% !TeX program = xelatex
% Compile this main file with XeLaTeX from its directory.
% Paragraph text uses Latin Modern. Mathematics uses default TeX math faces
% and standard AMS symbol faces. Page layouts position text at measured baselines.
% \Character uses a slot in a traditional TeX font. \Glyph uses a native font glyph.
% \EditionChange holds paired readings listed in the separate errata.
\newcommand\ActivateEdition{\CorrectedEditiontrue}
\newcommand\EditionTitle{Algebraic Number Theory — Corrected English Edition}
\documentclass{article}
\usepackage[paperwidth=612bp,paperheight=846bp,margin=0bp]{geometry}
\usepackage{graphicx}
\usepackage{eso-pic}
\usepackage[unicode,hidelinks,pageanchor=false,pdfpagelabels=false]{hyperref}
\usepackage{bookmark}
\pagestyle{empty}
\parindent=0pt
\newcommand\Face[1]{\csname face:#1\endcsname}
\newcommand\Font[3]{\expandafter\font\csname face:#1\endcsname=#2 at #3bp\relax}
\newcommand\Character[1]{\char#1\relax}
\newcommand\Glyph[1]{\XeTeXglyph#1\relax}
\newcommand\TextAt[4]{\put(#2,#3){\Face{#1}#4}}
\newcommand\Vector[1]{\put(0,0){\special{pdf:literal direct #1}}}
\newcommand\Destination[3]{\put(#1,#2){\special{pdf:dest (#3) [@thispage /XYZ @xpos @ypos null]}}}
\newcommand\Hotlink[5]{\put(#1,#2){\special{pdf:ann width #3bp height #4bp depth 0bp << /Subtype /Link /Border [0 0 0] /A << /S /GoTo /D (#5) >> >>}}}
\newcommand\Page[1]{\AddToShipoutPictureFG*{\AtPageLowerLeft{\setlength\unitlength{1bp}\begin{picture}(0,0)%
#1\end{picture}}}\null\newpage}
\providecommand\EditionTitle{Algebraic Number Theory}
\hypersetup{pdftitle={\EditionTitle},pdfauthor={Jürgen Neukirch}}


\usepackage{amsmath,amssymb,tikz}
\usetikzlibrary{patterns}
% Standard math constructions replace formerly positioned glyph composites.
\newcommand\NativeMathAt[4]{%
  \put(#1,#2){\begingroup\fontsize{#3bp}{#3bp}\selectfont\mathsurround=0pt\ensuremath{#4}\endgroup}}
\newcommand\NativeArrowAt[5]{%
  \put(#1,#2){\begingroup\fontsize{#3bp}{#3bp}\selectfont\mathsurround=0pt
  \resizebox{#4bp}{\height}{\ensuremath{#5}}\endgroup}}
\newcommand\NativeDiagram[1]{%
  \put(72,774){\begin{tikzpicture}[x=1bp,y=1bp,baseline=0pt]%
    \path[use as bounding box] (0,0) rectangle (1,1);%
    \begin{pgfinterruptboundingbox}#1\end{pgfinterruptboundingbox}%
  \end{tikzpicture}}}

\newif\ifCorrectedEdition
\providecommand\ActivateEdition{}
\ActivateEdition
\newcommand\EditionChange[2]{\ifCorrectedEdition#2\else#1\fi}
\newcommand\BodyAt[3]{\put(#1,#2){\fontsize{11.955bp}{16bp}\selectfont #3}}
\newcommand\SmallAt[3]{\put(#1,#2){\small #3}}
\newcommand\SeriesAt[3]{\put(#1,#2){\fontsize{9.3bp}{12bp}\selectfont #3}}
\newcommand\BoldAt[3]{\put(#1,#2){\Large\bfseries #3}}
\newcommand\CoverLine[3]{\put(306,#1){\makebox[0pt]{\fontsize{#2bp}{35bp}\selectfont #3}}}
\newcommand\NativeCover{%
  \CoverLine{681}{16}{Jürgen Neukirch}%
  \CoverLine{587}{37}{\bfseries Algebraic}%
  \CoverLine{543}{37}{\bfseries Number Theory}%
  \CoverLine{465}{13}{Translated from the German by Norbert Schappacher}%
  \CoverLine{425}{11}{With 16 Figures}%
  \CoverLine{100}{14}{Springer}%
  \CoverLine{76}{10}{\EditionChange{English edition}{Corrected English edition}}%
}

\Font{1}{"[lmroman12-bold.otf]:-kern"}{24.787}
\Font{2}{"[lmroman10-regular.otf]:-kern"}{10.909}
\Font{3}{cmsy8}{7.97}
\Font{4}{"[lmroman10-italic.otf]:-kern"}{10.909}
\Font{5}{"[lmmono10-regular.otf]:-kern"}{10.909}
\Font{6}{"[lmroman10-bold.otf]:-kern"}{10.909}
\Font{7}{"[lmroman8-regular.otf]:-kern"}{7.97}
\Font{8}{"[lmroman12-regular.otf]:-kern"}{11.955}
\Font{9}{cmr12}{11.955}
\Font{10}{"[lmroman10-regular.otf]:-kern"}{9.963}
\Font{11}{"[lmroman7-regular.otf]:-kern"}{6.974}
\Font{12}{"[lmroman12-bold.otf]:-kern"}{11.955}
\Font{13}{cmmi12}{11.955}
\Font{14}{eufm10}{11.955}
\Font{15}{"[lmroman12-bold.otf]:-kern"}{17.215}
\Font{16}{cmsy10}{11.955}
\Font{17}{cmr8}{7.97}
\Font{18}{msbm10}{11.955}
\Font{19}{"[lmroman12-italic.otf]:-kern"}{11.955}
\Font{20}{cmex10}{11.955}
\Font{21}{cmmi8}{7.97}
\Font{22}{cmmi6}{5.978}
\Font{23}{cmsy6}{5.978}
\Font{24}{cmr6}{5.978}
\Font{25}{msbm10}{7.97}
\Font{26}{"[lmromancaps10-regular.otf]:-kern"}{11.955}
\Font{27}{eufm10}{7.97}
\Font{28}{eufm7}{5.978}
\Font{29}{msbm7}{5.978}
\Font{30}{cmex8}{7.97}
\Font{31}{cmmi12}{17.215}
\Font{32}{eufm10}{17.215}
\Font{33}{cmmi10}{9.963}
\Font{34}{msbm10}{9.963}
\Font{35}{cmsy10}{9.963}
\Font{36}{cmmi7}{6.974}
\Font{37}{cmr10}{9.963}
\Font{38}{cmsy7}{6.974}
\Font{39}{cmr7}{6.974}
\Font{40}{"[lmroman6-regular.otf]:-kern"}{5.978}
\Font{41}{msam10}{11.955}
\Font{42}{cmex10}{9.963}
\Font{43}{cmsy5}{4.981}
\Font{44}{cmr5}{4.981}
\Font{45}{"[lmroman10-bold.otf]:-kern"}{9.963}
\Font{46}{msbm10}{17.215}
\Font{47}{cmmi10}{10.909}
\Font{48}{cmr10}{10.909}
\Font{49}{cmsy10}{10.909}
\Font{50}{cmmi12}{24.787}
\Font{51}{"[lmroman8-bold.otf]:-kern"}{7.97}
\Font{52}{eufm10}{9.963}
\Font{53}{eufm7}{6.974}
\Font{54}{cmmi5}{4.981}

% Place each verified correction at the original text baseline.
\newsavebox{\ErratumBox}
\newcommand\ErrataAt[4]{%
  \put(#1,#2){\begingroup\fontsize{11.955bp}{16bp}\selectfont
    \sbox{\ErratumBox}{#4}%
    \ifdim\wd\ErratumBox>#3bp
      \resizebox{#3bp}{\height}{\usebox{\ErratumBox}}%
    \else\usebox{\ErratumBox}\fi
  \endgroup}}

\begin{document}
\bookmark[level=0,dest={outline.1}]{Original Publication Data}
\bookmark[level=0,dest={outline.2}]{Foreword}
\bookmark[level=0,dest={outline.3}]{Translator's Note}
\bookmark[level=0,dest={outline.4}]{Preface to the German Edition}
\bookmark[level=0,dest={nav.chapterz2az.6}]{Contents}
\bookmark[level=0,dest={outline.5}]{I Algebraic Integers}
\bookmark[level=1,dest={outline.6}]{§1 The Gaussian Integers}
\bookmark[level=1,dest={outline.7}]{§2 Integrality}
\bookmark[level=1,dest={outline.8}]{§3 Ideals}
\bookmark[level=1,dest={outline.9}]{§4 Lattices}
\bookmark[level=1,dest={outline.10}]{§5 Minkowski Theory}
\bookmark[level=1,dest={outline.11}]{§6 The Class Number}
\bookmark[level=1,dest={outline.12}]{§7 Dirichlet's Unit Theorem}
\bookmark[level=1,dest={outline.13}]{§8 Extensions of Dedekind Domains}
\bookmark[level=1,dest={outline.14}]{§9 Hilbert's Ramification Theory}
\bookmark[level=1,dest={outline.15}]{§10 Cyclotomic Fields}
\bookmark[level=1,dest={outline.16}]{§11 Localization}
\bookmark[level=1,dest={outline.17}]{§12 Orders}
\bookmark[level=1,dest={outline.18}]{§13 One-dimensional Schemes}
\bookmark[level=1,dest={outline.19}]{§14 Function Fields}
\bookmark[level=0,dest={outline.20}]{II The Theory of Valuations}
\bookmark[level=1,dest={outline.21}]{§1 The p-adic Numbers}
\bookmark[level=1,dest={outline.22}]{§2 The p-adic Absolute Value}
\bookmark[level=1,dest={outline.23}]{§3 Valuations}
\bookmark[level=1,dest={outline.24}]{§4 Completions}
\bookmark[level=1,dest={outline.25}]{§5 Local Fields}
\bookmark[level=1,dest={outline.26}]{§6 Henselian Fields}
\bookmark[level=1,dest={outline.27}]{§7 Unramified and Tamely Ramified Extensions}
\bookmark[level=1,dest={outline.28}]{§8 Extensions of Valuations}
\bookmark[level=1,dest={outline.29}]{§9 Galois Theory of Valuations}
\bookmark[level=1,dest={outline.30}]{§10 Higher Ramification Groups}
\bookmark[level=0,dest={outline.31}]{III Riemann–Roch Theory}
\bookmark[level=1,dest={outline.32}]{§1 Primes}
\bookmark[level=1,dest={outline.33}]{§2 Different and Discriminant}
\bookmark[level=1,dest={outline.34}]{§3 Riemann–Roch}
\bookmark[level=1,dest={outline.35}]{§4 Metrized o-Modules}
\bookmark[level=1,dest={outline.36}]{§5 Grothendieck Groups}
\bookmark[level=1,dest={outline.37}]{§6 The Chern Character}
\bookmark[level=1,dest={outline.38}]{§7 Grothendieck–Riemann–Roch}
\bookmark[level=1,dest={outline.39}]{§8 The Euler–Minkowski Characteristic}
\bookmark[level=0,dest={outline.40}]{IV Abstract Class Field Theory}
\bookmark[level=1,dest={outline.41}]{§1 Infinite Galois Theory}
\bookmark[level=1,dest={outline.42}]{§2 Projective and Inductive Limits}
\bookmark[level=1,dest={outline.43}]{§3 Abstract Galois Theory}
\bookmark[level=1,dest={outline.44}]{§4 Abstract Valuation Theory}
\bookmark[level=1,dest={outline.45}]{§5 The Reciprocity Map}
\bookmark[level=1,dest={outline.46}]{§6 The General Reciprocity Law}
\bookmark[level=1,dest={outline.47}]{§7 The Herbrand Quotient}
\bookmark[level=0,dest={outline.48}]{V Local Class Field Theory}
\bookmark[level=1,dest={outline.49}]{§1 The Local Reciprocity Law}
\bookmark[level=1,dest={outline.50}]{§2 The Norm Residue Symbol over Qp}
\bookmark[level=1,dest={outline.51}]{§3 The Hilbert Symbol}
\bookmark[level=1,dest={outline.52}]{§4 Formal Groups}
\bookmark[level=1,dest={outline.53}]{§5 Generalized Cyclotomic Theory}
\bookmark[level=1,dest={outline.54}]{§6 Higher Ramification Groups}
\bookmark[level=0,dest={outline.55}]{VI Global Class Field Theory}
\bookmark[level=1,dest={outline.56}]{§1 Idèles and Idèle Classes}
\bookmark[level=1,dest={outline.57}]{§2 Idèles in Field Extensions}
\bookmark[level=1,dest={outline.58}]{§3 The Herbrand Quotient of the Idèle Class Group}
\bookmark[level=1,dest={outline.59}]{§4 The Class Field Axiom}
\bookmark[level=1,dest={outline.60}]{§5 The Global Reciprocity Law}
\bookmark[level=1,dest={outline.61}]{§6 Global Class Fields}
\bookmark[level=1,dest={outline.62}]{§7 The Ideal-Theoretic Version of Class Field Theory}
\bookmark[level=1,dest={outline.63}]{§8 The Reciprocity Law of the Power Residues}
\bookmark[level=0,dest={outline.64}]{VII Zeta Functions and L-series}
\bookmark[level=1,dest={outline.65}]{§1 The Riemann Zeta Function}
\bookmark[level=1,dest={outline.66}]{§2 Dirichlet L-series}
\bookmark[level=1,dest={outline.67}]{§3 Theta Series}
\bookmark[level=1,dest={outline.68}]{§4 The Higher-dimensional Gamma Function}
\bookmark[level=1,dest={outline.69}]{§5 The Dedekind Zeta Function}
\bookmark[level=1,dest={outline.70}]{§6 Hecke Characters}
\bookmark[level=1,dest={outline.71}]{§7 Theta Series of Algebraic Number Fields}
\bookmark[level=1,dest={outline.72}]{§8 Hecke L-series}
\bookmark[level=1,dest={outline.73}]{§9 Values of Dirichlet L-series at Integer Points}
\bookmark[level=1,dest={outline.74}]{§10 Artin L-series}
\bookmark[level=1,dest={outline.75}]{§11 The Artin Conductor}
\bookmark[level=1,dest={outline.76}]{§12 The Functional Equation of Artin L-series}
\bookmark[level=1,dest={outline.77}]{§13 Density Theorems}
\bookmark[level=0,dest={outline.78}]{Bibliography}
\bookmark[level=0,dest={outline.79}]{Index}
\bookmark[level=0,dest={publisher.series}]{Grundlehren der mathematischen Wissenschaften}
% Title
% Page 1 (body page front matter)
\Page{%
\NativeCover%
\Destination{79.23}{773.075}{nav.Doc-Start}%
\Destination{78.23}{809.815}{nav.page.i}%
}

% Copyright
% Page 2 (body page front matter)
\Page{%
\TextAt{1}{79.226}{673.054}{Original}%
\TextAt{1}{187.149}{673.054}{Publication}%
\TextAt{1}{336.02}{673.054}{Data}%
\TextAt{2}{79.226}{615.739}{Jürgen}%
\TextAt{2}{115.172}{615.739}{Neukirch}%
\TextAt{3}{157.935}{619.698}{\Character{121}}%
\TextAt{4}{79.226}{599.926}{Translator:}%
\TextAt{2}{142.837}{599.926}{Norbert}%
\TextAt{2}{184.117}{599.926}{Schappacher}%
\TextAt{2}{79.226}{584.113}{U.F.R.}%
\TextAt{2}{115.292}{584.113}{de}%
\TextAt{2}{129.833}{584.113}{Mathématique}%
\TextAt{2}{202.564}{584.113}{et}%
\TextAt{2}{215.284}{584.113}{d’Informatique}%
\TextAt{2}{79.226}{568.301}{Université}%
\TextAt{2}{131.132}{568.301}{Louis}%
\TextAt{2}{160.433}{568.301}{Pasteur}%
\TextAt{2}{79.226}{552.488}{7,}%
\TextAt{2}{91.346}{552.488}{rue}%
\TextAt{2}{110.164}{552.488}{René}%
\TextAt{2}{137.579}{552.488}{Descartes}%
\TextAt{2}{79.226}{536.675}{F-67084}%
\TextAt{2}{120.888}{536.675}{Strasbourg,}%
\TextAt{2}{179.513}{536.675}{France}%
\TextAt{2}{79.226}{520.862}{e-mail:}%
\TextAt{5}{116.186}{520.862}{schappa@math.u-strasbg.fr}%
\TextAt{2}{79.226}{505.05}{The}%
\TextAt{2}{101.644}{505.05}{original}%
\TextAt{2}{141.081}{505.05}{German}%
\TextAt{2}{183.004}{505.05}{edition}%
\TextAt{2}{219.375}{505.05}{was}%
\TextAt{2}{240.332}{505.05}{published}%
\TextAt{2}{289.499}{505.05}{in}%
\TextAt{2}{302.23}{505.05}{1992}%
\TextAt{2}{327.681}{505.05}{under}%
\TextAt{2}{358.63}{505.05}{the}%
\TextAt{2}{377.415}{505.05}{title}%
\TextAt{4}{400.444}{505.05}{Algebraische}%
\TextAt{4}{462.932}{505.05}{Zahlentheorie}%
\TextAt{2}{526.913}{505.05}{.}%
\TextAt{2}{79.226}{489.237}{ISBN}%
\TextAt{2}{108.768}{489.237}{3-540-54273-6}%
\TextAt{2}{177.844}{489.237}{Springer-Verlag}%
\TextAt{2}{255.528}{489.237}{Berlin}%
\TextAt{2}{288.135}{489.237}{Heidelberg}%
\TextAt{2}{342.713}{489.237}{New}%
\TextAt{2}{367.248}{489.237}{York.}%
\TextAt{6}{79.226}{473.424}{Library}%
\TextAt{6}{124.466}{473.424}{of}%
\TextAt{6}{138.746}{473.424}{Congress}%
\TextAt{6}{192.332}{473.424}{Cataloging-in-Publication}%
\TextAt{6}{336.953}{473.424}{Data}%
\TextAt{2}{79.226}{457.611}{Neukirch,}%
\TextAt{2}{130.718}{457.611}{Jürgen,}%
\TextAt{2}{171.758}{457.611}{1937–.}%
\TextAt{2}{211.857}{457.611}{[Algebraische}%
\TextAt{2}{280.26}{457.611}{zahlentheorie.}%
\TextAt{2}{355.846}{457.611}{English]}%
\TextAt{2}{399.54}{457.611}{Algebraic}%
\TextAt{2}{450.008}{457.611}{number}%
\TextAt{2}{491.389}{457.611}{theory}%
\TextAt{2}{527.315}{457.611}{/}%
\TextAt{2}{79.226}{443.791}{Jürgen}%
\TextAt{2}{114.869}{443.791}{Neukirch;}%
\TextAt{2}{164.096}{443.791}{translated}%
\TextAt{2}{215.404}{443.791}{from}%
\TextAt{2}{240.891}{443.791}{the}%
\TextAt{2}{259.374}{443.791}{German}%
\TextAt{2}{300.995}{443.791}{by}%
\TextAt{2}{315.845}{443.791}{Norbert}%
\TextAt{2}{356.822}{443.791}{Schappacher.}%
\TextAt{2}{424.03}{443.791}{p.}%
\TextAt{2}{437.871}{443.791}{cm.}%
\TextAt{2}{459.578}{443.791}{–}%
\TextAt{2}{468.362}{443.791}{(Grundlehren}%
\TextAt{2}{79.226}{429.971}{der}%
\TextAt{2}{97.821}{429.971}{mathematischen}%
\TextAt{2}{178.554}{429.971}{Wissenschaften;}%
\TextAt{2}{257.67}{429.971}{322).}%
\TextAt{2}{286.079}{429.971}{Includes}%
\TextAt{2}{328.642}{429.971}{bibliographical}%
\TextAt{2}{402.414}{429.971}{references}%
\TextAt{2}{452.296}{429.971}{and}%
\TextAt{2}{473.291}{429.971}{index.}%
\TextAt{2}{506.86}{429.971}{ISBN}%
\TextAt{2}{79.226}{416.151}{3-540-65399-6}%
\TextAt{2}{148.302}{416.151}{(hc.:}%
\TextAt{2}{174.364}{416.151}{alk.}%
\TextAt{2}{196.488}{416.151}{paper).}%
\TextAt{2}{235.619}{416.151}{1.}%
\TextAt{2}{248.95}{416.151}{Algebraic}%
\TextAt{2}{297.768}{416.151}{number}%
\TextAt{2}{337.499}{416.151}{theory.}%
\TextAt{2}{375.113}{416.151}{I.}%
\TextAt{2}{385.717}{416.151}{Title}%
\TextAt{2}{412.379}{416.151}{II.}%
\TextAt{2}{427.226}{416.151}{Series.}%
\TextAt{2}{79.226}{400.338}{QA247.N51713}%
\TextAt{2}{161.655}{400.338}{1999}%
\TextAt{2}{194.382}{400.338}{512’.74-dc21}%
%
%
\TextAt{2}{263.843}{400.338}{99-21030}%
\TextAt{2}{309.29}{400.338}{CIP.}%
\TextAt{2}{79.226}{384.525}{Mathematics}%
\TextAt{2}{144.43}{384.525}{Subject}%
\TextAt{2}{183.833}{384.525}{Classification}%
\TextAt{2}{250.619}{384.525}{(1991):}%
\TextAt{2}{288.801}{384.525}{11-XX,}%
\TextAt{2}{326.372}{384.525}{14-XX.}%
\TextAt{2}{79.226}{368.712}{ISSN}%
\TextAt{2}{107.11}{368.712}{0072-7830.}%
\TextAt{2}{79.226}{352.9}{ISBN}%
\TextAt{2}{108.768}{352.9}{3-540-65399-6}%
\TextAt{2}{177.844}{352.9}{Springer-Verlag}%
\TextAt{2}{255.528}{352.9}{Berlin}%
\TextAt{2}{288.135}{352.9}{Heidelberg}%
\TextAt{2}{342.713}{352.9}{New}%
\TextAt{2}{367.248}{352.9}{York.}%
\TextAt{2}{79.226}{337.087}{This}%
\TextAt{2}{105.175}{337.087}{work}%
\TextAt{2}{132.913}{337.087}{is}%
\TextAt{2}{144.92}{337.087}{subject}%
\TextAt{2}{183.6}{337.087}{to}%
\TextAt{2}{197.975}{337.087}{copyright.}%
\TextAt{2}{253.567}{337.087}{All}%
\TextAt{2}{272.491}{337.087}{rights}%
\TextAt{2}{304.233}{337.087}{are}%
\TextAt{2}{323.484}{337.087}{reserved,}%
\TextAt{2}{370.355}{337.087}{whether}%
\TextAt{2}{413.246}{337.087}{the}%
\TextAt{2}{433.075}{337.087}{whole}%
\TextAt{2}{465.024}{337.087}{or}%
\TextAt{2}{479.431}{337.087}{part}%
\TextAt{2}{504.148}{337.087}{of}%
\TextAt{2}{517.617}{337.087}{the}%
\TextAt{2}{79.226}{323.267}{material}%
\TextAt{2}{123.687}{323.267}{is}%
\TextAt{2}{136.053}{323.267}{concerned,}%
\TextAt{2}{191.774}{323.267}{specifically}%
\TextAt{2}{248.388}{323.267}{the}%
\TextAt{2}{268.576}{323.267}{rights}%
\TextAt{2}{300.677}{323.267}{of}%
\TextAt{2}{314.505}{323.267}{translation,}%
\TextAt{2}{374.545}{323.267}{reprinting,}%
\TextAt{2}{430.015}{323.267}{reuse}%
\TextAt{2}{459.378}{323.267}{of}%
\TextAt{2}{473.206}{323.267}{illustrations,}%
\TextAt{2}{79.226}{309.447}{recitation,}%
\TextAt{2}{132.604}{309.447}{broadcasting,}%
\TextAt{2}{201.2}{309.447}{reproduction}%
\TextAt{2}{266.803}{309.447}{on}%
\TextAt{2}{282.933}{309.447}{microfilm}%
\TextAt{2}{332.423}{309.447}{or}%
\TextAt{2}{346.764}{309.447}{in}%
\TextAt{2}{360.472}{309.447}{any}%
\TextAt{2}{382.057}{309.447}{other}%
\TextAt{2}{411.55}{309.447}{way,}%
\TextAt{2}{437.012}{309.447}{and}%
\TextAt{2}{459.208}{309.447}{storage}%
\TextAt{2}{497.843}{309.447}{in}%
\TextAt{2}{511.551}{309.447}{data}%
\TextAt{2}{79.226}{295.626}{banks.}%
\TextAt{2}{116.212}{295.626}{Duplication}%
\TextAt{2}{175.959}{295.626}{of}%
\TextAt{2}{188.873}{295.626}{this}%
\TextAt{2}{210.634}{295.626}{publication}%
\TextAt{2}{268.112}{295.626}{or}%
\TextAt{2}{281.964}{295.626}{parts}%
\TextAt{2}{310.423}{295.626}{thereof}%
\TextAt{2}{347.61}{295.626}{is}%
\TextAt{2}{359.062}{295.626}{permitted}%
\TextAt{2}{410.191}{295.626}{only}%
\TextAt{2}{434.625}{295.626}{under}%
\TextAt{2}{466.063}{295.626}{the}%
\TextAt{2}{485.337}{295.626}{provisions}%
\TextAt{2}{79.226}{281.806}{of}%
\TextAt{2}{91.908}{281.806}{the}%
\TextAt{2}{110.949}{281.806}{German}%
\TextAt{2}{153.129}{281.806}{Copyright}%
\TextAt{2}{204.636}{281.806}{Law}%
\TextAt{2}{228.369}{281.806}{of}%
\TextAt{2}{241.05}{281.806}{September}%
\TextAt{2}{295.274}{281.806}{9,}%
\TextAt{2}{307.714}{281.806}{1965,}%
\TextAt{2}{336.518}{281.806}{in}%
\TextAt{2}{349.505}{281.806}{its}%
\TextAt{2}{364.968}{281.806}{current}%
\TextAt{2}{403.166}{281.806}{version,}%
\TextAt{2}{443.577}{281.806}{and}%
\TextAt{2}{465.052}{281.806}{permission}%
\TextAt{2}{519.7}{281.806}{for}%
\TextAt{2}{79.226}{267.986}{use}%
\TextAt{2}{98.343}{267.986}{must}%
\TextAt{2}{125.642}{267.986}{always}%
\TextAt{2}{160.818}{267.986}{be}%
\TextAt{2}{175.942}{267.986}{obtained}%
\TextAt{2}{221.078}{267.986}{from}%
\TextAt{2}{247.144}{267.986}{Springer-Verlag.}%
\TextAt{2}{329.903}{267.986}{Violations}%
\TextAt{2}{382.064}{267.986}{are}%
\TextAt{2}{400.549}{267.986}{liable}%
\TextAt{2}{429.92}{267.986}{for}%
\TextAt{2}{446.899}{267.986}{prosecution}%
\TextAt{2}{505.453}{267.986}{under}%
\TextAt{2}{79.226}{254.166}{the}%
\TextAt{2}{98.012}{254.166}{German}%
\TextAt{2}{139.935}{254.166}{Copyright}%
\TextAt{2}{191.186}{254.166}{Law.}%
\TextAt{2}{79.226}{238.353}{©}%
\TextAt{2}{90.31}{238.353}{Springer-Verlag}%
\TextAt{2}{167.993}{238.353}{Berlin}%
\TextAt{2}{200.601}{238.353}{Heidelberg}%
\TextAt{2}{255.179}{238.353}{1999.}%
\TextAt{2}{79.226}{222.54}{Printed}%
\TextAt{2}{118.51}{222.54}{in}%
\TextAt{2}{131.241}{222.54}{Germany.}%
\TextAt{2}{79.226}{206.727}{Cover}%
\TextAt{2}{110.459}{206.727}{design:}%
\TextAt{2}{148.095}{206.727}{MetaDesign}%
\TextAt{2}{208.302}{206.727}{plus}%
\TextAt{2}{231.397}{206.727}{GmbH,}%
\TextAt{2}{269.655}{206.727}{Berlin.}%
\TextAt{2}{79.226}{190.915}{Photocomposed}%
\TextAt{2}{159.252}{190.915}{from}%
\TextAt{2}{186.064}{190.915}{the}%
\TextAt{2}{205.872}{190.915}{translator’s}%
\TextAt{2}{264.658}{190.915}{L}%
\TextAt{7}{267.549}{192.609}{A}%
\TextAt{2}{272.256}{190.915}{T}%
\TextAt{2}{278.314}{188.564}{E}%
\TextAt{2}{284.38}{190.915}{X}%
\TextAt{2}{297.217}{190.915}{files}%
\TextAt{2}{320.112}{190.915}{after}%
\TextAt{2}{346.923}{190.915}{editing}%
\TextAt{2}{384.317}{190.915}{and}%
\TextAt{2}{406.557}{190.915}{reformatting}%
\TextAt{2}{470.983}{190.915}{by}%
\TextAt{2}{487.159}{190.915}{Raymond}%
\TextAt{2}{79.226}{177.094}{Seroul,}%
\TextAt{2}{115.63}{177.094}{Strasbourg.}%
\TextAt{2}{79.226}{161.282}{SPIN:}%
\TextAt{2}{111.506}{161.282}{10425040}%
\TextAt{2}{176.961}{161.282}{41/3143-5}%
\TextAt{2}{235.139}{161.282}{4}%
\TextAt{2}{244.226}{161.282}{3}%
\TextAt{2}{253.313}{161.282}{2}%
\TextAt{2}{262.401}{161.282}{1}%
\TextAt{2}{271.488}{161.282}{0}%
\TextAt{2}{79.226}{145.469}{Printed}%
\TextAt{2}{118.51}{145.469}{on}%
\TextAt{2}{133.662}{145.469}{acid-free}%
\TextAt{2}{177.626}{145.469}{paper.}%
\TextAt{8}{304.372}{43.033}{i}%
\Destination{79.23}{773.075}{nav.chapterz2az.1}%
\Destination{79.23}{773.075}{outline.1}%
\Destination{79.226}{419}{errata.213}%
}

% Foreword
% Page 3 (body page front matter)
\Page{%
\TextAt{1}{79.226}{673.054}{Foreword}%
\TextAt{8}{79.226}{613.169}{It}%
\TextAt{8}{92.551}{613.169}{is}%
\TextAt{8}{104.967}{613.169}{a}%
\TextAt{8}{115.375}{613.169}{very}%
\TextAt{8}{141.719}{613.169}{sad}%
\TextAt{8}{163.245}{613.169}{moment}%
\TextAt{8}{209.1}{613.169}{for}%
\TextAt{8}{227.637}{613.169}{me}%
\TextAt{8}{247.143}{613.169}{to}%
\TextAt{8}{262.105}{613.169}{write}%
\TextAt{8}{292.669}{613.169}{this}%
\TextAt{8}{316.144}{613.169}{“Geleitwort”}%
\TextAt{8}{385.156}{613.169}{to}%
\TextAt{8}{400.119}{613.169}{the}%
\TextAt{8}{420.927}{613.169}{English}%
\TextAt{8}{463.423}{613.169}{translation}%
\TextAt{8}{523.337}{613.169}{of}%
\TextAt{8}{79.226}{596.123}{Jürgen}%
\TextAt{8}{116.888}{596.123}{Neukirch’s}%
\TextAt{8}{173.629}{596.123}{book}%
\TextAt{8}{201.703}{596.123}{on}%
\TextAt{8}{217.092}{596.123}{Algebraic}%
\TextAt{8}{268.573}{596.123}{Number}%
\TextAt{8}{312.894}{596.123}{Theory.}%
\TextAt{8}{356.822}{596.123}{It}%
\TextAt{8}{368.624}{596.123}{would}%
\TextAt{8}{401.898}{596.123}{have}%
\TextAt{8}{428.022}{596.123}{been}%
\TextAt{8}{454.781}{596.123}{so}%
\TextAt{8}{468.281}{596.123}{much}%
\TextAt{8}{498.626}{596.123}{better,}%
\TextAt{8}{79.226}{579.077}{if}%
\TextAt{8}{89.95}{579.077}{he}%
\TextAt{8}{105.551}{579.077}{could}%
\TextAt{8}{136.766}{579.077}{have}%
\TextAt{8}{163.761}{579.077}{done}%
\TextAt{8}{191.724}{579.077}{this}%
\TextAt{8}{214.547}{579.077}{himself.}%
\TextAt{8}{97.159}{559.042}{But}%
\TextAt{8}{122.073}{559.042}{it}%
\TextAt{8}{135.45}{559.042}{is}%
\TextAt{8}{148.887}{559.042}{also}%
\TextAt{8}{174.04}{559.042}{very}%
\TextAt{8}{201.405}{559.042}{difficult}%
\TextAt{8}{245.998}{559.042}{for}%
\TextAt{8}{265.556}{559.042}{me}%
\TextAt{8}{286.082}{559.042}{to}%
\TextAt{8}{302.066}{559.042}{write}%
\TextAt{8}{333.651}{559.042}{this}%
\TextAt{8}{358.146}{559.042}{“Geleitwort”:}%
\TextAt{8}{434.403}{559.042}{The}%
\TextAt{8}{460.13}{559.042}{book}%
\TextAt{8}{490.747}{559.042}{contains}%
\TextAt{8}{79.226}{541.996}{Neukirch’s}%
\TextAt{8}{136.809}{541.996}{Preface}%
\TextAt{8}{178.228}{541.996}{to}%
\TextAt{8}{192.509}{541.996}{the}%
\TextAt{8}{212.637}{541.996}{German}%
\TextAt{8}{257.559}{541.996}{edition.}%
\TextAt{8}{301.122}{541.996}{There}%
\TextAt{8}{334.902}{541.996}{he}%
\TextAt{8}{350.474}{541.996}{himself}%
\TextAt{8}{390.495}{541.996}{speaks}%
\TextAt{8}{427.658}{541.996}{about}%
\TextAt{8}{461.127}{541.996}{his}%
\TextAt{8}{479.366}{541.996}{intentions,}%
\TextAt{8}{79.226}{524.95}{the}%
\TextAt{8}{99.383}{524.95}{content}%
\TextAt{8}{141.01}{524.95}{of}%
\TextAt{8}{154.34}{524.95}{the}%
\TextAt{8}{174.497}{524.95}{book}%
\TextAt{8}{203.44}{524.95}{and}%
\TextAt{8}{226.203}{524.95}{his}%
\TextAt{8}{244.471}{524.95}{personal}%
\TextAt{8}{291.036}{524.95}{view}%
\TextAt{8}{318.019}{524.95}{of}%
\TextAt{8}{331.349}{524.95}{the}%
\TextAt{8}{351.505}{524.95}{subject.}%
\TextAt{8}{396.429}{524.95}{What}%
\TextAt{8}{429.27}{524.95}{else}%
\TextAt{8}{451.434}{524.95}{can}%
\TextAt{8}{472.894}{524.95}{be}%
\TextAt{8}{488.818}{524.95}{said?}%
\TextAt{8}{97.159}{504.915}{It}%
\TextAt{8}{109.08}{504.915}{becomes}%
\TextAt{8}{154.882}{504.915}{clear}%
\TextAt{8}{182.094}{504.915}{from}%
\TextAt{8}{208.984}{504.915}{his}%
\TextAt{8}{226.5}{504.915}{Preface}%
\TextAt{8}{267.197}{504.915}{that}%
\TextAt{8}{291.815}{504.915}{Number}%
\TextAt{8}{336.254}{504.915}{Theory}%
\TextAt{8}{376.151}{504.915}{was}%
\TextAt{8}{397.899}{504.915}{Neukirch’s}%
\TextAt{8}{454.76}{504.915}{favorite}%
\TextAt{8}{496.294}{504.915}{subject}%
\TextAt{8}{79.226}{487.869}{in}%
\TextAt{8}{94.608}{487.869}{mathematics.}%
\TextAt{8}{173.352}{487.869}{He}%
\TextAt{8}{192.954}{487.869}{was}%
\TextAt{8}{217.183}{487.869}{enthusiastic}%
\TextAt{8}{282.776}{487.869}{about}%
\TextAt{8}{318.004}{487.869}{it,}%
\TextAt{8}{335.121}{487.869}{and}%
\TextAt{8}{359.613}{487.869}{he}%
\TextAt{8}{376.944}{487.869}{was}%
\TextAt{8}{401.172}{487.869}{also}%
\TextAt{8}{426.382}{487.869}{able}%
\TextAt{8}{452.822}{487.869}{to}%
\TextAt{8}{468.861}{487.869}{implant}%
\TextAt{8}{513.844}{487.869}{this}%
\TextAt{8}{79.226}{470.823}{enthusiasm}%
\TextAt{8}{139.839}{470.823}{into}%
\TextAt{8}{163.582}{470.823}{the}%
\TextAt{8}{183.738}{470.823}{minds}%
\TextAt{8}{218.265}{470.823}{of}%
\TextAt{8}{231.595}{470.823}{his}%
\TextAt{8}{249.862}{470.823}{students.}%
\TextAt{8}{97.159}{450.789}{He}%
\TextAt{8}{116.536}{450.789}{attracted}%
\TextAt{8}{168.779}{450.789}{them,}%
\TextAt{8}{203.823}{450.789}{they}%
\TextAt{8}{231.665}{450.789}{gathered}%
\TextAt{8}{281.301}{450.789}{around}%
\TextAt{8}{322.485}{450.789}{him}%
\TextAt{8}{347.397}{450.789}{in}%
\TextAt{8}{362.555}{450.789}{Regensburg.}%
\TextAt{8}{434.934}{450.789}{He}%
\TextAt{8}{454.311}{450.789}{told}%
\TextAt{8}{479.882}{450.789}{them}%
\TextAt{8}{511.298}{450.789}{that}%
\TextAt{8}{79.226}{433.743}{the}%
\TextAt{8}{101.14}{433.743}{subject}%
\TextAt{8}{143.27}{433.743}{and}%
\TextAt{8}{167.79}{433.743}{its}%
\TextAt{8}{185.866}{433.743}{beauty}%
\TextAt{8}{226.323}{433.743}{justified}%
\TextAt{8}{272.685}{433.743}{the}%
\TextAt{8}{294.599}{433.743}{highest}%
\TextAt{8}{336.741}{433.743}{effort}%
\TextAt{8}{369.39}{433.743}{and}%
\TextAt{8}{393.91}{433.743}{so}%
\TextAt{8}{410.038}{433.743}{they}%
\TextAt{8}{438.132}{433.743}{were}%
\TextAt{8}{466.873}{433.743}{always}%
\TextAt{8}{506.097}{433.743}{eager}%
\TextAt{8}{79.226}{416.697}{and}%
\TextAt{8}{103.009}{416.697}{motivated}%
\TextAt{8}{158.998}{416.697}{to}%
\TextAt{8}{174.329}{416.697}{discuss}%
\TextAt{8}{214.549}{416.697}{and}%
\TextAt{8}{238.332}{416.697}{to}%
\TextAt{8}{253.662}{416.697}{learn}%
\TextAt{8}{283.948}{416.697}{the}%
\TextAt{8}{305.125}{416.697}{newest}%
\TextAt{8}{344.246}{416.697}{developments}%
\TextAt{8}{417.845}{416.697}{in}%
\TextAt{8}{432.518}{416.697}{number}%
\TextAt{8}{476.134}{416.697}{theory}%
\TextAt{8}{513.904}{416.697}{and}%
\TextAt{8}{79.226}{399.651}{arithmetic}%
\TextAt{8}{135.014}{399.651}{algebraic}%
\TextAt{8}{183.653}{399.651}{geometry.}%
\TextAt{8}{238.018}{399.651}{I}%
\TextAt{8}{245.339}{399.651}{remember}%
\TextAt{8}{299.166}{399.651}{very}%
\TextAt{8}{324.062}{399.651}{well}%
\TextAt{8}{346.997}{399.651}{the}%
\TextAt{8}{366.357}{399.651}{many}%
\TextAt{8}{397.434}{399.651}{occasions}%
\TextAt{8}{447.818}{399.651}{when}%
\TextAt{8}{477.579}{399.651}{this}%
\TextAt{8}{499.606}{399.651}{equipe}%
\TextAt{8}{79.226}{382.605}{showed}%
\TextAt{8}{119.769}{382.605}{up}%
\TextAt{8}{136.831}{382.605}{in}%
\TextAt{8}{150.642}{382.605}{the}%
\TextAt{8}{170.956}{382.605}{meetings}%
\TextAt{8}{219.95}{382.605}{of}%
\TextAt{8}{233.438}{382.605}{the}%
\TextAt{8}{253.753}{382.605}{“Oberwolfach}%
\TextAt{8}{327.076}{382.605}{Arbeitsgemeinschaft”}%
\TextAt{8}{440.186}{382.605}{and}%
\TextAt{8}{463.107}{382.605}{demonstrated}%
\TextAt{8}{79.226}{365.559}{their}%
\TextAt{8}{107.189}{365.559}{strength}%
\TextAt{8}{153.432}{365.559}{(mathematically}%
\TextAt{8}{240.92}{365.559}{and}%
\TextAt{8}{263.682}{365.559}{on}%
\TextAt{8}{279.941}{365.559}{the}%
\TextAt{8}{300.098}{365.559}{soccer}%
\TextAt{8}{334.947}{365.559}{field).}%
\TextAt{8}{97.159}{345.524}{During}%
\TextAt{8}{136.511}{345.524}{the}%
\TextAt{8}{156.508}{345.524}{meetings}%
\TextAt{8}{205.186}{345.524}{of}%
\TextAt{8}{218.356}{345.524}{the}%
\TextAt{8}{238.353}{345.524}{“Oberwolfach}%
\TextAt{8}{311.36}{345.524}{Arbeitsgemeinschaft”}%
\TextAt{8}{424.153}{345.524}{people}%
\TextAt{8}{460.732}{345.524}{come}%
\TextAt{8}{490.484}{345.524}{together}%
\TextAt{8}{79.226}{328.478}{to}%
\TextAt{8}{93.588}{328.478}{learn}%
\TextAt{8}{122.906}{328.478}{a}%
\TextAt{8}{132.713}{328.478}{subject}%
\TextAt{8}{173.137}{328.478}{which}%
\TextAt{8}{206.675}{328.478}{is}%
\TextAt{8}{218.49}{328.478}{not}%
\TextAt{8}{239.356}{328.478}{necessarily}%
\TextAt{8}{297.736}{328.478}{their}%
\TextAt{8}{325.751}{328.478}{own}%
\TextAt{8}{350.191}{328.478}{speciality.}%
\TextAt{8}{405.682}{328.478}{Always}%
\TextAt{8}{446.118}{328.478}{at}%
\TextAt{8}{460.48}{328.478}{the}%
\TextAt{8}{480.688}{328.478}{end,}%
\TextAt{8}{506.109}{328.478}{when}%
\TextAt{8}{79.226}{311.432}{the}%
\TextAt{8}{101.218}{311.432}{most}%
\TextAt{8}{131.735}{311.432}{difficult}%
\TextAt{8}{176.49}{311.432}{talks}%
\TextAt{8}{206.683}{311.432}{had}%
\TextAt{8}{231.282}{311.432}{to}%
\TextAt{8}{247.428}{311.432}{be}%
\TextAt{8}{265.188}{311.432}{delivered,}%
\TextAt{8}{320.157}{311.432}{the}%
\TextAt{8}{342.149}{311.432}{Regensburg}%
\TextAt{8}{407.3}{311.432}{crew}%
\TextAt{8}{436.441}{311.432}{took}%
\TextAt{8}{464.949}{311.432}{over.}%
\TextAt{8}{500.053}{311.432}{In}%
\TextAt{8}{516.51}{311.432}{the}%
\TextAt{8}{79.226}{294.386}{meantime}%
\TextAt{8}{132.881}{294.386}{many}%
\TextAt{8}{164.753}{294.386}{members}%
\TextAt{8}{214.236}{294.386}{of}%
\TextAt{8}{227.566}{294.386}{this}%
\TextAt{8}{250.388}{294.386}{team}%
\TextAt{8}{279.655}{294.386}{teach}%
\TextAt{8}{310.547}{294.386}{at}%
\TextAt{8}{324.857}{294.386}{German}%
\TextAt{8}{369.809}{294.386}{universities.}%
\TextAt{8}{97.159}{274.351}{We}%
\TextAt{8}{117.295}{274.351}{find}%
\TextAt{8}{140.695}{274.351}{this}%
\TextAt{8}{163.51}{274.351}{charisma}%
\TextAt{8}{212.673}{274.351}{of}%
\TextAt{8}{225.995}{274.351}{Jürgen}%
\TextAt{8}{264.519}{274.351}{Neukirch}%
\TextAt{8}{314.256}{274.351}{in}%
\TextAt{8}{327.901}{274.351}{the}%
\TextAt{8}{348.049}{274.351}{book.}%
\TextAt{8}{381.541}{274.351}{It}%
\TextAt{8}{394.205}{274.351}{will}%
\TextAt{8}{416.302}{274.351}{be}%
\TextAt{8}{432.219}{274.351}{a}%
\TextAt{8}{441.966}{274.351}{motivating}%
\TextAt{8}{500.837}{274.351}{source}%
\TextAt{8}{79.226}{257.305}{for}%
\TextAt{8}{96.979}{257.305}{young}%
\TextAt{8}{131.326}{257.305}{students}%
\TextAt{8}{177.82}{257.305}{to}%
\TextAt{8}{191.999}{257.305}{study}%
\TextAt{8}{224.122}{257.305}{Algebraic}%
\TextAt{8}{276.342}{257.305}{Number}%
\TextAt{8}{321.401}{257.305}{Theory,}%
\TextAt{8}{364.215}{257.305}{and}%
\TextAt{8}{386.846}{257.305}{I}%
\TextAt{8}{394.831}{257.305}{am}%
\TextAt{8}{414.211}{257.305}{sure}%
\TextAt{8}{438.85}{257.305}{that}%
\TextAt{8}{464.087}{257.305}{it}%
\TextAt{8}{475.66}{257.305}{will}%
\TextAt{8}{497.633}{257.305}{attract}%
\TextAt{8}{79.226}{240.259}{many}%
\TextAt{8}{111.099}{240.259}{of}%
\TextAt{8}{124.429}{240.259}{them.}%
\TextAt{8}{97.159}{220.225}{At}%
\TextAt{8}{113.683}{220.225}{Neukirch’s}%
\TextAt{8}{170.914}{220.225}{funeral}%
\TextAt{8}{209.878}{220.225}{his}%
\TextAt{8}{227.765}{220.225}{daughter}%
\TextAt{8}{276.496}{220.225}{Christiane}%
\TextAt{8}{332.759}{220.225}{recited}%
\TextAt{8}{370.743}{220.225}{the}%
\TextAt{8}{390.519}{220.225}{poem}%
\TextAt{8}{421.999}{220.225}{which}%
\TextAt{8}{455.105}{220.225}{she}%
\TextAt{8}{474.94}{220.225}{often}%
\TextAt{8}{504.149}{220.225}{heard}%
\TextAt{8}{79.226}{203.179}{from}%
\TextAt{8}{107.767}{203.179}{her}%
\TextAt{8}{128.825}{203.179}{father:}%
\TextAt{8}{169.321}{203.179}{Herr}%
\TextAt{8}{197.205}{203.179}{von}%
\TextAt{8}{220.223}{203.179}{Ribbeck}%
\TextAt{8}{266.482}{203.179}{auf}%
\TextAt{8}{287.216}{203.179}{Ribbeck}%
\TextAt{8}{333.475}{203.179}{im}%
\TextAt{8}{351.281}{203.179}{Havelland}%
\TextAt{8}{406.817}{203.179}{by}%
\TextAt{8}{423.976}{203.179}{Theodor}%
\TextAt{8}{472.028}{203.179}{Fontane.}%
\TextAt{8}{523.994}{203.179}{It}%
\TextAt{8}{79.226}{186.133}{tells}%
\TextAt{8}{103.549}{186.133}{the}%
\TextAt{8}{123.258}{186.133}{story}%
\TextAt{8}{152.471}{186.133}{of}%
\TextAt{8}{165.353}{186.133}{a}%
\TextAt{8}{174.66}{186.133}{nobleman}%
\TextAt{8}{227.544}{186.133}{who}%
\TextAt{8}{251.808}{186.133}{always}%
\TextAt{8}{288.827}{186.133}{generously}%
\TextAt{8}{346.003}{186.133}{gives}%
\TextAt{8}{374.236}{186.133}{away}%
\TextAt{8}{403.066}{186.133}{the}%
\TextAt{8}{422.775}{186.133}{pears}%
\TextAt{8}{453.279}{186.133}{from}%
\TextAt{8}{480.471}{186.133}{his}%
\TextAt{8}{498.291}{186.133}{garden}%
\TextAt{8}{79.226}{169.087}{to}%
\TextAt{8}{93.579}{169.087}{the}%
\TextAt{8}{113.778}{169.087}{children.}%
\TextAt{8}{163.002}{169.087}{When}%
\TextAt{8}{197.177}{169.087}{he}%
\TextAt{8}{212.821}{169.087}{dies}%
\TextAt{8}{236.331}{169.087}{he}%
\TextAt{8}{251.975}{169.087}{asks}%
\TextAt{8}{277.184}{169.087}{for}%
\TextAt{8}{295.111}{169.087}{a}%
\TextAt{8}{304.909}{169.087}{pear}%
\TextAt{8}{331.289}{169.087}{to}%
\TextAt{8}{345.642}{169.087}{be}%
\TextAt{8}{361.609}{169.087}{put}%
\TextAt{8}{383.111}{169.087}{in}%
\TextAt{8}{396.807}{169.087}{his}%
\TextAt{8}{415.117}{169.087}{grave,}%
\TextAt{8}{449.326}{169.087}{so}%
\TextAt{8}{463.739}{169.087}{that}%
\TextAt{8}{489.15}{169.087}{later}%
\TextAt{8}{516.51}{169.087}{the}%
\TextAt{8}{79.226}{152.041}{children}%
\TextAt{8}{123.771}{152.041}{can}%
\TextAt{8}{145.231}{152.041}{pick}%
\TextAt{8}{169.942}{152.041}{the}%
\TextAt{8}{190.098}{152.041}{pears}%
\TextAt{8}{221.05}{152.041}{from}%
\TextAt{8}{248.691}{152.041}{the}%
\TextAt{8}{268.847}{152.041}{growing}%
\TextAt{8}{312.758}{152.041}{tree.}%
\TextAt{8}{97.159}{132.006}{This}%
\TextAt{8}{124.065}{132.006}{is}%
\TextAt{8}{136.014}{132.006}{-}%
\TextAt{8}{143.995}{132.006}{I}%
\TextAt{8}{152.298}{132.006}{believe}%
\TextAt{8}{191.171}{132.006}{-}%
\TextAt{8}{199.151}{132.006}{a}%
\TextAt{8}{209.093}{132.006}{good}%
\TextAt{8}{237.899}{132.006}{way}%
\TextAt{8}{261.828}{132.006}{of}%
\TextAt{8}{275.344}{132.006}{thinking}%
\TextAt{8}{322.035}{132.006}{of}%
\TextAt{8}{335.551}{132.006}{Neukirch’s}%
\TextAt{8}{393.349}{132.006}{book:}%
\TextAt{8}{427.215}{132.006}{There}%
\TextAt{8}{461.21}{132.006}{are}%
\TextAt{8}{480.907}{132.006}{seeds}%
\TextAt{8}{511.124}{132.006}{in}%
\TextAt{8}{524.963}{132.006}{it}%
\TextAt{8}{79.226}{114.96}{for}%
\TextAt{8}{97.111}{114.96}{a}%
\TextAt{8}{106.866}{114.96}{tree}%
\TextAt{8}{130.275}{114.96}{to}%
\TextAt{8}{144.585}{114.96}{grow}%
\TextAt{8}{172.883}{114.96}{from}%
\TextAt{8}{200.523}{114.96}{which}%
\TextAt{8}{234.01}{114.96}{the}%
\TextAt{8}{254.166}{114.96}{“children”}%
\TextAt{8}{309.758}{114.96}{can}%
\TextAt{8}{331.217}{114.96}{pick}%
\TextAt{8}{355.929}{114.96}{fruits}%
\TextAt{8}{386.88}{114.96}{in}%
\TextAt{8}{400.533}{114.96}{the}%
\TextAt{8}{420.69}{114.96}{time}%
\TextAt{8}{447.35}{114.96}{to}%
\TextAt{8}{461.66}{114.96}{come.}%
\TextAt{8}{480.992}{84.963}{G.}%
\TextAt{8}{497.322}{84.963}{Harder}%
\TextAt{8}{301.12}{43.033}{iii}%
\Destination{79.23}{773.075}{nav.chapterz2az.3}%
\Destination{78.23}{809.815}{nav.page.iii}%
\Destination{79.23}{773.075}{outline.2}%
}

% Translator’s Note
% Page 4 (body page front matter)
\Page{%
\TextAt{1}{79.226}{673.054}{Translator’s}%
\TextAt{1}{233.327}{673.054}{Note}%
\TextAt{8}{79.226}{613.169}{When}%
\TextAt{8}{113.188}{613.169}{I}%
\TextAt{8}{121.134}{613.169}{first}%
\TextAt{8}{145.089}{613.169}{accepted}%
\TextAt{8}{193.038}{613.169}{Jürgen}%
\TextAt{8}{231.399}{613.169}{Neukirch’s}%
\TextAt{8}{288.84}{613.169}{request}%
\TextAt{8}{329.377}{613.169}{to}%
\TextAt{8}{343.517}{613.169}{translate}%
\TextAt{8}{392.195}{613.169}{his}%
\TextAt{8}{410.292}{613.169}{Algebraische}%
\TextAt{8}{478.469}{613.169}{Zahlenthe-}%
\TextAt{8}{79.226}{596.123}{orie,}%
\TextAt{8}{105.134}{596.123}{back}%
\TextAt{8}{132.319}{596.123}{in}%
\TextAt{8}{145.84}{596.123}{1991,}%
\TextAt{8}{176.315}{596.123}{no-one}%
\TextAt{8}{213.901}{596.123}{imagined}%
\TextAt{8}{263.848}{596.123}{that}%
\TextAt{8}{289.085}{596.123}{he}%
\TextAt{8}{304.553}{596.123}{would}%
\TextAt{8}{338.565}{596.123}{not}%
\TextAt{8}{359.246}{596.123}{live}%
\TextAt{8}{380.573}{596.123}{to}%
\TextAt{8}{394.751}{596.123}{see}%
\TextAt{8}{413.531}{596.123}{the}%
\TextAt{8}{433.555}{596.123}{English}%
\TextAt{8}{475.265}{596.123}{edition.}%
\TextAt{8}{518.794}{596.123}{He}%
\TextAt{8}{79.226}{579.077}{did}%
\TextAt{8}{98.647}{579.077}{see}%
\TextAt{8}{116.825}{579.077}{the}%
\TextAt{8}{136.246}{579.077}{raw}%
\TextAt{8}{157.951}{579.077}{version}%
\TextAt{8}{196.955}{579.077}{of}%
\TextAt{8}{209.55}{579.077}{the}%
\TextAt{8}{228.971}{579.077}{translation}%
\TextAt{8}{287.498}{579.077}{(I}%
\TextAt{8}{299.434}{579.077}{gave}%
\TextAt{8}{325.048}{579.077}{him}%
\TextAt{8}{347.721}{579.077}{the}%
\TextAt{8}{367.142}{579.077}{last}%
\TextAt{8}{388.584}{579.077}{chapters}%
\TextAt{8}{434.414}{579.077}{in}%
\TextAt{8}{447.332}{579.077}{the}%
\TextAt{8}{466.753}{579.077}{Fall}%
\TextAt{8}{488.936}{579.077}{of}%
\TextAt{8}{501.531}{579.077}{1996),}%
\TextAt{8}{79.226}{562.031}{and}%
\TextAt{8}{101.989}{562.031}{he}%
\TextAt{8}{117.59}{562.031}{still}%
\TextAt{8}{140.413}{562.031}{had}%
\TextAt{8}{163.175}{562.031}{time}%
\TextAt{8}{189.835}{562.031}{to}%
\TextAt{8}{204.146}{562.031}{go}%
\TextAt{8}{219.759}{562.031}{carefully}%
\TextAt{8}{267.233}{562.031}{through}%
\TextAt{8}{311.467}{562.031}{the}%
\TextAt{8}{331.624}{562.031}{first}%
\TextAt{8}{355.749}{562.031}{four}%
\TextAt{8}{380.138}{562.031}{chapters}%
\TextAt{8}{426.703}{562.031}{of}%
\TextAt{8}{440.033}{562.031}{it.}%
\TextAt{8}{97.159}{541.996}{The}%
\TextAt{8}{121.7}{541.996}{bulk}%
\TextAt{8}{148.525}{541.996}{of}%
\TextAt{8}{162.342}{541.996}{the}%
\TextAt{8}{182.986}{541.996}{text}%
\TextAt{8}{207.862}{541.996}{consists}%
\TextAt{8}{251.46}{541.996}{of}%
\TextAt{8}{265.277}{541.996}{detailed}%
\TextAt{8}{309.987}{541.996}{technical}%
\TextAt{8}{359.574}{541.996}{mathematical}%
\TextAt{8}{433.562}{541.996}{prose}%
\TextAt{8}{464.678}{541.996}{and}%
\TextAt{8}{487.928}{541.996}{was}%
\TextAt{8}{510.915}{541.996}{thus}%
\TextAt{8}{79.226}{524.95}{straightforward}%
\TextAt{8}{161.897}{524.95}{to}%
\TextAt{8}{176.518}{524.95}{translate,}%
\TextAt{8}{229.008}{524.95}{even}%
\TextAt{8}{255.98}{524.95}{though}%
\TextAt{8}{295.97}{524.95}{the}%
\TextAt{8}{316.438}{524.95}{author’s}%
\TextAt{8}{362.346}{524.95}{desire}%
\TextAt{8}{395.881}{524.95}{to}%
\TextAt{8}{410.503}{524.95}{integrate}%
\TextAt{8}{459.926}{524.95}{involved}%
\TextAt{8}{506.097}{524.95}{argu-}%
\TextAt{8}{79.226}{507.904}{ments}%
\TextAt{8}{114.01}{507.904}{and}%
\TextAt{8}{137.353}{507.904}{displayed}%
\TextAt{8}{189.054}{507.904}{formulae}%
\TextAt{8}{237.766}{507.904}{into}%
\TextAt{8}{262.089}{507.904}{comprehensive}%
\TextAt{8}{340.773}{507.904}{sentences}%
\TextAt{8}{392.521}{507.904}{could}%
\TextAt{8}{424.317}{507.904}{not}%
\TextAt{8}{445.711}{507.904}{simply}%
\TextAt{8}{483.747}{507.904}{be}%
\TextAt{8}{500.251}{507.904}{copied}%
\TextAt{8}{79.226}{490.858}{into}%
\TextAt{8}{104.035}{490.858}{English.}%
\TextAt{8}{153.627}{490.858}{However,}%
\TextAt{8}{205.362}{490.858}{Jürgen}%
\TextAt{8}{244.959}{490.858}{Neukirch}%
\TextAt{8}{295.771}{490.858}{had}%
\TextAt{8}{319.599}{490.858}{peppered}%
\TextAt{8}{371.379}{490.858}{his}%
\TextAt{8}{390.712}{490.858}{book}%
\TextAt{8}{420.722}{490.858}{with}%
\TextAt{8}{448.448}{490.858}{more}%
\TextAt{8}{478.78}{490.858}{meditative}%
\TextAt{8}{79.226}{473.812}{paragraphs}%
\TextAt{8}{140.518}{473.812}{which}%
\TextAt{8}{174.732}{473.812}{make}%
\TextAt{8}{206.028}{473.812}{rather}%
\TextAt{8}{241.88}{473.812}{serious}%
\TextAt{8}{281.103}{473.812}{use}%
\TextAt{8}{302.046}{473.812}{of}%
\TextAt{8}{316.103}{473.812}{the}%
\TextAt{8}{336.987}{473.812}{German}%
\TextAt{8}{382.666}{473.812}{language.}%
\TextAt{8}{438.187}{473.812}{When}%
\TextAt{8}{473.046}{473.812}{I}%
\TextAt{8}{481.89}{473.812}{started}%
\TextAt{8}{522.357}{473.812}{to}%
\TextAt{8}{79.226}{456.766}{work}%
\TextAt{8}{108.509}{456.766}{on}%
\TextAt{8}{125.43}{456.766}{the}%
\TextAt{8}{146.249}{456.766}{translation,}%
\TextAt{8}{209.59}{456.766}{he}%
\TextAt{8}{225.854}{456.766}{warned}%
\TextAt{8}{267.164}{456.766}{me}%
\TextAt{8}{286.68}{456.766}{that}%
\TextAt{8}{312.711}{456.766}{in}%
\TextAt{8}{327.026}{456.766}{every}%
\TextAt{8}{358.58}{456.766}{one}%
\TextAt{8}{380.702}{456.766}{of}%
\TextAt{8}{394.694}{456.766}{these}%
\TextAt{8}{425.328}{456.766}{passages,}%
\TextAt{8}{476.427}{456.766}{he}%
\TextAt{8}{492.691}{456.766}{was}%
\TextAt{8}{515.853}{456.766}{not}%
\TextAt{8}{79.226}{439.72}{seeking}%
\TextAt{8}{119.793}{439.72}{poetic}%
\TextAt{8}{154.766}{439.72}{beauty,}%
\TextAt{8}{195.624}{439.72}{but}%
\TextAt{8}{216.943}{439.72}{only}%
\TextAt{8}{242.495}{439.72}{the}%
\TextAt{8}{262.511}{439.72}{precisely}%
\TextAt{8}{310.227}{439.72}{adequate}%
\TextAt{8}{359.845}{439.72}{expression}%
\TextAt{8}{416.085}{439.72}{of}%
\TextAt{8}{429.275}{439.72}{an}%
\TextAt{8}{445.394}{439.72}{idea.}%
\TextAt{8}{474.61}{439.72}{It}%
\TextAt{8}{487.142}{439.72}{is}%
\TextAt{8}{498.766}{439.72}{for}%
\TextAt{8}{516.51}{439.72}{the}%
\TextAt{8}{79.226}{422.674}{reader}%
\TextAt{8}{114.996}{422.674}{to}%
\TextAt{8}{129.306}{422.674}{judge}%
\TextAt{8}{160.844}{422.674}{whether}%
\TextAt{8}{205.712}{422.674}{I}%
\TextAt{8}{213.829}{422.674}{managed}%
\TextAt{8}{263.264}{422.674}{to}%
\TextAt{8}{277.574}{422.674}{render}%
\TextAt{8}{313.99}{422.674}{his}%
\TextAt{8}{332.257}{422.674}{ideas}%
\TextAt{8}{361.583}{422.674}{faithfully.}%
\TextAt{8}{97.159}{402.64}{There}%
\TextAt{8}{132.587}{402.64}{is}%
\TextAt{8}{145.97}{402.64}{one}%
\TextAt{8}{169.049}{402.64}{neologism}%
\TextAt{8}{224.718}{402.64}{that}%
\TextAt{8}{251.706}{402.64}{I}%
\TextAt{8}{261.441}{402.64}{propose}%
\TextAt{8}{306.374}{402.64}{in}%
\TextAt{8}{321.646}{402.64}{this}%
\TextAt{8}{346.087}{402.64}{translation,}%
\TextAt{8}{410.625}{402.64}{with}%
\TextAt{8}{438.904}{402.64}{Jürgen}%
\TextAt{8}{479.055}{402.64}{Neukirch’s}%
\TextAt{8}{79.226}{385.594}{blessing:}%
\TextAt{8}{129.67}{385.594}{I}%
\TextAt{8}{138.886}{385.594}{call}%
\TextAt{8}{161.444}{385.594}{replete}%
\TextAt{8}{200.906}{385.594}{divisor,}%
\TextAt{8}{243.645}{385.594}{ideal,}%
\TextAt{8}{276.233}{385.594}{etc.,}%
\TextAt{8}{302.963}{385.594}{what}%
\TextAt{8}{333.328}{385.594}{is}%
\TextAt{8}{346.19}{385.594}{usually}%
\TextAt{8}{387.351}{385.594}{called}%
\TextAt{8}{421.613}{385.594}{Arakelov}%
\TextAt{8}{471.824}{385.594}{divisor,}%
\TextAt{8}{514.562}{385.594}{etc.}%
\TextAt{8}{79.226}{368.548}{(a}%
\TextAt{8}{93.825}{368.548}{terminology}%
\TextAt{8}{158.839}{368.548}{that}%
\TextAt{8}{184.497}{368.548}{Neukirch}%
\TextAt{8}{234.531}{368.548}{had}%
\TextAt{8}{257.583}{368.548}{avoided}%
\TextAt{8}{300.48}{368.548}{in}%
\TextAt{8}{314.421}{368.548}{the}%
\TextAt{8}{334.866}{368.548}{German}%
\TextAt{8}{380.107}{368.548}{edition).}%
\TextAt{8}{429.101}{368.548}{Time}%
\TextAt{8}{459.947}{368.548}{will}%
\TextAt{8}{482.341}{368.548}{deliver}%
\TextAt{8}{520.348}{368.548}{its}%
\TextAt{8}{79.226}{351.502}{verdict.}%
\TextAt{8}{97.159}{331.467}{I}%
\TextAt{8}{106.209}{331.467}{am}%
\TextAt{8}{126.652}{331.467}{much}%
\TextAt{8}{158.8}{331.467}{indebted}%
\TextAt{8}{207.853}{331.467}{to}%
\TextAt{8}{223.096}{331.467}{Frazer}%
\TextAt{8}{259.955}{331.467}{Jarvis}%
\TextAt{8}{295.259}{331.467}{for}%
\TextAt{8}{314.077}{331.467}{going}%
\TextAt{8}{346.237}{331.467}{through}%
\TextAt{8}{391.404}{331.467}{my}%
\TextAt{8}{411.848}{331.467}{entire}%
\TextAt{8}{445.622}{331.467}{manuscript,}%
\TextAt{8}{510.915}{331.467}{thus}%
\TextAt{8}{79.226}{314.421}{saving}%
\TextAt{8}{116.228}{314.421}{the}%
\TextAt{8}{137.544}{314.421}{English}%
\TextAt{8}{180.547}{314.421}{language}%
\TextAt{8}{230.495}{314.421}{from}%
\TextAt{8}{259.296}{314.421}{various}%
\TextAt{8}{300.529}{314.421}{infractions.}%
\TextAt{8}{366.183}{314.421}{But}%
\TextAt{8}{390.584}{314.421}{needless}%
\TextAt{8}{436.731}{314.421}{to}%
\TextAt{8}{452.201}{314.421}{say,}%
\TextAt{8}{476.151}{314.421}{I}%
\TextAt{8}{485.427}{314.421}{alone}%
\TextAt{8}{517.156}{314.421}{am}%
\TextAt{8}{79.226}{297.375}{responsible}%
\TextAt{8}{139.504}{297.375}{for}%
\TextAt{8}{157.389}{297.375}{all}%
\TextAt{8}{173.648}{297.375}{deficiencies}%
\TextAt{8}{234.177}{297.375}{that}%
\TextAt{8}{259.546}{297.375}{remain.}%
\TextAt{8}{97.159}{277.34}{After}%
\TextAt{8}{128.276}{277.34}{Jürgen}%
\TextAt{8}{167.366}{277.34}{Neukirch’s}%
\TextAt{8}{225.537}{277.34}{untimely}%
\TextAt{8}{274.874}{277.34}{death}%
\TextAt{8}{307.951}{277.34}{early}%
\TextAt{8}{337.454}{277.34}{in}%
\TextAt{8}{351.666}{277.34}{1997,}%
\TextAt{8}{382.946}{277.34}{it}%
\TextAt{8}{395.209}{277.34}{was}%
\TextAt{8}{418.268}{277.34}{Ms}%
\TextAt{8}{438.063}{277.34}{Eva-Maria}%
\TextAt{8}{496.019}{277.34}{Strobel}%
\TextAt{8}{79.226}{260.294}{who}%
\TextAt{8}{104.433}{260.294}{took}%
\TextAt{8}{131.6}{260.294}{it}%
\TextAt{8}{143.8}{260.294}{upon}%
\TextAt{8}{173.884}{260.294}{herself}%
\TextAt{8}{211.178}{260.294}{to}%
\TextAt{8}{225.983}{260.294}{finish}%
\TextAt{8}{257.753}{260.294}{as}%
\TextAt{8}{272.619}{260.294}{best}%
\TextAt{8}{298.208}{260.294}{she}%
\TextAt{8}{318.92}{260.294}{could}%
\TextAt{8}{350.63}{260.294}{what}%
\TextAt{8}{380.392}{260.294}{Jürgen}%
\TextAt{8}{419.419}{260.294}{Neukirch}%
\TextAt{8}{469.66}{260.294}{had}%
\TextAt{8}{492.918}{260.294}{to}%
\TextAt{8}{507.723}{260.294}{leave}%
\TextAt{8}{79.226}{243.248}{undone.}%
\TextAt{8}{126.668}{243.248}{She}%
\TextAt{8}{149.414}{243.248}{had}%
\TextAt{8}{172.817}{243.248}{already}%
\TextAt{8}{214.762}{243.248}{applied}%
\TextAt{8}{256.373}{243.248}{her}%
\TextAt{8}{277.169}{243.248}{infinite}%
\TextAt{8}{317.477}{243.248}{care}%
\TextAt{8}{342.828}{243.248}{and}%
\TextAt{8}{366.231}{243.248}{patience}%
\TextAt{8}{413.042}{243.248}{to}%
\TextAt{8}{427.993}{243.248}{the}%
\TextAt{8}{448.79}{243.248}{original}%
\TextAt{8}{491.715}{243.248}{German}%
\TextAt{8}{79.226}{226.202}{book,}%
\TextAt{8}{113.082}{226.202}{and}%
\TextAt{8}{137.172}{226.202}{she}%
\TextAt{8}{158.717}{226.202}{had}%
\TextAt{8}{182.808}{226.202}{assisted}%
\TextAt{8}{227.246}{226.202}{Jürgen}%
\TextAt{8}{267.106}{226.202}{Neukirch}%
\TextAt{8}{318.179}{226.202}{in}%
\TextAt{8}{333.16}{226.202}{proofreading}%
\TextAt{8}{402.788}{226.202}{the}%
\TextAt{8}{424.273}{226.202}{first}%
\TextAt{8}{449.727}{226.202}{four}%
\TextAt{8}{475.443}{226.202}{chapters}%
\TextAt{8}{523.337}{226.202}{of}%
\TextAt{8}{79.226}{209.156}{the}%
\TextAt{8}{99.594}{209.156}{translation.}%
\TextAt{8}{164.042}{209.156}{Without}%
\TextAt{8}{211.404}{209.156}{her}%
\TextAt{8}{231.772}{209.156}{knowledge,}%
\TextAt{8}{291.872}{209.156}{responsibility,}%
\TextAt{8}{366.402}{209.156}{and}%
\TextAt{8}{389.377}{209.156}{energy,}%
\TextAt{8}{429.308}{209.156}{this}%
\TextAt{8}{452.342}{209.156}{book}%
\TextAt{8}{481.497}{209.156}{would}%
\TextAt{8}{515.853}{209.156}{not}%
\TextAt{8}{79.226}{192.11}{be}%
\TextAt{8}{95.585}{192.11}{what}%
\TextAt{8}{125.285}{192.11}{it}%
\TextAt{8}{137.423}{192.11}{is.}%
\TextAt{8}{155.04}{192.11}{In}%
\TextAt{8}{170.095}{192.11}{particular,}%
\TextAt{8}{227.879}{192.11}{a}%
\TextAt{8}{238.069}{192.11}{fair}%
\TextAt{8}{259.64}{192.11}{number}%
\TextAt{8}{302.67}{192.11}{of}%
\TextAt{8}{316.434}{192.11}{small}%
\TextAt{8}{347.497}{192.11}{corrections}%
\TextAt{8}{407.181}{192.11}{and}%
\TextAt{8}{430.378}{192.11}{modifications}%
\TextAt{8}{502.746}{192.11}{of}%
\TextAt{8}{516.51}{192.11}{the}%
\TextAt{8}{79.226}{175.064}{German}%
\TextAt{8}{124.068}{175.064}{original}%
\TextAt{8}{166.244}{175.064}{that}%
\TextAt{8}{191.504}{175.064}{had}%
\TextAt{8}{214.157}{175.064}{been}%
\TextAt{8}{241.676}{175.064}{accumulated}%
\TextAt{8}{309.532}{175.064}{thanks}%
\TextAt{8}{347.535}{175.064}{to}%
\TextAt{8}{361.736}{175.064}{attentive}%
\TextAt{8}{410.739}{175.064}{readers,}%
\TextAt{8}{454.287}{175.064}{were}%
\TextAt{8}{481.161}{175.064}{taken}%
\TextAt{8}{512.924}{175.064}{into}%
\TextAt{8}{79.226}{158.018}{account}%
\TextAt{8}{122.01}{158.018}{for}%
\TextAt{8}{139.426}{158.018}{this}%
\TextAt{8}{161.779}{158.018}{English}%
\TextAt{8}{203.152}{158.018}{edition.}%
\TextAt{8}{246.568}{158.018}{Kay}%
\TextAt{8}{270.81}{158.018}{Wingberg}%
\TextAt{8}{324.318}{158.018}{graciously}%
\TextAt{8}{378.878}{158.018}{helped}%
\TextAt{8}{415.793}{158.018}{to}%
\TextAt{8}{429.633}{158.018}{check}%
\TextAt{8}{460.702}{158.018}{a}%
\TextAt{8}{469.987}{158.018}{few}%
\TextAt{8}{490.643}{158.018}{of}%
\TextAt{8}{503.503}{158.018}{them.}%
\TextAt{8}{79.226}{140.972}{We}%
\TextAt{8}{99.371}{140.972}{sincerely}%
\TextAt{8}{147.227}{140.972}{hope}%
\TextAt{8}{175.513}{140.972}{that}%
\TextAt{8}{200.882}{140.972}{the}%
\TextAt{8}{221.038}{140.972}{book}%
\TextAt{8}{249.982}{140.972}{published}%
\TextAt{8}{302.716}{140.972}{here}%
\TextAt{8}{328.073}{140.972}{would}%
\TextAt{8}{362.217}{140.972}{have}%
\TextAt{8}{389.212}{140.972}{made}%
\TextAt{8}{420.427}{140.972}{its}%
\TextAt{8}{436.746}{140.972}{author}%
\TextAt{8}{474.476}{140.972}{happy.}%
\TextAt{8}{97.159}{120.938}{Hearty}%
\TextAt{8}{134.975}{120.938}{thanks}%
\TextAt{8}{172.205}{120.938}{go}%
\TextAt{8}{186.935}{120.938}{to}%
\TextAt{8}{200.362}{120.938}{Raymond}%
\TextAt{8}{252.333}{120.938}{Seroul,}%
\TextAt{8}{290.648}{120.938}{Strasbourg,}%
\TextAt{8}{352.778}{120.938}{for}%
\TextAt{8}{369.78}{120.938}{applying}%
\TextAt{8}{416.706}{120.938}{his}%
\TextAt{8}{434.09}{120.938}{wonderful}%
\TextAt{8}{487.184}{120.938}{expertise}%
\TextAt{8}{79.226}{103.892}{of}%
\TextAt{8}{92.556}{103.892}{T}%
\TextAt{8}{99.016}{101.315}{E}%
\TextAt{8}{105.483}{103.892}{X}%
\TextAt{8}{118.156}{103.892}{to}%
\TextAt{8}{132.466}{103.892}{the}%
\TextAt{8}{152.622}{103.892}{final}%
\TextAt{8}{178.637}{103.892}{preparation}%
\TextAt{8}{241.736}{103.892}{of}%
\TextAt{8}{255.066}{103.892}{the}%
\TextAt{8}{275.223}{103.892}{camera-ready}%
\TextAt{8}{347.743}{103.892}{manuscript.}%
\TextAt{8}{97.159}{83.857}{Thanks}%
\TextAt{8}{140.484}{83.857}{go}%
\TextAt{8}{157.413}{83.857}{to}%
\TextAt{8}{173.038}{83.857}{the}%
\TextAt{8}{194.51}{83.857}{Springer}%
\TextAt{8}{242.653}{83.857}{staff}%
\TextAt{8}{269.72}{83.857}{for}%
\TextAt{8}{288.92}{83.857}{seeing}%
\TextAt{8}{324.761}{83.857}{this}%
\TextAt{8}{348.899}{83.857}{project}%
\TextAt{8}{389.881}{83.857}{through}%
\TextAt{8}{435.43}{83.857}{until}%
\TextAt{8}{464.386}{83.857}{it}%
\TextAt{8}{477.405}{83.857}{was}%
\TextAt{8}{501.22}{83.857}{finally}%
\TextAt{8}{301.281}{43.033}{iv}%
\Destination{79.23}{773.075}{nav.chapterz2az.4}%
\Destination{78.23}{809.815}{nav.page.iv}%
\Destination{79.23}{773.075}{outline.3}%
}

% Page 5 (body page front matter)
\Page{%
\TextAt{2}{79.226}{796.943}{Translator’s}%
\TextAt{2}{139.717}{796.943}{Note}%
\Vector{q 0 G 0 g 0.299 w 7.226 18.729 m 460.769 18.729 l S Q}%
\TextAt{8}{79.226}{761.124}{completed.}%
\TextAt{8}{139.552}{761.124}{Among}%
\TextAt{8}{179.744}{761.124}{them}%
\TextAt{8}{209.2}{761.124}{I}%
\TextAt{8}{216.862}{761.124}{want}%
\TextAt{8}{245.028}{761.124}{to}%
\TextAt{8}{258.883}{761.124}{thank}%
\TextAt{8}{291.926}{761.124}{especially}%
\TextAt{8}{344.205}{761.124}{Joachim}%
\TextAt{8}{389.765}{761.124}{Heinze}%
\TextAt{8}{427.339}{761.124}{for}%
\TextAt{8}{444.768}{761.124}{interfering}%
\TextAt{8}{500.897}{761.124}{rarely,}%
\TextAt{8}{79.226}{744.078}{but}%
\TextAt{8}{100.686}{744.078}{effectively,}%
\TextAt{8}{157.58}{744.078}{over}%
\TextAt{8}{182.626}{744.078}{the}%
\TextAt{8}{202.783}{744.078}{years,}%
\TextAt{8}{236.018}{744.078}{with}%
\TextAt{8}{262.678}{744.078}{the}%
\TextAt{8}{282.835}{744.078}{realization}%
\TextAt{8}{340.076}{744.078}{of}%
\TextAt{8}{353.406}{744.078}{this}%
\TextAt{8}{376.228}{744.078}{translation.}%
\TextAt{8}{410.086}{711.092}{Strasbourg,}%
\TextAt{8}{472.922}{711.092}{March}%
\TextAt{8}{509.337}{711.092}{1999}%
\TextAt{8}{424.802}{694.046}{Norbert}%
\TextAt{8}{469.025}{694.046}{Schappacher}%
\TextAt{8}{302.907}{43.033}{v}%
\Destination{78.23}{809.815}{nav.page.v}%
}

% Preface to the German Edition
% Page 6 (body page front matter)
\Page{%
\TextAt{1}{79.226}{673.054}{Preface}%
\TextAt{1}{178.895}{673.054}{to}%
\TextAt{1}{213.002}{673.054}{the}%
\TextAt{1}{261.361}{673.054}{German}%
\TextAt{1}{368.987}{673.054}{Edition}%
\TextAt{8}{79.226}{613.169}{Number}%
\TextAt{8}{123.824}{613.169}{Theory,}%
\TextAt{8}{166.269}{613.169}{among}%
\TextAt{8}{203.407}{613.169}{the}%
\TextAt{8}{222.971}{613.169}{mathematical}%
\TextAt{8}{295.878}{613.169}{disciplines,}%
\TextAt{8}{354.702}{613.169}{occupies}%
\TextAt{8}{400.663}{613.169}{an}%
\TextAt{8}{416.329}{613.169}{idealized}%
\TextAt{8}{463.856}{613.169}{position,}%
\TextAt{8}{511.25}{613.169}{sim-}%
\TextAt{8}{79.226}{596.123}{ilar}%
\TextAt{8}{100.01}{596.123}{to}%
\TextAt{8}{114.29}{596.123}{the}%
\TextAt{8}{134.416}{596.123}{one}%
\TextAt{8}{155.846}{596.123}{that}%
\TextAt{8}{181.184}{596.123}{mathematics}%
\TextAt{8}{250.159}{596.123}{holds}%
\TextAt{8}{280.758}{596.123}{among}%
\TextAt{8}{318.459}{596.123}{the}%
\TextAt{8}{338.585}{596.123}{sciences.}%
\TextAt{8}{386.81}{596.123}{Under}%
\TextAt{8}{422.215}{596.123}{no}%
\TextAt{8}{438.444}{596.123}{obligation}%
\TextAt{8}{493.061}{596.123}{to}%
\TextAt{8}{507.341}{596.123}{serve}%
\TextAt{8}{79.226}{579.077}{needs}%
\TextAt{8}{111.146}{579.077}{that}%
\TextAt{8}{136.514}{579.077}{do}%
\TextAt{8}{152.772}{579.077}{not}%
\TextAt{8}{173.585}{579.077}{originate}%
\TextAt{8}{222.373}{579.077}{within}%
\TextAt{8}{258.788}{579.077}{itself,}%
\TextAt{8}{290.384}{579.077}{it}%
\TextAt{8}{302.088}{579.077}{is}%
\TextAt{8}{313.851}{579.077}{essentially}%
\TextAt{8}{369.908}{579.077}{autonomous}%
\TextAt{8}{435.672}{579.077}{in}%
\TextAt{8}{449.324}{579.077}{setting}%
\TextAt{8}{487.759}{579.077}{its}%
\TextAt{8}{504.077}{579.077}{goals,}%
\TextAt{8}{79.226}{562.031}{and}%
\TextAt{8}{101.997}{562.031}{thus}%
\TextAt{8}{127.756}{562.031}{manages}%
\TextAt{8}{175.31}{562.031}{to}%
\TextAt{8}{189.628}{562.031}{protect}%
\TextAt{8}{229.961}{562.031}{its}%
\TextAt{8}{246.288}{562.031}{undisturbed}%
\TextAt{8}{311.714}{562.031}{harmony.}%
\TextAt{8}{364.098}{562.031}{The}%
\TextAt{8}{388.16}{562.031}{possibility}%
\TextAt{8}{443.903}{562.031}{of}%
\TextAt{8}{457.241}{562.031}{formulating}%
\TextAt{8}{520.348}{562.031}{its}%
\TextAt{8}{79.226}{544.985}{basic}%
\TextAt{8}{108.319}{544.985}{problems}%
\TextAt{8}{158.226}{544.985}{simply,}%
\TextAt{8}{197.767}{544.985}{the}%
\TextAt{8}{217.69}{544.985}{peculiar}%
\TextAt{8}{262.002}{544.985}{clarity}%
\TextAt{8}{298.197}{544.985}{of}%
\TextAt{8}{311.294}{544.985}{its}%
\TextAt{8}{327.38}{544.985}{statements,}%
\TextAt{8}{389.432}{544.985}{the}%
\TextAt{8}{409.356}{544.985}{arcane}%
\TextAt{8}{446.196}{544.985}{touch}%
\TextAt{8}{478.158}{544.985}{in}%
\TextAt{8}{491.578}{544.985}{its}%
\TextAt{8}{507.664}{544.985}{laws,}%
\TextAt{8}{79.226}{527.939}{be}%
\TextAt{8}{95.754}{527.939}{they}%
\TextAt{8}{122.694}{527.939}{discovered}%
\TextAt{8}{179.618}{527.939}{or}%
\TextAt{8}{194.532}{527.939}{undiscovered,}%
\TextAt{8}{267.866}{527.939}{merely}%
\TextAt{8}{306.511}{527.939}{divined;}%
\TextAt{8}{351.961}{527.939}{last}%
\TextAt{8}{374.741}{527.939}{but}%
\TextAt{8}{396.804}{527.939}{not}%
\TextAt{8}{418.221}{527.939}{least,}%
\TextAt{8}{449.605}{527.939}{the}%
\TextAt{8}{470.364}{527.939}{charm}%
\TextAt{8}{506.415}{527.939}{of}%
\TextAt{8}{520.348}{527.939}{its}%
\TextAt{8}{79.226}{510.893}{particularly}%
\TextAt{8}{142.656}{510.893}{satisfactory}%
\TextAt{8}{205.236}{510.893}{ways}%
\TextAt{8}{233.601}{510.893}{of}%
\TextAt{8}{246.939}{510.893}{reasoning}%
\TextAt{8}{299.047}{510.893}{-}%
\TextAt{8}{306.849}{510.893}{all}%
\TextAt{8}{323.115}{510.893}{these}%
\TextAt{8}{353.094}{510.893}{features}%
\TextAt{8}{397.061}{510.893}{have}%
\TextAt{8}{424.063}{510.893}{at}%
\TextAt{8}{438.381}{510.893}{all}%
\TextAt{8}{454.647}{510.893}{times}%
\TextAt{8}{485.929}{510.893}{attracted}%
\TextAt{8}{79.226}{493.847}{to}%
\TextAt{8}{93.536}{493.847}{number}%
\TextAt{8}{136.133}{493.847}{theory}%
\TextAt{8}{172.883}{493.847}{a}%
\TextAt{8}{182.638}{493.847}{community}%
\TextAt{8}{243.454}{493.847}{of}%
\TextAt{8}{256.784}{493.847}{dedicated}%
\TextAt{8}{309.459}{493.847}{followers.}%
\TextAt{8}{97.159}{473.812}{But}%
\TextAt{8}{121.35}{473.812}{different}%
\TextAt{8}{168.471}{473.812}{number}%
\TextAt{8}{212.018}{473.812}{theorists}%
\TextAt{8}{260.574}{473.812}{may}%
\TextAt{8}{286.893}{473.812}{dedicate}%
\TextAt{8}{334.015}{473.812}{themselves}%
\TextAt{8}{393.617}{473.812}{differently}%
\TextAt{8}{450.171}{473.812}{to}%
\TextAt{8}{465.432}{473.812}{their}%
\TextAt{8}{494.346}{473.812}{science.}%
\TextAt{8}{79.226}{456.766}{Some}%
\TextAt{8}{111.153}{456.766}{will}%
\TextAt{8}{133.971}{456.766}{push}%
\TextAt{8}{162.706}{456.766}{the}%
\TextAt{8}{183.574}{456.766}{theoretical}%
\TextAt{8}{242.173}{456.766}{development}%
\TextAt{8}{310.851}{456.766}{only}%
\TextAt{8}{337.255}{456.766}{as}%
\TextAt{8}{352.337}{456.766}{far}%
\TextAt{8}{370.934}{456.766}{as}%
\TextAt{8}{386.016}{456.766}{is}%
\TextAt{8}{398.492}{456.766}{necessary}%
\TextAt{8}{451.03}{456.766}{for}%
\TextAt{8}{469.627}{456.766}{the}%
\TextAt{8}{490.496}{456.766}{concrete}%
\TextAt{8}{79.226}{439.72}{result}%
\TextAt{8}{110.642}{439.72}{they}%
\TextAt{8}{135.818}{439.72}{desire.}%
\TextAt{8}{173.205}{439.72}{Others}%
\TextAt{8}{210.479}{439.72}{will}%
\TextAt{8}{231.423}{439.72}{strive}%
\TextAt{8}{262.194}{439.72}{for}%
\TextAt{8}{278.917}{439.72}{a}%
\TextAt{8}{287.511}{439.72}{more}%
\TextAt{8}{315.615}{439.72}{universal,}%
\TextAt{8}{367.432}{439.72}{conceptual}%
\TextAt{8}{424.803}{439.72}{clarity,}%
\TextAt{8}{462.573}{439.72}{never}%
\TextAt{8}{492.626}{439.72}{tiring}%
\TextAt{8}{523.337}{439.72}{of}%
\TextAt{8}{79.226}{422.674}{searching}%
\TextAt{8}{129.375}{422.674}{for}%
\TextAt{8}{146.288}{422.674}{the}%
\TextAt{8}{165.473}{422.674}{deep-lying}%
\TextAt{8}{220.75}{422.674}{reasons}%
\TextAt{8}{260.881}{422.674}{behind}%
\TextAt{8}{298.596}{422.674}{the}%
\TextAt{8}{317.781}{422.674}{apparent}%
\TextAt{8}{365.921}{422.674}{variety}%
\TextAt{8}{403.66}{422.674}{of}%
\TextAt{8}{416.019}{422.674}{arithmetic}%
\TextAt{8}{471.631}{422.674}{phenomena.}%
\TextAt{8}{79.226}{405.628}{Both}%
\TextAt{8}{110.199}{405.628}{attitudes}%
\TextAt{8}{161.566}{405.628}{are}%
\TextAt{8}{182.951}{405.628}{justified,}%
\TextAt{8}{233.149}{405.628}{and}%
\TextAt{8}{257.785}{405.628}{they}%
\TextAt{8}{285.996}{405.628}{grow}%
\TextAt{8}{316.168}{405.628}{particularly}%
\TextAt{8}{381.463}{405.628}{effective}%
\TextAt{8}{428.527}{405.628}{through}%
\TextAt{8}{474.635}{405.628}{the}%
\TextAt{8}{496.665}{405.628}{mutual}%
\TextAt{8}{79.226}{388.582}{inspirational}%
\TextAt{8}{147.486}{388.582}{influence}%
\TextAt{8}{196.796}{388.582}{they}%
\TextAt{8}{223.678}{388.582}{exert}%
\TextAt{8}{253.813}{388.582}{on}%
\TextAt{8}{270.617}{388.582}{one}%
\TextAt{8}{292.622}{388.582}{another.}%
\TextAt{8}{341.74}{388.582}{Several}%
\TextAt{8}{382.61}{388.582}{beautiful}%
\TextAt{8}{432.578}{388.582}{textbooks}%
\TextAt{8}{487.172}{388.582}{illustrate}%
\TextAt{8}{79.226}{371.536}{the}%
\TextAt{8}{100.895}{371.536}{success}%
\TextAt{8}{142.255}{371.536}{of}%
\TextAt{8}{157.098}{371.536}{the}%
\TextAt{8}{178.767}{371.536}{first}%
\TextAt{8}{204.406}{371.536}{attitude,}%
\TextAt{8}{254.428}{371.536}{which}%
\TextAt{8}{289.428}{371.536}{is}%
\TextAt{8}{302.705}{371.536}{oriented}%
\TextAt{8}{349.42}{371.536}{towards}%
\TextAt{8}{394.581}{371.536}{specific}%
\TextAt{8}{436.79}{371.536}{problems.}%
\TextAt{8}{496.019}{371.536}{Among}%
\TextAt{8}{79.226}{354.49}{them,}%
\TextAt{8}{113.921}{354.49}{let}%
\TextAt{8}{132.05}{354.49}{us}%
\TextAt{8}{148.291}{354.49}{pick}%
\TextAt{8}{174.227}{354.49}{out}%
\TextAt{8}{196.266}{354.49}{in}%
\TextAt{8}{211.144}{354.49}{particular}%
\TextAt{8}{266.358}{354.49}{Number}%
\TextAt{8}{312.774}{354.49}{Theory}%
\TextAt{8}{354.646}{354.49}{by}%
\TextAt{8}{372.13}{354.49}{S.I.}%
\TextAt{8}{394.48}{354.49}{Borevič}%
\TextAt{8}{438.133}{354.49}{and}%
\TextAt{8}{462.121}{354.49}{I.R.}%
\TextAt{8}{486.587}{354.49}{Šafarevič}%
\TextAt{8}{79.226}{337.445}{[14]:}%
\TextAt{8}{105.768}{337.445}{a}%
\TextAt{8}{115.27}{337.445}{book}%
\TextAt{8}{143.96}{337.445}{which}%
\TextAt{8}{177.194}{337.445}{is}%
\TextAt{8}{188.704}{337.445}{extremely}%
\TextAt{8}{242.429}{337.445}{rich}%
\TextAt{8}{265.261}{337.445}{in}%
\TextAt{8}{278.661}{337.445}{content,}%
\TextAt{8}{323.338}{337.445}{yet}%
\TextAt{8}{342.596}{337.445}{easy}%
\TextAt{8}{368.094}{337.445}{to}%
\TextAt{8}{382.152}{337.445}{read,}%
\TextAt{8}{411.215}{337.445}{and}%
\TextAt{8}{433.725}{337.445}{which}%
\TextAt{8}{466.958}{337.445}{we}%
\TextAt{8}{483.933}{337.445}{especially}%
\TextAt{8}{79.226}{320.399}{recommend}%
\TextAt{8}{141.656}{320.399}{to}%
\TextAt{8}{155.966}{320.399}{the}%
\TextAt{8}{176.123}{320.399}{reader.}%
\TextAt{8}{97.159}{300.364}{The}%
\TextAt{8}{120.437}{300.364}{present}%
\TextAt{8}{160.368}{300.364}{book}%
\TextAt{8}{188.535}{300.364}{was}%
\TextAt{8}{210.258}{300.364}{conceived}%
\TextAt{8}{262.157}{300.364}{with}%
\TextAt{8}{288.041}{300.364}{a}%
\TextAt{8}{297.02}{300.364}{different}%
\TextAt{8}{342.414}{300.364}{objective}%
\TextAt{8}{391.061}{300.364}{in}%
\TextAt{8}{403.938}{300.364}{mind.}%
\TextAt{8}{438.142}{300.364}{It}%
\TextAt{8}{450.038}{300.364}{does}%
\TextAt{8}{475.659}{300.364}{provide}%
\TextAt{8}{516.51}{300.364}{the}%
\TextAt{8}{79.226}{283.318}{student}%
\TextAt{8}{121.176}{283.318}{with}%
\TextAt{8}{147.775}{283.318}{an}%
\TextAt{8}{163.973}{283.318}{essentially}%
\TextAt{8}{219.97}{283.318}{self-contained}%
\TextAt{8}{293.457}{283.318}{introduction}%
\TextAt{8}{360.393}{283.318}{to}%
\TextAt{8}{374.643}{283.318}{the}%
\TextAt{8}{394.738}{283.318}{theory}%
\TextAt{8}{431.428}{283.318}{of}%
\TextAt{8}{444.697}{283.318}{algebraic}%
\TextAt{8}{494.071}{283.318}{number}%
\TextAt{8}{79.226}{266.272}{fields,}%
\TextAt{8}{112.216}{266.272}{presupposing}%
\TextAt{8}{182.615}{266.272}{only}%
\TextAt{8}{208.015}{266.272}{basic}%
\TextAt{8}{237.049}{266.272}{algebra}%
\TextAt{8}{277.74}{266.272}{(it}%
\TextAt{8}{293.707}{266.272}{starts}%
\TextAt{8}{326.065}{266.272}{with}%
\TextAt{8}{352.433}{266.272}{the}%
\TextAt{8}{372.298}{266.272}{equation}%
\TextAt{9}{419.815}{266.272}{2}%
\TextAt{9}{428.973}{266.272}{\Character{61}}%
\TextAt{9}{441.394}{266.272}{1}%
\TextAt{9}{449.297}{266.272}{\Character{43}}%
\TextAt{9}{460.463}{266.272}{1}%
\TextAt{8}{466.345}{266.272}{).}%
\TextAt{8}{479.251}{266.272}{But}%
\TextAt{8}{502.2}{266.272}{unlike}%
\TextAt{8}{79.226}{249.226}{the}%
\TextAt{8}{99.624}{249.226}{textbooks}%
\TextAt{8}{153.914}{249.226}{alluded}%
\TextAt{8}{195.126}{249.226}{to}%
\TextAt{8}{209.678}{249.226}{above,}%
\TextAt{8}{246.406}{249.226}{it}%
\TextAt{8}{258.352}{249.226}{progressively}%
\TextAt{8}{327.993}{249.226}{emphasizes}%
\TextAt{8}{388.835}{249.226}{theoretical}%
\TextAt{8}{446.963}{249.226}{aspects}%
\TextAt{8}{487.971}{249.226}{that}%
\TextAt{8}{513.581}{249.226}{rely}%
\TextAt{8}{79.226}{232.18}{on}%
\TextAt{8}{96.312}{232.18}{modern}%
\TextAt{8}{139.735}{232.18}{concepts.}%
\TextAt{8}{194.3}{232.18}{Still,}%
\TextAt{8}{223.296}{232.18}{in}%
\TextAt{8}{237.776}{232.18}{doing}%
\TextAt{8}{270.475}{232.18}{so,}%
\TextAt{8}{289.131}{232.18}{a}%
\TextAt{8}{299.713}{232.18}{special}%
\TextAt{8}{338.641}{232.18}{effort}%
\TextAt{8}{370.359}{232.18}{is}%
\TextAt{8}{382.95}{232.18}{made}%
\TextAt{8}{414.992}{232.18}{to}%
\TextAt{8}{430.129}{232.18}{limit}%
\TextAt{8}{458.919}{232.18}{the}%
\TextAt{8}{479.902}{232.18}{amount}%
\TextAt{8}{523.337}{232.18}{of}%
\TextAt{8}{79.226}{215.134}{abstraction}%
\TextAt{8}{141.046}{215.134}{used,}%
\TextAt{8}{171.779}{215.134}{in}%
\TextAt{8}{186.041}{215.134}{order}%
\TextAt{8}{217.219}{215.134}{that}%
\TextAt{8}{243.197}{215.134}{the}%
\TextAt{8}{263.963}{215.134}{reader}%
\TextAt{8}{300.342}{215.134}{should}%
\TextAt{8}{338.084}{215.134}{not}%
\TextAt{8}{359.507}{215.134}{lose}%
\TextAt{8}{382.938}{215.134}{sight}%
\TextAt{8}{411.905}{215.134}{of}%
\TextAt{8}{425.844}{215.134}{the}%
\TextAt{8}{446.61}{215.134}{concrete}%
\TextAt{8}{493.39}{215.134}{goals}%
\TextAt{8}{523.337}{215.134}{of}%
\TextAt{8}{79.226}{198.088}{number}%
\TextAt{8}{122.018}{198.088}{theory}%
\TextAt{8}{158.965}{198.088}{proper.}%
\TextAt{8}{201.499}{198.088}{The}%
\TextAt{8}{225.749}{198.088}{desire}%
\TextAt{8}{259.169}{198.088}{to}%
\TextAt{8}{273.675}{198.088}{present}%
\TextAt{8}{314.578}{198.088}{number}%
\TextAt{8}{357.371}{198.088}{theory}%
\TextAt{8}{394.317}{198.088}{as}%
\TextAt{8}{408.883}{198.088}{much}%
\TextAt{8}{440.294}{198.088}{as}%
\TextAt{8}{454.86}{198.088}{possible}%
\TextAt{8}{499.075}{198.088}{from}%
\TextAt{8}{526.911}{198.088}{a}%
\TextAt{8}{79.226}{181.042}{unified}%
\TextAt{8}{117.712}{181.042}{theoretical}%
\TextAt{8}{175.721}{181.042}{point}%
\TextAt{8}{206.412}{181.042}{of}%
\TextAt{8}{219.864}{181.042}{view}%
\TextAt{8}{246.968}{181.042}{seems}%
\TextAt{8}{280.373}{181.042}{imperative}%
\TextAt{8}{338.705}{181.042}{today,}%
\TextAt{8}{373.982}{181.042}{as}%
\TextAt{8}{388.473}{181.042}{a}%
\TextAt{8}{398.351}{181.042}{result}%
\TextAt{8}{431.05}{181.042}{of}%
\TextAt{8}{444.502}{181.042}{the}%
\TextAt{8}{464.78}{181.042}{revolutionary}%
\TextAt{8}{79.226}{163.996}{development}%
\TextAt{8}{147.576}{163.996}{that}%
\TextAt{8}{173.329}{163.996}{number}%
\TextAt{8}{216.31}{163.996}{theory}%
\TextAt{8}{253.445}{163.996}{has}%
\TextAt{8}{274.703}{163.996}{undergone}%
\TextAt{8}{331.672}{163.996}{in}%
\TextAt{8}{345.709}{163.996}{the}%
\TextAt{8}{366.25}{163.996}{last}%
\TextAt{8}{388.811}{163.996}{decades}%
\TextAt{8}{432.175}{163.996}{in}%
\TextAt{8}{446.212}{163.996}{conjunction}%
\TextAt{8}{510.007}{163.996}{with}%
\TextAt{8}{79.226}{146.95}{‘arithmetic}%
\TextAt{8}{138.214}{146.95}{algebraic}%
\TextAt{8}{186.801}{146.95}{geometry’}%
\TextAt{8}{235.602}{146.95}{.}%
\TextAt{8}{243.756}{146.95}{The}%
\TextAt{8}{266.962}{146.95}{immense}%
\TextAt{8}{314.294}{146.95}{success}%
\TextAt{8}{353.293}{146.95}{that}%
\TextAt{8}{377.815}{146.95}{this}%
\TextAt{8}{399.79}{146.95}{new}%
\TextAt{8}{422.996}{146.95}{geometric}%
\TextAt{8}{475.48}{146.95}{perspective}%
\TextAt{8}{79.226}{129.904}{has}%
\TextAt{8}{101.586}{129.904}{brought}%
\TextAt{8}{146.984}{129.904}{about}%
\TextAt{8}{181.969}{129.904}{-}%
\TextAt{8}{191.251}{129.904}{for}%
\TextAt{8}{210.622}{129.904}{instance,}%
\TextAt{8}{261.317}{129.904}{in}%
\TextAt{8}{276.457}{129.904}{the}%
\TextAt{8}{298.1}{129.904}{context}%
\TextAt{8}{341.214}{129.904}{of}%
\TextAt{8}{356.031}{129.904}{the}%
\TextAt{8}{377.674}{129.904}{Weil}%
\TextAt{8}{405.808}{129.904}{conjectures,}%
\TextAt{8}{471.782}{129.904}{the}%
\TextAt{8}{493.425}{129.904}{Mordell}%
\TextAt{8}{79.226}{112.858}{conjecture,}%
\TextAt{8}{141.06}{112.858}{of}%
\TextAt{8}{156.256}{112.858}{problems}%
\TextAt{8}{208.261}{112.858}{related}%
\TextAt{8}{249.149}{112.858}{to}%
\TextAt{8}{265.325}{112.858}{the}%
\TextAt{8}{287.348}{112.858}{conjectures}%
\TextAt{8}{350.078}{112.858}{of}%
\TextAt{8}{365.274}{112.858}{Birch}%
\TextAt{8}{398.51}{112.858}{and}%
\TextAt{8}{423.139}{112.858}{Swinnerton-Dyer}%
\TextAt{8}{515.242}{112.858}{-}%
\TextAt{8}{524.903}{112.858}{is}%
\TextAt{8}{79.226}{95.812}{largely}%
\TextAt{8}{117.139}{95.812}{based}%
\TextAt{8}{149.577}{95.812}{on}%
\TextAt{8}{165.696}{95.812}{the}%
\TextAt{8}{185.712}{95.812}{unconditional}%
\TextAt{8}{259.072}{95.812}{and}%
\TextAt{8}{281.695}{95.812}{universal}%
\TextAt{8}{331.049}{95.812}{application}%
\TextAt{8}{391.402}{95.812}{of}%
\TextAt{8}{404.592}{95.812}{the}%
\TextAt{8}{424.608}{95.812}{conceptual}%
\TextAt{8}{483}{95.812}{approach.}%
\TextAt{8}{97.159}{75.777}{It}%
\TextAt{8}{110.338}{75.777}{is}%
\TextAt{8}{122.609}{75.777}{true}%
\TextAt{8}{147.827}{75.777}{that}%
\TextAt{8}{173.703}{75.777}{those}%
\TextAt{8}{204.839}{75.777}{impressive}%
\TextAt{8}{262.049}{75.777}{results}%
\TextAt{8}{299.749}{75.777}{can}%
\TextAt{8}{321.715}{75.777}{hardly}%
\TextAt{8}{358.972}{75.777}{be}%
\TextAt{8}{375.404}{75.777}{touched}%
\TextAt{8}{419.81}{75.777}{upon}%
\TextAt{8}{449.906}{75.777}{in}%
\TextAt{8}{464.065}{75.777}{this}%
\TextAt{8}{487.395}{75.777}{book}%
\TextAt{8}{516.845}{75.777}{be-}%
\TextAt{8}{301.281}{43.033}{vi}%
\Destination{79.23}{773.075}{nav.chapterz2az.5}%
\Destination{78.23}{809.815}{nav.page.vi}%
\Destination{79.23}{773.075}{outline.4}%
\Hotlink{78.726}{335.255}{19.22}{11.047}{nav.cite.N014}%
}

% Page 7 (body page front matter)
\Page{%
\TextAt{2}{79.226}{796.943}{Preface}%
\TextAt{2}{117.888}{796.943}{to}%
\TextAt{2}{131.219}{796.943}{the}%
\TextAt{2}{150.004}{796.943}{German}%
\TextAt{2}{191.928}{796.943}{Edition}%
\Vector{q 0 G 0 g 0.299 w 7.226 18.729 m 460.769 18.729 l S Q}%
\TextAt{8}{79.226}{761.124}{cause}%
\TextAt{8}{111.67}{761.124}{they}%
\TextAt{8}{139.175}{761.124}{require}%
\TextAt{8}{179.689}{761.124}{higher}%
\TextAt{8}{216.627}{761.124}{dimensional}%
\TextAt{8}{282.246}{761.124}{theories,}%
\TextAt{8}{330.595}{761.124}{whereas}%
\TextAt{8}{376.046}{761.124}{the}%
\TextAt{8}{397.371}{761.124}{book}%
\TextAt{8}{427.483}{761.124}{deliberately}%
\TextAt{8}{492.385}{761.124}{confines}%
\TextAt{8}{79.226}{744.078}{itself}%
\TextAt{8}{109.259}{744.078}{to}%
\TextAt{8}{125.256}{744.078}{the}%
\TextAt{8}{147.1}{744.078}{theory}%
\TextAt{8}{185.537}{744.078}{of}%
\TextAt{8}{200.554}{744.078}{algebraic}%
\TextAt{8}{251.675}{744.078}{number}%
\TextAt{8}{295.959}{744.078}{fields,}%
\TextAt{8}{331.291}{744.078}{i.e.,}%
\TextAt{8}{355.505}{744.078}{to}%
\TextAt{8}{371.502}{744.078}{the}%
\TextAt{8}{393.346}{744.078}{1-dimensional}%
\TextAt{8}{469.238}{744.078}{case.}%
\TextAt{8}{503.621}{744.078}{But}%
\TextAt{8}{528.549}{744.078}{I}%
\TextAt{8}{79.226}{727.032}{thought}%
\TextAt{8}{123.22}{727.032}{it}%
\TextAt{8}{135.007}{727.032}{necessary}%
\TextAt{8}{186.916}{727.032}{to}%
\TextAt{8}{201.309}{727.032}{present}%
\TextAt{8}{242.099}{727.032}{the}%
\TextAt{8}{262.338}{727.032}{theory}%
\TextAt{8}{299.171}{727.032}{in}%
\TextAt{8}{312.907}{727.032}{a}%
\TextAt{8}{322.745}{727.032}{way}%
\TextAt{8}{346.571}{727.032}{which}%
\TextAt{8}{380.14}{727.032}{takes}%
\TextAt{8}{410.207}{727.032}{these}%
\TextAt{8}{440.261}{727.032}{developments}%
\TextAt{8}{512.924}{727.032}{into}%
\TextAt{8}{79.226}{709.986}{account,}%
\TextAt{8}{126.009}{709.986}{taking}%
\TextAt{8}{162.336}{709.986}{them}%
\TextAt{8}{192.469}{709.986}{as}%
\TextAt{8}{207.061}{709.986}{the}%
\TextAt{8}{227.44}{709.986}{distant}%
\TextAt{8}{267.078}{709.986}{focus,}%
\TextAt{8}{300.579}{709.986}{borrowing}%
\TextAt{8}{356.093}{709.986}{emphases}%
\TextAt{8}{408.463}{709.986}{and}%
\TextAt{8}{431.448}{709.986}{arguments}%
\TextAt{8}{488.648}{709.986}{from}%
\TextAt{8}{516.51}{709.986}{the}%
\TextAt{8}{79.226}{692.94}{higher}%
\TextAt{8}{114.297}{692.94}{point}%
\TextAt{8}{144.167}{692.94}{of}%
\TextAt{8}{156.798}{692.94}{view,}%
\TextAt{8}{186.473}{692.94}{thus}%
\TextAt{8}{211.526}{692.94}{integrating}%
\TextAt{8}{270.351}{692.94}{the}%
\TextAt{8}{289.809}{692.94}{theory}%
\TextAt{8}{325.86}{692.94}{of}%
\TextAt{8}{338.491}{692.94}{algebraic}%
\TextAt{8}{387.226}{692.94}{number}%
\TextAt{8}{429.123}{692.94}{fields}%
\TextAt{8}{458.396}{692.94}{into}%
\TextAt{8}{481.44}{692.94}{the}%
\TextAt{8}{500.897}{692.94}{higher}%
\TextAt{8}{79.226}{675.894}{dimensional}%
\TextAt{8}{144.623}{675.894}{theory}%
\TextAt{8}{182.319}{675.894}{-}%
\TextAt{8}{191.06}{675.894}{or}%
\TextAt{8}{206.317}{675.894}{at}%
\TextAt{8}{221.574}{675.894}{least}%
\TextAt{8}{249.897}{675.894}{avoiding}%
\TextAt{8}{297.361}{675.894}{any}%
\TextAt{8}{320.425}{675.894}{obstruction}%
\TextAt{8}{383.227}{675.894}{to}%
\TextAt{8}{398.484}{675.894}{such}%
\TextAt{8}{425.827}{675.894}{an}%
\TextAt{8}{443.032}{675.894}{integration.}%
\TextAt{8}{509.947}{675.894}{This}%
\TextAt{8}{79.226}{658.848}{is}%
\TextAt{8}{91.857}{658.848}{why}%
\TextAt{8}{117.435}{658.848}{I}%
\TextAt{8}{126.419}{658.848}{preferred,}%
\TextAt{8}{180.5}{658.848}{whenever}%
\TextAt{8}{232.738}{658.848}{it}%
\TextAt{8}{245.31}{658.848}{was}%
\TextAt{8}{268.676}{658.848}{feasible,}%
\TextAt{8}{314.365}{658.848}{the}%
\TextAt{8}{335.388}{658.848}{functorial}%
\TextAt{8}{389.264}{658.848}{point}%
\TextAt{8}{420.701}{658.848}{of}%
\TextAt{8}{434.898}{658.848}{view}%
\TextAt{8}{462.748}{658.848}{and}%
\TextAt{8}{486.377}{658.848}{the}%
\TextAt{8}{507.401}{658.848}{more}%
\TextAt{8}{79.226}{641.802}{far-reaching}%
\TextAt{8}{143.329}{641.802}{argument}%
\TextAt{8}{195.405}{641.802}{to}%
\TextAt{8}{209.428}{641.802}{the}%
\TextAt{8}{229.296}{641.802}{clever}%
\TextAt{8}{262.172}{641.802}{trick,}%
\TextAt{8}{292.512}{641.802}{and}%
\TextAt{8}{314.987}{641.802}{made}%
\TextAt{8}{345.914}{641.802}{a}%
\TextAt{8}{355.382}{641.802}{particular}%
\TextAt{8}{409.084}{641.802}{effort}%
\TextAt{8}{439.688}{641.802}{to}%
\TextAt{8}{453.711}{641.802}{place}%
\TextAt{8}{483.335}{641.802}{geometric}%
\TextAt{8}{79.226}{624.756}{interpretation}%
\TextAt{8}{153.707}{624.756}{to}%
\TextAt{8}{168.017}{624.756}{the}%
\TextAt{8}{188.174}{624.756}{fore,}%
\TextAt{8}{214.511}{624.756}{in}%
\TextAt{8}{228.164}{624.756}{the}%
\TextAt{8}{248.32}{624.756}{spirit}%
\TextAt{8}{278.949}{624.756}{of}%
\TextAt{8}{292.279}{624.756}{the}%
\TextAt{8}{312.436}{624.756}{theory}%
\TextAt{8}{349.186}{624.756}{of}%
\TextAt{8}{362.516}{624.756}{algebraic}%
\TextAt{8}{411.951}{624.756}{curves.}%
\TextAt{8}{97.159}{604.721}{Let}%
\TextAt{8}{118.332}{604.721}{me}%
\TextAt{8}{137.388}{604.721}{forego}%
\TextAt{8}{172.393}{604.721}{the}%
\TextAt{8}{192.752}{604.721}{usual}%
\TextAt{8}{223.585}{604.721}{habit}%
\TextAt{8}{254.357}{604.721}{of}%
\TextAt{8}{267.891}{604.721}{describing}%
\TextAt{8}{323.434}{604.721}{the}%
\TextAt{8}{343.794}{604.721}{content}%
\TextAt{8}{385.625}{604.721}{of}%
\TextAt{8}{399.158}{604.721}{each}%
\TextAt{8}{425.699}{604.721}{individual}%
\TextAt{8}{480.86}{604.721}{chapter}%
\TextAt{8}{523.014}{604.721}{in}%
\TextAt{8}{79.226}{587.675}{this}%
\TextAt{8}{101.328}{587.675}{foreword;}%
\TextAt{8}{152.232}{587.675}{simply}%
\TextAt{8}{188.967}{587.675}{turning}%
\TextAt{8}{229.875}{587.675}{pages}%
\TextAt{8}{261.087}{587.675}{will}%
\TextAt{8}{282.472}{587.675}{yield}%
\TextAt{8}{310.038}{587.675}{the}%
\TextAt{8}{329.474}{587.675}{same}%
\TextAt{8}{358.08}{587.675}{information}%
\TextAt{8}{420.782}{587.675}{in}%
\TextAt{8}{433.715}{587.675}{a}%
\TextAt{8}{442.75}{587.675}{more}%
\TextAt{8}{471.296}{587.675}{entertaining}%
\TextAt{8}{79.226}{570.629}{manner.}%
\TextAt{8}{127.896}{570.629}{I}%
\TextAt{8}{136.628}{570.629}{would}%
\TextAt{8}{171.388}{570.629}{however}%
\TextAt{8}{216.883}{570.629}{like}%
\TextAt{8}{238.957}{570.629}{to}%
\TextAt{8}{253.883}{570.629}{emphasize}%
\TextAt{8}{310.484}{570.629}{a}%
\TextAt{8}{320.855}{570.629}{few}%
\TextAt{8}{342.595}{570.629}{basic}%
\TextAt{8}{372.536}{570.629}{principles}%
\TextAt{8}{425.886}{570.629}{that}%
\TextAt{8}{451.87}{570.629}{have}%
\TextAt{8}{479.48}{570.629}{guided}%
\TextAt{8}{517.814}{570.629}{me}%
\TextAt{8}{79.226}{553.583}{while}%
\TextAt{8}{109.787}{553.583}{writing}%
\TextAt{8}{150.115}{553.583}{the}%
\TextAt{8}{170.274}{553.583}{book.}%
\TextAt{8}{203.779}{553.583}{The}%
\TextAt{8}{227.836}{553.583}{first}%
\TextAt{8}{251.964}{553.583}{chapter}%
\TextAt{8}{293.918}{553.583}{lays}%
\TextAt{8}{317.402}{553.583}{down}%
\TextAt{8}{348.297}{553.583}{the}%
\TextAt{8}{368.457}{553.583}{foundations}%
\TextAt{8}{431.942}{553.583}{of}%
\TextAt{8}{445.275}{553.583}{the}%
\TextAt{8}{465.435}{553.583}{global}%
\TextAt{8}{499.917}{553.583}{theory}%
\TextAt{8}{79.226}{536.537}{and}%
\TextAt{8}{100.989}{536.537}{the}%
\TextAt{8}{120.146}{536.537}{second}%
\TextAt{8}{156.925}{536.537}{of}%
\TextAt{8}{169.256}{536.537}{the}%
\TextAt{8}{188.413}{536.537}{local}%
\TextAt{8}{215.054}{536.537}{theory}%
\TextAt{8}{250.804}{536.537}{of}%
\TextAt{8}{263.135}{536.537}{algebraic}%
\TextAt{8}{311.57}{536.537}{number}%
\TextAt{8}{353.167}{536.537}{fields.}%
\TextAt{8}{387.356}{536.537}{These}%
\TextAt{8}{420.226}{536.537}{foundations}%
\TextAt{8}{482.708}{536.537}{are}%
\TextAt{8}{501.22}{536.537}{finally}%
\TextAt{8}{79.226}{519.491}{summed}%
\TextAt{8}{126.066}{519.491}{up}%
\TextAt{8}{143.579}{519.491}{in}%
\TextAt{8}{157.841}{519.491}{the}%
\TextAt{8}{178.607}{519.491}{first}%
\TextAt{8}{203.341}{519.491}{three}%
\TextAt{8}{233.862}{519.491}{sections}%
\TextAt{8}{278.167}{519.491}{of}%
\TextAt{8}{292.106}{519.491}{chapter}%
\TextAt{8}{334.666}{519.491}{III,}%
\TextAt{8}{355.73}{519.491}{the}%
\TextAt{8}{376.496}{519.491}{aim}%
\TextAt{8}{399.867}{519.491}{of}%
\TextAt{8}{413.806}{519.491}{which}%
\TextAt{8}{447.902}{519.491}{is}%
\TextAt{8}{460.275}{519.491}{to}%
\TextAt{8}{475.194}{519.491}{present}%
\TextAt{8}{516.51}{519.491}{the}%
\TextAt{8}{79.226}{502.445}{perfect}%
\TextAt{8}{117.576}{502.445}{analogy}%
\TextAt{8}{160.181}{502.445}{of}%
\TextAt{8}{172.851}{502.445}{the}%
\TextAt{8}{192.348}{502.445}{classical}%
\TextAt{8}{236.687}{502.445}{notions}%
\TextAt{8}{277.069}{502.445}{and}%
\TextAt{8}{299.171}{502.445}{results}%
\TextAt{8}{335.704}{502.445}{with}%
\TextAt{8}{361.704}{502.445}{the}%
\TextAt{8}{381.2}{502.445}{theory}%
\TextAt{8}{417.29}{502.445}{of}%
\TextAt{8}{429.96}{502.445}{algebraic}%
\TextAt{8}{478.735}{502.445}{curves}%
\TextAt{8}{513.904}{502.445}{and}%
\TextAt{8}{79.226}{485.399}{the}%
\TextAt{8}{100.163}{485.399}{idea}%
\TextAt{8}{125.654}{485.399}{of}%
\TextAt{8}{139.764}{485.399}{the}%
\TextAt{8}{160.701}{485.399}{Riemann-Roch}%
\TextAt{8}{241.15}{485.399}{theorem.}%
\TextAt{8}{293.567}{485.399}{The}%
\TextAt{8}{318.401}{485.399}{presentation}%
\TextAt{8}{385.915}{485.399}{is}%
\TextAt{8}{398.459}{485.399}{dominated}%
\TextAt{8}{457.126}{485.399}{by}%
\TextAt{8}{474.165}{485.399}{“Arakelov’s}%
\TextAt{8}{79.226}{468.353}{point}%
\TextAt{8}{108.578}{468.353}{of}%
\TextAt{8}{120.69}{468.353}{view”,}%
\TextAt{8}{155.474}{468.353}{which}%
\TextAt{8}{187.742}{468.353}{has}%
\TextAt{8}{207.398}{468.353}{acquired}%
\TextAt{8}{253.332}{468.353}{much}%
\TextAt{8}{283.329}{468.353}{importance}%
\TextAt{8}{343.573}{468.353}{in}%
\TextAt{8}{356.008}{468.353}{recent}%
\TextAt{8}{389.579}{468.353}{years.}%
\TextAt{8}{423.708}{468.353}{It}%
\TextAt{8}{435.163}{468.353}{is}%
\TextAt{8}{445.709}{468.353}{probably}%
\TextAt{8}{493.603}{468.353}{the}%
\TextAt{8}{512.541}{468.353}{first}%
\TextAt{8}{79.226}{451.307}{time}%
\TextAt{8}{106.581}{451.307}{that}%
\TextAt{8}{132.646}{451.307}{this}%
\TextAt{8}{156.163}{451.307}{approach,}%
\TextAt{8}{210.699}{451.307}{with}%
\TextAt{8}{238.055}{451.307}{all}%
\TextAt{8}{255.009}{451.307}{its}%
\TextAt{8}{272.023}{451.307}{intricate}%
\TextAt{8}{319.224}{451.307}{normalizations,}%
\TextAt{8}{402.118}{451.307}{has}%
\TextAt{8}{423.687}{451.307}{received}%
\TextAt{8}{469.25}{451.307}{an}%
\TextAt{8}{486.204}{451.307}{extensive}%
\TextAt{8}{79.226}{434.261}{treatment}%
\TextAt{8}{133.151}{434.261}{in}%
\TextAt{8}{146.415}{434.261}{a}%
\TextAt{8}{155.783}{434.261}{textbook.}%
\TextAt{8}{209.639}{434.261}{But}%
\TextAt{8}{232.492}{434.261}{I}%
\TextAt{8}{240.221}{434.261}{finally}%
\TextAt{8}{275.281}{434.261}{decided}%
\TextAt{8}{317.154}{434.261}{not}%
\TextAt{8}{337.58}{434.261}{to}%
\TextAt{8}{351.503}{434.261}{employ}%
\TextAt{8}{391.44}{434.261}{the}%
\TextAt{8}{411.208}{434.261}{term}%
\TextAt{8}{438.783}{434.261}{“Arakelov}%
\TextAt{8}{493.03}{434.261}{divisor”}%
\TextAt{8}{79.226}{417.215}{although}%
\TextAt{8}{127.677}{417.215}{it}%
\TextAt{8}{139.044}{417.215}{is}%
\TextAt{8}{150.47}{417.215}{now}%
\TextAt{8}{174.521}{417.215}{widely}%
\TextAt{8}{210.921}{417.215}{used.}%
\TextAt{8}{242.079}{417.215}{This}%
\TextAt{8}{268.462}{417.215}{would}%
\TextAt{8}{302.268}{417.215}{have}%
\TextAt{8}{328.925}{417.215}{entailed}%
\TextAt{8}{372.486}{417.215}{attaching}%
\TextAt{8}{423.867}{417.215}{the}%
\TextAt{8}{443.686}{417.215}{name}%
\TextAt{8}{474.563}{417.215}{of}%
\TextAt{8}{487.555}{417.215}{Arakelov}%
\TextAt{8}{79.226}{400.169}{to}%
\TextAt{8}{92.605}{400.169}{many}%
\TextAt{8}{123.547}{400.169}{other}%
\TextAt{8}{153.185}{400.169}{concepts,}%
\TextAt{8}{203.226}{400.169}{introducing}%
\TextAt{8}{264.737}{400.169}{too}%
\TextAt{8}{284.297}{400.169}{heavy}%
\TextAt{8}{316.864}{400.169}{a}%
\TextAt{8}{325.688}{400.169}{terminology}%
\TextAt{8}{389.482}{400.169}{for}%
\TextAt{8}{406.436}{400.169}{this}%
\TextAt{8}{428.328}{400.169}{elementary}%
\TextAt{8}{487.232}{400.169}{material.}%
\TextAt{8}{79.226}{383.124}{My}%
\TextAt{8}{99.486}{383.124}{decision}%
\TextAt{8}{143.225}{383.124}{seemed}%
\TextAt{8}{183.055}{383.124}{all}%
\TextAt{8}{198.771}{383.124}{the}%
\TextAt{8}{218.385}{383.124}{more}%
\TextAt{8}{247.109}{383.124}{justified}%
\TextAt{8}{291.171}{383.124}{as}%
\TextAt{8}{304.999}{383.124}{Arakelov}%
\TextAt{8}{353.568}{383.124}{himself}%
\TextAt{8}{393.075}{383.124}{introduced}%
\TextAt{8}{451.065}{383.124}{his}%
\TextAt{8}{468.79}{383.124}{divisors}%
\TextAt{8}{510.975}{383.124}{only}%
\TextAt{8}{79.226}{366.078}{for}%
\TextAt{8}{98.214}{366.078}{arithmetic}%
\TextAt{8}{155.9}{366.078}{surfaces.}%
\TextAt{8}{207.777}{366.078}{The}%
\TextAt{8}{232.934}{366.078}{corresponding}%
\TextAt{8}{309.222}{366.078}{idea}%
\TextAt{8}{335.036}{366.078}{in}%
\TextAt{8}{349.792}{366.078}{the}%
\TextAt{8}{371.051}{366.078}{number}%
\TextAt{8}{414.749}{366.078}{field}%
\TextAt{8}{441.209}{366.078}{case}%
\TextAt{8}{467.082}{366.078}{goes}%
\TextAt{8}{493.936}{366.078}{back}%
\TextAt{8}{522.357}{366.078}{to}%
\TextAt{8}{79.226}{349.032}{Hasse,}%
\TextAt{8}{115.438}{349.032}{and}%
\TextAt{8}{138.201}{349.032}{is}%
\TextAt{8}{149.965}{349.032}{clearly}%
\TextAt{8}{187.361}{349.032}{highlighted}%
\TextAt{8}{248.177}{349.032}{for}%
\TextAt{8}{266.062}{349.032}{instance}%
\TextAt{8}{311.647}{349.032}{in}%
\TextAt{8}{325.299}{349.032}{S.}%
\TextAt{8}{338.952}{349.032}{Lang’s}%
\TextAt{8}{376.252}{349.032}{textbook}%
\TextAt{8}{425.687}{349.032}{[94].}%
\TextAt{8}{97.159}{328.997}{It}%
\TextAt{8}{109.419}{328.997}{was}%
\TextAt{8}{131.506}{328.997}{not}%
\TextAt{8}{151.908}{328.997}{without}%
\TextAt{8}{195.072}{328.997}{hesitation}%
\TextAt{8}{248.709}{328.997}{that}%
\TextAt{8}{273.665}{328.997}{I}%
\TextAt{8}{281.37}{328.997}{decided}%
\TextAt{8}{323.219}{328.997}{to}%
\TextAt{8}{337.117}{328.997}{include}%
\TextAt{8}{377.017}{328.997}{Class}%
\TextAt{8}{407.294}{328.997}{Field}%
\TextAt{8}{436.626}{328.997}{Theory}%
\TextAt{8}{476.861}{328.997}{in}%
\TextAt{8}{490.101}{328.997}{chapters}%
\TextAt{8}{79.226}{311.951}{IV-VI.}%
\TextAt{8}{116.499}{311.951}{Since}%
\TextAt{8}{147.292}{311.951}{my}%
\TextAt{8}{167.039}{311.951}{book}%
\TextAt{8}{196.218}{311.951}{[107]}%
\TextAt{8}{224.429}{311.951}{on}%
\TextAt{8}{240.924}{311.951}{this}%
\TextAt{8}{263.982}{311.951}{subject}%
\TextAt{8}{304.59}{311.951}{had}%
\TextAt{8}{327.589}{311.951}{been}%
\TextAt{8}{355.453}{311.951}{published}%
\TextAt{8}{408.423}{311.951}{not}%
\TextAt{8}{429.472}{311.951}{long}%
\TextAt{8}{455.077}{311.951}{before,}%
\TextAt{8}{493.736}{311.951}{another}%
\TextAt{8}{79.226}{294.905}{treatment}%
\TextAt{8}{132.629}{294.905}{of}%
\TextAt{8}{145.05}{294.905}{this}%
\TextAt{8}{166.964}{294.905}{theory}%
\TextAt{8}{202.805}{294.905}{posed}%
\TextAt{8}{234.796}{294.905}{obvious}%
\TextAt{8}{276.232}{294.905}{questions.}%
\TextAt{8}{331.66}{294.905}{But}%
\TextAt{8}{353.992}{294.905}{in}%
\TextAt{8}{366.736}{294.905}{the}%
\TextAt{8}{385.983}{294.905}{end,}%
\TextAt{8}{410.613}{294.905}{after}%
\TextAt{8}{437.344}{294.905}{long}%
\TextAt{8}{461.804}{294.905}{consideration,}%
\TextAt{8}{79.226}{277.859}{there}%
\TextAt{8}{110.997}{277.859}{was}%
\TextAt{8}{135.355}{277.859}{simply}%
\TextAt{8}{174.669}{277.859}{no}%
\TextAt{8}{192.787}{277.859}{other}%
\TextAt{8}{225.215}{277.859}{choice.}%
\TextAt{8}{270.131}{277.859}{A}%
\TextAt{8}{284.66}{277.859}{sourcebook}%
\TextAt{8}{347.394}{277.859}{on}%
\TextAt{8}{365.512}{277.859}{algebraic}%
\TextAt{8}{416.805}{277.859}{number}%
\TextAt{8}{461.26}{277.859}{fields}%
\TextAt{8}{493.09}{277.859}{without}%
\TextAt{8}{79.226}{260.813}{the}%
\TextAt{8}{101.134}{260.813}{crowning}%
\TextAt{8}{152.643}{260.813}{conclusion}%
\TextAt{8}{211.038}{260.813}{of}%
\TextAt{8}{226.119}{260.813}{class}%
\TextAt{8}{255.308}{260.813}{field}%
\TextAt{8}{282.416}{260.813}{theory}%
\TextAt{8}{320.918}{260.813}{with}%
\TextAt{8}{349.329}{260.813}{its}%
\TextAt{8}{367.4}{260.813}{important}%
\TextAt{8}{424.444}{260.813}{consequences}%
\TextAt{8}{496.874}{260.813}{for}%
\TextAt{8}{516.51}{260.813}{the}%
\TextAt{8}{79.226}{243.767}{theory}%
\TextAt{8}{116.282}{243.767}{of}%
\TextAt{8}{129.918}{243.767}{L-series}%
\TextAt{8}{172.772}{243.767}{would}%
\TextAt{8}{207.222}{243.767}{have}%
\TextAt{8}{234.523}{243.767}{appeared}%
\TextAt{8}{285.232}{243.767}{like}%
\TextAt{8}{306.997}{243.767}{a}%
\TextAt{8}{317.059}{243.767}{torso,}%
\TextAt{8}{350.031}{243.767}{suffering}%
\TextAt{8}{397.547}{243.767}{from}%
\TextAt{8}{425.494}{243.767}{an}%
\TextAt{8}{442.059}{243.767}{unacceptable}%
\TextAt{8}{512.601}{243.767}{lack}%
\TextAt{8}{79.226}{226.721}{of}%
\TextAt{8}{93.9}{226.721}{completeness.}%
\TextAt{8}{172.839}{226.721}{This}%
\TextAt{8}{200.903}{226.721}{also}%
\TextAt{8}{225.727}{226.721}{gave}%
\TextAt{8}{253.42}{226.721}{me}%
\TextAt{8}{273.618}{226.721}{the}%
\TextAt{8}{295.118}{226.721}{opportunity}%
\TextAt{8}{361.188}{226.721}{to}%
\TextAt{8}{376.842}{226.721}{modify}%
\TextAt{8}{417.53}{226.721}{and}%
\TextAt{8}{441.637}{226.721}{emend}%
\TextAt{8}{480.042}{226.721}{my}%
\TextAt{8}{500.897}{226.721}{earlier}%
\TextAt{8}{79.226}{209.675}{treatment,}%
\TextAt{8}{136.856}{209.675}{to}%
\TextAt{8}{151.218}{209.675}{enrich}%
\TextAt{8}{186.06}{209.675}{that}%
\TextAt{8}{211.481}{209.675}{somewhat}%
\TextAt{8}{266.228}{209.675}{dry}%
\TextAt{8}{287.417}{209.675}{presentation}%
\TextAt{8}{354.203}{209.675}{with}%
\TextAt{8}{380.916}{209.675}{quite}%
\TextAt{8}{410.557}{209.675}{a}%
\TextAt{8}{420.365}{209.675}{few}%
\TextAt{8}{441.542}{209.675}{examples,}%
\TextAt{8}{495.321}{209.675}{to}%
\TextAt{8}{509.684}{209.675}{refer}%
\TextAt{8}{79.226}{192.629}{ahead}%
\TextAt{8}{113.047}{192.629}{with}%
\TextAt{8}{139.707}{192.629}{some}%
\TextAt{8}{169.033}{192.629}{remarks,}%
\TextAt{8}{216.902}{192.629}{and}%
\TextAt{8}{239.665}{192.629}{to}%
\TextAt{8}{253.975}{192.629}{add}%
\TextAt{8}{276.738}{192.629}{beneficial}%
\TextAt{8}{328.432}{192.629}{exercises.}%
\TextAt{8}{97.159}{172.594}{A}%
\TextAt{8}{108.833}{172.594}{lot}%
\TextAt{8}{125.398}{172.594}{of}%
\TextAt{8}{137.731}{172.594}{work}%
\TextAt{8}{165.354}{172.594}{went}%
\TextAt{8}{192.32}{172.594}{into}%
\TextAt{8}{215.066}{172.594}{the}%
\TextAt{8}{234.225}{172.594}{last}%
\TextAt{8}{255.405}{172.594}{chapter}%
\TextAt{8}{296.358}{172.594}{on}%
\TextAt{8}{311.62}{172.594}{zeta}%
\TextAt{8}{335.334}{172.594}{functions}%
\TextAt{8}{384.8}{172.594}{and}%
\TextAt{8}{406.565}{172.594}{L-series.}%
\TextAt{8}{453.332}{172.594}{These}%
\TextAt{8}{486.204}{172.594}{functions}%
\TextAt{8}{79.226}{155.548}{have}%
\TextAt{8}{107.419}{155.548}{gained}%
\TextAt{8}{145.691}{155.548}{central}%
\TextAt{8}{185.588}{155.548}{importance}%
\TextAt{8}{248.248}{155.548}{in}%
\TextAt{8}{263.099}{155.548}{recent}%
\TextAt{8}{299.087}{155.548}{decades,}%
\TextAt{8}{346.816}{155.548}{but}%
\TextAt{8}{369.474}{155.548}{textbooks}%
\TextAt{8}{424.722}{155.548}{do}%
\TextAt{8}{442.179}{155.548}{not}%
\TextAt{8}{464.192}{155.548}{pay}%
\TextAt{8}{487.507}{155.548}{sufficient}%
\TextAt{8}{79.226}{138.502}{attention}%
\TextAt{8}{129.094}{138.502}{to}%
\TextAt{8}{142.857}{138.502}{them.}%
\TextAt{8}{177.138}{138.502}{I}%
\TextAt{8}{184.708}{138.502}{did}%
\TextAt{8}{204.317}{138.502}{not,}%
\TextAt{8}{227.945}{138.502}{however,}%
\TextAt{8}{275.639}{138.502}{include}%
\TextAt{8}{315.405}{138.502}{Tate’s}%
\TextAt{8}{349.707}{138.502}{approach}%
\TextAt{8}{399.575}{138.502}{to}%
\TextAt{8}{413.339}{138.502}{Hecke}%
\TextAt{8}{446.601}{138.502}{L-series,}%
\TextAt{8}{491.964}{138.502}{which}%
\TextAt{8}{524.903}{138.502}{is}%
\TextAt{8}{79.226}{121.456}{based}%
\TextAt{8}{110.643}{121.456}{on}%
\TextAt{8}{125.741}{121.456}{harmonic}%
\TextAt{8}{175.963}{121.456}{analysis,}%
\TextAt{8}{222.317}{121.456}{although}%
\TextAt{8}{269.945}{121.456}{it}%
\TextAt{8}{280.488}{121.456}{would}%
\TextAt{8}{313.471}{121.456}{have}%
\TextAt{8}{339.305}{121.456}{suited}%
\TextAt{8}{372.67}{121.456}{the}%
\TextAt{8}{391.665}{121.456}{more}%
\TextAt{8}{419.771}{121.456}{conceptual}%
\TextAt{8}{477.142}{121.456}{orientation}%
\TextAt{8}{79.226}{104.41}{of}%
\TextAt{8}{92.383}{104.41}{the}%
\TextAt{8}{112.366}{104.41}{book}%
\TextAt{8}{141.137}{104.41}{perfectly}%
\TextAt{8}{189.406}{104.41}{well.}%
\TextAt{8}{217.63}{104.41}{In}%
\TextAt{8}{232.078}{104.41}{fact,}%
\TextAt{8}{258.277}{104.41}{the}%
\TextAt{8}{278.26}{104.41}{clarity}%
\TextAt{8}{314.515}{104.41}{of}%
\TextAt{8}{327.671}{104.41}{Tate’s}%
\TextAt{8}{362.348}{104.41}{own}%
\TextAt{8}{386.563}{104.41}{presentation}%
\TextAt{8}{453.124}{104.41}{could}%
\TextAt{8}{484.165}{104.41}{hardly}%
\TextAt{8}{520.743}{104.41}{be}%
\TextAt{8}{79.226}{87.364}{improved}%
\TextAt{8}{130.344}{87.364}{upon,}%
\TextAt{8}{163.257}{87.364}{and}%
\TextAt{8}{186.077}{87.364}{it}%
\TextAt{8}{197.839}{87.364}{has}%
\TextAt{8}{218.771}{87.364}{also}%
\TextAt{8}{242.308}{87.364}{been}%
\TextAt{8}{269.994}{87.364}{sufficiently}%
\TextAt{8}{328.644}{87.364}{repeated}%
\TextAt{8}{376.499}{87.364}{in}%
\TextAt{8}{390.209}{87.364}{other}%
\TextAt{8}{420.836}{87.364}{places.}%
\TextAt{8}{460.086}{87.364}{Instead}%
\TextAt{8}{501.497}{87.364}{I}%
\TextAt{8}{509.672}{87.364}{have}%
\TextAt{8}{299.656}{43.033}{vii}%
\Destination{78.23}{809.815}{nav.page.vii}%
\Hotlink{425.187}{346.855}{19.22}{11.047}{nav.cite.N094}%
\Hotlink{195.718}{309.755}{25.078}{11.047}{nav.cite.N107}%
}

% Page 8 (body page front matter)
\Page{%
\TextAt{2}{79.226}{796.943}{Preface}%
\TextAt{2}{117.888}{796.943}{to}%
\TextAt{2}{131.219}{796.943}{the}%
\TextAt{2}{150.004}{796.943}{German}%
\TextAt{2}{191.928}{796.943}{Edition}%
\Vector{q 0 G 0 g 0.299 w 7.226 18.729 m 460.769 18.729 l S Q}%
\TextAt{8}{79.226}{761.124}{preferred}%
\TextAt{8}{128.092}{761.124}{to}%
\TextAt{8}{141.523}{761.124}{turn}%
\TextAt{8}{166.657}{761.124}{back}%
\TextAt{8}{193.095}{761.124}{to}%
\TextAt{8}{206.526}{761.124}{Hecke’s}%
\TextAt{8}{247.322}{761.124}{approach,}%
\TextAt{8}{300.285}{761.124}{which}%
\TextAt{8}{332.892}{761.124}{is}%
\TextAt{8}{343.776}{761.124}{not}%
\TextAt{8}{363.71}{761.124}{easy}%
\TextAt{8}{388.582}{761.124}{to}%
\TextAt{8}{402.013}{761.124}{understand}%
\TextAt{8}{462.332}{761.124}{in}%
\TextAt{8}{475.105}{761.124}{the}%
\TextAt{8}{494.381}{761.124}{original}%
\TextAt{8}{79.226}{744.078}{version,}%
\TextAt{8}{122.19}{744.078}{but}%
\TextAt{8}{143.616}{744.078}{for}%
\TextAt{8}{161.468}{744.078}{all}%
\TextAt{8}{177.694}{744.078}{its}%
\TextAt{8}{193.979}{744.078}{various}%
\TextAt{8}{234.019}{744.078}{advantages}%
\TextAt{8}{293.905}{744.078}{cried}%
\TextAt{8}{322.481}{744.078}{out}%
\TextAt{8}{343.261}{744.078}{for}%
\TextAt{8}{361.113}{744.078}{a}%
\TextAt{8}{370.835}{744.078}{modern}%
\TextAt{8}{413.398}{744.078}{treatment.}%
\TextAt{8}{472.25}{744.078}{This}%
\TextAt{8}{498.936}{744.078}{having}%
\TextAt{8}{79.226}{727.032}{been}%
\TextAt{8}{106.191}{727.032}{done,}%
\TextAt{8}{136.876}{727.032}{there}%
\TextAt{8}{166.125}{727.032}{was}%
\TextAt{8}{187.961}{727.032}{the}%
\TextAt{8}{207.454}{727.032}{obvious}%
\TextAt{8}{249.136}{727.032}{opportunity}%
\TextAt{8}{313.199}{727.032}{of}%
\TextAt{8}{325.865}{727.032}{giving}%
\TextAt{8}{360.004}{727.032}{a}%
\TextAt{8}{369.096}{727.032}{thorough}%
\TextAt{8}{418.525}{727.032}{presentation}%
\TextAt{8}{484.596}{727.032}{of}%
\TextAt{8}{497.263}{727.032}{Artin’s}%
\TextAt{8}{79.226}{709.986}{L-series}%
\TextAt{8}{122.489}{709.986}{with}%
\TextAt{8}{149.863}{709.986}{their}%
\TextAt{8}{178.541}{709.986}{functional}%
\TextAt{8}{234.213}{709.986}{equation}%
\TextAt{8}{282.736}{709.986}{-}%
\TextAt{8}{291.245}{709.986}{which}%
\TextAt{8}{325.446}{709.986}{surprisingly}%
\TextAt{8}{389.702}{709.986}{has}%
\TextAt{8}{411.29}{709.986}{not}%
\TextAt{8}{432.818}{709.986}{been}%
\TextAt{8}{461.161}{709.986}{undertaken}%
\TextAt{8}{523.014}{709.986}{in}%
\TextAt{8}{79.226}{692.94}{any}%
\TextAt{8}{101.343}{692.94}{existing}%
\TextAt{8}{144.657}{692.94}{textbook.}%
\TextAt{8}{97.159}{672.905}{It}%
\TextAt{8}{109.222}{672.905}{was}%
\TextAt{8}{131.113}{672.905}{a}%
\TextAt{8}{140.26}{672.905}{difficult}%
\TextAt{8}{182.57}{672.905}{decision}%
\TextAt{8}{226.243}{672.905}{to}%
\TextAt{8}{239.944}{672.905}{exclude}%
\TextAt{8}{281.274}{672.905}{Iwasawa}%
\TextAt{8}{326.907}{672.905}{Theory,}%
\TextAt{8}{369.339}{672.905}{a}%
\TextAt{8}{378.486}{672.905}{relatively}%
\TextAt{8}{428.937}{672.905}{recent}%
\TextAt{8}{463.118}{672.905}{theory}%
\TextAt{8}{499.259}{672.905}{totally}%
\TextAt{8}{79.226}{655.859}{germane}%
\TextAt{8}{126.204}{655.859}{to}%
\TextAt{8}{140.663}{655.859}{algebraic}%
\TextAt{8}{190.247}{655.859}{number}%
\TextAt{8}{232.993}{655.859}{fields,}%
\TextAt{8}{266.403}{655.859}{the}%
\TextAt{8}{286.708}{655.859}{subject}%
\TextAt{8}{327.23}{655.859}{of}%
\TextAt{8}{340.709}{655.859}{this}%
\TextAt{8}{363.681}{655.859}{book.}%
\TextAt{8}{397.623}{655.859}{Since}%
\TextAt{8}{428.33}{655.859}{it}%
\TextAt{8}{440.183}{655.859}{mirrors}%
\TextAt{8}{481.374}{655.859}{important}%
\TextAt{8}{79.226}{638.813}{geometric}%
\TextAt{8}{132.594}{638.813}{properties}%
\TextAt{8}{187.648}{638.813}{of}%
\TextAt{8}{201.014}{638.813}{algebraic}%
\TextAt{8}{250.485}{638.813}{curves,}%
\TextAt{8}{289.612}{638.813}{it}%
\TextAt{8}{301.352}{638.813}{would}%
\TextAt{8}{335.532}{638.813}{have}%
\TextAt{8}{362.563}{638.813}{been}%
\TextAt{8}{390.227}{638.813}{a}%
\TextAt{8}{400.019}{638.813}{particularly}%
\TextAt{8}{463.477}{638.813}{beautiful}%
\TextAt{8}{512.936}{638.813}{vin-}%
\TextAt{8}{79.226}{621.767}{dication}%
\TextAt{8}{123.656}{621.767}{of}%
\TextAt{8}{136.536}{621.767}{our}%
\TextAt{8}{156.9}{621.767}{oft-repeated}%
\TextAt{8}{222.132}{621.767}{thesis}%
\TextAt{8}{254.32}{621.767}{that}%
\TextAt{8}{279.238}{621.767}{number}%
\TextAt{8}{321.385}{621.767}{theory}%
\TextAt{8}{357.685}{621.767}{is}%
\TextAt{8}{368.999}{621.767}{geometry.}%
\TextAt{8}{423.48}{621.767}{I}%
\TextAt{8}{431.147}{621.767}{do}%
\TextAt{8}{446.957}{621.767}{believe,}%
\TextAt{8}{488.535}{621.767}{however,}%
\TextAt{8}{79.226}{604.721}{that}%
\TextAt{8}{105.819}{604.721}{in}%
\TextAt{8}{120.695}{604.721}{this}%
\TextAt{8}{144.741}{604.721}{case}%
\TextAt{8}{170.736}{604.721}{the}%
\TextAt{8}{192.116}{604.721}{geometric}%
\TextAt{8}{246.672}{604.721}{aspect}%
\TextAt{8}{284.048}{604.721}{becomes}%
\TextAt{8}{331.825}{604.721}{truly}%
\TextAt{8}{361.992}{604.721}{convincing}%
\TextAt{8}{421.103}{604.721}{only}%
\TextAt{8}{448.018}{604.721}{if}%
\TextAt{8}{459.966}{604.721}{one}%
\TextAt{8}{482.649}{604.721}{uses}%
\TextAt{8}{508.704}{604.721}{étale}%
\TextAt{8}{79.226}{587.675}{cohomology}%
\TextAt{8}{143.292}{587.675}{-}%
\TextAt{8}{151.074}{587.675}{which}%
\TextAt{8}{184.547}{587.675}{can}%
\TextAt{8}{205.993}{587.675}{neither}%
\TextAt{8}{245.647}{587.675}{be}%
\TextAt{8}{261.558}{587.675}{assumed}%
\TextAt{8}{308.493}{587.675}{nor}%
\TextAt{8}{329.293}{587.675}{reasonably}%
\TextAt{8}{387.561}{587.675}{developed}%
\TextAt{8}{441.848}{587.675}{here.}%
\TextAt{8}{471.752}{587.675}{Perhaps}%
\TextAt{8}{516.51}{587.675}{the}%
\TextAt{8}{79.226}{570.629}{dissatisfaction}%
\TextAt{8}{155.895}{570.629}{with}%
\TextAt{8}{183.262}{570.629}{this}%
\TextAt{8}{206.79}{570.629}{exclusion}%
\TextAt{8}{257.959}{570.629}{will}%
\TextAt{8}{280.771}{570.629}{be}%
\TextAt{8}{297.401}{570.629}{strong}%
\TextAt{8}{333.949}{570.629}{enough}%
\TextAt{8}{374.98}{570.629}{to}%
\TextAt{8}{389.997}{570.629}{bring}%
\TextAt{8}{421.272}{570.629}{about}%
\TextAt{8}{455.477}{570.629}{a}%
\TextAt{8}{465.938}{570.629}{sequel}%
\TextAt{8}{501.494}{570.629}{to}%
\TextAt{8}{516.51}{570.629}{the}%
\TextAt{8}{79.226}{553.583}{present}%
\TextAt{8}{119.933}{553.583}{volume,}%
\TextAt{8}{163.51}{553.583}{devoted}%
\TextAt{8}{207.087}{553.583}{to}%
\TextAt{8}{221.397}{553.583}{the}%
\TextAt{8}{241.553}{553.583}{cohomology}%
\TextAt{8}{305.633}{553.583}{of}%
\TextAt{8}{318.963}{553.583}{algebraic}%
\TextAt{8}{368.398}{553.583}{number}%
\TextAt{8}{410.994}{553.583}{fields.}%
\TextAt{8}{97.159}{533.548}{From}%
\TextAt{8}{128.603}{533.548}{the}%
\TextAt{8}{149.478}{533.548}{very}%
\TextAt{8}{175.888}{533.548}{start}%
\TextAt{8}{204.642}{533.548}{the}%
\TextAt{8}{225.518}{533.548}{book}%
\TextAt{8}{255.18}{533.548}{was}%
\TextAt{8}{278.398}{533.548}{not}%
\TextAt{8}{299.931}{533.548}{just}%
\TextAt{8}{323.795}{533.548}{intended}%
\TextAt{8}{372.311}{533.548}{as}%
\TextAt{8}{387.4}{533.548}{a}%
\TextAt{8}{397.874}{533.548}{modern}%
\TextAt{8}{441.19}{533.548}{sourcebook}%
\TextAt{8}{502.784}{533.548}{on}%
\TextAt{8}{519.762}{533.548}{al-}%
\TextAt{8}{79.226}{516.503}{gebraic}%
\TextAt{8}{120.53}{516.503}{number}%
\TextAt{8}{164.106}{516.503}{theory,}%
\TextAt{8}{204.352}{516.503}{but}%
\TextAt{8}{226.791}{516.503}{also}%
\TextAt{8}{251.25}{516.503}{as}%
\TextAt{8}{266.6}{516.503}{a}%
\TextAt{8}{277.335}{516.503}{convenient}%
\TextAt{8}{336.201}{516.503}{textbook}%
\TextAt{8}{386.615}{516.503}{for}%
\TextAt{8}{405.479}{516.503}{a}%
\TextAt{8}{416.214}{516.503}{course.}%
\TextAt{8}{459.533}{516.503}{This}%
\TextAt{8}{487.232}{516.503}{intention}%
\TextAt{8}{79.226}{499.457}{was}%
\TextAt{8}{102.111}{499.457}{increasingly}%
\TextAt{8}{166.624}{499.457}{jeopardized}%
\TextAt{8}{229.116}{499.457}{by}%
\TextAt{8}{245.761}{499.457}{the}%
\TextAt{8}{266.303}{499.457}{unexpected}%
\TextAt{8}{328.46}{499.457}{growth}%
\TextAt{8}{368.202}{499.457}{of}%
\TextAt{8}{381.917}{499.457}{the}%
\TextAt{8}{402.459}{499.457}{material}%
\TextAt{8}{449.027}{499.457}{which}%
\TextAt{8}{482.899}{499.457}{had}%
\TextAt{8}{506.047}{499.457}{to}%
\TextAt{8}{520.743}{499.457}{be}%
\TextAt{8}{79.226}{482.411}{covered}%
\TextAt{8}{121.754}{482.411}{in}%
\TextAt{8}{135.984}{482.411}{view}%
\TextAt{8}{163.545}{482.411}{of}%
\TextAt{8}{177.452}{482.411}{the}%
\TextAt{8}{198.186}{482.411}{intrinsic}%
\TextAt{8}{244.026}{482.411}{necessities}%
\TextAt{8}{300.709}{482.411}{of}%
\TextAt{8}{314.616}{482.411}{the}%
\TextAt{8}{335.35}{482.411}{theory.}%
\TextAt{8}{377.403}{482.411}{Yet}%
\TextAt{8}{399.428}{482.411}{I}%
\TextAt{8}{408.122}{482.411}{think}%
\TextAt{8}{439.592}{482.411}{that}%
\TextAt{8}{465.538}{482.411}{the}%
\TextAt{8}{486.272}{482.411}{book}%
\TextAt{8}{515.793}{482.411}{has}%
\TextAt{8}{79.226}{465.365}{not}%
\TextAt{8}{100.644}{465.365}{lost}%
\TextAt{8}{123.424}{465.365}{that}%
\TextAt{8}{149.397}{465.365}{character.}%
\TextAt{8}{206.82}{465.365}{In}%
\TextAt{8}{222.045}{465.365}{fact,}%
\TextAt{8}{249.137}{465.365}{it}%
\TextAt{8}{261.445}{465.365}{has}%
\TextAt{8}{282.922}{465.365}{passed}%
\TextAt{8}{320.718}{465.365}{a}%
\TextAt{8}{331.078}{465.365}{first}%
\TextAt{8}{355.807}{465.365}{test}%
\TextAt{8}{379.233}{465.365}{in}%
\TextAt{8}{393.49}{465.365}{this}%
\TextAt{8}{416.916}{465.365}{respect.}%
\TextAt{8}{463.328}{465.365}{With}%
\TextAt{8}{494.166}{465.365}{a}%
\TextAt{8}{504.525}{465.365}{bit}%
\TextAt{8}{523.337}{465.365}{of}%
\TextAt{8}{79.226}{448.319}{careful}%
\TextAt{8}{117.767}{448.319}{planning,}%
\TextAt{8}{169.774}{448.319}{the}%
\TextAt{8}{190.43}{448.319}{basic}%
\TextAt{8}{220.255}{448.319}{content}%
\TextAt{8}{262.382}{448.319}{of}%
\TextAt{8}{276.211}{448.319}{the}%
\TextAt{8}{296.867}{448.319}{first}%
\TextAt{8}{321.491}{448.319}{three}%
\TextAt{8}{351.902}{448.319}{chapters}%
\TextAt{8}{398.967}{448.319}{can}%
\TextAt{8}{420.926}{448.319}{easily}%
\TextAt{8}{453.68}{448.319}{be}%
\TextAt{8}{470.103}{448.319}{presented}%
\TextAt{8}{523.014}{448.319}{in}%
\TextAt{8}{79.226}{431.273}{one}%
\TextAt{8}{101.034}{431.273}{academic}%
\TextAt{8}{152.107}{431.273}{year}%
\TextAt{8}{177.824}{431.273}{(if}%
\TextAt{8}{193.45}{431.273}{possible}%
\TextAt{8}{237.817}{431.273}{including}%
\TextAt{8}{288.891}{431.273}{infinite}%
\TextAt{8}{328.906}{431.273}{Galois}%
\TextAt{8}{365.167}{431.273}{theory).}%
\TextAt{8}{412.067}{431.273}{The}%
\TextAt{8}{436.468}{431.273}{following}%
\TextAt{8}{486.251}{431.273}{term}%
\TextAt{8}{514.562}{431.273}{will}%
\TextAt{8}{79.226}{414.227}{then}%
\TextAt{8}{105.886}{414.227}{provide}%
\TextAt{8}{147.514}{414.227}{scarce,}%
\TextAt{8}{185.292}{414.227}{yet}%
\TextAt{8}{204.803}{414.227}{sufficient}%
\TextAt{8}{253.963}{414.227}{room}%
\TextAt{8}{284.209}{414.227}{for}%
\TextAt{8}{302.094}{414.227}{the}%
\TextAt{8}{322.251}{414.227}{class}%
\TextAt{8}{349.688}{414.227}{field}%
\TextAt{8}{375.045}{414.227}{theory}%
\TextAt{8}{411.795}{414.227}{of}%
\TextAt{8}{425.125}{414.227}{chapters}%
\TextAt{8}{471.69}{414.227}{IV-VI.}%
\TextAt{8}{97.159}{394.192}{Sections}%
\TextAt{8}{144.668}{394.192}{11-14}%
\TextAt{8}{177.819}{394.192}{of}%
\TextAt{8}{193.073}{394.192}{chapter}%
\TextAt{8}{236.947}{394.192}{I}%
\TextAt{8}{246.987}{394.192}{may}%
\TextAt{8}{274.28}{394.192}{mostly}%
\TextAt{8}{314.317}{394.192}{be}%
\TextAt{8}{332.165}{394.192}{dropped}%
\TextAt{8}{379.937}{394.192}{from}%
\TextAt{8}{409.501}{394.192}{an}%
\TextAt{8}{427.684}{394.192}{introductory}%
\TextAt{8}{497.585}{394.192}{course.}%
\TextAt{8}{79.226}{377.146}{Although}%
\TextAt{8}{131.667}{377.146}{the}%
\TextAt{8}{152.558}{377.146}{results}%
\TextAt{8}{190.485}{377.146}{of}%
\TextAt{8}{204.55}{377.146}{section}%
\TextAt{8}{244.366}{377.146}{12}%
\TextAt{8}{260.714}{377.146}{on}%
\TextAt{8}{277.707}{377.146}{orders}%
\TextAt{8}{313.626}{377.146}{are}%
\TextAt{8}{333.872}{377.146}{irrelevant}%
\TextAt{8}{386.647}{377.146}{for}%
\TextAt{8}{405.267}{377.146}{the}%
\TextAt{8}{426.158}{377.146}{sequel,}%
\TextAt{8}{465.177}{377.146}{I}%
\TextAt{8}{474.028}{377.146}{consider}%
\TextAt{8}{520.348}{377.146}{its}%
\TextAt{8}{79.226}{360.1}{insertion}%
\TextAt{8}{127.997}{360.1}{in}%
\TextAt{8}{142.23}{360.1}{the}%
\TextAt{8}{162.967}{360.1}{book}%
\TextAt{8}{192.49}{360.1}{particularly}%
\TextAt{8}{256.492}{360.1}{important.}%
\TextAt{8}{318.076}{360.1}{For}%
\TextAt{8}{339.625}{360.1}{one}%
\TextAt{8}{361.665}{360.1}{thing,}%
\TextAt{8}{396.211}{360.1}{orders}%
\TextAt{8}{431.975}{360.1}{constitute}%
\TextAt{8}{487.25}{360.1}{the}%
\TextAt{8}{507.986}{360.1}{rings}%
\TextAt{8}{79.226}{343.054}{of}%
\TextAt{8}{94.238}{343.054}{multipliers,}%
\TextAt{8}{157.861}{343.054}{which}%
\TextAt{8}{193.029}{343.054}{play}%
\TextAt{8}{220.079}{343.054}{an}%
\TextAt{8}{238.02}{343.054}{eminent}%
\TextAt{8}{284.246}{343.054}{role}%
\TextAt{8}{308.691}{343.054}{in}%
\TextAt{8}{324.025}{343.054}{many}%
\TextAt{8}{357.579}{343.054}{diophantine}%
\TextAt{8}{423.328}{343.054}{problems.}%
\TextAt{8}{483.064}{343.054}{But}%
\TextAt{8}{507.986}{343.054}{most}%
\TextAt{8}{79.226}{326.008}{importantly,}%
\TextAt{8}{147.374}{326.008}{they}%
\TextAt{8}{174.632}{326.008}{represent}%
\TextAt{8}{226.015}{326.008}{the}%
\TextAt{8}{247.092}{326.008}{analogues}%
\TextAt{8}{301.416}{326.008}{of}%
\TextAt{8}{315.667}{326.008}{singular}%
\TextAt{8}{360.882}{326.008}{algebraic}%
\TextAt{8}{411.237}{326.008}{curves.}%
\TextAt{8}{454.379}{326.008}{As}%
\TextAt{8}{472.587}{326.008}{cohomology}%
\TextAt{8}{79.226}{308.962}{theory}%
\TextAt{8}{115.883}{308.962}{becomes}%
\TextAt{8}{162.343}{308.962}{increasingly}%
\TextAt{8}{226.378}{308.962}{important}%
\TextAt{8}{281.577}{308.962}{for}%
\TextAt{8}{299.369}{308.962}{algebraic}%
\TextAt{8}{348.71}{308.962}{number}%
\TextAt{8}{391.213}{308.962}{fields,}%
\TextAt{8}{424.362}{308.962}{and}%
\TextAt{8}{447.032}{308.962}{since}%
\TextAt{8}{475.607}{308.962}{this}%
\TextAt{8}{498.336}{308.962}{is}%
\TextAt{8}{510.007}{308.962}{even}%
\TextAt{8}{79.226}{291.916}{more}%
\TextAt{8}{108.667}{291.916}{true}%
\TextAt{8}{133.552}{291.916}{of}%
\TextAt{8}{147.057}{291.916}{algebraic}%
\TextAt{8}{196.666}{291.916}{K-theory,}%
\TextAt{8}{248.9}{291.916}{which}%
\TextAt{8}{282.561}{291.916}{cannot}%
\TextAt{8}{321.111}{291.916}{be}%
\TextAt{8}{337.21}{291.916}{constructed}%
\TextAt{8}{400.531}{291.916}{without}%
\TextAt{8}{444.282}{291.916}{singular}%
\TextAt{8}{488.751}{291.916}{schemes,}%
\TextAt{8}{79.226}{274.87}{the}%
\TextAt{8}{99.383}{274.87}{time}%
\TextAt{8}{126.043}{274.87}{has}%
\TextAt{8}{146.916}{274.87}{come}%
\TextAt{8}{176.828}{274.87}{to}%
\TextAt{8}{191.139}{274.87}{give}%
\TextAt{8}{215.204}{274.87}{orders}%
\TextAt{8}{250.388}{274.87}{an}%
\TextAt{8}{266.647}{274.87}{adequate}%
\TextAt{8}{316.405}{274.87}{treatment.}%
\TextAt{8}{97.159}{254.835}{In}%
\TextAt{8}{112.101}{254.835}{chapter}%
\TextAt{8}{154.373}{254.835}{II,}%
\TextAt{8}{170.606}{254.835}{the}%
\TextAt{8}{191.084}{254.835}{special}%
\TextAt{8}{229.506}{254.835}{treatment}%
\TextAt{8}{284.139}{254.835}{of}%
\TextAt{8}{297.791}{254.835}{henselian}%
\TextAt{8}{348.897}{254.835}{fields}%
\TextAt{8}{379.19}{254.835}{in}%
\TextAt{8}{393.164}{254.835}{section}%
\TextAt{8}{432.566}{254.835}{6}%
\TextAt{8}{442.643}{254.835}{may}%
\TextAt{8}{468.333}{254.835}{be}%
\TextAt{8}{484.578}{254.835}{restricted}%
\TextAt{8}{79.226}{237.789}{to}%
\TextAt{8}{93.828}{237.789}{complete}%
\TextAt{8}{143.543}{237.789}{valued}%
\TextAt{8}{180.585}{237.789}{fields,}%
\TextAt{8}{214.173}{237.789}{and}%
\TextAt{8}{237.228}{237.789}{thus}%
\TextAt{8}{263.271}{237.789}{joined}%
\TextAt{8}{298.353}{237.789}{with}%
\TextAt{8}{325.305}{237.789}{section}%
\TextAt{8}{364.678}{237.789}{4.}%
\TextAt{8}{379.86}{237.789}{If}%
\TextAt{8}{391.845}{237.789}{pressed}%
\TextAt{8}{433.226}{237.789}{for}%
\TextAt{8}{451.403}{237.789}{time,}%
\TextAt{8}{481.68}{237.789}{section}%
\TextAt{8}{521.053}{237.789}{10}%
\TextAt{8}{79.226}{220.743}{on}%
\TextAt{8}{95.485}{220.743}{higher}%
\TextAt{8}{131.255}{220.743}{ramification}%
\TextAt{8}{196.303}{220.743}{may}%
\TextAt{8}{221.672}{220.743}{be}%
\TextAt{8}{237.596}{220.743}{omitted}%
\TextAt{8}{281.173}{220.743}{completely.}%
\TextAt{8}{97.159}{200.709}{The}%
\TextAt{8}{121.782}{200.709}{first}%
\TextAt{8}{146.477}{200.709}{three}%
\TextAt{8}{176.959}{200.709}{sections}%
\TextAt{8}{221.225}{200.709}{of}%
\TextAt{8}{235.124}{200.709}{chapter}%
\TextAt{8}{277.645}{200.709}{III}%
\TextAt{8}{295.417}{200.709}{should}%
\TextAt{8}{333.12}{200.709}{be}%
\TextAt{8}{349.613}{200.709}{presented}%
\TextAt{8}{402.595}{200.709}{in}%
\TextAt{8}{416.817}{200.709}{the}%
\TextAt{8}{437.543}{200.709}{lectures}%
\TextAt{8}{481.091}{200.709}{since}%
\TextAt{8}{510.33}{200.709}{they}%
\TextAt{8}{79.226}{183.663}{highlight}%
\TextAt{8}{128.154}{183.663}{a}%
\TextAt{8}{137.725}{183.663}{new}%
\TextAt{8}{161.594}{183.663}{approach}%
\TextAt{8}{211.825}{183.663}{to}%
\TextAt{8}{225.951}{183.663}{classical}%
\TextAt{8}{270.766}{183.663}{results}%
\TextAt{8}{307.775}{183.663}{of}%
\TextAt{8}{320.92}{183.663}{algebraic}%
\TextAt{8}{370.171}{183.663}{number}%
\TextAt{8}{412.583}{183.663}{theory.}%
\TextAt{8}{452.842}{183.663}{The}%
\TextAt{8}{476.712}{183.663}{subsequent}%
\TextAt{8}{79.226}{166.617}{theory}%
\TextAt{8}{117.578}{166.617}{concerning}%
\TextAt{8}{177.713}{166.617}{the}%
\TextAt{8}{199.471}{166.617}{theorem}%
\TextAt{8}{246.599}{166.617}{of}%
\TextAt{8}{261.531}{166.617}{Grothendieck-Riemann-Roch}%
\TextAt{8}{415.071}{166.617}{is}%
\TextAt{8}{428.437}{166.617}{a}%
\TextAt{8}{439.794}{166.617}{nice}%
\TextAt{8}{465.45}{166.617}{subject}%
\TextAt{8}{507.425}{166.617}{for}%
\TextAt{8}{526.911}{166.617}{a}%
\TextAt{8}{79.226}{149.571}{student}%
\TextAt{8}{121.237}{149.571}{seminar}%
\TextAt{8}{164.873}{149.571}{rather}%
\TextAt{8}{199.997}{149.571}{than}%
\TextAt{8}{227.315}{149.571}{for}%
\TextAt{8}{245.2}{149.571}{an}%
\TextAt{8}{261.459}{149.571}{introductory}%
\TextAt{8}{329.436}{149.571}{course.}%
\TextAt{8}{97.159}{129.536}{Finally,}%
\TextAt{8}{139.646}{129.536}{in}%
\TextAt{8}{153.603}{129.536}{presenting}%
\TextAt{8}{210.229}{129.536}{class}%
\TextAt{8}{237.97}{129.536}{field}%
\TextAt{8}{263.632}{129.536}{theory,}%
\TextAt{8}{303.034}{129.536}{it}%
\TextAt{8}{315.043}{129.536}{saves}%
\TextAt{8}{345.068}{129.536}{considerable}%
\TextAt{8}{411.772}{129.536}{time}%
\TextAt{8}{438.737}{129.536}{if}%
\TextAt{8}{449.765}{129.536}{the}%
\TextAt{8}{470.226}{129.536}{students}%
\TextAt{8}{517.156}{129.536}{are}%
\TextAt{8}{79.226}{112.49}{already}%
\TextAt{8}{120.714}{112.49}{familiar}%
\TextAt{8}{164.151}{112.49}{with}%
\TextAt{8}{190.994}{112.49}{profinite}%
\TextAt{8}{238.006}{112.49}{groups}%
\TextAt{8}{275.979}{112.49}{and}%
\TextAt{8}{298.925}{112.49}{infinite}%
\TextAt{8}{338.775}{112.49}{Galois}%
\TextAt{8}{374.871}{112.49}{theory.}%
\TextAt{8}{415.741}{112.49}{Sections}%
\TextAt{8}{461.509}{112.49}{4--6}%
\TextAt{8}{481.203}{112.49}{of}%
\TextAt{8}{494.716}{112.49}{chapter}%
\TextAt{8}{79.226}{95.444}{V,}%
\TextAt{8}{106.391}{95.444}{on}%
\TextAt{8}{122.928}{95.444}{formal}%
\TextAt{8}{159.956}{95.444}{groups,}%
\TextAt{8}{201.346}{95.444}{Lubin-Tate}%
\TextAt{8}{262.583}{95.444}{theory}%
\TextAt{8}{299.611}{95.444}{and}%
\TextAt{8}{322.652}{95.444}{the}%
\TextAt{8}{343.086}{95.444}{theory}%
\TextAt{8}{380.114}{95.444}{of}%
\TextAt{8}{393.722}{95.444}{higher}%
\TextAt{8}{429.77}{95.444}{ramification}%
\TextAt{8}{495.096}{95.444}{may}%
\TextAt{8}{520.743}{95.444}{be}%
\TextAt{8}{298.03}{43.033}{viii}%
\Destination{78.23}{809.815}{nav.page.viii}%
\Destination{79.226}{126}{errata.212}%
}

% Page 9 (body page front matter)
\Page{%
\TextAt{2}{79.226}{796.943}{Preface}%
\TextAt{2}{117.888}{796.943}{to}%
\TextAt{2}{131.219}{796.943}{the}%
\TextAt{2}{150.004}{796.943}{German}%
\TextAt{2}{191.928}{796.943}{Edition}%
\Vector{q 0 G 0 g 0.299 w 7.226 18.729 m 460.769 18.729 l S Q}%
\TextAt{8}{79.226}{761.124}{omitted.}%
\TextAt{8}{128.376}{761.124}{Cutting}%
\TextAt{8}{172.293}{761.124}{out}%
\TextAt{8}{193.448}{761.124}{even}%
\TextAt{8}{220.448}{761.124}{more,}%
\TextAt{8}{253.392}{761.124}{chapter}%
\TextAt{8}{295.684}{761.124}{V,}%
\TextAt{8}{311.949}{761.124}{3,}%
\TextAt{8}{325.382}{761.124}{on}%
\TextAt{8}{341.982}{761.124}{the}%
\TextAt{8}{362.479}{761.124}{Hilbert}%
\TextAt{8}{403.132}{761.124}{symbol,}%
\TextAt{8}{446.872}{761.124}{and}%
\TextAt{8}{469.975}{761.124}{VI,}%
\TextAt{8}{490.46}{761.124}{7}%
\TextAt{8}{500.556}{761.124}{and}%
\TextAt{8}{523.66}{761.124}{8,}%
\TextAt{8}{79.226}{744.078}{still}%
\TextAt{8}{102.698}{744.078}{leaves}%
\TextAt{8}{136.905}{744.078}{a}%
\TextAt{8}{147.309}{744.078}{fully-fledged}%
\TextAt{8}{214.285}{744.078}{theory,}%
\TextAt{8}{254.118}{744.078}{which}%
\TextAt{8}{288.254}{744.078}{is}%
\TextAt{8}{300.667}{744.078}{however}%
\TextAt{8}{346.195}{744.078}{unsatisfactory}%
\TextAt{8}{422.425}{744.078}{because}%
\TextAt{8}{466.376}{744.078}{it}%
\TextAt{8}{478.729}{744.078}{remains}%
\TextAt{8}{523.014}{744.078}{in}%
\TextAt{8}{79.226}{727.032}{the}%
\TextAt{8}{99.383}{727.032}{abstract}%
\TextAt{8}{144.98}{727.032}{realm,}%
\TextAt{8}{180.749}{727.032}{and}%
\TextAt{8}{203.512}{727.032}{is}%
\TextAt{8}{215.276}{727.032}{never}%
\TextAt{8}{246.491}{727.032}{linked}%
\TextAt{8}{280.958}{727.032}{to}%
\TextAt{8}{295.268}{727.032}{classical}%
\TextAt{8}{340.267}{727.032}{problems.}%
\TextAt{8}{97.159}{706.997}{A}%
\TextAt{8}{110.614}{706.997}{word}%
\TextAt{8}{140.341}{706.997}{on}%
\TextAt{8}{157.384}{706.997}{the}%
\TextAt{8}{178.324}{706.997}{exercises}%
\TextAt{8}{227.024}{706.997}{at}%
\TextAt{8}{242.118}{706.997}{the}%
\TextAt{8}{263.058}{706.997}{end}%
\TextAt{8}{285.947}{706.997}{of}%
\TextAt{8}{300.06}{706.997}{the}%
\TextAt{8}{321}{706.997}{sections.}%
\TextAt{8}{371.598}{706.997}{Some}%
\TextAt{8}{403.597}{706.997}{of}%
\TextAt{8}{417.711}{706.997}{them}%
\TextAt{8}{448.406}{706.997}{are}%
\TextAt{8}{468.701}{706.997}{not}%
\TextAt{8}{490.298}{706.997}{so}%
\TextAt{8}{505.452}{706.997}{much}%
\TextAt{8}{79.226}{689.951}{exercises,}%
\TextAt{8}{130.285}{689.951}{but}%
\TextAt{8}{151.608}{689.951}{additional}%
\TextAt{8}{206.764}{689.951}{remarks}%
\TextAt{8}{251.244}{689.951}{which}%
\TextAt{8}{284.593}{689.951}{did}%
\TextAt{8}{304.613}{689.951}{not}%
\TextAt{8}{325.29}{689.951}{fit}%
\TextAt{8}{340.109}{689.951}{well}%
\TextAt{8}{363.703}{689.951}{into}%
\TextAt{8}{387.31}{689.951}{the}%
\TextAt{8}{407.329}{689.951}{main}%
\TextAt{8}{436.459}{689.951}{text.}%
\TextAt{8}{465.353}{689.951}{The}%
\TextAt{8}{489.27}{689.951}{reader}%
\TextAt{8}{524.903}{689.951}{is}%
\TextAt{8}{79.226}{672.905}{encouraged}%
\TextAt{8}{140.973}{672.905}{to}%
\TextAt{8}{155.891}{672.905}{prove}%
\TextAt{8}{188.048}{672.905}{his}%
\TextAt{8}{206.923}{672.905}{versatility}%
\TextAt{8}{262.238}{672.905}{in}%
\TextAt{8}{276.499}{672.905}{looking}%
\TextAt{8}{318.089}{672.905}{up}%
\TextAt{8}{335.601}{672.905}{the}%
\TextAt{8}{356.365}{672.905}{literature.}%
\TextAt{8}{414.123}{672.905}{I}%
\TextAt{8}{422.848}{672.905}{should}%
\TextAt{8}{460.588}{672.905}{also}%
\TextAt{8}{484.676}{672.905}{point}%
\TextAt{8}{515.853}{672.905}{out}%
\TextAt{8}{79.226}{655.859}{that}%
\TextAt{8}{104.932}{655.859}{I}%
\TextAt{8}{113.386}{655.859}{have}%
\TextAt{8}{140.718}{655.859}{not}%
\TextAt{8}{161.869}{655.859}{actually}%
\TextAt{8}{206.763}{655.859}{done}%
\TextAt{8}{235.064}{655.859}{all}%
\TextAt{8}{251.66}{655.859}{the}%
\TextAt{8}{272.153}{655.859}{exercises}%
\TextAt{8}{320.407}{655.859}{myself,}%
\TextAt{8}{360.233}{655.859}{so}%
\TextAt{8}{374.94}{655.859}{that}%
\TextAt{8}{400.646}{655.859}{there}%
\TextAt{8}{430.895}{655.859}{might}%
\TextAt{8}{464.73}{655.859}{be}%
\TextAt{8}{480.992}{655.859}{occasional}%
\TextAt{8}{79.226}{638.813}{mistakes}%
\TextAt{8}{127.113}{638.813}{in}%
\TextAt{8}{141.046}{638.813}{the}%
\TextAt{8}{161.484}{638.813}{way}%
\TextAt{8}{185.508}{638.813}{they}%
\TextAt{8}{212.126}{638.813}{are}%
\TextAt{8}{231.917}{638.813}{posed.}%
\TextAt{8}{270.212}{638.813}{If}%
\TextAt{8}{282.185}{638.813}{such}%
\TextAt{8}{308.863}{638.813}{a}%
\TextAt{8}{318.899}{638.813}{case}%
\TextAt{8}{343.951}{638.813}{arises,}%
\TextAt{8}{379.546}{638.813}{it}%
\TextAt{8}{391.531}{638.813}{is}%
\TextAt{8}{403.576}{638.813}{for}%
\TextAt{8}{421.742}{638.813}{the}%
\TextAt{8}{442.179}{638.813}{reader}%
\TextAt{8}{478.23}{638.813}{to}%
\TextAt{8}{492.821}{638.813}{find}%
\TextAt{8}{516.51}{638.813}{the}%
\TextAt{8}{79.226}{621.767}{correct}%
\TextAt{8}{119.036}{621.767}{formulation.}%
\TextAt{8}{189.052}{621.767}{May}%
\TextAt{8}{216.178}{621.767}{the}%
\TextAt{8}{237.123}{621.767}{reader’s}%
\TextAt{8}{281.548}{621.767}{reaction}%
\TextAt{8}{327.216}{621.767}{to}%
\TextAt{8}{342.315}{621.767}{such}%
\TextAt{8}{369.501}{621.767}{a}%
\TextAt{8}{380.045}{621.767}{possible}%
\TextAt{8}{424.852}{621.767}{slip}%
\TextAt{8}{447.16}{621.767}{of}%
\TextAt{8}{461.279}{621.767}{the}%
\TextAt{8}{482.223}{621.767}{author}%
\TextAt{8}{520.743}{621.767}{be}%
\TextAt{8}{79.226}{604.721}{mitigated}%
\TextAt{8}{131.913}{604.721}{by}%
\TextAt{8}{148.172}{604.721}{Goethe’s}%
\TextAt{8}{196.757}{604.721}{distich:}%
\TextAt{8}{109.114}{571.735}{“Irrtum}%
\TextAt{8}{152.679}{571.735}{verläßt}%
\TextAt{8}{191.712}{571.735}{uns}%
\TextAt{8}{213.232}{571.735}{nie,}%
\TextAt{8}{235.337}{571.735}{doch}%
\TextAt{8}{263.3}{571.735}{ziehet}%
\TextAt{8}{297.109}{571.735}{ein}%
\TextAt{8}{315.962}{571.735}{höher}%
\TextAt{8}{348.48}{571.735}{Bedürfnis}%
\TextAt{8}{109.114}{554.689}{Immer}%
\TextAt{8}{146.498}{554.689}{den}%
\TextAt{8}{168.603}{554.689}{strebenden}%
\TextAt{8}{228.164}{554.689}{Geist}%
\TextAt{8}{258.865}{554.689}{leise}%
\TextAt{8}{284.281}{554.689}{zur}%
\TextAt{8}{304.438}{554.689}{Wahrheit}%
\TextAt{8}{355.809}{554.689}{hinan.}%
\TextAt{8}{386.068}{554.689}{”}%
\TextAt{7}{391.579}{559.027}{*}%
\TextAt{8}{97.159}{521.703}{During}%
\TextAt{8}{135.84}{521.703}{the}%
\TextAt{8}{155.166}{521.703}{writing}%
\TextAt{8}{194.66}{521.703}{of}%
\TextAt{8}{207.159}{521.703}{this}%
\TextAt{8}{229.151}{521.703}{book}%
\TextAt{8}{257.263}{521.703}{I}%
\TextAt{8}{264.549}{521.703}{have}%
\TextAt{8}{290.713}{521.703}{been}%
\TextAt{8}{317.511}{521.703}{helped}%
\TextAt{8}{354.064}{521.703}{in}%
\TextAt{8}{366.886}{521.703}{many}%
\TextAt{8}{397.928}{521.703}{ways.}%
\TextAt{8}{430.559}{521.703}{I}%
\TextAt{8}{437.845}{521.703}{thank}%
\TextAt{8}{470.513}{521.703}{the}%
\TextAt{8}{489.838}{521.703}{Springer}%
\TextAt{8}{79.226}{504.657}{Verlag}%
\TextAt{8}{116.955}{504.657}{for}%
\TextAt{8}{136.154}{504.657}{considering}%
\TextAt{8}{198.666}{504.657}{my}%
\TextAt{8}{219.491}{504.657}{wishes}%
\TextAt{8}{257.34}{504.657}{with}%
\TextAt{8}{285.314}{504.657}{generosity.}%
\TextAt{8}{348.178}{504.657}{My}%
\TextAt{8}{370.294}{504.657}{students}%
\TextAt{8}{418.233}{504.657}{I.}%
\TextAt{8}{430.916}{504.657}{Kausz,}%
\TextAt{8}{470.982}{504.657}{B.}%
\TextAt{8}{487.73}{504.657}{Köck,}%
\TextAt{8}{522.536}{504.657}{P.}%
\TextAt{8}{79.226}{487.611}{Kölcze,}%
\TextAt{8}{121.771}{487.611}{TH.}%
\TextAt{8}{147.675}{487.611}{Moser,}%
\TextAt{8}{187.686}{487.611}{M.}%
\TextAt{8}{207.086}{487.611}{Spiess}%
\TextAt{8}{243.2}{487.611}{have}%
\TextAt{8}{271.722}{487.611}{critically}%
\TextAt{8}{321.704}{487.611}{examined}%
\TextAt{8}{375.583}{487.611}{larger}%
\TextAt{8}{410.286}{487.611}{or}%
\TextAt{8}{426.124}{487.611}{smaller}%
\TextAt{8}{468.036}{487.611}{parts,}%
\TextAt{8}{503.18}{487.611}{which}%
\TextAt{8}{79.226}{470.565}{led}%
\TextAt{8}{98.549}{470.565}{to}%
\TextAt{8}{113.328}{470.565}{numerous}%
\TextAt{8}{166.867}{470.565}{improvements}%
\TextAt{8}{242.199}{470.565}{and}%
\TextAt{8}{265.431}{470.565}{made}%
\TextAt{8}{297.115}{470.565}{it}%
\TextAt{8}{309.288}{470.565}{possible}%
\TextAt{8}{353.777}{470.565}{to}%
\TextAt{8}{368.556}{470.565}{avoid}%
\TextAt{8}{399.93}{470.565}{mistakes}%
\TextAt{8}{448.004}{470.565}{and}%
\TextAt{8}{471.236}{470.565}{ambiguities.}%
\TextAt{8}{79.226}{453.519}{To}%
\TextAt{8}{97.256}{453.519}{my}%
\TextAt{8}{117.569}{453.519}{friends}%
\TextAt{8}{156.472}{453.519}{W.-D.}%
\TextAt{8}{192.542}{453.519}{Geyer,}%
\TextAt{8}{230.69}{453.519}{G.}%
\TextAt{8}{247.823}{453.519}{Tamme,}%
\TextAt{8}{294.016}{453.519}{and}%
\TextAt{8}{317.581}{453.519}{K.}%
\TextAt{8}{334.631}{453.519}{Wingberg}%
\TextAt{8}{389.411}{453.519}{I}%
\TextAt{8}{398.33}{453.519}{owe}%
\TextAt{8}{421.895}{453.519}{much}%
\TextAt{8}{453.912}{453.519}{valuable}%
\TextAt{8}{500.574}{453.519}{advice}%
\TextAt{8}{79.226}{436.473}{from}%
\TextAt{8}{106.578}{436.473}{which}%
\TextAt{8}{139.776}{436.473}{the}%
\TextAt{8}{159.644}{436.473}{book}%
\TextAt{8}{188.299}{436.473}{has}%
\TextAt{8}{208.885}{436.473}{profited,}%
\TextAt{8}{255.482}{436.473}{and}%
\TextAt{8}{277.957}{436.473}{it}%
\TextAt{8}{289.372}{436.473}{was}%
\TextAt{8}{311.584}{436.473}{C.}%
\TextAt{8}{326.897}{436.473}{Deninger}%
\TextAt{8}{376.521}{436.473}{and}%
\TextAt{8}{398.995}{436.473}{U.}%
\TextAt{8}{414.631}{436.473}{Jannsen}%
\TextAt{8}{459.438}{436.473}{who}%
\TextAt{8}{483.861}{436.473}{suggested}%
\TextAt{8}{79.226}{419.427}{that}%
\TextAt{8}{106.5}{419.427}{I}%
\TextAt{8}{116.52}{419.427}{give}%
\TextAt{8}{142.49}{419.427}{a}%
\TextAt{8}{154.15}{419.427}{new}%
\TextAt{8}{180.109}{419.427}{treatment}%
\TextAt{8}{236.326}{419.427}{of}%
\TextAt{8}{251.56}{419.427}{Hecke’s}%
\TextAt{8}{295.141}{419.427}{theory}%
\TextAt{8}{333.796}{419.427}{of}%
\TextAt{8}{349.03}{419.427}{theta}%
\TextAt{8}{381.504}{419.427}{series}%
\TextAt{8}{414.743}{419.427}{and}%
\TextAt{8}{439.411}{419.427}{L-series.}%
\TextAt{8}{492.224}{419.427}{I}%
\TextAt{8}{502.244}{419.427}{owe}%
\TextAt{8}{526.911}{419.427}{a}%
\TextAt{8}{79.226}{402.381}{great}%
\TextAt{8}{110.105}{402.381}{debt}%
\TextAt{8}{137.721}{402.381}{of}%
\TextAt{8}{152.006}{402.381}{gratitude}%
\TextAt{8}{203.7}{402.381}{to}%
\TextAt{8}{218.965}{402.381}{Mrs.}%
\TextAt{8}{250.174}{402.381}{Eva-Maria}%
\TextAt{8}{308.526}{402.381}{Strobel.}%
\TextAt{8}{356.591}{402.381}{She}%
\TextAt{8}{379.651}{402.381}{drew}%
\TextAt{8}{409.216}{402.381}{the}%
\TextAt{8}{430.327}{402.381}{pictures}%
\TextAt{8}{475.565}{402.381}{and}%
\TextAt{8}{499.283}{402.381}{helped}%
\TextAt{8}{79.226}{385.335}{me}%
\TextAt{8}{98.713}{385.335}{with}%
\TextAt{8}{126.006}{385.335}{the}%
\TextAt{8}{146.796}{385.335}{proofreading}%
\TextAt{8}{215.73}{385.335}{and}%
\TextAt{8}{239.126}{385.335}{the}%
\TextAt{8}{259.915}{385.335}{formatting}%
\TextAt{8}{318.771}{385.335}{of}%
\TextAt{8}{332.734}{385.335}{the}%
\TextAt{8}{353.524}{385.335}{text,}%
\TextAt{8}{381.956}{385.335}{never}%
\TextAt{8}{413.804}{385.335}{tiring}%
\TextAt{8}{446.31}{385.335}{of}%
\TextAt{8}{460.274}{385.335}{going}%
\TextAt{8}{492.134}{385.335}{into}%
\TextAt{8}{516.51}{385.335}{the}%
\TextAt{8}{79.226}{368.289}{minutest}%
\TextAt{8}{128.779}{368.289}{detail.}%
\TextAt{8}{168.965}{368.289}{Let}%
\TextAt{8}{190.974}{368.289}{me}%
\TextAt{8}{210.866}{368.289}{heartily}%
\TextAt{8}{255.159}{368.289}{thank}%
\TextAt{8}{289.697}{368.289}{all}%
\TextAt{8}{306.995}{368.289}{those}%
\TextAt{8}{338.663}{368.289}{who}%
\TextAt{8}{364.413}{368.289}{assisted}%
\TextAt{8}{408.563}{368.289}{me,}%
\TextAt{8}{431.967}{368.289}{and}%
\TextAt{8}{455.768}{368.289}{also}%
\TextAt{8}{480.287}{368.289}{those}%
\TextAt{8}{511.956}{368.289}{who}%
\TextAt{8}{79.226}{351.243}{are}%
\TextAt{8}{99.88}{351.243}{not}%
\TextAt{8}{121.838}{351.243}{named}%
\TextAt{8}{160.7}{351.243}{here.}%
\TextAt{8}{194.038}{351.243}{Tremendous}%
\TextAt{8}{261.246}{351.243}{thanks}%
\TextAt{8}{300.503}{351.243}{are}%
\TextAt{8}{321.157}{351.243}{due}%
\TextAt{8}{344.406}{351.243}{to}%
\TextAt{8}{359.859}{351.243}{Mrs.}%
\TextAt{8}{391.632}{351.243}{Martina}%
\TextAt{8}{437.978}{351.243}{Hertl}%
\TextAt{8}{469.356}{351.243}{who}%
\TextAt{8}{495.211}{351.243}{did}%
\TextAt{8}{516.51}{351.243}{the}%
\TextAt{8}{79.226}{334.197}{typesetting}%
\TextAt{8}{139.787}{334.197}{of}%
\TextAt{8}{152.801}{334.197}{the}%
\TextAt{8}{172.643}{334.197}{manuscript}%
\TextAt{8}{233.203}{334.197}{in}%
\TextAt{8}{246.541}{334.197}{T}%
\TextAt{8}{253}{331.621}{E}%
\TextAt{8}{259.468}{334.197}{X.}%
\TextAt{8}{276.586}{334.197}{That}%
\TextAt{8}{305.537}{334.197}{the}%
\TextAt{8}{325.378}{334.197}{book}%
\TextAt{8}{354.006}{334.197}{can}%
\TextAt{8}{375.151}{334.197}{appear}%
\TextAt{8}{413.534}{334.197}{is}%
\TextAt{8}{424.983}{334.197}{essentially}%
\TextAt{8}{480.725}{334.197}{due}%
\TextAt{8}{502.515}{334.197}{to}%
\TextAt{8}{516.51}{334.197}{her}%
\TextAt{8}{79.226}{317.151}{competence,}%
\TextAt{8}{145.251}{317.151}{to}%
\TextAt{8}{158.779}{317.151}{the}%
\TextAt{8}{178.154}{317.151}{unfailing}%
\TextAt{8}{225.827}{317.151}{and}%
\TextAt{8}{247.807}{317.151}{kind}%
\TextAt{8}{273.363}{317.151}{willingness}%
\TextAt{8}{331.233}{317.151}{with}%
\TextAt{8}{357.112}{317.151}{which}%
\TextAt{8}{389.816}{317.151}{she}%
\TextAt{8}{409.251}{317.151}{worked}%
\TextAt{8}{448.471}{317.151}{through}%
\TextAt{8}{491.923}{317.151}{the}%
\TextAt{8}{511.298}{317.151}{long}%
\TextAt{8}{79.226}{300.105}{handwritten}%
\TextAt{8}{144.785}{300.105}{manuscript,}%
\TextAt{8}{208.289}{300.105}{and}%
\TextAt{8}{230.272}{300.105}{through}%
\TextAt{8}{273.726}{300.105}{the}%
\TextAt{8}{293.102}{300.105}{many}%
\TextAt{8}{324.195}{300.105}{modifications,}%
\TextAt{8}{398.757}{300.105}{additions,}%
\TextAt{8}{452.183}{300.105}{and}%
\TextAt{8}{474.165}{300.105}{corrections,}%
\TextAt{8}{79.226}{283.059}{always}%
\TextAt{8}{116.694}{283.059}{prepared}%
\TextAt{8}{165.471}{283.059}{to}%
\TextAt{8}{179.781}{283.059}{give}%
\TextAt{8}{203.847}{283.059}{her}%
\TextAt{8}{224.003}{283.059}{best.}%
\TextAt{8}{392.858}{250.073}{Regensburg,}%
\TextAt{8}{459.425}{250.073}{February}%
\TextAt{8}{509.337}{250.073}{1992}%
\TextAt{8}{448.39}{233.027}{Jürgen}%
\TextAt{8}{486.921}{233.027}{Neukirch}%
\Vector{q 0 G 0 g 0.399 w 7.226 -663.202 m 188.641 -663.202 l S Q}%
\ErrataAt{79.19}{99.64}{220.81}{}%
\TextAt{11}{92.693}{90.655}{*}%
\TextAt{10}{97.159}{87.039}{Error}%
\TextAt{10}{123.958}{87.039}{is}%
\TextAt{10}{133.971}{87.039}{ever}%
\TextAt{10}{155.022}{87.039}{with}%
\TextAt{10}{177.717}{87.039}{us.}%
\TextAt{10}{194.374}{87.039}{Yet}%
\TextAt{10}{212.636}{87.039}{some}%
\TextAt{10}{237.582}{87.039}{angelic}%
\TextAt{10}{270.788}{87.039}{need}%
\TextAt{10}{79.226}{72.932}{Gently}%
\TextAt{10}{111.953}{72.932}{coaxes}%
\TextAt{10}{143.266}{72.932}{our}%
\TextAt{10}{161.009}{72.932}{striving}%
\TextAt{10}{197.353}{72.932}{mind}%
\TextAt{10}{222.818}{72.932}{upwards,}%
\TextAt{10}{264.97}{72.932}{towards}%
\TextAt{10}{302.13}{72.932}{truth.}%
\TextAt{10}{332.058}{72.932}{(Translation}%
\TextAt{10}{388.885}{72.932}{suggested}%
\TextAt{10}{433.816}{72.932}{by}%
\TextAt{10}{447.655}{72.932}{Barry}%
\TextAt{10}{476.078}{72.932}{Mazur.)}%
\TextAt{8}{301.281}{43.033}{ix}%
\Destination{97.16}{101.145}{nav.Hfootnote.2}%
\Destination{78.23}{809.815}{nav.page.ix}%
\Hotlink{391.579}{556.7}{4.232}{11.294}{nav.Hfootnote.2}%
}

% Contents
% Page 10 (body page front matter)
\Page{%
\TextAt{1}{79.226}{673.054}{Contents}%
\TextAt{12}{79.226}{604.8}{Original}%
\TextAt{12}{131.279}{604.8}{Publication}%
\TextAt{12}{203.082}{604.8}{Data}%
\TextAt{12}{237.726}{604.8}{.}%
\TextAt{12}{247.446}{604.8}{.}%
\TextAt{12}{257.165}{604.8}{.}%
\TextAt{12}{266.885}{604.8}{.}%
\TextAt{12}{276.604}{604.8}{.}%
\TextAt{12}{286.323}{604.8}{.}%
\TextAt{12}{296.043}{604.8}{.}%
\TextAt{12}{305.762}{604.8}{.}%
\TextAt{12}{315.482}{604.8}{.}%
\TextAt{12}{325.201}{604.8}{.}%
\TextAt{12}{334.921}{604.8}{.}%
\TextAt{12}{344.64}{604.8}{.}%
\TextAt{12}{354.36}{604.8}{.}%
\TextAt{12}{364.079}{604.8}{.}%
\TextAt{12}{373.799}{604.8}{.}%
\TextAt{12}{383.518}{604.8}{.}%
\TextAt{12}{393.238}{604.8}{.}%
\TextAt{12}{402.957}{604.8}{.}%
\TextAt{12}{412.676}{604.8}{.}%
\TextAt{12}{422.396}{604.8}{.}%
\TextAt{12}{432.115}{604.8}{.}%
\TextAt{12}{441.835}{604.8}{.}%
\TextAt{12}{451.554}{604.8}{.}%
\TextAt{12}{461.274}{604.8}{.}%
\TextAt{12}{470.993}{604.8}{.}%
\TextAt{12}{480.713}{604.8}{.}%
\TextAt{12}{490.432}{604.8}{.}%
\TextAt{12}{500.152}{604.8}{.}%
\TextAt{8}{529.518}{604.8}{i}%
\TextAt{12}{79.226}{547.994}{Foreword}%
\TextAt{12}{140.532}{547.994}{.}%
\TextAt{12}{150.251}{547.994}{.}%
\TextAt{12}{159.971}{547.994}{.}%
\TextAt{12}{169.69}{547.994}{.}%
\TextAt{12}{179.409}{547.994}{.}%
\TextAt{12}{189.129}{547.994}{.}%
\TextAt{12}{198.848}{547.994}{.}%
\TextAt{12}{208.568}{547.994}{.}%
\TextAt{12}{218.287}{547.994}{.}%
\TextAt{12}{228.007}{547.994}{.}%
\TextAt{12}{237.726}{547.994}{.}%
\TextAt{12}{247.446}{547.994}{.}%
\TextAt{12}{257.165}{547.994}{.}%
\TextAt{12}{266.885}{547.994}{.}%
\TextAt{12}{276.604}{547.994}{.}%
\TextAt{12}{286.323}{547.994}{.}%
\TextAt{12}{296.043}{547.994}{.}%
\TextAt{12}{305.762}{547.994}{.}%
\TextAt{12}{315.482}{547.994}{.}%
\TextAt{12}{325.201}{547.994}{.}%
\TextAt{12}{334.921}{547.994}{.}%
\TextAt{12}{344.64}{547.994}{.}%
\TextAt{12}{354.36}{547.994}{.}%
\TextAt{12}{364.079}{547.994}{.}%
\TextAt{12}{373.799}{547.994}{.}%
\TextAt{12}{383.518}{547.994}{.}%
\TextAt{12}{393.238}{547.994}{.}%
\TextAt{12}{402.957}{547.994}{.}%
\TextAt{12}{412.676}{547.994}{.}%
\TextAt{12}{422.396}{547.994}{.}%
\TextAt{12}{432.115}{547.994}{.}%
\TextAt{12}{441.835}{547.994}{.}%
\TextAt{12}{451.554}{547.994}{.}%
\TextAt{12}{461.274}{547.994}{.}%
\TextAt{12}{470.993}{547.994}{.}%
\TextAt{12}{480.713}{547.994}{.}%
\TextAt{12}{490.432}{547.994}{.}%
\TextAt{12}{500.152}{547.994}{.}%
\TextAt{8}{523.014}{547.994}{iii}%
\TextAt{12}{79.226}{519.59}{Translator’s}%
\TextAt{12}{153.551}{519.59}{Note}%
\TextAt{12}{189.129}{519.59}{.}%
\TextAt{12}{198.848}{519.59}{.}%
\TextAt{12}{208.568}{519.59}{.}%
\TextAt{12}{218.287}{519.59}{.}%
\TextAt{12}{228.007}{519.59}{.}%
\TextAt{12}{237.726}{519.59}{.}%
\TextAt{12}{247.446}{519.59}{.}%
\TextAt{12}{257.165}{519.59}{.}%
\TextAt{12}{266.885}{519.59}{.}%
\TextAt{12}{276.604}{519.59}{.}%
\TextAt{12}{286.323}{519.59}{.}%
\TextAt{12}{296.043}{519.59}{.}%
\TextAt{12}{305.762}{519.59}{.}%
\TextAt{12}{315.482}{519.59}{.}%
\TextAt{12}{325.201}{519.59}{.}%
\TextAt{12}{334.921}{519.59}{.}%
\TextAt{12}{344.64}{519.59}{.}%
\TextAt{12}{354.36}{519.59}{.}%
\TextAt{12}{364.079}{519.59}{.}%
\TextAt{12}{373.799}{519.59}{.}%
\TextAt{12}{383.518}{519.59}{.}%
\TextAt{12}{393.238}{519.59}{.}%
\TextAt{12}{402.957}{519.59}{.}%
\TextAt{12}{412.676}{519.59}{.}%
\TextAt{12}{422.396}{519.59}{.}%
\TextAt{12}{432.115}{519.59}{.}%
\TextAt{12}{441.835}{519.59}{.}%
\TextAt{12}{451.554}{519.59}{.}%
\TextAt{12}{461.274}{519.59}{.}%
\TextAt{12}{470.993}{519.59}{.}%
\TextAt{12}{480.713}{519.59}{.}%
\TextAt{12}{490.432}{519.59}{.}%
\TextAt{12}{500.152}{519.59}{.}%
\TextAt{8}{523.337}{519.59}{iv}%
\TextAt{12}{79.226}{491.187}{Preface}%
\TextAt{12}{127.298}{491.187}{to}%
\TextAt{12}{143.748}{491.187}{the}%
\TextAt{12}{167.073}{491.187}{German}%
\TextAt{12}{218.982}{491.187}{Edition}%
\TextAt{12}{266.885}{491.187}{.}%
\TextAt{12}{276.604}{491.187}{.}%
\TextAt{12}{286.323}{491.187}{.}%
\TextAt{12}{296.043}{491.187}{.}%
\TextAt{12}{305.762}{491.187}{.}%
\TextAt{12}{315.482}{491.187}{.}%
\TextAt{12}{325.201}{491.187}{.}%
\TextAt{12}{334.921}{491.187}{.}%
\TextAt{12}{344.64}{491.187}{.}%
\TextAt{12}{354.36}{491.187}{.}%
\TextAt{12}{364.079}{491.187}{.}%
\TextAt{12}{373.799}{491.187}{.}%
\TextAt{12}{383.518}{491.187}{.}%
\TextAt{12}{393.238}{491.187}{.}%
\TextAt{12}{402.957}{491.187}{.}%
\TextAt{12}{412.676}{491.187}{.}%
\TextAt{12}{422.396}{491.187}{.}%
\TextAt{12}{432.115}{491.187}{.}%
\TextAt{12}{441.835}{491.187}{.}%
\TextAt{12}{451.554}{491.187}{.}%
\TextAt{12}{461.274}{491.187}{.}%
\TextAt{12}{470.993}{491.187}{.}%
\TextAt{12}{480.713}{491.187}{.}%
\TextAt{12}{490.432}{491.187}{.}%
\TextAt{12}{500.152}{491.187}{.}%
\TextAt{8}{523.337}{491.187}{vi}%
\TextAt{12}{79.226}{462.784}{I}%
\TextAt{12}{109.114}{462.784}{Algebraic}%
\TextAt{12}{169.595}{462.784}{Integers}%
\TextAt{12}{228.007}{462.784}{.}%
\TextAt{12}{237.726}{462.784}{.}%
\TextAt{12}{247.446}{462.784}{.}%
\TextAt{12}{257.165}{462.784}{.}%
\TextAt{12}{266.885}{462.784}{.}%
\TextAt{12}{276.604}{462.784}{.}%
\TextAt{12}{286.323}{462.784}{.}%
\TextAt{12}{296.043}{462.784}{.}%
\TextAt{12}{305.762}{462.784}{.}%
\TextAt{12}{315.482}{462.784}{.}%
\TextAt{12}{325.201}{462.784}{.}%
\TextAt{12}{334.921}{462.784}{.}%
\TextAt{12}{344.64}{462.784}{.}%
\TextAt{12}{354.36}{462.784}{.}%
\TextAt{12}{364.079}{462.784}{.}%
\TextAt{12}{373.799}{462.784}{.}%
\TextAt{12}{383.518}{462.784}{.}%
\TextAt{12}{393.238}{462.784}{.}%
\TextAt{12}{402.957}{462.784}{.}%
\TextAt{12}{412.676}{462.784}{.}%
\TextAt{12}{422.396}{462.784}{.}%
\TextAt{12}{432.115}{462.784}{.}%
\TextAt{12}{441.835}{462.784}{.}%
\TextAt{12}{451.554}{462.784}{.}%
\TextAt{12}{461.274}{462.784}{.}%
\TextAt{12}{470.993}{462.784}{.}%
\TextAt{12}{480.713}{462.784}{.}%
\TextAt{12}{490.432}{462.784}{.}%
\TextAt{12}{500.152}{462.784}{.}%
\TextAt{8}{526.911}{462.784}{1}%
\TextAt{8}{97.159}{442.749}{§}%
\TextAt{8}{104.831}{442.749}{1}%
\TextAt{8}{124.656}{442.749}{The}%
\TextAt{8}{148.71}{442.749}{Gaussian}%
\TextAt{8}{198.993}{442.749}{Integers}%
\TextAt{8}{248.342}{442.749}{.}%
\TextAt{8}{257.572}{442.749}{.}%
\TextAt{8}{266.801}{442.749}{.}%
\TextAt{8}{276.03}{442.749}{.}%
\TextAt{8}{285.26}{442.749}{.}%
\TextAt{8}{294.489}{442.749}{.}%
\TextAt{8}{303.718}{442.749}{.}%
\TextAt{8}{312.948}{442.749}{.}%
\TextAt{8}{322.177}{442.749}{.}%
\TextAt{8}{331.406}{442.749}{.}%
\TextAt{8}{340.635}{442.749}{.}%
\TextAt{8}{349.865}{442.749}{.}%
\TextAt{8}{359.094}{442.749}{.}%
\TextAt{8}{368.323}{442.749}{.}%
\TextAt{8}{377.553}{442.749}{.}%
\TextAt{8}{386.782}{442.749}{.}%
\TextAt{8}{396.011}{442.749}{.}%
\TextAt{8}{405.241}{442.749}{.}%
\TextAt{8}{414.47}{442.749}{.}%
\TextAt{8}{423.699}{442.749}{.}%
\TextAt{8}{432.929}{442.749}{.}%
\TextAt{8}{442.158}{442.749}{.}%
\TextAt{8}{451.387}{442.749}{.}%
\TextAt{8}{460.616}{442.749}{.}%
\TextAt{8}{469.846}{442.749}{.}%
\TextAt{8}{479.075}{442.749}{.}%
\TextAt{8}{488.304}{442.749}{.}%
\TextAt{8}{497.534}{442.749}{.}%
\TextAt{8}{506.763}{442.749}{.}%
\TextAt{8}{526.911}{442.749}{1}%
\TextAt{8}{97.159}{422.714}{§}%
\TextAt{8}{104.831}{422.714}{2}%
\TextAt{8}{124.656}{422.714}{Integrality}%
\TextAt{8}{183.737}{422.714}{.}%
\TextAt{8}{192.967}{422.714}{.}%
\TextAt{8}{202.196}{422.714}{.}%
\TextAt{8}{211.425}{422.714}{.}%
\TextAt{8}{220.654}{422.714}{.}%
\TextAt{8}{229.884}{422.714}{.}%
\TextAt{8}{239.113}{422.714}{.}%
\TextAt{8}{248.342}{422.714}{.}%
\TextAt{8}{257.572}{422.714}{.}%
\TextAt{8}{266.801}{422.714}{.}%
\TextAt{8}{276.03}{422.714}{.}%
\TextAt{8}{285.26}{422.714}{.}%
\TextAt{8}{294.489}{422.714}{.}%
\TextAt{8}{303.718}{422.714}{.}%
\TextAt{8}{312.948}{422.714}{.}%
\TextAt{8}{322.177}{422.714}{.}%
\TextAt{8}{331.406}{422.714}{.}%
\TextAt{8}{340.635}{422.714}{.}%
\TextAt{8}{349.865}{422.714}{.}%
\TextAt{8}{359.094}{422.714}{.}%
\TextAt{8}{368.323}{422.714}{.}%
\TextAt{8}{377.553}{422.714}{.}%
\TextAt{8}{386.782}{422.714}{.}%
\TextAt{8}{396.011}{422.714}{.}%
\TextAt{8}{405.241}{422.714}{.}%
\TextAt{8}{414.47}{422.714}{.}%
\TextAt{8}{423.699}{422.714}{.}%
\TextAt{8}{432.929}{422.714}{.}%
\TextAt{8}{442.158}{422.714}{.}%
\TextAt{8}{451.387}{422.714}{.}%
\TextAt{8}{460.616}{422.714}{.}%
\TextAt{8}{469.846}{422.714}{.}%
\TextAt{8}{479.075}{422.714}{.}%
\TextAt{8}{488.304}{422.714}{.}%
\TextAt{8}{497.534}{422.714}{.}%
\TextAt{8}{506.763}{422.714}{.}%
\TextAt{8}{526.911}{422.714}{5}%
\TextAt{8}{97.159}{402.679}{§}%
\TextAt{8}{104.831}{402.679}{3}%
\TextAt{8}{124.656}{402.679}{Ideals}%
\TextAt{8}{165.279}{402.679}{.}%
\TextAt{8}{174.508}{402.679}{.}%
\TextAt{8}{183.737}{402.679}{.}%
\TextAt{8}{192.967}{402.679}{.}%
\TextAt{8}{202.196}{402.679}{.}%
\TextAt{8}{211.425}{402.679}{.}%
\TextAt{8}{220.654}{402.679}{.}%
\TextAt{8}{229.884}{402.679}{.}%
\TextAt{8}{239.113}{402.679}{.}%
\TextAt{8}{248.342}{402.679}{.}%
\TextAt{8}{257.572}{402.679}{.}%
\TextAt{8}{266.801}{402.679}{.}%
\TextAt{8}{276.03}{402.679}{.}%
\TextAt{8}{285.26}{402.679}{.}%
\TextAt{8}{294.489}{402.679}{.}%
\TextAt{8}{303.718}{402.679}{.}%
\TextAt{8}{312.948}{402.679}{.}%
\TextAt{8}{322.177}{402.679}{.}%
\TextAt{8}{331.406}{402.679}{.}%
\TextAt{8}{340.635}{402.679}{.}%
\TextAt{8}{349.865}{402.679}{.}%
\TextAt{8}{359.094}{402.679}{.}%
\TextAt{8}{368.323}{402.679}{.}%
\TextAt{8}{377.553}{402.679}{.}%
\TextAt{8}{386.782}{402.679}{.}%
\TextAt{8}{396.011}{402.679}{.}%
\TextAt{8}{405.241}{402.679}{.}%
\TextAt{8}{414.47}{402.679}{.}%
\TextAt{8}{423.699}{402.679}{.}%
\TextAt{8}{432.929}{402.679}{.}%
\TextAt{8}{442.158}{402.679}{.}%
\TextAt{8}{451.387}{402.679}{.}%
\TextAt{8}{460.616}{402.679}{.}%
\TextAt{8}{469.846}{402.679}{.}%
\TextAt{8}{479.075}{402.679}{.}%
\TextAt{8}{488.304}{402.679}{.}%
\TextAt{8}{497.534}{402.679}{.}%
\TextAt{8}{506.763}{402.679}{.}%
\TextAt{8}{521.053}{402.679}{15}%
\TextAt{8}{97.159}{382.644}{§}%
\TextAt{8}{104.831}{382.644}{4}%
\TextAt{8}{124.656}{382.644}{Lattices}%
\TextAt{8}{174.508}{382.644}{.}%
\TextAt{8}{183.737}{382.644}{.}%
\TextAt{8}{192.967}{382.644}{.}%
\TextAt{8}{202.196}{382.644}{.}%
\TextAt{8}{211.425}{382.644}{.}%
\TextAt{8}{220.654}{382.644}{.}%
\TextAt{8}{229.884}{382.644}{.}%
\TextAt{8}{239.113}{382.644}{.}%
\TextAt{8}{248.342}{382.644}{.}%
\TextAt{8}{257.572}{382.644}{.}%
\TextAt{8}{266.801}{382.644}{.}%
\TextAt{8}{276.03}{382.644}{.}%
\TextAt{8}{285.26}{382.644}{.}%
\TextAt{8}{294.489}{382.644}{.}%
\TextAt{8}{303.718}{382.644}{.}%
\TextAt{8}{312.948}{382.644}{.}%
\TextAt{8}{322.177}{382.644}{.}%
\TextAt{8}{331.406}{382.644}{.}%
\TextAt{8}{340.635}{382.644}{.}%
\TextAt{8}{349.865}{382.644}{.}%
\TextAt{8}{359.094}{382.644}{.}%
\TextAt{8}{368.323}{382.644}{.}%
\TextAt{8}{377.553}{382.644}{.}%
\TextAt{8}{386.782}{382.644}{.}%
\TextAt{8}{396.011}{382.644}{.}%
\TextAt{8}{405.241}{382.644}{.}%
\TextAt{8}{414.47}{382.644}{.}%
\TextAt{8}{423.699}{382.644}{.}%
\TextAt{8}{432.929}{382.644}{.}%
\TextAt{8}{442.158}{382.644}{.}%
\TextAt{8}{451.387}{382.644}{.}%
\TextAt{8}{460.616}{382.644}{.}%
\TextAt{8}{469.846}{382.644}{.}%
\TextAt{8}{479.075}{382.644}{.}%
\TextAt{8}{488.304}{382.644}{.}%
\TextAt{8}{497.534}{382.644}{.}%
\TextAt{8}{506.763}{382.644}{.}%
\TextAt{8}{521.053}{382.644}{22}%
\TextAt{8}{97.159}{362.61}{§}%
\TextAt{8}{104.831}{362.61}{5}%
\TextAt{8}{124.656}{362.61}{Minkowski}%
\TextAt{8}{182.925}{362.61}{Theory}%
\TextAt{8}{229.884}{362.61}{.}%
\TextAt{8}{239.113}{362.61}{.}%
\TextAt{8}{248.342}{362.61}{.}%
\TextAt{8}{257.572}{362.61}{.}%
\TextAt{8}{266.801}{362.61}{.}%
\TextAt{8}{276.03}{362.61}{.}%
\TextAt{8}{285.26}{362.61}{.}%
\TextAt{8}{294.489}{362.61}{.}%
\TextAt{8}{303.718}{362.61}{.}%
\TextAt{8}{312.948}{362.61}{.}%
\TextAt{8}{322.177}{362.61}{.}%
\TextAt{8}{331.406}{362.61}{.}%
\TextAt{8}{340.635}{362.61}{.}%
\TextAt{8}{349.865}{362.61}{.}%
\TextAt{8}{359.094}{362.61}{.}%
\TextAt{8}{368.323}{362.61}{.}%
\TextAt{8}{377.553}{362.61}{.}%
\TextAt{8}{386.782}{362.61}{.}%
\TextAt{8}{396.011}{362.61}{.}%
\TextAt{8}{405.241}{362.61}{.}%
\TextAt{8}{414.47}{362.61}{.}%
\TextAt{8}{423.699}{362.61}{.}%
\TextAt{8}{432.929}{362.61}{.}%
\TextAt{8}{442.158}{362.61}{.}%
\TextAt{8}{451.387}{362.61}{.}%
\TextAt{8}{460.616}{362.61}{.}%
\TextAt{8}{469.846}{362.61}{.}%
\TextAt{8}{479.075}{362.61}{.}%
\TextAt{8}{488.304}{362.61}{.}%
\TextAt{8}{497.534}{362.61}{.}%
\TextAt{8}{506.763}{362.61}{.}%
\TextAt{8}{521.053}{362.61}{26}%
\TextAt{8}{97.159}{342.575}{§}%
\TextAt{8}{104.831}{342.575}{6}%
\TextAt{8}{124.656}{342.575}{The}%
\TextAt{8}{148.71}{342.575}{Class}%
\TextAt{8}{179.399}{342.575}{Number}%
\TextAt{8}{229.884}{342.575}{.}%
\TextAt{8}{239.113}{342.575}{.}%
\TextAt{8}{248.342}{342.575}{.}%
\TextAt{8}{257.572}{342.575}{.}%
\TextAt{8}{266.801}{342.575}{.}%
\TextAt{8}{276.03}{342.575}{.}%
\TextAt{8}{285.26}{342.575}{.}%
\TextAt{8}{294.489}{342.575}{.}%
\TextAt{8}{303.718}{342.575}{.}%
\TextAt{8}{312.948}{342.575}{.}%
\TextAt{8}{322.177}{342.575}{.}%
\TextAt{8}{331.406}{342.575}{.}%
\TextAt{8}{340.635}{342.575}{.}%
\TextAt{8}{349.865}{342.575}{.}%
\TextAt{8}{359.094}{342.575}{.}%
\TextAt{8}{368.323}{342.575}{.}%
\TextAt{8}{377.553}{342.575}{.}%
\TextAt{8}{386.782}{342.575}{.}%
\TextAt{8}{396.011}{342.575}{.}%
\TextAt{8}{405.241}{342.575}{.}%
\TextAt{8}{414.47}{342.575}{.}%
\TextAt{8}{423.699}{342.575}{.}%
\TextAt{8}{432.929}{342.575}{.}%
\TextAt{8}{442.158}{342.575}{.}%
\TextAt{8}{451.387}{342.575}{.}%
\TextAt{8}{460.616}{342.575}{.}%
\TextAt{8}{469.846}{342.575}{.}%
\TextAt{8}{479.075}{342.575}{.}%
\TextAt{8}{488.304}{342.575}{.}%
\TextAt{8}{497.534}{342.575}{.}%
\TextAt{8}{506.763}{342.575}{.}%
\TextAt{8}{521.053}{342.575}{32}%
\TextAt{8}{97.159}{322.54}{§}%
\TextAt{8}{104.831}{322.54}{7}%
\TextAt{8}{124.656}{322.54}{Dirichlet’s}%
\TextAt{8}{180.809}{322.54}{Unit}%
\TextAt{8}{207.792}{322.54}{Theorem}%
\TextAt{8}{257.572}{322.54}{.}%
\TextAt{8}{266.801}{322.54}{.}%
\TextAt{8}{276.03}{322.54}{.}%
\TextAt{8}{285.26}{322.54}{.}%
\TextAt{8}{294.489}{322.54}{.}%
\TextAt{8}{303.718}{322.54}{.}%
\TextAt{8}{312.948}{322.54}{.}%
\TextAt{8}{322.177}{322.54}{.}%
\TextAt{8}{331.406}{322.54}{.}%
\TextAt{8}{340.635}{322.54}{.}%
\TextAt{8}{349.865}{322.54}{.}%
\TextAt{8}{359.094}{322.54}{.}%
\TextAt{8}{368.323}{322.54}{.}%
\TextAt{8}{377.553}{322.54}{.}%
\TextAt{8}{386.782}{322.54}{.}%
\TextAt{8}{396.011}{322.54}{.}%
\TextAt{8}{405.241}{322.54}{.}%
\TextAt{8}{414.47}{322.54}{.}%
\TextAt{8}{423.699}{322.54}{.}%
\TextAt{8}{432.929}{322.54}{.}%
\TextAt{8}{442.158}{322.54}{.}%
\TextAt{8}{451.387}{322.54}{.}%
\TextAt{8}{460.616}{322.54}{.}%
\TextAt{8}{469.846}{322.54}{.}%
\TextAt{8}{479.075}{322.54}{.}%
\TextAt{8}{488.304}{322.54}{.}%
\TextAt{8}{497.534}{322.54}{.}%
\TextAt{8}{506.763}{322.54}{.}%
\TextAt{8}{521.053}{322.54}{36}%
\TextAt{8}{97.159}{302.505}{§}%
\TextAt{8}{104.831}{302.505}{8}%
\TextAt{8}{124.656}{302.505}{Extensions}%
\TextAt{8}{183.798}{302.505}{of}%
\TextAt{8}{197.128}{302.505}{Dedekind}%
\TextAt{8}{249.312}{302.505}{Domains}%
\TextAt{8}{303.718}{302.505}{.}%
\TextAt{8}{312.948}{302.505}{.}%
\TextAt{8}{322.177}{302.505}{.}%
\TextAt{8}{331.406}{302.505}{.}%
\TextAt{8}{340.635}{302.505}{.}%
\TextAt{8}{349.865}{302.505}{.}%
\TextAt{8}{359.094}{302.505}{.}%
\TextAt{8}{368.323}{302.505}{.}%
\TextAt{8}{377.553}{302.505}{.}%
\TextAt{8}{386.782}{302.505}{.}%
\TextAt{8}{396.011}{302.505}{.}%
\TextAt{8}{405.241}{302.505}{.}%
\TextAt{8}{414.47}{302.505}{.}%
\TextAt{8}{423.699}{302.505}{.}%
\TextAt{8}{432.929}{302.505}{.}%
\TextAt{8}{442.158}{302.505}{.}%
\TextAt{8}{451.387}{302.505}{.}%
\TextAt{8}{460.616}{302.505}{.}%
\TextAt{8}{469.846}{302.505}{.}%
\TextAt{8}{479.075}{302.505}{.}%
\TextAt{8}{488.304}{302.505}{.}%
\TextAt{8}{497.534}{302.505}{.}%
\TextAt{8}{506.763}{302.505}{.}%
\TextAt{8}{521.053}{302.505}{42}%
\TextAt{8}{97.159}{282.471}{§}%
\TextAt{8}{104.831}{282.471}{9}%
\TextAt{8}{124.656}{282.471}{Hilbert’s}%
\TextAt{8}{172.835}{282.471}{Ramification}%
\TextAt{8}{241.948}{282.471}{Theory}%
\TextAt{8}{285.26}{282.471}{.}%
\TextAt{8}{294.489}{282.471}{.}%
\TextAt{8}{303.718}{282.471}{.}%
\TextAt{8}{312.948}{282.471}{.}%
\TextAt{8}{322.177}{282.471}{.}%
\TextAt{8}{331.406}{282.471}{.}%
\TextAt{8}{340.635}{282.471}{.}%
\TextAt{8}{349.865}{282.471}{.}%
\TextAt{8}{359.094}{282.471}{.}%
\TextAt{8}{368.323}{282.471}{.}%
\TextAt{8}{377.553}{282.471}{.}%
\TextAt{8}{386.782}{282.471}{.}%
\TextAt{8}{396.011}{282.471}{.}%
\TextAt{8}{405.241}{282.471}{.}%
\TextAt{8}{414.47}{282.471}{.}%
\TextAt{8}{423.699}{282.471}{.}%
\TextAt{8}{432.929}{282.471}{.}%
\TextAt{8}{442.158}{282.471}{.}%
\TextAt{8}{451.387}{282.471}{.}%
\TextAt{8}{460.616}{282.471}{.}%
\TextAt{8}{469.846}{282.471}{.}%
\TextAt{8}{479.075}{282.471}{.}%
\TextAt{8}{488.304}{282.471}{.}%
\TextAt{8}{497.534}{282.471}{.}%
\TextAt{8}{506.763}{282.471}{.}%
\TextAt{8}{521.053}{282.471}{50}%
\TextAt{8}{97.159}{262.436}{§}%
\TextAt{8}{104.831}{262.436}{10}%
\TextAt{8}{124.656}{262.436}{Cyclotomic}%
\TextAt{8}{186.117}{262.436}{Fields}%
\TextAt{8}{220.654}{262.436}{.}%
\TextAt{8}{229.884}{262.436}{.}%
\TextAt{8}{239.113}{262.436}{.}%
\TextAt{8}{248.342}{262.436}{.}%
\TextAt{8}{257.572}{262.436}{.}%
\TextAt{8}{266.801}{262.436}{.}%
\TextAt{8}{276.03}{262.436}{.}%
\TextAt{8}{285.26}{262.436}{.}%
\TextAt{8}{294.489}{262.436}{.}%
\TextAt{8}{303.718}{262.436}{.}%
\TextAt{8}{312.948}{262.436}{.}%
\TextAt{8}{322.177}{262.436}{.}%
\TextAt{8}{331.406}{262.436}{.}%
\TextAt{8}{340.635}{262.436}{.}%
\TextAt{8}{349.865}{262.436}{.}%
\TextAt{8}{359.094}{262.436}{.}%
\TextAt{8}{368.323}{262.436}{.}%
\TextAt{8}{377.553}{262.436}{.}%
\TextAt{8}{386.782}{262.436}{.}%
\TextAt{8}{396.011}{262.436}{.}%
\TextAt{8}{405.241}{262.436}{.}%
\TextAt{8}{414.47}{262.436}{.}%
\TextAt{8}{423.699}{262.436}{.}%
\TextAt{8}{432.929}{262.436}{.}%
\TextAt{8}{442.158}{262.436}{.}%
\TextAt{8}{451.387}{262.436}{.}%
\TextAt{8}{460.616}{262.436}{.}%
\TextAt{8}{469.846}{262.436}{.}%
\TextAt{8}{479.075}{262.436}{.}%
\TextAt{8}{488.304}{262.436}{.}%
\TextAt{8}{497.534}{262.436}{.}%
\TextAt{8}{506.763}{262.436}{.}%
\TextAt{8}{521.053}{262.436}{55}%
\TextAt{8}{97.159}{242.401}{§}%
\TextAt{8}{104.831}{242.401}{11}%
\TextAt{8}{124.656}{242.401}{Localization}%
\TextAt{8}{192.967}{242.401}{.}%
\TextAt{8}{202.196}{242.401}{.}%
\TextAt{8}{211.425}{242.401}{.}%
\TextAt{8}{220.654}{242.401}{.}%
\TextAt{8}{229.884}{242.401}{.}%
\TextAt{8}{239.113}{242.401}{.}%
\TextAt{8}{248.342}{242.401}{.}%
\TextAt{8}{257.572}{242.401}{.}%
\TextAt{8}{266.801}{242.401}{.}%
\TextAt{8}{276.03}{242.401}{.}%
\TextAt{8}{285.26}{242.401}{.}%
\TextAt{8}{294.489}{242.401}{.}%
\TextAt{8}{303.718}{242.401}{.}%
\TextAt{8}{312.948}{242.401}{.}%
\TextAt{8}{322.177}{242.401}{.}%
\TextAt{8}{331.406}{242.401}{.}%
\TextAt{8}{340.635}{242.401}{.}%
\TextAt{8}{349.865}{242.401}{.}%
\TextAt{8}{359.094}{242.401}{.}%
\TextAt{8}{368.323}{242.401}{.}%
\TextAt{8}{377.553}{242.401}{.}%
\TextAt{8}{386.782}{242.401}{.}%
\TextAt{8}{396.011}{242.401}{.}%
\TextAt{8}{405.241}{242.401}{.}%
\TextAt{8}{414.47}{242.401}{.}%
\TextAt{8}{423.699}{242.401}{.}%
\TextAt{8}{432.929}{242.401}{.}%
\TextAt{8}{442.158}{242.401}{.}%
\TextAt{8}{451.387}{242.401}{.}%
\TextAt{8}{460.616}{242.401}{.}%
\TextAt{8}{469.846}{242.401}{.}%
\TextAt{8}{479.075}{242.401}{.}%
\TextAt{8}{488.304}{242.401}{.}%
\TextAt{8}{497.534}{242.401}{.}%
\TextAt{8}{506.763}{242.401}{.}%
\TextAt{8}{521.053}{242.401}{60}%
\TextAt{8}{97.159}{222.366}{§}%
\TextAt{8}{104.831}{222.366}{12}%
\TextAt{8}{124.656}{222.366}{Orders}%
\TextAt{8}{165.279}{222.366}{.}%
\TextAt{8}{174.508}{222.366}{.}%
\TextAt{8}{183.737}{222.366}{.}%
\TextAt{8}{192.967}{222.366}{.}%
\TextAt{8}{202.196}{222.366}{.}%
\TextAt{8}{211.425}{222.366}{.}%
\TextAt{8}{220.654}{222.366}{.}%
\TextAt{8}{229.884}{222.366}{.}%
\TextAt{8}{239.113}{222.366}{.}%
\TextAt{8}{248.342}{222.366}{.}%
\TextAt{8}{257.572}{222.366}{.}%
\TextAt{8}{266.801}{222.366}{.}%
\TextAt{8}{276.03}{222.366}{.}%
\TextAt{8}{285.26}{222.366}{.}%
\TextAt{8}{294.489}{222.366}{.}%
\TextAt{8}{303.718}{222.366}{.}%
\TextAt{8}{312.948}{222.366}{.}%
\TextAt{8}{322.177}{222.366}{.}%
\TextAt{8}{331.406}{222.366}{.}%
\TextAt{8}{340.635}{222.366}{.}%
\TextAt{8}{349.865}{222.366}{.}%
\TextAt{8}{359.094}{222.366}{.}%
\TextAt{8}{368.323}{222.366}{.}%
\TextAt{8}{377.553}{222.366}{.}%
\TextAt{8}{386.782}{222.366}{.}%
\TextAt{8}{396.011}{222.366}{.}%
\TextAt{8}{405.241}{222.366}{.}%
\TextAt{8}{414.47}{222.366}{.}%
\TextAt{8}{423.699}{222.366}{.}%
\TextAt{8}{432.929}{222.366}{.}%
\TextAt{8}{442.158}{222.366}{.}%
\TextAt{8}{451.387}{222.366}{.}%
\TextAt{8}{460.616}{222.366}{.}%
\TextAt{8}{469.846}{222.366}{.}%
\TextAt{8}{479.075}{222.366}{.}%
\TextAt{8}{488.304}{222.366}{.}%
\TextAt{8}{497.534}{222.366}{.}%
\TextAt{8}{506.763}{222.366}{.}%
\TextAt{8}{521.053}{222.366}{68}%
\TextAt{8}{97.159}{202.332}{§}%
\TextAt{8}{104.831}{202.332}{13}%
\TextAt{8}{124.656}{202.332}{One-dimensional}%
\TextAt{8}{213.817}{202.332}{Schemes}%
\TextAt{8}{266.801}{202.332}{.}%
\TextAt{8}{276.03}{202.332}{.}%
\TextAt{8}{285.26}{202.332}{.}%
\TextAt{8}{294.489}{202.332}{.}%
\TextAt{8}{303.718}{202.332}{.}%
\TextAt{8}{312.948}{202.332}{.}%
\TextAt{8}{322.177}{202.332}{.}%
\TextAt{8}{331.406}{202.332}{.}%
\TextAt{8}{340.635}{202.332}{.}%
\TextAt{8}{349.865}{202.332}{.}%
\TextAt{8}{359.094}{202.332}{.}%
\TextAt{8}{368.323}{202.332}{.}%
\TextAt{8}{377.553}{202.332}{.}%
\TextAt{8}{386.782}{202.332}{.}%
\TextAt{8}{396.011}{202.332}{.}%
\TextAt{8}{405.241}{202.332}{.}%
\TextAt{8}{414.47}{202.332}{.}%
\TextAt{8}{423.699}{202.332}{.}%
\TextAt{8}{432.929}{202.332}{.}%
\TextAt{8}{442.158}{202.332}{.}%
\TextAt{8}{451.387}{202.332}{.}%
\TextAt{8}{460.616}{202.332}{.}%
\TextAt{8}{469.846}{202.332}{.}%
\TextAt{8}{479.075}{202.332}{.}%
\TextAt{8}{488.304}{202.332}{.}%
\TextAt{8}{497.534}{202.332}{.}%
\TextAt{8}{506.763}{202.332}{.}%
\TextAt{8}{521.053}{202.332}{79}%
\TextAt{8}{97.159}{182.297}{§}%
\TextAt{8}{104.831}{182.297}{14}%
\TextAt{8}{124.656}{182.297}{Function}%
\TextAt{8}{173.588}{182.297}{Fields}%
\TextAt{8}{211.425}{182.297}{.}%
\TextAt{8}{220.654}{182.297}{.}%
\TextAt{8}{229.884}{182.297}{.}%
\TextAt{8}{239.113}{182.297}{.}%
\TextAt{8}{248.342}{182.297}{.}%
\TextAt{8}{257.572}{182.297}{.}%
\TextAt{8}{266.801}{182.297}{.}%
\TextAt{8}{276.03}{182.297}{.}%
\TextAt{8}{285.26}{182.297}{.}%
\TextAt{8}{294.489}{182.297}{.}%
\TextAt{8}{303.718}{182.297}{.}%
\TextAt{8}{312.948}{182.297}{.}%
\TextAt{8}{322.177}{182.297}{.}%
\TextAt{8}{331.406}{182.297}{.}%
\TextAt{8}{340.635}{182.297}{.}%
\TextAt{8}{349.865}{182.297}{.}%
\TextAt{8}{359.094}{182.297}{.}%
\TextAt{8}{368.323}{182.297}{.}%
\TextAt{8}{377.553}{182.297}{.}%
\TextAt{8}{386.782}{182.297}{.}%
\TextAt{8}{396.011}{182.297}{.}%
\TextAt{8}{405.241}{182.297}{.}%
\TextAt{8}{414.47}{182.297}{.}%
\TextAt{8}{423.699}{182.297}{.}%
\TextAt{8}{432.929}{182.297}{.}%
\TextAt{8}{442.158}{182.297}{.}%
\TextAt{8}{451.387}{182.297}{.}%
\TextAt{8}{460.616}{182.297}{.}%
\TextAt{8}{469.846}{182.297}{.}%
\TextAt{8}{479.075}{182.297}{.}%
\TextAt{8}{488.304}{182.297}{.}%
\TextAt{8}{497.534}{182.297}{.}%
\TextAt{8}{506.763}{182.297}{.}%
\TextAt{8}{521.053}{182.297}{87}%
\TextAt{12}{79.226}{153.894}{II}%
\TextAt{12}{109.114}{153.894}{The}%
\TextAt{12}{136.551}{153.894}{Theory}%
\TextAt{12}{183.32}{153.894}{of}%
\TextAt{12}{198.646}{153.894}{Valuations}%
\TextAt{12}{266.885}{153.894}{.}%
\TextAt{12}{276.604}{153.894}{.}%
\TextAt{12}{286.323}{153.894}{.}%
\TextAt{12}{296.043}{153.894}{.}%
\TextAt{12}{305.762}{153.894}{.}%
\TextAt{12}{315.482}{153.894}{.}%
\TextAt{12}{325.201}{153.894}{.}%
\TextAt{12}{334.921}{153.894}{.}%
\TextAt{12}{344.64}{153.894}{.}%
\TextAt{12}{354.36}{153.894}{.}%
\TextAt{12}{364.079}{153.894}{.}%
\TextAt{12}{373.799}{153.894}{.}%
\TextAt{12}{383.518}{153.894}{.}%
\TextAt{12}{393.238}{153.894}{.}%
\TextAt{12}{402.957}{153.894}{.}%
\TextAt{12}{412.676}{153.894}{.}%
\TextAt{12}{422.396}{153.894}{.}%
\TextAt{12}{432.115}{153.894}{.}%
\TextAt{12}{441.835}{153.894}{.}%
\TextAt{12}{451.554}{153.894}{.}%
\TextAt{12}{461.274}{153.894}{.}%
\TextAt{12}{470.993}{153.894}{.}%
\TextAt{12}{480.713}{153.894}{.}%
\TextAt{12}{490.432}{153.894}{.}%
\TextAt{12}{500.152}{153.894}{.}%
\TextAt{8}{521.053}{153.894}{91}%
\TextAt{8}{97.159}{133.859}{§}%
\TextAt{8}{104.831}{133.859}{1}%
\TextAt{8}{124.656}{133.859}{The}%
\TextAt{13}{148.71}{133.859}{p}%
\TextAt{8}{154.585}{133.859}{-adic}%
\TextAt{8}{183.193}{133.859}{Numbers}%
\TextAt{8}{239.113}{133.859}{.}%
\TextAt{8}{248.342}{133.859}{.}%
\TextAt{8}{257.572}{133.859}{.}%
\TextAt{8}{266.801}{133.859}{.}%
\TextAt{8}{276.03}{133.859}{.}%
\TextAt{8}{285.26}{133.859}{.}%
\TextAt{8}{294.489}{133.859}{.}%
\TextAt{8}{303.718}{133.859}{.}%
\TextAt{8}{312.948}{133.859}{.}%
\TextAt{8}{322.177}{133.859}{.}%
\TextAt{8}{331.406}{133.859}{.}%
\TextAt{8}{340.635}{133.859}{.}%
\TextAt{8}{349.865}{133.859}{.}%
\TextAt{8}{359.094}{133.859}{.}%
\TextAt{8}{368.323}{133.859}{.}%
\TextAt{8}{377.553}{133.859}{.}%
\TextAt{8}{386.782}{133.859}{.}%
\TextAt{8}{396.011}{133.859}{.}%
\TextAt{8}{405.241}{133.859}{.}%
\TextAt{8}{414.47}{133.859}{.}%
\TextAt{8}{423.699}{133.859}{.}%
\TextAt{8}{432.929}{133.859}{.}%
\TextAt{8}{442.158}{133.859}{.}%
\TextAt{8}{451.387}{133.859}{.}%
\TextAt{8}{460.616}{133.859}{.}%
\TextAt{8}{469.846}{133.859}{.}%
\TextAt{8}{479.075}{133.859}{.}%
\TextAt{8}{488.304}{133.859}{.}%
\TextAt{8}{497.534}{133.859}{.}%
\TextAt{8}{506.763}{133.859}{.}%
\TextAt{8}{521.053}{133.859}{91}%
\TextAt{8}{97.159}{113.824}{§}%
\TextAt{8}{104.831}{113.824}{2}%
\TextAt{8}{124.656}{113.824}{The}%
\TextAt{13}{148.71}{113.824}{p}%
\TextAt{8}{154.585}{113.824}{-adic}%
\TextAt{8}{183.193}{113.824}{Absolute}%
\TextAt{8}{232.353}{113.824}{Value}%
\TextAt{8}{266.801}{113.824}{.}%
\TextAt{8}{276.03}{113.824}{.}%
\TextAt{8}{285.26}{113.824}{.}%
\TextAt{8}{294.489}{113.824}{.}%
\TextAt{8}{303.718}{113.824}{.}%
\TextAt{8}{312.948}{113.824}{.}%
\TextAt{8}{322.177}{113.824}{.}%
\TextAt{8}{331.406}{113.824}{.}%
\TextAt{8}{340.635}{113.824}{.}%
\TextAt{8}{349.865}{113.824}{.}%
\TextAt{8}{359.094}{113.824}{.}%
\TextAt{8}{368.323}{113.824}{.}%
\TextAt{8}{377.553}{113.824}{.}%
\TextAt{8}{386.782}{113.824}{.}%
\TextAt{8}{396.011}{113.824}{.}%
\TextAt{8}{405.241}{113.824}{.}%
\TextAt{8}{414.47}{113.824}{.}%
\TextAt{8}{423.699}{113.824}{.}%
\TextAt{8}{432.929}{113.824}{.}%
\TextAt{8}{442.158}{113.824}{.}%
\TextAt{8}{451.387}{113.824}{.}%
\TextAt{8}{460.616}{113.824}{.}%
\TextAt{8}{469.846}{113.824}{.}%
\TextAt{8}{479.075}{113.824}{.}%
\TextAt{8}{488.304}{113.824}{.}%
\TextAt{8}{497.534}{113.824}{.}%
\TextAt{8}{506.763}{113.824}{.}%
\TextAt{8}{521.053}{113.824}{98}%
\TextAt{8}{97.159}{93.789}{§}%
\TextAt{8}{104.831}{93.789}{3}%
\TextAt{8}{124.656}{93.789}{Valuations}%
\TextAt{8}{183.737}{93.789}{.}%
\TextAt{8}{192.967}{93.789}{.}%
\TextAt{8}{202.196}{93.789}{.}%
\TextAt{8}{211.425}{93.789}{.}%
\TextAt{8}{220.654}{93.789}{.}%
\TextAt{8}{229.884}{93.789}{.}%
\TextAt{8}{239.113}{93.789}{.}%
\TextAt{8}{248.342}{93.789}{.}%
\TextAt{8}{257.572}{93.789}{.}%
\TextAt{8}{266.801}{93.789}{.}%
\TextAt{8}{276.03}{93.789}{.}%
\TextAt{8}{285.26}{93.789}{.}%
\TextAt{8}{294.489}{93.789}{.}%
\TextAt{8}{303.718}{93.789}{.}%
\TextAt{8}{312.948}{93.789}{.}%
\TextAt{8}{322.177}{93.789}{.}%
\TextAt{8}{331.406}{93.789}{.}%
\TextAt{8}{340.635}{93.789}{.}%
\TextAt{8}{349.865}{93.789}{.}%
\TextAt{8}{359.094}{93.789}{.}%
\TextAt{8}{368.323}{93.789}{.}%
\TextAt{8}{377.553}{93.789}{.}%
\TextAt{8}{386.782}{93.789}{.}%
\TextAt{8}{396.011}{93.789}{.}%
\TextAt{8}{405.241}{93.789}{.}%
\TextAt{8}{414.47}{93.789}{.}%
\TextAt{8}{423.699}{93.789}{.}%
\TextAt{8}{432.929}{93.789}{.}%
\TextAt{8}{442.158}{93.789}{.}%
\TextAt{8}{451.387}{93.789}{.}%
\TextAt{8}{460.616}{93.789}{.}%
\TextAt{8}{469.846}{93.789}{.}%
\TextAt{8}{479.075}{93.789}{.}%
\TextAt{8}{488.304}{93.789}{.}%
\TextAt{8}{497.534}{93.789}{.}%
\TextAt{8}{506.763}{93.789}{.}%
\TextAt{8}{515.195}{93.789}{106}%
\TextAt{8}{97.159}{73.754}{§}%
\TextAt{8}{104.831}{73.754}{4}%
\TextAt{8}{124.656}{73.754}{Completions}%
\TextAt{8}{192.967}{73.754}{.}%
\TextAt{8}{202.196}{73.754}{.}%
\TextAt{8}{211.425}{73.754}{.}%
\TextAt{8}{220.654}{73.754}{.}%
\TextAt{8}{229.884}{73.754}{.}%
\TextAt{8}{239.113}{73.754}{.}%
\TextAt{8}{248.342}{73.754}{.}%
\TextAt{8}{257.572}{73.754}{.}%
\TextAt{8}{266.801}{73.754}{.}%
\TextAt{8}{276.03}{73.754}{.}%
\TextAt{8}{285.26}{73.754}{.}%
\TextAt{8}{294.489}{73.754}{.}%
\TextAt{8}{303.718}{73.754}{.}%
\TextAt{8}{312.948}{73.754}{.}%
\TextAt{8}{322.177}{73.754}{.}%
\TextAt{8}{331.406}{73.754}{.}%
\TextAt{8}{340.635}{73.754}{.}%
\TextAt{8}{349.865}{73.754}{.}%
\TextAt{8}{359.094}{73.754}{.}%
\TextAt{8}{368.323}{73.754}{.}%
\TextAt{8}{377.553}{73.754}{.}%
\TextAt{8}{386.782}{73.754}{.}%
\TextAt{8}{396.011}{73.754}{.}%
\TextAt{8}{405.241}{73.754}{.}%
\TextAt{8}{414.47}{73.754}{.}%
\TextAt{8}{423.699}{73.754}{.}%
\TextAt{8}{432.929}{73.754}{.}%
\TextAt{8}{442.158}{73.754}{.}%
\TextAt{8}{451.387}{73.754}{.}%
\TextAt{8}{460.616}{73.754}{.}%
\TextAt{8}{469.846}{73.754}{.}%
\TextAt{8}{479.075}{73.754}{.}%
\TextAt{8}{488.304}{73.754}{.}%
\TextAt{8}{497.534}{73.754}{.}%
\TextAt{8}{506.763}{73.754}{.}%
\TextAt{8}{515.195}{73.754}{113}%
\TextAt{8}{302.907}{43.033}{x}%
\Destination{79.23}{773.075}{nav.chapterz2az.6}%
\Destination{78.23}{809.815}{nav.page.x}%
\Hotlink{79.226}{601.345}{152.476}{17.048}{nav.chapterz2az.1}%
\Hotlink{79.226}{544.538}{54.743}{17.048}{nav.chapterz2az.3}%
\Hotlink{79.226}{516.135}{102.946}{17.048}{nav.chapterz2az.4}%
\Hotlink{79.226}{487.732}{182.974}{17.048}{nav.chapterz2az.5}%
\Hotlink{79.226}{459.328}{137.52}{17.048}{nav.chapter.1}%
\Hotlink{97.159}{439.401}{142.219}{16.821}{nav.section.1.1}%
\Hotlink{97.159}{419.366}{80.841}{16.821}{nav.section.1.2}%
\Hotlink{97.159}{399.332}{57.146}{16.82}{nav.section.1.3}%
\Hotlink{97.159}{379.297}{68.049}{16.821}{nav.section.1.4}%
\Hotlink{97.159}{359.262}{122.517}{16.821}{nav.section.1.5}%
\Hotlink{97.159}{339.227}{123.533}{16.821}{nav.section.1.6}%
\Hotlink{97.159}{319.193}{156.158}{16.82}{nav.section.1.7}%
\Hotlink{97.159}{299.158}{196.937}{16.821}{nav.section.1.8}%
\Hotlink{97.159}{279.123}{181.539}{16.821}{nav.section.1.9}%
\Hotlink{97.159}{259.088}{119.42}{16.821}{nav.section.1.10}%
\Hotlink{97.159}{239.053}{89.783}{16.821}{nav.section.1.11}%
\Hotlink{97.159}{219.019}{62.035}{16.821}{nav.section.1.12}%
\Hotlink{97.159}{198.984}{159.315}{16.821}{nav.section.1.13}%
\Hotlink{97.159}{178.949}{106.891}{16.821}{nav.section.1.14}%
\Hotlink{79.226}{150.438}{181.24}{17.048}{nav.chapter.2}%
\Hotlink{97.159}{130.511}{131.942}{16.821}{nav.section.2.1}%
\Hotlink{97.159}{110.476}{163.803}{16.821}{nav.section.2.2}%
\Hotlink{97.159}{90.442}{81.546}{16.82}{nav.section.2.3}%
\Hotlink{97.159}{70.407}{91.302}{16.821}{nav.section.2.4}%
}

% Page 11 (body page front matter)
\Page{%
\TextAt{2}{79.226}{796.943}{Contents}%
\Vector{q 0 G 0 g 0.299 w 7.226 18.729 m 460.769 18.729 l S Q}%
\TextAt{8}{97.159}{761.124}{§}%
\TextAt{8}{104.831}{761.124}{5}%
\TextAt{8}{124.656}{761.124}{Local}%
\TextAt{8}{156.361}{761.124}{Fields}%
\TextAt{8}{192.967}{761.124}{.}%
\TextAt{8}{202.196}{761.124}{.}%
\TextAt{8}{211.425}{761.124}{.}%
\TextAt{8}{220.654}{761.124}{.}%
\TextAt{8}{229.884}{761.124}{.}%
\TextAt{8}{239.113}{761.124}{.}%
\TextAt{8}{248.342}{761.124}{.}%
\TextAt{8}{257.572}{761.124}{.}%
\TextAt{8}{266.801}{761.124}{.}%
\TextAt{8}{276.03}{761.124}{.}%
\TextAt{8}{285.26}{761.124}{.}%
\TextAt{8}{294.489}{761.124}{.}%
\TextAt{8}{303.718}{761.124}{.}%
\TextAt{8}{312.948}{761.124}{.}%
\TextAt{8}{322.177}{761.124}{.}%
\TextAt{8}{331.406}{761.124}{.}%
\TextAt{8}{340.635}{761.124}{.}%
\TextAt{8}{349.865}{761.124}{.}%
\TextAt{8}{359.094}{761.124}{.}%
\TextAt{8}{368.323}{761.124}{.}%
\TextAt{8}{377.553}{761.124}{.}%
\TextAt{8}{386.782}{761.124}{.}%
\TextAt{8}{396.011}{761.124}{.}%
\TextAt{8}{405.241}{761.124}{.}%
\TextAt{8}{414.47}{761.124}{.}%
\TextAt{8}{423.699}{761.124}{.}%
\TextAt{8}{432.929}{761.124}{.}%
\TextAt{8}{442.158}{761.124}{.}%
\TextAt{8}{451.387}{761.124}{.}%
\TextAt{8}{460.616}{761.124}{.}%
\TextAt{8}{469.846}{761.124}{.}%
\TextAt{8}{479.075}{761.124}{.}%
\TextAt{8}{488.304}{761.124}{.}%
\TextAt{8}{497.534}{761.124}{.}%
\TextAt{8}{506.763}{761.124}{.}%
\TextAt{8}{515.195}{761.124}{123}%
\TextAt{8}{97.159}{741.089}{§}%
\TextAt{8}{104.831}{741.089}{6}%
\TextAt{8}{124.656}{741.089}{Henselian}%
\TextAt{8}{177.713}{741.089}{Fields}%
\TextAt{8}{211.425}{741.089}{.}%
\TextAt{8}{220.654}{741.089}{.}%
\TextAt{8}{229.884}{741.089}{.}%
\TextAt{8}{239.113}{741.089}{.}%
\TextAt{8}{248.342}{741.089}{.}%
\TextAt{8}{257.572}{741.089}{.}%
\TextAt{8}{266.801}{741.089}{.}%
\TextAt{8}{276.03}{741.089}{.}%
\TextAt{8}{285.26}{741.089}{.}%
\TextAt{8}{294.489}{741.089}{.}%
\TextAt{8}{303.718}{741.089}{.}%
\TextAt{8}{312.948}{741.089}{.}%
\TextAt{8}{322.177}{741.089}{.}%
\TextAt{8}{331.406}{741.089}{.}%
\TextAt{8}{340.635}{741.089}{.}%
\TextAt{8}{349.865}{741.089}{.}%
\TextAt{8}{359.094}{741.089}{.}%
\TextAt{8}{368.323}{741.089}{.}%
\TextAt{8}{377.553}{741.089}{.}%
\TextAt{8}{386.782}{741.089}{.}%
\TextAt{8}{396.011}{741.089}{.}%
\TextAt{8}{405.241}{741.089}{.}%
\TextAt{8}{414.47}{741.089}{.}%
\TextAt{8}{423.699}{741.089}{.}%
\TextAt{8}{432.929}{741.089}{.}%
\TextAt{8}{442.158}{741.089}{.}%
\TextAt{8}{451.387}{741.089}{.}%
\TextAt{8}{460.616}{741.089}{.}%
\TextAt{8}{469.846}{741.089}{.}%
\TextAt{8}{479.075}{741.089}{.}%
\TextAt{8}{488.304}{741.089}{.}%
\TextAt{8}{497.534}{741.089}{.}%
\TextAt{8}{506.763}{741.089}{.}%
\TextAt{8}{515.195}{741.089}{131}%
\TextAt{8}{97.159}{721.054}{§}%
\TextAt{8}{104.831}{721.054}{7}%
\TextAt{8}{124.656}{721.054}{Unramified}%
\TextAt{8}{185.46}{721.054}{and}%
\TextAt{8}{208.222}{721.054}{Tamely}%
\TextAt{8}{249.838}{721.054}{Ramified}%
\TextAt{8}{299.428}{721.054}{Extensions}%
\TextAt{8}{359.094}{721.054}{.}%
\TextAt{8}{368.323}{721.054}{.}%
\TextAt{8}{377.553}{721.054}{.}%
\TextAt{8}{386.782}{721.054}{.}%
\TextAt{8}{396.011}{721.054}{.}%
\TextAt{8}{405.241}{721.054}{.}%
\TextAt{8}{414.47}{721.054}{.}%
\TextAt{8}{423.699}{721.054}{.}%
\TextAt{8}{432.929}{721.054}{.}%
\TextAt{8}{442.158}{721.054}{.}%
\TextAt{8}{451.387}{721.054}{.}%
\TextAt{8}{460.616}{721.054}{.}%
\TextAt{8}{469.846}{721.054}{.}%
\TextAt{8}{479.075}{721.054}{.}%
\TextAt{8}{488.304}{721.054}{.}%
\TextAt{8}{497.534}{721.054}{.}%
\TextAt{8}{506.763}{721.054}{.}%
\TextAt{8}{515.195}{721.054}{140}%
\TextAt{8}{97.159}{701.019}{§}%
\TextAt{8}{104.831}{701.019}{8}%
\TextAt{8}{124.656}{701.019}{Extensions}%
\TextAt{8}{183.798}{701.019}{of}%
\TextAt{8}{197.128}{701.019}{Valuations}%
\TextAt{8}{257.572}{701.019}{.}%
\TextAt{8}{266.801}{701.019}{.}%
\TextAt{8}{276.03}{701.019}{.}%
\TextAt{8}{285.26}{701.019}{.}%
\TextAt{8}{294.489}{701.019}{.}%
\TextAt{8}{303.718}{701.019}{.}%
\TextAt{8}{312.948}{701.019}{.}%
\TextAt{8}{322.177}{701.019}{.}%
\TextAt{8}{331.406}{701.019}{.}%
\TextAt{8}{340.635}{701.019}{.}%
\TextAt{8}{349.865}{701.019}{.}%
\TextAt{8}{359.094}{701.019}{.}%
\TextAt{8}{368.323}{701.019}{.}%
\TextAt{8}{377.553}{701.019}{.}%
\TextAt{8}{386.782}{701.019}{.}%
\TextAt{8}{396.011}{701.019}{.}%
\TextAt{8}{405.241}{701.019}{.}%
\TextAt{8}{414.47}{701.019}{.}%
\TextAt{8}{423.699}{701.019}{.}%
\TextAt{8}{432.929}{701.019}{.}%
\TextAt{8}{442.158}{701.019}{.}%
\TextAt{8}{451.387}{701.019}{.}%
\TextAt{8}{460.616}{701.019}{.}%
\TextAt{8}{469.846}{701.019}{.}%
\TextAt{8}{479.075}{701.019}{.}%
\TextAt{8}{488.304}{701.019}{.}%
\TextAt{8}{497.534}{701.019}{.}%
\TextAt{8}{506.763}{701.019}{.}%
\TextAt{8}{515.195}{701.019}{146}%
\TextAt{8}{97.159}{680.985}{§}%
\TextAt{8}{104.831}{680.985}{9}%
\TextAt{8}{124.656}{680.985}{Galois}%
\TextAt{8}{160.569}{680.985}{Theory}%
\TextAt{8}{201.217}{680.985}{of}%
\TextAt{8}{214.547}{680.985}{Valuations}%
\TextAt{8}{276.03}{680.985}{.}%
\TextAt{8}{285.26}{680.985}{.}%
\TextAt{8}{294.489}{680.985}{.}%
\TextAt{8}{303.718}{680.985}{.}%
\TextAt{8}{312.948}{680.985}{.}%
\TextAt{8}{322.177}{680.985}{.}%
\TextAt{8}{331.406}{680.985}{.}%
\TextAt{8}{340.635}{680.985}{.}%
\TextAt{8}{349.865}{680.985}{.}%
\TextAt{8}{359.094}{680.985}{.}%
\TextAt{8}{368.323}{680.985}{.}%
\TextAt{8}{377.553}{680.985}{.}%
\TextAt{8}{386.782}{680.985}{.}%
\TextAt{8}{396.011}{680.985}{.}%
\TextAt{8}{405.241}{680.985}{.}%
\TextAt{8}{414.47}{680.985}{.}%
\TextAt{8}{423.699}{680.985}{.}%
\TextAt{8}{432.929}{680.985}{.}%
\TextAt{8}{442.158}{680.985}{.}%
\TextAt{8}{451.387}{680.985}{.}%
\TextAt{8}{460.616}{680.985}{.}%
\TextAt{8}{469.846}{680.985}{.}%
\TextAt{8}{479.075}{680.985}{.}%
\TextAt{8}{488.304}{680.985}{.}%
\TextAt{8}{497.534}{680.985}{.}%
\TextAt{8}{506.763}{680.985}{.}%
\TextAt{8}{515.195}{680.985}{152}%
\TextAt{8}{97.159}{660.95}{§}%
\TextAt{8}{104.831}{660.95}{10}%
\TextAt{8}{124.656}{660.95}{Higher}%
\TextAt{8}{162.697}{660.95}{Ramification}%
\TextAt{8}{231.81}{660.95}{Groups}%
\TextAt{8}{276.03}{660.95}{.}%
\TextAt{8}{285.26}{660.95}{.}%
\TextAt{8}{294.489}{660.95}{.}%
\TextAt{8}{303.718}{660.95}{.}%
\TextAt{8}{312.948}{660.95}{.}%
\TextAt{8}{322.177}{660.95}{.}%
\TextAt{8}{331.406}{660.95}{.}%
\TextAt{8}{340.635}{660.95}{.}%
\TextAt{8}{349.865}{660.95}{.}%
\TextAt{8}{359.094}{660.95}{.}%
\TextAt{8}{368.323}{660.95}{.}%
\TextAt{8}{377.553}{660.95}{.}%
\TextAt{8}{386.782}{660.95}{.}%
\TextAt{8}{396.011}{660.95}{.}%
\TextAt{8}{405.241}{660.95}{.}%
\TextAt{8}{414.47}{660.95}{.}%
\TextAt{8}{423.699}{660.95}{.}%
\TextAt{8}{432.929}{660.95}{.}%
\TextAt{8}{442.158}{660.95}{.}%
\TextAt{8}{451.387}{660.95}{.}%
\TextAt{8}{460.616}{660.95}{.}%
\TextAt{8}{469.846}{660.95}{.}%
\TextAt{8}{479.075}{660.95}{.}%
\TextAt{8}{488.304}{660.95}{.}%
\TextAt{8}{497.534}{660.95}{.}%
\TextAt{8}{506.763}{660.95}{.}%
\TextAt{8}{515.195}{660.95}{161}%
\TextAt{12}{79.226}{632.547}{III}%
\TextAt{12}{109.114}{632.547}{Riemann–Roch}%
\TextAt{12}{203.141}{632.547}{Theory}%
\TextAt{12}{257.165}{632.547}{.}%
\TextAt{12}{266.885}{632.547}{.}%
\TextAt{12}{276.604}{632.547}{.}%
\TextAt{12}{286.323}{632.547}{.}%
\TextAt{12}{296.043}{632.547}{.}%
\TextAt{12}{305.762}{632.547}{.}%
\TextAt{12}{315.482}{632.547}{.}%
\TextAt{12}{325.201}{632.547}{.}%
\TextAt{12}{334.921}{632.547}{.}%
\TextAt{12}{344.64}{632.547}{.}%
\TextAt{12}{354.36}{632.547}{.}%
\TextAt{12}{364.079}{632.547}{.}%
\TextAt{12}{373.799}{632.547}{.}%
\TextAt{12}{383.518}{632.547}{.}%
\TextAt{12}{393.238}{632.547}{.}%
\TextAt{12}{402.957}{632.547}{.}%
\TextAt{12}{412.676}{632.547}{.}%
\TextAt{12}{422.396}{632.547}{.}%
\TextAt{12}{432.115}{632.547}{.}%
\TextAt{12}{441.835}{632.547}{.}%
\TextAt{12}{451.554}{632.547}{.}%
\TextAt{12}{461.274}{632.547}{.}%
\TextAt{12}{470.993}{632.547}{.}%
\TextAt{12}{480.713}{632.547}{.}%
\TextAt{12}{490.432}{632.547}{.}%
\TextAt{12}{500.152}{632.547}{.}%
\TextAt{8}{515.195}{632.547}{167}%
\TextAt{8}{97.159}{612.512}{§}%
\TextAt{8}{104.831}{612.512}{1}%
\TextAt{8}{124.656}{612.512}{Primes}%
\TextAt{8}{165.279}{612.512}{.}%
\TextAt{8}{174.508}{612.512}{.}%
\TextAt{8}{183.737}{612.512}{.}%
\TextAt{8}{192.967}{612.512}{.}%
\TextAt{8}{202.196}{612.512}{.}%
\TextAt{8}{211.425}{612.512}{.}%
\TextAt{8}{220.654}{612.512}{.}%
\TextAt{8}{229.884}{612.512}{.}%
\TextAt{8}{239.113}{612.512}{.}%
\TextAt{8}{248.342}{612.512}{.}%
\TextAt{8}{257.572}{612.512}{.}%
\TextAt{8}{266.801}{612.512}{.}%
\TextAt{8}{276.03}{612.512}{.}%
\TextAt{8}{285.26}{612.512}{.}%
\TextAt{8}{294.489}{612.512}{.}%
\TextAt{8}{303.718}{612.512}{.}%
\TextAt{8}{312.948}{612.512}{.}%
\TextAt{8}{322.177}{612.512}{.}%
\TextAt{8}{331.406}{612.512}{.}%
\TextAt{8}{340.635}{612.512}{.}%
\TextAt{8}{349.865}{612.512}{.}%
\TextAt{8}{359.094}{612.512}{.}%
\TextAt{8}{368.323}{612.512}{.}%
\TextAt{8}{377.553}{612.512}{.}%
\TextAt{8}{386.782}{612.512}{.}%
\TextAt{8}{396.011}{612.512}{.}%
\TextAt{8}{405.241}{612.512}{.}%
\TextAt{8}{414.47}{612.512}{.}%
\TextAt{8}{423.699}{612.512}{.}%
\TextAt{8}{432.929}{612.512}{.}%
\TextAt{8}{442.158}{612.512}{.}%
\TextAt{8}{451.387}{612.512}{.}%
\TextAt{8}{460.616}{612.512}{.}%
\TextAt{8}{469.846}{612.512}{.}%
\TextAt{8}{479.075}{612.512}{.}%
\TextAt{8}{488.304}{612.512}{.}%
\TextAt{8}{497.534}{612.512}{.}%
\TextAt{8}{506.763}{612.512}{.}%
\TextAt{8}{515.195}{612.512}{167}%
\TextAt{8}{97.159}{592.477}{§}%
\TextAt{8}{104.831}{592.477}{2}%
\TextAt{8}{124.656}{592.477}{Different}%
\TextAt{8}{173.266}{592.477}{and}%
\TextAt{8}{196.028}{592.477}{Discriminant}%
\TextAt{8}{266.801}{592.477}{.}%
\TextAt{8}{276.03}{592.477}{.}%
\TextAt{8}{285.26}{592.477}{.}%
\TextAt{8}{294.489}{592.477}{.}%
\TextAt{8}{303.718}{592.477}{.}%
\TextAt{8}{312.948}{592.477}{.}%
\TextAt{8}{322.177}{592.477}{.}%
\TextAt{8}{331.406}{592.477}{.}%
\TextAt{8}{340.635}{592.477}{.}%
\TextAt{8}{349.865}{592.477}{.}%
\TextAt{8}{359.094}{592.477}{.}%
\TextAt{8}{368.323}{592.477}{.}%
\TextAt{8}{377.553}{592.477}{.}%
\TextAt{8}{386.782}{592.477}{.}%
\TextAt{8}{396.011}{592.477}{.}%
\TextAt{8}{405.241}{592.477}{.}%
\TextAt{8}{414.47}{592.477}{.}%
\TextAt{8}{423.699}{592.477}{.}%
\TextAt{8}{432.929}{592.477}{.}%
\TextAt{8}{442.158}{592.477}{.}%
\TextAt{8}{451.387}{592.477}{.}%
\TextAt{8}{460.616}{592.477}{.}%
\TextAt{8}{469.846}{592.477}{.}%
\TextAt{8}{479.075}{592.477}{.}%
\TextAt{8}{488.304}{592.477}{.}%
\TextAt{8}{497.534}{592.477}{.}%
\TextAt{8}{506.763}{592.477}{.}%
\TextAt{8}{515.195}{592.477}{177}%
\TextAt{8}{97.159}{572.442}{§}%
\TextAt{8}{104.831}{572.442}{3}%
\TextAt{8}{124.656}{572.442}{Riemann–Roch}%
\TextAt{8}{211.425}{572.442}{.}%
\TextAt{8}{220.654}{572.442}{.}%
\TextAt{8}{229.884}{572.442}{.}%
\TextAt{8}{239.113}{572.442}{.}%
\TextAt{8}{248.342}{572.442}{.}%
\TextAt{8}{257.572}{572.442}{.}%
\TextAt{8}{266.801}{572.442}{.}%
\TextAt{8}{276.03}{572.442}{.}%
\TextAt{8}{285.26}{572.442}{.}%
\TextAt{8}{294.489}{572.442}{.}%
\TextAt{8}{303.718}{572.442}{.}%
\TextAt{8}{312.948}{572.442}{.}%
\TextAt{8}{322.177}{572.442}{.}%
\TextAt{8}{331.406}{572.442}{.}%
\TextAt{8}{340.635}{572.442}{.}%
\TextAt{8}{349.865}{572.442}{.}%
\TextAt{8}{359.094}{572.442}{.}%
\TextAt{8}{368.323}{572.442}{.}%
\TextAt{8}{377.553}{572.442}{.}%
\TextAt{8}{386.782}{572.442}{.}%
\TextAt{8}{396.011}{572.442}{.}%
\TextAt{8}{405.241}{572.442}{.}%
\TextAt{8}{414.47}{572.442}{.}%
\TextAt{8}{423.699}{572.442}{.}%
\TextAt{8}{432.929}{572.442}{.}%
\TextAt{8}{442.158}{572.442}{.}%
\TextAt{8}{451.387}{572.442}{.}%
\TextAt{8}{460.616}{572.442}{.}%
\TextAt{8}{469.846}{572.442}{.}%
\TextAt{8}{479.075}{572.442}{.}%
\TextAt{8}{488.304}{572.442}{.}%
\TextAt{8}{497.534}{572.442}{.}%
\TextAt{8}{506.763}{572.442}{.}%
\TextAt{8}{515.195}{572.442}{190}%
\TextAt{8}{97.159}{552.407}{§}%
\TextAt{8}{104.831}{552.407}{4}%
\TextAt{8}{124.656}{552.407}{Metrized}%
\TextAt{14}{173.744}{552.407}{o}%
\TextAt{8}{179.588}{552.407}{-Modules}%
\TextAt{8}{229.884}{552.407}{.}%
\TextAt{8}{239.113}{552.407}{.}%
\TextAt{8}{248.342}{552.407}{.}%
\TextAt{8}{257.572}{552.407}{.}%
\TextAt{8}{266.801}{552.407}{.}%
\TextAt{8}{276.03}{552.407}{.}%
\TextAt{8}{285.26}{552.407}{.}%
\TextAt{8}{294.489}{552.407}{.}%
\TextAt{8}{303.718}{552.407}{.}%
\TextAt{8}{312.948}{552.407}{.}%
\TextAt{8}{322.177}{552.407}{.}%
\TextAt{8}{331.406}{552.407}{.}%
\TextAt{8}{340.635}{552.407}{.}%
\TextAt{8}{349.865}{552.407}{.}%
\TextAt{8}{359.094}{552.407}{.}%
\TextAt{8}{368.323}{552.407}{.}%
\TextAt{8}{377.553}{552.407}{.}%
\TextAt{8}{386.782}{552.407}{.}%
\TextAt{8}{396.011}{552.407}{.}%
\TextAt{8}{405.241}{552.407}{.}%
\TextAt{8}{414.47}{552.407}{.}%
\TextAt{8}{423.699}{552.407}{.}%
\TextAt{8}{432.929}{552.407}{.}%
\TextAt{8}{442.158}{552.407}{.}%
\TextAt{8}{451.387}{552.407}{.}%
\TextAt{8}{460.616}{552.407}{.}%
\TextAt{8}{469.846}{552.407}{.}%
\TextAt{8}{479.075}{552.407}{.}%
\TextAt{8}{488.304}{552.407}{.}%
\TextAt{8}{497.534}{552.407}{.}%
\TextAt{8}{506.763}{552.407}{.}%
\TextAt{8}{515.195}{552.407}{206}%
\TextAt{8}{97.159}{532.373}{§}%
\TextAt{8}{104.831}{532.373}{5}%
\TextAt{8}{124.656}{532.373}{Grothendieck}%
\TextAt{8}{196.925}{532.373}{Groups}%
\TextAt{8}{239.113}{532.373}{.}%
\TextAt{8}{248.342}{532.373}{.}%
\TextAt{8}{257.572}{532.373}{.}%
\TextAt{8}{266.801}{532.373}{.}%
\TextAt{8}{276.03}{532.373}{.}%
\TextAt{8}{285.26}{532.373}{.}%
\TextAt{8}{294.489}{532.373}{.}%
\TextAt{8}{303.718}{532.373}{.}%
\TextAt{8}{312.948}{532.373}{.}%
\TextAt{8}{322.177}{532.373}{.}%
\TextAt{8}{331.406}{532.373}{.}%
\TextAt{8}{340.635}{532.373}{.}%
\TextAt{8}{349.865}{532.373}{.}%
\TextAt{8}{359.094}{532.373}{.}%
\TextAt{8}{368.323}{532.373}{.}%
\TextAt{8}{377.553}{532.373}{.}%
\TextAt{8}{386.782}{532.373}{.}%
\TextAt{8}{396.011}{532.373}{.}%
\TextAt{8}{405.241}{532.373}{.}%
\TextAt{8}{414.47}{532.373}{.}%
\TextAt{8}{423.699}{532.373}{.}%
\TextAt{8}{432.929}{532.373}{.}%
\TextAt{8}{442.158}{532.373}{.}%
\TextAt{8}{451.387}{532.373}{.}%
\TextAt{8}{460.616}{532.373}{.}%
\TextAt{8}{469.846}{532.373}{.}%
\TextAt{8}{479.075}{532.373}{.}%
\TextAt{8}{488.304}{532.373}{.}%
\TextAt{8}{497.534}{532.373}{.}%
\TextAt{8}{506.763}{532.373}{.}%
\TextAt{8}{515.195}{532.373}{214}%
\TextAt{8}{97.159}{512.338}{§}%
\TextAt{8}{104.831}{512.338}{6}%
\TextAt{8}{124.656}{512.338}{The}%
\TextAt{8}{148.71}{512.338}{Chern}%
\TextAt{8}{183.822}{512.338}{Character}%
\TextAt{8}{239.113}{512.338}{.}%
\TextAt{8}{248.342}{512.338}{.}%
\TextAt{8}{257.572}{512.338}{.}%
\TextAt{8}{266.801}{512.338}{.}%
\TextAt{8}{276.03}{512.338}{.}%
\TextAt{8}{285.26}{512.338}{.}%
\TextAt{8}{294.489}{512.338}{.}%
\TextAt{8}{303.718}{512.338}{.}%
\TextAt{8}{312.948}{512.338}{.}%
\TextAt{8}{322.177}{512.338}{.}%
\TextAt{8}{331.406}{512.338}{.}%
\TextAt{8}{340.635}{512.338}{.}%
\TextAt{8}{349.865}{512.338}{.}%
\TextAt{8}{359.094}{512.338}{.}%
\TextAt{8}{368.323}{512.338}{.}%
\TextAt{8}{377.553}{512.338}{.}%
\TextAt{8}{386.782}{512.338}{.}%
\TextAt{8}{396.011}{512.338}{.}%
\TextAt{8}{405.241}{512.338}{.}%
\TextAt{8}{414.47}{512.338}{.}%
\TextAt{8}{423.699}{512.338}{.}%
\TextAt{8}{432.929}{512.338}{.}%
\TextAt{8}{442.158}{512.338}{.}%
\TextAt{8}{451.387}{512.338}{.}%
\TextAt{8}{460.616}{512.338}{.}%
\TextAt{8}{469.846}{512.338}{.}%
\TextAt{8}{479.075}{512.338}{.}%
\TextAt{8}{488.304}{512.338}{.}%
\TextAt{8}{497.534}{512.338}{.}%
\TextAt{8}{506.763}{512.338}{.}%
\TextAt{8}{515.195}{512.338}{223}%
\TextAt{8}{97.159}{492.303}{§}%
\TextAt{8}{104.831}{492.303}{7}%
\TextAt{8}{124.656}{492.303}{Grothendieck–Riemann–Roch}%
\TextAt{8}{285.26}{492.303}{.}%
\TextAt{8}{294.489}{492.303}{.}%
\TextAt{8}{303.718}{492.303}{.}%
\TextAt{8}{312.948}{492.303}{.}%
\TextAt{8}{322.177}{492.303}{.}%
\TextAt{8}{331.406}{492.303}{.}%
\TextAt{8}{340.635}{492.303}{.}%
\TextAt{8}{349.865}{492.303}{.}%
\TextAt{8}{359.094}{492.303}{.}%
\TextAt{8}{368.323}{492.303}{.}%
\TextAt{8}{377.553}{492.303}{.}%
\TextAt{8}{386.782}{492.303}{.}%
\TextAt{8}{396.011}{492.303}{.}%
\TextAt{8}{405.241}{492.303}{.}%
\TextAt{8}{414.47}{492.303}{.}%
\TextAt{8}{423.699}{492.303}{.}%
\TextAt{8}{432.929}{492.303}{.}%
\TextAt{8}{442.158}{492.303}{.}%
\TextAt{8}{451.387}{492.303}{.}%
\TextAt{8}{460.616}{492.303}{.}%
\TextAt{8}{469.846}{492.303}{.}%
\TextAt{8}{479.075}{492.303}{.}%
\TextAt{8}{488.304}{492.303}{.}%
\TextAt{8}{497.534}{492.303}{.}%
\TextAt{8}{506.763}{492.303}{.}%
\TextAt{8}{515.195}{492.303}{226}%
\TextAt{8}{97.159}{472.268}{§}%
\TextAt{8}{104.831}{472.268}{8}%
\TextAt{8}{124.656}{472.268}{The}%
\TextAt{8}{148.71}{472.268}{Euler–Minkowski}%
\TextAt{8}{240.31}{472.268}{Characteristic}%
\TextAt{8}{322.177}{472.268}{.}%
\TextAt{8}{331.406}{472.268}{.}%
\TextAt{8}{340.635}{472.268}{.}%
\TextAt{8}{349.865}{472.268}{.}%
\TextAt{8}{359.094}{472.268}{.}%
\TextAt{8}{368.323}{472.268}{.}%
\TextAt{8}{377.553}{472.268}{.}%
\TextAt{8}{386.782}{472.268}{.}%
\TextAt{8}{396.011}{472.268}{.}%
\TextAt{8}{405.241}{472.268}{.}%
\TextAt{8}{414.47}{472.268}{.}%
\TextAt{8}{423.699}{472.268}{.}%
\TextAt{8}{432.929}{472.268}{.}%
\TextAt{8}{442.158}{472.268}{.}%
\TextAt{8}{451.387}{472.268}{.}%
\TextAt{8}{460.616}{472.268}{.}%
\TextAt{8}{469.846}{472.268}{.}%
\TextAt{8}{479.075}{472.268}{.}%
\TextAt{8}{488.304}{472.268}{.}%
\TextAt{8}{497.534}{472.268}{.}%
\TextAt{8}{506.763}{472.268}{.}%
\TextAt{8}{515.195}{472.268}{235}%
\TextAt{12}{79.226}{443.865}{IV}%
\TextAt{12}{109.114}{443.865}{Abstract}%
\TextAt{12}{165.028}{443.865}{Class}%
\TextAt{12}{200.129}{443.865}{Field}%
\TextAt{12}{234.153}{443.865}{Theory}%
\TextAt{12}{286.323}{443.865}{.}%
\TextAt{12}{296.043}{443.865}{.}%
\TextAt{12}{305.762}{443.865}{.}%
\TextAt{12}{315.482}{443.865}{.}%
\TextAt{12}{325.201}{443.865}{.}%
\TextAt{12}{334.921}{443.865}{.}%
\TextAt{12}{344.64}{443.865}{.}%
\TextAt{12}{354.36}{443.865}{.}%
\TextAt{12}{364.079}{443.865}{.}%
\TextAt{12}{373.799}{443.865}{.}%
\TextAt{12}{383.518}{443.865}{.}%
\TextAt{12}{393.238}{443.865}{.}%
\TextAt{12}{402.957}{443.865}{.}%
\TextAt{12}{412.676}{443.865}{.}%
\TextAt{12}{422.396}{443.865}{.}%
\TextAt{12}{432.115}{443.865}{.}%
\TextAt{12}{441.835}{443.865}{.}%
\TextAt{12}{451.554}{443.865}{.}%
\TextAt{12}{461.274}{443.865}{.}%
\TextAt{12}{470.993}{443.865}{.}%
\TextAt{12}{480.713}{443.865}{.}%
\TextAt{12}{490.432}{443.865}{.}%
\TextAt{12}{500.152}{443.865}{.}%
\TextAt{8}{515.195}{443.865}{239}%
\TextAt{8}{97.159}{423.83}{§}%
\TextAt{8}{104.831}{423.83}{1}%
\TextAt{8}{124.656}{423.83}{Infinite}%
\TextAt{8}{165.291}{423.83}{Galois}%
\TextAt{8}{201.205}{423.83}{Theory}%
\TextAt{8}{248.342}{423.83}{.}%
\TextAt{8}{257.572}{423.83}{.}%
\TextAt{8}{266.801}{423.83}{.}%
\TextAt{8}{276.03}{423.83}{.}%
\TextAt{8}{285.26}{423.83}{.}%
\TextAt{8}{294.489}{423.83}{.}%
\TextAt{8}{303.718}{423.83}{.}%
\TextAt{8}{312.948}{423.83}{.}%
\TextAt{8}{322.177}{423.83}{.}%
\TextAt{8}{331.406}{423.83}{.}%
\TextAt{8}{340.635}{423.83}{.}%
\TextAt{8}{349.865}{423.83}{.}%
\TextAt{8}{359.094}{423.83}{.}%
\TextAt{8}{368.323}{423.83}{.}%
\TextAt{8}{377.553}{423.83}{.}%
\TextAt{8}{386.782}{423.83}{.}%
\TextAt{8}{396.011}{423.83}{.}%
\TextAt{8}{405.241}{423.83}{.}%
\TextAt{8}{414.47}{423.83}{.}%
\TextAt{8}{423.699}{423.83}{.}%
\TextAt{8}{432.929}{423.83}{.}%
\TextAt{8}{442.158}{423.83}{.}%
\TextAt{8}{451.387}{423.83}{.}%
\TextAt{8}{460.616}{423.83}{.}%
\TextAt{8}{469.846}{423.83}{.}%
\TextAt{8}{479.075}{423.83}{.}%
\TextAt{8}{488.304}{423.83}{.}%
\TextAt{8}{497.534}{423.83}{.}%
\TextAt{8}{506.763}{423.83}{.}%
\TextAt{8}{515.195}{423.83}{239}%
\TextAt{8}{97.159}{403.796}{§}%
\TextAt{8}{104.831}{403.796}{2}%
\TextAt{8}{124.656}{403.796}{Projective}%
\TextAt{8}{180.092}{403.796}{and}%
\TextAt{8}{202.855}{403.796}{Inductive}%
\TextAt{8}{254.549}{403.796}{Limits}%
\TextAt{8}{294.489}{403.796}{.}%
\TextAt{8}{303.718}{403.796}{.}%
\TextAt{8}{312.948}{403.796}{.}%
\TextAt{8}{322.177}{403.796}{.}%
\TextAt{8}{331.406}{403.796}{.}%
\TextAt{8}{340.635}{403.796}{.}%
\TextAt{8}{349.865}{403.796}{.}%
\TextAt{8}{359.094}{403.796}{.}%
\TextAt{8}{368.323}{403.796}{.}%
\TextAt{8}{377.553}{403.796}{.}%
\TextAt{8}{386.782}{403.796}{.}%
\TextAt{8}{396.011}{403.796}{.}%
\TextAt{8}{405.241}{403.796}{.}%
\TextAt{8}{414.47}{403.796}{.}%
\TextAt{8}{423.699}{403.796}{.}%
\TextAt{8}{432.929}{403.796}{.}%
\TextAt{8}{442.158}{403.796}{.}%
\TextAt{8}{451.387}{403.796}{.}%
\TextAt{8}{460.616}{403.796}{.}%
\TextAt{8}{469.846}{403.796}{.}%
\TextAt{8}{479.075}{403.796}{.}%
\TextAt{8}{488.304}{403.796}{.}%
\TextAt{8}{497.534}{403.796}{.}%
\TextAt{8}{506.763}{403.796}{.}%
\TextAt{8}{515.195}{403.796}{243}%
\TextAt{8}{97.159}{383.761}{§}%
\TextAt{8}{104.831}{383.761}{3}%
\TextAt{8}{124.656}{383.761}{Abstract}%
\TextAt{8}{173.17}{383.761}{Galois}%
\TextAt{8}{209.083}{383.761}{Theory}%
\TextAt{8}{257.572}{383.761}{.}%
\TextAt{8}{266.801}{383.761}{.}%
\TextAt{8}{276.03}{383.761}{.}%
\TextAt{8}{285.26}{383.761}{.}%
\TextAt{8}{294.489}{383.761}{.}%
\TextAt{8}{303.718}{383.761}{.}%
\TextAt{8}{312.948}{383.761}{.}%
\TextAt{8}{322.177}{383.761}{.}%
\TextAt{8}{331.406}{383.761}{.}%
\TextAt{8}{340.635}{383.761}{.}%
\TextAt{8}{349.865}{383.761}{.}%
\TextAt{8}{359.094}{383.761}{.}%
\TextAt{8}{368.323}{383.761}{.}%
\TextAt{8}{377.553}{383.761}{.}%
\TextAt{8}{386.782}{383.761}{.}%
\TextAt{8}{396.011}{383.761}{.}%
\TextAt{8}{405.241}{383.761}{.}%
\TextAt{8}{414.47}{383.761}{.}%
\TextAt{8}{423.699}{383.761}{.}%
\TextAt{8}{432.929}{383.761}{.}%
\TextAt{8}{442.158}{383.761}{.}%
\TextAt{8}{451.387}{383.761}{.}%
\TextAt{8}{460.616}{383.761}{.}%
\TextAt{8}{469.846}{383.761}{.}%
\TextAt{8}{479.075}{383.761}{.}%
\TextAt{8}{488.304}{383.761}{.}%
\TextAt{8}{497.534}{383.761}{.}%
\TextAt{8}{506.763}{383.761}{.}%
\TextAt{8}{515.195}{383.761}{251}%
\TextAt{8}{97.159}{363.726}{§}%
\TextAt{8}{104.831}{363.726}{4}%
\TextAt{8}{124.656}{363.726}{Abstract}%
\TextAt{8}{173.17}{363.726}{Valuation}%
\TextAt{8}{226.502}{363.726}{Theory}%
\TextAt{8}{266.801}{363.726}{.}%
\TextAt{8}{276.03}{363.726}{.}%
\TextAt{8}{285.26}{363.726}{.}%
\TextAt{8}{294.489}{363.726}{.}%
\TextAt{8}{303.718}{363.726}{.}%
\TextAt{8}{312.948}{363.726}{.}%
\TextAt{8}{322.177}{363.726}{.}%
\TextAt{8}{331.406}{363.726}{.}%
\TextAt{8}{340.635}{363.726}{.}%
\TextAt{8}{349.865}{363.726}{.}%
\TextAt{8}{359.094}{363.726}{.}%
\TextAt{8}{368.323}{363.726}{.}%
\TextAt{8}{377.553}{363.726}{.}%
\TextAt{8}{386.782}{363.726}{.}%
\TextAt{8}{396.011}{363.726}{.}%
\TextAt{8}{405.241}{363.726}{.}%
\TextAt{8}{414.47}{363.726}{.}%
\TextAt{8}{423.699}{363.726}{.}%
\TextAt{8}{432.929}{363.726}{.}%
\TextAt{8}{442.158}{363.726}{.}%
\TextAt{8}{451.387}{363.726}{.}%
\TextAt{8}{460.616}{363.726}{.}%
\TextAt{8}{469.846}{363.726}{.}%
\TextAt{8}{479.075}{363.726}{.}%
\TextAt{8}{488.304}{363.726}{.}%
\TextAt{8}{497.534}{363.726}{.}%
\TextAt{8}{506.763}{363.726}{.}%
\TextAt{8}{515.195}{363.726}{260}%
\TextAt{8}{97.159}{343.691}{§}%
\TextAt{8}{104.831}{343.691}{5}%
\TextAt{8}{124.656}{343.691}{The}%
\TextAt{8}{148.71}{343.691}{Reciprocity}%
\TextAt{8}{210.984}{343.691}{Map}%
\TextAt{8}{239.113}{343.691}{.}%
\TextAt{8}{248.342}{343.691}{.}%
\TextAt{8}{257.572}{343.691}{.}%
\TextAt{8}{266.801}{343.691}{.}%
\TextAt{8}{276.03}{343.691}{.}%
\TextAt{8}{285.26}{343.691}{.}%
\TextAt{8}{294.489}{343.691}{.}%
\TextAt{8}{303.718}{343.691}{.}%
\TextAt{8}{312.948}{343.691}{.}%
\TextAt{8}{322.177}{343.691}{.}%
\TextAt{8}{331.406}{343.691}{.}%
\TextAt{8}{340.635}{343.691}{.}%
\TextAt{8}{349.865}{343.691}{.}%
\TextAt{8}{359.094}{343.691}{.}%
\TextAt{8}{368.323}{343.691}{.}%
\TextAt{8}{377.553}{343.691}{.}%
\TextAt{8}{386.782}{343.691}{.}%
\TextAt{8}{396.011}{343.691}{.}%
\TextAt{8}{405.241}{343.691}{.}%
\TextAt{8}{414.47}{343.691}{.}%
\TextAt{8}{423.699}{343.691}{.}%
\TextAt{8}{432.929}{343.691}{.}%
\TextAt{8}{442.158}{343.691}{.}%
\TextAt{8}{451.387}{343.691}{.}%
\TextAt{8}{460.616}{343.691}{.}%
\TextAt{8}{469.846}{343.691}{.}%
\TextAt{8}{479.075}{343.691}{.}%
\TextAt{8}{488.304}{343.691}{.}%
\TextAt{8}{497.534}{343.691}{.}%
\TextAt{8}{506.763}{343.691}{.}%
\TextAt{8}{515.195}{343.691}{265}%
\TextAt{8}{97.159}{323.656}{§}%
\TextAt{8}{104.831}{323.656}{6}%
\TextAt{8}{124.656}{323.656}{The}%
\TextAt{8}{148.71}{323.656}{General}%
\TextAt{8}{192.358}{323.656}{Reciprocity}%
\TextAt{8}{254.632}{323.656}{Law}%
\TextAt{8}{285.26}{323.656}{.}%
\TextAt{8}{294.489}{323.656}{.}%
\TextAt{8}{303.718}{323.656}{.}%
\TextAt{8}{312.948}{323.656}{.}%
\TextAt{8}{322.177}{323.656}{.}%
\TextAt{8}{331.406}{323.656}{.}%
\TextAt{8}{340.635}{323.656}{.}%
\TextAt{8}{349.865}{323.656}{.}%
\TextAt{8}{359.094}{323.656}{.}%
\TextAt{8}{368.323}{323.656}{.}%
\TextAt{8}{377.553}{323.656}{.}%
\TextAt{8}{386.782}{323.656}{.}%
\TextAt{8}{396.011}{323.656}{.}%
\TextAt{8}{405.241}{323.656}{.}%
\TextAt{8}{414.47}{323.656}{.}%
\TextAt{8}{423.699}{323.656}{.}%
\TextAt{8}{432.929}{323.656}{.}%
\TextAt{8}{442.158}{323.656}{.}%
\TextAt{8}{451.387}{323.656}{.}%
\TextAt{8}{460.616}{323.656}{.}%
\TextAt{8}{469.846}{323.656}{.}%
\TextAt{8}{479.075}{323.656}{.}%
\TextAt{8}{488.304}{323.656}{.}%
\TextAt{8}{497.534}{323.656}{.}%
\TextAt{8}{506.763}{323.656}{.}%
\TextAt{8}{515.195}{323.656}{274}%
\TextAt{8}{97.159}{303.622}{§}%
\TextAt{8}{104.831}{303.622}{7}%
\TextAt{8}{124.656}{303.622}{The}%
\TextAt{8}{148.71}{303.622}{Herbrand}%
\TextAt{8}{201.061}{303.622}{Quotient}%
\TextAt{8}{257.572}{303.622}{.}%
\TextAt{8}{266.801}{303.622}{.}%
\TextAt{8}{276.03}{303.622}{.}%
\TextAt{8}{285.26}{303.622}{.}%
\TextAt{8}{294.489}{303.622}{.}%
\TextAt{8}{303.718}{303.622}{.}%
\TextAt{8}{312.948}{303.622}{.}%
\TextAt{8}{322.177}{303.622}{.}%
\TextAt{8}{331.406}{303.622}{.}%
\TextAt{8}{340.635}{303.622}{.}%
\TextAt{8}{349.865}{303.622}{.}%
\TextAt{8}{359.094}{303.622}{.}%
\TextAt{8}{368.323}{303.622}{.}%
\TextAt{8}{377.553}{303.622}{.}%
\TextAt{8}{386.782}{303.622}{.}%
\TextAt{8}{396.011}{303.622}{.}%
\TextAt{8}{405.241}{303.622}{.}%
\TextAt{8}{414.47}{303.622}{.}%
\TextAt{8}{423.699}{303.622}{.}%
\TextAt{8}{432.929}{303.622}{.}%
\TextAt{8}{442.158}{303.622}{.}%
\TextAt{8}{451.387}{303.622}{.}%
\TextAt{8}{460.616}{303.622}{.}%
\TextAt{8}{469.846}{303.622}{.}%
\TextAt{8}{479.075}{303.622}{.}%
\TextAt{8}{488.304}{303.622}{.}%
\TextAt{8}{497.534}{303.622}{.}%
\TextAt{8}{506.763}{303.622}{.}%
\TextAt{8}{515.195}{303.622}{286}%
\TextAt{12}{79.226}{275.218}{V}%
\TextAt{12}{109.114}{275.218}{Local}%
\TextAt{12}{145.039}{275.218}{Class}%
\TextAt{12}{180.14}{275.218}{Field}%
\TextAt{12}{214.164}{275.218}{Theory}%
\TextAt{12}{266.885}{275.218}{.}%
\TextAt{12}{276.604}{275.218}{.}%
\TextAt{12}{286.323}{275.218}{.}%
\TextAt{12}{296.043}{275.218}{.}%
\TextAt{12}{305.762}{275.218}{.}%
\TextAt{12}{315.482}{275.218}{.}%
\TextAt{12}{325.201}{275.218}{.}%
\TextAt{12}{334.921}{275.218}{.}%
\TextAt{12}{344.64}{275.218}{.}%
\TextAt{12}{354.36}{275.218}{.}%
\TextAt{12}{364.079}{275.218}{.}%
\TextAt{12}{373.799}{275.218}{.}%
\TextAt{12}{383.518}{275.218}{.}%
\TextAt{12}{393.238}{275.218}{.}%
\TextAt{12}{402.957}{275.218}{.}%
\TextAt{12}{412.676}{275.218}{.}%
\TextAt{12}{422.396}{275.218}{.}%
\TextAt{12}{432.115}{275.218}{.}%
\TextAt{12}{441.835}{275.218}{.}%
\TextAt{12}{451.554}{275.218}{.}%
\TextAt{12}{461.274}{275.218}{.}%
\TextAt{12}{470.993}{275.218}{.}%
\TextAt{12}{480.713}{275.218}{.}%
\TextAt{12}{490.432}{275.218}{.}%
\TextAt{12}{500.152}{275.218}{.}%
\TextAt{8}{515.195}{275.218}{291}%
\TextAt{8}{97.159}{255.184}{§}%
\TextAt{8}{104.831}{255.184}{1}%
\TextAt{8}{124.656}{255.184}{The}%
\TextAt{8}{148.71}{255.184}{Local}%
\TextAt{8}{180.415}{255.184}{Reciprocity}%
\TextAt{8}{242.689}{255.184}{Law}%
\TextAt{8}{276.03}{255.184}{.}%
\TextAt{8}{285.26}{255.184}{.}%
\TextAt{8}{294.489}{255.184}{.}%
\TextAt{8}{303.718}{255.184}{.}%
\TextAt{8}{312.948}{255.184}{.}%
\TextAt{8}{322.177}{255.184}{.}%
\TextAt{8}{331.406}{255.184}{.}%
\TextAt{8}{340.635}{255.184}{.}%
\TextAt{8}{349.865}{255.184}{.}%
\TextAt{8}{359.094}{255.184}{.}%
\TextAt{8}{368.323}{255.184}{.}%
\TextAt{8}{377.553}{255.184}{.}%
\TextAt{8}{386.782}{255.184}{.}%
\TextAt{8}{396.011}{255.184}{.}%
\TextAt{8}{405.241}{255.184}{.}%
\TextAt{8}{414.47}{255.184}{.}%
\TextAt{8}{423.699}{255.184}{.}%
\TextAt{8}{432.929}{255.184}{.}%
\TextAt{8}{442.158}{255.184}{.}%
\TextAt{8}{451.387}{255.184}{.}%
\TextAt{8}{460.616}{255.184}{.}%
\TextAt{8}{469.846}{255.184}{.}%
\TextAt{8}{479.075}{255.184}{.}%
\TextAt{8}{488.304}{255.184}{.}%
\TextAt{8}{497.534}{255.184}{.}%
\TextAt{8}{506.763}{255.184}{.}%
\TextAt{8}{515.195}{255.184}{291}%
\TextAt{8}{97.159}{235.149}{§}%
\TextAt{8}{104.831}{235.149}{2}%
\TextAt{8}{124.656}{235.149}{The}%
\TextAt{8}{148.71}{235.149}{Norm}%
\TextAt{8}{181.55}{235.149}{Residue}%
\TextAt{8}{225.342}{235.149}{Symbol}%
\TextAt{8}{267.293}{235.149}{over}%
\TextAt{8}{292.339}{235.149}{Qp}%
\TextAt{8}{312.948}{235.149}{.}%
\TextAt{8}{322.177}{235.149}{.}%
\TextAt{8}{331.406}{235.149}{.}%
\TextAt{8}{340.635}{235.149}{.}%
\TextAt{8}{349.865}{235.149}{.}%
\TextAt{8}{359.094}{235.149}{.}%
\TextAt{8}{368.323}{235.149}{.}%
\TextAt{8}{377.553}{235.149}{.}%
\TextAt{8}{386.782}{235.149}{.}%
\TextAt{8}{396.011}{235.149}{.}%
\TextAt{8}{405.241}{235.149}{.}%
\TextAt{8}{414.47}{235.149}{.}%
\TextAt{8}{423.699}{235.149}{.}%
\TextAt{8}{432.929}{235.149}{.}%
\TextAt{8}{442.158}{235.149}{.}%
\TextAt{8}{451.387}{235.149}{.}%
\TextAt{8}{460.616}{235.149}{.}%
\TextAt{8}{469.846}{235.149}{.}%
\TextAt{8}{479.075}{235.149}{.}%
\TextAt{8}{488.304}{235.149}{.}%
\TextAt{8}{497.534}{235.149}{.}%
\TextAt{8}{506.763}{235.149}{.}%
\TextAt{8}{515.195}{235.149}{300}%
\TextAt{8}{97.159}{215.114}{§}%
\TextAt{8}{104.831}{215.114}{3}%
\TextAt{8}{124.656}{215.114}{The}%
\TextAt{8}{148.71}{215.114}{Hilbert}%
\TextAt{8}{189.022}{215.114}{Symbol}%
\TextAt{8}{239.113}{215.114}{.}%
\TextAt{8}{248.342}{215.114}{.}%
\TextAt{8}{257.572}{215.114}{.}%
\TextAt{8}{266.801}{215.114}{.}%
\TextAt{8}{276.03}{215.114}{.}%
\TextAt{8}{285.26}{215.114}{.}%
\TextAt{8}{294.489}{215.114}{.}%
\TextAt{8}{303.718}{215.114}{.}%
\TextAt{8}{312.948}{215.114}{.}%
\TextAt{8}{322.177}{215.114}{.}%
\TextAt{8}{331.406}{215.114}{.}%
\TextAt{8}{340.635}{215.114}{.}%
\TextAt{8}{349.865}{215.114}{.}%
\TextAt{8}{359.094}{215.114}{.}%
\TextAt{8}{368.323}{215.114}{.}%
\TextAt{8}{377.553}{215.114}{.}%
\TextAt{8}{386.782}{215.114}{.}%
\TextAt{8}{396.011}{215.114}{.}%
\TextAt{8}{405.241}{215.114}{.}%
\TextAt{8}{414.47}{215.114}{.}%
\TextAt{8}{423.699}{215.114}{.}%
\TextAt{8}{432.929}{215.114}{.}%
\TextAt{8}{442.158}{215.114}{.}%
\TextAt{8}{451.387}{215.114}{.}%
\TextAt{8}{460.616}{215.114}{.}%
\TextAt{8}{469.846}{215.114}{.}%
\TextAt{8}{479.075}{215.114}{.}%
\TextAt{8}{488.304}{215.114}{.}%
\TextAt{8}{497.534}{215.114}{.}%
\TextAt{8}{506.763}{215.114}{.}%
\TextAt{8}{515.195}{215.114}{305}%
\TextAt{8}{97.159}{195.079}{§}%
\TextAt{8}{104.831}{195.079}{4}%
\TextAt{8}{124.656}{195.079}{Formal}%
\TextAt{8}{164.49}{195.079}{Groups}%
\TextAt{8}{211.425}{195.079}{.}%
\TextAt{8}{220.654}{195.079}{.}%
\TextAt{8}{229.884}{195.079}{.}%
\TextAt{8}{239.113}{195.079}{.}%
\TextAt{8}{248.342}{195.079}{.}%
\TextAt{8}{257.572}{195.079}{.}%
\TextAt{8}{266.801}{195.079}{.}%
\TextAt{8}{276.03}{195.079}{.}%
\TextAt{8}{285.26}{195.079}{.}%
\TextAt{8}{294.489}{195.079}{.}%
\TextAt{8}{303.718}{195.079}{.}%
\TextAt{8}{312.948}{195.079}{.}%
\TextAt{8}{322.177}{195.079}{.}%
\TextAt{8}{331.406}{195.079}{.}%
\TextAt{8}{340.635}{195.079}{.}%
\TextAt{8}{349.865}{195.079}{.}%
\TextAt{8}{359.094}{195.079}{.}%
\TextAt{8}{368.323}{195.079}{.}%
\TextAt{8}{377.553}{195.079}{.}%
\TextAt{8}{386.782}{195.079}{.}%
\TextAt{8}{396.011}{195.079}{.}%
\TextAt{8}{405.241}{195.079}{.}%
\TextAt{8}{414.47}{195.079}{.}%
\TextAt{8}{423.699}{195.079}{.}%
\TextAt{8}{432.929}{195.079}{.}%
\TextAt{8}{442.158}{195.079}{.}%
\TextAt{8}{451.387}{195.079}{.}%
\TextAt{8}{460.616}{195.079}{.}%
\TextAt{8}{469.846}{195.079}{.}%
\TextAt{8}{479.075}{195.079}{.}%
\TextAt{8}{488.304}{195.079}{.}%
\TextAt{8}{497.534}{195.079}{.}%
\TextAt{8}{506.763}{195.079}{.}%
\TextAt{8}{515.195}{195.079}{313}%
\TextAt{8}{97.159}{175.045}{§}%
\TextAt{8}{104.831}{175.045}{5}%
\TextAt{8}{124.656}{175.045}{Generalized}%
\TextAt{8}{188.461}{175.045}{Cyclotomic}%
\TextAt{8}{249.922}{175.045}{Theory}%
\TextAt{8}{294.489}{175.045}{.}%
\TextAt{8}{303.718}{175.045}{.}%
\TextAt{8}{312.948}{175.045}{.}%
\TextAt{8}{322.177}{175.045}{.}%
\TextAt{8}{331.406}{175.045}{.}%
\TextAt{8}{340.635}{175.045}{.}%
\TextAt{8}{349.865}{175.045}{.}%
\TextAt{8}{359.094}{175.045}{.}%
\TextAt{8}{368.323}{175.045}{.}%
\TextAt{8}{377.553}{175.045}{.}%
\TextAt{8}{386.782}{175.045}{.}%
\TextAt{8}{396.011}{175.045}{.}%
\TextAt{8}{405.241}{175.045}{.}%
\TextAt{8}{414.47}{175.045}{.}%
\TextAt{8}{423.699}{175.045}{.}%
\TextAt{8}{432.929}{175.045}{.}%
\TextAt{8}{442.158}{175.045}{.}%
\TextAt{8}{451.387}{175.045}{.}%
\TextAt{8}{460.616}{175.045}{.}%
\TextAt{8}{469.846}{175.045}{.}%
\TextAt{8}{479.075}{175.045}{.}%
\TextAt{8}{488.304}{175.045}{.}%
\TextAt{8}{497.534}{175.045}{.}%
\TextAt{8}{506.763}{175.045}{.}%
\TextAt{8}{515.195}{175.045}{317}%
\TextAt{8}{97.159}{155.01}{§}%
\TextAt{8}{104.831}{155.01}{6}%
\TextAt{8}{124.656}{155.01}{Higher}%
\TextAt{8}{162.697}{155.01}{Ramification}%
\TextAt{8}{231.81}{155.01}{Groups}%
\TextAt{8}{276.03}{155.01}{.}%
\TextAt{8}{285.26}{155.01}{.}%
\TextAt{8}{294.489}{155.01}{.}%
\TextAt{8}{303.718}{155.01}{.}%
\TextAt{8}{312.948}{155.01}{.}%
\TextAt{8}{322.177}{155.01}{.}%
\TextAt{8}{331.406}{155.01}{.}%
\TextAt{8}{340.635}{155.01}{.}%
\TextAt{8}{349.865}{155.01}{.}%
\TextAt{8}{359.094}{155.01}{.}%
\TextAt{8}{368.323}{155.01}{.}%
\TextAt{8}{377.553}{155.01}{.}%
\TextAt{8}{386.782}{155.01}{.}%
\TextAt{8}{396.011}{155.01}{.}%
\TextAt{8}{405.241}{155.01}{.}%
\TextAt{8}{414.47}{155.01}{.}%
\TextAt{8}{423.699}{155.01}{.}%
\TextAt{8}{432.929}{155.01}{.}%
\TextAt{8}{442.158}{155.01}{.}%
\TextAt{8}{451.387}{155.01}{.}%
\TextAt{8}{460.616}{155.01}{.}%
\TextAt{8}{469.846}{155.01}{.}%
\TextAt{8}{479.075}{155.01}{.}%
\TextAt{8}{488.304}{155.01}{.}%
\TextAt{8}{497.534}{155.01}{.}%
\TextAt{8}{506.763}{155.01}{.}%
\TextAt{8}{515.195}{155.01}{323}%
\TextAt{12}{79.226}{126.606}{VI}%
\TextAt{12}{109.114}{126.606}{Global}%
\TextAt{12}{152.392}{126.606}{Class}%
\TextAt{12}{187.492}{126.606}{Field}%
\TextAt{12}{221.517}{126.606}{Theory}%
\TextAt{12}{266.885}{126.606}{.}%
\TextAt{12}{276.604}{126.606}{.}%
\TextAt{12}{286.323}{126.606}{.}%
\TextAt{12}{296.043}{126.606}{.}%
\TextAt{12}{305.762}{126.606}{.}%
\TextAt{12}{315.482}{126.606}{.}%
\TextAt{12}{325.201}{126.606}{.}%
\TextAt{12}{334.921}{126.606}{.}%
\TextAt{12}{344.64}{126.606}{.}%
\TextAt{12}{354.36}{126.606}{.}%
\TextAt{12}{364.079}{126.606}{.}%
\TextAt{12}{373.799}{126.606}{.}%
\TextAt{12}{383.518}{126.606}{.}%
\TextAt{12}{393.238}{126.606}{.}%
\TextAt{12}{402.957}{126.606}{.}%
\TextAt{12}{412.676}{126.606}{.}%
\TextAt{12}{422.396}{126.606}{.}%
\TextAt{12}{432.115}{126.606}{.}%
\TextAt{12}{441.835}{126.606}{.}%
\TextAt{12}{451.554}{126.606}{.}%
\TextAt{12}{461.274}{126.606}{.}%
\TextAt{12}{470.993}{126.606}{.}%
\TextAt{12}{480.713}{126.606}{.}%
\TextAt{12}{490.432}{126.606}{.}%
\TextAt{12}{500.152}{126.606}{.}%
\TextAt{8}{515.195}{126.606}{326}%
\TextAt{8}{97.159}{106.572}{§}%
\TextAt{8}{104.831}{106.572}{1}%
\TextAt{8}{124.656}{106.572}{Idèles}%
\TextAt{8}{157.544}{106.572}{and}%
\TextAt{8}{180.307}{106.572}{Idèle}%
\TextAt{8}{208.581}{106.572}{Classes}%
\TextAt{8}{248.342}{106.572}{.}%
\TextAt{8}{257.572}{106.572}{.}%
\TextAt{8}{266.801}{106.572}{.}%
\TextAt{8}{276.03}{106.572}{.}%
\TextAt{8}{285.26}{106.572}{.}%
\TextAt{8}{294.489}{106.572}{.}%
\TextAt{8}{303.718}{106.572}{.}%
\TextAt{8}{312.948}{106.572}{.}%
\TextAt{8}{322.177}{106.572}{.}%
\TextAt{8}{331.406}{106.572}{.}%
\TextAt{8}{340.635}{106.572}{.}%
\TextAt{8}{349.865}{106.572}{.}%
\TextAt{8}{359.094}{106.572}{.}%
\TextAt{8}{368.323}{106.572}{.}%
\TextAt{8}{377.553}{106.572}{.}%
\TextAt{8}{386.782}{106.572}{.}%
\TextAt{8}{396.011}{106.572}{.}%
\TextAt{8}{405.241}{106.572}{.}%
\TextAt{8}{414.47}{106.572}{.}%
\TextAt{8}{423.699}{106.572}{.}%
\TextAt{8}{432.929}{106.572}{.}%
\TextAt{8}{442.158}{106.572}{.}%
\TextAt{8}{451.387}{106.572}{.}%
\TextAt{8}{460.616}{106.572}{.}%
\TextAt{8}{469.846}{106.572}{.}%
\TextAt{8}{479.075}{106.572}{.}%
\TextAt{8}{488.304}{106.572}{.}%
\TextAt{8}{497.534}{106.572}{.}%
\TextAt{8}{506.763}{106.572}{.}%
\TextAt{8}{515.195}{106.572}{326}%
\TextAt{8}{97.159}{86.537}{§}%
\TextAt{8}{104.831}{86.537}{2}%
\TextAt{8}{124.656}{86.537}{Idèles}%
\TextAt{8}{157.544}{86.537}{in}%
\TextAt{8}{171.197}{86.537}{Field}%
\TextAt{8}{200.942}{86.537}{Extensions}%
\TextAt{8}{266.801}{86.537}{.}%
\TextAt{8}{276.03}{86.537}{.}%
\TextAt{8}{285.26}{86.537}{.}%
\TextAt{8}{294.489}{86.537}{.}%
\TextAt{8}{303.718}{86.537}{.}%
\TextAt{8}{312.948}{86.537}{.}%
\TextAt{8}{322.177}{86.537}{.}%
\TextAt{8}{331.406}{86.537}{.}%
\TextAt{8}{340.635}{86.537}{.}%
\TextAt{8}{349.865}{86.537}{.}%
\TextAt{8}{359.094}{86.537}{.}%
\TextAt{8}{368.323}{86.537}{.}%
\TextAt{8}{377.553}{86.537}{.}%
\TextAt{8}{386.782}{86.537}{.}%
\TextAt{8}{396.011}{86.537}{.}%
\TextAt{8}{405.241}{86.537}{.}%
\TextAt{8}{414.47}{86.537}{.}%
\TextAt{8}{423.699}{86.537}{.}%
\TextAt{8}{432.929}{86.537}{.}%
\TextAt{8}{442.158}{86.537}{.}%
\TextAt{8}{451.387}{86.537}{.}%
\TextAt{8}{460.616}{86.537}{.}%
\TextAt{8}{469.846}{86.537}{.}%
\TextAt{8}{479.075}{86.537}{.}%
\TextAt{8}{488.304}{86.537}{.}%
\TextAt{8}{497.534}{86.537}{.}%
\TextAt{8}{506.763}{86.537}{.}%
\TextAt{8}{515.195}{86.537}{336}%
\TextAt{8}{301.281}{43.033}{xi}%
\Destination{78.23}{809.815}{nav.page.xi}%
\Hotlink{97.159}{757.776}{89.664}{16.821}{nav.section.2.5}%
\Hotlink{97.159}{737.741}{111.016}{16.821}{nav.section.2.6}%
\Hotlink{97.159}{717.706}{257.514}{16.821}{nav.section.2.7}%
\Hotlink{97.159}{697.672}{154.018}{16.821}{nav.section.2.8}%
\Hotlink{97.159}{677.637}{171.437}{16.821}{nav.section.2.9}%
\Hotlink{97.159}{657.602}{171.867}{16.821}{nav.section.2.10}%
\Hotlink{79.226}{629.091}{166.201}{17.048}{nav.chapter.3}%
\Hotlink{97.159}{609.164}{62.836}{16.821}{nav.section.3.1}%
\Hotlink{97.159}{589.129}{164.79}{16.821}{nav.section.3.2}%
\Hotlink{97.159}{569.095}{105.229}{16.82}{nav.section.3.3}%
\Hotlink{97.159}{549.06}{129.305}{16.821}{nav.section.3.4}%
\Hotlink{97.159}{529.025}{136.982}{16.821}{nav.section.3.5}%
\Hotlink{97.159}{508.99}{137.401}{16.821}{nav.section.3.6}%
\Hotlink{97.159}{488.956}{179.459}{16.82}{nav.section.3.7}%
\Hotlink{97.159}{468.921}{214.763}{16.821}{nav.section.3.8}%
\Hotlink{79.226}{440.41}{197.213}{17.048}{nav.chapter.4}%
\Hotlink{97.159}{420.483}{140.796}{16.821}{nav.section.4.1}%
\Hotlink{97.159}{400.448}{190.135}{16.821}{nav.section.4.2}%
\Hotlink{97.159}{380.413}{148.674}{16.821}{nav.section.4.3}%
\Hotlink{97.159}{360.378}{166.093}{16.821}{nav.section.4.4}%
\Hotlink{97.159}{340.344}{136.911}{16.82}{nav.section.4.5}%
\Hotlink{97.159}{320.309}{178.778}{16.821}{nav.section.4.6}%
\Hotlink{97.159}{300.274}{149.117}{16.821}{nav.section.4.7}%
\Hotlink{79.226}{271.763}{177.224}{17.048}{nav.chapter.5}%
\Hotlink{97.159}{251.836}{166.834}{16.821}{nav.section.5.1}%
\Hotlink{97.159}{231.801}{210.793}{16.821}{nav.section.5.2}%
\Hotlink{97.159}{211.766}{129.917}{16.821}{nav.section.5.3}%
\Hotlink{97.159}{191.732}{104.548}{16.821}{nav.section.5.4}%
\Hotlink{97.159}{171.697}{189.513}{16.821}{nav.section.5.5}%
\Hotlink{97.159}{151.662}{171.867}{16.821}{nav.section.5.6}%
\Hotlink{79.226}{123.151}{184.576}{17.048}{nav.chapter.6}%
\Hotlink{97.159}{103.224}{148.029}{16.821}{nav.section.6.1}%
\Hotlink{97.159}{83.189}{159.028}{16.821}{nav.section.6.2}%
}

% Page 12 (body page front matter)
\Page{%
\TextAt{2}{79.226}{796.943}{Contents}%
\Vector{q 0 G 0 g 0.299 w 7.226 18.729 m 460.769 18.729 l S Q}%
\TextAt{8}{97.159}{761.124}{§}%
\TextAt{8}{104.831}{761.124}{3}%
\TextAt{8}{124.656}{761.124}{The}%
\TextAt{8}{148.71}{761.124}{Herbrand}%
\TextAt{8}{201.061}{761.124}{Quotient}%
\TextAt{8}{250.173}{761.124}{of}%
\TextAt{8}{263.503}{761.124}{the}%
\TextAt{8}{283.66}{761.124}{Idèle}%
\TextAt{8}{311.934}{761.124}{Class}%
\TextAt{8}{342.622}{761.124}{Group}%
\TextAt{8}{386.782}{761.124}{.}%
\TextAt{8}{396.011}{761.124}{.}%
\TextAt{8}{405.241}{761.124}{.}%
\TextAt{8}{414.47}{761.124}{.}%
\TextAt{8}{423.699}{761.124}{.}%
\TextAt{8}{432.929}{761.124}{.}%
\TextAt{8}{442.158}{761.124}{.}%
\TextAt{8}{451.387}{761.124}{.}%
\TextAt{8}{460.616}{761.124}{.}%
\TextAt{8}{469.846}{761.124}{.}%
\TextAt{8}{479.075}{761.124}{.}%
\TextAt{8}{488.304}{761.124}{.}%
\TextAt{8}{497.534}{761.124}{.}%
\TextAt{8}{506.763}{761.124}{.}%
\TextAt{8}{515.195}{761.124}{341}%
\TextAt{8}{97.159}{741.089}{§}%
\TextAt{8}{104.831}{741.089}{4}%
\TextAt{8}{124.656}{741.089}{The}%
\TextAt{8}{148.71}{741.089}{Class}%
\TextAt{8}{179.399}{741.089}{Field}%
\TextAt{8}{209.143}{741.089}{Axiom}%
\TextAt{8}{248.342}{741.089}{.}%
\TextAt{8}{257.572}{741.089}{.}%
\TextAt{8}{266.801}{741.089}{.}%
\TextAt{8}{276.03}{741.089}{.}%
\TextAt{8}{285.26}{741.089}{.}%
\TextAt{8}{294.489}{741.089}{.}%
\TextAt{8}{303.718}{741.089}{.}%
\TextAt{8}{312.948}{741.089}{.}%
\TextAt{8}{322.177}{741.089}{.}%
\TextAt{8}{331.406}{741.089}{.}%
\TextAt{8}{340.635}{741.089}{.}%
\TextAt{8}{349.865}{741.089}{.}%
\TextAt{8}{359.094}{741.089}{.}%
\TextAt{8}{368.323}{741.089}{.}%
\TextAt{8}{377.553}{741.089}{.}%
\TextAt{8}{386.782}{741.089}{.}%
\TextAt{8}{396.011}{741.089}{.}%
\TextAt{8}{405.241}{741.089}{.}%
\TextAt{8}{414.47}{741.089}{.}%
\TextAt{8}{423.699}{741.089}{.}%
\TextAt{8}{432.929}{741.089}{.}%
\TextAt{8}{442.158}{741.089}{.}%
\TextAt{8}{451.387}{741.089}{.}%
\TextAt{8}{460.616}{741.089}{.}%
\TextAt{8}{469.846}{741.089}{.}%
\TextAt{8}{479.075}{741.089}{.}%
\TextAt{8}{488.304}{741.089}{.}%
\TextAt{8}{497.534}{741.089}{.}%
\TextAt{8}{506.763}{741.089}{.}%
\TextAt{8}{515.195}{741.089}{346}%
\TextAt{8}{97.159}{721.054}{§}%
\TextAt{8}{104.831}{721.054}{5}%
\TextAt{8}{124.656}{721.054}{The}%
\TextAt{8}{148.71}{721.054}{Global}%
\TextAt{8}{186.512}{721.054}{Reciprocity}%
\TextAt{8}{248.786}{721.054}{Law}%
\TextAt{8}{276.03}{721.054}{.}%
\TextAt{8}{285.26}{721.054}{.}%
\TextAt{8}{294.489}{721.054}{.}%
\TextAt{8}{303.718}{721.054}{.}%
\TextAt{8}{312.948}{721.054}{.}%
\TextAt{8}{322.177}{721.054}{.}%
\TextAt{8}{331.406}{721.054}{.}%
\TextAt{8}{340.635}{721.054}{.}%
\TextAt{8}{349.865}{721.054}{.}%
\TextAt{8}{359.094}{721.054}{.}%
\TextAt{8}{368.323}{721.054}{.}%
\TextAt{8}{377.553}{721.054}{.}%
\TextAt{8}{386.782}{721.054}{.}%
\TextAt{8}{396.011}{721.054}{.}%
\TextAt{8}{405.241}{721.054}{.}%
\TextAt{8}{414.47}{721.054}{.}%
\TextAt{8}{423.699}{721.054}{.}%
\TextAt{8}{432.929}{721.054}{.}%
\TextAt{8}{442.158}{721.054}{.}%
\TextAt{8}{451.387}{721.054}{.}%
\TextAt{8}{460.616}{721.054}{.}%
\TextAt{8}{469.846}{721.054}{.}%
\TextAt{8}{479.075}{721.054}{.}%
\TextAt{8}{488.304}{721.054}{.}%
\TextAt{8}{497.534}{721.054}{.}%
\TextAt{8}{506.763}{721.054}{.}%
\TextAt{8}{515.195}{721.054}{351}%
\TextAt{8}{97.159}{701.019}{§}%
\TextAt{8}{104.831}{701.019}{6}%
\TextAt{8}{124.656}{701.019}{Global}%
\TextAt{8}{162.458}{701.019}{Class}%
\TextAt{8}{193.147}{701.019}{Fields}%
\TextAt{8}{229.884}{701.019}{.}%
\TextAt{8}{239.113}{701.019}{.}%
\TextAt{8}{248.342}{701.019}{.}%
\TextAt{8}{257.572}{701.019}{.}%
\TextAt{8}{266.801}{701.019}{.}%
\TextAt{8}{276.03}{701.019}{.}%
\TextAt{8}{285.26}{701.019}{.}%
\TextAt{8}{294.489}{701.019}{.}%
\TextAt{8}{303.718}{701.019}{.}%
\TextAt{8}{312.948}{701.019}{.}%
\TextAt{8}{322.177}{701.019}{.}%
\TextAt{8}{331.406}{701.019}{.}%
\TextAt{8}{340.635}{701.019}{.}%
\TextAt{8}{349.865}{701.019}{.}%
\TextAt{8}{359.094}{701.019}{.}%
\TextAt{8}{368.323}{701.019}{.}%
\TextAt{8}{377.553}{701.019}{.}%
\TextAt{8}{386.782}{701.019}{.}%
\TextAt{8}{396.011}{701.019}{.}%
\TextAt{8}{405.241}{701.019}{.}%
\TextAt{8}{414.47}{701.019}{.}%
\TextAt{8}{423.699}{701.019}{.}%
\TextAt{8}{432.929}{701.019}{.}%
\TextAt{8}{442.158}{701.019}{.}%
\TextAt{8}{451.387}{701.019}{.}%
\TextAt{8}{460.616}{701.019}{.}%
\TextAt{8}{469.846}{701.019}{.}%
\TextAt{8}{479.075}{701.019}{.}%
\TextAt{8}{488.304}{701.019}{.}%
\TextAt{8}{497.534}{701.019}{.}%
\TextAt{8}{506.763}{701.019}{.}%
\TextAt{8}{515.195}{701.019}{360}%
\TextAt{8}{97.159}{680.985}{§}%
\TextAt{8}{104.831}{680.985}{7}%
\TextAt{8}{124.656}{680.985}{The}%
\TextAt{8}{148.71}{680.985}{Ideal-Theoretic}%
\TextAt{8}{230.316}{680.985}{Version}%
\TextAt{8}{271.991}{680.985}{of}%
\TextAt{8}{285.321}{680.985}{Class}%
\TextAt{8}{316.01}{680.985}{Field}%
\TextAt{8}{345.755}{680.985}{Theory}%
\TextAt{8}{386.782}{680.985}{.}%
\TextAt{8}{396.011}{680.985}{.}%
\TextAt{8}{405.241}{680.985}{.}%
\TextAt{8}{414.47}{680.985}{.}%
\TextAt{8}{423.699}{680.985}{.}%
\TextAt{8}{432.929}{680.985}{.}%
\TextAt{8}{442.158}{680.985}{.}%
\TextAt{8}{451.387}{680.985}{.}%
\TextAt{8}{460.616}{680.985}{.}%
\TextAt{8}{469.846}{680.985}{.}%
\TextAt{8}{479.075}{680.985}{.}%
\TextAt{8}{488.304}{680.985}{.}%
\TextAt{8}{497.534}{680.985}{.}%
\TextAt{8}{506.763}{680.985}{.}%
\TextAt{8}{515.195}{680.985}{370}%
\TextAt{8}{97.159}{660.95}{§}%
\TextAt{8}{104.831}{660.95}{8}%
\TextAt{8}{124.656}{660.95}{The}%
\TextAt{8}{148.71}{660.95}{Reciprocity}%
\TextAt{8}{210.984}{660.95}{Law}%
\TextAt{8}{236.186}{660.95}{of}%
\TextAt{8}{249.516}{660.95}{the}%
\TextAt{8}{269.672}{660.95}{Power}%
\TextAt{8}{304.629}{660.95}{Residues}%
\TextAt{8}{359.094}{660.95}{.}%
\TextAt{8}{368.323}{660.95}{.}%
\TextAt{8}{377.553}{660.95}{.}%
\TextAt{8}{386.782}{660.95}{.}%
\TextAt{8}{396.011}{660.95}{.}%
\TextAt{8}{405.241}{660.95}{.}%
\TextAt{8}{414.47}{660.95}{.}%
\TextAt{8}{423.699}{660.95}{.}%
\TextAt{8}{432.929}{660.95}{.}%
\TextAt{8}{442.158}{660.95}{.}%
\TextAt{8}{451.387}{660.95}{.}%
\TextAt{8}{460.616}{660.95}{.}%
\TextAt{8}{469.846}{660.95}{.}%
\TextAt{8}{479.075}{660.95}{.}%
\TextAt{8}{488.304}{660.95}{.}%
\TextAt{8}{497.534}{660.95}{.}%
\TextAt{8}{506.763}{660.95}{.}%
\TextAt{8}{515.195}{660.95}{378}%
\TextAt{12}{79.226}{632.547}{VII}%
\TextAt{12}{109.114}{632.547}{Zeta}%
\TextAt{12}{139.731}{632.547}{Functions}%
\TextAt{12}{200.954}{632.547}{and}%
\TextAt{13}{226.92}{632.547}{L}%
\TextAt{12}{234.885}{632.547}{-series}%
\TextAt{12}{276.604}{632.547}{.}%
\TextAt{12}{286.323}{632.547}{.}%
\TextAt{12}{296.043}{632.547}{.}%
\TextAt{12}{305.762}{632.547}{.}%
\TextAt{12}{315.482}{632.547}{.}%
\TextAt{12}{325.201}{632.547}{.}%
\TextAt{12}{334.921}{632.547}{.}%
\TextAt{12}{344.64}{632.547}{.}%
\TextAt{12}{354.36}{632.547}{.}%
\TextAt{12}{364.079}{632.547}{.}%
\TextAt{12}{373.799}{632.547}{.}%
\TextAt{12}{383.518}{632.547}{.}%
\TextAt{12}{393.238}{632.547}{.}%
\TextAt{12}{402.957}{632.547}{.}%
\TextAt{12}{412.676}{632.547}{.}%
\TextAt{12}{422.396}{632.547}{.}%
\TextAt{12}{432.115}{632.547}{.}%
\TextAt{12}{441.835}{632.547}{.}%
\TextAt{12}{451.554}{632.547}{.}%
\TextAt{12}{461.274}{632.547}{.}%
\TextAt{12}{470.993}{632.547}{.}%
\TextAt{12}{480.713}{632.547}{.}%
\TextAt{12}{490.432}{632.547}{.}%
\TextAt{12}{500.152}{632.547}{.}%
\TextAt{8}{515.195}{632.547}{381}%
\TextAt{8}{97.159}{612.512}{§}%
\TextAt{8}{104.831}{612.512}{1}%
\TextAt{8}{124.656}{612.512}{The}%
\TextAt{8}{148.71}{612.512}{Riemann}%
\TextAt{8}{198.3}{612.512}{Zeta}%
\TextAt{8}{224.96}{612.512}{Function}%
\TextAt{8}{276.03}{612.512}{.}%
\TextAt{8}{285.26}{612.512}{.}%
\TextAt{8}{294.489}{612.512}{.}%
\TextAt{8}{303.718}{612.512}{.}%
\TextAt{8}{312.948}{612.512}{.}%
\TextAt{8}{322.177}{612.512}{.}%
\TextAt{8}{331.406}{612.512}{.}%
\TextAt{8}{340.635}{612.512}{.}%
\TextAt{8}{349.865}{612.512}{.}%
\TextAt{8}{359.094}{612.512}{.}%
\TextAt{8}{368.323}{612.512}{.}%
\TextAt{8}{377.553}{612.512}{.}%
\TextAt{8}{386.782}{612.512}{.}%
\TextAt{8}{396.011}{612.512}{.}%
\TextAt{8}{405.241}{612.512}{.}%
\TextAt{8}{414.47}{612.512}{.}%
\TextAt{8}{423.699}{612.512}{.}%
\TextAt{8}{432.929}{612.512}{.}%
\TextAt{8}{442.158}{612.512}{.}%
\TextAt{8}{451.387}{612.512}{.}%
\TextAt{8}{460.616}{612.512}{.}%
\TextAt{8}{469.846}{612.512}{.}%
\TextAt{8}{479.075}{612.512}{.}%
\TextAt{8}{488.304}{612.512}{.}%
\TextAt{8}{497.534}{612.512}{.}%
\TextAt{8}{506.763}{612.512}{.}%
\TextAt{8}{515.195}{612.512}{381}%
\TextAt{8}{97.159}{592.477}{§}%
\TextAt{8}{104.831}{592.477}{2}%
\TextAt{8}{124.656}{592.477}{Dirichlet}%
\TextAt{13}{172.943}{592.477}{L}%
\TextAt{8}{180.907}{592.477}{-series}%
\TextAt{8}{220.654}{592.477}{.}%
\TextAt{8}{229.884}{592.477}{.}%
\TextAt{8}{239.113}{592.477}{.}%
\TextAt{8}{248.342}{592.477}{.}%
\TextAt{8}{257.572}{592.477}{.}%
\TextAt{8}{266.801}{592.477}{.}%
\TextAt{8}{276.03}{592.477}{.}%
\TextAt{8}{285.26}{592.477}{.}%
\TextAt{8}{294.489}{592.477}{.}%
\TextAt{8}{303.718}{592.477}{.}%
\TextAt{8}{312.948}{592.477}{.}%
\TextAt{8}{322.177}{592.477}{.}%
\TextAt{8}{331.406}{592.477}{.}%
\TextAt{8}{340.635}{592.477}{.}%
\TextAt{8}{349.865}{592.477}{.}%
\TextAt{8}{359.094}{592.477}{.}%
\TextAt{8}{368.323}{592.477}{.}%
\TextAt{8}{377.553}{592.477}{.}%
\TextAt{8}{386.782}{592.477}{.}%
\TextAt{8}{396.011}{592.477}{.}%
\TextAt{8}{405.241}{592.477}{.}%
\TextAt{8}{414.47}{592.477}{.}%
\TextAt{8}{423.699}{592.477}{.}%
\TextAt{8}{432.929}{592.477}{.}%
\TextAt{8}{442.158}{592.477}{.}%
\TextAt{8}{451.387}{592.477}{.}%
\TextAt{8}{460.616}{592.477}{.}%
\TextAt{8}{469.846}{592.477}{.}%
\TextAt{8}{479.075}{592.477}{.}%
\TextAt{8}{488.304}{592.477}{.}%
\TextAt{8}{497.534}{592.477}{.}%
\TextAt{8}{506.763}{592.477}{.}%
\TextAt{8}{515.195}{592.477}{393}%
\TextAt{8}{97.159}{572.442}{§}%
\TextAt{8}{104.831}{572.442}{3}%
\TextAt{8}{124.656}{572.442}{Theta}%
\TextAt{8}{159.123}{572.442}{Series}%
\TextAt{8}{192.967}{572.442}{.}%
\TextAt{8}{202.196}{572.442}{.}%
\TextAt{8}{211.425}{572.442}{.}%
\TextAt{8}{220.654}{572.442}{.}%
\TextAt{8}{229.884}{572.442}{.}%
\TextAt{8}{239.113}{572.442}{.}%
\TextAt{8}{248.342}{572.442}{.}%
\TextAt{8}{257.572}{572.442}{.}%
\TextAt{8}{266.801}{572.442}{.}%
\TextAt{8}{276.03}{572.442}{.}%
\TextAt{8}{285.26}{572.442}{.}%
\TextAt{8}{294.489}{572.442}{.}%
\TextAt{8}{303.718}{572.442}{.}%
\TextAt{8}{312.948}{572.442}{.}%
\TextAt{8}{322.177}{572.442}{.}%
\TextAt{8}{331.406}{572.442}{.}%
\TextAt{8}{340.635}{572.442}{.}%
\TextAt{8}{349.865}{572.442}{.}%
\TextAt{8}{359.094}{572.442}{.}%
\TextAt{8}{368.323}{572.442}{.}%
\TextAt{8}{377.553}{572.442}{.}%
\TextAt{8}{386.782}{572.442}{.}%
\TextAt{8}{396.011}{572.442}{.}%
\TextAt{8}{405.241}{572.442}{.}%
\TextAt{8}{414.47}{572.442}{.}%
\TextAt{8}{423.699}{572.442}{.}%
\TextAt{8}{432.929}{572.442}{.}%
\TextAt{8}{442.158}{572.442}{.}%
\TextAt{8}{451.387}{572.442}{.}%
\TextAt{8}{460.616}{572.442}{.}%
\TextAt{8}{469.846}{572.442}{.}%
\TextAt{8}{479.075}{572.442}{.}%
\TextAt{8}{488.304}{572.442}{.}%
\TextAt{8}{497.534}{572.442}{.}%
\TextAt{8}{506.763}{572.442}{.}%
\TextAt{8}{515.195}{572.442}{401}%
\TextAt{8}{97.159}{552.407}{§}%
\TextAt{8}{104.831}{552.407}{4}%
\TextAt{8}{124.656}{552.407}{The}%
\TextAt{8}{148.71}{552.407}{Higher-dimensional}%
\TextAt{8}{251.201}{552.407}{Gamma}%
\TextAt{8}{295.507}{552.407}{Function}%
\TextAt{8}{349.865}{552.407}{.}%
\TextAt{8}{359.094}{552.407}{.}%
\TextAt{8}{368.323}{552.407}{.}%
\TextAt{8}{377.553}{552.407}{.}%
\TextAt{8}{386.782}{552.407}{.}%
\TextAt{8}{396.011}{552.407}{.}%
\TextAt{8}{405.241}{552.407}{.}%
\TextAt{8}{414.47}{552.407}{.}%
\TextAt{8}{423.699}{552.407}{.}%
\TextAt{8}{432.929}{552.407}{.}%
\TextAt{8}{442.158}{552.407}{.}%
\TextAt{8}{451.387}{552.407}{.}%
\TextAt{8}{460.616}{552.407}{.}%
\TextAt{8}{469.846}{552.407}{.}%
\TextAt{8}{479.075}{552.407}{.}%
\TextAt{8}{488.304}{552.407}{.}%
\TextAt{8}{497.534}{552.407}{.}%
\TextAt{8}{506.763}{552.407}{.}%
\TextAt{8}{515.195}{552.407}{410}%
\TextAt{8}{97.159}{532.373}{§}%
\TextAt{8}{104.831}{532.373}{5}%
\TextAt{8}{124.656}{532.373}{The}%
\TextAt{8}{148.71}{532.373}{Dedekind}%
\TextAt{8}{200.894}{532.373}{Zeta}%
\TextAt{8}{227.554}{532.373}{Function}%
\TextAt{8}{276.03}{532.373}{.}%
\TextAt{8}{285.26}{532.373}{.}%
\TextAt{8}{294.489}{532.373}{.}%
\TextAt{8}{303.718}{532.373}{.}%
\TextAt{8}{312.948}{532.373}{.}%
\TextAt{8}{322.177}{532.373}{.}%
\TextAt{8}{331.406}{532.373}{.}%
\TextAt{8}{340.635}{532.373}{.}%
\TextAt{8}{349.865}{532.373}{.}%
\TextAt{8}{359.094}{532.373}{.}%
\TextAt{8}{368.323}{532.373}{.}%
\TextAt{8}{377.553}{532.373}{.}%
\TextAt{8}{386.782}{532.373}{.}%
\TextAt{8}{396.011}{532.373}{.}%
\TextAt{8}{405.241}{532.373}{.}%
\TextAt{8}{414.47}{532.373}{.}%
\TextAt{8}{423.699}{532.373}{.}%
\TextAt{8}{432.929}{532.373}{.}%
\TextAt{8}{442.158}{532.373}{.}%
\TextAt{8}{451.387}{532.373}{.}%
\TextAt{8}{460.616}{532.373}{.}%
\TextAt{8}{469.846}{532.373}{.}%
\TextAt{8}{479.075}{532.373}{.}%
\TextAt{8}{488.304}{532.373}{.}%
\TextAt{8}{497.534}{532.373}{.}%
\TextAt{8}{506.763}{532.373}{.}%
\TextAt{8}{515.195}{532.373}{413}%
\TextAt{8}{97.159}{512.338}{§}%
\TextAt{8}{104.831}{512.338}{6}%
\TextAt{8}{124.656}{512.338}{Hecke}%
\TextAt{8}{158.465}{512.338}{Characters}%
\TextAt{8}{220.654}{512.338}{.}%
\TextAt{8}{229.884}{512.338}{.}%
\TextAt{8}{239.113}{512.338}{.}%
\TextAt{8}{248.342}{512.338}{.}%
\TextAt{8}{257.572}{512.338}{.}%
\TextAt{8}{266.801}{512.338}{.}%
\TextAt{8}{276.03}{512.338}{.}%
\TextAt{8}{285.26}{512.338}{.}%
\TextAt{8}{294.489}{512.338}{.}%
\TextAt{8}{303.718}{512.338}{.}%
\TextAt{8}{312.948}{512.338}{.}%
\TextAt{8}{322.177}{512.338}{.}%
\TextAt{8}{331.406}{512.338}{.}%
\TextAt{8}{340.635}{512.338}{.}%
\TextAt{8}{349.865}{512.338}{.}%
\TextAt{8}{359.094}{512.338}{.}%
\TextAt{8}{368.323}{512.338}{.}%
\TextAt{8}{377.553}{512.338}{.}%
\TextAt{8}{386.782}{512.338}{.}%
\TextAt{8}{396.011}{512.338}{.}%
\TextAt{8}{405.241}{512.338}{.}%
\TextAt{8}{414.47}{512.338}{.}%
\TextAt{8}{423.699}{512.338}{.}%
\TextAt{8}{432.929}{512.338}{.}%
\TextAt{8}{442.158}{512.338}{.}%
\TextAt{8}{451.387}{512.338}{.}%
\TextAt{8}{460.616}{512.338}{.}%
\TextAt{8}{469.846}{512.338}{.}%
\TextAt{8}{479.075}{512.338}{.}%
\TextAt{8}{488.304}{512.338}{.}%
\TextAt{8}{497.534}{512.338}{.}%
\TextAt{8}{506.763}{512.338}{.}%
\TextAt{8}{515.195}{512.338}{425}%
\TextAt{8}{97.159}{492.303}{§}%
\TextAt{8}{104.831}{492.303}{7}%
\TextAt{8}{124.656}{492.303}{Theta}%
\TextAt{8}{159.123}{492.303}{Series}%
\TextAt{8}{192.346}{492.303}{of}%
\TextAt{8}{205.676}{492.303}{Algebraic}%
\TextAt{8}{258.028}{492.303}{Number}%
\TextAt{8}{303.218}{492.303}{Fields}%
\TextAt{8}{340.635}{492.303}{.}%
\TextAt{8}{349.865}{492.303}{.}%
\TextAt{8}{359.094}{492.303}{.}%
\TextAt{8}{368.323}{492.303}{.}%
\TextAt{8}{377.553}{492.303}{.}%
\TextAt{8}{386.782}{492.303}{.}%
\TextAt{8}{396.011}{492.303}{.}%
\TextAt{8}{405.241}{492.303}{.}%
\TextAt{8}{414.47}{492.303}{.}%
\TextAt{8}{423.699}{492.303}{.}%
\TextAt{8}{432.929}{492.303}{.}%
\TextAt{8}{442.158}{492.303}{.}%
\TextAt{8}{451.387}{492.303}{.}%
\TextAt{8}{460.616}{492.303}{.}%
\TextAt{8}{469.846}{492.303}{.}%
\TextAt{8}{479.075}{492.303}{.}%
\TextAt{8}{488.304}{492.303}{.}%
\TextAt{8}{497.534}{492.303}{.}%
\TextAt{8}{506.763}{492.303}{.}%
\TextAt{8}{515.195}{492.303}{438}%
\TextAt{8}{97.159}{472.268}{§}%
\TextAt{8}{104.831}{472.268}{8}%
\TextAt{8}{124.656}{472.268}{Hecke}%
\TextAt{13}{158.465}{472.268}{L}%
\TextAt{8}{166.43}{472.268}{-series}%
\TextAt{8}{202.196}{472.268}{.}%
\TextAt{8}{211.425}{472.268}{.}%
\TextAt{8}{220.654}{472.268}{.}%
\TextAt{8}{229.884}{472.268}{.}%
\TextAt{8}{239.113}{472.268}{.}%
\TextAt{8}{248.342}{472.268}{.}%
\TextAt{8}{257.572}{472.268}{.}%
\TextAt{8}{266.801}{472.268}{.}%
\TextAt{8}{276.03}{472.268}{.}%
\TextAt{8}{285.26}{472.268}{.}%
\TextAt{8}{294.489}{472.268}{.}%
\TextAt{8}{303.718}{472.268}{.}%
\TextAt{8}{312.948}{472.268}{.}%
\TextAt{8}{322.177}{472.268}{.}%
\TextAt{8}{331.406}{472.268}{.}%
\TextAt{8}{340.635}{472.268}{.}%
\TextAt{8}{349.865}{472.268}{.}%
\TextAt{8}{359.094}{472.268}{.}%
\TextAt{8}{368.323}{472.268}{.}%
\TextAt{8}{377.553}{472.268}{.}%
\TextAt{8}{386.782}{472.268}{.}%
\TextAt{8}{396.011}{472.268}{.}%
\TextAt{8}{405.241}{472.268}{.}%
\TextAt{8}{414.47}{472.268}{.}%
\TextAt{8}{423.699}{472.268}{.}%
\TextAt{8}{432.929}{472.268}{.}%
\TextAt{8}{442.158}{472.268}{.}%
\TextAt{8}{451.387}{472.268}{.}%
\TextAt{8}{460.616}{472.268}{.}%
\TextAt{8}{469.846}{472.268}{.}%
\TextAt{8}{479.075}{472.268}{.}%
\TextAt{8}{488.304}{472.268}{.}%
\TextAt{8}{497.534}{472.268}{.}%
\TextAt{8}{506.763}{472.268}{.}%
\TextAt{8}{515.195}{472.268}{446}%
\TextAt{8}{97.159}{452.234}{§}%
\TextAt{8}{104.831}{452.234}{9}%
\TextAt{8}{124.656}{452.234}{Values}%
\TextAt{8}{161.777}{452.234}{of}%
\TextAt{8}{175.107}{452.234}{Dirichlet}%
\TextAt{13}{223.394}{452.234}{L}%
\TextAt{8}{231.358}{452.234}{-series}%
\TextAt{8}{266.59}{452.234}{at}%
\TextAt{8}{280.9}{452.234}{Integer}%
\TextAt{8}{320.568}{452.234}{Points}%
\TextAt{8}{359.094}{452.234}{.}%
\TextAt{8}{368.323}{452.234}{.}%
\TextAt{8}{377.553}{452.234}{.}%
\TextAt{8}{386.782}{452.234}{.}%
\TextAt{8}{396.011}{452.234}{.}%
\TextAt{8}{405.241}{452.234}{.}%
\TextAt{8}{414.47}{452.234}{.}%
\TextAt{8}{423.699}{452.234}{.}%
\TextAt{8}{432.929}{452.234}{.}%
\TextAt{8}{442.158}{452.234}{.}%
\TextAt{8}{451.387}{452.234}{.}%
\TextAt{8}{460.616}{452.234}{.}%
\TextAt{8}{469.846}{452.234}{.}%
\TextAt{8}{479.075}{452.234}{.}%
\TextAt{8}{488.304}{452.234}{.}%
\TextAt{8}{497.534}{452.234}{.}%
\TextAt{8}{506.763}{452.234}{.}%
\TextAt{8}{515.195}{452.234}{456}%
\TextAt{8}{97.159}{432.199}{§}%
\TextAt{8}{104.831}{432.199}{10}%
\TextAt{8}{124.656}{432.199}{Artin}%
\TextAt{13}{156.194}{432.199}{L}%
\TextAt{8}{164.158}{432.199}{-series}%
\TextAt{8}{202.196}{432.199}{.}%
\TextAt{8}{211.425}{432.199}{.}%
\TextAt{8}{220.654}{432.199}{.}%
\TextAt{8}{229.884}{432.199}{.}%
\TextAt{8}{239.113}{432.199}{.}%
\TextAt{8}{248.342}{432.199}{.}%
\TextAt{8}{257.572}{432.199}{.}%
\TextAt{8}{266.801}{432.199}{.}%
\TextAt{8}{276.03}{432.199}{.}%
\TextAt{8}{285.26}{432.199}{.}%
\TextAt{8}{294.489}{432.199}{.}%
\TextAt{8}{303.718}{432.199}{.}%
\TextAt{8}{312.948}{432.199}{.}%
\TextAt{8}{322.177}{432.199}{.}%
\TextAt{8}{331.406}{432.199}{.}%
\TextAt{8}{340.635}{432.199}{.}%
\TextAt{8}{349.865}{432.199}{.}%
\TextAt{8}{359.094}{432.199}{.}%
\TextAt{8}{368.323}{432.199}{.}%
\TextAt{8}{377.553}{432.199}{.}%
\TextAt{8}{386.782}{432.199}{.}%
\TextAt{8}{396.011}{432.199}{.}%
\TextAt{8}{405.241}{432.199}{.}%
\TextAt{8}{414.47}{432.199}{.}%
\TextAt{8}{423.699}{432.199}{.}%
\TextAt{8}{432.929}{432.199}{.}%
\TextAt{8}{442.158}{432.199}{.}%
\TextAt{8}{451.387}{432.199}{.}%
\TextAt{8}{460.616}{432.199}{.}%
\TextAt{8}{469.846}{432.199}{.}%
\TextAt{8}{479.075}{432.199}{.}%
\TextAt{8}{488.304}{432.199}{.}%
\TextAt{8}{497.534}{432.199}{.}%
\TextAt{8}{506.763}{432.199}{.}%
\TextAt{8}{515.195}{432.199}{468}%
\TextAt{8}{97.159}{412.164}{§}%
\TextAt{8}{104.831}{412.164}{11}%
\TextAt{8}{124.656}{412.164}{The}%
\TextAt{8}{148.71}{412.164}{Artin}%
\TextAt{8}{180.247}{412.164}{Conductor}%
\TextAt{8}{239.113}{412.164}{.}%
\TextAt{8}{248.342}{412.164}{.}%
\TextAt{8}{257.572}{412.164}{.}%
\TextAt{8}{266.801}{412.164}{.}%
\TextAt{8}{276.03}{412.164}{.}%
\TextAt{8}{285.26}{412.164}{.}%
\TextAt{8}{294.489}{412.164}{.}%
\TextAt{8}{303.718}{412.164}{.}%
\TextAt{8}{312.948}{412.164}{.}%
\TextAt{8}{322.177}{412.164}{.}%
\TextAt{8}{331.406}{412.164}{.}%
\TextAt{8}{340.635}{412.164}{.}%
\TextAt{8}{349.865}{412.164}{.}%
\TextAt{8}{359.094}{412.164}{.}%
\TextAt{8}{368.323}{412.164}{.}%
\TextAt{8}{377.553}{412.164}{.}%
\TextAt{8}{386.782}{412.164}{.}%
\TextAt{8}{396.011}{412.164}{.}%
\TextAt{8}{405.241}{412.164}{.}%
\TextAt{8}{414.47}{412.164}{.}%
\TextAt{8}{423.699}{412.164}{.}%
\TextAt{8}{432.929}{412.164}{.}%
\TextAt{8}{442.158}{412.164}{.}%
\TextAt{8}{451.387}{412.164}{.}%
\TextAt{8}{460.616}{412.164}{.}%
\TextAt{8}{469.846}{412.164}{.}%
\TextAt{8}{479.075}{412.164}{.}%
\TextAt{8}{488.304}{412.164}{.}%
\TextAt{8}{497.534}{412.164}{.}%
\TextAt{8}{506.763}{412.164}{.}%
\TextAt{8}{515.195}{412.164}{477}%
\TextAt{8}{97.159}{392.129}{§}%
\TextAt{8}{104.831}{392.129}{12}%
\TextAt{8}{124.656}{392.129}{The}%
\TextAt{8}{148.71}{392.129}{Functional}%
\TextAt{8}{206.752}{392.129}{Equation}%
\TextAt{8}{257.322}{392.129}{of}%
\TextAt{8}{270.652}{392.129}{Artin}%
\TextAt{13}{302.19}{392.129}{L}%
\TextAt{8}{310.155}{392.129}{-series}%
\TextAt{8}{349.865}{392.129}{.}%
\TextAt{8}{359.094}{392.129}{.}%
\TextAt{8}{368.323}{392.129}{.}%
\TextAt{8}{377.553}{392.129}{.}%
\TextAt{8}{386.782}{392.129}{.}%
\TextAt{8}{396.011}{392.129}{.}%
\TextAt{8}{405.241}{392.129}{.}%
\TextAt{8}{414.47}{392.129}{.}%
\TextAt{8}{423.699}{392.129}{.}%
\TextAt{8}{432.929}{392.129}{.}%
\TextAt{8}{442.158}{392.129}{.}%
\TextAt{8}{451.387}{392.129}{.}%
\TextAt{8}{460.616}{392.129}{.}%
\TextAt{8}{469.846}{392.129}{.}%
\TextAt{8}{479.075}{392.129}{.}%
\TextAt{8}{488.304}{392.129}{.}%
\TextAt{8}{497.534}{392.129}{.}%
\TextAt{8}{506.763}{392.129}{.}%
\TextAt{8}{515.195}{392.129}{485}%
\TextAt{8}{97.159}{372.095}{§}%
\TextAt{8}{104.831}{372.095}{13}%
\TextAt{8}{124.656}{372.095}{Density}%
\TextAt{8}{167.479}{372.095}{Theorems}%
\TextAt{8}{220.654}{372.095}{.}%
\TextAt{8}{229.884}{372.095}{.}%
\TextAt{8}{239.113}{372.095}{.}%
\TextAt{8}{248.342}{372.095}{.}%
\TextAt{8}{257.572}{372.095}{.}%
\TextAt{8}{266.801}{372.095}{.}%
\TextAt{8}{276.03}{372.095}{.}%
\TextAt{8}{285.26}{372.095}{.}%
\TextAt{8}{294.489}{372.095}{.}%
\TextAt{8}{303.718}{372.095}{.}%
\TextAt{8}{312.948}{372.095}{.}%
\TextAt{8}{322.177}{372.095}{.}%
\TextAt{8}{331.406}{372.095}{.}%
\TextAt{8}{340.635}{372.095}{.}%
\TextAt{8}{349.865}{372.095}{.}%
\TextAt{8}{359.094}{372.095}{.}%
\TextAt{8}{368.323}{372.095}{.}%
\TextAt{8}{377.553}{372.095}{.}%
\TextAt{8}{386.782}{372.095}{.}%
\TextAt{8}{396.011}{372.095}{.}%
\TextAt{8}{405.241}{372.095}{.}%
\TextAt{8}{414.47}{372.095}{.}%
\TextAt{8}{423.699}{372.095}{.}%
\TextAt{8}{432.929}{372.095}{.}%
\TextAt{8}{442.158}{372.095}{.}%
\TextAt{8}{451.387}{372.095}{.}%
\TextAt{8}{460.616}{372.095}{.}%
\TextAt{8}{469.846}{372.095}{.}%
\TextAt{8}{479.075}{372.095}{.}%
\TextAt{8}{488.304}{372.095}{.}%
\TextAt{8}{497.534}{372.095}{.}%
\TextAt{8}{506.763}{372.095}{.}%
\TextAt{8}{515.195}{372.095}{491}%
\TextAt{12}{79.226}{343.691}{Bibliography}%
\TextAt{12}{159.971}{343.691}{.}%
\TextAt{12}{169.69}{343.691}{.}%
\TextAt{12}{179.409}{343.691}{.}%
\TextAt{12}{189.129}{343.691}{.}%
\TextAt{12}{198.848}{343.691}{.}%
\TextAt{12}{208.568}{343.691}{.}%
\TextAt{12}{218.287}{343.691}{.}%
\TextAt{12}{228.007}{343.691}{.}%
\TextAt{12}{237.726}{343.691}{.}%
\TextAt{12}{247.446}{343.691}{.}%
\TextAt{12}{257.165}{343.691}{.}%
\TextAt{12}{266.885}{343.691}{.}%
\TextAt{12}{276.604}{343.691}{.}%
\TextAt{12}{286.323}{343.691}{.}%
\TextAt{12}{296.043}{343.691}{.}%
\TextAt{12}{305.762}{343.691}{.}%
\TextAt{12}{315.482}{343.691}{.}%
\TextAt{12}{325.201}{343.691}{.}%
\TextAt{12}{334.921}{343.691}{.}%
\TextAt{12}{344.64}{343.691}{.}%
\TextAt{12}{354.36}{343.691}{.}%
\TextAt{12}{364.079}{343.691}{.}%
\TextAt{12}{373.799}{343.691}{.}%
\TextAt{12}{383.518}{343.691}{.}%
\TextAt{12}{393.238}{343.691}{.}%
\TextAt{12}{402.957}{343.691}{.}%
\TextAt{12}{412.676}{343.691}{.}%
\TextAt{12}{422.396}{343.691}{.}%
\TextAt{12}{432.115}{343.691}{.}%
\TextAt{12}{441.835}{343.691}{.}%
\TextAt{12}{451.554}{343.691}{.}%
\TextAt{12}{461.274}{343.691}{.}%
\TextAt{12}{470.993}{343.691}{.}%
\TextAt{12}{480.713}{343.691}{.}%
\TextAt{12}{490.432}{343.691}{.}%
\TextAt{12}{500.152}{343.691}{.}%
\TextAt{8}{515.195}{343.691}{498}%
\TextAt{12}{79.226}{315.288}{Index}%
\TextAt{12}{121.093}{315.288}{.}%
\TextAt{12}{130.812}{315.288}{.}%
\TextAt{12}{140.532}{315.288}{.}%
\TextAt{12}{150.251}{315.288}{.}%
\TextAt{12}{159.971}{315.288}{.}%
\TextAt{12}{169.69}{315.288}{.}%
\TextAt{12}{179.409}{315.288}{.}%
\TextAt{12}{189.129}{315.288}{.}%
\TextAt{12}{198.848}{315.288}{.}%
\TextAt{12}{208.568}{315.288}{.}%
\TextAt{12}{218.287}{315.288}{.}%
\TextAt{12}{228.007}{315.288}{.}%
\TextAt{12}{237.726}{315.288}{.}%
\TextAt{12}{247.446}{315.288}{.}%
\TextAt{12}{257.165}{315.288}{.}%
\TextAt{12}{266.885}{315.288}{.}%
\TextAt{12}{276.604}{315.288}{.}%
\TextAt{12}{286.323}{315.288}{.}%
\TextAt{12}{296.043}{315.288}{.}%
\TextAt{12}{305.762}{315.288}{.}%
\TextAt{12}{315.482}{315.288}{.}%
\TextAt{12}{325.201}{315.288}{.}%
\TextAt{12}{334.921}{315.288}{.}%
\TextAt{12}{344.64}{315.288}{.}%
\TextAt{12}{354.36}{315.288}{.}%
\TextAt{12}{364.079}{315.288}{.}%
\TextAt{12}{373.799}{315.288}{.}%
\TextAt{12}{383.518}{315.288}{.}%
\TextAt{12}{393.238}{315.288}{.}%
\TextAt{12}{402.957}{315.288}{.}%
\TextAt{12}{412.676}{315.288}{.}%
\TextAt{12}{422.396}{315.288}{.}%
\TextAt{12}{432.115}{315.288}{.}%
\TextAt{12}{441.835}{315.288}{.}%
\TextAt{12}{451.554}{315.288}{.}%
\TextAt{12}{461.274}{315.288}{.}%
\TextAt{12}{470.993}{315.288}{.}%
\TextAt{12}{480.713}{315.288}{.}%
\TextAt{12}{490.432}{315.288}{.}%
\TextAt{12}{500.152}{315.288}{.}%
\TextAt{8}{515.195}{315.288}{507}%
\TextAt{8}{299.656}{43.033}{xii}%
\Destination{78.23}{809.815}{nav.page.xii}%
\Hotlink{97.159}{757.776}{278.065}{16.821}{nav.section.6.3}%
\Hotlink{97.159}{737.741}{145.805}{16.821}{nav.section.6.4}%
\Hotlink{97.159}{717.706}{172.931}{16.821}{nav.section.6.5}%
\Hotlink{97.159}{697.672}{126.45}{16.821}{nav.section.6.6}%
\Hotlink{97.159}{677.637}{285.346}{16.821}{nav.section.6.7}%
\Hotlink{97.159}{657.602}{251.979}{16.821}{nav.section.6.8}%
\Hotlink{79.226}{629.091}{192.265}{17.048}{nav.chapter.7}%
\Hotlink{97.159}{609.164}{172.836}{16.821}{nav.section.7.1}%
\Hotlink{97.159}{589.129}{115.083}{16.821}{nav.section.7.2}%
\Hotlink{97.159}{569.095}{91.29}{16.82}{nav.section.7.3}%
\Hotlink{97.159}{549.06}{243.383}{16.821}{nav.section.7.4}%
\Hotlink{97.159}{529.025}{175.43}{16.821}{nav.section.7.5}%
\Hotlink{97.159}{508.99}{116.658}{16.821}{nav.section.7.6}%
\Hotlink{97.159}{488.956}{236.521}{16.82}{nav.section.7.7}%
\Hotlink{97.159}{468.921}{100.605}{16.821}{nav.section.7.8}%
\Hotlink{97.159}{448.886}{255.508}{16.821}{nav.section.7.9}%
\Hotlink{97.159}{428.851}{98.334}{16.821}{nav.section.7.10}%
\Hotlink{97.159}{408.816}{137.078}{16.821}{nav.section.7.11}%
\Hotlink{97.159}{388.782}{244.33}{16.821}{nav.section.7.12}%
\Hotlink{97.159}{368.747}{120.46}{16.821}{nav.section.7.13}%
\Hotlink{79.226}{311.833}{33.188}{17.048}{nav.chapterz2az.108}%
\Hotlink{79.226}{340.238}{454.461}{11.047}{nav.chapterz2az.107}%
}


% Chapter 1
% Page 13 (body page 1)
\Page{%
\TextAt{1}{79.226}{673.054}{Chapter}%
\TextAt{1}{188.19}{673.054}{I}%
\TextAt{1}{79.226}{614.873}{Algebraic}%
\TextAt{1}{204.624}{614.873}{Integers}%
\TextAt{15}{79.226}{546.17}{§}%
\TextAt{15}{93.785}{546.17}{1}%
\TextAt{15}{120.693}{546.17}{The}%
\TextAt{15}{160.202}{546.17}{Gaussian}%
\TextAt{15}{242.905}{546.17}{Integers}%
\TextAt{8}{79.226}{514.284}{The}%
\TextAt{8}{103.28}{514.284}{equations}%
\TextAt{9}{99.344}{487.721}{2}%
\TextAt{9}{108.502}{487.721}{\Character{61}}%
\TextAt{9}{120.923}{487.721}{1}%
\TextAt{9}{129.423}{487.721}{\Character{43}}%
\TextAt{9}{141.187}{487.721}{1}%
\TextAt{13}{147.068}{487.721}{\Character{59}}%
\TextAt{9}{164.267}{487.721}{5}%
\TextAt{9}{173.425}{487.721}{\Character{61}}%
\TextAt{9}{185.846}{487.721}{1}%
\TextAt{9}{194.346}{487.721}{\Character{43}}%
\TextAt{9}{206.098}{487.721}{4}%
\TextAt{13}{211.99}{487.721}{\Character{59}}%
\TextAt{9}{229.189}{487.721}{13}%
\TextAt{9}{244.193}{487.721}{\Character{61}}%
\TextAt{9}{256.614}{487.721}{4}%
\TextAt{9}{265.114}{487.721}{\Character{43}}%
\TextAt{9}{276.878}{487.721}{9}%
\TextAt{13}{282.766}{487.721}{\Character{59}}%
\TextAt{9}{299.965}{487.721}{17}%
\TextAt{9}{314.969}{487.721}{\Character{61}}%
\TextAt{9}{327.39}{487.721}{1}%
\TextAt{9}{335.89}{487.721}{\Character{43}}%
\TextAt{9}{347.642}{487.721}{16}%
\TextAt{13}{359.394}{487.721}{\Character{59}}%
\TextAt{9}{376.594}{487.721}{29}%
\TextAt{9}{391.598}{487.721}{\Character{61}}%
\TextAt{9}{404.019}{487.721}{4}%
\TextAt{9}{412.519}{487.721}{\Character{43}}%
\TextAt{9}{424.271}{487.721}{25}%
\TextAt{13}{436.023}{487.721}{\Character{59}}%
\TextAt{9}{453.222}{487.721}{37}%
\TextAt{9}{468.226}{487.721}{\Character{61}}%
\TextAt{9}{480.647}{487.721}{1}%
\TextAt{9}{489.147}{487.721}{\Character{43}}%
\TextAt{9}{500.911}{487.721}{36}%
\TextAt{8}{79.226}{461.157}{show}%
\TextAt{8}{108.761}{461.157}{the}%
\TextAt{8}{129.45}{461.157}{first}%
\TextAt{8}{154.107}{461.157}{prime}%
\TextAt{8}{187.803}{461.157}{numbers}%
\TextAt{8}{235.545}{461.157}{that}%
\TextAt{8}{261.446}{461.157}{can}%
\TextAt{8}{283.438}{461.157}{be}%
\TextAt{8}{299.894}{461.157}{represented}%
\TextAt{8}{362.593}{461.157}{as}%
\TextAt{8}{377.495}{461.157}{a}%
\TextAt{8}{387.782}{461.157}{sum}%
\TextAt{8}{413.085}{461.157}{of}%
\TextAt{8}{426.947}{461.157}{two}%
\TextAt{8}{449.596}{461.157}{squares.}%
\TextAt{8}{497.167}{461.157}{Except}%
\TextAt{8}{79.226}{444.111}{for}%
\TextAt{9}{97.081}{444.111}{2}%
\TextAt{8}{102.934}{444.111}{,}%
\TextAt{8}{110.06}{444.111}{they}%
\TextAt{8}{136.367}{444.111}{are}%
\TextAt{8}{155.848}{444.111}{all}%
\NativeMathAt{172.077}{444.111}{11.955}{\equiv}%
\TextAt{9}{184.697}{444.112}{1}%
\TextAt{8}{193.871}{444.111}{mod}%
\TextAt{9}{219.309}{444.112}{4}%
\TextAt{8}{225.162}{444.111}{,}%
\TextAt{8}{232.287}{444.111}{and}%
\TextAt{8}{255.02}{444.111}{it}%
\TextAt{8}{266.694}{444.111}{is}%
\TextAt{8}{278.428}{444.111}{true}%
\TextAt{8}{303.11}{444.111}{in}%
\TextAt{8}{316.733}{444.111}{general}%
\TextAt{8}{357.028}{444.111}{that}%
\TextAt{8}{382.367}{444.111}{any}%
\TextAt{8}{404.454}{444.111}{odd}%
\TextAt{8}{427.51}{444.111}{prime}%
\TextAt{8}{460.644}{444.111}{number}%
\TextAt{8}{503.21}{444.111}{of}%
\TextAt{8}{516.51}{444.111}{the}%
\TextAt{8}{79.226}{427.066}{form}%
\TextAt{13}{107.626}{427.066}{p}%
\TextAt{9}{118.115}{427.066}{\Character{61}}%
\TextAt{13}{131.834}{427.066}{a}%
\TextAt{17}{137.979}{431.404}{2}%
\TextAt{9}{145.886}{427.066}{\Character{43}}%
\TextAt{13}{158.164}{427.066}{b}%
\TextAt{17}{163.141}{431.404}{2}%
\TextAt{8}{172.53}{427.066}{satisfies}%
\TextAt{13}{216.4}{427.066}{p}%
\NativeMathAt{226.889}{427.066}{11.955}{\equiv}%
\TextAt{9}{240.802}{427.066}{1}%
\TextAt{8}{249.976}{427.066}{mod}%
\TextAt{9}{275.414}{427.066}{4}%
\TextAt{8}{281.267}{427.066}{,}%
\TextAt{8}{289.365}{427.066}{because}%
\TextAt{8}{333.426}{427.066}{perfect}%
\TextAt{8}{373.195}{427.066}{squares}%
\TextAt{8}{415.378}{427.066}{are}%
\NativeMathAt{435.648}{427.066}{11.955}{\equiv}%
\TextAt{9}{449.561}{427.066}{0}%
\TextAt{8}{460.071}{427.066}{or}%
\NativeMathAt{475.14}{427.066}{11.955}{\equiv}%
\TextAt{9}{489.053}{427.066}{1}%
\TextAt{8}{498.227}{427.066}{mod}%
\TextAt{9}{523.665}{427.066}{4}%
\TextAt{8}{529.518}{427.066}{.}%
\TextAt{8}{79.226}{410.02}{This}%
\TextAt{8}{105.878}{410.02}{is}%
\TextAt{8}{117.574}{410.02}{obvious.}%
\TextAt{8}{164.447}{410.02}{What}%
\TextAt{8}{197.22}{410.02}{is}%
\TextAt{8}{208.916}{410.02}{not}%
\TextAt{8}{229.661}{410.02}{obvious}%
\TextAt{8}{271.938}{410.02}{is}%
\TextAt{8}{283.634}{410.02}{the}%
\TextAt{8}{303.723}{410.02}{remarkable}%
\TextAt{8}{364.47}{410.02}{fact}%
\TextAt{8}{387.488}{410.02}{that}%
\TextAt{8}{412.788}{410.02}{the}%
\TextAt{8}{432.877}{410.02}{converse}%
\TextAt{8}{479.374}{410.02}{also}%
\TextAt{8}{502.786}{410.02}{holds:}%
\TextAt{12}{79.226}{384.704}{Theorem}%
\TextAt{12}{136.24}{384.704}{(1.1).}%
\TextAt{8}{171.556}{384.704}{For}%
\TextAt{8}{192.525}{384.704}{all}%
\TextAt{8}{208.784}{384.704}{prime}%
\TextAt{8}{241.948}{384.704}{numbers}%
\TextAt{13}{289.159}{384.704}{p}%
\NativeMathAt{298.355}{384.704}{11.955}{\neq}%
\TextAt{9}{310.764}{384.704}{2}%
\TextAt{8}{316.633}{384.704}{,}%
\TextAt{8}{323.783}{384.704}{one}%
\TextAt{8}{345.242}{384.704}{has:}%
\TextAt{13}{188.68}{358.14}{p}%
\TextAt{9}{197.876}{358.14}{\Character{61}}%
\TextAt{13}{210.301}{358.14}{a}%
\TextAt{17}{216.446}{363.076}{2}%
\TextAt{9}{223.835}{358.14}{\Character{43}}%
\TextAt{13}{235.597}{358.14}{b}%
\TextAt{17}{240.574}{363.076}{2}%
\TextAt{9}{257.261}{358.14}{\Character{40}}%
\TextAt{13}{261.814}{358.14}{a\Character{59}}%
\TextAt{13}{273.183}{358.14}{b}%
\NativeMathAt{281.501}{358.14}{11.955}{\in}%
\TextAt{18}{292.792}{358.14}{Z}%
\TextAt{9}{300.762}{358.14}{\Character{41}}%
\NativeArrowAt{320.59}{358.14}{11.955}{21.926}{\Longleftrightarrow}%
\TextAt{13}{357.784}{358.14}{p}%
\NativeMathAt{366.98}{358.14}{11.955}{\equiv}%
\TextAt{9}{379.599}{358.14}{1}%
\TextAt{8}{388.773}{358.14}{mod}%
\TextAt{9}{414.211}{358.14}{4}%
\TextAt{13}{420.064}{358.14}{\Character{58}}%
\TextAt{8}{97.159}{323.307}{The}%
\TextAt{8}{121.109}{323.307}{natural}%
\TextAt{8}{161.988}{323.307}{explanation}%
\TextAt{8}{225.306}{323.307}{of}%
\TextAt{8}{238.533}{323.307}{this}%
\TextAt{8}{261.252}{323.307}{arithmetic}%
\TextAt{8}{317.732}{323.307}{law}%
\TextAt{8}{338.765}{323.307}{concerning}%
\TextAt{8}{397.194}{323.307}{the}%
\TextAt{8}{417.247}{323.307}{ring}%
\TextAt{18}{441.209}{323.307}{Z}%
\TextAt{8}{452.973}{323.307}{of}%
\TextAt{8}{466.199}{323.307}{rational}%
\TextAt{8}{509.684}{323.307}{inte-}%
\TextAt{8}{79.226}{306.261}{gers}%
\TextAt{8}{103.352}{306.261}{is}%
\TextAt{8}{115.116}{306.261}{found}%
\TextAt{8}{147.956}{306.261}{in}%
\TextAt{8}{161.609}{306.261}{the}%
\TextAt{8}{181.766}{306.261}{larger}%
\TextAt{8}{214.941}{306.261}{domain}%
\TextAt{8}{256.569}{306.261}{of}%
\TextAt{8}{269.899}{306.261}{the}%
\TextAt{12}{290.056}{306.261}{gaussian}%
\TextAt{12}{343.651}{306.261}{integers}%
\TextAt{18}{205.062}{279.697}{Z}%
\TextAt{9}{213.032}{279.697}{\Character{91}}%
\TextAt{13}{216.283}{279.697}{i}%
\TextAt{9}{220.277}{279.697}{\Character{93}}%
\TextAt{9}{226.828}{279.697}{\Character{61}}%
\TextAt{16}{239.275}{279.697}{\Character{102}}%
\TextAt{13}{245.252}{279.697}{a}%
\TextAt{9}{254.054}{279.697}{\Character{43}}%
\TextAt{13}{265.815}{279.697}{bi}%
\TextAt{16}{278.107}{279.697}{\Character{106}}%
\TextAt{13}{284.748}{279.697}{a\Character{59}}%
\TextAt{13}{296.117}{279.697}{b}%
\NativeMathAt{304.435}{279.697}{11.955}{\in}%
\TextAt{18}{315.726}{279.697}{Z}%
\TextAt{16}{323.697}{279.697}{\Character{103}}%
\TextAt{13}{329.674}{279.697}{\Character{59}}%
\TextAt{13}{358.809}{279.697}{i}%
\TextAt{9}{366.143}{279.697}{\Character{61}}%
\TextAt{16}{378.568}{289.672}{\Character{112}}%
\Vector{q 0 G 0 g 0.478 w 316.531 -484.088 m 331.683 -484.088 l S Q}%
\TextAt{16}{388.531}{279.697}{\Character{0}}%
\TextAt{9}{397.83}{279.697}{1}%
\TextAt{13}{403.683}{279.697}{\Character{58}}%
\TextAt{8}{79.226}{253.134}{In}%
\TextAt{8}{93.847}{253.134}{this}%
\TextAt{8}{116.67}{253.134}{ring,}%
\TextAt{8}{143.987}{253.134}{the}%
\TextAt{8}{164.144}{253.134}{equation}%
\TextAt{13}{211.952}{253.134}{p}%
\TextAt{9}{221.148}{253.134}{\Character{61}}%
\TextAt{13}{233.574}{253.134}{x}%
\TextAt{17}{240.226}{257.472}{2}%
\TextAt{9}{247.615}{253.134}{\Character{43}}%
\TextAt{13}{259.376}{253.134}{y}%
\TextAt{17}{265.513}{257.472}{2}%
\TextAt{8}{274.143}{253.134}{turns}%
\TextAt{8}{304.772}{253.134}{into}%
\TextAt{8}{328.515}{253.134}{the}%
\TextAt{8}{348.671}{253.134}{product}%
\TextAt{8}{392.571}{253.134}{decomposition}%
\TextAt{13}{253.16}{226.57}{p}%
\TextAt{9}{262.356}{226.57}{\Character{61}}%
\TextAt{9}{274.766}{226.57}{\Character{40}}%
\TextAt{13}{279.333}{226.57}{x}%
\TextAt{9}{288.642}{226.57}{\Character{43}}%
\TextAt{13}{300.403}{226.57}{iy}%
\TextAt{9}{310.533}{226.57}{\Character{41}\Character{40}}%
\TextAt{13}{319.638}{226.57}{x}%
\TextAt{16}{328.947}{226.57}{\Character{0}}%
\TextAt{13}{340.902}{226.57}{iy}%
\TextAt{9}{351.032}{226.57}{\Character{41}}%
\TextAt{13}{355.584}{226.57}{\Character{59}}%
\TextAt{8}{79.226}{200.007}{so}%
\TextAt{8}{94.416}{200.007}{that}%
\TextAt{8}{120.605}{200.007}{the}%
\TextAt{8}{141.581}{200.007}{problem}%
\TextAt{8}{187.927}{200.007}{is}%
\TextAt{8}{200.511}{200.007}{now}%
\TextAt{8}{225.719}{200.007}{when}%
\TextAt{8}{257.097}{200.007}{and}%
\TextAt{8}{280.679}{200.007}{how}%
\TextAt{8}{305.888}{200.007}{a}%
\TextAt{8}{316.463}{200.007}{prime}%
\TextAt{8}{350.447}{200.007}{number}%
\TextAt{13}{393.863}{200.007}{p}%
\NativeMathAt{404.456}{200.007}{11.955}{\in}%
\TextAt{18}{417.145}{200.007}{Z}%
\TextAt{8}{429.832}{200.007}{factors}%
\TextAt{8}{468.765}{200.007}{in}%
\TextAt{18}{483.238}{200.007}{Z}%
\TextAt{9}{491.208}{200.007}{\Character{91}}%
\TextAt{13}{494.46}{200.007}{i}%
\TextAt{9}{498.453}{200.007}{\Character{93}}%
\TextAt{8}{501.705}{200.007}{.}%
\TextAt{8}{512.613}{200.007}{The}%
\TextAt{8}{79.226}{182.961}{answer}%
\TextAt{8}{117.515}{182.961}{to}%
\TextAt{8}{131.356}{182.961}{this}%
\TextAt{8}{153.709}{182.961}{question}%
\TextAt{8}{199.805}{182.961}{is}%
\TextAt{8}{211.1}{182.961}{based}%
\TextAt{8}{243.208}{182.961}{on}%
\TextAt{8}{258.998}{182.961}{the}%
\TextAt{8}{278.685}{182.961}{following}%
\TextAt{8}{327.65}{182.961}{result}%
\TextAt{8}{359.759}{182.961}{about}%
\TextAt{8}{392.788}{182.961}{unique}%
\TextAt{8}{430.36}{182.961}{factorization}%
\TextAt{8}{497.867}{182.961}{in}%
\TextAt{18}{511.051}{182.961}{Z}%
\TextAt{9}{519.021}{182.961}{\Character{91}}%
\TextAt{13}{522.273}{182.961}{i}%
\TextAt{9}{526.266}{182.961}{\Character{93}}%
\TextAt{8}{529.518}{182.961}{.}%
\TextAt{12}{79.226}{157.645}{Proposition}%
\TextAt{12}{151.937}{157.645}{(1.2).}%
\TextAt{8}{187.253}{157.645}{The}%
\TextAt{8}{211.307}{157.645}{ring}%
\TextAt{18}{235.373}{157.645}{Z}%
\TextAt{9}{243.343}{157.645}{\Character{91}}%
\TextAt{13}{246.594}{157.645}{i}%
\TextAt{9}{250.588}{157.645}{\Character{93}}%
\TextAt{8}{257.737}{157.645}{is}%
\TextAt{8}{269.501}{157.645}{euclidean,}%
\TextAt{8}{324.124}{157.645}{therefore}%
\TextAt{8}{373.224}{157.645}{in}%
\TextAt{8}{386.877}{157.645}{particular}%
\TextAt{8}{440.866}{157.645}{factorial.}%
\TextAt{19}{79.226}{124.059}{Proof.}%
\TextAt{8}{116}{124.059}{We}%
\TextAt{8}{138.027}{124.059}{show}%
\TextAt{8}{168.913}{124.059}{that}%
\TextAt{18}{196.164}{124.059}{Z}%
\TextAt{9}{204.134}{124.059}{\Character{91}}%
\TextAt{13}{207.386}{124.059}{i}%
\TextAt{9}{211.379}{124.059}{\Character{93}}%
\TextAt{8}{220.41}{124.059}{is}%
\TextAt{8}{234.057}{124.059}{euclidean}%
\TextAt{8}{287.31}{124.059}{with}%
\TextAt{8}{315.853}{124.059}{respect}%
\TextAt{8}{357.785}{124.059}{to}%
\TextAt{8}{373.977}{124.059}{the}%
\TextAt{8}{396.016}{124.059}{function}%
\TextAt{18}{443.746}{124.059}{Z}%
\TextAt{9}{451.717}{124.059}{\Character{91}}%
\TextAt{13}{454.968}{124.059}{i}%
\TextAt{9}{458.962}{124.059}{\Character{93}}%
\NativeMathAt{468.742}{124.059}{11.955}{\rightarrow}%
\TextAt{18}{487.226}{124.059}{N}%
\TextAt{16}{499.8}{124.059}{\Character{91}}%
\TextAt{16}{511.695}{124.059}{\Character{102}}%
\TextAt{9}{517.687}{124.059}{0}%
\TextAt{16}{523.54}{124.059}{\Character{103}}%
\TextAt{8}{529.518}{124.059}{,}%
\TextAt{13}{79.226}{107.013}{\Character{11}}%
\NativeArrowAt{90.565}{107.013}{11.955}{11.955}{\mapsto}%
\TextAt{16}{106.334}{107.013}{\Character{106}}%
\TextAt{13}{109.659}{107.013}{\Character{11}}%
\TextAt{16}{117.181}{107.013}{\Character{106}}%
\TextAt{17}{120.502}{111.352}{2}%
\TextAt{8}{125.234}{107.013}{.}%
\TextAt{8}{134.557}{107.013}{So,}%
\TextAt{8}{154.432}{107.013}{for}%
\TextAt{13}{172.608}{107.013}{\Character{11}\Character{59}}%
\TextAt{13}{185.352}{107.013}{\Character{12}}%
\NativeMathAt{196.463}{107.013}{11.955}{\in}%
\TextAt{18}{208.25}{107.013}{Z}%
\TextAt{9}{216.221}{107.013}{\Character{91}}%
\TextAt{13}{219.472}{107.013}{i}%
\TextAt{9}{223.466}{107.013}{\Character{93}}%
\TextAt{8}{226.717}{107.013}{,}%
\TextAt{13}{234.231}{107.013}{\Character{12}}%
\NativeMathAt{245.32}{107.013}{11.955}{\neq}%
\TextAt{9}{258.232}{107.013}{0}%
\TextAt{8}{264.095}{107.013}{,}%
\TextAt{8}{271.608}{107.013}{one}%
\TextAt{8}{293.359}{107.013}{has}%
\TextAt{8}{314.525}{107.013}{to}%
\TextAt{8}{329.126}{107.013}{verify}%
\TextAt{8}{361.936}{107.013}{the}%
\TextAt{8}{382.384}{107.013}{existence}%
\TextAt{8}{432.48}{107.013}{of}%
\TextAt{8}{446.102}{107.013}{gaussian}%
\TextAt{8}{493.353}{107.013}{integers}%
\TextAt{13}{79.226}{89.967}{\Character{13}}%
\TextAt{13}{85.933}{89.967}{\Character{59}}%
\TextAt{13}{91.169}{89.967}{\Character{26}}%
\TextAt{8}{101.128}{89.967}{such}%
\TextAt{8}{127.525}{89.967}{that}%
\TextAt{13}{228.224}{72.921}{\Character{11}}%
\TextAt{9}{239.066}{72.921}{\Character{61}}%
\TextAt{13}{251.492}{72.921}{\Character{13}}%
\TextAt{13}{258.199}{72.921}{\Character{12}}%
\TextAt{9}{268.142}{72.921}{\Character{43}}%
\TextAt{13}{279.903}{72.921}{\Character{26}}%
\TextAt{8}{297.896}{72.921}{and}%
\TextAt{16}{328.717}{72.921}{\Character{106}}%
\TextAt{13}{332.038}{72.921}{\Character{26}}%
\TextAt{16}{338.075}{72.921}{\Character{106}}%
\TextAt{17}{341.396}{77.858}{2}%
\TextAt{13}{349.449}{72.921}{\Character{60}}%
\TextAt{16}{361.875}{72.921}{\Character{106}}%
\TextAt{13}{365.196}{72.921}{\Character{12}}%
\TextAt{16}{372.467}{72.921}{\Character{106}}%
\TextAt{17}{375.788}{77.858}{2}%
\TextAt{13}{380.52}{72.921}{\Character{58}}%
\TextAt{8}{303.069}{43.033}{1}%
\Destination{79.23}{773.075}{nav.chapter.1}%
\Destination{78.23}{809.815}{nav.page.1}%
\Destination{79.23}{570.035}{nav.section.1.1}%
\Destination{79.36}{629.855}{nav.print.1}%
\Destination{79.36}{398.655}{nav.number.1.1.1}%
\Destination{79.36}{171.655}{nav.number.1.1.2}%
\Destination{79.23}{773.075}{outline.5}%
\Destination{79.23}{570.035}{outline.6}%
}

% Page 14 (body page 2)
\Page{%
\TextAt{2}{79.226}{796.943}{I.}%
\TextAt{2}{89.83}{796.943}{Algebraic}%
\TextAt{2}{138.648}{796.943}{Integers}%
\Vector{q 0 G 0 g 0.299 w 7.226 18.729 m 460.769 18.729 l S Q}%
\TextAt{8}{79.226}{761.124}{It}%
\TextAt{8}{92.585}{761.124}{clearly}%
\TextAt{8}{130.667}{761.124}{suffices}%
\TextAt{8}{171.14}{761.124}{to}%
\TextAt{8}{186.136}{761.124}{find}%
\TextAt{13}{210.231}{761.124}{\Character{13}}%
\NativeMathAt{221.443}{761.124}{11.955}{\in}%
\TextAt{18}{233.904}{761.124}{Z}%
\TextAt{9}{241.874}{761.124}{\Character{91}}%
\TextAt{13}{245.125}{761.124}{i}%
\TextAt{9}{249.119}{761.124}{\Character{93}}%
\TextAt{8}{256.954}{761.124}{such}%
\TextAt{8}{284.037}{761.124}{that}%
\TextAt{16}{310.093}{761.124}{\Character{106}}%
\TextAt{13}{313.413}{761.124}{\Character{11}}%
\TextAt{9}{320.935}{761.124}{\Character{47}}%
\TextAt{13}{326.788}{761.124}{\Character{12}}%
\TextAt{16}{337.184}{761.124}{\Character{0}}%
\TextAt{13}{349.607}{761.124}{\Character{13}}%
\TextAt{16}{356.329}{761.124}{\Character{106}}%
\TextAt{13}{364.14}{761.124}{\Character{60}}%
\TextAt{9}{377.735}{761.124}{1}%
\TextAt{8}{383.588}{761.124}{.}%
\TextAt{8}{394.095}{761.124}{Now,}%
\TextAt{8}{424.865}{761.124}{the}%
\TextAt{8}{445.707}{761.124}{gaussian}%
\TextAt{8}{493.353}{761.124}{integers}%
\TextAt{8}{79.226}{744.078}{form}%
\TextAt{8}{106.361}{744.078}{a}%
\TextAt{8}{115.611}{744.078}{lattice}%
\TextAt{8}{150.875}{744.078}{in}%
\TextAt{8}{164.023}{744.078}{the}%
\TextAt{8}{183.674}{744.078}{complex}%
\TextAt{8}{229.016}{744.078}{plane}%
\TextAt{18}{259.726}{744.078}{C}%
\TextAt{8}{271.752}{744.078}{(the}%
\TextAt{8}{295.958}{744.078}{points}%
\TextAt{8}{330.636}{744.078}{with}%
\TextAt{8}{356.791}{744.078}{integer}%
\TextAt{8}{394.984}{744.078}{coordinates}%
\TextAt{8}{456.658}{744.078}{with}%
\TextAt{8}{482.812}{744.078}{respect}%
\TextAt{8}{522.357}{744.078}{to}%
\TextAt{8}{79.226}{727.032}{the}%
\TextAt{8}{99.143}{727.032}{basis}%
\TextAt{9}{127.643}{727.032}{1}%
\TextAt{13}{133.496}{727.032}{\Character{59}}%
\TextAt{13}{138.72}{727.032}{i}%
\TextAt{8}{142.733}{727.032}{).}%
\TextAt{8}{155.657}{727.032}{The}%
\TextAt{8}{179.47}{727.032}{complex}%
\TextAt{8}{225.078}{727.032}{number}%
\TextAt{13}{267.435}{727.032}{\Character{11}}%
\TextAt{9}{274.956}{727.032}{\Character{47}}%
\TextAt{13}{280.809}{727.032}{\Character{12}}%
\TextAt{8}{291.738}{727.032}{lies}%
\TextAt{8}{311.714}{727.032}{in}%
\TextAt{8}{325.127}{727.032}{some}%
\TextAt{8}{354.213}{727.032}{mesh}%
\TextAt{8}{383.944}{727.032}{of}%
\TextAt{8}{397.034}{727.032}{the}%
\TextAt{8}{416.951}{727.032}{lattice}%
\TextAt{8}{452.48}{727.032}{and}%
\TextAt{8}{475.003}{727.032}{its}%
\TextAt{8}{491.082}{727.032}{distance}%
\TextAt{8}{79.226}{709.986}{from}%
\TextAt{8}{107.524}{709.986}{the}%
\TextAt{8}{128.338}{709.986}{nearest}%
\TextAt{8}{169.38}{709.986}{lattice}%
\TextAt{8}{205.808}{709.986}{point}%
\TextAt{8}{237.034}{709.986}{is}%
\TextAt{8}{249.456}{709.986}{not}%
\TextAt{8}{270.927}{709.986}{greater}%
\TextAt{8}{311.264}{709.986}{than}%
\TextAt{8}{339.239}{709.986}{half}%
\TextAt{8}{362.982}{709.986}{the}%
\TextAt{8}{383.796}{709.986}{length}%
\TextAt{8}{420.223}{709.986}{of}%
\TextAt{8}{434.211}{709.986}{the}%
\TextAt{8}{455.025}{709.986}{diagonal}%
\TextAt{8}{502.523}{709.986}{of}%
\TextAt{8}{516.51}{709.986}{the}%
\TextAt{8}{79.226}{692.94}{mesh,}%
\TextAt{8}{112.45}{692.94}{i.e.,}%
\TextAt{16}{134.555}{702.829}{\Character{112}}%
\Vector{q 0 G 0 g 0.478 w 72.517 -70.932 m 78.37 -70.932 l S Q}%
\TextAt{9}{144.517}{692.94}{2\Character{47}2}%
\TextAt{8}{162.076}{692.94}{.}%
\TextAt{8}{170.525}{692.94}{Therefore}%
\TextAt{8}{223.522}{692.94}{there}%
\TextAt{8}{253.434}{692.94}{exists}%
\TextAt{8}{285.749}{692.94}{an}%
\TextAt{8}{302.008}{692.94}{element}%
\TextAt{13}{345.249}{692.94}{\Character{13}}%
\NativeMathAt{355.292}{692.94}{11.955}{\in}%
\TextAt{18}{366.583}{692.94}{Z}%
\TextAt{9}{374.554}{692.94}{\Character{91}}%
\TextAt{13}{377.805}{692.94}{i}%
\TextAt{9}{381.799}{692.94}{\Character{93}}%
\TextAt{8}{388.948}{692.94}{with}%
\TextAt{20}{256.89}{676.529}{\Character{12}}%
\TextAt{20}{256.89}{669.356}{\Character{12}}%
\TextAt{20}{256.89}{662.183}{\Character{12}}%
\TextAt{20}{256.89}{655.009}{\Character{12}}%
\TextAt{13}{262.071}{667.282}{\Character{11}}%
\Vector{q 0 G 0 g 0.478 w 190.071 -111.817 m 197.593 -111.817 l S Q}%
\TextAt{13}{262.196}{650.993}{\Character{12}}%
\TextAt{16}{273.445}{659.194}{\Character{0}}%
\TextAt{13}{285.4}{659.194}{\Character{13}}%
\TextAt{20}{292.122}{676.529}{\Character{12}}%
\TextAt{20}{292.122}{669.356}{\Character{12}}%
\TextAt{20}{292.122}{662.183}{\Character{12}}%
\TextAt{20}{292.122}{655.009}{\Character{12}}%
\NativeMathAt{299.428}{659.194}{11.955}{\leq}%
\TextAt{16}{313.243}{677.171}{\Character{112}}%
\Vector{q 0 G 0 g 0.478 w 251.206 -96.59 m 257.059 -96.59 l S Q}%
\TextAt{9}{323.206}{667.282}{2}%
\Vector{q 0 G 0 g 0.478 w 241.243 -111.817 m 257.059 -111.817 l S Q}%
\TextAt{9}{318.224}{650.993}{2}%
\TextAt{13}{333.575}{659.194}{\Character{60}}%
\TextAt{9}{346.001}{659.194}{1}%
\TextAt{13}{351.854}{659.194}{\Character{58}}%
\Vector{q 0 G 0 g 0.399 w 452.334 -146.057 m 452.334 -137.987 l S Q}%
\Vector{q 0 G 0 g 0.399 w 452.534 -138.187 m 459.707 -138.187 l S Q}%
\Vector{q 0 G 0 g 0.399 w 452.534 -145.858 m 459.707 -145.858 l S Q}%
\Vector{q 0 G 0 g 0.399 w 459.906 -146.057 m 459.906 -137.987 l S Q}%
\TextAt{8}{97.159}{599.264}{Based}%
\TextAt{8}{132.843}{599.264}{on}%
\TextAt{8}{150.427}{599.264}{this}%
\TextAt{8}{174.574}{599.264}{result}%
\TextAt{8}{208.476}{599.264}{about}%
\TextAt{8}{243.299}{599.264}{the}%
\TextAt{8}{264.781}{599.264}{ring}%
\TextAt{18}{290.171}{599.264}{Z}%
\TextAt{9}{298.141}{599.264}{\Character{91}}%
\TextAt{13}{301.393}{599.264}{i}%
\TextAt{9}{305.386}{599.264}{\Character{93}}%
\TextAt{8}{308.638}{599.264}{,}%
\TextAt{8}{317.443}{599.264}{theorem}%
\TextAt{8}{364.293}{599.264}{(1.1)}%
\TextAt{8}{393.593}{599.264}{now}%
\TextAt{8}{419.306}{599.264}{follows}%
\TextAt{8}{459.067}{599.264}{like}%
\TextAt{8}{481.851}{599.264}{this:}%
\TextAt{8}{511.874}{599.264}{it}%
\TextAt{8}{524.903}{599.264}{is}%
\TextAt{8}{79.226}{582.218}{sufficient}%
\TextAt{8}{127.74}{582.218}{to}%
\TextAt{8}{141.405}{582.218}{show}%
\TextAt{8}{169.762}{582.218}{that}%
\TextAt{8}{194.486}{582.218}{a}%
\TextAt{8}{203.595}{582.218}{prime}%
\TextAt{8}{236.113}{582.218}{number}%
\TextAt{13}{278.064}{582.218}{p}%
\NativeMathAt{287.26}{582.218}{11.955}{\equiv}%
\TextAt{9}{299.879}{582.218}{1}%
\TextAt{8}{309.053}{582.218}{mod}%
\TextAt{9}{334.491}{582.218}{4}%
\TextAt{8}{343.596}{582.218}{of}%
\TextAt{18}{356.28}{582.218}{Z}%
\TextAt{8}{367.502}{582.218}{does}%
\TextAt{8}{393.253}{582.218}{not}%
\TextAt{8}{413.421}{582.218}{remain}%
\TextAt{8}{451.797}{582.218}{a}%
\TextAt{8}{460.907}{582.218}{prime}%
\TextAt{8}{493.425}{582.218}{element}%
\TextAt{8}{79.226}{565.172}{in}%
\TextAt{8}{92.879}{565.172}{the}%
\TextAt{8}{113.035}{565.172}{ring}%
\TextAt{18}{137.101}{565.172}{Z}%
\TextAt{9}{145.071}{565.172}{\Character{91}}%
\TextAt{13}{148.323}{565.172}{i}%
\TextAt{9}{152.316}{565.172}{\Character{93}}%
\TextAt{8}{155.568}{565.172}{.}%
\TextAt{8}{164.016}{565.172}{Indeed,}%
\TextAt{8}{205.298}{565.172}{if}%
\TextAt{8}{216.021}{565.172}{this}%
\TextAt{8}{238.844}{565.172}{is}%
\TextAt{8}{250.608}{565.172}{proved,}%
\TextAt{8}{291.913}{565.172}{then}%
\TextAt{8}{318.573}{565.172}{there}%
\TextAt{8}{348.485}{565.172}{exists}%
\TextAt{8}{380.799}{565.172}{a}%
\TextAt{8}{390.555}{565.172}{decomposition}%
\TextAt{13}{283.473}{537.299}{p}%
\TextAt{9}{292.669}{537.299}{\Character{61}}%
\TextAt{13}{305.095}{537.299}{\Character{11}}%
\TextAt{16}{315.273}{537.299}{\Character{1}}%
\TextAt{13}{321.251}{537.299}{\Character{12}}%
\TextAt{8}{79.226}{509.426}{into}%
\TextAt{8}{102.969}{509.426}{two}%
\TextAt{8}{125.086}{509.426}{non-units}%
\TextAt{13}{177.175}{509.426}{\Character{11}\Character{59}}%
\TextAt{13}{189.919}{509.426}{\Character{12}}%
\TextAt{8}{201.109}{509.426}{of}%
\TextAt{18}{214.439}{509.426}{Z}%
\TextAt{9}{222.41}{509.426}{\Character{91}}%
\TextAt{13}{225.661}{509.426}{i}%
\TextAt{9}{229.655}{509.426}{\Character{93}}%
\TextAt{8}{232.906}{509.426}{.}%
\TextAt{8}{241.355}{509.426}{The}%
\TextAt{12}{265.408}{509.426}{norm}%
\TextAt{8}{300.222}{509.426}{of}%
\TextAt{13}{313.552}{509.426}{z}%
\TextAt{9}{322.843}{509.426}{\Character{61}}%
\TextAt{13}{335.269}{509.426}{x}%
\TextAt{9}{344.578}{509.426}{\Character{43}}%
\TextAt{13}{356.339}{509.426}{iy}%
\TextAt{8}{370.366}{509.426}{is}%
\TextAt{8}{382.13}{509.426}{defined}%
\TextAt{8}{422.443}{509.426}{by}%
\TextAt{13}{204.425}{481.552}{N}%
\TextAt{9}{215.047}{481.552}{\Character{40}}%
\TextAt{13}{219.6}{481.552}{x}%
\TextAt{9}{228.908}{481.552}{\Character{43}}%
\TextAt{13}{240.67}{481.552}{iy}%
\TextAt{9}{250.8}{481.552}{\Character{41}}%
\TextAt{9}{258.655}{481.552}{\Character{61}}%
\TextAt{9}{271.076}{481.552}{\Character{40}}%
\TextAt{13}{275.651}{481.552}{x}%
\TextAt{9}{284.959}{481.552}{\Character{43}}%
\TextAt{13}{296.721}{481.552}{iy}%
\TextAt{9}{306.851}{481.552}{\Character{41}\Character{40}}%
\TextAt{13}{315.955}{481.552}{x}%
\TextAt{16}{325.264}{481.552}{\Character{0}}%
\TextAt{13}{337.219}{481.552}{iy}%
\TextAt{9}{347.349}{481.552}{\Character{41}}%
\TextAt{9}{355.204}{481.552}{\Character{61}}%
\TextAt{13}{367.648}{481.552}{x}%
\TextAt{17}{374.3}{486.488}{2}%
\TextAt{9}{381.689}{481.552}{\Character{43}}%
\TextAt{13}{393.45}{481.552}{y}%
\TextAt{17}{399.587}{486.488}{2}%
\TextAt{13}{404.319}{481.552}{\Character{59}}%
\TextAt{8}{79.226}{453.679}{i.e.,}%
\TextAt{8}{101.331}{453.679}{by}%
\TextAt{13}{117.59}{453.679}{N}%
\TextAt{9}{128.213}{453.679}{\Character{40}}%
\TextAt{13}{132.765}{453.679}{z}%
\TextAt{9}{138.736}{453.679}{\Character{41}}%
\TextAt{9}{146.591}{453.679}{\Character{61}}%
\TextAt{16}{159.034}{453.679}{\Character{106}}%
\TextAt{13}{162.355}{453.679}{z}%
\TextAt{16}{168.326}{453.679}{\Character{106}}%
\TextAt{17}{171.647}{458.017}{2}%
\TextAt{8}{176.379}{453.679}{.}%
\TextAt{8}{184.827}{453.679}{It}%
\TextAt{8}{197.5}{453.679}{is}%
\TextAt{8}{209.264}{453.679}{multiplicative,}%
\TextAt{8}{286.339}{453.679}{so}%
\TextAt{8}{300.709}{453.679}{that}%
\TextAt{8}{326.078}{453.679}{one}%
\TextAt{8}{347.537}{453.679}{has}%
\ErrataAt{206}{425.855}{200}{\makebox[200bp][c]{$p^2=N(\alpha)\cdot N(\beta)$.}}%
\TextAt{8}{79.226}{397.932}{Since}%
\TextAt{13}{110.573}{397.932}{\Character{11}\Character{59}}%
\TextAt{13}{123.317}{397.932}{\Character{12}}%
\TextAt{8}{135.298}{397.932}{are}%
\TextAt{8}{155.598}{397.932}{not}%
\TextAt{8}{177.202}{397.932}{units,}%
\TextAt{8}{210.767}{397.932}{it}%
\TextAt{8}{223.261}{397.932}{follows}%
\TextAt{8}{262.487}{397.932}{that}%
\TextAt{13}{288.646}{397.932}{N}%
\TextAt{9}{299.268}{397.932}{\Character{40}}%
\TextAt{13}{303.821}{397.932}{\Character{11}}%
\TextAt{9}{311.342}{397.932}{\Character{41}}%
\TextAt{13}{315.895}{397.932}{\Character{59}}%
\TextAt{13}{321.119}{397.932}{N}%
\TextAt{9}{331.761}{397.932}{\Character{40}}%
\TextAt{13}{336.314}{397.932}{\Character{12}}%
\TextAt{9}{343.585}{397.932}{\Character{41}}%
\NativeMathAt{352.804}{397.932}{11.955}{\neq}%
\TextAt{9}{366.564}{397.932}{1}%
\TextAt{8}{377.116}{397.932}{(see}%
\TextAt{8}{401.374}{397.932}{exercise}%
\TextAt{8}{445.465}{397.932}{1),}%
\TextAt{8}{464.014}{397.932}{and}%
\TextAt{8}{487.567}{397.932}{therefore}%
\TextAt{13}{79.226}{380.886}{p}%
\TextAt{9}{88.422}{380.886}{\Character{61}}%
\TextAt{13}{100.848}{380.886}{N}%
\TextAt{9}{111.47}{380.886}{\Character{40}}%
\TextAt{13}{116.023}{380.886}{\Character{11}}%
\TextAt{9}{123.544}{380.886}{\Character{41}}%
\TextAt{9}{131.399}{380.886}{\Character{61}}%
\TextAt{13}{143.843}{380.886}{a}%
\TextAt{17}{149.988}{385.225}{2}%
\TextAt{9}{157.377}{380.886}{\Character{43}}%
\TextAt{13}{169.138}{380.886}{b}%
\TextAt{17}{174.115}{385.225}{2}%
\TextAt{8}{178.848}{380.886}{,}%
\TextAt{8}{185.997}{380.886}{where}%
\TextAt{8}{219.806}{380.886}{we}%
\TextAt{8}{237.033}{380.886}{put}%
\TextAt{13}{258.493}{380.886}{\Character{11}}%
\TextAt{9}{269.335}{380.886}{\Character{61}}%
\TextAt{13}{281.761}{380.886}{a}%
\TextAt{9}{290.563}{380.886}{\Character{43}}%
\TextAt{13}{302.324}{380.886}{bi}%
\TextAt{8}{311.294}{380.886}{.}%
\TextAt{8}{97.159}{360.852}{Finally,}%
\TextAt{8}{139.991}{360.852}{in}%
\TextAt{8}{154.225}{360.852}{order}%
\TextAt{8}{185.375}{360.852}{to}%
\TextAt{8}{200.266}{360.852}{prove}%
\TextAt{8}{232.397}{360.852}{that}%
\TextAt{8}{258.347}{360.852}{a}%
\TextAt{8}{268.683}{360.852}{rational}%
\TextAt{8}{312.852}{360.852}{prime}%
\TextAt{8}{346.597}{360.852}{of}%
\TextAt{8}{360.508}{360.852}{the}%
\TextAt{8}{381.245}{360.852}{form}%
\TextAt{13}{409.466}{360.852}{p}%
\TextAt{9}{419.652}{360.852}{\Character{61}}%
\TextAt{9}{433.054}{360.852}{1}%
\TextAt{9}{441.948}{360.852}{\Character{43}}%
\TextAt{9}{454.107}{360.852}{4}%
\TextAt{13}{459.983}{360.852}{n}%
\TextAt{8}{471.449}{360.852}{cannot}%
\TextAt{8}{510.406}{360.852}{be}%
\TextAt{8}{526.911}{360.852}{a}%
\TextAt{8}{79.226}{343.806}{prime}%
\TextAt{8}{112.39}{343.806}{element}%
\TextAt{8}{155.632}{343.806}{in}%
\TextAt{18}{169.284}{343.806}{Z}%
\TextAt{9}{177.255}{343.806}{\Character{91}}%
\TextAt{13}{180.506}{343.806}{i}%
\TextAt{9}{184.5}{343.806}{\Character{93}}%
\TextAt{8}{187.751}{343.806}{,}%
\TextAt{8}{194.901}{343.806}{we}%
\TextAt{8}{212.128}{343.806}{note}%
\TextAt{8}{238.142}{343.806}{that}%
\TextAt{8}{263.511}{343.806}{the}%
\TextAt{8}{283.668}{343.806}{congruence}%
\TextAt{16}{267.443}{315.932}{\Character{0}}%
\TextAt{9}{276.741}{315.932}{1}%
\NativeMathAt{285.915}{315.932}{11.955}{\equiv}%
\TextAt{13}{298.535}{315.932}{x}%
\TextAt{17}{305.187}{320.869}{2}%
\TextAt{8}{313.24}{315.932}{mod}%
\TextAt{13}{338.678}{315.932}{p}%
\TextAt{8}{79.226}{288.059}{admits}%
\TextAt{8}{118.842}{288.059}{a}%
\TextAt{8}{129.777}{288.059}{solution,}%
\TextAt{8}{178.797}{288.059}{namely}%
\TextAt{13}{220.624}{288.059}{x}%
\TextAt{9}{232.607}{288.059}{\Character{61}}%
\TextAt{9}{247.025}{288.059}{\Character{40}2}%
\TextAt{13}{257.448}{288.059}{n}%
\TextAt{9}{264.436}{288.059}{\Character{41}\Character{33}}%
\TextAt{8}{272.24}{288.059}{.}%
\TextAt{8}{284.227}{288.059}{Indeed,}%
\TextAt{8}{326.982}{288.059}{since}%
\TextAt{16}{356.83}{288.059}{\Character{0}}%
\TextAt{9}{366.129}{288.059}{1}%
\NativeMathAt{377.313}{288.059}{11.955}{\equiv}%
\TextAt{9}{391.942}{288.059}{\Character{40}}%
\TextAt{13}{396.495}{288.059}{p}%
\TextAt{16}{405.831}{288.059}{\Character{0}}%
\TextAt{9}{418.59}{288.059}{1\Character{41}\Character{33}}%
\TextAt{8}{435.568}{288.059}{mod}%
\TextAt{13}{461.005}{288.059}{p}%
\TextAt{8}{471.958}{288.059}{by}%
\TextAt{8}{489.396}{288.059}{Wilson’s}%
\TextAt{8}{79.226}{271.013}{theorem,}%
\TextAt{8}{128.003}{271.013}{one}%
\TextAt{8}{149.463}{271.013}{has}%
\TextAt{16}{160.991}{243.14}{\Character{0}}%
\TextAt{9}{170.29}{243.14}{1}%
\NativeMathAt{179.464}{243.14}{11.955}{\equiv}%
\TextAt{9}{192.083}{243.14}{\Character{40}}%
\TextAt{13}{196.635}{243.14}{p}%
\TextAt{16}{205.167}{243.14}{\Character{0}}%
\TextAt{9}{217.122}{243.14}{1\Character{41}\Character{33}}%
\TextAt{9}{234.062}{243.14}{\Character{61}}%
\TextAt{9}{246.484}{243.14}{\Character{91}1}%
\TextAt{16}{258.287}{243.14}{\Character{1}}%
\TextAt{9}{264.264}{243.14}{2}%
\TextAt{16}{272.11}{243.14}{\Character{1}}%
\TextAt{16}{277.406}{243.14}{\Character{1}}%
\TextAt{16}{282.714}{243.14}{\Character{1}}%
\TextAt{9}{288.05}{243.14}{\Character{40}2}%
\TextAt{13}{298.455}{243.14}{n}%
\TextAt{9}{305.443}{243.14}{\Character{41}\Character{93}\Character{91}\Character{40}}%
\TextAt{13}{321.051}{243.14}{p}%
\TextAt{16}{329.583}{243.14}{\Character{0}}%
\TextAt{9}{341.538}{243.14}{1\Character{41}\Character{40}}%
\TextAt{13}{356.496}{243.14}{p}%
\TextAt{16}{365.027}{243.14}{\Character{0}}%
\TextAt{9}{376.983}{243.14}{2\Character{41}}%
\TextAt{16}{389.38}{243.14}{\Character{1}}%
\TextAt{16}{394.676}{243.14}{\Character{1}}%
\TextAt{16}{399.984}{243.14}{\Character{1}}%
\TextAt{9}{405.32}{243.14}{\Character{40}}%
\TextAt{13}{409.873}{243.14}{p}%
\TextAt{16}{418.405}{243.14}{\Character{0}}%
\TextAt{9}{430.36}{243.14}{2}%
\TextAt{13}{436.213}{243.14}{n}%
\TextAt{9}{443.2}{243.14}{\Character{41}\Character{93}}%
\NativeMathAt{179.464}{223.105}{11.955}{\equiv}%
\TextAt{9}{192.083}{223.105}{\Character{91}\Character{40}2}%
\TextAt{13}{205.74}{223.105}{n}%
\TextAt{9}{212.728}{223.105}{\Character{41}\Character{33}\Character{93}\Character{91}\Character{40}}%
\TextAt{16}{231.587}{223.105}{\Character{0}}%
\TextAt{9}{240.886}{223.105}{1\Character{41}}%
\TextAt{17}{251.291}{228.041}{2}%
\TextAt{21}{255.525}{228.041}{n}%
\TextAt{9}{261.161}{223.105}{\Character{40}2}%
\TextAt{13}{271.567}{223.105}{n}%
\TextAt{9}{278.554}{223.105}{\Character{41}\Character{33}\Character{93}}%
\TextAt{9}{292.888}{223.105}{\Character{61}}%
\TextAt{9}{305.31}{223.105}{\Character{91}\Character{40}2}%
\TextAt{13}{319.013}{223.105}{n}%
\TextAt{9}{326.001}{223.105}{\Character{41}\Character{33}\Character{93}}%
\TextAt{17}{337.057}{228.041}{2}%
\TextAt{8}{345.11}{223.105}{mod}%
\TextAt{13}{370.548}{223.105}{p\Character{58}}%
\ErrataAt{79.22}{195.255}{453.64}{Thus we have $p\mid x^2+1=(x+i)(x-i)$. But since $\frac{x}{p}\pm\frac{i}{p}\notin\mathbb Z[i]$, $p$ does not}%
\TextAt{8}{79.226}{178.186}{divide}%
\TextAt{8}{113.839}{178.186}{any}%
\TextAt{8}{135.78}{178.186}{of}%
\TextAt{8}{148.934}{178.186}{the}%
\TextAt{8}{168.914}{178.186}{factors}%
\TextAt{13}{206.851}{178.186}{x}%
\TextAt{9}{215.799}{178.186}{\Character{43}}%
\TextAt{13}{227.2}{178.186}{i\Character{59}}%
\TextAt{13}{236.418}{178.186}{x}%
\TextAt{16}{245.386}{178.186}{\Character{0}}%
\TextAt{13}{256.981}{178.186}{i}%
\TextAt{8}{260.974}{178.186}{,}%
\TextAt{8}{267.983}{178.186}{and}%
\TextAt{8}{290.569}{178.186}{is}%
\TextAt{8}{302.157}{178.186}{therefore}%
\TextAt{8}{351.08}{178.186}{not}%
\TextAt{8}{371.718}{178.186}{a}%
\TextAt{8}{381.297}{178.186}{prime}%
\TextAt{8}{414.285}{178.186}{element}%
\TextAt{8}{457.35}{178.186}{in}%
\TextAt{8}{470.827}{178.186}{the}%
\TextAt{8}{490.807}{178.186}{factorial}%
\TextAt{8}{79.226}{161.14}{ring}%
\TextAt{18}{103.292}{161.14}{Z}%
\TextAt{9}{111.262}{161.14}{\Character{91}}%
\TextAt{13}{114.514}{161.14}{i}%
\TextAt{9}{118.507}{161.14}{\Character{93}}%
\TextAt{8}{121.759}{161.14}{.}%
\TextAt{8}{97.159}{141.105}{The}%
\TextAt{8}{121.204}{141.105}{example}%
\TextAt{8}{167.043}{141.105}{of}%
\TextAt{8}{180.365}{141.105}{the}%
\TextAt{8}{200.513}{141.105}{equation}%
\TextAt{13}{248.313}{141.105}{p}%
\TextAt{9}{257.509}{141.105}{\Character{61}}%
\TextAt{13}{269.934}{141.105}{x}%
\TextAt{17}{276.586}{145.444}{2}%
\TextAt{9}{283.957}{141.105}{\Character{43}}%
\TextAt{13}{295.701}{141.105}{y}%
\TextAt{17}{301.838}{145.444}{2}%
\TextAt{8}{310.459}{141.105}{shows}%
\TextAt{8}{344.068}{141.105}{that}%
\TextAt{8}{369.428}{141.105}{even}%
\TextAt{8}{396.079}{141.105}{quite}%
\TextAt{8}{425.66}{141.105}{elementary}%
\TextAt{8}{485.487}{141.105}{questions}%
\TextAt{8}{79.226}{124.059}{about}%
\TextAt{8}{112.798}{124.059}{rational}%
\TextAt{8}{156.46}{124.059}{integers}%
\TextAt{8}{199.847}{124.059}{may}%
\TextAt{8}{225.29}{124.059}{lead}%
\TextAt{8}{250.075}{124.059}{to}%
\TextAt{8}{264.459}{124.059}{the}%
\TextAt{8}{284.689}{124.059}{consideration}%
\TextAt{8}{356.374}{124.059}{of}%
\TextAt{8}{369.777}{124.059}{higher}%
\TextAt{8}{405.621}{124.059}{domains}%
\TextAt{8}{451.937}{124.059}{of}%
\TextAt{8}{465.341}{124.059}{integers.}%
\TextAt{8}{513.426}{124.059}{But}%
\TextAt{8}{79.226}{107.013}{it}%
\TextAt{8}{90.937}{107.013}{was}%
\TextAt{8}{113.444}{107.013}{not}%
\TextAt{8}{134.265}{107.013}{so}%
\TextAt{8}{148.643}{107.013}{much}%
\TextAt{8}{179.865}{107.013}{for}%
\TextAt{8}{197.757}{107.013}{this}%
\TextAt{8}{220.586}{107.013}{equation}%
\TextAt{8}{268.402}{107.013}{that}%
\TextAt{8}{293.778}{107.013}{we}%
\TextAt{8}{311.013}{107.013}{have}%
\TextAt{8}{338.015}{107.013}{introduced}%
\TextAt{8}{396.554}{107.013}{the}%
\TextAt{8}{416.718}{107.013}{ring}%
\TextAt{18}{440.791}{107.013}{Z}%
\TextAt{9}{448.761}{107.013}{\Character{91}}%
\TextAt{13}{452.013}{107.013}{i}%
\TextAt{9}{456.006}{107.013}{\Character{93}}%
\TextAt{8}{459.258}{107.013}{,}%
\TextAt{8}{466.416}{107.013}{but}%
\TextAt{8}{487.883}{107.013}{rather}%
\TextAt{8}{523.014}{107.013}{in}%
\TextAt{8}{79.226}{89.967}{order}%
\TextAt{8}{109.432}{89.967}{to}%
\TextAt{8}{123.378}{89.967}{preface}%
\TextAt{8}{163.005}{89.967}{the}%
\TextAt{8}{182.797}{89.967}{general}%
\TextAt{8}{222.758}{89.967}{theory}%
\TextAt{8}{259.145}{89.967}{of}%
\TextAt{8}{272.111}{89.967}{algebraic}%
\TextAt{8}{321.182}{89.967}{integers}%
\TextAt{8}{364.132}{89.967}{with}%
\TextAt{8}{390.428}{89.967}{a}%
\TextAt{8}{399.82}{89.967}{concrete}%
\TextAt{8}{445.627}{89.967}{example.}%
\TextAt{8}{495.905}{89.967}{For}%
\TextAt{8}{516.51}{89.967}{the}%
\TextAt{8}{79.226}{72.921}{same}%
\TextAt{8}{108.552}{72.921}{reason}%
\TextAt{8}{145.039}{72.921}{we}%
\TextAt{8}{162.267}{72.921}{will}%
\TextAt{8}{184.372}{72.921}{now}%
\TextAt{8}{208.76}{72.921}{look}%
\TextAt{8}{234.129}{72.921}{at}%
\TextAt{8}{248.44}{72.921}{this}%
\TextAt{8}{271.262}{72.921}{ring}%
\TextAt{8}{295.328}{72.921}{a}%
\TextAt{8}{305.083}{72.921}{bit}%
\TextAt{8}{323.291}{72.921}{more}%
\TextAt{8}{352.557}{72.921}{closely.}%
\TextAt{8}{303.069}{43.033}{2}%
\Destination{78.23}{809.815}{nav.page.2}%
\Destination{79.36}{759.055}{nav.print.2}%
\Hotlink{444.965}{395.755}{6.858}{11.047}{nav.ex.1.1.1}%
\Hotlink{363.793}{597.055}{25.078}{11.047}{nav.number.1.1.1}%
\Destination{79.226}{439.855}{errata.39}%
}

% Page 15 (body page 3)
\Page{%
\TextAt{2}{79.226}{796.943}{I.}%
\TextAt{2}{89.83}{796.943}{Algebraic}%
\TextAt{2}{138.648}{796.943}{Integers}%
\Vector{q 0 G 0 g 0.299 w 7.226 18.729 m 460.769 18.729 l S Q}%
\TextAt{8}{97.159}{761.124}{When}%
\TextAt{8}{132.415}{761.124}{developing}%
\TextAt{8}{191.426}{761.124}{the}%
\TextAt{8}{212.707}{761.124}{theory}%
\TextAt{8}{250.581}{761.124}{of}%
\TextAt{8}{265.036}{761.124}{divisibility}%
\TextAt{8}{323.784}{761.124}{for}%
\TextAt{8}{342.793}{761.124}{a}%
\TextAt{8}{353.673}{761.124}{ring,}%
\TextAt{8}{382.396}{761.124}{two}%
\TextAt{8}{405.637}{761.124}{basic}%
\TextAt{8}{436.087}{761.124}{problems}%
\TextAt{8}{487.351}{761.124}{are}%
\TextAt{8}{507.986}{761.124}{most}%
\TextAt{8}{79.226}{744.078}{prominent:}%
\TextAt{8}{139.924}{744.078}{on}%
\TextAt{8}{155.956}{744.078}{the}%
\TextAt{8}{175.884}{744.078}{one}%
\TextAt{8}{197.116}{744.078}{hand,}%
\TextAt{8}{229.452}{744.078}{to}%
\TextAt{8}{243.534}{744.078}{determine}%
\TextAt{8}{297.93}{744.078}{the}%
\TextAt{8}{317.858}{744.078}{units}%
\TextAt{8}{346.957}{744.078}{of}%
\TextAt{8}{360.059}{744.078}{the}%
\TextAt{8}{379.988}{744.078}{ring}%
\TextAt{8}{403.826}{744.078}{in}%
\TextAt{8}{417.251}{744.078}{question,}%
\TextAt{8}{466.886}{744.078}{on}%
\TextAt{8}{482.917}{744.078}{the}%
\TextAt{8}{502.846}{744.078}{other,}%
\TextAt{8}{79.226}{727.032}{its}%
\TextAt{8}{96.364}{727.032}{prime}%
\TextAt{8}{130.347}{727.032}{elements.}%
\TextAt{8}{185.212}{727.032}{The}%
\TextAt{8}{210.085}{727.032}{answer}%
\TextAt{8}{249.663}{727.032}{to}%
\TextAt{8}{264.792}{727.032}{the}%
\TextAt{8}{285.768}{727.032}{first}%
\TextAt{8}{310.713}{727.032}{question}%
\TextAt{8}{358.097}{727.032}{in}%
\TextAt{8}{372.569}{727.032}{the}%
\TextAt{8}{393.545}{727.032}{present}%
\TextAt{8}{435.071}{727.032}{case}%
\TextAt{8}{460.662}{727.032}{is}%
\TextAt{8}{473.245}{727.032}{particularly}%
\TextAt{8}{79.226}{709.986}{easy.}%
\TextAt{8}{108.548}{709.986}{A}%
\TextAt{8}{121.221}{709.986}{number}%
\TextAt{13}{163.817}{709.986}{\Character{11}}%
\TextAt{9}{174.66}{709.986}{\Character{61}}%
\TextAt{13}{187.085}{709.986}{a}%
\TextAt{9}{195.887}{709.986}{\Character{43}}%
\TextAt{13}{207.648}{709.986}{bi}%
\NativeMathAt{219.939}{709.986}{11.955}{\in}%
\TextAt{18}{231.23}{709.986}{Z}%
\TextAt{9}{239.2}{709.986}{\Character{91}}%
\TextAt{13}{242.452}{709.986}{i}%
\TextAt{9}{246.445}{709.986}{\Character{93}}%
\TextAt{8}{253.595}{709.986}{is}%
\TextAt{8}{265.358}{709.986}{a}%
\TextAt{8}{275.114}{709.986}{unit}%
\TextAt{8}{299.825}{709.986}{if}%
\TextAt{8}{310.549}{709.986}{and}%
\TextAt{8}{333.312}{709.986}{only}%
\TextAt{8}{359.003}{709.986}{if}%
\TextAt{8}{369.727}{709.986}{its}%
\TextAt{8}{386.046}{709.986}{norm}%
\TextAt{8}{416.615}{709.986}{is}%
\TextAt{9}{428.379}{709.986}{1}%
\TextAt{8}{434.232}{709.986}{:}%
\TextAt{13}{207.964}{682.404}{N}%
\TextAt{9}{218.587}{682.404}{\Character{40}}%
\TextAt{13}{223.139}{682.404}{\Character{11}}%
\TextAt{9}{230.661}{682.404}{\Character{41}}%
\TextAt{9}{238.516}{682.404}{\Character{58}\Character{61}}%
\TextAt{9}{254.177}{682.404}{\Character{40}}%
\TextAt{13}{258.763}{682.404}{a}%
\TextAt{9}{267.565}{682.404}{\Character{43}}%
\TextAt{13}{279.326}{682.404}{ib}%
\TextAt{9}{288.297}{682.404}{\Character{41}\Character{40}}%
\TextAt{13}{297.402}{682.404}{a}%
\TextAt{16}{306.203}{682.404}{\Character{0}}%
\TextAt{13}{318.158}{682.404}{ib}%
\TextAt{9}{327.129}{682.404}{\Character{41}}%
\TextAt{9}{334.984}{682.404}{\Character{61}}%
\TextAt{13}{347.428}{682.404}{a}%
\TextAt{17}{353.572}{687.34}{2}%
\TextAt{9}{360.961}{682.404}{\Character{43}}%
\TextAt{13}{372.723}{682.404}{b}%
\TextAt{17}{377.7}{687.34}{2}%
\TextAt{9}{385.753}{682.404}{\Character{61}}%
\TextAt{9}{398.163}{682.404}{1}%
\TextAt{8}{79.226}{654.822}{(exercise}%
\TextAt{8}{127.083}{654.822}{1),}%
\TextAt{8}{144.645}{654.822}{i.e.,}%
\TextAt{8}{166.75}{654.822}{if}%
\TextAt{8}{177.474}{654.822}{either}%
\TextAt{13}{210.637}{654.822}{a}%
\TextAt{17}{216.782}{659.16}{2}%
\TextAt{9}{224.835}{654.822}{\Character{61}}%
\TextAt{9}{237.245}{654.822}{1}%
\TextAt{13}{243.114}{654.822}{\Character{59}}%
\TextAt{13}{248.338}{654.822}{b}%
\TextAt{17}{253.335}{659.16}{2}%
\TextAt{9}{261.388}{654.822}{\Character{61}}%
\TextAt{9}{273.798}{654.822}{0}%
\TextAt{8}{279.667}{654.822}{,}%
\TextAt{8}{286.816}{654.822}{or}%
\TextAt{13}{301.126}{654.822}{a}%
\TextAt{17}{307.271}{659.16}{2}%
\TextAt{9}{315.324}{654.822}{\Character{61}}%
\TextAt{9}{327.733}{654.822}{0}%
\TextAt{13}{333.603}{654.822}{\Character{59}}%
\TextAt{13}{338.827}{654.822}{b}%
\TextAt{17}{343.824}{659.16}{2}%
\TextAt{9}{351.877}{654.822}{\Character{61}}%
\TextAt{9}{364.286}{654.822}{1}%
\TextAt{8}{370.156}{654.822}{.}%
\TextAt{8}{378.604}{654.822}{We}%
\TextAt{8}{398.749}{654.822}{thus}%
\TextAt{8}{424.5}{654.822}{obtain}%
\TextAt{8}{460.927}{654.822}{the}%
\TextAt{12}{79.226}{629.215}{Proposition}%
\TextAt{12}{153.044}{629.215}{(1.3).}%
\TextAt{8}{189.322}{629.215}{The}%
\TextAt{8}{214.338}{629.215}{group}%
\TextAt{8}{248.476}{629.215}{of}%
\TextAt{8}{262.768}{629.215}{units}%
\TextAt{8}{293.056}{629.215}{of}%
\TextAt{8}{307.348}{629.215}{the}%
\TextAt{8}{328.467}{629.215}{ring}%
\TextAt{18}{353.495}{629.215}{Z}%
\TextAt{9}{361.465}{629.215}{\Character{91}}%
\TextAt{13}{364.717}{629.215}{i}%
\TextAt{9}{368.71}{629.215}{\Character{93}}%
\TextAt{8}{376.821}{629.215}{consists}%
\TextAt{8}{420.894}{629.215}{of}%
\TextAt{8}{435.186}{629.215}{the}%
\TextAt{8}{456.305}{629.215}{fourth}%
\TextAt{8}{492.714}{629.215}{roots}%
\TextAt{8}{523.337}{629.215}{of}%
\TextAt{8}{79.226}{612.169}{unity,}%
\TextAt{18}{251.91}{595.123}{Z}%
\TextAt{9}{259.881}{595.123}{\Character{91}}%
\TextAt{13}{263.132}{595.123}{i}%
\TextAt{9}{267.126}{595.123}{\Character{93}}%
\TextAt{3}{270.377}{600.059}{\Character{3}}%
\TextAt{9}{278.431}{595.123}{\Character{61}}%
\TextAt{16}{290.856}{595.123}{\Character{102}}%
\TextAt{9}{296.834}{595.123}{1}%
\TextAt{13}{302.687}{595.123}{\Character{59}}%
\TextAt{16}{307.931}{595.123}{\Character{0}}%
\TextAt{9}{317.229}{595.123}{1}%
\TextAt{13}{323.082}{595.123}{\Character{59}}%
\TextAt{13}{328.306}{595.123}{i\Character{59}}%
\TextAt{16}{337.564}{595.123}{\Character{0}}%
\TextAt{13}{346.863}{595.123}{i}%
\TextAt{16}{350.856}{595.123}{\Character{103}}%
\TextAt{13}{356.834}{595.123}{\Character{58}}%
\TextAt{8}{97.159}{563.648}{In}%
\TextAt{8}{112.098}{563.648}{order}%
\TextAt{8}{142.985}{563.648}{to}%
\TextAt{8}{157.613}{563.648}{answer}%
\TextAt{8}{196.689}{563.648}{the}%
\TextAt{8}{217.163}{563.648}{question}%
\TextAt{8}{264.046}{563.648}{for}%
\TextAt{8}{282.249}{563.648}{primes,}%
\TextAt{8}{323.676}{563.648}{i.e.,}%
\TextAt{8}{346.179}{563.648}{irreducible}%
\TextAt{8}{404.371}{563.648}{elements}%
\TextAt{8}{452.545}{563.648}{of}%
\TextAt{8}{466.193}{563.648}{the}%
\TextAt{8}{486.667}{563.648}{ring}%
\TextAt{18}{511.051}{563.648}{Z}%
\TextAt{9}{519.021}{563.648}{\Character{91}}%
\TextAt{13}{522.273}{563.648}{i}%
\TextAt{9}{526.266}{563.648}{\Character{93}}%
\TextAt{8}{529.518}{563.648}{,}%
\TextAt{8}{79.226}{546.602}{we}%
\TextAt{8}{96.019}{546.602}{first}%
\TextAt{8}{119.71}{546.602}{recall}%
\TextAt{8}{150.49}{546.602}{that}%
\TextAt{8}{175.424}{546.602}{two}%
\TextAt{8}{197.107}{546.602}{elements}%
\TextAt{13}{244.529}{546.602}{\Character{11}\Character{59}}%
\TextAt{13}{257.261}{546.602}{\Character{12}}%
\TextAt{8}{268.028}{546.602}{in}%
\TextAt{8}{281.247}{546.602}{a}%
\TextAt{8}{290.567}{546.602}{ring}%
\TextAt{8}{314.199}{546.602}{are}%
\TextAt{8}{333.275}{546.602}{called}%
\TextAt{12}{366.004}{546.602}{associated}%
\TextAt{8}{425.361}{546.602}{,}%
\TextAt{8}{432.163}{546.602}{symbolically}%
\TextAt{13}{498.785}{546.602}{\Character{11}}%
\NativeMathAt{509.627}{546.602}{11.955}{\sim}%
\TextAt{13}{522.246}{546.602}{\Character{12}}%
\TextAt{8}{529.518}{546.602}{,}%
\TextAt{8}{79.226}{529.556}{if}%
\TextAt{8}{90.641}{529.556}{they}%
\TextAt{8}{117.668}{529.556}{differ}%
\TextAt{8}{148.594}{529.556}{only}%
\TextAt{8}{174.976}{529.556}{by}%
\TextAt{8}{191.926}{529.556}{a}%
\TextAt{8}{202.372}{529.556}{unit}%
\TextAt{8}{227.774}{529.556}{factor,}%
\TextAt{8}{265.387}{529.556}{and}%
\TextAt{8}{288.841}{529.556}{that}%
\TextAt{8}{314.9}{529.556}{every}%
\TextAt{8}{346.483}{529.556}{element}%
\TextAt{8}{390.416}{529.556}{associated}%
\TextAt{8}{446.841}{529.556}{to}%
\TextAt{8}{461.842}{529.556}{an}%
\TextAt{8}{478.792}{529.556}{irreducible}%
\TextAt{8}{79.226}{512.51}{element}%
\TextAt{13}{122.755}{512.51}{\Character{25}}%
\TextAt{8}{134.009}{512.51}{is}%
\TextAt{8}{146.059}{512.51}{also}%
\TextAt{8}{169.826}{512.51}{irreducible.}%
\TextAt{8}{233.113}{512.51}{Using}%
\TextAt{8}{266.301}{512.51}{theorem}%
\TextAt{8}{312.113}{512.51}{(1.1)}%
\TextAt{8}{340.375}{512.51}{we}%
\TextAt{8}{357.889}{512.51}{obtain}%
\TextAt{8}{394.603}{512.51}{the}%
\TextAt{8}{415.047}{512.51}{following}%
\TextAt{8}{464.768}{512.51}{precise}%
\TextAt{8}{503.479}{512.51}{list}%
\TextAt{8}{523.337}{512.51}{of}%
\TextAt{8}{79.226}{495.464}{all}%
\TextAt{8}{95.485}{495.464}{prime}%
\TextAt{8}{128.649}{495.464}{numbers}%
\TextAt{8}{175.86}{495.464}{of}%
\TextAt{18}{189.19}{495.464}{Z}%
\TextAt{9}{197.16}{495.464}{\Character{91}}%
\TextAt{13}{200.412}{495.464}{i}%
\TextAt{9}{204.405}{495.464}{\Character{93}}%
\TextAt{8}{207.657}{495.464}{.}%
\TextAt{12}{79.226}{469.858}{Theorem}%
\TextAt{12}{136.838}{469.858}{(1.4).}%
\TextAt{8}{172.674}{469.858}{The}%
\TextAt{8}{197.247}{469.858}{prime}%
\TextAt{8}{230.931}{469.858}{elements}%
\TextAt{13}{279.307}{469.858}{\Character{25}}%
\TextAt{8}{290.794}{469.858}{of}%
\TextAt{18}{304.643}{469.858}{Z}%
\TextAt{9}{312.614}{469.858}{\Character{91}}%
\TextAt{13}{315.865}{469.858}{i}%
\TextAt{9}{319.859}{469.858}{\Character{93}}%
\TextAt{8}{323.11}{469.858}{,}%
\TextAt{8}{330.909}{469.858}{up}%
\TextAt{8}{348.334}{469.858}{to}%
\TextAt{8}{363.164}{469.858}{associated}%
\TextAt{8}{419.419}{469.858}{elements,}%
\TextAt{8}{471.177}{469.858}{are}%
\TextAt{8}{491.208}{469.858}{given}%
\TextAt{8}{522.297}{469.858}{as}%
\TextAt{8}{79.226}{452.812}{follows.}%
\TextAt{8}{86.973}{432.777}{(1)}%
\TextAt{13}{107.918}{432.777}{\Character{25}}%
\TextAt{9}{118.309}{432.777}{\Character{61}}%
\TextAt{9}{130.719}{432.777}{1}%
\TextAt{9}{139.219}{432.777}{\Character{43}}%
\TextAt{13}{151.005}{432.777}{i}%
\TextAt{8}{154.998}{432.777}{,}%
\TextAt{8}{86.973}{415.731}{(2)}%
\TextAt{13}{107.918}{415.731}{\Character{25}}%
\TextAt{9}{118.309}{415.731}{\Character{61}}%
\TextAt{13}{130.734}{415.731}{a}%
\TextAt{9}{139.536}{415.731}{\Character{43}}%
\TextAt{13}{151.297}{415.731}{bi}%
\TextAt{8}{164.165}{415.731}{with}%
\TextAt{13}{190.825}{415.731}{a}%
\TextAt{17}{196.97}{420.069}{2}%
\TextAt{9}{204.359}{415.731}{\Character{43}}%
\TextAt{13}{216.12}{415.731}{b}%
\TextAt{17}{221.097}{420.069}{2}%
\TextAt{9}{229.15}{415.731}{\Character{61}}%
\TextAt{13}{241.576}{415.731}{p}%
\TextAt{8}{247.451}{415.731}{,}%
\TextAt{13}{254.6}{415.731}{p}%
\NativeMathAt{263.796}{415.731}{11.955}{\equiv}%
\TextAt{9}{276.416}{415.731}{1}%
\TextAt{8}{285.589}{415.731}{mod}%
\TextAt{9}{311.027}{415.731}{4}%
\TextAt{8}{316.88}{415.731}{,}%
\TextAt{13}{324.029}{415.731}{a}%
\TextAt{13}{333.486}{415.731}{\Character{62}}%
\TextAt{16}{345.921}{415.731}{\Character{106}}%
\TextAt{13}{349.242}{415.731}{b}%
\TextAt{16}{354.219}{415.731}{\Character{106}}%
\TextAt{13}{360.86}{415.731}{\Character{62}}%
\TextAt{9}{373.286}{415.731}{0}%
\TextAt{8}{379.139}{415.731}{,}%
\TextAt{8}{86.973}{398.685}{(3)}%
\TextAt{13}{107.918}{398.685}{\Character{25}}%
\TextAt{9}{118.309}{398.685}{\Character{61}}%
\TextAt{13}{130.734}{398.685}{p}%
\TextAt{8}{136.609}{398.685}{,}%
\TextAt{13}{143.758}{398.685}{p}%
\NativeMathAt{152.954}{398.685}{11.955}{\equiv}%
\TextAt{9}{165.574}{398.685}{3}%
\TextAt{8}{174.748}{398.685}{mod}%
\TextAt{9}{200.185}{398.685}{4}%
\TextAt{8}{206.038}{398.685}{.}%
\TextAt{8}{79.226}{378.65}{Here,}%
\TextAt{13}{110.106}{378.65}{p}%
\TextAt{8}{119.879}{378.65}{denotes}%
\TextAt{8}{162.212}{378.65}{a}%
\TextAt{8}{171.968}{378.65}{prime}%
\TextAt{8}{205.131}{378.65}{number}%
\TextAt{8}{247.727}{378.65}{of}%
\TextAt{18}{261.057}{378.65}{Z}%
\TextAt{8}{269.028}{378.65}{.}%
\TextAt{19}{79.226}{344.482}{Proof.}%
\TextAt{8}{116}{344.482}{Numbers}%
\TextAt{8}{167.031}{344.482}{as}%
\TextAt{8}{182.626}{344.482}{in}%
\TextAt{8}{197.504}{344.482}{(1)}%
\TextAt{8}{217.595}{344.482}{or}%
\TextAt{8}{233.13}{344.482}{(2)}%
\TextAt{8}{253.221}{344.482}{are}%
\TextAt{8}{273.957}{344.482}{prime}%
\TextAt{8}{308.345}{344.482}{because}%
\TextAt{8}{352.872}{344.482}{a}%
\TextAt{8}{363.853}{344.482}{decomposition}%
\TextAt{13}{442.213}{344.482}{\Character{25}}%
\TextAt{9}{454.691}{344.482}{\Character{61}}%
\TextAt{13}{469.204}{344.482}{\Character{11}}%
\TextAt{16}{480.218}{344.482}{\Character{1}}%
\TextAt{13}{487.031}{344.482}{\Character{12}}%
\TextAt{8}{499.425}{344.482}{in}%
\TextAt{18}{514.303}{344.482}{Z}%
\TextAt{9}{522.273}{344.482}{\Character{91}}%
\TextAt{13}{525.524}{344.482}{i}%
\TextAt{9}{529.518}{344.482}{\Character{93}}%
\TextAt{8}{79.226}{327.436}{implies}%
\TextAt{8}{118.953}{327.436}{an}%
\TextAt{8}{135.212}{327.436}{equation}%
\TextAt{13}{245.166}{310.39}{p}%
\TextAt{9}{254.362}{310.39}{\Character{61}}%
\TextAt{13}{266.788}{310.39}{N}%
\TextAt{9}{277.41}{310.39}{\Character{40}}%
\TextAt{13}{281.962}{310.39}{\Character{25}}%
\TextAt{9}{289.032}{310.39}{\Character{41}}%
\TextAt{9}{296.887}{310.39}{\Character{61}}%
\TextAt{13}{309.33}{310.39}{N}%
\TextAt{9}{319.953}{310.39}{\Character{40}}%
\TextAt{13}{324.505}{310.39}{\Character{11}}%
\TextAt{9}{332.027}{310.39}{\Character{41}}%
\TextAt{13}{336.579}{310.39}{N}%
\TextAt{9}{347.202}{310.39}{\Character{40}}%
\TextAt{13}{351.754}{310.39}{\Character{12}}%
\TextAt{9}{359.026}{310.39}{\Character{41}}%
\TextAt{13}{363.578}{310.39}{\Character{59}}%
\TextAt{8}{79.226}{287.477}{with}%
\TextAt{8}{105.824}{287.477}{some}%
\TextAt{8}{135.088}{287.477}{prime}%
\TextAt{8}{168.19}{287.477}{number}%
\TextAt{13}{210.724}{287.477}{p}%
\TextAt{8}{216.599}{287.477}{.}%
\TextAt{8}{225.027}{287.477}{Hence}%
\TextAt{8}{259.743}{287.477}{either}%
\TextAt{13}{292.844}{287.477}{N}%
\TextAt{9}{303.467}{287.477}{\Character{40}}%
\TextAt{13}{308.019}{287.477}{\Character{11}}%
\TextAt{9}{315.541}{287.477}{\Character{41}}%
\TextAt{9}{323.396}{287.477}{\Character{61}}%
\TextAt{9}{335.817}{287.477}{1}%
\TextAt{8}{345.528}{287.477}{or}%
\TextAt{13}{359.777}{287.477}{N}%
\TextAt{9}{370.399}{287.477}{\Character{40}}%
\TextAt{13}{374.951}{287.477}{\Character{12}}%
\TextAt{9}{382.223}{287.477}{\Character{41}}%
\TextAt{9}{390.078}{287.477}{\Character{61}}%
\TextAt{9}{402.499}{287.477}{1}%
\TextAt{8}{408.374}{287.477}{,}%
\TextAt{8}{415.474}{287.477}{so}%
\TextAt{8}{429.782}{287.477}{that}%
\TextAt{8}{455.089}{287.477}{either}%
\TextAt{13}{488.191}{287.477}{\Character{11}}%
\TextAt{8}{499.548}{287.477}{or}%
\TextAt{13}{513.796}{287.477}{\Character{12}}%
\TextAt{8}{524.903}{287.477}{is}%
\TextAt{8}{79.226}{270.431}{a}%
\TextAt{8}{88.982}{270.431}{unit.}%
\TextAt{8}{97.159}{250.396}{Numbers}%
\TextAt{13}{148.586}{250.396}{\Character{25}}%
\TextAt{9}{161.74}{250.396}{\Character{61}}%
\TextAt{13}{176.929}{250.396}{p}%
\TextAt{8}{182.804}{250.396}{,}%
\TextAt{8}{191.981}{250.396}{where}%
\TextAt{13}{227.412}{250.396}{p}%
\NativeMathAt{239.371}{250.396}{11.955}{\equiv}%
\TextAt{9}{254.754}{250.396}{3}%
\TextAt{8}{263.928}{250.396}{mod}%
\TextAt{9}{289.366}{250.396}{4}%
\TextAt{8}{295.219}{250.396}{,}%
\TextAt{8}{304.396}{250.396}{are}%
\TextAt{8}{325.528}{250.396}{prime}%
\TextAt{8}{360.314}{250.396}{in}%
\TextAt{18}{375.588}{250.396}{Z}%
\TextAt{9}{383.558}{250.396}{\Character{91}}%
\TextAt{13}{386.81}{250.396}{i}%
\TextAt{9}{390.803}{250.396}{\Character{93}}%
\TextAt{8}{394.055}{250.396}{,}%
\TextAt{8}{403.232}{250.396}{because}%
\TextAt{8}{448.155}{250.396}{a}%
\TextAt{8}{459.532}{250.396}{decomposition}%
\TextAt{13}{79.226}{233.35}{p}%
\TextAt{9}{90.137}{233.35}{\Character{61}}%
\TextAt{13}{104.278}{233.35}{\Character{11}}%
\TextAt{16}{115.142}{233.35}{\Character{1}}%
\TextAt{13}{121.805}{233.35}{\Character{12}}%
\TextAt{8}{133.98}{233.35}{into}%
\TextAt{8}{158.73}{233.35}{non-units}%
\TextAt{13}{211.825}{233.35}{\Character{11}\Character{59}}%
\TextAt{13}{224.569}{233.35}{\Character{12}}%
\TextAt{8}{236.766}{233.35}{would}%
\TextAt{8}{271.916}{233.35}{imply}%
\TextAt{8}{305.763}{233.35}{that}%
\TextAt{13}{332.139}{233.35}{p}%
\TextAt{17}{338.014}{237.689}{2}%
\TextAt{9}{347.782}{233.35}{\Character{61}}%
\TextAt{13}{361.922}{233.35}{N}%
\TextAt{9}{372.545}{233.35}{\Character{40}}%
\TextAt{13}{377.097}{233.35}{\Character{11}}%
\TextAt{9}{384.619}{233.35}{\Character{41}}%
\TextAt{13}{389.171}{233.35}{N}%
\TextAt{9}{399.794}{233.35}{\Character{40}}%
\TextAt{13}{404.346}{233.35}{\Character{12}}%
\TextAt{9}{411.617}{233.35}{\Character{41}}%
\TextAt{8}{416.17}{233.35}{,}%
\TextAt{8}{424.577}{233.35}{so}%
\TextAt{8}{439.953}{233.35}{that}%
\TextAt{13}{466.329}{233.35}{p}%
\TextAt{9}{477.24}{233.35}{\Character{61}}%
\TextAt{13}{491.38}{233.35}{N}%
\TextAt{9}{502.003}{233.35}{\Character{40}}%
\TextAt{13}{506.555}{233.35}{\Character{11}}%
\TextAt{9}{514.077}{233.35}{\Character{41}}%
\TextAt{9}{523.653}{233.35}{\Character{61}}%
\TextAt{13}{79.226}{216.304}{N}%
\TextAt{9}{89.849}{216.304}{\Character{40}}%
\TextAt{13}{94.401}{216.304}{a}%
\TextAt{9}{103.203}{216.304}{\Character{43}}%
\TextAt{13}{114.964}{216.304}{bi}%
\TextAt{9}{123.935}{216.304}{\Character{41}}%
\TextAt{9}{131.79}{216.304}{\Character{61}}%
\TextAt{13}{144.233}{216.304}{a}%
\TextAt{17}{150.378}{220.643}{2}%
\TextAt{9}{157.767}{216.304}{\Character{43}}%
\TextAt{13}{169.528}{216.304}{b}%
\TextAt{17}{174.506}{220.643}{2}%
\TextAt{8}{179.238}{216.304}{,}%
\TextAt{8}{186.387}{216.304}{which}%
\TextAt{8}{219.873}{216.304}{according}%
\TextAt{8}{272.56}{216.304}{to}%
\TextAt{8}{286.87}{216.304}{(1.1)}%
\TextAt{8}{314.845}{216.304}{would}%
\TextAt{8}{348.989}{216.304}{yield}%
\TextAt{13}{377.275}{216.304}{p}%
\NativeMathAt{386.471}{216.304}{11.955}{\equiv}%
\TextAt{9}{399.091}{216.304}{1}%
\TextAt{8}{408.264}{216.304}{mod}%
\TextAt{9}{433.702}{216.304}{4}%
\TextAt{8}{439.555}{216.304}{.}%
\TextAt{8}{97.159}{196.269}{This}%
\TextAt{8}{122.729}{196.269}{being}%
\TextAt{8}{153.118}{196.269}{said,}%
\TextAt{8}{179.576}{196.269}{we}%
\TextAt{8}{195.654}{196.269}{have}%
\TextAt{8}{221.499}{196.269}{to}%
\TextAt{8}{234.66}{196.269}{check}%
\TextAt{8}{265.048}{196.269}{that}%
\TextAt{8}{289.268}{196.269}{an}%
\TextAt{8}{304.378}{196.269}{arbitrary}%
\TextAt{8}{352.998}{196.269}{prime}%
\TextAt{8}{385.012}{196.269}{element}%
\TextAt{13}{427.104}{196.269}{\Character{25}}%
\TextAt{8}{436.922}{196.269}{of}%
\TextAt{18}{449.102}{196.269}{Z}%
\TextAt{9}{457.073}{196.269}{\Character{91}}%
\TextAt{13}{460.324}{196.269}{i}%
\TextAt{9}{464.318}{196.269}{\Character{93}}%
\TextAt{8}{470.317}{196.269}{is}%
\TextAt{8}{480.932}{196.269}{associated}%
\TextAt{8}{79.226}{179.223}{to}%
\TextAt{8}{93.536}{179.223}{one}%
\TextAt{8}{114.996}{179.223}{of}%
\TextAt{8}{128.326}{179.223}{those}%
\TextAt{8}{158.955}{179.223}{listed.}%
\TextAt{8}{194.781}{179.223}{First}%
\TextAt{8}{223.294}{179.223}{of}%
\TextAt{8}{236.624}{179.223}{all,}%
\TextAt{8}{256.135}{179.223}{the}%
\TextAt{8}{276.291}{179.223}{decomposition}%
\ErrataAt{206}{151.655}{200}{\makebox[200bp][c]{$N(\pi)=\pi\cdot\bar\pi=p_1\cdots p_r$,}}%
\TextAt{8}{79.226}{124.059}{with}%
\TextAt{8}{106.224}{124.059}{rational}%
\TextAt{8}{150.15}{124.059}{primes}%
\TextAt{13}{188.266}{124.059}{p}%
\TextAt{21}{194.141}{122.266}{i}%
\TextAt{8}{197.522}{124.059}{,}%
\TextAt{8}{205.094}{124.059}{shows}%
\TextAt{8}{239.049}{124.059}{that}%
\TextAt{13}{264.756}{124.059}{\Character{25}}%
\TextAt{16}{275.721}{124.059}{\Character{106}}%
\TextAt{13}{282.939}{124.059}{p}%
\TextAt{8}{293.049}{124.059}{for}%
\TextAt{8}{311.271}{124.059}{some}%
\TextAt{13}{340.935}{124.059}{p}%
\TextAt{9}{350.706}{124.059}{\Character{61}}%
\TextAt{13}{363.707}{124.059}{p}%
\TextAt{21}{369.582}{122.266}{i}%
\TextAt{8}{372.964}{124.059}{.}%
\TextAt{8}{382.425}{124.059}{This}%
\TextAt{8}{409.483}{124.059}{gives}%
\TextAt{13}{438.501}{124.059}{N}%
\TextAt{9}{449.123}{124.059}{\Character{40}}%
\TextAt{13}{453.676}{124.059}{\Character{25}}%
\TextAt{9}{460.745}{124.059}{\Character{41}}%
\TextAt{16}{469.193}{124.059}{\Character{106}}%
\TextAt{13}{476.411}{124.059}{N}%
\TextAt{9}{487.033}{124.059}{\Character{40}}%
\TextAt{13}{491.586}{124.059}{p}%
\TextAt{9}{497.461}{124.059}{\Character{41}}%
\TextAt{9}{505.889}{124.059}{\Character{61}}%
\TextAt{13}{518.91}{124.059}{p}%
\TextAt{17}{524.785}{128.398}{2}%
\TextAt{8}{529.518}{124.059}{,}%
\TextAt{8}{79.226}{107.013}{so}%
\TextAt{8}{94.423}{107.013}{that}%
\TextAt{8}{120.619}{107.013}{either}%
\TextAt{13}{154.61}{107.013}{N}%
\TextAt{9}{165.232}{107.013}{\Character{40}}%
\TextAt{13}{169.784}{107.013}{\Character{25}}%
\TextAt{9}{176.854}{107.013}{\Character{41}}%
\TextAt{9}{186.119}{107.013}{\Character{61}}%
\TextAt{13}{199.971}{107.013}{p}%
\TextAt{8}{210.57}{107.013}{or}%
\TextAt{13}{225.708}{107.013}{N}%
\TextAt{9}{236.33}{107.013}{\Character{40}}%
\TextAt{13}{240.882}{107.013}{\Character{25}}%
\TextAt{9}{247.952}{107.013}{\Character{41}}%
\TextAt{9}{257.217}{107.013}{\Character{61}}%
\TextAt{13}{271.069}{107.013}{p}%
\TextAt{17}{276.944}{111.352}{2}%
\TextAt{8}{281.676}{107.013}{.}%
\TextAt{8}{292.605}{107.013}{In}%
\TextAt{8}{308.053}{107.013}{the}%
\TextAt{8}{329.037}{107.013}{case}%
\TextAt{13}{354.635}{107.013}{N}%
\TextAt{9}{365.258}{107.013}{\Character{40}}%
\TextAt{13}{369.81}{107.013}{\Character{25}}%
\TextAt{9}{376.879}{107.013}{\Character{41}}%
\TextAt{9}{386.144}{107.013}{\Character{61}}%
\TextAt{13}{399.996}{107.013}{p}%
\TextAt{8}{410.596}{107.013}{we}%
\TextAt{8}{428.65}{107.013}{get}%
\TextAt{13}{448.988}{107.013}{\Character{25}}%
\TextAt{9}{460.787}{107.013}{\Character{61}}%
\TextAt{13}{474.622}{107.013}{a}%
\TextAt{9}{483.987}{107.013}{\Character{43}}%
\TextAt{13}{496.312}{107.013}{bi}%
\TextAt{8}{510.007}{107.013}{with}%
\TextAt{13}{79.226}{89.967}{a}%
\TextAt{17}{85.371}{94.306}{2}%
\TextAt{9}{92.898}{89.967}{\Character{43}}%
\TextAt{13}{104.798}{89.967}{b}%
\TextAt{17}{109.775}{94.306}{2}%
\TextAt{9}{118.174}{89.967}{\Character{61}}%
\TextAt{13}{130.945}{89.967}{p}%
\TextAt{8}{136.82}{89.967}{,}%
\TextAt{8}{144.222}{89.967}{so}%
\TextAt{13}{158.795}{89.967}{\Character{25}}%
\TextAt{8}{169.965}{89.967}{is}%
\TextAt{8}{181.931}{89.967}{of}%
\TextAt{8}{195.464}{89.967}{type}%
\TextAt{8}{222.004}{89.967}{(2)}%
\TextAt{8}{241.072}{89.967}{or,}%
\TextAt{8}{258.888}{89.967}{if}%
\TextAt{13}{269.814}{89.967}{p}%
\TextAt{9}{279.356}{89.967}{\Character{61}}%
\TextAt{9}{292.112}{89.967}{2}%
\TextAt{8}{297.98}{89.967}{,}%
\TextAt{8}{305.382}{89.967}{it}%
\TextAt{8}{317.289}{89.967}{is}%
\TextAt{8}{329.256}{89.967}{associated}%
\TextAt{8}{385.193}{89.967}{to}%
\TextAt{9}{399.707}{89.967}{1}%
\TextAt{9}{408.339}{89.967}{\Character{43}}%
\TextAt{13}{420.254}{89.967}{i}%
\TextAt{8}{424.247}{89.967}{.}%
\TextAt{8}{433.304}{89.967}{On}%
\TextAt{8}{453.018}{89.967}{the}%
\TextAt{8}{473.377}{89.967}{other}%
\TextAt{8}{504.149}{89.967}{hand,}%
\TextAt{8}{79.226}{72.921}{if}%
\TextAt{13}{90.05}{72.921}{N}%
\TextAt{9}{100.673}{72.921}{\Character{40}}%
\TextAt{13}{105.225}{72.921}{\Character{25}}%
\TextAt{9}{112.294}{72.921}{\Character{41}}%
\TextAt{9}{120.328}{72.921}{\Character{61}}%
\TextAt{13}{132.934}{72.921}{p}%
\TextAt{17}{138.809}{77.26}{2}%
\TextAt{8}{143.541}{72.921}{,}%
\TextAt{8}{150.815}{72.921}{then}%
\TextAt{13}{177.575}{72.921}{\Character{25}}%
\TextAt{8}{188.642}{72.921}{is}%
\TextAt{8}{200.506}{72.921}{associated}%
\TextAt{8}{256.341}{72.921}{to}%
\TextAt{13}{270.751}{72.921}{p}%
\TextAt{8}{280.623}{72.921}{since}%
\TextAt{13}{309.392}{72.921}{p}%
\TextAt{9}{315.267}{72.921}{\Character{47}}%
\TextAt{13}{321.12}{72.921}{\Character{25}}%
\TextAt{8}{332.187}{72.921}{is}%
\TextAt{8}{344.051}{72.921}{an}%
\TextAt{8}{360.41}{72.921}{integer}%
\TextAt{8}{399.208}{72.921}{with}%
\TextAt{8}{425.969}{72.921}{norm}%
\TextAt{8}{456.638}{72.921}{one}%
\TextAt{8}{478.197}{72.921}{and}%
\TextAt{8}{501.06}{72.921}{thus}%
\TextAt{8}{526.911}{72.921}{a}%
\TextAt{8}{303.069}{43.033}{3}%
\Destination{79.23}{449.705}{nav.Item.1}%
\Destination{79.23}{429.775}{nav.Item.2}%
\Destination{79.23}{412.725}{nav.Item.3}%
\Destination{78.23}{809.815}{nav.page.3}%
\Destination{79.36}{776.155}{nav.print.3}%
\Destination{79.36}{211.255}{nav.print.4}%
\Destination{79.36}{643.255}{nav.number.1.1.3}%
\Destination{79.36}{483.855}{nav.number.1.1.4}%
\Hotlink{126.583}{652.655}{6.858}{11.047}{nav.ex.1.1.1}%
\Hotlink{311.613}{510.355}{25.078}{11.047}{nav.number.1.1.1}%
\Hotlink{286.37}{214.055}{25.078}{11.047}{nav.number.1.1.1}%
\Destination{79.226}{165.655}{errata.40}%
}

% Page 16 (body page 4)
\Page{%
\TextAt{2}{79.226}{796.943}{I.}%
\TextAt{2}{89.83}{796.943}{Algebraic}%
\TextAt{2}{138.648}{796.943}{Integers}%
\Vector{q 0 G 0 g 0.299 w 7.226 18.729 m 460.769 18.729 l S Q}%
\TextAt{8}{79.226}{761.124}{unit.}%
\TextAt{8}{110.284}{761.124}{Moreover,}%
\TextAt{13}{165.667}{761.124}{p}%
\NativeMathAt{175.882}{761.124}{11.955}{\equiv}%
\TextAt{9}{189.521}{761.124}{3}%
\TextAt{8}{198.695}{761.124}{mod}%
\TextAt{9}{224.133}{761.124}{4}%
\TextAt{8}{234.482}{761.124}{has}%
\TextAt{8}{255.954}{761.124}{to}%
\TextAt{8}{270.863}{761.124}{hold}%
\TextAt{8}{297.475}{761.124}{in}%
\TextAt{8}{311.727}{761.124}{this}%
\TextAt{8}{335.147}{761.124}{case}%
\TextAt{8}{360.517}{761.124}{because}%
\TextAt{8}{404.417}{761.124}{otherwise}%
\TextAt{8}{457.104}{761.124}{we}%
\TextAt{8}{474.93}{761.124}{would}%
\TextAt{8}{509.672}{761.124}{have}%
\TextAt{13}{79.226}{744.078}{p}%
\TextAt{9}{89.438}{744.078}{\Character{61}}%
\TextAt{9}{102.864}{744.078}{2}%
\TextAt{8}{113.225}{744.078}{or}%
\TextAt{13}{128.131}{744.078}{p}%
\NativeMathAt{138.342}{744.078}{11.955}{\equiv}%
\TextAt{9}{151.977}{744.078}{1}%
\TextAt{8}{161.151}{744.078}{mod}%
\TextAt{9}{186.589}{744.078}{4}%
\TextAt{8}{196.935}{744.078}{and}%
\TextAt{8}{220.293}{744.078}{because}%
\TextAt{8}{264.191}{744.078}{of}%
\TextAt{8}{278.117}{744.078}{(1.1)}%
\TextAt{13}{306.688}{744.078}{p}%
\TextAt{9}{316.899}{744.078}{\Character{61}}%
\TextAt{13}{330.34}{744.078}{a}%
\TextAt{17}{336.485}{748.416}{2}%
\TextAt{9}{344.28}{744.078}{\Character{43}}%
\TextAt{13}{356.448}{744.078}{b}%
\TextAt{17}{361.425}{748.416}{2}%
\TextAt{9}{370.493}{744.078}{\Character{61}}%
\TextAt{9}{383.919}{744.078}{\Character{40}}%
\TextAt{13}{388.486}{744.078}{a}%
\TextAt{9}{397.694}{744.078}{\Character{43}}%
\TextAt{13}{409.862}{744.078}{bi}%
\TextAt{9}{418.832}{744.078}{\Character{41}\Character{40}}%
\TextAt{13}{427.937}{744.078}{a}%
\TextAt{16}{437.145}{744.078}{\Character{0}}%
\TextAt{13}{449.506}{744.078}{bi}%
\TextAt{9}{458.477}{744.078}{\Character{41}}%
\TextAt{8}{467.522}{744.078}{could}%
\TextAt{8}{499.333}{744.078}{not}%
\TextAt{8}{520.743}{744.078}{be}%
\TextAt{8}{79.226}{727.032}{prime.}%
\TextAt{8}{116.941}{727.032}{This}%
\TextAt{8}{143.661}{727.032}{completes}%
\TextAt{8}{197.698}{727.032}{the}%
\TextAt{8}{217.854}{727.032}{proof.}%
\Vector{q 0 G 0 g 0.399 w 452.334 -46.968 m 452.334 -38.898 l S Q}%
\Vector{q 0 G 0 g 0.399 w 452.534 -39.098 m 459.707 -39.098 l S Q}%
\Vector{q 0 G 0 g 0.399 w 452.534 -46.769 m 459.707 -46.769 l S Q}%
\Vector{q 0 G 0 g 0.399 w 459.906 -46.968 m 459.906 -38.898 l S Q}%
\TextAt{8}{97.159}{698.031}{The}%
\TextAt{8}{122.499}{698.031}{proposition}%
\TextAt{8}{185.318}{698.031}{also}%
\TextAt{8}{210.083}{698.031}{settles}%
\TextAt{8}{247.259}{698.031}{completely}%
\TextAt{8}{307.4}{698.031}{the}%
\TextAt{8}{328.842}{698.031}{question}%
\TextAt{8}{376.693}{698.031}{of}%
\TextAt{8}{391.309}{698.031}{how}%
\TextAt{8}{416.984}{698.031}{prime}%
\TextAt{8}{451.433}{698.031}{numbers}%
\TextAt{13}{499.93}{698.031}{p}%
\NativeMathAt{511.317}{698.031}{11.955}{\in}%
\TextAt{18}{524.799}{698.031}{Z}%
\TextAt{8}{79.226}{680.985}{decompose}%
\TextAt{8}{138.278}{680.985}{in}%
\TextAt{18}{152.067}{680.985}{Z}%
\TextAt{9}{160.037}{680.985}{\Character{91}}%
\TextAt{13}{163.289}{680.985}{i}%
\TextAt{9}{167.282}{680.985}{\Character{93}}%
\TextAt{8}{170.534}{680.985}{.}%
\TextAt{8}{179.392}{680.985}{The}%
\TextAt{8}{203.582}{680.985}{prime}%
\TextAt{9}{236.883}{680.985}{2}%
\TextAt{9}{246.28}{680.985}{\Character{61}}%
\TextAt{9}{258.928}{680.985}{\Character{40}1}%
\TextAt{9}{272.067}{680.985}{\Character{43}}%
\TextAt{13}{283.957}{680.985}{i}%
\TextAt{9}{287.95}{680.985}{\Character{41}\Character{40}1}%
\TextAt{16}{305.657}{680.985}{\Character{0}}%
\TextAt{13}{317.706}{680.985}{i}%
\TextAt{9}{321.699}{680.985}{\Character{41}}%
\TextAt{8}{330.285}{680.985}{is}%
\TextAt{8}{342.186}{680.985}{associated}%
\TextAt{8}{398.057}{680.985}{to}%
\TextAt{8}{412.504}{680.985}{the}%
\TextAt{8}{432.797}{680.985}{square}%
\TextAt{8}{469.744}{680.985}{of}%
\TextAt{8}{483.21}{680.985}{the}%
\TextAt{8}{503.503}{680.985}{prime}%
\TextAt{8}{79.226}{663.939}{element}%
\TextAt{9}{122.524}{663.939}{1}%
\TextAt{9}{131.06}{663.939}{\Character{43}}%
\TextAt{13}{142.871}{663.939}{i}%
\TextAt{8}{146.865}{663.939}{.}%
\TextAt{8}{155.481}{663.939}{Indeed,}%
\TextAt{8}{196.832}{663.939}{the}%
\TextAt{8}{217.045}{663.939}{identity}%
\TextAt{9}{260.354}{663.939}{1}%
\TextAt{16}{268.902}{663.939}{\Character{0}}%
\TextAt{13}{280.896}{663.939}{i}%
\TextAt{9}{288.305}{663.939}{\Character{61}}%
\TextAt{16}{300.826}{663.939}{\Character{0}}%
\TextAt{13}{310.125}{663.939}{i}%
\TextAt{9}{314.118}{663.939}{\Character{40}1}%
\TextAt{9}{327.197}{663.939}{\Character{43}}%
\TextAt{13}{339.018}{663.939}{i}%
\TextAt{9}{343.011}{663.939}{\Character{41}}%
\TextAt{8}{351.517}{663.939}{shows}%
\TextAt{8}{385.191}{663.939}{that}%
\TextAt{9}{410.616}{663.939}{2}%
\NativeMathAt{419.885}{663.939}{11.955}{\sim}%
\TextAt{9}{432.599}{663.939}{\Character{40}1}%
\TextAt{9}{445.678}{663.939}{\Character{43}}%
\TextAt{13}{457.499}{663.939}{i}%
\TextAt{9}{461.493}{663.939}{\Character{41}}%
\TextAt{17}{466.045}{668.277}{2}%
\TextAt{8}{470.777}{663.939}{.}%
\TextAt{8}{479.393}{663.939}{The}%
\TextAt{8}{503.503}{663.939}{prime}%
\TextAt{8}{79.226}{646.893}{numbers}%
\TextAt{13}{126.437}{646.893}{p}%
\NativeMathAt{135.633}{646.893}{11.955}{\equiv}%
\TextAt{9}{148.252}{646.893}{1}%
\TextAt{8}{157.426}{646.893}{mod}%
\TextAt{9}{182.864}{646.893}{4}%
\TextAt{8}{192.615}{646.893}{split}%
\TextAt{8}{218.689}{646.893}{into}%
\TextAt{8}{242.432}{646.893}{two}%
\TextAt{8}{264.549}{646.893}{conjugate}%
\TextAt{8}{317.558}{646.893}{prime}%
\TextAt{8}{350.722}{646.893}{factors}%
\TextAt{13}{254.826}{617.892}{p}%
\TextAt{9}{264.022}{617.892}{\Character{61}}%
\TextAt{9}{276.432}{617.892}{\Character{40}}%
\TextAt{13}{281}{617.892}{a}%
\TextAt{9}{289.802}{617.892}{\Character{43}}%
\TextAt{13}{301.563}{617.892}{bi}%
\TextAt{9}{310.533}{617.892}{\Character{41}\Character{40}}%
\TextAt{13}{319.638}{617.892}{a}%
\TextAt{16}{328.44}{617.892}{\Character{0}}%
\TextAt{13}{340.395}{617.892}{bi}%
\TextAt{9}{349.365}{617.892}{\Character{41}}%
\TextAt{13}{353.918}{617.892}{\Character{59}}%
\TextAt{8}{79.226}{588.89}{and}%
\TextAt{8}{101.989}{588.89}{the}%
\TextAt{8}{122.145}{588.89}{prime}%
\TextAt{8}{155.309}{588.89}{numbers}%
\TextAt{13}{202.52}{588.89}{p}%
\NativeMathAt{211.716}{588.89}{11.955}{\equiv}%
\TextAt{9}{224.335}{588.89}{3}%
\TextAt{8}{233.509}{588.89}{mod}%
\TextAt{9}{258.947}{588.89}{4}%
\TextAt{8}{268.697}{588.89}{remain}%
\TextAt{8}{307.719}{588.89}{prime}%
\TextAt{8}{340.883}{588.89}{in}%
\TextAt{18}{354.535}{588.89}{Z}%
\TextAt{9}{362.505}{588.89}{\Character{91}}%
\TextAt{13}{365.757}{588.89}{i}%
\TextAt{9}{369.751}{588.89}{\Character{93}}%
\TextAt{8}{373.002}{588.89}{.}%
\TextAt{8}{97.159}{568.856}{The}%
\TextAt{8}{121.213}{568.856}{gaussian}%
\TextAt{8}{168.173}{568.856}{integers}%
\TextAt{8}{211.486}{568.856}{play}%
\TextAt{8}{236.855}{568.856}{the}%
\TextAt{8}{257.011}{568.856}{same}%
\TextAt{8}{286.337}{568.856}{role}%
\TextAt{8}{309.1}{568.856}{in}%
\TextAt{8}{322.753}{568.856}{the}%
\TextAt{8}{342.909}{568.856}{field}%
\TextAt{18}{241.063}{539.855}{Q}%
\TextAt{9}{250.361}{539.855}{\Character{40}}%
\TextAt{13}{254.913}{539.855}{i}%
\TextAt{9}{258.907}{539.855}{\Character{41}}%
\TextAt{9}{266.762}{539.855}{\Character{61}}%
\TextAt{16}{279.205}{539.855}{\Character{102}}%
\TextAt{13}{285.183}{539.855}{a}%
\TextAt{9}{293.985}{539.855}{\Character{43}}%
\TextAt{13}{305.746}{539.855}{bi}%
\TextAt{16}{318.037}{539.855}{\Character{106}}%
\TextAt{13}{324.679}{539.855}{a\Character{59}}%
\TextAt{13}{336.048}{539.855}{b}%
\NativeMathAt{344.366}{539.855}{11.955}{\in}%
\TextAt{18}{355.657}{539.855}{Q}%
\TextAt{16}{364.956}{539.855}{\Character{103}}%
\TextAt{8}{79.226}{510.853}{as}%
\TextAt{8}{93.927}{510.853}{the}%
\TextAt{8}{114.414}{510.853}{rational}%
\TextAt{8}{158.333}{510.853}{integers}%
\TextAt{8}{201.978}{510.853}{do}%
\TextAt{8}{218.567}{510.853}{in}%
\TextAt{8}{232.551}{510.853}{the}%
\TextAt{8}{253.038}{510.853}{field}%
\TextAt{18}{278.725}{510.853}{Q}%
\TextAt{8}{288.024}{510.853}{.}%
\TextAt{8}{297.464}{510.853}{So}%
\TextAt{8}{314.054}{510.853}{they}%
\TextAt{8}{340.722}{510.853}{should}%
\TextAt{8}{378.185}{510.853}{be}%
\TextAt{8}{394.44}{510.853}{viewed}%
\TextAt{8}{433.135}{510.853}{as}%
\TextAt{8}{447.836}{510.853}{the}%
\TextAt{8}{468.323}{510.853}{“integers”}%
\TextAt{8}{523.014}{510.853}{in}%
\TextAt{18}{79.226}{493.807}{Q}%
\TextAt{9}{88.525}{493.807}{\Character{40}}%
\TextAt{13}{93.077}{493.807}{i}%
\TextAt{9}{97.07}{493.807}{\Character{41}}%
\TextAt{8}{101.623}{493.807}{.}%
\TextAt{8}{111.344}{493.807}{This}%
\TextAt{8}{138.489}{493.807}{notion}%
\TextAt{8}{175.34}{493.807}{of}%
\TextAt{8}{189.095}{493.807}{integrality}%
\TextAt{8}{245.792}{493.807}{is}%
\TextAt{8}{257.981}{493.807}{relative}%
\TextAt{8}{300.033}{493.807}{to}%
\TextAt{8}{314.768}{493.807}{the}%
\TextAt{8}{335.348}{493.807}{coordinates}%
\TextAt{8}{397.952}{493.807}{of}%
\TextAt{8}{411.706}{493.807}{the}%
\TextAt{8}{432.287}{493.807}{basis}%
\TextAt{9}{461.452}{493.807}{1}%
\TextAt{13}{467.305}{493.807}{\Character{59}}%
\TextAt{13}{472.529}{493.807}{i}%
\TextAt{8}{476.542}{493.807}{.}%
\TextAt{8}{486.264}{493.807}{However,}%
\TextAt{8}{79.226}{476.761}{we}%
\TextAt{8}{96.116}{476.761}{also}%
\TextAt{8}{119.259}{476.761}{have}%
\TextAt{8}{145.916}{476.761}{the}%
\TextAt{8}{165.735}{476.761}{following}%
\TextAt{8}{214.832}{476.761}{characterization}%
\TextAt{8}{300.034}{476.761}{of}%
\TextAt{8}{313.027}{476.761}{the}%
\TextAt{8}{332.846}{476.761}{gaussian}%
\TextAt{8}{379.469}{476.761}{integers,}%
\TextAt{8}{425.764}{476.761}{which}%
\TextAt{8}{458.913}{476.761}{is}%
\TextAt{8}{470.34}{476.761}{independent}%
\TextAt{8}{79.226}{459.715}{of}%
\TextAt{8}{92.556}{459.715}{a}%
\TextAt{8}{102.312}{459.715}{choice}%
\TextAt{8}{137.101}{459.715}{of}%
\TextAt{8}{150.431}{459.715}{basis.}%
\TextAt{12}{79.226}{433.703}{Proposition}%
\TextAt{12}{152.635}{433.703}{(1.5).}%
\TextAt{18}{188.556}{433.703}{Z}%
\TextAt{9}{196.526}{433.703}{\Character{91}}%
\TextAt{13}{199.778}{433.703}{i}%
\TextAt{9}{203.771}{433.703}{\Character{93}}%
\TextAt{8}{211.527}{433.703}{consists}%
\TextAt{8}{255.243}{433.703}{precisely}%
\TextAt{8}{303.706}{433.703}{of}%
\TextAt{8}{317.642}{433.703}{those}%
\TextAt{8}{348.877}{433.703}{elements}%
\TextAt{8}{397.339}{433.703}{of}%
\TextAt{8}{411.276}{433.703}{the}%
\TextAt{8}{432.038}{433.703}{extension}%
\TextAt{8}{484.41}{433.703}{field}%
\TextAt{18}{510.373}{433.703}{Q}%
\TextAt{9}{519.671}{433.703}{\Character{40}}%
\TextAt{13}{524.224}{433.703}{i}%
\TextAt{9}{528.217}{433.703}{\Character{41}}%
\TextAt{8}{79.226}{416.657}{of}%
\TextAt{18}{92.556}{416.657}{Q}%
\TextAt{8}{105.752}{416.657}{which}%
\TextAt{8}{139.238}{416.657}{satisfy}%
\TextAt{8}{175.785}{416.657}{a}%
\TextAt{8}{185.541}{416.657}{monic}%
\TextAt{8}{220.008}{416.657}{polynomial}%
\TextAt{8}{280.501}{416.657}{equation}%
\TextAt{13}{266.201}{387.656}{x}%
\TextAt{17}{272.853}{392.592}{2}%
\TextAt{9}{280.242}{387.656}{\Character{43}}%
\TextAt{13}{292.003}{387.656}{ax}%
\TextAt{9}{307.457}{387.656}{\Character{43}}%
\TextAt{13}{319.218}{387.656}{b}%
\TextAt{9}{327.516}{387.656}{\Character{61}}%
\TextAt{9}{339.925}{387.656}{0}%
\TextAt{8}{79.226}{358.655}{with}%
\TextAt{8}{105.886}{358.655}{coefficients}%
\TextAt{13}{165.124}{358.655}{a\Character{59}}%
\TextAt{13}{176.493}{358.655}{b}%
\NativeMathAt{184.811}{358.655}{11.955}{\in}%
\TextAt{18}{196.102}{358.655}{Z}%
\TextAt{8}{204.072}{358.655}{.}%
\TextAt{19}{79.226}{323.676}{Proof.}%
\TextAt{8}{116}{323.676}{An}%
\TextAt{8}{135.119}{323.676}{element}%
\TextAt{13}{178.303}{323.676}{\Character{11}}%
\TextAt{9}{189.145}{323.676}{\Character{61}}%
\TextAt{13}{201.571}{323.676}{c}%
\TextAt{9}{209.147}{323.676}{\Character{43}}%
\TextAt{13}{220.791}{323.676}{id}%
\NativeMathAt{234.188}{323.676}{11.955}{\in}%
\TextAt{18}{245.479}{323.676}{Q}%
\TextAt{9}{254.777}{323.676}{\Character{40}}%
\TextAt{13}{259.33}{323.676}{i}%
\TextAt{9}{263.323}{323.676}{\Character{41}}%
\TextAt{8}{271.715}{323.676}{is}%
\TextAt{8}{283.421}{323.676}{a}%
\TextAt{8}{293.119}{323.676}{zero}%
\TextAt{8}{317.773}{323.676}{of}%
\TextAt{8}{331.045}{323.676}{the}%
\TextAt{8}{351.144}{323.676}{polynomial}%
\TextAt{13}{411.579}{323.676}{x}%
\TextAt{17}{418.231}{328.015}{2}%
\TextAt{9}{425.502}{323.676}{\Character{43}}%
\TextAt{13}{437.145}{323.676}{ax}%
\TextAt{9}{452.481}{323.676}{\Character{43}}%
\TextAt{13}{464.124}{323.676}{b}%
\NativeMathAt{472.422}{323.676}{11.955}{\in}%
\TextAt{18}{483.713}{323.676}{Q}%
\TextAt{9}{493.012}{323.676}{\Character{91}}%
\TextAt{13}{496.263}{323.676}{x}%
\TextAt{9}{502.916}{323.676}{\Character{93}}%
\TextAt{8}{510.007}{323.676}{with}%
\TextAt{13}{240.891}{294.675}{a}%
\TextAt{9}{250.357}{294.675}{\Character{61}}%
\TextAt{16}{262.782}{294.675}{\Character{0}}%
\TextAt{9}{272.081}{294.675}{2}%
\TextAt{13}{277.934}{294.675}{c\Character{59}}%
\TextAt{13}{312.102}{294.675}{b}%
\TextAt{9}{320.424}{294.675}{\Character{61}}%
\TextAt{13}{332.85}{294.675}{c}%
\TextAt{17}{337.888}{299.611}{2}%
\TextAt{9}{345.277}{294.675}{\Character{43}}%
\TextAt{13}{357.038}{294.675}{d}%
\TextAt{17}{363.121}{299.611}{2}%
\TextAt{13}{367.853}{294.675}{\Character{58}}%
\ErrataAt{79.19}{265.674}{455.81}{If $c,d\in\mathbb Z$, then $a,b\in\mathbb Z$. Conversely, suppose $\alpha$ is a root of a monic quadratic in $\mathbb Z[x]$.}%
\ErrataAt{79.19}{248.628}{455.81}{If $d=0$, the rational-root theorem gives $c\in\mathbb Z$. If $d\ne0$, $\alpha$ and $\bar\alpha$ are distinct roots.}%
\ErrataAt{79.19}{231.582}{455.81}{Thus $-2c$ and $b=c^2+d^2$ are integers. The rational number $2d$ has integer square}%
\ErrataAt{79.19}{217}{455.81}{$4b-(2c)^2$, so it is integral. The sum $(2c)^2+(2d)^2=4b$ forces both integers even.}%
\ErrataAt{524.4}{217}{9.6}{$\square$}%
\TextAt{8}{97.159}{202.581}{The}%
\TextAt{8}{121.167}{202.581}{last}%
\TextAt{8}{143.298}{202.581}{proposition}%
\TextAt{8}{204.786}{202.581}{leads}%
\TextAt{8}{234.066}{202.581}{us}%
\TextAt{8}{249.036}{202.581}{to}%
\TextAt{8}{263.301}{202.581}{the}%
\TextAt{8}{283.412}{202.581}{general}%
\TextAt{8}{323.691}{202.581}{notion}%
\TextAt{8}{360.073}{202.581}{of}%
\TextAt{8}{373.357}{202.581}{an}%
\TextAt{12}{389.571}{202.581}{algebraic}%
\TextAt{12}{446.377}{202.581}{integer}%
\TextAt{8}{490.805}{202.581}{as}%
\TextAt{8}{505.129}{202.581}{being}%
\TextAt{8}{79.226}{185.535}{an}%
\TextAt{8}{95.442}{185.535}{element}%
\TextAt{8}{138.64}{185.535}{satisfying}%
\TextAt{8}{190.757}{185.535}{a}%
\TextAt{8}{200.469}{185.535}{monic}%
\TextAt{8}{234.892}{185.535}{polynomial}%
\TextAt{8}{295.341}{185.535}{equation}%
\TextAt{8}{343.106}{185.535}{with}%
\TextAt{8}{369.723}{185.535}{rational}%
\TextAt{8}{413.268}{185.535}{integer}%
\TextAt{8}{451.923}{185.535}{coefficients.}%
\TextAt{8}{515.697}{185.535}{For}%
\TextAt{8}{79.226}{168.489}{the}%
\TextAt{8}{98.832}{168.489}{domain}%
\TextAt{8}{139.909}{168.489}{of}%
\TextAt{8}{152.688}{168.489}{the}%
\TextAt{8}{172.294}{168.489}{gaussian}%
\TextAt{8}{218.703}{168.489}{integers}%
\TextAt{8}{261.466}{168.489}{we}%
\TextAt{8}{278.143}{168.489}{have}%
\TextAt{8}{304.587}{168.489}{obtained}%
\TextAt{8}{352.167}{168.489}{in}%
\TextAt{8}{365.27}{168.489}{this}%
\TextAt{8}{387.541}{168.489}{section}%
\TextAt{8}{426.072}{168.489}{a}%
\TextAt{8}{435.277}{168.489}{complete}%
\TextAt{8}{484.149}{168.489}{answer}%
\TextAt{8}{522.357}{168.489}{to}%
\TextAt{8}{79.226}{151.443}{the}%
\TextAt{8}{99.716}{151.443}{question}%
\TextAt{8}{146.614}{151.443}{of}%
\TextAt{8}{160.278}{151.443}{the}%
\TextAt{8}{180.767}{151.443}{units,}%
\TextAt{8}{213.762}{151.443}{the}%
\TextAt{8}{234.251}{151.443}{question}%
\TextAt{8}{281.15}{151.443}{of}%
\TextAt{8}{294.813}{151.443}{prime}%
\TextAt{8}{328.31}{151.443}{elements,}%
\TextAt{8}{379.835}{151.443}{and}%
\TextAt{8}{402.931}{151.443}{to}%
\TextAt{8}{417.574}{151.443}{the}%
\TextAt{8}{438.064}{151.443}{question}%
\TextAt{8}{484.962}{151.443}{of}%
\TextAt{8}{498.626}{151.443}{unique}%
\TextAt{8}{79.226}{134.397}{factorization.}%
\TextAt{8}{97.159}{114.362}{These}%
\TextAt{8}{131.644}{114.362}{questions}%
\TextAt{8}{183.44}{114.362}{indicate}%
\TextAt{8}{228.279}{114.362}{already}%
\TextAt{8}{270.2}{114.362}{the}%
\TextAt{8}{290.973}{114.362}{fundamental}%
\TextAt{8}{359.231}{114.362}{problems}%
\TextAt{8}{409.988}{114.362}{in}%
\TextAt{8}{424.257}{114.362}{the}%
\TextAt{8}{445.029}{114.362}{general}%
\TextAt{8}{485.97}{114.362}{theory}%
\TextAt{8}{523.337}{114.362}{of}%
\TextAt{8}{79.226}{97.316}{algebraic}%
\TextAt{8}{127.818}{97.316}{integers.}%
\TextAt{8}{175.402}{97.316}{But}%
\TextAt{8}{197.801}{97.316}{the}%
\TextAt{8}{217.115}{97.316}{answers}%
\TextAt{8}{259.645}{97.316}{we}%
\TextAt{8}{276.03}{97.316}{found}%
\TextAt{8}{308.029}{97.316}{in}%
\TextAt{8}{320.839}{97.316}{the}%
\TextAt{8}{340.153}{97.316}{special}%
\TextAt{8}{377.412}{97.316}{case}%
\TextAt{18}{401.341}{97.316}{Z}%
\TextAt{9}{409.311}{97.316}{\Character{91}}%
\TextAt{13}{412.563}{97.316}{i}%
\TextAt{9}{416.556}{97.316}{\Character{93}}%
\TextAt{8}{422.863}{97.316}{are}%
\TextAt{8}{441.531}{97.316}{not}%
\TextAt{8}{461.503}{97.316}{typical.}%
\TextAt{8}{504.149}{97.316}{Novel}%
\TextAt{8}{79.226}{80.27}{features}%
\TextAt{8}{123.185}{80.27}{will}%
\TextAt{8}{145.29}{80.27}{present}%
\TextAt{8}{185.998}{80.27}{themselves}%
\TextAt{8}{244.65}{80.27}{instead.}%
\TextAt{8}{303.069}{43.033}{4}%
\Destination{78.23}{809.815}{nav.page.4}%
\Destination{79.36}{431.655}{nav.print.5}%
\Destination{79.36}{447.655}{nav.number.1.1.5}%
\Hotlink{277.617}{741.855}{25.078}{11.047}{nav.number.1.1.1}%
\Destination{79.226}{282.674}{errata.45}%
}

% Page 17 (body page 5)
\Page{%
\TextAt{2}{79.226}{796.943}{I.}%
\TextAt{2}{89.83}{796.943}{Algebraic}%
\TextAt{2}{138.648}{796.943}{Integers}%
\Vector{q 0 G 0 g 0.299 w 7.226 18.729 m 460.769 18.729 l S Q}%
\TextAt{12}{79.226}{761.124}{Exercise}%
\TextAt{12}{132.427}{761.124}{I.1.1.}%
\TextAt{13}{166.021}{761.124}{\Character{11}}%
\NativeMathAt{176.863}{761.124}{11.955}{\in}%
\TextAt{18}{188.154}{761.124}{Z}%
\TextAt{9}{196.124}{761.124}{\Character{91}}%
\TextAt{13}{199.376}{761.124}{i}%
\TextAt{9}{203.369}{761.124}{\Character{93}}%
\TextAt{8}{210.518}{761.124}{is}%
\TextAt{8}{222.282}{761.124}{a}%
\TextAt{8}{232.038}{761.124}{unit}%
\TextAt{8}{256.749}{761.124}{if}%
\TextAt{8}{267.473}{761.124}{and}%
\TextAt{8}{290.236}{761.124}{only}%
\TextAt{8}{315.927}{761.124}{if}%
\TextAt{13}{326.651}{761.124}{N}%
\TextAt{9}{337.274}{761.124}{\Character{40}}%
\TextAt{13}{341.826}{761.124}{\Character{11}}%
\TextAt{9}{349.348}{761.124}{\Character{41}}%
\TextAt{9}{357.203}{761.124}{\Character{61}}%
\TextAt{9}{369.624}{761.124}{1}%
\TextAt{8}{375.499}{761.124}{.}%
\TextAt{12}{79.226}{729.134}{Exercise}%
\TextAt{12}{133.5}{729.134}{I.1.2.}%
\TextAt{8}{168.027}{729.134}{Show}%
\TextAt{8}{199.852}{729.134}{that,}%
\TextAt{8}{229.639}{729.134}{in}%
\TextAt{8}{244.225}{729.134}{the}%
\TextAt{8}{265.315}{729.134}{ring}%
\TextAt{18}{290.313}{729.134}{Z}%
\TextAt{9}{298.283}{729.134}{\Character{91}}%
\TextAt{13}{301.535}{729.134}{i}%
\TextAt{9}{305.529}{729.134}{\Character{93}}%
\TextAt{8}{308.78}{729.134}{,}%
\TextAt{8}{317.096}{729.134}{the}%
\TextAt{8}{338.185}{729.134}{relation}%
\TextAt{13}{382.049}{729.134}{\Character{11}\Character{12}}%
\TextAt{9}{401.753}{729.134}{\Character{61}}%
\TextAt{13}{415.769}{729.134}{\Character{34}\Character{13}}%
\TextAt{21}{427.969}{733.472}{n}%
\TextAt{8}{433.605}{729.134}{,}%
\TextAt{8}{441.921}{729.134}{for}%
\TextAt{13}{460.739}{729.134}{\Character{11}\Character{59}}%
\TextAt{13}{473.483}{729.134}{\Character{12}}%
\TextAt{8}{485.606}{729.134}{relatively}%
\TextAt{8}{79.226}{712.088}{prime}%
\TextAt{8}{112.39}{712.088}{numbers}%
\TextAt{8}{159.601}{712.088}{and}%
\TextAt{13}{182.363}{712.088}{\Character{34}}%
\TextAt{8}{191.739}{712.088}{a}%
\TextAt{8}{201.494}{712.088}{unit,}%
\TextAt{8}{229.457}{712.088}{implies}%
\TextAt{13}{269.184}{712.088}{\Character{11}}%
\TextAt{9}{280.027}{712.088}{\Character{61}}%
\TextAt{13}{292.452}{712.088}{\Character{34}}%
\TextAt{3}{297.93}{716.426}{\Character{48}}%
\TextAt{13}{300.726}{712.088}{\Character{24}}%
\TextAt{21}{306.416}{716.426}{n}%
\TextAt{8}{315.95}{712.088}{and}%
\TextAt{13}{338.712}{712.088}{\Character{12}}%
\TextAt{9}{349.304}{712.088}{\Character{61}}%
\TextAt{13}{361.73}{712.088}{\Character{34}}%
\TextAt{3}{367.208}{716.426}{\Character{48}\Character{48}}%
\TextAt{13}{372.3}{712.088}{\Character{17}}%
\TextAt{21}{378.504}{716.426}{n}%
\TextAt{8}{384.141}{712.088}{,}%
\TextAt{8}{391.29}{712.088}{with}%
\TextAt{13}{417.95}{712.088}{\Character{34}}%
\TextAt{3}{423.428}{716.426}{\Character{48}}%
\TextAt{13}{426.223}{712.088}{\Character{59}}%
\TextAt{13}{431.447}{712.088}{\Character{34}}%
\TextAt{3}{436.945}{716.426}{\Character{48}\Character{48}}%
\TextAt{8}{445.935}{712.088}{units.}%
\TextAt{12}{79.226}{680.098}{Exercise}%
\TextAt{12}{132.427}{680.098}{I.1.3.}%
\TextAt{8}{166.021}{680.098}{Show}%
\TextAt{8}{196.913}{680.098}{that}%
\TextAt{8}{222.282}{680.098}{the}%
\TextAt{8}{242.438}{680.098}{integer}%
\TextAt{8}{281.137}{680.098}{solutions}%
\TextAt{8}{330.046}{680.098}{of}%
\TextAt{8}{343.376}{680.098}{the}%
\TextAt{8}{363.532}{680.098}{equation}%
\TextAt{13}{274.438}{651.097}{x}%
\TextAt{17}{281.09}{656.033}{2}%
\TextAt{9}{288.479}{651.097}{\Character{43}}%
\TextAt{13}{300.24}{651.097}{y}%
\TextAt{17}{306.377}{656.033}{2}%
\TextAt{9}{314.43}{651.097}{\Character{61}}%
\TextAt{13}{326.855}{651.097}{z}%
\TextAt{17}{332.826}{656.033}{2}%
\TextAt{8}{79.226}{622.096}{such}%
\TextAt{8}{105.326}{622.096}{that}%
\TextAt{13}{130.399}{622.096}{x\Character{59}}%
\TextAt{13}{142.271}{622.096}{y}%
\TextAt{13}{148.404}{622.096}{\Character{59}}%
\TextAt{13}{153.64}{622.096}{z}%
\TextAt{13}{162.917}{622.096}{\Character{62}}%
\TextAt{9}{175.393}{622.096}{0}%
\TextAt{8}{184.846}{622.096}{and}%
\TextAt{9}{207.312}{622.096}{\Character{40}}%
\TextAt{13}{211.865}{622.096}{x\Character{59}}%
\TextAt{13}{223.737}{622.096}{y}%
\TextAt{13}{229.87}{622.096}{\Character{59}}%
\TextAt{13}{235.106}{622.096}{z}%
\TextAt{9}{241.112}{622.096}{\Character{41}}%
\TextAt{9}{248.967}{622.096}{\Character{61}}%
\TextAt{9}{261.388}{622.096}{1}%
\TextAt{8}{270.865}{622.096}{(“pythagorean}%
\TextAt{8}{347.655}{622.096}{triples”)}%
\TextAt{8}{393.266}{622.096}{are}%
\TextAt{8}{412.48}{622.096}{all}%
\TextAt{8}{428.443}{622.096}{given,}%
\TextAt{8}{462.026}{622.096}{up}%
\TextAt{8}{478.634}{622.096}{to}%
\TextAt{8}{492.648}{622.096}{possible}%
\TextAt{8}{79.226}{605.05}{permutation}%
\TextAt{8}{146.223}{605.05}{of}%
\TextAt{13}{159.553}{605.05}{x}%
\TextAt{8}{170.102}{605.05}{and}%
\TextAt{13}{192.865}{605.05}{y}%
\TextAt{8}{199.002}{605.05}{,}%
\TextAt{8}{206.151}{605.05}{by}%
\TextAt{8}{222.41}{605.05}{the}%
\TextAt{8}{242.566}{605.05}{formulae}%
\TextAt{13}{196.186}{576.048}{x}%
\TextAt{9}{206.159}{576.048}{\Character{61}}%
\TextAt{13}{218.585}{576.048}{u}%
\TextAt{17}{225.247}{580.985}{2}%
\TextAt{16}{232.636}{576.048}{\Character{0}}%
\TextAt{13}{244.591}{576.048}{v}%
\TextAt{17}{250.68}{580.985}{2}%
\TextAt{13}{255.412}{576.048}{\Character{59}}%
\TextAt{13}{284.547}{576.048}{y}%
\TextAt{9}{294.024}{576.048}{\Character{61}}%
\TextAt{9}{306.434}{576.048}{2}%
\TextAt{13}{312.302}{576.048}{uv}%
\TextAt{13}{325.034}{576.048}{\Character{59}}%
\TextAt{13}{354.181}{576.048}{z}%
\TextAt{9}{363.499}{576.048}{\Character{61}}%
\TextAt{13}{375.924}{576.048}{u}%
\TextAt{17}{382.587}{580.985}{2}%
\TextAt{9}{389.976}{576.048}{\Character{43}}%
\TextAt{13}{401.737}{576.048}{v}%
\TextAt{17}{407.825}{580.985}{2}%
\TextAt{13}{412.558}{576.048}{\Character{59}}%
\TextAt{8}{79.226}{547.047}{where}%
\TextAt{13}{113.035}{547.047}{u\Character{59}}%
\TextAt{13}{124.919}{547.047}{v}%
\NativeMathAt{134.351}{547.047}{11.955}{\in}%
\TextAt{18}{145.642}{547.047}{Z}%
\TextAt{8}{153.612}{547.047}{,}%
\TextAt{13}{160.761}{547.047}{u}%
\TextAt{13}{170.732}{547.047}{\Character{62}}%
\TextAt{13}{183.153}{547.047}{v}%
\TextAt{13}{192.562}{547.047}{\Character{62}}%
\TextAt{9}{205.005}{547.047}{0}%
\TextAt{8}{210.858}{547.047}{,}%
\TextAt{9}{218.007}{547.047}{\Character{40}}%
\TextAt{13}{222.559}{547.047}{u\Character{59}}%
\TextAt{13}{234.443}{547.047}{v}%
\TextAt{9}{240.554}{547.047}{\Character{41}}%
\TextAt{9}{248.409}{547.047}{\Character{61}}%
\TextAt{9}{260.83}{547.047}{1}%
\TextAt{8}{266.706}{547.047}{,}%
\TextAt{13}{273.855}{547.047}{u\Character{59}}%
\TextAt{13}{285.738}{547.047}{v}%
\TextAt{8}{295.747}{547.047}{not}%
\TextAt{8}{316.561}{547.047}{both}%
\TextAt{8}{344.201}{547.047}{odd.}%
\TextAt{12}{97.159}{527.013}{Hint:}%
\TextAt{8}{134.004}{527.013}{Use}%
\TextAt{8}{157.144}{527.013}{exercise}%
\TextAt{8}{201.099}{527.013}{2}%
\TextAt{8}{211.508}{527.013}{to}%
\TextAt{8}{226.471}{527.013}{show}%
\TextAt{8}{256.127}{527.013}{that}%
\TextAt{8}{282.149}{527.013}{necessarily}%
\TextAt{13}{341.132}{527.013}{x}%
\TextAt{9}{350.886}{527.013}{\Character{43}}%
\TextAt{13}{363.092}{527.013}{iy}%
\TextAt{9}{377.656}{527.013}{\Character{61}}%
\TextAt{13}{391.194}{527.013}{\Character{34}\Character{11}}%
\TextAt{17}{404.194}{531.351}{2}%
\TextAt{8}{413.477}{527.013}{with}%
\TextAt{8}{440.79}{527.013}{a}%
\TextAt{8}{451.198}{527.013}{unit}%
\TextAt{13}{476.563}{527.013}{\Character{34}}%
\TextAt{8}{486.591}{527.013}{and}%
\TextAt{8}{510.007}{527.013}{with}%
\TextAt{13}{79.226}{509.967}{\Character{11}}%
\TextAt{9}{90.069}{509.967}{\Character{61}}%
\TextAt{13}{102.494}{509.967}{u}%
\TextAt{9}{111.813}{509.967}{\Character{43}}%
\TextAt{13}{123.575}{509.967}{iv}%
\NativeMathAt{136.977}{509.967}{11.955}{\in}%
\TextAt{18}{148.268}{509.967}{Z}%
\TextAt{9}{156.238}{509.967}{\Character{91}}%
\TextAt{13}{159.49}{509.967}{i}%
\TextAt{9}{163.483}{509.967}{\Character{93}}%
\TextAt{8}{166.735}{509.967}{.}%
\TextAt{12}{79.226}{477.977}{Exercise}%
\TextAt{12}{132.427}{477.977}{I.1.4.}%
\TextAt{8}{166.021}{477.977}{Show}%
\TextAt{8}{196.913}{477.977}{that}%
\TextAt{8}{222.282}{477.977}{the}%
\TextAt{8}{242.438}{477.977}{ring}%
\TextAt{18}{266.504}{477.977}{Z}%
\TextAt{9}{274.474}{477.977}{\Character{91}}%
\TextAt{13}{277.726}{477.977}{i}%
\TextAt{9}{281.719}{477.977}{\Character{93}}%
\TextAt{8}{288.868}{477.977}{cannot}%
\TextAt{8}{327.244}{477.977}{be}%
\TextAt{8}{343.169}{477.977}{ordered.}%
\TextAt{12}{79.226}{445.987}{Exercise}%
\TextAt{12}{133.984}{445.987}{I.1.5.}%
\TextAt{8}{168.931}{445.987}{Show}%
\TextAt{8}{201.177}{445.987}{that}%
\TextAt{8}{227.899}{445.987}{the}%
\TextAt{8}{249.409}{445.987}{only}%
\TextAt{8}{276.454}{445.987}{units}%
\TextAt{8}{307.134}{445.987}{of}%
\TextAt{8}{321.817}{445.987}{the}%
\TextAt{8}{343.327}{445.987}{ring}%
\TextAt{18}{368.746}{445.987}{Z}%
\TextAt{9}{376.716}{445.987}{\Character{91}}%
\TextAt{16}{379.968}{455.677}{\Character{112}}%
\Vector{q 0 G 0 g 0.478 w 317.931 -318.084 m 333.312 -318.084 l S Q}%
\TextAt{16}{389.931}{445.987}{\Character{0}}%
\TextAt{13}{399.229}{445.987}{d}%
\TextAt{9}{405.312}{445.987}{\Character{93}}%
\TextAt{9}{414.171}{445.987}{\Character{61}}%
\TextAt{18}{428.923}{445.987}{Z}%
\TextAt{9}{440.472}{445.987}{\Character{43}}%
\TextAt{18}{453.156}{445.987}{Z}%
\TextAt{16}{461.127}{455.677}{\Character{112}}%
\Vector{q 0 G 0 g 0.478 w 399.089 -318.084 m 414.47 -318.084 l S Q}%
\TextAt{16}{471.089}{445.987}{\Character{0}}%
\TextAt{13}{480.388}{445.987}{d}%
\TextAt{8}{486.47}{445.987}{,}%
\TextAt{8}{495.311}{445.987}{for}%
\TextAt{8}{514.55}{445.987}{any}%
\TextAt{8}{79.226}{428.941}{rational}%
\TextAt{8}{122.815}{428.941}{integer}%
\TextAt{13}{161.514}{428.941}{d}%
\TextAt{13}{170.899}{428.941}{\Character{62}}%
\TextAt{9}{183.343}{428.941}{1}%
\TextAt{8}{189.196}{428.941}{,}%
\TextAt{8}{196.345}{428.941}{are}%
\NativeMathAt{215.856}{428.941}{11.955}{\pm}%
\TextAt{9}{225.154}{428.941}{1}%
\TextAt{8}{231.007}{428.941}{.}%
\TextAt{12}{79.226}{396.951}{Exercise}%
\TextAt{12}{131.869}{396.951}{I.1.6.}%
\TextAt{8}{164.979}{396.951}{Show}%
\TextAt{8}{195.387}{396.951}{that}%
\TextAt{8}{220.271}{396.951}{the}%
\TextAt{8}{239.944}{396.951}{ring}%
\TextAt{18}{263.525}{396.951}{Z}%
\TextAt{9}{271.495}{396.951}{\Character{91}}%
\TextAt{16}{274.747}{407.139}{\Character{112}}%
\Vector{q 0 G 0 g 0.478 w 212.709 -366.622 m 218.792 -366.622 l S Q}%
\TextAt{13}{284.709}{396.951}{d}%
\TextAt{9}{290.792}{396.951}{\Character{93}}%
\TextAt{9}{297.343}{396.951}{\Character{61}}%
\TextAt{18}{309.79}{396.951}{Z}%
\TextAt{9}{319.426}{396.951}{\Character{43}}%
\TextAt{18}{330.197}{396.951}{Z}%
\TextAt{16}{338.167}{407.139}{\Character{112}}%
\Vector{q 0 G 0 g 0.478 w 276.13 -366.622 m 282.213 -366.622 l S Q}%
\TextAt{13}{348.13}{396.951}{d}%
\TextAt{8}{354.213}{396.951}{,}%
\TextAt{8}{360.974}{396.951}{for}%
\TextAt{8}{378.375}{396.951}{any}%
\TextAt{8}{400.008}{396.951}{squarefree}%
\TextAt{8}{454.864}{396.951}{rational}%
\TextAt{8}{497.968}{396.951}{integer}%
\TextAt{13}{79.226}{379.905}{d}%
\TextAt{13}{88.611}{379.905}{\Character{62}}%
\TextAt{9}{101.055}{379.905}{1}%
\TextAt{8}{106.908}{379.905}{,}%
\TextAt{8}{114.057}{379.905}{has}%
\TextAt{8}{134.931}{379.905}{infinitely}%
\TextAt{8}{184.031}{379.905}{many}%
\TextAt{8}{215.903}{379.905}{units.}%
\TextAt{12}{79.226}{347.915}{Exercise}%
\TextAt{12}{133.158}{347.915}{I.1.7.}%
\TextAt{8}{167.389}{347.915}{Show}%
\TextAt{8}{198.917}{347.915}{that}%
\TextAt{8}{224.922}{347.915}{the}%
\TextAt{8}{245.714}{347.915}{ring}%
\TextAt{18}{270.416}{347.915}{Z}%
\TextAt{9}{278.386}{347.915}{\Character{91}}%
\TextAt{16}{281.638}{357.804}{\Character{112}}%
\Vector{q 0 G 0 g 0.478 w 219.601 -415.956 m 225.454 -415.956 l S Q}%
\TextAt{9}{291.601}{347.915}{2\Character{93}}%
\TextAt{9}{305.087}{347.915}{\Character{61}}%
\TextAt{18}{318.62}{347.915}{Z}%
\TextAt{9}{329.68}{347.915}{\Character{43}}%
\TextAt{18}{341.875}{347.915}{Z}%
\TextAt{16}{349.845}{357.804}{\Character{112}}%
\Vector{q 0 G 0 g 0.478 w 287.808 -415.956 m 293.661 -415.956 l S Q}%
\TextAt{9}{359.808}{347.915}{2}%
\TextAt{8}{370.195}{347.915}{is}%
\TextAt{8}{382.595}{347.915}{euclidean.}%
\TextAt{8}{440.425}{347.915}{Show}%
\TextAt{8}{471.954}{347.915}{furthermore}%
\TextAt{8}{79.226}{330.869}{that}%
\TextAt{8}{104.595}{330.869}{its}%
\TextAt{8}{120.914}{330.869}{units}%
\TextAt{8}{150.24}{330.869}{are}%
\TextAt{8}{169.751}{330.869}{given}%
\TextAt{8}{200.32}{330.869}{by}%
\NativeMathAt{216.579}{330.869}{11.955}{\pm}%
\TextAt{9}{225.878}{330.869}{\Character{40}1}%
\TextAt{9}{238.921}{330.869}{\Character{43}}%
\TextAt{16}{250.701}{340.759}{\Character{112}}%
\Vector{q 0 G 0 g 0.478 w 188.664 -433.002 m 194.517 -433.002 l S Q}%
\TextAt{9}{260.664}{330.869}{2\Character{41}}%
\TextAt{21}{271.069}{335.207}{n}%
\TextAt{8}{276.705}{330.869}{,}%
\TextAt{13}{283.854}{330.869}{n}%
\NativeMathAt{294.163}{330.869}{11.955}{\in}%
\TextAt{18}{305.454}{330.869}{Z}%
\TextAt{8}{313.424}{330.869}{,}%
\TextAt{8}{320.573}{330.869}{and}%
\TextAt{8}{343.336}{330.869}{determine}%
\TextAt{8}{397.959}{330.869}{its}%
\TextAt{8}{414.278}{330.869}{prime}%
\TextAt{8}{447.441}{330.869}{elements.}%
\TextAt{15}{79.226}{278.005}{§}%
\TextAt{15}{93.785}{278.005}{2}%
\TextAt{15}{120.693}{278.005}{Integrality}%
\TextAt{8}{79.226}{246.119}{An}%
\TextAt{12}{100.079}{246.119}{algebraic}%
\TextAt{12}{158.867}{246.119}{number}%
\TextAt{12}{210.171}{246.119}{field}%
\TextAt{8}{240.564}{246.119}{is}%
\TextAt{8}{254.005}{246.119}{a}%
\TextAt{8}{265.438}{246.119}{finite}%
\TextAt{8}{297.027}{246.119}{field}%
\TextAt{8}{324.061}{246.119}{extension}%
\TextAt{13}{377.503}{246.119}{K}%
\TextAt{8}{393.895}{246.119}{of}%
\TextAt{18}{408.903}{246.119}{Q}%
\TextAt{8}{418.201}{246.119}{.}%
\TextAt{8}{431.68}{246.119}{The}%
\TextAt{8}{457.411}{246.119}{elements}%
\TextAt{8}{506.945}{246.119}{of}%
\TextAt{13}{521.952}{246.119}{K}%
\TextAt{8}{79.226}{229.073}{are}%
\TextAt{8}{98.241}{229.073}{called}%
\TextAt{12}{130.91}{229.073}{algebraic}%
\TextAt{12}{187.198}{229.073}{numbers}%
\TextAt{8}{237.398}{229.073}{.}%
\TextAt{8}{245.681}{229.073}{An}%
\TextAt{8}{264.362}{229.073}{algebraic}%
\TextAt{8}{313.301}{229.073}{number}%
\TextAt{8}{355.402}{229.073}{is}%
\TextAt{8}{366.67}{229.073}{called}%
\TextAt{12}{399.338}{229.073}{integral}%
\TextAt{8}{444.063}{229.073}{,}%
\TextAt{8}{450.815}{229.073}{or}%
\TextAt{8}{464.63}{229.073}{an}%
\TextAt{12}{480.394}{229.073}{algebraic}%
\TextAt{12}{79.226}{212.027}{integer}%
\TextAt{8}{119.802}{212.027}{,}%
\TextAt{8}{126.056}{212.027}{if}%
\TextAt{8}{135.662}{212.027}{it}%
\TextAt{8}{146.247}{212.027}{is}%
\TextAt{8}{156.893}{212.027}{a}%
\TextAt{8}{165.529}{212.027}{zero}%
\TextAt{8}{189.122}{212.027}{of}%
\TextAt{8}{201.334}{212.027}{a}%
\TextAt{8}{209.971}{212.027}{monic}%
\TextAt{8}{243.319}{212.027}{polynomial}%
\TextAt{13}{302.693}{212.027}{f}%
\TextAt{9}{309.74}{212.027}{\Character{40}}%
\TextAt{13}{314.292}{212.027}{x}%
\TextAt{9}{320.944}{212.027}{\Character{41}}%
\NativeMathAt{328.817}{212.027}{11.955}{\in}%
\TextAt{18}{340.108}{212.027}{Z}%
\TextAt{9}{348.078}{212.027}{\Character{91}}%
\TextAt{13}{351.33}{212.027}{x}%
\TextAt{9}{357.982}{212.027}{\Character{93}}%
\TextAt{8}{361.234}{212.027}{.}%
\TextAt{8}{369.309}{212.027}{This}%
\TextAt{8}{394.911}{212.027}{notion}%
\TextAt{8}{430.22}{212.027}{of}%
\TextAt{8}{442.431}{212.027}{integrality}%
\TextAt{8}{497.585}{212.027}{applies}%
\TextAt{8}{79.226}{194.982}{not}%
\TextAt{8}{99.615}{194.982}{only}%
\TextAt{8}{124.881}{194.982}{to}%
\TextAt{8}{138.766}{194.982}{algebraic}%
\TextAt{8}{187.775}{194.982}{numbers,}%
\TextAt{8}{237.898}{194.982}{but}%
\TextAt{8}{258.932}{194.982}{occurs}%
\TextAt{8}{294.659}{194.982}{in}%
\TextAt{8}{307.886}{194.982}{many}%
\TextAt{8}{339.334}{194.982}{different}%
\TextAt{8}{385.079}{194.982}{contexts}%
\TextAt{8}{430.896}{194.982}{and}%
\TextAt{8}{453.234}{194.982}{therefore}%
\TextAt{8}{501.908}{194.982}{has}%
\TextAt{8}{522.357}{194.982}{to}%
\TextAt{8}{79.226}{177.936}{be}%
\TextAt{8}{95.15}{177.936}{treated}%
\TextAt{8}{135.475}{177.936}{in}%
\TextAt{8}{149.128}{177.936}{full}%
\TextAt{8}{169.607}{177.936}{generality.}%
\TextAt{8}{97.159}{157.901}{In}%
\TextAt{8}{111.78}{157.901}{what}%
\TextAt{8}{141.046}{157.901}{follows,}%
\TextAt{8}{182.734}{157.901}{rings}%
\TextAt{8}{211.414}{157.901}{are}%
\TextAt{8}{230.925}{157.901}{always}%
\TextAt{8}{268.393}{157.901}{understood}%
\TextAt{8}{329.591}{157.901}{to}%
\TextAt{8}{343.902}{157.901}{be}%
\TextAt{8}{359.826}{157.901}{commutative}%
\TextAt{8}{429.752}{157.901}{rings}%
\TextAt{8}{458.432}{157.901}{with}%
\TextAt{9}{485.092}{157.901}{1}%
\TextAt{8}{490.945}{157.901}{.}%
\TextAt{12}{79.226}{131.888}{Definition}%
\TextAt{12}{143.555}{131.888}{(2.1).}%
\TextAt{8}{180.21}{131.888}{Let}%
\TextAt{13}{202.519}{131.888}{A}%
\NativeMathAt{216.897}{131.888}{11.955}{\subseteq}%
\TextAt{13}{231.799}{131.888}{B}%
\TextAt{8}{246.532}{131.888}{be}%
\TextAt{8}{263.795}{131.888}{an}%
\TextAt{8}{281.394}{131.888}{extension}%
\TextAt{8}{334.499}{131.888}{of}%
\TextAt{8}{349.168}{131.888}{rings.}%
\TextAt{8}{386.417}{131.888}{An}%
\TextAt{8}{406.933}{131.888}{element}%
\TextAt{13}{451.514}{131.888}{b}%
\NativeMathAt{462.094}{131.888}{11.955}{\in}%
\TextAt{13}{475.667}{131.888}{B}%
\TextAt{8}{490.4}{131.888}{is}%
\TextAt{8}{503.503}{131.888}{called}%
\TextAt{12}{79.226}{114.842}{integral}%
\TextAt{12}{128.434}{114.842}{over}%
\TextAt{13}{157.64}{114.842}{A}%
\TextAt{8}{166.415}{114.842}{,}%
\TextAt{8}{173.565}{114.842}{if}%
\TextAt{8}{184.288}{114.842}{it}%
\TextAt{8}{195.993}{114.842}{satisfies}%
\TextAt{8}{239.103}{114.842}{a}%
\TextAt{8}{248.858}{114.842}{monic}%
\TextAt{8}{283.325}{114.842}{equation}%
\ErrataAt{181}{85.855}{250}{\makebox[250bp][c]{$x^n+a_1x^{n-1}+\cdots+a_n=0,\qquad n\ge1$,}}%
\TextAt{8}{303.069}{43.033}{5}%
\Destination{78.23}{809.815}{nav.page.5}%
\Destination{79.23}{303.865}{nav.section.1.2}%
\Destination{79.226}{775.155}{nav.ex.1.1.1}%
\Destination{79.226}{743.155}{nav.ex.1.1.2}%
\Destination{79.226}{694.055}{nav.ex.1.1.3}%
\Destination{79.226}{491.955}{nav.ex.1.1.4}%
\Destination{79.226}{459.955}{nav.ex.1.1.5}%
\Destination{79.226}{410.955}{nav.ex.1.1.6}%
\Destination{79.226}{361.955}{nav.ex.1.1.7}%
\Destination{79.36}{145.855}{nav.number.1.2.1}%
\Destination{79.23}{303.865}{outline.7}%
\Hotlink{200.599}{524.855}{6.858}{11.047}{nav.ex.1.1.2}%
\Hotlink{277.886}{345.755}{23.301}{11.047}{nav.cite.N002}%
\Destination{79.226}{99.855}{errata.44}%
}

% Page 18 (body page 6)
\Page{%
\TextAt{2}{79.226}{796.943}{I.}%
\TextAt{2}{89.83}{796.943}{Algebraic}%
\TextAt{2}{138.648}{796.943}{Integers}%
\Vector{q 0 G 0 g 0.299 w 7.226 18.729 m 460.769 18.729 l S Q}%
\TextAt{8}{79.226}{761.124}{with}%
\TextAt{8}{106.107}{761.124}{coefficients}%
\TextAt{13}{165.565}{761.124}{a}%
\TextAt{21}{171.71}{759.33}{i}%
\NativeMathAt{178.788}{761.124}{11.955}{\in}%
\TextAt{13}{190.455}{761.124}{A}%
\TextAt{8}{199.231}{761.124}{.}%
\TextAt{8}{208.341}{761.124}{The}%
\TextAt{8}{232.615}{761.124}{ring}%
\TextAt{13}{256.902}{761.124}{B}%
\TextAt{8}{270.516}{761.124}{is}%
\TextAt{8}{282.501}{761.124}{called}%
\TextAt{12}{315.885}{761.124}{integral}%
\TextAt{12}{365.346}{761.124}{over}%
\TextAt{13}{394.806}{761.124}{A}%
\TextAt{8}{407.7}{761.124}{if}%
\TextAt{8}{418.644}{761.124}{all}%
\TextAt{8}{435.124}{761.124}{elements}%
\TextAt{13}{483.201}{761.124}{b}%
\NativeMathAt{491.875}{761.124}{11.955}{\in}%
\TextAt{13}{503.542}{761.124}{B}%
\TextAt{8}{517.156}{761.124}{are}%
\TextAt{8}{79.226}{744.078}{integral}%
\TextAt{8}{121.834}{744.078}{over}%
\TextAt{13}{146.88}{744.078}{A}%
\TextAt{8}{155.656}{744.078}{.}%
\TextAt{8}{97.159}{718.656}{It}%
\TextAt{8}{110.16}{718.656}{is}%
\TextAt{8}{122.252}{718.656}{desirable,}%
\TextAt{8}{174.752}{718.656}{but}%
\TextAt{8}{196.54}{718.656}{strangely}%
\TextAt{8}{247.343}{718.656}{enough}%
\TextAt{8}{287.996}{718.656}{not}%
\TextAt{8}{309.139}{718.656}{immediately}%
\TextAt{8}{376.129}{718.656}{obvious,}%
\TextAt{8}{422.137}{718.656}{that}%
\TextAt{8}{447.835}{718.656}{the}%
\TextAt{8}{468.32}{718.656}{sum}%
\TextAt{8}{493.419}{718.656}{and}%
\TextAt{8}{516.51}{718.656}{the}%
\TextAt{8}{79.226}{701.61}{product}%
\TextAt{8}{124.567}{701.61}{of}%
\TextAt{8}{139.339}{701.61}{two}%
\TextAt{8}{162.898}{701.61}{elements}%
\TextAt{8}{212.197}{701.61}{which}%
\TextAt{8}{247.125}{701.61}{are}%
\TextAt{8}{268.078}{701.61}{integral}%
\TextAt{8}{312.128}{701.61}{over}%
\TextAt{13}{338.616}{701.61}{A}%
\TextAt{8}{352.73}{701.61}{are}%
\TextAt{8}{373.683}{701.61}{again}%
\TextAt{8}{406.352}{701.61}{integral.}%
\TextAt{8}{457.837}{701.61}{This}%
\TextAt{8}{485.998}{701.61}{will}%
\TextAt{8}{509.545}{701.61}{be}%
\TextAt{8}{526.911}{701.61}{a}%
\TextAt{8}{79.226}{684.564}{consequence}%
\TextAt{8}{145.29}{684.564}{of}%
\TextAt{8}{158.62}{684.564}{the}%
\TextAt{8}{178.777}{684.564}{following}%
\TextAt{8}{228.211}{684.564}{abstract}%
\TextAt{8}{273.809}{684.564}{reinterpretation}%
\TextAt{8}{358.045}{684.564}{of}%
\TextAt{8}{371.375}{684.564}{the}%
\TextAt{8}{391.531}{684.564}{notion}%
\TextAt{8}{427.958}{684.564}{of}%
\TextAt{8}{441.288}{684.564}{integrality.}%
\TextAt{12}{79.226}{659.143}{Proposition}%
\TextAt{12}{152.072}{659.143}{(2.2).}%
\TextAt{8}{187.505}{659.143}{Finitely}%
\TextAt{8}{231.354}{659.143}{many}%
\TextAt{8}{263.343}{659.143}{elements}%
\TextAt{13}{311.317}{659.143}{b}%
\TextAt{17}{316.294}{657.35}{1}%
\TextAt{13}{321.026}{659.143}{\Character{59}}%
\TextAt{13}{326.25}{659.143}{\Character{58}}%
\TextAt{13}{331.487}{659.143}{\Character{58}}%
\TextAt{13}{336.723}{659.143}{\Character{58}}%
\TextAt{13}{341.948}{659.143}{\Character{59}}%
\TextAt{13}{347.184}{659.143}{b}%
\TextAt{21}{352.224}{657.35}{n}%
\NativeMathAt{361.381}{659.143}{11.955}{\in}%
\TextAt{13}{372.871}{659.143}{B}%
\TextAt{8}{386.382}{659.143}{are}%
\TextAt{8}{406.009}{659.143}{all}%
\TextAt{8}{422.385}{659.143}{integral}%
\TextAt{8}{465.111}{659.143}{over}%
\TextAt{13}{490.274}{659.143}{A}%
\TextAt{8}{503.063}{659.143}{if}%
\TextAt{8}{513.904}{659.143}{and}%
\TextAt{8}{79.226}{642.097}{only}%
\TextAt{8}{104.918}{642.097}{if}%
\TextAt{8}{115.642}{642.097}{the}%
\TextAt{8}{135.798}{642.097}{ring}%
\TextAt{13}{159.864}{642.097}{A}%
\TextAt{9}{168.639}{642.097}{\Character{91}}%
\TextAt{13}{171.891}{642.097}{b}%
\TextAt{17}{176.868}{640.304}{1}%
\TextAt{13}{181.6}{642.097}{\Character{59}}%
\TextAt{13}{186.824}{642.097}{\Character{58}}%
\TextAt{13}{192.061}{642.097}{\Character{58}}%
\TextAt{13}{197.297}{642.097}{\Character{58}}%
\TextAt{13}{202.522}{642.097}{\Character{59}}%
\TextAt{13}{207.758}{642.097}{b}%
\TextAt{21}{212.798}{640.304}{n}%
\TextAt{9}{218.434}{642.097}{\Character{93}}%
\TextAt{8}{225.584}{642.097}{viewed}%
\TextAt{8}{263.948}{642.097}{as}%
\TextAt{8}{278.318}{642.097}{an}%
\TextAt{13}{294.577}{642.097}{A}%
\TextAt{8}{303.352}{642.097}{-module}%
\TextAt{8}{348.543}{642.097}{is}%
\TextAt{8}{360.307}{642.097}{finitely}%
\TextAt{8}{399.651}{642.097}{generated.}%
\TextAt{8}{97.159}{616.676}{To}%
\TextAt{8}{114.386}{616.676}{prove}%
\TextAt{8}{145.936}{616.676}{this}%
\TextAt{8}{168.758}{616.676}{we}%
\TextAt{8}{185.986}{616.676}{make}%
\TextAt{8}{216.555}{616.676}{use}%
\TextAt{8}{236.771}{616.676}{of}%
\TextAt{8}{250.101}{616.676}{the}%
\TextAt{8}{270.258}{616.676}{following}%
\TextAt{8}{319.692}{616.676}{result}%
\TextAt{8}{352.27}{616.676}{of}%
\TextAt{8}{365.6}{616.676}{linear}%
\TextAt{8}{398.118}{616.676}{algebra.}%
\TextAt{12}{79.226}{591.254}{Proposition}%
\TextAt{12}{152.931}{591.254}{(2.3)}%
\TextAt{12}{186.083}{591.254}{(Row-Column}%
\TextAt{12}{273.739}{591.254}{Expansion).}%
\TextAt{8}{348.139}{591.254}{Let}%
\TextAt{13}{369.972}{591.254}{A}%
\TextAt{9}{383.539}{591.254}{\Character{61}}%
\TextAt{9}{397.419}{591.254}{\Character{40}}%
\TextAt{13}{401.988}{591.254}{a}%
\TextAt{21}{408.133}{589.461}{ij}%
\TextAt{9}{415.399}{591.254}{\Character{41}}%
\TextAt{8}{424.712}{591.254}{be}%
\TextAt{8}{441.499}{591.254}{an}%
\TextAt{9}{458.622}{591.254}{\Character{40}}%
\TextAt{13}{463.174}{591.254}{r}%
\TextAt{16}{472.02}{591.254}{\Character{2}}%
\TextAt{13}{484.563}{591.254}{r}%
\TextAt{9}{490.164}{591.254}{\Character{41}}%
\TextAt{8}{494.716}{591.254}{-matrix}%
\TextAt{8}{79.226}{574.208}{with}%
\TextAt{8}{106.998}{574.208}{entries}%
\TextAt{8}{145.566}{574.208}{in}%
\TextAt{8}{160.331}{574.208}{an}%
\TextAt{8}{177.702}{574.208}{arbitrary}%
\TextAt{8}{228.584}{574.208}{ring,}%
\TextAt{8}{257.291}{574.208}{and}%
\TextAt{8}{281.166}{574.208}{let}%
\TextAt{13}{299.183}{574.208}{A}%
\TextAt{3}{307.958}{578.547}{\Character{3}}%
\TextAt{9}{317.907}{574.208}{\Character{61}}%
\TextAt{9}{332.217}{574.208}{\Character{40}}%
\TextAt{13}{336.78}{574.208}{a}%
\TextAt{3}{342.925}{578.547}{\Character{3}}%
\TextAt{21}{342.925}{571.253}{ij}%
\TextAt{9}{350.19}{574.208}{\Character{41}}%
\TextAt{8}{359.752}{574.208}{be}%
\TextAt{8}{376.788}{574.208}{the}%
\TextAt{8}{398.057}{574.208}{adjoint}%
\TextAt{8}{438.848}{574.208}{matrix,}%
\TextAt{8}{481.543}{574.208}{i.e.,}%
\TextAt{13}{505.039}{574.208}{a}%
\TextAt{3}{511.183}{578.547}{\Character{3}}%
\TextAt{21}{511.183}{571.253}{ij}%
\TextAt{9}{523.665}{574.208}{\Character{61}}%
\TextAt{9}{79.226}{557.162}{\Character{40}}%
\TextAt{16}{83.778}{557.162}{\Character{0}}%
\TextAt{9}{93.077}{557.162}{1\Character{41}}%
\TextAt{21}{103.482}{561.501}{i}%
\TextAt{17}{106.365}{561.501}{\Character{43}}%
\TextAt{21}{112.952}{561.501}{j}%
\TextAt{8}{119.327}{557.162}{det}%
\TextAt{9}{135.586}{557.162}{\Character{40}}%
\TextAt{13}{140.138}{557.162}{A}%
\TextAt{21}{148.913}{555.369}{ij}%
\TextAt{9}{156.179}{557.162}{\Character{41}}%
\TextAt{8}{160.731}{557.162}{,}%
\TextAt{8}{169.028}{557.162}{where}%
\TextAt{8}{203.755}{557.162}{the}%
\TextAt{8}{224.83}{557.162}{matrix}%
\TextAt{13}{263.802}{557.162}{A}%
\TextAt{21}{272.577}{555.369}{ij}%
\TextAt{8}{284.658}{557.162}{is}%
\TextAt{8}{297.34}{557.162}{obtained}%
\TextAt{8}{346.39}{557.162}{from}%
\TextAt{13}{374.948}{557.162}{A}%
\TextAt{8}{388.539}{557.162}{by}%
\TextAt{8}{405.717}{557.162}{deleting}%
\TextAt{8}{450.857}{557.162}{the}%
\TextAt{13}{471.932}{557.162}{i}%
\TextAt{8}{475.925}{557.162}{-th}%
\TextAt{8}{495.697}{557.162}{column}%
\TextAt{8}{79.226}{540.116}{and}%
\TextAt{8}{101.989}{540.116}{the}%
\TextAt{13}{122.145}{540.116}{j}%
\TextAt{8}{127.661}{540.116}{-th}%
\TextAt{8}{146.514}{540.116}{row.}%
\TextAt{8}{173.505}{540.116}{Then}%
\TextAt{8}{204.062}{540.116}{one}%
\TextAt{8}{225.522}{540.116}{has}%
\TextAt{13}{244.59}{513.183}{AA}%
\TextAt{3}{262.141}{518.12}{\Character{3}}%
\TextAt{9}{270.194}{513.183}{\Character{61}}%
\TextAt{13}{282.619}{513.183}{A}%
\TextAt{3}{291.395}{518.12}{\Character{3}}%
\TextAt{13}{296.127}{513.183}{A}%
\TextAt{9}{308.223}{513.183}{\Character{61}}%
\TextAt{8}{320.649}{513.183}{det}%
\TextAt{9}{336.908}{513.183}{\Character{40}}%
\TextAt{13}{341.46}{513.183}{A}%
\TextAt{9}{350.235}{513.183}{\Character{41}}%
\TextAt{13}{354.788}{513.183}{E}%
\TextAt{13}{364.137}{513.183}{\Character{59}}%
\TextAt{8}{79.226}{486.251}{where}%
\TextAt{13}{113.511}{486.251}{E}%
\TextAt{8}{127.249}{486.251}{denotes}%
\TextAt{8}{170.058}{486.251}{the}%
\TextAt{8}{190.69}{486.251}{unit}%
\TextAt{8}{215.876}{486.251}{matrix}%
\TextAt{8}{254.405}{486.251}{of}%
\TextAt{8}{268.21}{486.251}{rank}%
\TextAt{13}{295.68}{486.251}{r}%
\TextAt{8}{301.28}{486.251}{.}%
\TextAt{8}{311.154}{486.251}{For}%
\TextAt{8}{332.599}{486.251}{any}%
\TextAt{8}{355.191}{486.251}{vector}%
\TextAt{13}{390.791}{486.251}{x}%
\TextAt{9}{401.573}{486.251}{\Character{61}}%
\TextAt{9}{414.796}{486.251}{\Character{40}}%
\TextAt{13}{419.361}{486.251}{x}%
\TextAt{17}{426.013}{484.457}{1}%
\TextAt{13}{430.745}{486.251}{\Character{59}}%
\TextAt{13}{435.969}{486.251}{\Character{58}}%
\TextAt{13}{441.206}{486.251}{\Character{58}}%
\TextAt{13}{446.442}{486.251}{\Character{58}}%
\TextAt{13}{451.667}{486.251}{\Character{59}}%
\TextAt{13}{456.903}{486.251}{x}%
\TextAt{21}{463.618}{484.457}{r}%
\TextAt{9}{468.173}{486.251}{\Character{41}}%
\TextAt{8}{472.725}{486.251}{,}%
\TextAt{8}{480.469}{486.251}{this}%
\TextAt{8}{503.766}{486.251}{yields}%
\TextAt{8}{79.226}{469.205}{the}%
\TextAt{8}{99.383}{469.205}{implication}%
\TextAt{13}{228.857}{452.159}{Ax}%
\TextAt{9}{247.606}{452.159}{\Character{61}}%
\TextAt{9}{260.015}{452.159}{0}%
\NativeArrowAt{281.14}{452.159}{11.955}{19.087}{\Longrightarrow}%
\TextAt{9}{315.504}{452.159}{\Character{40}}%
\TextAt{8}{320.056}{452.159}{det}%
\TextAt{13}{338.307}{452.159}{A}%
\TextAt{9}{347.083}{452.159}{\Character{41}}%
\TextAt{13}{351.635}{452.159}{x}%
\TextAt{9}{361.608}{452.159}{\Character{61}}%
\TextAt{9}{374.018}{452.159}{0}%
\TextAt{13}{379.886}{452.159}{\Character{58}}%
\TextAt{19}{79.226}{412.773}{Proof}%
\TextAt{19}{110.62}{412.773}{of}%
\TextAt{19}{124.369}{412.773}{proposition}%
\TextAt{19}{184.599}{412.773}{(2.2).}%
\TextAt{8}{219.269}{412.773}{Let}%
\TextAt{13}{240.658}{412.773}{b}%
\NativeMathAt{249.671}{412.773}{11.955}{\in}%
\TextAt{13}{261.677}{412.773}{B}%
\TextAt{8}{275.49}{412.773}{be}%
\TextAt{8}{291.834}{412.773}{integral}%
\TextAt{8}{334.861}{412.773}{over}%
\TextAt{13}{360.327}{412.773}{A}%
\TextAt{8}{373.419}{412.773}{and}%
\TextAt{13}{396.601}{412.773}{f}%
\TextAt{9}{403.648}{412.773}{\Character{40}}%
\TextAt{13}{408.2}{412.773}{x}%
\TextAt{9}{414.852}{412.773}{\Character{41}}%
\NativeMathAt{423.44}{412.773}{11.955}{\in}%
\TextAt{13}{435.446}{412.773}{A}%
\TextAt{9}{444.222}{412.773}{\Character{91}}%
\TextAt{13}{447.473}{412.773}{x}%
\TextAt{9}{454.125}{412.773}{\Character{93}}%
\TextAt{8}{461.694}{412.773}{a}%
\TextAt{8}{471.869}{412.773}{monic}%
\TextAt{8}{506.755}{412.773}{poly-}%
\TextAt{8}{79.226}{395.727}{nomial}%
\TextAt{8}{117.881}{395.727}{of}%
\TextAt{8}{131.49}{395.727}{degree}%
\TextAt{13}{168.184}{395.727}{n}%
\NativeMathAt{178.967}{395.727}{11.955}{\geq}%
\TextAt{9}{192.061}{395.727}{1}%
\TextAt{8}{202.09}{395.727}{such}%
\TextAt{8}{228.766}{395.727}{that}%
\TextAt{13}{254.413}{395.727}{f}%
\TextAt{9}{261.46}{395.727}{\Character{40}}%
\TextAt{13}{266.012}{395.727}{b}%
\TextAt{9}{270.989}{395.727}{\Character{41}}%
\TextAt{9}{279.322}{395.727}{\Character{61}}%
\TextAt{9}{292.209}{395.727}{0}%
\TextAt{8}{298.09}{395.727}{.}%
\TextAt{8}{307.375}{395.727}{For}%
\TextAt{8}{328.623}{395.727}{an}%
\TextAt{8}{345.16}{395.727}{arbitrary}%
\TextAt{8}{395.208}{395.727}{polynomial}%
\TextAt{13}{455.98}{395.727}{g}%
\TextAt{9}{462.014}{395.727}{\Character{40}}%
\TextAt{13}{466.567}{395.727}{x}%
\TextAt{9}{473.219}{395.727}{\Character{41}}%
\NativeMathAt{481.567}{395.727}{11.955}{\in}%
\TextAt{13}{493.333}{395.727}{A}%
\TextAt{9}{502.108}{395.727}{\Character{91}}%
\TextAt{13}{505.36}{395.727}{x}%
\TextAt{9}{512.012}{395.727}{\Character{93}}%
\TextAt{8}{519.44}{395.727}{we}%
\TextAt{8}{79.226}{378.681}{may}%
\TextAt{8}{104.595}{378.681}{then}%
\TextAt{8}{131.255}{378.681}{write}%
\TextAt{13}{245.626}{361.635}{g}%
\TextAt{9}{251.66}{361.635}{\Character{40}}%
\TextAt{13}{256.213}{361.635}{x}%
\TextAt{9}{262.865}{361.635}{\Character{41}}%
\TextAt{9}{270.72}{361.635}{\Character{61}}%
\TextAt{13}{283.164}{361.635}{q}%
\TextAt{9}{288.783}{361.635}{\Character{40}}%
\TextAt{13}{293.335}{361.635}{x}%
\TextAt{9}{299.987}{361.635}{\Character{41}}%
\TextAt{13}{304.539}{361.635}{f}%
\TextAt{9}{311.586}{361.635}{\Character{40}}%
\TextAt{13}{316.138}{361.635}{x}%
\TextAt{9}{322.79}{361.635}{\Character{41}}%
\TextAt{9}{329.987}{361.635}{\Character{43}}%
\TextAt{13}{341.761}{361.635}{r}%
\TextAt{9}{347.361}{361.635}{\Character{40}}%
\TextAt{13}{351.913}{361.635}{x}%
\TextAt{9}{358.565}{361.635}{\Character{41}}%
\TextAt{13}{363.118}{361.635}{\Character{59}}%
\TextAt{8}{79.226}{338.999}{with}%
\TextAt{13}{105.886}{338.999}{q}%
\TextAt{9}{111.505}{338.999}{\Character{40}}%
\TextAt{13}{116.058}{338.999}{x}%
\TextAt{9}{122.71}{338.999}{\Character{41}}%
\TextAt{13}{127.262}{338.999}{\Character{59}}%
\TextAt{13}{132.486}{338.999}{r}%
\TextAt{9}{138.107}{338.999}{\Character{40}}%
\TextAt{13}{142.659}{338.999}{x}%
\TextAt{9}{149.311}{338.999}{\Character{41}}%
\NativeMathAt{157.184}{338.999}{11.955}{\in}%
\TextAt{13}{168.475}{338.999}{A}%
\TextAt{9}{177.251}{338.999}{\Character{91}}%
\TextAt{13}{180.502}{338.999}{x}%
\TextAt{9}{187.154}{338.999}{\Character{93}}%
\TextAt{8}{194.303}{338.999}{and}%
\TextAt{8}{217.066}{338.999}{deg}%
\TextAt{9}{234.628}{338.999}{\Character{40}}%
\TextAt{13}{239.181}{338.999}{r}%
\TextAt{9}{244.781}{338.999}{\Character{40}}%
\TextAt{13}{249.333}{338.999}{x}%
\TextAt{9}{255.985}{338.999}{\Character{41}\Character{41}}%
\TextAt{13}{268.411}{338.999}{\Character{60}}%
\TextAt{13}{280.821}{338.999}{n}%
\TextAt{8}{287.824}{338.999}{,}%
\TextAt{8}{294.973}{338.999}{so}%
\TextAt{8}{309.343}{338.999}{that}%
\TextAt{8}{334.712}{338.999}{one}%
\TextAt{8}{356.172}{338.999}{has}%
\TextAt{13}{204.742}{312.067}{g}%
\TextAt{9}{210.776}{312.067}{\Character{40}}%
\TextAt{13}{215.328}{312.067}{b}%
\TextAt{9}{220.306}{312.067}{\Character{41}}%
\TextAt{9}{228.161}{312.067}{\Character{61}}%
\TextAt{13}{240.604}{312.067}{r}%
\TextAt{9}{246.205}{312.067}{\Character{40}}%
\TextAt{13}{250.757}{312.067}{b}%
\TextAt{9}{255.734}{312.067}{\Character{41}}%
\TextAt{9}{263.589}{312.067}{\Character{61}}%
\TextAt{13}{276.033}{312.067}{a}%
\TextAt{17}{282.178}{310.273}{0}%
\TextAt{9}{289.567}{312.067}{\Character{43}}%
\TextAt{13}{301.328}{312.067}{a}%
\TextAt{17}{307.473}{310.273}{1}%
\TextAt{13}{312.205}{312.067}{b}%
\TextAt{9}{319.839}{312.067}{\Character{43}}%
\TextAt{16}{331.6}{312.067}{\Character{1}}%
\TextAt{16}{336.896}{312.067}{\Character{1}}%
\TextAt{16}{342.204}{312.067}{\Character{1}}%
\TextAt{9}{348.205}{312.067}{\Character{43}}%
\TextAt{13}{359.966}{312.067}{a}%
\TextAt{21}{366.111}{310.273}{n}%
\TextAt{3}{371.249}{310.273}{\Character{0}}%
\TextAt{17}{377.836}{310.273}{1}%
\TextAt{13}{382.568}{312.067}{b}%
\TextAt{21}{387.545}{317.003}{n}%
\TextAt{3}{392.683}{317.003}{\Character{0}}%
\TextAt{17}{399.27}{317.003}{1}%
\TextAt{13}{404.002}{312.067}{\Character{58}}%
\TextAt{8}{79.226}{285.134}{Thus}%
\TextAt{13}{108.875}{285.134}{A}%
\TextAt{9}{117.65}{285.134}{\Character{91}}%
\TextAt{13}{120.902}{285.134}{b}%
\TextAt{9}{125.879}{285.134}{\Character{93}}%
\TextAt{8}{133.028}{285.134}{is}%
\TextAt{8}{144.792}{285.134}{generated}%
\TextAt{8}{198.124}{285.134}{as}%
\TextAt{13}{212.494}{285.134}{A}%
\TextAt{8}{221.269}{285.134}{-module}%
\TextAt{8}{266.46}{285.134}{by}%
\TextAt{9}{282.719}{285.134}{1}%
\TextAt{13}{288.572}{285.134}{\Character{59}}%
\TextAt{13}{293.796}{285.134}{b\Character{59}}%
\TextAt{13}{304.006}{285.134}{\Character{58}}%
\TextAt{13}{309.231}{285.134}{\Character{58}}%
\TextAt{13}{314.467}{285.134}{\Character{58}}%
\TextAt{13}{319.703}{285.134}{\Character{59}}%
\TextAt{13}{324.928}{285.134}{b}%
\TextAt{21}{329.991}{289.472}{n}%
\TextAt{3}{335.129}{289.472}{\Character{0}}%
\TextAt{17}{341.716}{289.472}{1}%
\TextAt{8}{346.448}{285.134}{.}%
\TextAt{8}{97.159}{265.099}{More}%
\TextAt{8}{127.89}{265.099}{generally,}%
\TextAt{8}{180.539}{265.099}{if}%
\TextAt{13}{191.758}{265.099}{b}%
\TextAt{17}{196.736}{263.306}{1}%
\TextAt{13}{201.468}{265.099}{\Character{59}}%
\TextAt{13}{206.692}{265.099}{\Character{58}}%
\TextAt{13}{211.929}{265.099}{\Character{58}}%
\TextAt{13}{217.153}{265.099}{\Character{58}}%
\TextAt{13}{222.39}{265.099}{\Character{59}}%
\TextAt{13}{227.626}{265.099}{b}%
\TextAt{21}{232.666}{263.306}{n}%
\NativeMathAt{242.468}{265.099}{11.955}{\in}%
\TextAt{13}{254.605}{265.099}{B}%
\TextAt{8}{268.494}{265.099}{are}%
\TextAt{8}{288.501}{265.099}{integral}%
\TextAt{8}{331.605}{265.099}{over}%
\TextAt{13}{357.148}{265.099}{A}%
\TextAt{8}{365.923}{265.099}{,}%
\TextAt{8}{373.692}{265.099}{then}%
\TextAt{8}{400.848}{265.099}{the}%
\TextAt{8}{421.501}{265.099}{fact}%
\TextAt{8}{445.082}{265.099}{that}%
\TextAt{13}{470.947}{265.099}{A}%
\TextAt{9}{479.723}{265.099}{\Character{91}}%
\TextAt{13}{482.974}{265.099}{b}%
\TextAt{17}{487.951}{263.306}{1}%
\TextAt{13}{492.684}{265.099}{\Character{59}}%
\TextAt{13}{497.908}{265.099}{\Character{58}}%
\TextAt{13}{503.145}{265.099}{\Character{58}}%
\TextAt{13}{508.369}{265.099}{\Character{58}}%
\TextAt{13}{513.606}{265.099}{\Character{59}}%
\TextAt{13}{518.842}{265.099}{b}%
\TextAt{21}{523.882}{263.306}{n}%
\TextAt{9}{529.518}{265.099}{\Character{93}}%
\TextAt{8}{79.226}{248.053}{is}%
\TextAt{8}{91.503}{248.053}{of}%
\TextAt{8}{105.346}{248.053}{finite}%
\TextAt{8}{135.771}{248.053}{type}%
\TextAt{8}{162.621}{248.053}{over}%
\TextAt{13}{188.18}{248.053}{A}%
\TextAt{8}{201.366}{248.053}{follows}%
\TextAt{8}{240.315}{248.053}{by}%
\TextAt{8}{257.087}{248.053}{induction}%
\TextAt{8}{309.629}{248.053}{on}%
\TextAt{13}{326.401}{248.053}{n}%
\TextAt{8}{333.389}{248.053}{.}%
\TextAt{8}{343.376}{248.053}{Indeed,}%
\TextAt{8}{385.299}{248.053}{since}%
\TextAt{13}{414.48}{248.053}{b}%
\TextAt{21}{419.457}{246.26}{n}%
\TextAt{8}{429.504}{248.053}{is}%
\TextAt{8}{441.781}{248.053}{integral}%
\TextAt{8}{484.902}{248.053}{over}%
\TextAt{13}{510.461}{248.053}{R}%
\TextAt{9}{523.665}{248.053}{\Character{61}}%
\TextAt{13}{79.226}{231.007}{A}%
\TextAt{9}{88.002}{231.007}{\Character{91}}%
\TextAt{13}{91.253}{231.007}{b}%
\TextAt{17}{96.23}{229.214}{1}%
\TextAt{13}{100.963}{231.007}{\Character{59}}%
\TextAt{13}{106.187}{231.007}{\Character{58}}%
\TextAt{13}{111.424}{231.007}{\Character{58}}%
\TextAt{13}{116.648}{231.007}{\Character{58}}%
\TextAt{13}{121.885}{231.007}{\Character{59}}%
\TextAt{13}{127.121}{231.007}{b}%
\TextAt{21}{132.16}{229.214}{n}%
\TextAt{3}{137.299}{229.214}{\Character{0}}%
\TextAt{17}{143.885}{229.214}{1}%
\TextAt{9}{148.617}{231.007}{\Character{93}}%
\TextAt{8}{151.869}{231.007}{,}%
\TextAt{8}{161.005}{231.007}{what}%
\TextAt{8}{191.861}{231.007}{we}%
\TextAt{8}{210.679}{231.007}{have}%
\TextAt{8}{239.263}{231.007}{just}%
\TextAt{8}{263.998}{231.007}{shown}%
\TextAt{8}{301.094}{231.007}{implies}%
\TextAt{8}{342.411}{231.007}{that}%
\TextAt{13}{369.37}{231.007}{A}%
\TextAt{9}{378.145}{231.007}{\Character{91}}%
\TextAt{13}{381.397}{231.007}{b}%
\TextAt{17}{386.374}{229.214}{1}%
\TextAt{13}{391.106}{231.007}{\Character{59}}%
\TextAt{13}{396.33}{231.007}{\Character{58}}%
\TextAt{13}{401.567}{231.007}{\Character{58}}%
\TextAt{13}{406.803}{231.007}{\Character{58}}%
\TextAt{13}{412.028}{231.007}{\Character{59}}%
\TextAt{13}{417.264}{231.007}{b}%
\TextAt{21}{422.304}{229.214}{n}%
\TextAt{9}{427.94}{231.007}{\Character{93}}%
\TextAt{9}{437.205}{231.007}{\Character{61}}%
\TextAt{13}{452.357}{231.007}{R}%
\TextAt{9}{461.365}{231.007}{\Character{91}}%
\TextAt{13}{464.617}{231.007}{b}%
\TextAt{21}{469.594}{229.214}{n}%
\TextAt{9}{475.23}{231.007}{\Character{93}}%
\TextAt{8}{483.969}{231.007}{is}%
\TextAt{8}{497.322}{231.007}{finitely}%
\TextAt{8}{79.226}{213.961}{generated}%
\TextAt{8}{132.705}{213.961}{over}%
\TextAt{13}{157.898}{213.961}{R}%
\TextAt{8}{166.907}{213.961}{,}%
\TextAt{8}{174.239}{213.961}{hence}%
\TextAt{8}{206.892}{213.961}{also}%
\TextAt{8}{230.519}{213.961}{over}%
\TextAt{13}{255.712}{213.961}{A}%
\TextAt{8}{264.488}{213.961}{,}%
\TextAt{8}{271.821}{213.961}{if}%
\TextAt{8}{282.691}{213.961}{we}%
\TextAt{8}{300.066}{213.961}{assume,}%
\TextAt{8}{343.945}{213.961}{by}%
\TextAt{8}{360.351}{213.961}{induction,}%
\TextAt{8}{415.816}{213.961}{that}%
\TextAt{13}{441.331}{213.961}{R}%
\TextAt{8}{454.384}{213.961}{is}%
\TextAt{8}{466.295}{213.961}{an}%
\TextAt{13}{482.701}{213.961}{A}%
\TextAt{8}{491.476}{213.961}{-module}%
\TextAt{8}{79.226}{196.915}{of}%
\TextAt{8}{92.556}{196.915}{finite}%
\TextAt{8}{122.468}{196.915}{type.}%
\TextAt{8}{97.159}{176.88}{Conversely,}%
\TextAt{8}{161.09}{176.88}{assume}%
\TextAt{8}{203.462}{176.88}{that}%
\TextAt{8}{230.759}{176.88}{the}%
\TextAt{13}{252.843}{176.88}{A}%
\TextAt{8}{261.618}{176.88}{-module}%
\TextAt{13}{308.736}{176.88}{A}%
\TextAt{9}{317.512}{176.88}{\Character{91}}%
\TextAt{13}{320.763}{176.88}{b}%
\TextAt{17}{325.741}{175.087}{1}%
\TextAt{13}{330.473}{176.88}{\Character{59}}%
\TextAt{13}{335.697}{176.88}{\Character{58}}%
\TextAt{13}{340.934}{176.88}{\Character{58}}%
\TextAt{13}{346.158}{176.88}{\Character{58}}%
\TextAt{13}{351.395}{176.88}{\Character{59}}%
\TextAt{13}{356.631}{176.88}{b}%
\TextAt{21}{361.671}{175.087}{n}%
\TextAt{9}{367.307}{176.88}{\Character{93}}%
\TextAt{8}{376.384}{176.88}{is}%
\TextAt{8}{390.076}{176.88}{finitely}%
\TextAt{8}{431.348}{176.88}{generated}%
\TextAt{8}{486.608}{176.88}{and}%
\TextAt{8}{511.298}{176.88}{that}%
\TextAt{13}{79.226}{159.834}{\Character{33}}%
\TextAt{17}{86.515}{158.041}{1}%
\TextAt{13}{91.248}{159.834}{\Character{59}}%
\TextAt{13}{96.472}{159.834}{\Character{58}}%
\TextAt{13}{101.709}{159.834}{\Character{58}}%
\TextAt{13}{106.933}{159.834}{\Character{58}}%
\TextAt{13}{112.17}{159.834}{\Character{59}}%
\TextAt{13}{117.406}{159.834}{\Character{33}}%
\TextAt{21}{124.758}{158.041}{r}%
\TextAt{8}{133.963}{159.834}{is}%
\TextAt{8}{146.48}{159.834}{a}%
\TextAt{8}{156.989}{159.834}{system}%
\TextAt{8}{196.56}{159.834}{of}%
\TextAt{8}{210.643}{159.834}{generators.}%
\TextAt{8}{274.109}{159.834}{Then,}%
\TextAt{8}{308.86}{159.834}{for}%
\TextAt{8}{327.498}{159.834}{any}%
\TextAt{8}{350.368}{159.834}{element}%
\TextAt{13}{394.363}{159.834}{b}%
\NativeMathAt{403.944}{159.834}{11.955}{\in}%
\TextAt{13}{416.518}{159.834}{A}%
\TextAt{9}{425.294}{159.834}{\Character{91}}%
\TextAt{13}{428.545}{159.834}{b}%
\TextAt{17}{433.523}{158.041}{1}%
\TextAt{13}{438.255}{159.834}{\Character{59}}%
\TextAt{13}{443.479}{159.834}{\Character{58}}%
\TextAt{13}{448.716}{159.834}{\Character{58}}%
\TextAt{13}{453.94}{159.834}{\Character{58}}%
\TextAt{13}{459.177}{159.834}{\Character{59}}%
\TextAt{13}{464.413}{159.834}{b}%
\TextAt{21}{469.453}{158.041}{n}%
\TextAt{9}{475.089}{159.834}{\Character{93}}%
\TextAt{8}{478.341}{159.834}{,}%
\TextAt{8}{486.431}{159.834}{one}%
\TextAt{8}{508.644}{159.834}{finds}%
\TextAt{8}{79.226}{142.788}{that}%
\TextAt{13}{219.227}{122.101}{b}%
\TextAt{16}{226.861}{122.101}{\Character{1}}%
\TextAt{13}{232.839}{122.101}{\Character{33}}%
\TextAt{21}{240.128}{120.308}{i}%
\TextAt{9}{246.83}{122.101}{\Character{61}}%
\TextAt{21}{265.862}{137.045}{r}%
\TextAt{20}{259.256}{133.459}{\Character{88}}%
\TextAt{21}{260.538}{108.265}{j}%
\TextAt{17}{264.422}{108.265}{\Character{61}1}%
\TextAt{13}{278.517}{122.101}{a}%
\TextAt{21}{284.662}{120.308}{ij}%
\TextAt{13}{291.927}{122.101}{\Character{33}}%
\TextAt{21}{299.216}{120.308}{j}%
\TextAt{13}{303.598}{122.101}{\Character{59}}%
\TextAt{13}{332.733}{122.101}{i}%
\TextAt{9}{340.067}{122.101}{\Character{61}}%
\TextAt{9}{352.476}{122.101}{1}%
\TextAt{13}{358.346}{122.101}{\Character{59}}%
\TextAt{13}{363.57}{122.101}{\Character{58}}%
\TextAt{13}{368.807}{122.101}{\Character{58}}%
\TextAt{13}{374.031}{122.101}{\Character{58}}%
\TextAt{13}{379.268}{122.101}{\Character{59}}%
\TextAt{13}{384.504}{122.101}{r}%
\TextAt{13}{389.465}{122.101}{\Character{58}}%
\TextAt{8}{79.226}{89.967}{From}%
\TextAt{8}{110.346}{89.967}{(2.3)}%
\TextAt{8}{138.717}{89.967}{we}%
\TextAt{8}{156.34}{89.967}{see}%
\TextAt{8}{175.649}{89.967}{that}%
\TextAt{8}{201.413}{89.967}{det}%
\TextAt{9}{217.672}{89.967}{\Character{40}}%
\TextAt{13}{222.225}{89.967}{bE}%
\TextAt{16}{239.494}{89.967}{\Character{0}}%
\TextAt{9}{251.719}{89.967}{\Character{40}}%
\TextAt{13}{256.271}{89.967}{a}%
\TextAt{21}{262.416}{88.174}{ij}%
\TextAt{9}{269.682}{89.967}{\Character{41}\Character{41}}%
\TextAt{13}{278.786}{89.967}{\Character{33}}%
\TextAt{21}{286.075}{88.174}{i}%
\TextAt{9}{293.452}{89.967}{\Character{61}}%
\TextAt{9}{306.543}{89.967}{0}%
\TextAt{8}{312.404}{89.967}{,}%
\TextAt{13}{320.048}{89.967}{i}%
\TextAt{9}{328.036}{89.967}{\Character{61}}%
\TextAt{9}{341.127}{89.967}{1}%
\TextAt{13}{346.989}{89.967}{\Character{59}}%
\TextAt{13}{352.213}{89.967}{\Character{58}}%
\TextAt{13}{357.45}{89.967}{\Character{58}}%
\TextAt{13}{362.674}{89.967}{\Character{58}}%
\TextAt{13}{367.911}{89.967}{\Character{59}}%
\TextAt{13}{373.147}{89.967}{r}%
\TextAt{8}{383.103}{89.967}{(here}%
\TextAt{13}{413.411}{89.967}{E}%
\TextAt{8}{427.07}{89.967}{is}%
\TextAt{8}{439.229}{89.967}{the}%
\TextAt{8}{459.781}{89.967}{unit}%
\TextAt{8}{484.888}{89.967}{matrix}%
\TextAt{8}{523.337}{89.967}{of}%
\TextAt{8}{79.226}{72.921}{rank}%
\TextAt{13}{105.919}{72.921}{r}%
\TextAt{8}{111.52}{72.921}{),}%
\TextAt{8}{122.983}{72.921}{and}%
\TextAt{8}{145.444}{72.921}{since}%
\TextAt{9}{173.811}{72.921}{1}%
\TextAt{8}{183.26}{72.921}{can}%
\TextAt{8}{204.418}{72.921}{be}%
\TextAt{8}{220.041}{72.921}{written}%
\TextAt{9}{260.709}{72.921}{1}%
\TextAt{9}{269.867}{72.921}{\Character{61}}%
\TextAt{13}{282.309}{72.921}{c}%
\TextAt{17}{287.347}{71.128}{1}%
\TextAt{13}{292.079}{72.921}{\Character{33}}%
\TextAt{17}{299.368}{71.128}{1}%
\TextAt{9}{306.141}{72.921}{\Character{43}}%
\TextAt{16}{317.285}{72.921}{\Character{1}}%
\TextAt{16}{322.581}{72.921}{\Character{1}}%
\TextAt{16}{327.889}{72.921}{\Character{1}}%
\TextAt{9}{333.273}{72.921}{\Character{43}}%
\TextAt{13}{344.418}{72.921}{c}%
\TextAt{21}{349.456}{71.128}{r}%
\TextAt{13}{354.011}{72.921}{\Character{33}}%
\TextAt{21}{361.3}{71.128}{r}%
\TextAt{8}{365.855}{72.921}{,}%
\TextAt{8}{372.763}{72.921}{the}%
\TextAt{8}{392.617}{72.921}{identity}%
\TextAt{8}{435.57}{72.921}{det}%
\TextAt{9}{451.829}{72.921}{\Character{40}}%
\TextAt{13}{456.381}{72.921}{bE}%
\TextAt{16}{472.764}{72.921}{\Character{0}}%
\TextAt{9}{484.103}{72.921}{\Character{40}}%
\TextAt{13}{488.655}{72.921}{a}%
\TextAt{21}{494.8}{71.128}{ij}%
\TextAt{9}{502.066}{72.921}{\Character{41}\Character{41}}%
\TextAt{9}{514.464}{72.921}{\Character{61}}%
\TextAt{9}{526.885}{72.921}{0}%
\TextAt{8}{303.069}{43.033}{6}%
\Destination{78.23}{809.815}{nav.page.6}%
\Destination{79.36}{776.155}{nav.print.6}%
\Destination{79.36}{501.255}{nav.print.7}%
\Destination{79.36}{673.155}{nav.number.1.2.2}%
\Destination{79.36}{605.255}{nav.number.1.2.3}%
\Hotlink{184.099}{410.555}{26.106}{11.047}{nav.number.1.2.2}%
\Hotlink{109.846}{87.755}{25.078}{11.047}{nav.number.1.2.3}%
}

% Page 19 (body page 7)
\Page{%
\TextAt{2}{79.226}{796.943}{I.}%
\TextAt{2}{89.83}{796.943}{Algebraic}%
\TextAt{2}{138.648}{796.943}{Integers}%
\Vector{q 0 G 0 g 0.299 w 7.226 18.729 m 460.769 18.729 l S Q}%
\TextAt{8}{79.226}{761.124}{gives}%
\TextAt{8}{107.427}{761.124}{us}%
\TextAt{8}{121.963}{761.124}{a}%
\TextAt{8}{131.239}{761.124}{monic}%
\TextAt{8}{165.226}{761.124}{equation}%
\TextAt{8}{212.555}{761.124}{for}%
\TextAt{13}{229.96}{761.124}{b}%
\TextAt{8}{238.355}{761.124}{with}%
\TextAt{8}{264.535}{761.124}{coefficients}%
\TextAt{8}{323.294}{761.124}{in}%
\TextAt{13}{336.467}{761.124}{A}%
\TextAt{8}{345.242}{761.124}{.}%
\TextAt{8}{353.531}{761.124}{This}%
\TextAt{8}{379.771}{761.124}{shows}%
\TextAt{8}{412.909}{761.124}{that}%
\TextAt{13}{437.798}{761.124}{b}%
\TextAt{8}{446.193}{761.124}{is}%
\TextAt{8}{457.477}{761.124}{indeed}%
\TextAt{8}{494.059}{761.124}{integral}%
\TextAt{8}{79.226}{744.078}{over}%
\TextAt{13}{104.272}{744.078}{A}%
\TextAt{8}{113.048}{744.078}{.}%
\Vector{q 0 G 0 g 0.399 w 452.334 -29.922 m 452.334 -21.853 l S Q}%
\Vector{q 0 G 0 g 0.399 w 452.534 -22.052 m 459.707 -22.052 l S Q}%
\Vector{q 0 G 0 g 0.399 w 452.534 -29.723 m 459.707 -29.723 l S Q}%
\Vector{q 0 G 0 g 0.399 w 459.906 -29.922 m 459.906 -21.853 l S Q}%
\TextAt{8}{97.159}{715.077}{According}%
\TextAt{8}{153.565}{715.077}{to}%
\TextAt{8}{169.002}{715.077}{this}%
\TextAt{8}{192.95}{715.077}{proposition,}%
\TextAt{8}{259.142}{715.077}{if}%
\TextAt{13}{270.992}{715.077}{b}%
\TextAt{17}{275.969}{713.283}{1}%
\TextAt{13}{280.701}{715.077}{\Character{59}}%
\TextAt{13}{285.925}{715.077}{\Character{58}}%
\TextAt{13}{291.162}{715.077}{\Character{58}}%
\TextAt{13}{296.398}{715.077}{\Character{58}}%
\TextAt{13}{301.623}{715.077}{\Character{59}}%
\TextAt{13}{306.859}{715.077}{b}%
\TextAt{21}{311.899}{713.283}{n}%
\NativeMathAt{322.775}{715.077}{11.955}{\in}%
\TextAt{13}{335.984}{715.077}{B}%
\TextAt{8}{350.504}{715.077}{are}%
\TextAt{8}{371.141}{715.077}{integral}%
\TextAt{8}{414.875}{715.077}{over}%
\TextAt{13}{441.047}{715.077}{A}%
\TextAt{8}{449.822}{715.077}{,}%
\TextAt{8}{458.378}{715.077}{then}%
\TextAt{8}{486.164}{715.077}{so}%
\TextAt{8}{501.66}{715.077}{is}%
\TextAt{8}{514.55}{715.077}{any}%
\TextAt{8}{79.226}{698.031}{element}%
\TextAt{13}{123.417}{698.031}{b}%
\TextAt{8}{133.241}{698.031}{of}%
\TextAt{13}{147.52}{698.031}{A}%
\TextAt{9}{156.296}{698.031}{\Character{91}}%
\TextAt{13}{159.547}{698.031}{b}%
\TextAt{17}{164.524}{696.237}{1}%
\TextAt{13}{169.257}{698.031}{\Character{59}}%
\TextAt{13}{174.481}{698.031}{\Character{58}}%
\TextAt{13}{179.718}{698.031}{\Character{58}}%
\TextAt{13}{184.942}{698.031}{\Character{58}}%
\TextAt{13}{190.179}{698.031}{\Character{59}}%
\TextAt{13}{195.415}{698.031}{b}%
\TextAt{21}{200.455}{696.237}{n}%
\TextAt{9}{206.091}{698.031}{\Character{93}}%
\TextAt{8}{209.343}{698.031}{,}%
\TextAt{8}{217.678}{698.031}{because}%
\TextAt{13}{261.929}{698.031}{A}%
\TextAt{9}{270.705}{698.031}{\Character{91}}%
\TextAt{13}{273.956}{698.031}{b}%
\TextAt{17}{278.933}{696.237}{1}%
\TextAt{13}{283.666}{698.031}{\Character{59}}%
\TextAt{13}{288.89}{698.031}{\Character{58}}%
\TextAt{13}{294.127}{698.031}{\Character{58}}%
\TextAt{13}{299.351}{698.031}{\Character{58}}%
\TextAt{13}{304.588}{698.031}{\Character{59}}%
\TextAt{13}{309.824}{698.031}{b}%
\TextAt{21}{314.864}{696.237}{n}%
\TextAt{13}{320.5}{698.031}{\Character{59}}%
\TextAt{13}{325.724}{698.031}{b}%
\TextAt{9}{330.721}{698.031}{\Character{93}}%
\TextAt{9}{338.898}{698.031}{\Character{61}}%
\TextAt{13}{352.955}{698.031}{A}%
\TextAt{9}{361.73}{698.031}{\Character{91}}%
\TextAt{13}{364.982}{698.031}{b}%
\TextAt{17}{369.959}{696.237}{1}%
\TextAt{13}{374.691}{698.031}{\Character{59}}%
\TextAt{13}{379.915}{698.031}{\Character{58}}%
\TextAt{13}{385.152}{698.031}{\Character{58}}%
\TextAt{13}{390.376}{698.031}{\Character{58}}%
\TextAt{13}{395.613}{698.031}{\Character{59}}%
\TextAt{13}{400.849}{698.031}{b}%
\TextAt{21}{405.889}{696.237}{n}%
\TextAt{9}{411.525}{698.031}{\Character{93}}%
\TextAt{8}{419.623}{698.031}{is}%
\TextAt{8}{432.337}{698.031}{a}%
\TextAt{8}{443.041}{698.031}{finitely}%
\TextAt{8}{483.335}{698.031}{generated}%
\TextAt{13}{79.226}{680.985}{A}%
\TextAt{8}{88.001}{680.985}{-module.}%
\TextAt{8}{137.699}{680.985}{In}%
\TextAt{8}{152.187}{680.985}{particular,}%
\TextAt{8}{209.322}{680.985}{given}%
\TextAt{8}{239.759}{680.985}{two}%
\TextAt{8}{261.743}{680.985}{integral}%
\TextAt{8}{304.219}{680.985}{elements}%
\TextAt{13}{351.943}{680.985}{b}%
\TextAt{17}{356.92}{679.191}{1}%
\TextAt{13}{361.652}{680.985}{\Character{59}}%
\TextAt{13}{366.876}{680.985}{b}%
\TextAt{17}{371.873}{679.191}{2}%
\NativeMathAt{379.927}{680.985}{11.955}{\in}%
\TextAt{13}{391.217}{680.985}{B}%
\TextAt{8}{400.714}{680.985}{,}%
\TextAt{8}{407.757}{680.985}{then}%
\TextAt{13}{434.284}{680.985}{b}%
\TextAt{17}{439.261}{679.191}{1}%
\TextAt{9}{446.379}{680.985}{\Character{43}}%
\TextAt{13}{457.869}{680.985}{b}%
\TextAt{17}{462.846}{679.191}{2}%
\TextAt{8}{471.343}{680.985}{and}%
\TextAt{13}{493.973}{680.985}{b}%
\TextAt{17}{498.95}{679.191}{1}%
\TextAt{13}{503.682}{680.985}{b}%
\TextAt{17}{508.659}{679.191}{2}%
\TextAt{8}{517.156}{680.985}{are}%
\TextAt{8}{79.226}{663.939}{also}%
\TextAt{8}{102.706}{663.939}{integral}%
\TextAt{8}{145.314}{663.939}{over}%
\TextAt{13}{170.36}{663.939}{A}%
\TextAt{8}{179.136}{663.939}{.}%
\TextAt{8}{187.584}{663.939}{At}%
\TextAt{8}{204.489}{663.939}{the}%
\TextAt{8}{224.645}{663.939}{same}%
\TextAt{8}{253.971}{663.939}{time}%
\TextAt{8}{280.631}{663.939}{we}%
\TextAt{8}{297.859}{663.939}{obtain}%
\TextAt{8}{334.286}{663.939}{the}%
\TextAt{12}{79.226}{637.926}{Proposition}%
\TextAt{12}{152.024}{637.926}{(2.4).}%
\TextAt{8}{187.415}{637.926}{Let}%
\TextAt{13}{208.459}{637.926}{A}%
\NativeMathAt{220.683}{637.926}{11.955}{\subseteq}%
\TextAt{13}{233.431}{637.926}{B}%
\NativeMathAt{246.376}{637.926}{11.955}{\subseteq}%
\TextAt{13}{259.123}{637.926}{C}%
\TextAt{8}{272.329}{637.926}{be}%
\TextAt{8}{288.329}{637.926}{two}%
\TextAt{8}{310.521}{637.926}{ring}%
\TextAt{8}{334.662}{637.926}{extensions.}%
\TextAt{8}{395.819}{637.926}{If}%
\TextAt{13}{407.586}{637.926}{C}%
\TextAt{8}{420.792}{637.926}{is}%
\TextAt{8}{432.631}{637.926}{integral}%
\TextAt{8}{475.314}{637.926}{over}%
\TextAt{13}{500.436}{637.926}{B}%
\TextAt{8}{513.904}{637.926}{and}%
\TextAt{13}{79.226}{620.88}{B}%
\TextAt{8}{92.62}{620.88}{is}%
\TextAt{8}{104.384}{620.88}{integral}%
\TextAt{8}{146.992}{620.88}{over}%
\TextAt{13}{172.038}{620.88}{A}%
\TextAt{8}{180.813}{620.88}{,}%
\TextAt{8}{187.963}{620.88}{then}%
\TextAt{13}{214.623}{620.88}{C}%
\TextAt{8}{227.754}{620.88}{is}%
\TextAt{8}{239.517}{620.88}{integral}%
\TextAt{8}{282.126}{620.88}{over}%
\TextAt{13}{307.172}{620.88}{A}%
\TextAt{8}{315.947}{620.88}{.}%
\TextAt{19}{79.226}{585.902}{Proof.}%
\TextAt{8}{116}{585.902}{Take}%
\TextAt{13}{144.378}{585.902}{c}%
\NativeMathAt{152.894}{585.902}{11.955}{\in}%
\TextAt{13}{164.342}{585.902}{C}%
\TextAt{8}{173.575}{585.902}{,}%
\TextAt{8}{180.84}{585.902}{and}%
\TextAt{8}{203.694}{585.902}{let}%
\TextAt{13}{220.691}{585.902}{c}%
\TextAt{21}{225.729}{590.24}{n}%
\TextAt{9}{234.085}{585.902}{\Character{43}}%
\TextAt{13}{245.909}{585.902}{b}%
\TextAt{17}{250.886}{584.108}{1}%
\TextAt{13}{255.618}{585.902}{c}%
\TextAt{21}{260.656}{590.24}{n}%
\TextAt{3}{265.794}{590.24}{\Character{0}}%
\TextAt{17}{272.381}{590.24}{1}%
\TextAt{9}{279.833}{585.902}{\Character{43}}%
\TextAt{16}{291.657}{585.902}{\Character{1}}%
\TextAt{16}{296.953}{585.902}{\Character{1}}%
\TextAt{16}{302.261}{585.902}{\Character{1}}%
\TextAt{9}{308.324}{585.902}{\Character{43}}%
\TextAt{13}{320.148}{585.902}{b}%
\TextAt{21}{325.125}{584.108}{n}%
\TextAt{9}{334.239}{585.902}{\Character{61}}%
\TextAt{9}{346.804}{585.902}{0}%
\TextAt{8}{356.664}{585.902}{be}%
\TextAt{8}{372.68}{585.902}{an}%
\TextAt{8}{389.031}{585.902}{equation}%
\TextAt{8}{436.932}{585.902}{with}%
\TextAt{8}{463.684}{585.902}{coefficients}%
\TextAt{8}{523.014}{585.902}{in}%
\TextAt{13}{79.226}{568.856}{B}%
\TextAt{8}{88.722}{568.856}{.}%
\TextAt{8}{97.781}{568.856}{Write}%
\TextAt{13}{130.49}{568.856}{R}%
\TextAt{9}{143.166}{568.856}{\Character{61}}%
\TextAt{13}{155.938}{568.856}{A}%
\TextAt{9}{164.713}{568.856}{\Character{91}}%
\TextAt{13}{167.965}{568.856}{b}%
\TextAt{17}{172.942}{567.062}{1}%
\TextAt{13}{177.674}{568.856}{\Character{59}}%
\TextAt{13}{182.898}{568.856}{\Character{58}}%
\TextAt{13}{188.135}{568.856}{\Character{58}}%
\TextAt{13}{193.371}{568.856}{\Character{58}}%
\TextAt{13}{198.596}{568.856}{\Character{59}}%
\TextAt{13}{203.832}{568.856}{b}%
\TextAt{21}{208.872}{567.062}{n}%
\TextAt{9}{214.509}{568.856}{\Character{93}}%
\TextAt{8}{217.76}{568.856}{.}%
\TextAt{8}{226.819}{568.856}{Then}%
\TextAt{13}{257.579}{568.856}{R}%
\TextAt{9}{266.588}{568.856}{\Character{91}}%
\TextAt{13}{269.84}{568.856}{c}%
\TextAt{9}{274.878}{568.856}{\Character{93}}%
\TextAt{8}{282.23}{568.856}{is}%
\TextAt{8}{294.197}{568.856}{a}%
\TextAt{8}{304.156}{568.856}{finitely}%
\TextAt{8}{343.704}{568.856}{generated}%
\TextAt{13}{397.239}{568.856}{R}%
\TextAt{8}{406.248}{568.856}{-module.}%
\TextAt{8}{456.599}{568.856}{If}%
\TextAt{13}{468.495}{568.856}{B}%
\TextAt{8}{482.092}{568.856}{is}%
\TextAt{8}{494.059}{568.856}{integral}%
\TextAt{8}{79.226}{551.81}{over}%
\TextAt{13}{104.654}{551.81}{A}%
\TextAt{8}{113.43}{551.81}{,}%
\TextAt{8}{121.057}{551.81}{then}%
\TextAt{13}{148.099}{551.81}{R}%
\TextAt{9}{157.108}{551.81}{\Character{91}}%
\TextAt{13}{160.359}{551.81}{c}%
\TextAt{9}{165.397}{551.81}{\Character{93}}%
\TextAt{8}{172.928}{551.81}{is}%
\TextAt{8}{185.075}{551.81}{even}%
\TextAt{8}{212.117}{551.81}{finitely}%
\TextAt{8}{251.843}{551.81}{generated}%
\TextAt{8}{305.558}{551.81}{over}%
\TextAt{13}{330.986}{551.81}{A}%
\TextAt{8}{339.761}{551.81}{,}%
\TextAt{8}{347.388}{551.81}{since}%
\TextAt{13}{376.439}{551.81}{R}%
\TextAt{8}{389.727}{551.81}{is}%
\TextAt{8}{401.873}{551.81}{finitely}%
\TextAt{8}{441.6}{551.81}{generated}%
\TextAt{8}{495.314}{551.81}{over}%
\TextAt{13}{520.742}{551.81}{A}%
\TextAt{8}{529.518}{551.81}{.}%
\TextAt{8}{79.226}{534.764}{Thus}%
\TextAt{13}{108.875}{534.764}{c}%
\TextAt{8}{117.81}{534.764}{is}%
\TextAt{8}{129.574}{534.764}{integral}%
\TextAt{8}{172.182}{534.764}{over}%
\TextAt{13}{197.229}{534.764}{A}%
\TextAt{8}{206.004}{534.764}{.}%
\Vector{q 0 G 0 g 0.399 w 452.334 -239.236 m 452.334 -231.167 l S Q}%
\Vector{q 0 G 0 g 0.399 w 452.534 -231.366 m 459.707 -231.366 l S Q}%
\Vector{q 0 G 0 g 0.399 w 452.534 -239.037 m 459.707 -239.037 l S Q}%
\Vector{q 0 G 0 g 0.399 w 459.906 -239.236 m 459.906 -231.167 l S Q}%
\TextAt{8}{97.159}{505.763}{From}%
\TextAt{8}{127.884}{505.763}{what}%
\TextAt{8}{157.15}{505.763}{we}%
\TextAt{8}{174.377}{505.763}{have}%
\TextAt{8}{201.372}{505.763}{proven,}%
\TextAt{8}{242.677}{505.763}{the}%
\TextAt{8}{262.834}{505.763}{set}%
\TextAt{8}{281.101}{505.763}{of}%
\TextAt{8}{294.431}{505.763}{all}%
\TextAt{8}{310.69}{505.763}{elements}%
\TextAt{9}{228.671}{479.783}{\Character{22}}%
\TextAt{13}{225.583}{476.761}{A}%
\TextAt{9}{237.68}{476.761}{\Character{61}}%
\TextAt{16}{250.105}{476.761}{\Character{102}}%
\TextAt{13}{256.083}{476.761}{b}%
\NativeMathAt{264.381}{476.761}{11.955}{\in}%
\TextAt{13}{275.672}{476.761}{B}%
\TextAt{16}{288.489}{476.761}{\Character{106}}%
\TextAt{13}{295.13}{476.761}{b}%
\TextAt{8}{304.005}{476.761}{integral}%
\TextAt{8}{346.613}{476.761}{over}%
\TextAt{13}{371.659}{476.761}{A}%
\TextAt{16}{380.435}{476.761}{\Character{103}}%
\TextAt{8}{79.226}{447.76}{in}%
\TextAt{8}{92.638}{447.76}{a}%
\TextAt{8}{102.152}{447.76}{ring}%
\TextAt{8}{125.977}{447.76}{extension}%
\TextAt{13}{177.501}{447.76}{A}%
\NativeMathAt{189.598}{447.76}{11.955}{\subseteq}%
\TextAt{13}{202.217}{447.76}{B}%
\TextAt{8}{215.37}{447.76}{forms}%
\TextAt{8}{247.383}{447.76}{a}%
\TextAt{8}{256.898}{447.76}{ring.}%
\TextAt{8}{285.434}{447.76}{It}%
\TextAt{8}{297.865}{447.76}{is}%
\TextAt{8}{309.388}{447.76}{called}%
\TextAt{8}{342.311}{447.76}{the}%
\TextAt{12}{362.226}{447.76}{integral}%
\TextAt{12}{411.156}{447.76}{closure}%
\TextAt{8}{455.675}{447.76}{of}%
\TextAt{13}{468.764}{447.76}{A}%
\TextAt{8}{481.196}{447.76}{in}%
\TextAt{13}{494.607}{447.76}{B}%
\TextAt{8}{504.103}{447.76}{.}%
\TextAt{13}{512.471}{447.76}{A}%
\TextAt{8}{524.903}{447.76}{is}%
\TextAt{8}{79.226}{430.714}{said}%
\TextAt{8}{103.111}{430.714}{to}%
\TextAt{8}{117.182}{430.714}{be}%
\TextAt{12}{132.866}{430.714}{integrally}%
\TextAt{12}{192.64}{430.714}{closed}%
\TextAt{12}{232.21}{430.714}{in}%
\TextAt{13}{247.631}{430.714}{B}%
\TextAt{8}{260.785}{430.714}{if}%
\TextAt{13}{271.268}{430.714}{A}%
\TextAt{9}{283.365}{430.714}{\Character{61}}%
\TextAt{9}{298.877}{433.736}{\Character{22}}%
\TextAt{13}{295.79}{430.714}{A}%
\TextAt{8}{304.565}{430.714}{.}%
\TextAt{8}{312.934}{430.714}{It}%
\TextAt{8}{325.366}{430.714}{is}%
\TextAt{8}{336.89}{430.714}{immediate}%
\TextAt{8}{393.879}{430.714}{from}%
\TextAt{8}{421.279}{430.714}{(2.4)}%
\TextAt{8}{449.014}{430.714}{that}%
\TextAt{8}{474.142}{430.714}{the}%
\TextAt{8}{494.059}{430.714}{integral}%
\TextAt{8}{79.226}{413.668}{closure}%
\TextAt{9}{121.246}{416.69}{\Character{22}}%
\TextAt{13}{118.159}{413.668}{A}%
\TextAt{8}{130.683}{413.668}{is}%
\TextAt{8}{142.298}{413.668}{itself}%
\TextAt{8}{170.495}{413.668}{integrally}%
\TextAt{8}{222.388}{413.668}{closed}%
\TextAt{8}{256.766}{413.668}{in}%
\TextAt{13}{270.27}{413.668}{B}%
\TextAt{8}{279.766}{413.668}{.}%
\TextAt{8}{288.165}{413.668}{If}%
\TextAt{13}{299.708}{413.668}{A}%
\TextAt{8}{312.233}{413.668}{is}%
\TextAt{8}{323.848}{413.668}{an}%
\TextAt{8}{339.958}{413.668}{integral}%
\TextAt{8}{382.418}{413.668}{domain}%
\TextAt{8}{423.897}{413.668}{with}%
\TextAt{8}{450.409}{413.668}{field}%
\TextAt{8}{475.617}{413.668}{of}%
\TextAt{8}{488.798}{413.668}{fractions}%
\TextAt{13}{79.226}{396.622}{K}%
\TextAt{8}{90.044}{396.622}{,}%
\TextAt{8}{98.399}{396.622}{then}%
\TextAt{8}{126.024}{396.622}{the}%
\TextAt{8}{147.146}{396.622}{integral}%
\TextAt{8}{190.719}{396.622}{closure}%
\TextAt{9}{233.852}{399.644}{\Character{22}}%
\TextAt{13}{230.765}{396.622}{A}%
\TextAt{8}{244.403}{396.622}{of}%
\TextAt{13}{258.698}{396.622}{A}%
\TextAt{8}{272.335}{396.622}{in}%
\TextAt{13}{286.953}{396.622}{K}%
\TextAt{8}{302.633}{396.622}{is}%
\TextAt{8}{315.362}{396.622}{called}%
\TextAt{8}{349.491}{396.622}{the}%
\TextAt{12}{370.612}{396.622}{normalization}%
\TextAt{8}{456.112}{396.622}{of}%
\TextAt{13}{470.407}{396.622}{A}%
\TextAt{8}{479.182}{396.622}{,}%
\TextAt{8}{487.538}{396.622}{and}%
\TextAt{13}{511.265}{396.622}{A}%
\TextAt{8}{524.903}{396.622}{is}%
\TextAt{8}{79.226}{379.576}{simply}%
\TextAt{8}{116.716}{379.576}{called}%
\TextAt{12}{149.914}{379.576}{integrally}%
\TextAt{12}{210.004}{379.576}{closed}%
\TextAt{8}{249.3}{379.576}{if}%
\TextAt{13}{260.058}{379.576}{A}%
\TextAt{9}{272.212}{379.576}{\Character{61}}%
\TextAt{9}{287.783}{382.598}{\Character{22}}%
\TextAt{13}{284.696}{379.576}{A}%
\TextAt{8}{293.472}{379.576}{.}%
\TextAt{8}{302.023}{379.576}{For}%
\TextAt{8}{323.027}{379.576}{instance,}%
\TextAt{8}{371.907}{379.576}{every}%
\TextAt{8}{402.833}{379.576}{factorial}%
\TextAt{8}{448.728}{379.576}{ring}%
\TextAt{8}{472.828}{379.576}{is}%
\TextAt{8}{484.626}{379.576}{integrally}%
\TextAt{8}{79.226}{362.53}{closed.}%
\TextAt{8}{97.159}{342.496}{In}%
\TextAt{8}{111.78}{342.496}{fact,}%
\TextAt{8}{138.117}{342.496}{if}%
\TextAt{13}{148.841}{342.496}{a}%
\TextAt{9}{154.986}{342.496}{\Character{47}}%
\TextAt{13}{160.839}{342.496}{b}%
\NativeMathAt{169.137}{342.496}{11.955}{\in}%
\TextAt{13}{180.428}{342.496}{K}%
\TextAt{8}{195.143}{342.496}{(}%
\TextAt{13}{199.698}{342.496}{a\Character{59}}%
\TextAt{13}{211.067}{342.496}{b}%
\NativeMathAt{219.385}{342.496}{11.955}{\in}%
\TextAt{13}{230.676}{342.496}{A}%
\TextAt{8}{239.451}{342.496}{)}%
\TextAt{8}{247.904}{342.496}{is}%
\TextAt{8}{259.668}{342.496}{integral}%
\TextAt{8}{302.276}{342.496}{over}%
\TextAt{13}{327.322}{342.496}{A}%
\TextAt{8}{336.097}{342.496}{,}%
\TextAt{8}{343.246}{342.496}{i.e.,}%
\TextAt{9}{216.516}{313.495}{\Character{40}}%
\TextAt{13}{221.068}{313.495}{a}%
\TextAt{9}{227.213}{313.495}{\Character{47}}%
\TextAt{13}{233.066}{313.495}{b}%
\TextAt{9}{238.043}{313.495}{\Character{41}}%
\TextAt{21}{242.596}{318.431}{n}%
\TextAt{9}{250.889}{313.495}{\Character{43}}%
\TextAt{13}{262.65}{313.495}{a}%
\TextAt{17}{268.795}{311.701}{1}%
\TextAt{9}{273.527}{313.495}{\Character{40}}%
\TextAt{13}{278.079}{313.495}{a}%
\TextAt{9}{284.224}{313.495}{\Character{47}}%
\TextAt{13}{290.077}{313.495}{b}%
\TextAt{9}{295.055}{313.495}{\Character{41}}%
\TextAt{21}{299.607}{318.431}{n}%
\TextAt{3}{304.745}{318.431}{\Character{0}}%
\TextAt{17}{311.332}{318.431}{1}%
\TextAt{9}{318.721}{313.495}{\Character{43}}%
\TextAt{16}{330.482}{313.495}{\Character{1}}%
\TextAt{16}{335.778}{313.495}{\Character{1}}%
\TextAt{16}{341.086}{313.495}{\Character{1}}%
\TextAt{9}{347.086}{313.495}{\Character{43}}%
\TextAt{13}{358.847}{313.495}{a}%
\TextAt{21}{364.992}{311.701}{n}%
\TextAt{9}{373.95}{313.495}{\Character{61}}%
\TextAt{9}{386.36}{313.495}{0}%
\TextAt{13}{392.228}{313.495}{\Character{59}}%
\TextAt{8}{79.226}{284.493}{with}%
\TextAt{13}{105.886}{284.493}{a}%
\TextAt{21}{112.031}{282.7}{i}%
\NativeMathAt{118.733}{284.493}{11.955}{\in}%
\TextAt{13}{130.024}{284.493}{A}%
\TextAt{8}{138.8}{284.493}{,}%
\TextAt{8}{145.949}{284.493}{then}%
\TextAt{13}{228.655}{267.447}{a}%
\TextAt{21}{234.8}{272.384}{n}%
\TextAt{9}{243.093}{267.447}{\Character{43}}%
\TextAt{13}{254.855}{267.447}{a}%
\TextAt{17}{261}{265.654}{1}%
\TextAt{13}{265.732}{267.447}{ba}%
\TextAt{21}{276.854}{272.384}{n}%
\TextAt{3}{281.992}{272.384}{\Character{0}}%
\TextAt{17}{288.579}{272.384}{1}%
\TextAt{9}{295.968}{267.447}{\Character{43}}%
\TextAt{16}{307.729}{267.447}{\Character{1}}%
\TextAt{16}{313.025}{267.447}{\Character{1}}%
\TextAt{16}{318.333}{267.447}{\Character{1}}%
\TextAt{9}{324.333}{267.447}{\Character{43}}%
\TextAt{13}{336.095}{267.447}{a}%
\TextAt{21}{342.24}{265.654}{n}%
\TextAt{13}{347.876}{267.447}{b}%
\TextAt{21}{352.853}{272.384}{n}%
\TextAt{9}{361.81}{267.447}{\Character{61}}%
\TextAt{9}{374.22}{267.447}{0}%
\TextAt{13}{380.089}{267.447}{\Character{58}}%
\TextAt{8}{79.226}{243.926}{Therefore}%
\TextAt{8}{133.887}{243.926}{each}%
\TextAt{8}{161.887}{243.926}{prime}%
\TextAt{8}{196.714}{243.926}{element}%
\TextAt{13}{241.619}{243.926}{\Character{25}}%
\TextAt{8}{254.248}{243.926}{which}%
\TextAt{8}{289.398}{243.926}{divides}%
\TextAt{13}{330.465}{243.926}{b}%
\TextAt{8}{341.003}{243.926}{also}%
\TextAt{8}{366.146}{243.926}{divides}%
\TextAt{13}{407.213}{243.926}{a}%
\TextAt{8}{413.358}{243.926}{.}%
\TextAt{8}{426.796}{243.926}{Assuming}%
\TextAt{13}{482.234}{243.926}{a}%
\TextAt{9}{488.379}{243.926}{\Character{47}}%
\TextAt{13}{494.232}{243.926}{b}%
\TextAt{8}{504.769}{243.926}{to}%
\TextAt{8}{520.743}{243.926}{be}%
\TextAt{8}{79.226}{226.88}{reduced,}%
\TextAt{8}{126.043}{226.88}{this}%
\TextAt{8}{148.865}{226.88}{implies}%
\TextAt{13}{188.592}{226.88}{a}%
\TextAt{9}{194.737}{226.88}{\Character{47}}%
\TextAt{13}{200.59}{226.88}{b}%
\NativeMathAt{208.888}{226.88}{11.955}{\in}%
\TextAt{13}{220.179}{226.88}{A}%
\TextAt{8}{228.954}{226.88}{.}%
\TextAt{8}{97.159}{206.845}{We}%
\TextAt{8}{117.813}{206.845}{now}%
\TextAt{8}{142.711}{206.845}{turn}%
\TextAt{8}{169.234}{206.845}{to}%
\TextAt{8}{184.054}{206.845}{a}%
\TextAt{8}{194.319}{206.845}{more}%
\TextAt{8}{224.095}{206.845}{specialized}%
\TextAt{8}{282.862}{206.845}{situation.}%
\TextAt{8}{337.789}{206.845}{Let}%
\TextAt{13}{359.268}{206.845}{A}%
\TextAt{8}{372.45}{206.845}{be}%
\TextAt{8}{388.884}{206.845}{an}%
\TextAt{8}{405.652}{206.845}{integral}%
\TextAt{8}{448.77}{206.845}{domain}%
\TextAt{8}{490.907}{206.845}{which}%
\TextAt{8}{524.903}{206.845}{is}%
\TextAt{8}{79.226}{189.799}{integrally}%
\TextAt{8}{130.932}{189.799}{closed,}%
\TextAt{13}{168.442}{189.799}{K}%
\TextAt{8}{182.822}{189.799}{its}%
\TextAt{8}{198.806}{189.799}{field}%
\TextAt{8}{223.828}{189.799}{of}%
\TextAt{8}{236.823}{189.799}{fractions,}%
\TextAt{13}{287.675}{189.799}{L}%
\TextAt{16}{298.961}{189.799}{\Character{106}}%
\TextAt{13}{305.602}{189.799}{K}%
\TextAt{8}{319.982}{189.799}{a}%
\TextAt{8}{329.403}{189.799}{finite}%
\TextAt{8}{358.979}{189.799}{field}%
\TextAt{8}{384.001}{189.799}{extension,}%
\TextAt{8}{438.751}{189.799}{and}%
\TextAt{13}{461.179}{189.799}{B}%
\TextAt{8}{474.237}{189.799}{the}%
\TextAt{8}{494.059}{189.799}{integral}%
\TextAt{8}{79.226}{172.753}{closure}%
\TextAt{8}{118.229}{172.753}{of}%
\TextAt{13}{131.481}{172.753}{A}%
\TextAt{8}{144.075}{172.753}{in}%
\TextAt{13}{157.649}{172.753}{L}%
\TextAt{8}{165.614}{172.753}{.}%
\TextAt{8}{174.036}{172.753}{According}%
\TextAt{8}{229.238}{172.753}{to}%
\TextAt{8}{243.47}{172.753}{(2.4),}%
\TextAt{13}{274.634}{172.753}{B}%
\TextAt{8}{287.949}{172.753}{is}%
\TextAt{8}{299.634}{172.753}{automatically}%
\TextAt{8}{373.391}{172.753}{integrally}%
\TextAt{8}{425.353}{172.753}{closed.}%
\TextAt{8}{464.405}{172.753}{Each}%
\TextAt{8}{493.425}{172.753}{element}%
\TextAt{13}{79.226}{155.707}{\Character{12}}%
\NativeMathAt{89.818}{155.707}{11.955}{\in}%
\TextAt{13}{101.109}{155.707}{L}%
\TextAt{8}{112.971}{155.707}{is}%
\TextAt{8}{124.735}{155.707}{of}%
\TextAt{8}{138.065}{155.707}{the}%
\TextAt{8}{158.221}{155.707}{form}%
\TextAt{13}{249.398}{144.084}{b}%
\Vector{q 0 G 0 g 0.478 w 176.814 -635.015 m 182.959 -635.015 l S Q}%
\TextAt{13}{248.814}{127.796}{a}%
\TextAt{13}{256.154}{135.996}{\Character{59}}%
\TextAt{13}{285.289}{135.996}{b}%
\NativeMathAt{293.607}{135.996}{11.955}{\in}%
\TextAt{13}{304.898}{135.996}{B}%
\TextAt{13}{314.379}{135.996}{\Character{59}}%
\TextAt{13}{331.57}{135.996}{a}%
\NativeMathAt{341.059}{135.996}{11.955}{\in}%
\TextAt{13}{352.35}{135.996}{A\Character{59}}%
\TextAt{8}{79.226}{112.475}{because}%
\TextAt{8}{122.528}{112.475}{if}%
\TextAt{13}{181.79}{95.429}{a}%
\TextAt{21}{187.935}{93.635}{n}%
\TextAt{13}{193.572}{95.429}{\Character{12}}%
\TextAt{21}{200.843}{100.365}{n}%
\TextAt{9}{209.136}{95.429}{\Character{43}}%
\TextAt{16}{220.897}{95.429}{\Character{1}}%
\TextAt{16}{226.193}{95.429}{\Character{1}}%
\TextAt{16}{231.501}{95.429}{\Character{1}}%
\TextAt{9}{237.502}{95.429}{\Character{43}}%
\TextAt{13}{249.263}{95.429}{a}%
\TextAt{17}{255.408}{93.635}{1}%
\TextAt{13}{260.14}{95.429}{\Character{12}}%
\TextAt{9}{270.068}{95.429}{\Character{43}}%
\TextAt{13}{281.829}{95.429}{a}%
\TextAt{17}{287.974}{93.635}{0}%
\TextAt{9}{296.027}{95.429}{\Character{61}}%
\TextAt{9}{308.437}{95.429}{0}%
\TextAt{13}{314.306}{95.429}{\Character{59}}%
\TextAt{13}{343.441}{95.429}{a}%
\TextAt{21}{349.605}{93.635}{i}%
\NativeMathAt{356.307}{95.429}{11.955}{\in}%
\TextAt{13}{367.598}{95.429}{A\Character{59}}%
\TextAt{13}{393.553}{95.429}{a}%
\TextAt{21}{399.718}{93.635}{n}%
\NativeMathAt{408.675}{95.429}{11.955}{\neq}%
\TextAt{9}{421.084}{95.429}{0}%
\TextAt{13}{426.954}{95.429}{\Character{59}}%
\TextAt{8}{303.069}{43.033}{7}%
\Destination{78.23}{809.815}{nav.page.7}%
\Destination{79.36}{258.955}{nav.print.8}%
\Destination{79.36}{651.955}{nav.number.1.2.4}%
\Hotlink{420.779}{428.555}{25.078}{11.047}{nav.number.1.2.4}%
\Hotlink{242.97}{170.555}{25.078}{11.047}{nav.number.1.2.4}%
}

% Page 20 (body page 8)
\Page{%
\TextAt{2}{79.226}{796.943}{I.}%
\TextAt{2}{89.83}{796.943}{Algebraic}%
\TextAt{2}{138.648}{796.943}{Integers}%
\Vector{q 0 G 0 g 0.299 w 7.226 18.729 m 460.769 18.729 l S Q}%
\TextAt{8}{79.226}{761.124}{then}%
\TextAt{13}{105.886}{761.124}{b}%
\TextAt{9}{114.184}{761.124}{\Character{61}}%
\TextAt{13}{126.61}{761.124}{a}%
\TextAt{21}{132.755}{759.33}{n}%
\TextAt{13}{138.391}{761.124}{\Character{12}}%
\TextAt{8}{149.56}{761.124}{is}%
\TextAt{8}{161.323}{761.124}{integral}%
\TextAt{8}{203.932}{761.124}{over}%
\TextAt{13}{228.978}{761.124}{A}%
\TextAt{8}{237.753}{761.124}{,}%
\TextAt{8}{244.902}{761.124}{an}%
\TextAt{8}{261.161}{761.124}{integral}%
\TextAt{8}{303.77}{761.124}{equation}%
\TextAt{9}{161.326}{732.979}{\Character{40}}%
\TextAt{13}{165.878}{732.979}{a}%
\TextAt{21}{172.023}{731.186}{n}%
\TextAt{13}{177.659}{732.979}{\Character{12}}%
\TextAt{9}{184.931}{732.979}{\Character{41}}%
\TextAt{21}{189.483}{737.915}{n}%
\TextAt{9}{197.776}{732.979}{\Character{43}}%
\TextAt{13}{209.537}{732.979}{a}%
\TextAt{21}{215.682}{731.186}{n}%
\TextAt{3}{220.82}{731.186}{\Character{0}}%
\TextAt{17}{227.407}{731.186}{1}%
\TextAt{9}{232.139}{732.979}{\Character{40}}%
\TextAt{13}{236.691}{732.979}{a}%
\TextAt{21}{242.836}{731.186}{n}%
\TextAt{13}{248.473}{732.979}{\Character{12}}%
\TextAt{9}{255.744}{732.979}{\Character{41}}%
\TextAt{21}{260.296}{737.915}{n}%
\TextAt{3}{265.435}{737.915}{\Character{0}}%
\TextAt{17}{272.021}{737.915}{1}%
\TextAt{9}{279.41}{732.979}{\Character{43}}%
\TextAt{16}{291.171}{732.979}{\Character{1}}%
\TextAt{16}{296.467}{732.979}{\Character{1}}%
\TextAt{16}{301.775}{732.979}{\Character{1}}%
\TextAt{9}{307.776}{732.979}{\Character{43}}%
\TextAt{13}{319.537}{732.979}{a}%
\TextAt{17}{325.682}{731.186}{1}%
\TextAt{13}{330.414}{732.979}{a}%
\TextAt{21}{336.559}{737.915}{n}%
\TextAt{3}{341.697}{737.915}{\Character{0}}%
\TextAt{17}{348.284}{737.915}{2}%
\TextAt{21}{336.559}{730.023}{n}%
\TextAt{9}{353.016}{732.979}{\Character{40}}%
\TextAt{13}{357.569}{732.979}{a}%
\TextAt{21}{363.714}{731.186}{n}%
\TextAt{13}{369.35}{732.979}{\Character{12}}%
\TextAt{9}{376.621}{732.979}{\Character{41}}%
\TextAt{9}{383.818}{732.979}{\Character{43}}%
\TextAt{13}{395.591}{732.979}{a}%
\TextAt{17}{401.736}{731.186}{0}%
\TextAt{13}{406.469}{732.979}{a}%
\TextAt{21}{412.614}{737.915}{n}%
\TextAt{3}{417.752}{737.915}{\Character{0}}%
\TextAt{17}{424.338}{737.915}{1}%
\TextAt{21}{412.614}{730.023}{n}%
\TextAt{9}{432.391}{732.979}{\Character{61}}%
\TextAt{9}{444.8}{732.979}{0}%
\TextAt{8}{79.226}{704.834}{being}%
\TextAt{8}{110.772}{704.834}{obtained}%
\TextAt{8}{158.912}{704.834}{from}%
\TextAt{8}{186.561}{704.834}{the}%
\TextAt{8}{206.725}{704.834}{equation}%
\TextAt{8}{254.542}{704.834}{for}%
\TextAt{13}{272.436}{704.834}{\Character{12}}%
\TextAt{8}{283.612}{704.834}{by}%
\TextAt{8}{299.88}{704.834}{multiplication}%
\TextAt{8}{375.014}{704.834}{by}%
\TextAt{13}{391.282}{704.834}{a}%
\TextAt{21}{397.427}{709.173}{n}%
\TextAt{3}{402.565}{709.173}{\Character{0}}%
\TextAt{17}{409.151}{709.173}{1}%
\TextAt{21}{397.427}{701.879}{n}%
\TextAt{8}{413.884}{704.834}{.}%
\TextAt{8}{422.357}{704.834}{Furthermore,}%
\TextAt{8}{493.417}{704.834}{the}%
\TextAt{8}{513.581}{704.834}{fact}%
\TextAt{8}{79.226}{687.788}{that}%
\TextAt{13}{104.838}{687.788}{A}%
\TextAt{8}{117.754}{687.788}{is}%
\TextAt{8}{129.762}{687.788}{integrally}%
\TextAt{8}{182.046}{687.788}{closed}%
\TextAt{8}{216.816}{687.788}{has}%
\TextAt{8}{237.933}{687.788}{the}%
\TextAt{8}{258.332}{687.788}{effect}%
\TextAt{8}{289.456}{687.788}{that}%
\TextAt{8}{315.068}{687.788}{an}%
\TextAt{8}{331.57}{687.788}{element}%
\TextAt{13}{375.055}{687.788}{\Character{12}}%
\NativeMathAt{386.062}{687.788}{11.955}{\in}%
\TextAt{13}{397.768}{687.788}{L}%
\TextAt{8}{409.873}{687.788}{is}%
\TextAt{8}{421.88}{687.788}{integral}%
\TextAt{8}{464.732}{687.788}{over}%
\TextAt{13}{490.021}{687.788}{A}%
\TextAt{8}{502.937}{687.788}{if}%
\TextAt{8}{513.904}{687.788}{and}%
\TextAt{8}{79.226}{670.742}{only}%
\TextAt{8}{105.818}{670.742}{if}%
\TextAt{8}{117.441}{670.742}{its}%
\TextAt{8}{134.66}{670.742}{minimal}%
\TextAt{8}{181.085}{670.742}{polynomial}%
\TextAt{13}{242.477}{670.742}{p}%
\TextAt{9}{248.352}{670.742}{\Character{40}}%
\TextAt{13}{252.905}{670.742}{x}%
\TextAt{9}{259.557}{670.742}{\Character{41}}%
\TextAt{8}{268.906}{670.742}{takes}%
\TextAt{8}{299.79}{670.742}{its}%
\TextAt{8}{317.008}{670.742}{coefficients}%
\TextAt{8}{377.146}{670.742}{in}%
\TextAt{13}{391.698}{670.742}{A}%
\TextAt{8}{400.474}{670.742}{.}%
\TextAt{8}{411.621}{670.742}{In}%
\TextAt{8}{427.142}{670.742}{fact,}%
\TextAt{8}{454.604}{670.742}{let}%
\TextAt{13}{472.408}{670.742}{\Character{12}}%
\TextAt{8}{484.476}{670.742}{be}%
\TextAt{8}{501.3}{670.742}{a}%
\TextAt{8}{511.956}{670.742}{zero}%
\TextAt{8}{79.226}{653.696}{of}%
\TextAt{8}{92.831}{653.696}{the}%
\TextAt{8}{113.263}{653.696}{monic}%
\TextAt{8}{148.005}{653.696}{polynomial}%
\TextAt{13}{208.773}{653.696}{g}%
\TextAt{9}{214.808}{653.696}{\Character{40}}%
\TextAt{13}{219.36}{653.696}{x}%
\TextAt{9}{226.012}{653.696}{\Character{41}}%
\NativeMathAt{234.354}{653.696}{11.955}{\in}%
\TextAt{13}{246.114}{653.696}{A}%
\TextAt{9}{254.89}{653.696}{\Character{91}}%
\TextAt{13}{258.141}{653.696}{x}%
\TextAt{9}{264.793}{653.696}{\Character{93}}%
\TextAt{8}{268.045}{653.696}{.}%
\TextAt{8}{277.319}{653.696}{Then}%
\TextAt{13}{308.152}{653.696}{p}%
\TextAt{9}{314.027}{653.696}{\Character{40}}%
\TextAt{13}{318.579}{653.696}{x}%
\TextAt{9}{325.231}{653.696}{\Character{41}}%
\TextAt{8}{333.956}{653.696}{divides}%
\TextAt{13}{373.636}{653.696}{g}%
\TextAt{9}{379.67}{653.696}{\Character{40}}%
\TextAt{13}{384.222}{653.696}{x}%
\TextAt{9}{390.874}{653.696}{\Character{41}}%
\TextAt{8}{399.599}{653.696}{in}%
\TextAt{13}{413.527}{653.696}{K}%
\TextAt{9}{424.345}{653.696}{\Character{91}}%
\TextAt{13}{427.597}{653.696}{x}%
\TextAt{9}{434.249}{653.696}{\Character{93}}%
\TextAt{8}{437.501}{653.696}{,}%
\TextAt{8}{444.994}{653.696}{so}%
\TextAt{8}{459.639}{653.696}{that}%
\TextAt{8}{485.283}{653.696}{all}%
\TextAt{8}{501.818}{653.696}{zeroes}%
\TextAt{13}{79.226}{636.65}{\Character{12}}%
\TextAt{17}{85.835}{634.857}{1}%
\TextAt{13}{90.567}{636.65}{\Character{59}}%
\TextAt{13}{95.791}{636.65}{\Character{58}}%
\TextAt{13}{101.028}{636.65}{\Character{58}}%
\TextAt{13}{106.252}{636.65}{\Character{58}}%
\TextAt{13}{111.489}{636.65}{\Character{59}}%
\TextAt{13}{116.725}{636.65}{\Character{12}}%
\TextAt{21}{123.396}{634.857}{n}%
\TextAt{8}{132.171}{636.65}{of}%
\TextAt{13}{144.742}{636.65}{p}%
\TextAt{9}{150.617}{636.65}{\Character{40}}%
\TextAt{13}{155.169}{636.65}{x}%
\TextAt{9}{161.821}{636.65}{\Character{41}}%
\TextAt{8}{169.512}{636.65}{are}%
\TextAt{8}{188.264}{636.65}{integral}%
\TextAt{8}{230.113}{636.65}{over}%
\TextAt{13}{254.4}{636.65}{A}%
\TextAt{8}{263.175}{636.65}{,}%
\TextAt{8}{269.717}{636.65}{hence}%
\TextAt{8}{301.464}{636.65}{the}%
\TextAt{8}{320.861}{636.65}{same}%
\TextAt{8}{349.428}{636.65}{holds}%
\TextAt{8}{379.298}{636.65}{for}%
\TextAt{8}{396.424}{636.65}{all}%
\TextAt{8}{411.924}{636.65}{the}%
\TextAt{8}{431.321}{636.65}{coefficients,}%
\TextAt{8}{493.204}{636.65}{in}%
\TextAt{8}{506.097}{636.65}{other}%
\TextAt{8}{79.226}{619.604}{words}%
\TextAt{13}{112.784}{619.604}{p}%
\TextAt{9}{118.659}{619.604}{\Character{40}}%
\TextAt{13}{123.212}{619.604}{x}%
\TextAt{9}{129.864}{619.604}{\Character{41}}%
\NativeMathAt{137.737}{619.604}{11.955}{\in}%
\TextAt{13}{149.028}{619.604}{A}%
\TextAt{9}{157.803}{619.604}{\Character{91}}%
\TextAt{13}{161.055}{619.604}{x}%
\TextAt{9}{167.707}{619.604}{\Character{93}}%
\TextAt{8}{170.959}{619.604}{.}%
\TextAt{8}{97.159}{599.57}{The}%
\TextAt{12}{121.196}{599.57}{trace}%
\TextAt{8}{154.463}{599.57}{and}%
\TextAt{8}{177.209}{599.57}{the}%
\TextAt{12}{197.349}{599.57}{norm}%
\TextAt{8}{232.146}{599.57}{in}%
\TextAt{8}{245.782}{599.57}{the}%
\TextAt{8}{265.922}{599.57}{field}%
\TextAt{8}{291.262}{599.57}{extension}%
\TextAt{13}{343.012}{599.57}{L}%
\TextAt{16}{354.297}{599.57}{\Character{106}}%
\TextAt{13}{360.939}{599.57}{K}%
\TextAt{8}{375.637}{599.57}{furnish}%
\TextAt{8}{415.025}{599.57}{important}%
\TextAt{8}{470.301}{599.57}{tools}%
\TextAt{8}{498.642}{599.57}{for}%
\TextAt{8}{516.51}{599.57}{the}%
\TextAt{8}{79.226}{582.524}{study}%
\TextAt{8}{111.481}{582.524}{of}%
\TextAt{8}{124.811}{582.524}{the}%
\TextAt{8}{144.968}{582.524}{integral}%
\TextAt{8}{187.576}{582.524}{elements}%
\TextAt{8}{235.432}{582.524}{in}%
\TextAt{13}{249.085}{582.524}{L}%
\TextAt{8}{257.05}{582.524}{.}%
\TextAt{8}{265.498}{582.524}{We}%
\TextAt{8}{285.643}{582.524}{recall}%
\TextAt{8}{316.857}{582.524}{the}%
\TextAt{12}{79.226}{556.756}{Definition}%
\TextAt{12}{142.464}{556.756}{(2.5).}%
\TextAt{8}{178.17}{556.756}{The}%
\TextAt{8}{202.614}{556.756}{trace}%
\TextAt{8}{232.271}{556.756}{and}%
\TextAt{8}{255.424}{556.756}{norm}%
\TextAt{8}{286.384}{556.756}{of}%
\TextAt{8}{300.104}{556.756}{an}%
\TextAt{8}{316.754}{556.756}{element}%
\TextAt{13}{360.386}{556.756}{x}%
\NativeMathAt{371.025}{556.756}{11.955}{\in}%
\TextAt{13}{382.981}{556.756}{L}%
\TextAt{8}{395.233}{556.756}{are}%
\TextAt{8}{415.135}{556.756}{defined}%
\TextAt{8}{455.838}{556.756}{to}%
\TextAt{8}{470.539}{556.756}{be}%
\TextAt{8}{486.854}{556.756}{the}%
\TextAt{8}{507.401}{556.756}{trace}%
\TextAt{8}{79.226}{539.71}{and}%
\TextAt{8}{101.989}{539.71}{determinant,}%
\TextAt{8}{171.257}{539.71}{respectively,}%
\TextAt{8}{237.321}{539.71}{of}%
\TextAt{8}{250.651}{539.71}{the}%
\TextAt{8}{270.808}{539.71}{endomorphism}%
\TextAt{13}{228.533}{511.565}{T}%
\TextAt{21}{235.394}{509.772}{x}%
\TextAt{9}{243.98}{511.565}{\Character{58}}%
\TextAt{13}{250.552}{511.565}{L}%
\NativeArrowAt{261.838}{511.565}{11.955}{19.26}{\longrightarrow}%
\TextAt{13}{284.42}{511.565}{L\Character{59}}%
\TextAt{13}{321.517}{511.565}{T}%
\TextAt{21}{328.4}{509.772}{x}%
\TextAt{9}{333.665}{511.565}{\Character{40}}%
\TextAt{13}{338.217}{511.565}{\Character{11}}%
\TextAt{9}{345.739}{511.565}{\Character{41}}%
\TextAt{9}{353.594}{511.565}{\Character{61}}%
\TextAt{13}{366.037}{511.565}{x\Character{11}\Character{59}}%
\TextAt{8}{79.226}{483.421}{of}%
\TextAt{8}{92.556}{483.421}{the}%
\TextAt{13}{112.713}{483.421}{K}%
\TextAt{8}{123.53}{483.421}{-vector}%
\TextAt{8}{162.552}{483.421}{space}%
\TextAt{13}{193.827}{483.421}{L}%
\TextAt{8}{201.791}{483.421}{:}%
\TextAt{8}{194.923}{455.276}{Tr}%
\TextAt{21}{207.931}{453.416}{L}%
\TextAt{3}{213.69}{453.416}{\Character{106}}%
\TextAt{21}{216.042}{453.416}{K}%
\TextAt{9}{224.245}{455.276}{\Character{40}}%
\TextAt{13}{228.797}{455.276}{x}%
\TextAt{9}{235.449}{455.276}{\Character{41}}%
\TextAt{9}{243.304}{455.276}{\Character{61}}%
\TextAt{8}{255.748}{455.276}{Tr}%
\TextAt{9}{268.755}{455.276}{\Character{40}}%
\TextAt{13}{273.307}{455.276}{T}%
\TextAt{21}{280.168}{453.483}{x}%
\TextAt{9}{285.433}{455.276}{\Character{41}}%
\TextAt{13}{289.986}{455.276}{\Character{59}}%
\TextAt{13}{319.121}{455.276}{N}%
\TextAt{21}{328.514}{453.416}{L}%
\TextAt{3}{334.273}{453.416}{\Character{106}}%
\TextAt{21}{336.625}{453.416}{K}%
\TextAt{9}{344.828}{455.276}{\Character{40}}%
\TextAt{13}{349.38}{455.276}{x}%
\TextAt{9}{356.032}{455.276}{\Character{41}}%
\TextAt{9}{363.887}{455.276}{\Character{61}}%
\TextAt{8}{376.331}{455.276}{det}%
\TextAt{9}{392.59}{455.276}{\Character{40}}%
\TextAt{13}{397.142}{455.276}{T}%
\TextAt{21}{404.003}{453.483}{x}%
\TextAt{9}{409.268}{455.276}{\Character{41}}%
\TextAt{13}{413.821}{455.276}{\Character{58}}%
\TextAt{8}{97.159}{418.41}{In}%
\TextAt{8}{111.78}{418.41}{the}%
\TextAt{8}{131.937}{418.41}{characteristic}%
\TextAt{8}{203.871}{418.41}{polynomial}%
\TextAt{13}{156.367}{390.265}{f}%
\TextAt{21}{162.137}{388.472}{x}%
\TextAt{9}{167.402}{390.265}{\Character{40}}%
\TextAt{13}{171.954}{390.265}{t}%
\TextAt{9}{176.181}{390.265}{\Character{41}}%
\TextAt{9}{184.036}{390.265}{\Character{61}}%
\TextAt{8}{196.48}{390.265}{det}%
\TextAt{9}{212.739}{390.265}{\Character{40}}%
\TextAt{13}{217.291}{390.265}{t}%
\TextAt{8}{225.504}{390.265}{id}%
\TextAt{16}{237.251}{390.265}{\Character{0}}%
\TextAt{13}{246.55}{390.265}{T}%
\TextAt{21}{253.411}{388.472}{x}%
\TextAt{9}{258.676}{390.265}{\Character{41}}%
\TextAt{9}{266.531}{390.265}{\Character{61}}%
\TextAt{13}{278.975}{390.265}{t}%
\TextAt{21}{283.202}{395.201}{n}%
\TextAt{16}{291.495}{390.265}{\Character{0}}%
\TextAt{13}{303.45}{390.265}{a}%
\TextAt{17}{309.595}{388.472}{1}%
\TextAt{13}{314.327}{390.265}{t}%
\TextAt{21}{318.554}{395.201}{n}%
\TextAt{3}{323.693}{395.201}{\Character{0}}%
\TextAt{17}{330.279}{395.201}{1}%
\TextAt{9}{337.668}{390.265}{\Character{43}}%
\TextAt{16}{349.429}{390.265}{\Character{1}}%
\TextAt{16}{354.725}{390.265}{\Character{1}}%
\TextAt{16}{360.033}{390.265}{\Character{1}}%
\TextAt{9}{366.034}{390.265}{\Character{43}}%
\TextAt{9}{377.786}{390.265}{\Character{40}}%
\TextAt{16}{382.347}{390.265}{\Character{0}}%
\TextAt{9}{391.646}{390.265}{1\Character{41}}%
\TextAt{21}{402.051}{395.201}{n}%
\TextAt{13}{407.688}{390.265}{a}%
\TextAt{21}{413.832}{388.472}{n}%
\NativeMathAt{422.79}{390.265}{11.955}{\in}%
\TextAt{13}{434.081}{390.265}{K}%
\TextAt{9}{444.898}{390.265}{\Character{91}}%
\TextAt{13}{448.15}{390.265}{t}%
\TextAt{9}{452.377}{390.265}{\Character{93}}%
\TextAt{8}{79.226}{362.12}{of}%
\TextAt{13}{92.556}{362.12}{T}%
\TextAt{21}{99.417}{360.327}{x}%
\TextAt{8}{104.682}{362.12}{,}%
\TextAt{13}{111.831}{362.12}{n}%
\TextAt{9}{122.14}{362.12}{\Character{61}}%
\TextAt{9}{134.55}{362.12}{\Character{91}}%
\TextAt{13}{137.817}{362.12}{L}%
\TextAt{9}{149.102}{362.12}{\Character{58}}%
\TextAt{13}{155.675}{362.12}{K}%
\TextAt{9}{166.492}{362.12}{\Character{93}}%
\TextAt{8}{169.744}{362.12}{,}%
\TextAt{8}{176.893}{362.12}{we}%
\TextAt{8}{194.121}{362.12}{recognize}%
\TextAt{8}{244.847}{362.12}{the}%
\TextAt{8}{265.003}{362.12}{trace}%
\TextAt{8}{294.269}{362.12}{and}%
\TextAt{8}{317.032}{362.12}{the}%
\TextAt{8}{337.188}{362.12}{norm}%
\TextAt{8}{367.758}{362.12}{as}%
\TextAt{13}{212.647}{333.976}{a}%
\TextAt{17}{218.792}{332.182}{1}%
\TextAt{9}{226.845}{333.976}{\Character{61}}%
\TextAt{8}{239.271}{333.976}{Tr}%
\TextAt{21}{252.278}{332.116}{L}%
\TextAt{3}{258.037}{332.116}{\Character{106}}%
\TextAt{21}{260.389}{332.116}{K}%
\TextAt{9}{268.592}{333.976}{\Character{40}}%
\TextAt{13}{273.144}{333.976}{x}%
\TextAt{9}{279.797}{333.976}{\Character{41}}%
\TextAt{8}{296.304}{333.976}{and}%
\TextAt{13}{327.124}{333.976}{a}%
\TextAt{21}{333.269}{332.182}{n}%
\TextAt{9}{342.227}{333.976}{\Character{61}}%
\TextAt{13}{354.652}{333.976}{N}%
\TextAt{21}{364.026}{332.116}{L}%
\TextAt{3}{369.785}{332.116}{\Character{106}}%
\TextAt{21}{372.137}{332.116}{K}%
\TextAt{9}{380.34}{333.976}{\Character{40}}%
\TextAt{13}{384.892}{333.976}{x}%
\TextAt{9}{391.544}{333.976}{\Character{41}}%
\TextAt{13}{396.097}{333.976}{\Character{58}}%
\TextAt{8}{79.226}{305.831}{Since}%
\TextAt{13}{109.784}{305.831}{T}%
\TextAt{21}{116.645}{304.038}{x}%
\TextAt{17}{121.411}{304.038}{\Character{43}}%
\TextAt{21}{127.998}{304.038}{y}%
\TextAt{9}{136.289}{305.831}{\Character{61}}%
\TextAt{13}{148.714}{305.831}{T}%
\TextAt{21}{155.575}{304.038}{x}%
\TextAt{9}{163.497}{305.831}{\Character{43}}%
\TextAt{13}{175.258}{305.831}{T}%
\TextAt{21}{182.119}{304.038}{y}%
\TextAt{8}{190.986}{305.831}{and}%
\TextAt{13}{213.749}{305.831}{T}%
\TextAt{21}{220.61}{304.038}{xy}%
\TextAt{9}{233.668}{305.831}{\Character{61}}%
\TextAt{13}{246.093}{305.831}{T}%
\TextAt{21}{252.954}{304.038}{x}%
\TextAt{16}{260.876}{305.831}{\Character{14}}%
\TextAt{13}{269.51}{305.831}{T}%
\TextAt{21}{276.371}{304.038}{y}%
\TextAt{8}{281.341}{305.831}{,}%
\TextAt{8}{288.49}{305.831}{we}%
\TextAt{8}{305.718}{305.831}{obtain}%
\TextAt{8}{342.145}{305.831}{homomorphisms}%
\TextAt{8}{196.169}{277.686}{Tr}%
\TextAt{21}{209.176}{275.827}{L}%
\TextAt{3}{214.935}{275.827}{\Character{106}}%
\TextAt{21}{217.288}{275.827}{K}%
\TextAt{9}{228.811}{277.686}{\Character{58}}%
\TextAt{13}{235.384}{277.686}{L}%
\NativeArrowAt{246.669}{277.686}{11.955}{19.26}{\longrightarrow}%
\TextAt{13}{269.251}{277.686}{K}%
\TextAt{8}{292.024}{277.686}{and}%
\TextAt{13}{322.844}{277.686}{N}%
\TextAt{21}{332.218}{275.827}{L}%
\TextAt{3}{337.977}{275.827}{\Character{106}}%
\TextAt{21}{340.329}{275.827}{K}%
\TextAt{9}{351.853}{277.686}{\Character{58}}%
\TextAt{13}{358.425}{277.686}{L}%
\TextAt{3}{366.39}{282.623}{\Character{3}}%
\NativeArrowAt{374.443}{277.686}{11.955}{19.26}{\longrightarrow}%
\TextAt{13}{397.025}{277.686}{K}%
\TextAt{3}{407.843}{282.623}{\Character{3}}%
\TextAt{13}{412.575}{277.686}{\Character{58}}%
\TextAt{8}{79.226}{249.542}{In}%
\TextAt{8}{92.962}{249.542}{the}%
\TextAt{8}{112.232}{249.542}{case}%
\TextAt{8}{136.118}{249.542}{where}%
\TextAt{8}{169.042}{249.542}{the}%
\TextAt{8}{188.312}{249.542}{extension}%
\TextAt{13}{239.193}{249.542}{L}%
\TextAt{16}{250.478}{249.542}{\Character{106}}%
\TextAt{13}{257.12}{249.542}{K}%
\TextAt{8}{270.949}{249.542}{is}%
\TextAt{8}{281.827}{249.542}{separable,}%
\TextAt{8}{335.814}{249.542}{the}%
\TextAt{8}{355.085}{249.542}{trace}%
\TextAt{8}{383.465}{249.542}{and}%
\TextAt{8}{405.342}{249.542}{norm}%
\TextAt{8}{435.026}{249.542}{admit}%
\TextAt{8}{467.961}{249.542}{the}%
\TextAt{8}{487.232}{249.542}{following}%
\TextAt{8}{79.226}{232.496}{Galois-theoretic}%
\TextAt{8}{163.917}{232.496}{interpretation.}%
\TextAt{12}{79.226}{206.728}{Proposition}%
\TextAt{12}{153.028}{206.728}{(2.6).}%
\TextAt{8}{189.292}{206.728}{If}%
\TextAt{13}{201.933}{206.728}{L}%
\TextAt{16}{214.835}{206.728}{\Character{106}}%
\TextAt{13}{223.093}{206.728}{K}%
\TextAt{8}{238.756}{206.728}{is}%
\TextAt{8}{251.468}{206.728}{a}%
\TextAt{8}{262.172}{206.728}{separable}%
\TextAt{8}{314.564}{206.728}{extension}%
\TextAt{8}{367.278}{206.728}{and}%
\TextAt{13}{390.989}{206.728}{\Character{27}}%
\TextAt{9}{403.009}{206.728}{\Character{58}}%
\TextAt{13}{411.198}{206.728}{L}%
\NativeMathAt{424.099}{206.728}{11.955}{\rightarrow}%
\TextAt{9}{444.124}{209.75}{\Character{22}}%
\TextAt{13}{440.991}{206.728}{K}%
\TextAt{8}{456.655}{206.728}{varies}%
\TextAt{8}{490.516}{206.728}{over}%
\TextAt{8}{516.51}{206.728}{the}%
\TextAt{8}{79.226}{189.682}{different}%
\TextAt{13}{125.397}{189.682}{K}%
\TextAt{8}{136.215}{189.682}{-embeddings}%
\TextAt{8}{203.905}{189.682}{of}%
\TextAt{13}{217.235}{189.682}{L}%
\TextAt{8}{229.097}{189.682}{into}%
\TextAt{8}{252.84}{189.682}{an}%
\TextAt{8}{269.099}{189.682}{algebraic}%
\TextAt{8}{318.533}{189.682}{closure}%
\TextAt{9}{360.748}{192.704}{\Character{22}}%
\TextAt{13}{357.615}{189.682}{K}%
\TextAt{8}{372.33}{189.682}{of}%
\TextAt{13}{385.66}{189.682}{K}%
\TextAt{8}{396.478}{189.682}{,}%
\TextAt{8}{403.627}{189.682}{then}%
\TextAt{8}{430.287}{189.682}{we}%
\TextAt{8}{447.514}{189.682}{have}%
\TextAt{8}{89.579}{169.647}{(i)}%
\TextAt{13}{107.918}{169.647}{f}%
\TextAt{21}{113.688}{167.854}{x}%
\TextAt{9}{118.953}{169.647}{\Character{40}}%
\TextAt{13}{123.506}{169.647}{t}%
\TextAt{9}{127.733}{169.647}{\Character{41}}%
\TextAt{9}{135.588}{169.647}{\Character{61}}%
\TextAt{20}{148.032}{178.614}{\Character{81}}%
\TextAt{21}{159.323}{166.16}{\Character{27}}%
\TextAt{9}{164.946}{169.647}{\Character{40}}%
\TextAt{13}{169.498}{169.647}{t}%
\TextAt{16}{176.382}{169.647}{\Character{0}}%
\TextAt{13}{188.337}{169.647}{\Character{27}}%
\TextAt{13}{195.403}{169.647}{x}%
\TextAt{9}{202.072}{169.647}{\Character{41}}%
\TextAt{8}{206.624}{169.647}{,}%
\TextAt{8}{86.327}{152.601}{(ii)}%
\TextAt{8}{107.918}{152.601}{Tr}%
\TextAt{21}{120.926}{150.742}{L}%
\TextAt{3}{126.685}{150.742}{\Character{106}}%
\TextAt{21}{129.037}{150.742}{K}%
\TextAt{9}{137.24}{152.601}{\Character{40}}%
\TextAt{13}{141.792}{152.601}{x}%
\TextAt{9}{148.444}{152.601}{\Character{41}}%
\TextAt{9}{156.299}{152.601}{\Character{61}}%
\TextAt{20}{168.743}{161.568}{\Character{80}}%
\TextAt{21}{181.362}{149.114}{\Character{27}}%
\TextAt{13}{188.978}{152.601}{\Character{27}}%
\TextAt{13}{196.044}{152.601}{x}%
\TextAt{8}{202.713}{152.601}{,}%
\TextAt{8}{83.076}{135.555}{(iii)}%
\TextAt{13}{107.918}{135.555}{N}%
\TextAt{21}{117.292}{133.696}{L}%
\TextAt{3}{123.051}{133.696}{\Character{106}}%
\TextAt{21}{125.404}{133.696}{K}%
\TextAt{9}{133.606}{135.555}{\Character{40}}%
\TextAt{13}{138.159}{135.555}{x}%
\TextAt{9}{144.811}{135.555}{\Character{41}}%
\TextAt{9}{152.666}{135.555}{\Character{61}}%
\TextAt{20}{165.109}{144.522}{\Character{81}}%
\TextAt{21}{176.4}{132.068}{\Character{27}}%
\TextAt{13}{184.016}{135.555}{\Character{27}}%
\TextAt{13}{191.082}{135.555}{x}%
\TextAt{8}{197.751}{135.555}{.}%
\TextAt{19}{79.226}{101.066}{Proof.}%
\TextAt{8}{116}{101.066}{The}%
\TextAt{8}{140.054}{101.066}{characteristic}%
\TextAt{8}{211.988}{101.066}{polynomial}%
\TextAt{13}{272.481}{101.066}{f}%
\TextAt{21}{278.251}{99.273}{x}%
\TextAt{9}{283.516}{101.066}{\Character{40}}%
\TextAt{13}{288.069}{101.066}{t}%
\TextAt{9}{292.296}{101.066}{\Character{41}}%
\TextAt{8}{300.746}{101.066}{is}%
\TextAt{8}{312.51}{101.066}{a}%
\TextAt{8}{322.265}{101.066}{power}%
\TextAt{13}{220.318}{72.921}{f}%
\TextAt{21}{226.088}{71.128}{x}%
\TextAt{9}{231.353}{72.921}{\Character{40}}%
\TextAt{13}{235.905}{72.921}{t}%
\TextAt{9}{240.132}{72.921}{\Character{41}}%
\TextAt{9}{247.987}{72.921}{\Character{61}}%
\TextAt{13}{260.431}{72.921}{p}%
\TextAt{21}{266.306}{71.128}{x}%
\TextAt{9}{271.571}{72.921}{\Character{40}}%
\TextAt{13}{276.124}{72.921}{t}%
\TextAt{9}{280.351}{72.921}{\Character{41}}%
\TextAt{21}{284.903}{77.858}{d}%
\TextAt{13}{289.759}{72.921}{\Character{59}}%
\TextAt{13}{318.894}{72.921}{d}%
\TextAt{9}{328.317}{72.921}{\Character{61}}%
\TextAt{9}{340.726}{72.921}{\Character{91}}%
\TextAt{13}{343.994}{72.921}{L}%
\TextAt{9}{355.279}{72.921}{\Character{58}}%
\TextAt{13}{361.852}{72.921}{K}%
\TextAt{9}{372.669}{72.921}{\Character{40}}%
\TextAt{13}{377.222}{72.921}{x}%
\TextAt{9}{383.874}{72.921}{\Character{41}\Character{93}}%
\TextAt{8}{303.069}{43.033}{8}%
\Destination{79.23}{184.245}{nav.Item.4}%
\Destination{79.23}{166.165}{nav.Item.5}%
\Destination{79.23}{148.745}{nav.Item.6}%
\Destination{78.23}{809.815}{nav.page.8}%
\Destination{79.36}{614.555}{nav.print.9}%
\Destination{79.36}{570.755}{nav.number.1.2.5}%
\Destination{79.36}{220.755}{nav.number.1.2.6}%
}

% Page 21 (body page 9)
\Page{%
\TextAt{2}{79.226}{796.943}{I.}%
\TextAt{2}{89.83}{796.943}{Algebraic}%
\TextAt{2}{138.648}{796.943}{Integers}%
\Vector{q 0 G 0 g 0.299 w 7.226 18.729 m 460.769 18.729 l S Q}%
\TextAt{8}{79.226}{761.124}{of}%
\TextAt{8}{92.556}{761.124}{the}%
\TextAt{8}{112.713}{761.124}{minimal}%
\TextAt{8}{158.238}{761.124}{polynomial}%
\TextAt{13}{173.799}{737.189}{p}%
\TextAt{21}{179.674}{735.395}{x}%
\TextAt{9}{184.939}{737.189}{\Character{40}}%
\TextAt{13}{189.491}{737.189}{t}%
\TextAt{9}{193.718}{737.189}{\Character{41}}%
\TextAt{9}{201.573}{737.189}{\Character{61}}%
\TextAt{13}{214.017}{737.189}{t}%
\TextAt{21}{218.244}{742.125}{m}%
\TextAt{9}{228.89}{737.189}{\Character{43}}%
\TextAt{13}{240.651}{737.189}{c}%
\TextAt{17}{245.689}{735.395}{1}%
\TextAt{13}{250.421}{737.189}{t}%
\TextAt{21}{254.648}{742.125}{m}%
\TextAt{3}{262.139}{742.125}{\Character{0}}%
\TextAt{17}{268.725}{742.125}{1}%
\TextAt{9}{276.114}{737.189}{\Character{43}}%
\TextAt{16}{287.876}{737.189}{\Character{1}}%
\TextAt{16}{293.172}{737.189}{\Character{1}}%
\TextAt{16}{298.48}{737.189}{\Character{1}}%
\TextAt{9}{304.48}{737.189}{\Character{43}}%
\TextAt{13}{316.241}{737.189}{c}%
\TextAt{21}{321.279}{735.395}{m}%
\TextAt{13}{329.268}{737.189}{\Character{59}}%
\TextAt{13}{358.403}{737.189}{m}%
\TextAt{9}{371.983}{737.189}{\Character{61}}%
\TextAt{9}{384.393}{737.189}{\Character{91}}%
\TextAt{13}{387.66}{737.189}{K}%
\TextAt{9}{398.477}{737.189}{\Character{40}}%
\TextAt{13}{403.03}{737.189}{x}%
\TextAt{9}{409.682}{737.189}{\Character{41}}%
\TextAt{9}{417.537}{737.189}{\Character{58}}%
\TextAt{13}{424.128}{737.189}{K}%
\TextAt{9}{434.945}{737.189}{\Character{93}}%
\TextAt{8}{79.226}{713.254}{of}%
\TextAt{13}{92.136}{713.254}{x}%
\TextAt{8}{98.788}{713.254}{.}%
\TextAt{8}{107.096}{713.254}{In}%
\TextAt{8}{121.297}{713.254}{fact,}%
\TextAt{9}{147.299}{713.254}{1}%
\TextAt{13}{153.152}{713.254}{\Character{59}}%
\TextAt{13}{158.376}{713.254}{x\Character{59}}%
\TextAt{13}{170.26}{713.254}{\Character{58}}%
\TextAt{13}{175.484}{713.254}{\Character{58}}%
\TextAt{13}{180.721}{713.254}{\Character{58}}%
\TextAt{13}{185.957}{713.254}{\Character{59}}%
\TextAt{13}{191.181}{713.254}{x}%
\TextAt{21}{197.921}{717.592}{m}%
\TextAt{3}{205.411}{717.592}{\Character{0}}%
\TextAt{17}{211.998}{717.592}{1}%
\TextAt{8}{220.207}{713.254}{is}%
\TextAt{8}{231.551}{713.254}{a}%
\TextAt{8}{240.886}{713.254}{basis}%
\TextAt{8}{269.206}{713.254}{of}%
\TextAt{13}{282.116}{713.254}{K}%
\TextAt{9}{292.934}{713.254}{\Character{40}}%
\TextAt{13}{297.486}{713.254}{x}%
\TextAt{9}{304.138}{713.254}{\Character{41}}%
\TextAt{16}{312.012}{713.254}{\Character{106}}%
\TextAt{13}{318.653}{713.254}{K}%
\TextAt{8}{329.471}{713.254}{,}%
\TextAt{8}{336.284}{713.254}{and}%
\TextAt{8}{358.627}{713.254}{if}%
\TextAt{13}{368.93}{713.254}{\Character{11}}%
\TextAt{17}{376.376}{711.461}{1}%
\TextAt{13}{381.108}{713.254}{\Character{59}}%
\TextAt{13}{386.332}{713.254}{\Character{58}}%
\TextAt{13}{391.569}{713.254}{\Character{58}}%
\TextAt{13}{396.805}{713.254}{\Character{58}}%
\TextAt{13}{402.03}{713.254}{\Character{59}}%
\TextAt{13}{407.266}{713.254}{\Character{11}}%
\TextAt{21}{414.775}{711.461}{d}%
\TextAt{8}{423.107}{713.254}{is}%
\TextAt{8}{434.451}{713.254}{a}%
\TextAt{8}{443.786}{713.254}{basis}%
\TextAt{8}{472.106}{713.254}{of}%
\TextAt{13}{485.016}{713.254}{L}%
\TextAt{16}{496.302}{713.254}{\Character{106}}%
\TextAt{13}{502.943}{713.254}{K}%
\TextAt{9}{513.761}{713.254}{\Character{40}}%
\TextAt{13}{518.313}{713.254}{x}%
\TextAt{9}{524.965}{713.254}{\Character{41}}%
\TextAt{8}{529.518}{713.254}{,}%
\TextAt{8}{79.226}{696.208}{then}%
\TextAt{13}{188.693}{679.162}{\Character{11}}%
\TextAt{17}{196.139}{677.369}{1}%
\TextAt{13}{200.871}{679.162}{\Character{59}}%
\TextAt{13}{206.095}{679.162}{\Character{11}}%
\TextAt{17}{213.561}{677.369}{1}%
\TextAt{13}{218.293}{679.162}{x\Character{59}}%
\TextAt{13}{230.165}{679.162}{\Character{58}}%
\TextAt{13}{235.401}{679.162}{\Character{58}}%
\TextAt{13}{240.637}{679.162}{\Character{58}}%
\TextAt{13}{245.862}{679.162}{\Character{59}}%
\TextAt{13}{251.098}{679.162}{\Character{11}}%
\TextAt{17}{258.612}{677.369}{1}%
\TextAt{13}{263.344}{679.162}{x}%
\TextAt{21}{269.996}{684.098}{m}%
\TextAt{3}{277.487}{684.098}{\Character{0}}%
\TextAt{17}{284.073}{684.098}{1}%
\TextAt{9}{288.805}{679.162}{\Character{59}}%
\TextAt{13}{297.947}{679.162}{\Character{58}}%
\TextAt{13}{303.171}{679.162}{\Character{58}}%
\TextAt{13}{308.408}{679.162}{\Character{58}}%
\TextAt{9}{313.679}{679.162}{\Character{59}}%
\TextAt{13}{322.821}{679.162}{\Character{11}}%
\TextAt{21}{330.266}{677.369}{d}%
\TextAt{13}{335.122}{679.162}{\Character{59}}%
\TextAt{13}{340.346}{679.162}{\Character{11}}%
\TextAt{21}{347.812}{677.369}{d}%
\TextAt{13}{352.667}{679.162}{x\Character{59}}%
\TextAt{13}{364.538}{679.162}{\Character{58}}%
\TextAt{13}{369.775}{679.162}{\Character{58}}%
\TextAt{13}{375.011}{679.162}{\Character{58}}%
\TextAt{13}{380.236}{679.162}{\Character{59}}%
\TextAt{13}{385.472}{679.162}{\Character{11}}%
\TextAt{21}{392.986}{677.369}{d}%
\TextAt{13}{397.841}{679.162}{x}%
\TextAt{21}{404.493}{684.098}{m}%
\TextAt{3}{411.984}{684.098}{\Character{0}}%
\TextAt{17}{418.57}{684.098}{1}%
\TextAt{8}{79.226}{657.811}{is}%
\TextAt{8}{91.155}{657.811}{a}%
\TextAt{8}{101.076}{657.811}{basis}%
\TextAt{8}{129.982}{657.811}{of}%
\TextAt{13}{143.477}{657.811}{L}%
\TextAt{16}{155.044}{657.811}{\Character{106}}%
\TextAt{13}{161.967}{657.811}{K}%
\TextAt{8}{172.785}{657.811}{.}%
\TextAt{8}{181.729}{657.811}{The}%
\TextAt{8}{205.948}{657.811}{matrix}%
\TextAt{8}{244.167}{657.811}{of}%
\TextAt{8}{257.662}{657.811}{the}%
\TextAt{8}{277.984}{657.811}{linear}%
\TextAt{8}{310.667}{657.811}{transformation}%
\TextAt{13}{390.585}{657.811}{T}%
\TextAt{21}{397.446}{656.018}{x}%
\TextAt{9}{406.314}{657.811}{\Character{58}}%
\TextAt{13}{413.168}{657.811}{y}%
\NativeArrowAt{422.907}{657.811}{11.955}{11.955}{\mapsto}%
\TextAt{13}{438.465}{657.811}{xy}%
\TextAt{8}{455.316}{657.811}{with}%
\TextAt{8}{482.141}{657.811}{respect}%
\TextAt{8}{522.357}{657.811}{to}%
\TextAt{8}{79.226}{640.765}{this}%
\TextAt{8}{102.049}{640.765}{basis}%
\TextAt{8}{130.789}{640.765}{has}%
\TextAt{8}{151.663}{640.765}{obviously}%
\TextAt{8}{203.44}{640.765}{only}%
\TextAt{8}{229.132}{640.765}{blocks}%
\TextAt{8}{264.639}{640.765}{along}%
\TextAt{8}{295.866}{640.765}{the}%
\TextAt{8}{316.022}{640.765}{diagonal,}%
\TextAt{8}{366.114}{640.765}{each}%
\TextAt{8}{392.452}{640.765}{of}%
\TextAt{8}{405.782}{640.765}{them}%
\TextAt{8}{435.693}{640.765}{equal}%
\TextAt{8}{466.586}{640.765}{to}%
\TextAt{20}{240.663}{628.788}{\Character{48}}%
\TextAt{20}{240.663}{607.746}{\Character{66}}%
\TextAt{20}{240.663}{600.573}{\Character{66}}%
\TextAt{20}{240.663}{593.4}{\Character{66}}%
\TextAt{20}{240.663}{586.227}{\Character{66}}%
\TextAt{20}{240.663}{579.054}{\Character{66}}%
\TextAt{20}{240.663}{571.881}{\Character{66}}%
\TextAt{20}{240.663}{564.229}{\Character{64}}%
\TextAt{9}{251.124}{618.497}{0}%
\TextAt{9}{266.929}{618.497}{0}%
\TextAt{16}{282.756}{618.497}{\Character{1}}%
\TextAt{16}{288.052}{618.497}{\Character{1}}%
\TextAt{16}{293.36}{618.497}{\Character{1}}%
\TextAt{9}{306.666}{618.497}{0}%
\TextAt{16}{327.892}{618.497}{\Character{0}}%
\TextAt{13}{337.19}{618.497}{c}%
\TextAt{21}{342.228}{616.704}{m}%
\TextAt{9}{251.124}{601.451}{1}%
\TextAt{9}{266.929}{601.451}{0}%
\TextAt{16}{282.755}{601.451}{\Character{1}}%
\TextAt{16}{288.051}{601.451}{\Character{1}}%
\TextAt{16}{293.359}{601.451}{\Character{1}}%
\TextAt{9}{306.666}{601.451}{0}%
\TextAt{16}{322.481}{601.451}{\Character{0}}%
\TextAt{13}{331.78}{601.451}{c}%
\TextAt{21}{336.818}{599.658}{m}%
\TextAt{3}{344.308}{599.658}{\Character{0}}%
\TextAt{17}{350.895}{599.658}{1}%
\TextAt{9}{251.124}{584.405}{0}%
\TextAt{9}{266.929}{584.405}{1}%
\TextAt{16}{282.755}{584.405}{\Character{1}}%
\TextAt{16}{288.051}{584.405}{\Character{1}}%
\TextAt{16}{293.359}{584.405}{\Character{1}}%
\TextAt{9}{306.666}{584.405}{0}%
\TextAt{16}{322.481}{584.405}{\Character{0}}%
\TextAt{13}{331.78}{584.405}{c}%
\TextAt{21}{336.818}{582.612}{m}%
\TextAt{3}{344.308}{582.612}{\Character{0}}%
\TextAt{17}{350.895}{582.612}{2}%
\TextAt{8}{252.425}{572.154}{.}%
\TextAt{8}{252.425}{568.169}{.}%
\TextAt{8}{252.425}{564.184}{.}%
\TextAt{8}{268.24}{572.154}{.}%
\TextAt{8}{268.24}{568.169}{.}%
\TextAt{8}{268.24}{564.184}{.}%
\TextAt{8}{283.523}{571.158}{.}%
\TextAt{8}{288.103}{568.169}{.}%
\TextAt{8}{292.684}{565.18}{.}%
\TextAt{8}{307.966}{572.154}{.}%
\TextAt{8}{307.966}{568.169}{.}%
\TextAt{8}{307.966}{564.184}{.}%
\TextAt{8}{337.428}{572.154}{.}%
\TextAt{8}{337.428}{568.169}{.}%
\TextAt{8}{337.428}{564.184}{.}%
\TextAt{9}{251.124}{547.138}{0}%
\TextAt{9}{266.929}{547.138}{0}%
\TextAt{16}{282.756}{547.138}{\Character{1}}%
\TextAt{16}{288.052}{547.138}{\Character{1}}%
\TextAt{16}{293.36}{547.138}{\Character{1}}%
\TextAt{9}{306.666}{547.138}{1}%
\TextAt{16}{329.52}{547.138}{\Character{0}}%
\TextAt{13}{338.818}{547.138}{c}%
\TextAt{17}{343.856}{545.345}{1}%
\TextAt{20}{355.627}{628.788}{\Character{49}}%
\TextAt{20}{355.627}{607.746}{\Character{67}}%
\TextAt{20}{355.627}{600.573}{\Character{67}}%
\TextAt{20}{355.627}{593.4}{\Character{67}}%
\TextAt{20}{355.627}{586.227}{\Character{67}}%
\TextAt{20}{355.627}{579.054}{\Character{67}}%
\TextAt{20}{355.627}{571.881}{\Character{67}}%
\TextAt{20}{355.627}{564.229}{\Character{65}}%
\TextAt{13}{368.081}{583.238}{\Character{58}}%
\TextAt{8}{79.226}{525.842}{The}%
\TextAt{8}{103.28}{525.842}{corresponding}%
\TextAt{8}{178.466}{525.842}{characteristic}%
\TextAt{8}{250.4}{525.842}{polynomial}%
\TextAt{8}{310.893}{525.842}{is}%
\TextAt{8}{322.657}{525.842}{easily}%
\TextAt{8}{354.912}{525.842}{checked}%
\TextAt{8}{397.831}{525.842}{to}%
\TextAt{8}{412.142}{525.842}{be}%
\TextAt{13}{226.637}{501.907}{t}%
\TextAt{21}{230.865}{506.843}{m}%
\TextAt{9}{241.51}{501.907}{\Character{43}}%
\TextAt{13}{253.271}{501.907}{c}%
\TextAt{17}{258.309}{500.114}{1}%
\TextAt{13}{263.041}{501.907}{t}%
\TextAt{21}{267.269}{506.843}{m}%
\TextAt{3}{274.759}{506.843}{\Character{0}}%
\TextAt{17}{281.346}{506.843}{1}%
\TextAt{9}{288.735}{501.907}{\Character{43}}%
\TextAt{16}{300.496}{501.907}{\Character{1}}%
\TextAt{16}{305.792}{501.907}{\Character{1}}%
\TextAt{16}{311.1}{501.907}{\Character{1}}%
\TextAt{9}{317.1}{501.907}{\Character{43}}%
\TextAt{13}{328.862}{501.907}{c}%
\TextAt{21}{333.9}{500.114}{m}%
\TextAt{9}{345.209}{501.907}{\Character{61}}%
\TextAt{13}{357.635}{501.907}{p}%
\TextAt{21}{363.51}{500.114}{x}%
\TextAt{9}{368.775}{501.907}{\Character{40}}%
\TextAt{13}{373.327}{501.907}{t}%
\TextAt{9}{377.554}{501.907}{\Character{41}}%
\TextAt{13}{382.107}{501.907}{\Character{59}}%
\TextAt{8}{79.226}{477.972}{so}%
\TextAt{8}{93.596}{477.972}{that}%
\TextAt{8}{118.965}{477.972}{finally}%
\TextAt{13}{154.412}{477.972}{f}%
\TextAt{21}{160.182}{476.179}{x}%
\TextAt{9}{165.447}{477.972}{\Character{40}}%
\TextAt{13}{170}{477.972}{t}%
\TextAt{9}{174.227}{477.972}{\Character{41}}%
\TextAt{9}{182.082}{477.972}{\Character{61}}%
\TextAt{13}{194.525}{477.972}{p}%
\TextAt{21}{200.4}{476.179}{x}%
\TextAt{9}{205.666}{477.972}{\Character{40}}%
\TextAt{13}{210.218}{477.972}{t}%
\TextAt{9}{214.445}{477.972}{\Character{41}}%
\TextAt{21}{218.997}{482.311}{d}%
\TextAt{8}{223.853}{477.972}{.}%
\TextAt{8}{97.159}{457.938}{The}%
\TextAt{8}{120.994}{457.938}{set}%
\TextAt{8}{139.043}{457.938}{Hom}%
\TextAt{21}{163.432}{456.144}{K}%
\TextAt{9}{171.635}{457.938}{\Character{40}}%
\TextAt{13}{176.187}{457.938}{L\Character{59}}%
\TextAt{9}{192.528}{460.96}{\Character{22}}%
\TextAt{13}{189.396}{457.938}{K}%
\TextAt{9}{200.213}{457.938}{\Character{41}}%
\TextAt{8}{208.445}{457.938}{of}%
\TextAt{8}{221.556}{457.938}{all}%
\TextAt{13}{237.597}{457.938}{K}%
\TextAt{8}{248.415}{457.938}{-embeddings}%
\TextAt{8}{315.887}{457.938}{of}%
\TextAt{13}{328.998}{457.938}{L}%
\TextAt{8}{340.642}{457.938}{is}%
\TextAt{8}{352.187}{457.938}{partitioned}%
\TextAt{8}{412.462}{457.938}{by}%
\TextAt{8}{428.503}{457.938}{the}%
\TextAt{8}{448.441}{457.938}{equivalence}%
\TextAt{8}{510.007}{457.938}{rela-}%
\TextAt{8}{79.226}{440.892}{tion}%
\TextAt{13}{243.766}{423.846}{\Character{27}}%
\NativeMathAt{254.169}{423.846}{11.955}{\sim}%
\TextAt{13}{266.789}{423.846}{\Character{28}}%
\NativeArrowAt{288.484}{423.846}{11.955}{21.926}{\Longleftrightarrow}%
\TextAt{13}{325.678}{423.846}{\Character{27}}%
\TextAt{13}{332.744}{423.846}{x}%
\TextAt{9}{342.733}{423.846}{\Character{61}}%
\TextAt{13}{355.158}{423.846}{\Character{28}}%
\TextAt{13}{361.566}{423.846}{x}%
\TextAt{8}{79.226}{402.495}{into}%
\TextAt{13}{102.512}{402.495}{m}%
\TextAt{8}{116.192}{402.495}{equivalence}%
\TextAt{8}{177.52}{402.495}{classes}%
\TextAt{8}{214.316}{402.495}{of}%
\TextAt{13}{227.189}{402.495}{d}%
\TextAt{8}{236.712}{402.495}{elements}%
\TextAt{8}{284.112}{402.495}{each.}%
\TextAt{8}{314.848}{402.495}{If}%
\TextAt{13}{326.084}{402.495}{\Character{27}}%
\TextAt{17}{332.737}{400.702}{1}%
\TextAt{13}{337.47}{402.495}{\Character{59}}%
\TextAt{13}{342.694}{402.495}{\Character{58}}%
\TextAt{13}{347.931}{402.495}{\Character{58}}%
\TextAt{13}{353.155}{402.495}{\Character{58}}%
\TextAt{13}{358.392}{402.495}{\Character{59}}%
\TextAt{13}{363.628}{402.495}{\Character{27}}%
\TextAt{21}{370.344}{400.702}{m}%
\TextAt{8}{381.773}{402.495}{is}%
\TextAt{8}{393.08}{402.495}{a}%
\TextAt{8}{402.379}{402.495}{system}%
\TextAt{8}{440.741}{402.495}{of}%
\TextAt{8}{453.614}{402.495}{representatives,}%
\TextAt{8}{79.226}{385.449}{then}%
\TextAt{8}{105.886}{385.449}{we}%
\TextAt{8}{123.114}{385.449}{find}%
\TextAt{13}{254.31}{364.762}{p}%
\TextAt{21}{260.185}{362.969}{x}%
\TextAt{9}{265.45}{364.762}{\Character{40}}%
\TextAt{13}{270.002}{364.762}{t}%
\TextAt{9}{274.229}{364.762}{\Character{41}}%
\TextAt{9}{282.084}{364.762}{\Character{61}}%
\TextAt{21}{298.421}{379.706}{m}%
\TextAt{20}{294.528}{376.12}{\Character{89}}%
\TextAt{21}{295.314}{350.925}{i}%
\TextAt{17}{298.197}{350.925}{\Character{61}1}%
\TextAt{9}{309.804}{364.762}{\Character{40}}%
\TextAt{13}{314.356}{364.762}{t}%
\TextAt{16}{321.24}{364.762}{\Character{0}}%
\TextAt{13}{333.195}{364.762}{\Character{27}}%
\TextAt{21}{339.849}{362.969}{i}%
\TextAt{13}{343.23}{364.762}{x}%
\TextAt{9}{349.882}{364.762}{\Character{41}}%
\TextAt{13}{354.434}{364.762}{\Character{59}}%
\TextAt{8}{79.226}{334.556}{and}%
\TextAt{13}{102.695}{334.556}{f}%
\TextAt{21}{108.465}{332.763}{x}%
\TextAt{9}{113.73}{334.556}{\Character{40}}%
\TextAt{13}{118.282}{334.556}{t}%
\TextAt{9}{122.509}{334.556}{\Character{41}}%
\TextAt{9}{131.571}{334.556}{\Character{61}}%
\TextAt{20}{145.215}{343.522}{\Character{81}}%
\TextAt{21}{156.506}{340.367}{m}%
\TextAt{21}{156.506}{331.069}{i}%
\TextAt{17}{159.389}{331.069}{\Character{61}1}%
\TextAt{9}{170.708}{334.556}{\Character{40}}%
\TextAt{13}{175.26}{334.556}{t}%
\TextAt{16}{182.626}{334.556}{\Character{0}}%
\TextAt{13}{195.062}{334.556}{\Character{27}}%
\TextAt{21}{201.716}{332.763}{i}%
\TextAt{13}{205.097}{334.556}{x}%
\TextAt{9}{211.749}{334.556}{\Character{41}}%
\TextAt{21}{216.301}{338.894}{d}%
\TextAt{9}{225.681}{334.556}{\Character{61}}%
\TextAt{20}{239.31}{343.522}{\Character{81}}%
\TextAt{21}{250.601}{340.367}{m}%
\TextAt{21}{250.601}{331.069}{i}%
\TextAt{17}{253.484}{331.069}{\Character{61}1}%
\TextAt{20}{266.795}{343.522}{\Character{81}}%
\TextAt{21}{278.086}{331.069}{\Character{27}}%
\NativeMathAt{283.212}{331.069}{7.97}{\sim}%
\TextAt{21}{289.798}{331.069}{\Character{27}}%
\TextAt{22}{294.638}{329.854}{i}%
\TextAt{9}{298.297}{334.556}{\Character{40}}%
\TextAt{13}{302.85}{334.556}{t}%
\TextAt{16}{310.215}{334.556}{\Character{0}}%
\TextAt{13}{322.652}{334.556}{\Character{27}}%
\TextAt{13}{329.718}{334.556}{x}%
\TextAt{9}{336.386}{334.556}{\Character{41}}%
\TextAt{9}{345.448}{334.556}{\Character{61}}%
\TextAt{20}{359.092}{343.522}{\Character{81}}%
\TextAt{21}{370.383}{331.069}{\Character{27}}%
\TextAt{9}{376.006}{334.556}{\Character{40}}%
\TextAt{13}{380.558}{334.556}{t}%
\TextAt{16}{387.924}{334.556}{\Character{0}}%
\TextAt{13}{400.36}{334.556}{\Character{27}}%
\TextAt{13}{407.425}{334.556}{x}%
\TextAt{9}{414.095}{334.556}{\Character{41}}%
\TextAt{8}{418.647}{334.556}{.}%
\TextAt{8}{429.214}{334.556}{This}%
\TextAt{8}{456.64}{334.556}{proves}%
\TextAt{8}{493.511}{334.556}{(i),}%
\TextAt{8}{513.904}{334.556}{and}%
\TextAt{8}{79.226}{317.51}{therefore}%
\TextAt{8}{128.326}{317.51}{also}%
\TextAt{8}{151.806}{317.51}{(ii)}%
\TextAt{8}{171.317}{317.51}{and}%
\TextAt{8}{194.079}{317.51}{(iii),}%
\TextAt{8}{220.094}{317.51}{after}%
\TextAt{8}{247.734}{317.51}{Vièta.}%
\Vector{q 0 G 0 g 0.399 w 452.334 -456.49 m 452.334 -448.42 l S Q}%
\Vector{q 0 G 0 g 0.399 w 452.534 -448.62 m 459.707 -448.62 l S Q}%
\Vector{q 0 G 0 g 0.399 w 452.534 -456.291 m 459.707 -456.291 l S Q}%
\Vector{q 0 G 0 g 0.399 w 459.906 -456.49 m 459.906 -448.42 l S Q}%
\TextAt{12}{79.226}{285.426}{Corollary}%
\TextAt{12}{139.014}{285.426}{(2.7).}%
\TextAt{8}{174.33}{285.426}{In}%
\TextAt{8}{188.951}{285.426}{a}%
\TextAt{8}{198.706}{285.426}{tower}%
\TextAt{8}{230.579}{285.426}{of}%
\TextAt{8}{243.909}{285.426}{finite}%
\TextAt{8}{273.82}{285.426}{field}%
\TextAt{8}{299.177}{285.426}{extensions}%
\TextAt{13}{355.558}{285.426}{K}%
\NativeMathAt{369.696}{285.426}{11.955}{\subseteq}%
\TextAt{13}{382.316}{285.426}{L}%
\NativeMathAt{393.601}{285.426}{11.955}{\subseteq}%
\TextAt{13}{406.22}{285.426}{M}%
\TextAt{8}{418.794}{285.426}{,}%
\TextAt{8}{425.943}{285.426}{one}%
\TextAt{8}{447.403}{285.426}{has}%
\TextAt{8}{176.408}{261.491}{Tr}%
\TextAt{21}{189.415}{259.632}{L}%
\TextAt{3}{195.174}{259.631}{\Character{106}}%
\TextAt{21}{197.527}{259.632}{K}%
\TextAt{16}{207.722}{261.491}{\Character{14}}%
\TextAt{8}{215.692}{261.491}{Tr}%
\TextAt{21}{228.699}{259.632}{M}%
\TextAt{3}{237.681}{259.631}{\Character{106}}%
\TextAt{21}{240.033}{259.632}{L}%
\TextAt{9}{249.611}{261.491}{\Character{61}}%
\TextAt{8}{262.037}{261.491}{Tr}%
\TextAt{21}{275.044}{259.632}{M}%
\TextAt{3}{284.026}{259.631}{\Character{106}}%
\TextAt{21}{286.378}{259.632}{K}%
\TextAt{13}{294.581}{261.491}{\Character{59}}%
\TextAt{13}{323.716}{261.491}{N}%
\TextAt{21}{333.109}{259.632}{L}%
\TextAt{3}{338.868}{259.631}{\Character{106}}%
\TextAt{21}{341.22}{259.632}{K}%
\TextAt{16}{352.08}{261.491}{\Character{14}}%
\TextAt{13}{360.714}{261.491}{N}%
\TextAt{21}{370.088}{259.632}{M}%
\TextAt{3}{379.069}{259.631}{\Character{106}}%
\TextAt{21}{381.422}{259.632}{L}%
\TextAt{9}{391}{261.491}{\Character{61}}%
\TextAt{13}{403.425}{261.491}{N}%
\TextAt{21}{412.799}{259.632}{M}%
\TextAt{3}{421.781}{259.631}{\Character{106}}%
\TextAt{21}{424.133}{259.632}{K}%
\TextAt{13}{432.336}{261.491}{\Character{58}}%
\TextAt{19}{79.226}{222.518}{Proof.}%
\TextAt{8}{116}{222.518}{We}%
\TextAt{8}{136.671}{222.518}{assume}%
\TextAt{8}{177.641}{222.518}{that}%
\TextAt{13}{203.536}{222.518}{M}%
\TextAt{16}{220.327}{222.518}{\Character{106}}%
\TextAt{13}{227.865}{222.518}{K}%
\TextAt{8}{243.107}{222.518}{is}%
\TextAt{8}{255.397}{222.518}{separable.}%
\TextAt{8}{312.969}{222.518}{The}%
\TextAt{8}{337.549}{222.518}{set}%
\TextAt{8}{356.342}{222.518}{Hom}%
\TextAt{21}{380.731}{220.725}{K}%
\TextAt{9}{388.934}{222.518}{\Character{40}}%
\TextAt{13}{393.486}{222.518}{M}%
\TextAt{13}{405.405}{222.518}{\Character{59}}%
\TextAt{9}{413.786}{225.54}{\Character{22}}%
\TextAt{13}{410.653}{222.518}{K}%
\TextAt{9}{421.471}{222.518}{\Character{41}}%
\TextAt{8}{430.447}{222.518}{of}%
\TextAt{13}{444.303}{222.518}{K}%
\TextAt{8}{455.121}{222.518}{-embeddings}%
\TextAt{8}{523.337}{222.518}{of}%
\TextAt{13}{79.226}{205.473}{M}%
\TextAt{8}{95.697}{205.473}{is}%
\TextAt{8}{107.461}{205.473}{partitioned}%
\TextAt{8}{167.954}{205.473}{by}%
\TextAt{8}{184.213}{205.473}{the}%
\TextAt{8}{204.37}{205.473}{relation}%
\TextAt{13}{240.84}{181.538}{\Character{27}}%
\NativeMathAt{251.243}{181.538}{11.955}{\sim}%
\TextAt{13}{263.863}{181.538}{\Character{28}}%
\NativeArrowAt{285.558}{181.538}{11.955}{21.926}{\Longleftrightarrow}%
\TextAt{13}{322.752}{181.538}{\Character{27}}%
\TextAt{16}{329.834}{181.538}{\Character{106}}%
\TextAt{21}{333.155}{179.744}{L}%
\TextAt{9}{342.733}{181.538}{\Character{61}}%
\TextAt{13}{355.158}{181.538}{\Character{28}}%
\TextAt{16}{361.577}{181.538}{\Character{106}}%
\TextAt{21}{364.898}{179.744}{L}%
\TextAt{8}{79.226}{157.603}{into}%
\TextAt{13}{103.658}{157.603}{m}%
\TextAt{9}{118.393}{157.603}{\Character{61}}%
\TextAt{9}{131.986}{157.603}{\Character{91}}%
\TextAt{13}{135.245}{157.603}{L}%
\TextAt{9}{147.704}{157.603}{\Character{58}}%
\TextAt{13}{155.451}{157.603}{K}%
\TextAt{9}{166.269}{157.603}{\Character{93}}%
\TextAt{8}{174.107}{157.603}{equivalence}%
\TextAt{8}{236.581}{157.603}{classes.}%
\TextAt{8}{280.452}{157.603}{If}%
\TextAt{13}{292.833}{157.603}{\Character{27}}%
\TextAt{17}{299.486}{155.809}{1}%
\TextAt{13}{304.219}{157.603}{\Character{59}}%
\TextAt{13}{309.443}{157.603}{\Character{58}}%
\TextAt{13}{314.68}{157.603}{\Character{58}}%
\TextAt{13}{319.904}{157.603}{\Character{58}}%
\TextAt{13}{325.141}{157.603}{\Character{59}}%
\TextAt{13}{330.377}{157.603}{\Character{27}}%
\TextAt{21}{337.093}{155.809}{m}%
\TextAt{8}{349.668}{157.603}{is}%
\TextAt{8}{362.121}{157.603}{a}%
\TextAt{8}{372.566}{157.603}{system}%
\TextAt{8}{412.074}{157.603}{of}%
\TextAt{8}{426.093}{157.603}{representatives,}%
\TextAt{8}{510.007}{157.603}{then}%
\TextAt{8}{79.226}{140.557}{Hom}%
\TextAt{21}{103.615}{138.763}{K}%
\TextAt{9}{111.817}{140.557}{\Character{40}}%
\TextAt{13}{116.37}{140.557}{L\Character{59}}%
\TextAt{9}{132.711}{143.579}{\Character{22}}%
\TextAt{13}{129.578}{140.557}{K}%
\TextAt{9}{140.396}{140.557}{\Character{41}}%
\TextAt{9}{148.251}{140.557}{\Character{61}}%
\TextAt{16}{160.695}{140.557}{\Character{102}}%
\TextAt{13}{166.672}{140.557}{\Character{27}}%
\TextAt{21}{173.326}{138.763}{i}%
\TextAt{16}{176.707}{140.557}{\Character{106}}%
\TextAt{21}{180.028}{138.763}{L}%
\TextAt{16}{189.606}{140.557}{\Character{106}}%
\TextAt{13}{196.248}{140.557}{i}%
\TextAt{9}{203.562}{140.557}{\Character{61}}%
\TextAt{9}{215.972}{140.557}{1}%
\TextAt{13}{221.84}{140.557}{\Character{59}}%
\TextAt{13}{227.064}{140.557}{\Character{58}}%
\TextAt{13}{232.301}{140.557}{\Character{58}}%
\TextAt{13}{237.537}{140.557}{\Character{58}}%
\TextAt{13}{242.762}{140.557}{\Character{59}}%
\TextAt{13}{247.998}{140.557}{m}%
\TextAt{16}{258.3}{140.557}{\Character{103}}%
\TextAt{8}{264.278}{140.557}{,}%
\TextAt{8}{271.427}{140.557}{and}%
\TextAt{8}{294.19}{140.557}{we}%
\TextAt{8}{311.417}{140.557}{find}%
\TextAt{8}{146.229}{110.111}{Tr}%
\TextAt{21}{159.236}{108.252}{M}%
\TextAt{3}{168.218}{108.252}{\Character{106}}%
\TextAt{21}{170.57}{108.252}{K}%
\TextAt{9}{178.773}{110.111}{\Character{40}}%
\TextAt{13}{183.325}{110.111}{x}%
\TextAt{9}{189.977}{110.111}{\Character{41}}%
\TextAt{9}{197.832}{110.111}{\Character{61}}%
\TextAt{21}{215.165}{125.056}{m}%
\TextAt{20}{210.276}{121.469}{\Character{88}}%
\TextAt{21}{212.059}{96.275}{i}%
\TextAt{17}{214.942}{96.275}{\Character{61}1}%
\TextAt{20}{230.759}{121.469}{\Character{88}}%
\TextAt{21}{229.537}{96.363}{\Character{27}}%
\NativeMathAt{234.663}{96.363}{7.97}{\sim}%
\TextAt{21}{241.249}{96.363}{\Character{27}}%
\TextAt{22}{246.088}{95.148}{i}%
\TextAt{13}{251.243}{110.111}{\Character{27}}%
\TextAt{13}{258.308}{110.111}{x}%
\TextAt{9}{268.298}{110.112}{\Character{61}}%
\TextAt{21}{285.612}{125.056}{m}%
\TextAt{20}{280.723}{121.469}{\Character{88}}%
\TextAt{21}{282.506}{96.275}{i}%
\TextAt{17}{285.389}{96.275}{\Character{61}1}%
\TextAt{8}{299.985}{110.111}{Tr}%
\TextAt{21}{312.992}{108.252}{\Character{27}}%
\TextAt{22}{317.831}{107.037}{i}%
\TextAt{21}{320.993}{108.252}{M}%
\TextAt{3}{329.975}{108.252}{\Character{106}}%
\TextAt{21}{332.327}{108.252}{\Character{27}}%
\TextAt{22}{337.166}{107.037}{i}%
\TextAt{21}{340.328}{108.252}{L}%
\TextAt{9}{346.585}{110.111}{\Character{40}}%
\TextAt{13}{351.137}{110.112}{\Character{27}}%
\TextAt{21}{357.791}{108.318}{i}%
\TextAt{13}{361.172}{110.111}{x}%
\TextAt{9}{367.824}{110.111}{\Character{41}}%
\TextAt{9}{375.679}{110.111}{\Character{61}}%
\TextAt{21}{393.012}{125.056}{m}%
\TextAt{20}{388.123}{121.469}{\Character{88}}%
\TextAt{21}{389.905}{96.275}{i}%
\TextAt{17}{392.788}{96.275}{\Character{61}1}%
\TextAt{13}{407.384}{110.111}{\Character{27}}%
\TextAt{21}{414.037}{108.318}{i}%
\TextAt{8}{419.411}{110.112}{Tr}%
\TextAt{21}{432.418}{108.252}{M}%
\TextAt{3}{441.4}{108.252}{\Character{106}}%
\TextAt{21}{443.752}{108.252}{L}%
\TextAt{9}{450.01}{110.112}{\Character{40}}%
\TextAt{13}{454.562}{110.112}{x}%
\TextAt{9}{461.214}{110.112}{\Character{41}}%
\TextAt{9}{197.851}{78.035}{\Character{61}}%
\TextAt{8}{210.276}{78.035}{Tr}%
\TextAt{21}{223.283}{76.176}{L}%
\TextAt{3}{229.042}{76.176}{\Character{106}}%
\TextAt{21}{231.395}{76.176}{K}%
\TextAt{9}{239.597}{78.035}{\Character{40}}%
\TextAt{8}{244.15}{78.035}{Tr}%
\TextAt{21}{257.157}{76.176}{M}%
\TextAt{3}{266.139}{76.176}{\Character{106}}%
\TextAt{21}{268.491}{76.176}{L}%
\TextAt{9}{274.748}{78.035}{\Character{40}}%
\TextAt{13}{279.3}{78.035}{x}%
\TextAt{9}{285.953}{78.035}{\Character{41}\Character{41}}%
\TextAt{13}{295.057}{78.035}{\Character{58}}%
\TextAt{8}{303.069}{43.033}{9}%
\Destination{78.23}{809.815}{nav.page.9}%
\Destination{79.36}{540.855}{nav.print.10}%
\Destination{79.36}{299.455}{nav.number.1.2.7}%
}

% Page 22 (body page 10)
\Page{%
\TextAt{2}{79.226}{796.943}{I.}%
\TextAt{2}{89.83}{796.943}{Algebraic}%
\TextAt{2}{138.648}{796.943}{Integers}%
\Vector{q 0 G 0 g 0.299 w 7.226 18.729 m 460.769 18.729 l S Q}%
\TextAt{8}{79.226}{761.124}{Likewise}%
\TextAt{8}{126.27}{761.124}{for}%
\TextAt{8}{144.155}{761.124}{the}%
\TextAt{8}{164.311}{761.124}{norm.}%
\ErrataAt{97.1}{741.078}{437.9}{We will not need the inseparable case for the sequel. Put $q_{B/A}=[B:A]_i$ for a finite}%
\ErrataAt{79.19}{724.032}{455.81}{extension $B\mid A$. Each $K$-embedding of $L$ extends to $[M:L]_s$ embeddings of $M$.}%
\ErrataAt{79.19}{706.986}{455.81}{In \hyperlink{nav.number.1.2.6}{(2.6)}, the trace and the norm count each embedding with its inseparable multiplicity.}%
\ErrataAt{79.19}{689.94}{455.81}{Grouping the sums and products by restriction to $L$ proves both tower identities, since}%
\ErrataAt{79.19}{672.894}{455.81}{$q_{M/K}=q_{M/L}q_{L/K}$. If $M^s$ is the maximal separable subextension of $M\mid K$ and}%
\ErrataAt{79.16}{663.655}{455.84}{ }%
\ErrataAt{79.19}{655.848}{455.81}{$q=[M:M^s]$, then $N_{M\mid K}(x)=N_{M^s\mid K}(x^q)$ and $\operatorname{Tr}_{M\mid K}(x)=0$ for $q>1$}%
\ErrataAt{79.19}{637.255}{455.81}{(see \hyperlink{nav.cite.N143}{[143]}, vol. I, chap. II, \S10).}%
\Vector{q 0 G 0 g 0.399 w 452.334 -136.708 m 452.334 -128.639 l S Q}%
\Vector{q 0 G 0 g 0.399 w 452.534 -128.838 m 459.707 -128.838 l S Q}%
\Vector{q 0 G 0 g 0.399 w 452.534 -136.509 m 459.707 -136.509 l S Q}%
\Vector{q 0 G 0 g 0.399 w 459.906 -136.708 m 459.906 -128.639 l S Q}%
\TextAt{8}{97.159}{609.054}{The}%
\TextAt{12}{120.964}{609.054}{discriminant}%
\TextAt{12}{198.208}{609.054}{of}%
\TextAt{12}{213.249}{609.054}{a}%
\TextAt{12}{223.985}{609.054}{basis}%
\TextAt{13}{256.004}{609.054}{\Character{11}}%
\TextAt{17}{263.45}{607.26}{1}%
\TextAt{13}{268.182}{609.054}{\Character{59}}%
\TextAt{13}{273.406}{609.054}{\Character{58}}%
\TextAt{13}{278.643}{609.054}{\Character{58}}%
\TextAt{13}{283.867}{609.054}{\Character{58}}%
\TextAt{13}{289.104}{609.054}{\Character{59}}%
\TextAt{13}{294.34}{609.054}{\Character{11}}%
\TextAt{21}{301.848}{607.26}{n}%
\TextAt{8}{311.134}{609.054}{of}%
\TextAt{8}{324.215}{609.054}{a}%
\TextAt{8}{333.722}{609.054}{separable}%
\TextAt{8}{384.917}{609.054}{extension}%
\TextAt{13}{436.434}{609.054}{L}%
\TextAt{16}{447.72}{609.054}{\Character{106}}%
\TextAt{13}{454.361}{609.054}{K}%
\TextAt{8}{468.828}{609.054}{is}%
\TextAt{8}{480.343}{609.054}{defined}%
\TextAt{8}{520.408}{609.054}{by}%
\TextAt{13}{232.633}{582.724}{d}%
\TextAt{9}{238.716}{582.724}{\Character{40}}%
\TextAt{13}{243.268}{582.724}{\Character{11}}%
\TextAt{17}{250.714}{580.931}{1}%
\TextAt{13}{255.446}{582.724}{\Character{59}}%
\TextAt{13}{260.67}{582.724}{\Character{58}}%
\TextAt{13}{265.907}{582.724}{\Character{58}}%
\TextAt{13}{271.143}{582.724}{\Character{58}}%
\TextAt{13}{276.368}{582.724}{\Character{59}}%
\TextAt{13}{281.604}{582.724}{\Character{11}}%
\TextAt{21}{289.113}{580.931}{n}%
\TextAt{9}{294.749}{582.724}{\Character{41}}%
\TextAt{9}{302.604}{582.724}{\Character{61}}%
\TextAt{8}{315.048}{582.724}{det}%
\TextAt{9}{331.307}{582.724}{\Character{40}\Character{40}}%
\TextAt{13}{340.411}{582.724}{\Character{27}}%
\TextAt{21}{347.065}{580.931}{i}%
\TextAt{13}{350.446}{582.724}{\Character{11}}%
\TextAt{21}{357.892}{580.931}{j}%
\TextAt{9}{362.274}{582.724}{\Character{41}\Character{41}}%
\TextAt{17}{371.378}{587.66}{2}%
\TextAt{13}{376.111}{582.724}{\Character{59}}%
\TextAt{8}{79.226}{556.394}{where}%
\TextAt{13}{113.035}{556.394}{\Character{27}}%
\TextAt{21}{119.689}{554.601}{i}%
\TextAt{8}{123.07}{556.394}{,}%
\TextAt{13}{130.219}{556.394}{i}%
\TextAt{9}{137.534}{556.394}{\Character{61}}%
\TextAt{9}{149.944}{556.394}{1}%
\TextAt{13}{155.812}{556.394}{\Character{59}}%
\TextAt{13}{161.036}{556.394}{\Character{58}}%
\TextAt{13}{166.273}{556.394}{\Character{58}}%
\TextAt{13}{171.497}{556.394}{\Character{58}}%
\TextAt{13}{176.734}{556.394}{\Character{59}}%
\TextAt{13}{181.97}{556.394}{n}%
\TextAt{8}{189.02}{556.394}{,}%
\TextAt{8}{196.17}{556.394}{varies}%
\TextAt{8}{229.082}{556.394}{over}%
\TextAt{8}{254.128}{556.394}{the}%
\TextAt{13}{274.285}{556.394}{K}%
\TextAt{8}{285.102}{556.394}{-embeddings}%
\TextAt{13}{352.793}{556.394}{L}%
\NativeMathAt{364.078}{556.394}{11.955}{\rightarrow}%
\TextAt{9}{382.487}{559.416}{\Character{22}}%
\TextAt{13}{379.354}{556.394}{K}%
\TextAt{8}{390.172}{556.394}{.}%
\TextAt{8}{398.62}{556.394}{Because}%
\TextAt{8}{443.38}{556.394}{of}%
\TextAt{8}{456.71}{556.394}{the}%
\TextAt{8}{476.867}{556.394}{relation}%
\TextAt{8}{225.119}{530.065}{Tr}%
\TextAt{21}{238.126}{528.205}{L}%
\TextAt{3}{243.885}{528.205}{\Character{106}}%
\TextAt{21}{246.238}{528.205}{K}%
\TextAt{9}{254.44}{530.065}{\Character{40}}%
\TextAt{13}{258.993}{530.065}{\Character{11}}%
\TextAt{21}{266.438}{528.271}{i}%
\TextAt{13}{269.82}{530.065}{\Character{11}}%
\TextAt{21}{277.265}{528.271}{j}%
\TextAt{9}{281.647}{530.065}{\Character{41}}%
\TextAt{9}{289.502}{530.065}{\Character{61}}%
\TextAt{20}{301.946}{541.422}{\Character{88}}%
\TextAt{21}{308.269}{515.962}{k}%
\TextAt{9}{319.215}{530.065}{\Character{40}}%
\TextAt{13}{323.767}{530.065}{\Character{27}}%
\TextAt{21}{330.42}{528.271}{k}%
\TextAt{13}{335.54}{530.065}{\Character{11}}%
\TextAt{21}{342.986}{528.271}{i}%
\TextAt{9}{346.367}{530.065}{\Character{41}\Character{40}}%
\TextAt{13}{355.472}{530.065}{\Character{27}}%
\TextAt{21}{362.125}{528.271}{k}%
\TextAt{13}{367.245}{530.065}{\Character{11}}%
\TextAt{21}{374.69}{528.271}{j}%
\TextAt{9}{379.073}{530.065}{\Character{41}}%
\TextAt{13}{383.625}{530.065}{\Character{59}}%
\TextAt{8}{79.226}{495.246}{the}%
\TextAt{8}{100.367}{495.246}{matrix}%
\TextAt{9}{139.405}{495.246}{\Character{40}}%
\TextAt{8}{143.958}{495.246}{Tr}%
\TextAt{21}{156.965}{493.386}{L}%
\TextAt{3}{162.724}{493.386}{\Character{106}}%
\TextAt{21}{165.076}{493.386}{K}%
\TextAt{9}{173.279}{495.246}{\Character{40}}%
\TextAt{13}{177.831}{495.246}{\Character{11}}%
\TextAt{21}{185.277}{493.452}{i}%
\TextAt{13}{188.658}{495.246}{\Character{11}}%
\TextAt{21}{196.104}{493.452}{j}%
\TextAt{9}{200.486}{495.246}{\Character{41}\Character{41}}%
\TextAt{8}{214.473}{495.246}{is}%
\TextAt{8}{227.221}{495.246}{the}%
\TextAt{8}{248.362}{495.246}{product}%
\TextAt{8}{293.246}{495.246}{of}%
\TextAt{8}{307.561}{495.246}{the}%
\TextAt{8}{328.702}{495.246}{matrices}%
\TextAt{9}{376.575}{495.246}{\Character{40}}%
\TextAt{13}{381.128}{495.246}{\Character{27}}%
\TextAt{21}{387.781}{493.452}{k}%
\TextAt{13}{392.901}{495.246}{\Character{11}}%
\TextAt{21}{400.346}{493.452}{i}%
\TextAt{9}{403.728}{495.246}{\Character{41}}%
\TextAt{21}{408.28}{499.584}{t}%
\TextAt{8}{416.718}{495.246}{and}%
\TextAt{9}{440.466}{495.246}{\Character{40}}%
\TextAt{13}{445.018}{495.246}{\Character{27}}%
\TextAt{21}{451.671}{493.452}{k}%
\TextAt{13}{456.791}{495.246}{\Character{11}}%
\TextAt{21}{464.237}{493.452}{j}%
\TextAt{9}{468.619}{495.246}{\Character{41}}%
\TextAt{8}{473.171}{495.246}{.}%
\TextAt{8}{484.574}{495.246}{Thus}%
\TextAt{8}{515.207}{495.246}{one}%
\TextAt{8}{79.226}{478.2}{may}%
\TextAt{8}{104.595}{478.2}{also}%
\TextAt{8}{128.075}{478.2}{write}%
\TextAt{13}{219.943}{461.154}{d}%
\TextAt{9}{226.025}{461.154}{\Character{40}}%
\TextAt{13}{230.578}{461.154}{\Character{11}}%
\TextAt{17}{238.023}{459.36}{1}%
\TextAt{13}{242.756}{461.154}{\Character{59}}%
\TextAt{13}{247.98}{461.154}{\Character{58}}%
\TextAt{13}{253.217}{461.154}{\Character{58}}%
\TextAt{13}{258.441}{461.154}{\Character{58}}%
\TextAt{13}{263.678}{461.154}{\Character{59}}%
\TextAt{13}{268.914}{461.154}{\Character{11}}%
\TextAt{21}{276.422}{459.36}{n}%
\TextAt{9}{282.058}{461.154}{\Character{41}}%
\TextAt{9}{289.913}{461.154}{\Character{61}}%
\TextAt{8}{302.357}{461.154}{det}%
\TextAt{9}{318.616}{461.154}{\Character{40}}%
\TextAt{8}{323.168}{461.154}{Tr}%
\TextAt{21}{336.176}{459.294}{L}%
\TextAt{3}{341.935}{459.294}{\Character{106}}%
\TextAt{21}{344.287}{459.294}{K}%
\TextAt{9}{352.49}{461.154}{\Character{40}}%
\TextAt{13}{357.042}{461.154}{\Character{11}}%
\TextAt{21}{364.488}{459.36}{i}%
\TextAt{13}{367.869}{461.154}{\Character{11}}%
\TextAt{21}{375.314}{459.36}{j}%
\TextAt{9}{379.697}{461.154}{\Character{41}\Character{41}}%
\TextAt{13}{388.801}{461.154}{\Character{58}}%
\TextAt{8}{79.226}{438.777}{In}%
\TextAt{8}{93.847}{438.777}{the}%
\TextAt{8}{114.004}{438.777}{special}%
\TextAt{8}{152.105}{438.777}{case}%
\TextAt{8}{176.876}{438.777}{of}%
\TextAt{8}{190.206}{438.777}{a}%
\TextAt{8}{199.961}{438.777}{basis}%
\TextAt{8}{228.702}{438.777}{of}%
\TextAt{8}{242.032}{438.777}{type}%
\TextAt{9}{268.369}{438.777}{1}%
\TextAt{13}{274.222}{438.777}{\Character{59}}%
\TextAt{13}{279.446}{438.777}{\Character{18}}%
\TextAt{13}{285.221}{438.777}{\Character{59}}%
\TextAt{13}{290.457}{438.777}{\Character{58}}%
\TextAt{13}{295.694}{438.777}{\Character{58}}%
\TextAt{13}{300.918}{438.777}{\Character{58}}%
\TextAt{13}{306.154}{438.777}{\Character{59}}%
\TextAt{13}{311.391}{438.777}{\Character{18}}%
\TextAt{21}{317.248}{443.115}{n}%
\TextAt{3}{322.386}{443.115}{\Character{0}}%
\TextAt{17}{328.972}{443.115}{1}%
\TextAt{8}{337.602}{438.777}{one}%
\TextAt{8}{359.061}{438.777}{gets}%
\TextAt{13}{225.038}{412.447}{d}%
\TextAt{9}{231.121}{412.447}{\Character{40}1}%
\TextAt{13}{241.526}{412.447}{\Character{59}}%
\TextAt{13}{246.75}{412.447}{\Character{18}}%
\TextAt{13}{252.525}{412.447}{\Character{59}}%
\TextAt{13}{257.761}{412.447}{\Character{58}}%
\TextAt{13}{262.998}{412.447}{\Character{58}}%
\TextAt{13}{268.222}{412.447}{\Character{58}}%
\TextAt{13}{273.458}{412.447}{\Character{59}}%
\TextAt{13}{278.695}{412.447}{\Character{18}}%
\TextAt{21}{284.552}{417.383}{n}%
\TextAt{3}{289.69}{417.383}{\Character{0}}%
\TextAt{17}{296.276}{417.383}{1}%
\TextAt{9}{301.009}{412.447}{\Character{41}}%
\TextAt{9}{308.864}{412.447}{\Character{61}}%
\TextAt{20}{321.307}{423.805}{\Character{89}}%
\TextAt{21}{322.269}{398.61}{i\Character{60}j}%
\TextAt{9}{336.583}{412.447}{\Character{40}}%
\TextAt{13}{341.136}{412.447}{\Character{18}}%
\TextAt{21}{346.591}{410.654}{i}%
\TextAt{16}{352.629}{412.447}{\Character{0}}%
\TextAt{13}{364.584}{412.447}{\Character{18}}%
\TextAt{21}{370.039}{410.654}{j}%
\TextAt{9}{374.421}{412.447}{\Character{41}}%
\TextAt{17}{378.974}{417.383}{2}%
\TextAt{13}{383.706}{412.447}{\Character{59}}%
\TextAt{8}{79.226}{376.619}{where}%
\TextAt{13}{112.95}{376.619}{\Character{18}}%
\TextAt{21}{118.406}{374.826}{i}%
\TextAt{9}{125.108}{376.619}{\Character{61}}%
\TextAt{13}{137.533}{376.619}{\Character{27}}%
\TextAt{21}{144.187}{374.826}{i}%
\TextAt{13}{147.568}{376.619}{\Character{18}}%
\TextAt{8}{153.348}{376.619}{.}%
\TextAt{8}{161.768}{376.619}{This}%
\TextAt{8}{188.403}{376.619}{is}%
\TextAt{8}{200.082}{376.619}{seen}%
\TextAt{8}{225.414}{376.619}{by}%
\TextAt{8}{241.588}{376.619}{successively}%
\TextAt{8}{305.093}{376.619}{multiplying}%
\TextAt{8}{367.45}{376.619}{each}%
\TextAt{8}{393.702}{376.619}{of}%
\TextAt{8}{406.947}{376.619}{the}%
\TextAt{8}{427.019}{376.619}{first}%
\TextAt{9}{451.059}{376.619}{\Character{40}}%
\TextAt{13}{455.612}{376.619}{n}%
\TextAt{16}{465.082}{376.619}{\Character{0}}%
\TextAt{9}{476.864}{376.619}{1\Character{41}}%
\TextAt{8}{491.082}{376.619}{columns}%
\TextAt{8}{79.226}{359.573}{in}%
\TextAt{8}{92.879}{359.573}{the}%
\TextAt{8}{113.035}{359.573}{Vandermonde}%
\TextAt{8}{187.169}{359.573}{matrix}%
\TextAt{20}{243.501}{348.564}{\Character{48}}%
\TextAt{20}{243.501}{327.522}{\Character{66}}%
\TextAt{20}{243.501}{320.349}{\Character{66}}%
\TextAt{20}{243.501}{313.176}{\Character{66}}%
\TextAt{20}{243.501}{306.003}{\Character{66}}%
\TextAt{20}{243.501}{298.351}{\Character{64}}%
\TextAt{9}{253.962}{336.923}{1}%
\TextAt{13}{270.23}{336.923}{\Character{18}}%
\TextAt{17}{275.685}{335.13}{1}%
\TextAt{13}{291.121}{336.923}{\Character{18}}%
\TextAt{17}{296.902}{341.262}{2}%
\TextAt{17}{296.576}{333.968}{1}%
\TextAt{16}{311.886}{336.923}{\Character{1}}%
\TextAt{16}{317.182}{336.923}{\Character{1}}%
\TextAt{16}{322.49}{336.923}{\Character{1}}%
\TextAt{13}{335.796}{336.923}{\Character{18}}%
\TextAt{21}{341.577}{341.871}{n}%
\TextAt{3}{346.715}{341.871}{\Character{0}}%
\TextAt{17}{353.301}{341.871}{1}%
\TextAt{17}{341.251}{333.992}{1}%
\TextAt{9}{253.962}{319.877}{1}%
\TextAt{13}{270.23}{319.877}{\Character{18}}%
\TextAt{17}{275.685}{318.084}{2}%
\TextAt{13}{291.121}{319.877}{\Character{18}}%
\TextAt{17}{296.902}{324.216}{2}%
\TextAt{17}{296.576}{316.922}{2}%
\TextAt{16}{311.886}{319.877}{\Character{1}}%
\TextAt{16}{317.182}{319.877}{\Character{1}}%
\TextAt{16}{322.49}{319.877}{\Character{1}}%
\TextAt{13}{335.796}{319.877}{\Character{18}}%
\TextAt{21}{341.577}{324.826}{n}%
\TextAt{3}{346.715}{324.826}{\Character{0}}%
\TextAt{17}{353.301}{324.826}{1}%
\TextAt{17}{341.251}{316.946}{2}%
\TextAt{8}{255.263}{307.626}{.}%
\TextAt{8}{255.263}{303.641}{.}%
\TextAt{8}{255.263}{299.656}{.}%
\TextAt{8}{273.697}{307.626}{.}%
\TextAt{8}{273.697}{303.641}{.}%
\TextAt{8}{273.697}{299.656}{.}%
\TextAt{8}{294.752}{307.626}{.}%
\TextAt{8}{294.752}{303.641}{.}%
\TextAt{8}{294.752}{299.656}{.}%
\TextAt{8}{345.289}{307.626}{.}%
\TextAt{8}{345.289}{303.641}{.}%
\TextAt{8}{345.289}{299.656}{.}%
\TextAt{9}{253.962}{282.61}{1}%
\TextAt{13}{269.778}{282.61}{\Character{18}}%
\TextAt{21}{275.233}{280.817}{n}%
\TextAt{13}{290.832}{282.61}{\Character{18}}%
\TextAt{17}{296.612}{286.949}{2}%
\TextAt{21}{296.287}{279.655}{n}%
\TextAt{16}{311.886}{282.61}{\Character{1}}%
\TextAt{16}{317.182}{282.61}{\Character{1}}%
\TextAt{16}{322.49}{282.61}{\Character{1}}%
\TextAt{13}{335.796}{282.61}{\Character{18}}%
\TextAt{21}{341.577}{286.949}{n}%
\TextAt{3}{346.715}{286.949}{\Character{0}}%
\TextAt{17}{353.301}{286.949}{1}%
\TextAt{21}{341.251}{279.655}{n}%
\TextAt{20}{358.034}{348.564}{\Character{49}}%
\TextAt{20}{358.034}{327.522}{\Character{67}}%
\TextAt{20}{358.034}{320.349}{\Character{67}}%
\TextAt{20}{358.034}{313.176}{\Character{67}}%
\TextAt{20}{358.034}{306.003}{\Character{67}}%
\TextAt{20}{358.034}{298.351}{\Character{65}}%
\TextAt{8}{79.226}{258.614}{by}%
\TextAt{13}{95.485}{258.614}{\Character{18}}%
\TextAt{17}{100.94}{256.82}{1}%
\TextAt{8}{109.57}{258.614}{and}%
\TextAt{8}{132.333}{258.614}{subtracting}%
\TextAt{8}{194.189}{258.614}{it}%
\TextAt{8}{205.893}{258.614}{from}%
\TextAt{8}{233.533}{258.614}{the}%
\TextAt{8}{253.69}{258.614}{following.}%
\TextAt{12}{79.226}{233.365}{Proposition}%
\TextAt{12}{151.584}{233.365}{(2.8).}%
\TextAt{8}{186.593}{233.365}{If}%
\TextAt{13}{197.978}{233.365}{L}%
\TextAt{16}{209.264}{233.365}{\Character{106}}%
\TextAt{13}{215.905}{233.365}{K}%
\TextAt{8}{230.313}{233.365}{is}%
\TextAt{8}{241.77}{233.365}{separable}%
\TextAt{8}{292.906}{233.365}{and}%
\TextAt{13}{315.362}{233.365}{\Character{11}}%
\TextAt{17}{322.808}{231.571}{1}%
\TextAt{13}{327.54}{233.365}{\Character{59}}%
\TextAt{13}{332.764}{233.365}{\Character{58}}%
\TextAt{13}{338.001}{233.365}{\Character{58}}%
\TextAt{13}{343.225}{233.365}{\Character{58}}%
\TextAt{13}{348.462}{233.365}{\Character{59}}%
\TextAt{13}{353.698}{233.365}{\Character{11}}%
\TextAt{21}{361.206}{231.571}{n}%
\TextAt{8}{370.433}{233.365}{is}%
\TextAt{8}{381.89}{233.365}{a}%
\TextAt{8}{391.339}{233.365}{basis,}%
\TextAt{8}{423.085}{233.365}{then}%
\TextAt{8}{449.438}{233.365}{the}%
\TextAt{8}{469.288}{233.365}{discriminant}%
\TextAt{13}{260.238}{207.035}{d}%
\TextAt{9}{266.321}{207.035}{\Character{40}}%
\TextAt{13}{270.873}{207.035}{\Character{11}}%
\TextAt{17}{278.319}{205.242}{1}%
\TextAt{13}{283.051}{207.035}{\Character{59}}%
\TextAt{13}{288.275}{207.035}{\Character{58}}%
\TextAt{13}{293.512}{207.035}{\Character{58}}%
\TextAt{13}{298.748}{207.035}{\Character{58}}%
\TextAt{13}{303.973}{207.035}{\Character{59}}%
\TextAt{13}{309.209}{207.035}{\Character{11}}%
\TextAt{21}{316.718}{205.242}{n}%
\TextAt{9}{322.354}{207.035}{\Character{41}}%
\NativeMathAt{330.227}{207.035}{11.955}{\neq}%
\TextAt{9}{342.637}{207.035}{0}%
\TextAt{13}{348.506}{207.035}{\Character{59}}%
\TextAt{8}{79.226}{180.705}{and}%
\TextAt{9}{258.949}{163.659}{\Character{40}}%
\TextAt{13}{263.501}{163.659}{x\Character{59}}%
\TextAt{13}{275.373}{163.659}{y}%
\TextAt{9}{281.534}{163.659}{\Character{41}}%
\TextAt{9}{289.389}{163.659}{\Character{61}}%
\TextAt{8}{301.832}{163.659}{Tr}%
\TextAt{21}{314.84}{161.799}{L}%
\TextAt{3}{320.599}{161.799}{\Character{106}}%
\TextAt{21}{322.951}{161.799}{K}%
\TextAt{9}{331.154}{163.659}{\Character{40}}%
\TextAt{13}{335.706}{163.659}{xy}%
\TextAt{9}{348.495}{163.659}{\Character{41}}%
\TextAt{8}{79.226}{141.282}{is}%
\TextAt{8}{90.99}{141.282}{a}%
\TextAt{8}{100.745}{141.282}{nondegenerate}%
\TextAt{8}{178.143}{141.282}{bilinear}%
\TextAt{8}{220.417}{141.282}{form}%
\TextAt{8}{248.057}{141.282}{on}%
\TextAt{8}{264.316}{141.282}{the}%
\TextAt{13}{284.473}{141.282}{K}%
\TextAt{8}{295.29}{141.282}{-vector}%
\TextAt{8}{334.312}{141.282}{space}%
\TextAt{13}{365.587}{141.282}{L}%
\TextAt{8}{373.551}{141.282}{.}%
\TextAt{8}{300.14}{43.033}{10}%
\Destination{78.23}{809.815}{nav.page.10}%
\Destination{79.36}{624.055}{nav.print.11}%
\Destination{79.36}{247.355}{nav.number.1.2.8}%
\Destination{79.226}{703.955}{errata.6}%
\Destination{79.226}{680.655}{errata.61}%
}

% Page 23 (body page 11)
\Page{%
\TextAt{2}{79.226}{796.943}{I.}%
\TextAt{2}{89.83}{796.943}{Algebraic}%
\TextAt{2}{138.648}{796.943}{Integers}%
\Vector{q 0 G 0 g 0.299 w 7.226 18.729 m 460.769 18.729 l S Q}%
\TextAt{19}{79.226}{761.124}{Proof.}%
\TextAt{8}{116}{761.124}{We}%
\TextAt{8}{136.157}{761.124}{first}%
\TextAt{8}{160.296}{761.124}{show}%
\TextAt{8}{189.311}{761.124}{that}%
\TextAt{8}{214.693}{761.124}{the}%
\TextAt{8}{234.862}{761.124}{bilinear}%
\TextAt{8}{277.148}{761.124}{form}%
\TextAt{9}{304.801}{761.124}{\Character{40}}%
\TextAt{13}{309.354}{761.124}{x\Character{59}}%
\TextAt{13}{321.226}{761.124}{y}%
\TextAt{9}{327.386}{761.124}{\Character{41}}%
\TextAt{9}{335.264}{761.124}{\Character{61}}%
\TextAt{8}{347.728}{761.124}{Tr}%
\TextAt{9}{360.735}{761.124}{\Character{40}}%
\TextAt{13}{365.288}{761.124}{xy}%
\TextAt{9}{378.077}{761.124}{\Character{41}}%
\TextAt{8}{386.539}{761.124}{is}%
\TextAt{8}{398.315}{761.124}{nondegenerate.}%
\TextAt{8}{480.302}{761.124}{Let}%
\TextAt{13}{501.284}{761.124}{\Character{18}}%
\TextAt{8}{510.975}{761.124}{be}%
\TextAt{8}{526.911}{761.124}{a}%
\TextAt{8}{79.226}{744.078}{primitive}%
\TextAt{8}{129.945}{744.078}{element}%
\TextAt{8}{173.825}{744.078}{for}%
\TextAt{13}{192.349}{744.078}{L}%
\TextAt{16}{204.723}{744.078}{\Character{106}}%
\TextAt{13}{212.453}{744.078}{K}%
\TextAt{8}{223.27}{744.078}{,}%
\TextAt{8}{231.218}{744.078}{i.e.,}%
\TextAt{13}{254.121}{744.078}{L}%
\TextAt{9}{266.495}{744.078}{\Character{61}}%
\TextAt{13}{280.009}{744.078}{K}%
\TextAt{9}{290.826}{744.078}{\Character{40}}%
\TextAt{13}{295.379}{744.078}{\Character{18}}%
\TextAt{9}{301.159}{744.078}{\Character{41}}%
\TextAt{8}{305.711}{744.078}{.}%
\TextAt{8}{316.075}{744.078}{Then}%
\TextAt{9}{347.271}{744.078}{1}%
\TextAt{13}{353.124}{744.078}{\Character{59}}%
\TextAt{13}{358.348}{744.078}{\Character{18}}%
\TextAt{13}{364.123}{744.078}{\Character{59}}%
\TextAt{13}{369.359}{744.078}{\Character{58}}%
\TextAt{13}{374.596}{744.078}{\Character{58}}%
\TextAt{13}{379.82}{744.078}{\Character{58}}%
\TextAt{13}{385.056}{744.078}{\Character{59}}%
\TextAt{13}{390.293}{744.078}{\Character{18}}%
\TextAt{21}{396.15}{748.416}{n}%
\TextAt{3}{401.288}{748.416}{\Character{0}}%
\TextAt{17}{407.875}{748.416}{1}%
\TextAt{8}{417.143}{744.078}{is}%
\TextAt{8}{429.546}{744.078}{a}%
\TextAt{8}{439.94}{744.078}{basis}%
\TextAt{8}{469.318}{744.078}{with}%
\TextAt{8}{496.617}{744.078}{respect}%
\TextAt{8}{79.226}{727.032}{to}%
\TextAt{8}{95.086}{727.032}{which}%
\TextAt{8}{130.122}{727.032}{the}%
\TextAt{8}{151.827}{727.032}{form}%
\TextAt{9}{181.017}{727.032}{\Character{40}}%
\TextAt{13}{185.569}{727.032}{x\Character{59}}%
\TextAt{13}{197.441}{727.032}{y}%
\TextAt{9}{203.602}{727.032}{\Character{41}}%
\TextAt{8}{213.601}{727.032}{is}%
\TextAt{8}{226.915}{727.032}{given}%
\TextAt{8}{259.033}{727.032}{by}%
\TextAt{8}{276.842}{727.032}{the}%
\TextAt{8}{298.547}{727.032}{matrix}%
\TextAt{13}{338.15}{727.032}{M}%
\TextAt{9}{356.685}{727.032}{\Character{61}}%
\TextAt{9}{371.737}{727.032}{\Character{40}}%
\TextAt{8}{376.303}{727.032}{Tr}%
\TextAt{21}{389.31}{725.172}{L}%
\TextAt{3}{395.069}{725.172}{\Character{106}}%
\TextAt{21}{397.422}{725.172}{K}%
\TextAt{9}{405.624}{727.032}{\Character{40}}%
\TextAt{13}{410.177}{727.032}{\Character{18}}%
\TextAt{21}{415.957}{731.37}{i}%
\TextAt{3}{418.84}{731.37}{\Character{0}}%
\TextAt{17}{425.427}{731.37}{1}%
\TextAt{13}{430.159}{727.032}{\Character{18}}%
\TextAt{21}{435.939}{731.37}{j}%
\TextAt{3}{439.823}{731.37}{\Character{0}}%
\TextAt{17}{446.41}{731.37}{1}%
\TextAt{9}{451.142}{727.032}{\Character{41}\Character{41}}%
\TextAt{21}{460.247}{725.238}{i\Character{59}j}%
\TextAt{17}{469.366}{725.238}{\Character{61}1}%
\TextAt{21}{480.187}{725.238}{\Character{59}\Character{58}\Character{58}\Character{58}\Character{59}n}%
\TextAt{8}{497.585}{727.032}{.}%
\TextAt{8}{510.681}{727.032}{It}%
\TextAt{8}{524.903}{727.032}{is}%
\TextAt{8}{79.226}{709.986}{nondegenerate}%
\TextAt{8}{156.624}{709.986}{because,}%
\TextAt{8}{203.177}{709.986}{for}%
\TextAt{13}{221.062}{709.986}{\Character{18}}%
\TextAt{21}{226.517}{708.193}{i}%
\TextAt{9}{233.22}{709.986}{\Character{61}}%
\TextAt{13}{245.645}{709.986}{\Character{27}}%
\TextAt{21}{252.298}{708.193}{i}%
\TextAt{13}{255.68}{709.986}{\Character{18}}%
\TextAt{8}{261.46}{709.986}{,}%
\TextAt{8}{268.609}{709.986}{we}%
\TextAt{8}{285.837}{709.986}{have}%
\TextAt{8}{187.397}{681.138}{det}%
\TextAt{9}{203.656}{681.138}{\Character{40}}%
\TextAt{13}{208.208}{681.138}{M}%
\TextAt{9}{220.782}{681.138}{\Character{41}}%
\TextAt{9}{228.637}{681.138}{\Character{61}}%
\TextAt{13}{241.08}{681.138}{d}%
\TextAt{9}{247.163}{681.138}{\Character{40}1}%
\TextAt{13}{257.568}{681.138}{\Character{59}}%
\TextAt{13}{262.792}{681.138}{\Character{18}}%
\TextAt{13}{268.567}{681.138}{\Character{59}}%
\TextAt{13}{273.803}{681.138}{\Character{58}}%
\TextAt{13}{279.04}{681.138}{\Character{58}}%
\TextAt{13}{284.264}{681.138}{\Character{58}}%
\TextAt{13}{289.5}{681.138}{\Character{59}}%
\TextAt{13}{294.737}{681.138}{\Character{18}}%
\TextAt{21}{300.594}{686.074}{n}%
\TextAt{3}{305.732}{686.074}{\Character{0}}%
\TextAt{17}{312.319}{686.074}{1}%
\TextAt{9}{317.051}{681.138}{\Character{41}}%
\TextAt{9}{324.906}{681.138}{\Character{61}}%
\TextAt{20}{337.349}{692.495}{\Character{89}}%
\TextAt{21}{338.311}{667.301}{i\Character{60}j}%
\TextAt{9}{352.626}{681.138}{\Character{40}}%
\TextAt{13}{357.178}{681.138}{\Character{18}}%
\TextAt{21}{362.633}{679.345}{i}%
\TextAt{16}{368.671}{681.138}{\Character{0}}%
\TextAt{13}{380.626}{681.138}{\Character{18}}%
\TextAt{21}{386.081}{679.345}{j}%
\TextAt{9}{390.463}{681.138}{\Character{41}}%
\TextAt{17}{395.016}{686.074}{2}%
\NativeMathAt{403.069}{681.138}{11.955}{\neq}%
\TextAt{9}{415.478}{681.138}{0}%
\TextAt{13}{421.347}{681.138}{\Character{58}}%
\TextAt{8}{79.226}{642.791}{If}%
\TextAt{13}{92.059}{642.791}{\Character{11}}%
\TextAt{17}{99.504}{640.998}{1}%
\TextAt{13}{104.237}{642.791}{\Character{59}}%
\TextAt{13}{109.461}{642.791}{\Character{58}}%
\TextAt{13}{114.698}{642.791}{\Character{58}}%
\TextAt{13}{119.922}{642.791}{\Character{58}}%
\TextAt{13}{125.159}{642.791}{\Character{59}}%
\TextAt{13}{130.395}{642.791}{\Character{11}}%
\TextAt{21}{137.903}{640.998}{n}%
\TextAt{8}{148.577}{642.791}{is}%
\TextAt{8}{161.482}{642.791}{an}%
\TextAt{8}{178.881}{642.791}{arbitrary}%
\TextAt{8}{229.791}{642.791}{basis}%
\TextAt{8}{259.672}{642.791}{of}%
\TextAt{13}{274.142}{642.791}{L}%
\TextAt{16}{287.371}{642.791}{\Character{106}}%
\TextAt{13}{295.956}{642.791}{K}%
\TextAt{8}{306.774}{642.791}{,}%
\TextAt{8}{315.349}{642.791}{then}%
\TextAt{8}{343.149}{642.791}{the}%
\TextAt{8}{364.446}{642.791}{bilinear}%
\TextAt{8}{407.86}{642.791}{form}%
\TextAt{9}{436.641}{642.791}{\Character{40}}%
\TextAt{13}{441.193}{642.791}{x\Character{59}}%
\TextAt{13}{453.065}{642.791}{y}%
\TextAt{9}{459.226}{642.791}{\Character{41}}%
\TextAt{8}{468.817}{642.791}{with}%
\TextAt{8}{496.617}{642.791}{respect}%
\TextAt{8}{79.226}{625.745}{to}%
\TextAt{8}{93.961}{625.745}{this}%
\TextAt{8}{117.209}{625.745}{basis}%
\TextAt{8}{146.374}{625.745}{is}%
\TextAt{8}{158.563}{625.745}{given}%
\TextAt{8}{189.557}{625.745}{by}%
\TextAt{8}{206.241}{625.745}{the}%
\TextAt{8}{226.822}{625.745}{matrix}%
\TextAt{13}{265.3}{625.745}{M}%
\TextAt{9}{281.919}{625.745}{\Character{61}}%
\TextAt{9}{295.058}{625.745}{\Character{40}}%
\TextAt{8}{299.621}{625.745}{Tr}%
\TextAt{21}{312.628}{623.886}{L}%
\TextAt{3}{318.387}{623.886}{\Character{106}}%
\TextAt{21}{320.739}{623.886}{K}%
\TextAt{9}{328.942}{625.745}{\Character{40}}%
\TextAt{13}{333.494}{625.745}{\Character{11}}%
\TextAt{21}{340.94}{623.952}{i}%
\TextAt{13}{344.321}{625.745}{\Character{11}}%
\TextAt{21}{351.767}{623.952}{j}%
\TextAt{9}{356.149}{625.745}{\Character{41}\Character{41}}%
\TextAt{8}{365.254}{625.745}{.}%
\TextAt{8}{374.977}{625.745}{From}%
\TextAt{8}{406.126}{625.745}{the}%
\TextAt{8}{426.708}{625.745}{above}%
\TextAt{8}{460.308}{625.745}{it}%
\TextAt{8}{472.437}{625.745}{follows}%
\TextAt{8}{511.298}{625.745}{that}%
\TextAt{13}{79.226}{608.699}{d}%
\TextAt{9}{85.309}{608.699}{\Character{40}}%
\TextAt{13}{89.861}{608.699}{\Character{11}}%
\TextAt{17}{97.307}{606.906}{1}%
\TextAt{13}{102.039}{608.699}{\Character{59}}%
\TextAt{13}{107.263}{608.699}{\Character{58}}%
\TextAt{13}{112.5}{608.699}{\Character{58}}%
\TextAt{13}{117.736}{608.699}{\Character{58}}%
\TextAt{13}{122.961}{608.699}{\Character{59}}%
\TextAt{13}{128.197}{608.699}{\Character{11}}%
\TextAt{21}{135.705}{606.906}{n}%
\TextAt{9}{141.342}{608.699}{\Character{41}}%
\TextAt{9}{149.197}{608.699}{\Character{61}}%
\TextAt{8}{161.64}{608.699}{det}%
\TextAt{9}{177.899}{608.699}{\Character{40}}%
\TextAt{13}{182.452}{608.699}{M}%
\TextAt{9}{195.025}{608.699}{\Character{41}}%
\NativeMathAt{202.899}{608.699}{11.955}{\neq}%
\TextAt{9}{215.309}{608.699}{0}%
\TextAt{8}{221.177}{608.699}{.}%
\Vector{q 0 G 0 g 0.399 w 452.334 -165.301 m 452.334 -157.231 l S Q}%
\Vector{q 0 G 0 g 0.399 w 452.534 -157.43 m 459.707 -157.43 l S Q}%
\Vector{q 0 G 0 g 0.399 w 452.534 -165.101 m 459.707 -165.101 l S Q}%
\Vector{q 0 G 0 g 0.399 w 459.906 -165.301 m 459.906 -157.231 l S Q}%
\TextAt{8}{97.159}{579.742}{After}%
\TextAt{8}{127.832}{579.742}{this}%
\TextAt{8}{150.77}{579.742}{review}%
\TextAt{8}{187.624}{579.742}{from}%
\TextAt{8}{215.38}{579.742}{the}%
\TextAt{8}{235.653}{579.742}{theory}%
\TextAt{8}{272.519}{579.742}{of}%
\TextAt{8}{285.964}{579.742}{fields,}%
\TextAt{8}{319.333}{579.742}{we}%
\TextAt{8}{336.676}{579.742}{return}%
\TextAt{8}{372.561}{579.742}{to}%
\TextAt{8}{386.987}{579.742}{the}%
\TextAt{8}{407.26}{579.742}{integrally}%
\TextAt{8}{459.416}{579.742}{closed}%
\TextAt{8}{494.059}{579.742}{integral}%
\TextAt{8}{79.226}{562.696}{domain}%
\TextAt{13}{121.366}{562.696}{A}%
\TextAt{8}{134.551}{562.696}{with}%
\TextAt{8}{161.723}{562.696}{field}%
\TextAt{8}{187.591}{562.696}{of}%
\TextAt{8}{201.433}{562.696}{fractions}%
\TextAt{13}{249.814}{562.696}{K}%
\TextAt{8}{260.631}{562.696}{,}%
\TextAt{8}{268.42}{562.696}{and}%
\TextAt{8}{291.695}{562.696}{to}%
\TextAt{8}{306.517}{562.696}{its}%
\TextAt{8}{323.348}{562.696}{integral}%
\TextAt{8}{366.468}{562.696}{closure}%
\TextAt{13}{406.061}{562.696}{B}%
\TextAt{8}{419.967}{562.696}{in}%
\TextAt{8}{434.132}{562.696}{the}%
\TextAt{8}{454.8}{562.696}{finite}%
\TextAt{8}{485.224}{562.696}{separable}%
\TextAt{8}{79.226}{545.65}{extension}%
\TextAt{13}{130.31}{545.65}{L}%
\TextAt{16}{141.596}{545.65}{\Character{106}}%
\TextAt{13}{148.237}{545.65}{K}%
\TextAt{8}{159.055}{545.65}{.}%
\TextAt{8}{167.276}{545.65}{If}%
\TextAt{13}{178.286}{545.65}{x}%
\NativeMathAt{188.259}{545.65}{11.955}{\in}%
\TextAt{13}{199.55}{545.65}{B}%
\TextAt{8}{212.262}{545.65}{is}%
\TextAt{8}{223.344}{545.65}{an}%
\TextAt{8}{238.922}{545.65}{integral}%
\TextAt{8}{280.848}{545.65}{element}%
\TextAt{8}{323.408}{545.65}{of}%
\TextAt{13}{336.056}{545.65}{L}%
\TextAt{8}{344.021}{545.65}{,}%
\TextAt{8}{350.625}{545.65}{then}%
\TextAt{8}{376.603}{545.65}{all}%
\TextAt{8}{392.18}{545.65}{of}%
\TextAt{8}{404.829}{545.65}{its}%
\TextAt{8}{420.466}{545.65}{conjugates}%
\TextAt{13}{477.408}{545.65}{\Character{27}}%
\TextAt{13}{484.474}{545.65}{x}%
\TextAt{8}{494.358}{545.65}{are}%
\TextAt{8}{513.187}{545.65}{also}%
\TextAt{8}{79.226}{528.604}{integral.}%
\TextAt{8}{127.855}{528.604}{Taking}%
\TextAt{8}{167.366}{528.604}{into}%
\TextAt{8}{191.599}{528.604}{account}%
\TextAt{8}{235.343}{528.604}{that}%
\TextAt{13}{261.202}{528.604}{A}%
\TextAt{8}{274.364}{528.604}{is}%
\TextAt{8}{286.618}{528.604}{integrally}%
\TextAt{8}{339.149}{528.604}{closed,}%
\TextAt{8}{377.539}{528.604}{i.e.,}%
\TextAt{13}{400.257}{528.604}{A}%
\TextAt{9}{413.188}{528.604}{\Character{61}}%
\TextAt{13}{426.448}{528.604}{B}%
\TextAt{16}{438.935}{528.604}{\Character{92}}%
\TextAt{13}{449.896}{528.604}{K}%
\TextAt{8}{460.713}{528.604}{,}%
\TextAt{8}{468.475}{528.604}{(2.6)}%
\TextAt{8}{496.94}{528.604}{implies}%
\TextAt{8}{79.226}{511.558}{that}%
\TextAt{8}{240.817}{494.512}{Tr}%
\TextAt{21}{253.825}{492.652}{L}%
\TextAt{3}{259.584}{492.652}{\Character{106}}%
\TextAt{21}{261.936}{492.652}{K}%
\TextAt{9}{270.139}{494.512}{\Character{40}}%
\TextAt{13}{274.691}{494.512}{x}%
\TextAt{9}{281.343}{494.512}{\Character{41}}%
\TextAt{13}{285.895}{494.512}{\Character{59}}%
\TextAt{13}{303.075}{494.512}{N}%
\TextAt{21}{312.469}{492.652}{L}%
\TextAt{3}{318.228}{492.652}{\Character{106}}%
\TextAt{21}{320.58}{492.652}{K}%
\TextAt{9}{328.783}{494.512}{\Character{40}}%
\TextAt{13}{333.335}{494.512}{x}%
\TextAt{9}{339.987}{494.512}{\Character{41}}%
\NativeMathAt{347.86}{494.512}{11.955}{\in}%
\TextAt{13}{359.151}{494.512}{A\Character{58}}%
\TextAt{8}{79.226}{471.056}{Furthermore,}%
\TextAt{8}{150.276}{471.056}{for}%
\TextAt{8}{168.161}{471.056}{the}%
\TextAt{12}{188.317}{471.056}{group}%
\TextAt{12}{226.705}{471.056}{of}%
\TextAt{12}{242.032}{471.056}{units}%
\TextAt{8}{275.159}{471.056}{of}%
\TextAt{13}{288.489}{471.056}{B}%
\TextAt{8}{301.883}{471.056}{over}%
\TextAt{13}{326.929}{471.056}{A}%
\TextAt{8}{335.705}{471.056}{,}%
\TextAt{8}{342.854}{471.056}{we}%
\TextAt{8}{360.081}{471.056}{obtain}%
\TextAt{8}{396.509}{471.056}{the}%
\TextAt{8}{416.665}{471.056}{relation}%
\TextAt{13}{225.609}{442.208}{x}%
\NativeMathAt{235.582}{442.208}{11.955}{\in}%
\TextAt{13}{246.873}{442.208}{B}%
\TextAt{3}{256.369}{447.144}{\Character{3}}%
\NativeArrowAt{276.377}{442.208}{11.955}{21.926}{\Longleftrightarrow}%
\TextAt{13}{313.571}{442.208}{N}%
\TextAt{21}{322.945}{440.349}{L}%
\TextAt{3}{328.704}{440.349}{\Character{106}}%
\TextAt{21}{331.056}{440.349}{K}%
\TextAt{9}{339.259}{442.208}{\Character{40}}%
\TextAt{13}{343.811}{442.208}{x}%
\TextAt{9}{350.463}{442.208}{\Character{41}}%
\NativeMathAt{358.337}{442.208}{11.955}{\in}%
\TextAt{13}{369.628}{442.208}{A}%
\TextAt{3}{378.403}{447.144}{\Character{3}}%
\TextAt{13}{383.135}{442.208}{\Character{58}}%
\TextAt{8}{79.226}{413.36}{For}%
\TextAt{8}{100.266}{413.36}{if}%
\TextAt{13}{111.061}{413.36}{aN}%
\TextAt{21}{126.579}{411.501}{L}%
\TextAt{3}{132.338}{411.501}{\Character{106}}%
\TextAt{21}{134.691}{411.501}{K}%
\TextAt{9}{142.893}{413.36}{\Character{40}}%
\TextAt{13}{147.446}{413.36}{x}%
\TextAt{9}{154.098}{413.36}{\Character{41}}%
\TextAt{9}{162.072}{413.36}{\Character{61}}%
\TextAt{9}{174.613}{413.36}{1}%
\TextAt{8}{180.49}{413.36}{,}%
\TextAt{13}{187.727}{413.36}{a}%
\NativeMathAt{197.314}{413.36}{11.955}{\in}%
\TextAt{13}{208.725}{413.36}{A}%
\TextAt{8}{217.5}{413.36}{,}%
\TextAt{8}{224.738}{413.36}{then}%
\TextAt{9}{251.468}{413.36}{1}%
\TextAt{9}{260.745}{413.36}{\Character{61}}%
\TextAt{13}{273.308}{413.36}{a}%
\TextAt{20}{281.446}{422.327}{\Character{81}}%
\TextAt{21}{292.737}{409.873}{\Character{27}}%
\TextAt{13}{300.353}{413.36}{\Character{27}}%
\TextAt{13}{307.418}{413.36}{x}%
\TextAt{9}{317.528}{413.36}{\Character{61}}%
\TextAt{13}{330.074}{413.36}{y}%
\TextAt{13}{336.195}{413.36}{x}%
\TextAt{8}{346.831}{413.36}{for}%
\TextAt{8}{364.787}{413.36}{some}%
\TextAt{13}{394.183}{413.36}{y}%
\NativeMathAt{403.761}{413.36}{11.955}{\in}%
\TextAt{13}{415.172}{413.36}{B}%
\TextAt{8}{424.669}{413.36}{.}%
\TextAt{8}{433.329}{413.36}{The}%
\TextAt{8}{457.453}{413.36}{discriminant}%
\TextAt{8}{524.903}{413.36}{is}%
\TextAt{8}{79.226}{396.314}{often}%
\TextAt{8}{108.815}{396.314}{useful}%
\TextAt{8}{142.361}{396.314}{because}%
\TextAt{8}{185.663}{396.314}{of}%
\TextAt{8}{198.993}{396.314}{the}%
\TextAt{8}{219.149}{396.314}{following}%
\TextAt{12}{79.226}{370.346}{Lemma}%
\TextAt{12}{126.003}{370.346}{(2.9).}%
\TextAt{8}{160.546}{370.346}{Let}%
\TextAt{13}{180.743}{370.346}{\Character{11}}%
\TextAt{17}{188.189}{368.553}{1}%
\TextAt{13}{192.921}{370.346}{\Character{59}}%
\TextAt{13}{198.145}{370.346}{\Character{58}}%
\TextAt{13}{203.382}{370.346}{\Character{58}}%
\TextAt{13}{208.618}{370.346}{\Character{58}}%
\TextAt{13}{213.843}{370.346}{\Character{59}}%
\TextAt{13}{219.079}{370.346}{\Character{11}}%
\TextAt{21}{226.588}{368.553}{n}%
\TextAt{8}{235.349}{370.346}{be}%
\TextAt{8}{250.501}{370.346}{a}%
\TextAt{8}{259.484}{370.346}{basis}%
\TextAt{8}{287.452}{370.346}{of}%
\TextAt{13}{300.01}{370.346}{L}%
\TextAt{16}{311.295}{370.346}{\Character{106}}%
\TextAt{13}{317.937}{370.346}{K}%
\TextAt{8}{331.879}{370.346}{which}%
\TextAt{8}{364.593}{370.346}{is}%
\TextAt{8}{375.585}{370.346}{contained}%
\TextAt{8}{427.822}{370.346}{in}%
\TextAt{13}{440.702}{370.346}{B}%
\TextAt{8}{450.199}{370.346}{,}%
\TextAt{8}{456.73}{370.346}{of}%
\TextAt{8}{469.288}{370.346}{discriminant}%
\TextAt{13}{79.226}{353.3}{d}%
\TextAt{9}{88.63}{353.3}{\Character{61}}%
\TextAt{13}{101.055}{353.3}{d}%
\TextAt{9}{107.138}{353.3}{\Character{40}}%
\TextAt{13}{111.69}{353.3}{\Character{11}}%
\TextAt{17}{119.136}{351.507}{1}%
\TextAt{13}{123.868}{353.3}{\Character{59}}%
\TextAt{13}{129.092}{353.3}{\Character{58}}%
\TextAt{13}{134.329}{353.3}{\Character{58}}%
\TextAt{13}{139.565}{353.3}{\Character{58}}%
\TextAt{13}{144.79}{353.3}{\Character{59}}%
\TextAt{13}{150.026}{353.3}{\Character{11}}%
\TextAt{21}{157.534}{351.507}{n}%
\TextAt{9}{163.171}{353.3}{\Character{41}}%
\TextAt{8}{167.723}{353.3}{.}%
\TextAt{8}{176.171}{353.3}{Then}%
\TextAt{8}{206.729}{353.3}{one}%
\TextAt{8}{228.188}{353.3}{has}%
\TextAt{13}{245.815}{324.452}{dB}%
\NativeMathAt{264.715}{324.452}{11.955}{\subseteq}%
\TextAt{13}{277.334}{324.452}{A\Character{11}}%
\TextAt{17}{293.555}{322.659}{1}%
\TextAt{9}{300.944}{324.452}{\Character{43}}%
\TextAt{16}{312.706}{324.452}{\Character{1}}%
\TextAt{16}{318.002}{324.452}{\Character{1}}%
\TextAt{16}{323.31}{324.452}{\Character{1}}%
\TextAt{9}{329.31}{324.452}{\Character{43}}%
\TextAt{13}{341.071}{324.452}{A\Character{11}}%
\TextAt{21}{357.292}{322.659}{n}%
\TextAt{13}{362.929}{324.452}{\Character{58}}%
\TextAt{19}{79.226}{277.759}{Proof.}%
\TextAt{8}{116}{277.759}{If}%
\TextAt{13}{127.74}{277.759}{\Character{11}}%
\TextAt{9}{138.663}{277.759}{\Character{61}}%
\TextAt{13}{151.169}{277.759}{a}%
\TextAt{17}{157.314}{275.966}{1}%
\TextAt{13}{162.047}{277.759}{\Character{11}}%
\TextAt{17}{169.492}{275.966}{1}%
\TextAt{9}{176.914}{277.759}{\Character{43}}%
\TextAt{16}{188.707}{277.759}{\Character{1}}%
\TextAt{16}{194.003}{277.759}{\Character{1}}%
\TextAt{16}{199.311}{277.759}{\Character{1}}%
\TextAt{9}{205.344}{277.759}{\Character{43}}%
\TextAt{13}{217.137}{277.759}{a}%
\TextAt{21}{223.282}{275.966}{n}%
\TextAt{13}{228.919}{277.759}{\Character{11}}%
\TextAt{21}{236.364}{275.966}{n}%
\NativeMathAt{245.402}{277.759}{11.955}{\in}%
\TextAt{13}{256.774}{277.759}{B}%
\TextAt{8}{266.27}{277.759}{,}%
\TextAt{13}{273.479}{277.759}{a}%
\TextAt{21}{279.624}{275.966}{i}%
\NativeMathAt{286.406}{277.759}{11.955}{\in}%
\TextAt{13}{297.778}{277.759}{K}%
\TextAt{8}{308.596}{277.759}{,}%
\TextAt{8}{315.804}{277.759}{then}%
\TextAt{8}{342.512}{277.759}{the}%
\TextAt{13}{362.716}{277.759}{a}%
\TextAt{21}{368.861}{275.966}{j}%
\TextAt{8}{377.187}{277.759}{are}%
\TextAt{8}{396.746}{277.759}{a}%
\TextAt{8}{406.549}{277.759}{solution}%
\TextAt{8}{450.89}{277.759}{of}%
\TextAt{8}{464.267}{277.759}{the}%
\TextAt{8}{484.471}{277.759}{system}%
\TextAt{8}{523.337}{277.759}{of}%
\TextAt{8}{79.226}{260.713}{linear}%
\TextAt{8}{111.744}{260.713}{equations}%
\TextAt{8}{222.677}{243.667}{Tr}%
\TextAt{21}{235.684}{241.807}{L}%
\TextAt{3}{241.443}{241.807}{\Character{106}}%
\TextAt{21}{243.796}{241.807}{K}%
\TextAt{9}{251.998}{243.667}{\Character{40}}%
\TextAt{13}{256.551}{243.667}{\Character{11}}%
\TextAt{21}{263.996}{241.874}{i}%
\TextAt{13}{267.378}{243.667}{\Character{11}}%
\TextAt{9}{274.899}{243.667}{\Character{41}}%
\TextAt{9}{282.754}{243.667}{\Character{61}}%
\TextAt{20}{295.198}{255.025}{\Character{88}}%
\TextAt{21}{301.89}{229.83}{j}%
\TextAt{8}{314.459}{243.667}{Tr}%
\TextAt{21}{327.466}{241.807}{L}%
\TextAt{3}{333.225}{241.807}{\Character{106}}%
\TextAt{21}{335.578}{241.807}{K}%
\TextAt{9}{343.78}{243.667}{\Character{40}}%
\TextAt{13}{348.333}{243.667}{\Character{11}}%
\TextAt{21}{355.778}{241.874}{i}%
\TextAt{13}{359.16}{243.667}{\Character{11}}%
\TextAt{21}{366.605}{241.874}{j}%
\TextAt{9}{370.987}{243.667}{\Character{41}}%
\TextAt{13}{375.54}{243.667}{a}%
\TextAt{21}{381.685}{241.874}{j}%
\TextAt{13}{386.067}{243.667}{\Character{59}}%
\TextAt{8}{79.226}{210.712}{and,}%
\TextAt{8}{105.464}{210.712}{as}%
\TextAt{8}{120.013}{210.712}{Tr}%
\TextAt{21}{133.02}{208.853}{L}%
\TextAt{3}{138.779}{208.853}{\Character{106}}%
\TextAt{21}{141.132}{208.853}{K}%
\TextAt{9}{149.334}{210.712}{\Character{40}}%
\TextAt{13}{153.887}{210.712}{\Character{11}}%
\TextAt{21}{161.332}{208.919}{i}%
\TextAt{13}{164.714}{210.712}{\Character{11}}%
\TextAt{9}{172.235}{210.712}{\Character{41}}%
\NativeMathAt{180.413}{210.712}{11.955}{\in}%
\TextAt{13}{192.009}{210.712}{A}%
\TextAt{8}{200.784}{210.712}{,}%
\TextAt{8}{208.157}{210.712}{they}%
\TextAt{8}{234.673}{210.712}{are}%
\TextAt{8}{254.363}{210.712}{given}%
\TextAt{8}{285.111}{210.712}{as}%
\TextAt{8}{299.66}{210.712}{the}%
\TextAt{8}{319.995}{210.712}{quotient}%
\TextAt{8}{366.357}{210.712}{of}%
\TextAt{8}{379.866}{210.712}{an}%
\TextAt{8}{396.303}{210.712}{element}%
\TextAt{8}{439.724}{210.712}{of}%
\TextAt{13}{453.233}{210.712}{A}%
\TextAt{8}{466.085}{210.712}{by}%
\TextAt{8}{482.522}{210.712}{the}%
\TextAt{8}{502.858}{210.712}{deter-}%
\TextAt{8}{79.226}{193.666}{minant}%
\TextAt{8}{119.228}{193.666}{det}%
\TextAt{9}{135.487}{193.666}{\Character{40}}%
\TextAt{8}{140.04}{193.666}{Tr}%
\TextAt{21}{153.047}{191.807}{L}%
\TextAt{3}{158.806}{191.807}{\Character{106}}%
\TextAt{21}{161.158}{191.807}{K}%
\TextAt{9}{169.361}{193.666}{\Character{40}}%
\TextAt{13}{173.913}{193.666}{\Character{11}}%
\TextAt{21}{181.359}{191.873}{i}%
\TextAt{13}{184.74}{193.666}{\Character{11}}%
\TextAt{21}{192.186}{191.873}{j}%
\TextAt{9}{196.568}{193.666}{\Character{41}\Character{41}}%
\TextAt{9}{208.966}{193.666}{\Character{61}}%
\TextAt{13}{221.419}{193.666}{d}%
\TextAt{8}{227.501}{193.666}{.}%
\TextAt{8}{235.95}{193.666}{Therefore}%
\TextAt{13}{288.947}{193.666}{da}%
\TextAt{21}{301.175}{191.873}{j}%
\NativeMathAt{308.878}{193.666}{11.955}{\in}%
\TextAt{13}{320.169}{193.666}{A}%
\TextAt{8}{328.944}{193.666}{,}%
\TextAt{8}{336.093}{193.666}{and}%
\TextAt{8}{358.856}{193.666}{thus}%
\TextAt{13}{247.467}{164.818}{d\Character{11}}%
\NativeMathAt{264.392}{164.819}{11.955}{\in}%
\TextAt{13}{275.683}{164.818}{A\Character{11}}%
\TextAt{17}{291.904}{163.025}{1}%
\TextAt{9}{299.293}{164.819}{\Character{43}}%
\TextAt{16}{311.054}{164.819}{\Character{1}}%
\TextAt{16}{316.35}{164.819}{\Character{1}}%
\TextAt{16}{321.658}{164.819}{\Character{1}}%
\TextAt{9}{327.659}{164.819}{\Character{43}}%
\TextAt{13}{339.42}{164.819}{A\Character{11}}%
\TextAt{21}{355.641}{163.025}{n}%
\TextAt{13}{361.277}{164.819}{\Character{58}}%
\Vector{q 0 G 0 g 0.399 w 452.334 -638.029 m 452.334 -629.96 l S Q}%
\Vector{q 0 G 0 g 0.399 w 452.534 -630.159 m 459.707 -630.159 l S Q}%
\Vector{q 0 G 0 g 0.399 w 452.534 -637.83 m 459.707 -637.83 l S Q}%
\Vector{q 0 G 0 g 0.399 w 459.906 -638.029 m 459.906 -629.96 l S Q}%
\TextAt{8}{97.159}{107.013}{A}%
\TextAt{8}{109.552}{107.013}{system}%
\TextAt{8}{148.092}{107.013}{of}%
\TextAt{8}{161.143}{107.013}{elements}%
\TextAt{13}{208.72}{107.013}{\Character{33}}%
\TextAt{17}{216.009}{105.22}{1}%
\TextAt{13}{220.742}{107.013}{\Character{59}}%
\TextAt{13}{225.966}{107.013}{\Character{58}}%
\TextAt{13}{231.203}{107.013}{\Character{58}}%
\TextAt{13}{236.427}{107.013}{\Character{58}}%
\TextAt{13}{241.664}{107.013}{\Character{59}}%
\TextAt{13}{246.9}{107.013}{\Character{33}}%
\TextAt{21}{254.252}{105.22}{n}%
\NativeMathAt{263.209}{107.013}{11.955}{\in}%
\TextAt{13}{274.5}{107.013}{B}%
\TextAt{8}{287.615}{107.013}{such}%
\TextAt{8}{313.733}{107.013}{that}%
\TextAt{8}{338.823}{107.013}{each}%
\TextAt{13}{364.881}{107.013}{b}%
\NativeMathAt{373.179}{107.013}{11.955}{\in}%
\TextAt{13}{384.47}{107.013}{B}%
\TextAt{8}{397.585}{107.013}{can}%
\TextAt{8}{418.765}{107.013}{be}%
\TextAt{8}{434.41}{107.013}{written}%
\TextAt{8}{475.102}{107.013}{uniquely}%
\TextAt{8}{522.297}{107.013}{as}%
\TextAt{8}{79.226}{89.967}{a}%
\TextAt{8}{88.982}{89.967}{linear}%
\TextAt{8}{121.5}{89.967}{combination}%
\TextAt{13}{250.441}{72.921}{b}%
\TextAt{9}{258.739}{72.921}{\Character{61}}%
\TextAt{13}{271.165}{72.921}{a}%
\TextAt{17}{277.31}{71.128}{1}%
\TextAt{13}{282.042}{72.921}{\Character{33}}%
\TextAt{17}{289.331}{71.128}{1}%
\TextAt{9}{296.72}{72.921}{\Character{43}}%
\TextAt{16}{308.482}{72.921}{\Character{1}}%
\TextAt{16}{313.778}{72.921}{\Character{1}}%
\TextAt{16}{319.086}{72.921}{\Character{1}}%
\TextAt{9}{325.086}{72.921}{\Character{43}}%
\TextAt{13}{336.847}{72.921}{a}%
\TextAt{21}{342.992}{71.128}{n}%
\TextAt{13}{348.629}{72.921}{\Character{33}}%
\TextAt{21}{355.918}{71.128}{n}%
\TextAt{8}{300.14}{43.033}{11}%
\Destination{78.23}{809.815}{nav.page.11}%
\Destination{79.36}{594.755}{nav.print.12}%
\Destination{79.36}{384.355}{nav.number.1.2.9}%
\Hotlink{467.975}{526.355}{25.078}{11.047}{nav.number.1.2.6}%
}

% Page 24 (body page 12)
\Page{%
\TextAt{2}{79.226}{796.943}{I.}%
\TextAt{2}{89.83}{796.943}{Algebraic}%
\TextAt{2}{138.648}{796.943}{Integers}%
\Vector{q 0 G 0 g 0.299 w 7.226 18.729 m 460.769 18.729 l S Q}%
\TextAt{8}{79.226}{761.124}{with}%
\TextAt{8}{106.307}{761.124}{coefficients}%
\TextAt{13}{165.966}{761.124}{a}%
\TextAt{21}{172.111}{759.33}{i}%
\NativeMathAt{179.53}{761.124}{11.955}{\in}%
\TextAt{13}{191.538}{761.124}{A}%
\TextAt{8}{200.314}{761.124}{,}%
\TextAt{8}{207.989}{761.124}{is}%
\TextAt{8}{220.174}{761.124}{called}%
\TextAt{8}{253.758}{761.124}{an}%
\TextAt{12}{270.438}{761.124}{integral}%
\TextAt{12}{320.13}{761.124}{basis}%
\TextAt{8}{352.818}{761.124}{of}%
\TextAt{13}{366.569}{761.124}{B}%
\TextAt{8}{380.383}{761.124}{over}%
\TextAt{13}{405.85}{761.124}{A}%
\TextAt{8}{418.944}{761.124}{(or:}%
\TextAt{8}{443.202}{761.124}{an}%
\TextAt{13}{459.882}{761.124}{A}%
\TextAt{8}{468.657}{761.124}{-basis}%
\TextAt{8}{501.716}{761.124}{of}%
\TextAt{13}{515.467}{761.124}{B}%
\TextAt{8}{524.963}{761.124}{).}%
\TextAt{8}{79.226}{744.078}{Since}%
\TextAt{8}{109.977}{744.078}{such}%
\TextAt{8}{136.568}{744.078}{an}%
\TextAt{8}{153.02}{744.078}{integral}%
\TextAt{8}{195.822}{744.078}{basis}%
\TextAt{8}{224.756}{744.078}{is}%
\TextAt{8}{236.714}{744.078}{automatically}%
\TextAt{8}{310.742}{744.078}{a}%
\TextAt{8}{320.691}{744.078}{basis}%
\TextAt{8}{349.625}{744.078}{of}%
\TextAt{13}{363.149}{744.078}{L}%
\TextAt{16}{374.764}{744.078}{\Character{106}}%
\TextAt{13}{381.736}{744.078}{K}%
\TextAt{8}{392.553}{744.078}{,}%
\TextAt{8}{399.944}{744.078}{its}%
\TextAt{8}{416.457}{744.078}{length}%
\TextAt{13}{452.42}{744.078}{n}%
\TextAt{8}{463.499}{744.078}{always}%
\TextAt{8}{501.16}{744.078}{equals}%
\TextAt{8}{79.226}{727.032}{the}%
\TextAt{8}{99.525}{727.032}{degree}%
\TextAt{9}{136.084}{727.032}{\Character{91}}%
\TextAt{13}{139.335}{727.032}{L}%
\TextAt{9}{150.864}{727.032}{\Character{58}}%
\TextAt{13}{157.68}{727.032}{K}%
\TextAt{9}{168.497}{727.032}{\Character{93}}%
\TextAt{8}{175.789}{727.032}{of}%
\TextAt{8}{189.262}{727.032}{the}%
\TextAt{8}{209.561}{727.032}{field}%
\TextAt{8}{235.061}{727.032}{extension.}%
\TextAt{8}{291.806}{727.032}{The}%
\TextAt{8}{316.002}{727.032}{existence}%
\TextAt{8}{365.95}{727.032}{of}%
\TextAt{8}{379.423}{727.032}{an}%
\TextAt{8}{395.825}{727.032}{integral}%
\TextAt{8}{438.576}{727.032}{basis}%
\TextAt{8}{467.459}{727.032}{signifies}%
\TextAt{8}{511.298}{727.032}{that}%
\TextAt{13}{79.226}{709.986}{B}%
\TextAt{8}{93.452}{709.986}{is}%
\TextAt{8}{106.049}{709.986}{a}%
\TextAt{8}{116.637}{709.986}{free}%
\TextAt{13}{139.897}{709.986}{A}%
\TextAt{8}{148.672}{709.986}{-module}%
\TextAt{8}{194.695}{709.986}{of}%
\TextAt{8}{208.858}{709.986}{rank}%
\TextAt{13}{236.685}{709.986}{n}%
\TextAt{9}{248.412}{709.986}{\Character{61}}%
\TextAt{9}{262.244}{709.986}{\Character{91}}%
\TextAt{13}{265.508}{709.986}{L}%
\TextAt{9}{278.212}{709.986}{\Character{58}}%
\TextAt{13}{286.203}{709.986}{K}%
\TextAt{9}{297.021}{709.986}{\Character{93}}%
\TextAt{8}{300.272}{709.986}{.}%
\TextAt{8}{311.218}{709.986}{In}%
\TextAt{8}{326.672}{709.986}{general,}%
\TextAt{8}{371.289}{709.986}{such}%
\TextAt{8}{398.519}{709.986}{an}%
\TextAt{8}{415.61}{709.986}{integral}%
\TextAt{8}{459.051}{709.986}{basis}%
\TextAt{8}{488.623}{709.986}{does}%
\TextAt{8}{515.853}{709.986}{not}%
\TextAt{8}{79.226}{692.94}{exist.}%
\TextAt{8}{111.432}{692.94}{If,}%
\TextAt{8}{126.267}{692.94}{however,}%
\TextAt{13}{174.289}{692.94}{A}%
\TextAt{8}{186.825}{692.94}{is}%
\TextAt{8}{198.453}{692.94}{a}%
\TextAt{8}{208.072}{692.94}{principal}%
\TextAt{8}{256.713}{692.94}{ideal}%
\TextAt{8}{284.539}{692.94}{domain,}%
\TextAt{8}{329.31}{692.94}{then}%
\TextAt{8}{355.834}{692.94}{one}%
\TextAt{8}{377.157}{692.94}{has}%
\TextAt{8}{397.894}{692.94}{the}%
\TextAt{8}{417.914}{692.94}{following}%
\TextAt{8}{467.212}{692.94}{more}%
\TextAt{8}{496.342}{692.94}{general}%
\TextAt{12}{79.226}{667.983}{Proposition}%
\TextAt{12}{151.856}{667.983}{(2.10).}%
\TextAt{8}{193.832}{667.983}{If}%
\TextAt{13}{205.454}{667.983}{L}%
\TextAt{16}{216.739}{667.983}{\Character{106}}%
\TextAt{13}{223.381}{667.983}{K}%
\TextAt{8}{238.025}{667.983}{is}%
\TextAt{8}{249.718}{667.983}{separable}%
\TextAt{8}{301.091}{667.983}{and}%
\TextAt{13}{323.783}{667.983}{A}%
\TextAt{8}{336.385}{667.983}{is}%
\TextAt{8}{348.078}{667.983}{a}%
\TextAt{8}{357.763}{667.983}{principal}%
\TextAt{8}{406.47}{667.983}{ideal}%
\TextAt{8}{434.362}{667.983}{domain,}%
\TextAt{8}{479.185}{667.983}{then}%
\TextAt{8}{505.775}{667.983}{every}%
\TextAt{8}{79.226}{650.937}{finitely}%
\TextAt{8}{120.07}{650.937}{generated}%
\TextAt{13}{174.901}{650.937}{B}%
\TextAt{8}{184.397}{650.937}{-submodule}%
\TextAt{13}{248.709}{650.937}{M}%
\NativeMathAt{267.158}{650.937}{11.955}{\neq}%
\TextAt{9}{282.126}{650.937}{0}%
\TextAt{8}{293.388}{650.937}{of}%
\TextAt{13}{308.217}{650.937}{L}%
\TextAt{8}{321.578}{650.937}{is}%
\TextAt{8}{334.841}{650.937}{a}%
\TextAt{8}{346.096}{650.937}{free}%
\TextAt{13}{370.023}{650.937}{A}%
\TextAt{8}{378.798}{650.937}{-module}%
\TextAt{8}{425.488}{650.937}{of}%
\TextAt{8}{440.317}{650.937}{rank}%
\TextAt{9}{468.811}{650.937}{\Character{91}}%
\TextAt{13}{472.063}{650.937}{L}%
\TextAt{9}{485.903}{650.937}{\Character{58}}%
\TextAt{13}{495.03}{650.937}{K}%
\TextAt{9}{505.848}{650.937}{\Character{93}}%
\TextAt{8}{509.1}{650.937}{.}%
\TextAt{8}{522.046}{650.937}{In}%
\TextAt{8}{79.226}{633.891}{particular,}%
\TextAt{13}{136.467}{633.891}{B}%
\TextAt{8}{149.861}{633.891}{admits}%
\TextAt{8}{188.297}{633.891}{an}%
\TextAt{8}{204.556}{633.891}{integral}%
\TextAt{8}{247.164}{633.891}{basis}%
\TextAt{8}{275.905}{633.891}{over}%
\TextAt{13}{300.951}{633.891}{A}%
\TextAt{8}{309.726}{633.891}{.}%
\TextAt{19}{79.226}{601.024}{Proof.}%
\TextAt{8}{116}{601.024}{Let}%
\TextAt{13}{137.933}{601.024}{M}%
\NativeMathAt{155.469}{601.024}{11.955}{\neq}%
\TextAt{9}{169.528}{601.024}{0}%
\TextAt{8}{180.249}{601.024}{be}%
\TextAt{8}{197.136}{601.024}{a}%
\TextAt{8}{207.855}{601.024}{finitely}%
\TextAt{8}{248.163}{601.024}{generated}%
\TextAt{13}{302.458}{601.024}{B}%
\TextAt{8}{311.954}{601.024}{-submodule}%
\TextAt{8}{375.73}{601.024}{of}%
\TextAt{13}{390.023}{601.024}{L}%
\TextAt{8}{402.848}{601.024}{and}%
\TextAt{13}{426.574}{601.024}{\Character{11}}%
\TextAt{17}{434.019}{599.23}{1}%
\TextAt{13}{438.752}{601.024}{\Character{59}}%
\TextAt{13}{443.976}{601.024}{\Character{58}}%
\TextAt{13}{449.213}{601.024}{\Character{58}}%
\TextAt{13}{454.437}{601.024}{\Character{58}}%
\TextAt{13}{459.674}{601.024}{\Character{59}}%
\TextAt{13}{464.91}{601.024}{\Character{11}}%
\TextAt{21}{472.418}{599.23}{n}%
\TextAt{8}{482.915}{601.024}{a}%
\TextAt{8}{493.633}{601.024}{basis}%
\TextAt{8}{523.337}{601.024}{of}%
\TextAt{13}{79.226}{583.978}{L}%
\TextAt{16}{90.694}{583.978}{\Character{106}}%
\TextAt{13}{97.519}{583.978}{K}%
\TextAt{8}{108.337}{583.978}{.}%
\TextAt{8}{117.107}{583.978}{Multiplying}%
\TextAt{8}{180.947}{583.978}{by}%
\TextAt{8}{197.314}{583.978}{an}%
\TextAt{8}{213.68}{583.978}{element}%
\TextAt{8}{257.029}{583.978}{of}%
\TextAt{13}{270.466}{583.978}{A}%
\TextAt{8}{279.242}{583.978}{,}%
\TextAt{8}{286.525}{583.978}{we}%
\TextAt{8}{303.86}{583.978}{may}%
\TextAt{8}{329.336}{583.978}{arrange}%
\TextAt{8}{371.729}{583.978}{for}%
\TextAt{8}{389.721}{583.978}{the}%
\TextAt{13}{409.985}{583.978}{\Character{11}}%
\TextAt{21}{417.43}{582.185}{i}%
\TextAt{8}{424.816}{583.978}{to}%
\TextAt{8}{439.234}{583.978}{lie}%
\TextAt{8}{454.943}{583.978}{in}%
\TextAt{13}{468.703}{583.978}{B}%
\TextAt{8}{478.199}{583.978}{.}%
\TextAt{8}{486.97}{583.978}{By}%
\TextAt{8}{505.44}{583.978}{(2.9),}%
\TextAt{8}{79.226}{566.932}{we}%
\TextAt{8}{95.966}{566.932}{then}%
\TextAt{8}{122.139}{566.932}{have}%
\TextAt{13}{148.646}{566.932}{dB}%
\NativeMathAt{167.546}{566.932}{11.955}{\subseteq}%
\TextAt{13}{180.165}{566.932}{A\Character{11}}%
\TextAt{17}{196.386}{565.139}{1}%
\TextAt{9}{202.778}{566.932}{\Character{43}}%
\TextAt{16}{213.543}{566.932}{\Character{1}}%
\TextAt{16}{218.839}{566.932}{\Character{1}}%
\TextAt{16}{224.147}{566.932}{\Character{1}}%
\TextAt{9}{229.15}{566.932}{\Character{43}}%
\TextAt{13}{239.915}{566.932}{A\Character{11}}%
\TextAt{21}{256.136}{565.139}{n}%
\TextAt{8}{261.772}{566.932}{,}%
\TextAt{8}{268.531}{566.932}{in}%
\TextAt{8}{281.696}{566.932}{particular,}%
\TextAt{8}{338.548}{566.932}{rank}%
\TextAt{9}{361.645}{566.932}{\Character{40}}%
\TextAt{13}{366.198}{566.932}{B}%
\TextAt{9}{375.694}{566.932}{\Character{41}}%
\NativeMathAt{383.567}{566.932}{11.955}{\leq}%
\TextAt{9}{396.186}{566.932}{\Character{91}}%
\TextAt{13}{399.438}{566.932}{L}%
\TextAt{9}{410.723}{566.932}{\Character{58}}%
\TextAt{13}{417.296}{566.932}{K}%
\TextAt{9}{428.113}{566.932}{\Character{93}}%
\TextAt{8}{431.365}{566.932}{,}%
\TextAt{8}{438.124}{566.932}{and}%
\TextAt{8}{460.4}{566.932}{since}%
\TextAt{8}{488.581}{566.932}{a}%
\TextAt{8}{497.849}{566.932}{system}%
\TextAt{8}{79.226}{549.886}{of}%
\TextAt{8}{93.231}{549.886}{generators}%
\TextAt{8}{150.562}{549.886}{of}%
\TextAt{8}{164.567}{549.886}{the}%
\TextAt{13}{185.398}{549.886}{A}%
\TextAt{8}{194.174}{549.886}{-module}%
\TextAt{13}{240.039}{549.886}{B}%
\TextAt{8}{254.108}{549.886}{is}%
\TextAt{8}{266.547}{549.886}{also}%
\TextAt{8}{290.702}{549.886}{a}%
\TextAt{8}{301.133}{549.886}{system}%
\TextAt{8}{340.626}{549.886}{of}%
\TextAt{8}{354.631}{549.886}{generators}%
\TextAt{8}{411.962}{549.886}{of}%
\TextAt{8}{425.967}{549.886}{the}%
\TextAt{13}{446.799}{549.886}{K}%
\TextAt{8}{457.616}{549.886}{-module}%
\TextAt{13}{503.482}{549.886}{L}%
\TextAt{8}{511.446}{549.886}{,}%
\TextAt{8}{519.439}{549.886}{we}%
\TextAt{8}{79.226}{532.84}{have}%
\TextAt{8}{106.439}{532.84}{rank}%
\TextAt{9}{129.537}{532.84}{\Character{40}}%
\TextAt{13}{134.089}{532.84}{B}%
\TextAt{9}{143.585}{532.84}{\Character{41}}%
\TextAt{9}{151.81}{532.84}{\Character{61}}%
\TextAt{9}{164.602}{532.84}{\Character{91}}%
\TextAt{13}{167.88}{532.84}{L}%
\TextAt{9}{179.537}{532.84}{\Character{58}}%
\TextAt{13}{186.482}{532.84}{K}%
\TextAt{9}{197.299}{532.84}{\Character{93}}%
\TextAt{8}{200.551}{532.84}{.}%
\TextAt{8}{209.655}{532.84}{Let}%
\TextAt{13}{230.842}{532.84}{\Character{22}}%
\TextAt{17}{237.885}{531.047}{1}%
\TextAt{13}{242.618}{532.84}{\Character{59}}%
\TextAt{13}{247.842}{532.84}{\Character{58}}%
\TextAt{13}{253.079}{532.84}{\Character{58}}%
\TextAt{13}{258.303}{532.84}{\Character{58}}%
\TextAt{13}{263.54}{532.84}{\Character{59}}%
\TextAt{13}{268.776}{532.84}{\Character{22}}%
\TextAt{21}{275.881}{531.047}{r}%
\NativeMathAt{284.129}{532.84}{11.955}{\in}%
\TextAt{13}{295.792}{532.84}{M}%
\TextAt{8}{312.481}{532.84}{be}%
\TextAt{8}{328.624}{532.84}{a}%
\TextAt{8}{338.598}{532.84}{system}%
\TextAt{8}{377.635}{532.84}{of}%
\TextAt{8}{391.183}{532.84}{generators}%
\TextAt{8}{448.057}{532.84}{of}%
\TextAt{8}{461.605}{532.84}{the}%
\TextAt{13}{481.98}{532.84}{B}%
\TextAt{8}{491.476}{532.84}{-module}%
\TextAt{13}{79.226}{515.794}{M}%
\TextAt{8}{91.8}{515.794}{.}%
\TextAt{8}{100.107}{515.794}{There}%
\TextAt{8}{133.494}{515.794}{exists}%
\TextAt{8}{165.386}{515.794}{an}%
\TextAt{13}{181.222}{515.794}{a}%
\NativeMathAt{190.688}{515.794}{11.955}{\in}%
\TextAt{13}{201.979}{515.794}{A}%
\TextAt{8}{210.754}{515.794}{,}%
\TextAt{13}{217.565}{515.794}{a}%
\NativeMathAt{227.031}{515.794}{11.955}{\neq}%
\TextAt{9}{239.441}{515.794}{0}%
\TextAt{8}{245.309}{515.794}{,}%
\TextAt{8}{252.12}{515.794}{such}%
\TextAt{8}{278.094}{515.794}{that}%
\TextAt{13}{303.041}{515.794}{a\Character{22}}%
\TextAt{21}{316.228}{514.001}{i}%
\NativeMathAt{322.931}{515.794}{11.955}{\in}%
\TextAt{13}{334.222}{515.794}{B}%
\TextAt{8}{343.718}{515.794}{,}%
\TextAt{13}{350.529}{515.794}{i}%
\TextAt{9}{357.843}{515.794}{\Character{61}}%
\TextAt{9}{370.252}{515.794}{1}%
\TextAt{13}{376.121}{515.794}{\Character{59}}%
\TextAt{13}{381.345}{515.794}{\Character{58}}%
\TextAt{13}{386.582}{515.794}{\Character{58}}%
\TextAt{13}{391.818}{515.794}{\Character{58}}%
\TextAt{13}{397.043}{515.794}{\Character{59}}%
\TextAt{13}{402.279}{515.794}{r}%
\TextAt{8}{407.943}{515.794}{,}%
\TextAt{8}{414.754}{515.794}{so}%
\TextAt{8}{428.701}{515.794}{that}%
\TextAt{13}{453.647}{515.794}{aM}%
\NativeMathAt{475.687}{515.794}{11.955}{\subseteq}%
\TextAt{13}{488.306}{515.794}{B}%
\TextAt{8}{497.802}{515.794}{.}%
\TextAt{8}{506.11}{515.794}{Then}%
\TextAt{13}{209.543}{490.488}{adM}%
\NativeMathAt{237.665}{490.488}{11.955}{\subseteq}%
\TextAt{13}{250.284}{490.488}{dB}%
\NativeMathAt{269.184}{490.488}{11.955}{\subseteq}%
\TextAt{13}{281.803}{490.488}{A\Character{11}}%
\TextAt{17}{298.024}{488.694}{1}%
\TextAt{9}{305.413}{490.488}{\Character{43}}%
\TextAt{16}{317.175}{490.488}{\Character{1}}%
\TextAt{16}{322.471}{490.488}{\Character{1}}%
\TextAt{16}{327.779}{490.488}{\Character{1}}%
\TextAt{9}{333.779}{490.488}{\Character{43}}%
\TextAt{13}{345.54}{490.488}{A\Character{11}}%
\TextAt{21}{361.761}{488.694}{n}%
\TextAt{9}{370.718}{490.488}{\Character{61}}%
\TextAt{13}{383.144}{490.488}{M}%
\TextAt{17}{394.469}{488.694}{0}%
\TextAt{13}{399.201}{490.488}{\Character{58}}%
\TextAt{8}{79.226}{465.181}{According}%
\TextAt{8}{133.405}{465.181}{to}%
\TextAt{8}{146.613}{465.181}{the}%
\TextAt{8}{165.668}{465.181}{main}%
\TextAt{8}{193.832}{465.181}{theorem}%
\TextAt{8}{238.256}{465.181}{on}%
\TextAt{8}{253.413}{465.181}{finitely}%
\TextAt{8}{291.656}{465.181}{generated}%
\TextAt{8}{343.886}{465.181}{modules}%
\TextAt{8}{388.692}{465.181}{over}%
\TextAt{8}{412.636}{465.181}{principal}%
\TextAt{8}{460.311}{465.181}{ideal}%
\TextAt{8}{487.172}{465.181}{domains,}%
\TextAt{8}{79.226}{448.135}{since}%
\TextAt{13}{107.895}{448.135}{M}%
\TextAt{17}{119.219}{446.342}{0}%
\TextAt{8}{127.849}{448.135}{is}%
\TextAt{8}{139.613}{448.135}{a}%
\TextAt{8}{149.368}{448.135}{free}%
\TextAt{13}{171.796}{448.135}{A}%
\TextAt{8}{180.572}{448.135}{-module,}%
\TextAt{8}{229.014}{448.135}{so}%
\TextAt{8}{243.384}{448.135}{is}%
\TextAt{13}{255.148}{448.135}{adM}%
\TextAt{8}{279.949}{448.135}{,}%
\TextAt{8}{287.098}{448.135}{and}%
\TextAt{8}{309.861}{448.135}{hence}%
\TextAt{8}{342.367}{448.135}{also}%
\TextAt{13}{365.847}{448.135}{M}%
\TextAt{8}{378.421}{448.135}{.}%
\TextAt{8}{386.869}{448.135}{Finally,}%
\TextAt{9}{133.765}{422.829}{\Character{91}}%
\TextAt{13}{137.017}{422.829}{L}%
\TextAt{9}{148.302}{422.829}{\Character{58}}%
\TextAt{13}{154.875}{422.829}{K}%
\TextAt{9}{165.693}{422.829}{\Character{93}}%
\TextAt{9}{172.244}{422.829}{\Character{61}}%
\TextAt{8}{184.69}{422.829}{rank}%
\TextAt{9}{207.788}{422.829}{\Character{40}}%
\TextAt{13}{212.34}{422.829}{B}%
\TextAt{9}{221.836}{422.829}{\Character{41}}%
\NativeMathAt{229.71}{422.829}{11.955}{\leq}%
\TextAt{8}{242.329}{422.829}{rank}%
\TextAt{9}{265.426}{422.829}{\Character{40}}%
\TextAt{13}{269.979}{422.829}{M}%
\TextAt{9}{282.552}{422.829}{\Character{41}}%
\TextAt{9}{290.407}{422.829}{\Character{61}}%
\TextAt{8}{302.851}{422.829}{rank}%
\TextAt{9}{325.948}{422.829}{\Character{40}}%
\TextAt{13}{330.501}{422.829}{adM}%
\TextAt{9}{355.302}{422.829}{\Character{41}}%
\NativeMathAt{363.175}{422.829}{11.955}{\leq}%
\TextAt{8}{375.794}{422.829}{rank}%
\TextAt{9}{398.892}{422.829}{\Character{40}}%
\TextAt{13}{403.444}{422.829}{M}%
\TextAt{17}{414.769}{421.036}{0}%
\TextAt{9}{419.501}{422.829}{\Character{41}}%
\TextAt{9}{427.356}{422.829}{\Character{61}}%
\TextAt{9}{439.777}{422.829}{\Character{91}}%
\TextAt{13}{443.052}{422.829}{L}%
\TextAt{9}{454.337}{422.829}{\Character{58}}%
\TextAt{13}{460.909}{422.829}{K}%
\TextAt{9}{471.727}{422.829}{\Character{93}}%
\TextAt{13}{474.979}{422.829}{\Character{59}}%
\TextAt{8}{79.226}{397.523}{hence}%
\TextAt{8}{111.732}{397.523}{rank}%
\TextAt{9}{134.83}{397.523}{\Character{40}}%
\TextAt{13}{139.382}{397.523}{M}%
\TextAt{9}{151.956}{397.523}{\Character{41}}%
\TextAt{9}{159.811}{397.523}{\Character{61}}%
\TextAt{9}{172.232}{397.523}{\Character{91}}%
\TextAt{13}{175.506}{397.523}{L}%
\TextAt{9}{186.791}{397.523}{\Character{58}}%
\TextAt{13}{193.364}{397.523}{K}%
\TextAt{9}{204.181}{397.523}{\Character{93}}%
\TextAt{8}{207.433}{397.523}{.}%
\Vector{q 0 G 0 g 0.399 w 452.334 -376.477 m 452.334 -368.408 l S Q}%
\Vector{q 0 G 0 g 0.399 w 452.534 -368.607 m 459.707 -368.607 l S Q}%
\Vector{q 0 G 0 g 0.399 w 452.534 -376.278 m 459.707 -376.278 l S Q}%
\Vector{q 0 G 0 g 0.399 w 459.906 -376.477 m 459.906 -368.408 l S Q}%
\TextAt{8}{97.159}{369.577}{It}%
\TextAt{8}{109.846}{369.577}{is}%
\TextAt{8}{121.624}{369.577}{in}%
\TextAt{8}{135.291}{369.577}{general}%
\TextAt{8}{175.631}{369.577}{a}%
\TextAt{8}{185.401}{369.577}{difficult}%
\TextAt{8}{228.334}{369.577}{problem}%
\TextAt{8}{273.874}{369.577}{to}%
\TextAt{8}{288.199}{369.577}{produce}%
\TextAt{8}{332.758}{369.577}{integral}%
\TextAt{8}{375.381}{369.577}{bases.}%
\TextAt{8}{410.664}{369.577}{In}%
\TextAt{8}{425.299}{369.577}{concrete}%
\TextAt{8}{471.485}{369.577}{situations}%
\TextAt{8}{524.963}{369.577}{it}%
\TextAt{8}{79.226}{352.531}{can}%
\TextAt{8}{99.849}{352.531}{also}%
\TextAt{8}{122.493}{352.531}{be}%
\TextAt{8}{137.58}{352.531}{an}%
\TextAt{8}{153.003}{352.531}{important}%
\TextAt{8}{207.459}{352.531}{one.}%
\TextAt{8}{233.191}{352.531}{This}%
\TextAt{8}{259.074}{352.531}{is}%
\TextAt{8}{270.001}{352.531}{why}%
\TextAt{8}{293.876}{352.531}{the}%
\TextAt{8}{313.196}{352.531}{following}%
\TextAt{8}{361.794}{352.531}{proposition}%
\TextAt{8}{422.491}{352.531}{is}%
\TextAt{8}{433.418}{352.531}{interesting.}%
\TextAt{8}{495.314}{352.531}{Instead}%
\TextAt{8}{79.226}{335.485}{of}%
\TextAt{8}{92.432}{335.485}{integral}%
\TextAt{8}{134.915}{335.485}{bases}%
\TextAt{8}{165.479}{335.485}{of}%
\TextAt{8}{178.685}{335.485}{the}%
\TextAt{8}{198.717}{335.485}{integral}%
\TextAt{8}{241.2}{335.485}{closure}%
\TextAt{13}{280.157}{335.485}{B}%
\TextAt{8}{293.426}{335.485}{of}%
\TextAt{13}{306.632}{335.485}{A}%
\TextAt{8}{319.18}{335.485}{in}%
\TextAt{13}{332.708}{335.485}{L}%
\TextAt{8}{340.672}{335.485}{,}%
\TextAt{8}{347.722}{335.485}{we}%
\TextAt{8}{364.825}{335.485}{will}%
\TextAt{8}{386.805}{335.485}{now}%
\TextAt{8}{411.069}{335.485}{simply}%
\TextAt{8}{448.4}{335.485}{speak}%
\TextAt{8}{480.853}{335.485}{of}%
\TextAt{8}{494.059}{335.485}{integral}%
\TextAt{8}{79.226}{318.439}{bases}%
\TextAt{8}{109.915}{318.439}{of}%
\TextAt{8}{123.245}{318.439}{the}%
\TextAt{8}{143.402}{318.439}{extension}%
\TextAt{13}{195.167}{318.439}{L}%
\TextAt{16}{206.453}{318.439}{\Character{106}}%
\TextAt{13}{213.094}{318.439}{K}%
\TextAt{8}{223.912}{318.439}{.}%
\TextAt{12}{79.226}{293.483}{Proposition}%
\TextAt{12}{152.429}{293.483}{(2.11).}%
\TextAt{8}{194.902}{293.483}{Let}%
\TextAt{13}{216.298}{293.483}{L}%
\TextAt{16}{228.311}{293.482}{\Character{106}}%
\TextAt{13}{235.68}{293.483}{K}%
\TextAt{8}{250.822}{293.483}{and}%
\TextAt{13}{274.012}{293.483}{L}%
\TextAt{3}{281.976}{297.821}{\Character{48}}%
\TextAt{16}{288.82}{293.483}{\Character{106}}%
\TextAt{13}{296.189}{293.482}{K}%
\TextAt{8}{311.331}{293.483}{be}%
\TextAt{8}{327.682}{293.483}{two}%
\TextAt{8}{350.226}{293.483}{Galois}%
\TextAt{8}{386.567}{293.483}{extensions}%
\TextAt{8}{443.374}{293.483}{of}%
\TextAt{8}{457.131}{293.483}{degree}%
\TextAt{13}{493.973}{293.483}{n}%
\TextAt{8}{500.961}{293.483}{,}%
\TextAt{8}{508.644}{293.483}{resp.}%
\TextAt{13}{79.226}{276.437}{n}%
\TextAt{3}{86.214}{280.775}{\Character{48}}%
\TextAt{8}{89.009}{276.437}{,}%
\TextAt{8}{96.124}{276.437}{such}%
\TextAt{8}{122.479}{276.437}{that}%
\TextAt{13}{147.805}{276.437}{L}%
\TextAt{16}{158.34}{276.437}{\Character{92}}%
\TextAt{13}{168.88}{276.437}{L}%
\TextAt{3}{176.844}{280.775}{\Character{48}}%
\TextAt{9}{182.96}{276.437}{\Character{61}}%
\TextAt{13}{195.386}{276.437}{K}%
\TextAt{8}{206.203}{276.437}{.}%
\TextAt{8}{214.637}{276.437}{Let}%
\TextAt{13}{235.564}{276.437}{\Character{33}}%
\TextAt{17}{242.854}{274.643}{1}%
\TextAt{13}{247.586}{276.437}{\Character{59}}%
\TextAt{13}{252.81}{276.437}{\Character{58}}%
\TextAt{13}{258.047}{276.437}{\Character{58}}%
\TextAt{13}{263.271}{276.437}{\Character{58}}%
\TextAt{13}{268.508}{276.437}{\Character{59}}%
\TextAt{13}{273.744}{276.437}{\Character{33}}%
\TextAt{21}{281.096}{274.643}{n}%
\TextAt{8}{286.732}{276.437}{,}%
\TextAt{8}{293.848}{276.437}{resp.}%
\TextAt{13}{323.156}{276.437}{\Character{33}}%
\TextAt{3}{330.874}{280.775}{\Character{48}}%
\TextAt{17}{330.445}{273.481}{1}%
\TextAt{13}{335.177}{276.437}{\Character{59}}%
\TextAt{13}{340.401}{276.437}{\Character{58}}%
\TextAt{13}{345.638}{276.437}{\Character{58}}%
\TextAt{13}{350.874}{276.437}{\Character{58}}%
\TextAt{13}{356.099}{276.437}{\Character{59}}%
\TextAt{13}{361.335}{276.437}{\Character{33}}%
\TextAt{3}{369.116}{280.775}{\Character{48}}%
\TextAt{21}{368.687}{273.272}{n}%
\TextAt{23}{373.825}{275.541}{\Character{48}}%
\TextAt{8}{377.008}{276.437}{,}%
\TextAt{8}{384.123}{276.437}{be}%
\TextAt{8}{400.005}{276.437}{an}%
\TextAt{8}{416.222}{276.437}{integral}%
\TextAt{8}{458.788}{276.437}{basis}%
\TextAt{8}{487.485}{276.437}{of}%
\TextAt{13}{500.773}{276.437}{L}%
\TextAt{16}{512.058}{276.437}{\Character{106}}%
\TextAt{13}{518.7}{276.437}{K}%
\TextAt{8}{529.518}{276.437}{,}%
\TextAt{8}{79.226}{259.391}{resp.}%
\TextAt{13}{108.708}{259.391}{L}%
\TextAt{3}{116.673}{263.729}{\Character{48}}%
\TextAt{16}{122.879}{259.391}{\Character{106}}%
\TextAt{13}{129.612}{259.391}{K}%
\TextAt{8}{140.43}{259.391}{,}%
\TextAt{8}{147.645}{259.391}{with}%
\TextAt{8}{174.359}{259.391}{discriminant}%
\TextAt{13}{241.791}{259.391}{d}%
\TextAt{8}{247.874}{259.391}{,}%
\TextAt{8}{255.09}{259.391}{resp.}%
\TextAt{13}{284.572}{259.391}{d}%
\TextAt{3}{290.655}{263.729}{\Character{48}}%
\TextAt{8}{293.45}{259.391}{.}%
\TextAt{8}{302.058}{259.391}{Suppose}%
\TextAt{8}{348.019}{259.391}{that}%
\TextAt{13}{373.441}{259.391}{d}%
\TextAt{8}{383.475}{259.391}{and}%
\TextAt{13}{406.291}{259.391}{d}%
\TextAt{3}{412.373}{263.729}{\Character{48}}%
\TextAt{8}{419.119}{259.391}{are}%
\TextAt{8}{438.683}{259.391}{relatively}%
\TextAt{8}{489.797}{259.391}{prime}%
\TextAt{8}{523.014}{259.391}{in}%
\TextAt{8}{79.226}{242.345}{the}%
\TextAt{8}{99.242}{242.345}{sense}%
\TextAt{8}{129.134}{242.345}{that}%
\TextAt{13}{154.362}{242.345}{xd}%
\TextAt{9}{169.467}{242.345}{\Character{43}}%
\TextAt{13}{180.942}{242.345}{x}%
\TextAt{3}{187.594}{246.683}{\Character{48}}%
\TextAt{13}{190.389}{242.345}{d}%
\TextAt{3}{196.472}{246.683}{\Character{48}}%
\TextAt{9}{202.588}{242.345}{\Character{61}}%
\TextAt{9}{214.998}{242.345}{1}%
\TextAt{8}{220.866}{242.345}{,}%
\TextAt{8}{227.903}{242.345}{for}%
\TextAt{8}{245.648}{242.345}{suitable}%
\TextAt{13}{289.144}{242.345}{x\Character{59}}%
\TextAt{13}{301.016}{242.345}{x}%
\TextAt{3}{307.693}{246.683}{\Character{48}}%
\NativeMathAt{313.809}{242.345}{11.955}{\in}%
\TextAt{13}{325.1}{242.345}{A}%
\TextAt{8}{333.875}{242.345}{.}%
\TextAt{8}{342.276}{242.345}{Then}%
\TextAt{13}{372.694}{242.345}{\Character{33}}%
\TextAt{21}{379.983}{240.551}{i}%
\TextAt{13}{383.364}{242.345}{\Character{33}}%
\TextAt{3}{391.082}{246.683}{\Character{48}}%
\TextAt{21}{390.654}{239.389}{j}%
\TextAt{8}{398.793}{242.345}{is}%
\TextAt{8}{410.417}{242.345}{an}%
\TextAt{8}{426.536}{242.345}{integral}%
\TextAt{8}{469.004}{242.345}{basis}%
\TextAt{8}{497.604}{242.345}{of}%
\TextAt{13}{510.794}{242.345}{LL}%
\TextAt{3}{526.723}{246.683}{\Character{48}}%
\TextAt{8}{529.518}{242.345}{,}%
\TextAt{8}{79.226}{225.299}{of}%
\TextAt{8}{92.556}{225.299}{discriminant}%
\TextAt{13}{159.935}{225.299}{d}%
\TextAt{21}{166.018}{229.637}{n}%
\TextAt{23}{171.156}{232.45}{\Character{48}}%
\TextAt{13}{174.339}{225.299}{d}%
\TextAt{3}{180.422}{229.637}{\Character{48}}%
\TextAt{21}{182.719}{229.637}{n}%
\TextAt{8}{188.355}{225.299}{.}%
\TextAt{19}{79.226}{192.431}{Proof.}%
\TextAt{8}{116}{192.431}{As}%
\TextAt{13}{133.017}{192.431}{L}%
\TextAt{16}{143.086}{192.431}{\Character{92}}%
\TextAt{13}{153.16}{192.431}{L}%
\TextAt{3}{161.124}{196.77}{\Character{48}}%
\TextAt{9}{167.24}{192.431}{\Character{61}}%
\TextAt{13}{179.666}{192.431}{K}%
\TextAt{8}{190.483}{192.431}{,}%
\TextAt{8}{197.416}{192.431}{we}%
\TextAt{8}{214.374}{192.431}{have}%
\TextAt{9}{241.098}{192.431}{\Character{91}}%
\TextAt{13}{244.35}{192.431}{LL}%
\TextAt{3}{260.279}{196.77}{\Character{48}}%
\TextAt{9}{266.395}{192.431}{\Character{58}}%
\TextAt{13}{272.967}{192.431}{K}%
\TextAt{9}{283.785}{192.431}{\Character{93}}%
\TextAt{9}{290.336}{192.431}{\Character{61}}%
\TextAt{13}{302.783}{192.431}{nn}%
\TextAt{3}{316.758}{196.77}{\Character{48}}%
\TextAt{8}{319.553}{192.431}{,}%
\TextAt{8}{326.486}{192.431}{so}%
\TextAt{8}{340.586}{192.431}{the}%
\TextAt{13}{360.472}{192.431}{nn}%
\TextAt{3}{374.447}{196.77}{\Character{48}}%
\TextAt{8}{380.87}{192.431}{products}%
\TextAt{13}{429.113}{192.431}{\Character{33}}%
\TextAt{21}{436.403}{190.638}{i}%
\TextAt{13}{439.784}{192.431}{\Character{33}}%
\TextAt{3}{447.502}{196.77}{\Character{48}}%
\TextAt{21}{447.073}{189.476}{j}%
\TextAt{8}{455.083}{192.431}{do}%
\TextAt{8}{471.071}{192.431}{form}%
\TextAt{8}{498.441}{192.431}{a}%
\TextAt{8}{507.927}{192.431}{basis}%
\TextAt{8}{79.226}{175.385}{of}%
\TextAt{13}{92.556}{175.385}{LL}%
\TextAt{3}{108.485}{179.724}{\Character{48}}%
\TextAt{16}{114.601}{175.385}{\Character{106}}%
\TextAt{13}{121.243}{175.385}{K}%
\TextAt{8}{132.06}{175.385}{.}%
\TextAt{8}{140.509}{175.385}{Now}%
\TextAt{8}{167.169}{175.385}{let}%
\TextAt{13}{184.073}{175.385}{\Character{11}}%
\TextAt{8}{195.493}{175.385}{be}%
\TextAt{8}{211.417}{175.385}{an}%
\TextAt{8}{227.676}{175.385}{integral}%
\TextAt{8}{270.284}{175.385}{element}%
\TextAt{8}{313.526}{175.385}{of}%
\TextAt{13}{326.856}{175.385}{LL}%
\TextAt{3}{342.785}{179.724}{\Character{48}}%
\TextAt{8}{345.58}{175.385}{,}%
\TextAt{8}{352.729}{175.385}{and}%
\TextAt{8}{375.492}{175.385}{write}%
\TextAt{13}{231.559}{150.079}{\Character{11}}%
\TextAt{9}{242.402}{150.079}{\Character{61}}%
\TextAt{20}{254.827}{161.436}{\Character{88}}%
\TextAt{21}{258.902}{136.242}{i\Character{59}j}%
\TextAt{13}{274.089}{150.079}{a}%
\TextAt{21}{280.233}{148.286}{ij}%
\TextAt{13}{287.499}{150.079}{\Character{33}}%
\TextAt{21}{294.788}{148.286}{i}%
\TextAt{13}{298.169}{150.079}{\Character{33}}%
\TextAt{3}{305.888}{155.015}{\Character{48}}%
\TextAt{21}{305.459}{147.123}{j}%
\TextAt{13}{309.841}{150.079}{\Character{59}}%
\TextAt{13}{338.976}{150.079}{a}%
\TextAt{21}{345.14}{148.286}{ij}%
\NativeMathAt{355.726}{150.079}{11.955}{\in}%
\TextAt{13}{367.017}{150.079}{K}%
\TextAt{13}{377.179}{150.079}{\Character{58}}%
\TextAt{8}{79.226}{115.274}{We}%
\TextAt{8}{99.848}{115.274}{have}%
\TextAt{8}{127.32}{115.274}{to}%
\TextAt{8}{142.107}{115.274}{show}%
\TextAt{8}{171.587}{115.274}{that}%
\TextAt{13}{197.433}{115.274}{a}%
\TextAt{21}{203.578}{113.48}{ij}%
\NativeMathAt{214.978}{115.274}{11.955}{\in}%
\TextAt{13}{227.082}{115.274}{A}%
\TextAt{8}{235.857}{115.274}{.}%
\TextAt{8}{245.737}{115.274}{Put}%
\TextAt{13}{269.132}{115.274}{\Character{12}}%
\TextAt{21}{275.74}{113.48}{j}%
\TextAt{9}{284.256}{115.274}{\Character{61}}%
\TextAt{20}{297.495}{124.24}{\Character{80}}%
\TextAt{21}{310.114}{111.787}{i}%
\TextAt{13}{315.488}{115.274}{a}%
\TextAt{21}{321.633}{113.48}{ij}%
\TextAt{13}{328.898}{115.274}{\Character{33}}%
\TextAt{21}{336.188}{113.48}{i}%
\TextAt{8}{339.569}{115.274}{.}%
\TextAt{8}{349.449}{115.274}{Let}%
\TextAt{13}{370.895}{115.274}{G}%
\TextAt{9}{380.129}{115.274}{\Character{40}}%
\TextAt{13}{384.681}{115.274}{LL}%
\TextAt{3}{400.61}{119.612}{\Character{48}}%
\TextAt{16}{407.539}{115.274}{\Character{106}}%
\TextAt{13}{414.994}{115.274}{L}%
\TextAt{3}{422.958}{119.612}{\Character{48}}%
\TextAt{9}{425.753}{115.274}{\Character{41}}%
\TextAt{9}{434.421}{115.274}{\Character{61}}%
\TextAt{16}{447.678}{115.274}{\Character{102}}%
\TextAt{13}{453.656}{115.274}{\Character{27}}%
\TextAt{17}{460.309}{113.48}{1}%
\TextAt{13}{465.042}{115.274}{\Character{59}}%
\TextAt{13}{470.266}{115.274}{\Character{58}}%
\TextAt{13}{475.503}{115.274}{\Character{58}}%
\TextAt{13}{480.727}{115.274}{\Character{58}}%
\TextAt{13}{485.964}{115.274}{\Character{59}}%
\TextAt{13}{491.2}{115.274}{\Character{27}}%
\TextAt{21}{497.916}{113.48}{n}%
\TextAt{16}{503.552}{115.274}{\Character{103}}%
\TextAt{8}{513.904}{115.274}{and}%
\TextAt{13}{79.226}{98.228}{G}%
\TextAt{9}{88.46}{98.228}{\Character{40}}%
\TextAt{13}{93.012}{98.228}{LL}%
\TextAt{3}{108.941}{102.566}{\Character{48}}%
\TextAt{16}{115.057}{98.228}{\Character{106}}%
\TextAt{13}{121.699}{98.228}{L}%
\TextAt{9}{129.663}{98.228}{\Character{41}}%
\TextAt{9}{137.518}{98.228}{\Character{61}}%
\TextAt{16}{149.962}{98.228}{\Character{102}}%
\TextAt{13}{155.939}{98.228}{\Character{27}}%
\TextAt{3}{163.022}{102.566}{\Character{48}}%
\TextAt{17}{162.593}{95.272}{1}%
\TextAt{13}{167.325}{98.228}{\Character{59}}%
\TextAt{13}{172.549}{98.228}{\Character{58}}%
\TextAt{13}{177.786}{98.228}{\Character{58}}%
\TextAt{13}{183.022}{98.228}{\Character{58}}%
\TextAt{13}{188.247}{98.228}{\Character{59}}%
\TextAt{13}{193.483}{98.228}{\Character{27}}%
\TextAt{3}{200.628}{102.566}{\Character{48}}%
\TextAt{21}{200.199}{95.063}{n}%
\TextAt{23}{205.338}{97.333}{\Character{48}}%
\TextAt{16}{208.52}{98.228}{\Character{103}}%
\TextAt{8}{214.498}{98.228}{.}%
\TextAt{8}{222.946}{98.228}{Thus}%
\TextAt{13}{178.95}{72.921}{G}%
\TextAt{9}{188.183}{72.921}{\Character{40}}%
\TextAt{13}{192.736}{72.921}{LL}%
\TextAt{3}{208.665}{77.858}{\Character{48}}%
\TextAt{16}{214.781}{72.921}{\Character{106}}%
\TextAt{13}{221.422}{72.921}{K}%
\TextAt{9}{232.24}{72.921}{\Character{41}}%
\TextAt{9}{240.095}{72.921}{\Character{61}}%
\TextAt{16}{252.539}{72.921}{\Character{102}}%
\TextAt{13}{258.516}{72.921}{\Character{27}}%
\TextAt{21}{265.17}{71.128}{k}%
\TextAt{13}{270.29}{72.921}{\Character{27}}%
\TextAt{3}{277.372}{77.858}{\Character{48}}%
\TextAt{21}{276.943}{69.966}{\Character{96}}%
\TextAt{16}{284.221}{72.921}{\Character{106}}%
\TextAt{13}{290.863}{72.921}{k}%
\TextAt{9}{300.673}{72.921}{\Character{61}}%
\TextAt{9}{313.083}{72.921}{1}%
\TextAt{13}{318.952}{72.921}{\Character{59}}%
\TextAt{13}{324.176}{72.921}{\Character{58}}%
\TextAt{13}{329.413}{72.921}{\Character{58}}%
\TextAt{13}{334.637}{72.921}{\Character{58}}%
\TextAt{13}{339.874}{72.921}{\Character{59}}%
\TextAt{13}{345.11}{72.921}{n\Character{59}}%
\TextAt{13}{361.214}{72.921}{\Character{96}}%
\TextAt{9}{369.535}{72.921}{\Character{61}}%
\TextAt{9}{381.945}{72.921}{1}%
\TextAt{13}{387.813}{72.921}{\Character{59}}%
\TextAt{13}{393.037}{72.921}{\Character{58}}%
\TextAt{13}{398.274}{72.921}{\Character{58}}%
\TextAt{13}{403.498}{72.921}{\Character{58}}%
\TextAt{13}{408.735}{72.921}{\Character{59}}%
\TextAt{13}{413.971}{72.921}{n}%
\TextAt{3}{421.021}{77.858}{\Character{48}}%
\TextAt{16}{423.817}{72.921}{\Character{103}}%
\TextAt{13}{429.794}{72.921}{\Character{58}}%
\TextAt{8}{300.14}{43.033}{12}%
\Destination{78.23}{809.815}{nav.page.12}%
\Destination{79.36}{616.055}{nav.print.13}%
\Destination{79.36}{681.955}{nav.number.1.2.10}%
\Destination{79.36}{307.455}{nav.number.1.2.11}%
\Hotlink{504.94}{581.755}{25.078}{11.047}{nav.number.1.2.9}%
}

% Page 25 (body page 13)
\Page{%
\TextAt{2}{79.226}{796.943}{I.}%
\TextAt{2}{89.83}{796.943}{Algebraic}%
\TextAt{2}{138.648}{796.943}{Integers}%
\Vector{q 0 G 0 g 0.299 w 7.226 18.729 m 460.769 18.729 l S Q}%
\TextAt{8}{79.226}{761.124}{Putting}%
\TextAt{13}{198.264}{744.078}{T}%
\TextAt{9}{210.072}{744.078}{\Character{61}}%
\TextAt{9}{222.482}{744.078}{\Character{40}}%
\TextAt{13}{227.05}{744.078}{\Character{27}}%
\TextAt{3}{234.132}{749.014}{\Character{48}}%
\TextAt{21}{233.703}{741.122}{\Character{96}}%
\TextAt{13}{237.66}{744.078}{\Character{33}}%
\TextAt{3}{245.379}{749.014}{\Character{48}}%
\TextAt{21}{244.95}{741.122}{j}%
\TextAt{9}{249.332}{744.078}{\Character{41}}%
\TextAt{13}{253.884}{744.078}{\Character{59}}%
\TextAt{13}{283.019}{744.078}{a}%
\TextAt{9}{292.504}{744.078}{\Character{61}}%
\TextAt{9}{304.913}{744.078}{\Character{40}}%
\TextAt{13}{309.482}{744.078}{\Character{27}}%
\TextAt{3}{316.565}{749.014}{\Character{48}}%
\TextAt{21}{316.136}{741.122}{\Character{96}}%
\TextAt{13}{320.093}{744.078}{\Character{11}}%
\TextAt{9}{327.615}{744.078}{\Character{41}}%
\TextAt{21}{332.167}{742.284}{\Character{96}}%
\TextAt{13}{336.125}{744.078}{\Character{59}}%
\TextAt{13}{365.26}{744.078}{b}%
\TextAt{9}{373.577}{744.078}{\Character{61}}%
\TextAt{9}{385.987}{744.078}{\Character{40}}%
\TextAt{13}{390.555}{744.078}{\Character{12}}%
\TextAt{21}{397.163}{742.284}{j}%
\TextAt{9}{401.545}{744.078}{\Character{41}}%
\TextAt{21}{406.098}{742.284}{j}%
\TextAt{13}{410.48}{744.078}{\Character{59}}%
\TextAt{8}{79.226}{720.556}{one}%
\TextAt{8}{101.561}{720.556}{finds}%
\TextAt{8}{130.46}{720.556}{det}%
\TextAt{9}{146.719}{720.556}{\Character{40}}%
\TextAt{13}{151.271}{720.556}{T}%
\TextAt{9}{159.758}{720.556}{\Character{41}}%
\TextAt{17}{164.31}{724.894}{2}%
\TextAt{9}{173.856}{720.556}{\Character{61}}%
\TextAt{13}{187.774}{720.556}{d}%
\TextAt{3}{193.856}{724.894}{\Character{48}}%
\TextAt{8}{201.425}{720.556}{and}%
\TextAt{13}{225.063}{720.556}{a}%
\TextAt{9}{236.021}{720.556}{\Character{61}}%
\TextAt{13}{249.939}{720.556}{T}%
\TextAt{13}{258.403}{720.556}{b}%
\TextAt{8}{263.403}{720.556}{.}%
\TextAt{8}{274.478}{720.556}{Write}%
\TextAt{13}{307.86}{720.556}{T}%
\TextAt{3}{316.347}{724.894}{\Character{3}}%
\TextAt{8}{325.852}{720.556}{for}%
\TextAt{8}{344.613}{720.556}{the}%
\TextAt{8}{365.645}{720.556}{adjoint}%
\TextAt{8}{406.2}{720.556}{matrix}%
\TextAt{8}{445.129}{720.556}{of}%
\TextAt{13}{459.334}{720.556}{T}%
\TextAt{8}{467.821}{720.556}{.}%
\TextAt{8}{478.897}{720.556}{Then}%
\TextAt{8}{510.33}{720.556}{row-}%
\TextAt{8}{79.226}{703.51}{column}%
\TextAt{8}{120.197}{703.51}{expansion}%
\TextAt{8}{174.569}{703.51}{(2.3)}%
\TextAt{8}{202.544}{703.51}{gives}%
\TextAt{8}{267.403}{686.464}{det}%
\TextAt{9}{283.662}{686.464}{\Character{40}}%
\TextAt{13}{288.214}{686.464}{T}%
\TextAt{9}{296.701}{686.464}{\Character{41}}%
\TextAt{13}{301.253}{686.464}{b}%
\TextAt{9}{309.551}{686.464}{\Character{61}}%
\TextAt{13}{321.977}{686.464}{T}%
\TextAt{3}{330.464}{691.4}{\Character{3}}%
\TextAt{13}{335.196}{686.464}{a\Character{58}}%
\TextAt{8}{79.226}{662.942}{Since}%
\TextAt{13}{110.825}{662.942}{T}%
\TextAt{3}{119.311}{667.281}{\Character{3}}%
\TextAt{8}{128.982}{662.942}{and}%
\TextAt{13}{152.786}{662.942}{a}%
\TextAt{8}{163.869}{662.942}{have}%
\TextAt{8}{191.905}{662.942}{integral}%
\TextAt{8}{235.554}{662.942}{entries}%
\TextAt{8}{274.051}{662.942}{in}%
\TextAt{13}{288.744}{662.942}{LL}%
\TextAt{3}{304.674}{667.281}{\Character{48}}%
\TextAt{8}{307.469}{662.942}{,}%
\TextAt{8}{315.919}{662.942}{the}%
\TextAt{8}{337.116}{662.942}{multiple}%
\TextAt{13}{384.006}{662.942}{d}%
\TextAt{3}{390.088}{667.281}{\Character{48}}%
\TextAt{13}{392.883}{662.942}{b}%
\TextAt{8}{402.799}{662.942}{has}%
\TextAt{8}{424.714}{662.942}{integral}%
\TextAt{8}{468.363}{662.942}{entries}%
\TextAt{8}{506.859}{662.942}{in}%
\TextAt{13}{521.553}{662.942}{L}%
\TextAt{8}{529.518}{662.942}{,}%
\TextAt{8}{79.226}{645.896}{namely}%
\TextAt{13}{121.157}{645.896}{d}%
\TextAt{3}{127.239}{650.235}{\Character{48}}%
\TextAt{13}{130.035}{645.896}{\Character{12}}%
\TextAt{21}{136.643}{644.103}{j}%
\TextAt{9}{146.533}{645.896}{\Character{61}}%
\TextAt{20}{161.145}{654.863}{\Character{80}}%
\TextAt{21}{173.764}{642.409}{i}%
\TextAt{13}{179.138}{645.896}{d}%
\TextAt{3}{185.22}{650.235}{\Character{48}}%
\TextAt{13}{188.016}{645.896}{a}%
\TextAt{21}{194.16}{644.103}{ij}%
\TextAt{13}{201.426}{645.896}{\Character{33}}%
\TextAt{21}{208.715}{644.103}{i}%
\TextAt{8}{212.096}{645.896}{.}%
\TextAt{8}{224.394}{645.896}{Thus}%
\TextAt{13}{255.326}{645.896}{d}%
\TextAt{3}{261.408}{650.235}{\Character{48}}%
\TextAt{13}{264.204}{645.896}{a}%
\TextAt{21}{270.348}{644.103}{ij}%
\NativeMathAt{283.121}{645.896}{11.955}{\in}%
\TextAt{13}{296.599}{645.896}{A}%
\TextAt{8}{305.374}{645.896}{.}%
\TextAt{8}{317.672}{645.896}{Swapping}%
\TextAt{8}{371.964}{645.896}{the}%
\TextAt{8}{393.403}{645.896}{roles}%
\TextAt{8}{422.064}{645.896}{of}%
\TextAt{9}{436.677}{645.896}{\Character{40}}%
\TextAt{13}{441.229}{645.896}{\Character{33}}%
\TextAt{21}{448.518}{644.103}{i}%
\TextAt{9}{451.9}{645.896}{\Character{41}}%
\TextAt{8}{461.633}{645.896}{and}%
\TextAt{9}{485.678}{645.896}{\Character{40}}%
\TextAt{13}{490.231}{645.896}{\Character{33}}%
\TextAt{3}{497.949}{650.235}{\Character{48}}%
\TextAt{21}{497.52}{642.941}{j}%
\TextAt{9}{501.902}{645.896}{\Character{41}}%
\TextAt{8}{506.454}{645.896}{,}%
\TextAt{8}{515.207}{645.896}{one}%
\TextAt{8}{79.226}{628.85}{checks}%
\TextAt{8}{115.379}{628.85}{in}%
\TextAt{8}{129.031}{628.85}{the}%
\TextAt{8}{149.188}{628.85}{same}%
\TextAt{8}{178.514}{628.85}{manner}%
\TextAt{8}{220.787}{628.85}{that}%
\TextAt{13}{246.156}{628.85}{da}%
\TextAt{21}{258.384}{627.057}{ij}%
\NativeMathAt{268.97}{628.85}{11.955}{\in}%
\TextAt{13}{280.261}{628.85}{A}%
\TextAt{8}{289.036}{628.85}{,}%
\TextAt{8}{296.185}{628.85}{so}%
\TextAt{8}{310.556}{628.85}{that}%
\TextAt{13}{241.951}{599.849}{a}%
\TextAt{21}{248.096}{598.056}{ij}%
\TextAt{9}{258.682}{599.849}{\Character{61}}%
\TextAt{13}{271.108}{599.849}{xda}%
\TextAt{21}{289.987}{598.056}{ij}%
\TextAt{9}{299.909}{599.849}{\Character{43}}%
\TextAt{13}{311.671}{599.849}{x}%
\TextAt{3}{318.323}{604.786}{\Character{48}}%
\TextAt{13}{321.118}{599.849}{d}%
\TextAt{3}{327.2}{604.786}{\Character{48}}%
\TextAt{13}{329.996}{599.849}{a}%
\TextAt{21}{336.14}{598.056}{ij}%
\NativeMathAt{346.727}{599.849}{11.955}{\in}%
\TextAt{13}{358.018}{599.849}{A\Character{58}}%
\TextAt{8}{79.226}{570.848}{Therefore}%
\TextAt{13}{132.858}{570.848}{\Character{33}}%
\TextAt{21}{140.147}{569.055}{i}%
\TextAt{13}{143.529}{570.848}{\Character{33}}%
\TextAt{3}{151.247}{575.187}{\Character{48}}%
\TextAt{21}{150.818}{567.893}{j}%
\TextAt{8}{159.732}{570.848}{is}%
\TextAt{8}{172.131}{570.848}{indeed}%
\TextAt{8}{209.826}{570.848}{an}%
\TextAt{8}{226.72}{570.848}{integral}%
\TextAt{8}{269.963}{570.848}{basis}%
\TextAt{8}{299.338}{570.848}{of}%
\TextAt{13}{313.302}{570.848}{LL}%
\TextAt{3}{329.231}{575.187}{\Character{48}}%
\TextAt{16}{336.429}{570.848}{\Character{106}}%
\TextAt{13}{344.152}{570.848}{K}%
\TextAt{8}{354.97}{570.848}{.}%
\TextAt{8}{365.322}{570.848}{We}%
\TextAt{8}{386.101}{570.848}{compute}%
\TextAt{8}{434.209}{570.848}{the}%
\TextAt{8}{455.001}{570.848}{discriminant}%
\TextAt{9}{523.015}{570.848}{\Character{1}}%
\TextAt{8}{79.226}{553.802}{of}%
\TextAt{8}{93.023}{553.802}{this}%
\TextAt{8}{116.313}{553.802}{integral}%
\TextAt{8}{159.388}{553.802}{basis.}%
\TextAt{8}{194.08}{553.802}{Since}%
\TextAt{13}{225.105}{553.802}{G}%
\TextAt{9}{234.338}{553.802}{\Character{40}}%
\TextAt{13}{238.891}{553.802}{LL}%
\TextAt{3}{254.82}{558.141}{\Character{48}}%
\TextAt{16}{261.732}{553.802}{\Character{106}}%
\TextAt{13}{269.169}{553.802}{K}%
\TextAt{9}{279.987}{553.802}{\Character{41}}%
\TextAt{9}{288.643}{553.802}{\Character{61}}%
\TextAt{16}{301.877}{553.802}{\Character{102}}%
\TextAt{13}{307.855}{553.802}{\Character{27}}%
\TextAt{21}{314.508}{552.009}{k}%
\TextAt{13}{319.628}{553.802}{\Character{27}}%
\TextAt{3}{326.711}{558.141}{\Character{48}}%
\TextAt{21}{326.282}{550.693}{\Character{96}}%
\TextAt{16}{334.356}{553.802}{\Character{106}}%
\TextAt{13}{341.793}{553.802}{k}%
\TextAt{9}{352.4}{553.802}{\Character{61}}%
\TextAt{9}{365.611}{553.802}{1}%
\TextAt{13}{371.474}{553.802}{\Character{59}}%
\TextAt{13}{376.698}{553.802}{\Character{58}}%
\TextAt{13}{381.935}{553.802}{\Character{58}}%
\TextAt{13}{387.159}{553.802}{\Character{58}}%
\TextAt{13}{392.396}{553.802}{\Character{59}}%
\TextAt{13}{397.632}{553.802}{n\Character{59}}%
\TextAt{13}{414.214}{553.802}{\Character{96}}%
\TextAt{9}{423.32}{553.802}{\Character{61}}%
\TextAt{9}{436.531}{553.802}{1}%
\TextAt{13}{442.394}{553.802}{\Character{59}}%
\TextAt{13}{447.618}{553.802}{\Character{58}}%
\TextAt{13}{452.855}{553.802}{\Character{58}}%
\TextAt{13}{458.091}{553.802}{\Character{58}}%
\TextAt{13}{463.316}{553.802}{\Character{59}}%
\TextAt{13}{468.552}{553.802}{n}%
\TextAt{3}{475.603}{558.141}{\Character{48}}%
\TextAt{16}{478.398}{553.802}{\Character{103}}%
\TextAt{8}{484.375}{553.802}{,}%
\TextAt{8}{492.108}{553.802}{it}%
\TextAt{8}{504.28}{553.802}{is}%
\TextAt{8}{516.51}{553.802}{the}%
\TextAt{8}{79.226}{536.756}{square}%
\TextAt{8}{116.036}{536.756}{of}%
\TextAt{8}{129.366}{536.756}{the}%
\TextAt{8}{149.523}{536.756}{determinant}%
\TextAt{8}{215.539}{536.756}{of}%
\TextAt{8}{228.869}{536.756}{the}%
\TextAt{9}{249.025}{536.756}{\Character{40}}%
\TextAt{13}{253.578}{536.756}{nn}%
\TextAt{3}{267.553}{541.095}{\Character{48}}%
\TextAt{16}{273.005}{536.756}{\Character{2}}%
\TextAt{13}{284.96}{536.756}{nn}%
\TextAt{3}{298.935}{541.095}{\Character{48}}%
\TextAt{9}{301.73}{536.756}{\Character{41}}%
\TextAt{8}{306.282}{536.756}{-matrix}%
\TextAt{13}{227.512}{507.755}{M}%
\TextAt{9}{243.406}{507.755}{\Character{61}}%
\TextAt{9}{255.816}{507.755}{\Character{40}}%
\TextAt{13}{260.384}{507.755}{\Character{27}}%
\TextAt{21}{267.038}{505.962}{k}%
\TextAt{13}{272.157}{507.755}{\Character{27}}%
\TextAt{3}{279.24}{512.691}{\Character{48}}%
\TextAt{21}{278.811}{504.8}{\Character{96}}%
\TextAt{13}{282.768}{507.755}{\Character{33}}%
\TextAt{21}{290.058}{505.962}{i}%
\TextAt{13}{293.439}{507.755}{\Character{33}}%
\TextAt{3}{301.157}{512.691}{\Character{48}}%
\TextAt{21}{300.728}{504.8}{j}%
\TextAt{9}{305.11}{507.755}{\Character{41}}%
\TextAt{9}{312.965}{507.755}{\Character{61}}%
\TextAt{9}{325.386}{507.755}{\Character{40}}%
\TextAt{13}{329.961}{507.755}{\Character{27}}%
\TextAt{21}{336.615}{505.962}{k}%
\TextAt{13}{341.734}{507.755}{\Character{33}}%
\TextAt{21}{349.024}{505.962}{i}%
\TextAt{13}{354.397}{507.755}{\Character{27}}%
\TextAt{3}{361.48}{512.691}{\Character{48}}%
\TextAt{21}{361.051}{504.8}{\Character{96}}%
\TextAt{13}{365.008}{507.755}{\Character{33}}%
\TextAt{3}{372.726}{512.691}{\Character{48}}%
\TextAt{21}{372.298}{504.8}{j}%
\TextAt{9}{376.68}{507.755}{\Character{41}}%
\TextAt{13}{381.232}{507.755}{\Character{58}}%
\TextAt{8}{79.226}{478.754}{This}%
\TextAt{8}{105.893}{478.754}{matrix}%
\TextAt{8}{143.893}{478.754}{is}%
\TextAt{8}{155.604}{478.754}{itself}%
\TextAt{8}{183.897}{478.754}{an}%
\TextAt{9}{200.103}{478.754}{\Character{40}}%
\TextAt{13}{204.655}{478.754}{n}%
\TextAt{3}{211.643}{483.092}{\Character{48}}%
\TextAt{16}{216.986}{478.754}{\Character{2}}%
\TextAt{13}{228.833}{478.754}{n}%
\TextAt{3}{235.82}{483.092}{\Character{48}}%
\TextAt{9}{238.615}{478.754}{\Character{41}}%
\TextAt{8}{243.168}{478.754}{-matrix}%
\TextAt{8}{285.065}{478.754}{with}%
\TextAt{8}{311.672}{478.754}{entries}%
\TextAt{9}{349.075}{478.754}{\Character{40}}%
\TextAt{13}{353.627}{478.754}{n}%
\TextAt{16}{363.163}{478.754}{\Character{2}}%
\TextAt{13}{375.01}{478.754}{n}%
\TextAt{9}{381.997}{478.754}{\Character{41}}%
\TextAt{8}{386.549}{478.754}{-matrices}%
\TextAt{8}{437.282}{478.754}{of}%
\TextAt{8}{450.559}{478.754}{which}%
\TextAt{8}{483.992}{478.754}{the}%
\TextAt{9}{504.096}{478.754}{\Character{40}}%
\TextAt{13}{508.648}{478.754}{\Character{96}\Character{59}}%
\TextAt{13}{518.774}{478.754}{j}%
\TextAt{9}{524.32}{478.754}{\Character{41}}%
\TextAt{8}{528.872}{478.754}{-}%
\TextAt{8}{79.226}{461.708}{entry}%
\TextAt{8}{109.796}{461.708}{is}%
\TextAt{8}{121.559}{461.708}{the}%
\TextAt{8}{141.716}{461.708}{matrix}%
\TextAt{13}{179.769}{461.708}{Q\Character{27}}%
\TextAt{3}{196.123}{466.046}{\Character{48}}%
\TextAt{21}{195.694}{458.599}{\Character{96}}%
\TextAt{13}{199.651}{461.708}{\Character{33}}%
\TextAt{3}{207.369}{466.046}{\Character{48}}%
\TextAt{21}{206.94}{458.752}{j}%
\TextAt{8}{211.323}{461.708}{,}%
\TextAt{8}{218.472}{461.708}{where}%
\TextAt{13}{252.281}{461.708}{Q}%
\TextAt{9}{264.873}{461.708}{\Character{61}}%
\TextAt{9}{277.283}{461.708}{\Character{40}}%
\TextAt{13}{281.851}{461.708}{\Character{27}}%
\TextAt{21}{288.504}{459.915}{k}%
\TextAt{13}{293.624}{461.708}{\Character{33}}%
\TextAt{21}{300.913}{459.915}{i}%
\TextAt{9}{304.295}{461.708}{\Character{41}}%
\TextAt{8}{308.847}{461.708}{.}%
\TextAt{8}{317.295}{461.708}{In}%
\TextAt{8}{331.917}{461.708}{other}%
\TextAt{8}{362.486}{461.708}{words,}%
\TextAt{13}{186.101}{414.106}{M}%
\TextAt{9}{201.996}{414.106}{\Character{61}}%
\TextAt{20}{214.421}{441.723}{\Character{48}}%
\TextAt{20}{214.421}{420.681}{\Character{66}}%
\TextAt{20}{214.421}{413.03}{\Character{64}}%
\TextAt{13}{224.882}{432.319}{Q}%
\TextAt{9}{269.528}{432.319}{0}%
\TextAt{8}{244.78}{419.072}{.}%
\TextAt{8}{249.36}{416.083}{.}%
\TextAt{8}{253.94}{413.094}{.}%
\TextAt{9}{226.591}{395.052}{0}%
\TextAt{13}{267.819}{395.052}{Q}%
\TextAt{20}{277.09}{441.723}{\Character{49}}%
\TextAt{20}{277.09}{420.681}{\Character{67}}%
\TextAt{20}{277.09}{413.03}{\Character{65}}%
\TextAt{20}{289.543}{441.723}{\Character{48}}%
\TextAt{20}{289.543}{420.681}{\Character{66}}%
\TextAt{20}{289.543}{413.03}{\Character{64}}%
\TextAt{13}{301.798}{432.319}{E}%
\TextAt{13}{311.147}{432.319}{\Character{27}}%
\TextAt{3}{318.247}{436.658}{\Character{48}}%
\TextAt{17}{317.818}{429.364}{1}%
\TextAt{13}{322.55}{432.319}{\Character{33}}%
\TextAt{3}{330.269}{436.658}{\Character{48}}%
\TextAt{17}{329.84}{429.364}{1}%
\TextAt{16}{346.329}{432.319}{\Character{1}}%
\TextAt{16}{351.625}{432.319}{\Character{1}}%
\TextAt{16}{356.933}{432.319}{\Character{1}}%
\ErrataAt{371.99}{432.319}{56.01}{$E\sigma_1^{\prime}\omega_{n^{\prime}}^{\prime}$}%
\TextAt{8}{316.559}{420.068}{.}%
\TextAt{8}{316.559}{416.083}{.}%
\TextAt{8}{316.559}{412.098}{.}%
\TextAt{8}{388.588}{420.068}{.}%
\TextAt{8}{388.588}{416.083}{.}%
\TextAt{8}{388.588}{412.098}{.}%
\ErrataAt{299.99}{395.052}{51.01}{$E\sigma_{n^{\prime}}^{\prime}\omega_1^{\prime}$}%
\TextAt{16}{351.625}{395.052}{\Character{1}}%
\TextAt{16}{356.933}{395.052}{\Character{1}}%
\TextAt{13}{370.239}{395.052}{E}%
\TextAt{13}{379.588}{395.052}{\Character{27}}%
\TextAt{3}{386.688}{399.39}{\Character{48}}%
\TextAt{21}{386.259}{391.888}{n}%
\TextAt{23}{391.397}{394.157}{\Character{48}}%
\TextAt{13}{394.58}{395.052}{\Character{33}}%
\TextAt{3}{402.298}{399.39}{\Character{48}}%
\TextAt{21}{401.869}{391.888}{n}%
\TextAt{23}{407.007}{394.157}{\Character{48}}%
\TextAt{20}{410.189}{441.723}{\Character{49}}%
\TextAt{20}{410.189}{420.681}{\Character{67}}%
\TextAt{20}{410.189}{413.03}{\Character{65}}%
\TextAt{13}{422.643}{414.106}{\Character{58}}%
\TextAt{8}{79.226}{368.02}{Here}%
\TextAt{13}{106.975}{368.02}{E}%
\TextAt{8}{120.359}{368.02}{denotes}%
\TextAt{8}{162.813}{368.02}{the}%
\TextAt{9}{183.09}{368.02}{\Character{40}}%
\TextAt{13}{187.643}{368.02}{n}%
\TextAt{16}{197.369}{368.02}{\Character{2}}%
\TextAt{13}{209.407}{368.02}{n}%
\TextAt{9}{216.394}{368.02}{\Character{41}}%
\TextAt{8}{220.947}{368.02}{-unit}%
\TextAt{8}{249.676}{368.02}{matrix.}%
\TextAt{8}{292.642}{368.02}{By}%
\TextAt{8}{311.126}{368.02}{changing}%
\TextAt{8}{360.358}{368.02}{indices}%
\TextAt{8}{398.903}{368.02}{the}%
\TextAt{8}{419.18}{368.02}{second}%
\TextAt{8}{457.079}{368.02}{matrix}%
\TextAt{8}{495.253}{368.02}{may}%
\TextAt{8}{520.743}{368.02}{be}%
\TextAt{8}{79.226}{350.974}{transformed}%
\TextAt{8}{144.657}{350.974}{to}%
\TextAt{8}{158.967}{350.974}{look}%
\TextAt{8}{184.336}{350.974}{like}%
\TextAt{8}{205.796}{350.974}{the}%
\TextAt{8}{225.952}{350.974}{first}%
\TextAt{8}{250.077}{350.974}{one.}%
\TextAt{8}{276.088}{350.974}{This}%
\TextAt{8}{302.808}{350.974}{yields}%
\TextAt{9}{182.429}{321.973}{\Character{1}}%
\TextAt{9}{195.484}{321.973}{\Character{61}}%
\TextAt{8}{207.93}{321.973}{det}%
\TextAt{9}{224.189}{321.973}{\Character{40}}%
\TextAt{13}{228.741}{321.973}{M}%
\TextAt{9}{241.315}{321.973}{\Character{41}}%
\TextAt{17}{245.867}{326.909}{2}%
\TextAt{9}{253.921}{321.973}{\Character{61}}%
\TextAt{8}{266.346}{321.973}{det}%
\TextAt{9}{282.605}{321.973}{\Character{40}}%
\TextAt{13}{287.157}{321.973}{Q}%
\TextAt{9}{296.429}{321.973}{\Character{41}}%
\TextAt{17}{300.981}{326.909}{2}%
\TextAt{21}{305.215}{326.909}{n}%
\TextAt{23}{310.353}{329.722}{\Character{48}}%
\TextAt{8}{315.528}{321.973}{det}%
\TextAt{9}{331.787}{321.973}{\Character{40}\Character{40}}%
\TextAt{13}{340.892}{321.973}{\Character{27}}%
\TextAt{3}{347.974}{326.909}{\Character{48}}%
\TextAt{21}{347.545}{319.018}{\Character{96}}%
\TextAt{13}{351.503}{321.973}{\Character{33}}%
\TextAt{3}{359.221}{326.909}{\Character{48}}%
\TextAt{21}{358.792}{319.018}{j}%
\TextAt{9}{363.174}{321.973}{\Character{41}\Character{41}}%
\TextAt{17}{372.279}{326.909}{2}%
\TextAt{21}{376.513}{326.909}{n}%
\TextAt{9}{385.47}{321.973}{\Character{61}}%
\TextAt{13}{397.896}{321.973}{d}%
\TextAt{21}{403.978}{326.909}{n}%
\TextAt{23}{409.117}{329.722}{\Character{48}}%
\TextAt{13}{412.299}{321.973}{d}%
\TextAt{3}{418.382}{326.909}{\Character{48}}%
\TextAt{21}{420.679}{326.909}{n}%
\TextAt{13}{426.315}{321.973}{\Character{58}}%
\Vector{q 0 G 0 g 0.399 w 452.334 -481.028 m 452.334 -472.958 l S Q}%
\Vector{q 0 G 0 g 0.399 w 452.534 -473.157 m 459.707 -473.157 l S Q}%
\Vector{q 0 G 0 g 0.399 w 452.534 -480.829 m 459.707 -480.829 l S Q}%
\Vector{q 0 G 0 g 0.399 w 459.906 -481.028 m 459.906 -472.958 l S Q}%
\TextAt{12}{79.226}{260.982}{Remark.}%
\TextAt{8}{133.967}{260.982}{It}%
\TextAt{8}{147.223}{260.982}{follows}%
\TextAt{8}{186.243}{260.982}{from}%
\TextAt{8}{214.468}{260.982}{the}%
\TextAt{8}{235.208}{260.982}{proof}%
\TextAt{8}{266.361}{260.982}{that}%
\TextAt{8}{292.314}{260.982}{the}%
\TextAt{8}{313.055}{260.982}{proposition}%
\TextAt{8}{375.172}{260.982}{is}%
\TextAt{8}{387.52}{260.982}{valid}%
\TextAt{8}{416.402}{260.982}{for}%
\TextAt{8}{434.87}{260.982}{arbitrary}%
\TextAt{8}{485.224}{260.982}{separable}%
\TextAt{8}{79.226}{243.936}{extensions}%
\TextAt{8}{135.916}{243.936}{(not}%
\TextAt{8}{161.593}{243.936}{necessarily}%
\TextAt{8}{220.231}{243.936}{Galois),}%
\TextAt{8}{264.337}{243.936}{if}%
\TextAt{8}{275.37}{243.936}{one}%
\TextAt{8}{297.138}{243.936}{assumes}%
\TextAt{8}{342.506}{243.936}{instead}%
\TextAt{8}{383.2}{243.936}{of}%
\TextAt{13}{396.839}{243.936}{L}%
\TextAt{16}{407.67}{243.936}{\Character{92}}%
\TextAt{13}{418.508}{243.936}{L}%
\TextAt{3}{426.472}{248.275}{\Character{48}}%
\TextAt{9}{433.114}{243.936}{\Character{61}}%
\TextAt{13}{446.066}{243.936}{K}%
\TextAt{8}{461.09}{243.936}{that}%
\TextAt{13}{486.768}{243.936}{L}%
\TextAt{8}{498.938}{243.936}{and}%
\TextAt{13}{522.01}{243.936}{L}%
\TextAt{3}{529.974}{248.275}{\Character{48}}%
\TextAt{8}{79.226}{226.89}{are}%
\TextAt{8}{98.737}{226.89}{linearly}%
\TextAt{8}{140.688}{226.89}{disjoint.}%
\TextAt{8}{97.159}{203.867}{The}%
\TextAt{8}{122.755}{203.867}{chief}%
\TextAt{8}{151.603}{203.867}{application}%
\TextAt{8}{213.638}{203.867}{of}%
\TextAt{8}{228.51}{203.867}{our}%
\TextAt{8}{250.867}{203.867}{considerations}%
\TextAt{8}{328.635}{203.867}{on}%
\TextAt{8}{346.436}{203.867}{integrality}%
\TextAt{8}{404.251}{203.867}{will}%
\TextAt{8}{427.899}{203.867}{concern}%
\TextAt{8}{472.36}{203.867}{the}%
\TextAt{8}{494.059}{203.867}{integral}%
\TextAt{8}{79.226}{186.821}{closure}%
\TextAt{14}{118.706}{186.821}{o}%
\TextAt{21}{124.55}{185.027}{K}%
\NativeMathAt{136.751}{186.821}{11.955}{\subseteq}%
\TextAt{13}{150.049}{186.821}{K}%
\TextAt{8}{165.162}{186.821}{of}%
\TextAt{18}{178.89}{186.821}{Z}%
\NativeMathAt{190.859}{186.821}{11.955}{\subseteq}%
\TextAt{18}{204.157}{186.821}{Q}%
\TextAt{8}{217.751}{186.821}{in}%
\TextAt{8}{231.801}{186.821}{an}%
\TextAt{8}{248.458}{186.821}{algebraic}%
\TextAt{8}{298.291}{186.821}{number}%
\TextAt{8}{341.285}{186.821}{field}%
\TextAt{13}{367.04}{186.821}{K}%
\TextAt{8}{377.858}{186.821}{.}%
\TextAt{8}{387.5}{186.821}{By}%
\TextAt{8}{406.261}{186.821}{proposition}%
\TextAt{8}{468.192}{186.821}{(2.10),}%
\TextAt{8}{505.775}{186.821}{every}%
\ErrataAt{79.19}{169.787}{455.81}{nonzero finitely generated $\mathfrak o_K$-submodule $\mathfrak a$ of $K$ admits a $\mathbb Z$-basis $\alpha_1,\ldots,\alpha_n$,}%
\TextAt{14}{251.515}{140.774}{a}%
\TextAt{9}{260.819}{140.774}{\Character{61}}%
\TextAt{18}{273.245}{140.774}{Z}%
\TextAt{13}{281.215}{140.774}{\Character{11}}%
\TextAt{17}{288.661}{138.98}{1}%
\TextAt{9}{296.05}{140.774}{\Character{43}}%
\TextAt{16}{307.811}{140.774}{\Character{1}}%
\TextAt{16}{313.107}{140.774}{\Character{1}}%
\TextAt{16}{318.415}{140.774}{\Character{1}}%
\TextAt{9}{324.415}{140.774}{\Character{43}}%
\TextAt{18}{336.177}{140.774}{Z}%
\TextAt{13}{344.147}{140.774}{\Character{11}}%
\TextAt{21}{351.592}{138.98}{n}%
\TextAt{13}{357.229}{140.774}{\Character{58}}%
\TextAt{8}{79.226}{111.772}{The}%
\TextAt{8}{103.28}{111.772}{discriminant}%
\TextAt{13}{234.259}{94.727}{d}%
\TextAt{9}{240.342}{94.727}{\Character{40}}%
\TextAt{13}{244.894}{94.727}{\Character{11}}%
\TextAt{17}{252.34}{92.933}{1}%
\TextAt{13}{257.072}{94.727}{\Character{59}}%
\TextAt{13}{262.296}{94.727}{\Character{58}}%
\TextAt{13}{267.533}{94.727}{\Character{58}}%
\TextAt{13}{272.757}{94.727}{\Character{58}}%
\TextAt{13}{277.994}{94.727}{\Character{59}}%
\TextAt{13}{283.23}{94.727}{\Character{11}}%
\TextAt{21}{290.738}{92.933}{n}%
\TextAt{9}{296.375}{94.727}{\Character{41}}%
\TextAt{9}{304.23}{94.727}{\Character{61}}%
\TextAt{8}{316.673}{94.727}{det}%
\TextAt{9}{332.932}{94.727}{\Character{40}\Character{40}}%
\TextAt{13}{342.037}{94.727}{\Character{27}}%
\TextAt{21}{348.691}{92.933}{i}%
\TextAt{13}{352.072}{94.727}{\Character{11}}%
\TextAt{21}{359.517}{92.933}{j}%
\TextAt{9}{363.9}{94.727}{\Character{41}\Character{41}}%
\TextAt{17}{373.004}{99.663}{2}%
\TextAt{8}{300.14}{43.033}{13}%
\Destination{78.23}{809.815}{nav.page.13}%
\Destination{79.36}{735.555}{nav.print.14}%
\Hotlink{174.069}{701.355}{25.078}{11.047}{nav.number.1.2.3}%
\Hotlink{467.692}{184.655}{30.936}{11.047}{nav.number.1.2.10}%
\Destination{79.226}{449.319}{errata.67}%
\Destination{79.226}{193}{errata.186}%
}

% Page 26 (body page 14)
\Page{%
\TextAt{2}{79.226}{796.943}{I.}%
\TextAt{2}{89.83}{796.943}{Algebraic}%
\TextAt{2}{138.648}{796.943}{Integers}%
\Vector{q 0 G 0 g 0.299 w 7.226 18.729 m 460.769 18.729 l S Q}%
\TextAt{8}{79.226}{761.124}{is}%
\TextAt{8}{92.161}{761.124}{independent}%
\TextAt{8}{159.66}{761.124}{of}%
\TextAt{8}{174.161}{761.124}{the}%
\TextAt{8}{195.489}{761.124}{choice}%
\TextAt{8}{231.45}{761.124}{of}%
\TextAt{8}{245.952}{761.124}{a}%
\TextAt{18}{256.878}{761.124}{Z}%
\TextAt{8}{264.848}{761.124}{-basis;}%
\TextAt{8}{302.495}{761.124}{if}%
\TextAt{13}{314.39}{761.124}{\Character{11}}%
\TextAt{3}{321.912}{765.462}{\Character{48}}%
\TextAt{17}{321.836}{758.168}{1}%
\TextAt{13}{326.568}{761.124}{\Character{59}}%
\TextAt{13}{331.792}{761.124}{\Character{58}}%
\TextAt{13}{337.029}{761.124}{\Character{58}}%
\TextAt{13}{342.253}{761.124}{\Character{58}}%
\TextAt{13}{347.49}{761.124}{\Character{59}}%
\TextAt{13}{352.726}{761.124}{\Character{11}}%
\TextAt{3}{360.311}{765.462}{\Character{48}}%
\TextAt{21}{360.234}{758.168}{n}%
\TextAt{8}{370.94}{761.124}{is}%
\TextAt{8}{383.875}{761.124}{another}%
\TextAt{8}{427.977}{761.124}{basis,}%
\TextAt{8}{461.433}{761.124}{then}%
\TextAt{8}{489.265}{761.124}{the}%
\TextAt{8}{510.593}{761.124}{base}%
\TextAt{8}{79.226}{744.078}{change}%
\TextAt{8}{118.847}{744.078}{matrix}%
\TextAt{13}{157.822}{744.078}{T}%
\TextAt{9}{171.2}{744.078}{\Character{61}}%
\TextAt{9}{185.188}{744.078}{\Character{40}}%
\TextAt{13}{189.748}{744.078}{a}%
\TextAt{21}{195.893}{742.284}{ij}%
\TextAt{9}{203.158}{744.078}{\Character{41}}%
\TextAt{8}{207.711}{744.078}{,}%
\TextAt{13}{216.012}{744.078}{\Character{11}}%
\TextAt{3}{223.533}{748.416}{\Character{48}}%
\TextAt{21}{223.457}{741.122}{i}%
\TextAt{9}{231.73}{744.078}{\Character{61}}%
\TextAt{20}{245.726}{753.044}{\Character{80}}%
\TextAt{21}{258.345}{740.591}{j}%
\TextAt{13}{264.72}{744.078}{a}%
\TextAt{21}{270.865}{742.284}{ij}%
\TextAt{13}{278.13}{744.078}{\Character{11}}%
\TextAt{21}{285.576}{742.284}{j}%
\TextAt{8}{289.958}{744.078}{,}%
\TextAt{8}{298.259}{744.078}{as}%
\TextAt{8}{313.551}{744.078}{well}%
\TextAt{8}{338.203}{744.078}{as}%
\TextAt{8}{353.495}{744.078}{its}%
\TextAt{8}{370.736}{744.078}{inverse,}%
\TextAt{8}{413.898}{744.078}{has}%
\TextAt{8}{435.693}{744.078}{integral}%
\TextAt{8}{479.223}{744.078}{entries.}%
\TextAt{8}{523.994}{744.078}{It}%
\TextAt{8}{79.226}{727.032}{therefore}%
\TextAt{8}{128.326}{727.032}{has}%
\TextAt{8}{149.2}{727.032}{determinant}%
\NativeMathAt{215.216}{727.032}{11.955}{\pm}%
\TextAt{9}{224.515}{727.032}{1}%
\TextAt{8}{230.368}{727.032}{,}%
\TextAt{8}{237.517}{727.032}{so}%
\TextAt{8}{251.887}{727.032}{that}%
\TextAt{8}{277.256}{727.032}{indeed}%
\TextAt{13}{169.332}{698.031}{d}%
\TextAt{9}{175.415}{698.031}{\Character{40}}%
\TextAt{13}{179.967}{698.031}{\Character{11}}%
\TextAt{3}{187.489}{702.967}{\Character{48}}%
\TextAt{17}{187.413}{695.075}{1}%
\TextAt{13}{192.145}{698.031}{\Character{59}}%
\TextAt{13}{197.369}{698.031}{\Character{58}}%
\TextAt{13}{202.606}{698.031}{\Character{58}}%
\TextAt{13}{207.842}{698.031}{\Character{58}}%
\TextAt{13}{213.067}{698.031}{\Character{59}}%
\TextAt{13}{218.303}{698.031}{\Character{11}}%
\TextAt{3}{225.888}{702.967}{\Character{48}}%
\TextAt{21}{225.812}{695.075}{n}%
\TextAt{9}{231.448}{698.031}{\Character{41}}%
\TextAt{9}{239.303}{698.031}{\Character{61}}%
\TextAt{8}{251.747}{698.031}{det}%
\TextAt{9}{268.006}{698.031}{\Character{40}}%
\TextAt{13}{272.558}{698.031}{T}%
\TextAt{9}{281.045}{698.031}{\Character{41}}%
\TextAt{17}{285.597}{702.967}{2}%
\TextAt{13}{290.329}{698.031}{d}%
\TextAt{9}{296.412}{698.031}{\Character{40}}%
\TextAt{13}{300.964}{698.031}{\Character{11}}%
\TextAt{17}{308.41}{696.237}{1}%
\TextAt{13}{313.142}{698.031}{\Character{59}}%
\TextAt{13}{318.366}{698.031}{\Character{58}}%
\TextAt{13}{323.603}{698.031}{\Character{58}}%
\TextAt{13}{328.839}{698.031}{\Character{58}}%
\TextAt{13}{334.064}{698.031}{\Character{59}}%
\TextAt{13}{339.3}{698.031}{\Character{11}}%
\TextAt{21}{346.809}{696.237}{n}%
\TextAt{9}{352.445}{698.031}{\Character{41}}%
\TextAt{9}{360.3}{698.031}{\Character{61}}%
\TextAt{13}{372.744}{698.031}{d}%
\TextAt{9}{378.826}{698.031}{\Character{40}}%
\TextAt{13}{383.379}{698.031}{\Character{11}}%
\TextAt{17}{390.824}{696.237}{1}%
\TextAt{13}{395.557}{698.031}{\Character{59}}%
\TextAt{13}{400.781}{698.031}{\Character{58}}%
\TextAt{13}{406.018}{698.031}{\Character{58}}%
\TextAt{13}{411.242}{698.031}{\Character{58}}%
\TextAt{13}{416.479}{698.031}{\Character{59}}%
\TextAt{13}{421.715}{698.031}{\Character{11}}%
\TextAt{21}{429.223}{696.237}{n}%
\TextAt{9}{434.859}{698.031}{\Character{41}}%
\TextAt{13}{439.412}{698.031}{\Character{58}}%
\TextAt{8}{79.226}{669.029}{We}%
\TextAt{8}{99.371}{669.029}{may}%
\TextAt{8}{124.739}{669.029}{therefore}%
\TextAt{8}{173.839}{669.029}{write}%
\TextAt{13}{252.58}{651.983}{d}%
\TextAt{9}{258.662}{651.983}{\Character{40}}%
\TextAt{14}{263.215}{651.983}{a}%
\TextAt{9}{269.198}{651.983}{\Character{41}}%
\TextAt{9}{277.053}{651.983}{\Character{61}}%
\TextAt{13}{289.496}{651.983}{d}%
\TextAt{9}{295.579}{651.983}{\Character{40}}%
\TextAt{13}{300.131}{651.983}{\Character{11}}%
\TextAt{17}{307.577}{650.19}{1}%
\TextAt{13}{312.309}{651.983}{\Character{59}}%
\TextAt{13}{317.533}{651.983}{\Character{58}}%
\TextAt{13}{322.77}{651.983}{\Character{58}}%
\TextAt{13}{328.006}{651.983}{\Character{58}}%
\TextAt{13}{333.231}{651.983}{\Character{59}}%
\TextAt{13}{338.467}{651.983}{\Character{11}}%
\TextAt{21}{345.976}{650.19}{n}%
\TextAt{9}{351.612}{651.983}{\Character{41}}%
\TextAt{13}{356.164}{651.983}{\Character{58}}%
\TextAt{8}{79.226}{628.462}{In}%
\TextAt{8}{94.272}{628.462}{the}%
\TextAt{8}{114.853}{628.462}{special}%
\TextAt{8}{153.379}{628.462}{case}%
\TextAt{8}{178.574}{628.462}{of}%
\TextAt{8}{192.329}{628.462}{an}%
\TextAt{8}{209.013}{628.462}{integral}%
\TextAt{8}{252.045}{628.462}{basis}%
\TextAt{13}{281.21}{628.462}{\Character{33}}%
\TextAt{17}{288.499}{626.669}{1}%
\TextAt{13}{293.232}{628.462}{\Character{59}}%
\TextAt{13}{298.456}{628.462}{\Character{58}}%
\TextAt{13}{303.693}{628.462}{\Character{58}}%
\TextAt{13}{308.917}{628.462}{\Character{58}}%
\TextAt{13}{314.154}{628.462}{\Character{59}}%
\TextAt{13}{319.39}{628.462}{\Character{33}}%
\TextAt{21}{326.742}{626.669}{n}%
\TextAt{8}{336.7}{628.462}{of}%
\TextAt{14}{350.455}{628.462}{o}%
\TextAt{21}{356.299}{626.669}{K}%
\TextAt{8}{368.824}{628.462}{we}%
\TextAt{8}{386.475}{628.462}{obtain}%
\TextAt{8}{423.327}{628.462}{the}%
\TextAt{12}{443.908}{628.462}{discriminant}%
\TextAt{12}{521.926}{628.462}{of}%
\TextAt{12}{79.226}{611.416}{the}%
\TextAt{12}{102.551}{611.416}{algebraic}%
\TextAt{12}{159.409}{611.416}{number}%
\TextAt{12}{208.784}{611.416}{field}%
\TextAt{13}{238.086}{611.416}{K}%
\TextAt{8}{248.904}{611.416}{,}%
\TextAt{13}{233.688}{582.415}{d}%
\TextAt{21}{239.771}{580.621}{K}%
\TextAt{9}{251.294}{582.415}{\Character{61}}%
\TextAt{13}{263.72}{582.415}{d}%
\TextAt{9}{269.803}{582.415}{\Character{40}}%
\TextAt{14}{274.355}{582.415}{o}%
\TextAt{21}{280.199}{580.621}{K}%
\TextAt{9}{288.402}{582.415}{\Character{41}}%
\TextAt{9}{296.257}{582.415}{\Character{61}}%
\TextAt{13}{308.7}{582.415}{d}%
\TextAt{9}{314.783}{582.415}{\Character{40}}%
\TextAt{13}{319.335}{582.415}{\Character{33}}%
\TextAt{17}{326.625}{580.621}{1}%
\TextAt{13}{331.357}{582.415}{\Character{59}}%
\TextAt{13}{336.581}{582.415}{\Character{58}}%
\TextAt{13}{341.818}{582.415}{\Character{58}}%
\TextAt{13}{347.042}{582.415}{\Character{58}}%
\TextAt{13}{352.279}{582.415}{\Character{59}}%
\TextAt{13}{357.515}{582.415}{\Character{33}}%
\TextAt{21}{364.867}{580.621}{n}%
\TextAt{9}{370.503}{582.415}{\Character{41}}%
\TextAt{13}{375.056}{582.415}{\Character{58}}%
\TextAt{8}{79.226}{553.414}{In}%
\TextAt{8}{93.847}{553.414}{general,}%
\TextAt{8}{137.424}{553.414}{one}%
\TextAt{8}{158.883}{553.414}{has}%
\TextAt{8}{179.757}{553.414}{the}%
\TextAt{12}{79.226}{527.401}{Proposition}%
\TextAt{12}{152.391}{527.401}{(2.12).}%
\TextAt{8}{194.832}{527.401}{If}%
\TextAt{14}{206.918}{527.401}{a}%
\NativeMathAt{216.894}{527.401}{11.955}{\subseteq}%
\TextAt{14}{230.185}{527.401}{a}%
\TextAt{3}{236.168}{531.74}{\Character{48}}%
\TextAt{8}{243.255}{527.401}{are}%
\TextAt{8}{263.16}{527.401}{two}%
\TextAt{8}{285.671}{527.401}{nonzero}%
\TextAt{8}{329.642}{527.401}{finitely}%
\TextAt{8}{369.381}{527.401}{generated}%
\TextAt{14}{423.107}{527.401}{o}%
\TextAt{21}{428.952}{525.608}{K}%
\TextAt{8}{437.154}{527.401}{-submodules}%
\TextAt{8}{504.976}{527.401}{of}%
\TextAt{13}{518.7}{527.401}{K}%
\TextAt{8}{529.518}{527.401}{,}%
\TextAt{8}{79.226}{510.355}{then}%
\TextAt{8}{105.886}{510.355}{the}%
\TextAt{8}{126.043}{510.355}{index}%
\TextAt{9}{157.58}{510.355}{\Character{40}}%
\TextAt{14}{162.133}{510.355}{a}%
\TextAt{3}{168.116}{514.694}{\Character{48}}%
\TextAt{9}{174.232}{510.355}{\Character{58}}%
\TextAt{14}{180.804}{510.355}{a}%
\TextAt{9}{186.787}{510.355}{\Character{41}}%
\TextAt{8}{195.237}{510.355}{is}%
\TextAt{8}{207.001}{510.355}{finite}%
\TextAt{8}{236.913}{510.355}{and}%
\TextAt{8}{259.675}{510.355}{satisfies}%
\TextAt{13}{254.685}{481.354}{d}%
\TextAt{9}{260.768}{481.354}{\Character{40}}%
\TextAt{14}{265.32}{481.354}{a}%
\TextAt{9}{271.303}{481.354}{\Character{41}}%
\TextAt{9}{279.158}{481.354}{\Character{61}}%
\TextAt{9}{291.579}{481.354}{\Character{40}}%
\TextAt{14}{296.154}{481.354}{a}%
\TextAt{3}{302.137}{486.29}{\Character{48}}%
\TextAt{9}{308.253}{481.354}{\Character{58}}%
\TextAt{14}{314.826}{481.354}{a}%
\TextAt{9}{320.809}{481.354}{\Character{41}}%
\TextAt{17}{325.361}{486.29}{2}%
\TextAt{13}{330.093}{481.354}{d}%
\TextAt{9}{336.176}{481.354}{\Character{40}}%
\TextAt{14}{340.728}{481.354}{a}%
\TextAt{3}{346.711}{486.29}{\Character{48}}%
\TextAt{9}{349.507}{481.354}{\Character{41}}%
\TextAt{13}{354.059}{481.354}{\Character{58}}%
\ErrataAt{79.19}{443.387}{455.81}{Let $(\alpha_i)$ and $(\alpha_i^{\prime})$ be $\mathbb Z$-bases of $\mathfrak a$ and $\mathfrak a^{\prime}$, respectively, and write}%
\ErrataAt{79.19}{426.341}{455.81}{$\alpha_i=\sum_jt_{ij}\alpha_j^{\prime}$. Then $(\mathfrak a^{\prime}:\mathfrak a)=|\det(t_{ij})|$ by the index formula for lattices.}%
\ErrataAt{79.19}{409.295}{455.81}{Changing the basis in the discriminant multiplies it by $\det(t_{ij})^2$, as required.}%
\TextAt{12}{79.226}{377.947}{Exercise}%
\TextAt{12}{132.427}{377.947}{I.2.1.}%
\TextAt{8}{166.021}{377.947}{Is}%
\TextAt{9}{179.948}{386.035}{3}%
\TextAt{9}{188.448}{386.035}{\Character{43}}%
\TextAt{9}{200.2}{386.035}{2}%
\TextAt{16}{206.072}{395.925}{\Character{112}}%
\Vector{q 0 G 0 g 0.478 w 144.035 -377.836 m 149.888 -377.836 l S Q}%
\TextAt{9}{216.035}{386.035}{6}%
\Vector{q 0 G 0 g 0.478 w 107.948 -393.064 m 149.888 -393.064 l S Q}%
\TextAt{9}{182.778}{368.417}{1}%
\TextAt{16}{191.288}{368.417}{\Character{0}}%
\TextAt{16}{203.243}{378.306}{\Character{112}}%
\Vector{q 0 G 0 g 0.478 w 141.206 -395.455 m 147.059 -395.455 l S Q}%
\TextAt{9}{213.206}{368.417}{6}%
\TextAt{8}{226.981}{377.947}{an}%
\TextAt{8}{243.24}{377.947}{algebraic}%
\TextAt{8}{292.675}{377.947}{integer?}%
\TextAt{12}{79.226}{342.473}{Exercise}%
\TextAt{12}{133.074}{342.473}{I.2.2.}%
\TextAt{8}{167.231}{342.473}{Show}%
\TextAt{8}{198.686}{342.473}{that,}%
\TextAt{8}{228.011}{342.473}{if}%
\TextAt{8}{239.298}{342.473}{the}%
\TextAt{8}{260.017}{342.473}{integral}%
\TextAt{8}{303.188}{342.473}{domain}%
\TextAt{13}{345.379}{342.473}{A}%
\TextAt{8}{358.615}{342.473}{is}%
\TextAt{8}{370.942}{342.473}{integrally}%
\TextAt{8}{423.545}{342.473}{closed,}%
\TextAt{8}{462.027}{342.473}{then}%
\TextAt{8}{489.251}{342.473}{so}%
\TextAt{8}{504.184}{342.473}{is}%
\TextAt{8}{516.51}{342.473}{the}%
\TextAt{8}{79.226}{325.427}{polynomial}%
\TextAt{8}{139.719}{325.427}{ring}%
\TextAt{13}{163.785}{325.427}{A}%
\TextAt{9}{172.56}{325.427}{\Character{91}}%
\TextAt{13}{175.812}{325.427}{t}%
\TextAt{9}{180.039}{325.427}{\Character{93}}%
\TextAt{8}{183.291}{325.427}{.}%
\TextAt{12}{79.226}{293.437}{Exercise}%
\TextAt{12}{133.43}{293.437}{I.2.3.}%
\TextAt{8}{167.897}{293.437}{In}%
\TextAt{8}{183.39}{293.437}{the}%
\TextAt{8}{204.419}{293.437}{polynomial}%
\TextAt{8}{265.785}{293.437}{ring}%
\TextAt{13}{290.723}{293.437}{A}%
\TextAt{9}{304.306}{293.437}{\Character{61}}%
\TextAt{18}{318.219}{293.437}{Q}%
\TextAt{9}{327.517}{293.437}{\Character{91}}%
\TextAt{13}{330.769}{293.437}{X}%
\TextAt{13}{340.764}{293.437}{\Character{59}}%
\TextAt{13}{345.988}{293.437}{Y}%
\TextAt{9}{355.411}{293.437}{\Character{93}}%
\TextAt{8}{358.663}{293.437}{,}%
\TextAt{8}{366.902}{293.437}{consider}%
\TextAt{8}{413.36}{293.437}{the}%
\TextAt{8}{434.389}{293.437}{principal}%
\TextAt{8}{484.038}{293.437}{ideal}%
\TextAt{14}{512.874}{293.437}{p}%
\TextAt{9}{523.665}{293.437}{\Character{61}}%
\TextAt{9}{79.226}{276.391}{\Character{40}}%
\TextAt{13}{83.778}{276.391}{X}%
\TextAt{17}{94.434}{280.729}{2}%
\TextAt{16}{101.823}{276.391}{\Character{0}}%
\TextAt{13}{113.778}{276.391}{Y}%
\TextAt{17}{123.171}{280.729}{3}%
\TextAt{9}{127.903}{276.391}{\Character{41}}%
\TextAt{8}{132.456}{276.391}{.}%
\TextAt{8}{140.904}{276.391}{Show}%
\TextAt{8}{171.796}{276.391}{that}%
\TextAt{14}{197.165}{276.391}{p}%
\TextAt{8}{207.045}{276.391}{is}%
\TextAt{8}{218.809}{276.391}{a}%
\TextAt{8}{228.565}{276.391}{prime}%
\TextAt{8}{261.728}{276.391}{ideal,}%
\TextAt{8}{292.943}{276.391}{but}%
\TextAt{13}{314.403}{276.391}{A}%
\TextAt{9}{323.178}{276.391}{\Character{47}}%
\TextAt{14}{329.031}{276.391}{p}%
\TextAt{8}{338.912}{276.391}{is}%
\TextAt{8}{350.676}{276.391}{not}%
\TextAt{8}{371.489}{276.391}{integrally}%
\TextAt{8}{423.53}{276.391}{closed.}%
\TextAt{12}{79.226}{244.401}{Exercise}%
\TextAt{12}{132.829}{244.401}{I.2.4.}%
\TextAt{8}{166.773}{244.401}{Let}%
\TextAt{13}{188.092}{244.401}{D}%
\TextAt{8}{202.384}{244.401}{be}%
\TextAt{8}{218.658}{244.401}{a}%
\TextAt{8}{228.763}{244.401}{squarefree}%
\TextAt{8}{284.453}{244.401}{rational}%
\TextAt{8}{328.392}{244.401}{integer}%
\NativeMathAt{367.44}{244.401}{11.955}{\neq}%
\TextAt{9}{380.447}{244.401}{0}%
\TextAt{13}{386.315}{244.401}{\Character{59}}%
\TextAt{9}{391.559}{244.401}{1}%
\TextAt{8}{401.659}{244.401}{and}%
\TextAt{13}{424.772}{244.401}{d}%
\TextAt{8}{435.101}{244.401}{the}%
\TextAt{8}{455.608}{244.401}{discriminant}%
\TextAt{8}{523.337}{244.401}{of}%
\TextAt{8}{79.226}{227.355}{the}%
\TextAt{8}{99.383}{227.355}{quadratic}%
\TextAt{8}{151.746}{227.355}{number}%
\TextAt{8}{194.342}{227.355}{field}%
\TextAt{13}{219.699}{227.355}{K}%
\TextAt{9}{233.838}{227.355}{\Character{61}}%
\TextAt{18}{246.263}{227.355}{Q}%
\TextAt{9}{255.562}{227.355}{\Character{40}}%
\TextAt{16}{260.114}{237.477}{\Character{112}}%
\Vector{q 0 G 0 g 0.478 w 198.077 -536.284 m 208.121 -536.284 l S Q}%
\TextAt{13}{270.077}{227.355}{D}%
\TextAt{9}{280.121}{227.355}{\Character{41}}%
\TextAt{8}{284.674}{227.355}{.}%
\TextAt{8}{293.122}{227.355}{Show}%
\TextAt{8}{324.014}{227.355}{that}%
\TextAt{13}{148.592}{198.354}{d}%
\TextAt{9}{157.995}{198.354}{\Character{61}}%
\TextAt{13}{170.421}{198.354}{D}%
\TextAt{13}{180.463}{198.354}{\Character{59}}%
\TextAt{8}{197.665}{198.354}{if}%
\TextAt{13}{208.389}{198.354}{D}%
\NativeMathAt{221.754}{198.354}{11.955}{\equiv}%
\TextAt{9}{234.373}{198.354}{1}%
\TextAt{8}{243.547}{198.354}{mod}%
\TextAt{9}{268.985}{198.354}{4}%
\TextAt{13}{274.838}{198.354}{\Character{59}}%
\TextAt{13}{303.973}{198.354}{d}%
\TextAt{9}{313.396}{198.354}{\Character{61}}%
\TextAt{9}{325.805}{198.354}{4}%
\TextAt{13}{331.674}{198.354}{D}%
\TextAt{13}{341.716}{198.354}{\Character{59}}%
\TextAt{8}{358.918}{198.354}{if}%
\TextAt{13}{369.642}{198.354}{D}%
\NativeMathAt{383.007}{198.354}{11.955}{\equiv}%
\TextAt{9}{395.627}{198.354}{2}%
\TextAt{8}{405.377}{198.354}{or}%
\TextAt{9}{419.687}{198.354}{3}%
\TextAt{8}{428.861}{198.354}{mod}%
\TextAt{9}{454.299}{198.354}{4}%
\TextAt{13}{460.152}{198.354}{\Character{59}}%
\TextAt{8}{79.226}{169.353}{and}%
\TextAt{8}{101.45}{169.353}{that}%
\TextAt{8}{126.28}{169.353}{an}%
\TextAt{8}{142}{169.353}{integral}%
\TextAt{8}{184.069}{169.353}{basis}%
\TextAt{8}{212.27}{169.353}{of}%
\TextAt{13}{225.061}{169.353}{K}%
\TextAt{8}{239.238}{169.353}{is}%
\TextAt{8}{250.462}{169.353}{given}%
\TextAt{8}{280.493}{169.353}{by}%
\TextAt{16}{296.213}{169.353}{\Character{102}}%
\TextAt{9}{302.191}{169.353}{1}%
\TextAt{13}{308.044}{169.353}{\Character{59}}%
\TextAt{16}{313.288}{179.475}{\Character{112}}%
\Vector{q 0 G 0 g 0.478 w 251.25 -594.286 m 261.295 -594.286 l S Q}%
\TextAt{13}{323.25}{169.353}{D}%
\TextAt{16}{333.295}{169.353}{\Character{103}}%
\TextAt{8}{342.631}{169.353}{in}%
\TextAt{8}{355.745}{169.353}{the}%
\TextAt{8}{375.362}{169.353}{second}%
\TextAt{8}{412.602}{169.353}{case,}%
\TextAt{8}{440.193}{169.353}{by}%
\TextAt{16}{455.913}{169.353}{\Character{102}}%
\TextAt{9}{461.891}{169.353}{1}%
\TextAt{13}{467.744}{169.353}{\Character{59}}%
\TextAt{17}{474.184}{174.06}{1}%
\Vector{q 0 G 0 g 0.478 w 402.184 -601.659 m 406.418 -601.659 l S Q}%
\TextAt{17}{474.184}{165.23}{2}%
\TextAt{9}{479.613}{169.353}{\Character{40}1}%
\TextAt{9}{491.556}{169.353}{\Character{43}}%
\TextAt{16}{502.232}{179.475}{\Character{112}}%
\Vector{q 0 G 0 g 0.478 w 440.195 -594.286 m 450.24 -594.286 l S Q}%
\TextAt{13}{512.195}{169.353}{D}%
\TextAt{9}{522.24}{169.353}{\Character{41}}%
\TextAt{16}{526.792}{169.353}{\Character{103}}%
\TextAt{8}{79.226}{152.307}{in}%
\TextAt{8}{92.879}{152.307}{the}%
\TextAt{8}{113.035}{152.307}{first}%
\TextAt{8}{137.161}{152.307}{case,}%
\TextAt{8}{165.184}{152.307}{and}%
\TextAt{8}{187.946}{152.307}{by}%
\TextAt{16}{204.206}{152.307}{\Character{102}}%
\TextAt{9}{210.183}{152.307}{1}%
\TextAt{13}{216.036}{152.307}{\Character{59}}%
\TextAt{17}{222.476}{157.014}{1}%
\Vector{q 0 G 0 g 0.478 w 150.476 -618.705 m 154.71 -618.705 l S Q}%
\TextAt{17}{222.476}{148.184}{2}%
\TextAt{9}{227.905}{152.307}{\Character{40}}%
\TextAt{13}{232.458}{152.307}{d}%
\TextAt{9}{241.197}{152.307}{\Character{43}}%
\TextAt{16}{252.958}{162.495}{\Character{112}}%
\Vector{q 0 G 0 g 0.478 w 190.921 -611.266 m 197.004 -611.266 l S Q}%
\TextAt{13}{262.921}{152.307}{d}%
\TextAt{9}{269.004}{152.307}{\Character{41}}%
\TextAt{16}{273.556}{152.307}{\Character{103}}%
\TextAt{8}{283.431}{152.307}{in}%
\TextAt{8}{297.084}{152.307}{both}%
\TextAt{8}{324.724}{152.307}{cases.}%
\TextAt{12}{79.226}{120.317}{Exercise}%
\TextAt{12}{132.427}{120.317}{I.2.5.}%
\TextAt{8}{166.021}{120.317}{Show}%
\TextAt{8}{196.913}{120.317}{that}%
\TextAt{16}{222.282}{120.317}{\Character{102}}%
\TextAt{9}{228.259}{120.317}{1}%
\TextAt{13}{234.112}{120.317}{\Character{59}}%
\TextAt{24}{242.677}{125.872}{3}%
\TextAt{16}{239.689}{130.206}{\Character{112}}%
\Vector{q 0 G 0 g 0.478 w 177.651 -643.555 m 183.504 -643.555 l S Q}%
\TextAt{9}{249.651}{120.317}{2}%
\TextAt{13}{255.504}{120.317}{\Character{59}}%
\TextAt{24}{264.069}{126.403}{3}%
\TextAt{16}{261.08}{130.649}{\Character{112}}%
\Vector{q 0 G 0 g 0.478 w 199.043 -643.112 m 209.628 -643.112 l S Q}%
\TextAt{9}{271.043}{120.317}{2}%
\TextAt{17}{276.896}{123.77}{2}%
\TextAt{16}{281.628}{120.317}{\Character{103}}%
\TextAt{8}{291.503}{120.317}{is}%
\TextAt{8}{303.267}{120.317}{an}%
\TextAt{8}{319.526}{120.317}{integral}%
\TextAt{8}{362.135}{120.317}{basis}%
\TextAt{8}{390.875}{120.317}{of}%
\TextAt{18}{404.205}{120.317}{Q}%
\TextAt{9}{413.503}{120.317}{\Character{40}}%
\TextAt{24}{421.376}{125.872}{3}%
\TextAt{16}{418.388}{130.206}{\Character{112}}%
\Vector{q 0 G 0 g 0.478 w 356.35 -643.555 m 362.203 -643.555 l S Q}%
\TextAt{9}{428.35}{120.317}{2\Character{41}}%
\TextAt{8}{438.756}{120.317}{.}%
\TextAt{12}{79.226}{88.327}{Exercise}%
\TextAt{12}{132.427}{88.327}{I.2.6.}%
\TextAt{8}{166.021}{88.327}{Show}%
\TextAt{8}{196.913}{88.327}{that}%
\TextAt{16}{222.282}{88.327}{\Character{102}}%
\TextAt{9}{228.259}{88.327}{1}%
\TextAt{13}{234.112}{88.327}{\Character{59}}%
\TextAt{13}{239.336}{88.327}{\Character{18}}%
\TextAt{13}{245.111}{88.327}{\Character{59}}%
\TextAt{17}{251.576}{93.034}{1}%
\Vector{q 0 G 0 g 0.478 w 179.576 -682.684 m 183.811 -682.684 l S Q}%
\TextAt{17}{251.576}{84.204}{2}%
\TextAt{9}{257.006}{88.327}{\Character{40}}%
\TextAt{13}{261.559}{88.327}{\Character{18}}%
\TextAt{9}{269.996}{88.327}{\Character{43}}%
\TextAt{13}{281.757}{88.327}{\Character{18}}%
\TextAt{17}{287.537}{92.665}{2}%
\TextAt{9}{292.269}{88.327}{\Character{41}}%
\TextAt{16}{296.822}{88.327}{\Character{103}}%
\TextAt{8}{306.697}{88.327}{is}%
\TextAt{8}{318.461}{88.327}{an}%
\TextAt{8}{334.72}{88.327}{integral}%
\TextAt{8}{377.328}{88.327}{basis}%
\TextAt{8}{406.068}{88.327}{of}%
\TextAt{18}{419.398}{88.327}{Q}%
\TextAt{9}{428.697}{88.327}{\Character{40}}%
\TextAt{13}{433.249}{88.327}{\Character{18}}%
\TextAt{9}{439.029}{88.327}{\Character{41}}%
\TextAt{8}{443.582}{88.327}{,}%
\TextAt{13}{450.731}{88.327}{\Character{18}}%
\TextAt{17}{456.511}{92.665}{3}%
\TextAt{16}{463.9}{88.327}{\Character{0}}%
\TextAt{13}{475.855}{88.327}{\Character{18}}%
\TextAt{16}{484.292}{88.327}{\Character{0}}%
\TextAt{9}{496.248}{88.327}{4}%
\TextAt{9}{505.406}{88.327}{\Character{61}}%
\TextAt{9}{517.827}{88.327}{0}%
\TextAt{8}{523.7}{88.327}{.}%
\TextAt{8}{300.14}{43.033}{14}%
\Destination{78.23}{809.815}{nav.page.14}%
\Destination{79.226}{391.955}{nav.ex.1.2.1}%
\Destination{79.226}{356.455}{nav.ex.1.2.2}%
\Destination{79.226}{307.455}{nav.ex.1.2.3}%
\Destination{79.226}{258.355}{nav.ex.1.2.4}%
\Destination{79.226}{134.355}{nav.ex.1.2.5}%
\Destination{79.226}{102.355}{nav.ex.1.2.6}%
\Destination{79.36}{776.155}{nav.print.15}%
\Destination{79.36}{541.355}{nav.number.1.2.12}%
\Destination{79.226}{460.387}{errata.71}%
}

% Page 27 (body page 15)
\Page{%
\TextAt{2}{79.226}{796.943}{I.}%
\TextAt{2}{89.83}{796.943}{Algebraic}%
\TextAt{2}{138.648}{796.943}{Integers}%
\Vector{q 0 G 0 g 0.299 w 7.226 18.729 m 460.769 18.729 l S Q}%
\TextAt{12}{79.226}{761.124}{Exercise}%
\TextAt{12}{131.825}{761.124}{I.2.7.}%
\TextAt{8}{164.897}{761.124}{The}%
\TextAt{8}{188.428}{761.124}{discriminant}%
\TextAt{13}{255.284}{761.124}{d}%
\TextAt{21}{261.367}{759.33}{K}%
\TextAt{8}{272.944}{761.124}{of}%
\TextAt{8}{285.751}{761.124}{an}%
\TextAt{8}{301.488}{761.124}{algebraic}%
\TextAt{8}{350.4}{761.124}{number}%
\TextAt{8}{392.473}{761.124}{field}%
\TextAt{13}{417.307}{761.124}{K}%
\TextAt{8}{431.5}{761.124}{is}%
\TextAt{8}{442.741}{761.124}{always}%
\NativeMathAt{479.685}{761.124}{11.955}{\equiv}%
\TextAt{9}{492.305}{761.124}{0}%
\TextAt{8}{501.479}{761.124}{mod}%
\TextAt{9}{526.916}{761.124}{4}%
\TextAt{8}{79.226}{744.078}{or}%
\NativeMathAt{93.536}{744.078}{11.955}{\equiv}%
\TextAt{9}{106.156}{744.078}{1}%
\TextAt{8}{115.33}{744.078}{mod}%
\TextAt{9}{140.768}{744.078}{4}%
\TextAt{8}{150.518}{744.078}{(}%
\TextAt{12}{155.073}{744.078}{Stickelberger’s}%
\TextAt{12}{245.107}{744.078}{discriminant}%
\TextAt{12}{322.636}{744.078}{relation}%
\TextAt{8}{367.731}{744.078}{).}%
\TextAt{12}{97.159}{724.043}{Hint:}%
\TextAt{8}{135.929}{724.043}{The}%
\TextAt{8}{161.598}{724.043}{determinant}%
\TextAt{8}{229.23}{724.043}{det}%
\TextAt{9}{245.489}{724.043}{\Character{40}}%
\TextAt{13}{250.042}{724.043}{\Character{27}}%
\TextAt{21}{256.695}{722.25}{i}%
\TextAt{13}{260.076}{724.043}{\Character{33}}%
\TextAt{21}{267.366}{722.25}{j}%
\TextAt{9}{271.748}{724.043}{\Character{41}}%
\TextAt{8}{281.813}{724.043}{of}%
\TextAt{8}{296.759}{724.043}{an}%
\TextAt{8}{314.633}{724.043}{integral}%
\TextAt{8}{358.857}{724.043}{basis}%
\TextAt{13}{389.213}{724.043}{\Character{33}}%
\TextAt{21}{396.502}{722.25}{j}%
\TextAt{8}{406.398}{724.043}{is}%
\TextAt{8}{419.777}{724.043}{a}%
\TextAt{8}{431.148}{724.043}{sum}%
\TextAt{8}{457.535}{724.043}{of}%
\TextAt{8}{472.48}{724.043}{terms,}%
\TextAt{8}{510.33}{724.043}{each}%
\TextAt{8}{79.226}{706.997}{prefixed}%
\TextAt{8}{123.473}{706.997}{by}%
\TextAt{8}{139.433}{706.997}{a}%
\TextAt{8}{148.891}{706.997}{positive}%
\TextAt{8}{191.906}{706.997}{or}%
\TextAt{8}{205.918}{706.997}{a}%
\TextAt{8}{215.375}{706.997}{negative}%
\TextAt{8}{261.259}{706.997}{sign.}%
\TextAt{8}{289.836}{706.997}{Writing}%
\TextAt{13}{332.457}{706.997}{P}%
\TextAt{8}{341.628}{706.997}{,}%
\TextAt{8}{348.539}{706.997}{resp.}%
\TextAt{13}{377.761}{706.997}{N}%
\TextAt{8}{388.384}{706.997}{,}%
\TextAt{8}{395.294}{706.997}{for}%
\TextAt{8}{412.881}{706.997}{the}%
\TextAt{8}{432.739}{706.997}{sum}%
\TextAt{8}{457.212}{706.997}{of}%
\TextAt{8}{470.243}{706.997}{the}%
\TextAt{8}{490.101}{706.997}{positive,}%
\TextAt{8}{79.226}{689.951}{resp.}%
\TextAt{8}{108.548}{689.951}{negative}%
\TextAt{8}{154.731}{689.951}{terms,}%
\TextAt{8}{190.561}{689.951}{one}%
\TextAt{8}{212.02}{689.951}{finds}%
\TextAt{13}{240.043}{689.951}{d}%
\TextAt{21}{246.126}{688.158}{K}%
\TextAt{9}{257.649}{689.951}{\Character{61}}%
\TextAt{9}{270.059}{689.951}{\Character{40}}%
\TextAt{13}{274.627}{689.951}{P}%
\TextAt{16}{286.455}{689.951}{\Character{0}}%
\TextAt{13}{298.41}{689.951}{N}%
\TextAt{9}{309.033}{689.951}{\Character{41}}%
\TextAt{17}{313.585}{694.289}{2}%
\TextAt{9}{321.638}{689.951}{\Character{61}}%
\TextAt{9}{334.048}{689.951}{\Character{40}}%
\TextAt{13}{338.616}{689.951}{P}%
\TextAt{9}{350.444}{689.951}{\Character{43}}%
\TextAt{13}{362.205}{689.951}{N}%
\TextAt{9}{372.828}{689.951}{\Character{41}}%
\TextAt{17}{377.38}{694.289}{2}%
\TextAt{16}{384.769}{689.951}{\Character{0}}%
\TextAt{9}{396.724}{689.951}{4}%
\TextAt{13}{402.577}{689.951}{P}%
\TextAt{13}{411.735}{689.951}{N}%
\TextAt{8}{422.371}{689.951}{.}%
\TextAt{15}{79.226}{637.087}{§}%
\TextAt{15}{93.785}{637.087}{3}%
\TextAt{15}{120.693}{637.087}{Ideals}%
\TextAt{8}{79.226}{605.201}{Being}%
\TextAt{8}{112.305}{605.201}{a}%
\TextAt{8}{122.143}{605.201}{generalization}%
\TextAt{8}{197.03}{605.201}{of}%
\TextAt{8}{210.443}{605.201}{the}%
\TextAt{8}{230.682}{605.201}{ring}%
\TextAt{18}{254.83}{605.201}{Z}%
\NativeMathAt{266.263}{605.202}{11.955}{\subseteq}%
\TextAt{18}{279.023}{605.201}{Q}%
\TextAt{8}{288.322}{605.201}{,}%
\TextAt{8}{295.574}{605.201}{the}%
\TextAt{8}{315.814}{605.201}{ring}%
\TextAt{14}{339.962}{605.201}{o}%
\TextAt{21}{345.806}{603.408}{K}%
\TextAt{8}{357.989}{605.201}{of}%
\TextAt{8}{371.402}{605.201}{integers}%
\TextAt{8}{414.798}{605.201}{of}%
\TextAt{8}{428.211}{605.201}{an}%
\TextAt{8}{444.553}{605.201}{algebraic}%
\TextAt{8}{494.071}{605.201}{number}%
\TextAt{8}{79.226}{588.156}{field}%
\TextAt{13}{105.039}{588.156}{K}%
\TextAt{8}{120.211}{588.156}{is}%
\TextAt{8}{132.431}{588.156}{at}%
\TextAt{8}{147.198}{588.156}{the}%
\TextAt{8}{167.811}{588.156}{center}%
\TextAt{8}{203.057}{588.156}{of}%
\TextAt{8}{216.843}{588.156}{all}%
\TextAt{8}{233.558}{588.156}{our}%
\TextAt{8}{254.829}{588.156}{considerations.}%
\TextAt{8}{336.975}{588.156}{As}%
\TextAt{8}{354.718}{588.156}{in}%
\TextAt{18}{368.828}{588.156}{Z}%
\TextAt{8}{376.798}{588.156}{,}%
\TextAt{8}{384.518}{588.156}{every}%
\TextAt{8}{415.866}{588.156}{non-unit}%
\TextAt{13}{463.796}{588.156}{\Character{11}}%
\NativeMathAt{475.417}{588.156}{11.955}{\neq}%
\TextAt{9}{488.604}{588.156}{0}%
\TextAt{8}{498.827}{588.156}{can}%
\TextAt{8}{520.743}{588.156}{be}%
\TextAt{8}{79.226}{571.11}{factored}%
\TextAt{8}{125.28}{571.11}{in}%
\TextAt{14}{139.784}{571.11}{o}%
\TextAt{21}{145.628}{569.316}{K}%
\TextAt{8}{158.579}{571.11}{into}%
\TextAt{8}{183.173}{571.11}{a}%
\TextAt{8}{193.78}{571.11}{product}%
\TextAt{8}{238.53}{571.11}{of}%
\TextAt{8}{252.711}{571.11}{irreducible}%
\TextAt{8}{311.437}{571.11}{elements.}%
\TextAt{8}{366.398}{571.11}{For}%
\TextAt{8}{388.218}{571.11}{if}%
\TextAt{13}{399.793}{571.11}{\Character{11}}%
\TextAt{8}{412.063}{571.11}{is}%
\TextAt{8}{424.678}{571.11}{not}%
\TextAt{8}{446.343}{571.11}{itself}%
\TextAt{8}{475.54}{571.11}{irreducible,}%
\TextAt{8}{79.226}{554.064}{then}%
\TextAt{8}{105.886}{554.064}{it}%
\TextAt{8}{117.59}{554.064}{can}%
\TextAt{8}{139.05}{554.064}{be}%
\TextAt{8}{154.974}{554.064}{written}%
\TextAt{8}{195.944}{554.064}{as}%
\TextAt{8}{210.315}{554.064}{a}%
\TextAt{8}{220.07}{554.064}{product}%
\TextAt{8}{263.969}{554.064}{of}%
\TextAt{8}{277.299}{554.064}{two}%
\TextAt{8}{299.416}{554.064}{non-units}%
\TextAt{13}{351.505}{554.064}{\Character{11}}%
\TextAt{9}{362.348}{554.064}{\Character{61}}%
\TextAt{13}{374.773}{554.064}{\Character{12}}%
\TextAt{13}{382.042}{554.064}{\Character{13}}%
\TextAt{8}{388.767}{554.064}{.}%
\TextAt{8}{397.215}{554.064}{Then}%
\TextAt{8}{427.772}{554.064}{by}%
\TextAt{8}{444.031}{554.064}{§2,}%
\TextAt{8}{462.717}{554.064}{one}%
\TextAt{8}{484.177}{554.064}{has}%
\TextAt{9}{154.182}{525.062}{1}%
\TextAt{13}{163.356}{525.062}{\Character{60}}%
\TextAt{16}{175.781}{525.062}{\Character{106}}%
\TextAt{13}{179.102}{525.062}{N}%
\TextAt{21}{188.476}{523.203}{K}%
\TextAt{3}{196.181}{523.203}{\Character{106}}%
\TextAt{25}{198.533}{523.203}{Q}%
\TextAt{9}{205.23}{525.062}{\Character{40}}%
\TextAt{13}{209.783}{525.062}{\Character{12}}%
\TextAt{9}{217.054}{525.062}{\Character{41}}%
\TextAt{16}{221.606}{525.062}{\Character{106}}%
\TextAt{13}{228.248}{525.062}{\Character{60}}%
\TextAt{16}{240.673}{525.062}{\Character{106}}%
\TextAt{13}{243.994}{525.062}{N}%
\TextAt{21}{253.368}{523.203}{K}%
\TextAt{3}{261.073}{523.203}{\Character{106}}%
\TextAt{25}{263.425}{523.203}{Q}%
\TextAt{9}{270.122}{525.062}{\Character{40}}%
\TextAt{13}{274.674}{525.062}{\Character{11}}%
\TextAt{9}{282.196}{525.062}{\Character{41}}%
\TextAt{16}{286.748}{525.062}{\Character{106}}%
\TextAt{13}{290.069}{525.062}{\Character{59}}%
\TextAt{9}{319.224}{525.062}{1}%
\TextAt{13}{328.398}{525.062}{\Character{60}}%
\TextAt{16}{340.823}{525.062}{\Character{106}}%
\TextAt{13}{344.144}{525.062}{N}%
\TextAt{21}{353.518}{523.203}{K}%
\TextAt{3}{361.222}{523.203}{\Character{106}}%
\TextAt{25}{363.575}{523.203}{Q}%
\TextAt{9}{370.272}{525.062}{\Character{40}}%
\TextAt{13}{374.824}{525.062}{\Character{13}}%
\TextAt{9}{381.546}{525.062}{\Character{41}}%
\TextAt{16}{386.099}{525.062}{\Character{106}}%
\TextAt{13}{392.74}{525.062}{\Character{60}}%
\TextAt{16}{405.166}{525.062}{\Character{106}}%
\TextAt{13}{408.487}{525.062}{N}%
\TextAt{21}{417.861}{523.203}{K}%
\TextAt{3}{425.565}{523.203}{\Character{106}}%
\TextAt{25}{427.917}{523.203}{Q}%
\TextAt{9}{434.615}{525.062}{\Character{40}}%
\TextAt{13}{439.167}{525.062}{\Character{11}}%
\TextAt{9}{446.689}{525.062}{\Character{41}}%
\TextAt{16}{451.241}{525.062}{\Character{106}}%
\TextAt{13}{454.562}{525.062}{\Character{59}}%
\TextAt{8}{79.226}{496.061}{and}%
\TextAt{8}{102.086}{496.061}{the}%
\TextAt{8}{122.339}{496.061}{prime}%
\TextAt{8}{155.599}{496.061}{decomposition}%
\TextAt{8}{232.831}{496.061}{of}%
\TextAt{13}{246.258}{496.061}{\Character{11}}%
\TextAt{8}{257.774}{496.061}{follows}%
\TextAt{8}{296.306}{496.061}{by}%
\TextAt{8}{312.662}{496.061}{induction}%
\TextAt{8}{364.788}{496.061}{from}%
\TextAt{8}{392.525}{496.061}{those}%
\TextAt{8}{423.251}{496.061}{of}%
\TextAt{13}{436.678}{496.061}{\Character{12}}%
\TextAt{8}{447.943}{496.061}{and}%
\TextAt{13}{470.803}{496.061}{\Character{13}}%
\TextAt{8}{477.525}{496.061}{.}%
\TextAt{8}{486.264}{496.061}{However,}%
\TextAt{8}{79.226}{479.015}{contrary}%
\TextAt{8}{126.402}{479.015}{to}%
\TextAt{8}{141.047}{479.015}{what}%
\TextAt{8}{170.649}{479.015}{happens}%
\TextAt{8}{216.892}{479.015}{in}%
\TextAt{8}{230.88}{479.015}{the}%
\TextAt{8}{251.372}{479.015}{rings}%
\TextAt{18}{280.388}{479.015}{Z}%
\TextAt{8}{292.591}{479.015}{and}%
\TextAt{18}{315.688}{479.015}{Z}%
\TextAt{9}{323.659}{479.015}{\Character{91}}%
\TextAt{13}{326.91}{479.015}{i}%
\TextAt{9}{330.904}{479.015}{\Character{93}}%
\TextAt{8}{334.155}{479.015}{,}%
\TextAt{8}{341.724}{479.015}{the}%
\TextAt{8}{362.215}{479.015}{uniqueness}%
\TextAt{8}{421.526}{479.015}{of}%
\TextAt{8}{435.191}{479.015}{prime}%
\TextAt{8}{468.69}{479.015}{factorization}%
\TextAt{8}{79.226}{461.969}{does}%
\TextAt{8}{105.623}{461.969}{not}%
\TextAt{8}{126.437}{461.969}{hold}%
\TextAt{8}{152.452}{461.969}{in}%
\TextAt{8}{166.104}{461.969}{general.}%
\TextAt{12}{97.159}{441.935}{Example:}%
\TextAt{8}{158.02}{441.935}{The}%
\TextAt{8}{182.522}{441.935}{ring}%
\TextAt{8}{207.037}{441.935}{of}%
\TextAt{8}{220.815}{441.935}{integers}%
\TextAt{8}{264.578}{441.935}{of}%
\TextAt{8}{278.356}{441.935}{the}%
\TextAt{8}{298.961}{441.935}{field}%
\TextAt{13}{324.767}{441.935}{K}%
\TextAt{9}{339.67}{441.935}{\Character{61}}%
\TextAt{18}{352.861}{441.935}{Q}%
\TextAt{9}{362.159}{441.935}{\Character{40}}%
\TextAt{16}{366.711}{451.326}{\Character{112}}%
\Vector{q 0 G 0 g 0.478 w 304.674 -322.435 m 319.826 -322.435 l S Q}%
\TextAt{16}{376.674}{441.935}{\Character{0}}%
\TextAt{9}{385.973}{441.935}{5\Character{41}}%
\TextAt{8}{400.724}{441.935}{is}%
\TextAt{8}{412.937}{441.935}{given}%
\TextAt{8}{443.955}{441.935}{by}%
\TextAt{8}{460.662}{441.935}{§2,}%
\TextAt{8}{479.909}{441.935}{exercise}%
\TextAt{8}{523.66}{441.935}{4,}%
\TextAt{8}{79.226}{424.889}{as}%
\TextAt{14}{93.385}{424.889}{o}%
\TextAt{21}{99.229}{423.095}{K}%
\TextAt{9}{110.752}{424.889}{\Character{61}}%
\TextAt{18}{123.178}{424.889}{Z}%
\TextAt{9}{133.372}{424.889}{\Character{43}}%
\TextAt{18}{144.701}{424.889}{Z}%
\TextAt{16}{152.671}{434.28}{\Character{112}}%
\Vector{q 0 G 0 g 0.478 w 90.634 -339.481 m 105.785 -339.481 l S Q}%
\TextAt{16}{162.634}{424.889}{\Character{0}}%
\TextAt{9}{171.932}{424.889}{5}%
\TextAt{8}{177.785}{424.889}{.}%
\TextAt{8}{186.163}{424.889}{In}%
\TextAt{8}{200.573}{424.889}{this}%
\TextAt{8}{223.183}{424.889}{ring,}%
\TextAt{8}{250.332}{424.889}{the}%
\TextAt{8}{270.277}{424.889}{rational}%
\TextAt{8}{313.654}{424.889}{integer}%
\TextAt{8}{352.141}{424.889}{21}%
\TextAt{8}{367.543}{424.889}{can}%
\TextAt{8}{388.791}{424.889}{be}%
\TextAt{8}{404.504}{424.889}{decomposed}%
\TextAt{8}{469.711}{424.889}{in}%
\TextAt{8}{483.152}{424.889}{two}%
\TextAt{8}{505.057}{424.889}{ways,}%
\TextAt{9}{212.163}{395.887}{21}%
\TextAt{9}{227.167}{395.887}{\Character{61}}%
\TextAt{9}{239.588}{395.887}{3}%
\TextAt{16}{248.125}{395.887}{\Character{1}}%
\TextAt{9}{254.102}{395.887}{7}%
\TextAt{9}{263.26}{395.887}{\Character{61}}%
\TextAt{9}{275.681}{395.887}{\Character{40}1}%
\TextAt{9}{288.724}{395.887}{\Character{43}}%
\TextAt{9}{300.488}{395.887}{2}%
\TextAt{16}{306.378}{405.863}{\Character{112}}%
\Vector{q 0 G 0 g 0.478 w 244.341 -367.898 m 259.492 -367.898 l S Q}%
\TextAt{16}{316.341}{395.887}{\Character{0}}%
\TextAt{9}{325.639}{395.887}{5\Character{41}\Character{40}1}%
\TextAt{16}{349.106}{395.887}{\Character{0}}%
\TextAt{9}{361.062}{395.887}{2}%
\TextAt{16}{366.915}{405.863}{\Character{112}}%
\Vector{q 0 G 0 g 0.478 w 304.877 -367.898 m 320.029 -367.898 l S Q}%
\TextAt{16}{376.877}{395.887}{\Character{0}}%
\TextAt{9}{386.176}{395.887}{5\Character{41}}%
\TextAt{13}{396.581}{395.887}{\Character{58}}%
\TextAt{8}{79.226}{366.886}{All}%
\TextAt{8}{97.926}{366.886}{factors}%
\TextAt{8}{135.564}{366.886}{occurring}%
\TextAt{8}{186.794}{366.886}{here}%
\TextAt{8}{211.675}{366.886}{are}%
\TextAt{8}{230.71}{366.886}{irreducible}%
\TextAt{8}{288.109}{366.886}{in}%
\TextAt{14}{301.286}{366.886}{o}%
\TextAt{21}{307.13}{365.093}{K}%
\TextAt{8}{315.333}{366.886}{.}%
\TextAt{8}{323.622}{366.886}{For}%
\TextAt{8}{344.116}{366.886}{if}%
\TextAt{8}{354.364}{366.886}{one}%
\TextAt{8}{375.347}{366.886}{had,}%
\TextAt{8}{400.981}{366.886}{for}%
\TextAt{8}{418.39}{366.886}{instance,}%
\TextAt{9}{466.846}{366.886}{3}%
\TextAt{9}{476.004}{366.886}{\Character{61}}%
\TextAt{13}{488.445}{366.886}{\Character{11}\Character{12}}%
\TextAt{8}{503.238}{366.886}{,}%
\TextAt{8}{510.007}{366.886}{with}%
\TextAt{13}{79.226}{349.84}{\Character{11}\Character{59}}%
\TextAt{13}{91.97}{349.84}{\Character{12}}%
\TextAt{8}{103.161}{349.84}{non-units,}%
\TextAt{8}{158.501}{349.84}{then}%
\TextAt{9}{185.161}{349.84}{9}%
\TextAt{9}{194.319}{349.84}{\Character{61}}%
\TextAt{13}{206.76}{349.84}{N}%
\TextAt{21}{216.134}{347.981}{K}%
\TextAt{3}{223.839}{347.981}{\Character{106}}%
\TextAt{25}{226.191}{347.981}{Q}%
\TextAt{9}{232.888}{349.84}{\Character{40}}%
\TextAt{13}{237.441}{349.84}{\Character{11}}%
\TextAt{9}{244.962}{349.84}{\Character{41}}%
\TextAt{13}{249.515}{349.84}{N}%
\TextAt{21}{258.889}{347.981}{K}%
\TextAt{3}{266.593}{347.981}{\Character{106}}%
\TextAt{25}{268.945}{347.981}{Q}%
\TextAt{9}{275.642}{349.84}{\Character{40}}%
\TextAt{13}{280.195}{349.84}{\Character{12}}%
\TextAt{9}{287.466}{349.84}{\Character{41}}%
\TextAt{8}{295.916}{349.84}{would}%
\TextAt{8}{330.06}{349.84}{imply}%
\TextAt{13}{362.901}{349.84}{N}%
\TextAt{21}{372.274}{347.981}{K}%
\TextAt{3}{379.979}{347.981}{\Character{106}}%
\TextAt{25}{382.331}{347.981}{Q}%
\TextAt{9}{389.028}{349.84}{\Character{40}}%
\TextAt{13}{393.581}{349.84}{\Character{11}}%
\TextAt{9}{401.102}{349.84}{\Character{41}}%
\TextAt{9}{408.957}{349.84}{\Character{61}}%
\NativeMathAt{421.401}{349.84}{11.955}{\pm}%
\TextAt{9}{430.7}{349.84}{3}%
\TextAt{8}{436.553}{349.84}{.}%
\TextAt{8}{445.001}{349.84}{But}%
\TextAt{8}{468.242}{349.84}{the}%
\TextAt{8}{488.398}{349.84}{equation}%
\TextAt{13}{217.637}{320.839}{N}%
\TextAt{21}{227.011}{318.98}{K}%
\TextAt{3}{234.715}{318.98}{\Character{106}}%
\TextAt{25}{237.068}{318.98}{Q}%
\TextAt{9}{243.765}{320.839}{\Character{40}}%
\TextAt{13}{248.317}{320.839}{x}%
\TextAt{9}{257.626}{320.839}{\Character{43}}%
\TextAt{13}{269.387}{320.839}{y}%
\TextAt{16}{275.524}{330.814}{\Character{112}}%
\Vector{q 0 G 0 g 0.478 w 213.486 -442.947 m 228.638 -442.947 l S Q}%
\TextAt{16}{285.486}{320.839}{\Character{0}}%
\TextAt{9}{294.785}{320.839}{5\Character{41}}%
\TextAt{9}{308.486}{320.839}{\Character{61}}%
\TextAt{13}{320.937}{320.839}{x}%
\TextAt{17}{327.589}{325.775}{2}%
\TextAt{9}{334.978}{320.839}{\Character{43}}%
\TextAt{9}{346.73}{320.839}{5}%
\TextAt{13}{352.592}{320.839}{y}%
\TextAt{17}{358.729}{325.775}{2}%
\TextAt{9}{366.782}{320.839}{\Character{61}}%
\NativeMathAt{379.207}{320.839}{11.955}{\pm}%
\TextAt{9}{388.506}{320.839}{3}%
\TextAt{8}{79.226}{291.838}{has}%
\TextAt{8}{100.56}{291.838}{no}%
\TextAt{8}{117.28}{291.838}{solutions}%
\TextAt{8}{166.649}{291.838}{in}%
\TextAt{18}{180.762}{291.838}{Z}%
\TextAt{8}{188.733}{291.838}{.}%
\TextAt{8}{198.562}{291.838}{In}%
\TextAt{8}{213.644}{291.838}{the}%
\TextAt{8}{234.261}{291.838}{same}%
\TextAt{8}{264.048}{291.838}{way}%
\TextAt{8}{288.251}{291.838}{it}%
\TextAt{8}{300.416}{291.838}{is}%
\TextAt{8}{312.64}{291.838}{seen}%
\TextAt{8}{338.517}{291.838}{that}%
\TextAt{9}{364.347}{291.838}{7}%
\TextAt{8}{370.2}{291.838}{,}%
\TextAt{9}{377.924}{291.838}{1}%
\TextAt{9}{386.735}{291.838}{\Character{43}}%
\TextAt{9}{398.798}{291.838}{2}%
\TextAt{16}{404.676}{301.23}{\Character{112}}%
\Vector{q 0 G 0 g 0.478 w 342.639 -472.531 m 357.79 -472.531 l S Q}%
\TextAt{16}{414.639}{291.838}{\Character{0}}%
\TextAt{9}{423.937}{291.838}{5}%
\TextAt{8}{429.79}{291.838}{,}%
\TextAt{8}{437.515}{291.838}{and}%
\TextAt{9}{460.738}{291.838}{1}%
\TextAt{16}{469.562}{291.838}{\Character{0}}%
\TextAt{9}{481.831}{291.838}{2}%
\TextAt{16}{487.684}{301.23}{\Character{112}}%
\Vector{q 0 G 0 g 0.478 w 425.647 -472.531 m 440.798 -472.531 l S Q}%
\TextAt{16}{497.647}{291.838}{\Character{0}}%
\TextAt{9}{506.945}{291.838}{5}%
\TextAt{8}{517.156}{291.838}{are}%
\TextAt{8}{79.226}{274.792}{irreducible.}%
\TextAt{8}{141.652}{274.792}{As}%
\TextAt{8}{158.939}{274.792}{the}%
\TextAt{8}{179.096}{274.792}{fractions}%
\TextAt{9}{238.987}{251.373}{1}%
\NativeMathAt{247.487}{251.373}{11.955}{\pm}%
\TextAt{9}{259.239}{251.373}{2}%
\TextAt{16}{265.111}{260.765}{\Character{112}}%
\Vector{q 0 G 0 g 0.478 w 203.074 -512.996 m 218.225 -512.996 l S Q}%
\TextAt{16}{275.074}{251.373}{\Character{0}}%
\TextAt{9}{284.372}{251.373}{5}%
\Vector{q 0 G 0 g 0.478 w 166.987 -527.726 m 218.225 -527.726 l S Q}%
\TextAt{9}{261.68}{235.085}{3}%
\TextAt{13}{291.421}{243.286}{\Character{59}}%
\TextAt{9}{321.771}{251.373}{1}%
\NativeMathAt{330.271}{251.373}{11.955}{\pm}%
\TextAt{9}{342.023}{251.373}{2}%
\TextAt{16}{347.895}{260.765}{\Character{112}}%
\Vector{q 0 G 0 g 0.478 w 285.857 -512.996 m 301.009 -512.996 l S Q}%
\TextAt{16}{357.857}{251.373}{\Character{0}}%
\TextAt{9}{367.156}{251.373}{5}%
\Vector{q 0 G 0 g 0.478 w 249.771 -527.726 m 301.009 -527.726 l S Q}%
\TextAt{9}{344.463}{235.085}{7}%
\TextAt{8}{79.226}{211.786}{do}%
\TextAt{8}{95.461}{211.786}{not}%
\TextAt{8}{116.251}{211.786}{belong}%
\TextAt{8}{153.623}{211.786}{to}%
\TextAt{14}{167.91}{211.786}{o}%
\TextAt{21}{173.754}{209.992}{K}%
\TextAt{8}{181.957}{211.786}{,}%
\TextAt{8}{189.087}{211.786}{the}%
\TextAt{8}{209.219}{211.786}{numbers}%
\TextAt{8}{256.406}{211.786}{3}%
\TextAt{8}{266.138}{211.786}{and}%
\TextAt{8}{288.877}{211.786}{7}%
\TextAt{8}{298.608}{211.786}{are}%
\TextAt{8}{318.095}{211.786}{not}%
\TextAt{8}{338.885}{211.786}{associated}%
\TextAt{8}{394.596}{211.786}{to}%
\TextAt{9}{408.883}{211.786}{1}%
\TextAt{9}{417.335}{211.786}{\Character{43}}%
\TextAt{9}{429.04}{211.786}{2}%
\TextAt{16}{434.909}{221.177}{\Character{112}}%
\Vector{q 0 G 0 g 0.478 w 372.872 -552.584 m 388.023 -552.584 l S Q}%
\TextAt{16}{444.872}{211.786}{\Character{0}}%
\TextAt{9}{454.17}{211.786}{5}%
\TextAt{8}{463.897}{211.786}{or}%
\TextAt{9}{478.183}{211.786}{1}%
\TextAt{16}{486.644}{211.786}{\Character{0}}%
\TextAt{9}{498.551}{211.786}{2}%
\TextAt{16}{504.404}{221.177}{\Character{112}}%
\Vector{q 0 G 0 g 0.478 w 442.366 -552.584 m 457.518 -552.584 l S Q}%
\TextAt{16}{514.366}{211.786}{\Character{0}}%
\TextAt{9}{523.665}{211.786}{5}%
\TextAt{8}{529.518}{211.786}{.}%
\TextAt{8}{79.226}{194.74}{The}%
\TextAt{8}{103.28}{194.74}{two}%
\TextAt{8}{125.397}{194.74}{prime}%
\TextAt{8}{158.561}{194.74}{factorizations}%
\TextAt{8}{231.152}{194.74}{of}%
\TextAt{8}{244.482}{194.74}{21}%
\TextAt{8}{260.096}{194.74}{are}%
\TextAt{8}{279.607}{194.74}{therefore}%
\TextAt{8}{328.707}{194.74}{essentially}%
\TextAt{8}{384.764}{194.74}{different.}%
\TextAt{8}{97.159}{174.705}{Realizing}%
\TextAt{8}{146.81}{174.705}{the}%
\TextAt{8}{165.725}{174.705}{failure}%
\TextAt{8}{200.576}{174.705}{of}%
\TextAt{8}{212.664}{174.705}{unique}%
\TextAt{8}{249.464}{174.705}{factorization}%
\TextAt{8}{316.199}{174.705}{in}%
\TextAt{8}{328.61}{174.705}{general}%
\TextAt{8}{367.693}{174.705}{has}%
\TextAt{8}{387.325}{174.705}{led}%
\TextAt{8}{404.937}{174.705}{to}%
\TextAt{8}{418.005}{174.705}{one}%
\TextAt{8}{438.223}{174.705}{of}%
\TextAt{8}{450.311}{174.705}{the}%
\TextAt{8}{469.226}{174.705}{grand}%
\TextAt{8}{501.16}{174.705}{events}%
\TextAt{8}{79.226}{157.659}{in}%
\TextAt{8}{94.189}{157.659}{the}%
\TextAt{8}{115.655}{157.659}{history}%
\TextAt{8}{156.381}{157.659}{of}%
\TextAt{8}{171.02}{157.659}{number}%
\TextAt{8}{214.926}{157.659}{theory,}%
\TextAt{8}{255.585}{157.659}{the}%
\TextAt{8}{277.051}{157.659}{discovery}%
\TextAt{8}{329.158}{157.659}{of}%
\TextAt{8}{343.798}{157.659}{ideal}%
\TextAt{8}{373.071}{157.659}{theory}%
\TextAt{8}{411.131}{157.659}{by}%
\TextAt{26}{428.699}{157.659}{Eduard}%
\TextAt{26}{480.382}{157.659}{Kummer}%
\TextAt{8}{529.518}{157.659}{.}%
\TextAt{8}{79.226}{140.613}{Inspired}%
\TextAt{8}{124.086}{140.613}{by}%
\TextAt{8}{139.954}{140.613}{the}%
\TextAt{8}{159.72}{140.613}{discovery}%
\TextAt{8}{210.127}{140.613}{of}%
\TextAt{8}{223.067}{140.613}{complex}%
\TextAt{8}{268.524}{140.613}{numbers,}%
\TextAt{8}{318.674}{140.613}{Kummer’s}%
\TextAt{8}{374.915}{140.613}{idea}%
\TextAt{8}{399.236}{140.613}{was}%
\TextAt{8}{421.345}{140.613}{that}%
\TextAt{8}{446.324}{140.613}{the}%
\TextAt{8}{466.089}{140.613}{integers}%
\TextAt{8}{509.012}{140.613}{of}%
\TextAt{13}{521.952}{140.613}{K}%
\TextAt{8}{79.226}{123.567}{would}%
\TextAt{8}{112.733}{123.567}{have}%
\TextAt{8}{139.091}{123.567}{to}%
\TextAt{8}{152.764}{123.567}{admit}%
\TextAt{8}{185.948}{123.567}{an}%
\TextAt{8}{201.57}{123.567}{embedding}%
\TextAt{8}{260.111}{123.567}{into}%
\TextAt{8}{283.217}{123.567}{a}%
\TextAt{8}{292.336}{123.567}{bigger}%
\TextAt{8}{326.823}{123.567}{domain}%
\TextAt{8}{367.814}{123.567}{of}%
\TextAt{8}{380.507}{123.567}{“ideal}%
\TextAt{8}{413.356}{123.567}{numbers”}%
\TextAt{8}{465.453}{123.567}{where}%
\TextAt{8}{498.626}{123.567}{unique}%
\TextAt{8}{79.226}{106.521}{factorization}%
\TextAt{8}{147.203}{106.521}{into}%
\TextAt{8}{170.946}{106.521}{“ideal}%
\TextAt{8}{204.433}{106.521}{prime}%
\TextAt{8}{237.596}{106.521}{numbers”}%
\TextAt{8}{290.331}{106.521}{would}%
\TextAt{8}{324.474}{106.521}{hold.}%
\TextAt{8}{355.04}{106.521}{For}%
\TextAt{8}{376.009}{106.521}{instance,}%
\TextAt{8}{424.846}{106.521}{in}%
\TextAt{8}{438.499}{106.521}{the}%
\TextAt{8}{458.655}{106.521}{example}%
\TextAt{8}{504.503}{106.521}{of}%
\TextAt{9}{212.163}{77.52}{21}%
\TextAt{9}{227.167}{77.52}{\Character{61}}%
\TextAt{9}{239.588}{77.52}{3}%
\TextAt{16}{248.125}{77.52}{\Character{1}}%
\TextAt{9}{254.102}{77.52}{7}%
\TextAt{9}{263.26}{77.52}{\Character{61}}%
\TextAt{9}{275.681}{77.52}{\Character{40}1}%
\TextAt{9}{288.724}{77.52}{\Character{43}}%
\TextAt{9}{300.488}{77.52}{2}%
\TextAt{16}{306.378}{87.495}{\Character{112}}%
\Vector{q 0 G 0 g 0.478 w 244.341 -686.266 m 259.492 -686.266 l S Q}%
\TextAt{16}{316.341}{77.52}{\Character{0}}%
\TextAt{9}{325.639}{77.52}{5\Character{41}\Character{40}1}%
\TextAt{16}{349.106}{77.52}{\Character{0}}%
\TextAt{9}{361.062}{77.52}{2}%
\TextAt{16}{366.915}{87.495}{\Character{112}}%
\Vector{q 0 G 0 g 0.478 w 304.877 -686.266 m 320.029 -686.266 l S Q}%
\TextAt{16}{376.877}{77.52}{\Character{0}}%
\TextAt{9}{386.176}{77.52}{5\Character{41}}%
\TextAt{13}{396.581}{77.52}{\Character{59}}%
\TextAt{8}{300.14}{43.033}{15}%
\Destination{78.23}{809.815}{nav.page.15}%
\Destination{79.23}{662.955}{nav.section.1.3}%
\Destination{79.226}{775.155}{nav.ex.1.2.7}%
\Destination{79.36}{620.155}{nav.print.16}%
\Destination{79.23}{662.955}{outline.8}%
\Hotlink{449.21}{551.855}{6.858}{11.047}{nav.section.1.2}%
\Hotlink{523.16}{439.755}{6.858}{11.047}{nav.ex.1.2.4}%
\Hotlink{465.841}{439.755}{6.858}{11.047}{nav.section.1.2}%
\Destination{79.226}{268.373}{errata.65}%
}

% Page 28 (body page 16)
\Page{%
\TextAt{2}{79.226}{796.943}{I.}%
\TextAt{2}{89.83}{796.943}{Algebraic}%
\TextAt{2}{138.648}{796.943}{Integers}%
\Vector{q 0 G 0 g 0.299 w 7.226 18.729 m 460.769 18.729 l S Q}%
\TextAt{8}{79.226}{761.124}{the}%
\TextAt{8}{99.601}{761.124}{factors}%
\TextAt{8}{137.932}{761.124}{on}%
\TextAt{8}{154.409}{761.124}{the}%
\TextAt{8}{174.784}{761.124}{right}%
\TextAt{8}{203.3}{761.124}{would}%
\TextAt{8}{237.662}{761.124}{be}%
\TextAt{8}{253.805}{761.124}{composed}%
\TextAt{8}{307.738}{761.124}{of}%
\TextAt{8}{321.286}{761.124}{ideal}%
\TextAt{8}{349.467}{761.124}{prime}%
\TextAt{8}{382.849}{761.124}{numbers}%
\TextAt{14}{430.278}{761.124}{p}%
\TextAt{17}{436.261}{759.33}{1}%
\TextAt{13}{440.994}{761.124}{\Character{59}}%
\TextAt{14}{446.238}{761.124}{p}%
\TextAt{17}{452.221}{759.33}{2}%
\TextAt{13}{456.953}{761.124}{\Character{59}}%
\TextAt{14}{462.197}{761.124}{p}%
\TextAt{17}{468.18}{759.33}{3}%
\TextAt{13}{472.913}{761.124}{\Character{59}}%
\TextAt{14}{478.157}{761.124}{p}%
\TextAt{17}{484.14}{759.33}{4}%
\TextAt{8}{488.872}{761.124}{,}%
\TextAt{8}{496.294}{761.124}{subject}%
\TextAt{8}{79.226}{744.078}{to}%
\TextAt{8}{93.536}{744.078}{the}%
\TextAt{8}{113.693}{744.078}{rules}%
\TextAt{9}{147.031}{717.537}{3}%
\TextAt{9}{156.189}{717.537}{\Character{61}}%
\TextAt{14}{168.63}{717.537}{p}%
\TextAt{17}{174.613}{715.744}{1}%
\TextAt{14}{179.346}{717.537}{p}%
\TextAt{17}{185.329}{715.744}{2}%
\TextAt{13}{190.061}{717.537}{\Character{59}}%
\TextAt{9}{207.26}{717.537}{7}%
\TextAt{9}{216.418}{717.537}{\Character{61}}%
\TextAt{14}{228.86}{717.537}{p}%
\TextAt{17}{234.843}{715.744}{3}%
\TextAt{14}{239.575}{717.537}{p}%
\TextAt{17}{245.558}{715.744}{4}%
\TextAt{13}{250.29}{717.537}{\Character{59}}%
\TextAt{9}{267.49}{717.537}{1}%
\TextAt{9}{275.99}{717.537}{\Character{43}}%
\TextAt{9}{287.742}{717.537}{2}%
\TextAt{16}{293.614}{727.512}{\Character{112}}%
\Vector{q 0 G 0 g 0.478 w 231.576 -46.249 m 246.728 -46.249 l S Q}%
\TextAt{16}{303.576}{717.537}{\Character{0}}%
\TextAt{9}{312.875}{717.537}{5}%
\TextAt{9}{322.033}{717.537}{\Character{61}}%
\TextAt{14}{334.474}{717.537}{p}%
\TextAt{17}{340.457}{715.744}{1}%
\TextAt{14}{345.189}{717.537}{p}%
\TextAt{17}{351.173}{715.744}{3}%
\TextAt{13}{355.905}{717.537}{\Character{59}}%
\TextAt{9}{373.104}{717.537}{1}%
\TextAt{16}{381.614}{717.537}{\Character{0}}%
\TextAt{9}{393.569}{717.537}{2}%
\TextAt{16}{399.422}{727.512}{\Character{112}}%
\Vector{q 0 G 0 g 0.478 w 337.385 -46.249 m 352.536 -46.249 l S Q}%
\TextAt{16}{409.385}{717.537}{\Character{0}}%
\TextAt{9}{418.683}{717.537}{5}%
\TextAt{9}{427.841}{717.537}{\Character{61}}%
\TextAt{14}{440.282}{717.537}{p}%
\TextAt{17}{446.266}{715.744}{2}%
\TextAt{14}{450.998}{717.537}{p}%
\TextAt{17}{456.981}{715.744}{4}%
\TextAt{13}{461.713}{717.537}{\Character{58}}%
\TextAt{8}{79.226}{690.996}{This}%
\TextAt{8}{105.946}{690.996}{would}%
\TextAt{8}{140.09}{690.996}{resolve}%
\TextAt{8}{178.526}{690.996}{the}%
\TextAt{8}{198.682}{690.996}{above}%
\TextAt{8}{231.858}{690.996}{non-uniqueness}%
\TextAt{8}{313.595}{690.996}{into}%
\TextAt{8}{337.338}{690.996}{the}%
\TextAt{8}{357.495}{690.996}{wonderfully}%
\TextAt{8}{420.905}{690.996}{unique}%
\TextAt{8}{458.946}{690.996}{factorization}%
\TextAt{9}{221.702}{664.455}{21}%
\TextAt{9}{236.706}{664.455}{\Character{61}}%
\TextAt{9}{249.127}{664.455}{\Character{40}}%
\TextAt{14}{253.706}{664.455}{p}%
\TextAt{17}{259.69}{662.662}{1}%
\TextAt{14}{264.422}{664.455}{p}%
\TextAt{17}{270.405}{662.662}{2}%
\TextAt{9}{275.137}{664.455}{\Character{41}\Character{40}}%
\TextAt{14}{284.242}{664.455}{p}%
\TextAt{17}{290.225}{662.662}{3}%
\TextAt{14}{294.957}{664.455}{p}%
\TextAt{17}{300.94}{662.662}{4}%
\TextAt{9}{305.673}{664.455}{\Character{41}}%
\TextAt{9}{313.528}{664.455}{\Character{61}}%
\TextAt{9}{325.949}{664.455}{\Character{40}}%
\TextAt{14}{330.524}{664.455}{p}%
\TextAt{17}{336.507}{662.662}{1}%
\TextAt{14}{341.239}{664.455}{p}%
\TextAt{17}{347.222}{662.662}{3}%
\TextAt{9}{351.954}{664.455}{\Character{41}\Character{40}}%
\TextAt{14}{361.059}{664.455}{p}%
\TextAt{17}{367.042}{662.662}{2}%
\TextAt{14}{371.775}{664.455}{p}%
\TextAt{17}{377.758}{662.662}{4}%
\TextAt{9}{382.49}{664.455}{\Character{41}}%
\TextAt{13}{387.042}{664.455}{\Character{58}}%
\TextAt{8}{79.226}{637.915}{Kummer’s}%
\TextAt{8}{135.596}{637.915}{concept}%
\TextAt{8}{178.252}{637.915}{of}%
\TextAt{8}{191.32}{637.915}{“ideal}%
\TextAt{8}{224.545}{637.915}{numbers”}%
\TextAt{8}{277.017}{637.915}{was}%
\TextAt{8}{299.254}{637.915}{later}%
\TextAt{8}{326.309}{637.915}{replaced}%
\TextAt{8}{372.218}{637.915}{by}%
\TextAt{8}{388.215}{637.915}{that}%
\TextAt{8}{413.322}{637.915}{of}%
\TextAt{8}{426.389}{637.915}{ideals}%
\TextAt{8}{458.705}{637.915}{of}%
\TextAt{8}{471.773}{637.915}{the}%
\TextAt{8}{491.667}{637.915}{ring}%
\TextAt{14}{515.471}{637.915}{o}%
\TextAt{21}{521.315}{636.121}{K}%
\TextAt{8}{529.518}{637.915}{.}%
\TextAt{8}{79.226}{620.869}{The}%
\TextAt{8}{103.694}{620.869}{reason}%
\TextAt{8}{140.594}{620.869}{for}%
\TextAt{8}{158.893}{620.869}{this}%
\TextAt{8}{182.129}{620.869}{is}%
\TextAt{8}{194.307}{620.869}{easily}%
\TextAt{8}{226.975}{620.869}{seen:}%
\TextAt{8}{257.77}{620.869}{whatever}%
\TextAt{8}{308.264}{620.869}{an}%
\TextAt{8}{324.937}{620.869}{ideal}%
\TextAt{8}{353.313}{620.869}{number}%
\TextAt{13}{396.323}{620.869}{\Character{11}}%
\Vector{q 0 G 0 g 0.478 w 324.323 -154.805 m 331.845 -154.805 l S Q}%
\TextAt{8}{408.156}{620.869}{should}%
\TextAt{8}{445.702}{620.869}{be}%
\TextAt{8}{462.04}{620.869}{defined}%
\TextAt{8}{502.767}{620.869}{to}%
\TextAt{8}{517.491}{620.869}{be,}%
\TextAt{8}{79.226}{603.823}{it}%
\TextAt{8}{91.059}{603.823}{ought}%
\TextAt{8}{124.041}{603.823}{to}%
\TextAt{8}{138.48}{603.823}{be}%
\TextAt{8}{154.533}{603.823}{linked}%
\TextAt{8}{189.129}{603.823}{to}%
\TextAt{8}{203.568}{603.823}{certain}%
\TextAt{8}{242.718}{603.823}{numbers}%
\TextAt{13}{290.058}{603.823}{a}%
\NativeMathAt{299.743}{603.823}{11.955}{\in}%
\TextAt{14}{311.254}{603.823}{o}%
\TextAt{21}{317.098}{602.03}{K}%
\TextAt{8}{329.327}{603.823}{by}%
\TextAt{8}{345.715}{603.823}{a}%
\TextAt{8}{355.599}{603.823}{divisibility}%
\TextAt{8}{413.352}{603.823}{relation}%
\TextAt{13}{456.412}{603.823}{\Character{11}}%
\Vector{q 0 G 0 g 0.478 w 384.412 -171.851 m 391.934 -171.851 l S Q}%
\TextAt{16}{467.474}{603.823}{\Character{106}}%
\TextAt{13}{474.335}{603.823}{a}%
\TextAt{8}{484.506}{603.823}{satisfying}%
\TextAt{8}{79.226}{586.777}{the}%
\TextAt{8}{99.383}{586.777}{following}%
\TextAt{8}{148.817}{586.777}{rules,}%
\TextAt{8}{180.092}{586.777}{for}%
\TextAt{13}{197.977}{586.777}{a\Character{59}}%
\TextAt{13}{209.346}{586.777}{b\Character{59}}%
\TextAt{13}{219.556}{586.777}{\Character{21}}%
\NativeMathAt{229.737}{586.777}{11.955}{\in}%
\TextAt{14}{241.028}{586.777}{o}%
\TextAt{21}{246.872}{584.984}{K}%
\TextAt{8}{255.074}{586.777}{,}%
\TextAt{13}{171.263}{560.236}{\Character{11}}%
\Vector{q 0 G 0 g 0.478 w 99.263 -215.438 m 106.785 -215.438 l S Q}%
\TextAt{16}{182.106}{560.236}{\Character{106}}%
\TextAt{13}{188.747}{560.236}{a}%
\TextAt{8}{198.79}{560.236}{and}%
\TextAt{13}{221.552}{560.236}{\Character{11}}%
\Vector{q 0 G 0 g 0.478 w 149.552 -215.438 m 157.074 -215.438 l S Q}%
\TextAt{16}{232.395}{560.236}{\Character{106}}%
\TextAt{13}{239.037}{560.236}{b}%
\NativeArrowAt{251.232}{560.236}{11.955}{19.067}{\Longrightarrow}%
\TextAt{13}{277.518}{560.236}{\Character{11}}%
\Vector{q 0 G 0 g 0.478 w 205.518 -215.438 m 213.039 -215.438 l S Q}%
\TextAt{16}{288.36}{560.236}{\Character{106}}%
\TextAt{13}{295.002}{560.236}{a}%
\NativeMathAt{303.803}{560.236}{11.955}{\pm}%
\TextAt{13}{315.759}{560.236}{b}%
\TextAt{9}{320.736}{560.236}{\Character{59}}%
\TextAt{13}{349.89}{560.236}{\Character{11}}%
\Vector{q 0 G 0 g 0.478 w 277.89 -215.438 m 285.412 -215.438 l S Q}%
\TextAt{16}{360.733}{560.236}{\Character{106}}%
\TextAt{13}{367.374}{560.236}{a}%
\NativeArrowAt{380.738}{560.236}{11.955}{19.067}{\Longrightarrow}%
\TextAt{13}{407.023}{560.236}{\Character{11}}%
\Vector{q 0 G 0 g 0.478 w 335.023 -215.438 m 342.545 -215.438 l S Q}%
\TextAt{16}{417.866}{560.236}{\Character{106}}%
\TextAt{13}{424.507}{560.236}{\Character{21}a\Character{58}}%
\TextAt{8}{79.226}{533.695}{And}%
\TextAt{8}{104.906}{533.695}{an}%
\TextAt{8}{121.165}{533.695}{ideal}%
\TextAt{8}{149.128}{533.695}{number}%
\TextAt{13}{191.724}{533.695}{\Character{11}}%
\Vector{q 0 G 0 g 0.478 w 119.724 -241.978 m 127.246 -241.978 l S Q}%
\TextAt{8}{203.143}{533.695}{should}%
\TextAt{8}{240.276}{533.695}{be}%
\TextAt{8}{256.2}{533.695}{determined}%
\TextAt{8}{317.327}{533.695}{by}%
\TextAt{8}{333.586}{533.695}{the}%
\TextAt{8}{353.743}{533.695}{totality}%
\TextAt{8}{395.383}{533.695}{of}%
\TextAt{8}{408.713}{533.695}{its}%
\ErrataAt{425.02}{533.695}{109.98}{multiples in $\mathfrak o_K$}%
\TextAt{14}{253.332}{507.155}{a}%
\TextAt{9}{262.636}{507.155}{\Character{61}}%
\TextAt{16}{275.061}{507.155}{\Character{102}}%
\TextAt{13}{281.039}{507.155}{a}%
\NativeMathAt{290.505}{507.155}{11.955}{\in}%
\TextAt{14}{301.796}{507.155}{o}%
\TextAt{21}{307.64}{505.361}{K}%
\TextAt{16}{319.163}{507.155}{\Character{106}}%
\TextAt{13}{325.805}{507.155}{\Character{11}}%
\Vector{q 0 G 0 g 0.478 w 253.805 -268.519 m 261.327 -268.519 l S Q}%
\TextAt{16}{336.648}{507.155}{\Character{106}}%
\TextAt{13}{343.289}{507.155}{a}%
\TextAt{16}{349.434}{507.155}{\Character{103}}%
\TextAt{13}{355.412}{507.155}{\Character{58}}%
\TextAt{8}{79.226}{480.614}{But}%
\TextAt{8}{102.539}{480.614}{in}%
\TextAt{8}{116.264}{480.614}{view}%
\TextAt{8}{143.319}{480.614}{of}%
\TextAt{8}{156.722}{480.614}{the}%
\TextAt{8}{176.95}{480.614}{rules}%
\TextAt{8}{205.045}{480.614}{for}%
\TextAt{8}{223.002}{480.614}{divisibility,}%
\TextAt{8}{282.988}{480.614}{this}%
\TextAt{8}{305.883}{480.614}{set}%
\TextAt{8}{324.222}{480.614}{is}%
\TextAt{8}{336.059}{480.614}{an}%
\TextAt{8}{352.39}{480.614}{ideal}%
\TextAt{8}{380.425}{480.614}{of}%
\TextAt{14}{393.827}{480.614}{o}%
\TextAt{21}{399.672}{478.821}{K}%
\TextAt{8}{407.874}{480.614}{.}%
\TextAt{8}{416.539}{480.614}{This}%
\TextAt{8}{443.331}{480.614}{is}%
\TextAt{8}{455.167}{480.614}{the}%
\TextAt{8}{475.396}{480.614}{reason}%
\TextAt{8}{511.956}{480.614}{why}%
\TextAt{26}{79.226}{463.568}{Richard}%
\TextAt{26}{132.733}{463.568}{Dedekind}%
\TextAt{8}{192.857}{463.568}{re-introduced}%
\TextAt{8}{264.554}{463.568}{Kummer’s}%
\TextAt{8}{320.697}{463.568}{“ideal}%
\TextAt{8}{353.695}{463.568}{numbers”}%
\TextAt{8}{405.941}{463.568}{as}%
\TextAt{8}{419.823}{463.568}{being}%
\TextAt{8}{450.872}{463.568}{the}%
\TextAt{8}{470.54}{463.568}{ideals}%
\TextAt{8}{502.629}{463.568}{of}%
\TextAt{14}{515.471}{463.568}{o}%
\TextAt{21}{521.315}{461.775}{K}%
\TextAt{8}{529.518}{463.568}{.}%
\TextAt{8}{79.226}{446.522}{Once}%
\TextAt{8}{109.966}{446.522}{this}%
\TextAt{8}{133.617}{446.522}{is}%
\TextAt{8}{146.209}{446.522}{done,}%
\TextAt{8}{178.459}{446.522}{the}%
\TextAt{8}{199.443}{446.522}{divisibility}%
\TextAt{8}{257.896}{446.522}{relation}%
\TextAt{14}{301.655}{446.522}{a}%
\TextAt{16}{312.37}{446.522}{\Character{106}}%
\TextAt{13}{320.423}{446.522}{a}%
\TextAt{8}{331.293}{446.522}{can}%
\TextAt{8}{353.581}{446.522}{simply}%
\TextAt{8}{391.865}{446.522}{be}%
\TextAt{8}{408.617}{446.522}{defined}%
\TextAt{8}{449.758}{446.522}{by}%
\TextAt{8}{466.845}{446.522}{the}%
\TextAt{8}{487.83}{446.522}{inclusion}%
\TextAt{13}{79.226}{429.476}{a}%
\NativeMathAt{89.214}{429.476}{11.955}{\in}%
\TextAt{14}{101.027}{429.476}{a}%
\TextAt{8}{107.01}{429.476}{,}%
\TextAt{8}{114.542}{429.476}{and}%
\TextAt{8}{137.611}{429.476}{more}%
\TextAt{8}{167.184}{429.476}{generally}%
\TextAt{8}{217.248}{429.476}{the}%
\TextAt{8}{237.71}{429.476}{divisibility}%
\TextAt{8}{295.641}{429.476}{relation}%
\TextAt{14}{338.878}{429.476}{a}%
\TextAt{16}{348.704}{429.476}{\Character{106}}%
\TextAt{14}{355.868}{429.476}{b}%
\TextAt{8}{366.206}{429.476}{between}%
\TextAt{8}{411.703}{429.476}{two}%
\TextAt{8}{434.127}{429.476}{ideals}%
\TextAt{8}{467.011}{429.476}{by}%
\TextAt{14}{483.576}{429.476}{b}%
\NativeMathAt{493.554}{429.476}{11.955}{\subseteq}%
\TextAt{14}{506.695}{429.476}{a}%
\TextAt{8}{512.678}{429.476}{.}%
\TextAt{8}{522.046}{429.476}{In}%
\TextAt{8}{79.226}{412.43}{what}%
\TextAt{8}{109.509}{412.43}{follows,}%
\TextAt{8}{152.467}{412.43}{we}%
\TextAt{8}{170.71}{412.43}{will}%
\TextAt{8}{193.832}{412.43}{study}%
\TextAt{8}{227.103}{412.43}{this}%
\TextAt{8}{250.942}{412.43}{notion}%
\TextAt{8}{288.385}{412.43}{of}%
\TextAt{8}{302.732}{412.43}{divisibility}%
\TextAt{8}{361.372}{412.43}{more}%
\TextAt{8}{391.654}{412.43}{closely.}%
\TextAt{8}{435.729}{412.43}{The}%
\TextAt{8}{460.799}{412.43}{basic}%
\TextAt{8}{491.142}{412.43}{theorem}%
\TextAt{8}{79.226}{395.384}{here}%
\TextAt{8}{104.583}{395.384}{is}%
\TextAt{8}{116.347}{395.384}{the}%
\TextAt{8}{136.503}{395.384}{following.}%
\TextAt{12}{79.226}{370.075}{Theorem}%
\TextAt{12}{135.315}{370.075}{(3.1).}%
\TextAt{8}{169.826}{370.075}{The}%
\TextAt{8}{193.075}{370.075}{ring}%
\TextAt{14}{216.336}{370.075}{o}%
\TextAt{21}{222.18}{368.281}{K}%
\TextAt{8}{233.476}{370.075}{is}%
\TextAt{8}{244.435}{370.075}{noetherian,}%
\TextAt{8}{305.253}{370.075}{integrally}%
\TextAt{8}{356.489}{370.075}{closed,}%
\TextAt{8}{393.624}{370.075}{and}%
\TextAt{8}{415.582}{370.075}{every}%
\TextAt{8}{445.67}{370.075}{prime}%
\TextAt{8}{478.029}{370.075}{ideal}%
\TextAt{14}{505.187}{370.075}{p}%
\NativeMathAt{514.491}{370.075}{11.955}{\neq}%
\TextAt{9}{526.9}{370.075}{0}%
\TextAt{8}{79.226}{353.029}{is}%
\TextAt{8}{90.99}{353.029}{a}%
\TextAt{8}{100.745}{353.029}{maximal}%
\TextAt{8}{148.554}{353.029}{ideal.}%
\ErrataAt{79.19}{319.515}{455.81}{By \hyperlink{nav.number.1.2.10}{(2.10)}, $\mathfrak o_K$ is a finitely generated $\mathbb Z$-module. Its subgroup $\mathfrak a$ is therefore also finitely}%
\ErrataAt{79.19}{302.469}{455.81}{generated over $\mathbb Z$ and hence over $\mathfrak o_K$.}%
\TextAt{8}{79.226}{285.364}{being}%
\TextAt{8}{110.535}{285.364}{the}%
\TextAt{8}{130.462}{285.364}{integral}%
\TextAt{8}{172.84}{285.364}{closure}%
\TextAt{8}{211.693}{285.364}{of}%
\TextAt{18}{224.793}{285.364}{Z}%
\TextAt{8}{236.431}{285.364}{in}%
\TextAt{13}{249.855}{285.364}{K}%
\TextAt{8}{260.673}{285.364}{.}%
\TextAt{8}{269.044}{285.364}{It}%
\TextAt{8}{281.488}{285.364}{thus}%
\TextAt{8}{307.01}{285.364}{remains}%
\TextAt{8}{350.417}{285.364}{to}%
\TextAt{8}{364.498}{285.364}{show}%
\TextAt{8}{393.272}{285.364}{that}%
\TextAt{8}{418.411}{285.364}{each}%
\TextAt{8}{444.519}{285.364}{prime}%
\TextAt{8}{477.453}{285.364}{ideal}%
\TextAt{14}{505.187}{285.364}{p}%
\NativeMathAt{514.491}{285.364}{11.955}{\neq}%
\TextAt{9}{526.9}{285.364}{0}%
\TextAt{8}{79.226}{268.318}{is}%
\TextAt{8}{91.436}{268.318}{maximal.}%
\TextAt{8}{145.132}{268.318}{Now,}%
\TextAt{14}{175.601}{268.318}{p}%
\TextAt{16}{184.544}{268.318}{\Character{92}}%
\TextAt{18}{195.474}{268.318}{Z}%
\TextAt{8}{207.788}{268.318}{is}%
\TextAt{8}{219.997}{268.318}{a}%
\TextAt{8}{230.198}{268.318}{nonzero}%
\TextAt{8}{274.22}{268.318}{prime}%
\TextAt{8}{307.829}{268.318}{ideal}%
\TextAt{9}{336.238}{268.318}{\Character{40}}%
\TextAt{13}{340.79}{268.318}{p}%
\TextAt{9}{346.665}{268.318}{\Character{41}}%
\TextAt{8}{355.561}{268.318}{in}%
\TextAt{18}{369.659}{268.318}{Z}%
\TextAt{8}{377.629}{268.318}{:}%
\TextAt{8}{386.968}{268.318}{the}%
\TextAt{8}{407.57}{268.318}{primality}%
\TextAt{8}{458.754}{268.318}{is}%
\TextAt{8}{470.963}{268.318}{clear,}%
\TextAt{8}{502.735}{268.318}{and}%
\TextAt{8}{525.943}{268.318}{if}%
\TextAt{13}{79.226}{251.272}{y}%
\NativeMathAt{88.684}{251.272}{11.955}{\in}%
\TextAt{14}{99.975}{251.272}{p}%
\TextAt{8}{105.958}{251.272}{,}%
\TextAt{13}{113.107}{251.272}{y}%
\NativeMathAt{122.564}{251.272}{11.955}{\neq}%
\TextAt{9}{134.974}{251.272}{0}%
\TextAt{8}{140.843}{251.272}{,}%
\TextAt{8}{147.992}{251.272}{and}%
\TextAt{13}{238.085}{234.226}{y}%
\TextAt{21}{244.221}{239.162}{n}%
\TextAt{9}{252.514}{234.226}{\Character{43}}%
\TextAt{13}{264.276}{234.226}{a}%
\TextAt{17}{270.421}{232.433}{1}%
\TextAt{13}{275.153}{234.226}{y}%
\TextAt{21}{281.29}{239.162}{n}%
\TextAt{3}{286.428}{239.162}{\Character{0}}%
\TextAt{17}{293.014}{239.162}{1}%
\TextAt{9}{300.403}{234.226}{\Character{43}}%
\TextAt{16}{312.165}{234.226}{\Character{1}}%
\TextAt{16}{317.461}{234.226}{\Character{1}}%
\TextAt{16}{322.769}{234.226}{\Character{1}}%
\TextAt{9}{328.769}{234.226}{\Character{43}}%
\TextAt{13}{340.53}{234.226}{a}%
\TextAt{21}{346.675}{232.433}{n}%
\TextAt{9}{355.632}{234.226}{\Character{61}}%
\TextAt{9}{368.042}{234.226}{0}%
\TextAt{8}{79.226}{211.759}{is}%
\TextAt{8}{90.918}{211.759}{an}%
\TextAt{8}{107.105}{211.759}{equation}%
\TextAt{8}{154.841}{211.759}{for}%
\TextAt{13}{172.654}{211.759}{y}%
\TextAt{8}{182.616}{211.759}{with}%
\TextAt{13}{209.204}{211.759}{a}%
\TextAt{21}{215.349}{209.966}{i}%
\NativeMathAt{222.051}{211.759}{11.955}{\in}%
\TextAt{18}{233.342}{211.759}{Z}%
\TextAt{8}{241.312}{211.759}{,}%
\TextAt{13}{248.404}{211.759}{a}%
\TextAt{21}{254.549}{209.966}{n}%
\NativeMathAt{263.506}{211.759}{11.955}{\neq}%
\TextAt{9}{275.916}{211.759}{0}%
\TextAt{8}{281.784}{211.759}{,}%
\TextAt{8}{288.876}{211.759}{then}%
\TextAt{13}{315.464}{211.759}{a}%
\TextAt{21}{321.608}{209.966}{n}%
\NativeMathAt{330.566}{211.759}{11.955}{\in}%
\TextAt{14}{341.857}{211.759}{p}%
\TextAt{16}{350.349}{211.759}{\Character{92}}%
\TextAt{18}{360.828}{211.759}{Z}%
\TextAt{8}{368.798}{211.759}{.}%
\TextAt{8}{377.223}{211.759}{The}%
\TextAt{8}{401.204}{211.759}{integral}%
\TextAt{8}{443.74}{211.759}{domain}%
\TextAt{9}{485.292}{212.294}{\Character{22}}%
\TextAt{14}{485.296}{211.759}{o}%
\TextAt{9}{494.461}{211.759}{\Character{61}}%
\TextAt{14}{506.887}{211.759}{o}%
\TextAt{21}{512.731}{209.966}{K}%
\TextAt{9}{520.934}{211.759}{\Character{47}}%
\TextAt{14}{526.786}{211.759}{p}%
\TextAt{8}{79.226}{194.713}{arises}%
\TextAt{8}{110.711}{194.713}{from}%
\TextAt{13}{137.845}{194.713}{\Character{20}}%
\TextAt{9}{147.904}{194.713}{\Character{61}}%
\TextAt{18}{160.33}{194.713}{Z}%
\TextAt{9}{168.3}{194.713}{\Character{47}}%
\TextAt{13}{174.153}{194.713}{p}%
\TextAt{18}{180.028}{194.713}{Z}%
\TextAt{8}{191.389}{194.713}{by}%
\TextAt{8}{207.141}{194.713}{adjoining}%
\TextAt{8}{257.695}{194.713}{algebraic}%
\TextAt{8}{306.622}{194.713}{elements}%
\TextAt{8}{353.972}{194.713}{and}%
\TextAt{8}{376.228}{194.713}{is}%
\TextAt{8}{387.485}{194.713}{therefore}%
\TextAt{8}{436.078}{194.713}{again}%
\TextAt{8}{466.798}{194.713}{a}%
\TextAt{8}{476.047}{194.713}{field}%
\TextAt{8}{500.897}{194.713}{(recall}%
\TextAt{8}{79.226}{177.667}{the}%
\TextAt{8}{99.383}{177.667}{fact}%
\TextAt{8}{122.468}{177.667}{that}%
\TextAt{13}{147.837}{177.667}{\Character{20}}%
\TextAt{9}{154.575}{177.667}{\Character{91}}%
\TextAt{13}{157.827}{177.667}{\Character{11}}%
\TextAt{9}{165.349}{177.667}{\Character{93}}%
\TextAt{9}{171.9}{177.667}{\Character{61}}%
\TextAt{13}{184.347}{177.667}{\Character{20}}%
\TextAt{9}{191.085}{177.667}{\Character{40}}%
\TextAt{13}{195.638}{177.667}{\Character{11}}%
\TextAt{9}{203.159}{177.667}{\Character{41}}%
\TextAt{8}{207.712}{177.667}{,}%
\TextAt{8}{214.861}{177.667}{if}%
\TextAt{13}{225.585}{177.667}{\Character{11}}%
\TextAt{8}{237.004}{177.667}{is}%
\TextAt{8}{248.768}{177.667}{algebraic).}%
\TextAt{8}{307.308}{177.667}{Therefore}%
\TextAt{14}{360.305}{177.667}{p}%
\TextAt{8}{370.186}{177.667}{is}%
\TextAt{8}{381.95}{177.667}{a}%
\TextAt{8}{391.705}{177.667}{maximal}%
\TextAt{8}{439.514}{177.667}{ideal.}%
\Vector{q 0 G 0 g 0.399 w 452.334 -596.333 m 452.334 -588.263 l S Q}%
\Vector{q 0 G 0 g 0.399 w 452.534 -588.463 m 459.707 -588.463 l S Q}%
\Vector{q 0 G 0 g 0.399 w 452.534 -596.134 m 459.707 -596.134 l S Q}%
\Vector{q 0 G 0 g 0.399 w 459.906 -596.333 m 459.906 -588.263 l S Q}%
\TextAt{8}{97.159}{149.369}{The}%
\TextAt{8}{121.54}{149.369}{three}%
\TextAt{8}{151.778}{149.369}{properties}%
\TextAt{8}{207.123}{149.369}{of}%
\TextAt{8}{220.779}{149.369}{the}%
\TextAt{8}{241.263}{149.369}{ring}%
\TextAt{14}{265.655}{149.369}{o}%
\TextAt{21}{271.499}{147.575}{K}%
\TextAt{8}{283.926}{149.369}{which}%
\TextAt{8}{317.74}{149.369}{we}%
\TextAt{8}{335.294}{149.369}{have}%
\TextAt{8}{362.615}{149.369}{just}%
\TextAt{8}{386.087}{149.369}{proven}%
\TextAt{8}{424.467}{149.369}{lay}%
\TextAt{8}{443.66}{149.369}{the}%
\TextAt{8}{464.143}{149.369}{foundation}%
\TextAt{8}{523.337}{149.369}{of}%
\TextAt{8}{79.226}{132.323}{the}%
\TextAt{8}{100.561}{132.323}{whole}%
\TextAt{8}{134.903}{132.323}{theory}%
\TextAt{8}{172.831}{132.323}{of}%
\TextAt{8}{187.339}{132.323}{divisibility}%
\TextAt{8}{246.142}{132.323}{of}%
\TextAt{8}{260.65}{132.323}{its}%
\TextAt{8}{278.147}{132.323}{ideals.}%
\TextAt{8}{318.81}{132.323}{This}%
\TextAt{8}{346.709}{132.323}{theory}%
\TextAt{8}{384.637}{132.323}{was}%
\TextAt{8}{408.315}{132.323}{developed}%
\TextAt{8}{463.793}{132.323}{by}%
\TextAt{8}{481.231}{132.323}{Dedekind,}%
\TextAt{8}{79.226}{115.277}{which}%
\TextAt{8}{112.713}{115.277}{suggested}%
\TextAt{8}{165.519}{115.277}{the}%
\TextAt{8}{185.675}{115.277}{following}%
\TextAt{12}{79.226}{89.967}{Definition}%
\TextAt{12}{141.257}{89.967}{(3.2).}%
\TextAt{8}{175.915}{89.967}{A}%
\TextAt{8}{187.928}{89.967}{noetherian,}%
\TextAt{8}{248.863}{89.967}{integrally}%
\TextAt{8}{300.246}{89.967}{closed}%
\TextAt{8}{334.114}{89.967}{integral}%
\TextAt{8}{376.064}{89.967}{domain}%
\TextAt{8}{417.034}{89.967}{in}%
\TextAt{8}{430.028}{89.967}{which}%
\TextAt{8}{462.856}{89.967}{every}%
\TextAt{8}{493.09}{89.967}{nonzero}%
\TextAt{8}{79.226}{72.921}{prime}%
\TextAt{8}{112.39}{72.921}{ideal}%
\TextAt{8}{140.353}{72.921}{is}%
\TextAt{8}{152.117}{72.921}{maximal}%
\TextAt{8}{199.926}{72.921}{is}%
\TextAt{8}{211.689}{72.921}{called}%
\TextAt{8}{244.853}{72.921}{a}%
\TextAt{12}{254.608}{72.921}{Dedekind}%
\TextAt{12}{314.922}{72.921}{domain}%
\TextAt{8}{358.092}{72.921}{.}%
\TextAt{8}{300.14}{43.033}{16}%
\Destination{78.23}{809.815}{nav.page.16}%
\Destination{79.36}{776.155}{nav.print.17}%
\Destination{79.36}{164.355}{nav.print.18}%
\Destination{79.36}{384.055}{nav.number.1.3.1}%
\Destination{79.36}{103.955}{nav.number.1.3.2}%
\Hotlink{383.738}{300.255}{6.858}{11.047}{nav.section.1.2}%
\Destination{79.226}{550.695}{errata.66}%
\Destination{79.226}{332}{errata.191}%
}

% Page 29 (body page 17)
\Page{%
\TextAt{2}{79.226}{796.943}{I.}%
\TextAt{2}{89.83}{796.943}{Algebraic}%
\TextAt{2}{138.648}{796.943}{Integers}%
\Vector{q 0 G 0 g 0.299 w 7.226 18.729 m 460.769 18.729 l S Q}%
\TextAt{8}{97.159}{761.124}{Just}%
\TextAt{8}{123.054}{761.124}{as}%
\TextAt{8}{137.736}{761.124}{the}%
\TextAt{8}{158.204}{761.124}{rings}%
\TextAt{8}{187.196}{761.124}{of}%
\TextAt{8}{200.837}{761.124}{the}%
\TextAt{8}{221.305}{761.124}{form}%
\TextAt{14}{249.257}{761.124}{o}%
\TextAt{21}{255.101}{759.33}{K}%
\TextAt{8}{267.513}{761.124}{may}%
\TextAt{8}{293.193}{761.124}{be}%
\TextAt{8}{309.429}{761.124}{viewed}%
\TextAt{8}{348.104}{761.124}{as}%
\TextAt{8}{362.786}{761.124}{generalizations}%
\TextAt{8}{442.515}{761.124}{of}%
\TextAt{8}{456.157}{761.124}{the}%
\TextAt{8}{476.625}{761.124}{ring}%
\TextAt{18}{501.002}{761.124}{Z}%
\TextAt{8}{508.972}{761.124}{,}%
\TextAt{8}{516.51}{761.124}{the}%
\TextAt{8}{79.226}{744.078}{Dedekind}%
\TextAt{8}{131.445}{744.078}{domains}%
\TextAt{8}{177.723}{744.078}{may}%
\TextAt{8}{203.126}{744.078}{be}%
\TextAt{8}{219.085}{744.078}{viewed}%
\TextAt{8}{257.484}{744.078}{as}%
\TextAt{8}{271.889}{744.078}{generalized}%
\TextAt{8}{332.405}{744.078}{principal}%
\TextAt{8}{381.217}{744.078}{ideal}%
\TextAt{8}{409.215}{744.078}{domains.}%
\TextAt{8}{460.112}{744.078}{Indeed,}%
\TextAt{8}{501.437}{744.078}{if}%
\TextAt{13}{512.196}{744.078}{A}%
\TextAt{8}{524.903}{744.078}{is}%
\TextAt{8}{79.226}{727.032}{a}%
\TextAt{8}{89.567}{727.032}{principal}%
\TextAt{8}{138.93}{727.032}{ideal}%
\TextAt{8}{167.478}{727.032}{domain}%
\TextAt{8}{209.692}{727.032}{with}%
\TextAt{8}{236.937}{727.032}{field}%
\TextAt{8}{262.88}{727.032}{of}%
\TextAt{8}{276.795}{727.032}{fractions}%
\TextAt{13}{325.249}{727.032}{K}%
\TextAt{8}{336.067}{727.032}{,}%
\TextAt{8}{343.948}{727.032}{and}%
\TextAt{13}{367.296}{727.032}{L}%
\TextAt{16}{379.579}{727.032}{\Character{106}}%
\TextAt{13}{387.219}{727.032}{K}%
\TextAt{8}{402.519}{727.032}{is}%
\TextAt{8}{414.869}{727.032}{a}%
\TextAt{8}{425.209}{727.032}{finite}%
\TextAt{8}{455.707}{727.032}{field}%
\TextAt{8}{481.649}{727.032}{extension,}%
\TextAt{8}{79.226}{709.986}{then}%
\TextAt{8}{106.908}{709.986}{the}%
\TextAt{8}{128.085}{709.986}{integral}%
\TextAt{8}{171.715}{709.986}{closure}%
\TextAt{13}{211.818}{709.986}{B}%
\TextAt{8}{226.233}{709.986}{of}%
\TextAt{13}{240.584}{709.986}{A}%
\TextAt{8}{254.278}{709.986}{in}%
\TextAt{13}{268.953}{709.986}{L}%
\TextAt{8}{281.836}{709.986}{is,}%
\TextAt{8}{298.128}{709.986}{in}%
\TextAt{8}{312.802}{709.986}{general,}%
\TextAt{8}{357.656}{709.986}{not}%
\TextAt{8}{379.491}{709.986}{a}%
\TextAt{8}{390.268}{709.986}{principal}%
\TextAt{8}{440.066}{709.986}{ideal}%
\TextAt{8}{469.051}{709.986}{domain,}%
\TextAt{8}{515.207}{709.986}{but}%
\TextAt{8}{79.226}{692.94}{always}%
\TextAt{8}{116.694}{692.94}{a}%
\TextAt{8}{126.449}{692.94}{Dedekind}%
\TextAt{8}{178.633}{692.94}{domain,}%
\TextAt{8}{223.513}{692.94}{as}%
\TextAt{8}{237.883}{692.94}{we}%
\TextAt{8}{255.111}{692.94}{shall}%
\TextAt{8}{282.488}{692.94}{show}%
\TextAt{8}{311.491}{692.94}{further}%
\TextAt{8}{350.836}{692.94}{on.}%
\TextAt{8}{97.159}{672.905}{Instead}%
\TextAt{8}{138.546}{672.905}{of}%
\TextAt{8}{151.911}{672.905}{the}%
\TextAt{8}{172.102}{672.905}{ring}%
\TextAt{14}{196.203}{672.905}{o}%
\TextAt{21}{202.047}{671.112}{K}%
\TextAt{8}{214.182}{672.905}{we}%
\TextAt{8}{231.444}{672.905}{will}%
\TextAt{8}{253.583}{672.905}{now}%
\TextAt{8}{278.007}{672.905}{consider}%
\TextAt{8}{323.626}{672.905}{an}%
\TextAt{8}{339.92}{672.905}{arbitrary}%
\TextAt{8}{389.724}{672.905}{Dedekind}%
\TextAt{8}{441.943}{672.905}{domain}%
\TextAt{14}{483.605}{672.905}{o}%
\TextAt{8}{489.45}{672.905}{,}%
\TextAt{8}{496.642}{672.905}{and}%
\TextAt{8}{519.439}{672.905}{we}%
\TextAt{8}{79.226}{655.859}{denote}%
\TextAt{8}{117.049}{655.859}{by}%
\TextAt{13}{133.412}{655.859}{K}%
\TextAt{8}{148.232}{655.859}{the}%
\TextAt{8}{168.492}{655.859}{field}%
\TextAt{8}{193.954}{655.859}{of}%
\TextAt{8}{207.388}{655.859}{fractions}%
\TextAt{8}{255.361}{655.859}{of}%
\TextAt{14}{268.795}{655.859}{o}%
\TextAt{8}{274.639}{655.859}{.}%
\TextAt{8}{283.4}{655.859}{Given}%
\TextAt{8}{317.398}{655.859}{two}%
\TextAt{8}{339.619}{655.859}{ideals}%
\TextAt{14}{372.301}{655.859}{a}%
\TextAt{8}{382.286}{655.859}{and}%
\TextAt{14}{405.153}{655.859}{b}%
\TextAt{8}{415.289}{655.859}{of}%
\TextAt{14}{428.723}{655.859}{o}%
\TextAt{8}{438.569}{655.859}{(or}%
\TextAt{8}{457.539}{655.859}{more}%
\TextAt{8}{486.909}{655.859}{generally}%
\TextAt{8}{79.226}{638.813}{of}%
\TextAt{8}{93.087}{638.813}{an}%
\TextAt{8}{109.876}{638.813}{arbitrary}%
\TextAt{8}{160.176}{638.813}{ring),}%
\TextAt{8}{192.711}{638.813}{the}%
\TextAt{8}{213.398}{638.813}{divisibility}%
\TextAt{8}{271.553}{638.813}{relation}%
\TextAt{14}{315.014}{638.813}{a}%
\TextAt{16}{325.222}{638.813}{\Character{106}}%
\TextAt{14}{332.767}{638.813}{b}%
\TextAt{8}{343.33}{638.813}{is}%
\TextAt{8}{355.624}{638.813}{defined}%
\TextAt{8}{396.468}{638.813}{by}%
\TextAt{14}{413.257}{638.813}{b}%
\NativeMathAt{423.617}{638.813}{11.955}{\subseteq}%
\TextAt{14}{437.14}{638.813}{a}%
\TextAt{8}{443.123}{638.813}{,}%
\TextAt{8}{450.935}{638.813}{and}%
\TextAt{8}{474.228}{638.813}{the}%
\TextAt{12}{494.915}{638.813}{sum}%
\TextAt{8}{523.337}{638.813}{of}%
\TextAt{8}{79.226}{621.767}{the}%
\TextAt{8}{99.383}{621.767}{ideals}%
\TextAt{8}{131.96}{621.767}{by}%
\TextAt{14}{230.648}{604.721}{a}%
\TextAt{9}{239.288}{604.721}{\Character{43}}%
\TextAt{14}{251.049}{604.721}{b}%
\TextAt{9}{260.505}{604.721}{\Character{61}}%
\TextAt{16}{272.93}{604.721}{\Character{102}}%
\TextAt{13}{278.908}{604.721}{a}%
\TextAt{9}{287.71}{604.721}{\Character{43}}%
\TextAt{13}{299.471}{604.721}{b}%
\TextAt{16}{307.769}{604.721}{\Character{106}}%
\TextAt{13}{314.411}{604.721}{a}%
\NativeMathAt{323.876}{604.721}{11.955}{\in}%
\TextAt{14}{335.167}{604.721}{a}%
\TextAt{13}{341.15}{604.721}{\Character{59}}%
\TextAt{13}{346.374}{604.721}{b}%
\NativeMathAt{354.693}{604.721}{11.955}{\in}%
\TextAt{14}{365.984}{604.721}{b}%
\TextAt{16}{372.118}{604.721}{\Character{103}}%
\TextAt{13}{378.096}{604.721}{\Character{58}}%
\TextAt{8}{79.226}{581.199}{This}%
\TextAt{8}{106.97}{581.199}{is}%
\TextAt{8}{119.757}{581.199}{the}%
\TextAt{8}{140.937}{581.199}{smallest}%
\TextAt{8}{186.96}{581.199}{ideal}%
\TextAt{8}{215.947}{581.199}{containing}%
\TextAt{14}{273.889}{581.199}{a}%
\TextAt{8}{284.793}{581.199}{as}%
\TextAt{8}{300.187}{581.199}{well}%
\TextAt{8}{324.941}{581.199}{as}%
\TextAt{14}{340.335}{581.199}{b}%
\TextAt{8}{346.47}{581.199}{,}%
\TextAt{8}{354.899}{581.199}{in}%
\TextAt{8}{369.575}{581.199}{other}%
\TextAt{8}{401.168}{581.199}{words,}%
\TextAt{8}{439.258}{581.199}{it}%
\TextAt{8}{451.985}{581.199}{is}%
\TextAt{8}{464.773}{581.199}{the}%
\TextAt{12}{485.953}{581.199}{greatest}%
\TextAt{12}{79.226}{564.153}{common}%
\TextAt{12}{132.946}{564.153}{divisor}%
\TextAt{8}{176.349}{564.153}{gcd}%
\TextAt{9}{193.911}{564.153}{\Character{40}}%
\TextAt{14}{198.464}{564.153}{a}%
\TextAt{13}{204.447}{564.153}{\Character{59}}%
\TextAt{14}{209.691}{564.153}{b}%
\TextAt{9}{215.826}{564.153}{\Character{41}}%
\TextAt{8}{224.186}{564.153}{of}%
\TextAt{14}{237.427}{564.153}{a}%
\TextAt{8}{247.218}{564.153}{and}%
\TextAt{14}{269.891}{564.153}{b}%
\TextAt{8}{276.026}{564.153}{.}%
\TextAt{8}{284.444}{564.153}{By}%
\TextAt{8}{302.718}{564.153}{the}%
\TextAt{8}{322.785}{564.153}{same}%
\TextAt{8}{352.022}{564.153}{token}%
\TextAt{8}{383.805}{564.153}{the}%
\TextAt{12}{403.872}{564.153}{intersection}%
\TextAt{14}{475.992}{564.153}{a}%
\TextAt{16}{484.449}{564.153}{\Character{92}}%
\TextAt{14}{494.893}{564.153}{b}%
\TextAt{8}{504.836}{564.153}{is}%
\TextAt{8}{516.51}{564.153}{the}%
\TextAt{8}{79.226}{547.108}{lcm}%
\TextAt{8}{101.331}{547.108}{(least}%
\TextAt{8}{133.264}{547.108}{common}%
\TextAt{8}{180.092}{547.108}{multiple)}%
\TextAt{8}{230.495}{547.108}{of}%
\TextAt{14}{243.825}{547.108}{a}%
\TextAt{8}{253.705}{547.108}{and}%
\TextAt{14}{276.468}{547.108}{b}%
\TextAt{8}{282.603}{547.108}{.}%
\TextAt{8}{291.051}{547.108}{We}%
\TextAt{8}{311.196}{547.108}{define}%
\TextAt{8}{345.005}{547.108}{the}%
\TextAt{12}{365.161}{547.108}{product}%
\TextAt{8}{415.289}{547.108}{of}%
\TextAt{14}{428.619}{547.108}{a}%
\TextAt{8}{438.5}{547.108}{and}%
\TextAt{14}{461.262}{547.108}{b}%
\TextAt{8}{471.295}{547.108}{by}%
\TextAt{14}{223.115}{510.234}{ab}%
\TextAt{9}{238.554}{510.234}{\Character{61}}%
\TextAt{20}{250.979}{530.677}{\Character{40}}%
\TextAt{20}{260.61}{521.591}{\Character{88}}%
\TextAt{21}{267.803}{496.397}{i}%
\TextAt{13}{279.871}{510.234}{a}%
\TextAt{21}{286.016}{508.44}{i}%
\TextAt{13}{289.397}{510.234}{b}%
\TextAt{21}{294.374}{508.44}{i}%
\TextAt{20}{301.653}{531.155}{\Character{12}}%
\TextAt{20}{301.653}{523.982}{\Character{12}}%
\TextAt{20}{301.653}{516.809}{\Character{12}}%
\TextAt{20}{301.653}{509.636}{\Character{12}}%
\TextAt{20}{301.653}{502.463}{\Character{12}}%
\TextAt{13}{309.536}{510.234}{a}%
\TextAt{21}{315.68}{508.44}{i}%
\NativeMathAt{322.383}{510.234}{11.955}{\in}%
\TextAt{14}{333.674}{510.234}{a}%
\TextAt{13}{339.657}{510.234}{\Character{59}}%
\TextAt{13}{344.881}{510.234}{b}%
\TextAt{21}{349.878}{508.44}{i}%
\NativeMathAt{356.58}{510.234}{11.955}{\in}%
\TextAt{14}{367.871}{510.234}{b}%
\TextAt{20}{374.006}{530.677}{\Character{41}}%
\TextAt{13}{385.629}{510.234}{\Character{58}}%
\TextAt{8}{79.226}{473.833}{With}%
\TextAt{8}{109.675}{473.833}{respect}%
\TextAt{8}{149.94}{473.833}{to}%
\TextAt{8}{164.464}{473.833}{this}%
\TextAt{8}{187.501}{473.833}{multiplication}%
\TextAt{8}{262.842}{473.833}{the}%
\TextAt{8}{283.213}{473.833}{ideals}%
\TextAt{8}{316.006}{473.833}{of}%
\TextAt{14}{329.55}{473.833}{o}%
\TextAt{8}{339.506}{473.833}{will}%
\TextAt{8}{361.826}{473.833}{grant}%
\TextAt{8}{392.944}{473.833}{us}%
\TextAt{8}{408.175}{473.833}{what}%
\TextAt{8}{437.655}{473.833}{the}%
\TextAt{8}{458.026}{473.833}{elements}%
\TextAt{8}{506.097}{473.833}{alone}%
\TextAt{8}{79.226}{456.787}{may}%
\TextAt{8}{104.595}{456.787}{refuse}%
\TextAt{8}{138.141}{456.787}{to}%
\TextAt{8}{152.452}{456.787}{provide:}%
\TextAt{8}{198.63}{456.787}{the}%
\TextAt{8}{218.787}{456.787}{unique}%
\TextAt{8}{256.828}{456.787}{prime}%
\TextAt{8}{289.992}{456.787}{factorization.}%
\TextAt{12}{79.226}{430.775}{Theorem}%
\TextAt{12}{136.24}{430.775}{(3.3).}%
\TextAt{8}{171.556}{430.775}{Every}%
\TextAt{8}{205.21}{430.775}{ideal}%
\TextAt{14}{233.173}{430.775}{a}%
\TextAt{8}{243.053}{430.775}{of}%
\TextAt{14}{256.383}{430.775}{o}%
\TextAt{8}{266.125}{430.775}{different}%
\TextAt{8}{312.296}{430.775}{from}%
\TextAt{9}{339.936}{430.775}{\Character{40}0\Character{41}}%
\TextAt{8}{358.791}{430.775}{and}%
\TextAt{9}{381.554}{430.775}{\Character{40}1\Character{41}}%
\TextAt{8}{400.409}{430.775}{admits}%
\TextAt{8}{438.845}{430.775}{a}%
\TextAt{8}{448.6}{430.775}{factorization}%
\TextAt{14}{275.54}{401.774}{a}%
\TextAt{9}{284.844}{401.774}{\Character{61}}%
\TextAt{14}{297.27}{401.774}{p}%
\TextAt{17}{303.253}{399.981}{1}%
\TextAt{16}{309.978}{401.774}{\Character{1}}%
\TextAt{16}{315.274}{401.774}{\Character{1}}%
\TextAt{16}{320.582}{401.774}{\Character{1}}%
\TextAt{14}{325.918}{401.774}{p}%
\TextAt{21}{331.901}{399.981}{r}%
\TextAt{8}{79.226}{372.773}{into}%
\TextAt{8}{102.969}{372.773}{nonzero}%
\TextAt{8}{146.546}{372.773}{prime}%
\TextAt{8}{179.709}{372.773}{ideals}%
\TextAt{14}{212.287}{372.773}{p}%
\TextAt{21}{218.27}{370.98}{i}%
\TextAt{8}{225.549}{372.773}{of}%
\TextAt{14}{238.879}{372.773}{o}%
\TextAt{8}{248.621}{372.773}{which}%
\TextAt{8}{282.107}{372.773}{is}%
\TextAt{8}{293.871}{372.773}{unique}%
\TextAt{8}{331.912}{372.773}{up}%
\TextAt{8}{348.817}{372.773}{to}%
\TextAt{8}{363.127}{372.773}{the}%
\TextAt{8}{383.284}{372.773}{order}%
\TextAt{8}{413.853}{372.773}{of}%
\TextAt{8}{427.183}{372.773}{the}%
\TextAt{8}{447.339}{372.773}{factors.}%
\TextAt{8}{97.159}{346.76}{This}%
\TextAt{8}{123.44}{346.76}{theorem}%
\TextAt{8}{168.526}{346.76}{is}%
\TextAt{8}{179.851}{346.76}{of}%
\TextAt{8}{192.742}{346.76}{course}%
\TextAt{8}{228.132}{346.76}{perfectly}%
\TextAt{8}{276.136}{346.76}{in}%
\TextAt{8}{289.349}{346.76}{line}%
\TextAt{8}{311.015}{346.76}{with}%
\TextAt{8}{337.236}{346.76}{the}%
\TextAt{8}{356.954}{346.76}{invention}%
\TextAt{8}{407.252}{346.76}{of}%
\TextAt{8}{420.143}{346.76}{“ideal}%
\TextAt{8}{453.191}{346.76}{numbers”.}%
\TextAt{8}{508.704}{346.76}{Still,}%
\TextAt{8}{79.226}{329.714}{the}%
\TextAt{8}{100.65}{329.714}{fact}%
\TextAt{8}{125.003}{329.714}{that}%
\TextAt{8}{151.64}{329.714}{it}%
\TextAt{8}{164.612}{329.714}{holds}%
\TextAt{8}{196.508}{329.714}{is}%
\TextAt{8}{209.54}{329.714}{remarkable}%
\TextAt{8}{271.624}{329.714}{because}%
\TextAt{8}{316.193}{329.714}{its}%
\TextAt{8}{333.779}{329.714}{proof}%
\TextAt{8}{365.616}{329.714}{is}%
\TextAt{8}{378.648}{329.714}{far}%
\TextAt{8}{397.801}{329.714}{from}%
\TextAt{8}{426.709}{329.714}{straightforward,}%
\TextAt{8}{513.904}{329.714}{and}%
\TextAt{8}{79.226}{312.668}{unveils}%
\TextAt{8}{117.904}{312.668}{a}%
\TextAt{8}{127.578}{312.668}{deeper}%
\TextAt{8}{164.881}{312.668}{principle}%
\TextAt{8}{212.919}{312.668}{governing}%
\TextAt{8}{265.859}{312.668}{the}%
\TextAt{8}{285.934}{312.668}{arithmetic}%
\TextAt{8}{342.437}{312.668}{in}%
\TextAt{14}{356.009}{312.668}{o}%
\TextAt{8}{361.853}{312.668}{.}%
\TextAt{8}{370.274}{312.668}{We}%
\TextAt{8}{390.337}{312.668}{prepare}%
\TextAt{8}{432.53}{312.668}{the}%
\TextAt{8}{452.605}{312.668}{proof}%
\TextAt{8}{483.093}{312.668}{proper}%
\TextAt{8}{520.408}{312.668}{by}%
\TextAt{8}{79.226}{295.623}{two}%
\TextAt{8}{101.343}{295.623}{lemmas.}%
\TextAt{12}{79.226}{269.61}{Lemma}%
\TextAt{12}{127.296}{269.61}{(3.4).}%
\TextAt{8}{162.963}{269.61}{For}%
\TextAt{8}{184.284}{269.61}{every}%
\TextAt{8}{215.528}{269.61}{ideal}%
\TextAt{14}{243.843}{269.61}{a}%
\NativeMathAt{253.746}{269.61}{11.955}{\neq}%
\TextAt{9}{266.753}{269.61}{0}%
\TextAt{8}{276.873}{269.61}{of}%
\TextAt{14}{290.554}{269.61}{o}%
\TextAt{8}{300.648}{269.61}{there}%
\TextAt{8}{330.911}{269.61}{exist}%
\TextAt{8}{358.963}{269.61}{nonzero}%
\TextAt{8}{402.891}{269.61}{prime}%
\TextAt{8}{436.406}{269.61}{ideals}%
\TextAt{14}{469.336}{269.61}{p}%
\TextAt{17}{475.319}{267.817}{1}%
\TextAt{13}{480.051}{269.61}{\Character{59}}%
\TextAt{14}{485.296}{269.61}{p}%
\TextAt{17}{491.279}{267.817}{2}%
\TextAt{13}{496.011}{269.61}{\Character{59}}%
\TextAt{13}{501.235}{269.61}{\Character{58}}%
\TextAt{13}{506.472}{269.61}{\Character{58}}%
\TextAt{13}{511.696}{269.61}{\Character{58}}%
\TextAt{13}{516.933}{269.61}{\Character{59}}%
\TextAt{14}{522.232}{269.61}{p}%
\TextAt{21}{528.215}{267.817}{r}%
\TextAt{8}{79.226}{252.564}{such}%
\TextAt{8}{105.623}{252.564}{that}%
\TextAt{14}{268.46}{235.518}{a}%
\NativeMathAt{277.764}{235.518}{11.955}{\supseteq}%
\TextAt{14}{290.383}{235.518}{p}%
\TextAt{17}{296.366}{233.725}{1}%
\TextAt{14}{301.098}{235.518}{p}%
\TextAt{17}{307.082}{233.725}{2}%
\TextAt{16}{313.806}{235.518}{\Character{1}}%
\TextAt{16}{319.102}{235.518}{\Character{1}}%
\TextAt{16}{324.41}{235.518}{\Character{1}}%
\TextAt{14}{329.746}{235.518}{p}%
\TextAt{21}{335.73}{233.725}{r}%
\TextAt{13}{340.284}{235.518}{\Character{58}}%
\TextAt{19}{79.226}{194.064}{Proof.}%
\TextAt{8}{116}{194.064}{Suppose}%
\TextAt{8}{161.885}{194.064}{the}%
\TextAt{8}{182.019}{194.064}{set}%
\TextAt{14}{200.264}{194.064}{M}%
\TextAt{8}{216.686}{194.064}{of}%
\TextAt{8}{229.993}{194.064}{those}%
\TextAt{8}{260.6}{194.064}{ideals}%
\TextAt{8}{293.155}{194.064}{which}%
\TextAt{8}{326.619}{194.064}{do}%
\TextAt{8}{342.855}{194.064}{not}%
\TextAt{8}{363.647}{194.064}{fulfill}%
\TextAt{8}{393.859}{194.064}{this}%
\TextAt{8}{416.658}{194.064}{condition}%
\TextAt{8}{468.019}{194.064}{is}%
\TextAt{8}{479.76}{194.064}{nonempty.}%
\TextAt{8}{79.226}{177.018}{As}%
\TextAt{14}{97.109}{177.018}{o}%
\TextAt{8}{107.446}{177.018}{is}%
\TextAt{8}{119.806}{177.018}{noetherian,}%
\TextAt{8}{182.012}{177.018}{every}%
\TextAt{8}{213.5}{177.018}{ascending}%
\TextAt{8}{267.488}{177.018}{chain}%
\TextAt{8}{298.976}{177.018}{of}%
\TextAt{8}{312.902}{177.018}{ideals}%
\TextAt{8}{346.075}{177.018}{becomes}%
\TextAt{8}{393.224}{177.018}{stationary.}%
\TextAt{8}{454.269}{177.018}{Therefore}%
\TextAt{14}{507.862}{177.018}{M}%
\TextAt{8}{524.903}{177.018}{is}%
\TextAt{8}{79.226}{159.972}{inductively}%
\TextAt{8}{139.08}{159.972}{ordered}%
\TextAt{8}{181.049}{159.972}{with}%
\TextAt{8}{207.405}{159.972}{respect}%
\TextAt{8}{247.15}{159.972}{to}%
\TextAt{8}{261.156}{159.972}{inclusion}%
\TextAt{8}{309.688}{159.972}{and}%
\TextAt{8}{332.146}{159.972}{thus}%
\TextAt{8}{357.593}{159.972}{admits}%
\TextAt{8}{395.725}{159.972}{a}%
\TextAt{8}{405.176}{159.972}{maximal}%
\TextAt{8}{452.68}{159.972}{element}%
\TextAt{14}{495.617}{159.972}{a}%
\TextAt{8}{501.6}{159.972}{.}%
\TextAt{8}{509.947}{159.972}{This}%
\TextAt{8}{79.226}{142.926}{cannot}%
\TextAt{8}{116.673}{142.926}{be}%
\TextAt{8}{131.667}{142.926}{a}%
\TextAt{8}{140.493}{142.926}{prime}%
\TextAt{8}{172.727}{142.926}{ideal,}%
\TextAt{8}{203.198}{142.926}{so}%
\TextAt{8}{216.638}{142.926}{there}%
\TextAt{8}{245.62}{142.926}{exist}%
\TextAt{8}{272.391}{142.926}{elements}%
\TextAt{13}{319.318}{142.926}{b}%
\TextAt{17}{324.295}{141.133}{1}%
\TextAt{13}{329.027}{142.926}{\Character{59}}%
\TextAt{13}{334.251}{142.926}{b}%
\TextAt{17}{339.248}{141.133}{2}%
\NativeMathAt{347.302}{142.926}{11.955}{\in}%
\TextAt{14}{358.593}{142.926}{o}%
\TextAt{8}{367.404}{142.926}{such}%
\TextAt{8}{392.872}{142.926}{that}%
\TextAt{13}{417.311}{142.926}{b}%
\TextAt{17}{422.288}{141.133}{1}%
\TextAt{13}{427.02}{142.926}{b}%
\TextAt{17}{431.997}{141.133}{2}%
\NativeMathAt{440.051}{142.926}{11.955}{\in}%
\TextAt{14}{451.342}{142.926}{a}%
\TextAt{8}{457.325}{142.926}{,}%
\TextAt{8}{463.73}{142.926}{but}%
\TextAt{13}{484.26}{142.926}{b}%
\TextAt{17}{489.237}{141.133}{1}%
\TextAt{13}{493.969}{142.926}{\Character{59}}%
\TextAt{13}{499.193}{142.926}{b}%
\TextAt{17}{504.191}{141.133}{2}%
\NativeMathAt{512.244}{142.926}{11.955}{\notin}%
\TextAt{14}{523.535}{142.926}{a}%
\TextAt{8}{529.518}{142.926}{.}%
\TextAt{8}{79.226}{125.88}{Put}%
\TextAt{14}{102.323}{125.88}{a}%
\TextAt{17}{108.306}{124.087}{1}%
\TextAt{9}{116.664}{125.88}{\Character{61}}%
\TextAt{9}{129.384}{125.88}{\Character{40}}%
\TextAt{13}{133.947}{125.88}{b}%
\TextAt{17}{138.924}{124.087}{1}%
\TextAt{9}{143.657}{125.88}{\Character{41}}%
\TextAt{9}{150.974}{125.88}{\Character{43}}%
\TextAt{14}{162.871}{125.88}{a}%
\TextAt{8}{168.854}{125.88}{,}%
\TextAt{14}{176.227}{125.88}{a}%
\TextAt{17}{182.21}{124.087}{2}%
\TextAt{9}{190.568}{125.88}{\Character{61}}%
\TextAt{9}{203.288}{125.88}{\Character{40}}%
\TextAt{13}{207.851}{125.88}{b}%
\TextAt{17}{212.828}{124.087}{2}%
\TextAt{9}{217.56}{125.88}{\Character{41}}%
\TextAt{9}{224.877}{125.88}{\Character{43}}%
\TextAt{14}{236.774}{125.88}{a}%
\TextAt{8}{242.757}{125.88}{.}%
\TextAt{8}{251.743}{125.88}{Then}%
\TextAt{14}{282.479}{125.88}{a}%
\TextAt{18}{292.088}{125.88}{\Character{40}}%
\TextAt{14}{305.012}{125.88}{a}%
\TextAt{17}{310.995}{124.087}{1}%
\TextAt{8}{315.728}{125.88}{,}%
\TextAt{14}{323.1}{125.88}{a}%
\TextAt{18}{332.709}{125.88}{\Character{40}}%
\TextAt{14}{345.634}{125.88}{a}%
\TextAt{17}{351.617}{124.087}{2}%
\TextAt{8}{360.425}{125.88}{and}%
\TextAt{14}{383.367}{125.88}{a}%
\TextAt{17}{389.35}{124.087}{1}%
\TextAt{14}{394.082}{125.88}{a}%
\TextAt{17}{400.065}{124.087}{2}%
\NativeMathAt{408.424}{125.88}{11.955}{\subseteq}%
\TextAt{14}{421.348}{125.88}{a}%
\TextAt{8}{427.331}{125.88}{.}%
\TextAt{8}{436.316}{125.88}{By}%
\TextAt{8}{454.858}{125.88}{the}%
\TextAt{8}{475.194}{125.88}{maximality}%
\TextAt{8}{79.226}{108.834}{of}%
\TextAt{14}{92.281}{108.834}{a}%
\TextAt{8}{98.264}{108.834}{,}%
\TextAt{8}{105.193}{108.834}{both}%
\TextAt{14}{132.558}{108.834}{a}%
\TextAt{17}{138.541}{107.041}{1}%
\TextAt{8}{146.896}{108.834}{and}%
\TextAt{14}{169.384}{108.834}{a}%
\TextAt{17}{175.367}{107.041}{2}%
\TextAt{8}{183.721}{108.834}{contain}%
\TextAt{8}{224.751}{108.834}{a}%
\TextAt{8}{234.231}{108.834}{product}%
\TextAt{8}{277.856}{108.834}{of}%
\TextAt{8}{290.91}{108.834}{prime}%
\TextAt{8}{323.799}{108.834}{ideals,}%
\TextAt{8}{359.409}{108.834}{and}%
\TextAt{8}{381.896}{108.834}{the}%
\TextAt{8}{401.777}{108.834}{product}%
\TextAt{8}{445.402}{108.834}{of}%
\TextAt{8}{458.456}{108.834}{these}%
\TextAt{8}{488.153}{108.834}{products}%
\TextAt{8}{79.226}{91.788}{is}%
\TextAt{8}{90.99}{91.788}{contained}%
\TextAt{8}{143.999}{91.788}{in}%
\TextAt{14}{157.652}{91.788}{a}%
\TextAt{8}{163.635}{91.788}{,}%
\TextAt{8}{170.784}{91.788}{a}%
\TextAt{8}{180.54}{91.788}{contradiction.}%
\Vector{q 0 G 0 g 0.399 w 452.334 -682.212 m 452.334 -674.142 l S Q}%
\Vector{q 0 G 0 g 0.399 w 452.534 -674.342 m 459.707 -674.342 l S Q}%
\Vector{q 0 G 0 g 0.399 w 452.534 -682.013 m 459.707 -682.013 l S Q}%
\Vector{q 0 G 0 g 0.399 w 459.906 -682.212 m 459.906 -674.142 l S Q}%
\TextAt{8}{300.14}{43.033}{17}%
\Destination{78.23}{809.815}{nav.page.17}%
\Destination{79.36}{284.655}{nav.print.19}%
\Destination{79.36}{444.755}{nav.number.1.3.3}%
\Destination{79.36}{283.655}{nav.number.1.3.4}%
}

% Page 30 (body page 18)
\Page{%
\TextAt{2}{79.226}{796.943}{I.}%
\TextAt{2}{89.83}{796.943}{Algebraic}%
\TextAt{2}{138.648}{796.943}{Integers}%
\Vector{q 0 G 0 g 0.299 w 7.226 18.729 m 460.769 18.729 l S Q}%
\TextAt{12}{79.226}{761.124}{Lemma}%
\TextAt{12}{126.891}{761.124}{(3.5).}%
\ErrataAt{162.2}{761.124}{372.8}{Let $\mathfrak p$ be a nonzero prime ideal of $\mathfrak o$ and define}%
\TextAt{14}{243.638}{736.197}{p}%
\TextAt{3}{249.622}{741.133}{\Character{0}}%
\TextAt{17}{256.208}{741.133}{1}%
\TextAt{9}{264.261}{736.197}{\Character{61}}%
\TextAt{16}{276.687}{736.197}{\Character{102}}%
\TextAt{13}{282.664}{736.197}{x}%
\NativeMathAt{292.637}{736.197}{11.955}{\in}%
\TextAt{13}{303.928}{736.197}{K}%
\TextAt{16}{318.067}{736.197}{\Character{106}}%
\TextAt{13}{324.708}{736.197}{x}%
\TextAt{14}{331.36}{736.197}{p}%
\NativeMathAt{340.664}{736.197}{11.955}{\subseteq}%
\TextAt{14}{353.284}{736.197}{o}%
\TextAt{16}{359.128}{736.197}{\Character{103}}%
\TextAt{13}{365.106}{736.197}{\Character{58}}%
\TextAt{8}{79.226}{711.27}{Then}%
\TextAt{8}{109.784}{711.27}{one}%
\TextAt{8}{131.243}{711.27}{has}%
\TextAt{14}{152.117}{711.27}{ap}%
\TextAt{3}{164.083}{715.608}{\Character{0}}%
\TextAt{17}{170.669}{715.608}{1}%
\TextAt{9}{178.723}{711.27}{\Character{58}\Character{61}}%
\TextAt{16}{194.4}{711.27}{\Character{102}}%
\TextAt{20}{200.377}{720.236}{\Character{80}}%
\TextAt{21}{212.997}{707.783}{i}%
\TextAt{13}{218.371}{711.27}{a}%
\TextAt{21}{224.515}{709.477}{i}%
\TextAt{13}{227.897}{711.27}{x}%
\TextAt{21}{234.549}{709.477}{i}%
\TextAt{16}{241.251}{711.27}{\Character{106}}%
\TextAt{13}{247.893}{711.27}{a}%
\TextAt{21}{254.038}{709.477}{i}%
\NativeMathAt{260.74}{711.27}{11.955}{\in}%
\TextAt{14}{272.031}{711.27}{a}%
\TextAt{13}{278.014}{711.27}{\Character{59}}%
\TextAt{13}{283.238}{711.27}{x}%
\TextAt{21}{289.91}{709.477}{i}%
\NativeMathAt{296.612}{711.27}{11.955}{\in}%
\TextAt{14}{307.903}{711.27}{p}%
\TextAt{3}{313.886}{715.608}{\Character{0}}%
\TextAt{17}{320.473}{715.608}{1}%
\TextAt{16}{325.205}{711.27}{\Character{103}}%
\TextAt{18}{334.503}{711.27}{\Character{41}}%
\TextAt{14}{347.123}{711.27}{a}%
\TextAt{8}{353.106}{711.27}{,}%
\TextAt{8}{360.255}{711.27}{for}%
\TextAt{8}{378.14}{711.27}{every}%
\TextAt{8}{409.032}{711.27}{ideal}%
\TextAt{14}{436.995}{711.27}{a}%
\NativeMathAt{446.299}{711.27}{11.955}{\neq}%
\TextAt{9}{458.709}{711.27}{0}%
\TextAt{8}{464.578}{711.27}{.}%
\TextAt{19}{79.226}{678.619}{Proof.}%
\TextAt{8}{116}{678.619}{Let}%
\TextAt{13}{137.036}{678.619}{a}%
\NativeMathAt{146.614}{678.619}{11.955}{\in}%
\TextAt{14}{158.018}{678.619}{p}%
\TextAt{8}{164.001}{678.619}{,}%
\TextAt{13}{171.232}{678.619}{a}%
\NativeMathAt{180.811}{678.619}{11.955}{\neq}%
\TextAt{9}{193.34}{678.619}{0}%
\TextAt{8}{199.202}{678.619}{,}%
\TextAt{8}{206.433}{678.619}{and}%
\TextAt{14}{229.262}{678.619}{p}%
\TextAt{17}{235.245}{676.826}{1}%
\TextAt{14}{239.978}{678.619}{p}%
\TextAt{17}{245.961}{676.826}{2}%
\TextAt{16}{252.685}{678.619}{\Character{1}}%
\TextAt{16}{257.981}{678.619}{\Character{1}}%
\TextAt{16}{263.289}{678.619}{\Character{1}}%
\TextAt{14}{268.626}{678.619}{p}%
\TextAt{21}{274.609}{676.826}{r}%
\NativeMathAt{282.597}{678.619}{11.955}{\subseteq}%
\TextAt{9}{295.329}{678.619}{\Character{40}}%
\TextAt{13}{299.881}{678.619}{a}%
\TextAt{9}{306.026}{678.619}{\Character{41}}%
\NativeMathAt{314.012}{678.619}{11.955}{\subseteq}%
\TextAt{14}{326.744}{678.619}{p}%
\TextAt{8}{332.727}{678.619}{,}%
\TextAt{8}{339.958}{678.619}{with}%
\TextAt{13}{366.685}{678.619}{r}%
\TextAt{8}{376.248}{678.619}{as}%
\TextAt{8}{390.685}{678.619}{small}%
\TextAt{8}{421.38}{678.619}{as}%
\TextAt{8}{435.816}{678.619}{possible.}%
\TextAt{8}{484.584}{678.619}{Then}%
\TextAt{8}{515.207}{678.619}{one}%
\TextAt{8}{79.226}{661.573}{of}%
\TextAt{8}{92.48}{661.573}{the}%
\TextAt{14}{112.56}{661.573}{p}%
\TextAt{21}{118.543}{659.78}{i}%
\TextAt{8}{121.924}{661.573}{,}%
\TextAt{8}{129.012}{661.573}{say}%
\TextAt{14}{149.164}{661.573}{p}%
\TextAt{17}{155.147}{659.78}{1}%
\TextAt{8}{159.88}{661.573}{,}%
\TextAt{8}{166.968}{661.573}{is}%
\TextAt{8}{178.655}{661.573}{contained}%
\TextAt{8}{231.588}{661.573}{in}%
\TextAt{14}{245.164}{661.573}{p}%
\TextAt{8}{251.148}{661.573}{,}%
\TextAt{8}{258.236}{661.573}{and}%
\TextAt{8}{280.922}{661.573}{so}%
\TextAt{14}{295.216}{661.573}{p}%
\TextAt{17}{301.199}{659.78}{1}%
\TextAt{9}{309.252}{661.573}{\Character{61}}%
\TextAt{14}{321.677}{661.573}{p}%
\TextAt{8}{331.482}{661.573}{because}%
\TextAt{14}{374.707}{661.573}{p}%
\TextAt{17}{380.69}{659.78}{1}%
\TextAt{8}{389.243}{661.573}{is}%
\TextAt{8}{400.931}{661.573}{a}%
\TextAt{8}{410.61}{661.573}{maximal}%
\TextAt{8}{458.342}{661.573}{ideal.}%
\TextAt{8}{490.831}{661.573}{(Indeed,}%
\TextAt{8}{79.226}{644.527}{if}%
\TextAt{8}{90.362}{644.527}{none}%
\TextAt{8}{118.737}{644.527}{of}%
\TextAt{8}{132.479}{644.527}{the}%
\TextAt{14}{153.048}{644.527}{p}%
\TextAt{21}{159.031}{642.734}{i}%
\TextAt{8}{166.721}{644.527}{were}%
\TextAt{8}{194.116}{644.527}{contained}%
\TextAt{8}{247.538}{644.527}{in}%
\TextAt{14}{261.602}{644.527}{p}%
\TextAt{8}{267.586}{644.527}{,}%
\TextAt{8}{275.25}{644.527}{then}%
\TextAt{8}{302.322}{644.527}{for}%
\TextAt{8}{320.619}{644.527}{every}%
\TextAt{13}{351.923}{644.527}{i}%
\TextAt{8}{360.226}{644.527}{there}%
\TextAt{8}{390.55}{644.527}{would}%
\TextAt{8}{425.106}{644.527}{exist}%
\TextAt{13}{453.218}{644.527}{a}%
\TextAt{21}{459.363}{642.734}{i}%
\NativeMathAt{466.767}{644.527}{11.955}{\in}%
\TextAt{14}{478.76}{644.527}{p}%
\TextAt{21}{484.743}{642.734}{i}%
\TextAt{16}{491.062}{644.527}{\Character{110}}%
\TextAt{14}{499.977}{644.527}{p}%
\TextAt{8}{510.27}{644.527}{such}%
\TextAt{8}{79.226}{627.481}{that}%
\TextAt{13}{105.262}{627.481}{a}%
\TextAt{17}{111.407}{625.688}{1}%
\TextAt{16}{118.132}{627.481}{\Character{1}}%
\TextAt{16}{123.428}{627.481}{\Character{1}}%
\TextAt{16}{128.736}{627.481}{\Character{1}}%
\TextAt{13}{134.072}{627.481}{a}%
\TextAt{21}{140.217}{625.688}{r}%
\NativeMathAt{149.229}{627.481}{11.955}{\in}%
\TextAt{14}{161.656}{627.481}{p}%
\TextAt{8}{167.639}{627.481}{.}%
\TextAt{8}{178.088}{627.481}{But}%
\TextAt{14}{201.995}{627.481}{p}%
\TextAt{8}{212.542}{627.481}{is}%
\TextAt{8}{224.973}{627.481}{prime.)}%
\TextAt{8}{269.243}{627.481}{Since}%
\TextAt{14}{300.467}{627.481}{p}%
\TextAt{17}{306.45}{625.688}{2}%
\TextAt{16}{313.175}{627.481}{\Character{1}}%
\TextAt{16}{318.471}{627.481}{\Character{1}}%
\TextAt{16}{323.779}{627.481}{\Character{1}}%
\TextAt{14}{329.115}{627.481}{p}%
\TextAt{21}{335.098}{625.688}{r}%
\NativeMathAt{344.11}{627.481}{11.955}{\not\subseteq}%
\TextAt{9}{357.866}{627.481}{\Character{40}}%
\TextAt{13}{362.418}{627.481}{a}%
\TextAt{9}{368.563}{627.481}{\Character{41}}%
\TextAt{8}{373.115}{627.481}{,}%
\TextAt{8}{381.098}{627.481}{there}%
\TextAt{8}{411.677}{627.481}{exists}%
\TextAt{13}{444.658}{627.481}{b}%
\NativeMathAt{454.093}{627.481}{11.955}{\in}%
\TextAt{14}{466.52}{627.481}{p}%
\TextAt{17}{472.503}{625.688}{2}%
\TextAt{16}{479.228}{627.481}{\Character{1}}%
\TextAt{16}{484.524}{627.481}{\Character{1}}%
\TextAt{16}{489.832}{627.481}{\Character{1}}%
\TextAt{14}{495.168}{627.481}{p}%
\TextAt{21}{501.151}{625.688}{r}%
\TextAt{8}{510.27}{627.481}{such}%
\TextAt{8}{79.226}{610.435}{that}%
\TextAt{13}{104.505}{610.435}{b}%
\NativeMathAt{112.803}{610.435}{11.955}{\notin}%
\TextAt{13}{124.094}{610.435}{a}%
\TextAt{14}{130.239}{610.435}{o}%
\TextAt{8}{136.083}{610.435}{,}%
\TextAt{8}{143.161}{610.435}{i.e.,}%
\TextAt{13}{165.194}{610.435}{a}%
\TextAt{3}{171.339}{614.774}{\Character{0}}%
\TextAt{17}{177.926}{614.774}{1}%
\TextAt{13}{182.658}{610.435}{b}%
\NativeMathAt{190.956}{610.435}{11.955}{\notin}%
\TextAt{14}{202.247}{610.435}{o}%
\TextAt{8}{208.091}{610.435}{.}%
\TextAt{8}{216.509}{610.435}{On}%
\TextAt{8}{235.931}{610.435}{the}%
\TextAt{8}{255.997}{610.435}{other}%
\TextAt{8}{286.477}{610.435}{hand}%
\TextAt{8}{315.653}{610.435}{we}%
\TextAt{8}{332.791}{610.435}{have}%
\TextAt{13}{359.696}{610.435}{b}%
\TextAt{14}{364.673}{610.435}{p}%
\NativeMathAt{373.977}{610.435}{11.955}{\subseteq}%
\TextAt{9}{386.597}{610.435}{\Character{40}}%
\TextAt{13}{391.149}{610.435}{a}%
\TextAt{9}{397.294}{610.435}{\Character{41}}%
\TextAt{8}{401.846}{610.435}{,}%
\TextAt{8}{408.924}{610.435}{i.e.,}%
\TextAt{13}{430.957}{610.435}{a}%
\TextAt{3}{437.102}{614.774}{\Character{0}}%
\TextAt{17}{443.688}{614.774}{1}%
\TextAt{13}{448.421}{610.435}{b}%
\TextAt{14}{453.398}{610.435}{p}%
\NativeMathAt{462.702}{610.435}{11.955}{\subseteq}%
\TextAt{14}{475.321}{610.435}{o}%
\TextAt{8}{481.165}{610.435}{,}%
\TextAt{8}{488.243}{610.435}{and}%
\TextAt{8}{510.916}{610.435}{thus}%
\TextAt{13}{79.226}{593.389}{a}%
\TextAt{3}{85.371}{597.728}{\Character{0}}%
\TextAt{17}{91.958}{597.728}{1}%
\TextAt{13}{96.69}{593.389}{b}%
\NativeMathAt{104.988}{593.389}{11.955}{\in}%
\TextAt{14}{116.279}{593.389}{p}%
\TextAt{3}{122.262}{597.728}{\Character{0}}%
\TextAt{17}{128.848}{597.728}{1}%
\TextAt{8}{133.581}{593.389}{.}%
\TextAt{8}{142.029}{593.389}{It}%
\TextAt{8}{154.702}{593.389}{follows}%
\TextAt{8}{193.137}{593.389}{that}%
\TextAt{14}{218.506}{593.389}{p}%
\TextAt{3}{224.489}{597.728}{\Character{0}}%
\TextAt{17}{231.076}{597.728}{1}%
\TextAt{18}{239.129}{593.389}{\Character{41}}%
\TextAt{14}{251.748}{593.389}{o}%
\TextAt{8}{257.593}{593.389}{.}%
\TextAt{8}{97.159}{573.355}{Now}%
\TextAt{8}{124.107}{573.355}{let}%
\TextAt{14}{141.3}{573.355}{a}%
\NativeMathAt{151.094}{573.355}{11.955}{\neq}%
\TextAt{9}{163.994}{573.355}{0}%
\TextAt{8}{174.049}{573.355}{be}%
\TextAt{8}{190.261}{573.355}{an}%
\TextAt{8}{206.809}{573.355}{ideal}%
\TextAt{8}{235.06}{573.355}{of}%
\TextAt{14}{248.678}{573.355}{o}%
\TextAt{8}{258.707}{573.355}{and}%
\TextAt{13}{281.758}{573.355}{a}%
\TextAt{17}{287.903}{571.561}{1}%
\TextAt{13}{292.635}{573.355}{\Character{59}}%
\TextAt{13}{297.859}{573.355}{\Character{58}}%
\TextAt{13}{303.096}{573.355}{\Character{58}}%
\TextAt{13}{308.332}{573.355}{\Character{58}}%
\TextAt{13}{313.557}{573.355}{\Character{59}}%
\TextAt{13}{318.793}{573.355}{a}%
\TextAt{21}{325.001}{571.561}{n}%
\TextAt{8}{334.823}{573.355}{a}%
\TextAt{8}{344.866}{573.355}{system}%
\TextAt{8}{383.973}{573.355}{of}%
\TextAt{8}{397.591}{573.355}{generators.}%
\TextAt{8}{459.661}{573.355}{Let}%
\TextAt{8}{480.919}{573.355}{us}%
\TextAt{8}{496.223}{573.355}{assume}%
\TextAt{8}{79.226}{556.309}{that}%
\TextAt{14}{104.595}{556.309}{ap}%
\TextAt{3}{116.561}{560.647}{\Character{0}}%
\TextAt{17}{123.148}{560.647}{1}%
\TextAt{9}{131.201}{556.309}{\Character{61}}%
\TextAt{14}{143.626}{556.309}{a}%
\TextAt{8}{149.609}{556.309}{.}%
\TextAt{8}{158.058}{556.309}{Then}%
\TextAt{8}{188.615}{556.309}{for}%
\TextAt{8}{206.5}{556.309}{every}%
\TextAt{13}{237.392}{556.309}{x}%
\NativeMathAt{247.365}{556.309}{11.955}{\in}%
\TextAt{14}{258.656}{556.309}{p}%
\TextAt{3}{264.639}{560.647}{\Character{0}}%
\TextAt{17}{271.226}{560.647}{1}%
\TextAt{8}{275.958}{556.309}{,}%
\TextAt{13}{235.3}{531.382}{xa}%
\TextAt{21}{248.097}{529.589}{i}%
\TextAt{9}{254.799}{531.382}{\Character{61}}%
\TextAt{20}{267.225}{542.739}{\Character{88}}%
\TextAt{21}{273.917}{517.545}{j}%
\TextAt{13}{286.486}{531.382}{a}%
\TextAt{21}{292.631}{529.589}{ij}%
\TextAt{13}{299.896}{531.382}{a}%
\TextAt{21}{306.041}{529.589}{j}%
\TextAt{13}{310.423}{531.382}{\Character{59}}%
\TextAt{13}{339.558}{531.382}{a}%
\TextAt{21}{345.723}{529.589}{ij}%
\NativeMathAt{356.309}{531.382}{11.955}{\in}%
\TextAt{14}{367.6}{531.382}{o}%
\TextAt{13}{373.444}{531.382}{\Character{58}}%
\TextAt{8}{79.226}{496.682}{Writing}%
\TextAt{13}{121.82}{496.682}{A}%
\TextAt{8}{134.167}{496.682}{for}%
\TextAt{8}{151.727}{496.682}{the}%
\TextAt{8}{171.558}{496.682}{matrix}%
\TextAt{9}{209.286}{496.682}{\Character{40}}%
\TextAt{13}{213.838}{496.682}{x\Character{14}}%
\TextAt{21}{225.675}{494.889}{ij}%
\TextAt{16}{234.932}{496.682}{\Character{0}}%
\TextAt{13}{246.221}{496.682}{a}%
\TextAt{21}{252.366}{494.889}{ij}%
\TextAt{9}{259.632}{496.682}{\Character{41}}%
\TextAt{8}{267.756}{496.682}{we}%
\TextAt{8}{284.658}{496.682}{obtain}%
\TextAt{13}{320.76}{496.682}{A}%
\TextAt{9}{329.536}{496.682}{\Character{40}}%
\TextAt{13}{334.088}{496.682}{a}%
\TextAt{17}{340.233}{494.889}{1}%
\TextAt{13}{344.965}{496.682}{\Character{59}}%
\TextAt{13}{350.189}{496.682}{\Character{58}}%
\TextAt{13}{355.426}{496.682}{\Character{58}}%
\TextAt{13}{360.662}{496.682}{\Character{58}}%
\TextAt{13}{365.887}{496.682}{\Character{59}}%
\TextAt{13}{371.123}{496.682}{a}%
\TextAt{21}{377.331}{494.889}{n}%
\TextAt{9}{382.967}{496.682}{\Character{41}}%
\TextAt{21}{387.52}{501.02}{t}%
\TextAt{9}{394.397}{496.682}{\Character{61}}%
\TextAt{9}{406.807}{496.682}{0}%
\TextAt{8}{412.675}{496.682}{.}%
\TextAt{8}{421.015}{496.682}{By}%
\TextAt{8}{439.053}{496.682}{(2.3),}%
\TextAt{8}{470.019}{496.682}{the}%
\TextAt{8}{489.85}{496.682}{determi-}%
\TextAt{8}{79.226}{479.636}{nant}%
\TextAt{13}{105.839}{479.636}{d}%
\TextAt{9}{115.242}{479.636}{\Character{61}}%
\TextAt{8}{127.668}{479.636}{det}%
\TextAt{9}{143.927}{479.636}{\Character{40}}%
\TextAt{13}{148.479}{479.636}{A}%
\TextAt{9}{157.254}{479.636}{\Character{41}}%
\TextAt{8}{165.322}{479.636}{satisfies}%
\TextAt{13}{208.049}{479.636}{da}%
\TextAt{17}{220.277}{477.843}{1}%
\TextAt{9}{228.33}{479.636}{\Character{61}}%
\TextAt{16}{240.756}{479.636}{\Character{1}}%
\TextAt{16}{246.052}{479.636}{\Character{1}}%
\TextAt{16}{251.36}{479.636}{\Character{1}}%
\TextAt{9}{258.024}{479.636}{\Character{61}}%
\TextAt{13}{270.45}{479.636}{da}%
\TextAt{21}{282.677}{477.843}{n}%
\TextAt{9}{291.634}{479.636}{\Character{61}}%
\TextAt{9}{304.043}{479.636}{0}%
\TextAt{8}{313.428}{479.636}{and}%
\TextAt{8}{335.808}{479.636}{thus}%
\TextAt{13}{361.177}{479.636}{d}%
\TextAt{9}{370.581}{479.636}{\Character{61}}%
\TextAt{9}{382.99}{479.636}{0}%
\TextAt{8}{388.859}{479.636}{.}%
\TextAt{8}{397.18}{479.636}{It}%
\TextAt{8}{409.47}{479.636}{follows}%
\TextAt{8}{447.524}{479.636}{that}%
\TextAt{13}{472.51}{479.636}{x}%
\TextAt{8}{482.677}{479.636}{is}%
\TextAt{8}{494.059}{479.636}{integral}%
\TextAt{8}{79.226}{462.59}{over}%
\TextAt{14}{104.607}{462.59}{o}%
\TextAt{8}{110.452}{462.59}{,}%
\TextAt{8}{118.02}{462.59}{being}%
\TextAt{8}{149.893}{462.59}{a}%
\TextAt{8}{159.983}{462.59}{zero}%
\TextAt{8}{185.03}{462.59}{of}%
\TextAt{8}{198.695}{462.59}{the}%
\TextAt{8}{219.186}{462.59}{monic}%
\TextAt{8}{253.988}{462.59}{polynomial}%
\TextAt{13}{314.817}{462.59}{f}%
\TextAt{9}{321.863}{462.59}{\Character{40}}%
\TextAt{13}{326.415}{462.59}{X}%
\TextAt{9}{337.071}{462.59}{\Character{41}}%
\TextAt{9}{345.499}{462.59}{\Character{61}}%
\TextAt{8}{358.511}{462.59}{det}%
\TextAt{9}{374.77}{462.59}{\Character{40}}%
\TextAt{13}{379.323}{462.59}{X}%
\TextAt{13}{389.963}{462.59}{\Character{14}}%
\TextAt{21}{395.163}{460.797}{ij}%
\TextAt{16}{405.313}{462.59}{\Character{0}}%
\TextAt{13}{417.497}{462.59}{a}%
\TextAt{21}{423.642}{460.797}{ij}%
\TextAt{9}{430.907}{462.59}{\Character{41}}%
\NativeMathAt{439.351}{462.59}{11.955}{\in}%
\TextAt{14}{451.213}{462.59}{o}%
\TextAt{9}{457.057}{462.59}{\Character{91}}%
\TextAt{13}{460.309}{462.59}{X}%
\TextAt{9}{470.964}{462.59}{\Character{93}}%
\TextAt{8}{474.216}{462.59}{.}%
\TextAt{8}{483.67}{462.59}{Therefore}%
\TextAt{13}{79.226}{445.544}{x}%
\NativeMathAt{89.199}{445.544}{11.955}{\in}%
\TextAt{14}{100.49}{445.544}{o}%
\TextAt{8}{106.334}{445.544}{.}%
\TextAt{8}{114.783}{445.544}{This}%
\TextAt{8}{141.502}{445.544}{means}%
\TextAt{8}{177.332}{445.544}{that}%
\TextAt{14}{202.701}{445.544}{p}%
\TextAt{3}{208.684}{449.882}{\Character{0}}%
\TextAt{17}{215.271}{449.882}{1}%
\TextAt{9}{223.324}{445.544}{\Character{61}}%
\TextAt{14}{235.749}{445.544}{o}%
\TextAt{8}{241.593}{445.544}{,}%
\TextAt{8}{248.743}{445.544}{a}%
\TextAt{8}{258.498}{445.544}{contradiction.}%
\Vector{q 0 G 0 g 0.399 w 452.334 -328.456 m 452.334 -320.386 l S Q}%
\Vector{q 0 G 0 g 0.399 w 452.534 -320.586 m 459.707 -320.586 l S Q}%
\Vector{q 0 G 0 g 0.399 w 452.534 -328.257 m 459.707 -328.257 l S Q}%
\Vector{q 0 G 0 g 0.399 w 459.906 -328.456 m 459.906 -320.386 l S Q}%
\TextAt{19}{79.226}{417.707}{Proof}%
\TextAt{19}{110.62}{417.707}{of}%
\TextAt{19}{124.369}{417.707}{(3.3).}%
\TextAt{12}{159.039}{417.707}{I.}%
\TextAt{12}{172.822}{417.707}{Existence}%
\TextAt{12}{233.781}{417.707}{of}%
\TextAt{12}{249.656}{417.707}{the}%
\TextAt{12}{273.529}{417.707}{prime}%
\TextAt{12}{312.622}{417.707}{ideal}%
\TextAt{12}{345.282}{417.707}{factorization.}%
\TextAt{8}{429.189}{417.707}{Let}%
\TextAt{14}{450.635}{417.707}{M}%
\TextAt{8}{467.557}{417.707}{be}%
\TextAt{8}{483.959}{417.707}{the}%
\TextAt{8}{504.592}{417.707}{set}%
\TextAt{8}{523.337}{417.707}{of}%
\TextAt{8}{79.226}{400.661}{all}%
\TextAt{8}{95.819}{400.661}{ideals}%
\TextAt{8}{128.731}{400.661}{different}%
\TextAt{8}{175.236}{400.661}{from}%
\TextAt{9}{203.211}{400.661}{\Character{40}0\Character{41}}%
\TextAt{8}{222.4}{400.661}{and}%
\TextAt{9}{245.496}{400.661}{\Character{40}1\Character{41}}%
\TextAt{8}{264.685}{400.661}{which}%
\TextAt{8}{298.506}{400.661}{do}%
\TextAt{8}{315.099}{400.661}{not}%
\TextAt{8}{336.247}{400.661}{admit}%
\TextAt{8}{370.402}{400.661}{a}%
\TextAt{8}{380.492}{400.661}{prime}%
\TextAt{8}{413.99}{400.661}{ideal}%
\TextAt{8}{442.287}{400.661}{decomposition.}%
\TextAt{8}{524.975}{400.661}{If}%
\TextAt{14}{79.226}{383.615}{M}%
\TextAt{8}{95.944}{383.615}{is}%
\TextAt{8}{107.98}{383.615}{nonempty,}%
\TextAt{8}{165.228}{383.615}{then}%
\TextAt{8}{192.16}{383.615}{we}%
\TextAt{8}{209.66}{383.615}{argue}%
\TextAt{8}{241.806}{383.615}{as}%
\TextAt{8}{256.448}{383.615}{for}%
\TextAt{8}{274.606}{383.615}{(3.4)}%
\TextAt{8}{302.854}{383.615}{that}%
\TextAt{8}{328.495}{383.615}{there}%
\TextAt{8}{358.68}{383.615}{exists}%
\TextAt{8}{391.268}{383.615}{a}%
\TextAt{8}{401.296}{383.615}{maximal}%
\TextAt{8}{449.377}{383.615}{element}%
\TextAt{14}{492.892}{383.615}{a}%
\TextAt{8}{503.045}{383.615}{in}%
\TextAt{14}{516.97}{383.615}{M}%
\TextAt{8}{529.518}{383.615}{.}%
\TextAt{8}{79.226}{366.569}{It}%
\TextAt{8}{91.899}{366.569}{is}%
\TextAt{8}{103.663}{366.569}{contained}%
\TextAt{8}{156.672}{366.569}{in}%
\TextAt{8}{170.325}{366.569}{a}%
\TextAt{8}{180.08}{366.569}{maximal}%
\TextAt{8}{227.889}{366.569}{ideal}%
\TextAt{14}{255.852}{366.569}{p}%
\TextAt{8}{261.835}{366.569}{,}%
\TextAt{8}{268.984}{366.569}{and}%
\TextAt{8}{291.747}{366.569}{the}%
\TextAt{8}{311.903}{366.569}{inclusion}%
\TextAt{14}{360.74}{366.569}{o}%
\NativeMathAt{369.905}{366.569}{11.955}{\subseteq}%
\TextAt{14}{382.524}{366.569}{p}%
\TextAt{3}{388.508}{370.907}{\Character{0}}%
\TextAt{17}{395.094}{370.907}{1}%
\TextAt{8}{403.724}{366.569}{gives}%
\TextAt{8}{432.404}{366.569}{us}%
\TextAt{14}{251.263}{341.642}{a}%
\NativeMathAt{260.567}{341.642}{11.955}{\subseteq}%
\TextAt{14}{273.186}{341.642}{ap}%
\TextAt{3}{285.153}{346.578}{\Character{0}}%
\TextAt{17}{291.739}{346.578}{1}%
\NativeMathAt{299.792}{341.642}{11.955}{\subseteq}%
\TextAt{14}{312.412}{341.642}{pp}%
\TextAt{3}{324.378}{346.578}{\Character{0}}%
\TextAt{17}{330.964}{346.578}{1}%
\NativeMathAt{339.017}{341.642}{11.955}{\subseteq}%
\TextAt{14}{351.637}{341.642}{o}%
\TextAt{13}{357.481}{341.642}{\Character{58}}%
\TextAt{8}{79.226}{316.715}{By}%
\TextAt{8}{98.603}{316.715}{(3.5),}%
\TextAt{8}{131.097}{316.715}{one}%
\TextAt{8}{153.57}{316.715}{has}%
\TextAt{14}{175.457}{316.715}{a}%
\TextAt{18}{186.488}{316.715}{\Character{40}}%
\TextAt{14}{200.835}{316.715}{ap}%
\TextAt{3}{212.801}{321.054}{\Character{0}}%
\TextAt{17}{219.388}{321.054}{1}%
\TextAt{8}{229.031}{316.715}{and}%
\TextAt{14}{252.807}{316.715}{p}%
\TextAt{18}{263.838}{316.715}{\Character{40}}%
\TextAt{14}{278.185}{316.715}{pp}%
\TextAt{3}{290.151}{321.054}{\Character{0}}%
\TextAt{17}{296.738}{321.054}{1}%
\NativeMathAt{306.518}{316.715}{11.955}{\subseteq}%
\TextAt{14}{320.865}{316.715}{o}%
\TextAt{8}{326.709}{316.715}{.}%
\TextAt{8}{338.198}{316.715}{Since}%
\TextAt{14}{369.769}{316.715}{p}%
\TextAt{8}{380.663}{316.715}{is}%
\TextAt{8}{393.44}{316.715}{a}%
\TextAt{8}{404.209}{316.715}{maximal}%
\TextAt{8}{453.032}{316.715}{ideal,}%
\TextAt{8}{485.513}{316.715}{it}%
\TextAt{8}{498.231}{316.715}{follows}%
\TextAt{8}{79.226}{299.669}{that}%
\TextAt{14}{104.881}{299.669}{pp}%
\TextAt{3}{116.848}{304.008}{\Character{0}}%
\TextAt{17}{123.434}{304.008}{1}%
\TextAt{9}{131.975}{299.669}{\Character{61}}%
\TextAt{14}{144.889}{299.669}{o}%
\TextAt{8}{150.733}{299.669}{.}%
\TextAt{8}{160.04}{299.669}{In}%
\TextAt{8}{174.948}{299.669}{view}%
\TextAt{8}{202.217}{299.669}{of}%
\TextAt{8}{215.833}{299.669}{the}%
\TextAt{8}{236.276}{299.669}{maximality}%
\TextAt{8}{298.036}{299.669}{of}%
\TextAt{14}{311.652}{299.669}{a}%
\TextAt{8}{321.819}{299.669}{in}%
\TextAt{14}{335.758}{299.669}{M}%
\TextAt{8}{352.489}{299.669}{and}%
\TextAt{8}{375.538}{299.669}{since}%
\TextAt{14}{404.493}{299.669}{a}%
\NativeMathAt{414.285}{299.669}{11.955}{\neq}%
\TextAt{14}{427.199}{299.669}{p}%
\TextAt{8}{433.182}{299.669}{,}%
\TextAt{8}{440.689}{299.669}{i.e.,}%
\TextAt{14}{463.152}{299.669}{ap}%
\TextAt{3}{475.118}{304.008}{\Character{0}}%
\TextAt{17}{481.705}{304.008}{1}%
\NativeMathAt{490.246}{299.669}{11.955}{\neq}%
\TextAt{14}{503.159}{299.669}{o}%
\TextAt{8}{509.003}{299.669}{,}%
\TextAt{8}{516.511}{299.669}{the}%
\TextAt{8}{79.226}{282.623}{ideal}%
\TextAt{14}{107.155}{282.623}{ap}%
\TextAt{3}{119.121}{286.962}{\Character{0}}%
\TextAt{17}{125.707}{286.962}{1}%
\TextAt{8}{134.302}{282.623}{admits}%
\TextAt{8}{172.704}{282.623}{a}%
\TextAt{8}{182.424}{282.623}{prime}%
\TextAt{8}{215.553}{282.623}{ideal}%
\TextAt{8}{243.482}{282.623}{decomposition}%
\TextAt{14}{320.582}{282.623}{ap}%
\TextAt{3}{332.548}{286.962}{\Character{0}}%
\TextAt{17}{339.135}{286.962}{1}%
\TextAt{9}{347.188}{282.623}{\Character{61}}%
\TextAt{14}{359.613}{282.623}{p}%
\TextAt{17}{365.596}{280.83}{1}%
\TextAt{16}{372.321}{282.623}{\Character{1}}%
\TextAt{16}{377.617}{282.623}{\Character{1}}%
\TextAt{16}{382.925}{282.623}{\Character{1}}%
\TextAt{14}{388.261}{282.623}{p}%
\TextAt{21}{394.245}{280.83}{r}%
\TextAt{8}{398.799}{282.623}{,}%
\TextAt{8}{405.921}{282.623}{and}%
\TextAt{8}{428.649}{282.623}{so}%
\TextAt{8}{442.984}{282.623}{does}%
\TextAt{14}{469.347}{282.623}{a}%
\TextAt{9}{478.651}{282.623}{\Character{61}}%
\TextAt{14}{491.076}{282.623}{ap}%
\TextAt{3}{503.042}{286.962}{\Character{0}}%
\TextAt{17}{509.629}{286.962}{1}%
\TextAt{14}{514.361}{282.623}{p}%
\TextAt{9}{523.665}{282.623}{\Character{61}}%
\TextAt{14}{79.226}{265.577}{p}%
\TextAt{17}{85.209}{263.784}{1}%
\TextAt{16}{91.934}{265.577}{\Character{1}}%
\TextAt{16}{97.23}{265.577}{\Character{1}}%
\TextAt{16}{102.538}{265.577}{\Character{1}}%
\TextAt{14}{107.874}{265.577}{p}%
\TextAt{21}{113.857}{263.784}{r}%
\TextAt{14}{118.412}{265.577}{p}%
\TextAt{8}{124.395}{265.577}{,}%
\TextAt{8}{131.544}{265.577}{a}%
\TextAt{8}{141.3}{265.577}{contradiction.}%
\TextAt{12}{97.159}{245.542}{II.}%
\TextAt{12}{115.773}{245.542}{Uniqueness}%
\TextAt{12}{186.739}{245.542}{of}%
\TextAt{12}{202.066}{245.542}{the}%
\TextAt{12}{225.39}{245.542}{prime}%
\TextAt{12}{263.934}{245.542}{ideal}%
\TextAt{12}{296.045}{245.542}{factorization.}%
\TextAt{8}{378.52}{245.542}{For}%
\TextAt{8}{399.489}{245.542}{a}%
\TextAt{8}{409.245}{245.542}{prime}%
\TextAt{8}{442.408}{245.542}{ideal}%
\TextAt{14}{470.371}{245.542}{p}%
\TextAt{8}{480.252}{245.542}{one}%
\TextAt{8}{501.711}{245.542}{has:}%
\TextAt{14}{141.963}{220.616}{ab}%
\NativeMathAt{157.401}{220.616}{11.955}{\subseteq}%
\TextAt{14}{170.021}{220.616}{p}%
\NativeArrowAt{183.222}{220.616}{11.955}{19.067}{\Longrightarrow}%
\TextAt{14}{209.508}{220.616}{a}%
\NativeMathAt{218.812}{220.616}{11.955}{\subseteq}%
\TextAt{14}{231.431}{220.616}{p}%
\TextAt{8}{241.311}{220.616}{or}%
\TextAt{14}{255.622}{220.616}{b}%
\NativeMathAt{265.077}{220.616}{11.955}{\subseteq}%
\TextAt{14}{277.697}{220.616}{p}%
\TextAt{13}{283.68}{220.616}{\Character{59}}%
\TextAt{8}{312.834}{220.616}{i.e.,}%
\TextAt{14}{342.997}{220.616}{p}%
\TextAt{16}{352.301}{220.616}{\Character{106}}%
\TextAt{14}{358.943}{220.616}{ab}%
\NativeArrowAt{378.279}{220.616}{11.955}{19.067}{\Longrightarrow}%
\TextAt{14}{404.565}{220.616}{p}%
\TextAt{16}{413.868}{220.616}{\Character{106}}%
\TextAt{14}{420.51}{220.616}{a}%
\TextAt{8}{430.391}{220.616}{or}%
\TextAt{14}{444.701}{220.616}{p}%
\TextAt{16}{454.005}{220.616}{\Character{106}}%
\TextAt{14}{460.647}{220.616}{b}%
\TextAt{13}{466.781}{220.616}{\Character{58}}%
\TextAt{8}{79.226}{195.689}{Let}%
\TextAt{14}{237.632}{178.643}{a}%
\TextAt{9}{246.936}{178.643}{\Character{61}}%
\TextAt{14}{259.362}{178.643}{p}%
\TextAt{17}{265.345}{176.849}{1}%
\TextAt{14}{270.077}{178.643}{p}%
\TextAt{17}{276.06}{176.849}{2}%
\TextAt{16}{282.785}{178.643}{\Character{1}}%
\TextAt{16}{288.081}{178.643}{\Character{1}}%
\TextAt{16}{293.389}{178.643}{\Character{1}}%
\TextAt{14}{298.725}{178.643}{p}%
\TextAt{21}{304.708}{176.849}{r}%
\TextAt{9}{312.584}{178.643}{\Character{61}}%
\TextAt{14}{325.01}{178.643}{q}%
\TextAt{17}{330.857}{176.849}{1}%
\TextAt{14}{335.589}{178.643}{q}%
\TextAt{17}{341.437}{176.849}{2}%
\TextAt{16}{348.162}{178.643}{\Character{1}}%
\TextAt{16}{353.458}{178.643}{\Character{1}}%
\TextAt{16}{358.766}{178.643}{\Character{1}}%
\TextAt{14}{364.102}{178.643}{q}%
\TextAt{21}{369.949}{176.849}{s}%
\TextAt{8}{79.226}{156.867}{be}%
\TextAt{8}{96.256}{156.867}{two}%
\TextAt{8}{119.479}{156.867}{prime}%
\TextAt{8}{153.748}{156.867}{ideal}%
\TextAt{8}{182.817}{156.867}{factorizations}%
\TextAt{8}{256.514}{156.867}{of}%
\TextAt{14}{270.95}{156.867}{a}%
\TextAt{8}{276.933}{156.867}{.}%
\TextAt{8}{288.698}{156.867}{Then}%
\TextAt{14}{320.362}{156.867}{p}%
\TextAt{17}{326.345}{155.074}{1}%
\TextAt{8}{336.08}{156.867}{divides}%
\TextAt{8}{376.59}{156.867}{a}%
\TextAt{8}{387.451}{156.867}{factor}%
\TextAt{14}{422.055}{156.867}{q}%
\TextAt{21}{427.903}{155.074}{j}%
\TextAt{8}{432.285}{156.867}{,}%
\TextAt{8}{440.816}{156.867}{say}%
\TextAt{14}{462.15}{156.867}{q}%
\TextAt{17}{467.997}{155.074}{1}%
\TextAt{8}{472.73}{156.867}{,}%
\TextAt{8}{481.261}{156.867}{and}%
\TextAt{8}{505.129}{156.867}{being}%
\TextAt{8}{79.226}{139.821}{maximal}%
\TextAt{8}{127.035}{139.821}{equals}%
\TextAt{14}{162.542}{139.821}{q}%
\TextAt{17}{168.389}{138.028}{1}%
\TextAt{8}{173.121}{139.821}{.}%
\TextAt{8}{181.57}{139.821}{We}%
\TextAt{8}{201.714}{139.821}{multiply}%
\TextAt{8}{248.543}{139.821}{by}%
\TextAt{14}{264.802}{139.821}{p}%
\TextAt{3}{270.785}{144.769}{\Character{0}}%
\TextAt{17}{277.371}{144.769}{1}%
\TextAt{17}{270.785}{136.89}{1}%
\TextAt{8}{286.001}{139.821}{and}%
\TextAt{8}{308.764}{139.821}{obtain,}%
\TextAt{8}{348.443}{139.821}{in}%
\TextAt{8}{362.096}{139.821}{view}%
\TextAt{8}{389.078}{139.821}{of}%
\TextAt{14}{402.408}{139.821}{p}%
\TextAt{17}{408.392}{138.028}{1}%
\TextAt{14}{413.124}{139.821}{p}%
\TextAt{3}{419.107}{144.769}{\Character{0}}%
\TextAt{17}{425.694}{144.769}{1}%
\TextAt{17}{419.107}{136.89}{1}%
\TextAt{9}{433.747}{139.821}{\Character{61}}%
\TextAt{14}{446.172}{139.821}{o}%
\TextAt{8}{452.016}{139.821}{,}%
\TextAt{8}{459.166}{139.821}{that}%
\TextAt{14}{257.519}{114.894}{p}%
\TextAt{17}{263.502}{113.101}{2}%
\TextAt{16}{270.227}{114.894}{\Character{1}}%
\TextAt{16}{275.523}{114.894}{\Character{1}}%
\TextAt{16}{280.831}{114.894}{\Character{1}}%
\TextAt{14}{286.167}{114.894}{p}%
\TextAt{21}{292.15}{113.101}{r}%
\TextAt{9}{300.026}{114.894}{\Character{61}}%
\TextAt{14}{312.451}{114.894}{q}%
\TextAt{17}{318.299}{113.101}{2}%
\TextAt{16}{325.023}{114.894}{\Character{1}}%
\TextAt{16}{330.319}{114.894}{\Character{1}}%
\TextAt{16}{335.627}{114.894}{\Character{1}}%
\TextAt{14}{340.964}{114.894}{q}%
\TextAt{21}{346.811}{113.101}{s}%
\TextAt{13}{351.225}{114.894}{\Character{58}}%
\TextAt{8}{79.226}{89.967}{Continuing}%
\TextAt{8}{140.444}{89.967}{like}%
\TextAt{8}{162.627}{89.967}{this}%
\TextAt{8}{186.174}{89.967}{we}%
\TextAt{8}{204.126}{89.967}{see}%
\TextAt{8}{223.763}{89.967}{that}%
\TextAt{13}{249.856}{89.967}{r}%
\TextAt{9}{260.011}{89.967}{\Character{61}}%
\TextAt{13}{273.671}{89.967}{s}%
\TextAt{8}{283.807}{89.967}{and,}%
\TextAt{8}{310.726}{89.967}{possibly}%
\TextAt{8}{356.45}{89.967}{after}%
\TextAt{8}{384.814}{89.967}{renumbering,}%
\TextAt{14}{456.937}{89.967}{p}%
\TextAt{21}{462.92}{88.174}{i}%
\TextAt{9}{470.856}{89.967}{\Character{61}}%
\TextAt{14}{484.516}{89.967}{q}%
\TextAt{21}{490.363}{88.174}{i}%
\TextAt{8}{493.744}{89.967}{,}%
\TextAt{8}{501.799}{89.967}{for}%
\TextAt{8}{520.408}{89.967}{all}%
\TextAt{13}{79.226}{72.921}{i}%
\TextAt{9}{86.54}{72.921}{\Character{61}}%
\TextAt{9}{98.95}{72.921}{1}%
\TextAt{13}{104.819}{72.921}{\Character{59}}%
\TextAt{13}{110.043}{72.921}{\Character{58}}%
\TextAt{13}{115.28}{72.921}{\Character{58}}%
\TextAt{13}{120.504}{72.921}{\Character{58}}%
\TextAt{13}{125.741}{72.921}{\Character{59}}%
\TextAt{13}{130.977}{72.921}{r}%
\TextAt{8}{136.64}{72.921}{.}%
\Vector{q 0 G 0 g 0.399 w 452.334 -701.079 m 452.334 -693.009 l S Q}%
\Vector{q 0 G 0 g 0.399 w 452.534 -693.208 m 459.707 -693.208 l S Q}%
\Vector{q 0 G 0 g 0.399 w 452.534 -700.879 m 459.707 -700.879 l S Q}%
\Vector{q 0 G 0 g 0.399 w 459.906 -701.079 m 459.906 -693.009 l S Q}%
\TextAt{8}{300.14}{43.033}{18}%
\Destination{78.23}{809.815}{nav.page.18}%
\Destination{79.36}{415.655}{nav.print.20}%
\Destination{79.36}{775.155}{nav.number.1.3.5}%
\Hotlink{438.553}{494.455}{25.078}{11.047}{nav.number.1.2.3}%
\Hotlink{123.869}{415.555}{26.106}{11.047}{nav.number.1.3.3}%
\Hotlink{274.106}{381.455}{25.078}{11.047}{nav.number.1.3.4}%
\Hotlink{98.103}{314.555}{25.078}{11.047}{nav.number.1.3.5}%
\Destination{79.226}{778.124}{errata.46}%
}

% Page 31 (body page 19)
\Page{%
\TextAt{2}{79.226}{796.943}{I.}%
\TextAt{2}{89.83}{796.943}{Algebraic}%
\TextAt{2}{138.648}{796.943}{Integers}%
\Vector{q 0 G 0 g 0.299 w 7.226 18.729 m 460.769 18.729 l S Q}%
\TextAt{8}{97.159}{761.124}{Grouping}%
\TextAt{8}{149.068}{761.124}{together}%
\TextAt{8}{195.048}{761.124}{the}%
\TextAt{8}{215.001}{761.124}{occurrences}%
\TextAt{8}{277.61}{761.124}{of}%
\TextAt{8}{290.737}{761.124}{the}%
\TextAt{8}{310.69}{761.124}{same}%
\TextAt{8}{339.813}{761.124}{prime}%
\TextAt{8}{372.773}{761.124}{ideals}%
\TextAt{8}{405.148}{761.124}{in}%
\TextAt{8}{418.598}{761.124}{the}%
\TextAt{8}{438.551}{761.124}{prime}%
\TextAt{8}{471.511}{761.124}{ideal}%
\TextAt{8}{499.271}{761.124}{factor-}%
\TextAt{8}{79.226}{744.078}{ization}%
\TextAt{8}{117.602}{744.078}{of}%
\TextAt{8}{130.932}{744.078}{an}%
\TextAt{8}{147.191}{744.078}{ideal}%
\TextAt{14}{175.154}{744.078}{a}%
\NativeMathAt{184.458}{744.078}{11.955}{\neq}%
\TextAt{9}{196.868}{744.078}{0}%
\TextAt{8}{206.634}{744.078}{of}%
\TextAt{14}{219.964}{744.078}{o}%
\TextAt{8}{225.808}{744.078}{,}%
\TextAt{8}{232.958}{744.078}{gives}%
\TextAt{8}{261.638}{744.078}{a}%
\TextAt{8}{271.394}{744.078}{product}%
\TextAt{8}{315.293}{744.078}{representation}%
\TextAt{14}{239.783}{719.499}{a}%
\TextAt{9}{249.087}{719.499}{\Character{61}}%
\TextAt{14}{261.513}{719.499}{p}%
\TextAt{21}{267.496}{724.724}{\Character{23}}%
\TextAt{24}{271.658}{723.617}{1}%
\TextAt{17}{267.496}{716.568}{1}%
\TextAt{16}{278.3}{719.499}{\Character{1}}%
\TextAt{16}{283.596}{719.499}{\Character{1}}%
\TextAt{16}{288.904}{719.499}{\Character{1}}%
\TextAt{14}{294.24}{719.499}{p}%
\TextAt{21}{300.223}{724.435}{\Character{23}}%
\TextAt{22}{304.386}{723.439}{r}%
\TextAt{21}{300.223}{716.544}{r}%
\TextAt{13}{309.042}{719.499}{\Character{59}}%
\TextAt{13}{338.177}{719.499}{\Character{23}}%
\TextAt{21}{343.98}{717.706}{i}%
\TextAt{13}{350.682}{719.499}{\Character{62}}%
\TextAt{9}{363.108}{719.499}{0}%
\TextAt{13}{368.961}{719.499}{\Character{58}}%
\TextAt{8}{79.226}{694.921}{In}%
\TextAt{8}{93.787}{694.921}{the}%
\TextAt{8}{113.884}{694.921}{sequel}%
\TextAt{8}{148.673}{694.921}{such}%
\TextAt{8}{175.01}{694.921}{an}%
\TextAt{8}{191.209}{694.921}{identity}%
\TextAt{8}{234.403}{694.921}{will}%
\TextAt{8}{256.448}{694.921}{be}%
\TextAt{8}{272.312}{694.921}{automatically}%
\TextAt{8}{346.087}{694.921}{understood}%
\TextAt{8}{407.226}{694.921}{to}%
\TextAt{8}{421.476}{694.921}{signify}%
\TextAt{8}{458.549}{694.921}{that}%
\TextAt{8}{483.858}{694.921}{the}%
\TextAt{14}{503.954}{694.921}{p}%
\TextAt{21}{509.937}{693.127}{i}%
\TextAt{8}{517.156}{694.921}{are}%
\TextAt{8}{79.226}{677.875}{pairwise}%
\TextAt{8}{124.959}{677.875}{distinct.}%
\TextAt{8}{172.286}{677.875}{If}%
\TextAt{8}{184.126}{677.875}{in}%
\TextAt{8}{197.926}{677.875}{particular}%
\TextAt{14}{252.063}{677.875}{a}%
\TextAt{8}{262.091}{677.875}{is}%
\TextAt{8}{274.003}{677.875}{a}%
\TextAt{8}{283.906}{677.875}{principal}%
\TextAt{8}{332.831}{677.875}{ideal}%
\TextAt{9}{360.941}{677.875}{\Character{40}}%
\TextAt{13}{365.494}{677.875}{a}%
\TextAt{9}{371.639}{677.875}{\Character{41}}%
\TextAt{8}{376.191}{677.875}{,}%
\TextAt{8}{383.525}{677.875}{then—following}%
\TextAt{8}{467.574}{677.875}{the}%
\TextAt{8}{487.878}{677.875}{tradition}%
\TextAt{8}{79.226}{660.829}{which}%
\TextAt{8}{113.109}{660.829}{tends}%
\TextAt{8}{144.781}{660.829}{to}%
\TextAt{8}{159.488}{660.829}{attribute}%
\TextAt{8}{209.32}{660.829}{to}%
\TextAt{8}{224.027}{660.829}{the}%
\TextAt{8}{244.58}{660.829}{ideals}%
\TextAt{8}{277.555}{660.829}{the}%
\TextAt{8}{298.109}{660.829}{role}%
\TextAt{8}{321.268}{660.829}{of}%
\TextAt{8}{334.995}{660.829}{“ideal}%
\TextAt{8}{368.878}{660.829}{numbers”—we}%
\TextAt{8}{447.044}{660.829}{will}%
\TextAt{8}{469.546}{660.829}{write}%
\TextAt{8}{499.854}{660.829}{with}%
\TextAt{8}{526.911}{660.829}{a}%
\TextAt{8}{79.226}{643.783}{slight}%
\TextAt{8}{110.836}{643.783}{abuse}%
\TextAt{8}{143.413}{643.783}{of}%
\TextAt{8}{156.743}{643.783}{notation}%
\TextAt{13}{269.662}{626.737}{a}%
\TextAt{9}{279.127}{626.737}{\Character{61}}%
\TextAt{14}{291.553}{626.737}{p}%
\TextAt{21}{297.536}{631.962}{\Character{23}}%
\TextAt{24}{301.699}{630.855}{1}%
\TextAt{17}{297.536}{623.806}{1}%
\TextAt{16}{308.34}{626.737}{\Character{1}}%
\TextAt{16}{313.636}{626.737}{\Character{1}}%
\TextAt{16}{318.944}{626.737}{\Character{1}}%
\TextAt{14}{324.28}{626.737}{p}%
\TextAt{21}{330.264}{631.673}{\Character{23}}%
\TextAt{22}{334.426}{630.677}{r}%
\TextAt{21}{330.264}{623.781}{r}%
\TextAt{13}{339.082}{626.737}{\Character{58}}%
\TextAt{8}{79.226}{605.11}{Similarly,}%
\TextAt{8}{130.611}{605.11}{the}%
\TextAt{8}{149.962}{605.11}{notation}%
\TextAt{14}{195.998}{605.11}{a}%
\TextAt{16}{205.302}{605.11}{\Character{106}}%
\TextAt{13}{211.943}{605.11}{a}%
\TextAt{8}{221.18}{605.11}{is}%
\TextAt{8}{232.139}{605.11}{often}%
\TextAt{8}{260.923}{605.11}{used}%
\TextAt{8}{286.838}{605.11}{instead}%
\TextAt{8}{326.417}{605.11}{of}%
\TextAt{14}{338.942}{605.11}{a}%
\TextAt{16}{348.246}{605.11}{\Character{106}}%
\TextAt{9}{354.888}{605.11}{\Character{40}}%
\TextAt{13}{359.44}{605.11}{a}%
\TextAt{9}{365.585}{605.11}{\Character{41}}%
\TextAt{8}{373.23}{605.11}{and}%
\TextAt{9}{395.188}{605.11}{\Character{40}}%
\TextAt{14}{399.74}{605.11}{a}%
\TextAt{13}{405.723}{605.11}{\Character{59}}%
\TextAt{14}{410.967}{605.11}{b}%
\TextAt{9}{417.102}{605.11}{\Character{41}}%
\TextAt{9}{424.957}{605.11}{\Character{61}}%
\TextAt{9}{437.378}{605.11}{1}%
\TextAt{8}{446.346}{605.11}{is}%
\TextAt{8}{457.305}{605.11}{written}%
\TextAt{8}{497.47}{605.11}{for}%
\TextAt{8}{514.55}{605.11}{two}%
\TextAt{8}{79.226}{588.064}{relatively}%
\TextAt{8}{130.667}{588.064}{prime}%
\TextAt{8}{164.211}{588.064}{ideals,}%
\TextAt{8}{200.516}{588.064}{instead}%
\TextAt{8}{241.281}{588.064}{of}%
\TextAt{8}{254.991}{588.064}{the}%
\TextAt{8}{275.528}{588.064}{correct}%
\TextAt{8}{314.93}{588.064}{formula}%
\TextAt{9}{358.242}{588.064}{\Character{40}}%
\TextAt{14}{362.794}{588.064}{a}%
\TextAt{13}{368.777}{588.064}{\Character{59}}%
\TextAt{14}{374.021}{588.064}{b}%
\TextAt{9}{380.156}{588.064}{\Character{41}}%
\TextAt{9}{388.656}{588.064}{\Character{61}}%
\TextAt{14}{401.751}{588.064}{a}%
\TextAt{9}{410.65}{588.064}{\Character{43}}%
\TextAt{14}{422.671}{588.064}{b}%
\TextAt{9}{432.775}{588.064}{\Character{61}}%
\TextAt{14}{445.848}{588.064}{o}%
\TextAt{8}{451.693}{588.064}{.}%
\TextAt{8}{461.282}{588.064}{For}%
\TextAt{8}{482.632}{588.064}{a}%
\TextAt{8}{492.767}{588.064}{product}%
\TextAt{14}{79.226}{571.018}{a}%
\TextAt{9}{88.97}{571.019}{\Character{61}}%
\TextAt{14}{101.834}{571.018}{a}%
\TextAt{17}{107.818}{569.225}{1}%
\TextAt{16}{114.542}{571.019}{\Character{1}}%
\TextAt{16}{119.838}{571.019}{\Character{1}}%
\TextAt{16}{125.146}{571.019}{\Character{1}}%
\TextAt{14}{130.483}{571.019}{a}%
\TextAt{21}{136.466}{569.225}{n}%
\TextAt{8}{146.257}{571.018}{of}%
\TextAt{8}{159.845}{571.018}{relatively}%
\TextAt{8}{211.164}{571.018}{prime}%
\TextAt{8}{244.585}{571.018}{ideals}%
\TextAt{14}{277.421}{571.018}{a}%
\TextAt{17}{283.404}{569.225}{1}%
\TextAt{13}{288.136}{571.019}{\Character{59}}%
\TextAt{13}{293.36}{571.019}{\Character{58}}%
\TextAt{13}{298.597}{571.019}{\Character{58}}%
\TextAt{13}{303.833}{571.019}{\Character{58}}%
\TextAt{13}{309.058}{571.019}{\Character{59}}%
\TextAt{14}{314.357}{571.019}{a}%
\TextAt{21}{320.34}{569.225}{n}%
\TextAt{8}{325.976}{571.018}{,}%
\TextAt{8}{333.448}{571.018}{one}%
\TextAt{8}{355.165}{571.018}{has}%
\TextAt{8}{376.297}{571.018}{an}%
\TextAt{8}{392.814}{571.018}{analogue}%
\TextAt{8}{441.861}{571.018}{of}%
\TextAt{8}{455.449}{571.018}{the}%
\TextAt{8}{475.863}{571.018}{well-known}%
\TextAt{8}{79.226}{553.973}{“Chinese}%
\TextAt{8}{128.589}{553.973}{Remainder}%
\TextAt{8}{188.15}{553.973}{Theorem”}%
\TextAt{8}{243.312}{553.973}{from}%
\TextAt{8}{271.168}{553.973}{elementary}%
\TextAt{8}{331.219}{553.973}{number}%
\TextAt{8}{374.031}{553.973}{theory.}%
\TextAt{8}{414.998}{553.973}{We}%
\TextAt{8}{435.358}{553.973}{may}%
\TextAt{8}{460.942}{553.973}{formulate}%
\TextAt{8}{513.844}{553.973}{this}%
\TextAt{8}{79.226}{536.927}{result}%
\TextAt{8}{111.804}{536.927}{for}%
\TextAt{8}{129.689}{536.927}{an}%
\TextAt{8}{145.948}{536.927}{arbitrary}%
\TextAt{8}{195.717}{536.927}{ring}%
\TextAt{8}{219.783}{536.927}{taking}%
\TextAt{8}{255.888}{536.927}{into}%
\TextAt{8}{279.631}{536.927}{account}%
\TextAt{8}{322.884}{536.927}{that}%
\TextAt{14}{280.977}{506.376}{a}%
\TextAt{9}{290.281}{506.376}{\Character{61}}%
\TextAt{21}{306.989}{521.32}{n}%
\TextAt{20}{302.916}{517.733}{\Character{92}}%
\TextAt{21}{302.706}{492.539}{i}%
\TextAt{17}{305.589}{492.539}{\Character{61}1}%
\TextAt{14}{318.403}{506.376}{a}%
\TextAt{21}{324.386}{504.583}{i}%
\TextAt{13}{327.767}{506.376}{\Character{58}}%
\ErrataAt{79.19}{473.571}{455.81}{For comaximal ideals $A,B$, choose $u\in A$, $v\in B$ with $u+v=1$. If $x\in A\cap B$, then}%
\ErrataAt{79.19}{456.525}{455.81}{$x=xu+xv\in AB$. The reverse inclusion is clear. For pairwise comaximal $\mathfrak a_1,\ldots,\mathfrak a_n$,}%
\ErrataAt{79.19}{439.479}{455.81}{take $b_i\in\mathfrak a_i$, $b_i\equiv1\pmod{\mathfrak a_1}$ for $i>1$. Thus $\mathfrak a_1+\prod_{i>1}\mathfrak a_i=\mathfrak o$. Now induct.}%
\TextAt{12}{79.226}{414.731}{Chinese}%
\TextAt{12}{130.974}{414.731}{Remainder}%
\TextAt{12}{200.977}{414.731}{Theorem}%
\TextAt{12}{259.276}{414.731}{(3.6).}%
\TextAt{8}{295.708}{414.731}{Let}%
\TextAt{14}{317.795}{414.731}{a}%
\TextAt{17}{323.778}{412.937}{1}%
\TextAt{13}{328.51}{414.731}{\Character{59}}%
\TextAt{13}{333.734}{414.731}{\Character{58}}%
\TextAt{13}{338.971}{414.731}{\Character{58}}%
\TextAt{13}{344.207}{414.731}{\Character{58}}%
\TextAt{13}{349.432}{414.731}{\Character{59}}%
\TextAt{14}{354.731}{414.731}{a}%
\TextAt{21}{360.714}{412.937}{n}%
\TextAt{8}{371.365}{414.731}{be}%
\TextAt{8}{388.406}{414.731}{ideals}%
\TextAt{8}{422.101}{414.731}{in}%
\TextAt{8}{436.87}{414.731}{a}%
\TextAt{8}{447.743}{414.731}{ring}%
\TextAt{14}{472.925}{414.731}{o}%
\TextAt{8}{483.784}{414.731}{such}%
\TextAt{8}{511.298}{414.731}{that}%
\TextAt{14}{79.226}{397.685}{a}%
\TextAt{21}{85.209}{395.891}{i}%
\TextAt{9}{91.247}{397.685}{\Character{43}}%
\TextAt{14}{103.008}{397.685}{a}%
\TextAt{21}{108.992}{395.891}{j}%
\TextAt{9}{116.695}{397.685}{\Character{61}}%
\TextAt{14}{129.12}{397.685}{o}%
\TextAt{8}{138.862}{397.685}{for}%
\TextAt{13}{156.747}{397.685}{i}%
\NativeMathAt{164.061}{397.685}{11.955}{\neq}%
\TextAt{13}{176.486}{397.685}{j}%
\TextAt{8}{182.002}{397.685}{.}%
\TextAt{8}{190.45}{397.685}{Then,}%
\TextAt{8}{224.259}{397.685}{if}%
\TextAt{14}{234.983}{397.685}{a}%
\TextAt{9}{244.287}{397.685}{\Character{61}}%
\TextAt{20}{256.713}{406.651}{\Character{81}}%
\TextAt{21}{268.004}{403.496}{n}%
\TextAt{21}{268.004}{394.198}{i}%
\TextAt{17}{270.887}{394.198}{\Character{61}1}%
\TextAt{14}{284.198}{397.685}{a}%
\TextAt{21}{290.181}{395.891}{i}%
\TextAt{8}{293.562}{397.685}{,}%
\TextAt{8}{300.712}{397.685}{one}%
\TextAt{8}{322.171}{397.685}{has}%
\TextAt{14}{268.397}{366.098}{o}%
\TextAt{9}{274.241}{366.098}{\Character{47}}%
\TextAt{14}{280.094}{366.098}{a}%
\NativeMathAt{289.398}{366.098}{11.955}{\simeq}%
\TextAt{21}{307.086}{381.042}{n}%
\TextAt{20}{302.017}{377.455}{\Character{89}}%
\TextAt{21}{302.803}{352.261}{i}%
\TextAt{17}{305.686}{352.261}{\Character{61}1}%
\TextAt{14}{319.286}{366.098}{o}%
\TextAt{9}{325.13}{366.098}{\Character{47}}%
\TextAt{14}{330.983}{366.098}{a}%
\TextAt{21}{336.966}{364.305}{i}%
\TextAt{13}{340.347}{366.098}{\Character{58}}%
\TextAt{19}{79.226}{318.715}{Proof.}%
\TextAt{8}{116}{318.715}{The}%
\TextAt{8}{140.054}{318.715}{canonical}%
\TextAt{8}{191.437}{318.715}{homomorphism}%
\TextAt{14}{196.495}{288.176}{o}%
\NativeArrowAt{205.66}{288.176}{11.955}{19.26}{\longrightarrow}%
\TextAt{21}{233.311}{303.12}{n}%
\TextAt{20}{228.242}{299.533}{\Character{89}}%
\TextAt{21}{229.028}{274.339}{i}%
\TextAt{17}{231.911}{274.339}{\Character{61}1}%
\TextAt{14}{245.511}{288.176}{o}%
\TextAt{9}{251.355}{288.176}{\Character{47}}%
\TextAt{14}{257.208}{288.176}{a}%
\TextAt{21}{263.191}{286.382}{i}%
\TextAt{13}{266.572}{288.176}{\Character{59}}%
\TextAt{13}{295.707}{288.176}{a}%
\NativeArrowAt{305.193}{288.176}{11.955}{19.26}{\longmapsto}%
\TextAt{9}{327.775}{288.176}{\Character{40}}%
\TextAt{13}{332.327}{288.176}{a}%
\TextAt{8}{341.793}{288.176}{mod}%
\TextAt{14}{367.231}{288.176}{a}%
\TextAt{21}{373.214}{286.382}{i}%
\TextAt{9}{376.595}{288.176}{\Character{41}}%
\TextAt{21}{381.147}{286.382}{i}%
\TextAt{17}{384.03}{286.382}{\Character{61}1}%
\TextAt{21}{394.851}{286.382}{\Character{59}\Character{58}\Character{58}\Character{58}\Character{59}n}%
\TextAt{13}{412.249}{288.176}{\Character{59}}%
\TextAt{8}{79.226}{255.648}{has}%
\TextAt{8}{101.605}{255.648}{kernel}%
\TextAt{14}{137.577}{255.648}{a}%
\TextAt{9}{149.446}{255.648}{\Character{61}}%
\TextAt{20}{164.437}{264.614}{\Character{84}}%
\TextAt{21}{174.4}{252.161}{i}%
\TextAt{14}{179.773}{255.648}{a}%
\TextAt{21}{185.756}{253.855}{i}%
\TextAt{8}{189.138}{255.648}{.}%
\TextAt{8}{202.102}{255.648}{It}%
\TextAt{8}{216.28}{255.648}{therefore}%
\TextAt{8}{266.885}{255.648}{suffices}%
\TextAt{8}{308.177}{255.648}{to}%
\TextAt{8}{323.992}{255.648}{show}%
\TextAt{8}{354.501}{255.648}{that}%
\TextAt{8}{381.375}{255.648}{it}%
\TextAt{8}{394.584}{255.648}{is}%
\TextAt{8}{407.854}{255.648}{surjective.}%
\TextAt{8}{469.332}{255.648}{For}%
\TextAt{8}{491.806}{255.648}{this,}%
\TextAt{8}{519.762}{255.648}{let}%
\TextAt{13}{79.226}{238.602}{x}%
\TextAt{21}{85.878}{236.809}{i}%
\TextAt{8}{92.58}{238.602}{mod}%
\TextAt{14}{118.018}{238.602}{a}%
\TextAt{21}{124.001}{236.809}{i}%
\NativeMathAt{130.766}{238.602}{11.955}{\in}%
\TextAt{14}{142.119}{238.602}{o}%
\TextAt{9}{147.963}{238.602}{\Character{47}}%
\TextAt{14}{153.816}{238.602}{a}%
\TextAt{21}{159.799}{236.809}{i}%
\TextAt{8}{163.181}{238.602}{,}%
\TextAt{13}{170.376}{238.602}{i}%
\TextAt{9}{177.752}{238.602}{\Character{61}}%
\TextAt{9}{190.233}{238.602}{1}%
\TextAt{13}{196.093}{238.602}{\Character{59}}%
\TextAt{13}{201.317}{238.602}{\Character{58}}%
\TextAt{13}{206.554}{238.602}{\Character{58}}%
\TextAt{13}{211.79}{238.602}{\Character{58}}%
\TextAt{13}{217.015}{238.602}{\Character{59}}%
\TextAt{13}{222.251}{238.602}{n}%
\TextAt{8}{229.302}{238.602}{,}%
\TextAt{8}{236.497}{238.602}{be}%
\TextAt{8}{252.457}{238.602}{given.}%
\TextAt{8}{287.688}{238.602}{If}%
\TextAt{13}{299.416}{238.602}{n}%
\TextAt{9}{309.787}{238.602}{\Character{61}}%
\TextAt{9}{322.256}{238.602}{2}%
\TextAt{8}{328.128}{238.602}{,}%
\TextAt{8}{335.323}{238.602}{we}%
\TextAt{8}{352.587}{238.602}{may}%
\TextAt{8}{377.992}{238.602}{write}%
\TextAt{9}{407.941}{238.602}{1}%
\TextAt{9}{417.159}{238.602}{\Character{61}}%
\TextAt{13}{429.665}{238.602}{a}%
\TextAt{17}{435.81}{236.809}{1}%
\TextAt{9}{443.224}{238.602}{\Character{43}}%
\TextAt{13}{455.01}{238.602}{a}%
\TextAt{17}{461.155}{236.809}{2}%
\TextAt{8}{465.887}{238.602}{,}%
\TextAt{13}{473.082}{238.602}{a}%
\TextAt{21}{479.227}{236.809}{i}%
\NativeMathAt{485.992}{238.602}{11.955}{\in}%
\TextAt{14}{497.345}{238.602}{a}%
\TextAt{21}{503.328}{236.809}{i}%
\TextAt{8}{506.709}{238.602}{,}%
\TextAt{8}{513.904}{238.602}{and}%
\TextAt{8}{79.226}{221.556}{putting}%
\TextAt{13}{120.854}{221.556}{x}%
\TextAt{9}{130.827}{221.556}{\Character{61}}%
\TextAt{13}{143.252}{221.556}{x}%
\TextAt{17}{149.905}{219.763}{2}%
\TextAt{13}{154.637}{221.556}{a}%
\TextAt{17}{160.782}{219.763}{1}%
\TextAt{9}{168.171}{221.556}{\Character{43}}%
\TextAt{13}{179.932}{221.556}{x}%
\TextAt{17}{186.584}{219.763}{1}%
\TextAt{13}{191.317}{221.556}{a}%
\TextAt{17}{197.461}{219.763}{2}%
\TextAt{8}{206.091}{221.556}{we}%
\TextAt{8}{223.319}{221.556}{get}%
\TextAt{13}{242.829}{221.556}{x}%
\NativeMathAt{252.802}{221.556}{11.955}{\equiv}%
\TextAt{13}{265.422}{221.556}{x}%
\TextAt{21}{272.074}{219.763}{i}%
\TextAt{8}{278.776}{221.556}{mod}%
\TextAt{14}{304.214}{221.556}{a}%
\TextAt{21}{310.197}{219.763}{i}%
\TextAt{8}{313.578}{221.556}{,}%
\TextAt{13}{320.727}{221.556}{i}%
\TextAt{9}{328.042}{221.556}{\Character{61}}%
\TextAt{9}{340.452}{221.556}{1}%
\TextAt{13}{346.32}{221.556}{\Character{59}}%
\TextAt{9}{351.564}{221.556}{2}%
\TextAt{8}{357.417}{221.556}{.}%
\TextAt{8}{97.159}{201.521}{If}%
\TextAt{13}{108.851}{201.521}{n}%
\TextAt{13}{119.144}{201.521}{\Character{62}}%
\TextAt{9}{131.585}{201.521}{2}%
\TextAt{8}{137.438}{201.521}{,}%
\TextAt{8}{144.587}{201.521}{we}%
\TextAt{8}{161.815}{201.521}{may}%
\TextAt{8}{187.183}{201.521}{find}%
\TextAt{8}{210.592}{201.521}{as}%
\TextAt{8}{224.962}{201.521}{before}%
\TextAt{8}{260.074}{201.521}{an}%
\TextAt{8}{276.333}{201.521}{element}%
\TextAt{13}{319.575}{201.521}{y}%
\TextAt{17}{325.283}{199.728}{1}%
\NativeMathAt{333.336}{201.521}{11.955}{\in}%
\TextAt{14}{344.627}{201.521}{o}%
\TextAt{8}{354.368}{201.521}{such}%
\TextAt{8}{380.765}{201.521}{that}%
\TextAt{13}{210.129}{170.983}{y}%
\TextAt{17}{215.836}{169.189}{1}%
\NativeMathAt{223.89}{170.983}{11.955}{\equiv}%
\TextAt{9}{236.509}{170.983}{1}%
\TextAt{8}{245.683}{170.983}{mod}%
\TextAt{14}{271.121}{170.983}{a}%
\TextAt{17}{277.104}{169.189}{1}%
\TextAt{13}{281.836}{170.983}{\Character{59}}%
\TextAt{13}{310.971}{170.983}{y}%
\TextAt{17}{316.698}{169.189}{1}%
\NativeMathAt{324.751}{170.983}{11.955}{\equiv}%
\TextAt{9}{337.371}{170.983}{0}%
\TextAt{8}{346.545}{170.983}{mod}%
\TextAt{21}{377.051}{185.927}{n}%
\TextAt{20}{371.982}{182.34}{\Character{89}}%
\TextAt{21}{372.769}{157.146}{i}%
\TextAt{17}{375.652}{157.146}{\Character{61}2}%
\TextAt{14}{389.251}{170.983}{a}%
\TextAt{21}{395.234}{169.189}{i}%
\TextAt{13}{398.615}{170.983}{\Character{59}}%
\TextAt{8}{79.226}{139.124}{and,}%
\TextAt{8}{105.241}{139.124}{by}%
\TextAt{8}{121.5}{139.124}{the}%
\TextAt{8}{141.656}{139.124}{same}%
\TextAt{8}{170.982}{139.124}{token,}%
\TextAt{8}{206.106}{139.124}{elements}%
\TextAt{13}{253.963}{139.124}{y}%
\TextAt{17}{259.671}{137.331}{2}%
\TextAt{13}{264.403}{139.124}{\Character{59}}%
\TextAt{13}{269.627}{139.124}{\Character{58}}%
\TextAt{13}{274.864}{139.124}{\Character{58}}%
\TextAt{13}{280.088}{139.124}{\Character{58}}%
\TextAt{13}{285.325}{139.124}{\Character{59}}%
\TextAt{13}{290.561}{139.124}{y}%
\TextAt{21}{296.331}{137.331}{n}%
\TextAt{8}{305.865}{139.124}{such}%
\TextAt{8}{332.262}{139.124}{that}%
\TextAt{13}{193.067}{114.546}{y}%
\TextAt{21}{198.774}{112.753}{i}%
\NativeMathAt{205.476}{114.546}{11.955}{\equiv}%
\TextAt{9}{218.096}{114.546}{1}%
\TextAt{8}{227.27}{114.546}{mod}%
\TextAt{14}{252.707}{114.546}{a}%
\TextAt{21}{258.691}{112.753}{i}%
\TextAt{13}{262.072}{114.546}{\Character{59}}%
\TextAt{13}{291.207}{114.546}{y}%
\TextAt{21}{296.934}{112.753}{i}%
\NativeMathAt{303.636}{114.546}{11.955}{\equiv}%
\TextAt{9}{316.255}{114.546}{0}%
\TextAt{8}{325.429}{114.546}{mod}%
\TextAt{14}{350.867}{114.546}{a}%
\TextAt{21}{356.85}{112.753}{j}%
\TextAt{8}{373.188}{114.546}{for}%
\TextAt{13}{391.073}{114.546}{i}%
\NativeMathAt{398.387}{114.546}{11.955}{\neq}%
\TextAt{13}{410.812}{114.546}{j\Character{58}}%
\TextAt{8}{79.226}{89.967}{Putting}%
\TextAt{13}{123.648}{89.967}{x}%
\TextAt{9}{135.898}{89.967}{\Character{61}}%
\TextAt{13}{150.6}{89.967}{x}%
\TextAt{17}{157.252}{88.174}{1}%
\TextAt{13}{161.984}{89.967}{y}%
\TextAt{17}{167.692}{88.174}{1}%
\TextAt{9}{175.992}{89.967}{\Character{43}}%
\TextAt{16}{188.664}{89.967}{\Character{1}}%
\TextAt{16}{193.96}{89.967}{\Character{1}}%
\TextAt{16}{199.268}{89.967}{\Character{1}}%
\TextAt{9}{206.178}{89.967}{\Character{43}}%
\TextAt{13}{218.85}{89.967}{x}%
\TextAt{21}{225.503}{88.174}{n}%
\TextAt{13}{231.139}{89.967}{y}%
\TextAt{21}{236.847}{88.174}{n}%
\TextAt{8}{247.716}{89.967}{we}%
\TextAt{8}{266.279}{89.967}{find}%
\TextAt{13}{291.024}{89.967}{x}%
\NativeMathAt{303.273}{89.967}{11.955}{\equiv}%
\TextAt{13}{318.169}{89.967}{x}%
\TextAt{21}{324.821}{88.174}{i}%
\TextAt{8}{331.523}{89.967}{mod}%
\TextAt{14}{356.961}{89.967}{a}%
\TextAt{21}{362.944}{88.174}{i}%
\TextAt{8}{366.325}{89.967}{,}%
\TextAt{13}{375.144}{89.967}{i}%
\TextAt{9}{384.735}{89.967}{\Character{61}}%
\TextAt{9}{399.428}{89.967}{1}%
\TextAt{13}{405.29}{89.967}{\Character{59}}%
\TextAt{13}{410.514}{89.967}{\Character{58}}%
\TextAt{13}{415.751}{89.967}{\Character{58}}%
\TextAt{13}{420.987}{89.967}{\Character{58}}%
\TextAt{13}{426.212}{89.967}{\Character{59}}%
\TextAt{13}{431.448}{89.967}{n}%
\TextAt{8}{438.499}{89.967}{.}%
\TextAt{8}{450.955}{89.967}{This}%
\TextAt{8}{479.01}{89.967}{proves}%
\TextAt{8}{516.51}{89.967}{the}%
\TextAt{8}{79.226}{72.921}{surjectivity.}%
\Vector{q 0 G 0 g 0.399 w 452.334 -701.079 m 452.334 -693.009 l S Q}%
\Vector{q 0 G 0 g 0.399 w 452.534 -693.208 m 459.707 -693.208 l S Q}%
\Vector{q 0 G 0 g 0.399 w 452.534 -700.879 m 459.707 -700.879 l S Q}%
\Vector{q 0 G 0 g 0.399 w 459.906 -701.079 m 459.906 -693.009 l S Q}%
\TextAt{8}{300.14}{43.033}{19}%
\Destination{78.23}{809.815}{nav.page.19}%
\Destination{79.36}{429.755}{nav.print.21}%
\Destination{79.226}{490.571}{errata.47}%
\Destination{79.23}{421.7}{nav.number.1.3.6}%
}

% Page 32 (body page 20)
\Page{%
\TextAt{2}{79.226}{796.943}{I.}%
\TextAt{2}{89.83}{796.943}{Algebraic}%
\TextAt{2}{138.648}{796.943}{Integers}%
\Vector{q 0 G 0 g 0.299 w 7.226 18.729 m 460.769 18.729 l S Q}%
\TextAt{8}{97.159}{761.124}{Now}%
\TextAt{8}{123.821}{761.124}{let}%
\TextAt{14}{140.727}{761.124}{o}%
\TextAt{8}{150.47}{761.124}{be}%
\TextAt{8}{166.396}{761.124}{again}%
\TextAt{8}{197.625}{761.124}{a}%
\TextAt{8}{207.382}{761.124}{Dedekind}%
\TextAt{8}{259.568}{761.124}{domain.}%
\TextAt{8}{305.752}{761.124}{Just}%
\TextAt{8}{331.338}{761.124}{as}%
\TextAt{8}{345.71}{761.124}{for}%
\TextAt{8}{363.596}{761.124}{nonzero}%
\TextAt{8}{407.175}{761.124}{numbers,}%
\TextAt{8}{457.64}{761.124}{we}%
\TextAt{8}{474.869}{761.124}{may}%
\TextAt{8}{500.239}{761.124}{obtain}%
\TextAt{8}{79.226}{744.078}{inverses}%
\TextAt{8}{123.536}{744.078}{for}%
\TextAt{8}{142.358}{744.078}{the}%
\TextAt{8}{163.451}{744.078}{nonzero}%
\TextAt{8}{207.964}{744.078}{ideals}%
\TextAt{8}{241.478}{744.078}{of}%
\TextAt{14}{255.745}{744.078}{o}%
\TextAt{8}{266.423}{744.078}{by}%
\TextAt{8}{283.619}{744.078}{introducing}%
\TextAt{8}{346.997}{744.078}{the}%
\TextAt{8}{368.09}{744.078}{notion}%
\TextAt{8}{405.454}{744.078}{of}%
\TextAt{8}{419.721}{744.078}{fractional}%
\TextAt{8}{473.021}{744.078}{ideal}%
\TextAt{8}{501.921}{744.078}{in}%
\TextAt{8}{516.51}{744.078}{the}%
\TextAt{8}{79.226}{727.032}{field}%
\TextAt{8}{104.583}{727.032}{of}%
\TextAt{8}{117.913}{727.032}{fractions}%
\TextAt{13}{165.782}{727.032}{K}%
\TextAt{8}{176.599}{727.032}{.}%
\TextAt{12}{79.226}{701.019}{Definition}%
\TextAt{12}{141.959}{701.019}{(3.7).}%
\TextAt{8}{177.226}{701.019}{A}%
\TextAt{12}{189.85}{701.019}{fractional}%
\TextAt{12}{249.869}{701.019}{ideal}%
\TextAt{8}{281.347}{701.019}{of}%
\TextAt{13}{294.628}{701.019}{K}%
\TextAt{8}{309.295}{701.019}{is}%
\TextAt{8}{321.01}{701.019}{a}%
\TextAt{8}{330.717}{701.019}{finitely}%
\TextAt{8}{370.014}{701.019}{generated}%
\TextAt{14}{423.297}{701.019}{o}%
\TextAt{8}{429.141}{701.019}{-submodule}%
\TextAt{14}{491.905}{701.019}{a}%
\NativeMathAt{501.209}{701.019}{11.955}{\neq}%
\TextAt{9}{513.619}{701.019}{0}%
\TextAt{8}{523.337}{701.019}{of}%
\TextAt{13}{79.226}{683.973}{K}%
\TextAt{8}{90.044}{683.973}{.}%
\TextAt{8}{97.159}{657.961}{For}%
\TextAt{8}{119.366}{657.961}{instance,}%
\TextAt{8}{169.751}{657.961}{an}%
\TextAt{8}{187.248}{657.961}{element}%
\TextAt{13}{231.728}{657.961}{a}%
\NativeMathAt{243.304}{657.961}{11.955}{\in}%
\TextAt{13}{256.705}{657.961}{K}%
\TextAt{3}{267.523}{662.3}{\Character{3}}%
\TextAt{8}{277.391}{657.961}{defines}%
\TextAt{8}{317.053}{657.961}{the}%
\TextAt{8}{338.447}{657.961}{fractional}%
\TextAt{8}{392.049}{657.961}{“principal}%
\TextAt{8}{447.588}{657.961}{ideal”}%
\TextAt{9}{482.312}{657.961}{\Character{40}}%
\TextAt{13}{486.865}{657.961}{a}%
\TextAt{9}{493.01}{657.961}{\Character{41}}%
\TextAt{9}{502.981}{657.961}{\Character{61}}%
\TextAt{13}{517.528}{657.961}{a}%
\TextAt{14}{523.673}{657.961}{o}%
\TextAt{8}{529.518}{657.961}{.}%
\TextAt{8}{79.226}{640.915}{Obviously,}%
\TextAt{8}{137.104}{640.915}{since}%
\TextAt{14}{166.235}{640.915}{o}%
\TextAt{8}{176.438}{640.915}{is}%
\TextAt{8}{188.664}{640.915}{noetherian,}%
\TextAt{8}{250.702}{640.915}{an}%
\TextAt{14}{267.423}{640.915}{o}%
\TextAt{8}{273.267}{640.915}{-submodule}%
\TextAt{14}{336.542}{640.915}{a}%
\NativeMathAt{346.632}{640.915}{11.955}{\neq}%
\TextAt{9}{359.831}{640.915}{0}%
\TextAt{8}{370.057}{640.915}{of}%
\TextAt{13}{383.849}{640.915}{K}%
\TextAt{8}{399.026}{640.915}{is}%
\TextAt{8}{411.251}{640.915}{a}%
\TextAt{8}{421.468}{640.915}{fractional}%
\TextAt{8}{474.294}{640.915}{ideal}%
\TextAt{8}{502.719}{640.915}{if}%
\TextAt{8}{513.904}{640.915}{and}%
\TextAt{8}{79.226}{623.869}{only}%
\TextAt{8}{104.167}{623.869}{if}%
\TextAt{8}{114.141}{623.869}{there}%
\TextAt{8}{143.302}{623.869}{exists}%
\TextAt{13}{174.867}{623.869}{c}%
\NativeMathAt{183.225}{623.869}{11.955}{\in}%
\TextAt{14}{194.516}{623.869}{o}%
\TextAt{8}{200.361}{623.869}{,}%
\TextAt{13}{206.91}{623.869}{c}%
\NativeMathAt{215.268}{623.869}{11.955}{\neq}%
\TextAt{9}{227.678}{623.869}{0}%
\TextAt{8}{233.547}{623.869}{,}%
\TextAt{8}{240.096}{623.869}{such}%
\TextAt{8}{265.742}{623.869}{that}%
\TextAt{13}{290.361}{623.869}{c}%
\TextAt{14}{295.399}{623.869}{a}%
\NativeMathAt{304.703}{623.869}{11.955}{\subseteq}%
\TextAt{14}{317.322}{623.869}{o}%
\TextAt{8}{326.313}{623.869}{is}%
\TextAt{8}{337.327}{623.869}{an}%
\TextAt{8}{352.835}{623.869}{ideal}%
\TextAt{8}{380.048}{623.869}{of}%
\TextAt{8}{392.628}{623.869}{the}%
\TextAt{8}{412.034}{623.869}{ring}%
\TextAt{14}{435.349}{623.869}{o}%
\TextAt{8}{441.193}{623.869}{.}%
\TextAt{8}{449.391}{623.869}{Fractional}%
\TextAt{8}{504.089}{623.869}{ideals}%
\TextAt{8}{79.226}{606.823}{are}%
\TextAt{8}{99.13}{606.823}{multiplied}%
\TextAt{8}{155.127}{606.823}{in}%
\TextAt{8}{169.174}{606.823}{the}%
\TextAt{8}{189.724}{606.823}{same}%
\TextAt{8}{219.443}{606.823}{way}%
\TextAt{8}{243.579}{606.823}{as}%
\TextAt{8}{258.343}{606.823}{ideals}%
\TextAt{8}{291.314}{606.823}{in}%
\TextAt{14}{305.361}{606.823}{o}%
\TextAt{8}{311.205}{606.823}{.}%
\TextAt{8}{320.834}{606.823}{For}%
\TextAt{8}{342.196}{606.823}{distinction}%
\TextAt{8}{400.537}{606.823}{the}%
\TextAt{8}{421.086}{606.823}{latter}%
\TextAt{8}{453.352}{606.823}{may}%
\TextAt{8}{479.115}{606.823}{henceforth}%
\TextAt{8}{79.226}{589.777}{be}%
\TextAt{8}{95.15}{589.777}{called}%
\TextAt{12}{128.314}{589.777}{integral}%
\TextAt{12}{177.522}{589.777}{ideals}%
\TextAt{8}{214.355}{589.777}{of}%
\TextAt{13}{227.685}{589.777}{K}%
\TextAt{8}{238.503}{589.777}{.}%
\TextAt{12}{79.226}{563.765}{Proposition}%
\TextAt{12}{151.716}{563.765}{(3.8).}%
\TextAt{8}{186.838}{563.765}{The}%
\TextAt{8}{210.699}{563.765}{fractional}%
\TextAt{8}{262.87}{563.765}{ideals}%
\TextAt{8}{295.255}{563.765}{form}%
\TextAt{8}{322.703}{563.765}{an}%
\TextAt{8}{338.769}{563.765}{abelian}%
\TextAt{8}{379.224}{563.765}{group,}%
\TextAt{8}{415.497}{563.765}{the}%
\TextAt{12}{435.46}{563.765}{ideal}%
\TextAt{12}{467.35}{563.765}{group}%
\TextAt{13}{504.959}{563.765}{J}%
\TextAt{21}{511.43}{561.972}{K}%
\TextAt{8}{523.337}{563.765}{of}%
\TextAt{13}{79.226}{546.719}{K}%
\TextAt{8}{90.044}{546.719}{.}%
\TextAt{8}{98.492}{546.719}{The}%
\TextAt{8}{122.546}{546.719}{identity}%
\TextAt{8}{165.8}{546.719}{element}%
\TextAt{8}{209.042}{546.719}{is}%
\TextAt{9}{220.805}{546.719}{\Character{40}1\Character{41}}%
\TextAt{9}{239.049}{546.719}{\Character{61}}%
\TextAt{14}{251.509}{546.719}{o}%
\TextAt{8}{257.354}{546.719}{,}%
\TextAt{8}{264.503}{546.719}{and}%
\TextAt{8}{287.266}{546.719}{the}%
\TextAt{8}{307.422}{546.719}{inverse}%
\TextAt{8}{346.181}{546.719}{of}%
\TextAt{14}{359.511}{546.719}{a}%
\TextAt{8}{369.391}{546.719}{is}%
\TextAt{14}{243.638}{517.718}{a}%
\TextAt{3}{249.622}{522.654}{\Character{0}}%
\TextAt{17}{256.208}{522.654}{1}%
\TextAt{9}{264.261}{517.718}{\Character{61}}%
\TextAt{16}{276.687}{517.718}{\Character{102}}%
\TextAt{13}{282.664}{517.718}{x}%
\NativeMathAt{292.637}{517.718}{11.955}{\in}%
\TextAt{13}{303.928}{517.718}{K}%
\TextAt{16}{318.067}{517.718}{\Character{106}}%
\TextAt{13}{324.708}{517.718}{x}%
\TextAt{14}{331.36}{517.718}{a}%
\NativeMathAt{340.664}{517.718}{11.955}{\subseteq}%
\TextAt{14}{353.284}{517.718}{o}%
\TextAt{16}{359.128}{517.718}{\Character{103}}%
\TextAt{13}{365.106}{517.718}{\Character{58}}%
\TextAt{19}{79.226}{470.784}{Proof.}%
\TextAt{8}{116}{470.784}{One}%
\TextAt{8}{140.349}{470.784}{obviously}%
\TextAt{8}{191.764}{470.784}{has}%
\TextAt{8}{212.276}{470.784}{associativity,}%
\TextAt{8}{281.385}{470.784}{commutativity}%
\TextAt{8}{359.736}{470.784}{and}%
\TextAt{14}{382.136}{470.784}{a}%
\TextAt{9}{388.119}{470.784}{\Character{40}1\Character{41}}%
\TextAt{9}{406.363}{470.784}{\Character{61}}%
\TextAt{14}{418.823}{470.784}{a}%
\TextAt{8}{424.806}{470.784}{.}%
\TextAt{8}{433.133}{470.784}{For}%
\TextAt{8}{453.74}{470.784}{a}%
\TextAt{8}{463.133}{470.784}{prime}%
\TextAt{8}{495.934}{470.784}{ideal}%
\TextAt{14}{523.535}{470.784}{p}%
\TextAt{8}{529.518}{470.784}{,}%
\TextAt{8}{79.226}{453.738}{(3.5)}%
\TextAt{8}{107.516}{453.738}{says}%
\TextAt{8}{132.674}{453.738}{that}%
\TextAt{14}{158.357}{453.738}{p}%
\TextAt{18}{168.198}{453.738}{\Character{40}}%
\TextAt{14}{181.354}{453.738}{pp}%
\TextAt{3}{193.32}{458.076}{\Character{0}}%
\TextAt{17}{199.906}{458.076}{1}%
\TextAt{8}{208.851}{453.738}{and}%
\TextAt{8}{231.928}{453.738}{therefore}%
\TextAt{14}{281.343}{453.738}{pp}%
\TextAt{3}{293.309}{458.076}{\Character{0}}%
\TextAt{17}{299.896}{458.076}{1}%
\TextAt{9}{308.485}{453.738}{\Character{61}}%
\TextAt{14}{321.447}{453.738}{o}%
\TextAt{8}{331.504}{453.738}{because}%
\TextAt{14}{375.12}{453.738}{p}%
\TextAt{8}{385.316}{453.738}{is}%
\TextAt{8}{397.394}{453.738}{maximal.}%
\TextAt{8}{450.698}{453.738}{Consequently,}%
\TextAt{8}{525.943}{453.738}{if}%
\TextAt{14}{79.226}{436.692}{a}%
\TextAt{9}{89.35}{436.692}{\Character{61}}%
\TextAt{14}{102.595}{436.692}{p}%
\TextAt{17}{108.578}{434.899}{1}%
\TextAt{16}{115.303}{436.692}{\Character{1}}%
\TextAt{16}{120.599}{436.692}{\Character{1}}%
\TextAt{16}{125.907}{436.692}{\Character{1}}%
\TextAt{14}{131.243}{436.692}{p}%
\TextAt{21}{137.226}{434.899}{r}%
\TextAt{8}{146.159}{436.692}{is}%
\TextAt{8}{158.404}{436.692}{an}%
\TextAt{8}{175.144}{436.692}{integral}%
\TextAt{8}{218.233}{436.692}{ideal,}%
\TextAt{8}{250.05}{436.692}{then}%
\TextAt{14}{277.191}{436.692}{b}%
\TextAt{9}{287.466}{436.692}{\Character{61}}%
\TextAt{14}{300.711}{436.692}{p}%
\TextAt{3}{306.694}{441.64}{\Character{0}}%
\TextAt{17}{313.281}{441.64}{1}%
\TextAt{17}{306.694}{433.761}{1}%
\TextAt{16}{320.006}{436.692}{\Character{1}}%
\TextAt{16}{325.302}{436.692}{\Character{1}}%
\TextAt{16}{330.61}{436.692}{\Character{1}}%
\TextAt{14}{335.946}{436.692}{p}%
\TextAt{3}{341.929}{441.03}{\Character{0}}%
\TextAt{17}{348.515}{441.03}{1}%
\TextAt{21}{341.929}{433.736}{r}%
\TextAt{8}{357.626}{436.692}{is}%
\TextAt{8}{369.871}{436.692}{an}%
\TextAt{8}{386.611}{436.692}{inverse.}%
\TextAt{14}{431.364}{436.692}{ba}%
\TextAt{9}{447.622}{436.692}{\Character{61}}%
\TextAt{14}{460.867}{436.692}{o}%
\TextAt{8}{471.09}{436.692}{implies}%
\TextAt{8}{511.298}{436.692}{that}%
\TextAt{14}{79.226}{419.646}{b}%
\NativeMathAt{89.496}{419.646}{11.955}{\subseteq}%
\TextAt{14}{102.93}{419.646}{a}%
\TextAt{3}{108.913}{423.984}{\Character{0}}%
\TextAt{17}{115.5}{423.984}{1}%
\TextAt{8}{120.232}{419.646}{.}%
\TextAt{8}{130.114}{419.646}{Conversely,}%
\TextAt{8}{192.233}{419.646}{if}%
\TextAt{13}{203.435}{419.646}{x}%
\TextAt{14}{210.087}{419.646}{a}%
\NativeMathAt{220.206}{419.646}{11.955}{\subseteq}%
\TextAt{14}{233.639}{419.646}{o}%
\TextAt{8}{239.484}{419.646}{,}%
\TextAt{8}{247.23}{419.646}{then}%
\TextAt{13}{274.368}{419.646}{x}%
\TextAt{14}{281.02}{419.646}{ab}%
\NativeMathAt{297.274}{419.646}{11.955}{\subseteq}%
\TextAt{14}{310.708}{419.646}{b}%
\TextAt{8}{316.843}{419.646}{,}%
\TextAt{8}{324.589}{419.646}{so}%
\TextAt{13}{339.437}{419.646}{x}%
\NativeMathAt{350.225}{419.646}{11.955}{\in}%
\TextAt{14}{362.33}{419.646}{b}%
\TextAt{8}{372.841}{419.646}{because}%
\TextAt{14}{416.62}{419.646}{ab}%
\TextAt{9}{432.873}{419.646}{\Character{61}}%
\TextAt{14}{446.113}{419.646}{o}%
\TextAt{8}{451.958}{419.646}{.}%
\TextAt{8}{461.84}{419.646}{Thus}%
\TextAt{8}{491.967}{419.646}{we}%
\TextAt{8}{509.672}{419.646}{have}%
\TextAt{14}{79.226}{402.6}{b}%
\TextAt{9}{88.682}{402.6}{\Character{61}}%
\TextAt{14}{101.107}{402.6}{a}%
\TextAt{3}{107.09}{406.938}{\Character{0}}%
\TextAt{17}{113.677}{406.938}{1}%
\TextAt{8}{118.409}{402.6}{.}%
\TextAt{8}{126.768}{402.6}{Finally,}%
\TextAt{8}{168.659}{402.6}{if}%
\TextAt{14}{179.113}{402.6}{a}%
\TextAt{8}{188.724}{402.6}{is}%
\TextAt{8}{200.219}{402.6}{an}%
\TextAt{8}{216.209}{402.6}{arbitrary}%
\TextAt{8}{265.709}{402.6}{fractional}%
\TextAt{8}{317.804}{402.6}{ideal}%
\TextAt{8}{345.498}{402.6}{and}%
\TextAt{13}{367.991}{402.6}{c}%
\NativeMathAt{376.35}{402.6}{11.955}{\in}%
\TextAt{14}{387.641}{402.6}{o}%
\TextAt{8}{393.485}{402.6}{,}%
\TextAt{13}{400.419}{402.6}{c}%
\NativeMathAt{408.778}{402.6}{11.955}{\neq}%
\TextAt{9}{421.188}{402.6}{0}%
\TextAt{8}{427.056}{402.6}{,}%
\TextAt{8}{433.99}{402.6}{is}%
\TextAt{8}{445.485}{402.6}{such}%
\TextAt{8}{471.613}{402.6}{that}%
\TextAt{13}{496.712}{402.6}{c}%
\TextAt{14}{501.75}{402.6}{a}%
\NativeMathAt{511.054}{402.6}{11.955}{\subseteq}%
\TextAt{14}{523.674}{402.6}{o}%
\TextAt{8}{529.518}{402.6}{,}%
\TextAt{8}{79.226}{385.554}{then}%
\TextAt{9}{105.886}{385.554}{\Character{40}}%
\TextAt{13}{110.439}{385.554}{c}%
\TextAt{14}{115.477}{385.554}{a}%
\TextAt{9}{121.46}{385.554}{\Character{41}}%
\TextAt{3}{126.012}{389.892}{\Character{0}}%
\TextAt{17}{132.598}{389.892}{1}%
\TextAt{9}{140.652}{385.554}{\Character{61}}%
\TextAt{13}{153.077}{385.554}{c}%
\TextAt{3}{158.115}{389.892}{\Character{0}}%
\TextAt{17}{164.702}{389.892}{1}%
\TextAt{14}{169.434}{385.554}{a}%
\TextAt{3}{175.417}{389.892}{\Character{0}}%
\TextAt{17}{182.003}{389.892}{1}%
\TextAt{8}{190.633}{385.554}{is}%
\TextAt{8}{202.397}{385.554}{the}%
\TextAt{8}{222.553}{385.554}{inverse}%
\TextAt{8}{261.312}{385.554}{of}%
\TextAt{13}{274.642}{385.554}{c}%
\TextAt{14}{279.68}{385.554}{a}%
\TextAt{8}{285.663}{385.554}{,}%
\TextAt{8}{292.812}{385.554}{so}%
\TextAt{14}{307.183}{385.554}{aa}%
\TextAt{3}{319.149}{389.892}{\Character{0}}%
\TextAt{17}{325.735}{389.892}{1}%
\TextAt{9}{333.788}{385.554}{\Character{61}}%
\TextAt{14}{346.214}{385.554}{o}%
\TextAt{8}{352.058}{385.554}{.}%
\Vector{q 0 G 0 g 0.399 w 452.334 -388.446 m 452.334 -380.376 l S Q}%
\Vector{q 0 G 0 g 0.399 w 452.534 -380.575 m 459.707 -380.575 l S Q}%
\Vector{q 0 G 0 g 0.399 w 452.534 -388.247 m 459.707 -388.247 l S Q}%
\Vector{q 0 G 0 g 0.399 w 459.906 -388.446 m 459.906 -380.376 l S Q}%
\TextAt{12}{79.226}{350.575}{Corollary}%
\TextAt{12}{139.014}{350.575}{(3.9).}%
\TextAt{8}{174.33}{350.575}{Every}%
\TextAt{8}{207.983}{350.575}{fractional}%
\TextAt{8}{260.347}{350.575}{ideal}%
\TextAt{14}{288.31}{350.575}{a}%
\TextAt{8}{298.191}{350.575}{admits}%
\TextAt{8}{336.626}{350.575}{a}%
\TextAt{8}{346.382}{350.575}{unique}%
\TextAt{8}{384.423}{350.575}{representation}%
\TextAt{8}{460.912}{350.575}{as}%
\TextAt{8}{475.282}{350.575}{a}%
\TextAt{8}{485.038}{350.575}{product}%
\TextAt{14}{279.136}{321.574}{a}%
\TextAt{9}{288.439}{321.574}{\Character{61}}%
\TextAt{20}{300.865}{332.932}{\Character{89}}%
\TextAt{27}{306.509}{307.826}{p}%
\TextAt{14}{318.133}{321.574}{p}%
\TextAt{21}{324.117}{326.51}{\Character{23}}%
\TextAt{28}{328.279}{325.514}{p}%
\TextAt{8}{79.226}{283.755}{with}%
\TextAt{13}{106.045}{283.755}{\Character{23}}%
\TextAt{27}{111.829}{281.962}{p}%
\NativeMathAt{119.907}{283.755}{11.955}{\in}%
\TextAt{18}{131.468}{283.755}{Z}%
\TextAt{8}{143.494}{283.755}{and}%
\TextAt{13}{166.415}{283.755}{\Character{23}}%
\TextAt{27}{172.199}{281.962}{p}%
\TextAt{9}{180.277}{283.755}{\Character{61}}%
\TextAt{9}{192.962}{283.755}{0}%
\TextAt{8}{202.882}{283.755}{for}%
\TextAt{8}{220.925}{283.755}{almost}%
\TextAt{8}{258.874}{283.755}{all}%
\TextAt{14}{275.292}{283.755}{p}%
\TextAt{8}{281.275}{283.755}{.}%
\TextAt{8}{290.199}{283.755}{In}%
\TextAt{8}{304.979}{283.755}{other}%
\TextAt{8}{335.707}{283.755}{words,}%
\TextAt{13}{372.715}{283.755}{J}%
\TextAt{21}{379.185}{281.962}{K}%
\TextAt{8}{391.444}{283.755}{is}%
\TextAt{8}{403.366}{283.755}{the}%
\TextAt{8}{423.681}{283.755}{free}%
\TextAt{8}{446.268}{283.755}{abelian}%
\TextAt{8}{487.074}{283.755}{group}%
\TextAt{8}{520.408}{283.755}{on}%
\TextAt{8}{79.226}{266.709}{the}%
\TextAt{8}{99.383}{266.709}{set}%
\TextAt{8}{117.65}{266.709}{of}%
\TextAt{8}{130.98}{266.709}{nonzero}%
\TextAt{8}{174.557}{266.709}{prime}%
\TextAt{8}{207.72}{266.709}{ideals}%
\TextAt{14}{240.298}{266.709}{p}%
\TextAt{8}{250.179}{266.709}{of}%
\TextAt{14}{263.509}{266.709}{o}%
\TextAt{8}{269.353}{266.709}{.}%
\TextAt{19}{79.226}{231.731}{Proof.}%
\TextAt{8}{116}{231.731}{Every}%
\TextAt{8}{149.527}{231.731}{fractional}%
\TextAt{8}{201.763}{231.731}{ideal}%
\TextAt{14}{229.598}{231.731}{a}%
\TextAt{8}{239.351}{231.731}{is}%
\TextAt{8}{250.988}{231.731}{a}%
\TextAt{8}{260.616}{231.731}{quotient}%
\TextAt{14}{306.671}{231.731}{a}%
\TextAt{9}{315.975}{231.731}{\Character{61}}%
\TextAt{14}{328.4}{231.731}{b}%
\TextAt{9}{334.535}{231.731}{\Character{47}}%
\TextAt{14}{340.388}{231.731}{c}%
\TextAt{8}{348.814}{231.731}{of}%
\TextAt{8}{362.017}{231.731}{two}%
\TextAt{8}{384.006}{231.731}{integral}%
\TextAt{8}{426.487}{231.731}{ideals}%
\TextAt{14}{458.937}{231.731}{b}%
\TextAt{8}{468.842}{231.731}{and}%
\TextAt{14}{491.477}{231.731}{c}%
\TextAt{8}{496.133}{231.731}{,}%
\TextAt{8}{503.18}{231.731}{which}%
\TextAt{8}{79.226}{214.685}{by}%
\TextAt{8}{95.334}{214.685}{(3.3)}%
\TextAt{8}{123.157}{214.685}{have}%
\TextAt{8}{150}{214.685}{a}%
\TextAt{8}{159.604}{214.685}{prime}%
\TextAt{8}{192.616}{214.685}{decomposition.}%
\TextAt{8}{274.251}{214.685}{Therefore}%
\TextAt{14}{327.097}{214.685}{a}%
\TextAt{8}{336.825}{214.685}{has}%
\TextAt{8}{357.548}{214.685}{a}%
\TextAt{8}{367.151}{214.685}{prime}%
\TextAt{8}{400.163}{214.685}{decomposition}%
\TextAt{8}{477.146}{214.685}{of}%
\TextAt{8}{490.325}{214.685}{the}%
\TextAt{8}{510.33}{214.685}{type}%
\TextAt{8}{79.226}{197.639}{stated}%
\TextAt{8}{114.446}{197.639}{in}%
\TextAt{8}{128.134}{197.639}{the}%
\TextAt{8}{148.326}{197.639}{corollary.}%
\TextAt{8}{200.469}{197.639}{By}%
\TextAt{8}{218.868}{197.639}{(3.3),}%
\TextAt{8}{250.139}{197.639}{it}%
\TextAt{8}{261.879}{197.639}{is}%
\TextAt{8}{273.678}{197.639}{unique}%
\TextAt{8}{311.755}{197.639}{if}%
\TextAt{14}{322.514}{197.639}{a}%
\TextAt{8}{332.43}{197.639}{is}%
\TextAt{8}{344.229}{197.639}{integral,}%
\TextAt{8}{390.134}{197.639}{and}%
\TextAt{8}{412.932}{197.639}{therefore}%
\TextAt{8}{462.067}{197.639}{clearly}%
\TextAt{8}{499.499}{197.639}{also}%
\TextAt{8}{523.014}{197.639}{in}%
\TextAt{8}{79.226}{180.593}{general.}%
\Vector{q 0 G 0 g 0.399 w 452.334 -593.407 m 452.334 -585.337 l S Q}%
\Vector{q 0 G 0 g 0.399 w 452.534 -585.537 m 459.707 -585.537 l S Q}%
\Vector{q 0 G 0 g 0.399 w 452.534 -593.208 m 459.707 -593.208 l S Q}%
\Vector{q 0 G 0 g 0.399 w 459.906 -593.407 m 459.906 -585.337 l S Q}%
\TextAt{8}{97.159}{151.592}{The}%
\TextAt{8}{122.246}{151.592}{fractional}%
\TextAt{8}{175.643}{151.592}{principal}%
\TextAt{8}{225.453}{151.592}{ideals}%
\TextAt{9}{259.065}{151.592}{\Character{40}}%
\TextAt{13}{263.617}{151.592}{a}%
\TextAt{9}{269.762}{151.592}{\Character{41}}%
\TextAt{9}{279.386}{151.592}{\Character{61}}%
\TextAt{13}{293.582}{151.592}{a}%
\TextAt{14}{299.727}{151.592}{o}%
\TextAt{8}{305.571}{151.592}{,}%
\TextAt{13}{314.012}{151.592}{a}%
\NativeMathAt{325.239}{151.592}{11.955}{\in}%
\TextAt{13}{338.291}{151.592}{K}%
\TextAt{3}{349.108}{155.93}{\Character{3}}%
\TextAt{8}{353.841}{151.592}{,}%
\TextAt{8}{362.282}{151.592}{form}%
\TextAt{8}{390.955}{151.592}{a}%
\TextAt{8}{401.744}{151.592}{subgroup}%
\TextAt{8}{453.575}{151.592}{of}%
\TextAt{8}{467.938}{151.592}{the}%
\TextAt{8}{489.128}{151.592}{group}%
\TextAt{8}{523.337}{151.592}{of}%
\TextAt{8}{79.226}{134.546}{ideals}%
\TextAt{13}{111.804}{134.546}{J}%
\TextAt{21}{118.274}{132.752}{K}%
\TextAt{8}{126.477}{134.546}{,}%
\TextAt{8}{133.626}{134.546}{which}%
\TextAt{8}{167.112}{134.546}{will}%
\TextAt{8}{189.218}{134.546}{be}%
\TextAt{8}{205.142}{134.546}{denoted}%
\TextAt{13}{249.364}{134.546}{P}%
\TextAt{21}{256.909}{132.752}{K}%
\TextAt{8}{265.112}{134.546}{.}%
\TextAt{8}{273.56}{134.546}{The}%
\TextAt{8}{297.614}{134.546}{quotient}%
\TextAt{8}{343.797}{134.546}{group}%
\TextAt{13}{269.526}{105.545}{C}%
\TextAt{13}{278.744}{105.545}{l}%
\TextAt{21}{282.247}{103.751}{K}%
\TextAt{9}{293.77}{105.545}{\Character{61}}%
\TextAt{13}{306.196}{105.545}{J}%
\TextAt{21}{312.666}{103.751}{K}%
\TextAt{9}{320.869}{105.545}{\Character{47}}%
\TextAt{13}{326.721}{105.545}{P}%
\TextAt{21}{334.267}{103.751}{K}%
\TextAt{8}{300.14}{43.033}{20}%
\Destination{78.23}{809.815}{nav.page.20}%
\Destination{79.36}{485.755}{nav.print.22}%
\Destination{79.36}{715.055}{nav.number.1.3.7}%
\Destination{79.36}{577.755}{nav.number.1.3.8}%
\Destination{79.36}{364.555}{nav.number.1.3.9}%
\Hotlink{78.726}{451.555}{25.078}{11.047}{nav.number.1.3.5}%
\Hotlink{94.834}{212.455}{25.078}{11.047}{nav.number.1.3.3}%
\Hotlink{218.368}{195.455}{25.078}{11.047}{nav.number.1.3.3}%
}

% Page 33 (body page 21)
\Page{%
\TextAt{2}{79.226}{796.943}{I.}%
\TextAt{2}{89.83}{796.943}{Algebraic}%
\TextAt{2}{138.648}{796.943}{Integers}%
\Vector{q 0 G 0 g 0.299 w 7.226 18.729 m 460.769 18.729 l S Q}%
\TextAt{8}{79.226}{761.124}{is}%
\TextAt{8}{91.334}{761.124}{called}%
\TextAt{8}{124.841}{761.124}{the}%
\TextAt{12}{145.341}{761.124}{ideal}%
\TextAt{12}{177.847}{761.124}{class}%
\TextAt{12}{209.601}{761.124}{group}%
\TextAt{8}{243.506}{761.124}{,}%
\TextAt{8}{251.084}{761.124}{or}%
\TextAt{12}{265.738}{761.124}{class}%
\TextAt{12}{297.492}{761.124}{group}%
\TextAt{8}{335.637}{761.124}{for}%
\TextAt{8}{353.866}{761.124}{short,}%
\TextAt{8}{387.53}{761.124}{of}%
\TextAt{13}{401.204}{761.124}{K}%
\TextAt{8}{412.022}{761.124}{.}%
\TextAt{8}{421.5}{761.124}{Along}%
\TextAt{8}{455.988}{761.124}{with}%
\TextAt{8}{482.991}{761.124}{the}%
\TextAt{8}{503.491}{761.124}{group}%
\TextAt{8}{79.226}{744.078}{of}%
\TextAt{8}{92.556}{744.078}{units}%
\TextAt{14}{121.882}{744.078}{o}%
\TextAt{3}{127.726}{748.416}{\Character{3}}%
\TextAt{8}{136.356}{744.078}{of}%
\TextAt{14}{149.686}{744.078}{o}%
\TextAt{8}{155.53}{744.078}{,}%
\TextAt{8}{162.68}{744.078}{it}%
\TextAt{8}{174.384}{744.078}{fits}%
\TextAt{8}{193.954}{744.078}{into}%
\TextAt{8}{217.697}{744.078}{the}%
\TextAt{8}{237.854}{744.078}{exact}%
\TextAt{8}{268.746}{744.078}{sequence}%
\TextAt{9}{202.901}{715.177}{1}%
\NativeArrowAt{212.074}{715.177}{11.955}{19.26}{\longrightarrow}%
\TextAt{14}{234.656}{715.177}{o}%
\TextAt{3}{240.501}{720.113}{\Character{3}}%
\NativeArrowAt{248.554}{715.177}{11.955}{19.26}{\longrightarrow}%
\TextAt{13}{271.136}{715.177}{K}%
\TextAt{3}{281.953}{720.113}{\Character{3}}%
\NativeArrowAt{290.007}{715.177}{11.955}{19.26}{\longrightarrow}%
\TextAt{13}{312.589}{715.177}{J}%
\TextAt{21}{319.059}{713.384}{K}%
\NativeArrowAt{330.582}{715.177}{11.955}{19.26}{\longrightarrow}%
\TextAt{13}{353.164}{715.177}{C}%
\TextAt{13}{362.382}{715.177}{l}%
\TextAt{21}{365.885}{713.384}{K}%
\NativeArrowAt{377.408}{715.177}{11.955}{19.26}{\longrightarrow}%
\TextAt{9}{399.99}{715.177}{1}%
\TextAt{13}{405.843}{715.177}{\Character{59}}%
\TextAt{8}{79.226}{686.277}{where}%
\TextAt{8}{112.984}{686.277}{the}%
\TextAt{8}{133.089}{686.277}{arrow}%
\TextAt{8}{165.89}{686.277}{in}%
\TextAt{8}{179.492}{686.277}{the}%
\TextAt{8}{199.597}{686.277}{middle}%
\TextAt{8}{237.91}{686.277}{is}%
\TextAt{8}{249.622}{686.277}{given}%
\TextAt{8}{280.14}{686.277}{by}%
\TextAt{13}{296.348}{686.277}{a}%
\NativeArrowAt{305.813}{686.277}{11.955}{11.955}{\mapsto}%
\TextAt{9}{321.09}{686.277}{\Character{40}}%
\TextAt{13}{325.642}{686.277}{a}%
\TextAt{9}{331.787}{686.277}{\Character{41}}%
\TextAt{8}{336.339}{686.277}{.}%
\TextAt{8}{344.77}{686.277}{So}%
\TextAt{8}{360.978}{686.277}{the}%
\TextAt{8}{381.083}{686.277}{class}%
\TextAt{8}{408.469}{686.277}{group}%
\TextAt{13}{441.593}{686.277}{C}%
\TextAt{13}{450.81}{686.277}{l}%
\TextAt{21}{454.313}{684.483}{K}%
\TextAt{8}{466.362}{686.277}{measures}%
\TextAt{8}{516.51}{686.277}{the}%
\TextAt{8}{79.226}{669.231}{expansion}%
\TextAt{8}{133.396}{669.231}{that}%
\TextAt{8}{158.563}{669.231}{takes}%
\TextAt{8}{188.345}{669.231}{place}%
\TextAt{8}{218.055}{669.231}{when}%
\TextAt{8}{248.411}{669.231}{we}%
\TextAt{8}{265.436}{669.231}{pass}%
\TextAt{8}{290.723}{669.231}{from}%
\TextAt{8}{318.161}{669.231}{numbers}%
\TextAt{8}{365.17}{669.231}{to}%
\TextAt{8}{379.279}{669.231}{ideals,}%
\TextAt{8}{414.947}{669.231}{whereas}%
\TextAt{8}{459.027}{669.231}{the}%
\TextAt{8}{478.982}{669.231}{unit}%
\TextAt{8}{503.491}{669.231}{group}%
\TextAt{14}{79.226}{652.185}{o}%
\TextAt{3}{85.07}{656.523}{\Character{3}}%
\TextAt{8}{93.562}{652.185}{measures}%
\TextAt{8}{143.624}{652.185}{the}%
\TextAt{8}{163.642}{652.185}{contraction}%
\TextAt{8}{224.978}{652.185}{in}%
\TextAt{8}{238.493}{652.185}{the}%
\TextAt{8}{258.511}{652.185}{same}%
\TextAt{8}{287.699}{652.185}{process.}%
\TextAt{8}{332.971}{652.185}{This}%
\TextAt{8}{359.553}{652.185}{immediately}%
\TextAt{8}{426.077}{652.185}{raises}%
\TextAt{8}{457.931}{652.185}{the}%
\TextAt{8}{477.95}{652.185}{problem}%
\TextAt{8}{523.337}{652.185}{of}%
\TextAt{8}{79.226}{635.139}{understanding}%
\TextAt{8}{155.915}{635.139}{these}%
\TextAt{8}{185.763}{635.139}{groups}%
\TextAt{14}{223.43}{635.139}{o}%
\TextAt{3}{229.275}{639.477}{\Character{3}}%
\TextAt{8}{237.781}{635.139}{and}%
\TextAt{13}{260.421}{635.139}{C}%
\TextAt{13}{269.639}{635.139}{l}%
\TextAt{21}{273.141}{633.345}{K}%
\TextAt{8}{285.118}{635.139}{more}%
\TextAt{8}{314.261}{635.139}{thoroughly.}%
\TextAt{8}{377.315}{635.139}{For}%
\TextAt{8}{398.162}{635.139}{general}%
\TextAt{8}{438.363}{635.139}{Dedekind}%
\TextAt{8}{490.424}{635.139}{domains}%
\TextAt{8}{79.226}{618.093}{they}%
\TextAt{8}{106.37}{618.093}{may}%
\TextAt{8}{132.545}{618.093}{turn}%
\TextAt{8}{159.366}{618.093}{out}%
\TextAt{8}{180.987}{618.093}{to}%
\TextAt{8}{196.104}{618.093}{be}%
\TextAt{8}{212.835}{618.093}{completely}%
\TextAt{8}{272.496}{618.093}{arbitrary}%
\TextAt{8}{323.072}{618.093}{groups.}%
\TextAt{8}{367.833}{618.093}{For}%
\TextAt{8}{389.609}{618.093}{the}%
\TextAt{8}{410.572}{618.093}{ring}%
\TextAt{14}{435.445}{618.093}{o}%
\TextAt{21}{441.289}{616.299}{K}%
\TextAt{8}{454.195}{618.093}{of}%
\TextAt{8}{468.332}{618.093}{integers}%
\TextAt{8}{512.452}{618.093}{in}%
\TextAt{8}{526.911}{618.093}{a}%
\TextAt{8}{79.226}{601.047}{number}%
\TextAt{8}{122.812}{601.047}{field}%
\TextAt{13}{149.158}{601.047}{K}%
\TextAt{8}{159.975}{601.047}{,}%
\TextAt{8}{168.361}{601.047}{however,}%
\TextAt{8}{217.729}{601.047}{one}%
\TextAt{8}{240.177}{601.047}{obtains}%
\TextAt{8}{282.209}{601.047}{important}%
\TextAt{8}{338.49}{601.047}{finiteness}%
\TextAt{8}{390.325}{601.047}{theorems,}%
\TextAt{8}{444.953}{601.047}{which}%
\TextAt{8}{479.429}{601.047}{are}%
\TextAt{8}{499.929}{601.047}{funda-}%
\TextAt{8}{79.226}{584.001}{mental}%
\TextAt{8}{118.57}{584.001}{for}%
\TextAt{8}{137.1}{584.001}{the}%
\TextAt{8}{157.902}{584.001}{further}%
\TextAt{8}{197.891}{584.001}{development}%
\TextAt{8}{266.501}{584.001}{of}%
\TextAt{8}{280.476}{584.001}{number}%
\TextAt{8}{323.718}{584.001}{theory.}%
\TextAt{8}{365.974}{584.001}{But}%
\TextAt{8}{389.859}{584.001}{these}%
\TextAt{8}{420.476}{584.001}{results}%
\TextAt{8}{458.314}{584.001}{cannot}%
\TextAt{8}{497.335}{584.001}{be}%
\TextAt{8}{513.904}{584.001}{had}%
\TextAt{8}{79.226}{566.955}{for}%
\TextAt{8}{96.483}{566.955}{nothing.}%
\TextAt{8}{143.756}{566.955}{They}%
\TextAt{8}{173.363}{566.955}{will}%
\TextAt{8}{194.84}{566.955}{be}%
\TextAt{8}{210.137}{566.955}{obtained}%
\TextAt{8}{257.64}{566.955}{by}%
\TextAt{8}{273.272}{566.955}{viewing}%
\TextAt{8}{315.24}{566.955}{the}%
\TextAt{8}{334.769}{566.955}{numbers}%
\TextAt{8}{381.352}{566.955}{geometrically}%
\TextAt{8}{452.598}{566.955}{as}%
\TextAt{8}{466.341}{566.955}{lattice}%
\TextAt{8}{501.483}{566.955}{points}%
\TextAt{8}{79.226}{549.909}{in}%
\TextAt{8}{92.567}{549.909}{space.}%
\TextAt{8}{128.289}{549.909}{For}%
\TextAt{8}{148.946}{549.909}{this}%
\TextAt{8}{171.457}{549.909}{we}%
\TextAt{8}{188.372}{549.909}{will}%
\TextAt{8}{210.166}{549.909}{now}%
\TextAt{8}{234.242}{549.909}{prepare}%
\TextAt{8}{276.204}{549.909}{the}%
\TextAt{8}{296.049}{549.909}{necessary}%
\TextAt{8}{347.562}{549.909}{concepts,}%
\TextAt{8}{398.099}{549.909}{which}%
\TextAt{8}{431.273}{549.909}{all}%
\TextAt{8}{447.22}{549.909}{come}%
\TextAt{8}{476.82}{549.909}{from}%
\TextAt{8}{504.149}{549.909}{linear}%
\TextAt{8}{79.226}{532.863}{algebra.}%
\TextAt{12}{79.226}{506.879}{Exercise}%
\TextAt{12}{132.427}{506.879}{I.3.1.}%
\TextAt{8}{166.021}{506.879}{Decompose}%
\TextAt{9}{227.375}{506.879}{33}%
\TextAt{9}{241.721}{506.879}{\Character{43}}%
\TextAt{9}{253.473}{506.879}{11}%
\TextAt{16}{265.205}{516.271}{\Character{112}}%
\Vector{q 0 G 0 g 0.478 w 203.167 -257.49 m 218.319 -257.49 l S Q}%
\TextAt{16}{275.167}{506.879}{\Character{0}}%
\TextAt{9}{284.466}{506.879}{7}%
\TextAt{8}{294.216}{506.879}{into}%
\TextAt{8}{317.959}{506.879}{irreducible}%
\TextAt{8}{375.834}{506.879}{integral}%
\TextAt{8}{418.442}{506.879}{elements}%
\TextAt{8}{466.299}{506.879}{of}%
\TextAt{18}{479.629}{506.879}{Q}%
\TextAt{9}{488.927}{506.879}{\Character{40}}%
\TextAt{16}{493.48}{516.271}{\Character{112}}%
\Vector{q 0 G 0 g 0.478 w 431.442 -257.49 m 446.594 -257.49 l S Q}%
\TextAt{16}{503.442}{506.879}{\Character{0}}%
\TextAt{9}{512.741}{506.879}{7\Character{41}}%
\TextAt{8}{523.146}{506.879}{.}%
\TextAt{12}{79.226}{474.947}{Exercise}%
\TextAt{12}{132.427}{474.947}{I.3.2.}%
\TextAt{8}{166.021}{474.947}{Show}%
\TextAt{8}{196.913}{474.947}{that}%
\TextAt{9}{207.966}{443.409}{54}%
\TextAt{9}{222.97}{443.409}{\Character{61}}%
\TextAt{9}{235.391}{443.409}{2}%
\TextAt{16}{243.928}{443.409}{\Character{1}}%
\TextAt{9}{249.906}{443.409}{3}%
\TextAt{17}{255.759}{448.345}{3}%
\TextAt{9}{263.812}{443.409}{\Character{61}}%
\TextAt{9}{277.433}{451.497}{13}%
\TextAt{9}{291.779}{451.497}{\Character{43}}%
\TextAt{16}{303.557}{460.888}{\Character{112}}%
\Vector{q 0 G 0 g 0.478 w 241.519 -312.873 m 262.524 -312.873 l S Q}%
\TextAt{16}{313.519}{451.497}{\Character{0}}%
\TextAt{9}{322.818}{451.497}{47}%
\Vector{q 0 G 0 g 0.478 w 205.433 -327.602 m 262.524 -327.602 l S Q}%
\TextAt{9}{303.052}{435.209}{2}%
\TextAt{16}{338.376}{443.409}{\Character{1}}%
\TextAt{9}{345.549}{451.497}{13}%
\TextAt{16}{359.912}{451.497}{\Character{0}}%
\TextAt{16}{371.867}{460.888}{\Character{112}}%
\Vector{q 0 G 0 g 0.478 w 309.83 -312.873 m 330.834 -312.873 l S Q}%
\TextAt{16}{381.83}{451.497}{\Character{0}}%
\TextAt{9}{391.128}{451.497}{47}%
\Vector{q 0 G 0 g 0.478 w 273.549 -327.602 m 330.834 -327.602 l S Q}%
\TextAt{9}{371.265}{435.209}{2}%
\TextAt{8}{79.226}{412.01}{are}%
\TextAt{8}{97.727}{412.01}{two}%
\TextAt{8}{118.835}{412.01}{essentially}%
\TextAt{8}{173.883}{412.01}{different}%
\TextAt{8}{219.044}{412.01}{decompositions}%
\TextAt{8}{299.783}{412.01}{into}%
\TextAt{8}{322.517}{412.01}{irreducible}%
\TextAt{8}{379.382}{412.01}{integral}%
\TextAt{8}{420.98}{412.01}{elements}%
\TextAt{8}{467.827}{412.01}{of}%
\TextAt{18}{480.147}{412.01}{Q}%
\TextAt{9}{489.446}{412.01}{\Character{40}}%
\TextAt{16}{493.998}{421.401}{\Character{112}}%
\Vector{q 0 G 0 g 0.478 w 431.961 -352.36 m 452.965 -352.36 l S Q}%
\TextAt{16}{503.961}{412.01}{\Character{0}}%
\TextAt{9}{513.259}{412.01}{47\Character{41}}%
\TextAt{8}{529.518}{412.01}{.}%
\TextAt{12}{79.226}{380.077}{Exercise}%
\TextAt{12}{132.231}{380.077}{I.3.3.}%
\ErrataAt{165.64}{380.077}{369.36}{Let $d\ne1$ be squarefree and $p$ a prime not dividing $2d$. Let $\mathfrak o$ be the}%
\TextAt{8}{79.226}{363.031}{ring}%
\TextAt{8}{104.237}{363.031}{of}%
\TextAt{8}{118.512}{363.031}{integers}%
\TextAt{8}{162.771}{363.031}{of}%
\TextAt{18}{177.046}{363.031}{Q}%
\TextAt{9}{186.345}{363.031}{\Character{40}}%
\TextAt{16}{190.897}{373.22}{\Character{112}}%
\Vector{q 0 G 0 g 0.478 w 128.86 -400.541 m 134.942 -400.541 l S Q}%
\TextAt{13}{200.86}{363.031}{d}%
\TextAt{9}{206.942}{363.031}{\Character{41}}%
\TextAt{8}{211.495}{363.031}{.}%
\TextAt{8}{222.778}{363.031}{Show}%
\TextAt{8}{254.616}{363.031}{that}%
\TextAt{9}{280.93}{363.031}{\Character{40}}%
\TextAt{13}{285.482}{363.031}{p}%
\TextAt{9}{291.357}{363.031}{\Character{41}}%
\TextAt{9}{300.826}{363.031}{\Character{61}}%
\TextAt{13}{314.877}{363.031}{p}%
\TextAt{14}{320.752}{363.031}{o}%
\TextAt{8}{331.439}{363.031}{is}%
\TextAt{8}{344.148}{363.031}{a}%
\TextAt{8}{354.849}{363.031}{prime}%
\TextAt{8}{388.958}{363.031}{ideal}%
\TextAt{8}{417.866}{363.031}{of}%
\TextAt{14}{432.141}{363.031}{o}%
\TextAt{8}{442.828}{363.031}{if}%
\TextAt{8}{454.497}{363.031}{and}%
\TextAt{8}{478.205}{363.031}{only}%
\TextAt{8}{504.841}{363.031}{if}%
\TextAt{8}{516.51}{363.031}{the}%
\TextAt{8}{79.226}{345.985}{congruence}%
\TextAt{13}{139.707}{345.985}{x}%
\TextAt{17}{146.359}{350.324}{2}%
\NativeMathAt{154.413}{345.985}{11.955}{\equiv}%
\TextAt{13}{167.032}{345.985}{d}%
\TextAt{8}{176.435}{345.985}{mod}%
\TextAt{13}{201.873}{345.985}{p}%
\TextAt{8}{211.646}{345.985}{has}%
\TextAt{8}{232.52}{345.985}{no}%
\TextAt{8}{248.779}{345.985}{solution.}%
\TextAt{12}{79.226}{314.053}{Exercise}%
\TextAt{12}{132.904}{314.053}{I.3.4.}%
\TextAt{8}{166.913}{314.053}{A}%
\TextAt{8}{180}{314.053}{Dedekind}%
\TextAt{8}{232.599}{314.053}{domain}%
\TextAt{8}{274.642}{314.053}{with}%
\TextAt{8}{301.717}{314.053}{a}%
\TextAt{8}{311.887}{314.053}{finite}%
\TextAt{8}{342.214}{314.053}{number}%
\TextAt{8}{385.225}{314.053}{of}%
\TextAt{8}{398.97}{314.053}{prime}%
\TextAt{8}{432.548}{314.053}{ideals}%
\TextAt{8}{465.541}{314.053}{is}%
\TextAt{8}{477.719}{314.053}{a}%
\TextAt{8}{487.89}{314.053}{principal}%
\TextAt{8}{79.226}{297.007}{ideal}%
\TextAt{8}{107.189}{297.007}{domain.}%
\TextAt{12}{97.159}{276.972}{Hint:}%
\ErrataAt{133.675}{276.972}{401.325}{Write $\mathfrak a=\prod_{i=1}^r\mathfrak p_i^{\nu_i}\ne0$, where $\mathfrak p_1,\ldots,\mathfrak p_r$ are all nonzero primes and $\nu_i\ge0$.}%
\ErrataAt{79.19}{259.926}{455.81}{Choose $\pi_i\in\mathfrak p_i\setminus\mathfrak p_i^2$ and use CRT for the cosets $\pi_i^{\nu_i}\bmod\mathfrak p_i^{\nu_i+1}$.}%
\TextAt{12}{79.226}{227.994}{Exercise}%
\TextAt{12}{133.938}{227.994}{I.3.5.}%
\TextAt{8}{168.846}{227.994}{The}%
\TextAt{8}{194.214}{227.994}{quotient}%
\TextAt{8}{241.711}{227.994}{ring}%
\TextAt{14}{267.091}{227.994}{o}%
\TextAt{9}{272.935}{227.994}{\Character{47}}%
\TextAt{14}{278.788}{227.994}{a}%
\TextAt{8}{289.983}{227.994}{of}%
\TextAt{8}{304.627}{227.994}{a}%
\TextAt{8}{315.697}{227.994}{Dedekind}%
\TextAt{8}{369.195}{227.994}{domain}%
\TextAt{8}{412.137}{227.994}{by}%
\TextAt{8}{429.71}{227.994}{an}%
\TextAt{8}{447.283}{227.994}{ideal}%
\TextAt{14}{476.561}{227.994}{a}%
\NativeMathAt{488.104}{227.994}{11.955}{\neq}%
\TextAt{9}{502.761}{227.994}{0}%
\TextAt{8}{513.833}{227.994}{is}%
\TextAt{8}{526.911}{227.994}{a}%
\EditionChange{\BodyAt{79.226}{210.955}{principal ideal domain.}}{\BodyAt{79.226}{210.955}{principal ideal ring.}}%
\ErrataAt{97.13}{190.913}{437.87}{\textbf{Hint:} For $\mathfrak a=\mathfrak p^n$, the proper ideals of $\mathfrak o/\mathfrak a$ are $\mathfrak p^\nu/\mathfrak p^n$ for $1\le \nu\le n$. Choose}%
\TextAt{13}{79.226}{173.867}{\Character{25}}%
\NativeMathAt{89.616}{173.867}{11.955}{\in}%
\TextAt{14}{100.907}{173.867}{p}%
\TextAt{16}{109.547}{173.867}{\Character{110}}%
\TextAt{14}{118.181}{173.867}{p}%
\TextAt{17}{124.164}{178.205}{2}%
\TextAt{8}{132.794}{173.867}{and}%
\TextAt{8}{155.557}{173.867}{show}%
\TextAt{8}{184.56}{173.867}{that}%
\TextAt{14}{209.929}{173.867}{p}%
\TextAt{21}{215.912}{178.205}{\Character{23}}%
\TextAt{9}{224.415}{173.867}{\Character{61}}%
\TextAt{14}{236.84}{173.867}{o}%
\TextAt{13}{242.684}{173.867}{\Character{25}}%
\TextAt{21}{249.754}{178.205}{\Character{23}}%
\TextAt{9}{257.592}{173.867}{\Character{43}}%
\TextAt{14}{269.354}{173.867}{p}%
\TextAt{21}{275.337}{178.205}{n}%
\TextAt{8}{280.973}{173.867}{.}%
\TextAt{12}{79.226}{141.935}{Exercise}%
\TextAt{12}{132.427}{141.935}{I.3.6.}%
\TextAt{8}{166.021}{141.935}{Every}%
\TextAt{8}{199.674}{141.935}{ideal}%
\TextAt{8}{227.638}{141.935}{of}%
\TextAt{8}{240.968}{141.935}{a}%
\TextAt{8}{250.723}{141.935}{Dedekind}%
\TextAt{8}{302.907}{141.935}{domain}%
\TextAt{8}{344.535}{141.935}{can}%
\TextAt{8}{365.995}{141.935}{be}%
\TextAt{8}{381.919}{141.935}{generated}%
\TextAt{8}{435.251}{141.935}{by}%
\TextAt{8}{451.51}{141.935}{two}%
\TextAt{8}{473.627}{141.935}{elements.}%
\TextAt{12}{97.159}{121.9}{Hint:}%
\TextAt{8}{132.698}{121.9}{Use}%
\TextAt{8}{155.185}{121.9}{exercise}%
\TextAt{8}{198.487}{121.9}{5.}%
\TextAt{12}{79.226}{89.967}{Exercise}%
\TextAt{12}{134.245}{89.967}{I.3.7.}%
\TextAt{8}{169.421}{89.967}{In}%
\TextAt{8}{185.623}{89.967}{a}%
\TextAt{8}{196.959}{89.967}{noetherian}%
\TextAt{8}{256.75}{89.967}{ring}%
\TextAt{13}{282.397}{89.967}{R}%
\TextAt{8}{296.884}{89.967}{in}%
\TextAt{8}{312.118}{89.967}{which}%
\TextAt{8}{347.185}{89.967}{every}%
\TextAt{8}{379.659}{89.967}{prime}%
\TextAt{8}{414.403}{89.967}{ideal}%
\TextAt{8}{443.948}{89.967}{is}%
\TextAt{8}{457.293}{89.967}{maximal,}%
\TextAt{8}{510.33}{89.967}{each}%
\TextAt{8}{79.226}{72.921}{descending}%
\TextAt{8}{138.464}{72.921}{chain}%
\TextAt{8}{169.356}{72.921}{of}%
\TextAt{8}{182.686}{72.921}{ideals}%
\TextAt{14}{215.264}{72.921}{a}%
\TextAt{17}{221.247}{71.128}{1}%
\NativeMathAt{229.3}{72.921}{11.955}{\supseteq}%
\TextAt{14}{241.92}{72.921}{a}%
\TextAt{17}{247.903}{71.128}{2}%
\NativeMathAt{255.956}{72.921}{11.955}{\supseteq}%
\TextAt{16}{268.557}{72.921}{\Character{1}}%
\TextAt{16}{273.865}{72.921}{\Character{1}}%
\TextAt{16}{279.173}{72.921}{\Character{1}}%
\TextAt{8}{288.413}{72.921}{becomes}%
\TextAt{8}{334.966}{72.921}{stationary.}%
\TextAt{8}{300.14}{43.033}{21}%
\Destination{79.36}{229.955}{errata.1}%
\Destination{78.23}{809.815}{nav.page.21}%
\Destination{79.226}{520.855}{nav.ex.1.3.1}%
\Destination{79.226}{488.955}{nav.ex.1.3.2}%
\Destination{79.226}{394.055}{nav.ex.1.3.3}%
\Destination{79.226}{328.055}{nav.ex.1.3.4}%
\Destination{79.226}{241.955}{nav.ex.1.3.5}%
\Destination{79.226}{155.955}{nav.ex.1.3.6}%
\Destination{79.226}{103.955}{nav.ex.1.3.7}%
\Destination{79.36}{521.855}{nav.print.23}%
\Destination{79.226}{207.913}{errata.60}%
\Destination{79.226}{397.077}{errata.68}%
\Destination{79.226}{293.972}{errata.69}%
}

% Page 34 (body page 22)
\Page{%
\TextAt{2}{79.226}{796.943}{I.}%
\TextAt{2}{89.83}{796.943}{Algebraic}%
\TextAt{2}{138.648}{796.943}{Integers}%
\Vector{q 0 G 0 g 0.299 w 7.226 18.729 m 460.769 18.729 l S Q}%
\TextAt{12}{97.159}{761.124}{Hint:}%
\TextAt{8}{134.697}{761.124}{Show}%
\TextAt{8}{166.589}{761.124}{as}%
\TextAt{8}{181.958}{761.124}{in}%
\TextAt{8}{196.611}{761.124}{(3.4)}%
\TextAt{8}{225.586}{761.124}{that}%
\TextAt{9}{251.954}{761.124}{\Character{40}0\Character{41}}%
\TextAt{8}{271.809}{761.124}{is}%
\TextAt{8}{284.572}{761.124}{a}%
\TextAt{8}{295.327}{761.124}{product}%
\TextAt{14}{340.226}{761.124}{p}%
\TextAt{17}{346.209}{759.33}{1}%
\TextAt{16}{352.934}{761.124}{\Character{1}}%
\TextAt{16}{358.23}{761.124}{\Character{1}}%
\TextAt{16}{363.538}{761.124}{\Character{1}}%
\TextAt{14}{368.874}{761.124}{p}%
\TextAt{21}{374.858}{759.33}{r}%
\TextAt{8}{384.309}{761.124}{of}%
\TextAt{8}{398.639}{761.124}{prime}%
\TextAt{8}{432.802}{761.124}{ideals}%
\TextAt{8}{466.38}{761.124}{and}%
\TextAt{8}{490.142}{761.124}{that}%
\TextAt{8}{516.51}{761.124}{the}%
\TextAt{8}{79.226}{744.078}{chain}%
\TextAt{13}{110.118}{744.078}{R}%
\NativeMathAt{122.448}{744.078}{11.955}{\supseteq}%
\TextAt{14}{135.067}{744.078}{p}%
\TextAt{17}{141.05}{742.284}{1}%
\NativeMathAt{149.103}{744.078}{11.955}{\supseteq}%
\TextAt{14}{161.723}{744.078}{p}%
\TextAt{17}{167.706}{742.284}{1}%
\TextAt{14}{172.438}{744.078}{p}%
\TextAt{17}{178.421}{742.284}{2}%
\NativeMathAt{186.474}{744.078}{11.955}{\supseteq}%
\TextAt{16}{199.075}{744.078}{\Character{1}}%
\TextAt{16}{204.383}{744.078}{\Character{1}}%
\TextAt{16}{209.691}{744.078}{\Character{1}}%
\NativeMathAt{216.314}{744.078}{11.955}{\supseteq}%
\TextAt{14}{228.981}{744.078}{p}%
\TextAt{17}{234.964}{742.284}{1}%
\TextAt{16}{241.689}{744.078}{\Character{1}}%
\TextAt{16}{246.985}{744.078}{\Character{1}}%
\TextAt{16}{252.293}{744.078}{\Character{1}}%
\TextAt{14}{257.629}{744.078}{p}%
\TextAt{21}{263.613}{742.284}{r}%
\TextAt{9}{271.488}{744.078}{\Character{61}}%
\TextAt{9}{283.898}{744.078}{\Character{40}0\Character{41}}%
\TextAt{8}{302.769}{744.078}{can}%
\TextAt{8}{324.228}{744.078}{be}%
\TextAt{8}{340.152}{744.078}{refined}%
\TextAt{8}{378.517}{744.078}{into}%
\TextAt{8}{402.26}{744.078}{a}%
\TextAt{8}{412.015}{744.078}{composition}%
\TextAt{8}{477.446}{744.078}{series.}%
\TextAt{12}{79.226}{712.088}{Exercise}%
\TextAt{12}{132.342}{712.088}{I.3.8.}%
\TextAt{8}{165.862}{712.088}{Let}%
\TextAt{14}{186.758}{712.088}{m}%
\TextAt{8}{199.748}{712.088}{be}%
\TextAt{8}{215.599}{712.088}{a}%
\TextAt{8}{225.281}{712.088}{nonzero}%
\TextAt{8}{268.784}{712.088}{integral}%
\TextAt{8}{311.318}{712.088}{ideal}%
\TextAt{8}{339.208}{712.088}{of}%
\TextAt{8}{352.464}{712.088}{the}%
\TextAt{8}{372.547}{712.088}{Dedekind}%
\TextAt{8}{424.657}{712.088}{domain}%
\TextAt{14}{466.212}{712.088}{o}%
\TextAt{8}{472.056}{712.088}{.}%
\TextAt{8}{480.48}{712.088}{Show}%
\TextAt{8}{511.298}{712.088}{that}%
\TextAt{8}{79.226}{695.042}{in}%
\TextAt{8}{92.879}{695.042}{every}%
\TextAt{8}{123.771}{695.042}{ideal}%
\TextAt{8}{151.734}{695.042}{class}%
\TextAt{8}{179.171}{695.042}{of}%
\TextAt{13}{192.501}{695.042}{C}%
\TextAt{13}{201.719}{695.042}{l}%
\TextAt{21}{205.222}{693.249}{K}%
\TextAt{8}{213.425}{695.042}{,}%
\TextAt{8}{220.574}{695.042}{there}%
\TextAt{8}{250.486}{695.042}{exists}%
\TextAt{8}{282.8}{695.042}{an}%
\TextAt{8}{299.059}{695.042}{integral}%
\TextAt{8}{341.668}{695.042}{ideal}%
\TextAt{8}{369.631}{695.042}{prime}%
\TextAt{8}{402.794}{695.042}{to}%
\TextAt{14}{417.105}{695.042}{m}%
\TextAt{8}{426.272}{695.042}{.}%
\TextAt{12}{79.226}{663.052}{Exercise}%
\TextAt{12}{132.605}{663.052}{I.3.9.}%
\TextAt{8}{166.354}{663.052}{Let}%
\TextAt{14}{187.478}{663.052}{o}%
\TextAt{8}{197.374}{663.052}{be}%
\TextAt{8}{213.453}{663.052}{an}%
\TextAt{8}{229.867}{663.052}{integral}%
\TextAt{8}{272.63}{663.052}{domain}%
\TextAt{8}{314.413}{663.052}{in}%
\TextAt{8}{328.22}{663.052}{which}%
\TextAt{8}{361.862}{663.052}{all}%
\TextAt{8}{378.275}{663.052}{nonzero}%
\TextAt{8}{422.007}{663.052}{ideals}%
\TextAt{8}{454.739}{663.052}{admit}%
\TextAt{8}{488.715}{663.052}{a}%
\TextAt{8}{498.626}{663.052}{unique}%
\TextAt{8}{79.226}{646.006}{factorization}%
\TextAt{8}{147.203}{646.006}{into}%
\TextAt{8}{170.946}{646.006}{prime}%
\TextAt{8}{204.11}{646.006}{ideals.}%
\TextAt{8}{241.239}{646.006}{Show}%
\TextAt{8}{272.131}{646.006}{that}%
\TextAt{14}{297.5}{646.006}{o}%
\TextAt{8}{307.241}{646.006}{is}%
\TextAt{8}{319.005}{646.006}{a}%
\TextAt{8}{328.761}{646.006}{Dedekind}%
\TextAt{8}{380.945}{646.006}{domain.}%
\TextAt{12}{79.226}{614.016}{Exercise}%
\TextAt{12}{131.203}{614.016}{I.3.10.}%
\TextAt{8}{170.464}{614.016}{The}%
\TextAt{8}{193.454}{614.016}{fractional}%
\TextAt{8}{244.754}{614.016}{ideals}%
\TextAt{14}{276.268}{614.016}{a}%
\TextAt{8}{285.084}{614.016}{of}%
\TextAt{8}{297.351}{614.016}{a}%
\TextAt{8}{306.042}{614.016}{Dedekind}%
\TextAt{8}{357.163}{614.016}{domain}%
\TextAt{14}{397.727}{614.016}{o}%
\TextAt{8}{406.405}{614.016}{are}%
\TextAt{8}{424.852}{614.016}{projective}%
\TextAt{14}{477.766}{614.016}{o}%
\TextAt{8}{483.61}{614.016}{-modules,}%
\TextAt{8}{79.226}{596.97}{i.e.,}%
\TextAt{8}{103.33}{596.97}{given}%
\TextAt{8}{135.498}{596.97}{any}%
\TextAt{8}{159.214}{596.97}{surjective}%
\TextAt{8}{213.224}{596.97}{homomorphism}%
\TextAt{13}{297.493}{596.97}{M}%
\TextAt{21}{318.911}{604.234}{f}%
\NativeArrowAt{316.112}{596.97}{11.955}{11.955}{\longrightarrow}%
\TextAt{13}{334.113}{596.97}{N}%
\TextAt{8}{350.232}{596.97}{of}%
\TextAt{14}{365.161}{596.97}{o}%
\TextAt{8}{371.005}{596.97}{-modules,}%
\TextAt{8}{426.061}{596.97}{each}%
\TextAt{8}{453.997}{596.97}{homomorphism}%
\TextAt{14}{79.226}{579.924}{a}%
\TextAt{21}{91.644}{587.188}{g}%
\NativeArrowAt{88.53}{579.924}{11.955}{11.955}{\longrightarrow}%
\TextAt{13}{103.806}{579.924}{N}%
\TextAt{8}{118.326}{579.924}{can}%
\TextAt{8}{139.786}{579.924}{be}%
\TextAt{8}{155.71}{579.924}{lifted}%
\TextAt{8}{185.945}{579.924}{to}%
\TextAt{8}{200.255}{579.924}{a}%
\TextAt{8}{210.01}{579.924}{homomorphism}%
\TextAt{13}{292.68}{579.924}{h}%
\TextAt{9}{302.74}{579.924}{\Character{58}}%
\TextAt{14}{309.312}{579.924}{a}%
\NativeMathAt{318.616}{579.924}{11.955}{\rightarrow}%
\TextAt{13}{333.892}{579.924}{M}%
\TextAt{8}{350.363}{579.924}{such}%
\TextAt{8}{376.76}{579.924}{that}%
\TextAt{13}{402.129}{579.924}{f}%
\TextAt{16}{411.832}{579.924}{\Character{14}}%
\TextAt{13}{420.467}{579.924}{h}%
\TextAt{9}{430.526}{579.924}{\Character{61}}%
\TextAt{13}{442.951}{579.924}{g}%
\TextAt{8}{448.986}{579.924}{.}%
\TextAt{15}{79.226}{527.06}{§}%
\TextAt{15}{93.785}{527.06}{4}%
\TextAt{15}{120.693}{527.06}{Lattices}%
\TextAt{8}{79.226}{495.175}{In}%
\TextAt{8}{94.834}{495.175}{§1,}%
\TextAt{8}{114.754}{495.175}{when}%
\TextAt{8}{146.298}{495.175}{solving}%
\TextAt{8}{186.701}{495.175}{the}%
\TextAt{8}{207.845}{495.175}{basic}%
\TextAt{8}{238.158}{495.175}{problems}%
\TextAt{8}{289.284}{495.175}{concerning}%
\TextAt{8}{348.804}{495.175}{the}%
\TextAt{8}{369.947}{495.175}{gaussian}%
\TextAt{8}{417.894}{495.175}{integers,}%
\TextAt{8}{465.693}{495.175}{we}%
\TextAt{8}{483.907}{495.175}{used}%
\TextAt{8}{511.614}{495.175}{at}%
\TextAt{8}{526.911}{495.175}{a}%
\TextAt{8}{79.226}{478.129}{crucial}%
\TextAt{8}{116.945}{478.129}{place}%
\TextAt{8}{146.857}{478.129}{the}%
\TextAt{8}{167.013}{478.129}{inclusion}%
\TextAt{18}{284.477}{461.083}{Z}%
\TextAt{9}{292.447}{461.083}{\Character{91}}%
\TextAt{13}{295.699}{461.083}{i}%
\TextAt{9}{299.692}{461.083}{\Character{93}}%
\NativeMathAt{306.265}{461.083}{11.955}{\subseteq}%
\TextAt{18}{318.884}{461.083}{C}%
\TextAt{8}{79.226}{437.561}{and}%
\TextAt{8}{102.161}{437.561}{considered}%
\TextAt{8}{159.622}{437.561}{the}%
\TextAt{8}{179.951}{437.561}{integers}%
\TextAt{8}{223.437}{437.561}{of}%
\TextAt{18}{236.939}{437.561}{Q}%
\TextAt{9}{246.238}{437.561}{\Character{40}}%
\TextAt{13}{250.79}{437.561}{i}%
\TextAt{9}{254.784}{437.561}{\Character{41}}%
\TextAt{8}{263.406}{437.561}{as}%
\TextAt{8}{277.948}{437.561}{lattice}%
\TextAt{8}{313.89}{437.561}{points}%
\TextAt{8}{349.246}{437.561}{in}%
\TextAt{8}{363.071}{437.561}{the}%
\TextAt{8}{383.4}{437.561}{complex}%
\TextAt{8}{429.42}{437.561}{plane.}%
\TextAt{8}{465.703}{437.561}{This}%
\TextAt{8}{492.595}{437.561}{point}%
\TextAt{8}{523.337}{437.561}{of}%
\TextAt{8}{79.226}{420.515}{view}%
\TextAt{8}{106.546}{420.515}{has}%
\TextAt{8}{127.758}{420.515}{been}%
\TextAt{8}{155.723}{420.515}{generalized}%
\TextAt{8}{216.542}{420.515}{to}%
\TextAt{8}{231.19}{420.515}{arbitrary}%
\TextAt{8}{281.297}{420.515}{number}%
\TextAt{8}{324.231}{420.515}{fields}%
\TextAt{8}{354.54}{420.515}{by}%
\TextAt{26}{371.136}{420.515}{Hermann}%
\TextAt{26}{430.562}{420.515}{Minkowski}%
\TextAt{8}{498.924}{420.515}{(1864–}%
\TextAt{8}{79.226}{403.469}{1909)}%
\TextAt{8}{110.024}{403.469}{and}%
\TextAt{8}{131.7}{403.469}{has}%
\TextAt{8}{151.488}{403.469}{led}%
\TextAt{8}{169.255}{403.469}{to}%
\TextAt{8}{182.478}{403.469}{results}%
\TextAt{8}{218.585}{403.469}{which}%
\TextAt{8}{250.985}{403.469}{make}%
\TextAt{8}{280.468}{403.469}{up}%
\TextAt{8}{296.286}{403.469}{an}%
\TextAt{8}{311.458}{403.469}{essential}%
\TextAt{8}{356.997}{403.469}{part}%
\TextAt{8}{381.28}{403.469}{of}%
\TextAt{8}{393.523}{403.469}{the}%
\TextAt{8}{412.593}{403.469}{foundations}%
\TextAt{8}{474.989}{403.469}{of}%
\TextAt{8}{487.232}{403.469}{algebraic}%
\TextAt{8}{79.226}{386.423}{number}%
\TextAt{8}{122.824}{386.423}{theory.}%
\TextAt{8}{166.148}{386.423}{In}%
\TextAt{8}{181.771}{386.423}{order}%
\TextAt{8}{213.342}{386.423}{to}%
\TextAt{8}{228.653}{386.423}{develop}%
\TextAt{8}{271.928}{386.423}{Minkowski’s}%
\TextAt{8}{339.065}{386.423}{theory}%
\TextAt{8}{376.817}{386.423}{we}%
\TextAt{8}{395.046}{386.423}{first}%
\TextAt{8}{420.172}{386.423}{have}%
\TextAt{8}{448.169}{386.423}{to}%
\TextAt{8}{463.48}{386.423}{introduce}%
\TextAt{8}{516.51}{386.423}{the}%
\TextAt{8}{79.226}{369.377}{general}%
\TextAt{8}{119.551}{369.377}{notion}%
\TextAt{8}{155.978}{369.377}{of}%
\TextAt{8}{169.308}{369.377}{lattice}%
\TextAt{8}{205.078}{369.377}{and}%
\TextAt{8}{227.841}{369.377}{study}%
\TextAt{8}{260.096}{369.377}{some}%
\TextAt{8}{289.422}{369.377}{of}%
\TextAt{8}{302.752}{369.377}{its}%
\TextAt{8}{319.071}{369.377}{basic}%
\TextAt{8}{348.397}{369.377}{properties.}%
\TextAt{12}{79.226}{343.365}{Definition}%
\TextAt{12}{140.551}{343.365}{(4.1).}%
\TextAt{8}{174.595}{343.365}{Let}%
\TextAt{13}{194.292}{343.365}{V}%
\TextAt{8}{206.343}{343.365}{be}%
\TextAt{8}{220.995}{343.365}{an}%
\TextAt{13}{235.982}{343.365}{n}%
\TextAt{8}{242.97}{343.365}{-dimensional}%
\TextAt{18}{310.046}{343.365}{R}%
\TextAt{8}{318.68}{343.365}{-vector}%
\TextAt{8}{356.43}{343.365}{space.}%
\TextAt{8}{391.831}{343.365}{A}%
\TextAt{12}{403.23}{343.365}{lattice}%
\TextAt{8}{442.462}{343.365}{in}%
\TextAt{13}{454.843}{343.365}{V}%
\TextAt{8}{466.894}{343.365}{is}%
\TextAt{8}{477.386}{343.365}{a}%
\TextAt{8}{485.869}{343.365}{subgroup}%
\TextAt{8}{79.226}{326.319}{of}%
\TextAt{8}{92.556}{326.319}{the}%
\TextAt{8}{112.713}{326.319}{form}%
\TextAt{9}{253.086}{309.273}{\Character{0}}%
\TextAt{9}{263.702}{309.273}{\Character{61}}%
\TextAt{18}{276.146}{309.273}{Z}%
\TextAt{13}{284.117}{309.273}{v}%
\TextAt{17}{289.776}{307.479}{1}%
\TextAt{9}{297.165}{309.273}{\Character{43}}%
\TextAt{16}{308.926}{309.273}{\Character{1}}%
\TextAt{16}{314.222}{309.273}{\Character{1}}%
\TextAt{16}{319.53}{309.273}{\Character{1}}%
\TextAt{9}{325.53}{309.273}{\Character{43}}%
\TextAt{18}{337.292}{309.273}{Z}%
\TextAt{13}{345.262}{309.273}{v}%
\TextAt{21}{350.921}{307.479}{m}%
\TextAt{8}{79.226}{285.751}{with}%
\TextAt{8}{106.697}{285.751}{linearly}%
\TextAt{8}{149.459}{285.751}{independent}%
\TextAt{8}{216.597}{285.751}{vectors}%
\TextAt{13}{257.147}{285.751}{v}%
\TextAt{17}{262.807}{283.958}{1}%
\TextAt{13}{267.539}{285.751}{\Character{59}}%
\TextAt{13}{272.763}{285.751}{\Character{58}}%
\TextAt{13}{278}{285.751}{\Character{58}}%
\TextAt{13}{283.224}{285.751}{\Character{58}}%
\TextAt{13}{288.461}{285.751}{\Character{59}}%
\TextAt{13}{293.697}{285.751}{v}%
\TextAt{21}{299.419}{283.958}{m}%
\TextAt{8}{312.116}{285.751}{of}%
\TextAt{13}{326.257}{285.751}{V}%
\TextAt{8}{335.683}{285.751}{.}%
\TextAt{8}{346.564}{285.751}{The}%
\TextAt{13}{371.429}{285.751}{m}%
\TextAt{8}{381.668}{285.751}{-tuple}%
\TextAt{9}{416.289}{285.751}{\Character{40}}%
\TextAt{13}{420.841}{285.751}{v}%
\TextAt{17}{426.5}{283.958}{1}%
\TextAt{13}{431.232}{285.751}{\Character{59}}%
\TextAt{13}{436.456}{285.751}{\Character{58}}%
\TextAt{13}{441.693}{285.751}{\Character{58}}%
\TextAt{13}{446.929}{285.751}{\Character{58}}%
\TextAt{13}{452.154}{285.751}{\Character{59}}%
\TextAt{13}{457.39}{285.751}{v}%
\TextAt{21}{463.112}{283.958}{m}%
\TextAt{9}{471.101}{285.751}{\Character{41}}%
\TextAt{8}{480.362}{285.751}{is}%
\TextAt{8}{492.937}{285.751}{called}%
\TextAt{8}{526.912}{285.751}{a}%
\TextAt{12}{79.226}{268.705}{basis}%
\TextAt{8}{111.493}{268.705}{and}%
\TextAt{8}{134.256}{268.705}{the}%
\TextAt{8}{154.412}{268.705}{set}%
\TextAt{9}{188.591}{239.704}{\Character{8}}%
\TextAt{9}{200.355}{239.704}{\Character{61}}%
\TextAt{16}{212.792}{239.704}{\Character{102}}%
\TextAt{13}{218.769}{239.704}{x}%
\TextAt{17}{225.421}{237.911}{1}%
\TextAt{13}{230.154}{239.704}{v}%
\TextAt{17}{235.813}{237.911}{1}%
\TextAt{9}{243.202}{239.704}{\Character{43}}%
\TextAt{16}{254.963}{239.704}{\Character{1}}%
\TextAt{16}{260.259}{239.704}{\Character{1}}%
\TextAt{16}{265.567}{239.704}{\Character{1}}%
\TextAt{9}{271.568}{239.704}{\Character{43}}%
\TextAt{13}{283.329}{239.704}{x}%
\TextAt{21}{289.981}{237.911}{m}%
\TextAt{13}{297.97}{239.704}{v}%
\TextAt{21}{303.629}{237.911}{m}%
\TextAt{16}{314.938}{239.704}{\Character{106}}%
\TextAt{13}{321.58}{239.704}{x}%
\TextAt{21}{328.232}{237.911}{i}%
\NativeMathAt{334.934}{239.704}{11.955}{\in}%
\TextAt{18}{346.225}{239.704}{R}%
\TextAt{13}{354.86}{239.704}{\Character{59}}%
\TextAt{9}{364.001}{239.704}{0}%
\NativeMathAt{373.175}{239.704}{11.955}{\leq}%
\TextAt{13}{385.794}{239.704}{x}%
\TextAt{21}{392.446}{237.911}{i}%
\TextAt{13}{399.149}{239.704}{\Character{60}}%
\TextAt{9}{411.574}{239.704}{1}%
\TextAt{16}{417.427}{239.704}{\Character{103}}%
\TextAt{8}{79.226}{210.703}{a}%
\TextAt{12}{89.43}{210.703}{fundamental}%
\TextAt{12}{167.464}{210.703}{mesh}%
\TextAt{8}{201.937}{210.703}{of}%
\TextAt{8}{215.716}{210.703}{the}%
\TextAt{8}{236.321}{210.703}{lattice.}%
\TextAt{8}{277.988}{210.703}{The}%
\TextAt{8}{302.49}{210.703}{lattice}%
\TextAt{8}{338.709}{210.703}{is}%
\TextAt{8}{350.922}{210.703}{called}%
\TextAt{12}{384.534}{210.703}{complete}%
\TextAt{8}{441.519}{210.703}{or}%
\TextAt{8}{456.278}{210.703}{a}%
\TextAt{18}{466.482}{210.703}{Z}%
\TextAt{12}{474.452}{210.703}{-structure}%
\TextAt{8}{79.226}{193.657}{of}%
\TextAt{13}{92.556}{193.657}{V}%
\TextAt{8}{101.982}{193.657}{,}%
\TextAt{8}{109.131}{193.657}{if}%
\TextAt{13}{119.855}{193.657}{m}%
\TextAt{9}{133.415}{193.657}{\Character{61}}%
\TextAt{13}{145.84}{193.657}{n}%
\TextAt{8}{152.828}{193.657}{.}%
\TextAt{8}{97.159}{167.644}{The}%
\TextAt{8}{121.053}{167.644}{completeness}%
\TextAt{8}{191.25}{167.644}{of}%
\TextAt{8}{204.421}{167.644}{the}%
\TextAt{8}{224.418}{167.644}{lattice}%
\TextAt{8}{260.028}{167.644}{is}%
\TextAt{8}{271.633}{167.644}{obviously}%
\TextAt{8}{323.251}{167.644}{tantamount}%
\TextAt{8}{386.849}{167.644}{to}%
\TextAt{8}{401}{167.644}{the}%
\TextAt{8}{420.997}{167.644}{fact}%
\TextAt{8}{443.923}{167.644}{that}%
\TextAt{8}{469.132}{167.644}{the}%
\TextAt{8}{489.129}{167.644}{set}%
\TextAt{8}{507.237}{167.644}{of}%
\TextAt{8}{520.408}{167.644}{all}%
\TextAt{8}{79.226}{150.598}{translates}%
\TextAt{9}{132.69}{150.598}{\Character{8}}%
\TextAt{9}{143.796}{150.598}{\Character{43}}%
\TextAt{13}{155.562}{150.598}{\Character{13}}%
\TextAt{8}{162.284}{150.598}{,}%
\TextAt{13}{169.433}{150.598}{\Character{13}}%
\NativeMathAt{179.476}{150.598}{11.955}{\in}%
\TextAt{9}{190.767}{150.598}{\Character{0}}%
\TextAt{8}{198.082}{150.598}{,}%
\TextAt{8}{205.231}{150.598}{of}%
\TextAt{8}{218.561}{150.598}{the}%
\TextAt{8}{238.717}{150.598}{fundamental}%
\TextAt{8}{306.36}{150.598}{mesh}%
\TextAt{8}{336.331}{150.598}{covers}%
\TextAt{8}{371.192}{150.598}{the}%
\TextAt{8}{391.349}{150.598}{entire}%
\TextAt{8}{424.19}{150.598}{space}%
\TextAt{13}{455.464}{150.598}{V}%
\TextAt{8}{464.89}{150.598}{.}%
\TextAt{8}{97.159}{130.564}{The}%
\TextAt{8}{121.92}{130.564}{above}%
\TextAt{8}{155.803}{130.564}{definition}%
\TextAt{8}{208.539}{130.564}{makes}%
\TextAt{8}{244.43}{130.564}{use}%
\TextAt{8}{265.354}{130.564}{of}%
\TextAt{8}{279.391}{130.564}{a}%
\TextAt{8}{289.853}{130.564}{choice}%
\TextAt{8}{325.35}{130.564}{of}%
\TextAt{8}{339.387}{130.564}{linearly}%
\TextAt{8}{382.045}{130.564}{independent}%
\TextAt{8}{449.08}{130.564}{vectors.}%
\TextAt{8}{495.491}{130.564}{But}%
\TextAt{8}{519.439}{130.564}{we}%
\TextAt{8}{79.226}{113.518}{will}%
\TextAt{8}{101.438}{113.518}{need}%
\TextAt{8}{128.851}{113.518}{a}%
\TextAt{8}{138.714}{113.518}{characterization}%
\TextAt{8}{224.36}{113.518}{of}%
\TextAt{8}{237.797}{113.518}{lattices}%
\TextAt{8}{278.289}{113.518}{which}%
\TextAt{8}{311.882}{113.518}{is}%
\TextAt{8}{323.753}{113.518}{independent}%
\TextAt{8}{390.188}{113.518}{of}%
\TextAt{8}{403.625}{113.518}{such}%
\TextAt{8}{430.129}{113.518}{a}%
\TextAt{8}{439.991}{113.518}{choice.}%
\TextAt{8}{479.653}{113.518}{Note}%
\TextAt{8}{508.046}{113.518}{that,}%
\TextAt{8}{79.226}{96.472}{first}%
\TextAt{8}{102.765}{96.472}{of}%
\TextAt{8}{115.509}{96.472}{all,}%
\TextAt{8}{134.551}{96.472}{a}%
\TextAt{8}{143.72}{96.472}{lattice}%
\TextAt{8}{178.904}{96.472}{is}%
\TextAt{8}{190.082}{96.472}{a}%
\TextAt{8}{199.251}{96.472}{finitely}%
\TextAt{8}{238.009}{96.472}{generated}%
\TextAt{8}{290.755}{96.472}{subgroup}%
\TextAt{8}{340.966}{96.472}{of}%
\TextAt{13}{353.71}{96.472}{V}%
\TextAt{8}{363.135}{96.472}{.}%
\TextAt{8}{371.388}{96.472}{But}%
\TextAt{8}{394.043}{96.472}{not}%
\TextAt{8}{414.271}{96.472}{every}%
\TextAt{8}{444.577}{96.472}{finitely}%
\TextAt{8}{483.335}{96.472}{generated}%
\TextAt{8}{79.226}{79.426}{subgroup}%
\TextAt{8}{128.751}{79.426}{is}%
\TextAt{8}{139.242}{79.426}{a}%
\TextAt{8}{147.725}{79.426}{lattice—for}%
\TextAt{8}{207.914}{79.426}{instance}%
\TextAt{18}{252.226}{79.426}{Z}%
\TextAt{9}{260.25}{79.426}{\Character{43}}%
\TextAt{18}{269.409}{79.426}{Z}%
\TextAt{16}{277.379}{89.315}{\Character{112}}%
\Vector{q 0 G 0 g 0.478 w 215.342 -684.446 m 221.195 -684.446 l S Q}%
\TextAt{9}{287.342}{79.426}{2}%
\NativeMathAt{296.516}{79.426}{11.955}{\subseteq}%
\TextAt{18}{309.135}{79.426}{R}%
\TextAt{8}{320.394}{79.426}{is}%
\TextAt{8}{330.885}{79.426}{not.}%
\TextAt{8}{355.826}{79.426}{But}%
\TextAt{8}{377.794}{79.426}{each}%
\TextAt{8}{402.859}{79.426}{lattice}%
\TextAt{9}{437.356}{79.426}{\Character{0}}%
\TextAt{9}{447.972}{79.426}{\Character{61}}%
\TextAt{18}{460.416}{79.426}{Z}%
\TextAt{13}{468.387}{79.426}{v}%
\TextAt{17}{474.046}{77.633}{1}%
\TextAt{9}{478.832}{79.426}{\Character{43}}%
\TextAt{16}{487.991}{79.426}{\Character{1}}%
\TextAt{16}{493.287}{79.426}{\Character{1}}%
\TextAt{16}{498.595}{79.426}{\Character{1}}%
\TextAt{9}{501.993}{79.426}{\Character{43}}%
\TextAt{18}{511.152}{79.426}{Z}%
\TextAt{13}{519.122}{79.426}{v}%
\TextAt{21}{524.781}{77.633}{m}%
\TextAt{8}{300.14}{43.033}{22}%
\Destination{78.23}{809.815}{nav.page.22}%
\Destination{79.23}{553.585}{nav.section.1.4}%
\Destination{79.226}{726.055}{nav.ex.1.3.8}%
\Destination{79.226}{677.055}{nav.ex.1.3.9}%
\Destination{79.226}{628.055}{nav.ex.1.3.10}%
\Destination{79.36}{452.555}{nav.print.24}%
\Destination{79.36}{357.355}{nav.number.1.4.1}%
\Destination{79.23}{553.585}{outline.9}%
\Hotlink{100.013}{492.955}{6.858}{11.047}{nav.section.1.1}%
\Hotlink{196.111}{758.955}{25.078}{11.047}{nav.number.1.3.4}%
}

% Page 35 (body page 23)
\Page{%
\TextAt{2}{79.226}{796.943}{I.}%
\TextAt{2}{89.83}{796.943}{Algebraic}%
\TextAt{2}{138.648}{796.943}{Integers}%
\Vector{q 0 G 0 g 0.299 w 7.226 18.729 m 460.769 18.729 l S Q}%
\TextAt{8}{79.226}{761.124}{has}%
\TextAt{8}{100.944}{761.124}{the}%
\TextAt{8}{121.945}{761.124}{special}%
\TextAt{8}{160.89}{761.124}{property}%
\TextAt{8}{209.543}{761.124}{of}%
\TextAt{8}{223.717}{761.124}{being}%
\TextAt{8}{256.099}{761.124}{a}%
\TextAt{12}{266.698}{761.124}{discrete}%
\TextAt{8}{316.941}{761.124}{subgroup}%
\TextAt{8}{368.583}{761.124}{of}%
\TextAt{13}{382.757}{761.124}{V}%
\TextAt{8}{392.182}{761.124}{.}%
\TextAt{8}{403.163}{761.124}{This}%
\TextAt{8}{430.727}{761.124}{is}%
\TextAt{8}{443.335}{761.124}{to}%
\TextAt{8}{458.489}{761.124}{say}%
\TextAt{8}{479.562}{761.124}{that}%
\TextAt{8}{505.775}{761.124}{every}%
\TextAt{8}{79.226}{744.078}{point}%
\TextAt{13}{110.833}{744.078}{\Character{13}}%
\NativeMathAt{122.644}{744.078}{11.955}{\in}%
\TextAt{9}{135.703}{744.078}{\Character{0}}%
\TextAt{8}{147.952}{744.078}{is}%
\TextAt{8}{160.753}{744.078}{an}%
\TextAt{8}{178.05}{744.078}{isolated}%
\TextAt{8}{222.078}{744.078}{point}%
\TextAt{8}{253.685}{744.078}{in}%
\TextAt{8}{268.375}{744.078}{the}%
\TextAt{8}{289.569}{744.078}{sense}%
\TextAt{8}{320.638}{744.078}{that}%
\TextAt{8}{347.044}{744.078}{there}%
\TextAt{8}{377.994}{744.078}{exists}%
\TextAt{8}{411.346}{744.078}{a}%
\TextAt{8}{422.139}{744.078}{neighbourhood}%
\TextAt{8}{503.18}{744.078}{which}%
\TextAt{8}{79.226}{727.032}{contains}%
\TextAt{8}{125.146}{727.032}{no}%
\TextAt{8}{141.405}{727.032}{other}%
\TextAt{8}{171.974}{727.032}{points}%
\TextAt{8}{207.158}{727.032}{of}%
\TextAt{9}{220.488}{727.032}{\Character{0}}%
\TextAt{8}{227.803}{727.032}{.}%
\TextAt{8}{236.251}{727.032}{In}%
\TextAt{8}{250.872}{727.032}{fact,}%
\TextAt{8}{277.209}{727.032}{if}%
\TextAt{13}{236.258}{698.662}{\Character{13}}%
\TextAt{9}{246.301}{698.662}{\Character{61}}%
\TextAt{13}{258.726}{698.662}{a}%
\TextAt{17}{264.871}{696.869}{1}%
\TextAt{13}{269.604}{698.662}{v}%
\TextAt{17}{275.263}{696.869}{1}%
\TextAt{9}{282.652}{698.662}{\Character{43}}%
\TextAt{16}{294.413}{698.662}{\Character{1}}%
\TextAt{16}{299.709}{698.662}{\Character{1}}%
\TextAt{16}{305.017}{698.662}{\Character{1}}%
\TextAt{9}{311.017}{698.662}{\Character{43}}%
\TextAt{13}{322.779}{698.662}{a}%
\TextAt{21}{328.924}{696.869}{m}%
\TextAt{13}{336.912}{698.662}{v}%
\TextAt{21}{342.572}{696.869}{m}%
\NativeMathAt{353.881}{698.662}{11.955}{\in}%
\TextAt{9}{365.172}{698.662}{\Character{0}}%
\TextAt{13}{372.486}{698.662}{\Character{59}}%
\TextAt{8}{79.226}{670.292}{then,}%
\TextAt{8}{109.138}{670.292}{extending}%
\TextAt{13}{162.793}{670.292}{v}%
\TextAt{17}{168.452}{668.499}{1}%
\TextAt{13}{173.184}{670.292}{\Character{59}}%
\TextAt{13}{178.408}{670.292}{\Character{58}}%
\TextAt{13}{183.645}{670.292}{\Character{58}}%
\TextAt{13}{188.881}{670.292}{\Character{58}}%
\TextAt{13}{194.106}{670.292}{\Character{59}}%
\TextAt{13}{199.342}{670.292}{v}%
\TextAt{21}{205.064}{668.499}{m}%
\TextAt{8}{216.95}{670.292}{to}%
\TextAt{8}{231.261}{670.292}{a}%
\TextAt{8}{241.016}{670.292}{basis}%
\TextAt{13}{269.756}{670.292}{v}%
\TextAt{17}{275.416}{668.499}{1}%
\TextAt{13}{280.148}{670.292}{\Character{59}}%
\TextAt{13}{285.372}{670.292}{\Character{58}}%
\TextAt{13}{290.609}{670.292}{\Character{58}}%
\TextAt{13}{295.833}{670.292}{\Character{58}}%
\TextAt{13}{301.07}{670.292}{\Character{59}}%
\TextAt{13}{306.306}{670.292}{v}%
\TextAt{21}{312.028}{668.499}{n}%
\TextAt{8}{321.562}{670.292}{of}%
\TextAt{13}{334.892}{670.292}{V}%
\TextAt{8}{344.318}{670.292}{,}%
\TextAt{8}{351.467}{670.292}{the}%
\TextAt{8}{371.623}{670.292}{set}%
\TextAt{16}{156.633}{641.922}{\Character{102}}%
\TextAt{13}{162.61}{641.922}{x}%
\TextAt{17}{169.263}{640.129}{1}%
\TextAt{13}{173.995}{641.922}{v}%
\TextAt{17}{179.654}{640.129}{1}%
\TextAt{9}{187.043}{641.922}{\Character{43}}%
\TextAt{16}{198.804}{641.922}{\Character{1}}%
\TextAt{16}{204.1}{641.922}{\Character{1}}%
\TextAt{16}{209.408}{641.922}{\Character{1}}%
\TextAt{9}{215.409}{641.922}{\Character{43}}%
\TextAt{13}{227.17}{641.922}{x}%
\TextAt{21}{233.822}{640.129}{n}%
\TextAt{13}{239.459}{641.922}{v}%
\TextAt{21}{245.118}{640.129}{n}%
\TextAt{16}{254.075}{641.922}{\Character{106}}%
\TextAt{13}{260.717}{641.922}{x}%
\TextAt{21}{267.369}{640.129}{i}%
\NativeMathAt{274.071}{641.922}{11.955}{\in}%
\TextAt{18}{285.362}{641.922}{R}%
\TextAt{13}{293.996}{641.922}{\Character{59}}%
\TextAt{16}{303.138}{641.922}{\Character{106}}%
\TextAt{13}{306.459}{641.922}{a}%
\TextAt{21}{312.604}{640.129}{i}%
\TextAt{16}{318.641}{641.922}{\Character{0}}%
\TextAt{13}{330.597}{641.922}{x}%
\TextAt{21}{337.249}{640.129}{i}%
\TextAt{16}{340.63}{641.922}{\Character{106}}%
\TextAt{13}{347.272}{641.922}{\Character{60}}%
\TextAt{9}{359.697}{641.922}{1}%
\TextAt{8}{369.448}{641.922}{for}%
\TextAt{13}{387.332}{641.922}{i}%
\TextAt{9}{394.647}{641.922}{\Character{61}}%
\TextAt{9}{407.057}{641.922}{1}%
\TextAt{13}{412.925}{641.922}{\Character{59}}%
\TextAt{13}{418.149}{641.922}{\Character{58}}%
\TextAt{13}{423.386}{641.922}{\Character{58}}%
\TextAt{13}{428.622}{641.922}{\Character{58}}%
\TextAt{13}{433.847}{641.922}{\Character{59}}%
\TextAt{13}{439.083}{641.922}{m}%
\TextAt{16}{449.385}{641.922}{\Character{103}}%
\TextAt{8}{79.226}{613.553}{clearly}%
\TextAt{8}{116.622}{613.553}{is}%
\TextAt{8}{128.386}{613.553}{such}%
\TextAt{8}{154.783}{613.553}{a}%
\TextAt{8}{164.538}{613.553}{neighbourhood.}%
\TextAt{8}{249.093}{613.553}{This}%
\TextAt{8}{275.813}{613.553}{property}%
\TextAt{8}{323.622}{613.553}{is}%
\TextAt{8}{335.386}{613.553}{indeed}%
\TextAt{8}{372.447}{613.553}{characteristic.}%
\TextAt{12}{79.226}{587.721}{Proposition}%
\TextAt{12}{151.937}{587.721}{(4.2).}%
\TextAt{8}{187.253}{587.721}{A}%
\TextAt{8}{199.926}{587.721}{subgroup}%
\TextAt{9}{250.723}{587.721}{\Character{0}}%
\NativeMathAt{261.358}{587.721}{11.955}{\subseteq}%
\TextAt{13}{273.977}{587.721}{V}%
\TextAt{8}{287.3}{587.721}{is}%
\TextAt{8}{299.064}{587.721}{a}%
\TextAt{8}{308.82}{587.721}{lattice}%
\TextAt{8}{344.59}{587.721}{if}%
\TextAt{8}{355.313}{587.721}{and}%
\TextAt{8}{378.076}{587.721}{only}%
\TextAt{8}{403.768}{587.721}{if}%
\TextAt{8}{414.492}{587.721}{it}%
\TextAt{8}{426.196}{587.721}{is}%
\TextAt{8}{437.96}{587.721}{discrete.}%
\TextAt{19}{79.226}{553.103}{Proof.}%
\TextAt{8}{116}{553.103}{Let}%
\TextAt{9}{137.555}{553.103}{\Character{0}}%
\TextAt{8}{149.352}{553.103}{be}%
\TextAt{8}{165.861}{553.103}{a}%
\TextAt{8}{176.202}{553.103}{discrete}%
\TextAt{8}{219.766}{553.103}{subgroup}%
\TextAt{8}{271.149}{553.103}{of}%
\TextAt{13}{285.065}{553.103}{V}%
\TextAt{8}{294.49}{553.103}{.}%
\TextAt{8}{304.695}{553.103}{Then}%
\TextAt{9}{335.837}{553.103}{\Character{0}}%
\TextAt{8}{347.634}{553.103}{is}%
\TextAt{8}{359.984}{553.103}{closed.}%
\TextAt{8}{400.817}{553.103}{For}%
\TextAt{8}{422.372}{553.103}{let}%
\TextAt{13}{439.862}{553.103}{U}%
\TextAt{8}{453.543}{553.103}{be}%
\TextAt{8}{470.053}{553.103}{an}%
\TextAt{8}{486.898}{553.103}{arbitrary}%
\TextAt{8}{79.226}{536.057}{neighbourhood}%
\TextAt{8}{160.343}{536.057}{of}%
\TextAt{8}{174.786}{536.057}{0.}%
\TextAt{8}{192.43}{536.057}{Then}%
\TextAt{8}{224.1}{536.057}{there}%
\TextAt{8}{255.125}{536.057}{exists}%
\TextAt{8}{288.552}{536.057}{a}%
\TextAt{8}{299.42}{536.057}{neighbourhood}%
\TextAt{13}{380.537}{536.057}{U}%
\TextAt{3}{389.736}{540.395}{\Character{48}}%
\NativeMathAt{397.748}{536.057}{11.955}{\subseteq}%
\TextAt{13}{412.264}{536.057}{U}%
\TextAt{8}{426.473}{536.057}{of}%
\TextAt{8}{440.915}{536.057}{0}%
\TextAt{8}{451.783}{536.057}{such}%
\TextAt{8}{479.293}{536.057}{that}%
\TextAt{8}{505.775}{536.057}{every}%
\TextAt{8}{79.226}{519.011}{difference}%
\TextAt{8}{131.076}{519.011}{of}%
\TextAt{8}{143.916}{519.011}{elements}%
\TextAt{8}{191.282}{519.011}{of}%
\TextAt{13}{204.122}{519.011}{U}%
\TextAt{3}{213.321}{523.349}{\Character{48}}%
\TextAt{8}{219.523}{519.011}{lies}%
\TextAt{8}{239.249}{519.011}{in}%
\TextAt{13}{252.412}{519.011}{U}%
\TextAt{8}{261.611}{519.011}{.}%
\TextAt{8}{269.896}{519.011}{If}%
\TextAt{8}{281.098}{519.011}{there}%
\TextAt{8}{310.52}{519.011}{were}%
\TextAt{8}{337.012}{519.011}{an}%
\TextAt{13}{352.781}{519.011}{x}%
\NativeMathAt{362.754}{519.011}{11.955}{\notin}%
\TextAt{9}{374.045}{519.011}{\Character{0}}%
\TextAt{8}{384.767}{519.011}{belonging}%
\TextAt{8}{437.286}{519.011}{to}%
\TextAt{8}{451.106}{519.011}{the}%
\TextAt{8}{470.772}{519.011}{closure}%
\TextAt{8}{509.364}{519.011}{of}%
\TextAt{9}{522.204}{519.011}{\Character{0}}%
\TextAt{8}{529.518}{519.011}{,}%
\TextAt{8}{79.226}{501.965}{then}%
\TextAt{8}{105.974}{501.965}{we}%
\TextAt{8}{123.289}{501.965}{could}%
\TextAt{8}{154.592}{501.965}{find}%
\TextAt{8}{178.088}{501.965}{in}%
\TextAt{8}{191.829}{501.965}{the}%
\TextAt{8}{212.073}{501.965}{neighbourhood}%
\TextAt{13}{292.165}{501.965}{x}%
\TextAt{9}{301.534}{501.965}{\Character{43}}%
\TextAt{13}{313.355}{501.965}{U}%
\TextAt{3}{322.554}{506.303}{\Character{48}}%
\TextAt{8}{329.334}{501.965}{of}%
\TextAt{13}{342.752}{501.965}{x}%
\TextAt{8}{353.389}{501.965}{two}%
\TextAt{8}{375.594}{501.965}{distinct}%
\TextAt{8}{418.015}{501.965}{elements}%
\TextAt{13}{465.96}{501.965}{\Character{13}}%
\TextAt{17}{472.032}{500.172}{1}%
\TextAt{13}{476.764}{501.965}{\Character{59}}%
\TextAt{13}{481.988}{501.965}{\Character{13}}%
\TextAt{17}{488.08}{500.172}{2}%
\NativeMathAt{496.283}{501.965}{11.955}{\in}%
\TextAt{9}{507.724}{501.965}{\Character{0}}%
\TextAt{8}{515.038}{501.965}{,}%
\TextAt{8}{522.297}{501.965}{so}%
\TextAt{8}{79.226}{484.919}{that}%
\TextAt{9}{104.587}{484.919}{0}%
\NativeMathAt{113.761}{484.919}{11.955}{\neq}%
\TextAt{13}{126.187}{484.919}{\Character{13}}%
\TextAt{17}{132.259}{483.126}{1}%
\TextAt{16}{139.632}{484.919}{\Character{0}}%
\TextAt{13}{151.571}{484.919}{\Character{13}}%
\TextAt{17}{157.643}{483.126}{2}%
\NativeMathAt{165.696}{484.919}{11.955}{\in}%
\TextAt{13}{176.987}{484.919}{U}%
\TextAt{3}{186.186}{489.257}{\Character{48}}%
\TextAt{16}{191.622}{484.919}{\Character{0}}%
\TextAt{13}{203.562}{484.919}{U}%
\TextAt{3}{212.761}{489.257}{\Character{48}}%
\NativeMathAt{218.877}{484.919}{11.955}{\subseteq}%
\TextAt{13}{231.496}{484.919}{U}%
\TextAt{8}{240.695}{484.919}{.}%
\TextAt{8}{249.141}{484.919}{Thus}%
\TextAt{8}{278.782}{484.919}{0}%
\TextAt{8}{288.53}{484.919}{would}%
\TextAt{8}{322.666}{484.919}{not}%
\TextAt{8}{343.472}{484.919}{be}%
\TextAt{8}{359.389}{484.919}{an}%
\TextAt{8}{375.64}{484.919}{isolated}%
\TextAt{8}{418.623}{484.919}{point,}%
\TextAt{8}{452.439}{484.919}{a}%
\TextAt{8}{462.186}{484.919}{contradiction.}%
\TextAt{8}{97.159}{464.884}{Let}%
\TextAt{13}{118.799}{464.884}{V}%
\TextAt{17}{125.623}{463.091}{0}%
\TextAt{8}{134.923}{464.884}{be}%
\TextAt{8}{151.518}{464.884}{the}%
\TextAt{8}{172.345}{464.884}{linear}%
\TextAt{8}{205.534}{464.884}{subspace}%
\TextAt{8}{255.101}{464.884}{of}%
\TextAt{13}{269.101}{464.884}{V}%
\TextAt{8}{283.095}{464.884}{which}%
\TextAt{8}{317.252}{464.884}{is}%
\TextAt{8}{329.686}{464.884}{spanned}%
\TextAt{8}{375.942}{464.884}{by}%
\TextAt{8}{392.871}{464.884}{the}%
\TextAt{8}{413.698}{464.884}{set}%
\TextAt{9}{432.636}{464.884}{\Character{0}}%
\TextAt{8}{439.95}{464.884}{,}%
\TextAt{8}{447.938}{464.884}{and}%
\TextAt{8}{471.371}{464.884}{let}%
\TextAt{13}{488.946}{464.884}{m}%
\TextAt{8}{503.753}{464.884}{be}%
\TextAt{8}{520.348}{464.884}{its}%
\TextAt{8}{79.226}{447.838}{dimension.}%
\TextAt{8}{141.671}{447.838}{Then}%
\TextAt{8}{173.079}{447.838}{we}%
\TextAt{8}{191.158}{447.838}{may}%
\TextAt{8}{217.378}{447.838}{choose}%
\TextAt{8}{255.361}{447.838}{a}%
\TextAt{8}{265.968}{447.838}{basis}%
\TextAt{13}{295.559}{447.838}{u}%
\TextAt{17}{302.222}{446.045}{1}%
\TextAt{13}{306.954}{447.838}{\Character{59}}%
\TextAt{13}{312.178}{447.838}{\Character{58}}%
\TextAt{13}{317.415}{447.838}{\Character{58}}%
\TextAt{13}{322.639}{447.838}{\Character{58}}%
\TextAt{13}{327.876}{447.838}{\Character{59}}%
\TextAt{13}{333.112}{447.838}{u}%
\TextAt{21}{339.837}{446.045}{m}%
\TextAt{8}{352.574}{447.838}{of}%
\TextAt{13}{366.755}{447.838}{V}%
\TextAt{17}{373.58}{446.045}{0}%
\TextAt{8}{383.06}{447.838}{which}%
\TextAt{8}{417.398}{447.838}{is}%
\TextAt{8}{430.013}{447.838}{contained}%
\TextAt{8}{483.873}{447.838}{in}%
\TextAt{9}{498.377}{447.838}{\Character{0}}%
\TextAt{8}{505.691}{447.838}{,}%
\TextAt{8}{513.904}{447.838}{and}%
\TextAt{8}{79.226}{430.792}{form}%
\TextAt{8}{106.866}{430.792}{the}%
\TextAt{8}{127.023}{430.792}{complete}%
\TextAt{8}{176.446}{430.792}{lattice}%
\TextAt{9}{238.089}{413.746}{\Character{0}}%
\TextAt{17}{245.404}{411.953}{0}%
\TextAt{9}{253.457}{413.746}{\Character{61}}%
\TextAt{18}{265.882}{413.746}{Z}%
\TextAt{13}{273.852}{413.746}{u}%
\TextAt{17}{280.515}{411.953}{1}%
\TextAt{9}{287.904}{413.746}{\Character{43}}%
\TextAt{16}{299.665}{413.746}{\Character{1}}%
\TextAt{16}{304.961}{413.746}{\Character{1}}%
\TextAt{16}{310.269}{413.746}{\Character{1}}%
\TextAt{9}{316.269}{413.746}{\Character{43}}%
\TextAt{18}{328.031}{413.746}{Z}%
\TextAt{13}{336.001}{413.746}{u}%
\TextAt{21}{342.663}{411.953}{m}%
\NativeMathAt{353.973}{413.746}{11.955}{\subseteq}%
\TextAt{9}{366.592}{413.746}{\Character{0}}%
\TextAt{8}{79.226}{390.495}{of}%
\TextAt{13}{92.141}{390.495}{V}%
\TextAt{17}{98.965}{388.702}{0}%
\TextAt{8}{103.697}{390.495}{.}%
\TextAt{8}{112.007}{390.495}{We}%
\TextAt{8}{131.736}{390.495}{claim}%
\TextAt{8}{162.536}{390.495}{that}%
\TextAt{8}{187.489}{390.495}{the}%
\TextAt{8}{207.23}{390.495}{index}%
\TextAt{9}{238.352}{390.495}{\Character{40}\Character{0}}%
\TextAt{9}{253.511}{390.495}{\Character{58}}%
\TextAt{9}{260.075}{390.495}{\Character{0}}%
\TextAt{17}{267.426}{388.702}{0}%
\TextAt{9}{272.159}{390.495}{\Character{41}}%
\TextAt{8}{280.193}{390.495}{is}%
\TextAt{8}{291.541}{390.495}{finite.}%
\TextAt{8}{325.865}{390.495}{To}%
\TextAt{8}{342.677}{390.495}{see}%
\TextAt{8}{361.175}{390.495}{this,}%
\TextAt{8}{386.917}{390.495}{let}%
\TextAt{13}{403.406}{390.495}{\Character{13}}%
\TextAt{21}{409.478}{388.702}{i}%
\NativeMathAt{416.18}{390.495}{11.955}{\in}%
\TextAt{9}{427.471}{390.495}{\Character{0}}%
\TextAt{8}{438.267}{390.495}{vary}%
\TextAt{8}{463.878}{390.495}{over}%
\TextAt{8}{488.509}{390.495}{a}%
\TextAt{8}{497.848}{390.495}{system}%
\TextAt{8}{79.226}{373.449}{of}%
\TextAt{8}{92.25}{373.449}{representatives}%
\TextAt{8}{171.745}{373.449}{of}%
\TextAt{8}{184.769}{373.449}{the}%
\TextAt{8}{204.62}{373.449}{cosets}%
\TextAt{8}{238.255}{373.449}{in}%
\TextAt{9}{251.601}{373.449}{\Character{0}\Character{47}\Character{0}}%
\TextAt{17}{272.083}{371.656}{0}%
\TextAt{8}{276.815}{373.449}{.}%
\TextAt{8}{285.162}{373.449}{Since}%
\TextAt{9}{315.413}{373.449}{\Character{0}}%
\TextAt{17}{322.727}{371.656}{0}%
\TextAt{8}{331.051}{373.449}{is}%
\TextAt{8}{342.509}{373.449}{complete}%
\TextAt{8}{391.626}{373.449}{in}%
\TextAt{13}{404.973}{373.449}{V}%
\TextAt{17}{411.797}{371.656}{0}%
\TextAt{8}{416.529}{373.449}{,}%
\TextAt{8}{423.434}{373.449}{the}%
\TextAt{8}{443.284}{373.449}{translates}%
\TextAt{9}{496.442}{373.449}{\Character{8}}%
\TextAt{17}{504.896}{371.656}{0}%
\TextAt{9}{511.66}{373.449}{\Character{43}}%
\TextAt{13}{522.795}{373.449}{\Character{13}}%
\TextAt{8}{529.518}{373.449}{,}%
\TextAt{13}{79.226}{356.403}{\Character{13}}%
\NativeMathAt{89.269}{356.403}{11.955}{\in}%
\TextAt{9}{100.56}{356.403}{\Character{0}}%
\TextAt{17}{107.874}{354.61}{0}%
\TextAt{8}{112.607}{356.403}{,}%
\TextAt{8}{119.756}{356.403}{of}%
\TextAt{8}{133.086}{356.403}{the}%
\TextAt{8}{153.242}{356.403}{fundamental}%
\TextAt{8}{220.885}{356.403}{mesh}%
\TextAt{9}{185.222}{328.033}{\Character{8}}%
\TextAt{17}{193.676}{326.24}{0}%
\TextAt{9}{201.729}{328.033}{\Character{61}}%
\TextAt{16}{214.155}{328.033}{\Character{102}}%
\TextAt{13}{220.132}{328.033}{x}%
\TextAt{17}{226.784}{326.24}{1}%
\TextAt{13}{231.517}{328.033}{u}%
\TextAt{17}{238.179}{326.24}{1}%
\TextAt{9}{245.568}{328.033}{\Character{43}}%
\TextAt{16}{257.329}{328.033}{\Character{1}}%
\TextAt{16}{262.625}{328.033}{\Character{1}}%
\TextAt{16}{267.933}{328.033}{\Character{1}}%
\TextAt{9}{273.934}{328.033}{\Character{43}}%
\TextAt{13}{285.695}{328.033}{x}%
\TextAt{21}{292.347}{326.24}{m}%
\TextAt{13}{300.336}{328.033}{u}%
\TextAt{21}{306.998}{326.24}{m}%
\TextAt{16}{318.308}{328.033}{\Character{106}}%
\TextAt{13}{324.95}{328.033}{x}%
\TextAt{21}{331.602}{326.24}{i}%
\NativeMathAt{338.304}{328.033}{11.955}{\in}%
\TextAt{18}{349.595}{328.033}{R}%
\TextAt{13}{358.229}{328.033}{\Character{59}}%
\TextAt{9}{367.371}{328.033}{0}%
\NativeMathAt{376.544}{328.033}{11.955}{\leq}%
\TextAt{13}{389.164}{328.033}{x}%
\TextAt{21}{395.816}{326.24}{i}%
\TextAt{13}{402.518}{328.033}{\Character{60}}%
\TextAt{9}{414.943}{328.033}{1}%
\TextAt{16}{420.796}{328.033}{\Character{103}}%
\TextAt{8}{79.226}{299.664}{cover}%
\TextAt{8}{109.473}{299.664}{the}%
\TextAt{8}{129.629}{299.664}{entire}%
\TextAt{8}{162.47}{299.664}{space}%
\TextAt{13}{193.745}{299.664}{V}%
\TextAt{17}{200.569}{297.87}{0}%
\TextAt{8}{205.301}{299.664}{.}%
\TextAt{8}{213.75}{299.664}{We}%
\TextAt{8}{233.894}{299.664}{may}%
\TextAt{8}{259.263}{299.664}{therefore}%
\TextAt{8}{308.363}{299.664}{write}%
\TextAt{13}{196.298}{271.294}{\Character{13}}%
\TextAt{21}{202.37}{269.501}{i}%
\TextAt{9}{209.072}{271.294}{\Character{61}}%
\TextAt{13}{221.497}{271.294}{\Character{22}}%
\TextAt{21}{228.54}{269.501}{i}%
\TextAt{9}{234.578}{271.294}{\Character{43}}%
\TextAt{13}{246.34}{271.294}{\Character{13}}%
\TextAt{17}{252.412}{269.501}{0}%
\TextAt{21}{256.646}{269.501}{i}%
\TextAt{13}{260.027}{271.294}{\Character{59}}%
\TextAt{13}{289.162}{271.294}{\Character{22}}%
\TextAt{21}{296.225}{269.501}{i}%
\NativeMathAt{302.927}{271.294}{11.955}{\in}%
\TextAt{9}{314.218}{271.294}{\Character{8}}%
\TextAt{17}{322.672}{269.501}{0}%
\TextAt{13}{327.404}{271.294}{\Character{59}}%
\TextAt{13}{344.584}{271.294}{\Character{13}}%
\TextAt{17}{350.675}{269.501}{0}%
\TextAt{21}{354.91}{269.501}{i}%
\NativeMathAt{361.612}{271.294}{11.955}{\in}%
\TextAt{9}{372.903}{271.294}{\Character{0}}%
\TextAt{17}{380.217}{269.501}{0}%
\NativeMathAt{388.27}{271.294}{11.955}{\subseteq}%
\TextAt{13}{400.889}{271.294}{V}%
\TextAt{17}{407.714}{269.501}{0}%
\TextAt{13}{412.446}{271.294}{\Character{58}}%
\TextAt{8}{79.226}{242.924}{As}%
\TextAt{8}{97.149}{242.924}{the}%
\TextAt{13}{117.941}{242.924}{\Character{22}}%
\TextAt{21}{124.984}{241.131}{i}%
\TextAt{9}{132.77}{242.924}{\Character{61}}%
\TextAt{13}{146.279}{242.924}{\Character{13}}%
\TextAt{21}{152.351}{241.131}{i}%
\TextAt{16}{158.822}{242.924}{\Character{0}}%
\TextAt{13}{171.211}{242.924}{\Character{13}}%
\TextAt{17}{177.283}{241.131}{0}%
\TextAt{21}{181.517}{241.131}{i}%
\NativeMathAt{189.303}{242.924}{11.955}{\in}%
\TextAt{9}{201.677}{242.924}{\Character{0}}%
\TextAt{8}{213.524}{242.924}{lie}%
\TextAt{8}{229.762}{242.924}{discretely}%
\TextAt{8}{282.809}{242.924}{in}%
\TextAt{8}{297.098}{242.924}{the}%
\TextAt{8}{317.89}{242.924}{bounded}%
\TextAt{8}{366.323}{242.924}{set}%
\TextAt{9}{385.226}{242.924}{\Character{8}}%
\TextAt{17}{393.68}{241.131}{0}%
\TextAt{8}{398.412}{242.924}{,}%
\TextAt{8}{406.356}{242.924}{they}%
\TextAt{8}{433.33}{242.924}{have}%
\TextAt{8}{460.96}{242.924}{to}%
\TextAt{8}{475.906}{242.924}{be}%
\TextAt{8}{492.466}{242.924}{finite}%
\TextAt{8}{523.014}{242.924}{in}%
\TextAt{8}{79.226}{225.878}{number.}%
\TextAt{8}{129.489}{225.878}{In}%
\TextAt{8}{145.149}{225.878}{fact,}%
\TextAt{8}{172.784}{225.878}{the}%
\TextAt{8}{193.979}{225.878}{intersection}%
\TextAt{8}{257.842}{225.878}{of}%
\TextAt{9}{272.211}{225.878}{\Character{0}}%
\TextAt{8}{284.461}{225.878}{with}%
\TextAt{8}{312.159}{225.878}{the}%
\TextAt{8}{333.354}{225.878}{closure}%
\TextAt{8}{373.474}{225.878}{of}%
\TextAt{9}{387.843}{225.878}{\Character{8}}%
\TextAt{17}{396.297}{224.085}{0}%
\TextAt{8}{405.966}{225.878}{is}%
\TextAt{8}{418.768}{225.878}{compact}%
\TextAt{8}{466.635}{225.878}{and}%
\TextAt{8}{490.436}{225.878}{discrete,}%
\TextAt{8}{79.226}{208.832}{hence}%
\TextAt{8}{111.732}{208.832}{finite.}%
\TextAt{8}{97.159}{188.797}{Putting}%
\TextAt{8}{140.245}{188.797}{now}%
\TextAt{13}{164.634}{188.797}{q}%
\TextAt{9}{173.574}{188.797}{\Character{61}}%
\TextAt{9}{185.984}{188.797}{\Character{40}\Character{0}}%
\TextAt{9}{201.155}{188.797}{\Character{58}}%
\TextAt{9}{207.718}{188.797}{\Character{0}}%
\TextAt{17}{215.073}{187.004}{0}%
\TextAt{9}{219.806}{188.797}{\Character{41}}%
\TextAt{8}{224.358}{188.797}{,}%
\TextAt{8}{231.507}{188.797}{we}%
\TextAt{8}{248.735}{188.797}{have}%
\TextAt{13}{275.729}{188.797}{q}%
\TextAt{9}{281.349}{188.797}{\Character{0}}%
\NativeMathAt{291.984}{188.797}{11.955}{\subseteq}%
\TextAt{9}{304.603}{188.797}{\Character{0}}%
\TextAt{17}{311.917}{187.004}{0}%
\TextAt{8}{316.649}{188.797}{,}%
\TextAt{8}{323.799}{188.797}{whence}%
\TextAt{9}{203.508}{156.153}{\Character{0}}%
\NativeMathAt{214.143}{156.153}{11.955}{\subseteq}%
\TextAt{9}{227.958}{164.241}{1}%
\Vector{q 0 G 0 g 0.478 w 155.958 -614.858 m 161.811 -614.858 l S Q}%
\TextAt{13}{228.075}{147.953}{q}%
\TextAt{9}{235.006}{156.153}{\Character{0}}%
\TextAt{17}{242.321}{154.36}{0}%
\TextAt{9}{250.374}{156.153}{\Character{61}}%
\TextAt{18}{262.799}{156.153}{Z}%
\TextAt{20}{272.762}{173.01}{\Character{18}}%
\TextAt{9}{282.758}{164.241}{1}%
\Vector{q 0 G 0 g 0.478 w 210.758 -614.858 m 216.611 -614.858 l S Q}%
\TextAt{13}{282.875}{147.953}{q}%
\TextAt{13}{289.806}{156.153}{u}%
\TextAt{17}{296.469}{154.36}{1}%
\TextAt{20}{301.201}{173.01}{\Character{19}}%
\TextAt{9}{312.658}{156.153}{\Character{43}}%
\TextAt{16}{324.419}{156.153}{\Character{1}}%
\TextAt{16}{329.715}{156.153}{\Character{1}}%
\TextAt{16}{335.023}{156.153}{\Character{1}}%
\TextAt{9}{341.024}{156.153}{\Character{43}}%
\TextAt{18}{352.785}{156.153}{Z}%
\TextAt{20}{362.748}{173.01}{\Character{18}}%
\TextAt{9}{372.744}{164.241}{1}%
\Vector{q 0 G 0 g 0.478 w 300.743 -614.858 m 306.596 -614.858 l S Q}%
\TextAt{13}{372.86}{147.953}{q}%
\TextAt{13}{379.792}{156.153}{u}%
\TextAt{21}{386.454}{154.36}{m}%
\TextAt{20}{394.443}{173.01}{\Character{19}}%
\TextAt{13}{405.236}{156.153}{\Character{58}}%
\TextAt{8}{79.226}{124.059}{By}%
\TextAt{8}{97.921}{124.059}{the}%
\TextAt{8}{118.409}{124.059}{main}%
\TextAt{8}{148.007}{124.059}{theorem}%
\TextAt{8}{193.864}{124.059}{on}%
\TextAt{8}{210.454}{124.059}{finitely}%
\TextAt{8}{250.13}{124.059}{generated}%
\TextAt{8}{303.794}{124.059}{abelian}%
\TextAt{8}{344.773}{124.059}{groups,}%
\TextAt{9}{386.23}{124.059}{\Character{0}}%
\TextAt{8}{397.773}{124.059}{therefore}%
\TextAt{8}{447.205}{124.059}{admits}%
\TextAt{8}{485.972}{124.059}{a}%
\TextAt{18}{496.059}{124.059}{Z}%
\TextAt{8}{504.029}{124.059}{-basis}%
\TextAt{13}{79.226}{107.013}{v}%
\TextAt{17}{84.885}{105.22}{1}%
\TextAt{13}{89.618}{107.013}{\Character{59}}%
\TextAt{13}{94.842}{107.013}{\Character{58}}%
\TextAt{13}{100.079}{107.013}{\Character{58}}%
\TextAt{13}{105.303}{107.013}{\Character{58}}%
\TextAt{13}{110.54}{107.013}{\Character{59}}%
\TextAt{13}{115.776}{107.013}{v}%
\TextAt{21}{121.498}{105.22}{r}%
\TextAt{8}{126.053}{107.013}{,}%
\TextAt{13}{134.353}{107.013}{r}%
\NativeMathAt{144.843}{107.013}{11.955}{\leq}%
\TextAt{13}{159.031}{107.013}{m}%
\TextAt{8}{169.27}{107.013}{,}%
\TextAt{8}{177.57}{107.013}{i.e.,}%
\TextAt{9}{200.826}{107.013}{\Character{0}}%
\TextAt{9}{213.02}{107.013}{\Character{61}}%
\TextAt{18}{227.024}{107.013}{Z}%
\TextAt{13}{234.994}{107.013}{v}%
\TextAt{17}{240.654}{105.22}{1}%
\TextAt{9}{248.67}{107.013}{\Character{43}}%
\TextAt{16}{261.059}{107.013}{\Character{1}}%
\TextAt{16}{266.355}{107.013}{\Character{1}}%
\TextAt{16}{271.663}{107.013}{\Character{1}}%
\TextAt{9}{278.291}{107.013}{\Character{43}}%
\TextAt{18}{290.68}{107.013}{Z}%
\TextAt{13}{298.65}{107.013}{v}%
\TextAt{21}{304.309}{105.22}{r}%
\TextAt{8}{308.864}{107.013}{.}%
\TextAt{8}{320.074}{107.013}{The}%
\TextAt{8}{345.049}{107.013}{vectors}%
\TextAt{13}{385.708}{107.013}{v}%
\TextAt{17}{391.367}{105.22}{1}%
\TextAt{13}{396.1}{107.013}{\Character{59}}%
\TextAt{13}{401.324}{107.013}{\Character{58}}%
\TextAt{13}{406.561}{107.013}{\Character{58}}%
\TextAt{13}{411.785}{107.013}{\Character{58}}%
\TextAt{13}{417.022}{107.013}{\Character{59}}%
\TextAt{13}{422.258}{107.013}{v}%
\TextAt{21}{427.98}{105.22}{r}%
\TextAt{8}{437.353}{107.013}{are}%
\TextAt{8}{457.784}{107.013}{also}%
\TextAt{18}{482.185}{107.013}{R}%
\TextAt{8}{490.819}{107.013}{-linearly}%
\TextAt{8}{79.226}{89.967}{independent}%
\TextAt{8}{147.263}{89.967}{because}%
\TextAt{8}{192.273}{89.967}{they}%
\TextAt{8}{220.32}{89.967}{span}%
\TextAt{8}{249.406}{89.967}{the}%
\TextAt{13}{271.272}{89.967}{m}%
\TextAt{8}{281.511}{89.967}{-dimensional}%
\TextAt{8}{351.568}{89.967}{space}%
\TextAt{13}{384.552}{89.967}{V}%
\TextAt{17}{391.376}{88.174}{0}%
\TextAt{8}{396.108}{89.967}{.}%
\TextAt{8}{409.684}{89.967}{This}%
\TextAt{8}{438.113}{89.967}{shows}%
\TextAt{8}{473.44}{89.967}{that}%
\TextAt{9}{500.518}{89.967}{\Character{0}}%
\TextAt{8}{513.438}{89.967}{is}%
\TextAt{8}{526.911}{89.967}{a}%
\TextAt{8}{79.226}{72.921}{lattice.}%
\Vector{q 0 G 0 g 0.399 w 452.334 -701.079 m 452.334 -693.009 l S Q}%
\Vector{q 0 G 0 g 0.399 w 452.534 -693.208 m 459.707 -693.208 l S Q}%
\Vector{q 0 G 0 g 0.399 w 452.534 -700.879 m 459.707 -700.879 l S Q}%
\Vector{q 0 G 0 g 0.399 w 459.906 -701.079 m 459.906 -693.009 l S Q}%
\TextAt{8}{300.14}{43.033}{23}%
\Destination{78.23}{809.815}{nav.page.23}%
\Destination{79.36}{568.055}{nav.print.25}%
\Destination{79.36}{601.755}{nav.number.1.4.2}%
}

% Page 36 (body page 24)
\Page{%
\TextAt{2}{79.226}{796.943}{I.}%
\TextAt{2}{89.83}{796.943}{Algebraic}%
\TextAt{2}{138.648}{796.943}{Integers}%
\Vector{q 0 G 0 g 0.299 w 7.226 18.729 m 460.769 18.729 l S Q}%
\TextAt{8}{97.159}{761.124}{Next}%
\TextAt{8}{125.524}{761.124}{we}%
\TextAt{8}{142.509}{761.124}{prove}%
\TextAt{8}{173.815}{761.124}{a}%
\TextAt{8}{183.327}{761.124}{criterion}%
\TextAt{8}{229.913}{761.124}{which}%
\TextAt{8}{263.156}{761.124}{will}%
\TextAt{8}{285.018}{761.124}{tell}%
\TextAt{8}{304.931}{761.124}{us}%
\TextAt{8}{319.703}{761.124}{when}%
\TextAt{8}{350.018}{761.124}{a}%
\TextAt{8}{359.53}{761.124}{lattice}%
\TextAt{8}{395.056}{761.124}{in}%
\TextAt{8}{408.466}{761.124}{the}%
\TextAt{8}{428.379}{761.124}{space}%
\TextAt{13}{459.411}{761.124}{V}%
\TextAt{8}{468.837}{761.124}{—given,}%
\TextAt{8}{514.167}{761.124}{say,}%
\TextAt{8}{79.226}{744.078}{as}%
\TextAt{8}{93.596}{744.078}{a}%
\TextAt{8}{103.352}{744.078}{discrete}%
\TextAt{8}{146.331}{744.078}{subgroup}%
\TextAt{9}{197.128}{744.078}{\Character{0}}%
\NativeMathAt{207.763}{744.078}{11.955}{\subseteq}%
\TextAt{13}{220.382}{744.078}{V}%
\TextAt{8}{229.808}{744.078}{—is}%
\TextAt{8}{253.276}{744.078}{complete.}%
\TextAt{12}{79.226}{718.435}{Lemma}%
\TextAt{12}{126.963}{718.435}{(4.3).}%
\TextAt{8}{162.342}{718.435}{A}%
\TextAt{8}{175.077}{718.435}{lattice}%
\TextAt{9}{210.909}{718.435}{\Character{0}}%
\TextAt{8}{222.183}{718.435}{in}%
\TextAt{13}{235.899}{718.435}{V}%
\TextAt{8}{249.284}{718.435}{is}%
\TextAt{8}{261.111}{718.435}{complete}%
\TextAt{8}{310.596}{718.435}{if}%
\TextAt{8}{321.382}{718.435}{and}%
\TextAt{8}{344.208}{718.435}{only}%
\TextAt{8}{369.962}{718.435}{if}%
\TextAt{8}{380.748}{718.435}{there}%
\TextAt{8}{410.723}{718.435}{exists}%
\TextAt{8}{443.1}{718.435}{a}%
\TextAt{8}{452.918}{718.435}{bounded}%
\TextAt{8}{500.777}{718.435}{subset}%
\TextAt{13}{79.226}{701.389}{M}%
\NativeMathAt{95.121}{701.389}{11.955}{\subseteq}%
\TextAt{13}{107.74}{701.389}{V}%
\TextAt{8}{120.85}{701.389}{such}%
\TextAt{8}{147.034}{701.389}{that}%
\TextAt{8}{172.19}{701.389}{the}%
\TextAt{8}{192.133}{701.389}{collection}%
\TextAt{8}{243.949}{701.389}{of}%
\TextAt{8}{257.066}{701.389}{all}%
\TextAt{8}{273.112}{701.389}{translates}%
\TextAt{13}{326.363}{701.389}{M}%
\TextAt{9}{341.157}{701.389}{\Character{43}}%
\TextAt{13}{352.483}{701.389}{\Character{13}}%
\TextAt{8}{359.205}{701.389}{,}%
\TextAt{13}{366.184}{701.389}{\Character{13}}%
\NativeMathAt{376.227}{701.389}{11.955}{\in}%
\TextAt{9}{387.518}{701.389}{\Character{0}}%
\TextAt{8}{394.832}{701.389}{,}%
\TextAt{8}{401.811}{701.389}{covers}%
\TextAt{8}{436.459}{701.389}{the}%
\TextAt{8}{456.403}{701.389}{entire}%
\TextAt{8}{489.03}{701.389}{space}%
\TextAt{13}{520.092}{701.389}{V}%
\TextAt{8}{529.518}{701.389}{.}%
\TextAt{19}{79.226}{667.15}{Proof.}%
\TextAt{8}{116}{667.15}{If}%
\TextAt{9}{127.899}{667.15}{\Character{0}}%
\TextAt{9}{138.874}{667.15}{\Character{61}}%
\TextAt{18}{151.665}{667.15}{Z}%
\TextAt{13}{159.635}{667.15}{v}%
\TextAt{17}{165.295}{665.356}{1}%
\TextAt{9}{172.825}{667.15}{\Character{43}}%
\TextAt{16}{184.727}{667.15}{\Character{1}}%
\TextAt{16}{190.023}{667.15}{\Character{1}}%
\TextAt{16}{195.331}{667.15}{\Character{1}}%
\TextAt{9}{201.473}{667.15}{\Character{43}}%
\TextAt{18}{213.375}{667.15}{Z}%
\TextAt{13}{221.345}{667.15}{v}%
\TextAt{21}{227.004}{665.356}{n}%
\TextAt{8}{236.745}{667.15}{is}%
\TextAt{8}{248.716}{667.15}{complete,}%
\TextAt{8}{301.649}{667.15}{then}%
\TextAt{8}{328.516}{667.15}{one}%
\TextAt{8}{350.183}{667.15}{may}%
\TextAt{8}{375.759}{667.15}{take}%
\TextAt{13}{401.335}{667.15}{M}%
\TextAt{8}{418.013}{667.15}{to}%
\TextAt{8}{432.53}{667.15}{be}%
\TextAt{8}{448.661}{667.15}{the}%
\TextAt{8}{469.025}{667.15}{fundamental}%
\TextAt{8}{79.226}{650.104}{mesh}%
\TextAt{9}{109.198}{650.104}{\Character{8}}%
\TextAt{9}{120.962}{650.104}{\Character{61}}%
\TextAt{16}{133.398}{650.104}{\Character{102}}%
\TextAt{13}{139.376}{650.104}{x}%
\TextAt{17}{146.028}{648.31}{1}%
\TextAt{13}{150.76}{650.104}{v}%
\TextAt{17}{156.42}{648.31}{1}%
\TextAt{9}{163.809}{650.104}{\Character{43}}%
\TextAt{16}{175.57}{650.104}{\Character{1}}%
\TextAt{16}{180.866}{650.104}{\Character{1}}%
\TextAt{16}{186.174}{650.104}{\Character{1}}%
\TextAt{9}{192.174}{650.104}{\Character{43}}%
\TextAt{13}{203.936}{650.104}{x}%
\TextAt{21}{210.588}{648.31}{n}%
\TextAt{13}{216.224}{650.104}{v}%
\TextAt{21}{221.883}{648.31}{n}%
\TextAt{16}{230.84}{650.104}{\Character{106}}%
\TextAt{9}{237.482}{650.104}{0}%
\NativeMathAt{246.656}{650.104}{11.955}{\leq}%
\TextAt{13}{259.275}{650.104}{x}%
\TextAt{21}{265.927}{648.31}{i}%
\TextAt{13}{272.63}{650.104}{\Character{60}}%
\TextAt{9}{285.055}{650.104}{1}%
\TextAt{16}{290.908}{650.104}{\Character{103}}%
\TextAt{8}{296.886}{650.104}{.}%
\TextAt{8}{97.159}{630.069}{Conversely,}%
\TextAt{8}{158.489}{630.069}{let}%
\TextAt{13}{175.154}{630.069}{M}%
\TextAt{8}{191.385}{630.069}{be}%
\TextAt{8}{207.07}{630.069}{a}%
\TextAt{8}{216.586}{630.069}{bounded}%
\TextAt{8}{264.143}{630.069}{subset}%
\TextAt{8}{299.793}{630.069}{of}%
\TextAt{13}{312.883}{630.069}{V}%
\TextAt{8}{325.967}{630.069}{whose}%
\TextAt{8}{360.254}{630.069}{translates}%
\TextAt{13}{413.478}{630.069}{M}%
\TextAt{9}{428.218}{630.069}{\Character{43}}%
\TextAt{13}{439.489}{630.069}{\Character{13}}%
\TextAt{8}{446.212}{630.069}{,}%
\TextAt{8}{453.169}{630.069}{for}%
\TextAt{13}{470.815}{630.069}{\Character{13}}%
\NativeMathAt{480.858}{630.069}{11.955}{\in}%
\TextAt{9}{492.149}{630.069}{\Character{0}}%
\TextAt{8}{499.463}{630.069}{,}%
\TextAt{8}{506.42}{630.069}{cover}%
\TextAt{13}{79.226}{613.023}{V}%
\TextAt{8}{88.652}{613.023}{.}%
\TextAt{8}{98.018}{613.023}{Let}%
\TextAt{13}{119.293}{613.023}{V}%
\TextAt{17}{126.118}{611.23}{0}%
\TextAt{8}{135.053}{613.023}{be}%
\TextAt{8}{151.284}{613.023}{the}%
\TextAt{8}{171.746}{613.023}{subspace}%
\TextAt{8}{220.948}{613.023}{spanned}%
\TextAt{8}{266.84}{613.023}{by}%
\TextAt{9}{283.404}{613.023}{\Character{0}}%
\TextAt{8}{290.719}{613.023}{.}%
\TextAt{8}{300.085}{613.023}{We}%
\TextAt{8}{320.535}{613.023}{have}%
\TextAt{8}{347.836}{613.023}{to}%
\TextAt{8}{362.452}{613.023}{show}%
\TextAt{8}{391.761}{613.023}{that}%
\TextAt{13}{417.436}{613.023}{V}%
\TextAt{9}{430.704}{613.023}{\Character{61}}%
\TextAt{13}{443.651}{613.023}{V}%
\TextAt{17}{450.475}{611.23}{0}%
\TextAt{8}{455.208}{613.023}{.}%
\TextAt{8}{464.574}{613.023}{So}%
\TextAt{8}{481.139}{613.023}{let}%
\TextAt{13}{498.349}{613.023}{v}%
\NativeMathAt{508.28}{613.023}{11.955}{\in}%
\TextAt{13}{520.092}{613.023}{V}%
\TextAt{8}{529.518}{613.023}{.}%
\TextAt{8}{79.226}{595.977}{Since}%
\TextAt{13}{109.784}{595.977}{V}%
\TextAt{9}{122.53}{595.977}{\Character{61}}%
\TextAt{20}{134.956}{604.943}{\Character{83}}%
\TextAt{21}{144.918}{592.49}{\Character{13}}%
\NativeMathAt{149.75}{592.49}{7.97}{\in}%
\TextAt{17}{155.396}{592.49}{\Character{0}}%
\TextAt{9}{161.183}{595.977}{\Character{40}}%
\TextAt{13}{165.735}{595.977}{M}%
\TextAt{9}{180.965}{595.977}{\Character{43}}%
\TextAt{13}{192.727}{595.977}{\Character{13}}%
\TextAt{9}{199.449}{595.977}{\Character{41}}%
\TextAt{8}{207.898}{595.977}{we}%
\TextAt{8}{225.126}{595.977}{may}%
\TextAt{8}{250.495}{595.977}{write,}%
\TextAt{8}{283.658}{595.977}{for}%
\TextAt{8}{301.543}{595.977}{each}%
\TextAt{13}{327.881}{595.977}{\Character{23}}%
\NativeMathAt{337.739}{595.977}{11.955}{\in}%
\TextAt{18}{349.03}{595.977}{N}%
\TextAt{8}{357.665}{595.977}{,}%
\TextAt{13}{212.988}{568.27}{\Character{23}}%
\TextAt{13}{219.516}{568.27}{v}%
\TextAt{9}{228.935}{568.27}{\Character{61}}%
\TextAt{13}{241.361}{568.27}{a}%
\TextAt{21}{247.506}{566.476}{\Character{23}}%
\TextAt{9}{255.345}{568.27}{\Character{43}}%
\TextAt{13}{267.106}{568.27}{\Character{13}}%
\TextAt{21}{273.178}{566.476}{\Character{23}}%
\TextAt{13}{278.36}{568.27}{\Character{59}}%
\TextAt{13}{307.495}{568.27}{a}%
\TextAt{21}{313.659}{566.476}{\Character{23}}%
\NativeMathAt{322.162}{568.27}{11.955}{\in}%
\TextAt{13}{333.453}{568.27}{M}%
\TextAt{13}{345.372}{568.27}{\Character{59}}%
\TextAt{13}{362.552}{568.27}{\Character{13}}%
\TextAt{21}{368.648}{566.476}{\Character{23}}%
\NativeMathAt{377.15}{568.27}{11.955}{\in}%
\TextAt{9}{388.441}{568.27}{\Character{0}}%
\TextAt{13}{395.756}{568.27}{\Character{58}}%
\TextAt{8}{79.226}{540.562}{Since}%
\TextAt{13}{109.784}{540.562}{M}%
\TextAt{8}{126.255}{540.562}{is}%
\TextAt{8}{138.018}{540.562}{bounded,}%
\TextAt{13}{189.067}{540.562}{\Character{23}}%
\TextAt{3}{195.605}{544.901}{\Character{0}}%
\TextAt{17}{202.191}{544.901}{1}%
\TextAt{13}{206.924}{540.562}{a}%
\TextAt{21}{213.069}{538.769}{\Character{23}}%
\TextAt{8}{222.148}{540.562}{converges}%
\TextAt{8}{274.572}{540.562}{to}%
\TextAt{8}{288.882}{540.562}{zero,}%
\TextAt{8}{316.845}{540.562}{and}%
\TextAt{8}{339.608}{540.562}{since}%
\TextAt{13}{368.276}{540.562}{V}%
\TextAt{17}{375.101}{538.769}{0}%
\TextAt{8}{383.73}{540.562}{is}%
\TextAt{8}{395.494}{540.562}{closed,}%
\TextAt{13}{199.558}{510.662}{v}%
\TextAt{9}{208.967}{510.662}{\Character{61}}%
\TextAt{8}{224.073}{510.662}{lim}%
\TextAt{21}{221.392}{503.489}{\Character{23}}%
\NativeMathAt{226.076}{503.489}{7.97}{\rightarrow}%
\NativeMathAt{234.54}{503.489}{7.97}{\infty}%
\TextAt{9}{246.543}{518.75}{1}%
\Vector{q 0 G 0 g 0.478 w 174.201 -260.349 m 180.739 -260.349 l S Q}%
\TextAt{13}{246.201}{502.461}{\Character{23}}%
\TextAt{13}{253.934}{510.662}{a}%
\TextAt{21}{260.079}{508.869}{\Character{23}}%
\TextAt{9}{267.918}{510.662}{\Character{43}}%
\TextAt{8}{282.36}{510.662}{lim}%
\TextAt{21}{279.679}{503.489}{\Character{23}}%
\NativeMathAt{284.363}{503.489}{7.97}{\rightarrow}%
\NativeMathAt{292.827}{503.489}{7.97}{\infty}%
\TextAt{9}{304.83}{518.75}{1}%
\Vector{q 0 G 0 g 0.478 w 232.488 -260.349 m 239.026 -260.349 l S Q}%
\TextAt{13}{304.488}{502.461}{\Character{23}}%
\TextAt{13}{312.221}{510.662}{\Character{13}}%
\TextAt{21}{318.293}{508.869}{\Character{23}}%
\TextAt{9}{326.796}{510.662}{\Character{61}}%
\TextAt{8}{341.902}{510.662}{lim}%
\TextAt{21}{339.222}{503.489}{\Character{23}}%
\NativeMathAt{343.906}{503.489}{7.97}{\rightarrow}%
\NativeMathAt{352.37}{503.489}{7.97}{\infty}%
\TextAt{9}{364.373}{518.75}{1}%
\Vector{q 0 G 0 g 0.478 w 292.03 -260.349 m 298.568 -260.349 l S Q}%
\TextAt{13}{364.03}{502.461}{\Character{23}}%
\TextAt{13}{371.764}{510.662}{\Character{13}}%
\TextAt{21}{377.836}{508.869}{\Character{23}}%
\NativeMathAt{386.338}{510.662}{11.955}{\in}%
\TextAt{13}{397.629}{510.662}{V}%
\TextAt{17}{404.454}{508.869}{0}%
\TextAt{13}{409.186}{510.662}{\Character{58}}%
\Vector{q 0 G 0 g 0.399 w 452.334 -291.046 m 452.334 -282.976 l S Q}%
\Vector{q 0 G 0 g 0.399 w 452.534 -283.175 m 459.707 -283.175 l S Q}%
\Vector{q 0 G 0 g 0.399 w 452.534 -290.846 m 459.707 -290.846 l S Q}%
\Vector{q 0 G 0 g 0.399 w 459.906 -291.046 m 459.906 -282.976 l S Q}%
\TextAt{8}{97.159}{454.323}{Now}%
\TextAt{8}{123.99}{454.323}{let}%
\TextAt{13}{141.065}{454.323}{V}%
\TextAt{8}{154.559}{454.323}{be}%
\TextAt{8}{170.654}{454.323}{a}%
\TextAt{12}{180.58}{454.323}{euclidean}%
\TextAt{12}{239.942}{454.323}{vector}%
\TextAt{12}{280.93}{454.323}{space}%
\TextAt{8}{312.36}{454.323}{,}%
\TextAt{8}{319.723}{454.323}{i.e.,}%
\TextAt{8}{342.041}{454.323}{an}%
\TextAt{18}{358.471}{454.323}{R}%
\TextAt{8}{367.105}{454.323}{-vector}%
\TextAt{8}{406.298}{454.323}{space}%
\TextAt{8}{437.743}{454.323}{of}%
\TextAt{8}{451.244}{454.323}{finite}%
\TextAt{8}{481.326}{454.323}{dimension}%
\TextAt{13}{79.226}{437.277}{n}%
\TextAt{8}{90.111}{437.277}{equipped}%
\TextAt{8}{140.179}{437.277}{with}%
\TextAt{8}{166.839}{437.277}{a}%
\TextAt{8}{176.595}{437.277}{symmetric,}%
\TextAt{8}{236.813}{437.277}{positive}%
\TextAt{8}{280.127}{437.277}{definite}%
\TextAt{8}{321.743}{437.277}{bilinear}%
\TextAt{8}{364.016}{437.277}{form}%
\TextAt{9}{254.353}{409.57}{\Character{40}}%
\TextAt{13}{262.803}{409.57}{\Character{59}}%
\TextAt{9}{271.945}{409.57}{\Character{41}}%
\TextAt{9}{279.8}{409.57}{\Character{58}}%
\TextAt{13}{286.39}{409.57}{V}%
\TextAt{16}{298.473}{409.57}{\Character{2}}%
\TextAt{13}{310.428}{409.57}{V}%
\NativeArrowAt{323.174}{409.57}{11.955}{19.26}{\longrightarrow}%
\TextAt{18}{345.756}{409.57}{R}%
\TextAt{13}{354.391}{409.57}{\Character{58}}%
\TextAt{8}{79.226}{381.862}{Then}%
\TextAt{8}{108.542}{381.862}{we}%
\TextAt{8}{124.529}{381.862}{have}%
\TextAt{8}{150.282}{381.862}{on}%
\TextAt{13}{165.3}{381.862}{V}%
\TextAt{8}{177.382}{381.862}{a}%
\TextAt{8}{185.896}{381.862}{notion}%
\TextAt{8}{221.082}{381.862}{of}%
\TextAt{8}{233.171}{381.862}{volume—more}%
\TextAt{8}{309.328}{381.862}{precisely}%
\TextAt{8}{355.943}{381.862}{a}%
\TextAt{8}{364.457}{381.862}{Haar}%
\TextAt{8}{392.16}{381.862}{measure.}%
\TextAt{8}{441.882}{381.862}{The}%
\TextAt{8}{464.695}{381.862}{cube}%
\TextAt{8}{491.082}{381.862}{spanned}%
\TextAt{8}{79.226}{364.816}{by}%
\TextAt{8}{95.829}{364.816}{an}%
\TextAt{8}{112.432}{364.816}{orthonormal}%
\TextAt{8}{179.785}{364.816}{basis}%
\TextAt{13}{208.87}{364.816}{e}%
\TextAt{17}{214.295}{363.023}{1}%
\TextAt{13}{219.027}{364.816}{\Character{59}}%
\TextAt{13}{224.251}{364.816}{\Character{58}}%
\TextAt{13}{229.488}{364.816}{\Character{58}}%
\TextAt{13}{234.724}{364.816}{\Character{58}}%
\TextAt{13}{239.949}{364.816}{\Character{59}}%
\TextAt{13}{245.185}{364.816}{e}%
\TextAt{21}{250.674}{363.023}{n}%
\TextAt{8}{260.552}{364.816}{has}%
\TextAt{8}{281.769}{364.816}{volume}%
\TextAt{8}{322.438}{364.816}{1,}%
\TextAt{8}{335.876}{364.816}{and}%
\TextAt{8}{358.983}{364.816}{more}%
\TextAt{8}{388.593}{364.816}{generally,}%
\TextAt{8}{441.052}{364.816}{the}%
\TextAt{8}{461.553}{364.816}{parallelepiped}%
\TextAt{8}{79.226}{347.77}{spanned}%
\TextAt{8}{124.811}{347.77}{by}%
\TextAt{13}{141.07}{347.77}{n}%
\TextAt{8}{151.955}{347.77}{linearly}%
\TextAt{8}{193.906}{347.77}{independent}%
\TextAt{8}{260.233}{347.77}{vectors}%
\TextAt{13}{299.972}{347.77}{v}%
\TextAt{17}{305.631}{345.977}{1}%
\TextAt{13}{310.364}{347.77}{\Character{59}}%
\TextAt{13}{315.588}{347.77}{\Character{58}}%
\TextAt{13}{320.825}{347.77}{\Character{58}}%
\TextAt{13}{326.049}{347.77}{\Character{58}}%
\TextAt{13}{331.286}{347.77}{\Character{59}}%
\TextAt{13}{336.522}{347.77}{v}%
\TextAt{21}{342.244}{345.977}{n}%
\TextAt{8}{347.88}{347.77}{,}%
\TextAt{9}{190.943}{320.063}{\Character{8}}%
\TextAt{9}{202.707}{320.063}{\Character{61}}%
\TextAt{16}{215.144}{320.063}{\Character{102}}%
\TextAt{13}{221.122}{320.063}{x}%
\TextAt{17}{227.774}{318.27}{1}%
\TextAt{13}{232.506}{320.063}{v}%
\TextAt{17}{238.165}{318.27}{1}%
\TextAt{9}{245.554}{320.063}{\Character{43}}%
\TextAt{16}{257.316}{320.063}{\Character{1}}%
\TextAt{16}{262.612}{320.063}{\Character{1}}%
\TextAt{16}{267.92}{320.063}{\Character{1}}%
\TextAt{9}{273.92}{320.063}{\Character{43}}%
\TextAt{13}{285.681}{320.063}{x}%
\TextAt{21}{292.333}{318.27}{n}%
\TextAt{13}{297.97}{320.063}{v}%
\TextAt{21}{303.629}{318.27}{n}%
\TextAt{16}{312.586}{320.063}{\Character{106}}%
\TextAt{13}{319.228}{320.063}{x}%
\TextAt{21}{325.88}{318.27}{i}%
\NativeMathAt{332.582}{320.063}{11.955}{\in}%
\TextAt{18}{343.873}{320.063}{R}%
\TextAt{13}{352.507}{320.063}{\Character{59}}%
\TextAt{9}{361.649}{320.063}{0}%
\NativeMathAt{370.823}{320.063}{11.955}{\leq}%
\TextAt{13}{383.442}{320.063}{x}%
\TextAt{21}{390.094}{318.27}{i}%
\TextAt{13}{396.796}{320.063}{\Character{60}}%
\TextAt{9}{409.222}{320.063}{1}%
\TextAt{16}{415.075}{320.063}{\Character{103}}%
\TextAt{8}{79.226}{292.356}{has}%
\TextAt{8}{100.1}{292.356}{volume}%
\TextAt{8}{262.243}{275.31}{vol}%
\TextAt{9}{277.534}{275.31}{\Character{40}\Character{8}\Character{41}}%
\TextAt{9}{298.384}{275.31}{\Character{61}}%
\TextAt{16}{310.839}{275.31}{\Character{106}}%
\TextAt{8}{316.153}{275.31}{det}%
\TextAt{13}{334.404}{275.31}{A}%
\TextAt{16}{343.18}{275.31}{\Character{106}}%
\TextAt{13}{346.501}{275.31}{\Character{59}}%
\TextAt{8}{79.226}{252.343}{where}%
\TextAt{13}{113.921}{252.343}{A}%
\TextAt{9}{127.527}{252.343}{\Character{61}}%
\TextAt{9}{141.455}{252.343}{\Character{40}}%
\TextAt{13}{146.014}{252.343}{a}%
\TextAt{21}{152.159}{250.549}{ik}%
\TextAt{9}{160.162}{252.343}{\Character{41}}%
\TextAt{8}{169.498}{252.343}{is}%
\TextAt{8}{182.147}{252.343}{the}%
\TextAt{8}{203.19}{252.343}{matrix}%
\TextAt{8}{242.129}{252.343}{of}%
\TextAt{8}{256.345}{252.343}{the}%
\TextAt{8}{277.387}{252.343}{base}%
\TextAt{8}{304.347}{252.343}{change}%
\TextAt{8}{343.932}{252.343}{from}%
\TextAt{13}{372.458}{252.343}{e}%
\TextAt{17}{377.883}{250.549}{1}%
\TextAt{13}{382.616}{252.343}{\Character{59}}%
\TextAt{13}{387.84}{252.343}{\Character{58}}%
\TextAt{13}{393.077}{252.343}{\Character{58}}%
\TextAt{13}{398.301}{252.343}{\Character{58}}%
\TextAt{13}{403.538}{252.343}{\Character{59}}%
\TextAt{13}{408.774}{252.343}{e}%
\TextAt{21}{414.262}{250.549}{n}%
\TextAt{8}{424.681}{252.343}{to}%
\TextAt{13}{439.878}{252.343}{v}%
\TextAt{17}{445.537}{250.549}{1}%
\TextAt{13}{450.269}{252.343}{\Character{59}}%
\TextAt{13}{455.493}{252.343}{\Character{58}}%
\TextAt{13}{460.73}{252.343}{\Character{58}}%
\TextAt{13}{465.966}{252.343}{\Character{58}}%
\TextAt{13}{471.191}{252.343}{\Character{59}}%
\TextAt{13}{476.427}{252.343}{v}%
\TextAt{21}{482.149}{250.549}{n}%
\TextAt{8}{487.786}{252.343}{,}%
\TextAt{8}{496.042}{252.343}{so}%
\TextAt{8}{511.298}{252.343}{that}%
\TextAt{13}{79.226}{235.297}{v}%
\TextAt{21}{84.885}{233.503}{i}%
\TextAt{9}{91.588}{235.297}{\Character{61}}%
\TextAt{20}{104.013}{244.263}{\Character{80}}%
\TextAt{21}{116.632}{231.81}{k}%
\TextAt{13}{123.745}{235.297}{a}%
\TextAt{21}{129.89}{233.503}{ik}%
\TextAt{13}{137.892}{235.297}{e}%
\TextAt{21}{143.318}{233.503}{k}%
\TextAt{8}{148.438}{235.297}{.}%
\TextAt{8}{156.886}{235.297}{Since}%
\TextAt{9}{164.887}{199.23}{\Character{40}\Character{40}}%
\TextAt{13}{173.992}{199.23}{v}%
\TextAt{21}{179.651}{197.437}{i}%
\TextAt{13}{183.033}{199.23}{\Character{59}}%
\TextAt{13}{188.257}{199.23}{v}%
\TextAt{21}{193.936}{197.437}{j}%
\TextAt{9}{198.318}{199.23}{\Character{41}\Character{41}}%
\TextAt{9}{210.716}{199.23}{\Character{61}}%
\TextAt{20}{223.169}{219.674}{\Character{32}}%
\TextAt{20}{232.634}{210.588}{\Character{88}}%
\TextAt{21}{236.051}{185.128}{k}%
\TextAt{21}{240.666}{185.128}{\Character{59}\Character{96}}%
\TextAt{13}{251.895}{199.23}{a}%
\TextAt{21}{258.04}{197.437}{ik}%
\TextAt{13}{266.043}{199.23}{a}%
\TextAt{21}{272.188}{197.437}{j}%
\TextAt{21}{276.062}{197.437}{\Character{96}}%
\TextAt{9}{280.029}{199.23}{\Character{40}}%
\TextAt{13}{284.581}{199.23}{e}%
\TextAt{21}{290.007}{197.437}{k}%
\TextAt{13}{295.127}{199.23}{\Character{59}}%
\TextAt{13}{300.351}{199.23}{e}%
\TextAt{21}{305.796}{197.437}{\Character{96}}%
\TextAt{9}{309.754}{199.23}{\Character{41}}%
\TextAt{20}{314.306}{219.674}{\Character{33}}%
\TextAt{9}{327.091}{199.23}{\Character{61}}%
\TextAt{20}{339.517}{219.674}{\Character{32}}%
\TextAt{20}{348.981}{210.588}{\Character{88}}%
\TextAt{21}{355.305}{185.128}{k}%
\TextAt{13}{368.242}{199.23}{a}%
\TextAt{21}{374.387}{197.437}{ik}%
\TextAt{13}{382.39}{199.23}{a}%
\TextAt{21}{388.535}{197.437}{j}%
\TextAt{21}{392.409}{197.437}{k}%
\TextAt{20}{397.539}{219.674}{\Character{33}}%
\TextAt{9}{410.324}{199.23}{\Character{61}}%
\TextAt{13}{422.75}{199.23}{AA}%
\TextAt{21}{440.3}{204.167}{t}%
\TextAt{13}{443.857}{199.23}{\Character{59}}%
\TextAt{8}{79.226}{162.428}{we}%
\TextAt{8}{96.454}{162.428}{also}%
\TextAt{8}{119.934}{162.428}{have}%
\TextAt{8}{146.928}{162.428}{the}%
\TextAt{8}{167.085}{162.428}{invariant}%
\TextAt{8}{216.208}{162.428}{notation}%
\TextAt{8}{239.759}{134.721}{vol}%
\TextAt{9}{255.05}{134.721}{\Character{40}\Character{8}\Character{41}}%
\TextAt{9}{275.9}{134.721}{\Character{61}}%
\TextAt{16}{288.355}{134.721}{\Character{106}}%
\TextAt{8}{293.669}{134.721}{det}%
\TextAt{9}{309.928}{134.721}{\Character{40}\Character{40}}%
\TextAt{13}{319.032}{134.721}{v}%
\TextAt{21}{324.692}{132.927}{i}%
\TextAt{13}{328.073}{134.721}{\Character{59}}%
\TextAt{13}{333.297}{134.721}{v}%
\TextAt{21}{338.976}{132.927}{j}%
\TextAt{9}{343.358}{134.721}{\Character{41}\Character{41}}%
\TextAt{16}{352.463}{134.721}{\Character{106}}%
\TextAt{17}{355.784}{139.657}{1\Character{47}2}%
\TextAt{13}{368.985}{134.721}{\Character{58}}%
\TextAt{8}{79.226}{107.013}{Let}%
\TextAt{9}{100.475}{107.013}{\Character{0}}%
\TextAt{8}{111.966}{107.013}{be}%
\TextAt{8}{128.17}{107.013}{the}%
\TextAt{8}{148.606}{107.013}{lattice}%
\TextAt{8}{184.656}{107.013}{spanned}%
\TextAt{8}{230.521}{107.013}{by}%
\TextAt{13}{247.059}{107.013}{v}%
\TextAt{17}{252.718}{105.22}{1}%
\TextAt{13}{257.451}{107.013}{\Character{59}}%
\TextAt{13}{262.675}{107.013}{\Character{58}}%
\TextAt{13}{267.912}{107.013}{\Character{58}}%
\TextAt{13}{273.136}{107.013}{\Character{58}}%
\TextAt{13}{278.373}{107.013}{\Character{59}}%
\TextAt{13}{283.609}{107.013}{v}%
\TextAt{21}{289.331}{105.22}{n}%
\TextAt{8}{294.967}{107.013}{.}%
\TextAt{8}{304.254}{107.013}{Then}%
\TextAt{9}{335.091}{107.013}{\Character{8}}%
\TextAt{8}{347.723}{107.013}{is}%
\TextAt{8}{359.766}{107.013}{a}%
\TextAt{8}{369.801}{107.013}{fundamental}%
\TextAt{8}{437.723}{107.013}{mesh}%
\TextAt{8}{467.975}{107.013}{of}%
\TextAt{9}{481.584}{107.013}{\Character{0}}%
\TextAt{8}{488.898}{107.013}{,}%
\TextAt{8}{496.397}{107.013}{and}%
\TextAt{8}{519.439}{107.013}{we}%
\TextAt{8}{79.226}{89.967}{write}%
\TextAt{8}{109.138}{89.967}{for}%
\TextAt{8}{127.023}{89.967}{short}%
\TextAt{8}{264.219}{72.921}{vol}%
\TextAt{9}{279.51}{72.921}{\Character{40}\Character{0}\Character{41}}%
\TextAt{9}{299.212}{72.921}{\Character{61}}%
\TextAt{8}{311.675}{72.921}{vol}%
\TextAt{9}{326.966}{72.921}{\Character{40}\Character{8}\Character{41}}%
\TextAt{13}{344.525}{72.921}{\Character{58}}%
\TextAt{8}{300.14}{43.033}{24}%
\Destination{78.23}{809.815}{nav.page.24}%
\Destination{79.36}{682.155}{nav.print.26}%
\Destination{79.36}{732.455}{nav.number.1.4.3}%
}

% Page 37 (body page 25)
\Page{%
\TextAt{2}{79.226}{796.943}{I.}%
\TextAt{2}{89.83}{796.943}{Algebraic}%
\TextAt{2}{138.648}{796.943}{Integers}%
\Vector{q 0 G 0 g 0.299 w 7.226 18.729 m 460.769 18.729 l S Q}%
\TextAt{8}{79.226}{761.124}{This}%
\TextAt{8}{104.816}{761.124}{does}%
\TextAt{8}{130.083}{761.124}{not}%
\TextAt{8}{149.767}{761.124}{depend}%
\TextAt{8}{189.273}{761.124}{on}%
\TextAt{8}{204.402}{761.124}{the}%
\TextAt{8}{223.428}{761.124}{choice}%
\TextAt{8}{257.088}{761.124}{of}%
\TextAt{8}{269.288}{761.124}{a}%
\TextAt{8}{277.913}{761.124}{basis}%
\TextAt{13}{305.523}{761.124}{v}%
\TextAt{17}{311.183}{759.33}{1}%
\TextAt{13}{315.915}{761.124}{\Character{59}}%
\TextAt{13}{321.139}{761.124}{\Character{58}}%
\TextAt{13}{326.376}{761.124}{\Character{58}}%
\TextAt{13}{331.6}{761.124}{\Character{58}}%
\TextAt{13}{336.837}{761.124}{\Character{59}}%
\TextAt{13}{342.073}{761.124}{v}%
\TextAt{21}{347.795}{759.33}{n}%
\TextAt{8}{356.199}{761.124}{of}%
\TextAt{8}{368.399}{761.124}{the}%
\TextAt{8}{387.425}{761.124}{lattice}%
\TextAt{8}{422.065}{761.124}{because}%
\TextAt{8}{464.237}{761.124}{the}%
\TextAt{8}{483.263}{761.124}{transition}%
\TextAt{8}{79.226}{744.078}{matrix}%
\TextAt{8}{118.665}{744.078}{passing}%
\TextAt{8}{161.152}{744.078}{to}%
\TextAt{8}{176.847}{744.078}{a}%
\TextAt{8}{187.988}{744.078}{different}%
\TextAt{8}{235.544}{744.078}{basis,}%
\TextAt{8}{269.267}{744.078}{as}%
\TextAt{8}{285.023}{744.078}{well}%
\TextAt{8}{310.139}{744.078}{as}%
\TextAt{8}{325.894}{744.078}{its}%
\TextAt{8}{343.598}{744.078}{inverse,}%
\TextAt{8}{387.34}{744.078}{has}%
\TextAt{8}{409.599}{744.078}{integer}%
\TextAt{8}{449.683}{744.078}{coefficients,}%
\TextAt{8}{513.904}{744.078}{and}%
\TextAt{8}{79.226}{727.032}{therefore}%
\TextAt{8}{129.583}{727.032}{has}%
\TextAt{8}{151.714}{727.032}{determinant}%
\NativeMathAt{218.988}{727.032}{11.955}{\pm}%
\TextAt{9}{228.286}{727.032}{1}%
\TextAt{8}{239.294}{727.032}{so}%
\TextAt{8}{254.921}{727.032}{that}%
\TextAt{8}{281.547}{727.032}{the}%
\TextAt{8}{302.96}{727.032}{set}%
\TextAt{9}{322.485}{727.032}{\Character{8}}%
\TextAt{8}{336.094}{727.032}{is}%
\TextAt{8}{349.115}{727.032}{transformed}%
\TextAt{8}{415.803}{727.032}{into}%
\TextAt{8}{440.803}{727.032}{a}%
\TextAt{8}{451.815}{727.032}{set}%
\TextAt{8}{471.34}{727.032}{of}%
\TextAt{8}{485.927}{727.032}{the}%
\TextAt{8}{507.341}{727.032}{same}%
\TextAt{8}{79.226}{709.986}{volume.}%
\TextAt{8}{97.159}{689.951}{We}%
\TextAt{8}{118.547}{689.951}{now}%
\TextAt{8}{144.178}{689.951}{come}%
\TextAt{8}{175.333}{689.951}{to}%
\TextAt{8}{190.887}{689.951}{the}%
\TextAt{8}{212.286}{689.951}{most}%
\TextAt{8}{242.21}{689.951}{important}%
\TextAt{8}{298.746}{689.951}{theorem}%
\TextAt{8}{345.514}{689.951}{about}%
\TextAt{8}{380.256}{689.951}{lattices.}%
\TextAt{8}{428.921}{689.951}{A}%
\TextAt{8}{442.835}{689.951}{subset}%
\TextAt{13}{479.968}{689.951}{X}%
\TextAt{8}{495.764}{689.951}{of}%
\TextAt{13}{510.337}{689.951}{V}%
\TextAt{8}{524.903}{689.951}{is}%
\TextAt{8}{79.226}{672.905}{called}%
\TextAt{12}{114.093}{672.905}{centrally}%
\TextAt{12}{171.609}{672.905}{symmetric}%
\TextAt{8}{233.034}{672.905}{,}%
\TextAt{8}{242.313}{672.905}{if,}%
\TextAt{8}{258.418}{672.905}{given}%
\TextAt{8}{290.691}{672.905}{any}%
\TextAt{8}{314.512}{672.905}{point}%
\TextAt{13}{346.785}{672.905}{x}%
\NativeMathAt{359.661}{672.905}{11.955}{\in}%
\TextAt{13}{373.855}{672.905}{X}%
\TextAt{8}{384.51}{672.905}{,}%
\TextAt{8}{393.788}{672.905}{the}%
\TextAt{8}{415.648}{672.905}{point}%
\TextAt{16}{447.921}{672.905}{\Character{0}}%
\TextAt{13}{457.22}{672.905}{x}%
\TextAt{8}{469.473}{672.905}{also}%
\TextAt{8}{494.656}{672.905}{belongs}%
\TextAt{8}{79.226}{655.859}{to}%
\TextAt{13}{94.859}{655.859}{X}%
\TextAt{8}{105.514}{655.859}{.}%
\TextAt{8}{117.931}{655.859}{It}%
\TextAt{8}{131.926}{655.859}{is}%
\TextAt{8}{145.012}{655.859}{called}%
\TextAt{12}{179.499}{655.859}{convex}%
\TextAt{8}{224.494}{655.859}{if,}%
\TextAt{8}{240.123}{655.859}{given}%
\TextAt{8}{272.015}{655.859}{any}%
\TextAt{8}{295.454}{655.859}{two}%
\TextAt{8}{318.894}{655.859}{points}%
\TextAt{13}{355.401}{655.859}{x\Character{59}}%
\TextAt{13}{367.273}{655.859}{y}%
\NativeMathAt{379.009}{655.859}{11.955}{\in}%
\TextAt{13}{392.554}{655.859}{X}%
\TextAt{8}{403.209}{655.859}{,}%
\TextAt{8}{412.011}{655.859}{the}%
\TextAt{8}{433.491}{655.859}{whole}%
\TextAt{8}{467.977}{655.859}{line}%
\TextAt{8}{491.405}{655.859}{segment}%
\TextAt{16}{79.226}{638.813}{\Character{102}}%
\TextAt{13}{85.204}{638.813}{ty}%
\TextAt{9}{98.25}{638.813}{\Character{43}}%
\TextAt{9}{110.026}{638.813}{\Character{40}1}%
\TextAt{16}{123.126}{638.813}{\Character{0}}%
\TextAt{13}{135.107}{638.813}{t}%
\TextAt{9}{139.334}{638.813}{\Character{41}}%
\TextAt{13}{143.887}{638.813}{x}%
\TextAt{16}{153.925}{638.813}{\Character{106}}%
\TextAt{9}{160.632}{638.813}{0}%
\NativeMathAt{169.871}{638.813}{11.955}{\leq}%
\TextAt{13}{182.556}{638.813}{t}%
\NativeMathAt{190.169}{638.813}{11.955}{\leq}%
\TextAt{9}{202.854}{638.813}{1}%
\TextAt{16}{208.707}{638.813}{\Character{103}}%
\TextAt{8}{218.62}{638.813}{joining}%
\TextAt{13}{257.357}{638.813}{x}%
\TextAt{8}{267.945}{638.813}{with}%
\TextAt{13}{294.643}{638.813}{y}%
\TextAt{8}{304.716}{638.813}{is}%
\TextAt{8}{316.518}{638.813}{contained}%
\TextAt{8}{369.565}{638.813}{in}%
\TextAt{13}{383.256}{638.813}{X}%
\TextAt{8}{393.912}{638.813}{.}%
\TextAt{8}{402.475}{638.813}{With}%
\TextAt{8}{432.748}{638.813}{these}%
\TextAt{8}{462.758}{638.813}{definitions}%
\TextAt{8}{519.44}{638.813}{we}%
\TextAt{8}{79.226}{621.767}{have}%
\TextAt{12}{79.226}{595.814}{Minkowski’s}%
\TextAt{12}{156.6}{595.814}{Lattice}%
\TextAt{12}{202.937}{595.814}{Point}%
\TextAt{12}{239.962}{595.814}{Theorem}%
\TextAt{12}{297.884}{595.814}{(4.4).}%
\TextAt{8}{333.99}{595.814}{Let}%
\TextAt{9}{355.748}{595.814}{\Character{0}}%
\TextAt{8}{367.749}{595.814}{be}%
\TextAt{8}{384.463}{595.814}{a}%
\TextAt{8}{395.008}{595.814}{complete}%
\TextAt{8}{445.22}{595.814}{lattice}%
\TextAt{8}{481.78}{595.814}{in}%
\TextAt{8}{496.222}{595.814}{the}%
\TextAt{8}{517.168}{595.814}{eu-}%
\TextAt{8}{79.226}{578.768}{clidean}%
\TextAt{8}{118.893}{578.768}{vector}%
\TextAt{8}{154.018}{578.768}{space}%
\TextAt{13}{185.292}{578.768}{V}%
\TextAt{8}{198.615}{578.768}{and}%
\TextAt{13}{221.378}{578.768}{X}%
\TextAt{8}{235.931}{578.768}{a}%
\TextAt{8}{245.686}{578.768}{centrally}%
\TextAt{8}{293.818}{578.768}{symmetric,}%
\TextAt{8}{354.036}{578.768}{convex}%
\TextAt{8}{392.412}{578.768}{subset}%
\TextAt{8}{428.301}{578.768}{of}%
\TextAt{13}{441.631}{578.768}{V}%
\TextAt{8}{451.057}{578.768}{.}%
\TextAt{8}{459.505}{578.768}{Suppose}%
\TextAt{8}{505.413}{578.768}{that}%
\TextAt{8}{256.378}{549.975}{vol}%
\TextAt{9}{271.669}{549.975}{\Character{40}}%
\TextAt{13}{276.221}{549.975}{X}%
\TextAt{9}{286.876}{549.975}{\Character{41}}%
\TextAt{13}{294.749}{549.975}{\Character{62}}%
\TextAt{9}{307.175}{549.975}{2}%
\TextAt{21}{313.028}{554.911}{n}%
\TextAt{8}{320.656}{549.975}{vol}%
\TextAt{9}{335.947}{549.975}{\Character{40}\Character{0}\Character{41}}%
\TextAt{13}{352.366}{549.975}{\Character{58}}%
\TextAt{8}{79.226}{521.182}{Then}%
\TextAt{13}{109.784}{521.182}{X}%
\TextAt{8}{124.336}{521.182}{contains}%
\TextAt{8}{170.256}{521.182}{at}%
\TextAt{8}{184.566}{521.182}{least}%
\TextAt{8}{211.944}{521.182}{one}%
\TextAt{8}{233.403}{521.182}{nonzero}%
\TextAt{8}{276.98}{521.182}{lattice}%
\TextAt{8}{312.75}{521.182}{point}%
\TextAt{13}{343.319}{521.182}{\Character{13}}%
\NativeMathAt{353.362}{521.182}{11.955}{\in}%
\TextAt{9}{364.653}{521.182}{\Character{0}}%
\TextAt{8}{371.967}{521.182}{.}%
\TextAt{19}{79.226}{486.322}{Proof.}%
\TextAt{8}{116}{486.322}{It}%
\TextAt{8}{129.595}{486.322}{is}%
\TextAt{8}{142.281}{486.322}{enough}%
\TextAt{8}{183.528}{486.322}{to}%
\TextAt{8}{198.761}{486.322}{show}%
\TextAt{8}{228.687}{486.322}{that}%
\TextAt{8}{254.978}{486.322}{there}%
\TextAt{8}{285.812}{486.322}{exist}%
\TextAt{8}{314.434}{486.322}{two}%
\TextAt{8}{337.474}{486.322}{distinct}%
\TextAt{8}{380.729}{486.322}{lattice}%
\TextAt{8}{417.422}{486.322}{points}%
\TextAt{13}{453.528}{486.322}{\Character{13}}%
\TextAt{17}{459.6}{484.529}{1}%
\TextAt{13}{464.332}{486.322}{\Character{59}}%
\TextAt{13}{469.556}{486.322}{\Character{13}}%
\TextAt{17}{475.648}{484.529}{2}%
\NativeMathAt{485.273}{486.322}{11.955}{\in}%
\TextAt{9}{498.136}{486.322}{\Character{0}}%
\TextAt{8}{510.27}{486.322}{such}%
\TextAt{8}{79.226}{469.276}{that}%
\TextAt{20}{231.746}{461.914}{\Character{0}}%
\TextAt{17}{238.421}{456.937}{1}%
\Vector{q 0 G 0 g 0.478 w 166.421 -318.781 m 170.656 -318.781 l S Q}%
\TextAt{17}{238.421}{448.107}{2}%
\TextAt{13}{243.851}{452.23}{X}%
\TextAt{9}{257.163}{452.23}{\Character{43}}%
\TextAt{13}{268.924}{452.23}{\Character{13}}%
\TextAt{17}{274.996}{450.437}{1}%
\TextAt{20}{279.728}{461.914}{\Character{1}}%
\TextAt{16}{287.865}{452.23}{\Character{92}}%
\TextAt{20}{298.491}{461.914}{\Character{0}}%
\TextAt{17}{305.166}{456.937}{1}%
\Vector{q 0 G 0 g 0.478 w 233.166 -318.781 m 237.4 -318.781 l S Q}%
\TextAt{17}{305.166}{448.107}{2}%
\TextAt{13}{310.596}{452.23}{X}%
\TextAt{9}{323.908}{452.23}{\Character{43}}%
\TextAt{13}{335.669}{452.23}{\Character{13}}%
\TextAt{17}{341.741}{450.437}{2}%
\TextAt{20}{346.473}{461.914}{\Character{1}}%
\NativeMathAt{355.274}{452.23}{11.955}{\neq}%
\TextAt{18}{367.699}{452.23}{\Character{63}}%
\TextAt{13}{376.998}{452.23}{\Character{58}}%
\TextAt{8}{79.226}{428.797}{In}%
\TextAt{8}{93.847}{428.797}{fact,}%
\TextAt{8}{120.185}{428.797}{choosing}%
\TextAt{8}{167.73}{428.797}{a}%
\TextAt{8}{177.486}{428.797}{point}%
\TextAt{8}{208.055}{428.797}{in}%
\TextAt{8}{221.708}{428.797}{this}%
\TextAt{8}{244.53}{428.797}{intersection,}%
\TextAt{17}{213.571}{404.711}{1}%
\Vector{q 0 G 0 g 0.478 w 141.571 -371.007 m 145.805 -371.007 l S Q}%
\TextAt{17}{213.571}{395.882}{2}%
\TextAt{13}{219}{400.004}{x}%
\TextAt{17}{225.652}{398.211}{1}%
\TextAt{9}{233.041}{400.004}{\Character{43}}%
\TextAt{13}{244.803}{400.004}{\Character{13}}%
\TextAt{17}{250.875}{398.211}{1}%
\TextAt{9}{258.928}{400.004}{\Character{61}}%
\TextAt{17}{272.549}{404.711}{1}%
\Vector{q 0 G 0 g 0.478 w 200.549 -371.007 m 204.783 -371.007 l S Q}%
\TextAt{17}{272.549}{395.882}{2}%
\TextAt{13}{277.978}{400.004}{x}%
\TextAt{17}{284.63}{398.211}{2}%
\TextAt{9}{292.019}{400.004}{\Character{43}}%
\TextAt{13}{303.781}{400.004}{\Character{13}}%
\TextAt{17}{309.853}{398.211}{2}%
\TextAt{13}{314.585}{400.004}{\Character{59}}%
\TextAt{13}{343.72}{400.004}{x}%
\TextAt{17}{350.392}{398.211}{1}%
\TextAt{13}{355.124}{400.004}{\Character{59}}%
\TextAt{13}{360.348}{400.004}{x}%
\TextAt{17}{367.02}{398.211}{2}%
\NativeMathAt{375.073}{400.004}{11.955}{\in}%
\TextAt{13}{386.364}{400.004}{X}%
\TextAt{13}{396.358}{400.004}{\Character{59}}%
\TextAt{8}{79.226}{371.211}{we}%
\TextAt{8}{96.454}{371.211}{obtain}%
\TextAt{8}{132.881}{371.211}{an}%
\TextAt{8}{149.14}{371.211}{element}%
\TextAt{13}{241.839}{354.165}{\Character{13}}%
\TextAt{9}{251.882}{354.165}{\Character{61}}%
\TextAt{13}{264.307}{354.165}{\Character{13}}%
\TextAt{17}{270.379}{352.372}{1}%
\TextAt{16}{277.768}{354.165}{\Character{0}}%
\TextAt{13}{289.724}{354.165}{\Character{13}}%
\TextAt{17}{295.795}{352.372}{2}%
\TextAt{9}{303.849}{354.165}{\Character{61}}%
\TextAt{17}{317.47}{358.872}{1}%
\Vector{q 0 G 0 g 0.478 w 245.47 -416.846 m 249.704 -416.846 l S Q}%
\TextAt{17}{317.47}{350.042}{2}%
\TextAt{13}{322.899}{354.165}{x}%
\TextAt{17}{329.551}{352.372}{2}%
\TextAt{16}{336.94}{354.165}{\Character{0}}%
\TextAt{17}{350.091}{358.872}{1}%
\Vector{q 0 G 0 g 0.478 w 278.091 -416.846 m 282.325 -416.846 l S Q}%
\TextAt{17}{350.091}{350.042}{2}%
\TextAt{13}{355.521}{354.165}{x}%
\TextAt{17}{362.173}{352.372}{1}%
\TextAt{13}{366.905}{354.165}{\Character{59}}%
\TextAt{8}{79.226}{330.732}{which}%
\TextAt{8}{111.97}{330.732}{is}%
\TextAt{8}{122.991}{330.732}{the}%
\TextAt{8}{142.405}{330.732}{center}%
\TextAt{8}{176.452}{330.732}{of}%
\TextAt{8}{189.039}{330.732}{the}%
\TextAt{8}{208.453}{330.732}{line}%
\TextAt{8}{229.816}{330.732}{segment}%
\TextAt{8}{274.336}{330.732}{joining}%
\TextAt{13}{312.292}{330.732}{x}%
\TextAt{17}{318.944}{328.939}{2}%
\TextAt{8}{326.831}{330.732}{and}%
\TextAt{16}{348.851}{330.732}{\Character{0}}%
\TextAt{13}{358.15}{330.733}{x}%
\TextAt{17}{364.802}{328.939}{1}%
\TextAt{8}{369.534}{330.732}{,}%
\TextAt{8}{376.089}{330.732}{and}%
\TextAt{8}{398.109}{330.732}{therefore}%
\TextAt{8}{446.466}{330.732}{belongs}%
\TextAt{8}{487.734}{330.732}{to}%
\TextAt{13}{501.302}{330.732}{X}%
\TextAt{16}{513.095}{330.733}{\Character{92}}%
\TextAt{9}{522.204}{330.732}{\Character{0}}%
\TextAt{8}{529.518}{330.732}{.}%
\TextAt{8}{97.159}{310.698}{Now,}%
\TextAt{8}{127.452}{310.698}{if}%
\TextAt{8}{138.482}{310.698}{the}%
\TextAt{8}{158.943}{310.698}{sets}%
\TextAt{17}{183.326}{315.405}{1}%
\Vector{q 0 G 0 g 0.478 w 111.326 -460.313 m 115.561 -460.313 l S Q}%
\TextAt{17}{183.326}{306.575}{2}%
\TextAt{13}{188.756}{310.698}{X}%
\TextAt{9}{202.276}{310.698}{\Character{43}}%
\TextAt{13}{214.245}{310.698}{\Character{13}}%
\TextAt{8}{220.968}{310.698}{,}%
\TextAt{13}{228.499}{310.698}{\Character{13}}%
\NativeMathAt{239.062}{310.698}{11.955}{\in}%
\TextAt{9}{250.873}{310.698}{\Character{0}}%
\TextAt{8}{258.187}{310.698}{,}%
\TextAt{8}{265.718}{310.698}{were}%
\TextAt{8}{293.006}{310.698}{pairwise}%
\TextAt{8}{338.897}{310.698}{disjoint,}%
\TextAt{8}{384.218}{310.698}{then}%
\TextAt{8}{411.183}{310.698}{the}%
\TextAt{8}{431.645}{310.698}{same}%
\TextAt{8}{461.277}{310.698}{would}%
\TextAt{8}{495.726}{310.698}{be}%
\TextAt{8}{511.956}{310.698}{true}%
\TextAt{8}{79.226}{293.652}{of}%
\TextAt{8}{92.556}{293.652}{their}%
\TextAt{8}{120.519}{293.652}{intersections}%
\TextAt{9}{187.958}{293.652}{\Character{8}}%
\TextAt{16}{199.069}{293.652}{\Character{92}}%
\TextAt{9}{209.696}{293.652}{\Character{40}}%
\TextAt{17}{215.444}{298.359}{1}%
\Vector{q 0 G 0 g 0.478 w 143.444 -477.359 m 147.678 -477.359 l S Q}%
\TextAt{17}{215.444}{289.529}{2}%
\TextAt{13}{220.874}{293.652}{X}%
\TextAt{9}{234.186}{293.652}{\Character{43}}%
\TextAt{13}{245.947}{293.652}{\Character{13}}%
\TextAt{9}{252.669}{293.652}{\Character{41}}%
\TextAt{8}{261.119}{293.652}{with}%
\TextAt{8}{287.779}{293.652}{a}%
\TextAt{8}{297.534}{293.652}{fundamental}%
\TextAt{8}{365.177}{293.652}{mesh}%
\TextAt{9}{395.148}{293.652}{\Character{8}}%
\TextAt{8}{407.5}{293.652}{of}%
\TextAt{9}{420.83}{293.652}{\Character{0}}%
\TextAt{8}{428.144}{293.652}{,}%
\TextAt{8}{435.293}{293.652}{i.e.,}%
\TextAt{8}{457.398}{293.652}{we}%
\TextAt{8}{474.626}{293.652}{would}%
\TextAt{8}{508.77}{293.652}{have}%
\TextAt{8}{220.598}{264.859}{vol}%
\TextAt{9}{235.888}{264.859}{\Character{40}\Character{8}\Character{41}}%
\NativeMathAt{256.768}{264.859}{11.955}{\geq}%
\TextAt{20}{269.388}{276.216}{\Character{88}}%
\TextAt{21}{270.139}{250.844}{\Character{13}}%
\NativeMathAt{274.971}{250.844}{7.97}{\in}%
\TextAt{17}{280.616}{250.844}{\Character{0}}%
\TextAt{8}{288.649}{264.859}{vol}%
\TextAt{20}{305.932}{274.542}{\Character{0}}%
\TextAt{9}{311.411}{264.859}{\Character{8}}%
\TextAt{16}{322.522}{264.859}{\Character{92}}%
\TextAt{9}{333.149}{264.859}{\Character{40}}%
\TextAt{17}{338.897}{269.566}{1}%
\Vector{q 0 G 0 g 0.478 w 266.897 -506.153 m 271.131 -506.153 l S Q}%
\TextAt{17}{338.897}{260.736}{2}%
\TextAt{13}{344.327}{264.859}{X}%
\TextAt{9}{357.638}{264.859}{\Character{43}}%
\TextAt{13}{369.4}{264.859}{\Character{13}}%
\TextAt{9}{376.122}{264.859}{\Character{41}}%
\TextAt{20}{380.674}{274.542}{\Character{1}}%
\TextAt{13}{388.146}{264.859}{\Character{58}}%
\TextAt{8}{79.226}{225.512}{But}%
\TextAt{8}{102.153}{225.512}{translation}%
\TextAt{8}{161.101}{225.512}{of}%
\TextAt{9}{174.117}{225.512}{\Character{8}}%
\TextAt{16}{184.587}{225.512}{\Character{92}}%
\TextAt{9}{194.572}{225.512}{\Character{40}}%
\TextAt{17}{200.32}{230.219}{1}%
\Vector{q 0 G 0 g 0.478 w 128.32 -545.499 m 132.554 -545.499 l S Q}%
\TextAt{17}{200.32}{221.39}{2}%
\TextAt{13}{205.749}{225.512}{X}%
\TextAt{9}{218.419}{225.512}{\Character{43}}%
\TextAt{13}{229.539}{225.512}{\Character{13}}%
\TextAt{9}{236.261}{225.512}{\Character{41}}%
\TextAt{8}{244.397}{225.512}{by}%
\TextAt{16}{260.342}{225.512}{\Character{0}}%
\TextAt{13}{269.641}{225.512}{\Character{13}}%
\TextAt{8}{279.947}{225.512}{creates}%
\TextAt{8}{318.714}{225.512}{the}%
\TextAt{8}{338.557}{225.512}{set}%
\TextAt{9}{356.511}{225.512}{\Character{40}\Character{8}}%
\TextAt{16}{371.532}{225.512}{\Character{0}}%
\TextAt{13}{382.846}{225.512}{\Character{13}}%
\TextAt{9}{389.568}{225.512}{\Character{41}}%
\TextAt{16}{396.135}{225.512}{\Character{92}}%
\TextAt{17}{407.316}{230.219}{1}%
\Vector{q 0 G 0 g 0.478 w 335.316 -545.499 m 339.55 -545.499 l S Q}%
\TextAt{17}{407.316}{221.39}{2}%
\TextAt{13}{412.745}{225.512}{X}%
\TextAt{8}{426.984}{225.512}{of}%
\TextAt{8}{440}{225.512}{equal}%
\TextAt{8}{470.579}{225.512}{volume,}%
\TextAt{8}{513.904}{225.512}{and}%
\TextAt{8}{79.226}{208.466}{the}%
\TextAt{9}{99.704}{208.466}{\Character{8}}%
\TextAt{16}{111.034}{208.466}{\Character{0}}%
\TextAt{13}{123.208}{208.466}{\Character{13}}%
\TextAt{8}{129.93}{208.466}{,}%
\TextAt{13}{137.481}{208.466}{\Character{13}}%
\NativeMathAt{148.072}{208.466}{11.955}{\in}%
\TextAt{9}{159.911}{208.466}{\Character{0}}%
\TextAt{8}{167.225}{208.466}{,}%
\TextAt{8}{174.776}{208.466}{cover}%
\TextAt{8}{205.344}{208.466}{the}%
\TextAt{8}{225.821}{208.466}{entire}%
\TextAt{8}{258.984}{208.466}{space}%
\TextAt{13}{290.58}{208.466}{V}%
\TextAt{8}{300.005}{208.466}{,}%
\TextAt{8}{307.556}{208.466}{therefore}%
\TextAt{8}{356.978}{208.466}{also}%
\TextAt{8}{380.779}{208.466}{the}%
\TextAt{8}{401.257}{208.466}{set}%
\TextAt{17}{421.041}{213.173}{1}%
\Vector{q 0 G 0 g 0.478 w 349.041 -562.545 m 353.275 -562.545 l S Q}%
\TextAt{17}{421.041}{204.344}{2}%
\TextAt{13}{426.471}{208.466}{X}%
\TextAt{8}{437.126}{208.466}{.}%
\TextAt{8}{446.538}{208.466}{Consequently}%
\TextAt{8}{519.44}{208.466}{we}%
\TextAt{8}{79.226}{191.42}{would}%
\TextAt{8}{113.37}{191.42}{obtain}%
\TextAt{8}{159.451}{162.627}{vol}%
\TextAt{9}{174.742}{162.627}{\Character{40}\Character{8}\Character{41}}%
\NativeMathAt{195.622}{162.627}{11.955}{\geq}%
\TextAt{20}{208.241}{173.985}{\Character{88}}%
\TextAt{21}{208.992}{148.613}{\Character{13}}%
\NativeMathAt{213.824}{148.613}{7.97}{\in}%
\TextAt{17}{219.47}{148.613}{\Character{0}}%
\TextAt{8}{227.502}{162.627}{vol}%
\TextAt{20}{244.785}{172.311}{\Character{0}}%
\TextAt{9}{250.265}{162.627}{\Character{40}\Character{8}}%
\TextAt{16}{265.928}{162.627}{\Character{0}}%
\TextAt{13}{277.883}{162.627}{\Character{13}}%
\TextAt{9}{284.606}{162.627}{\Character{41}}%
\TextAt{16}{291.815}{162.627}{\Character{92}}%
\TextAt{17}{303.637}{167.334}{1}%
\Vector{q 0 G 0 g 0.478 w 231.637 -608.384 m 235.871 -608.384 l S Q}%
\TextAt{17}{303.637}{158.504}{2}%
\TextAt{13}{309.067}{162.627}{X}%
\TextAt{20}{319.722}{172.311}{\Character{1}}%
\TextAt{9}{328.522}{162.627}{\Character{61}}%
\TextAt{8}{340.947}{162.627}{vol}%
\TextAt{9}{356.238}{162.627}{\Character{40}}%
\TextAt{17}{361.986}{167.334}{1}%
\Vector{q 0 G 0 g 0.478 w 289.986 -608.384 m 294.22 -608.384 l S Q}%
\TextAt{17}{361.986}{158.504}{2}%
\TextAt{13}{367.416}{162.627}{X}%
\TextAt{9}{378.071}{162.627}{\Character{41}}%
\TextAt{9}{385.926}{162.627}{\Character{61}}%
\TextAt{9}{402.383}{170.715}{1}%
\Vector{q 0 G 0 g 0.478 w 327.565 -608.384 m 339.054 -608.384 l S Q}%
\TextAt{9}{399.565}{154.426}{2}%
\TextAt{21}{405.418}{157.88}{n}%
\TextAt{8}{414.242}{162.627}{vol}%
\TextAt{9}{429.533}{162.627}{\Character{40}}%
\TextAt{13}{434.085}{162.627}{X}%
\TextAt{9}{444.74}{162.627}{\Character{41}}%
\TextAt{13}{449.293}{162.627}{\Character{59}}%
\TextAt{8}{79.226}{124.827}{which}%
\TextAt{8}{112.713}{124.827}{contradicts}%
\TextAt{8}{172.943}{124.827}{the}%
\TextAt{8}{193.099}{124.827}{hypothesis.}%
\Vector{q 0 G 0 g 0.399 w 452.334 -649.173 m 452.334 -641.103 l S Q}%
\Vector{q 0 G 0 g 0.399 w 452.534 -641.302 m 459.707 -641.302 l S Q}%
\Vector{q 0 G 0 g 0.399 w 452.534 -648.974 m 459.707 -648.974 l S Q}%
\Vector{q 0 G 0 g 0.399 w 459.906 -649.173 m 459.906 -641.103 l S Q}%
\TextAt{12}{79.226}{89.967}{Exercise}%
\TextAt{12}{132.23}{89.967}{I.4.1.}%
\TextAt{8}{165.653}{89.967}{Show}%
\TextAt{8}{196.374}{89.967}{that}%
\TextAt{8}{221.571}{89.967}{a}%
\TextAt{8}{231.156}{89.967}{lattice}%
\TextAt{9}{266.754}{89.967}{\Character{0}}%
\TextAt{8}{277.795}{89.967}{in}%
\TextAt{18}{291.276}{89.967}{R}%
\TextAt{21}{299.911}{94.306}{n}%
\TextAt{8}{309.273}{89.967}{is}%
\TextAt{8}{320.866}{89.967}{complete}%
\TextAt{8}{370.117}{89.967}{if}%
\TextAt{8}{380.67}{89.967}{and}%
\TextAt{8}{403.262}{89.967}{only}%
\TextAt{8}{428.782}{89.967}{if}%
\TextAt{8}{439.335}{89.967}{the}%
\TextAt{8}{459.32}{89.967}{quotient}%
\TextAt{18}{505.332}{89.967}{R}%
\TextAt{21}{513.966}{94.306}{n}%
\TextAt{9}{519.602}{89.967}{\Character{47}\Character{0}}%
\TextAt{8}{79.226}{72.921}{is}%
\TextAt{8}{90.99}{72.921}{compact.}%
\TextAt{8}{300.14}{43.033}{25}%
\Destination{78.23}{809.815}{nav.page.25}%
\Destination{79.226}{103.955}{nav.ex.1.4.1}%
\Destination{79.36}{610.855}{nav.print.27}%
\Destination{79.22}{610.2}{nav.number.1.4.4}%
}

% Page 38 (body page 26)
\Page{%
\TextAt{2}{79.226}{796.943}{I.}%
\TextAt{2}{89.83}{796.943}{Algebraic}%
\TextAt{2}{138.648}{796.943}{Integers}%
\Vector{q 0 G 0 g 0.299 w 7.226 18.729 m 460.769 18.729 l S Q}%
\TextAt{12}{79.226}{761.124}{Exercise}%
\TextAt{12}{133.588}{761.124}{I.4.2.}%
\TextAt{8}{168.192}{761.124}{Show}%
\TextAt{8}{200.094}{761.124}{that}%
\TextAt{8}{226.473}{761.124}{Minkowski’s}%
\TextAt{8}{293.619}{761.124}{lattice}%
\TextAt{8}{330.399}{761.124}{point}%
\TextAt{8}{361.978}{761.124}{theorem}%
\TextAt{8}{408.513}{761.124}{cannot}%
\TextAt{8}{447.899}{761.124}{be}%
\TextAt{8}{464.833}{761.124}{improved,}%
\TextAt{8}{520.408}{761.124}{by}%
\TextAt{8}{79.226}{744.078}{giving}%
\TextAt{8}{115.935}{744.078}{an}%
\TextAt{8}{134.102}{744.078}{example}%
\TextAt{8}{181.858}{744.078}{of}%
\TextAt{8}{197.096}{744.078}{a}%
\TextAt{8}{208.759}{744.078}{centrally}%
\TextAt{8}{258.799}{744.078}{symmetric}%
\TextAt{8}{317.673}{744.078}{convex}%
\TextAt{8}{357.957}{744.078}{set}%
\TextAt{13}{378.132}{744.078}{X}%
\NativeMathAt{395.359}{744.078}{11.955}{\subseteq}%
\TextAt{13}{411.23}{744.078}{V}%
\TextAt{8}{426.461}{744.078}{such}%
\TextAt{8}{454.766}{744.078}{that}%
\TextAt{8}{482.042}{744.078}{vol}%
\TextAt{9}{497.333}{744.078}{\Character{40}}%
\TextAt{13}{501.885}{744.078}{X}%
\TextAt{9}{512.54}{744.078}{\Character{41}}%
\TextAt{9}{523.646}{744.078}{\Character{61}}%
\TextAt{9}{79.226}{727.032}{2}%
\TextAt{21}{85.079}{731.37}{n}%
\TextAt{8}{92.708}{727.032}{vol}%
\TextAt{9}{107.999}{727.032}{\Character{40}\Character{0}\Character{41}}%
\TextAt{8}{129.314}{727.032}{which}%
\TextAt{8}{163.799}{727.032}{does}%
\TextAt{8}{191.195}{727.032}{not}%
\TextAt{8}{213.007}{727.032}{contain}%
\TextAt{8}{255.311}{727.032}{any}%
\TextAt{8}{278.427}{727.032}{nonzero}%
\TextAt{8}{323.002}{727.032}{point}%
\TextAt{8}{354.57}{727.032}{of}%
\TextAt{8}{368.899}{727.032}{the}%
\TextAt{8}{390.054}{727.032}{lattice}%
\TextAt{9}{426.823}{727.032}{\Character{0}}%
\TextAt{8}{434.137}{727.032}{.}%
\TextAt{8}{445.582}{727.032}{If}%
\TextAt{13}{458.273}{727.032}{X}%
\TextAt{8}{473.824}{727.032}{is}%
\TextAt{8}{486.587}{727.032}{compact,}%
\TextAt{8}{79.226}{709.986}{however,}%
\TextAt{8}{127.358}{709.986}{then}%
\TextAt{8}{154.018}{709.986}{the}%
\TextAt{8}{174.174}{709.986}{statement}%
\TextAt{8}{228.546}{709.986}{(4.4)}%
\TextAt{8}{256.521}{709.986}{does}%
\TextAt{8}{282.918}{709.986}{remain}%
\TextAt{8}{321.94}{709.986}{true}%
\TextAt{8}{346.651}{709.986}{in}%
\TextAt{8}{360.304}{709.986}{the}%
\TextAt{8}{380.461}{709.986}{case}%
\TextAt{8}{405.232}{709.986}{of}%
\TextAt{8}{418.562}{709.986}{equality.}%
\TextAt{12}{79.226}{679.473}{Exercise}%
\TextAt{12}{132.427}{679.473}{I.4.3.}%
\TextAt{8}{166.021}{679.473}{(}%
\TextAt{12}{170.576}{679.473}{Minkowski’s}%
\TextAt{12}{247.041}{679.473}{Theorem}%
\TextAt{12}{304.055}{679.473}{on}%
\TextAt{12}{322.741}{679.473}{Linear}%
\TextAt{12}{364.692}{679.473}{Forms}%
\TextAt{8}{400.772}{679.473}{).}%
\TextAt{8}{413.776}{679.473}{Let}%
\TextAt{13}{195.496}{646.536}{L}%
\TextAt{21}{203.461}{644.742}{i}%
\TextAt{9}{206.842}{646.536}{\Character{40}}%
\TextAt{13}{211.394}{646.536}{x}%
\TextAt{17}{218.046}{644.742}{1}%
\TextAt{13}{222.779}{646.536}{\Character{59}}%
\TextAt{13}{228.003}{646.536}{\Character{58}}%
\TextAt{13}{233.24}{646.536}{\Character{58}}%
\TextAt{13}{238.464}{646.536}{\Character{58}}%
\TextAt{13}{243.701}{646.536}{\Character{59}}%
\TextAt{13}{248.937}{646.536}{x}%
\TextAt{21}{255.652}{644.742}{n}%
\TextAt{9}{261.288}{646.536}{\Character{41}}%
\TextAt{9}{269.143}{646.536}{\Character{61}}%
\TextAt{21}{287.652}{661.48}{n}%
\TextAt{20}{281.587}{657.893}{\Character{88}}%
\TextAt{21}{282.868}{632.699}{j}%
\TextAt{17}{286.753}{632.699}{\Character{61}1}%
\TextAt{13}{300.848}{646.536}{a}%
\TextAt{21}{306.993}{644.742}{ij}%
\TextAt{13}{314.258}{646.536}{x}%
\TextAt{21}{320.91}{644.742}{j}%
\TextAt{13}{325.292}{646.536}{\Character{59}}%
\TextAt{13}{354.427}{646.536}{i}%
\TextAt{9}{361.761}{646.536}{\Character{61}}%
\TextAt{9}{374.171}{646.536}{1}%
\TextAt{13}{380.039}{646.536}{\Character{59}}%
\TextAt{13}{385.263}{646.536}{\Character{58}}%
\TextAt{13}{390.5}{646.536}{\Character{58}}%
\TextAt{13}{395.736}{646.536}{\Character{58}}%
\TextAt{13}{400.961}{646.536}{\Character{59}}%
\TextAt{13}{406.197}{646.536}{n\Character{59}}%
\TextAt{8}{79.226}{610.621}{be}%
\TextAt{8}{96.066}{610.621}{real}%
\TextAt{8}{119.744}{610.621}{linear}%
\TextAt{8}{153.178}{610.621}{forms}%
\TextAt{8}{186.348}{610.621}{such}%
\TextAt{8}{213.661}{610.621}{that}%
\TextAt{8}{239.945}{610.621}{det}%
\TextAt{9}{256.204}{610.621}{\Character{40}}%
\TextAt{13}{260.757}{610.621}{a}%
\TextAt{21}{266.901}{608.828}{ij}%
\TextAt{9}{274.167}{610.621}{\Character{41}}%
\NativeMathAt{283.6}{610.621}{11.955}{\neq}%
\TextAt{9}{297.576}{610.621}{0}%
\TextAt{8}{303.439}{610.621}{,}%
\TextAt{8}{311.732}{610.621}{and}%
\TextAt{8}{335.41}{610.621}{let}%
\TextAt{13}{353.231}{610.621}{c}%
\TextAt{17}{358.269}{608.828}{1}%
\TextAt{13}{363.001}{610.621}{\Character{59}}%
\TextAt{13}{368.225}{610.621}{\Character{58}}%
\TextAt{13}{373.462}{610.621}{\Character{58}}%
\TextAt{13}{378.686}{610.621}{\Character{58}}%
\TextAt{13}{383.923}{610.621}{\Character{59}}%
\TextAt{13}{389.159}{610.621}{c}%
\TextAt{21}{394.26}{608.828}{n}%
\TextAt{8}{404.709}{610.621}{be}%
\TextAt{8}{421.549}{610.621}{positive}%
\TextAt{8}{465.778}{610.621}{real}%
\TextAt{8}{489.456}{610.621}{numbers}%
\TextAt{8}{79.226}{593.575}{such}%
\TextAt{8}{105.623}{593.575}{that}%
\TextAt{13}{130.992}{593.575}{c}%
\TextAt{17}{136.03}{591.782}{1}%
\TextAt{16}{142.755}{593.575}{\Character{1}}%
\TextAt{16}{148.051}{593.575}{\Character{1}}%
\TextAt{16}{153.359}{593.575}{\Character{1}}%
\TextAt{13}{158.695}{593.575}{c}%
\TextAt{21}{163.733}{591.782}{n}%
\TextAt{13}{172.69}{593.575}{\Character{62}}%
\TextAt{16}{185.116}{593.575}{\Character{106}}%
\TextAt{8}{190.429}{593.575}{det}%
\TextAt{9}{206.688}{593.575}{\Character{40}}%
\TextAt{13}{211.24}{593.575}{a}%
\TextAt{21}{217.385}{591.782}{ij}%
\TextAt{9}{224.651}{593.575}{\Character{41}}%
\TextAt{16}{229.203}{593.575}{\Character{106}}%
\TextAt{8}{232.524}{593.575}{.}%
\ErrataAt{240.91}{593.575}{294.09}{Find $0\ne(m_1,\ldots,m_n)\in\mathbb Z^n$ such that}%
\TextAt{16}{206.231}{567.16}{\Character{106}}%
\TextAt{13}{209.552}{567.16}{L}%
\TextAt{21}{217.517}{565.366}{i}%
\TextAt{9}{220.898}{567.16}{\Character{40}}%
\TextAt{13}{225.45}{567.16}{m}%
\TextAt{17}{235.69}{565.366}{1}%
\TextAt{13}{240.422}{567.16}{\Character{59}}%
\TextAt{13}{245.646}{567.16}{\Character{58}}%
\TextAt{13}{250.883}{567.16}{\Character{58}}%
\TextAt{13}{256.107}{567.16}{\Character{58}}%
\TextAt{13}{261.344}{567.16}{\Character{59}}%
\TextAt{13}{266.58}{567.16}{m}%
\TextAt{21}{276.882}{565.366}{n}%
\TextAt{9}{282.518}{567.16}{\Character{41}}%
\TextAt{16}{287.071}{567.16}{\Character{106}}%
\TextAt{13}{293.712}{567.16}{\Character{60}}%
\TextAt{13}{306.122}{567.16}{c}%
\TextAt{21}{311.176}{565.366}{i}%
\TextAt{13}{314.557}{567.16}{\Character{59}}%
\TextAt{13}{343.692}{567.16}{i}%
\TextAt{9}{351.026}{567.16}{\Character{61}}%
\TextAt{9}{363.435}{567.16}{1}%
\TextAt{13}{369.304}{567.16}{\Character{59}}%
\TextAt{13}{374.528}{567.16}{\Character{58}}%
\TextAt{13}{379.765}{567.16}{\Character{58}}%
\TextAt{13}{385.001}{567.16}{\Character{58}}%
\TextAt{13}{390.226}{567.16}{\Character{59}}%
\TextAt{13}{395.462}{567.16}{n\Character{58}}%
\TextAt{12}{97.159}{537.755}{Hint:}%
\TextAt{8}{132.698}{537.755}{Use}%
\TextAt{8}{155.185}{537.755}{Minkowski’s}%
\TextAt{8}{221.321}{537.755}{lattice}%
\TextAt{8}{257.091}{537.755}{point}%
\TextAt{8}{287.661}{537.755}{theorem.}%
\TextAt{15}{79.226}{486.013}{§}%
\TextAt{15}{93.785}{486.013}{5}%
\TextAt{15}{120.693}{486.013}{Minkowski}%
\TextAt{15}{217.771}{486.013}{Theory}%
\TextAt{8}{79.226}{454.127}{The}%
\TextAt{8}{104.167}{454.127}{basic}%
\TextAt{8}{134.379}{454.127}{idea}%
\TextAt{8}{159.978}{454.127}{in}%
\TextAt{8}{174.517}{454.127}{Minkowski’s}%
\TextAt{8}{241.54}{454.127}{treatment}%
\TextAt{8}{296.739}{454.127}{of}%
\TextAt{8}{310.956}{454.127}{an}%
\TextAt{8}{328.101}{454.127}{algebraic}%
\TextAt{8}{378.423}{454.127}{number}%
\TextAt{8}{421.906}{454.127}{field}%
\TextAt{13}{448.15}{454.127}{K}%
\TextAt{16}{463.799}{454.127}{\Character{106}}%
\TextAt{18}{471.952}{454.127}{Q}%
\TextAt{8}{486.035}{454.127}{of}%
\TextAt{8}{500.251}{454.127}{degree}%
\TextAt{13}{79.226}{437.081}{n}%
\TextAt{8}{91.309}{437.081}{is}%
\TextAt{8}{104.271}{437.081}{to}%
\TextAt{8}{119.779}{437.081}{interpret}%
\TextAt{8}{169.432}{437.081}{its}%
\TextAt{8}{186.949}{437.081}{numbers}%
\TextAt{8}{235.358}{437.081}{as}%
\TextAt{8}{250.926}{437.081}{points}%
\TextAt{8}{287.308}{437.081}{in}%
\TextAt{13}{302.159}{437.081}{n}%
\TextAt{8}{309.146}{437.081}{-dimensional}%
\TextAt{8}{378.692}{437.081}{space.}%
\TextAt{8}{418.112}{437.081}{This}%
\TextAt{8}{446.03}{437.081}{explains}%
\TextAt{8}{492.49}{437.081}{why}%
\TextAt{8}{518.399}{437.081}{his}%
\TextAt{8}{79.226}{420.035}{theory}%
\TextAt{8}{115.882}{420.035}{has}%
\TextAt{8}{136.661}{420.035}{been}%
\TextAt{8}{164.194}{420.035}{called}%
\TextAt{8}{197.263}{420.035}{“Geometry}%
\TextAt{8}{257.076}{420.035}{of}%
\TextAt{8}{270.311}{420.035}{Numbers.}%
\TextAt{8}{317.845}{420.035}{”}%
\TextAt{8}{328.533}{420.035}{It}%
\TextAt{8}{341.111}{420.035}{seems}%
\TextAt{8}{374.299}{420.035}{appropriate,}%
\TextAt{8}{440.574}{420.035}{however,}%
\TextAt{8}{488.63}{420.035}{to}%
\TextAt{8}{502.846}{420.035}{follow}%
\TextAt{8}{79.226}{402.989}{the}%
\TextAt{8}{100.543}{402.989}{current}%
\TextAt{8}{142.351}{402.989}{trend}%
\TextAt{8}{174.726}{402.989}{and}%
\TextAt{8}{198.649}{402.989}{call}%
\TextAt{8}{221.269}{402.989}{it}%
\TextAt{8}{234.133}{402.989}{“Minkowski}%
\TextAt{8}{299.086}{402.989}{Theory”}%
\TextAt{8}{346.418}{402.989}{instead,}%
\TextAt{8}{391.504}{402.989}{because}%
\TextAt{8}{435.966}{402.989}{in}%
\TextAt{8}{450.78}{402.989}{the}%
\TextAt{8}{472.096}{402.989}{meantime}%
\TextAt{8}{526.911}{402.989}{a}%
\TextAt{8}{79.226}{385.943}{geometric}%
\TextAt{8}{131.443}{385.943}{approach}%
\TextAt{8}{180.742}{385.943}{to}%
\TextAt{8}{193.937}{385.943}{number}%
\TextAt{8}{235.418}{385.943}{theory}%
\TextAt{8}{271.052}{385.943}{has}%
\TextAt{8}{290.811}{385.943}{been}%
\TextAt{8}{317.323}{385.943}{developed}%
\TextAt{8}{370.508}{385.943}{which}%
\TextAt{8}{402.879}{385.943}{is}%
\TextAt{8}{413.528}{385.943}{quite}%
\TextAt{8}{442.001}{385.943}{different}%
\TextAt{8}{487.057}{385.943}{in}%
\TextAt{8}{499.594}{385.943}{nature}%
\TextAt{8}{79.226}{368.897}{and}%
\TextAt{8}{102.607}{368.897}{much}%
\TextAt{8}{134.44}{368.897}{more}%
\TextAt{8}{164.325}{368.897}{comprehensive.}%
\TextAt{8}{248.833}{368.897}{We}%
\TextAt{8}{269.596}{368.897}{will}%
\TextAt{8}{292.319}{368.897}{explain}%
\TextAt{8}{333.585}{368.897}{this}%
\TextAt{8}{357.025}{368.897}{in}%
\TextAt{8}{371.296}{368.897}{§13.}%
\TextAt{8}{398.994}{368.897}{In}%
\TextAt{8}{414.233}{368.897}{the}%
\TextAt{8}{435.008}{368.897}{present}%
\TextAt{8}{476.333}{368.897}{section,}%
\TextAt{8}{519.439}{368.897}{we}%
\TextAt{8}{79.226}{351.851}{consider}%
\TextAt{8}{124.811}{351.851}{the}%
\TextAt{8}{144.968}{351.851}{canonical}%
\TextAt{8}{196.351}{351.851}{mapping}%
\TextAt{13}{192.611}{325.435}{j}%
\TextAt{9}{201.448}{325.435}{\Character{58}}%
\TextAt{13}{208.02}{325.435}{K}%
\NativeArrowAt{222.159}{325.435}{11.955}{19.26}{\longrightarrow}%
\TextAt{13}{244.741}{325.435}{K}%
\TextAt{25}{254.7}{323.642}{C}%
\TextAt{9}{264.275}{325.435}{\Character{58}\Character{61}}%
\TextAt{20}{279.952}{336.793}{\Character{89}}%
\TextAt{21}{285.251}{311.687}{\Character{28}}%
\TextAt{18}{297.221}{325.435}{C}%
\TextAt{13}{305.855}{325.435}{\Character{59}}%
\TextAt{13}{334.99}{325.435}{a}%
\NativeArrowAt{344.475}{325.435}{11.955}{19.26}{\longmapsto}%
\TextAt{13}{367.057}{325.435}{j}%
\TextAt{13}{372.568}{325.435}{a}%
\TextAt{9}{382.039}{325.435}{\Character{61}}%
\TextAt{9}{394.449}{325.435}{\Character{40}}%
\TextAt{13}{399.016}{325.435}{\Character{28}}%
\TextAt{13}{405.424}{325.435}{a}%
\TextAt{9}{411.58}{325.435}{\Character{41}}%
\TextAt{13}{416.133}{325.435}{\Character{59}}%
\TextAt{8}{79.226}{291.709}{which}%
\TextAt{8}{113.71}{291.709}{results}%
\TextAt{8}{151.901}{291.709}{from}%
\TextAt{8}{180.539}{291.709}{the}%
\TextAt{13}{201.693}{291.709}{n}%
\TextAt{8}{213.576}{291.709}{complex}%
\TextAt{8}{260.422}{291.709}{embeddings}%
\TextAt{13}{325.212}{291.709}{\Character{28}}%
\TextAt{9}{336.652}{291.709}{\Character{58}}%
\TextAt{13}{344.925}{291.709}{K}%
\NativeMathAt{360.764}{291.709}{11.955}{\rightarrow}%
\TextAt{18}{377.74}{291.709}{C}%
\TextAt{8}{386.375}{291.709}{.}%
\TextAt{8}{397.816}{291.709}{The}%
\TextAt{18}{422.868}{291.709}{C}%
\TextAt{8}{431.502}{291.709}{-vector}%
\TextAt{8}{471.522}{291.709}{space}%
\TextAt{13}{503.794}{291.709}{K}%
\TextAt{25}{513.753}{289.916}{C}%
\TextAt{8}{524.903}{291.709}{is}%
\TextAt{8}{79.226}{274.663}{equipped}%
\TextAt{8}{129.294}{274.663}{with}%
\TextAt{8}{155.954}{274.663}{the}%
\TextAt{12}{176.111}{274.663}{hermitian}%
\TextAt{12}{237.644}{274.663}{scalar}%
\TextAt{12}{275.733}{274.663}{product}%
\TextAt{16}{261.846}{248.247}{\Character{104}}%
\TextAt{13}{266.495}{248.247}{x\Character{59}}%
\TextAt{13}{278.367}{248.247}{y}%
\TextAt{16}{284.528}{248.247}{\Character{105}}%
\TextAt{9}{292.498}{248.247}{\Character{61}}%
\TextAt{20}{304.924}{259.605}{\Character{88}}%
\TextAt{21}{311.219}{234.499}{\Character{28}}%
\TextAt{13}{324.185}{248.247}{x}%
\TextAt{21}{330.837}{246.454}{\Character{28}}%
\TextAt{9}{336.806}{248.247}{\Character{22}}%
\TextAt{13}{336.014}{248.247}{y}%
\TextAt{21}{341.721}{246.454}{\Character{28}}%
\TextAt{13}{346.898}{248.247}{\Character{58}}%
\TextAt{8}{517.802}{248.247}{(*)}%
\TextAt{8}{79.226}{213.971}{Let}%
\TextAt{8}{99.936}{213.971}{us}%
\TextAt{8}{114.692}{213.971}{recall}%
\TextAt{8}{145.647}{213.971}{that}%
\TextAt{8}{170.756}{213.971}{a}%
\TextAt{8}{180.251}{213.971}{hermitian}%
\TextAt{8}{233.323}{213.971}{scalar}%
\TextAt{8}{266.299}{213.971}{product}%
\TextAt{8}{309.939}{213.971}{is}%
\TextAt{8}{321.443}{213.971}{given}%
\TextAt{8}{351.752}{213.971}{by}%
\TextAt{8}{367.751}{213.971}{a}%
\TextAt{8}{377.247}{213.971}{form}%
\TextAt{13}{404.627}{213.971}{H}%
\TextAt{9}{415.25}{213.971}{\Character{40}}%
\TextAt{13}{419.802}{213.971}{x\Character{59}}%
\TextAt{13}{431.673}{213.971}{y}%
\TextAt{9}{437.835}{213.971}{\Character{41}}%
\TextAt{8}{446.025}{213.971}{which}%
\TextAt{8}{479.252}{213.971}{is}%
\TextAt{8}{490.756}{213.971}{linear}%
\TextAt{8}{523.014}{213.971}{in}%
\TextAt{8}{79.226}{196.925}{the}%
\TextAt{8}{99.486}{196.925}{first}%
\TextAt{8}{123.715}{196.925}{variable}%
\TextAt{8}{167.729}{196.925}{and}%
\TextAt{8}{190.595}{196.925}{satisfies}%
\Vector{q 0 G 0 g 0.478 w 161.809 -566.435 m 199.569 -566.435 l S Q}%
\TextAt{13}{233.809}{196.925}{H}%
\TextAt{9}{244.431}{196.925}{\Character{40}}%
\TextAt{13}{248.984}{196.925}{x\Character{59}}%
\TextAt{13}{260.856}{196.925}{y}%
\TextAt{9}{267.017}{196.925}{\Character{41}}%
\TextAt{9}{275.051}{196.925}{\Character{61}}%
\TextAt{13}{287.667}{196.925}{H}%
\TextAt{9}{298.29}{196.925}{\Character{40}}%
\TextAt{13}{302.842}{196.925}{y}%
\TextAt{13}{308.963}{196.925}{\Character{59}}%
\TextAt{13}{314.199}{196.925}{x}%
\TextAt{9}{320.875}{196.925}{\Character{41}}%
\TextAt{8}{329.428}{196.925}{as}%
\TextAt{8}{343.901}{196.925}{well}%
\TextAt{8}{367.736}{196.925}{as}%
\TextAt{13}{382.209}{196.925}{H}%
\TextAt{9}{392.832}{196.925}{\Character{40}}%
\TextAt{13}{397.384}{196.925}{x\Character{59}}%
\TextAt{13}{409.256}{196.925}{x}%
\TextAt{9}{415.932}{196.925}{\Character{41}}%
\TextAt{13}{423.981}{196.925}{\Character{62}}%
\TextAt{9}{436.583}{196.925}{0}%
\TextAt{8}{446.436}{196.925}{for}%
\TextAt{13}{464.425}{196.925}{x}%
\NativeMathAt{474.573}{196.925}{11.955}{\neq}%
\TextAt{9}{487.162}{196.925}{0}%
\TextAt{8}{493.028}{196.925}{.}%
\TextAt{8}{501.786}{196.925}{In}%
\TextAt{8}{516.51}{196.925}{the}%
\TextAt{8}{79.226}{179.879}{sequel}%
\TextAt{8}{114.871}{179.879}{we}%
\TextAt{8}{132.893}{179.879}{always}%
\TextAt{8}{171.156}{179.879}{view}%
\TextAt{13}{198.934}{179.879}{K}%
\TextAt{25}{208.893}{178.086}{C}%
\TextAt{8}{219.84}{179.879}{as}%
\TextAt{8}{235.005}{179.879}{a}%
\TextAt{8}{245.556}{179.879}{hermitian}%
\TextAt{8}{299.683}{179.879}{space,}%
\TextAt{8}{335.204}{179.879}{with}%
\TextAt{8}{362.659}{179.879}{respect}%
\TextAt{8}{403.504}{179.879}{to}%
\TextAt{8}{418.609}{179.879}{the}%
\TextAt{8}{439.561}{179.879}{“standard}%
\TextAt{8}{494.728}{179.879}{metric”}%
\TextAt{8}{79.226}{162.833}{(*).}%
\TextAt{8}{97.159}{142.798}{The}%
\TextAt{8}{121.213}{142.798}{Galois}%
\TextAt{8}{157.126}{142.798}{group}%
\TextAt{13}{190.302}{142.798}{G}%
\TextAt{9}{199.535}{142.798}{\Character{40}}%
\TextAt{18}{204.088}{142.798}{C}%
\TextAt{16}{216.043}{142.798}{\Character{106}}%
\TextAt{18}{222.684}{142.798}{R}%
\TextAt{9}{231.319}{142.798}{\Character{41}}%
\TextAt{8}{239.768}{142.798}{is}%
\TextAt{8}{251.532}{142.798}{generated}%
\TextAt{8}{304.864}{142.798}{by}%
\TextAt{8}{321.123}{142.798}{complex}%
\TextAt{8}{366.971}{142.798}{conjugation}%
\TextAt{13}{275.901}{116.383}{F}%
\TextAt{9}{288.426}{116.383}{\Character{58}}%
\TextAt{13}{294.998}{116.383}{z}%
\NativeArrowAt{304.29}{116.383}{11.955}{19.26}{\longmapsto}%
\TextAt{9}{327.581}{116.383}{\Character{22}}%
\TextAt{13}{326.872}{116.383}{z}%
\TextAt{13}{332.826}{116.383}{\Character{58}}%
\TextAt{8}{79.226}{89.967}{The}%
\TextAt{8}{103.888}{89.967}{notation}%
\TextAt{13}{151.337}{89.967}{F}%
\TextAt{8}{165.047}{89.967}{will}%
\TextAt{8}{187.761}{89.967}{be}%
\TextAt{8}{204.294}{89.967}{justified}%
\TextAt{8}{249.507}{89.967}{only}%
\TextAt{8}{275.807}{89.967}{later}%
\TextAt{8}{303.733}{89.967}{(see}%
\TextAt{8}{327.81}{89.967}{chap.}%
\TextAt{8}{361.827}{89.967}{III,}%
\TextAt{8}{382.891}{89.967}{§4).}%
\TextAt{13}{409.256}{89.967}{F}%
\TextAt{8}{422.966}{89.967}{acts}%
\TextAt{8}{447.7}{89.967}{on}%
\TextAt{8}{464.568}{89.967}{the}%
\TextAt{8}{485.333}{89.967}{one}%
\TextAt{8}{507.401}{89.967}{hand}%
\TextAt{8}{79.226}{72.921}{on}%
\TextAt{8}{96.105}{72.921}{the}%
\TextAt{8}{116.882}{72.921}{factors}%
\TextAt{8}{155.615}{72.921}{of}%
\TextAt{8}{169.565}{72.921}{the}%
\TextAt{8}{190.341}{72.921}{product}%
\TextAt{20}{234.861}{81.888}{\Character{81}}%
\TextAt{21}{246.152}{69.434}{\Character{28}}%
\TextAt{18}{253.321}{72.921}{C}%
\TextAt{8}{261.955}{72.921}{,}%
\TextAt{8}{269.879}{72.921}{but}%
\TextAt{8}{291.959}{72.921}{on}%
\TextAt{8}{308.838}{72.921}{the}%
\TextAt{8}{329.615}{72.921}{other}%
\TextAt{8}{360.804}{72.921}{hand}%
\TextAt{8}{390.69}{72.921}{it}%
\TextAt{8}{403.015}{72.921}{also}%
\TextAt{8}{427.115}{72.921}{acts}%
\TextAt{8}{451.86}{72.921}{on}%
\TextAt{8}{468.739}{72.921}{the}%
\TextAt{8}{489.516}{72.921}{indexing}%
\TextAt{8}{300.14}{43.033}{26}%
\Destination{522.36}{265.295}{nav.AMS.7}%
\Destination{78.23}{809.815}{nav.page.26}%
\Destination{79.23}{512.545}{nav.section.1.5}%
\Destination{79.226}{775.155}{nav.ex.1.4.2}%
\Destination{79.226}{693.455}{nav.ex.1.4.3}%
\Destination{79.36}{694.455}{nav.print.28}%
\Destination{79.23}{512.545}{outline.10}%
\Hotlink{376.475}{366.655}{12.716}{11.047}{nav.section.1.13}%
\Hotlink{361.327}{87.755}{14.294}{11.047}{nav.chapter.3}%
\Hotlink{388.07}{87.755}{6.858}{11.047}{nav.section.3.4}%
\Destination{79.226}{610.575}{errata.63}%
}

% Page 39 (body page 27)
\Page{%
\TextAt{2}{79.226}{796.943}{I.}%
\TextAt{2}{89.83}{796.943}{Algebraic}%
\TextAt{2}{138.648}{796.943}{Integers}%
\Vector{q 0 G 0 g 0.299 w 7.226 18.729 m 460.769 18.729 l S Q}%
\TextAt{8}{79.226}{761.124}{set}%
\TextAt{8}{97.711}{761.124}{of}%
\TextAt{13}{111.257}{761.124}{\Character{28}}%
\TextAt{8}{117.676}{761.124}{’s;}%
\TextAt{8}{133.017}{761.124}{to}%
\TextAt{8}{147.545}{761.124}{each}%
\TextAt{8}{174.099}{761.124}{embedding}%
\TextAt{13}{233.494}{761.124}{\Character{28}}%
\TextAt{9}{243.603}{761.124}{\Character{58}}%
\TextAt{13}{250.545}{761.124}{K}%
\NativeMathAt{265.053}{761.124}{11.955}{\rightarrow}%
\TextAt{18}{280.699}{761.124}{C}%
\TextAt{8}{293.447}{761.124}{corresponds}%
\TextAt{8}{357.851}{761.124}{its}%
\TextAt{8}{374.387}{761.124}{complex}%
\TextAt{8}{420.452}{761.124}{conjugate}%
\TextAt{9}{474.286}{761.124}{\Character{22}}%
\TextAt{13}{473.678}{761.124}{\Character{28}}%
\TextAt{9}{483.788}{761.124}{\Character{58}}%
\TextAt{13}{490.73}{761.124}{K}%
\NativeMathAt{505.238}{761.124}{11.955}{\rightarrow}%
\TextAt{18}{520.883}{761.124}{C}%
\TextAt{8}{529.518}{761.124}{.}%
\TextAt{8}{79.226}{744.078}{Altogether,}%
\TextAt{8}{140.688}{744.078}{this}%
\TextAt{8}{163.51}{744.078}{defines}%
\TextAt{8}{201.934}{744.078}{an}%
\TextAt{8}{218.193}{744.078}{involution}%
\TextAt{13}{267.284}{715.214}{F}%
\TextAt{9}{279.809}{715.214}{\Character{58}}%
\TextAt{13}{286.381}{715.214}{K}%
\TextAt{25}{296.341}{713.421}{C}%
\NativeArrowAt{305.916}{715.214}{11.955}{19.26}{\longrightarrow}%
\TextAt{13}{328.498}{715.214}{K}%
\TextAt{25}{338.457}{713.421}{C}%
\TextAt{8}{79.226}{686.351}{which,}%
\TextAt{8}{115.964}{686.351}{on}%
\TextAt{8}{132.223}{686.351}{the}%
\TextAt{8}{152.38}{686.351}{points}%
\TextAt{13}{187.564}{686.351}{z}%
\TextAt{9}{196.855}{686.351}{\Character{61}}%
\TextAt{9}{209.265}{686.351}{\Character{40}}%
\TextAt{13}{213.833}{686.351}{z}%
\TextAt{21}{219.271}{684.557}{\Character{28}}%
\TextAt{9}{224.448}{686.351}{\Character{41}}%
\NativeMathAt{232.321}{686.351}{11.955}{\in}%
\TextAt{13}{243.612}{686.351}{K}%
\TextAt{25}{253.571}{684.557}{C}%
\TextAt{8}{259.825}{686.351}{,}%
\TextAt{8}{266.974}{686.351}{is}%
\TextAt{8}{278.738}{686.351}{given}%
\TextAt{8}{309.308}{686.351}{by}%
\TextAt{9}{276.464}{657.487}{\Character{40}}%
\TextAt{13}{281.016}{657.487}{F}%
\TextAt{13}{290.209}{657.487}{z}%
\TextAt{9}{296.191}{657.487}{\Character{41}}%
\TextAt{21}{300.743}{655.694}{\Character{28}}%
\TextAt{9}{309.24}{657.487}{\Character{61}}%
\TextAt{9}{322.367}{657.487}{\Character{22}}%
\TextAt{13}{321.666}{657.487}{z}%
\TextAt{17}{327.561}{655.694}{\Character{22}}%
\TextAt{21}{327.104}{655.694}{\Character{28}}%
\TextAt{13}{332.28}{657.487}{\Character{58}}%
\TextAt{8}{79.226}{628.624}{The}%
\TextAt{8}{103.28}{628.624}{scalar}%
\TextAt{8}{136.515}{628.624}{product}%
\TextAt{16}{180.415}{628.624}{\Character{104}}%
\TextAt{13}{188.961}{628.624}{\Character{59}}%
\TextAt{16}{198.103}{628.624}{\Character{105}}%
\TextAt{8}{206.65}{628.624}{is}%
\TextAt{8}{218.413}{628.624}{equivariant}%
\TextAt{8}{279.241}{628.624}{under}%
\TextAt{13}{312.405}{628.624}{F}%
\TextAt{8}{321.609}{628.624}{,}%
\TextAt{8}{328.758}{628.624}{that}%
\TextAt{8}{354.127}{628.624}{is}%
\TextAt{16}{255.362}{599.76}{\Character{104}}%
\TextAt{13}{260.011}{599.76}{F}%
\TextAt{13}{269.204}{599.76}{x\Character{59}}%
\TextAt{13}{281.076}{599.76}{F}%
\TextAt{13}{290.27}{599.76}{y}%
\TextAt{16}{296.451}{599.76}{\Character{105}}%
\TextAt{9}{304.422}{599.76}{\Character{61}}%
\TextAt{13}{316.847}{599.76}{F}%
\TextAt{16}{326.051}{599.76}{\Character{104}}%
\TextAt{13}{330.7}{599.76}{x\Character{59}}%
\TextAt{13}{342.572}{599.76}{y}%
\TextAt{16}{348.733}{599.76}{\Character{105}}%
\TextAt{13}{353.382}{599.76}{\Character{58}}%
\TextAt{8}{79.226}{570.897}{Finally,}%
\TextAt{8}{121.332}{570.897}{we}%
\TextAt{8}{138.56}{570.897}{have}%
\TextAt{8}{165.554}{570.897}{on}%
\TextAt{8}{181.813}{570.897}{the}%
\TextAt{18}{201.97}{570.897}{C}%
\TextAt{8}{210.604}{570.897}{-vector}%
\TextAt{8}{249.626}{570.897}{space}%
\TextAt{13}{280.901}{570.897}{K}%
\TextAt{25}{290.86}{569.103}{C}%
\TextAt{9}{300.435}{570.897}{\Character{61}}%
\TextAt{20}{312.86}{579.863}{\Character{81}}%
\TextAt{21}{324.151}{567.41}{\Character{28}}%
\TextAt{18}{331.32}{570.897}{C}%
\TextAt{8}{343.852}{570.897}{the}%
\TextAt{8}{364.008}{570.897}{linear}%
\TextAt{8}{396.527}{570.897}{map}%
\TextAt{8}{267.546}{542.033}{Tr}%
\TextAt{9}{283.874}{542.033}{\Character{58}}%
\TextAt{13}{290.447}{542.033}{K}%
\TextAt{25}{300.406}{540.24}{C}%
\NativeArrowAt{309.981}{542.033}{11.955}{19.26}{\longrightarrow}%
\TextAt{18}{332.563}{542.033}{C}%
\TextAt{13}{341.198}{542.033}{\Character{59}}%
\TextAt{8}{79.226}{513.17}{given}%
\TextAt{8}{109.796}{513.17}{as}%
\TextAt{8}{124.166}{513.17}{the}%
\TextAt{8}{144.322}{513.17}{sum}%
\TextAt{8}{169.093}{513.17}{of}%
\TextAt{8}{182.423}{513.17}{the}%
\TextAt{8}{202.58}{513.17}{coordinates.}%
\TextAt{8}{269.309}{513.17}{It}%
\TextAt{8}{281.982}{513.17}{is}%
\TextAt{8}{293.746}{513.17}{also}%
\TextAt{13}{317.226}{513.17}{F}%
\TextAt{8}{326.429}{513.17}{-invariant.}%
\TextAt{8}{384.001}{513.17}{The}%
\TextAt{8}{408.055}{513.17}{composite}%
\TextAt{13}{267.546}{484.306}{K}%
\TextAt{21}{285.015}{491.57}{j}%
\NativeArrowAt{281.685}{484.306}{11.955}{11.955}{\longrightarrow}%
\TextAt{13}{296.961}{484.306}{K}%
\TextAt{25}{306.92}{482.513}{C}%
\TextAt{7}{319.083}{491.083}{Tr}%
\NativeArrowAt{316.495}{484.306}{11.955}{15.996}{\longrightarrow}%
\TextAt{18}{335.815}{484.306}{C}%
\TextAt{8}{79.226}{455.443}{gives}%
\TextAt{8}{107.907}{455.443}{the}%
\TextAt{8}{128.063}{455.443}{usual}%
\TextAt{8}{158.692}{455.443}{trace}%
\TextAt{8}{187.958}{455.443}{of}%
\TextAt{13}{201.288}{455.443}{K}%
\TextAt{16}{215.427}{455.443}{\Character{106}}%
\TextAt{18}{222.069}{455.443}{Q}%
\TextAt{8}{235.265}{455.443}{(see}%
\TextAt{8}{258.733}{455.443}{(2.6),}%
\TextAt{8}{289.959}{455.443}{(ii)),}%
\TextAt{8}{257.107}{426.579}{Tr}%
\TextAt{21}{270.114}{424.719}{K}%
\TextAt{3}{277.819}{424.719}{\Character{106}}%
\TextAt{25}{280.171}{424.719}{Q}%
\TextAt{9}{286.868}{426.579}{\Character{40}}%
\TextAt{13}{291.421}{426.579}{a}%
\TextAt{9}{297.566}{426.579}{\Character{41}}%
\TextAt{9}{305.421}{426.579}{\Character{61}}%
\TextAt{8}{317.864}{426.579}{Tr}%
\TextAt{9}{330.872}{426.579}{\Character{40}}%
\TextAt{13}{335.424}{426.579}{j}%
\TextAt{13}{340.935}{426.579}{a}%
\TextAt{9}{347.084}{426.579}{\Character{41}}%
\TextAt{13}{351.637}{426.579}{\Character{58}}%
\TextAt{8}{79.226}{397.716}{We}%
\TextAt{8}{99.371}{397.716}{now}%
\TextAt{8}{123.759}{397.716}{concentrate}%
\TextAt{8}{186.524}{397.716}{on}%
\TextAt{8}{202.783}{397.716}{the}%
\TextAt{18}{222.939}{397.716}{R}%
\TextAt{8}{231.574}{397.716}{-vector}%
\TextAt{8}{270.595}{397.716}{space}%
\TextAt{13}{249.726}{360}{K}%
\TextAt{25}{259.685}{358.207}{R}%
\TextAt{9}{269.26}{360}{\Character{61}}%
\TextAt{13}{281.686}{360}{K}%
\TextAt{17}{292.503}{364.948}{\Character{43}}%
\TextAt{25}{291.645}{356.715}{C}%
\TextAt{9}{302.909}{360}{\Character{61}}%
\TextAt{20}{315.334}{380.444}{\Character{34}}%
\TextAt{20}{322.308}{371.358}{\Character{89}}%
\TextAt{21}{327.607}{346.252}{\Character{28}}%
\TextAt{18}{339.577}{360}{C}%
\TextAt{20}{348.211}{380.444}{\Character{35}}%
\TextAt{17}{355.185}{377.767}{\Character{43}}%
\TextAt{8}{79.226}{323.276}{consisting}%
\TextAt{8}{133.198}{323.276}{of}%
\TextAt{8}{146.39}{323.276}{the}%
\TextAt{13}{166.409}{323.276}{G}%
\TextAt{9}{175.643}{323.276}{\Character{40}}%
\TextAt{18}{180.195}{323.276}{C}%
\TextAt{16}{192.15}{323.276}{\Character{106}}%
\TextAt{18}{198.792}{323.276}{R}%
\TextAt{9}{207.426}{323.276}{\Character{41}}%
\TextAt{8}{211.979}{323.276}{-invariant,}%
\TextAt{8}{268.142}{323.276}{i.e.,}%
\TextAt{13}{290.137}{323.276}{F}%
\TextAt{8}{299.341}{323.276}{-invariant,}%
\TextAt{8}{355.504}{323.276}{points}%
\TextAt{8}{390.55}{323.276}{of}%
\TextAt{13}{403.743}{323.276}{K}%
\TextAt{25}{413.702}{321.483}{C}%
\TextAt{8}{419.956}{323.276}{.}%
\TextAt{8}{428.359}{323.276}{These}%
\TextAt{8}{462.09}{323.276}{are}%
\TextAt{8}{481.464}{323.276}{the}%
\TextAt{8}{501.483}{323.276}{points}%
\TextAt{9}{79.226}{306.23}{\Character{40}}%
\TextAt{13}{83.778}{306.23}{z}%
\TextAt{21}{89.216}{304.437}{\Character{28}}%
\TextAt{9}{94.393}{306.23}{\Character{41}}%
\TextAt{8}{103.073}{306.23}{such}%
\TextAt{8}{129.701}{306.23}{that}%
\TextAt{13}{155.301}{306.23}{z}%
\TextAt{17}{161.196}{304.437}{\Character{22}}%
\TextAt{21}{160.739}{304.437}{\Character{28}}%
\TextAt{9}{169.63}{306.23}{\Character{61}}%
\TextAt{9}{183.139}{306.23}{\Character{22}}%
\TextAt{13}{182.449}{306.23}{z}%
\TextAt{21}{187.887}{304.437}{\Character{28}}%
\TextAt{8}{193.063}{306.23}{.}%
\TextAt{8}{202.204}{306.23}{An}%
\TextAt{8}{221.611}{306.23}{explicit}%
\TextAt{8}{263.135}{306.23}{description}%
\TextAt{8}{323.262}{306.23}{of}%
\TextAt{13}{336.823}{306.23}{K}%
\TextAt{25}{346.782}{304.437}{R}%
\TextAt{8}{357.164}{306.23}{will}%
\TextAt{8}{379.5}{306.23}{be}%
\TextAt{8}{395.655}{306.23}{given}%
\TextAt{8}{426.456}{306.23}{anon.}%
\TextAt{8}{460.32}{306.23}{Since}%
\Vector{q 0 G 0 g 0.478 w 419.108 -460.949 m 431.672 -460.949 l S Q}%
\TextAt{13}{491.108}{306.23}{\Character{28}}%
\TextAt{13}{497.516}{306.23}{a}%
\TextAt{9}{507.387}{306.23}{\Character{61}}%
\TextAt{9}{520.801}{306.23}{\Character{22}}%
\TextAt{13}{520.206}{306.23}{\Character{28}}%
\TextAt{13}{526.614}{306.23}{a}%
\TextAt{8}{79.226}{289.184}{for}%
\TextAt{13}{97.111}{289.184}{a}%
\NativeMathAt{106.577}{289.184}{11.955}{\in}%
\TextAt{13}{117.868}{289.184}{K}%
\TextAt{8}{128.685}{289.184}{,}%
\TextAt{8}{135.835}{289.184}{one}%
\TextAt{8}{157.294}{289.184}{has}%
\TextAt{13}{178.168}{289.184}{F}%
\TextAt{9}{187.372}{289.184}{\Character{40}}%
\TextAt{13}{191.924}{289.184}{j}%
\TextAt{13}{197.435}{289.184}{a}%
\TextAt{9}{203.584}{289.184}{\Character{41}}%
\TextAt{9}{211.439}{289.184}{\Character{61}}%
\TextAt{13}{223.883}{289.184}{j}%
\TextAt{13}{229.394}{289.184}{a}%
\TextAt{8}{235.543}{289.184}{.}%
\TextAt{8}{243.992}{289.184}{This}%
\TextAt{8}{270.712}{289.184}{yields}%
\TextAt{8}{303.612}{289.184}{a}%
\TextAt{8}{313.368}{289.184}{mapping}%
\TextAt{13}{270.201}{260.32}{j}%
\TextAt{9}{279.037}{260.32}{\Character{58}}%
\TextAt{13}{285.609}{260.32}{K}%
\NativeArrowAt{299.748}{260.32}{11.955}{19.26}{\longrightarrow}%
\TextAt{13}{322.33}{260.32}{K}%
\TextAt{25}{332.289}{258.527}{R}%
\TextAt{13}{338.543}{260.32}{\Character{58}}%
\TextAt{8}{79.226}{231.457}{The}%
\TextAt{8}{102.66}{231.457}{restriction}%
\TextAt{8}{158.038}{231.457}{of}%
\TextAt{8}{170.748}{231.457}{the}%
\TextAt{8}{190.284}{231.457}{hermitian}%
\TextAt{8}{242.996}{231.457}{scalar}%
\TextAt{8}{275.611}{231.457}{product}%
\TextAt{16}{318.89}{231.457}{\Character{104}}%
\TextAt{13}{326.817}{231.457}{\Character{59}}%
\TextAt{16}{335.338}{231.457}{\Character{105}}%
\TextAt{8}{343.265}{231.457}{from}%
\TextAt{13}{370.285}{231.457}{K}%
\TextAt{25}{380.244}{229.664}{C}%
\TextAt{8}{389.776}{231.457}{to}%
\TextAt{13}{403.466}{231.457}{K}%
\TextAt{25}{413.425}{229.664}{R}%
\TextAt{8}{422.957}{231.457}{gives}%
\TextAt{8}{451.017}{231.457}{a}%
\TextAt{8}{460.152}{231.457}{scalar}%
\TextAt{8}{492.767}{231.457}{product}%
\TextAt{9}{249.191}{202.593}{\Character{40}}%
\TextAt{13}{257.641}{202.593}{\Character{59}}%
\TextAt{9}{266.783}{202.593}{\Character{41}}%
\TextAt{9}{274.638}{202.593}{\Character{58}}%
\TextAt{13}{281.228}{202.593}{K}%
\TextAt{25}{291.187}{200.8}{R}%
\TextAt{16}{300.098}{202.593}{\Character{2}}%
\TextAt{13}{312.054}{202.593}{K}%
\TextAt{25}{322.013}{200.8}{R}%
\NativeArrowAt{331.588}{202.593}{11.955}{19.26}{\longrightarrow}%
\TextAt{18}{354.17}{202.593}{R}%
\TextAt{8}{79.226}{173.73}{on}%
\TextAt{8}{96.829}{173.73}{the}%
\TextAt{18}{118.329}{173.73}{R}%
\TextAt{8}{126.963}{173.73}{-vector}%
\TextAt{8}{167.329}{173.73}{space}%
\TextAt{13}{199.947}{173.73}{K}%
\TextAt{25}{209.906}{171.937}{R}%
\TextAt{8}{216.16}{173.73}{.}%
\TextAt{8}{228.64}{173.73}{Indeed,}%
\TextAt{8}{271.6}{173.73}{for}%
\TextAt{13}{290.829}{173.73}{x\Character{59}}%
\TextAt{13}{302.701}{173.73}{y}%
\NativeMathAt{314.472}{173.73}{11.955}{\in}%
\TextAt{13}{328.053}{173.73}{K}%
\TextAt{25}{338.012}{171.937}{R}%
\TextAt{8}{344.267}{173.73}{,}%
\TextAt{8}{353.095}{173.73}{one}%
\TextAt{8}{375.899}{173.73}{has}%
\TextAt{16}{398.116}{173.73}{\Character{104}}%
\TextAt{13}{402.765}{173.73}{x\Character{59}}%
\TextAt{13}{414.637}{173.73}{y}%
\TextAt{16}{420.798}{173.73}{\Character{105}}%
\NativeMathAt{431.044}{173.73}{11.955}{\in}%
\TextAt{18}{444.639}{173.73}{R}%
\TextAt{8}{458.514}{173.73}{in}%
\TextAt{8}{473.51}{173.73}{view}%
\TextAt{8}{501.837}{173.73}{of}%
\TextAt{8}{516.511}{173.73}{the}%
\TextAt{8}{79.226}{156.684}{relations}%
\TextAt{13}{126.427}{156.684}{F}%
\TextAt{16}{135.631}{156.684}{\Character{104}}%
\TextAt{13}{140.28}{156.684}{x\Character{59}}%
\TextAt{13}{152.152}{156.684}{y}%
\TextAt{16}{158.313}{156.684}{\Character{105}}%
\TextAt{9}{166.283}{156.684}{\Character{61}}%
\TextAt{16}{178.708}{156.684}{\Character{104}}%
\TextAt{13}{183.358}{156.684}{F}%
\TextAt{13}{192.54}{156.684}{x\Character{59}}%
\TextAt{13}{204.423}{156.684}{F}%
\TextAt{13}{213.617}{156.684}{y}%
\TextAt{16}{219.798}{156.684}{\Character{105}}%
\TextAt{9}{227.768}{156.684}{\Character{61}}%
\TextAt{16}{240.193}{156.684}{\Character{104}}%
\TextAt{13}{244.843}{156.684}{x\Character{59}}%
\TextAt{13}{256.715}{156.684}{y}%
\TextAt{16}{262.875}{156.684}{\Character{105}}%
\TextAt{8}{267.525}{156.684}{,}%
\TextAt{16}{274.398}{156.684}{\Character{104}}%
\TextAt{13}{279.047}{156.684}{x\Character{59}}%
\TextAt{13}{290.918}{156.684}{y}%
\TextAt{16}{297.08}{156.684}{\Character{105}}%
\TextAt{9}{305.05}{156.684}{\Character{61}}%
\Vector{q 0 G 0 g 0.478 w 245.476 -606.676 m 272.807 -606.676 l S Q}%
\TextAt{16}{317.476}{156.684}{\Character{104}}%
\TextAt{13}{322.125}{156.684}{x\Character{59}}%
\TextAt{13}{333.997}{156.684}{y}%
\TextAt{16}{340.158}{156.684}{\Character{105}}%
\TextAt{9}{348.128}{156.684}{\Character{61}}%
\TextAt{16}{360.553}{156.684}{\Character{104}}%
\TextAt{13}{365.203}{156.684}{y}%
\TextAt{13}{371.324}{156.684}{\Character{59}}%
\TextAt{13}{376.56}{156.684}{x}%
\TextAt{16}{383.235}{156.684}{\Character{105}}%
\TextAt{8}{387.885}{156.684}{,}%
\TextAt{8}{394.758}{156.684}{and,}%
\TextAt{8}{420.497}{156.684}{in}%
\TextAt{8}{433.804}{156.684}{any}%
\TextAt{8}{455.577}{156.684}{case,}%
\TextAt{16}{483.324}{156.684}{\Character{104}}%
\TextAt{13}{487.973}{156.684}{x\Character{59}}%
\TextAt{13}{499.845}{156.684}{x}%
\TextAt{16}{506.521}{156.684}{\Character{105}}%
\TextAt{13}{514.491}{156.684}{\Character{62}}%
\TextAt{9}{526.917}{156.684}{0}%
\TextAt{8}{79.226}{139.638}{for}%
\TextAt{13}{97.111}{139.638}{x}%
\NativeMathAt{107.084}{139.638}{11.955}{\neq}%
\TextAt{9}{119.494}{139.638}{0}%
\TextAt{8}{125.362}{139.638}{.}%
\TextAt{8}{97.159}{119.603}{We}%
\TextAt{8}{117.303}{119.603}{call}%
\TextAt{8}{138.763}{119.603}{the}%
\TextAt{8}{158.919}{119.603}{euclidean}%
\TextAt{8}{210.291}{119.603}{vector}%
\TextAt{8}{245.415}{119.603}{space}%
\TextAt{13}{266.55}{81.888}{K}%
\TextAt{25}{276.509}{80.095}{R}%
\TextAt{9}{286.085}{81.888}{\Character{61}}%
\TextAt{20}{298.51}{102.331}{\Character{34}}%
\TextAt{20}{305.484}{93.245}{\Character{89}}%
\TextAt{21}{310.783}{68.139}{\Character{28}}%
\TextAt{18}{322.753}{81.888}{C}%
\TextAt{20}{331.387}{102.331}{\Character{35}}%
\TextAt{17}{338.361}{99.655}{\Character{43}}%
\TextAt{8}{300.14}{43.033}{27}%
\Destination{78.23}{809.815}{nav.page.27}%
\Destination{79.36}{643.655}{nav.print.29}%
\Hotlink{258.233}{453.255}{25.078}{11.047}{nav.number.1.2.6}%
}

% Page 40 (body page 28)
\Page{%
\TextAt{2}{79.226}{796.943}{I.}%
\TextAt{2}{89.83}{796.943}{Algebraic}%
\TextAt{2}{138.648}{796.943}{Integers}%
\Vector{q 0 G 0 g 0.299 w 7.226 18.729 m 460.769 18.729 l S Q}%
\TextAt{8}{79.226}{761.124}{the}%
\TextAt{12}{98.617}{761.124}{Minkowski}%
\TextAt{12}{165.151}{761.124}{space}%
\TextAt{8}{196.582}{761.124}{,}%
\TextAt{8}{203.118}{761.124}{its}%
\TextAt{8}{218.671}{761.124}{scalar}%
\TextAt{8}{251.141}{761.124}{product}%
\TextAt{9}{294.275}{761.124}{\Character{40}}%
\TextAt{13}{301.959}{761.124}{\Character{59}}%
\TextAt{9}{310.335}{761.124}{\Character{41}}%
\TextAt{8}{318.019}{761.124}{the}%
\TextAt{12}{337.41}{761.124}{canonical}%
\TextAt{12}{395.205}{761.124}{metric}%
\TextAt{8}{433.007}{761.124}{,}%
\TextAt{8}{439.544}{761.124}{and}%
\TextAt{8}{461.541}{761.124}{the}%
\TextAt{8}{480.932}{761.124}{associated}%
\TextAt{8}{79.226}{744.078}{Haar}%
\TextAt{8}{108.216}{744.078}{measure}%
\TextAt{8}{153.847}{744.078}{(see}%
\TextAt{8}{177.362}{744.078}{§4,}%
\TextAt{8}{196.106}{744.078}{p.26)}%
\TextAt{8}{226.076}{744.078}{the}%
\TextAt{12}{246.279}{744.078}{canonical}%
\TextAt{12}{305.008}{744.078}{measure}%
\TextAt{8}{353.307}{744.078}{.}%
\TextAt{8}{361.895}{744.078}{Since}%
\TextAt{8}{392.499}{744.078}{Tr}%
\TextAt{16}{407.498}{744.078}{\Character{14}}%
\TextAt{13}{413.476}{744.078}{F}%
\TextAt{9}{426.079}{744.078}{\Character{61}}%
\TextAt{13}{438.584}{744.078}{F}%
\TextAt{16}{450.476}{744.078}{\Character{14}}%
\TextAt{8}{459.142}{744.078}{Tr}%
\TextAt{8}{476.093}{744.078}{we}%
\TextAt{8}{493.367}{744.078}{have}%
\TextAt{8}{520.408}{744.078}{on}%
\TextAt{13}{79.226}{727.032}{K}%
\TextAt{25}{89.185}{725.238}{R}%
\TextAt{8}{99.337}{727.032}{the}%
\TextAt{18}{119.493}{727.032}{R}%
\TextAt{8}{128.128}{727.032}{-linear}%
\TextAt{8}{164.543}{727.032}{map}%
\TextAt{8}{267.546}{709.986}{Tr}%
\TextAt{9}{283.874}{709.986}{\Character{58}}%
\TextAt{13}{290.447}{709.986}{K}%
\TextAt{25}{300.406}{708.193}{R}%
\NativeArrowAt{309.981}{709.986}{11.955}{19.26}{\longrightarrow}%
\TextAt{18}{332.563}{709.986}{R}%
\TextAt{13}{341.198}{709.986}{\Character{59}}%
\TextAt{8}{79.226}{688.163}{and}%
\TextAt{8}{101.989}{688.163}{its}%
\TextAt{8}{118.308}{688.163}{composite}%
\TextAt{8}{173.325}{688.163}{with}%
\TextAt{13}{199.985}{688.163}{j}%
\TextAt{9}{208.822}{688.163}{\Character{58}}%
\TextAt{13}{215.394}{688.163}{K}%
\NativeMathAt{229.533}{688.163}{11.955}{\rightarrow}%
\TextAt{13}{244.809}{688.163}{K}%
\TextAt{25}{254.768}{686.369}{R}%
\TextAt{8}{264.92}{688.163}{is}%
\TextAt{8}{276.683}{688.163}{again}%
\TextAt{8}{307.91}{688.163}{the}%
\TextAt{8}{328.067}{688.163}{usual}%
\TextAt{8}{358.696}{688.163}{trace}%
\TextAt{8}{387.962}{688.163}{of}%
\TextAt{13}{401.292}{688.163}{K}%
\TextAt{16}{415.431}{688.163}{\Character{106}}%
\TextAt{18}{422.072}{688.163}{Q}%
\TextAt{8}{431.371}{688.163}{,}%
\TextAt{8}{257.107}{663.125}{Tr}%
\TextAt{21}{270.114}{661.265}{K}%
\TextAt{3}{277.819}{661.265}{\Character{106}}%
\TextAt{25}{280.171}{661.265}{Q}%
\TextAt{9}{286.868}{663.125}{\Character{40}}%
\TextAt{13}{291.421}{663.125}{a}%
\TextAt{9}{297.566}{663.125}{\Character{41}}%
\TextAt{9}{305.421}{663.125}{\Character{61}}%
\TextAt{8}{317.864}{663.125}{Tr}%
\TextAt{9}{330.872}{663.125}{\Character{40}}%
\TextAt{13}{335.424}{663.125}{j}%
\TextAt{13}{340.935}{663.125}{a}%
\TextAt{9}{347.084}{663.125}{\Character{41}}%
\TextAt{13}{351.637}{663.125}{\Character{58}}%
\TextAt{12}{79.226}{632.675}{Remark.}%
\TextAt{8}{133.997}{632.675}{We}%
\TextAt{8}{154.756}{632.675}{mention}%
\TextAt{8}{200.573}{632.675}{in}%
\TextAt{8}{214.84}{632.675}{passing—it}%
\TextAt{8}{276.067}{632.675}{will}%
\TextAt{8}{298.787}{632.675}{not}%
\TextAt{8}{320.215}{632.675}{be}%
\TextAt{8}{336.753}{632.675}{used}%
\TextAt{8}{364.088}{632.675}{in}%
\TextAt{8}{378.355}{632.675}{the}%
\TextAt{8}{399.125}{632.675}{sequel—that}%
\TextAt{8}{467.765}{632.675}{the}%
\TextAt{8}{488.535}{632.675}{mapping}%
\TextAt{13}{79.226}{615.629}{j}%
\TextAt{9}{88.062}{615.629}{\Character{58}}%
\TextAt{13}{94.635}{615.629}{K}%
\NativeMathAt{108.773}{615.629}{11.955}{\rightarrow}%
\TextAt{13}{124.049}{615.629}{K}%
\TextAt{25}{134.009}{613.836}{R}%
\TextAt{8}{144.16}{615.629}{identifies}%
\TextAt{8}{193.32}{615.629}{the}%
\TextAt{8}{213.476}{615.629}{vector}%
\TextAt{8}{248.601}{615.629}{space}%
\TextAt{13}{279.875}{615.629}{K}%
\TextAt{25}{289.835}{613.836}{R}%
\TextAt{8}{299.986}{615.629}{with}%
\TextAt{8}{326.646}{615.629}{the}%
\TextAt{8}{346.803}{615.629}{tensor}%
\TextAt{8}{381.987}{615.629}{product}%
\TextAt{13}{425.886}{615.629}{K}%
\TextAt{16}{439.361}{615.629}{\Character{10}}%
\TextAt{25}{448.659}{613.836}{Q}%
\TextAt{18}{458.013}{615.629}{R}%
\TextAt{8}{466.647}{615.629}{,}%
\TextAt{13}{211.036}{590.591}{K}%
\TextAt{16}{224.51}{590.591}{\Character{10}}%
\TextAt{25}{233.809}{588.798}{Q}%
\TextAt{18}{243.162}{590.591}{R}%
\NativeMathAt{257.705}{597.368}{7.97}{\sim}%
\NativeArrowAt{255.118}{590.591}{11.955}{13.175}{\longrightarrow}%
\TextAt{13}{271.611}{590.591}{K}%
\TextAt{25}{281.57}{588.798}{R}%
\TextAt{13}{287.825}{590.591}{\Character{59}}%
\TextAt{13}{316.96}{590.591}{a}%
\TextAt{16}{325.781}{590.591}{\Character{10}}%
\TextAt{13}{337.736}{590.591}{x}%
\NativeArrowAt{347.709}{590.591}{11.955}{19.26}{\longmapsto}%
\TextAt{9}{370.291}{590.591}{\Character{40}}%
\TextAt{13}{374.843}{590.591}{j}%
\TextAt{13}{380.354}{590.591}{a}%
\TextAt{9}{386.504}{590.591}{\Character{41}}%
\TextAt{13}{391.056}{590.591}{x\Character{58}}%
\TextAt{8}{79.226}{565.553}{Likewise,}%
\TextAt{13}{130.414}{565.553}{K}%
\TextAt{16}{144.375}{565.553}{\Character{10}}%
\TextAt{25}{153.673}{563.76}{Q}%
\TextAt{18}{163.513}{565.553}{C}%
\NativeMathAt{176.685}{565.553}{11.955}{\simeq}%
\TextAt{13}{190.521}{565.553}{K}%
\TextAt{25}{200.48}{563.76}{C}%
\TextAt{8}{206.734}{565.553}{.}%
\TextAt{8}{217.324}{565.553}{In}%
\TextAt{8}{232.659}{565.553}{this}%
\TextAt{8}{256.195}{565.553}{approach,}%
\TextAt{8}{310.754}{565.553}{the}%
\TextAt{8}{331.624}{565.553}{inclusion}%
\TextAt{13}{381.174}{565.553}{K}%
\TextAt{25}{391.134}{563.76}{R}%
\NativeMathAt{401.925}{565.553}{11.955}{\subseteq}%
\TextAt{13}{415.761}{565.553}{K}%
\TextAt{25}{425.72}{563.76}{C}%
\TextAt{8}{436.585}{565.553}{corresponds}%
\TextAt{8}{501.486}{565.553}{to}%
\TextAt{8}{516.51}{565.553}{the}%
\TextAt{8}{79.226}{548.507}{canonical}%
\TextAt{8}{131.651}{548.507}{mapping}%
\TextAt{13}{180.823}{548.507}{K}%
\TextAt{16}{195.007}{548.507}{\Character{10}}%
\TextAt{25}{204.306}{546.714}{Q}%
\TextAt{18}{214.369}{548.507}{R}%
\NativeMathAt{228.098}{548.507}{11.955}{\rightarrow}%
\TextAt{13}{245.148}{548.507}{K}%
\TextAt{16}{259.332}{548.507}{\Character{10}}%
\TextAt{25}{268.631}{546.714}{Q}%
\TextAt{18}{278.694}{548.507}{C}%
\TextAt{8}{292.267}{548.507}{which}%
\TextAt{8}{326.795}{548.507}{is}%
\TextAt{8}{339.599}{548.507}{induced}%
\TextAt{8}{384.205}{548.507}{by}%
\TextAt{8}{401.505}{548.507}{the}%
\TextAt{8}{422.703}{548.507}{inclusion}%
\TextAt{18}{472.581}{548.507}{R}%
\NativeMathAt{486.31}{548.507}{11.955}{\rightarrow}%
\TextAt{18}{503.36}{548.507}{C}%
\TextAt{8}{511.994}{548.507}{.}%
\TextAt{13}{523.566}{548.507}{F}%
\TextAt{8}{79.226}{531.461}{corresponds}%
\TextAt{8}{143.413}{531.461}{to}%
\TextAt{13}{157.724}{531.461}{F}%
\TextAt{9}{166.927}{531.461}{\Character{40}}%
\TextAt{13}{171.48}{531.461}{a}%
\TextAt{16}{180.281}{531.461}{\Character{10}}%
\TextAt{13}{192.237}{531.461}{z}%
\TextAt{9}{198.207}{531.461}{\Character{41}}%
\TextAt{9}{206.062}{531.461}{\Character{61}}%
\TextAt{13}{218.506}{531.461}{a}%
\TextAt{16}{227.307}{531.461}{\Character{10}}%
\TextAt{9}{239.972}{531.461}{\Character{22}}%
\TextAt{13}{239.263}{531.461}{z}%
\TextAt{8}{245.233}{531.461}{.}%
\TextAt{8}{97.159}{509.004}{An}%
\TextAt{8}{118.243}{509.004}{explicit}%
\TextAt{8}{161.444}{509.004}{description}%
\TextAt{8}{223.248}{509.004}{of}%
\TextAt{8}{238.486}{509.004}{the}%
\TextAt{8}{260.55}{509.004}{Minkowski}%
\TextAt{8}{320.728}{509.004}{space}%
\TextAt{13}{353.911}{509.004}{K}%
\TextAt{25}{363.87}{507.211}{R}%
\TextAt{8}{375.93}{509.004}{can}%
\TextAt{8}{399.297}{509.004}{be}%
\TextAt{8}{417.129}{509.004}{given}%
\TextAt{8}{449.607}{509.004}{in}%
\TextAt{8}{465.168}{509.004}{the}%
\TextAt{8}{487.232}{509.004}{following}%
\TextAt{8}{79.226}{491.958}{manner.}%
\TextAt{8}{125.901}{491.958}{Some}%
\TextAt{8}{156.668}{491.958}{of}%
\TextAt{8}{169.55}{491.958}{the}%
\TextAt{8}{189.258}{491.958}{embeddings}%
\TextAt{13}{252.602}{491.958}{\Character{28}}%
\TextAt{9}{262.342}{491.958}{\Character{58}}%
\TextAt{13}{268.914}{491.958}{K}%
\NativeMathAt{283.053}{491.958}{11.955}{\rightarrow}%
\TextAt{18}{298.329}{491.958}{C}%
\TextAt{8}{310.412}{491.958}{are}%
\TextAt{12}{329.475}{491.958}{real}%
\TextAt{8}{354.838}{491.958}{in}%
\TextAt{8}{368.042}{491.958}{that}%
\TextAt{8}{392.963}{491.958}{they}%
\TextAt{8}{418.852}{491.958}{land}%
\TextAt{8}{444.418}{491.958}{already}%
\TextAt{8}{485.275}{491.958}{in}%
\TextAt{18}{498.479}{491.958}{R}%
\TextAt{8}{507.114}{491.958}{,}%
\TextAt{8}{513.904}{491.958}{and}%
\TextAt{8}{79.226}{474.912}{others}%
\TextAt{8}{114.41}{474.912}{are}%
\TextAt{12}{133.921}{474.912}{complex}%
\TextAt{8}{182.292}{474.912}{,}%
\TextAt{8}{189.441}{474.912}{i.e.,}%
\TextAt{8}{211.546}{474.912}{not}%
\TextAt{8}{232.36}{474.912}{real.}%
\TextAt{8}{259.673}{474.912}{Let}%
\TextAt{13}{254.582}{449.874}{\Character{26}}%
\TextAt{17}{260.62}{448.081}{1}%
\TextAt{13}{265.352}{449.874}{\Character{59}}%
\TextAt{13}{270.576}{449.874}{\Character{58}}%
\TextAt{13}{275.813}{449.874}{\Character{58}}%
\TextAt{13}{281.049}{449.874}{\Character{58}}%
\TextAt{13}{286.274}{449.874}{\Character{59}}%
\TextAt{13}{291.51}{449.874}{\Character{26}}%
\TextAt{21}{297.611}{448.081}{r}%
\TextAt{9}{305.486}{449.874}{\Character{58}}%
\TextAt{13}{312.059}{449.874}{K}%
\NativeArrowAt{326.197}{449.874}{11.955}{19.26}{\longrightarrow}%
\TextAt{18}{348.779}{449.874}{R}%
\TextAt{8}{79.226}{424.836}{be}%
\TextAt{8}{95.15}{424.836}{the}%
\TextAt{8}{115.307}{424.836}{real}%
\TextAt{8}{138.07}{424.836}{embeddings.}%
\TextAt{8}{206.413}{424.836}{The}%
\TextAt{8}{230.467}{424.836}{complex}%
\TextAt{8}{276.315}{424.836}{ones}%
\TextAt{8}{302.389}{424.836}{come}%
\TextAt{8}{332.301}{424.836}{in}%
\TextAt{8}{345.954}{424.836}{pairs}%
\TextAt{13}{237.566}{399.798}{\Character{27}}%
\TextAt{17}{244.219}{398.005}{1}%
\TextAt{13}{248.952}{399.798}{\Character{59}}%
\TextAt{9}{254.811}{399.798}{\Character{22}}%
\TextAt{13}{254.196}{399.798}{\Character{27}}%
\TextAt{17}{260.849}{398.005}{1}%
\TextAt{13}{265.582}{399.798}{\Character{59}}%
\TextAt{13}{270.806}{399.798}{\Character{58}}%
\TextAt{13}{276.043}{399.798}{\Character{58}}%
\TextAt{13}{281.267}{399.798}{\Character{58}}%
\TextAt{13}{286.504}{399.798}{\Character{59}}%
\TextAt{13}{291.74}{399.798}{\Character{27}}%
\TextAt{21}{298.456}{398.005}{s}%
\TextAt{13}{302.87}{399.798}{\Character{59}}%
\TextAt{9}{308.729}{399.798}{\Character{22}}%
\TextAt{13}{308.114}{399.798}{\Character{27}}%
\TextAt{21}{314.768}{398.005}{s}%
\TextAt{9}{322.502}{399.798}{\Character{58}}%
\TextAt{13}{329.075}{399.798}{K}%
\NativeArrowAt{343.213}{399.798}{11.955}{19.26}{\longrightarrow}%
\TextAt{18}{365.795}{399.798}{C}%
\TextAt{8}{79.226}{374.76}{of}%
\TextAt{8}{93.57}{374.76}{complex}%
\TextAt{8}{140.432}{374.76}{conjugate}%
\TextAt{8}{194.455}{374.76}{embeddings.}%
\TextAt{8}{265.841}{374.76}{Thus}%
\TextAt{13}{296.503}{374.76}{n}%
\TextAt{9}{308.54}{374.76}{\Character{61}}%
\TextAt{13}{322.693}{374.76}{r}%
\TextAt{9}{331.641}{374.76}{\Character{43}}%
\TextAt{9}{344.086}{374.76}{2}%
\TextAt{13}{349.947}{374.76}{s}%
\TextAt{8}{355.461}{374.76}{.}%
\TextAt{8}{366.951}{374.76}{We}%
\TextAt{8}{388.109}{374.76}{choose}%
\TextAt{8}{426.256}{374.76}{from}%
\TextAt{8}{454.91}{374.76}{each}%
\TextAt{8}{482.261}{374.76}{pair}%
\TextAt{8}{507.341}{374.76}{some}%
\TextAt{8}{79.226}{357.714}{fixed}%
\TextAt{8}{107.879}{357.714}{complex}%
\TextAt{8}{154.095}{357.714}{embedding,}%
\TextAt{8}{216.984}{357.714}{and}%
\TextAt{8}{240.114}{357.714}{let}%
\TextAt{13}{257.386}{357.714}{\Character{26}}%
\TextAt{8}{267.688}{357.714}{vary}%
\TextAt{8}{294.082}{357.714}{over}%
\TextAt{8}{319.495}{357.714}{the}%
\TextAt{8}{340.019}{357.714}{family}%
\TextAt{8}{376.156}{357.714}{of}%
\TextAt{8}{389.854}{357.714}{real}%
\TextAt{8}{412.984}{357.714}{embeddings}%
\TextAt{8}{477.144}{357.714}{and}%
\TextAt{13}{500.274}{357.714}{\Character{27}}%
\TextAt{8}{511.621}{357.714}{over}%
\TextAt{8}{79.226}{340.668}{the}%
\TextAt{8}{99.428}{340.668}{family}%
\TextAt{8}{135.244}{340.668}{of}%
\TextAt{8}{148.62}{340.668}{chosen}%
\TextAt{8}{186.121}{340.668}{complex}%
\TextAt{8}{232.015}{340.668}{embeddings.}%
\TextAt{8}{300.496}{340.668}{Since}%
\TextAt{13}{331.1}{340.668}{F}%
\TextAt{8}{344.247}{340.668}{leaves}%
\TextAt{8}{377.851}{340.668}{the}%
\TextAt{13}{398.053}{340.668}{\Character{26}}%
\TextAt{8}{408.034}{340.668}{invariant,}%
\TextAt{8}{460.467}{340.668}{but}%
\TextAt{8}{481.972}{340.668}{exchanges}%
\TextAt{8}{79.226}{323.622}{the}%
\TextAt{13}{99.383}{323.622}{\Character{27}}%
\TextAt{13}{105.815}{323.622}{\Character{59}}%
\TextAt{9}{111.674}{323.622}{\Character{22}}%
\TextAt{13}{111.059}{323.622}{\Character{27}}%
\TextAt{8}{118.141}{323.622}{,}%
\TextAt{8}{125.29}{323.622}{we}%
\TextAt{8}{142.518}{323.622}{have}%
\TextAt{13}{197.465}{299.38}{K}%
\TextAt{25}{207.424}{297.586}{R}%
\TextAt{9}{216.999}{299.38}{\Character{61}}%
\TextAt{20}{229.425}{319.823}{\Character{40}}%
\TextAt{9}{239.055}{299.38}{\Character{40}}%
\TextAt{13}{243.608}{299.38}{z}%
\TextAt{21}{249.046}{297.586}{\Character{28}}%
\TextAt{9}{254.222}{299.38}{\Character{41}}%
\NativeMathAt{262.095}{299.38}{11.955}{\in}%
\TextAt{20}{273.386}{310.737}{\Character{89}}%
\TextAt{21}{278.685}{285.631}{\Character{28}}%
\TextAt{18}{290.655}{299.38}{C}%
\TextAt{20}{303.186}{320.301}{\Character{12}}%
\TextAt{20}{303.186}{313.128}{\Character{12}}%
\TextAt{20}{303.186}{305.955}{\Character{12}}%
\TextAt{20}{303.186}{298.782}{\Character{12}}%
\TextAt{20}{303.186}{291.609}{\Character{12}}%
\TextAt{13}{311.069}{299.38}{z}%
\TextAt{21}{316.507}{297.586}{\Character{26}}%
\NativeMathAt{324.662}{299.38}{11.955}{\in}%
\TextAt{18}{335.953}{299.38}{R}%
\TextAt{13}{344.588}{299.38}{\Character{59}}%
\TextAt{13}{361.768}{299.38}{z}%
\TextAt{21}{367.225}{297.586}{\Character{27}}%
\TextAt{9}{376.169}{299.38}{\Character{61}}%
\TextAt{9}{389.296}{299.38}{\Character{22}}%
\TextAt{13}{388.595}{299.38}{z}%
\TextAt{17}{394.478}{297.586}{\Character{22}}%
\TextAt{21}{394.032}{297.586}{\Character{27}}%
\TextAt{20}{399.656}{319.823}{\Character{41}}%
\TextAt{13}{411.279}{299.38}{\Character{58}}%
\TextAt{8}{79.226}{270.365}{This}%
\TextAt{8}{105.946}{270.365}{gives}%
\TextAt{8}{134.626}{270.365}{the}%
\TextAt{12}{79.226}{245.485}{Proposition}%
\TextAt{12}{151.937}{245.485}{(5.1).}%
\TextAt{8}{187.253}{245.485}{There}%
\TextAt{8}{221.062}{245.485}{is}%
\TextAt{8}{232.826}{245.485}{an}%
\TextAt{8}{249.085}{245.485}{isomorphism}%
\TextAt{13}{244.726}{220.447}{f}%
\TextAt{9}{255.094}{220.447}{\Character{58}}%
\TextAt{13}{261.666}{220.447}{K}%
\TextAt{25}{271.625}{218.654}{R}%
\NativeMathAt{283.788}{227.224}{7.97}{\sim}%
\NativeArrowAt{281.2}{220.447}{11.955}{13.175}{\longrightarrow}%
\TextAt{20}{297.694}{231.805}{\Character{89}}%
\TextAt{21}{302.993}{206.699}{\Character{28}}%
\TextAt{18}{314.963}{220.447}{R}%
\TextAt{9}{326.918}{220.447}{\Character{61}}%
\TextAt{18}{339.343}{220.447}{R}%
\TextAt{21}{347.978}{225.383}{r}%
\TextAt{17}{352.034}{225.383}{\Character{43}2}%
\TextAt{21}{362.855}{225.383}{s}%
\TextAt{8}{79.226}{187.549}{given}%
\TextAt{8}{109.796}{187.549}{by}%
\TextAt{8}{126.055}{187.549}{the}%
\TextAt{8}{146.211}{187.549}{rule}%
\TextAt{9}{169.619}{187.549}{\Character{40}}%
\TextAt{13}{174.172}{187.549}{z}%
\TextAt{21}{179.609}{185.755}{\Character{28}}%
\TextAt{9}{184.786}{187.549}{\Character{41}}%
\NativeArrowAt{192.659}{187.549}{11.955}{11.955}{\mapsto}%
\TextAt{9}{207.935}{187.549}{\Character{40}}%
\TextAt{13}{212.487}{187.549}{x}%
\TextAt{21}{219.139}{185.755}{\Character{28}}%
\TextAt{9}{224.316}{187.549}{\Character{41}}%
\TextAt{8}{232.766}{187.549}{where}%
\TextAt{13}{199.643}{162.511}{x}%
\TextAt{21}{206.295}{160.717}{\Character{26}}%
\TextAt{9}{214.45}{162.511}{\Character{61}}%
\TextAt{13}{226.876}{162.511}{z}%
\TextAt{21}{232.314}{160.717}{\Character{26}}%
\TextAt{13}{237.148}{162.511}{\Character{59}}%
\TextAt{13}{266.283}{162.511}{x}%
\TextAt{21}{272.955}{160.717}{\Character{27}}%
\TextAt{9}{281.899}{162.511}{\Character{61}}%
\TextAt{16}{294.325}{162.511}{\Character{60}}%
\TextAt{9}{302.959}{162.511}{\Character{40}}%
\TextAt{13}{307.511}{162.511}{z}%
\TextAt{21}{312.949}{160.717}{\Character{27}}%
\TextAt{9}{318.573}{162.511}{\Character{41}}%
\TextAt{13}{323.125}{162.511}{\Character{59}}%
\TextAt{13}{352.26}{162.511}{x}%
\TextAt{17}{359.377}{160.717}{\Character{22}}%
\TextAt{21}{358.931}{160.717}{\Character{27}}%
\TextAt{9}{367.876}{162.511}{\Character{61}}%
\TextAt{16}{380.301}{162.511}{\Character{61}}%
\TextAt{9}{388.935}{162.511}{\Character{40}}%
\TextAt{13}{393.488}{162.511}{z}%
\TextAt{21}{398.926}{160.717}{\Character{27}}%
\TextAt{9}{404.549}{162.511}{\Character{41}}%
\TextAt{13}{409.101}{162.511}{\Character{58}}%
\TextAt{8}{79.226}{137.473}{This}%
\TextAt{8}{105.946}{137.473}{isomorphism}%
\TextAt{8}{174.365}{137.473}{transforms}%
\TextAt{8}{232.707}{137.473}{the}%
\TextAt{8}{252.863}{137.473}{canonical}%
\TextAt{8}{304.246}{137.473}{metric}%
\TextAt{9}{340.662}{137.473}{\Character{40}}%
\TextAt{13}{349.112}{137.473}{\Character{59}}%
\TextAt{9}{358.253}{137.473}{\Character{41}}%
\TextAt{8}{366.703}{137.473}{into}%
\TextAt{8}{390.446}{137.473}{the}%
\TextAt{8}{410.602}{137.473}{scalar}%
\TextAt{8}{443.838}{137.473}{product}%
\TextAt{9}{256.282}{112.435}{\Character{40}}%
\TextAt{13}{260.835}{112.435}{x\Character{59}}%
\TextAt{13}{272.707}{112.435}{y}%
\TextAt{9}{278.868}{112.435}{\Character{41}}%
\TextAt{9}{286.723}{112.435}{\Character{61}}%
\TextAt{20}{299.166}{123.792}{\Character{88}}%
\TextAt{21}{305.461}{98.686}{\Character{28}}%
\TextAt{13}{318.427}{112.435}{a}%
\TextAt{21}{324.572}{110.641}{\Character{28}}%
\TextAt{13}{329.749}{112.435}{x}%
\TextAt{21}{336.401}{110.641}{\Character{28}}%
\TextAt{13}{341.577}{112.435}{y}%
\TextAt{21}{347.285}{110.641}{\Character{28}}%
\TextAt{13}{352.462}{112.435}{\Character{59}}%
\TextAt{8}{79.226}{80.086}{where}%
\TextAt{13}{113.035}{80.086}{a}%
\TextAt{21}{119.18}{78.293}{\Character{28}}%
\TextAt{9}{127.678}{80.086}{\Character{61}}%
\TextAt{9}{140.088}{80.086}{1}%
\TextAt{8}{145.956}{80.086}{,}%
\TextAt{8}{153.105}{80.086}{resp.}%
\TextAt{13}{182.427}{80.086}{a}%
\TextAt{21}{188.572}{78.293}{\Character{28}}%
\TextAt{9}{197.07}{80.086}{\Character{61}}%
\TextAt{9}{209.48}{80.086}{2}%
\TextAt{8}{215.348}{80.086}{,}%
\TextAt{8}{222.497}{80.086}{if}%
\TextAt{13}{233.221}{80.086}{\Character{28}}%
\TextAt{8}{243.537}{80.086}{is}%
\TextAt{8}{255.301}{80.086}{real,}%
\TextAt{8}{281.316}{80.086}{resp.}%
\TextAt{8}{310.638}{80.086}{complex.}%
\TextAt{8}{300.14}{43.033}{28}%
\Destination{78.23}{809.815}{nav.page.28}%
\Destination{79.36}{703.155}{nav.print.30}%
\Destination{79.36}{259.455}{nav.number.1.5.1}%
\Hotlink{182.541}{741.855}{6.858}{11.047}{nav.section.1.4}%
\Hotlink{205.361}{741.855}{12.716}{11.047}{nav.print.26}%
}

% Page 41 (body page 29)
\Page{%
\TextAt{2}{79.226}{796.943}{I.}%
\TextAt{2}{89.83}{796.943}{Algebraic}%
\TextAt{2}{138.648}{796.943}{Integers}%
\Vector{q 0 G 0 g 0.299 w 7.226 18.729 m 460.769 18.729 l S Q}%
\ErrataAt{79.19}{761.124}{455.81}{\textbf{Proof.} The map is clearly an isomorphism. Write $z_\sigma=x_\sigma+iy_\sigma$ and $z_\sigma^{\prime}=x_\sigma^{\prime}+iy_\sigma^{\prime}$}%
\ErrataAt{79.19}{744.078}{455.81}{for the chosen complex embeddings $\sigma$. At real $\rho$, $z_\rho\bar z_\rho^{\prime}=x_\rho x_\rho^{\prime}$. Using $y_\sigma=x_{\bar\sigma}$ and $y_\sigma^{\prime}=x_{\bar\sigma}^{\prime}$ gives}%
\TextAt{13}{159.993}{715.077}{z}%
\TextAt{21}{165.431}{713.283}{\Character{27}}%
\TextAt{9}{171.764}{715.077}{\Character{22}}%
\TextAt{13}{171.054}{715.077}{z}%
\TextAt{3}{177.025}{720.013}{\Character{48}}%
\TextAt{21}{176.492}{712.121}{\Character{27}}%
\TextAt{9}{184.772}{715.077}{\Character{43}}%
\TextAt{13}{196.534}{715.077}{z}%
\TextAt{17}{202.417}{713.283}{\Character{22}}%
\TextAt{21}{201.972}{713.283}{\Character{27}}%
\TextAt{9}{208.304}{715.077}{\Character{22}}%
\TextAt{13}{207.595}{715.077}{z}%
\TextAt{3}{213.566}{720.013}{\Character{48}}%
\TextAt{17}{213.479}{712.121}{\Character{22}}%
\TextAt{21}{213.033}{712.121}{\Character{27}}%
\TextAt{9}{221.977}{715.077}{\Character{61}}%
\TextAt{13}{234.403}{715.077}{z}%
\TextAt{21}{239.841}{713.283}{\Character{27}}%
\TextAt{9}{246.173}{715.077}{\Character{22}}%
\TextAt{13}{245.464}{715.077}{z}%
\TextAt{3}{251.435}{720.013}{\Character{48}}%
\TextAt{21}{250.902}{712.121}{\Character{27}}%
\TextAt{9}{259.182}{715.077}{\Character{43}}%
\TextAt{9}{271.639}{715.077}{\Character{22}}%
\TextAt{13}{270.943}{715.077}{z}%
\TextAt{21}{276.381}{713.283}{\Character{27}}%
\TextAt{13}{282.005}{715.077}{z}%
\TextAt{3}{287.975}{720.013}{\Character{48}}%
\TextAt{21}{287.443}{712.121}{\Character{27}}%
\TextAt{9}{296.387}{715.077}{\Character{61}}%
\TextAt{9}{308.797}{715.077}{2}%
\TextAt{16}{314.665}{715.077}{\Character{60}}%
\TextAt{9}{323.3}{715.077}{\Character{40}}%
\TextAt{13}{327.852}{715.077}{z}%
\TextAt{21}{333.29}{713.283}{\Character{27}}%
\TextAt{9}{339.622}{715.077}{\Character{22}}%
\TextAt{13}{338.913}{715.077}{z}%
\TextAt{3}{344.884}{720.013}{\Character{48}}%
\TextAt{21}{344.351}{712.121}{\Character{27}}%
\TextAt{9}{349.975}{715.077}{\Character{41}}%
\TextAt{9}{357.83}{715.077}{\Character{61}}%
\TextAt{9}{370.251}{715.077}{2\Character{40}}%
\TextAt{13}{380.678}{715.077}{x}%
\TextAt{21}{387.331}{713.283}{\Character{27}}%
\TextAt{13}{392.954}{715.077}{x}%
\TextAt{3}{399.606}{720.013}{\Character{48}}%
\TextAt{21}{399.606}{712.121}{\Character{27}}%
\TextAt{9}{407.886}{715.077}{\Character{43}}%
\TextAt{13}{419.647}{715.077}{x}%
\TextAt{17}{426.745}{713.283}{\Character{22}}%
\TextAt{21}{426.3}{713.283}{\Character{27}}%
\TextAt{13}{431.923}{715.077}{x}%
\TextAt{3}{438.575}{720.013}{\Character{48}}%
\TextAt{17}{439.021}{712.121}{\Character{22}}%
\TextAt{21}{438.575}{712.121}{\Character{27}}%
\TextAt{9}{444.199}{715.077}{\Character{41}}%
\TextAt{13}{448.751}{715.077}{\Character{58}}%
\TextAt{8}{79.226}{686.075}{This}%
\TextAt{8}{105.946}{686.075}{proves}%
\TextAt{8}{142.11}{686.075}{the}%
\TextAt{8}{162.267}{686.075}{claim}%
\TextAt{8}{193.482}{686.075}{concerning}%
\TextAt{8}{252.014}{686.075}{the}%
\TextAt{8}{272.171}{686.075}{scalar}%
\TextAt{8}{305.406}{686.075}{products.}%
\Vector{q 0 G 0 g 0.399 w 452.334 -87.925 m 452.334 -79.855 l S Q}%
\Vector{q 0 G 0 g 0.399 w 452.534 -80.054 m 459.707 -80.054 l S Q}%
\Vector{q 0 G 0 g 0.399 w 452.534 -87.725 m 459.707 -87.725 l S Q}%
\Vector{q 0 G 0 g 0.399 w 459.906 -87.925 m 459.906 -79.855 l S Q}%
\TextAt{8}{97.159}{657.074}{The}%
\TextAt{8}{122.003}{657.074}{scalar}%
\TextAt{8}{156.029}{657.074}{product}%
\TextAt{9}{200.719}{657.074}{\Character{40}}%
\TextAt{13}{205.272}{657.074}{x\Character{59}}%
\TextAt{13}{217.144}{657.074}{y}%
\TextAt{9}{223.305}{657.074}{\Character{41}}%
\TextAt{9}{232.511}{657.074}{\Character{61}}%
\TextAt{20}{246.298}{666.041}{\Character{80}}%
\TextAt{21}{258.917}{653.587}{\Character{28}}%
\TextAt{13}{266.086}{657.074}{a}%
\TextAt{21}{272.231}{655.281}{\Character{28}}%
\TextAt{13}{277.408}{657.074}{x}%
\TextAt{21}{284.06}{655.281}{\Character{28}}%
\TextAt{13}{289.236}{657.074}{y}%
\TextAt{21}{294.944}{655.281}{\Character{28}}%
\TextAt{8}{304.808}{657.074}{transfers}%
\TextAt{8}{353.527}{657.074}{the}%
\TextAt{8}{374.474}{657.074}{canonical}%
\TextAt{8}{426.648}{657.074}{measure}%
\TextAt{8}{473.024}{657.074}{from}%
\TextAt{13}{501.455}{657.074}{K}%
\TextAt{25}{511.414}{655.281}{R}%
\TextAt{8}{522.357}{657.074}{to}%
\TextAt{18}{79.226}{640.028}{R}%
\TextAt{21}{87.86}{644.367}{r}%
\TextAt{17}{91.917}{644.367}{\Character{43}2}%
\TextAt{21}{102.738}{644.367}{s}%
\TextAt{8}{107.152}{640.028}{.}%
\TextAt{8}{115.6}{640.028}{It}%
\TextAt{8}{128.273}{640.028}{obviously}%
\TextAt{8}{180.05}{640.028}{differs}%
\TextAt{8}{214.9}{640.028}{from}%
\TextAt{8}{242.54}{640.028}{the}%
\TextAt{8}{262.697}{640.028}{standard}%
\TextAt{8}{311.545}{640.028}{Lebesgue}%
\TextAt{8}{362.164}{640.028}{measure}%
\TextAt{8}{407.749}{640.028}{by}%
\TextAt{8}{212.804}{611.027}{vol}%
\TextAt{7}{228.095}{609.234}{canonical}%
\TextAt{9}{262.92}{611.027}{\Character{40}}%
\TextAt{13}{267.473}{611.027}{X}%
\TextAt{9}{278.128}{611.027}{\Character{41}}%
\TextAt{9}{285.983}{611.027}{\Character{61}}%
\TextAt{9}{298.404}{611.027}{2}%
\TextAt{21}{304.279}{615.963}{s}%
\TextAt{8}{310.686}{611.027}{vol}%
\TextAt{7}{325.977}{609.234}{Lebesgue}%
\TextAt{9}{360.029}{611.027}{\Character{40}}%
\TextAt{13}{364.581}{611.027}{f}%
\TextAt{9}{371.628}{611.027}{\Character{40}}%
\TextAt{13}{376.18}{611.027}{X}%
\TextAt{9}{386.835}{611.027}{\Character{41}\Character{41}}%
\TextAt{13}{395.94}{611.027}{\Character{58}}%
\TextAt{8}{79.226}{582.026}{Minkowski}%
\TextAt{8}{136.263}{582.026}{himself}%
\TextAt{8}{175.08}{582.026}{worked}%
\TextAt{8}{213.849}{582.026}{with}%
\TextAt{8}{239.276}{582.026}{the}%
\TextAt{8}{258.2}{582.026}{Lebesgue}%
\TextAt{8}{307.585}{582.026}{measure}%
\TextAt{8}{351.937}{582.026}{on}%
\TextAt{18}{366.963}{582.026}{R}%
\TextAt{21}{375.598}{586.364}{r}%
\TextAt{17}{379.654}{586.364}{\Character{43}2}%
\TextAt{21}{390.475}{586.364}{s}%
\TextAt{8}{394.889}{582.026}{,}%
\TextAt{8}{401.052}{582.026}{and}%
\TextAt{8}{422.582}{582.026}{most}%
\TextAt{8}{450.029}{582.026}{textbooks}%
\TextAt{8}{502.846}{582.026}{follow}%
\TextAt{8}{79.226}{564.98}{suit.}%
\TextAt{8}{106.599}{564.98}{The}%
\TextAt{8}{130.653}{564.98}{corresponding}%
\TextAt{8}{205.839}{564.98}{measure}%
\TextAt{8}{251.424}{564.98}{on}%
\TextAt{13}{267.683}{564.98}{K}%
\TextAt{25}{277.643}{563.187}{R}%
\TextAt{8}{287.794}{564.98}{is}%
\TextAt{8}{299.558}{564.98}{the}%
\TextAt{8}{319.715}{564.98}{one}%
\TextAt{8}{341.174}{564.98}{determined}%
\TextAt{8}{402.301}{564.98}{by}%
\TextAt{8}{418.56}{564.98}{the}%
\TextAt{8}{438.716}{564.98}{scalar}%
\TextAt{8}{471.952}{564.98}{product}%
\TextAt{9}{255.087}{533.786}{\Character{40}}%
\TextAt{13}{259.639}{533.786}{x\Character{59}}%
\TextAt{13}{271.511}{533.786}{y}%
\TextAt{9}{277.672}{533.786}{\Character{41}}%
\TextAt{9}{285.527}{533.786}{\Character{61}}%
\TextAt{20}{297.971}{545.143}{\Character{88}}%
\TextAt{21}{304.266}{520.037}{\Character{28}}%
\TextAt{9}{321.162}{541.873}{1}%
\Vector{q 0 G 0 g 0.478 w 246.427 -237.226 m 257.749 -237.226 l S Q}%
\TextAt{13}{318.427}{525.585}{a}%
\TextAt{21}{324.572}{523.792}{\Character{28}}%
\TextAt{13}{330.944}{533.786}{x}%
\TextAt{21}{337.596}{531.992}{\Character{28}}%
\TextAt{9}{343.565}{533.786}{\Character{22}}%
\TextAt{13}{342.773}{533.786}{y}%
\TextAt{21}{348.481}{531.992}{\Character{28}}%
\TextAt{13}{353.657}{533.786}{\Character{58}}%
\TextAt{8}{79.226}{497.474}{This}%
\TextAt{8}{107.224}{497.474}{scalar}%
\TextAt{8}{141.738}{497.474}{product}%
\TextAt{8}{186.916}{497.474}{may}%
\TextAt{8}{213.563}{497.474}{therefore}%
\TextAt{8}{263.942}{497.474}{be}%
\TextAt{8}{281.145}{497.474}{called}%
\TextAt{8}{315.587}{497.474}{the}%
\TextAt{12}{337.022}{497.474}{Minkowski}%
\TextAt{12}{405.907}{497.474}{metric}%
\TextAt{8}{448.885}{497.474}{on}%
\TextAt{13}{466.423}{497.474}{K}%
\TextAt{25}{476.382}{495.68}{R}%
\TextAt{8}{482.636}{497.474}{.}%
\TextAt{8}{494.92}{497.474}{But}%
\TextAt{8}{519.439}{497.474}{we}%
\TextAt{8}{79.226}{480.428}{will}%
\TextAt{8}{101.442}{480.428}{systematically}%
\TextAt{8}{177.778}{480.428}{work}%
\TextAt{8}{206.509}{480.428}{with}%
\TextAt{8}{233.279}{480.428}{the}%
\TextAt{8}{253.546}{480.428}{canonical}%
\TextAt{8}{305.04}{480.428}{metric,}%
\TextAt{8}{344.845}{480.428}{and}%
\TextAt{8}{367.718}{480.428}{denote}%
\TextAt{8}{405.546}{480.428}{by}%
\TextAt{8}{421.916}{480.428}{vol}%
\TextAt{8}{441.214}{480.428}{the}%
\TextAt{8}{461.481}{480.428}{corresponding}%
\TextAt{8}{79.226}{463.382}{canonical}%
\TextAt{8}{130.609}{463.382}{measure.}%
\TextAt{8}{97.159}{443.347}{The}%
\TextAt{8}{121.213}{443.347}{mapping}%
\TextAt{13}{169.344}{443.347}{j}%
\TextAt{9}{178.181}{443.347}{\Character{58}}%
\TextAt{13}{184.753}{443.347}{K}%
\NativeMathAt{198.892}{443.347}{11.955}{\rightarrow}%
\TextAt{13}{214.168}{443.347}{K}%
\TextAt{25}{224.127}{441.554}{R}%
\TextAt{8}{234.278}{443.347}{gives}%
\TextAt{8}{262.959}{443.347}{us}%
\TextAt{8}{277.975}{443.347}{the}%
\TextAt{8}{298.131}{443.347}{following}%
\TextAt{8}{347.566}{443.347}{lattices}%
\TextAt{8}{387.95}{443.347}{in}%
\TextAt{8}{401.603}{443.347}{Minkowski}%
\TextAt{8}{459.872}{443.347}{space}%
\TextAt{13}{491.147}{443.347}{K}%
\TextAt{25}{501.106}{441.554}{R}%
\TextAt{8}{507.361}{443.347}{.}%
\TextAt{12}{79.226}{417.335}{Proposition}%
\TextAt{12}{152.279}{417.335}{(5.2).}%
\TextAt{8}{187.891}{417.335}{If}%
\TextAt{14}{199.88}{417.335}{a}%
\NativeMathAt{209.689}{417.335}{11.955}{\neq}%
\TextAt{9}{222.613}{417.335}{0}%
\TextAt{8}{232.667}{417.335}{is}%
\TextAt{8}{244.727}{417.335}{an}%
\TextAt{8}{261.283}{417.335}{ideal}%
\TextAt{8}{289.543}{417.335}{of}%
\TextAt{14}{303.169}{417.335}{o}%
\TextAt{21}{309.013}{415.541}{K}%
\TextAt{8}{317.216}{417.335}{,}%
\TextAt{8}{324.736}{417.335}{then}%
\TextAt{9}{351.693}{417.335}{\Character{0}}%
\TextAt{9}{362.823}{417.335}{\Character{61}}%
\TextAt{13}{375.764}{417.335}{j}%
\TextAt{14}{381.28}{417.335}{a}%
\TextAt{8}{391.457}{417.335}{is}%
\TextAt{8}{403.517}{417.335}{a}%
\TextAt{8}{413.569}{417.335}{complete}%
\TextAt{8}{463.288}{417.335}{lattice}%
\TextAt{8}{499.355}{417.335}{in}%
\TextAt{13}{513.304}{417.335}{K}%
\TextAt{25}{523.263}{415.541}{R}%
\TextAt{8}{529.518}{417.335}{.}%
\TextAt{8}{79.226}{400.289}{Its}%
\TextAt{8}{96.513}{400.289}{fundamental}%
\TextAt{8}{164.156}{400.289}{mesh}%
\TextAt{8}{194.127}{400.289}{has}%
\TextAt{8}{215.001}{400.289}{volume}%
\TextAt{8}{243.693}{371.288}{vol}%
\TextAt{9}{258.983}{371.288}{\Character{40}\Character{0}\Character{41}}%
\TextAt{9}{278.685}{371.288}{\Character{61}}%
\TextAt{20}{291.148}{382.093}{\Character{112}}%
\Vector{q 0 G 0 g 0.478 w 231.104 -391.668 m 252.031 -391.668 l S Q}%
\TextAt{16}{303.104}{371.288}{\Character{106}}%
\TextAt{13}{306.425}{371.288}{d}%
\TextAt{21}{312.507}{369.494}{K}%
\TextAt{16}{320.71}{371.288}{\Character{106}}%
\TextAt{9}{326.023}{371.288}{\Character{40}}%
\TextAt{14}{330.576}{371.288}{o}%
\TextAt{21}{336.42}{369.494}{K}%
\TextAt{9}{347.943}{371.288}{\Character{58}}%
\TextAt{14}{354.516}{371.288}{a}%
\TextAt{9}{360.499}{371.288}{\Character{41}}%
\TextAt{13}{365.051}{371.288}{\Character{58}}%
\TextAt{19}{79.226}{324.354}{Proof.}%
\TextAt{8}{116}{324.354}{Let}%
\TextAt{13}{137.866}{324.354}{\Character{11}}%
\TextAt{17}{145.312}{322.56}{1}%
\TextAt{13}{150.044}{324.354}{\Character{59}}%
\TextAt{13}{155.268}{324.354}{\Character{58}}%
\TextAt{13}{160.505}{324.354}{\Character{58}}%
\TextAt{13}{165.729}{324.354}{\Character{58}}%
\TextAt{13}{170.966}{324.354}{\Character{59}}%
\TextAt{13}{176.202}{324.354}{\Character{11}}%
\TextAt{21}{183.71}{322.56}{n}%
\TextAt{8}{194.141}{324.354}{be}%
\TextAt{8}{210.961}{324.354}{a}%
\TextAt{18}{221.613}{324.354}{Z}%
\TextAt{8}{229.583}{324.354}{-basis}%
\TextAt{8}{263.117}{324.354}{of}%
\TextAt{14}{277.344}{324.354}{a}%
\TextAt{8}{283.327}{324.354}{,}%
\TextAt{8}{291.597}{324.354}{so}%
\TextAt{8}{306.863}{324.354}{that}%
\TextAt{9}{333.129}{324.354}{\Character{0}}%
\TextAt{9}{345.276}{324.354}{\Character{61}}%
\TextAt{18}{359.245}{324.354}{Z}%
\TextAt{13}{367.215}{324.354}{j}%
\TextAt{13}{372.726}{324.354}{\Character{11}}%
\TextAt{17}{380.176}{322.56}{1}%
\TextAt{9}{388.176}{324.354}{\Character{43}}%
\TextAt{16}{400.548}{324.354}{\Character{1}}%
\TextAt{16}{405.844}{324.354}{\Character{1}}%
\TextAt{16}{411.152}{324.354}{\Character{1}}%
\TextAt{9}{417.764}{324.354}{\Character{43}}%
\TextAt{18}{430.136}{324.354}{Z}%
\TextAt{13}{438.106}{324.354}{j}%
\TextAt{13}{443.617}{324.354}{\Character{11}}%
\TextAt{21}{451.067}{322.56}{n}%
\TextAt{8}{456.704}{324.354}{.}%
\TextAt{8}{467.841}{324.354}{We}%
\TextAt{8}{488.882}{324.354}{choose}%
\TextAt{8}{526.911}{324.354}{a}%
\TextAt{8}{79.226}{307.308}{numbering}%
\TextAt{8}{138.27}{307.308}{of}%
\TextAt{8}{152.433}{307.308}{the}%
\TextAt{8}{173.424}{307.308}{embeddings}%
\TextAt{13}{238.05}{307.308}{\Character{28}}%
\TextAt{9}{249.211}{307.308}{\Character{58}}%
\TextAt{13}{257.204}{307.308}{K}%
\NativeMathAt{272.763}{307.308}{11.955}{\rightarrow}%
\TextAt{18}{289.46}{307.308}{C}%
\TextAt{8}{298.094}{307.308}{,}%
\TextAt{13}{306.286}{307.308}{\Character{28}}%
\TextAt{17}{311.37}{305.514}{1}%
\TextAt{13}{316.102}{307.308}{\Character{59}}%
\TextAt{13}{321.326}{307.308}{\Character{58}}%
\TextAt{13}{326.563}{307.308}{\Character{58}}%
\TextAt{13}{331.787}{307.308}{\Character{58}}%
\TextAt{13}{337.024}{307.308}{\Character{59}}%
\TextAt{13}{342.26}{307.308}{\Character{28}}%
\TextAt{21}{347.406}{305.514}{n}%
\TextAt{8}{353.043}{307.308}{,}%
\TextAt{8}{361.234}{307.308}{and}%
\TextAt{8}{384.83}{307.308}{form}%
\TextAt{8}{413.305}{307.308}{the}%
\TextAt{8}{434.295}{307.308}{matrix}%
\TextAt{13}{473.182}{307.308}{A}%
\TextAt{9}{486.699}{307.308}{\Character{61}}%
\TextAt{9}{500.531}{307.308}{\Character{40}}%
\TextAt{13}{505.097}{307.308}{\Character{28}}%
\TextAt{21}{510.181}{305.514}{\Character{96}}%
\TextAt{13}{514.138}{307.308}{\Character{11}}%
\TextAt{21}{521.584}{305.514}{i}%
\TextAt{9}{524.965}{307.308}{\Character{41}}%
\TextAt{8}{529.518}{307.308}{.}%
\TextAt{8}{79.226}{290.262}{Then,}%
\TextAt{8}{113.035}{290.262}{according}%
\TextAt{8}{165.722}{290.262}{to}%
\TextAt{8}{180.032}{290.262}{(2.12),}%
\TextAt{8}{217.117}{290.262}{we}%
\TextAt{8}{234.344}{290.262}{have}%
\TextAt{13}{143.008}{261.261}{d}%
\TextAt{9}{149.091}{261.261}{\Character{40}}%
\TextAt{14}{153.643}{261.261}{a}%
\TextAt{9}{159.626}{261.261}{\Character{41}}%
\TextAt{9}{167.481}{261.261}{\Character{61}}%
\TextAt{13}{179.925}{261.261}{d}%
\TextAt{9}{186.008}{261.261}{\Character{40}}%
\TextAt{13}{190.56}{261.261}{\Character{11}}%
\TextAt{17}{198.006}{259.467}{1}%
\TextAt{13}{202.738}{261.261}{\Character{59}}%
\TextAt{13}{207.962}{261.261}{\Character{58}}%
\TextAt{13}{213.199}{261.261}{\Character{58}}%
\TextAt{13}{218.423}{261.261}{\Character{58}}%
\TextAt{13}{223.66}{261.261}{\Character{59}}%
\TextAt{13}{228.896}{261.261}{\Character{11}}%
\TextAt{21}{236.404}{259.467}{n}%
\TextAt{9}{242.041}{261.261}{\Character{41}}%
\TextAt{9}{249.896}{261.261}{\Character{61}}%
\TextAt{9}{262.317}{261.261}{\Character{40}}%
\TextAt{8}{266.892}{261.261}{det}%
\TextAt{13}{285.143}{261.261}{A}%
\TextAt{9}{293.918}{261.261}{\Character{41}}%
\TextAt{17}{298.471}{266.197}{2}%
\TextAt{9}{306.524}{261.261}{\Character{61}}%
\TextAt{9}{318.934}{261.261}{\Character{40}}%
\TextAt{14}{323.502}{261.261}{o}%
\TextAt{21}{329.346}{259.467}{K}%
\TextAt{9}{340.869}{261.261}{\Character{58}}%
\TextAt{14}{347.442}{261.261}{a}%
\TextAt{9}{353.425}{261.261}{\Character{41}}%
\TextAt{17}{357.977}{266.197}{2}%
\TextAt{13}{362.71}{261.261}{d}%
\TextAt{9}{368.792}{261.261}{\Character{40}}%
\TextAt{14}{373.345}{261.261}{o}%
\TextAt{21}{379.189}{259.467}{K}%
\TextAt{9}{387.392}{261.261}{\Character{41}}%
\TextAt{9}{395.247}{261.261}{\Character{61}}%
\TextAt{9}{407.668}{261.261}{\Character{40}}%
\TextAt{14}{412.242}{261.261}{o}%
\TextAt{21}{418.087}{259.467}{K}%
\TextAt{9}{429.61}{261.261}{\Character{58}}%
\TextAt{14}{436.183}{261.261}{a}%
\TextAt{9}{442.166}{261.261}{\Character{41}}%
\TextAt{17}{446.718}{266.197}{2}%
\TextAt{13}{451.45}{261.261}{d}%
\TextAt{21}{457.533}{259.467}{K}%
\TextAt{13}{465.736}{261.261}{\Character{59}}%
\TextAt{8}{79.226}{232.259}{and}%
\TextAt{8}{101.989}{232.259}{on}%
\TextAt{8}{118.248}{232.259}{the}%
\TextAt{8}{138.404}{232.259}{other}%
\TextAt{8}{168.974}{232.259}{hand}%
\TextAt{9}{208.305}{198.267}{\Character{40}\Character{40}}%
\TextAt{13}{217.41}{198.267}{j}%
\TextAt{13}{222.921}{198.267}{\Character{11}}%
\TextAt{21}{230.371}{196.474}{i}%
\TextAt{13}{233.752}{198.267}{\Character{59}}%
\TextAt{13}{238.976}{198.267}{j}%
\TextAt{13}{244.488}{198.267}{\Character{11}}%
\TextAt{21}{251.958}{196.474}{k}%
\TextAt{9}{257.077}{198.267}{\Character{41}\Character{41}}%
\TextAt{9}{269.475}{198.267}{\Character{61}}%
\TextAt{20}{281.928}{218.71}{\Character{32}}%
\TextAt{21}{297.458}{213.211}{n}%
\TextAt{20}{291.393}{209.624}{\Character{88}}%
\TextAt{21}{292.887}{184.164}{\Character{96}}%
\TextAt{17}{296.347}{184.164}{\Character{61}1}%
\TextAt{13}{310.654}{198.267}{\Character{28}}%
\TextAt{21}{315.738}{196.474}{\Character{96}}%
\TextAt{13}{319.695}{198.267}{\Character{11}}%
\TextAt{21}{327.141}{196.474}{i}%
\Vector{q 0 G 0 g 0.478 w 260.514 -568.912 m 282.121 -568.912 l S Q}%
\TextAt{13}{332.514}{198.267}{\Character{28}}%
\TextAt{21}{337.598}{196.474}{\Character{96}}%
\TextAt{13}{341.556}{198.267}{\Character{11}}%
\TextAt{21}{349.001}{196.474}{k}%
\TextAt{20}{354.121}{218.71}{\Character{33}}%
\TextAt{9}{366.906}{198.267}{\Character{61}}%
\ErrataAt{379.3}{198.267}{55.7}{$A^t\overline A$.}%
\TextAt{8}{79.226}{161.72}{This}%
\TextAt{8}{105.946}{161.72}{indeed}%
\TextAt{8}{143.007}{161.72}{yields}%
\TextAt{8}{162.13}{132.719}{vol}%
\TextAt{9}{177.421}{132.719}{\Character{40}\Character{0}\Character{41}}%
\TextAt{9}{197.123}{132.719}{\Character{61}}%
\TextAt{16}{209.586}{132.719}{\Character{106}}%
\TextAt{8}{214.9}{132.719}{det}%
\TextAt{9}{231.159}{132.719}{\Character{40}\Character{40}}%
\TextAt{13}{240.263}{132.719}{j}%
\TextAt{13}{245.774}{132.719}{\Character{11}}%
\TextAt{21}{253.224}{130.926}{i}%
\TextAt{13}{256.606}{132.719}{\Character{59}}%
\TextAt{13}{261.83}{132.719}{j}%
\TextAt{13}{267.342}{132.719}{\Character{11}}%
\TextAt{21}{274.811}{130.926}{k}%
\TextAt{9}{279.931}{132.719}{\Character{41}\Character{41}}%
\TextAt{16}{289.035}{132.719}{\Character{106}}%
\TextAt{17}{292.356}{137.655}{1\Character{47}2}%
\TextAt{9}{308.878}{132.719}{\Character{61}}%
\TextAt{16}{321.303}{132.719}{\Character{106}}%
\TextAt{8}{326.617}{132.719}{det}%
\TextAt{13}{344.868}{132.719}{A}%
\TextAt{16}{353.644}{132.719}{\Character{106}}%
\TextAt{9}{360.285}{132.719}{\Character{61}}%
\TextAt{20}{372.711}{143.524}{\Character{112}}%
\Vector{q 0 G 0 g 0.478 w 312.666 -630.236 m 333.593 -630.236 l S Q}%
\TextAt{16}{384.666}{132.719}{\Character{106}}%
\TextAt{13}{387.987}{132.719}{d}%
\TextAt{21}{394.07}{130.926}{K}%
\TextAt{16}{402.272}{132.719}{\Character{106}}%
\TextAt{9}{407.586}{132.719}{\Character{40}}%
\TextAt{14}{412.138}{132.719}{o}%
\TextAt{21}{417.982}{130.926}{K}%
\TextAt{9}{429.506}{132.719}{\Character{58}}%
\TextAt{14}{436.078}{132.719}{a}%
\TextAt{9}{442.061}{132.719}{\Character{41}}%
\TextAt{13}{446.614}{132.719}{\Character{58}}%
\Vector{q 0 G 0 g 0.399 w 452.334 -670.282 m 452.334 -662.212 l S Q}%
\Vector{q 0 G 0 g 0.399 w 452.534 -662.412 m 459.707 -662.412 l S Q}%
\Vector{q 0 G 0 g 0.399 w 452.534 -670.083 m 459.707 -670.083 l S Q}%
\Vector{q 0 G 0 g 0.399 w 459.906 -670.282 m 459.906 -662.212 l S Q}%
\TextAt{8}{300.14}{43.033}{29}%
\Destination{78.23}{809.815}{nav.page.29}%
\Destination{79.36}{776.155}{nav.print.31}%
\Destination{79.36}{431.355}{nav.number.1.5.2}%
\Hotlink{179.532}{288.055}{30.936}{11.047}{nav.number.1.2.12}%
\Destination{79.226}{215.267}{errata.48}%
\Destination{79.226}{774}{errata.187}%
}

% Page 42 (body page 30)
\Page{%
\TextAt{2}{79.226}{796.943}{I.}%
\TextAt{2}{89.83}{796.943}{Algebraic}%
\TextAt{2}{138.648}{796.943}{Integers}%
\Vector{q 0 G 0 g 0.299 w 7.226 18.729 m 460.769 18.729 l S Q}%
\TextAt{8}{97.159}{761.124}{Using}%
\TextAt{8}{131.849}{761.124}{this}%
\TextAt{8}{156.46}{761.124}{proposition,}%
\TextAt{8}{223.482}{761.124}{Minkowski’s}%
\TextAt{8}{291.407}{761.124}{lattice}%
\TextAt{8}{328.966}{761.124}{point}%
\TextAt{8}{361.325}{761.124}{theorem}%
\TextAt{8}{408.639}{761.124}{now}%
\TextAt{8}{434.817}{761.124}{gives}%
\TextAt{8}{465.287}{761.124}{the}%
\TextAt{8}{487.232}{761.124}{following}%
\TextAt{8}{79.226}{744.078}{result,}%
\TextAt{8}{115.056}{744.078}{which}%
\TextAt{8}{148.542}{744.078}{is}%
\TextAt{8}{160.306}{744.078}{what}%
\TextAt{8}{189.572}{744.078}{we}%
\TextAt{8}{206.8}{744.078}{chiefly}%
\TextAt{8}{243.215}{744.078}{intend}%
\TextAt{8}{279.308}{744.078}{to}%
\TextAt{8}{293.618}{744.078}{use}%
\TextAt{8}{313.834}{744.078}{in}%
\TextAt{8}{327.487}{744.078}{our}%
\TextAt{8}{348.301}{744.078}{applications}%
\TextAt{8}{413.409}{744.078}{to}%
\TextAt{8}{427.719}{744.078}{number}%
\TextAt{8}{470.316}{744.078}{theory.}%
\TextAt{12}{79.226}{718.409}{Theorem}%
\TextAt{12}{135.937}{718.409}{(5.3).}%
\TextAt{8}{170.989}{718.409}{Let}%
\TextAt{14}{191.695}{718.409}{a}%
\NativeMathAt{200.999}{718.409}{11.955}{\neq}%
\TextAt{9}{213.409}{718.409}{0}%
\TextAt{8}{222.911}{718.409}{be}%
\TextAt{8}{238.572}{718.409}{an}%
\TextAt{8}{254.567}{718.409}{integral}%
\TextAt{8}{296.912}{718.409}{ideal}%
\TextAt{8}{324.612}{718.409}{of}%
\TextAt{13}{337.678}{718.409}{K}%
\TextAt{8}{348.496}{718.409}{,}%
\TextAt{8}{355.434}{718.409}{and}%
\TextAt{8}{377.933}{718.409}{let}%
\TextAt{13}{394.574}{718.409}{c}%
\TextAt{21}{399.612}{716.616}{\Character{28}}%
\TextAt{13}{408.11}{718.409}{\Character{62}}%
\TextAt{9}{420.535}{718.409}{0}%
\TextAt{8}{426.388}{718.409}{,}%
\TextAt{8}{433.327}{718.409}{for}%
\TextAt{13}{450.948}{718.409}{\Character{28}}%
\NativeMathAt{460.688}{718.409}{11.955}{\in}%
\TextAt{8}{471.979}{718.409}{Hom}%
\TextAt{9}{496.367}{718.409}{\Character{40}}%
\TextAt{13}{500.92}{718.409}{K}%
\TextAt{13}{511.082}{718.409}{\Character{59}}%
\TextAt{18}{516.331}{718.409}{C}%
\TextAt{9}{524.965}{718.409}{\Character{41}}%
\TextAt{8}{529.518}{718.409}{,}%
\TextAt{8}{79.226}{701.363}{be}%
\TextAt{8}{95.15}{701.363}{real}%
\TextAt{8}{117.913}{701.363}{numbers}%
\TextAt{8}{165.124}{701.363}{such}%
\TextAt{8}{191.521}{701.363}{that}%
\TextAt{13}{216.89}{701.363}{c}%
\TextAt{21}{221.928}{699.57}{\Character{28}}%
\TextAt{9}{230.425}{701.363}{\Character{61}}%
\TextAt{13}{242.851}{701.363}{c}%
\TextAt{17}{248.346}{699.57}{\Character{22}}%
\TextAt{21}{247.889}{699.57}{\Character{28}}%
\TextAt{8}{256.963}{701.363}{and}%
\TextAt{20}{258.856}{684.923}{\Character{89}}%
\TextAt{21}{264.155}{659.817}{\Character{28}}%
\TextAt{13}{276.124}{673.565}{c}%
\TextAt{21}{281.162}{671.772}{\Character{28}}%
\TextAt{13}{289.66}{673.565}{\Character{62}}%
\TextAt{13}{302.07}{673.565}{A}%
\TextAt{9}{310.86}{673.565}{\Character{40}}%
\TextAt{14}{315.413}{673.565}{o}%
\TextAt{21}{321.257}{671.772}{K}%
\TextAt{9}{332.78}{673.565}{\Character{58}}%
\TextAt{14}{339.353}{673.565}{a}%
\TextAt{9}{345.336}{673.565}{\Character{41}}%
\TextAt{13}{349.888}{673.565}{\Character{59}}%
\TextAt{8}{79.226}{635.695}{where}%
\TextAt{13}{113.035}{635.695}{A}%
\TextAt{9}{125.132}{635.695}{\Character{61}}%
\TextAt{9}{137.542}{635.695}{\Character{40}2\Character{47}}%
\TextAt{13}{153.815}{635.695}{\Character{25}}%
\TextAt{9}{160.885}{635.695}{\Character{41}}%
\TextAt{21}{165.437}{640.034}{s}%
\TextAt{20}{169.851}{645.917}{\Character{112}}%
\Vector{q 0 G 0 g 0.478 w 109.806 -127.844 m 130.733 -127.844 l S Q}%
\TextAt{16}{181.806}{635.695}{\Character{106}}%
\TextAt{13}{185.127}{635.695}{d}%
\TextAt{21}{191.21}{633.902}{K}%
\TextAt{16}{199.412}{635.695}{\Character{106}}%
\TextAt{8}{202.733}{635.695}{.}%
\TextAt{8}{211.182}{635.695}{Then}%
\TextAt{8}{241.739}{635.695}{there}%
\TextAt{8}{271.651}{635.695}{exists}%
\TextAt{13}{303.966}{635.695}{a}%
\NativeMathAt{313.431}{635.695}{11.955}{\in}%
\TextAt{14}{324.722}{635.695}{a}%
\TextAt{8}{330.706}{635.695}{,}%
\TextAt{13}{337.855}{635.695}{a}%
\NativeMathAt{347.321}{635.695}{11.955}{\neq}%
\TextAt{9}{359.731}{635.695}{0}%
\TextAt{8}{365.599}{635.695}{,}%
\TextAt{8}{372.748}{635.695}{such}%
\TextAt{8}{399.145}{635.695}{that}%
\TextAt{16}{213.477}{607.897}{\Character{106}}%
\TextAt{13}{216.798}{607.897}{\Character{28}}%
\TextAt{13}{223.206}{607.897}{a}%
\TextAt{16}{229.362}{607.897}{\Character{106}}%
\TextAt{13}{236.003}{607.897}{\Character{60}}%
\TextAt{13}{248.413}{607.897}{c}%
\TextAt{21}{253.467}{606.104}{\Character{28}}%
\TextAt{8}{282.554}{607.897}{for}%
\TextAt{8}{300.438}{607.897}{all}%
\TextAt{13}{316.697}{607.897}{\Character{28}}%
\NativeMathAt{326.437}{607.897}{11.955}{\in}%
\TextAt{8}{337.728}{607.897}{Hom}%
\TextAt{9}{362.117}{607.897}{\Character{40}}%
\TextAt{13}{366.669}{607.897}{K}%
\TextAt{13}{376.831}{607.897}{\Character{59}}%
\TextAt{18}{382.081}{607.897}{C}%
\TextAt{9}{390.715}{607.897}{\Character{41}}%
\TextAt{13}{395.267}{607.897}{\Character{58}}%
\TextAt{19}{79.226}{562.854}{Proof.}%
\TextAt{8}{116}{562.854}{The}%
\TextAt{8}{141.43}{562.854}{set}%
\TextAt{13}{161.074}{562.854}{X}%
\TextAt{9}{177.396}{562.854}{\Character{61}}%
\TextAt{16}{192.167}{562.854}{\Character{102}}%
\TextAt{9}{198.145}{562.854}{\Character{40}}%
\TextAt{13}{202.697}{562.854}{z}%
\TextAt{21}{208.135}{561.061}{\Character{28}}%
\TextAt{9}{213.311}{562.854}{\Character{41}}%
\NativeMathAt{223.53}{562.854}{11.955}{\in}%
\TextAt{13}{237.167}{562.854}{K}%
\TextAt{25}{247.126}{561.061}{R}%
\TextAt{16}{259.047}{562.854}{\Character{106}}%
\TextAt{16}{268.013}{562.854}{\Character{106}}%
\TextAt{13}{271.355}{562.854}{z}%
\TextAt{21}{276.793}{561.061}{\Character{28}}%
\TextAt{16}{281.969}{562.854}{\Character{106}}%
\TextAt{13}{290.956}{562.854}{\Character{60}}%
\TextAt{13}{305.721}{562.854}{c}%
\TextAt{21}{310.765}{561.061}{\Character{28}}%
\TextAt{16}{315.942}{562.854}{\Character{103}}%
\TextAt{8}{327.193}{562.854}{is}%
\TextAt{8}{340.334}{562.854}{centrally}%
\TextAt{8}{389.842}{562.854}{symmetric}%
\TextAt{8}{448.184}{562.854}{and}%
\TextAt{8}{472.323}{562.854}{convex.}%
\TextAt{8}{519.38}{562.854}{Its}%
\TextAt{8}{79.226}{545.808}{volume}%
\TextAt{8}{119.551}{545.808}{vol}%
\TextAt{9}{134.842}{545.808}{\Character{40}}%
\TextAt{13}{139.394}{545.808}{X}%
\TextAt{9}{150.049}{545.808}{\Character{41}}%
\TextAt{8}{158.499}{545.808}{can}%
\TextAt{8}{179.958}{545.808}{be}%
\TextAt{8}{195.883}{545.808}{computed}%
\TextAt{8}{249.86}{545.808}{via}%
\TextAt{8}{269.048}{545.808}{the}%
\TextAt{8}{289.205}{545.808}{map}%
\TextAt{8}{315.219}{545.808}{(5.1)}%
\TextAt{13}{217.082}{518.011}{f}%
\TextAt{9}{227.449}{518.01}{\Character{58}}%
\TextAt{13}{234.022}{518.011}{K}%
\TextAt{25}{243.981}{516.217}{R}%
\NativeMathAt{256.143}{524.788}{7.97}{\sim}%
\NativeArrowAt{253.556}{518.01}{11.955}{13.175}{\longrightarrow}%
\TextAt{20}{270.05}{529.368}{\Character{89}}%
\TextAt{21}{275.349}{504.262}{\Character{28}}%
\TextAt{18}{287.318}{518.011}{R}%
\TextAt{13}{295.953}{518.011}{\Character{59}}%
\TextAt{9}{325.107}{518.011}{\Character{40}}%
\TextAt{13}{329.659}{518.011}{z}%
\TextAt{21}{335.097}{516.217}{\Character{28}}%
\TextAt{9}{340.274}{518.011}{\Character{41}}%
\NativeArrowAt{348.147}{518.011}{11.955}{19.26}{\longmapsto}%
\TextAt{9}{370.729}{518.011}{\Character{40}}%
\TextAt{13}{375.281}{518.011}{x}%
\TextAt{21}{381.933}{516.217}{\Character{28}}%
\TextAt{9}{387.11}{518.011}{\Character{41}}%
\TextAt{13}{391.662}{518.011}{\Character{59}}%
\TextAt{8}{79.226}{482.352}{given}%
\TextAt{8}{110.509}{482.352}{by}%
\TextAt{13}{127.482}{482.352}{x}%
\TextAt{21}{134.134}{480.559}{\Character{26}}%
\TextAt{9}{143.505}{482.352}{\Character{61}}%
\TextAt{13}{157.147}{482.352}{z}%
\TextAt{21}{162.585}{480.559}{\Character{26}}%
\TextAt{8}{167.42}{482.352}{,}%
\TextAt{13}{175.461}{482.352}{x}%
\TextAt{21}{182.113}{480.559}{\Character{27}}%
\TextAt{9}{192.273}{482.352}{\Character{61}}%
\TextAt{16}{205.915}{482.352}{\Character{60}}%
\TextAt{9}{214.549}{482.352}{\Character{40}}%
\TextAt{13}{219.101}{482.352}{z}%
\TextAt{21}{224.539}{480.559}{\Character{27}}%
\TextAt{9}{230.163}{482.352}{\Character{41}}%
\TextAt{8}{234.715}{482.352}{,}%
\TextAt{13}{242.756}{482.352}{x}%
\TextAt{17}{249.854}{480.559}{\Character{22}}%
\TextAt{21}{249.408}{480.559}{\Character{27}}%
\TextAt{9}{259.569}{482.352}{\Character{61}}%
\TextAt{16}{273.21}{482.352}{\Character{61}}%
\TextAt{9}{281.844}{482.352}{\Character{40}}%
\TextAt{13}{286.397}{482.352}{z}%
\TextAt{21}{291.835}{480.559}{\Character{27}}%
\TextAt{9}{297.458}{482.352}{\Character{41}}%
\TextAt{8}{302.01}{482.352}{.}%
\TextAt{8}{312.599}{482.352}{It}%
\TextAt{8}{325.985}{482.352}{comes}%
\TextAt{8}{361.226}{482.352}{out}%
\TextAt{8}{382.753}{482.352}{to}%
\TextAt{8}{397.777}{482.352}{be}%
\TextAt{9}{414.415}{482.352}{2}%
\TextAt{21}{420.268}{486.69}{s}%
\TextAt{8}{429.293}{482.352}{times}%
\TextAt{8}{461.281}{482.352}{the}%
\TextAt{8}{482.151}{482.352}{Lebesgue-}%
\TextAt{8}{79.226}{465.306}{volume}%
\TextAt{8}{119.551}{465.306}{of}%
\TextAt{8}{132.881}{465.306}{the}%
\TextAt{8}{153.037}{465.306}{image}%
\TextAt{13}{172.693}{430.185}{f}%
\TextAt{9}{179.74}{430.185}{\Character{40}}%
\TextAt{13}{184.292}{430.185}{X}%
\TextAt{9}{194.947}{430.185}{\Character{41}}%
\TextAt{9}{202.802}{430.185}{\Character{61}}%
\TextAt{20}{215.246}{450.629}{\Character{40}}%
\TextAt{9}{224.876}{430.185}{\Character{40}}%
\TextAt{13}{229.429}{430.185}{x}%
\TextAt{21}{236.081}{428.392}{\Character{28}}%
\TextAt{9}{241.257}{430.185}{\Character{41}}%
\NativeMathAt{249.13}{430.185}{11.955}{\in}%
\TextAt{20}{260.421}{441.543}{\Character{89}}%
\TextAt{21}{265.72}{416.437}{\Character{28}}%
\TextAt{18}{277.69}{430.185}{R}%
\TextAt{20}{290.222}{451.107}{\Character{12}}%
\TextAt{20}{290.222}{443.934}{\Character{12}}%
\TextAt{20}{290.222}{436.761}{\Character{12}}%
\TextAt{20}{290.222}{429.588}{\Character{12}}%
\TextAt{20}{290.222}{422.414}{\Character{12}}%
\TextAt{16}{298.104}{430.185}{\Character{106}}%
\TextAt{13}{301.425}{430.185}{x}%
\TextAt{21}{308.077}{428.392}{\Character{26}}%
\TextAt{16}{312.912}{430.185}{\Character{106}}%
\TextAt{13}{319.553}{430.185}{\Character{60}}%
\TextAt{13}{331.963}{430.185}{c}%
\TextAt{21}{337.017}{428.392}{\Character{26}}%
\TextAt{13}{341.852}{430.185}{\Character{59}}%
\TextAt{13}{359.032}{430.185}{x}%
\TextAt{17}{365.703}{435.122}{2}%
\TextAt{21}{365.703}{427.23}{\Character{27}}%
\TextAt{9}{373.983}{430.185}{\Character{43}}%
\TextAt{13}{385.744}{430.185}{x}%
\TextAt{17}{392.397}{435.122}{2}%
\TextAt{17}{392.842}{427.23}{\Character{22}}%
\TextAt{21}{392.397}{427.23}{\Character{27}}%
\TextAt{13}{401.341}{430.185}{\Character{60}}%
\TextAt{13}{413.75}{430.185}{c}%
\TextAt{17}{418.804}{435.122}{2}%
\TextAt{21}{418.804}{427.23}{\Character{27}}%
\TextAt{20}{424.428}{450.629}{\Character{41}}%
\TextAt{13}{436.051}{430.185}{\Character{58}}%
\TextAt{8}{79.226}{395.196}{This}%
\TextAt{8}{105.946}{395.196}{gives}%
\TextAt{8}{135.301}{367.398}{vol}%
\TextAt{9}{150.592}{367.398}{\Character{40}}%
\TextAt{13}{155.144}{367.398}{X}%
\TextAt{9}{165.799}{367.398}{\Character{41}}%
\TextAt{9}{173.654}{367.398}{\Character{61}}%
\TextAt{9}{186.075}{367.398}{2}%
\TextAt{21}{191.951}{372.335}{s}%
\TextAt{8}{198.357}{367.398}{vol}%
\TextAt{7}{213.648}{365.605}{Lebesgue}%
\TextAt{9}{247.7}{367.398}{\Character{40}}%
\TextAt{13}{252.253}{367.398}{f}%
\TextAt{9}{259.299}{367.398}{\Character{40}}%
\TextAt{13}{263.852}{367.398}{X}%
\TextAt{9}{274.507}{367.398}{\Character{41}\Character{41}}%
\TextAt{9}{286.905}{367.398}{\Character{61}}%
\TextAt{9}{299.326}{367.398}{2}%
\TextAt{21}{305.211}{372.335}{s}%
\TextAt{20}{311.617}{378.756}{\Character{89}}%
\TextAt{21}{317.087}{353.65}{\Character{26}}%
\TextAt{9}{326.893}{367.398}{\Character{40}2}%
\TextAt{13}{337.299}{367.398}{c}%
\TextAt{21}{342.337}{365.605}{\Character{26}}%
\TextAt{9}{347.171}{367.398}{\Character{41}}%
\TextAt{20}{353.716}{378.756}{\Character{89}}%
\TextAt{21}{358.791}{353.65}{\Character{27}}%
\TextAt{9}{368.992}{367.398}{\Character{40}}%
\TextAt{13}{373.544}{367.398}{\Character{25}}%
\TextAt{13}{380.598}{367.398}{c}%
\TextAt{17}{385.652}{372.335}{2}%
\TextAt{21}{385.652}{364.443}{\Character{27}}%
\TextAt{9}{391.275}{367.398}{\Character{41}}%
\TextAt{9}{399.13}{367.398}{\Character{61}}%
\TextAt{9}{411.551}{367.398}{2}%
\TextAt{21}{417.427}{372.335}{r}%
\TextAt{17}{421.483}{372.335}{\Character{43}}%
\TextAt{21}{428.07}{372.335}{s}%
\TextAt{13}{432.484}{367.398}{\Character{25}}%
\TextAt{21}{439.553}{372.335}{s}%
\TextAt{20}{445.96}{378.756}{\Character{89}}%
\TextAt{21}{451.259}{353.65}{\Character{28}}%
\TextAt{13}{463.228}{367.398}{c}%
\TextAt{21}{468.266}{365.605}{\Character{28}}%
\TextAt{13}{473.443}{367.398}{\Character{58}}%
\TextAt{8}{79.226}{330.19}{Now}%
\TextAt{8}{105.886}{330.19}{using}%
\TextAt{8}{136.515}{330.19}{(5.2),}%
\TextAt{8}{167.742}{330.19}{we}%
\TextAt{8}{184.97}{330.19}{obtain}%
\TextAt{8}{177.627}{297.829}{vol}%
\TextAt{9}{192.917}{297.829}{\Character{40}}%
\TextAt{13}{197.47}{297.829}{X}%
\TextAt{9}{208.125}{297.829}{\Character{41}}%
\TextAt{13}{215.998}{297.829}{\Character{62}}%
\TextAt{9}{228.423}{297.829}{2}%
\TextAt{21}{234.276}{302.766}{r}%
\TextAt{17}{238.333}{302.766}{\Character{43}}%
\TextAt{21}{244.919}{302.766}{s}%
\TextAt{13}{249.334}{297.829}{\Character{25}}%
\TextAt{21}{256.403}{302.766}{s}%
\TextAt{20}{262.809}{314.686}{\Character{18}}%
\TextAt{9}{273.413}{305.917}{2}%
\Vector{q 0 G 0 g 0.478 w 200.805 -473.182 m 207.875 -473.182 l S Q}%
\TextAt{13}{272.805}{289.629}{\Character{25}}%
\TextAt{20}{281.07}{314.686}{\Character{19}}%
\TextAt{21}{289.87}{312.01}{s}%
\TextAt{20}{296.277}{308.635}{\Character{112}}%
\Vector{q 0 G 0 g 0.478 w 236.232 -465.126 m 257.159 -465.126 l S Q}%
\TextAt{16}{308.232}{297.829}{\Character{106}}%
\TextAt{13}{311.553}{297.829}{d}%
\TextAt{21}{317.636}{296.036}{K}%
\TextAt{16}{325.838}{297.829}{\Character{106}}%
\TextAt{9}{331.152}{297.829}{\Character{40}}%
\TextAt{14}{335.704}{297.829}{o}%
\TextAt{21}{341.548}{296.036}{K}%
\TextAt{9}{353.072}{297.829}{\Character{58}}%
\TextAt{14}{359.644}{297.829}{a}%
\TextAt{9}{365.627}{297.829}{\Character{41}}%
\TextAt{9}{373.482}{297.829}{\Character{61}}%
\TextAt{9}{385.903}{297.829}{2}%
\TextAt{21}{391.779}{302.766}{n}%
\TextAt{8}{399.408}{297.829}{vol}%
\TextAt{9}{414.699}{297.829}{\Character{40}\Character{0}\Character{41}}%
\TextAt{13}{431.117}{297.829}{\Character{58}}%
\TextAt{8}{79.226}{266.307}{Thus}%
\TextAt{8}{107.876}{266.307}{the}%
\TextAt{8}{127.033}{266.307}{hypothesis}%
\TextAt{8}{183.718}{266.307}{of}%
\TextAt{8}{196.049}{266.307}{Minkowski’s}%
\TextAt{8}{261.186}{266.307}{lattice}%
\TextAt{8}{295.957}{266.307}{point}%
\TextAt{8}{325.527}{266.307}{theorem}%
\TextAt{8}{370.053}{266.307}{is}%
\TextAt{8}{380.818}{266.307}{satisfied.}%
\TextAt{8}{430.035}{266.307}{So}%
\TextAt{8}{445.295}{266.307}{there}%
\TextAt{8}{474.208}{266.307}{does}%
\TextAt{8}{499.606}{266.307}{indeed}%
\TextAt{8}{79.226}{249.261}{exist}%
\TextAt{8}{106.926}{249.261}{a}%
\TextAt{8}{116.682}{249.261}{lattice}%
\TextAt{8}{152.452}{249.261}{point}%
\TextAt{13}{183.021}{249.261}{j}%
\TextAt{13}{188.532}{249.261}{a}%
\NativeMathAt{198.002}{249.261}{11.955}{\in}%
\TextAt{13}{209.293}{249.261}{X}%
\TextAt{8}{219.948}{249.261}{,}%
\TextAt{13}{227.097}{249.261}{a}%
\NativeMathAt{236.563}{249.261}{11.955}{\neq}%
\TextAt{9}{248.973}{249.261}{0}%
\TextAt{8}{254.842}{249.261}{,}%
\TextAt{13}{261.991}{249.261}{a}%
\NativeMathAt{271.457}{249.261}{11.955}{\in}%
\TextAt{14}{282.748}{249.261}{a}%
\TextAt{8}{288.731}{249.261}{;}%
\TextAt{8}{295.88}{249.261}{in}%
\TextAt{8}{309.533}{249.261}{other}%
\TextAt{8}{340.102}{249.261}{words}%
\TextAt{16}{373.66}{249.261}{\Character{106}}%
\TextAt{13}{376.981}{249.261}{\Character{28}}%
\TextAt{13}{383.389}{249.261}{a}%
\TextAt{16}{389.545}{249.261}{\Character{106}}%
\TextAt{13}{396.187}{249.261}{\Character{60}}%
\TextAt{13}{408.596}{249.261}{c}%
\TextAt{21}{413.65}{247.468}{\Character{28}}%
\TextAt{8}{418.827}{249.261}{.}%
\Vector{q 0 G 0 g 0.399 w 452.334 -524.739 m 452.334 -516.669 l S Q}%
\Vector{q 0 G 0 g 0.399 w 452.534 -516.868 m 459.707 -516.868 l S Q}%
\Vector{q 0 G 0 g 0.399 w 452.534 -524.539 m 459.707 -524.539 l S Q}%
\Vector{q 0 G 0 g 0.399 w 459.906 -524.739 m 459.906 -516.669 l S Q}%
\TextAt{8}{97.159}{220.604}{There}%
\TextAt{8}{131.344}{220.604}{is}%
\TextAt{8}{143.484}{220.604}{also}%
\TextAt{8}{167.34}{220.604}{a}%
\TextAt{8}{177.472}{220.604}{multiplicative}%
\TextAt{8}{251.671}{220.604}{version}%
\TextAt{8}{291.786}{220.604}{of}%
\TextAt{8}{305.492}{220.604}{Minkowski}%
\TextAt{8}{364.138}{220.604}{theory.}%
\TextAt{8}{405.587}{220.604}{It}%
\TextAt{8}{418.635}{220.604}{is}%
\TextAt{8}{430.775}{220.604}{based}%
\TextAt{8}{463.729}{220.604}{on}%
\TextAt{8}{480.364}{220.604}{the}%
\TextAt{8}{500.897}{220.604}{homo-}%
\TextAt{8}{79.226}{203.558}{morphism}%
\TextAt{13}{244.644}{186.512}{j}%
\TextAt{9}{253.48}{186.512}{\Character{58}}%
\TextAt{13}{260.052}{186.512}{K}%
\TextAt{3}{270.87}{191.448}{\Character{3}}%
\NativeArrowAt{278.923}{186.512}{11.955}{19.26}{\longrightarrow}%
\TextAt{13}{301.505}{186.512}{K}%
\TextAt{3}{312.323}{191.448}{\Character{3}}%
\TextAt{25}{311.464}{183.557}{C}%
\TextAt{9}{321.04}{186.512}{\Character{61}}%
\TextAt{20}{333.465}{197.87}{\Character{89}}%
\TextAt{21}{338.764}{172.764}{\Character{28}}%
\TextAt{18}{350.734}{186.512}{C}%
\TextAt{3}{359.368}{191.448}{\Character{3}}%
\TextAt{13}{364.1}{186.512}{\Character{58}}%
\TextAt{8}{79.226}{156.315}{The}%
\TextAt{8}{103.28}{156.315}{multiplicative}%
\TextAt{8}{177.103}{156.315}{group}%
\TextAt{13}{210.279}{156.315}{K}%
\TextAt{3}{221.096}{160.653}{\Character{3}}%
\TextAt{25}{220.238}{153.25}{C}%
\TextAt{8}{230.39}{156.315}{admits}%
\TextAt{8}{268.825}{156.315}{the}%
\TextAt{8}{288.982}{156.315}{homomorphism}%
\TextAt{13}{267.998}{128.517}{N}%
\TextAt{9}{281.942}{128.517}{\Character{58}}%
\TextAt{13}{288.514}{128.517}{K}%
\TextAt{3}{299.332}{133.453}{\Character{3}}%
\TextAt{25}{298.473}{125.561}{C}%
\NativeArrowAt{308.049}{128.517}{11.955}{19.26}{\longrightarrow}%
\TextAt{18}{330.631}{128.517}{C}%
\TextAt{3}{339.265}{133.453}{\Character{3}}%
\TextAt{8}{79.226}{100.719}{given}%
\TextAt{8}{109.796}{100.719}{by}%
\TextAt{8}{126.055}{100.719}{the}%
\TextAt{8}{146.211}{100.719}{product}%
\TextAt{8}{190.11}{100.719}{of}%
\TextAt{8}{203.44}{100.719}{the}%
\TextAt{8}{223.597}{100.719}{coordinates.}%
\TextAt{8}{290.327}{100.719}{The}%
\TextAt{8}{314.38}{100.719}{composite}%
\TextAt{13}{263.735}{72.921}{K}%
\TextAt{3}{274.553}{77.858}{\Character{3}}%
\TextAt{21}{285.936}{80.185}{j}%
\NativeArrowAt{282.606}{72.921}{11.955}{11.955}{\longrightarrow}%
\TextAt{13}{297.882}{72.921}{K}%
\TextAt{3}{308.7}{77.858}{\Character{3}}%
\TextAt{25}{307.841}{69.966}{C}%
\TextAt{21}{320.004}{79.698}{N}%
\NativeArrowAt{317.416}{72.921}{11.955}{14.155}{\longrightarrow}%
\TextAt{18}{334.894}{72.921}{C}%
\TextAt{3}{343.528}{77.858}{\Character{3}}%
\TextAt{8}{300.14}{43.033}{30}%
\Destination{78.23}{809.815}{nav.page.30}%
\Destination{79.36}{733.455}{nav.print.32}%
\Destination{79.36}{732.455}{nav.number.1.5.3}%
\Hotlink{314.719}{543.655}{25.078}{11.047}{nav.number.1.5.1}%
\Hotlink{136.015}{327.955}{25.078}{11.047}{nav.number.1.5.2}%
}

% Page 43 (body page 31)
\Page{%
\TextAt{2}{79.226}{796.943}{I.}%
\TextAt{2}{89.83}{796.943}{Algebraic}%
\TextAt{2}{138.648}{796.943}{Integers}%
\Vector{q 0 G 0 g 0.299 w 7.226 18.729 m 460.769 18.729 l S Q}%
\TextAt{8}{79.226}{761.124}{is}%
\TextAt{8}{90.99}{761.124}{the}%
\TextAt{8}{111.146}{761.124}{usual}%
\TextAt{8}{141.776}{761.124}{norm}%
\TextAt{8}{172.345}{761.124}{of}%
\TextAt{13}{185.675}{761.124}{K}%
\TextAt{16}{199.813}{761.124}{\Character{106}}%
\TextAt{18}{206.455}{761.124}{Q}%
\TextAt{8}{215.754}{761.124}{,}%
\TextAt{13}{260.116}{744.078}{N}%
\TextAt{21}{269.49}{742.218}{K}%
\TextAt{3}{277.195}{742.218}{\Character{106}}%
\TextAt{25}{279.547}{742.218}{Q}%
\TextAt{9}{286.244}{744.078}{\Character{40}}%
\TextAt{13}{290.796}{744.078}{a}%
\TextAt{9}{296.941}{744.078}{\Character{41}}%
\TextAt{9}{304.796}{744.078}{\Character{61}}%
\TextAt{13}{317.24}{744.078}{N}%
\TextAt{9}{327.863}{744.078}{\Character{40}}%
\TextAt{13}{332.415}{744.078}{j}%
\TextAt{13}{337.926}{744.078}{a}%
\TextAt{9}{344.075}{744.078}{\Character{41}}%
\TextAt{13}{348.628}{744.078}{\Character{58}}%
\TextAt{8}{79.226}{720.556}{In}%
\TextAt{8}{93.61}{720.556}{order}%
\TextAt{8}{123.942}{720.556}{to}%
\TextAt{8}{138.015}{720.556}{produce}%
\TextAt{8}{182.322}{720.556}{a}%
\TextAt{8}{191.84}{720.556}{lattice}%
\TextAt{8}{227.372}{720.556}{from}%
\TextAt{8}{254.775}{720.556}{the}%
\TextAt{8}{274.694}{720.556}{multiplicative}%
\TextAt{8}{348.28}{720.556}{theory,}%
\TextAt{8}{387.112}{720.556}{we}%
\TextAt{8}{404.101}{720.556}{use}%
\TextAt{8}{424.08}{720.556}{the}%
\TextAt{8}{443.999}{720.556}{logarithm}%
\TextAt{8}{497.106}{720.556}{to}%
\TextAt{8}{511.178}{720.556}{pass}%
\TextAt{8}{79.226}{703.51}{from}%
\TextAt{8}{106.866}{703.51}{multiplicative}%
\TextAt{8}{180.69}{703.51}{to}%
\TextAt{8}{195}{703.51}{additive}%
\TextAt{8}{239.88}{703.51}{groups}%
\TextAt{13}{227.717}{674.509}{\Character{96}}%
\TextAt{9}{235.95}{674.509}{\Character{58}}%
\TextAt{18}{242.522}{674.509}{C}%
\TextAt{3}{251.157}{679.445}{\Character{3}}%
\NativeArrowAt{259.21}{674.509}{11.955}{19.26}{\longrightarrow}%
\TextAt{18}{281.792}{674.509}{R}%
\TextAt{13}{290.426}{674.509}{\Character{59}}%
\TextAt{13}{319.561}{674.509}{z}%
\NativeArrowAt{328.872}{674.509}{11.955}{19.26}{\longmapsto}%
\TextAt{8}{351.454}{674.509}{log}%
\TextAt{16}{368.415}{674.509}{\Character{106}}%
\TextAt{13}{371.735}{674.509}{z}%
\TextAt{16}{377.706}{674.509}{\Character{106}}%
\TextAt{13}{381.027}{674.509}{\Character{58}}%
\TextAt{8}{79.226}{645.508}{It}%
\TextAt{8}{91.899}{645.508}{induces}%
\TextAt{8}{133.574}{645.508}{a}%
\TextAt{8}{143.33}{645.508}{surjective}%
\TextAt{8}{195.741}{645.508}{homomorphism}%
\TextAt{13}{262.96}{616.507}{\Character{96}}%
\TextAt{9}{271.193}{616.507}{\Character{58}}%
\TextAt{13}{277.765}{616.507}{K}%
\TextAt{3}{288.583}{621.443}{\Character{3}}%
\TextAt{25}{287.724}{613.551}{C}%
\NativeArrowAt{297.299}{616.507}{11.955}{19.26}{\longrightarrow}%
\TextAt{20}{319.881}{627.864}{\Character{89}}%
\TextAt{21}{325.18}{602.758}{\Character{28}}%
\TextAt{18}{337.15}{616.507}{R}%
\TextAt{13}{345.784}{616.507}{\Character{59}}%
\TextAt{8}{79.226}{580.314}{and}%
\TextAt{8}{101.989}{580.314}{we}%
\TextAt{8}{119.216}{580.314}{obtain}%
\TextAt{8}{155.644}{580.314}{the}%
\TextAt{8}{175.8}{580.314}{commutative}%
\TextAt{8}{245.726}{580.314}{diagram}%
\TextAt{13}{229.797}{549.712}{K}%
\TextAt{3}{240.615}{554.051}{\Character{3}}%
\TextAt{13}{298.691}{549.712}{K}%
\TextAt{3}{309.508}{554.051}{\Character{3}}%
\TextAt{25}{308.65}{546.648}{C}%
\TextAt{20}{368.249}{558.679}{\Character{81}}%
\TextAt{21}{379.54}{546.225}{\Character{28}}%
\TextAt{18}{386.709}{549.712}{R}%
\TextAt{18}{230.557}{496.951}{Q}%
\TextAt{3}{239.855}{501.29}{\Character{3}}%
\TextAt{18}{300.115}{496.951}{C}%
\TextAt{3}{308.749}{501.29}{\Character{3}}%
\TextAt{18}{375.853}{496.951}{R}%
\TextAt{13}{384.487}{496.951}{\Character{58}}%
\NativeDiagram{%
\path[draw=black,line width=0.478bp] (178.699,-221.298) -- (220.862,-221.298);%
\path[draw=black,line width=0.478bp,line cap=round,line join=round] (218.614,-218.429) .. controls (219.069,-220.151) and (220.081,-220.963) .. (221.101,-221.298) .. controls (220.081,-221.633) and (219.069,-222.446) .. (218.614,-224.167);%
}%
\TextAt{21}{270.078}{557.067}{j}%
\NativeDiagram{%
\path[draw=black,line width=0.478bp] (165.572,-228.866) -- (165.572,-263.756);%
\path[draw=black,line width=0.478bp,line cap=round,line join=round] (168.441,-261.508) .. controls (166.72,-261.963) and (165.907,-262.975) .. (165.572,-263.995) .. controls (165.237,-262.975) and (164.425,-261.963) .. (162.703,-261.508);%
}%
\TextAt{21}{214.074}{526.343}{N}%
\TextAt{22}{220.796}{524.605}{K}%
\TextAt{23}{227.341}{524.605}{\Character{106}}%
\TextAt{29}{229.61}{524.605}{Q}%
\NativeDiagram{%
\path[draw=black,line width=0.478bp] (248.258,-221.298) -- (290.42,-221.298);%
\path[draw=black,line width=0.478bp,line cap=round,line join=round] (288.173,-218.429) .. controls (288.627,-220.151) and (289.639,-220.963) .. (290.659,-221.298) .. controls (289.639,-221.633) and (288.627,-222.446) .. (288.173,-224.167);%
}%
\TextAt{21}{339.848}{555.517}{\Character{96}}%
\NativeDiagram{%
\path[draw=black,line width=0.478bp] (234.799,-231.931) -- (234.799,-263.756);%
\path[draw=black,line width=0.478bp,line cap=round,line join=round] (237.668,-261.508) .. controls (235.946,-261.963) and (235.134,-262.975) .. (234.799,-263.995) .. controls (234.464,-262.975) and (233.651,-261.963) .. (231.93,-261.508);%
}%
\TextAt{21}{309.614}{523.194}{N}%
\NativeDiagram{%
\path[draw=black,line width=0.478bp] (309.798,-232.353) -- (309.798,-263.756);%
\path[draw=black,line width=0.478bp,line cap=round,line join=round] (312.667,-261.508) .. controls (310.945,-261.963) and (310.132,-262.975) .. (309.798,-263.995) .. controls (309.463,-262.975) and (308.65,-261.963) .. (306.928,-261.508);%
}%
\TextAt{7}{384.613}{523.012}{Tr}%
\NativeDiagram{%
\path[draw=black,line width=0.478bp] (177.94,-274.06) -- (222.285,-274.06);%
\path[draw=black,line width=0.478bp,line cap=round,line join=round] (220.038,-271.191) .. controls (220.492,-272.912) and (221.504,-273.725) .. (222.524,-274.06) .. controls (221.504,-274.395) and (220.492,-275.208) .. (220.038,-276.929);%
\path[draw=black,line width=0.478bp] (246.834,-274.06) -- (298.024,-274.06);%
\path[draw=black,line width=0.478bp,line cap=round,line join=round] (295.777,-271.191) .. controls (296.231,-272.912) and (297.244,-273.725) .. (298.264,-274.06) .. controls (297.244,-274.395) and (296.231,-275.208) .. (295.777,-276.929);%
}%
\TextAt{21}{342.939}{502.755}{\Character{96}}%
\TextAt{8}{79.226}{468.661}{The}%
\TextAt{8}{103.556}{468.661}{involution}%
\TextAt{13}{158.801}{468.661}{F}%
\NativeMathAt{171.795}{468.661}{11.955}{\in}%
\TextAt{13}{183.556}{468.661}{G}%
\TextAt{9}{192.789}{468.661}{\Character{40}}%
\TextAt{18}{197.342}{468.661}{C}%
\TextAt{16}{209.766}{468.661}{\Character{106}}%
\TextAt{18}{216.878}{468.661}{R}%
\TextAt{9}{225.512}{468.661}{\Character{41}}%
\TextAt{8}{234.237}{468.661}{acts}%
\TextAt{8}{258.638}{468.661}{on}%
\TextAt{8}{275.173}{468.661}{all}%
\TextAt{8}{291.708}{468.661}{groups}%
\TextAt{8}{329.774}{468.661}{in}%
\TextAt{8}{343.702}{468.661}{this}%
\TextAt{8}{366.8}{468.661}{diagram,}%
\TextAt{8}{415.933}{468.661}{trivially}%
\TextAt{8}{460.443}{468.661}{on}%
\TextAt{13}{476.978}{468.661}{K}%
\TextAt{3}{487.795}{473}{\Character{3}}%
\TextAt{8}{492.528}{468.661}{,}%
\TextAt{8}{500.021}{468.661}{on}%
\TextAt{13}{516.556}{468.661}{K}%
\TextAt{3}{527.374}{473}{\Character{3}}%
\TextAt{25}{526.515}{465.597}{C}%
\TextAt{8}{79.226}{451.615}{as}%
\TextAt{8}{93.596}{451.615}{before,}%
\TextAt{8}{131.96}{451.615}{and}%
\TextAt{8}{154.723}{451.615}{on}%
\TextAt{8}{170.982}{451.615}{the}%
\TextAt{8}{191.139}{451.615}{points}%
\TextAt{13}{226.323}{451.615}{x}%
\TextAt{9}{236.295}{451.615}{\Character{61}}%
\TextAt{9}{248.705}{451.615}{\Character{40}}%
\TextAt{13}{253.273}{451.615}{x}%
\TextAt{21}{259.925}{449.822}{\Character{28}}%
\TextAt{9}{265.102}{451.615}{\Character{41}}%
\NativeMathAt{272.975}{451.615}{11.955}{\in}%
\TextAt{20}{284.266}{460.582}{\Character{81}}%
\TextAt{21}{295.557}{448.128}{\Character{28}}%
\TextAt{18}{302.726}{451.615}{R}%
\TextAt{8}{315.258}{451.615}{by}%
\TextAt{9}{331.517}{451.615}{\Character{40}}%
\TextAt{13}{336.069}{451.615}{F}%
\TextAt{13}{345.251}{451.615}{x}%
\TextAt{9}{351.925}{451.615}{\Character{41}}%
\TextAt{21}{356.477}{449.822}{\Character{28}}%
\TextAt{9}{364.974}{451.615}{\Character{61}}%
\TextAt{13}{377.4}{451.615}{x}%
\TextAt{17}{384.509}{449.822}{\Character{22}}%
\TextAt{21}{384.052}{449.822}{\Character{28}}%
\TextAt{8}{389.228}{451.615}{.}%
\TextAt{8}{397.677}{451.615}{One}%
\TextAt{8}{422.388}{451.615}{clearly}%
\TextAt{8}{459.784}{451.615}{has}%
\TextAt{13}{136.568}{422.614}{F}%
\TextAt{16}{148.428}{422.614}{\Character{14}}%
\TextAt{13}{157.062}{422.614}{j}%
\TextAt{9}{165.899}{422.614}{\Character{61}}%
\TextAt{13}{178.324}{422.614}{j\Character{59}}%
\TextAt{13}{212.325}{422.614}{F}%
\TextAt{16}{224.204}{422.614}{\Character{14}}%
\TextAt{13}{232.838}{422.614}{\Character{96}}%
\TextAt{9}{241.071}{422.614}{\Character{61}}%
\TextAt{13}{253.497}{422.614}{\Character{96}}%
\TextAt{16}{261.065}{422.614}{\Character{14}}%
\TextAt{13}{269.7}{422.614}{F}%
\TextAt{13}{277.59}{422.614}{\Character{59}}%
\TextAt{13}{306.725}{422.614}{N}%
\TextAt{16}{320.036}{422.614}{\Character{14}}%
\TextAt{13}{328.671}{422.614}{F}%
\TextAt{9}{341.195}{422.614}{\Character{61}}%
\TextAt{13}{353.621}{422.614}{F}%
\TextAt{16}{365.481}{422.614}{\Character{14}}%
\TextAt{13}{374.115}{422.614}{N}%
\TextAt{13}{384.086}{422.614}{\Character{59}}%
\TextAt{8}{413.242}{422.614}{Tr}%
\TextAt{16}{428.242}{422.614}{\Character{14}}%
\TextAt{13}{434.219}{422.614}{F}%
\TextAt{9}{446.744}{422.614}{\Character{61}}%
\TextAt{8}{459.169}{422.614}{Tr}%
\TextAt{13}{472.176}{422.614}{\Character{59}}%
\TextAt{8}{79.226}{393.613}{i.e.,}%
\TextAt{8}{102.771}{393.613}{the}%
\TextAt{8}{124.079}{393.613}{homomorphisms}%
\TextAt{8}{212.516}{393.613}{of}%
\TextAt{8}{226.998}{393.613}{the}%
\TextAt{8}{248.306}{393.613}{diagram}%
\TextAt{8}{294.995}{393.613}{are}%
\TextAt{13}{315.658}{393.613}{G}%
\TextAt{9}{324.891}{393.613}{\Character{40}}%
\TextAt{18}{329.444}{393.613}{C}%
\TextAt{16}{343.362}{393.613}{\Character{106}}%
\TextAt{18}{351.966}{393.613}{R}%
\TextAt{9}{360.601}{393.613}{\Character{41}}%
\TextAt{8}{365.153}{393.613}{-homomorphisms.}%
\TextAt{8}{464.342}{393.613}{We}%
\TextAt{8}{485.638}{393.613}{now}%
\TextAt{8}{511.178}{393.613}{pass}%
\TextAt{8}{79.226}{376.567}{everywhere}%
\TextAt{8}{140.03}{376.567}{to}%
\TextAt{8}{154.34}{376.567}{the}%
\TextAt{8}{174.497}{376.567}{fixed}%
\TextAt{8}{202.783}{376.567}{modules}%
\TextAt{8}{248.691}{376.567}{under}%
\TextAt{13}{281.854}{376.567}{G}%
\TextAt{9}{291.088}{376.567}{\Character{40}}%
\TextAt{18}{295.64}{376.567}{C}%
\TextAt{16}{307.595}{376.567}{\Character{106}}%
\TextAt{18}{314.237}{376.567}{R}%
\TextAt{9}{322.871}{376.567}{\Character{41}}%
\TextAt{8}{331.321}{376.567}{and}%
\TextAt{8}{354.084}{376.567}{obtain}%
\TextAt{8}{390.511}{376.567}{the}%
\TextAt{8}{410.668}{376.567}{diagram}%
\TextAt{13}{223.003}{345.428}{K}%
\TextAt{3}{233.821}{349.766}{\Character{3}}%
\TextAt{13}{291.897}{345.428}{K}%
\TextAt{3}{302.714}{349.766}{\Character{3}}%
\TextAt{25}{301.856}{342.363}{R}%
\TextAt{9}{361.455}{345.428}{\Character{91}}%
\TextAt{20}{364.706}{354.394}{\Character{81}}%
\TextAt{21}{375.997}{341.941}{\Character{28}}%
\TextAt{18}{383.166}{345.428}{R}%
\TextAt{9}{391.801}{345.428}{\Character{93}}%
\TextAt{17}{395.052}{351.239}{\Character{43}}%
\TextAt{18}{223.763}{292.666}{Q}%
\TextAt{3}{233.061}{297.005}{\Character{3}}%
\TextAt{18}{293.321}{292.666}{R}%
\TextAt{3}{301.955}{297.005}{\Character{3}}%
\TextAt{18}{375.853}{292.666}{R}%
\TextAt{13}{384.487}{292.666}{\Character{58}}%
\NativeDiagram{%
\path[draw=black,line width=0.478bp] (171.905,-425.583) -- (214.068,-425.583);%
\path[draw=black,line width=0.478bp,line cap=round,line join=round] (211.82,-422.714) .. controls (212.274,-424.436) and (213.287,-425.249) .. (214.307,-425.584) .. controls (213.287,-425.918) and (212.274,-426.731) .. (211.82,-428.453);%
}%
\TextAt{21}{263.284}{352.782}{j}%
\NativeDiagram{%
\path[draw=black,line width=0.478bp] (158.778,-433.151) -- (158.778,-468.041);%
\path[draw=black,line width=0.478bp,line cap=round,line join=round] (161.647,-465.793) .. controls (159.926,-466.248) and (159.113,-467.26) .. (158.778,-468.28) .. controls (158.443,-467.26) and (157.63,-466.248) .. (155.909,-465.793);%
}%
\TextAt{21}{207.28}{322.058}{N}%
\TextAt{22}{214.002}{320.32}{K}%
\TextAt{23}{220.547}{320.32}{\Character{106}}%
\TextAt{29}{222.816}{320.32}{Q}%
\NativeDiagram{%
\path[draw=black,line width=0.478bp] (241.464,-425.583) -- (283.626,-425.583);%
\path[draw=black,line width=0.478bp,line cap=round,line join=round] (281.379,-422.714) .. controls (281.833,-424.436) and (282.845,-425.249) .. (283.865,-425.584) .. controls (282.845,-425.918) and (281.833,-426.731) .. (281.379,-428.453);%
}%
\TextAt{21}{333.054}{351.232}{\Character{96}}%
\NativeDiagram{%
\path[draw=black,line width=0.478bp] (228.005,-436.216) -- (228.005,-468.041);%
\path[draw=black,line width=0.478bp,line cap=round,line join=round] (230.874,-465.793) .. controls (229.152,-466.248) and (228.339,-467.26) .. (228.005,-468.28) .. controls (227.67,-467.26) and (226.857,-466.248) .. (225.136,-465.793);%
}%
\TextAt{21}{302.82}{318.909}{N}%
\NativeDiagram{%
\path[draw=black,line width=0.478bp] (309.798,-436.638) -- (309.798,-468.041);%
\path[draw=black,line width=0.478bp,line cap=round,line join=round] (312.667,-465.793) .. controls (310.945,-466.248) and (310.132,-467.26) .. (309.798,-468.28) .. controls (309.463,-467.26) and (308.65,-466.248) .. (306.928,-465.793);%
}%
\TextAt{7}{384.613}{318.727}{Tr}%
\NativeDiagram{%
\path[draw=black,line width=0.478bp] (171.146,-478.345) -- (215.491,-478.345);%
\path[draw=black,line width=0.478bp,line cap=round,line join=round] (213.244,-475.476) .. controls (213.698,-477.198) and (214.71,-478.011) .. (215.73,-478.345) .. controls (214.71,-478.68) and (213.698,-479.493) .. (213.244,-481.215);%
\path[draw=black,line width=0.478bp] (240.04,-478.345) -- (298.024,-478.345);%
\path[draw=black,line width=0.478bp,line cap=round,line join=round] (295.777,-475.476) .. controls (296.231,-477.198) and (297.244,-478.011) .. (298.264,-478.345) .. controls (297.244,-478.68) and (296.231,-479.493) .. (295.777,-481.215);%
}%
\TextAt{21}{339.542}{298.47}{\Character{96}}%
\TextAt{8}{79.226}{264.189}{The}%
\TextAt{18}{104.543}{264.189}{R}%
\TextAt{8}{113.177}{264.189}{-vector}%
\TextAt{8}{153.461}{264.189}{space}%
\TextAt{9}{185.999}{264.189}{\Character{91}}%
\TextAt{20}{189.25}{273.155}{\Character{81}}%
\TextAt{21}{200.541}{260.702}{\Character{28}}%
\TextAt{18}{207.71}{264.189}{R}%
\TextAt{9}{216.345}{264.189}{\Character{93}}%
\TextAt{17}{219.596}{268.527}{\Character{43}}%
\TextAt{8}{231.841}{264.189}{is}%
\TextAt{8}{244.867}{264.189}{explicitly}%
\TextAt{8}{296.856}{264.189}{given}%
\TextAt{8}{328.688}{264.189}{as}%
\TextAt{8}{344.321}{264.189}{follows.}%
\TextAt{8}{391.095}{264.189}{Separate}%
\TextAt{8}{440.489}{264.189}{as}%
\TextAt{8}{456.122}{264.189}{before}%
\TextAt{8}{492.497}{264.189}{the}%
\TextAt{8}{513.916}{264.189}{em-}%
\TextAt{8}{79.226}{247.143}{beddings}%
\TextAt{13}{129.833}{247.143}{\Character{28}}%
\TextAt{9}{142.038}{247.143}{\Character{58}}%
\TextAt{13}{151.076}{247.143}{K}%
\NativeMathAt{167.681}{247.143}{11.955}{\rightarrow}%
\TextAt{18}{185.422}{247.143}{C}%
\TextAt{8}{199.401}{247.143}{into}%
\TextAt{8}{224.591}{247.143}{real}%
\TextAt{8}{248.8}{247.143}{ones,}%
\TextAt{13}{279.935}{247.143}{\Character{26}}%
\TextAt{17}{285.973}{245.35}{1}%
\TextAt{13}{290.705}{247.143}{\Character{59}}%
\TextAt{13}{295.929}{247.143}{\Character{58}}%
\TextAt{13}{301.166}{247.143}{\Character{58}}%
\TextAt{13}{306.39}{247.143}{\Character{58}}%
\TextAt{13}{311.627}{247.143}{\Character{59}}%
\TextAt{13}{316.863}{247.143}{\Character{26}}%
\TextAt{21}{322.963}{245.35}{r}%
\TextAt{8}{327.518}{247.143}{,}%
\TextAt{8}{336.476}{247.143}{and}%
\TextAt{8}{360.686}{247.143}{pairs}%
\TextAt{8}{390.813}{247.143}{of}%
\TextAt{8}{405.59}{247.143}{complex}%
\TextAt{8}{452.885}{247.143}{conjugate}%
\TextAt{8}{507.341}{247.143}{ones,}%
\TextAt{13}{79.226}{230.097}{\Character{27}}%
\TextAt{17}{85.88}{228.304}{1}%
\TextAt{13}{90.612}{230.097}{\Character{59}}%
\TextAt{9}{96.471}{230.097}{\Character{22}}%
\TextAt{13}{95.856}{230.097}{\Character{27}}%
\TextAt{17}{102.51}{228.304}{1}%
\TextAt{13}{107.242}{230.097}{\Character{59}}%
\TextAt{13}{112.466}{230.097}{\Character{58}}%
\TextAt{13}{117.703}{230.097}{\Character{58}}%
\TextAt{13}{122.927}{230.097}{\Character{58}}%
\TextAt{13}{128.164}{230.097}{\Character{59}}%
\TextAt{13}{133.4}{230.097}{\Character{27}}%
\TextAt{21}{140.116}{228.304}{s}%
\TextAt{13}{144.53}{230.097}{\Character{59}}%
\TextAt{9}{150.389}{230.097}{\Character{22}}%
\TextAt{13}{149.774}{230.097}{\Character{27}}%
\TextAt{21}{156.428}{228.304}{s}%
\TextAt{8}{160.842}{230.097}{.}%
\TextAt{8}{169.176}{230.097}{We}%
\TextAt{8}{188.978}{230.097}{obtain}%
\TextAt{8}{225.063}{230.097}{a}%
\TextAt{8}{234.475}{230.097}{decomposition}%
\TextAt{8}{311.267}{230.097}{which}%
\TextAt{8}{344.411}{230.097}{is}%
\TextAt{8}{355.832}{230.097}{analogous}%
\TextAt{8}{409.551}{230.097}{to}%
\TextAt{8}{423.519}{230.097}{the}%
\TextAt{8}{443.333}{230.097}{one}%
\TextAt{8}{464.449}{230.097}{we}%
\TextAt{8}{481.334}{230.097}{saw}%
\TextAt{8}{503.491}{230.097}{above}%
\TextAt{8}{79.226}{213.051}{for}%
\TextAt{9}{97.111}{213.051}{\Character{91}}%
\TextAt{20}{100.363}{222.017}{\Character{81}}%
\TextAt{21}{111.654}{209.564}{\Character{28}}%
\TextAt{18}{118.823}{213.051}{C}%
\TextAt{9}{127.457}{213.051}{\Character{93}}%
\TextAt{17}{130.709}{217.389}{\Character{43}}%
\TextAt{20}{222.402}{206.429}{\Character{34}}%
\TextAt{20}{229.375}{197.343}{\Character{89}}%
\TextAt{21}{234.674}{172.237}{\Character{28}}%
\TextAt{18}{246.644}{185.986}{R}%
\TextAt{20}{255.278}{206.429}{\Character{35}}%
\TextAt{17}{262.252}{203.752}{\Character{43}}%
\TextAt{9}{272.658}{185.986}{\Character{61}}%
\TextAt{20}{285.083}{197.343}{\Character{89}}%
\TextAt{21}{290.553}{172.237}{\Character{26}}%
\TextAt{18}{302.352}{185.986}{R}%
\TextAt{16}{313.643}{185.986}{\Character{2}}%
\TextAt{20}{325.598}{197.343}{\Character{89}}%
\TextAt{21}{330.673}{172.237}{\Character{27}}%
\TextAt{9}{340.874}{185.986}{\Character{91}}%
\TextAt{18}{344.126}{185.986}{R}%
\TextAt{16}{355.417}{185.986}{\Character{2}}%
\TextAt{18}{367.372}{185.986}{R}%
\TextAt{9}{376.006}{185.986}{\Character{93}}%
\TextAt{17}{379.258}{190.922}{\Character{43}}%
\TextAt{13}{386.342}{185.986}{\Character{58}}%
\TextAt{8}{79.226}{152.866}{The}%
\TextAt{8}{103.828}{152.866}{factor}%
\TextAt{9}{137.874}{152.866}{\Character{91}}%
\TextAt{18}{141.126}{152.866}{R}%
\TextAt{16}{152.791}{152.866}{\Character{2}}%
\TextAt{18}{165.12}{152.866}{R}%
\TextAt{9}{173.754}{152.866}{\Character{93}}%
\TextAt{17}{177.005}{157.204}{\Character{43}}%
\TextAt{8}{188.536}{152.866}{now}%
\TextAt{8}{213.472}{152.866}{consists}%
\TextAt{8}{257.131}{152.866}{of}%
\TextAt{8}{271.009}{152.866}{the}%
\TextAt{8}{291.713}{152.866}{points}%
\TextAt{9}{327.445}{152.866}{\Character{40}}%
\TextAt{13}{331.998}{152.866}{x\Character{59}}%
\TextAt{13}{343.87}{152.866}{x}%
\TextAt{9}{350.546}{152.866}{\Character{41}}%
\TextAt{8}{355.098}{152.866}{,}%
\TextAt{8}{362.933}{152.866}{and}%
\TextAt{8}{386.243}{152.866}{we}%
\TextAt{8}{404.019}{152.866}{identify}%
\TextAt{8}{447.163}{152.866}{it}%
\TextAt{8}{459.415}{152.866}{with}%
\TextAt{18}{486.624}{152.866}{R}%
\TextAt{8}{499.703}{152.866}{by}%
\TextAt{8}{516.51}{152.866}{the}%
\TextAt{8}{79.226}{135.82}{map}%
\TextAt{9}{105.241}{135.82}{\Character{40}}%
\TextAt{13}{109.793}{135.82}{x\Character{59}}%
\TextAt{13}{121.665}{135.82}{x}%
\TextAt{9}{128.341}{135.82}{\Character{41}}%
\NativeArrowAt{136.214}{135.82}{11.955}{11.955}{\mapsto}%
\TextAt{9}{151.49}{135.82}{2}%
\TextAt{13}{157.343}{135.82}{x}%
\TextAt{8}{163.996}{135.82}{.}%
\TextAt{8}{172.444}{135.82}{In}%
\TextAt{8}{187.065}{135.82}{this}%
\TextAt{8}{209.887}{135.82}{way}%
\TextAt{8}{233.63}{135.82}{we}%
\TextAt{8}{250.858}{135.82}{obtain}%
\TextAt{8}{287.285}{135.82}{an}%
\TextAt{8}{303.544}{135.82}{isomorphism}%
\TextAt{20}{261.089}{117.741}{\Character{34}}%
\TextAt{20}{268.062}{108.655}{\Character{89}}%
\TextAt{21}{273.361}{83.549}{\Character{28}}%
\TextAt{18}{285.331}{97.298}{R}%
\TextAt{20}{293.965}{117.741}{\Character{35}}%
\TextAt{17}{300.939}{115.065}{\Character{43}}%
\NativeMathAt{311.345}{97.298}{11.955}{\simeq}%
\TextAt{18}{323.964}{97.298}{R}%
\TextAt{21}{332.598}{102.234}{r}%
\TextAt{17}{336.655}{102.234}{\Character{43}}%
\TextAt{21}{343.241}{102.234}{s}%
\TextAt{13}{347.655}{97.298}{\Character{59}}%
\TextAt{8}{300.14}{43.033}{31}%
\Destination{78.23}{809.815}{nav.page.31}%
\Destination{79.36}{735.555}{nav.print.33}%
}

% Page 44 (body page 32)
\Page{%
\TextAt{2}{79.226}{796.943}{I.}%
\TextAt{2}{89.83}{796.943}{Algebraic}%
\TextAt{2}{138.648}{796.943}{Integers}%
\Vector{q 0 G 0 g 0.299 w 7.226 18.729 m 460.769 18.729 l S Q}%
\TextAt{8}{79.226}{761.124}{which}%
\TextAt{8}{112.713}{761.124}{again}%
\TextAt{8}{143.939}{761.124}{transforms}%
\TextAt{8}{202.281}{761.124}{the}%
\TextAt{8}{222.437}{761.124}{map}%
\TextAt{8}{248.452}{761.124}{Tr}%
\TextAt{9}{264.78}{761.124}{\Character{58}}%
\TextAt{9}{271.331}{761.124}{\Character{91}}%
\TextAt{20}{274.604}{770.09}{\Character{81}}%
\TextAt{21}{285.895}{757.637}{\Character{28}}%
\TextAt{18}{293.064}{761.124}{R}%
\TextAt{9}{301.698}{761.124}{\Character{93}}%
\TextAt{17}{304.95}{765.462}{\Character{43}}%
\NativeMathAt{315.355}{761.124}{11.955}{\rightarrow}%
\TextAt{18}{330.631}{761.124}{R}%
\TextAt{8}{343.163}{761.124}{into}%
\TextAt{8}{366.906}{761.124}{the}%
\TextAt{8}{387.062}{761.124}{usual}%
\TextAt{8}{417.691}{761.124}{map}%
\TextAt{8}{265.433}{732.123}{Tr}%
\TextAt{9}{281.761}{732.123}{\Character{58}}%
\TextAt{18}{288.334}{732.123}{R}%
\TextAt{21}{296.968}{737.059}{r}%
\TextAt{17}{301.025}{737.059}{\Character{43}}%
\TextAt{21}{307.611}{737.059}{s}%
\NativeArrowAt{315.346}{732.123}{11.955}{19.26}{\longrightarrow}%
\TextAt{18}{337.928}{732.123}{R}%
\TextAt{8}{79.226}{703.121}{given}%
\TextAt{8}{109.796}{703.121}{by}%
\TextAt{8}{126.055}{703.121}{the}%
\TextAt{8}{146.211}{703.121}{sum}%
\TextAt{8}{170.982}{703.121}{of}%
\TextAt{8}{184.312}{703.121}{the}%
\TextAt{8}{204.469}{703.121}{coordinates.}%
\TextAt{8}{271.198}{703.121}{Identifying}%
\TextAt{9}{330.376}{703.121}{\Character{91}}%
\TextAt{20}{333.628}{712.088}{\Character{81}}%
\TextAt{21}{344.919}{699.634}{\Character{28}}%
\TextAt{18}{352.088}{703.121}{R}%
\TextAt{9}{360.722}{703.121}{\Character{93}}%
\TextAt{17}{363.974}{707.46}{\Character{43}}%
\TextAt{8}{374.956}{703.121}{with}%
\TextAt{18}{401.616}{703.121}{R}%
\TextAt{21}{410.25}{707.46}{r}%
\TextAt{17}{414.307}{707.46}{\Character{43}}%
\TextAt{21}{420.893}{707.46}{s}%
\TextAt{8}{425.308}{703.121}{,}%
\TextAt{8}{432.457}{703.121}{the}%
\TextAt{8}{452.613}{703.121}{homomorphism}%
\TextAt{13}{265.691}{674.12}{\Character{96}}%
\TextAt{9}{273.924}{674.12}{\Character{58}}%
\TextAt{13}{280.497}{674.12}{K}%
\TextAt{3}{291.314}{679.056}{\Character{3}}%
\TextAt{25}{290.456}{671.165}{R}%
\NativeArrowAt{300.031}{674.12}{11.955}{19.26}{\longrightarrow}%
\TextAt{18}{322.613}{674.12}{R}%
\TextAt{21}{331.247}{679.056}{r}%
\TextAt{17}{335.304}{679.056}{\Character{43}}%
\TextAt{21}{341.89}{679.056}{s}%
\TextAt{8}{79.226}{645.119}{is}%
\TextAt{8}{90.99}{645.119}{given}%
\TextAt{8}{121.559}{645.119}{by}%
\TextAt{13}{169.13}{628.073}{\Character{96}}%
\TextAt{9}{174.042}{628.073}{\Character{40}}%
\TextAt{13}{178.595}{628.073}{x}%
\TextAt{9}{185.247}{628.073}{\Character{41}}%
\TextAt{9}{193.102}{628.073}{\Character{61}}%
\TextAt{9}{205.523}{628.073}{\Character{40}}%
\TextAt{8}{210.098}{628.073}{log}%
\TextAt{16}{227.058}{628.073}{\Character{106}}%
\TextAt{13}{230.379}{628.073}{x}%
\TextAt{21}{237.031}{626.28}{\Character{26}}%
\TextAt{24}{241.368}{625.173}{1}%
\TextAt{16}{246.017}{628.073}{\Character{106}}%
\TextAt{13}{249.338}{628.073}{\Character{59}}%
\TextAt{13}{254.562}{628.073}{\Character{58}}%
\TextAt{13}{259.799}{628.073}{\Character{58}}%
\TextAt{13}{265.023}{628.073}{\Character{58}}%
\TextAt{13}{270.26}{628.073}{\Character{59}}%
\TextAt{8}{275.558}{628.073}{log}%
\TextAt{16}{292.519}{628.073}{\Character{106}}%
\TextAt{13}{295.84}{628.073}{x}%
\TextAt{21}{302.492}{626.28}{\Character{26}}%
\TextAt{22}{306.828}{625.284}{r}%
\TextAt{16}{311.485}{628.073}{\Character{106}}%
\TextAt{13}{314.805}{628.073}{\Character{59}}%
\TextAt{8}{320.05}{628.073}{log}%
\TextAt{16}{337.01}{628.073}{\Character{106}}%
\TextAt{13}{340.331}{628.073}{x}%
\TextAt{21}{346.983}{626.28}{\Character{27}}%
\TextAt{24}{351.822}{625.173}{1}%
\TextAt{16}{356.471}{628.073}{\Character{106}}%
\TextAt{17}{359.792}{633.009}{2}%
\TextAt{13}{364.525}{628.073}{\Character{59}}%
\TextAt{13}{369.749}{628.073}{\Character{58}}%
\TextAt{13}{374.986}{628.073}{\Character{58}}%
\TextAt{13}{380.21}{628.073}{\Character{58}}%
\TextAt{13}{385.447}{628.073}{\Character{59}}%
\TextAt{8}{390.745}{628.073}{log}%
\TextAt{16}{407.706}{628.073}{\Character{106}}%
\TextAt{13}{411.027}{628.073}{x}%
\TextAt{21}{417.679}{626.28}{\Character{27}}%
\TextAt{22}{422.518}{625.284}{s}%
\TextAt{16}{427.008}{628.073}{\Character{106}}%
\TextAt{17}{430.329}{633.009}{2}%
\TextAt{9}{435.061}{628.073}{\Character{41}}%
\TextAt{13}{439.614}{628.073}{\Character{59}}%
\TextAt{8}{79.226}{604.551}{where}%
\TextAt{8}{113.035}{604.551}{we}%
\TextAt{8}{130.263}{604.551}{write}%
\TextAt{13}{160.175}{604.551}{x}%
\NativeMathAt{170.148}{604.551}{11.955}{\in}%
\TextAt{13}{181.439}{604.551}{K}%
\TextAt{3}{192.256}{608.89}{\Character{3}}%
\TextAt{25}{191.398}{601.487}{R}%
\NativeMathAt{200.973}{604.551}{11.955}{\subseteq}%
\TextAt{20}{213.592}{613.518}{\Character{81}}%
\TextAt{21}{224.883}{601.064}{\Character{28}}%
\TextAt{18}{232.052}{604.551}{C}%
\TextAt{3}{240.686}{608.89}{\Character{3}}%
\TextAt{8}{249.316}{604.551}{as}%
\TextAt{13}{263.686}{604.551}{x}%
\TextAt{9}{273.659}{604.551}{\Character{61}}%
\TextAt{9}{286.069}{604.551}{\Character{40}}%
\TextAt{13}{290.637}{604.551}{x}%
\TextAt{21}{297.289}{602.758}{\Character{28}}%
\TextAt{9}{302.466}{604.551}{\Character{41}}%
\TextAt{8}{307.018}{604.551}{.}%
\TextAt{12}{79.226}{578.539}{Exercise}%
\TextAt{12}{134.015}{578.539}{I.5.1.}%
\TextAt{8}{168.989}{578.539}{Write}%
\TextAt{8}{202.876}{578.539}{down}%
\TextAt{8}{235.149}{578.539}{a}%
\TextAt{8}{246.285}{578.539}{constant}%
\TextAt{13}{294.889}{578.539}{A}%
\TextAt{8}{308.942}{578.539}{which}%
\TextAt{8}{343.809}{578.539}{depends}%
\TextAt{8}{390.44}{578.539}{only}%
\TextAt{8}{417.512}{578.539}{on}%
\TextAt{13}{435.152}{578.539}{K}%
\TextAt{8}{451.248}{578.539}{such}%
\TextAt{8}{479.025}{578.539}{that}%
\TextAt{8}{505.775}{578.539}{every}%
\TextAt{8}{79.226}{561.493}{integral}%
\TextAt{8}{121.834}{561.493}{ideal}%
\TextAt{14}{149.798}{561.493}{a}%
\NativeMathAt{159.101}{561.493}{11.955}{\neq}%
\TextAt{9}{171.511}{561.493}{0}%
\TextAt{8}{181.277}{561.493}{of}%
\TextAt{13}{194.607}{561.493}{K}%
\TextAt{8}{209.322}{561.493}{contains}%
\TextAt{8}{255.242}{561.493}{an}%
\TextAt{8}{271.501}{561.493}{element}%
\TextAt{13}{314.743}{561.493}{a}%
\NativeMathAt{324.209}{561.493}{11.955}{\neq}%
\TextAt{9}{336.619}{561.493}{0}%
\TextAt{8}{346.385}{561.493}{satisfying}%
\TextAt{16}{155.385}{532.492}{\Character{106}}%
\TextAt{13}{158.705}{532.492}{\Character{28}}%
\TextAt{13}{165.113}{532.492}{a}%
\TextAt{16}{171.269}{532.492}{\Character{106}}%
\TextAt{13}{177.911}{532.492}{\Character{60}}%
\TextAt{13}{190.321}{532.492}{A}%
\TextAt{9}{199.112}{532.492}{\Character{40}}%
\TextAt{14}{203.664}{532.492}{o}%
\TextAt{21}{209.508}{530.699}{K}%
\TextAt{9}{221.032}{532.492}{\Character{58}}%
\TextAt{14}{227.604}{532.492}{a}%
\TextAt{9}{233.587}{532.492}{\Character{41}}%
\TextAt{17}{238.14}{537.428}{1\Character{47}}%
\TextAt{21}{246.608}{537.428}{n}%
\TextAt{8}{264.2}{532.492}{for}%
\TextAt{8}{282.085}{532.492}{all}%
\TextAt{13}{298.344}{532.492}{\Character{28}}%
\NativeMathAt{308.083}{532.492}{11.955}{\in}%
\TextAt{8}{319.374}{532.492}{Hom}%
\TextAt{9}{343.763}{532.492}{\Character{40}}%
\TextAt{13}{348.315}{532.492}{K}%
\TextAt{13}{358.477}{532.492}{\Character{59}}%
\TextAt{18}{363.727}{532.492}{C}%
\TextAt{9}{372.361}{532.492}{\Character{41}}%
\TextAt{13}{376.913}{532.492}{\Character{59}}%
\TextAt{13}{394.093}{532.492}{n}%
\TextAt{9}{404.421}{532.492}{\Character{61}}%
\TextAt{9}{416.831}{532.492}{\Character{91}}%
\TextAt{13}{420.098}{532.492}{K}%
\TextAt{9}{434.237}{532.492}{\Character{58}}%
\TextAt{18}{440.809}{532.492}{Q}%
\TextAt{9}{450.108}{532.492}{\Character{93}}%
\TextAt{13}{453.359}{532.492}{\Character{58}}%
\TextAt{12}{79.226}{488.547}{Exercise}%
\TextAt{12}{132.427}{488.547}{I.5.2.}%
\ErrataAt{166.01}{488.547}{368.99}{Show that for $t>0$ the convex, centrally symmetric set}%
\TextAt{13}{223.759}{452.235}{X}%
\TextAt{9}{237.735}{452.235}{\Character{61}}%
\TextAt{20}{250.161}{472.678}{\Character{40}}%
\TextAt{9}{259.791}{452.235}{\Character{40}}%
\TextAt{13}{264.343}{452.235}{z}%
\TextAt{21}{269.781}{450.442}{\Character{28}}%
\TextAt{9}{274.958}{452.235}{\Character{41}}%
\NativeMathAt{282.831}{452.235}{11.955}{\in}%
\TextAt{13}{294.122}{452.235}{K}%
\TextAt{25}{304.081}{450.442}{R}%
\TextAt{20}{314.233}{473.157}{\Character{12}}%
\TextAt{20}{314.233}{465.983}{\Character{12}}%
\TextAt{20}{314.233}{458.81}{\Character{12}}%
\TextAt{20}{314.233}{451.637}{\Character{12}}%
\TextAt{20}{314.233}{444.464}{\Character{12}}%
\TextAt{20}{322.115}{463.592}{\Character{88}}%
\TextAt{21}{328.41}{438.486}{\Character{28}}%
\TextAt{16}{341.376}{452.235}{\Character{106}}%
\TextAt{13}{344.697}{452.235}{z}%
\TextAt{21}{350.135}{450.442}{\Character{28}}%
\TextAt{16}{355.312}{452.235}{\Character{106}}%
\TextAt{13}{361.953}{452.235}{\Character{60}}%
\TextAt{13}{374.363}{452.235}{t}%
\TextAt{20}{378.606}{472.678}{\Character{41}}%
\TextAt{8}{79.226}{415.373}{has}%
\TextAt{8}{100.1}{415.373}{volume}%
\TextAt{8}{140.425}{415.373}{vol}%
\TextAt{9}{155.715}{415.373}{\Character{40}}%
\TextAt{13}{160.268}{415.373}{X}%
\TextAt{9}{170.923}{415.373}{\Character{41}}%
\TextAt{9}{178.778}{415.373}{\Character{61}}%
\TextAt{9}{191.199}{415.373}{2}%
\TextAt{21}{197.074}{419.711}{r}%
\TextAt{13}{201.629}{415.373}{\Character{25}}%
\TextAt{21}{208.698}{419.711}{s}%
\TextAt{13}{213.112}{415.373}{t}%
\TextAt{21}{217.34}{419.711}{n}%
\TextAt{9}{222.976}{415.373}{\Character{47}}%
\TextAt{13}{228.829}{415.373}{n}%
\TextAt{9}{235.817}{415.373}{\Character{33}}%
\TextAt{8}{242.966}{415.373}{(see}%
\TextAt{8}{266.434}{415.373}{chap.}%
\TextAt{8}{298.625}{415.373}{III,}%
\TextAt{8}{319.08}{415.373}{(2.15)).}%
\TextAt{12}{79.226}{383.383}{Exercise}%
\TextAt{12}{132.427}{383.383}{I.5.3.}%
\TextAt{8}{166.021}{383.383}{Show}%
\TextAt{8}{196.913}{383.383}{that}%
\TextAt{8}{222.282}{383.383}{in}%
\TextAt{8}{235.935}{383.383}{every}%
\TextAt{8}{266.827}{383.383}{ideal}%
\TextAt{14}{294.79}{383.383}{a}%
\NativeMathAt{304.094}{383.383}{11.955}{\neq}%
\TextAt{9}{316.504}{383.383}{0}%
\TextAt{8}{326.27}{383.383}{of}%
\TextAt{14}{339.6}{383.383}{o}%
\TextAt{21}{345.444}{381.59}{K}%
\TextAt{8}{357.544}{383.383}{there}%
\TextAt{8}{387.456}{383.383}{exists}%
\TextAt{8}{419.771}{383.383}{an}%
\TextAt{13}{436.03}{383.383}{a}%
\NativeMathAt{445.495}{383.383}{11.955}{\neq}%
\TextAt{9}{457.905}{383.383}{0}%
\TextAt{8}{467.671}{383.383}{such}%
\TextAt{8}{494.068}{383.383}{that}%
\TextAt{16}{246.592}{354.382}{\Character{106}}%
\TextAt{13}{249.912}{354.382}{N}%
\TextAt{21}{259.286}{352.522}{K}%
\TextAt{3}{266.991}{352.522}{\Character{106}}%
\TextAt{25}{269.343}{352.522}{Q}%
\TextAt{9}{276.04}{354.382}{\Character{40}}%
\TextAt{13}{280.593}{354.382}{a}%
\TextAt{9}{286.738}{354.382}{\Character{41}}%
\TextAt{16}{291.29}{354.382}{\Character{106}}%
\NativeMathAt{297.913}{354.382}{11.955}{\leq}%
\TextAt{13}{310.551}{354.382}{M}%
\TextAt{9}{323.125}{354.382}{\Character{40}}%
\TextAt{14}{327.677}{354.382}{o}%
\TextAt{21}{333.521}{352.589}{K}%
\TextAt{9}{345.045}{354.382}{\Character{58}}%
\TextAt{14}{351.617}{354.382}{a}%
\TextAt{9}{357.6}{354.382}{\Character{41}}%
\TextAt{13}{362.152}{354.382}{\Character{59}}%
\TextAt{8}{79.226}{319.966}{where}%
\TextAt{13}{113.035}{319.966}{M}%
\TextAt{9}{128.93}{319.966}{\Character{61}}%
\TextAt{13}{143.743}{328.054}{n}%
\TextAt{9}{150.731}{328.054}{\Character{33}}%
\Vector{q 0 G 0 g 0.478 w 70.551 -451.045 m 83.175 -451.045 l S Q}%
\TextAt{13}{142.551}{311.766}{n}%
\TextAt{21}{149.538}{315.219}{n}%
\TextAt{20}{158.363}{336.823}{\Character{18}}%
\TextAt{9}{168.967}{328.054}{4}%
\Vector{q 0 G 0 g 0.478 w 96.359 -451.045 m 103.428 -451.045 l S Q}%
\TextAt{13}{168.359}{311.766}{\Character{25}}%
\TextAt{20}{176.623}{336.823}{\Character{19}}%
\TextAt{21}{185.424}{334.147}{s}%
\TextAt{20}{191.83}{330.188}{\Character{112}}%
\Vector{q 0 G 0 g 0.478 w 131.786 -443.573 m 152.713 -443.573 l S Q}%
\TextAt{16}{203.786}{319.966}{\Character{106}}%
\TextAt{13}{207.106}{319.966}{d}%
\TextAt{21}{213.189}{318.173}{K}%
\TextAt{16}{221.392}{319.966}{\Character{106}}%
\TextAt{8}{228.61}{319.966}{(the}%
\TextAt{8}{253.321}{319.966}{so-called}%
\TextAt{12}{300.855}{319.966}{Minkowski}%
\TextAt{12}{368.27}{319.966}{bound}%
\TextAt{8}{405.26}{319.966}{).}%
\TextAt{12}{97.159}{295.657}{Hint:}%
\TextAt{8}{133.323}{295.657}{Use}%
\TextAt{8}{156.124}{295.657}{exercise}%
\TextAt{8}{199.739}{295.657}{2}%
\TextAt{8}{209.807}{295.657}{to}%
\TextAt{8}{224.43}{295.657}{proceed}%
\TextAt{8}{267.985}{295.657}{as}%
\TextAt{8}{282.668}{295.657}{in}%
\TextAt{8}{296.634}{295.657}{(5.3),}%
\TextAt{8}{328.252}{295.657}{and}%
\TextAt{8}{351.327}{295.657}{make}%
\TextAt{8}{382.21}{295.657}{use}%
\TextAt{8}{402.739}{295.657}{of}%
\TextAt{8}{416.382}{295.657}{the}%
\TextAt{8}{436.851}{295.657}{inequality}%
\TextAt{8}{491.476}{295.657}{between}%
\TextAt{8}{79.226}{278.611}{arithmetic}%
\TextAt{8}{135.81}{278.611}{and}%
\TextAt{8}{158.573}{278.611}{geometric}%
\TextAt{8}{211.905}{278.611}{means,}%
\TextAt{9}{239.445}{247.552}{1}%
\Vector{q 0 G 0 g 0.478 w 166.878 -531.546 m 173.865 -531.546 l S Q}%
\TextAt{13}{238.878}{231.264}{n}%
\TextAt{20}{249.053}{250.822}{\Character{88}}%
\TextAt{21}{255.348}{225.716}{\Character{28}}%
\TextAt{16}{268.314}{239.465}{\Character{106}}%
\TextAt{13}{271.635}{239.465}{z}%
\TextAt{21}{277.073}{237.671}{\Character{28}}%
\TextAt{16}{282.25}{239.465}{\Character{106}}%
\NativeMathAt{288.873}{239.465}{11.955}{\geq}%
\TextAt{20}{301.511}{259.908}{\Character{32}}%
\TextAt{20}{310.975}{250.822}{\Character{89}}%
\TextAt{21}{316.274}{225.716}{\Character{28}}%
\TextAt{16}{328.244}{239.465}{\Character{106}}%
\TextAt{13}{331.565}{239.465}{z}%
\TextAt{21}{337.003}{237.671}{\Character{28}}%
\TextAt{16}{342.179}{239.465}{\Character{106}}%
\TextAt{20}{345.5}{259.908}{\Character{33}}%
\TextAt{17}{354.965}{257.232}{1\Character{47}}%
\TextAt{21}{363.433}{257.232}{n}%
\TextAt{13}{371.062}{239.465}{\Character{58}}%
\TextAt{15}{79.226}{172.294}{§}%
\TextAt{15}{93.785}{172.294}{6}%
\TextAt{15}{120.693}{172.294}{The}%
\TextAt{15}{160.202}{172.294}{Class}%
\TextAt{15}{210.747}{172.294}{Number}%
\TextAt{8}{79.226}{140.408}{As}%
\TextAt{8}{95.866}{140.408}{a}%
\TextAt{8}{104.975}{140.408}{first}%
\TextAt{8}{128.453}{140.408}{application}%
\TextAt{8}{188.3}{140.408}{of}%
\TextAt{8}{200.983}{140.408}{Minkowski}%
\TextAt{8}{258.605}{140.408}{theory,}%
\TextAt{8}{297.109}{140.408}{we}%
\TextAt{8}{313.69}{140.408}{are}%
\TextAt{8}{332.554}{140.408}{going}%
\TextAt{8}{363.134}{140.408}{to}%
\TextAt{8}{376.797}{140.408}{show}%
\TextAt{8}{405.153}{140.408}{that}%
\TextAt{8}{429.875}{140.408}{the}%
\TextAt{8}{449.385}{140.408}{ideal}%
\TextAt{8}{476.701}{140.408}{class}%
\TextAt{8}{503.491}{140.408}{group}%
\TextAt{13}{79.226}{123.362}{C}%
\TextAt{13}{88.444}{123.362}{l}%
\TextAt{21}{91.947}{121.569}{K}%
\TextAt{9}{103.47}{123.362}{\Character{61}}%
\TextAt{13}{115.896}{123.362}{J}%
\TextAt{21}{122.366}{121.569}{K}%
\TextAt{9}{130.568}{123.362}{\Character{47}}%
\TextAt{13}{136.421}{123.362}{P}%
\TextAt{21}{143.967}{121.569}{K}%
\TextAt{8}{155.842}{123.362}{of}%
\TextAt{8}{168.946}{123.362}{an}%
\TextAt{8}{184.98}{123.362}{algebraic}%
\TextAt{8}{234.189}{123.362}{number}%
\TextAt{8}{276.56}{123.362}{field}%
\TextAt{13}{301.692}{123.362}{K}%
\TextAt{8}{316.182}{123.362}{is}%
\TextAt{8}{327.721}{123.362}{finite.}%
\TextAt{8}{362.108}{123.362}{In}%
\TextAt{8}{376.504}{123.362}{order}%
\TextAt{8}{406.848}{123.362}{to}%
\TextAt{8}{420.933}{123.362}{count}%
\TextAt{8}{452.903}{123.362}{the}%
\TextAt{8}{472.835}{123.362}{ideals}%
\TextAt{14}{505.187}{123.362}{a}%
\NativeMathAt{514.491}{123.362}{11.955}{\neq}%
\TextAt{9}{526.9}{123.362}{0}%
\TextAt{8}{79.226}{106.316}{of}%
\TextAt{8}{92.556}{106.316}{the}%
\TextAt{8}{112.713}{106.316}{ring}%
\TextAt{14}{136.778}{106.316}{o}%
\TextAt{21}{142.623}{104.523}{K}%
\TextAt{8}{154.723}{106.316}{we}%
\TextAt{8}{171.95}{106.316}{consider}%
\TextAt{8}{217.535}{106.316}{their}%
\TextAt{12}{245.498}{106.316}{absolute}%
\TextAt{12}{298.615}{106.316}{norm}%
\TextAt{14}{264.467}{77.315}{N}%
\TextAt{9}{274.415}{77.315}{\Character{40}}%
\TextAt{14}{278.967}{77.315}{a}%
\TextAt{9}{284.95}{77.315}{\Character{41}}%
\TextAt{9}{292.805}{77.315}{\Character{61}}%
\TextAt{9}{305.226}{77.315}{\Character{40}}%
\TextAt{14}{309.801}{77.315}{o}%
\TextAt{21}{315.646}{75.522}{K}%
\TextAt{9}{327.169}{77.315}{\Character{58}}%
\TextAt{14}{333.742}{77.315}{a}%
\TextAt{9}{339.725}{77.315}{\Character{41}}%
\TextAt{13}{344.277}{77.315}{\Character{58}}%
\TextAt{8}{300.14}{43.033}{32}%
\Destination{78.23}{809.815}{nav.page.32}%
\Destination{79.23}{188.555}{nav.section.1.6}%
\Destination{79.226}{592.555}{nav.ex.1.5.1}%
\Destination{79.226}{502.555}{nav.ex.1.5.2}%
\Destination{79.226}{397.355}{nav.ex.1.5.3}%
\Destination{79.36}{776.155}{nav.print.34}%
\Destination{79.23}{188.555}{outline.11}%
\Hotlink{298.125}{413.155}{14.294}{11.047}{nav.chapter.3}%
\Hotlink{199.239}{293.455}{6.858}{11.047}{nav.ex.1.5.2}%
\Hotlink{318.58}{413.155}{30.936}{11.047}{nav.number.3.2.15}%
\Hotlink{296.134}{293.455}{25.078}{11.047}{nav.number.1.5.3}%
\Destination{79.226}{505.547}{errata.70}%
}

% Page 45 (body page 33)
\Page{%
\TextAt{2}{79.226}{796.943}{I.}%
\TextAt{2}{89.83}{796.943}{Algebraic}%
\TextAt{2}{138.648}{796.943}{Integers}%
\Vector{q 0 G 0 g 0.299 w 7.226 18.729 m 460.769 18.729 l S Q}%
\TextAt{8}{79.226}{761.124}{(Throughout}%
\TextAt{8}{148.615}{761.124}{this}%
\TextAt{8}{171.223}{761.124}{book}%
\TextAt{8}{199.952}{761.124}{the}%
\TextAt{8}{219.895}{761.124}{case}%
\TextAt{8}{244.452}{761.124}{of}%
\TextAt{8}{257.567}{761.124}{the}%
\TextAt{8}{277.51}{761.124}{zero}%
\TextAt{8}{302.007}{761.124}{ideal}%
\TextAt{14}{329.756}{761.124}{a}%
\TextAt{9}{339.06}{761.124}{\Character{61}}%
\TextAt{9}{351.47}{761.124}{0}%
\TextAt{8}{361.021}{761.124}{is}%
\TextAt{8}{372.571}{761.124}{often}%
\TextAt{8}{401.946}{761.124}{tacitly}%
\TextAt{8}{438.482}{761.124}{excluded,}%
\TextAt{8}{490.005}{761.124}{when}%
\TextAt{8}{520.348}{761.124}{its}%
\TextAt{8}{79.226}{744.078}{consideration}%
\TextAt{8}{150.91}{744.078}{would}%
\TextAt{8}{185.126}{744.078}{visibly}%
\TextAt{8}{222.331}{744.078}{make}%
\TextAt{8}{252.973}{744.078}{no}%
\TextAt{8}{269.304}{744.078}{sense.)}%
\TextAt{8}{308.659}{744.078}{This}%
\TextAt{8}{335.451}{744.078}{index}%
\TextAt{8}{367.061}{744.078}{is}%
\TextAt{8}{378.897}{744.078}{finite}%
\TextAt{8}{408.881}{744.078}{by}%
\TextAt{8}{425.213}{744.078}{(2.12),}%
\TextAt{8}{462.388}{744.078}{and}%
\TextAt{8}{485.223}{744.078}{the}%
\TextAt{8}{505.452}{744.078}{name}%
\TextAt{8}{79.226}{727.032}{is}%
\TextAt{8}{90.99}{727.032}{justified}%
\TextAt{8}{135.595}{727.032}{by}%
\TextAt{8}{151.854}{727.032}{the}%
\TextAt{8}{172.01}{727.032}{special}%
\TextAt{8}{210.111}{727.032}{case}%
\TextAt{8}{234.882}{727.032}{of}%
\TextAt{8}{248.212}{727.032}{a}%
\TextAt{8}{257.968}{727.032}{principal}%
\TextAt{8}{306.745}{727.032}{ideal}%
\TextAt{9}{334.708}{727.032}{\Character{40}}%
\TextAt{13}{339.26}{727.032}{a}%
\TextAt{9}{345.405}{727.032}{\Character{41}}%
\TextAt{8}{353.855}{727.032}{of}%
\TextAt{14}{367.185}{727.032}{o}%
\TextAt{21}{373.029}{725.238}{K}%
\TextAt{8}{381.232}{727.032}{,}%
\TextAt{8}{388.381}{727.032}{where}%
\TextAt{8}{422.19}{727.032}{we}%
\TextAt{8}{439.418}{727.032}{have}%
\TextAt{8}{466.413}{727.032}{the}%
\TextAt{8}{486.569}{727.032}{identity}%
\TextAt{14}{255.338}{700.053}{N}%
\TextAt{9}{265.286}{700.053}{\Character{40}\Character{40}}%
\TextAt{13}{274.391}{700.053}{a}%
\TextAt{9}{280.536}{700.053}{\Character{41}\Character{41}}%
\TextAt{9}{292.934}{700.053}{\Character{61}}%
\TextAt{16}{305.387}{700.053}{\Character{106}}%
\TextAt{13}{308.708}{700.053}{N}%
\TextAt{21}{318.081}{698.194}{K}%
\TextAt{3}{325.786}{698.194}{\Character{106}}%
\TextAt{25}{328.138}{698.194}{Q}%
\TextAt{9}{334.835}{700.053}{\Character{40}}%
\TextAt{13}{339.388}{700.053}{a}%
\TextAt{9}{345.533}{700.053}{\Character{41}}%
\TextAt{16}{350.085}{700.053}{\Character{106}}%
\TextAt{13}{353.406}{700.053}{\Character{58}}%
\TextAt{8}{79.226}{673.075}{Indeed,}%
\TextAt{8}{120.114}{673.075}{if}%
\TextAt{13}{130.345}{673.075}{\Character{33}}%
\TextAt{17}{137.635}{671.282}{1}%
\TextAt{13}{142.367}{673.075}{\Character{59}}%
\TextAt{13}{147.591}{673.075}{\Character{58}}%
\TextAt{13}{152.828}{673.075}{\Character{58}}%
\TextAt{13}{158.052}{673.075}{\Character{58}}%
\TextAt{13}{163.289}{673.075}{\Character{59}}%
\TextAt{13}{168.525}{673.075}{\Character{33}}%
\TextAt{21}{175.877}{671.282}{n}%
\TextAt{8}{184.919}{673.075}{is}%
\TextAt{8}{196.19}{673.075}{a}%
\TextAt{18}{205.454}{673.075}{Z}%
\TextAt{8}{213.424}{673.075}{-basis}%
\TextAt{8}{245.569}{673.075}{of}%
\TextAt{14}{258.407}{673.075}{o}%
\TextAt{21}{264.251}{671.282}{K}%
\TextAt{8}{272.454}{673.075}{,}%
\TextAt{8}{279.21}{673.075}{then}%
\TextAt{13}{305.378}{673.075}{a\Character{33}}%
\TextAt{17}{318.812}{671.282}{1}%
\TextAt{13}{323.544}{673.075}{\Character{59}}%
\TextAt{13}{328.768}{673.075}{\Character{58}}%
\TextAt{13}{334.005}{673.075}{\Character{58}}%
\TextAt{13}{339.241}{673.075}{\Character{58}}%
\TextAt{13}{344.466}{673.075}{\Character{59}}%
\TextAt{13}{349.702}{673.075}{a\Character{33}}%
\TextAt{21}{363.199}{671.282}{n}%
\TextAt{8}{372.241}{673.075}{is}%
\TextAt{8}{383.512}{673.075}{a}%
\TextAt{18}{392.776}{673.075}{Z}%
\TextAt{8}{400.746}{673.075}{-basis}%
\TextAt{8}{432.891}{673.075}{of}%
\TextAt{9}{445.729}{673.075}{\Character{40}}%
\TextAt{13}{450.282}{673.075}{a}%
\TextAt{9}{456.427}{673.075}{\Character{41}}%
\TextAt{9}{464.282}{673.075}{\Character{61}}%
\TextAt{13}{476.725}{673.075}{a}%
\TextAt{14}{482.87}{673.075}{o}%
\TextAt{21}{488.714}{671.282}{K}%
\TextAt{8}{496.917}{673.075}{,}%
\TextAt{8}{503.673}{673.075}{and}%
\TextAt{8}{525.943}{673.075}{if}%
\TextAt{13}{79.226}{656.029}{A}%
\TextAt{9}{91.322}{656.029}{\Character{61}}%
\TextAt{9}{103.732}{656.029}{\Character{40}}%
\TextAt{13}{108.3}{656.029}{a}%
\TextAt{21}{114.445}{654.236}{ij}%
\TextAt{9}{121.71}{656.029}{\Character{41}}%
\TextAt{8}{129.617}{656.029}{denotes}%
\TextAt{8}{171.406}{656.029}{the}%
\TextAt{8}{191.019}{656.029}{transition}%
\TextAt{8}{243.879}{656.029}{matrix,}%
\TextAt{13}{284.75}{656.029}{a\Character{33}}%
\TextAt{21}{298.184}{654.236}{i}%
\TextAt{9}{304.886}{656.029}{\Character{61}}%
\TextAt{20}{317.312}{664.995}{\Character{80}}%
\TextAt{21}{329.931}{652.542}{j}%
\TextAt{13}{336.306}{656.029}{a}%
\TextAt{21}{342.451}{654.236}{ij}%
\TextAt{13}{349.716}{656.029}{\Character{33}}%
\TextAt{21}{357.005}{654.236}{j}%
\TextAt{8}{361.387}{656.029}{,}%
\TextAt{8}{368.102}{656.029}{then,}%
\TextAt{8}{397.579}{656.029}{as}%
\TextAt{8}{411.405}{656.029}{was}%
\TextAt{8}{433.361}{656.029}{pointed}%
\TextAt{8}{475.091}{656.029}{out}%
\TextAt{8}{495.362}{656.029}{already}%
\TextAt{8}{79.226}{638.983}{in}%
\TextAt{8}{92.879}{638.983}{§2,}%
\TextAt{8}{111.565}{638.983}{one}%
\TextAt{8}{133.024}{638.983}{has}%
\TextAt{16}{153.898}{638.983}{\Character{106}}%
\TextAt{8}{159.212}{638.983}{det}%
\TextAt{9}{175.471}{638.983}{\Character{40}}%
\TextAt{13}{180.023}{638.983}{A}%
\TextAt{9}{188.798}{638.983}{\Character{41}}%
\TextAt{16}{193.351}{638.983}{\Character{106}}%
\TextAt{9}{199.992}{638.983}{\Character{61}}%
\TextAt{9}{212.402}{638.983}{\Character{40}}%
\TextAt{14}{216.97}{638.983}{o}%
\TextAt{21}{222.814}{637.19}{K}%
\TextAt{9}{234.338}{638.983}{\Character{58}}%
\TextAt{9}{240.889}{638.983}{\Character{40}}%
\TextAt{13}{245.463}{638.983}{a}%
\TextAt{9}{251.608}{638.983}{\Character{41}\Character{41}}%
\TextAt{8}{264.61}{638.983}{as}%
\TextAt{8}{278.98}{638.983}{well}%
\TextAt{8}{302.711}{638.983}{as}%
\TextAt{8}{317.081}{638.983}{det}%
\TextAt{9}{333.34}{638.983}{\Character{40}}%
\TextAt{13}{337.892}{638.983}{A}%
\TextAt{9}{346.668}{638.983}{\Character{41}}%
\TextAt{9}{354.523}{638.983}{\Character{61}}%
\TextAt{13}{366.966}{638.983}{N}%
\TextAt{21}{376.34}{637.123}{K}%
\TextAt{3}{384.045}{637.123}{\Character{106}}%
\TextAt{25}{386.397}{637.123}{Q}%
\TextAt{9}{393.094}{638.983}{\Character{40}}%
\TextAt{13}{397.646}{638.983}{a}%
\TextAt{9}{403.791}{638.983}{\Character{41}}%
\TextAt{8}{412.241}{638.983}{by}%
\TextAt{8}{428.5}{638.983}{definition.}%
\TextAt{12}{79.226}{613.548}{Proposition}%
\TextAt{12}{152.689}{613.548}{(6.1).}%
\TextAt{8}{188.659}{613.548}{If}%
\TextAt{14}{201.004}{613.548}{a}%
\TextAt{9}{211.422}{613.548}{\Character{61}}%
\TextAt{14}{224.961}{613.548}{p}%
\TextAt{21}{230.945}{618.773}{\Character{23}}%
\TextAt{24}{235.107}{617.666}{1}%
\TextAt{17}{230.945}{610.617}{1}%
\TextAt{16}{241.749}{613.548}{\Character{1}}%
\TextAt{16}{247.045}{613.548}{\Character{1}}%
\TextAt{16}{252.353}{613.548}{\Character{1}}%
\TextAt{14}{257.689}{613.548}{p}%
\TextAt{21}{263.672}{617.887}{\Character{23}}%
\TextAt{22}{267.835}{616.891}{r}%
\TextAt{21}{263.672}{610.593}{r}%
\TextAt{8}{277.042}{613.548}{is}%
\TextAt{8}{289.459}{613.548}{the}%
\TextAt{8}{310.27}{613.548}{prime}%
\TextAt{8}{344.087}{613.548}{factorization}%
\TextAt{8}{412.717}{613.548}{of}%
\TextAt{8}{426.701}{613.548}{an}%
\TextAt{8}{443.614}{613.548}{ideal}%
\TextAt{14}{472.231}{613.548}{a}%
\NativeMathAt{482.648}{613.548}{11.955}{\neq}%
\TextAt{9}{496.169}{613.548}{0}%
\TextAt{8}{502.041}{613.548}{,}%
\TextAt{8}{510.007}{613.548}{then}%
\TextAt{8}{79.226}{596.502}{one}%
\TextAt{8}{100.686}{596.502}{has}%
\TextAt{14}{236.52}{579.456}{N}%
\TextAt{9}{246.468}{579.456}{\Character{40}}%
\TextAt{14}{251.02}{579.456}{a}%
\TextAt{9}{257.003}{579.456}{\Character{41}}%
\TextAt{9}{264.858}{579.456}{\Character{61}}%
\TextAt{14}{277.302}{579.456}{N}%
\TextAt{9}{287.25}{579.456}{\Character{40}}%
\TextAt{14}{291.802}{579.456}{p}%
\TextAt{17}{297.786}{577.663}{1}%
\TextAt{9}{302.518}{579.456}{\Character{41}}%
\TextAt{21}{307.07}{584.393}{\Character{23}}%
\TextAt{24}{311.233}{583.286}{1}%
\TextAt{16}{317.875}{579.456}{\Character{1}}%
\TextAt{16}{323.171}{579.456}{\Character{1}}%
\TextAt{16}{328.479}{579.456}{\Character{1}}%
\TextAt{14}{333.815}{579.456}{N}%
\TextAt{9}{343.763}{579.456}{\Character{40}}%
\TextAt{14}{348.315}{579.456}{p}%
\TextAt{21}{354.298}{577.663}{r}%
\TextAt{9}{358.853}{579.456}{\Character{41}}%
\TextAt{21}{363.405}{584.393}{\Character{23}}%
\TextAt{22}{367.568}{583.396}{r}%
\TextAt{13}{372.224}{579.456}{\Character{58}}%
\TextAt{19}{79.226}{540.025}{Proof.}%
\TextAt{8}{116}{540.025}{By}%
\TextAt{8}{134.363}{540.025}{the}%
\TextAt{8}{154.52}{540.025}{Chinese}%
\TextAt{8}{198.144}{540.025}{remainder}%
\TextAt{8}{253.425}{540.025}{theorem}%
\TextAt{8}{298.95}{540.025}{(3.6),}%
\TextAt{8}{330.177}{540.025}{one}%
\TextAt{8}{351.637}{540.025}{has}%
\TextAt{14}{227.176}{513.046}{o}%
\TextAt{21}{233.021}{511.253}{K}%
\TextAt{9}{241.223}{513.046}{\Character{47}}%
\TextAt{14}{247.076}{513.046}{a}%
\NativeMathAt{256.38}{513.046}{11.955}{\simeq}%
\TextAt{14}{269}{513.046}{o}%
\TextAt{21}{274.844}{511.253}{K}%
\TextAt{9}{283.047}{513.046}{\Character{47}}%
\TextAt{14}{288.9}{513.046}{p}%
\TextAt{21}{294.883}{518.271}{\Character{23}}%
\TextAt{24}{299.045}{517.164}{1}%
\TextAt{17}{294.883}{510.115}{1}%
\TextAt{16}{306.351}{513.046}{\Character{8}}%
\TextAt{16}{318.294}{513.046}{\Character{1}}%
\TextAt{16}{323.59}{513.046}{\Character{1}}%
\TextAt{16}{328.899}{513.046}{\Character{1}}%
\TextAt{16}{334.864}{513.046}{\Character{8}}%
\TextAt{14}{346.866}{513.046}{o}%
\TextAt{21}{352.71}{511.253}{K}%
\TextAt{9}{360.913}{513.046}{\Character{47}}%
\TextAt{14}{366.766}{513.046}{p}%
\TextAt{21}{372.749}{517.982}{\Character{23}}%
\TextAt{22}{376.911}{516.986}{r}%
\TextAt{21}{372.749}{510.091}{r}%
\TextAt{13}{381.568}{513.046}{\Character{58}}%
\TextAt{8}{79.226}{486.068}{We}%
\TextAt{8}{99.371}{486.068}{are}%
\TextAt{8}{118.881}{486.068}{thus}%
\TextAt{8}{144.633}{486.068}{reduced}%
\TextAt{8}{188.198}{486.068}{to}%
\TextAt{8}{202.508}{486.068}{considering}%
\TextAt{8}{263.706}{486.068}{the}%
\TextAt{8}{283.863}{486.068}{case}%
\TextAt{8}{308.634}{486.068}{where}%
\TextAt{14}{342.443}{486.068}{a}%
\TextAt{8}{352.324}{486.068}{is}%
\TextAt{8}{364.087}{486.068}{a}%
\TextAt{8}{373.843}{486.068}{prime}%
\TextAt{8}{407.007}{486.068}{power}%
\TextAt{14}{441.151}{486.068}{p}%
\TextAt{21}{447.134}{490.406}{\Character{23}}%
\TextAt{8}{452.316}{486.068}{.}%
\TextAt{8}{460.764}{486.068}{In}%
\TextAt{8}{475.385}{486.068}{the}%
\TextAt{8}{495.542}{486.068}{chain}%
\TextAt{14}{261.182}{459.089}{p}%
\NativeMathAt{270.486}{459.089}{11.955}{\supseteq}%
\TextAt{14}{283.105}{459.089}{p}%
\TextAt{17}{289.088}{464.025}{2}%
\NativeMathAt{297.141}{459.089}{11.955}{\supseteq}%
\TextAt{16}{309.742}{459.089}{\Character{1}}%
\TextAt{16}{315.05}{459.089}{\Character{1}}%
\TextAt{16}{320.358}{459.089}{\Character{1}}%
\NativeMathAt{326.981}{459.089}{11.955}{\supseteq}%
\TextAt{14}{339.649}{459.089}{p}%
\TextAt{21}{345.632}{464.025}{\Character{23}}%
\TextAt{8}{79.226}{432.111}{one}%
\TextAt{8}{100.706}{432.111}{has}%
\TextAt{14}{121.601}{432.111}{p}%
\TextAt{21}{127.584}{436.449}{i}%
\NativeMathAt{134.322}{432.111}{11.955}{\neq}%
\TextAt{14}{146.782}{432.111}{p}%
\TextAt{21}{152.766}{436.449}{i}%
\TextAt{17}{155.649}{436.449}{\Character{43}1}%
\TextAt{8}{170.886}{432.111}{because}%
\TextAt{8}{214.208}{432.111}{of}%
\TextAt{8}{227.559}{432.111}{the}%
\TextAt{8}{247.736}{432.111}{unique}%
\TextAt{8}{285.798}{432.111}{prime}%
\TextAt{8}{318.983}{432.111}{factorization,}%
\TextAt{8}{390.237}{432.111}{and}%
\TextAt{8}{413.021}{432.111}{each}%
\TextAt{8}{439.379}{432.111}{quotient}%
\TextAt{14}{485.582}{432.111}{p}%
\TextAt{21}{491.566}{436.449}{i}%
\TextAt{9}{494.947}{432.111}{\Character{47}}%
\TextAt{14}{500.8}{432.111}{p}%
\TextAt{21}{506.783}{436.449}{i}%
\TextAt{17}{509.666}{436.449}{\Character{43}1}%
\TextAt{8}{524.903}{432.111}{is}%
\TextAt{8}{79.226}{415.065}{an}%
\TextAt{14}{96.173}{415.065}{o}%
\TextAt{21}{102.017}{413.272}{K}%
\TextAt{9}{110.22}{415.065}{\Character{47}}%
\TextAt{14}{116.073}{415.065}{p}%
\TextAt{8}{122.056}{415.065}{-vector}%
\TextAt{8}{161.766}{415.065}{space}%
\TextAt{8}{193.729}{415.065}{of}%
\TextAt{8}{207.747}{415.065}{dimension}%
\TextAt{8}{263.775}{415.065}{1.}%
\TextAt{8}{280.146}{415.065}{In}%
\TextAt{8}{295.455}{415.065}{fact,}%
\TextAt{8}{322.652}{415.065}{if}%
\TextAt{13}{334.064}{415.065}{a}%
\NativeMathAt{344.703}{415.065}{11.955}{\in}%
\TextAt{14}{357.166}{415.065}{p}%
\TextAt{21}{363.149}{419.403}{i}%
\TextAt{16}{369.656}{415.065}{\Character{110}}%
\TextAt{14}{378.76}{415.065}{p}%
\TextAt{21}{384.743}{419.403}{i}%
\TextAt{17}{387.626}{419.403}{\Character{43}1}%
\TextAt{8}{403.53}{415.065}{and}%
\TextAt{14}{426.981}{415.065}{b}%
\TextAt{9}{437.609}{415.065}{\Character{61}}%
\TextAt{9}{451.19}{415.065}{\Character{40}}%
\TextAt{13}{455.759}{415.065}{a}%
\TextAt{9}{461.904}{415.065}{\Character{41}}%
\TextAt{9}{469.567}{415.065}{\Character{43}}%
\TextAt{14}{481.813}{415.065}{p}%
\TextAt{21}{487.796}{419.403}{i}%
\TextAt{17}{490.679}{419.403}{\Character{43}1}%
\TextAt{8}{501.998}{415.065}{,}%
\TextAt{8}{510.007}{415.065}{then}%
\TextAt{14}{79.226}{398.019}{p}%
\TextAt{21}{85.209}{402.357}{i}%
\NativeMathAt{92.784}{398.019}{11.955}{\supseteq}%
\TextAt{14}{106.275}{398.019}{b}%
\TextAt{18}{116.603}{398.019}{\Character{41}}%
\TextAt{14}{130.095}{398.019}{p}%
\TextAt{21}{136.078}{402.357}{i}%
\TextAt{17}{138.961}{402.357}{\Character{43}1}%
\TextAt{8}{154.689}{398.019}{and}%
\TextAt{8}{177.963}{398.019}{consequently}%
\TextAt{14}{247.803}{398.019}{p}%
\TextAt{21}{253.786}{402.357}{i}%
\TextAt{9}{261.361}{398.019}{\Character{61}}%
\TextAt{14}{274.658}{398.019}{b}%
\TextAt{8}{280.793}{398.019}{,}%
\TextAt{8}{288.582}{398.019}{because}%
\TextAt{8}{332.396}{398.019}{otherwise}%
\TextAt{14}{384.996}{398.019}{b}%
\TextAt{3}{391.131}{402.357}{\Character{48}}%
\TextAt{9}{398.119}{398.019}{\Character{61}}%
\TextAt{14}{411.417}{398.019}{bp}%
\TextAt{3}{423.535}{402.357}{\Character{0}}%
\TextAt{21}{430.121}{402.357}{i}%
\TextAt{8}{437.912}{398.019}{would}%
\TextAt{8}{472.568}{398.019}{be}%
\TextAt{8}{489.004}{398.019}{a}%
\TextAt{8}{499.271}{398.019}{proper}%
\TextAt{8}{79.226}{380.973}{divisor}%
\TextAt{8}{117.66}{380.973}{of}%
\TextAt{14}{131.311}{380.973}{p}%
\TextAt{9}{141.162}{380.973}{\Character{61}}%
\TextAt{14}{154.134}{380.973}{p}%
\TextAt{21}{160.117}{385.311}{i}%
\TextAt{17}{163.001}{385.311}{\Character{43}1}%
\TextAt{14}{174.319}{380.973}{p}%
\TextAt{3}{180.303}{385.311}{\Character{0}}%
\TextAt{21}{186.889}{385.311}{i}%
\TextAt{8}{190.27}{380.973}{.}%
\TextAt{8}{199.681}{380.973}{Thus}%
\TextAt{9}{229.797}{380.973}{\Character{22}}%
\TextAt{13}{229.651}{380.973}{a}%
\TextAt{9}{239.664}{380.973}{\Character{61}}%
\TextAt{13}{252.636}{380.973}{a}%
\TextAt{8}{262.102}{380.973}{mod}%
\TextAt{14}{287.54}{380.973}{p}%
\TextAt{21}{293.523}{385.311}{i}%
\TextAt{17}{296.406}{385.311}{\Character{43}1}%
\TextAt{8}{311.943}{380.973}{is}%
\TextAt{8}{324.028}{380.973}{a}%
\TextAt{8}{334.105}{380.973}{basis}%
\TextAt{8}{363.166}{380.973}{of}%
\TextAt{8}{376.817}{380.973}{the}%
\TextAt{14}{397.294}{380.973}{o}%
\TextAt{21}{403.138}{379.18}{K}%
\TextAt{9}{411.341}{380.973}{\Character{47}}%
\TextAt{14}{417.194}{380.973}{p}%
\TextAt{8}{423.177}{380.973}{-vector}%
\TextAt{8}{462.52}{380.973}{space}%
\TextAt{14}{494.115}{380.973}{p}%
\TextAt{21}{500.098}{385.311}{i}%
\TextAt{9}{503.48}{380.973}{\Character{47}}%
\TextAt{14}{509.333}{380.973}{p}%
\TextAt{21}{515.316}{385.311}{i}%
\TextAt{17}{518.199}{385.311}{\Character{43}1}%
\TextAt{8}{529.518}{380.973}{.}%
\TextAt{8}{79.226}{363.927}{So}%
\TextAt{8}{95.485}{363.927}{we}%
\TextAt{8}{112.713}{363.927}{have}%
\TextAt{14}{139.707}{363.927}{p}%
\TextAt{21}{145.69}{368.265}{i}%
\TextAt{9}{149.072}{363.927}{\Character{47}}%
\TextAt{14}{154.925}{363.927}{p}%
\TextAt{21}{160.908}{368.265}{i}%
\TextAt{17}{163.791}{368.265}{\Character{43}1}%
\NativeMathAt{178.431}{363.927}{11.955}{\simeq}%
\TextAt{14}{191.05}{363.927}{o}%
\TextAt{21}{196.894}{362.134}{K}%
\TextAt{9}{205.097}{363.927}{\Character{47}}%
\TextAt{14}{210.95}{363.927}{p}%
\TextAt{8}{220.83}{363.927}{and}%
\TextAt{8}{243.593}{363.927}{therefore}%
\TextAt{14}{156.027}{336.948}{N}%
\TextAt{9}{165.975}{336.948}{\Character{40}}%
\TextAt{14}{170.527}{336.948}{p}%
\TextAt{21}{176.51}{341.885}{\Character{23}}%
\TextAt{9}{181.692}{336.948}{\Character{41}}%
\TextAt{9}{189.547}{336.948}{\Character{61}}%
\TextAt{9}{201.968}{336.948}{\Character{40}}%
\TextAt{14}{206.543}{336.948}{o}%
\TextAt{21}{212.387}{335.155}{K}%
\TextAt{9}{223.911}{336.948}{\Character{58}}%
\TextAt{14}{230.483}{336.948}{p}%
\TextAt{21}{236.467}{341.885}{\Character{23}}%
\TextAt{9}{241.649}{336.948}{\Character{41}}%
\TextAt{9}{249.504}{336.948}{\Character{61}}%
\TextAt{9}{261.925}{336.948}{\Character{40}}%
\TextAt{14}{266.5}{336.948}{o}%
\TextAt{21}{272.344}{335.155}{K}%
\TextAt{9}{283.867}{336.948}{\Character{58}}%
\TextAt{14}{290.44}{336.948}{p}%
\TextAt{9}{296.423}{336.948}{\Character{41}\Character{40}}%
\TextAt{14}{305.527}{336.948}{p}%
\TextAt{9}{314.831}{336.948}{\Character{58}}%
\TextAt{14}{321.404}{336.948}{p}%
\TextAt{17}{327.387}{341.885}{2}%
\TextAt{9}{332.119}{336.948}{\Character{41}}%
\TextAt{16}{338.664}{336.948}{\Character{1}}%
\TextAt{16}{343.96}{336.948}{\Character{1}}%
\TextAt{16}{349.268}{336.948}{\Character{1}}%
\TextAt{9}{354.604}{336.948}{\Character{40}}%
\TextAt{14}{359.157}{336.948}{p}%
\TextAt{21}{365.14}{341.885}{\Character{23}}%
\TextAt{3}{369.824}{341.885}{\Character{0}}%
\TextAt{17}{376.41}{341.885}{1}%
\TextAt{9}{384.463}{336.948}{\Character{58}}%
\TextAt{14}{391.036}{336.948}{p}%
\TextAt{21}{397.019}{341.885}{\Character{23}}%
\TextAt{9}{402.201}{336.948}{\Character{41}}%
\TextAt{9}{410.056}{336.948}{\Character{61}}%
\TextAt{14}{422.499}{336.948}{N}%
\TextAt{9}{432.448}{336.948}{\Character{40}}%
\TextAt{14}{437}{336.948}{p}%
\TextAt{9}{442.983}{336.948}{\Character{41}}%
\TextAt{21}{447.535}{341.885}{\Character{23}}%
\TextAt{13}{452.717}{336.948}{\Character{58}}%
\Vector{q 0 G 0 g 0.399 w 452.334 -464.03 m 452.334 -455.96 l S Q}%
\Vector{q 0 G 0 g 0.399 w 452.534 -456.159 m 459.707 -456.159 l S Q}%
\Vector{q 0 G 0 g 0.399 w 452.534 -463.831 m 459.707 -463.831 l S Q}%
\Vector{q 0 G 0 g 0.399 w 459.906 -464.03 m 459.906 -455.96 l S Q}%
\TextAt{8}{97.159}{281.547}{The}%
\TextAt{8}{121.213}{281.547}{proposition}%
\TextAt{8}{182.746}{281.547}{immediately}%
\TextAt{8}{249.408}{281.547}{implies}%
\TextAt{8}{289.135}{281.547}{the}%
\TextAt{8}{309.291}{281.547}{multiplicativity}%
\TextAt{14}{257.428}{254.568}{N}%
\TextAt{9}{267.376}{254.568}{\Character{40}}%
\TextAt{14}{271.928}{254.568}{ab}%
\TextAt{9}{284.046}{254.568}{\Character{41}}%
\TextAt{9}{291.901}{254.568}{\Character{61}}%
\TextAt{14}{304.345}{254.568}{N}%
\TextAt{9}{314.293}{254.568}{\Character{40}}%
\TextAt{14}{318.845}{254.568}{a}%
\TextAt{9}{324.828}{254.568}{\Character{41}}%
\TextAt{14}{329.38}{254.568}{N}%
\TextAt{9}{339.329}{254.568}{\Character{40}}%
\TextAt{14}{343.881}{254.568}{b}%
\TextAt{9}{350.016}{254.568}{\Character{41}}%
\TextAt{8}{79.226}{227.59}{of}%
\TextAt{8}{92.556}{227.59}{the}%
\TextAt{8}{112.713}{227.59}{absolute}%
\TextAt{8}{158.955}{227.59}{norm.}%
\TextAt{8}{194.075}{227.59}{It}%
\TextAt{8}{206.748}{227.59}{may}%
\TextAt{8}{232.117}{227.59}{therefore}%
\TextAt{8}{281.217}{227.59}{be}%
\TextAt{8}{297.141}{227.59}{extended}%
\TextAt{8}{346.886}{227.59}{to}%
\TextAt{8}{361.197}{227.59}{a}%
\TextAt{8}{370.952}{227.59}{homomorphism}%
\TextAt{14}{267.93}{200.611}{N}%
\TextAt{9}{281.199}{200.611}{\Character{58}}%
\TextAt{13}{287.771}{200.611}{J}%
\TextAt{21}{294.241}{198.818}{K}%
\NativeArrowAt{305.765}{200.611}{11.955}{19.26}{\longrightarrow}%
\TextAt{18}{328.347}{200.611}{R}%
\TextAt{3}{336.981}{205.548}{\Character{3}}%
\TextAt{17}{336.981}{197.656}{\Character{43}}%
\TextAt{8}{79.226}{173.633}{defined}%
\TextAt{8}{119.869}{173.633}{on}%
\TextAt{8}{136.459}{173.633}{all}%
\TextAt{8}{153.048}{173.633}{fractional}%
\TextAt{8}{205.742}{173.633}{ideals}%
\TextAt{14}{238.651}{173.633}{a}%
\TextAt{9}{248.518}{173.633}{\Character{61}}%
\TextAt{20}{261.506}{182.599}{\Character{81}}%
\TextAt{27}{272.797}{170.146}{p}%
\TextAt{14}{279.277}{173.633}{p}%
\TextAt{21}{285.26}{177.971}{\Character{23}}%
\TextAt{28}{289.422}{176.975}{p}%
\TextAt{8}{294.003}{173.633}{,}%
\TextAt{13}{301.566}{173.633}{\Character{23}}%
\TextAt{27}{307.349}{171.84}{p}%
\NativeMathAt{315.72}{173.633}{11.955}{\in}%
\TextAt{18}{327.574}{173.633}{Z}%
\TextAt{8}{335.544}{173.633}{.}%
\TextAt{8}{344.984}{173.633}{The}%
\TextAt{8}{369.368}{173.633}{following}%
\TextAt{8}{419.133}{173.633}{lemma,}%
\TextAt{8}{460.517}{173.633}{a}%
\TextAt{8}{470.603}{173.633}{consequence}%
\TextAt{8}{79.226}{156.587}{of}%
\TextAt{8}{92.556}{156.587}{(5.3),}%
\TextAt{8}{123.783}{156.587}{is}%
\TextAt{8}{135.547}{156.587}{crucial}%
\TextAt{8}{173.266}{156.587}{for}%
\TextAt{8}{191.15}{156.587}{the}%
\TextAt{8}{211.307}{156.587}{finiteness}%
\TextAt{8}{262.152}{156.587}{of}%
\TextAt{8}{275.482}{156.587}{the}%
\TextAt{8}{295.639}{156.587}{ideal}%
\TextAt{8}{323.602}{156.587}{class}%
\TextAt{8}{351.039}{156.587}{group.}%
\TextAt{12}{79.226}{131.152}{Lemma}%
\TextAt{12}{126.891}{131.152}{(6.2).}%
\TextAt{8}{162.207}{131.152}{In}%
\TextAt{8}{176.828}{131.152}{every}%
\TextAt{8}{207.72}{131.152}{ideal}%
\TextAt{14}{235.683}{131.152}{a}%
\NativeMathAt{244.987}{131.152}{11.955}{\neq}%
\TextAt{9}{257.397}{131.152}{0}%
\TextAt{8}{267.163}{131.152}{of}%
\TextAt{14}{280.493}{131.152}{o}%
\TextAt{21}{286.337}{129.359}{K}%
\TextAt{8}{298.438}{131.152}{there}%
\TextAt{8}{328.349}{131.152}{exists}%
\TextAt{8}{360.664}{131.152}{an}%
\TextAt{13}{376.923}{131.152}{a}%
\NativeMathAt{386.389}{131.152}{11.955}{\in}%
\TextAt{14}{397.68}{131.152}{a}%
\TextAt{8}{403.663}{131.152}{,}%
\TextAt{13}{410.812}{131.152}{a}%
\NativeMathAt{420.278}{131.152}{11.955}{\neq}%
\TextAt{9}{432.688}{131.152}{0}%
\TextAt{8}{438.557}{131.152}{,}%
\TextAt{8}{445.706}{131.152}{such}%
\TextAt{8}{472.103}{131.152}{that}%
\TextAt{16}{225.703}{99.611}{\Character{106}}%
\TextAt{13}{229.024}{99.611}{N}%
\TextAt{21}{238.398}{97.751}{K}%
\TextAt{3}{246.102}{97.751}{\Character{106}}%
\TextAt{25}{248.455}{97.751}{Q}%
\TextAt{9}{255.152}{99.611}{\Character{40}}%
\TextAt{13}{259.704}{99.611}{a}%
\TextAt{9}{265.849}{99.611}{\Character{41}}%
\TextAt{16}{270.401}{99.611}{\Character{106}}%
\NativeMathAt{277.024}{99.611}{11.955}{\leq}%
\TextAt{20}{289.663}{116.468}{\Character{18}}%
\TextAt{9}{300.267}{107.699}{2}%
\Vector{q 0 G 0 g 0.478 w 227.658 -671.4 m 234.728 -671.4 l S Q}%
\TextAt{13}{299.658}{91.41}{\Character{25}}%
\TextAt{20}{307.923}{116.468}{\Character{19}}%
\TextAt{21}{316.724}{113.791}{s}%
\TextAt{20}{323.13}{110.416}{\Character{112}}%
\Vector{q 0 G 0 g 0.478 w 263.085 -663.344 m 284.012 -663.344 l S Q}%
\TextAt{16}{335.085}{99.611}{\Character{106}}%
\TextAt{13}{338.406}{99.611}{d}%
\TextAt{21}{344.489}{97.818}{K}%
\TextAt{16}{352.692}{99.611}{\Character{106}}%
\TextAt{14}{358.005}{99.611}{N}%
\TextAt{9}{367.953}{99.611}{\Character{40}}%
\TextAt{14}{372.505}{99.611}{a}%
\TextAt{9}{378.488}{99.611}{\Character{41}}%
\TextAt{13}{383.041}{99.611}{\Character{58}}%
\TextAt{8}{300.14}{43.033}{33}%
\Destination{78.23}{809.815}{nav.page.33}%
\Destination{79.36}{759.055}{nav.print.35}%
\Destination{79.36}{627.555}{nav.number.1.6.1}%
\Destination{79.36}{145.155}{nav.number.1.6.2}%
\Hotlink{98.058}{636.755}{6.858}{11.047}{nav.section.1.2}%
\Hotlink{424.713}{741.855}{30.936}{11.047}{nav.number.1.2.12}%
\Hotlink{92.056}{154.355}{25.078}{11.047}{nav.number.1.5.3}%
\Hotlink{300.118}{537.8}{21.722}{11.047}{nav.number.1.3.6}%
}

% Page 46 (body page 34)
\Page{%
\TextAt{2}{79.226}{796.943}{I.}%
\TextAt{2}{89.83}{796.943}{Algebraic}%
\TextAt{2}{138.648}{796.943}{Integers}%
\Vector{q 0 G 0 g 0.299 w 7.226 18.729 m 460.769 18.729 l S Q}%
\TextAt{19}{79.226}{761.124}{Proof.}%
\TextAt{8}{116}{761.124}{Given}%
\TextAt{13}{150.548}{761.124}{\Character{34}}%
\TextAt{13}{160.459}{761.124}{\Character{62}}%
\TextAt{9}{174.006}{761.124}{0}%
\TextAt{8}{179.859}{761.124}{,}%
\TextAt{8}{187.827}{761.124}{we}%
\TextAt{8}{205.71}{761.124}{choose}%
\TextAt{8}{243.498}{761.124}{positive}%
\TextAt{8}{287.467}{761.124}{real}%
\TextAt{8}{310.885}{761.124}{numbers}%
\TextAt{13}{358.751}{761.124}{c}%
\TextAt{21}{363.789}{759.33}{\Character{28}}%
\TextAt{8}{368.966}{761.124}{,}%
\TextAt{8}{376.934}{761.124}{for}%
\TextAt{13}{395.474}{761.124}{\Character{28}}%
\NativeMathAt{406.331}{761.124}{11.955}{\in}%
\TextAt{8}{418.738}{761.124}{Hom}%
\TextAt{9}{443.127}{761.124}{\Character{40}}%
\TextAt{13}{447.679}{761.124}{K}%
\TextAt{13}{457.841}{761.124}{\Character{59}}%
\TextAt{18}{463.091}{761.124}{C}%
\TextAt{9}{471.725}{761.124}{\Character{41}}%
\TextAt{8}{476.277}{761.124}{,}%
\TextAt{8}{484.246}{761.124}{such}%
\TextAt{8}{511.298}{761.124}{that}%
\TextAt{13}{79.226}{744.078}{c}%
\TextAt{21}{84.264}{742.284}{\Character{28}}%
\TextAt{9}{92.761}{744.078}{\Character{61}}%
\TextAt{13}{105.187}{744.078}{c}%
\TextAt{17}{110.682}{742.284}{\Character{22}}%
\TextAt{21}{110.225}{742.284}{\Character{28}}%
\TextAt{8}{119.299}{744.078}{and}%
\TextAt{20}{226.12}{735.034}{\Character{89}}%
\TextAt{21}{231.419}{709.928}{\Character{28}}%
\TextAt{13}{243.389}{723.676}{c}%
\TextAt{21}{248.427}{721.883}{\Character{28}}%
\TextAt{9}{256.924}{723.676}{\Character{61}}%
\TextAt{20}{269.35}{740.533}{\Character{18}}%
\TextAt{9}{279.954}{731.764}{2}%
\Vector{q 0 G 0 g 0.478 w 207.345 -47.335 m 214.415 -47.335 l S Q}%
\TextAt{13}{279.345}{715.476}{\Character{25}}%
\TextAt{20}{287.61}{740.533}{\Character{19}}%
\TextAt{21}{296.411}{737.857}{s}%
\TextAt{20}{302.817}{734.482}{\Character{112}}%
\Vector{q 0 G 0 g 0.478 w 242.772 -39.279 m 263.699 -39.279 l S Q}%
\TextAt{16}{314.772}{723.676}{\Character{106}}%
\TextAt{13}{318.093}{723.676}{d}%
\TextAt{21}{324.176}{721.883}{K}%
\TextAt{16}{332.379}{723.676}{\Character{106}}%
\TextAt{14}{337.692}{723.676}{N}%
\TextAt{9}{347.64}{723.676}{\Character{40}}%
\TextAt{14}{352.192}{723.676}{a}%
\TextAt{9}{358.175}{723.676}{\Character{41}}%
\TextAt{9}{365.372}{723.676}{\Character{43}}%
\TextAt{13}{377.146}{723.676}{\Character{34}\Character{58}}%
\TextAt{8}{79.226}{692.486}{Then}%
\TextAt{8}{109.784}{692.486}{by}%
\TextAt{8}{126.043}{692.486}{(5.3)}%
\TextAt{8}{154.018}{692.486}{we}%
\TextAt{8}{171.245}{692.486}{find}%
\TextAt{8}{194.653}{692.486}{an}%
\TextAt{8}{210.912}{692.486}{element}%
\TextAt{13}{254.154}{692.486}{a}%
\NativeMathAt{263.62}{692.486}{11.955}{\in}%
\TextAt{14}{274.911}{692.486}{a}%
\TextAt{8}{280.894}{692.486}{,}%
\TextAt{13}{288.043}{692.486}{a}%
\NativeMathAt{297.509}{692.486}{11.955}{\neq}%
\TextAt{9}{309.918}{692.486}{0}%
\TextAt{8}{315.787}{692.486}{,}%
\TextAt{8}{322.937}{692.486}{satisfying}%
\TextAt{16}{375.097}{692.486}{\Character{106}}%
\TextAt{13}{378.418}{692.486}{\Character{28}}%
\TextAt{13}{384.826}{692.486}{a}%
\TextAt{16}{390.982}{692.486}{\Character{106}}%
\TextAt{13}{397.624}{692.486}{\Character{60}}%
\TextAt{13}{410.034}{692.486}{c}%
\TextAt{21}{415.087}{690.692}{\Character{28}}%
\TextAt{8}{420.264}{692.486}{.}%
\TextAt{8}{428.712}{692.486}{Thus}%
\TextAt{16}{189.742}{659.369}{\Character{106}}%
\TextAt{13}{193.063}{659.369}{N}%
\TextAt{21}{202.437}{657.509}{K}%
\TextAt{3}{210.141}{657.509}{\Character{106}}%
\TextAt{25}{212.493}{657.509}{Q}%
\TextAt{9}{219.191}{659.369}{\Character{40}}%
\TextAt{13}{223.743}{659.369}{a}%
\TextAt{9}{229.888}{659.369}{\Character{41}}%
\TextAt{16}{234.44}{659.369}{\Character{106}}%
\TextAt{9}{241.082}{659.369}{\Character{61}}%
\TextAt{20}{253.507}{670.726}{\Character{89}}%
\TextAt{21}{258.806}{645.62}{\Character{28}}%
\TextAt{16}{270.776}{659.369}{\Character{106}}%
\TextAt{13}{274.097}{659.369}{\Character{28}}%
\TextAt{13}{280.505}{659.369}{a}%
\TextAt{16}{286.661}{659.369}{\Character{106}}%
\TextAt{13}{293.302}{659.369}{\Character{60}}%
\TextAt{20}{305.728}{676.226}{\Character{18}}%
\TextAt{9}{316.332}{667.457}{2}%
\Vector{q 0 G 0 g 0.478 w 243.724 -111.642 m 250.793 -111.642 l S Q}%
\TextAt{13}{315.724}{651.168}{\Character{25}}%
\TextAt{20}{323.989}{676.226}{\Character{19}}%
\TextAt{21}{332.789}{673.549}{s}%
\TextAt{20}{339.195}{670.174}{\Character{112}}%
\Vector{q 0 G 0 g 0.478 w 279.151 -103.587 m 300.078 -103.587 l S Q}%
\TextAt{16}{351.151}{659.369}{\Character{106}}%
\TextAt{13}{354.472}{659.369}{d}%
\TextAt{21}{360.554}{657.576}{K}%
\TextAt{16}{368.757}{659.369}{\Character{106}}%
\TextAt{14}{374.07}{659.369}{N}%
\TextAt{9}{384.018}{659.369}{\Character{40}}%
\TextAt{14}{388.571}{659.369}{a}%
\TextAt{9}{394.554}{659.369}{\Character{41}}%
\TextAt{9}{401.751}{659.369}{\Character{43}}%
\TextAt{13}{413.524}{659.369}{\Character{34}\Character{58}}%
\TextAt{8}{79.226}{622.954}{This}%
\TextAt{8}{105.734}{622.954}{being}%
\TextAt{8}{137.061}{622.954}{true}%
\TextAt{8}{161.561}{622.954}{for}%
\TextAt{8}{179.234}{622.954}{all}%
\TextAt{13}{195.282}{622.954}{\Character{34}}%
\TextAt{13}{204.069}{622.954}{\Character{62}}%
\TextAt{9}{216.506}{622.954}{0}%
\TextAt{8}{226.045}{622.954}{and}%
\TextAt{8}{248.596}{622.954}{since}%
\TextAt{16}{277.053}{622.954}{\Character{106}}%
\TextAt{13}{280.374}{622.954}{N}%
\TextAt{21}{289.748}{621.095}{K}%
\TextAt{3}{297.452}{621.095}{\Character{106}}%
\TextAt{25}{299.804}{621.095}{Q}%
\TextAt{9}{306.502}{622.954}{\Character{40}}%
\TextAt{13}{311.054}{622.954}{a}%
\TextAt{9}{317.199}{622.954}{\Character{41}}%
\TextAt{16}{321.751}{622.954}{\Character{106}}%
\TextAt{8}{328.758}{622.954}{is}%
\TextAt{8}{340.31}{622.954}{always}%
\TextAt{8}{377.566}{622.954}{a}%
\TextAt{8}{387.11}{622.954}{positive}%
\TextAt{8}{430.212}{622.954}{integer,}%
\TextAt{8}{471.994}{622.954}{there}%
\TextAt{8}{501.694}{622.954}{has}%
\TextAt{8}{522.357}{622.954}{to}%
\TextAt{8}{79.226}{605.908}{exist}%
\TextAt{8}{106.926}{605.908}{an}%
\TextAt{13}{123.185}{605.908}{a}%
\NativeMathAt{132.651}{605.908}{11.955}{\in}%
\TextAt{14}{143.942}{605.908}{a}%
\TextAt{8}{149.925}{605.908}{,}%
\TextAt{13}{157.074}{605.908}{a}%
\NativeMathAt{166.54}{605.908}{11.955}{\neq}%
\TextAt{9}{178.95}{605.908}{0}%
\TextAt{8}{184.819}{605.908}{,}%
\TextAt{8}{191.968}{605.908}{such}%
\TextAt{8}{218.365}{605.908}{that}%
\TextAt{16}{225.703}{573.468}{\Character{106}}%
\TextAt{13}{229.024}{573.468}{N}%
\TextAt{21}{238.398}{571.608}{K}%
\TextAt{3}{246.102}{571.608}{\Character{106}}%
\TextAt{25}{248.455}{571.608}{Q}%
\TextAt{9}{255.152}{573.468}{\Character{40}}%
\TextAt{13}{259.704}{573.468}{a}%
\TextAt{9}{265.849}{573.468}{\Character{41}}%
\TextAt{16}{270.401}{573.468}{\Character{106}}%
\NativeMathAt{277.024}{573.468}{11.955}{\leq}%
\TextAt{20}{289.663}{590.325}{\Character{18}}%
\TextAt{9}{300.267}{581.556}{2}%
\Vector{q 0 G 0 g 0.478 w 227.658 -197.543 m 234.728 -197.543 l S Q}%
\TextAt{13}{299.658}{565.267}{\Character{25}}%
\TextAt{20}{307.923}{590.325}{\Character{19}}%
\TextAt{21}{316.724}{587.648}{s}%
\TextAt{20}{323.13}{584.273}{\Character{112}}%
\Vector{q 0 G 0 g 0.478 w 263.085 -189.488 m 284.012 -189.488 l S Q}%
\TextAt{16}{335.085}{573.468}{\Character{106}}%
\TextAt{13}{338.406}{573.468}{d}%
\TextAt{21}{344.489}{571.675}{K}%
\TextAt{16}{352.692}{573.468}{\Character{106}}%
\TextAt{14}{358.005}{573.468}{N}%
\TextAt{9}{367.953}{573.468}{\Character{40}}%
\TextAt{14}{372.505}{573.468}{a}%
\TextAt{9}{378.488}{573.468}{\Character{41}}%
\TextAt{13}{383.041}{573.468}{\Character{58}}%
\Vector{q 0 G 0 g 0.399 w 452.334 -232.464 m 452.334 -224.394 l S Q}%
\Vector{q 0 G 0 g 0.399 w 452.534 -224.593 m 459.707 -224.593 l S Q}%
\Vector{q 0 G 0 g 0.399 w 452.534 -232.264 m 459.707 -232.264 l S Q}%
\Vector{q 0 G 0 g 0.399 w 459.906 -232.464 m 459.906 -224.394 l S Q}%
\TextAt{12}{79.226}{506.813}{Theorem}%
\TextAt{12}{136.24}{506.813}{(6.3).}%
\TextAt{8}{171.556}{506.813}{The}%
\TextAt{8}{195.61}{506.813}{ideal}%
\TextAt{8}{223.573}{506.813}{class}%
\TextAt{8}{251.01}{506.813}{group}%
\TextAt{13}{284.186}{506.813}{C}%
\TextAt{13}{293.404}{506.813}{l}%
\TextAt{21}{296.906}{505.02}{K}%
\TextAt{9}{308.43}{506.813}{\Character{61}}%
\TextAt{13}{320.855}{506.813}{J}%
\TextAt{21}{327.325}{505.02}{K}%
\TextAt{9}{335.528}{506.813}{\Character{47}}%
\TextAt{13}{341.381}{506.813}{P}%
\TextAt{21}{348.926}{505.02}{K}%
\TextAt{8}{361.026}{506.813}{is}%
\TextAt{8}{372.79}{506.813}{finite.}%
\TextAt{8}{407.253}{506.813}{Its}%
\TextAt{8}{424.54}{506.813}{order}%
\TextAt{13}{265.945}{478.259}{h}%
\TextAt{21}{272.683}{476.466}{K}%
\TextAt{9}{284.207}{478.259}{\Character{61}}%
\TextAt{9}{296.617}{478.259}{\Character{40}}%
\TextAt{13}{301.185}{478.259}{J}%
\TextAt{21}{307.655}{476.466}{K}%
\TextAt{9}{319.178}{478.259}{\Character{58}}%
\TextAt{13}{325.751}{478.259}{P}%
\TextAt{21}{333.296}{476.466}{K}%
\TextAt{9}{341.499}{478.259}{\Character{41}}%
\TextAt{8}{79.226}{449.706}{is}%
\TextAt{8}{90.99}{449.706}{called}%
\TextAt{8}{124.154}{449.706}{the}%
\TextAt{12}{144.31}{449.706}{class}%
\TextAt{12}{175.669}{449.706}{number}%
\TextAt{8}{224.458}{449.706}{of}%
\TextAt{13}{237.788}{449.706}{K}%
\TextAt{8}{248.605}{449.706}{.}%
\TextAt{19}{79.226}{414.982}{Proof.}%
\TextAt{8}{116}{414.982}{If}%
\TextAt{14}{127.282}{414.982}{p}%
\NativeMathAt{136.586}{414.982}{11.955}{\neq}%
\TextAt{9}{148.996}{414.982}{0}%
\TextAt{8}{158.352}{414.982}{is}%
\TextAt{8}{169.705}{414.982}{a}%
\TextAt{8}{179.05}{414.982}{prime}%
\TextAt{8}{211.804}{414.982}{ideal}%
\TextAt{8}{239.357}{414.982}{of}%
\TextAt{14}{252.276}{414.982}{o}%
\TextAt{21}{258.121}{413.189}{K}%
\TextAt{8}{269.81}{414.982}{and}%
\TextAt{14}{292.163}{414.982}{p}%
\TextAt{16}{299.964}{414.982}{\Character{92}}%
\TextAt{18}{309.751}{414.982}{Z}%
\TextAt{9}{321.042}{414.982}{\Character{61}}%
\TextAt{13}{333.468}{414.982}{p}%
\TextAt{18}{339.343}{414.982}{Z}%
\TextAt{8}{347.313}{414.982}{,}%
\TextAt{8}{354.134}{414.982}{then}%
\TextAt{14}{380.384}{414.982}{o}%
\TextAt{21}{386.228}{413.189}{K}%
\TextAt{9}{394.431}{414.982}{\Character{47}}%
\TextAt{14}{400.284}{414.982}{p}%
\TextAt{8}{409.754}{414.982}{is}%
\TextAt{8}{421.108}{414.982}{a}%
\TextAt{8}{430.453}{414.982}{finite}%
\TextAt{8}{459.954}{414.982}{field}%
\TextAt{8}{484.901}{414.982}{extension}%
\TextAt{8}{79.226}{397.936}{of}%
\TextAt{18}{92.556}{397.936}{Z}%
\TextAt{9}{100.526}{397.936}{\Character{47}}%
\TextAt{13}{106.379}{397.936}{p}%
\TextAt{18}{112.254}{397.936}{Z}%
\TextAt{8}{124.122}{397.936}{of}%
\TextAt{8}{137.452}{397.936}{degree,}%
\TextAt{8}{177.119}{397.936}{say,}%
\TextAt{13}{199.619}{397.936}{f}%
\NativeMathAt{209.986}{397.936}{11.955}{\geq}%
\TextAt{9}{222.605}{397.936}{1}%
\TextAt{8}{228.458}{397.936}{,}%
\TextAt{8}{235.608}{397.936}{and}%
\TextAt{8}{258.37}{397.936}{we}%
\TextAt{8}{275.598}{397.936}{have}%
\TextAt{14}{278.321}{369.383}{N}%
\TextAt{9}{288.269}{369.383}{\Character{40}}%
\TextAt{14}{292.821}{369.383}{p}%
\TextAt{9}{298.804}{369.383}{\Character{41}}%
\TextAt{9}{306.659}{369.383}{\Character{61}}%
\TextAt{13}{319.103}{369.383}{p}%
\TextAt{21}{324.978}{374.319}{f}%
\TextAt{13}{330.423}{369.383}{\Character{58}}%
\TextAt{8}{79.226}{340.829}{Given}%
\TextAt{13}{113.467}{340.829}{p}%
\TextAt{8}{119.342}{340.829}{,}%
\TextAt{8}{126.926}{340.829}{there}%
\TextAt{8}{157.186}{340.829}{are}%
\TextAt{8}{177.044}{340.829}{only}%
\TextAt{8}{203.084}{340.829}{finitely}%
\TextAt{8}{242.776}{340.829}{many}%
\TextAt{8}{274.996}{340.829}{prime}%
\TextAt{8}{308.508}{340.829}{ideals}%
\TextAt{14}{341.433}{340.829}{p}%
\TextAt{8}{351.662}{340.829}{such}%
\TextAt{8}{378.407}{340.829}{that}%
\TextAt{14}{404.123}{340.829}{p}%
\TextAt{16}{413}{340.829}{\Character{92}}%
\TextAt{18}{423.864}{340.829}{Z}%
\TextAt{9}{435.748}{340.829}{\Character{61}}%
\TextAt{13}{448.766}{340.829}{p}%
\TextAt{18}{454.641}{340.829}{Z}%
\TextAt{8}{462.611}{340.829}{,}%
\TextAt{8}{470.195}{340.829}{because}%
\TextAt{8}{513.845}{340.829}{this}%
\TextAt{8}{79.226}{323.783}{means}%
\TextAt{8}{114.868}{323.783}{that}%
\TextAt{14}{140.049}{323.783}{p}%
\TextAt{16}{149.353}{323.783}{\Character{106}}%
\TextAt{9}{155.995}{323.783}{\Character{40}}%
\TextAt{13}{160.547}{323.783}{p}%
\TextAt{9}{166.423}{323.783}{\Character{41}}%
\TextAt{8}{170.975}{323.783}{.}%
\TextAt{8}{179.361}{323.783}{It}%
\TextAt{8}{191.846}{323.783}{follows}%
\TextAt{8}{230.094}{323.783}{that}%
\TextAt{8}{255.275}{323.783}{there}%
\TextAt{8}{284.999}{323.783}{are}%
\TextAt{8}{304.323}{323.783}{only}%
\TextAt{8}{329.827}{323.783}{finitely}%
\TextAt{8}{368.984}{323.783}{many}%
\TextAt{8}{400.669}{323.783}{prime}%
\TextAt{8}{433.645}{323.783}{ideals}%
\TextAt{14}{466.035}{323.783}{p}%
\TextAt{8}{475.728}{323.783}{of}%
\TextAt{8}{488.87}{323.783}{bounded}%
\TextAt{8}{79.226}{306.737}{absolute}%
\TextAt{8}{126.274}{306.737}{norm.}%
\TextAt{8}{163.811}{306.737}{Since}%
\TextAt{8}{195.174}{306.737}{every}%
\TextAt{8}{226.871}{306.737}{integral}%
\TextAt{8}{270.285}{306.737}{ideal}%
\TextAt{8}{299.054}{306.737}{admits}%
\TextAt{8}{338.295}{306.737}{a}%
\TextAt{8}{348.856}{306.737}{representation}%
\TextAt{14}{426.151}{306.737}{a}%
\TextAt{9}{436.827}{306.737}{\Character{61}}%
\TextAt{14}{450.625}{306.737}{p}%
\TextAt{21}{456.608}{311.962}{\Character{23}}%
\TextAt{24}{460.771}{310.855}{1}%
\TextAt{17}{456.608}{303.806}{1}%
\TextAt{16}{467.413}{306.737}{\Character{1}}%
\TextAt{16}{472.709}{306.737}{\Character{1}}%
\TextAt{16}{478.017}{306.737}{\Character{1}}%
\TextAt{14}{483.353}{306.737}{p}%
\TextAt{21}{489.336}{311.075}{\Character{23}}%
\TextAt{22}{493.499}{310.079}{r}%
\TextAt{21}{489.336}{303.781}{r}%
\TextAt{8}{502.858}{306.737}{where}%
\TextAt{13}{79.226}{289.691}{\Character{23}}%
\TextAt{21}{85.01}{287.898}{i}%
\TextAt{13}{91.712}{289.691}{\Character{62}}%
\TextAt{9}{104.138}{289.691}{0}%
\TextAt{8}{113.888}{289.691}{and}%
\TextAt{14}{236.52}{272.645}{N}%
\TextAt{9}{246.468}{272.645}{\Character{40}}%
\TextAt{14}{251.02}{272.645}{a}%
\TextAt{9}{257.003}{272.645}{\Character{41}}%
\TextAt{9}{264.858}{272.645}{\Character{61}}%
\TextAt{14}{277.302}{272.645}{N}%
\TextAt{9}{287.25}{272.645}{\Character{40}}%
\TextAt{14}{291.802}{272.645}{p}%
\TextAt{17}{297.786}{270.852}{1}%
\TextAt{9}{302.518}{272.645}{\Character{41}}%
\TextAt{21}{307.07}{277.581}{\Character{23}}%
\TextAt{24}{311.233}{276.474}{1}%
\TextAt{16}{317.875}{272.645}{\Character{1}}%
\TextAt{16}{323.171}{272.645}{\Character{1}}%
\TextAt{16}{328.479}{272.645}{\Character{1}}%
\TextAt{14}{333.815}{272.645}{N}%
\TextAt{9}{343.763}{272.645}{\Character{40}}%
\TextAt{14}{348.315}{272.645}{p}%
\TextAt{21}{354.298}{270.852}{r}%
\TextAt{9}{358.853}{272.645}{\Character{41}}%
\TextAt{21}{363.405}{277.581}{\Character{23}}%
\TextAt{22}{367.568}{276.585}{r}%
\TextAt{13}{372.224}{272.645}{\Character{59}}%
\TextAt{8}{79.226}{249.315}{there}%
\TextAt{8}{109.567}{249.315}{are}%
\TextAt{8}{129.507}{249.315}{altogether}%
\TextAt{8}{185.229}{249.315}{only}%
\TextAt{8}{211.35}{249.315}{a}%
\TextAt{8}{221.535}{249.315}{finite}%
\TextAt{8}{251.876}{249.315}{number}%
\TextAt{8}{294.901}{249.315}{of}%
\TextAt{8}{308.661}{249.315}{ideals}%
\TextAt{14}{341.668}{249.315}{a}%
\TextAt{8}{351.978}{249.315}{of}%
\TextAt{14}{365.737}{249.315}{o}%
\TextAt{21}{371.581}{247.522}{K}%
\TextAt{8}{384.11}{249.315}{with}%
\TextAt{8}{411.2}{249.315}{bounded}%
\TextAt{8}{459.426}{249.315}{absolute}%
\TextAt{8}{506.097}{249.315}{norm}%
\TextAt{14}{79.226}{232.269}{N}%
\TextAt{9}{89.174}{232.269}{\Character{40}}%
\TextAt{14}{93.727}{232.269}{a}%
\TextAt{9}{99.71}{232.269}{\Character{41}}%
\NativeMathAt{107.583}{232.269}{11.955}{\leq}%
\TextAt{13}{120.202}{232.269}{M}%
\TextAt{8}{132.776}{232.269}{.}%
\TextAt{8}{97.159}{212.234}{It}%
\TextAt{8}{110.819}{212.234}{therefore}%
\TextAt{8}{160.907}{212.234}{suffices}%
\TextAt{8}{201.682}{212.234}{to}%
\TextAt{8}{216.98}{212.234}{show}%
\TextAt{8}{246.971}{212.234}{that}%
\TextAt{8}{273.328}{212.234}{each}%
\TextAt{8}{300.653}{212.234}{class}%
\TextAt{9}{329.078}{212.234}{\Character{91}}%
\TextAt{14}{332.33}{212.234}{a}%
\TextAt{9}{338.313}{212.234}{\Character{93}}%
\NativeMathAt{346.569}{212.234}{11.955}{\in}%
\TextAt{13}{359.544}{212.234}{C}%
\TextAt{13}{368.762}{212.234}{l}%
\TextAt{21}{372.264}{210.441}{K}%
\TextAt{8}{385.352}{212.234}{contains}%
\TextAt{8}{432.26}{212.234}{an}%
\TextAt{8}{449.507}{212.234}{integral}%
\TextAt{8}{493.103}{212.234}{ideal}%
\TextAt{14}{522.054}{212.234}{a}%
\TextAt{17}{528.037}{210.441}{1}%
\TextAt{8}{79.226}{195.188}{satisfying}%
\TextAt{14}{234.183}{174.129}{N}%
\TextAt{9}{244.131}{174.129}{\Character{40}}%
\TextAt{14}{248.683}{174.129}{a}%
\TextAt{17}{254.666}{172.336}{1}%
\TextAt{9}{259.399}{174.129}{\Character{41}}%
\NativeMathAt{267.272}{174.129}{11.955}{\leq}%
\TextAt{13}{279.891}{174.129}{M}%
\TextAt{9}{295.786}{174.129}{\Character{61}}%
\TextAt{20}{308.211}{190.986}{\Character{18}}%
\TextAt{9}{318.815}{182.217}{2}%
\Vector{q 0 G 0 g 0.478 w 246.207 -596.882 m 253.276 -596.882 l S Q}%
\TextAt{13}{318.207}{165.929}{\Character{25}}%
\TextAt{20}{326.472}{190.986}{\Character{19}}%
\TextAt{21}{335.272}{188.31}{s}%
\TextAt{20}{341.679}{184.935}{\Character{112}}%
\Vector{q 0 G 0 g 0.478 w 281.634 -588.826 m 302.561 -588.826 l S Q}%
\TextAt{16}{353.634}{174.129}{\Character{106}}%
\TextAt{13}{356.955}{174.129}{d}%
\TextAt{21}{363.038}{172.336}{K}%
\TextAt{16}{371.24}{174.129}{\Character{106}}%
\TextAt{13}{374.561}{174.129}{\Character{58}}%
\TextAt{8}{79.226}{147.075}{For}%
\TextAt{8}{99.893}{147.075}{this,}%
\TextAt{8}{125.726}{147.075}{choose}%
\TextAt{8}{162.556}{147.075}{an}%
\TextAt{8}{178.513}{147.075}{arbitrary}%
\TextAt{8}{227.981}{147.075}{representative}%
\TextAt{14}{302.864}{147.075}{a}%
\TextAt{8}{312.443}{147.075}{of}%
\TextAt{8}{325.471}{147.075}{the}%
\TextAt{8}{345.325}{147.075}{class,}%
\TextAt{8}{375.772}{147.075}{and}%
\TextAt{8}{398.233}{147.075}{a}%
\TextAt{13}{407.686}{147.075}{\Character{13}}%
\NativeMathAt{417.729}{147.075}{11.955}{\in}%
\TextAt{14}{429.02}{147.075}{o}%
\TextAt{21}{434.864}{145.282}{K}%
\TextAt{8}{443.067}{147.075}{,}%
\TextAt{13}{449.974}{147.075}{\Character{13}}%
\NativeMathAt{460.017}{147.075}{11.955}{\neq}%
\TextAt{9}{472.427}{147.075}{0}%
\TextAt{8}{478.296}{147.075}{,}%
\TextAt{8}{485.203}{147.075}{such}%
\TextAt{8}{511.298}{147.075}{that}%
\TextAt{14}{79.226}{130.029}{b}%
\TextAt{9}{88.682}{130.029}{\Character{61}}%
\TextAt{13}{101.107}{130.029}{\Character{13}}%
\TextAt{14}{107.83}{130.029}{a}%
\TextAt{3}{113.813}{134.367}{\Character{0}}%
\TextAt{17}{120.399}{134.367}{1}%
\NativeMathAt{128.452}{130.029}{11.955}{\subseteq}%
\TextAt{14}{141.072}{130.029}{o}%
\TextAt{21}{146.916}{128.236}{K}%
\TextAt{8}{155.118}{130.029}{.}%
\TextAt{8}{163.567}{130.029}{By}%
\TextAt{8}{181.93}{130.029}{(6.2),}%
\TextAt{8}{213.157}{130.029}{there}%
\TextAt{8}{243.069}{130.029}{exists}%
\TextAt{13}{275.384}{130.029}{a}%
\NativeMathAt{284.849}{130.029}{11.955}{\in}%
\TextAt{14}{296.14}{130.029}{b}%
\TextAt{8}{302.275}{130.029}{,}%
\TextAt{13}{309.424}{130.029}{a}%
\NativeMathAt{318.89}{130.029}{11.955}{\neq}%
\TextAt{9}{331.3}{130.029}{0}%
\TextAt{8}{337.169}{130.029}{,}%
\TextAt{8}{344.318}{130.029}{such}%
\TextAt{8}{370.715}{130.029}{that}%
\TextAt{16}{184.231}{101.475}{\Character{106}}%
\TextAt{13}{187.552}{101.475}{N}%
\TextAt{21}{196.926}{99.616}{K}%
\TextAt{3}{204.63}{99.615}{\Character{106}}%
\TextAt{25}{206.983}{99.616}{Q}%
\TextAt{9}{213.68}{101.475}{\Character{40}}%
\TextAt{13}{218.232}{101.475}{a}%
\TextAt{9}{224.377}{101.475}{\Character{41}}%
\TextAt{16}{228.93}{101.475}{\Character{106}}%
\TextAt{14}{234.243}{101.475}{N}%
\TextAt{9}{244.191}{101.475}{\Character{40}}%
\TextAt{14}{248.743}{101.475}{b}%
\TextAt{9}{254.878}{101.475}{\Character{41}}%
\TextAt{3}{259.431}{106.411}{\Character{0}}%
\TextAt{17}{266.017}{106.411}{1}%
\TextAt{9}{274.07}{101.475}{\Character{61}}%
\TextAt{14}{286.496}{101.475}{N}%
\TextAt{9}{296.444}{101.475}{\Character{40}\Character{40}}%
\TextAt{13}{305.548}{101.475}{a}%
\TextAt{9}{311.693}{101.475}{\Character{41}}%
\TextAt{14}{316.246}{101.475}{b}%
\TextAt{3}{322.381}{106.411}{\Character{0}}%
\TextAt{17}{328.967}{106.411}{1}%
\TextAt{9}{333.699}{101.475}{\Character{41}}%
\TextAt{9}{341.554}{101.475}{\Character{61}}%
\TextAt{14}{353.998}{101.475}{N}%
\TextAt{9}{363.946}{101.475}{\Character{40}}%
\TextAt{13}{368.498}{101.475}{a}%
\TextAt{14}{374.643}{101.475}{b}%
\TextAt{3}{380.778}{106.411}{\Character{0}}%
\TextAt{17}{387.365}{106.411}{1}%
\TextAt{9}{392.097}{101.475}{\Character{41}}%
\NativeMathAt{399.97}{101.475}{11.955}{\leq}%
\TextAt{13}{412.589}{101.475}{M}%
\TextAt{13}{424.508}{101.475}{\Character{58}}%
\TextAt{8}{79.226}{72.921}{The}%
\TextAt{8}{103.28}{72.921}{ideal}%
\TextAt{14}{131.243}{72.921}{a}%
\TextAt{17}{137.226}{71.128}{1}%
\TextAt{9}{145.279}{72.921}{\Character{61}}%
\TextAt{13}{157.705}{72.921}{a}%
\TextAt{14}{163.85}{72.921}{b}%
\TextAt{3}{169.985}{77.26}{\Character{0}}%
\TextAt{17}{176.571}{77.26}{1}%
\TextAt{9}{184.624}{72.921}{\Character{61}}%
\TextAt{13}{197.05}{72.921}{a\Character{13}}%
\TextAt{3}{209.917}{77.26}{\Character{0}}%
\TextAt{17}{216.503}{77.26}{1}%
\TextAt{14}{221.236}{72.921}{a}%
\NativeMathAt{230.54}{72.921}{11.955}{\in}%
\TextAt{9}{241.831}{72.921}{\Character{91}}%
\TextAt{14}{245.082}{72.921}{a}%
\TextAt{9}{251.065}{72.921}{\Character{93}}%
\TextAt{8}{258.214}{72.921}{therefore}%
\TextAt{8}{307.314}{72.921}{has}%
\TextAt{8}{328.188}{72.921}{the}%
\TextAt{8}{348.344}{72.921}{required}%
\TextAt{8}{394.192}{72.921}{property.}%
\Vector{q 0 G 0 g 0.399 w 452.334 -701.079 m 452.334 -693.009 l S Q}%
\Vector{q 0 G 0 g 0.399 w 452.534 -693.208 m 459.707 -693.208 l S Q}%
\Vector{q 0 G 0 g 0.399 w 452.534 -700.879 m 459.707 -700.879 l S Q}%
\Vector{q 0 G 0 g 0.399 w 459.906 -701.079 m 459.906 -693.009 l S Q}%
\TextAt{8}{300.14}{43.033}{34}%
\Destination{78.23}{809.815}{nav.page.34}%
\Destination{79.36}{776.155}{nav.print.36}%
\Destination{79.36}{520.855}{nav.number.1.6.3}%
\Hotlink{125.543}{690.255}{25.078}{11.047}{nav.number.1.5.3}%
\Hotlink{181.43}{127.855}{25.078}{11.047}{nav.number.1.6.2}%
}

% Page 47 (body page 35)
\Page{%
\TextAt{2}{79.226}{796.943}{I.}%
\TextAt{2}{89.83}{796.943}{Algebraic}%
\TextAt{2}{138.648}{796.943}{Integers}%
\Vector{q 0 G 0 g 0.299 w 7.226 18.729 m 460.769 18.729 l S Q}%
\TextAt{8}{97.159}{761.124}{The}%
\TextAt{8}{120.219}{761.124}{theorem}%
\TextAt{8}{164.75}{761.124}{of}%
\TextAt{8}{177.087}{761.124}{the}%
\TextAt{8}{196.249}{761.124}{finiteness}%
\TextAt{8}{246.101}{761.124}{of}%
\TextAt{8}{258.437}{761.124}{the}%
\TextAt{8}{277.6}{761.124}{class}%
\TextAt{8}{304.043}{761.124}{number}%
\TextAt{13}{345.645}{761.124}{h}%
\TextAt{21}{352.384}{759.33}{K}%
\TextAt{8}{363.49}{761.124}{means}%
\TextAt{8}{398.326}{761.124}{that}%
\TextAt{8}{422.701}{761.124}{passing}%
\TextAt{8}{462.809}{761.124}{from}%
\TextAt{8}{489.456}{761.124}{numbers}%
\TextAt{8}{79.226}{744.078}{to}%
\TextAt{8}{93.327}{744.078}{ideals}%
\TextAt{8}{125.695}{744.078}{has}%
\TextAt{8}{146.36}{744.078}{not}%
\TextAt{8}{166.964}{744.078}{thrust}%
\TextAt{8}{201.939}{744.078}{us}%
\TextAt{8}{216.745}{744.078}{into}%
\TextAt{8}{240.278}{744.078}{unlimited}%
\TextAt{8}{292.743}{744.078}{new}%
\TextAt{8}{316.587}{744.078}{territory.}%
\TextAt{8}{367.251}{744.078}{The}%
\TextAt{8}{391.096}{744.078}{most}%
\TextAt{8}{419.567}{744.078}{favourable}%
\TextAt{8}{475.953}{744.078}{case}%
\TextAt{8}{500.514}{744.078}{occurs}%
\TextAt{8}{79.226}{727.032}{of}%
\TextAt{8}{92.704}{727.032}{course}%
\TextAt{8}{128.682}{727.032}{when}%
\TextAt{13}{159.387}{727.032}{h}%
\TextAt{21}{166.126}{725.238}{K}%
\TextAt{9}{177.902}{727.032}{\Character{61}}%
\TextAt{9}{190.563}{727.032}{1}%
\TextAt{8}{196.432}{727.032}{.}%
\TextAt{8}{205.325}{727.032}{This}%
\TextAt{8}{232.193}{727.032}{means}%
\TextAt{8}{268.17}{727.032}{that}%
\TextAt{14}{293.687}{727.032}{o}%
\TextAt{21}{299.532}{725.238}{K}%
\TextAt{8}{311.78}{727.032}{is}%
\TextAt{8}{323.692}{727.032}{a}%
\TextAt{8}{333.595}{727.032}{principal}%
\TextAt{8}{382.52}{727.032}{ideal}%
\TextAt{8}{410.631}{727.032}{domain,}%
\TextAt{8}{455.696}{727.032}{i.e.,}%
\TextAt{8}{477.986}{727.032}{that}%
\TextAt{8}{503.503}{727.032}{prime}%
\TextAt{8}{79.226}{709.986}{factorization}%
\TextAt{8}{146.525}{709.986}{of}%
\TextAt{8}{159.177}{709.986}{elements}%
\TextAt{8}{206.355}{709.986}{in}%
\TextAt{8}{219.329}{709.986}{the}%
\TextAt{8}{238.807}{709.986}{classical}%
\TextAt{8}{283.128}{709.986}{sense}%
\TextAt{8}{312.481}{709.986}{holds.}%
\TextAt{8}{347.435}{709.986}{In}%
\TextAt{8}{361.378}{709.986}{general,}%
\TextAt{8}{404.412}{709.986}{however,}%
\TextAt{8}{452.001}{709.986}{one}%
\TextAt{8}{472.782}{709.986}{has}%
\TextAt{13}{492.977}{709.986}{h}%
\TextAt{21}{499.716}{708.193}{K}%
\TextAt{13}{511.239}{709.986}{\Character{62}}%
\TextAt{9}{523.665}{709.986}{1}%
\TextAt{8}{529.518}{709.986}{.}%
\TextAt{8}{79.226}{692.94}{For}%
\TextAt{8}{100.196}{692.94}{instance,}%
\TextAt{8}{149.034}{692.94}{we}%
\TextAt{8}{166.262}{692.94}{know}%
\TextAt{8}{196.832}{692.94}{now}%
\TextAt{8}{221.221}{692.94}{that}%
\TextAt{8}{246.591}{692.94}{the}%
\TextAt{8}{266.748}{692.94}{only}%
\TextAt{8}{292.44}{692.94}{imaginary}%
\TextAt{8}{347.411}{692.94}{quadratic}%
\TextAt{8}{399.775}{692.94}{fields}%
\TextAt{18}{429.747}{692.94}{Q}%
\TextAt{9}{439.046}{692.94}{\Character{40}}%
\TextAt{16}{443.598}{703.128}{\Character{112}}%
\Vector{q 0 G 0 g 0.478 w 381.561 -70.633 m 387.643 -70.633 l S Q}%
\TextAt{13}{453.561}{692.94}{d}%
\TextAt{9}{459.643}{692.94}{\Character{41}}%
\TextAt{8}{464.196}{692.94}{,}%
\TextAt{13}{471.346}{692.94}{d}%
\TextAt{8}{481.326}{692.94}{squarefree}%
\TextAt{8}{79.226}{675.894}{and}%
\TextAt{13}{101.989}{675.894}{\Character{60}}%
\TextAt{9}{114.414}{675.894}{0}%
\TextAt{8}{120.267}{675.894}{,}%
\TextAt{8}{127.416}{675.894}{which}%
\TextAt{8}{160.903}{675.894}{have}%
\TextAt{8}{187.898}{675.894}{class}%
\TextAt{8}{215.335}{675.894}{number}%
\TextAt{8}{257.931}{675.894}{1}%
\TextAt{8}{267.686}{675.894}{are}%
\TextAt{8}{287.197}{675.894}{those}%
\TextAt{8}{317.826}{675.894}{with}%
\TextAt{13}{186.74}{646.893}{d}%
\TextAt{9}{196.144}{646.893}{\Character{61}}%
\TextAt{16}{208.569}{646.893}{\Character{0}}%
\TextAt{9}{217.868}{646.893}{1}%
\TextAt{13}{223.721}{646.893}{\Character{59}}%
\TextAt{16}{228.965}{646.893}{\Character{0}}%
\TextAt{9}{238.263}{646.893}{2}%
\TextAt{13}{244.116}{646.893}{\Character{59}}%
\TextAt{16}{249.36}{646.893}{\Character{0}}%
\TextAt{9}{258.659}{646.893}{3}%
\TextAt{13}{264.512}{646.893}{\Character{59}}%
\TextAt{16}{269.756}{646.893}{\Character{0}}%
\TextAt{9}{279.055}{646.893}{7}%
\TextAt{13}{284.908}{646.893}{\Character{59}}%
\TextAt{16}{290.152}{646.893}{\Character{0}}%
\TextAt{9}{299.45}{646.893}{11}%
\TextAt{13}{311.156}{646.893}{\Character{59}}%
\TextAt{16}{316.4}{646.893}{\Character{0}}%
\TextAt{9}{325.699}{646.893}{19}%
\TextAt{13}{337.405}{646.893}{\Character{59}}%
\TextAt{16}{342.649}{646.893}{\Character{0}}%
\TextAt{9}{351.948}{646.893}{43}%
\TextAt{13}{363.654}{646.893}{\Character{59}}%
\TextAt{16}{368.898}{646.893}{\Character{0}}%
\TextAt{9}{378.196}{646.893}{67}%
\TextAt{13}{389.902}{646.893}{\Character{59}}%
\TextAt{16}{395.146}{646.893}{\Character{0}}%
\TextAt{9}{404.445}{646.893}{163}%
\TextAt{13}{422.004}{646.893}{\Character{58}}%
\TextAt{8}{79.226}{617.892}{Among}%
\TextAt{8}{120.078}{617.892}{real}%
\TextAt{8}{143.046}{617.892}{quadratic}%
\TextAt{8}{195.614}{617.892}{fields,}%
\TextAt{8}{229.093}{617.892}{class}%
\TextAt{8}{256.735}{617.892}{number}%
\TextAt{8}{299.536}{617.892}{1}%
\TextAt{8}{309.496}{617.892}{is}%
\TextAt{8}{321.465}{617.892}{more}%
\TextAt{8}{350.936}{617.892}{common.}%
\TextAt{8}{402.929}{617.892}{In}%
\TextAt{8}{417.755}{617.892}{the}%
\TextAt{8}{438.116}{617.892}{range}%
\TextAt{9}{470.193}{617.892}{2}%
\NativeMathAt{479.716}{617.892}{11.955}{\leq}%
\TextAt{13}{492.684}{617.892}{d}%
\TextAt{13}{502.416}{617.892}{\Character{60}}%
\TextAt{9}{515.211}{617.892}{100}%
\TextAt{8}{79.226}{600.846}{for}%
\TextAt{8}{97.111}{600.846}{instance,}%
\TextAt{8}{145.948}{600.846}{it}%
\TextAt{8}{157.652}{600.846}{occurs}%
\TextAt{8}{193.804}{600.846}{for}%
\TextAt{13}{192.061}{571.844}{d}%
\TextAt{9}{201.465}{571.844}{\Character{61}}%
\TextAt{9}{213.875}{571.844}{2}%
\TextAt{13}{219.743}{571.844}{\Character{59}}%
\TextAt{9}{224.987}{571.844}{3}%
\TextAt{13}{230.84}{571.844}{\Character{59}}%
\TextAt{9}{236.084}{571.844}{5}%
\TextAt{13}{241.937}{571.844}{\Character{59}}%
\TextAt{9}{247.181}{571.844}{6}%
\TextAt{13}{253.034}{571.844}{\Character{59}}%
\TextAt{9}{258.279}{571.844}{7}%
\TextAt{13}{264.132}{571.844}{\Character{59}}%
\TextAt{9}{269.376}{571.844}{11}%
\TextAt{13}{281.082}{571.844}{\Character{59}}%
\TextAt{9}{286.326}{571.844}{13}%
\TextAt{13}{298.032}{571.844}{\Character{59}}%
\TextAt{9}{303.276}{571.844}{14}%
\TextAt{13}{314.982}{571.844}{\Character{59}}%
\TextAt{9}{320.226}{571.844}{17}%
\TextAt{13}{331.932}{571.844}{\Character{59}}%
\TextAt{9}{337.176}{571.844}{19}%
\TextAt{13}{348.882}{571.844}{\Character{59}}%
\TextAt{9}{354.126}{571.844}{21}%
\TextAt{13}{365.832}{571.844}{\Character{59}}%
\TextAt{9}{371.077}{571.844}{22}%
\TextAt{13}{382.783}{571.844}{\Character{59}}%
\TextAt{9}{388.027}{571.844}{23}%
\TextAt{13}{399.733}{571.844}{\Character{59}}%
\TextAt{9}{404.977}{571.844}{29}%
\TextAt{13}{416.683}{571.844}{\Character{59}}%
\TextAt{9}{213.89}{551.81}{31}%
\TextAt{13}{225.596}{551.81}{\Character{59}}%
\TextAt{9}{230.84}{551.81}{33}%
\TextAt{13}{242.546}{551.81}{\Character{59}}%
\TextAt{9}{247.79}{551.81}{37}%
\TextAt{13}{259.496}{551.81}{\Character{59}}%
\TextAt{9}{264.74}{551.81}{38}%
\TextAt{13}{276.446}{551.81}{\Character{59}}%
\TextAt{9}{281.691}{551.81}{41}%
\TextAt{13}{293.397}{551.81}{\Character{59}}%
\TextAt{9}{298.641}{551.81}{43}%
\TextAt{13}{310.347}{551.81}{\Character{59}}%
\TextAt{9}{315.591}{551.81}{46}%
\TextAt{13}{327.297}{551.81}{\Character{59}}%
\TextAt{9}{332.541}{551.81}{47}%
\TextAt{13}{344.247}{551.81}{\Character{59}}%
\TextAt{9}{349.491}{551.81}{53}%
\TextAt{13}{361.197}{551.81}{\Character{59}}%
\TextAt{9}{366.441}{551.81}{57}%
\TextAt{13}{378.147}{551.81}{\Character{59}}%
\TextAt{9}{383.391}{551.81}{59}%
\TextAt{13}{395.097}{551.81}{\Character{59}}%
\TextAt{9}{400.342}{551.81}{61}%
\TextAt{13}{412.048}{551.81}{\Character{59}}%
\TextAt{9}{213.89}{531.775}{62}%
\TextAt{13}{225.596}{531.775}{\Character{59}}%
\TextAt{9}{230.84}{531.775}{67}%
\TextAt{13}{242.546}{531.775}{\Character{59}}%
\TextAt{9}{247.79}{531.775}{69}%
\TextAt{13}{259.496}{531.775}{\Character{59}}%
\TextAt{9}{264.74}{531.775}{71}%
\TextAt{13}{276.446}{531.775}{\Character{59}}%
\TextAt{9}{281.691}{531.775}{73}%
\TextAt{13}{293.397}{531.775}{\Character{59}}%
\TextAt{9}{298.641}{531.775}{77}%
\TextAt{13}{310.347}{531.775}{\Character{59}}%
\TextAt{9}{315.591}{531.775}{83}%
\TextAt{13}{327.297}{531.775}{\Character{59}}%
\TextAt{9}{332.541}{531.775}{86}%
\TextAt{13}{344.247}{531.775}{\Character{59}}%
\TextAt{9}{349.491}{531.775}{89}%
\TextAt{13}{361.197}{531.775}{\Character{59}}%
\TextAt{9}{366.441}{531.775}{93}%
\TextAt{13}{378.147}{531.775}{\Character{59}}%
\TextAt{9}{383.391}{531.775}{94}%
\TextAt{13}{395.097}{531.775}{\Character{59}}%
\TextAt{9}{400.342}{531.775}{97}%
\TextAt{13}{412.048}{531.775}{\Character{58}}%
\TextAt{8}{79.226}{502.774}{It}%
\TextAt{8}{92.829}{502.774}{is}%
\TextAt{8}{105.522}{502.774}{conjectured}%
\TextAt{8}{169.205}{502.774}{that}%
\TextAt{8}{195.504}{502.774}{there}%
\TextAt{8}{226.346}{502.774}{are}%
\TextAt{8}{246.786}{502.774}{infinitely}%
\TextAt{8}{296.816}{502.774}{many}%
\TextAt{8}{329.619}{502.774}{real}%
\TextAt{8}{353.311}{502.774}{quadratic}%
\TextAt{8}{406.605}{502.774}{fields}%
\TextAt{8}{437.506}{502.774}{of}%
\TextAt{8}{451.766}{502.774}{class}%
\TextAt{8}{480.133}{502.774}{number}%
\TextAt{8}{523.66}{502.774}{1.}%
\TextAt{8}{79.226}{485.728}{But}%
\TextAt{8}{102.888}{485.728}{we}%
\TextAt{8}{120.536}{485.728}{do}%
\TextAt{8}{137.216}{485.728}{not}%
\TextAt{8}{158.451}{485.728}{even}%
\TextAt{8}{185.532}{485.728}{yet}%
\TextAt{8}{205.464}{485.728}{know}%
\TextAt{8}{236.455}{485.728}{whether}%
\TextAt{8}{281.743}{485.728}{there}%
\TextAt{8}{312.076}{485.728}{are}%
\TextAt{8}{332.008}{485.728}{infinitely}%
\TextAt{8}{381.529}{485.728}{many}%
\TextAt{8}{413.822}{485.728}{algebraic}%
\TextAt{8}{463.678}{485.728}{number}%
\TextAt{8}{506.695}{485.728}{fields}%
\TextAt{8}{79.226}{468.682}{(of}%
\TextAt{8}{97.664}{468.682}{arbitrary}%
\TextAt{8}{147.985}{468.682}{degree)}%
\TextAt{8}{189.508}{468.682}{with}%
\TextAt{8}{216.721}{468.682}{class}%
\TextAt{8}{244.71}{468.682}{number}%
\TextAt{8}{287.859}{468.682}{1.}%
\TextAt{8}{303.822}{468.682}{It}%
\TextAt{8}{317.047}{468.682}{was}%
\TextAt{8}{340.099}{468.682}{found}%
\TextAt{8}{373.493}{468.682}{time}%
\TextAt{8}{400.705}{468.682}{and}%
\TextAt{8}{424.02}{468.682}{again}%
\TextAt{8}{455.8}{468.682}{in}%
\TextAt{8}{470.005}{468.682}{innumerable}%
\TextAt{8}{79.226}{451.636}{investigations}%
\TextAt{8}{152.952}{451.636}{that}%
\TextAt{8}{178.737}{451.636}{the}%
\TextAt{8}{199.31}{451.636}{ideal}%
\TextAt{8}{227.689}{451.636}{class}%
\TextAt{8}{255.543}{451.636}{groups}%
\TextAt{13}{293.75}{451.636}{C}%
\TextAt{13}{302.968}{451.636}{l}%
\TextAt{21}{306.47}{449.843}{K}%
\TextAt{8}{318.987}{451.636}{behave}%
\TextAt{8}{358.425}{451.636}{completely}%
\TextAt{8}{417.696}{451.636}{unpredictably,}%
\TextAt{8}{494.957}{451.636}{both}%
\TextAt{8}{523.014}{451.636}{in}%
\TextAt{8}{79.226}{434.59}{their}%
\TextAt{8}{107.312}{434.59}{size}%
\TextAt{8}{129.599}{434.59}{and}%
\TextAt{8}{152.484}{434.59}{their}%
\TextAt{8}{180.569}{434.59}{structure.}%
\TextAt{8}{235.627}{434.59}{An}%
\TextAt{8}{254.926}{434.59}{exception}%
\TextAt{8}{307.4}{434.59}{to}%
\TextAt{8}{321.832}{434.59}{this}%
\TextAt{8}{344.777}{434.59}{lack}%
\TextAt{8}{368.965}{434.59}{of}%
\TextAt{8}{382.417}{434.59}{rule}%
\TextAt{8}{405.948}{434.59}{is}%
\TextAt{26}{417.834}{434.59}{Kenkichi}%
\TextAt{26}{475.791}{434.59}{Iwasawa}%
\TextAt{8}{524.903}{434.59}{’s}%
\TextAt{8}{79.226}{417.544}{discovery}%
\TextAt{8}{130.26}{417.544}{that}%
\TextAt{8}{155.866}{417.544}{the}%
\TextAt{13}{176.259}{417.544}{p}%
\TextAt{8}{182.134}{417.544}{-part}%
\TextAt{8}{211.637}{417.544}{of}%
\TextAt{8}{225.204}{417.544}{the}%
\TextAt{8}{245.597}{417.544}{class}%
\TextAt{8}{273.271}{417.544}{number}%
\TextAt{8}{316.104}{417.544}{of}%
\TextAt{8}{329.671}{417.544}{the}%
\TextAt{8}{350.064}{417.544}{field}%
\TextAt{8}{375.658}{417.544}{of}%
\TextAt{13}{389.225}{417.544}{p}%
\TextAt{21}{395.1}{421.882}{n}%
\TextAt{8}{400.736}{417.544}{-th}%
\TextAt{8}{419.826}{417.544}{roots}%
\TextAt{8}{449.724}{417.544}{of}%
\TextAt{8}{463.291}{417.544}{unity}%
\TextAt{8}{494.097}{417.544}{obeys}%
\TextAt{8}{526.911}{417.544}{a}%
\TextAt{8}{79.226}{400.498}{very}%
\TextAt{8}{104.918}{400.498}{strict}%
\TextAt{8}{135.547}{400.498}{law}%
\TextAt{8}{156.684}{400.498}{when}%
\TextAt{13}{187.241}{400.498}{n}%
\TextAt{8}{198.126}{400.498}{varies}%
\TextAt{8}{231.039}{400.498}{(see}%
\TextAt{8}{254.507}{400.498}{[136],}%
\TextAt{8}{285.734}{400.498}{th.13.13).}%
\TextAt{8}{97.159}{380.463}{In}%
\TextAt{8}{111.404}{380.463}{the}%
\TextAt{8}{131.184}{380.463}{case}%
\TextAt{8}{155.579}{380.463}{of}%
\TextAt{8}{168.533}{380.463}{the}%
\TextAt{8}{188.314}{380.463}{field}%
\TextAt{8}{213.295}{380.463}{of}%
\TextAt{13}{226.249}{380.463}{p}%
\TextAt{8}{232.124}{380.463}{-th}%
\TextAt{8}{250.601}{380.463}{roots}%
\TextAt{8}{279.886}{380.463}{of}%
\TextAt{8}{292.84}{380.463}{unity,}%
\TextAt{8}{325.38}{380.463}{the}%
\TextAt{8}{345.16}{380.463}{question}%
\TextAt{8}{391.35}{380.463}{whether}%
\TextAt{8}{435.841}{380.463}{the}%
\TextAt{8}{455.622}{380.463}{class}%
\TextAt{8}{482.683}{380.463}{number}%
\TextAt{8}{524.903}{380.463}{is}%
\TextAt{8}{79.226}{363.417}{divisible}%
\TextAt{8}{125.453}{363.417}{by}%
\TextAt{13}{142.032}{363.417}{p}%
\TextAt{8}{152.123}{363.417}{has}%
\TextAt{8}{173.316}{363.417}{played}%
\TextAt{8}{210.386}{363.417}{a}%
\TextAt{8}{220.46}{363.417}{very}%
\TextAt{8}{246.471}{363.417}{important}%
\TextAt{8}{302.083}{363.417}{special}%
\TextAt{8}{340.504}{363.417}{role}%
\TextAt{8}{363.586}{363.417}{because}%
\TextAt{8}{407.206}{363.417}{it}%
\TextAt{8}{419.23}{363.417}{is}%
\TextAt{8}{431.313}{363.417}{intimately}%
\TextAt{8}{487.57}{363.417}{linked}%
\TextAt{8}{522.357}{363.417}{to}%
\TextAt{8}{79.226}{346.371}{the}%
\TextAt{8}{99.383}{346.371}{celebrated}%
\TextAt{12}{155.309}{346.371}{Fermat’s}%
\TextAt{12}{210.793}{346.371}{Last}%
\TextAt{12}{240.442}{346.371}{Theorem}%
\TextAt{8}{296.87}{346.371}{according}%
\TextAt{8}{349.556}{346.371}{to}%
\TextAt{8}{363.867}{346.371}{which}%
\TextAt{8}{397.353}{346.371}{the}%
\TextAt{8}{417.51}{346.371}{equation}%
\TextAt{13}{274.395}{317.37}{x}%
\TextAt{21}{281.047}{322.306}{p}%
\TextAt{9}{288.464}{317.37}{\Character{43}}%
\TextAt{13}{300.226}{317.37}{y}%
\TextAt{21}{306.362}{322.306}{p}%
\TextAt{9}{314.444}{317.37}{\Character{61}}%
\TextAt{13}{326.869}{317.37}{z}%
\TextAt{21}{332.84}{322.306}{p}%
\TextAt{8}{79.226}{288.369}{for}%
\TextAt{13}{97.564}{288.369}{p}%
\NativeMathAt{107.531}{288.369}{11.955}{\geq}%
\TextAt{9}{120.922}{288.369}{3}%
\TextAt{8}{131.125}{288.369}{has}%
\TextAt{8}{152.452}{288.369}{no}%
\TextAt{8}{169.163}{288.369}{solutions}%
\TextAt{8}{218.525}{288.369}{in}%
\TextAt{8}{232.63}{288.369}{integers}%
\NativeMathAt{276.396}{288.369}{11.955}{\neq}%
\TextAt{9}{289.583}{288.369}{0}%
\TextAt{8}{295.446}{288.369}{.}%
\TextAt{8}{305.253}{288.369}{In}%
\TextAt{8}{320.327}{288.369}{a}%
\TextAt{8}{330.535}{288.369}{similar}%
\TextAt{8}{369.424}{288.369}{way}%
\TextAt{8}{393.619}{288.369}{as}%
\TextAt{8}{408.442}{288.369}{the}%
\TextAt{8}{429.051}{288.369}{sums}%
\TextAt{8}{458.89}{288.369}{of}%
\TextAt{8}{472.672}{288.369}{two}%
\TextAt{8}{495.242}{288.369}{squares}%
\TextAt{13}{79.226}{271.323}{x}%
\TextAt{17}{85.878}{275.661}{2}%
\TextAt{9}{93.639}{271.323}{\Character{43}}%
\TextAt{13}{105.772}{271.323}{y}%
\TextAt{17}{111.908}{275.661}{2}%
\TextAt{9}{120.89}{271.323}{\Character{61}}%
\TextAt{9}{134.232}{271.323}{\Character{40}}%
\TextAt{13}{138.797}{271.323}{x}%
\TextAt{9}{148.477}{271.323}{\Character{43}}%
\TextAt{13}{160.61}{271.323}{iy}%
\TextAt{9}{170.74}{271.323}{\Character{41}\Character{40}}%
\TextAt{13}{179.845}{271.323}{x}%
\TextAt{16}{189.525}{271.323}{\Character{0}}%
\TextAt{13}{201.852}{271.323}{iy}%
\TextAt{9}{211.982}{271.323}{\Character{41}}%
\TextAt{8}{220.977}{271.323}{lead}%
\TextAt{8}{246.233}{271.323}{to}%
\TextAt{8}{261.089}{271.323}{studying}%
\TextAt{8}{309.502}{271.323}{the}%
\TextAt{8}{330.204}{271.323}{gaussian}%
\TextAt{8}{377.709}{271.323}{integers,}%
\TextAt{8}{424.956}{271.323}{the}%
\TextAt{8}{445.657}{271.323}{decomposition}%
\TextAt{8}{523.337}{271.323}{of}%
\TextAt{13}{79.226}{254.277}{x}%
\TextAt{21}{85.878}{258.615}{p}%
\TextAt{9}{93.953}{254.277}{\Character{43}}%
\TextAt{13}{106.371}{254.277}{y}%
\TextAt{21}{112.508}{258.615}{p}%
\TextAt{8}{122.13}{254.277}{by}%
\TextAt{8}{139.352}{254.277}{means}%
\TextAt{8}{176.146}{254.277}{of}%
\TextAt{8}{190.44}{254.277}{a}%
\TextAt{13}{201.159}{254.277}{p}%
\TextAt{8}{207.034}{254.277}{-th}%
\TextAt{8}{226.851}{254.277}{root}%
\TextAt{8}{252.861}{254.277}{of}%
\TextAt{8}{267.155}{254.277}{unity}%
\TextAt{13}{298.688}{254.277}{\Character{16}}%
\NativeMathAt{309.666}{254.277}{11.955}{\neq}%
\TextAt{9}{323.725}{254.277}{1}%
\TextAt{8}{334.448}{254.277}{leads}%
\TextAt{8}{364.738}{254.277}{to}%
\TextAt{8}{380.012}{254.277}{a}%
\TextAt{8}{390.731}{254.277}{problem}%
\TextAt{8}{437.22}{254.277}{in}%
\TextAt{8}{451.837}{254.277}{the}%
\TextAt{8}{472.957}{254.277}{ring}%
\TextAt{18}{497.987}{254.277}{Z}%
\TextAt{9}{505.957}{254.277}{\Character{91}}%
\TextAt{13}{509.208}{254.277}{\Character{16}}%
\TextAt{9}{515.224}{254.277}{\Character{93}}%
\TextAt{8}{523.337}{254.277}{of}%
\TextAt{8}{79.226}{237.231}{integers}%
\TextAt{8}{122.54}{237.231}{of}%
\TextAt{18}{135.87}{237.231}{Q}%
\TextAt{9}{145.168}{237.231}{\Character{40}}%
\TextAt{13}{149.721}{237.231}{\Character{16}}%
\TextAt{9}{155.736}{237.231}{\Character{41}}%
\TextAt{8}{160.288}{237.231}{.}%
\TextAt{8}{168.737}{237.231}{The}%
\TextAt{8}{192.791}{237.231}{equation}%
\TextAt{13}{240.599}{237.231}{y}%
\TextAt{21}{246.736}{241.569}{p}%
\TextAt{9}{254.818}{237.231}{\Character{61}}%
\TextAt{13}{267.243}{237.231}{z}%
\TextAt{21}{273.214}{241.569}{p}%
\TextAt{16}{280.631}{237.231}{\Character{0}}%
\TextAt{13}{292.587}{237.231}{x}%
\TextAt{21}{299.239}{241.569}{p}%
\TextAt{8}{307.897}{237.231}{there}%
\TextAt{8}{337.809}{237.231}{turns}%
\TextAt{8}{368.438}{237.231}{into}%
\TextAt{8}{392.181}{237.231}{the}%
\TextAt{8}{412.337}{237.231}{identity}%
\TextAt{13}{196.729}{208.23}{y}%
\TextAt{16}{205.522}{208.23}{\Character{1}}%
\TextAt{13}{211.5}{208.23}{y}%
\TextAt{16}{219.629}{208.23}{\Character{1}}%
\TextAt{16}{224.925}{208.23}{\Character{1}}%
\TextAt{16}{230.233}{208.23}{\Character{1}}%
\TextAt{13}{235.569}{208.23}{y}%
\TextAt{9}{245.027}{208.23}{\Character{61}}%
\TextAt{9}{257.437}{208.23}{\Character{40}}%
\TextAt{13}{262.004}{208.23}{z}%
\TextAt{16}{270.632}{208.23}{\Character{0}}%
\TextAt{13}{282.587}{208.23}{x}%
\TextAt{9}{289.239}{208.23}{\Character{41}\Character{40}}%
\TextAt{13}{298.344}{208.23}{z}%
\TextAt{16}{306.971}{208.23}{\Character{0}}%
\TextAt{13}{318.926}{208.23}{\Character{16}}%
\TextAt{13}{324.928}{208.23}{x}%
\TextAt{9}{331.594}{208.23}{\Character{41}}%
\TextAt{16}{338.139}{208.23}{\Character{1}}%
\TextAt{16}{343.435}{208.23}{\Character{1}}%
\TextAt{16}{348.743}{208.23}{\Character{1}}%
\TextAt{9}{354.079}{208.23}{\Character{40}}%
\TextAt{13}{358.631}{208.23}{z}%
\TextAt{16}{367.258}{208.23}{\Character{0}}%
\TextAt{13}{379.213}{208.23}{\Character{16}}%
\TextAt{21}{385.229}{213.166}{p}%
\TextAt{3}{389.492}{213.166}{\Character{0}}%
\TextAt{17}{396.078}{213.166}{1}%
\TextAt{13}{400.811}{208.23}{x}%
\TextAt{9}{407.463}{208.23}{\Character{41}}%
\TextAt{13}{412.015}{208.23}{\Character{58}}%
\TextAt{8}{79.226}{179.229}{Thus,}%
\TextAt{8}{111.591}{179.229}{assuming}%
\TextAt{8}{161.778}{179.229}{the}%
\TextAt{8}{181.265}{179.229}{existence}%
\TextAt{8}{230.4}{179.229}{of}%
\TextAt{8}{243.06}{179.229}{a}%
\TextAt{8}{252.146}{179.229}{solution,}%
\TextAt{8}{299.155}{179.229}{one}%
\TextAt{8}{319.945}{179.229}{obtains}%
\TextAt{8}{360.317}{179.229}{two}%
\TextAt{8}{381.764}{179.229}{multiplicative}%
\TextAt{8}{454.917}{179.229}{decompositions}%
\TextAt{8}{79.226}{162.183}{of}%
\TextAt{8}{91.654}{162.183}{the}%
\TextAt{8}{110.908}{162.183}{same}%
\TextAt{8}{139.331}{162.183}{number}%
\TextAt{8}{181.025}{162.183}{in}%
\TextAt{18}{193.775}{162.183}{Z}%
\TextAt{9}{201.746}{162.183}{\Character{91}}%
\TextAt{13}{204.997}{162.183}{\Character{16}}%
\TextAt{9}{211.013}{162.183}{\Character{93}}%
\TextAt{8}{214.264}{162.183}{.}%
\TextAt{8}{222.412}{162.183}{One}%
\TextAt{8}{246.221}{162.183}{can}%
\TextAt{8}{266.778}{162.183}{show}%
\TextAt{8}{294.879}{162.183}{that}%
\TextAt{8}{319.345}{162.183}{this}%
\TextAt{8}{341.265}{162.183}{contradicts}%
\TextAt{8}{400.593}{162.183}{the}%
\TextAt{8}{419.847}{162.183}{unique}%
\TextAt{8}{456.986}{162.183}{factorization—}%
\TextAt{8}{79.226}{145.137}{provided}%
\TextAt{8}{127.955}{145.137}{that}%
\TextAt{8}{153.921}{145.137}{this}%
\TextAt{8}{177.341}{145.137}{holds}%
\TextAt{8}{208.568}{145.137}{in}%
\TextAt{8}{222.818}{145.137}{the}%
\TextAt{8}{243.572}{145.137}{ring}%
\TextAt{18}{268.235}{145.137}{Z}%
\TextAt{9}{276.205}{145.137}{\Character{91}}%
\TextAt{13}{279.457}{145.137}{\Character{16}}%
\TextAt{9}{285.472}{145.137}{\Character{93}}%
\TextAt{8}{288.724}{145.137}{.}%
\TextAt{8}{298.964}{145.137}{Supposing}%
\TextAt{8}{355.883}{145.137}{erroneously}%
\TextAt{8}{418.659}{145.137}{that}%
\TextAt{8}{444.625}{145.137}{this}%
\TextAt{8}{468.045}{145.137}{was}%
\TextAt{8}{491.142}{145.137}{the}%
\TextAt{8}{511.896}{145.137}{case}%
\TextAt{8}{79.226}{128.091}{in}%
\TextAt{8}{93.174}{128.091}{general—in}%
\TextAt{8}{155.252}{128.091}{other}%
\TextAt{8}{186.116}{128.091}{words}%
\TextAt{8}{219.969}{128.091}{that}%
\TextAt{8}{245.633}{128.091}{the}%
\TextAt{8}{266.084}{128.091}{class}%
\TextAt{8}{293.815}{128.091}{number}%
\TextAt{13}{336.706}{128.091}{h}%
\TextAt{21}{343.445}{126.298}{p}%
\TextAt{8}{352.397}{128.091}{of}%
\TextAt{8}{366.022}{128.091}{the}%
\TextAt{8}{386.473}{128.091}{field}%
\TextAt{18}{412.125}{128.091}{Q}%
\TextAt{9}{421.423}{128.091}{\Character{40}}%
\TextAt{13}{425.975}{128.091}{\Character{16}}%
\TextAt{9}{431.991}{128.091}{\Character{41}}%
\TextAt{8}{440.735}{128.091}{were}%
\TextAt{8}{468.013}{128.091}{always}%
\TextAt{8}{505.775}{128.091}{equal}%
\TextAt{8}{79.226}{111.045}{to}%
\TextAt{8}{95.459}{111.045}{1—some}%
\TextAt{8}{144.27}{111.045}{actually}%
\TextAt{8}{190.75}{111.045}{thought}%
\TextAt{8}{236.584}{111.045}{they}%
\TextAt{8}{264.844}{111.045}{had}%
\TextAt{8}{289.53}{111.045}{proved}%
\TextAt{8}{329.506}{111.045}{“Fermat’s}%
\TextAt{8}{385.298}{111.045}{Last}%
\TextAt{8}{413.463}{111.045}{Theorem”}%
\TextAt{8}{470.332}{111.045}{in}%
\TextAt{8}{485.907}{111.045}{this}%
\TextAt{8}{510.652}{111.045}{way.}%
\TextAt{26}{79.226}{93.999}{Kummer}%
\TextAt{8}{128.362}{93.999}{,}%
\TextAt{8}{136.699}{93.999}{however,}%
\TextAt{8}{186.019}{93.999}{did}%
\TextAt{8}{207.126}{93.999}{not}%
\TextAt{8}{228.891}{93.999}{fall}%
\TextAt{8}{249.675}{93.999}{into}%
\TextAt{8}{274.368}{93.999}{this}%
\TextAt{8}{298.141}{93.999}{trap.}%
\TextAt{8}{330.913}{93.999}{Instead,}%
\TextAt{8}{376.706}{93.999}{he}%
\TextAt{8}{393.258}{93.999}{proved}%
\TextAt{8}{432.262}{93.999}{that}%
\TextAt{8}{458.582}{93.999}{the}%
\TextAt{8}{479.689}{93.999}{arguments}%
\TextAt{8}{79.226}{76.953}{we}%
\TextAt{8}{96.943}{76.953}{have}%
\TextAt{8}{124.427}{76.953}{indicated}%
\TextAt{8}{175.643}{76.953}{can}%
\TextAt{8}{197.592}{76.953}{be}%
\TextAt{8}{214.005}{76.953}{salvaged}%
\TextAt{8}{261.072}{76.953}{if}%
\TextAt{8}{272.285}{76.953}{one}%
\TextAt{8}{294.234}{76.953}{only}%
\TextAt{8}{320.416}{76.953}{assumes}%
\TextAt{13}{365.964}{76.953}{p}%
\TextAt{18}{375.994}{76.953}{\Character{45}}%
\TextAt{13}{383.47}{76.953}{h}%
\TextAt{21}{390.209}{75.16}{p}%
\TextAt{8}{399.356}{76.953}{instead}%
\TextAt{8}{440.23}{76.953}{of}%
\TextAt{13}{454.05}{76.953}{h}%
\TextAt{21}{460.789}{75.16}{p}%
\TextAt{9}{469.704}{76.953}{\Character{61}}%
\TextAt{9}{482.95}{76.953}{1}%
\TextAt{8}{488.817}{76.953}{.}%
\TextAt{8}{498.734}{76.953}{In}%
\TextAt{8}{513.844}{76.953}{this}%
\TextAt{8}{300.14}{43.033}{35}%
\Destination{78.23}{809.815}{nav.page.35}%
\Destination{79.36}{776.155}{nav.print.37}%
\Hotlink{254.007}{398.255}{25.078}{11.047}{nav.cite.N136}%
}

% Page 48 (body page 36)
\Page{%
\TextAt{2}{79.226}{796.943}{I.}%
\TextAt{2}{89.83}{796.943}{Algebraic}%
\TextAt{2}{138.648}{796.943}{Integers}%
\Vector{q 0 G 0 g 0.299 w 7.226 18.729 m 460.769 18.729 l S Q}%
\TextAt{8}{79.226}{761.124}{case}%
\TextAt{8}{103.759}{761.124}{he}%
\TextAt{8}{119.122}{761.124}{called}%
\TextAt{8}{152.048}{761.124}{a}%
\TextAt{8}{161.565}{761.124}{prime}%
\TextAt{8}{194.491}{761.124}{number}%
\TextAt{13}{236.849}{761.124}{p}%
\TextAt{12}{246.383}{761.124}{regular}%
\TextAt{8}{287.999}{761.124}{,}%
\TextAt{8}{294.958}{761.124}{otherwise}%
\TextAt{12}{346.808}{761.124}{irregular}%
\TextAt{8}{397.666}{761.124}{.}%
\TextAt{8}{406.035}{761.124}{He}%
\TextAt{8}{423.67}{761.124}{even}%
\TextAt{8}{450.091}{761.124}{showed}%
\TextAt{8}{490.238}{761.124}{that}%
\TextAt{13}{515.369}{761.124}{p}%
\TextAt{8}{524.903}{761.124}{is}%
\TextAt{8}{79.226}{744.078}{regular}%
\TextAt{8}{118.342}{744.078}{if}%
\TextAt{8}{128.503}{744.078}{and}%
\TextAt{8}{150.702}{744.078}{only}%
\TextAt{8}{175.831}{744.078}{if}%
\TextAt{8}{185.991}{744.078}{the}%
\TextAt{8}{205.584}{744.078}{numerators}%
\TextAt{8}{266.555}{744.078}{of}%
\TextAt{8}{279.321}{744.078}{the}%
\TextAt{12}{298.915}{744.078}{Bernoulli}%
\TextAt{12}{356.847}{744.078}{numbers}%
\TextAt{13}{410.381}{744.078}{B}%
\TextAt{17}{419.274}{742.284}{2}%
\TextAt{13}{424.006}{744.078}{\Character{59}}%
\TextAt{13}{429.23}{744.078}{B}%
\TextAt{17}{438.143}{742.284}{4}%
\TextAt{13}{442.876}{744.078}{\Character{59}}%
\TextAt{13}{448.1}{744.078}{\Character{58}}%
\TextAt{13}{453.337}{744.078}{\Character{58}}%
\TextAt{13}{458.561}{744.078}{\Character{58}}%
\TextAt{13}{463.798}{744.078}{\Character{59}}%
\TextAt{13}{469.034}{744.078}{B}%
\TextAt{21}{477.99}{742.284}{p}%
\TextAt{3}{482.252}{742.284}{\Character{0}}%
\TextAt{17}{488.839}{742.284}{3}%
\TextAt{8}{496.905}{744.078}{are}%
\TextAt{8}{515.853}{744.078}{not}%
\TextAt{8}{79.226}{727.032}{divisible}%
\TextAt{8}{124.19}{727.032}{by}%
\TextAt{13}{139.505}{727.032}{p}%
\TextAt{8}{145.38}{727.032}{.}%
\TextAt{8}{153.514}{727.032}{Among}%
\TextAt{8}{193.217}{727.032}{the}%
\TextAt{8}{212.429}{727.032}{first}%
\TextAt{8}{235.611}{727.032}{25}%
\TextAt{8}{250.28}{727.032}{prime}%
\TextAt{8}{282.5}{727.032}{numbers}%
\TextAt{13}{328.766}{727.032}{\Character{60}}%
\TextAt{9}{341.192}{727.032}{100}%
\TextAt{8}{361.704}{727.032}{only}%
\TextAt{8}{386.452}{727.032}{three}%
\TextAt{8}{415.419}{727.032}{are}%
\TextAt{8}{433.986}{727.032}{irregular:37,59,}%
\TextAt{8}{513.904}{727.032}{and}%
\TextAt{8}{79.226}{709.986}{67.}%
\TextAt{8}{99.272}{709.986}{We}%
\TextAt{8}{119.06}{709.986}{still}%
\TextAt{8}{141.525}{709.986}{do}%
\TextAt{8}{157.428}{709.986}{not}%
\TextAt{8}{177.885}{709.986}{know}%
\TextAt{8}{208.098}{709.986}{today}%
\TextAt{8}{240.594}{709.986}{whether}%
\TextAt{8}{285.106}{709.986}{there}%
\TextAt{8}{314.661}{709.986}{are}%
\TextAt{8}{333.815}{709.986}{infinitely}%
\TextAt{8}{382.558}{709.986}{many}%
\TextAt{8}{414.074}{709.986}{regular}%
\TextAt{8}{453.397}{709.986}{prime}%
\TextAt{8}{486.204}{709.986}{numbers.}%
\TextAt{8}{97.159}{689.951}{The}%
\TextAt{8}{121.438}{689.951}{connection}%
\TextAt{8}{180.195}{689.951}{with}%
\TextAt{8}{207.08}{689.951}{Fermat’s}%
\TextAt{8}{255.651}{689.951}{last}%
\TextAt{8}{278.053}{689.951}{theorem}%
\TextAt{8}{323.803}{689.951}{has}%
\TextAt{8}{344.902}{689.951}{at}%
\TextAt{8}{359.437}{689.951}{last}%
\TextAt{8}{381.839}{689.951}{become}%
\TextAt{8}{424.002}{689.951}{obsolete.}%
\TextAt{8}{474.167}{689.951}{Following}%
\TextAt{8}{526.911}{689.951}{a}%
\TextAt{8}{79.226}{672.905}{surprising}%
\TextAt{8}{133.08}{672.905}{discovery}%
\TextAt{8}{183.622}{672.905}{by}%
\TextAt{8}{199.626}{672.905}{the}%
\TextAt{8}{219.527}{672.905}{mathematician}%
\TextAt{26}{299.276}{672.905}{Gerhard}%
\TextAt{26}{356.855}{672.905}{Frey}%
\TextAt{8}{386.575}{672.905}{,}%
\TextAt{8}{393.52}{672.905}{who}%
\TextAt{8}{417.977}{672.905}{established}%
\TextAt{8}{477.676}{672.905}{a}%
\TextAt{8}{487.177}{672.905}{link}%
\TextAt{8}{510.007}{672.905}{with}%
\TextAt{8}{79.226}{655.859}{the}%
\TextAt{8}{99.654}{655.859}{theory}%
\TextAt{8}{136.675}{655.859}{of}%
\TextAt{8}{150.277}{655.859}{elliptic}%
\TextAt{8}{188.912}{655.859}{curves,}%
\TextAt{8}{228.333}{655.859}{it}%
\TextAt{8}{240.308}{655.859}{was}%
\TextAt{26}{263.079}{655.859}{Kenneth}%
\TextAt{26}{320.384}{655.859}{Ribet}%
\TextAt{8}{354.229}{655.859}{,}%
\TextAt{8}{361.718}{655.859}{who}%
\TextAt{8}{386.7}{655.859}{managed}%
\TextAt{8}{436.406}{655.859}{to}%
\TextAt{8}{450.988}{655.859}{reduce}%
\TextAt{8}{488.32}{655.859}{Fermat’s}%
\TextAt{8}{79.226}{638.813}{statement}%
\TextAt{8}{133.816}{638.813}{to}%
\TextAt{8}{148.345}{638.813}{another,}%
\TextAt{8}{194.8}{638.813}{much}%
\TextAt{8}{226.233}{638.813}{more}%
\TextAt{8}{255.717}{638.813}{important}%
\TextAt{8}{311.228}{638.813}{conjecture,}%
\TextAt{8}{371.001}{638.813}{the}%
\TextAt{12}{391.376}{638.813}{Taniyama-Shimura-Weil}%
\TextAt{12}{79.226}{621.767}{Conjecture}%
\TextAt{8}{143.712}{621.767}{.}%
\TextAt{8}{152.241}{621.767}{This}%
\TextAt{8}{178.988}{621.767}{was}%
\TextAt{8}{201.514}{621.767}{proved}%
\TextAt{8}{239.594}{621.767}{in}%
\TextAt{8}{253.274}{621.767}{sufficient}%
\TextAt{8}{302.46}{621.767}{generality}%
\TextAt{8}{356.477}{621.767}{in}%
\TextAt{8}{370.157}{621.767}{1994}%
\TextAt{8}{397.513}{621.767}{by}%
\TextAt{26}{413.799}{621.767}{Andrew}%
\TextAt{26}{466.731}{621.767}{Wiles}%
\TextAt{8}{501.844}{621.767}{,}%
\TextAt{8}{509.027}{621.767}{after}%
\TextAt{8}{79.226}{604.721}{many}%
\TextAt{8}{111.658}{604.721}{years}%
\TextAt{8}{142.2}{604.721}{of}%
\TextAt{8}{156.089}{604.721}{work,}%
\TextAt{8}{188.66}{604.721}{and}%
\TextAt{8}{211.982}{604.721}{with}%
\TextAt{8}{239.201}{604.721}{a}%
\TextAt{8}{249.515}{604.721}{helping}%
\TextAt{8}{291.044}{604.721}{hand}%
\TextAt{8}{320.869}{604.721}{from}%
\TextAt{26}{349.069}{604.721}{Richard}%
\TextAt{26}{403.79}{604.721}{Taylor}%
\TextAt{8}{446.924}{604.721}{.}%
\TextAt{8}{457.049}{604.721}{See}%
\TextAt{8}{478.41}{604.721}{[144].}%
\TextAt{8}{512.613}{604.721}{The}%
\TextAt{8}{79.226}{587.675}{regular}%
\TextAt{8}{118.905}{587.675}{and}%
\TextAt{8}{141.668}{587.675}{irregular}%
\TextAt{8}{189.154}{587.675}{prime}%
\TextAt{8}{222.318}{587.675}{numbers}%
\TextAt{8}{269.529}{587.675}{do}%
\TextAt{8}{285.788}{587.675}{however}%
\TextAt{8}{330.667}{587.675}{continue}%
\TextAt{8}{377.496}{587.675}{to}%
\TextAt{8}{391.806}{587.675}{be}%
\TextAt{8}{407.73}{587.675}{important.}%
\TextAt{12}{79.226}{561.663}{Exercise}%
\TextAt{12}{132.427}{561.663}{I.6.1.}%
\TextAt{8}{166.021}{561.663}{How}%
\TextAt{8}{192.681}{561.663}{many}%
\TextAt{8}{224.553}{561.663}{integral}%
\TextAt{8}{267.161}{561.663}{ideals}%
\TextAt{14}{299.739}{561.663}{a}%
\TextAt{8}{309.62}{561.663}{are}%
\TextAt{8}{329.131}{561.663}{there}%
\TextAt{8}{359.042}{561.663}{with}%
\TextAt{8}{385.702}{561.663}{the}%
\TextAt{8}{405.859}{561.663}{given}%
\TextAt{8}{436.428}{561.663}{norm}%
\TextAt{14}{466.998}{561.663}{N}%
\TextAt{9}{476.946}{561.663}{\Character{40}}%
\TextAt{14}{481.498}{561.663}{a}%
\TextAt{9}{487.481}{561.663}{\Character{41}}%
\TextAt{9}{495.336}{561.663}{\Character{61}}%
\TextAt{13}{507.78}{561.663}{n}%
\TextAt{8}{514.767}{561.663}{?}%
\TextAt{12}{79.226}{529.673}{Exercise}%
\TextAt{12}{131.826}{529.673}{I.6.2.}%
\TextAt{8}{164.898}{529.673}{Show}%
\ErrataAt{165.7}{529.674}{369.3}{Which of $5,8,11,-3,-4,-7,-8,-11$ occur as quadratic discriminants? Show that the}%
\EditionChange{\BodyAt{436.229}{529.655}{11}}{\BodyAt{436.229}{529.655}{13}}%
\ErrataAt{79.19}{512.628}{455.81}{associated quadratic fields have class number $1$.}%
\TextAt{12}{79.226}{480.637}{Exercise}%
\TextAt{12}{132.655}{480.637}{I.6.3.}%
\TextAt{8}{166.447}{480.637}{Show}%
\TextAt{8}{197.537}{480.637}{that}%
\TextAt{8}{223.104}{480.637}{in}%
\TextAt{8}{236.955}{480.637}{every}%
\TextAt{8}{268.046}{480.637}{ideal}%
\TextAt{8}{296.207}{480.637}{class}%
\TextAt{8}{323.842}{480.637}{of}%
\TextAt{8}{337.37}{480.637}{an}%
\TextAt{8}{353.828}{480.637}{algebraic}%
\TextAt{8}{403.46}{480.637}{number}%
\TextAt{8}{446.255}{480.637}{field}%
\TextAt{13}{471.81}{480.637}{K}%
\TextAt{8}{486.723}{480.637}{of}%
\TextAt{8}{500.251}{480.637}{degree}%
\TextAt{13}{79.226}{463.591}{n}%
\TextAt{8}{86.214}{463.591}{,}%
\TextAt{8}{93.363}{463.591}{there}%
\TextAt{8}{123.275}{463.591}{exists}%
\TextAt{8}{155.59}{463.591}{an}%
\TextAt{8}{171.849}{463.591}{integral}%
\TextAt{8}{214.457}{463.591}{ideal}%
\TextAt{14}{242.42}{463.591}{a}%
\TextAt{8}{252.3}{463.591}{such}%
\TextAt{8}{278.698}{463.591}{that}%
\TextAt{14}{242.205}{430.577}{N}%
\TextAt{9}{252.153}{430.577}{\Character{40}}%
\TextAt{14}{256.706}{430.577}{a}%
\TextAt{9}{262.689}{430.577}{\Character{41}}%
\NativeMathAt{270.562}{430.577}{11.955}{\leq}%
\TextAt{13}{285.569}{438.665}{n}%
\TextAt{9}{292.557}{438.665}{\Character{33}}%
\Vector{q 0 G 0 g 0.478 w 212.377 -340.434 m 225.001 -340.434 l S Q}%
\TextAt{13}{284.377}{422.376}{n}%
\TextAt{21}{291.364}{425.83}{n}%
\TextAt{20}{300.189}{447.434}{\Character{18}}%
\TextAt{9}{310.793}{438.665}{4}%
\Vector{q 0 G 0 g 0.478 w 238.185 -340.434 m 245.254 -340.434 l S Q}%
\TextAt{13}{310.185}{422.376}{\Character{25}}%
\TextAt{20}{318.449}{447.434}{\Character{19}}%
\TextAt{21}{327.25}{444.757}{s}%
\TextAt{20}{333.656}{441.382}{\Character{112}}%
\Vector{q 0 G 0 g 0.478 w 273.612 -332.379 m 294.539 -332.379 l S Q}%
\TextAt{16}{345.612}{430.577}{\Character{106}}%
\TextAt{13}{348.932}{430.577}{d}%
\TextAt{21}{355.015}{428.784}{K}%
\TextAt{16}{363.218}{430.577}{\Character{106}}%
\TextAt{13}{366.539}{430.577}{\Character{58}}%
\TextAt{12}{97.159}{394.313}{Hint:}%
\TextAt{8}{132.698}{394.313}{Using}%
\TextAt{8}{165.598}{394.313}{exercise}%
\TextAt{8}{208.9}{394.313}{3,}%
\TextAt{8}{221.907}{394.313}{§5,}%
\TextAt{8}{240.593}{394.313}{proceed}%
\TextAt{8}{283.835}{394.313}{as}%
\TextAt{8}{298.205}{394.313}{in}%
\TextAt{8}{311.858}{394.313}{the}%
\TextAt{8}{332.014}{394.313}{proof}%
\TextAt{8}{362.584}{394.313}{of}%
\TextAt{8}{375.914}{394.313}{(6.3).}%
\TextAt{12}{79.226}{362.323}{Exercise}%
\TextAt{12}{133.104}{362.323}{I.6.4.}%
\TextAt{8}{167.286}{362.323}{Show}%
\TextAt{8}{198.767}{362.323}{that}%
\TextAt{8}{224.725}{362.323}{the}%
\TextAt{8}{245.47}{362.323}{absolute}%
\TextAt{8}{292.301}{362.323}{value}%
\TextAt{8}{323.136}{362.323}{of}%
\TextAt{8}{337.055}{362.323}{the}%
\TextAt{8}{357.8}{362.323}{discriminant}%
\TextAt{16}{425.768}{362.323}{\Character{106}}%
\TextAt{13}{429.089}{362.323}{d}%
\TextAt{21}{435.171}{360.53}{K}%
\TextAt{16}{443.374}{362.323}{\Character{106}}%
\TextAt{8}{451.181}{362.323}{is}%
\TextAt{13}{463.534}{362.323}{\Character{62}}%
\TextAt{9}{476.962}{362.323}{1}%
\TextAt{8}{487.301}{362.323}{for}%
\TextAt{8}{505.775}{362.323}{every}%
\TextAt{8}{79.226}{345.277}{algebraic}%
\TextAt{8}{128.989}{345.277}{number}%
\TextAt{8}{171.914}{345.277}{field}%
\TextAt{13}{197.599}{345.277}{K}%
\NativeMathAt{212.297}{345.277}{11.955}{\neq}%
\TextAt{18}{225.282}{345.277}{Q}%
\TextAt{8}{238.807}{345.277}{(Minkowski’s}%
\TextAt{8}{309.826}{345.277}{theorem}%
\TextAt{8}{355.68}{345.277}{on}%
\TextAt{8}{372.267}{345.277}{the}%
\TextAt{8}{392.752}{345.277}{discriminant,}%
\TextAt{8}{463.794}{345.277}{see}%
\TextAt{8}{483.035}{345.277}{chap.}%
\TextAt{8}{516.212}{345.277}{III,}%
\TextAt{8}{79.226}{328.231}{(2.17)).}%
\TextAt{12}{79.226}{296.241}{Exercise}%
\TextAt{12}{132.748}{296.241}{I.6.5.}%
\TextAt{8}{166.621}{296.241}{Show}%
\TextAt{8}{197.793}{296.241}{that}%
\TextAt{8}{223.441}{296.241}{the}%
\TextAt{8}{243.876}{296.241}{absolute}%
\TextAt{8}{290.398}{296.241}{value}%
\TextAt{8}{320.924}{296.241}{of}%
\TextAt{8}{334.533}{296.241}{the}%
\TextAt{8}{354.969}{296.241}{discriminant}%
\TextAt{16}{422.628}{296.241}{\Character{106}}%
\TextAt{13}{425.949}{296.241}{d}%
\TextAt{21}{432.031}{294.448}{K}%
\TextAt{16}{440.234}{296.241}{\Character{106}}%
\TextAt{8}{447.731}{296.241}{tends}%
\TextAt{8}{479.285}{296.241}{to}%
\NativeMathAt{493.875}{296.241}{11.955}{\infty}%
\TextAt{8}{510.007}{296.241}{with}%
\TextAt{8}{79.226}{279.195}{the}%
\TextAt{8}{99.383}{279.195}{degree}%
\TextAt{13}{135.798}{279.195}{n}%
\TextAt{8}{146.683}{279.195}{of}%
\TextAt{8}{160.013}{279.195}{the}%
\TextAt{8}{180.169}{279.195}{field.}%
\TextAt{12}{79.226}{247.205}{Exercise}%
\TextAt{12}{132.98}{247.205}{I.6.6.}%
\TextAt{8}{167.055}{247.205}{Let}%
\TextAt{14}{188.505}{247.205}{a}%
\TextAt{8}{198.866}{247.205}{be}%
\TextAt{8}{215.272}{247.205}{an}%
\TextAt{8}{232.012}{247.205}{integral}%
\TextAt{8}{275.101}{247.205}{ideal}%
\TextAt{8}{303.545}{247.205}{of}%
\TextAt{13}{317.356}{247.205}{K}%
\TextAt{8}{332.552}{247.205}{and}%
\TextAt{14}{355.795}{247.205}{a}%
\TextAt{21}{361.778}{251.544}{m}%
\TextAt{9}{373.907}{247.205}{\Character{61}}%
\TextAt{9}{387.141}{247.205}{\Character{40}}%
\TextAt{13}{391.705}{247.205}{a}%
\TextAt{9}{397.85}{247.205}{\Character{41}}%
\TextAt{8}{402.402}{247.205}{.}%
\TextAt{8}{412.293}{247.205}{Show}%
\TextAt{8}{443.666}{247.205}{that}%
\TextAt{14}{469.516}{247.205}{a}%
\TextAt{8}{479.877}{247.205}{becomes}%
\TextAt{8}{526.911}{247.205}{a}%
\TextAt{8}{79.226}{230.159}{principal}%
\TextAt{8}{128.003}{230.159}{ideal}%
\TextAt{8}{155.966}{230.159}{in}%
\TextAt{8}{169.619}{230.159}{the}%
\TextAt{8}{189.776}{230.159}{field}%
\TextAt{13}{215.133}{230.159}{L}%
\TextAt{9}{226.418}{230.159}{\Character{61}}%
\TextAt{13}{238.843}{230.159}{K}%
\TextAt{9}{249.661}{230.159}{\Character{40}}%
\TextAt{22}{257.534}{234.18}{m}%
\TextAt{16}{257.43}{238.77}{\Character{112}}%
\Vector{q 0 G 0 g 0.478 w 195.393 -534.991 m 201.538 -534.991 l S Q}%
\TextAt{13}{267.393}{230.159}{a}%
\TextAt{9}{273.538}{230.159}{\Character{41}}%
\TextAt{8}{278.09}{230.159}{,}%
\TextAt{8}{285.24}{230.159}{in}%
\TextAt{8}{298.892}{230.159}{the}%
\TextAt{8}{319.049}{230.159}{sense}%
\TextAt{8}{349.08}{230.159}{that}%
\TextAt{14}{374.449}{230.159}{ao}%
\TextAt{21}{386.276}{228.366}{L}%
\TextAt{9}{395.854}{230.159}{\Character{61}}%
\TextAt{9}{408.264}{230.159}{\Character{40}}%
\TextAt{13}{412.832}{230.159}{\Character{11}}%
\TextAt{9}{420.354}{230.159}{\Character{41}}%
\TextAt{8}{424.906}{230.159}{.}%
\TextAt{12}{79.226}{198.169}{Exercise}%
\TextAt{12}{133.284}{198.169}{I.6.7.}%
\TextAt{8}{167.623}{198.169}{Show}%
\TextAt{8}{199.26}{198.169}{that,}%
\TextAt{8}{228.812}{198.169}{for}%
\TextAt{8}{247.442}{198.169}{every}%
\TextAt{8}{279.079}{198.169}{number}%
\TextAt{8}{322.42}{198.169}{field}%
\TextAt{13}{348.522}{198.169}{K}%
\TextAt{8}{359.34}{198.169}{,}%
\TextAt{8}{367.42}{198.169}{there}%
\TextAt{8}{398.077}{198.169}{exists}%
\TextAt{8}{431.137}{198.169}{a}%
\TextAt{8}{441.637}{198.169}{finite}%
\TextAt{8}{472.294}{198.169}{extension}%
\TextAt{13}{524.805}{198.169}{L}%
\TextAt{8}{79.226}{181.123}{such}%
\TextAt{8}{105.623}{181.123}{that}%
\TextAt{8}{130.992}{181.123}{every}%
\TextAt{8}{161.884}{181.123}{ideal}%
\TextAt{8}{189.847}{181.123}{of}%
\TextAt{13}{203.177}{181.123}{K}%
\TextAt{8}{217.892}{181.123}{becomes}%
\TextAt{8}{264.446}{181.123}{a}%
\TextAt{8}{274.201}{181.123}{principal}%
\TextAt{8}{322.978}{181.123}{ideal.}%
\TextAt{15}{79.226}{128.26}{§}%
\TextAt{15}{93.785}{128.26}{7}%
\TextAt{15}{120.693}{128.26}{Dirichlet’s}%
\TextAt{15}{214.311}{128.26}{Unit}%
\TextAt{15}{259.346}{128.26}{Theorem}%
\TextAt{8}{79.226}{96.374}{After}%
\TextAt{8}{110.781}{96.374}{considering}%
\TextAt{8}{172.977}{96.374}{the}%
\TextAt{8}{194.131}{96.374}{ideal}%
\TextAt{8}{223.091}{96.374}{class}%
\TextAt{8}{251.526}{96.374}{group}%
\TextAt{13}{285.699}{96.374}{C}%
\TextAt{13}{294.917}{96.374}{l}%
\TextAt{21}{298.419}{94.58}{K}%
\TextAt{8}{306.622}{96.374}{,}%
\TextAt{8}{315.018}{96.374}{we}%
\TextAt{8}{333.243}{96.374}{now}%
\TextAt{8}{358.629}{96.374}{turn}%
\TextAt{8}{385.641}{96.374}{to}%
\TextAt{8}{400.948}{96.374}{the}%
\TextAt{8}{422.102}{96.374}{second}%
\TextAt{8}{460.878}{96.374}{main}%
\TextAt{8}{491.142}{96.374}{problem}%
\TextAt{8}{79.226}{79.328}{posed}%
\TextAt{8}{112.552}{79.328}{by}%
\TextAt{8}{129.237}{79.328}{the}%
\TextAt{8}{149.819}{79.328}{ring}%
\TextAt{14}{174.311}{79.328}{o}%
\TextAt{21}{180.155}{77.534}{K}%
\TextAt{8}{192.681}{79.328}{of}%
\TextAt{8}{206.436}{79.328}{integers}%
\TextAt{8}{250.176}{79.328}{of}%
\TextAt{8}{263.931}{79.328}{an}%
\TextAt{8}{280.616}{79.328}{algebraic}%
\TextAt{8}{330.476}{79.328}{number}%
\TextAt{8}{373.498}{79.328}{field}%
\TextAt{13}{399.281}{79.328}{K}%
\TextAt{8}{410.099}{79.328}{,}%
\TextAt{8}{417.78}{79.328}{the}%
\TextAt{8}{438.362}{79.328}{group}%
\TextAt{8}{471.963}{79.328}{of}%
\TextAt{8}{485.719}{79.328}{units}%
\TextAt{14}{515.471}{79.328}{o}%
\TextAt{3}{521.315}{83.666}{\Character{3}}%
\TextAt{21}{521.315}{76.307}{K}%
\TextAt{8}{529.518}{79.328}{.}%
\TextAt{8}{300.14}{43.033}{36}%
\Destination{79.36}{548.655}{errata.2}%
\Destination{78.23}{809.815}{nav.page.36}%
\Destination{79.23}{154.675}{nav.section.1.7}%
\Destination{79.226}{575.655}{nav.ex.1.6.1}%
\Destination{79.226}{543.655}{nav.ex.1.6.2}%
\Destination{79.226}{494.655}{nav.ex.1.6.3}%
\Destination{79.226}{376.355}{nav.ex.1.6.4}%
\Destination{79.226}{310.255}{nav.ex.1.6.5}%
\Destination{79.226}{261.255}{nav.ex.1.6.6}%
\Destination{79.226}{212.155}{nav.ex.1.6.7}%
\Destination{79.36}{776.155}{nav.print.38}%
\Destination{79.23}{154.675}{outline.12}%
\Hotlink{477.91}{602.555}{25.078}{11.047}{nav.cite.N144}%
\Hotlink{208.4}{392.155}{6.858}{11.047}{nav.ex.1.5.3}%
\Hotlink{227.086}{392.155}{6.858}{11.047}{nav.section.1.5}%
\Hotlink{515.712}{343.055}{14.294}{11.047}{nav.chapter.3}%
\Hotlink{375.414}{392.155}{25.078}{11.047}{nav.number.1.6.3}%
\Hotlink{80.397}{326}{27.58}{11.047}{nav.number.3.2.17}%
}

% Page 49 (body page 37)
\Page{%
\TextAt{2}{79.226}{796.943}{I.}%
\TextAt{2}{89.83}{796.943}{Algebraic}%
\TextAt{2}{138.648}{796.943}{Integers}%
\Vector{q 0 G 0 g 0.299 w 7.226 18.729 m 460.769 18.729 l S Q}%
\TextAt{8}{79.226}{761.124}{It}%
\TextAt{8}{91.744}{761.124}{contains}%
\TextAt{8}{137.509}{761.124}{the}%
\TextAt{8}{157.511}{761.124}{finite}%
\TextAt{8}{187.268}{761.124}{group}%
\TextAt{13}{220.288}{761.124}{\Character{22}}%
\TextAt{9}{227.331}{761.124}{\Character{40}}%
\TextAt{13}{231.884}{761.124}{K}%
\TextAt{9}{242.701}{761.124}{\Character{41}}%
\TextAt{8}{250.996}{761.124}{of}%
\TextAt{8}{264.172}{761.124}{the}%
\TextAt{8}{284.173}{761.124}{roots}%
\TextAt{8}{313.679}{761.124}{of}%
\TextAt{8}{326.854}{761.124}{unity}%
\TextAt{8}{357.269}{761.124}{that}%
\TextAt{8}{382.483}{761.124}{lie}%
\TextAt{8}{397.93}{761.124}{in}%
\TextAt{13}{411.428}{761.124}{K}%
\TextAt{8}{422.246}{761.124}{,}%
\TextAt{8}{429.271}{761.124}{but}%
\TextAt{8}{450.576}{761.124}{in}%
\TextAt{8}{464.074}{761.124}{general}%
\TextAt{8}{504.244}{761.124}{is}%
\TextAt{8}{515.853}{761.124}{not}%
\TextAt{8}{79.226}{744.078}{itself}%
\TextAt{8}{107.522}{744.078}{finite.}%
\TextAt{8}{141.968}{744.078}{Its}%
\TextAt{8}{159.204}{744.078}{size}%
\TextAt{8}{181.319}{744.078}{is}%
\TextAt{8}{193.033}{744.078}{in}%
\TextAt{8}{206.635}{744.078}{fact}%
\TextAt{8}{229.67}{744.078}{determined}%
\TextAt{8}{290.747}{744.078}{by}%
\TextAt{8}{306.955}{744.078}{the}%
\TextAt{8}{327.062}{744.078}{number}%
\TextAt{13}{369.608}{744.078}{r}%
\TextAt{8}{379.055}{744.078}{of}%
\TextAt{8}{392.335}{744.078}{real}%
\TextAt{8}{415.047}{744.078}{embeddings}%
\TextAt{13}{478.79}{744.078}{\Character{26}}%
\TextAt{9}{488.148}{744.078}{\Character{58}}%
\TextAt{13}{494.721}{744.078}{K}%
\NativeMathAt{508.859}{744.078}{11.955}{\rightarrow}%
\TextAt{18}{524.135}{744.078}{R}%
\TextAt{8}{79.226}{727.032}{and}%
\TextAt{8}{102.441}{727.032}{the}%
\TextAt{8}{123.05}{727.032}{number}%
\TextAt{13}{166.099}{727.032}{s}%
\TextAt{8}{175.963}{727.032}{of}%
\TextAt{8}{189.746}{727.032}{pairs}%
\TextAt{13}{218.879}{727.032}{\Character{27}}%
\TextAt{13}{225.311}{727.032}{\Character{59}}%
\TextAt{9}{231.17}{727.032}{\Character{22}}%
\TextAt{13}{230.555}{727.032}{\Character{27}}%
\TextAt{9}{241.729}{727.032}{\Character{58}}%
\TextAt{13}{249.073}{727.032}{K}%
\NativeMathAt{263.983}{727.032}{11.955}{\rightarrow}%
\TextAt{18}{280.03}{727.032}{C}%
\TextAt{8}{293.014}{727.032}{of}%
\TextAt{8}{306.797}{727.032}{complex}%
\TextAt{8}{353.098}{727.032}{conjugate}%
\TextAt{8}{406.559}{727.032}{embeddings.}%
\TextAt{8}{476.261}{727.032}{In}%
\TextAt{8}{491.335}{727.032}{order}%
\TextAt{8}{522.357}{727.032}{to}%
\TextAt{8}{79.226}{709.986}{describe}%
\TextAt{8}{124.476}{709.986}{the}%
\TextAt{8}{144.633}{709.986}{group,}%
\TextAt{8}{181.06}{709.986}{we}%
\TextAt{8}{198.288}{709.986}{use}%
\TextAt{8}{218.504}{709.986}{the}%
\TextAt{8}{238.66}{709.986}{diagram}%
\TextAt{8}{284.198}{709.986}{which}%
\TextAt{8}{317.684}{709.986}{was}%
\TextAt{8}{340.184}{709.986}{set}%
\TextAt{8}{358.451}{709.986}{up}%
\TextAt{8}{375.356}{709.986}{in}%
\TextAt{8}{389.009}{709.986}{§5:}%
\TextAt{13}{223.003}{680.256}{K}%
\TextAt{3}{233.821}{684.595}{\Character{3}}%
\TextAt{13}{291.897}{680.256}{K}%
\TextAt{3}{302.714}{684.595}{\Character{3}}%
\TextAt{25}{301.856}{677.192}{R}%
\TextAt{9}{361.455}{680.256}{\Character{91}}%
\TextAt{20}{364.706}{689.223}{\Character{81}}%
\TextAt{21}{375.997}{676.769}{\Character{28}}%
\TextAt{18}{383.166}{680.256}{R}%
\TextAt{9}{391.801}{680.256}{\Character{93}}%
\TextAt{17}{395.052}{686.068}{\Character{43}}%
\TextAt{18}{223.763}{627.495}{Q}%
\TextAt{3}{233.061}{631.834}{\Character{3}}%
\TextAt{18}{293.321}{627.495}{R}%
\TextAt{3}{301.955}{631.834}{\Character{3}}%
\TextAt{18}{375.853}{627.495}{R}%
\TextAt{13}{384.487}{627.495}{\Character{58}}%
\NativeDiagram{%
\path[draw=black,line width=0.478bp] (171.905,-90.754) -- (214.068,-90.754);%
\path[draw=black,line width=0.478bp,line cap=round,line join=round] (211.82,-87.885) .. controls (212.274,-89.607) and (213.287,-90.42) .. (214.307,-90.754) .. controls (213.287,-91.089) and (212.274,-91.902) .. (211.82,-93.624);%
}%
\TextAt{21}{263.284}{687.611}{j}%
\NativeDiagram{%
\path[draw=black,line width=0.478bp] (158.778,-98.322) -- (158.778,-133.212);%
\path[draw=black,line width=0.478bp,line cap=round,line join=round] (161.647,-130.965) .. controls (159.926,-131.419) and (159.113,-132.431) .. (158.778,-133.451) .. controls (158.443,-132.431) and (157.63,-131.419) .. (155.909,-130.965);%
}%
\TextAt{21}{207.28}{656.887}{N}%
\TextAt{22}{214.002}{655.149}{K}%
\TextAt{23}{220.547}{655.149}{\Character{106}}%
\TextAt{29}{222.816}{655.149}{Q}%
\NativeDiagram{%
\path[draw=black,line width=0.478bp] (241.464,-90.754) -- (283.626,-90.754);%
\path[draw=black,line width=0.478bp,line cap=round,line join=round] (281.379,-87.885) .. controls (281.833,-89.607) and (282.845,-90.42) .. (283.865,-90.754) .. controls (282.845,-91.089) and (281.833,-91.902) .. (281.379,-93.624);%
}%
\TextAt{21}{333.054}{686.061}{\Character{96}}%
\NativeDiagram{%
\path[draw=black,line width=0.478bp] (228.005,-101.387) -- (228.005,-133.212);%
\path[draw=black,line width=0.478bp,line cap=round,line join=round] (230.874,-130.965) .. controls (229.152,-131.419) and (228.339,-132.431) .. (228.005,-133.451) .. controls (227.67,-132.431) and (226.857,-131.419) .. (225.136,-130.965);%
}%
\TextAt{21}{302.82}{653.738}{N}%
\NativeDiagram{%
\path[draw=black,line width=0.478bp] (309.798,-101.809) -- (309.798,-133.212);%
\path[draw=black,line width=0.478bp,line cap=round,line join=round] (312.667,-130.965) .. controls (310.945,-131.419) and (310.132,-132.431) .. (309.798,-133.451) .. controls (309.463,-132.431) and (308.65,-131.419) .. (306.928,-130.965);%
}%
\TextAt{7}{384.613}{653.556}{Tr}%
\NativeDiagram{%
\path[draw=black,line width=0.478bp] (171.146,-143.516) -- (215.491,-143.516);%
\path[draw=black,line width=0.478bp,line cap=round,line join=round] (213.244,-140.647) .. controls (213.698,-142.369) and (214.71,-143.182) .. (215.73,-143.516) .. controls (214.71,-143.851) and (213.698,-144.664) .. (213.244,-146.386);%
\path[draw=black,line width=0.478bp] (240.04,-143.516) -- (298.024,-143.516);%
\path[draw=black,line width=0.478bp,line cap=round,line join=round] (295.777,-140.647) .. controls (296.231,-142.369) and (297.244,-143.182) .. (298.264,-143.516) .. controls (297.244,-143.851) and (296.231,-144.664) .. (295.777,-146.386);%
}%
\TextAt{21}{339.542}{633.299}{\Character{96}}%
\TextAt{8}{79.226}{600.723}{In}%
\TextAt{8}{93.847}{600.723}{the}%
\TextAt{8}{114.004}{600.723}{upper}%
\TextAt{8}{147.49}{600.723}{part}%
\TextAt{8}{172.859}{600.723}{of}%
\TextAt{8}{186.189}{600.723}{the}%
\TextAt{8}{206.345}{600.723}{diagram}%
\TextAt{8}{251.883}{600.723}{we}%
\TextAt{8}{269.11}{600.723}{consider}%
\TextAt{8}{314.695}{600.723}{the}%
\TextAt{8}{334.852}{600.723}{subgroups}%
\TextAt{14}{120.629}{572.69}{o}%
\TextAt{3}{126.473}{577.626}{\Character{3}}%
\TextAt{21}{126.473}{569.734}{K}%
\TextAt{9}{137.997}{572.69}{\Character{61}}%
\TextAt{16}{150.422}{572.69}{\Character{102}}%
\TextAt{13}{156.4}{572.69}{\Character{34}}%
\NativeMathAt{165.199}{572.69}{11.955}{\in}%
\TextAt{14}{176.49}{572.69}{o}%
\TextAt{21}{182.334}{570.897}{K}%
\TextAt{16}{193.858}{572.69}{\Character{106}}%
\TextAt{13}{200.499}{572.69}{N}%
\TextAt{21}{209.873}{570.83}{K}%
\TextAt{3}{217.578}{570.83}{\Character{106}}%
\TextAt{25}{219.93}{570.83}{Q}%
\TextAt{9}{226.627}{572.69}{\Character{40}}%
\TextAt{13}{231.179}{572.69}{\Character{34}}%
\TextAt{9}{236.657}{572.69}{\Character{41}}%
\TextAt{9}{244.512}{572.69}{\Character{61}}%
\NativeMathAt{256.956}{572.69}{11.955}{\pm}%
\TextAt{9}{266.255}{572.69}{1}%
\TextAt{16}{272.108}{572.69}{\Character{103}}%
\TextAt{13}{278.085}{572.69}{\Character{59}}%
\TextAt{8}{346.673}{572.69}{the}%
\TextAt{8}{366.83}{572.69}{group}%
\TextAt{8}{400.005}{572.69}{of}%
\TextAt{8}{413.335}{572.69}{units}%
\TextAt{13}{438.764}{572.69}{\Character{59}}%
\TextAt{13}{126.781}{552.655}{S}%
\TextAt{9}{137.997}{552.655}{\Character{61}}%
\TextAt{16}{150.422}{552.655}{\Character{102}}%
\TextAt{13}{156.4}{552.655}{y}%
\NativeMathAt{165.858}{552.655}{11.955}{\in}%
\TextAt{13}{177.149}{552.655}{K}%
\TextAt{3}{187.966}{557.591}{\Character{3}}%
\TextAt{25}{187.108}{549.7}{R}%
\TextAt{16}{196.683}{552.655}{\Character{106}}%
\TextAt{13}{203.325}{552.655}{N}%
\TextAt{9}{213.947}{552.655}{\Character{40}}%
\TextAt{13}{218.499}{552.655}{y}%
\TextAt{9}{224.636}{552.655}{\Character{41}}%
\TextAt{9}{232.491}{552.655}{\Character{61}}%
\NativeMathAt{244.935}{552.655}{11.955}{\pm}%
\TextAt{9}{254.233}{552.655}{1}%
\TextAt{16}{260.086}{552.655}{\Character{103}}%
\TextAt{13}{266.064}{552.655}{\Character{59}}%
\TextAt{8}{346.673}{552.655}{the}%
\TextAt{8}{366.83}{552.655}{“norm-one}%
\TextAt{8}{424.382}{552.655}{surface”}%
\TextAt{13}{465.412}{552.655}{\Character{59}}%
\TextAt{13}{124.054}{520.974}{H}%
\TextAt{9}{137.997}{520.974}{\Character{61}}%
\TextAt{20}{150.422}{541.418}{\Character{40}}%
\TextAt{13}{160.053}{520.974}{x}%
\NativeMathAt{170.026}{520.974}{11.955}{\in}%
\TextAt{20}{181.317}{541.418}{\Character{34}}%
\TextAt{20}{188.291}{532.332}{\Character{89}}%
\TextAt{21}{193.59}{507.226}{\Character{28}}%
\TextAt{18}{205.559}{520.974}{R}%
\TextAt{20}{214.194}{541.418}{\Character{35}}%
\TextAt{17}{221.168}{538.741}{\Character{43}}%
\TextAt{20}{232.15}{541.896}{\Character{12}}%
\TextAt{20}{232.15}{534.723}{\Character{12}}%
\TextAt{20}{232.15}{527.55}{\Character{12}}%
\TextAt{20}{232.15}{520.376}{\Character{12}}%
\TextAt{20}{232.15}{513.203}{\Character{12}}%
\TextAt{8}{240.032}{520.974}{Tr}%
\TextAt{9}{253.039}{520.974}{\Character{40}}%
\TextAt{13}{257.592}{520.974}{x}%
\TextAt{9}{264.244}{520.974}{\Character{41}}%
\TextAt{9}{272.099}{520.974}{\Character{61}}%
\TextAt{9}{284.52}{520.974}{0}%
\TextAt{20}{290.395}{541.418}{\Character{41}}%
\TextAt{13}{302.018}{520.974}{\Character{59}}%
\TextAt{8}{346.673}{520.974}{the}%
\TextAt{8}{366.83}{520.974}{“trace-zero}%
\TextAt{8}{426.33}{520.974}{hyperplane”}%
\TextAt{13}{488.115}{520.974}{\Character{58}}%
\TextAt{8}{79.226}{485.75}{We}%
\TextAt{8}{99.371}{485.75}{obtain}%
\TextAt{8}{135.798}{485.75}{the}%
\TextAt{8}{155.954}{485.75}{homomorphisms}%
\TextAt{14}{271.118}{468.704}{o}%
\TextAt{3}{276.963}{473.64}{\Character{3}}%
\TextAt{21}{276.963}{465.748}{K}%
\TextAt{21}{291.816}{475.968}{j}%
\NativeArrowAt{288.486}{468.704}{11.955}{11.955}{\longrightarrow}%
\TextAt{13}{303.762}{468.704}{S}%
\TextAt{21}{318.521}{475.481}{\Character{96}}%
\NativeArrowAt{314.978}{468.704}{11.955}{11.955}{\longrightarrow}%
\TextAt{13}{330.255}{468.704}{H}%
\TextAt{8}{79.226}{445.597}{and}%
\TextAt{8}{101.989}{445.597}{the}%
\TextAt{8}{122.145}{445.597}{composite}%
\TextAt{13}{177.163}{445.597}{\Character{21}}%
\TextAt{9}{187.312}{445.597}{\Character{58}\Character{61}}%
\TextAt{13}{202.989}{445.597}{\Character{96}}%
\TextAt{16}{210.558}{445.597}{\Character{14}}%
\TextAt{13}{219.192}{445.597}{j}%
\TextAt{9}{228.029}{445.597}{\Character{58}}%
\TextAt{14}{234.601}{445.597}{o}%
\TextAt{3}{240.445}{449.935}{\Character{3}}%
\TextAt{21}{240.445}{442.576}{K}%
\NativeMathAt{251.969}{445.597}{11.955}{\rightarrow}%
\TextAt{13}{267.245}{445.597}{H}%
\TextAt{8}{277.868}{445.597}{.}%
\TextAt{8}{286.316}{445.597}{The}%
\TextAt{8}{310.37}{445.597}{image}%
\TextAt{8}{344.191}{445.597}{will}%
\TextAt{8}{366.296}{445.597}{be}%
\TextAt{8}{382.22}{445.597}{denoted}%
\TextAt{8}{426.442}{445.597}{by}%
\TextAt{9}{264.896}{417.564}{\Character{0}}%
\TextAt{9}{275.512}{417.564}{\Character{61}}%
\TextAt{13}{287.956}{417.564}{\Character{21}}%
\TextAt{9}{294.784}{417.564}{\Character{40}}%
\TextAt{14}{299.337}{417.564}{o}%
\TextAt{3}{305.181}{422.5}{\Character{3}}%
\TextAt{21}{305.181}{414.608}{K}%
\TextAt{9}{313.384}{417.564}{\Character{41}}%
\NativeMathAt{321.257}{417.564}{11.955}{\subseteq}%
\TextAt{13}{333.876}{417.564}{H}%
\TextAt{13}{343.835}{417.564}{\Character{59}}%
\TextAt{8}{79.226}{389.53}{and}%
\TextAt{8}{101.989}{389.53}{we}%
\TextAt{8}{119.216}{389.53}{obtain}%
\TextAt{8}{155.644}{389.53}{the}%
\TextAt{12}{79.226}{363.795}{Proposition}%
\TextAt{12}{151.937}{363.795}{(7.1).}%
\TextAt{8}{187.253}{363.795}{The}%
\TextAt{8}{211.307}{363.795}{sequence}%
\TextAt{9}{227.829}{335.762}{1}%
\NativeArrowAt{237.003}{335.762}{11.955}{19.26}{\longrightarrow}%
\TextAt{13}{259.585}{335.762}{\Character{22}}%
\TextAt{9}{266.628}{335.762}{\Character{40}}%
\TextAt{13}{271.18}{335.762}{K}%
\TextAt{9}{281.998}{335.762}{\Character{41}}%
\NativeArrowAt{289.871}{335.762}{11.955}{19.26}{\longrightarrow}%
\TextAt{14}{312.453}{335.762}{o}%
\TextAt{3}{318.297}{340.698}{\Character{3}}%
\TextAt{21}{318.297}{332.806}{K}%
\TextAt{21}{332.623}{342.539}{\Character{21}}%
\NativeArrowAt{329.821}{335.762}{11.955}{11.955}{\longrightarrow}%
\TextAt{9}{345.097}{335.762}{\Character{0}}%
\NativeArrowAt{355.732}{335.762}{11.955}{19.26}{\longrightarrow}%
\TextAt{9}{378.314}{335.762}{0}%
\TextAt{8}{79.226}{307.728}{is}%
\TextAt{8}{90.99}{307.728}{exact.}%
\TextAt{19}{79.226}{273.303}{Proof.}%
\TextAt{8}{116}{273.303}{We}%
\TextAt{8}{136.269}{273.303}{have}%
\TextAt{8}{163.389}{273.303}{to}%
\TextAt{8}{177.823}{273.303}{show}%
\TextAt{8}{206.951}{273.303}{that}%
\TextAt{13}{232.445}{273.303}{\Character{22}}%
\TextAt{9}{239.488}{273.303}{\Character{40}}%
\TextAt{13}{244.04}{273.303}{K}%
\TextAt{9}{254.858}{273.303}{\Character{41}}%
\TextAt{8}{263.432}{273.303}{is}%
\TextAt{8}{275.32}{273.303}{the}%
\TextAt{8}{295.601}{273.303}{kernel}%
\TextAt{8}{330.193}{273.303}{of}%
\TextAt{13}{343.647}{273.303}{\Character{21}}%
\TextAt{8}{350.476}{273.303}{.}%
\TextAt{8}{359.298}{273.303}{For}%
\TextAt{13}{380.391}{273.303}{\Character{16}}%
\NativeMathAt{389.94}{273.303}{11.955}{\in}%
\TextAt{13}{401.443}{273.303}{\Character{22}}%
\TextAt{9}{408.486}{273.303}{\Character{40}}%
\TextAt{13}{413.039}{273.303}{K}%
\TextAt{9}{423.856}{273.303}{\Character{41}}%
\TextAt{8}{432.431}{273.303}{and}%
\TextAt{13}{455.318}{273.303}{\Character{28}}%
\TextAt{9}{465.27}{273.303}{\Character{58}}%
\TextAt{13}{472.055}{273.303}{K}%
\NativeMathAt{486.405}{273.303}{11.955}{\rightarrow}%
\TextAt{18}{501.894}{273.303}{C}%
\TextAt{8}{514.55}{273.303}{any}%
\TextAt{8}{79.226}{256.257}{embedding,}%
\TextAt{8}{141.724}{256.257}{we}%
\TextAt{8}{159.005}{256.257}{find}%
\TextAt{8}{182.467}{256.257}{log}%
\TextAt{16}{199.428}{256.257}{\Character{106}}%
\TextAt{13}{202.749}{256.257}{\Character{28}}%
\TextAt{13}{209.157}{256.257}{\Character{16}}%
\TextAt{16}{215.183}{256.257}{\Character{106}}%
\TextAt{9}{221.917}{256.257}{\Character{61}}%
\TextAt{8}{234.435}{256.257}{log}%
\TextAt{9}{251.395}{256.257}{1}%
\TextAt{9}{260.648}{256.257}{\Character{61}}%
\TextAt{9}{273.154}{256.257}{0}%
\TextAt{8}{279.032}{256.257}{,}%
\TextAt{8}{286.248}{256.257}{so}%
\TextAt{8}{300.673}{256.257}{that}%
\TextAt{8}{326.095}{256.257}{certainly}%
\TextAt{13}{374.604}{256.257}{\Character{22}}%
\TextAt{9}{381.647}{256.257}{\Character{40}}%
\TextAt{13}{386.199}{256.257}{K}%
\TextAt{9}{397.017}{256.257}{\Character{41}}%
\NativeMathAt{404.982}{256.257}{11.955}{\subseteq}%
\TextAt{8}{417.693}{256.257}{ker}%
\TextAt{9}{433.63}{256.257}{\Character{40}}%
\TextAt{13}{438.182}{256.257}{\Character{21}}%
\TextAt{9}{445.011}{256.257}{\Character{41}}%
\TextAt{8}{449.563}{256.257}{.}%
\TextAt{8}{458.173}{256.257}{Conversely,}%
\TextAt{8}{519.762}{256.257}{let}%
\TextAt{13}{79.226}{239.211}{\Character{34}}%
\NativeMathAt{88.025}{239.211}{11.955}{\in}%
\TextAt{14}{99.316}{239.211}{o}%
\TextAt{3}{105.16}{243.549}{\Character{3}}%
\TextAt{21}{105.16}{236.19}{K}%
\TextAt{8}{116.549}{239.211}{be}%
\TextAt{8}{131.762}{239.211}{an}%
\TextAt{8}{147.31}{239.211}{element}%
\TextAt{8}{189.841}{239.211}{in}%
\TextAt{8}{202.783}{239.211}{the}%
\TextAt{8}{222.228}{239.211}{kernel,}%
\TextAt{8}{259.378}{239.211}{so}%
\TextAt{8}{273.037}{239.211}{that}%
\TextAt{13}{297.694}{239.211}{\Character{21}}%
\TextAt{9}{304.523}{239.211}{\Character{40}}%
\TextAt{13}{309.075}{239.211}{\Character{34}}%
\TextAt{9}{314.553}{239.211}{\Character{41}}%
\TextAt{9}{322.408}{239.211}{\Character{61}}%
\TextAt{13}{334.852}{239.211}{\Character{96}}%
\TextAt{9}{339.764}{239.211}{\Character{40}}%
\TextAt{13}{344.316}{239.211}{j}%
\TextAt{13}{349.827}{239.211}{\Character{34}}%
\TextAt{9}{355.31}{239.211}{\Character{41}}%
\TextAt{9}{363.165}{239.211}{\Character{61}}%
\TextAt{9}{375.586}{239.211}{0}%
\TextAt{8}{381.461}{239.211}{.}%
\TextAt{8}{389.673}{239.211}{This}%
\TextAt{8}{415.681}{239.211}{means}%
\TextAt{8}{450.8}{239.211}{that}%
\TextAt{16}{475.458}{239.211}{\Character{106}}%
\TextAt{13}{478.779}{239.211}{\Character{28}}%
\TextAt{13}{485.187}{239.211}{\Character{34}}%
\TextAt{16}{490.675}{239.211}{\Character{106}}%
\TextAt{9}{497.317}{239.211}{\Character{61}}%
\TextAt{9}{509.726}{239.211}{1}%
\TextAt{8}{518.782}{239.211}{for}%
\TextAt{8}{79.226}{222.165}{each}%
\TextAt{8}{105.77}{222.165}{embedding}%
\TextAt{13}{165.154}{222.165}{\Character{28}}%
\TextAt{9}{175.246}{222.165}{\Character{58}}%
\TextAt{13}{182.17}{222.165}{K}%
\NativeMathAt{196.66}{222.165}{11.955}{\rightarrow}%
\TextAt{18}{212.288}{222.165}{C}%
\TextAt{8}{220.923}{222.165}{,}%
\TextAt{8}{228.33}{222.165}{so}%
\TextAt{8}{242.906}{222.165}{that}%
\TextAt{13}{268.482}{222.165}{j}%
\TextAt{13}{273.993}{222.165}{\Character{34}}%
\TextAt{9}{283.148}{222.165}{\Character{61}}%
\TextAt{9}{295.916}{222.165}{\Character{40}}%
\TextAt{13}{300.477}{222.165}{\Character{28}}%
\TextAt{13}{306.885}{222.165}{\Character{34}}%
\TextAt{9}{312.374}{222.165}{\Character{41}}%
\TextAt{8}{321.03}{222.165}{lies}%
\TextAt{8}{341.453}{222.165}{in}%
\TextAt{8}{355.312}{222.165}{a}%
\TextAt{8}{365.274}{222.165}{bounded}%
\TextAt{8}{413.277}{222.165}{domain}%
\TextAt{8}{455.112}{222.165}{of}%
\TextAt{8}{468.648}{222.165}{the}%
\TextAt{18}{489.011}{222.165}{R}%
\TextAt{8}{497.645}{222.165}{-vector}%
\TextAt{8}{79.226}{205.119}{space}%
\TextAt{13}{110.019}{205.119}{K}%
\TextAt{25}{119.978}{203.326}{R}%
\TextAt{8}{126.233}{205.119}{.}%
\TextAt{8}{134.52}{205.119}{On}%
\TextAt{8}{153.549}{205.119}{the}%
\TextAt{8}{173.224}{205.119}{other}%
\TextAt{8}{203.312}{205.119}{hand,}%
\TextAt{13}{235.444}{205.119}{j}%
\TextAt{13}{240.955}{205.119}{\Character{34}}%
\TextAt{8}{249.853}{205.119}{is}%
\TextAt{8}{261.135}{205.119}{a}%
\TextAt{8}{270.409}{205.119}{point}%
\TextAt{8}{300.497}{205.119}{of}%
\TextAt{8}{313.345}{205.119}{the}%
\TextAt{8}{333.02}{205.119}{lattice}%
\TextAt{13}{368.308}{205.119}{j}%
\TextAt{14}{373.823}{205.119}{o}%
\TextAt{21}{379.667}{203.326}{K}%
\TextAt{8}{391.286}{205.119}{of}%
\TextAt{13}{404.134}{205.119}{K}%
\TextAt{25}{414.093}{203.326}{R}%
\TextAt{8}{423.763}{205.119}{(see}%
\TextAt{8}{446.749}{205.119}{(5.2)).}%
\TextAt{8}{483.67}{205.119}{Therefore}%
\TextAt{8}{79.226}{188.073}{the}%
\TextAt{8}{100.633}{188.073}{kernel}%
\TextAt{8}{136.35}{188.073}{of}%
\TextAt{13}{150.931}{188.073}{\Character{21}}%
\TextAt{8}{162.907}{188.073}{can}%
\TextAt{8}{185.617}{188.073}{contain}%
\TextAt{8}{228.173}{188.073}{only}%
\TextAt{8}{255.115}{188.073}{a}%
\TextAt{8}{266.121}{188.073}{finite}%
\TextAt{8}{297.283}{188.073}{number}%
\TextAt{8}{341.13}{188.073}{of}%
\TextAt{8}{355.71}{188.073}{elements,}%
\TextAt{8}{408.382}{188.073}{and}%
\TextAt{8}{432.395}{188.073}{thus,}%
\TextAt{8}{462.961}{188.073}{being}%
\TextAt{8}{495.749}{188.073}{a}%
\TextAt{8}{506.755}{188.073}{finite}%
\TextAt{8}{79.226}{171.027}{group,}%
\TextAt{8}{115.654}{171.027}{contains}%
\TextAt{8}{161.573}{171.027}{only}%
\TextAt{8}{187.265}{171.027}{roots}%
\TextAt{8}{216.926}{171.027}{of}%
\TextAt{8}{230.256}{171.027}{unity}%
\TextAt{8}{260.825}{171.027}{in}%
\TextAt{13}{274.478}{171.027}{K}%
\TextAt{3}{285.296}{175.365}{\Character{3}}%
\TextAt{8}{290.028}{171.027}{.}%
\Vector{q 0 G 0 g 0.399 w 452.334 -602.973 m 452.334 -594.903 l S Q}%
\Vector{q 0 G 0 g 0.399 w 452.534 -595.103 m 459.707 -595.103 l S Q}%
\Vector{q 0 G 0 g 0.399 w 452.534 -602.774 m 459.707 -602.774 l S Q}%
\Vector{q 0 G 0 g 0.399 w 459.906 -602.973 m 459.906 -594.903 l S Q}%
\TextAt{8}{97.159}{142.302}{Given}%
\TextAt{8}{131.834}{142.302}{this}%
\TextAt{8}{155.439}{142.302}{proposition,}%
\TextAt{8}{221.201}{142.302}{it}%
\TextAt{8}{233.688}{142.302}{remains}%
\TextAt{8}{278.106}{142.302}{to}%
\TextAt{8}{293.199}{142.302}{determine}%
\TextAt{8}{348.604}{142.302}{the}%
\TextAt{8}{369.543}{142.302}{group}%
\TextAt{9}{403.5}{142.302}{\Character{0}}%
\TextAt{8}{410.815}{142.302}{.}%
\TextAt{8}{421.61}{142.302}{For}%
\TextAt{8}{443.361}{142.302}{this,}%
\TextAt{8}{470.413}{142.302}{we}%
\TextAt{8}{488.423}{142.302}{need}%
\TextAt{8}{516.51}{142.302}{the}%
\TextAt{8}{79.226}{125.256}{following}%
\TextAt{12}{79.226}{99.521}{Lemma}%
\TextAt{12}{125.666}{99.521}{(7.2).}%
\TextAt{8}{159.917}{99.521}{Up}%
\TextAt{8}{178.028}{99.521}{to}%
\TextAt{8}{191.274}{99.521}{multiplication}%
\TextAt{8}{265.335}{99.521}{by}%
\TextAt{8}{280.529}{99.521}{units}%
\TextAt{8}{308.791}{99.521}{there}%
\TextAt{8}{337.637}{99.521}{are}%
\TextAt{8}{356.083}{99.521}{only}%
\TextAt{8}{380.71}{99.521}{finitely}%
\TextAt{8}{418.99}{99.521}{many}%
\TextAt{8}{449.797}{99.521}{elements}%
\TextAt{13}{496.589}{99.521}{\Character{11}}%
\NativeMathAt{507.432}{99.521}{11.955}{\in}%
\TextAt{14}{518.723}{99.521}{o}%
\TextAt{21}{524.567}{97.727}{K}%
\TextAt{8}{79.226}{82.475}{of}%
\TextAt{8}{92.556}{82.475}{given}%
\TextAt{8}{123.126}{82.475}{norm}%
\TextAt{13}{153.695}{82.475}{N}%
\TextAt{21}{163.069}{80.615}{K}%
\TextAt{3}{170.773}{80.615}{\Character{106}}%
\TextAt{25}{173.126}{80.615}{Q}%
\TextAt{9}{179.823}{82.475}{\Character{40}}%
\TextAt{13}{184.375}{82.475}{\Character{11}}%
\TextAt{9}{191.897}{82.475}{\Character{41}}%
\TextAt{9}{199.752}{82.475}{\Character{61}}%
\TextAt{13}{212.195}{82.475}{a}%
\TextAt{8}{218.34}{82.475}{.}%
\TextAt{8}{300.14}{43.033}{37}%
\Destination{78.23}{809.815}{nav.page.37}%
\Destination{79.36}{776.155}{nav.print.39}%
\Destination{79.36}{237.155}{nav.print.40}%
\Destination{79.36}{377.755}{nav.number.1.7.1}%
\Destination{79.36}{113.555}{nav.number.1.7.2}%
\Hotlink{394.188}{707.755}{6.858}{11.047}{nav.section.1.5}%
\Hotlink{446.249}{202.955}{25.078}{11.047}{nav.number.1.5.2}%
}

% Page 50 (body page 38)
\Page{%
\TextAt{2}{79.226}{796.943}{I.}%
\TextAt{2}{89.83}{796.943}{Algebraic}%
\TextAt{2}{138.648}{796.943}{Integers}%
\Vector{q 0 G 0 g 0.299 w 7.226 18.729 m 460.769 18.729 l S Q}%
\TextAt{19}{79.226}{761.124}{Proof.}%
\TextAt{8}{116}{761.124}{Let}%
\TextAt{13}{136.895}{761.124}{a}%
\NativeMathAt{146.361}{761.124}{11.955}{\in}%
\TextAt{18}{157.652}{761.124}{Z}%
\TextAt{8}{165.622}{761.124}{,}%
\TextAt{13}{172.711}{761.124}{a}%
\NativeMathAt{182.177}{761.124}{11.955}{\geq}%
\TextAt{9}{194.797}{761.124}{1}%
\TextAt{8}{200.649}{761.124}{.}%
\TextAt{8}{209.073}{761.124}{In}%
\TextAt{8}{223.62}{761.124}{every}%
\TextAt{8}{254.437}{761.124}{one}%
\TextAt{8}{275.822}{761.124}{of}%
\TextAt{8}{289.077}{761.124}{the}%
\TextAt{8}{309.159}{761.124}{finitely}%
\TextAt{8}{348.429}{761.124}{many}%
\TextAt{8}{380.227}{761.124}{cosets}%
\TextAt{8}{414.093}{761.124}{of}%
\TextAt{14}{427.349}{761.124}{o}%
\TextAt{21}{433.193}{759.33}{K}%
\TextAt{9}{441.395}{761.124}{\Character{47}}%
\TextAt{13}{447.248}{761.124}{a}%
\TextAt{14}{453.393}{761.124}{o}%
\TextAt{21}{459.238}{759.33}{K}%
\TextAt{8}{471.263}{761.124}{there}%
\TextAt{8}{501.1}{761.124}{exists,}%
\TextAt{8}{79.226}{744.078}{up}%
\TextAt{8}{96.304}{744.078}{to}%
\TextAt{8}{110.788}{744.078}{multiplication}%
\TextAt{8}{186.088}{744.078}{by}%
\TextAt{8}{202.521}{744.078}{units,}%
\TextAt{8}{235.316}{744.078}{at}%
\TextAt{8}{249.8}{744.078}{most}%
\TextAt{8}{278.654}{744.078}{one}%
\TextAt{8}{300.288}{744.078}{element}%
\TextAt{13}{343.703}{744.078}{\Character{11}}%
\TextAt{8}{355.296}{744.078}{such}%
\TextAt{8}{381.867}{744.078}{that}%
\TextAt{16}{407.409}{744.078}{\Character{106}}%
\TextAt{13}{410.73}{744.078}{N}%
\TextAt{9}{421.353}{744.078}{\Character{40}}%
\TextAt{13}{425.905}{744.078}{\Character{11}}%
\TextAt{9}{433.427}{744.078}{\Character{41}}%
\TextAt{16}{437.979}{744.078}{\Character{106}}%
\TextAt{9}{444.917}{744.078}{\Character{61}}%
\TextAt{16}{457.638}{744.078}{\Character{106}}%
\TextAt{13}{460.959}{744.078}{N}%
\TextAt{21}{470.333}{742.218}{K}%
\TextAt{3}{478.038}{742.218}{\Character{106}}%
\TextAt{25}{480.39}{742.218}{Q}%
\TextAt{9}{487.087}{744.078}{\Character{40}}%
\TextAt{13}{491.639}{744.078}{\Character{11}}%
\TextAt{9}{499.161}{744.078}{\Character{41}}%
\TextAt{16}{503.713}{744.078}{\Character{106}}%
\TextAt{9}{510.651}{744.078}{\Character{61}}%
\TextAt{13}{523.373}{744.078}{a}%
\TextAt{8}{529.518}{744.078}{.}%
\TextAt{8}{79.226}{727.032}{For}%
\TextAt{8}{100.196}{727.032}{if}%
\TextAt{13}{110.919}{727.032}{\Character{12}}%
\TextAt{9}{121.511}{727.032}{\Character{61}}%
\TextAt{13}{133.937}{727.032}{\Character{11}}%
\TextAt{9}{144.115}{727.032}{\Character{43}}%
\TextAt{13}{155.877}{727.032}{a\Character{13}}%
\TextAt{8}{168.744}{727.032}{,}%
\TextAt{13}{175.893}{727.032}{\Character{13}}%
\NativeMathAt{185.936}{727.032}{11.955}{\in}%
\TextAt{14}{197.227}{727.032}{o}%
\TextAt{21}{203.071}{725.238}{K}%
\TextAt{8}{211.274}{727.032}{,}%
\TextAt{8}{218.423}{727.032}{is}%
\TextAt{8}{230.187}{727.032}{another}%
\TextAt{8}{273.118}{727.032}{one,}%
\TextAt{8}{297.829}{727.032}{then}%
\TextAt{13}{251.746}{705.813}{\Character{11}}%
\Vector{q 0 G 0 g 0.478 w 179.746 -73.286 m 187.268 -73.286 l S Q}%
\TextAt{13}{251.871}{689.524}{\Character{12}}%
\TextAt{9}{263.784}{697.725}{\Character{61}}%
\TextAt{9}{276.194}{697.725}{1}%
\NativeMathAt{284.719}{697.725}{11.955}{\pm}%
\TextAt{13}{297.87}{705.813}{N}%
\TextAt{9}{308.493}{705.813}{\Character{40}}%
\TextAt{13}{313.045}{705.813}{\Character{12}}%
\TextAt{9}{320.316}{705.813}{\Character{41}}%
\Vector{q 0 G 0 g 0.478 w 225.87 -73.286 m 252.869 -73.286 l S Q}%
\TextAt{13}{307.734}{689.524}{\Character{12}}%
\TextAt{13}{326.064}{697.725}{\Character{13}}%
\NativeMathAt{336.107}{697.725}{11.955}{\in}%
\TextAt{14}{347.398}{697.725}{o}%
\TextAt{21}{353.242}{695.932}{K}%
\TextAt{8}{79.226}{668.305}{because}%
\TextAt{13}{122.528}{668.305}{N}%
\TextAt{9}{133.15}{668.305}{\Character{40}}%
\TextAt{13}{137.703}{668.305}{\Character{12}}%
\TextAt{9}{144.974}{668.305}{\Character{41}\Character{47}}%
\TextAt{13}{155.379}{668.305}{\Character{12}}%
\NativeMathAt{165.971}{668.305}{11.955}{\in}%
\TextAt{14}{177.262}{668.305}{o}%
\TextAt{21}{183.107}{666.512}{K}%
\TextAt{8}{191.309}{668.305}{,}%
\TextAt{8}{198.458}{668.305}{and}%
\TextAt{8}{221.221}{668.305}{by}%
\TextAt{8}{237.48}{668.305}{the}%
\TextAt{8}{257.637}{668.305}{same}%
\TextAt{8}{286.963}{668.305}{token}%
\TextAt{13}{250.12}{646.421}{\Character{12}}%
\Vector{q 0 G 0 g 0.478 w 177.995 -132.677 m 185.517 -132.677 l S Q}%
\TextAt{13}{249.995}{630.133}{\Character{11}}%
\TextAt{9}{262.033}{638.334}{\Character{61}}%
\TextAt{9}{274.442}{638.334}{1}%
\NativeMathAt{282.968}{638.334}{11.955}{\pm}%
\TextAt{13}{296.119}{646.422}{N}%
\TextAt{9}{306.742}{646.421}{\Character{40}}%
\TextAt{13}{311.294}{646.422}{\Character{11}}%
\TextAt{9}{318.816}{646.421}{\Character{41}}%
\Vector{q 0 G 0 g 0.478 w 224.119 -132.677 m 251.368 -132.677 l S Q}%
\TextAt{13}{305.983}{630.133}{\Character{11}}%
\TextAt{13}{324.563}{638.334}{\Character{13}}%
\NativeMathAt{334.607}{638.334}{11.955}{\in}%
\TextAt{14}{345.898}{638.334}{o}%
\TextAt{21}{351.742}{636.54}{K}%
\TextAt{13}{359.944}{638.334}{\Character{59}}%
\TextAt{8}{79.226}{612.356}{i.e.,}%
\TextAt{13}{102.919}{612.356}{\Character{12}}%
\TextAt{8}{115.358}{612.356}{is}%
\TextAt{8}{128.392}{612.356}{associated}%
\TextAt{8}{185.397}{612.356}{to}%
\TextAt{13}{200.978}{612.356}{\Character{11}}%
\TextAt{8}{208.5}{612.356}{.}%
\TextAt{8}{220.759}{612.356}{Therefore,}%
\TextAt{8}{278.596}{612.356}{up}%
\TextAt{8}{296.771}{612.356}{to}%
\TextAt{8}{312.351}{612.356}{multiplication}%
\TextAt{8}{388.748}{612.356}{by}%
\TextAt{8}{406.277}{612.356}{units,}%
\TextAt{8}{440.443}{612.356}{there}%
\TextAt{8}{471.625}{612.356}{are}%
\TextAt{8}{492.406}{612.356}{at}%
\TextAt{8}{507.986}{612.356}{most}%
\TextAt{9}{79.226}{595.31}{\Character{40}}%
\TextAt{14}{83.778}{595.31}{o}%
\TextAt{21}{89.623}{593.517}{K}%
\TextAt{9}{101.146}{595.31}{\Character{58}}%
\TextAt{13}{107.719}{595.31}{a}%
\TextAt{14}{113.864}{595.31}{o}%
\TextAt{21}{119.708}{593.517}{K}%
\TextAt{9}{127.911}{595.31}{\Character{41}}%
\TextAt{8}{136.36}{595.31}{elements}%
\TextAt{8}{184.217}{595.31}{of}%
\TextAt{8}{197.547}{595.31}{norm}%
\NativeMathAt{228.116}{595.31}{11.955}{\pm}%
\TextAt{13}{237.415}{595.31}{a}%
\TextAt{8}{243.56}{595.31}{.}%
\Vector{q 0 G 0 g 0.399 w 452.334 -178.69 m 452.334 -170.62 l S Q}%
\Vector{q 0 G 0 g 0.399 w 452.534 -170.82 m 459.707 -170.82 l S Q}%
\Vector{q 0 G 0 g 0.399 w 452.534 -178.491 m 459.707 -178.491 l S Q}%
\Vector{q 0 G 0 g 0.399 w 459.906 -178.69 m 459.906 -170.62 l S Q}%
\TextAt{12}{79.226}{562.059}{Theorem}%
\TextAt{12}{136.103}{562.059}{(7.3).}%
\TextAt{8}{171.299}{562.059}{The}%
\TextAt{8}{195.234}{562.059}{group}%
\TextAt{9}{228.29}{562.059}{\Character{0}}%
\TextAt{8}{239.383}{562.059}{is}%
\TextAt{8}{251.027}{562.059}{a}%
\TextAt{8}{260.663}{562.059}{complete}%
\TextAt{8}{309.967}{562.059}{lattice}%
\TextAt{8}{345.617}{562.059}{in}%
\TextAt{8}{359.151}{562.059}{the}%
\TextAt{9}{379.188}{562.059}{\Character{40}}%
\TextAt{13}{383.74}{562.059}{r}%
\TextAt{9}{391.753}{562.059}{\Character{43}}%
\TextAt{13}{403.271}{562.059}{s}%
\TextAt{16}{411.198}{562.059}{\Character{0}}%
\TextAt{9}{422.909}{562.059}{1\Character{41}}%
\TextAt{8}{433.314}{562.059}{-dimensional}%
\TextAt{8}{501.543}{562.059}{vector}%
\TextAt{8}{79.226}{545.013}{space}%
\TextAt{13}{110.501}{545.013}{H}%
\TextAt{8}{121.123}{545.013}{,}%
\TextAt{8}{128.273}{545.013}{and}%
\TextAt{8}{151.035}{545.013}{is}%
\TextAt{8}{162.799}{545.013}{therefore}%
\TextAt{8}{211.899}{545.013}{isomorphic}%
\TextAt{8}{271.149}{545.013}{to}%
\TextAt{18}{287.448}{519.035}{Z}%
\TextAt{21}{295.418}{523.971}{r}%
\TextAt{17}{299.475}{523.971}{\Character{43}}%
\TextAt{21}{306.061}{523.971}{s}%
\TextAt{3}{309.977}{523.971}{\Character{0}}%
\TextAt{17}{316.564}{523.971}{1}%
\TextAt{13}{321.296}{519.035}{\Character{58}}%
\TextAt{19}{79.226}{476.852}{Proof.}%
\TextAt{8}{116}{476.852}{We}%
\TextAt{8}{137.256}{476.852}{first}%
\TextAt{8}{162.494}{476.852}{show}%
\TextAt{8}{192.608}{476.852}{that}%
\TextAt{9}{219.089}{476.852}{\Character{0}}%
\TextAt{9}{231.606}{476.852}{\Character{61}}%
\TextAt{13}{245.938}{476.852}{\Character{21}}%
\TextAt{9}{252.767}{476.852}{\Character{40}}%
\TextAt{14}{257.319}{476.852}{o}%
\TextAt{3}{263.163}{481.19}{\Character{3}}%
\TextAt{21}{263.163}{473.831}{K}%
\TextAt{9}{271.366}{476.852}{\Character{41}}%
\TextAt{8}{280.927}{476.852}{is}%
\TextAt{8}{293.803}{476.852}{a}%
\TextAt{8}{304.67}{476.852}{lattice}%
\TextAt{8}{341.552}{476.852}{in}%
\TextAt{13}{356.316}{476.852}{H}%
\TextAt{8}{366.939}{476.852}{,}%
\TextAt{8}{375.477}{476.852}{i.e.,}%
\TextAt{8}{398.972}{476.852}{a}%
\TextAt{8}{409.839}{476.852}{discrete}%
\TextAt{8}{453.93}{476.852}{subgroup.}%
\TextAt{8}{512.613}{476.852}{The}%
\TextAt{8}{79.226}{459.806}{mapping}%
\TextAt{13}{127.358}{459.806}{\Character{21}}%
\TextAt{9}{137.507}{459.806}{\Character{58}}%
\TextAt{14}{144.079}{459.806}{o}%
\TextAt{3}{149.924}{464.144}{\Character{3}}%
\TextAt{21}{149.924}{456.785}{K}%
\NativeMathAt{161.447}{459.806}{11.955}{\rightarrow}%
\TextAt{13}{176.723}{459.806}{H}%
\TextAt{8}{191.243}{459.806}{arises}%
\TextAt{8}{223.235}{459.806}{by}%
\TextAt{8}{239.494}{459.806}{restricting}%
\TextAt{8}{295.492}{459.806}{the}%
\TextAt{8}{315.649}{459.806}{mapping}%
\TextAt{13}{249.731}{433.828}{K}%
\TextAt{3}{260.549}{438.764}{\Character{3}}%
\TextAt{21}{271.932}{441.092}{j}%
\NativeArrowAt{268.602}{433.828}{11.955}{11.955}{\longrightarrow}%
\TextAt{20}{283.878}{445.185}{\Character{89}}%
\TextAt{21}{289.177}{420.079}{\Character{28}}%
\TextAt{18}{301.147}{433.828}{C}%
\TextAt{3}{309.781}{438.764}{\Character{3}}%
\TextAt{21}{321.376}{440.605}{\Character{96}}%
\NativeArrowAt{317.834}{433.828}{11.955}{11.955}{\longrightarrow}%
\TextAt{20}{333.11}{445.185}{\Character{89}}%
\TextAt{21}{338.409}{420.079}{\Character{28}}%
\TextAt{18}{350.379}{433.828}{R}%
\TextAt{13}{359.013}{433.828}{\Character{59}}%
\TextAt{8}{79.226}{399.989}{and}%
\TextAt{8}{101.987}{399.989}{it}%
\TextAt{8}{113.69}{399.989}{suffices}%
\TextAt{8}{153.475}{399.989}{to}%
\TextAt{8}{167.783}{399.989}{show}%
\TextAt{8}{196.785}{399.989}{that,}%
\TextAt{8}{225.404}{399.989}{for}%
\TextAt{8}{243.287}{399.989}{any}%
\TextAt{13}{265.403}{399.989}{c}%
\TextAt{13}{273.748}{399.989}{\Character{62}}%
\TextAt{9}{286.187}{399.989}{0}%
\TextAt{8}{292.04}{399.989}{,}%
\TextAt{8}{299.188}{399.989}{the}%
\TextAt{8}{319.343}{399.989}{bounded}%
\TextAt{8}{367.138}{399.989}{domain}%
\TextAt{16}{408.764}{399.989}{\Character{102}}%
\TextAt{9}{414.742}{399.989}{\Character{40}}%
\TextAt{13}{419.294}{399.989}{x}%
\TextAt{21}{425.946}{398.196}{\Character{28}}%
\TextAt{9}{431.122}{399.989}{\Character{41}}%
\NativeMathAt{438.996}{399.989}{11.955}{\in}%
\TextAt{20}{450.287}{408.955}{\Character{81}}%
\TextAt{21}{461.578}{396.502}{\Character{28}}%
\TextAt{18}{468.747}{399.989}{R}%
\TextAt{16}{480.702}{399.989}{\Character{106}}%
\TextAt{16}{487.325}{399.989}{\Character{106}}%
\TextAt{13}{490.664}{399.989}{x}%
\TextAt{21}{497.316}{398.196}{\Character{28}}%
\TextAt{16}{502.493}{399.989}{\Character{106}}%
\NativeMathAt{509.116}{399.989}{11.955}{\leq}%
\TextAt{13}{521.754}{399.989}{c}%
\TextAt{16}{526.792}{399.989}{\Character{103}}%
\TextAt{8}{79.226}{382.943}{contains}%
\TextAt{8}{125.16}{382.943}{only}%
\TextAt{8}{150.865}{382.943}{finitely}%
\TextAt{8}{190.224}{382.943}{many}%
\TextAt{8}{222.11}{382.943}{points}%
\TextAt{8}{257.308}{382.943}{of}%
\TextAt{9}{270.652}{382.943}{\Character{0}}%
\TextAt{9}{281.292}{382.943}{\Character{61}}%
\TextAt{13}{293.759}{382.943}{\Character{96}}%
\TextAt{9}{298.671}{382.943}{\Character{40}}%
\TextAt{13}{303.224}{382.943}{j}%
\TextAt{14}{308.739}{382.943}{o}%
\TextAt{3}{314.583}{387.281}{\Character{3}}%
\TextAt{21}{314.583}{379.922}{K}%
\TextAt{9}{322.786}{382.943}{\Character{41}}%
\TextAt{8}{327.338}{382.943}{.}%
\TextAt{8}{335.828}{382.943}{Since}%
\TextAt{13}{366.399}{382.943}{\Character{96}}%
\TextAt{9}{371.312}{382.943}{\Character{40}\Character{40}}%
\TextAt{13}{380.416}{382.943}{z}%
\TextAt{21}{385.854}{381.15}{\Character{28}}%
\TextAt{9}{391.031}{382.943}{\Character{41}\Character{41}}%
\TextAt{9}{403.452}{382.943}{\Character{61}}%
\TextAt{9}{415.898}{382.943}{\Character{40}}%
\TextAt{8}{420.481}{382.943}{log}%
\TextAt{16}{437.441}{382.943}{\Character{106}}%
\TextAt{13}{440.762}{382.943}{z}%
\TextAt{21}{446.2}{381.15}{\Character{28}}%
\TextAt{16}{451.377}{382.943}{\Character{106}}%
\TextAt{9}{454.698}{382.943}{\Character{41}}%
\TextAt{8}{459.25}{382.943}{,}%
\TextAt{8}{466.416}{382.943}{the}%
\TextAt{8}{486.587}{382.943}{preimage}%
\TextAt{8}{79.226}{365.897}{of}%
\TextAt{8}{92.556}{365.897}{this}%
\TextAt{8}{115.379}{365.897}{domain}%
\TextAt{8}{157.006}{365.897}{with}%
\TextAt{8}{183.667}{365.897}{respect}%
\TextAt{8}{223.716}{365.897}{to}%
\TextAt{13}{238.027}{365.897}{\Character{96}}%
\TextAt{8}{246.836}{365.897}{is}%
\TextAt{8}{258.6}{365.897}{the}%
\TextAt{8}{278.756}{365.897}{bounded}%
\TextAt{8}{326.553}{365.897}{domain}%
\TextAt{20}{217.919}{353.171}{\Character{40}}%
\TextAt{9}{227.55}{332.728}{\Character{40}}%
\TextAt{13}{232.102}{332.728}{z}%
\TextAt{21}{237.54}{330.934}{\Character{28}}%
\TextAt{9}{242.717}{332.728}{\Character{41}}%
\NativeMathAt{250.59}{332.728}{11.955}{\in}%
\TextAt{20}{261.881}{344.085}{\Character{89}}%
\TextAt{21}{267.18}{318.979}{\Character{28}}%
\TextAt{18}{279.149}{332.728}{C}%
\TextAt{3}{287.784}{337.664}{\Character{3}}%
\TextAt{20}{296.413}{353.649}{\Character{12}}%
\TextAt{20}{296.413}{346.476}{\Character{12}}%
\TextAt{20}{296.413}{339.303}{\Character{12}}%
\TextAt{20}{296.413}{332.13}{\Character{12}}%
\TextAt{20}{296.413}{324.957}{\Character{12}}%
\TextAt{13}{304.296}{332.728}{e}%
\TextAt{3}{309.721}{337.664}{\Character{0}}%
\TextAt{21}{316.308}{337.664}{c}%
\NativeMathAt{323.795}{332.728}{11.955}{\leq}%
\TextAt{16}{336.396}{332.728}{\Character{106}}%
\TextAt{13}{339.735}{332.728}{z}%
\TextAt{21}{345.173}{330.935}{\Character{28}}%
\TextAt{16}{350.349}{332.728}{\Character{106}}%
\NativeMathAt{356.972}{332.728}{11.955}{\leq}%
\TextAt{13}{369.61}{332.728}{e}%
\TextAt{21}{375.036}{337.664}{c}%
\TextAt{20}{379.202}{353.171}{\Character{41}}%
\TextAt{13}{390.825}{332.728}{\Character{58}}%
\TextAt{8}{79.226}{299.439}{It}%
\TextAt{8}{91.08}{299.439}{contains}%
\TextAt{8}{136.181}{299.439}{only}%
\TextAt{8}{161.054}{299.439}{finitely}%
\TextAt{8}{199.579}{299.439}{many}%
\TextAt{8}{230.633}{299.439}{elements}%
\TextAt{8}{277.671}{299.439}{of}%
\TextAt{8}{290.182}{299.439}{the}%
\TextAt{8}{309.52}{299.439}{set}%
\TextAt{13}{326.968}{299.439}{j}%
\TextAt{14}{332.484}{299.439}{o}%
\TextAt{3}{338.328}{303.777}{\Character{3}}%
\TextAt{21}{338.328}{296.418}{K}%
\TextAt{8}{349.609}{299.439}{because}%
\TextAt{8}{392.092}{299.439}{this}%
\TextAt{8}{414.096}{299.439}{is}%
\TextAt{8}{425.041}{299.439}{a}%
\TextAt{8}{433.978}{299.439}{subset}%
\TextAt{8}{469.048}{299.439}{of}%
\TextAt{8}{481.559}{299.439}{the}%
\TextAt{8}{500.897}{299.439}{lattice}%
\TextAt{13}{79.226}{282.393}{j}%
\TextAt{14}{84.742}{282.393}{o}%
\TextAt{21}{90.586}{280.6}{K}%
\TextAt{8}{102.686}{282.393}{in}%
\TextAt{9}{116.339}{282.393}{\Character{91}}%
\TextAt{20}{119.59}{291.359}{\Character{81}}%
\TextAt{21}{130.881}{278.906}{\Character{28}}%
\TextAt{18}{138.05}{282.393}{C}%
\TextAt{9}{146.685}{282.393}{\Character{93}}%
\TextAt{17}{149.936}{286.731}{\Character{43}}%
\TextAt{8}{160.918}{282.393}{(see}%
\TextAt{8}{184.386}{282.393}{(5.2)).}%
\TextAt{8}{221.467}{282.393}{Therefore}%
\TextAt{9}{274.465}{282.393}{\Character{0}}%
\TextAt{8}{285.676}{282.393}{is}%
\TextAt{8}{297.44}{282.393}{a}%
\TextAt{8}{307.195}{282.393}{lattice.}%
\TextAt{8}{97.159}{262.358}{We}%
\TextAt{8}{118.242}{262.358}{now}%
\TextAt{8}{143.569}{262.358}{show}%
\TextAt{8}{173.512}{262.358}{that}%
\TextAt{9}{199.819}{262.358}{\Character{0}}%
\TextAt{8}{211.97}{262.358}{is}%
\TextAt{8}{224.672}{262.358}{a}%
\TextAt{8}{235.366}{262.358}{complete}%
\TextAt{8}{285.728}{262.358}{lattice}%
\TextAt{8}{322.437}{262.358}{in}%
\TextAt{13}{337.028}{262.358}{H}%
\TextAt{8}{347.651}{262.358}{.}%
\TextAt{8}{358.915}{262.358}{This}%
\TextAt{8}{386.574}{262.358}{is}%
\TextAt{8}{399.277}{262.358}{the}%
\TextAt{8}{420.372}{262.358}{principal}%
\TextAt{8}{470.088}{262.358}{claim}%
\TextAt{8}{502.242}{262.358}{of}%
\TextAt{8}{516.51}{262.358}{the}%
\TextAt{8}{79.226}{245.312}{theorem.}%
\TextAt{8}{130.917}{245.312}{We}%
\TextAt{8}{151.6}{245.312}{apply}%
\TextAt{8}{184.334}{245.312}{the}%
\TextAt{8}{205.029}{245.312}{criterion}%
\TextAt{8}{252.395}{245.312}{(4.3).}%
\TextAt{8}{286.537}{245.312}{So}%
\TextAt{8}{303.334}{245.312}{we}%
\TextAt{8}{321.1}{245.312}{have}%
\TextAt{8}{348.633}{245.312}{to}%
\TextAt{8}{363.482}{245.312}{find}%
\TextAt{8}{387.428}{245.312}{a}%
\TextAt{8}{397.722}{245.312}{bounded}%
\TextAt{8}{446.057}{245.312}{set}%
\TextAt{13}{464.863}{245.312}{M}%
\NativeMathAt{481.675}{245.312}{11.955}{\subseteq}%
\TextAt{13}{495.212}{245.312}{H}%
\TextAt{8}{510.27}{245.312}{such}%
\TextAt{8}{79.226}{228.266}{that}%
\TextAt{13}{261.895}{211.22}{H}%
\TextAt{9}{275.839}{211.22}{\Character{61}}%
\TextAt{20}{289.505}{222.578}{\Character{91}}%
\TextAt{21}{288.264}{197.206}{\Character{13}}%
\NativeMathAt{293.096}{197.206}{7.97}{\in}%
\TextAt{17}{298.742}{197.206}{\Character{0}}%
\TextAt{9}{304.03}{211.22}{\Character{40}}%
\TextAt{13}{308.583}{211.22}{M}%
\TextAt{9}{323.813}{211.22}{\Character{43}}%
\TextAt{13}{335.574}{211.22}{\Character{13}}%
\TextAt{9}{342.296}{211.22}{\Character{41}}%
\TextAt{13}{346.849}{211.22}{\Character{58}}%
\TextAt{8}{79.226}{179.988}{We}%
\TextAt{8}{99.193}{179.988}{construct}%
\TextAt{8}{150.459}{179.988}{this}%
\TextAt{8}{173.104}{179.988}{set}%
\TextAt{8}{191.194}{179.988}{through}%
\TextAt{8}{235.251}{179.988}{its}%
\TextAt{8}{251.393}{179.988}{preimage}%
\TextAt{8}{301.295}{179.988}{with}%
\TextAt{8}{327.778}{179.988}{respect}%
\TextAt{8}{367.651}{179.988}{to}%
\TextAt{8}{381.784}{179.988}{the}%
\TextAt{8}{401.763}{179.988}{surjective}%
\TextAt{8}{453.997}{179.988}{homomorphism}%
\TextAt{13}{275.084}{154.01}{\Character{96}}%
\TextAt{9}{283.317}{154.01}{\Character{58}}%
\TextAt{13}{289.889}{154.01}{S}%
\NativeArrowAt{301.106}{154.01}{11.955}{19.26}{\longrightarrow}%
\TextAt{13}{323.688}{154.01}{H}%
\TextAt{13}{333.647}{154.01}{\Character{58}}%
\TextAt{8}{79.226}{128.032}{More}%
\TextAt{8}{109.646}{128.032}{precisely,}%
\TextAt{8}{160.005}{128.032}{we}%
\TextAt{8}{177.417}{128.032}{will}%
\TextAt{8}{199.707}{128.032}{construct}%
\TextAt{8}{251.335}{128.032}{a}%
\TextAt{8}{261.275}{128.032}{bounded}%
\TextAt{8}{309.257}{128.032}{set}%
\TextAt{13}{327.709}{128.032}{T}%
\TextAt{8}{340.278}{128.032}{in}%
\TextAt{8}{354.116}{128.032}{the}%
\TextAt{8}{374.457}{128.032}{norm-one}%
\TextAt{8}{426.671}{128.032}{surface}%
\TextAt{13}{466.26}{128.032}{S}%
\TextAt{8}{474.156}{128.032}{,}%
\TextAt{8}{481.536}{128.032}{the}%
\TextAt{8}{501.877}{128.032}{multi-}%
\TextAt{8}{79.226}{110.986}{plicative}%
\TextAt{8}{126.055}{110.986}{translations}%
\TextAt{13}{189.931}{110.986}{T}%
\TextAt{13}{198.395}{110.986}{j}%
\TextAt{13}{203.919}{110.986}{\Character{34}}%
\TextAt{8}{209.411}{110.986}{,}%
\TextAt{13}{216.561}{110.986}{\Character{34}}%
\NativeMathAt{225.359}{110.986}{11.955}{\in}%
\TextAt{14}{236.65}{110.986}{o}%
\TextAt{3}{242.495}{115.324}{\Character{3}}%
\TextAt{21}{242.495}{107.965}{K}%
\TextAt{8}{250.697}{110.986}{,}%
\TextAt{8}{257.846}{110.986}{of}%
\TextAt{8}{271.176}{110.986}{which}%
\TextAt{8}{304.663}{110.986}{cover}%
\TextAt{8}{334.909}{110.986}{all}%
\TextAt{8}{351.169}{110.986}{of}%
\TextAt{13}{364.499}{110.986}{S}%
\TextAt{8}{372.394}{110.986}{:}%
\TextAt{13}{271.567}{85.008}{S}%
\TextAt{9}{282.783}{85.008}{\Character{61}}%
\TextAt{20}{298.815}{96.365}{\Character{91}}%
\TextAt{21}{295.209}{70.913}{\Character{34}}%
\NativeMathAt{299.119}{70.913}{7.97}{\in}%
\TextAt{27}{304.765}{70.913}{o}%
\TextAt{23}{308.661}{73.659}{\Character{3}}%
\TextAt{22}{308.661}{68.139}{K}%
\TextAt{13}{317.696}{85.008}{T}%
\TextAt{13}{326.172}{85.008}{j}%
\TextAt{13}{331.684}{85.008}{\Character{34}\Character{58}}%
\TextAt{8}{300.14}{43.033}{38}%
\Destination{78.23}{809.815}{nav.page.38}%
\Destination{79.36}{277.355}{nav.print.41}%
\Destination{79.36}{576.055}{nav.number.1.7.3}%
\Hotlink{183.886}{280.155}{25.078}{11.047}{nav.number.1.5.2}%
\Hotlink{251.895}{243.155}{25.078}{11.047}{nav.number.1.4.3}%
}

% Page 51 (body page 39)
\Page{%
\TextAt{2}{79.226}{796.943}{I.}%
\TextAt{2}{89.83}{796.943}{Algebraic}%
\TextAt{2}{138.648}{796.943}{Integers}%
\Vector{q 0 G 0 g 0.299 w 7.226 18.729 m 460.769 18.729 l S Q}%
\TextAt{8}{79.226}{761.124}{For}%
\TextAt{13}{101.1}{761.124}{x}%
\TextAt{9}{112.613}{761.124}{\Character{61}}%
\TextAt{9}{126.565}{761.124}{\Character{40}}%
\TextAt{13}{131.132}{761.124}{x}%
\TextAt{21}{137.784}{759.33}{\Character{28}}%
\TextAt{9}{142.96}{761.124}{\Character{41}}%
\NativeMathAt{152.374}{761.124}{11.955}{\in}%
\TextAt{13}{165.205}{761.124}{T}%
\TextAt{8}{173.692}{761.124}{,}%
\TextAt{8}{181.972}{761.124}{it}%
\TextAt{8}{194.58}{761.124}{will}%
\TextAt{8}{217.589}{761.124}{follow}%
\TextAt{8}{252.314}{761.124}{that}%
\TextAt{8}{278.587}{761.124}{the}%
\TextAt{8}{299.647}{761.124}{absolute}%
\TextAt{8}{346.794}{761.124}{values}%
\TextAt{16}{382.559}{761.124}{\Character{106}}%
\TextAt{13}{385.88}{761.124}{x}%
\TextAt{21}{392.532}{759.33}{\Character{28}}%
\TextAt{16}{397.709}{761.124}{\Character{106}}%
\TextAt{8}{405.831}{761.124}{are}%
\TextAt{8}{426.246}{761.124}{bounded}%
\TextAt{8}{474.947}{761.124}{from}%
\TextAt{8}{503.491}{761.124}{above}%
\TextAt{8}{79.226}{744.078}{and}%
\TextAt{8}{101.785}{744.078}{also}%
\TextAt{8}{125.062}{744.078}{away}%
\TextAt{8}{154.136}{744.078}{from}%
\TextAt{8}{181.573}{744.078}{zero,}%
\TextAt{8}{209.373}{744.078}{because}%
\TextAt{20}{252.471}{753.044}{\Character{81}}%
\TextAt{21}{263.762}{740.591}{\Character{28}}%
\TextAt{16}{270.931}{744.078}{\Character{106}}%
\TextAt{13}{274.252}{744.078}{x}%
\TextAt{21}{280.904}{742.284}{\Character{28}}%
\TextAt{16}{286.081}{744.078}{\Character{106}}%
\TextAt{9}{292.722}{744.078}{\Character{61}}%
\TextAt{9}{305.132}{744.078}{1}%
\TextAt{8}{311.001}{744.078}{.}%
\TextAt{8}{319.381}{744.078}{Thus}%
\TextAt{13}{348.827}{744.078}{M}%
\TextAt{9}{364.721}{744.078}{\Character{61}}%
\TextAt{13}{377.147}{744.078}{\Character{96}}%
\TextAt{9}{382.059}{744.078}{\Character{40}}%
\TextAt{13}{386.611}{744.078}{T}%
\TextAt{9}{395.098}{744.078}{\Character{41}}%
\TextAt{8}{403.344}{744.078}{will}%
\TextAt{8}{425.245}{744.078}{also}%
\TextAt{8}{448.522}{744.078}{be}%
\TextAt{8}{464.243}{744.078}{bounded.}%
\TextAt{8}{516.522}{744.078}{We}%
\TextAt{8}{79.226}{727.032}{choose}%
\TextAt{8}{116.359}{727.032}{real}%
\TextAt{8}{139.122}{727.032}{numbers}%
\TextAt{13}{186.333}{727.032}{c}%
\TextAt{21}{191.371}{725.238}{\Character{28}}%
\TextAt{13}{199.868}{727.032}{\Character{62}}%
\TextAt{9}{212.293}{727.032}{0}%
\TextAt{8}{218.146}{727.032}{,}%
\TextAt{8}{225.295}{727.032}{for}%
\TextAt{13}{243.18}{727.032}{\Character{28}}%
\NativeMathAt{252.92}{727.032}{11.955}{\in}%
\TextAt{8}{264.211}{727.032}{Hom}%
\TextAt{9}{288.6}{727.032}{\Character{40}}%
\TextAt{13}{293.152}{727.032}{K}%
\TextAt{13}{303.314}{727.032}{\Character{59}}%
\TextAt{18}{308.564}{727.032}{C}%
\TextAt{9}{317.198}{727.032}{\Character{41}}%
\TextAt{8}{321.75}{727.032}{,}%
\TextAt{8}{328.899}{727.032}{satisfying}%
\TextAt{13}{381.06}{727.032}{c}%
\TextAt{21}{386.098}{725.238}{\Character{28}}%
\TextAt{9}{394.595}{727.032}{\Character{61}}%
\TextAt{13}{407.021}{727.032}{c}%
\TextAt{17}{412.516}{725.238}{\Character{22}}%
\TextAt{21}{412.059}{725.238}{\Character{28}}%
\TextAt{8}{421.132}{727.032}{and}%
\TextAt{13}{237.092}{698.97}{C}%
\TextAt{9}{249.647}{698.97}{\Character{61}}%
\TextAt{20}{262.072}{710.328}{\Character{89}}%
\TextAt{21}{267.371}{685.222}{\Character{28}}%
\TextAt{13}{279.341}{698.97}{c}%
\TextAt{21}{284.379}{697.177}{\Character{28}}%
\TextAt{13}{292.876}{698.97}{\Character{62}}%
\TextAt{20}{305.302}{715.827}{\Character{18}}%
\TextAt{9}{315.906}{707.058}{2}%
\Vector{q 0 G 0 g 0.478 w 243.298 -72.041 m 250.367 -72.041 l S Q}%
\TextAt{13}{315.298}{690.77}{\Character{25}}%
\TextAt{20}{323.562}{715.827}{\Character{19}}%
\TextAt{21}{332.363}{713.151}{s}%
\TextAt{20}{338.769}{709.776}{\Character{112}}%
\Vector{q 0 G 0 g 0.478 w 278.724 -63.985 m 299.652 -63.985 l S Q}%
\TextAt{16}{350.724}{698.97}{\Character{106}}%
\TextAt{13}{354.045}{698.97}{d}%
\TextAt{21}{360.128}{697.177}{K}%
\TextAt{16}{368.331}{698.97}{\Character{106}}%
\TextAt{13}{371.652}{698.97}{\Character{59}}%
\TextAt{8}{79.226}{668.269}{and}%
\TextAt{8}{101.989}{668.269}{we}%
\TextAt{8}{119.216}{668.269}{consider}%
\TextAt{8}{164.801}{668.269}{the}%
\TextAt{8}{184.958}{668.269}{set}%
\TextAt{13}{233.332}{651.223}{X}%
\TextAt{9}{247.308}{651.223}{\Character{61}}%
\TextAt{16}{259.733}{651.223}{\Character{102}}%
\TextAt{9}{265.711}{651.223}{\Character{40}}%
\TextAt{13}{270.263}{651.223}{z}%
\TextAt{21}{275.701}{649.43}{\Character{28}}%
\TextAt{9}{280.878}{651.223}{\Character{41}}%
\NativeMathAt{288.751}{651.223}{11.955}{\in}%
\TextAt{13}{300.042}{651.223}{K}%
\TextAt{25}{310.001}{649.43}{R}%
\TextAt{16}{319.576}{651.223}{\Character{106}}%
\TextAt{16}{326.199}{651.223}{\Character{106}}%
\TextAt{13}{329.539}{651.223}{z}%
\TextAt{21}{334.976}{649.43}{\Character{28}}%
\TextAt{16}{340.153}{651.223}{\Character{106}}%
\TextAt{13}{346.795}{651.223}{\Character{60}}%
\TextAt{13}{359.204}{651.223}{c}%
\TextAt{21}{364.258}{649.43}{\Character{28}}%
\TextAt{16}{369.435}{651.223}{\Character{103}}%
\TextAt{13}{375.412}{651.223}{\Character{58}}%
\TextAt{8}{79.226}{630.054}{For}%
\TextAt{8}{100.196}{630.054}{an}%
\TextAt{8}{116.455}{630.054}{arbitrary}%
\TextAt{8}{166.224}{630.054}{point}%
\TextAt{13}{196.793}{630.054}{y}%
\TextAt{9}{206.251}{630.054}{\Character{61}}%
\TextAt{9}{218.661}{630.054}{\Character{40}}%
\TextAt{13}{223.229}{630.054}{y}%
\TextAt{21}{228.936}{628.261}{\Character{28}}%
\TextAt{9}{234.113}{630.054}{\Character{41}}%
\NativeMathAt{241.986}{630.054}{11.955}{\in}%
\TextAt{13}{253.277}{630.054}{S}%
\TextAt{8}{261.172}{630.054}{,}%
\TextAt{8}{268.321}{630.054}{it}%
\TextAt{8}{280.026}{630.054}{follows}%
\TextAt{8}{318.461}{630.054}{that}%
\TextAt{13}{231.889}{606.544}{X}%
\TextAt{13}{242.529}{606.544}{y}%
\TextAt{9}{252.002}{606.544}{\Character{61}}%
\TextAt{16}{264.427}{606.544}{\Character{102}}%
\TextAt{9}{270.405}{606.544}{\Character{40}}%
\TextAt{13}{274.957}{606.544}{z}%
\TextAt{21}{280.395}{604.751}{\Character{28}}%
\TextAt{9}{285.572}{606.544}{\Character{41}}%
\NativeMathAt{293.445}{606.544}{11.955}{\in}%
\TextAt{13}{304.736}{606.544}{K}%
\TextAt{25}{314.695}{604.751}{R}%
\TextAt{16}{324.27}{606.544}{\Character{106}}%
\TextAt{16}{330.893}{606.544}{\Character{106}}%
\TextAt{13}{334.233}{606.544}{z}%
\TextAt{21}{339.671}{604.751}{\Character{28}}%
\TextAt{16}{344.847}{606.544}{\Character{106}}%
\TextAt{13}{351.489}{606.544}{\Character{60}}%
\TextAt{13}{363.899}{606.544}{c}%
\TextAt{3}{368.952}{611.48}{\Character{48}}%
\TextAt{21}{368.952}{603.588}{\Character{28}}%
\TextAt{16}{374.129}{606.544}{\Character{103}}%
\TextAt{8}{79.226}{583.034}{where}%
\TextAt{13}{114.551}{583.034}{c}%
\TextAt{3}{119.589}{587.372}{\Character{48}}%
\TextAt{21}{119.589}{580.078}{\Character{28}}%
\TextAt{9}{130.669}{583.034}{\Character{61}}%
\TextAt{13}{145.678}{583.034}{c}%
\TextAt{21}{150.716}{581.24}{\Character{28}}%
\TextAt{16}{155.892}{583.034}{\Character{106}}%
\TextAt{13}{159.213}{583.034}{y}%
\TextAt{21}{164.921}{581.24}{\Character{28}}%
\TextAt{16}{170.097}{583.034}{\Character{106}}%
\TextAt{8}{173.418}{583.034}{,}%
\TextAt{8}{182.462}{583.034}{and}%
\TextAt{8}{206.74}{583.034}{one}%
\TextAt{8}{229.716}{583.034}{has}%
\TextAt{13}{252.105}{583.034}{c}%
\TextAt{3}{257.143}{587.372}{\Character{48}}%
\TextAt{21}{257.143}{580.078}{\Character{28}}%
\TextAt{9}{268.223}{583.034}{\Character{61}}%
\TextAt{13}{283.232}{583.034}{c}%
\TextAt{3}{288.27}{587.372}{\Character{48}}%
\TextAt{17}{288.727}{580.078}{\Character{22}}%
\TextAt{21}{288.27}{580.078}{\Character{28}}%
\TextAt{8}{298.859}{583.034}{and}%
\TextAt{20}{323.138}{592}{\Character{81}}%
\TextAt{21}{334.429}{579.547}{\Character{28}}%
\TextAt{13}{341.598}{583.034}{c}%
\TextAt{3}{346.636}{587.372}{\Character{48}}%
\TextAt{21}{346.636}{580.078}{\Character{28}}%
\TextAt{9}{357.716}{583.034}{\Character{61}}%
\TextAt{20}{372.724}{592}{\Character{81}}%
\TextAt{21}{384.015}{579.547}{\Character{28}}%
\TextAt{13}{391.184}{583.034}{c}%
\TextAt{21}{396.222}{581.24}{\Character{28}}%
\TextAt{9}{407.303}{583.034}{\Character{61}}%
\TextAt{13}{422.311}{583.034}{C}%
\TextAt{8}{436.958}{583.034}{because}%
\TextAt{20}{481.775}{592}{\Character{81}}%
\TextAt{21}{493.066}{579.547}{\Character{28}}%
\TextAt{16}{500.235}{583.034}{\Character{106}}%
\TextAt{13}{503.556}{583.034}{y}%
\TextAt{21}{509.264}{581.24}{\Character{28}}%
\TextAt{16}{514.44}{583.034}{\Character{106}}%
\TextAt{9}{523.665}{583.034}{\Character{61}}%
\TextAt{16}{79.226}{565.988}{\Character{106}}%
\TextAt{13}{82.547}{565.988}{N}%
\TextAt{9}{93.17}{565.988}{\Character{40}}%
\TextAt{13}{97.722}{565.988}{y}%
\TextAt{9}{103.859}{565.988}{\Character{41}}%
\TextAt{16}{108.411}{565.988}{\Character{106}}%
\TextAt{9}{115.053}{565.988}{\Character{61}}%
\TextAt{9}{127.463}{565.988}{1}%
\TextAt{8}{133.331}{565.988}{.}%
\TextAt{8}{141.779}{565.988}{Then,}%
\TextAt{8}{175.589}{565.988}{by}%
\TextAt{8}{191.848}{565.988}{(5.3),}%
\TextAt{8}{223.075}{565.988}{there}%
\TextAt{8}{252.986}{565.988}{is}%
\TextAt{8}{264.75}{565.988}{a}%
\TextAt{8}{274.506}{565.988}{point}%
\TextAt{13}{209.682}{542.477}{j}%
\TextAt{13}{215.193}{542.477}{a}%
\TextAt{9}{224.663}{542.477}{\Character{61}}%
\TextAt{9}{237.073}{542.477}{\Character{40}}%
\TextAt{13}{241.641}{542.477}{\Character{28}}%
\TextAt{13}{248.049}{542.477}{a}%
\TextAt{9}{254.205}{542.477}{\Character{41}}%
\NativeMathAt{262.078}{542.477}{11.955}{\in}%
\TextAt{13}{273.369}{542.477}{X}%
\TextAt{13}{284.009}{542.477}{y}%
\TextAt{13}{290.142}{542.477}{\Character{59}}%
\TextAt{13}{319.289}{542.477}{a}%
\NativeMathAt{328.781}{542.477}{11.955}{\in}%
\TextAt{14}{340.072}{542.477}{o}%
\TextAt{21}{345.916}{540.684}{K}%
\TextAt{13}{354.119}{542.477}{\Character{59}}%
\TextAt{13}{371.299}{542.477}{a}%
\NativeMathAt{380.784}{542.477}{11.955}{\neq}%
\TextAt{9}{393.194}{542.477}{0}%
\TextAt{13}{399.062}{542.477}{\Character{58}}%
\TextAt{8}{79.226}{518.967}{Now,}%
\TextAt{8}{109.52}{518.967}{according}%
\TextAt{8}{162.512}{518.967}{to}%
\TextAt{8}{177.128}{518.967}{lemma}%
\TextAt{8}{215.152}{518.967}{(7.2),}%
\TextAt{8}{246.76}{518.967}{we}%
\TextAt{8}{264.293}{518.967}{may}%
\TextAt{8}{289.968}{518.967}{pick}%
\TextAt{8}{314.985}{518.967}{a}%
\TextAt{8}{325.045}{518.967}{system}%
\TextAt{13}{364.169}{518.967}{\Character{11}}%
\TextAt{17}{371.615}{517.174}{1}%
\TextAt{13}{376.347}{518.967}{\Character{59}}%
\TextAt{13}{381.571}{518.967}{\Character{58}}%
\TextAt{13}{386.808}{518.967}{\Character{58}}%
\TextAt{13}{392.044}{518.967}{\Character{58}}%
\TextAt{13}{397.269}{518.967}{\Character{59}}%
\TextAt{13}{402.505}{518.967}{\Character{11}}%
\TextAt{21}{410.014}{517.174}{N}%
\NativeMathAt{421.924}{518.967}{11.955}{\in}%
\TextAt{14}{433.735}{518.967}{o}%
\TextAt{21}{439.579}{517.174}{K}%
\TextAt{8}{447.782}{518.967}{,}%
\TextAt{13}{455.313}{518.967}{\Character{11}}%
\TextAt{21}{462.759}{517.174}{i}%
\NativeMathAt{469.981}{518.967}{11.955}{\neq}%
\TextAt{9}{482.917}{518.967}{0}%
\TextAt{8}{488.78}{518.967}{,}%
\TextAt{8}{496.312}{518.967}{in}%
\TextAt{8}{510.27}{518.967}{such}%
\TextAt{8}{79.226}{501.921}{a}%
\TextAt{8}{89.266}{501.921}{way}%
\TextAt{8}{113.293}{501.921}{that}%
\TextAt{8}{138.946}{501.921}{every}%
\TextAt{13}{170.122}{501.921}{a}%
\NativeMathAt{180.072}{501.921}{11.955}{\in}%
\TextAt{14}{191.847}{501.921}{o}%
\TextAt{21}{197.691}{500.128}{K}%
\TextAt{8}{210.075}{501.921}{with}%
\TextAt{9}{237.019}{501.921}{0}%
\TextAt{13}{246.677}{501.921}{\Character{60}}%
\TextAt{16}{259.587}{501.921}{\Character{106}}%
\TextAt{13}{262.908}{501.921}{N}%
\TextAt{21}{272.281}{500.061}{K}%
\TextAt{3}{279.986}{500.061}{\Character{106}}%
\TextAt{25}{282.338}{500.061}{Q}%
\TextAt{9}{289.035}{501.921}{\Character{40}}%
\TextAt{13}{293.588}{501.921}{a}%
\TextAt{9}{299.733}{501.921}{\Character{41}}%
\TextAt{16}{304.285}{501.921}{\Character{106}}%
\NativeMathAt{311.398}{501.921}{11.955}{\leq}%
\TextAt{13}{324.514}{501.921}{C}%
\TextAt{8}{337.929}{501.921}{is}%
\TextAt{8}{349.977}{501.921}{associated}%
\TextAt{8}{405.996}{501.921}{to}%
\TextAt{8}{420.591}{501.921}{one}%
\TextAt{8}{442.334}{501.921}{of}%
\TextAt{8}{455.948}{501.921}{these}%
\TextAt{8}{486.204}{501.921}{numbers.}%
\TextAt{8}{79.226}{484.875}{The}%
\TextAt{8}{103.28}{484.875}{set}%
\TextAt{13}{251.733}{462.173}{T}%
\TextAt{9}{263.541}{462.173}{\Character{61}}%
\TextAt{13}{275.966}{462.173}{S}%
\TextAt{16}{286.518}{462.173}{\Character{92}}%
\TextAt{21}{300.212}{477.117}{N}%
\TextAt{20}{297.355}{473.531}{\Character{91}}%
\TextAt{21}{297.145}{448.336}{i}%
\TextAt{17}{300.028}{448.336}{\Character{61}1}%
\TextAt{13}{312.841}{462.173}{X}%
\TextAt{9}{323.497}{462.173}{\Character{40}}%
\TextAt{13}{328.049}{462.173}{j}%
\TextAt{13}{333.56}{462.173}{\Character{11}}%
\TextAt{21}{341.01}{460.38}{i}%
\TextAt{9}{344.391}{462.173}{\Character{41}}%
\TextAt{3}{348.944}{467.109}{\Character{0}}%
\TextAt{17}{355.53}{467.109}{1}%
\TextAt{8}{79.226}{432.547}{then}%
\TextAt{8}{105.682}{432.547}{has}%
\TextAt{8}{126.351}{432.547}{the}%
\TextAt{8}{146.303}{432.547}{required}%
\TextAt{8}{191.947}{432.547}{property:}%
\TextAt{8}{244.205}{432.547}{since}%
\TextAt{13}{272.669}{432.547}{X}%
\TextAt{8}{287.017}{432.547}{is}%
\TextAt{8}{298.577}{432.547}{bounded,}%
\TextAt{8}{349.462}{432.547}{so}%
\TextAt{8}{363.628}{432.547}{is}%
\TextAt{13}{375.187}{432.547}{X}%
\TextAt{9}{385.842}{432.547}{\Character{40}}%
\TextAt{13}{390.395}{432.547}{j}%
\TextAt{13}{395.906}{432.547}{\Character{11}}%
\TextAt{21}{403.356}{430.754}{i}%
\TextAt{9}{406.737}{432.547}{\Character{41}}%
\TextAt{3}{411.289}{436.886}{\Character{0}}%
\TextAt{17}{417.876}{436.886}{1}%
\TextAt{8}{426.301}{432.547}{and}%
\TextAt{8}{448.86}{432.547}{therefore}%
\TextAt{8}{497.755}{432.547}{also}%
\TextAt{13}{521.031}{432.547}{T}%
\TextAt{8}{529.518}{432.547}{,}%
\TextAt{8}{79.226}{415.501}{and}%
\TextAt{8}{101.989}{415.501}{we}%
\TextAt{8}{119.216}{415.501}{have}%
\TextAt{13}{271.567}{398.455}{S}%
\TextAt{9}{282.783}{398.455}{\Character{61}}%
\TextAt{20}{298.815}{409.813}{\Character{91}}%
\TextAt{21}{295.209}{384.361}{\Character{34}}%
\NativeMathAt{299.119}{384.361}{7.97}{\in}%
\TextAt{27}{304.765}{384.361}{o}%
\TextAt{23}{308.661}{387.106}{\Character{3}}%
\TextAt{22}{308.661}{381.587}{K}%
\TextAt{13}{317.696}{398.455}{T}%
\TextAt{13}{326.172}{398.455}{j}%
\TextAt{13}{331.684}{398.455}{\Character{34}\Character{58}}%
\TextAt{8}{79.226}{365.798}{In}%
\TextAt{8}{95.151}{365.798}{fact,}%
\TextAt{8}{123.118}{365.798}{if}%
\TextAt{13}{135.146}{365.798}{y}%
\NativeMathAt{146.826}{365.798}{11.955}{\in}%
\TextAt{13}{160.339}{365.798}{S}%
\TextAt{8}{168.234}{365.798}{,}%
\TextAt{8}{177.013}{365.798}{we}%
\TextAt{8}{195.545}{365.798}{find}%
\TextAt{8}{220.257}{365.798}{by}%
\TextAt{8}{237.82}{365.798}{the}%
\TextAt{8}{259.28}{365.798}{above}%
\TextAt{8}{293.76}{365.798}{an}%
\TextAt{13}{311.323}{365.798}{a}%
\NativeMathAt{323.01}{365.798}{11.955}{\in}%
\TextAt{14}{336.523}{365.798}{o}%
\TextAt{21}{342.368}{364.005}{K}%
\TextAt{8}{350.57}{365.798}{,}%
\TextAt{13}{359.349}{365.798}{a}%
\NativeMathAt{371.037}{365.798}{11.955}{\neq}%
\TextAt{9}{385.67}{365.798}{0}%
\TextAt{8}{391.538}{365.798}{,}%
\TextAt{8}{400.317}{365.798}{such}%
\TextAt{8}{428.018}{365.798}{that}%
\TextAt{13}{454.691}{365.798}{j}%
\TextAt{13}{460.202}{365.798}{a}%
\NativeMathAt{471.894}{365.798}{11.955}{\in}%
\TextAt{13}{485.407}{365.798}{X}%
\TextAt{13}{496.047}{365.798}{y}%
\TextAt{3}{502.199}{370.136}{\Character{0}}%
\TextAt{17}{508.785}{370.136}{1}%
\TextAt{8}{513.518}{365.798}{,}%
\TextAt{8}{522.297}{365.798}{so}%
\TextAt{13}{79.226}{348.752}{j}%
\TextAt{13}{84.737}{348.752}{a}%
\TextAt{9}{94.207}{348.752}{\Character{61}}%
\TextAt{13}{106.633}{348.752}{xy}%
\TextAt{3}{119.422}{353.09}{\Character{0}}%
\TextAt{17}{126.008}{353.09}{1}%
\TextAt{8}{134.638}{348.752}{for}%
\TextAt{8}{152.523}{348.752}{some}%
\TextAt{13}{181.849}{348.752}{x}%
\NativeMathAt{191.822}{348.752}{11.955}{\in}%
\TextAt{13}{203.113}{348.752}{X}%
\TextAt{8}{213.768}{348.752}{.}%
\TextAt{8}{222.216}{348.752}{Since}%
\TextAt{16}{189.088}{325.241}{\Character{106}}%
\TextAt{13}{192.409}{325.241}{N}%
\TextAt{21}{201.782}{323.382}{K}%
\TextAt{3}{209.487}{323.382}{\Character{106}}%
\TextAt{25}{211.839}{323.382}{Q}%
\TextAt{9}{218.536}{325.241}{\Character{40}}%
\TextAt{13}{223.089}{325.241}{a}%
\TextAt{9}{229.234}{325.241}{\Character{41}}%
\TextAt{16}{233.786}{325.241}{\Character{106}}%
\TextAt{9}{240.428}{325.241}{\Character{61}}%
\TextAt{16}{252.853}{325.241}{\Character{106}}%
\TextAt{13}{256.174}{325.241}{N}%
\TextAt{9}{266.797}{325.241}{\Character{40}}%
\TextAt{13}{271.349}{325.241}{xy}%
\TextAt{3}{284.138}{330.178}{\Character{0}}%
\TextAt{17}{290.724}{330.178}{1}%
\TextAt{9}{295.457}{325.241}{\Character{41}}%
\TextAt{16}{300.009}{325.241}{\Character{106}}%
\TextAt{9}{306.651}{325.241}{\Character{61}}%
\TextAt{16}{319.076}{325.241}{\Character{106}}%
\TextAt{13}{322.397}{325.241}{N}%
\TextAt{9}{333.02}{325.241}{\Character{40}}%
\TextAt{13}{337.572}{325.241}{x}%
\TextAt{9}{344.224}{325.241}{\Character{41}}%
\TextAt{16}{348.776}{325.241}{\Character{106}}%
\TextAt{13}{355.418}{325.241}{\Character{60}}%
\TextAt{20}{367.844}{336.599}{\Character{89}}%
\TextAt{21}{373.142}{311.493}{\Character{28}}%
\TextAt{13}{385.112}{325.241}{c}%
\TextAt{21}{390.15}{323.448}{\Character{28}}%
\TextAt{9}{398.647}{325.241}{\Character{61}}%
\TextAt{13}{411.073}{325.241}{C}%
\TextAt{13}{419.645}{325.241}{\Character{59}}%
\TextAt{13}{79.226}{294.42}{a}%
\TextAt{8}{89.268}{294.42}{is}%
\TextAt{8}{101.032}{294.42}{associated}%
\TextAt{8}{156.767}{294.42}{to}%
\TextAt{8}{171.078}{294.42}{some}%
\TextAt{13}{200.404}{294.42}{\Character{11}}%
\TextAt{21}{207.849}{292.627}{i}%
\TextAt{8}{211.231}{294.42}{,}%
\TextAt{13}{218.38}{294.42}{\Character{11}}%
\TextAt{21}{225.825}{292.627}{i}%
\TextAt{9}{232.527}{294.42}{\Character{61}}%
\TextAt{13}{244.953}{294.42}{\Character{34}a}%
\TextAt{8}{256.576}{294.42}{,}%
\TextAt{13}{263.725}{294.42}{\Character{34}}%
\NativeMathAt{272.524}{294.42}{11.955}{\in}%
\TextAt{14}{283.815}{294.42}{o}%
\TextAt{3}{289.659}{298.759}{\Character{3}}%
\TextAt{21}{289.659}{291.4}{K}%
\TextAt{8}{297.862}{294.42}{.}%
\TextAt{8}{306.31}{294.42}{Consequently}%
\TextAt{13}{247.946}{270.91}{y}%
\TextAt{9}{257.404}{270.91}{\Character{61}}%
\TextAt{13}{269.829}{270.91}{xj}%
\TextAt{13}{281.987}{270.91}{a}%
\TextAt{3}{288.142}{275.846}{\Character{0}}%
\TextAt{17}{294.728}{275.846}{1}%
\TextAt{9}{302.781}{270.91}{\Character{61}}%
\TextAt{13}{315.207}{270.91}{xj}%
\TextAt{9}{327.374}{270.91}{\Character{40}}%
\TextAt{13}{331.927}{270.91}{\Character{11}}%
\TextAt{3}{339.449}{275.858}{\Character{0}}%
\TextAt{17}{346.035}{275.858}{1}%
\TextAt{21}{339.372}{267.846}{i}%
\TextAt{13}{350.767}{270.91}{\Character{34}}%
\TextAt{9}{356.245}{270.91}{\Character{41}}%
\TextAt{13}{360.798}{270.91}{\Character{58}}%
\TextAt{8}{79.226}{247.4}{Since}%
\TextAt{13}{109.784}{247.4}{y}%
\TextAt{13}{115.905}{247.4}{\Character{59}}%
\TextAt{13}{121.141}{247.4}{j}%
\TextAt{13}{126.653}{247.4}{\Character{34}}%
\NativeMathAt{135.479}{247.4}{11.955}{\in}%
\TextAt{13}{146.77}{247.4}{S}%
\TextAt{8}{154.665}{247.4}{,}%
\TextAt{8}{161.814}{247.4}{one}%
\TextAt{8}{183.274}{247.4}{finds}%
\TextAt{13}{211.297}{247.4}{xj}%
\TextAt{13}{223.455}{247.4}{\Character{11}}%
\TextAt{3}{230.986}{252.348}{\Character{0}}%
\TextAt{17}{237.572}{252.348}{1}%
\TextAt{21}{230.91}{244.336}{i}%
\NativeMathAt{245.626}{247.4}{11.955}{\in}%
\TextAt{13}{256.917}{247.4}{S}%
\TextAt{16}{267.469}{247.4}{\Character{92}}%
\TextAt{13}{278.095}{247.4}{X}%
\TextAt{13}{288.735}{247.4}{j}%
\TextAt{13}{294.247}{247.4}{\Character{11}}%
\TextAt{3}{301.788}{252.348}{\Character{0}}%
\TextAt{17}{308.374}{252.348}{1}%
\TextAt{21}{301.712}{244.336}{i}%
\NativeMathAt{316.427}{247.4}{11.955}{\subseteq}%
\TextAt{13}{329.047}{247.4}{T}%
\TextAt{8}{337.534}{247.4}{,}%
\TextAt{8}{344.683}{247.4}{and}%
\TextAt{8}{367.445}{247.4}{thus}%
\TextAt{13}{393.197}{247.4}{y}%
\NativeMathAt{402.654}{247.4}{11.955}{\in}%
\TextAt{13}{413.945}{247.4}{T}%
\TextAt{13}{422.421}{247.4}{j}%
\TextAt{13}{427.933}{247.4}{\Character{34}}%
\TextAt{8}{433.426}{247.4}{.}%
\Vector{q 0 G 0 g 0.399 w 452.334 -526.6 m 452.334 -518.531 l S Q}%
\Vector{q 0 G 0 g 0.399 w 452.534 -518.73 m 459.707 -518.73 l S Q}%
\Vector{q 0 G 0 g 0.399 w 452.534 -526.401 m 459.707 -526.401 l S Q}%
\Vector{q 0 G 0 g 0.399 w 459.906 -526.6 m 459.906 -518.531 l S Q}%
\TextAt{8}{97.159}{219.967}{From}%
\TextAt{8}{128.601}{219.967}{proposition}%
\TextAt{8}{190.852}{219.967}{(7.1)}%
\TextAt{8}{219.544}{219.967}{and}%
\TextAt{8}{243.024}{219.967}{theorem}%
\TextAt{8}{289.266}{219.967}{(7.3)}%
\TextAt{8}{317.959}{219.967}{we}%
\TextAt{8}{335.904}{219.967}{immediately}%
\TextAt{8}{403.283}{219.967}{deduce}%
\TextAt{12}{443.01}{219.967}{Dirichlet’s}%
\TextAt{12}{508.847}{219.967}{unit}%
\TextAt{12}{79.226}{202.921}{theorem}%
\TextAt{8}{131.542}{202.921}{in}%
\TextAt{8}{145.195}{202.921}{its}%
\TextAt{8}{161.514}{202.921}{classical}%
\TextAt{8}{206.513}{202.921}{form.}%
\TextAt{12}{79.226}{178.478}{Theorem}%
\TextAt{12}{137.055}{178.478}{(7.4).}%
\TextAt{8}{173.079}{178.478}{The}%
\TextAt{8}{197.842}{178.478}{group}%
\TextAt{8}{231.726}{178.478}{of}%
\TextAt{8}{245.764}{178.478}{units}%
\TextAt{14}{275.799}{178.478}{o}%
\TextAt{3}{281.643}{182.816}{\Character{3}}%
\TextAt{21}{281.643}{175.457}{K}%
\TextAt{8}{294.451}{178.478}{of}%
\TextAt{14}{308.49}{178.478}{o}%
\TextAt{21}{314.334}{176.684}{K}%
\TextAt{8}{327.143}{178.478}{is}%
\TextAt{8}{339.615}{178.478}{the}%
\TextAt{8}{360.48}{178.478}{direct}%
\TextAt{8}{394.352}{178.478}{product}%
\TextAt{8}{438.96}{178.478}{of}%
\TextAt{8}{452.998}{178.478}{the}%
\TextAt{8}{473.863}{178.478}{finite}%
\TextAt{8}{504.484}{178.478}{cyclic}%
\TextAt{8}{79.226}{161.432}{group}%
\TextAt{13}{112.402}{161.432}{\Character{22}}%
\TextAt{9}{119.445}{161.432}{\Character{40}}%
\TextAt{13}{123.997}{161.432}{K}%
\TextAt{9}{134.815}{161.432}{\Character{41}}%
\TextAt{8}{143.264}{161.432}{and}%
\TextAt{8}{166.027}{161.432}{a}%
\TextAt{8}{175.783}{161.432}{free}%
\TextAt{8}{198.21}{161.432}{abelian}%
\TextAt{8}{238.858}{161.432}{group}%
\TextAt{8}{272.034}{161.432}{of}%
\TextAt{8}{285.364}{161.432}{rank}%
\TextAt{13}{312.358}{161.432}{r}%
\TextAt{9}{320.616}{161.432}{\Character{43}}%
\TextAt{13}{332.377}{161.432}{s}%
\TextAt{16}{340.547}{161.432}{\Character{0}}%
\TextAt{9}{352.503}{161.432}{1}%
\TextAt{8}{358.356}{161.432}{.}%
\TextAt{8}{97.159}{136.988}{In}%
\TextAt{8}{111.578}{136.988}{other}%
\TextAt{8}{141.945}{136.988}{words:}%
\TextAt{8}{179.953}{136.988}{there}%
\TextAt{8}{209.663}{136.988}{exist}%
\TextAt{8}{237.161}{136.988}{units}%
\TextAt{13}{266.284}{136.988}{\Character{34}}%
\TextAt{17}{271.762}{135.195}{1}%
\TextAt{13}{276.495}{136.988}{\Character{59}}%
\TextAt{13}{281.719}{136.988}{\Character{58}}%
\TextAt{13}{286.956}{136.988}{\Character{58}}%
\TextAt{13}{292.18}{136.988}{\Character{58}}%
\TextAt{13}{297.417}{136.988}{\Character{59}}%
\TextAt{13}{302.653}{136.988}{\Character{34}}%
\TextAt{21}{308.193}{135.195}{t}%
\TextAt{8}{311.75}{136.988}{,}%
\TextAt{13}{318.737}{136.988}{t}%
\TextAt{9}{326.285}{136.988}{\Character{61}}%
\TextAt{13}{338.711}{136.988}{r}%
\TextAt{9}{346.554}{136.988}{\Character{43}}%
\TextAt{13}{357.902}{136.988}{s}%
\TextAt{16}{365.659}{136.988}{\Character{0}}%
\TextAt{9}{377.201}{136.988}{1}%
\TextAt{8}{383.054}{136.988}{,}%
\TextAt{8}{390.041}{136.988}{called}%
\TextAt{12}{423.003}{136.988}{fundamental}%
\TextAt{12}{500.287}{136.988}{units}%
\TextAt{8}{529.518}{136.988}{,}%
\TextAt{8}{79.226}{119.942}{such}%
\TextAt{8}{105.623}{119.942}{that}%
\TextAt{8}{130.992}{119.942}{any}%
\TextAt{8}{153.109}{119.942}{other}%
\TextAt{8}{183.678}{119.942}{unit}%
\TextAt{13}{208.39}{119.942}{\Character{34}}%
\TextAt{8}{217.765}{119.942}{can}%
\TextAt{8}{239.225}{119.942}{be}%
\TextAt{8}{255.149}{119.942}{written}%
\TextAt{8}{296.119}{119.942}{uniquely}%
\TextAt{8}{343.593}{119.942}{as}%
\TextAt{8}{357.963}{119.942}{a}%
\TextAt{8}{367.719}{119.942}{product}%
\TextAt{13}{269.523}{96.432}{\Character{34}}%
\TextAt{9}{278.322}{96.432}{\Character{61}}%
\TextAt{13}{290.747}{96.432}{\Character{16}}%
\TextAt{13}{296.749}{96.432}{\Character{34}}%
\TextAt{21}{302.241}{101.657}{\Character{23}}%
\TextAt{24}{306.404}{100.55}{1}%
\TextAt{17}{302.241}{93.501}{1}%
\TextAt{16}{313.045}{96.432}{\Character{1}}%
\TextAt{16}{318.341}{96.432}{\Character{1}}%
\TextAt{16}{323.649}{96.432}{\Character{1}}%
\TextAt{13}{328.985}{96.432}{\Character{34}}%
\TextAt{21}{334.463}{101.368}{\Character{23}}%
\TextAt{22}{338.626}{100.372}{t}%
\TextAt{21}{334.463}{93.476}{t}%
\TextAt{8}{79.226}{72.921}{with}%
\TextAt{8}{105.886}{72.921}{a}%
\TextAt{8}{115.642}{72.921}{root}%
\TextAt{8}{140.688}{72.921}{of}%
\TextAt{8}{154.018}{72.921}{unity}%
\TextAt{13}{184.587}{72.921}{\Character{16}}%
\TextAt{8}{194.5}{72.921}{and}%
\TextAt{8}{217.263}{72.921}{integers}%
\TextAt{13}{260.576}{72.921}{\Character{23}}%
\TextAt{21}{266.36}{71.128}{i}%
\TextAt{8}{269.741}{72.921}{.}%
\TextAt{8}{300.14}{43.033}{39}%
\Destination{78.23}{809.815}{nav.page.39}%
\Destination{79.36}{380.755}{nav.print.42}%
\Destination{79.36}{192.455}{nav.number.1.7.4}%
\Hotlink{191.348}{563.755}{25.078}{11.047}{nav.number.1.5.3}%
\Hotlink{214.652}{516.755}{25.078}{11.047}{nav.number.1.7.2}%
\Hotlink{190.352}{217.755}{25.078}{11.047}{nav.number.1.7.1}%
\Hotlink{288.766}{217.755}{25.078}{11.047}{nav.number.1.7.3}%
}

% Page 52 (body page 40)
\Page{%
\TextAt{2}{79.226}{796.943}{I.}%
\TextAt{2}{89.83}{796.943}{Algebraic}%
\TextAt{2}{138.648}{796.943}{Integers}%
\Vector{q 0 G 0 g 0.299 w 7.226 18.729 m 460.769 18.729 l S Q}%
\TextAt{19}{79.226}{761.124}{Proof.}%
\TextAt{8}{116}{761.124}{In}%
\TextAt{8}{130.621}{761.124}{the}%
\TextAt{8}{150.778}{761.124}{exact}%
\TextAt{8}{181.67}{761.124}{sequence}%
\TextAt{9}{227.829}{733.29}{1}%
\NativeArrowAt{237.003}{733.29}{11.955}{19.26}{\longrightarrow}%
\TextAt{13}{259.585}{733.29}{\Character{22}}%
\TextAt{9}{266.628}{733.29}{\Character{40}}%
\TextAt{13}{271.18}{733.29}{K}%
\TextAt{9}{281.998}{733.29}{\Character{41}}%
\NativeArrowAt{289.871}{733.29}{11.955}{19.26}{\longrightarrow}%
\TextAt{14}{312.453}{733.29}{o}%
\TextAt{3}{318.297}{738.226}{\Character{3}}%
\TextAt{21}{318.297}{730.334}{K}%
\TextAt{21}{332.623}{740.067}{\Character{21}}%
\NativeArrowAt{329.821}{733.29}{11.955}{11.955}{\longrightarrow}%
\TextAt{9}{345.097}{733.29}{\Character{0}}%
\NativeArrowAt{355.732}{733.29}{11.955}{19.26}{\longrightarrow}%
\TextAt{9}{378.314}{733.29}{0}%
\TextAt{9}{79.226}{705.455}{\Character{0}}%
\TextAt{8}{90.625}{705.455}{is}%
\TextAt{8}{102.576}{705.455}{a}%
\TextAt{8}{112.519}{705.455}{free}%
\TextAt{8}{135.134}{705.455}{abelian}%
\TextAt{8}{175.969}{705.455}{group}%
\TextAt{8}{209.331}{705.455}{of}%
\TextAt{8}{222.849}{705.455}{rank}%
\TextAt{13}{250.031}{705.455}{t}%
\TextAt{9}{257.898}{705.455}{\Character{61}}%
\TextAt{13}{270.642}{705.455}{r}%
\TextAt{9}{279.027}{705.455}{\Character{43}}%
\TextAt{13}{290.916}{705.455}{s}%
\TextAt{16}{299.214}{705.455}{\Character{0}}%
\TextAt{9}{311.297}{705.455}{1}%
\TextAt{8}{321.234}{705.455}{by}%
\TextAt{8}{337.681}{705.455}{(7.3).}%
\TextAt{8}{370.768}{705.455}{Let}%
\TextAt{13}{391.925}{705.455}{v}%
\TextAt{17}{397.584}{703.662}{1}%
\TextAt{13}{402.316}{705.455}{\Character{59}}%
\TextAt{13}{407.54}{705.455}{\Character{58}}%
\TextAt{13}{412.777}{705.455}{\Character{58}}%
\TextAt{13}{418.013}{705.455}{\Character{58}}%
\TextAt{13}{423.238}{705.455}{\Character{59}}%
\TextAt{13}{428.474}{705.455}{v}%
\TextAt{21}{434.196}{703.662}{t}%
\TextAt{8}{441.837}{705.455}{be}%
\TextAt{8}{457.949}{705.455}{a}%
\TextAt{18}{467.891}{705.455}{Z}%
\TextAt{8}{475.862}{705.455}{-basis}%
\TextAt{8}{508.686}{705.455}{of}%
\TextAt{9}{522.204}{705.455}{\Character{0}}%
\TextAt{8}{529.518}{705.455}{,}%
\TextAt{8}{79.226}{688.409}{let}%
\TextAt{13}{96.293}{688.409}{\Character{34}}%
\TextAt{17}{101.771}{686.616}{1}%
\TextAt{13}{106.504}{688.409}{\Character{59}}%
\TextAt{13}{111.728}{688.409}{\Character{58}}%
\TextAt{13}{116.965}{688.409}{\Character{58}}%
\TextAt{13}{122.189}{688.409}{\Character{58}}%
\TextAt{13}{127.426}{688.409}{\Character{59}}%
\TextAt{13}{132.662}{688.409}{\Character{34}}%
\TextAt{21}{138.202}{686.616}{t}%
\NativeMathAt{145.356}{688.409}{11.955}{\in}%
\TextAt{14}{156.924}{688.409}{o}%
\TextAt{3}{162.768}{692.748}{\Character{3}}%
\TextAt{21}{162.768}{685.389}{K}%
\TextAt{8}{175.031}{688.409}{be}%
\TextAt{8}{191.118}{688.409}{preimages}%
\TextAt{8}{245.975}{688.409}{of}%
\TextAt{8}{259.468}{688.409}{the}%
\TextAt{13}{279.787}{688.409}{v}%
\TextAt{21}{285.446}{686.616}{i}%
\TextAt{8}{288.827}{688.409}{,}%
\TextAt{8}{296.179}{688.409}{and}%
\TextAt{8}{319.105}{688.409}{let}%
\TextAt{13}{336.172}{688.409}{A}%
\NativeMathAt{348.545}{688.409}{11.955}{\subseteq}%
\TextAt{14}{361.441}{688.409}{o}%
\TextAt{3}{367.285}{692.748}{\Character{3}}%
\TextAt{21}{367.285}{685.389}{K}%
\TextAt{8}{379.548}{688.409}{be}%
\TextAt{8}{395.634}{688.409}{the}%
\TextAt{8}{415.953}{688.409}{subgroup}%
\TextAt{8}{466.913}{688.409}{generated}%
\TextAt{8}{520.408}{688.409}{by}%
\TextAt{8}{79.226}{671.363}{the}%
\TextAt{13}{99.435}{671.363}{\Character{34}}%
\TextAt{21}{104.913}{669.57}{i}%
\TextAt{8}{108.294}{671.363}{.}%
\TextAt{8}{116.9}{671.363}{Then}%
\TextAt{13}{147.51}{671.363}{A}%
\TextAt{8}{160.236}{671.363}{is}%
\TextAt{8}{172.052}{671.363}{mapped}%
\TextAt{8}{216.65}{671.363}{isomorphically}%
\TextAt{8}{294.495}{671.363}{onto}%
\TextAt{9}{320.896}{671.363}{\Character{0}}%
\TextAt{8}{332.161}{671.363}{by}%
\TextAt{13}{348.472}{671.363}{\Character{21}}%
\TextAt{8}{355.301}{671.363}{,}%
\TextAt{8}{362.516}{671.363}{i.e.,}%
\TextAt{8}{384.686}{671.363}{one}%
\TextAt{8}{406.198}{671.363}{has}%
\TextAt{13}{427.125}{671.363}{\Character{22}}%
\TextAt{9}{434.168}{671.363}{\Character{40}}%
\TextAt{13}{438.72}{671.363}{K}%
\TextAt{9}{449.538}{671.363}{\Character{41}}%
\TextAt{16}{456.783}{671.363}{\Character{92}}%
\TextAt{13}{467.445}{671.363}{A}%
\TextAt{9}{479.631}{671.363}{\Character{61}}%
\TextAt{16}{492.146}{671.363}{\Character{102}}%
\TextAt{9}{498.124}{671.363}{1}%
\TextAt{16}{503.977}{671.363}{\Character{103}}%
\TextAt{8}{513.904}{671.363}{and}%
\TextAt{8}{79.226}{654.317}{therefore}%
\TextAt{14}{128.326}{654.317}{o}%
\TextAt{3}{134.17}{658.656}{\Character{3}}%
\TextAt{21}{134.17}{651.297}{K}%
\TextAt{9}{145.694}{654.318}{\Character{61}}%
\TextAt{13}{158.119}{654.317}{\Character{22}}%
\TextAt{9}{165.162}{654.318}{\Character{40}}%
\TextAt{13}{169.715}{654.317}{K}%
\TextAt{9}{180.532}{654.318}{\Character{41}}%
\TextAt{16}{187.741}{654.317}{\Character{2}}%
\TextAt{13}{199.696}{654.318}{A}%
\TextAt{8}{208.472}{654.317}{.}%
\Vector{q 0 G 0 g 0.399 w 452.334 -119.683 m 452.334 -111.613 l S Q}%
\Vector{q 0 G 0 g 0.399 w 452.534 -111.812 m 459.707 -111.812 l S Q}%
\Vector{q 0 G 0 g 0.399 w 452.534 -119.483 m 459.707 -119.483 l S Q}%
\Vector{q 0 G 0 g 0.399 w 459.906 -119.683 m 459.906 -111.613 l S Q}%
\TextAt{8}{97.159}{625.65}{Identifying}%
\TextAt{9}{157.437}{625.65}{\Character{91}}%
\TextAt{20}{160.689}{634.616}{\Character{81}}%
\TextAt{21}{171.98}{622.163}{\Character{28}}%
\TextAt{18}{179.149}{625.65}{R}%
\TextAt{9}{187.783}{625.65}{\Character{93}}%
\TextAt{17}{191.035}{629.988}{\Character{43}}%
\TextAt{9}{203.315}{625.65}{\Character{61}}%
\TextAt{18}{217.615}{625.65}{R}%
\TextAt{21}{226.249}{629.988}{r}%
\TextAt{17}{230.306}{629.988}{\Character{43}}%
\TextAt{21}{236.892}{629.988}{s}%
\TextAt{8}{246.304}{625.65}{(see}%
\TextAt{8}{270.872}{625.65}{§5,}%
\TextAt{8}{290.933}{625.65}{p.33),}%
\TextAt{13}{325.483}{625.65}{H}%
\TextAt{8}{341.103}{625.65}{becomes}%
\TextAt{8}{388.757}{625.65}{a}%
\TextAt{8}{399.612}{625.65}{subspace}%
\TextAt{8}{449.609}{625.65}{of}%
\TextAt{8}{464.039}{625.65}{the}%
\TextAt{8}{485.295}{625.65}{euclidean}%
\TextAt{8}{79.226}{608.604}{space}%
\TextAt{18}{111.763}{608.604}{R}%
\TextAt{21}{120.397}{612.942}{r}%
\TextAt{17}{124.454}{612.942}{\Character{43}}%
\TextAt{21}{131.04}{612.942}{s}%
\TextAt{8}{135.455}{608.604}{,}%
\TextAt{8}{144.181}{608.604}{and}%
\TextAt{8}{168.206}{608.604}{thus}%
\TextAt{8}{195.22}{608.604}{itself}%
\TextAt{8}{224.828}{608.604}{a}%
\TextAt{8}{235.846}{608.604}{euclidean}%
\TextAt{8}{288.479}{608.604}{space.}%
\TextAt{8}{328.091}{608.604}{We}%
\TextAt{8}{349.498}{608.604}{may}%
\TextAt{8}{376.129}{608.604}{therefore}%
\TextAt{8}{426.491}{608.604}{speak}%
\TextAt{8}{460.331}{608.604}{of}%
\TextAt{8}{474.923}{608.604}{the}%
\TextAt{8}{496.342}{608.604}{volume}%
\TextAt{8}{79.226}{591.558}{of}%
\TextAt{8}{93.505}{591.558}{the}%
\TextAt{8}{114.61}{591.558}{fundamental}%
\TextAt{8}{183.201}{591.558}{mesh}%
\TextAt{8}{214.122}{591.558}{vol}%
\TextAt{9}{229.413}{591.558}{\Character{40}}%
\TextAt{13}{233.965}{591.558}{\Character{21}}%
\TextAt{9}{240.793}{591.558}{\Character{40}}%
\TextAt{14}{245.346}{591.558}{o}%
\TextAt{3}{251.19}{595.896}{\Character{3}}%
\TextAt{21}{251.19}{588.537}{K}%
\TextAt{9}{259.393}{591.558}{\Character{41}\Character{41}}%
\TextAt{8}{273.343}{591.558}{of}%
\TextAt{8}{287.622}{591.558}{the}%
\TextAt{8}{308.728}{591.558}{unit}%
\TextAt{8}{334.388}{591.558}{lattice}%
\TextAt{9}{371.107}{591.558}{\Character{0}}%
\TextAt{9}{383.337}{591.558}{\Character{61}}%
\TextAt{13}{397.401}{591.558}{\Character{21}}%
\TextAt{9}{404.229}{591.558}{\Character{40}}%
\TextAt{14}{408.782}{591.558}{o}%
\TextAt{3}{414.626}{595.896}{\Character{3}}%
\TextAt{21}{414.626}{588.537}{K}%
\TextAt{9}{422.829}{591.558}{\Character{41}}%
\NativeMathAt{432.319}{591.558}{11.955}{\subseteq}%
\TextAt{13}{446.555}{591.558}{H}%
\TextAt{8}{457.178}{591.558}{,}%
\TextAt{8}{465.513}{591.558}{and}%
\TextAt{8}{489.224}{591.558}{will}%
\TextAt{8}{512.278}{591.558}{now}%
\TextAt{8}{79.226}{574.512}{compute}%
\TextAt{8}{127.483}{574.512}{it.}%
\TextAt{8}{146.087}{574.512}{Let}%
\TextAt{13}{167.84}{574.512}{\Character{34}}%
\TextAt{17}{173.318}{572.719}{1}%
\TextAt{13}{178.05}{574.512}{\Character{59}}%
\TextAt{13}{183.274}{574.512}{\Character{58}}%
\TextAt{13}{188.511}{574.512}{\Character{58}}%
\TextAt{13}{193.735}{574.512}{\Character{58}}%
\TextAt{13}{198.972}{574.512}{\Character{59}}%
\TextAt{13}{204.208}{574.512}{\Character{34}}%
\TextAt{21}{209.749}{572.719}{t}%
\TextAt{8}{213.305}{574.512}{,}%
\TextAt{13}{221.433}{574.512}{t}%
\TextAt{9}{230.315}{574.512}{\Character{61}}%
\TextAt{13}{244.075}{574.512}{r}%
\TextAt{9}{252.866}{574.512}{\Character{43}}%
\TextAt{13}{265.161}{574.512}{s}%
\TextAt{16}{273.866}{574.512}{\Character{0}}%
\TextAt{9}{286.354}{574.512}{1}%
\TextAt{8}{292.207}{574.512}{,}%
\TextAt{8}{300.335}{574.512}{be}%
\TextAt{8}{317.043}{574.512}{a}%
\TextAt{8}{327.581}{574.512}{system}%
\TextAt{8}{367.183}{574.512}{of}%
\TextAt{8}{381.296}{574.512}{fundamental}%
\TextAt{8}{449.721}{574.512}{units}%
\TextAt{8}{479.83}{574.512}{and}%
\TextAt{9}{503.376}{574.512}{\Character{8}}%
\TextAt{8}{516.51}{574.512}{the}%
\TextAt{8}{79.226}{557.466}{fundamental}%
\TextAt{8}{146.616}{557.466}{mesh}%
\TextAt{8}{176.336}{557.466}{of}%
\TextAt{8}{189.413}{557.466}{the}%
\TextAt{8}{209.317}{557.466}{unit}%
\TextAt{8}{233.776}{557.466}{lattice}%
\TextAt{13}{269.294}{557.466}{\Character{21}}%
\TextAt{9}{276.122}{557.466}{\Character{40}}%
\TextAt{14}{280.675}{557.466}{o}%
\TextAt{3}{286.519}{561.804}{\Character{3}}%
\TextAt{21}{286.519}{554.445}{K}%
\TextAt{9}{294.722}{557.466}{\Character{41}}%
\TextAt{8}{299.274}{557.466}{,}%
\TextAt{8}{306.221}{557.466}{spanned}%
\TextAt{8}{351.554}{557.466}{by}%
\TextAt{8}{367.561}{557.466}{the}%
\TextAt{8}{387.465}{557.466}{vectors}%
\TextAt{13}{426.952}{557.466}{\Character{21}}%
\TextAt{9}{433.78}{557.466}{\Character{40}}%
\TextAt{13}{438.332}{557.466}{\Character{34}}%
\TextAt{17}{443.81}{555.673}{1}%
\TextAt{9}{448.543}{557.466}{\Character{41}}%
\TextAt{13}{453.095}{557.466}{\Character{59}}%
\TextAt{13}{458.319}{557.466}{\Character{58}}%
\TextAt{13}{463.556}{557.466}{\Character{58}}%
\TextAt{13}{468.792}{557.466}{\Character{58}}%
\TextAt{13}{474.017}{557.466}{\Character{59}}%
\TextAt{13}{479.253}{557.466}{\Character{21}}%
\TextAt{9}{486.144}{557.466}{\Character{40}}%
\TextAt{13}{490.697}{557.466}{\Character{34}}%
\TextAt{21}{496.175}{555.673}{t}%
\TextAt{9}{499.731}{557.466}{\Character{41}}%
\NativeMathAt{507.604}{557.466}{11.955}{\in}%
\TextAt{13}{518.895}{557.466}{H}%
\TextAt{8}{529.518}{557.466}{.}%
\TextAt{8}{79.226}{540.42}{The}%
\TextAt{8}{103.28}{540.42}{vector}%
\TextAt{13}{230.734}{523.374}{\Character{21}}%
\TextAt{17}{237.562}{521.581}{0}%
\TextAt{9}{245.615}{523.374}{\Character{61}}%
\TextAt{9}{274.058}{531.462}{1}%
\Vector{q 0 G 0 g 0.478 w 187.236 -247.637 m 222.732 -247.637 l S Q}%
\TextAt{16}{259.236}{523.733}{\Character{112}}%
\Vector{q 0 G 0 g 0.478 w 197.199 -250.028 m 222.732 -250.028 l S Q}%
\TextAt{13}{269.199}{514.707}{r}%
\TextAt{9}{277.456}{514.707}{\Character{43}}%
\TextAt{13}{289.218}{514.707}{s}%
\TextAt{9}{295.927}{523.374}{\Character{40}1}%
\TextAt{13}{306.332}{523.374}{\Character{59}}%
\TextAt{13}{311.556}{523.374}{\Character{58}}%
\TextAt{13}{316.793}{523.374}{\Character{58}}%
\TextAt{13}{322.029}{523.374}{\Character{58}}%
\TextAt{13}{327.254}{523.374}{\Character{59}}%
\TextAt{9}{332.553}{523.374}{1\Character{41}}%
\NativeMathAt{346.279}{523.374}{11.955}{\in}%
\TextAt{18}{357.57}{523.374}{R}%
\TextAt{21}{366.205}{528.31}{r}%
\TextAt{17}{370.261}{528.31}{\Character{43}}%
\TextAt{21}{376.848}{528.31}{s}%
\TextAt{8}{79.226}{496.867}{is}%
\TextAt{8}{91.057}{496.867}{obviously}%
\TextAt{8}{142.901}{496.867}{orthogonal}%
\TextAt{8}{201.524}{496.867}{to}%
\TextAt{13}{215.901}{496.867}{H}%
\TextAt{8}{230.487}{496.867}{and}%
\TextAt{8}{253.316}{496.867}{has}%
\TextAt{8}{274.257}{496.867}{length}%
\TextAt{8}{310.093}{496.867}{1.}%
\TextAt{8}{324.599}{496.867}{The}%
\TextAt{13}{348.719}{496.867}{t}%
\TextAt{8}{352.947}{496.867}{-dimensional}%
\TextAt{8}{421.361}{496.867}{volume}%
\TextAt{8}{461.752}{496.867}{of}%
\TextAt{9}{475.149}{496.867}{\Character{8}}%
\TextAt{8}{487.567}{496.867}{therefore}%
\TextAt{8}{79.226}{479.822}{equals}%
\TextAt{8}{114.047}{479.822}{the}%
\TextAt{9}{133.518}{479.822}{\Character{40}}%
\TextAt{13}{138.07}{479.822}{t}%
\TextAt{9}{143.552}{479.822}{\Character{43}}%
\TextAt{9}{153.893}{479.822}{1\Character{41}}%
\TextAt{8}{164.316}{479.822}{-dimensional}%
\TextAt{8}{231.978}{479.822}{volume}%
\TextAt{8}{271.617}{479.822}{of}%
\TextAt{8}{284.261}{479.822}{the}%
\TextAt{8}{303.732}{479.822}{parallelepiped}%
\TextAt{8}{378.16}{479.822}{spanned}%
\TextAt{8}{423.06}{479.822}{by}%
\TextAt{13}{438.633}{479.822}{\Character{21}}%
\TextAt{17}{445.461}{478.028}{0}%
\TextAt{13}{450.194}{479.822}{\Character{59}}%
\TextAt{13}{455.418}{479.822}{\Character{21}}%
\TextAt{9}{462.266}{479.822}{\Character{40}}%
\TextAt{13}{466.819}{479.822}{\Character{34}}%
\TextAt{17}{472.297}{478.028}{1}%
\TextAt{9}{477.029}{479.822}{\Character{41}}%
\TextAt{13}{481.581}{479.822}{\Character{59}}%
\TextAt{13}{486.805}{479.822}{\Character{58}}%
\TextAt{13}{492.042}{479.822}{\Character{58}}%
\TextAt{13}{497.278}{479.822}{\Character{58}}%
\TextAt{13}{502.503}{479.822}{\Character{59}}%
\TextAt{13}{507.739}{479.822}{\Character{21}}%
\TextAt{9}{514.631}{479.822}{\Character{40}}%
\TextAt{13}{519.183}{479.822}{\Character{34}}%
\TextAt{21}{524.661}{478.028}{t}%
\TextAt{9}{528.217}{479.822}{\Character{41}}%
\TextAt{8}{79.226}{462.776}{in}%
\TextAt{18}{92.879}{462.776}{R}%
\TextAt{21}{101.513}{467.114}{t}%
\TextAt{17}{104.571}{467.114}{\Character{43}1}%
\TextAt{8}{115.89}{462.776}{.}%
\TextAt{8}{124.338}{462.776}{But}%
\TextAt{8}{147.579}{462.776}{this}%
\TextAt{8}{170.402}{462.776}{has}%
\TextAt{8}{191.275}{462.776}{volume}%
\NativeMathAt{202.398}{420.726}{11.955}{\pm}%
\TextAt{8}{213.689}{420.726}{det}%
\TextAt{20}{231.94}{448.343}{\Character{48}}%
\TextAt{20}{231.94}{427.302}{\Character{66}}%
\TextAt{20}{231.94}{419.65}{\Character{64}}%
\TextAt{13}{248.399}{438.939}{\Character{21}}%
\TextAt{17}{255.228}{437.146}{01}%
\TextAt{13}{284.978}{438.939}{\Character{21}}%
\TextAt{17}{291.806}{437.146}{1}%
\TextAt{9}{296.539}{438.939}{\Character{40}}%
\TextAt{13}{301.091}{438.939}{\Character{34}}%
\TextAt{17}{306.569}{437.146}{1}%
\TextAt{9}{311.301}{438.939}{\Character{41}}%
\TextAt{16}{330.639}{438.939}{\Character{1}}%
\TextAt{16}{335.935}{438.939}{\Character{1}}%
\TextAt{16}{341.243}{438.939}{\Character{1}}%
\TextAt{13}{359.371}{438.939}{\Character{21}}%
\TextAt{17}{366.2}{437.146}{1}%
\TextAt{9}{370.932}{438.939}{\Character{40}}%
\TextAt{13}{375.484}{438.939}{\Character{34}}%
\TextAt{21}{380.962}{437.146}{t}%
\TextAt{9}{384.518}{438.939}{\Character{41}}%
\TextAt{8}{254.671}{426.688}{.}%
\TextAt{8}{254.671}{422.703}{.}%
\TextAt{8}{254.671}{418.718}{.}%
\TextAt{8}{298.79}{426.688}{.}%
\TextAt{8}{298.79}{422.703}{.}%
\TextAt{8}{298.79}{418.718}{.}%
\TextAt{8}{372.595}{426.688}{.}%
\TextAt{8}{372.595}{422.703}{.}%
\TextAt{8}{372.595}{418.718}{.}%
\TextAt{13}{242.401}{401.672}{\Character{21}}%
\TextAt{17}{249.23}{399.879}{0}%
\TextAt{21}{253.464}{399.879}{\Character{59}t}%
\TextAt{17}{258.874}{399.879}{\Character{43}1}%
\TextAt{13}{280.156}{401.672}{\Character{21}}%
\TextAt{21}{286.984}{399.879}{t}%
\TextAt{17}{290.042}{399.879}{\Character{43}1}%
\TextAt{9}{301.361}{401.672}{\Character{40}}%
\TextAt{13}{305.913}{401.672}{\Character{34}}%
\TextAt{17}{311.391}{399.879}{1}%
\TextAt{9}{316.124}{401.672}{\Character{41}}%
\TextAt{16}{330.639}{401.672}{\Character{1}}%
\TextAt{16}{335.935}{401.672}{\Character{1}}%
\TextAt{16}{341.243}{401.672}{\Character{1}}%
\TextAt{13}{354.549}{401.672}{\Character{21}}%
\TextAt{21}{361.377}{399.879}{t}%
\TextAt{17}{364.435}{399.879}{\Character{43}1}%
\TextAt{9}{375.754}{401.672}{\Character{40}}%
\TextAt{13}{380.306}{401.672}{\Character{34}}%
\TextAt{21}{385.784}{399.879}{t}%
\TextAt{9}{389.341}{401.672}{\Character{41}}%
\TextAt{20}{393.893}{448.343}{\Character{49}}%
\TextAt{20}{393.893}{427.302}{\Character{67}}%
\TextAt{20}{393.893}{419.65}{\Character{65}}%
\TextAt{13}{406.346}{420.726}{\Character{58}}%
\TextAt{8}{79.226}{376.238}{Adding}%
\TextAt{8}{120.328}{376.238}{all}%
\TextAt{8}{136.718}{376.238}{rows}%
\TextAt{8}{163.904}{376.238}{to}%
\TextAt{8}{178.346}{376.238}{a}%
\TextAt{8}{188.233}{376.238}{fixed}%
\TextAt{8}{216.65}{376.238}{one,}%
\TextAt{8}{241.525}{376.238}{say}%
\TextAt{8}{261.885}{376.238}{the}%
\TextAt{13}{282.172}{376.238}{i}%
\TextAt{8}{286.166}{376.238}{-th}%
\TextAt{8}{305.151}{376.238}{row,}%
\TextAt{8}{331.006}{376.238}{this}%
\TextAt{8}{353.96}{376.238}{row}%
\TextAt{8}{376.531}{376.238}{has}%
\TextAt{8}{397.536}{376.238}{only}%
\TextAt{8}{423.359}{376.238}{zeroes,}%
\TextAt{8}{461.625}{376.238}{except}%
\TextAt{8}{498.494}{376.238}{for}%
\TextAt{8}{516.51}{376.238}{the}%
\TextAt{8}{79.226}{359.192}{first}%
\TextAt{8}{103.352}{359.192}{entry,}%
\TextAt{8}{136.193}{359.192}{which}%
\TextAt{8}{169.679}{359.192}{equals}%
\TextAt{16}{205.186}{368.218}{\Character{112}}%
\Vector{q 0 G 0 g 0.478 w 143.148 -405.543 m 168.681 -405.543 l S Q}%
\TextAt{13}{215.148}{359.192}{r}%
\TextAt{9}{223.406}{359.192}{\Character{43}}%
\TextAt{13}{235.167}{359.192}{s}%
\TextAt{8}{240.681}{359.192}{.}%
\TextAt{8}{249.129}{359.192}{We}%
\TextAt{8}{269.274}{359.192}{therefore}%
\TextAt{8}{318.374}{359.192}{get}%
\TextAt{8}{337.884}{359.192}{the}%
\TextAt{12}{79.226}{333.513}{Proposition}%
\TextAt{12}{151.768}{333.513}{(7.5).}%
\TextAt{8}{186.937}{333.513}{The}%
\TextAt{8}{210.843}{333.513}{volume}%
\TextAt{8}{251.021}{333.513}{of}%
\TextAt{8}{264.203}{333.513}{the}%
\TextAt{8}{284.213}{333.513}{fundamental}%
\TextAt{8}{351.708}{333.513}{mesh}%
\TextAt{8}{381.532}{333.513}{of}%
\TextAt{8}{394.715}{333.513}{the}%
\TextAt{8}{414.724}{333.513}{unit}%
\TextAt{8}{439.288}{333.513}{lattice}%
\TextAt{13}{474.911}{333.513}{\Character{21}}%
\TextAt{9}{481.74}{333.513}{\Character{40}}%
\TextAt{14}{486.292}{333.513}{o}%
\TextAt{3}{492.136}{337.852}{\Character{3}}%
\TextAt{21}{492.136}{330.493}{K}%
\TextAt{9}{500.339}{333.513}{\Character{41}}%
\TextAt{8}{508.641}{333.513}{in}%
\TextAt{13}{522.147}{333.513}{H}%
\TextAt{8}{79.226}{316.467}{is}%
\TextAt{8}{246.063}{299.421}{vol}%
\TextAt{9}{261.354}{299.421}{\Character{40}}%
\TextAt{13}{265.906}{299.421}{\Character{21}}%
\TextAt{9}{272.735}{299.421}{\Character{40}}%
\TextAt{14}{277.287}{299.421}{o}%
\TextAt{3}{283.131}{304.357}{\Character{3}}%
\TextAt{21}{283.131}{296.466}{K}%
\TextAt{9}{291.334}{299.421}{\Character{41}\Character{41}}%
\TextAt{9}{303.732}{299.421}{\Character{61}}%
\TextAt{16}{316.185}{309.031}{\Character{112}}%
\Vector{q 0 G 0 g 0.478 w 254.147 -464.73 m 279.68 -464.73 l S Q}%
\TextAt{13}{326.147}{299.421}{r}%
\TextAt{9}{334.405}{299.421}{\Character{43}}%
\TextAt{13}{346.166}{299.421}{s}%
\TextAt{13}{353.662}{299.421}{R\Character{59}}%
\TextAt{8}{79.226}{276.4}{where}%
\TextAt{13}{112.148}{276.4}{R}%
\TextAt{8}{124.167}{276.4}{is}%
\TextAt{8}{135.044}{276.4}{the}%
\TextAt{8}{154.314}{276.4}{absolute}%
\TextAt{8}{199.669}{276.4}{value}%
\TextAt{8}{229.029}{276.4}{of}%
\TextAt{8}{241.472}{276.4}{the}%
\TextAt{8}{260.742}{276.4}{determinant}%
\TextAt{8}{325.871}{276.4}{of}%
\TextAt{8}{338.314}{276.4}{an}%
\TextAt{8}{353.686}{276.4}{arbitrary}%
\TextAt{8}{402.569}{276.4}{minor}%
\TextAt{8}{435.503}{276.4}{of}%
\TextAt{8}{447.946}{276.4}{rank}%
\TextAt{13}{474.054}{276.4}{t}%
\TextAt{9}{481.602}{276.4}{\Character{61}}%
\TextAt{13}{494.027}{276.4}{r}%
\TextAt{9}{500.471}{276.4}{\Character{43}}%
\TextAt{13}{510.418}{276.4}{s}%
\TextAt{16}{516.775}{276.4}{\Character{0}}%
\TextAt{9}{526.916}{276.4}{1}%
\TextAt{8}{79.226}{259.354}{of}%
\TextAt{8}{92.556}{259.354}{the}%
\TextAt{8}{112.713}{259.354}{following}%
\TextAt{8}{162.147}{259.354}{matrix:}%
\TextAt{20}{236.046}{253.378}{\Character{48}}%
\TextAt{20}{236.046}{232.337}{\Character{66}}%
\TextAt{20}{236.046}{224.685}{\Character{64}}%
\TextAt{13}{251.329}{243.975}{\Character{21}}%
\TextAt{17}{258.158}{242.181}{1}%
\TextAt{9}{262.89}{243.975}{\Character{40}}%
\TextAt{13}{267.442}{243.975}{\Character{34}}%
\TextAt{17}{272.92}{242.181}{1}%
\TextAt{9}{277.653}{243.975}{\Character{41}}%
\TextAt{16}{296.99}{243.975}{\Character{1}}%
\TextAt{16}{302.286}{243.975}{\Character{1}}%
\TextAt{16}{307.594}{243.975}{\Character{1}}%
\TextAt{13}{325.723}{243.975}{\Character{21}}%
\TextAt{17}{332.551}{242.181}{1}%
\TextAt{9}{337.283}{243.975}{\Character{40}}%
\TextAt{13}{341.836}{243.975}{\Character{34}}%
\TextAt{21}{347.314}{242.181}{t}%
\TextAt{9}{350.87}{243.975}{\Character{41}}%
\TextAt{8}{265.141}{231.723}{.}%
\TextAt{8}{265.141}{227.738}{.}%
\TextAt{8}{265.141}{223.753}{.}%
\TextAt{8}{338.946}{231.723}{.}%
\TextAt{8}{338.946}{227.738}{.}%
\TextAt{8}{338.946}{223.753}{.}%
\TextAt{13}{246.507}{206.707}{\Character{21}}%
\TextAt{21}{253.336}{204.914}{t}%
\TextAt{17}{256.394}{204.914}{\Character{43}1}%
\TextAt{9}{267.712}{206.707}{\Character{40}}%
\TextAt{13}{272.265}{206.707}{\Character{34}}%
\TextAt{17}{277.743}{204.914}{1}%
\TextAt{9}{282.475}{206.707}{\Character{41}}%
\TextAt{16}{296.99}{206.707}{\Character{1}}%
\TextAt{16}{302.286}{206.707}{\Character{1}}%
\TextAt{16}{307.594}{206.707}{\Character{1}}%
\TextAt{13}{320.9}{206.707}{\Character{21}}%
\TextAt{21}{327.729}{204.914}{t}%
\TextAt{17}{330.787}{204.914}{\Character{43}1}%
\TextAt{9}{342.106}{206.707}{\Character{40}}%
\TextAt{13}{346.658}{206.707}{\Character{34}}%
\TextAt{21}{352.136}{204.914}{t}%
\TextAt{9}{355.692}{206.707}{\Character{41}}%
\TextAt{20}{360.244}{253.378}{\Character{49}}%
\TextAt{20}{360.244}{232.337}{\Character{67}}%
\TextAt{20}{360.244}{224.685}{\Character{65}}%
\TextAt{13}{372.698}{225.761}{\Character{58}}%
\TextAt{8}{79.226}{186.205}{This}%
\TextAt{8}{105.946}{186.205}{absolute}%
\TextAt{8}{152.189}{186.205}{value}%
\TextAt{13}{182.435}{186.205}{R}%
\TextAt{8}{195.341}{186.205}{is}%
\TextAt{8}{207.105}{186.205}{called}%
\TextAt{8}{240.269}{186.205}{the}%
\TextAt{12}{260.425}{186.205}{regulator}%
\TextAt{8}{317.905}{186.205}{of}%
\TextAt{8}{331.235}{186.205}{the}%
\TextAt{8}{351.392}{186.205}{field}%
\TextAt{13}{376.749}{186.205}{K}%
\TextAt{8}{387.566}{186.205}{.}%
\TextAt{8}{97.159}{160.526}{The}%
\TextAt{8}{121.213}{160.526}{importance}%
\TextAt{8}{182.674}{160.526}{of}%
\TextAt{8}{196.004}{160.526}{the}%
\TextAt{8}{216.161}{160.526}{regulator}%
\TextAt{8}{266.253}{160.526}{will}%
\TextAt{8}{288.358}{160.526}{only}%
\TextAt{8}{314.05}{160.526}{be}%
\TextAt{8}{329.974}{160.526}{demonstrated}%
\TextAt{8}{403.534}{160.526}{later}%
\TextAt{8}{430.852}{160.526}{(see}%
\TextAt{8}{454.32}{160.526}{chap.}%
\TextAt{8}{486.511}{160.526}{VII,}%
\TextAt{8}{511.198}{160.526}{§5).}%
\TextAt{12}{79.226}{134.847}{Exercise}%
\TextAt{12}{133.599}{134.847}{I.7.1.}%
\TextAt{8}{168.211}{134.847}{Let}%
\TextAt{13}{190.2}{134.847}{D}%
\TextAt{13}{205.299}{134.847}{\Character{62}}%
\TextAt{9}{219.463}{134.847}{1}%
\TextAt{8}{230.232}{134.847}{be}%
\TextAt{8}{247.175}{134.847}{a}%
\TextAt{8}{257.949}{134.847}{squarefree}%
\TextAt{8}{314.309}{134.847}{integer}%
\TextAt{8}{354.026}{134.847}{and}%
\TextAt{13}{377.808}{134.847}{d}%
\TextAt{8}{388.807}{134.847}{the}%
\TextAt{8}{409.982}{134.847}{discriminant}%
\TextAt{8}{478.38}{134.847}{of}%
\TextAt{8}{492.729}{134.847}{the}%
\TextAt{8}{513.904}{134.847}{real}%
\TextAt{8}{79.226}{117.801}{quadratic}%
\TextAt{8}{133.201}{117.801}{number}%
\TextAt{8}{177.408}{117.801}{field}%
\TextAt{13}{204.376}{117.801}{K}%
\TextAt{9}{221.259}{117.801}{\Character{61}}%
\TextAt{18}{236.43}{117.801}{Q}%
\TextAt{9}{245.728}{117.801}{\Character{40}}%
\TextAt{16}{250.281}{127.924}{\Character{112}}%
\Vector{q 0 G 0 g 0.478 w 188.243 -645.837 m 198.288 -645.837 l S Q}%
\TextAt{13}{260.243}{117.801}{D}%
\TextAt{9}{270.288}{117.801}{\Character{41}}%
\TextAt{8}{280.349}{117.801}{(see}%
\TextAt{8}{305.427}{117.801}{§2,}%
\TextAt{8}{326.127}{117.801}{exercise}%
\TextAt{8}{371.039}{117.801}{4).}%
\TextAt{8}{394.733}{117.801}{Let}%
\TextAt{13}{417.314}{117.801}{x}%
\TextAt{17}{423.966}{116.008}{1}%
\TextAt{13}{428.698}{117.801}{\Character{59}}%
\TextAt{13}{433.922}{117.801}{y}%
\TextAt{17}{439.65}{116.008}{1}%
\TextAt{8}{449.89}{117.801}{be}%
\TextAt{8}{467.426}{117.801}{the}%
\TextAt{8}{489.193}{117.801}{uniquely}%
\TextAt{8}{79.226}{100.756}{determined}%
\TextAt{8}{140.353}{100.756}{rational}%
\TextAt{8}{183.941}{100.756}{integer}%
\TextAt{8}{222.64}{100.756}{solution}%
\TextAt{8}{266.934}{100.756}{of}%
\TextAt{8}{280.264}{100.756}{the}%
\TextAt{8}{300.421}{100.756}{equation}%
\TextAt{13}{267.449}{72.921}{x}%
\TextAt{17}{274.101}{77.858}{2}%
\TextAt{16}{281.49}{72.921}{\Character{0}}%
\TextAt{13}{293.445}{72.921}{dy}%
\TextAt{17}{305.665}{77.858}{2}%
\TextAt{9}{313.718}{72.921}{\Character{61}}%
\TextAt{16}{326.143}{72.921}{\Character{0}}%
\TextAt{9}{335.442}{72.921}{4}%
\TextAt{13}{341.295}{72.921}{\Character{59}}%
\TextAt{8}{300.14}{43.033}{40}%
\Destination{78.23}{809.815}{nav.page.40}%
\Destination{79.226}{148.855}{nav.ex.1.7.1}%
\Destination{79.36}{606.555}{nav.print.43}%
\Destination{79.36}{347.555}{nav.number.1.7.5}%
\Hotlink{276.051}{623.455}{6.858}{11.047}{nav.section.1.5}%
\Hotlink{300.188}{623.455}{12.716}{11.047}{nav.print.33}%
\Hotlink{486.011}{158.355}{18.526}{11.047}{nav.chapter.7}%
\Hotlink{516.377}{158.355}{6.858}{11.047}{nav.section.7.5}%
\Hotlink{370.539}{115.555}{6.858}{11.047}{nav.ex.1.2.4}%
\Hotlink{310.606}{115.555}{6.858}{11.047}{nav.section.1.2}%
\Hotlink{337.181}{703.255}{25.078}{11.047}{nav.number.1.7.3}%
}

% Page 53 (body page 41)
\Page{%
\TextAt{2}{79.226}{796.943}{I.}%
\TextAt{2}{89.83}{796.943}{Algebraic}%
\TextAt{2}{138.648}{796.943}{Integers}%
\Vector{q 0 G 0 g 0.299 w 7.226 18.729 m 460.769 18.729 l S Q}%
\TextAt{8}{79.226}{761.124}{or—in}%
\TextAt{8}{114.996}{761.124}{case}%
\TextAt{8}{139.767}{761.124}{this}%
\TextAt{8}{162.59}{761.124}{equation}%
\TextAt{8}{210.398}{761.124}{has}%
\TextAt{8}{231.272}{761.124}{no}%
\TextAt{8}{247.531}{761.124}{rational}%
\TextAt{8}{291.12}{761.124}{integer}%
\TextAt{8}{329.818}{761.124}{solutions—of}%
\TextAt{8}{399.864}{761.124}{the}%
\TextAt{8}{420.02}{761.124}{equation}%
\TextAt{13}{272.098}{735.994}{x}%
\TextAt{17}{278.75}{740.93}{2}%
\TextAt{16}{286.139}{735.994}{\Character{0}}%
\TextAt{13}{298.095}{735.994}{dy}%
\TextAt{17}{310.314}{740.93}{2}%
\TextAt{9}{318.367}{735.994}{\Character{61}}%
\TextAt{9}{330.776}{735.994}{4}%
\TextAt{13}{336.646}{735.994}{\Character{59}}%
\TextAt{8}{79.226}{710.864}{for}%
\TextAt{8}{97.111}{710.864}{which}%
\TextAt{13}{130.597}{710.864}{x}%
\TextAt{17}{137.25}{709.071}{1}%
\TextAt{13}{141.982}{710.864}{\Character{59}}%
\TextAt{13}{147.206}{710.864}{y}%
\TextAt{17}{152.934}{709.071}{1}%
\TextAt{13}{160.987}{710.864}{\Character{62}}%
\TextAt{9}{173.412}{710.864}{0}%
\TextAt{8}{183.163}{710.864}{are}%
\TextAt{8}{202.674}{710.864}{as}%
\TextAt{8}{217.044}{710.864}{small}%
\TextAt{8}{247.673}{710.864}{as}%
\TextAt{8}{262.043}{710.864}{possible.}%
\TextAt{8}{310.613}{710.864}{Then}%
\TextAt{13}{265.68}{680.227}{\Character{34}}%
\TextAt{17}{271.158}{678.434}{1}%
\TextAt{9}{279.211}{680.227}{\Character{61}}%
\TextAt{13}{292.832}{688.315}{x}%
\TextAt{17}{299.484}{686.522}{1}%
\TextAt{9}{306.873}{688.315}{\Character{43}}%
\TextAt{13}{318.635}{688.315}{y}%
\TextAt{17}{324.342}{686.522}{1}%
\TextAt{16}{329.075}{698.503}{\Character{112}}%
\Vector{q 0 G 0 g 0.478 w 267.037 -75.257 m 273.12 -75.257 l S Q}%
\TextAt{13}{339.037}{688.315}{d}%
\Vector{q 0 G 0 g 0.478 w 220.832 -90.784 m 273.12 -90.784 l S Q}%
\TextAt{9}{316.05}{672.027}{2}%
\TextAt{8}{79.226}{653.472}{is}%
\TextAt{8}{92.734}{653.472}{a}%
\TextAt{8}{104.234}{653.472}{fundamental}%
\TextAt{8}{173.62}{653.472}{unit}%
\TextAt{8}{200.076}{653.472}{of}%
\TextAt{13}{215.15}{653.472}{K}%
\TextAt{8}{225.968}{653.472}{.}%
\TextAt{8}{239.649}{653.472}{(The}%
\TextAt{8}{270.002}{653.472}{pair}%
\TextAt{8}{295.812}{653.472}{of}%
\TextAt{8}{310.886}{653.472}{equations}%
\TextAt{13}{365.054}{653.472}{x}%
\TextAt{17}{371.706}{657.81}{2}%
\TextAt{16}{380.284}{653.472}{\Character{0}}%
\TextAt{13}{393.428}{653.472}{dy}%
\TextAt{17}{405.647}{657.81}{2}%
\TextAt{9}{416.673}{653.472}{\Character{61}}%
\NativeMathAt{432.071}{653.472}{11.955}{\pm}%
\TextAt{9}{441.369}{653.472}{4}%
\TextAt{8}{452.864}{653.472}{is}%
\TextAt{8}{466.372}{653.472}{called}%
\TextAt{12}{501.28}{653.472}{Pell’s}%
\TextAt{12}{79.226}{636.426}{equation}%
\TextAt{8}{129.653}{636.426}{.)}%
\TextAt{12}{79.226}{606.648}{Exercise}%
\TextAt{12}{132.427}{606.648}{I.7.2.}%
\TextAt{8}{166.021}{606.648}{Check}%
\TextAt{8}{201.133}{606.648}{the}%
\TextAt{8}{221.289}{606.648}{following}%
\TextAt{8}{270.724}{606.648}{table}%
\TextAt{8}{299.99}{606.648}{of}%
\TextAt{8}{313.32}{606.648}{fundamental}%
\TextAt{8}{380.963}{606.648}{units}%
\TextAt{13}{410.289}{606.648}{\Character{34}}%
\TextAt{17}{415.767}{604.855}{1}%
\TextAt{8}{424.396}{606.648}{for}%
\TextAt{18}{442.281}{606.648}{Q}%
\TextAt{9}{451.58}{606.648}{\Character{40}}%
\TextAt{16}{456.132}{616.77}{\Character{112}}%
\Vector{q 0 G 0 g 0.478 w 394.095 -156.991 m 404.139 -156.991 l S Q}%
\TextAt{13}{466.095}{606.648}{D}%
\TextAt{9}{476.139}{606.648}{\Character{41}}%
\TextAt{8}{480.692}{606.648}{:}%
\TextAt{13}{143.444}{582.647}{D}%
\Vector{q 0 G 0 g 0.399 w 86.752 -196.466 m 86.752 -179.42 l S Q}%
\TextAt{9}{179.049}{582.647}{2}%
\TextAt{9}{225.098}{582.647}{3}%
\TextAt{9}{281.553}{582.647}{5}%
\TextAt{9}{340.934}{582.647}{6}%
\TextAt{9}{392.836}{582.647}{7}%
\TextAt{9}{441.812}{582.647}{10}%
\Vector{q 0 G 0 g 0.399 w 66.379 -196.666 m 401.616 -196.666 l S Q}%
\TextAt{13}{143.361}{565.203}{\Character{34}}%
\TextAt{17}{148.839}{563.41}{1}%
\Vector{q 0 G 0 g 0.399 w 86.752 -213.911 m 86.752 -196.865 l S Q}%
\TextAt{9}{163.932}{565.203}{1}%
\TextAt{9}{172.432}{565.203}{\Character{43}}%
\TextAt{16}{184.203}{575.093}{\Character{112}}%
\Vector{q 0 G 0 g 0.478 w 122.166 -198.668 m 128.019 -198.668 l S Q}%
\TextAt{9}{194.166}{565.203}{2}%
\TextAt{9}{209.971}{565.203}{2}%
\TextAt{9}{218.471}{565.203}{\Character{43}}%
\TextAt{16}{230.252}{575.093}{\Character{112}}%
\Vector{q 0 G 0 g 0.478 w 168.215 -198.668 m 174.068 -198.668 l S Q}%
\TextAt{9}{240.215}{565.203}{3}%
\TextAt{9}{256.02}{565.203}{\Character{40}1}%
\TextAt{9}{269.063}{565.203}{\Character{43}}%
\TextAt{16}{280.854}{575.093}{\Character{112}}%
\Vector{q 0 G 0 g 0.478 w 218.817 -198.668 m 224.67 -198.668 l S Q}%
\TextAt{9}{290.817}{565.203}{5\Character{41}\Character{47}2}%
\TextAt{9}{322.857}{565.203}{5}%
\TextAt{9}{331.357}{565.203}{\Character{43}}%
\TextAt{9}{343.109}{565.203}{2}%
\TextAt{16}{349.015}{575.093}{\Character{112}}%
\Vector{q 0 G 0 g 0.478 w 286.977 -198.668 m 292.83 -198.668 l S Q}%
\TextAt{9}{358.977}{565.203}{6}%
\TextAt{9}{374.782}{565.203}{8}%
\TextAt{9}{383.282}{565.203}{\Character{43}}%
\TextAt{9}{395.034}{565.203}{3}%
\TextAt{16}{400.917}{575.093}{\Character{112}}%
\Vector{q 0 G 0 g 0.478 w 338.88 -198.668 m 344.733 -198.668 l S Q}%
\TextAt{9}{410.88}{565.203}{7}%
\TextAt{9}{426.685}{565.203}{3}%
\TextAt{9}{435.185}{565.203}{\Character{43}}%
\TextAt{16}{446.966}{575.093}{\Character{112}}%
\Vector{q 0 G 0 g 0.478 w 384.929 -198.668 m 396.635 -198.668 l S Q}%
\TextAt{9}{456.929}{565.203}{10}%
\TextAt{12}{97.159}{538.546}{Hint:}%
\TextAt{8}{133.284}{538.546}{Check}%
\TextAt{8}{168.69}{538.546}{one}%
\TextAt{8}{190.443}{538.546}{by}%
\TextAt{8}{206.995}{538.546}{one}%
\TextAt{8}{228.748}{538.546}{for}%
\TextAt{13}{246.926}{538.546}{y}%
\TextAt{9}{256.883}{538.546}{\Character{61}}%
\TextAt{9}{269.795}{538.546}{1}%
\TextAt{13}{275.662}{538.546}{\Character{59}}%
\TextAt{9}{280.906}{538.546}{2}%
\TextAt{13}{286.759}{538.546}{\Character{59}}%
\TextAt{9}{292.003}{538.546}{3}%
\TextAt{13}{297.856}{538.546}{\Character{59}}%
\TextAt{13}{303.08}{538.546}{\Character{58}}%
\TextAt{13}{308.317}{538.546}{\Character{58}}%
\TextAt{13}{313.541}{538.546}{\Character{58}}%
\TextAt{8}{316.84}{538.546}{,}%
\TextAt{8}{324.356}{538.546}{whether}%
\TextAt{8}{369.517}{538.546}{one}%
\TextAt{8}{391.27}{538.546}{of}%
\TextAt{8}{404.893}{538.546}{the}%
\TextAt{8}{425.343}{538.546}{numbers}%
\TextAt{13}{472.847}{538.546}{dy}%
\TextAt{17}{485.067}{542.884}{2}%
\NativeMathAt{492.656}{538.546}{11.955}{\mp}%
\TextAt{9}{504.811}{538.546}{4}%
\TextAt{8}{514.854}{538.546}{is}%
\TextAt{8}{526.911}{538.546}{a}%
\TextAt{8}{79.226}{521.5}{square}%
\TextAt{13}{116.651}{521.5}{x}%
\TextAt{17}{123.303}{525.838}{2}%
\TextAt{8}{128.036}{521.5}{.}%
\TextAt{8}{138.329}{521.5}{By}%
\TextAt{8}{157.307}{521.5}{the}%
\TextAt{8}{178.079}{521.5}{unit}%
\TextAt{8}{203.405}{521.5}{theorem}%
\TextAt{8}{249.546}{521.5}{this}%
\TextAt{8}{272.983}{521.5}{is}%
\TextAt{8}{285.362}{521.5}{bound}%
\TextAt{8}{322.07}{521.5}{to}%
\TextAt{8}{336.995}{521.5}{happen,}%
\TextAt{8}{382.309}{521.5}{with}%
\TextAt{8}{409.584}{521.5}{the}%
\TextAt{8}{430.356}{521.5}{plus}%
\TextAt{8}{455.742}{521.5}{sign.}%
\TextAt{8}{486.264}{521.5}{However,}%
\TextAt{8}{79.226}{504.454}{for}%
\TextAt{8}{97.546}{504.454}{fixed}%
\TextAt{13}{126.268}{504.454}{y}%
\TextAt{8}{132.404}{504.454}{,}%
\TextAt{8}{140.098}{504.454}{let}%
\TextAt{8}{157.438}{504.454}{preference}%
\TextAt{8}{213.464}{504.454}{be}%
\TextAt{8}{229.824}{504.454}{given}%
\TextAt{8}{260.829}{504.454}{to}%
\TextAt{8}{275.574}{504.454}{the}%
\TextAt{8}{296.166}{504.454}{minus}%
\TextAt{8}{330.805}{504.454}{sign.}%
\TextAt{8}{360.787}{504.454}{Then}%
\TextAt{8}{391.78}{504.454}{the}%
\TextAt{8}{412.372}{504.454}{first}%
\TextAt{8}{436.933}{504.454}{case,}%
\TextAt{8}{465.5}{504.454}{in}%
\TextAt{8}{479.588}{504.454}{this}%
\TextAt{8}{502.846}{504.454}{order,}%
\TextAt{8}{79.226}{487.408}{where}%
\TextAt{13}{113.035}{487.408}{dy}%
\TextAt{17}{125.255}{491.746}{2}%
\TextAt{17}{124.826}{484.452}{1}%
\NativeMathAt{132.644}{487.408}{11.955}{\mp}%
\TextAt{9}{144.599}{487.408}{4}%
\TextAt{9}{153.757}{487.408}{\Character{61}}%
\TextAt{13}{166.198}{487.408}{x}%
\TextAt{17}{172.85}{491.746}{2}%
\TextAt{17}{172.85}{484.452}{1}%
\TextAt{8}{177.583}{487.408}{,}%
\TextAt{8}{184.732}{487.408}{gives}%
\TextAt{8}{213.412}{487.408}{the}%
\TextAt{8}{233.569}{487.408}{fundamental}%
\TextAt{8}{301.211}{487.408}{unit}%
\TextAt{13}{325.922}{487.408}{\Character{34}}%
\TextAt{17}{331.4}{485.614}{1}%
\TextAt{9}{339.453}{487.408}{\Character{61}}%
\TextAt{9}{351.863}{487.408}{\Character{40}}%
\TextAt{13}{356.431}{487.408}{x}%
\TextAt{17}{363.083}{485.614}{1}%
\TextAt{9}{370.472}{487.408}{\Character{43}}%
\TextAt{13}{382.234}{487.408}{y}%
\TextAt{17}{387.941}{485.614}{1}%
\TextAt{16}{392.674}{497.596}{\Character{112}}%
\Vector{q 0 G 0 g 0.478 w 330.636 -276.165 m 336.719 -276.165 l S Q}%
\TextAt{13}{402.636}{487.408}{d}%
\TextAt{9}{408.719}{487.408}{\Character{41}\Character{47}2}%
\TextAt{8}{424.977}{487.408}{.}%
\TextAt{12}{79.226}{457.63}{Exercise}%
\TextAt{12}{132.427}{457.63}{I.7.3.}%
\TextAt{12}{166.021}{457.63}{The}%
\TextAt{12}{193.458}{457.63}{Battle}%
\TextAt{12}{234.392}{457.63}{of}%
\TextAt{12}{249.719}{457.63}{Hastings}%
\TextAt{12}{305.059}{457.63}{(October}%
\TextAt{12}{362.301}{457.63}{14,}%
\TextAt{12}{383.987}{457.63}{1066).}%
\TextAt{8}{97.159}{437.595}{“The}%
\TextAt{8}{128.544}{437.595}{men}%
\TextAt{8}{155.71}{437.595}{of}%
\TextAt{8}{170.848}{437.595}{Harold}%
\TextAt{8}{211.355}{437.595}{stood}%
\TextAt{8}{245.096}{437.595}{well}%
\TextAt{8}{270.635}{437.595}{together,}%
\TextAt{8}{322.33}{437.595}{as}%
\TextAt{8}{338.509}{437.595}{their}%
\TextAt{8}{368.28}{437.595}{wont}%
\TextAt{8}{398.709}{437.595}{was,}%
\TextAt{8}{426.721}{437.595}{and}%
\TextAt{8}{451.292}{437.595}{formed}%
\TextAt{8}{492.445}{437.595}{thirteen}%
\TextAt{8}{79.226}{420.549}{squares,}%
\TextAt{8}{123.406}{420.549}{with}%
\TextAt{8}{149.445}{420.549}{a}%
\TextAt{8}{158.58}{420.549}{like}%
\TextAt{8}{179.419}{420.549}{number}%
\TextAt{8}{221.394}{420.549}{of}%
\TextAt{8}{234.104}{420.549}{men}%
\TextAt{8}{258.84}{420.549}{in}%
\TextAt{8}{271.872}{420.549}{every}%
\TextAt{8}{302.143}{420.549}{square}%
\TextAt{8}{338.332}{420.549}{thereof,}%
\TextAt{8}{380.432}{420.549}{and}%
\TextAt{8}{402.574}{420.549}{woe}%
\TextAt{8}{425.361}{420.549}{to}%
\TextAt{8}{439.051}{420.549}{the}%
\TextAt{8}{458.587}{420.549}{hardy}%
\TextAt{8}{491.464}{420.549}{Norman}%
\TextAt{8}{79.226}{403.503}{who}%
\TextAt{8}{104.954}{403.503}{ventured}%
\TextAt{8}{154.425}{403.503}{to}%
\TextAt{8}{169.752}{403.503}{enter}%
\TextAt{8}{200.358}{403.503}{their}%
\TextAt{8}{229.337}{403.503}{redoubts;}%
\TextAt{8}{282.305}{403.503}{for}%
\TextAt{8}{301.207}{403.503}{a}%
\TextAt{8}{311.979}{403.503}{single}%
\TextAt{8}{345.573}{403.503}{blow}%
\TextAt{8}{374.23}{403.503}{of}%
\TextAt{8}{388.577}{403.503}{a}%
\TextAt{8}{399.349}{403.503}{Saxon}%
\TextAt{8}{435.167}{403.503}{war-hatched}%
\TextAt{8}{502.523}{403.503}{would}%
\TextAt{8}{79.226}{386.457}{break}%
\TextAt{8}{110.518}{386.457}{his}%
\TextAt{8}{127.882}{386.457}{lance}%
\TextAt{8}{156.891}{386.457}{and}%
\TextAt{8}{178.75}{386.457}{cut}%
\TextAt{8}{198.003}{386.457}{through}%
\TextAt{8}{241.334}{386.457}{his}%
\TextAt{8}{258.698}{386.457}{coat}%
\TextAt{8}{283.163}{386.457}{of}%
\TextAt{8}{295.59}{386.457}{mail…}%
\TextAt{8}{328.4}{386.457}{When}%
\TextAt{8}{361.629}{386.457}{Harold}%
\TextAt{8}{399.424}{386.457}{threw}%
\TextAt{8}{431.685}{386.457}{himself}%
\TextAt{8}{470.831}{386.457}{into}%
\TextAt{8}{493.671}{386.457}{the}%
\TextAt{8}{512.924}{386.457}{fray}%
\TextAt{8}{79.226}{369.411}{the}%
\TextAt{8}{99.934}{369.411}{Saxons}%
\TextAt{8}{139.901}{369.411}{were}%
\TextAt{8}{167.436}{369.411}{one}%
\TextAt{8}{189.447}{369.411}{mighty}%
\TextAt{8}{229.354}{369.411}{square}%
\TextAt{8}{266.716}{369.411}{of}%
\TextAt{8}{280.597}{369.411}{men,}%
\TextAt{8}{309.895}{369.411}{shouting}%
\TextAt{8}{357.992}{369.411}{the}%
\TextAt{8}{378.7}{369.411}{battle-cries,}%
\TextAt{8}{443.182}{369.411}{‘Ut!’,}%
\TextAt{8}{474.105}{369.411}{‘Olicrosse!’,}%
\TextAt{8}{79.226}{352.365}{‘Godemite!’}%
\TextAt{8}{137.197}{352.365}{.}%
\TextAt{8}{138.823}{352.365}{”}%
\TextAt{8}{152.329}{352.365}{[Fictitious}%
\TextAt{8}{209.093}{352.365}{historical}%
\TextAt{8}{260.823}{352.365}{text,}%
\TextAt{8}{289.63}{352.365}{following}%
\TextAt{8}{339.997}{352.365}{essentially}%
\TextAt{8}{396.988}{352.365}{problem}%
\TextAt{8}{443.446}{352.365}{no.}%
\TextAt{8}{467.055}{352.365}{129}%
\TextAt{8}{489.459}{352.365}{in:}%
\TextAt{8}{509.529}{352.365}{H.E.}%
\TextAt{8}{79.226}{335.319}{Dundeney,}%
\TextAt{19}{136.934}{335.319}{Amusements}%
\TextAt{19}{205.401}{335.319}{in}%
\TextAt{19}{219.747}{335.319}{Mathematics}%
\TextAt{8}{284.628}{335.319}{,}%
\TextAt{8}{291.777}{335.319}{1917}%
\TextAt{8}{319.107}{335.319}{(Dover}%
\TextAt{8}{357.65}{335.319}{reprints}%
\TextAt{8}{400.964}{335.319}{1958}%
\TextAt{8}{428.293}{335.319}{and}%
\TextAt{8}{451.056}{335.319}{1970).]}%
\TextAt{12}{97.159}{315.285}{Question.}%
\TextAt{8}{157.898}{315.285}{How}%
\TextAt{8}{183.37}{315.285}{many}%
\TextAt{8}{214.055}{315.285}{troops}%
\TextAt{8}{249.032}{315.285}{does}%
\TextAt{8}{274.241}{315.285}{this}%
\TextAt{8}{295.876}{315.285}{suggest}%
\TextAt{8}{335.79}{315.285}{Harold}%
\TextAt{8}{373.301}{315.285}{II}%
\TextAt{8}{384.773}{315.285}{had}%
\TextAt{8}{406.348}{315.285}{at}%
\TextAt{8}{419.471}{315.285}{the}%
\TextAt{8}{438.439}{315.285}{battle}%
\TextAt{8}{471.073}{315.285}{of}%
\TextAt{8}{483.215}{315.285}{Hastings?}%
\TextAt{12}{79.226}{285.507}{Exercise}%
\TextAt{12}{132.91}{285.507}{I.7.4.}%
\TextAt{8}{166.925}{285.507}{Let}%
\TextAt{13}{188.315}{285.507}{\Character{16}}%
\TextAt{8}{198.649}{285.507}{be}%
\TextAt{8}{214.993}{285.507}{a}%
\TextAt{8}{225.169}{285.507}{primitive}%
\TextAt{13}{275.67}{285.507}{p}%
\TextAt{8}{281.545}{285.507}{-th}%
\TextAt{8}{300.819}{285.507}{root}%
\TextAt{8}{326.286}{285.507}{of}%
\TextAt{8}{340.036}{285.507}{unity,}%
\TextAt{13}{373.403}{285.507}{p}%
\TextAt{8}{383.596}{285.507}{an}%
\TextAt{8}{400.276}{285.507}{odd}%
\TextAt{8}{423.782}{285.507}{prime}%
\TextAt{8}{457.366}{285.507}{number.}%
\TextAt{8}{505.775}{285.507}{Show}%
\TextAt{8}{79.226}{268.461}{that}%
\TextAt{18}{105.336}{268.461}{Z}%
\TextAt{9}{113.306}{268.461}{\Character{91}}%
\TextAt{13}{116.558}{268.461}{\Character{16}}%
\TextAt{9}{122.574}{268.461}{\Character{93}}%
\TextAt{3}{125.825}{272.8}{\Character{3}}%
\TextAt{9}{135.142}{268.461}{\Character{61}}%
\TextAt{16}{148.83}{268.461}{\Character{104}}%
\TextAt{13}{153.479}{268.461}{\Character{16}}%
\TextAt{16}{159.495}{268.461}{\Character{105}}%
\TextAt{18}{164.144}{268.461}{Z}%
\TextAt{9}{172.115}{268.461}{\Character{91}}%
\TextAt{13}{175.366}{268.461}{\Character{16}}%
\TextAt{9}{184.544}{268.461}{\Character{43}}%
\TextAt{13}{196.81}{268.461}{\Character{16}}%
\TextAt{3}{202.826}{272.8}{\Character{0}}%
\TextAt{17}{209.412}{272.8}{1}%
\TextAt{9}{214.145}{268.461}{\Character{93}}%
\TextAt{3}{217.396}{272.8}{\Character{3}}%
\TextAt{8}{222.129}{268.461}{.}%
\TextAt{8}{232.801}{268.461}{Show}%
\TextAt{8}{264.434}{268.461}{that}%
\TextAt{18}{290.544}{268.461}{Z}%
\TextAt{9}{298.514}{268.461}{\Character{91}}%
\TextAt{13}{301.766}{268.461}{\Character{16}}%
\TextAt{9}{307.782}{268.461}{\Character{93}}%
\TextAt{3}{311.033}{272.8}{\Character{3}}%
\TextAt{9}{320.35}{268.461}{\Character{61}}%
\TextAt{16}{334.038}{268.461}{\Character{102}}%
\NativeMathAt{340.016}{268.461}{11.955}{\pm}%
\TextAt{13}{349.314}{268.461}{\Character{16}}%
\TextAt{21}{355.33}{272.8}{k}%
\TextAt{9}{360.45}{268.461}{\Character{40}1}%
\TextAt{9}{373.995}{268.461}{\Character{43}}%
\TextAt{13}{386.283}{268.461}{\Character{16}}%
\TextAt{9}{392.299}{268.461}{\Character{41}}%
\TextAt{21}{396.851}{272.8}{n}%
\TextAt{16}{407.072}{268.461}{\Character{106}}%
\TextAt{9}{414.977}{268.461}{0}%
\NativeMathAt{425.414}{268.461}{11.955}{\leq}%
\TextAt{13}{439.296}{268.461}{k}%
\TextAt{13}{450.367}{268.461}{\Character{60}}%
\TextAt{9}{464.058}{268.461}{5}%
\TextAt{13}{469.911}{268.461}{\Character{59}}%
\TextAt{13}{479.774}{268.461}{n}%
\NativeMathAt{491.366}{268.461}{11.955}{\in}%
\TextAt{18}{503.92}{268.461}{Z}%
\TextAt{16}{511.89}{268.461}{\Character{103}}%
\TextAt{8}{517.867}{268.461}{,}%
\TextAt{8}{525.943}{268.461}{if}%
\TextAt{13}{79.226}{251.415}{p}%
\TextAt{9}{88.422}{251.415}{\Character{61}}%
\TextAt{9}{100.832}{251.415}{5}%
\TextAt{8}{106.701}{251.415}{.}%
\TextAt{12}{79.226}{221.637}{Exercise}%
\TextAt{12}{132.393}{221.637}{I.7.5.}%
\TextAt{8}{165.958}{221.637}{Let}%
\TextAt{13}{186.899}{221.637}{\Character{16}}%
\TextAt{8}{196.783}{221.637}{be}%
\TextAt{8}{212.678}{221.637}{a}%
\TextAt{8}{222.404}{221.637}{primitive}%
\TextAt{13}{272.455}{221.637}{m}%
\TextAt{8}{282.695}{221.637}{-th}%
\TextAt{8}{301.519}{221.637}{root}%
\TextAt{8}{326.536}{221.637}{of}%
\TextAt{8}{339.837}{221.637}{unity,}%
\TextAt{13}{372.654}{221.637}{m}%
\NativeMathAt{386.215}{221.637}{11.955}{\geq}%
\TextAt{9}{398.834}{221.637}{3}%
\TextAt{8}{404.687}{221.637}{.}%
\TextAt{8}{413.126}{221.637}{Show}%
\TextAt{8}{443.989}{221.637}{that}%
\TextAt{8}{469.329}{221.637}{the}%
\TextAt{8}{489.456}{221.637}{numbers}%
\TextAt{9}{80.422}{208.329}{1}%
\TextAt{16}{88.931}{208.329}{\Character{0}}%
\TextAt{13}{100.886}{208.329}{\Character{16}}%
\TextAt{21}{106.902}{212.668}{k}%
\Vector{q 0 G 0 g 0.478 w 8.422 -570.77 m 40.022 -570.77 l S Q}%
\TextAt{9}{82.982}{192.041}{1}%
\TextAt{16}{91.491}{192.041}{\Character{0}}%
\TextAt{13}{103.446}{192.041}{\Character{16}}%
\TextAt{8}{117.201}{200.241}{for}%
\TextAt{9}{135.172}{200.241}{\Character{40}}%
\TextAt{13}{139.724}{200.241}{k}%
\TextAt{13}{146.204}{200.241}{\Character{59}}%
\TextAt{13}{151.44}{200.241}{m}%
\TextAt{9}{161.697}{200.241}{\Character{41}}%
\TextAt{9}{169.707}{200.241}{\Character{61}}%
\TextAt{9}{182.272}{200.241}{1}%
\TextAt{8}{192.126}{200.241}{are}%
\TextAt{8}{211.723}{200.241}{units}%
\TextAt{8}{241.136}{200.241}{in}%
\TextAt{8}{254.875}{200.241}{the}%
\TextAt{8}{275.117}{200.241}{ring}%
\TextAt{8}{299.269}{200.241}{of}%
\TextAt{8}{312.686}{200.241}{integers}%
\TextAt{8}{356.085}{200.241}{of}%
\TextAt{8}{369.502}{200.241}{the}%
\TextAt{8}{389.744}{200.241}{field}%
\TextAt{18}{415.187}{200.241}{Q}%
\TextAt{9}{424.486}{200.241}{\Character{40}}%
\TextAt{13}{429.038}{200.241}{\Character{16}}%
\TextAt{9}{435.054}{200.241}{\Character{41}}%
\TextAt{8}{439.606}{200.241}{.}%
\TextAt{8}{448.313}{200.241}{The}%
\TextAt{8}{472.453}{200.241}{subgroup}%
\TextAt{8}{523.337}{200.241}{of}%
\TextAt{8}{79.226}{180.303}{the}%
\TextAt{8}{99.383}{180.303}{group}%
\TextAt{8}{132.558}{180.303}{of}%
\TextAt{8}{145.888}{180.303}{units}%
\TextAt{8}{175.214}{180.303}{they}%
\TextAt{8}{201.551}{180.303}{generate}%
\TextAt{8}{248.38}{180.303}{is}%
\TextAt{8}{260.144}{180.303}{called}%
\TextAt{8}{293.307}{180.303}{the}%
\TextAt{12}{313.464}{180.303}{group}%
\TextAt{12}{351.852}{180.303}{of}%
\TextAt{12}{367.178}{180.303}{cyclotomic}%
\TextAt{12}{434.091}{180.303}{units}%
\TextAt{8}{463.322}{180.303}{.}%
\TextAt{12}{79.226}{150.526}{Exercise}%
\TextAt{12}{131.509}{150.526}{I.7.6.}%
\TextAt{8}{164.306}{150.526}{Let}%
\TextAt{13}{184.478}{150.526}{K}%
\TextAt{8}{198.395}{150.526}{be}%
\TextAt{8}{213.522}{150.526}{a}%
\TextAt{8}{222.48}{150.526}{totally}%
\TextAt{8}{259.09}{150.526}{real}%
\TextAt{8}{281.055}{150.526}{number}%
\TextAt{8}{322.854}{150.526}{field,}%
\TextAt{8}{350.825}{150.526}{i.e.,}%
\TextAt{13}{372.292}{150.526}{X}%
\TextAt{9}{386.268}{150.526}{\Character{61}}%
\TextAt{8}{398.693}{150.526}{Hom}%
\TextAt{9}{423.082}{150.526}{\Character{40}}%
\TextAt{13}{427.634}{150.526}{K}%
\TextAt{13}{437.796}{150.526}{\Character{59}}%
\TextAt{18}{443.046}{150.526}{C}%
\TextAt{9}{451.68}{150.526}{\Character{41}}%
\TextAt{9}{459.535}{150.526}{\Character{61}}%
\TextAt{8}{471.979}{150.526}{Hom}%
\TextAt{9}{496.367}{150.526}{\Character{40}}%
\TextAt{13}{500.92}{150.526}{K}%
\TextAt{13}{511.082}{150.526}{\Character{59}}%
\TextAt{18}{516.331}{150.526}{R}%
\TextAt{9}{524.965}{150.526}{\Character{41}}%
\TextAt{8}{529.518}{150.526}{,}%
\TextAt{8}{79.226}{133.48}{and}%
\TextAt{8}{101.989}{133.48}{let}%
\TextAt{13}{118.893}{133.48}{T}%
\TextAt{8}{131.278}{133.48}{be}%
\TextAt{8}{147.202}{133.48}{a}%
\TextAt{8}{156.957}{133.48}{proper}%
\TextAt{8}{194.353}{133.48}{nonempty}%
\TextAt{8}{248.988}{133.48}{subset}%
\TextAt{8}{284.878}{133.48}{of}%
\TextAt{13}{298.208}{133.48}{X}%
\TextAt{8}{308.863}{133.48}{.}%
\TextAt{8}{317.311}{133.48}{Then}%
\TextAt{8}{347.869}{133.48}{there}%
\TextAt{8}{377.78}{133.48}{exists}%
\TextAt{8}{410.095}{133.48}{a}%
\TextAt{8}{419.851}{133.48}{unit}%
\TextAt{13}{444.562}{133.48}{\Character{34}}%
\TextAt{8}{453.937}{133.48}{satisfying}%
\TextAt{9}{173.461}{108.35}{0}%
\TextAt{13}{182.635}{108.35}{\Character{60}}%
\TextAt{13}{195.045}{108.35}{\Character{28}}%
\TextAt{13}{201.464}{108.35}{\Character{34}}%
\TextAt{13}{210.264}{108.35}{\Character{60}}%
\TextAt{9}{222.704}{108.35}{1}%
\TextAt{8}{232.454}{108.35}{for}%
\TextAt{13}{250.339}{108.35}{\Character{28}}%
\NativeMathAt{260.079}{108.35}{11.955}{\in}%
\TextAt{13}{271.37}{108.35}{T}%
\TextAt{13}{279.189}{108.35}{\Character{59}}%
\TextAt{8}{308.361}{108.35}{and}%
\TextAt{13}{351.137}{108.35}{\Character{28}}%
\TextAt{13}{357.545}{108.35}{\Character{34}}%
\TextAt{13}{366.344}{108.35}{\Character{62}}%
\TextAt{9}{378.78}{108.35}{1}%
\TextAt{8}{388.53}{108.35}{for}%
\TextAt{13}{406.415}{108.35}{\Character{28}}%
\NativeMathAt{416.155}{108.35}{11.955}{\notin}%
\TextAt{13}{427.446}{108.35}{T}%
\TextAt{13}{435.265}{108.35}{\Character{58}}%
\TextAt{12}{97.159}{80.232}{Hint:}%
\TextAt{8}{132.426}{80.232}{Apply}%
\TextAt{8}{166.995}{80.232}{Minkowski’s}%
\TextAt{8}{232.588}{80.232}{lattice}%
\TextAt{8}{267.814}{80.232}{point}%
\TextAt{8}{297.84}{80.232}{theorem}%
\TextAt{8}{342.822}{80.232}{to}%
\TextAt{8}{356.589}{80.232}{the}%
\TextAt{8}{376.202}{80.232}{unit}%
\TextAt{8}{400.37}{80.232}{lattice}%
\TextAt{8}{435.597}{80.232}{in}%
\TextAt{8}{448.706}{80.232}{trace-zero}%
\TextAt{8}{502.14}{80.232}{space.}%
\TextAt{8}{300.14}{43.033}{41}%
\Destination{78.23}{809.815}{nav.page.41}%
\Destination{79.226}{620.655}{nav.ex.1.7.2}%
\Destination{79.226}{471.655}{nav.ex.1.7.3}%
\Destination{79.226}{299.555}{nav.ex.1.7.4}%
\Destination{79.226}{235.655}{nav.ex.1.7.5}%
\Destination{79.226}{164.555}{nav.ex.1.7.6}%
\Destination{79.36}{553.555}{nav.print.44}%
}

% Page 54 (body page 42)
\Page{%
\TextAt{2}{79.226}{796.943}{I.}%
\TextAt{2}{89.83}{796.943}{Algebraic}%
\TextAt{2}{138.648}{796.943}{Integers}%
\Vector{q 0 G 0 g 0.299 w 7.226 18.729 m 460.769 18.729 l S Q}%
\TextAt{15}{79.226}{760.804}{§}%
\TextAt{15}{93.785}{760.804}{8}%
\TextAt{15}{120.693}{760.804}{Extensions}%
\TextAt{15}{218.339}{760.804}{of}%
\TextAt{15}{240.409}{760.804}{Dedekind}%
\TextAt{15}{327.261}{760.804}{Domains}%
\TextAt{8}{79.226}{728.918}{Having}%
\TextAt{8}{119.431}{728.918}{studied}%
\TextAt{8}{160.665}{728.918}{the}%
\TextAt{8}{181.025}{728.918}{ideal}%
\TextAt{8}{209.191}{728.918}{class}%
\TextAt{8}{236.831}{728.918}{group}%
\TextAt{8}{270.21}{728.918}{and}%
\TextAt{8}{293.176}{728.918}{the}%
\TextAt{8}{313.536}{728.918}{group}%
\TextAt{8}{346.915}{728.918}{of}%
\TextAt{8}{360.448}{728.918}{units}%
\TextAt{8}{389.977}{728.918}{of}%
\TextAt{8}{403.51}{728.918}{the}%
\TextAt{8}{423.87}{728.918}{ring}%
\TextAt{14}{448.139}{728.918}{o}%
\TextAt{21}{453.983}{727.125}{K}%
\TextAt{8}{466.287}{728.918}{of}%
\TextAt{8}{479.82}{728.918}{integers}%
\TextAt{8}{523.337}{728.918}{of}%
\TextAt{8}{79.226}{711.872}{a}%
\TextAt{8}{88.764}{711.872}{number}%
\TextAt{8}{131.143}{711.872}{field}%
\TextAt{13}{156.283}{711.872}{K}%
\TextAt{8}{167.1}{711.872}{,}%
\TextAt{8}{174.076}{711.872}{we}%
\TextAt{8}{191.086}{711.872}{now}%
\TextAt{8}{215.257}{711.872}{propose}%
\TextAt{8}{258.353}{711.872}{to}%
\TextAt{8}{272.446}{711.872}{make}%
\TextAt{8}{302.798}{711.872}{a}%
\TextAt{8}{312.336}{711.872}{first}%
\TextAt{8}{336.244}{711.872}{survey}%
\TextAt{8}{372.837}{711.872}{of}%
\TextAt{8}{385.95}{711.872}{the}%
\TextAt{8}{405.889}{711.872}{set}%
\TextAt{8}{423.939}{711.872}{of}%
\TextAt{8}{437.051}{711.872}{prime}%
\TextAt{8}{469.998}{711.872}{ideals}%
\TextAt{8}{502.358}{711.872}{of}%
\TextAt{14}{515.471}{711.872}{o}%
\TextAt{21}{521.315}{710.079}{K}%
\TextAt{8}{529.518}{711.872}{.}%
\TextAt{8}{79.226}{694.826}{They}%
\TextAt{8}{109.511}{694.826}{are}%
\TextAt{8}{129.073}{694.826}{often}%
\TextAt{8}{158.712}{694.826}{referred}%
\TextAt{8}{202.005}{694.826}{to}%
\TextAt{8}{216.366}{694.826}{as}%
\TextAt{8}{230.786}{694.826}{the}%
\TextAt{12}{250.993}{694.826}{prime}%
\TextAt{12}{289.595}{694.826}{ideals}%
\TextAt{12}{327.073}{694.826}{of}%
\TextAt{13}{342.458}{694.826}{K}%
\TextAt{8}{353.275}{694.826}{—an}%
\TextAt{8}{381.289}{694.826}{imprecise}%
\TextAt{8}{432.771}{694.826}{manner}%
\TextAt{8}{475.095}{694.826}{of}%
\TextAt{8}{488.476}{694.826}{speaking}%
\TextAt{8}{79.226}{677.78}{which}%
\TextAt{8}{112.713}{677.78}{is,}%
\TextAt{8}{127.728}{677.78}{however,}%
\TextAt{8}{175.86}{677.78}{not}%
\TextAt{8}{196.674}{677.78}{likely}%
\TextAt{8}{227.566}{677.78}{to}%
\TextAt{8}{241.876}{677.78}{cause}%
\TextAt{8}{273.151}{677.78}{any}%
\TextAt{8}{295.268}{677.78}{misunderstanding.}%
\TextAt{8}{97.159}{657.746}{Every}%
\TextAt{8}{130.02}{657.746}{prime}%
\TextAt{8}{162.39}{657.746}{ideal}%
\TextAt{14}{189.56}{657.746}{p}%
\NativeMathAt{198.864}{657.746}{11.955}{\neq}%
\TextAt{9}{211.274}{657.746}{0}%
\TextAt{8}{220.247}{657.746}{of}%
\TextAt{14}{232.784}{657.746}{o}%
\TextAt{21}{238.628}{655.952}{K}%
\TextAt{8}{249.935}{657.746}{contains}%
\TextAt{8}{295.061}{657.746}{a}%
\TextAt{8}{304.024}{657.746}{rational}%
\TextAt{8}{346.819}{657.746}{prime}%
\TextAt{8}{379.19}{657.746}{number}%
\TextAt{13}{420.993}{657.746}{p}%
\TextAt{8}{429.972}{657.746}{(see}%
\TextAt{8}{452.647}{657.746}{§3,}%
\TextAt{8}{470.699}{657.746}{p.}%
\TextAt{8}{483.558}{657.746}{17)}%
\TextAt{8}{502.933}{657.746}{and}%
\TextAt{8}{524.903}{657.746}{is}%
\TextAt{8}{79.226}{640.7}{therefore}%
\TextAt{8}{127.641}{640.7}{a}%
\TextAt{8}{136.711}{640.7}{divisor}%
\TextAt{8}{174.138}{640.7}{of}%
\TextAt{8}{186.783}{640.7}{the}%
\TextAt{8}{206.254}{640.7}{ideal}%
\TextAt{13}{233.531}{640.7}{p}%
\TextAt{14}{239.406}{640.7}{o}%
\TextAt{21}{245.251}{638.906}{K}%
\TextAt{8}{253.453}{640.7}{.}%
\TextAt{8}{261.673}{640.7}{Hence}%
\TextAt{8}{295.765}{640.7}{the}%
\TextAt{8}{315.236}{640.7}{question}%
\TextAt{8}{361.116}{640.7}{arises}%
\TextAt{8}{392.423}{640.7}{as}%
\TextAt{8}{406.107}{640.7}{to}%
\TextAt{8}{419.732}{640.7}{how}%
\TextAt{8}{443.435}{640.7}{a}%
\TextAt{8}{452.505}{640.7}{prime}%
\TextAt{8}{484.984}{640.7}{number}%
\TextAt{13}{526.894}{640.7}{p}%
\TextAt{8}{79.226}{623.654}{factors}%
\TextAt{8}{117}{623.654}{into}%
\TextAt{8}{140.404}{623.654}{prime}%
\TextAt{8}{173.229}{623.654}{ideals}%
\TextAt{8}{205.468}{623.654}{of}%
\TextAt{8}{218.459}{623.654}{the}%
\TextAt{8}{238.277}{623.654}{ring}%
\TextAt{14}{262.004}{623.654}{o}%
\TextAt{21}{267.848}{621.86}{K}%
\TextAt{8}{276.051}{623.654}{.}%
\TextAt{8}{284.386}{623.654}{We}%
\TextAt{8}{304.192}{623.654}{treat}%
\TextAt{8}{332.473}{623.654}{this}%
\TextAt{8}{354.957}{623.654}{problem}%
\TextAt{8}{400.143}{623.654}{in}%
\TextAt{8}{413.457}{623.654}{a}%
\TextAt{8}{422.874}{623.654}{more}%
\TextAt{8}{451.801}{623.654}{general}%
\TextAt{8}{491.787}{623.654}{context,}%
\TextAt{8}{79.226}{606.608}{starting}%
\TextAt{8}{122.06}{606.608}{from}%
\TextAt{8}{148.886}{606.608}{an}%
\TextAt{8}{164.33}{606.608}{arbitrary}%
\TextAt{8}{213.285}{606.608}{Dedekind}%
\TextAt{8}{264.654}{606.608}{domain}%
\TextAt{14}{305.468}{606.608}{o}%
\TextAt{8}{314.394}{606.608}{at}%
\TextAt{8}{327.89}{606.608}{the}%
\TextAt{8}{347.232}{606.608}{base}%
\TextAt{8}{372.491}{606.608}{instead}%
\TextAt{8}{412.061}{606.608}{of}%
\TextAt{18}{424.577}{606.608}{Z}%
\TextAt{8}{432.547}{606.608}{,}%
\TextAt{8}{439.044}{606.608}{and}%
\TextAt{8}{460.992}{606.608}{taking}%
\TextAt{8}{496.282}{606.608}{instead}%
\TextAt{8}{79.226}{589.562}{of}%
\TextAt{14}{92.556}{589.562}{o}%
\TextAt{21}{98.4}{587.769}{K}%
\TextAt{8}{110.5}{589.562}{the}%
\TextAt{8}{130.657}{589.562}{integral}%
\TextAt{8}{173.265}{589.562}{closure}%
\TextAt{16}{212.347}{589.562}{\Character{79}}%
\TextAt{8}{226.094}{589.562}{of}%
\TextAt{14}{239.424}{589.562}{o}%
\TextAt{8}{249.165}{589.562}{in}%
\TextAt{8}{262.818}{589.562}{a}%
\TextAt{8}{272.574}{589.562}{finite}%
\TextAt{8}{302.485}{589.562}{extension}%
\TextAt{8}{354.251}{589.562}{of}%
\TextAt{8}{367.581}{589.562}{its}%
\TextAt{8}{383.9}{589.562}{field}%
\TextAt{8}{409.257}{589.562}{of}%
\TextAt{8}{422.587}{589.562}{fractions.}%
\TextAt{12}{79.226}{564.4}{Proposition}%
\TextAt{12}{152.094}{564.4}{(8.1).}%
\TextAt{8}{187.545}{564.4}{Let}%
\TextAt{14}{208.651}{564.4}{o}%
\TextAt{8}{218.528}{564.4}{be}%
\TextAt{8}{234.589}{564.4}{a}%
\TextAt{8}{244.48}{564.4}{Dedekind}%
\TextAt{8}{296.8}{564.4}{domain}%
\TextAt{8}{338.564}{564.4}{with}%
\TextAt{8}{365.36}{564.4}{field}%
\TextAt{8}{390.853}{564.4}{of}%
\TextAt{8}{404.319}{564.4}{fractions}%
\TextAt{13}{452.324}{564.4}{K}%
\TextAt{8}{463.141}{564.4}{,}%
\TextAt{8}{470.46}{564.4}{let}%
\TextAt{13}{487.501}{564.4}{L}%
\TextAt{16}{499.018}{564.4}{\Character{106}}%
\TextAt{13}{505.892}{564.4}{K}%
\TextAt{8}{520.743}{564.4}{be}%
\TextAt{8}{79.226}{547.354}{a}%
\TextAt{8}{88.831}{547.354}{finite}%
\TextAt{8}{118.592}{547.354}{extension}%
\TextAt{8}{170.207}{547.354}{of}%
\TextAt{13}{183.386}{547.354}{K}%
\TextAt{8}{197.95}{547.354}{and}%
\TextAt{16}{220.562}{547.354}{\Character{79}}%
\TextAt{8}{234.158}{547.354}{the}%
\TextAt{8}{254.164}{547.354}{integral}%
\TextAt{8}{296.621}{547.354}{closure}%
\TextAt{8}{335.552}{547.354}{of}%
\TextAt{14}{348.731}{547.354}{o}%
\TextAt{8}{358.322}{547.354}{in}%
\TextAt{13}{371.824}{547.354}{L}%
\TextAt{8}{379.788}{547.354}{.}%
\TextAt{8}{388.186}{547.354}{Then}%
\TextAt{16}{418.593}{547.354}{\Character{79}}%
\TextAt{8}{432.189}{547.354}{is}%
\TextAt{8}{443.802}{547.354}{again}%
\TextAt{8}{474.878}{547.354}{a}%
\TextAt{8}{484.483}{547.354}{Dedekind}%
\TextAt{8}{79.226}{530.308}{domain.}%
\TextAt{19}{79.226}{497.03}{Proof.}%
\TextAt{8}{116}{497.03}{Being}%
\TextAt{8}{149.405}{497.03}{the}%
\TextAt{8}{169.969}{497.03}{integral}%
\TextAt{8}{212.985}{497.03}{closure}%
\TextAt{8}{252.475}{497.03}{of}%
\TextAt{14}{266.213}{497.03}{o}%
\TextAt{8}{272.057}{497.03}{,}%
\TextAt{16}{279.716}{497.03}{\Character{79}}%
\TextAt{8}{293.871}{497.03}{is}%
\TextAt{8}{306.043}{497.03}{integrally}%
\TextAt{8}{358.492}{497.03}{closed.}%
\TextAt{8}{398.794}{497.03}{The}%
\TextAt{8}{423.255}{497.03}{fact}%
\TextAt{8}{446.749}{497.03}{that}%
\TextAt{8}{472.526}{497.03}{the}%
\TextAt{8}{493.09}{497.03}{nonzero}%
\TextAt{8}{79.226}{479.984}{prime}%
\TextAt{8}{113.429}{479.984}{ideals}%
\TextAt{14}{147.046}{479.984}{P}%
\TextAt{8}{161.876}{479.984}{of}%
\TextAt{16}{176.245}{479.984}{\Character{79}}%
\TextAt{8}{191.031}{479.984}{are}%
\TextAt{8}{211.581}{479.984}{maximal}%
\TextAt{8}{260.429}{479.984}{is}%
\TextAt{8}{273.232}{479.984}{proved}%
\TextAt{8}{312.324}{479.984}{similarly}%
\TextAt{8}{361.232}{479.984}{as}%
\TextAt{8}{376.641}{479.984}{in}%
\TextAt{8}{391.333}{479.984}{the}%
\TextAt{8}{412.529}{479.984}{case}%
\TextAt{14}{438.339}{479.984}{o}%
\TextAt{9}{449.275}{479.984}{\Character{61}}%
\TextAt{18}{463.471}{479.984}{Z}%
\TextAt{8}{476.378}{479.984}{(see}%
\TextAt{8}{500.885}{479.984}{(3.1)):}%
\TextAt{14}{79.226}{462.938}{p}%
\TextAt{9}{88.97}{462.938}{\Character{61}}%
\TextAt{14}{101.836}{462.938}{P}%
\TextAt{16}{114.562}{462.938}{\Character{92}}%
\TextAt{14}{125.365}{462.938}{o}%
\TextAt{8}{135.365}{462.938}{is}%
\TextAt{8}{147.387}{462.938}{a}%
\TextAt{8}{157.401}{462.938}{nonzero}%
\TextAt{8}{201.236}{462.938}{prime}%
\TextAt{8}{234.658}{462.938}{ideal}%
\TextAt{8}{262.879}{462.938}{of}%
\TextAt{14}{276.468}{462.938}{o}%
\TextAt{8}{282.312}{462.938}{.}%
\TextAt{8}{291.535}{462.938}{Thus}%
\TextAt{8}{321.442}{462.938}{the}%
\TextAt{8}{341.857}{462.938}{integral}%
\TextAt{8}{384.724}{462.938}{domain}%
\TextAt{16}{426.61}{462.938}{\Character{79}}%
\TextAt{9}{436.46}{462.938}{\Character{47}}%
\TextAt{14}{442.313}{462.938}{P}%
\TextAt{8}{456.361}{462.938}{is}%
\TextAt{8}{468.384}{462.938}{an}%
\TextAt{8}{484.901}{462.938}{extension}%
\TextAt{8}{79.226}{445.892}{of}%
\TextAt{8}{92.274}{445.892}{the}%
\TextAt{8}{112.147}{445.892}{field}%
\TextAt{14}{137.222}{445.892}{o}%
\TextAt{9}{143.066}{445.892}{\Character{47}}%
\TextAt{14}{148.919}{445.892}{p}%
\TextAt{8}{154.902}{445.892}{,}%
\TextAt{8}{161.825}{445.892}{and}%
\TextAt{8}{184.305}{445.892}{therefore}%
\TextAt{8}{233.123}{445.892}{has}%
\TextAt{8}{253.714}{445.892}{itself}%
\TextAt{8}{281.777}{445.892}{to}%
\TextAt{8}{295.805}{445.892}{be}%
\TextAt{8}{311.447}{445.892}{a}%
\TextAt{8}{320.919}{445.892}{field,}%
\TextAt{8}{349.302}{445.892}{because}%
\TextAt{8}{392.321}{445.892}{if}%
\TextAt{8}{402.762}{445.892}{it}%
\TextAt{8}{414.184}{445.892}{were}%
\TextAt{8}{440.884}{445.892}{not,}%
\TextAt{8}{464.724}{445.892}{then}%
\TextAt{8}{491.101}{445.892}{it}%
\TextAt{8}{502.523}{445.892}{would}%
\TextAt{8}{79.226}{428.846}{admit}%
\TextAt{8}{112.746}{428.846}{a}%
\TextAt{8}{122.2}{428.846}{nonzero}%
\TextAt{8}{165.475}{428.846}{prime}%
\TextAt{8}{198.338}{428.846}{ideal}%
\TextAt{8}{225.999}{428.846}{whose}%
\TextAt{8}{260.225}{428.846}{intersection}%
\TextAt{8}{322.748}{428.846}{with}%
\TextAt{14}{349.106}{428.846}{o}%
\TextAt{9}{354.951}{428.846}{\Character{47}}%
\TextAt{14}{360.804}{428.846}{p}%
\TextAt{8}{370.383}{428.846}{would}%
\TextAt{8}{404.225}{428.846}{again}%
\TextAt{8}{435.151}{428.846}{be}%
\TextAt{8}{450.774}{428.846}{a}%
\TextAt{8}{460.228}{428.846}{nonzero}%
\TextAt{8}{503.503}{428.846}{prime}%
\TextAt{8}{79.226}{411.8}{ideal}%
\TextAt{8}{107.016}{411.8}{in}%
\TextAt{14}{120.495}{411.8}{o}%
\TextAt{9}{126.339}{411.8}{\Character{47}}%
\TextAt{14}{132.192}{411.8}{p}%
\TextAt{8}{138.175}{411.8}{.}%
\TextAt{8}{146.566}{411.8}{It}%
\TextAt{8}{159.064}{411.8}{remains}%
\TextAt{8}{202.527}{411.8}{to}%
\TextAt{8}{216.664}{411.8}{show}%
\TextAt{8}{245.493}{411.8}{that}%
\TextAt{16}{270.688}{411.8}{\Character{79}}%
\TextAt{8}{284.262}{411.8}{is}%
\TextAt{8}{295.852}{411.8}{noetherian.}%
\TextAt{8}{358.555}{411.8}{In}%
\TextAt{8}{373.003}{411.8}{the}%
\TextAt{8}{392.985}{411.8}{case}%
\TextAt{8}{417.583}{411.8}{that}%
\TextAt{8}{442.778}{411.8}{is}%
\TextAt{8}{454.368}{411.8}{of}%
\TextAt{8}{467.524}{411.8}{chief}%
\TextAt{8}{494.656}{411.8}{interest}%
\TextAt{8}{79.226}{394.754}{to}%
\TextAt{8}{92.969}{394.754}{us,}%
\TextAt{8}{110.782}{394.754}{namely,}%
\TextAt{8}{153.247}{394.754}{if}%
\TextAt{13}{163.403}{394.754}{L}%
\TextAt{16}{174.688}{394.754}{\Character{106}}%
\TextAt{13}{181.33}{394.754}{K}%
\TextAt{8}{195.477}{394.754}{is}%
\TextAt{8}{206.673}{394.754}{a}%
\TextAt{8}{215.861}{394.754}{separable}%
\TextAt{8}{266.736}{394.754}{extension,}%
\TextAt{8}{321.3}{394.754}{the}%
\TextAt{8}{340.889}{394.754}{proof}%
\TextAt{8}{370.89}{394.754}{is}%
\TextAt{8}{382.086}{394.754}{very}%
\TextAt{8}{407.21}{394.754}{easy.}%
\TextAt{8}{436.343}{394.754}{Let}%
\TextAt{13}{456.745}{394.754}{\Character{11}}%
\TextAt{17}{464.19}{392.961}{1}%
\TextAt{13}{468.923}{394.754}{\Character{59}}%
\TextAt{13}{474.147}{394.754}{\Character{58}}%
\TextAt{13}{479.384}{394.754}{\Character{58}}%
\TextAt{13}{484.608}{394.754}{\Character{58}}%
\TextAt{13}{489.845}{394.754}{\Character{59}}%
\TextAt{13}{495.081}{394.754}{\Character{11}}%
\TextAt{21}{502.589}{392.961}{n}%
\TextAt{8}{511.555}{394.754}{be}%
\TextAt{8}{526.911}{394.754}{a}%
\TextAt{8}{79.226}{377.708}{basis}%
\TextAt{8}{107.047}{377.708}{of}%
\TextAt{13}{119.457}{377.708}{L}%
\TextAt{16}{130.743}{377.708}{\Character{106}}%
\TextAt{13}{137.384}{377.708}{K}%
\TextAt{8}{151.18}{377.708}{contained}%
\TextAt{8}{203.27}{377.708}{in}%
\TextAt{16}{216.003}{377.708}{\Character{79}}%
\TextAt{8}{225.853}{377.708}{,}%
\TextAt{8}{232.266}{377.708}{of}%
\TextAt{8}{244.677}{377.708}{discriminant}%
\TextAt{13}{311.137}{377.708}{d}%
\TextAt{9}{320.54}{377.708}{\Character{61}}%
\TextAt{13}{332.966}{377.708}{d}%
\TextAt{9}{339.048}{377.708}{\Character{40}}%
\TextAt{13}{343.601}{377.708}{\Character{11}}%
\TextAt{17}{351.046}{375.915}{1}%
\TextAt{13}{355.779}{377.708}{\Character{59}}%
\TextAt{13}{361.003}{377.708}{\Character{58}}%
\TextAt{13}{366.24}{377.708}{\Character{58}}%
\TextAt{13}{371.464}{377.708}{\Character{58}}%
\TextAt{13}{376.701}{377.708}{\Character{59}}%
\TextAt{13}{381.937}{377.708}{\Character{11}}%
\TextAt{21}{389.445}{375.915}{n}%
\TextAt{9}{395.081}{377.708}{\Character{41}}%
\TextAt{8}{399.634}{377.708}{.}%
\TextAt{8}{407.776}{377.708}{Then}%
\TextAt{13}{437.414}{377.708}{d}%
\NativeMathAt{446.817}{377.708}{11.955}{\neq}%
\TextAt{9}{459.226}{377.708}{0}%
\TextAt{8}{468.073}{377.708}{by}%
\TextAt{8}{483.413}{377.708}{(2.8),}%
\TextAt{8}{513.904}{377.708}{and}%
\TextAt{8}{79.226}{360.662}{(2.9)}%
\TextAt{8}{107.078}{360.662}{tells}%
\TextAt{8}{131.726}{360.662}{us}%
\TextAt{8}{146.619}{360.662}{that}%
\TextAt{16}{171.865}{360.662}{\Character{79}}%
\TextAt{8}{185.489}{360.662}{is}%
\TextAt{8}{197.13}{360.662}{contained}%
\TextAt{8}{250.016}{360.662}{in}%
\TextAt{8}{263.546}{360.662}{the}%
\TextAt{8}{283.579}{360.662}{finitely}%
\TextAt{8}{322.8}{360.662}{generated}%
\TextAt{14}{376.009}{360.662}{o}%
\TextAt{8}{381.853}{360.662}{-module}%
\TextAt{14}{426.921}{360.662}{o}%
\TextAt{13}{432.765}{360.662}{\Character{11}}%
\TextAt{17}{440.211}{358.869}{1}%
\TextAt{9}{444.943}{360.662}{\Character{47}}%
\TextAt{13}{450.796}{360.662}{d}%
\TextAt{9}{459.284}{360.662}{\Character{43}}%
\TextAt{16}{470.793}{360.662}{\Character{1}}%
\TextAt{16}{476.089}{360.662}{\Character{1}}%
\TextAt{16}{481.397}{360.662}{\Character{1}}%
\TextAt{9}{487.146}{360.662}{\Character{43}}%
\TextAt{14}{498.656}{360.662}{o}%
\TextAt{13}{504.5}{360.662}{\Character{11}}%
\TextAt{21}{511.946}{358.869}{n}%
\TextAt{9}{517.582}{360.662}{\Character{47}}%
\TextAt{13}{523.435}{360.662}{d}%
\TextAt{8}{529.518}{360.662}{.}%
\TextAt{8}{79.226}{343.616}{Every}%
\TextAt{8}{113.822}{343.616}{ideal}%
\TextAt{8}{142.728}{343.616}{of}%
\TextAt{16}{157}{343.616}{\Character{79}}%
\TextAt{8}{171.69}{343.616}{is}%
\TextAt{8}{184.396}{343.616}{also}%
\TextAt{8}{208.818}{343.616}{contained}%
\TextAt{8}{262.77}{343.616}{in}%
\TextAt{8}{277.365}{343.616}{this}%
\TextAt{8}{301.13}{343.616}{finitely}%
\TextAt{8}{341.417}{343.616}{generated}%
\TextAt{14}{395.691}{343.616}{o}%
\TextAt{8}{401.535}{343.616}{-module,}%
\TextAt{8}{451.156}{343.616}{and}%
\TextAt{8}{474.861}{343.616}{therefore}%
\TextAt{8}{524.903}{343.616}{is}%
\TextAt{8}{79.226}{326.57}{itself}%
\TextAt{8}{108.792}{326.57}{an}%
\TextAt{14}{126.271}{326.57}{o}%
\TextAt{8}{132.116}{326.57}{-module}%
\TextAt{8}{178.527}{326.57}{of}%
\TextAt{8}{193.077}{326.57}{finite}%
\TextAt{8}{224.209}{326.57}{type,}%
\TextAt{8}{255.323}{326.57}{hence}%
\TextAt{8}{289.05}{326.57}{a}%
\TextAt{8}{300.025}{326.57}{fortiori}%
\TextAt{8}{340.602}{326.57}{a}%
\TextAt{8}{351.578}{326.57}{finitely}%
\TextAt{8}{392.143}{326.57}{generated}%
\TextAt{16}{446.695}{326.57}{\Character{79}}%
\TextAt{8}{456.545}{326.57}{-module.}%
\TextAt{8}{509.947}{326.57}{This}%
\TextAt{8}{79.226}{309.524}{shows}%
\TextAt{8}{112.851}{309.524}{that}%
\TextAt{16}{138.227}{309.524}{\Character{79}}%
\TextAt{8}{151.981}{309.524}{is}%
\TextAt{8}{163.752}{309.524}{noetherian,}%
\TextAt{8}{225.223}{309.524}{provided}%
\TextAt{13}{273.361}{309.524}{L}%
\TextAt{16}{284.659}{309.524}{\Character{106}}%
\TextAt{13}{291.313}{309.524}{K}%
\TextAt{8}{306.035}{309.524}{is}%
\TextAt{8}{317.806}{309.524}{separable.}%
\TextAt{8}{373.821}{309.524}{We}%
\TextAt{8}{393.973}{309.524}{ask}%
\TextAt{8}{414.531}{309.524}{the}%
\TextAt{8}{434.694}{309.524}{reader’s}%
\TextAt{8}{478.338}{309.524}{permission}%
\TextAt{8}{79.226}{292.478}{to}%
\TextAt{8}{93.364}{292.478}{content}%
\TextAt{8}{134.82}{292.478}{ourselves}%
\TextAt{8}{184.202}{292.478}{for}%
\TextAt{8}{201.915}{292.478}{the}%
\TextAt{8}{221.9}{292.478}{time}%
\TextAt{8}{248.388}{292.478}{being}%
\TextAt{8}{279.753}{292.478}{with}%
\TextAt{8}{306.241}{292.478}{this}%
\TextAt{8}{328.892}{292.478}{case.}%
\TextAt{8}{358.156}{292.478}{We}%
\TextAt{8}{378.129}{292.478}{shall}%
\TextAt{8}{405.334}{292.478}{come}%
\TextAt{8}{435.074}{292.478}{back}%
\TextAt{8}{462.219}{292.478}{to}%
\TextAt{8}{476.358}{292.478}{the}%
\TextAt{8}{496.342}{292.478}{general}%
\TextAt{8}{79.226}{275.432}{case}%
\TextAt{8}{104.821}{275.432}{on}%
\TextAt{8}{121.904}{275.432}{a}%
\TextAt{8}{132.483}{275.432}{more}%
\TextAt{8}{162.573}{275.432}{convenient}%
\TextAt{8}{221.283}{275.432}{occasion.}%
\TextAt{8}{274.87}{275.432}{In}%
\TextAt{8}{290.315}{275.432}{fact,}%
\TextAt{8}{317.682}{275.432}{we}%
\TextAt{8}{335.733}{275.432}{shall}%
\TextAt{8}{363.934}{275.432}{give}%
\TextAt{8}{388.824}{275.432}{the}%
\TextAt{8}{409.804}{275.432}{proof}%
\TextAt{8}{441.197}{275.432}{in}%
\TextAt{8}{455.673}{275.432}{a}%
\TextAt{8}{466.252}{275.432}{more}%
\TextAt{8}{496.342}{275.432}{general}%
\TextAt{8}{79.226}{258.386}{framework}%
\TextAt{8}{136.79}{258.386}{in}%
\TextAt{8}{150.443}{258.386}{§12}%
\TextAt{8}{171.735}{258.386}{(see}%
\TextAt{8}{195.203}{258.386}{(12.8)).}%
\Vector{q 0 G 0 g 0.399 w 452.334 -515.614 m 452.334 -507.544 l S Q}%
\Vector{q 0 G 0 g 0.399 w 452.534 -507.743 m 459.707 -507.743 l S Q}%
\Vector{q 0 G 0 g 0.399 w 452.534 -515.414 m 459.707 -515.414 l S Q}%
\Vector{q 0 G 0 g 0.399 w 459.906 -515.614 m 459.906 -507.544 l S Q}%
\TextAt{8}{97.159}{230.236}{For}%
\TextAt{8}{118.128}{230.236}{a}%
\TextAt{8}{127.884}{230.236}{prime}%
\TextAt{8}{161.047}{230.236}{ideal}%
\TextAt{14}{189.01}{230.236}{p}%
\TextAt{8}{198.891}{230.236}{of}%
\TextAt{14}{212.221}{230.236}{o}%
\TextAt{8}{221.963}{230.236}{one}%
\TextAt{8}{243.422}{230.236}{always}%
\TextAt{8}{280.89}{230.236}{has}%
\TextAt{14}{283.658}{204.211}{p}%
\TextAt{16}{289.641}{204.211}{\Character{79}}%
\NativeMathAt{302.804}{204.211}{11.955}{\neq}%
\TextAt{16}{315.237}{204.211}{\Character{79}}%
\TextAt{13}{325.086}{204.211}{\Character{58}}%
\TextAt{8}{79.226}{178.186}{In}%
\TextAt{8}{94.726}{178.186}{fact,}%
\TextAt{8}{122.161}{178.186}{let}%
\TextAt{13}{139.944}{178.186}{\Character{25}}%
\NativeMathAt{151.831}{178.186}{11.955}{\in}%
\TextAt{14}{164.619}{178.186}{p}%
\TextAt{16}{173.857}{178.186}{\Character{110}}%
\TextAt{14}{183.09}{178.186}{p}%
\TextAt{17}{189.073}{182.524}{2}%
\TextAt{8}{198.581}{178.186}{(}%
\TextAt{14}{203.136}{178.186}{p}%
\NativeMathAt{213.937}{178.186}{11.955}{\neq}%
\TextAt{9}{227.841}{178.186}{0}%
\TextAt{8}{233.713}{178.186}{),}%
\TextAt{8}{246.515}{178.186}{so}%
\TextAt{8}{261.763}{178.186}{that}%
\TextAt{13}{288.01}{178.186}{\Character{25}}%
\TextAt{14}{295.08}{178.186}{o}%
\TextAt{9}{305.742}{178.186}{\Character{61}}%
\TextAt{14}{319.664}{178.186}{pa}%
\TextAt{8}{336.406}{178.186}{with}%
\TextAt{14}{363.944}{178.186}{p}%
\TextAt{18}{374.745}{178.186}{\Character{45}}%
\TextAt{14}{382.884}{178.186}{a}%
\TextAt{8}{388.867}{178.186}{,}%
\TextAt{8}{397.114}{178.186}{hence}%
\TextAt{14}{430.499}{178.186}{p}%
\TextAt{9}{439.737}{178.186}{\Character{43}}%
\TextAt{14}{452.097}{178.186}{a}%
\TextAt{9}{462.898}{178.186}{\Character{61}}%
\TextAt{14}{476.82}{178.186}{o}%
\TextAt{8}{482.664}{178.186}{.}%
\TextAt{8}{493.748}{178.186}{Writing}%
\TextAt{9}{79.226}{161.14}{1}%
\TextAt{9}{89.113}{161.14}{\Character{61}}%
\TextAt{13}{102.267}{161.14}{b}%
\TextAt{9}{110.189}{161.14}{\Character{43}}%
\TextAt{13}{122.239}{161.14}{s}%
\TextAt{8}{127.753}{161.14}{,}%
\TextAt{8}{135.431}{161.14}{with}%
\TextAt{13}{162.514}{161.14}{b}%
\NativeMathAt{171.533}{161.14}{11.955}{\in}%
\TextAt{14}{183.544}{161.14}{p}%
\TextAt{8}{193.848}{161.14}{and}%
\TextAt{13}{217.033}{161.14}{s}%
\NativeMathAt{226.589}{161.14}{11.955}{\in}%
\TextAt{14}{238.601}{161.14}{a}%
\TextAt{8}{244.584}{161.14}{,}%
\TextAt{8}{252.262}{161.14}{we}%
\TextAt{8}{269.912}{161.14}{find}%
\TextAt{13}{293.744}{161.14}{s}%
\NativeMathAt{303.299}{161.14}{11.955}{\notin}%
\TextAt{14}{315.311}{161.14}{p}%
\TextAt{8}{325.614}{161.14}{and}%
\TextAt{13}{348.8}{161.14}{s}%
\TextAt{14}{354.314}{161.14}{p}%
\NativeMathAt{364.339}{161.14}{11.955}{\subseteq}%
\TextAt{14}{377.679}{161.14}{pa}%
\TextAt{9}{393.687}{161.14}{\Character{61}}%
\TextAt{13}{406.833}{161.14}{\Character{25}}%
\TextAt{14}{413.902}{161.14}{o}%
\TextAt{8}{419.747}{161.14}{.}%
\TextAt{8}{429.464}{161.14}{If}%
\TextAt{8}{441.579}{161.14}{one}%
\TextAt{8}{463.462}{161.14}{had}%
\TextAt{14}{486.647}{161.14}{p}%
\TextAt{16}{492.63}{161.14}{\Character{79}}%
\TextAt{9}{506.522}{161.14}{\Character{61}}%
\TextAt{16}{519.668}{161.14}{\Character{79}}%
\TextAt{8}{529.518}{161.14}{,}%
\TextAt{8}{79.226}{144.094}{then}%
\TextAt{8}{105.931}{144.094}{it}%
\TextAt{8}{117.68}{144.094}{would}%
\TextAt{8}{151.869}{144.094}{follow}%
\TextAt{8}{185.735}{144.094}{that}%
\TextAt{13}{211.149}{144.094}{s}%
\TextAt{16}{216.663}{144.094}{\Character{79}}%
\TextAt{9}{229.91}{144.094}{\Character{61}}%
\TextAt{13}{242.412}{144.094}{s}%
\TextAt{14}{247.926}{144.094}{p}%
\TextAt{16}{253.909}{144.094}{\Character{79}}%
\NativeMathAt{267.143}{144.094}{11.955}{\subseteq}%
\TextAt{13}{279.852}{144.094}{\Character{25}}%
\TextAt{16}{286.922}{144.094}{\Character{79}}%
\TextAt{8}{296.771}{144.094}{,}%
\TextAt{8}{303.977}{144.094}{so}%
\TextAt{8}{318.392}{144.094}{that}%
\TextAt{13}{343.806}{144.094}{s}%
\TextAt{9}{352.717}{144.094}{\Character{61}}%
\TextAt{13}{365.219}{144.094}{\Character{25}}%
\TextAt{13}{372.273}{144.094}{x}%
\TextAt{8}{382.883}{144.094}{for}%
\TextAt{8}{400.813}{144.094}{some}%
\TextAt{13}{430.184}{144.094}{x}%
\NativeMathAt{440.233}{144.094}{11.955}{\in}%
\TextAt{16}{451.59}{144.094}{\Character{79}}%
\TextAt{16}{464.12}{144.094}{\Character{92}}%
\TextAt{13}{474.795}{144.094}{K}%
\TextAt{9}{489.01}{144.094}{\Character{61}}%
\TextAt{14}{501.512}{144.094}{o}%
\TextAt{8}{507.356}{144.094}{,}%
\TextAt{8}{514.562}{144.094}{i.e.,}%
\TextAt{13}{79.226}{127.048}{s}%
\NativeMathAt{88.061}{127.048}{11.955}{\in}%
\TextAt{14}{99.352}{127.048}{p}%
\TextAt{8}{105.335}{127.048}{,}%
\TextAt{8}{112.484}{127.048}{a}%
\TextAt{8}{122.24}{127.048}{contradiction.}%
\TextAt{8}{97.159}{107.013}{A}%
\TextAt{8}{109.835}{107.013}{prime}%
\TextAt{8}{143.003}{107.013}{ideal}%
\TextAt{14}{170.97}{107.013}{p}%
\NativeMathAt{180.281}{107.013}{11.955}{\neq}%
\TextAt{9}{192.703}{107.013}{0}%
\TextAt{8}{202.467}{107.013}{of}%
\TextAt{8}{215.801}{107.013}{the}%
\TextAt{8}{235.962}{107.013}{ring}%
\TextAt{14}{260.031}{107.013}{o}%
\TextAt{8}{269.777}{107.013}{decomposes}%
\TextAt{8}{333.311}{107.013}{in}%
\TextAt{16}{346.968}{107.013}{\Character{79}}%
\TextAt{8}{360.719}{107.013}{in}%
\TextAt{8}{374.376}{107.013}{a}%
\TextAt{8}{384.135}{107.013}{unique}%
\TextAt{8}{422.18}{107.013}{way}%
\TextAt{8}{445.927}{107.013}{into}%
\TextAt{8}{469.674}{107.013}{a}%
\TextAt{8}{479.434}{107.013}{product}%
\TextAt{8}{523.337}{107.013}{of}%
\TextAt{8}{79.226}{89.967}{prime}%
\TextAt{8}{112.39}{89.967}{ideals,}%
\TextAt{14}{261.142}{72.921}{p}%
\TextAt{16}{267.125}{72.921}{\Character{79}}%
\TextAt{9}{280.296}{72.921}{\Character{61}}%
\TextAt{14}{292.721}{72.921}{P}%
\TextAt{21}{302.614}{78.146}{e}%
\TextAt{24}{306.543}{77.039}{1}%
\TextAt{17}{302.614}{69.99}{1}%
\TextAt{16}{313.184}{72.921}{\Character{1}}%
\TextAt{16}{318.48}{72.921}{\Character{1}}%
\TextAt{16}{323.788}{72.921}{\Character{1}}%
\TextAt{14}{329.124}{72.921}{P}%
\TextAt{21}{339.018}{77.858}{e}%
\TextAt{22}{342.946}{76.861}{r}%
\TextAt{21}{339.018}{69.966}{r}%
\TextAt{13}{347.602}{72.921}{\Character{58}}%
\TextAt{8}{300.14}{43.033}{42}%
\Destination{78.23}{809.815}{nav.page.42}%
\Destination{79.23}{773.075}{nav.section.1.8}%
\Destination{79.36}{579.355}{nav.print.45}%
\Destination{79.36}{578.355}{nav.number.1.8.1}%
\Destination{79.23}{773.075}{outline.13}%
\Hotlink{457.826}{655.555}{6.858}{11.047}{nav.section.1.3}%
\Hotlink{483.058}{655.555}{12.716}{11.047}{nav.print.17}%
\Hotlink{155.622}{256.155}{12.716}{11.047}{nav.section.1.12}%
\Hotlink{500.385}{477.755}{25.078}{11.047}{nav.number.1.3.1}%
\Hotlink{482.913}{375.555}{25.078}{11.047}{nav.number.1.2.8}%
\Hotlink{78.726}{358.455}{25.078}{11.047}{nav.number.1.2.9}%
\Hotlink{194.703}{256.155}{30.936}{11.047}{nav.number.1.12.8}%
}

% Page 55 (body page 43)
\Page{%
\TextAt{2}{79.226}{796.943}{I.}%
\TextAt{2}{89.83}{796.943}{Algebraic}%
\TextAt{2}{138.648}{796.943}{Integers}%
\Vector{q 0 G 0 g 0.299 w 7.226 18.729 m 460.769 18.729 l S Q}%
\TextAt{8}{79.226}{761.124}{Instead}%
\TextAt{8}{120.547}{761.124}{of}%
\TextAt{14}{133.845}{761.124}{p}%
\TextAt{16}{139.829}{761.124}{\Character{79}}%
\TextAt{8}{153.544}{761.124}{we}%
\TextAt{8}{170.739}{761.124}{will}%
\TextAt{8}{192.813}{761.124}{often}%
\TextAt{8}{222.37}{761.124}{write}%
\TextAt{8}{252.25}{761.124}{simply}%
\TextAt{14}{289.674}{761.124}{p}%
\TextAt{8}{295.657}{761.124}{.}%
\TextAt{8}{304.094}{761.124}{The}%
\TextAt{8}{328.116}{761.124}{prime}%
\TextAt{8}{361.248}{761.124}{ideals}%
\TextAt{14}{393.794}{761.124}{P}%
\TextAt{21}{403.687}{759.33}{i}%
\TextAt{8}{410.934}{761.124}{occurring}%
\TextAt{8}{462.609}{761.124}{in}%
\TextAt{8}{476.229}{761.124}{the}%
\TextAt{8}{496.354}{761.124}{decom-}%
\TextAt{8}{79.226}{744.078}{position}%
\TextAt{8}{123.435}{744.078}{are}%
\TextAt{8}{142.538}{744.078}{precisely}%
\TextAt{8}{189.987}{744.078}{those}%
\TextAt{8}{220.208}{744.078}{prime}%
\TextAt{8}{252.964}{744.078}{ideals}%
\TextAt{14}{285.135}{744.078}{P}%
\TextAt{8}{298.517}{744.078}{of}%
\TextAt{16}{311.44}{744.078}{\Character{79}}%
\TextAt{8}{324.779}{744.078}{which}%
\TextAt{8}{357.858}{744.078}{lie}%
\TextAt{8}{373.052}{744.078}{over}%
\TextAt{14}{397.69}{744.078}{p}%
\TextAt{8}{407.163}{744.078}{in}%
\TextAt{8}{420.408}{744.078}{the}%
\TextAt{8}{440.157}{744.078}{sense}%
\TextAt{8}{469.78}{744.078}{that}%
\TextAt{8}{494.741}{744.078}{one}%
\TextAt{8}{515.793}{744.078}{has}%
\TextAt{8}{79.226}{727.032}{the}%
\TextAt{8}{99.383}{727.032}{relation}%
\TextAt{14}{278.997}{709.986}{p}%
\TextAt{9}{288.301}{709.986}{\Character{61}}%
\TextAt{14}{300.726}{709.986}{P}%
\TextAt{16}{313.276}{709.986}{\Character{92}}%
\TextAt{14}{323.903}{709.986}{o}%
\TextAt{13}{329.747}{709.986}{\Character{58}}%
\TextAt{8}{79.226}{686.464}{This}%
\TextAt{8}{105.132}{686.464}{we}%
\TextAt{8}{121.545}{686.464}{also}%
\TextAt{8}{144.211}{686.464}{denote}%
\TextAt{8}{181.115}{686.464}{for}%
\TextAt{8}{198.186}{686.464}{short}%
\TextAt{8}{227.356}{686.464}{by}%
\TextAt{14}{242.801}{686.464}{P}%
\TextAt{16}{256.015}{686.464}{\Character{106}}%
\TextAt{14}{262.656}{686.464}{p}%
\TextAt{8}{268.64}{686.464}{,}%
\TextAt{8}{275.137}{686.464}{and}%
\TextAt{8}{297.086}{686.464}{we}%
\TextAt{8}{313.499}{686.464}{call}%
\TextAt{14}{334.145}{686.464}{P}%
\TextAt{8}{347.121}{686.464}{a}%
\TextAt{12}{356.062}{686.464}{prime}%
\TextAt{12}{393.669}{686.464}{divisor}%
\TextAt{8}{436.348}{686.464}{of}%
\TextAt{14}{448.864}{686.464}{p}%
\TextAt{8}{454.847}{686.464}{.}%
\TextAt{8}{463.024}{686.464}{The}%
\TextAt{8}{486.264}{686.464}{exponent}%
\TextAt{13}{79.226}{669.418}{e}%
\TextAt{21}{84.652}{667.625}{i}%
\TextAt{8}{91.93}{669.418}{is}%
\TextAt{8}{103.694}{669.418}{called}%
\TextAt{8}{136.858}{669.418}{the}%
\TextAt{12}{157.014}{669.418}{ramification}%
\TextAt{12}{231.662}{669.418}{index}%
\TextAt{8}{263.583}{669.418}{,}%
\TextAt{8}{270.732}{669.418}{and}%
\TextAt{8}{293.494}{669.418}{the}%
\TextAt{8}{313.651}{669.418}{degree}%
\TextAt{8}{350.066}{669.418}{of}%
\TextAt{8}{363.396}{669.418}{the}%
\TextAt{8}{383.553}{669.418}{field}%
\TextAt{8}{408.91}{669.418}{extension}%
\TextAt{13}{262.022}{640.417}{f}%
\TextAt{21}{267.792}{638.624}{i}%
\TextAt{9}{274.494}{640.417}{\Character{61}}%
\TextAt{9}{286.904}{640.417}{\Character{91}}%
\TextAt{16}{290.171}{640.417}{\Character{79}}%
\TextAt{9}{300.021}{640.417}{\Character{47}}%
\TextAt{14}{305.874}{640.417}{P}%
\TextAt{21}{315.767}{638.624}{i}%
\TextAt{9}{322.469}{640.417}{\Character{58}}%
\TextAt{14}{329.042}{640.417}{o}%
\TextAt{9}{334.886}{640.417}{\Character{47}}%
\TextAt{14}{340.739}{640.417}{p}%
\TextAt{9}{346.722}{640.417}{\Character{93}}%
\TextAt{8}{79.226}{611.416}{is}%
\TextAt{8}{90.169}{611.416}{called}%
\TextAt{8}{122.511}{611.416}{the}%
\TextAt{12}{141.846}{611.416}{inertia}%
\TextAt{12}{183.749}{611.416}{degree}%
\TextAt{8}{224.926}{611.416}{of}%
\TextAt{14}{237.435}{611.416}{P}%
\TextAt{21}{247.328}{609.623}{i}%
\TextAt{8}{253.786}{611.416}{over}%
\TextAt{14}{278.011}{611.416}{p}%
\TextAt{8}{283.994}{611.416}{.}%
\TextAt{8}{292.168}{611.416}{If}%
\TextAt{8}{303.039}{611.416}{the}%
\TextAt{8}{322.375}{611.416}{extension}%
\TextAt{13}{373.319}{611.416}{L}%
\TextAt{16}{384.605}{611.416}{\Character{106}}%
\TextAt{13}{391.246}{611.416}{K}%
\TextAt{8}{405.14}{611.416}{is}%
\TextAt{8}{416.083}{611.416}{separable,}%
\TextAt{8}{470.121}{611.416}{the}%
\TextAt{8}{489.456}{611.416}{numbers}%
\TextAt{13}{79.226}{594.37}{e}%
\TextAt{21}{84.652}{592.577}{i}%
\TextAt{13}{88.033}{594.37}{\Character{59}}%
\TextAt{13}{93.257}{594.37}{f}%
\TextAt{21}{99.047}{592.577}{i}%
\TextAt{8}{106.326}{594.37}{and}%
\TextAt{8}{129.088}{594.37}{the}%
\TextAt{8}{149.245}{594.37}{degree}%
\TextAt{13}{185.66}{594.37}{n}%
\TextAt{9}{195.969}{594.37}{\Character{61}}%
\TextAt{9}{208.379}{594.37}{\Character{91}}%
\TextAt{13}{211.646}{594.37}{L}%
\TextAt{9}{222.931}{594.37}{\Character{58}}%
\TextAt{13}{229.504}{594.37}{K}%
\TextAt{9}{240.321}{594.37}{\Character{93}}%
\TextAt{8}{247.47}{594.37}{are}%
\TextAt{8}{266.981}{594.37}{connected}%
\TextAt{8}{321.604}{594.37}{by}%
\TextAt{8}{337.863}{594.37}{the}%
\TextAt{8}{358.02}{594.37}{following}%
\TextAt{8}{407.454}{594.37}{law.}%
\TextAt{12}{79.226}{568.358}{Proposition}%
\TextAt{12}{151.937}{568.358}{(8.2).}%
\TextAt{8}{187.253}{568.358}{Let}%
\TextAt{13}{208.222}{568.358}{L}%
\TextAt{16}{219.508}{568.358}{\Character{106}}%
\TextAt{13}{226.149}{568.358}{K}%
\TextAt{8}{240.865}{568.358}{be}%
\TextAt{8}{256.789}{568.358}{separable.}%
\TextAt{8}{312.783}{568.358}{Then}%
\TextAt{8}{343.34}{568.358}{we}%
\TextAt{8}{360.568}{568.358}{have}%
\TextAt{8}{387.562}{568.358}{the}%
\TextAt{12}{407.719}{568.358}{fundamental}%
\TextAt{12}{485.236}{568.358}{identity}%
\TextAt{21}{281.001}{547.778}{r}%
\TextAt{20}{274.395}{544.192}{\Character{88}}%
\TextAt{21}{276.178}{518.997}{i}%
\TextAt{17}{279.061}{518.997}{\Character{61}1}%
\TextAt{13}{293.657}{532.834}{e}%
\TextAt{21}{299.082}{531.041}{i}%
\TextAt{13}{302.463}{532.834}{f}%
\TextAt{21}{308.233}{531.041}{i}%
\TextAt{9}{314.935}{532.834}{\Character{61}}%
\TextAt{13}{327.361}{532.834}{n\Character{58}}%
\TextAt{19}{79.226}{478.501}{Proof.}%
\TextAt{8}{116}{478.501}{The}%
\TextAt{8}{140.054}{478.501}{proof}%
\TextAt{8}{170.623}{478.501}{is}%
\TextAt{8}{182.387}{478.501}{based}%
\TextAt{8}{214.965}{478.501}{on}%
\TextAt{8}{231.224}{478.501}{the}%
\TextAt{8}{251.381}{478.501}{Chinese}%
\TextAt{8}{295.005}{478.501}{remainder}%
\TextAt{8}{350.286}{478.501}{theorem}%
\TextAt{16}{254.013}{443.54}{\Character{79}}%
\TextAt{9}{263.863}{443.54}{\Character{47}}%
\TextAt{14}{269.716}{443.54}{p}%
\TextAt{16}{275.699}{443.54}{\Character{79}}%
\NativeMathAt{288.869}{443.54}{11.955}{\cong}%
\TextAt{21}{308.493}{458.484}{r}%
\TextAt{20}{301.489}{454.897}{\Character{77}}%
\TextAt{21}{303.67}{429.703}{i}%
\TextAt{17}{306.553}{429.703}{\Character{61}1}%
\TextAt{16}{321.547}{443.54}{\Character{79}}%
\TextAt{9}{331.397}{443.54}{\Character{47}}%
\TextAt{14}{337.25}{443.54}{P}%
\TextAt{21}{347.143}{448.873}{e}%
\TextAt{22}{351.071}{447.658}{i}%
\TextAt{21}{347.143}{440.476}{i}%
\TextAt{13}{354.731}{443.54}{\Character{58}}%
\TextAt{16}{79.226}{406.59}{\Character{79}}%
\TextAt{9}{89.076}{406.59}{\Character{47}}%
\TextAt{14}{94.929}{406.59}{p}%
\TextAt{16}{100.912}{406.59}{\Character{79}}%
\TextAt{8}{114.659}{406.59}{and}%
\TextAt{16}{137.422}{406.59}{\Character{79}}%
\TextAt{9}{147.272}{406.59}{\Character{47}}%
\TextAt{14}{153.125}{406.59}{P}%
\TextAt{21}{163.018}{411.923}{e}%
\TextAt{22}{166.946}{410.708}{i}%
\TextAt{21}{163.018}{403.526}{i}%
\TextAt{8}{174.503}{406.59}{are}%
\TextAt{8}{194.014}{406.59}{vector}%
\TextAt{8}{229.139}{406.59}{spaces}%
\TextAt{8}{265.028}{406.59}{over}%
\TextAt{8}{290.074}{406.59}{the}%
\TextAt{8}{310.231}{406.59}{field}%
\TextAt{13}{335.587}{406.59}{\Character{20}}%
\TextAt{9}{345.647}{406.59}{\Character{61}}%
\TextAt{14}{358.072}{406.59}{o}%
\TextAt{9}{363.917}{406.59}{\Character{47}}%
\TextAt{14}{369.77}{406.59}{p}%
\TextAt{8}{375.753}{406.59}{,}%
\TextAt{8}{382.902}{406.59}{and}%
\TextAt{8}{405.664}{406.59}{it}%
\TextAt{8}{417.369}{406.59}{suffices}%
\TextAt{8}{457.155}{406.59}{to}%
\TextAt{8}{471.466}{406.59}{show}%
\TextAt{8}{500.469}{406.59}{that}%
\TextAt{8}{175.46}{377.588}{dim}%
\TextAt{21}{194.971}{375.795}{\Character{20}}%
\TextAt{9}{200.349}{377.589}{\Character{40}}%
\TextAt{16}{204.901}{377.589}{\Character{79}}%
\TextAt{9}{214.751}{377.589}{\Character{47}}%
\TextAt{14}{220.604}{377.589}{p}%
\TextAt{16}{226.587}{377.589}{\Character{79}}%
\TextAt{9}{236.437}{377.589}{\Character{41}}%
\TextAt{9}{244.292}{377.589}{\Character{61}}%
\TextAt{13}{256.736}{377.589}{n}%
\TextAt{8}{287.634}{377.588}{and}%
\TextAt{8}{332.402}{377.588}{dim}%
\TextAt{21}{351.913}{375.795}{\Character{20}}%
\TextAt{9}{357.291}{377.589}{\Character{40}}%
\TextAt{16}{361.843}{377.589}{\Character{79}}%
\TextAt{9}{371.693}{377.589}{\Character{47}}%
\TextAt{14}{377.546}{377.589}{P}%
\TextAt{21}{387.439}{382.921}{e}%
\TextAt{22}{391.367}{381.706}{i}%
\TextAt{21}{387.439}{374.525}{i}%
\TextAt{9}{395.027}{377.589}{\Character{41}}%
\TextAt{9}{402.882}{377.589}{\Character{61}}%
\TextAt{13}{415.326}{377.589}{e}%
\TextAt{21}{420.751}{375.795}{i}%
\TextAt{13}{424.133}{377.589}{f}%
\TextAt{21}{429.903}{375.795}{i}%
\TextAt{13}{433.284}{377.589}{\Character{58}}%
\TextAt{8}{79.226}{348.587}{In}%
\TextAt{8}{95.663}{348.587}{order}%
\TextAt{8}{128.048}{348.587}{to}%
\TextAt{8}{144.174}{348.587}{prove}%
\TextAt{8}{177.54}{348.587}{the}%
\TextAt{8}{199.512}{348.587}{first}%
\TextAt{8}{225.454}{348.587}{identity,}%
\TextAt{8}{273.249}{348.587}{let}%
\TextAt{13}{291.969}{348.587}{\Character{33}}%
\TextAt{17}{299.258}{346.794}{1}%
\TextAt{13}{303.991}{348.587}{\Character{59}}%
\TextAt{13}{309.215}{348.587}{\Character{58}}%
\TextAt{13}{314.452}{348.587}{\Character{58}}%
\TextAt{13}{319.676}{348.587}{\Character{58}}%
\TextAt{13}{324.913}{348.587}{\Character{59}}%
\TextAt{13}{330.149}{348.587}{\Character{33}}%
\TextAt{21}{337.501}{346.794}{m}%
\NativeMathAt{351.905}{348.587}{11.955}{\in}%
\TextAt{16}{366.275}{348.587}{\Character{79}}%
\TextAt{8}{381.853}{348.587}{be}%
\TextAt{8}{399.593}{348.587}{representatives}%
\TextAt{8}{481.21}{348.587}{of}%
\TextAt{8}{496.355}{348.587}{a}%
\TextAt{8}{507.927}{348.587}{basis}%
\Vector{q 0 G 0 g 0.478 w 7.226 -435.638 m 14.944 -435.638 l S Q}%
\TextAt{13}{79.226}{331.541}{\Character{33}}%
\TextAt{17}{86.944}{329.748}{1}%
\TextAt{13}{91.677}{331.541}{\Character{59}}%
\TextAt{13}{96.901}{331.541}{\Character{58}}%
\TextAt{13}{102.138}{331.541}{\Character{58}}%
\TextAt{13}{107.362}{331.541}{\Character{58}}%
\TextAt{13}{112.599}{331.541}{\Character{59}}%
\Vector{q 0 G 0 g 0.478 w 45.897 -435.638 m 53.616 -435.638 l S Q}%
\TextAt{13}{117.897}{331.541}{\Character{33}}%
\TextAt{21}{125.616}{329.748}{m}%
\TextAt{8}{136.476}{331.541}{of}%
\TextAt{16}{148.78}{331.541}{\Character{79}}%
\TextAt{9}{158.629}{331.541}{\Character{47}}%
\TextAt{14}{164.482}{331.541}{p}%
\TextAt{16}{170.465}{331.541}{\Character{79}}%
\TextAt{8}{183.187}{331.541}{over}%
\TextAt{13}{207.207}{331.541}{\Character{20}}%
\TextAt{8}{216.816}{331.541}{(we}%
\TextAt{8}{237.573}{331.541}{have}%
\TextAt{8}{263.541}{331.541}{seen}%
\TextAt{8}{287.932}{331.541}{in}%
\TextAt{8}{300.559}{331.541}{the}%
\TextAt{8}{319.689}{331.541}{proof}%
\TextAt{8}{349.232}{331.541}{of}%
\TextAt{8}{361.536}{331.541}{(8.1)}%
\TextAt{8}{388.485}{331.541}{that}%
\TextAt{16}{412.828}{331.541}{\Character{79}}%
\TextAt{8}{425.549}{331.541}{is}%
\TextAt{8}{436.287}{331.541}{a}%
\TextAt{8}{445.017}{331.541}{finitely}%
\TextAt{8}{483.335}{331.541}{generated}%
\TextAt{14}{79.226}{314.495}{o}%
\TextAt{8}{85.07}{314.495}{-module,}%
\TextAt{8}{134.425}{314.495}{so}%
\TextAt{8}{149.524}{314.495}{certainly}%
\TextAt{8}{198.708}{314.495}{dim}%
\TextAt{21}{218.219}{312.702}{\Character{20}}%
\TextAt{9}{223.597}{314.495}{\Character{40}}%
\TextAt{16}{228.149}{314.495}{\Character{79}}%
\TextAt{9}{237.999}{314.495}{\Character{47}}%
\TextAt{14}{243.852}{314.495}{p}%
\TextAt{16}{249.835}{314.495}{\Character{79}}%
\TextAt{9}{259.685}{314.495}{\Character{41}}%
\TextAt{13}{268.801}{314.495}{\Character{60}}%
\NativeMathAt{282.469}{314.495}{11.955}{\infty}%
\TextAt{8}{294.425}{314.495}{).}%
\TextAt{8}{309.616}{314.495}{It}%
\TextAt{8}{323.018}{314.495}{is}%
\TextAt{8}{335.511}{314.495}{sufficient}%
\TextAt{8}{385.4}{314.495}{to}%
\TextAt{8}{400.44}{314.495}{show}%
\TextAt{8}{430.173}{314.495}{that}%
\TextAt{13}{456.271}{314.495}{\Character{33}}%
\TextAt{17}{463.56}{312.702}{1}%
\TextAt{13}{468.293}{314.495}{\Character{59}}%
\TextAt{13}{473.517}{314.495}{\Character{58}}%
\TextAt{13}{478.754}{314.495}{\Character{58}}%
\TextAt{13}{483.978}{314.495}{\Character{58}}%
\TextAt{13}{489.215}{314.495}{\Character{59}}%
\TextAt{13}{494.451}{314.495}{\Character{33}}%
\TextAt{21}{501.803}{312.702}{m}%
\TextAt{8}{514.418}{314.495}{is}%
\TextAt{8}{526.912}{314.495}{a}%
\TextAt{8}{79.226}{297.449}{basis}%
\TextAt{8}{107.388}{297.449}{of}%
\TextAt{13}{120.139}{297.449}{L}%
\TextAt{16}{131.424}{297.449}{\Character{106}}%
\TextAt{13}{138.066}{297.449}{K}%
\TextAt{8}{148.884}{297.449}{.}%
\TextAt{8}{157.139}{297.449}{Assume}%
\TextAt{8}{199.922}{297.449}{the}%
\TextAt{13}{219.5}{297.449}{\Character{33}}%
\TextAt{17}{226.789}{295.656}{1}%
\TextAt{13}{231.521}{297.449}{\Character{59}}%
\TextAt{13}{236.745}{297.449}{\Character{58}}%
\TextAt{13}{241.982}{297.449}{\Character{58}}%
\TextAt{13}{247.218}{297.449}{\Character{58}}%
\TextAt{13}{252.443}{297.449}{\Character{59}}%
\TextAt{13}{257.679}{297.449}{\Character{33}}%
\TextAt{21}{265.032}{295.656}{m}%
\TextAt{8}{276.339}{297.449}{are}%
\TextAt{8}{295.271}{297.449}{linearly}%
\TextAt{8}{336.643}{297.449}{dependent}%
\TextAt{8}{392.637}{297.449}{over}%
\TextAt{13}{417.104}{297.449}{K}%
\TextAt{8}{427.922}{297.449}{,}%
\TextAt{8}{434.608}{297.449}{and}%
\TextAt{8}{456.792}{297.449}{hence}%
\TextAt{8}{488.719}{297.449}{also}%
\TextAt{8}{511.621}{297.449}{over}%
\TextAt{14}{79.226}{280.403}{o}%
\TextAt{8}{85.07}{280.403}{.}%
\TextAt{8}{93.519}{280.403}{Then}%
\TextAt{8}{124.076}{280.403}{there}%
\TextAt{8}{153.988}{280.403}{are}%
\TextAt{8}{173.499}{280.403}{elements}%
\TextAt{13}{221.355}{280.403}{a}%
\TextAt{17}{227.5}{278.61}{1}%
\TextAt{13}{232.233}{280.403}{\Character{59}}%
\TextAt{13}{237.457}{280.403}{\Character{58}}%
\TextAt{13}{242.694}{280.403}{\Character{58}}%
\TextAt{13}{247.918}{280.403}{\Character{58}}%
\TextAt{13}{253.155}{280.403}{\Character{59}}%
\TextAt{13}{258.391}{280.403}{a}%
\TextAt{21}{264.598}{278.61}{m}%
\NativeMathAt{275.908}{280.403}{11.955}{\in}%
\TextAt{14}{287.199}{280.403}{o}%
\TextAt{8}{296.94}{280.403}{not}%
\TextAt{8}{317.754}{280.403}{all}%
\TextAt{8}{334.013}{280.403}{zero}%
\TextAt{8}{358.725}{280.403}{such}%
\TextAt{8}{385.122}{280.403}{that}%
\TextAt{13}{246.025}{251.402}{a}%
\TextAt{17}{252.17}{249.609}{1}%
\TextAt{13}{256.903}{251.402}{\Character{33}}%
\TextAt{17}{264.192}{249.609}{1}%
\TextAt{9}{271.581}{251.402}{\Character{43}}%
\TextAt{16}{283.342}{251.402}{\Character{1}}%
\TextAt{16}{288.638}{251.402}{\Character{1}}%
\TextAt{16}{293.946}{251.402}{\Character{1}}%
\TextAt{9}{299.947}{251.402}{\Character{43}}%
\TextAt{13}{311.708}{251.402}{a}%
\TextAt{21}{317.853}{249.609}{m}%
\TextAt{13}{325.841}{251.402}{\Character{33}}%
\TextAt{21}{333.131}{249.609}{m}%
\TextAt{9}{344.44}{251.402}{\Character{61}}%
\TextAt{9}{356.85}{251.402}{0}%
\TextAt{13}{362.719}{251.402}{\Character{58}}%
\TextAt{8}{79.226}{222.401}{Consider}%
\TextAt{8}{129.303}{222.401}{the}%
\TextAt{8}{150.699}{222.401}{ideal}%
\TextAt{14}{179.901}{222.401}{a}%
\TextAt{9}{191.318}{222.401}{\Character{61}}%
\TextAt{9}{205.844}{222.401}{\Character{40}}%
\TextAt{13}{210.408}{222.401}{a}%
\TextAt{17}{216.553}{220.608}{1}%
\TextAt{13}{221.285}{222.401}{\Character{59}}%
\TextAt{13}{226.509}{222.401}{\Character{58}}%
\TextAt{13}{231.746}{222.401}{\Character{58}}%
\TextAt{13}{236.982}{222.401}{\Character{58}}%
\TextAt{13}{242.207}{222.401}{\Character{59}}%
\TextAt{13}{247.443}{222.401}{a}%
\TextAt{21}{253.651}{220.608}{m}%
\TextAt{9}{261.64}{222.401}{\Character{41}}%
\TextAt{8}{271.329}{222.401}{of}%
\TextAt{14}{285.899}{222.401}{o}%
\TextAt{8}{296.88}{222.401}{and}%
\TextAt{8}{320.882}{222.401}{find}%
\TextAt{13}{345.53}{222.401}{a}%
\NativeMathAt{357.108}{222.401}{11.955}{\in}%
\TextAt{14}{370.512}{222.401}{a}%
\TextAt{3}{376.495}{226.74}{\Character{0}}%
\TextAt{17}{383.081}{226.74}{1}%
\TextAt{8}{392.95}{222.401}{such}%
\TextAt{8}{420.587}{222.401}{that}%
\TextAt{13}{447.196}{222.401}{a}%
\NativeMathAt{458.774}{222.401}{11.955}{\notin}%
\TextAt{14}{472.177}{222.401}{a}%
\TextAt{3}{478.16}{226.74}{\Character{0}}%
\TextAt{17}{484.747}{226.74}{1}%
\TextAt{14}{489.479}{222.401}{p}%
\TextAt{8}{495.462}{222.401}{,}%
\TextAt{8}{504.161}{222.401}{hence}%
\TextAt{13}{79.226}{205.355}{a}%
\TextAt{14}{85.371}{205.355}{a}%
\TextAt{18}{94.675}{205.355}{\Character{42}}%
\TextAt{14}{107.294}{205.355}{p}%
\TextAt{8}{113.277}{205.355}{.}%
\TextAt{8}{121.664}{205.355}{Then}%
\TextAt{8}{152.036}{205.355}{the}%
\TextAt{8}{172.006}{205.355}{elements}%
\TextAt{13}{219.677}{205.355}{aa}%
\TextAt{17}{231.967}{203.562}{1}%
\TextAt{13}{236.699}{205.355}{\Character{59}}%
\TextAt{13}{241.923}{205.355}{\Character{58}}%
\TextAt{13}{247.16}{205.355}{\Character{58}}%
\TextAt{13}{252.396}{205.355}{\Character{58}}%
\TextAt{13}{257.621}{205.355}{\Character{59}}%
\TextAt{13}{262.857}{205.355}{aa}%
\TextAt{21}{275.21}{203.562}{m}%
\TextAt{8}{286.91}{205.355}{lie}%
\TextAt{8}{302.326}{205.355}{in}%
\TextAt{14}{315.793}{205.355}{o}%
\TextAt{8}{321.637}{205.355}{,}%
\TextAt{8}{328.638}{205.355}{but}%
\TextAt{8}{349.912}{205.355}{not}%
\TextAt{8}{370.54}{205.355}{all}%
\TextAt{8}{386.613}{205.355}{belong}%
\TextAt{8}{423.823}{205.355}{to}%
\TextAt{14}{437.948}{205.355}{p}%
\TextAt{8}{443.931}{205.355}{.}%
\TextAt{8}{452.318}{205.355}{The}%
\TextAt{8}{476.186}{205.355}{congruence}%
\TextAt{13}{214.837}{176.354}{aa}%
\TextAt{17}{227.127}{174.561}{1}%
\TextAt{13}{231.859}{176.354}{\Character{33}}%
\TextAt{17}{239.148}{174.561}{1}%
\TextAt{9}{246.537}{176.354}{\Character{43}}%
\TextAt{16}{258.299}{176.354}{\Character{1}}%
\TextAt{16}{263.595}{176.354}{\Character{1}}%
\TextAt{16}{268.903}{176.354}{\Character{1}}%
\TextAt{9}{274.903}{176.354}{\Character{43}}%
\TextAt{13}{286.664}{176.354}{aa}%
\TextAt{21}{298.954}{174.561}{m}%
\TextAt{13}{306.943}{176.354}{\Character{33}}%
\TextAt{21}{314.232}{174.561}{m}%
\NativeMathAt{325.542}{176.354}{11.955}{\equiv}%
\TextAt{9}{338.161}{176.354}{0}%
\TextAt{9}{355.95}{176.354}{\Character{40}}%
\TextAt{8}{360.521}{176.354}{mod}%
\TextAt{14}{386.623}{176.354}{p}%
\TextAt{9}{392.606}{176.354}{\Character{41}}%
\TextAt{8}{79.226}{147.353}{thus}%
\TextAt{8}{106.177}{147.353}{gives}%
\TextAt{8}{136.057}{147.353}{us}%
\TextAt{8}{152.272}{147.353}{a}%
\TextAt{8}{163.227}{147.353}{linear}%
\TextAt{8}{196.945}{147.353}{dependence}%
\TextAt{8}{260.885}{147.353}{among}%
\TextAt{8}{299.815}{147.353}{the}%
\Vector{q 0 G 0 g 0.478 w 249.171 -619.826 m 256.89 -619.826 l S Q}%
\TextAt{13}{321.171}{147.353}{\Character{33}}%
\TextAt{17}{328.89}{145.56}{1}%
\TextAt{13}{333.622}{147.353}{\Character{59}}%
\TextAt{13}{338.846}{147.353}{\Character{58}}%
\TextAt{13}{344.083}{147.353}{\Character{58}}%
\TextAt{13}{349.307}{147.353}{\Character{58}}%
\TextAt{13}{354.544}{147.353}{\Character{59}}%
\Vector{q 0 G 0 g 0.478 w 287.843 -619.826 m 295.561 -619.826 l S Q}%
\TextAt{13}{359.843}{147.353}{\Character{33}}%
\TextAt{21}{367.561}{145.56}{m}%
\TextAt{8}{380.647}{147.353}{over}%
\TextAt{13}{406.892}{147.353}{\Character{20}}%
\TextAt{8}{413.631}{147.353}{,}%
\TextAt{8}{422.279}{147.353}{a}%
\TextAt{8}{433.234}{147.353}{contradiction.}%
\TextAt{8}{512.613}{147.353}{The}%
\TextAt{13}{79.226}{130.307}{\Character{33}}%
\TextAt{17}{86.515}{128.514}{1}%
\TextAt{13}{91.248}{130.307}{\Character{59}}%
\TextAt{13}{96.472}{130.307}{\Character{58}}%
\TextAt{13}{101.709}{130.307}{\Character{58}}%
\TextAt{13}{106.933}{130.307}{\Character{58}}%
\TextAt{13}{112.17}{130.307}{\Character{59}}%
\TextAt{13}{117.406}{130.307}{\Character{33}}%
\TextAt{21}{124.758}{128.514}{m}%
\TextAt{8}{136.644}{130.307}{are}%
\TextAt{8}{156.155}{130.307}{therefore}%
\TextAt{8}{205.255}{130.307}{linearly}%
\TextAt{8}{247.205}{130.307}{independent}%
\TextAt{8}{313.532}{130.307}{over}%
\TextAt{13}{338.579}{130.307}{K}%
\TextAt{8}{349.396}{130.307}{.}%
\TextAt{8}{97.159}{110.272}{In}%
\TextAt{8}{112.537}{110.272}{order}%
\TextAt{8}{143.863}{110.272}{to}%
\TextAt{8}{158.93}{110.272}{show}%
\TextAt{8}{188.69}{110.272}{that}%
\TextAt{8}{214.815}{110.272}{the}%
\TextAt{13}{235.729}{110.272}{\Character{33}}%
\TextAt{21}{243.018}{108.479}{i}%
\TextAt{8}{251.053}{110.272}{are}%
\TextAt{8}{271.321}{110.272}{a}%
\TextAt{8}{281.833}{110.272}{basis}%
\TextAt{8}{311.33}{110.272}{of}%
\TextAt{13}{325.417}{110.272}{L}%
\TextAt{16}{337.992}{110.272}{\Character{106}}%
\TextAt{13}{345.923}{110.272}{K}%
\TextAt{8}{356.74}{110.272}{,}%
\TextAt{8}{364.836}{110.272}{we}%
\TextAt{8}{382.82}{110.272}{consider}%
\TextAt{8}{429.162}{110.272}{the}%
\TextAt{14}{450.075}{110.272}{o}%
\TextAt{8}{455.919}{110.272}{-modules}%
\TextAt{13}{506.481}{110.272}{M}%
\TextAt{9}{523.665}{110.272}{\Character{61}}%
\TextAt{14}{79.226}{93.226}{o}%
\TextAt{13}{85.07}{93.226}{\Character{33}}%
\TextAt{17}{92.36}{91.433}{1}%
\TextAt{9}{100.227}{93.226}{\Character{43}}%
\TextAt{16}{112.467}{93.226}{\Character{1}}%
\TextAt{16}{117.763}{93.226}{\Character{1}}%
\TextAt{16}{123.071}{93.226}{\Character{1}}%
\TextAt{9}{129.55}{93.226}{\Character{43}}%
\TextAt{14}{141.79}{93.226}{o}%
\TextAt{13}{147.634}{93.226}{\Character{33}}%
\TextAt{21}{154.924}{91.433}{m}%
\TextAt{8}{167.512}{93.226}{and}%
\TextAt{13}{190.977}{93.226}{N}%
\TextAt{9}{206.117}{93.226}{\Character{61}}%
\TextAt{16}{219.739}{93.226}{\Character{79}}%
\TextAt{9}{229.588}{93.226}{\Character{47}}%
\TextAt{13}{235.441}{93.226}{M}%
\TextAt{8}{248.015}{93.226}{.}%
\TextAt{8}{258.57}{93.226}{Since}%
\TextAt{16}{289.829}{93.226}{\Character{79}}%
\TextAt{9}{304.197}{93.226}{\Character{61}}%
\TextAt{13}{317.819}{93.226}{M}%
\TextAt{9}{333.528}{93.226}{\Character{43}}%
\TextAt{14}{345.767}{93.226}{p}%
\TextAt{16}{351.751}{93.226}{\Character{79}}%
\TextAt{8}{361.6}{93.226}{,}%
\TextAt{8}{369.627}{93.226}{we}%
\TextAt{8}{387.557}{93.226}{have}%
\TextAt{14}{415.254}{93.226}{p}%
\TextAt{13}{421.237}{93.226}{N}%
\TextAt{9}{436.377}{93.226}{\Character{61}}%
\TextAt{13}{449.999}{93.226}{N}%
\TextAt{8}{460.622}{93.226}{.}%
\TextAt{8}{471.176}{93.226}{As}%
\TextAt{13}{489.166}{93.226}{L}%
\TextAt{16}{501.648}{93.226}{\Character{106}}%
\TextAt{13}{509.486}{93.226}{K}%
\TextAt{8}{524.903}{93.226}{is}%
\TextAt{8}{300.14}{43.033}{43}%
\Destination{78.23}{809.815}{nav.page.43}%
\Destination{79.36}{626.455}{nav.print.46}%
\Destination{79.36}{582.355}{nav.number.1.8.2}%
\Hotlink{361.036}{329.355}{25.078}{11.047}{nav.number.1.8.1}%
}

% Page 56 (body page 44)
\Page{%
\TextAt{2}{79.226}{796.943}{I.}%
\TextAt{2}{89.83}{796.943}{Algebraic}%
\TextAt{2}{138.648}{796.943}{Integers}%
\Vector{q 0 G 0 g 0.299 w 7.226 18.729 m 460.769 18.729 l S Q}%
\TextAt{8}{79.226}{761.124}{separable,}%
\TextAt{16}{133.525}{761.124}{\Character{79}}%
\TextAt{8}{143.375}{761.124}{,}%
\TextAt{8}{150.129}{761.124}{and}%
\TextAt{8}{172.397}{761.124}{hence}%
\TextAt{8}{204.408}{761.124}{also}%
\TextAt{13}{227.393}{761.124}{N}%
\TextAt{8}{238.016}{761.124}{,}%
\TextAt{8}{244.769}{761.124}{are}%
\TextAt{8}{263.786}{761.124}{finitely}%
\TextAt{8}{302.635}{761.124}{generated}%
\TextAt{14}{355.473}{761.124}{o}%
\TextAt{8}{361.317}{761.124}{-modules}%
\TextAt{8}{410.628}{761.124}{(see}%
\TextAt{8}{433.601}{761.124}{p.}%
\TextAt{8}{446.759}{761.124}{45).}%
\TextAt{8}{471.314}{761.124}{If}%
\TextAt{13}{482.511}{761.124}{\Character{11}}%
\TextAt{17}{489.957}{759.33}{1}%
\TextAt{13}{494.689}{761.124}{\Character{59}}%
\TextAt{13}{499.913}{761.124}{\Character{58}}%
\TextAt{13}{505.15}{761.124}{\Character{58}}%
\TextAt{13}{510.374}{761.124}{\Character{58}}%
\TextAt{13}{515.611}{761.124}{\Character{59}}%
\TextAt{13}{520.847}{761.124}{\Character{11}}%
\TextAt{21}{528.355}{759.33}{s}%
\TextAt{8}{79.226}{744.078}{is}%
\TextAt{8}{90.99}{744.078}{a}%
\TextAt{8}{100.745}{744.078}{system}%
\TextAt{8}{139.564}{744.078}{of}%
\TextAt{8}{152.894}{744.078}{generators}%
\TextAt{8}{209.549}{744.078}{of}%
\TextAt{13}{222.879}{744.078}{N}%
\TextAt{8}{233.502}{744.078}{,}%
\TextAt{8}{240.651}{744.078}{then}%
\TextAt{13}{230.936}{715.077}{\Character{11}}%
\TextAt{21}{238.381}{713.283}{i}%
\TextAt{9}{245.083}{715.077}{\Character{61}}%
\TextAt{20}{257.509}{726.434}{\Character{88}}%
\TextAt{21}{264.201}{701.24}{j}%
\TextAt{13}{276.77}{715.077}{a}%
\TextAt{21}{282.915}{713.283}{ij}%
\TextAt{13}{290.18}{715.077}{\Character{11}}%
\TextAt{21}{297.626}{713.283}{j}%
\TextAt{8}{325.918}{715.077}{for}%
\TextAt{13}{343.803}{715.077}{a}%
\TextAt{21}{349.948}{713.283}{ij}%
\NativeMathAt{360.534}{715.077}{11.955}{\in}%
\TextAt{14}{371.825}{715.077}{p}%
\TextAt{13}{377.808}{715.077}{\Character{58}}%
\TextAt{8}{79.226}{676.577}{Let}%
\TextAt{13}{101.09}{676.577}{A}%
\TextAt{8}{114.657}{676.577}{be}%
\TextAt{8}{131.476}{676.577}{the}%
\TextAt{8}{152.527}{676.577}{matrix}%
\TextAt{9}{191.475}{676.577}{\Character{40}}%
\TextAt{13}{196.027}{676.577}{a}%
\TextAt{21}{202.172}{674.783}{ij}%
\TextAt{9}{209.437}{676.577}{\Character{41}}%
\TextAt{16}{217.256}{676.577}{\Character{0}}%
\TextAt{13}{229.821}{676.577}{I}%
\TextAt{8}{235.924}{676.577}{,}%
\TextAt{8}{244.191}{676.577}{where}%
\TextAt{13}{278.895}{676.577}{I}%
\TextAt{8}{289.789}{676.577}{is}%
\TextAt{8}{302.448}{676.577}{the}%
\TextAt{8}{323.499}{676.577}{unit}%
\TextAt{8}{349.104}{676.577}{matrix}%
\TextAt{8}{388.052}{676.577}{of}%
\TextAt{8}{402.277}{676.577}{rank}%
\TextAt{13}{430.166}{676.577}{s}%
\TextAt{8}{435.68}{676.577}{,}%
\TextAt{8}{443.947}{676.577}{and}%
\TextAt{8}{467.604}{676.577}{let}%
\TextAt{13}{485.404}{676.577}{B}%
\TextAt{8}{499.692}{676.577}{be}%
\TextAt{8}{516.511}{676.577}{the}%
\ErrataAt{79.19}{659.531}{455.81}{adjugate matrix of $A$, the transpose of its cofactor matrix. Then one has}%
\TextAt{13}{79.226}{642.485}{A}%
\TextAt{9}{88.002}{642.485}{\Character{40}}%
\TextAt{13}{92.554}{642.485}{\Character{11}}%
\TextAt{17}{99.999}{640.691}{1}%
\TextAt{13}{104.732}{642.485}{\Character{59}}%
\TextAt{13}{109.956}{642.485}{\Character{58}}%
\TextAt{13}{115.193}{642.485}{\Character{58}}%
\TextAt{13}{120.417}{642.485}{\Character{58}}%
\TextAt{13}{125.654}{642.485}{\Character{59}}%
\TextAt{13}{130.89}{642.485}{\Character{11}}%
\TextAt{21}{138.398}{640.691}{s}%
\TextAt{9}{142.812}{642.485}{\Character{41}}%
\TextAt{21}{147.365}{646.823}{t}%
\TextAt{9}{154.241}{642.485}{\Character{61}}%
\TextAt{9}{166.651}{642.485}{0}%
\TextAt{8}{176.417}{642.485}{and}%
\TextAt{13}{199.18}{642.485}{B}%
\TextAt{13}{208.661}{642.485}{A}%
\TextAt{9}{220.772}{642.485}{\Character{61}}%
\TextAt{13}{233.198}{642.485}{dI}%
\TextAt{8}{245.383}{642.485}{,}%
\TextAt{8}{252.533}{642.485}{with}%
\TextAt{13}{279.193}{642.485}{d}%
\TextAt{9}{288.596}{642.485}{\Character{61}}%
\TextAt{8}{301.022}{642.485}{det}%
\TextAt{9}{317.281}{642.485}{\Character{40}}%
\TextAt{13}{321.833}{642.485}{A}%
\TextAt{9}{330.608}{642.485}{\Character{41}}%
\TextAt{8}{339.058}{642.485}{(see}%
\TextAt{8}{362.526}{642.485}{(2.3)).}%
\TextAt{8}{399.607}{642.485}{Hence}%
\TextAt{9}{207.562}{613.484}{0}%
\TextAt{9}{216.72}{613.484}{\Character{61}}%
\TextAt{13}{229.161}{613.484}{B}%
\TextAt{13}{238.642}{613.484}{A}%
\TextAt{9}{247.432}{613.484}{\Character{40}}%
\TextAt{13}{251.985}{613.484}{\Character{11}}%
\TextAt{17}{259.43}{611.69}{1}%
\TextAt{13}{264.163}{613.484}{\Character{59}}%
\TextAt{13}{269.387}{613.484}{\Character{58}}%
\TextAt{13}{274.624}{613.484}{\Character{58}}%
\TextAt{13}{279.848}{613.484}{\Character{58}}%
\TextAt{13}{285.085}{613.484}{\Character{59}}%
\TextAt{13}{290.321}{613.484}{\Character{11}}%
\TextAt{21}{297.829}{611.69}{s}%
\TextAt{9}{302.243}{613.484}{\Character{41}}%
\TextAt{21}{306.795}{618.42}{t}%
\TextAt{9}{313.672}{613.484}{\Character{61}}%
\TextAt{9}{326.082}{613.484}{\Character{40}}%
\TextAt{13}{330.65}{613.484}{d\Character{11}}%
\TextAt{17}{344.179}{611.69}{1}%
\TextAt{13}{348.911}{613.484}{\Character{59}}%
\TextAt{13}{354.135}{613.484}{\Character{58}}%
\TextAt{13}{359.372}{613.484}{\Character{58}}%
\TextAt{13}{364.596}{613.484}{\Character{58}}%
\TextAt{13}{369.833}{613.484}{\Character{59}}%
\TextAt{13}{375.069}{613.484}{d\Character{11}}%
\TextAt{21}{388.66}{611.69}{s}%
\TextAt{9}{393.074}{613.484}{\Character{41}}%
\TextAt{21}{397.626}{618.42}{t}%
\TextAt{13}{401.182}{613.484}{\Character{59}}%
\TextAt{8}{79.226}{584.482}{and}%
\TextAt{8}{101.114}{584.482}{therefore}%
\TextAt{13}{149.338}{584.482}{dN}%
\TextAt{9}{169.365}{584.482}{\Character{61}}%
\TextAt{9}{181.775}{584.482}{0}%
\TextAt{8}{187.643}{584.482}{,}%
\TextAt{8}{194.092}{584.482}{i.e.,}%
\TextAt{13}{215.497}{584.482}{d}%
\TextAt{16}{221.58}{584.482}{\Character{79}}%
\NativeMathAt{234.743}{584.482}{11.955}{\subseteq}%
\TextAt{13}{247.37}{584.482}{M}%
\TextAt{9}{263.264}{584.482}{\Character{61}}%
\TextAt{14}{275.69}{584.482}{o}%
\TextAt{13}{281.534}{584.482}{\Character{33}}%
\TextAt{17}{288.823}{582.689}{1}%
\TextAt{9}{294.423}{584.482}{\Character{43}}%
\TextAt{16}{304.394}{584.482}{\Character{1}}%
\TextAt{16}{309.69}{584.482}{\Character{1}}%
\TextAt{16}{314.998}{584.482}{\Character{1}}%
\TextAt{9}{319.209}{584.482}{\Character{43}}%
\TextAt{14}{329.181}{584.482}{o}%
\TextAt{13}{335.025}{584.482}{\Character{33}}%
\TextAt{21}{342.314}{582.689}{m}%
\TextAt{8}{350.303}{584.482}{.}%
\TextAt{8}{358.459}{584.482}{We}%
\TextAt{8}{377.729}{584.482}{have}%
\TextAt{13}{403.848}{584.482}{d}%
\NativeMathAt{413.252}{584.482}{11.955}{\neq}%
\TextAt{9}{425.661}{584.482}{0}%
\TextAt{8}{431.53}{584.482}{,}%
\TextAt{8}{437.979}{584.482}{because}%
\TextAt{8}{480.406}{584.482}{expanding}%
\TextAt{8}{79.226}{567.436}{the}%
\TextAt{8}{99.25}{567.436}{determinant}%
\TextAt{13}{165.134}{567.436}{d}%
\TextAt{9}{174.538}{567.436}{\Character{61}}%
\TextAt{8}{186.963}{567.436}{det}%
\TextAt{9}{203.222}{567.436}{\Character{40}\Character{40}}%
\TextAt{13}{212.327}{567.436}{a}%
\TextAt{21}{218.472}{565.643}{ij}%
\TextAt{9}{225.737}{567.436}{\Character{41}}%
\TextAt{16}{232.676}{567.436}{\Character{0}}%
\TextAt{13}{244.36}{567.436}{I}%
\TextAt{9}{250.463}{567.436}{\Character{41}}%
\TextAt{8}{258.78}{567.436}{we}%
\TextAt{8}{275.875}{567.436}{find}%
\TextAt{13}{299.151}{567.436}{d}%
\NativeMathAt{308.555}{567.436}{11.955}{\equiv}%
\TextAt{9}{321.174}{567.436}{\Character{40}}%
\TextAt{16}{325.726}{567.436}{\Character{0}}%
\TextAt{9}{335.025}{567.436}{1\Character{41}}%
\TextAt{21}{345.43}{571.775}{s}%
\TextAt{9}{355.157}{567.436}{\Character{40}}%
\TextAt{8}{359.71}{567.436}{mod}%
\TextAt{14}{385.812}{567.436}{p}%
\TextAt{9}{391.795}{567.436}{\Character{41}}%
\TextAt{8}{400.112}{567.436}{because}%
\TextAt{13}{443.282}{567.436}{a}%
\TextAt{21}{449.427}{565.643}{ij}%
\NativeMathAt{460.013}{567.436}{11.955}{\in}%
\TextAt{14}{471.304}{567.436}{p}%
\TextAt{8}{477.287}{567.436}{.}%
\TextAt{8}{485.691}{567.436}{It}%
\TextAt{8}{498.231}{567.436}{follows}%
\TextAt{8}{79.226}{550.39}{that}%
\TextAt{13}{104.595}{550.39}{L}%
\TextAt{9}{115.88}{550.39}{\Character{61}}%
\TextAt{13}{128.306}{550.39}{dL}%
\TextAt{9}{145.674}{550.39}{\Character{61}}%
\TextAt{13}{158.099}{550.39}{K}%
\TextAt{13}{168.907}{550.39}{\Character{33}}%
\TextAt{17}{176.206}{548.597}{1}%
\TextAt{9}{183.595}{550.39}{\Character{43}}%
\TextAt{16}{195.357}{550.39}{\Character{1}}%
\TextAt{16}{200.653}{550.39}{\Character{1}}%
\TextAt{16}{205.961}{550.39}{\Character{1}}%
\TextAt{9}{211.961}{550.39}{\Character{43}}%
\TextAt{13}{223.722}{550.39}{K}%
\TextAt{13}{234.53}{550.39}{\Character{33}}%
\TextAt{21}{241.829}{548.597}{m}%
\TextAt{8}{249.818}{550.39}{.}%
\TextAt{13}{258.266}{550.39}{\Character{33}}%
\TextAt{17}{265.555}{548.597}{1}%
\TextAt{13}{270.288}{550.39}{\Character{59}}%
\TextAt{13}{275.512}{550.39}{\Character{58}}%
\TextAt{13}{280.749}{550.39}{\Character{58}}%
\TextAt{13}{285.973}{550.39}{\Character{58}}%
\TextAt{13}{291.21}{550.39}{\Character{59}}%
\TextAt{13}{296.446}{550.39}{\Character{33}}%
\TextAt{21}{303.798}{548.597}{m}%
\TextAt{8}{315.684}{550.39}{is}%
\TextAt{8}{327.448}{550.39}{therefore}%
\TextAt{8}{376.548}{550.39}{indeed}%
\TextAt{8}{413.609}{550.39}{a}%
\TextAt{8}{423.364}{550.39}{basis}%
\TextAt{8}{452.104}{550.39}{of}%
\TextAt{13}{465.434}{550.39}{L}%
\TextAt{16}{476.72}{550.39}{\Character{106}}%
\TextAt{13}{483.361}{550.39}{K}%
\TextAt{8}{494.179}{550.39}{.}%
\TextAt{8}{97.159}{530.356}{In}%
\TextAt{8}{111.78}{530.356}{order}%
\TextAt{8}{142.349}{530.356}{to}%
\TextAt{8}{156.66}{530.356}{prove}%
\TextAt{8}{188.209}{530.356}{the}%
\TextAt{8}{208.366}{530.356}{second}%
\TextAt{8}{246.144}{530.356}{identity,}%
\TextAt{8}{291.67}{530.356}{let}%
\TextAt{8}{308.574}{530.356}{us}%
\TextAt{8}{323.59}{530.356}{consider}%
\TextAt{8}{369.175}{530.356}{the}%
\TextAt{8}{389.331}{530.356}{descending}%
\TextAt{8}{448.569}{530.356}{chain}%
\TextAt{16}{198.95}{501.355}{\Character{79}}%
\TextAt{9}{208.8}{501.355}{\Character{47}}%
\TextAt{14}{214.653}{501.355}{P}%
\TextAt{21}{224.546}{506.687}{e}%
\TextAt{22}{228.474}{505.472}{i}%
\TextAt{21}{224.546}{498.291}{i}%
\NativeMathAt{235.455}{501.355}{11.955}{\supseteq}%
\TextAt{14}{248.074}{501.355}{P}%
\TextAt{21}{257.968}{499.561}{i}%
\TextAt{9}{261.349}{501.355}{\Character{47}}%
\TextAt{14}{267.202}{501.355}{P}%
\TextAt{21}{277.095}{506.688}{e}%
\TextAt{22}{281.023}{505.472}{i}%
\TextAt{21}{277.095}{498.291}{i}%
\NativeMathAt{288.004}{501.355}{11.955}{\supseteq}%
\TextAt{16}{300.605}{501.355}{\Character{1}}%
\TextAt{16}{305.913}{501.355}{\Character{1}}%
\TextAt{16}{311.221}{501.355}{\Character{1}}%
\NativeMathAt{317.844}{501.355}{11.955}{\supseteq}%
\TextAt{14}{330.511}{501.355}{P}%
\TextAt{21}{340.405}{506.688}{e}%
\TextAt{22}{344.333}{505.472}{i}%
\TextAt{3}{347.495}{506.688}{\Character{0}}%
\TextAt{17}{354.081}{506.688}{1}%
\TextAt{21}{340.405}{498.291}{i}%
\TextAt{9}{358.814}{501.355}{\Character{47}}%
\TextAt{14}{364.667}{501.355}{P}%
\TextAt{21}{374.56}{506.688}{e}%
\TextAt{22}{378.488}{505.472}{i}%
\TextAt{21}{374.56}{498.291}{i}%
\NativeMathAt{385.469}{501.355}{11.955}{\supseteq}%
\TextAt{9}{398.088}{501.355}{\Character{40}0\Character{41}}%
\TextAt{8}{79.226}{472.353}{of}%
\TextAt{13}{93.823}{472.353}{\Character{20}}%
\TextAt{8}{100.562}{472.353}{-vector}%
\TextAt{8}{140.851}{472.353}{spaces.}%
\TextAt{8}{185.093}{472.353}{The}%
\TextAt{8}{210.414}{472.353}{successive}%
\TextAt{8}{265.838}{472.353}{quotients}%
\TextAt{14}{317.903}{472.353}{P}%
\TextAt{21}{327.796}{476.692}{\Character{23}}%
\TextAt{21}{327.796}{469.398}{i}%
\TextAt{9}{332.978}{472.353}{\Character{47}}%
\TextAt{14}{338.831}{472.353}{P}%
\TextAt{21}{348.725}{477.302}{\Character{23}}%
\TextAt{17}{353.409}{477.302}{\Character{43}1}%
\TextAt{21}{348.725}{469.29}{i}%
\TextAt{8}{369.892}{472.353}{in}%
\TextAt{8}{384.812}{472.353}{this}%
\TextAt{8}{408.902}{472.353}{chain}%
\TextAt{8}{441.061}{472.353}{are}%
\TextAt{8}{461.839}{472.353}{isomorphic}%
\TextAt{8}{522.357}{472.353}{to}%
\TextAt{16}{79.226}{455.307}{\Character{79}}%
\TextAt{9}{89.076}{455.307}{\Character{47}}%
\TextAt{14}{94.929}{455.307}{P}%
\TextAt{21}{104.822}{453.514}{i}%
\TextAt{8}{108.203}{455.307}{,}%
\TextAt{8}{115.353}{455.307}{for}%
\TextAt{8}{133.238}{455.307}{if}%
\TextAt{13}{143.961}{455.307}{\Character{11}}%
\NativeMathAt{154.804}{455.307}{11.955}{\in}%
\TextAt{14}{166.095}{455.307}{P}%
\TextAt{21}{175.988}{459.646}{\Character{23}}%
\TextAt{21}{175.988}{452.352}{i}%
\TextAt{16}{183.827}{455.307}{\Character{110}}%
\TextAt{14}{192.461}{455.307}{P}%
\TextAt{21}{202.354}{460.256}{\Character{23}}%
\TextAt{17}{207.038}{460.256}{\Character{43}1}%
\TextAt{21}{202.354}{452.244}{i}%
\TextAt{8}{218.357}{455.307}{,}%
\TextAt{8}{225.506}{455.307}{then}%
\TextAt{8}{252.166}{455.307}{the}%
\TextAt{8}{272.323}{455.307}{homomorphism}%
\TextAt{16}{225.649}{426.306}{\Character{79}}%
\NativeArrowAt{238.812}{426.306}{11.955}{19.26}{\longrightarrow}%
\TextAt{14}{261.402}{426.306}{P}%
\TextAt{21}{271.295}{431.242}{\Character{23}}%
\TextAt{21}{271.295}{423.351}{i}%
\TextAt{9}{276.477}{426.306}{\Character{47}}%
\TextAt{14}{282.33}{426.306}{P}%
\TextAt{21}{292.223}{431.254}{\Character{23}}%
\TextAt{17}{296.907}{431.254}{\Character{43}1}%
\TextAt{21}{292.223}{423.242}{i}%
\TextAt{13}{308.226}{426.306}{\Character{59}}%
\TextAt{13}{337.361}{426.306}{a}%
\NativeArrowAt{346.846}{426.306}{11.955}{19.26}{\longmapsto}%
\TextAt{13}{369.428}{426.306}{a\Character{11}\Character{59}}%
\TextAt{8}{79.226}{397.305}{has}%
\TextAt{8}{101.289}{397.305}{kernel}%
\TextAt{14}{136.944}{397.305}{P}%
\TextAt{21}{146.837}{395.512}{i}%
\TextAt{8}{155.304}{397.305}{and}%
\TextAt{8}{179.256}{397.305}{is}%
\TextAt{8}{192.208}{397.305}{surjective}%
\TextAt{8}{245.808}{397.305}{because}%
\TextAt{14}{290.299}{397.305}{P}%
\TextAt{21}{300.192}{401.644}{\Character{23}}%
\TextAt{21}{300.192}{394.35}{i}%
\TextAt{8}{310.46}{397.305}{is}%
\TextAt{8}{323.413}{397.305}{the}%
\TextAt{8}{344.758}{397.305}{gcd}%
\TextAt{8}{367.406}{397.305}{of}%
\TextAt{14}{381.925}{397.305}{P}%
\TextAt{21}{391.818}{402.253}{\Character{23}}%
\TextAt{17}{396.502}{402.253}{\Character{43}1}%
\TextAt{21}{391.818}{394.241}{i}%
\TextAt{8}{412.906}{397.305}{and}%
\TextAt{9}{436.858}{397.305}{\Character{40}}%
\TextAt{13}{441.41}{397.305}{\Character{11}}%
\TextAt{9}{448.932}{397.305}{\Character{41}}%
\TextAt{9}{458.819}{397.305}{\Character{61}}%
\TextAt{13}{473.282}{397.305}{\Character{11}}%
\TextAt{16}{480.803}{397.305}{\Character{79}}%
\TextAt{8}{495.739}{397.305}{so}%
\TextAt{8}{511.298}{397.305}{that}%
\TextAt{14}{79.226}{380.259}{P}%
\TextAt{21}{89.119}{384.598}{\Character{23}}%
\TextAt{21}{89.119}{377.304}{i}%
\TextAt{9}{97.622}{380.259}{\Character{61}}%
\TextAt{13}{110.048}{380.259}{\Character{11}}%
\TextAt{16}{117.569}{380.259}{\Character{79}}%
\TextAt{9}{130.076}{380.259}{\Character{43}}%
\TextAt{14}{141.837}{380.259}{P}%
\TextAt{21}{151.73}{385.207}{\Character{23}}%
\TextAt{17}{156.414}{385.207}{\Character{43}1}%
\TextAt{21}{151.73}{377.195}{i}%
\TextAt{8}{167.733}{380.259}{.}%
\TextAt{8}{176.181}{380.259}{Since}%
\TextAt{13}{206.739}{380.259}{f}%
\TextAt{21}{212.509}{378.466}{i}%
\TextAt{9}{219.211}{380.259}{\Character{61}}%
\TextAt{9}{231.621}{380.259}{\Character{91}}%
\TextAt{16}{234.888}{380.259}{\Character{79}}%
\TextAt{9}{244.738}{380.259}{\Character{47}}%
\TextAt{14}{250.591}{380.259}{P}%
\TextAt{21}{260.484}{378.466}{i}%
\TextAt{9}{267.186}{380.259}{\Character{58}}%
\TextAt{13}{273.759}{380.259}{\Character{20}}%
\TextAt{9}{280.497}{380.259}{\Character{93}}%
\TextAt{8}{283.749}{380.259}{,}%
\TextAt{8}{290.898}{380.259}{we}%
\TextAt{8}{308.125}{380.259}{obtain}%
\TextAt{8}{344.553}{380.259}{dim}%
\TextAt{21}{364.064}{378.466}{\Character{20}}%
\TextAt{9}{369.442}{380.259}{\Character{40}}%
\TextAt{14}{373.994}{380.259}{P}%
\TextAt{21}{383.887}{384.598}{\Character{23}}%
\TextAt{21}{383.887}{377.304}{i}%
\TextAt{9}{389.069}{380.259}{\Character{47}}%
\TextAt{14}{394.922}{380.259}{P}%
\TextAt{21}{404.816}{385.207}{\Character{23}}%
\TextAt{17}{409.499}{385.207}{\Character{43}1}%
\TextAt{21}{404.816}{377.195}{i}%
\TextAt{9}{420.818}{380.259}{\Character{41}}%
\TextAt{9}{428.673}{380.259}{\Character{61}}%
\TextAt{13}{441.117}{380.259}{f}%
\TextAt{21}{446.887}{378.466}{i}%
\TextAt{8}{454.166}{380.259}{and}%
\TextAt{8}{476.928}{380.259}{therefore}%
\TextAt{8}{195.697}{342.816}{dim}%
\TextAt{21}{215.208}{341.022}{\Character{20}}%
\TextAt{9}{220.586}{342.816}{\Character{40}}%
\TextAt{16}{225.139}{342.816}{\Character{79}}%
\TextAt{9}{234.988}{342.816}{\Character{47}}%
\TextAt{14}{240.841}{342.816}{P}%
\TextAt{21}{250.735}{348.149}{e}%
\TextAt{22}{254.663}{346.934}{i}%
\TextAt{21}{250.734}{339.752}{i}%
\TextAt{9}{258.323}{342.816}{\Character{41}}%
\TextAt{9}{266.178}{342.816}{\Character{61}}%
\TextAt{21}{278.621}{357.912}{e}%
\TextAt{22}{282.55}{356.697}{i}%
\TextAt{3}{285.712}{357.912}{\Character{0}}%
\TextAt{17}{292.298}{357.912}{1}%
\TextAt{20}{278.943}{354.173}{\Character{88}}%
\TextAt{21}{279.825}{329.067}{\Character{23}}%
\TextAt{17}{284.508}{329.067}{\Character{61}0}%
\TextAt{8}{298.525}{342.816}{dim}%
\TextAt{21}{318.036}{341.022}{\Character{20}}%
\TextAt{9}{323.414}{342.816}{\Character{40}}%
\TextAt{14}{327.966}{342.816}{P}%
\TextAt{21}{337.859}{347.752}{\Character{23}}%
\TextAt{21}{337.859}{339.86}{i}%
\TextAt{9}{343.041}{342.816}{\Character{47}}%
\TextAt{14}{348.894}{342.816}{P}%
\TextAt{21}{358.787}{347.764}{\Character{23}}%
\TextAt{17}{363.471}{347.764}{\Character{43}1}%
\TextAt{21}{358.787}{339.752}{i}%
\TextAt{9}{374.79}{342.816}{\Character{41}}%
\TextAt{9}{382.645}{342.816}{\Character{61}}%
\TextAt{13}{395.089}{342.816}{e}%
\TextAt{21}{400.514}{341.022}{i}%
\TextAt{13}{403.895}{342.816}{f}%
\TextAt{21}{409.665}{341.022}{i}%
\TextAt{13}{413.047}{342.816}{\Character{58}}%
\Vector{q 0 G 0 g 0.399 w 452.334 -467.15 m 452.334 -459.08 l S Q}%
\Vector{q 0 G 0 g 0.399 w 452.534 -459.279 m 459.707 -459.279 l S Q}%
\Vector{q 0 G 0 g 0.399 w 452.534 -466.95 m 459.707 -466.95 l S Q}%
\Vector{q 0 G 0 g 0.399 w 459.906 -467.15 m 459.906 -459.08 l S Q}%
\TextAt{8}{97.159}{277.849}{Suppose}%
\TextAt{8}{142.453}{277.849}{now}%
\TextAt{8}{166.229}{277.849}{that}%
\TextAt{8}{190.984}{277.849}{the}%
\TextAt{8}{210.527}{277.849}{separable}%
\TextAt{8}{261.357}{277.849}{extension}%
\TextAt{13}{312.509}{277.849}{L}%
\TextAt{16}{323.795}{277.849}{\Character{106}}%
\TextAt{13}{330.436}{277.849}{K}%
\TextAt{8}{344.538}{277.849}{is}%
\TextAt{8}{355.689}{277.849}{given}%
\TextAt{8}{385.645}{277.849}{by}%
\TextAt{8}{401.29}{277.849}{a}%
\TextAt{8}{410.432}{277.849}{primitive}%
\TextAt{8}{459.899}{277.849}{element}%
\TextAt{13}{502.528}{277.849}{\Character{18}}%
\NativeMathAt{511.629}{277.849}{11.955}{\in}%
\TextAt{16}{522.903}{277.849}{\Character{79}}%
\TextAt{8}{79.226}{260.803}{with}%
\TextAt{8}{105.886}{260.803}{minimal}%
\TextAt{8}{151.411}{260.803}{polynomial}%
\TextAt{13}{272.747}{243.757}{p}%
\TextAt{9}{278.622}{243.757}{\Character{40}}%
\TextAt{13}{283.175}{243.757}{X}%
\TextAt{9}{293.83}{243.757}{\Character{41}}%
\NativeMathAt{301.703}{243.757}{11.955}{\in}%
\TextAt{14}{312.994}{243.757}{o}%
\TextAt{9}{318.838}{243.757}{\Character{91}}%
\TextAt{13}{322.09}{243.757}{X}%
\TextAt{9}{332.745}{243.757}{\Character{93}}%
\TextAt{13}{335.997}{243.757}{\Character{59}}%
\TextAt{8}{79.226}{220.236}{so}%
\TextAt{8}{93.853}{220.236}{that}%
\TextAt{13}{119.478}{220.236}{L}%
\TextAt{9}{131.2}{220.236}{\Character{61}}%
\TextAt{13}{144.062}{220.236}{K}%
\TextAt{9}{154.88}{220.236}{\Character{40}}%
\TextAt{13}{159.432}{220.236}{\Character{18}}%
\TextAt{9}{165.212}{220.236}{\Character{41}}%
\TextAt{8}{169.765}{220.236}{.}%
\TextAt{8}{178.982}{220.236}{We}%
\TextAt{8}{199.383}{220.236}{may}%
\TextAt{8}{225.008}{220.236}{then}%
\TextAt{8}{251.924}{220.236}{deduce}%
\TextAt{8}{291.19}{220.236}{a}%
\TextAt{8}{301.202}{220.236}{result}%
\TextAt{8}{334.036}{220.236}{about}%
\TextAt{8}{367.791}{220.236}{the}%
\TextAt{8}{388.204}{220.236}{nature}%
\TextAt{8}{425.533}{220.236}{of}%
\TextAt{8}{439.119}{220.236}{the}%
\TextAt{8}{459.532}{220.236}{decomposition}%
\TextAt{8}{79.226}{203.19}{of}%
\TextAt{14}{92.516}{203.19}{p}%
\TextAt{8}{102.357}{203.19}{in}%
\TextAt{16}{115.97}{203.19}{\Character{79}}%
\TextAt{8}{129.677}{203.19}{which,}%
\TextAt{8}{166.384}{203.19}{albeit}%
\TextAt{8}{199.185}{203.19}{not}%
\TextAt{8}{219.959}{203.19}{complete,}%
\TextAt{8}{272.601}{203.19}{does}%
\TextAt{8}{298.958}{203.19}{show}%
\TextAt{8}{327.922}{203.19}{characteristic}%
\TextAt{8}{399.816}{203.19}{phenomena}%
\TextAt{8}{461.561}{203.19}{and}%
\TextAt{8}{484.283}{203.19}{a}%
\TextAt{8}{493.999}{203.19}{striking}%
\TextAt{8}{79.226}{186.144}{simplicity.}%
\TextAt{8}{135.926}{186.144}{It}%
\TextAt{8}{147.812}{186.144}{is}%
\TextAt{8}{158.788}{186.144}{incomplete}%
\TextAt{8}{217.18}{186.144}{in}%
\TextAt{8}{230.045}{186.144}{that}%
\TextAt{8}{254.627}{186.144}{a}%
\TextAt{8}{263.595}{186.144}{finite}%
\TextAt{8}{292.72}{186.144}{number}%
\TextAt{8}{334.529}{186.144}{of}%
\TextAt{8}{347.072}{186.144}{prime}%
\TextAt{8}{379.449}{186.144}{ideals}%
\TextAt{8}{411.24}{186.144}{are}%
\TextAt{8}{429.964}{186.144}{excluded;}%
\TextAt{8}{481.133}{186.144}{only}%
\TextAt{8}{506.038}{186.144}{those}%
\TextAt{8}{79.226}{169.098}{relatively}%
\TextAt{8}{130.559}{169.098}{prime}%
\TextAt{8}{163.996}{169.098}{to}%
\TextAt{8}{178.579}{169.098}{the}%
\TextAt{12}{199.008}{169.098}{conductor}%
\TextAt{8}{261.747}{169.098}{of}%
\TextAt{8}{275.349}{169.098}{the}%
\TextAt{8}{295.778}{169.098}{ring}%
\TextAt{14}{320.117}{169.098}{o}%
\TextAt{9}{325.961}{169.098}{\Character{91}}%
\TextAt{13}{329.213}{169.098}{\Character{18}}%
\TextAt{9}{334.993}{169.098}{\Character{93}}%
\TextAt{8}{342.415}{169.098}{can}%
\TextAt{8}{364.147}{169.098}{be}%
\TextAt{8}{380.344}{169.098}{considered.}%
\TextAt{8}{443.003}{169.098}{This}%
\TextAt{8}{469.995}{169.098}{conductor}%
\TextAt{8}{524.903}{169.098}{is}%
\TextAt{8}{79.226}{152.052}{defined}%
\TextAt{8}{119.539}{152.052}{to}%
\TextAt{8}{133.849}{152.052}{be}%
\TextAt{8}{149.774}{152.052}{the}%
\TextAt{8}{169.93}{152.052}{biggest}%
\TextAt{8}{209.669}{152.052}{ideal}%
\TextAt{14}{237.632}{152.052}{F}%
\TextAt{8}{248.836}{152.052}{of}%
\TextAt{16}{262.166}{152.052}{\Character{79}}%
\TextAt{8}{275.914}{152.052}{which}%
\TextAt{8}{309.4}{152.052}{is}%
\TextAt{8}{321.164}{152.052}{contained}%
\TextAt{8}{374.173}{152.052}{in}%
\TextAt{14}{387.826}{152.052}{o}%
\TextAt{9}{393.67}{152.052}{\Character{91}}%
\TextAt{13}{396.922}{152.052}{\Character{18}}%
\TextAt{9}{402.702}{152.052}{\Character{93}}%
\TextAt{8}{405.954}{152.052}{.}%
\TextAt{8}{414.402}{152.052}{In}%
\TextAt{8}{429.023}{152.052}{other}%
\TextAt{8}{459.593}{152.052}{words}%
\TextAt{14}{240.175}{123.051}{F}%
\TextAt{9}{250.803}{123.051}{\Character{61}}%
\TextAt{16}{263.228}{123.051}{\Character{102}}%
\TextAt{13}{269.206}{123.051}{\Character{11}}%
\NativeMathAt{280.048}{123.051}{11.955}{\in}%
\TextAt{16}{291.322}{123.051}{\Character{79}}%
\TextAt{16}{304.496}{123.051}{\Character{106}}%
\TextAt{13}{311.152}{123.051}{\Character{11}}%
\TextAt{16}{318.674}{123.051}{\Character{79}}%
\NativeMathAt{331.837}{123.051}{11.955}{\subseteq}%
\TextAt{14}{344.463}{123.051}{o}%
\TextAt{9}{350.308}{123.051}{\Character{91}}%
\TextAt{13}{353.559}{123.051}{\Character{18}}%
\TextAt{9}{359.34}{123.051}{\Character{93}}%
\TextAt{16}{362.591}{123.051}{\Character{103}}%
\TextAt{13}{368.569}{123.051}{\Character{58}}%
\TextAt{8}{79.226}{94.05}{Since}%
\TextAt{16}{109.784}{94.05}{\Character{79}}%
\TextAt{8}{123.531}{94.05}{is}%
\TextAt{8}{135.295}{94.05}{a}%
\TextAt{8}{145.05}{94.05}{finitely}%
\TextAt{8}{184.394}{94.05}{generated}%
\TextAt{14}{237.726}{94.05}{o}%
\TextAt{8}{243.571}{94.05}{-module}%
\TextAt{8}{288.761}{94.05}{(see}%
\TextAt{8}{312.229}{94.05}{proof}%
\TextAt{8}{342.799}{94.05}{of}%
\TextAt{8}{356.129}{94.05}{(8.1)),}%
\TextAt{8}{391.91}{94.05}{one}%
\TextAt{8}{413.37}{94.05}{has}%
\TextAt{14}{434.244}{94.05}{F}%
\NativeMathAt{444.871}{94.05}{11.955}{\neq}%
\TextAt{9}{457.281}{94.05}{0}%
\TextAt{8}{463.15}{94.05}{.}%
\TextAt{8}{300.14}{43.033}{44}%
\Destination{78.23}{809.815}{nav.page.44}%
\Destination{79.36}{657.455}{nav.print.47}%
\Hotlink{446.259}{758.955}{12.716}{11.047}{nav.print.45}%
\Hotlink{362.026}{640.255}{25.078}{11.047}{nav.number.1.2.3}%
\Hotlink{355.629}{91.855}{25.078}{11.047}{nav.number.1.8.1}%
\Destination{79.226}{676.531}{errata.55}%
}

% Page 57 (body page 45)
\Page{%
\TextAt{2}{79.226}{796.943}{I.}%
\TextAt{2}{89.83}{796.943}{Algebraic}%
\TextAt{2}{138.648}{796.943}{Integers}%
\Vector{q 0 G 0 g 0.299 w 7.226 18.729 m 460.769 18.729 l S Q}%
\TextAt{12}{79.226}{761.124}{Proposition}%
\TextAt{12}{151.233}{761.124}{(8.3).}%
\TextAt{8}{185.936}{761.124}{Let}%
\TextAt{14}{206.293}{761.124}{p}%
\TextAt{8}{215.561}{761.124}{be}%
\TextAt{8}{230.873}{761.124}{a}%
\TextAt{8}{240.016}{761.124}{prime}%
\TextAt{8}{272.567}{761.124}{ideal}%
\TextAt{8}{299.918}{761.124}{of}%
\TextAt{14}{312.636}{761.124}{o}%
\TextAt{8}{321.765}{761.124}{which}%
\TextAt{8}{354.639}{761.124}{is}%
\TextAt{8}{365.79}{761.124}{relatively}%
\TextAt{8}{416.239}{761.124}{prime}%
\TextAt{8}{448.79}{761.124}{to}%
\TextAt{8}{462.488}{761.124}{the}%
\TextAt{8}{482.032}{761.124}{conductor}%
\TextAt{14}{79.226}{744.078}{F}%
\TextAt{8}{90.43}{744.078}{of}%
\TextAt{14}{103.76}{744.078}{o}%
\TextAt{9}{109.605}{744.078}{\Character{91}}%
\TextAt{13}{112.856}{744.078}{\Character{18}}%
\TextAt{9}{118.637}{744.078}{\Character{93}}%
\TextAt{8}{121.888}{744.078}{,}%
\TextAt{8}{129.038}{744.078}{and}%
\TextAt{8}{151.8}{744.078}{let}%
\Vector{q 0 G 0 g 0.478 w 165.481 -40.147 m 171.357 -40.147 l S Q}%
\TextAt{13}{237.481}{727.032}{p}%
\TextAt{9}{243.357}{727.032}{\Character{40}}%
\TextAt{13}{247.909}{727.032}{X}%
\TextAt{9}{258.564}{727.032}{\Character{41}}%
\TextAt{9}{266.419}{727.032}{\Character{61}}%
\Vector{q 0 G 0 g 0.478 w 206.863 -40.147 m 212.738 -40.147 l S Q}%
\TextAt{13}{278.863}{727.032}{p}%
\TextAt{17}{284.738}{724.209}{1}%
\TextAt{9}{289.47}{727.032}{\Character{40}}%
\TextAt{13}{294.022}{727.032}{X}%
\TextAt{9}{304.678}{727.032}{\Character{41}}%
\TextAt{21}{309.23}{731.968}{e}%
\TextAt{24}{313.158}{730.861}{1}%
\TextAt{16}{319.8}{727.032}{\Character{1}}%
\TextAt{16}{325.096}{727.032}{\Character{1}}%
\TextAt{16}{330.404}{727.032}{\Character{1}}%
\Vector{q 0 G 0 g 0.478 w 263.74 -40.147 m 269.615 -40.147 l S Q}%
\TextAt{13}{335.74}{727.032}{p}%
\TextAt{21}{341.615}{724.209}{r}%
\TextAt{9}{346.17}{727.032}{\Character{40}}%
\TextAt{13}{350.722}{727.032}{X}%
\TextAt{9}{361.377}{727.032}{\Character{41}}%
\TextAt{21}{365.93}{731.968}{e}%
\TextAt{22}{369.858}{730.972}{r}%
\TextAt{8}{79.226}{703.51}{be}%
\TextAt{8}{96.722}{703.51}{the}%
\TextAt{8}{118.45}{703.51}{factorization}%
\TextAt{8}{187.998}{703.51}{of}%
\TextAt{8}{202.9}{703.51}{the}%
\TextAt{8}{224.627}{703.51}{polynomial}%
\Vector{q 0 G 0 g 0.478 w 214.692 -63.669 m 220.567 -63.669 l S Q}%
\TextAt{13}{286.692}{703.51}{p}%
\TextAt{9}{292.567}{703.51}{\Character{40}}%
\TextAt{13}{297.119}{703.51}{X}%
\TextAt{9}{307.775}{703.51}{\Character{41}}%
\TextAt{9}{318.308}{703.51}{\Character{61}}%
\TextAt{13}{333.429}{703.51}{p}%
\TextAt{9}{339.304}{703.51}{\Character{40}}%
\TextAt{13}{343.856}{703.51}{X}%
\TextAt{9}{354.512}{703.51}{\Character{41}}%
\TextAt{8}{362.385}{703.51}{mod}%
\TextAt{14}{387.823}{703.51}{p}%
\TextAt{8}{399.275}{703.51}{into}%
\TextAt{8}{424.589}{703.51}{irreducibles}%
\Vector{q 0 G 0 g 0.478 w 416.65 -63.669 m 422.525 -63.669 l S Q}%
\TextAt{13}{488.65}{703.51}{p}%
\TextAt{21}{494.525}{700.687}{i}%
\TextAt{9}{497.906}{703.51}{\Character{40}}%
\TextAt{13}{502.459}{703.51}{X}%
\TextAt{9}{513.114}{703.51}{\Character{41}}%
\TextAt{9}{523.647}{703.51}{\Character{61}}%
\TextAt{13}{79.226}{686.464}{p}%
\TextAt{21}{85.101}{684.671}{i}%
\TextAt{9}{88.483}{686.464}{\Character{40}}%
\TextAt{13}{93.035}{686.464}{X}%
\TextAt{9}{103.69}{686.464}{\Character{41}}%
\TextAt{8}{111.563}{686.464}{mod}%
\TextAt{14}{137.001}{686.464}{p}%
\TextAt{8}{146.882}{686.464}{over}%
\TextAt{8}{171.928}{686.464}{the}%
\TextAt{8}{192.084}{686.464}{residue}%
\TextAt{8}{231.811}{686.464}{class}%
\TextAt{8}{259.248}{686.464}{field}%
\TextAt{14}{284.605}{686.464}{o}%
\TextAt{9}{290.449}{686.464}{\Character{47}}%
\TextAt{14}{296.302}{686.464}{p}%
\TextAt{8}{302.285}{686.464}{,}%
\TextAt{8}{309.435}{686.464}{with}%
\TextAt{8}{336.095}{686.464}{all}%
\TextAt{13}{352.354}{686.464}{p}%
\TextAt{21}{358.229}{684.671}{i}%
\TextAt{9}{361.61}{686.464}{\Character{40}}%
\TextAt{13}{366.162}{686.464}{X}%
\TextAt{9}{376.818}{686.464}{\Character{41}}%
\NativeMathAt{384.691}{686.464}{11.955}{\in}%
\TextAt{14}{395.982}{686.464}{o}%
\TextAt{9}{401.826}{686.464}{\Character{91}}%
\TextAt{13}{405.078}{686.464}{X}%
\TextAt{9}{415.733}{686.464}{\Character{93}}%
\TextAt{8}{422.882}{686.464}{monic.}%
\TextAt{8}{461.899}{686.464}{Then}%
\TextAt{14}{214.782}{657.463}{P}%
\TextAt{21}{224.675}{655.67}{i}%
\TextAt{9}{231.377}{657.463}{\Character{61}}%
\TextAt{14}{243.802}{657.463}{p}%
\TextAt{16}{249.785}{657.463}{\Character{79}}%
\TextAt{9}{262.292}{657.463}{\Character{43}}%
\TextAt{13}{274.053}{657.463}{p}%
\TextAt{21}{279.928}{655.67}{i}%
\TextAt{9}{283.31}{657.463}{\Character{40}}%
\TextAt{13}{287.862}{657.463}{\Character{18}}%
\TextAt{9}{293.642}{657.463}{\Character{41}}%
\TextAt{16}{298.195}{657.463}{\Character{79}}%
\TextAt{13}{308.044}{657.463}{\Character{59}}%
\TextAt{13}{337.179}{657.463}{i}%
\TextAt{9}{344.513}{657.463}{\Character{61}}%
\TextAt{9}{356.923}{657.463}{1}%
\TextAt{13}{362.792}{657.463}{\Character{59}}%
\TextAt{13}{368.016}{657.463}{\Character{58}}%
\TextAt{13}{373.253}{657.463}{\Character{58}}%
\TextAt{13}{378.477}{657.463}{\Character{58}}%
\TextAt{13}{383.714}{657.463}{\Character{59}}%
\TextAt{13}{388.95}{657.463}{r}%
\TextAt{13}{393.911}{657.463}{\Character{59}}%
\TextAt{8}{79.226}{628.462}{are}%
\TextAt{8}{99.158}{628.462}{the}%
\TextAt{8}{119.735}{628.462}{different}%
\TextAt{8}{166.327}{628.462}{prime}%
\TextAt{8}{199.912}{628.462}{ideals}%
\TextAt{8}{232.91}{628.462}{of}%
\TextAt{16}{246.661}{628.462}{\Character{79}}%
\TextAt{8}{260.829}{628.462}{above}%
\TextAt{14}{294.426}{628.462}{p}%
\TextAt{8}{300.409}{628.462}{.}%
\TextAt{8}{310.12}{628.462}{The}%
\TextAt{8}{334.595}{628.462}{inertia}%
\TextAt{8}{372.089}{628.462}{degree}%
\TextAt{13}{408.925}{628.462}{f}%
\TextAt{21}{414.695}{626.669}{i}%
\TextAt{8}{422.395}{628.462}{of}%
\TextAt{14}{436.146}{628.462}{P}%
\TextAt{21}{446.039}{626.669}{i}%
\TextAt{8}{453.738}{628.462}{is}%
\TextAt{8}{465.923}{628.462}{the}%
\TextAt{8}{486.501}{628.462}{degree}%
\TextAt{8}{523.337}{628.462}{of}%
\Vector{q 0 G 0 g 0.478 w 7.226 -155.763 m 13.101 -155.763 l S Q}%
\TextAt{13}{79.226}{611.416}{p}%
\TextAt{21}{85.101}{608.593}{i}%
\TextAt{9}{88.483}{611.416}{\Character{40}}%
\TextAt{13}{93.035}{611.416}{X}%
\TextAt{9}{103.69}{611.416}{\Character{41}}%
\TextAt{8}{108.242}{611.416}{,}%
\TextAt{8}{115.392}{611.416}{and}%
\TextAt{8}{138.154}{611.416}{one}%
\TextAt{8}{159.614}{611.416}{has}%
\TextAt{14}{266.067}{594.37}{p}%
\TextAt{9}{275.371}{594.37}{\Character{61}}%
\TextAt{14}{287.796}{594.37}{P}%
\TextAt{21}{297.689}{599.595}{e}%
\TextAt{24}{301.618}{598.488}{1}%
\TextAt{17}{297.689}{591.439}{1}%
\TextAt{16}{308.259}{594.37}{\Character{1}}%
\TextAt{16}{313.555}{594.37}{\Character{1}}%
\TextAt{16}{318.863}{594.37}{\Character{1}}%
\TextAt{14}{324.2}{594.37}{P}%
\TextAt{21}{334.093}{599.306}{e}%
\TextAt{22}{338.021}{598.31}{r}%
\TextAt{21}{334.093}{591.414}{r}%
\TextAt{13}{342.677}{594.37}{\Character{58}}%
\TextAt{19}{79.226}{552.915}{Proof.}%
\TextAt{8}{116}{552.915}{Writing}%
\TextAt{16}{158.919}{552.915}{\Character{79}}%
\TextAt{3}{168.769}{557.254}{\Character{48}}%
\TextAt{9}{174.885}{552.915}{\Character{61}}%
\TextAt{14}{187.31}{552.915}{o}%
\TextAt{9}{193.155}{552.915}{\Character{91}}%
\TextAt{13}{196.406}{552.915}{\Character{18}}%
\TextAt{9}{202.187}{552.915}{\Character{93}}%
\TextAt{8}{209.336}{552.915}{and}%
\Vector{q 0 G 0 g 0.478 w 160.098 -213.728 m 165.943 -213.728 l S Q}%
\TextAt{14}{232.098}{552.915}{o}%
\TextAt{9}{241.263}{552.915}{\Character{61}}%
\TextAt{14}{253.689}{552.915}{o}%
\TextAt{9}{259.533}{552.915}{\Character{47}}%
\TextAt{14}{265.386}{552.915}{p}%
\TextAt{8}{271.369}{552.915}{,}%
\TextAt{8}{278.518}{552.915}{we}%
\TextAt{8}{295.746}{552.915}{have}%
\TextAt{8}{322.741}{552.915}{a}%
\TextAt{8}{332.496}{552.915}{canonical}%
\TextAt{8}{383.879}{552.915}{isomorphism}%
\TextAt{16}{222.304}{523.914}{\Character{79}}%
\TextAt{9}{232.153}{523.914}{\Character{47}}%
\TextAt{14}{238.006}{523.914}{p}%
\TextAt{16}{243.989}{523.914}{\Character{79}}%
\NativeMathAt{257.16}{523.914}{11.955}{\cong}%
\TextAt{16}{269.779}{523.914}{\Character{79}}%
\TextAt{3}{279.629}{528.85}{\Character{48}}%
\TextAt{9}{282.424}{523.914}{\Character{47}}%
\TextAt{14}{288.277}{523.914}{p}%
\TextAt{16}{294.26}{523.914}{\Character{79}}%
\TextAt{3}{304.11}{528.85}{\Character{48}}%
\NativeMathAt{310.226}{523.914}{11.955}{\cong}%
\Vector{q 0 G 0 g 0.478 w 250.845 -242.729 m 256.69 -242.729 l S Q}%
\TextAt{14}{322.845}{523.914}{o}%
\TextAt{9}{328.69}{523.914}{\Character{91}}%
\TextAt{13}{331.941}{523.914}{X}%
\TextAt{9}{342.596}{523.914}{\Character{93}\Character{47}\Character{40}}%
\Vector{q 0 G 0 g 0.478 w 284.253 -243.265 m 290.128 -243.265 l S Q}%
\TextAt{13}{356.253}{523.914}{p}%
\TextAt{9}{362.128}{523.914}{\Character{40}}%
\TextAt{13}{366.681}{523.914}{X}%
\TextAt{9}{377.336}{523.914}{\Character{41}\Character{41}}%
\TextAt{13}{386.441}{523.914}{\Character{58}}%
\TextAt{8}{79.226}{494.913}{The}%
\TextAt{8}{104.148}{494.913}{first}%
\TextAt{8}{129.142}{494.913}{isomorphism}%
\TextAt{8}{198.43}{494.913}{follows}%
\TextAt{8}{237.734}{494.913}{from}%
\TextAt{8}{266.243}{494.913}{the}%
\TextAt{8}{287.268}{494.913}{relative}%
\TextAt{8}{329.764}{494.913}{primality}%
\TextAt{14}{381.37}{494.913}{p}%
\TextAt{16}{387.353}{494.913}{\Character{79}}%
\TextAt{9}{400.452}{494.913}{\Character{43}}%
\TextAt{14}{412.805}{494.913}{F}%
\TextAt{9}{424.912}{494.913}{\Character{61}}%
\TextAt{16}{438.818}{494.913}{\Character{79}}%
\TextAt{8}{448.668}{494.913}{.}%
\TextAt{8}{459.721}{494.913}{As}%
\TextAt{14}{477.877}{494.913}{F}%
\NativeMathAt{489.984}{494.913}{11.955}{\subseteq}%
\TextAt{16}{504.067}{494.913}{\Character{79}}%
\TextAt{3}{513.933}{499.252}{\Character{48}}%
\TextAt{8}{516.728}{494.913}{,}%
\TextAt{8}{524.963}{494.913}{it}%
\TextAt{8}{79.226}{477.867}{follows}%
\TextAt{8}{116.924}{477.867}{that}%
\TextAt{16}{141.555}{477.867}{\Character{79}}%
\TextAt{9}{154.725}{477.867}{\Character{61}}%
\TextAt{14}{167.151}{477.867}{p}%
\TextAt{16}{173.134}{477.867}{\Character{79}}%
\TextAt{9}{184.131}{477.867}{\Character{43}}%
\TextAt{16}{194.383}{477.867}{\Character{79}}%
\TextAt{3}{204.233}{482.206}{\Character{48}}%
\TextAt{8}{207.028}{477.867}{,}%
\TextAt{8}{213.586}{477.867}{i.e.,}%
\TextAt{8}{235.101}{477.867}{the}%
\TextAt{8}{254.519}{477.867}{homomorphism}%
\TextAt{16}{336.451}{477.867}{\Character{79}}%
\TextAt{3}{346.301}{482.206}{\Character{48}}%
\NativeMathAt{352.417}{477.867}{11.955}{\rightarrow}%
\TextAt{16}{367.684}{477.867}{\Character{79}}%
\TextAt{9}{377.543}{477.867}{\Character{47}}%
\TextAt{14}{383.396}{477.867}{p}%
\TextAt{16}{389.379}{477.867}{\Character{79}}%
\TextAt{8}{402.388}{477.867}{is}%
\TextAt{8}{413.414}{477.867}{surjective.}%
\TextAt{8}{470.13}{477.867}{It}%
\TextAt{8}{482.065}{477.867}{has}%
\TextAt{8}{502.2}{477.867}{kernel}%
\TextAt{14}{79.226}{460.821}{p}%
\TextAt{16}{85.209}{460.821}{\Character{79}}%
\TextAt{16}{97.714}{460.821}{\Character{92}}%
\TextAt{16}{108.33}{460.821}{\Character{79}}%
\TextAt{3}{118.192}{465.16}{\Character{48}}%
\TextAt{8}{120.987}{460.821}{,}%
\TextAt{8}{128.137}{460.821}{which}%
\TextAt{8}{161.623}{460.821}{equals}%
\TextAt{14}{197.13}{460.821}{p}%
\TextAt{16}{203.113}{460.821}{\Character{79}}%
\TextAt{3}{212.963}{465.16}{\Character{48}}%
\TextAt{8}{215.758}{460.821}{.}%
\TextAt{8}{224.206}{460.821}{Since}%
\TextAt{9}{254.763}{460.821}{\Character{40}}%
\TextAt{14}{259.316}{460.821}{p}%
\TextAt{13}{265.299}{460.821}{\Character{59}}%
\TextAt{14}{270.543}{460.821}{F}%
\TextAt{16}{280.507}{460.821}{\Character{92}}%
\TextAt{14}{291.133}{460.821}{o}%
\TextAt{9}{296.978}{460.821}{\Character{41}}%
\TextAt{9}{304.833}{460.821}{\Character{61}}%
\TextAt{9}{317.254}{460.821}{1}%
\TextAt{8}{323.129}{460.821}{,}%
\TextAt{8}{330.279}{460.821}{it}%
\TextAt{8}{341.983}{460.821}{follows}%
\TextAt{8}{380.418}{460.821}{that}%
\TextAt{14}{214.495}{431.82}{p}%
\TextAt{16}{220.478}{431.82}{\Character{79}}%
\TextAt{16}{232.983}{431.82}{\Character{92}}%
\TextAt{16}{243.599}{431.82}{\Character{79}}%
\TextAt{3}{253.461}{436.756}{\Character{48}}%
\TextAt{9}{259.577}{431.82}{\Character{61}}%
\TextAt{9}{271.987}{431.82}{\Character{40}}%
\TextAt{14}{276.555}{431.82}{p}%
\TextAt{9}{285.195}{431.82}{\Character{43}}%
\TextAt{14}{296.956}{431.82}{F}%
\TextAt{9}{304.263}{431.82}{\Character{41}\Character{40}}%
\TextAt{14}{313.368}{431.82}{p}%
\TextAt{16}{319.351}{431.82}{\Character{79}}%
\TextAt{16}{331.844}{431.82}{\Character{92}}%
\TextAt{16}{342.472}{431.82}{\Character{79}}%
\TextAt{3}{352.334}{436.756}{\Character{48}}%
\TextAt{9}{355.129}{431.82}{\Character{41}}%
\NativeMathAt{363.002}{431.82}{11.955}{\subseteq}%
\TextAt{14}{375.621}{431.82}{p}%
\TextAt{16}{381.604}{431.82}{\Character{79}}%
\TextAt{3}{391.454}{436.756}{\Character{48}}%
\TextAt{13}{394.249}{431.82}{\Character{58}}%
\TextAt{8}{79.226}{402.819}{The}%
\TextAt{8}{103.28}{402.819}{second}%
\TextAt{8}{141.058}{402.819}{isomorphism}%
\TextAt{8}{209.478}{402.819}{is}%
\TextAt{8}{221.242}{402.819}{deduced}%
\TextAt{8}{266.755}{402.819}{from}%
\TextAt{8}{294.395}{402.819}{the}%
\TextAt{8}{314.552}{402.819}{surjective}%
\TextAt{8}{366.963}{402.819}{homomorphism}%
\TextAt{14}{248.122}{373.818}{o}%
\TextAt{9}{253.966}{373.818}{\Character{91}}%
\TextAt{13}{257.218}{373.818}{X}%
\TextAt{9}{267.873}{373.818}{\Character{93}}%
\NativeArrowAt{274.445}{373.818}{11.955}{19.26}{\longrightarrow}%
\Vector{q 0 G 0 g 0.478 w 225.027 -392.826 m 230.871 -392.826 l S Q}%
\TextAt{14}{297.027}{373.818}{o}%
\TextAt{9}{302.871}{373.818}{\Character{91}}%
\TextAt{13}{306.123}{373.818}{X}%
\TextAt{9}{316.778}{373.818}{\Character{93}\Character{47}\Character{40}}%
\Vector{q 0 G 0 g 0.478 w 258.435 -393.361 m 264.31 -393.361 l S Q}%
\TextAt{13}{330.435}{373.818}{p}%
\TextAt{9}{336.31}{373.818}{\Character{40}}%
\TextAt{13}{340.863}{373.818}{X}%
\TextAt{9}{351.518}{373.818}{\Character{41}\Character{41}}%
\TextAt{13}{360.622}{373.818}{\Character{58}}%
\TextAt{8}{79.226}{344.817}{Its}%
\TextAt{8}{96.563}{344.817}{kernel}%
\TextAt{8}{131.079}{344.817}{is}%
\TextAt{8}{142.893}{344.817}{the}%
\TextAt{8}{163.099}{344.817}{ideal}%
\TextAt{8}{191.112}{344.817}{generated}%
\TextAt{8}{244.494}{344.817}{by}%
\TextAt{14}{260.802}{344.817}{p}%
\TextAt{8}{270.733}{344.817}{and}%
\TextAt{13}{293.545}{344.817}{p}%
\TextAt{9}{299.42}{344.817}{\Character{40}}%
\TextAt{13}{303.972}{344.817}{X}%
\TextAt{9}{314.628}{344.817}{\Character{41}}%
\TextAt{8}{319.18}{344.817}{,}%
\TextAt{8}{326.391}{344.817}{and}%
\TextAt{8}{349.204}{344.817}{in}%
\TextAt{8}{362.906}{344.817}{view}%
\TextAt{8}{389.939}{344.817}{of}%
\TextAt{16}{403.318}{344.817}{\Character{79}}%
\TextAt{3}{413.168}{349.155}{\Character{48}}%
\TextAt{9}{419.369}{344.817}{\Character{61}}%
\TextAt{14}{431.879}{344.817}{o}%
\TextAt{9}{437.723}{344.817}{\Character{91}}%
\TextAt{13}{440.975}{344.817}{\Character{18}}%
\TextAt{9}{446.755}{344.817}{\Character{93}}%
\TextAt{9}{453.39}{344.817}{\Character{61}}%
\TextAt{14}{465.923}{344.817}{o}%
\TextAt{9}{471.767}{344.817}{\Character{91}}%
\TextAt{13}{475.018}{344.817}{X}%
\TextAt{9}{485.673}{344.817}{\Character{93}\Character{47}\Character{40}}%
\TextAt{13}{499.33}{344.817}{p}%
\TextAt{9}{505.206}{344.817}{\Character{40}}%
\TextAt{13}{509.758}{344.817}{X}%
\TextAt{9}{520.413}{344.817}{\Character{41}\Character{41}}%
\TextAt{8}{529.518}{344.817}{,}%
\TextAt{8}{79.226}{327.771}{we}%
\TextAt{8}{96.454}{327.771}{have}%
\TextAt{16}{123.448}{327.771}{\Character{79}}%
\TextAt{3}{133.298}{332.109}{\Character{48}}%
\TextAt{9}{136.093}{327.771}{\Character{47}}%
\TextAt{14}{141.946}{327.771}{p}%
\TextAt{16}{147.929}{327.771}{\Character{79}}%
\TextAt{3}{157.779}{332.109}{\Character{48}}%
\NativeMathAt{163.895}{327.771}{11.955}{\cong}%
\Vector{q 0 G 0 g 0.478 w 104.514 -438.873 m 110.359 -438.873 l S Q}%
\TextAt{14}{176.514}{327.771}{o}%
\TextAt{9}{182.359}{327.771}{\Character{91}}%
\TextAt{13}{185.61}{327.771}{X}%
\TextAt{9}{196.265}{327.771}{\Character{93}\Character{47}\Character{40}}%
\Vector{q 0 G 0 g 0.478 w 137.922 -439.408 m 143.797 -439.408 l S Q}%
\TextAt{13}{209.922}{327.771}{p}%
\TextAt{9}{215.797}{327.771}{\Character{40}}%
\TextAt{13}{220.35}{327.771}{X}%
\TextAt{9}{231.005}{327.771}{\Character{41}\Character{41}}%
\TextAt{8}{240.109}{327.771}{.}%
\TextAt{8}{97.159}{307.736}{Since}%
\Vector{q 0 G 0 g 0.478 w 56.172 -459.443 m 62.047 -459.443 l S Q}%
\TextAt{13}{128.172}{307.736}{p}%
\TextAt{9}{134.047}{307.736}{\Character{40}}%
\TextAt{13}{138.6}{307.736}{X}%
\TextAt{9}{149.255}{307.736}{\Character{41}}%
\TextAt{9}{157.887}{307.736}{\Character{61}}%
\TextAt{20}{171.108}{316.702}{\Character{81}}%
\TextAt{21}{182.399}{313.548}{r}%
\TextAt{21}{182.399}{304.249}{i}%
\TextAt{17}{185.282}{304.249}{\Character{61}1}%
\Vector{q 0 G 0 g 0.478 w 126.593 -459.443 m 132.468 -459.443 l S Q}%
\TextAt{13}{198.593}{307.736}{p}%
\TextAt{21}{204.468}{304.913}{i}%
\TextAt{9}{207.849}{307.736}{\Character{40}}%
\TextAt{13}{212.402}{307.736}{X}%
\TextAt{9}{223.057}{307.736}{\Character{41}}%
\TextAt{21}{227.609}{312.074}{e}%
\TextAt{22}{231.538}{310.859}{i}%
\TextAt{8}{235.197}{307.736}{,}%
\TextAt{8}{242.917}{307.736}{the}%
\TextAt{8}{263.529}{307.736}{Chinese}%
\TextAt{8}{307.609}{307.736}{remainder}%
\TextAt{8}{363.346}{307.736}{theorem}%
\TextAt{8}{409.327}{307.736}{finally}%
\TextAt{8}{445.23}{307.736}{gives}%
\TextAt{8}{474.367}{307.736}{the}%
\TextAt{8}{494.979}{307.736}{isomor-}%
\TextAt{8}{79.226}{290.69}{phism}%
\Vector{q 0 G 0 g 0.478 w 145.293 -498.84 m 151.137 -498.84 l S Q}%
\TextAt{14}{217.293}{267.803}{o}%
\TextAt{9}{223.137}{267.803}{\Character{91}}%
\TextAt{13}{226.389}{267.803}{X}%
\TextAt{9}{237.044}{267.803}{\Character{93}\Character{47}\Character{40}}%
\Vector{q 0 G 0 g 0.478 w 178.701 -499.376 m 184.576 -499.376 l S Q}%
\TextAt{13}{250.701}{267.803}{p}%
\TextAt{9}{256.576}{267.803}{\Character{40}}%
\TextAt{13}{261.128}{267.803}{X}%
\TextAt{9}{271.783}{267.803}{\Character{41}\Character{41}}%
\NativeMathAt{284.209}{267.803}{11.955}{\cong}%
\TextAt{21}{303.833}{282.747}{r}%
\TextAt{20}{296.828}{279.161}{\Character{77}}%
\TextAt{21}{299.009}{253.966}{i}%
\TextAt{17}{301.892}{253.966}{\Character{61}1}%
\Vector{q 0 G 0 g 0.478 w 244.886 -498.84 m 250.731 -498.84 l S Q}%
\TextAt{14}{316.886}{267.803}{o}%
\TextAt{9}{322.731}{267.803}{\Character{91}}%
\TextAt{13}{325.982}{267.803}{X}%
\TextAt{9}{336.637}{267.803}{\Character{93}\Character{47}\Character{40}}%
\Vector{q 0 G 0 g 0.478 w 278.294 -499.376 m 284.17 -499.376 l S Q}%
\TextAt{13}{350.294}{267.803}{p}%
\TextAt{21}{356.17}{264.981}{i}%
\TextAt{9}{359.551}{267.803}{\Character{40}}%
\TextAt{13}{364.103}{267.803}{X}%
\TextAt{9}{374.758}{267.803}{\Character{41}\Character{41}}%
\TextAt{21}{383.863}{272.739}{e}%
\TextAt{22}{387.791}{271.524}{i}%
\TextAt{13}{391.451}{267.803}{\Character{58}}%
\TextAt{8}{79.226}{236.333}{This}%
\TextAt{8}{105.429}{236.333}{shows}%
\TextAt{8}{138.529}{236.333}{that}%
\TextAt{8}{163.38}{236.333}{the}%
\TextAt{8}{183.019}{236.333}{prime}%
\TextAt{8}{215.666}{236.333}{ideals}%
\TextAt{8}{247.726}{236.333}{of}%
\TextAt{8}{260.539}{236.333}{the}%
\TextAt{8}{280.178}{236.333}{ring}%
\TextAt{13}{303.726}{236.333}{R}%
\TextAt{9}{316.055}{236.333}{\Character{61}}%
\Vector{q 0 G 0 g 0.478 w 256.481 -530.311 m 262.325 -530.311 l S Q}%
\TextAt{14}{328.481}{236.333}{o}%
\TextAt{9}{334.325}{236.333}{\Character{91}}%
\TextAt{13}{337.577}{236.333}{X}%
\TextAt{9}{348.232}{236.333}{\Character{93}\Character{47}\Character{40}}%
\Vector{q 0 G 0 g 0.478 w 289.889 -530.846 m 295.764 -530.846 l S Q}%
\TextAt{13}{361.889}{236.333}{p}%
\TextAt{9}{367.764}{236.333}{\Character{40}}%
\TextAt{13}{372.316}{236.333}{X}%
\TextAt{9}{382.971}{236.333}{\Character{41}\Character{41}}%
\TextAt{8}{395.456}{236.333}{are}%
\TextAt{8}{414.449}{236.333}{the}%
\TextAt{8}{434.088}{236.333}{principal}%
\TextAt{8}{482.348}{236.333}{ideals}%
\TextAt{9}{514.408}{236.333}{\Character{40}}%
\Vector{q 0 G 0 g 0.478 w 446.961 -530.846 m 452.836 -530.846 l S Q}%
\TextAt{13}{518.961}{236.333}{p}%
\TextAt{21}{524.836}{233.51}{i}%
\TextAt{9}{528.217}{236.333}{\Character{41}}%
\TextAt{8}{79.226}{219.287}{generated}%
\TextAt{8}{132.343}{219.287}{by}%
\TextAt{8}{148.387}{219.287}{the}%
\Vector{q 0 G 0 g 0.478 w 96.329 -547.892 m 102.204 -547.892 l S Q}%
\TextAt{13}{168.329}{219.287}{p}%
\TextAt{21}{174.204}{216.464}{i}%
\TextAt{9}{177.585}{219.287}{\Character{40}}%
\TextAt{13}{182.137}{219.287}{X}%
\TextAt{9}{192.793}{219.287}{\Character{41}}%
\TextAt{8}{200.666}{219.287}{mod}%
\Vector{q 0 G 0 g 0.478 w 154.104 -547.892 m 159.979 -547.892 l S Q}%
\TextAt{13}{226.104}{219.287}{p}%
\TextAt{9}{231.979}{219.287}{\Character{40}}%
\TextAt{13}{236.531}{219.287}{X}%
\TextAt{9}{247.186}{219.287}{\Character{41}}%
\TextAt{8}{251.738}{219.287}{,}%
\TextAt{8}{258.716}{219.287}{for}%
\TextAt{13}{276.386}{219.287}{i}%
\TextAt{9}{283.7}{219.287}{\Character{61}}%
\TextAt{9}{296.11}{219.287}{1}%
\TextAt{13}{301.978}{219.287}{\Character{59}}%
\TextAt{13}{307.202}{219.287}{\Character{58}}%
\TextAt{13}{312.439}{219.287}{\Character{58}}%
\TextAt{13}{317.675}{219.287}{\Character{58}}%
\TextAt{13}{322.9}{219.287}{\Character{59}}%
\TextAt{13}{328.136}{219.287}{r}%
\TextAt{8}{333.8}{219.287}{,}%
\TextAt{8}{340.777}{219.287}{that}%
\TextAt{8}{365.931}{219.287}{the}%
\TextAt{8}{385.872}{219.287}{degree}%
\TextAt{9}{422.073}{219.287}{\Character{91}}%
\TextAt{13}{425.324}{219.287}{R}%
\TextAt{9}{434.333}{219.287}{\Character{47}\Character{40}}%
\Vector{q 0 G 0 g 0.478 w 372.738 -547.892 m 378.613 -547.892 l S Q}%
\TextAt{13}{444.738}{219.287}{p}%
\TextAt{21}{450.613}{216.464}{i}%
\TextAt{9}{453.995}{219.287}{\Character{41}}%
\TextAt{9}{461.85}{219.287}{\Character{58}}%
\Vector{q 0 G 0 g 0.478 w 396.44 -547.357 m 402.285 -547.357 l S Q}%
\TextAt{14}{468.44}{219.287}{o}%
\TextAt{9}{474.285}{219.287}{\Character{93}}%
\TextAt{8}{481.219}{219.287}{equals}%
\TextAt{8}{516.51}{219.287}{the}%
\TextAt{8}{79.226}{202.241}{degree}%
\TextAt{8}{115.642}{202.241}{of}%
\TextAt{8}{128.972}{202.241}{the}%
\TextAt{8}{149.128}{202.241}{polynomial}%
\Vector{q 0 G 0 g 0.478 w 137.621 -564.938 m 143.496 -564.938 l S Q}%
\TextAt{13}{209.621}{202.241}{p}%
\TextAt{21}{215.496}{199.418}{i}%
\TextAt{9}{218.878}{202.241}{\Character{40}}%
\TextAt{13}{223.43}{202.241}{X}%
\TextAt{9}{234.085}{202.241}{\Character{41}}%
\TextAt{8}{238.637}{202.241}{,}%
\TextAt{8}{245.787}{202.241}{and}%
\TextAt{8}{268.549}{202.241}{that}%
\TextAt{9}{253.83}{166.729}{\Character{40}0\Character{41}}%
\TextAt{9}{272.074}{166.729}{\Character{61}}%
\TextAt{9}{284.495}{166.729}{\Character{40}}%
\Vector{q 0 G 0 g 0.478 w 217.087 -600.45 m 222.962 -600.45 l S Q}%
\TextAt{13}{289.087}{166.729}{p}%
\TextAt{9}{294.962}{166.729}{\Character{41}}%
\TextAt{9}{302.817}{166.729}{\Character{61}}%
\TextAt{21}{320.084}{181.673}{r}%
\TextAt{20}{315.471}{178.087}{\Character{92}}%
\TextAt{21}{315.26}{152.892}{i}%
\TextAt{17}{318.144}{152.892}{\Character{61}1}%
\TextAt{9}{328.964}{166.729}{\Character{40}}%
\Vector{q 0 G 0 g 0.478 w 261.517 -600.45 m 267.392 -600.45 l S Q}%
\TextAt{13}{333.517}{166.729}{p}%
\TextAt{21}{339.392}{163.907}{i}%
\TextAt{9}{342.773}{166.729}{\Character{41}}%
\TextAt{21}{347.325}{171.665}{e}%
\TextAt{22}{351.254}{170.45}{i}%
\TextAt{13}{354.914}{166.729}{\Character{58}}%
\TextAt{8}{79.226}{129.779}{In}%
\TextAt{8}{95.251}{129.779}{view}%
\TextAt{8}{123.637}{129.779}{of}%
\TextAt{8}{138.37}{129.779}{the}%
\TextAt{8}{159.929}{129.779}{isomorphism}%
\Vector{q 0 G 0 g 0.478 w 157.752 -636.865 m 163.596 -636.865 l S Q}%
\TextAt{14}{229.752}{129.779}{o}%
\TextAt{9}{235.596}{129.779}{\Character{91}}%
\TextAt{13}{238.848}{129.779}{X}%
\TextAt{9}{249.503}{129.779}{\Character{93}\Character{47}\Character{40}}%
\Vector{q 0 G 0 g 0.478 w 191.16 -637.4 m 197.035 -637.4 l S Q}%
\TextAt{13}{263.16}{129.779}{p}%
\TextAt{9}{269.035}{129.779}{\Character{40}}%
\TextAt{13}{273.588}{129.779}{X}%
\TextAt{9}{284.243}{129.779}{\Character{41}\Character{41}}%
\NativeMathAt{299.059}{129.779}{11.955}{\cong}%
\TextAt{16}{314.07}{129.779}{\Character{79}}%
\TextAt{9}{323.92}{129.779}{\Character{47}}%
\TextAt{14}{329.773}{129.779}{p}%
\TextAt{16}{335.756}{129.779}{\Character{79}}%
\TextAt{8}{345.606}{129.779}{,}%
\TextAt{13}{354.509}{129.779}{f}%
\TextAt{9}{361.555}{129.779}{\Character{40}}%
\TextAt{13}{366.108}{129.779}{X}%
\TextAt{9}{376.763}{129.779}{\Character{41}}%
\NativeArrowAt{387.027}{129.779}{11.955}{11.955}{\mapsto}%
\TextAt{13}{404.695}{129.779}{f}%
\TextAt{9}{411.741}{129.779}{\Character{40}}%
\TextAt{13}{416.293}{129.779}{\Character{18}}%
\TextAt{9}{422.074}{129.779}{\Character{41}}%
\TextAt{8}{426.626}{129.779}{,}%
\TextAt{8}{435.529}{129.779}{the}%
\TextAt{8}{457.089}{129.779}{same}%
\TextAt{8}{487.818}{129.779}{situation}%
\TextAt{8}{79.226}{112.733}{holds}%
\TextAt{8}{110.657}{112.733}{in}%
\TextAt{8}{125.111}{112.733}{the}%
\TextAt{8}{146.069}{112.733}{ring}%
\Vector{q 0 G 0 g 0.478 w 98.936 -651.424 m 108.786 -651.424 l S Q}%
\TextAt{16}{170.936}{112.733}{\Character{79}}%
\TextAt{9}{185.473}{112.733}{\Character{61}}%
\TextAt{16}{199.264}{112.733}{\Character{79}}%
\TextAt{9}{209.114}{112.733}{\Character{47}}%
\TextAt{14}{214.967}{112.733}{p}%
\TextAt{16}{220.95}{112.733}{\Character{79}}%
\TextAt{8}{230.8}{112.733}{.}%
\TextAt{8}{241.653}{112.733}{Thus}%
\TextAt{8}{272.103}{112.733}{the}%
\TextAt{8}{293.061}{112.733}{prime}%
\TextAt{8}{327.026}{112.733}{ideals}%
\Vector{q 0 G 0 g 0.478 w 288.406 -651.327 m 298.299 -651.327 l S Q}%
\TextAt{14}{360.406}{112.733}{P}%
\TextAt{21}{370.299}{109.975}{i}%
\TextAt{8}{378.379}{112.733}{of}%
\Vector{q 0 G 0 g 0.478 w 320.511 -651.424 m 330.36 -651.424 l S Q}%
\TextAt{16}{392.511}{112.733}{\Character{79}}%
\TextAt{8}{407.059}{112.733}{correspond}%
\TextAt{8}{467.433}{112.733}{to}%
\TextAt{8}{482.545}{112.733}{the}%
\TextAt{8}{503.503}{112.733}{prime}%
\TextAt{8}{79.226}{95.687}{ideals}%
\TextAt{9}{111.924}{95.687}{\Character{40}}%
\Vector{q 0 G 0 g 0.478 w 44.476 -671.492 m 50.352 -671.492 l S Q}%
\TextAt{13}{116.476}{95.687}{p}%
\TextAt{21}{122.352}{92.864}{i}%
\TextAt{9}{125.733}{95.687}{\Character{41}}%
\TextAt{8}{130.285}{95.687}{,}%
\TextAt{8}{137.585}{95.687}{and}%
\TextAt{8}{160.467}{95.687}{they}%
\TextAt{8}{186.925}{95.687}{are}%
\TextAt{8}{206.556}{95.687}{the}%
\TextAt{8}{226.832}{95.687}{principal}%
\TextAt{8}{275.73}{95.687}{ideals}%
\TextAt{8}{308.428}{95.687}{generated}%
\TextAt{8}{361.88}{95.687}{by}%
\TextAt{8}{378.259}{95.687}{the}%
\TextAt{13}{398.536}{95.687}{p}%
\TextAt{21}{404.411}{93.894}{i}%
\TextAt{9}{407.792}{95.687}{\Character{40}}%
\TextAt{13}{412.344}{95.687}{\Character{18}}%
\TextAt{9}{418.125}{95.687}{\Character{41}}%
\TextAt{8}{425.998}{95.687}{mod}%
\TextAt{14}{451.436}{95.687}{p}%
\TextAt{16}{457.419}{95.687}{\Character{79}}%
\TextAt{8}{467.269}{95.687}{.}%
\TextAt{8}{476.077}{95.687}{The}%
\TextAt{8}{500.251}{95.687}{degree}%
\TextAt{8}{300.14}{43.033}{45}%
\Destination{78.23}{809.815}{nav.page.45}%
\Destination{79.36}{718.555}{nav.print.48}%
\Destination{79.36}{775.155}{nav.number.1.8.3}%
}

% Page 58 (body page 46)
\Page{%
\TextAt{2}{79.226}{796.943}{I.}%
\TextAt{2}{89.83}{796.943}{Algebraic}%
\TextAt{2}{138.648}{796.943}{Integers}%
\Vector{q 0 G 0 g 0.299 w 7.226 18.729 m 460.769 18.729 l S Q}%
\TextAt{9}{79.226}{761.124}{\Character{91}}%
\TextAt{16}{82.478}{761.124}{\Character{79}}%
\TextAt{9}{92.328}{761.124}{\Character{47}}%
\TextAt{14}{98.181}{761.124}{P}%
\TextAt{21}{108.074}{759.33}{i}%
\TextAt{9}{115.587}{761.124}{\Character{58}}%
\Vector{q 0 G 0 g 0.478 w 50.971 -5.52 m 56.815 -5.52 l S Q}%
\TextAt{14}{122.971}{761.124}{o}%
\TextAt{9}{128.815}{761.124}{\Character{93}}%
\TextAt{8}{136.44}{761.124}{is}%
\TextAt{8}{148.68}{761.124}{the}%
\TextAt{8}{169.312}{761.124}{degree}%
\TextAt{8}{206.204}{761.124}{of}%
\TextAt{8}{220.01}{761.124}{the}%
\TextAt{8}{240.642}{761.124}{polynomial}%
\Vector{q 0 G 0 g 0.478 w 229.612 -6.055 m 235.487 -6.055 l S Q}%
\TextAt{13}{301.612}{761.124}{p}%
\TextAt{21}{307.487}{758.301}{i}%
\TextAt{9}{310.868}{761.124}{\Character{40}}%
\TextAt{13}{315.42}{761.124}{X}%
\TextAt{9}{326.075}{761.124}{\Character{41}}%
\TextAt{8}{330.628}{761.124}{,}%
\TextAt{8}{338.372}{761.124}{and}%
\TextAt{8}{361.611}{761.124}{we}%
\TextAt{8}{379.314}{761.124}{have}%
\TextAt{9}{406.785}{761.124}{\Character{40}0\Character{41}}%
\TextAt{9}{425.842}{761.124}{\Character{61}}%
\TextAt{20}{439.111}{770.09}{\Character{84}}%
\TextAt{21}{449.074}{766.935}{r}%
\TextAt{21}{449.074}{757.637}{i}%
\TextAt{17}{451.957}{757.637}{\Character{61}1}%
\Vector{q 0 G 0 g 0.478 w 393.268 -2.937 m 403.162 -2.937 l S Q}%
\TextAt{14}{465.268}{761.124}{P}%
\TextAt{21}{475.162}{768.626}{e}%
\TextAt{22}{479.09}{767.411}{i}%
\TextAt{21}{475.162}{758.168}{i}%
\TextAt{8}{482.75}{761.124}{.}%
\TextAt{8}{492.626}{761.124}{Now}%
\TextAt{8}{519.762}{761.124}{let}%
\TextAt{14}{79.226}{744.078}{P}%
\TextAt{21}{89.119}{742.284}{i}%
\TextAt{9}{95.821}{744.078}{\Character{61}}%
\TextAt{14}{108.247}{744.078}{p}%
\TextAt{16}{114.23}{744.078}{\Character{79}}%
\TextAt{9}{126.737}{744.078}{\Character{43}}%
\TextAt{13}{138.498}{744.078}{p}%
\TextAt{21}{144.373}{742.284}{i}%
\TextAt{9}{147.754}{744.078}{\Character{40}}%
\TextAt{13}{152.307}{744.078}{\Character{18}}%
\TextAt{9}{158.087}{744.078}{\Character{41}}%
\TextAt{16}{162.639}{744.078}{\Character{79}}%
\TextAt{8}{176.386}{744.078}{be}%
\TextAt{8}{192.311}{744.078}{the}%
\TextAt{8}{212.467}{744.078}{preimage}%
\TextAt{8}{262.547}{744.078}{of}%
\Vector{q 0 G 0 g 0.478 w 203.877 -19.983 m 213.771 -19.983 l S Q}%
\TextAt{14}{275.877}{744.078}{P}%
\TextAt{21}{285.771}{741.319}{i}%
\TextAt{8}{293.049}{744.078}{with}%
\TextAt{8}{319.709}{744.078}{respect}%
\TextAt{8}{359.759}{744.078}{to}%
\TextAt{8}{374.069}{744.078}{the}%
\TextAt{8}{394.226}{744.078}{canonical}%
\TextAt{8}{445.609}{744.078}{homomorphism}%
\TextAt{16}{270.728}{718.38}{\Character{79}}%
\NativeArrowAt{283.891}{718.38}{11.955}{19.26}{\longrightarrow}%
\TextAt{16}{306.462}{718.38}{\Character{79}}%
\TextAt{9}{316.33}{718.38}{\Character{47}}%
\TextAt{14}{322.183}{718.38}{p}%
\TextAt{16}{328.166}{718.38}{\Character{79}}%
\TextAt{13}{338.016}{718.38}{\Character{58}}%
\TextAt{8}{79.226}{692.683}{Then}%
\TextAt{14}{110.313}{692.683}{P}%
\TextAt{21}{120.206}{690.89}{i}%
\TextAt{8}{123.588}{692.683}{,}%
\TextAt{8}{131.399}{692.683}{for}%
\TextAt{13}{149.813}{692.683}{i}%
\TextAt{9}{158.03}{692.683}{\Character{61}}%
\TextAt{9}{171.348}{692.683}{1}%
\TextAt{13}{177.211}{692.683}{\Character{59}}%
\TextAt{13}{182.435}{692.683}{\Character{58}}%
\TextAt{13}{187.672}{692.683}{\Character{58}}%
\TextAt{13}{192.896}{692.683}{\Character{58}}%
\TextAt{13}{198.133}{692.683}{\Character{59}}%
\TextAt{13}{203.369}{692.683}{r}%
\TextAt{8}{209.032}{692.683}{,}%
\TextAt{8}{216.843}{692.683}{varies}%
\TextAt{8}{250.285}{692.683}{over}%
\TextAt{8}{275.861}{692.683}{the}%
\TextAt{8}{296.547}{692.683}{prime}%
\TextAt{8}{330.24}{692.683}{ideals}%
\TextAt{8}{363.347}{692.683}{of}%
\TextAt{16}{377.207}{692.683}{\Character{79}}%
\TextAt{8}{391.483}{692.683}{above}%
\TextAt{14}{425.188}{692.683}{p}%
\TextAt{8}{431.171}{692.683}{.}%
\TextAt{13}{441.208}{692.683}{f}%
\TextAt{21}{446.978}{690.89}{i}%
\TextAt{9}{454.583}{692.683}{\Character{61}}%
\TextAt{9}{467.901}{692.683}{\Character{91}}%
\TextAt{16}{471.162}{692.683}{\Character{79}}%
\TextAt{9}{481.012}{692.683}{\Character{47}}%
\TextAt{14}{486.865}{692.683}{P}%
\TextAt{21}{496.758}{690.89}{i}%
\TextAt{9}{504.363}{692.683}{\Character{58}}%
\TextAt{14}{511.838}{692.683}{o}%
\TextAt{9}{517.682}{692.683}{\Character{47}}%
\TextAt{14}{523.535}{692.683}{p}%
\TextAt{9}{529.518}{692.683}{\Character{93}}%
\TextAt{8}{79.226}{675.637}{is}%
\TextAt{8}{91.235}{675.637}{the}%
\TextAt{8}{111.636}{675.637}{degree}%
\TextAt{8}{148.297}{675.637}{of}%
\TextAt{8}{161.871}{675.637}{the}%
\TextAt{8}{182.273}{675.637}{polynomial}%
\Vector{q 0 G 0 g 0.478 w 171.011 -91.542 m 176.886 -91.542 l S Q}%
\TextAt{13}{243.011}{675.637}{p}%
\TextAt{21}{248.886}{672.814}{i}%
\TextAt{9}{252.267}{675.637}{\Character{40}}%
\TextAt{13}{256.82}{675.637}{X}%
\TextAt{9}{267.475}{675.637}{\Character{41}}%
\TextAt{8}{272.027}{675.637}{.}%
\TextAt{8}{281.21}{675.637}{Furthermore}%
\TextAt{14}{349.253}{675.637}{P}%
\TextAt{21}{359.146}{680.97}{e}%
\TextAt{22}{363.074}{679.755}{i}%
\TextAt{21}{359.146}{672.573}{i}%
\TextAt{8}{370.876}{675.637}{is}%
\TextAt{8}{382.885}{675.637}{the}%
\TextAt{8}{403.287}{675.637}{preimage}%
\TextAt{8}{453.612}{675.637}{of}%
\Vector{q 0 G 0 g 0.478 w 395.187 -88.423 m 405.08 -88.423 l S Q}%
\TextAt{14}{467.187}{675.637}{P}%
\TextAt{21}{477.08}{683.139}{e}%
\TextAt{22}{481.008}{681.924}{i}%
\TextAt{21}{477.08}{672.682}{i}%
\TextAt{8}{488.81}{675.637}{(because}%
\TextAt{13}{79.226}{658.591}{e}%
\TextAt{21}{84.652}{656.798}{i}%
\TextAt{9}{92.662}{658.591}{\Character{61}}%
\TextAt{8}{106.395}{658.591}{\char35\relax }%
\TextAt{16}{116.15}{658.591}{\Character{102}}%
\Vector{q 0 G 0 g 0.478 w 50.128 -105.469 m 60.021 -105.469 l S Q}%
\TextAt{14}{122.128}{658.591}{P}%
\TextAt{21}{132.021}{666.093}{\Character{23}}%
\TextAt{21}{132.021}{655.636}{i}%
\TextAt{16}{141.832}{658.591}{\Character{106}}%
\TextAt{13}{149.782}{658.591}{\Character{23}}%
\NativeMathAt{160.948}{658.591}{11.955}{\in}%
\TextAt{18}{173.547}{658.591}{N}%
\TextAt{16}{182.182}{658.591}{\Character{103}}%
\TextAt{8}{188.159}{658.591}{),}%
\TextAt{8}{200.823}{658.591}{and}%
\TextAt{14}{224.353}{658.591}{p}%
\TextAt{16}{230.336}{658.591}{\Character{79}}%
\NativeMathAt{244.802}{658.591}{11.955}{\supseteq}%
\TextAt{20}{258.742}{667.558}{\Character{84}}%
\TextAt{21}{268.704}{664.403}{r}%
\TextAt{21}{268.704}{655.104}{i}%
\TextAt{17}{271.587}{655.104}{\Character{61}1}%
\TextAt{14}{284.899}{658.591}{P}%
\TextAt{21}{294.792}{663.924}{e}%
\TextAt{22}{298.72}{662.709}{i}%
\TextAt{21}{294.792}{655.527}{i}%
\TextAt{8}{302.38}{658.591}{,}%
\TextAt{8}{310.489}{658.591}{so}%
\TextAt{8}{325.626}{658.591}{that}%
\TextAt{14}{351.763}{658.591}{p}%
\TextAt{16}{357.746}{658.591}{\Character{79}}%
\TextAt{16}{372.212}{658.591}{\Character{106}}%
\TextAt{20}{380.174}{667.558}{\Character{81}}%
\TextAt{21}{391.465}{664.403}{r}%
\TextAt{21}{391.465}{655.104}{i}%
\TextAt{17}{394.348}{655.104}{\Character{61}1}%
\TextAt{14}{407.66}{658.591}{P}%
\TextAt{21}{417.553}{663.924}{e}%
\TextAt{22}{421.481}{662.709}{i}%
\TextAt{21}{417.553}{655.527}{i}%
\TextAt{8}{429.806}{658.591}{and}%
\TextAt{8}{453.336}{658.591}{therefore}%
\TextAt{14}{503.203}{658.591}{p}%
\TextAt{16}{509.186}{658.591}{\Character{79}}%
\TextAt{9}{523.665}{658.591}{\Character{61}}%
\TextAt{20}{79.226}{650.512}{\Character{81}}%
\TextAt{21}{90.517}{647.357}{r}%
\TextAt{21}{90.517}{638.058}{i}%
\TextAt{17}{93.4}{638.058}{\Character{61}1}%
\TextAt{14}{106.712}{641.545}{P}%
\TextAt{21}{116.605}{646.878}{e}%
\TextAt{22}{120.533}{645.663}{i}%
\TextAt{21}{116.605}{638.481}{i}%
\TextAt{8}{128.091}{641.545}{because}%
\TextAt{20}{171.392}{650.512}{\Character{80}}%
\TextAt{13}{186.004}{641.545}{e}%
\TextAt{21}{191.429}{639.752}{i}%
\TextAt{13}{194.811}{641.545}{f}%
\TextAt{21}{200.581}{639.752}{i}%
\TextAt{9}{207.283}{641.545}{\Character{61}}%
\TextAt{13}{219.708}{641.545}{n}%
\TextAt{8}{226.696}{641.545}{.}%
\Vector{q 0 G 0 g 0.399 w 452.334 -132.455 m 452.334 -124.385 l S Q}%
\Vector{q 0 G 0 g 0.399 w 452.534 -124.584 m 459.707 -124.584 l S Q}%
\Vector{q 0 G 0 g 0.399 w 452.534 -132.256 m 459.707 -132.256 l S Q}%
\Vector{q 0 G 0 g 0.399 w 459.906 -132.455 m 459.906 -124.385 l S Q}%
\TextAt{8}{97.159}{613.488}{The}%
\TextAt{8}{121.087}{613.488}{prime}%
\TextAt{8}{154.124}{613.488}{ideal}%
\TextAt{14}{181.962}{613.488}{p}%
\TextAt{8}{191.716}{613.488}{is}%
\TextAt{8}{203.354}{613.488}{said}%
\TextAt{8}{227.354}{613.488}{to}%
\TextAt{12}{241.538}{613.488}{split}%
\TextAt{12}{271.377}{613.488}{completely}%
\TextAt{8}{338.63}{613.488}{(or}%
\TextAt{8}{357.37}{613.488}{to}%
\TextAt{8}{371.554}{613.488}{be}%
\TextAt{12}{387.352}{613.488}{totally}%
\TextAt{12}{430.019}{613.488}{split}%
\TextAt{8}{455.519}{613.488}{)}%
\TextAt{8}{463.846}{613.488}{in}%
\TextAt{13}{477.373}{613.488}{L}%
\TextAt{8}{485.337}{613.488}{,}%
\TextAt{8}{492.386}{613.488}{if}%
\TextAt{8}{502.984}{613.488}{in}%
\TextAt{8}{516.51}{613.488}{the}%
\TextAt{8}{79.226}{596.442}{decomposition}%
\TextAt{14}{267.693}{579.396}{p}%
\TextAt{9}{276.996}{579.396}{\Character{61}}%
\TextAt{14}{289.422}{579.396}{P}%
\TextAt{21}{299.315}{584.621}{e}%
\TextAt{24}{303.244}{583.514}{1}%
\TextAt{17}{299.315}{576.465}{1}%
\TextAt{16}{309.885}{579.396}{\Character{1}}%
\TextAt{16}{315.181}{579.396}{\Character{1}}%
\TextAt{16}{320.489}{579.396}{\Character{1}}%
\TextAt{14}{325.825}{579.396}{P}%
\TextAt{21}{335.719}{584.332}{e}%
\TextAt{22}{339.647}{583.336}{r}%
\TextAt{21}{335.719}{576.441}{r}%
\TextAt{8}{79.226}{557.29}{one}%
\TextAt{8}{100.861}{557.29}{has}%
\TextAt{13}{121.91}{557.29}{r}%
\TextAt{9}{131.131}{557.29}{\Character{61}}%
\TextAt{13}{143.855}{557.29}{n}%
\TextAt{9}{154.463}{557.29}{\Character{61}}%
\TextAt{9}{167.171}{557.29}{\Character{91}}%
\TextAt{13}{170.439}{557.29}{L}%
\TextAt{9}{182.024}{557.29}{\Character{58}}%
\TextAt{13}{188.895}{557.29}{K}%
\TextAt{9}{199.713}{557.29}{\Character{93}}%
\TextAt{8}{202.965}{557.29}{,}%
\TextAt{8}{210.333}{557.29}{so}%
\TextAt{8}{224.879}{557.29}{that}%
\TextAt{13}{250.423}{557.29}{e}%
\TextAt{21}{255.849}{555.497}{i}%
\TextAt{9}{262.85}{557.29}{\Character{61}}%
\TextAt{13}{275.574}{557.29}{f}%
\TextAt{21}{281.344}{555.497}{i}%
\TextAt{9}{288.346}{557.29}{\Character{61}}%
\TextAt{9}{301.054}{557.29}{1}%
\TextAt{8}{310.996}{557.29}{for}%
\TextAt{8}{329.057}{557.29}{all}%
\TextAt{13}{345.491}{557.29}{i}%
\TextAt{9}{353.105}{557.29}{\Character{61}}%
\TextAt{9}{365.813}{557.29}{1}%
\TextAt{13}{371.682}{557.29}{\Character{59}}%
\TextAt{13}{376.906}{557.29}{\Character{58}}%
\TextAt{13}{382.143}{557.29}{\Character{58}}%
\TextAt{13}{387.379}{557.29}{\Character{58}}%
\TextAt{13}{392.604}{557.29}{\Character{59}}%
\TextAt{13}{397.84}{557.29}{r}%
\TextAt{8}{403.503}{557.29}{.}%
\TextAt{14}{412.478}{557.29}{p}%
\TextAt{8}{422.534}{557.29}{is}%
\TextAt{8}{434.474}{557.29}{called}%
\TextAt{12}{467.813}{557.29}{nonsplit}%
\TextAt{8}{514.988}{557.29}{,}%
\TextAt{8}{522.357}{557.29}{or}%
\TextAt{12}{79.226}{540.244}{indecomposed}%
\TextAt{8}{161.454}{540.244}{,}%
\TextAt{8}{168.807}{540.244}{if}%
\TextAt{13}{179.694}{540.244}{r}%
\TextAt{9}{188.894}{540.244}{\Character{61}}%
\TextAt{9}{201.59}{540.244}{1}%
\TextAt{8}{207.451}{540.244}{,}%
\TextAt{8}{214.805}{540.244}{i.e.,}%
\TextAt{8}{237.114}{540.244}{if}%
\TextAt{8}{248.001}{540.244}{there}%
\TextAt{8}{278.076}{540.244}{is}%
\TextAt{8}{290.004}{540.244}{only}%
\TextAt{8}{315.859}{540.244}{a}%
\TextAt{8}{325.778}{540.244}{single}%
\TextAt{8}{358.519}{540.244}{prime}%
\TextAt{8}{391.846}{540.244}{ideal}%
\TextAt{8}{419.972}{540.244}{of}%
\TextAt{13}{433.466}{540.244}{L}%
\TextAt{8}{445.491}{540.244}{over}%
\TextAt{14}{470.701}{540.244}{p}%
\TextAt{8}{476.684}{540.244}{.}%
\TextAt{8}{485.622}{540.244}{From}%
\TextAt{8}{516.51}{540.244}{the}%
\TextAt{8}{79.226}{523.198}{fundamental}%
\TextAt{8}{146.868}{523.198}{identity}%
\TextAt{21}{282.627}{515.136}{r}%
\TextAt{20}{276.021}{511.55}{\Character{88}}%
\TextAt{21}{277.804}{486.355}{i}%
\TextAt{17}{280.687}{486.355}{\Character{61}1}%
\TextAt{13}{295.282}{500.192}{e}%
\TextAt{21}{300.708}{498.399}{i}%
\TextAt{13}{304.089}{500.192}{f}%
\TextAt{21}{309.859}{498.399}{i}%
\TextAt{9}{316.561}{500.192}{\Character{61}}%
\TextAt{13}{328.987}{500.192}{n}%
\TextAt{8}{79.226}{470.687}{we}%
\TextAt{8}{96.916}{470.687}{now}%
\TextAt{8}{121.767}{470.687}{understand}%
\TextAt{8}{183.427}{470.687}{the}%
\TextAt{8}{204.046}{470.687}{name}%
\TextAt{8}{235.723}{470.687}{of}%
\TextAt{8}{249.516}{470.687}{inertia}%
\TextAt{8}{287.051}{470.687}{degree:}%
\TextAt{8}{328.942}{470.687}{the}%
\TextAt{8}{349.56}{470.687}{smaller}%
\TextAt{8}{390.407}{470.687}{this}%
\TextAt{8}{413.692}{470.687}{degree}%
\TextAt{8}{450.57}{470.687}{is,}%
\TextAt{8}{466.163}{470.687}{the}%
\TextAt{8}{486.782}{470.687}{more}%
\TextAt{8}{516.51}{470.687}{the}%
\TextAt{8}{79.226}{453.641}{ideal}%
\TextAt{14}{107.189}{453.641}{p}%
\TextAt{8}{117.07}{453.641}{will}%
\TextAt{8}{139.175}{453.641}{tend}%
\TextAt{8}{165.835}{453.641}{to}%
\TextAt{8}{180.145}{453.641}{factor}%
\TextAt{8}{213.644}{453.641}{into}%
\TextAt{8}{237.387}{453.641}{different}%
\TextAt{8}{283.557}{453.641}{prime}%
\TextAt{8}{316.721}{453.641}{ideals.}%
\TextAt{8}{97.159}{433.606}{The}%
\TextAt{8}{121.796}{433.606}{prime}%
\TextAt{8}{155.543}{433.606}{ideal}%
\TextAt{14}{184.09}{433.606}{P}%
\TextAt{21}{193.983}{431.813}{i}%
\TextAt{8}{201.845}{433.606}{in}%
\TextAt{8}{216.081}{433.606}{the}%
\TextAt{8}{236.821}{433.606}{decomposition}%
\TextAt{14}{314.539}{433.606}{p}%
\TextAt{9}{324.837}{433.606}{\Character{61}}%
\TextAt{20}{338.257}{442.573}{\Character{81}}%
\TextAt{21}{349.548}{439.418}{r}%
\TextAt{21}{349.548}{430.119}{i}%
\TextAt{17}{352.431}{430.119}{\Character{61}1}%
\TextAt{14}{365.743}{433.606}{P}%
\TextAt{21}{375.636}{438.939}{e}%
\TextAt{22}{379.564}{437.724}{i}%
\TextAt{21}{375.636}{430.543}{i}%
\TextAt{8}{387.705}{433.606}{is}%
\TextAt{8}{400.052}{433.606}{called}%
\TextAt{12}{433.799}{433.606}{unramified}%
\TextAt{8}{501.296}{433.606}{over}%
\TextAt{14}{526.925}{433.606}{o}%
\TextAt{8}{79.226}{416.56}{(or}%
\TextAt{8}{98.567}{416.56}{over}%
\TextAt{13}{124.089}{416.56}{K}%
\TextAt{8}{134.907}{416.56}{)}%
\TextAt{8}{143.835}{416.56}{if}%
\TextAt{13}{155.035}{416.56}{e}%
\TextAt{21}{160.46}{414.767}{i}%
\TextAt{9}{167.974}{416.56}{\Character{61}}%
\TextAt{9}{181.197}{416.56}{1}%
\TextAt{8}{191.436}{416.56}{and}%
\TextAt{8}{214.675}{416.56}{if}%
\TextAt{8}{225.875}{416.56}{the}%
\TextAt{8}{246.507}{416.56}{residue}%
\TextAt{8}{286.71}{416.56}{class}%
\TextAt{8}{314.623}{416.56}{field}%
\TextAt{8}{340.456}{416.56}{extension}%
\TextAt{16}{392.698}{416.56}{\Character{79}}%
\TextAt{9}{402.547}{416.56}{\Character{47}}%
\TextAt{14}{408.4}{416.56}{P}%
\TextAt{21}{418.294}{414.767}{i}%
\TextAt{16}{425.807}{416.56}{\Character{106}}%
\TextAt{14}{433.26}{416.56}{o}%
\TextAt{9}{439.104}{416.56}{\Character{47}}%
\TextAt{14}{444.957}{416.56}{p}%
\TextAt{8}{455.313}{416.56}{is}%
\TextAt{8}{467.553}{416.56}{separable.}%
\TextAt{8}{524.975}{416.56}{If}%
\TextAt{8}{79.226}{399.515}{not,}%
\TextAt{8}{103.565}{399.515}{it}%
\TextAt{8}{115.487}{399.515}{is}%
\TextAt{8}{127.47}{399.515}{called}%
\TextAt{12}{160.852}{399.515}{ramified}%
\TextAt{8}{208.924}{399.515}{,}%
\TextAt{8}{216.346}{399.515}{and}%
\TextAt{12}{239.327}{399.515}{totally}%
\TextAt{12}{282.389}{399.515}{ramified}%
\TextAt{8}{334.577}{399.515}{if}%
\TextAt{8}{345.519}{399.515}{furthermore}%
\TextAt{13}{410.451}{399.515}{f}%
\TextAt{21}{416.221}{397.721}{i}%
\TextAt{9}{423.295}{399.515}{\Character{61}}%
\TextAt{9}{436.075}{399.515}{1}%
\TextAt{8}{441.946}{399.515}{.}%
\TextAt{8}{451.049}{399.515}{The}%
\TextAt{8}{475.322}{399.515}{prime}%
\TextAt{8}{508.704}{399.515}{ideal}%
\TextAt{14}{79.226}{382.469}{p}%
\TextAt{8}{90.397}{382.469}{is}%
\TextAt{8}{103.452}{382.469}{called}%
\TextAt{12}{137.907}{382.469}{unramified}%
\TextAt{8}{206.11}{382.469}{if}%
\TextAt{8}{218.125}{382.469}{all}%
\TextAt{14}{235.675}{382.469}{P}%
\TextAt{21}{245.568}{380.675}{i}%
\TextAt{8}{254.138}{382.469}{are}%
\TextAt{8}{274.939}{382.469}{unramified,}%
\TextAt{8}{338.337}{382.469}{otherwise}%
\TextAt{8}{391.717}{382.469}{it}%
\TextAt{8}{404.711}{382.469}{is}%
\TextAt{8}{417.766}{382.469}{called}%
\TextAt{12}{452.221}{382.469}{ramified}%
\TextAt{8}{500.292}{382.469}{.}%
\TextAt{8}{512.613}{382.469}{The}%
\TextAt{8}{79.226}{365.423}{extension}%
\TextAt{13}{130.704}{365.423}{L}%
\TextAt{16}{141.989}{365.423}{\Character{106}}%
\TextAt{13}{148.631}{365.423}{K}%
\TextAt{8}{163.057}{365.423}{itself}%
\TextAt{8}{191.115}{365.423}{is}%
\TextAt{8}{202.59}{365.423}{called}%
\TextAt{12}{235.466}{365.423}{unramified}%
\TextAt{8}{302.09}{365.423}{if}%
\TextAt{8}{312.526}{365.423}{all}%
\TextAt{8}{328.497}{365.423}{prime}%
\TextAt{8}{361.372}{365.423}{ideals}%
\TextAt{14}{393.661}{365.423}{p}%
\TextAt{8}{403.254}{365.423}{of}%
\TextAt{13}{416.295}{365.423}{K}%
\TextAt{8}{430.722}{365.423}{are}%
\TextAt{8}{449.945}{365.423}{unramified}%
\TextAt{8}{508.189}{365.423}{in}%
\TextAt{13}{521.553}{365.423}{L}%
\TextAt{8}{529.518}{365.423}{.}%
\TextAt{8}{97.159}{345.388}{The}%
\TextAt{8}{121.478}{345.388}{case}%
\TextAt{8}{146.514}{345.388}{where}%
\TextAt{8}{180.588}{345.388}{a}%
\TextAt{8}{190.608}{345.388}{prime}%
\TextAt{8}{224.036}{345.388}{ideal}%
\TextAt{14}{252.264}{345.388}{p}%
\TextAt{8}{262.41}{345.388}{of}%
\TextAt{13}{276.005}{345.388}{K}%
\TextAt{8}{290.985}{345.388}{is}%
\TextAt{8}{303.013}{345.388}{ramified}%
\TextAt{8}{348.804}{345.388}{in}%
\TextAt{13}{362.721}{345.388}{L}%
\TextAt{8}{374.848}{345.388}{is}%
\TextAt{8}{386.877}{345.388}{an}%
\TextAt{8}{403.401}{345.388}{exceptional}%
\TextAt{8}{465.127}{345.388}{phenomenon.}%
\TextAt{8}{79.226}{328.342}{In}%
\TextAt{8}{93.847}{328.342}{fact,}%
\TextAt{8}{120.185}{328.342}{we}%
\TextAt{8}{137.412}{328.342}{have}%
\TextAt{8}{164.407}{328.342}{the}%
\TextAt{12}{79.226}{303.273}{Proposition}%
\TextAt{12}{151.928}{303.273}{(8.4).}%
\TextAt{8}{187.234}{303.273}{If}%
\TextAt{13}{198.918}{303.273}{L}%
\TextAt{16}{210.203}{303.273}{\Character{106}}%
\TextAt{13}{216.845}{303.273}{K}%
\TextAt{8}{231.551}{303.273}{is}%
\TextAt{8}{243.307}{303.273}{separable,}%
\TextAt{8}{297.995}{303.273}{then}%
\TextAt{8}{324.646}{303.273}{there}%
\TextAt{8}{354.549}{303.273}{are}%
\TextAt{8}{374.051}{303.273}{only}%
\TextAt{8}{399.734}{303.273}{finitely}%
\TextAt{8}{439.07}{303.273}{many}%
\TextAt{8}{470.934}{303.273}{prime}%
\TextAt{8}{504.089}{303.273}{ideals}%
\TextAt{8}{79.226}{286.227}{of}%
\TextAt{13}{92.556}{286.227}{K}%
\TextAt{8}{107.271}{286.227}{which}%
\TextAt{8}{140.758}{286.227}{are}%
\TextAt{8}{160.268}{286.227}{ramified}%
\TextAt{8}{205.794}{286.227}{in}%
\TextAt{13}{219.447}{286.227}{L}%
\TextAt{8}{227.411}{286.227}{.}%
\TextAt{19}{79.226}{253.137}{Proof.}%
\TextAt{8}{116}{253.137}{Let}%
\TextAt{13}{136.659}{253.137}{\Character{18}}%
\NativeMathAt{145.76}{253.137}{11.955}{\in}%
\TextAt{16}{157.034}{253.137}{\Character{79}}%
\TextAt{8}{170.487}{253.137}{be}%
\TextAt{8}{186.1}{253.137}{a}%
\TextAt{8}{195.545}{253.137}{primitive}%
\TextAt{8}{245.314}{253.137}{element}%
\TextAt{8}{288.245}{253.137}{for}%
\TextAt{13}{305.819}{253.137}{L}%
\TextAt{8}{313.783}{253.137}{,}%
\TextAt{8}{320.684}{253.137}{i.e.,}%
\TextAt{13}{342.54}{253.137}{L}%
\TextAt{9}{353.825}{253.137}{\Character{61}}%
\TextAt{13}{366.251}{253.137}{K}%
\TextAt{9}{377.068}{253.137}{\Character{40}}%
\TextAt{13}{381.621}{253.137}{\Character{18}}%
\TextAt{9}{387.401}{253.137}{\Character{41}}%
\TextAt{8}{391.953}{253.137}{,}%
\TextAt{8}{398.854}{253.137}{and}%
\TextAt{8}{421.305}{253.137}{let}%
\TextAt{13}{437.899}{253.137}{p}%
\TextAt{9}{443.774}{253.137}{\Character{40}}%
\TextAt{13}{448.327}{253.137}{X}%
\TextAt{9}{458.982}{253.137}{\Character{41}}%
\NativeMathAt{466.855}{253.137}{11.955}{\in}%
\TextAt{14}{478.146}{253.137}{o}%
\TextAt{9}{483.99}{253.137}{\Character{91}}%
\TextAt{13}{487.242}{253.137}{X}%
\TextAt{9}{497.897}{253.137}{\Character{93}}%
\TextAt{8}{504.735}{253.137}{be}%
\TextAt{8}{520.348}{253.137}{its}%
\TextAt{8}{79.226}{236.091}{minimal}%
\TextAt{8}{124.751}{236.091}{polynomial.}%
\TextAt{8}{189.796}{236.091}{Let}%
\TextAt{13}{205.521}{210.393}{d}%
\TextAt{9}{214.925}{210.393}{\Character{61}}%
\TextAt{13}{227.35}{210.393}{d}%
\TextAt{9}{233.433}{210.393}{\Character{40}1}%
\TextAt{13}{243.838}{210.393}{\Character{59}}%
\TextAt{13}{249.062}{210.393}{\Character{18}}%
\TextAt{13}{254.837}{210.393}{\Character{59}}%
\TextAt{13}{260.073}{210.393}{\Character{58}}%
\TextAt{13}{265.31}{210.393}{\Character{58}}%
\TextAt{13}{270.534}{210.393}{\Character{58}}%
\TextAt{13}{275.77}{210.393}{\Character{59}}%
\TextAt{13}{281.007}{210.393}{\Character{18}}%
\TextAt{21}{286.864}{215.33}{n}%
\TextAt{3}{292.002}{215.33}{\Character{0}}%
\TextAt{17}{298.589}{215.33}{1}%
\TextAt{9}{303.321}{210.393}{\Character{41}}%
\TextAt{9}{311.176}{210.393}{\Character{61}}%
\TextAt{20}{323.62}{221.751}{\Character{89}}%
\TextAt{21}{324.581}{196.556}{i\Character{60}j}%
\TextAt{9}{338.896}{210.393}{\Character{40}}%
\TextAt{13}{343.448}{210.393}{\Character{18}}%
\TextAt{21}{348.903}{208.6}{i}%
\TextAt{16}{354.941}{210.393}{\Character{0}}%
\TextAt{13}{366.896}{210.393}{\Character{18}}%
\TextAt{21}{372.351}{208.6}{j}%
\TextAt{9}{376.734}{210.393}{\Character{41}}%
\TextAt{17}{381.286}{215.33}{2}%
\NativeMathAt{389.339}{210.393}{11.955}{\in}%
\TextAt{14}{400.63}{210.393}{o}%
\TextAt{8}{79.226}{175.197}{be}%
\TextAt{8}{96.441}{175.197}{the}%
\TextAt{8}{117.888}{175.197}{discriminant}%
\TextAt{8}{186.557}{175.197}{of}%
\TextAt{13}{201.178}{175.197}{p}%
\TextAt{9}{207.053}{175.197}{\Character{40}}%
\TextAt{13}{211.605}{175.197}{X}%
\TextAt{9}{222.26}{175.197}{\Character{41}}%
\TextAt{8}{232}{175.197}{(see}%
\TextAt{8}{256.759}{175.197}{§2,}%
\TextAt{8}{277.058}{175.197}{p.}%
\TextAt{8}{292.001}{175.197}{11).}%
\TextAt{8}{320.591}{175.197}{Then}%
\TextAt{8}{352.439}{175.197}{every}%
\TextAt{8}{384.622}{175.197}{prime}%
\TextAt{8}{419.076}{175.197}{ideal}%
\TextAt{14}{448.329}{175.197}{p}%
\TextAt{8}{459.5}{175.197}{of}%
\TextAt{13}{474.121}{175.197}{K}%
\TextAt{8}{490.126}{175.197}{which}%
\TextAt{8}{524.903}{175.197}{is}%
\TextAt{8}{79.226}{158.151}{relatively}%
\TextAt{8}{131.778}{158.151}{prime}%
\TextAt{8}{166.432}{158.151}{to}%
\TextAt{13}{182.234}{158.151}{d}%
\TextAt{8}{193.705}{158.151}{and}%
\TextAt{8}{217.958}{158.151}{to}%
\TextAt{8}{233.76}{158.151}{the}%
\TextAt{8}{255.407}{158.151}{conductor}%
\TextAt{14}{311.533}{158.151}{F}%
\TextAt{8}{324.228}{158.151}{of}%
\TextAt{14}{339.049}{158.151}{o}%
\TextAt{9}{344.894}{158.151}{\Character{91}}%
\TextAt{13}{348.145}{158.151}{\Character{18}}%
\TextAt{9}{353.926}{158.151}{\Character{93}}%
\TextAt{8}{362.566}{158.151}{is}%
\TextAt{8}{375.82}{158.151}{unramified.}%
\TextAt{8}{443.377}{158.151}{In}%
\TextAt{8}{459.489}{158.151}{fact,}%
\TextAt{8}{487.69}{158.151}{by}%
\TextAt{8}{505.44}{158.151}{(8.3),}%
\TextAt{8}{79.226}{141.105}{the}%
\TextAt{8}{100.763}{141.105}{ramification}%
\TextAt{8}{167.191}{141.105}{indices}%
\TextAt{13}{206.995}{141.105}{e}%
\TextAt{21}{212.421}{139.312}{i}%
\TextAt{8}{221.08}{141.105}{equal}%
\TextAt{9}{253.352}{141.105}{1}%
\TextAt{8}{264.483}{141.105}{as}%
\TextAt{8}{280.233}{141.105}{soon}%
\TextAt{8}{308.668}{141.105}{as}%
\TextAt{8}{324.418}{141.105}{they}%
\TextAt{8}{352.135}{141.105}{are}%
\TextAt{8}{373.027}{141.105}{equal}%
\TextAt{8}{405.299}{141.105}{to}%
\TextAt{9}{420.989}{141.105}{1}%
\TextAt{8}{432.12}{141.105}{in}%
\TextAt{8}{447.153}{141.105}{the}%
\TextAt{8}{468.69}{141.105}{factorization}%
\TextAt{8}{79.226}{124.059}{of}%
\Vector{q 0 G 0 g 0.478 w 21.669 -643.12 m 27.544 -643.12 l S Q}%
\TextAt{13}{93.669}{124.059}{p}%
\TextAt{9}{99.544}{124.059}{\Character{40}}%
\TextAt{13}{104.096}{124.059}{X}%
\TextAt{9}{114.751}{124.059}{\Character{41}}%
\TextAt{9}{124.506}{124.059}{\Character{61}}%
\TextAt{13}{138.842}{124.059}{p}%
\TextAt{9}{144.717}{124.059}{\Character{40}}%
\TextAt{13}{149.27}{124.059}{X}%
\TextAt{9}{159.925}{124.059}{\Character{41}}%
\TextAt{8}{167.798}{124.059}{mod}%
\TextAt{14}{193.236}{124.059}{p}%
\TextAt{8}{204.229}{124.059}{in}%
\TextAt{14}{218.995}{124.059}{o}%
\TextAt{9}{224.839}{124.059}{\Character{47}}%
\TextAt{14}{230.692}{124.059}{p}%
\TextAt{8}{236.675}{124.059}{,}%
\TextAt{8}{245.215}{124.059}{so}%
\TextAt{8}{260.698}{124.059}{certainly}%
\TextAt{8}{310.265}{124.059}{if}%
\Vector{q 0 G 0 g 0.478 w 250.101 -643.12 m 255.976 -643.12 l S Q}%
\TextAt{13}{322.101}{124.059}{p}%
\TextAt{9}{327.976}{124.059}{\Character{40}}%
\TextAt{13}{332.528}{124.059}{X}%
\TextAt{9}{343.184}{124.059}{\Character{41}}%
\TextAt{8}{352.746}{124.059}{has}%
\TextAt{8}{374.732}{124.059}{no}%
\TextAt{8}{392.104}{124.059}{multiple}%
\TextAt{8}{439.065}{124.059}{roots.}%
\TextAt{8}{476.614}{124.059}{But}%
\TextAt{8}{500.968}{124.059}{this}%
\TextAt{8}{524.903}{124.059}{is}%
\TextAt{8}{79.226}{107.013}{the}%
\TextAt{8}{99.949}{107.013}{case}%
\TextAt{8}{125.287}{107.013}{since}%
\TextAt{8}{154.522}{107.013}{the}%
\TextAt{8}{175.245}{107.013}{discriminant}%
\Vector{q 0 G 0 g 0.478 w 171.191 -657.011 m 177.273 -657.011 l S Q}%
\TextAt{13}{243.191}{107.013}{d}%
\TextAt{9}{253.56}{107.013}{\Character{61}}%
\TextAt{13}{266.951}{107.013}{d}%
\TextAt{8}{276.354}{107.013}{mod}%
\TextAt{14}{301.792}{107.013}{p}%
\TextAt{8}{312.239}{107.013}{of}%
\Vector{q 0 G 0 g 0.478 w 254.136 -660.166 m 260.011 -660.166 l S Q}%
\TextAt{13}{326.136}{107.013}{p}%
\TextAt{9}{332.011}{107.013}{\Character{40}}%
\TextAt{13}{336.563}{107.013}{X}%
\TextAt{9}{347.218}{107.013}{\Character{41}}%
\TextAt{8}{356.235}{107.013}{is}%
\TextAt{8}{368.565}{107.013}{nonzero.}%
\TextAt{8}{418.392}{107.013}{The}%
\TextAt{8}{443.013}{107.013}{residue}%
\TextAt{8}{483.306}{107.013}{class}%
\TextAt{8}{511.31}{107.013}{field}%
\TextAt{8}{79.226}{89.967}{extensions}%
\TextAt{16}{135.909}{89.967}{\Character{79}}%
\TextAt{9}{145.759}{89.967}{\Character{47}}%
\TextAt{14}{151.612}{89.967}{P}%
\TextAt{21}{161.505}{88.174}{i}%
\TextAt{16}{168.722}{89.967}{\Character{106}}%
\TextAt{14}{175.879}{89.967}{o}%
\TextAt{9}{181.723}{89.967}{\Character{47}}%
\TextAt{14}{187.576}{89.967}{p}%
\TextAt{8}{197.759}{89.967}{are}%
\TextAt{8}{217.572}{89.967}{generated}%
\TextAt{8}{271.207}{89.967}{by}%
\Vector{q 0 G 0 g 0.478 w 215.768 -674.057 m 221.548 -674.057 l S Q}%
\TextAt{13}{287.768}{89.967}{\Character{18}}%
\TextAt{9}{297.384}{89.967}{\Character{61}}%
\TextAt{13}{310.325}{89.967}{\Character{18}}%
\TextAt{8}{319.426}{89.967}{mod}%
\TextAt{14}{344.864}{89.967}{P}%
\TextAt{21}{354.757}{88.174}{i}%
\TextAt{8}{362.338}{89.967}{and}%
\TextAt{8}{385.403}{89.967}{therefore}%
\TextAt{8}{434.806}{89.967}{separable.}%
\TextAt{8}{491.707}{89.967}{Hence}%
\TextAt{14}{526.786}{89.967}{p}%
\TextAt{8}{79.226}{72.921}{is}%
\TextAt{8}{90.99}{72.921}{unramified.}%
\Vector{q 0 G 0 g 0.399 w 452.334 -701.079 m 452.334 -693.009 l S Q}%
\Vector{q 0 G 0 g 0.399 w 452.534 -693.208 m 459.707 -693.208 l S Q}%
\Vector{q 0 G 0 g 0.399 w 452.534 -700.879 m 459.707 -700.879 l S Q}%
\Vector{q 0 G 0 g 0.399 w 459.906 -701.079 m 459.906 -693.009 l S Q}%
\TextAt{8}{300.14}{43.033}{46}%
\Destination{78.23}{809.815}{nav.page.46}%
\Destination{79.36}{628.455}{nav.print.49}%
\Destination{79.36}{317.255}{nav.number.1.8.4}%
\Hotlink{261.938}{172.955}{6.858}{11.047}{nav.section.1.2}%
\Hotlink{291.501}{172.955}{12.716}{11.047}{nav.print.11}%
\Hotlink{504.94}{155.955}{25.078}{11.047}{nav.number.1.8.3}%
}

% Page 59 (body page 47)
\Page{%
\TextAt{2}{79.226}{796.943}{I.}%
\TextAt{2}{89.83}{796.943}{Algebraic}%
\TextAt{2}{138.648}{796.943}{Integers}%
\Vector{q 0 G 0 g 0.299 w 7.226 18.729 m 460.769 18.729 l S Q}%
\TextAt{8}{97.159}{761.124}{The}%
\TextAt{8}{122.603}{761.124}{precise}%
\TextAt{8}{162.418}{761.124}{description}%
\TextAt{8}{223.704}{761.124}{of}%
\TextAt{8}{238.425}{761.124}{the}%
\TextAt{8}{259.972}{761.124}{ramified}%
\TextAt{8}{306.888}{761.124}{prime}%
\TextAt{8}{341.443}{761.124}{ideals}%
\TextAt{8}{375.412}{761.124}{is}%
\TextAt{8}{388.566}{761.124}{given}%
\TextAt{8}{420.526}{761.124}{by}%
\TextAt{8}{438.176}{761.124}{the}%
\TextAt{12}{459.723}{761.124}{discriminant}%
\TextAt{8}{79.226}{744.078}{of}%
\TextAt{16}{93.76}{744.078}{\Character{79}}%
\TextAt{16}{108.979}{744.078}{\Character{106}}%
\TextAt{14}{117.676}{744.078}{o}%
\TextAt{8}{123.521}{744.078}{.}%
\TextAt{8}{135.581}{744.078}{It}%
\TextAt{8}{149.458}{744.078}{is}%
\TextAt{8}{162.426}{744.078}{defined}%
\TextAt{8}{203.943}{744.078}{to}%
\TextAt{8}{219.457}{744.078}{be}%
\TextAt{8}{236.585}{744.078}{the}%
\TextAt{8}{257.946}{744.078}{ideal}%
\TextAt{14}{287.113}{744.078}{d}%
\TextAt{8}{298.175}{744.078}{of}%
\TextAt{14}{312.709}{744.078}{o}%
\TextAt{8}{323.655}{744.078}{which}%
\TextAt{8}{358.345}{744.078}{is}%
\TextAt{8}{371.313}{744.078}{generated}%
\TextAt{8}{425.849}{744.078}{by}%
\TextAt{8}{443.312}{744.078}{the}%
\TextAt{8}{464.673}{744.078}{discriminants}%
\TextAt{13}{79.226}{727.032}{d}%
\TextAt{9}{85.309}{727.032}{\Character{40}}%
\TextAt{13}{89.861}{727.032}{\Character{33}}%
\TextAt{17}{97.15}{725.238}{1}%
\TextAt{13}{101.883}{727.032}{\Character{59}}%
\TextAt{13}{107.107}{727.032}{\Character{58}}%
\TextAt{13}{112.344}{727.032}{\Character{58}}%
\TextAt{13}{117.568}{727.032}{\Character{58}}%
\TextAt{13}{122.805}{727.032}{\Character{59}}%
\TextAt{13}{128.041}{727.032}{\Character{33}}%
\TextAt{21}{135.393}{725.238}{n}%
\TextAt{9}{141.029}{727.032}{\Character{41}}%
\TextAt{8}{149.032}{727.032}{of}%
\TextAt{8}{161.916}{727.032}{all}%
\TextAt{8}{177.729}{727.032}{bases}%
\TextAt{13}{207.971}{727.032}{\Character{33}}%
\TextAt{17}{215.26}{725.238}{1}%
\TextAt{13}{219.993}{727.032}{\Character{59}}%
\TextAt{13}{225.217}{727.032}{\Character{58}}%
\TextAt{13}{230.454}{727.032}{\Character{58}}%
\TextAt{13}{235.678}{727.032}{\Character{58}}%
\TextAt{13}{240.915}{727.032}{\Character{59}}%
\TextAt{13}{246.151}{727.032}{\Character{33}}%
\TextAt{21}{253.503}{725.238}{n}%
\TextAt{8}{262.59}{727.032}{of}%
\TextAt{13}{275.473}{727.032}{L}%
\TextAt{16}{286.759}{727.032}{\Character{106}}%
\TextAt{13}{293.401}{727.032}{K}%
\TextAt{8}{307.669}{727.032}{contained}%
\TextAt{8}{360.232}{727.032}{in}%
\TextAt{16}{373.438}{727.032}{\Character{79}}%
\TextAt{8}{383.288}{727.032}{.}%
\TextAt{8}{391.588}{727.032}{We}%
\TextAt{8}{411.286}{727.032}{will}%
\TextAt{8}{432.944}{727.032}{show}%
\TextAt{8}{461.501}{727.032}{in}%
\TextAt{8}{474.707}{727.032}{chapter}%
\TextAt{8}{516.212}{727.032}{III,}%
\TextAt{8}{79.226}{709.986}{§2}%
\TextAt{8}{94.66}{709.986}{that}%
\TextAt{8}{120.029}{709.986}{the}%
\TextAt{8}{140.186}{709.986}{prime}%
\TextAt{8}{173.349}{709.986}{divisors}%
\TextAt{8}{216.077}{709.986}{of}%
\TextAt{14}{229.407}{709.986}{d}%
\TextAt{8}{239.265}{709.986}{are}%
\TextAt{8}{258.776}{709.986}{exactly}%
\TextAt{8}{299.1}{709.986}{the}%
\TextAt{8}{319.257}{709.986}{prime}%
\TextAt{8}{352.421}{709.986}{ideals}%
\TextAt{8}{384.998}{709.986}{which}%
\TextAt{8}{418.485}{709.986}{ramify}%
\TextAt{8}{455.558}{709.986}{in}%
\TextAt{13}{469.211}{709.986}{L}%
\TextAt{8}{477.175}{709.986}{.}%
\TextAt{12}{97.159}{689.951}{Example:}%
\TextAt{8}{160.513}{689.951}{The}%
\TextAt{8}{186.263}{689.951}{law}%
\TextAt{8}{209.095}{689.951}{of}%
\TextAt{8}{224.121}{689.951}{decomposition}%
\TextAt{8}{302.952}{689.951}{of}%
\TextAt{8}{317.977}{689.951}{prime}%
\TextAt{8}{352.837}{689.951}{numbers}%
\TextAt{13}{401.743}{689.951}{p}%
\TextAt{8}{413.212}{689.951}{in}%
\TextAt{8}{428.56}{689.951}{a}%
\TextAt{8}{440.011}{689.951}{quadratic}%
\TextAt{8}{494.071}{689.951}{number}%
\TextAt{8}{79.226}{672.905}{field}%
\TextAt{18}{104.394}{672.905}{Q}%
\TextAt{9}{113.692}{672.905}{\Character{40}}%
\TextAt{16}{118.245}{681.516}{\Character{112}}%
\Vector{q 0 G 0 g 0.478 w 56.207 -92.245 m 62.352 -92.245 l S Q}%
\TextAt{13}{128.207}{672.905}{a}%
\TextAt{9}{134.352}{672.905}{\Character{41}}%
\TextAt{8}{142.613}{672.905}{is}%
\TextAt{8}{154.188}{672.905}{intimately}%
\TextAt{8}{209.937}{672.905}{related}%
\TextAt{8}{248.769}{672.905}{to}%
\TextAt{8}{262.89}{672.905}{Gauss’s}%
\TextAt{8}{305.238}{672.905}{famous}%
\TextAt{8}{345.11}{672.905}{quadratic}%
\TextAt{8}{397.285}{672.905}{reciprocity}%
\TextAt{8}{455.305}{672.905}{law.}%
\TextAt{8}{480.93}{672.905}{The}%
\TextAt{8}{504.794}{672.905}{latter}%
\TextAt{8}{79.226}{655.859}{concerns}%
\TextAt{8}{126.76}{655.859}{the}%
\TextAt{8}{146.916}{655.859}{problem}%
\TextAt{8}{192.442}{655.859}{of}%
\TextAt{8}{205.772}{655.859}{integer}%
\TextAt{8}{244.471}{655.859}{solutions}%
\TextAt{8}{293.379}{655.859}{of}%
\TextAt{8}{306.709}{655.859}{the}%
\TextAt{8}{326.866}{655.859}{equation}%
\TextAt{13}{236.365}{629.894}{x}%
\TextAt{17}{243.017}{634.83}{2}%
\TextAt{9}{250.406}{629.894}{\Character{43}}%
\TextAt{13}{262.167}{629.894}{by}%
\TextAt{9}{276.602}{629.894}{\Character{61}}%
\TextAt{13}{289.027}{629.894}{a\Character{59}}%
\TextAt{9}{324.327}{629.894}{\Character{40}}%
\TextAt{13}{328.879}{629.894}{a\Character{59}}%
\TextAt{13}{340.248}{629.894}{b}%
\NativeMathAt{348.566}{629.894}{11.955}{\in}%
\TextAt{18}{359.857}{629.894}{Z}%
\TextAt{9}{367.827}{629.894}{\Character{41}}%
\TextAt{13}{372.379}{629.894}{\Character{59}}%
\TextAt{8}{79.226}{603.929}{the}%
\TextAt{8}{100.506}{603.929}{simplest}%
\TextAt{8}{147.274}{603.929}{among}%
\TextAt{8}{186.127}{603.929}{the}%
\TextAt{8}{207.407}{603.929}{nontrivial}%
\TextAt{8}{261.874}{603.929}{diophantine}%
\TextAt{8}{327.065}{603.929}{equations.}%
\TextAt{8}{387.409}{603.929}{The}%
\TextAt{8}{412.586}{603.929}{theory}%
\TextAt{8}{450.459}{603.929}{of}%
\TextAt{8}{464.913}{603.929}{this}%
\TextAt{8}{488.858}{603.929}{equation}%
\TextAt{8}{79.226}{586.883}{reduces}%
\TextAt{8}{120.106}{586.883}{immediately}%
\TextAt{8}{185.971}{586.883}{to}%
\TextAt{8}{199.485}{586.883}{the}%
\TextAt{8}{218.846}{586.883}{case}%
\TextAt{8}{242.82}{586.883}{where}%
\TextAt{13}{275.833}{586.883}{b}%
\TextAt{8}{283.912}{586.883}{is}%
\TextAt{8}{294.879}{586.883}{an}%
\TextAt{8}{310.342}{586.883}{odd}%
\TextAt{8}{332.631}{586.883}{prime}%
\TextAt{8}{364.998}{586.883}{number}%
\TextAt{13}{406.798}{586.883}{p}%
\TextAt{8}{415.775}{586.883}{and}%
\TextAt{9}{437.741}{586.883}{\Character{40}}%
\TextAt{13}{442.293}{586.883}{a\Character{59}}%
\TextAt{13}{453.662}{586.883}{p}%
\TextAt{9}{459.558}{586.883}{\Character{41}}%
\TextAt{9}{467.413}{586.883}{\Character{61}}%
\TextAt{9}{479.834}{586.883}{1}%
\TextAt{8}{488.81}{586.883}{(exercise}%
\TextAt{8}{79.226}{569.837}{6).}%
\TextAt{8}{98.074}{569.837}{Let}%
\TextAt{8}{119.005}{569.837}{us}%
\TextAt{8}{133.982}{569.837}{assume}%
\TextAt{8}{174.387}{569.837}{this}%
\TextAt{8}{197.17}{569.837}{for}%
\TextAt{8}{215.016}{569.837}{the}%
\TextAt{8}{235.134}{569.837}{sequel.}%
\TextAt{8}{274.521}{569.837}{We}%
\TextAt{8}{294.627}{569.837}{are}%
\TextAt{8}{314.098}{569.837}{then}%
\TextAt{8}{340.719}{569.837}{facing}%
\TextAt{8}{374.824}{569.837}{the}%
\TextAt{8}{394.942}{569.837}{question}%
\TextAt{8}{441.468}{569.837}{as}%
\TextAt{8}{455.799}{569.837}{to}%
\TextAt{8}{470.071}{569.837}{whether}%
\TextAt{13}{514.9}{569.837}{a}%
\TextAt{8}{524.903}{569.837}{is}%
\TextAt{8}{79.226}{552.791}{a}%
\TextAt{12}{88.982}{552.791}{quadratic}%
\TextAt{12}{149.044}{552.791}{residue}%
\TextAt{8}{194.701}{552.791}{mod}%
\TextAt{13}{221.038}{552.791}{p}%
\TextAt{8}{226.913}{552.791}{,}%
\TextAt{8}{234.063}{552.791}{i.e.,}%
\TextAt{8}{256.168}{552.791}{whether}%
\TextAt{8}{301.036}{552.791}{the}%
\TextAt{8}{321.192}{552.791}{congruence}%
\TextAt{13}{262.745}{526.826}{x}%
\TextAt{17}{269.397}{531.763}{2}%
\NativeMathAt{277.45}{526.826}{11.955}{\equiv}%
\TextAt{13}{290.069}{526.826}{a}%
\TextAt{9}{308.169}{526.826}{\Character{40}}%
\TextAt{8}{312.721}{526.826}{mod}%
\TextAt{13}{338.824}{526.826}{p}%
\TextAt{9}{344.699}{526.826}{\Character{41}}%
\TextAt{8}{79.226}{500.861}{does}%
\TextAt{8}{105.133}{500.861}{or}%
\TextAt{8}{118.952}{500.861}{does}%
\TextAt{8}{144.859}{500.861}{not}%
\TextAt{8}{165.182}{500.861}{have}%
\TextAt{8}{191.686}{500.861}{a}%
\TextAt{8}{200.951}{500.861}{solution.}%
\TextAt{8}{249.632}{500.861}{In}%
\TextAt{8}{263.763}{500.861}{other}%
\TextAt{8}{293.842}{500.861}{words,}%
\TextAt{8}{330.259}{500.861}{we}%
\TextAt{8}{346.996}{500.861}{want}%
\TextAt{8}{375.126}{500.861}{to}%
\TextAt{8}{388.946}{500.861}{know}%
\TextAt{8}{419.025}{500.861}{if}%
\TextAt{8}{429.258}{500.861}{the}%
\TextAt{8}{448.924}{500.861}{equation}%
\ErrataAt{496.2}{500.861}{38.8}{$\bar x^2=\bar a$,}%
\ErrataAt{79.19}{483.816}{455.81}{for a given element $\bar a=a\bmod p\in\mathbb F_p^*$, admits a solution in the field $\mathbb F_p$ or not. For this}%
\TextAt{8}{79.226}{465.728}{one}%
\TextAt{8}{100.214}{465.728}{introduces}%
\TextAt{8}{156.386}{465.728}{the}%
\TextAt{12}{176.07}{465.728}{Legendre}%
\TextAt{12}{233.666}{465.728}{symbol}%
\TextAt{20}{278.659}{478.998}{\Character{16}}%
\TextAt{21}{286.995}{470.435}{a}%
\Vector{q 0 G 0 g 0.478 w 214.995 -305.283 m 219.493 -305.283 l S Q}%
\TextAt{21}{287.112}{461.605}{p}%
\TextAt{20}{292.688}{478.998}{\Character{17}}%
\TextAt{8}{299.828}{465.728}{,}%
\TextAt{8}{306.6}{465.728}{which,}%
\TextAt{8}{342.961}{465.728}{for}%
\TextAt{8}{360.374}{465.728}{every}%
\TextAt{8}{390.794}{465.728}{rational}%
\TextAt{8}{433.911}{465.728}{number}%
\TextAt{13}{476.036}{465.728}{a}%
\TextAt{8}{485.606}{465.728}{relatively}%
\TextAt{8}{79.226}{443.212}{prime}%
\TextAt{8}{112.452}{443.212}{to}%
\TextAt{13}{126.824}{443.212}{p}%
\TextAt{8}{132.699}{443.212}{,}%
\TextAt{8}{139.926}{443.212}{is}%
\TextAt{8}{151.751}{443.212}{defined}%
\TextAt{20}{192.126}{456.482}{\Character{16}}%
\TextAt{21}{200.461}{447.919}{a}%
\Vector{q 0 G 0 g 0.478 w 128.461 -327.799 m 132.959 -327.799 l S Q}%
\TextAt{21}{200.579}{439.089}{p}%
\TextAt{20}{206.155}{456.482}{\Character{17}}%
\TextAt{9}{216.721}{443.212}{\Character{61}}%
\TextAt{9}{229.238}{443.212}{1}%
\TextAt{8}{239.064}{443.212}{or}%
\TextAt{16}{253.436}{443.212}{\Character{0}}%
\TextAt{9}{262.735}{443.212}{1}%
\TextAt{8}{268.588}{443.212}{,}%
\TextAt{8}{275.815}{443.212}{according}%
\TextAt{8}{328.563}{443.212}{as}%
\TextAt{13}{342.995}{443.212}{x}%
\TextAt{17}{349.647}{447.55}{2}%
\NativeMathAt{357.805}{443.212}{11.955}{\equiv}%
\TextAt{13}{370.53}{443.212}{a}%
\TextAt{9}{381.988}{443.212}{\Character{40}}%
\TextAt{8}{386.541}{443.212}{mod}%
\TextAt{13}{412.643}{443.212}{p}%
\TextAt{9}{418.518}{443.212}{\Character{41}}%
\TextAt{8}{427.03}{443.212}{has}%
\TextAt{8}{447.965}{443.212}{or}%
\TextAt{8}{462.337}{443.212}{does}%
\TextAt{8}{488.796}{443.212}{not}%
\TextAt{8}{509.672}{443.212}{have}%
\TextAt{8}{79.226}{426.166}{a}%
\TextAt{8}{88.982}{426.166}{solution.}%
\TextAt{8}{137.826}{426.166}{This}%
\TextAt{8}{164.546}{426.166}{symbol}%
\TextAt{8}{204.608}{426.166}{is}%
\TextAt{8}{216.372}{426.166}{multiplicative,}%
\TextAt{20}{252.948}{413.334}{\Character{18}}%
\TextAt{13}{262.944}{404.565}{ab}%
\Vector{q 0 G 0 g 0.478 w 190.944 -374.534 m 202.066 -374.534 l S Q}%
\TextAt{13}{265.567}{388.276}{p}%
\TextAt{20}{275.261}{413.334}{\Character{19}}%
\TextAt{9}{287.382}{396.477}{\Character{61}}%
\TextAt{20}{299.808}{413.334}{\Character{18}}%
\TextAt{13}{309.804}{404.565}{a}%
\Vector{q 0 G 0 g 0.478 w 237.804 -374.534 m 243.949 -374.534 l S Q}%
\TextAt{13}{309.939}{388.276}{p}%
\TextAt{20}{317.144}{413.334}{\Character{19}}%
\TextAt{20}{327.928}{413.334}{\Character{18}}%
\TextAt{13}{338.382}{404.565}{b}%
\Vector{q 0 G 0 g 0.478 w 265.933 -374.534 m 271.808 -374.534 l S Q}%
\TextAt{13}{337.933}{388.276}{p}%
\TextAt{20}{345.003}{413.334}{\Character{19}}%
\TextAt{13}{355.796}{396.477}{\Character{58}}%
\TextAt{8}{79.226}{365.729}{This}%
\TextAt{8}{105.887}{365.729}{is}%
\TextAt{8}{117.592}{365.729}{because}%
\TextAt{8}{160.835}{365.729}{the}%
\TextAt{8}{180.933}{365.729}{group}%
\TextAt{18}{214.05}{365.729}{F}%
\TextAt{3}{221.356}{370.068}{\Character{3}}%
\TextAt{21}{221.356}{362.774}{p}%
\TextAt{8}{229.955}{365.729}{is}%
\TextAt{8}{241.66}{365.729}{cyclic}%
\TextAt{8}{273.785}{365.729}{of}%
\TextAt{8}{287.056}{365.729}{order}%
\TextAt{13}{317.567}{365.729}{p}%
\TextAt{16}{325.978}{365.729}{\Character{0}}%
\TextAt{9}{337.813}{365.729}{1}%
\TextAt{8}{347.505}{365.729}{and}%
\TextAt{8}{370.209}{365.729}{the}%
\TextAt{8}{390.306}{365.729}{subgroup}%
\TextAt{18}{441.045}{365.729}{F}%
\TextAt{3}{448.351}{370.068}{\Character{3}}%
\TextAt{17}{452.585}{370.068}{2}%
\TextAt{21}{448.351}{362.774}{p}%
\TextAt{8}{461.156}{365.729}{of}%
\TextAt{8}{474.427}{365.729}{squares}%
\TextAt{8}{515.793}{365.729}{has}%
\TextAt{8}{79.226}{346.479}{index}%
\TextAt{9}{110.764}{346.479}{2}%
\TextAt{8}{116.617}{346.479}{,}%
\TextAt{8}{123.766}{346.479}{i.e.,}%
\TextAt{18}{145.871}{346.479}{F}%
\TextAt{3}{153.177}{350.818}{\Character{3}}%
\TextAt{21}{153.177}{343.524}{p}%
\TextAt{9}{157.938}{346.479}{\Character{47}}%
\TextAt{18}{163.791}{346.479}{F}%
\TextAt{3}{171.097}{350.818}{\Character{3}}%
\TextAt{17}{175.331}{350.818}{2}%
\TextAt{21}{171.097}{343.524}{p}%
\NativeMathAt{183.384}{346.479}{11.955}{\cong}%
\TextAt{18}{196.004}{346.479}{Z}%
\TextAt{9}{203.974}{346.479}{\Character{47}2}%
\TextAt{18}{215.68}{346.479}{Z}%
\TextAt{8}{223.65}{346.479}{.}%
\TextAt{8}{232.098}{346.479}{Since}%
\TextAt{20}{262.656}{359.75}{\Character{16}}%
\TextAt{21}{270.991}{351.186}{a}%
\Vector{q 0 G 0 g 0.478 w 198.991 -424.532 m 203.489 -424.532 l S Q}%
\TextAt{21}{271.109}{342.357}{p}%
\TextAt{20}{276.685}{359.75}{\Character{17}}%
\TextAt{9}{287.145}{346.479}{\Character{61}}%
\TextAt{9}{299.555}{346.479}{1}%
\NativeArrowAt{308.745}{346.479}{11.955}{21.926}{\Longleftrightarrow}%
\Vector{q 0 G 0 g 0.478 w 261.983 -420.7 m 268.128 -420.7 l S Q}%
\TextAt{13}{333.983}{346.479}{a}%
\NativeMathAt{343.449}{346.479}{11.955}{\in}%
\TextAt{18}{354.74}{346.479}{F}%
\TextAt{3}{362.046}{350.818}{\Character{3}}%
\TextAt{17}{366.28}{350.818}{2}%
\TextAt{21}{362.046}{343.524}{p}%
\TextAt{8}{374.91}{346.479}{one}%
\TextAt{8}{396.369}{346.479}{also}%
\TextAt{8}{419.849}{346.479}{has}%
\TextAt{20}{238.424}{328.315}{\Character{18}}%
\TextAt{13}{248.42}{319.546}{a}%
\Vector{q 0 G 0 g 0.478 w 176.42 -459.553 m 182.565 -459.553 l S Q}%
\TextAt{13}{248.555}{303.257}{p}%
\TextAt{20}{255.761}{328.315}{\Character{19}}%
\NativeMathAt{267.882}{311.458}{11.955}{\equiv}%
\TextAt{13}{280.501}{311.458}{a}%
\TextAt{17}{286.646}{316.394}{\Character{40}}%
\TextAt{21}{289.939}{316.394}{p}%
\TextAt{3}{294.202}{316.394}{\Character{0}}%
\TextAt{17}{300.789}{316.394}{1\Character{41}\Character{47}2}%
\TextAt{9}{329.238}{311.458}{\Character{40}}%
\TextAt{8}{333.79}{311.458}{mod}%
\TextAt{13}{359.892}{311.458}{p}%
\TextAt{9}{365.767}{311.458}{\Character{41}}%
\TextAt{13}{370.32}{311.458}{\Character{58}}%
\TextAt{8}{79.226}{276.437}{In}%
\TextAt{8}{94.568}{276.437}{the}%
\TextAt{8}{115.446}{276.437}{case}%
\TextAt{8}{140.938}{276.437}{of}%
\TextAt{8}{154.989}{276.437}{squarefree}%
\TextAt{13}{211.05}{276.437}{a}%
\TextAt{8}{217.195}{276.437}{,}%
\TextAt{8}{225.245}{276.437}{the}%
\TextAt{8}{246.123}{276.437}{Legendre}%
\TextAt{8}{297.079}{276.437}{symbol}%
\TextAt{20}{337.862}{289.707}{\Character{16}}%
\TextAt{21}{346.197}{281.144}{a}%
\Vector{q 0 G 0 g 0.478 w 274.197 -494.575 m 278.695 -494.575 l S Q}%
\TextAt{21}{346.315}{272.314}{p}%
\TextAt{20}{351.891}{289.707}{\Character{17}}%
\TextAt{8}{363.649}{276.437}{bears}%
\TextAt{8}{395.322}{276.437}{the}%
\TextAt{8}{416.199}{276.437}{following}%
\TextAt{8}{466.355}{276.437}{relation}%
\TextAt{8}{510.007}{276.437}{with}%
\TextAt{8}{79.226}{253.921}{prime}%
\TextAt{8}{112.39}{253.921}{factorization.}%
\TextAt{20}{184.918}{267.191}{\Character{16}}%
\TextAt{21}{193.253}{258.628}{a}%
\Vector{q 0 G 0 g 0.478 w 121.253 -517.09 m 125.751 -517.09 l S Q}%
\TextAt{21}{193.371}{249.798}{p}%
\TextAt{20}{198.947}{267.191}{\Character{17}}%
\TextAt{9}{209.408}{253.921}{\Character{61}}%
\TextAt{9}{221.818}{253.921}{1}%
\TextAt{8}{231.583}{253.921}{signifies}%
\TextAt{8}{275.28}{253.921}{that}%
\TextAt{13}{217.645}{226.162}{x}%
\TextAt{17}{224.297}{231.098}{2}%
\TextAt{16}{231.686}{226.162}{\Character{0}}%
\TextAt{13}{243.642}{226.162}{a}%
\NativeMathAt{253.107}{226.162}{11.955}{\equiv}%
\TextAt{9}{265.727}{226.162}{\Character{40}}%
\TextAt{13}{270.279}{226.162}{x}%
\TextAt{16}{279.588}{226.162}{\Character{0}}%
\TextAt{13}{291.543}{226.162}{\Character{11}}%
\TextAt{9}{299.065}{226.162}{\Character{41}\Character{40}}%
\TextAt{13}{308.169}{226.162}{x}%
\TextAt{9}{317.478}{226.162}{\Character{43}}%
\TextAt{13}{329.239}{226.162}{\Character{11}}%
\TextAt{9}{336.761}{226.162}{\Character{41}}%
\TextAt{9}{353.247}{226.162}{\Character{40}}%
\TextAt{8}{357.821}{226.162}{mod}%
\TextAt{13}{383.923}{226.162}{p}%
\TextAt{9}{389.798}{226.162}{\Character{41}}%
\TextAt{8}{79.226}{200.197}{for}%
\TextAt{8}{96.892}{200.197}{some}%
\TextAt{13}{126}{200.197}{\Character{11}}%
\NativeMathAt{136.842}{200.197}{11.955}{\in}%
\TextAt{18}{148.133}{200.197}{Z}%
\TextAt{8}{156.104}{200.197}{.}%
\TextAt{8}{164.479}{200.197}{The}%
\TextAt{8}{188.314}{200.197}{conductor}%
\TextAt{8}{242.731}{200.197}{of}%
\TextAt{18}{255.842}{200.197}{Z}%
\TextAt{9}{263.812}{200.197}{\Character{91}}%
\TextAt{16}{267.064}{208.808}{\Character{112}}%
\Vector{q 0 G 0 g 0.478 w 205.027 -564.953 m 211.172 -564.953 l S Q}%
\TextAt{13}{277.027}{200.197}{a}%
\TextAt{9}{283.172}{200.197}{\Character{93}}%
\TextAt{8}{290.102}{200.197}{in}%
\TextAt{8}{303.536}{200.197}{the}%
\TextAt{8}{323.474}{200.197}{ring}%
\TextAt{8}{347.321}{200.197}{of}%
\TextAt{8}{360.433}{200.197}{integers}%
\TextAt{8}{403.528}{200.197}{of}%
\TextAt{18}{416.639}{200.197}{Q}%
\TextAt{9}{425.938}{200.197}{\Character{40}}%
\TextAt{16}{430.49}{208.808}{\Character{112}}%
\Vector{q 0 G 0 g 0.478 w 368.452 -564.953 m 374.597 -564.953 l S Q}%
\TextAt{13}{440.452}{200.197}{a}%
\TextAt{9}{446.597}{200.197}{\Character{41}}%
\TextAt{8}{454.829}{200.197}{is}%
\TextAt{8}{466.374}{200.197}{a}%
\TextAt{8}{475.911}{200.197}{divisor}%
\TextAt{8}{513.805}{200.197}{of}%
\TextAt{9}{526.917}{200.197}{2}%
\TextAt{8}{79.226}{183.151}{(see}%
\TextAt{8}{102.694}{183.151}{§2,}%
\TextAt{8}{121.38}{183.151}{exercise}%
\TextAt{8}{164.682}{183.151}{4).}%
\TextAt{8}{183.543}{183.151}{We}%
\TextAt{8}{203.687}{183.151}{may}%
\TextAt{8}{229.056}{183.151}{therefore}%
\TextAt{8}{278.156}{183.151}{apply}%
\TextAt{8}{310.351}{183.151}{proposition}%
\TextAt{8}{371.885}{183.151}{(8.3)}%
\TextAt{8}{399.86}{183.151}{and}%
\TextAt{8}{422.622}{183.151}{obtain}%
\TextAt{8}{459.05}{183.151}{the}%
\TextAt{12}{79.226}{158.006}{Proposition}%
\TextAt{12}{151.937}{158.006}{(8.5).}%
\TextAt{8}{187.253}{158.006}{For}%
\TextAt{8}{208.222}{158.006}{squarefree}%
\TextAt{13}{263.563}{158.006}{a}%
\TextAt{8}{273.605}{158.006}{and}%
\TextAt{9}{296.368}{158.006}{\Character{40}}%
\TextAt{13}{300.92}{158.006}{p\Character{59}}%
\TextAt{9}{312.039}{158.006}{2}%
\TextAt{13}{317.892}{158.006}{a}%
\TextAt{9}{324.037}{158.006}{\Character{41}}%
\TextAt{9}{331.892}{158.006}{\Character{61}}%
\TextAt{9}{344.313}{158.006}{1}%
\TextAt{8}{350.189}{158.006}{,}%
\TextAt{8}{357.338}{158.006}{we}%
\TextAt{8}{374.566}{158.006}{have}%
\TextAt{8}{401.56}{158.006}{the}%
\TextAt{8}{421.717}{158.006}{equivalence}%
\TextAt{20}{187.678}{144.612}{\Character{18}}%
\TextAt{13}{197.674}{135.843}{a}%
\Vector{q 0 G 0 g 0.478 w 125.674 -643.256 m 131.819 -643.256 l S Q}%
\TextAt{13}{197.809}{119.555}{p}%
\TextAt{20}{205.014}{144.612}{\Character{19}}%
\TextAt{9}{217.136}{127.755}{\Character{61}}%
\TextAt{9}{229.546}{127.755}{1}%
\NativeArrowAt{250.69}{127.755}{11.955}{21.926}{\Longleftrightarrow}%
\TextAt{13}{287.884}{127.755}{p}%
\TextAt{8}{297.656}{127.755}{is}%
\TextAt{8}{309.42}{127.755}{totally}%
\TextAt{8}{346.828}{127.755}{split}%
\TextAt{8}{372.902}{127.755}{in}%
\TextAt{18}{386.555}{127.755}{Q}%
\TextAt{9}{395.854}{127.755}{\Character{40}}%
\TextAt{16}{400.406}{136.95}{\Character{112}}%
\Vector{q 0 G 0 g 0.478 w 338.369 -636.811 m 344.514 -636.811 l S Q}%
\TextAt{13}{410.369}{127.755}{a}%
\TextAt{9}{416.514}{127.755}{\Character{41}}%
\TextAt{13}{421.066}{127.755}{\Character{58}}%
\TextAt{8}{97.159}{89.967}{For}%
\TextAt{8}{117.67}{89.967}{the}%
\TextAt{8}{137.368}{89.967}{Legendre}%
\TextAt{8}{187.145}{89.967}{symbol,}%
\TextAt{8}{230.092}{89.967}{one}%
\TextAt{8}{251.094}{89.967}{has}%
\TextAt{8}{271.509}{89.967}{the}%
\TextAt{8}{291.207}{89.967}{following}%
\TextAt{8}{340.184}{89.967}{remarkable}%
\TextAt{8}{400.541}{89.967}{law,}%
\TextAt{8}{424.563}{89.967}{which}%
\TextAt{8}{457.591}{89.967}{like}%
\TextAt{8}{478.593}{89.967}{none}%
\TextAt{8}{506.097}{89.967}{other}%
\TextAt{8}{79.226}{72.921}{has}%
\TextAt{8}{100.1}{72.921}{left}%
\TextAt{8}{120.579}{72.921}{its}%
\TextAt{8}{136.898}{72.921}{mark}%
\TextAt{8}{167.144}{72.921}{on}%
\TextAt{8}{183.404}{72.921}{the}%
\TextAt{8}{203.56}{72.921}{development}%
\TextAt{8}{271.525}{72.921}{of}%
\TextAt{8}{284.855}{72.921}{algebraic}%
\TextAt{8}{334.29}{72.921}{number}%
\TextAt{8}{376.886}{72.921}{theory.}%
\TextAt{8}{300.14}{43.033}{47}%
\Destination{78.23}{809.815}{nav.page.47}%
\Destination{79.36}{724.955}{nav.print.50}%
\Destination{79.36}{172.055}{nav.number.1.8.5}%
\Hotlink{515.712}{724.855}{14.294}{11.047}{nav.chapter.3}%
\Hotlink{84.405}{707.755}{6.858}{11.047}{nav.section.3.2}%
\Hotlink{164.182}{180.955}{6.858}{11.047}{nav.ex.1.2.4}%
\Hotlink{107.873}{180.955}{6.858}{11.047}{nav.section.1.2}%
\Hotlink{371.385}{180.955}{25.078}{11.047}{nav.number.1.8.3}%
\Destination{79.226}{516}{errata.185}%
\Hotlink{79.715}{567.6}{4.866}{11.047}{nav.ex.1.8.6}%
}

% Page 60 (body page 48)
\Page{%
\TextAt{2}{79.226}{796.943}{I.}%
\TextAt{2}{89.83}{796.943}{Algebraic}%
\TextAt{2}{138.648}{796.943}{Integers}%
\Vector{q 0 G 0 g 0.299 w 7.226 18.729 m 460.769 18.729 l S Q}%
\TextAt{12}{79.226}{761.124}{Theorem}%
\TextAt{12}{136.998}{761.124}{(8.6)}%
\TextAt{12}{169.915}{761.124}{(Gauss’s}%
\TextAt{12}{224.638}{761.124}{Reciprocity}%
\TextAt{12}{297.521}{761.124}{Law).}%
\TextAt{8}{335.026}{761.124}{For}%
\TextAt{8}{356.654}{761.124}{two}%
\TextAt{8}{379.43}{761.124}{distinct}%
\TextAt{8}{422.421}{761.124}{odd}%
\TextAt{8}{446.166}{761.124}{prime}%
\TextAt{8}{479.988}{761.124}{numbers}%
\TextAt{13}{527.857}{761.124}{\Character{96}}%
\TextAt{8}{79.226}{744.078}{and}%
\TextAt{13}{101.989}{744.078}{p}%
\TextAt{8}{107.864}{744.078}{,}%
\TextAt{8}{115.013}{744.078}{the}%
\TextAt{8}{135.17}{744.078}{following}%
\TextAt{8}{184.604}{744.078}{identity}%
\TextAt{8}{227.858}{744.078}{holds:}%
\TextAt{20}{243.649}{731.979}{\Character{18}}%
\TextAt{13}{254.126}{723.21}{\Character{96}}%
\Vector{q 0 G 0 g 0.478 w 181.645 -55.889 m 187.52 -55.889 l S Q}%
\TextAt{13}{253.645}{706.922}{p}%
\TextAt{20}{260.715}{731.979}{\Character{19}}%
\TextAt{20}{271.508}{728.393}{\Character{16}}%
\TextAt{13}{279.844}{723.21}{p}%
\Vector{q 0 G 0 g 0.478 w 207.844 -55.889 m 213.719 -55.889 l S Q}%
\TextAt{13}{280.325}{706.922}{\Character{96}}%
\TextAt{20}{286.914}{728.393}{\Character{17}}%
\TextAt{9}{297.375}{715.122}{\Character{61}}%
\TextAt{9}{309.785}{715.122}{\Character{40}}%
\TextAt{16}{314.353}{715.123}{\Character{0}}%
\TextAt{9}{323.651}{715.122}{1\Character{41}}%
\TextAt{22}{335.252}{723.419}{\Character{96}}%
\TextAt{23}{338.352}{723.419}{\Character{0}}%
\TextAt{24}{344.108}{723.419}{1}%
\Vector{q 0 G 0 g 0.359 w 263.252 -51.949 m 275.761 -51.949 l S Q}%
\TextAt{24}{339.68}{716.933}{2}%
\TextAt{22}{350.152}{723.752}{p}%
\TextAt{23}{353.992}{723.752}{\Character{0}}%
\TextAt{24}{359.748}{723.752}{1}%
\Vector{q 0 G 0 g 0.359 w 278.152 -51.949 m 291.401 -51.949 l S Q}%
\TextAt{24}{354.95}{716.933}{2}%
\TextAt{13}{365.095}{715.122}{\Character{58}}%
\TextAt{8}{79.226}{686.179}{One}%
\TextAt{8}{103.937}{686.179}{also}%
\TextAt{8}{127.417}{686.179}{has}%
\TextAt{8}{148.291}{686.179}{the}%
\TextAt{8}{168.448}{686.179}{two}%
\TextAt{8}{190.565}{686.179}{“supplementary}%
\TextAt{8}{274.849}{686.179}{theorems”}%
\TextAt{20}{186.575}{674.093}{\Character{18}}%
\TextAt{16}{196.571}{665.324}{\Character{0}}%
\TextAt{9}{205.87}{665.324}{1}%
\Vector{q 0 G 0 g 0.478 w 124.571 -113.775 m 139.723 -113.775 l S Q}%
\TextAt{13}{201.209}{649.035}{p}%
\TextAt{20}{212.918}{674.093}{\Character{19}}%
\TextAt{9}{225.039}{657.236}{\Character{61}}%
\TextAt{9}{237.449}{657.236}{\Character{40}}%
\TextAt{16}{242.017}{657.236}{\Character{0}}%
\TextAt{9}{251.316}{657.236}{1\Character{41}}%
\TextAt{17}{261.721}{662.172}{\Character{40}}%
\TextAt{21}{265.014}{662.172}{p}%
\TextAt{3}{269.277}{662.172}{\Character{0}}%
\TextAt{17}{275.863}{662.172}{1\Character{41}\Character{47}2}%
\TextAt{13}{292.357}{657.236}{\Character{59}}%
\TextAt{20}{321.512}{674.093}{\Character{18}}%
\TextAt{9}{331.519}{665.324}{2}%
\Vector{q 0 G 0 g 0.478 w 259.508 -113.775 m 265.383 -113.775 l S Q}%
\TextAt{13}{331.508}{649.035}{p}%
\TextAt{20}{338.578}{674.093}{\Character{19}}%
\TextAt{9}{350.7}{657.236}{\Character{61}}%
\TextAt{9}{363.11}{657.236}{\Character{40}}%
\TextAt{16}{367.677}{657.236}{\Character{0}}%
\TextAt{9}{376.976}{657.236}{1\Character{41}}%
\TextAt{17}{387.381}{662.172}{\Character{40}}%
\TextAt{21}{390.675}{662.172}{p}%
\TextAt{24}{394.937}{664.985}{2}%
\TextAt{3}{399.088}{662.172}{\Character{0}}%
\TextAt{17}{405.675}{662.172}{1\Character{41}\Character{47}8}%
\TextAt{13}{422.169}{657.236}{\Character{58}}%
\TextAt{19}{79.226}{607.189}{Proof.}%
\TextAt{20}{116}{620.459}{\Character{16}}%
\TextAt{3}{124.336}{611.896}{\Character{0}}%
\TextAt{17}{130.922}{611.896}{1}%
\Vector{q 0 G 0 g 0.478 w 52.336 -163.822 m 63.156 -163.822 l S Q}%
\TextAt{21}{127.615}{603.066}{p}%
\TextAt{20}{136.352}{620.459}{\Character{17}}%
\NativeMathAt{146.813}{607.189}{11.955}{\equiv}%
\TextAt{9}{159.432}{607.189}{\Character{40}}%
\TextAt{16}{163.984}{607.189}{\Character{0}}%
\TextAt{9}{173.283}{607.189}{1\Character{41}}%
\TextAt{17}{183.688}{611.528}{\Character{40}}%
\TextAt{21}{186.981}{611.528}{p}%
\TextAt{3}{191.244}{611.528}{\Character{0}}%
\TextAt{17}{197.831}{611.528}{1\Character{41}\Character{47}2}%
\TextAt{9}{219.638}{607.189}{\Character{40}}%
\TextAt{8}{224.19}{607.189}{mod}%
\TextAt{13}{250.292}{607.189}{p}%
\TextAt{9}{256.167}{607.189}{\Character{41}}%
\TextAt{8}{264.617}{607.189}{implies}%
\TextAt{20}{304.344}{620.459}{\Character{16}}%
\TextAt{3}{312.68}{611.896}{\Character{0}}%
\TextAt{17}{319.266}{611.896}{1}%
\Vector{q 0 G 0 g 0.478 w 240.68 -163.822 m 251.5 -163.822 l S Q}%
\TextAt{21}{315.959}{603.066}{p}%
\TextAt{20}{324.696}{620.459}{\Character{17}}%
\TextAt{9}{335.157}{607.189}{\Character{61}}%
\TextAt{9}{347.567}{607.189}{\Character{40}}%
\TextAt{16}{352.134}{607.189}{\Character{0}}%
\TextAt{9}{361.433}{607.189}{1\Character{41}}%
\TextAt{17}{371.838}{611.528}{\Character{40}}%
\TextAt{21}{375.131}{611.528}{p}%
\TextAt{3}{379.394}{611.528}{\Character{0}}%
\TextAt{17}{385.981}{611.528}{1\Character{41}\Character{47}2}%
\TextAt{8}{406.372}{607.189}{since}%
\TextAt{13}{435.041}{607.189}{p}%
\NativeMathAt{444.236}{607.189}{11.955}{\neq}%
\TextAt{9}{456.647}{607.189}{2}%
\TextAt{8}{462.515}{607.189}{.}%
\TextAt{8}{97.159}{581.685}{In}%
\TextAt{8}{110.577}{581.685}{order}%
\TextAt{8}{139.943}{581.685}{to}%
\TextAt{8}{153.051}{581.685}{determine}%
\TextAt{20}{206.471}{594.955}{\Character{16}}%
\TextAt{17}{214.821}{586.392}{2}%
\Vector{q 0 G 0 g 0.478 w 142.806 -189.327 m 147.069 -189.327 l S Q}%
\TextAt{21}{214.806}{577.562}{p}%
\TextAt{20}{220.265}{594.955}{\Character{17}}%
\TextAt{8}{227.405}{581.685}{,}%
\TextAt{8}{233.591}{581.685}{we}%
\TextAt{8}{249.616}{581.685}{work}%
\TextAt{8}{277.034}{581.685}{in}%
\TextAt{8}{289.483}{581.685}{the}%
\TextAt{8}{308.437}{581.685}{ring}%
\TextAt{18}{331.3}{581.685}{Z}%
\TextAt{9}{339.27}{581.685}{\Character{91}}%
\TextAt{13}{342.521}{581.685}{i}%
\TextAt{9}{346.515}{581.685}{\Character{93}}%
\TextAt{8}{352.461}{581.685}{of}%
\TextAt{8}{364.588}{581.685}{gaussian}%
\TextAt{8}{410.345}{581.685}{integers.}%
\TextAt{8}{457.808}{581.685}{Since}%
\TextAt{9}{487.163}{581.685}{\Character{40}1}%
\TextAt{9}{497.743}{581.685}{\Character{43}}%
\TextAt{13}{507.066}{581.685}{i}%
\TextAt{9}{511.059}{581.685}{\Character{41}}%
\TextAt{17}{515.612}{586.023}{2}%
\TextAt{9}{523.665}{581.685}{\Character{61}}%
\TextAt{9}{79.226}{564.639}{2}%
\TextAt{13}{85.079}{564.639}{i}%
\TextAt{8}{89.073}{564.639}{,}%
\TextAt{8}{96.222}{564.639}{we}%
\TextAt{8}{113.449}{564.639}{find}%
\TextAt{9}{161.711}{547.593}{\Character{40}1}%
\TextAt{9}{174.754}{547.593}{\Character{43}}%
\TextAt{13}{186.534}{547.593}{i}%
\TextAt{9}{190.527}{547.593}{\Character{41}}%
\TextAt{21}{195.08}{552.529}{p}%
\TextAt{9}{203.161}{547.593}{\Character{61}}%
\TextAt{9}{215.571}{547.593}{\Character{40}1}%
\TextAt{9}{228.626}{547.593}{\Character{43}}%
\TextAt{13}{240.41}{547.593}{i}%
\TextAt{9}{244.404}{547.593}{\Character{41}\Character{40}\Character{40}1}%
\TextAt{9}{266.533}{547.593}{\Character{43}}%
\TextAt{13}{278.332}{547.593}{i}%
\TextAt{9}{282.325}{547.593}{\Character{41}}%
\TextAt{17}{286.877}{552.529}{2}%
\TextAt{9}{291.61}{547.593}{\Character{41}}%
\TextAt{17}{296.162}{552.529}{\Character{40}}%
\TextAt{21}{299.455}{552.529}{p}%
\TextAt{3}{303.718}{552.529}{\Character{0}}%
\TextAt{17}{310.305}{552.529}{1\Character{41}\Character{47}2}%
\TextAt{9}{330.119}{547.593}{\Character{61}}%
\TextAt{9}{342.529}{547.593}{\Character{40}1}%
\TextAt{9}{355.584}{547.593}{\Character{43}}%
\TextAt{13}{367.368}{547.593}{i}%
\TextAt{9}{371.362}{547.593}{\Character{41}}%
\TextAt{13}{375.914}{547.593}{i}%
\TextAt{17}{379.907}{552.529}{\Character{40}}%
\TextAt{21}{383.201}{552.529}{p}%
\TextAt{3}{387.463}{552.529}{\Character{0}}%
\TextAt{17}{394.05}{552.529}{1\Character{41}\Character{47}2}%
\TextAt{9}{410.544}{547.593}{2}%
\TextAt{17}{416.397}{552.529}{\Character{40}}%
\TextAt{21}{419.69}{552.529}{p}%
\TextAt{3}{423.953}{552.529}{\Character{0}}%
\TextAt{17}{430.539}{552.529}{1\Character{41}\Character{47}2}%
\TextAt{13}{447.033}{547.593}{\Character{59}}%
\TextAt{8}{79.226}{525.692}{and}%
\TextAt{8}{101.989}{525.692}{since}%
\TextAt{9}{130.657}{525.692}{\Character{40}1}%
\TextAt{9}{143.7}{525.692}{\Character{43}}%
\TextAt{13}{155.481}{525.692}{i}%
\TextAt{9}{159.474}{525.692}{\Character{41}}%
\TextAt{21}{164.026}{530.03}{p}%
\NativeMathAt{172.108}{525.692}{11.955}{\equiv}%
\TextAt{9}{184.727}{525.692}{1}%
\TextAt{9}{193.227}{525.692}{\Character{43}}%
\TextAt{13}{204.998}{525.692}{i}%
\TextAt{21}{208.992}{530.03}{p}%
\TextAt{9}{219.066}{525.692}{\Character{40}}%
\TextAt{8}{223.618}{525.692}{mod}%
\TextAt{13}{249.72}{525.692}{p}%
\TextAt{9}{255.596}{525.692}{\Character{41}}%
\TextAt{8}{264.045}{525.692}{and}%
\TextAt{20}{286.808}{538.962}{\Character{16}}%
\TextAt{17}{295.158}{530.399}{2}%
\Vector{q 0 G 0 g 0.478 w 223.143 -245.319 m 227.406 -245.319 l S Q}%
\TextAt{21}{295.143}{521.569}{p}%
\TextAt{20}{300.602}{538.962}{\Character{17}}%
\NativeMathAt{311.062}{525.692}{11.955}{\equiv}%
\TextAt{9}{323.682}{525.692}{2}%
\TextAt{17}{329.535}{530.03}{\Character{40}}%
\TextAt{21}{332.828}{530.03}{p}%
\TextAt{3}{337.091}{530.03}{\Character{0}}%
\TextAt{17}{343.677}{530.03}{1\Character{41}\Character{47}2}%
\TextAt{9}{365.485}{525.692}{\Character{40}}%
\TextAt{8}{370.037}{525.692}{mod}%
\TextAt{13}{396.139}{525.692}{p}%
\TextAt{9}{402.014}{525.692}{\Character{41}}%
\TextAt{8}{406.566}{525.692}{,}%
\TextAt{8}{413.716}{525.692}{it}%
\TextAt{8}{425.42}{525.692}{follows}%
\TextAt{8}{463.856}{525.692}{that}%
\TextAt{20}{182.376}{508.273}{\Character{18}}%
\TextAt{9}{192.383}{499.504}{2}%
\Vector{q 0 G 0 g 0.478 w 120.372 -279.595 m 126.247 -279.595 l S Q}%
\TextAt{13}{192.372}{483.216}{p}%
\TextAt{20}{199.442}{508.273}{\Character{19}}%
\TextAt{9}{210.235}{491.416}{\Character{40}1}%
\TextAt{9}{223.278}{491.416}{\Character{43}}%
\TextAt{13}{235.058}{491.416}{i}%
\TextAt{9}{239.052}{491.416}{\Character{41}}%
\TextAt{13}{243.604}{491.416}{i}%
\TextAt{17}{247.598}{496.353}{\Character{40}}%
\TextAt{21}{250.891}{496.353}{p}%
\TextAt{3}{255.154}{496.353}{\Character{0}}%
\TextAt{17}{261.74}{496.353}{1\Character{41}\Character{47}2}%
\NativeMathAt{281.555}{491.416}{11.955}{\equiv}%
\TextAt{9}{294.174}{491.416}{1}%
\TextAt{9}{302.674}{491.416}{\Character{43}}%
\TextAt{13}{314.445}{491.416}{i}%
\TextAt{9}{318.439}{491.416}{\Character{40}}%
\TextAt{16}{322.991}{491.416}{\Character{0}}%
\TextAt{9}{332.29}{491.416}{1\Character{41}}%
\TextAt{17}{342.695}{496.353}{\Character{40}}%
\TextAt{21}{345.988}{496.353}{p}%
\TextAt{3}{350.251}{496.353}{\Character{0}}%
\TextAt{17}{356.837}{496.353}{1\Character{41}\Character{47}2}%
\TextAt{9}{385.286}{491.416}{\Character{40}}%
\TextAt{8}{389.839}{491.416}{mod}%
\TextAt{13}{415.941}{491.416}{p}%
\TextAt{9}{421.816}{491.416}{\Character{41}}%
\TextAt{13}{426.368}{491.416}{\Character{58}}%
\TextAt{8}{79.226}{462.593}{From}%
\TextAt{8}{109.951}{462.593}{this,}%
\TextAt{8}{136.025}{462.593}{an}%
\TextAt{8}{152.284}{462.593}{easy}%
\TextAt{8}{178.036}{462.593}{computation}%
\TextAt{8}{246.335}{462.593}{yields}%
\TextAt{20}{126.999}{450.506}{\Character{18}}%
\TextAt{9}{137.006}{441.737}{2}%
\Vector{q 0 G 0 g 0.478 w 64.995 -337.362 m 70.87 -337.362 l S Q}%
\TextAt{13}{136.995}{425.449}{p}%
\TextAt{20}{144.066}{450.506}{\Character{19}}%
\NativeMathAt{156.187}{433.649}{11.955}{\equiv}%
\TextAt{9}{168.806}{433.649}{\Character{40}}%
\TextAt{16}{173.359}{433.649}{\Character{0}}%
\TextAt{9}{182.657}{433.649}{1\Character{41}}%
\TextAt{17}{193.062}{438.586}{\Character{40}}%
\TextAt{21}{196.356}{438.586}{p}%
\TextAt{3}{200.618}{438.586}{\Character{0}}%
\TextAt{17}{207.205}{438.586}{1\Character{41}\Character{47}4}%
\TextAt{9}{235.654}{433.649}{\Character{40}}%
\TextAt{8}{240.206}{433.649}{mod}%
\TextAt{13}{266.308}{433.649}{p}%
\TextAt{9}{272.183}{433.649}{\Character{41}}%
\TextAt{13}{276.736}{433.649}{\Character{59}}%
\TextAt{8}{293.935}{433.649}{resp.}%
\TextAt{20}{332.008}{450.506}{\Character{18}}%
\TextAt{9}{342.015}{441.737}{2}%
\Vector{q 0 G 0 g 0.478 w 270.004 -337.362 m 275.879 -337.362 l S Q}%
\TextAt{13}{342.004}{425.449}{p}%
\TextAt{20}{349.075}{450.506}{\Character{19}}%
\NativeMathAt{361.196}{433.649}{11.955}{\equiv}%
\TextAt{9}{373.815}{433.649}{\Character{40}}%
\TextAt{16}{378.368}{433.649}{\Character{0}}%
\TextAt{9}{387.666}{433.649}{1\Character{41}}%
\TextAt{17}{398.071}{438.586}{\Character{40}}%
\TextAt{21}{401.365}{438.586}{p}%
\TextAt{17}{405.628}{438.586}{\Character{43}1\Character{41}\Character{47}4}%
\TextAt{9}{440.663}{433.649}{\Character{40}}%
\TextAt{8}{445.215}{433.649}{mod}%
\TextAt{13}{471.317}{433.649}{p}%
\TextAt{9}{477.192}{433.649}{\Character{41}}%
\TextAt{13}{481.745}{433.649}{\Character{59}}%
\TextAt{8}{79.226}{399.374}{if}%
\TextAt{21}{91.145}{404.63}{p}%
\TextAt{3}{95.408}{404.63}{\Character{0}}%
\TextAt{17}{101.995}{404.63}{1}%
\Vector{q 0 G 0 g 0.478 w 19.145 -371.637 m 34.229 -371.637 l S Q}%
\TextAt{17}{96.57}{395.251}{2}%
\TextAt{8}{111.322}{399.374}{is}%
\TextAt{8}{123.086}{399.374}{even,}%
\TextAt{8}{152.998}{399.374}{resp.}%
\TextAt{8}{182.32}{399.374}{odd.}%
\TextAt{8}{209.956}{399.374}{Since}%
\TextAt{21}{241.709}{404.63}{p}%
\TextAt{24}{245.972}{407.443}{2}%
\TextAt{3}{250.123}{404.63}{\Character{0}}%
\TextAt{17}{256.709}{404.63}{1}%
\Vector{q 0 G 0 g 0.478 w 169.709 -371.637 m 188.943 -371.637 l S Q}%
\TextAt{17}{249.209}{395.251}{8}%
\TextAt{9}{265.46}{399.374}{\Character{61}}%
\TextAt{21}{279.081}{404.63}{p}%
\TextAt{3}{283.344}{404.63}{\Character{0}}%
\TextAt{17}{289.93}{404.63}{1}%
\Vector{q 0 G 0 g 0.478 w 207.081 -371.637 m 222.164 -371.637 l S Q}%
\TextAt{17}{284.505}{395.251}{4}%
\TextAt{21}{296.555}{404.63}{p}%
\TextAt{17}{300.818}{404.63}{\Character{43}1}%
\Vector{q 0 G 0 g 0.478 w 224.555 -371.637 m 239.639 -371.637 l S Q}%
\TextAt{17}{301.98}{395.251}{2}%
\TextAt{9}{316.155}{399.374}{\Character{61}}%
\TextAt{21}{329.776}{404.63}{p}%
\TextAt{17}{334.039}{404.63}{\Character{43}1}%
\Vector{q 0 G 0 g 0.478 w 257.776 -371.637 m 272.859 -371.637 l S Q}%
\TextAt{17}{335.201}{395.251}{4}%
\TextAt{21}{347.251}{404.63}{p}%
\TextAt{3}{351.513}{404.63}{\Character{0}}%
\TextAt{17}{358.1}{404.63}{1}%
\Vector{q 0 G 0 g 0.478 w 275.251 -371.637 m 290.334 -371.637 l S Q}%
\TextAt{17}{352.675}{395.251}{2}%
\TextAt{8}{363.529}{399.374}{,}%
\TextAt{8}{370.679}{399.374}{we}%
\TextAt{8}{387.906}{399.374}{deduce}%
\TextAt{20}{426.916}{412.644}{\Character{16}}%
\TextAt{17}{435.266}{404.081}{2}%
\Vector{q 0 G 0 g 0.478 w 363.251 -371.637 m 367.514 -371.637 l S Q}%
\TextAt{21}{435.251}{395.251}{p}%
\TextAt{20}{440.71}{412.644}{\Character{17}}%
\TextAt{9}{451.17}{399.374}{\Character{61}}%
\TextAt{9}{463.579}{399.374}{\Character{40}}%
\TextAt{16}{468.148}{399.374}{\Character{0}}%
\TextAt{9}{477.447}{399.374}{1\Character{41}}%
\TextAt{17}{487.852}{403.713}{\Character{40}}%
\TextAt{21}{491.145}{403.713}{p}%
\TextAt{24}{495.408}{406.525}{2}%
\TextAt{3}{499.559}{403.713}{\Character{0}}%
\TextAt{17}{506.145}{403.713}{1\Character{41}\Character{47}8}%
\TextAt{8}{522.639}{399.374}{.}%
\TextAt{8}{97.159}{379.339}{In}%
\TextAt{8}{112.426}{379.339}{order}%
\TextAt{8}{143.642}{379.339}{to}%
\TextAt{8}{158.599}{379.339}{prove}%
\TextAt{8}{190.795}{379.339}{the}%
\TextAt{8}{211.597}{379.339}{first}%
\TextAt{8}{236.369}{379.339}{formula,}%
\TextAt{8}{283.36}{379.339}{we}%
\TextAt{8}{301.233}{379.339}{work}%
\TextAt{8}{330.5}{379.339}{in}%
\TextAt{8}{344.799}{379.339}{the}%
\TextAt{8}{365.602}{379.339}{ring}%
\TextAt{18}{390.314}{379.339}{Z}%
\TextAt{9}{398.284}{379.339}{\Character{91}}%
\TextAt{13}{401.536}{379.339}{\Character{16}}%
\TextAt{9}{407.551}{379.339}{\Character{93}}%
\TextAt{8}{410.803}{379.339}{,}%
\TextAt{8}{418.76}{379.339}{where}%
\TextAt{13}{453.216}{379.339}{\Character{16}}%
\TextAt{8}{463.775}{379.339}{is}%
\TextAt{8}{476.185}{379.339}{a}%
\TextAt{8}{486.587}{379.339}{primitive}%
\TextAt{13}{79.226}{362.293}{\Character{96}}%
\TextAt{8}{84.138}{362.293}{-th}%
\TextAt{8}{102.992}{362.293}{root}%
\TextAt{8}{128.038}{362.293}{of}%
\TextAt{8}{141.368}{362.293}{unity.}%
\TextAt{8}{175.508}{362.293}{We}%
\TextAt{8}{195.652}{362.293}{consider}%
\TextAt{8}{241.237}{362.293}{the}%
\TextAt{12}{261.394}{362.293}{Gauss}%
\TextAt{12}{301.073}{362.293}{sum}%
\TextAt{13}{256.325}{337.074}{\Character{28}}%
\TextAt{9}{266.065}{337.074}{\Character{61}}%
\TextAt{20}{289.54}{348.432}{\Character{88}}%
\TextAt{21}{278.49}{322.529}{a}%
\NativeMathAt{282.988}{322.529}{7.97}{\in}%
\TextAt{17}{288.634}{322.529}{\Character{40}}%
\TextAt{25}{291.927}{322.529}{Z}%
\TextAt{17}{297.241}{322.529}{\Character{47}}%
\TextAt{21}{301.475}{322.529}{\Character{96}}%
\TextAt{25}{304.934}{322.529}{Z}%
\TextAt{17}{310.248}{322.529}{\Character{41}}%
\TextAt{23}{313.541}{324.798}{\Character{3}}%
\TextAt{20}{319.85}{350.345}{\Character{16}}%
\TextAt{13}{328.186}{345.162}{a}%
\Vector{q 0 G 0 g 0.478 w 256.186 -433.937 m 262.331 -433.937 l S Q}%
\TextAt{13}{328.802}{328.874}{\Character{96}}%
\TextAt{20}{335.526}{350.345}{\Character{17}}%
\TextAt{13}{344.659}{337.074}{\Character{16}}%
\TextAt{21}{350.674}{342.011}{a}%
\TextAt{8}{79.226}{301.875}{and}%
\TextAt{8}{101.989}{301.875}{show}%
\TextAt{8}{130.992}{301.875}{that}%
\TextAt{13}{269.899}{283.424}{\Character{28}}%
\TextAt{17}{276.318}{288.36}{2}%
\TextAt{9}{284.371}{283.424}{\Character{61}}%
\TextAt{20}{296.797}{300.281}{\Character{18}}%
\TextAt{16}{306.793}{291.511}{\Character{0}}%
\TextAt{9}{316.091}{291.511}{1}%
\Vector{q 0 G 0 g 0.478 w 234.793 -487.588 m 249.944 -487.588 l S Q}%
\TextAt{13}{311.912}{275.223}{\Character{96}}%
\TextAt{20}{323.14}{300.281}{\Character{19}}%
\TextAt{13}{333.933}{283.424}{\Character{96}\Character{58}}%
\TextAt{8}{79.226}{256.74}{For}%
\TextAt{8}{101.279}{256.74}{this,}%
\TextAt{8}{128.707}{256.74}{let}%
\TextAt{13}{146.695}{256.74}{a}%
\TextAt{8}{157.821}{256.74}{and}%
\TextAt{13}{181.666}{256.74}{b}%
\TextAt{8}{191.624}{256.74}{vary}%
\TextAt{8}{218.734}{256.74}{over}%
\TextAt{8}{244.863}{256.74}{the}%
\TextAt{8}{266.103}{256.74}{group}%
\TextAt{9}{300.362}{256.74}{\Character{40}}%
\TextAt{18}{304.914}{256.74}{Z}%
\TextAt{9}{312.884}{256.74}{\Character{47}}%
\TextAt{13}{318.737}{256.74}{\Character{96}}%
\TextAt{18}{323.649}{256.74}{Z}%
\TextAt{9}{331.619}{256.74}{\Character{41}}%
\TextAt{3}{336.172}{261.079}{\Character{3}}%
\TextAt{8}{340.904}{256.74}{,}%
\TextAt{8}{349.407}{256.74}{put}%
\TextAt{13}{371.95}{256.74}{c}%
\TextAt{9}{382.155}{256.74}{\Character{61}}%
\TextAt{13}{396.426}{256.74}{ab}%
\TextAt{3}{407.549}{261.079}{\Character{0}}%
\TextAt{17}{414.135}{261.079}{1}%
\TextAt{8}{423.848}{256.74}{and}%
\TextAt{8}{447.694}{256.74}{deduce}%
\TextAt{8}{487.787}{256.74}{from}%
\TextAt{8}{516.511}{256.74}{the}%
\TextAt{8}{79.226}{239.694}{identity}%
\TextAt{20}{122.48}{249.378}{\Character{0}}%
\TextAt{21}{129.155}{244.401}{b}%
\Vector{q 0 G 0 g 0.478 w 57.155 -531.317 m 60.778 -531.317 l S Q}%
\TextAt{21}{129.237}{235.572}{\Character{96}}%
\TextAt{20}{133.973}{249.378}{\Character{1}}%
\TextAt{9}{142.773}{239.694}{\Character{61}}%
\TextAt{20}{155.199}{252.965}{\Character{16}}%
\TextAt{21}{163.534}{244.402}{b}%
\TextAt{23}{167.157}{247.214}{\Character{0}}%
\TextAt{24}{172.913}{247.214}{1}%
\Vector{q 0 G 0 g 0.478 w 91.534 -531.317 m 105.064 -531.317 l S Q}%
\TextAt{21}{168.569}{235.572}{\Character{96}}%
\TextAt{20}{178.26}{252.965}{\Character{17}}%
\TextAt{8}{189.297}{239.694}{that}%
\TextAt{20}{140.85}{222.049}{\Character{18}}%
\TextAt{16}{150.846}{213.279}{\Character{0}}%
\TextAt{9}{160.145}{213.279}{1}%
\Vector{q 0 G 0 g 0.478 w 78.846 -565.82 m 93.997 -565.82 l S Q}%
\TextAt{13}{155.966}{196.991}{\Character{96}}%
\TextAt{20}{167.193}{222.049}{\Character{19}}%
\TextAt{13}{177.986}{205.192}{\Character{28}}%
\TextAt{17}{184.405}{210.128}{2}%
\TextAt{9}{192.458}{205.192}{\Character{61}}%
\TextAt{20}{204.883}{216.549}{\Character{88}}%
\TextAt{21}{208.281}{191.089}{a\Character{59}b}%
\TextAt{20}{224.145}{222.049}{\Character{18}}%
\TextAt{16}{234.14}{213.279}{\Character{0}}%
\TextAt{13}{243.439}{213.279}{ab}%
\Vector{q 0 G 0 g 0.478 w 162.14 -565.82 m 182.561 -565.82 l S Q}%
\TextAt{13}{241.895}{196.991}{\Character{96}}%
\TextAt{20}{255.756}{222.049}{\Character{19}}%
\TextAt{13}{266.549}{205.192}{\Character{16}}%
\TextAt{21}{272.565}{210.128}{a}%
\TextAt{17}{277.063}{210.128}{\Character{43}}%
\TextAt{21}{283.649}{210.128}{b}%
\TextAt{9}{291.091}{205.192}{\Character{61}}%
\TextAt{20}{303.516}{216.549}{\Character{88}}%
\TextAt{21}{306.914}{191.089}{a\Character{59}b}%
\TextAt{20}{322.778}{222.049}{\Character{18}}%
\TextAt{13}{332.773}{213.279}{ab}%
\TextAt{3}{343.896}{217.618}{\Character{0}}%
\TextAt{17}{350.482}{217.618}{1}%
\Vector{q 0 G 0 g 0.478 w 260.773 -565.82 m 283.214 -565.82 l S Q}%
\TextAt{13}{341.538}{196.991}{\Character{96}}%
\TextAt{20}{356.41}{222.049}{\Character{19}}%
\TextAt{13}{367.203}{205.192}{\Character{16}}%
\TextAt{21}{373.218}{210.128}{a}%
\TextAt{3}{377.716}{210.128}{\Character{0}}%
\TextAt{21}{384.303}{210.128}{b}%
\TextAt{9}{391.744}{205.192}{\Character{61}}%
\TextAt{20}{404.17}{216.549}{\Character{88}}%
\TextAt{21}{407.983}{191.089}{b\Character{59}c}%
\TextAt{20}{423.431}{218.462}{\Character{16}}%
\TextAt{13}{431.766}{213.279}{c}%
\Vector{q 0 G 0 g 0.478 w 359.766 -565.82 m 364.804 -565.82 l S Q}%
\TextAt{13}{431.829}{196.991}{\Character{96}}%
\TextAt{20}{438}{218.462}{\Character{17}}%
\TextAt{13}{447.132}{205.192}{\Character{16}}%
\TextAt{21}{453.148}{210.128}{bc}%
\TextAt{3}{460.438}{210.128}{\Character{0}}%
\TextAt{21}{467.025}{210.128}{b}%
\TextAt{9}{192.458}{167.023}{\Character{61}}%
\TextAt{20}{204.883}{178.381}{\Character{88}}%
\TextAt{21}{206.274}{152.921}{c}%
\NativeMathAt{209.941}{152.921}{7.97}{\neq}%
\TextAt{17}{216.524}{152.921}{1}%
\TextAt{20}{224.145}{180.294}{\Character{16}}%
\TextAt{13}{232.48}{175.111}{c}%
\Vector{q 0 G 0 g 0.478 w 160.48 -603.988 m 165.518 -603.988 l S Q}%
\TextAt{13}{232.543}{158.823}{\Character{96}}%
\TextAt{20}{238.713}{180.294}{\Character{17}}%
\TextAt{20}{247.846}{178.381}{\Character{88}}%
\TextAt{21}{254.669}{152.921}{b}%
\TextAt{13}{267.107}{167.023}{\Character{16}}%
\TextAt{21}{273.123}{171.96}{b}%
\TextAt{17}{276.745}{171.96}{\Character{40}}%
\TextAt{21}{280.038}{171.96}{c}%
\TextAt{3}{283.706}{171.96}{\Character{0}}%
\TextAt{17}{290.293}{171.96}{1\Character{41}}%
\TextAt{9}{300.975}{167.023}{\Character{43}}%
\TextAt{20}{312.736}{178.381}{\Character{88}}%
\TextAt{21}{319.559}{152.921}{b}%
\TextAt{20}{331.997}{183.88}{\Character{18}}%
\TextAt{9}{341.993}{175.111}{1}%
\Vector{q 0 G 0 g 0.478 w 269.993 -603.988 m 275.846 -603.988 l S Q}%
\TextAt{13}{342.464}{158.823}{\Character{96}}%
\TextAt{20}{349.042}{183.88}{\Character{19}}%
\TextAt{13}{359.835}{167.023}{\Character{58}}%
\TextAt{8}{79.226}{130.844}{Now}%
\TextAt{20}{107.257}{139.811}{\Character{80}}%
\TextAt{21}{119.876}{127.357}{c}%
\TextAt{20}{126.035}{140.528}{\Character{0}}%
\TextAt{21}{132.71}{135.551}{c}%
\Vector{q 0 G 0 g 0.478 w 60.71 -640.167 m 64.377 -640.167 l S Q}%
\TextAt{21}{132.814}{126.722}{\Character{96}}%
\TextAt{20}{137.573}{140.528}{\Character{1}}%
\TextAt{9}{148.709}{130.844}{\Character{61}}%
\TextAt{9}{163.462}{130.844}{0}%
\TextAt{8}{169.323}{130.844}{,}%
\TextAt{8}{178.186}{130.844}{as}%
\TextAt{8}{193.927}{130.844}{one}%
\TextAt{8}{216.757}{130.844}{sees}%
\TextAt{8}{241.655}{130.844}{by}%
\TextAt{8}{259.285}{130.844}{multiplying}%
\TextAt{8}{323.097}{130.844}{the}%
\TextAt{8}{344.624}{130.844}{sum}%
\TextAt{8}{370.766}{130.844}{with}%
\TextAt{8}{398.797}{130.844}{a}%
\TextAt{8}{409.923}{130.844}{symbol}%
\TextAt{20}{451.355}{140.528}{\Character{0}}%
\TextAt{21}{458.03}{135.551}{x}%
\Vector{q 0 G 0 g 0.478 w 386.03 -640.167 m 390.797 -640.167 l S Q}%
\TextAt{21}{458.684}{126.722}{\Character{96}}%
\TextAt{20}{463.993}{140.528}{\Character{1}}%
\TextAt{9}{475.129}{130.844}{\Character{61}}%
\TextAt{16}{489.89}{130.844}{\Character{0}}%
\TextAt{9}{499.189}{130.844}{1}%
\TextAt{8}{505.042}{130.844}{,}%
\TextAt{8}{513.904}{130.844}{and}%
\TextAt{8}{79.226}{113.798}{putting}%
\TextAt{13}{121.051}{113.798}{\Character{24}}%
\TextAt{9}{130.399}{113.798}{\Character{61}}%
\TextAt{13}{143.161}{113.798}{\Character{16}}%
\TextAt{21}{149.176}{118.137}{c}%
\TextAt{3}{152.844}{118.137}{\Character{0}}%
\TextAt{17}{159.43}{118.137}{1}%
\TextAt{8}{168.257}{113.798}{gives}%
\TextAt{20}{197.135}{122.765}{\Character{80}}%
\TextAt{21}{209.755}{110.311}{b}%
\TextAt{13}{215.868}{113.798}{\Character{16}}%
\TextAt{21}{221.883}{118.137}{b}%
\TextAt{17}{225.506}{118.137}{\Character{40}}%
\TextAt{21}{228.799}{118.137}{c}%
\TextAt{3}{232.467}{118.137}{\Character{0}}%
\TextAt{17}{239.054}{118.137}{1\Character{41}}%
\TextAt{9}{250.736}{113.798}{\Character{61}}%
\TextAt{13}{263.498}{113.798}{\Character{24}}%
\TextAt{9}{271.979}{113.798}{\Character{43}}%
\TextAt{13}{283.875}{113.798}{\Character{24}}%
\TextAt{17}{289.566}{118.137}{2}%
\TextAt{9}{297.089}{113.798}{\Character{43}}%
\TextAt{16}{308.985}{113.798}{\Character{1}}%
\TextAt{16}{314.281}{113.798}{\Character{1}}%
\TextAt{16}{319.589}{113.798}{\Character{1}}%
\TextAt{9}{325.724}{113.798}{\Character{43}}%
\TextAt{13}{337.62}{113.798}{\Character{24}}%
\TextAt{21}{343.31}{118.137}{\Character{96}}%
\TextAt{3}{346.769}{118.137}{\Character{0}}%
\TextAt{17}{353.356}{118.137}{1}%
\TextAt{9}{361.745}{113.798}{\Character{61}}%
\TextAt{16}{374.507}{113.798}{\Character{0}}%
\TextAt{9}{383.805}{113.798}{1}%
\TextAt{8}{389.658}{113.798}{,}%
\TextAt{8}{397.054}{113.798}{from}%
\TextAt{8}{424.892}{113.798}{which}%
\TextAt{8}{458.576}{113.798}{we}%
\TextAt{8}{476}{113.798}{indeed}%
\TextAt{8}{513.259}{113.798}{find}%
\TextAt{8}{79.226}{96.752}{that}%
\TextAt{20}{217.873}{95.158}{\Character{18}}%
\TextAt{16}{227.868}{86.389}{\Character{0}}%
\TextAt{9}{237.167}{86.389}{1}%
\Vector{q 0 G 0 g 0.478 w 155.868 -692.71 m 171.02 -692.71 l S Q}%
\TextAt{13}{232.988}{70.101}{\Character{96}}%
\TextAt{20}{244.215}{95.158}{\Character{19}}%
\TextAt{13}{255.008}{78.301}{\Character{28}}%
\TextAt{17}{261.427}{83.238}{2}%
\TextAt{9}{269.48}{78.301}{\Character{61}}%
\TextAt{9}{281.889}{78.301}{\Character{40}}%
\TextAt{16}{286.458}{78.301}{\Character{0}}%
\TextAt{9}{295.757}{78.301}{1\Character{41}\Character{40}}%
\TextAt{16}{310.714}{78.301}{\Character{0}}%
\TextAt{9}{320.013}{78.301}{1\Character{41}}%
\TextAt{9}{333.056}{78.301}{\Character{43}}%
\TextAt{13}{344.836}{78.301}{\Character{96}}%
\TextAt{16}{352.405}{78.301}{\Character{0}}%
\TextAt{9}{364.36}{78.301}{1}%
\TextAt{9}{373.518}{78.301}{\Character{61}}%
\TextAt{13}{385.959}{78.301}{\Character{96}\Character{58}}%
\TextAt{8}{300.14}{43.033}{48}%
\Destination{78.23}{809.815}{nav.page.48}%
\Destination{79.36}{776.155}{nav.print.51}%
\Destination{79.36}{775.155}{nav.number.1.8.6}%
}

% Page 61 (body page 49)
\Page{%
\TextAt{2}{79.226}{796.943}{I.}%
\TextAt{2}{89.83}{796.943}{Algebraic}%
\TextAt{2}{138.648}{796.943}{Integers}%
\Vector{q 0 G 0 g 0.299 w 7.226 18.729 m 460.769 18.729 l S Q}%
\TextAt{8}{79.226}{759.33}{This,}%
\TextAt{8}{111.339}{759.33}{together}%
\TextAt{8}{159.236}{759.33}{with}%
\TextAt{8}{187.609}{759.33}{the}%
\TextAt{8}{209.479}{759.33}{congruence}%
\TextAt{20}{271.673}{772.601}{\Character{16}}%
\TextAt{21}{280.41}{764.037}{\Character{96}}%
\Vector{q 0 G 0 g 0.478 w 208.009 -11.681 m 212.271 -11.681 l S Q}%
\TextAt{21}{280.009}{755.208}{p}%
\TextAt{20}{285.467}{772.601}{\Character{17}}%
\NativeMathAt{298.847}{759.33}{11.955}{\equiv}%
\TextAt{13}{314.386}{759.33}{\Character{96}}%
\TextAt{17}{319.298}{763.669}{\Character{40}}%
\TextAt{21}{322.592}{763.669}{p}%
\TextAt{3}{326.854}{763.669}{\Character{0}}%
\TextAt{17}{333.441}{763.669}{1\Character{41}\Character{47}2}%
\TextAt{9}{355.248}{759.33}{\Character{40}}%
\TextAt{8}{359.801}{759.33}{mod}%
\TextAt{13}{385.903}{759.33}{p}%
\TextAt{9}{391.778}{759.33}{\Character{41}}%
\TextAt{8}{401.941}{759.33}{and}%
\TextAt{8}{426.417}{759.33}{the}%
\TextAt{8}{448.286}{759.33}{identity}%
\TextAt{20}{493.254}{769.014}{\Character{0}}%
\TextAt{3}{499.929}{764.037}{\Character{0}}%
\TextAt{17}{506.515}{764.037}{1}%
\Vector{q 0 G 0 g 0.478 w 427.929 -11.681 m 438.749 -11.681 l S Q}%
\TextAt{21}{503.609}{755.208}{\Character{96}}%
\TextAt{20}{511.945}{769.014}{\Character{1}}%
\TextAt{9}{523.665}{759.33}{\Character{61}}%
\TextAt{9}{79.226}{740.247}{\Character{40}}%
\TextAt{16}{83.778}{740.247}{\Character{0}}%
\TextAt{9}{93.077}{740.247}{1\Character{41}}%
\TextAt{17}{103.482}{744.585}{\Character{40}}%
\TextAt{21}{106.776}{744.586}{\Character{96}}%
\TextAt{3}{110.235}{744.585}{\Character{0}}%
\TextAt{17}{116.821}{744.586}{1\Character{41}\Character{47}2}%
\TextAt{8}{133.315}{740.247}{,}%
\TextAt{8}{140.464}{740.247}{implies}%
\TextAt{13}{181.982}{709.153}{\Character{28}}%
\TextAt{21}{188.401}{714.089}{p}%
\TextAt{9}{196.483}{709.153}{\Character{61}}%
\TextAt{13}{208.908}{709.153}{\Character{28}}%
\TextAt{9}{215.327}{709.153}{\Character{40}}%
\TextAt{13}{219.88}{709.153}{\Character{28}}%
\TextAt{17}{226.299}{714.089}{2}%
\TextAt{9}{231.031}{709.153}{\Character{41}}%
\TextAt{17}{235.583}{714.089}{\Character{40}}%
\TextAt{21}{238.876}{714.089}{p}%
\TextAt{3}{243.139}{714.089}{\Character{0}}%
\TextAt{17}{249.726}{714.089}{1\Character{41}\Character{47}2}%
\NativeMathAt{269.54}{709.153}{11.955}{\equiv}%
\TextAt{13}{282.16}{709.153}{\Character{28}}%
\TextAt{9}{288.579}{709.153}{\Character{40}}%
\TextAt{16}{293.131}{709.153}{\Character{0}}%
\TextAt{9}{302.43}{709.153}{1\Character{41}}%
\TextAt{22}{314.03}{717.45}{\Character{96}}%
\TextAt{23}{317.13}{717.45}{\Character{0}}%
\TextAt{24}{322.886}{717.45}{1}%
\Vector{q 0 G 0 g 0.359 w 242.03 -57.918 m 254.539 -57.918 l S Q}%
\TextAt{24}{318.458}{710.964}{2}%
\TextAt{22}{328.93}{717.782}{p}%
\TextAt{23}{332.77}{717.782}{\Character{0}}%
\TextAt{24}{338.527}{717.782}{1}%
\Vector{q 0 G 0 g 0.359 w 256.93 -57.918 m 270.179 -57.918 l S Q}%
\TextAt{24}{333.728}{710.964}{2}%
\TextAt{20}{345.866}{726.01}{\Character{18}}%
\TextAt{13}{356.343}{717.241}{\Character{96}}%
\Vector{q 0 G 0 g 0.478 w 283.861 -61.858 m 289.737 -61.858 l S Q}%
\TextAt{13}{355.861}{700.952}{p}%
\TextAt{20}{362.932}{726.01}{\Character{19}}%
\TextAt{9}{385.68}{709.153}{\Character{40}}%
\TextAt{8}{390.232}{709.153}{mod}%
\TextAt{13}{416.334}{709.153}{p}%
\TextAt{9}{422.21}{709.153}{\Character{41}}%
\TextAt{13}{426.762}{709.153}{\Character{58}}%
\TextAt{8}{79.226}{678.609}{On}%
\TextAt{8}{98.737}{678.609}{the}%
\TextAt{8}{118.893}{678.609}{other}%
\TextAt{8}{149.463}{678.609}{hand}%
\TextAt{8}{178.729}{678.609}{one}%
\TextAt{8}{200.189}{678.609}{has}%
\TextAt{13}{158.518}{651.789}{\Character{28}}%
\TextAt{21}{164.937}{656.726}{p}%
\NativeMathAt{173.018}{651.789}{11.955}{\equiv}%
\TextAt{20}{185.638}{663.147}{\Character{88}}%
\TextAt{21}{192.023}{638.041}{a}%
\TextAt{20}{204.899}{665.06}{\Character{16}}%
\TextAt{13}{213.234}{659.877}{a}%
\Vector{q 0 G 0 g 0.478 w 141.234 -119.222 m 147.379 -119.222 l S Q}%
\TextAt{13}{213.851}{643.589}{\Character{96}}%
\TextAt{20}{220.575}{665.06}{\Character{17}}%
\TextAt{13}{229.707}{651.789}{\Character{16}}%
\TextAt{21}{235.723}{656.726}{ap}%
\TextAt{9}{248.303}{651.789}{\Character{61}}%
\TextAt{20}{260.728}{665.06}{\Character{16}}%
\TextAt{13}{269.063}{659.877}{p}%
\Vector{q 0 G 0 g 0.478 w 197.063 -119.222 m 202.939 -119.222 l S Q}%
\TextAt{13}{269.545}{643.589}{\Character{96}}%
\TextAt{20}{276.134}{665.06}{\Character{17}}%
\TextAt{20}{285.267}{663.147}{\Character{88}}%
\TextAt{21}{291.652}{638.041}{a}%
\TextAt{20}{304.528}{665.06}{\Character{16}}%
\TextAt{13}{312.863}{659.877}{ap}%
\Vector{q 0 G 0 g 0.478 w 240.863 -119.222 m 252.883 -119.222 l S Q}%
\TextAt{13}{316.417}{643.589}{\Character{96}}%
\TextAt{20}{326.079}{665.06}{\Character{17}}%
\TextAt{13}{335.211}{651.789}{\Character{16}}%
\TextAt{21}{341.227}{656.726}{ap}%
\TextAt{9}{353.806}{651.789}{\Character{61}}%
\TextAt{20}{366.232}{665.06}{\Character{16}}%
\TextAt{13}{374.567}{659.877}{p}%
\Vector{q 0 G 0 g 0.478 w 302.567 -119.222 m 308.442 -119.222 l S Q}%
\TextAt{13}{375.049}{643.589}{\Character{96}}%
\TextAt{20}{381.638}{665.06}{\Character{17}}%
\TextAt{13}{390.77}{651.789}{\Character{28}}%
\TextAt{9}{409.144}{651.789}{\Character{40}}%
\TextAt{8}{413.697}{651.789}{mod}%
\TextAt{13}{439.799}{651.789}{p}%
\TextAt{9}{445.674}{651.789}{\Character{41}}%
\TextAt{13}{450.226}{651.789}{\Character{59}}%
\TextAt{8}{79.226}{617.778}{so}%
\TextAt{8}{93.596}{617.778}{that}%
\TextAt{13}{208.622}{599.327}{\Character{28}}%
\TextAt{20}{217.034}{612.598}{\Character{16}}%
\TextAt{13}{225.369}{607.415}{p}%
\Vector{q 0 G 0 g 0.478 w 153.369 -171.684 m 159.244 -171.684 l S Q}%
\TextAt{13}{225.851}{591.127}{\Character{96}}%
\TextAt{20}{232.44}{612.598}{\Character{17}}%
\NativeMathAt{242.9}{599.327}{11.955}{\equiv}%
\TextAt{13}{255.52}{599.327}{\Character{28}}%
\TextAt{9}{261.939}{599.327}{\Character{40}}%
\TextAt{16}{266.491}{599.327}{\Character{0}}%
\TextAt{9}{275.79}{599.327}{1\Character{41}}%
\TextAt{22}{287.39}{607.624}{\Character{96}}%
\TextAt{23}{290.49}{607.624}{\Character{0}}%
\TextAt{24}{296.246}{607.624}{1}%
\Vector{q 0 G 0 g 0.359 w 215.39 -167.744 m 227.899 -167.744 l S Q}%
\TextAt{24}{291.818}{601.138}{2}%
\TextAt{22}{302.29}{607.956}{p}%
\TextAt{23}{306.13}{607.956}{\Character{0}}%
\TextAt{24}{311.887}{607.956}{1}%
\Vector{q 0 G 0 g 0.359 w 230.29 -167.744 m 243.539 -167.744 l S Q}%
\TextAt{24}{307.088}{601.138}{2}%
\TextAt{20}{319.226}{616.184}{\Character{18}}%
\TextAt{13}{329.703}{607.415}{\Character{96}}%
\Vector{q 0 G 0 g 0.478 w 257.221 -171.684 m 263.097 -171.684 l S Q}%
\TextAt{13}{329.221}{591.127}{p}%
\TextAt{20}{336.292}{616.184}{\Character{19}}%
\TextAt{9}{359.04}{599.327}{\Character{40}}%
\TextAt{8}{363.592}{599.327}{mod}%
\TextAt{13}{389.694}{599.327}{p}%
\TextAt{9}{395.569}{599.327}{\Character{41}}%
\TextAt{13}{400.122}{599.327}{\Character{58}}%
\TextAt{8}{79.226}{573.131}{Multiplying}%
\TextAt{8}{142.959}{573.131}{by}%
\TextAt{13}{159.218}{573.131}{\Character{28}}%
\TextAt{8}{169.535}{573.131}{and}%
\TextAt{8}{192.297}{573.131}{dividing}%
\TextAt{8}{237.5}{573.131}{by}%
\NativeMathAt{253.759}{573.131}{11.955}{\pm}%
\TextAt{13}{263.057}{573.131}{\Character{96}}%
\TextAt{8}{271.867}{573.131}{yields}%
\TextAt{8}{304.767}{573.131}{the}%
\TextAt{8}{324.924}{573.131}{claim.}%
\Vector{q 0 G 0 g 0.399 w 452.334 -200.869 m 452.334 -192.8 l S Q}%
\Vector{q 0 G 0 g 0.399 w 452.534 -192.999 m 459.707 -192.999 l S Q}%
\Vector{q 0 G 0 g 0.399 w 452.534 -200.67 m 459.707 -200.67 l S Q}%
\Vector{q 0 G 0 g 0.399 w 459.906 -200.869 m 459.906 -192.8 l S Q}%
\TextAt{8}{97.159}{544.753}{We}%
\TextAt{8}{118.541}{544.753}{have}%
\TextAt{8}{146.773}{544.753}{proved}%
\TextAt{8}{186.064}{544.753}{Gauss’s}%
\TextAt{8}{229.839}{544.753}{reciprocity}%
\TextAt{8}{289.286}{544.753}{law}%
\TextAt{8}{311.66}{544.753}{by}%
\TextAt{8}{329.157}{544.753}{a}%
\TextAt{8}{340.15}{544.753}{rather}%
\TextAt{8}{376.512}{544.753}{contrived}%
\TextAt{8}{428.81}{544.753}{calculation.}%
\TextAt{8}{496.264}{544.753}{In}%
\TextAt{8}{512.123}{544.753}{§10,}%
\TextAt{8}{79.226}{527.707}{however,}%
\TextAt{8}{128.35}{527.707}{we}%
\TextAt{8}{146.372}{527.707}{will}%
\TextAt{8}{169.271}{527.707}{recognize}%
\TextAt{8}{220.791}{527.707}{the}%
\TextAt{8}{241.742}{527.707}{true}%
\TextAt{8}{267.247}{527.707}{reason}%
\TextAt{8}{304.528}{527.707}{why}%
\TextAt{8}{330.034}{527.707}{it}%
\TextAt{8}{342.532}{527.707}{holds}%
\TextAt{8}{373.955}{527.707}{in}%
\TextAt{8}{388.402}{527.707}{the}%
\TextAt{8}{409.353}{527.707}{law}%
\TextAt{8}{431.284}{527.707}{of}%
\TextAt{8}{445.408}{527.707}{decomposition}%
\TextAt{8}{523.337}{527.707}{of}%
\TextAt{8}{79.226}{510.661}{primes}%
\TextAt{8}{115.975}{510.661}{in}%
\TextAt{8}{128.597}{510.661}{the}%
\TextAt{8}{147.724}{510.661}{field}%
\TextAt{18}{172.051}{510.661}{Q}%
\TextAt{9}{181.35}{510.661}{\Character{40}}%
\TextAt{13}{185.902}{510.661}{\Character{16}}%
\TextAt{9}{191.917}{510.661}{\Character{41}}%
\TextAt{8}{199.337}{510.661}{of}%
\TextAt{13}{211.637}{510.661}{\Character{96}}%
\TextAt{8}{216.549}{510.661}{-th}%
\TextAt{8}{234.373}{510.661}{roots}%
\TextAt{8}{263.004}{510.661}{of}%
\TextAt{8}{275.304}{510.661}{unity.}%
\TextAt{8}{309.101}{510.661}{The}%
\TextAt{8}{332.124}{510.661}{Gauss}%
\TextAt{8}{365.765}{510.661}{sums}%
\TextAt{8}{394.12}{510.661}{do}%
\TextAt{8}{409.35}{510.661}{have}%
\TextAt{8}{435.314}{510.661}{a}%
\TextAt{8}{444.04}{510.661}{higher}%
\TextAt{8}{478.78}{510.661}{theoretical}%
\TextAt{8}{79.226}{493.615}{significance,}%
\TextAt{8}{144.322}{493.615}{though,}%
\TextAt{8}{187.253}{493.615}{as}%
\TextAt{8}{201.623}{493.615}{will}%
\TextAt{8}{223.728}{493.615}{become}%
\TextAt{8}{265.667}{493.615}{apparent}%
\TextAt{8}{314.779}{493.615}{later}%
\TextAt{8}{342.096}{493.615}{(see}%
\TextAt{8}{365.564}{493.615}{VII,}%
\TextAt{8}{390.252}{493.615}{§2}%
\TextAt{8}{405.686}{493.615}{and}%
\TextAt{8}{428.449}{493.615}{§6).}%
\TextAt{12}{79.226}{468.226}{Exercise}%
\TextAt{12}{131.578}{468.226}{I.8.1.}%
\TextAt{8}{164.434}{468.226}{If}%
\TextAt{14}{175.389}{468.226}{a}%
\TextAt{8}{184.532}{468.226}{and}%
\TextAt{14}{206.557}{468.226}{b}%
\TextAt{8}{215.851}{468.226}{are}%
\TextAt{8}{234.624}{468.226}{ideals}%
\TextAt{8}{266.464}{468.226}{of}%
\TextAt{14}{279.057}{468.226}{o}%
\TextAt{8}{284.901}{468.226}{,}%
\TextAt{8}{291.46}{468.226}{then}%
\TextAt{8}{317.382}{468.226}{one}%
\TextAt{8}{338.104}{468.226}{has}%
\TextAt{14}{358.24}{468.226}{a}%
\TextAt{9}{367.544}{468.226}{\Character{61}}%
\TextAt{14}{379.97}{468.226}{a}%
\TextAt{16}{385.953}{468.226}{\Character{79}}%
\TextAt{16}{396.94}{468.226}{\Character{92}}%
\TextAt{14}{406.069}{468.226}{o}%
\TextAt{8}{415.073}{468.226}{and}%
\TextAt{14}{437.098}{468.226}{a}%
\TextAt{16}{446.402}{468.226}{\Character{106}}%
\TextAt{14}{453.043}{468.226}{b}%
\NativeArrowAt{462.499}{468.226}{11.955}{21.926}{\Longleftrightarrow}%
\TextAt{14}{487.738}{468.226}{a}%
\TextAt{16}{493.721}{468.226}{\Character{79}}%
\TextAt{16}{506.884}{468.226}{\Character{106}}%
\TextAt{14}{513.533}{468.226}{b}%
\TextAt{16}{519.668}{468.226}{\Character{79}}%
\TextAt{8}{529.518}{468.226}{.}%
\ErrataAt{79.22}{437.455}{453.64}{\textbf{Exercise I.8.2.} Let $\mathfrak A\ne0$ be an integral ideal of $\mathcal O$. Suppose that $\mathcal O/\mathfrak p\mathcal O$ is}%
\ErrataAt{79.22}{420.455}{453.64}{generated by one element over $\mathfrak o/\mathfrak p$ for each base prime $\mathfrak p$ below $\mathfrak A$. Show that}%
\ErrataAt{79.22}{406.055}{453.64}{some $\theta\in\mathcal O$ generates $L$ over $K$ and has conductor prime to $\mathfrak A$.}%
\TextAt{12}{79.226}{389.693}{Exercise}%
\TextAt{12}{132.332}{389.693}{I.8.3.}%
\TextAt{8}{165.843}{389.693}{If}%
\TextAt{8}{177.452}{389.693}{a}%
\TextAt{8}{187.125}{389.693}{prime}%
\TextAt{8}{220.206}{389.693}{ideal}%
\TextAt{14}{248.087}{389.693}{p}%
\TextAt{8}{257.884}{389.693}{of}%
\TextAt{13}{271.132}{389.693}{K}%
\TextAt{8}{285.764}{389.693}{is}%
\TextAt{8}{297.445}{389.693}{totally}%
\TextAt{8}{334.77}{389.693}{split}%
\TextAt{8}{360.762}{389.693}{in}%
\TextAt{8}{374.332}{389.693}{two}%
\TextAt{8}{396.366}{389.693}{separable}%
\TextAt{8}{447.727}{389.693}{extensions}%
\TextAt{13}{504.025}{389.693}{L}%
\TextAt{16}{515.31}{389.693}{\Character{106}}%
\TextAt{13}{521.952}{389.693}{K}%
\TextAt{8}{79.226}{372.647}{and}%
\TextAt{13}{101.989}{372.647}{L}%
\TextAt{3}{109.953}{376.985}{\Character{48}}%
\TextAt{16}{116.069}{372.647}{\Character{106}}%
\TextAt{13}{122.711}{372.647}{K}%
\TextAt{8}{133.529}{372.647}{,}%
\TextAt{8}{140.678}{372.647}{then}%
\TextAt{8}{167.338}{372.647}{it}%
\TextAt{8}{179.042}{372.647}{is}%
\TextAt{8}{190.806}{372.647}{also}%
\TextAt{8}{214.286}{372.647}{totally}%
\TextAt{8}{251.694}{372.647}{split}%
\TextAt{8}{277.768}{372.647}{in}%
\TextAt{8}{291.421}{372.647}{the}%
\TextAt{8}{311.577}{372.647}{composite}%
\TextAt{8}{366.595}{372.647}{extension.}%
\TextAt{12}{79.226}{341.904}{Exercise}%
\TextAt{12}{132.62}{341.904}{I.8.4.}%
\TextAt{8}{166.383}{341.904}{A}%
\TextAt{8}{179.223}{341.904}{prime}%
\TextAt{8}{212.555}{341.904}{ideal}%
\TextAt{14}{240.687}{341.904}{p}%
\TextAt{8}{250.735}{341.904}{of}%
\TextAt{13}{264.234}{341.904}{K}%
\TextAt{8}{279.117}{341.904}{is}%
\TextAt{8}{291.049}{341.904}{totally}%
\TextAt{8}{328.625}{341.904}{split}%
\TextAt{8}{354.868}{341.904}{in}%
\TextAt{8}{368.689}{341.904}{the}%
\TextAt{8}{389.014}{341.904}{separable}%
\TextAt{8}{440.625}{341.904}{extension}%
\TextAt{13}{492.559}{341.904}{L}%
\TextAt{16}{504.131}{341.904}{\Character{106}}%
\TextAt{13}{511.06}{341.904}{K}%
\TextAt{8}{525.943}{341.904}{if}%
\TextAt{8}{79.226}{324.858}{and}%
\TextAt{8}{101.989}{324.858}{only}%
\TextAt{8}{127.68}{324.858}{if}%
\TextAt{8}{138.404}{324.858}{it}%
\TextAt{8}{150.108}{324.858}{is}%
\TextAt{8}{161.872}{324.858}{totally}%
\TextAt{8}{199.28}{324.858}{split}%
\TextAt{8}{225.354}{324.858}{in}%
\TextAt{8}{239.007}{324.858}{the}%
\TextAt{8}{259.163}{324.858}{Galois}%
\TextAt{8}{295.077}{324.858}{closure}%
\TextAt{13}{334.158}{324.858}{N}%
\TextAt{16}{348.102}{324.858}{\Character{106}}%
\TextAt{13}{354.743}{324.858}{K}%
\TextAt{8}{369.458}{324.858}{of}%
\TextAt{13}{382.788}{324.858}{L}%
\TextAt{16}{394.074}{324.858}{\Character{106}}%
\TextAt{13}{400.716}{324.858}{K}%
\TextAt{8}{411.533}{324.858}{.}%
\TextAt{12}{79.226}{294.114}{Exercise}%
\TextAt{12}{132.531}{294.114}{I.8.5.}%
\TextAt{8}{166.216}{294.114}{For}%
\TextAt{8}{187.277}{294.114}{a}%
\TextAt{8}{197.123}{294.114}{number}%
\TextAt{8}{239.811}{294.114}{field}%
\TextAt{13}{265.259}{294.114}{K}%
\TextAt{8}{280.065}{294.114}{the}%
\TextAt{8}{300.312}{294.114}{statement}%
\TextAt{8}{354.775}{294.114}{of}%
\TextAt{8}{368.196}{294.114}{proposition}%
\TextAt{8}{429.821}{294.114}{(8.3)}%
\TextAt{8}{457.887}{294.114}{concerning}%
\TextAt{8}{516.51}{294.114}{the}%
\TextAt{8}{79.226}{277.068}{prime}%
\TextAt{8}{112.39}{277.068}{decomposition}%
\TextAt{8}{189.525}{277.068}{in}%
\TextAt{8}{203.177}{277.068}{the}%
\TextAt{8}{223.334}{277.068}{extension}%
\TextAt{13}{275.1}{277.068}{K}%
\TextAt{9}{285.917}{277.068}{\Character{40}}%
\TextAt{13}{290.47}{277.068}{\Character{18}}%
\TextAt{9}{296.25}{277.068}{\Character{41}}%
\TextAt{8}{304.7}{277.068}{holds}%
\TextAt{8}{335.329}{277.068}{for}%
\TextAt{8}{353.214}{277.068}{all}%
\TextAt{8}{369.473}{277.068}{prime}%
\TextAt{8}{402.636}{277.068}{ideals}%
\TextAt{14}{435.214}{277.068}{p}%
\TextAt{18}{444.518}{277.068}{\Character{45}}%
\TextAt{9}{451.16}{277.068}{\Character{40}}%
\TextAt{16}{455.712}{277.068}{\Character{79}}%
\TextAt{9}{468.883}{277.068}{\Character{58}}%
\TextAt{14}{475.455}{277.068}{o}%
\TextAt{9}{481.3}{277.068}{\Character{91}}%
\TextAt{13}{484.551}{277.068}{\Character{18}}%
\TextAt{9}{490.332}{277.068}{\Character{93}\Character{41}}%
\TextAt{8}{498.136}{277.068}{.}%
\TextAt{12}{79.226}{246.325}{Exercise}%
\TextAt{12}{133.487}{246.325}{I.8.6.}%
\TextAt{8}{168.002}{246.325}{Given}%
\TextAt{8}{202.816}{246.325}{a}%
\TextAt{8}{213.493}{246.325}{positive}%
\TextAt{8}{257.728}{246.325}{integer}%
\TextAt{13}{297.349}{246.325}{b}%
\TextAt{13}{307.212}{246.325}{\Character{62}}%
\TextAt{9}{321.213}{246.325}{1}%
\TextAt{8}{327.066}{246.325}{,}%
\TextAt{8}{335.367}{246.325}{an}%
\TextAt{8}{352.547}{246.325}{integer}%
\TextAt{13}{392.168}{246.325}{a}%
\TextAt{8}{403.131}{246.325}{relatively}%
\TextAt{8}{455.113}{246.325}{prime}%
\TextAt{8}{489.198}{246.325}{to}%
\TextAt{13}{504.43}{246.325}{b}%
\TextAt{8}{514.226}{246.325}{is}%
\TextAt{8}{526.911}{246.325}{a}%
\TextAt{8}{79.226}{229.279}{quadratic}%
\TextAt{8}{131.3}{229.279}{residue}%
\TextAt{8}{170.737}{229.279}{mod}%
\TextAt{13}{196.784}{229.279}{b}%
\TextAt{8}{205.368}{229.279}{if}%
\TextAt{8}{215.802}{229.279}{and}%
\TextAt{8}{238.274}{229.279}{only}%
\TextAt{8}{263.676}{229.279}{if}%
\TextAt{8}{274.109}{229.279}{it}%
\TextAt{8}{285.523}{229.279}{is}%
\TextAt{8}{296.997}{229.279}{a}%
\TextAt{8}{306.462}{229.279}{quadratic}%
\TextAt{8}{358.536}{229.279}{residue}%
\TextAt{8}{397.973}{229.279}{modulo}%
\TextAt{8}{439.633}{229.279}{each}%
\TextAt{8}{465.68}{229.279}{prime}%
\TextAt{8}{498.554}{229.279}{divisor}%
\TextAt{13}{79.226}{212.233}{p}%
\TextAt{8}{88.999}{212.233}{of}%
\TextAt{13}{102.329}{212.233}{b}%
\TextAt{8}{107.306}{212.233}{,}%
\TextAt{8}{114.455}{212.233}{and}%
\TextAt{8}{137.218}{212.233}{if}%
\TextAt{13}{147.941}{212.233}{a}%
\NativeMathAt{157.407}{212.233}{11.955}{\equiv}%
\TextAt{9}{170.027}{212.233}{1}%
\TextAt{9}{181.181}{212.233}{\Character{40}}%
\TextAt{8}{185.745}{212.233}{mod}%
\TextAt{9}{211.847}{212.233}{4\Character{41}}%
\TextAt{8}{226.15}{212.233}{when}%
\TextAt{9}{256.707}{212.233}{4}%
\TextAt{16}{265.881}{212.233}{\Character{106}}%
\TextAt{13}{272.523}{212.233}{b}%
\TextAt{8}{277.5}{212.233}{,}%
\TextAt{9}{284.649}{212.233}{8}%
\TextAt{18}{293.823}{212.233}{\Character{45}}%
\TextAt{13}{300.465}{212.233}{b}%
\TextAt{8}{305.442}{212.233}{,}%
\TextAt{8}{312.591}{212.233}{resp.}%
\TextAt{13}{341.913}{212.233}{a}%
\NativeMathAt{351.379}{212.233}{11.955}{\equiv}%
\TextAt{9}{363.998}{212.233}{1}%
\TextAt{9}{375.152}{212.233}{\Character{40}}%
\TextAt{8}{379.717}{212.233}{mod}%
\TextAt{9}{405.819}{212.233}{8\Character{41}}%
\TextAt{8}{420.122}{212.233}{when}%
\TextAt{9}{450.679}{212.233}{8}%
\TextAt{16}{459.853}{212.233}{\Character{106}}%
\TextAt{13}{466.495}{212.233}{b}%
\TextAt{8}{471.472}{212.233}{.}%
\TextAt{12}{79.226}{181.489}{Exercise}%
\TextAt{12}{132.559}{181.489}{I.8.7.}%
\TextAt{8}{166.268}{181.489}{Let}%
\TextAt{9}{187.352}{181.489}{\Character{40}}%
\TextAt{13}{191.905}{181.489}{a\Character{59}}%
\TextAt{13}{203.274}{181.489}{p}%
\TextAt{9}{209.169}{181.489}{\Character{41}}%
\TextAt{9}{217.227}{181.489}{\Character{61}}%
\TextAt{9}{229.84}{181.489}{1}%
\TextAt{8}{239.725}{181.489}{and}%
\TextAt{13}{262.602}{181.489}{av}%
\NativeMathAt{278.352}{181.489}{11.955}{\equiv}%
\TextAt{13}{291.168}{181.489}{r}%
\TextAt{21}{296.443}{179.696}{v}%
\TextAt{9}{306.664}{181.489}{\Character{40}}%
\TextAt{8}{311.216}{181.489}{mod}%
\TextAt{13}{337.318}{181.489}{p}%
\TextAt{9}{343.193}{181.489}{\Character{41}}%
\TextAt{8}{347.746}{181.489}{,}%
\TextAt{13}{355.039}{181.489}{v}%
\TextAt{9}{364.644}{181.489}{\Character{61}}%
\TextAt{9}{377.257}{181.489}{1}%
\TextAt{13}{383.118}{181.489}{\Character{59}}%
\TextAt{13}{388.342}{181.489}{\Character{58}}%
\TextAt{13}{393.579}{181.489}{\Character{58}}%
\TextAt{13}{398.815}{181.489}{\Character{58}}%
\TextAt{13}{404.04}{181.489}{\Character{59}}%
\TextAt{13}{409.276}{181.489}{p}%
\TextAt{16}{417.949}{181.489}{\Character{0}}%
\TextAt{9}{429.983}{181.489}{1}%
\TextAt{8}{435.836}{181.489}{,}%
\TextAt{9}{443.129}{181.489}{0}%
\TextAt{13}{452.498}{181.489}{\Character{60}}%
\TextAt{13}{465.111}{181.489}{r}%
\TextAt{21}{470.395}{179.696}{v}%
\TextAt{13}{478.82}{181.489}{\Character{60}}%
\TextAt{13}{491.433}{181.489}{p}%
\TextAt{8}{497.316}{181.489}{.}%
\TextAt{8}{506.11}{181.489}{Then}%
\TextAt{8}{79.226}{164.443}{the}%
\TextAt{13}{99.383}{164.443}{r}%
\TextAt{21}{104.658}{162.65}{v}%
\TextAt{8}{113.463}{164.443}{give}%
\TextAt{8}{137.529}{164.443}{a}%
\TextAt{8}{147.284}{164.443}{permutation}%
\TextAt{13}{214.281}{164.443}{\Character{25}}%
\TextAt{8}{225.247}{164.443}{of}%
\TextAt{8}{238.577}{164.443}{the}%
\TextAt{8}{258.734}{164.443}{numbers}%
\TextAt{9}{305.945}{164.443}{1}%
\TextAt{13}{311.798}{164.443}{\Character{59}}%
\TextAt{13}{317.022}{164.443}{\Character{58}}%
\TextAt{13}{322.259}{164.443}{\Character{58}}%
\TextAt{13}{327.483}{164.443}{\Character{58}}%
\TextAt{13}{332.72}{164.443}{\Character{59}}%
\TextAt{13}{337.956}{164.443}{p}%
\TextAt{16}{346.55}{164.443}{\Character{0}}%
\TextAt{9}{358.506}{164.443}{1}%
\TextAt{8}{364.359}{164.443}{.}%
\TextAt{8}{372.807}{164.443}{Show}%
\TextAt{8}{403.699}{164.443}{that}%
\TextAt{8}{429.068}{164.443}{sgn}%
\TextAt{13}{448.037}{164.443}{\Character{25}}%
\TextAt{9}{458.427}{164.443}{\Character{61}}%
\TextAt{20}{470.852}{177.714}{\Character{16}}%
\TextAt{21}{479.188}{169.151}{a}%
\Vector{q 0 G 0 g 0.478 w 407.188 -606.568 m 411.686 -606.568 l S Q}%
\TextAt{21}{479.305}{160.321}{p}%
\TextAt{20}{484.881}{177.714}{\Character{17}}%
\TextAt{8}{492.021}{164.443}{.}%
\TextAt{12}{79.226}{123.045}{Exercise}%
\TextAt{12}{133.982}{123.045}{I.8.8.}%
\TextAt{8}{168.928}{123.045}{Let}%
\TextAt{13}{191.25}{123.045}{a}%
\TextAt{21}{197.395}{121.252}{n}%
\TextAt{9}{208.656}{123.045}{\Character{61}}%
\TextAt{13}{224.582}{131.133}{\Character{34}}%
\TextAt{21}{230.06}{135.471}{n}%
\TextAt{16}{238.353}{131.133}{\Character{0}}%
\TextAt{13}{250.308}{131.133}{\Character{34}}%
\TextAt{3}{255.786}{135.471}{\Character{48}}%
\TextAt{21}{258.083}{135.471}{n}%
\Vector{q 0 G 0 g 0.478 w 152.582 -647.966 m 191.719 -647.966 l S Q}%
\TextAt{16}{236.243}{123.404}{\Character{112}}%
\Vector{q 0 G 0 g 0.478 w 174.205 -650.357 m 180.058 -650.357 l S Q}%
\TextAt{9}{246.205}{113.514}{5}%
\TextAt{8}{264.915}{123.045}{,}%
\TextAt{8}{273.754}{123.045}{where}%
\TextAt{13}{308.916}{123.045}{\Character{34}}%
\TextAt{9}{320.019}{123.045}{\Character{61}}%
\TextAt{9}{335.944}{131.133}{1}%
\TextAt{9}{344.444}{131.133}{\Character{43}}%
\TextAt{16}{356.215}{141.022}{\Character{112}}%
\Vector{q 0 G 0 g 0.478 w 294.178 -632.739 m 300.031 -632.739 l S Q}%
\TextAt{9}{366.178}{131.133}{5}%
\Vector{q 0 G 0 g 0.478 w 263.944 -647.966 m 300.031 -647.966 l S Q}%
\TextAt{9}{351.061}{114.844}{2}%
\TextAt{8}{373.227}{123.045}{,}%
\TextAt{13}{382.066}{123.045}{\Character{34}}%
\TextAt{3}{387.544}{127.383}{\Character{48}}%
\TextAt{9}{395.964}{123.045}{\Character{61}}%
\TextAt{9}{411.89}{131.133}{1}%
\TextAt{16}{420.399}{131.133}{\Character{0}}%
\TextAt{16}{432.355}{141.022}{\Character{112}}%
\Vector{q 0 G 0 g 0.478 w 370.317 -632.739 m 376.17 -632.739 l S Q}%
\TextAt{9}{442.317}{131.133}{5}%
\Vector{q 0 G 0 g 0.478 w 339.89 -647.966 m 376.17 -647.966 l S Q}%
\TextAt{9}{427.104}{114.844}{2}%
\TextAt{8}{454.615}{123.045}{(}%
\TextAt{13}{459.17}{123.045}{a}%
\TextAt{21}{465.315}{121.252}{n}%
\TextAt{8}{476.201}{123.045}{is}%
\TextAt{8}{489.317}{123.045}{the}%
\TextAt{13}{510.826}{123.045}{n}%
\TextAt{8}{517.814}{123.045}{-th}%
\TextAt{8}{79.226}{101.964}{Fibonacci}%
\TextAt{8}{132.714}{101.964}{number).}%
\TextAt{8}{184.416}{101.964}{If}%
\TextAt{13}{196.108}{101.964}{p}%
\TextAt{8}{205.88}{101.964}{is}%
\TextAt{8}{217.644}{101.964}{a}%
\TextAt{8}{227.4}{101.964}{prime}%
\TextAt{8}{260.563}{101.964}{number}%
\NativeMathAt{303.16}{101.964}{11.955}{\neq}%
\TextAt{9}{315.57}{101.964}{2}%
\TextAt{13}{321.438}{101.964}{\Character{59}}%
\TextAt{9}{326.682}{101.964}{5}%
\TextAt{8}{332.535}{101.964}{,}%
\TextAt{8}{339.684}{101.964}{then}%
\TextAt{8}{366.344}{101.964}{one}%
\TextAt{8}{387.804}{101.964}{has}%
\TextAt{13}{252.161}{75.144}{a}%
\TextAt{21}{258.306}{73.351}{p}%
\NativeMathAt{266.388}{75.144}{11.955}{\equiv}%
\TextAt{20}{279.007}{88.415}{\Character{16}}%
\TextAt{13}{287.343}{83.232}{p}%
\Vector{q 0 G 0 g 0.478 w 215.343 -695.867 m 221.218 -695.867 l S Q}%
\TextAt{9}{287.354}{66.944}{5}%
\TextAt{20}{294.413}{88.415}{\Character{17}}%
\TextAt{9}{315.501}{75.144}{\Character{40}}%
\TextAt{8}{320.053}{75.144}{mod}%
\TextAt{13}{346.155}{75.144}{p}%
\TextAt{9}{352.03}{75.144}{\Character{41}}%
\TextAt{13}{356.583}{75.144}{\Character{58}}%
\TextAt{8}{300.14}{43.033}{49}%
\Destination{78.23}{809.815}{nav.page.49}%
\Destination{79.226}{482.255}{nav.ex.1.8.1}%
\Destination{79.226}{451.455}{nav.ex.1.8.2}%
\Destination{79.226}{403.655}{nav.ex.1.8.3}%
\Destination{79.226}{355.855}{nav.ex.1.8.4}%
\Destination{79.226}{308.155}{nav.ex.1.8.5}%
\Destination{79.226}{260.355}{nav.ex.1.8.6}%
\Destination{79.226}{195.455}{nav.ex.1.8.7}%
\Destination{79.226}{137.055}{nav.ex.1.8.8}%
\Destination{79.36}{774.355}{nav.print.52}%
\Destination{79.36}{404.655}{nav.print.53}%
\Hotlink{517.302}{542.555}{12.716}{11.047}{nav.section.1.10}%
\Hotlink{395.431}{491.455}{6.858}{11.047}{nav.section.7.2}%
\Hotlink{433.628}{491.455}{6.858}{11.047}{nav.section.7.6}%
\Hotlink{429.321}{291.955}{25.078}{11.047}{nav.number.1.8.3}%
\Destination{79.226}{451.455}{errata.7}%
\Hotlink{365.774}{491.4}{16.57}{11.047}{nav.chapter.7}%
}

% Page 62 (body page 50)
\Page{%
\TextAt{2}{79.226}{796.943}{I.}%
\TextAt{2}{89.83}{796.943}{Algebraic}%
\TextAt{2}{138.648}{796.943}{Integers}%
\Vector{q 0 G 0 g 0.299 w 7.226 18.729 m 460.769 18.729 l S Q}%
\TextAt{12}{79.226}{759.33}{Exercise}%
\TextAt{12}{132.754}{759.33}{I.8.9.}%
\TextAt{8}{166.633}{759.33}{Study}%
\TextAt{8}{201.062}{759.33}{the}%
\TextAt{8}{221.504}{759.33}{Legendre}%
\TextAt{8}{272.024}{759.33}{symbol}%
\TextAt{20}{312.371}{772.601}{\Character{16}}%
\TextAt{17}{320.721}{764.037}{3}%
\Vector{q 0 G 0 g 0.478 w 248.707 -11.681 m 252.969 -11.681 l S Q}%
\TextAt{21}{320.707}{755.208}{p}%
\TextAt{20}{326.165}{772.601}{\Character{17}}%
\TextAt{8}{337.487}{759.33}{as}%
\TextAt{8}{352.142}{759.33}{a}%
\TextAt{8}{362.183}{759.33}{function}%
\TextAt{8}{408.316}{759.33}{of}%
\TextAt{13}{421.931}{759.33}{p}%
\TextAt{13}{431.603}{759.33}{\Character{62}}%
\TextAt{9}{444.523}{759.33}{3}%
\TextAt{8}{450.376}{759.33}{.}%
\TextAt{8}{459.679}{759.33}{Show}%
\TextAt{8}{490.857}{759.33}{that}%
\TextAt{8}{516.51}{759.33}{the}%
\TextAt{8}{79.226}{742.284}{property}%
\TextAt{8}{126.787}{742.284}{of}%
\TextAt{9}{139.869}{742.284}{3}%
\TextAt{8}{149.371}{742.284}{being}%
\TextAt{8}{180.661}{742.284}{a}%
\TextAt{8}{190.168}{742.284}{quadratic}%
\TextAt{8}{242.284}{742.284}{residue}%
\TextAt{8}{281.763}{742.284}{or}%
\TextAt{8}{295.825}{742.284}{nonresidue}%
\TextAt{8}{354.169}{742.284}{mod}%
\TextAt{13}{380.258}{742.284}{p}%
\TextAt{8}{389.783}{742.284}{depends}%
\TextAt{8}{434.785}{742.284}{only}%
\TextAt{8}{460.228}{742.284}{on}%
\TextAt{8}{476.239}{742.284}{the}%
\TextAt{8}{496.148}{742.284}{class}%
\TextAt{8}{523.337}{742.284}{of}%
\TextAt{13}{79.226}{725.238}{p}%
\TextAt{8}{88.999}{725.238}{mod}%
\TextAt{9}{115.336}{725.238}{12}%
\TextAt{8}{127.042}{725.238}{.}%
\ErrataAt{79.19}{710.8}{455.81}{In Exercises I.8.10 and I.8.11, take $p$ odd and set $(u/p)=0$ when $p\mid u$.}%
\TextAt{12}{79.226}{694.441}{Exercise}%
\TextAt{12}{132.427}{694.441}{I.8.10.}%
\TextAt{8}{172.751}{694.441}{Show}%
\TextAt{8}{203.644}{694.441}{that}%
\TextAt{8}{229.012}{694.441}{the}%
\TextAt{8}{249.169}{694.441}{number}%
\TextAt{8}{291.765}{694.441}{of}%
\TextAt{8}{305.095}{694.441}{solutions}%
\TextAt{8}{354.004}{694.441}{of}%
\TextAt{13}{367.334}{694.441}{x}%
\TextAt{17}{373.986}{698.779}{2}%
\NativeMathAt{382.039}{694.441}{11.955}{\equiv}%
\TextAt{13}{394.658}{694.441}{a}%
\TextAt{9}{406.117}{694.441}{\Character{40}}%
\TextAt{8}{410.669}{694.441}{mod}%
\TextAt{13}{436.771}{694.441}{p}%
\TextAt{9}{442.646}{694.441}{\Character{41}}%
\TextAt{8}{451.096}{694.441}{equals}%
\TextAt{9}{486.603}{694.441}{1}%
\TextAt{9}{495.103}{694.441}{\Character{43}}%
\TextAt{20}{506.874}{707.711}{\Character{16}}%
\TextAt{21}{515.209}{699.148}{a}%
\Vector{q 0 G 0 g 0.478 w 443.209 -76.57 m 447.707 -76.57 l S Q}%
\TextAt{21}{515.327}{690.318}{p}%
\TextAt{20}{520.903}{707.711}{\Character{17}}%
\TextAt{8}{528.043}{694.441}{.}%
\TextAt{12}{79.226}{662.447}{Exercise}%
\TextAt{12}{132.425}{662.447}{I.8.11.}%
\TextAt{8}{172.748}{662.447}{Show}%
\TextAt{8}{203.639}{662.447}{that}%
\TextAt{8}{229.007}{662.447}{the}%
\TextAt{8}{249.162}{662.447}{number}%
\TextAt{8}{291.756}{662.447}{of}%
\TextAt{8}{305.085}{662.447}{solutions}%
\TextAt{8}{353.992}{662.447}{of}%
\TextAt{8}{367.321}{662.447}{the}%
\TextAt{8}{387.476}{662.447}{congruence}%
\TextAt{13}{447.956}{662.447}{ax}%
\TextAt{17}{460.753}{666.786}{2}%
\TextAt{9}{468.139}{662.447}{\Character{43}}%
\TextAt{13}{479.897}{662.447}{bx}%
\TextAt{9}{494.18}{662.447}{\Character{43}}%
\TextAt{13}{505.938}{662.447}{c}%
\NativeMathAt{514.297}{662.447}{11.955}{\equiv}%
\TextAt{9}{526.917}{662.447}{0}%
\TextAt{9}{79.226}{645.401}{\Character{40}}%
\TextAt{8}{83.778}{645.401}{mod}%
\TextAt{13}{109.881}{645.401}{p}%
\TextAt{9}{115.756}{645.401}{\Character{41}}%
\TextAt{8}{120.308}{645.401}{,}%
\TextAt{8}{127.457}{645.401}{where}%
\TextAt{9}{161.266}{645.401}{\Character{40}}%
\TextAt{13}{165.819}{645.401}{a\Character{59}}%
\TextAt{13}{177.188}{645.401}{p}%
\TextAt{9}{183.083}{645.401}{\Character{41}}%
\TextAt{9}{190.938}{645.401}{\Character{61}}%
\TextAt{9}{203.359}{645.401}{1}%
\TextAt{8}{209.235}{645.401}{,}%
\TextAt{8}{216.384}{645.401}{equals}%
\TextAt{9}{251.891}{645.401}{1}%
\TextAt{9}{260.391}{645.401}{\Character{43}}%
\TextAt{20}{272.162}{658.671}{\Character{16}}%
\TextAt{21}{280.497}{650.108}{b}%
\TextAt{24}{284.12}{652.921}{2}%
\TextAt{3}{288.271}{650.108}{\Character{0}}%
\TextAt{17}{294.857}{650.108}{4}%
\TextAt{21}{299.091}{650.108}{ac}%
\Vector{q 0 G 0 g 0.478 w 208.497 -125.61 m 235.257 -125.61 l S Q}%
\TextAt{21}{291.746}{641.278}{p}%
\TextAt{20}{308.453}{658.671}{\Character{17}}%
\TextAt{8}{315.593}{645.401}{.}%
\TextAt{15}{79.226}{593.442}{§}%
\TextAt{15}{93.785}{593.442}{9}%
\TextAt{15}{120.693}{593.442}{Hilbert’s}%
\TextAt{15}{201.692}{593.442}{Ramification}%
\TextAt{15}{315.71}{593.442}{Theory}%
\TextAt{8}{79.226}{561.556}{The}%
\TextAt{8}{104.727}{561.556}{question}%
\TextAt{8}{152.74}{561.556}{of}%
\TextAt{8}{167.518}{561.556}{prime}%
\TextAt{8}{202.129}{561.556}{decomposition}%
\TextAt{8}{280.711}{561.556}{in}%
\TextAt{8}{295.812}{561.556}{a}%
\TextAt{8}{307.015}{561.556}{finite}%
\TextAt{8}{338.374}{561.556}{extension}%
\TextAt{13}{391.587}{561.556}{L}%
\TextAt{16}{405.34}{561.556}{\Character{106}}%
\TextAt{13}{414.448}{561.556}{K}%
\TextAt{8}{430.611}{561.556}{takes}%
\TextAt{8}{462.042}{561.556}{a}%
\TextAt{8}{473.245}{561.556}{particularly}%
\TextAt{8}{79.226}{544.51}{interesting}%
\TextAt{8}{138.072}{544.51}{and}%
\TextAt{8}{162.057}{544.51}{important}%
\TextAt{8}{218.572}{544.51}{turn}%
\TextAt{8}{245.809}{544.51}{once}%
\TextAt{8}{273.691}{544.51}{we}%
\TextAt{8}{292.141}{544.51}{assume}%
\TextAt{13}{333.807}{544.51}{L}%
\TextAt{16}{347.175}{544.51}{\Character{106}}%
\TextAt{13}{355.9}{544.51}{K}%
\TextAt{8}{371.837}{544.51}{to}%
\TextAt{8}{387.37}{544.51}{be}%
\TextAt{8}{404.516}{544.51}{a}%
\TextAt{8}{415.494}{544.51}{Galois}%
\TextAt{8}{452.63}{544.51}{extension.}%
\TextAt{8}{512.613}{544.51}{The}%
\TextAt{8}{79.226}{527.464}{prime}%
\TextAt{8}{112.39}{527.464}{ideals}%
\TextAt{8}{144.968}{527.464}{are}%
\TextAt{8}{164.478}{527.464}{then}%
\TextAt{8}{191.138}{527.464}{subject}%
\TextAt{8}{231.511}{527.464}{to}%
\TextAt{8}{245.821}{527.464}{the}%
\TextAt{8}{265.978}{527.464}{action}%
\TextAt{8}{301.102}{527.464}{of}%
\TextAt{8}{314.432}{527.464}{the}%
\TextAt{8}{334.589}{527.464}{Galois}%
\TextAt{8}{370.502}{527.464}{group}%
\TextAt{13}{268.341}{500.549}{G}%
\TextAt{9}{280.895}{500.549}{\Character{61}}%
\TextAt{13}{293.32}{500.549}{G}%
\TextAt{9}{302.554}{500.549}{\Character{40}}%
\TextAt{13}{307.106}{500.549}{L}%
\TextAt{16}{318.392}{500.549}{\Character{106}}%
\TextAt{13}{325.033}{500.549}{K}%
\TextAt{9}{335.851}{500.549}{\Character{41}}%
\TextAt{13}{340.403}{500.549}{\Character{58}}%
\TextAt{8}{79.226}{473.635}{The}%
\TextAt{8}{105.1}{473.635}{“ramification}%
\TextAt{8}{177.492}{473.635}{theory”}%
\TextAt{8}{221.586}{473.635}{that}%
\TextAt{8}{248.775}{473.635}{arises}%
\TextAt{8}{282.588}{473.635}{from}%
\TextAt{8}{312.049}{473.635}{this}%
\TextAt{8}{336.692}{473.635}{assumption}%
\TextAt{8}{400.428}{473.635}{has}%
\TextAt{8}{423.122}{473.635}{been}%
\TextAt{8}{452.571}{473.635}{introduced}%
\TextAt{8}{512.924}{473.635}{into}%
\TextAt{8}{79.226}{456.589}{number}%
\TextAt{8}{120.564}{456.589}{theory}%
\TextAt{8}{156.055}{456.589}{by}%
\TextAt{26}{171.056}{456.589}{David}%
\TextAt{26}{208.343}{456.589}{Hilbert}%
\TextAt{8}{257.523}{456.589}{(1862–1943).}%
\TextAt{8}{327.384}{456.589}{Given}%
\TextAt{13}{360.018}{456.589}{a}%
\TextAt{8}{368.802}{456.589}{in}%
\TextAt{8}{381.196}{456.589}{the}%
\TextAt{8}{400.094}{456.589}{ring}%
\TextAt{16}{422.901}{456.589}{\Character{79}}%
\TextAt{8}{435.389}{456.589}{of}%
\TextAt{8}{447.461}{456.589}{integral}%
\TextAt{8}{488.81}{456.589}{elements}%
\TextAt{8}{79.226}{439.543}{of}%
\TextAt{13}{92.964}{439.543}{L}%
\TextAt{8}{100.929}{439.543}{,}%
\TextAt{8}{108.588}{439.543}{the}%
\TextAt{8}{129.153}{439.543}{conjugate}%
\TextAt{13}{182.57}{439.543}{\Character{27}}%
\TextAt{13}{189.636}{439.543}{a}%
\TextAt{8}{195.797}{439.543}{,}%
\TextAt{8}{203.457}{439.543}{for}%
\TextAt{8}{221.75}{439.543}{every}%
\TextAt{13}{253.05}{439.543}{\Character{27}}%
\NativeMathAt{264.149}{439.543}{11.955}{\in}%
\TextAt{13}{276.135}{439.543}{G}%
\TextAt{8}{285.369}{439.543}{,}%
\TextAt{8}{293.028}{439.543}{also}%
\TextAt{8}{316.916}{439.543}{belongs}%
\TextAt{8}{359.335}{439.543}{to}%
\TextAt{16}{374.054}{439.543}{\Character{79}}%
\TextAt{8}{383.903}{439.543}{,}%
\TextAt{8}{391.563}{439.543}{i.e.,}%
\TextAt{13}{414.178}{439.543}{G}%
\TextAt{8}{427.717}{439.543}{acts}%
\TextAt{8}{452.251}{439.543}{on}%
\TextAt{16}{468.918}{439.543}{\Character{79}}%
\TextAt{8}{478.768}{439.543}{.}%
\TextAt{8}{488.44}{439.543}{If}%
\TextAt{14}{500.541}{439.543}{P}%
\TextAt{8}{514.739}{439.543}{is}%
\TextAt{8}{526.911}{439.543}{a}%
\TextAt{8}{79.226}{422.497}{prime}%
\TextAt{8}{112.39}{422.497}{ideal}%
\TextAt{8}{140.353}{422.497}{of}%
\TextAt{16}{153.683}{422.497}{\Character{79}}%
\TextAt{8}{167.43}{422.497}{above}%
\TextAt{14}{200.606}{422.497}{p}%
\TextAt{8}{206.589}{422.497}{,}%
\TextAt{8}{213.738}{422.497}{then}%
\TextAt{8}{240.398}{422.497}{so}%
\TextAt{8}{254.768}{422.497}{is}%
\TextAt{13}{266.532}{422.497}{\Character{27}}%
\TextAt{14}{273.614}{422.497}{P}%
\TextAt{8}{283.508}{422.497}{,}%
\TextAt{8}{290.657}{422.497}{for}%
\TextAt{8}{308.542}{422.497}{each}%
\TextAt{13}{334.879}{422.497}{\Character{27}}%
\NativeMathAt{345.282}{422.497}{11.955}{\in}%
\TextAt{13}{356.573}{422.497}{G}%
\TextAt{8}{365.807}{422.497}{,}%
\TextAt{8}{372.956}{422.497}{because}%
\TextAt{13}{230.573}{395.582}{\Character{27}}%
\TextAt{14}{237.655}{395.582}{P}%
\TextAt{16}{250.205}{395.582}{\Character{92}}%
\TextAt{14}{260.832}{395.582}{o}%
\TextAt{9}{269.997}{395.582}{\Character{61}}%
\TextAt{13}{282.422}{395.582}{\Character{27}}%
\TextAt{9}{289.505}{395.582}{\Character{40}}%
\TextAt{14}{294.057}{395.582}{P}%
\TextAt{16}{306.607}{395.582}{\Character{92}}%
\TextAt{14}{317.234}{395.582}{o}%
\TextAt{9}{323.078}{395.582}{\Character{41}}%
\TextAt{9}{330.933}{395.582}{\Character{61}}%
\TextAt{13}{343.377}{395.582}{\Character{27}}%
\TextAt{14}{350.459}{395.582}{p}%
\TextAt{9}{359.763}{395.582}{\Character{61}}%
\TextAt{14}{372.188}{395.582}{p}%
\TextAt{13}{378.171}{395.582}{\Character{58}}%
\TextAt{8}{79.226}{368.668}{The}%
\TextAt{8}{103.28}{368.668}{ideals}%
\TextAt{13}{135.858}{368.668}{\Character{27}}%
\TextAt{14}{142.94}{368.668}{P}%
\TextAt{8}{152.833}{368.668}{,}%
\TextAt{8}{159.983}{368.668}{for}%
\TextAt{13}{177.868}{368.668}{\Character{27}}%
\NativeMathAt{188.271}{368.668}{11.955}{\in}%
\TextAt{13}{199.562}{368.668}{G}%
\TextAt{8}{208.795}{368.668}{,}%
\TextAt{8}{215.945}{368.668}{are}%
\TextAt{8}{235.455}{368.668}{called}%
\TextAt{8}{268.619}{368.668}{the}%
\TextAt{8}{288.775}{368.668}{prime}%
\TextAt{8}{321.939}{368.668}{ideals}%
\TextAt{12}{354.517}{368.668}{conjugate}%
\TextAt{8}{414.819}{368.668}{to}%
\TextAt{14}{429.129}{368.668}{P}%
\TextAt{8}{439.022}{368.668}{.}%
\TextAt{12}{79.226}{343.251}{Proposition}%
\TextAt{12}{152.545}{343.251}{(9.1).}%
\TextAt{8}{188.388}{343.251}{The}%
\TextAt{8}{212.969}{343.251}{Galois}%
\TextAt{8}{249.41}{343.251}{group}%
\TextAt{13}{283.114}{343.251}{G}%
\TextAt{8}{296.773}{343.251}{acts}%
\TextAt{8}{321.426}{343.251}{transitively}%
\TextAt{8}{383.487}{343.251}{on}%
\TextAt{8}{400.274}{343.251}{the}%
\TextAt{8}{420.958}{343.251}{set}%
\TextAt{8}{439.753}{343.251}{of}%
\TextAt{8}{453.611}{343.251}{all}%
\TextAt{8}{470.398}{343.251}{prime}%
\TextAt{8}{504.089}{343.251}{ideals}%
\TextAt{14}{79.226}{326.205}{P}%
\TextAt{8}{93.017}{326.205}{of}%
\TextAt{16}{106.347}{326.205}{\Character{79}}%
\TextAt{8}{120.094}{326.205}{lying}%
\TextAt{8}{149.037}{326.205}{above}%
\TextAt{14}{182.213}{326.205}{p}%
\TextAt{8}{188.196}{326.205}{,}%
\TextAt{8}{195.345}{326.205}{i.e.,}%
\TextAt{8}{217.45}{326.205}{these}%
\TextAt{8}{247.422}{326.205}{prime}%
\TextAt{8}{280.586}{326.205}{ideals}%
\TextAt{8}{313.163}{326.205}{are}%
\TextAt{8}{332.674}{326.205}{all}%
\TextAt{8}{348.933}{326.205}{conjugates}%
\TextAt{8}{406.557}{326.205}{of}%
\TextAt{8}{419.887}{326.205}{each}%
\TextAt{8}{446.225}{326.205}{other.}%
\TextAt{19}{79.226}{292.419}{Proof.}%
\TextAt{8}{116}{292.419}{Let}%
\TextAt{14}{137.154}{292.419}{P}%
\TextAt{8}{151.129}{292.419}{and}%
\TextAt{14}{174.076}{292.419}{P}%
\TextAt{3}{183.969}{296.757}{\Character{48}}%
\TextAt{8}{190.846}{292.419}{be}%
\TextAt{8}{206.954}{292.419}{two}%
\TextAt{8}{229.256}{292.419}{prime}%
\TextAt{8}{262.604}{292.419}{ideals}%
\TextAt{8}{295.366}{292.419}{above}%
\TextAt{14}{328.726}{292.419}{p}%
\TextAt{8}{334.709}{292.419}{.}%
\TextAt{8}{343.71}{292.419}{Assume}%
\TextAt{14}{387.256}{292.419}{P}%
\TextAt{3}{397.149}{296.757}{\Character{48}}%
\NativeMathAt{403.579}{292.419}{11.955}{\neq}%
\TextAt{13}{416.319}{292.419}{\Character{27}}%
\TextAt{14}{423.401}{292.419}{P}%
\TextAt{8}{437.376}{292.419}{for}%
\TextAt{8}{455.445}{292.419}{any}%
\TextAt{13}{477.747}{292.419}{\Character{27}}%
\NativeMathAt{488.464}{292.419}{11.955}{\in}%
\TextAt{13}{500.069}{292.419}{G}%
\TextAt{8}{509.303}{292.419}{.}%
\TextAt{8}{518.304}{292.419}{By}%
\TextAt{8}{79.226}{275.373}{the}%
\TextAt{8}{99.383}{275.373}{Chinese}%
\TextAt{8}{143.007}{275.373}{remainder}%
\TextAt{8}{198.288}{275.373}{theorem}%
\TextAt{8}{243.813}{275.373}{there}%
\TextAt{8}{273.725}{275.373}{exists}%
\TextAt{13}{306.04}{275.373}{x}%
\NativeMathAt{316.013}{275.373}{11.955}{\in}%
\TextAt{16}{327.287}{275.373}{\Character{79}}%
\TextAt{8}{341.051}{275.373}{such}%
\TextAt{8}{367.448}{275.373}{that}%
\TextAt{13}{154.032}{248.458}{x}%
\NativeMathAt{164.005}{248.458}{11.955}{\equiv}%
\TextAt{9}{176.624}{248.458}{0}%
\TextAt{9}{194.413}{248.458}{\Character{40}}%
\TextAt{8}{198.984}{248.458}{mod}%
\TextAt{14}{225.086}{248.458}{P}%
\TextAt{3}{234.98}{253.395}{\Character{48}}%
\TextAt{9}{237.775}{248.458}{\Character{41}}%
\TextAt{8}{254.282}{248.458}{and}%
\TextAt{13}{285.103}{248.458}{x}%
\NativeMathAt{295.076}{248.458}{11.955}{\equiv}%
\TextAt{9}{307.695}{248.458}{1}%
\TextAt{9}{325.484}{248.458}{\Character{40}}%
\TextAt{8}{330.055}{248.458}{mod}%
\TextAt{13}{356.157}{248.458}{\Character{27}}%
\TextAt{14}{363.24}{248.458}{P}%
\TextAt{9}{373.133}{248.458}{\Character{41}}%
\TextAt{8}{389.64}{248.458}{for}%
\TextAt{8}{407.525}{248.458}{all}%
\TextAt{13}{423.784}{248.458}{\Character{27}}%
\NativeMathAt{434.188}{248.458}{11.955}{\in}%
\TextAt{13}{445.479}{248.458}{G\Character{58}}%
\TextAt{8}{79.226}{221.544}{Then}%
\TextAt{8}{110.064}{221.544}{the}%
\TextAt{8}{130.502}{221.544}{norm}%
\TextAt{13}{161.352}{221.544}{N}%
\TextAt{21}{170.726}{219.684}{L}%
\TextAt{3}{176.485}{219.684}{\Character{106}}%
\TextAt{21}{178.837}{219.684}{K}%
\TextAt{9}{187.04}{221.544}{\Character{40}}%
\TextAt{13}{191.592}{221.544}{x}%
\TextAt{9}{198.244}{221.544}{\Character{41}}%
\TextAt{9}{206.577}{221.544}{\Character{61}}%
\TextAt{20}{219.5}{230.51}{\Character{81}}%
\TextAt{21}{230.791}{218.057}{\Character{27}}%
\NativeMathAt{235.917}{218.057}{7.97}{\in}%
\TextAt{21}{241.562}{218.057}{G}%
\TextAt{13}{250.652}{221.544}{\Character{27}}%
\TextAt{13}{257.717}{221.544}{x}%
\TextAt{8}{268.564}{221.544}{belongs}%
\TextAt{8}{310.856}{221.544}{to}%
\TextAt{14}{325.447}{221.544}{P}%
\TextAt{3}{335.34}{225.882}{\Character{48}}%
\TextAt{16}{340.984}{221.544}{\Character{92}}%
\TextAt{14}{351.802}{221.544}{o}%
\TextAt{9}{361.446}{221.544}{\Character{61}}%
\TextAt{14}{374.35}{221.544}{p}%
\TextAt{8}{380.333}{221.544}{.}%
\TextAt{8}{389.624}{221.544}{On}%
\TextAt{8}{409.416}{221.544}{the}%
\TextAt{8}{429.853}{221.544}{other}%
\TextAt{8}{460.703}{221.544}{hand,}%
\TextAt{13}{493.573}{221.544}{x}%
\NativeMathAt{504.024}{221.544}{11.955}{\notin}%
\TextAt{13}{515.794}{221.544}{\Character{27}}%
\TextAt{14}{522.876}{221.544}{P}%
\TextAt{8}{79.226}{204.498}{for}%
\TextAt{8}{97.94}{204.498}{any}%
\TextAt{13}{120.885}{204.498}{\Character{27}}%
\NativeMathAt{132.701}{204.498}{11.955}{\in}%
\TextAt{13}{145.404}{204.498}{G}%
\TextAt{8}{154.637}{204.498}{,}%
\TextAt{8}{162.822}{204.498}{hence}%
\TextAt{13}{196.157}{204.498}{\Character{27}}%
\TextAt{13}{203.223}{204.498}{x}%
\NativeMathAt{214.625}{204.498}{11.955}{\notin}%
\TextAt{14}{227.328}{204.498}{P}%
\TextAt{8}{241.947}{204.498}{for}%
\TextAt{8}{260.66}{204.498}{any}%
\TextAt{13}{283.606}{204.498}{\Character{27}}%
\NativeMathAt{295.421}{204.498}{11.955}{\in}%
\TextAt{13}{308.124}{204.498}{G}%
\TextAt{8}{317.358}{204.498}{.}%
\TextAt{8}{328.292}{204.498}{Consequently}%
\TextAt{20}{401.701}{213.464}{\Character{81}}%
\TextAt{21}{412.992}{201.011}{\Character{27}}%
\NativeMathAt{418.117}{201.011}{7.97}{\in}%
\TextAt{21}{423.763}{201.011}{G}%
\TextAt{13}{432.852}{204.498}{\Character{27}}%
\TextAt{13}{439.918}{204.498}{x}%
\NativeMathAt{451.319}{204.498}{11.955}{\notin}%
\TextAt{14}{464.022}{204.498}{P}%
\TextAt{16}{477.137}{204.498}{\Character{92}}%
\TextAt{14}{488.329}{204.498}{o}%
\TextAt{9}{498.906}{204.498}{\Character{61}}%
\TextAt{14}{512.743}{204.498}{p}%
\TextAt{8}{518.727}{204.498}{,}%
\TextAt{8}{526.912}{204.498}{a}%
\TextAt{8}{79.226}{187.452}{contradiction.}%
\Vector{q 0 G 0 g 0.399 w 452.334 -586.548 m 452.334 -578.478 l S Q}%
\Vector{q 0 G 0 g 0.399 w 452.534 -578.678 m 459.707 -578.678 l S Q}%
\Vector{q 0 G 0 g 0.399 w 452.534 -586.349 m 459.707 -586.349 l S Q}%
\Vector{q 0 G 0 g 0.399 w 459.906 -586.548 m 459.906 -578.478 l S Q}%
\TextAt{12}{79.226}{153.665}{Definition}%
\TextAt{12}{142.015}{153.665}{(9.2).}%
\TextAt{8}{177.33}{153.665}{If}%
\TextAt{14}{189.022}{153.665}{P}%
\TextAt{8}{202.813}{153.665}{is}%
\TextAt{8}{214.577}{153.665}{a}%
\TextAt{8}{224.332}{153.665}{prime}%
\TextAt{8}{257.496}{153.665}{ideal}%
\TextAt{8}{285.459}{153.665}{of}%
\TextAt{16}{298.789}{153.665}{\Character{79}}%
\TextAt{8}{308.639}{153.665}{,}%
\TextAt{8}{315.788}{153.665}{then}%
\TextAt{8}{342.448}{153.665}{the}%
\TextAt{8}{362.605}{153.665}{subgroup}%
\TextAt{13}{242.231}{126.751}{G}%
\TextAt{27}{251.464}{124.957}{P}%
\TextAt{9}{261.879}{126.751}{\Character{61}}%
\TextAt{16}{274.304}{126.751}{\Character{102}}%
\TextAt{13}{280.282}{126.751}{\Character{27}}%
\NativeMathAt{290.685}{126.751}{11.955}{\in}%
\TextAt{13}{301.976}{126.751}{G}%
\TextAt{16}{314.531}{126.751}{\Character{106}}%
\TextAt{13}{321.172}{126.751}{\Character{27}}%
\TextAt{14}{328.255}{126.751}{P}%
\TextAt{9}{341.469}{126.751}{\Character{61}}%
\TextAt{14}{353.894}{126.751}{P}%
\TextAt{16}{363.787}{126.751}{\Character{103}}%
\TextAt{8}{79.226}{99.836}{is}%
\TextAt{8}{90.99}{99.836}{called}%
\TextAt{8}{124.154}{99.836}{the}%
\TextAt{12}{144.31}{99.836}{decomposition}%
\TextAt{12}{233.125}{99.836}{group}%
\TextAt{8}{270.927}{99.836}{of}%
\TextAt{14}{284.257}{99.836}{P}%
\TextAt{8}{298.048}{99.836}{over}%
\TextAt{13}{323.094}{99.836}{K}%
\TextAt{8}{333.912}{99.836}{.}%
\TextAt{8}{342.36}{99.836}{The}%
\TextAt{8}{366.414}{99.836}{fixed}%
\TextAt{8}{394.7}{99.836}{field}%
\TextAt{13}{208.901}{72.921}{Z}%
\TextAt{27}{216.913}{71.128}{P}%
\TextAt{9}{227.328}{72.921}{\Character{61}}%
\TextAt{16}{239.753}{72.921}{\Character{102}}%
\TextAt{13}{245.731}{72.921}{x}%
\NativeMathAt{255.704}{72.921}{11.955}{\in}%
\TextAt{13}{266.995}{72.921}{L}%
\TextAt{16}{278.28}{72.921}{\Character{106}}%
\TextAt{13}{284.922}{72.921}{\Character{27}}%
\TextAt{13}{291.988}{72.921}{x}%
\TextAt{9}{301.977}{72.921}{\Character{61}}%
\TextAt{13}{314.402}{72.921}{x}%
\TextAt{8}{324.952}{72.921}{for}%
\TextAt{8}{342.837}{72.921}{all}%
\TextAt{13}{359.096}{72.921}{\Character{27}}%
\NativeMathAt{369.499}{72.921}{11.955}{\in}%
\TextAt{13}{380.79}{72.921}{G}%
\TextAt{27}{390.024}{71.128}{P}%
\TextAt{16}{397.117}{72.921}{\Character{103}}%
\TextAt{8}{300.14}{43.033}{50}%
\Destination{78.23}{809.815}{nav.page.50}%
\Destination{79.23}{614.525}{nav.section.1.9}%
\Destination{79.226}{773.355}{nav.ex.1.8.9}%
\Destination{79.226}{708.455}{nav.ex.1.8.10}%
\Destination{79.226}{676.455}{nav.ex.1.8.11}%
\Destination{79.36}{358.255}{nav.print.54}%
\Destination{79.36}{357.255}{nav.number.1.9.1}%
\Destination{79.36}{167.655}{nav.number.1.9.2}%
\Destination{79.23}{614.525}{outline.14}%
\Destination{79.226}{727.8}{errata.56}%
}

% Page 63 (body page 51)
\Page{%
\TextAt{2}{79.226}{796.943}{I.}%
\TextAt{2}{89.83}{796.943}{Algebraic}%
\TextAt{2}{138.648}{796.943}{Integers}%
\Vector{q 0 G 0 g 0.299 w 7.226 18.729 m 460.769 18.729 l S Q}%
\TextAt{8}{79.226}{761.124}{is}%
\TextAt{8}{90.99}{761.124}{called}%
\TextAt{8}{124.154}{761.124}{the}%
\TextAt{12}{144.31}{761.124}{decomposition}%
\TextAt{12}{233.125}{761.124}{field}%
\TextAt{8}{261.841}{761.124}{of}%
\TextAt{14}{275.171}{761.124}{P}%
\TextAt{8}{288.962}{761.124}{over}%
\TextAt{13}{314.008}{761.124}{K}%
\TextAt{8}{324.826}{761.124}{.}%
\TextAt{8}{97.159}{735.591}{The}%
\TextAt{8}{120.732}{735.591}{decomposition}%
\TextAt{8}{197.386}{735.591}{group}%
\TextAt{8}{230.081}{735.591}{encodes}%
\TextAt{8}{272.902}{735.591}{in}%
\TextAt{8}{286.075}{735.591}{group-theoretic}%
\TextAt{8}{367.547}{735.591}{language}%
\TextAt{8}{415.855}{735.591}{the}%
\TextAt{8}{435.531}{735.591}{number}%
\TextAt{8}{477.647}{735.591}{of}%
\TextAt{8}{490.496}{735.591}{different}%
\TextAt{8}{79.226}{718.545}{prime}%
\TextAt{8}{113.019}{718.545}{ideals}%
\TextAt{8}{146.225}{718.545}{into}%
\TextAt{8}{170.597}{718.545}{which}%
\TextAt{8}{204.712}{718.545}{a}%
\TextAt{8}{215.097}{718.545}{prime}%
\TextAt{8}{248.889}{718.545}{ideal}%
\TextAt{14}{277.481}{718.545}{p}%
\TextAt{8}{287.99}{718.545}{of}%
\TextAt{14}{301.949}{718.545}{o}%
\TextAt{8}{312.32}{718.545}{decomposes}%
\TextAt{8}{376.478}{718.545}{in}%
\TextAt{16}{390.76}{718.545}{\Character{79}}%
\TextAt{8}{400.61}{718.545}{.}%
\TextAt{8}{410.945}{718.545}{For}%
\TextAt{8}{432.543}{718.545}{if}%
\TextAt{14}{443.896}{718.545}{P}%
\TextAt{8}{458.315}{718.545}{is}%
\TextAt{8}{470.708}{718.545}{one}%
\TextAt{8}{492.796}{718.545}{of}%
\TextAt{8}{506.755}{718.545}{them}%
\TextAt{8}{79.226}{701.499}{and}%
\TextAt{13}{101.919}{701.499}{\Character{27}}%
\TextAt{8}{112.829}{701.499}{varies}%
\TextAt{8}{145.672}{701.499}{over}%
\TextAt{8}{170.649}{701.499}{a}%
\TextAt{8}{180.335}{701.499}{system}%
\TextAt{8}{219.083}{701.499}{of}%
\TextAt{8}{232.344}{701.499}{representatives}%
\TextAt{8}{312.075}{701.499}{of}%
\TextAt{8}{325.335}{701.499}{the}%
\TextAt{8}{345.422}{701.499}{cosets}%
\TextAt{8}{379.293}{701.499}{in}%
\TextAt{13}{392.877}{701.499}{G}%
\TextAt{9}{402.11}{701.499}{\Character{47}}%
\TextAt{13}{407.963}{701.499}{G}%
\TextAt{27}{417.197}{699.706}{P}%
\TextAt{8}{424.29}{701.499}{,}%
\TextAt{8}{431.384}{701.499}{then}%
\TextAt{13}{457.974}{701.499}{\Character{27}}%
\TextAt{14}{465.057}{701.499}{P}%
\TextAt{8}{478.778}{701.499}{varies}%
\TextAt{8}{511.621}{701.499}{over}%
\TextAt{8}{79.226}{684.453}{the}%
\TextAt{8}{100.046}{684.453}{different}%
\TextAt{8}{146.881}{684.453}{prime}%
\TextAt{8}{180.709}{684.453}{ideals}%
\TextAt{8}{213.95}{684.453}{above}%
\TextAt{14}{247.79}{684.453}{p}%
\TextAt{8}{253.773}{684.453}{,}%
\TextAt{8}{261.752}{684.453}{each}%
\TextAt{8}{288.753}{684.453}{one}%
\TextAt{8}{310.877}{684.453}{occurring}%
\TextAt{8}{363.247}{684.453}{precisely}%
\TextAt{8}{411.767}{684.453}{once,}%
\TextAt{8}{442.509}{684.453}{i.e.,}%
\TextAt{8}{465.444}{684.453}{their}%
\TextAt{8}{494.071}{684.453}{number}%
\TextAt{8}{79.226}{667.407}{equals}%
\TextAt{8}{114.733}{667.407}{the}%
\TextAt{8}{134.889}{667.407}{index}%
\TextAt{9}{166.427}{667.407}{\Character{40}}%
\TextAt{13}{170.979}{667.407}{G}%
\TextAt{9}{183.534}{667.407}{\Character{58}}%
\TextAt{13}{190.106}{667.407}{G}%
\TextAt{27}{199.34}{665.614}{P}%
\TextAt{9}{206.434}{667.407}{\Character{41}}%
\TextAt{8}{210.986}{667.407}{.}%
\TextAt{8}{219.434}{667.407}{In}%
\TextAt{8}{234.055}{667.407}{particular,}%
\TextAt{8}{291.297}{667.407}{one}%
\TextAt{8}{312.756}{667.407}{has}%
\TextAt{13}{198.517}{640.086}{G}%
\TextAt{27}{207.75}{638.293}{P}%
\TextAt{9}{218.165}{640.086}{\Character{61}}%
\TextAt{9}{230.575}{640.086}{1}%
\NativeArrowAt{239.764}{640.086}{11.955}{21.926}{\Longleftrightarrow}%
\TextAt{13}{265.003}{640.086}{Z}%
\TextAt{27}{273.015}{638.293}{P}%
\TextAt{9}{283.429}{640.086}{\Character{61}}%
\TextAt{13}{295.855}{640.086}{L}%
\NativeArrowAt{307.14}{640.086}{11.955}{21.926}{\Longleftrightarrow}%
\TextAt{14}{332.379}{640.086}{p}%
\TextAt{8}{342.259}{640.086}{is}%
\TextAt{8}{354.023}{640.086}{totally}%
\TextAt{8}{391.431}{640.086}{split}%
\TextAt{13}{413.608}{640.086}{\Character{59}}%
\TextAt{13}{195.136}{620.051}{G}%
\TextAt{27}{204.37}{618.258}{P}%
\TextAt{9}{214.784}{620.051}{\Character{61}}%
\TextAt{13}{227.21}{620.051}{G}%
\NativeArrowAt{239.764}{620.051}{11.955}{21.926}{\Longleftrightarrow}%
\TextAt{13}{265.003}{620.051}{Z}%
\TextAt{27}{273.015}{618.258}{P}%
\TextAt{9}{283.429}{620.051}{\Character{61}}%
\TextAt{13}{295.855}{620.051}{K}%
\NativeArrowAt{309.993}{620.051}{11.955}{21.926}{\Longleftrightarrow}%
\TextAt{14}{335.232}{620.051}{p}%
\TextAt{8}{345.113}{620.051}{is}%
\TextAt{8}{356.877}{620.051}{nonsplit}%
\TextAt{13}{397.919}{620.051}{\Character{58}}%
\TextAt{8}{79.226}{592.73}{The}%
\TextAt{8}{103.28}{592.73}{decomposition}%
\TextAt{8}{180.415}{592.73}{group}%
\TextAt{8}{213.59}{592.73}{of}%
\TextAt{8}{226.92}{592.73}{a}%
\TextAt{8}{236.676}{592.73}{prime}%
\TextAt{8}{269.839}{592.73}{ideal}%
\TextAt{13}{297.803}{592.73}{\Character{27}}%
\TextAt{14}{304.885}{592.73}{P}%
\TextAt{8}{318.676}{592.73}{conjugate}%
\TextAt{8}{371.685}{592.73}{to}%
\TextAt{14}{385.995}{592.73}{P}%
\TextAt{8}{399.786}{592.73}{is}%
\TextAt{8}{411.55}{592.73}{the}%
\TextAt{8}{431.706}{592.73}{conjugate}%
\TextAt{8}{484.715}{592.73}{subgroup}%
\TextAt{13}{264.867}{565.409}{G}%
\TextAt{21}{274.101}{563.615}{\Character{27}}%
\TextAt{27}{279.226}{563.615}{P}%
\TextAt{9}{289.64}{565.409}{\Character{61}}%
\TextAt{13}{302.066}{565.409}{\Character{27}}%
\TextAt{13}{309.132}{565.409}{G}%
\TextAt{27}{318.382}{563.615}{P}%
\TextAt{13}{325.476}{565.409}{\Character{27}}%
\TextAt{3}{332.558}{570.345}{\Character{0}}%
\TextAt{17}{339.145}{570.345}{1}%
\TextAt{13}{343.877}{565.409}{\Character{58}}%
\TextAt{8}{79.226}{538.087}{In}%
\TextAt{8}{93.847}{538.087}{fact,}%
\TextAt{8}{120.185}{538.087}{for}%
\TextAt{13}{138.069}{538.087}{\Character{28}}%
\NativeMathAt{147.809}{538.087}{11.955}{\in}%
\TextAt{13}{159.1}{538.087}{G}%
\TextAt{8}{168.334}{538.087}{,}%
\TextAt{8}{175.483}{538.087}{one}%
\TextAt{8}{196.943}{538.087}{has}%
\TextAt{8}{217.816}{538.087}{the}%
\TextAt{8}{237.973}{538.087}{equivalences}%
\TextAt{13}{191.729}{510.766}{\Character{28}}%
\NativeMathAt{201.469}{510.766}{11.955}{\in}%
\TextAt{13}{212.76}{510.766}{G}%
\TextAt{21}{221.994}{508.973}{\Character{27}}%
\TextAt{27}{227.119}{508.973}{P}%
\NativeArrowAt{237.533}{510.766}{11.955}{21.926}{\Longleftrightarrow}%
\TextAt{13}{262.772}{510.766}{\Character{28}}%
\TextAt{13}{269.18}{510.766}{\Character{27}}%
\TextAt{14}{276.273}{510.766}{P}%
\TextAt{9}{289.487}{510.766}{\Character{61}}%
\TextAt{13}{301.913}{510.766}{\Character{27}}%
\TextAt{14}{308.995}{510.766}{P}%
\NativeArrowAt{322.209}{510.766}{11.955}{21.926}{\Longleftrightarrow}%
\TextAt{13}{347.448}{510.766}{\Character{27}}%
\TextAt{3}{354.53}{515.702}{\Character{0}}%
\TextAt{17}{361.117}{515.702}{1}%
\TextAt{13}{365.849}{510.766}{\Character{28}}%
\TextAt{13}{372.257}{510.766}{\Character{27}}%
\TextAt{14}{379.351}{510.766}{P}%
\TextAt{9}{392.565}{510.766}{\Character{61}}%
\TextAt{14}{404.99}{510.766}{P}%
\NativeArrowAt{237.533}{490.731}{11.955}{21.926}{\Longleftrightarrow}%
\TextAt{13}{262.772}{490.731}{\Character{27}}%
\TextAt{3}{269.854}{495.667}{\Character{0}}%
\TextAt{17}{276.441}{495.667}{1}%
\TextAt{13}{281.173}{490.731}{\Character{28}}%
\TextAt{13}{287.581}{490.731}{\Character{27}}%
\NativeMathAt{297.995}{490.731}{11.955}{\in}%
\TextAt{13}{309.286}{490.731}{G}%
\TextAt{27}{318.52}{488.938}{P}%
\NativeArrowAt{328.934}{490.731}{11.955}{21.926}{\Longleftrightarrow}%
\TextAt{13}{354.173}{490.731}{\Character{28}}%
\NativeMathAt{363.913}{490.731}{11.955}{\in}%
\TextAt{13}{375.204}{490.731}{\Character{27}}%
\TextAt{13}{382.269}{490.731}{G}%
\TextAt{27}{391.52}{488.938}{P}%
\TextAt{13}{398.614}{490.731}{\Character{27}}%
\TextAt{3}{405.696}{495.667}{\Character{0}}%
\TextAt{17}{412.282}{495.667}{1}%
\TextAt{13}{417.015}{490.731}{\Character{58}}%
\TextAt{12}{79.226}{463.41}{Remark:}%
\TextAt{8}{135.617}{463.41}{The}%
\TextAt{8}{160.139}{463.41}{decomposition}%
\TextAt{8}{237.741}{463.41}{group}%
\TextAt{8}{271.384}{463.41}{regulates}%
\TextAt{8}{321.346}{463.41}{the}%
\TextAt{8}{341.97}{463.41}{prime}%
\TextAt{8}{375.601}{463.41}{decomposition}%
\TextAt{8}{453.204}{463.41}{also}%
\TextAt{8}{477.151}{463.41}{in}%
\TextAt{8}{491.272}{463.41}{the}%
\TextAt{8}{511.896}{463.41}{case}%
\TextAt{8}{79.226}{446.364}{of}%
\TextAt{8}{92.105}{446.364}{a}%
\TextAt{8}{101.41}{446.364}{non-Galois}%
\TextAt{8}{159.635}{446.364}{extension.}%
\TextAt{8}{215.802}{446.364}{For}%
\TextAt{8}{236.32}{446.364}{subgroups}%
\TextAt{13}{291.282}{446.364}{U}%
\TextAt{8}{303.928}{446.364}{and}%
\TextAt{13}{326.239}{446.364}{V}%
\TextAt{8}{339.112}{446.364}{of}%
\TextAt{8}{351.991}{446.364}{a}%
\TextAt{8}{361.296}{446.364}{group}%
\TextAt{13}{394.02}{446.364}{G}%
\TextAt{8}{403.254}{446.364}{,}%
\TextAt{8}{410.043}{446.364}{consider}%
\TextAt{8}{455.177}{446.364}{the}%
\TextAt{8}{474.883}{446.364}{equivalence}%
\TextAt{8}{79.226}{429.318}{relation}%
\TextAt{8}{122.157}{429.318}{in}%
\TextAt{13}{135.81}{429.318}{G}%
\TextAt{8}{148.941}{429.318}{defined}%
\TextAt{8}{189.254}{429.318}{by}%
\TextAt{13}{190.794}{401.997}{\Character{27}}%
\NativeMathAt{201.197}{401.997}{11.955}{\sim}%
\TextAt{13}{213.817}{401.997}{\Character{27}}%
\TextAt{3}{220.899}{406.933}{\Character{48}}%
\NativeArrowAt{238.97}{401.997}{11.955}{21.926}{\Longleftrightarrow}%
\TextAt{13}{276.164}{401.997}{\Character{27}}%
\TextAt{3}{283.246}{406.933}{\Character{48}}%
\TextAt{9}{289.362}{401.997}{\Character{61}}%
\TextAt{13}{301.788}{401.997}{u\Character{27}}%
\TextAt{13}{315.513}{401.997}{v}%
\TextAt{8}{333.576}{401.997}{for}%
\TextAt{13}{351.461}{401.997}{u}%
\NativeMathAt{361.444}{401.997}{11.955}{\in}%
\TextAt{13}{372.735}{401.997}{U\Character{59}}%
\TextAt{13}{389.759}{401.997}{v}%
\NativeMathAt{399.184}{401.997}{11.955}{\in}%
\TextAt{13}{410.475}{401.997}{V}%
\TextAt{13}{417.935}{401.997}{\Character{58}}%
\TextAt{8}{79.226}{374.675}{The}%
\TextAt{8}{103.28}{374.675}{corresponding}%
\TextAt{8}{178.466}{374.675}{equivalence}%
\TextAt{8}{240.25}{374.675}{classes}%
\TextAt{13}{230.176}{347.354}{U}%
\TextAt{13}{239.37}{347.354}{\Character{27}}%
\TextAt{13}{246.447}{347.354}{V}%
\TextAt{9}{259.204}{347.354}{\Character{61}}%
\TextAt{16}{271.629}{347.354}{\Character{102}}%
\TextAt{13}{277.607}{347.354}{u\Character{27}}%
\TextAt{13}{291.332}{347.354}{v}%
\TextAt{16}{300.761}{347.354}{\Character{106}}%
\TextAt{13}{307.402}{347.354}{u}%
\NativeMathAt{317.386}{347.354}{11.955}{\in}%
\TextAt{13}{328.677}{347.354}{U\Character{59}}%
\TextAt{13}{345.701}{347.354}{v}%
\NativeMathAt{355.126}{347.354}{11.955}{\in}%
\TextAt{13}{366.417}{347.354}{V}%
\TextAt{16}{375.842}{347.354}{\Character{103}}%
\TextAt{8}{79.226}{320.033}{are}%
\TextAt{8}{98.118}{320.033}{called}%
\TextAt{8}{130.662}{320.033}{the}%
\TextAt{12}{150.199}{320.033}{double}%
\TextAt{12}{192.992}{320.033}{cosets}%
\TextAt{8}{230.963}{320.033}{of}%
\TextAt{13}{243.674}{320.033}{G}%
\TextAt{8}{256.186}{320.033}{modd}%
\TextAt{13}{288.407}{320.033}{U\Character{59}}%
\TextAt{13}{301.534}{320.033}{V}%
\TextAt{8}{310.976}{320.033}{.}%
\TextAt{8}{319.218}{320.033}{The}%
\TextAt{8}{342.652}{320.033}{set}%
\TextAt{8}{360.3}{320.033}{of}%
\TextAt{8}{373.011}{320.033}{these}%
\TextAt{8}{402.363}{320.033}{double}%
\TextAt{8}{439.462}{320.033}{cosets,}%
\TextAt{8}{476.159}{320.033}{which}%
\TextAt{8}{509.027}{320.033}{form}%
\TextAt{8}{79.226}{302.987}{a}%
\TextAt{8}{88.982}{302.987}{partition}%
\TextAt{8}{137.771}{302.987}{of}%
\TextAt{13}{151.101}{302.987}{G}%
\TextAt{8}{160.334}{302.987}{,}%
\TextAt{8}{167.483}{302.987}{is}%
\TextAt{8}{179.247}{302.987}{denoted}%
\TextAt{13}{223.47}{302.987}{U}%
\TextAt{16}{232.669}{302.987}{\Character{110}}%
\TextAt{13}{238.646}{302.987}{G}%
\TextAt{9}{247.88}{302.987}{\Character{47}}%
\TextAt{13}{253.733}{302.987}{V}%
\TextAt{8}{263.158}{302.987}{.}%
\TextAt{8}{97.159}{282.952}{Now}%
\TextAt{8}{122.528}{282.952}{let}%
\TextAt{13}{138.141}{282.952}{L}%
\TextAt{16}{149.427}{282.952}{\Character{106}}%
\TextAt{13}{156.068}{282.952}{K}%
\TextAt{8}{169.492}{282.952}{be}%
\TextAt{8}{184.125}{282.952}{an}%
\TextAt{8}{199.093}{282.952}{arbitrary}%
\TextAt{8}{247.572}{282.952}{separable}%
\TextAt{8}{297.724}{282.952}{extension,}%
\TextAt{8}{351.708}{282.952}{and}%
\TextAt{8}{373.18}{282.952}{embed}%
\TextAt{8}{408.95}{282.952}{it}%
\TextAt{8}{419.363}{282.952}{into}%
\TextAt{8}{441.814}{282.952}{a}%
\TextAt{8}{450.279}{282.952}{Galois}%
\TextAt{8}{484.901}{282.952}{extension}%
\TextAt{13}{79.226}{265.906}{N}%
\TextAt{16}{94.542}{265.906}{\Character{106}}%
\TextAt{13}{102.556}{265.906}{K}%
\TextAt{8}{118.077}{265.906}{with}%
\TextAt{8}{145.542}{265.906}{Galois}%
\TextAt{8}{182.261}{265.906}{group}%
\TextAt{13}{216.241}{265.906}{G}%
\TextAt{8}{225.475}{265.906}{.}%
\TextAt{8}{236.339}{265.906}{In}%
\TextAt{13}{251.766}{265.906}{G}%
\TextAt{8}{261}{265.906}{,}%
\TextAt{8}{269.155}{265.906}{consider}%
\TextAt{8}{315.546}{265.906}{the}%
\TextAt{8}{336.508}{265.906}{subgroup}%
\TextAt{13}{388.11}{265.906}{H}%
\TextAt{9}{403.426}{265.906}{\Character{61}}%
\TextAt{13}{417.224}{265.906}{G}%
\TextAt{9}{426.458}{265.906}{\Character{40}}%
\TextAt{13}{431.01}{265.906}{N}%
\TextAt{16}{446.326}{265.906}{\Character{106}}%
\TextAt{13}{454.34}{265.906}{L}%
\TextAt{9}{462.305}{265.906}{\Character{41}}%
\TextAt{8}{466.857}{265.906}{.}%
\TextAt{8}{477.721}{265.906}{Let}%
\TextAt{14}{499.496}{265.906}{p}%
\TextAt{8}{510.182}{265.906}{be}%
\TextAt{8}{526.911}{265.906}{a}%
\TextAt{8}{79.226}{248.86}{prime}%
\TextAt{8}{112.527}{248.86}{ideal}%
\TextAt{8}{140.628}{248.86}{of}%
\TextAt{13}{154.096}{248.86}{K}%
\TextAt{8}{168.949}{248.86}{and}%
\TextAt{13}{191.849}{248.86}{P}%
\TextAt{27}{199.394}{247.067}{p}%
\TextAt{8}{207.916}{248.86}{the}%
\TextAt{8}{228.21}{248.86}{set}%
\TextAt{8}{246.615}{248.86}{of}%
\TextAt{8}{260.083}{248.86}{prime}%
\TextAt{8}{293.384}{248.86}{ideals}%
\TextAt{8}{326.1}{248.86}{of}%
\TextAt{13}{339.567}{248.86}{L}%
\TextAt{8}{351.567}{248.86}{above}%
\TextAt{14}{384.88}{248.86}{p}%
\TextAt{8}{390.863}{248.86}{.}%
\TextAt{8}{399.724}{248.86}{If}%
\TextAt{14}{411.554}{248.86}{P}%
\TextAt{8}{425.483}{248.86}{is}%
\TextAt{8}{437.384}{248.86}{a}%
\TextAt{8}{447.277}{248.86}{prime}%
\TextAt{8}{480.578}{248.86}{ideal}%
\TextAt{8}{508.679}{248.86}{of}%
\TextAt{13}{522.147}{248.86}{N}%
\TextAt{8}{79.226}{231.814}{above}%
\TextAt{14}{112.402}{231.814}{p}%
\TextAt{8}{118.385}{231.814}{,}%
\TextAt{8}{125.534}{231.814}{then}%
\TextAt{8}{152.194}{231.814}{the}%
\TextAt{8}{172.35}{231.814}{rule}%
\TextAt{13}{197.741}{204.493}{H}%
\TextAt{16}{208.363}{204.493}{\Character{110}}%
\TextAt{13}{214.341}{204.493}{G}%
\TextAt{9}{223.575}{204.493}{\Character{47}}%
\TextAt{13}{229.428}{204.493}{G}%
\TextAt{27}{238.661}{202.699}{P}%
\NativeArrowAt{249.076}{204.493}{11.955}{19.26}{\longrightarrow}%
\TextAt{13}{271.658}{204.493}{P}%
\TextAt{27}{279.203}{202.699}{p}%
\TextAt{13}{283.69}{204.493}{\Character{59}}%
\TextAt{13}{312.825}{204.493}{H}%
\TextAt{13}{323.441}{204.493}{\Character{27}}%
\TextAt{13}{330.519}{204.493}{G}%
\TextAt{27}{339.783}{202.699}{P}%
\NativeArrowAt{350.197}{204.493}{11.955}{19.26}{\longmapsto}%
\TextAt{13}{372.78}{204.493}{\Character{27}}%
\TextAt{14}{379.862}{204.493}{P}%
\TextAt{16}{392.412}{204.493}{\Character{92}}%
\TextAt{13}{403.039}{204.493}{L\Character{59}}%
\TextAt{8}{79.226}{177.171}{gives}%
\TextAt{8}{107.907}{177.171}{a}%
\TextAt{8}{117.662}{177.171}{well-defined}%
\TextAt{8}{181.706}{177.171}{bijection.}%
\TextAt{8}{234.054}{177.171}{The}%
\TextAt{8}{258.107}{177.171}{proof}%
\TextAt{8}{288.677}{177.171}{is}%
\TextAt{8}{300.441}{177.171}{left}%
\TextAt{8}{320.92}{177.171}{to}%
\TextAt{8}{335.23}{177.171}{the}%
\TextAt{8}{355.387}{177.171}{reader.}%
\TextAt{8}{97.159}{157.136}{In}%
\TextAt{8}{111.351}{157.136}{the}%
\TextAt{8}{131.079}{157.136}{Galois}%
\TextAt{8}{166.564}{157.136}{case,}%
\TextAt{8}{194.244}{157.136}{the}%
\TextAt{8}{213.972}{157.136}{inertia}%
\TextAt{8}{250.616}{157.136}{degrees}%
\TextAt{13}{291.218}{157.136}{f}%
\TextAt{17}{296.988}{155.343}{1}%
\TextAt{13}{301.72}{157.136}{\Character{59}}%
\TextAt{13}{306.944}{157.136}{\Character{58}}%
\TextAt{13}{312.181}{157.136}{\Character{58}}%
\TextAt{13}{317.405}{157.136}{\Character{58}}%
\TextAt{13}{322.642}{157.136}{\Character{59}}%
\TextAt{13}{327.878}{157.136}{f}%
\TextAt{21}{333.711}{155.343}{r}%
\TextAt{8}{341.734}{157.136}{and}%
\TextAt{8}{364.068}{157.136}{the}%
\TextAt{8}{383.796}{157.136}{ramification}%
\TextAt{8}{448.415}{157.136}{indices}%
\TextAt{13}{486.411}{157.136}{e}%
\TextAt{17}{491.836}{155.343}{1}%
\TextAt{13}{496.569}{157.136}{\Character{59}}%
\TextAt{13}{501.793}{157.136}{\Character{58}}%
\TextAt{13}{507.03}{157.136}{\Character{58}}%
\TextAt{13}{512.254}{157.136}{\Character{58}}%
\TextAt{13}{517.491}{157.136}{\Character{59}}%
\TextAt{13}{522.727}{157.136}{e}%
\TextAt{21}{528.215}{155.343}{r}%
\TextAt{8}{79.226}{140.09}{in}%
\TextAt{8}{92.879}{140.09}{the}%
\TextAt{8}{113.035}{140.09}{prime}%
\TextAt{8}{146.199}{140.09}{decomposition}%
\TextAt{14}{267.693}{123.044}{p}%
\TextAt{9}{276.996}{123.044}{\Character{61}}%
\TextAt{14}{289.422}{123.044}{P}%
\TextAt{21}{299.315}{128.269}{e}%
\TextAt{24}{303.244}{127.162}{1}%
\TextAt{17}{299.315}{120.113}{1}%
\TextAt{16}{309.885}{123.044}{\Character{1}}%
\TextAt{16}{315.181}{123.044}{\Character{1}}%
\TextAt{16}{320.489}{123.044}{\Character{1}}%
\TextAt{14}{325.825}{123.044}{P}%
\TextAt{21}{335.719}{127.981}{e}%
\TextAt{22}{339.647}{126.984}{r}%
\TextAt{21}{335.719}{120.089}{r}%
\TextAt{8}{79.226}{100.243}{of}%
\TextAt{8}{92.556}{100.243}{a}%
\TextAt{8}{102.312}{100.243}{prime}%
\TextAt{8}{135.475}{100.243}{ideal}%
\TextAt{14}{163.438}{100.243}{p}%
\TextAt{8}{173.319}{100.243}{of}%
\TextAt{13}{186.649}{100.243}{K}%
\TextAt{8}{201.364}{100.243}{are}%
\TextAt{8}{220.875}{100.243}{both}%
\TextAt{8}{248.515}{100.243}{independent}%
\TextAt{8}{314.842}{100.243}{of}%
\TextAt{13}{328.172}{100.243}{i}%
\TextAt{8}{332.166}{100.243}{,}%
\TextAt{13}{202.215}{72.921}{f}%
\TextAt{17}{207.985}{71.128}{1}%
\TextAt{9}{216.038}{72.921}{\Character{61}}%
\TextAt{16}{228.463}{72.921}{\Character{1}}%
\TextAt{16}{233.759}{72.921}{\Character{1}}%
\TextAt{16}{239.067}{72.921}{\Character{1}}%
\TextAt{9}{245.732}{72.921}{\Character{61}}%
\TextAt{13}{258.157}{72.921}{f}%
\TextAt{21}{263.927}{71.128}{r}%
\TextAt{9}{271.803}{72.921}{\Character{61}}%
\TextAt{13}{284.229}{72.921}{f}%
\TextAt{13}{290.613}{72.921}{\Character{59}}%
\TextAt{13}{319.748}{72.921}{e}%
\TextAt{17}{325.205}{71.128}{1}%
\TextAt{9}{333.258}{72.921}{\Character{61}}%
\TextAt{16}{345.683}{72.921}{\Character{1}}%
\TextAt{16}{350.979}{72.921}{\Character{1}}%
\TextAt{16}{356.287}{72.921}{\Character{1}}%
\TextAt{9}{362.952}{72.921}{\Character{61}}%
\TextAt{13}{375.377}{72.921}{e}%
\TextAt{21}{380.803}{71.128}{r}%
\TextAt{9}{388.678}{72.921}{\Character{61}}%
\TextAt{13}{401.104}{72.921}{e\Character{58}}%
\TextAt{8}{300.14}{43.033}{51}%
\Destination{78.23}{809.815}{nav.page.51}%
\Destination{79.36}{389.655}{nav.print.55}%
}

% Page 64 (body page 52)
\Page{%
\TextAt{2}{79.226}{796.943}{I.}%
\TextAt{2}{89.83}{796.943}{Algebraic}%
\TextAt{2}{138.648}{796.943}{Integers}%
\Vector{q 0 G 0 g 0.299 w 7.226 18.729 m 460.769 18.729 l S Q}%
\TextAt{8}{79.226}{761.124}{In}%
\TextAt{8}{94.778}{761.124}{fact,}%
\TextAt{8}{122.278}{761.124}{writing}%
\TextAt{14}{163.533}{761.124}{P}%
\TextAt{9}{178.332}{761.124}{\Character{61}}%
\TextAt{14}{192.343}{761.124}{P}%
\TextAt{17}{202.236}{759.33}{1}%
\TextAt{8}{206.968}{761.124}{,}%
\TextAt{8}{215.28}{761.124}{we}%
\TextAt{8}{233.438}{761.124}{find}%
\TextAt{14}{257.776}{761.124}{P}%
\TextAt{21}{267.67}{759.33}{i}%
\TextAt{9}{275.957}{761.124}{\Character{61}}%
\TextAt{13}{289.968}{761.124}{\Character{27}}%
\TextAt{21}{296.621}{759.33}{i}%
\TextAt{14}{300.003}{761.124}{P}%
\TextAt{8}{314.723}{761.124}{for}%
\TextAt{8}{333.539}{761.124}{suitable}%
\TextAt{13}{378.105}{761.124}{\Character{27}}%
\TextAt{21}{384.759}{759.33}{i}%
\NativeMathAt{393.046}{761.124}{11.955}{\in}%
\TextAt{13}{405.922}{761.124}{G}%
\TextAt{8}{415.156}{761.124}{,}%
\TextAt{8}{423.468}{761.124}{and}%
\TextAt{8}{447.161}{761.124}{the}%
\TextAt{8}{468.247}{761.124}{isomorphism}%
\TextAt{13}{79.226}{744.078}{\Character{27}}%
\TextAt{21}{85.88}{742.284}{i}%
\TextAt{9}{92.582}{744.078}{\Character{58}}%
\TextAt{16}{99.154}{744.078}{\Character{79}}%
\NativeMathAt{118.662}{750.855}{7.97}{\sim}%
\NativeArrowAt{112.325}{744.078}{11.955}{19.26}{\longrightarrow}%
\TextAt{16}{134.896}{744.078}{\Character{79}}%
\TextAt{8}{148.654}{744.078}{induces}%
\TextAt{8}{190.33}{744.078}{an}%
\TextAt{8}{206.589}{744.078}{isomorphism}%
\TextAt{16}{178.447}{715.352}{\Character{79}}%
\TextAt{9}{188.297}{715.352}{\Character{47}}%
\TextAt{14}{194.15}{715.352}{P}%
\NativeMathAt{213.701}{722.129}{7.97}{\sim}%
\NativeArrowAt{207.364}{715.352}{11.955}{19.26}{\longrightarrow}%
\TextAt{16}{229.935}{715.352}{\Character{79}}%
\TextAt{9}{239.796}{715.352}{\Character{47}}%
\TextAt{13}{245.649}{715.352}{\Character{27}}%
\TextAt{21}{252.302}{713.558}{i}%
\TextAt{14}{255.683}{715.352}{P}%
\TextAt{13}{265.577}{715.352}{\Character{59}}%
\TextAt{13}{294.712}{715.352}{a}%
\TextAt{8}{304.197}{715.352}{mod}%
\TextAt{14}{329.635}{715.352}{P}%
\NativeArrowAt{342.849}{715.352}{11.955}{19.26}{\longmapsto}%
\TextAt{13}{365.431}{715.352}{\Character{27}}%
\TextAt{21}{372.084}{713.558}{i}%
\TextAt{13}{375.465}{715.352}{a}%
\TextAt{8}{384.931}{715.352}{mod}%
\TextAt{13}{410.369}{715.352}{\Character{27}}%
\TextAt{21}{417.023}{713.558}{i}%
\TextAt{14}{420.404}{715.352}{P}%
\TextAt{13}{430.297}{715.352}{\Character{59}}%
\TextAt{8}{79.226}{686.625}{so}%
\TextAt{8}{93.596}{686.625}{that}%
\TextAt{13}{176.401}{669.579}{f}%
\TextAt{21}{182.171}{667.786}{i}%
\TextAt{9}{188.873}{669.579}{\Character{61}}%
\TextAt{9}{201.283}{669.579}{\Character{91}}%
\TextAt{16}{204.55}{669.579}{\Character{79}}%
\TextAt{9}{214.4}{669.579}{\Character{47}}%
\TextAt{13}{220.253}{669.579}{\Character{27}}%
\TextAt{21}{226.906}{667.786}{i}%
\TextAt{14}{230.287}{669.579}{P}%
\TextAt{9}{243.501}{669.579}{\Character{58}}%
\TextAt{14}{250.074}{669.579}{o}%
\TextAt{9}{255.918}{669.579}{\Character{47}}%
\TextAt{14}{261.771}{669.579}{p}%
\TextAt{9}{267.754}{669.579}{\Character{93}}%
\TextAt{9}{274.305}{669.579}{\Character{61}}%
\TextAt{9}{286.727}{669.579}{\Character{91}}%
\TextAt{16}{290.004}{669.579}{\Character{79}}%
\TextAt{9}{299.854}{669.579}{\Character{47}}%
\TextAt{14}{305.707}{669.579}{P}%
\TextAt{9}{318.921}{669.579}{\Character{58}}%
\TextAt{14}{325.493}{669.579}{o}%
\TextAt{9}{331.337}{669.579}{\Character{47}}%
\TextAt{14}{337.19}{669.579}{p}%
\TextAt{9}{343.174}{669.579}{\Character{93}}%
\TextAt{13}{346.425}{669.579}{\Character{59}}%
\TextAt{13}{375.56}{669.579}{i}%
\TextAt{9}{382.894}{669.579}{\Character{61}}%
\TextAt{9}{395.303}{669.579}{1}%
\TextAt{13}{401.172}{669.579}{\Character{59}}%
\TextAt{13}{406.396}{669.579}{\Character{58}}%
\TextAt{13}{411.633}{669.579}{\Character{58}}%
\TextAt{13}{416.869}{669.579}{\Character{58}}%
\TextAt{13}{422.094}{669.579}{\Character{59}}%
\TextAt{13}{427.33}{669.579}{r}%
\TextAt{13}{432.291}{669.579}{\Character{58}}%
\TextAt{8}{79.226}{646.176}{Furthermore,}%
\TextAt{8}{150.276}{646.176}{since}%
\TextAt{13}{178.944}{646.176}{\Character{27}}%
\TextAt{21}{185.598}{644.382}{i}%
\TextAt{9}{188.979}{646.176}{\Character{40}}%
\TextAt{14}{193.531}{646.176}{p}%
\TextAt{16}{199.514}{646.176}{\Character{79}}%
\TextAt{9}{209.364}{646.176}{\Character{41}}%
\TextAt{9}{217.219}{646.176}{\Character{61}}%
\TextAt{14}{229.663}{646.176}{p}%
\TextAt{16}{235.646}{646.176}{\Character{79}}%
\TextAt{8}{245.496}{646.176}{,}%
\TextAt{8}{252.645}{646.176}{we}%
\TextAt{8}{269.872}{646.176}{deduce}%
\TextAt{8}{308.882}{646.176}{from}%
\TextAt{14}{187.423}{617.449}{P}%
\TextAt{21}{197.316}{622.386}{\Character{23}}%
\TextAt{16}{205.819}{617.449}{\Character{106}}%
\TextAt{14}{212.461}{617.449}{p}%
\TextAt{16}{218.444}{617.449}{\Character{79}}%
\NativeArrowAt{231.607}{617.449}{11.955}{21.926}{\Longleftrightarrow}%
\TextAt{13}{256.853}{617.449}{\Character{27}}%
\TextAt{21}{263.507}{615.656}{i}%
\TextAt{9}{266.888}{617.449}{\Character{40}}%
\TextAt{14}{271.441}{617.449}{P}%
\TextAt{21}{281.334}{622.386}{\Character{23}}%
\TextAt{9}{286.516}{617.449}{\Character{41}}%
\TextAt{16}{294.389}{617.449}{\Character{106}}%
\TextAt{13}{301.031}{617.449}{\Character{27}}%
\TextAt{21}{307.684}{615.656}{i}%
\TextAt{9}{311.065}{617.449}{\Character{40}}%
\TextAt{14}{315.618}{617.449}{p}%
\TextAt{16}{321.601}{617.449}{\Character{79}}%
\TextAt{9}{331.451}{617.449}{\Character{41}}%
\NativeArrowAt{339.324}{617.449}{11.955}{21.926}{\Longleftrightarrow}%
\TextAt{9}{364.562}{617.449}{\Character{40}}%
\TextAt{13}{369.115}{617.449}{\Character{27}}%
\TextAt{21}{375.768}{615.656}{i}%
\TextAt{14}{379.15}{617.449}{P}%
\TextAt{9}{389.043}{617.449}{\Character{41}}%
\TextAt{21}{393.595}{622.386}{\Character{23}}%
\TextAt{16}{402.098}{617.449}{\Character{106}}%
\TextAt{14}{408.74}{617.449}{p}%
\TextAt{16}{414.723}{617.449}{\Character{79}}%
\TextAt{8}{79.226}{588.723}{the}%
\TextAt{8}{99.397}{588.723}{equality}%
\TextAt{8}{143.968}{588.723}{of}%
\TextAt{8}{157.312}{588.723}{the}%
\TextAt{13}{177.483}{588.723}{e}%
\TextAt{21}{182.909}{586.93}{i}%
\TextAt{8}{186.29}{588.723}{,}%
\TextAt{13}{193.457}{588.723}{i}%
\TextAt{9}{200.795}{588.723}{\Character{61}}%
\TextAt{9}{213.228}{588.723}{1}%
\TextAt{13}{219.098}{588.723}{\Character{59}}%
\TextAt{13}{224.322}{588.723}{\Character{58}}%
\TextAt{13}{229.559}{588.723}{\Character{58}}%
\TextAt{13}{234.795}{588.723}{\Character{58}}%
\TextAt{13}{240.02}{588.723}{\Character{59}}%
\TextAt{13}{245.256}{588.723}{r}%
\TextAt{8}{250.92}{588.723}{.}%
\TextAt{8}{259.411}{588.723}{Thus}%
\TextAt{8}{289.074}{588.723}{the}%
\TextAt{8}{309.245}{588.723}{prime}%
\TextAt{8}{342.423}{588.723}{decomposition}%
\TextAt{8}{419.572}{588.723}{of}%
\TextAt{14}{432.916}{588.723}{p}%
\TextAt{8}{442.811}{588.723}{in}%
\TextAt{16}{456.478}{588.723}{\Character{79}}%
\TextAt{8}{470.239}{588.723}{takes}%
\TextAt{8}{500.237}{588.723}{on}%
\TextAt{8}{516.51}{588.723}{the}%
\TextAt{8}{79.226}{571.677}{following}%
\TextAt{8}{128.661}{571.677}{simple}%
\TextAt{8}{165.136}{571.677}{form}%
\TextAt{8}{192.776}{571.677}{in}%
\TextAt{8}{206.429}{571.677}{the}%
\TextAt{8}{226.586}{571.677}{Galois}%
\TextAt{8}{262.499}{571.677}{case:}%
\TextAt{14}{263.711}{535.352}{p}%
\TextAt{9}{273.015}{535.352}{\Character{61}}%
\TextAt{20}{285.441}{555.795}{\Character{32}}%
\TextAt{20}{294.905}{546.709}{\Character{89}}%
\TextAt{21}{299.98}{521.603}{\Character{27}}%
\TextAt{13}{312.174}{535.352}{\Character{27}}%
\TextAt{14}{319.256}{535.352}{P}%
\TextAt{20}{329.149}{555.795}{\Character{33}}%
\TextAt{21}{338.614}{553.119}{e}%
\TextAt{13}{345.033}{535.352}{\Character{59}}%
\TextAt{8}{79.226}{498.765}{where}%
\TextAt{13}{113.093}{498.765}{\Character{27}}%
\TextAt{8}{124.13}{498.765}{varies}%
\TextAt{8}{157.1}{498.765}{over}%
\TextAt{8}{182.203}{498.765}{a}%
\TextAt{8}{192.016}{498.765}{system}%
\TextAt{8}{230.892}{498.765}{of}%
\TextAt{8}{244.279}{498.765}{representatives}%
\TextAt{8}{324.137}{498.765}{of}%
\TextAt{13}{337.525}{498.765}{G}%
\TextAt{9}{346.758}{498.765}{\Character{47}}%
\TextAt{13}{352.611}{498.765}{G}%
\TextAt{27}{361.845}{496.971}{P}%
\TextAt{8}{368.938}{498.765}{.}%
\TextAt{8}{377.559}{498.765}{The}%
\TextAt{8}{401.67}{498.765}{decomposition}%
\TextAt{8}{478.862}{498.765}{field}%
\TextAt{13}{504.276}{498.765}{Z}%
\TextAt{27}{512.288}{496.971}{P}%
\TextAt{8}{523.337}{498.765}{of}%
\TextAt{14}{79.226}{481.719}{P}%
\TextAt{8}{93.498}{481.719}{over}%
\TextAt{13}{119.025}{481.719}{K}%
\TextAt{8}{134.221}{481.719}{has}%
\TextAt{8}{155.575}{481.719}{the}%
\TextAt{8}{176.212}{481.719}{following}%
\TextAt{8}{226.128}{481.719}{significance}%
\TextAt{8}{288.453}{481.719}{for}%
\TextAt{8}{306.819}{481.719}{the}%
\TextAt{8}{327.456}{481.719}{decomposition}%
\TextAt{8}{405.072}{481.719}{of}%
\TextAt{14}{418.883}{481.719}{p}%
\TextAt{8}{429.244}{481.719}{and}%
\TextAt{8}{452.487}{481.719}{the}%
\TextAt{8}{473.125}{481.719}{invariants}%
\TextAt{13}{527.344}{481.719}{e}%
\TextAt{8}{79.226}{464.673}{and}%
\TextAt{13}{101.989}{464.673}{f}%
\TextAt{8}{109.035}{464.673}{.}%
\TextAt{12}{79.226}{438.739}{Proposition}%
\TextAt{12}{151.594}{438.739}{(9.3).}%
\TextAt{8}{186.61}{438.739}{Let}%
\TextAt{14}{207.281}{438.739}{P}%
\TextAt{21}{217.174}{436.946}{Z}%
\TextAt{9}{227.295}{438.739}{\Character{61}}%
\TextAt{14}{239.72}{438.739}{P}%
\TextAt{16}{251.659}{438.739}{\Character{92}}%
\TextAt{13}{261.674}{438.739}{Z}%
\TextAt{27}{269.687}{436.946}{P}%
\TextAt{8}{280.379}{438.739}{be}%
\TextAt{8}{296.004}{438.739}{the}%
\TextAt{8}{315.862}{438.739}{prime}%
\TextAt{8}{348.727}{438.739}{ideal}%
\TextAt{8}{376.391}{438.739}{of}%
\TextAt{13}{389.422}{438.739}{Z}%
\TextAt{27}{397.434}{436.946}{P}%
\TextAt{8}{408.126}{438.739}{below}%
\TextAt{14}{440.991}{438.739}{P}%
\TextAt{8}{450.885}{438.739}{.}%
\TextAt{8}{459.233}{438.739}{Then}%
\TextAt{8}{489.492}{438.739}{we}%
\TextAt{8}{506.42}{438.739}{have:}%
\TextAt{8}{89.579}{418.704}{(i)}%
\TextAt{14}{107.918}{418.704}{P}%
\TextAt{21}{117.812}{416.911}{Z}%
\TextAt{8}{128.509}{418.704}{is}%
\TextAt{8}{140.273}{418.704}{nonsplit}%
\TextAt{8}{185.212}{418.704}{in}%
\TextAt{13}{198.865}{418.704}{L}%
\TextAt{8}{206.829}{418.704}{,}%
\TextAt{8}{213.979}{418.704}{i.e.,}%
\TextAt{14}{236.084}{418.704}{P}%
\TextAt{8}{249.874}{418.704}{is}%
\TextAt{8}{261.638}{418.704}{the}%
\TextAt{8}{281.795}{418.704}{only}%
\TextAt{8}{307.486}{418.704}{prime}%
\TextAt{8}{340.65}{418.704}{ideal}%
\TextAt{8}{368.613}{418.704}{of}%
\TextAt{13}{381.943}{418.704}{L}%
\TextAt{8}{393.805}{418.704}{above}%
\TextAt{14}{426.981}{418.704}{P}%
\TextAt{21}{436.874}{416.911}{Z}%
\TextAt{8}{443.673}{418.704}{.}%
\TextAt{8}{86.327}{401.658}{(ii)}%
\TextAt{14}{107.918}{401.658}{P}%
\TextAt{8}{121.709}{401.658}{over}%
\TextAt{13}{146.755}{401.658}{Z}%
\TextAt{27}{154.767}{399.865}{P}%
\TextAt{8}{165.758}{401.658}{has}%
\TextAt{8}{186.632}{401.658}{ramification}%
\TextAt{8}{251.68}{401.658}{index}%
\TextAt{13}{283.218}{401.658}{e}%
\TextAt{8}{292.541}{401.658}{and}%
\TextAt{8}{315.303}{401.658}{inertia}%
\TextAt{8}{352.376}{401.658}{degree}%
\TextAt{13}{388.792}{401.658}{f}%
\TextAt{8}{395.838}{401.658}{.}%
\TextAt{8}{83.076}{384.612}{(iii)}%
\TextAt{8}{107.918}{384.612}{The}%
\TextAt{8}{131.972}{384.612}{ramification}%
\TextAt{8}{197.02}{384.612}{index}%
\TextAt{8}{228.558}{384.612}{and}%
\TextAt{8}{251.321}{384.612}{the}%
\TextAt{8}{271.477}{384.612}{inertia}%
\TextAt{8}{308.55}{384.612}{degree}%
\TextAt{8}{344.966}{384.612}{of}%
\TextAt{14}{358.296}{384.612}{P}%
\TextAt{21}{368.189}{382.819}{Z}%
\TextAt{8}{378.886}{384.612}{over}%
\TextAt{13}{403.932}{384.612}{K}%
\TextAt{8}{418.647}{384.612}{both}%
\TextAt{8}{446.287}{384.612}{equal}%
\TextAt{9}{477.179}{384.612}{1}%
\TextAt{8}{483.032}{384.612}{.}%
\TextAt{19}{79.226}{349.791}{Proof.}%
\TextAt{8}{116}{349.791}{(i)}%
\TextAt{8}{132.584}{349.791}{Since}%
\TextAt{13}{163.466}{349.791}{G}%
\TextAt{9}{172.7}{349.791}{\Character{40}}%
\TextAt{13}{177.252}{349.791}{L}%
\TextAt{16}{189.091}{349.791}{\Character{106}}%
\TextAt{13}{196.286}{349.791}{Z}%
\TextAt{27}{204.298}{347.997}{P}%
\TextAt{9}{211.392}{349.791}{\Character{41}}%
\TextAt{9}{219.809}{349.791}{\Character{61}}%
\TextAt{13}{232.797}{349.791}{G}%
\TextAt{27}{242.03}{347.997}{P}%
\TextAt{8}{249.124}{349.791}{,}%
\TextAt{8}{256.679}{349.791}{the}%
\TextAt{8}{277.16}{349.791}{prime}%
\TextAt{8}{310.649}{349.791}{ideals}%
\TextAt{8}{343.551}{349.791}{above}%
\TextAt{14}{377.051}{349.791}{P}%
\TextAt{21}{386.945}{347.997}{Z}%
\TextAt{8}{397.966}{349.791}{are}%
\TextAt{8}{417.802}{349.791}{the}%
\TextAt{13}{438.283}{349.791}{\Character{27}}%
\TextAt{14}{445.365}{349.791}{P}%
\TextAt{8}{455.259}{349.791}{,}%
\TextAt{8}{462.814}{349.791}{for}%
\TextAt{13}{481.023}{349.791}{\Character{27}}%
\NativeMathAt{491.98}{349.791}{11.955}{\in}%
\TextAt{13}{503.824}{349.791}{G}%
\TextAt{9}{513.058}{349.791}{\Character{40}}%
\TextAt{13}{517.61}{349.791}{L}%
\TextAt{16}{529.449}{349.791}{\Character{106}}%
\TextAt{13}{79.226}{332.745}{Z}%
\TextAt{27}{87.238}{330.951}{P}%
\TextAt{9}{94.332}{332.745}{\Character{41}}%
\TextAt{8}{98.884}{332.745}{,}%
\TextAt{8}{106.034}{332.745}{and}%
\TextAt{8}{128.796}{332.745}{they}%
\TextAt{8}{155.133}{332.745}{are}%
\TextAt{8}{174.644}{332.745}{all}%
\TextAt{8}{190.903}{332.745}{equal}%
\TextAt{8}{221.795}{332.745}{to}%
\TextAt{14}{236.106}{332.745}{P}%
\TextAt{8}{245.999}{332.745}{.}%
\TextAt{8}{97.159}{312.71}{(ii)}%
\TextAt{8}{116.794}{312.71}{Since}%
\TextAt{8}{147.475}{312.71}{in}%
\TextAt{8}{161.252}{312.71}{the}%
\TextAt{8}{181.532}{312.71}{Galois}%
\TextAt{8}{217.569}{312.71}{case,}%
\TextAt{8}{245.747}{312.71}{ramification}%
\TextAt{8}{310.919}{312.71}{indices}%
\TextAt{8}{349.467}{312.71}{and}%
\TextAt{8}{372.354}{312.71}{inertia}%
\TextAt{8}{409.551}{312.71}{degrees}%
\TextAt{8}{450.705}{312.71}{are}%
\TextAt{8}{470.34}{312.71}{independent}%
\TextAt{8}{79.226}{295.664}{of}%
\TextAt{8}{92.556}{295.664}{the}%
\TextAt{8}{112.713}{295.664}{prime}%
\TextAt{8}{145.876}{295.664}{divisor,}%
\TextAt{8}{187.241}{295.664}{the}%
\TextAt{8}{207.398}{295.664}{fundamental}%
\TextAt{8}{275.04}{295.664}{identity}%
\TextAt{8}{318.294}{295.664}{in}%
\TextAt{8}{331.947}{295.664}{this}%
\TextAt{8}{354.769}{295.664}{case}%
\TextAt{8}{379.54}{295.664}{reads}%
\TextAt{13}{284.294}{266.938}{n}%
\TextAt{9}{294.602}{266.938}{\Character{61}}%
\TextAt{13}{307.028}{266.938}{ef}%
\TextAt{13}{319.473}{266.938}{r}%
\TextAt{13}{324.435}{266.938}{\Character{59}}%
\TextAt{8}{79.226}{238.212}{where}%
\TextAt{13}{113.906}{238.212}{n}%
\TextAt{9}{125.697}{238.212}{\Character{58}\Character{61}}%
\TextAt{8}{142.857}{238.212}{\char35\relax }%
\TextAt{13}{152.613}{238.212}{G}%
\TextAt{8}{161.846}{238.212}{,}%
\TextAt{13}{170.083}{238.212}{r}%
\TextAt{9}{180.487}{238.212}{\Character{61}}%
\TextAt{9}{194.379}{238.212}{\Character{40}}%
\TextAt{13}{198.948}{238.212}{G}%
\TextAt{9}{212.986}{238.212}{\Character{58}}%
\TextAt{13}{221.041}{238.212}{G}%
\TextAt{27}{230.275}{236.418}{P}%
\TextAt{9}{237.368}{238.212}{\Character{41}}%
\TextAt{8}{241.921}{238.212}{.}%
\TextAt{8}{252.98}{238.212}{We}%
\TextAt{8}{273.994}{238.212}{see}%
\TextAt{8}{293.777}{238.212}{therefore}%
\TextAt{8}{343.748}{238.212}{that}%
\TextAt{8}{369.987}{238.212}{\char35\relax }%
\TextAt{13}{379.742}{238.212}{G}%
\TextAt{27}{388.976}{236.418}{P}%
\TextAt{9}{400.873}{238.212}{\Character{61}}%
\TextAt{9}{414.765}{238.212}{\Character{91}}%
\TextAt{13}{418.033}{238.212}{L}%
\TextAt{9}{430.801}{238.212}{\Character{58}}%
\TextAt{13}{438.857}{238.212}{Z}%
\TextAt{27}{446.869}{236.418}{P}%
\TextAt{9}{453.963}{238.212}{\Character{93}}%
\TextAt{9}{461.997}{238.212}{\Character{61}}%
\TextAt{13}{475.927}{238.212}{ef}%
\TextAt{8}{488.399}{238.212}{.}%
\TextAt{8}{499.458}{238.212}{Let}%
\TextAt{13}{521.297}{238.212}{e}%
\TextAt{3}{526.723}{242.55}{\Character{48}}%
\TextAt{8}{529.518}{238.212}{,}%
\TextAt{8}{79.226}{221.166}{resp.}%
\TextAt{13}{108.838}{221.166}{e}%
\TextAt{3}{114.263}{225.504}{\Character{48}\Character{48}}%
\TextAt{8}{119.355}{221.166}{,}%
\TextAt{8}{126.625}{221.166}{be}%
\TextAt{8}{142.646}{221.166}{the}%
\TextAt{8}{162.899}{221.166}{ramification}%
\TextAt{8}{228.044}{221.166}{index}%
\TextAt{8}{259.678}{221.166}{of}%
\TextAt{14}{273.105}{221.166}{P}%
\TextAt{8}{286.992}{221.166}{over}%
\TextAt{13}{312.135}{221.166}{Z}%
\TextAt{27}{320.147}{219.372}{P}%
\TextAt{8}{327.24}{221.166}{,}%
\TextAt{8}{334.51}{221.166}{resp.}%
\TextAt{8}{364.122}{221.166}{of}%
\TextAt{14}{377.549}{221.166}{P}%
\TextAt{21}{387.442}{219.372}{Z}%
\TextAt{8}{398.236}{221.166}{over}%
\TextAt{13}{423.378}{221.166}{K}%
\TextAt{8}{434.196}{221.166}{.}%
\TextAt{8}{442.934}{221.166}{Then}%
\TextAt{14}{473.588}{221.166}{p}%
\TextAt{9}{483.056}{221.166}{\Character{61}}%
\TextAt{14}{495.646}{221.166}{P}%
\TextAt{21}{505.54}{225.504}{e}%
\TextAt{23}{509.468}{228.317}{\Character{48}\Character{48}}%
\TextAt{21}{505.54}{218.145}{Z}%
\TextAt{16}{516.829}{221.166}{\Character{1}}%
\TextAt{16}{522.125}{221.166}{\Character{1}}%
\TextAt{16}{527.433}{221.166}{\Character{1}}%
\TextAt{8}{79.226}{204.12}{in}%
\TextAt{13}{93.361}{204.12}{Z}%
\TextAt{27}{101.373}{202.327}{P}%
\TextAt{8}{112.846}{204.12}{and}%
\TextAt{14}{136.091}{204.12}{P}%
\TextAt{21}{145.984}{202.327}{Z}%
\TextAt{9}{156.926}{204.12}{\Character{61}}%
\TextAt{14}{170.173}{204.12}{P}%
\TextAt{21}{180.066}{208.458}{e}%
\TextAt{23}{183.995}{211.271}{\Character{48}}%
\TextAt{8}{191.557}{204.12}{in}%
\TextAt{13}{205.692}{204.12}{L}%
\TextAt{8}{213.656}{204.12}{,}%
\TextAt{8}{221.408}{204.12}{so}%
\TextAt{8}{236.26}{204.12}{that}%
\TextAt{14}{262.111}{204.12}{p}%
\TextAt{9}{272.236}{204.12}{\Character{61}}%
\TextAt{14}{285.483}{204.12}{P}%
\TextAt{21}{295.377}{208.458}{e}%
\TextAt{23}{299.305}{211.271}{\Character{48}\Character{48}}%
\TextAt{21}{304.176}{208.458}{e}%
\TextAt{23}{308.104}{211.271}{\Character{48}}%
\TextAt{16}{313.279}{204.12}{\Character{1}}%
\TextAt{16}{318.575}{204.12}{\Character{1}}%
\TextAt{16}{323.883}{204.12}{\Character{1}}%
\TextAt{8}{329.219}{204.12}{,}%
\TextAt{8}{336.971}{204.12}{i.e.,}%
\TextAt{13}{359.679}{204.12}{e}%
\TextAt{9}{369.247}{204.12}{\Character{61}}%
\TextAt{13}{382.494}{204.12}{e}%
\TextAt{3}{387.919}{208.458}{\Character{48}}%
\TextAt{13}{390.714}{204.12}{e}%
\TextAt{3}{396.14}{208.458}{\Character{48}\Character{48}}%
\TextAt{8}{401.232}{204.12}{.}%
\TextAt{8}{411.126}{204.12}{One}%
\TextAt{8}{436.319}{204.12}{also}%
\TextAt{8}{460.281}{204.12}{obviously}%
\TextAt{8}{512.541}{204.12}{gets}%
\TextAt{8}{79.226}{187.074}{the}%
\TextAt{8}{99.172}{187.074}{analogous}%
\TextAt{8}{153.022}{187.074}{identity}%
\TextAt{8}{196.066}{187.074}{for}%
\TextAt{8}{213.74}{187.074}{the}%
\TextAt{8}{233.686}{187.074}{inertia}%
\TextAt{8}{270.548}{187.074}{degrees}%
\TextAt{13}{311.367}{187.074}{f}%
\TextAt{9}{321.735}{187.074}{\Character{61}}%
\TextAt{13}{334.16}{187.074}{f}%
\TextAt{3}{341.207}{191.412}{\Character{48}}%
\TextAt{13}{344.002}{187.074}{f}%
\TextAt{3}{351.048}{191.412}{\Character{48}\Character{48}}%
\TextAt{8}{356.14}{187.074}{.}%
\TextAt{8}{364.518}{187.074}{The}%
\TextAt{8}{388.361}{187.074}{fundamental}%
\TextAt{8}{455.793}{187.074}{identity}%
\TextAt{8}{498.836}{187.074}{for}%
\TextAt{8}{516.51}{187.074}{the}%
\TextAt{8}{79.226}{170.028}{decomposition}%
\TextAt{8}{155.678}{170.028}{of}%
\TextAt{14}{168.326}{170.028}{P}%
\TextAt{21}{178.219}{168.235}{Z}%
\TextAt{8}{188.234}{170.028}{in}%
\TextAt{13}{201.204}{170.028}{L}%
\TextAt{8}{212.384}{170.028}{then}%
\TextAt{8}{238.361}{170.028}{reads}%
\TextAt{9}{268.308}{170.028}{\Character{91}}%
\TextAt{13}{271.56}{170.028}{L}%
\TextAt{9}{282.845}{170.028}{\Character{58}}%
\TextAt{13}{289.418}{170.028}{Z}%
\TextAt{27}{297.43}{168.235}{P}%
\TextAt{9}{304.523}{170.028}{\Character{93}}%
\TextAt{9}{311.074}{170.028}{\Character{61}}%
\TextAt{13}{323.521}{170.028}{e}%
\TextAt{3}{328.947}{174.366}{\Character{48}}%
\TextAt{13}{331.742}{170.028}{f}%
\TextAt{3}{338.788}{174.366}{\Character{48}}%
\TextAt{8}{341.583}{170.028}{,}%
\TextAt{8}{348.187}{170.028}{i.e.,}%
\TextAt{8}{369.746}{170.028}{we}%
\TextAt{8}{386.291}{170.028}{have}%
\TextAt{13}{412.603}{170.028}{e}%
\TextAt{3}{418.029}{174.366}{\Character{48}}%
\TextAt{13}{420.824}{170.028}{f}%
\TextAt{3}{427.87}{174.366}{\Character{48}}%
\TextAt{9}{433.986}{170.028}{\Character{61}}%
\TextAt{13}{446.412}{170.028}{ef}%
\TextAt{8}{458.883}{170.028}{,}%
\TextAt{8}{465.487}{170.028}{and}%
\TextAt{8}{487.567}{170.028}{therefore}%
\TextAt{13}{79.226}{152.982}{e}%
\TextAt{3}{84.652}{157.32}{\Character{48}}%
\TextAt{9}{90.767}{152.982}{\Character{61}}%
\TextAt{13}{103.193}{152.982}{e}%
\TextAt{8}{108.618}{152.982}{,}%
\TextAt{13}{115.768}{152.982}{f}%
\TextAt{3}{122.814}{157.32}{\Character{48}}%
\TextAt{9}{128.93}{152.982}{\Character{61}}%
\TextAt{13}{141.355}{152.982}{f}%
\TextAt{8}{148.402}{152.982}{,}%
\TextAt{13}{155.551}{152.982}{e}%
\TextAt{3}{160.977}{157.32}{\Character{48}\Character{48}}%
\TextAt{9}{169.389}{152.982}{\Character{61}}%
\TextAt{13}{181.815}{152.982}{f}%
\TextAt{3}{188.861}{157.32}{\Character{48}\Character{48}}%
\TextAt{9}{197.274}{152.982}{\Character{61}}%
\TextAt{9}{209.684}{152.982}{1}%
\TextAt{8}{215.553}{152.982}{.}%
\Vector{q 0 G 0 g 0.399 w 452.334 -621.018 m 452.334 -612.948 l S Q}%
\Vector{q 0 G 0 g 0.399 w 452.534 -613.148 m 459.707 -613.148 l S Q}%
\Vector{q 0 G 0 g 0.399 w 452.534 -620.819 m 459.707 -620.819 l S Q}%
\Vector{q 0 G 0 g 0.399 w 459.906 -621.018 m 459.906 -612.948 l S Q}%
\TextAt{8}{97.159}{124.059}{The}%
\TextAt{8}{121.565}{124.059}{ramification}%
\TextAt{8}{186.965}{124.059}{index}%
\TextAt{13}{218.855}{124.059}{e}%
\TextAt{8}{228.529}{124.059}{and}%
\TextAt{8}{251.644}{124.059}{the}%
\TextAt{8}{272.153}{124.059}{inertia}%
\TextAt{8}{309.578}{124.059}{degree}%
\TextAt{13}{346.345}{124.059}{f}%
\TextAt{8}{357.641}{124.059}{admit}%
\TextAt{8}{391.814}{124.059}{a}%
\TextAt{8}{401.921}{124.059}{further}%
\TextAt{8}{441.618}{124.059}{interesting}%
\TextAt{8}{499.594}{124.059}{group-}%
\TextAt{8}{79.226}{107.013}{theoretic}%
\TextAt{8}{127.854}{107.013}{interpretation.}%
\TextAt{8}{206.836}{107.013}{Since}%
\TextAt{13}{237.244}{107.013}{\Character{27}}%
\TextAt{16}{244.327}{107.013}{\Character{79}}%
\TextAt{9}{257.497}{107.013}{\Character{61}}%
\TextAt{16}{269.923}{107.013}{\Character{79}}%
\TextAt{8}{283.521}{107.013}{and}%
\TextAt{13}{306.135}{107.013}{\Character{27}}%
\TextAt{14}{313.217}{107.013}{P}%
\TextAt{9}{326.431}{107.013}{\Character{61}}%
\TextAt{14}{338.857}{107.013}{P}%
\TextAt{8}{348.75}{107.013}{,}%
\TextAt{8}{355.78}{107.013}{every}%
\TextAt{13}{386.523}{107.013}{\Character{27}}%
\NativeMathAt{396.926}{107.013}{11.955}{\in}%
\TextAt{13}{408.217}{107.013}{G}%
\TextAt{27}{417.451}{105.22}{P}%
\TextAt{8}{428.293}{107.013}{induces}%
\TextAt{8}{469.819}{107.013}{an}%
\TextAt{8}{485.929}{107.013}{automor-}%
\TextAt{8}{79.226}{89.967}{phism}%
\Vector{q 0 G 0 g 0.478 w 109.47 -694.258 m 116.552 -694.258 l S Q}%
\TextAt{13}{181.47}{72.921}{\Character{27}}%
\TextAt{9}{191.873}{72.921}{\Character{58}}%
\TextAt{16}{198.446}{72.921}{\Character{79}}%
\TextAt{9}{208.295}{72.921}{\Character{47}}%
\TextAt{14}{214.148}{72.921}{P}%
\NativeArrowAt{227.362}{72.921}{11.955}{19.26}{\longrightarrow}%
\TextAt{16}{249.933}{72.921}{\Character{79}}%
\TextAt{9}{259.794}{72.921}{\Character{47}}%
\TextAt{14}{265.647}{72.921}{P}%
\TextAt{13}{275.541}{72.921}{\Character{59}}%
\TextAt{13}{304.676}{72.921}{a}%
\TextAt{8}{314.161}{72.921}{mod}%
\TextAt{14}{339.599}{72.921}{P}%
\NativeArrowAt{352.813}{72.921}{11.955}{19.26}{\longmapsto}%
\TextAt{13}{375.395}{72.921}{\Character{27}}%
\TextAt{13}{382.461}{72.921}{a}%
\TextAt{8}{391.943}{72.921}{mod}%
\TextAt{14}{417.381}{72.921}{P}%
\TextAt{13}{427.274}{72.921}{\Character{59}}%
\TextAt{8}{300.14}{43.033}{52}%
\Destination{79.23}{432.445}{nav.Item.7}%
\Destination{79.23}{415.705}{nav.Item.8}%
\Destination{79.23}{398.355}{nav.Item.9}%
\Destination{78.23}{809.815}{nav.page.52}%
\Destination{79.36}{364.755}{nav.print.56}%
\Destination{79.36}{452.755}{nav.number.1.9.3}%
}

% Page 65 (body page 53)
\Page{%
\TextAt{2}{79.226}{796.943}{I.}%
\TextAt{2}{89.83}{796.943}{Algebraic}%
\TextAt{2}{138.648}{796.943}{Integers}%
\Vector{q 0 G 0 g 0.299 w 7.226 18.729 m 460.769 18.729 l S Q}%
\TextAt{8}{79.226}{761.124}{of}%
\TextAt{8}{92.556}{761.124}{the}%
\TextAt{8}{112.713}{761.124}{residue}%
\TextAt{8}{152.44}{761.124}{class}%
\TextAt{8}{179.877}{761.124}{field}%
\TextAt{16}{205.234}{761.124}{\Character{79}}%
\TextAt{9}{215.083}{761.124}{\Character{47}}%
\TextAt{14}{220.936}{761.124}{P}%
\TextAt{8}{230.83}{761.124}{.}%
\TextAt{8}{239.278}{761.124}{Putting}%
\TextAt{13}{282.364}{761.124}{\Character{20}}%
\TextAt{9}{289.103}{761.124}{\Character{40}}%
\TextAt{14}{293.655}{761.124}{P}%
\TextAt{9}{303.549}{761.124}{\Character{41}}%
\TextAt{9}{311.404}{761.124}{\Character{61}}%
\TextAt{16}{323.847}{761.124}{\Character{79}}%
\TextAt{9}{333.697}{761.124}{\Character{47}}%
\TextAt{14}{339.55}{761.124}{P}%
\TextAt{8}{353.341}{761.124}{and}%
\TextAt{13}{376.103}{761.124}{\Character{20}}%
\TextAt{9}{382.842}{761.124}{\Character{40}}%
\TextAt{14}{387.394}{761.124}{p}%
\TextAt{9}{393.377}{761.124}{\Character{41}}%
\TextAt{9}{401.232}{761.124}{\Character{61}}%
\TextAt{14}{413.676}{761.124}{o}%
\TextAt{9}{419.52}{761.124}{\Character{47}}%
\TextAt{14}{425.373}{761.124}{p}%
\TextAt{8}{431.356}{761.124}{,}%
\TextAt{8}{438.505}{761.124}{we}%
\TextAt{8}{455.733}{761.124}{obtain}%
\TextAt{8}{492.16}{761.124}{the}%
\TextAt{12}{79.226}{736.639}{Proposition}%
\TextAt{12}{152.063}{736.639}{(9.4).}%
\TextAt{8}{187.488}{736.639}{The}%
\TextAt{8}{211.652}{736.639}{extension}%
\TextAt{13}{263.527}{736.639}{\Character{20}}%
\TextAt{9}{270.265}{736.639}{\Character{40}}%
\TextAt{14}{274.818}{736.639}{P}%
\TextAt{9}{284.711}{736.639}{\Character{41}}%
\TextAt{16}{292.771}{736.639}{\Character{106}}%
\TextAt{13}{299.599}{736.639}{\Character{20}}%
\TextAt{9}{306.337}{736.639}{\Character{40}}%
\TextAt{14}{310.89}{736.639}{p}%
\TextAt{9}{316.873}{736.639}{\Character{41}}%
\TextAt{8}{325.432}{736.639}{is}%
\TextAt{8}{337.305}{736.639}{normal}%
\TextAt{8}{377.094}{736.639}{and}%
\TextAt{8}{399.966}{736.639}{admits}%
\TextAt{8}{438.511}{736.639}{a}%
\TextAt{8}{448.376}{736.639}{surjective}%
\TextAt{8}{500.897}{736.639}{homo-}%
\TextAt{8}{79.226}{719.593}{morphism}%
\TextAt{13}{245.325}{702.547}{G}%
\TextAt{27}{254.559}{700.754}{P}%
\NativeArrowAt{264.973}{702.547}{11.955}{19.26}{\longrightarrow}%
\TextAt{13}{287.555}{702.547}{G}%
\TextAt{9}{296.789}{702.547}{\Character{40}}%
\TextAt{13}{301.341}{702.547}{\Character{20}}%
\TextAt{9}{308.08}{702.547}{\Character{40}}%
\TextAt{14}{312.632}{702.547}{P}%
\TextAt{9}{322.525}{702.547}{\Character{41}}%
\TextAt{16}{330.398}{702.547}{\Character{106}}%
\TextAt{13}{337.04}{702.547}{\Character{20}}%
\TextAt{9}{343.779}{702.547}{\Character{40}}%
\TextAt{14}{348.331}{702.547}{p}%
\TextAt{9}{354.314}{702.547}{\Character{41}\Character{41}}%
\TextAt{13}{363.419}{702.547}{\Character{58}}%
\TextAt{19}{79.226}{666.439}{Proof.}%
\TextAt{8}{116}{666.439}{The}%
\TextAt{8}{141.493}{666.439}{inertia}%
\TextAt{8}{180.005}{666.439}{degree}%
\TextAt{8}{217.86}{666.439}{of}%
\TextAt{14}{232.629}{666.439}{P}%
\TextAt{21}{242.522}{664.646}{Z}%
\TextAt{8}{254.658}{666.439}{over}%
\TextAt{13}{281.143}{666.439}{K}%
\TextAt{8}{297.297}{666.439}{equals}%
\TextAt{9}{334.243}{666.439}{1}%
\TextAt{8}{340.096}{666.439}{,}%
\TextAt{8}{349.044}{666.439}{i.e.,}%
\TextAt{13}{372.948}{666.439}{Z}%
\TextAt{27}{380.96}{664.646}{P}%
\TextAt{8}{393.39}{666.439}{has}%
\TextAt{8}{415.703}{666.439}{the}%
\TextAt{8}{437.299}{666.439}{same}%
\TextAt{8}{468.064}{666.439}{residue}%
\TextAt{8}{509.23}{666.439}{class}%
\TextAt{8}{79.226}{649.393}{field}%
\TextAt{13}{104.806}{649.393}{\Character{20}}%
\TextAt{9}{111.545}{649.393}{\Character{40}}%
\TextAt{14}{116.097}{649.393}{p}%
\TextAt{9}{122.08}{649.393}{\Character{41}}%
\TextAt{8}{130.753}{649.393}{as}%
\TextAt{13}{145.347}{649.393}{K}%
\TextAt{8}{160.285}{649.393}{with}%
\TextAt{8}{187.168}{649.393}{respect}%
\TextAt{8}{227.441}{649.393}{to}%
\TextAt{14}{241.975}{649.393}{p}%
\TextAt{8}{247.958}{649.393}{.}%
\TextAt{8}{257.076}{649.393}{Therefore}%
\TextAt{8}{310.297}{649.393}{we}%
\TextAt{8}{327.748}{649.393}{may,}%
\TextAt{8}{355.667}{649.393}{and}%
\TextAt{8}{378.653}{649.393}{do,}%
\TextAt{8}{398.443}{649.393}{assume}%
\TextAt{8}{439.11}{649.393}{that}%
\TextAt{13}{464.703}{649.393}{Z}%
\TextAt{27}{472.715}{647.6}{P}%
\TextAt{9}{483.51}{649.393}{\Character{61}}%
\TextAt{13}{496.316}{649.393}{K}%
\TextAt{8}{507.133}{649.393}{,}%
\TextAt{8}{514.562}{649.393}{i.e.,}%
\TextAt{13}{79.226}{632.347}{G}%
\TextAt{27}{88.46}{630.554}{P}%
\TextAt{9}{99.486}{632.347}{\Character{61}}%
\TextAt{13}{112.524}{632.347}{G}%
\TextAt{8}{121.758}{632.347}{.}%
\TextAt{8}{131.284}{632.347}{Let}%
\TextAt{13}{152.613}{632.347}{\Character{18}}%
\NativeMathAt{162.326}{632.347}{11.955}{\in}%
\TextAt{16}{174.21}{632.347}{\Character{79}}%
\TextAt{8}{188.336}{632.347}{be}%
\TextAt{8}{204.619}{632.347}{a}%
\TextAt{8}{214.734}{632.347}{representative}%
\TextAt{8}{290.279}{632.347}{of}%
\TextAt{8}{303.968}{632.347}{an}%
\TextAt{8}{320.587}{632.347}{element}%
\Vector{q 0 G 0 g 0.478 w 292.188 -131.677 m 297.968 -131.677 l S Q}%
\TextAt{13}{364.188}{632.347}{\Character{18}}%
\NativeMathAt{373.901}{632.347}{11.955}{\in}%
\TextAt{13}{385.804}{632.347}{\Character{20}}%
\TextAt{9}{392.543}{632.347}{\Character{40}}%
\TextAt{14}{397.095}{632.347}{P}%
\TextAt{9}{406.989}{632.347}{\Character{41}}%
\TextAt{8}{415.797}{632.347}{and}%
\TextAt{13}{438.919}{632.347}{f}%
\TextAt{9}{445.966}{632.347}{\Character{40}}%
\TextAt{13}{450.518}{632.347}{X}%
\TextAt{9}{461.173}{632.347}{\Character{41}}%
\TextAt{8}{465.726}{632.347}{,}%
\TextAt{8}{473.324}{632.347}{resp.}%
\TextAt{13}{503.724}{632.347}{g}%
\TextAt{9}{509.758}{632.347}{\Character{40}}%
\TextAt{13}{514.31}{632.347}{X}%
\TextAt{9}{524.965}{632.347}{\Character{41}}%
\TextAt{8}{529.518}{632.347}{,}%
\TextAt{8}{79.226}{615.301}{the}%
\TextAt{8}{99.845}{615.301}{minimal}%
\TextAt{8}{145.834}{615.301}{polynomial}%
\TextAt{8}{206.79}{615.301}{of}%
\TextAt{13}{220.582}{615.301}{\Character{18}}%
\TextAt{8}{230.723}{615.301}{over}%
\TextAt{13}{256.232}{615.301}{K}%
\TextAt{8}{267.05}{615.301}{,}%
\TextAt{8}{274.777}{615.301}{resp.}%
\TextAt{8}{305.488}{615.301}{of}%
\Vector{q 0 G 0 g 0.478 w 247.281 -148.723 m 253.061 -148.723 l S Q}%
\TextAt{13}{319.281}{615.301}{\Character{18}}%
\TextAt{8}{329.421}{615.301}{over}%
\TextAt{13}{354.93}{615.301}{\Character{20}}%
\TextAt{9}{361.669}{615.301}{\Character{40}}%
\TextAt{14}{366.221}{615.301}{p}%
\TextAt{9}{372.204}{615.301}{\Character{41}}%
\TextAt{8}{376.757}{615.301}{.}%
\TextAt{8}{386.594}{615.301}{Then}%
\Vector{q 0 G 0 g 0.478 w 345.614 -148.723 m 351.394 -148.723 l S Q}%
\TextAt{13}{417.614}{615.301}{\Character{18}}%
\TextAt{9}{427.504}{615.301}{\Character{61}}%
\TextAt{13}{440.718}{615.301}{\Character{18}}%
\TextAt{8}{449.819}{615.301}{mod}%
\TextAt{14}{475.257}{615.301}{P}%
\TextAt{8}{489.511}{615.301}{is}%
\TextAt{8}{501.737}{615.301}{a}%
\TextAt{8}{511.956}{615.301}{zero}%
\TextAt{8}{79.226}{598.255}{of}%
\TextAt{8}{92.943}{598.255}{the}%
\TextAt{8}{113.487}{598.255}{polynomial}%
\Vector{q 0 G 0 g 0.478 w 102.367 -165.769 m 109.413 -165.769 l S Q}%
\TextAt{13}{174.367}{598.255}{f}%
\TextAt{9}{181.413}{598.255}{\Character{40}}%
\TextAt{13}{185.965}{598.255}{X}%
\TextAt{9}{196.62}{598.255}{\Character{41}}%
\TextAt{9}{205.132}{598.255}{\Character{61}}%
\TextAt{13}{218.238}{598.255}{f}%
\TextAt{9}{225.284}{598.255}{\Character{40}}%
\TextAt{13}{229.837}{598.255}{X}%
\TextAt{9}{240.492}{598.255}{\Character{41}}%
\TextAt{8}{248.365}{598.255}{mod}%
\TextAt{14}{273.803}{598.255}{p}%
\TextAt{8}{279.786}{598.255}{,}%
\TextAt{8}{287.419}{598.255}{i.e.,}%
\Vector{q 0 G 0 g 0.478 w 238.008 -168.924 m 244.042 -168.924 l S Q}%
\TextAt{13}{310.008}{598.255}{g}%
\TextAt{9}{316.042}{598.255}{\Character{40}}%
\TextAt{13}{320.594}{598.255}{X}%
\TextAt{9}{331.249}{598.255}{\Character{41}}%
\TextAt{8}{340.086}{598.255}{divides}%
\Vector{q 0 G 0 g 0.478 w 307.877 -165.769 m 314.924 -165.769 l S Q}%
\TextAt{13}{379.877}{598.255}{f}%
\TextAt{9}{386.924}{598.255}{\Character{40}}%
\TextAt{13}{391.476}{598.255}{X}%
\TextAt{9}{402.131}{598.255}{\Character{41}}%
\TextAt{8}{406.684}{598.255}{.}%
\TextAt{8}{416.293}{598.255}{Since}%
\TextAt{13}{447.237}{598.255}{L}%
\TextAt{16}{459.182}{598.255}{\Character{106}}%
\TextAt{13}{466.483}{598.255}{K}%
\TextAt{8}{481.585}{598.255}{is}%
\TextAt{8}{493.736}{598.255}{normal,}%
\TextAt{13}{79.226}{581.209}{f}%
\TextAt{9}{86.273}{581.209}{\Character{40}}%
\TextAt{13}{90.825}{581.209}{X}%
\TextAt{9}{101.48}{581.209}{\Character{41}}%
\TextAt{8}{110.546}{581.209}{splits}%
\TextAt{8}{141.85}{581.209}{over}%
\TextAt{16}{167.512}{581.209}{\Character{79}}%
\TextAt{8}{181.875}{581.209}{into}%
\TextAt{8}{206.234}{581.209}{linear}%
\TextAt{8}{239.368}{581.209}{factors.}%
\TextAt{8}{283.88}{581.209}{Hence}%
\Vector{q 0 G 0 g 0.478 w 247.273 -182.815 m 254.32 -182.815 l S Q}%
\TextAt{13}{319.273}{581.209}{f}%
\TextAt{9}{326.32}{581.209}{\Character{40}}%
\TextAt{13}{330.872}{581.209}{X}%
\TextAt{9}{341.527}{581.209}{\Character{41}}%
\TextAt{8}{350.593}{581.209}{splits}%
\TextAt{8}{381.898}{581.209}{into}%
\TextAt{8}{406.256}{581.209}{linear}%
\TextAt{8}{439.39}{581.209}{factors}%
\TextAt{8}{478.119}{581.209}{over}%
\TextAt{13}{503.781}{581.209}{\Character{20}}%
\TextAt{9}{510.52}{581.209}{\Character{40}}%
\TextAt{14}{515.072}{581.209}{P}%
\TextAt{9}{524.965}{581.209}{\Character{41}}%
\TextAt{8}{529.518}{581.209}{,}%
\TextAt{8}{79.226}{564.163}{and}%
\TextAt{8}{101.989}{564.163}{the}%
\TextAt{8}{122.145}{564.163}{same}%
\TextAt{8}{151.471}{564.163}{is}%
\TextAt{8}{163.235}{564.163}{true}%
\TextAt{8}{187.946}{564.163}{of}%
\Vector{q 0 G 0 g 0.478 w 129.276 -203.016 m 135.311 -203.016 l S Q}%
\TextAt{13}{201.276}{564.163}{g}%
\TextAt{9}{207.311}{564.163}{\Character{40}}%
\TextAt{13}{211.863}{564.163}{X}%
\TextAt{9}{222.518}{564.163}{\Character{41}}%
\TextAt{8}{227.071}{564.163}{.}%
\TextAt{8}{235.519}{564.163}{In}%
\TextAt{8}{250.14}{564.163}{other}%
\TextAt{8}{280.709}{564.163}{words,}%
\TextAt{13}{317.519}{564.163}{\Character{20}}%
\TextAt{9}{324.258}{564.163}{\Character{40}}%
\TextAt{14}{328.81}{564.163}{P}%
\TextAt{9}{338.703}{564.163}{\Character{41}}%
\TextAt{16}{346.577}{564.163}{\Character{106}}%
\TextAt{13}{353.218}{564.163}{\Character{20}}%
\TextAt{9}{359.957}{564.163}{\Character{40}}%
\TextAt{14}{364.509}{564.163}{p}%
\TextAt{9}{370.492}{564.163}{\Character{41}}%
\TextAt{8}{378.942}{564.163}{is}%
\TextAt{8}{390.706}{564.163}{a}%
\TextAt{8}{400.461}{564.163}{normal}%
\TextAt{8}{440.141}{564.163}{extension.}%
\TextAt{8}{97.159}{544.129}{Now}%
\TextAt{8}{122.704}{544.129}{let}%
\Vector{q 0 G 0 g 0.478 w 66.494 -219.896 m 72.274 -219.896 l S Q}%
\TextAt{13}{138.494}{544.129}{\Character{18}}%
\TextAt{8}{147.057}{544.129}{be}%
\TextAt{8}{161.866}{544.129}{a}%
\TextAt{8}{170.507}{544.129}{primitive}%
\TextAt{8}{219.472}{544.129}{element}%
\TextAt{8}{261.599}{544.129}{for}%
\TextAt{8}{278.369}{544.129}{the}%
\TextAt{8}{297.411}{544.129}{maximal}%
\TextAt{8}{344.105}{544.129}{separable}%
\TextAt{8}{394.433}{544.129}{subextension}%
\TextAt{8}{463.029}{544.129}{of}%
\TextAt{13}{475.244}{544.129}{\Character{20}}%
\TextAt{9}{481.983}{544.129}{\Character{40}}%
\TextAt{14}{486.535}{544.129}{P}%
\TextAt{9}{496.428}{544.129}{\Character{41}}%
\TextAt{16}{504.301}{544.129}{\Character{106}}%
\TextAt{13}{510.943}{544.129}{\Character{20}}%
\TextAt{9}{517.682}{544.129}{\Character{40}}%
\TextAt{14}{522.234}{544.129}{p}%
\TextAt{9}{528.217}{544.129}{\Character{41}}%
\TextAt{8}{79.226}{527.083}{and}%
\Vector{q 0 G 0 g 0.478 w 132.301 -257.142 m 139.383 -257.142 l S Q}%
\TextAt{13}{204.301}{510.037}{\Character{27}}%
\NativeMathAt{214.704}{510.037}{11.955}{\in}%
\TextAt{13}{225.995}{510.037}{G}%
\TextAt{9}{235.229}{510.037}{\Character{40}}%
\TextAt{13}{239.781}{510.037}{\Character{20}}%
\TextAt{9}{246.519}{510.037}{\Character{40}}%
\TextAt{14}{251.072}{510.037}{P}%
\TextAt{9}{260.965}{510.037}{\Character{41}}%
\TextAt{16}{268.838}{510.037}{\Character{106}}%
\TextAt{13}{275.48}{510.037}{\Character{20}}%
\TextAt{9}{282.218}{510.037}{\Character{40}}%
\TextAt{14}{286.771}{510.037}{p}%
\TextAt{9}{292.754}{510.037}{\Character{41}\Character{41}}%
\TextAt{9}{305.151}{510.037}{\Character{61}}%
\TextAt{13}{317.605}{510.037}{G}%
\TextAt{9}{326.838}{510.037}{\Character{40}}%
\TextAt{13}{331.391}{510.037}{\Character{20}}%
\TextAt{9}{338.129}{510.037}{\Character{40}}%
\TextAt{14}{342.682}{510.037}{p}%
\TextAt{9}{348.665}{510.037}{\Character{41}\Character{40}}%
\Vector{q 0 G 0 g 0.478 w 285.769 -253.988 m 291.55 -253.988 l S Q}%
\TextAt{13}{357.769}{510.037}{\Character{18}}%
\TextAt{9}{363.55}{510.037}{\Character{41}}%
\TextAt{16}{371.423}{510.037}{\Character{106}}%
\TextAt{13}{378.065}{510.037}{\Character{20}}%
\TextAt{9}{384.803}{510.037}{\Character{40}}%
\TextAt{14}{389.356}{510.037}{p}%
\TextAt{9}{395.339}{510.037}{\Character{41}\Character{41}}%
\TextAt{13}{404.443}{510.037}{\Character{58}}%
\TextAt{8}{79.226}{488.806}{Then}%
\Vector{q 0 G 0 g 0.478 w 38.443 -278.373 m 45.526 -278.373 l S Q}%
\TextAt{13}{110.443}{488.806}{\Character{27}}%
\Vector{q 0 G 0 g 0.478 w 45.526 -275.218 m 51.306 -275.218 l S Q}%
\TextAt{13}{117.526}{488.806}{\Character{18}}%
\TextAt{8}{127.863}{488.806}{is}%
\TextAt{8}{140.287}{488.806}{a}%
\TextAt{8}{150.702}{488.806}{root}%
\TextAt{8}{176.407}{488.806}{of}%
\Vector{q 0 G 0 g 0.478 w 118.397 -278.373 m 124.431 -278.373 l S Q}%
\TextAt{13}{190.397}{488.806}{g}%
\TextAt{9}{196.431}{488.806}{\Character{40}}%
\TextAt{13}{200.984}{488.806}{X}%
\TextAt{9}{211.639}{488.806}{\Character{41}}%
\TextAt{8}{216.191}{488.806}{,}%
\TextAt{8}{224.165}{488.806}{and}%
\TextAt{8}{247.587}{488.806}{hence}%
\TextAt{8}{280.753}{488.806}{of}%
\Vector{q 0 G 0 g 0.478 w 222.743 -275.218 m 229.789 -275.218 l S Q}%
\TextAt{13}{294.743}{488.806}{f}%
\TextAt{9}{301.789}{488.806}{\Character{40}}%
\TextAt{13}{306.341}{488.806}{X}%
\TextAt{9}{316.997}{488.806}{\Character{41}}%
\TextAt{8}{321.549}{488.806}{,}%
\TextAt{8}{329.523}{488.806}{i.e.,}%
\TextAt{8}{352.452}{488.806}{there}%
\TextAt{8}{383.024}{488.806}{exists}%
\TextAt{8}{415.998}{488.806}{a}%
\TextAt{8}{426.413}{488.806}{zero}%
\TextAt{13}{451.784}{488.806}{\Character{18}}%
\TextAt{3}{457.565}{493.145}{\Character{48}}%
\TextAt{8}{464.917}{488.806}{of}%
\TextAt{13}{478.907}{488.806}{f}%
\TextAt{9}{485.953}{488.806}{\Character{40}}%
\TextAt{13}{490.505}{488.806}{X}%
\TextAt{9}{501.161}{488.806}{\Character{41}}%
\TextAt{8}{510.27}{488.806}{such}%
\TextAt{8}{79.226}{471.76}{that}%
\TextAt{13}{104.736}{471.76}{\Character{18}}%
\TextAt{3}{110.516}{476.099}{\Character{48}}%
\NativeMathAt{116.873}{471.76}{11.955}{\equiv}%
\Vector{q 0 G 0 g 0.478 w 57.732 -295.419 m 64.815 -295.419 l S Q}%
\TextAt{13}{129.732}{471.76}{\Character{27}}%
\Vector{q 0 G 0 g 0.478 w 64.815 -292.264 m 70.595 -292.264 l S Q}%
\TextAt{13}{136.815}{471.76}{\Character{18}}%
\TextAt{9}{147.908}{471.76}{\Character{40}}%
\TextAt{8}{152.461}{471.76}{mod}%
\TextAt{14}{178.563}{471.76}{P}%
\TextAt{9}{188.456}{471.76}{\Character{41}}%
\TextAt{8}{193.008}{471.76}{.}%
\TextAt{13}{201.88}{471.76}{\Character{18}}%
\TextAt{3}{207.66}{476.099}{\Character{48}}%
\TextAt{8}{214.493}{471.76}{is}%
\TextAt{8}{226.398}{471.76}{a}%
\TextAt{8}{236.295}{471.76}{conjugate}%
\TextAt{8}{289.445}{471.76}{of}%
\TextAt{13}{302.916}{471.76}{\Character{18}}%
\TextAt{8}{308.696}{471.76}{,}%
\TextAt{8}{316.022}{471.76}{i.e.,}%
\TextAt{13}{338.303}{471.76}{\Character{18}}%
\TextAt{3}{344.084}{476.099}{\Character{48}}%
\TextAt{9}{350.44}{471.76}{\Character{61}}%
\TextAt{13}{363.106}{471.76}{\Character{27}}%
\TextAt{13}{370.172}{471.76}{\Character{18}}%
\TextAt{8}{380.007}{471.76}{for}%
\TextAt{8}{398.033}{471.76}{some}%
\TextAt{13}{427.5}{471.76}{\Character{27}}%
\NativeMathAt{438.143}{471.76}{11.955}{\in}%
\TextAt{13}{449.675}{471.76}{G}%
\TextAt{9}{458.908}{471.76}{\Character{40}}%
\TextAt{13}{463.46}{471.76}{L}%
\TextAt{16}{474.986}{471.76}{\Character{106}}%
\TextAt{13}{481.868}{471.76}{K}%
\TextAt{9}{492.686}{471.76}{\Character{41}}%
\TextAt{8}{497.238}{471.76}{.}%
\TextAt{8}{506.109}{471.76}{Since}%
\TextAt{13}{79.226}{454.714}{\Character{27}}%
\TextAt{13}{86.292}{454.714}{\Character{18}}%
\NativeMathAt{95.41}{454.714}{11.955}{\equiv}%
\Vector{q 0 G 0 g 0.478 w 36.029 -312.465 m 43.111 -312.465 l S Q}%
\TextAt{13}{108.029}{454.714}{\Character{27}}%
\Vector{q 0 G 0 g 0.478 w 43.111 -309.31 m 48.892 -309.31 l S Q}%
\TextAt{13}{115.111}{454.714}{\Character{18}}%
\TextAt{9}{126.205}{454.714}{\Character{40}}%
\TextAt{8}{130.757}{454.714}{mod}%
\TextAt{14}{156.86}{454.714}{P}%
\TextAt{9}{166.753}{454.714}{\Character{41}}%
\TextAt{8}{171.305}{454.714}{,}%
\TextAt{8}{178.431}{454.714}{the}%
\TextAt{8}{198.558}{454.714}{automorphism}%
\TextAt{13}{275.998}{454.714}{\Character{27}}%
\TextAt{8}{286.949}{454.714}{is}%
\TextAt{8}{298.683}{454.714}{mapped}%
\TextAt{8}{343.199}{454.714}{by}%
\TextAt{8}{359.429}{454.714}{the}%
\TextAt{8}{379.556}{454.714}{homomorphism}%
\TextAt{8}{462.197}{454.714}{in}%
\TextAt{8}{475.82}{454.714}{question}%
\TextAt{8}{522.357}{454.714}{to}%
\Vector{q 0 G 0 g 0.478 w 7.226 -329.51 m 14.309 -329.51 l S Q}%
\TextAt{13}{79.226}{437.668}{\Character{27}}%
\TextAt{8}{86.309}{437.668}{.}%
\TextAt{8}{94.757}{437.668}{This}%
\TextAt{8}{121.477}{437.668}{proves}%
\TextAt{8}{157.641}{437.668}{the}%
\TextAt{8}{177.797}{437.668}{surjectivity.}%
\Vector{q 0 G 0 g 0.399 w 452.334 -336.332 m 452.334 -328.262 l S Q}%
\Vector{q 0 G 0 g 0.399 w 452.534 -328.461 m 459.707 -328.461 l S Q}%
\Vector{q 0 G 0 g 0.399 w 452.534 -336.132 m 459.707 -336.132 l S Q}%
\Vector{q 0 G 0 g 0.399 w 459.906 -336.332 m 459.906 -328.262 l S Q}%
\TextAt{12}{79.226}{405.745}{Definition}%
\TextAt{12}{142.015}{405.745}{(9.5).}%
\TextAt{8}{177.33}{405.745}{The}%
\TextAt{8}{201.384}{405.745}{kernel}%
\TextAt{13}{235.851}{405.745}{I}%
\TextAt{27}{241.014}{403.952}{P}%
\NativeMathAt{251.428}{405.745}{11.955}{\subseteq}%
\TextAt{13}{264.047}{405.745}{G}%
\TextAt{27}{273.281}{403.952}{P}%
\TextAt{8}{284.272}{405.745}{of}%
\TextAt{8}{297.602}{405.745}{the}%
\TextAt{8}{317.759}{405.745}{homomorphism}%
\TextAt{13}{246.951}{382.09}{G}%
\TextAt{27}{256.185}{380.297}{P}%
\NativeArrowAt{266.599}{382.09}{11.955}{19.26}{\longrightarrow}%
\TextAt{13}{289.181}{382.09}{G}%
\TextAt{9}{298.415}{382.09}{\Character{40}}%
\TextAt{13}{302.967}{382.09}{\Character{20}}%
\TextAt{9}{309.706}{382.09}{\Character{40}}%
\TextAt{14}{314.258}{382.09}{P}%
\TextAt{9}{324.151}{382.09}{\Character{41}}%
\TextAt{16}{332.024}{382.09}{\Character{106}}%
\TextAt{13}{338.666}{382.09}{\Character{20}}%
\TextAt{9}{345.405}{382.09}{\Character{40}}%
\TextAt{14}{349.957}{382.09}{p}%
\TextAt{9}{355.94}{382.09}{\Character{41}\Character{41}}%
\TextAt{8}{79.226}{358.436}{is}%
\TextAt{8}{90.99}{358.436}{called}%
\TextAt{8}{124.154}{358.436}{the}%
\TextAt{12}{144.31}{358.436}{inertia}%
\TextAt{12}{187.157}{358.436}{group}%
\TextAt{8}{224.96}{358.436}{of}%
\TextAt{14}{238.29}{358.436}{P}%
\TextAt{8}{252.08}{358.436}{over}%
\TextAt{13}{277.126}{358.436}{K}%
\TextAt{8}{287.944}{358.436}{.}%
\TextAt{8}{296.392}{358.436}{The}%
\TextAt{8}{320.446}{358.436}{fixed}%
\TextAt{8}{348.732}{358.436}{field}%
\TextAt{13}{211.512}{334.781}{T}%
\TextAt{27}{218.373}{332.988}{P}%
\TextAt{9}{228.787}{334.781}{\Character{61}}%
\TextAt{16}{241.213}{334.781}{\Character{102}}%
\TextAt{13}{247.19}{334.781}{x}%
\NativeMathAt{257.163}{334.781}{11.955}{\in}%
\TextAt{13}{268.454}{334.781}{L}%
\TextAt{16}{279.74}{334.781}{\Character{106}}%
\TextAt{13}{286.381}{334.781}{\Character{27}}%
\TextAt{13}{293.447}{334.781}{x}%
\TextAt{9}{303.437}{334.781}{\Character{61}}%
\TextAt{13}{315.862}{334.781}{x}%
\TextAt{8}{326.412}{334.781}{for}%
\TextAt{8}{344.297}{334.781}{all}%
\TextAt{13}{360.556}{334.781}{\Character{27}}%
\NativeMathAt{370.959}{334.781}{11.955}{\in}%
\TextAt{13}{382.25}{334.781}{I}%
\TextAt{27}{387.413}{332.988}{P}%
\TextAt{16}{394.506}{334.781}{\Character{103}}%
\TextAt{8}{79.226}{311.127}{is}%
\TextAt{8}{90.99}{311.127}{called}%
\TextAt{8}{124.154}{311.127}{the}%
\TextAt{12}{144.31}{311.127}{inertia}%
\TextAt{12}{187.157}{311.127}{field}%
\TextAt{8}{215.874}{311.127}{of}%
\TextAt{14}{229.204}{311.127}{P}%
\TextAt{8}{242.994}{311.127}{over}%
\TextAt{13}{268.04}{311.127}{K}%
\TextAt{8}{278.858}{311.127}{.}%
\TextAt{8}{97.159}{286.642}{This}%
\TextAt{8}{123.879}{286.642}{inertia}%
\TextAt{8}{160.952}{286.642}{field}%
\TextAt{13}{186.309}{286.642}{T}%
\TextAt{27}{193.17}{284.849}{P}%
\TextAt{8}{204.161}{286.642}{appears}%
\TextAt{8}{247.474}{286.642}{in}%
\TextAt{8}{261.127}{286.642}{the}%
\TextAt{8}{281.283}{286.642}{tower}%
\TextAt{8}{313.156}{286.642}{of}%
\TextAt{8}{326.486}{286.642}{fields}%
\TextAt{13}{256.54}{262.988}{K}%
\NativeMathAt{270.679}{262.988}{11.955}{\subseteq}%
\TextAt{13}{283.298}{262.988}{Z}%
\TextAt{27}{291.311}{261.194}{P}%
\NativeMathAt{301.725}{262.988}{11.955}{\subseteq}%
\TextAt{13}{314.344}{262.988}{T}%
\TextAt{27}{321.205}{261.194}{P}%
\NativeMathAt{331.62}{262.988}{11.955}{\subseteq}%
\TextAt{13}{344.239}{262.988}{L\Character{59}}%
\TextAt{8}{79.226}{239.333}{and}%
\TextAt{8}{101.989}{239.333}{we}%
\TextAt{8}{119.216}{239.333}{have}%
\TextAt{8}{146.211}{239.333}{the}%
\TextAt{8}{166.367}{239.333}{exact}%
\TextAt{8}{197.26}{239.333}{sequence}%
\TextAt{9}{194.49}{215.679}{1}%
\NativeArrowAt{203.663}{215.679}{11.955}{19.26}{\longrightarrow}%
\TextAt{13}{226.246}{215.679}{I}%
\TextAt{27}{231.408}{213.885}{P}%
\NativeArrowAt{241.823}{215.679}{11.955}{19.26}{\longrightarrow}%
\TextAt{13}{264.405}{215.679}{G}%
\TextAt{27}{273.639}{213.885}{P}%
\NativeArrowAt{284.053}{215.679}{11.955}{19.26}{\longrightarrow}%
\TextAt{13}{306.635}{215.679}{G}%
\TextAt{9}{315.869}{215.679}{\Character{40}}%
\TextAt{13}{320.421}{215.679}{\Character{20}}%
\TextAt{9}{327.159}{215.679}{\Character{40}}%
\TextAt{14}{331.712}{215.679}{P}%
\TextAt{9}{341.605}{215.679}{\Character{41}}%
\TextAt{16}{349.478}{215.679}{\Character{106}}%
\TextAt{13}{356.12}{215.679}{\Character{20}}%
\TextAt{9}{362.858}{215.679}{\Character{40}}%
\TextAt{14}{367.411}{215.679}{p}%
\TextAt{9}{373.394}{215.679}{\Character{41}\Character{41}}%
\NativeArrowAt{385.819}{215.679}{11.955}{19.26}{\longrightarrow}%
\TextAt{9}{408.401}{215.679}{1}%
\TextAt{13}{414.254}{215.679}{\Character{58}}%
\TextAt{8}{79.226}{192.024}{Its}%
\TextAt{8}{96.513}{192.024}{properties}%
\TextAt{8}{151.531}{192.024}{are}%
\TextAt{8}{171.042}{192.024}{expressed}%
\TextAt{8}{223.513}{192.024}{in}%
\TextAt{8}{237.166}{192.024}{the}%
\TextAt{12}{79.226}{167.539}{Proposition}%
\TextAt{12}{151.937}{167.539}{(9.6).}%
\TextAt{8}{187.253}{167.539}{The}%
\TextAt{8}{211.307}{167.539}{extension}%
\TextAt{13}{263.073}{167.539}{T}%
\TextAt{27}{269.934}{165.746}{P}%
\TextAt{16}{280.348}{167.539}{\Character{106}}%
\TextAt{13}{286.99}{167.539}{Z}%
\TextAt{27}{295.002}{165.746}{P}%
\TextAt{8}{305.993}{167.539}{is}%
\TextAt{8}{317.757}{167.539}{normal,}%
\TextAt{8}{360.688}{167.539}{and}%
\TextAt{8}{383.451}{167.539}{one}%
\TextAt{8}{404.91}{167.539}{has}%
\TextAt{13}{176.101}{143.885}{G}%
\TextAt{9}{185.335}{143.885}{\Character{40}}%
\TextAt{13}{189.887}{143.885}{T}%
\TextAt{27}{196.748}{142.092}{P}%
\TextAt{16}{207.162}{143.885}{\Character{106}}%
\TextAt{13}{213.804}{143.885}{Z}%
\TextAt{27}{221.816}{142.092}{P}%
\TextAt{9}{228.91}{143.885}{\Character{41}}%
\NativeMathAt{236.783}{143.885}{11.955}{\cong}%
\TextAt{13}{249.402}{143.885}{G}%
\TextAt{9}{258.636}{143.885}{\Character{40}}%
\TextAt{13}{263.188}{143.885}{\Character{20}}%
\TextAt{9}{269.927}{143.885}{\Character{40}}%
\TextAt{14}{274.479}{143.885}{P}%
\TextAt{9}{284.372}{143.885}{\Character{41}}%
\TextAt{16}{292.246}{143.885}{\Character{106}}%
\TextAt{13}{298.887}{143.885}{\Character{20}}%
\TextAt{9}{305.626}{143.885}{\Character{40}}%
\TextAt{14}{310.178}{143.885}{p}%
\TextAt{9}{316.161}{143.885}{\Character{41}\Character{41}}%
\TextAt{13}{325.266}{143.885}{\Character{59}}%
\TextAt{13}{354.401}{143.885}{G}%
\TextAt{9}{363.654}{143.885}{\Character{40}}%
\TextAt{13}{368.206}{143.885}{L}%
\TextAt{16}{379.492}{143.885}{\Character{106}}%
\TextAt{13}{386.133}{143.885}{T}%
\TextAt{27}{392.994}{142.092}{P}%
\TextAt{9}{400.088}{143.885}{\Character{41}}%
\TextAt{9}{407.943}{143.885}{\Character{61}}%
\TextAt{13}{420.387}{143.885}{I}%
\TextAt{27}{425.549}{142.092}{P}%
\TextAt{13}{432.643}{143.885}{\Character{58}}%
\TextAt{8}{79.226}{120.23}{If}%
\TextAt{8}{90.918}{120.23}{the}%
\TextAt{8}{111.075}{120.23}{residue}%
\TextAt{8}{150.802}{120.23}{field}%
\TextAt{8}{176.159}{120.23}{extension}%
\TextAt{13}{227.925}{120.23}{\Character{20}}%
\TextAt{9}{234.663}{120.23}{\Character{40}}%
\TextAt{14}{239.215}{120.23}{P}%
\TextAt{9}{249.109}{120.23}{\Character{41}}%
\TextAt{16}{256.982}{120.23}{\Character{106}}%
\TextAt{13}{263.624}{120.23}{\Character{20}}%
\TextAt{9}{270.362}{120.23}{\Character{40}}%
\TextAt{14}{274.914}{120.23}{p}%
\TextAt{9}{280.898}{120.23}{\Character{41}}%
\TextAt{8}{289.347}{120.23}{is}%
\TextAt{8}{301.111}{120.23}{separable,}%
\TextAt{8}{355.806}{120.23}{then}%
\TextAt{8}{382.466}{120.23}{one}%
\TextAt{8}{403.926}{120.23}{has}%
\TextAt{8}{175.708}{96.576}{\char35\relax }%
\TextAt{13}{185.464}{96.576}{I}%
\TextAt{27}{190.626}{94.783}{P}%
\TextAt{9}{201.041}{96.576}{\Character{61}}%
\TextAt{9}{213.451}{96.576}{\Character{91}}%
\TextAt{13}{216.718}{96.576}{L}%
\TextAt{9}{228.003}{96.576}{\Character{58}}%
\TextAt{13}{234.576}{96.576}{T}%
\TextAt{27}{241.437}{94.783}{P}%
\TextAt{9}{248.53}{96.576}{\Character{93}}%
\TextAt{9}{255.081}{96.576}{\Character{61}}%
\TextAt{13}{267.528}{96.576}{e\Character{59}}%
\TextAt{9}{302.108}{96.576}{\Character{40}}%
\TextAt{13}{306.661}{96.576}{G}%
\TextAt{27}{315.894}{94.783}{P}%
\TextAt{9}{326.309}{96.576}{\Character{58}}%
\TextAt{13}{332.881}{96.576}{I}%
\TextAt{27}{338.044}{94.783}{P}%
\TextAt{9}{345.138}{96.576}{\Character{41}}%
\TextAt{9}{352.993}{96.576}{\Character{61}}%
\TextAt{9}{365.414}{96.576}{\Character{91}}%
\TextAt{13}{368.688}{96.576}{T}%
\TextAt{27}{375.549}{94.783}{P}%
\TextAt{9}{385.963}{96.576}{\Character{58}}%
\TextAt{13}{392.536}{96.576}{Z}%
\TextAt{27}{400.548}{94.783}{P}%
\TextAt{9}{407.642}{96.576}{\Character{93}}%
\TextAt{9}{414.193}{96.576}{\Character{61}}%
\TextAt{13}{426.64}{96.576}{f}%
\TextAt{13}{433.024}{96.576}{\Character{58}}%
\TextAt{8}{79.226}{72.921}{In}%
\TextAt{8}{93.847}{72.921}{this}%
\TextAt{8}{116.67}{72.921}{case}%
\TextAt{8}{141.441}{72.921}{one}%
\TextAt{8}{162.9}{72.921}{finds}%
\TextAt{8}{190.923}{72.921}{for}%
\TextAt{8}{208.808}{72.921}{the}%
\TextAt{8}{228.965}{72.921}{prime}%
\TextAt{8}{262.128}{72.921}{ideal}%
\TextAt{14}{290.091}{72.921}{P}%
\TextAt{21}{299.985}{71.128}{T}%
\TextAt{8}{310.487}{72.921}{of}%
\TextAt{13}{323.817}{72.921}{T}%
\TextAt{27}{330.678}{71.128}{P}%
\TextAt{8}{341.669}{72.921}{below}%
\TextAt{14}{374.832}{72.921}{P}%
\TextAt{8}{384.725}{72.921}{:}%
\TextAt{8}{300.14}{43.033}{53}%
\Destination{78.23}{809.815}{nav.page.53}%
\Destination{79.36}{751.655}{nav.print.57}%
\Destination{79.36}{750.655}{nav.number.1.9.4}%
\Destination{79.36}{419.755}{nav.number.1.9.5}%
\Destination{79.36}{181.555}{nav.number.1.9.6}%
}

% Page 66 (body page 54)
\Page{%
\TextAt{2}{79.226}{796.943}{I.}%
\TextAt{2}{89.83}{796.943}{Algebraic}%
\TextAt{2}{138.648}{796.943}{Integers}%
\Vector{q 0 G 0 g 0.299 w 7.226 18.729 m 460.769 18.729 l S Q}%
\TextAt{8}{89.579}{761.124}{(i)}%
\TextAt{8}{107.918}{761.124}{The}%
\TextAt{8}{131.972}{761.124}{ramification}%
\TextAt{8}{197.02}{761.124}{index}%
\TextAt{8}{228.558}{761.124}{of}%
\TextAt{14}{241.888}{761.124}{P}%
\TextAt{8}{255.679}{761.124}{over}%
\TextAt{14}{280.725}{761.124}{P}%
\TextAt{21}{290.618}{759.33}{T}%
\TextAt{8}{301.12}{761.124}{is}%
\TextAt{13}{312.884}{761.124}{e}%
\TextAt{8}{322.207}{761.124}{and}%
\TextAt{8}{344.969}{761.124}{the}%
\TextAt{8}{365.126}{761.124}{inertia}%
\TextAt{8}{402.199}{761.124}{degree}%
\TextAt{8}{438.614}{761.124}{is}%
\TextAt{9}{450.378}{761.124}{1}%
\TextAt{8}{456.231}{761.124}{.}%
\TextAt{8}{86.327}{744.078}{(ii)}%
\TextAt{8}{107.918}{744.078}{The}%
\TextAt{8}{131.972}{744.078}{ramification}%
\TextAt{8}{197.02}{744.078}{index}%
\TextAt{8}{228.558}{744.078}{of}%
\TextAt{14}{241.888}{744.078}{P}%
\TextAt{21}{251.781}{742.284}{T}%
\TextAt{8}{262.283}{744.078}{over}%
\TextAt{14}{287.329}{744.078}{P}%
\TextAt{21}{297.223}{742.284}{Z}%
\TextAt{8}{307.919}{744.078}{is}%
\TextAt{9}{319.683}{744.078}{1}%
\TextAt{8}{325.536}{744.078}{,}%
\TextAt{8}{332.686}{744.078}{and}%
\TextAt{8}{355.448}{744.078}{the}%
\TextAt{8}{375.605}{744.078}{inertia}%
\TextAt{8}{412.678}{744.078}{degree}%
\TextAt{8}{449.093}{744.078}{is}%
\TextAt{13}{460.857}{744.078}{f}%
\TextAt{8}{467.903}{744.078}{.}%
\TextAt{19}{79.226}{709.099}{Proof.}%
\TextAt{8}{116}{709.099}{The}%
\TextAt{8}{139.862}{709.099}{first}%
\TextAt{8}{163.795}{709.099}{two}%
\TextAt{8}{185.719}{709.099}{claims}%
\TextAt{8}{221.356}{709.099}{follow}%
\TextAt{8}{254.985}{709.099}{from}%
\TextAt{8}{282.433}{709.099}{the}%
\TextAt{8}{302.397}{709.099}{identity}%
\TextAt{8}{345.458}{709.099}{\char35\relax }%
\TextAt{13}{355.213}{709.099}{G}%
\TextAt{27}{364.447}{707.306}{P}%
\TextAt{9}{374.861}{709.099}{\Character{61}}%
\TextAt{13}{387.287}{709.099}{ef}%
\TextAt{8}{399.759}{709.099}{.}%
\TextAt{8}{408.143}{709.099}{So}%
\TextAt{8}{424.209}{709.099}{we}%
\TextAt{8}{441.244}{709.099}{only}%
\TextAt{8}{466.744}{709.099}{have}%
\TextAt{8}{493.546}{709.099}{to}%
\TextAt{8}{507.664}{709.099}{show}%
\TextAt{8}{79.226}{692.053}{statements}%
\TextAt{8}{139.242}{692.053}{(i)}%
\TextAt{8}{156.531}{692.053}{and}%
\TextAt{8}{180.323}{692.053}{(ii).}%
\TextAt{8}{207.473}{692.053}{Using}%
\TextAt{8}{241.403}{692.053}{the}%
\TextAt{8}{262.589}{692.053}{fundamental}%
\TextAt{8}{331.261}{692.053}{identity,}%
\TextAt{8}{378.073}{692.053}{they}%
\TextAt{8}{405.44}{692.053}{all}%
\TextAt{8}{422.728}{692.053}{follow}%
\TextAt{8}{457.579}{692.053}{from}%
\TextAt{13}{486.249}{692.053}{\Character{20}}%
\TextAt{9}{492.987}{692.053}{\Character{40}}%
\TextAt{14}{497.54}{692.053}{P}%
\TextAt{21}{507.433}{690.26}{T}%
\TextAt{9}{514.037}{692.053}{\Character{41}}%
\TextAt{9}{523.649}{692.053}{\Character{61}}%
\TextAt{13}{79.226}{675.007}{\Character{20}}%
\TextAt{9}{85.965}{675.007}{\Character{40}}%
\TextAt{14}{90.517}{675.007}{P}%
\TextAt{9}{100.41}{675.007}{\Character{41}}%
\TextAt{8}{104.963}{675.007}{.}%
\TextAt{8}{113.123}{675.007}{As}%
\TextAt{8}{129.546}{675.007}{the}%
\TextAt{8}{148.839}{675.007}{inertia}%
\TextAt{8}{185.048}{675.007}{group}%
\TextAt{13}{217.359}{675.007}{I}%
\TextAt{27}{222.522}{673.214}{P}%
\TextAt{8}{232.649}{675.007}{of}%
\TextAt{14}{245.115}{675.007}{P}%
\TextAt{8}{258.042}{675.007}{over}%
\TextAt{13}{282.224}{675.007}{K}%
\TextAt{8}{296.075}{675.007}{is}%
\TextAt{8}{306.975}{675.007}{also}%
\TextAt{8}{329.591}{675.007}{the}%
\TextAt{8}{348.883}{675.007}{inertia}%
\TextAt{8}{385.092}{675.007}{group}%
\TextAt{8}{417.404}{675.007}{of}%
\TextAt{14}{429.87}{675.007}{P}%
\TextAt{8}{442.796}{675.007}{over}%
\TextAt{13}{466.978}{675.007}{T}%
\TextAt{27}{473.839}{673.214}{P}%
\TextAt{8}{480.933}{675.007}{,}%
\TextAt{8}{487.391}{675.007}{it}%
\TextAt{8}{498.231}{675.007}{follows}%
\TextAt{8}{79.226}{657.961}{from}%
\TextAt{8}{105.749}{657.961}{an}%
\TextAt{8}{120.89}{657.961}{application}%
\TextAt{8}{180.265}{657.961}{of}%
\TextAt{8}{192.477}{657.961}{proposition}%
\TextAt{8}{252.892}{657.961}{(9.4)}%
\TextAt{8}{279.749}{657.961}{to}%
\TextAt{8}{292.942}{657.961}{the}%
\TextAt{8}{311.98}{657.961}{extension}%
\TextAt{13}{362.628}{657.961}{L}%
\TextAt{16}{373.913}{657.961}{\Character{106}}%
\TextAt{13}{380.555}{657.961}{T}%
\TextAt{27}{387.416}{656.168}{P}%
\TextAt{8}{397.289}{657.961}{that}%
\TextAt{13}{421.54}{657.961}{G}%
\TextAt{9}{430.774}{657.961}{\Character{40}}%
\TextAt{13}{435.326}{657.961}{\Character{20}}%
\TextAt{9}{442.065}{657.961}{\Character{40}}%
\TextAt{14}{446.617}{657.961}{P}%
\TextAt{9}{456.51}{657.961}{\Character{41}}%
\TextAt{16}{464.383}{657.961}{\Character{106}}%
\TextAt{13}{471.025}{657.961}{\Character{20}}%
\TextAt{9}{477.764}{657.961}{\Character{40}}%
\TextAt{14}{482.316}{657.961}{P}%
\TextAt{21}{492.209}{656.168}{T}%
\TextAt{9}{498.814}{657.961}{\Character{41}\Character{41}}%
\TextAt{9}{511.212}{657.961}{\Character{61}}%
\TextAt{9}{523.633}{657.961}{1}%
\TextAt{8}{529.518}{657.961}{,}%
\TextAt{8}{79.226}{640.915}{hence}%
\TextAt{13}{111.732}{640.915}{\Character{20}}%
\TextAt{9}{118.471}{640.915}{\Character{40}}%
\TextAt{14}{123.023}{640.915}{P}%
\TextAt{21}{132.916}{639.122}{T}%
\TextAt{9}{139.521}{640.915}{\Character{41}}%
\TextAt{9}{147.376}{640.915}{\Character{61}}%
\TextAt{13}{159.819}{640.915}{\Character{20}}%
\TextAt{9}{166.558}{640.915}{\Character{40}}%
\TextAt{14}{171.11}{640.915}{P}%
\TextAt{9}{181.004}{640.915}{\Character{41}}%
\TextAt{8}{185.556}{640.915}{.}%
\Vector{q 0 G 0 g 0.399 w 452.334 -133.085 m 452.334 -125.015 l S Q}%
\Vector{q 0 G 0 g 0.399 w 452.534 -125.214 m 459.707 -125.214 l S Q}%
\Vector{q 0 G 0 g 0.399 w 452.534 -132.886 m 459.707 -132.886 l S Q}%
\Vector{q 0 G 0 g 0.399 w 459.906 -133.085 m 459.906 -125.015 l S Q}%
\TextAt{8}{97.159}{611.914}{In}%
\TextAt{8}{111.78}{611.914}{the}%
\TextAt{8}{131.937}{611.914}{diagram}%
\TextAt{13}{202.059}{582.912}{K}%
\TextAt{13}{266.221}{582.912}{Z}%
\TextAt{27}{274.233}{581.119}{P}%
\TextAt{13}{334.671}{582.912}{T}%
\TextAt{27}{341.532}{581.119}{P}%
\TextAt{13}{401.97}{582.912}{L}%
\NativeDiagram{%
\path[draw=black,line width=0.478bp] (146.229,-188.099) -- (188.392,-188.099);%
\path[draw=black,line width=0.478bp,line cap=round,line join=round] (186.144,-185.23) .. controls (186.599,-186.951) and (187.611,-187.764) .. (188.631,-188.099) .. controls (187.611,-188.434) and (186.599,-189.247) .. (186.144,-190.968);%
}%
\TextAt{17}{237.432}{588.717}{1}%
\TextAt{17}{237.432}{577.949}{1}%
\NativeDiagram{%
\path[draw=black,line width=0.478bp] (214.68,-188.099) -- (256.842,-188.099);%
\path[draw=black,line width=0.478bp,line cap=round,line join=round] (254.595,-185.23) .. controls (255.049,-186.951) and (256.062,-187.764) .. (257.082,-188.099) .. controls (256.062,-188.434) and (255.049,-189.247) .. (254.595,-190.968);%
}%
\TextAt{17}{305.883}{588.717}{1}%
\TextAt{21}{305.527}{577.551}{f}%
\NativeDiagram{%
\path[draw=black,line width=0.478bp] (281.979,-188.099) -- (324.142,-188.099);%
\path[draw=black,line width=0.478bp,line cap=round,line join=round] (321.894,-185.23) .. controls (322.349,-186.951) and (323.361,-187.764) .. (324.381,-188.099) .. controls (323.361,-188.434) and (322.349,-189.247) .. (321.894,-190.968);%
}%
\TextAt{21}{373.336}{588.717}{e}%
\TextAt{17}{373.183}{577.949}{1}%
\TextAt{8}{79.226}{552.057}{we}%
\TextAt{8}{96.876}{552.057}{have}%
\TextAt{8}{124.293}{552.057}{indicated}%
\TextAt{8}{175.441}{552.057}{the}%
\TextAt{8}{196.02}{552.057}{ramification}%
\TextAt{8}{261.49}{552.057}{indices}%
\TextAt{8}{300.336}{552.057}{of}%
\TextAt{8}{314.088}{552.057}{the}%
\TextAt{8}{334.667}{552.057}{individual}%
\TextAt{8}{390.047}{552.057}{field}%
\TextAt{8}{415.826}{552.057}{extensions}%
\TextAt{8}{472.629}{552.057}{on}%
\TextAt{8}{489.311}{552.057}{top,}%
\TextAt{8}{513.904}{552.057}{and}%
\TextAt{8}{79.226}{535.011}{the}%
\TextAt{8}{99.91}{535.011}{inertia}%
\TextAt{8}{137.511}{535.011}{degrees}%
\TextAt{8}{179.069}{535.011}{on}%
\TextAt{8}{195.856}{535.011}{the}%
\TextAt{8}{216.54}{535.011}{bottom.}%
\TextAt{8}{263.979}{535.011}{In}%
\TextAt{8}{279.128}{535.011}{the}%
\TextAt{8}{299.813}{535.011}{special}%
\TextAt{8}{338.441}{535.011}{case}%
\TextAt{8}{363.74}{535.011}{where}%
\TextAt{8}{398.077}{535.011}{the}%
\TextAt{8}{418.761}{535.011}{residue}%
\TextAt{8}{459.016}{535.011}{field}%
\TextAt{8}{484.901}{535.011}{extension}%
\TextAt{13}{79.226}{517.965}{\Character{20}}%
\TextAt{9}{85.965}{517.965}{\Character{40}}%
\TextAt{14}{90.517}{517.965}{P}%
\TextAt{9}{100.41}{517.965}{\Character{41}}%
\TextAt{16}{108.283}{517.965}{\Character{106}}%
\TextAt{13}{114.925}{517.965}{\Character{20}}%
\TextAt{9}{121.664}{517.965}{\Character{40}}%
\TextAt{14}{126.216}{517.965}{p}%
\TextAt{9}{132.199}{517.965}{\Character{41}}%
\TextAt{8}{140.649}{517.965}{is}%
\TextAt{8}{152.413}{517.965}{separable}%
\TextAt{8}{203.856}{517.965}{we}%
\TextAt{8}{221.083}{517.965}{find}%
\TextAt{13}{165.244}{488.964}{I}%
\TextAt{27}{170.407}{487.171}{P}%
\TextAt{9}{180.822}{488.964}{\Character{61}}%
\TextAt{9}{193.232}{488.964}{1}%
\NativeArrowAt{214.376}{488.964}{11.955}{21.926}{\Longleftrightarrow}%
\TextAt{13}{251.57}{488.964}{T}%
\TextAt{27}{258.431}{487.171}{P}%
\TextAt{9}{268.846}{488.964}{\Character{61}}%
\TextAt{13}{281.271}{488.964}{L}%
\NativeArrowAt{304.512}{488.964}{11.955}{21.926}{\Longleftrightarrow}%
\TextAt{14}{341.705}{488.964}{p}%
\TextAt{8}{351.586}{488.964}{is}%
\TextAt{8}{363.35}{488.964}{unramified}%
\TextAt{8}{421.882}{488.964}{in}%
\TextAt{13}{435.535}{488.964}{L\Character{58}}%
\ErrataAt{79.19}{459.997}{455.81}{When $\mathfrak p$ is unramified in $L$, we have $G(\kappa(\mathfrak P)\mid\kappa(\mathfrak p))\cong G_{\mathfrak P}$, so this residue Galois group}%
\TextAt{8}{79.226}{442.917}{may}%
\TextAt{8}{104.595}{442.917}{be}%
\TextAt{8}{120.519}{442.917}{viewed}%
\TextAt{8}{158.883}{442.917}{as}%
\TextAt{8}{173.254}{442.917}{a}%
\TextAt{8}{183.009}{442.917}{subgroup}%
\TextAt{8}{233.807}{442.917}{of}%
\TextAt{13}{247.137}{442.917}{G}%
\TextAt{9}{259.691}{442.917}{\Character{61}}%
\TextAt{13}{272.116}{442.917}{G}%
\TextAt{9}{281.35}{442.917}{\Character{40}}%
\TextAt{13}{285.902}{442.917}{L}%
\TextAt{16}{297.188}{442.917}{\Character{106}}%
\TextAt{13}{303.829}{442.917}{K}%
\TextAt{9}{314.647}{442.917}{\Character{41}}%
\TextAt{8}{319.199}{442.917}{.}%
\TextAt{8}{97.159}{422.882}{Hilbert’s}%
\TextAt{8}{144.71}{422.882}{ramification}%
\TextAt{8}{209.13}{422.882}{theory,}%
\TextAt{8}{247.649}{422.882}{with}%
\TextAt{8}{273.68}{422.882}{its}%
\TextAt{8}{289.371}{422.882}{various}%
\TextAt{8}{328.816}{422.882}{refinements}%
\TextAt{8}{390.355}{422.882}{and}%
\TextAt{8}{412.489}{422.882}{generalizations,}%
\TextAt{8}{494.656}{422.882}{belongs}%
\TextAt{8}{79.226}{405.837}{naturally}%
\TextAt{8}{131.247}{405.837}{to}%
\TextAt{8}{147.163}{405.837}{the}%
\TextAt{8}{168.925}{405.837}{theory}%
\TextAt{8}{207.281}{405.837}{of}%
\TextAt{8}{222.217}{405.837}{valuations,}%
\TextAt{8}{283.163}{405.837}{which}%
\TextAt{8}{318.255}{405.837}{we}%
\TextAt{8}{337.088}{405.837}{will}%
\TextAt{8}{360.799}{405.837}{develop}%
\TextAt{8}{404.679}{405.837}{in}%
\TextAt{8}{419.937}{405.837}{the}%
\TextAt{8}{441.699}{405.837}{next}%
\TextAt{8}{469.642}{405.837}{chapter}%
\TextAt{8}{513.199}{405.837}{(see}%
\TextAt{8}{79.226}{388.791}{chap.}%
\TextAt{8}{110.118}{388.791}{II,}%
\TextAt{8}{126.031}{388.791}{§9).}%
\TextAt{12}{79.226}{362.778}{Exercise}%
\TextAt{12}{132.71}{362.778}{I.9.1.}%
\TextAt{8}{166.549}{362.778}{If}%
\TextAt{13}{178.488}{362.778}{L}%
\TextAt{16}{190.192}{362.778}{\Character{106}}%
\TextAt{13}{197.253}{362.778}{K}%
\TextAt{8}{212.214}{362.778}{is}%
\TextAt{8}{224.224}{362.778}{a}%
\TextAt{8}{234.225}{362.778}{Galois}%
\TextAt{8}{270.384}{362.778}{extension}%
\TextAt{8}{322.396}{362.778}{of}%
\TextAt{8}{335.972}{362.778}{algebraic}%
\TextAt{8}{385.653}{362.778}{number}%
\TextAt{8}{428.495}{362.778}{fields}%
\TextAt{8}{458.712}{362.778}{with}%
\TextAt{8}{485.618}{362.778}{noncyclic}%
\TextAt{8}{79.226}{345.732}{Galois}%
\TextAt{8}{115.139}{345.732}{group,}%
\TextAt{8}{151.567}{345.732}{then}%
\TextAt{8}{178.227}{345.732}{there}%
\TextAt{8}{208.139}{345.732}{are}%
\TextAt{8}{227.65}{345.732}{at}%
\TextAt{8}{241.96}{345.732}{most}%
\TextAt{8}{270.64}{345.732}{finitely}%
\TextAt{8}{309.985}{345.732}{many}%
\TextAt{8}{341.857}{345.732}{nonsplit}%
\TextAt{8}{386.797}{345.732}{prime}%
\TextAt{8}{419.96}{345.732}{ideals}%
\TextAt{8}{452.538}{345.732}{of}%
\TextAt{13}{465.868}{345.732}{K}%
\TextAt{8}{476.686}{345.732}{.}%
\TextAt{12}{79.226}{313.742}{Exercise}%
\TextAt{12}{132.354}{313.742}{I.9.2.}%
\TextAt{8}{165.885}{313.742}{If}%
\TextAt{13}{177.514}{313.742}{L}%
\TextAt{16}{188.799}{313.742}{\Character{106}}%
\TextAt{13}{195.441}{313.742}{K}%
\TextAt{8}{210.093}{313.742}{is}%
\TextAt{8}{221.794}{313.742}{a}%
\TextAt{8}{231.486}{313.742}{Galois}%
\TextAt{8}{267.336}{313.742}{extension}%
\TextAt{8}{319.039}{313.742}{of}%
\TextAt{8}{332.306}{313.742}{algebraic}%
\TextAt{8}{381.678}{313.742}{number}%
\TextAt{8}{424.211}{313.742}{fields,}%
\TextAt{8}{457.384}{313.742}{and}%
\TextAt{14}{480.083}{313.742}{P}%
\TextAt{8}{493.811}{313.742}{a}%
\TextAt{8}{503.503}{313.742}{prime}%
\TextAt{8}{79.226}{296.696}{ideal}%
\TextAt{8}{107.497}{296.696}{which}%
\TextAt{8}{141.291}{296.696}{is}%
\TextAt{8}{153.363}{296.696}{unramified}%
\TextAt{8}{212.203}{296.696}{over}%
\TextAt{13}{237.557}{296.696}{K}%
\TextAt{8}{252.58}{296.696}{(i.e.,}%
\TextAt{14}{279.625}{296.696}{p}%
\TextAt{9}{289.453}{296.696}{\Character{61}}%
\TextAt{14}{302.403}{296.696}{P}%
\TextAt{16}{315.163}{296.696}{\Character{92}}%
\TextAt{13}{325.999}{296.696}{K}%
\TextAt{8}{341.022}{296.696}{is}%
\TextAt{8}{353.094}{296.696}{unramified}%
\TextAt{8}{411.934}{296.696}{in}%
\TextAt{13}{425.895}{296.696}{L}%
\TextAt{8}{433.859}{296.696}{),}%
\TextAt{8}{445.948}{296.696}{then}%
\TextAt{8}{472.916}{296.696}{there}%
\TextAt{8}{503.136}{296.696}{is}%
\TextAt{8}{515.207}{296.696}{one}%
\TextAt{8}{79.226}{279.65}{and}%
\TextAt{8}{101.989}{279.65}{only}%
\TextAt{8}{127.68}{279.65}{one}%
\TextAt{8}{149.14}{279.65}{automorphism}%
\TextAt{13}{226.609}{279.65}{\Character{39}}%
\TextAt{27}{234.278}{277.857}{P}%
\NativeMathAt{244.692}{279.65}{11.955}{\in}%
\TextAt{13}{255.983}{279.65}{G}%
\TextAt{9}{265.217}{279.65}{\Character{40}}%
\TextAt{13}{269.769}{279.65}{L}%
\TextAt{16}{281.055}{279.65}{\Character{106}}%
\TextAt{13}{287.696}{279.65}{K}%
\TextAt{9}{298.514}{279.65}{\Character{41}}%
\TextAt{8}{306.964}{279.65}{such}%
\TextAt{8}{333.361}{279.65}{that}%
\TextAt{13}{213.713}{250.649}{\Character{39}}%
\TextAt{27}{221.382}{248.856}{P}%
\TextAt{13}{228.475}{250.649}{a}%
\NativeMathAt{237.941}{250.649}{11.955}{\equiv}%
\TextAt{13}{250.56}{250.649}{a}%
\TextAt{21}{256.705}{255.585}{q}%
\TextAt{9}{273.225}{250.649}{\Character{40}}%
\TextAt{8}{277.777}{250.649}{mod}%
\TextAt{14}{303.88}{250.649}{P}%
\TextAt{9}{313.773}{250.649}{\Character{41}}%
\TextAt{8}{330.28}{250.649}{for}%
\TextAt{8}{348.165}{250.649}{all}%
\TextAt{13}{364.424}{250.649}{a}%
\NativeMathAt{373.89}{250.649}{11.955}{\in}%
\TextAt{16}{385.164}{250.649}{\Character{79}}%
\TextAt{13}{395.031}{250.649}{\Character{59}}%
\EditionChange{\BodyAt{79.226}{221.655}{where $q=[\kappa(\mathfrak P):\kappa(\mathfrak p)]$.}}{\BodyAt{79.226}{221.655}{where $q=|\kappa(\mathfrak p)|$.}}%
\TextAt{8}{206.096}{221.648}{It}%
\TextAt{8}{218.234}{221.648}{is}%
\TextAt{8}{229.464}{221.648}{called}%
\TextAt{8}{262.093}{221.648}{the}%
\TextAt{12}{281.715}{221.648}{Frobenius}%
\TextAt{12}{343.112}{221.648}{automorphism}%
\TextAt{8}{427.742}{221.648}{.}%
\TextAt{8}{436.013}{221.648}{The}%
\TextAt{8}{459.532}{221.648}{decomposition}%
\TextAt{8}{79.226}{204.602}{group}%
\TextAt{13}{112.402}{204.602}{G}%
\TextAt{27}{121.635}{202.809}{P}%
\TextAt{8}{132.626}{204.602}{is}%
\TextAt{8}{144.39}{204.602}{cyclic}%
\TextAt{8}{176.574}{204.602}{and}%
\TextAt{13}{199.336}{204.602}{\Character{39}}%
\TextAt{27}{207.005}{202.809}{P}%
\TextAt{8}{217.996}{204.602}{is}%
\TextAt{8}{229.759}{204.602}{a}%
\TextAt{8}{239.515}{204.602}{generator}%
\TextAt{8}{291.556}{204.602}{of}%
\TextAt{13}{304.886}{204.602}{G}%
\TextAt{27}{314.119}{202.809}{P}%
\TextAt{8}{321.213}{204.602}{.}%
\TextAt{12}{79.226}{172.612}{Exercise}%
\TextAt{12}{133.763}{172.612}{I.9.3.}%
\TextAt{8}{168.519}{172.612}{Let}%
\TextAt{13}{190.65}{172.612}{L}%
\TextAt{16}{203.915}{172.612}{\Character{106}}%
\TextAt{13}{212.537}{172.612}{K}%
\TextAt{8}{228.414}{172.612}{be}%
\TextAt{8}{245.5}{172.612}{a}%
\TextAt{8}{256.417}{172.612}{solvable}%
\TextAt{8}{301.55}{172.612}{extension}%
\TextAt{8}{354.478}{172.612}{of}%
\TextAt{8}{368.97}{172.612}{prime}%
\TextAt{8}{403.295}{172.612}{degree}%
\TextAt{13}{440.873}{172.612}{p}%
\TextAt{8}{451.807}{172.612}{(not}%
\TextAt{8}{478.338}{172.612}{necessarily}%
\ErrataAt{79.19}{155.551}{455.81}{Galois). If a prime $\mathfrak p$ of $K$ unramified in $L$ has two degree-one factors $\mathfrak P$ and $\mathfrak P^{\prime}$,}%
\TextAt{8}{79.226}{138.52}{then}%
\TextAt{8}{105.886}{138.52}{it}%
\TextAt{8}{117.59}{138.52}{is}%
\TextAt{8}{129.354}{138.52}{already}%
\TextAt{8}{170.659}{138.52}{totally}%
\TextAt{8}{208.067}{138.52}{split}%
\TextAt{8}{234.141}{138.52}{(theorem}%
\TextAt{8}{284.221}{138.52}{of}%
\TextAt{26}{298.066}{138.52}{F.K.}%
\TextAt{26}{328.814}{138.52}{Schmidt}%
\TextAt{8}{377.579}{138.52}{).}%
\TextAt{12}{97.159}{118.485}{Hint:}%
\TextAt{8}{133.091}{118.485}{Use}%
\TextAt{8}{155.775}{118.485}{the}%
\TextAt{8}{176.129}{118.485}{following}%
\TextAt{8}{225.76}{118.485}{result}%
\TextAt{8}{258.535}{118.485}{of}%
\TextAt{26}{272.575}{118.485}{Galois}%
\TextAt{8}{316.994}{118.485}{(see}%
\TextAt{8}{340.659}{118.485}{[75],}%
\TextAt{8}{366.274}{118.485}{chap.}%
\TextAt{8}{397.363}{118.485}{II,}%
\TextAt{8}{413.472}{118.485}{§3):}%
\TextAt{8}{438.405}{118.485}{if}%
\TextAt{13}{449.326}{118.485}{G}%
\TextAt{8}{462.653}{118.485}{is}%
\TextAt{8}{474.614}{118.485}{a}%
\TextAt{8}{484.566}{118.485}{transitive}%
\TextAt{8}{79.226}{101.439}{solvable}%
\TextAt{8}{124.212}{101.439}{permutation}%
\TextAt{8}{192.223}{101.439}{group}%
\TextAt{8}{226.413}{101.439}{of}%
\TextAt{8}{240.758}{101.439}{prime}%
\TextAt{8}{274.936}{101.439}{degree}%
\TextAt{13}{312.366}{101.439}{p}%
\TextAt{8}{318.241}{101.439}{,}%
\TextAt{8}{326.659}{101.439}{then}%
\TextAt{8}{354.333}{101.439}{there}%
\TextAt{8}{385.26}{101.439}{is}%
\TextAt{8}{398.038}{101.439}{no}%
\TextAt{8}{415.312}{101.439}{nontrivial}%
\TextAt{8}{469.67}{101.439}{permutation}%
\TextAt{13}{79.226}{84.394}{\Character{27}}%
\NativeMathAt{89.629}{84.394}{11.955}{\in}%
\TextAt{13}{100.92}{84.394}{G}%
\TextAt{8}{114.051}{84.394}{which}%
\TextAt{8}{147.538}{84.394}{fixes}%
\TextAt{8}{173.935}{84.394}{two}%
\TextAt{8}{196.052}{84.394}{distinct}%
\TextAt{8}{238.385}{84.394}{letters.}%
\TextAt{8}{300.14}{43.033}{54}%
\Destination{79.36}{240.655}{errata.3}%
\Destination{79.23}{773.075}{nav.Item.10}%
\Destination{79.23}{758.125}{nav.Item.11}%
\Destination{78.23}{809.815}{nav.page.54}%
\Destination{79.226}{376.755}{nav.ex.1.9.1}%
\Destination{79.226}{327.755}{nav.ex.1.9.2}%
\Destination{79.226}{186.655}{nav.ex.1.9.3}%
\Destination{79.36}{567.055}{nav.print.58}%
\Hotlink{109.618}{386.555}{9.763}{11.047}{nav.chapter.2}%
\Hotlink{131.21}{386.555}{6.858}{11.047}{nav.section.2.9}%
%
%
\Hotlink{340.159}{116.255}{19.22}{11.047}{nav.cite.N075}%
\Hotlink{252.392}{655.755}{25.078}{11.047}{nav.number.1.9.4}%
\Destination{79.226}{473}{errata.182}%
}

% Page 67 (body page 55)
\Page{%
\TextAt{2}{79.226}{796.943}{I.}%
\TextAt{2}{89.83}{796.943}{Algebraic}%
\TextAt{2}{138.648}{796.943}{Integers}%
\Vector{q 0 G 0 g 0.299 w 7.226 18.729 m 460.769 18.729 l S Q}%
\TextAt{12}{79.226}{761.124}{Exercise}%
\TextAt{12}{134.108}{761.124}{I.9.4.}%
\TextAt{8}{169.164}{761.124}{Let}%
\TextAt{13}{191.596}{761.124}{L}%
\TextAt{16}{205.372}{761.124}{\Character{106}}%
\TextAt{13}{214.506}{761.124}{K}%
\TextAt{8}{230.683}{761.124}{be}%
\TextAt{8}{248.069}{761.124}{a}%
\TextAt{8}{259.286}{761.124}{finite}%
\TextAt{8}{290.66}{761.124}{(not}%
\TextAt{8}{317.491}{761.124}{necessarily}%
\TextAt{8}{377.282}{761.124}{Galois)}%
\TextAt{8}{419.212}{761.124}{extension}%
\TextAt{8}{472.44}{761.124}{of}%
\TextAt{8}{487.232}{761.124}{algebraic}%
\TextAt{8}{79.226}{744.078}{number}%
\TextAt{8}{121.843}{744.078}{fields}%
\TextAt{8}{151.836}{744.078}{and}%
\TextAt{13}{174.619}{744.078}{N}%
\TextAt{16}{188.598}{744.078}{\Character{106}}%
\TextAt{13}{195.275}{744.078}{K}%
\TextAt{8}{210.011}{744.078}{the}%
\TextAt{8}{230.188}{744.078}{normal}%
\TextAt{8}{269.888}{744.078}{closure}%
\TextAt{8}{308.991}{744.078}{of}%
\TextAt{13}{322.341}{744.078}{L}%
\TextAt{16}{333.662}{744.078}{\Character{106}}%
\TextAt{13}{340.339}{744.078}{K}%
\TextAt{8}{351.157}{744.078}{.}%
\TextAt{8}{359.668}{744.078}{Show}%
\TextAt{8}{390.581}{744.078}{that}%
\TextAt{8}{415.97}{744.078}{a}%
\TextAt{8}{425.747}{744.078}{prime}%
\TextAt{8}{458.931}{744.078}{ideal}%
\TextAt{14}{486.915}{744.078}{p}%
\TextAt{8}{496.816}{744.078}{of}%
\TextAt{13}{510.167}{744.078}{K}%
\TextAt{8}{524.903}{744.078}{is}%
\TextAt{8}{79.226}{727.032}{totally}%
\TextAt{8}{116.634}{727.032}{split}%
\TextAt{8}{142.708}{727.032}{in}%
\TextAt{13}{156.361}{727.032}{L}%
\TextAt{8}{168.223}{727.032}{if}%
\TextAt{8}{178.947}{727.032}{and}%
\TextAt{8}{201.709}{727.032}{only}%
\TextAt{8}{227.401}{727.032}{if}%
\TextAt{8}{238.125}{727.032}{it}%
\TextAt{8}{249.829}{727.032}{is}%
\TextAt{8}{261.593}{727.032}{totally}%
\TextAt{8}{299}{727.032}{split}%
\TextAt{8}{325.075}{727.032}{in}%
\TextAt{13}{338.727}{727.032}{N}%
\TextAt{8}{349.35}{727.032}{.}%
\TextAt{12}{97.159}{706.997}{Hint:}%
\TextAt{8}{134.733}{706.997}{Use}%
\TextAt{8}{158.238}{706.997}{the}%
\TextAt{8}{179.413}{706.997}{double}%
\TextAt{8}{218.149}{706.997}{coset}%
\TextAt{8}{248.493}{706.997}{decomposition}%
\TextAt{13}{326.645}{706.997}{H}%
\TextAt{16}{337.268}{706.997}{\Character{110}}%
\TextAt{13}{343.245}{706.997}{G}%
\TextAt{9}{352.479}{706.997}{\Character{47}}%
\TextAt{13}{358.332}{706.997}{G}%
\TextAt{27}{367.565}{705.204}{P}%
\TextAt{8}{374.659}{706.997}{,}%
\TextAt{8}{383.08}{706.997}{where}%
\TextAt{13}{417.907}{706.997}{G}%
\TextAt{9}{432.196}{706.997}{\Character{61}}%
\TextAt{13}{446.356}{706.997}{G}%
\TextAt{9}{455.59}{706.997}{\Character{40}}%
\TextAt{13}{460.142}{706.997}{N}%
\TextAt{16}{475.82}{706.997}{\Character{106}}%
\TextAt{13}{484.196}{706.997}{K}%
\TextAt{9}{495.013}{706.997}{\Character{41}}%
\TextAt{8}{499.566}{706.997}{,}%
\TextAt{13}{507.987}{706.997}{H}%
\TextAt{9}{523.665}{706.997}{\Character{61}}%
\TextAt{13}{79.226}{689.951}{G}%
\TextAt{9}{88.46}{689.951}{\Character{40}}%
\TextAt{13}{93.012}{689.951}{N}%
\TextAt{16}{106.956}{689.951}{\Character{106}}%
\TextAt{13}{113.597}{689.951}{L}%
\TextAt{9}{121.562}{689.951}{\Character{41}}%
\TextAt{8}{130.011}{689.951}{and}%
\TextAt{13}{152.774}{689.951}{G}%
\TextAt{27}{162.008}{688.158}{P}%
\TextAt{8}{172.999}{689.951}{is}%
\TextAt{8}{184.763}{689.951}{the}%
\TextAt{8}{204.919}{689.951}{decomposition}%
\TextAt{8}{282.054}{689.951}{group}%
\TextAt{8}{315.229}{689.951}{of}%
\TextAt{8}{328.559}{689.951}{a}%
\TextAt{8}{338.315}{689.951}{prime}%
\TextAt{8}{371.478}{689.951}{ideal}%
\TextAt{14}{399.442}{689.951}{P}%
\TextAt{8}{413.232}{689.951}{over}%
\TextAt{14}{438.278}{689.951}{p}%
\TextAt{8}{444.261}{689.951}{.}%
\TextAt{15}{79.226}{637.886}{§}%
\TextAt{15}{93.785}{637.886}{10}%
\TextAt{15}{130.385}{637.886}{Cyclotomic}%
\TextAt{15}{232.129}{637.886}{Fields}%
\TextAt{8}{79.226}{606}{The}%
\TextAt{8}{102.738}{606}{concepts}%
\TextAt{8}{149.731}{606}{and}%
\TextAt{8}{171.952}{606}{results}%
\TextAt{8}{208.603}{606}{of}%
\TextAt{8}{221.391}{606}{the}%
\TextAt{8}{241.006}{606}{theory}%
\TextAt{8}{277.215}{606}{as}%
\TextAt{8}{291.044}{606}{far}%
\TextAt{8}{308.387}{606}{as}%
\TextAt{8}{322.216}{606}{it}%
\TextAt{8}{333.378}{606}{has}%
\TextAt{8}{353.711}{606}{now}%
\TextAt{8}{377.558}{606}{been}%
\TextAt{8}{404.644}{606}{developed}%
\TextAt{8}{458.403}{606}{have}%
\TextAt{8}{484.857}{606}{reached}%
\TextAt{8}{526.911}{606}{a}%
\TextAt{8}{79.226}{588.954}{degree}%
\TextAt{8}{114.97}{588.954}{of}%
\TextAt{8}{127.629}{588.954}{abstraction}%
\TextAt{8}{188.168}{588.954}{which}%
\TextAt{8}{220.983}{588.954}{we}%
\TextAt{8}{237.539}{588.954}{will}%
\TextAt{8}{258.973}{588.954}{now}%
\TextAt{8}{282.691}{588.954}{balance}%
\TextAt{8}{324.293}{588.954}{by}%
\TextAt{8}{339.881}{588.954}{something}%
\TextAt{8}{395.207}{588.954}{more}%
\TextAt{8}{423.802}{588.954}{concrete.}%
\TextAt{8}{474.3}{588.954}{We}%
\TextAt{8}{493.773}{588.954}{will}%
\TextAt{8}{515.207}{588.954}{put}%
\TextAt{8}{79.226}{571.908}{the}%
\TextAt{8}{98.818}{571.908}{insights}%
\TextAt{8}{140.981}{571.908}{of}%
\TextAt{8}{153.746}{571.908}{the}%
\TextAt{8}{173.337}{571.908}{general}%
\TextAt{8}{213.097}{571.908}{theory}%
\TextAt{8}{249.282}{571.908}{to}%
\TextAt{8}{263.028}{571.908}{the}%
\TextAt{8}{282.619}{571.908}{task}%
\TextAt{8}{307.16}{571.908}{and}%
\TextAt{8}{329.358}{571.908}{make}%
\TextAt{8}{359.363}{571.908}{them}%
\TextAt{8}{388.71}{571.908}{more}%
\TextAt{8}{417.411}{571.908}{explicit}%
\TextAt{8}{458.139}{571.908}{in}%
\TextAt{8}{471.227}{571.908}{the}%
\TextAt{8}{490.819}{571.908}{example}%
\TextAt{8}{79.226}{554.862}{of}%
\TextAt{8}{92.029}{554.862}{the}%
\TextAt{13}{111.658}{554.862}{n}%
\TextAt{12}{118.645}{554.862}{-th}%
\TextAt{12}{139.713}{554.862}{cyclotomic}%
\TextAt{12}{206.02}{554.862}{field}%
\TextAt{18}{234.208}{554.862}{Q}%
\TextAt{9}{243.507}{554.862}{\Character{40}}%
\TextAt{13}{248.059}{554.862}{\Character{16}}%
\TextAt{9}{254.075}{554.862}{\Character{41}}%
\TextAt{8}{258.627}{554.862}{,}%
\TextAt{8}{265.354}{554.862}{where}%
\TextAt{13}{298.636}{554.862}{\Character{16}}%
\TextAt{8}{308.022}{554.862}{is}%
\TextAt{8}{319.258}{554.862}{a}%
\TextAt{12}{328.486}{554.862}{primitive}%
\TextAt{13}{385.874}{554.862}{n}%
\TextAt{12}{392.862}{554.862}{-th}%
\TextAt{12}{413.93}{554.862}{root}%
\TextAt{12}{442.374}{554.862}{of}%
\TextAt{12}{457.094}{554.862}{unity}%
\TextAt{8}{487.747}{554.862}{.}%
\TextAt{8}{496.019}{554.862}{Among}%
\TextAt{8}{79.226}{537.816}{all}%
\TextAt{8}{95.315}{537.816}{number}%
\TextAt{8}{137.741}{537.816}{fields,}%
\TextAt{8}{170.828}{537.816}{this}%
\TextAt{8}{193.48}{537.816}{field}%
\TextAt{8}{218.666}{537.816}{occupies}%
\TextAt{8}{265.049}{537.816}{a}%
\TextAt{8}{274.634}{537.816}{special,}%
\TextAt{8}{315.851}{537.816}{central}%
\TextAt{8}{354.379}{537.816}{place.}%
\TextAt{8}{388.785}{537.816}{So}%
\TextAt{8}{404.873}{537.816}{studying}%
\TextAt{8}{452.572}{537.816}{it}%
\TextAt{8}{464.105}{537.816}{does}%
\TextAt{8}{490.332}{537.816}{not}%
\TextAt{8}{510.975}{537.816}{only}%
\TextAt{8}{79.226}{520.77}{furnish}%
\TextAt{8}{117.823}{520.77}{a}%
\TextAt{8}{126.77}{520.77}{worthwhile}%
\TextAt{8}{185.798}{520.77}{example}%
\TextAt{8}{230.838}{520.77}{but}%
\TextAt{8}{251.49}{520.77}{in}%
\TextAt{8}{264.335}{520.77}{fact}%
\TextAt{8}{286.612}{520.77}{an}%
\TextAt{8}{302.064}{520.77}{essential}%
\TextAt{8}{347.881}{520.77}{building}%
\TextAt{8}{392.598}{520.77}{block}%
\TextAt{8}{422.683}{520.77}{for}%
\TextAt{8}{439.76}{520.77}{the}%
\TextAt{8}{459.109}{520.77}{further}%
\TextAt{8}{497.645}{520.77}{theory.}%
\TextAt{8}{97.159}{500.735}{It}%
\TextAt{8}{110.344}{500.735}{will}%
\TextAt{8}{132.962}{500.735}{be}%
\TextAt{8}{149.4}{500.735}{our}%
\TextAt{8}{170.726}{500.735}{first}%
\TextAt{8}{195.365}{500.735}{goal}%
\TextAt{8}{220.601}{500.735}{to}%
\TextAt{8}{235.424}{500.735}{determine}%
\TextAt{8}{290.561}{500.735}{explicitly}%
\TextAt{8}{341.799}{500.735}{the}%
\TextAt{8}{362.469}{500.735}{ring}%
\TextAt{8}{387.047}{500.735}{of}%
\TextAt{8}{400.89}{500.735}{integers}%
\TextAt{8}{444.717}{500.735}{of}%
\TextAt{8}{458.56}{500.735}{the}%
\TextAt{8}{479.229}{500.735}{field}%
\TextAt{18}{505.099}{500.735}{Q}%
\TextAt{9}{514.397}{500.735}{\Character{40}}%
\TextAt{13}{518.95}{500.735}{\Character{16}}%
\TextAt{9}{524.965}{500.735}{\Character{41}}%
\TextAt{8}{529.518}{500.735}{.}%
\TextAt{8}{79.226}{483.689}{For}%
\TextAt{8}{100.196}{483.689}{this}%
\TextAt{8}{123.018}{483.689}{we}%
\TextAt{8}{140.245}{483.689}{need}%
\TextAt{8}{167.551}{483.689}{the}%
\TextAt{12}{79.226}{458.203}{Lemma}%
\TextAt{12}{126.874}{458.203}{(10.1).}%
\TextAt{8}{168.905}{458.203}{Let}%
\TextAt{13}{189.859}{458.203}{n}%
\TextAt{8}{200.729}{458.203}{be}%
\TextAt{8}{216.638}{458.203}{a}%
\TextAt{8}{226.378}{458.203}{prime}%
\TextAt{8}{259.526}{458.203}{power}%
\TextAt{13}{293.655}{458.203}{\Character{96}}%
\TextAt{21}{298.567}{462.541}{\Character{23}}%
\TextAt{8}{307.631}{458.203}{and}%
\TextAt{8}{330.378}{458.203}{put}%
\TextAt{13}{351.823}{458.203}{\Character{21}}%
\TextAt{9}{361.972}{458.203}{\Character{61}}%
\TextAt{9}{374.382}{458.203}{1}%
\TextAt{16}{382.876}{458.203}{\Character{0}}%
\TextAt{13}{394.8}{458.203}{\Character{16}}%
\TextAt{8}{400.815}{458.203}{.}%
\TextAt{8}{409.259}{458.203}{Then}%
\TextAt{8}{439.801}{458.203}{the}%
\TextAt{8}{459.942}{458.203}{principal}%
\TextAt{8}{508.704}{458.203}{ideal}%
\TextAt{9}{79.226}{441.157}{\Character{40}}%
\TextAt{13}{83.778}{441.157}{\Character{21}}%
\TextAt{9}{90.607}{441.157}{\Character{41}}%
\TextAt{8}{99.057}{441.157}{in}%
\TextAt{8}{112.709}{441.157}{the}%
\TextAt{8}{132.866}{441.157}{ring}%
\TextAt{14}{156.932}{441.157}{o}%
\TextAt{8}{166.673}{441.157}{of}%
\TextAt{8}{180.003}{441.157}{integers}%
\TextAt{8}{223.317}{441.157}{of}%
\TextAt{18}{236.647}{441.157}{Q}%
\TextAt{9}{245.945}{441.157}{\Character{40}}%
\TextAt{13}{250.498}{441.157}{\Character{16}}%
\TextAt{9}{256.513}{441.157}{\Character{41}}%
\TextAt{8}{264.963}{441.157}{is}%
\TextAt{8}{276.727}{441.157}{a}%
\TextAt{8}{286.482}{441.157}{prime}%
\TextAt{8}{319.646}{441.157}{ideal}%
\TextAt{8}{347.609}{441.157}{of}%
\TextAt{8}{360.939}{441.157}{degree}%
\TextAt{9}{397.355}{441.157}{1}%
\TextAt{8}{403.208}{441.157}{,}%
\TextAt{8}{410.357}{441.157}{and}%
\TextAt{8}{433.119}{441.157}{we}%
\TextAt{8}{450.347}{441.157}{have}%
\TextAt{13}{191.966}{413.997}{\Character{96}}%
\TextAt{14}{196.878}{413.997}{o}%
\TextAt{9}{206.044}{413.997}{\Character{61}}%
\TextAt{9}{218.454}{413.997}{\Character{40}}%
\TextAt{13}{223.021}{413.997}{\Character{21}}%
\TextAt{9}{229.85}{413.997}{\Character{41}}%
\TextAt{21}{234.402}{418.933}{d}%
\TextAt{13}{239.258}{413.997}{\Character{59}}%
\TextAt{8}{268.412}{413.997}{where}%
\TextAt{13}{302.221}{413.997}{d}%
\TextAt{9}{311.625}{413.997}{\Character{61}}%
\TextAt{13}{324.05}{413.997}{\Character{39}}%
\TextAt{9}{331.719}{413.997}{\Character{40}}%
\TextAt{13}{336.271}{413.997}{\Character{96}}%
\TextAt{21}{341.183}{418.933}{\Character{23}}%
\TextAt{9}{346.365}{413.997}{\Character{41}}%
\TextAt{9}{354.22}{413.997}{\Character{61}}%
\TextAt{9}{366.641}{413.997}{\Character{91}}%
\TextAt{18}{369.915}{413.997}{Q}%
\TextAt{9}{379.214}{413.997}{\Character{40}}%
\TextAt{13}{383.766}{413.997}{\Character{16}}%
\TextAt{9}{389.782}{413.997}{\Character{41}}%
\TextAt{9}{397.637}{413.997}{\Character{58}}%
\TextAt{18}{404.227}{413.997}{Q}%
\TextAt{9}{413.526}{413.997}{\Character{93}}%
\TextAt{13}{416.778}{413.997}{\Character{58}}%
\TextAt{8}{79.226}{386.837}{Furthermore,}%
\TextAt{8}{150.276}{386.837}{the}%
\TextAt{8}{170.432}{386.837}{basis}%
\TextAt{9}{199.172}{386.837}{1}%
\TextAt{13}{205.025}{386.837}{\Character{59}}%
\TextAt{13}{210.249}{386.837}{\Character{16}}%
\TextAt{13}{216.263}{386.837}{\Character{59}}%
\TextAt{13}{221.487}{386.837}{\Character{58}}%
\TextAt{13}{226.724}{386.837}{\Character{58}}%
\TextAt{13}{231.96}{386.837}{\Character{58}}%
\TextAt{13}{237.185}{386.837}{\Character{59}}%
\TextAt{13}{242.421}{386.837}{\Character{16}}%
\TextAt{21}{248.522}{391.176}{d}%
\TextAt{3}{252.879}{391.176}{\Character{0}}%
\TextAt{17}{259.465}{391.176}{1}%
\TextAt{8}{268.095}{386.837}{of}%
\TextAt{18}{281.425}{386.837}{Q}%
\TextAt{9}{290.724}{386.837}{\Character{40}}%
\TextAt{13}{295.276}{386.837}{\Character{16}}%
\TextAt{9}{301.291}{386.837}{\Character{41}}%
\TextAt{16}{309.165}{386.837}{\Character{106}}%
\TextAt{18}{315.806}{386.837}{Q}%
\TextAt{8}{329.002}{386.837}{has}%
\TextAt{8}{349.876}{386.837}{the}%
\TextAt{8}{370.032}{386.837}{discriminant}%
\TextAt{13}{180.331}{359.677}{d}%
\TextAt{9}{186.413}{359.677}{\Character{40}1}%
\TextAt{13}{196.819}{359.677}{\Character{59}}%
\TextAt{13}{202.043}{359.677}{\Character{16}}%
\TextAt{13}{208.057}{359.677}{\Character{59}}%
\TextAt{13}{213.281}{359.677}{\Character{58}}%
\TextAt{13}{218.518}{359.677}{\Character{58}}%
\TextAt{13}{223.754}{359.677}{\Character{58}}%
\TextAt{13}{228.979}{359.677}{\Character{59}}%
\TextAt{13}{234.215}{359.677}{\Character{16}}%
\TextAt{21}{240.315}{364.614}{d}%
\TextAt{3}{244.672}{364.614}{\Character{0}}%
\TextAt{17}{251.259}{364.614}{1}%
\TextAt{9}{255.991}{359.677}{\Character{41}}%
\TextAt{9}{263.846}{359.677}{\Character{61}}%
\NativeMathAt{276.29}{359.677}{11.955}{\pm}%
\TextAt{13}{285.588}{359.677}{\Character{96}}%
\TextAt{21}{290.5}{364.614}{s}%
\TextAt{13}{294.914}{359.677}{\Character{59}}%
\TextAt{13}{324.049}{359.677}{s}%
\TextAt{9}{332.904}{359.677}{\Character{61}}%
\TextAt{13}{345.329}{359.677}{\Character{96}}%
\TextAt{21}{350.241}{364.614}{\Character{23}}%
\TextAt{3}{354.925}{364.614}{\Character{0}}%
\TextAt{17}{361.512}{364.614}{1}%
\TextAt{9}{366.244}{359.677}{\Character{40}}%
\TextAt{13}{370.796}{359.677}{\Character{23}}%
\TextAt{13}{377.324}{359.677}{\Character{96}}%
\TextAt{16}{384.903}{359.677}{\Character{0}}%
\TextAt{13}{396.858}{359.677}{\Character{23}}%
\TextAt{16}{406.053}{359.677}{\Character{0}}%
\TextAt{9}{418.008}{359.677}{1\Character{41}}%
\TextAt{13}{428.413}{359.677}{\Character{58}}%
\TextAt{19}{79.226}{315.637}{Proof.}%
\TextAt{8}{116}{315.637}{The}%
\TextAt{8}{140.054}{315.637}{minimal}%
\TextAt{8}{185.579}{315.637}{polynomial}%
\TextAt{8}{246.072}{315.637}{of}%
\TextAt{13}{259.402}{315.637}{\Character{16}}%
\TextAt{8}{269.315}{315.637}{over}%
\TextAt{18}{294.362}{315.637}{Q}%
\TextAt{8}{307.557}{315.637}{is}%
\TextAt{8}{319.321}{315.637}{the}%
\TextAt{13}{339.478}{315.637}{n}%
\TextAt{8}{346.465}{315.637}{-th}%
\TextAt{8}{365.319}{315.637}{cyclotomic}%
\TextAt{8}{423.528}{315.637}{polynomial}%
\TextAt{13}{176.11}{284.127}{\Character{30}}%
\TextAt{21}{183.039}{282.334}{n}%
\TextAt{9}{188.675}{284.127}{\Character{40}}%
\TextAt{13}{193.227}{284.127}{X}%
\TextAt{9}{203.882}{284.127}{\Character{41}}%
\TextAt{9}{211.737}{284.127}{\Character{61}}%
\TextAt{13}{230.081}{292.215}{X}%
\TextAt{21}{240.736}{296.553}{\Character{96}}%
\TextAt{22}{244.195}{299.366}{\Character{23}}%
\TextAt{16}{251.972}{292.215}{\Character{0}}%
\TextAt{9}{263.927}{292.215}{1}%
\Vector{q 0 G 0 g 0.478 w 153.377 -486.884 m 202.484 -486.884 l S Q}%
\TextAt{13}{225.377}{275.867}{X}%
\TextAt{21}{236.032}{279.321}{\Character{96}}%
\TextAt{22}{239.491}{281.59}{\Character{23}}%
\TextAt{23}{243.614}{281.59}{\Character{0}}%
\TextAt{24}{249.37}{281.59}{1}%
\TextAt{16}{256.676}{275.867}{\Character{0}}%
\TextAt{9}{268.631}{275.867}{1}%
\TextAt{9}{279.001}{284.127}{\Character{61}}%
\TextAt{13}{291.426}{284.127}{X}%
\TextAt{21}{302.081}{289.063}{\Character{96}}%
\TextAt{22}{305.541}{291.876}{\Character{23}}%
\TextAt{23}{309.664}{291.876}{\Character{0}}%
\TextAt{24}{315.42}{291.876}{1}%
\TextAt{17}{319.571}{289.063}{\Character{40}}%
\TextAt{21}{322.864}{289.063}{\Character{96}}%
\TextAt{3}{326.324}{289.063}{\Character{0}}%
\TextAt{17}{332.91}{289.063}{1\Character{41}}%
\TextAt{9}{343.592}{284.127}{\Character{43}}%
\TextAt{16}{355.354}{284.127}{\Character{1}}%
\TextAt{16}{360.65}{284.127}{\Character{1}}%
\TextAt{16}{365.958}{284.127}{\Character{1}}%
\TextAt{9}{371.958}{284.127}{\Character{43}}%
\TextAt{13}{383.719}{284.127}{X}%
\TextAt{21}{394.375}{289.063}{\Character{96}}%
\TextAt{22}{397.834}{291.876}{\Character{23}}%
\TextAt{23}{401.957}{291.876}{\Character{0}}%
\TextAt{24}{407.713}{291.876}{1}%
\TextAt{9}{415.019}{284.127}{\Character{43}}%
\TextAt{9}{426.771}{284.127}{1}%
\TextAt{13}{432.634}{284.127}{\Character{58}}%
\TextAt{8}{79.226}{255.464}{Putting}%
\TextAt{13}{122.313}{255.464}{X}%
\TextAt{9}{136.289}{255.464}{\Character{61}}%
\TextAt{9}{148.699}{255.464}{1}%
\TextAt{8}{154.567}{255.464}{,}%
\TextAt{8}{161.716}{255.464}{we}%
\TextAt{8}{178.944}{255.464}{obtain}%
\TextAt{8}{215.371}{255.464}{the}%
\TextAt{8}{235.527}{255.464}{identity}%
\TextAt{13}{253.411}{228.304}{\Character{96}}%
\TextAt{9}{261.644}{228.304}{\Character{61}}%
\TextAt{20}{286.864}{239.661}{\Character{89}}%
\TextAt{21}{274.07}{213.758}{g}%
\NativeMathAt{278.385}{213.758}{7.97}{\in}%
\TextAt{17}{284.031}{213.758}{\Character{40}}%
\TextAt{25}{287.324}{213.758}{Z}%
\TextAt{17}{292.638}{213.758}{\Character{47}}%
\TextAt{21}{296.872}{213.758}{n}%
\TextAt{25}{302.01}{213.758}{Z}%
\TextAt{17}{307.323}{213.758}{\Character{41}}%
\TextAt{23}{310.617}{216.028}{\Character{3}}%
\TextAt{9}{314.934}{228.304}{\Character{40}1}%
\TextAt{16}{327.996}{228.304}{\Character{0}}%
\TextAt{13}{339.951}{228.304}{\Character{16}}%
\TextAt{21}{345.967}{233.24}{g}%
\TextAt{9}{350.781}{228.304}{\Character{41}}%
\TextAt{13}{355.333}{228.304}{\Character{58}}%
\TextAt{8}{79.226}{188.669}{But}%
\TextAt{9}{102.785}{188.669}{1}%
\TextAt{16}{111.511}{188.669}{\Character{0}}%
\TextAt{13}{123.683}{188.669}{\Character{16}}%
\TextAt{21}{129.699}{193.008}{g}%
\TextAt{9}{138.376}{188.669}{\Character{61}}%
\TextAt{13}{151.343}{188.669}{\Character{34}}%
\TextAt{21}{156.821}{186.876}{g}%
\TextAt{9}{161.635}{188.669}{\Character{40}1}%
\TextAt{16}{174.914}{188.669}{\Character{0}}%
\TextAt{13}{187.086}{188.669}{\Character{16}}%
\TextAt{9}{193.101}{188.669}{\Character{41}}%
\TextAt{8}{197.653}{188.669}{,}%
\TextAt{8}{205.2}{188.669}{for}%
\TextAt{8}{223.403}{188.669}{the}%
\TextAt{8}{243.877}{188.669}{algebraic}%
\TextAt{8}{293.63}{188.669}{integer}%
\TextAt{13}{332.647}{188.669}{\Character{34}}%
\TextAt{21}{338.125}{186.876}{g}%
\TextAt{9}{346.802}{188.669}{\Character{61}}%
\TextAt{17}{360.964}{193.925}{1}%
\TextAt{3}{365.199}{193.925}{\Character{0}}%
\TextAt{21}{371.785}{193.925}{\Character{16}}%
\TextAt{22}{376.054}{196.738}{g}%
\Vector{q 0 G 0 g 0.478 w 288.964 -582.342 m 308.399 -582.342 l S Q}%
\TextAt{17}{363.137}{184.547}{1}%
\TextAt{3}{367.371}{184.547}{\Character{0}}%
\TextAt{21}{373.958}{184.547}{\Character{16}}%
\TextAt{9}{385.458}{188.669}{\Character{61}}%
\TextAt{9}{398.417}{188.669}{1}%
\TextAt{9}{407.133}{188.669}{\Character{43}}%
\TextAt{13}{419.129}{188.669}{\Character{16}}%
\TextAt{9}{428.018}{188.669}{\Character{43}}%
\TextAt{16}{439.997}{188.669}{\Character{1}}%
\TextAt{16}{445.293}{188.669}{\Character{1}}%
\TextAt{16}{450.601}{188.669}{\Character{1}}%
\TextAt{9}{456.818}{188.669}{\Character{43}}%
\TextAt{13}{468.796}{188.669}{\Character{16}}%
\TextAt{21}{474.811}{193.008}{g}%
\TextAt{3}{479.127}{193.008}{\Character{0}}%
\TextAt{17}{485.714}{193.008}{1}%
\TextAt{8}{490.446}{188.669}{.}%
\TextAt{8}{499.848}{188.669}{If}%
\TextAt{13}{511.858}{188.669}{g}%
\TextAt{3}{517.893}{193.008}{\Character{48}}%
\TextAt{8}{524.903}{188.669}{is}%
\TextAt{8}{79.226}{171.624}{an}%
\TextAt{8}{95.485}{171.624}{integer}%
\TextAt{8}{134.184}{171.624}{such}%
\TextAt{8}{160.581}{171.624}{that}%
\TextAt{13}{185.95}{171.624}{g}%
\TextAt{13}{191.964}{171.624}{g}%
\TextAt{3}{198.018}{175.962}{\Character{48}}%
\NativeMathAt{204.134}{171.624}{11.955}{\equiv}%
\TextAt{9}{216.754}{171.624}{1}%
\TextAt{9}{227.908}{171.624}{\Character{40}}%
\TextAt{8}{232.472}{171.624}{mod}%
\TextAt{13}{258.574}{171.624}{\Character{96}}%
\TextAt{21}{263.487}{175.962}{\Character{23}}%
\TextAt{9}{268.669}{171.624}{\Character{41}}%
\TextAt{8}{273.221}{171.624}{,}%
\TextAt{8}{280.37}{171.624}{then}%
\TextAt{9}{195.798}{147.052}{1}%
\TextAt{16}{204.308}{147.052}{\Character{0}}%
\TextAt{13}{216.263}{147.052}{\Character{16}}%
\Vector{q 0 G 0 g 0.478 w 121.391 -632.046 m 152.686 -632.046 l S Q}%
\TextAt{9}{193.391}{130.764}{1}%
\TextAt{16}{201.901}{130.764}{\Character{0}}%
\TextAt{13}{213.856}{130.764}{\Character{16}}%
\TextAt{21}{219.872}{134.218}{g}%
\TextAt{9}{229.202}{138.965}{\Character{61}}%
\TextAt{9}{242.823}{147.052}{1}%
\TextAt{16}{251.333}{147.052}{\Character{0}}%
\TextAt{9}{263.288}{147.052}{\Character{40}}%
\TextAt{13}{267.84}{147.052}{\Character{16}}%
\TextAt{21}{273.856}{151.391}{g}%
\TextAt{9}{278.67}{147.052}{\Character{41}}%
\TextAt{21}{283.222}{151.391}{g}%
\TextAt{23}{287.538}{154.204}{\Character{48}}%
\Vector{q 0 G 0 g 0.478 w 170.823 -632.046 m 218.72 -632.046 l S Q}%
\TextAt{9}{251.125}{130.764}{1}%
\TextAt{16}{259.634}{130.764}{\Character{0}}%
\TextAt{13}{271.589}{130.764}{\Character{16}}%
\TextAt{21}{277.605}{134.218}{g}%
\TextAt{9}{295.237}{138.965}{\Character{61}}%
\TextAt{9}{307.647}{138.965}{1}%
\TextAt{9}{316.147}{138.965}{\Character{43}}%
\TextAt{13}{327.933}{138.965}{\Character{16}}%
\TextAt{21}{333.949}{143.901}{g}%
\TextAt{9}{341.42}{138.965}{\Character{43}}%
\TextAt{16}{353.181}{138.965}{\Character{1}}%
\TextAt{16}{358.477}{138.965}{\Character{1}}%
\TextAt{16}{363.785}{138.965}{\Character{1}}%
\TextAt{9}{369.785}{138.965}{\Character{43}}%
\TextAt{9}{381.537}{138.965}{\Character{40}}%
\TextAt{13}{386.099}{138.965}{\Character{16}}%
\TextAt{21}{392.114}{143.901}{g}%
\TextAt{9}{396.928}{138.965}{\Character{41}}%
\TextAt{21}{401.481}{143.901}{g}%
\TextAt{23}{405.796}{146.714}{\Character{48}}%
\TextAt{3}{408.481}{143.901}{\Character{0}}%
\TextAt{17}{415.067}{143.901}{1}%
\TextAt{8}{79.226}{107.013}{is}%
\TextAt{8}{90.052}{107.013}{integral}%
\TextAt{8}{131.722}{107.013}{as}%
\TextAt{8}{145.153}{107.013}{well,}%
\TextAt{8}{171.386}{107.013}{i.e.,}%
\TextAt{13}{192.74}{107.013}{\Character{34}}%
\TextAt{21}{198.218}{105.22}{g}%
\TextAt{8}{205.991}{107.013}{is}%
\TextAt{8}{216.817}{107.013}{a}%
\TextAt{8}{225.634}{107.013}{unit.}%
\TextAt{8}{254.583}{107.013}{Consequently}%
\TextAt{13}{326.225}{107.013}{\Character{96}}%
\TextAt{9}{334.458}{107.013}{\Character{61}}%
\TextAt{13}{346.883}{107.013}{\Character{34}}%
\TextAt{9}{352.361}{107.013}{\Character{40}1}%
\TextAt{16}{363.504}{107.013}{\Character{0}}%
\TextAt{13}{373.541}{107.013}{\Character{16}}%
\TextAt{9}{379.556}{107.013}{\Character{41}}%
\TextAt{21}{384.109}{111.352}{\Character{39}}%
\TextAt{17}{389.627}{111.352}{\Character{40}}%
\TextAt{21}{392.92}{111.352}{\Character{96}}%
\TextAt{22}{396.38}{114.165}{\Character{23}}%
\TextAt{17}{401.001}{111.352}{\Character{41}}%
\TextAt{8}{404.792}{107.013}{,}%
\TextAt{8}{411.191}{107.013}{with}%
\TextAt{8}{436.913}{107.013}{the}%
\TextAt{8}{456.131}{107.013}{unit}%
\TextAt{13}{479.904}{107.013}{\Character{34}}%
\TextAt{9}{488.703}{107.013}{\Character{61}}%
\TextAt{20}{501.128}{115.98}{\Character{81}}%
\TextAt{21}{512.419}{103.526}{g}%
\TextAt{13}{519.226}{107.013}{\Character{34}}%
\TextAt{21}{524.704}{105.22}{g}%
\TextAt{8}{529.518}{107.013}{,}%
\TextAt{8}{79.226}{89.967}{hence}%
\TextAt{13}{111.961}{89.967}{\Character{96}}%
\TextAt{14}{116.873}{89.967}{o}%
\TextAt{9}{126.429}{89.967}{\Character{61}}%
\TextAt{9}{139.233}{89.967}{\Character{40}}%
\TextAt{13}{143.797}{89.967}{\Character{21}}%
\TextAt{9}{150.625}{89.967}{\Character{41}}%
\TextAt{21}{155.178}{94.306}{\Character{39}}%
\TextAt{17}{160.696}{94.306}{\Character{40}}%
\TextAt{21}{163.989}{94.306}{\Character{96}}%
\TextAt{22}{167.449}{97.119}{\Character{23}}%
\TextAt{17}{172.07}{94.306}{\Character{41}}%
\TextAt{8}{175.861}{89.967}{.}%
\TextAt{8}{184.997}{89.967}{Since}%
\TextAt{9}{215.783}{89.967}{\Character{91}}%
\TextAt{18}{219.035}{89.967}{Q}%
\TextAt{9}{228.333}{89.967}{\Character{40}}%
\TextAt{13}{232.886}{89.967}{\Character{16}}%
\TextAt{9}{238.901}{89.967}{\Character{41}}%
\TextAt{9}{247.15}{89.967}{\Character{58}}%
\TextAt{18}{254.128}{89.967}{Q}%
\TextAt{9}{263.426}{89.967}{\Character{93}}%
\TextAt{9}{270.372}{89.967}{\Character{61}}%
\TextAt{13}{283.205}{89.967}{\Character{39}}%
\TextAt{9}{290.873}{89.967}{\Character{40}}%
\TextAt{13}{295.425}{89.967}{\Character{96}}%
\TextAt{21}{300.337}{94.306}{\Character{23}}%
\TextAt{9}{305.519}{89.967}{\Character{41}}%
\TextAt{8}{310.072}{89.967}{,}%
\TextAt{8}{317.507}{89.967}{the}%
\TextAt{8}{337.893}{89.967}{fundamental}%
\TextAt{8}{405.764}{89.967}{identity}%
\TextAt{8}{449.247}{89.967}{(8.2)}%
\TextAt{8}{477.451}{89.967}{shows}%
\TextAt{8}{511.298}{89.967}{that}%
\TextAt{9}{79.226}{72.921}{\Character{40}}%
\TextAt{13}{83.778}{72.921}{\Character{21}}%
\TextAt{9}{90.607}{72.921}{\Character{41}}%
\TextAt{8}{99.057}{72.921}{is}%
\TextAt{8}{110.821}{72.921}{a}%
\TextAt{8}{120.576}{72.921}{prime}%
\TextAt{8}{153.74}{72.921}{ideal}%
\TextAt{8}{181.703}{72.921}{of}%
\TextAt{8}{195.033}{72.921}{degree}%
\TextAt{9}{231.448}{72.921}{1}%
\TextAt{8}{237.301}{72.921}{.}%
\TextAt{8}{300.14}{43.033}{55}%
\Destination{78.23}{809.815}{nav.page.55}%
\Destination{79.23}{663.435}{nav.section.1.10}%
\Destination{79.226}{775.155}{nav.ex.1.9.4}%
\Destination{79.36}{603.955}{nav.print.59}%
\Destination{79.36}{472.155}{nav.number.1.10.1}%
\Destination{79.23}{663.435}{outline.15}%
\Hotlink{448.747}{87.755}{25.078}{11.047}{nav.number.1.8.2}%
}

% Page 68 (body page 56)
\Page{%
\TextAt{2}{79.226}{796.943}{I.}%
\TextAt{2}{89.83}{796.943}{Algebraic}%
\TextAt{2}{138.648}{796.943}{Integers}%
\Vector{q 0 G 0 g 0.299 w 7.226 18.729 m 460.769 18.729 l S Q}%
\TextAt{8}{97.159}{761.124}{Let}%
\TextAt{13}{117.733}{761.124}{\Character{16}}%
\TextAt{9}{127.069}{761.124}{\Character{61}}%
\TextAt{13}{139.495}{761.124}{\Character{16}}%
\TextAt{17}{144.608}{759.33}{1}%
\TextAt{13}{149.34}{761.124}{\Character{59}}%
\TextAt{13}{154.564}{761.124}{\Character{58}}%
\TextAt{13}{159.801}{761.124}{\Character{58}}%
\TextAt{13}{165.025}{761.124}{\Character{58}}%
\TextAt{13}{170.262}{761.124}{\Character{59}}%
\TextAt{13}{175.498}{761.124}{\Character{16}}%
\TextAt{21}{180.673}{759.33}{d}%
\TextAt{8}{189.031}{761.124}{be}%
\TextAt{8}{204.56}{761.124}{the}%
\TextAt{8}{224.321}{761.124}{conjugates}%
\TextAt{8}{281.549}{761.124}{of}%
\TextAt{13}{294.484}{761.124}{\Character{16}}%
\TextAt{8}{300.5}{761.124}{.}%
\TextAt{8}{308.816}{761.124}{Then}%
\TextAt{8}{338.978}{761.124}{the}%
\TextAt{8}{358.739}{761.124}{cyclotomic}%
\TextAt{8}{416.554}{761.124}{polynomial}%
\TextAt{8}{476.651}{761.124}{is}%
\TextAt{13}{488.02}{761.124}{\Character{30}}%
\TextAt{21}{494.948}{759.33}{n}%
\TextAt{9}{500.584}{761.124}{\Character{40}}%
\TextAt{13}{505.137}{761.124}{X}%
\TextAt{9}{515.792}{761.124}{\Character{41}}%
\TextAt{9}{523.647}{761.124}{\Character{61}}%
\TextAt{20}{79.226}{753.044}{\Character{81}}%
\TextAt{21}{90.517}{749.889}{d}%
\TextAt{21}{90.517}{740.591}{i}%
\TextAt{17}{93.4}{740.591}{\Character{61}1}%
\TextAt{9}{104.719}{744.078}{\Character{40}}%
\TextAt{13}{109.271}{744.078}{X}%
\TextAt{16}{122.583}{744.078}{\Character{0}}%
\TextAt{13}{134.538}{744.078}{\Character{16}}%
\TextAt{21}{139.651}{742.284}{i}%
\TextAt{9}{143.032}{744.078}{\Character{41}}%
\TextAt{8}{151.482}{744.078}{and}%
\TextAt{8}{174.245}{744.078}{(see}%
\TextAt{8}{197.713}{744.078}{§2,}%
\TextAt{8}{216.399}{744.078}{p.}%
\TextAt{8}{230.051}{744.078}{11)}%
\NativeMathAt{147.657}{706.226}{11.955}{\pm}%
\TextAt{13}{156.955}{706.226}{d}%
\TextAt{9}{163.038}{706.226}{\Character{40}1}%
\TextAt{13}{173.443}{706.226}{\Character{59}}%
\TextAt{13}{178.667}{706.226}{\Character{16}}%
\TextAt{13}{184.681}{706.226}{\Character{59}}%
\TextAt{13}{189.905}{706.226}{\Character{58}}%
\TextAt{13}{195.142}{706.226}{\Character{58}}%
\TextAt{13}{200.378}{706.226}{\Character{58}}%
\TextAt{13}{205.603}{706.226}{\Character{59}}%
\TextAt{13}{210.839}{706.226}{\Character{16}}%
\TextAt{21}{216.94}{711.162}{d}%
\TextAt{3}{221.297}{711.162}{\Character{0}}%
\TextAt{17}{227.883}{711.162}{1}%
\TextAt{9}{232.616}{706.226}{\Character{41}}%
\TextAt{9}{240.471}{706.226}{\Character{61}}%
\TextAt{20}{252.914}{717.583}{\Character{89}}%
\TextAt{21}{253.876}{692.123}{i}%
\NativeMathAt{256.759}{692.123}{7.97}{\neq}%
\TextAt{21}{263.345}{692.123}{j}%
\TextAt{9}{268.191}{706.226}{\Character{40}}%
\TextAt{13}{272.743}{706.226}{\Character{16}}%
\TextAt{21}{277.856}{704.432}{i}%
\TextAt{16}{283.893}{706.226}{\Character{0}}%
\TextAt{13}{295.849}{706.226}{\Character{16}}%
\TextAt{21}{300.961}{704.432}{j}%
\TextAt{9}{305.344}{706.226}{\Character{41}}%
\TextAt{9}{313.199}{706.226}{\Character{61}}%
\TextAt{21}{331.102}{721.17}{d}%
\TextAt{20}{325.642}{717.583}{\Character{89}}%
\TextAt{21}{326.428}{692.389}{i}%
\TextAt{17}{329.311}{692.389}{\Character{61}1}%
\TextAt{13}{342.911}{706.226}{\Character{30}}%
\TextAt{3}{349.839}{711.162}{\Character{48}}%
\TextAt{21}{349.839}{703.27}{n}%
\TextAt{9}{355.475}{706.226}{\Character{40}}%
\TextAt{13}{360.028}{706.226}{\Character{16}}%
\TextAt{21}{365.14}{704.432}{i}%
\TextAt{9}{368.522}{706.226}{\Character{41}}%
\TextAt{9}{376.377}{706.226}{\Character{61}}%
\TextAt{13}{388.82}{706.226}{N}%
\TextAt{25}{398.194}{704.366}{Q}%
\TextAt{17}{404.393}{704.366}{\Character{40}}%
\TextAt{21}{407.686}{704.366}{\Character{16}}%
\TextAt{17}{411.955}{704.366}{\Character{41}}%
\TextAt{3}{415.248}{704.366}{\Character{106}}%
\TextAt{25}{417.601}{704.366}{Q}%
\TextAt{9}{424.298}{706.226}{\Character{40}}%
\TextAt{13}{428.85}{706.226}{\Character{30}}%
\TextAt{3}{435.778}{711.162}{\Character{48}}%
\TextAt{21}{435.778}{703.27}{n}%
\TextAt{9}{441.415}{706.226}{\Character{40}}%
\TextAt{13}{445.967}{706.226}{\Character{16}}%
\TextAt{9}{451.982}{706.226}{\Character{41}\Character{41}}%
\TextAt{13}{461.087}{706.226}{\Character{58}}%
\TextAt{8}{79.226}{668.27}{Differentiating}%
\TextAt{8}{157.114}{668.27}{the}%
\TextAt{8}{177.27}{668.27}{equation}%
\TextAt{9}{233.007}{639.53}{\Character{40}}%
\TextAt{13}{237.559}{639.53}{X}%
\TextAt{21}{248.214}{644.466}{\Character{96}}%
\TextAt{22}{251.674}{647.279}{\Character{23}}%
\TextAt{23}{255.797}{647.279}{\Character{0}}%
\TextAt{24}{261.553}{647.279}{1}%
\TextAt{16}{268.859}{639.53}{\Character{0}}%
\TextAt{9}{280.814}{639.53}{1\Character{41}}%
\TextAt{13}{291.219}{639.53}{\Character{30}}%
\TextAt{21}{298.148}{637.737}{n}%
\TextAt{9}{303.784}{639.53}{\Character{40}}%
\TextAt{13}{308.336}{639.53}{X}%
\TextAt{9}{318.991}{639.53}{\Character{41}}%
\TextAt{9}{326.846}{639.53}{\Character{61}}%
\TextAt{13}{339.29}{639.53}{X}%
\TextAt{21}{349.945}{644.466}{\Character{96}}%
\TextAt{22}{353.404}{647.279}{\Character{23}}%
\TextAt{16}{361.181}{639.53}{\Character{0}}%
\TextAt{9}{373.136}{639.53}{1}%
\TextAt{8}{79.226}{610.79}{and}%
\TextAt{8}{101.989}{610.79}{substituting}%
\TextAt{13}{167.156}{610.79}{\Character{16}}%
\TextAt{8}{177.069}{610.79}{for}%
\TextAt{13}{194.954}{610.79}{X}%
\TextAt{8}{209.507}{610.79}{yields}%
\TextAt{9}{251.312}{582.049}{\Character{40}}%
\TextAt{13}{255.865}{582.049}{\Character{24}}%
\TextAt{16}{264.212}{582.049}{\Character{0}}%
\TextAt{9}{276.167}{582.049}{1\Character{41}}%
\TextAt{13}{286.572}{582.049}{\Character{30}}%
\TextAt{3}{293.5}{586.985}{\Character{48}}%
\TextAt{21}{293.5}{579.094}{n}%
\TextAt{9}{299.137}{582.049}{\Character{40}}%
\TextAt{13}{303.689}{582.049}{\Character{16}}%
\TextAt{9}{309.705}{582.049}{\Character{41}}%
\TextAt{9}{317.56}{582.049}{\Character{61}}%
\TextAt{13}{330.003}{582.049}{\Character{96}}%
\TextAt{21}{334.915}{586.985}{\Character{23}}%
\TextAt{13}{340.097}{582.049}{\Character{16}}%
\TextAt{3}{346.113}{586.985}{\Character{0}}%
\TextAt{17}{352.699}{586.985}{1}%
\TextAt{13}{357.432}{582.049}{\Character{59}}%
\TextAt{8}{79.226}{553.309}{with}%
\TextAt{8}{105.886}{553.309}{the}%
\TextAt{8}{126.043}{553.309}{primitive}%
\TextAt{13}{176.123}{553.309}{\Character{96}}%
\TextAt{8}{181.035}{553.309}{-th}%
\TextAt{8}{199.888}{553.309}{root}%
\TextAt{8}{224.934}{553.309}{of}%
\TextAt{8}{238.264}{553.309}{unity}%
\TextAt{13}{268.834}{553.309}{\Character{24}}%
\TextAt{9}{277.845}{553.309}{\Character{61}}%
\TextAt{13}{290.27}{553.309}{\Character{16}}%
\TextAt{21}{296.286}{557.647}{\Character{96}}%
\TextAt{22}{299.745}{560.46}{\Character{23}}%
\TextAt{23}{303.869}{560.46}{\Character{0}}%
\TextAt{24}{309.625}{560.46}{1}%
\TextAt{8}{314.274}{553.309}{.}%
\TextAt{8}{322.722}{553.309}{But}%
\TextAt{13}{345.963}{553.309}{N}%
\TextAt{25}{355.337}{551.449}{Q}%
\TextAt{17}{361.536}{551.449}{\Character{40}}%
\TextAt{21}{364.829}{551.449}{\Character{24}}%
\TextAt{17}{368.863}{551.449}{\Character{41}}%
\TextAt{3}{372.156}{551.449}{\Character{106}}%
\TextAt{25}{374.508}{551.449}{Q}%
\TextAt{9}{381.206}{553.309}{\Character{40}}%
\TextAt{13}{385.758}{553.309}{\Character{24}}%
\TextAt{16}{394.105}{553.309}{\Character{0}}%
\TextAt{9}{406.06}{553.309}{1\Character{41}}%
\TextAt{9}{419.761}{553.309}{\Character{61}}%
\NativeMathAt{432.212}{553.309}{11.955}{\pm}%
\TextAt{13}{441.51}{553.309}{\Character{96}}%
\TextAt{8}{446.422}{553.309}{,}%
\TextAt{8}{453.572}{553.309}{so}%
\TextAt{8}{467.942}{553.309}{that}%
\TextAt{13}{192.913}{524.568}{N}%
\TextAt{25}{202.286}{522.708}{Q}%
\TextAt{17}{208.485}{522.708}{\Character{40}}%
\TextAt{21}{211.779}{522.708}{\Character{16}}%
\TextAt{17}{216.047}{522.708}{\Character{41}}%
\TextAt{3}{219.341}{522.708}{\Character{106}}%
\TextAt{25}{221.693}{522.708}{Q}%
\TextAt{9}{228.39}{524.568}{\Character{40}}%
\TextAt{13}{232.942}{524.568}{\Character{24}}%
\TextAt{16}{241.29}{524.568}{\Character{0}}%
\TextAt{9}{253.245}{524.568}{1\Character{41}}%
\TextAt{9}{266.946}{524.568}{\Character{61}}%
\TextAt{13}{279.396}{524.568}{N}%
\TextAt{25}{288.77}{522.708}{Q}%
\TextAt{17}{294.969}{522.708}{\Character{40}}%
\TextAt{21}{298.262}{522.708}{\Character{24}}%
\TextAt{17}{302.296}{522.708}{\Character{41}}%
\TextAt{3}{305.589}{522.708}{\Character{106}}%
\TextAt{25}{307.942}{522.708}{Q}%
\TextAt{9}{314.639}{524.568}{\Character{40}}%
\TextAt{13}{319.191}{524.568}{\Character{24}}%
\TextAt{16}{327.538}{524.568}{\Character{0}}%
\TextAt{9}{339.493}{524.568}{1\Character{41}}%
\TextAt{21}{349.899}{529.504}{\Character{96}}%
\TextAt{22}{353.358}{532.317}{\Character{23}}%
\TextAt{23}{357.481}{532.317}{\Character{0}}%
\TextAt{24}{363.237}{532.317}{1}%
\TextAt{9}{371.207}{524.568}{\Character{61}}%
\NativeMathAt{383.633}{524.568}{11.955}{\pm}%
\TextAt{13}{392.931}{524.568}{\Character{96}}%
\TextAt{21}{397.843}{529.504}{\Character{96}}%
\TextAt{22}{401.303}{532.317}{\Character{23}}%
\TextAt{23}{405.426}{532.317}{\Character{0}}%
\TextAt{24}{411.182}{532.317}{1}%
\TextAt{13}{415.831}{524.568}{\Character{58}}%
\TextAt{8}{79.226}{495.828}{Observing}%
\TextAt{8}{134.901}{495.828}{that}%
\TextAt{13}{160.27}{495.828}{\Character{16}}%
\TextAt{3}{166.286}{500.166}{\Character{0}}%
\TextAt{17}{172.872}{500.166}{1}%
\TextAt{8}{181.502}{495.828}{has}%
\TextAt{8}{202.376}{495.828}{norm}%
\NativeMathAt{232.945}{495.828}{11.955}{\pm}%
\TextAt{9}{242.244}{495.828}{1}%
\TextAt{8}{251.994}{495.828}{we}%
\TextAt{8}{269.221}{495.828}{obtain}%
\TextAt{13}{151.264}{467.087}{d}%
\TextAt{9}{157.347}{467.087}{\Character{40}1}%
\TextAt{13}{167.752}{467.087}{\Character{59}}%
\TextAt{13}{172.976}{467.087}{\Character{16}}%
\TextAt{13}{178.99}{467.087}{\Character{59}}%
\TextAt{13}{184.214}{467.087}{\Character{58}}%
\TextAt{13}{189.451}{467.087}{\Character{58}}%
\TextAt{13}{194.687}{467.087}{\Character{58}}%
\TextAt{13}{199.912}{467.087}{\Character{59}}%
\TextAt{13}{205.148}{467.087}{\Character{16}}%
\TextAt{21}{211.249}{472.023}{d}%
\TextAt{3}{215.606}{472.023}{\Character{0}}%
\TextAt{17}{222.192}{472.023}{1}%
\TextAt{9}{226.925}{467.087}{\Character{41}}%
\TextAt{9}{234.78}{467.087}{\Character{61}}%
\NativeMathAt{247.223}{467.087}{11.955}{\pm}%
\TextAt{13}{256.522}{467.087}{N}%
\TextAt{25}{265.896}{465.228}{Q}%
\TextAt{17}{272.095}{465.227}{\Character{40}}%
\TextAt{21}{275.388}{465.228}{\Character{16}}%
\TextAt{17}{279.657}{465.227}{\Character{41}}%
\TextAt{3}{282.95}{465.228}{\Character{106}}%
\TextAt{25}{285.302}{465.227}{Q}%
\TextAt{9}{291.999}{467.087}{\Character{40}}%
\TextAt{13}{296.552}{467.087}{\Character{30}}%
\TextAt{3}{303.48}{472.023}{\Character{48}}%
\TextAt{21}{303.48}{464.132}{n}%
\TextAt{9}{309.116}{467.087}{\Character{40}}%
\TextAt{13}{313.668}{467.087}{\Character{16}}%
\TextAt{9}{319.684}{467.087}{\Character{41}\Character{41}}%
\TextAt{9}{332.082}{467.087}{\Character{61}}%
\NativeMathAt{344.535}{467.087}{11.955}{\pm}%
\TextAt{13}{353.834}{467.087}{\Character{96}}%
\TextAt{21}{358.746}{472.023}{\Character{23}}%
\TextAt{21}{363.424}{472.023}{\Character{96}}%
\TextAt{22}{366.889}{474.836}{\Character{23}}%
\TextAt{23}{371.012}{474.836}{\Character{0}}%
\TextAt{24}{376.768}{474.836}{1}%
\TextAt{17}{380.919}{472.023}{\Character{40}}%
\TextAt{21}{384.213}{472.023}{\Character{96}}%
\TextAt{3}{387.672}{472.023}{\Character{0}}%
\TextAt{17}{394.258}{472.023}{1\Character{41}}%
\TextAt{3}{401.786}{472.023}{\Character{0}}%
\TextAt{21}{408.372}{472.023}{\Character{96}}%
\TextAt{22}{411.832}{474.836}{\Character{23}}%
\TextAt{23}{415.955}{474.836}{\Character{0}}%
\TextAt{24}{421.711}{474.836}{1}%
\TextAt{9}{429.681}{467.087}{\Character{61}}%
\NativeMathAt{442.107}{467.087}{11.955}{\pm}%
\TextAt{13}{451.405}{467.087}{\Character{96}}%
\TextAt{21}{456.317}{472.023}{s}%
\TextAt{8}{79.226}{438.347}{with}%
\TextAt{13}{105.886}{438.347}{s}%
\TextAt{9}{114.721}{438.347}{\Character{61}}%
\TextAt{13}{127.146}{438.347}{\Character{96}}%
\TextAt{21}{132.059}{442.685}{\Character{23}}%
\TextAt{3}{136.742}{442.685}{\Character{0}}%
\TextAt{17}{143.329}{442.685}{1}%
\TextAt{9}{148.061}{438.347}{\Character{40}}%
\TextAt{13}{152.614}{438.347}{\Character{23}}%
\TextAt{13}{159.142}{438.347}{\Character{96}}%
\TextAt{16}{166.72}{438.347}{\Character{0}}%
\TextAt{13}{178.676}{438.347}{\Character{23}}%
\TextAt{16}{187.87}{438.347}{\Character{0}}%
\TextAt{9}{199.825}{438.347}{1\Character{41}}%
\TextAt{8}{210.231}{438.347}{.}%
\Vector{q 0 G 0 g 0.399 w 452.334 -335.653 m 452.334 -327.583 l S Q}%
\Vector{q 0 G 0 g 0.399 w 452.534 -327.783 m 459.707 -327.783 l S Q}%
\Vector{q 0 G 0 g 0.399 w 452.534 -335.454 m 459.707 -335.454 l S Q}%
\Vector{q 0 G 0 g 0.399 w 459.906 -335.653 m 459.906 -327.583 l S Q}%
\TextAt{8}{97.159}{409.42}{The}%
\TextAt{8}{121.213}{409.42}{ring}%
\TextAt{8}{145.278}{409.42}{of}%
\TextAt{8}{158.608}{409.42}{integers}%
\TextAt{8}{201.922}{409.42}{of}%
\TextAt{18}{215.252}{409.42}{Q}%
\TextAt{9}{224.551}{409.42}{\Character{40}}%
\TextAt{13}{229.103}{409.42}{\Character{16}}%
\TextAt{9}{235.118}{409.42}{\Character{41}}%
\TextAt{8}{243.568}{409.42}{is}%
\TextAt{8}{255.332}{409.42}{now}%
\TextAt{8}{279.721}{409.42}{determined,}%
\TextAt{8}{344.099}{409.42}{for}%
\TextAt{8}{361.984}{409.42}{arbitrary}%
\TextAt{13}{411.754}{409.42}{n}%
\TextAt{8}{418.741}{409.42}{,}%
\TextAt{8}{425.89}{409.42}{as}%
\TextAt{8}{440.26}{409.42}{follows.}%
\TextAt{12}{79.226}{383.482}{Proposition}%
\TextAt{12}{151.41}{383.482}{(10.2).}%
\TextAt{8}{192.998}{383.482}{A}%
\TextAt{18}{205.211}{383.482}{Z}%
\TextAt{8}{213.181}{383.482}{-basis}%
\TextAt{8}{245.361}{383.482}{of}%
\TextAt{8}{258.232}{383.482}{the}%
\TextAt{8}{277.93}{383.482}{ring}%
\TextAt{14}{301.537}{383.482}{o}%
\TextAt{8}{310.82}{383.482}{of}%
\TextAt{8}{323.692}{383.482}{integers}%
\TextAt{8}{366.547}{383.482}{of}%
\TextAt{18}{379.418}{383.482}{Q}%
\TextAt{9}{388.717}{383.482}{\Character{40}}%
\TextAt{13}{393.269}{383.482}{\Character{16}}%
\TextAt{9}{399.285}{383.482}{\Character{41}}%
\TextAt{8}{407.276}{383.482}{is}%
\TextAt{8}{418.581}{383.482}{given}%
\TextAt{8}{448.692}{383.482}{by}%
\TextAt{9}{464.492}{383.482}{1}%
\TextAt{13}{470.345}{383.482}{\Character{59}}%
\TextAt{13}{475.569}{383.482}{\Character{16}}%
\TextAt{13}{481.583}{383.482}{\Character{59}}%
\TextAt{13}{486.807}{383.482}{\Character{58}}%
\TextAt{13}{492.044}{383.482}{\Character{58}}%
\TextAt{13}{497.28}{383.482}{\Character{58}}%
\TextAt{13}{502.505}{383.482}{\Character{59}}%
\TextAt{13}{507.741}{383.482}{\Character{16}}%
\TextAt{21}{513.842}{387.821}{d}%
\TextAt{3}{518.199}{387.821}{\Character{0}}%
\TextAt{17}{524.785}{387.821}{1}%
\TextAt{8}{529.518}{383.482}{,}%
\TextAt{8}{79.226}{366.436}{with}%
\TextAt{13}{105.886}{366.436}{d}%
\TextAt{9}{115.29}{366.436}{\Character{61}}%
\TextAt{13}{127.715}{366.436}{\Character{39}}%
\TextAt{9}{135.384}{366.436}{\Character{40}}%
\TextAt{13}{139.936}{366.436}{n}%
\TextAt{9}{146.923}{366.436}{\Character{41}}%
\TextAt{8}{151.476}{366.436}{,}%
\TextAt{8}{158.625}{366.436}{in}%
\TextAt{8}{172.278}{366.436}{other}%
\TextAt{8}{202.847}{366.436}{words,}%
\TextAt{14}{221.049}{337.696}{o}%
\TextAt{9}{230.214}{337.696}{\Character{61}}%
\TextAt{18}{242.64}{337.696}{Z}%
\TextAt{9}{253.267}{337.696}{\Character{43}}%
\TextAt{18}{265.028}{337.696}{Z}%
\TextAt{13}{272.998}{337.696}{\Character{16}}%
\TextAt{9}{281.67}{337.696}{\Character{43}}%
\TextAt{16}{293.432}{337.696}{\Character{1}}%
\TextAt{16}{298.728}{337.696}{\Character{1}}%
\TextAt{16}{304.036}{337.696}{\Character{1}}%
\TextAt{9}{310.036}{337.696}{\Character{43}}%
\TextAt{18}{321.797}{337.696}{Z}%
\TextAt{13}{329.768}{337.696}{\Character{16}}%
\TextAt{21}{335.783}{342.632}{d}%
\TextAt{3}{340.14}{342.632}{\Character{0}}%
\TextAt{17}{346.727}{342.632}{1}%
\TextAt{9}{354.78}{337.696}{\Character{61}}%
\TextAt{18}{367.206}{337.696}{Z}%
\TextAt{9}{375.176}{337.696}{\Character{91}}%
\TextAt{13}{378.427}{337.696}{\Character{16}}%
\TextAt{9}{384.443}{337.696}{\Character{93}}%
\TextAt{13}{387.695}{337.696}{\Character{58}}%
\TextAt{19}{79.226}{291.172}{Proof.}%
\TextAt{8}{116}{291.172}{We}%
\TextAt{8}{137.436}{291.172}{first}%
\TextAt{8}{162.853}{291.172}{prove}%
\TextAt{8}{195.694}{291.172}{the}%
\TextAt{8}{217.142}{291.172}{proposition}%
\TextAt{8}{279.967}{291.172}{in}%
\TextAt{8}{294.912}{291.172}{the}%
\TextAt{8}{316.36}{291.172}{case}%
\TextAt{8}{342.422}{291.172}{where}%
\TextAt{13}{377.523}{291.172}{n}%
\TextAt{8}{389.699}{291.172}{is}%
\TextAt{8}{402.755}{291.172}{a}%
\TextAt{8}{413.802}{291.172}{prime}%
\TextAt{8}{448.257}{291.172}{power}%
\TextAt{13}{483.692}{291.172}{\Character{96}}%
\TextAt{21}{488.605}{295.51}{\Character{23}}%
\TextAt{8}{493.787}{291.172}{.}%
\TextAt{8}{506.109}{291.172}{Since}%
\TextAt{13}{79.226}{274.126}{d}%
\TextAt{9}{85.309}{274.126}{\Character{40}1}%
\TextAt{13}{95.714}{274.126}{\Character{59}}%
\TextAt{13}{100.938}{274.126}{\Character{16}}%
\TextAt{13}{106.952}{274.126}{\Character{59}}%
\TextAt{13}{112.176}{274.126}{\Character{58}}%
\TextAt{13}{117.413}{274.126}{\Character{58}}%
\TextAt{13}{122.649}{274.126}{\Character{58}}%
\TextAt{13}{127.874}{274.126}{\Character{59}}%
\TextAt{13}{133.11}{274.126}{\Character{16}}%
\TextAt{21}{139.21}{278.464}{d}%
\TextAt{3}{143.568}{278.464}{\Character{0}}%
\TextAt{17}{150.154}{278.464}{1}%
\TextAt{9}{154.886}{274.126}{\Character{41}}%
\TextAt{9}{162.741}{274.126}{\Character{61}}%
\NativeMathAt{175.185}{274.126}{11.955}{\pm}%
\TextAt{13}{184.484}{274.126}{\Character{96}}%
\TextAt{21}{189.396}{278.464}{s}%
\TextAt{8}{193.81}{274.126}{,}%
\TextAt{8}{200.959}{274.126}{(2.9)}%
\TextAt{8}{228.934}{274.126}{gives}%
\TextAt{8}{257.615}{274.126}{us}%
\TextAt{13}{267.68}{245.385}{\Character{96}}%
\TextAt{21}{272.592}{250.321}{s}%
\TextAt{14}{277.006}{245.385}{o}%
\NativeMathAt{286.171}{245.385}{11.955}{\subseteq}%
\TextAt{18}{298.791}{245.385}{Z}%
\TextAt{9}{306.761}{245.385}{\Character{91}}%
\TextAt{13}{310.012}{245.385}{\Character{16}}%
\TextAt{9}{316.028}{245.385}{\Character{93}}%
\NativeMathAt{322.6}{245.385}{11.955}{\subseteq}%
\TextAt{14}{335.22}{245.385}{o}%
\TextAt{13}{341.064}{245.385}{\Character{58}}%
\TextAt{8}{79.226}{216.645}{Putting}%
\TextAt{13}{122.749}{216.645}{\Character{21}}%
\TextAt{9}{133.642}{216.645}{\Character{61}}%
\TextAt{9}{146.805}{216.645}{1}%
\TextAt{16}{155.619}{216.645}{\Character{0}}%
\TextAt{13}{167.872}{216.645}{\Character{16}}%
\TextAt{8}{173.887}{216.645}{,}%
\TextAt{8}{181.582}{216.645}{lemma}%
\TextAt{8}{219.737}{216.645}{(10.1)}%
\TextAt{8}{254.007}{216.645}{tells}%
\TextAt{8}{279.215}{216.645}{us}%
\TextAt{8}{294.667}{216.645}{that}%
\TextAt{14}{320.473}{216.645}{o}%
\TextAt{9}{326.317}{216.645}{\Character{47}}%
\TextAt{13}{332.17}{216.645}{\Character{21}}%
\TextAt{14}{338.998}{216.645}{o}%
\NativeMathAt{348.907}{216.645}{11.955}{\cong}%
\TextAt{18}{362.271}{216.645}{Z}%
\TextAt{9}{370.241}{216.645}{\Character{47}}%
\TextAt{13}{376.094}{216.645}{\Character{96}}%
\TextAt{18}{381.006}{216.645}{Z}%
\TextAt{8}{388.976}{216.645}{,}%
\TextAt{8}{396.671}{216.645}{so}%
\TextAt{8}{411.477}{216.645}{that}%
\TextAt{14}{437.283}{216.645}{o}%
\TextAt{9}{447.192}{216.645}{\Character{61}}%
\TextAt{18}{460.361}{216.645}{Z}%
\TextAt{9}{471.286}{216.645}{\Character{43}}%
\TextAt{13}{483.345}{216.645}{\Character{21}}%
\TextAt{14}{490.173}{216.645}{o}%
\TextAt{8}{496.017}{216.645}{,}%
\TextAt{8}{503.712}{216.645}{and}%
\TextAt{8}{526.911}{216.645}{a}%
\TextAt{8}{79.226}{199.599}{fortiori}%
\TextAt{13}{269.787}{182.553}{\Character{21}}%
\TextAt{14}{276.615}{182.553}{o}%
\TextAt{9}{285.116}{182.553}{\Character{43}}%
\TextAt{18}{296.878}{182.553}{Z}%
\TextAt{9}{304.848}{182.553}{\Character{91}}%
\TextAt{13}{308.099}{182.553}{\Character{16}}%
\TextAt{9}{314.115}{182.553}{\Character{93}}%
\TextAt{9}{320.666}{182.553}{\Character{61}}%
\TextAt{14}{333.113}{182.553}{o}%
\TextAt{13}{338.957}{182.553}{\Character{58}}%
\TextAt{8}{79.226}{159.143}{Multiplying}%
\TextAt{8}{142.959}{159.143}{this}%
\TextAt{8}{165.782}{159.143}{by}%
\TextAt{13}{182.041}{159.143}{\Character{21}}%
\TextAt{8}{192.766}{159.143}{and}%
\TextAt{8}{215.529}{159.143}{substituting}%
\TextAt{8}{280.697}{159.143}{the}%
\TextAt{8}{300.853}{159.143}{result}%
\TextAt{13}{333.431}{159.143}{\Character{21}}%
\TextAt{14}{340.259}{159.143}{o}%
\TextAt{9}{349.425}{159.143}{\Character{61}}%
\TextAt{13}{361.85}{159.143}{\Character{21}}%
\TextAt{17}{368.679}{163.481}{2}%
\TextAt{14}{373.411}{159.143}{o}%
\TextAt{9}{381.912}{159.143}{\Character{43}}%
\TextAt{13}{393.673}{159.143}{\Character{21}}%
\TextAt{18}{400.502}{159.143}{Z}%
\TextAt{9}{408.472}{159.143}{\Character{91}}%
\TextAt{13}{411.723}{159.143}{\Character{16}}%
\TextAt{9}{417.739}{159.143}{\Character{93}}%
\TextAt{8}{420.991}{159.143}{,}%
\TextAt{8}{428.14}{159.143}{we}%
\TextAt{8}{445.367}{159.143}{obtain}%
\TextAt{13}{267.421}{130.402}{\Character{21}}%
\TextAt{17}{274.249}{135.338}{2}%
\TextAt{14}{278.981}{130.402}{o}%
\TextAt{9}{287.482}{130.402}{\Character{43}}%
\TextAt{18}{299.244}{130.402}{Z}%
\TextAt{9}{307.214}{130.402}{\Character{91}}%
\TextAt{13}{310.465}{130.402}{\Character{16}}%
\TextAt{9}{316.481}{130.402}{\Character{93}}%
\TextAt{9}{323.032}{130.402}{\Character{61}}%
\TextAt{14}{335.479}{130.402}{o}%
\TextAt{13}{341.323}{130.402}{\Character{58}}%
\TextAt{8}{79.226}{101.662}{Iterating}%
\TextAt{8}{127.68}{101.662}{this}%
\TextAt{8}{150.503}{101.662}{procedure,}%
\TextAt{8}{208.055}{101.662}{we}%
\TextAt{8}{225.282}{101.662}{find}%
\TextAt{13}{225.971}{72.921}{\Character{21}}%
\TextAt{21}{232.8}{77.858}{t}%
\TextAt{14}{236.356}{72.921}{o}%
\TextAt{9}{244.857}{72.921}{\Character{43}}%
\TextAt{18}{256.618}{72.921}{Z}%
\TextAt{9}{264.588}{72.921}{\Character{91}}%
\TextAt{13}{267.84}{72.921}{\Character{16}}%
\TextAt{9}{273.856}{72.921}{\Character{93}}%
\TextAt{9}{280.407}{72.921}{\Character{61}}%
\TextAt{14}{292.854}{72.921}{o}%
\TextAt{8}{322.608}{72.921}{for}%
\TextAt{8}{340.493}{72.921}{all}%
\TextAt{13}{356.752}{72.921}{t}%
\NativeMathAt{364.3}{72.921}{11.955}{\geq}%
\TextAt{9}{376.92}{72.921}{0}%
\TextAt{13}{382.773}{72.921}{\Character{58}}%
\TextAt{8}{300.14}{43.033}{56}%
\Destination{78.23}{809.815}{nav.page.56}%
\Destination{79.36}{683.255}{nav.print.60}%
\Destination{79.36}{397.455}{nav.number.1.10.2}%
\Hotlink{202.892}{741.855}{6.858}{11.047}{nav.section.1.2}%
\Hotlink{229.551}{741.855}{12.716}{11.047}{nav.print.11}%
\Hotlink{200.459}{271.955}{25.078}{11.047}{nav.number.1.2.9}%
\Hotlink{219.237}{214.455}{30.936}{11.047}{nav.number.1.10.1}%
\Destination{79.226}{89.921}{errata.57}%
}

% Page 69 (body page 57)
\Page{%
\TextAt{2}{79.226}{796.943}{I.}%
\TextAt{2}{89.83}{796.943}{Algebraic}%
\TextAt{2}{138.648}{796.943}{Integers}%
\Vector{q 0 G 0 g 0.299 w 7.226 18.729 m 460.769 18.729 l S Q}%
\TextAt{8}{79.226}{761.124}{For}%
\TextAt{13}{100.196}{761.124}{t}%
\TextAt{9}{107.744}{761.124}{\Character{61}}%
\TextAt{13}{120.169}{761.124}{s\Character{39}}%
\TextAt{9}{133.351}{761.124}{\Character{40}}%
\TextAt{13}{137.904}{761.124}{\Character{96}}%
\TextAt{21}{142.816}{765.462}{\Character{23}}%
\TextAt{9}{147.998}{761.124}{\Character{41}}%
\TextAt{8}{156.448}{761.124}{this}%
\TextAt{8}{179.27}{761.124}{implies,}%
\TextAt{8}{222.249}{761.124}{in}%
\TextAt{8}{235.902}{761.124}{view}%
\TextAt{8}{262.884}{761.124}{of}%
\TextAt{13}{276.214}{761.124}{\Character{96}}%
\TextAt{14}{281.127}{761.124}{o}%
\TextAt{9}{290.292}{761.124}{\Character{61}}%
\TextAt{13}{302.717}{761.124}{\Character{21}}%
\TextAt{21}{309.546}{765.462}{\Character{39}}%
\TextAt{17}{315.064}{765.462}{\Character{40}}%
\TextAt{21}{318.357}{765.462}{\Character{96}}%
\TextAt{22}{321.816}{768.275}{\Character{23}}%
\TextAt{17}{326.438}{765.462}{\Character{41}}%
\TextAt{14}{330.229}{761.124}{o}%
\TextAt{8}{339.971}{761.124}{(see}%
\TextAt{8}{363.439}{761.124}{(10.1)),}%
\TextAt{8}{405.079}{761.124}{that}%
\TextAt{14}{216.979}{732.123}{o}%
\TextAt{9}{226.144}{732.123}{\Character{61}}%
\TextAt{13}{238.57}{732.123}{\Character{21}}%
\TextAt{21}{245.398}{737.059}{t}%
\TextAt{14}{248.954}{732.123}{o}%
\TextAt{9}{257.455}{732.123}{\Character{43}}%
\TextAt{18}{269.217}{732.123}{Z}%
\TextAt{9}{277.187}{732.123}{\Character{91}}%
\TextAt{13}{280.438}{732.123}{\Character{16}}%
\TextAt{9}{286.454}{732.123}{\Character{93}}%
\TextAt{9}{293.005}{732.123}{\Character{61}}%
\TextAt{13}{305.452}{732.123}{\Character{96}}%
\TextAt{21}{310.364}{737.059}{s}%
\TextAt{14}{314.778}{732.123}{o}%
\TextAt{9}{323.279}{732.123}{\Character{43}}%
\TextAt{18}{335.04}{732.123}{Z}%
\TextAt{9}{343.01}{732.123}{\Character{91}}%
\TextAt{13}{346.262}{732.123}{\Character{16}}%
\TextAt{9}{352.278}{732.123}{\Character{93}}%
\TextAt{9}{358.829}{732.123}{\Character{61}}%
\TextAt{18}{371.276}{732.123}{Z}%
\TextAt{9}{379.246}{732.123}{\Character{91}}%
\TextAt{13}{382.498}{732.123}{\Character{16}}%
\TextAt{9}{388.513}{732.123}{\Character{93}}%
\TextAt{13}{391.765}{732.123}{\Character{58}}%
\TextAt{8}{79.226}{703.121}{In}%
\TextAt{8}{93.892}{703.121}{the}%
\TextAt{8}{114.092}{703.121}{general}%
\TextAt{8}{154.461}{703.121}{case,}%
\TextAt{8}{182.54}{703.121}{let}%
\TextAt{13}{199.488}{703.121}{n}%
\TextAt{9}{209.872}{703.121}{\Character{61}}%
\TextAt{13}{222.373}{703.121}{\Character{96}}%
\TextAt{21}{227.285}{708.346}{\Character{23}}%
\TextAt{24}{231.448}{707.239}{1}%
\TextAt{17}{227.285}{700.19}{1}%
\TextAt{16}{238.09}{703.121}{\Character{1}}%
\TextAt{16}{243.386}{703.121}{\Character{1}}%
\TextAt{16}{248.694}{703.121}{\Character{1}}%
\TextAt{13}{254.03}{703.121}{\Character{96}}%
\TextAt{21}{258.942}{707.46}{\Character{23}}%
\TextAt{22}{263.105}{706.464}{r}%
\TextAt{21}{258.942}{700.166}{r}%
\TextAt{8}{267.761}{703.121}{.}%
\TextAt{8}{276.342}{703.121}{Then}%
\TextAt{13}{306.944}{703.121}{\Character{16}}%
\TextAt{21}{312.056}{701.328}{i}%
\TextAt{9}{318.834}{703.121}{\Character{61}}%
\TextAt{13}{331.335}{703.121}{\Character{16}}%
\TextAt{21}{337.35}{707.46}{n}%
\TextAt{17}{342.489}{707.46}{\Character{47}}%
\TextAt{21}{346.723}{707.46}{\Character{96}}%
\TextAt{22}{350.182}{712.107}{\Character{23}}%
\TextAt{22}{353.897}{710.205}{i}%
\TextAt{22}{350.182}{704.81}{i}%
\TextAt{8}{361.997}{703.121}{is}%
\TextAt{8}{373.805}{703.121}{a}%
\TextAt{8}{383.605}{703.121}{primitive}%
\TextAt{13}{433.729}{703.121}{\Character{96}}%
\TextAt{21}{438.641}{708.454}{\Character{23}}%
\TextAt{22}{442.804}{707.239}{i}%
\TextAt{21}{438.641}{700.058}{i}%
\TextAt{8}{446.464}{703.121}{-th}%
\TextAt{8}{465.361}{703.121}{root}%
\TextAt{8}{490.452}{703.121}{of}%
\TextAt{8}{503.826}{703.121}{unity,}%
\TextAt{8}{79.226}{686.075}{and}%
\TextAt{8}{101.989}{686.075}{one}%
\TextAt{8}{123.448}{686.075}{has}%
\TextAt{18}{248.79}{669.029}{Q}%
\TextAt{9}{258.088}{669.029}{\Character{40}}%
\TextAt{13}{262.64}{669.029}{\Character{16}}%
\TextAt{9}{268.656}{669.029}{\Character{41}}%
\TextAt{9}{276.511}{669.029}{\Character{61}}%
\TextAt{18}{288.955}{669.029}{Q}%
\TextAt{9}{298.253}{669.029}{\Character{40}}%
\TextAt{13}{302.805}{669.029}{\Character{16}}%
\TextAt{17}{307.918}{667.236}{1}%
\TextAt{9}{312.65}{669.029}{\Character{41}}%
\TextAt{16}{319.195}{669.029}{\Character{1}}%
\TextAt{16}{324.491}{669.029}{\Character{1}}%
\TextAt{16}{329.799}{669.029}{\Character{1}}%
\TextAt{18}{335.135}{669.029}{Q}%
\TextAt{9}{344.434}{669.029}{\Character{40}}%
\TextAt{13}{348.986}{669.029}{\Character{16}}%
\TextAt{21}{354.099}{667.236}{r}%
\TextAt{9}{358.654}{669.029}{\Character{41}}%
\TextAt{8}{79.226}{645.508}{and}%
\TextAt{18}{102.196}{645.508}{Q}%
\TextAt{9}{111.495}{645.508}{\Character{40}}%
\TextAt{13}{116.047}{645.508}{\Character{16}}%
\TextAt{17}{121.16}{643.715}{1}%
\TextAt{9}{125.892}{645.508}{\Character{41}}%
\TextAt{16}{132.437}{645.508}{\Character{1}}%
\TextAt{16}{137.733}{645.508}{\Character{1}}%
\TextAt{16}{143.041}{645.508}{\Character{1}}%
\TextAt{18}{148.377}{645.508}{Q}%
\TextAt{9}{157.676}{645.508}{\Character{40}}%
\TextAt{13}{162.228}{645.508}{\Character{16}}%
\TextAt{21}{167.341}{643.715}{i}%
\TextAt{3}{170.224}{643.715}{\Character{0}}%
\TextAt{17}{176.81}{643.715}{1}%
\TextAt{9}{181.543}{645.508}{\Character{41}}%
\TextAt{16}{188.893}{645.508}{\Character{92}}%
\TextAt{18}{199.661}{645.508}{Q}%
\TextAt{9}{208.96}{645.508}{\Character{40}}%
\TextAt{13}{213.512}{645.508}{\Character{16}}%
\TextAt{21}{218.625}{643.715}{i}%
\TextAt{9}{222.006}{645.508}{\Character{41}}%
\TextAt{9}{230.219}{645.508}{\Character{61}}%
\TextAt{18}{243.012}{645.508}{Q}%
\TextAt{8}{252.311}{645.508}{.}%
\TextAt{8}{261.381}{645.508}{By}%
\TextAt{8}{279.952}{645.508}{what}%
\TextAt{8}{309.426}{645.508}{we}%
\TextAt{8}{326.861}{645.508}{have}%
\TextAt{8}{354.063}{645.508}{just}%
\TextAt{8}{377.416}{645.508}{seen,}%
\TextAt{8}{406.344}{645.508}{for}%
\TextAt{8}{424.436}{645.508}{each}%
\TextAt{13}{450.981}{645.508}{i}%
\TextAt{9}{458.649}{645.508}{\Character{61}}%
\TextAt{9}{471.417}{645.508}{1}%
\TextAt{13}{477.281}{645.508}{\Character{59}}%
\TextAt{13}{482.505}{645.508}{\Character{58}}%
\TextAt{13}{487.742}{645.508}{\Character{58}}%
\TextAt{13}{492.966}{645.508}{\Character{58}}%
\TextAt{13}{498.203}{645.508}{\Character{59}}%
\TextAt{13}{503.439}{645.508}{r}%
\TextAt{8}{509.102}{645.508}{,}%
\TextAt{8}{516.511}{645.508}{the}%
\TextAt{8}{79.226}{628.462}{elements}%
\TextAt{9}{128.24}{628.462}{1}%
\TextAt{13}{134.093}{628.462}{\Character{59}}%
\TextAt{13}{139.317}{628.462}{\Character{16}}%
\TextAt{21}{144.45}{626.669}{i}%
\TextAt{13}{147.831}{628.462}{\Character{59}}%
\TextAt{13}{153.055}{628.462}{\Character{58}}%
\TextAt{13}{158.292}{628.462}{\Character{58}}%
\TextAt{13}{163.516}{628.462}{\Character{58}}%
\TextAt{13}{168.753}{628.462}{\Character{59}}%
\TextAt{13}{173.989}{628.462}{\Character{16}}%
\TextAt{21}{180.067}{633.795}{d}%
\TextAt{22}{184.425}{632.58}{i}%
\TextAt{3}{187.586}{633.795}{\Character{0}}%
\TextAt{17}{194.173}{633.795}{1}%
\TextAt{21}{179.164}{625.398}{i}%
\TextAt{8}{198.905}{628.462}{,}%
\TextAt{8}{207.501}{628.462}{where}%
\TextAt{13}{242.467}{628.462}{d}%
\TextAt{21}{248.55}{626.669}{i}%
\TextAt{9}{257.224}{628.462}{\Character{61}}%
\TextAt{13}{271.621}{628.462}{\Character{39}}%
\TextAt{9}{279.29}{628.462}{\Character{40}}%
\TextAt{13}{283.842}{628.462}{\Character{96}}%
\TextAt{21}{288.754}{633.795}{\Character{23}}%
\TextAt{22}{292.917}{632.58}{i}%
\TextAt{21}{288.754}{625.398}{i}%
\TextAt{9}{296.577}{628.462}{\Character{41}}%
\TextAt{8}{301.129}{628.462}{,}%
\TextAt{8}{309.724}{628.462}{form}%
\TextAt{8}{338.522}{628.462}{an}%
\TextAt{8}{355.938}{628.462}{integral}%
\TextAt{8}{399.703}{628.462}{basis}%
\TextAt{8}{429.601}{628.462}{of}%
\TextAt{18}{444.088}{628.462}{Q}%
\TextAt{9}{453.386}{628.462}{\Character{40}}%
\TextAt{13}{457.939}{628.462}{\Character{16}}%
\TextAt{21}{463.051}{626.669}{i}%
\TextAt{9}{466.433}{628.462}{\Character{41}}%
\TextAt{16}{476.278}{628.462}{\Character{106}}%
\TextAt{18}{484.891}{628.462}{Q}%
\TextAt{8}{494.19}{628.462}{.}%
\TextAt{8}{506.109}{628.462}{Since}%
\TextAt{8}{79.226}{611.416}{the}%
\TextAt{8}{101.106}{611.416}{discriminants}%
\TextAt{13}{174.823}{611.416}{d}%
\TextAt{9}{180.906}{611.416}{\Character{40}1}%
\TextAt{13}{191.311}{611.416}{\Character{59}}%
\TextAt{13}{196.535}{611.416}{\Character{16}}%
\TextAt{21}{201.668}{609.623}{i}%
\TextAt{13}{205.049}{611.416}{\Character{59}}%
\TextAt{13}{210.273}{611.416}{\Character{58}}%
\TextAt{13}{215.51}{611.416}{\Character{58}}%
\TextAt{13}{220.734}{611.416}{\Character{58}}%
\TextAt{13}{225.971}{611.416}{\Character{59}}%
\TextAt{13}{231.207}{611.416}{\Character{16}}%
\TextAt{21}{237.285}{616.749}{d}%
\TextAt{22}{241.643}{615.534}{i}%
\TextAt{3}{244.804}{616.749}{\Character{0}}%
\TextAt{17}{251.391}{616.749}{1}%
\TextAt{21}{236.383}{608.352}{i}%
\TextAt{9}{256.123}{611.416}{\Character{41}}%
\TextAt{9}{266.918}{611.416}{\Character{61}}%
\NativeMathAt{282.295}{611.416}{11.955}{\pm}%
\TextAt{13}{291.593}{611.416}{\Character{96}}%
\TextAt{21}{296.505}{616.749}{s}%
\TextAt{22}{300.421}{615.534}{i}%
\TextAt{21}{296.505}{608.352}{i}%
\TextAt{8}{309.702}{611.416}{are}%
\TextAt{8}{330.936}{611.416}{pairwise}%
\TextAt{8}{378.244}{611.416}{relatively}%
\TextAt{8}{431.027}{611.416}{prime,}%
\TextAt{8}{469.597}{611.416}{we}%
\TextAt{8}{488.547}{611.416}{conclude}%
\TextAt{8}{79.226}{594.37}{successively}%
\TextAt{8}{144.259}{594.37}{from}%
\TextAt{8}{173.343}{594.37}{(2.11)}%
\TextAt{8}{208.62}{594.37}{that}%
\TextAt{8}{235.433}{594.37}{the}%
\TextAt{8}{257.033}{594.37}{elements}%
\TextAt{13}{306.333}{594.37}{\Character{16}}%
\TextAt{21}{312.348}{600.037}{j}%
\TextAt{24}{315.773}{598.931}{1}%
\TextAt{17}{311.445}{591.439}{1}%
\TextAt{16}{322.415}{594.37}{\Character{1}}%
\TextAt{16}{327.711}{594.37}{\Character{1}}%
\TextAt{16}{333.019}{594.37}{\Character{1}}%
\TextAt{13}{338.355}{594.37}{\Character{16}}%
\TextAt{21}{344.37}{598.708}{j}%
\TextAt{22}{347.795}{597.712}{r}%
\TextAt{21}{343.468}{591.414}{r}%
\TextAt{8}{352.451}{594.37}{,}%
\TextAt{8}{361.405}{594.37}{with}%
\TextAt{13}{389.509}{594.37}{j}%
\TextAt{21}{394.352}{592.577}{i}%
\TextAt{9}{403.514}{594.37}{\Character{61}}%
\TextAt{9}{418.386}{594.37}{0}%
\TextAt{13}{424.252}{594.37}{\Character{59}}%
\TextAt{13}{429.476}{594.37}{\Character{58}}%
\TextAt{13}{434.713}{594.37}{\Character{58}}%
\TextAt{13}{439.949}{594.37}{\Character{58}}%
\TextAt{13}{445.174}{594.37}{\Character{59}}%
\TextAt{13}{450.41}{594.37}{d}%
\TextAt{21}{456.556}{592.577}{i}%
\TextAt{16}{463.578}{594.37}{\Character{0}}%
\TextAt{9}{476.517}{594.37}{1}%
\TextAt{8}{482.37}{594.37}{,}%
\TextAt{8}{491.324}{594.37}{form}%
\TextAt{8}{520.408}{594.37}{an}%
\TextAt{8}{79.226}{577.324}{integral}%
\TextAt{8}{121.227}{577.324}{basis}%
\TextAt{8}{149.36}{577.324}{of}%
\TextAt{18}{162.082}{577.324}{Q}%
\TextAt{9}{171.381}{577.324}{\Character{40}}%
\TextAt{13}{175.933}{577.324}{\Character{16}}%
\TextAt{9}{181.949}{577.324}{\Character{41}}%
\TextAt{16}{189.822}{577.324}{\Character{106}}%
\TextAt{18}{196.464}{577.324}{Q}%
\TextAt{8}{205.762}{577.324}{.}%
\TextAt{8}{214.008}{577.324}{But}%
\TextAt{8}{236.642}{577.324}{each}%
\TextAt{8}{262.371}{577.324}{one}%
\TextAt{8}{283.224}{577.324}{of}%
\TextAt{8}{295.946}{577.324}{these}%
\TextAt{8}{325.31}{577.324}{elements}%
\TextAt{8}{372.56}{577.324}{is}%
\TextAt{8}{383.716}{577.324}{a}%
\TextAt{8}{392.864}{577.324}{power}%
\TextAt{8}{426.401}{577.324}{of}%
\TextAt{13}{439.123}{577.324}{\Character{16}}%
\TextAt{8}{445.139}{577.324}{.}%
\TextAt{8}{453.385}{577.324}{Therefore}%
\TextAt{8}{505.775}{577.324}{every}%
\TextAt{13}{79.226}{560.278}{\Character{11}}%
\NativeMathAt{90.069}{560.278}{11.955}{\in}%
\TextAt{14}{101.36}{560.278}{o}%
\TextAt{8}{110.796}{560.278}{may}%
\TextAt{8}{135.859}{560.278}{be}%
\TextAt{8}{151.478}{560.278}{written}%
\TextAt{8}{192.143}{560.278}{as}%
\TextAt{8}{206.208}{560.278}{a}%
\TextAt{8}{215.658}{560.278}{polynomial}%
\TextAt{13}{275.846}{560.278}{\Character{11}}%
\TextAt{9}{286.688}{560.278}{\Character{61}}%
\TextAt{13}{299.114}{560.278}{f}%
\TextAt{9}{306.16}{560.278}{\Character{40}}%
\TextAt{13}{310.712}{560.278}{\Character{16}}%
\TextAt{9}{316.728}{560.278}{\Character{41}}%
\TextAt{8}{324.872}{560.278}{with}%
\TextAt{8}{351.227}{560.278}{coefficients}%
\TextAt{8}{410.159}{560.278}{in}%
\TextAt{18}{423.507}{560.278}{Z}%
\TextAt{8}{431.477}{560.278}{.}%
\TextAt{8}{439.824}{560.278}{Since}%
\TextAt{13}{470.076}{560.278}{\Character{16}}%
\TextAt{8}{479.683}{560.278}{has}%
\TextAt{8}{500.251}{560.278}{degree}%
\TextAt{13}{79.226}{543.232}{\Character{39}}%
\TextAt{9}{86.895}{543.232}{\Character{40}}%
\TextAt{13}{91.447}{543.232}{n}%
\TextAt{9}{98.434}{543.232}{\Character{41}}%
\TextAt{8}{106.504}{543.232}{over}%
\TextAt{18}{131.169}{543.232}{Q}%
\TextAt{8}{140.468}{543.232}{,}%
\TextAt{8}{147.313}{543.232}{the}%
\TextAt{8}{167.089}{543.232}{degree}%
\TextAt{8}{203.124}{543.232}{of}%
\TextAt{8}{216.073}{543.232}{the}%
\TextAt{8}{235.849}{543.232}{polynomial}%
\TextAt{13}{295.962}{543.232}{f}%
\TextAt{8}{306.525}{543.232}{may}%
\TextAt{8}{331.513}{543.232}{be}%
\TextAt{8}{347.057}{543.232}{reduced}%
\TextAt{8}{390.241}{543.232}{to}%
\TextAt{13}{404.171}{543.232}{\Character{39}}%
\TextAt{9}{411.84}{543.232}{\Character{40}}%
\TextAt{13}{416.392}{543.232}{n}%
\TextAt{9}{423.38}{543.232}{\Character{41}}%
\TextAt{16}{429.811}{543.232}{\Character{0}}%
\TextAt{9}{440.988}{543.232}{1}%
\TextAt{8}{446.841}{543.232}{.}%
\TextAt{8}{455.162}{543.232}{In}%
\TextAt{8}{469.403}{543.232}{this}%
\TextAt{8}{491.845}{543.232}{way}%
\TextAt{8}{515.207}{543.232}{one}%
\TextAt{8}{79.226}{526.186}{obtains}%
\TextAt{8}{120.268}{526.186}{a}%
\TextAt{8}{130.024}{526.186}{representation}%
\TextAt{13}{215.61}{497.185}{\Character{11}}%
\TextAt{9}{226.453}{497.185}{\Character{61}}%
\TextAt{13}{238.878}{497.185}{a}%
\TextAt{17}{245.023}{495.392}{0}%
\TextAt{9}{252.412}{497.185}{\Character{43}}%
\TextAt{13}{264.173}{497.185}{a}%
\TextAt{17}{270.318}{495.392}{1}%
\TextAt{13}{275.051}{497.185}{\Character{16}}%
\TextAt{9}{283.723}{497.185}{\Character{43}}%
\TextAt{16}{295.484}{497.185}{\Character{1}}%
\TextAt{16}{300.78}{497.185}{\Character{1}}%
\TextAt{16}{306.088}{497.185}{\Character{1}}%
\TextAt{9}{312.088}{497.185}{\Character{43}}%
\TextAt{13}{323.85}{497.185}{a}%
\TextAt{21}{329.995}{495.325}{\Character{39}}%
\TextAt{17}{335.513}{495.325}{\Character{40}}%
\TextAt{21}{338.806}{495.325}{n}%
\TextAt{17}{343.944}{495.325}{\Character{41}}%
\TextAt{3}{347.238}{495.325}{\Character{0}}%
\TextAt{17}{353.824}{495.325}{1}%
\TextAt{13}{358.557}{497.185}{\Character{16}}%
\TextAt{21}{364.572}{502.121}{\Character{39}}%
\TextAt{17}{370.09}{502.121}{\Character{40}}%
\TextAt{21}{373.384}{502.121}{n}%
\TextAt{17}{378.522}{502.121}{\Character{41}}%
\TextAt{3}{381.815}{502.121}{\Character{0}}%
\TextAt{17}{388.402}{502.121}{1}%
\TextAt{13}{393.134}{497.185}{\Character{58}}%
\TextAt{8}{79.226}{468.184}{Thus}%
\TextAt{9}{108.875}{468.184}{1}%
\TextAt{13}{114.728}{468.184}{\Character{59}}%
\TextAt{13}{119.952}{468.184}{\Character{16}}%
\TextAt{13}{125.966}{468.184}{\Character{59}}%
\TextAt{13}{131.19}{468.184}{\Character{58}}%
\TextAt{13}{136.427}{468.184}{\Character{58}}%
\TextAt{13}{141.663}{468.184}{\Character{58}}%
\TextAt{13}{146.888}{468.184}{\Character{59}}%
\TextAt{13}{152.124}{468.184}{\Character{16}}%
\TextAt{21}{158.224}{472.522}{\Character{39}}%
\TextAt{17}{163.742}{472.522}{\Character{40}}%
\TextAt{21}{167.036}{472.522}{n}%
\TextAt{17}{172.174}{472.522}{\Character{41}}%
\TextAt{3}{175.467}{472.522}{\Character{0}}%
\TextAt{17}{182.054}{472.522}{1}%
\TextAt{8}{190.683}{468.184}{is}%
\TextAt{8}{202.447}{468.184}{indeed}%
\TextAt{8}{239.508}{468.184}{an}%
\TextAt{8}{255.767}{468.184}{integral}%
\TextAt{8}{298.375}{468.184}{basis.}%
\Vector{q 0 G 0 g 0.399 w 452.334 -305.816 m 452.334 -297.747 l S Q}%
\Vector{q 0 G 0 g 0.399 w 452.534 -297.946 m 459.707 -297.946 l S Q}%
\Vector{q 0 G 0 g 0.399 w 452.534 -305.617 m 459.707 -305.617 l S Q}%
\Vector{q 0 G 0 g 0.399 w 459.906 -305.816 m 459.906 -297.747 l S Q}%
\TextAt{8}{97.159}{439.183}{Knowing}%
\TextAt{8}{146.114}{439.183}{that}%
\TextAt{18}{171.339}{439.183}{Z}%
\TextAt{9}{179.309}{439.183}{\Character{91}}%
\TextAt{13}{182.561}{439.183}{\Character{16}}%
\TextAt{9}{188.576}{439.183}{\Character{93}}%
\TextAt{8}{195.581}{439.183}{is}%
\TextAt{8}{207.201}{439.183}{the}%
\TextAt{8}{227.213}{439.183}{ring}%
\TextAt{8}{251.134}{439.183}{of}%
\TextAt{8}{264.32}{439.183}{integers}%
\TextAt{8}{307.489}{439.183}{of}%
\TextAt{8}{320.675}{439.183}{the}%
\TextAt{8}{340.687}{439.183}{field}%
\TextAt{18}{365.899}{439.183}{Q}%
\TextAt{9}{375.198}{439.183}{\Character{40}}%
\TextAt{13}{379.75}{439.183}{\Character{16}}%
\TextAt{9}{385.766}{439.183}{\Character{41}}%
\TextAt{8}{394.071}{439.183}{we}%
\TextAt{8}{411.154}{439.183}{are}%
\TextAt{8}{430.521}{439.183}{now}%
\TextAt{8}{454.765}{439.183}{in}%
\TextAt{8}{468.273}{439.183}{a}%
\TextAt{8}{477.884}{439.183}{position}%
\TextAt{8}{522.357}{439.183}{to}%
\TextAt{8}{79.226}{422.137}{state}%
\TextAt{8}{107.519}{422.137}{explicitly}%
\TextAt{8}{157.857}{422.137}{the}%
\TextAt{8}{177.626}{422.137}{law}%
\TextAt{8}{198.375}{422.137}{of}%
\TextAt{8}{211.318}{422.137}{decomposition}%
\TextAt{8}{288.065}{422.137}{of}%
\TextAt{8}{301.008}{422.137}{prime}%
\TextAt{8}{333.784}{422.137}{numbers}%
\TextAt{13}{380.607}{422.137}{p}%
\TextAt{8}{389.992}{422.137}{into}%
\TextAt{8}{413.348}{422.137}{prime}%
\TextAt{8}{446.124}{422.137}{ideals}%
\TextAt{8}{478.314}{422.137}{of}%
\TextAt{18}{491.256}{422.137}{Q}%
\TextAt{9}{500.555}{422.137}{\Character{40}}%
\TextAt{13}{505.107}{422.137}{\Character{16}}%
\TextAt{9}{511.123}{422.137}{\Character{41}}%
\TextAt{8}{515.675}{422.137}{.}%
\TextAt{8}{523.994}{422.137}{It}%
\TextAt{8}{79.226}{405.091}{is}%
\TextAt{8}{90.99}{405.091}{of}%
\TextAt{8}{104.32}{405.091}{the}%
\TextAt{8}{124.476}{405.091}{most}%
\TextAt{8}{153.157}{405.091}{beautiful}%
\TextAt{8}{202.58}{405.091}{simplicity.}%
\TextAt{12}{79.226}{379.078}{Proposition}%
\TextAt{12}{151.295}{379.078}{(10.3).}%
\TextAt{8}{192.784}{379.078}{Let}%
\TextAt{13}{213.195}{379.078}{n}%
\TextAt{9}{223.503}{379.078}{\Character{61}}%
\TextAt{20}{235.929}{388.045}{\Character{81}}%
\TextAt{21}{247.22}{375.591}{p}%
\TextAt{13}{253.973}{379.078}{p}%
\TextAt{21}{259.848}{383.417}{\Character{23}}%
\TextAt{22}{264.011}{382.42}{p}%
\TextAt{8}{272.187}{379.078}{be}%
\TextAt{8}{287.553}{379.078}{the}%
\TextAt{8}{307.151}{379.078}{prime}%
\TextAt{8}{339.757}{379.078}{factorization}%
\TextAt{8}{407.176}{379.078}{of}%
\TextAt{13}{419.948}{379.078}{n}%
\TextAt{8}{430.275}{379.078}{and,}%
\TextAt{8}{455.842}{379.078}{for}%
\TextAt{8}{473.169}{379.078}{every}%
\TextAt{8}{503.503}{379.078}{prime}%
\TextAt{8}{79.226}{362.032}{number}%
\TextAt{13}{121.822}{362.032}{p}%
\TextAt{8}{127.698}{362.032}{,}%
\TextAt{8}{134.847}{362.032}{let}%
\TextAt{13}{151.751}{362.032}{f}%
\TextAt{21}{157.521}{360.239}{p}%
\TextAt{8}{166.18}{362.032}{be}%
\TextAt{8}{182.104}{362.032}{the}%
\TextAt{8}{202.26}{362.032}{smallest}%
\TextAt{8}{247.26}{362.032}{positive}%
\TextAt{8}{290.573}{362.032}{integer}%
\TextAt{8}{329.272}{362.032}{such}%
\TextAt{8}{355.669}{362.032}{that}%
\TextAt{13}{248.647}{333.031}{p}%
\TextAt{21}{254.522}{337.967}{f}%
\TextAt{22}{258.59}{336.971}{p}%
\NativeMathAt{266.748}{333.031}{11.955}{\equiv}%
\TextAt{9}{279.367}{333.031}{1}%
\TextAt{9}{297.156}{333.031}{\Character{40}}%
\TextAt{8}{301.727}{333.031}{mod}%
\TextAt{13}{327.83}{333.031}{n}%
\TextAt{9}{334.817}{333.031}{\Character{47}}%
\TextAt{13}{340.67}{333.031}{p}%
\TextAt{21}{346.545}{337.967}{\Character{23}}%
\TextAt{22}{350.708}{336.971}{p}%
\TextAt{9}{355.545}{333.031}{\Character{41}}%
\TextAt{13}{360.097}{333.031}{\Character{58}}%
\TextAt{8}{79.226}{304.03}{Then}%
\TextAt{8}{109.784}{304.03}{one}%
\TextAt{8}{131.243}{304.03}{has}%
\TextAt{8}{152.117}{304.03}{in}%
\TextAt{18}{165.77}{304.03}{Q}%
\TextAt{9}{175.068}{304.03}{\Character{40}}%
\TextAt{13}{179.62}{304.03}{\Character{16}}%
\TextAt{9}{185.636}{304.03}{\Character{41}}%
\TextAt{8}{194.086}{304.03}{the}%
\TextAt{8}{214.242}{304.03}{factorization}%
\TextAt{13}{256.707}{275.029}{p}%
\TextAt{9}{265.903}{275.029}{\Character{61}}%
\TextAt{9}{278.312}{275.029}{\Character{40}}%
\TextAt{14}{282.881}{275.029}{p}%
\TextAt{17}{288.864}{273.236}{1}%
\TextAt{16}{295.589}{275.029}{\Character{1}}%
\TextAt{16}{300.885}{275.029}{\Character{1}}%
\TextAt{16}{306.193}{275.029}{\Character{1}}%
\TextAt{14}{311.529}{275.029}{p}%
\TextAt{21}{317.512}{273.236}{r}%
\TextAt{9}{322.067}{275.029}{\Character{41}}%
\TextAt{21}{326.619}{279.965}{\Character{39}}%
\TextAt{17}{332.137}{279.965}{\Character{40}}%
\TextAt{21}{335.431}{279.965}{p}%
\TextAt{22}{339.693}{282.982}{\Character{23}}%
\TextAt{22}{343.409}{281.985}{p}%
\TextAt{17}{348.245}{279.965}{\Character{41}}%
\TextAt{13}{352.037}{275.029}{\Character{59}}%
\TextAt{8}{79.226}{246.028}{where}%
\TextAt{14}{113.035}{246.028}{p}%
\TextAt{17}{119.018}{244.234}{1}%
\TextAt{13}{123.751}{246.028}{\Character{59}}%
\TextAt{13}{128.975}{246.028}{\Character{58}}%
\TextAt{13}{134.212}{246.028}{\Character{58}}%
\TextAt{13}{139.436}{246.028}{\Character{58}}%
\TextAt{13}{144.673}{246.028}{\Character{59}}%
\TextAt{14}{149.972}{246.028}{p}%
\TextAt{21}{155.955}{244.234}{r}%
\TextAt{8}{164.407}{246.028}{are}%
\TextAt{8}{183.918}{246.028}{distinct}%
\TextAt{8}{226.251}{246.028}{prime}%
\TextAt{8}{259.414}{246.028}{ideals,}%
\TextAt{8}{295.244}{246.028}{all}%
\TextAt{8}{311.503}{246.028}{of}%
\TextAt{8}{324.833}{246.028}{degree}%
\TextAt{13}{361.249}{246.028}{f}%
\TextAt{21}{367.019}{244.234}{p}%
\TextAt{8}{371.78}{246.028}{.}%
\TextAt{19}{79.226}{211.049}{Proof.}%
\TextAt{8}{116}{211.049}{Since}%
\TextAt{14}{147.864}{211.049}{o}%
\TextAt{9}{159.254}{211.049}{\Character{61}}%
\TextAt{18}{173.905}{211.049}{Z}%
\TextAt{9}{181.875}{211.049}{\Character{91}}%
\TextAt{13}{185.127}{211.049}{\Character{16}}%
\TextAt{9}{191.143}{211.049}{\Character{93}}%
\TextAt{8}{194.394}{211.049}{,}%
\TextAt{8}{203.176}{211.049}{the}%
\TextAt{8}{224.638}{211.049}{conductor}%
\TextAt{8}{280.58}{211.049}{of}%
\TextAt{18}{295.216}{211.049}{Z}%
\TextAt{9}{303.186}{211.049}{\Character{91}}%
\TextAt{13}{306.437}{211.049}{\Character{16}}%
\TextAt{9}{312.453}{211.049}{\Character{93}}%
\TextAt{8}{320.908}{211.049}{equals}%
\TextAt{9}{357.721}{211.049}{1}%
\TextAt{8}{363.574}{211.049}{,}%
\TextAt{8}{372.356}{211.049}{and}%
\TextAt{8}{396.424}{211.049}{we}%
\TextAt{8}{414.958}{211.049}{may}%
\TextAt{8}{441.632}{211.049}{apply}%
\TextAt{8}{475.134}{211.049}{proposition}%
\TextAt{8}{79.226}{194.003}{(8.3)}%
\TextAt{8}{107.996}{194.003}{to}%
\TextAt{8}{123.101}{194.003}{any}%
\TextAt{8}{146.013}{194.003}{prime}%
\TextAt{8}{179.972}{194.003}{number}%
\TextAt{13}{223.363}{194.003}{p}%
\TextAt{8}{229.238}{194.003}{.}%
\TextAt{8}{240.071}{194.003}{As}%
\TextAt{8}{258.153}{194.003}{a}%
\TextAt{8}{268.704}{194.003}{consequence,}%
\TextAt{8}{339.013}{194.003}{every}%
\TextAt{13}{370.7}{194.003}{p}%
\TextAt{8}{381.268}{194.003}{decomposes}%
\TextAt{8}{445.593}{194.003}{into}%
\TextAt{8}{470.13}{194.003}{prime}%
\TextAt{8}{504.089}{194.003}{ideals}%
\TextAt{8}{79.226}{176.957}{in}%
\TextAt{8}{93.629}{176.957}{exactly}%
\TextAt{8}{134.703}{176.957}{the}%
\TextAt{8}{155.61}{176.957}{same}%
\TextAt{8}{185.685}{176.957}{way}%
\TextAt{8}{210.178}{176.957}{as}%
\TextAt{8}{225.298}{176.957}{the}%
\TextAt{8}{246.204}{176.957}{minimal}%
\TextAt{8}{292.479}{176.957}{polynomial}%
\TextAt{13}{353.722}{176.957}{\Character{30}}%
\TextAt{21}{360.65}{175.164}{n}%
\TextAt{9}{366.287}{176.957}{\Character{40}}%
\TextAt{13}{370.839}{176.957}{X}%
\TextAt{9}{381.494}{176.957}{\Character{41}}%
\TextAt{8}{390.694}{176.957}{of}%
\TextAt{13}{404.773}{176.957}{\Character{16}}%
\TextAt{8}{415.436}{176.957}{factors}%
\TextAt{8}{454.299}{176.957}{into}%
\TextAt{8}{478.792}{176.957}{irreducible}%
\TextAt{8}{79.226}{159.911}{polynomials}%
\TextAt{8}{144.334}{159.911}{mod}%
\TextAt{13}{170.671}{159.911}{p}%
\TextAt{8}{176.546}{159.911}{.}%
\TextAt{8}{184.995}{159.911}{All}%
\TextAt{8}{204.171}{159.911}{we}%
\TextAt{8}{221.398}{159.911}{have}%
\TextAt{8}{248.393}{159.911}{to}%
\TextAt{8}{262.703}{159.911}{show}%
\TextAt{8}{291.707}{159.911}{is}%
\TextAt{8}{303.47}{159.911}{therefore}%
\TextAt{8}{352.57}{159.911}{that}%
\TextAt{13}{197.216}{130.91}{\Character{30}}%
\TextAt{21}{204.144}{129.117}{n}%
\TextAt{9}{209.78}{130.91}{\Character{40}}%
\TextAt{13}{214.332}{130.91}{X}%
\TextAt{9}{224.987}{130.91}{\Character{41}}%
\NativeMathAt{232.861}{130.91}{11.955}{\equiv}%
\TextAt{9}{245.48}{130.91}{\Character{40}}%
\TextAt{13}{250.032}{130.91}{p}%
\TextAt{17}{255.907}{129.117}{1}%
\TextAt{9}{260.64}{130.91}{\Character{40}}%
\TextAt{13}{265.192}{130.91}{X}%
\TextAt{9}{275.847}{130.91}{\Character{41}}%
\TextAt{16}{282.392}{130.91}{\Character{1}}%
\TextAt{16}{287.688}{130.91}{\Character{1}}%
\TextAt{16}{292.996}{130.91}{\Character{1}}%
\TextAt{13}{298.332}{130.91}{p}%
\TextAt{21}{304.207}{129.117}{r}%
\TextAt{9}{308.762}{130.91}{\Character{40}}%
\TextAt{13}{313.314}{130.91}{X}%
\TextAt{9}{323.969}{130.91}{\Character{41}\Character{41}}%
\TextAt{21}{333.074}{135.846}{\Character{39}}%
\TextAt{17}{338.592}{135.846}{\Character{40}}%
\TextAt{21}{341.886}{135.846}{p}%
\TextAt{22}{346.148}{138.863}{\Character{23}}%
\TextAt{22}{349.864}{137.866}{p}%
\TextAt{17}{354.7}{135.846}{\Character{41}}%
\TextAt{9}{370.447}{130.91}{\Character{40}}%
\TextAt{8}{374.999}{130.91}{mod}%
\TextAt{13}{401.101}{130.91}{p}%
\TextAt{9}{406.976}{130.91}{\Character{41}}%
\TextAt{13}{411.528}{130.91}{\Character{59}}%
\TextAt{8}{79.226}{101.909}{where}%
\TextAt{13}{113.664}{101.909}{p}%
\TextAt{17}{119.539}{100.116}{1}%
\TextAt{9}{124.271}{101.909}{\Character{40}}%
\TextAt{13}{128.824}{101.909}{X}%
\TextAt{9}{139.479}{101.909}{\Character{41}}%
\TextAt{13}{144.031}{101.909}{\Character{59}}%
\TextAt{13}{149.255}{101.909}{\Character{58}}%
\TextAt{13}{154.492}{101.909}{\Character{58}}%
\TextAt{13}{159.728}{101.909}{\Character{58}}%
\TextAt{13}{164.953}{101.909}{\Character{59}}%
\TextAt{13}{170.189}{101.909}{p}%
\TextAt{21}{176.127}{100.116}{r}%
\TextAt{9}{180.682}{101.909}{\Character{40}}%
\TextAt{13}{185.234}{101.909}{X}%
\TextAt{9}{195.889}{101.909}{\Character{41}}%
\TextAt{8}{204.968}{101.909}{are}%
\TextAt{8}{225.107}{101.909}{distinct}%
\TextAt{8}{268.069}{101.909}{irreducible}%
\TextAt{8}{326.572}{101.909}{polynomials}%
\TextAt{8}{392.309}{101.909}{over}%
\TextAt{18}{417.984}{101.909}{Z}%
\TextAt{9}{425.954}{101.909}{\Character{47}}%
\TextAt{13}{431.807}{101.909}{p}%
\TextAt{18}{437.682}{101.909}{Z}%
\TextAt{8}{450.178}{101.909}{of}%
\TextAt{8}{464.137}{101.909}{degree}%
\TextAt{13}{501.181}{101.909}{f}%
\TextAt{21}{506.951}{100.116}{p}%
\TextAt{8}{511.712}{101.909}{.}%
\TextAt{8}{522.046}{101.909}{In}%
\TextAt{8}{79.226}{84.863}{order}%
\TextAt{8}{110.43}{84.863}{to}%
\TextAt{8}{125.375}{84.863}{see}%
\TextAt{8}{144.923}{84.863}{this,}%
\TextAt{8}{171.79}{84.863}{put}%
\TextAt{13}{193.884}{84.863}{n}%
\TextAt{9}{205.274}{84.863}{\Character{61}}%
\TextAt{13}{218.781}{84.863}{p}%
\TextAt{21}{224.656}{89.201}{\Character{23}}%
\TextAt{22}{228.819}{88.205}{p}%
\TextAt{13}{233.655}{84.863}{m}%
\TextAt{8}{243.895}{84.863}{.}%
\TextAt{8}{254.247}{84.863}{As}%
\TextAt{13}{272.168}{84.863}{\Character{24}}%
\TextAt{21}{277.281}{83.07}{i}%
\TextAt{8}{280.662}{84.863}{,}%
\TextAt{8}{288.605}{84.863}{resp.}%
\TextAt{13}{319.83}{84.863}{\Character{17}}%
\TextAt{21}{325.606}{83.07}{j}%
\TextAt{8}{329.988}{84.863}{,}%
\TextAt{8}{337.93}{84.863}{varies}%
\TextAt{8}{371.478}{84.863}{over}%
\TextAt{8}{397.158}{84.863}{primitive}%
\TextAt{8}{447.873}{84.863}{roots}%
\TextAt{8}{478.168}{84.863}{of}%
\TextAt{8}{492.133}{84.863}{unity}%
\TextAt{8}{523.337}{84.863}{of}%
\TextAt{8}{300.14}{43.033}{57}%
\Destination{78.23}{809.815}{nav.page.57}%
\Destination{79.36}{776.155}{nav.print.61}%
\Destination{79.36}{208.955}{nav.print.62}%
\Destination{79.36}{393.055}{nav.number.1.10.3}%
\Hotlink{362.939}{758.955}{30.936}{11.047}{nav.number.1.10.1}%
\Hotlink{172.843}{592.155}{30.936}{11.047}{nav.number.1.2.11}%
\Hotlink{78.726}{191.755}{25.078}{11.047}{nav.number.1.8.3}%
}

% Page 70 (body page 58)
\Page{%
\TextAt{2}{79.226}{796.943}{I.}%
\TextAt{2}{89.83}{796.943}{Algebraic}%
\TextAt{2}{138.648}{796.943}{Integers}%
\Vector{q 0 G 0 g 0.299 w 7.226 18.729 m 460.769 18.729 l S Q}%
\TextAt{8}{79.226}{761.124}{order}%
\TextAt{13}{109.514}{761.124}{m}%
\TextAt{8}{119.753}{761.124}{,}%
\TextAt{8}{126.677}{761.124}{resp.}%
\TextAt{13}{155.906}{761.124}{p}%
\TextAt{21}{161.781}{765.462}{\Character{23}}%
\TextAt{22}{165.944}{764.466}{p}%
\TextAt{8}{170.78}{761.124}{,}%
\TextAt{8}{177.704}{761.124}{the}%
\TextAt{8}{197.579}{761.124}{products}%
\TextAt{13}{245.812}{761.124}{\Character{24}}%
\TextAt{21}{250.925}{759.33}{i}%
\TextAt{13}{254.306}{761.124}{\Character{17}}%
\TextAt{21}{260.082}{759.33}{j}%
\TextAt{8}{268.08}{761.124}{vary}%
\TextAt{8}{293.825}{761.124}{precisely}%
\TextAt{8}{341.4}{761.124}{over}%
\TextAt{8}{366.165}{761.124}{the}%
\TextAt{8}{386.04}{761.124}{primitive}%
\TextAt{13}{435.838}{761.124}{n}%
\TextAt{8}{442.826}{761.124}{-th}%
\TextAt{8}{461.398}{761.124}{roots}%
\TextAt{8}{490.777}{761.124}{of}%
\TextAt{8}{503.826}{761.124}{unity,}%
\TextAt{8}{79.226}{744.078}{i.e.,}%
\TextAt{8}{101.331}{744.078}{one}%
\TextAt{8}{122.791}{744.078}{has}%
\TextAt{8}{143.665}{744.078}{the}%
\TextAt{8}{163.821}{744.078}{decomposition}%
\TextAt{8}{240.956}{744.078}{over}%
\TextAt{14}{266.002}{744.078}{o}%
\TextAt{8}{271.846}{744.078}{:}%
\TextAt{13}{246.187}{717.069}{\Character{30}}%
\TextAt{21}{253.115}{715.276}{n}%
\TextAt{9}{258.752}{717.069}{\Character{40}}%
\TextAt{13}{263.304}{717.069}{X}%
\TextAt{9}{273.959}{717.069}{\Character{41}}%
\TextAt{9}{281.814}{717.069}{\Character{61}}%
\TextAt{20}{294.258}{728.427}{\Character{89}}%
\TextAt{21}{297.336}{703.232}{i\Character{59}j}%
\TextAt{9}{309.534}{717.069}{\Character{40}}%
\TextAt{13}{314.086}{717.069}{X}%
\TextAt{16}{327.398}{717.069}{\Character{0}}%
\TextAt{13}{339.353}{717.069}{\Character{24}}%
\TextAt{21}{344.466}{715.276}{i}%
\TextAt{13}{347.847}{717.069}{\Character{17}}%
\TextAt{21}{353.622}{715.276}{j}%
\TextAt{9}{358.005}{717.069}{\Character{41}}%
\TextAt{13}{362.557}{717.069}{\Character{58}}%
\TextAt{8}{79.226}{679.6}{Since}%
\TextAt{13}{109.841}{679.6}{X}%
\TextAt{21}{120.496}{683.938}{p}%
\TextAt{22}{124.759}{686.955}{\Character{23}}%
\TextAt{22}{128.474}{685.958}{p}%
\TextAt{16}{136.504}{679.6}{\Character{0}}%
\TextAt{9}{148.499}{679.6}{1}%
\NativeMathAt{157.77}{679.6}{11.955}{\equiv}%
\TextAt{9}{170.487}{679.6}{\Character{40}}%
\TextAt{13}{175.039}{679.6}{X}%
\TextAt{16}{188.39}{679.6}{\Character{0}}%
\TextAt{9}{200.384}{679.6}{1\Character{41}}%
\TextAt{21}{210.79}{683.938}{p}%
\TextAt{22}{215.052}{686.955}{\Character{23}}%
\TextAt{22}{218.768}{685.958}{p}%
\TextAt{9}{229.416}{679.6}{\Character{40}}%
\TextAt{8}{233.968}{679.6}{mod}%
\TextAt{14}{260.07}{679.6}{p}%
\TextAt{9}{266.053}{679.6}{\Character{41}}%
\TextAt{8}{270.606}{679.6}{,}%
\TextAt{8}{277.826}{679.6}{one}%
\TextAt{8}{299.343}{679.6}{has}%
\TextAt{13}{320.274}{679.6}{\Character{17}}%
\TextAt{21}{326.05}{677.807}{j}%
\NativeMathAt{333.85}{679.6}{11.955}{\equiv}%
\TextAt{9}{346.567}{679.6}{1}%
\TextAt{9}{357.721}{679.6}{\Character{40}}%
\TextAt{8}{362.286}{679.6}{mod}%
\TextAt{14}{388.388}{679.6}{p}%
\TextAt{9}{394.371}{679.6}{\Character{41}}%
\TextAt{8}{398.923}{679.6}{,}%
\TextAt{8}{406.144}{679.6}{for}%
\TextAt{8}{424.086}{679.6}{any}%
\TextAt{8}{446.261}{679.6}{prime}%
\TextAt{8}{479.481}{679.6}{ideal}%
\TextAt{14}{507.502}{679.6}{p}%
\TextAt{16}{516.903}{679.6}{\Character{106}}%
\TextAt{13}{523.643}{679.6}{p}%
\TextAt{8}{529.518}{679.6}{.}%
\TextAt{8}{79.226}{662.554}{In}%
\TextAt{8}{93.847}{662.554}{other}%
\TextAt{8}{124.417}{662.554}{words,}%
\TextAt{13}{173.968}{635.545}{\Character{30}}%
\TextAt{21}{180.896}{633.752}{n}%
\TextAt{9}{186.532}{635.545}{\Character{40}}%
\TextAt{13}{191.084}{635.545}{X}%
\TextAt{9}{201.739}{635.545}{\Character{41}}%
\NativeMathAt{209.613}{635.545}{11.955}{\equiv}%
\TextAt{20}{222.232}{646.903}{\Character{89}}%
\TextAt{21}{228.428}{621.709}{i}%
\TextAt{9}{237.508}{635.545}{\Character{40}}%
\TextAt{13}{242.06}{635.545}{X}%
\TextAt{16}{255.372}{635.545}{\Character{0}}%
\TextAt{13}{267.327}{635.545}{\Character{24}}%
\TextAt{21}{272.44}{633.752}{i}%
\TextAt{9}{275.821}{635.545}{\Character{41}}%
\TextAt{21}{280.374}{640.482}{\Character{39}}%
\TextAt{17}{285.892}{640.482}{\Character{40}}%
\TextAt{21}{289.185}{640.482}{p}%
\TextAt{22}{293.448}{643.498}{\Character{23}}%
\TextAt{22}{297.163}{642.502}{p}%
\TextAt{17}{302}{640.482}{\Character{41}}%
\TextAt{9}{309.112}{635.545}{\Character{61}}%
\TextAt{13}{321.538}{635.545}{\Character{30}}%
\TextAt{21}{328.466}{633.752}{m}%
\TextAt{9}{336.454}{635.545}{\Character{40}}%
\TextAt{13}{341.007}{635.545}{X}%
\TextAt{9}{351.662}{635.545}{\Character{41}}%
\TextAt{21}{356.214}{640.482}{\Character{39}}%
\TextAt{17}{361.732}{640.482}{\Character{40}}%
\TextAt{21}{365.026}{640.482}{p}%
\TextAt{22}{369.288}{643.498}{\Character{23}}%
\TextAt{22}{373.004}{642.502}{p}%
\TextAt{17}{377.84}{640.482}{\Character{41}}%
\TextAt{9}{393.587}{635.545}{\Character{40}}%
\TextAt{8}{398.139}{635.545}{mod}%
\TextAt{14}{424.241}{635.545}{p}%
\TextAt{9}{430.224}{635.545}{\Character{41}}%
\TextAt{13}{434.776}{635.545}{\Character{58}}%
\TextAt{8}{79.226}{601.257}{This}%
\TextAt{8}{105.946}{601.257}{implies}%
\TextAt{8}{145.673}{601.257}{the}%
\TextAt{8}{165.829}{601.257}{congruence}%
\TextAt{13}{223.674}{574.249}{\Character{30}}%
\TextAt{21}{230.602}{572.456}{n}%
\TextAt{9}{236.239}{574.249}{\Character{40}}%
\TextAt{13}{240.791}{574.249}{X}%
\TextAt{9}{251.446}{574.249}{\Character{41}}%
\NativeMathAt{259.319}{574.249}{11.955}{\equiv}%
\TextAt{13}{271.939}{574.249}{\Character{30}}%
\TextAt{21}{278.867}{572.456}{m}%
\TextAt{9}{286.855}{574.249}{\Character{40}}%
\TextAt{13}{291.408}{574.249}{X}%
\TextAt{9}{302.063}{574.249}{\Character{41}}%
\TextAt{21}{306.615}{579.185}{\Character{39}}%
\TextAt{17}{312.133}{579.185}{\Character{40}}%
\TextAt{21}{315.427}{579.185}{p}%
\TextAt{22}{319.69}{582.202}{\Character{23}}%
\TextAt{22}{323.405}{581.205}{p}%
\TextAt{17}{328.241}{579.185}{\Character{41}}%
\TextAt{9}{343.988}{574.249}{\Character{40}}%
\TextAt{8}{348.54}{574.249}{mod}%
\TextAt{13}{374.642}{574.249}{p}%
\TextAt{9}{380.517}{574.249}{\Character{41}}%
\TextAt{13}{385.07}{574.249}{\Character{58}}%
\TextAt{8}{79.226}{547.241}{Observing}%
\TextAt{8}{134.196}{547.241}{that}%
\TextAt{13}{158.86}{547.241}{f}%
\TextAt{21}{164.63}{545.447}{p}%
\TextAt{8}{172.583}{547.241}{is}%
\TextAt{8}{183.642}{547.241}{the}%
\TextAt{8}{203.094}{547.241}{smallest}%
\TextAt{8}{247.388}{547.241}{positive}%
\TextAt{8}{289.996}{547.241}{integer}%
\TextAt{8}{327.99}{547.241}{such}%
\TextAt{8}{353.682}{547.241}{that}%
\TextAt{13}{378.346}{547.241}{p}%
\TextAt{21}{384.221}{551.579}{f}%
\TextAt{22}{388.289}{550.583}{p}%
\NativeMathAt{396.447}{547.241}{11.955}{\equiv}%
\TextAt{9}{409.066}{547.241}{1}%
\TextAt{9}{420.22}{547.241}{\Character{40}}%
\TextAt{8}{424.785}{547.241}{mod}%
\TextAt{13}{450.887}{547.241}{m}%
\TextAt{9}{461.126}{547.241}{\Character{41}}%
\TextAt{8}{465.679}{547.241}{,}%
\TextAt{8}{472.264}{547.241}{it}%
\TextAt{8}{483.263}{547.241}{is}%
\TextAt{8}{494.322}{547.241}{obvious}%
\TextAt{8}{79.226}{530.195}{that}%
\TextAt{8}{104.595}{530.195}{this}%
\TextAt{8}{127.417}{530.195}{congruence}%
\TextAt{8}{187.899}{530.195}{reduces}%
\TextAt{8}{229.574}{530.195}{us}%
\TextAt{8}{244.59}{530.195}{to}%
\TextAt{8}{258.9}{530.195}{the}%
\TextAt{8}{279.057}{530.195}{case}%
\TextAt{8}{303.828}{530.195}{where}%
\TextAt{13}{337.637}{530.195}{p}%
\TextAt{18}{346.833}{530.195}{\Character{45}}%
\TextAt{13}{353.475}{530.195}{n}%
\TextAt{8}{360.462}{530.195}{,}%
\TextAt{8}{367.612}{530.195}{and}%
\TextAt{8}{390.374}{530.195}{hence}%
\TextAt{13}{422.88}{530.195}{\Character{39}}%
\TextAt{9}{430.549}{530.195}{\Character{40}}%
\TextAt{13}{435.101}{530.195}{p}%
\TextAt{21}{440.976}{534.533}{\Character{23}}%
\TextAt{22}{445.139}{533.537}{p}%
\TextAt{9}{449.976}{530.195}{\Character{41}}%
\TextAt{9}{457.831}{530.195}{\Character{61}}%
\TextAt{13}{470.274}{530.195}{\Character{39}}%
\TextAt{9}{477.943}{530.195}{\Character{40}1\Character{41}}%
\TextAt{9}{496.187}{530.195}{\Character{61}}%
\TextAt{9}{508.608}{530.195}{1}%
\TextAt{8}{514.5}{530.195}{.}%
\TextAt{8}{97.159}{510.16}{As}%
\TextAt{8}{114.615}{510.16}{the}%
\TextAt{8}{134.94}{510.16}{characteristic}%
\TextAt{13}{207.043}{510.16}{p}%
\TextAt{8}{216.984}{510.16}{of}%
\TextAt{14}{230.483}{510.16}{o}%
\TextAt{9}{236.327}{510.16}{\Character{47}}%
\TextAt{14}{242.18}{510.16}{p}%
\TextAt{8}{252.23}{510.16}{does}%
\TextAt{8}{278.795}{510.16}{not}%
\TextAt{8}{299.778}{510.16}{divide}%
\TextAt{13}{334.736}{510.16}{n}%
\TextAt{8}{341.724}{510.16}{,}%
\TextAt{8}{349.084}{510.16}{the}%
\TextAt{8}{369.409}{510.16}{polynomials}%
\TextAt{13}{434.686}{510.16}{X}%
\TextAt{21}{445.341}{514.498}{n}%
\TextAt{16}{453.749}{510.16}{\Character{0}}%
\TextAt{9}{465.819}{510.16}{1}%
\TextAt{8}{475.738}{510.16}{and}%
\TextAt{13}{498.67}{510.16}{nX}%
\TextAt{21}{516.313}{514.498}{n}%
\TextAt{3}{521.451}{514.498}{\Character{0}}%
\TextAt{17}{528.037}{514.498}{1}%
\TextAt{8}{79.226}{493.114}{have}%
\TextAt{8}{106.506}{493.114}{no}%
\TextAt{8}{123.051}{493.114}{common}%
\TextAt{8}{170.165}{493.114}{root}%
\TextAt{8}{195.496}{493.114}{in}%
\TextAt{14}{209.434}{493.114}{o}%
\TextAt{9}{215.279}{493.114}{\Character{47}}%
\TextAt{14}{221.132}{493.114}{p}%
\TextAt{8}{227.115}{493.114}{.}%
\TextAt{8}{236.419}{493.114}{So}%
\TextAt{13}{252.964}{493.114}{X}%
\TextAt{21}{263.619}{497.452}{n}%
\TextAt{16}{272.106}{493.114}{\Character{0}}%
\TextAt{9}{284.256}{493.114}{1}%
\TextAt{8}{293.43}{493.114}{mod}%
\TextAt{14}{318.868}{493.114}{p}%
\TextAt{8}{329.034}{493.114}{has}%
\TextAt{8}{350.193}{493.114}{no}%
\TextAt{8}{366.737}{493.114}{multiple}%
\TextAt{8}{412.871}{493.114}{roots.}%
\TextAt{8}{447.939}{493.114}{We}%
\TextAt{8}{468.369}{493.114}{therefore}%
\TextAt{8}{517.754}{493.114}{see}%
\TextAt{8}{79.226}{476.068}{that}%
\TextAt{8}{104.269}{476.068}{passing}%
\TextAt{8}{145.044}{476.068}{to}%
\TextAt{8}{159.029}{476.068}{the}%
\TextAt{8}{178.859}{476.068}{quotient}%
\TextAt{14}{224.715}{476.068}{o}%
\NativeMathAt{233.88}{476.068}{11.955}{\rightarrow}%
\TextAt{14}{249.156}{476.068}{o}%
\TextAt{9}{255.001}{476.068}{\Character{47}}%
\TextAt{14}{260.854}{476.068}{p}%
\TextAt{8}{270.408}{476.068}{maps}%
\TextAt{8}{300.711}{476.068}{the}%
\TextAt{8}{320.541}{476.068}{group}%
\TextAt{13}{353.39}{476.068}{\Character{22}}%
\TextAt{21}{360.433}{474.275}{n}%
\TextAt{8}{369.641}{476.068}{of}%
\TextAt{13}{382.644}{476.068}{n}%
\TextAt{8}{389.632}{476.068}{-th}%
\TextAt{8}{408.159}{476.068}{roots}%
\TextAt{8}{437.494}{476.068}{of}%
\TextAt{8}{450.497}{476.068}{unity}%
\TextAt{8}{480.741}{476.068}{bijectively}%
\TextAt{8}{79.226}{459.022}{onto}%
\TextAt{8}{105.714}{459.022}{the}%
\TextAt{8}{126.01}{459.022}{group}%
\Vector{q 0 G 0 g 0.478 w 87.325 -308.157 m 94.368 -308.157 l S Q}%
\TextAt{13}{159.325}{459.022}{\Character{22}}%
\TextAt{21}{166.367}{456.199}{n}%
\TextAt{8}{176.04}{459.022}{of}%
\TextAt{13}{189.509}{459.022}{n}%
\TextAt{8}{196.497}{459.022}{-th}%
\TextAt{8}{215.489}{459.022}{roots}%
\TextAt{8}{245.289}{459.022}{of}%
\TextAt{8}{258.758}{459.022}{unity}%
\TextAt{8}{289.467}{459.022}{of}%
\TextAt{14}{302.936}{459.022}{o}%
\TextAt{9}{308.78}{459.022}{\Character{47}}%
\TextAt{14}{314.633}{459.022}{p}%
\TextAt{8}{320.616}{459.022}{.}%
\TextAt{8}{329.482}{459.022}{In}%
\TextAt{8}{344.242}{459.022}{particular,}%
\TextAt{8}{401.657}{459.022}{the}%
\TextAt{8}{421.952}{459.022}{primitive}%
\TextAt{13}{472.172}{459.022}{n}%
\TextAt{8}{479.159}{459.022}{-th}%
\TextAt{8}{498.152}{459.022}{root}%
\TextAt{8}{523.337}{459.022}{of}%
\TextAt{8}{79.226}{441.976}{unity}%
\TextAt{13}{110.019}{441.976}{\Character{16}}%
\TextAt{8}{120.155}{441.976}{modulo}%
\TextAt{14}{162.328}{441.976}{p}%
\TextAt{8}{172.432}{441.976}{remains}%
\TextAt{8}{216.291}{441.976}{a}%
\TextAt{8}{226.27}{441.976}{primitive}%
\TextAt{13}{276.573}{441.976}{n}%
\TextAt{8}{283.561}{441.976}{-th}%
\TextAt{8}{302.637}{441.976}{root}%
\TextAt{8}{327.906}{441.976}{of}%
\TextAt{8}{341.459}{441.976}{unity.}%
\TextAt{8}{376.269}{441.976}{The}%
\TextAt{8}{400.546}{441.976}{smallest}%
\TextAt{8}{445.768}{441.976}{extension}%
\TextAt{8}{497.757}{441.976}{field}%
\TextAt{8}{523.337}{441.976}{of}%
\TextAt{18}{79.226}{424.93}{F}%
\TextAt{21}{86.532}{423.137}{p}%
\TextAt{9}{95.592}{424.93}{\Character{61}}%
\TextAt{18}{108.995}{424.93}{Z}%
\TextAt{9}{116.966}{424.93}{\Character{47}}%
\TextAt{13}{122.819}{424.93}{p}%
\TextAt{18}{128.694}{424.93}{Z}%
\TextAt{8}{141.135}{424.93}{containing}%
\TextAt{8}{198.628}{424.93}{it}%
\TextAt{8}{210.906}{424.93}{is}%
\TextAt{8}{223.243}{424.93}{the}%
\TextAt{8}{243.974}{424.93}{field}%
\TextAt{18}{269.905}{424.93}{F}%
\TextAt{21}{277.211}{421.88}{p}%
\TextAt{22}{281.473}{424.897}{f}%
\TextAt{22}{285.064}{423.9}{p}%
\TextAt{8}{290.399}{424.93}{,}%
\TextAt{8}{298.265}{424.93}{because}%
\TextAt{8}{342.141}{424.93}{its}%
\TextAt{8}{359.034}{424.93}{multiplicative}%
\TextAt{8}{433.431}{424.93}{group}%
\TextAt{18}{467.18}{424.93}{F}%
\TextAt{3}{474.486}{429.268}{\Character{3}}%
\TextAt{21}{474.486}{420.188}{p}%
\TextAt{22}{478.749}{423.205}{f}%
\TextAt{22}{482.34}{422.208}{p}%
\TextAt{8}{492.146}{424.93}{is}%
\TextAt{8}{504.484}{424.93}{cyclic}%
\TextAt{8}{79.226}{407.884}{of}%
\TextAt{8}{93.372}{407.884}{order}%
\TextAt{13}{124.758}{407.884}{p}%
\TextAt{21}{130.633}{412.222}{f}%
\TextAt{22}{134.701}{411.226}{p}%
\TextAt{16}{142.751}{407.884}{\Character{0}}%
\TextAt{9}{155.262}{407.884}{1}%
\TextAt{8}{161.115}{407.884}{.}%
\TextAt{18}{172.012}{407.884}{F}%
\TextAt{21}{179.318}{404.834}{p}%
\TextAt{22}{183.58}{407.851}{f}%
\TextAt{22}{187.171}{406.855}{p}%
\TextAt{8}{197.219}{407.884}{is}%
\TextAt{8}{209.799}{407.884}{therefore}%
\TextAt{8}{259.715}{407.884}{the}%
\TextAt{8}{280.688}{407.884}{field}%
\TextAt{8}{306.861}{407.884}{of}%
\TextAt{8}{321.007}{407.884}{decomposition}%
\TextAt{8}{398.958}{407.884}{of}%
\TextAt{8}{413.104}{407.884}{the}%
\TextAt{8}{434.076}{407.884}{reduced}%
\TextAt{8}{478.457}{407.884}{cyclotomic}%
\TextAt{8}{79.226}{390.838}{polynomial}%
\Vector{q 0 G 0 g 0.478 w 174.858 -390.232 m 181.786 -390.232 l S Q}%
\TextAt{13}{246.858}{373.792}{\Character{30}}%
\TextAt{21}{253.786}{370.969}{n}%
\TextAt{9}{259.422}{373.792}{\Character{40}}%
\TextAt{13}{263.974}{373.792}{X}%
\TextAt{9}{274.63}{373.792}{\Character{41}}%
\TextAt{9}{282.485}{373.792}{\Character{61}}%
\TextAt{13}{294.928}{373.792}{\Character{30}}%
\TextAt{21}{301.856}{371.999}{n}%
\TextAt{9}{307.493}{373.792}{\Character{40}}%
\TextAt{13}{312.045}{373.792}{X}%
\TextAt{9}{322.7}{373.792}{\Character{41}}%
\TextAt{8}{330.573}{373.792}{mod}%
\TextAt{13}{356.011}{373.792}{p\Character{58}}%
\TextAt{8}{79.226}{351.124}{Being}%
\TextAt{8}{112.222}{351.124}{a}%
\TextAt{8}{121.978}{351.124}{divisor}%
\TextAt{8}{160.091}{351.124}{of}%
\TextAt{13}{173.421}{351.124}{X}%
\TextAt{21}{184.076}{355.463}{n}%
\TextAt{16}{192.369}{351.124}{\Character{0}}%
\TextAt{9}{204.324}{351.124}{1}%
\TextAt{8}{213.498}{351.124}{mod}%
\TextAt{13}{238.936}{351.124}{p}%
\TextAt{8}{244.811}{351.124}{,}%
\TextAt{8}{251.96}{351.124}{this}%
\TextAt{8}{274.783}{351.124}{polynomial}%
\TextAt{8}{335.276}{351.124}{has}%
\TextAt{8}{356.15}{351.124}{no}%
\TextAt{8}{372.409}{351.124}{multiple}%
\TextAt{8}{418.257}{351.124}{roots,}%
\TextAt{8}{451.169}{351.124}{and}%
\TextAt{8}{473.932}{351.124}{if}%
\Vector{q 0 G 0 g 0.478 w 170.718 -439.908 m 177.646 -439.908 l S Q}%
\TextAt{13}{242.718}{324.116}{\Character{30}}%
\TextAt{21}{249.646}{321.293}{n}%
\TextAt{9}{255.282}{324.116}{\Character{40}}%
\TextAt{13}{259.835}{324.116}{X}%
\TextAt{9}{270.49}{324.116}{\Character{41}}%
\TextAt{9}{278.345}{324.116}{\Character{61}}%
\Vector{q 0 G 0 g 0.478 w 218.788 -443.063 m 224.663 -443.063 l S Q}%
\TextAt{13}{290.788}{324.116}{p}%
\TextAt{17}{296.663}{321.293}{1}%
\TextAt{9}{301.396}{324.116}{\Character{40}}%
\TextAt{13}{305.948}{324.116}{X}%
\TextAt{9}{316.603}{324.116}{\Character{41}}%
\TextAt{16}{323.148}{324.116}{\Character{1}}%
\TextAt{16}{328.444}{324.116}{\Character{1}}%
\TextAt{16}{333.752}{324.116}{\Character{1}}%
\Vector{q 0 G 0 g 0.478 w 267.088 -443.063 m 272.963 -443.063 l S Q}%
\TextAt{13}{339.088}{324.116}{p}%
\TextAt{21}{344.963}{321.293}{r}%
\TextAt{9}{349.518}{324.116}{\Character{40}}%
\TextAt{13}{354.07}{324.116}{X}%
\TextAt{9}{364.726}{324.116}{\Character{41}}%
\TextAt{8}{79.226}{297.108}{is}%
\TextAt{8}{91.177}{297.108}{its}%
\TextAt{8}{107.683}{297.108}{factorization}%
\TextAt{8}{175.848}{297.108}{into}%
\TextAt{8}{199.778}{297.108}{irreducibles}%
\TextAt{8}{262.455}{297.108}{over}%
\TextAt{18}{287.688}{297.108}{F}%
\TextAt{21}{294.994}{295.314}{p}%
\TextAt{8}{299.755}{297.108}{,}%
\TextAt{8}{307.139}{297.108}{then}%
\TextAt{8}{333.986}{297.108}{every}%
\Vector{q 0 G 0 g 0.478 w 293.065 -470.071 m 298.941 -470.071 l S Q}%
\TextAt{13}{365.065}{297.108}{p}%
\TextAt{21}{370.941}{294.285}{i}%
\TextAt{9}{374.322}{297.108}{\Character{40}}%
\TextAt{13}{378.874}{297.108}{X}%
\TextAt{9}{389.529}{297.108}{\Character{41}}%
\TextAt{8}{398.166}{297.108}{is}%
\TextAt{8}{410.117}{297.108}{the}%
\TextAt{8}{430.461}{297.108}{minimal}%
\TextAt{8}{476.174}{297.108}{polynomial}%
\TextAt{8}{79.226}{280.062}{of}%
\TextAt{8}{93.597}{280.062}{a}%
\TextAt{8}{104.394}{280.062}{primitive}%
\TextAt{13}{155.516}{280.062}{n}%
\TextAt{8}{162.503}{280.062}{-th}%
\TextAt{8}{182.398}{280.062}{root}%
\TextAt{8}{208.485}{280.062}{of}%
\TextAt{8}{222.856}{280.062}{unity}%
\Vector{q 0 G 0 g 0.478 w 182.467 -483.963 m 188.482 -483.963 l S Q}%
\TextAt{13}{254.467}{280.062}{\Character{16}}%
\NativeMathAt{265.578}{280.062}{11.955}{\in}%
\TextAt{18}{278.643}{280.062}{F}%
\TextAt{3}{285.949}{284.4}{\Character{3}}%
\TextAt{21}{285.949}{275.32}{p}%
\TextAt{22}{290.212}{278.336}{f}%
\TextAt{22}{293.803}{277.34}{p}%
\TextAt{8}{299.137}{280.062}{.}%
\TextAt{8}{310.709}{280.062}{Its}%
\TextAt{8}{329.038}{280.062}{degree}%
\TextAt{8}{366.495}{280.062}{is}%
\TextAt{8}{379.3}{280.062}{therefore}%
\TextAt{13}{429.441}{280.062}{f}%
\TextAt{21}{435.211}{278.268}{p}%
\TextAt{8}{439.972}{280.062}{.}%
\TextAt{8}{451.544}{280.062}{This}%
\TextAt{8}{479.305}{280.062}{proves}%
\TextAt{8}{516.51}{280.062}{the}%
\TextAt{8}{79.226}{263.016}{proposition.}%
\Vector{q 0 G 0 g 0.399 w 452.334 -510.984 m 452.334 -502.915 l S Q}%
\Vector{q 0 G 0 g 0.399 w 452.534 -503.114 m 459.707 -503.114 l S Q}%
\Vector{q 0 G 0 g 0.399 w 452.534 -510.785 m 459.707 -510.785 l S Q}%
\Vector{q 0 G 0 g 0.399 w 459.906 -510.984 m 459.906 -502.915 l S Q}%
\TextAt{8}{97.159}{234.584}{Let}%
\TextAt{8}{118.128}{234.584}{us}%
\TextAt{8}{133.144}{234.584}{emphasize}%
\TextAt{8}{189.13}{234.584}{two}%
\TextAt{8}{211.247}{234.584}{special}%
\TextAt{8}{249.348}{234.584}{cases}%
\TextAt{8}{278.734}{234.584}{of}%
\TextAt{8}{292.064}{234.584}{the}%
\TextAt{8}{312.22}{234.584}{above}%
\TextAt{8}{345.396}{234.584}{law}%
\TextAt{8}{366.533}{234.584}{of}%
\TextAt{8}{379.863}{234.584}{decomposition:}%
\TextAt{12}{79.226}{209.141}{Corollary}%
\TextAt{12}{139.014}{209.141}{(10.4).}%
\TextAt{8}{181.06}{209.141}{A}%
\TextAt{8}{193.733}{209.141}{prime}%
\TextAt{8}{226.896}{209.141}{number}%
\TextAt{13}{269.493}{209.141}{p}%
\TextAt{8}{279.265}{209.141}{is}%
\TextAt{8}{291.029}{209.141}{ramified}%
\TextAt{8}{336.554}{209.141}{in}%
\TextAt{18}{350.207}{209.141}{Q}%
\TextAt{9}{359.506}{209.141}{\Character{40}}%
\TextAt{13}{364.058}{209.141}{\Character{16}}%
\TextAt{9}{370.074}{209.141}{\Character{41}}%
\TextAt{8}{378.523}{209.141}{if}%
\TextAt{8}{389.247}{209.141}{and}%
\TextAt{8}{412.01}{209.141}{only}%
\TextAt{8}{437.701}{209.141}{if}%
\TextAt{13}{263.463}{182.132}{n}%
\NativeMathAt{273.772}{182.132}{11.955}{\equiv}%
\TextAt{9}{286.391}{182.132}{0}%
\TextAt{9}{304.18}{182.132}{\Character{40}}%
\TextAt{8}{308.751}{182.132}{mod}%
\TextAt{13}{334.853}{182.132}{p}%
\TextAt{9}{340.728}{182.132}{\Character{41}}%
\TextAt{13}{345.281}{182.132}{\Character{59}}%
\TextAt{8}{79.226}{155.124}{except}%
\TextAt{8}{116.027}{155.124}{in}%
\TextAt{8}{129.743}{155.124}{the}%
\TextAt{8}{149.962}{155.124}{case}%
\TextAt{8}{174.796}{155.124}{where}%
\TextAt{13}{208.668}{155.124}{p}%
\TextAt{9}{217.971}{155.124}{\Character{61}}%
\TextAt{9}{230.488}{155.124}{2}%
\TextAt{9}{239.765}{155.124}{\Character{61}}%
\TextAt{9}{252.294}{155.124}{\Character{40}4}%
\TextAt{13}{262.723}{155.124}{\Character{59}}%
\TextAt{13}{267.947}{155.124}{n}%
\TextAt{9}{274.955}{155.124}{\Character{41}}%
\TextAt{8}{279.507}{155.124}{.}%
\TextAt{8}{288.144}{155.124}{A}%
\TextAt{8}{300.879}{155.124}{prime}%
\TextAt{8}{334.106}{155.124}{number}%
\TextAt{13}{376.765}{155.124}{p}%
\NativeMathAt{386.068}{155.124}{11.955}{\neq}%
\TextAt{9}{398.585}{155.124}{2}%
\TextAt{8}{408.414}{155.124}{is}%
\TextAt{8}{420.241}{155.124}{totally}%
\TextAt{8}{457.711}{155.124}{split}%
\TextAt{8}{483.848}{155.124}{in}%
\TextAt{18}{497.564}{155.124}{Q}%
\TextAt{9}{506.863}{155.124}{\Character{40}}%
\TextAt{13}{511.415}{155.124}{\Character{16}}%
\TextAt{9}{517.43}{155.124}{\Character{41}}%
\TextAt{8}{525.943}{155.124}{if}%
\TextAt{8}{79.226}{138.078}{and}%
\TextAt{8}{101.989}{138.078}{only}%
\TextAt{8}{127.68}{138.078}{if}%
\TextAt{13}{263.463}{121.032}{p}%
\NativeMathAt{272.659}{121.032}{11.955}{\equiv}%
\TextAt{9}{285.278}{121.032}{1}%
\TextAt{9}{303.079}{121.032}{\Character{40}}%
\TextAt{8}{307.639}{121.032}{mod}%
\TextAt{13}{333.741}{121.032}{n}%
\TextAt{9}{340.728}{121.032}{\Character{41}}%
\TextAt{13}{345.281}{121.032}{\Character{58}}%
\TextAt{8}{97.159}{89.967}{The}%
\TextAt{8}{120.714}{89.967}{completeness}%
\TextAt{8}{190.571}{89.967}{of}%
\TextAt{8}{203.402}{89.967}{these}%
\TextAt{8}{232.875}{89.967}{results}%
\TextAt{8}{269.569}{89.967}{concerning}%
\TextAt{8}{327.602}{89.967}{the}%
\TextAt{8}{347.26}{89.967}{integral}%
\TextAt{8}{389.369}{89.967}{basis}%
\TextAt{8}{417.611}{89.967}{and}%
\TextAt{8}{439.875}{89.967}{the}%
\TextAt{8}{459.532}{89.967}{decomposition}%
\TextAt{8}{79.226}{72.921}{of}%
\TextAt{8}{93.192}{72.921}{primes}%
\TextAt{8}{131.606}{72.921}{in}%
\TextAt{8}{145.895}{72.921}{the}%
\TextAt{8}{166.687}{72.921}{field}%
\TextAt{18}{192.68}{72.921}{Q}%
\TextAt{9}{201.979}{72.921}{\Character{40}}%
\TextAt{13}{206.531}{72.921}{\Character{16}}%
\TextAt{9}{212.547}{72.921}{\Character{41}}%
\TextAt{8}{221.632}{72.921}{will}%
\TextAt{8}{244.373}{72.921}{not}%
\TextAt{8}{265.823}{72.921}{be}%
\TextAt{8}{282.384}{72.921}{matched}%
\TextAt{8}{330.171}{72.921}{by}%
\TextAt{8}{347.066}{72.921}{our}%
\TextAt{8}{368.516}{72.921}{study}%
\TextAt{8}{401.406}{72.921}{of}%
\TextAt{8}{415.372}{72.921}{the}%
\TextAt{8}{436.165}{72.921}{group}%
\TextAt{8}{469.976}{72.921}{of}%
\TextAt{8}{483.942}{72.921}{units}%
\TextAt{8}{513.904}{72.921}{and}%
\TextAt{8}{300.14}{43.033}{58}%
\Destination{78.23}{809.815}{nav.page.58}%
\Destination{79.36}{224.155}{nav.print.63}%
\Destination{79.36}{223.155}{nav.number.1.10.4}%
}

% Page 71 (body page 59)
\Page{%
\TextAt{2}{79.226}{796.943}{I.}%
\TextAt{2}{89.83}{796.943}{Algebraic}%
\TextAt{2}{138.648}{796.943}{Integers}%
\Vector{q 0 G 0 g 0.299 w 7.226 18.729 m 460.769 18.729 l S Q}%
\TextAt{8}{79.226}{761.124}{the}%
\TextAt{8}{100.045}{761.124}{ideal}%
\TextAt{8}{128.67}{761.124}{class}%
\TextAt{8}{156.769}{761.124}{group.}%
\TextAt{8}{196.481}{761.124}{The}%
\TextAt{8}{221.197}{761.124}{problems}%
\TextAt{8}{271.999}{761.124}{arising}%
\TextAt{8}{310.451}{761.124}{in}%
\TextAt{8}{324.766}{761.124}{this}%
\TextAt{8}{348.251}{761.124}{context}%
\TextAt{8}{390.54}{761.124}{are}%
\TextAt{8}{410.713}{761.124}{in}%
\TextAt{8}{425.028}{761.124}{fact}%
\TextAt{8}{448.775}{761.124}{among}%
\TextAt{8}{487.168}{761.124}{the}%
\TextAt{8}{507.986}{761.124}{most}%
\TextAt{8}{79.226}{744.078}{difficult}%
\TextAt{8}{122.866}{744.078}{problems}%
\TextAt{8}{173.727}{744.078}{posed}%
\TextAt{8}{207.349}{744.078}{by}%
\TextAt{8}{224.329}{744.078}{algebraic}%
\TextAt{8}{274.485}{744.078}{number}%
\TextAt{8}{317.802}{744.078}{theory.}%
\TextAt{8}{360.286}{744.078}{At}%
\TextAt{8}{377.911}{744.078}{the}%
\TextAt{8}{398.789}{744.078}{same}%
\TextAt{8}{428.836}{744.078}{time}%
\TextAt{8}{456.217}{744.078}{one}%
\TextAt{8}{478.397}{744.078}{encounters}%
\TextAt{8}{79.226}{727.032}{here}%
\TextAt{8}{103.692}{727.032}{plenty}%
\TextAt{8}{138.248}{727.032}{of}%
\TextAt{8}{150.687}{727.032}{astonishing}%
\TextAt{8}{211.066}{727.032}{laws}%
\TextAt{8}{235.927}{727.032}{which}%
\TextAt{8}{268.522}{727.032}{are}%
\TextAt{8}{287.142}{727.032}{the}%
\TextAt{8}{306.407}{727.032}{subject}%
\TextAt{8}{345.889}{727.032}{of}%
\TextAt{8}{358.328}{727.032}{a}%
\TextAt{8}{367.192}{727.032}{theory}%
\TextAt{8}{403.051}{727.032}{which}%
\TextAt{8}{435.646}{727.032}{has}%
\TextAt{8}{455.629}{727.032}{been}%
\TextAt{8}{482.366}{727.032}{developed}%
\TextAt{8}{79.226}{709.986}{only}%
\TextAt{8}{104.918}{709.986}{recently,}%
\TextAt{12}{151.411}{709.986}{Iwasawa}%
\TextAt{12}{204.158}{709.986}{theory}%
\TextAt{8}{242.331}{709.986}{.}%
\TextAt{8}{97.159}{689.951}{The}%
\TextAt{8}{119.998}{689.951}{law}%
\TextAt{8}{139.921}{689.951}{of}%
\TextAt{8}{152.036}{689.951}{decomposition}%
\TextAt{8}{227.956}{689.951}{(10.3)}%
\TextAt{8}{260.575}{689.951}{in}%
\TextAt{8}{273.013}{689.951}{the}%
\TextAt{8}{291.955}{689.951}{cyclotomic}%
\TextAt{8}{348.951}{689.951}{field}%
\TextAt{8}{373.093}{689.951}{provides}%
\TextAt{8}{418.121}{689.951}{the}%
\TextAt{8}{437.063}{689.951}{proper}%
\TextAt{8}{473.245}{689.951}{explanation}%
\TextAt{8}{79.226}{672.905}{of}%
\TextAt{8}{92.556}{672.905}{Gauss’s}%
\TextAt{8}{135.093}{672.905}{reciprocity}%
\TextAt{8}{193.302}{672.905}{law}%
\TextAt{8}{214.439}{672.905}{(8.6).}%
\TextAt{8}{246.965}{672.905}{This}%
\TextAt{8}{273.685}{672.905}{is}%
\TextAt{8}{285.449}{672.905}{based}%
\TextAt{8}{318.027}{672.905}{on}%
\TextAt{8}{334.286}{672.905}{the}%
\TextAt{8}{354.442}{672.905}{following}%
\TextAt{12}{79.226}{646.893}{Proposition}%
\TextAt{12}{153.54}{646.893}{(10.5).}%
\TextAt{8}{196.979}{646.893}{Let}%
\TextAt{13}{219.341}{646.893}{\Character{96}}%
\TextAt{8}{229.544}{646.893}{and}%
\TextAt{13}{253.699}{646.893}{p}%
\TextAt{8}{264.865}{646.893}{be}%
\TextAt{8}{282.182}{646.893}{odd}%
\TextAt{8}{306.66}{646.893}{prime}%
\TextAt{8}{341.217}{646.893}{numbers,}%
\TextAt{13}{393.42}{646.893}{\Character{96}}%
\TextAt{3}{398.333}{651.231}{\Character{3}}%
\TextAt{9}{408.759}{646.893}{\Character{61}}%
\TextAt{9}{423.548}{646.893}{\Character{40}}%
\TextAt{16}{428.111}{646.893}{\Character{0}}%
\TextAt{9}{437.409}{646.893}{1\Character{41}}%
\TextAt{17}{447.815}{651.231}{\Character{40}}%
\TextAt{21}{451.108}{651.231}{\Character{96}}%
\TextAt{3}{454.567}{651.231}{\Character{0}}%
\TextAt{17}{461.154}{651.231}{1\Character{41}\Character{47}2}%
\TextAt{13}{477.648}{646.893}{\Character{96}}%
\TextAt{8}{482.56}{646.893}{,}%
\TextAt{8}{491.45}{646.893}{and}%
\TextAt{13}{515.606}{646.893}{\Character{16}}%
\TextAt{8}{526.911}{646.893}{a}%
\TextAt{8}{79.226}{629.847}{primitive}%
\TextAt{13}{129.306}{629.847}{\Character{96}}%
\TextAt{8}{134.218}{629.847}{-th}%
\TextAt{8}{153.072}{629.847}{root}%
\TextAt{8}{178.118}{629.847}{of}%
\TextAt{8}{191.448}{629.847}{unity.}%
\TextAt{8}{225.588}{629.847}{Then}%
\TextAt{8}{256.145}{629.847}{one}%
\TextAt{8}{277.605}{629.847}{has:}%
\TextAt{13}{219.06}{602.524}{p}%
\TextAt{8}{228.833}{602.524}{is}%
\TextAt{8}{240.597}{602.524}{totally}%
\TextAt{8}{278.004}{602.524}{split}%
\TextAt{8}{304.079}{602.524}{in}%
\TextAt{18}{317.731}{602.524}{Q}%
\TextAt{9}{327.03}{602.524}{\Character{40}}%
\TextAt{16}{331.582}{613.296}{\Character{112}}%
\Vector{q 0 G 0 g 0.478 w 269.545 -160.464 m 279.189 -160.464 l S Q}%
\TextAt{13}{341.545}{602.524}{\Character{96}}%
\TextAt{3}{346.457}{605.978}{\Character{3}}%
\TextAt{9}{351.189}{602.524}{\Character{41}}%
\NativeArrowAt{371.018}{602.524}{11.955}{21.926}{\Longleftrightarrow}%
\TextAt{13}{170.941}{582.49}{p}%
\TextAt{8}{180.713}{582.49}{splits}%
\TextAt{8}{211.402}{582.49}{in}%
\TextAt{18}{225.055}{582.49}{Q}%
\TextAt{9}{234.353}{582.49}{\Character{40}}%
\TextAt{13}{238.906}{582.49}{\Character{16}}%
\TextAt{9}{244.921}{582.49}{\Character{41}}%
\TextAt{8}{253.371}{582.49}{into}%
\TextAt{8}{277.114}{582.49}{an}%
\TextAt{8}{293.373}{582.49}{even}%
\TextAt{8}{320.033}{582.49}{number}%
\TextAt{8}{362.629}{582.49}{of}%
\TextAt{8}{375.959}{582.49}{prime}%
\TextAt{8}{409.123}{582.49}{ideals}%
\TextAt{13}{437.803}{582.49}{\Character{58}}%
\TextAt{19}{79.226}{537.017}{Proof.}%
\TextAt{8}{116}{537.017}{The}%
\TextAt{8}{142.893}{537.017}{little}%
\TextAt{8}{173.696}{537.017}{computation}%
\TextAt{8}{244.835}{537.017}{in}%
\TextAt{8}{261.328}{537.017}{§8,}%
\TextAt{8}{283.563}{537.017}{p.}%
\TextAt{8}{300.055}{537.017}{51}%
\TextAt{8}{318.508}{537.017}{has}%
\TextAt{8}{342.221}{537.017}{shown}%
\TextAt{8}{380.567}{537.017}{us}%
\TextAt{8}{398.422}{537.017}{that}%
\TextAt{13}{426.631}{537.017}{\Character{96}}%
\TextAt{3}{431.543}{541.355}{\Character{3}}%
\TextAt{9}{444.435}{537.017}{\Character{61}}%
\TextAt{13}{461.699}{537.017}{\Character{28}}%
\TextAt{17}{468.118}{541.355}{2}%
\TextAt{8}{479.587}{537.017}{with}%
\TextAt{13}{509.086}{537.017}{\Character{28}}%
\TextAt{9}{523.665}{537.017}{\Character{61}}%
\TextAt{20}{79.226}{528.937}{\Character{80}}%
\TextAt{21}{91.846}{516.484}{a}%
\NativeMathAt{96.344}{516.484}{7.97}{\in}%
\TextAt{17}{101.989}{516.484}{\Character{40}}%
\TextAt{25}{105.282}{516.484}{Z}%
\TextAt{17}{110.596}{516.484}{\Character{47}}%
\TextAt{21}{114.83}{516.484}{\Character{96}}%
\TextAt{25}{118.289}{516.484}{Z}%
\TextAt{17}{123.603}{516.484}{\Character{41}}%
\TextAt{23}{126.896}{518.753}{\Character{3}}%
\TextAt{20}{133.704}{529.655}{\Character{0}}%
\TextAt{21}{140.379}{524.678}{a}%
\Vector{q 0 G 0 g 0.478 w 68.379 -251.04 m 72.877 -251.04 l S Q}%
\TextAt{21}{140.898}{515.848}{\Character{96}}%
\TextAt{20}{146.072}{529.655}{\Character{1}}%
\TextAt{13}{153.544}{519.971}{\Character{16}}%
\TextAt{21}{159.56}{524.309}{a}%
\TextAt{8}{164.556}{519.971}{,}%
\TextAt{8}{172.435}{519.971}{so}%
\TextAt{8}{187.39}{519.971}{that}%
\TextAt{18}{213.343}{519.971}{Q}%
\TextAt{9}{222.641}{519.971}{\Character{40}}%
\TextAt{16}{227.194}{530.159}{\Character{112}}%
\Vector{q 0 G 0 g 0.478 w 165.156 -243.602 m 174.801 -243.602 l S Q}%
\TextAt{13}{237.156}{519.971}{\Character{96}}%
\TextAt{3}{242.068}{523.425}{\Character{3}}%
\TextAt{9}{246.801}{519.971}{\Character{41}}%
\NativeMathAt{255.669}{519.971}{11.955}{\subseteq}%
\TextAt{18}{269.284}{519.971}{Q}%
\TextAt{9}{278.583}{519.971}{\Character{40}}%
\TextAt{13}{283.135}{519.971}{\Character{16}}%
\TextAt{9}{289.151}{519.971}{\Character{41}}%
\TextAt{8}{293.703}{519.971}{.}%
\TextAt{8}{303.904}{519.971}{If}%
\TextAt{13}{316.18}{519.971}{p}%
\TextAt{8}{326.537}{519.971}{is}%
\TextAt{8}{338.885}{519.971}{totally}%
\TextAt{8}{376.877}{519.971}{split}%
\TextAt{8}{403.535}{519.971}{in}%
\TextAt{18}{417.772}{519.971}{Q}%
\TextAt{9}{427.071}{519.971}{\Character{40}}%
\TextAt{16}{431.623}{530.159}{\Character{112}}%
\Vector{q 0 G 0 g 0.478 w 369.586 -243.602 m 379.23 -243.602 l S Q}%
\TextAt{13}{441.586}{519.971}{\Character{96}}%
\TextAt{3}{446.498}{523.425}{\Character{3}}%
\TextAt{9}{451.23}{519.971}{\Character{41}}%
\TextAt{8}{455.783}{519.971}{,}%
\TextAt{8}{463.662}{519.971}{say}%
\TextAt{13}{484.474}{519.971}{p}%
\TextAt{9}{494.666}{519.971}{\Character{61}}%
\TextAt{14}{508.087}{519.971}{p}%
\TextAt{17}{514.07}{518.178}{1}%
\TextAt{14}{518.802}{519.971}{p}%
\TextAt{17}{524.785}{518.178}{2}%
\TextAt{8}{529.518}{519.971}{,}%
\TextAt{8}{79.226}{502.925}{then}%
\TextAt{8}{106.571}{502.925}{some}%
\TextAt{8}{136.582}{502.925}{automorphism}%
\TextAt{13}{214.736}{502.925}{\Character{27}}%
\TextAt{8}{226.4}{502.925}{of}%
\TextAt{18}{240.415}{502.925}{Q}%
\TextAt{9}{249.713}{502.925}{\Character{40}}%
\TextAt{13}{254.266}{502.925}{\Character{16}}%
\TextAt{9}{260.281}{502.925}{\Character{41}}%
\TextAt{8}{269.416}{502.925}{such}%
\TextAt{8}{296.497}{502.925}{that}%
\TextAt{13}{322.551}{502.925}{\Character{27}}%
\TextAt{14}{329.633}{502.925}{p}%
\TextAt{17}{335.616}{501.132}{1}%
\TextAt{9}{344.836}{502.925}{\Character{61}}%
\TextAt{14}{358.428}{502.925}{p}%
\TextAt{17}{364.411}{501.132}{2}%
\TextAt{8}{373.726}{502.925}{transforms}%
\TextAt{8}{432.752}{502.925}{the}%
\TextAt{8}{453.593}{502.925}{set}%
\TextAt{8}{472.545}{502.925}{of}%
\TextAt{8}{486.56}{502.925}{all}%
\TextAt{8}{503.503}{502.925}{prime}%
\TextAt{8}{79.226}{485.879}{ideals}%
\TextAt{8}{113.396}{485.879}{lying}%
\TextAt{8}{143.932}{485.879}{above}%
\TextAt{14}{178.699}{485.879}{p}%
\TextAt{17}{184.682}{484.086}{1}%
\TextAt{8}{194.904}{485.879}{bijectively}%
\TextAt{8}{252.423}{485.879}{into}%
\TextAt{8}{277.758}{485.879}{the}%
\TextAt{8}{299.506}{485.879}{set}%
\TextAt{8}{319.366}{485.879}{of}%
\TextAt{8}{334.288}{485.879}{prime}%
\TextAt{8}{369.044}{485.879}{ideals}%
\TextAt{8}{403.213}{485.879}{above}%
\TextAt{14}{437.981}{485.879}{p}%
\TextAt{17}{443.964}{484.086}{2}%
\TextAt{8}{448.697}{485.879}{.}%
\TextAt{8}{461.921}{485.879}{Therefore}%
\TextAt{8}{516.51}{485.879}{the}%
\TextAt{8}{79.226}{468.833}{number}%
\TextAt{8}{123.302}{468.833}{of}%
\TextAt{8}{138.111}{468.833}{prime}%
\TextAt{8}{172.754}{468.833}{ideals}%
\TextAt{8}{206.811}{468.833}{of}%
\TextAt{18}{221.62}{468.833}{Q}%
\TextAt{9}{230.918}{468.833}{\Character{40}}%
\TextAt{13}{235.471}{468.833}{\Character{16}}%
\TextAt{9}{241.486}{468.833}{\Character{41}}%
\TextAt{8}{251.415}{468.833}{above}%
\TextAt{13}{286.07}{468.833}{p}%
\TextAt{8}{297.322}{468.833}{is}%
\TextAt{8}{310.565}{468.833}{even.}%
\TextAt{8}{346.213}{468.833}{Now}%
\TextAt{8}{374.353}{468.833}{assume}%
\TextAt{8}{416.276}{468.833}{conversely}%
\TextAt{8}{473.753}{468.833}{that}%
\TextAt{8}{500.601}{468.833}{this}%
\TextAt{8}{524.903}{468.833}{is}%
\TextAt{8}{79.226}{451.787}{the}%
\TextAt{8}{99.904}{451.787}{case.}%
\TextAt{8}{130.791}{451.787}{Then}%
\TextAt{8}{161.871}{451.787}{the}%
\TextAt{8}{182.549}{451.787}{index}%
\TextAt{8}{214.608}{451.787}{of}%
\TextAt{8}{228.46}{451.787}{the}%
\TextAt{8}{249.138}{451.787}{decomposition}%
\TextAt{8}{326.794}{451.787}{group}%
\TextAt{13}{360.492}{451.787}{G}%
\TextAt{27}{369.725}{449.994}{p}%
\TextAt{8}{374.212}{451.787}{,}%
\TextAt{8}{382.014}{451.787}{or}%
\TextAt{8}{396.846}{451.787}{in}%
\TextAt{8}{411.02}{451.787}{other}%
\TextAt{8}{442.111}{451.787}{words,}%
\TextAt{8}{479.573}{451.787}{the}%
\TextAt{8}{500.251}{451.787}{degree}%
\TextAt{9}{79.226}{434.741}{\Character{91}}%
\TextAt{13}{82.478}{434.741}{Z}%
\TextAt{27}{90.49}{432.948}{p}%
\TextAt{9}{100.818}{434.741}{\Character{58}}%
\TextAt{18}{109.911}{434.741}{Q}%
\TextAt{9}{119.209}{434.741}{\Character{93}}%
\TextAt{8}{127.837}{434.741}{of}%
\TextAt{8}{142.646}{434.741}{the}%
\TextAt{8}{164.281}{434.741}{decomposition}%
\TextAt{8}{242.895}{434.741}{field}%
\TextAt{8}{269.731}{434.741}{of}%
\TextAt{8}{284.539}{434.741}{a}%
\TextAt{8}{295.774}{434.741}{prime}%
\TextAt{8}{330.416}{434.741}{ideal}%
\TextAt{14}{359.858}{434.741}{p}%
\TextAt{8}{371.218}{434.741}{of}%
\TextAt{18}{386.026}{434.741}{Q}%
\TextAt{9}{395.325}{434.741}{\Character{40}}%
\TextAt{13}{399.877}{434.741}{\Character{16}}%
\TextAt{9}{405.893}{434.741}{\Character{41}}%
\TextAt{8}{415.821}{434.741}{over}%
\TextAt{13}{442.346}{434.741}{p}%
\TextAt{8}{448.221}{434.741}{,}%
\TextAt{8}{457.219}{434.741}{is}%
\TextAt{8}{470.462}{434.741}{even.}%
\TextAt{8}{506.109}{434.741}{Since}%
\ErrataAt{79.19}{417.676}{455.81}{$G(\mathbb Q(\zeta)\mid\mathbb Q)$ is cyclic, it follows that $\mathbb Q(\sqrt{\ell^*})\subseteq Z_{\mathfrak p}$. By \hyperlink{nav.number.1.9.3}{(9.3)}, the ramification index}%
\TextAt{16}{285.469}{427.883}{\Character{112}}%
\Vector{q 0 G 0 g 0.478 w 223.432 -345.877 m 233.076 -345.877 l S Q}%
\ErrataAt{79.19}{400.63}{455.81}{and inertia degree of $\mathfrak p\cap Z_{\mathfrak p}$ over $\mathbb Q$ both equal $1$. The same holds in $\mathbb Q(\sqrt{\ell^*})$,}%
\TextAt{16}{349.631}{410.837}{\Character{112}}%
\Vector{q 0 G 0 g 0.478 w 287.593 -362.923 m 297.238 -362.923 l S Q}%
\ErrataAt{79.19}{383.584}{455.81}{so $p$ is totally split in that field.}%
\TextAt{16}{132.804}{393.792}{\Character{112}}%
\Vector{q 0 G 0 g 0.478 w 70.767 -379.969 m 80.411 -379.969 l S Q}%
\Vector{q 0 G 0 g 0.399 w 452.334 -390.397 m 452.334 -382.327 l S Q}%
\Vector{q 0 G 0 g 0.399 w 452.534 -382.526 m 459.707 -382.526 l S Q}%
\Vector{q 0 G 0 g 0.399 w 459.906 -390.397 m 459.906 -382.327 l S Q}%
\ErrataAt{97.1}{354.602}{437.9}{From this proposition we obtain the reciprocity law for two distinct odd prime numbers}%
\ErrataAt{79.19}{337.556}{455.81}{$\ell$ and $p$,}%
\TextAt{13}{254.126}{324.988}{\Character{96}}%
\Vector{q 0 G 0 g 0.478 w 181.645 -454.111 m 187.52 -454.111 l S Q}%
\TextAt{13}{253.645}{308.699}{p}%
\TextAt{20}{271.508}{330.17}{\Character{16}}%
\TextAt{13}{279.844}{324.988}{p}%
\Vector{q 0 G 0 g 0.478 w 207.844 -454.111 m 213.719 -454.111 l S Q}%
\TextAt{13}{280.325}{308.699}{\Character{96}}%
\TextAt{20}{286.914}{330.17}{\Character{17}}%
\TextAt{9}{297.375}{316.9}{\Character{61}}%
\TextAt{9}{309.785}{316.9}{\Character{40}}%
\TextAt{16}{314.353}{316.9}{\Character{0}}%
\TextAt{9}{323.651}{316.9}{1\Character{41}}%
\TextAt{22}{335.252}{325.197}{\Character{96}}%
\TextAt{23}{338.352}{325.197}{\Character{0}}%
\TextAt{24}{344.108}{325.197}{1}%
\Vector{q 0 G 0 g 0.359 w 263.252 -450.171 m 275.761 -450.171 l S Q}%
\TextAt{24}{339.68}{318.711}{2}%
\TextAt{22}{350.152}{325.529}{p}%
\TextAt{23}{353.992}{325.529}{\Character{0}}%
\TextAt{24}{359.748}{325.529}{1}%
\Vector{q 0 G 0 g 0.359 w 278.152 -450.171 m 291.401 -450.171 l S Q}%
\TextAt{24}{354.95}{318.711}{2}%
\TextAt{13}{365.095}{316.9}{\Character{59}}%
\TextAt{8}{79.226}{289.654}{as}%
\TextAt{8}{93.596}{289.654}{follows.}%
\TextAt{8}{136.583}{289.654}{It}%
\TextAt{8}{149.256}{289.654}{suffices}%
\TextAt{8}{189.042}{289.654}{to}%
\TextAt{8}{203.353}{289.654}{show}%
\TextAt{8}{232.356}{289.654}{that}%
\TextAt{20}{269.411}{276.105}{\Character{18}}%
\TextAt{13}{279.407}{267.336}{\Character{96}}%
\TextAt{3}{284.319}{271.674}{\Character{3}}%
\Vector{q 0 G 0 g 0.478 w 207.407 -511.763 m 217.052 -511.763 l S Q}%
\TextAt{13}{281.292}{251.047}{p}%
\TextAt{20}{290.247}{276.105}{\Character{19}}%
\TextAt{9}{302.368}{259.248}{\Character{61}}%
\TextAt{20}{314.794}{272.518}{\Character{16}}%
\TextAt{13}{323.129}{267.336}{p}%
\Vector{q 0 G 0 g 0.478 w 251.129 -511.763 m 257.005 -511.763 l S Q}%
\TextAt{13}{323.611}{251.047}{\Character{96}}%
\TextAt{20}{330.2}{272.518}{\Character{17}}%
\TextAt{13}{339.332}{259.248}{\Character{58}}%
\TextAt{8}{79.226}{221.19}{In}%
\TextAt{8}{93.847}{221.19}{fact,}%
\TextAt{8}{120.185}{221.19}{the}%
\TextAt{8}{140.341}{221.19}{completely}%
\TextAt{8}{199.196}{221.19}{elementary}%
\TextAt{8}{259.032}{221.19}{result}%
\TextAt{20}{291.61}{234.461}{\Character{16}}%
\TextAt{3}{299.945}{225.898}{\Character{0}}%
\TextAt{17}{306.532}{225.897}{1}%
\Vector{q 0 G 0 g 0.478 w 227.945 -549.821 m 238.766 -549.821 l S Q}%
\TextAt{21}{303.224}{217.068}{p}%
\TextAt{20}{311.961}{234.461}{\Character{17}}%
\TextAt{9}{322.422}{221.19}{\Character{61}}%
\TextAt{9}{334.832}{221.19}{\Character{40}}%
\TextAt{16}{339.4}{221.19}{\Character{0}}%
\TextAt{9}{348.698}{221.19}{1\Character{41}}%
\TextAt{17}{359.104}{225.529}{\Character{40}}%
\TextAt{21}{362.397}{225.529}{p}%
\TextAt{3}{366.66}{225.529}{\Character{0}}%
\TextAt{17}{373.246}{225.529}{1\Character{41}\Character{47}2}%
\TextAt{8}{393.638}{221.19}{(see}%
\TextAt{8}{417.106}{221.19}{§8,}%
\TextAt{8}{435.792}{221.19}{p.}%
\TextAt{8}{449.444}{221.19}{51)}%
\TextAt{8}{469.613}{221.19}{then}%
\TextAt{8}{496.273}{221.19}{gives}%
\TextAt{20}{166.667}{193.58}{\Character{16}}%
\TextAt{13}{175.002}{188.398}{p}%
\Vector{q 0 G 0 g 0.478 w 103.002 -590.701 m 108.877 -590.701 l S Q}%
\TextAt{13}{175.484}{172.109}{\Character{96}}%
\TextAt{20}{182.073}{193.58}{\Character{17}}%
\TextAt{9}{192.534}{180.31}{\Character{61}}%
\TextAt{20}{204.959}{197.167}{\Character{18}}%
\TextAt{13}{214.955}{188.398}{\Character{96}}%
\TextAt{3}{219.867}{192.736}{\Character{3}}%
\Vector{q 0 G 0 g 0.478 w 142.955 -590.701 m 152.599 -590.701 l S Q}%
\TextAt{13}{216.84}{172.109}{p}%
\TextAt{20}{225.795}{197.167}{\Character{19}}%
\TextAt{9}{237.916}{180.31}{\Character{61}}%
\TextAt{20}{250.342}{197.167}{\Character{18}}%
\TextAt{16}{260.337}{188.398}{\Character{0}}%
\TextAt{9}{269.636}{188.398}{1}%
\Vector{q 0 G 0 g 0.478 w 188.337 -590.701 m 203.489 -590.701 l S Q}%
\TextAt{13}{264.976}{172.109}{p}%
\TextAt{20}{276.684}{197.167}{\Character{19}}%
\TextAt{17}{285.485}{194.49}{\Character{40}}%
\TextAt{21}{288.778}{194.49}{\Character{96}}%
\TextAt{3}{292.237}{194.49}{\Character{0}}%
\TextAt{17}{298.824}{194.49}{1\Character{41}\Character{47}2}%
\TextAt{20}{317.31}{197.167}{\Character{18}}%
\TextAt{13}{327.788}{188.398}{\Character{96}}%
\Vector{q 0 G 0 g 0.478 w 255.306 -590.701 m 261.181 -590.701 l S Q}%
\TextAt{13}{327.306}{172.109}{p}%
\TextAt{20}{334.377}{197.167}{\Character{19}}%
\TextAt{9}{346.498}{180.31}{\Character{61}}%
\TextAt{20}{358.923}{197.167}{\Character{18}}%
\TextAt{13}{369.401}{188.398}{\Character{96}}%
\Vector{q 0 G 0 g 0.478 w 296.919 -590.701 m 302.795 -590.701 l S Q}%
\TextAt{13}{368.919}{172.109}{p}%
\TextAt{20}{375.99}{197.167}{\Character{19}}%
\TextAt{9}{386.783}{180.31}{\Character{40}}%
\TextAt{16}{391.335}{180.31}{\Character{0}}%
\TextAt{9}{400.634}{180.31}{1\Character{41}}%
\TextAt{22}{412.235}{188.939}{p}%
\TextAt{23}{416.075}{188.939}{\Character{0}}%
\TextAt{24}{421.831}{188.939}{1}%
\Vector{q 0 G 0 g 0.359 w 340.235 -586.761 m 353.484 -586.761 l S Q}%
\TextAt{24}{417.033}{182.121}{2}%
\TextAt{22}{427.875}{188.607}{\Character{96}}%
\TextAt{23}{430.975}{188.607}{\Character{0}}%
\TextAt{24}{436.731}{188.607}{1}%
\Vector{q 0 G 0 g 0.359 w 355.875 -586.761 m 368.384 -586.761 l S Q}%
\TextAt{24}{432.303}{182.121}{2}%
\TextAt{13}{442.077}{180.31}{\Character{58}}%
\TextAt{8}{79.226}{142.253}{By}%
\TextAt{8}{97.438}{142.253}{(8.5)}%
\TextAt{8}{125.262}{142.253}{and}%
\TextAt{8}{147.873}{142.253}{(10.5),}%
\TextAt{8}{184.837}{142.253}{we}%
\TextAt{8}{201.913}{142.253}{know}%
\TextAt{8}{232.331}{142.253}{that}%
\TextAt{20}{257.549}{155.523}{\Character{16}}%
\TextAt{21}{265.884}{146.96}{\Character{96}}%
\TextAt{23}{269.344}{149.772}{\Character{3}}%
\Vector{q 0 G 0 g 0.478 w 193.884 -628.759 m 201.661 -628.759 l S Q}%
\TextAt{21}{267.641}{138.13}{p}%
\TextAt{20}{274.856}{155.523}{\Character{17}}%
\TextAt{9}{285.317}{142.253}{\Character{61}}%
\TextAt{9}{297.726}{142.253}{1}%
\TextAt{8}{307.342}{142.253}{if}%
\TextAt{8}{317.914}{142.253}{and}%
\TextAt{8}{340.526}{142.253}{only}%
\TextAt{8}{366.066}{142.253}{if}%
\TextAt{13}{376.639}{142.253}{p}%
\TextAt{8}{386.26}{142.253}{decomposes}%
\TextAt{8}{449.638}{142.253}{in}%
\TextAt{8}{463.14}{142.253}{the}%
\TextAt{8}{483.145}{142.253}{field}%
\TextAt{18}{508.351}{142.253}{Q}%
\TextAt{9}{517.649}{142.253}{\Character{40}}%
\TextAt{13}{522.202}{142.253}{\Character{16}}%
\TextAt{9}{528.217}{142.253}{\Character{41}}%
\TextAt{8}{79.226}{125.207}{of}%
\TextAt{13}{93.722}{125.207}{\Character{96}}%
\TextAt{8}{98.634}{125.207}{-th}%
\TextAt{8}{118.653}{125.207}{roots}%
\TextAt{8}{149.479}{125.207}{of}%
\TextAt{8}{163.975}{125.207}{unity}%
\TextAt{8}{195.71}{125.207}{into}%
\TextAt{8}{220.619}{125.207}{an}%
\TextAt{8}{238.044}{125.207}{even}%
\TextAt{8}{265.869}{125.207}{number}%
\TextAt{8}{309.631}{125.207}{of}%
\TextAt{8}{324.127}{125.207}{prime}%
\TextAt{8}{358.456}{125.207}{ideals.}%
\TextAt{8}{399.082}{125.207}{By}%
\TextAt{8}{418.611}{125.207}{(10.3),}%
\TextAt{8}{457.153}{125.207}{this}%
\TextAt{8}{481.141}{125.207}{number}%
\TextAt{8}{524.903}{125.207}{is}%
\TextAt{13}{79.226}{108.161}{r}%
\TextAt{9}{88.234}{108.161}{\Character{61}}%
\TextAt{21}{101.942}{112.868}{\Character{96}}%
\TextAt{3}{105.402}{112.868}{\Character{0}}%
\TextAt{17}{111.988}{112.868}{1}%
\Vector{q 0 G 0 g 0.478 w 29.942 -662.851 m 44.222 -662.851 l S Q}%
\TextAt{21}{106.609}{104.038}{f}%
\TextAt{8}{117.418}{108.161}{,}%
\TextAt{8}{124.631}{108.161}{where}%
\TextAt{13}{158.491}{108.161}{f}%
\TextAt{8}{169.486}{108.161}{is}%
\TextAt{8}{181.301}{108.161}{the}%
\TextAt{8}{201.509}{108.161}{smallest}%
\TextAt{8}{246.559}{108.161}{positive}%
\TextAt{8}{289.924}{108.161}{integer}%
\TextAt{8}{328.674}{108.161}{such}%
\TextAt{8}{355.122}{108.161}{that}%
\TextAt{13}{380.542}{108.161}{p}%
\TextAt{21}{386.417}{112.499}{f}%
\NativeMathAt{395.269}{108.161}{11.955}{\equiv}%
\TextAt{9}{407.976}{108.161}{1}%
\TextAt{9}{419.13}{108.161}{\Character{40}}%
\TextAt{8}{423.694}{108.161}{mod}%
\TextAt{13}{449.797}{108.161}{\Character{96}}%
\TextAt{9}{454.709}{108.161}{\Character{41}}%
\TextAt{8}{459.261}{108.161}{,}%
\TextAt{8}{466.474}{108.161}{i.e.,}%
\TextAt{13}{488.643}{108.161}{r}%
\TextAt{8}{498.192}{108.161}{is}%
\TextAt{8}{510.007}{108.161}{even}%
\TextAt{8}{79.226}{91.115}{if}%
\TextAt{8}{90.382}{91.115}{and}%
\TextAt{8}{113.577}{91.115}{only}%
\TextAt{8}{139.7}{91.115}{if}%
\TextAt{13}{150.856}{91.115}{f}%
\TextAt{8}{162.232}{91.115}{is}%
\TextAt{8}{174.428}{91.115}{a}%
\TextAt{8}{184.615}{91.115}{divisor}%
\TextAt{8}{223.16}{91.115}{of}%
\TextAt{21}{238.118}{95.822}{\Character{96}}%
\TextAt{3}{241.577}{95.822}{\Character{0}}%
\TextAt{17}{248.164}{95.822}{1}%
\Vector{q 0 G 0 g 0.478 w 166.118 -679.897 m 180.398 -679.897 l S Q}%
\TextAt{17}{243.141}{86.992}{2}%
\TextAt{8}{253.593}{91.115}{.}%
\TextAt{8}{263.338}{91.115}{But}%
\TextAt{8}{287.011}{91.115}{this}%
\TextAt{8}{310.265}{91.115}{is}%
\TextAt{8}{322.461}{91.115}{tantamount}%
\TextAt{8}{386.65}{91.115}{to}%
\TextAt{8}{401.392}{91.115}{the}%
\TextAt{8}{421.981}{91.115}{condition}%
\TextAt{13}{473.796}{91.115}{p}%
\TextAt{17}{479.671}{95.453}{\Character{40}}%
\TextAt{21}{482.964}{95.453}{\Character{96}}%
\TextAt{3}{486.424}{95.453}{\Character{0}}%
\TextAt{17}{493.01}{95.453}{1\Character{41}\Character{47}2}%
\NativeMathAt{513.561}{91.115}{11.955}{\equiv}%
\TextAt{9}{526.916}{91.115}{1}%
\TextAt{9}{79.226}{74.069}{\Character{40}}%
\TextAt{8}{83.778}{74.069}{mod}%
\TextAt{13}{109.881}{74.069}{\Character{96}}%
\TextAt{9}{114.793}{74.069}{\Character{41}}%
\TextAt{8}{119.345}{74.069}{.}%
\TextAt{8}{127.846}{74.069}{Since}%
\TextAt{8}{158.422}{74.069}{an}%
\TextAt{8}{174.698}{74.069}{element}%
\TextAt{8}{217.958}{74.069}{in}%
\TextAt{8}{231.628}{74.069}{the}%
\TextAt{8}{251.803}{74.069}{cyclic}%
\TextAt{8}{284.004}{74.069}{group}%
\TextAt{18}{317.197}{74.069}{F}%
\TextAt{3}{324.503}{78.407}{\Character{3}}%
\TextAt{21}{324.503}{70.96}{\Character{96}}%
\TextAt{8}{333.15}{74.069}{has}%
\TextAt{8}{354.042}{74.069}{an}%
\TextAt{8}{370.318}{74.069}{order}%
\TextAt{8}{400.906}{74.069}{dividing}%
\TextAt{21}{447.321}{78.776}{\Character{96}}%
\TextAt{3}{450.781}{78.776}{\Character{0}}%
\TextAt{17}{457.367}{78.776}{1}%
\Vector{q 0 G 0 g 0.478 w 375.321 -696.943 m 389.601 -696.943 l S Q}%
\TextAt{17}{452.344}{69.946}{2}%
\TextAt{8}{466.712}{74.069}{if}%
\TextAt{8}{477.453}{74.069}{and}%
\TextAt{8}{500.234}{74.069}{only}%
\TextAt{8}{525.943}{74.069}{if}%
\TextAt{8}{300.14}{43.033}{59}%
\Destination{78.23}{809.815}{nav.page.59}%
\Destination{79.36}{597.455}{nav.print.64}%
\Destination{79.36}{660.855}{nav.number.1.10.5}%
\Hotlink{266.507}{534.855}{6.858}{11.047}{nav.section.1.8}%
\Hotlink{299.555}{534.855}{12.716}{11.047}{nav.print.51}%
\Hotlink{422.285}{218.955}{6.858}{11.047}{nav.section.1.8}%
\Hotlink{448.944}{218.955}{12.716}{11.047}{nav.print.51}%
\Hotlink{227.456}{687.755}{30.936}{11.047}{nav.number.1.10.3}%
\Hotlink{213.939}{670.755}{25.078}{11.047}{nav.number.1.8.6}%
\Hotlink{96.938}{140.055}{25.078}{11.047}{nav.number.1.8.5}%
\Hotlink{147.373}{140.055}{30.936}{11.047}{nav.number.1.10.5}%
\Hotlink{418.111}{123.055}{30.936}{11.047}{nav.number.1.10.3}%
\Destination{79.226}{371.602}{errata.58}%
\Destination{79.226}{435}{errata.183}%
}

% Page 72 (body page 60)
\Page{%
\TextAt{2}{79.226}{796.943}{I.}%
\TextAt{2}{89.83}{796.943}{Algebraic}%
\TextAt{2}{138.648}{796.943}{Integers}%
\Vector{q 0 G 0 g 0.299 w 7.226 18.729 m 460.769 18.729 l S Q}%
\TextAt{8}{79.226}{759.33}{it}%
\TextAt{8}{90.604}{759.33}{belongs}%
\TextAt{8}{132.288}{759.33}{to}%
\TextAt{18}{146.272}{759.33}{F}%
\TextAt{3}{153.578}{763.669}{\Character{3}}%
\TextAt{17}{157.812}{763.669}{2}%
\TextAt{21}{153.578}{756.221}{\Character{96}}%
\TextAt{8}{162.544}{759.33}{,}%
\TextAt{8}{169.432}{759.33}{the}%
\TextAt{8}{189.262}{759.33}{last}%
\TextAt{8}{211.112}{759.33}{congruence}%
\TextAt{8}{271.267}{759.33}{is}%
\TextAt{8}{282.704}{759.33}{equivalent}%
\TextAt{8}{337.993}{759.33}{to}%
\TextAt{20}{351.977}{769.014}{\Character{0}}%
\TextAt{21}{358.652}{764.586}{p}%
\Vector{q 0 G 0 g 0.478 w 286.652 -11.681 m 290.915 -11.681 l S Q}%
\TextAt{21}{359.054}{755.208}{\Character{96}}%
\TextAt{20}{364.111}{769.014}{\Character{1}}%
\TextAt{9}{372.911}{759.33}{\Character{61}}%
\TextAt{9}{385.321}{759.33}{1}%
\TextAt{8}{391.189}{759.33}{.}%
\TextAt{8}{399.529}{759.33}{So}%
\TextAt{8}{415.461}{759.33}{we}%
\TextAt{8}{432.362}{759.33}{do}%
\TextAt{8}{448.295}{759.33}{have}%
\TextAt{20}{474.963}{772.601}{\Character{16}}%
\TextAt{21}{483.299}{764.037}{\Character{96}}%
\TextAt{23}{486.758}{766.85}{\Character{3}}%
\Vector{q 0 G 0 g 0.478 w 411.299 -11.681 m 419.075 -11.681 l S Q}%
\TextAt{21}{485.055}{755.208}{p}%
\TextAt{20}{492.27}{772.601}{\Character{17}}%
\TextAt{9}{502.731}{759.33}{\Character{61}}%
\TextAt{20}{515.157}{769.014}{\Character{0}}%
\TextAt{21}{521.832}{764.586}{p}%
\Vector{q 0 G 0 g 0.478 w 449.832 -11.681 m 454.094 -11.681 l S Q}%
\TextAt{21}{522.233}{755.208}{\Character{96}}%
\TextAt{20}{527.29}{769.014}{\Character{1}}%
\TextAt{8}{79.226}{742.284}{as}%
\TextAt{8}{93.596}{742.284}{claimed.}%
\TextAt{8}{97.159}{722.25}{Historically,}%
\TextAt{8}{161.363}{722.25}{Gauss’s}%
\TextAt{8}{203.188}{722.25}{reciprocity}%
\TextAt{8}{260.687}{722.25}{law}%
\TextAt{8}{281.112}{722.25}{marked}%
\TextAt{8}{322.029}{722.25}{the}%
\TextAt{8}{341.474}{722.25}{beginning}%
\TextAt{8}{394.418}{722.25}{of}%
\TextAt{8}{407.037}{722.25}{algebraic}%
\TextAt{8}{455.76}{722.25}{number}%
\TextAt{8}{497.645}{722.25}{theory.}%
\TextAt{8}{79.226}{705.204}{It}%
\TextAt{8}{93.462}{705.204}{was}%
\TextAt{8}{117.524}{705.204}{discovered}%
\TextAt{8}{175.408}{705.204}{by}%
\TextAt{26}{193.23}{705.204}{Euler}%
\TextAt{8}{229.43}{705.204}{,}%
\TextAt{8}{238.533}{705.204}{but}%
\TextAt{8}{261.555}{705.204}{first}%
\TextAt{8}{287.243}{705.204}{proven}%
\TextAt{8}{326.86}{705.204}{by}%
\TextAt{26}{344.682}{705.204}{Gauss}%
\TextAt{8}{380.2}{705.204}{.}%
\TextAt{8}{393.337}{705.204}{The}%
\TextAt{8}{418.954}{705.204}{quest}%
\TextAt{8}{451.469}{705.204}{for}%
\TextAt{8}{470.917}{705.204}{similar}%
\TextAt{8}{510.915}{705.204}{laws}%
\TextAt{8}{79.226}{688.158}{concerning}%
\TextAt{8}{139.191}{688.158}{higher}%
\TextAt{8}{176.393}{688.158}{power}%
\TextAt{8}{211.97}{688.158}{residues,}%
\TextAt{8}{261.354}{688.158}{i.e.,}%
\TextAt{8}{285.249}{688.158}{the}%
\TextAt{8}{306.838}{688.158}{congruences}%
\TextAt{13}{373.366}{688.158}{x}%
\TextAt{21}{380.018}{692.496}{n}%
\NativeMathAt{391.416}{688.158}{11.955}{\equiv}%
\TextAt{13}{406.477}{688.158}{a}%
\TextAt{9}{417.935}{688.158}{\Character{40}}%
\TextAt{8}{422.487}{688.158}{mod}%
\TextAt{13}{448.589}{688.158}{p}%
\TextAt{9}{454.464}{688.158}{\Character{41}}%
\TextAt{8}{459.017}{688.158}{,}%
\TextAt{8}{467.956}{688.158}{with}%
\TextAt{13}{496.049}{688.158}{n}%
\TextAt{13}{508.781}{688.158}{\Character{62}}%
\TextAt{9}{523.665}{688.158}{2}%
\TextAt{8}{529.518}{688.158}{,}%
\TextAt{8}{79.226}{671.112}{dominated}%
\TextAt{8}{138.566}{671.112}{number}%
\TextAt{8}{182.615}{671.112}{theory}%
\TextAt{8}{220.818}{671.112}{for}%
\TextAt{8}{240.156}{671.112}{a}%
\TextAt{8}{251.364}{671.112}{long}%
\TextAt{8}{278.186}{671.112}{time.}%
\TextAt{8}{313.755}{671.112}{Since}%
\TextAt{8}{345.766}{671.112}{this}%
\TextAt{8}{370.041}{671.112}{problem}%
\TextAt{8}{417.019}{671.112}{required}%
\TextAt{8}{464.32}{671.112}{working}%
\TextAt{8}{510.007}{671.112}{with}%
\TextAt{8}{79.226}{654.066}{the}%
\TextAt{13}{100.513}{654.066}{n}%
\TextAt{8}{107.5}{654.066}{-th}%
\TextAt{8}{127.484}{654.066}{cyclotomic}%
\TextAt{8}{186.824}{654.066}{field,}%
\TextAt{26}{216.845}{654.066}{Kummer}%
\TextAt{8}{265.981}{654.066}{’s}%
\TextAt{8}{278.875}{654.066}{attempts}%
\TextAt{8}{329.499}{654.066}{to}%
\TextAt{8}{344.94}{654.066}{solve}%
\TextAt{8}{374.75}{654.066}{it}%
\TextAt{8}{387.584}{654.066}{led}%
\TextAt{8}{407.568}{654.066}{to}%
\TextAt{8}{423.008}{654.066}{his}%
\TextAt{8}{442.406}{654.066}{seminal}%
\TextAt{8}{485.869}{654.066}{discovery}%
\TextAt{8}{79.226}{637.02}{of}%
\TextAt{8}{93.632}{637.02}{ideal}%
\TextAt{8}{122.67}{637.02}{theory.}%
\TextAt{8}{166.217}{637.02}{We}%
\TextAt{8}{187.437}{637.02}{have}%
\TextAt{8}{215.508}{637.02}{developed}%
\TextAt{8}{270.883}{637.02}{the}%
\TextAt{8}{292.115}{637.02}{basics}%
\TextAt{8}{327.132}{637.02}{of}%
\TextAt{8}{341.537}{637.02}{this}%
\TextAt{8}{365.435}{637.02}{theory}%
\TextAt{8}{403.261}{637.02}{in}%
\TextAt{8}{417.989}{637.02}{the}%
\TextAt{8}{439.221}{637.02}{preceding}%
\TextAt{8}{492.971}{637.02}{sections}%
\TextAt{8}{79.226}{619.974}{and}%
\TextAt{8}{101.618}{619.974}{tested}%
\TextAt{8}{135.773}{619.974}{it}%
\TextAt{8}{147.107}{619.974}{successfully}%
\TextAt{8}{209.345}{619.974}{in}%
\TextAt{8}{222.627}{619.974}{the}%
\TextAt{8}{242.412}{619.974}{example}%
\TextAt{8}{287.889}{619.974}{of}%
\TextAt{8}{300.849}{619.974}{cyclotomic}%
\TextAt{8}{358.687}{619.974}{fields.}%
\TextAt{8}{393.086}{619.974}{The}%
\TextAt{8}{416.769}{619.974}{further}%
\TextAt{8}{455.743}{619.974}{development}%
\TextAt{8}{523.337}{619.974}{of}%
\TextAt{8}{79.226}{602.928}{this}%
\TextAt{8}{102.49}{602.928}{theory}%
\TextAt{8}{139.681}{602.928}{has}%
\TextAt{8}{160.997}{602.928}{led}%
\TextAt{8}{180.291}{602.928}{to}%
\TextAt{8}{195.043}{602.928}{a}%
\TextAt{8}{205.24}{602.928}{totally}%
\TextAt{8}{243.089}{602.928}{comprehensive}%
\TextAt{8}{321.633}{602.928}{generalization}%
\TextAt{8}{396.878}{602.928}{of}%
\TextAt{8}{410.649}{602.928}{Gauss’s}%
\TextAt{8}{453.627}{602.928}{reciprocity}%
\TextAt{8}{512.278}{602.928}{law,}%
\TextAt{12}{79.226}{585.882}{Artin’s}%
\TextAt{12}{125.675}{585.882}{reciprocity}%
\TextAt{12}{194.073}{585.882}{law}%
\TextAt{8}{213.703}{585.882}{,}%
\TextAt{8}{221.726}{585.882}{one}%
\TextAt{8}{243.884}{585.882}{of}%
\TextAt{8}{257.913}{585.882}{the}%
\TextAt{8}{278.768}{585.882}{high}%
\TextAt{8}{305.481}{585.882}{points}%
\TextAt{8}{341.364}{585.882}{in}%
\TextAt{8}{355.715}{585.882}{the}%
\TextAt{8}{376.571}{585.882}{history}%
\TextAt{8}{416.685}{585.882}{of}%
\TextAt{8}{430.714}{585.882}{number}%
\TextAt{8}{474.009}{585.882}{theory,}%
\TextAt{8}{513.904}{585.882}{and}%
\TextAt{8}{79.226}{568.836}{of}%
\TextAt{8}{93.294}{568.836}{compelling}%
\TextAt{8}{152.887}{568.836}{charm.}%
\TextAt{8}{195.098}{568.836}{This}%
\TextAt{8}{222.555}{568.836}{law}%
\TextAt{8}{244.43}{568.836}{is}%
\TextAt{8}{256.931}{568.836}{the}%
\TextAt{8}{277.825}{568.836}{main}%
\TextAt{8}{307.829}{568.836}{theorem}%
\TextAt{8}{354.092}{568.836}{of}%
\TextAt{8}{368.16}{568.836}{class}%
\TextAt{8}{396.334}{568.836}{field}%
\TextAt{8}{422.429}{568.836}{theory,}%
\TextAt{8}{462.373}{568.836}{which}%
\TextAt{8}{496.597}{568.836}{we}%
\TextAt{8}{514.562}{568.836}{will}%
\TextAt{8}{79.226}{551.79}{develop}%
\TextAt{8}{121.5}{551.79}{in}%
\TextAt{8}{135.152}{551.79}{chapters}%
\TextAt{8}{181.718}{551.79}{IV–VI.}%
\TextAt{12}{79.226}{525.892}{Exercise}%
\TextAt{12}{132.229}{525.892}{I.10.1.}%
\TextAt{12}{172.381}{525.892}{(Dirichlet’s}%
\TextAt{12}{242.432}{525.892}{Prime}%
\TextAt{12}{282.499}{525.892}{Number}%
\TextAt{12}{335.095}{525.892}{Theorem).}%
\TextAt{8}{401.743}{525.892}{For}%
\TextAt{8}{422.54}{525.892}{every}%
\TextAt{8}{453.26}{525.892}{natural}%
\TextAt{8}{494.071}{525.892}{number}%
\TextAt{13}{79.226}{508.846}{n}%
\TextAt{8}{90.111}{508.846}{there}%
\TextAt{8}{120.023}{508.846}{are}%
\TextAt{8}{139.534}{508.846}{infinitely}%
\TextAt{8}{188.634}{508.846}{many}%
\TextAt{8}{220.506}{508.846}{prime}%
\TextAt{8}{253.67}{508.846}{numbers}%
\TextAt{13}{300.881}{508.846}{p}%
\NativeMathAt{310.077}{508.846}{11.955}{\equiv}%
\TextAt{9}{322.696}{508.846}{1}%
\TextAt{9}{333.85}{508.846}{\Character{40}}%
\TextAt{8}{338.415}{508.846}{mod}%
\TextAt{13}{364.517}{508.846}{n}%
\TextAt{9}{371.504}{508.846}{\Character{41}}%
\TextAt{8}{376.057}{508.846}{.}%
\TextAt{12}{97.159}{488.811}{Hint:}%
\TextAt{8}{132.742}{488.811}{Assume}%
\TextAt{8}{176.126}{488.811}{there}%
\TextAt{8}{206.06}{488.811}{are}%
\TextAt{8}{225.593}{488.811}{only}%
\TextAt{8}{251.306}{488.811}{finitely}%
\TextAt{8}{290.673}{488.811}{many.}%
\TextAt{8}{326.183}{488.811}{Let}%
\TextAt{13}{347.174}{488.811}{P}%
\TextAt{8}{360.265}{488.811}{be}%
\TextAt{8}{376.211}{488.811}{their}%
\TextAt{8}{404.197}{488.811}{product}%
\TextAt{8}{448.118}{488.811}{and}%
\TextAt{8}{470.903}{488.811}{consider}%
\TextAt{8}{516.51}{488.811}{the}%
\TextAt{13}{79.226}{471.765}{n}%
\TextAt{8}{86.214}{471.765}{-th}%
\TextAt{8}{105.632}{471.765}{cyclotomic}%
\TextAt{8}{164.406}{471.765}{polynomial}%
\TextAt{13}{225.464}{471.765}{\Character{30}}%
\TextAt{21}{232.392}{469.972}{n}%
\TextAt{8}{238.029}{471.765}{.}%
\TextAt{8}{248.172}{471.765}{Not}%
\TextAt{8}{271.822}{471.765}{all}%
\TextAt{8}{288.646}{471.765}{numbers}%
\TextAt{13}{336.421}{471.765}{\Character{30}}%
\TextAt{21}{343.35}{469.972}{n}%
\TextAt{9}{348.986}{471.765}{\Character{40}}%
\TextAt{13}{353.538}{471.765}{xnP}%
\TextAt{9}{376.349}{471.765}{\Character{41}}%
\TextAt{8}{380.901}{471.765}{,}%
\TextAt{8}{388.756}{471.765}{for}%
\TextAt{13}{407.206}{471.765}{x}%
\NativeMathAt{418.142}{471.765}{11.955}{\in}%
\TextAt{18}{430.395}{471.765}{Z}%
\TextAt{8}{438.365}{471.765}{,}%
\TextAt{8}{446.22}{471.765}{can}%
\TextAt{8}{468.245}{471.765}{equal}%
\TextAt{9}{499.702}{471.765}{1}%
\TextAt{8}{505.555}{471.765}{.}%
\TextAt{8}{515.698}{471.765}{Let}%
\TextAt{13}{79.226}{454.719}{p}%
\TextAt{16}{88.671}{454.719}{\Character{106}}%
\TextAt{13}{95.562}{454.719}{\Character{30}}%
\TextAt{21}{102.49}{452.926}{n}%
\TextAt{9}{108.126}{454.719}{\Character{40}}%
\TextAt{13}{112.679}{454.719}{xnP}%
\TextAt{9}{135.489}{454.719}{\Character{41}}%
\TextAt{8}{144.085}{454.719}{for}%
\TextAt{8}{162.116}{454.719}{suitable}%
\TextAt{13}{205.899}{454.719}{x}%
\TextAt{8}{212.551}{454.719}{.}%
\TextAt{8}{221.437}{454.719}{Deduce}%
\TextAt{8}{263.032}{454.719}{a}%
\TextAt{8}{272.933}{454.719}{contradiction}%
\TextAt{8}{344.308}{454.719}{from}%
\TextAt{8}{372.095}{454.719}{this.}%
\TextAt{8}{399.906}{454.719}{(Dirichlet’s}%
\TextAt{8}{460.761}{454.719}{prime}%
\TextAt{8}{494.071}{454.719}{number}%
\TextAt{8}{79.226}{437.673}{theorem}%
\TextAt{8}{125.387}{437.673}{is}%
\TextAt{8}{137.786}{437.673}{valid}%
\TextAt{8}{166.719}{437.673}{more}%
\TextAt{8}{196.621}{437.673}{generally}%
\TextAt{8}{247.014}{437.673}{for}%
\TextAt{8}{265.534}{437.673}{prime}%
\TextAt{8}{299.333}{437.673}{numbers}%
\TextAt{13}{347.179}{437.673}{p}%
\NativeMathAt{357.458}{437.673}{11.955}{\equiv}%
\TextAt{13}{371.16}{437.673}{a}%
\TextAt{9}{382.619}{437.673}{\Character{40}}%
\TextAt{8}{387.171}{437.673}{mod}%
\TextAt{13}{413.273}{437.673}{n}%
\TextAt{9}{420.261}{437.673}{\Character{41}}%
\TextAt{8}{424.813}{437.673}{,}%
\TextAt{8}{432.756}{437.673}{provided}%
\TextAt{9}{481.523}{437.673}{\Character{40}}%
\TextAt{13}{486.076}{437.673}{a\Character{59}}%
\TextAt{13}{497.445}{437.673}{n}%
\TextAt{9}{504.452}{437.673}{\Character{41}}%
\TextAt{9}{513.394}{437.673}{\Character{61}}%
\TextAt{9}{526.892}{437.673}{1}%
\TextAt{8}{79.226}{420.628}{(see}%
\TextAt{8}{102.694}{420.628}{VII,}%
\TextAt{8}{127.382}{420.628}{(5.14)}%
\TextAt{8}{161.215}{420.628}{and}%
\TextAt{8}{183.977}{420.628}{VII,}%
\TextAt{8}{208.665}{420.628}{§13)).}%
\TextAt{12}{79.226}{388.867}{Exercise}%
\TextAt{12}{132.764}{388.867}{I.10.2.}%
\TextAt{8}{173.383}{388.867}{For}%
\TextAt{8}{194.646}{388.867}{every}%
\TextAt{8}{225.832}{388.867}{finite}%
\TextAt{8}{256.037}{388.867}{abelian}%
\TextAt{8}{296.978}{388.867}{group}%
\TextAt{13}{330.447}{388.867}{A}%
\TextAt{8}{343.414}{388.867}{there}%
\TextAt{8}{373.619}{388.867}{exists}%
\TextAt{8}{406.228}{388.867}{a}%
\TextAt{8}{416.277}{388.867}{Galois}%
\TextAt{8}{452.484}{388.867}{extension}%
\TextAt{13}{504.543}{388.867}{L}%
\TextAt{16}{516.329}{388.867}{\Character{106}}%
\TextAt{18}{523.471}{388.867}{Q}%
\TextAt{8}{79.226}{371.821}{with}%
\TextAt{8}{105.886}{371.821}{Galois}%
\TextAt{8}{141.8}{371.821}{group}%
\TextAt{13}{174.975}{371.821}{G}%
\TextAt{9}{184.209}{371.821}{\Character{40}}%
\TextAt{13}{188.761}{371.821}{L}%
\TextAt{16}{200.046}{371.821}{\Character{106}}%
\TextAt{18}{206.688}{371.821}{Q}%
\TextAt{9}{215.987}{371.821}{\Character{41}}%
\NativeMathAt{223.86}{371.821}{11.955}{\cong}%
\TextAt{13}{236.479}{371.821}{A}%
\TextAt{8}{245.254}{371.821}{.}%
\TextAt{12}{97.159}{351.786}{Hint:}%
\TextAt{8}{132.698}{351.786}{Use}%
\TextAt{8}{155.185}{351.786}{exercise}%
\TextAt{8}{198.487}{351.786}{1.}%
\TextAt{12}{79.226}{320.025}{Exercise}%
\TextAt{12}{133.008}{320.025}{I.10.3.}%
\TextAt{8}{173.839}{320.025}{Every}%
\TextAt{8}{207.998}{320.025}{quadratic}%
\TextAt{8}{260.867}{320.025}{number}%
\TextAt{8}{303.969}{320.025}{field}%
\TextAt{18}{329.831}{320.025}{Q}%
\TextAt{9}{339.13}{320.025}{\Character{40}}%
\TextAt{16}{343.682}{330.214}{\Character{112}}%
\Vector{q 0 G 0 g 0.478 w 281.645 -443.547 m 287.728 -443.547 l S Q}%
\TextAt{13}{353.645}{320.025}{d}%
\TextAt{9}{359.728}{320.025}{\Character{41}}%
\TextAt{8}{368.683}{320.025}{is}%
\TextAt{8}{380.952}{320.025}{contained}%
\TextAt{8}{434.467}{320.025}{in}%
\TextAt{8}{448.626}{320.025}{some}%
\TextAt{8}{478.457}{320.025}{cyclotomic}%
\TextAt{8}{79.226}{302.979}{field}%
\TextAt{18}{104.583}{302.979}{Q}%
\TextAt{9}{113.882}{302.979}{\Character{40}}%
\TextAt{13}{118.434}{302.979}{\Character{16}}%
\TextAt{21}{123.547}{301.186}{n}%
\TextAt{9}{129.183}{302.979}{\Character{41}}%
\TextAt{8}{133.735}{302.979}{,}%
\TextAt{13}{140.884}{302.979}{\Character{16}}%
\TextAt{21}{145.997}{301.186}{n}%
\TextAt{8}{155.531}{302.979}{a}%
\TextAt{8}{165.286}{302.979}{primitive}%
\TextAt{13}{215.367}{302.979}{n}%
\TextAt{8}{222.354}{302.979}{-th}%
\TextAt{8}{241.207}{302.979}{root}%
\TextAt{8}{266.254}{302.979}{of}%
\TextAt{8}{279.584}{302.979}{unity.}%
\TextAt{12}{79.226}{271.219}{Exercise}%
\TextAt{12}{133.26}{271.219}{I.10.4.}%
\TextAt{8}{174.31}{271.219}{Describe}%
\TextAt{8}{222.724}{271.219}{the}%
\TextAt{8}{243.605}{271.219}{quadratic}%
\TextAt{8}{296.693}{271.219}{subfields}%
\TextAt{8}{345.012}{271.219}{of}%
\TextAt{18}{359.066}{271.219}{Q}%
\TextAt{9}{368.365}{271.219}{\Character{40}}%
\TextAt{13}{372.917}{271.219}{\Character{16}}%
\TextAt{21}{378.03}{269.425}{n}%
\TextAt{9}{383.666}{271.219}{\Character{41}}%
\TextAt{16}{392.774}{271.219}{\Character{106}}%
\TextAt{18}{400.651}{271.219}{Q}%
\TextAt{8}{409.95}{271.219}{,}%
\TextAt{8}{418.005}{271.219}{in}%
\TextAt{8}{432.382}{271.219}{the}%
\TextAt{8}{453.263}{271.219}{case}%
\TextAt{8}{478.759}{271.219}{where}%
\TextAt{13}{513.293}{271.219}{n}%
\TextAt{8}{524.903}{271.219}{is}%
\TextAt{8}{79.226}{254.173}{odd.}%
\TextAt{12}{79.226}{222.412}{Exercise}%
\TextAt{12}{133.932}{222.412}{I.10.5.}%
\TextAt{8}{175.566}{222.412}{Show}%
\TextAt{8}{207.767}{222.412}{that}%
\TextAt{18}{234.445}{222.412}{Q}%
\TextAt{9}{243.744}{222.412}{\Character{40}}%
\TextAt{16}{248.296}{231.803}{\Character{112}}%
\Vector{q 0 G 0 g 0.478 w 186.259 -541.957 m 201.41 -541.957 l S Q}%
\TextAt{16}{258.259}{222.412}{\Character{0}}%
\TextAt{9}{267.557}{222.412}{1\Character{41}}%
\TextAt{8}{277.963}{222.412}{,}%
\TextAt{18}{286.748}{222.412}{Q}%
\TextAt{9}{296.047}{222.412}{\Character{40}}%
\TextAt{16}{300.599}{232.302}{\Character{112}}%
\Vector{q 0 G 0 g 0.478 w 238.562 -541.459 m 244.415 -541.459 l S Q}%
\TextAt{9}{310.562}{222.412}{2\Character{41}}%
\TextAt{8}{320.967}{222.412}{,}%
\TextAt{18}{329.753}{222.412}{Q}%
\TextAt{9}{339.051}{222.412}{\Character{40}}%
\TextAt{16}{343.603}{231.803}{\Character{112}}%
\Vector{q 0 G 0 g 0.478 w 281.566 -541.957 m 296.718 -541.957 l S Q}%
\TextAt{16}{353.566}{222.412}{\Character{0}}%
\TextAt{9}{362.865}{222.412}{2\Character{41}}%
\TextAt{8}{378.476}{222.412}{are}%
\TextAt{8}{399.296}{222.412}{the}%
\TextAt{8}{420.762}{222.412}{quadratic}%
\TextAt{8}{474.434}{222.412}{subfields}%
\TextAt{8}{523.337}{222.412}{of}%
\TextAt{18}{79.226}{205.366}{Q}%
\TextAt{9}{88.525}{205.366}{\Character{40}}%
\TextAt{13}{93.077}{205.366}{\Character{16}}%
\TextAt{21}{98.19}{203.573}{n}%
\TextAt{9}{103.826}{205.366}{\Character{41}}%
\TextAt{16}{111.699}{205.366}{\Character{106}}%
\TextAt{18}{118.341}{205.366}{Q}%
\TextAt{8}{131.537}{205.366}{for}%
\TextAt{13}{149.422}{205.366}{n}%
\TextAt{9}{159.73}{205.366}{\Character{61}}%
\TextAt{9}{172.14}{205.366}{2}%
\TextAt{21}{178.009}{209.704}{q}%
\TextAt{8}{182.574}{205.366}{,}%
\TextAt{13}{189.723}{205.366}{q}%
\NativeMathAt{198.663}{205.366}{11.955}{\geq}%
\TextAt{9}{211.282}{205.366}{3}%
\TextAt{8}{217.135}{205.366}{.}%
\TextAt{15}{79.226}{152.676}{§}%
\TextAt{15}{93.785}{152.676}{11}%
\TextAt{15}{130.385}{152.676}{Localization}%
\TextAt{8}{79.226}{120.79}{To}%
\TextAt{8}{96.537}{120.79}{“localize”}%
\TextAt{8}{148.959}{120.79}{means}%
\TextAt{8}{184.872}{120.79}{to}%
\TextAt{8}{199.265}{120.79}{form}%
\TextAt{8}{226.989}{120.79}{quotients,}%
\TextAt{8}{281.142}{120.79}{the}%
\TextAt{8}{301.382}{120.79}{most}%
\TextAt{8}{330.145}{120.79}{familiar}%
\TextAt{8}{373.482}{120.79}{case}%
\TextAt{8}{398.336}{120.79}{being}%
\TextAt{8}{429.957}{120.79}{the}%
\TextAt{8}{450.196}{120.79}{passage}%
\TextAt{8}{492.684}{120.79}{from}%
\TextAt{8}{520.408}{120.79}{an}%
\TextAt{8}{79.226}{103.744}{integral}%
\TextAt{8}{121.834}{103.744}{domain}%
\TextAt{13}{163.462}{103.744}{A}%
\TextAt{8}{176.135}{103.744}{to}%
\TextAt{8}{190.445}{103.744}{its}%
\TextAt{8}{206.764}{103.744}{field}%
\TextAt{8}{232.121}{103.744}{of}%
\TextAt{8}{245.451}{103.744}{fractions}%
\TextAt{13}{223.897}{75.144}{K}%
\TextAt{9}{238.036}{75.144}{\Character{61}}%
\TextAt{20}{250.461}{88.415}{\Character{110}}%
\TextAt{13}{259.627}{83.232}{a}%
\Vector{q 0 G 0 g 0.478 w 187.627 -695.867 m 193.772 -695.867 l S Q}%
\TextAt{13}{260.211}{66.944}{b}%
\TextAt{20}{270.865}{88.893}{\Character{12}}%
\TextAt{20}{270.865}{81.72}{\Character{12}}%
\TextAt{20}{270.865}{74.547}{\Character{12}}%
\TextAt{13}{278.747}{75.144}{a}%
\NativeMathAt{288.213}{75.144}{11.955}{\in}%
\TextAt{13}{299.504}{75.144}{A\Character{59}}%
\TextAt{13}{317.401}{75.144}{b}%
\NativeMathAt{325.719}{75.144}{11.955}{\in}%
\TextAt{13}{337.01}{75.144}{A}%
\TextAt{16}{348.442}{75.144}{\Character{110}}%
\TextAt{16}{357.074}{75.144}{\Character{102}}%
\TextAt{9}{363.054}{75.144}{0}%
\TextAt{16}{368.907}{75.144}{\Character{103}}%
\TextAt{20}{374.884}{88.415}{\Character{111}}%
\TextAt{13}{384.847}{75.144}{\Character{58}}%
\TextAt{8}{300.14}{43.033}{60}%
\Destination{78.23}{809.815}{nav.page.60}%
\Destination{79.23}{178.535}{nav.section.1.11}%
\Destination{79.226}{539.855}{nav.ex.1.10.1}%
\Destination{79.226}{402.855}{nav.ex.1.10.2}%
\Destination{79.226}{334.055}{nav.ex.1.10.3}%
\Destination{79.226}{285.255}{nav.ex.1.10.4}%
\Destination{79.226}{236.455}{nav.ex.1.10.5}%
\Destination{79.36}{403.855}{nav.print.65}%
\Destination{79.23}{178.535}{outline.16}%
\Hotlink{213.844}{418.455}{12.716}{11.047}{nav.section.7.13}%
\Hotlink{102.908}{418.4}{16.57}{11.047}{nav.chapter.7}%
\Hotlink{184.19}{418.4}{16.57}{11.047}{nav.chapter.7}%
\Hotlink{128.551}{418.4}{27.58}{11.047}{nav.number.7.5.14}%
\Hotlink{199.572}{349.6}{3.814}{11.047}{nav.ex.1.10.1}%
}

% Page 73 (body page 61)
\Page{%
\TextAt{2}{79.226}{796.943}{I.}%
\TextAt{2}{89.83}{796.943}{Algebraic}%
\TextAt{2}{138.648}{796.943}{Integers}%
\Vector{q 0 G 0 g 0.299 w 7.226 18.729 m 460.769 18.729 l S Q}%
\TextAt{8}{79.226}{761.124}{More}%
\TextAt{8}{110.036}{761.124}{generally,}%
\TextAt{8}{162.785}{761.124}{choosing}%
\TextAt{8}{210.906}{761.124}{instead}%
\TextAt{8}{251.867}{761.124}{of}%
\TextAt{13}{265.772}{761.124}{A}%
\TextAt{16}{277.597}{761.124}{\Character{110}}%
\TextAt{16}{286.623}{761.124}{\Character{102}}%
\TextAt{9}{292.601}{761.124}{0}%
\TextAt{16}{298.454}{761.124}{\Character{103}}%
\TextAt{8}{308.905}{761.124}{any}%
\TextAt{8}{331.598}{761.124}{nonempty}%
\TextAt{13}{386.809}{761.124}{S}%
\NativeMathAt{399.006}{761.124}{11.955}{\subseteq}%
\TextAt{13}{412.606}{761.124}{A}%
\TextAt{16}{424.431}{761.124}{\Character{110}}%
\TextAt{16}{433.457}{761.124}{\Character{102}}%
\TextAt{9}{439.435}{761.124}{0}%
\TextAt{16}{445.288}{761.124}{\Character{103}}%
\TextAt{8}{455.739}{761.124}{which}%
\TextAt{8}{489.801}{761.124}{is}%
\TextAt{8}{502.14}{761.124}{closed}%
\TextAt{8}{79.226}{744.078}{under}%
\TextAt{8}{112.39}{744.078}{multiplication,}%
\TextAt{8}{190.768}{744.078}{one}%
\TextAt{8}{212.227}{744.078}{again}%
\TextAt{8}{243.454}{744.078}{obtains}%
\TextAt{8}{284.496}{744.078}{a}%
\TextAt{8}{294.252}{744.078}{ring}%
\TextAt{8}{318.318}{744.078}{structure}%
\TextAt{8}{368.458}{744.078}{on}%
\TextAt{8}{384.717}{744.078}{the}%
\TextAt{8}{404.873}{744.078}{set}%
\TextAt{13}{217.318}{715.077}{AS}%
\TextAt{3}{233.988}{720.013}{\Character{0}}%
\TextAt{17}{240.575}{720.013}{1}%
\TextAt{9}{248.628}{715.077}{\Character{61}}%
\TextAt{20}{261.054}{728.347}{\Character{110}}%
\TextAt{13}{270.219}{723.164}{a}%
\Vector{q 0 G 0 g 0.478 w 198.219 -55.935 m 204.364 -55.935 l S Q}%
\TextAt{13}{270.535}{706.876}{s}%
\NativeMathAt{280.88}{715.077}{11.955}{\in}%
\TextAt{13}{292.171}{715.077}{K}%
\TextAt{20}{306.886}{728.825}{\Character{12}}%
\TextAt{20}{306.886}{721.652}{\Character{12}}%
\TextAt{20}{306.886}{714.479}{\Character{12}}%
\TextAt{13}{314.769}{715.077}{a}%
\NativeMathAt{324.235}{715.077}{11.955}{\in}%
\TextAt{13}{335.526}{715.077}{A\Character{59}}%
\TextAt{13}{353.423}{715.077}{s}%
\NativeMathAt{362.277}{715.077}{11.955}{\in}%
\TextAt{13}{373.568}{715.077}{S}%
\TextAt{20}{381.464}{728.347}{\Character{111}}%
\TextAt{13}{391.426}{715.077}{\Character{58}}%
\TextAt{8}{79.226}{686.075}{The}%
\TextAt{8}{103.678}{686.075}{most}%
\TextAt{8}{132.756}{686.075}{important}%
\TextAt{8}{188.446}{686.075}{special}%
\TextAt{8}{226.945}{686.075}{case}%
\TextAt{8}{252.114}{686.075}{of}%
\TextAt{8}{265.842}{686.075}{such}%
\TextAt{8}{292.637}{686.075}{a}%
\TextAt{8}{302.79}{686.075}{multiplicative}%
\TextAt{8}{377.011}{686.075}{subset}%
\TextAt{8}{413.298}{686.075}{is}%
\TextAt{8}{425.46}{686.075}{the}%
\TextAt{8}{446.014}{686.075}{complement}%
\TextAt{13}{511.771}{686.075}{S}%
\TextAt{9}{523.665}{686.075}{\Character{61}}%
\TextAt{13}{79.226}{669.029}{A}%
\TextAt{16}{90.909}{669.029}{\Character{110}}%
\TextAt{14}{99.795}{669.029}{p}%
\TextAt{8}{110.044}{669.029}{of}%
\TextAt{8}{123.743}{669.029}{a}%
\TextAt{8}{133.867}{669.029}{prime}%
\TextAt{8}{167.399}{669.029}{ideal}%
\TextAt{14}{195.731}{669.029}{p}%
\TextAt{8}{205.98}{669.029}{of}%
\TextAt{13}{219.679}{669.029}{A}%
\TextAt{8}{228.455}{669.029}{.}%
\TextAt{8}{238.009}{669.029}{In}%
\TextAt{8}{252.999}{669.029}{this}%
\TextAt{8}{276.19}{669.029}{case}%
\TextAt{8}{301.33}{669.029}{one}%
\TextAt{8}{323.158}{669.029}{writes}%
\TextAt{13}{358.053}{669.029}{A}%
\TextAt{27}{366.829}{667.236}{p}%
\TextAt{8}{375.582}{669.029}{instead}%
\TextAt{8}{416.335}{669.029}{of}%
\TextAt{13}{430.034}{669.029}{AS}%
\TextAt{3}{446.704}{673.368}{\Character{0}}%
\TextAt{17}{453.291}{673.368}{1}%
\TextAt{8}{458.023}{669.029}{,}%
\TextAt{8}{465.633}{669.029}{and}%
\TextAt{8}{488.765}{669.029}{one}%
\TextAt{8}{510.593}{669.029}{calls}%
\TextAt{8}{79.226}{651.983}{the}%
\TextAt{8}{98.887}{651.983}{ring}%
\TextAt{13}{122.457}{651.983}{A}%
\TextAt{27}{131.233}{650.19}{p}%
\TextAt{8}{139.122}{651.983}{the}%
\TextAt{12}{158.782}{651.983}{localization}%
\TextAt{8}{228.727}{651.983}{of}%
\TextAt{13}{241.561}{651.983}{A}%
\TextAt{8}{253.739}{651.983}{at}%
\TextAt{14}{267.554}{651.983}{p}%
\TextAt{8}{273.537}{651.983}{.}%
\TextAt{8}{281.82}{651.983}{When}%
\TextAt{8}{315.456}{651.983}{dealing}%
\TextAt{8}{355.286}{651.983}{with}%
\TextAt{8}{381.45}{651.983}{problems}%
\TextAt{8}{431.095}{651.983}{that}%
\TextAt{8}{455.968}{651.983}{involve}%
\TextAt{8}{494.829}{651.983}{a}%
\TextAt{8}{504.089}{651.983}{single}%
\TextAt{8}{79.226}{634.938}{prime}%
\TextAt{8}{111.846}{634.938}{ideal}%
\TextAt{14}{139.266}{634.938}{p}%
\TextAt{8}{148.603}{634.938}{of}%
\TextAt{13}{161.389}{634.938}{A}%
\TextAt{8}{173.519}{634.938}{at}%
\TextAt{8}{187.285}{634.938}{a}%
\TextAt{8}{196.497}{634.938}{time}%
\TextAt{8}{222.614}{634.938}{it}%
\TextAt{8}{233.774}{634.938}{is}%
\TextAt{8}{244.995}{634.938}{often}%
\TextAt{8}{274.04}{634.938}{expedient}%
\TextAt{8}{326.494}{634.938}{to}%
\TextAt{8}{340.261}{634.938}{replace}%
\TextAt{13}{379.385}{634.938}{A}%
\TextAt{8}{391.514}{634.938}{by}%
\TextAt{8}{407.229}{634.938}{the}%
\TextAt{8}{426.842}{634.938}{localization}%
\TextAt{13}{488.418}{634.938}{A}%
\TextAt{27}{497.193}{633.144}{p}%
\TextAt{8}{501.68}{634.938}{.}%
\TextAt{8}{509.947}{634.938}{This}%
\TextAt{8}{79.226}{617.892}{procedure}%
\TextAt{8}{133.702}{617.892}{forgets}%
\TextAt{8}{171.991}{617.892}{everything}%
\TextAt{8}{229.731}{617.892}{that}%
\TextAt{8}{255.276}{617.892}{has}%
\TextAt{8}{276.325}{617.892}{nothing}%
\TextAt{8}{319.432}{617.892}{to}%
\TextAt{8}{333.918}{617.892}{do}%
\TextAt{8}{350.353}{617.892}{with}%
\TextAt{14}{377.189}{617.892}{p}%
\TextAt{8}{383.172}{617.892}{,}%
\TextAt{8}{390.541}{617.892}{and}%
\TextAt{8}{413.479}{617.892}{brings}%
\TextAt{8}{448.839}{617.892}{out}%
\TextAt{8}{469.829}{617.892}{more}%
\TextAt{8}{499.271}{617.892}{clearly}%
\TextAt{8}{79.226}{600.846}{all}%
\TextAt{8}{95.485}{600.846}{the}%
\TextAt{8}{115.642}{600.846}{properties}%
\TextAt{8}{170.659}{600.846}{concerning}%
\TextAt{14}{229.192}{600.846}{p}%
\TextAt{8}{235.175}{600.846}{.}%
\TextAt{8}{243.623}{600.846}{For}%
\TextAt{8}{264.593}{600.846}{instance,}%
\TextAt{8}{313.429}{600.846}{the}%
\TextAt{8}{333.586}{600.846}{mapping}%
\TextAt{14}{280.568}{571.844}{q}%
\NativeArrowAt{289.736}{571.844}{11.955}{19.26}{\longmapsto}%
\TextAt{14}{312.318}{571.844}{q}%
\TextAt{13}{318.166}{571.844}{A}%
\TextAt{27}{326.941}{570.051}{p}%
\TextAt{8}{79.226}{542.843}{gives}%
\TextAt{8}{108.2}{542.843}{a}%
\TextAt{8}{118.249}{542.843}{1–1-correspondence}%
\TextAt{8}{221.691}{542.843}{between}%
\TextAt{8}{267.175}{542.843}{the}%
\TextAt{8}{287.625}{542.843}{prime}%
\TextAt{8}{321.082}{542.843}{ideals}%
\TextAt{14}{353.953}{542.843}{q}%
\NativeMathAt{363.621}{542.843}{11.955}{\subseteq}%
\TextAt{14}{376.74}{542.843}{p}%
\TextAt{8}{386.914}{542.843}{of}%
\TextAt{13}{400.537}{542.843}{A}%
\TextAt{8}{413.503}{542.843}{and}%
\TextAt{8}{436.559}{542.843}{the}%
\TextAt{8}{457.009}{542.843}{prime}%
\TextAt{8}{490.466}{542.843}{ideals}%
\TextAt{8}{523.337}{542.843}{of}%
\TextAt{13}{79.226}{525.797}{A}%
\TextAt{27}{88.002}{524.004}{p}%
\TextAt{8}{92.488}{525.797}{.}%
\TextAt{8}{100.937}{525.797}{More}%
\TextAt{8}{131.171}{525.797}{generally}%
\TextAt{8}{180.929}{525.797}{for}%
\TextAt{8}{198.814}{525.797}{any}%
\TextAt{8}{220.931}{525.797}{multiplicative}%
\TextAt{8}{294.754}{525.797}{set}%
\TextAt{13}{313.021}{525.797}{S}%
\TextAt{8}{320.917}{525.797}{,}%
\TextAt{8}{328.066}{525.797}{one}%
\TextAt{8}{349.525}{525.797}{has}%
\TextAt{8}{370.399}{525.797}{the}%
\TextAt{12}{79.226}{499.785}{Proposition}%
\TextAt{12}{151.937}{499.785}{(11.1).}%
\TextAt{8}{193.984}{499.785}{The}%
\TextAt{8}{218.038}{499.785}{mappings}%
\TextAt{14}{210.381}{470.784}{q}%
\NativeArrowAt{219.55}{470.784}{11.955}{19.26}{\longmapsto}%
\TextAt{14}{242.132}{470.784}{q}%
\TextAt{13}{247.979}{470.784}{S}%
\TextAt{3}{255.874}{475.72}{\Character{0}}%
\TextAt{17}{262.461}{475.72}{1}%
\TextAt{8}{291.104}{470.784}{and}%
\TextAt{14}{333.879}{470.784}{Q}%
\NativeArrowAt{347.087}{470.784}{11.955}{19.26}{\longmapsto}%
\TextAt{14}{369.669}{470.784}{Q}%
\TextAt{16}{382.212}{470.784}{\Character{92}}%
\TextAt{13}{392.839}{470.784}{A}%
\TextAt{8}{79.226}{441.783}{are}%
\TextAt{8}{99.274}{441.783}{mutually}%
\TextAt{8}{149.246}{441.783}{inverse}%
\TextAt{8}{188.542}{441.783}{1–1-correspondences}%
\TextAt{8}{296.843}{441.783}{between}%
\TextAt{8}{342.571}{441.783}{the}%
\TextAt{8}{363.264}{441.783}{prime}%
\TextAt{8}{396.965}{441.783}{ideals}%
\TextAt{14}{430.08}{441.783}{q}%
\NativeMathAt{440.164}{441.783}{11.955}{\subseteq}%
\TextAt{13}{453.699}{441.783}{A}%
\TextAt{16}{465.497}{441.783}{\Character{110}}%
\TextAt{13}{474.497}{441.783}{S}%
\TextAt{8}{486.827}{441.783}{of}%
\TextAt{13}{500.694}{441.783}{A}%
\TextAt{8}{513.904}{441.783}{and}%
\TextAt{8}{79.226}{424.737}{the}%
\TextAt{8}{99.383}{424.737}{prime}%
\TextAt{8}{132.546}{424.737}{ideals}%
\TextAt{14}{165.124}{424.737}{Q}%
\TextAt{8}{178.908}{424.737}{of}%
\TextAt{13}{192.238}{424.737}{AS}%
\TextAt{3}{208.909}{429.075}{\Character{0}}%
\TextAt{17}{215.495}{429.075}{1}%
\TextAt{8}{220.228}{424.737}{.}%
\TextAt{19}{79.226}{389.758}{Proof.}%
\TextAt{8}{116}{389.758}{If}%
\TextAt{14}{127.692}{389.758}{q}%
\NativeMathAt{136.861}{389.758}{11.955}{\subseteq}%
\TextAt{13}{149.48}{389.758}{A}%
\TextAt{16}{160.912}{389.758}{\Character{110}}%
\TextAt{13}{169.546}{389.758}{S}%
\TextAt{8}{181.339}{389.758}{is}%
\TextAt{8}{193.103}{389.758}{a}%
\TextAt{8}{202.858}{389.758}{prime}%
\TextAt{8}{236.022}{389.758}{ideal}%
\TextAt{8}{263.985}{389.758}{of}%
\TextAt{13}{277.315}{389.758}{A}%
\TextAt{8}{286.09}{389.758}{,}%
\TextAt{8}{293.24}{389.758}{then}%
\TextAt{14}{223.292}{360.757}{Q}%
\TextAt{9}{236.499}{360.757}{\Character{61}}%
\TextAt{14}{248.925}{360.757}{q}%
\TextAt{13}{254.772}{360.757}{S}%
\TextAt{3}{262.668}{365.693}{\Character{0}}%
\TextAt{17}{269.254}{365.693}{1}%
\TextAt{9}{277.307}{360.757}{\Character{61}}%
\TextAt{20}{289.733}{374.027}{\Character{110}}%
\TextAt{13}{298.898}{368.845}{q}%
\Vector{q 0 G 0 g 0.478 w 226.898 -410.254 m 232.517 -410.254 l S Q}%
\TextAt{13}{298.951}{352.556}{s}%
\TextAt{20}{309.61}{374.505}{\Character{12}}%
\TextAt{20}{309.61}{367.332}{\Character{12}}%
\TextAt{20}{309.61}{360.159}{\Character{12}}%
\TextAt{13}{317.493}{360.757}{q}%
\NativeMathAt{326.433}{360.757}{11.955}{\in}%
\TextAt{14}{337.724}{360.757}{q}%
\TextAt{13}{343.571}{360.757}{\Character{59}}%
\TextAt{13}{352.693}{360.757}{s}%
\NativeMathAt{361.548}{360.757}{11.955}{\in}%
\TextAt{13}{372.839}{360.757}{S}%
\TextAt{20}{380.734}{374.027}{\Character{111}}%
\TextAt{8}{79.226}{328.764}{is}%
\TextAt{8}{90.687}{328.764}{a}%
\TextAt{8}{100.14}{328.764}{prime}%
\TextAt{8}{133.001}{328.764}{ideal}%
\TextAt{8}{160.662}{328.764}{of}%
\TextAt{13}{173.69}{328.764}{AS}%
\TextAt{3}{190.36}{333.102}{\Character{0}}%
\TextAt{17}{196.947}{333.102}{1}%
\TextAt{8}{201.679}{328.764}{.}%
\TextAt{8}{210.027}{328.764}{Indeed,}%
\TextAt{8}{251.066}{328.764}{in}%
\TextAt{8}{264.416}{328.764}{obvious}%
\TextAt{8}{306.459}{328.764}{notation,}%
\TextAt{8}{356.309}{328.764}{the}%
\TextAt{8}{376.163}{328.764}{relation}%
\TextAt{21}{419.987}{333.471}{a}%
\Vector{q 0 G 0 g 0.478 w 347.987 -442.247 m 352.485 -442.247 l S Q}%
\TextAt{21}{420.278}{324.641}{s}%
\TextAt{21}{426.876}{333.471}{a}%
\TextAt{23}{431.374}{336.284}{\Character{48}}%
\Vector{q 0 G 0 g 0.478 w 354.876 -442.247 m 362.058 -442.247 l S Q}%
\TextAt{21}{427.167}{324.641}{s}%
\TextAt{23}{431.083}{326.911}{\Character{48}}%
\NativeMathAt{438.574}{328.764}{11.955}{\in}%
\TextAt{14}{449.865}{328.764}{Q}%
\TextAt{8}{459.752}{328.764}{,}%
\TextAt{8}{466.659}{328.764}{i.e.,}%
\TextAt{21}{489.718}{333.471}{aa}%
\TextAt{23}{498.714}{336.284}{\Character{48}}%
\Vector{q 0 G 0 g 0.478 w 417.718 -442.247 m 429.398 -442.247 l S Q}%
\TextAt{21}{490.3}{324.641}{ss}%
\TextAt{23}{498.132}{326.911}{\Character{48}}%
\TextAt{9}{505.915}{328.764}{\Character{61}}%
\TextAt{21}{521.896}{334.02}{q}%
\Vector{q 0 G 0 g 0.478 w 447.536 -442.247 m 456.322 -442.247 l S Q}%
\TextAt{21}{519.536}{324.641}{s}%
\TextAt{23}{523.452}{326.911}{\Character{48}\Character{48}}%
\TextAt{8}{529.518}{328.764}{,}%
\TextAt{8}{79.226}{311.718}{implies}%
\TextAt{8}{118.541}{311.718}{that}%
\TextAt{13}{143.498}{311.718}{s}%
\TextAt{3}{149.012}{316.056}{\Character{48}\Character{48}}%
\TextAt{13}{154.104}{311.718}{aa}%
\TextAt{3}{166.394}{316.056}{\Character{48}}%
\TextAt{9}{172.51}{311.718}{\Character{61}}%
\TextAt{13}{184.935}{311.718}{q}%
\TextAt{13}{190.542}{311.718}{ss}%
\TextAt{3}{201.583}{316.056}{\Character{48}}%
\NativeMathAt{207.699}{311.718}{11.955}{\in}%
\TextAt{14}{218.99}{311.718}{q}%
\TextAt{8}{224.837}{311.718}{.}%
\TextAt{8}{233.148}{311.718}{Therefore}%
\TextAt{13}{285.733}{311.718}{aa}%
\TextAt{3}{298.023}{316.056}{\Character{48}}%
\NativeMathAt{304.139}{311.718}{11.955}{\in}%
\TextAt{14}{315.43}{311.718}{q}%
\TextAt{8}{321.278}{311.718}{,}%
\TextAt{8}{328.097}{311.718}{because}%
\TextAt{13}{370.987}{311.718}{s}%
\TextAt{3}{376.501}{316.056}{\Character{48}\Character{48}}%
\NativeMathAt{384.914}{311.718}{11.955}{\notin}%
\TextAt{14}{396.205}{311.718}{q}%
\TextAt{8}{402.052}{311.718}{,}%
\TextAt{8}{408.872}{311.718}{and}%
\TextAt{8}{431.223}{311.718}{hence}%
\TextAt{13}{463.317}{311.718}{a}%
\TextAt{8}{472.947}{311.718}{or}%
\TextAt{13}{486.846}{311.718}{a}%
\TextAt{3}{492.991}{316.056}{\Character{48}}%
\TextAt{8}{499.271}{311.718}{belong}%
\TextAt{8}{79.226}{294.672}{to}%
\TextAt{14}{93.536}{294.672}{q}%
\TextAt{8}{99.384}{294.672}{,}%
\TextAt{8}{106.533}{294.672}{which}%
\TextAt{8}{140.02}{294.672}{shows}%
\TextAt{8}{173.637}{294.672}{that}%
\TextAt{21}{200.202}{299.379}{a}%
\Vector{q 0 G 0 g 0.478 w 128.202 -476.339 m 132.7 -476.339 l S Q}%
\TextAt{21}{200.493}{290.549}{s}%
\TextAt{8}{209.793}{294.672}{or}%
\TextAt{21}{225.299}{299.379}{a}%
\TextAt{23}{229.797}{302.192}{\Character{48}}%
\Vector{q 0 G 0 g 0.478 w 153.299 -476.339 m 160.481 -476.339 l S Q}%
\TextAt{21}{225.59}{290.549}{s}%
\TextAt{23}{229.506}{292.819}{\Character{48}}%
\TextAt{8}{237.574}{294.672}{belong}%
\TextAt{8}{274.97}{294.672}{to}%
\TextAt{14}{289.28}{294.672}{Q}%
\TextAt{8}{299.167}{294.672}{.}%
\TextAt{8}{307.615}{294.672}{Furthermore}%
\TextAt{8}{375.413}{294.672}{one}%
\TextAt{8}{396.872}{294.672}{has}%
\TextAt{14}{277.602}{265.671}{q}%
\TextAt{9}{286.771}{265.671}{\Character{61}}%
\TextAt{14}{299.196}{265.671}{Q}%
\TextAt{16}{311.74}{265.671}{\Character{92}}%
\TextAt{13}{322.366}{265.671}{A\Character{59}}%
\TextAt{8}{79.226}{236.67}{since}%
\TextAt{21}{109.09}{241.926}{q}%
\Vector{q 0 G 0 g 0.478 w 37.09 -534.341 m 41.157 -534.341 l S Q}%
\TextAt{21}{109.166}{232.547}{s}%
\TextAt{9}{117.673}{236.67}{\Character{61}}%
\TextAt{13}{130.099}{236.67}{a}%
\NativeMathAt{139.565}{236.67}{11.955}{\in}%
\TextAt{14}{150.855}{236.67}{Q}%
\TextAt{16}{163.399}{236.67}{\Character{92}}%
\TextAt{13}{174.026}{236.67}{A}%
\TextAt{8}{186.698}{236.67}{implies}%
\TextAt{13}{226.425}{236.67}{q}%
\TextAt{9}{235.365}{236.67}{\Character{61}}%
\TextAt{13}{247.791}{236.67}{as}%
\NativeMathAt{262.771}{236.67}{11.955}{\in}%
\TextAt{14}{274.062}{236.67}{q}%
\TextAt{8}{279.909}{236.67}{,}%
\TextAt{8}{287.058}{236.67}{whence}%
\TextAt{13}{328.017}{236.67}{a}%
\NativeMathAt{337.483}{236.67}{11.955}{\in}%
\TextAt{14}{348.774}{236.67}{q}%
\TextAt{8}{358.518}{236.67}{because}%
\TextAt{13}{401.82}{236.67}{s}%
\NativeMathAt{410.655}{236.67}{11.955}{\notin}%
\TextAt{14}{421.946}{236.67}{q}%
\TextAt{8}{427.793}{236.67}{.}%
\TextAt{8}{97.159}{216.635}{Conversely,}%
\TextAt{8}{158.862}{216.635}{let}%
\TextAt{14}{175.912}{216.635}{Q}%
\TextAt{8}{189.842}{216.635}{be}%
\TextAt{8}{205.912}{216.635}{an}%
\TextAt{8}{222.317}{216.635}{arbitrary}%
\TextAt{8}{272.232}{216.635}{prime}%
\TextAt{8}{305.541}{216.635}{ideal}%
\TextAt{8}{333.649}{216.635}{of}%
\TextAt{13}{347.125}{216.635}{AS}%
\TextAt{3}{363.796}{220.973}{\Character{0}}%
\TextAt{17}{370.382}{220.973}{1}%
\TextAt{8}{375.115}{216.635}{.}%
\TextAt{8}{384}{216.635}{Then}%
\TextAt{14}{414.703}{216.635}{q}%
\TextAt{9}{424.119}{216.635}{\Character{61}}%
\TextAt{14}{436.793}{216.635}{Q}%
\TextAt{16}{449.435}{216.635}{\Character{92}}%
\TextAt{13}{460.161}{216.635}{A}%
\TextAt{8}{472.98}{216.635}{is}%
\TextAt{8}{484.889}{216.635}{obviously}%
\TextAt{8}{79.226}{199.589}{a}%
\TextAt{8}{88.737}{199.589}{prime}%
\TextAt{8}{121.657}{199.589}{ideal}%
\TextAt{8}{149.375}{199.589}{of}%
\TextAt{13}{162.461}{199.589}{A}%
\TextAt{8}{171.236}{199.589}{,}%
\TextAt{8}{178.19}{199.589}{and}%
\TextAt{8}{200.709}{199.589}{one}%
\TextAt{8}{221.924}{199.589}{has}%
\TextAt{14}{242.553}{199.589}{q}%
\NativeMathAt{251.722}{199.589}{11.955}{\subseteq}%
\TextAt{13}{264.341}{199.589}{A}%
\TextAt{16}{275.273}{199.589}{\Character{110}}%
\TextAt{13}{283.408}{199.589}{S}%
\TextAt{8}{291.303}{199.589}{.}%
\TextAt{8}{299.67}{199.589}{In}%
\TextAt{8}{314.047}{199.589}{fact,}%
\TextAt{8}{340.189}{199.589}{if}%
\TextAt{14}{350.668}{199.589}{q}%
\TextAt{8}{360.169}{199.589}{were}%
\TextAt{8}{386.907}{199.589}{to}%
\TextAt{8}{400.973}{199.589}{contain}%
\TextAt{8}{442.034}{199.589}{an}%
\TextAt{13}{458.049}{199.589}{s}%
\NativeMathAt{466.884}{199.589}{11.955}{\in}%
\TextAt{13}{478.175}{199.589}{S}%
\TextAt{8}{486.07}{199.589}{,}%
\TextAt{8}{493.024}{199.589}{then}%
\TextAt{8}{519.44}{199.589}{we}%
\TextAt{8}{79.226}{182.543}{would}%
\TextAt{8}{113.37}{182.543}{have}%
\TextAt{9}{140.365}{182.543}{1}%
\TextAt{9}{149.523}{182.543}{\Character{61}}%
\TextAt{13}{161.964}{182.543}{s}%
\TextAt{17}{168.674}{187.25}{1}%
\Vector{q 0 G 0 g 0.478 w 96.674 -588.468 m 100.908 -588.468 l S Q}%
\TextAt{21}{168.833}{178.42}{s}%
\NativeMathAt{177.424}{182.543}{11.955}{\in}%
\TextAt{14}{188.715}{182.543}{Q}%
\TextAt{8}{202.499}{182.543}{because}%
\TextAt{17}{246.996}{187.25}{1}%
\Vector{q 0 G 0 g 0.478 w 174.996 -588.468 m 179.231 -588.468 l S Q}%
\TextAt{21}{247.156}{178.42}{s}%
\NativeMathAt{255.747}{182.543}{11.955}{\in}%
\TextAt{13}{267.038}{182.543}{AS}%
\TextAt{3}{283.709}{186.881}{\Character{0}}%
\TextAt{17}{290.295}{186.881}{1}%
\TextAt{8}{295.027}{182.543}{.}%
\TextAt{8}{303.476}{182.543}{Furthermore}%
\TextAt{8}{371.274}{182.543}{one}%
\TextAt{8}{392.733}{182.543}{has}%
\TextAt{14}{279.025}{153.542}{Q}%
\TextAt{9}{292.232}{153.542}{\Character{61}}%
\TextAt{14}{304.658}{153.542}{q}%
\TextAt{13}{310.505}{153.542}{S}%
\TextAt{3}{318.401}{158.478}{\Character{0}}%
\TextAt{17}{324.987}{158.478}{1}%
\TextAt{13}{329.719}{153.542}{\Character{58}}%
\TextAt{8}{79.226}{124.541}{For}%
\TextAt{8}{99.953}{124.541}{if}%
\TextAt{21}{111.629}{129.248}{a}%
\Vector{q 0 G 0 g 0.478 w 39.629 -646.47 m 44.127 -646.47 l S Q}%
\TextAt{21}{111.92}{120.418}{s}%
\NativeMathAt{120.643}{124.541}{11.955}{\in}%
\TextAt{14}{131.934}{124.541}{Q}%
\TextAt{8}{141.821}{124.541}{,}%
\TextAt{8}{148.776}{124.541}{then}%
\TextAt{13}{175.193}{124.541}{a}%
\TextAt{9}{184.659}{124.541}{\Character{61}}%
\TextAt{21}{198.28}{129.248}{a}%
\Vector{q 0 G 0 g 0.478 w 126.28 -646.47 m 130.778 -646.47 l S Q}%
\TextAt{21}{198.571}{120.418}{s}%
\TextAt{16}{206.133}{124.541}{\Character{1}}%
\TextAt{13}{211.614}{124.541}{s}%
\NativeMathAt{220.449}{124.541}{11.955}{\in}%
\TextAt{14}{231.74}{124.541}{Q}%
\TextAt{16}{243.786}{124.541}{\Character{92}}%
\TextAt{13}{253.916}{124.541}{A}%
\TextAt{9}{266.012}{124.541}{\Character{61}}%
\TextAt{14}{278.438}{124.541}{q}%
\TextAt{8}{284.285}{124.541}{,}%
\TextAt{8}{291.24}{124.541}{hence}%
\TextAt{21}{324.699}{129.248}{a}%
\Vector{q 0 G 0 g 0.478 w 252.699 -646.47 m 257.197 -646.47 l S Q}%
\TextAt{21}{324.99}{120.418}{s}%
\TextAt{9}{333.713}{124.541}{\Character{61}}%
\TextAt{13}{346.139}{124.541}{a}%
\TextAt{17}{353.479}{129.248}{1}%
\Vector{q 0 G 0 g 0.478 w 281.479 -646.47 m 285.713 -646.47 l S Q}%
\TextAt{21}{353.638}{120.418}{s}%
\NativeMathAt{362.23}{124.541}{11.955}{\in}%
\TextAt{14}{373.521}{124.541}{q}%
\TextAt{13}{379.368}{124.541}{S}%
\TextAt{3}{387.263}{128.879}{\Character{0}}%
\TextAt{17}{393.85}{128.879}{1}%
\TextAt{8}{398.582}{124.541}{.}%
\TextAt{8}{406.95}{124.541}{The}%
\TextAt{8}{430.76}{124.541}{mappings}%
\TextAt{14}{483.264}{124.541}{q}%
\NativeArrowAt{492.432}{124.541}{11.955}{11.955}{\mapsto}%
\TextAt{14}{507.708}{124.541}{q}%
\TextAt{13}{513.555}{124.541}{S}%
\TextAt{3}{521.451}{128.879}{\Character{0}}%
\TextAt{17}{528.037}{128.879}{1}%
\TextAt{8}{79.226}{107.495}{and}%
\TextAt{14}{101.989}{107.495}{Q}%
\NativeArrowAt{115.196}{107.495}{11.955}{11.955}{\mapsto}%
\TextAt{14}{130.472}{107.495}{Q}%
\TextAt{16}{143.016}{107.495}{\Character{92}}%
\TextAt{13}{153.643}{107.495}{A}%
\TextAt{8}{166.315}{107.495}{are}%
\TextAt{8}{185.826}{107.495}{therefore}%
\TextAt{8}{234.926}{107.495}{inverses}%
\TextAt{8}{278.299}{107.495}{of}%
\TextAt{8}{291.629}{107.495}{each}%
\TextAt{8}{317.967}{107.495}{other,}%
\TextAt{8}{351.788}{107.495}{which}%
\TextAt{8}{385.274}{107.495}{proves}%
\TextAt{8}{421.439}{107.495}{the}%
\TextAt{8}{441.595}{107.495}{proposition.}%
\Vector{q 0 G 0 g 0.399 w 452.334 -666.505 m 452.334 -658.435 l S Q}%
\Vector{q 0 G 0 g 0.399 w 452.534 -658.635 m 459.707 -658.635 l S Q}%
\Vector{q 0 G 0 g 0.399 w 452.534 -666.306 m 459.707 -666.306 l S Q}%
\Vector{q 0 G 0 g 0.399 w 459.906 -666.505 m 459.906 -658.435 l S Q}%
\TextAt{8}{300.14}{43.033}{61}%
\Destination{78.23}{809.815}{nav.page.61}%
\Destination{79.36}{404.755}{nav.print.66}%
\Destination{79.36}{513.755}{nav.number.1.11.1}%
}

% Page 74 (body page 62)
\Page{%
\TextAt{2}{79.226}{796.943}{I.}%
\TextAt{2}{89.83}{796.943}{Algebraic}%
\TextAt{2}{138.648}{796.943}{Integers}%
\Vector{q 0 G 0 g 0.299 w 7.226 18.729 m 460.769 18.729 l S Q}%
\TextAt{8}{97.159}{761.124}{Usually}%
\TextAt{13}{139.884}{761.124}{S}%
\TextAt{8}{152.068}{761.124}{will}%
\TextAt{8}{174.565}{761.124}{be}%
\TextAt{8}{190.88}{761.124}{the}%
\TextAt{8}{211.428}{761.124}{complement}%
\TextAt{8}{277.179}{761.124}{of}%
\TextAt{8}{290.9}{761.124}{a}%
\TextAt{8}{301.047}{761.124}{union}%
\TextAt{20}{333.957}{770.09}{\Character{83}}%
\TextAt{27}{343.919}{757.637}{p}%
\NativeMathAt{347.908}{757.637}{7.97}{\in}%
\TextAt{21}{353.554}{757.637}{X}%
\TextAt{14}{363.631}{761.124}{p}%
\TextAt{8}{373.903}{761.124}{over}%
\TextAt{8}{399.341}{761.124}{a}%
\TextAt{8}{409.488}{761.124}{set}%
\TextAt{13}{428.147}{761.124}{X}%
\TextAt{8}{443.091}{761.124}{of}%
\TextAt{8}{456.812}{761.124}{prime}%
\TextAt{8}{490.368}{761.124}{ideals}%
\TextAt{8}{523.337}{761.124}{of}%
\TextAt{13}{79.226}{744.078}{A}%
\TextAt{8}{88.001}{744.078}{.}%
\TextAt{8}{96.45}{744.078}{In}%
\TextAt{8}{111.071}{744.078}{this}%
\TextAt{8}{133.893}{744.078}{case}%
\TextAt{8}{158.665}{744.078}{one}%
\TextAt{8}{180.124}{744.078}{writes}%
\TextAt{13}{172.178}{713.978}{A}%
\TextAt{9}{180.954}{713.978}{\Character{40}}%
\TextAt{13}{185.506}{713.978}{X}%
\TextAt{9}{196.161}{713.978}{\Character{41}}%
\TextAt{9}{204.016}{713.978}{\Character{61}}%
\TextAt{20}{216.46}{730.835}{\Character{26}}%
\TextAt{13}{226.622}{722.065}{f}%
\Vector{q 0 G 0 g 0.478 w 154.622 -57.033 m 161.668 -57.033 l S Q}%
\TextAt{13}{227.128}{705.777}{g}%
\TextAt{20}{238.761}{731.313}{\Character{12}}%
\TextAt{20}{238.761}{724.14}{\Character{12}}%
\TextAt{20}{238.761}{716.967}{\Character{12}}%
\TextAt{20}{238.761}{709.793}{\Character{12}}%
\TextAt{13}{246.643}{713.978}{f}%
\TextAt{13}{253.027}{713.978}{\Character{59}}%
\TextAt{13}{258.264}{713.978}{g}%
\NativeMathAt{267.639}{713.978}{11.955}{\in}%
\TextAt{13}{278.93}{713.978}{A\Character{59}}%
\TextAt{13}{296.827}{713.978}{g}%
\NativeMathAt{306.202}{713.978}{11.955}{\not\equiv}%
\TextAt{9}{318.821}{713.978}{0}%
\TextAt{9}{336.61}{713.978}{\Character{40}}%
\TextAt{8}{341.181}{713.978}{mod}%
\TextAt{14}{367.283}{713.978}{p}%
\TextAt{9}{373.266}{713.978}{\Character{41}}%
\TextAt{8}{381.716}{713.978}{for}%
\TextAt{14}{399.601}{713.978}{p}%
\NativeMathAt{408.905}{713.978}{11.955}{\in}%
\TextAt{13}{420.196}{713.978}{X}%
\TextAt{20}{430.851}{730.835}{\Character{27}}%
\TextAt{8}{79.226}{680.5}{instead}%
\TextAt{8}{119.633}{680.5}{of}%
\TextAt{13}{132.986}{680.5}{AS}%
\TextAt{3}{149.657}{684.839}{\Character{0}}%
\TextAt{17}{156.243}{684.839}{1}%
\TextAt{8}{160.976}{680.5}{.}%
\TextAt{8}{169.492}{680.5}{The}%
\TextAt{8}{193.569}{680.5}{prime}%
\TextAt{8}{226.755}{680.5}{ideals}%
\TextAt{8}{259.355}{680.5}{of}%
\TextAt{13}{272.708}{680.5}{A}%
\TextAt{9}{281.484}{680.5}{\Character{40}}%
\TextAt{13}{286.036}{680.5}{X}%
\TextAt{9}{296.691}{680.5}{\Character{41}}%
\TextAt{8}{305.163}{680.5}{correspond}%
\TextAt{8}{364.759}{680.5}{by}%
\TextAt{8}{381.041}{680.5}{(11.1)}%
\TextAt{8}{414.896}{680.5}{1–1}%
\TextAt{8}{436.391}{680.5}{to}%
\TextAt{8}{450.724}{680.5}{the}%
\TextAt{8}{470.903}{680.5}{prime}%
\TextAt{8}{504.089}{680.5}{ideals}%
\TextAt{8}{79.226}{663.454}{of}%
\TextAt{13}{91.867}{663.454}{A}%
\TextAt{8}{103.85}{663.454}{which}%
\TextAt{8}{136.647}{663.454}{are}%
\TextAt{8}{155.468}{663.454}{contained}%
\TextAt{8}{207.788}{663.454}{in}%
\TextAt{20}{220.751}{672.421}{\Character{83}}%
\TextAt{27}{230.714}{659.967}{p}%
\NativeMathAt{234.702}{659.967}{7.97}{\in}%
\TextAt{21}{240.348}{659.967}{X}%
\TextAt{14}{250.425}{663.454}{p}%
\TextAt{8}{256.409}{663.454}{,}%
\TextAt{8}{263.006}{663.454}{all}%
\TextAt{8}{278.576}{663.454}{the}%
\TextAt{8}{298.042}{663.454}{others}%
\TextAt{8}{332.537}{663.454}{are}%
\TextAt{8}{351.358}{663.454}{being}%
\TextAt{8}{382.206}{663.454}{eliminated}%
\TextAt{8}{438.746}{663.454}{when}%
\TextAt{8}{468.614}{663.454}{passing}%
\TextAt{8}{509.027}{663.454}{from}%
\ErrataAt{79.22}{646.455}{453.64}{$A$ to $A(X)$. If $A$ is a Dedekind domain and $X$ is a finite set of nonzero primes,}%
\ErrataAt{79.22}{629.355}{453.64}{the nonzero primes of $A(X)$ are exactly the extensions of the primes in $X$.}%
\TextAt{8}{97.159}{609.327}{In}%
\TextAt{8}{111.274}{609.327}{the}%
\TextAt{8}{130.924}{609.327}{case}%
\TextAt{8}{155.189}{609.327}{that}%
\TextAt{13}{180.051}{609.327}{X}%
\TextAt{8}{194.097}{609.327}{consists}%
\TextAt{8}{236.701}{609.327}{of}%
\TextAt{8}{249.525}{609.327}{only}%
\TextAt{8}{274.71}{609.327}{one}%
\TextAt{8}{295.663}{609.327}{prime}%
\TextAt{8}{328.321}{609.327}{ideal}%
\TextAt{14}{355.777}{609.327}{p}%
\TextAt{8}{361.761}{609.327}{,}%
\TextAt{8}{368.505}{609.327}{the}%
\TextAt{8}{388.155}{609.327}{ring}%
\TextAt{13}{411.714}{609.327}{A}%
\TextAt{9}{420.489}{609.327}{\Character{40}}%
\TextAt{13}{425.042}{609.327}{X}%
\TextAt{9}{435.697}{609.327}{\Character{41}}%
\TextAt{8}{443.64}{609.327}{is}%
\TextAt{8}{454.898}{609.327}{the}%
\TextAt{8}{474.548}{609.327}{localization}%
\TextAt{13}{206.331}{576.358}{A}%
\TextAt{27}{215.106}{574.565}{p}%
\TextAt{9}{222.914}{576.358}{\Character{61}}%
\TextAt{20}{235.339}{593.215}{\Character{26}}%
\TextAt{13}{245.501}{584.446}{f}%
\Vector{q 0 G 0 g 0.478 w 173.501 -194.653 m 180.548 -194.653 l S Q}%
\TextAt{13}{246.007}{568.158}{g}%
\TextAt{20}{257.641}{593.693}{\Character{12}}%
\TextAt{20}{257.641}{586.52}{\Character{12}}%
\TextAt{20}{257.641}{579.347}{\Character{12}}%
\TextAt{20}{257.641}{572.174}{\Character{12}}%
\TextAt{13}{265.523}{576.358}{f}%
\TextAt{13}{271.907}{576.358}{\Character{59}}%
\TextAt{13}{277.144}{576.358}{g}%
\NativeMathAt{286.518}{576.358}{11.955}{\in}%
\TextAt{13}{297.809}{576.358}{A\Character{59}}%
\TextAt{13}{315.706}{576.358}{g}%
\NativeMathAt{325.081}{576.358}{11.955}{\not\equiv}%
\TextAt{9}{337.701}{576.358}{0}%
\TextAt{9}{355.49}{576.358}{\Character{40}}%
\TextAt{8}{360.061}{576.358}{mod}%
\TextAt{14}{386.163}{576.358}{p}%
\TextAt{9}{392.146}{576.358}{\Character{41}}%
\TextAt{20}{396.698}{593.215}{\Character{27}}%
\TextAt{8}{79.226}{543.939}{of}%
\TextAt{13}{92.556}{543.939}{A}%
\TextAt{8}{105.229}{543.939}{at}%
\TextAt{14}{119.539}{543.939}{p}%
\TextAt{8}{125.522}{543.939}{.}%
\TextAt{8}{133.971}{543.939}{Here}%
\TextAt{8}{161.599}{543.939}{we}%
\TextAt{8}{178.826}{543.939}{have}%
\TextAt{8}{205.821}{543.939}{the}%
\TextAt{12}{79.226}{518.014}{Corollary}%
\TextAt{12}{138.236}{518.014}{(11.2).}%
\TextAt{8}{179.606}{518.014}{If}%
\TextAt{14}{190.623}{518.014}{p}%
\TextAt{8}{199.827}{518.014}{is}%
\TextAt{8}{210.915}{518.014}{a}%
\TextAt{8}{219.994}{518.014}{prime}%
\TextAt{8}{252.482}{518.014}{ideal}%
\TextAt{8}{279.769}{518.014}{of}%
\TextAt{13}{292.422}{518.014}{A}%
\TextAt{8}{301.198}{518.014}{,}%
\TextAt{8}{307.806}{518.014}{then}%
\TextAt{13}{333.79}{518.014}{A}%
\TextAt{27}{342.565}{516.221}{p}%
\TextAt{8}{350.274}{518.014}{is}%
\TextAt{8}{361.361}{518.014}{a}%
\TextAt{12}{370.441}{518.014}{local}%
\TextAt{12}{401.248}{518.014}{ring}%
\TextAt{8}{424.693}{518.014}{,}%
\TextAt{8}{431.301}{518.014}{i.e.,}%
\TextAt{13}{452.865}{518.014}{A}%
\TextAt{27}{461.64}{516.221}{p}%
\TextAt{8}{469.349}{518.014}{has}%
\TextAt{8}{489.546}{518.014}{a}%
\TextAt{8}{498.626}{518.014}{unique}%
\TextAt{8}{79.226}{500.968}{maximal}%
\TextAt{8}{127.035}{500.968}{ideal,}%
\TextAt{8}{158.25}{500.968}{namely}%
\TextAt{14}{198.897}{500.968}{m}%
\TextAt{27}{208.064}{499.175}{p}%
\TextAt{9}{215.872}{500.968}{\Character{61}}%
\TextAt{14}{228.297}{500.968}{p}%
\TextAt{13}{234.28}{500.968}{A}%
\TextAt{27}{243.056}{499.175}{p}%
\TextAt{8}{247.543}{500.968}{.}%
\TextAt{8}{255.991}{500.968}{There}%
\TextAt{8}{289.8}{500.968}{is}%
\TextAt{8}{301.564}{500.968}{a}%
\TextAt{8}{311.319}{500.968}{canonical}%
\TextAt{8}{362.703}{500.968}{embedding}%
\TextAt{13}{264.101}{472.273}{A}%
\TextAt{9}{272.876}{472.273}{\Character{47}}%
\TextAt{14}{278.729}{472.273}{p}%
\TextAt{13}{288.033}{472.273}{\Character{44}}%
\NativeArrowAt{289.292}{472.273}{11.955}{19.26}{\longrightarrow}%
\TextAt{13}{311.874}{472.273}{A}%
\TextAt{27}{320.65}{470.48}{p}%
\TextAt{9}{325.137}{472.273}{\Character{47}}%
\TextAt{14}{330.99}{472.273}{m}%
\TextAt{27}{340.156}{470.48}{p}%
\TextAt{13}{344.643}{472.273}{\Character{59}}%
\TextAt{8}{79.226}{443.578}{identifying}%
\TextAt{13}{137.241}{443.578}{A}%
\TextAt{27}{146.016}{441.785}{p}%
\TextAt{9}{150.503}{443.578}{\Character{47}}%
\TextAt{14}{156.356}{443.578}{m}%
\TextAt{27}{165.523}{441.785}{p}%
\TextAt{8}{173.712}{443.578}{with}%
\TextAt{8}{200.176}{443.578}{the}%
\TextAt{8}{220.138}{443.578}{field}%
\TextAt{8}{245.299}{443.578}{of}%
\TextAt{8}{258.434}{443.578}{fractions}%
\TextAt{8}{306.107}{443.578}{of}%
\TextAt{13}{319.242}{443.578}{A}%
\TextAt{9}{328.017}{443.578}{\Character{47}}%
\TextAt{14}{333.87}{443.578}{p}%
\TextAt{8}{339.853}{443.578}{.}%
\TextAt{8}{348.237}{443.578}{In}%
\TextAt{8}{362.663}{443.578}{particular,}%
\TextAt{8}{419.748}{443.578}{if}%
\TextAt{14}{430.276}{443.578}{p}%
\TextAt{8}{439.961}{443.578}{is}%
\TextAt{8}{451.53}{443.578}{a}%
\TextAt{8}{461.09}{443.578}{maximal}%
\TextAt{8}{508.704}{443.578}{ideal}%
\TextAt{8}{79.226}{426.532}{of}%
\TextAt{13}{92.556}{426.532}{A}%
\TextAt{8}{101.332}{426.532}{,}%
\TextAt{8}{108.481}{426.532}{then}%
\TextAt{8}{135.141}{426.532}{one}%
\TextAt{8}{156.6}{426.532}{has}%
\TextAt{13}{231.031}{409.486}{A}%
\TextAt{9}{239.806}{409.486}{\Character{47}}%
\TextAt{14}{245.659}{409.486}{p}%
\TextAt{21}{251.642}{414.422}{n}%
\NativeMathAt{260.599}{409.486}{11.955}{\cong}%
\TextAt{13}{273.219}{409.486}{A}%
\TextAt{27}{281.994}{407.693}{p}%
\TextAt{9}{286.481}{409.486}{\Character{47}}%
\TextAt{14}{292.334}{409.486}{m}%
\TextAt{21}{301.501}{414.422}{n}%
\TextAt{27}{301.501}{406.53}{p}%
\TextAt{8}{331.047}{409.486}{for}%
\TextAt{13}{348.932}{409.486}{n}%
\NativeMathAt{359.241}{409.486}{11.955}{\geq}%
\TextAt{9}{371.86}{409.486}{1}%
\TextAt{13}{377.713}{409.486}{\Character{58}}%
\ErrataAt{79.22}{368.355}{453.64}{\textbf{Proof.} The prime ideals of $A_{\mathfrak p}$ correspond to the prime ideals of $A$ contained in $\mathfrak p$. The ideal}%
\TextAt{14}{79.226}{351.292}{m}%
\TextAt{27}{88.393}{349.498}{p}%
\TextAt{9}{96.201}{351.292}{\Character{61}}%
\TextAt{14}{108.626}{351.292}{p}%
\TextAt{13}{114.609}{351.292}{A}%
\TextAt{27}{123.385}{349.498}{p}%
\TextAt{8}{131.769}{351.292}{is}%
\TextAt{8}{143.533}{351.292}{the}%
\TextAt{8}{163.689}{351.292}{unique}%
\TextAt{8}{201.73}{351.292}{maximal}%
\TextAt{8}{249.539}{351.292}{ideal.}%
\TextAt{8}{282.053}{351.292}{Let}%
\TextAt{8}{303.023}{351.292}{us}%
\TextAt{8}{318.038}{351.292}{consider}%
\TextAt{8}{363.623}{351.292}{the}%
\TextAt{8}{383.78}{351.292}{homomorphism}%
\TextAt{13}{177.224}{322.597}{f}%
\TextAt{9}{187.591}{322.597}{\Character{58}}%
\TextAt{13}{194.164}{322.597}{A}%
\TextAt{9}{202.939}{322.597}{\Character{47}}%
\TextAt{14}{208.792}{322.597}{p}%
\TextAt{21}{214.775}{327.533}{n}%
\NativeArrowAt{223.732}{322.597}{11.955}{19.26}{\longrightarrow}%
\TextAt{13}{246.314}{322.597}{A}%
\TextAt{27}{255.09}{320.803}{p}%
\TextAt{9}{259.577}{322.597}{\Character{47}}%
\TextAt{14}{265.43}{322.597}{m}%
\TextAt{21}{274.596}{327.533}{n}%
\TextAt{27}{274.596}{319.641}{p}%
\TextAt{13}{280.233}{322.597}{\Character{59}}%
\TextAt{13}{309.368}{322.597}{a}%
\TextAt{8}{318.853}{322.597}{mod}%
\TextAt{14}{344.291}{322.597}{p}%
\TextAt{21}{350.274}{327.533}{n}%
\NativeArrowAt{359.231}{322.597}{11.955}{19.26}{\longmapsto}%
\TextAt{13}{381.813}{322.597}{a}%
\TextAt{8}{391.279}{322.597}{mod}%
\TextAt{14}{416.717}{322.597}{m}%
\TextAt{21}{425.884}{327.533}{n}%
\TextAt{27}{425.884}{319.641}{p}%
\TextAt{13}{431.52}{322.597}{\Character{58}}%
\TextAt{8}{79.226}{293.902}{For}%
\TextAt{13}{101.633}{293.902}{n}%
\TextAt{9}{114.391}{293.902}{\Character{61}}%
\TextAt{9}{129.251}{293.902}{1}%
\TextAt{8}{135.119}{293.902}{,}%
\TextAt{13}{144.065}{293.902}{f}%
\TextAt{8}{156.446}{293.902}{is}%
\TextAt{8}{169.647}{293.902}{injective}%
\TextAt{8}{217.578}{293.902}{because}%
\TextAt{14}{262.317}{293.902}{p}%
\TextAt{9}{274.07}{293.902}{\Character{61}}%
\TextAt{14}{288.945}{293.902}{m}%
\TextAt{27}{298.112}{292.108}{p}%
\TextAt{16}{306.235}{293.902}{\Character{92}}%
\TextAt{13}{317.842}{293.902}{A}%
\TextAt{8}{326.617}{293.902}{.}%
\TextAt{8}{339.377}{293.902}{Hence}%
\TextAt{13}{375.592}{293.902}{A}%
\TextAt{27}{384.368}{292.108}{p}%
\TextAt{9}{388.855}{293.902}{\Character{47}}%
\TextAt{14}{394.708}{293.902}{m}%
\TextAt{27}{403.874}{292.108}{p}%
\TextAt{13}{408.361}{293.902}{A}%
\TextAt{27}{417.137}{292.108}{p}%
\TextAt{8}{426.958}{293.902}{becomes}%
\TextAt{8}{474.949}{293.902}{the}%
\TextAt{8}{496.543}{293.902}{field}%
\TextAt{8}{523.337}{293.902}{of}%
\TextAt{8}{79.226}{276.856}{fractions}%
\TextAt{8}{126.919}{276.856}{of}%
\TextAt{13}{140.074}{276.856}{A}%
\TextAt{9}{148.849}{276.856}{\Character{47}}%
\TextAt{14}{154.702}{276.856}{p}%
\TextAt{8}{160.685}{276.856}{.}%
\TextAt{8}{169.075}{276.856}{Let}%
\TextAt{14}{189.869}{276.856}{p}%
\TextAt{8}{199.574}{276.856}{be}%
\TextAt{8}{215.323}{276.856}{maximal}%
\TextAt{8}{262.956}{276.856}{and}%
\TextAt{13}{285.543}{276.856}{n}%
\NativeMathAt{295.852}{276.856}{11.955}{\geq}%
\TextAt{9}{308.471}{276.856}{1}%
\TextAt{8}{314.324}{276.856}{.}%
\TextAt{8}{322.714}{276.856}{For}%
\TextAt{8}{343.508}{276.856}{every}%
\TextAt{13}{374.225}{276.856}{s}%
\NativeMathAt{383.06}{276.856}{11.955}{\in}%
\TextAt{13}{394.351}{276.856}{A}%
\TextAt{16}{405.424}{276.856}{\Character{110}}%
\TextAt{14}{413.699}{276.856}{p}%
\TextAt{8}{423.404}{276.856}{one}%
\TextAt{8}{444.689}{276.856}{has}%
\TextAt{14}{465.387}{276.856}{p}%
\TextAt{21}{471.37}{281.194}{n}%
\TextAt{9}{479.304}{276.856}{\Character{43}}%
\TextAt{13}{490.707}{276.856}{sA}%
\TextAt{9}{508.317}{276.856}{\Character{61}}%
\TextAt{13}{520.742}{276.856}{A}%
\TextAt{8}{529.518}{276.856}{,}%
\TextAt{8}{300.14}{43.033}{62}%
\Destination{78.23}{809.815}{nav.page.62}%
\Destination{79.36}{383.355}{nav.print.67}%
\Destination{79.36}{532.055}{nav.number.1.11.2}%
\Hotlink{380.541}{678.255}{30.936}{11.047}{nav.number.1.11.1}%
\Hotlink{396.565}{98.955}{26.464}{9.572}{nav.number.1.11.1}%
\Destination{79.226}{660.455}{errata.8}%
\Destination{79.226}{382.355}{errata.9}%
}

% Page 75 (body page 63)
\Page{%
\TextAt{2}{79.226}{796.943}{I.}%
\TextAt{2}{89.83}{796.943}{Algebraic}%
\TextAt{2}{138.648}{796.943}{Integers}%
\Vector{q 0 G 0 g 0.299 w 7.226 18.729 m 460.769 18.729 l S Q}%
\TextAt{8}{79.226}{761.124}{i.e.,}%
\Vector{q 0 G 0 g 0.478 w 29.091 -6.055 m 34.605 -6.055 l S Q}%
\TextAt{13}{101.091}{761.124}{s}%
\TextAt{9}{109.925}{761.124}{\Character{61}}%
\TextAt{13}{122.351}{761.124}{s}%
\TextAt{8}{131.186}{761.124}{mod}%
\TextAt{14}{156.624}{761.124}{p}%
\TextAt{21}{162.607}{765.462}{n}%
\TextAt{8}{171.84}{761.124}{is}%
\TextAt{8}{183.303}{761.124}{a}%
\TextAt{8}{192.758}{761.124}{unit}%
\TextAt{8}{217.168}{761.124}{in}%
\TextAt{13}{230.52}{761.124}{A}%
\TextAt{9}{239.296}{761.124}{\Character{47}}%
\TextAt{14}{245.149}{761.124}{p}%
\TextAt{21}{251.132}{765.462}{n}%
\TextAt{8}{256.768}{761.124}{.}%
\TextAt{8}{265.116}{761.124}{For}%
\TextAt{13}{285.785}{761.124}{n}%
\TextAt{9}{296.093}{761.124}{\Character{61}}%
\TextAt{9}{308.502}{761.124}{1}%
\TextAt{8}{317.968}{761.124}{this}%
\TextAt{8}{340.49}{761.124}{is}%
\TextAt{8}{351.953}{761.124}{clear}%
\TextAt{8}{379.615}{761.124}{from}%
\TextAt{8}{406.955}{761.124}{the}%
\TextAt{8}{426.811}{761.124}{maximality}%
\TextAt{8}{487.983}{761.124}{of}%
\TextAt{14}{501.013}{761.124}{p}%
\TextAt{8}{506.996}{761.124}{,}%
\TextAt{8}{513.904}{761.124}{and}%
\TextAt{8}{79.226}{744.078}{for}%
\TextAt{13}{97.111}{744.078}{n}%
\NativeMathAt{107.42}{744.078}{11.955}{\geq}%
\TextAt{9}{120.039}{744.078}{1}%
\TextAt{8}{129.789}{744.078}{it}%
\TextAt{8}{141.493}{744.078}{follows}%
\TextAt{8}{179.929}{744.078}{by}%
\TextAt{8}{196.188}{744.078}{induction:}%
\TextAt{13}{121.036}{719.166}{A}%
\TextAt{9}{133.132}{719.166}{\Character{61}}%
\TextAt{14}{145.558}{719.166}{p}%
\TextAt{21}{151.541}{724.102}{n}%
\TextAt{3}{156.679}{724.102}{\Character{0}}%
\TextAt{17}{163.266}{724.102}{1}%
\TextAt{9}{170.655}{719.166}{\Character{43}}%
\TextAt{13}{182.416}{719.166}{sA}%
\NativeArrowAt{200.026}{719.166}{11.955}{19.067}{\Longrightarrow}%
\TextAt{14}{222.414}{719.166}{p}%
\TextAt{9}{231.718}{719.166}{\Character{61}}%
\TextAt{14}{244.144}{719.166}{p}%
\TextAt{13}{250.127}{719.166}{A}%
\TextAt{9}{262.223}{719.166}{\Character{61}}%
\TextAt{14}{274.648}{719.166}{p}%
\TextAt{9}{280.632}{719.166}{\Character{40}}%
\TextAt{14}{285.184}{719.166}{p}%
\TextAt{21}{291.167}{724.102}{n}%
\TextAt{3}{296.305}{724.102}{\Character{0}}%
\TextAt{17}{302.892}{724.102}{1}%
\TextAt{9}{310.281}{719.166}{\Character{43}}%
\TextAt{13}{322.042}{719.166}{sA}%
\TextAt{9}{336.331}{719.166}{\Character{41}}%
\NativeMathAt{344.205}{719.166}{11.955}{\subseteq}%
\TextAt{14}{356.824}{719.166}{p}%
\TextAt{21}{362.807}{724.102}{n}%
\TextAt{9}{371.1}{719.166}{\Character{43}}%
\TextAt{13}{382.861}{719.166}{sA}%
\NativeArrowAt{400.471}{719.166}{11.955}{19.068}{\Longrightarrow}%
\TextAt{14}{422.86}{719.166}{p}%
\TextAt{21}{428.843}{724.102}{n}%
\TextAt{9}{437.136}{719.166}{\Character{43}}%
\TextAt{13}{448.897}{719.166}{sA}%
\TextAt{9}{466.507}{719.166}{\Character{61}}%
\TextAt{13}{478.933}{719.166}{A\Character{58}}%
\TextAt{8}{79.226}{694.254}{Injectivity}%
\TextAt{8}{135.971}{694.254}{of}%
\TextAt{13}{149.798}{694.254}{f}%
\TextAt{8}{156.844}{694.254}{:}%
\TextAt{8}{166.285}{694.254}{let}%
\TextAt{13}{183.686}{694.254}{a}%
\NativeMathAt{193.997}{694.254}{11.955}{\in}%
\TextAt{13}{206.134}{694.254}{A}%
\TextAt{8}{219.303}{694.254}{be}%
\TextAt{8}{235.723}{694.254}{such}%
\TextAt{8}{262.616}{694.254}{that}%
\TextAt{13}{288.481}{694.254}{a}%
\NativeMathAt{298.793}{694.254}{11.955}{\in}%
\TextAt{14}{310.929}{694.254}{m}%
\TextAt{21}{320.096}{698.593}{n}%
\TextAt{27}{320.096}{691.299}{p}%
\TextAt{8}{325.732}{694.254}{,}%
\TextAt{8}{333.502}{694.254}{i.e.,}%
\TextAt{13}{356.227}{694.254}{a}%
\TextAt{9}{366.539}{694.254}{\Character{61}}%
\TextAt{13}{379.81}{694.254}{b}%
\TextAt{9}{384.787}{694.254}{\Character{47}}%
\TextAt{13}{390.64}{694.254}{s}%
\TextAt{8}{400.547}{694.254}{with}%
\TextAt{13}{427.704}{694.254}{b}%
\NativeMathAt{436.847}{694.254}{11.955}{\in}%
\TextAt{14}{448.984}{694.254}{p}%
\TextAt{21}{454.967}{698.593}{n}%
\TextAt{8}{460.603}{694.254}{,}%
\TextAt{13}{468.373}{694.254}{s}%
\NativeMathAt{478.053}{694.254}{11.955}{\notin}%
\TextAt{14}{490.19}{694.254}{p}%
\TextAt{8}{496.173}{694.254}{.}%
\TextAt{8}{506.11}{694.254}{Then}%
\TextAt{13}{79.226}{677.208}{as}%
\TextAt{9}{94.206}{677.208}{\Character{61}}%
\TextAt{13}{106.631}{677.208}{b}%
\NativeMathAt{114.929}{677.208}{11.955}{\in}%
\TextAt{14}{126.22}{677.208}{p}%
\TextAt{21}{132.203}{681.547}{n}%
\TextAt{8}{137.84}{677.208}{,}%
\TextAt{8}{144.989}{677.208}{so}%
\TextAt{8}{159.359}{677.208}{that}%
\Vector{q 0 G 0 g 0.478 w 112.728 -89.971 m 118.873 -89.971 l S Q}%
\TextAt{13}{184.728}{677.208}{a}%
\Vector{q 0 G 0 g 0.478 w 118.873 -89.971 m 124.387 -89.971 l S Q}%
\TextAt{13}{190.873}{677.208}{s}%
\TextAt{9}{199.708}{677.208}{\Character{61}}%
\TextAt{9}{212.118}{677.208}{0}%
\TextAt{8}{221.884}{677.208}{in}%
\TextAt{13}{235.536}{677.208}{A}%
\TextAt{9}{244.312}{677.208}{\Character{47}}%
\TextAt{14}{250.165}{677.208}{p}%
\TextAt{21}{256.148}{681.547}{n}%
\TextAt{8}{261.784}{677.208}{,}%
\TextAt{8}{268.933}{677.208}{and}%
\TextAt{8}{291.696}{677.208}{hence}%
\Vector{q 0 G 0 g 0.478 w 252.202 -89.971 m 258.347 -89.971 l S Q}%
\TextAt{13}{324.202}{677.208}{a}%
\TextAt{9}{333.668}{677.208}{\Character{61}}%
\TextAt{9}{346.077}{677.208}{0}%
\TextAt{8}{355.844}{677.208}{in}%
\TextAt{13}{369.497}{677.208}{A}%
\TextAt{9}{378.272}{677.208}{\Character{47}}%
\TextAt{14}{384.125}{677.208}{p}%
\TextAt{21}{390.108}{681.547}{n}%
\TextAt{8}{395.744}{677.208}{.}%
\TextAt{8}{97.159}{657.173}{Surjectivity}%
\TextAt{8}{161.174}{657.173}{of}%
\TextAt{13}{175.432}{657.173}{f}%
\TextAt{8}{182.479}{657.173}{:}%
\TextAt{8}{192.783}{657.173}{let}%
\TextAt{13}{210.616}{657.173}{a}%
\TextAt{9}{216.76}{657.173}{\Character{47}}%
\TextAt{13}{222.613}{657.173}{s}%
\NativeMathAt{233.03}{657.173}{11.955}{\in}%
\TextAt{13}{245.902}{657.173}{A}%
\TextAt{27}{254.677}{655.38}{p}%
\TextAt{8}{259.164}{657.173}{,}%
\TextAt{13}{267.473}{657.173}{a}%
\NativeMathAt{278.521}{657.173}{11.955}{\in}%
\TextAt{13}{291.393}{657.173}{A}%
\TextAt{8}{300.168}{657.173}{,}%
\TextAt{13}{308.477}{657.173}{s}%
\NativeMathAt{318.894}{657.173}{11.955}{\notin}%
\TextAt{14}{331.766}{657.173}{p}%
\TextAt{8}{337.749}{657.173}{.}%
\TextAt{8}{348.981}{657.173}{Then}%
\TextAt{8}{380.467}{657.173}{by}%
\TextAt{8}{397.654}{657.173}{the}%
\TextAt{8}{418.738}{657.173}{above,}%
\TextAt{8}{456.325}{657.173}{there}%
\TextAt{8}{487.165}{657.173}{exists}%
\TextAt{8}{520.408}{657.173}{an}%
\TextAt{13}{79.226}{640.127}{a}%
\TextAt{3}{85.371}{644.466}{\Character{48}}%
\NativeMathAt{91.487}{640.127}{11.955}{\in}%
\TextAt{13}{102.778}{640.127}{A}%
\TextAt{8}{115.031}{640.127}{such}%
\TextAt{8}{141.009}{640.127}{that}%
\TextAt{13}{165.958}{640.127}{a}%
\NativeMathAt{175.424}{640.127}{11.955}{\equiv}%
\TextAt{13}{188.043}{640.127}{a}%
\TextAt{3}{194.188}{644.466}{\Character{48}}%
\TextAt{13}{196.983}{640.127}{s}%
\TextAt{9}{207.81}{640.127}{\Character{40}}%
\TextAt{8}{212.363}{640.127}{mod}%
\TextAt{14}{238.465}{640.127}{p}%
\TextAt{21}{244.448}{644.466}{n}%
\TextAt{9}{250.084}{640.127}{\Character{41}}%
\TextAt{8}{254.636}{640.127}{.}%
\TextAt{8}{262.945}{640.127}{Therefore}%
\TextAt{13}{315.523}{640.127}{a}%
\TextAt{9}{321.667}{640.127}{\Character{47}}%
\TextAt{13}{327.52}{640.127}{s}%
\NativeMathAt{336.355}{640.127}{11.955}{\equiv}%
\TextAt{13}{348.975}{640.127}{a}%
\TextAt{3}{355.12}{644.466}{\Character{48}}%
\TextAt{9}{363.228}{640.127}{\Character{40}}%
\TextAt{8}{367.78}{640.127}{mod}%
\TextAt{14}{393.882}{640.127}{p}%
\TextAt{21}{399.865}{644.466}{n}%
\TextAt{13}{405.502}{640.127}{A}%
\TextAt{27}{414.277}{638.334}{p}%
\TextAt{9}{418.764}{640.127}{\Character{41}}%
\TextAt{8}{423.316}{640.127}{,}%
\TextAt{8}{430.13}{640.127}{i.e.,}%
\TextAt{13}{451.899}{640.127}{a}%
\TextAt{9}{458.044}{640.127}{\Character{47}}%
\TextAt{13}{463.897}{640.127}{s}%
\TextAt{8}{472.732}{640.127}{mod}%
\TextAt{14}{498.17}{640.127}{m}%
\TextAt{21}{507.337}{644.466}{n}%
\TextAt{27}{507.337}{637.172}{p}%
\TextAt{8}{516.451}{640.127}{lies}%
\TextAt{8}{79.226}{623.081}{in}%
\TextAt{8}{92.879}{623.081}{the}%
\TextAt{8}{113.035}{623.081}{image}%
\TextAt{8}{146.857}{623.081}{of}%
\TextAt{13}{160.187}{623.081}{f}%
\TextAt{8}{167.233}{623.081}{.}%
\Vector{q 0 G 0 g 0.399 w 452.334 -150.919 m 452.334 -142.849 l S Q}%
\Vector{q 0 G 0 g 0.399 w 452.534 -143.048 m 459.707 -143.048 l S Q}%
\Vector{q 0 G 0 g 0.399 w 452.534 -150.719 m 459.707 -150.719 l S Q}%
\Vector{q 0 G 0 g 0.399 w 459.906 -150.919 m 459.906 -142.849 l S Q}%
\TextAt{8}{97.159}{595.249}{In}%
\TextAt{8}{112.545}{595.249}{a}%
\TextAt{8}{123.065}{595.249}{local}%
\TextAt{8}{151.47}{595.249}{ring}%
\TextAt{8}{176.301}{595.249}{with}%
\TextAt{8}{203.726}{595.249}{maximal}%
\TextAt{8}{252.299}{595.249}{ideal}%
\TextAt{14}{281.027}{595.249}{m}%
\TextAt{8}{290.194}{595.249}{,}%
\TextAt{8}{298.299}{595.249}{every}%
\TextAt{8}{329.956}{595.249}{element}%
\TextAt{13}{373.963}{595.249}{a}%
\NativeMathAt{384.732}{595.249}{11.955}{\notin}%
\TextAt{14}{397.326}{595.249}{m}%
\TextAt{8}{411.155}{595.249}{is}%
\TextAt{8}{423.684}{595.249}{a}%
\TextAt{8}{434.204}{595.249}{unit.}%
\TextAt{8}{465.761}{595.249}{Indeed,}%
\TextAt{8}{507.998}{595.249}{since}%
\TextAt{8}{79.226}{578.203}{the}%
\TextAt{8}{99.044}{578.203}{principal}%
\TextAt{8}{147.483}{578.203}{ideal}%
\TextAt{9}{175.108}{578.203}{\Character{40}}%
\TextAt{13}{179.66}{578.203}{a}%
\TextAt{9}{185.805}{578.203}{\Character{41}}%
\TextAt{8}{193.916}{578.203}{is}%
\TextAt{8}{205.342}{578.203}{not}%
\TextAt{8}{225.817}{578.203}{contained}%
\TextAt{8}{278.488}{578.203}{in}%
\TextAt{8}{291.802}{578.203}{any}%
\TextAt{8}{313.581}{578.203}{other}%
\TextAt{8}{343.812}{578.203}{maximal}%
\TextAt{8}{391.282}{578.203}{ideal,}%
\TextAt{8}{422.226}{578.203}{it}%
\TextAt{8}{433.592}{578.203}{has}%
\TextAt{8}{454.127}{578.203}{to}%
\TextAt{8}{468.099}{578.203}{be}%
\TextAt{8}{483.685}{578.203}{the}%
\TextAt{8}{503.503}{578.203}{whole}%
\TextAt{8}{79.226}{561.157}{ring.}%
\TextAt{8}{107.843}{561.157}{So}%
\TextAt{8}{124.102}{561.157}{we}%
\TextAt{8}{141.329}{561.157}{have}%
\TextAt{13}{275.129}{544.111}{A}%
\TextAt{3}{283.904}{549.047}{\Character{3}}%
\TextAt{9}{291.957}{544.111}{\Character{61}}%
\TextAt{13}{304.382}{544.111}{A}%
\TextAt{16}{315.814}{544.111}{\Character{110}}%
\TextAt{14}{324.449}{544.111}{m}%
\TextAt{13}{333.616}{544.111}{\Character{58}}%
\TextAt{8}{79.226}{522.342}{The}%
\TextAt{8}{103.28}{522.342}{simplest}%
\TextAt{8}{148.925}{522.342}{local}%
\TextAt{8}{176.565}{522.342}{rings,}%
\TextAt{8}{208.497}{522.342}{except}%
\TextAt{8}{245.236}{522.342}{for}%
\TextAt{8}{263.121}{522.342}{fields,}%
\TextAt{8}{296.344}{522.342}{are}%
\TextAt{12}{315.855}{522.342}{discrete}%
\TextAt{12}{365.839}{522.342}{valuation}%
\TextAt{12}{424.145}{522.342}{rings}%
\TextAt{8}{452.897}{522.342}{.}%
\TextAt{12}{79.226}{497.498}{Definition}%
\TextAt{12}{141.83}{497.498}{(11.3).}%
\TextAt{8}{183.716}{497.498}{A}%
\TextAt{8}{196.228}{497.498}{discrete}%
\TextAt{8}{239.047}{497.498}{valuation}%
\TextAt{8}{289.959}{497.498}{ring}%
\TextAt{8}{313.864}{497.498}{is}%
\TextAt{8}{325.468}{497.498}{a}%
\TextAt{8}{335.063}{497.498}{principal}%
\TextAt{8}{383.679}{497.498}{ideal}%
\TextAt{8}{411.482}{497.498}{domain}%
\TextAt{14}{452.95}{497.498}{o}%
\TextAt{8}{462.531}{497.498}{with}%
\TextAt{8}{489.03}{497.498}{a}%
\TextAt{8}{498.626}{497.498}{unique}%
\TextAt{8}{79.226}{480.452}{maximal}%
\TextAt{8}{127.035}{480.452}{ideal}%
\TextAt{14}{154.998}{480.452}{p}%
\NativeMathAt{164.302}{480.452}{11.955}{\neq}%
\TextAt{9}{176.712}{480.452}{0}%
\TextAt{8}{182.58}{480.452}{.}%
\TextAt{8}{97.159}{455.608}{The}%
\TextAt{8}{120.641}{455.608}{maximal}%
\TextAt{8}{167.879}{455.608}{ideal}%
\TextAt{8}{195.27}{455.608}{is}%
\TextAt{8}{206.463}{455.608}{of}%
\TextAt{8}{219.221}{455.608}{the}%
\TextAt{8}{238.807}{455.608}{form}%
\TextAt{14}{265.875}{455.608}{p}%
\TextAt{9}{275.179}{455.608}{\Character{61}}%
\TextAt{9}{287.589}{455.608}{\Character{40}}%
\TextAt{13}{292.157}{455.608}{\Character{25}}%
\TextAt{9}{299.226}{455.608}{\Character{41}}%
\TextAt{9}{307.081}{455.608}{\Character{61}}%
\TextAt{13}{319.525}{455.608}{\Character{25}}%
\TextAt{14}{326.594}{455.608}{o}%
\TextAt{8}{332.439}{455.608}{,}%
\TextAt{8}{339.131}{455.608}{for}%
\TextAt{8}{356.444}{455.608}{some}%
\TextAt{8}{385.199}{455.608}{prime}%
\TextAt{8}{417.791}{455.608}{element}%
\TextAt{13}{460.462}{455.608}{\Character{25}}%
\TextAt{8}{467.531}{455.608}{.}%
\TextAt{8}{475.789}{455.608}{Since}%
\TextAt{8}{505.775}{455.608}{every}%
\TextAt{8}{79.226}{438.562}{element}%
\TextAt{8}{122.899}{438.562}{not}%
\TextAt{8}{144.144}{438.562}{contained}%
\TextAt{8}{197.585}{438.562}{in}%
\TextAt{14}{211.669}{438.562}{p}%
\TextAt{8}{221.981}{438.562}{is}%
\TextAt{8}{234.176}{438.562}{a}%
\TextAt{8}{244.362}{438.562}{unit,}%
\TextAt{8}{272.864}{438.562}{it}%
\TextAt{8}{285}{438.562}{follows}%
\TextAt{8}{323.867}{438.562}{that,}%
\TextAt{8}{353.027}{438.562}{up}%
\TextAt{8}{370.362}{438.562}{to}%
\TextAt{8}{385.104}{438.562}{associated}%
\TextAt{8}{441.27}{438.562}{elements,}%
\TextAt{13}{492.917}{438.562}{\Character{25}}%
\TextAt{8}{504.315}{438.562}{is}%
\TextAt{8}{516.51}{438.562}{the}%
\TextAt{8}{79.226}{421.516}{only}%
\TextAt{8}{104.602}{421.516}{prime}%
\TextAt{8}{137.451}{421.516}{element}%
\TextAt{8}{180.377}{421.516}{of}%
\TextAt{14}{193.391}{421.516}{o}%
\TextAt{8}{199.236}{421.516}{.}%
\TextAt{8}{207.579}{421.516}{Every}%
\TextAt{8}{240.917}{421.516}{nonzero}%
\TextAt{8}{284.178}{421.516}{element}%
\TextAt{8}{327.105}{421.516}{of}%
\TextAt{14}{340.119}{421.516}{o}%
\TextAt{8}{349.545}{421.516}{may}%
\TextAt{8}{374.599}{421.516}{therefore}%
\TextAt{8}{423.383}{421.516}{be}%
\TextAt{8}{438.992}{421.516}{written}%
\TextAt{8}{479.647}{421.516}{as}%
\TextAt{13}{493.701}{421.516}{\Character{34}\Character{25}}%
\TextAt{21}{506.249}{425.854}{n}%
\TextAt{8}{511.885}{421.516}{,}%
\TextAt{8}{518.782}{421.516}{for}%
\TextAt{8}{79.226}{404.47}{some}%
\TextAt{13}{108.833}{404.47}{\Character{34}}%
\NativeMathAt{118.111}{404.47}{11.955}{\in}%
\TextAt{14}{129.881}{404.47}{o}%
\TextAt{3}{135.725}{408.808}{\Character{3}}%
\TextAt{8}{140.458}{404.47}{,}%
\TextAt{8}{147.958}{404.47}{and}%
\TextAt{13}{171.002}{404.47}{n}%
\NativeMathAt{181.789}{404.47}{11.955}{\geq}%
\TextAt{9}{194.887}{404.47}{0}%
\TextAt{8}{200.74}{404.47}{.}%
\TextAt{8}{210.032}{404.47}{More}%
\TextAt{8}{240.548}{404.47}{generally,}%
\TextAt{8}{292.928}{404.47}{every}%
\TextAt{8}{324.101}{404.47}{element}%
\TextAt{13}{367.624}{404.47}{a}%
\NativeMathAt{377.569}{404.47}{11.955}{\neq}%
\TextAt{9}{390.457}{404.47}{0}%
\TextAt{8}{400.505}{404.47}{of}%
\TextAt{8}{414.116}{404.47}{the}%
\TextAt{8}{434.553}{404.47}{field}%
\TextAt{8}{460.191}{404.47}{of}%
\TextAt{8}{473.802}{404.47}{fractions}%
\TextAt{13}{521.952}{404.47}{K}%
\TextAt{8}{79.226}{387.424}{may}%
\TextAt{8}{104.595}{387.424}{be}%
\TextAt{8}{120.519}{387.424}{uniquely}%
\TextAt{8}{167.993}{387.424}{written}%
\TextAt{8}{208.964}{387.424}{as}%
\TextAt{13}{231.04}{362.512}{a}%
\TextAt{9}{240.505}{362.512}{\Character{61}}%
\TextAt{13}{252.931}{362.512}{\Character{34}\Character{25}}%
\TextAt{21}{265.478}{367.448}{n}%
\TextAt{13}{271.115}{362.512}{\Character{59}}%
\TextAt{13}{300.25}{362.512}{\Character{34}}%
\NativeMathAt{309.068}{362.512}{11.955}{\in}%
\TextAt{14}{320.359}{362.512}{o}%
\TextAt{3}{326.203}{367.448}{\Character{3}}%
\TextAt{13}{330.935}{362.512}{\Character{59}}%
\TextAt{13}{348.115}{362.512}{n}%
\NativeMathAt{358.443}{362.512}{11.955}{\in}%
\TextAt{18}{369.734}{362.512}{Z}%
\TextAt{13}{377.704}{362.512}{\Character{58}}%
\TextAt{8}{79.226}{337.6}{The}%
\TextAt{8}{104.717}{337.6}{exponent}%
\TextAt{13}{156.556}{337.6}{n}%
\TextAt{8}{168.878}{337.6}{is}%
\TextAt{8}{182.078}{337.6}{called}%
\TextAt{8}{216.679}{337.6}{the}%
\TextAt{12}{238.272}{337.6}{valuation}%
\TextAt{8}{297.428}{337.6}{of}%
\TextAt{13}{312.194}{337.6}{a}%
\TextAt{8}{318.339}{337.6}{.}%
\TextAt{8}{331.097}{337.6}{It}%
\TextAt{8}{345.207}{337.6}{is}%
\TextAt{8}{358.407}{337.6}{denoted}%
\TextAt{13}{404.066}{337.6}{v}%
\TextAt{9}{410.154}{337.6}{\Character{40}}%
\TextAt{13}{414.706}{337.6}{a}%
\TextAt{9}{420.851}{337.6}{\Character{41}}%
\TextAt{8}{425.404}{337.6}{,}%
\TextAt{8}{434.349}{337.6}{and}%
\TextAt{8}{458.548}{337.6}{it}%
\TextAt{8}{471.689}{337.6}{is}%
\TextAt{8}{484.889}{337.6}{obviously}%
\TextAt{8}{79.226}{320.554}{characterized}%
\TextAt{8}{150.443}{320.554}{by}%
\TextAt{8}{166.702}{320.554}{the}%
\TextAt{8}{186.859}{320.554}{equation}%
\TextAt{9}{277.886}{303.508}{\Character{40}}%
\TextAt{13}{282.439}{303.508}{a}%
\TextAt{9}{288.584}{303.508}{\Character{41}}%
\TextAt{9}{296.439}{303.508}{\Character{61}}%
\TextAt{14}{308.882}{303.508}{p}%
\TextAt{21}{314.865}{308.444}{v}%
\TextAt{17}{319.275}{308.444}{\Character{40}}%
\TextAt{21}{322.568}{308.444}{a}%
\TextAt{17}{327.066}{308.444}{\Character{41}}%
\TextAt{13}{330.858}{303.508}{\Character{58}}%
\TextAt{8}{79.226}{281.739}{The}%
\TextAt{8}{103.28}{281.739}{valuation}%
\TextAt{8}{154.352}{281.739}{is}%
\TextAt{8}{166.116}{281.739}{a}%
\TextAt{8}{175.872}{281.739}{function}%
\TextAt{13}{271.67}{264.693}{v}%
\TextAt{9}{281.079}{264.693}{\Character{58}}%
\TextAt{13}{287.651}{264.693}{K}%
\TextAt{3}{298.469}{269.629}{\Character{3}}%
\NativeArrowAt{306.522}{264.693}{11.955}{19.26}{\longrightarrow}%
\TextAt{18}{329.104}{264.693}{Z}%
\TextAt{13}{337.074}{264.693}{\Character{58}}%
\TextAt{8}{79.226}{242.924}{Extending}%
\TextAt{8}{134.934}{242.924}{it}%
\TextAt{8}{145.929}{242.924}{to}%
\TextAt{13}{159.53}{242.924}{K}%
\TextAt{8}{173.536}{242.924}{by}%
\TextAt{8}{189.086}{242.924}{the}%
\TextAt{8}{208.534}{242.924}{convention}%
\TextAt{13}{266.369}{242.924}{v}%
\TextAt{9}{272.458}{242.924}{\Character{40}0\Character{41}}%
\TextAt{9}{290.702}{242.924}{\Character{61}}%
\NativeMathAt{303.161}{242.924}{11.955}{\infty}%
\TextAt{8}{315.117}{242.924}{,}%
\TextAt{8}{321.699}{242.924}{a}%
\TextAt{8}{330.745}{242.924}{simple}%
\TextAt{8}{366.511}{242.924}{calculation}%
\TextAt{8}{424.992}{242.924}{shows}%
\TextAt{8}{457.901}{242.924}{that}%
\TextAt{8}{482.561}{242.924}{it}%
\TextAt{8}{493.557}{242.924}{satisfies}%
\TextAt{8}{79.226}{225.878}{the}%
\TextAt{8}{99.383}{225.878}{conditions}%
\TextAt{13}{173.356}{200.966}{v}%
\TextAt{9}{179.444}{200.966}{\Character{40}}%
\TextAt{13}{183.996}{200.966}{ab}%
\TextAt{9}{195.118}{200.966}{\Character{41}}%
\TextAt{9}{202.973}{200.966}{\Character{61}}%
\TextAt{13}{215.417}{200.966}{v}%
\TextAt{9}{221.505}{200.966}{\Character{40}}%
\TextAt{13}{226.058}{200.966}{a}%
\TextAt{9}{232.202}{200.966}{\Character{41}}%
\TextAt{9}{239.399}{200.966}{\Character{43}}%
\TextAt{13}{251.173}{200.966}{v}%
\TextAt{9}{257.261}{200.966}{\Character{40}}%
\TextAt{13}{261.813}{200.966}{b}%
\TextAt{9}{266.79}{200.966}{\Character{41}}%
\TextAt{13}{271.343}{200.966}{\Character{59}}%
\TextAt{13}{300.478}{200.966}{v}%
\TextAt{9}{306.586}{200.966}{\Character{40}}%
\TextAt{13}{311.138}{200.966}{a}%
\TextAt{9}{319.939}{200.966}{\Character{43}}%
\TextAt{13}{331.701}{200.966}{b}%
\TextAt{9}{336.678}{200.966}{\Character{41}}%
\NativeMathAt{344.551}{200.966}{11.955}{\geq}%
\TextAt{8}{357.17}{200.966}{min}%
\TextAt{16}{376.681}{200.966}{\Character{102}}%
\TextAt{13}{382.659}{200.966}{v}%
\TextAt{9}{388.747}{200.966}{\Character{40}}%
\TextAt{13}{393.299}{200.966}{a}%
\TextAt{9}{399.444}{200.966}{\Character{41}}%
\TextAt{13}{403.997}{200.966}{\Character{59}}%
\TextAt{13}{409.221}{200.966}{v}%
\TextAt{9}{415.329}{200.966}{\Character{40}}%
\TextAt{13}{419.881}{200.966}{b}%
\TextAt{9}{424.858}{200.966}{\Character{41}}%
\TextAt{16}{429.411}{200.966}{\Character{103}}%
\TextAt{13}{435.388}{200.966}{\Character{58}}%
\TextAt{8}{79.226}{176.055}{This}%
\TextAt{8}{106.402}{176.055}{innocuous}%
\TextAt{8}{161.875}{176.055}{looking}%
\TextAt{8}{203.314}{176.055}{function}%
\TextAt{8}{249.618}{176.055}{gives}%
\TextAt{8}{278.754}{176.055}{rise}%
\TextAt{8}{300.729}{176.055}{to}%
\TextAt{8}{315.495}{176.055}{a}%
\TextAt{8}{325.707}{176.055}{theory}%
\TextAt{8}{362.913}{176.055}{which}%
\TextAt{8}{396.855}{176.055}{will}%
\TextAt{8}{419.416}{176.055}{occupy}%
\TextAt{8}{459.216}{176.055}{all}%
\TextAt{8}{475.931}{176.055}{of}%
\TextAt{8}{489.717}{176.055}{the}%
\TextAt{8}{510.33}{176.055}{next}%
\TextAt{8}{79.226}{159.009}{chapter.}%
\TextAt{8}{97.159}{138.974}{The}%
\TextAt{8}{122.619}{138.974}{discrete}%
\TextAt{8}{167.005}{138.974}{valuation}%
\TextAt{8}{219.484}{138.974}{rings}%
\TextAt{8}{249.571}{138.974}{arise}%
\TextAt{8}{278.355}{138.974}{as}%
\TextAt{8}{294.132}{138.974}{localizations}%
\TextAt{8}{362.273}{138.974}{of}%
\TextAt{8}{377.01}{138.974}{Dedekind}%
\TextAt{8}{430.601}{138.974}{domains.}%
\TextAt{8}{485.614}{138.974}{This}%
\TextAt{8}{513.741}{138.974}{is}%
\TextAt{8}{526.911}{138.974}{a}%
\TextAt{8}{79.226}{121.928}{consequence}%
\TextAt{8}{145.29}{121.928}{of}%
\TextAt{8}{158.62}{121.928}{the}%
\TextAt{12}{79.226}{97.084}{Proposition}%
\TextAt{12}{153.403}{97.084}{(11.4).}%
\TextAt{8}{196.724}{97.084}{If}%
\TextAt{14}{209.691}{97.084}{o}%
\TextAt{8}{220.707}{97.084}{is}%
\TextAt{8}{233.745}{97.084}{a}%
\TextAt{8}{244.775}{97.084}{Dedekind}%
\TextAt{8}{298.234}{97.084}{domain,}%
\TextAt{8}{344.707}{97.084}{and}%
\TextAt{13}{368.744}{97.084}{S}%
\NativeMathAt{382.132}{97.084}{11.955}{\subseteq}%
\TextAt{14}{396.923}{97.084}{o}%
\TextAt{16}{406.292}{97.084}{\Character{110}}%
\TextAt{16}{415.784}{97.084}{\Character{102}}%
\TextAt{9}{421.773}{97.084}{0}%
\TextAt{16}{427.626}{97.084}{\Character{103}}%
\TextAt{8}{438.776}{97.084}{is}%
\TextAt{8}{451.814}{97.084}{a}%
\TextAt{8}{462.844}{97.084}{multiplicative}%
\TextAt{8}{79.226}{80.038}{subset,}%
\TextAt{8}{118.367}{80.038}{then}%
\TextAt{14}{145.027}{80.038}{o}%
\TextAt{13}{150.872}{80.038}{S}%
\TextAt{3}{158.767}{84.376}{\Character{0}}%
\TextAt{17}{165.353}{84.376}{1}%
\TextAt{8}{173.983}{80.038}{is}%
\TextAt{8}{185.747}{80.038}{also}%
\TextAt{8}{209.227}{80.038}{a}%
\TextAt{8}{218.982}{80.038}{Dedekind}%
\TextAt{8}{271.167}{80.038}{domain.}%
\TextAt{8}{300.14}{43.033}{63}%
\Destination{78.23}{809.815}{nav.page.63}%
\Destination{79.36}{257.955}{nav.print.68}%
\Destination{79.36}{511.455}{nav.number.1.11.3}%
\Destination{79.36}{111.055}{nav.number.1.11.4}%
}

% Page 76 (body page 64)
\Page{%
\TextAt{2}{79.226}{796.943}{I.}%
\TextAt{2}{89.83}{796.943}{Algebraic}%
\TextAt{2}{138.648}{796.943}{Integers}%
\Vector{q 0 G 0 g 0.299 w 7.226 18.729 m 460.769 18.729 l S Q}%
\TextAt{19}{79.226}{761.124}{Proof.}%
\TextAt{8}{116}{761.124}{Let}%
\TextAt{14}{137.651}{761.124}{A}%
\TextAt{8}{150.809}{761.124}{be}%
\TextAt{8}{167.414}{761.124}{an}%
\TextAt{8}{184.355}{761.124}{ideal}%
\TextAt{8}{212.999}{761.124}{of}%
\TextAt{14}{227.011}{761.124}{o}%
\TextAt{13}{232.855}{761.124}{S}%
\TextAt{3}{240.75}{765.462}{\Character{0}}%
\TextAt{17}{247.337}{765.462}{1}%
\TextAt{8}{256.648}{761.124}{and}%
\TextAt{14}{280.092}{761.124}{a}%
\TextAt{9}{290.557}{761.124}{\Character{61}}%
\TextAt{14}{304.144}{761.124}{A}%
\TextAt{16}{315.844}{761.124}{\Character{92}}%
\TextAt{14}{326.935}{761.124}{o}%
\TextAt{8}{332.779}{761.124}{.}%
\TextAt{8}{343.272}{761.124}{Then}%
\TextAt{14}{374.511}{761.124}{A}%
\TextAt{9}{387.572}{761.124}{\Character{61}}%
\TextAt{14}{401.158}{761.124}{a}%
\TextAt{13}{407.141}{761.124}{S}%
\TextAt{3}{415.037}{765.462}{\Character{0}}%
\TextAt{17}{421.623}{765.462}{1}%
\TextAt{8}{426.355}{761.124}{,}%
\TextAt{8}{434.356}{761.124}{because}%
\TextAt{8}{478.339}{761.124}{if}%
\TextAt{13}{489.744}{761.124}{a}%
\TextAt{9}{495.889}{761.124}{\Character{47}}%
\TextAt{13}{501.742}{761.124}{s}%
\NativeMathAt{511.738}{761.124}{11.955}{\in}%
\TextAt{14}{524.191}{761.124}{A}%
\TextAt{8}{79.226}{744.078}{and}%
\TextAt{13}{102.097}{744.078}{s}%
\NativeMathAt{111.117}{744.078}{11.955}{\in}%
\TextAt{13}{122.593}{744.078}{S}%
\TextAt{8}{130.488}{744.078}{,}%
\TextAt{8}{137.773}{744.078}{then}%
\TextAt{8}{164.542}{744.078}{one}%
\TextAt{8}{186.11}{744.078}{has}%
\TextAt{13}{207.092}{744.078}{a}%
\TextAt{9}{216.743}{744.078}{\Character{61}}%
\TextAt{21}{230.548}{748.785}{a}%
\Vector{q 0 G 0 g 0.478 w 158.548 -26.933 m 163.046 -26.933 l S Q}%
\TextAt{21}{230.84}{739.955}{s}%
\TextAt{16}{238.973}{744.078}{\Character{1}}%
\TextAt{13}{245.024}{744.078}{s}%
\NativeMathAt{254.044}{744.078}{11.955}{\in}%
\TextAt{14}{265.52}{744.078}{A}%
\TextAt{16}{276.83}{744.078}{\Character{92}}%
\TextAt{14}{287.53}{744.078}{o}%
\TextAt{9}{296.88}{744.078}{\Character{61}}%
\TextAt{14}{309.491}{744.078}{a}%
\TextAt{8}{315.474}{744.078}{,}%
\TextAt{8}{322.759}{744.078}{so}%
\TextAt{8}{337.237}{744.078}{that}%
\TextAt{13}{362.715}{744.078}{a}%
\TextAt{9}{368.859}{744.078}{\Character{47}}%
\TextAt{13}{374.712}{744.078}{s}%
\TextAt{9}{383.732}{744.078}{\Character{61}}%
\TextAt{13}{396.343}{744.078}{a}%
\TextAt{17}{403.683}{748.785}{1}%
\Vector{q 0 G 0 g 0.478 w 331.683 -26.933 m 335.917 -26.933 l S Q}%
\TextAt{21}{403.842}{739.955}{s}%
\NativeMathAt{412.619}{744.078}{11.955}{\in}%
\TextAt{14}{424.094}{744.078}{a}%
\TextAt{13}{430.077}{744.078}{S}%
\TextAt{3}{437.973}{748.416}{\Character{0}}%
\TextAt{17}{444.559}{748.416}{1}%
\TextAt{8}{449.292}{744.078}{.}%
\TextAt{8}{458.065}{744.078}{As}%
\TextAt{14}{475.461}{744.078}{a}%
\TextAt{8}{485.45}{744.078}{is}%
\TextAt{8}{497.323}{744.078}{finitely}%
\ErrataAt{79.19}{727.017}{455.81}{generated, so is $\mathfrak A$, i.e., $\mathfrak oS^{-1}$ is noetherian. It follows from \hyperlink{nav.number.1.11.1}{(11.1)} that every nonzero prime ideal}%
\TextAt{8}{79.226}{709.986}{of}%
\TextAt{14}{91.905}{709.986}{o}%
\TextAt{13}{97.749}{709.986}{S}%
\TextAt{3}{105.644}{714.324}{\Character{0}}%
\TextAt{17}{112.231}{714.324}{1}%
\TextAt{8}{120.209}{709.986}{is}%
\TextAt{8}{131.321}{709.986}{maximal,}%
\TextAt{8}{181.86}{709.986}{because}%
\TextAt{8}{224.51}{709.986}{this}%
\TextAt{8}{246.681}{709.986}{holds}%
\TextAt{8}{276.658}{709.986}{in}%
\TextAt{14}{289.66}{709.986}{o}%
\TextAt{8}{295.504}{709.986}{.}%
\TextAt{8}{303.735}{709.986}{Finally,}%
\TextAt{14}{345.32}{709.986}{o}%
\TextAt{13}{351.164}{709.986}{S}%
\TextAt{3}{359.059}{714.324}{\Character{0}}%
\TextAt{17}{365.646}{714.324}{1}%
\TextAt{8}{373.624}{709.986}{is}%
\TextAt{8}{384.736}{709.986}{integrally}%
\TextAt{8}{436.125}{709.986}{closed,}%
\TextAt{8}{473.382}{709.986}{for}%
\TextAt{8}{490.616}{709.986}{if}%
\TextAt{13}{500.688}{709.986}{x}%
\NativeMathAt{510.661}{709.986}{11.955}{\in}%
\TextAt{13}{521.952}{709.986}{K}%
\TextAt{8}{79.226}{692.94}{satisfies}%
\TextAt{8}{122.336}{692.94}{the}%
\TextAt{8}{142.493}{692.94}{equation}%
\TextAt{13}{235.178}{675.894}{x}%
\TextAt{21}{241.83}{680.83}{n}%
\TextAt{9}{250.123}{675.894}{\Character{43}}%
\TextAt{13}{263.08}{683.982}{a}%
\TextAt{17}{269.225}{682.188}{1}%
\Vector{q 0 G 0 g 0.478 w 191.08 -95.117 m 201.957 -95.117 l S Q}%
\TextAt{13}{263.396}{667.693}{s}%
\TextAt{17}{268.91}{665.9}{1}%
\TextAt{13}{275.153}{675.894}{x}%
\TextAt{21}{281.805}{680.83}{n}%
\TextAt{3}{286.943}{680.83}{\Character{0}}%
\TextAt{17}{293.53}{680.83}{1}%
\TextAt{9}{300.919}{675.894}{\Character{43}}%
\TextAt{16}{312.68}{675.894}{\Character{1}}%
\TextAt{16}{317.976}{675.894}{\Character{1}}%
\TextAt{16}{323.284}{675.894}{\Character{1}}%
\TextAt{9}{329.284}{675.894}{\Character{43}}%
\TextAt{13}{342.241}{683.982}{a}%
\TextAt{21}{348.386}{682.188}{n}%
\Vector{q 0 G 0 g 0.478 w 270.241 -95.117 m 282.023 -95.117 l S Q}%
\TextAt{13}{342.557}{667.693}{s}%
\TextAt{21}{348.071}{665.9}{n}%
\TextAt{9}{358.539}{675.894}{\Character{61}}%
\TextAt{9}{370.949}{675.894}{0}%
\TextAt{8}{79.226}{649.313}{with}%
\TextAt{8}{106.718}{649.313}{coefficients}%
\TextAt{13}{166.787}{649.313}{a}%
\TextAt{21}{172.932}{647.52}{i}%
\TextAt{9}{176.314}{649.313}{\Character{47}}%
\TextAt{13}{182.167}{649.313}{s}%
\TextAt{21}{187.681}{647.52}{i}%
\NativeMathAt{195.8}{649.313}{11.955}{\in}%
\TextAt{14}{208.508}{649.313}{o}%
\TextAt{13}{214.352}{649.313}{S}%
\TextAt{3}{222.248}{653.652}{\Character{0}}%
\TextAt{17}{228.834}{653.652}{1}%
\TextAt{8}{233.566}{649.313}{,}%
\TextAt{8}{241.755}{649.313}{then}%
\TextAt{8}{269.247}{649.313}{multiplying}%
\TextAt{8}{332.52}{649.313}{it}%
\TextAt{8}{345.056}{649.313}{with}%
\TextAt{8}{372.548}{649.313}{the}%
\TextAt{13}{393.536}{649.313}{n}%
\TextAt{8}{400.523}{649.313}{-th}%
\TextAt{8}{420.208}{649.313}{power}%
\TextAt{8}{455.184}{649.313}{of}%
\TextAt{13}{469.345}{649.313}{s}%
\TextAt{9}{479.597}{649.313}{\Character{61}}%
\TextAt{13}{493.44}{649.313}{s}%
\TextAt{17}{498.954}{647.52}{1}%
\TextAt{16}{505.679}{649.313}{\Character{1}}%
\TextAt{16}{510.975}{649.313}{\Character{1}}%
\TextAt{16}{516.283}{649.313}{\Character{1}}%
\TextAt{13}{521.619}{649.313}{s}%
\TextAt{21}{527.133}{647.52}{n}%
\TextAt{8}{79.226}{632.267}{shows}%
\TextAt{8}{112.851}{632.267}{that}%
\TextAt{13}{138.227}{632.267}{sx}%
\TextAt{8}{154.297}{632.267}{is}%
\TextAt{8}{166.068}{632.267}{integral}%
\TextAt{8}{208.683}{632.267}{over}%
\TextAt{14}{233.737}{632.267}{o}%
\TextAt{8}{239.581}{632.267}{,}%
\TextAt{8}{246.739}{632.267}{whence}%
\TextAt{13}{287.704}{632.267}{sx}%
\NativeMathAt{303.203}{632.267}{11.955}{\in}%
\TextAt{14}{314.506}{632.267}{o}%
\TextAt{8}{324.254}{632.267}{and}%
\TextAt{8}{347.024}{632.267}{therefore}%
\TextAt{13}{396.131}{632.267}{x}%
\NativeMathAt{406.116}{632.267}{11.955}{\in}%
\TextAt{14}{417.419}{632.267}{o}%
\TextAt{13}{423.263}{632.267}{S}%
\TextAt{3}{431.158}{636.606}{\Character{0}}%
\TextAt{17}{437.745}{636.606}{1}%
\TextAt{8}{442.477}{632.267}{.}%
\TextAt{8}{450.946}{632.267}{This}%
\TextAt{8}{477.673}{632.267}{shows}%
\TextAt{8}{511.298}{632.267}{that}%
\TextAt{14}{79.226}{615.221}{o}%
\TextAt{13}{85.07}{615.221}{S}%
\TextAt{3}{92.966}{619.56}{\Character{0}}%
\TextAt{17}{99.552}{619.56}{1}%
\TextAt{8}{108.182}{615.221}{is}%
\TextAt{8}{119.946}{615.221}{a}%
\TextAt{8}{129.701}{615.221}{Dedekind}%
\TextAt{8}{181.886}{615.221}{domain.}%
\Vector{q 0 G 0 g 0.399 w 452.334 -158.779 m 452.334 -150.709 l S Q}%
\Vector{q 0 G 0 g 0.399 w 452.534 -150.908 m 459.707 -150.908 l S Q}%
\Vector{q 0 G 0 g 0.399 w 452.534 -158.579 m 459.707 -158.579 l S Q}%
\Vector{q 0 G 0 g 0.399 w 459.906 -158.779 m 459.906 -150.709 l S Q}%
\TextAt{12}{79.226}{580.723}{Proposition}%
\TextAt{12}{152.187}{580.723}{(11.5).}%
\TextAt{8}{194.45}{580.723}{Let}%
\TextAt{14}{215.636}{580.723}{o}%
\TextAt{8}{225.594}{580.723}{be}%
\TextAt{8}{241.735}{580.723}{a}%
\TextAt{8}{251.707}{580.723}{noetherian}%
\TextAt{8}{310.133}{580.723}{integral}%
\TextAt{8}{352.958}{580.723}{domain.}%
\TextAt{14}{399.787}{580.723}{o}%
\TextAt{8}{409.745}{580.723}{is}%
\TextAt{8}{421.726}{580.723}{a}%
\TextAt{8}{431.698}{580.723}{Dedekind}%
\TextAt{8}{484.099}{580.723}{domain}%
\TextAt{8}{525.943}{580.723}{if}%
\TextAt{8}{79.226}{563.677}{and}%
\TextAt{8}{101.989}{563.677}{only}%
\TextAt{8}{127.68}{563.677}{if,}%
\TextAt{8}{141.656}{563.677}{for}%
\TextAt{8}{159.541}{563.677}{all}%
\TextAt{8}{175.8}{563.677}{prime}%
\TextAt{8}{208.964}{563.677}{ideals}%
\TextAt{14}{241.541}{563.677}{p}%
\NativeMathAt{250.845}{563.677}{11.955}{\neq}%
\TextAt{9}{263.255}{563.677}{0}%
\TextAt{8}{269.124}{563.677}{,}%
\TextAt{8}{276.273}{563.677}{the}%
\TextAt{8}{296.43}{563.677}{localizations}%
\TextAt{14}{363.163}{563.677}{o}%
\TextAt{27}{369.007}{561.884}{p}%
\TextAt{8}{377.392}{563.677}{are}%
\TextAt{8}{396.903}{563.677}{discrete}%
\TextAt{8}{439.881}{563.677}{valuation}%
\TextAt{8}{490.954}{563.677}{rings.}%
\TextAt{19}{79.226}{529.178}{Proof.}%
\TextAt{8}{116}{529.178}{If}%
\TextAt{14}{128.592}{529.178}{o}%
\TextAt{8}{139.233}{529.178}{is}%
\TextAt{8}{151.896}{529.178}{a}%
\TextAt{8}{162.55}{529.178}{Dedekind}%
\TextAt{8}{215.634}{529.178}{domain,}%
\TextAt{8}{261.638}{529.178}{then}%
\TextAt{8}{289.197}{529.178}{so}%
\TextAt{8}{304.467}{529.178}{are}%
\TextAt{8}{324.877}{529.178}{the}%
\TextAt{8}{345.932}{529.178}{localizations}%
\TextAt{14}{413.565}{529.178}{o}%
\TextAt{27}{419.41}{527.385}{p}%
\TextAt{8}{423.896}{529.178}{.}%
\TextAt{8}{435.043}{529.178}{The}%
\TextAt{8}{459.996}{529.178}{maximal}%
\TextAt{8}{508.704}{529.178}{ideal}%
\TextAt{14}{79.226}{512.132}{m}%
\TextAt{9}{93.901}{512.132}{\Character{61}}%
\TextAt{14}{108.514}{512.132}{po}%
\TextAt{27}{120.342}{510.339}{p}%
\TextAt{8}{130.01}{512.132}{is}%
\TextAt{8}{143.057}{512.132}{the}%
\TextAt{8}{164.497}{512.132}{only}%
\TextAt{8}{191.473}{512.132}{nonzero}%
\TextAt{8}{236.333}{512.132}{prime}%
\TextAt{8}{270.78}{512.132}{ideal}%
\TextAt{8}{300.027}{512.132}{of}%
\TextAt{14}{314.641}{512.132}{o}%
\TextAt{27}{320.485}{510.339}{p}%
\TextAt{8}{324.972}{512.132}{.}%
\TextAt{8}{337.271}{512.132}{Therefore,}%
\TextAt{8}{395.125}{512.132}{choosing}%
\TextAt{8}{443.954}{512.132}{any}%
\TextAt{13}{467.355}{512.132}{\Character{25}}%
\NativeMathAt{479.932}{512.132}{11.955}{\in}%
\TextAt{14}{493.411}{512.132}{m}%
\TextAt{16}{506.109}{512.132}{\Character{110}}%
\TextAt{14}{515.619}{512.132}{m}%
\TextAt{17}{524.785}{516.471}{2}%
\TextAt{8}{529.518}{512.132}{,}%
\TextAt{8}{79.226}{495.086}{one}%
\TextAt{8}{101.207}{495.086}{necessarily}%
\TextAt{8}{160.057}{495.086}{finds}%
\TextAt{9}{188.6}{495.086}{\Character{40}}%
\TextAt{13}{193.153}{495.086}{\Character{25}}%
\TextAt{9}{200.222}{495.086}{\Character{41}}%
\TextAt{9}{208.973}{495.086}{\Character{61}}%
\TextAt{14}{222.296}{495.086}{m}%
\TextAt{8}{231.462}{495.086}{,}%
\TextAt{8}{239.263}{495.086}{and}%
\TextAt{8}{262.546}{495.086}{furthermore}%
\TextAt{14}{327.78}{495.086}{m}%
\TextAt{21}{336.947}{499.425}{n}%
\TextAt{9}{346.792}{495.086}{\Character{61}}%
\TextAt{9}{360.086}{495.086}{\Character{40}}%
\TextAt{13}{364.657}{495.086}{\Character{25}}%
\TextAt{21}{371.726}{499.425}{n}%
\TextAt{9}{377.363}{495.086}{\Character{41}}%
\TextAt{8}{381.915}{495.086}{.}%
\TextAt{8}{391.926}{495.086}{Thus}%
\TextAt{14}{422.096}{495.086}{o}%
\TextAt{27}{427.94}{493.293}{p}%
\TextAt{8}{436.845}{495.086}{is}%
\TextAt{8}{449.13}{495.086}{a}%
\TextAt{8}{459.406}{495.086}{principal}%
\TextAt{8}{508.704}{495.086}{ideal}%
\TextAt{8}{79.226}{478.04}{domain,}%
\TextAt{8}{124.106}{478.04}{and}%
\TextAt{8}{146.868}{478.04}{hence}%
\TextAt{8}{179.375}{478.04}{a}%
\TextAt{8}{189.13}{478.04}{discrete}%
\TextAt{8}{232.109}{478.04}{valuation}%
\TextAt{8}{283.181}{478.04}{ring.}%
\TextAt{8}{97.159}{458.006}{Letting}%
\TextAt{14}{138.297}{458.006}{p}%
\TextAt{8}{148.177}{458.006}{vary}%
\TextAt{8}{174.204}{458.006}{over}%
\TextAt{8}{199.25}{458.006}{all}%
\TextAt{8}{215.509}{458.006}{prime}%
\TextAt{8}{248.672}{458.006}{ideals}%
\NativeMathAt{281.25}{458.006}{11.955}{\neq}%
\TextAt{9}{293.66}{458.006}{0}%
\TextAt{8}{303.426}{458.006}{of}%
\TextAt{14}{316.756}{458.006}{o}%
\TextAt{8}{322.6}{458.006}{,}%
\TextAt{8}{329.749}{458.006}{we}%
\TextAt{8}{346.977}{458.006}{find}%
\TextAt{8}{370.385}{458.006}{in}%
\TextAt{8}{384.038}{458.006}{any}%
\TextAt{8}{406.155}{458.006}{case}%
\TextAt{8}{430.926}{458.006}{that}%
\TextAt{14}{280.773}{429.845}{o}%
\TextAt{9}{289.938}{429.845}{\Character{61}}%
\TextAt{20}{302.364}{441.202}{\Character{92}}%
\TextAt{27}{307.011}{416.096}{p}%
\TextAt{14}{317.64}{429.845}{o}%
\TextAt{27}{323.484}{428.052}{p}%
\TextAt{13}{327.971}{429.845}{\Character{58}}%
\TextAt{8}{79.226}{392.316}{For}%
\TextAt{8}{100.196}{392.316}{if}%
\TextAt{13}{110.919}{392.316}{a}%
\TextAt{9}{117.064}{392.316}{\Character{47}}%
\TextAt{13}{122.917}{392.316}{b}%
\NativeMathAt{131.215}{392.316}{11.955}{\in}%
\TextAt{20}{142.506}{401.283}{\Character{84}}%
\TextAt{27}{152.469}{388.829}{p}%
\TextAt{14}{158.948}{392.316}{o}%
\TextAt{27}{164.792}{390.523}{p}%
\TextAt{8}{169.279}{392.316}{,}%
\TextAt{8}{176.428}{392.316}{with}%
\TextAt{13}{203.088}{392.316}{a\Character{59}}%
\TextAt{13}{214.457}{392.316}{b}%
\NativeMathAt{222.776}{392.316}{11.955}{\in}%
\TextAt{14}{234.067}{392.316}{o}%
\TextAt{8}{239.911}{392.316}{,}%
\TextAt{8}{247.06}{392.316}{then}%
\TextAt{14}{251.505}{364.156}{a}%
\TextAt{9}{260.809}{364.156}{\Character{61}}%
\TextAt{16}{273.235}{364.156}{\Character{102}}%
\TextAt{13}{279.212}{364.156}{x}%
\NativeMathAt{289.185}{364.156}{11.955}{\in}%
\TextAt{14}{300.476}{364.156}{o}%
\TextAt{16}{309.641}{364.156}{\Character{106}}%
\TextAt{13}{316.283}{364.156}{xa}%
\NativeMathAt{332.401}{364.156}{11.955}{\in}%
\TextAt{13}{343.692}{364.156}{b}%
\TextAt{14}{348.669}{364.156}{o}%
\TextAt{16}{354.513}{364.156}{\Character{103}}%
\TextAt{8}{79.226}{335.995}{is}%
\TextAt{8}{90.959}{335.995}{an}%
\TextAt{8}{107.188}{335.995}{ideal}%
\TextAt{8}{135.12}{335.995}{which}%
\TextAt{8}{168.576}{335.995}{cannot}%
\TextAt{8}{206.921}{335.995}{be}%
\TextAt{8}{222.815}{335.995}{contained}%
\TextAt{8}{275.794}{335.995}{in}%
\TextAt{8}{289.416}{335.995}{any}%
\TextAt{8}{311.502}{335.995}{prime}%
\TextAt{8}{344.635}{335.995}{ideal}%
\TextAt{8}{372.568}{335.995}{of}%
\TextAt{14}{385.867}{335.995}{o}%
\TextAt{8}{391.711}{335.995}{.}%
\TextAt{8}{400.15}{335.995}{In}%
\TextAt{8}{414.74}{335.995}{fact,}%
\TextAt{8}{441.053}{335.995}{for}%
\TextAt{8}{458.907}{335.995}{any}%
\TextAt{14}{480.993}{335.995}{p}%
\TextAt{8}{486.977}{335.995}{,}%
\TextAt{8}{494.101}{335.995}{we}%
\TextAt{8}{511.298}{335.995}{may}%
\TextAt{8}{79.226}{318.949}{write}%
\TextAt{13}{109.153}{318.949}{a}%
\TextAt{9}{115.298}{318.949}{\Character{47}}%
\TextAt{13}{121.151}{318.949}{b}%
\TextAt{9}{129.476}{318.949}{\Character{61}}%
\TextAt{13}{141.927}{318.949}{c}%
\TextAt{9}{146.965}{318.949}{\Character{47}}%
\TextAt{13}{152.818}{318.949}{s}%
\TextAt{8}{162.245}{318.949}{with}%
\TextAt{13}{188.921}{318.949}{c}%
\NativeMathAt{197.306}{318.949}{11.955}{\in}%
\TextAt{14}{208.623}{318.949}{o}%
\TextAt{8}{214.467}{318.949}{,}%
\TextAt{13}{221.636}{318.949}{s}%
\NativeMathAt{230.497}{318.949}{11.955}{\notin}%
\TextAt{14}{241.814}{318.949}{p}%
\TextAt{8}{247.798}{318.949}{,}%
\TextAt{8}{254.966}{318.949}{so}%
\TextAt{8}{269.352}{318.949}{that}%
\TextAt{13}{294.736}{318.949}{sa}%
\TextAt{9}{309.742}{318.949}{\Character{61}}%
\TextAt{13}{322.194}{318.949}{bc}%
\TextAt{8}{332.209}{318.949}{,}%
\TextAt{8}{339.378}{318.949}{hence}%
\TextAt{13}{371.899}{318.949}{s}%
\NativeMathAt{380.76}{318.949}{11.955}{\in}%
\TextAt{14}{392.078}{318.949}{a}%
\TextAt{16}{400.728}{318.949}{\Character{110}}%
\TextAt{14}{409.373}{318.949}{p}%
\TextAt{8}{415.356}{318.949}{.}%
\TextAt{8}{423.851}{318.949}{As}%
\TextAt{14}{441.153}{318.949}{a}%
\TextAt{8}{451.049}{318.949}{is}%
\TextAt{8}{462.828}{318.949}{not}%
\TextAt{8}{483.658}{318.949}{contained}%
\TextAt{8}{79.226}{301.903}{in}%
\TextAt{8}{92.879}{301.903}{any}%
\TextAt{8}{114.996}{301.903}{maximal}%
\TextAt{8}{162.805}{301.903}{ideal,}%
\TextAt{8}{194.02}{301.903}{it}%
\TextAt{8}{205.724}{301.903}{follows}%
\TextAt{8}{244.16}{301.903}{that}%
\TextAt{14}{269.529}{301.903}{a}%
\TextAt{9}{278.832}{301.903}{\Character{61}}%
\TextAt{14}{291.258}{301.903}{o}%
\TextAt{8}{297.102}{301.903}{,}%
\TextAt{8}{304.251}{301.903}{hence}%
\TextAt{13}{336.757}{301.903}{a}%
\TextAt{9}{346.223}{301.903}{\Character{61}}%
\TextAt{9}{358.632}{301.903}{1}%
\TextAt{16}{367.158}{301.903}{\Character{1}}%
\TextAt{13}{373.136}{301.903}{a}%
\NativeMathAt{382.602}{301.903}{11.955}{\in}%
\TextAt{13}{393.893}{301.903}{b}%
\TextAt{14}{398.87}{301.903}{o}%
\TextAt{8}{404.714}{301.903}{,}%
\TextAt{8}{411.863}{301.903}{i.e.,}%
\TextAt{13}{433.968}{301.903}{a}%
\TextAt{9}{440.113}{301.903}{\Character{47}}%
\TextAt{13}{445.966}{301.903}{b}%
\NativeMathAt{454.264}{301.903}{11.955}{\in}%
\TextAt{14}{465.555}{301.903}{o}%
\TextAt{8}{471.399}{301.903}{.}%
\TextAt{8}{97.159}{281.868}{Suppose}%
\TextAt{8}{143.312}{281.868}{now}%
\TextAt{8}{167.946}{281.868}{that}%
\TextAt{8}{193.56}{281.868}{the}%
\TextAt{14}{213.962}{281.868}{o}%
\TextAt{27}{219.806}{280.075}{p}%
\TextAt{8}{228.435}{281.868}{are}%
\TextAt{8}{248.191}{281.868}{discrete}%
\TextAt{8}{291.415}{281.868}{valuation}%
\TextAt{8}{342.733}{281.868}{rings.}%
\TextAt{8}{376.7}{281.868}{Being}%
\TextAt{8}{409.942}{281.868}{principal}%
\TextAt{8}{458.964}{281.868}{ideal}%
\TextAt{8}{487.172}{281.868}{domains,}%
\ErrataAt{79.19}{264.829}{455.81}{they are integrally closed (see \hyperlink{nav.section.1.2}{\S2}), so $\mathfrak o=\bigcap_{\mathfrak p}\mathfrak o_{\mathfrak p}$ is integrally closed. Finally, for}%
\ErrataAt{79.19}{247.783}{455.81}{$\mathfrak p\ne0$ prime choose a maximal $\mathfrak q\supseteq\mathfrak p$. The DVR $\mathfrak o_{\mathfrak q}$ has only one nonzero prime,}%
\ErrataAt{79.19}{230.737}{455.81}{so $\mathfrak p\mathfrak o_{\mathfrak q}=\mathfrak q\mathfrak o_{\mathfrak q}$. Contracting gives $\mathfrak p=\mathfrak q$. Therefore $\mathfrak o$ is a Dedekind domain.}%
\Vector{q 0 G 0 g 0.399 w 452.334 -543.27 m 452.334 -535.2 l S Q}%
\Vector{q 0 G 0 g 0.399 w 452.534 -535.399 m 459.707 -535.399 l S Q}%
\Vector{q 0 G 0 g 0.399 w 452.534 -543.071 m 459.707 -543.071 l S Q}%
\Vector{q 0 G 0 g 0.399 w 459.906 -543.27 m 459.906 -535.2 l S Q}%
\TextAt{8}{97.159}{201.969}{For}%
\TextAt{8}{118.509}{201.969}{a}%
\TextAt{8}{128.645}{201.969}{Dedekind}%
\TextAt{8}{181.209}{201.969}{domain}%
\TextAt{14}{223.218}{201.969}{o}%
\TextAt{8}{229.062}{201.969}{,}%
\TextAt{8}{236.687}{201.969}{we}%
\TextAt{8}{254.295}{201.969}{have}%
\TextAt{8}{281.67}{201.969}{for}%
\TextAt{8}{299.936}{201.969}{each}%
\TextAt{8}{326.653}{201.969}{prime}%
\TextAt{8}{360.197}{201.969}{ideal}%
\TextAt{14}{388.541}{201.969}{p}%
\NativeMathAt{398.493}{201.969}{11.955}{\neq}%
\TextAt{9}{411.56}{201.969}{0}%
\TextAt{8}{421.698}{201.969}{the}%
\TextAt{8}{442.235}{201.969}{discrete}%
\TextAt{8}{485.594}{201.969}{valuation}%
\TextAt{8}{79.226}{184.923}{ring}%
\TextAt{14}{103.292}{184.923}{o}%
\TextAt{27}{109.136}{183.13}{p}%
\TextAt{8}{117.52}{184.923}{and}%
\TextAt{8}{140.283}{184.923}{the}%
\TextAt{8}{160.439}{184.923}{corresponding}%
\TextAt{8}{235.626}{184.923}{valuation}%
\TextAt{13}{271.267}{156.762}{v}%
\TextAt{27}{276.926}{154.969}{p}%
\TextAt{9}{284.734}{156.762}{\Character{58}}%
\TextAt{13}{291.306}{156.762}{K}%
\TextAt{3}{302.124}{161.698}{\Character{3}}%
\NativeArrowAt{310.177}{156.762}{11.955}{19.26}{\longrightarrow}%
\TextAt{18}{332.759}{156.762}{Z}%
\TextAt{8}{79.226}{128.601}{of}%
\TextAt{8}{93.07}{128.601}{the}%
\TextAt{8}{113.74}{128.601}{field}%
\TextAt{8}{139.61}{128.601}{of}%
\TextAt{8}{153.454}{128.601}{fractions.}%
\TextAt{8}{207.414}{128.601}{The}%
\TextAt{8}{231.982}{128.601}{significance}%
\TextAt{8}{294.339}{128.601}{of}%
\TextAt{8}{308.183}{128.601}{these}%
\TextAt{8}{338.668}{128.601}{valuations}%
\TextAt{8}{394.869}{128.601}{lies}%
\TextAt{8}{415.599}{128.601}{in}%
\TextAt{8}{429.765}{128.601}{their}%
\TextAt{8}{458.242}{128.601}{relation}%
\TextAt{8}{501.686}{128.601}{to}%
\TextAt{8}{516.51}{128.601}{the}%
\TextAt{8}{79.226}{111.555}{prime}%
\TextAt{8}{112.39}{111.555}{ideal}%
\TextAt{8}{140.353}{111.555}{factorization.}%
\TextAt{8}{212.881}{111.555}{If}%
\TextAt{13}{224.573}{111.555}{x}%
\NativeMathAt{234.546}{111.555}{11.955}{\in}%
\TextAt{13}{245.837}{111.555}{K}%
\TextAt{3}{256.655}{115.894}{\Character{3}}%
\TextAt{8}{265.284}{111.555}{and}%
\TextAt{9}{274.249}{83.395}{\Character{40}}%
\TextAt{13}{278.801}{83.395}{x}%
\TextAt{9}{285.453}{83.395}{\Character{41}}%
\TextAt{9}{293.308}{83.395}{\Character{61}}%
\TextAt{20}{305.752}{94.752}{\Character{89}}%
\TextAt{27}{311.395}{69.646}{p}%
\TextAt{14}{323.02}{83.395}{p}%
\TextAt{21}{329.003}{88.331}{\Character{23}}%
\TextAt{28}{333.166}{87.335}{p}%
\TextAt{8}{300.14}{43.033}{64}%
\Destination{78.23}{809.815}{nav.page.64}%
\Destination{79.36}{742.055}{nav.print.69}%
\Destination{79.36}{594.755}{nav.number.1.11.5}%
\Hotlink{78.726}{245.555}{30.936}{11.047}{nav.number.1.11.1}%
\Destination{79.226}{744.017}{errata.49}%
\Destination{79.226}{281.829}{errata.50}%
}

% Page 77 (body page 65)
\Page{%
\TextAt{2}{79.226}{796.943}{I.}%
\TextAt{2}{89.83}{796.943}{Algebraic}%
\TextAt{2}{138.648}{796.943}{Integers}%
\Vector{q 0 G 0 g 0.299 w 7.226 18.729 m 460.769 18.729 l S Q}%
\TextAt{8}{79.226}{761.124}{is}%
\TextAt{8}{90.99}{761.124}{the}%
\TextAt{8}{111.146}{761.124}{prime}%
\TextAt{8}{144.31}{761.124}{factorization}%
\TextAt{8}{212.287}{761.124}{of}%
\TextAt{8}{225.617}{761.124}{the}%
\TextAt{8}{245.774}{761.124}{principal}%
\TextAt{8}{294.551}{761.124}{ideal}%
\TextAt{9}{322.514}{761.124}{\Character{40}}%
\TextAt{13}{327.066}{761.124}{x}%
\TextAt{9}{333.718}{761.124}{\Character{41}}%
\TextAt{8}{338.271}{761.124}{,}%
\TextAt{8}{345.42}{761.124}{then,}%
\TextAt{8}{375.332}{761.124}{for}%
\TextAt{8}{393.217}{761.124}{each}%
\TextAt{14}{419.554}{761.124}{p}%
\TextAt{8}{425.537}{761.124}{,}%
\TextAt{8}{432.686}{761.124}{one}%
\TextAt{8}{454.146}{761.124}{has}%
\TextAt{13}{278.412}{732.139}{\Character{23}}%
\TextAt{27}{284.196}{730.346}{p}%
\TextAt{9}{292.004}{732.139}{\Character{61}}%
\TextAt{13}{304.429}{732.139}{v}%
\TextAt{27}{310.088}{730.346}{p}%
\TextAt{9}{314.575}{732.139}{\Character{40}}%
\TextAt{13}{319.128}{732.139}{x}%
\TextAt{9}{325.78}{732.139}{\Character{41}}%
\TextAt{13}{330.332}{732.139}{\Character{58}}%
\TextAt{8}{79.226}{703.154}{In}%
\TextAt{8}{95.376}{703.154}{fact,}%
\TextAt{8}{123.625}{703.154}{for}%
\TextAt{8}{143.039}{703.154}{a}%
\TextAt{8}{154.323}{703.154}{fixed}%
\TextAt{8}{184.138}{703.154}{prime}%
\TextAt{8}{218.831}{703.154}{ideal}%
\TextAt{14}{248.323}{703.154}{q}%
\NativeMathAt{260.097}{703.154}{11.955}{\neq}%
\TextAt{9}{275.113}{703.154}{0}%
\TextAt{8}{286.408}{703.154}{of}%
\TextAt{14}{301.267}{703.154}{o}%
\TextAt{8}{307.111}{703.154}{,}%
\TextAt{8}{316.172}{703.154}{the}%
\TextAt{8}{337.857}{703.154}{first}%
\TextAt{8}{363.512}{703.154}{equation}%
\TextAt{8}{412.85}{703.154}{above}%
\TextAt{8}{447.554}{703.154}{implies}%
\TextAt{8}{488.81}{703.154}{(because}%
\TextAt{14}{79.226}{686.108}{po}%
\TextAt{27}{91.053}{684.315}{q}%
\TextAt{9}{98.771}{686.108}{\Character{61}}%
\TextAt{14}{111.196}{686.108}{o}%
\TextAt{27}{117.04}{684.315}{q}%
\TextAt{8}{125.334}{686.108}{for}%
\TextAt{14}{143.219}{686.108}{p}%
\NativeMathAt{152.523}{686.108}{11.955}{\neq}%
\TextAt{14}{164.949}{686.108}{q}%
\TextAt{8}{170.796}{686.108}{)}%
\TextAt{8}{179.248}{686.108}{that}%
\TextAt{13}{219.43}{648.951}{x}%
\TextAt{14}{226.082}{648.951}{o}%
\TextAt{27}{231.926}{647.158}{q}%
\TextAt{9}{239.643}{648.951}{\Character{61}}%
\TextAt{20}{252.069}{669.394}{\Character{32}}%
\TextAt{20}{261.533}{660.308}{\Character{89}}%
\TextAt{27}{267.177}{635.202}{p}%
\TextAt{14}{278.802}{648.951}{p}%
\TextAt{21}{284.785}{653.887}{\Character{23}}%
\TextAt{28}{288.948}{652.891}{p}%
\TextAt{20}{293.529}{669.394}{\Character{33}}%
\TextAt{14}{304.986}{648.951}{o}%
\TextAt{27}{310.83}{647.158}{q}%
\TextAt{9}{318.547}{648.951}{\Character{61}}%
\TextAt{14}{330.973}{648.951}{q}%
\TextAt{21}{336.82}{653.887}{\Character{23}}%
\TextAt{28}{340.983}{652.891}{q}%
\TextAt{14}{345.49}{648.951}{o}%
\TextAt{27}{351.335}{647.158}{q}%
\TextAt{9}{359.052}{648.951}{\Character{61}}%
\TextAt{14}{371.477}{648.951}{m}%
\TextAt{21}{380.644}{653.887}{\Character{23}}%
\TextAt{28}{384.807}{652.891}{q}%
\TextAt{27}{380.644}{645.995}{q}%
\TextAt{13}{389.314}{648.951}{\Character{58}}%
\TextAt{8}{79.226}{610.598}{Hence}%
\TextAt{8}{115.903}{610.598}{indeed}%
\TextAt{13}{154.863}{610.598}{v}%
\TextAt{27}{160.522}{608.805}{q}%
\TextAt{9}{164.919}{610.599}{\Character{40}}%
\TextAt{13}{169.471}{610.599}{x}%
\TextAt{9}{176.123}{610.599}{\Character{41}}%
\TextAt{9}{187.217}{610.599}{\Character{61}}%
\TextAt{13}{202.894}{610.599}{\Character{23}}%
\TextAt{27}{208.678}{608.805}{q}%
\TextAt{8}{213.075}{610.598}{.}%
\TextAt{8}{227.22}{610.598}{In}%
\TextAt{8}{243.741}{610.598}{view}%
\TextAt{8}{272.622}{610.598}{of}%
\TextAt{8}{287.852}{610.598}{this}%
\TextAt{8}{312.573}{610.598}{relation,}%
\TextAt{8}{361.13}{610.598}{the}%
\TextAt{8}{383.185}{610.598}{valuations}%
\TextAt{13}{440.772}{610.598}{v}%
\TextAt{27}{446.431}{608.805}{p}%
\TextAt{8}{456.714}{610.598}{are}%
\TextAt{8}{478.124}{610.598}{also}%
\TextAt{8}{503.503}{610.598}{called}%
\TextAt{12}{79.226}{593.553}{exponential}%
\TextAt{12}{151.483}{593.553}{valuations}%
\TextAt{8}{210.613}{593.553}{.}%
\TextAt{8}{97.159}{573.518}{The}%
\TextAt{8}{120.539}{573.518}{reader}%
\TextAt{8}{155.634}{573.518}{should}%
\TextAt{8}{192.093}{573.518}{check}%
\TextAt{8}{222.956}{573.518}{that}%
\TextAt{8}{247.651}{573.518}{the}%
\TextAt{8}{267.133}{573.518}{localization}%
\TextAt{8}{328.578}{573.518}{of}%
\TextAt{8}{341.234}{573.518}{the}%
\TextAt{8}{360.716}{573.518}{ring}%
\TextAt{18}{384.108}{573.518}{Z}%
\TextAt{8}{395.301}{573.518}{at}%
\TextAt{8}{408.937}{573.518}{the}%
\TextAt{8}{428.42}{573.518}{prime}%
\TextAt{8}{460.909}{573.518}{ideal}%
\TextAt{9}{488.198}{573.518}{\Character{40}}%
\TextAt{13}{492.75}{573.518}{p}%
\TextAt{9}{498.626}{573.518}{\Character{41}}%
\TextAt{9}{506.481}{573.518}{\Character{61}}%
\TextAt{13}{518.924}{573.518}{p}%
\TextAt{18}{524.799}{573.518}{Z}%
\TextAt{8}{79.226}{556.472}{is}%
\TextAt{8}{90.99}{556.472}{given}%
\TextAt{8}{121.559}{556.472}{by}%
\TextAt{18}{233.264}{539.426}{Z}%
\TextAt{17}{241.234}{537.566}{\Character{40}}%
\TextAt{21}{244.527}{537.566}{p}%
\TextAt{17}{248.79}{537.566}{\Character{41}}%
\TextAt{9}{255.902}{539.426}{\Character{61}}%
\TextAt{20}{268.327}{552.696}{\Character{110}}%
\TextAt{13}{277.493}{547.514}{a}%
\Vector{q 0 G 0 g 0.478 w 205.493 -231.585 m 211.638 -231.585 l S Q}%
\TextAt{13}{278.077}{531.225}{b}%
\TextAt{20}{288.731}{553.174}{\Character{12}}%
\TextAt{20}{288.731}{546.001}{\Character{12}}%
\TextAt{20}{288.731}{538.828}{\Character{12}}%
\TextAt{13}{296.613}{539.426}{a\Character{59}}%
\TextAt{13}{307.982}{539.426}{b}%
\NativeMathAt{316.3}{539.426}{11.955}{\in}%
\TextAt{18}{327.591}{539.426}{Z}%
\TextAt{13}{335.561}{539.426}{\Character{59}}%
\TextAt{13}{344.683}{539.426}{p}%
\TextAt{18}{353.899}{539.426}{\Character{45}}%
\TextAt{13}{360.541}{539.426}{b}%
\TextAt{20}{365.518}{552.696}{\Character{111}}%
\TextAt{13}{375.48}{539.426}{\Character{58}}%
\TextAt{8}{79.226}{514.794}{The}%
\TextAt{8}{103.361}{514.794}{maximal}%
\TextAt{8}{151.251}{514.794}{ideal}%
\TextAt{13}{179.296}{514.794}{p}%
\TextAt{18}{185.171}{514.794}{Z}%
\TextAt{17}{193.141}{512.934}{\Character{40}}%
\TextAt{21}{196.434}{512.934}{p}%
\TextAt{17}{200.697}{512.934}{\Character{41}}%
\TextAt{8}{208.467}{514.794}{consists}%
\TextAt{8}{251.659}{514.794}{of}%
\TextAt{8}{265.07}{514.794}{all}%
\TextAt{8}{281.411}{514.794}{fractions}%
\TextAt{13}{329.36}{514.794}{a}%
\TextAt{9}{335.505}{514.794}{\Character{47}}%
\TextAt{13}{341.358}{514.794}{b}%
\TextAt{8}{350.314}{514.794}{satisfying}%
\TextAt{13}{402.556}{514.794}{p}%
\TextAt{16}{411.891}{514.794}{\Character{106}}%
\TextAt{13}{418.671}{514.794}{a}%
\TextAt{8}{424.816}{514.794}{,}%
\TextAt{13}{432.067}{514.794}{p}%
\TextAt{18}{441.401}{514.794}{\Character{45}}%
\TextAt{13}{448.182}{514.794}{b}%
\TextAt{8}{453.159}{514.794}{,}%
\TextAt{8}{460.41}{514.794}{and}%
\TextAt{8}{483.254}{514.794}{the}%
\TextAt{8}{503.491}{514.794}{group}%
\TextAt{8}{79.226}{497.748}{of}%
\TextAt{8}{92.556}{497.748}{units}%
\TextAt{8}{121.882}{497.748}{consists}%
\TextAt{8}{164.993}{497.748}{of}%
\TextAt{8}{178.323}{497.748}{all}%
\TextAt{8}{194.582}{497.748}{fractions}%
\TextAt{13}{242.45}{497.748}{a}%
\TextAt{9}{248.595}{497.748}{\Character{47}}%
\TextAt{13}{254.448}{497.748}{b}%
\TextAt{8}{263.323}{497.748}{satisfying}%
\TextAt{13}{315.483}{497.748}{p}%
\TextAt{18}{324.679}{497.748}{\Character{45}}%
\TextAt{13}{331.321}{497.748}{ab}%
\TextAt{8}{342.443}{497.748}{.}%
\TextAt{8}{350.891}{497.748}{The}%
\TextAt{8}{374.945}{497.748}{valuation}%
\TextAt{8}{426.017}{497.748}{associated}%
\TextAt{8}{481.752}{497.748}{to}%
\TextAt{18}{496.063}{497.748}{Z}%
\TextAt{17}{504.033}{495.888}{\Character{40}}%
\TextAt{21}{507.326}{495.888}{p}%
\TextAt{17}{511.589}{495.888}{\Character{41}}%
\TextAt{8}{515.38}{497.748}{,}%
\TextAt{13}{254.033}{468.763}{v}%
\TextAt{21}{259.692}{466.97}{p}%
\TextAt{9}{267.774}{468.763}{\Character{58}}%
\TextAt{18}{274.346}{468.763}{Q}%
\NativeArrowAt{286.965}{468.763}{11.955}{19.26}{\longrightarrow}%
\TextAt{18}{309.547}{468.763}{Z}%
\TextAt{16}{320.174}{468.763}{\Character{91}}%
\TextAt{16}{330.79}{468.763}{\Character{102}}%
\NativeMathAt{336.768}{468.763}{11.955}{\infty}%
\TextAt{16}{348.723}{468.763}{\Character{103}}%
\TextAt{13}{354.711}{468.763}{\Character{59}}%
\TextAt{8}{79.226}{439.778}{is}%
\TextAt{8}{91.254}{439.778}{called}%
\TextAt{8}{124.682}{439.778}{the}%
\TextAt{13}{145.103}{439.778}{p}%
\TextAt{12}{150.978}{439.778}{-adic}%
\TextAt{12}{183.979}{439.778}{valuation}%
\TextAt{8}{241.963}{439.778}{of}%
\TextAt{18}{255.557}{439.778}{Q}%
\TextAt{8}{264.855}{439.778}{.}%
\TextAt{8}{274.096}{439.778}{The}%
\TextAt{8}{298.414}{439.778}{valuation}%
\TextAt{13}{349.751}{439.778}{v}%
\TextAt{21}{355.41}{437.985}{p}%
\TextAt{9}{360.171}{439.778}{\Character{40}}%
\TextAt{13}{364.723}{439.778}{x}%
\TextAt{9}{371.376}{439.778}{\Character{41}}%
\TextAt{8}{380.089}{439.778}{of}%
\TextAt{8}{393.684}{439.778}{an}%
\TextAt{8}{410.207}{439.778}{element}%
\TextAt{13}{453.713}{439.778}{x}%
\NativeMathAt{464.136}{439.778}{11.955}{\in}%
\TextAt{18}{475.877}{439.778}{Q}%
\TextAt{3}{485.176}{444.116}{\Character{3}}%
\TextAt{8}{494.069}{439.778}{is}%
\TextAt{8}{506.098}{439.778}{given}%
\TextAt{8}{79.226}{422.732}{by}%
\TextAt{13}{280.467}{405.686}{v}%
\TextAt{21}{286.126}{403.893}{p}%
\TextAt{9}{290.887}{405.686}{\Character{40}}%
\TextAt{13}{295.439}{405.686}{x}%
\TextAt{9}{302.091}{405.686}{\Character{41}}%
\TextAt{9}{309.946}{405.686}{\Character{61}}%
\TextAt{13}{322.39}{405.686}{\Character{23}\Character{59}}%
\TextAt{8}{79.226}{382.171}{where}%
\TextAt{13}{113.035}{382.171}{x}%
\TextAt{9}{123.008}{382.171}{\Character{61}}%
\TextAt{13}{135.434}{382.171}{p}%
\TextAt{21}{141.309}{386.51}{\Character{23}}%
\TextAt{13}{146.491}{382.171}{a}%
\TextAt{9}{152.636}{382.171}{\Character{47}}%
\TextAt{13}{158.489}{382.171}{b}%
\TextAt{8}{167.363}{382.171}{with}%
\TextAt{8}{194.023}{382.171}{integers}%
\TextAt{13}{237.337}{382.171}{a\Character{59}}%
\TextAt{13}{248.706}{382.171}{b}%
\TextAt{8}{257.601}{382.171}{relatively}%
\TextAt{8}{308.661}{382.171}{prime}%
\TextAt{8}{341.825}{382.171}{to}%
\TextAt{13}{356.135}{382.171}{p}%
\TextAt{8}{362.01}{382.171}{.}%
\TextAt{8}{97.159}{362.137}{To}%
\TextAt{8}{114.386}{362.137}{end}%
\TextAt{8}{136.491}{362.137}{this}%
\TextAt{8}{159.314}{362.137}{section,}%
\TextAt{8}{201.647}{362.137}{we}%
\TextAt{8}{218.874}{362.137}{now}%
\TextAt{8}{243.263}{362.137}{want}%
\TextAt{8}{271.884}{362.137}{to}%
\TextAt{8}{286.194}{362.137}{compare}%
\TextAt{8}{333.022}{362.137}{a}%
\TextAt{8}{342.778}{362.137}{Dedekind}%
\TextAt{8}{394.962}{362.137}{domain}%
\TextAt{14}{436.59}{362.137}{o}%
\TextAt{8}{446.332}{362.137}{to}%
\TextAt{8}{460.642}{362.137}{the}%
\TextAt{8}{480.798}{362.137}{ring}%
\TextAt{14}{172.487}{329.415}{o}%
\TextAt{9}{178.332}{329.415}{\Character{40}}%
\TextAt{13}{182.884}{329.415}{X}%
\TextAt{9}{193.539}{329.415}{\Character{41}}%
\TextAt{9}{201.394}{329.415}{\Character{61}}%
\TextAt{20}{213.838}{346.272}{\Character{26}}%
\TextAt{13}{223.999}{337.503}{f}%
\Vector{q 0 G 0 g 0.478 w 151.999 -441.596 m 159.046 -441.596 l S Q}%
\TextAt{13}{224.506}{321.215}{g}%
\TextAt{20}{236.139}{346.751}{\Character{12}}%
\TextAt{20}{236.139}{339.577}{\Character{12}}%
\TextAt{20}{236.139}{332.404}{\Character{12}}%
\TextAt{20}{236.139}{325.231}{\Character{12}}%
\TextAt{13}{244.021}{329.415}{f}%
\TextAt{13}{250.405}{329.415}{\Character{59}}%
\TextAt{13}{255.642}{329.415}{g}%
\NativeMathAt{265.017}{329.415}{11.955}{\in}%
\TextAt{14}{276.308}{329.415}{o}%
\TextAt{13}{282.152}{329.415}{\Character{59}}%
\TextAt{13}{291.274}{329.415}{g}%
\NativeMathAt{300.648}{329.415}{11.955}{\not\equiv}%
\TextAt{9}{313.268}{329.415}{0}%
\TextAt{9}{331.057}{329.415}{\Character{40}}%
\TextAt{8}{335.628}{329.415}{mod}%
\TextAt{14}{361.73}{329.415}{p}%
\TextAt{9}{367.713}{329.415}{\Character{41}}%
\TextAt{8}{376.163}{329.415}{for}%
\TextAt{14}{394.048}{329.415}{p}%
\NativeMathAt{403.352}{329.415}{11.955}{\in}%
\TextAt{13}{414.643}{329.415}{X}%
\TextAt{20}{425.298}{346.272}{\Character{27}}%
\TextAt{13}{436.257}{329.415}{\Character{59}}%
\ErrataAt{79.22}{296.755}{453.64}{where $X$ consists of nonzero primes of $\mathfrak o$ and contains almost all of them. Assume that}%
\ErrataAt{79.22}{279.655}{453.64}{each omitted prime has an element outside all primes in $X$. By \hyperlink{nav.number.1.11.1}{(11.1)}, the nonzero primes of $\mathfrak o(X)$ are}%
\ErrataAt{79.22}{262.655}{453.64}{$\mathfrak p_X=\mathfrak p\mathfrak o(X)$ for $\mathfrak p\in X$, and $\mathfrak o$ and $\mathfrak o(X)$ have the same localizations}%
\TextAt{14}{272.815}{233.629}{o}%
\TextAt{27}{278.659}{231.836}{p}%
\TextAt{9}{286.467}{233.629}{\Character{61}}%
\TextAt{14}{298.892}{233.629}{o}%
\TextAt{9}{304.736}{233.629}{\Character{40}}%
\TextAt{13}{309.289}{233.629}{X}%
\TextAt{9}{319.944}{233.629}{\Character{41}}%
\TextAt{27}{324.496}{231.836}{p}%
\TextAt{22}{328.485}{230.497}{X}%
\TextAt{13}{335.929}{233.629}{\Character{58}}%
\TextAt{8}{79.226}{204.644}{We}%
\TextAt{8}{100.014}{204.644}{denote}%
\TextAt{8}{138.376}{204.644}{by}%
\TextAt{13}{155.278}{204.644}{C}%
\TextAt{13}{164.496}{204.644}{l}%
\TextAt{9}{168.262}{204.644}{\Character{40}}%
\TextAt{14}{172.814}{204.644}{o}%
\TextAt{9}{178.658}{204.644}{\Character{41}}%
\TextAt{8}{183.21}{204.644}{,}%
\TextAt{8}{191.164}{204.644}{resp.}%
\TextAt{13}{222.416}{204.644}{C}%
\TextAt{13}{231.634}{204.644}{l}%
\TextAt{9}{235.399}{204.644}{\Character{40}}%
\TextAt{14}{239.951}{204.644}{o}%
\TextAt{9}{245.796}{204.644}{\Character{40}}%
\TextAt{13}{250.348}{204.644}{X}%
\TextAt{9}{261.003}{204.644}{\Character{41}\Character{41}}%
\TextAt{8}{270.108}{204.644}{,}%
\TextAt{8}{278.061}{204.644}{the}%
\TextAt{8}{298.861}{204.644}{ideal}%
\TextAt{8}{327.467}{204.644}{class}%
\TextAt{8}{355.548}{204.644}{groups}%
\TextAt{8}{393.981}{204.644}{of}%
\TextAt{14}{407.955}{204.644}{o}%
\TextAt{8}{413.799}{204.644}{,}%
\TextAt{8}{421.752}{204.644}{resp.}%
\TextAt{14}{453.004}{204.644}{o}%
\TextAt{9}{458.849}{204.644}{\Character{40}}%
\TextAt{13}{463.401}{204.644}{X}%
\TextAt{9}{474.056}{204.644}{\Character{41}}%
\TextAt{8}{478.608}{204.644}{.}%
\TextAt{8}{488.987}{204.644}{They,}%
\TextAt{8}{522.297}{204.644}{as}%
\TextAt{8}{79.226}{187.598}{well}%
\TextAt{8}{102.957}{187.598}{as}%
\TextAt{8}{117.327}{187.598}{the}%
\TextAt{8}{137.484}{187.598}{groups}%
\TextAt{8}{175.274}{187.598}{of}%
\TextAt{8}{188.604}{187.598}{units}%
\TextAt{14}{217.93}{187.598}{o}%
\TextAt{3}{223.774}{191.937}{\Character{3}}%
\TextAt{8}{232.404}{187.598}{and}%
\TextAt{14}{255.167}{187.598}{o}%
\TextAt{9}{261.011}{187.598}{\Character{40}}%
\TextAt{13}{265.563}{187.598}{X}%
\TextAt{9}{276.218}{187.598}{\Character{41}}%
\TextAt{3}{280.771}{191.937}{\Character{3}}%
\TextAt{8}{285.503}{187.598}{,}%
\TextAt{8}{292.652}{187.598}{are}%
\TextAt{8}{312.163}{187.598}{related}%
\TextAt{8}{351.185}{187.598}{by}%
\TextAt{8}{367.444}{187.598}{the}%
\TextAt{8}{387.6}{187.598}{following}%
\TextAt{12}{79.226}{161.591}{Proposition}%
\TextAt{12}{151.937}{161.591}{(11.6).}%
\TextAt{8}{193.984}{161.591}{There}%
\TextAt{8}{227.793}{161.591}{is}%
\TextAt{8}{239.557}{161.591}{a}%
\TextAt{8}{249.312}{161.591}{canonical}%
\TextAt{8}{300.696}{161.591}{exact}%
\TextAt{8}{331.588}{161.591}{sequence}%
\TextAt{9}{175.493}{133.102}{1}%
\NativeArrowAt{184.667}{133.102}{11.955}{19.26}{\longrightarrow}%
\TextAt{14}{207.249}{133.102}{o}%
\TextAt{3}{213.093}{138.038}{\Character{3}}%
\NativeArrowAt{221.146}{133.102}{11.955}{19.26}{\longrightarrow}%
\TextAt{14}{243.728}{133.102}{o}%
\TextAt{9}{249.573}{133.102}{\Character{40}}%
\TextAt{13}{254.125}{133.102}{X}%
\TextAt{9}{264.78}{133.102}{\Character{41}}%
\TextAt{3}{269.332}{138.038}{\Character{3}}%
\NativeArrowAt{277.386}{133.102}{11.955}{19.26}{\longrightarrow}%
\TextAt{20}{299.968}{144.459}{\Character{77}}%
\TextAt{27}{300.39}{118.556}{p}%
\TextAt{17}{305.319}{118.556}{\Character{47}}%
\NativeMathAt{304.379}{118.556}{7.97}{\in}%
\TextAt{21}{310.024}{118.556}{X}%
\TextAt{13}{320.026}{133.102}{K}%
\TextAt{3}{330.843}{138.038}{\Character{3}}%
\TextAt{9}{335.576}{133.102}{\Character{47}}%
\TextAt{14}{341.429}{133.102}{o}%
\TextAt{3}{347.273}{138.038}{\Character{3}}%
\TextAt{27}{347.273}{130.146}{p}%
\NativeArrowAt{277.386}{99.937}{11.955}{19.26}{\longrightarrow}%
\TextAt{13}{299.968}{99.937}{C}%
\TextAt{13}{309.185}{99.937}{l}%
\TextAt{9}{312.951}{99.937}{\Character{40}}%
\TextAt{14}{317.503}{99.937}{o}%
\TextAt{9}{323.348}{99.937}{\Character{41}}%
\NativeArrowAt{331.221}{99.937}{11.955}{19.26}{\longrightarrow}%
\TextAt{13}{353.803}{99.937}{C}%
\TextAt{13}{363.021}{99.937}{l}%
\TextAt{9}{366.786}{99.937}{\Character{40}}%
\TextAt{14}{371.338}{99.937}{o}%
\TextAt{9}{377.183}{99.937}{\Character{40}}%
\TextAt{13}{381.735}{99.937}{X}%
\TextAt{9}{392.39}{99.937}{\Character{41}\Character{41}}%
\NativeArrowAt{404.816}{99.937}{11.955}{19.26}{\longrightarrow}%
\TextAt{9}{427.398}{99.937}{1}%
\TextAt{13}{433.251}{99.937}{\Character{59}}%
\TextAt{8}{79.226}{72.921}{and}%
\TextAt{8}{101.989}{72.921}{one}%
\TextAt{8}{123.448}{72.921}{has}%
\TextAt{13}{144.322}{72.921}{K}%
\TextAt{3}{155.14}{77.26}{\Character{3}}%
\TextAt{9}{159.872}{72.921}{\Character{47}}%
\TextAt{14}{165.725}{72.921}{o}%
\TextAt{3}{171.569}{77.26}{\Character{3}}%
\TextAt{27}{171.569}{69.966}{p}%
\NativeMathAt{179.622}{72.921}{11.955}{\cong}%
\TextAt{18}{192.242}{72.921}{Z}%
\TextAt{8}{200.212}{72.921}{.}%
\TextAt{8}{300.14}{43.033}{65}%
\Destination{78.23}{809.815}{nav.page.65}%
\Destination{79.36}{377.155}{nav.print.70}%
\Destination{79.36}{175.555}{nav.number.1.11.6}%
\Destination{79.226}{310.755}{errata.10}%
}

% Page 78 (body page 66)
\Page{%
\TextAt{2}{79.226}{796.943}{I.}%
\TextAt{2}{89.83}{796.943}{Algebraic}%
\TextAt{2}{138.648}{796.943}{Integers}%
\Vector{q 0 G 0 g 0.299 w 7.226 18.729 m 460.769 18.729 l S Q}%
\TextAt{19}{79.226}{732.123}{Proof.}%
\TextAt{8}{116}{732.123}{The}%
\TextAt{8}{139.11}{732.123}{first}%
\TextAt{8}{162.292}{732.123}{arrow}%
\TextAt{8}{194.201}{732.123}{is}%
\TextAt{8}{205.021}{732.123}{inclusion}%
\TextAt{8}{252.914}{732.123}{and}%
\TextAt{8}{274.733}{732.123}{the}%
\TextAt{8}{293.945}{732.123}{second}%
\TextAt{8}{330.78}{732.123}{one}%
\TextAt{8}{351.295}{732.123}{is}%
\TextAt{8}{362.116}{732.123}{induced}%
\TextAt{8}{404.736}{732.123}{by}%
\TextAt{8}{420.052}{732.123}{the}%
\TextAt{8}{439.264}{732.123}{inclusion}%
\TextAt{14}{487.157}{732.123}{o}%
\TextAt{9}{493.001}{732.123}{\Character{40}}%
\TextAt{13}{497.554}{732.123}{X}%
\TextAt{9}{508.209}{732.123}{\Character{41}}%
\TextAt{3}{512.761}{736.461}{\Character{3}}%
\NativeMathAt{520.814}{732.123}{11.955}{\rightarrow}%
\TextAt{13}{79.226}{715.077}{K}%
\TextAt{3}{90.044}{719.415}{\Character{3}}%
\TextAt{8}{94.776}{715.077}{,}%
\TextAt{8}{102.577}{715.077}{followed}%
\TextAt{8}{148.301}{715.077}{by}%
\TextAt{8}{165.082}{715.077}{the}%
\TextAt{8}{185.759}{715.077}{projections}%
\TextAt{13}{246.176}{715.077}{K}%
\TextAt{3}{256.994}{719.415}{\Character{3}}%
\NativeMathAt{265.936}{715.077}{11.955}{\rightarrow}%
\TextAt{13}{282.1}{715.077}{K}%
\TextAt{3}{292.918}{719.415}{\Character{3}}%
\TextAt{9}{297.65}{715.077}{\Character{47}}%
\TextAt{14}{303.503}{715.077}{o}%
\TextAt{3}{309.347}{719.415}{\Character{3}}%
\TextAt{27}{309.347}{712.121}{p}%
\TextAt{8}{314.08}{715.077}{.}%
\TextAt{8}{324.092}{715.077}{If}%
\TextAt{13}{336.306}{715.077}{a}%
\NativeMathAt{346.66}{715.077}{11.955}{\in}%
\TextAt{14}{358.84}{715.077}{o}%
\TextAt{9}{364.684}{715.077}{\Character{40}}%
\TextAt{13}{369.236}{715.077}{X}%
\TextAt{9}{379.892}{715.077}{\Character{41}}%
\TextAt{3}{384.444}{719.415}{\Character{3}}%
\TextAt{8}{393.595}{715.077}{belongs}%
\TextAt{8}{436.127}{715.077}{to}%
\TextAt{8}{450.959}{715.077}{the}%
\TextAt{8}{471.637}{715.077}{kernel,}%
\TextAt{8}{510.007}{715.077}{then}%
\TextAt{13}{79.226}{698.031}{a}%
\NativeMathAt{88.692}{698.031}{11.955}{\in}%
\TextAt{14}{99.983}{698.031}{o}%
\TextAt{3}{105.827}{702.369}{\Character{3}}%
\TextAt{27}{105.827}{695.075}{p}%
\TextAt{8}{114.367}{698.031}{for}%
\TextAt{14}{132.162}{698.031}{p}%
\NativeMathAt{141.466}{698.031}{11.955}{\notin}%
\TextAt{13}{152.757}{698.031}{X}%
\TextAt{8}{163.412}{698.031}{,}%
\TextAt{8}{170.49}{698.031}{and}%
\TextAt{8}{193.163}{698.031}{also}%
\TextAt{8}{216.553}{698.031}{for}%
\TextAt{14}{234.348}{698.031}{p}%
\NativeMathAt{243.652}{698.031}{11.955}{\in}%
\TextAt{13}{254.943}{698.031}{X}%
\TextAt{8}{269.406}{698.031}{because}%
\TextAt{14}{312.618}{698.031}{o}%
\TextAt{27}{318.462}{696.237}{p}%
\TextAt{9}{326.27}{698.031}{\Character{61}}%
\TextAt{14}{338.695}{698.031}{o}%
\TextAt{9}{344.54}{698.031}{\Character{40}}%
\TextAt{13}{349.092}{698.031}{X}%
\TextAt{9}{359.747}{698.031}{\Character{41}}%
\TextAt{27}{364.299}{696.237}{p}%
\TextAt{22}{368.288}{694.898}{X}%
\TextAt{8}{375.732}{698.031}{,}%
\TextAt{8}{382.81}{698.031}{hence}%
\TextAt{13}{415.226}{698.031}{a}%
\NativeMathAt{424.692}{698.031}{11.955}{\in}%
\TextAt{20}{435.983}{706.997}{\Character{84}}%
\TextAt{27}{445.946}{694.544}{p}%
\TextAt{14}{452.425}{698.031}{o}%
\TextAt{3}{458.269}{702.369}{\Character{3}}%
\TextAt{27}{458.269}{695.075}{p}%
\TextAt{9}{466.322}{698.031}{\Character{61}}%
\TextAt{14}{478.748}{698.031}{o}%
\TextAt{3}{484.592}{702.369}{\Character{3}}%
\TextAt{8}{493.132}{698.031}{(see}%
\TextAt{8}{516.511}{698.031}{the}%
\TextAt{8}{79.226}{680.985}{argument}%
\TextAt{8}{131.59}{680.985}{in}%
\TextAt{8}{145.243}{680.985}{the}%
\TextAt{8}{165.399}{680.985}{proof}%
\TextAt{8}{195.968}{680.985}{of}%
\TextAt{8}{209.298}{680.985}{(11.5)).}%
\TextAt{8}{252.237}{680.985}{This}%
\TextAt{8}{278.957}{680.985}{shows}%
\TextAt{8}{312.575}{680.985}{the}%
\TextAt{8}{332.732}{680.985}{exactness}%
\TextAt{8}{384.557}{680.985}{at}%
\TextAt{14}{398.868}{680.985}{o}%
\TextAt{9}{404.712}{680.985}{\Character{40}}%
\TextAt{13}{409.264}{680.985}{X}%
\TextAt{9}{419.919}{680.985}{\Character{41}}%
\TextAt{3}{424.472}{685.323}{\Character{3}}%
\TextAt{8}{429.204}{680.985}{.}%
\TextAt{8}{437.652}{680.985}{The}%
\TextAt{8}{461.706}{680.985}{arrow}%
\TextAt{20}{253.061}{663.341}{\Character{77}}%
\TextAt{27}{253.484}{637.438}{p}%
\TextAt{17}{258.413}{637.438}{\Character{47}}%
\NativeMathAt{257.472}{637.438}{7.97}{\in}%
\TextAt{21}{263.118}{637.438}{X}%
\TextAt{13}{273.12}{651.983}{K}%
\TextAt{3}{283.937}{656.92}{\Character{3}}%
\TextAt{9}{288.67}{651.983}{\Character{47}}%
\TextAt{14}{294.523}{651.983}{o}%
\TextAt{3}{300.367}{656.92}{\Character{3}}%
\TextAt{27}{300.367}{649.028}{p}%
\NativeArrowAt{308.42}{651.983}{11.955}{19.26}{\longrightarrow}%
\TextAt{13}{331.002}{651.983}{C}%
\TextAt{13}{340.22}{651.983}{l}%
\TextAt{9}{343.985}{651.983}{\Character{40}}%
\TextAt{14}{348.538}{651.983}{o}%
\TextAt{9}{354.382}{651.983}{\Character{41}}%
\TextAt{8}{79.226}{613.487}{is}%
\TextAt{8}{90.99}{613.487}{induced}%
\TextAt{8}{134.555}{613.487}{by}%
\TextAt{8}{150.814}{613.487}{mapping}%
\TextAt{20}{230.753}{607.799}{\Character{77}}%
\TextAt{27}{231.175}{581.896}{p}%
\TextAt{17}{236.105}{581.896}{\Character{47}}%
\NativeMathAt{235.164}{581.896}{7.97}{\in}%
\TextAt{21}{240.809}{581.896}{X}%
\TextAt{13}{250.811}{596.441}{\Character{11}}%
\TextAt{27}{258.257}{594.648}{p}%
\TextAt{8}{266.064}{596.441}{mod}%
\TextAt{14}{291.502}{596.441}{o}%
\TextAt{3}{297.347}{601.378}{\Character{3}}%
\TextAt{27}{297.347}{593.486}{p}%
\NativeArrowAt{305.4}{596.441}{11.955}{19.26}{\longmapsto}%
\TextAt{20}{328.954}{607.799}{\Character{89}}%
\TextAt{27}{327.982}{581.896}{p}%
\TextAt{17}{332.911}{581.896}{\Character{47}}%
\NativeMathAt{331.97}{581.896}{7.97}{\in}%
\TextAt{21}{337.616}{581.896}{X}%
\TextAt{14}{347.195}{596.441}{p}%
\TextAt{21}{353.179}{601.378}{v}%
\TextAt{28}{357.302}{600.381}{p}%
\TextAt{17}{361.385}{601.378}{\Character{40}}%
\TextAt{21}{364.678}{601.378}{\Character{11}}%
\TextAt{28}{370.117}{600.381}{p}%
\TextAt{17}{374.2}{601.378}{\Character{41}}%
\TextAt{13}{377.991}{596.441}{\Character{59}}%
\TextAt{8}{79.226}{562.755}{where}%
\TextAt{13}{111.784}{562.755}{v}%
\TextAt{27}{117.444}{560.962}{p}%
\TextAt{9}{125.251}{562.755}{\Character{58}}%
\TextAt{13}{131.824}{562.755}{K}%
\TextAt{3}{142.641}{567.093}{\Character{3}}%
\NativeMathAt{150.695}{562.755}{11.955}{\rightarrow}%
\TextAt{18}{165.971}{562.755}{Z}%
\TextAt{8}{176.587}{562.755}{is}%
\TextAt{8}{187.1}{562.755}{the}%
\TextAt{8}{206.005}{562.755}{exponential}%
\TextAt{8}{267.518}{562.755}{valuation}%
\TextAt{8}{317.34}{562.755}{of}%
\TextAt{13}{329.419}{562.755}{K}%
\TextAt{8}{342.883}{562.755}{associated}%
\TextAt{8}{397.367}{562.755}{to}%
\TextAt{14}{410.426}{562.755}{o}%
\TextAt{27}{416.27}{560.962}{p}%
\TextAt{8}{420.757}{562.755}{.}%
\TextAt{8}{428.788}{562.755}{Let}%
\TextAt{20}{448.506}{571.722}{\Character{76}}%
\TextAt{27}{461.79}{559.268}{p}%
\TextAt{17}{466.72}{559.268}{\Character{47}}%
\NativeMathAt{465.779}{559.268}{7.97}{\in}%
\TextAt{21}{471.424}{559.268}{X}%
\TextAt{13}{481.502}{562.755}{\Character{11}}%
\TextAt{27}{488.947}{560.962}{p}%
\TextAt{8}{496.755}{562.755}{mod}%
\TextAt{14}{522.193}{562.755}{o}%
\TextAt{3}{528.037}{567.093}{\Character{3}}%
\TextAt{27}{528.037}{559.8}{p}%
\TextAt{8}{79.226}{545.709}{be}%
\TextAt{8}{95.15}{545.709}{an}%
\TextAt{8}{111.409}{545.709}{element}%
\TextAt{8}{154.651}{545.709}{in}%
\TextAt{8}{168.304}{545.709}{the}%
\TextAt{8}{188.461}{545.709}{kernel,}%
\TextAt{8}{226.179}{545.709}{i.e.,}%
\TextAt{20}{235.916}{528.065}{\Character{89}}%
\TextAt{27}{234.943}{502.162}{p}%
\TextAt{17}{239.873}{502.162}{\Character{47}}%
\NativeMathAt{238.932}{502.162}{7.97}{\in}%
\TextAt{21}{244.577}{502.162}{X}%
\TextAt{14}{254.157}{516.708}{p}%
\TextAt{21}{260.14}{521.644}{v}%
\TextAt{28}{264.263}{520.648}{p}%
\TextAt{17}{268.346}{521.644}{\Character{40}}%
\TextAt{21}{271.639}{521.644}{\Character{11}}%
\TextAt{28}{277.078}{520.648}{p}%
\TextAt{17}{281.161}{521.644}{\Character{41}}%
\TextAt{9}{288.273}{516.708}{\Character{61}}%
\TextAt{9}{300.683}{516.708}{\Character{40}}%
\TextAt{13}{305.251}{516.708}{\Character{11}}%
\TextAt{9}{312.773}{516.708}{\Character{41}}%
\TextAt{9}{320.628}{516.708}{\Character{61}}%
\TextAt{20}{333.071}{528.065}{\Character{89}}%
\TextAt{27}{338.715}{502.959}{p}%
\TextAt{14}{350.34}{516.708}{p}%
\TextAt{21}{356.323}{521.644}{v}%
\TextAt{28}{360.446}{520.648}{p}%
\TextAt{17}{364.529}{521.644}{\Character{40}}%
\TextAt{21}{367.822}{521.644}{\Character{11}}%
\TextAt{17}{373.261}{521.644}{\Character{41}}%
\TextAt{8}{79.226}{477.542}{for}%
\TextAt{8}{97.557}{477.542}{some}%
\TextAt{13}{127.329}{477.542}{\Character{11}}%
\NativeMathAt{138.932}{477.542}{11.955}{\in}%
\TextAt{13}{150.984}{477.542}{K}%
\TextAt{3}{161.801}{481.881}{\Character{3}}%
\TextAt{8}{166.534}{477.542}{.}%
\TextAt{8}{176.32}{477.542}{Because}%
\TextAt{8}{221.527}{477.542}{of}%
\TextAt{8}{235.303}{477.542}{unique}%
\TextAt{8}{273.79}{477.542}{prime}%
\TextAt{8}{307.4}{477.542}{factorization,}%
\TextAt{8}{379.187}{477.542}{this}%
\TextAt{8}{402.455}{477.542}{means}%
\TextAt{8}{438.731}{477.542}{that}%
\TextAt{13}{464.546}{477.542}{v}%
\TextAt{27}{470.205}{475.749}{p}%
\TextAt{9}{474.692}{477.542}{\Character{40}}%
\TextAt{13}{479.244}{477.542}{\Character{11}}%
\TextAt{9}{486.766}{477.542}{\Character{41}}%
\TextAt{9}{495.386}{477.542}{\Character{61}}%
\TextAt{9}{508.56}{477.542}{0}%
\TextAt{8}{518.782}{477.542}{for}%
\TextAt{14}{79.226}{460.496}{p}%
\NativeMathAt{89.062}{460.496}{11.955}{\in}%
\TextAt{13}{100.884}{460.496}{X}%
\TextAt{8}{111.539}{460.496}{,}%
\TextAt{8}{119.078}{460.496}{and}%
\TextAt{13}{142.153}{460.496}{v}%
\TextAt{27}{147.812}{458.703}{p}%
\TextAt{9}{152.299}{460.496}{\Character{40}}%
\TextAt{13}{156.851}{460.496}{\Character{11}}%
\TextAt{27}{164.297}{458.703}{p}%
\TextAt{9}{168.784}{460.496}{\Character{41}}%
\TextAt{9}{177.177}{460.496}{\Character{61}}%
\TextAt{13}{190.146}{460.496}{v}%
\TextAt{27}{195.805}{458.703}{p}%
\TextAt{9}{200.292}{460.496}{\Character{40}}%
\TextAt{13}{204.844}{460.496}{\Character{11}}%
\TextAt{9}{212.366}{460.496}{\Character{41}}%
\TextAt{8}{221.128}{460.496}{for}%
\TextAt{14}{239.324}{460.496}{p}%
\NativeMathAt{249.16}{460.496}{11.955}{\notin}%
\TextAt{13}{260.982}{460.496}{X}%
\TextAt{8}{271.638}{460.496}{.}%
\TextAt{8}{281.022}{460.496}{It}%
\TextAt{8}{294.006}{460.496}{follows}%
\TextAt{8}{332.754}{460.496}{that}%
\TextAt{13}{358.435}{460.496}{\Character{11}}%
\NativeMathAt{369.809}{460.496}{11.955}{\in}%
\TextAt{20}{381.631}{469.463}{\Character{84}}%
\TextAt{27}{391.594}{457.009}{p}%
\NativeMathAt{395.583}{457.009}{7.97}{\in}%
\TextAt{21}{401.228}{457.009}{X}%
\TextAt{14}{411.306}{460.496}{o}%
\TextAt{3}{417.15}{464.835}{\Character{3}}%
\TextAt{27}{417.15}{457.541}{p}%
\TextAt{9}{425.735}{460.496}{\Character{61}}%
\TextAt{14}{438.692}{460.496}{o}%
\TextAt{9}{444.536}{460.496}{\Character{40}}%
\TextAt{13}{449.088}{460.496}{X}%
\TextAt{9}{459.744}{460.496}{\Character{41}}%
\TextAt{3}{464.296}{464.835}{\Character{3}}%
\TextAt{8}{473.238}{460.496}{and}%
\TextAt{13}{496.312}{460.496}{\Character{11}}%
\NativeMathAt{507.686}{460.496}{11.955}{\equiv}%
\TextAt{13}{520.837}{460.496}{\Character{11}}%
\TextAt{27}{528.283}{458.703}{p}%
\TextAt{9}{79.226}{443.45}{\Character{40}}%
\TextAt{8}{83.778}{443.45}{mod}%
\TextAt{14}{109.881}{443.45}{o}%
\TextAt{3}{115.725}{447.789}{\Character{3}}%
\TextAt{27}{115.725}{440.495}{p}%
\TextAt{9}{120.457}{443.45}{\Character{41}}%
\TextAt{8}{125.009}{443.45}{.}%
\TextAt{8}{133.458}{443.45}{This}%
\TextAt{8}{160.178}{443.45}{shows}%
\TextAt{8}{193.795}{443.45}{exactness}%
\TextAt{8}{245.621}{443.45}{in}%
\TextAt{8}{259.274}{443.45}{the}%
\TextAt{8}{279.43}{443.45}{middle.}%
\TextAt{8}{322.345}{443.45}{The}%
\TextAt{8}{346.399}{443.45}{arrow}%
\TextAt{13}{255.234}{414.449}{C}%
\TextAt{13}{264.452}{414.449}{l}%
\TextAt{9}{268.218}{414.449}{\Character{40}}%
\TextAt{14}{272.77}{414.449}{o}%
\TextAt{9}{278.614}{414.449}{\Character{41}}%
\NativeArrowAt{286.487}{414.449}{11.955}{19.26}{\longrightarrow}%
\TextAt{13}{309.069}{414.449}{C}%
\TextAt{13}{318.286}{414.449}{l}%
\TextAt{9}{322.053}{414.449}{\Character{40}}%
\TextAt{14}{326.605}{414.449}{o}%
\TextAt{9}{332.449}{414.449}{\Character{40}}%
\TextAt{13}{337.002}{414.449}{X}%
\TextAt{9}{347.657}{414.449}{\Character{41}\Character{41}}%
\TextAt{8}{79.226}{385.448}{comes}%
\TextAt{8}{113.464}{385.448}{from}%
\TextAt{8}{140.816}{385.448}{mapping}%
\TextAt{14}{188.659}{385.448}{a}%
\NativeArrowAt{197.963}{385.448}{11.955}{11.955}{\mapsto}%
\TextAt{14}{213.239}{385.448}{ao}%
\TextAt{9}{225.066}{385.448}{\Character{40}}%
\TextAt{13}{229.618}{385.448}{X}%
\TextAt{9}{240.273}{385.448}{\Character{41}}%
\TextAt{8}{244.826}{385.448}{.}%
\TextAt{8}{253.178}{385.448}{The}%
\TextAt{8}{276.943}{385.448}{classes}%
\TextAt{8}{313.907}{385.448}{of}%
\TextAt{8}{326.948}{385.448}{prime}%
\TextAt{8}{359.823}{385.448}{ideals}%
\TextAt{14}{392.113}{385.448}{p}%
\NativeMathAt{401.416}{385.448}{11.955}{\in}%
\TextAt{13}{412.707}{385.448}{X}%
\TextAt{8}{426.971}{385.448}{are}%
\TextAt{8}{446.194}{385.448}{mapped}%
\TextAt{8}{490.45}{385.448}{onto}%
\TextAt{8}{516.511}{385.448}{the}%
\TextAt{8}{79.226}{368.402}{classes}%
\TextAt{8}{116.124}{368.402}{of}%
\TextAt{8}{129.099}{368.402}{prime}%
\TextAt{8}{161.907}{368.402}{ideals}%
\TextAt{8}{194.13}{368.402}{of}%
\TextAt{14}{207.105}{368.402}{o}%
\TextAt{9}{212.95}{368.402}{\Character{40}}%
\TextAt{13}{217.502}{368.402}{X}%
\TextAt{9}{228.157}{368.402}{\Character{41}}%
\TextAt{8}{232.709}{368.402}{.}%
\TextAt{8}{241.039}{368.402}{Since}%
\TextAt{13}{271.242}{368.402}{C}%
\TextAt{13}{280.459}{368.402}{l}%
\TextAt{9}{284.225}{368.402}{\Character{40}}%
\TextAt{14}{288.778}{368.402}{o}%
\TextAt{9}{294.622}{368.402}{\Character{40}}%
\TextAt{13}{299.174}{368.402}{X}%
\TextAt{9}{309.829}{368.402}{\Character{41}\Character{41}}%
\TextAt{8}{322.476}{368.402}{is}%
\TextAt{8}{333.885}{368.402}{generated}%
\TextAt{8}{386.863}{368.402}{by}%
\TextAt{8}{402.767}{368.402}{these}%
\TextAt{8}{432.383}{368.402}{classes,}%
\TextAt{8}{472.604}{368.402}{the}%
\TextAt{8}{492.405}{368.402}{arrow}%
\TextAt{8}{524.903}{368.402}{is}%
\TextAt{8}{79.226}{351.356}{surjective.}%
\TextAt{8}{136.003}{351.356}{For}%
\TextAt{14}{156.416}{351.356}{p}%
\NativeMathAt{165.72}{351.356}{11.955}{\notin}%
\TextAt{13}{177.011}{351.356}{X}%
\TextAt{8}{191.008}{351.356}{we}%
\TextAt{8}{207.679}{351.356}{have}%
\TextAt{14}{234.117}{351.356}{po}%
\TextAt{9}{245.945}{351.356}{\Character{40}}%
\TextAt{13}{250.497}{351.356}{X}%
\TextAt{9}{261.152}{351.356}{\Character{41}}%
\TextAt{9}{269.007}{351.356}{\Character{61}}%
\TextAt{9}{281.428}{351.356}{\Character{40}1\Character{41}}%
\TextAt{8}{296.408}{351.356}{,}%
\TextAt{8}{303.113}{351.356}{and}%
\TextAt{8}{325.319}{351.356}{this}%
\TextAt{8}{347.585}{351.356}{means}%
\TextAt{8}{382.859}{351.356}{that}%
\TextAt{8}{407.672}{351.356}{the}%
\TextAt{8}{427.272}{351.356}{kernel}%
\TextAt{8}{461.182}{351.356}{consists}%
\TextAt{8}{503.737}{351.356}{of}%
\TextAt{8}{516.51}{351.356}{the}%
\TextAt{8}{79.226}{334.31}{classes}%
\TextAt{8}{116.224}{334.31}{of}%
\TextAt{8}{129.299}{334.31}{the}%
\TextAt{8}{149.2}{334.31}{ideals}%
\TextAt{20}{181.523}{343.277}{\Character{81}}%
\TextAt{27}{192.814}{330.823}{p}%
\TextAt{17}{197.744}{330.823}{\Character{47}}%
\NativeMathAt{196.803}{330.823}{7.97}{\in}%
\TextAt{21}{202.448}{330.823}{X}%
\TextAt{14}{212.526}{334.31}{p}%
\TextAt{21}{218.509}{338.649}{\Character{23}}%
\TextAt{28}{222.672}{337.652}{p}%
\TextAt{8}{227.253}{334.31}{.}%
\TextAt{8}{235.616}{334.31}{This,}%
\TextAt{8}{265.384}{334.31}{however,}%
\TextAt{8}{313.311}{334.31}{is}%
\TextAt{8}{324.82}{334.31}{visibly}%
\TextAt{8}{361.698}{334.31}{the}%
\TextAt{8}{381.6}{334.31}{image}%
\TextAt{8}{415.166}{334.31}{of}%
\TextAt{8}{428.241}{334.31}{the}%
\TextAt{8}{448.143}{334.31}{preceding}%
\TextAt{8}{500.562}{334.31}{arrow.}%
\TextAt{8}{79.226}{317.264}{Therefore}%
\TextAt{8}{132.558}{317.264}{the}%
\TextAt{8}{153.049}{317.264}{whole}%
\TextAt{8}{186.547}{317.264}{sequence}%
\TextAt{8}{235.383}{317.264}{is}%
\TextAt{8}{247.481}{317.264}{exact.}%
\TextAt{8}{283.927}{317.264}{Finally,}%
\TextAt{8}{326.451}{317.264}{the}%
\TextAt{8}{346.942}{317.264}{valuation}%
\TextAt{13}{398.349}{317.264}{v}%
\TextAt{27}{404.008}{315.471}{p}%
\TextAt{9}{412.386}{317.264}{\Character{58}}%
\TextAt{13}{419.528}{317.264}{K}%
\TextAt{3}{430.346}{321.603}{\Character{3}}%
\NativeMathAt{438.969}{317.264}{11.955}{\rightarrow}%
\TextAt{18}{454.815}{317.264}{Z}%
\TextAt{8}{467.016}{317.264}{produces}%
\TextAt{8}{516.51}{317.264}{the}%
\TextAt{8}{79.226}{300.218}{isomorphism}%
\TextAt{13}{147.646}{300.218}{K}%
\TextAt{3}{158.463}{304.557}{\Character{3}}%
\TextAt{9}{163.196}{300.218}{\Character{47}}%
\TextAt{14}{169.049}{300.218}{o}%
\TextAt{3}{174.893}{304.557}{\Character{3}}%
\TextAt{27}{174.893}{297.263}{p}%
\NativeMathAt{182.946}{300.218}{11.955}{\cong}%
\TextAt{18}{195.565}{300.218}{Z}%
\TextAt{8}{203.535}{300.218}{.}%
\Vector{q 0 G 0 g 0.399 w 452.334 -473.782 m 452.334 -465.712 l S Q}%
\Vector{q 0 G 0 g 0.399 w 452.534 -465.911 m 459.707 -465.911 l S Q}%
\Vector{q 0 G 0 g 0.399 w 452.534 -473.583 m 459.707 -473.583 l S Q}%
\Vector{q 0 G 0 g 0.399 w 459.906 -473.782 m 459.906 -465.712 l S Q}%
\TextAt{8}{97.159}{271.217}{For}%
\TextAt{8}{119.545}{271.217}{the}%
\TextAt{8}{141.119}{271.217}{ring}%
\TextAt{8}{166.602}{271.217}{of}%
\TextAt{8}{181.349}{271.217}{integers}%
\TextAt{14}{226.08}{271.217}{o}%
\TextAt{21}{231.924}{269.424}{K}%
\TextAt{8}{245.441}{271.217}{of}%
\TextAt{8}{260.188}{271.217}{an}%
\TextAt{8}{277.865}{271.217}{algebraic}%
\TextAt{8}{328.716}{271.217}{number}%
\TextAt{8}{372.73}{271.217}{field}%
\TextAt{13}{399.504}{271.217}{K}%
\TextAt{8}{410.322}{271.217}{,}%
\TextAt{8}{419.242}{271.217}{the}%
\TextAt{8}{440.816}{271.217}{proposition}%
\TextAt{8}{503.766}{271.217}{yields}%
\TextAt{8}{79.226}{254.171}{the}%
\TextAt{8}{100.479}{254.171}{following}%
\TextAt{8}{151.011}{254.171}{results.}%
\TextAt{8}{196.044}{254.171}{Let}%
\TextAt{13}{218.11}{254.171}{S}%
\TextAt{8}{231}{254.171}{denote}%
\TextAt{8}{269.815}{254.171}{a}%
\TextAt{8}{280.667}{254.171}{finite}%
\TextAt{8}{311.676}{254.171}{set}%
\TextAt{8}{331.04}{254.171}{of}%
\TextAt{8}{345.467}{254.171}{prime}%
\TextAt{8}{379.727}{254.171}{ideals}%
\TextAt{8}{413.401}{254.171}{of}%
\TextAt{14}{427.828}{254.171}{o}%
\TextAt{21}{433.672}{252.378}{K}%
\TextAt{8}{446.869}{254.171}{(not}%
\TextAt{8}{473.335}{254.171}{any}%
\TextAt{8}{496.548}{254.171}{more}%
\TextAt{8}{526.911}{254.171}{a}%
\TextAt{8}{79.226}{237.125}{multiplicative}%
\TextAt{8}{153.445}{237.125}{subset),}%
\TextAt{8}{197.636}{237.125}{and}%
\TextAt{8}{220.795}{237.125}{let}%
\TextAt{13}{238.096}{237.125}{X}%
\TextAt{8}{253.044}{237.125}{be}%
\TextAt{8}{269.364}{237.125}{the}%
\TextAt{8}{289.917}{237.125}{set}%
\TextAt{8}{308.58}{237.125}{of}%
\TextAt{8}{322.306}{237.125}{all}%
\TextAt{8}{338.961}{237.125}{prime}%
\TextAt{8}{372.521}{237.125}{ideals}%
\TextAt{8}{405.495}{237.125}{that}%
\TextAt{8}{431.259}{237.125}{do}%
\TextAt{8}{447.914}{237.125}{not}%
\TextAt{8}{469.124}{237.125}{belong}%
\TextAt{8}{506.916}{237.125}{to}%
\TextAt{13}{521.622}{237.125}{S}%
\TextAt{8}{529.518}{237.125}{.}%
\TextAt{8}{79.226}{220.079}{We}%
\TextAt{8}{99.371}{220.079}{put}%
\TextAt{14}{272.572}{203.033}{o}%
\TextAt{21}{278.416}{207.969}{S}%
\TextAt{21}{278.416}{200.078}{K}%
\TextAt{9}{289.94}{203.033}{\Character{61}}%
\TextAt{14}{302.365}{203.033}{o}%
\TextAt{21}{308.209}{201.24}{K}%
\TextAt{9}{316.412}{203.033}{\Character{40}}%
\TextAt{13}{320.964}{203.033}{X}%
\TextAt{9}{331.62}{203.033}{\Character{41}}%
\TextAt{13}{336.172}{203.033}{\Character{58}}%
\TextAt{8}{79.226}{179.512}{The}%
\TextAt{8}{103.628}{179.512}{units}%
\TextAt{8}{133.302}{179.512}{of}%
\TextAt{8}{146.981}{179.512}{this}%
\TextAt{8}{170.151}{179.512}{ring}%
\TextAt{8}{194.565}{179.512}{are}%
\TextAt{8}{214.425}{179.512}{called}%
\TextAt{8}{247.936}{179.512}{the}%
\TextAt{13}{268.441}{179.512}{S}%
\TextAt{12}{276.336}{179.512}{-units}%
\TextAt{8}{310.05}{179.512}{,}%
\TextAt{8}{317.635}{179.512}{and}%
\TextAt{8}{340.745}{179.512}{the}%
\TextAt{8}{361.25}{179.512}{group}%
\TextAt{13}{394.774}{179.512}{C}%
\TextAt{13}{403.992}{179.512}{l}%
\TextAt{21}{407.757}{183.85}{S}%
\TextAt{21}{407.494}{176.491}{K}%
\TextAt{9}{419.611}{179.512}{\Character{61}}%
\TextAt{13}{432.63}{179.511}{C}%
\TextAt{13}{441.848}{179.511}{l}%
\TextAt{9}{445.614}{179.512}{\Character{40}}%
\TextAt{14}{450.166}{179.511}{o}%
\TextAt{21}{456.01}{183.85}{S}%
\TextAt{21}{456.01}{176.491}{K}%
\TextAt{9}{464.213}{179.512}{\Character{41}}%
\TextAt{8}{473.011}{179.512}{the}%
\TextAt{13}{493.516}{179.512}{S}%
\TextAt{12}{501.411}{179.512}{-class}%
\TextAt{12}{79.226}{162.466}{group}%
\TextAt{8}{117.028}{162.466}{of}%
\TextAt{13}{130.358}{162.466}{K}%
\TextAt{8}{141.176}{162.466}{.}%
\TextAt{8}{300.14}{43.033}{66}%
\Destination{78.23}{809.815}{nav.page.66}%
\Destination{79.36}{400.455}{nav.print.71}%
\Hotlink{208.798}{678.755}{30.936}{11.047}{nav.number.1.11.5}%
}

% Page 79 (body page 67)
\Page{%
\TextAt{2}{79.226}{796.943}{I.}%
\TextAt{2}{89.83}{796.943}{Algebraic}%
\TextAt{2}{138.648}{796.943}{Integers}%
\Vector{q 0 G 0 g 0.299 w 7.226 18.729 m 460.769 18.729 l S Q}%
\TextAt{12}{79.226}{761.124}{Corollary}%
\TextAt{12}{139.014}{761.124}{(11.7).}%
\TextAt{8}{181.06}{761.124}{For}%
\TextAt{8}{202.03}{761.124}{the}%
\TextAt{8}{222.186}{761.124}{group}%
\TextAt{13}{255.362}{761.124}{K}%
\TextAt{21}{266.179}{765.462}{S}%
\TextAt{9}{275.594}{761.124}{\Character{61}}%
\TextAt{9}{288.004}{761.124}{\Character{40}}%
\TextAt{14}{292.572}{761.124}{o}%
\TextAt{21}{298.416}{765.462}{S}%
\TextAt{21}{298.416}{758.103}{K}%
\TextAt{9}{306.619}{761.124}{\Character{41}}%
\TextAt{3}{311.171}{765.462}{\Character{3}}%
\TextAt{8}{319.801}{761.124}{of}%
\TextAt{13}{333.131}{761.124}{S}%
\TextAt{8}{341.026}{761.124}{-units}%
\TextAt{8}{374.25}{761.124}{of}%
\TextAt{13}{387.58}{761.124}{K}%
\TextAt{8}{402.295}{761.124}{there}%
\TextAt{8}{432.206}{761.124}{is}%
\TextAt{8}{443.97}{761.124}{an}%
\TextAt{8}{460.229}{761.124}{isomorphism}%
\TextAt{13}{240.616}{732.706}{K}%
\TextAt{21}{251.433}{737.642}{S}%
\NativeMathAt{260.848}{732.706}{11.955}{\cong}%
\TextAt{13}{273.467}{732.706}{\Character{22}}%
\TextAt{9}{280.51}{732.706}{\Character{40}}%
\TextAt{13}{285.063}{732.706}{K}%
\TextAt{9}{295.88}{732.706}{\Character{41}}%
\TextAt{16}{303.089}{732.706}{\Character{2}}%
\TextAt{18}{315.045}{732.706}{Z}%
\TextAt{7}{323.015}{737.642}{\char35\relax }%
\TextAt{21}{330.068}{737.642}{S}%
\TextAt{17}{335.664}{737.642}{\Character{43}}%
\TextAt{21}{342.25}{737.642}{r}%
\TextAt{17}{346.307}{737.642}{\Character{43}}%
\TextAt{21}{352.894}{737.642}{s}%
\TextAt{3}{356.809}{737.642}{\Character{0}}%
\TextAt{17}{363.396}{737.642}{1}%
\TextAt{13}{368.128}{732.706}{\Character{59}}%
\TextAt{8}{79.226}{704.289}{where}%
\TextAt{13}{113.035}{704.289}{r}%
\TextAt{8}{122.533}{704.289}{and}%
\TextAt{13}{145.296}{704.289}{s}%
\TextAt{8}{154.707}{704.289}{are}%
\TextAt{8}{174.218}{704.289}{defined}%
\TextAt{8}{214.531}{704.289}{as}%
\TextAt{8}{228.901}{704.289}{in}%
\TextAt{8}{242.554}{704.289}{§5,}%
\TextAt{8}{261.24}{704.289}{p.}%
\TextAt{8}{274.893}{704.289}{30.}%
\TextAt{19}{79.226}{669.644}{Proof.}%
\TextAt{8}{116}{669.644}{The}%
\TextAt{8}{139.691}{669.644}{torsion}%
\TextAt{8}{178.421}{669.644}{subgroup}%
\TextAt{8}{228.855}{669.644}{of}%
\TextAt{13}{241.821}{669.644}{K}%
\TextAt{21}{252.639}{673.982}{S}%
\TextAt{8}{262.267}{669.644}{is}%
\TextAt{8}{273.668}{669.644}{the}%
\TextAt{8}{293.461}{669.644}{group}%
\TextAt{13}{326.273}{669.644}{\Character{22}}%
\TextAt{9}{333.316}{669.644}{\Character{40}}%
\TextAt{13}{337.868}{669.644}{K}%
\TextAt{9}{348.686}{669.644}{\Character{41}}%
\TextAt{8}{356.772}{669.644}{of}%
\TextAt{8}{369.739}{669.644}{roots}%
\TextAt{8}{399.036}{669.644}{of}%
\TextAt{8}{412.003}{669.644}{unity}%
\TextAt{8}{442.209}{669.644}{in}%
\TextAt{13}{455.498}{669.644}{K}%
\TextAt{8}{466.316}{669.644}{.}%
\TextAt{8}{474.643}{669.644}{Since}%
\TextAt{13}{504.837}{669.644}{C}%
\TextAt{13}{514.054}{669.644}{l}%
\TextAt{9}{517.821}{669.644}{\Character{40}}%
\TextAt{14}{522.373}{669.644}{o}%
\TextAt{9}{528.217}{669.644}{\Character{41}}%
\TextAt{8}{79.226}{652.598}{is}%
\TextAt{8}{90.711}{652.598}{finite,}%
\TextAt{8}{123.652}{652.598}{we}%
\TextAt{8}{140.6}{652.598}{obtain}%
\TextAt{8}{176.748}{652.598}{the}%
\TextAt{8}{196.626}{652.598}{following}%
\TextAt{8}{245.782}{652.598}{identities}%
\TextAt{8}{295.965}{652.598}{from}%
\TextAt{8}{323.327}{652.598}{the}%
\TextAt{8}{343.204}{652.598}{exact}%
\TextAt{8}{373.818}{652.598}{sequence}%
\TextAt{8}{422.041}{652.598}{(11.6)}%
\TextAt{8}{455.595}{652.598}{and}%
\TextAt{8}{478.079}{652.598}{from}%
\TextAt{8}{505.44}{652.598}{(7.4):}%
\TextAt{8}{158.679}{616.307}{rank}%
\TextAt{9}{181.776}{616.307}{\Character{40}}%
\TextAt{13}{186.329}{616.307}{K}%
\TextAt{21}{197.146}{621.244}{S}%
\TextAt{9}{203.24}{616.307}{\Character{41}}%
\TextAt{9}{211.095}{616.307}{\Character{61}}%
\TextAt{8}{223.539}{616.307}{rank}%
\TextAt{9}{246.636}{616.307}{\Character{40}}%
\TextAt{14}{251.188}{616.307}{o}%
\TextAt{3}{257.033}{621.244}{\Character{3}}%
\TextAt{21}{257.033}{613.352}{K}%
\TextAt{9}{265.235}{616.307}{\Character{41}}%
\TextAt{9}{272.432}{616.307}{\Character{43}}%
\TextAt{8}{284.206}{616.307}{rank}%
\TextAt{20}{309.296}{636.751}{\Character{32}}%
\TextAt{20}{318.76}{627.665}{\Character{77}}%
\TextAt{27}{320.178}{602.293}{p}%
\NativeMathAt{324.167}{602.293}{7.97}{\in}%
\TextAt{21}{329.812}{602.293}{S}%
\TextAt{18}{338.818}{616.307}{Z}%
\TextAt{20}{346.788}{636.751}{\Character{33}}%
\TextAt{9}{359.574}{616.307}{\Character{61}}%
\TextAt{8}{371.999}{616.307}{\char35\relax }%
\TextAt{13}{381.755}{616.307}{S}%
\TextAt{9}{392.307}{616.307}{\Character{43}}%
\TextAt{13}{404.068}{616.307}{r}%
\TextAt{9}{412.325}{616.307}{\Character{43}}%
\TextAt{13}{424.086}{616.307}{s}%
\TextAt{16}{432.257}{616.307}{\Character{0}}%
\TextAt{9}{444.212}{616.307}{1}%
\TextAt{13}{450.065}{616.307}{\Character{58}}%
\TextAt{8}{79.226}{578.926}{This}%
\TextAt{8}{105.946}{578.926}{proves}%
\TextAt{8}{142.11}{578.926}{the}%
\TextAt{8}{162.267}{578.926}{corollary.}%
\Vector{q 0 G 0 g 0.399 w 452.334 -195.074 m 452.334 -187.004 l S Q}%
\Vector{q 0 G 0 g 0.399 w 452.534 -187.203 m 459.707 -187.203 l S Q}%
\Vector{q 0 G 0 g 0.399 w 452.534 -194.874 m 459.707 -194.874 l S Q}%
\Vector{q 0 G 0 g 0.399 w 459.906 -195.074 m 459.906 -187.004 l S Q}%
\TextAt{12}{79.226}{544.281}{Corollary}%
\TextAt{12}{139.014}{544.281}{(11.8).}%
\TextAt{8}{181.06}{544.281}{The}%
\TextAt{13}{205.114}{544.281}{S}%
\TextAt{8}{213.009}{544.281}{-class}%
\TextAt{8}{244.344}{544.281}{group}%
\TextAt{13}{277.52}{544.281}{C}%
\TextAt{13}{286.738}{544.281}{l}%
\TextAt{21}{290.503}{548.62}{S}%
\TextAt{21}{290.24}{541.261}{K}%
\TextAt{9}{301.763}{544.281}{\Character{61}}%
\TextAt{13}{314.189}{544.281}{C}%
\TextAt{13}{323.406}{544.281}{l}%
\TextAt{9}{327.172}{544.281}{\Character{40}}%
\TextAt{14}{331.725}{544.281}{o}%
\TextAt{21}{337.569}{548.62}{S}%
\TextAt{21}{337.569}{541.261}{K}%
\TextAt{9}{345.772}{544.281}{\Character{41}}%
\TextAt{8}{354.221}{544.281}{is}%
\TextAt{8}{365.985}{544.281}{finite.}%
\TextAt{12}{79.226}{512.625}{Exercise}%
\TextAt{12}{132.799}{512.625}{I.11.1.}%
\TextAt{8}{173.447}{512.625}{Let}%
\TextAt{13}{194.739}{512.625}{A}%
\TextAt{8}{207.736}{512.625}{be}%
\TextAt{8}{223.983}{512.625}{an}%
\TextAt{8}{240.566}{512.625}{arbitrary}%
\TextAt{8}{290.658}{512.625}{ring,}%
\TextAt{8}{318.38}{512.625}{not}%
\TextAt{8}{339.517}{512.625}{necessarily}%
\TextAt{8}{398.17}{512.625}{an}%
\TextAt{8}{414.752}{512.625}{integral}%
\TextAt{8}{457.684}{512.625}{domain,}%
\TextAt{8}{502.968}{512.625}{let}%
\TextAt{13}{520.196}{512.625}{M}%
\TextAt{8}{79.226}{495.579}{be}%
\TextAt{8}{95.386}{495.579}{an}%
\TextAt{13}{111.881}{495.579}{A}%
\TextAt{8}{120.656}{495.579}{-module}%
\TextAt{8}{166.083}{495.579}{and}%
\TextAt{13}{189.081}{495.579}{S}%
\TextAt{8}{201.109}{495.579}{a}%
\TextAt{8}{211.101}{495.579}{multiplicatively}%
\TextAt{8}{294.592}{495.579}{closed}%
\TextAt{8}{329.354}{495.579}{subset}%
\TextAt{8}{365.48}{495.579}{of}%
\TextAt{13}{379.045}{495.579}{A}%
\TextAt{8}{391.954}{495.579}{such}%
\TextAt{8}{418.587}{495.579}{that}%
\TextAt{9}{444.191}{495.579}{0}%
\NativeMathAt{453.767}{495.579}{11.955}{\notin}%
\TextAt{13}{465.46}{495.579}{S}%
\TextAt{8}{473.355}{495.579}{.}%
\TextAt{8}{482.51}{495.579}{In}%
\TextAt{13}{497.367}{495.579}{M}%
\TextAt{16}{512.758}{495.579}{\Character{2}}%
\TextAt{13}{524.874}{495.579}{S}%
\TextAt{8}{79.226}{478.533}{consider}%
\TextAt{8}{124.811}{478.533}{the}%
\TextAt{8}{144.968}{478.533}{equivalence}%
\TextAt{8}{206.752}{478.533}{relation}%
\TextAt{9}{143.051}{450.115}{\Character{40}}%
\TextAt{13}{147.604}{450.115}{m\Character{59}}%
\TextAt{13}{163.062}{450.115}{s}%
\TextAt{9}{168.601}{450.115}{\Character{41}}%
\NativeMathAt{176.474}{450.115}{11.955}{\sim}%
\TextAt{9}{189.094}{450.115}{\Character{40}}%
\TextAt{13}{193.646}{450.115}{m}%
\TextAt{3}{203.885}{455.051}{\Character{48}}%
\TextAt{13}{206.68}{450.115}{\Character{59}}%
\TextAt{13}{211.904}{450.115}{s}%
\TextAt{3}{217.439}{455.051}{\Character{48}}%
\TextAt{9}{220.234}{450.115}{\Character{41}}%
\NativeArrowAt{240.062}{450.115}{11.955}{21.926}{\Longleftrightarrow}%
\TextAt{16}{277.255}{450.115}{\Character{57}}%
\TextAt{13}{283.898}{450.115}{s}%
\TextAt{3}{289.412}{455.051}{\Character{48}\Character{48}}%
\NativeMathAt{297.824}{450.115}{11.955}{\in}%
\TextAt{13}{309.115}{450.115}{S}%
\TextAt{8}{320.908}{450.115}{such}%
\TextAt{8}{347.305}{450.115}{that}%
\TextAt{13}{372.674}{450.115}{s}%
\TextAt{3}{378.188}{455.051}{\Character{48}\Character{48}}%
\TextAt{9}{383.28}{450.115}{\Character{40}}%
\TextAt{13}{387.832}{450.115}{s}%
\TextAt{3}{393.346}{455.051}{\Character{48}}%
\TextAt{13}{396.141}{450.115}{m}%
\TextAt{16}{409.037}{450.115}{\Character{0}}%
\TextAt{13}{420.993}{450.115}{sm}%
\TextAt{3}{436.746}{455.051}{\Character{48}}%
\TextAt{9}{439.541}{450.115}{\Character{41}}%
\TextAt{9}{447.396}{450.115}{\Character{61}}%
\TextAt{9}{459.817}{450.115}{0}%
\TextAt{13}{465.693}{450.115}{\Character{58}}%
\TextAt{8}{79.226}{421.698}{Show}%
\TextAt{8}{109.066}{421.698}{that}%
\TextAt{8}{133.382}{421.698}{the}%
\TextAt{8}{152.485}{421.698}{set}%
\TextAt{13}{169.7}{421.698}{M}%
\TextAt{21}{181.025}{419.905}{S}%
\TextAt{8}{189.963}{421.698}{of}%
\TextAt{8}{202.241}{421.698}{equivalence}%
\TextAt{8}{262.972}{421.698}{classes}%
\Vector{q 0 G 0 g 0.478 w 227.172 -341.662 m 257.274 -341.662 l S Q}%
\TextAt{9}{299.172}{421.698}{\Character{40}}%
\TextAt{13}{303.724}{421.698}{m\Character{59}}%
\TextAt{13}{319.182}{421.698}{s}%
\TextAt{9}{324.721}{421.698}{\Character{41}}%
\TextAt{8}{332.118}{421.698}{forms}%
\TextAt{8}{363.321}{421.698}{an}%
\TextAt{13}{378.527}{421.698}{A}%
\TextAt{8}{387.302}{421.698}{-module,}%
\TextAt{8}{434.903}{421.698}{and}%
\TextAt{8}{456.612}{421.698}{that}%
\TextAt{13}{480.929}{421.698}{M}%
\NativeMathAt{496.823}{421.698}{11.955}{\rightarrow}%
\TextAt{13}{512.099}{421.698}{M}%
\TextAt{21}{523.424}{419.905}{S}%
\TextAt{8}{529.518}{421.698}{,}%
\TextAt{13}{79.226}{404.652}{a}%
\NativeArrowAt{88.809}{404.652}{11.955}{11.955}{\mapsto}%
\Vector{q 0 G 0 g 0.478 w 32.202 -358.708 m 58.548 -358.708 l S Q}%
\TextAt{9}{104.202}{404.652}{\Character{40}}%
\TextAt{13}{108.754}{404.652}{a\Character{59}}%
\TextAt{9}{120.143}{404.652}{1\Character{41}}%
\TextAt{8}{130.548}{404.652}{,}%
\TextAt{8}{137.783}{404.652}{is}%
\TextAt{8}{149.616}{404.652}{a}%
\TextAt{8}{159.439}{404.652}{homomorphism.}%
\TextAt{8}{246.866}{404.652}{In}%
\TextAt{8}{261.556}{404.652}{particular,}%
\TextAt{13}{318.883}{404.652}{A}%
\TextAt{21}{327.658}{402.859}{S}%
\TextAt{8}{337.718}{404.652}{is}%
\TextAt{8}{349.55}{404.652}{a}%
\TextAt{8}{359.374}{404.652}{ring.}%
\TextAt{8}{388.196}{404.652}{It}%
\TextAt{8}{400.937}{404.652}{is}%
\TextAt{8}{412.77}{404.652}{called}%
\TextAt{8}{446.002}{404.652}{the}%
\TextAt{12}{466.227}{404.652}{localization}%
\TextAt{8}{79.226}{387.606}{of}%
\TextAt{13}{92.556}{387.606}{A}%
\TextAt{8}{105.229}{387.606}{with}%
\TextAt{8}{131.889}{387.606}{respect}%
\TextAt{8}{171.939}{387.606}{to}%
\TextAt{13}{186.249}{387.606}{S}%
\TextAt{8}{194.144}{387.606}{.}%
\TextAt{12}{79.226}{355.949}{Exercise}%
\TextAt{12}{133.08}{355.949}{I.11.2.}%
\TextAt{8}{173.972}{355.949}{Show}%
\TextAt{8}{205.432}{355.949}{that,}%
\TextAt{8}{234.763}{355.949}{in}%
\TextAt{8}{248.983}{355.949}{the}%
\TextAt{8}{269.707}{355.949}{above}%
\TextAt{8}{303.451}{355.949}{situation,}%
\TextAt{8}{356.261}{355.949}{the}%
\TextAt{8}{376.985}{355.949}{prime}%
\TextAt{8}{410.717}{355.949}{ideals}%
\TextAt{8}{443.862}{355.949}{of}%
\TextAt{13}{457.76}{355.949}{A}%
\TextAt{21}{466.535}{354.156}{S}%
\TextAt{8}{477.094}{355.949}{correspond}%
\TextAt{8}{79.226}{338.904}{1–1}%
\TextAt{8}{100.416}{338.904}{to}%
\TextAt{8}{114.445}{338.904}{the}%
\TextAt{8}{134.32}{338.904}{prime}%
\TextAt{8}{167.202}{338.904}{ideals}%
\TextAt{8}{199.498}{338.904}{of}%
\TextAt{13}{212.546}{338.904}{A}%
\TextAt{8}{224.937}{338.904}{which}%
\TextAt{8}{258.142}{338.904}{are}%
\TextAt{8}{277.372}{338.904}{disjoint}%
\TextAt{8}{318.778}{338.904}{from}%
\TextAt{13}{346.136}{338.904}{S}%
\TextAt{8}{354.032}{338.904}{.}%
\TextAt{8}{362.386}{338.904}{If}%
\TextAt{14}{373.797}{338.904}{p}%
\NativeMathAt{383.101}{338.904}{11.955}{\subseteq}%
\TextAt{13}{395.72}{338.904}{A}%
\TextAt{8}{408.111}{338.904}{and}%
\TextAt{14}{430.592}{338.904}{p}%
\TextAt{21}{436.575}{337.11}{S}%
\NativeMathAt{445.99}{338.904}{11.955}{\subseteq}%
\TextAt{13}{458.609}{338.904}{A}%
\TextAt{21}{467.385}{337.11}{S}%
\TextAt{8}{477.094}{338.904}{correspond}%
\TextAt{8}{79.226}{321.858}{in}%
\TextAt{8}{92.879}{321.858}{this}%
\TextAt{8}{115.701}{321.858}{way,}%
\TextAt{8}{141.716}{321.858}{then}%
\TextAt{13}{168.376}{321.858}{A}%
\TextAt{21}{177.151}{320.064}{S}%
\TextAt{9}{183.245}{321.858}{\Character{47}}%
\TextAt{14}{189.098}{321.858}{p}%
\TextAt{21}{195.081}{320.064}{S}%
\TextAt{8}{205.072}{321.858}{is}%
\TextAt{8}{216.836}{321.858}{the}%
\TextAt{8}{236.993}{321.858}{localization}%
\TextAt{8}{299.112}{321.858}{of}%
\TextAt{13}{312.442}{321.858}{A}%
\TextAt{9}{321.217}{321.858}{\Character{47}}%
\TextAt{14}{327.07}{321.858}{p}%
\TextAt{8}{336.951}{321.858}{with}%
\TextAt{8}{363.611}{321.858}{respect}%
\TextAt{8}{403.66}{321.858}{to}%
\TextAt{8}{417.971}{321.858}{the}%
\TextAt{8}{438.127}{321.858}{image}%
\TextAt{8}{471.948}{321.858}{of}%
\TextAt{13}{485.278}{321.858}{S}%
\TextAt{8}{493.174}{321.858}{.}%
\TextAt{12}{79.226}{290.201}{Exercise}%
\TextAt{12}{131.917}{290.201}{I.11.3.}%
\TextAt{8}{171.799}{290.201}{Let}%
\TextAt{13}{192.325}{290.201}{f}%
\TextAt{9}{202.693}{290.201}{\Character{58}}%
\TextAt{13}{209.265}{290.201}{M}%
\NativeMathAt{225.16}{290.201}{11.955}{\rightarrow}%
\TextAt{13}{240.436}{290.201}{N}%
\TextAt{8}{254.513}{290.201}{be}%
\TextAt{8}{269.994}{290.201}{a}%
\TextAt{8}{279.306}{290.201}{homomorphism}%
\TextAt{8}{361.533}{290.201}{of}%
\TextAt{13}{374.42}{290.201}{A}%
\TextAt{8}{383.196}{290.201}{-modules.}%
\TextAt{8}{437.404}{290.201}{Then}%
\TextAt{8}{467.519}{290.201}{the}%
\TextAt{8}{487.232}{290.201}{following}%
\TextAt{8}{79.226}{273.155}{conditions}%
\TextAt{8}{135.224}{273.155}{are}%
\TextAt{8}{154.735}{273.155}{equivalent:}%
\TextAt{8}{89.579}{253.12}{(i)}%
\TextAt{13}{107.918}{253.12}{f}%
\TextAt{8}{118.862}{253.12}{is}%
\TextAt{8}{130.626}{253.12}{injective}%
\TextAt{8}{177.12}{253.12}{(surjective).}%
\TextAt{8}{86.327}{236.074}{(ii)}%
\TextAt{13}{107.918}{236.074}{f}%
\TextAt{27}{113.688}{234.281}{p}%
\TextAt{9}{121.496}{236.074}{\Character{58}}%
\TextAt{13}{128.069}{236.074}{M}%
\TextAt{27}{139.393}{234.281}{p}%
\NativeMathAt{147.201}{236.074}{11.955}{\rightarrow}%
\TextAt{13}{162.477}{236.074}{N}%
\TextAt{27}{171.851}{234.281}{p}%
\TextAt{8}{180.235}{236.074}{is}%
\TextAt{8}{191.999}{236.074}{injective}%
\TextAt{8}{238.493}{236.074}{(surjective)}%
\TextAt{8}{300.014}{236.074}{for}%
\TextAt{8}{317.899}{236.074}{every}%
\TextAt{8}{348.791}{236.074}{prime}%
\TextAt{8}{381.955}{236.074}{ideal}%
\TextAt{14}{409.918}{236.074}{p}%
\TextAt{8}{415.901}{236.074}{.}%
\TextAt{8}{83.076}{219.028}{(iii)}%
\TextAt{13}{107.918}{219.028}{f}%
\TextAt{27}{113.688}{217.235}{m}%
\TextAt{9}{123.619}{219.028}{\Character{58}}%
\TextAt{13}{130.191}{219.028}{M}%
\TextAt{27}{141.516}{217.235}{m}%
\NativeMathAt{151.446}{219.028}{11.955}{\rightarrow}%
\TextAt{13}{166.722}{219.028}{N}%
\TextAt{27}{176.096}{217.235}{m}%
\TextAt{8}{186.602}{219.028}{is}%
\TextAt{8}{198.366}{219.028}{injective}%
\TextAt{8}{244.86}{219.028}{(surjective)}%
\TextAt{8}{306.381}{219.028}{for}%
\TextAt{8}{324.266}{219.028}{every}%
\TextAt{8}{355.158}{219.028}{maximal}%
\TextAt{8}{402.967}{219.028}{ideal}%
\TextAt{14}{430.93}{219.028}{m}%
\TextAt{8}{440.097}{219.028}{.}%
\TextAt{12}{79.226}{187.372}{Exercise}%
\TextAt{12}{132.075}{187.372}{I.11.4.}%
\TextAt{8}{172.093}{187.372}{Let}%
\TextAt{13}{192.757}{187.372}{S}%
\TextAt{8}{204.243}{187.372}{and}%
\TextAt{13}{226.7}{187.372}{T}%
\TextAt{8}{238.778}{187.372}{be}%
\TextAt{8}{254.396}{187.372}{two}%
\TextAt{8}{276.207}{187.372}{multiplicative}%
\TextAt{8}{349.724}{187.372}{subsets}%
\TextAt{8}{389.922}{187.372}{of}%
\TextAt{13}{402.946}{187.372}{A}%
\TextAt{8}{411.722}{187.372}{,}%
\TextAt{8}{418.626}{187.372}{and}%
\TextAt{13}{441.083}{187.372}{T}%
\TextAt{3}{449.57}{191.711}{\Character{3}}%
\TextAt{8}{457.893}{187.372}{the}%
\TextAt{8}{477.744}{187.372}{image}%
\TextAt{8}{511.259}{187.372}{of}%
\TextAt{13}{524.283}{187.372}{T}%
\TextAt{8}{79.226}{170.326}{in}%
\TextAt{13}{92.879}{170.326}{A}%
\TextAt{21}{101.654}{168.533}{S}%
\TextAt{8}{107.748}{170.326}{.}%
\TextAt{8}{116.197}{170.326}{Then}%
\TextAt{8}{146.754}{170.326}{one}%
\TextAt{8}{168.213}{170.326}{has}%
\TextAt{13}{189.087}{170.326}{A}%
\TextAt{21}{197.863}{168.533}{S}%
\TextAt{21}{203.45}{168.533}{T}%
\NativeMathAt{213.384}{170.326}{11.955}{\cong}%
\TextAt{9}{226.003}{170.326}{\Character{40}}%
\TextAt{13}{230.555}{170.326}{A}%
\TextAt{21}{239.331}{168.533}{S}%
\TextAt{9}{245.424}{170.326}{\Character{41}}%
\TextAt{21}{249.977}{168.533}{T}%
\TextAt{23}{256.083}{170.802}{\Character{3}}%
\TextAt{8}{260.898}{170.326}{.}%
\TextAt{12}{79.226}{138.67}{Exercise}%
\TextAt{12}{132.87}{138.67}{I.11.5.}%
\TextAt{8}{173.581}{138.67}{Let}%
\TextAt{13}{194.936}{138.67}{f}%
\TextAt{9}{205.96}{138.67}{\Character{58}}%
\TextAt{13}{213.19}{138.67}{A}%
\NativeMathAt{225.943}{138.67}{11.955}{\rightarrow}%
\TextAt{13}{241.876}{138.67}{B}%
\TextAt{8}{255.655}{138.67}{be}%
\TextAt{8}{271.965}{138.67}{a}%
\TextAt{8}{282.106}{138.67}{homomorphism}%
\TextAt{8}{365.162}{138.67}{of}%
\TextAt{8}{378.877}{138.67}{rings}%
\TextAt{8}{407.943}{138.67}{and}%
\TextAt{13}{431.092}{138.67}{S}%
\TextAt{8}{443.27}{138.67}{a}%
\TextAt{8}{453.411}{138.67}{multiplicatively}%
\TextAt{8}{79.226}{121.624}{closed}%
\TextAt{8}{113.753}{121.624}{subset}%
\TextAt{8}{149.642}{121.624}{such}%
\TextAt{8}{176.039}{121.624}{that}%
\TextAt{13}{201.408}{121.624}{f}%
\TextAt{9}{208.454}{121.624}{\Character{40}}%
\TextAt{13}{213.007}{121.624}{S}%
\TextAt{9}{220.902}{121.624}{\Character{41}}%
\NativeMathAt{228.775}{121.624}{11.955}{\subseteq}%
\TextAt{13}{241.395}{121.624}{B}%
\TextAt{3}{250.891}{125.962}{\Character{3}}%
\TextAt{8}{255.623}{121.624}{.}%
\TextAt{8}{264.071}{121.624}{Then}%
\TextAt{13}{294.629}{121.624}{f}%
\TextAt{8}{305.573}{121.624}{induces}%
\TextAt{8}{347.248}{121.624}{a}%
\TextAt{8}{357.004}{121.624}{homomorphism}%
\TextAt{13}{439.674}{121.624}{A}%
\TextAt{21}{448.449}{119.83}{S}%
\NativeMathAt{457.864}{121.624}{11.955}{\rightarrow}%
\TextAt{13}{473.14}{121.624}{B}%
\TextAt{8}{482.636}{121.624}{.}%
\TextAt{12}{79.226}{89.967}{Exercise}%
\TextAt{12}{132.509}{89.967}{I.11.6.}%
\TextAt{8}{172.905}{89.967}{Let}%
\TextAt{13}{193.946}{89.967}{A}%
\TextAt{8}{206.69}{89.967}{be}%
\TextAt{8}{222.686}{89.967}{an}%
\TextAt{8}{239.017}{89.967}{integral}%
\TextAt{8}{281.696}{89.967}{domain.}%
\TextAt{8}{328.09}{89.967}{If}%
\TextAt{8}{339.853}{89.967}{the}%
\TextAt{8}{360.081}{89.967}{localization}%
\TextAt{13}{422.272}{89.967}{A}%
\TextAt{21}{431.047}{88.174}{S}%
\TextAt{8}{441.11}{89.967}{is}%
\TextAt{8}{452.945}{89.967}{integral}%
\TextAt{8}{495.625}{89.967}{over}%
\TextAt{13}{520.742}{89.967}{A}%
\TextAt{8}{529.518}{89.967}{,}%
\TextAt{8}{79.226}{72.921}{then}%
\TextAt{13}{105.886}{72.921}{A}%
\TextAt{21}{114.662}{71.128}{S}%
\TextAt{9}{124.076}{72.921}{\Character{61}}%
\TextAt{13}{136.502}{72.921}{A}%
\TextAt{8}{145.277}{72.921}{.}%
\TextAt{8}{300.14}{43.033}{67}%
\Destination{79.23}{267.845}{nav.Item.12}%
\Destination{79.23}{250.115}{nav.Item.13}%
\Destination{79.23}{232.775}{nav.Item.14}%
\Destination{78.23}{809.815}{nav.page.67}%
\Destination{79.226}{526.655}{nav.ex.1.11.1}%
\Destination{79.226}{369.955}{nav.ex.1.11.2}%
\Destination{79.226}{304.155}{nav.ex.1.11.3}%
\Destination{79.226}{201.355}{nav.ex.1.11.4}%
\Destination{79.226}{152.655}{nav.ex.1.11.5}%
\Destination{79.226}{103.955}{nav.ex.1.11.6}%
\Destination{79.36}{370.955}{nav.print.72}%
\Destination{79.36}{775.155}{nav.number.1.11.7}%
\Destination{79.36}{558.255}{nav.number.1.11.8}%
\Hotlink{247.733}{702.055}{6.858}{11.047}{nav.section.1.5}%
\Hotlink{274.393}{702.055}{12.716}{11.047}{nav.print.30}%
\Hotlink{421.541}{650.355}{30.936}{11.047}{nav.number.1.11.6}%
\Hotlink{504.94}{650.355}{25.078}{11.047}{nav.number.1.7.4}%
}

% Page 80 (body page 68)
\Page{%
\TextAt{2}{79.226}{796.943}{I.}%
\TextAt{2}{89.83}{796.943}{Algebraic}%
\TextAt{2}{138.648}{796.943}{Integers}%
\Vector{q 0 G 0 g 0.299 w 7.226 18.729 m 460.769 18.729 l S Q}%
\TextAt{12}{79.226}{761.124}{Exercise}%
\TextAt{12}{132.122}{761.124}{I.11.7.}%
\TextAt{12}{172.183}{761.124}{(Nakayama’s}%
\TextAt{12}{252.002}{761.124}{Lemma).}%
\TextAt{8}{309.271}{761.124}{Let}%
\TextAt{13}{329.976}{761.124}{A}%
\TextAt{8}{342.384}{761.124}{be}%
\TextAt{8}{358.044}{761.124}{a}%
\TextAt{8}{367.535}{761.124}{local}%
\TextAt{8}{394.911}{761.124}{ring}%
\TextAt{8}{418.712}{761.124}{with}%
\TextAt{8}{445.108}{761.124}{maximal}%
\TextAt{8}{492.652}{761.124}{ideal}%
\TextAt{14}{520.351}{761.124}{m}%
\TextAt{8}{529.518}{761.124}{,}%
\TextAt{8}{79.226}{744.078}{let}%
\TextAt{13}{96.765}{744.078}{M}%
\TextAt{8}{113.871}{744.078}{be}%
\TextAt{8}{130.429}{744.078}{an}%
\TextAt{13}{147.323}{744.078}{A}%
\TextAt{8}{156.098}{744.078}{-module}%
\TextAt{8}{201.923}{744.078}{and}%
\TextAt{13}{225.32}{744.078}{N}%
\NativeMathAt{240.345}{744.078}{11.955}{\subseteq}%
\TextAt{13}{254.045}{744.078}{M}%
\TextAt{8}{271.151}{744.078}{a}%
\TextAt{8}{281.541}{744.078}{submodule}%
\TextAt{8}{341.09}{744.078}{such}%
\TextAt{8}{368.122}{744.078}{that}%
\TextAt{13}{394.125}{744.078}{M}%
\TextAt{9}{406.698}{744.078}{\Character{47}}%
\TextAt{13}{412.551}{744.078}{N}%
\TextAt{8}{427.706}{744.078}{is}%
\TextAt{8}{440.104}{744.078}{finitely}%
\TextAt{8}{480.083}{744.078}{generated.}%
\TextAt{8}{79.226}{727.032}{Then}%
\TextAt{8}{109.784}{727.032}{one}%
\TextAt{8}{131.243}{727.032}{has}%
\TextAt{8}{152.117}{727.032}{the}%
\TextAt{8}{172.273}{727.032}{implication:}%
\TextAt{13}{222.866}{698.031}{M}%
\TextAt{9}{238.76}{698.031}{\Character{61}}%
\TextAt{13}{251.186}{698.031}{N}%
\TextAt{9}{264.465}{698.031}{\Character{43}}%
\TextAt{14}{276.226}{698.031}{m}%
\TextAt{13}{285.393}{698.031}{M}%
\NativeArrowAt{313.243}{698.031}{11.955}{19.067}{\Longrightarrow}%
\TextAt{13}{347.586}{698.031}{M}%
\TextAt{9}{363.48}{698.031}{\Character{61}}%
\TextAt{13}{375.906}{698.031}{N}%
\TextAt{13}{385.877}{698.031}{\Character{58}}%
\TextAt{15}{79.226}{633.212}{§}%
\TextAt{15}{93.785}{633.212}{12}%
\TextAt{15}{130.385}{633.212}{Orders}%
\TextAt{8}{79.226}{601.326}{The}%
\TextAt{8}{103.174}{601.326}{ring}%
\TextAt{14}{127.134}{601.326}{o}%
\TextAt{21}{132.978}{599.533}{K}%
\TextAt{8}{144.972}{601.326}{of}%
\TextAt{8}{158.196}{601.326}{integers}%
\TextAt{8}{201.403}{601.326}{of}%
\TextAt{8}{214.627}{601.326}{an}%
\TextAt{8}{230.78}{601.326}{algebraic}%
\TextAt{8}{280.109}{601.326}{number}%
\TextAt{8}{322.599}{601.326}{field}%
\TextAt{13}{347.85}{601.326}{K}%
\TextAt{8}{362.459}{601.326}{is}%
\TextAt{8}{374.117}{601.326}{our}%
\TextAt{8}{394.825}{601.326}{chief}%
\TextAt{8}{422.024}{601.326}{interest}%
\TextAt{8}{463.929}{601.326}{because}%
\TextAt{8}{507.124}{601.326}{of}%
\TextAt{8}{520.348}{601.326}{its}%
\TextAt{8}{79.226}{584.28}{excellent}%
\TextAt{8}{127.614}{584.28}{property}%
\TextAt{8}{175.691}{584.28}{of}%
\TextAt{8}{189.289}{584.28}{being}%
\TextAt{8}{221.094}{584.28}{a}%
\TextAt{8}{231.118}{584.28}{Dedekind}%
\TextAt{8}{283.57}{584.28}{domain.}%
\TextAt{8}{330.553}{584.28}{Due}%
\TextAt{8}{355.365}{584.28}{to}%
\TextAt{8}{369.944}{584.28}{important}%
\TextAt{8}{425.504}{584.28}{theoretical}%
\TextAt{8}{483.659}{584.28}{as}%
\TextAt{8}{498.298}{584.28}{well}%
\TextAt{8}{522.297}{584.28}{as}%
\TextAt{8}{79.226}{567.234}{practical}%
\TextAt{8}{127.95}{567.234}{circumstances,}%
\TextAt{8}{206.853}{567.234}{however,}%
\TextAt{8}{255.724}{567.234}{one}%
\TextAt{8}{277.776}{567.234}{is}%
\TextAt{8}{290.132}{567.234}{pushed}%
\TextAt{8}{330.451}{567.234}{to}%
\TextAt{8}{345.353}{567.234}{devise}%
\TextAt{8}{380.794}{567.234}{a}%
\TextAt{8}{391.142}{567.234}{theory}%
\TextAt{8}{428.484}{567.234}{of}%
\TextAt{8}{442.406}{567.234}{greater}%
\TextAt{8}{482.677}{567.234}{generality}%
\TextAt{8}{79.226}{550.188}{which}%
\TextAt{8}{112.713}{550.188}{comprises}%
\TextAt{8}{166.164}{550.188}{also}%
\TextAt{8}{189.644}{550.188}{the}%
\TextAt{8}{209.801}{550.188}{theory}%
\TextAt{8}{246.551}{550.188}{of}%
\TextAt{8}{259.881}{550.188}{rings}%
\TextAt{8}{288.561}{550.188}{of}%
\TextAt{8}{301.891}{550.188}{algebraic}%
\TextAt{8}{351.326}{550.188}{integers}%
\TextAt{8}{394.639}{550.188}{which,}%
\TextAt{8}{431.378}{550.188}{like}%
\TextAt{8}{452.837}{550.188}{the}%
\TextAt{8}{472.994}{550.188}{ring}%
\TextAt{14}{245.41}{521.187}{o}%
\TextAt{9}{254.575}{521.187}{\Character{61}}%
\TextAt{18}{267.001}{521.187}{Z}%
\TextAt{9}{277.628}{521.187}{\Character{43}}%
\TextAt{18}{289.389}{521.187}{Z}%
\TextAt{16}{297.359}{531.66}{\Character{112}}%
\Vector{q 0 G 0 g 0.478 w 235.322 -242.101 m 241.175 -242.101 l S Q}%
\TextAt{9}{307.322}{521.187}{5}%
\NativeMathAt{316.496}{521.187}{11.955}{\subseteq}%
\TextAt{18}{329.115}{521.187}{Q}%
\TextAt{9}{338.413}{521.187}{\Character{40}}%
\TextAt{16}{342.966}{531.66}{\Character{112}}%
\Vector{q 0 G 0 g 0.478 w 280.928 -242.101 m 286.781 -242.101 l S Q}%
\TextAt{9}{352.928}{521.187}{5\Character{41}}%
\TextAt{13}{363.334}{521.187}{\Character{59}}%
\TextAt{8}{79.226}{492.186}{are}%
\TextAt{8}{98.737}{492.186}{not}%
\TextAt{8}{119.551}{492.186}{necessarily}%
\TextAt{8}{177.88}{492.186}{integrally}%
\TextAt{8}{229.921}{492.186}{closed.}%
\TextAt{8}{268.999}{492.186}{These}%
\TextAt{8}{302.868}{492.186}{rings}%
\TextAt{8}{331.548}{492.186}{are}%
\TextAt{8}{351.059}{492.186}{the}%
\TextAt{8}{371.215}{492.186}{so-called}%
\TextAt{12}{418.749}{492.186}{orders}%
\TextAt{8}{455.392}{492.186}{.}%
\TextAt{12}{79.226}{466.173}{Definition}%
\TextAt{12}{142.04}{466.173}{(12.1).}%
\TextAt{8}{184.109}{466.173}{Let}%
\TextAt{13}{205.1}{466.173}{K}%
\TextAt{16}{219.276}{466.173}{\Character{106}}%
\TextAt{18}{225.956}{466.173}{Q}%
\TextAt{8}{239.174}{466.173}{be}%
\TextAt{8}{255.12}{466.173}{an}%
\TextAt{8}{271.401}{466.173}{algebraic}%
\TextAt{8}{320.858}{466.173}{number}%
\TextAt{8}{363.476}{466.173}{field}%
\TextAt{8}{388.855}{466.173}{of}%
\TextAt{8}{402.208}{466.173}{degree}%
\TextAt{13}{438.645}{466.173}{n}%
\TextAt{8}{445.633}{466.173}{.}%
\TextAt{8}{454.147}{466.173}{An}%
\TextAt{12}{473.346}{466.173}{order}%
\TextAt{8}{508.6}{466.173}{of}%
\TextAt{13}{521.952}{466.173}{K}%
\TextAt{8}{79.226}{449.127}{is}%
\TextAt{8}{91.524}{449.127}{a}%
\TextAt{8}{101.813}{449.127}{subring}%
\TextAt{14}{144.035}{449.127}{o}%
\TextAt{8}{154.31}{449.127}{of}%
\TextAt{14}{168.174}{449.127}{o}%
\TextAt{21}{174.018}{447.334}{K}%
\TextAt{8}{186.652}{449.127}{which}%
\TextAt{8}{220.672}{449.127}{contains}%
\TextAt{8}{267.126}{449.127}{an}%
\TextAt{8}{283.919}{449.127}{integral}%
\TextAt{8}{327.061}{449.127}{basis}%
\TextAt{8}{356.335}{449.127}{of}%
\TextAt{8}{370.199}{449.127}{length}%
\TextAt{13}{406.503}{449.127}{n}%
\TextAt{8}{413.49}{449.127}{.}%
\TextAt{8}{423.54}{449.127}{The}%
\TextAt{8}{448.128}{449.127}{ring}%
\TextAt{14}{472.727}{449.127}{o}%
\TextAt{21}{478.572}{447.334}{K}%
\TextAt{8}{491.205}{449.127}{is}%
\TextAt{8}{503.503}{449.127}{called}%
\TextAt{8}{79.226}{432.081}{the}%
\TextAt{12}{99.383}{432.081}{maximal}%
\TextAt{12}{153.958}{432.081}{order}%
\TextAt{8}{189.19}{432.081}{of}%
\TextAt{13}{202.52}{432.081}{K}%
\TextAt{8}{213.337}{432.081}{.}%
\TextAt{8}{97.159}{406.069}{In}%
\TextAt{8}{111.78}{406.069}{concrete}%
\TextAt{8}{157.951}{406.069}{terms,}%
\TextAt{8}{193.781}{406.069}{orders}%
\TextAt{8}{228.965}{406.069}{are}%
\TextAt{8}{248.475}{406.069}{obtained}%
\TextAt{8}{296.607}{406.069}{as}%
\TextAt{8}{310.977}{406.069}{rings}%
\TextAt{8}{339.658}{406.069}{of}%
\TextAt{8}{352.988}{406.069}{the}%
\TextAt{8}{373.144}{406.069}{form}%
\TextAt{14}{261.14}{377.068}{o}%
\TextAt{9}{270.306}{377.068}{\Character{61}}%
\TextAt{18}{282.731}{377.068}{Z}%
\TextAt{9}{290.701}{377.068}{\Character{91}}%
\TextAt{13}{293.953}{377.068}{\Character{11}}%
\TextAt{17}{301.398}{375.275}{1}%
\TextAt{13}{306.131}{377.068}{\Character{59}}%
\TextAt{13}{311.355}{377.068}{\Character{58}}%
\TextAt{13}{316.592}{377.068}{\Character{58}}%
\TextAt{13}{321.816}{377.068}{\Character{58}}%
\TextAt{13}{327.053}{377.068}{\Character{59}}%
\TextAt{13}{332.289}{377.068}{\Character{11}}%
\TextAt{21}{339.797}{375.275}{r}%
\TextAt{9}{344.352}{377.068}{\Character{93}}%
\TextAt{13}{347.604}{377.068}{\Character{59}}%
\TextAt{8}{79.226}{348.067}{where}%
\TextAt{13}{113.997}{348.067}{\Character{11}}%
\TextAt{17}{121.442}{346.274}{1}%
\TextAt{13}{126.175}{348.067}{\Character{59}}%
\TextAt{13}{131.399}{348.067}{\Character{58}}%
\TextAt{13}{136.636}{348.067}{\Character{58}}%
\TextAt{13}{141.86}{348.067}{\Character{58}}%
\TextAt{13}{147.097}{348.067}{\Character{59}}%
\TextAt{13}{152.333}{348.067}{\Character{11}}%
\TextAt{21}{159.841}{346.274}{r}%
\TextAt{8}{169.255}{348.067}{are}%
\TextAt{8}{189.727}{348.067}{integers}%
\TextAt{8}{234.002}{348.067}{such}%
\TextAt{8}{261.361}{348.067}{that}%
\TextAt{13}{287.691}{348.067}{K}%
\TextAt{9}{303.468}{348.067}{\Character{61}}%
\TextAt{18}{317.532}{348.067}{Q}%
\TextAt{9}{326.831}{348.067}{\Character{40}}%
\TextAt{13}{331.383}{348.067}{\Character{11}}%
\TextAt{17}{338.829}{346.274}{1}%
\TextAt{13}{343.561}{348.067}{\Character{59}}%
\TextAt{13}{348.785}{348.067}{\Character{58}}%
\TextAt{13}{354.022}{348.067}{\Character{58}}%
\TextAt{13}{359.258}{348.067}{\Character{58}}%
\TextAt{13}{364.483}{348.067}{\Character{59}}%
\TextAt{13}{369.719}{348.067}{\Character{11}}%
\TextAt{21}{377.228}{346.274}{r}%
\TextAt{9}{381.782}{348.067}{\Character{41}}%
\TextAt{8}{386.335}{348.067}{.}%
\TextAt{8}{397.668}{348.067}{Being}%
\TextAt{8}{431.625}{348.067}{a}%
\TextAt{8}{442.342}{348.067}{submodule}%
\TextAt{8}{502.219}{348.067}{of}%
\TextAt{8}{516.51}{348.067}{the}%
\TextAt{8}{79.226}{331.021}{free}%
\TextAt{18}{101.92}{331.021}{Z}%
\TextAt{8}{109.89}{331.021}{-module}%
\TextAt{14}{155.347}{331.021}{o}%
\TextAt{21}{161.192}{329.228}{K}%
\TextAt{8}{169.394}{331.021}{,}%
\TextAt{14}{176.876}{331.021}{o}%
\TextAt{8}{186.884}{331.021}{does}%
\TextAt{8}{213.547}{331.021}{of}%
\TextAt{8}{227.144}{331.021}{course}%
\TextAt{8}{263.24}{331.021}{admit}%
\TextAt{8}{297.327}{331.021}{a}%
\TextAt{18}{307.349}{331.021}{Z}%
\TextAt{8}{315.319}{331.021}{-basis}%
\TextAt{8}{348.223}{331.021}{which,}%
\TextAt{8}{385.294}{331.021}{as}%
\TextAt{18}{399.93}{331.021}{Q}%
\TextAt{14}{409.229}{331.021}{o}%
\TextAt{9}{418.848}{331.021}{\Character{61}}%
\TextAt{13}{431.727}{331.021}{K}%
\TextAt{8}{442.545}{331.021}{,}%
\TextAt{8}{450.027}{331.021}{has}%
\TextAt{8}{471.167}{331.021}{to}%
\TextAt{8}{485.743}{331.021}{be}%
\TextAt{8}{501.934}{331.021}{at}%
\TextAt{8}{516.51}{331.021}{the}%
\TextAt{8}{79.226}{313.975}{same}%
\TextAt{8}{109.322}{313.975}{time}%
\TextAt{8}{136.752}{313.975}{a}%
\TextAt{8}{147.278}{313.975}{basis}%
\TextAt{8}{176.788}{313.975}{of}%
\TextAt{13}{190.888}{313.975}{K}%
\TextAt{16}{206.338}{313.975}{\Character{106}}%
\TextAt{18}{214.292}{313.975}{Q}%
\TextAt{8}{223.591}{313.975}{,}%
\TextAt{8}{231.702}{313.975}{and}%
\TextAt{8}{255.235}{313.975}{therefore}%
\TextAt{8}{305.105}{313.975}{has}%
\TextAt{8}{326.749}{313.975}{length}%
\TextAt{13}{363.288}{313.975}{n}%
\TextAt{8}{370.276}{313.975}{.}%
\TextAt{8}{381.034}{313.975}{Orders}%
\TextAt{8}{420.24}{313.975}{arise}%
\TextAt{8}{448.387}{313.975}{often}%
\TextAt{8}{478.746}{313.975}{as}%
\TextAt{8}{493.886}{313.975}{rings}%
\TextAt{8}{523.337}{313.975}{of}%
\TextAt{8}{79.226}{296.929}{multipliers,}%
\TextAt{8}{141.779}{296.929}{and}%
\TextAt{8}{165.367}{296.929}{as}%
\TextAt{8}{180.563}{296.929}{such}%
\TextAt{8}{207.785}{296.929}{have}%
\TextAt{8}{235.606}{296.929}{their}%
\TextAt{8}{264.394}{296.929}{practical}%
\TextAt{8}{313.351}{296.929}{applications.}%
\TextAt{8}{385.487}{296.929}{For}%
\TextAt{8}{407.281}{296.929}{instance,}%
\TextAt{8}{457.15}{296.929}{if}%
\TextAt{13}{468.699}{296.929}{\Character{11}}%
\TextAt{17}{476.145}{295.136}{1}%
\TextAt{13}{480.877}{296.929}{\Character{59}}%
\TextAt{13}{486.101}{296.929}{\Character{58}}%
\TextAt{13}{491.338}{296.929}{\Character{58}}%
\TextAt{13}{496.574}{296.929}{\Character{58}}%
\TextAt{13}{501.799}{296.929}{\Character{59}}%
\TextAt{13}{507.035}{296.929}{\Character{11}}%
\TextAt{21}{514.544}{295.136}{n}%
\TextAt{8}{524.903}{296.929}{is}%
\TextAt{8}{79.226}{279.883}{any}%
\TextAt{8}{101.343}{279.883}{basis}%
\TextAt{8}{130.083}{279.883}{of}%
\TextAt{13}{143.413}{279.883}{K}%
\TextAt{16}{157.552}{279.883}{\Character{106}}%
\TextAt{18}{164.194}{279.883}{Q}%
\TextAt{8}{177.39}{279.883}{and}%
\TextAt{13}{200.152}{279.883}{M}%
\TextAt{9}{216.047}{279.883}{\Character{61}}%
\TextAt{18}{228.472}{279.883}{Z}%
\TextAt{13}{236.442}{279.883}{\Character{11}}%
\TextAt{17}{243.888}{278.09}{1}%
\TextAt{9}{251.277}{279.883}{\Character{43}}%
\TextAt{16}{263.038}{279.883}{\Character{1}}%
\TextAt{16}{268.334}{279.883}{\Character{1}}%
\TextAt{16}{273.642}{279.883}{\Character{1}}%
\TextAt{9}{279.643}{279.883}{\Character{43}}%
\TextAt{18}{291.404}{279.883}{Z}%
\TextAt{13}{299.374}{279.883}{\Character{11}}%
\TextAt{21}{306.82}{278.09}{n}%
\TextAt{8}{312.456}{279.883}{,}%
\TextAt{8}{319.605}{279.883}{then}%
\TextAt{14}{238.724}{250.882}{o}%
\TextAt{21}{244.568}{249.089}{M}%
\TextAt{9}{257.368}{250.882}{\Character{61}}%
\TextAt{16}{269.794}{250.882}{\Character{102}}%
\TextAt{13}{275.772}{250.882}{\Character{11}}%
\NativeMathAt{286.614}{250.882}{11.955}{\in}%
\TextAt{13}{297.905}{250.882}{K}%
\TextAt{16}{312.044}{250.882}{\Character{106}}%
\TextAt{13}{318.685}{250.882}{\Character{11}M}%
\NativeMathAt{342.101}{250.882}{11.955}{\subseteq}%
\TextAt{13}{354.721}{250.882}{M}%
\TextAt{16}{367.294}{250.882}{\Character{103}}%
\TextAt{8}{79.226}{221.881}{is}%
\TextAt{8}{90.127}{221.881}{an}%
\TextAt{8}{105.523}{221.881}{order.}%
\TextAt{8}{140.356}{221.881}{The}%
\TextAt{8}{163.547}{221.881}{theoretical}%
\TextAt{8}{220.571}{221.881}{significance}%
\TextAt{8}{281.552}{221.881}{of}%
\TextAt{8}{294.019}{221.881}{orders,}%
\TextAt{8}{331.765}{221.881}{however,}%
\TextAt{8}{379.206}{221.881}{lies}%
\TextAt{8}{398.559}{221.881}{in}%
\TextAt{8}{411.349}{221.881}{the}%
\TextAt{8}{430.643}{221.881}{fact}%
\TextAt{8}{452.865}{221.881}{that}%
\TextAt{8}{477.371}{221.881}{they}%
\TextAt{8}{502.846}{221.881}{admit}%
\TextAt{8}{79.226}{204.835}{“singularities”,}%
\TextAt{8}{158.215}{204.835}{which}%
\TextAt{8}{191.104}{204.835}{are}%
\TextAt{8}{210.019}{204.835}{excluded}%
\TextAt{8}{257.864}{204.835}{as}%
\TextAt{8}{271.638}{204.835}{long}%
\TextAt{8}{296.41}{204.835}{as}%
\TextAt{8}{310.183}{204.835}{only}%
\TextAt{8}{335.278}{204.835}{Dedekind}%
\TextAt{8}{386.866}{204.835}{domains}%
\TextAt{8}{432.511}{204.835}{with}%
\TextAt{8}{458.575}{204.835}{their}%
\TextAt{8}{485.941}{204.835}{“regular”}%
\TextAt{8}{79.226}{187.789}{localizations}%
\TextAt{14}{145.96}{187.789}{o}%
\TextAt{27}{151.804}{185.995}{p}%
\TextAt{8}{160.188}{187.789}{are}%
\TextAt{8}{179.699}{187.789}{considered.}%
\TextAt{8}{241.539}{187.789}{We}%
\TextAt{8}{261.684}{187.789}{will}%
\TextAt{8}{283.789}{187.789}{explain}%
\TextAt{8}{324.436}{187.789}{what}%
\TextAt{8}{353.703}{187.789}{this}%
\TextAt{8}{376.525}{187.789}{means}%
\TextAt{8}{412.355}{187.789}{in}%
\TextAt{8}{426.008}{187.789}{the}%
\TextAt{8}{446.164}{187.789}{next}%
\TextAt{8}{472.501}{187.789}{section.}%
\TextAt{8}{97.159}{167.754}{In}%
\TextAt{8}{111.634}{167.754}{the}%
\TextAt{8}{131.644}{167.754}{preceding}%
\TextAt{8}{184.173}{167.754}{section}%
\TextAt{8}{223.108}{167.754}{we}%
\TextAt{8}{240.19}{167.754}{studied}%
\TextAt{8}{281.074}{167.754}{the}%
\TextAt{8}{301.084}{167.754}{localizations}%
\TextAt{8}{367.672}{167.754}{of}%
\TextAt{8}{380.856}{167.754}{a}%
\TextAt{8}{390.465}{167.754}{Dedekind}%
\TextAt{8}{442.504}{167.754}{domain}%
\TextAt{14}{483.986}{167.754}{o}%
\TextAt{21}{489.83}{165.961}{K}%
\TextAt{8}{498.033}{167.754}{.}%
\TextAt{8}{506.432}{167.754}{They}%
\TextAt{8}{79.226}{150.708}{are}%
\TextAt{8}{97.715}{150.708}{extension}%
\TextAt{8}{148.459}{150.708}{rings}%
\TextAt{8}{176.118}{150.708}{of}%
\TextAt{14}{188.427}{150.708}{o}%
\TextAt{21}{194.271}{148.915}{K}%
\TextAt{8}{205.349}{150.708}{which}%
\TextAt{8}{237.814}{150.708}{are}%
\TextAt{8}{256.303}{150.708}{integrally}%
\TextAt{8}{307.322}{150.708}{closed,}%
\TextAt{8}{344.283}{150.708}{yet}%
\TextAt{8}{362.772}{150.708}{no}%
\TextAt{8}{378.01}{150.708}{longer}%
\TextAt{8}{412.112}{150.708}{integral}%
\TextAt{8}{453.699}{150.708}{over}%
\TextAt{18}{477.723}{150.708}{Z}%
\TextAt{8}{485.693}{150.708}{.}%
\TextAt{8}{493.801}{150.708}{Now}%
\TextAt{8}{519.439}{150.708}{we}%
\TextAt{8}{79.226}{133.662}{study}%
\TextAt{8}{111.44}{133.662}{orders.}%
\TextAt{8}{151.162}{133.662}{They}%
\TextAt{8}{181.356}{133.662}{are}%
\TextAt{8}{200.825}{133.662}{subrings}%
\TextAt{8}{247.087}{133.662}{of}%
\TextAt{14}{260.376}{133.662}{o}%
\TextAt{21}{266.22}{131.869}{K}%
\TextAt{8}{278.28}{133.662}{which}%
\TextAt{8}{311.725}{133.662}{are}%
\TextAt{8}{331.195}{133.662}{integral}%
\TextAt{8}{373.763}{133.662}{over}%
\TextAt{18}{398.768}{133.662}{Z}%
\TextAt{8}{406.738}{133.662}{,}%
\TextAt{8}{413.854}{133.662}{yet}%
\TextAt{8}{433.324}{133.662}{no}%
\TextAt{8}{449.543}{133.662}{longer}%
\TextAt{8}{484.626}{133.662}{integrally}%
\TextAt{8}{79.226}{116.616}{closed.}%
\TextAt{8}{118.108}{116.616}{As}%
\TextAt{8}{134.808}{116.616}{a}%
\TextAt{8}{143.977}{116.616}{common}%
\TextAt{8}{190.218}{116.616}{generalization}%
\TextAt{8}{264.435}{116.616}{of}%
\TextAt{8}{277.178}{116.616}{both}%
\TextAt{8}{304.232}{116.616}{types}%
\TextAt{8}{334.597}{116.616}{of}%
\TextAt{8}{347.34}{116.616}{rings}%
\TextAt{8}{375.434}{116.616}{let}%
\TextAt{8}{391.751}{116.616}{us}%
\TextAt{8}{406.18}{116.616}{consider}%
\TextAt{8}{451.179}{116.616}{for}%
\TextAt{8}{468.477}{116.616}{now}%
\TextAt{8}{492.278}{116.616}{all}%
\TextAt{12}{507.951}{116.616}{one-}%
\TextAt{12}{79.226}{99.57}{dimensional}%
\TextAt{12}{151.885}{99.57}{noetherian}%
\TextAt{12}{217.778}{99.57}{integral}%
\TextAt{12}{265.594}{99.57}{domains}%
\TextAt{8}{314.072}{99.57}{.}%
\TextAt{8}{322.117}{99.57}{These}%
\TextAt{8}{354.777}{99.57}{are}%
\TextAt{8}{373.078}{99.57}{the}%
\TextAt{8}{392.025}{99.57}{noetherian}%
\TextAt{8}{449.025}{99.57}{integral}%
\TextAt{8}{490.424}{99.57}{domains}%
\TextAt{8}{79.226}{82.524}{in}%
\TextAt{8}{93.117}{82.524}{which}%
\TextAt{8}{126.842}{82.524}{every}%
\TextAt{8}{157.973}{82.524}{prime}%
\TextAt{8}{191.375}{82.524}{ideal}%
\TextAt{14}{219.577}{82.524}{p}%
\NativeMathAt{229.287}{82.524}{11.955}{\neq}%
\TextAt{9}{242.103}{82.524}{0}%
\TextAt{8}{252.108}{82.524}{is}%
\TextAt{8}{264.11}{82.524}{a}%
\TextAt{8}{274.104}{82.524}{maximal}%
\TextAt{8}{322.151}{82.524}{ideal.}%
\TextAt{8}{355.381}{82.524}{The}%
\TextAt{8}{379.673}{82.524}{term}%
\TextAt{8}{407.874}{82.524}{“one-dimensional”}%
\TextAt{8}{505.069}{82.524}{refers}%
\TextAt{8}{300.14}{43.033}{68}%
\Destination{78.23}{809.815}{nav.page.68}%
\Destination{79.23}{661.155}{nav.section.1.12}%
\Destination{79.226}{775.155}{nav.ex.1.11.7}%
\Destination{79.36}{421.055}{nav.print.73}%
\Destination{79.36}{480.155}{nav.number.1.12.1}%
\Destination{79.23}{661.155}{outline.17}%
}

% Page 81 (body page 69)
\Page{%
\TextAt{2}{79.226}{796.943}{I.}%
\TextAt{2}{89.83}{796.943}{Algebraic}%
\TextAt{2}{138.648}{796.943}{Integers}%
\Vector{q 0 G 0 g 0.299 w 7.226 18.729 m 460.769 18.729 l S Q}%
\TextAt{8}{79.226}{761.124}{to}%
\TextAt{8}{93.648}{761.124}{the}%
\TextAt{8}{113.915}{761.124}{general}%
\TextAt{8}{154.351}{761.124}{definition}%
\TextAt{8}{206.492}{761.124}{of}%
\TextAt{8}{219.933}{761.124}{the}%
\TextAt{12}{240.201}{761.124}{Krull}%
\TextAt{12}{275.8}{761.124}{dimension}%
\TextAt{8}{339.094}{761.124}{of}%
\TextAt{8}{352.535}{761.124}{a}%
\TextAt{8}{362.402}{761.124}{ring}%
\TextAt{8}{386.579}{761.124}{as}%
\TextAt{8}{401.06}{761.124}{being}%
\TextAt{8}{432.709}{761.124}{the}%
\TextAt{8}{452.977}{761.124}{maximal}%
\TextAt{8}{500.897}{761.124}{length}%
\TextAt{13}{79.226}{744.078}{d}%
\TextAt{8}{89.206}{744.078}{of}%
\TextAt{8}{102.536}{744.078}{a}%
\TextAt{8}{112.292}{744.078}{chain}%
\TextAt{8}{143.184}{744.078}{of}%
\TextAt{8}{156.514}{744.078}{prime}%
\TextAt{8}{189.677}{744.078}{ideals}%
\TextAt{14}{222.255}{744.078}{p}%
\TextAt{17}{228.238}{742.284}{0}%
\TextAt{18}{236.292}{744.078}{\Character{40}}%
\TextAt{14}{248.911}{744.078}{p}%
\TextAt{17}{254.894}{742.284}{1}%
\TextAt{18}{262.947}{744.078}{\Character{40}}%
\TextAt{16}{275.566}{744.078}{\Character{1}}%
\TextAt{16}{280.862}{744.078}{\Character{1}}%
\TextAt{16}{286.17}{744.078}{\Character{1}}%
\TextAt{18}{292.835}{744.078}{\Character{40}}%
\TextAt{14}{305.454}{744.078}{p}%
\TextAt{21}{311.437}{742.284}{d}%
\TextAt{8}{316.293}{744.078}{.}%
\TextAt{12}{79.226}{718.58}{Proposition}%
\TextAt{12}{151.937}{718.58}{(12.2).}%
\TextAt{8}{193.984}{718.58}{An}%
\TextAt{8}{213.16}{718.58}{order}%
\TextAt{14}{243.729}{718.58}{o}%
\TextAt{8}{253.471}{718.58}{of}%
\TextAt{13}{266.801}{718.58}{K}%
\TextAt{8}{281.516}{718.58}{is}%
\TextAt{8}{293.28}{718.58}{a}%
\TextAt{8}{303.035}{718.58}{one-dimensional}%
\TextAt{8}{388.945}{718.58}{noetherian}%
\TextAt{8}{447.155}{718.58}{integral}%
\TextAt{8}{489.763}{718.58}{domain.}%
\TextAt{19}{79.226}{684.63}{Proof.}%
\TextAt{8}{116}{684.63}{Since}%
\TextAt{14}{146.667}{684.63}{o}%
\TextAt{8}{156.519}{684.63}{is}%
\TextAt{8}{168.392}{684.63}{a}%
\TextAt{8}{178.257}{684.63}{finitely}%
\TextAt{8}{217.711}{684.63}{generated}%
\TextAt{18}{271.153}{684.63}{Z}%
\TextAt{8}{279.123}{684.63}{-module}%
\TextAt{8}{324.424}{684.63}{of}%
\TextAt{8}{337.863}{684.63}{rank}%
\TextAt{13}{364.968}{684.63}{n}%
\TextAt{9}{375.463}{684.63}{\Character{61}}%
\TextAt{9}{388.064}{684.63}{\Character{91}}%
\TextAt{13}{391.327}{684.63}{K}%
\TextAt{9}{405.653}{684.63}{\Character{58}}%
\TextAt{18}{412.412}{684.63}{Q}%
\TextAt{9}{421.711}{684.63}{\Character{93}}%
\TextAt{8}{424.962}{684.63}{,}%
\TextAt{8}{432.249}{684.63}{every}%
\TextAt{8}{463.25}{684.63}{ideal}%
\TextAt{14}{491.323}{684.63}{a}%
\TextAt{8}{501.313}{684.63}{is}%
\TextAt{8}{513.187}{684.63}{also}%
\TextAt{8}{79.226}{667.584}{a}%
\TextAt{8}{89.569}{667.584}{finitely}%
\TextAt{8}{129.5}{667.584}{generated}%
\TextAt{18}{183.419}{667.584}{Z}%
\TextAt{8}{191.39}{667.584}{-module,}%
\TextAt{8}{240.566}{667.584}{and}%
\TextAt{8}{263.916}{667.584}{a}%
\TextAt{8}{274.258}{667.584}{fortiori}%
\TextAt{8}{314.202}{667.584}{a}%
\TextAt{8}{324.544}{667.584}{finitely}%
\TextAt{8}{364.476}{667.584}{generated}%
\TextAt{14}{418.395}{667.584}{o}%
\TextAt{8}{424.239}{667.584}{-module.}%
\TextAt{8}{475.742}{667.584}{This}%
\TextAt{8}{503.049}{667.584}{shows}%
\TextAt{8}{79.226}{650.538}{that}%
\TextAt{14}{104.459}{650.538}{o}%
\TextAt{8}{114.065}{650.538}{is}%
\TextAt{8}{125.693}{650.538}{noetherian.}%
\TextAt{8}{188.408}{650.538}{If}%
\TextAt{14}{199.964}{650.538}{p}%
\NativeMathAt{209.268}{650.538}{11.955}{\neq}%
\TextAt{9}{221.678}{650.538}{0}%
\TextAt{8}{231.308}{650.538}{is}%
\TextAt{8}{242.936}{650.538}{a}%
\TextAt{8}{252.556}{650.538}{prime}%
\TextAt{8}{285.584}{650.538}{ideal}%
\TextAt{8}{313.411}{650.538}{and}%
\TextAt{13}{336.038}{650.538}{a}%
\NativeMathAt{345.503}{650.538}{11.955}{\in}%
\TextAt{14}{356.794}{650.538}{p}%
\TextAt{16}{365.156}{650.538}{\Character{92}}%
\TextAt{18}{375.505}{650.538}{Z}%
\TextAt{8}{383.475}{650.538}{,}%
\TextAt{13}{390.516}{650.538}{a}%
\NativeMathAt{399.981}{650.538}{11.955}{\neq}%
\TextAt{9}{412.391}{650.538}{0}%
\TextAt{8}{418.26}{650.538}{,}%
\TextAt{8}{425.3}{650.538}{then}%
\TextAt{13}{451.825}{650.538}{a}%
\TextAt{14}{457.97}{650.538}{o}%
\NativeMathAt{467.135}{650.538}{11.955}{\subseteq}%
\TextAt{14}{479.754}{650.538}{p}%
\NativeMathAt{489.058}{650.538}{11.955}{\subseteq}%
\TextAt{14}{501.677}{650.538}{o}%
\TextAt{8}{507.521}{650.538}{,}%
\TextAt{8}{514.562}{650.538}{i.e.,}%
\TextAt{14}{79.226}{633.492}{p}%
\TextAt{8}{89.719}{633.492}{and}%
\TextAt{14}{113.094}{633.492}{o}%
\TextAt{8}{123.447}{633.492}{have}%
\TextAt{8}{151.054}{633.492}{the}%
\TextAt{8}{171.823}{633.492}{same}%
\TextAt{8}{201.761}{633.492}{rank}%
\TextAt{13}{229.368}{633.492}{n}%
\TextAt{8}{236.356}{633.492}{.}%
\TextAt{8}{246.64}{633.492}{Therefore}%
\TextAt{14}{300.25}{633.492}{o}%
\TextAt{9}{306.094}{633.492}{\Character{47}}%
\TextAt{14}{311.947}{633.492}{p}%
\TextAt{8}{322.44}{633.492}{is}%
\TextAt{8}{334.816}{633.492}{a}%
\TextAt{8}{345.183}{633.492}{finite}%
\TextAt{8}{375.707}{633.492}{integral}%
\TextAt{8}{418.927}{633.492}{domain,}%
\TextAt{8}{464.572}{633.492}{hence}%
\TextAt{8}{497.691}{633.492}{a}%
\TextAt{8}{508.058}{633.492}{field,}%
\TextAt{8}{79.226}{616.446}{and}%
\TextAt{8}{101.989}{616.446}{thus}%
\TextAt{14}{127.74}{616.446}{p}%
\TextAt{8}{137.621}{616.446}{is}%
\TextAt{8}{149.385}{616.446}{a}%
\TextAt{8}{159.14}{616.446}{maximal}%
\TextAt{8}{206.949}{616.446}{ideal.}%
\Vector{q 0 G 0 g 0.399 w 452.334 -157.554 m 452.334 -149.484 l S Q}%
\Vector{q 0 G 0 g 0.399 w 452.534 -149.683 m 459.707 -149.683 l S Q}%
\Vector{q 0 G 0 g 0.399 w 452.534 -157.354 m 459.707 -157.354 l S Q}%
\Vector{q 0 G 0 g 0.399 w 459.906 -157.554 m 459.906 -149.484 l S Q}%
\TextAt{8}{97.159}{587.96}{In}%
\TextAt{8}{111.451}{587.96}{what}%
\TextAt{8}{140.389}{587.96}{follows,}%
\TextAt{8}{181.814}{587.96}{we}%
\TextAt{8}{198.713}{587.96}{always}%
\TextAt{8}{235.852}{587.96}{let}%
\TextAt{14}{252.428}{587.96}{o}%
\TextAt{8}{261.841}{587.96}{be}%
\TextAt{8}{277.436}{587.96}{a}%
\TextAt{8}{286.863}{587.96}{one-dimensional}%
\TextAt{8}{372.444}{587.96}{noetherian}%
\TextAt{8}{430.325}{587.96}{integral}%
\TextAt{8}{472.605}{587.96}{domain}%
\TextAt{8}{513.904}{587.96}{and}%
\TextAt{13}{79.226}{570.914}{K}%
\TextAt{8}{92.995}{570.914}{its}%
\TextAt{8}{108.368}{570.914}{field}%
\TextAt{8}{132.778}{570.914}{of}%
\TextAt{8}{145.162}{570.914}{fractions.}%
\TextAt{8}{197.266}{570.914}{We}%
\TextAt{8}{216.465}{570.914}{set}%
\TextAt{8}{233.786}{570.914}{out}%
\TextAt{8}{253.654}{570.914}{by}%
\TextAt{8}{268.967}{570.914}{proving}%
\TextAt{8}{310.306}{570.914}{the}%
\TextAt{8}{329.516}{570.914}{following}%
\TextAt{8}{378.005}{570.914}{stronger}%
\TextAt{8}{422.656}{570.914}{version}%
\TextAt{8}{461.448}{570.914}{of}%
\TextAt{8}{473.832}{570.914}{the}%
\TextAt{8}{493.042}{570.914}{Chinese}%
\TextAt{8}{79.226}{553.868}{remainder}%
\TextAt{8}{134.507}{553.868}{theorem.}%
\TextAt{12}{79.226}{528.37}{Proposition}%
\TextAt{12}{151.937}{528.37}{(12.3).}%
\TextAt{8}{193.984}{528.37}{If}%
\TextAt{14}{205.676}{528.37}{a}%
\NativeMathAt{214.98}{528.37}{11.955}{\neq}%
\TextAt{9}{227.39}{528.37}{0}%
\TextAt{8}{237.156}{528.37}{is}%
\TextAt{8}{248.92}{528.37}{an}%
\TextAt{8}{265.179}{528.37}{ideal}%
\TextAt{8}{293.142}{528.37}{of}%
\TextAt{14}{306.472}{528.37}{o}%
\TextAt{8}{312.316}{528.37}{,}%
\TextAt{8}{319.465}{528.37}{then}%
\TextAt{14}{227.035}{501.17}{o}%
\TextAt{9}{232.88}{501.17}{\Character{47}}%
\TextAt{14}{238.733}{501.17}{a}%
\NativeMathAt{248.036}{501.17}{11.955}{\cong}%
\TextAt{20}{260.656}{512.527}{\Character{77}}%
\TextAt{27}{267.694}{487.421}{p}%
\TextAt{14}{280.714}{501.17}{o}%
\TextAt{27}{286.558}{499.377}{p}%
\TextAt{9}{291.045}{501.17}{\Character{47}}%
\TextAt{14}{296.898}{501.17}{ao}%
\TextAt{27}{308.725}{499.377}{p}%
\NativeMathAt{316.533}{501.17}{11.955}{\cong}%
\TextAt{20}{329.152}{512.527}{\Character{77}}%
\TextAt{27}{330.903}{487.273}{p}%
\NativeMathAt{334.892}{487.273}{7.97}{\supseteq}%
\TextAt{27}{341.478}{487.273}{a}%
\TextAt{14}{349.21}{501.17}{o}%
\TextAt{27}{355.055}{499.377}{p}%
\TextAt{9}{359.542}{501.17}{\Character{47}}%
\TextAt{14}{365.394}{501.17}{ao}%
\TextAt{27}{377.222}{499.377}{p}%
\TextAt{13}{381.709}{501.17}{\Character{58}}%
\TextAt{19}{79.226}{447.026}{Proof.}%
\TextAt{8}{116}{447.026}{Let}%
\NativeMathAt{137.614}{447.026}{11.955}{\bar{\mathfrak a}}%
\TextAt{27}{143.597}{445.233}{p}%
\TextAt{9}{152.502}{447.026}{\Character{61}}%
\TextAt{14}{166.026}{447.026}{o}%
\TextAt{16}{174.966}{447.026}{\Character{92}}%
\TextAt{14}{186.032}{447.026}{ao}%
\TextAt{27}{197.859}{445.233}{p}%
\TextAt{8}{202.346}{447.026}{.}%
\TextAt{8}{212.727}{447.026}{For}%
\TextAt{8}{234.34}{447.026}{almost}%
\TextAt{8}{272.775}{447.026}{all}%
\TextAt{14}{289.678}{447.026}{p}%
\TextAt{8}{300.203}{447.026}{one}%
\TextAt{8}{322.307}{447.026}{has}%
\TextAt{14}{343.825}{447.026}{p}%
\TextAt{18}{354.226}{447.026}{\Character{43}}%
\TextAt{14}{367.944}{447.026}{a}%
\TextAt{8}{378.468}{447.026}{and}%
\TextAt{8}{401.875}{447.026}{therefore}%
\TextAt{14}{451.619}{447.026}{ao}%
\TextAt{27}{463.446}{445.233}{p}%
\TextAt{9}{472.352}{447.026}{\Character{61}}%
\TextAt{14}{485.875}{447.026}{o}%
\TextAt{27}{491.72}{445.233}{p}%
\TextAt{8}{496.206}{447.026}{,}%
\TextAt{8}{504.161}{447.026}{hence}%
\NativeMathAt{79.226}{429.98}{11.955}{\bar{\mathfrak a}}%
\TextAt{27}{85.209}{428.187}{p}%
\TextAt{9}{93.558}{429.98}{\Character{61}}%
\TextAt{14}{106.524}{429.98}{o}%
\TextAt{8}{112.368}{429.98}{.}%
\TextAt{8}{121.768}{429.98}{Furthermore,}%
\TextAt{8}{193.214}{429.98}{one}%
\TextAt{8}{214.991}{429.98}{has}%
\TextAt{14}{236.182}{429.98}{a}%
\TextAt{9}{246.027}{429.98}{\Character{61}}%
\TextAt{20}{258.993}{438.947}{\Character{84}}%
\TextAt{27}{268.956}{426.493}{p}%
\NativeMathAt{275.435}{429.98}{11.955}{\bar{\mathfrak a}}%
\TextAt{27}{281.418}{428.187}{p}%
\TextAt{9}{289.766}{429.98}{\Character{61}}%
\TextAt{20}{302.733}{438.947}{\Character{84}}%
\TextAt{27}{312.695}{426.493}{p}%
\NativeMathAt{316.684}{426.493}{7.97}{\supseteq}%
\TextAt{27}{323.27}{426.493}{a}%
\NativeMathAt{329.75}{429.98}{11.955}{\bar{\mathfrak a}}%
\TextAt{27}{335.733}{428.187}{p}%
\TextAt{8}{340.22}{429.98}{.}%
\TextAt{8}{349.62}{429.98}{Indeed,}%
\TextAt{8}{391.298}{429.98}{for}%
\TextAt{8}{409.5}{429.98}{any}%
\TextAt{13}{431.934}{429.98}{a}%
\NativeMathAt{441.941}{429.98}{11.955}{\in}%
\TextAt{20}{453.772}{438.947}{\Character{84}}%
\TextAt{27}{463.735}{426.493}{p}%
\NativeMathAt{470.214}{429.98}{11.955}{\bar{\mathfrak a}}%
\TextAt{27}{476.197}{428.187}{p}%
\TextAt{8}{480.684}{429.98}{,}%
\TextAt{8}{488.23}{429.98}{the}%
\TextAt{8}{508.704}{429.98}{ideal}%
\TextAt{14}{79.226}{412.934}{b}%
\TextAt{9}{89.756}{412.934}{\Character{61}}%
\TextAt{16}{103.256}{412.934}{\Character{102}}%
\TextAt{13}{109.234}{412.934}{x}%
\NativeMathAt{120.281}{412.934}{11.955}{\in}%
\TextAt{14}{132.646}{412.934}{o}%
\TextAt{16}{142.886}{412.934}{\Character{106}}%
\TextAt{13}{150.602}{412.934}{xa}%
\NativeMathAt{167.794}{412.934}{11.955}{\in}%
\TextAt{14}{180.159}{412.934}{a}%
\TextAt{16}{186.142}{412.934}{\Character{103}}%
\TextAt{8}{196.648}{412.934}{does}%
\TextAt{8}{223.675}{412.934}{not}%
\TextAt{8}{245.12}{412.934}{belong}%
\TextAt{8}{283.146}{412.934}{to}%
\TextAt{8}{298.087}{412.934}{any}%
\TextAt{8}{320.834}{412.934}{of}%
\TextAt{8}{334.795}{412.934}{the}%
\TextAt{8}{355.581}{412.934}{maximal}%
\TextAt{8}{404.021}{412.934}{ideals}%
\TextAt{14}{437.229}{412.934}{p}%
\TextAt{8}{447.74}{412.934}{(in}%
\TextAt{8}{466.578}{412.934}{fact,}%
\TextAt{8}{493.703}{412.934}{one}%
\TextAt{8}{515.793}{412.934}{has}%
\TextAt{13}{79.226}{395.888}{s}%
\TextAt{27}{84.74}{394.095}{p}%
\TextAt{13}{89.227}{395.888}{a}%
\NativeMathAt{99.012}{395.888}{11.955}{\in}%
\TextAt{14}{110.622}{395.888}{a}%
\TextAt{8}{120.689}{395.888}{for}%
\TextAt{8}{138.761}{395.888}{an}%
\TextAt{13}{155.208}{395.888}{s}%
\TextAt{27}{160.722}{394.095}{p}%
\NativeMathAt{168.848}{395.888}{11.955}{\notin}%
\TextAt{14}{180.458}{395.888}{p}%
\TextAt{8}{186.441}{395.888}{),}%
\TextAt{8}{198.379}{395.888}{consequently,}%
\TextAt{14}{270.213}{395.888}{b}%
\TextAt{9}{279.987}{395.888}{\Character{61}}%
\TextAt{14}{292.732}{395.888}{o}%
\TextAt{8}{298.576}{395.888}{,}%
\TextAt{8}{305.959}{395.888}{i.e.,}%
\TextAt{13}{328.298}{395.888}{a}%
\TextAt{9}{338.083}{395.888}{\Character{61}}%
\TextAt{9}{350.815}{395.888}{1}%
\TextAt{16}{359.465}{395.888}{\Character{1}}%
\TextAt{13}{365.57}{395.888}{a}%
\NativeMathAt{375.355}{395.888}{11.955}{\in}%
\TextAt{14}{386.965}{395.888}{a}%
\TextAt{8}{392.948}{395.888}{,}%
\TextAt{8}{400.331}{395.888}{as}%
\TextAt{8}{414.888}{395.888}{claimed.}%
\TextAt{8}{462.92}{395.888}{(11.1)}%
\TextAt{8}{496.94}{395.888}{implies}%
\TextAt{8}{79.226}{378.842}{that,}%
\TextAt{8}{107.886}{378.842}{if}%
\TextAt{14}{118.64}{378.842}{p}%
\NativeMathAt{127.997}{378.842}{11.955}{\supseteq}%
\TextAt{14}{140.669}{378.842}{a}%
\TextAt{8}{146.652}{378.842}{,}%
\TextAt{8}{153.84}{378.842}{then}%
\TextAt{14}{180.531}{378.842}{p}%
\TextAt{8}{190.443}{378.842}{is}%
\TextAt{8}{202.238}{378.842}{the}%
\TextAt{8}{222.425}{378.842}{only}%
\TextAt{8}{248.148}{378.842}{prime}%
\TextAt{8}{281.342}{378.842}{ideal}%
\TextAt{8}{309.337}{378.842}{containing}%
\NativeMathAt{366.286}{378.842}{11.955}{\bar{\mathfrak a}}%
\TextAt{27}{372.269}{377.049}{p}%
\TextAt{8}{376.756}{378.842}{.}%
\TextAt{8}{385.297}{378.842}{Therefore,}%
\TextAt{8}{441.585}{378.842}{given}%
\TextAt{8}{472.186}{378.842}{two}%
\TextAt{8}{494.334}{378.842}{distinct}%
\TextAt{8}{79.226}{361.796}{prime}%
\TextAt{8}{113.162}{361.796}{ideals}%
\TextAt{14}{146.512}{361.796}{p}%
\TextAt{8}{157.165}{361.796}{and}%
\TextAt{14}{180.699}{361.796}{q}%
\TextAt{8}{191.216}{361.796}{of}%
\TextAt{14}{205.319}{361.796}{o}%
\TextAt{8}{211.163}{361.796}{,}%
\TextAt{8}{219.277}{361.796}{the}%
\TextAt{8}{240.206}{361.796}{ideal}%
\NativeMathAt{268.941}{361.796}{11.955}{\bar{\mathfrak a}}%
\TextAt{27}{274.924}{360.003}{p}%
\TextAt{9}{282.594}{361.796}{\Character{43}}%
\NativeMathAt{294.882}{361.796}{11.955}{\bar{\mathfrak a}}%
\TextAt{27}{300.865}{360.003}{q}%
\TextAt{8}{309.931}{361.796}{cannot}%
\TextAt{8}{349.079}{361.796}{be}%
\TextAt{8}{365.775}{361.796}{contained}%
\TextAt{8}{419.557}{361.796}{in}%
\TextAt{8}{433.982}{361.796}{any}%
\TextAt{8}{456.871}{361.796}{maximal}%
\TextAt{8}{505.452}{361.796}{ideal,}%
\TextAt{8}{79.226}{344.75}{whence}%
\NativeMathAt{120.185}{344.75}{11.955}{\bar{\mathfrak a}}%
\TextAt{27}{126.168}{342.957}{p}%
\TextAt{9}{133.311}{344.75}{\Character{43}}%
\NativeMathAt{145.072}{344.75}{11.955}{\bar{\mathfrak a}}%
\TextAt{27}{151.056}{342.957}{q}%
\TextAt{9}{158.773}{344.75}{\Character{61}}%
\TextAt{14}{171.198}{344.75}{o}%
\TextAt{8}{177.043}{344.75}{.}%
\TextAt{8}{185.491}{344.75}{The}%
\TextAt{8}{209.545}{344.75}{Chinese}%
\TextAt{8}{253.169}{344.75}{remainder}%
\TextAt{8}{308.45}{344.75}{theorem}%
\TextAt{8}{353.975}{344.75}{(3.6)}%
\TextAt{8}{381.95}{344.75}{now}%
\TextAt{8}{406.339}{344.75}{gives}%
\TextAt{8}{435.019}{344.75}{the}%
\TextAt{8}{455.176}{344.75}{isomorphism}%
\TextAt{14}{266.449}{317.55}{o}%
\TextAt{9}{272.293}{317.55}{\Character{47}}%
\TextAt{14}{278.146}{317.55}{a}%
\NativeMathAt{287.45}{317.55}{11.955}{\cong}%
\TextAt{20}{300.07}{328.908}{\Character{77}}%
\TextAt{27}{301.82}{303.653}{p}%
\NativeMathAt{305.809}{303.653}{7.97}{\supseteq}%
\TextAt{27}{312.396}{303.653}{a}%
\TextAt{14}{320.128}{317.55}{o}%
\TextAt{9}{325.972}{317.55}{\Character{47}}%
\NativeMathAt{331.825}{317.55}{11.955}{\bar{\mathfrak a}}%
\TextAt{27}{337.808}{315.757}{p}%
\TextAt{13}{342.295}{317.55}{\Character{59}}%
\TextAt{8}{79.226}{280.834}{and}%
\TextAt{8}{101.989}{280.834}{we}%
\TextAt{8}{119.216}{280.834}{have}%
\TextAt{14}{146.211}{280.834}{o}%
\TextAt{9}{152.055}{280.834}{\Character{47}}%
\NativeMathAt{157.908}{280.834}{11.955}{\bar{\mathfrak a}}%
\TextAt{27}{163.891}{279.041}{p}%
\NativeMathAt{171.699}{280.834}{11.955}{\cong}%
\TextAt{14}{184.318}{280.834}{o}%
\TextAt{27}{190.163}{279.041}{p}%
\TextAt{9}{194.649}{280.834}{\Character{47}}%
\TextAt{14}{200.502}{280.834}{ao}%
\TextAt{27}{212.33}{279.041}{p}%
\TextAt{8}{216.817}{280.834}{,}%
\TextAt{8}{223.966}{280.834}{because}%
\TextAt{14}{267.267}{280.834}{p}%
\TextAt{8}{276.571}{280.834}{mod}%
\NativeMathAt{302.009}{280.834}{11.955}{\bar{\mathfrak a}}%
\TextAt{27}{307.992}{279.041}{p}%
\TextAt{8}{316.377}{280.834}{is}%
\TextAt{8}{328.14}{280.834}{the}%
\TextAt{8}{348.297}{280.834}{only}%
\TextAt{8}{373.988}{280.834}{maximal}%
\TextAt{8}{421.797}{280.834}{ideal}%
\TextAt{8}{449.76}{280.834}{of}%
\TextAt{14}{463.09}{280.834}{o}%
\TextAt{9}{468.935}{280.834}{\Character{47}}%
\NativeMathAt{474.788}{280.834}{11.955}{\bar{\mathfrak a}}%
\TextAt{27}{480.771}{279.041}{p}%
\TextAt{8}{485.258}{280.834}{.}%
\Vector{q 0 G 0 g 0.399 w 452.334 -493.166 m 452.334 -485.096 l S Q}%
\Vector{q 0 G 0 g 0.399 w 452.534 -485.295 m 459.707 -485.295 l S Q}%
\Vector{q 0 G 0 g 0.399 w 452.534 -492.967 m 459.707 -492.967 l S Q}%
\Vector{q 0 G 0 g 0.399 w 459.906 -493.166 m 459.906 -485.096 l S Q}%
\TextAt{8}{97.159}{252.347}{For}%
\TextAt{8}{119.744}{252.347}{the}%
\TextAt{8}{141.515}{252.347}{ring}%
\TextAt{14}{167.197}{252.347}{o}%
\TextAt{8}{173.041}{252.347}{,}%
\TextAt{8}{182.209}{252.347}{the}%
\TextAt{12}{203.981}{252.347}{fractional}%
\TextAt{12}{265.914}{252.347}{ideals}%
\TextAt{8}{304.363}{252.347}{of}%
\TextAt{14}{319.309}{252.347}{o}%
\TextAt{8}{325.153}{252.347}{,}%
\TextAt{8}{334.321}{252.347}{in}%
\TextAt{8}{349.589}{252.347}{other}%
\TextAt{8}{381.774}{252.347}{words,}%
\TextAt{8}{420.603}{252.347}{the}%
\TextAt{8}{442.375}{252.347}{finitely}%
\TextAt{8}{483.335}{252.347}{generated}%
\TextAt{8}{79.226}{235.302}{nonzero}%
\TextAt{14}{121.541}{235.302}{o}%
\TextAt{8}{127.385}{235.302}{-submodules}%
\TextAt{8}{193.55}{235.302}{of}%
\TextAt{8}{205.619}{235.302}{the}%
\TextAt{8}{224.513}{235.302}{field}%
\TextAt{8}{248.608}{235.302}{of}%
\TextAt{8}{260.677}{235.302}{fractions}%
\TextAt{13}{307.283}{235.302}{K}%
\TextAt{8}{318.101}{235.302}{,}%
\TextAt{8}{324.241}{235.302}{no}%
\TextAt{8}{339.238}{235.302}{longer}%
\TextAt{8}{373.1}{235.302}{form}%
\TextAt{8}{399.479}{235.302}{a}%
\TextAt{8}{407.972}{235.302}{group—unless}%
\TextAt{14}{482.279}{235.302}{o}%
\TextAt{8}{490.759}{235.302}{happens}%
\TextAt{8}{79.226}{218.256}{to}%
\TextAt{8}{94.33}{218.256}{be}%
\TextAt{8}{111.048}{218.256}{Dedekind.}%
\TextAt{8}{170.163}{218.256}{The}%
\TextAt{8}{195.01}{218.256}{way}%
\TextAt{8}{219.547}{218.256}{out}%
\TextAt{8}{241.154}{218.256}{is}%
\TextAt{8}{253.711}{218.256}{to}%
\TextAt{8}{268.815}{218.256}{restrict}%
\TextAt{8}{309.993}{218.256}{attention}%
\TextAt{8}{361.201}{218.256}{to}%
\TextAt{8}{376.305}{218.256}{the}%
\TextAt{12}{397.255}{218.256}{invertible}%
\TextAt{12}{458.182}{218.256}{ideals}%
\TextAt{8}{491.119}{218.256}{,}%
\TextAt{8}{499.26}{218.256}{i.e.,}%
\TextAt{8}{522.357}{218.256}{to}%
\TextAt{8}{79.226}{201.21}{those}%
\TextAt{8}{109.855}{201.21}{fractional}%
\TextAt{8}{162.219}{201.21}{ideals}%
\TextAt{14}{194.797}{201.21}{a}%
\TextAt{8}{204.677}{201.21}{of}%
\TextAt{14}{218.007}{201.21}{o}%
\TextAt{8}{227.749}{201.21}{for}%
\TextAt{8}{245.634}{201.21}{which}%
\TextAt{8}{279.12}{201.21}{there}%
\TextAt{8}{309.032}{201.21}{exists}%
\TextAt{8}{341.347}{201.21}{a}%
\TextAt{8}{351.102}{201.21}{fractional}%
\TextAt{8}{403.466}{201.21}{ideal}%
\TextAt{14}{431.429}{201.21}{b}%
\TextAt{8}{441.461}{201.21}{such}%
\TextAt{8}{467.858}{201.21}{that}%
\TextAt{14}{287.518}{174.009}{ab}%
\TextAt{9}{302.957}{174.009}{\Character{61}}%
\TextAt{14}{315.382}{174.009}{o}%
\TextAt{13}{321.226}{174.009}{\Character{58}}%
\TextAt{8}{79.226}{146.809}{These}%
\TextAt{8}{113.74}{146.809}{form}%
\TextAt{8}{142.026}{146.809}{an}%
\TextAt{8}{158.93}{146.809}{abelian}%
\TextAt{8}{200.223}{146.809}{group,}%
\TextAt{8}{237.457}{146.809}{for}%
\TextAt{8}{255.987}{146.809}{trivial}%
\TextAt{8}{291.433}{146.809}{reasons.}%
\TextAt{8}{339.022}{146.809}{The}%
\TextAt{8}{363.72}{146.809}{inverse}%
\TextAt{8}{403.124}{146.809}{of}%
\TextAt{14}{417.099}{146.809}{a}%
\TextAt{8}{427.625}{146.809}{is}%
\TextAt{8}{440.034}{146.809}{still}%
\TextAt{8}{463.502}{146.809}{the}%
\TextAt{8}{484.303}{146.809}{fractional}%
\TextAt{8}{79.226}{129.763}{ideal}%
\TextAt{14}{243.638}{112.717}{a}%
\TextAt{3}{249.622}{117.653}{\Character{0}}%
\TextAt{17}{256.208}{117.653}{1}%
\TextAt{9}{264.261}{112.717}{\Character{61}}%
\TextAt{16}{276.687}{112.717}{\Character{102}}%
\TextAt{13}{282.664}{112.717}{x}%
\NativeMathAt{292.637}{112.717}{11.955}{\in}%
\TextAt{13}{303.928}{112.717}{K}%
\TextAt{16}{318.067}{112.717}{\Character{106}}%
\TextAt{13}{324.708}{112.717}{x}%
\TextAt{14}{331.36}{112.717}{a}%
\NativeMathAt{340.664}{112.717}{11.955}{\subseteq}%
\TextAt{14}{353.284}{112.717}{o}%
\TextAt{16}{359.128}{112.717}{\Character{103}}%
\TextAt{13}{365.106}{112.717}{\Character{59}}%
\TextAt{8}{79.226}{89.967}{because}%
\TextAt{8}{123.744}{89.967}{it}%
\TextAt{8}{136.665}{89.967}{is}%
\TextAt{8}{149.645}{89.967}{the}%
\TextAt{8}{171.018}{89.967}{biggest}%
\TextAt{8}{211.974}{89.967}{ideal}%
\TextAt{8}{241.153}{89.967}{such}%
\TextAt{8}{268.767}{89.967}{that}%
\TextAt{14}{295.352}{89.967}{aa}%
\TextAt{3}{307.318}{94.306}{\Character{0}}%
\TextAt{17}{313.905}{94.306}{1}%
\NativeMathAt{324.031}{89.967}{11.955}{\subseteq}%
\TextAt{14}{338.723}{89.967}{o}%
\TextAt{8}{344.568}{89.967}{.}%
\TextAt{8}{356.665}{89.967}{The}%
\TextAt{8}{381.936}{89.967}{invertible}%
\TextAt{8}{434.858}{89.967}{ideals}%
\TextAt{8}{468.653}{89.967}{of}%
\TextAt{14}{483.199}{89.967}{o}%
\TextAt{8}{494.157}{89.967}{may}%
\TextAt{8}{520.743}{89.967}{be}%
\TextAt{8}{79.226}{72.921}{characterized}%
\TextAt{8}{150.443}{72.921}{as}%
\TextAt{8}{164.813}{72.921}{those}%
\TextAt{8}{195.442}{72.921}{fractional}%
\TextAt{8}{247.806}{72.921}{ideals}%
\TextAt{8}{280.384}{72.921}{which}%
\TextAt{8}{313.87}{72.921}{are}%
\TextAt{8}{333.381}{72.921}{“locally”}%
\TextAt{8}{381.501}{72.921}{principal:}%
\TextAt{8}{300.14}{43.033}{69}%
\Destination{78.23}{809.815}{nav.page.69}%
\Destination{79.36}{543.355}{nav.print.74}%
\Destination{79.36}{732.555}{nav.number.1.12.2}%
\Destination{79.36}{542.355}{nav.number.1.12.3}%
\Hotlink{462.42}{393.655}{30.936}{11.047}{nav.number.1.11.1}%
\Destination{79.226}{413}{errata.62}%
\Destination{79.226}{460}{errata.180}%
\Hotlink{355.135}{341.757}{21.722}{11.047}{nav.number.1.3.6}%
}

% Page 82 (body page 70)
\Page{%
\TextAt{2}{79.226}{796.943}{I.}%
\TextAt{2}{89.83}{796.943}{Algebraic}%
\TextAt{2}{138.648}{796.943}{Integers}%
\Vector{q 0 G 0 g 0.299 w 7.226 18.729 m 460.769 18.729 l S Q}%
\TextAt{12}{79.226}{761.124}{Proposition}%
\TextAt{12}{151.805}{761.124}{(12.4).}%
\TextAt{8}{193.737}{761.124}{A}%
\TextAt{8}{206.294}{761.124}{fractional}%
\TextAt{8}{258.543}{761.124}{ideal}%
\TextAt{14}{286.391}{761.124}{a}%
\TextAt{8}{296.157}{761.124}{of}%
\TextAt{14}{309.372}{761.124}{o}%
\TextAt{8}{318.999}{761.124}{is}%
\TextAt{8}{330.648}{761.124}{invertible}%
\TextAt{8}{382.239}{761.124}{if}%
\TextAt{8}{392.848}{761.124}{and}%
\TextAt{8}{415.495}{761.124}{only}%
\TextAt{8}{441.072}{761.124}{if,}%
\TextAt{8}{454.956}{761.124}{for}%
\TextAt{8}{472.726}{761.124}{every}%
\TextAt{8}{503.503}{761.124}{prime}%
\TextAt{8}{79.226}{744.078}{ideal}%
\TextAt{14}{107.189}{744.078}{p}%
\NativeMathAt{116.493}{744.078}{11.955}{\neq}%
\TextAt{9}{128.903}{744.078}{0}%
\TextAt{8}{134.772}{744.078}{,}%
\TextAt{14}{284.733}{727.032}{a}%
\TextAt{27}{290.716}{725.238}{p}%
\TextAt{9}{298.523}{727.032}{\Character{61}}%
\TextAt{14}{310.949}{727.032}{ao}%
\TextAt{27}{322.776}{725.238}{p}%
\TextAt{8}{79.226}{703.51}{is}%
\TextAt{8}{90.99}{703.51}{a}%
\TextAt{8}{100.745}{703.51}{fractional}%
\TextAt{8}{153.109}{703.51}{principal}%
\TextAt{8}{201.886}{703.51}{ideal}%
\TextAt{8}{229.849}{703.51}{of}%
\TextAt{14}{243.179}{703.51}{o}%
\TextAt{27}{249.024}{701.717}{p}%
\TextAt{8}{253.51}{703.51}{.}%
\TextAt{19}{79.226}{668.531}{Proof.}%
\TextAt{8}{116}{668.531}{Let}%
\TextAt{14}{137.332}{668.531}{a}%
\TextAt{8}{147.575}{668.531}{be}%
\TextAt{8}{163.862}{668.531}{an}%
\TextAt{8}{180.483}{668.531}{invertible}%
\TextAt{8}{232.552}{668.531}{ideal}%
\TextAt{8}{260.878}{668.531}{and}%
\TextAt{14}{284.003}{668.531}{ab}%
\TextAt{9}{300.059}{668.531}{\Character{61}}%
\TextAt{14}{313.103}{668.531}{o}%
\TextAt{8}{318.947}{668.531}{.}%
\TextAt{8}{328.483}{668.531}{Then}%
\TextAt{9}{359.403}{668.531}{1}%
\TextAt{9}{369.182}{668.531}{\Character{61}}%
\TextAt{20}{382.237}{677.498}{\Character{80}}%
\TextAt{21}{394.857}{674.343}{r}%
\TextAt{21}{394.857}{665.044}{i}%
\TextAt{17}{397.74}{665.044}{\Character{61}1}%
\TextAt{13}{411.051}{668.531}{a}%
\TextAt{21}{417.196}{666.738}{i}%
\TextAt{13}{420.577}{668.531}{b}%
\TextAt{21}{425.555}{666.738}{i}%
\TextAt{8}{433.196}{668.531}{with}%
\TextAt{13}{460.218}{668.531}{a}%
\TextAt{21}{466.363}{666.738}{i}%
\NativeMathAt{473.683}{668.531}{11.955}{\in}%
\TextAt{14}{485.592}{668.531}{a}%
\TextAt{8}{491.575}{668.531}{,}%
\TextAt{13}{499.177}{668.531}{b}%
\TextAt{21}{504.154}{666.738}{i}%
\NativeMathAt{511.474}{668.531}{11.955}{\in}%
\TextAt{14}{523.383}{668.531}{b}%
\TextAt{8}{529.518}{668.531}{,}%
\TextAt{8}{79.226}{651.485}{and}%
\TextAt{8}{101.981}{651.485}{not}%
\TextAt{8}{122.786}{651.485}{all}%
\TextAt{13}{139.037}{651.485}{a}%
\TextAt{21}{145.182}{649.692}{i}%
\TextAt{13}{148.563}{651.485}{b}%
\TextAt{21}{153.541}{649.692}{i}%
\NativeMathAt{160.243}{651.485}{11.955}{\in}%
\TextAt{14}{171.534}{651.485}{o}%
\TextAt{27}{177.378}{649.692}{p}%
\TextAt{8}{185.754}{651.485}{can}%
\TextAt{8}{207.205}{651.485}{lie}%
\TextAt{8}{222.799}{651.485}{in}%
\TextAt{8}{236.443}{651.485}{the}%
\TextAt{8}{256.591}{651.485}{maximal}%
\TextAt{8}{304.392}{651.485}{ideal}%
\TextAt{14}{332.347}{651.485}{po}%
\TextAt{27}{344.174}{649.692}{p}%
\TextAt{8}{348.661}{651.485}{.}%
\TextAt{8}{357.107}{651.485}{Suppose}%
\TextAt{13}{403.006}{651.485}{a}%
\TextAt{17}{409.151}{649.692}{1}%
\TextAt{13}{413.884}{651.485}{b}%
\TextAt{17}{418.861}{649.692}{1}%
\TextAt{8}{427.482}{651.485}{is}%
\TextAt{8}{439.238}{651.485}{a}%
\TextAt{8}{448.985}{651.485}{unit}%
\TextAt{8}{473.688}{651.485}{in}%
\TextAt{14}{487.333}{651.485}{o}%
\TextAt{27}{493.177}{649.692}{p}%
\TextAt{8}{497.664}{651.485}{.}%
\TextAt{8}{506.11}{651.485}{Then}%
\TextAt{14}{79.226}{634.439}{a}%
\TextAt{27}{85.209}{632.646}{p}%
\TextAt{9}{93.017}{634.439}{\Character{61}}%
\TextAt{13}{105.442}{634.439}{a}%
\TextAt{17}{111.587}{632.646}{1}%
\TextAt{14}{116.32}{634.439}{o}%
\TextAt{27}{122.164}{632.646}{p}%
\TextAt{8}{130.548}{634.439}{because,}%
\TextAt{8}{177.102}{634.439}{for}%
\TextAt{13}{194.986}{634.439}{x}%
\NativeMathAt{204.959}{634.439}{11.955}{\in}%
\TextAt{14}{216.25}{634.439}{a}%
\TextAt{27}{222.233}{632.646}{p}%
\TextAt{8}{226.72}{634.439}{,}%
\TextAt{13}{233.87}{634.439}{xb}%
\TextAt{17}{245.499}{632.646}{1}%
\NativeMathAt{253.552}{634.439}{11.955}{\in}%
\TextAt{14}{264.843}{634.439}{a}%
\TextAt{27}{270.826}{632.646}{p}%
\TextAt{14}{275.313}{634.439}{b}%
\TextAt{9}{284.768}{634.439}{\Character{61}}%
\TextAt{14}{297.194}{634.439}{o}%
\TextAt{27}{303.038}{632.646}{p}%
\TextAt{8}{307.525}{634.439}{,}%
\TextAt{8}{314.674}{634.439}{hence}%
\TextAt{13}{347.18}{634.439}{x}%
\TextAt{9}{357.153}{634.439}{\Character{61}}%
\TextAt{13}{369.579}{634.439}{xb}%
\TextAt{17}{381.208}{632.646}{1}%
\TextAt{9}{385.94}{634.439}{\Character{40}}%
\TextAt{13}{390.493}{634.439}{b}%
\TextAt{17}{395.47}{632.646}{1}%
\TextAt{13}{400.202}{634.439}{a}%
\TextAt{17}{406.347}{632.646}{1}%
\TextAt{9}{411.079}{634.439}{\Character{41}}%
\TextAt{3}{415.632}{638.778}{\Character{0}}%
\TextAt{17}{422.218}{638.778}{1}%
\TextAt{13}{426.95}{634.439}{a}%
\TextAt{17}{433.095}{632.646}{1}%
\NativeMathAt{441.148}{634.439}{11.955}{\in}%
\TextAt{13}{452.439}{634.439}{a}%
\TextAt{17}{458.584}{632.646}{1}%
\TextAt{14}{463.317}{634.439}{o}%
\TextAt{27}{469.161}{632.646}{p}%
\TextAt{8}{473.648}{634.439}{.}%
\TextAt{8}{97.159}{614.405}{Conversely,}%
\TextAt{8}{159.002}{614.405}{assume}%
\TextAt{14}{199.704}{614.405}{a}%
\TextAt{27}{205.687}{612.611}{p}%
\TextAt{9}{213.933}{614.405}{\Character{61}}%
\TextAt{14}{226.797}{614.405}{ao}%
\TextAt{27}{238.624}{612.611}{p}%
\TextAt{8}{247.266}{614.405}{is}%
\TextAt{8}{259.287}{614.405}{a}%
\TextAt{8}{269.3}{614.405}{principal}%
\TextAt{8}{318.334}{614.405}{ideal}%
\TextAt{13}{346.554}{614.405}{a}%
\TextAt{27}{352.699}{612.611}{p}%
\TextAt{14}{357.186}{614.405}{o}%
\TextAt{27}{363.031}{612.611}{p}%
\TextAt{8}{367.517}{614.405}{,}%
\TextAt{13}{374.988}{614.405}{a}%
\TextAt{27}{381.133}{612.611}{p}%
\NativeMathAt{389.379}{614.405}{11.955}{\in}%
\TextAt{13}{401.109}{614.405}{K}%
\TextAt{3}{411.927}{618.743}{\Character{3}}%
\TextAt{8}{416.659}{614.405}{,}%
\TextAt{8}{424.13}{614.405}{for}%
\TextAt{8}{442.272}{614.405}{every}%
\TextAt{14}{473.421}{614.405}{p}%
\TextAt{8}{479.404}{614.405}{.}%
\TextAt{8}{488.625}{614.405}{Then}%
\TextAt{8}{519.44}{614.405}{we}%
\TextAt{8}{79.226}{597.359}{may}%
\TextAt{8}{104.599}{597.359}{and}%
\TextAt{8}{127.366}{597.359}{do}%
\TextAt{8}{143.63}{597.359}{assume}%
\TextAt{8}{184.079}{597.359}{that}%
\TextAt{13}{209.452}{597.359}{a}%
\TextAt{27}{215.597}{595.565}{p}%
\NativeMathAt{223.412}{597.359}{11.955}{\in}%
\TextAt{14}{234.71}{597.359}{a}%
\TextAt{8}{240.693}{597.359}{.}%
\TextAt{8}{249.155}{597.359}{We}%
\TextAt{8}{269.304}{597.359}{claim}%
\TextAt{8}{300.523}{597.359}{that}%
\TextAt{8}{325.896}{597.359}{the}%
\TextAt{8}{346.057}{597.359}{fractional}%
\TextAt{8}{398.425}{597.359}{ideal}%
\TextAt{14}{426.393}{597.359}{a}%
\TextAt{3}{432.376}{601.697}{\Character{0}}%
\TextAt{17}{438.962}{601.697}{1}%
\TextAt{9}{447.023}{597.359}{\Character{61}}%
\TextAt{16}{459.456}{597.359}{\Character{102}}%
\TextAt{13}{465.434}{597.359}{x}%
\NativeMathAt{475.414}{597.359}{11.955}{\in}%
\TextAt{13}{486.712}{597.359}{K}%
\TextAt{16}{500.858}{597.359}{\Character{106}}%
\TextAt{13}{507.508}{597.359}{x}%
\TextAt{14}{514.16}{597.359}{a}%
\NativeMathAt{523.471}{597.359}{11.955}{\subseteq}%
\TextAt{14}{79.226}{580.313}{o}%
\TextAt{16}{85.07}{580.313}{\Character{103}}%
\TextAt{8}{95.496}{580.313}{is}%
\TextAt{8}{107.81}{580.313}{an}%
\TextAt{8}{124.62}{580.313}{inverse}%
\TextAt{8}{163.929}{580.313}{for}%
\TextAt{14}{182.365}{580.313}{a}%
\TextAt{8}{188.348}{580.313}{.}%
\TextAt{8}{198.448}{580.313}{If}%
\TextAt{8}{210.69}{580.313}{this}%
\TextAt{8}{234.063}{580.313}{were}%
\TextAt{8}{261.597}{580.313}{not}%
\TextAt{8}{282.961}{580.313}{the}%
\TextAt{8}{303.668}{580.313}{case,}%
\TextAt{8}{332.379}{580.313}{then}%
\TextAt{8}{359.59}{580.313}{we}%
\TextAt{8}{377.368}{580.313}{would}%
\TextAt{8}{412.062}{580.313}{have}%
\TextAt{8}{439.608}{580.313}{a}%
\TextAt{8}{449.913}{580.313}{maximal}%
\TextAt{8}{498.273}{580.313}{ideal}%
\TextAt{14}{526.786}{580.313}{p}%
\TextAt{8}{79.226}{563.267}{such}%
\TextAt{8}{106.136}{563.267}{that}%
\TextAt{14}{132.019}{563.267}{aa}%
\TextAt{3}{143.985}{567.605}{\Character{0}}%
\TextAt{17}{150.571}{567.605}{1}%
\NativeMathAt{159.499}{563.267}{11.955}{\subseteq}%
\TextAt{14}{172.993}{563.267}{p}%
\NativeMathAt{183.171}{563.267}{11.955}{\subset}%
\TextAt{14}{196.665}{563.267}{o}%
\TextAt{8}{202.51}{563.267}{.}%
\TextAt{8}{212.498}{563.267}{Let}%
\TextAt{13}{233.98}{563.267}{a}%
\TextAt{17}{240.125}{561.473}{1}%
\TextAt{13}{244.858}{563.267}{\Character{59}}%
\TextAt{13}{250.082}{563.267}{\Character{58}}%
\TextAt{13}{255.319}{563.267}{\Character{58}}%
\TextAt{13}{260.543}{563.267}{\Character{58}}%
\TextAt{13}{265.78}{563.267}{\Character{59}}%
\TextAt{13}{271.016}{563.267}{a}%
\TextAt{21}{277.223}{561.473}{n}%
\TextAt{8}{287.27}{563.267}{be}%
\TextAt{8}{303.708}{563.267}{generators}%
\TextAt{8}{360.877}{563.267}{of}%
\TextAt{14}{374.72}{563.267}{a}%
\TextAt{8}{380.703}{563.267}{.}%
\TextAt{8}{390.691}{563.267}{As}%
\TextAt{13}{408.491}{563.267}{a}%
\TextAt{21}{414.636}{561.473}{i}%
\NativeMathAt{422.213}{563.267}{11.955}{\in}%
\TextAt{13}{434.379}{563.267}{a}%
\TextAt{27}{440.524}{561.473}{p}%
\TextAt{14}{445.011}{563.267}{o}%
\TextAt{27}{450.855}{561.473}{p}%
\TextAt{8}{455.342}{563.267}{,}%
\TextAt{8}{463.132}{563.267}{we}%
\TextAt{8}{480.873}{563.267}{may}%
\TextAt{8}{506.755}{563.267}{write}%
\TextAt{13}{79.226}{546.221}{a}%
\TextAt{21}{85.371}{544.427}{i}%
\TextAt{9}{93.664}{546.221}{\Character{61}}%
\TextAt{13}{107.681}{546.221}{a}%
\TextAt{27}{113.826}{544.427}{p}%
\TextAt{13}{118.313}{546.221}{b}%
\TextAt{21}{123.29}{544.427}{i}%
\TextAt{9}{126.671}{546.221}{\Character{47}}%
\TextAt{13}{132.524}{546.221}{s}%
\TextAt{21}{138.038}{544.427}{i}%
\TextAt{8}{141.419}{546.221}{,}%
\TextAt{8}{149.735}{546.221}{with}%
\TextAt{13}{177.329}{546.221}{b}%
\TextAt{21}{182.306}{544.427}{i}%
\NativeMathAt{190.599}{546.221}{11.955}{\in}%
\TextAt{14}{203.481}{546.221}{o}%
\TextAt{8}{209.326}{546.221}{,}%
\TextAt{13}{217.642}{546.221}{s}%
\TextAt{21}{223.156}{544.427}{i}%
\NativeMathAt{231.449}{546.221}{11.955}{\in}%
\TextAt{14}{244.331}{546.221}{o}%
\TextAt{16}{253.468}{546.221}{\Character{110}}%
\TextAt{14}{262.739}{546.221}{p}%
\TextAt{8}{268.722}{546.221}{.}%
\TextAt{8}{279.971}{546.221}{Then}%
\TextAt{13}{311.462}{546.221}{s}%
\TextAt{21}{316.976}{544.427}{i}%
\TextAt{13}{320.358}{546.221}{a}%
\TextAt{21}{326.503}{544.427}{i}%
\NativeMathAt{334.796}{546.221}{11.955}{\in}%
\TextAt{13}{347.678}{546.221}{a}%
\TextAt{27}{353.823}{544.427}{p}%
\TextAt{14}{358.31}{546.221}{o}%
\TextAt{8}{364.154}{546.221}{.}%
\TextAt{8}{375.403}{546.221}{Putting}%
\TextAt{13}{419.423}{546.221}{s}%
\TextAt{9}{429.849}{546.221}{\Character{61}}%
\TextAt{13}{443.866}{546.221}{s}%
\TextAt{17}{449.38}{544.427}{1}%
\TextAt{16}{456.104}{546.221}{\Character{1}}%
\TextAt{16}{461.4}{546.221}{\Character{1}}%
\TextAt{16}{466.708}{546.221}{\Character{1}}%
\TextAt{13}{472.045}{546.221}{s}%
\TextAt{21}{477.559}{544.427}{n}%
\TextAt{8}{483.195}{546.221}{,}%
\TextAt{8}{491.511}{546.221}{we}%
\TextAt{8}{509.672}{546.221}{have}%
\TextAt{13}{79.226}{529.175}{sa}%
\TextAt{21}{90.885}{527.382}{i}%
\NativeMathAt{99.469}{529.175}{11.955}{\in}%
\TextAt{13}{112.642}{529.175}{a}%
\TextAt{27}{118.787}{527.382}{p}%
\TextAt{14}{123.274}{529.175}{o}%
\TextAt{8}{134.12}{529.175}{for}%
\TextAt{13}{153.11}{529.175}{i}%
\TextAt{9}{162.306}{529.175}{\Character{61}}%
\TextAt{9}{176.604}{529.175}{1}%
\TextAt{13}{182.467}{529.175}{\Character{59}}%
\TextAt{13}{187.691}{529.175}{\Character{58}}%
\TextAt{13}{192.928}{529.175}{\Character{58}}%
\TextAt{13}{198.152}{529.175}{\Character{58}}%
\TextAt{13}{203.389}{529.175}{\Character{59}}%
\TextAt{13}{208.625}{529.175}{n}%
\TextAt{8}{215.675}{529.175}{,}%
\TextAt{8}{224.205}{529.175}{hence}%
\TextAt{13}{257.815}{529.175}{sa}%
\TextAt{3}{269.474}{533.513}{\Character{0}}%
\TextAt{17}{276.061}{533.513}{1}%
\TextAt{27}{269.474}{526.219}{p}%
\TextAt{14}{280.793}{529.175}{a}%
\NativeMathAt{291.979}{529.175}{11.955}{\subseteq}%
\TextAt{14}{306.481}{529.175}{o}%
\TextAt{8}{317.327}{529.175}{and}%
\TextAt{8}{341.194}{529.175}{therefore}%
\TextAt{13}{391.398}{529.175}{sa}%
\TextAt{3}{403.057}{533.513}{\Character{0}}%
\TextAt{17}{409.644}{533.513}{1}%
\TextAt{27}{403.057}{526.219}{p}%
\NativeMathAt{419.579}{529.175}{11.955}{\in}%
\TextAt{14}{432.752}{529.175}{a}%
\TextAt{3}{438.735}{533.513}{\Character{0}}%
\TextAt{17}{445.322}{533.513}{1}%
\TextAt{8}{450.054}{529.175}{.}%
\TextAt{8}{461.816}{529.175}{Consequently,}%
\TextAt{13}{79.226}{512.129}{s}%
\TextAt{9}{88.061}{512.129}{\Character{61}}%
\TextAt{13}{100.486}{512.129}{sa}%
\TextAt{3}{112.145}{516.467}{\Character{0}}%
\TextAt{17}{118.732}{516.467}{1}%
\TextAt{27}{112.145}{509.173}{p}%
\TextAt{13}{123.464}{512.129}{a}%
\TextAt{27}{129.609}{510.336}{p}%
\NativeMathAt{137.417}{512.129}{11.955}{\in}%
\TextAt{14}{148.708}{512.129}{a}%
\TextAt{3}{154.691}{516.467}{\Character{0}}%
\TextAt{17}{161.277}{516.467}{1}%
\TextAt{14}{166.01}{512.129}{a}%
\NativeMathAt{175.314}{512.129}{11.955}{\subseteq}%
\TextAt{14}{187.933}{512.129}{p}%
\TextAt{8}{193.916}{512.129}{,}%
\TextAt{8}{201.065}{512.129}{a}%
\TextAt{8}{210.821}{512.129}{contradiction.}%
\Vector{q 0 G 0 g 0.399 w 452.334 -261.871 m 452.334 -253.801 l S Q}%
\Vector{q 0 G 0 g 0.399 w 452.534 -254.001 m 459.707 -254.001 l S Q}%
\Vector{q 0 G 0 g 0.399 w 452.534 -261.672 m 459.707 -261.672 l S Q}%
\Vector{q 0 G 0 g 0.399 w 459.906 -261.871 m 459.906 -253.801 l S Q}%
\TextAt{8}{97.159}{483.128}{We}%
\TextAt{8}{117.699}{483.128}{denote}%
\TextAt{8}{155.813}{483.128}{the}%
\TextAt{8}{176.365}{483.128}{group}%
\TextAt{8}{209.937}{483.128}{of}%
\TextAt{8}{223.662}{483.128}{invertible}%
\TextAt{8}{275.764}{483.128}{ideals}%
\TextAt{8}{308.738}{483.128}{of}%
\TextAt{14}{322.463}{483.128}{o}%
\TextAt{8}{332.601}{483.128}{by}%
\TextAt{13}{349.255}{483.128}{J}%
\TextAt{9}{356.832}{483.128}{\Character{40}}%
\TextAt{14}{361.385}{483.128}{o}%
\TextAt{9}{367.229}{483.128}{\Character{41}}%
\TextAt{8}{371.781}{483.128}{.}%
\TextAt{8}{381.417}{483.128}{It}%
\TextAt{8}{394.485}{483.128}{contains}%
\TextAt{8}{440.8}{483.128}{the}%
\TextAt{8}{461.353}{483.128}{group}%
\TextAt{13}{494.924}{483.128}{P}%
\TextAt{9}{504.095}{483.128}{\Character{40}}%
\TextAt{14}{508.647}{483.128}{o}%
\TextAt{9}{514.491}{483.128}{\Character{41}}%
\TextAt{8}{523.337}{483.128}{of}%
\TextAt{8}{79.226}{466.082}{fractional}%
\TextAt{8}{131.59}{466.082}{principal}%
\TextAt{8}{180.367}{466.082}{ideals}%
\TextAt{13}{212.945}{466.082}{a}%
\TextAt{14}{219.09}{466.082}{o}%
\TextAt{8}{224.934}{466.082}{,}%
\TextAt{13}{232.083}{466.082}{a}%
\NativeMathAt{241.549}{466.082}{11.955}{\in}%
\TextAt{13}{252.84}{466.082}{K}%
\TextAt{3}{263.657}{470.42}{\Character{3}}%
\TextAt{8}{268.39}{466.082}{.}%
\TextAt{12}{79.226}{440.069}{Definition}%
\TextAt{12}{142.015}{440.069}{(12.5).}%
\TextAt{8}{184.061}{440.069}{The}%
\TextAt{8}{208.115}{440.069}{quotient}%
\TextAt{8}{254.298}{440.069}{group}%
\TextAt{13}{255.299}{411.068}{P}%
\TextAt{13}{264.469}{411.068}{ic}%
\TextAt{9}{273.502}{411.068}{\Character{40}}%
\TextAt{14}{278.054}{411.068}{o}%
\TextAt{9}{283.899}{411.068}{\Character{41}}%
\TextAt{9}{291.754}{411.068}{\Character{61}}%
\TextAt{13}{304.197}{411.068}{J}%
\TextAt{9}{311.774}{411.068}{\Character{40}}%
\TextAt{14}{316.327}{411.068}{o}%
\TextAt{9}{322.171}{411.068}{\Character{41}\Character{47}}%
\TextAt{13}{332.576}{411.068}{P}%
\TextAt{9}{341.747}{411.068}{\Character{40}}%
\TextAt{14}{346.3}{411.068}{o}%
\TextAt{9}{352.144}{411.068}{\Character{41}}%
\TextAt{8}{79.226}{382.067}{is}%
\TextAt{8}{90.99}{382.067}{called}%
\TextAt{8}{124.154}{382.067}{the}%
\TextAt{12}{144.31}{382.067}{Picard}%
\TextAt{12}{187.217}{382.067}{group}%
\TextAt{8}{225.019}{382.067}{of}%
\TextAt{8}{238.349}{382.067}{the}%
\TextAt{8}{258.506}{382.067}{ring}%
\TextAt{14}{282.572}{382.067}{o}%
\TextAt{8}{288.416}{382.067}{.}%
\TextAt{8}{97.159}{356.055}{In}%
\TextAt{8}{111.9}{356.055}{the}%
\TextAt{8}{132.177}{356.055}{case}%
\TextAt{8}{157.068}{356.055}{where}%
\TextAt{14}{190.997}{356.055}{o}%
\TextAt{8}{200.859}{356.055}{is}%
\TextAt{8}{212.743}{356.055}{a}%
\TextAt{8}{222.618}{356.055}{Dedekind}%
\TextAt{8}{274.922}{356.055}{domain,}%
\TextAt{8}{319.952}{356.055}{the}%
\TextAt{8}{340.229}{356.055}{Picard}%
\TextAt{8}{377.577}{356.055}{group}%
\TextAt{8}{410.873}{356.055}{is}%
\TextAt{8}{422.757}{356.055}{of}%
\TextAt{8}{436.207}{356.055}{course}%
\TextAt{8}{472.156}{356.055}{nothing}%
\TextAt{8}{515.207}{356.055}{but}%
\TextAt{8}{79.226}{339.009}{the}%
\TextAt{8}{100.635}{339.009}{ideal}%
\TextAt{8}{129.852}{339.009}{class}%
\TextAt{8}{158.542}{339.009}{group}%
\TextAt{13}{192.97}{339.009}{C}%
\TextAt{13}{202.188}{339.009}{l}%
\TextAt{21}{205.691}{337.215}{K}%
\TextAt{8}{213.893}{339.009}{.}%
\TextAt{8}{226.1}{339.009}{In}%
\TextAt{8}{241.974}{339.009}{general,}%
\TextAt{8}{287.117}{339.009}{we}%
\TextAt{8}{305.597}{339.009}{have}%
\TextAt{8}{333.845}{339.009}{the}%
\TextAt{8}{355.254}{339.009}{following}%
\TextAt{8}{405.942}{339.009}{description}%
\TextAt{8}{467.09}{339.009}{for}%
\TextAt{13}{486.228}{339.009}{J}%
\TextAt{9}{493.805}{339.009}{\Character{40}}%
\TextAt{14}{498.357}{339.009}{o}%
\TextAt{9}{504.202}{339.009}{\Character{41}}%
\TextAt{8}{513.904}{339.009}{and}%
\TextAt{13}{79.226}{321.963}{P}%
\TextAt{13}{88.396}{321.963}{ic}%
\TextAt{9}{97.429}{321.963}{\Character{40}}%
\TextAt{14}{101.981}{321.963}{o}%
\TextAt{9}{107.825}{321.963}{\Character{41}}%
\TextAt{8}{112.378}{321.963}{.}%
\TextAt{12}{79.226}{295.95}{Proposition}%
\TextAt{12}{151.937}{295.95}{(12.6).}%
\TextAt{8}{193.984}{295.95}{The}%
\TextAt{8}{218.038}{295.95}{correspondence}%
\TextAt{14}{299.715}{295.95}{a}%
\NativeArrowAt{309.019}{295.95}{11.955}{11.955}{\mapsto}%
\TextAt{9}{324.295}{295.95}{\Character{40}}%
\TextAt{14}{328.848}{295.95}{a}%
\TextAt{27}{334.831}{294.157}{p}%
\TextAt{9}{339.318}{295.95}{\Character{41}}%
\TextAt{9}{347.173}{295.95}{\Character{61}}%
\TextAt{9}{359.594}{295.95}{\Character{40}}%
\TextAt{14}{364.169}{295.95}{ao}%
\TextAt{27}{375.996}{294.157}{p}%
\TextAt{9}{380.483}{295.95}{\Character{41}}%
\TextAt{8}{388.932}{295.95}{yields}%
\TextAt{8}{421.833}{295.95}{an}%
\TextAt{8}{438.092}{295.95}{isomorphism}%
\TextAt{13}{260.806}{266.949}{J}%
\TextAt{9}{268.384}{266.949}{\Character{40}}%
\TextAt{14}{272.936}{266.949}{o}%
\TextAt{9}{278.78}{266.949}{\Character{41}}%
\NativeMathAt{286.653}{266.949}{11.955}{\cong}%
\TextAt{20}{299.273}{278.307}{\Character{77}}%
\TextAt{27}{306.311}{253.201}{p}%
\TextAt{13}{319.331}{266.949}{P}%
\TextAt{9}{328.502}{266.949}{\Character{40}}%
\TextAt{14}{333.054}{266.949}{o}%
\TextAt{27}{338.898}{265.156}{p}%
\TextAt{9}{343.385}{266.949}{\Character{41}}%
\TextAt{13}{347.938}{266.949}{\Character{58}}%
\TextAt{8}{79.226}{228.581}{Identifying}%
\TextAt{8}{138.404}{228.581}{the}%
\TextAt{8}{158.561}{228.581}{subgroup}%
\TextAt{13}{209.358}{228.581}{P}%
\TextAt{9}{218.529}{228.581}{\Character{40}}%
\TextAt{14}{223.082}{228.581}{o}%
\TextAt{9}{228.926}{228.581}{\Character{41}}%
\TextAt{8}{237.376}{228.581}{with}%
\TextAt{8}{264.036}{228.581}{its}%
\TextAt{8}{280.354}{228.581}{image}%
\TextAt{8}{314.176}{228.581}{in}%
\TextAt{8}{327.828}{228.581}{the}%
\TextAt{8}{347.985}{228.581}{direct}%
\TextAt{8}{381.148}{228.581}{sum}%
\TextAt{8}{405.92}{228.581}{one}%
\TextAt{8}{427.379}{228.581}{gets}%
\TextAt{13}{230.046}{191.719}{P}%
\TextAt{13}{239.216}{191.719}{ic}%
\TextAt{9}{248.249}{191.719}{\Character{40}}%
\TextAt{14}{252.801}{191.719}{o}%
\TextAt{9}{258.646}{191.719}{\Character{41}}%
\NativeMathAt{266.519}{191.719}{11.955}{\cong}%
\TextAt{20}{279.138}{212.162}{\Character{32}}%
\TextAt{20}{288.603}{203.076}{\Character{77}}%
\TextAt{27}{295.641}{177.97}{p}%
\TextAt{13}{308.661}{191.719}{P}%
\TextAt{9}{317.832}{191.719}{\Character{40}}%
\TextAt{14}{322.384}{191.719}{o}%
\TextAt{27}{328.228}{189.925}{p}%
\TextAt{9}{332.715}{191.719}{\Character{41}}%
\TextAt{20}{337.267}{212.162}{\Character{33}}%
\TextAt{9}{348.725}{191.719}{\Character{47}}%
\TextAt{13}{354.578}{191.719}{P}%
\TextAt{9}{363.749}{191.719}{\Character{40}}%
\TextAt{14}{368.301}{191.719}{o}%
\TextAt{9}{374.145}{191.719}{\Character{41}}%
\TextAt{13}{378.697}{191.719}{\Character{58}}%
\TextAt{19}{79.226}{135.417}{Proof.}%
\TextAt{8}{116}{135.417}{For}%
\TextAt{8}{136.616}{135.417}{every}%
\TextAt{14}{167.154}{135.417}{a}%
\NativeMathAt{176.458}{135.417}{11.955}{\in}%
\TextAt{13}{187.749}{135.417}{J}%
\TextAt{9}{195.326}{135.417}{\Character{40}}%
\TextAt{14}{199.879}{135.417}{o}%
\TextAt{9}{205.723}{135.417}{\Character{41}}%
\TextAt{8}{210.275}{135.417}{,}%
\TextAt{14}{217.141}{135.417}{a}%
\TextAt{27}{223.125}{133.624}{p}%
\TextAt{9}{230.932}{135.417}{\Character{61}}%
\TextAt{14}{243.358}{135.417}{ao}%
\TextAt{27}{255.185}{133.624}{p}%
\TextAt{8}{263.216}{135.417}{is}%
\TextAt{8}{274.626}{135.417}{a}%
\TextAt{8}{284.027}{135.417}{principal}%
\TextAt{8}{332.451}{135.417}{ideal}%
\TextAt{8}{360.06}{135.417}{by}%
\TextAt{8}{375.965}{135.417}{(12.4),}%
\TextAt{8}{412.767}{135.417}{and}%
\TextAt{8}{435.176}{135.417}{we}%
\TextAt{8}{452.05}{135.417}{have}%
\TextAt{14}{478.691}{135.417}{a}%
\TextAt{27}{484.674}{133.624}{p}%
\TextAt{9}{492.482}{135.417}{\Character{61}}%
\TextAt{14}{504.907}{135.417}{o}%
\TextAt{27}{510.752}{133.624}{p}%
\TextAt{8}{518.782}{135.417}{for}%
\ErrataAt{79.19}{118.371}{455.81}{almost all $\mathfrak p$. Choose $a,c\in\mathfrak o\setminus\{0\}$ with $a\mathfrak o\subseteq\mathfrak a\subseteq c^{-1}\mathfrak o$. The primes dividing $ac$}%
\ErrataAt{79.19}{101.325}{455.81}{can be exceptions. We therefore obtain a homomorphism}%
\TextAt{13}{215.518}{84.279}{J}%
\TextAt{9}{223.095}{84.279}{\Character{40}}%
\TextAt{14}{227.647}{84.279}{o}%
\TextAt{9}{233.491}{84.279}{\Character{41}}%
\NativeArrowAt{241.364}{84.279}{11.955}{19.26}{\longrightarrow}%
\TextAt{20}{263.946}{95.637}{\Character{77}}%
\TextAt{27}{270.985}{70.531}{p}%
\TextAt{13}{284.005}{84.279}{P}%
\TextAt{9}{293.176}{84.279}{\Character{40}}%
\TextAt{14}{297.728}{84.279}{o}%
\TextAt{27}{303.572}{82.486}{p}%
\TextAt{9}{308.059}{84.279}{\Character{41}}%
\TextAt{13}{312.611}{84.279}{\Character{59}}%
\TextAt{14}{341.766}{84.279}{a}%
\NativeArrowAt{351.07}{84.279}{11.955}{19.26}{\longmapsto}%
\TextAt{9}{373.652}{84.279}{\Character{40}}%
\TextAt{14}{378.204}{84.279}{a}%
\TextAt{27}{384.187}{82.486}{p}%
\TextAt{9}{388.674}{84.279}{\Character{41}}%
\TextAt{13}{393.226}{84.279}{\Character{58}}%
\TextAt{8}{300.14}{43.033}{70}%
\Destination{78.23}{809.815}{nav.page.70}%
\Destination{79.36}{683.555}{nav.print.75}%
\Destination{79.36}{775.155}{nav.number.1.12.4}%
\Destination{79.36}{454.055}{nav.number.1.12.5}%
\Destination{79.36}{309.955}{nav.number.1.12.6}%
\Hotlink{375.465}{133.255}{30.936}{11.047}{nav.number.1.12.4}%
\Destination{79.226}{135}{errata.190}%
}

% Page 83 (body page 71)
\Page{%
\TextAt{2}{79.226}{796.943}{I.}%
\TextAt{2}{89.83}{796.943}{Algebraic}%
\TextAt{2}{138.648}{796.943}{Integers}%
\Vector{q 0 G 0 g 0.299 w 7.226 18.729 m 460.769 18.729 l S Q}%
\TextAt{8}{79.226}{761.124}{It}%
\TextAt{8}{92.424}{761.124}{is}%
\TextAt{8}{104.712}{761.124}{injective,}%
\TextAt{8}{155.114}{761.124}{for}%
\TextAt{8}{173.524}{761.124}{if}%
\TextAt{14}{184.772}{761.124}{a}%
\TextAt{27}{190.755}{759.33}{p}%
\TextAt{9}{199.458}{761.124}{\Character{61}}%
\TextAt{14}{212.778}{761.124}{o}%
\TextAt{27}{218.622}{759.33}{p}%
\TextAt{8}{227.531}{761.124}{for}%
\TextAt{8}{245.941}{761.124}{all}%
\TextAt{14}{262.725}{761.124}{p}%
\TextAt{8}{268.708}{761.124}{,}%
\TextAt{8}{276.513}{761.124}{then}%
\TextAt{14}{303.698}{761.124}{a}%
\NativeMathAt{313.896}{761.124}{11.955}{\subseteq}%
\TextAt{20}{327.41}{770.09}{\Character{84}}%
\TextAt{27}{337.373}{757.637}{p}%
\TextAt{14}{343.852}{761.124}{o}%
\TextAt{27}{349.696}{759.33}{p}%
\TextAt{9}{358.399}{761.124}{\Character{61}}%
\TextAt{14}{371.718}{761.124}{o}%
\TextAt{8}{381.985}{761.124}{(see}%
\TextAt{8}{405.978}{761.124}{the}%
\TextAt{8}{426.659}{761.124}{proof}%
\TextAt{8}{457.753}{761.124}{of}%
\TextAt{8}{471.608}{761.124}{(11.5)),}%
\TextAt{8}{513.904}{761.124}{and}%
\TextAt{8}{79.226}{744.078}{one}%
\TextAt{8}{101.121}{744.078}{has}%
\TextAt{8}{122.43}{744.078}{to}%
\TextAt{8}{137.175}{744.078}{have}%
\TextAt{14}{164.605}{744.078}{a}%
\TextAt{9}{174.651}{744.078}{\Character{61}}%
\TextAt{14}{187.818}{744.078}{o}%
\TextAt{8}{197.995}{744.078}{because}%
\TextAt{8}{241.732}{744.078}{otherwise}%
\TextAt{8}{294.256}{744.078}{there}%
\TextAt{8}{324.603}{744.078}{would}%
\TextAt{8}{359.182}{744.078}{exist}%
\TextAt{8}{387.317}{744.078}{a}%
\TextAt{8}{397.508}{744.078}{maximal}%
\TextAt{8}{445.752}{744.078}{ideal}%
\TextAt{14}{474.15}{744.078}{p}%
\TextAt{8}{484.466}{744.078}{such}%
\TextAt{8}{511.298}{744.078}{that}%
\TextAt{14}{79.226}{727.032}{a}%
\NativeMathAt{89.5}{727.032}{11.955}{\subseteq}%
\TextAt{14}{103.089}{727.032}{p}%
\NativeMathAt{113.362}{727.032}{11.955}{\subset}%
\TextAt{14}{126.951}{727.032}{o}%
\TextAt{8}{132.795}{727.032}{,}%
\TextAt{8}{140.656}{727.032}{i.e.,}%
\TextAt{14}{163.472}{727.032}{a}%
\TextAt{27}{169.455}{725.238}{p}%
\NativeMathAt{178.232}{727.032}{11.955}{\subseteq}%
\TextAt{14}{191.821}{727.032}{po}%
\TextAt{27}{203.648}{725.238}{p}%
\NativeMathAt{212.426}{727.032}{11.955}{\neq}%
\TextAt{14}{225.821}{727.032}{o}%
\TextAt{27}{231.665}{725.238}{p}%
\TextAt{8}{236.152}{727.032}{.}%
\TextAt{8}{246.307}{727.032}{In}%
\TextAt{8}{261.497}{727.032}{order}%
\TextAt{8}{292.635}{727.032}{to}%
\TextAt{8}{307.515}{727.032}{prove}%
\TextAt{8}{339.633}{727.032}{surjectivity,}%
\TextAt{8}{403.815}{727.032}{let}%
\TextAt{9}{421.288}{727.032}{\Character{40}}%
\TextAt{13}{425.84}{727.032}{a}%
\TextAt{27}{431.985}{725.238}{p}%
\TextAt{14}{436.472}{727.032}{o}%
\TextAt{27}{442.316}{725.238}{p}%
\TextAt{9}{446.803}{727.032}{\Character{41}}%
\NativeMathAt{455.646}{727.032}{11.955}{\in}%
\TextAt{20}{467.907}{735.998}{\Character{76}}%
\TextAt{27}{481.19}{723.545}{p}%
\TextAt{13}{487.669}{727.032}{P}%
\TextAt{9}{496.841}{727.032}{\Character{40}}%
\TextAt{14}{501.393}{727.032}{o}%
\TextAt{27}{507.237}{725.238}{p}%
\TextAt{9}{511.724}{727.032}{\Character{41}}%
\TextAt{8}{520.743}{727.032}{be}%
\TextAt{8}{79.226}{709.986}{given.}%
\TextAt{8}{114.346}{709.986}{Then}%
\TextAt{8}{144.904}{709.986}{the}%
\TextAt{14}{165.06}{709.986}{o}%
\TextAt{8}{170.905}{709.986}{-submodule}%
\TextAt{14}{277.014}{692.94}{a}%
\TextAt{9}{286.318}{692.94}{\Character{61}}%
\TextAt{20}{298.743}{704.297}{\Character{92}}%
\TextAt{27}{303.39}{679.191}{p}%
\TextAt{13}{314.019}{692.94}{a}%
\TextAt{27}{320.164}{691.147}{p}%
\TextAt{14}{324.651}{692.94}{o}%
\TextAt{27}{330.495}{691.147}{p}%
\TextAt{8}{79.226}{661.563}{of}%
\TextAt{13}{93.214}{661.563}{K}%
\TextAt{8}{108.587}{661.563}{is}%
\TextAt{8}{121.008}{661.563}{a}%
\TextAt{8}{131.421}{661.563}{fractional}%
\TextAt{8}{184.443}{661.563}{ideal.}%
\TextAt{8}{218.93}{661.563}{Indeed,}%
\TextAt{8}{261.034}{661.563}{since}%
\TextAt{13}{290.36}{661.563}{a}%
\TextAt{27}{296.505}{659.77}{p}%
\TextAt{14}{300.992}{661.563}{o}%
\TextAt{27}{306.836}{659.77}{p}%
\TextAt{9}{315.764}{661.563}{\Character{61}}%
\TextAt{14}{329.311}{661.563}{o}%
\TextAt{27}{335.155}{659.77}{p}%
\TextAt{8}{344.197}{661.563}{for}%
\TextAt{8}{362.74}{661.563}{almost}%
\TextAt{8}{401.188}{661.563}{all}%
\TextAt{14}{418.104}{661.563}{p}%
\TextAt{8}{424.087}{661.563}{,}%
\TextAt{8}{432.059}{661.563}{there}%
\TextAt{8}{462.628}{661.563}{is}%
\TextAt{8}{475.05}{661.563}{some}%
\TextAt{13}{505.034}{661.563}{c}%
\NativeMathAt{514.513}{661.563}{11.955}{\in}%
\TextAt{14}{526.925}{661.563}{o}%
\TextAt{8}{79.226}{644.517}{such}%
\TextAt{8}{105.623}{644.517}{that}%
\TextAt{13}{130.992}{644.517}{ca}%
\TextAt{27}{142.175}{642.724}{p}%
\NativeMathAt{149.983}{644.517}{11.955}{\in}%
\TextAt{14}{161.274}{644.517}{o}%
\TextAt{27}{167.118}{642.724}{p}%
\TextAt{8}{175.502}{644.517}{for}%
\TextAt{8}{193.387}{644.517}{all}%
\TextAt{14}{209.646}{644.517}{p}%
\TextAt{8}{215.629}{644.517}{,}%
\TextAt{8}{222.778}{644.517}{i.e.,}%
\TextAt{13}{244.883}{644.517}{c}%
\TextAt{14}{249.921}{644.517}{a}%
\NativeMathAt{259.225}{644.517}{11.955}{\subseteq}%
\TextAt{20}{271.845}{653.484}{\Character{84}}%
\TextAt{27}{281.807}{641.03}{p}%
\TextAt{14}{288.287}{644.517}{o}%
\TextAt{27}{294.131}{642.724}{p}%
\TextAt{9}{301.939}{644.517}{\Character{61}}%
\TextAt{14}{314.364}{644.517}{o}%
\TextAt{8}{320.208}{644.517}{.}%
\TextAt{8}{328.657}{644.517}{We}%
\TextAt{8}{348.801}{644.517}{have}%
\TextAt{8}{375.796}{644.517}{to}%
\TextAt{8}{390.106}{644.517}{show}%
\TextAt{8}{419.109}{644.517}{that}%
\TextAt{8}{444.478}{644.517}{one}%
\TextAt{8}{465.938}{644.517}{has}%
\TextAt{14}{279.486}{617.763}{ao}%
\TextAt{27}{291.313}{615.97}{p}%
\TextAt{9}{299.121}{617.763}{\Character{61}}%
\TextAt{13}{311.547}{617.763}{a}%
\TextAt{27}{317.692}{615.97}{p}%
\TextAt{14}{322.178}{617.763}{o}%
\TextAt{27}{328.023}{615.97}{p}%
\TextAt{8}{79.226}{591.009}{for}%
\TextAt{8}{97.692}{591.009}{every}%
\TextAt{14}{129.165}{591.009}{p}%
\TextAt{8}{135.148}{591.009}{.}%
\TextAt{8}{145.338}{591.009}{The}%
\TextAt{8}{169.972}{591.009}{inclusion}%
\NativeMathAt{219.39}{591.009}{11.955}{\subseteq}%
\TextAt{8}{233.166}{591.009}{is}%
\TextAt{8}{245.511}{591.009}{trivial.}%
\TextAt{8}{286.605}{591.009}{In}%
\TextAt{8}{301.807}{591.009}{order}%
\TextAt{8}{332.957}{591.009}{to}%
\TextAt{8}{347.848}{591.009}{show}%
\TextAt{8}{377.432}{591.009}{that}%
\TextAt{13}{403.381}{591.009}{a}%
\TextAt{27}{409.526}{589.215}{p}%
\TextAt{14}{414.013}{591.009}{o}%
\TextAt{27}{419.857}{589.215}{p}%
\NativeMathAt{428.655}{591.009}{11.955}{\subseteq}%
\TextAt{14}{442.263}{591.009}{ao}%
\TextAt{27}{454.091}{589.215}{p}%
\TextAt{8}{458.578}{591.009}{,}%
\TextAt{8}{466.453}{591.009}{let}%
\TextAt{8}{483.938}{591.009}{us}%
\TextAt{8}{499.534}{591.009}{choose}%
\TextAt{13}{79.226}{573.963}{c}%
\NativeMathAt{88.354}{573.963}{11.955}{\in}%
\TextAt{14}{100.414}{573.963}{o}%
\TextAt{8}{106.258}{573.963}{,}%
\TextAt{13}{113.971}{573.963}{c}%
\NativeMathAt{123.099}{573.963}{11.955}{\neq}%
\TextAt{9}{136.286}{573.963}{0}%
\TextAt{8}{142.146}{573.963}{,}%
\TextAt{8}{149.859}{573.963}{such}%
\TextAt{8}{176.707}{573.963}{that}%
\TextAt{13}{202.527}{573.963}{ca}%
\TextAt{3}{213.71}{578.301}{\Character{0}}%
\TextAt{17}{220.297}{578.301}{1}%
\TextAt{27}{213.71}{571.007}{p}%
\TextAt{13}{225.029}{573.963}{a}%
\TextAt{27}{231.174}{572.169}{q}%
\NativeMathAt{239.66}{573.963}{11.955}{\in}%
\TextAt{14}{251.72}{573.963}{o}%
\TextAt{8}{261.912}{573.963}{for}%
\TextAt{8}{280.248}{573.963}{the}%
\TextAt{8}{300.856}{573.963}{finitely}%
\TextAt{8}{340.652}{573.963}{many}%
\TextAt{14}{372.975}{573.963}{q}%
\TextAt{8}{383.171}{573.963}{which}%
\TextAt{8}{417.109}{573.963}{satisfy}%
\TextAt{13}{454.107}{573.963}{a}%
\TextAt{3}{460.252}{578.301}{\Character{0}}%
\TextAt{17}{466.838}{578.301}{1}%
\TextAt{27}{460.252}{571.007}{p}%
\TextAt{13}{471.571}{573.963}{a}%
\TextAt{27}{477.716}{572.169}{q}%
\NativeMathAt{486.202}{573.963}{11.955}{\notin}%
\TextAt{14}{498.261}{573.963}{o}%
\TextAt{27}{504.106}{572.169}{q}%
\TextAt{8}{508.502}{573.963}{.}%
\TextAt{8}{518.304}{573.963}{By}%
\TextAt{8}{79.226}{556.917}{the}%
\TextAt{8}{99.383}{556.917}{Chinese}%
\TextAt{8}{143.007}{556.917}{remainder}%
\TextAt{8}{198.288}{556.917}{theorem}%
\TextAt{8}{243.813}{556.917}{(12.3),}%
\TextAt{8}{280.898}{556.917}{we}%
\TextAt{8}{298.125}{556.917}{may}%
\TextAt{8}{323.494}{556.917}{find}%
\TextAt{13}{346.902}{556.917}{a}%
\NativeMathAt{356.368}{556.917}{11.955}{\in}%
\TextAt{14}{367.659}{556.917}{o}%
\TextAt{8}{377.401}{556.917}{such}%
\TextAt{8}{403.798}{556.917}{that}%
\ErrataAt{182}{530.176}{318}{$a-c\in c\mathfrak p\mathfrak o_{\mathfrak p}\quad\text{and}\quad a\in ca_{\mathfrak p}^{-1}a_{\mathfrak q}\mathfrak o_{\mathfrak q}\text{ for }\mathfrak q\ne\mathfrak p.$}%
\TextAt{8}{79.226}{503.408}{Then}%
\TextAt{13}{109.784}{503.408}{\Character{34}}%
\TextAt{9}{118.582}{503.408}{\Character{61}}%
\TextAt{13}{131.008}{503.408}{ac}%
\TextAt{3}{142.191}{507.746}{\Character{0}}%
\TextAt{17}{148.777}{507.746}{1}%
\TextAt{8}{157.407}{503.408}{is}%
\TextAt{8}{169.171}{503.408}{a}%
\TextAt{8}{178.926}{503.408}{unit}%
\TextAt{8}{203.638}{503.408}{in}%
\TextAt{14}{217.29}{503.408}{o}%
\TextAt{27}{223.135}{501.615}{p}%
\TextAt{8}{242.482}{503.408}{and}%
\TextAt{13}{265.244}{503.408}{a}%
\TextAt{27}{271.389}{501.615}{p}%
\TextAt{13}{275.876}{503.408}{\Character{34}}%
\NativeMathAt{284.675}{503.408}{11.955}{\in}%
\TextAt{20}{295.966}{512.374}{\Character{84}}%
\TextAt{27}{305.929}{499.921}{q}%
\TextAt{13}{312.318}{503.408}{a}%
\TextAt{27}{318.463}{501.615}{q}%
\TextAt{14}{322.859}{503.408}{o}%
\TextAt{27}{328.703}{501.615}{q}%
\TextAt{9}{336.421}{503.408}{\Character{61}}%
\TextAt{14}{348.846}{503.408}{a}%
\TextAt{8}{354.829}{503.408}{,}%
\TextAt{8}{361.978}{503.408}{hence}%
\TextAt{13}{252.117}{476.654}{a}%
\TextAt{27}{258.262}{474.86}{p}%
\TextAt{14}{262.749}{476.654}{o}%
\TextAt{27}{268.593}{474.86}{p}%
\TextAt{9}{276.401}{476.654}{\Character{61}}%
\TextAt{9}{288.81}{476.654}{\Character{40}}%
\TextAt{13}{293.379}{476.654}{a}%
\TextAt{27}{299.524}{474.86}{p}%
\TextAt{13}{304.011}{476.654}{\Character{34}}%
\TextAt{9}{309.489}{476.654}{\Character{41}}%
\TextAt{14}{314.041}{476.654}{o}%
\TextAt{27}{319.885}{474.86}{p}%
\NativeMathAt{327.693}{476.654}{11.955}{\subseteq}%
\TextAt{14}{340.312}{476.654}{ao}%
\TextAt{27}{352.14}{474.86}{p}%
\TextAt{13}{356.627}{476.654}{\Character{58}}%
\Vector{q 0 G 0 g 0.399 w 452.334 -324.101 m 452.334 -316.031 l S Q}%
\Vector{q 0 G 0 g 0.399 w 452.534 -316.23 m 459.707 -316.23 l S Q}%
\Vector{q 0 G 0 g 0.399 w 452.534 -323.902 m 459.707 -323.902 l S Q}%
\Vector{q 0 G 0 g 0.399 w 459.906 -324.101 m 459.906 -316.031 l S Q}%
\TextAt{8}{97.159}{421.54}{Passing}%
\TextAt{8}{139.069}{421.54}{from}%
\TextAt{8}{166.381}{421.54}{the}%
\TextAt{8}{186.209}{421.54}{ring}%
\TextAt{14}{209.947}{421.54}{o}%
\TextAt{8}{219.361}{421.54}{to}%
\TextAt{8}{233.343}{421.54}{its}%
\TextAt{12}{249.334}{421.54}{normalization}%
\TextAt{20}{333.142}{422.075}{\Character{101}}%
\TextAt{14}{333.541}{421.54}{o}%
\TextAt{8}{339.385}{421.54}{,}%
\TextAt{8}{346.272}{421.54}{i.e.,}%
\TextAt{8}{368.115}{421.54}{to}%
\TextAt{8}{382.097}{421.54}{the}%
\TextAt{8}{401.926}{421.54}{integral}%
\TextAt{8}{444.206}{421.54}{closure}%
\TextAt{8}{482.959}{421.54}{of}%
\TextAt{14}{495.961}{421.54}{o}%
\TextAt{8}{505.375}{421.54}{in}%
\TextAt{13}{518.7}{421.54}{K}%
\TextAt{8}{529.518}{421.54}{,}%
\TextAt{8}{79.226}{404.494}{one}%
\TextAt{8}{100.558}{404.494}{obtains}%
\TextAt{8}{141.472}{404.494}{a}%
\TextAt{8}{151.1}{404.494}{Dedekind}%
\TextAt{8}{203.157}{404.494}{domain.}%
\TextAt{8}{249.293}{404.494}{This}%
\TextAt{8}{275.885}{404.494}{is}%
\TextAt{8}{287.521}{404.494}{not}%
\TextAt{8}{308.208}{404.494}{all}%
\TextAt{8}{324.339}{404.494}{that}%
\TextAt{8}{349.58}{404.494}{easy}%
\TextAt{8}{375.204}{404.494}{to}%
\TextAt{8}{389.386}{404.494}{prove,}%
\TextAt{8}{424.086}{404.494}{however,}%
\TextAt{8}{472.115}{404.494}{because}%
\TextAt{20}{514.89}{405.029}{\Character{101}}%
\TextAt{14}{515.289}{404.494}{o}%
\TextAt{8}{524.903}{404.494}{is}%
\TextAt{8}{79.226}{387.448}{in}%
\TextAt{8}{92.879}{387.448}{general}%
\TextAt{8}{133.204}{387.448}{not}%
\TextAt{8}{154.018}{387.448}{a}%
\TextAt{8}{163.773}{387.448}{finitely}%
\TextAt{8}{203.118}{387.448}{generated}%
\TextAt{14}{256.45}{387.448}{o}%
\TextAt{8}{262.294}{387.448}{-module.}%
\TextAt{8}{312.035}{387.448}{But}%
\TextAt{8}{335.276}{387.448}{at}%
\TextAt{8}{349.586}{387.448}{any}%
\TextAt{8}{371.704}{387.448}{rate}%
\TextAt{8}{395.769}{387.448}{we}%
\TextAt{8}{412.997}{387.448}{have}%
\TextAt{8}{439.991}{387.448}{the}%
\TextAt{12}{79.226}{362.078}{Lemma}%
\TextAt{12}{127.718}{362.078}{(12.7).}%
\TextAt{8}{170.483}{362.078}{Let}%
\TextAt{14}{192.171}{362.078}{o}%
\TextAt{8}{202.631}{362.078}{be}%
\TextAt{8}{219.273}{362.078}{a}%
\TextAt{8}{229.747}{362.078}{one-dimensional}%
\TextAt{8}{316.376}{362.078}{noetherian}%
\TextAt{8}{375.304}{362.078}{integral}%
\TextAt{8}{418.631}{362.078}{domain}%
\TextAt{8}{460.977}{362.078}{and}%
\TextAt{20}{484.059}{362.613}{\Character{101}}%
\TextAt{14}{484.458}{362.078}{o}%
\TextAt{8}{494.918}{362.078}{its}%
\TextAt{8}{511.956}{362.078}{nor-}%
\TextAt{8}{79.226}{345.032}{malization.}%
\TextAt{8}{144.619}{345.032}{Then,}%
\TextAt{8}{179.928}{345.032}{for}%
\TextAt{8}{199.014}{345.032}{each}%
\TextAt{8}{226.551}{345.032}{ideal}%
\TextAt{14}{255.714}{345.032}{a}%
\NativeMathAt{267.064}{345.032}{11.955}{\neq}%
\TextAt{9}{281.518}{345.032}{0}%
\TextAt{8}{292.485}{345.032}{of}%
\TextAt{14}{307.015}{345.032}{o}%
\TextAt{8}{312.859}{345.032}{,}%
\TextAt{8}{321.509}{345.032}{the}%
\TextAt{8}{342.865}{345.032}{quotient}%
\TextAt{20}{389.85}{345.567}{\Character{101}}%
\TextAt{14}{390.248}{345.032}{o}%
\TextAt{9}{396.093}{345.032}{\Character{47}}%
\TextAt{14}{401.946}{345.032}{a}%
\TextAt{20}{407.53}{345.567}{\Character{101}}%
\TextAt{14}{407.929}{345.032}{o}%
\TextAt{8}{418.871}{345.032}{is}%
\TextAt{8}{431.835}{345.032}{a}%
\TextAt{8}{442.79}{345.032}{finitely}%
\TextAt{8}{483.335}{345.032}{generated}%
\TextAt{14}{79.226}{327.986}{o}%
\TextAt{8}{85.07}{327.986}{-module.}%
\TextAt{19}{79.226}{294.291}{Proof.}%
\TextAt{8}{116}{294.291}{Let}%
\TextAt{13}{137.414}{294.291}{a}%
\NativeMathAt{147.636}{294.291}{11.955}{\in}%
\TextAt{14}{159.684}{294.291}{a}%
\TextAt{8}{165.667}{294.291}{,}%
\TextAt{13}{173.372}{294.291}{a}%
\NativeMathAt{183.594}{294.291}{11.955}{\neq}%
\TextAt{9}{196.769}{294.291}{0}%
\TextAt{8}{202.629}{294.291}{.}%
\TextAt{8}{212.41}{294.291}{Then}%
\TextAt{20}{243.013}{294.826}{\Character{101}}%
\TextAt{14}{243.412}{294.291}{o}%
\TextAt{9}{249.256}{294.291}{\Character{47}}%
\TextAt{14}{255.109}{294.291}{a}%
\TextAt{20}{260.693}{294.826}{\Character{101}}%
\TextAt{14}{261.092}{294.291}{o}%
\TextAt{8}{271.278}{294.291}{is}%
\TextAt{8}{283.486}{294.291}{a}%
\TextAt{8}{293.685}{294.291}{quotient}%
\TextAt{8}{340.312}{294.291}{of}%
\TextAt{20}{353.687}{294.826}{\Character{101}}%
\TextAt{14}{354.086}{294.291}{o}%
\TextAt{9}{359.93}{294.291}{\Character{47}}%
\TextAt{13}{365.783}{294.291}{a}%
\TextAt{20}{371.53}{294.826}{\Character{101}}%
\TextAt{14}{371.928}{294.291}{o}%
\TextAt{8}{377.773}{294.291}{.}%
\TextAt{8}{387.553}{294.291}{It}%
\TextAt{8}{400.67}{294.291}{thus}%
\TextAt{8}{426.865}{294.291}{suffices}%
\TextAt{8}{467.096}{294.291}{to}%
\TextAt{8}{481.851}{294.291}{show}%
\TextAt{8}{511.298}{294.291}{that}%
\TextAt{20}{78.827}{277.78}{\Character{101}}%
\TextAt{14}{79.226}{277.245}{o}%
\TextAt{9}{85.07}{277.245}{\Character{47}}%
\TextAt{13}{90.923}{277.245}{a}%
\TextAt{20}{96.67}{277.78}{\Character{101}}%
\TextAt{14}{97.068}{277.245}{o}%
\TextAt{8}{106.957}{277.245}{is}%
\TextAt{8}{118.869}{277.245}{a}%
\TextAt{8}{128.771}{277.245}{finitely}%
\TextAt{8}{168.263}{277.245}{generated}%
\TextAt{14}{221.743}{277.245}{o}%
\TextAt{8}{227.587}{277.245}{-module.}%
\TextAt{8}{277.771}{277.245}{With}%
\TextAt{8}{308.153}{277.245}{this}%
\TextAt{8}{331.123}{277.245}{end,}%
\TextAt{8}{356.664}{277.245}{consider}%
\TextAt{8}{402.396}{277.245}{in}%
\TextAt{14}{416.197}{277.245}{o}%
\TextAt{8}{426.086}{277.245}{the}%
\TextAt{8}{446.389}{277.245}{descending}%
\TextAt{8}{505.775}{277.245}{chain}%
\TextAt{8}{79.226}{260.199}{of}%
\TextAt{8}{92.556}{260.199}{ideals}%
\TextAt{8}{125.134}{260.199}{containing}%
\TextAt{13}{182.053}{260.199}{a}%
\TextAt{14}{188.197}{260.199}{o}%
\TextAt{14}{256.791}{243.153}{a}%
\TextAt{21}{262.774}{241.36}{m}%
\TextAt{9}{274.084}{243.153}{\Character{61}}%
\TextAt{9}{286.494}{243.153}{\Character{40}}%
\TextAt{13}{291.062}{243.153}{a}%
\TextAt{21}{297.207}{248.089}{m}%
\TextAt{20}{304.796}{243.688}{\Character{101}}%
\TextAt{14}{305.195}{243.153}{o}%
\TextAt{16}{313.696}{243.153}{\Character{92}}%
\TextAt{14}{324.323}{243.153}{o}%
\TextAt{13}{330.167}{243.153}{\Character{59}}%
\TextAt{13}{335.391}{243.153}{a}%
\TextAt{14}{341.556}{243.153}{o}%
\TextAt{9}{347.4}{243.153}{\Character{41}}%
\TextAt{13}{351.953}{243.153}{\Character{58}}%
\TextAt{8}{79.226}{220.594}{This}%
\TextAt{8}{106.874}{220.594}{chain}%
\TextAt{8}{138.694}{220.594}{becomes}%
\TextAt{8}{186.176}{220.594}{stationary.}%
\TextAt{8}{248.218}{220.594}{In}%
\TextAt{8}{263.767}{220.594}{fact,}%
\TextAt{8}{291.264}{220.594}{the}%
\TextAt{8}{312.349}{220.594}{prime}%
\TextAt{8}{346.44}{220.594}{ideals}%
\TextAt{8}{379.946}{220.594}{of}%
\TextAt{8}{394.204}{220.594}{the}%
\TextAt{8}{415.289}{220.594}{ring}%
\TextAt{14}{440.282}{220.594}{o}%
\TextAt{9}{446.127}{220.594}{\Character{47}}%
\TextAt{13}{451.98}{220.594}{a}%
\TextAt{14}{458.125}{220.594}{o}%
\TextAt{8}{468.794}{220.594}{are}%
\TextAt{8}{489.233}{220.594}{not}%
\TextAt{8}{510.975}{220.594}{only}%
\TextAt{8}{79.226}{203.548}{maximal}%
\TextAt{8}{127.509}{203.548}{but}%
\TextAt{8}{149.443}{203.548}{also}%
\TextAt{8}{173.397}{203.548}{minimal}%
\TextAt{8}{219.397}{203.548}{in}%
\TextAt{8}{233.524}{203.548}{the}%
\TextAt{8}{254.155}{203.548}{sense}%
\TextAt{8}{284.66}{203.548}{that}%
\TextAt{14}{310.503}{203.548}{o}%
\TextAt{9}{316.348}{203.548}{\Character{47}}%
\TextAt{13}{322.201}{203.548}{a}%
\TextAt{14}{328.346}{203.548}{o}%
\TextAt{8}{338.561}{203.548}{is}%
\TextAt{8}{350.8}{203.548}{a}%
\TextAt{8}{361.029}{203.548}{zero-dimensional}%
\TextAt{8}{450.665}{203.548}{noetherian}%
\TextAt{8}{509.349}{203.548}{ring.}%
\TextAt{8}{79.226}{186.502}{In}%
\TextAt{8}{93.665}{186.502}{such}%
\TextAt{8}{119.879}{186.502}{a}%
\TextAt{8}{129.452}{186.502}{ring}%
\TextAt{8}{153.336}{186.502}{every}%
\TextAt{8}{184.046}{186.502}{descending}%
\TextAt{8}{243.101}{186.502}{chain}%
\TextAt{8}{273.811}{186.502}{of}%
\TextAt{8}{286.958}{186.502}{ideals}%
\TextAt{8}{319.354}{186.502}{becomes}%
\TextAt{8}{365.725}{186.502}{stationary}%
\TextAt{8}{421.229}{186.502}{(see}%
\TextAt{8}{444.515}{186.502}{§3,}%
\TextAt{8}{463.055}{186.502}{exercise}%
\TextAt{8}{506.174}{186.502}{7).}%
\TextAt{8}{524.975}{186.502}{If}%
\TextAt{8}{79.226}{169.456}{the}%
\TextAt{8}{99.394}{169.456}{chain}%
\Vector{q 0 G 0 g 0.478 w 58.297 -597.187 m 64.28 -597.187 l S Q}%
\TextAt{14}{130.297}{169.456}{a}%
\TextAt{21}{136.28}{167.663}{m}%
\TextAt{9}{147.608}{169.456}{\Character{61}}%
\TextAt{14}{160.053}{169.456}{a}%
\TextAt{21}{166.036}{167.663}{m}%
\TextAt{8}{177.345}{169.456}{mod}%
\TextAt{13}{202.783}{169.456}{a}%
\TextAt{14}{208.928}{169.456}{o}%
\TextAt{8}{218.681}{169.456}{is}%
\TextAt{8}{230.456}{169.456}{stationary}%
\TextAt{8}{286.154}{169.456}{at}%
\TextAt{13}{300.476}{169.456}{n}%
\TextAt{8}{307.463}{169.456}{,}%
\TextAt{8}{314.626}{169.456}{then}%
\TextAt{8}{341.297}{169.456}{so}%
\TextAt{8}{355.678}{169.456}{is}%
\TextAt{8}{367.453}{169.456}{the}%
\TextAt{8}{387.621}{169.456}{chain}%
\TextAt{14}{418.524}{169.456}{a}%
\TextAt{21}{424.507}{167.663}{m}%
\TextAt{8}{432.496}{169.456}{.}%
\TextAt{8}{440.978}{169.456}{We}%
\TextAt{8}{461.133}{169.456}{show}%
\TextAt{8}{490.147}{169.456}{that,}%
\TextAt{8}{518.782}{169.456}{for}%
\TextAt{8}{79.226}{152.41}{this}%
\TextAt{13}{102.049}{152.41}{n}%
\TextAt{8}{109.036}{152.41}{,}%
\TextAt{8}{116.185}{152.41}{we}%
\TextAt{8}{133.413}{152.41}{have}%
\TextAt{20}{267.771}{135.9}{\Character{101}}%
\TextAt{14}{268.17}{135.364}{o}%
\NativeMathAt{277.335}{135.364}{11.955}{\subseteq}%
\TextAt{13}{289.955}{135.364}{a}%
\TextAt{3}{296.1}{140.3}{\Character{0}}%
\TextAt{21}{302.686}{140.3}{n}%
\TextAt{14}{308.322}{135.364}{o}%
\TextAt{9}{316.823}{135.364}{\Character{43}}%
\TextAt{13}{328.585}{135.364}{a}%
\TextAt{20}{334.331}{135.9}{\Character{101}}%
\TextAt{14}{334.73}{135.364}{o}%
\TextAt{13}{340.574}{135.364}{\Character{58}}%
\TextAt{8}{300.14}{43.033}{71}%
\Destination{78.23}{809.815}{nav.page.71}%
\Destination{79.36}{776.155}{nav.print.76}%
\Destination{79.36}{376.055}{nav.number.1.12.7}%
\Hotlink{449.694}{184.255}{6.858}{11.047}{nav.section.1.3}%
\Hotlink{471.108}{758.955}{30.936}{11.047}{nav.number.1.11.5}%
\Hotlink{243.313}{554.755}{30.936}{11.047}{nav.number.1.12.3}%
\Hotlink{304.205}{84.855}{26.464}{9.572}{nav.number.1.12.3}%
\Destination{79.226}{547.176}{errata.51}%
}

% Page 84 (body page 72)
\Page{%
\TextAt{2}{79.226}{796.943}{I.}%
\TextAt{2}{89.83}{796.943}{Algebraic}%
\TextAt{2}{138.648}{796.943}{Integers}%
\Vector{q 0 G 0 g 0.299 w 7.226 18.729 m 460.769 18.729 l S Q}%
\TextAt{8}{79.226}{761.124}{Let}%
\TextAt{13}{100.239}{761.124}{\Character{12}}%
\TextAt{9}{110.905}{761.124}{\Character{61}}%
\TextAt{13}{123.405}{761.124}{b}%
\TextAt{9}{128.382}{761.124}{\Character{47}}%
\TextAt{13}{134.235}{761.124}{c}%
\NativeMathAt{142.667}{761.124}{11.955}{\in}%
\TextAt{20}{153.634}{761.659}{\Character{101}}%
\TextAt{14}{154.032}{761.124}{o}%
\TextAt{8}{159.877}{761.124}{,}%
\TextAt{13}{167.08}{761.124}{b\Character{59}}%
\TextAt{13}{177.278}{761.124}{c}%
\NativeMathAt{185.734}{761.124}{11.955}{\in}%
\TextAt{14}{197.099}{761.124}{o}%
\TextAt{8}{202.944}{761.124}{.}%
\TextAt{8}{211.522}{761.124}{Apply}%
\TextAt{8}{246.678}{761.124}{the}%
\TextAt{8}{266.878}{761.124}{descending}%
\TextAt{8}{326.159}{761.124}{chain}%
\TextAt{8}{357.095}{761.124}{condition}%
\TextAt{8}{408.521}{761.124}{to}%
\TextAt{8}{422.875}{761.124}{the}%
\TextAt{8}{443.075}{761.124}{ring}%
\TextAt{14}{467.184}{761.124}{o}%
\TextAt{9}{473.028}{761.124}{\Character{47}}%
\TextAt{13}{478.881}{761.124}{c}%
\TextAt{14}{483.919}{761.124}{o}%
\TextAt{8}{493.704}{761.124}{and}%
\TextAt{8}{516.511}{761.124}{the}%
\TextAt{8}{79.226}{744.078}{chain}%
\TextAt{8}{110.587}{744.078}{of}%
\TextAt{8}{124.387}{744.078}{ideals}%
\TextAt{9}{157.434}{744.078}{\Character{40}}%
\Vector{q 0 G 0 g 0.478 w 89.986 -23.101 m 96.131 -23.101 l S Q}%
\TextAt{13}{161.986}{744.078}{a}%
\TextAt{21}{168.131}{748.461}{m}%
\TextAt{9}{176.12}{744.078}{\Character{41}}%
\TextAt{8}{180.672}{744.078}{,}%
\TextAt{8}{188.408}{744.078}{where}%
\Vector{q 0 G 0 g 0.478 w 150.686 -23.101 m 156.831 -23.101 l S Q}%
\TextAt{13}{222.686}{744.078}{a}%
\TextAt{9}{232.951}{744.078}{\Character{61}}%
\TextAt{13}{246.176}{744.078}{a}%
\TextAt{8}{255.642}{744.078}{mod}%
\TextAt{13}{281.08}{744.078}{c}%
\TextAt{14}{286.118}{744.078}{o}%
\TextAt{8}{291.962}{744.078}{.}%
\TextAt{8}{301.818}{744.078}{Then}%
\TextAt{9}{332.845}{744.078}{\Character{40}}%
\Vector{q 0 G 0 g 0.478 w 265.397 -23.101 m 271.542 -23.101 l S Q}%
\TextAt{13}{337.397}{744.078}{a}%
\TextAt{21}{343.542}{748.461}{h}%
\TextAt{9}{348.92}{744.078}{\Character{41}}%
\TextAt{9}{357.576}{744.078}{\Character{61}}%
\TextAt{9}{370.798}{744.078}{\Character{40}}%
\Vector{q 0 G 0 g 0.478 w 303.37 -23.101 m 309.515 -23.101 l S Q}%
\TextAt{13}{375.37}{744.078}{a}%
\TextAt{21}{381.515}{748.461}{h}%
\TextAt{17}{386.395}{748.461}{\Character{43}1}%
\TextAt{9}{397.714}{744.078}{\Character{41}}%
\TextAt{8}{402.266}{744.078}{,}%
\TextAt{8}{410.002}{744.078}{i.e.,}%
\TextAt{8}{432.693}{744.078}{we}%
\TextAt{8}{450.39}{744.078}{find}%
\TextAt{8}{474.267}{744.078}{some}%
\TextAt{13}{504.062}{744.078}{x}%
\NativeMathAt{514.835}{744.078}{11.955}{\in}%
\TextAt{14}{526.925}{744.078}{o}%
\TextAt{8}{79.226}{727.032}{such}%
\TextAt{8}{105.623}{727.032}{that}%
\TextAt{13}{130.992}{727.032}{a}%
\TextAt{21}{137.137}{731.37}{h}%
\NativeMathAt{145.836}{727.032}{11.955}{\equiv}%
\TextAt{13}{158.455}{727.032}{xa}%
\TextAt{21}{171.252}{731.37}{h}%
\TextAt{17}{176.132}{731.37}{\Character{43}1}%
\TextAt{9}{192.764}{727.032}{\Character{40}}%
\TextAt{8}{197.317}{727.032}{mod}%
\TextAt{13}{223.419}{727.032}{c}%
\TextAt{14}{228.457}{727.032}{o}%
\TextAt{9}{234.301}{727.032}{\Character{41}}%
\TextAt{8}{238.853}{727.032}{,}%
\TextAt{8}{246.002}{727.032}{hence}%
\TextAt{9}{278.509}{727.032}{\Character{40}1}%
\TextAt{16}{291.57}{727.032}{\Character{0}}%
\TextAt{13}{303.526}{727.032}{xa}%
\TextAt{9}{316.323}{727.032}{\Character{41}}%
\TextAt{13}{320.875}{727.032}{a}%
\TextAt{21}{327.02}{731.37}{h}%
\NativeMathAt{335.719}{727.032}{11.955}{\in}%
\TextAt{13}{347.01}{727.032}{c}%
\TextAt{14}{352.048}{727.032}{o}%
\TextAt{8}{357.892}{727.032}{,}%
\TextAt{8}{365.041}{727.032}{and}%
\TextAt{8}{387.804}{727.032}{therefore}%
\TextAt{13}{158.022}{693.131}{\Character{12}}%
\TextAt{9}{168.615}{693.131}{\Character{61}}%
\TextAt{13}{182.266}{701.218}{b}%
\Vector{q 0 G 0 g 0.478 w 110.236 -77.881 m 115.273 -77.881 l S Q}%
\TextAt{13}{182.236}{684.93}{c}%
\TextAt{9}{188.469}{693.131}{\Character{40}1}%
\TextAt{16}{201.531}{693.131}{\Character{0}}%
\TextAt{13}{213.486}{693.131}{xa}%
\TextAt{9}{226.283}{693.131}{\Character{41}}%
\TextAt{9}{233.48}{693.131}{\Character{43}}%
\TextAt{13}{245.253}{693.131}{\Character{12}}%
\TextAt{13}{252.522}{693.131}{xa}%
\TextAt{9}{268.643}{693.131}{\Character{61}}%
\TextAt{13}{285.537}{701.218}{b}%
\Vector{q 0 G 0 g 0.478 w 210.264 -77.881 m 221.787 -77.881 l S Q}%
\TextAt{13}{282.264}{684.93}{a}%
\TextAt{21}{288.409}{688.384}{h}%
\TextAt{9}{296.178}{701.218}{\Character{40}1}%
\TextAt{16}{309.24}{701.218}{\Character{0}}%
\TextAt{13}{321.195}{701.218}{xa}%
\TextAt{9}{333.992}{701.218}{\Character{41}}%
\TextAt{13}{338.544}{701.218}{a}%
\TextAt{21}{344.689}{705.557}{h}%
\Vector{q 0 G 0 g 0.478 w 224.178 -77.881 m 278.067 -77.881 l S Q}%
\TextAt{13}{320.603}{684.93}{c}%
\TextAt{9}{353.919}{693.131}{\Character{43}}%
\TextAt{13}{365.681}{693.131}{\Character{12}}%
\TextAt{13}{372.95}{693.131}{xa}%
\NativeMathAt{389.07}{693.131}{11.955}{\in}%
\TextAt{13}{400.361}{693.131}{a}%
\TextAt{3}{406.506}{698.067}{\Character{0}}%
\TextAt{21}{413.092}{698.067}{h}%
\TextAt{14}{418.47}{693.131}{o}%
\TextAt{9}{426.971}{693.131}{\Character{43}}%
\TextAt{13}{438.732}{693.131}{a}%
\TextAt{20}{444.479}{693.666}{\Character{101}}%
\TextAt{14}{444.877}{693.131}{o}%
\TextAt{13}{450.722}{693.131}{\Character{58}}%
\TextAt{8}{79.226}{662.105}{Let}%
\TextAt{13}{100.574}{662.105}{h}%
\TextAt{8}{111.587}{662.105}{be}%
\TextAt{8}{127.89}{662.105}{the}%
\TextAt{8}{148.424}{662.105}{smallest}%
\TextAt{8}{193.801}{662.105}{positive}%
\TextAt{8}{237.493}{662.105}{integer}%
\TextAt{8}{276.57}{662.105}{such}%
\TextAt{8}{303.345}{662.105}{that}%
\TextAt{13}{329.092}{662.105}{\Character{12}}%
\NativeMathAt{340.328}{662.105}{11.955}{\in}%
\TextAt{13}{352.263}{662.105}{a}%
\TextAt{3}{358.408}{666.444}{\Character{0}}%
\TextAt{21}{364.995}{666.444}{h}%
\TextAt{14}{370.373}{662.105}{o}%
\TextAt{9}{379.131}{662.105}{\Character{43}}%
\TextAt{13}{391.15}{662.105}{a}%
\TextAt{20}{396.897}{662.641}{\Character{101}}%
\TextAt{14}{397.295}{662.105}{o}%
\TextAt{8}{403.14}{662.105}{.}%
\TextAt{8}{412.722}{662.105}{It}%
\TextAt{8}{425.772}{662.105}{then}%
\TextAt{8}{452.81}{662.105}{suffices}%
\TextAt{8}{492.975}{662.105}{to}%
\TextAt{8}{507.664}{662.105}{show}%
\TextAt{8}{79.226}{645.059}{that}%
\TextAt{13}{104.595}{645.059}{h}%
\NativeMathAt{114.654}{645.059}{11.955}{\leq}%
\TextAt{13}{127.274}{645.059}{n}%
\TextAt{8}{134.261}{645.059}{.}%
\TextAt{8}{142.71}{645.059}{Assume}%
\TextAt{13}{186.071}{645.059}{h}%
\TextAt{13}{196.113}{645.059}{\Character{62}}%
\TextAt{13}{208.535}{645.059}{n}%
\TextAt{8}{215.544}{645.059}{.}%
\TextAt{8}{223.992}{645.059}{Writing}%
\TextAt{13}{215.319}{616.058}{\Character{12}}%
\TextAt{9}{225.911}{616.058}{\Character{61}}%
\TextAt{13}{241.962}{624.146}{u}%
\Vector{q 0 G 0 g 0.478 w 167.532 -154.953 m 179.055 -154.953 l S Q}%
\TextAt{13}{239.532}{607.857}{a}%
\TextAt{21}{245.677}{611.311}{h}%
\TextAt{9}{254.907}{616.058}{\Character{43}}%
\TextAt{13}{266.669}{616.058}{a}%
\TextAt{20}{273.149}{616.058}{\Character{101}}%
\TextAt{13}{272.814}{616.058}{u}%
\TextAt{8}{303.386}{616.058}{with}%
\TextAt{13}{330.046}{616.058}{u}%
\NativeMathAt{340.03}{616.058}{11.955}{\in}%
\TextAt{14}{351.321}{616.058}{o}%
\TextAt{13}{357.165}{616.058}{\Character{59}}%
\TextAt{20}{366.642}{616.058}{\Character{101}}%
\TextAt{13}{366.306}{616.058}{u}%
\NativeMathAt{376.29}{616.058}{11.955}{\in}%
\TextAt{20}{387.182}{616.594}{\Character{101}}%
\TextAt{14}{387.581}{616.058}{o}%
\TextAt{13}{393.425}{616.058}{\Character{59}}%
\TextAt{8}{517.802}{616.058}{(*)}%
\TextAt{8}{79.226}{585.033}{we}%
\TextAt{8}{96.825}{585.033}{have}%
\TextAt{13}{124.192}{585.033}{u}%
\TextAt{9}{134.809}{585.033}{\Character{61}}%
\TextAt{13}{147.868}{585.033}{a}%
\TextAt{21}{154.013}{589.371}{h}%
\TextAt{9}{159.391}{585.033}{\Character{40}}%
\TextAt{13}{163.943}{585.033}{\Character{12}}%
\TextAt{16}{174.125}{585.033}{\Character{0}}%
\TextAt{13}{186.334}{585.033}{a}%
\TextAt{20}{192.814}{585.033}{\Character{101}}%
\TextAt{13}{192.478}{585.033}{u}%
\TextAt{9}{199.141}{585.033}{\Character{41}}%
\NativeMathAt{207.648}{585.033}{11.955}{\in}%
\TextAt{13}{219.572}{585.033}{a}%
\TextAt{21}{225.717}{589.371}{h}%
\TextAt{20}{230.697}{585.568}{\Character{101}}%
\TextAt{14}{231.095}{585.033}{o}%
\TextAt{16}{239.85}{585.033}{\Character{92}}%
\TextAt{14}{250.73}{585.033}{o}%
\NativeMathAt{260.529}{585.033}{11.955}{\subseteq}%
\TextAt{14}{273.782}{585.033}{a}%
\TextAt{21}{279.765}{583.239}{h}%
\TextAt{9}{289.097}{585.033}{\Character{61}}%
\TextAt{14}{302.157}{585.033}{a}%
\TextAt{21}{308.14}{583.239}{h}%
\TextAt{17}{313.02}{583.239}{\Character{43}1}%
\TextAt{8}{328.608}{585.033}{because}%
\TextAt{13}{372.281}{585.033}{h}%
\TextAt{13}{382.957}{585.033}{\Character{62}}%
\TextAt{13}{396.012}{585.033}{n}%
\TextAt{8}{403.021}{585.033}{,}%
\TextAt{8}{410.635}{585.033}{hence}%
\TextAt{13}{443.513}{585.033}{u}%
\TextAt{9}{454.13}{585.033}{\Character{61}}%
\TextAt{13}{467.189}{585.033}{a}%
\TextAt{21}{473.334}{589.371}{h}%
\TextAt{17}{478.214}{589.371}{\Character{43}1}%
\TextAt{20}{489.868}{585.033}{\Character{101}}%
\TextAt{13}{489.533}{585.033}{u}%
\TextAt{3}{496.195}{589.371}{\Character{48}}%
\TextAt{9}{501.901}{585.033}{\Character{43}}%
\TextAt{13}{513.915}{585.033}{au}%
\TextAt{3}{526.723}{589.371}{\Character{48}}%
\TextAt{8}{529.518}{585.033}{,}%
\TextAt{13}{79.226}{567.987}{u}%
\TextAt{3}{85.889}{572.325}{\Character{48}}%
\NativeMathAt{92.004}{567.987}{11.955}{\in}%
\TextAt{14}{103.295}{567.987}{o}%
\TextAt{8}{109.14}{567.987}{,}%
\TextAt{20}{116.624}{567.987}{\Character{101}}%
\TextAt{13}{116.289}{567.987}{u}%
\TextAt{3}{122.951}{572.325}{\Character{48}}%
\NativeMathAt{129.067}{567.987}{11.955}{\in}%
\TextAt{20}{139.959}{568.522}{\Character{101}}%
\TextAt{14}{140.358}{567.987}{o}%
\TextAt{8}{146.202}{567.987}{.}%
\TextAt{8}{154.651}{567.987}{Substituting}%
\TextAt{8}{221.707}{567.987}{this}%
\TextAt{8}{244.53}{567.987}{into}%
\TextAt{8}{268.273}{567.987}{(*)}%
\TextAt{8}{287.138}{567.987}{gives}%
\TextAt{13}{215.79}{535.181}{\Character{12}}%
\TextAt{9}{226.382}{535.181}{\Character{61}}%
\TextAt{13}{246.446}{543.268}{u}%
\TextAt{3}{253.108}{547.607}{\Character{48}}%
\Vector{q 0 G 0 g 0.478 w 168.003 -235.831 m 190.346 -235.831 l S Q}%
\TextAt{13}{240.003}{526.98}{a}%
\TextAt{21}{246.148}{530.434}{h}%
\TextAt{3}{251.028}{530.434}{\Character{0}}%
\TextAt{17}{257.614}{530.434}{1}%
\TextAt{9}{266.199}{535.181}{\Character{43}}%
\TextAt{13}{277.96}{535.181}{a}%
\TextAt{9}{284.105}{535.181}{\Character{40}}%
\TextAt{20}{288.993}{535.181}{\Character{101}}%
\TextAt{13}{288.657}{535.181}{u}%
\TextAt{9}{297.976}{535.181}{\Character{43}}%
\TextAt{20}{310.073}{535.181}{\Character{101}}%
\TextAt{13}{309.738}{535.181}{u}%
\TextAt{3}{316.4}{540.117}{\Character{48}}%
\TextAt{9}{319.195}{535.181}{\Character{41}}%
\NativeMathAt{327.068}{535.181}{11.955}{\in}%
\TextAt{13}{338.359}{535.181}{a}%
\TextAt{17}{344.504}{540.117}{1}%
\TextAt{3}{348.738}{540.117}{\Character{0}}%
\TextAt{21}{355.325}{540.117}{h}%
\TextAt{14}{360.703}{535.181}{o}%
\TextAt{9}{369.204}{535.181}{\Character{43}}%
\TextAt{13}{380.965}{535.181}{a}%
\TextAt{20}{386.711}{535.716}{\Character{101}}%
\TextAt{14}{387.11}{535.181}{o}%
\TextAt{13}{392.954}{535.181}{\Character{58}}%
\TextAt{8}{79.226}{504.859}{This}%
\TextAt{8}{105.671}{504.859}{contradicts}%
\TextAt{8}{165.626}{504.859}{the}%
\TextAt{8}{185.507}{504.859}{minimality}%
\TextAt{8}{244.422}{504.859}{of}%
\TextAt{13}{257.477}{504.859}{h}%
\TextAt{8}{264.216}{504.859}{.}%
\TextAt{8}{272.573}{504.859}{So}%
\TextAt{8}{288.557}{504.859}{we}%
\TextAt{8}{305.509}{504.859}{do}%
\TextAt{8}{321.493}{504.859}{have}%
\TextAt{20}{347.814}{505.394}{\Character{101}}%
\TextAt{14}{348.213}{504.859}{o}%
\NativeMathAt{357.378}{504.859}{11.955}{\subseteq}%
\TextAt{13}{369.997}{504.859}{a}%
\TextAt{3}{376.142}{509.197}{\Character{0}}%
\TextAt{21}{382.728}{509.197}{n}%
\TextAt{14}{388.365}{504.859}{o}%
\TextAt{9}{396.303}{504.859}{\Character{43}}%
\TextAt{13}{407.502}{504.859}{a}%
\TextAt{20}{413.248}{505.394}{\Character{101}}%
\TextAt{14}{413.647}{504.859}{o}%
\TextAt{8}{419.491}{504.859}{.}%
\TextAt{20}{427.449}{505.394}{\Character{101}}%
\TextAt{14}{427.848}{504.859}{o}%
\TextAt{9}{433.692}{504.859}{\Character{47}}%
\TextAt{13}{439.545}{504.859}{a}%
\TextAt{20}{445.291}{505.394}{\Character{101}}%
\TextAt{14}{445.69}{504.859}{o}%
\TextAt{8}{455.157}{504.859}{thus}%
\TextAt{8}{480.633}{504.859}{becomes}%
\TextAt{8}{526.911}{504.859}{a}%
\TextAt{8}{79.226}{487.813}{submodule}%
\TextAt{8}{137.691}{487.813}{of}%
\TextAt{8}{150.57}{487.813}{the}%
\TextAt{14}{170.277}{487.813}{o}%
\TextAt{8}{176.121}{487.813}{-module}%
\TextAt{9}{220.861}{487.813}{\Character{40}}%
\TextAt{13}{225.413}{487.813}{a}%
\TextAt{3}{231.558}{492.151}{\Character{0}}%
\TextAt{21}{238.145}{492.151}{n}%
\TextAt{14}{243.781}{487.813}{o}%
\TextAt{9}{251.361}{487.813}{\Character{43}}%
\TextAt{13}{262.201}{487.813}{a}%
\TextAt{20}{267.947}{488.348}{\Character{101}}%
\TextAt{14}{268.346}{487.813}{o}%
\TextAt{9}{274.19}{487.813}{\Character{41}\Character{47}}%
\TextAt{13}{284.596}{487.813}{a}%
\TextAt{20}{290.342}{488.348}{\Character{101}}%
\TextAt{14}{290.741}{487.813}{o}%
\TextAt{8}{300.032}{487.813}{generated}%
\TextAt{8}{352.913}{487.813}{by}%
\TextAt{13}{368.722}{487.813}{a}%
\TextAt{3}{374.867}{492.151}{\Character{0}}%
\TextAt{21}{381.454}{492.151}{n}%
\TextAt{8}{390.411}{487.813}{mod}%
\TextAt{13}{415.849}{487.813}{a}%
\TextAt{20}{421.595}{488.348}{\Character{101}}%
\TextAt{14}{421.994}{487.813}{o}%
\TextAt{8}{427.838}{487.813}{.}%
\TextAt{8}{436.136}{487.813}{It}%
\TextAt{8}{448.358}{487.813}{is}%
\TextAt{8}{459.672}{487.813}{therefore}%
\TextAt{8}{508.321}{487.813}{itself}%
\TextAt{8}{79.226}{470.767}{a}%
\TextAt{8}{88.982}{470.767}{finitely}%
\TextAt{8}{128.326}{470.767}{generated}%
\TextAt{14}{181.658}{470.767}{o}%
\TextAt{8}{187.502}{470.767}{-module,}%
\TextAt{8}{235.945}{470.767}{q.e.d.}%
\Vector{q 0 G 0 g 0.399 w 452.334 -303.233 m 452.334 -295.163 l S Q}%
\Vector{q 0 G 0 g 0.399 w 452.534 -295.363 m 459.707 -295.363 l S Q}%
\Vector{q 0 G 0 g 0.399 w 452.534 -303.034 m 459.707 -303.034 l S Q}%
\Vector{q 0 G 0 g 0.399 w 459.906 -303.233 m 459.906 -295.163 l S Q}%
\TextAt{12}{79.226}{435.788}{Proposition}%
\TextAt{12}{152.134}{435.788}{(12.8)}%
\TextAt{12}{191.221}{435.788}{(Krull–Akizuki).}%
\TextAt{8}{292.521}{435.788}{Let}%
\TextAt{14}{313.662}{435.788}{o}%
\TextAt{8}{323.575}{435.788}{be}%
\TextAt{8}{339.67}{435.788}{a}%
\TextAt{8}{349.597}{435.788}{one-dimensional}%
\TextAt{8}{435.678}{435.788}{noetherian}%
\TextAt{8}{494.059}{435.788}{integral}%
\TextAt{8}{79.226}{418.742}{domain}%
\TextAt{8}{121.985}{418.742}{with}%
\TextAt{8}{149.777}{418.742}{field}%
\TextAt{8}{176.265}{418.742}{of}%
\TextAt{8}{190.726}{418.742}{fractions}%
\TextAt{13}{239.726}{418.742}{K}%
\TextAt{8}{250.544}{418.742}{.}%
\TextAt{8}{262.386}{418.742}{Let}%
\TextAt{13}{284.487}{418.742}{L}%
\TextAt{16}{297.7}{418.742}{\Character{106}}%
\TextAt{13}{306.269}{418.742}{K}%
\TextAt{8}{322.116}{418.742}{be}%
\TextAt{8}{339.171}{418.742}{a}%
\TextAt{8}{350.058}{418.742}{finite}%
\TextAt{8}{381.101}{418.742}{extension}%
\TextAt{8}{433.999}{418.742}{and}%
\TextAt{16}{457.892}{418.742}{\Character{79}}%
\TextAt{8}{472.771}{418.742}{the}%
\TextAt{8}{494.059}{418.742}{integral}%
\TextAt{8}{79.226}{401.696}{closure}%
\TextAt{8}{118.308}{401.696}{of}%
\TextAt{14}{131.638}{401.696}{o}%
\TextAt{8}{141.379}{401.696}{in}%
\TextAt{13}{155.032}{401.696}{L}%
\TextAt{8}{162.997}{401.696}{.}%
\TextAt{8}{171.445}{401.696}{Then}%
\TextAt{16}{202.002}{401.696}{\Character{79}}%
\TextAt{8}{215.749}{401.696}{is}%
\TextAt{8}{227.513}{401.696}{a}%
\TextAt{8}{237.269}{401.696}{Dedekind}%
\TextAt{8}{289.453}{401.696}{domain.}%
\TextAt{19}{79.226}{366.717}{Proof.}%
\TextAt{8}{116}{366.717}{The}%
\TextAt{8}{138.976}{366.717}{facts}%
\TextAt{8}{165.599}{366.717}{that}%
\TextAt{16}{189.89}{366.717}{\Character{79}}%
\TextAt{8}{202.56}{366.717}{is}%
\TextAt{8}{213.246}{366.717}{integrally}%
\TextAt{8}{264.209}{366.717}{closed}%
\TextAt{8}{297.658}{366.717}{and}%
\TextAt{8}{319.344}{366.717}{that}%
\TextAt{8}{343.635}{366.717}{every}%
\TextAt{8}{373.449}{366.717}{nonzero}%
\TextAt{8}{415.948}{366.717}{prime}%
\TextAt{8}{448.034}{366.717}{ideal}%
\TextAt{8}{474.92}{366.717}{is}%
\TextAt{8}{485.606}{366.717}{maximal,}%
\TextAt{8}{79.226}{349.671}{are}%
\TextAt{8}{97.981}{349.671}{deduced}%
\TextAt{8}{142.738}{349.671}{as}%
\TextAt{8}{156.351}{349.671}{in}%
\TextAt{8}{169.248}{349.671}{(3.1).}%
\TextAt{8}{201.522}{349.671}{It}%
\TextAt{8}{213.438}{349.671}{remains}%
\TextAt{8}{256.318}{349.671}{to}%
\TextAt{8}{269.872}{349.671}{show}%
\TextAt{8}{298.118}{349.671}{that}%
\TextAt{16}{322.731}{349.671}{\Character{79}}%
\TextAt{8}{335.722}{349.671}{is}%
\TextAt{8}{346.729}{349.671}{noetherian.}%
\TextAt{8}{409.238}{349.671}{Let}%
\TextAt{13}{429.451}{349.671}{\Character{33}}%
\TextAt{17}{436.74}{347.878}{1}%
\TextAt{13}{441.472}{349.671}{\Character{59}}%
\TextAt{13}{446.696}{349.671}{\Character{58}}%
\TextAt{13}{451.933}{349.671}{\Character{58}}%
\TextAt{13}{457.169}{349.671}{\Character{58}}%
\TextAt{13}{462.394}{349.671}{\Character{59}}%
\TextAt{13}{467.63}{349.671}{\Character{33}}%
\TextAt{21}{474.982}{347.878}{n}%
\TextAt{8}{483.76}{349.671}{be}%
\TextAt{8}{498.928}{349.671}{a}%
\TextAt{8}{507.927}{349.671}{basis}%
\TextAt{8}{79.226}{332.625}{of}%
\TextAt{13}{92.172}{332.625}{L}%
\TextAt{16}{103.457}{332.625}{\Character{106}}%
\TextAt{13}{110.099}{332.625}{K}%
\TextAt{8}{124.429}{332.625}{which}%
\TextAt{8}{157.531}{332.625}{is}%
\TextAt{8}{168.911}{332.625}{contained}%
\TextAt{8}{221.535}{332.625}{in}%
\TextAt{16}{234.804}{332.625}{\Character{79}}%
\TextAt{8}{244.654}{332.625}{.}%
\TextAt{8}{252.974}{332.625}{Then}%
\TextAt{8}{283.147}{332.625}{the}%
\TextAt{8}{302.919}{332.625}{ring}%
\TextAt{16}{326.6}{332.625}{\Character{79}}%
\TextAt{17}{336.118}{330.832}{0}%
\TextAt{9}{344.171}{332.625}{\Character{61}}%
\TextAt{14}{356.596}{332.625}{o}%
\TextAt{9}{362.44}{332.625}{\Character{91}}%
\TextAt{13}{365.692}{332.625}{\Character{33}}%
\TextAt{17}{372.981}{330.832}{1}%
\TextAt{13}{377.714}{332.625}{\Character{59}}%
\TextAt{13}{382.938}{332.625}{\Character{58}}%
\TextAt{13}{388.175}{332.625}{\Character{58}}%
\TextAt{13}{393.399}{332.625}{\Character{58}}%
\TextAt{13}{398.636}{332.625}{\Character{59}}%
\TextAt{13}{403.872}{332.625}{\Character{33}}%
\TextAt{21}{411.224}{330.832}{n}%
\TextAt{9}{416.86}{332.625}{\Character{93}}%
\TextAt{8}{423.625}{332.625}{is}%
\TextAt{8}{435.004}{332.625}{a}%
\TextAt{8}{444.375}{332.625}{finitely}%
\TextAt{8}{483.335}{332.625}{generated}%
\TextAt{14}{79.226}{315.579}{o}%
\TextAt{8}{85.07}{315.579}{-module}%
\TextAt{8}{130.65}{315.579}{and}%
\TextAt{8}{153.802}{315.579}{in}%
\TextAt{8}{167.843}{315.579}{particular}%
\TextAt{8}{222.222}{315.579}{is}%
\TextAt{8}{234.375}{315.579}{noetherian}%
\TextAt{8}{292.974}{315.579}{since}%
\TextAt{14}{322.031}{315.579}{o}%
\TextAt{8}{332.162}{315.579}{is}%
\TextAt{8}{344.315}{315.579}{noetherian.}%
\TextAt{8}{408.243}{315.579}{We}%
\TextAt{8}{428.776}{315.579}{argue}%
\TextAt{8}{461.037}{315.579}{as}%
\TextAt{8}{475.797}{315.579}{before}%
\TextAt{8}{511.298}{315.579}{that}%
\TextAt{16}{79.226}{298.534}{\Character{79}}%
\TextAt{17}{88.744}{296.74}{0}%
\TextAt{8}{97.082}{298.534}{is}%
\TextAt{8}{108.555}{298.534}{one-dimensional}%
\TextAt{8}{194.174}{298.534}{and}%
\TextAt{8}{216.645}{298.534}{are}%
\TextAt{8}{235.864}{298.534}{thus}%
\TextAt{8}{261.325}{298.534}{reduced}%
\TextAt{8}{304.598}{298.534}{to}%
\TextAt{8}{318.617}{298.534}{the}%
\TextAt{8}{338.482}{298.534}{case}%
\TextAt{13}{362.962}{298.534}{L}%
\TextAt{9}{374.247}{298.534}{\Character{61}}%
\TextAt{13}{386.673}{298.534}{K}%
\TextAt{8}{397.491}{298.534}{.}%
\TextAt{8}{405.842}{298.534}{So}%
\TextAt{8}{421.81}{298.534}{let}%
\TextAt{14}{438.423}{298.534}{A}%
\TextAt{8}{450.608}{298.534}{be}%
\TextAt{8}{466.241}{298.534}{an}%
\TextAt{8}{482.209}{298.534}{ideal}%
\TextAt{8}{509.881}{298.534}{of}%
\TextAt{16}{522.92}{298.534}{\Character{79}}%
\TextAt{8}{79.226}{281.488}{and}%
\TextAt{13}{102.268}{281.488}{a}%
\NativeMathAt{112.21}{281.488}{11.955}{\in}%
\TextAt{14}{123.977}{281.488}{A}%
\TextAt{16}{135.403}{281.488}{\Character{92}}%
\TextAt{14}{146.22}{281.488}{o}%
\TextAt{8}{152.064}{281.488}{,}%
\TextAt{13}{159.563}{281.488}{a}%
\NativeMathAt{169.504}{281.488}{11.955}{\neq}%
\TextAt{9}{182.392}{281.488}{0}%
\TextAt{8}{188.259}{281.488}{;}%
\TextAt{8}{195.827}{281.488}{then}%
\TextAt{8}{222.766}{281.488}{by}%
\TextAt{8}{239.305}{281.488}{the}%
\TextAt{8}{259.74}{281.488}{above}%
\TextAt{8}{293.195}{281.488}{lemma}%
\TextAt{16}{331.193}{281.488}{\Character{79}}%
\TextAt{9}{341.043}{281.488}{\Character{47}}%
\TextAt{13}{346.896}{281.488}{a}%
\TextAt{16}{353.041}{281.488}{\Character{79}}%
\TextAt{8}{367.067}{281.488}{is}%
\TextAt{8}{379.111}{281.488}{a}%
\TextAt{8}{389.145}{281.488}{finitely}%
\TextAt{8}{428.769}{281.488}{generated}%
\TextAt{14}{482.38}{281.488}{o}%
\TextAt{8}{488.225}{281.488}{-module.}%
\TextAt{8}{79.226}{264.442}{Since}%
\TextAt{14}{109.784}{264.442}{o}%
\TextAt{8}{119.525}{264.442}{is}%
\TextAt{8}{131.289}{264.442}{noetherian,}%
\TextAt{8}{192.751}{264.442}{so}%
\TextAt{8}{207.121}{264.442}{is}%
\TextAt{8}{218.885}{264.442}{the}%
\TextAt{14}{239.041}{264.442}{o}%
\TextAt{8}{244.885}{264.442}{-submodule}%
\TextAt{14}{307.698}{264.442}{A}%
\TextAt{9}{316.277}{264.442}{\Character{47}}%
\TextAt{13}{322.13}{264.442}{a}%
\TextAt{16}{328.275}{264.442}{\Character{79}}%
\TextAt{8}{338.124}{264.442}{,}%
\TextAt{8}{345.274}{264.442}{and}%
\TextAt{8}{368.036}{264.442}{also}%
\TextAt{8}{391.516}{264.442}{the}%
\TextAt{16}{411.673}{264.442}{\Character{79}}%
\TextAt{8}{421.522}{264.442}{-module}%
\TextAt{14}{466.713}{264.442}{A}%
\TextAt{8}{475.292}{264.442}{.}%
\Vector{q 0 G 0 g 0.399 w 452.334 -509.558 m 452.334 -501.489 l S Q}%
\Vector{q 0 G 0 g 0.399 w 452.534 -501.688 m 459.707 -501.688 l S Q}%
\Vector{q 0 G 0 g 0.399 w 452.534 -509.359 m 459.707 -509.359 l S Q}%
\Vector{q 0 G 0 g 0.399 w 459.906 -509.558 m 459.906 -501.489 l S Q}%
\TextAt{12}{97.159}{235.44}{Remark:}%
\TextAt{8}{152.637}{235.44}{The}%
\TextAt{8}{176.702}{235.44}{above}%
\TextAt{8}{209.888}{235.44}{proof}%
\TextAt{8}{240.469}{235.44}{is}%
\TextAt{8}{252.243}{235.44}{taken}%
\TextAt{8}{284.127}{235.44}{from}%
\TextAt{26}{311.778}{235.44}{Kaplansky}%
\TextAt{8}{377.053}{235.44}{’s}%
\TextAt{8}{388.828}{235.44}{book}%
\TextAt{8}{417.783}{235.44}{[82]}%
\TextAt{8}{439.911}{235.44}{(see}%
\TextAt{8}{463.39}{235.44}{also}%
\TextAt{8}{486.881}{235.44}{[101]).}%
\TextAt{8}{523.994}{235.44}{It}%
\TextAt{8}{79.226}{218.394}{shows}%
\TextAt{8}{112.627}{218.394}{at}%
\TextAt{8}{126.72}{218.394}{the}%
\TextAt{8}{146.66}{218.394}{same}%
\TextAt{8}{175.769}{218.394}{time}%
\TextAt{8}{202.212}{218.394}{that}%
\TextAt{8}{227.364}{218.394}{proposition}%
\TextAt{8}{288.68}{218.394}{(8.1),}%
\TextAt{8}{319.733}{218.394}{which}%
\TextAt{8}{353.002}{218.394}{we}%
\TextAt{8}{370.013}{218.394}{had}%
\TextAt{8}{392.558}{218.394}{proved}%
\TextAt{8}{430.395}{218.394}{only}%
\TextAt{8}{455.869}{218.394}{in}%
\TextAt{8}{469.305}{218.394}{the}%
\TextAt{8}{489.244}{218.394}{case}%
\TextAt{8}{513.798}{218.394}{of}%
\TextAt{8}{526.911}{218.394}{a}%
\TextAt{8}{79.226}{201.349}{separable}%
\TextAt{8}{130.692}{201.349}{extension}%
\TextAt{13}{182.481}{201.349}{L}%
\TextAt{16}{193.805}{201.348}{\Character{106}}%
\TextAt{13}{200.485}{201.349}{K}%
\TextAt{8}{211.303}{201.349}{,}%
\TextAt{8}{218.481}{201.349}{is}%
\TextAt{8}{230.267}{201.349}{valid}%
\TextAt{8}{258.588}{201.349}{for}%
\TextAt{8}{276.496}{201.349}{general}%
\TextAt{8}{316.843}{201.349}{finite}%
\TextAt{8}{346.778}{201.349}{extensions}%
\TextAt{8}{403.181}{201.349}{of}%
\TextAt{8}{416.534}{201.349}{the}%
\TextAt{8}{436.713}{201.349}{field}%
\TextAt{8}{462.093}{201.349}{of}%
\TextAt{8}{475.446}{201.349}{fractions}%
\TextAt{8}{523.337}{201.349}{of}%
\TextAt{8}{79.226}{184.303}{a}%
\TextAt{8}{88.982}{184.303}{Dedekind}%
\TextAt{8}{141.166}{184.303}{domain.}%
\TextAt{8}{97.159}{164.268}{Next}%
\TextAt{8}{125.991}{164.268}{we}%
\TextAt{8}{143.442}{164.268}{want}%
\TextAt{8}{172.286}{164.268}{to}%
\TextAt{8}{186.819}{164.268}{compare}%
\TextAt{8}{233.871}{164.268}{the}%
\TextAt{8}{254.251}{164.268}{one-dimensional}%
\TextAt{8}{340.384}{164.268}{noetherian}%
\TextAt{8}{398.817}{164.268}{integral}%
\TextAt{8}{441.649}{164.268}{domain}%
\TextAt{14}{483.5}{164.268}{o}%
\TextAt{8}{493.465}{164.268}{with}%
\TextAt{8}{520.348}{164.268}{its}%
\TextAt{8}{79.226}{147.222}{normalization}%
\TextAt{20}{153.536}{147.757}{\Character{101}}%
\TextAt{14}{153.934}{147.222}{o}%
\TextAt{8}{159.779}{147.222}{.}%
\TextAt{8}{169.878}{147.222}{The}%
\TextAt{8}{194.482}{147.222}{fact}%
\TextAt{8}{218.118}{147.222}{that}%
\TextAt{20}{243.639}{147.757}{\Character{101}}%
\TextAt{14}{244.037}{147.222}{o}%
\TextAt{8}{254.329}{147.222}{is}%
\TextAt{8}{266.644}{147.222}{a}%
\TextAt{8}{276.95}{147.222}{Dedekind}%
\TextAt{8}{329.684}{147.222}{domain}%
\TextAt{8}{371.863}{147.222}{is}%
\TextAt{8}{384.177}{147.222}{evident}%
\TextAt{8}{425.698}{147.222}{and}%
\TextAt{8}{449.011}{147.222}{does}%
\TextAt{8}{475.958}{147.222}{not}%
\TextAt{8}{497.322}{147.222}{require}%
\TextAt{8}{79.226}{130.176}{the}%
\TextAt{8}{99.383}{130.176}{lengthy}%
\TextAt{8}{141.01}{130.176}{proof}%
\TextAt{8}{171.58}{130.176}{of}%
\TextAt{8}{184.91}{130.176}{(12.8)}%
\TextAt{8}{218.743}{130.176}{provided}%
\TextAt{8}{266.875}{130.176}{we}%
\TextAt{8}{284.102}{130.176}{make}%
\TextAt{8}{314.671}{130.176}{the}%
\TextAt{8}{334.828}{130.176}{following}%
\TextAt{8}{384.262}{130.176}{hypothesis:}%
\TextAt{8}{109.114}{100.178}{(*)}%
\TextAt{14}{128.79}{100.178}{o}%
\TextAt{8}{139.342}{100.178}{is}%
\TextAt{8}{151.917}{100.178}{an}%
\TextAt{8}{168.986}{100.178}{integral}%
\TextAt{8}{212.405}{100.178}{domain}%
\TextAt{8}{254.844}{100.178}{whose}%
\TextAt{8}{290.181}{100.178}{normalization}%
\TextAt{20}{364.751}{100.714}{\Character{101}}%
\TextAt{14}{365.149}{100.178}{o}%
\TextAt{8}{375.702}{100.178}{is}%
\TextAt{8}{388.276}{100.178}{a}%
\TextAt{8}{398.842}{100.178}{finitely}%
\TextAt{8}{438.997}{100.178}{generated}%
\TextAt{14}{493.14}{100.178}{o}%
\TextAt{8}{498.984}{100.178}{-}%
\TextAt{8}{109.114}{83.132}{module.}%
\TextAt{8}{300.14}{43.033}{72}%
\Destination{522.36}{633.105}{nav.AMS.8}%
\Destination{78.23}{809.815}{nav.page.72}%
\Destination{79.36}{776.155}{nav.print.77}%
\Destination{79.36}{162.255}{nav.print.78}%
\Destination{79.36}{449.755}{nav.number.1.12.8}%
\Hotlink{417.283}{233.255}{19.22}{11.047}{nav.cite.N082}%
\Hotlink{486.381}{233.255}{25.078}{11.047}{nav.cite.N101}%
\Hotlink{168.748}{347.455}{25.078}{11.047}{nav.number.1.3.1}%
\Hotlink{288.18}{216.155}{25.078}{11.047}{nav.number.1.8.1}%
\Hotlink{184.41}{127.955}{30.936}{11.047}{nav.number.1.12.8}%
}

% Page 85 (body page 73)
\Page{%
\TextAt{2}{79.226}{796.943}{I.}%
\TextAt{2}{89.83}{796.943}{Algebraic}%
\TextAt{2}{138.648}{796.943}{Integers}%
\Vector{q 0 G 0 g 0.299 w 7.226 18.729 m 460.769 18.729 l S Q}%
\TextAt{8}{79.226}{761.124}{This}%
\TextAt{8}{106.062}{761.124}{condition}%
\TextAt{8}{157.561}{761.124}{will}%
\TextAt{8}{179.782}{761.124}{be}%
\TextAt{8}{195.822}{761.124}{assumed}%
\TextAt{8}{242.886}{761.124}{for}%
\TextAt{8}{260.887}{761.124}{all}%
\TextAt{8}{277.261}{761.124}{that}%
\TextAt{8}{302.746}{761.124}{follows.}%
\TextAt{8}{346.08}{761.124}{It}%
\TextAt{8}{358.869}{761.124}{avoids}%
\TextAt{8}{394.503}{761.124}{pathological}%
\TextAt{8}{460.325}{761.124}{situations}%
\TextAt{8}{513.904}{761.124}{and}%
\TextAt{8}{79.226}{744.078}{is}%
\TextAt{8}{90.253}{744.078}{satisfied}%
\TextAt{8}{134.514}{744.078}{in}%
\TextAt{8}{147.43}{744.078}{all}%
\TextAt{8}{162.951}{744.078}{interesting}%
\TextAt{8}{219.838}{744.078}{cases,}%
\TextAt{8}{251.885}{744.078}{in}%
\TextAt{8}{264.8}{744.078}{particular}%
\TextAt{8}{318.052}{744.078}{for}%
\TextAt{8}{335.2}{744.078}{the}%
\TextAt{8}{354.619}{744.078}{orders}%
\TextAt{8}{389.065}{744.078}{in}%
\TextAt{8}{401.981}{744.078}{an}%
\TextAt{8}{417.502}{744.078}{algebraic}%
\TextAt{8}{466.199}{744.078}{number}%
\TextAt{8}{508.058}{744.078}{field.}%
\TextAt{8}{97.159}{724.043}{The}%
\TextAt{8}{121.53}{724.043}{groups}%
\TextAt{8}{159.639}{724.043}{of}%
\TextAt{8}{173.286}{724.043}{units}%
\TextAt{8}{202.93}{724.043}{and}%
\TextAt{8}{226.011}{724.043}{the}%
\TextAt{8}{246.485}{724.043}{Picard}%
\TextAt{8}{284.031}{724.043}{groups}%
\TextAt{8}{322.139}{724.043}{of}%
\TextAt{14}{335.787}{724.043}{o}%
\TextAt{8}{345.846}{724.043}{and}%
\TextAt{20}{368.528}{724.578}{\Character{101}}%
\TextAt{14}{368.927}{724.043}{o}%
\TextAt{8}{378.986}{724.043}{are}%
\TextAt{8}{398.815}{724.043}{compared}%
\TextAt{8}{452.465}{724.043}{with}%
\TextAt{8}{479.442}{724.043}{each}%
\TextAt{8}{506.097}{724.043}{other}%
\TextAt{8}{79.226}{706.997}{by}%
\TextAt{8}{95.485}{706.997}{the}%
\TextAt{8}{115.642}{706.997}{following}%
\TextAt{12}{79.226}{682.063}{Proposition}%
\TextAt{12}{151.937}{682.063}{(12.9).}%
\TextAt{8}{193.984}{682.063}{One}%
\TextAt{8}{218.695}{682.063}{has}%
\TextAt{8}{239.569}{682.063}{the}%
\TextAt{8}{259.725}{682.063}{canonical}%
\TextAt{8}{311.109}{682.063}{exact}%
\TextAt{8}{342.001}{682.063}{sequence}%
\TextAt{9}{190.034}{657.331}{1}%
\NativeArrowAt{199.208}{657.332}{11.955}{19.26}{\longrightarrow}%
\TextAt{14}{221.79}{657.331}{o}%
\TextAt{3}{227.634}{662.268}{\Character{3}}%
\NativeArrowAt{235.687}{657.331}{11.955}{19.26}{\longrightarrow}%
\TextAt{20}{257.87}{657.867}{\Character{101}}%
\TextAt{14}{258.269}{657.332}{o}%
\TextAt{3}{264.113}{662.268}{\Character{3}}%
\NativeArrowAt{272.167}{657.331}{11.955}{19.26}{\longrightarrow}%
\TextAt{20}{294.749}{668.689}{\Character{77}}%
\TextAt{27}{301.787}{643.583}{p}%
\TextAt{20}{314.408}{657.867}{\Character{101}}%
\TextAt{14}{314.807}{657.332}{o}%
\TextAt{3}{320.651}{662.268}{\Character{3}}%
\TextAt{27}{320.651}{654.376}{p}%
\TextAt{9}{325.383}{657.332}{\Character{47}}%
\TextAt{14}{331.236}{657.331}{o}%
\TextAt{3}{337.08}{662.268}{\Character{3}}%
\TextAt{27}{337.08}{654.376}{p}%
\NativeArrowAt{272.166}{624.963}{11.955}{19.26}{\longrightarrow}%
\TextAt{13}{294.749}{624.963}{P}%
\TextAt{13}{303.907}{624.963}{ic}%
\TextAt{9}{312.951}{624.963}{\Character{40}}%
\TextAt{14}{317.503}{624.963}{o}%
\TextAt{9}{323.348}{624.963}{\Character{41}}%
\NativeArrowAt{331.221}{624.963}{11.955}{19.26}{\longrightarrow}%
\TextAt{13}{353.803}{624.963}{P}%
\TextAt{13}{362.961}{624.963}{ic}%
\TextAt{9}{372.005}{624.963}{\Character{40}}%
\TextAt{20}{376.159}{625.499}{\Character{101}}%
\TextAt{14}{376.558}{624.963}{o}%
\TextAt{9}{382.402}{624.963}{\Character{41}}%
\NativeArrowAt{390.275}{624.963}{11.955}{19.26}{\longrightarrow}%
\TextAt{9}{412.857}{624.963}{1}%
\TextAt{13}{418.71}{624.963}{\Character{58}}%
\TextAt{8}{79.226}{601.502}{In}%
\TextAt{8}{93.905}{601.502}{the}%
\TextAt{8}{114.12}{601.502}{sum,}%
\TextAt{14}{142.215}{601.502}{p}%
\TextAt{8}{152.153}{601.502}{varies}%
\TextAt{8}{185.124}{601.502}{over}%
\TextAt{8}{210.228}{601.502}{the}%
\TextAt{8}{230.442}{601.502}{prime}%
\TextAt{8}{263.663}{601.502}{ideals}%
\NativeMathAt{296.299}{601.502}{11.955}{\neq}%
\TextAt{9}{308.816}{601.502}{0}%
\TextAt{8}{318.632}{601.502}{of}%
\TextAt{14}{332.02}{601.502}{o}%
\TextAt{8}{341.819}{601.502}{and}%
\TextAt{20}{364.241}{602.038}{\Character{101}}%
\TextAt{14}{364.64}{601.502}{o}%
\TextAt{27}{370.484}{599.709}{p}%
\TextAt{8}{378.926}{601.502}{denotes}%
\TextAt{8}{421.317}{601.502}{the}%
\TextAt{8}{441.531}{601.502}{integral}%
\TextAt{8}{484.198}{601.502}{closure}%
\TextAt{8}{523.337}{601.502}{of}%
\TextAt{14}{79.226}{584.456}{o}%
\TextAt{27}{85.07}{582.663}{p}%
\TextAt{8}{93.455}{584.456}{in}%
\TextAt{13}{107.107}{584.456}{K}%
\TextAt{8}{117.925}{584.456}{.}%
\TextAt{19}{79.226}{551.634}{Proof.}%
\TextAt{8}{116}{551.634}{If}%
\TextAt{20}{127.363}{552.751}{\Character{101}}%
\TextAt{14}{127.692}{551.634}{p}%
\TextAt{8}{137.573}{551.634}{varies}%
\TextAt{8}{170.485}{551.634}{over}%
\TextAt{8}{195.532}{551.634}{the}%
\TextAt{8}{215.688}{551.634}{prime}%
\TextAt{8}{248.852}{551.634}{ideals}%
\TextAt{8}{281.429}{551.634}{of}%
\TextAt{20}{294.361}{552.169}{\Character{101}}%
\TextAt{14}{294.759}{551.634}{o}%
\TextAt{8}{300.604}{551.634}{,}%
\TextAt{8}{307.753}{551.634}{we}%
\TextAt{8}{324.98}{551.634}{know}%
\TextAt{8}{355.55}{551.634}{from}%
\TextAt{8}{383.19}{551.634}{(12.6)}%
\TextAt{8}{417.023}{551.634}{that}%
\TextAt{13}{260.806}{526.407}{J}%
\TextAt{9}{268.384}{526.407}{\Character{40}}%
\TextAt{20}{272.537}{526.942}{\Character{101}}%
\TextAt{14}{272.936}{526.407}{o}%
\TextAt{9}{278.78}{526.407}{\Character{41}}%
\NativeMathAt{286.653}{526.407}{11.955}{\cong}%
\TextAt{20}{299.273}{537.764}{\Character{77}}%
\TextAt{30}{305.953}{512.083}{\Character{101}}%
\TextAt{27}{306.311}{511.338}{p}%
\TextAt{13}{319.331}{526.407}{P}%
\TextAt{9}{328.502}{526.407}{\Character{40}}%
\TextAt{20}{332.655}{526.942}{\Character{101}}%
\TextAt{14}{333.054}{526.407}{o}%
\TextAt{30}{338.54}{524.768}{\Character{101}}%
\TextAt{27}{338.898}{524.024}{p}%
\TextAt{9}{343.385}{526.407}{\Character{41}}%
\TextAt{13}{347.938}{526.407}{\Character{58}}%
\TextAt{8}{79.226}{490.288}{If}%
\TextAt{14}{90.918}{490.288}{p}%
\TextAt{8}{100.799}{490.288}{is}%
\TextAt{8}{112.563}{490.288}{a}%
\TextAt{8}{122.318}{490.288}{prime}%
\TextAt{8}{155.482}{490.288}{ideal}%
\TextAt{8}{183.445}{490.288}{of}%
\TextAt{14}{196.775}{490.288}{o}%
\TextAt{8}{202.619}{490.288}{,}%
\TextAt{8}{209.768}{490.288}{then}%
\TextAt{14}{236.428}{490.288}{p}%
\TextAt{20}{242.013}{490.824}{\Character{101}}%
\TextAt{14}{242.411}{490.288}{o}%
\TextAt{8}{252.153}{490.288}{splits}%
\TextAt{8}{282.842}{490.288}{in}%
\TextAt{8}{296.495}{490.288}{the}%
\TextAt{8}{316.651}{490.288}{Dedekind}%
\TextAt{8}{368.836}{490.288}{domain}%
\TextAt{20}{410.065}{490.824}{\Character{101}}%
\TextAt{14}{410.463}{490.288}{o}%
\TextAt{8}{420.205}{490.288}{into}%
\TextAt{8}{443.948}{490.288}{a}%
\TextAt{8}{453.703}{490.288}{product}%
\TextAt{14}{267.055}{465.061}{p}%
\TextAt{20}{272.639}{465.596}{\Character{101}}%
\TextAt{14}{273.038}{465.061}{o}%
\TextAt{9}{282.203}{465.061}{\Character{61}}%
\TextAt{20}{294.299}{466.178}{\Character{101}}%
\TextAt{14}{294.628}{465.061}{p}%
\TextAt{21}{300.611}{470.286}{e}%
\TextAt{24}{304.54}{469.179}{1}%
\TextAt{17}{300.611}{462.13}{1}%
\TextAt{16}{311.181}{465.061}{\Character{1}}%
\TextAt{16}{316.477}{465.061}{\Character{1}}%
\TextAt{16}{321.785}{465.061}{\Character{1}}%
\TextAt{20}{326.792}{466.178}{\Character{101}}%
\TextAt{14}{327.122}{465.061}{p}%
\TextAt{21}{333.105}{469.997}{e}%
\TextAt{22}{337.033}{469.001}{r}%
\TextAt{21}{333.105}{462.106}{r}%
\TextAt{13}{341.689}{465.061}{\Character{59}}%
\TextAt{8}{79.226}{439.834}{i.e.,}%
\TextAt{8}{100.733}{439.834}{there}%
\TextAt{8}{129.896}{439.834}{are}%
\TextAt{8}{148.658}{439.834}{only}%
\TextAt{8}{173.602}{439.834}{finitely}%
\TextAt{8}{212.198}{439.834}{many}%
\TextAt{8}{243.322}{439.834}{prime}%
\TextAt{8}{275.737}{439.834}{ideals}%
\TextAt{8}{307.566}{439.834}{of}%
\TextAt{20}{319.749}{440.369}{\Character{101}}%
\TextAt{14}{320.148}{439.834}{o}%
\TextAt{8}{329.141}{439.834}{above}%
\TextAt{14}{361.568}{439.834}{p}%
\TextAt{8}{367.552}{439.834}{.}%
\TextAt{8}{375.75}{439.834}{The}%
\TextAt{8}{399.056}{439.834}{same}%
\TextAt{8}{427.633}{439.834}{holds}%
\TextAt{8}{457.514}{439.834}{for}%
\TextAt{8}{474.651}{439.834}{the}%
\TextAt{8}{494.059}{439.834}{integral}%
\TextAt{8}{79.226}{422.788}{closure}%
\TextAt{20}{117.203}{423.323}{\Character{101}}%
\TextAt{14}{117.602}{422.788}{o}%
\TextAt{27}{123.446}{420.994}{p}%
\TextAt{8}{131.124}{422.788}{of}%
\TextAt{14}{143.748}{422.788}{o}%
\TextAt{27}{149.592}{420.994}{p}%
\TextAt{8}{154.079}{422.788}{.}%
\TextAt{8}{162.292}{422.788}{Since}%
\TextAt{8}{192.144}{422.788}{every}%
\TextAt{8}{222.33}{422.788}{nonzero}%
\TextAt{8}{265.2}{422.788}{prime}%
\TextAt{8}{297.658}{422.788}{ideal}%
\TextAt{8}{324.915}{422.788}{of}%
\TextAt{20}{337.14}{423.323}{\Character{101}}%
\TextAt{14}{337.539}{422.788}{o}%
\TextAt{27}{343.383}{420.994}{p}%
\TextAt{8}{351.062}{422.788}{has}%
\TextAt{8}{371.229}{422.788}{to}%
\TextAt{8}{384.834}{422.788}{lie}%
\TextAt{8}{399.729}{422.788}{above}%
\TextAt{14}{432.199}{422.788}{po}%
\TextAt{27}{444.026}{420.994}{p}%
\TextAt{8}{448.513}{422.788}{,}%
\TextAt{8}{455.097}{422.788}{the}%
\TextAt{8}{474.548}{422.788}{localization}%
\TextAt{20}{78.827}{406.277}{\Character{101}}%
\TextAt{14}{79.226}{405.742}{o}%
\TextAt{27}{85.07}{403.948}{p}%
\TextAt{8}{93.46}{405.742}{has}%
\TextAt{8}{114.34}{405.742}{only}%
\TextAt{8}{140.038}{405.742}{a}%
\TextAt{8}{149.799}{405.742}{finite}%
\TextAt{8}{179.717}{405.742}{number}%
\TextAt{8}{222.319}{405.742}{of}%
\TextAt{8}{235.655}{405.742}{prime}%
\TextAt{8}{268.824}{405.742}{ideals}%
\TextAt{8}{301.408}{405.742}{and}%
\TextAt{8}{324.176}{405.742}{is}%
\TextAt{8}{335.946}{405.742}{therefore}%
\TextAt{8}{385.052}{405.742}{a}%
\TextAt{8}{394.813}{405.742}{principal}%
\TextAt{8}{443.596}{405.742}{ideal}%
\TextAt{8}{471.565}{405.742}{domain}%
\TextAt{8}{513.199}{405.742}{(see}%
\TextAt{8}{79.226}{388.696}{§3,}%
\TextAt{8}{97.912}{388.696}{exercise}%
\TextAt{8}{141.214}{388.696}{4).}%
\TextAt{8}{160.075}{388.696}{In}%
\TextAt{8}{174.696}{388.696}{view}%
\TextAt{8}{201.679}{388.696}{of}%
\TextAt{8}{215.009}{388.696}{(12.6),}%
\TextAt{8}{252.094}{388.696}{it}%
\TextAt{8}{263.798}{388.696}{follows}%
\TextAt{8}{302.234}{388.696}{that}%
\TextAt{13}{237.915}{363.468}{P}%
\TextAt{9}{247.086}{363.468}{\Character{40}}%
\TextAt{20}{251.24}{364.004}{\Character{101}}%
\TextAt{14}{251.639}{363.468}{o}%
\TextAt{27}{257.483}{361.675}{p}%
\TextAt{9}{261.97}{363.468}{\Character{41}}%
\NativeMathAt{269.843}{363.468}{11.955}{\cong}%
\TextAt{13}{282.462}{363.468}{J}%
\TextAt{9}{290.039}{363.468}{\Character{40}}%
\TextAt{20}{294.193}{364.004}{\Character{101}}%
\TextAt{14}{294.592}{363.468}{o}%
\TextAt{27}{300.436}{361.675}{p}%
\TextAt{9}{304.923}{363.468}{\Character{41}}%
\NativeMathAt{312.796}{363.468}{11.955}{\cong}%
\TextAt{20}{325.415}{374.826}{\Character{77}}%
\TextAt{30}{326.808}{349.144}{\Character{101}}%
\TextAt{27}{327.166}{348.4}{p}%
\NativeMathAt{331.155}{348.4}{7.97}{\supseteq}%
\TextAt{27}{337.741}{348.4}{p}%
\TextAt{13}{345.473}{363.468}{P}%
\TextAt{9}{354.645}{363.468}{\Character{40}}%
\TextAt{20}{358.798}{364.004}{\Character{101}}%
\TextAt{14}{359.197}{363.468}{o}%
\TextAt{30}{364.683}{361.83}{\Character{101}}%
\TextAt{27}{365.041}{361.086}{p}%
\TextAt{9}{369.528}{363.468}{\Character{41}}%
\TextAt{8}{79.226}{328.103}{and}%
\TextAt{8}{101.989}{328.103}{therefore}%
\TextAt{13}{218.475}{311.057}{J}%
\TextAt{9}{226.052}{311.057}{\Character{40}}%
\TextAt{20}{230.205}{311.593}{\Character{101}}%
\TextAt{14}{230.604}{311.057}{o}%
\TextAt{9}{236.448}{311.057}{\Character{41}}%
\NativeMathAt{244.322}{311.057}{11.955}{\cong}%
\TextAt{20}{256.941}{322.415}{\Character{77}}%
\TextAt{27}{263.979}{297.309}{p}%
\TextAt{20}{276.999}{322.415}{\Character{77}}%
\TextAt{30}{278.392}{296.733}{\Character{101}}%
\TextAt{27}{278.75}{295.989}{p}%
\NativeMathAt{282.739}{295.989}{7.97}{\supseteq}%
\TextAt{27}{289.325}{295.989}{p}%
\TextAt{13}{297.057}{311.057}{P}%
\TextAt{9}{306.228}{311.057}{\Character{40}}%
\TextAt{20}{310.382}{311.593}{\Character{101}}%
\TextAt{14}{310.781}{311.057}{o}%
\TextAt{30}{316.267}{309.419}{\Character{101}}%
\TextAt{27}{316.625}{308.675}{p}%
\TextAt{9}{321.112}{311.057}{\Character{41}}%
\NativeMathAt{328.985}{311.057}{11.955}{\cong}%
\TextAt{20}{341.604}{322.415}{\Character{77}}%
\TextAt{27}{348.643}{297.309}{p}%
\TextAt{13}{361.662}{311.057}{P}%
\TextAt{9}{370.833}{311.057}{\Character{40}}%
\TextAt{20}{374.987}{311.593}{\Character{101}}%
\TextAt{14}{375.386}{311.057}{o}%
\TextAt{27}{381.23}{309.264}{p}%
\TextAt{9}{385.717}{311.057}{\Character{41}}%
\TextAt{13}{390.269}{311.057}{\Character{58}}%
\TextAt{8}{79.226}{278.466}{Observing}%
\TextAt{8}{135.416}{278.466}{that}%
\TextAt{13}{161.3}{278.466}{P}%
\TextAt{9}{170.471}{278.466}{\Character{40}}%
\TextAt{13}{175.023}{278.466}{R}%
\TextAt{9}{184.032}{278.466}{\Character{41}}%
\NativeMathAt{192.783}{278.465}{11.955}{\cong}%
\TextAt{13}{206.279}{278.466}{K}%
\TextAt{3}{217.097}{282.804}{\Character{3}}%
\TextAt{9}{221.829}{278.466}{\Character{47}}%
\TextAt{13}{227.682}{278.466}{R}%
\TextAt{3}{236.691}{282.804}{\Character{3}}%
\TextAt{8}{245.835}{278.466}{for}%
\TextAt{8}{264.235}{278.466}{any}%
\TextAt{8}{286.867}{278.466}{integral}%
\TextAt{8}{329.99}{278.466}{domain}%
\TextAt{13}{372.133}{278.466}{R}%
\TextAt{8}{385.554}{278.466}{with}%
\TextAt{8}{412.729}{278.466}{field}%
\TextAt{8}{438.601}{278.466}{of}%
\TextAt{8}{452.446}{278.466}{fractions}%
\TextAt{13}{500.829}{278.466}{K}%
\TextAt{8}{511.647}{278.466}{,}%
\TextAt{8}{519.439}{278.466}{we}%
\TextAt{8}{79.226}{261.42}{obtain}%
\TextAt{8}{115.654}{261.42}{the}%
\TextAt{8}{135.81}{261.42}{commutative}%
\TextAt{8}{205.736}{261.42}{exact}%
\TextAt{8}{236.628}{261.42}{diagram}%
\TextAt{9}{193.02}{232.858}{1}%
\TextAt{13}{222.329}{232.858}{K}%
\TextAt{3}{233.147}{237.196}{\Character{3}}%
\TextAt{9}{237.879}{232.858}{\Character{47}}%
\TextAt{14}{243.732}{232.858}{o}%
\TextAt{3}{249.577}{237.196}{\Character{3}}%
\TextAt{20}{277.765}{244.215}{\Character{77}}%
\TextAt{27}{284.804}{219.109}{p}%
\TextAt{13}{297.823}{232.858}{K}%
\TextAt{3}{308.641}{237.794}{\Character{3}}%
\TextAt{9}{313.373}{232.858}{\Character{47}}%
\TextAt{14}{319.226}{232.858}{o}%
\TextAt{3}{325.07}{237.794}{\Character{3}}%
\TextAt{27}{325.07}{229.902}{p}%
\TextAt{13}{353.259}{232.858}{P}%
\TextAt{13}{362.417}{232.858}{ic}%
\TextAt{9}{371.461}{232.858}{\Character{40}}%
\TextAt{14}{376.014}{232.858}{o}%
\TextAt{9}{381.858}{232.858}{\Character{41}}%
\TextAt{9}{411.492}{232.858}{1}%
\TextAt{9}{193.02}{173.575}{1}%
\TextAt{13}{222.329}{173.575}{K}%
\TextAt{3}{233.147}{177.913}{\Character{3}}%
\TextAt{9}{237.879}{173.575}{\Character{47}}%
\TextAt{20}{243.334}{174.11}{\Character{101}}%
\TextAt{14}{243.732}{173.575}{o}%
\TextAt{3}{249.577}{177.913}{\Character{3}}%
\TextAt{20}{277.765}{184.933}{\Character{77}}%
\TextAt{27}{284.804}{159.827}{p}%
\TextAt{13}{297.823}{173.575}{K}%
\TextAt{3}{308.641}{178.511}{\Character{3}}%
\TextAt{9}{313.373}{173.575}{\Character{47}}%
\TextAt{20}{318.827}{174.11}{\Character{101}}%
\TextAt{14}{319.226}{173.575}{o}%
\TextAt{3}{325.07}{178.511}{\Character{3}}%
\TextAt{27}{325.07}{170.62}{p}%
\TextAt{13}{353.259}{173.575}{P}%
\TextAt{13}{362.417}{173.575}{ic}%
\TextAt{9}{371.461}{173.575}{\Character{40}}%
\TextAt{20}{375.615}{174.11}{\Character{101}}%
\TextAt{14}{376.014}{173.575}{o}%
\TextAt{9}{381.858}{173.575}{\Character{41}}%
\TextAt{9}{409.867}{173.575}{1}%
\TextAt{13}{415.72}{173.575}{\Character{58}}%
\NativeDiagram{%
\path[draw=black,line width=0.478bp] (132.226,-538.152) -- (144.5,-538.152);%
\path[draw=black,line width=0.478bp,line cap=round,line join=round] (142.253,-535.283) .. controls (142.707,-537.005) and (143.72,-537.818) .. (144.74,-538.153) .. controls (143.72,-538.487) and (142.707,-539.3) .. (142.253,-541.022);%
\path[draw=black,line width=0.478bp] (187.663,-538.152) -- (199.937,-538.152);%
\path[draw=black,line width=0.478bp,line cap=round,line join=round] (197.689,-535.283) .. controls (198.144,-537.005) and (199.156,-537.818) .. (200.176,-538.153) .. controls (199.156,-538.487) and (198.144,-539.3) .. (197.689,-541.022);%
\path[draw=black,line width=0.478bp] (166.321,-548.709) -- (166.321,-586.198);%
\path[draw=black,line width=0.478bp,line cap=round,line join=round] (169.19,-583.95) .. controls (167.468,-584.404) and (166.655,-585.417) .. (166.321,-586.437) .. controls (165.986,-585.417) and (165.173,-584.404) .. (163.452,-583.95);%
}%
\TextAt{21}{230.066}{204.592}{\Character{11}}%
\NativeDiagram{%
\path[draw=black,line width=0.478bp] (263.157,-538.152) -- (275.432,-538.152);%
\path[draw=black,line width=0.478bp,line cap=round,line join=round] (273.184,-535.283) .. controls (273.638,-537.005) and (274.651,-537.818) .. (275.671,-538.153) .. controls (274.651,-538.487) and (273.638,-539.3) .. (273.184,-541.022);%
\path[draw=black,line width=0.478bp] (231.786,-562.171) -- (231.786,-582.814);%
\path[draw=black,line width=0.478bp,line cap=round,line join=round] (234.655,-580.567) .. controls (232.934,-581.021) and (232.121,-582.034) .. (231.786,-583.053) .. controls (231.451,-582.034) and (230.639,-581.021) .. (228.917,-580.567);%
}%
\TextAt{21}{295.809}{199.275}{\Character{12}}%
\NativeDiagram{%
\path[draw=black,line width=0.478bp] (319.766,-538.152) -- (333.666,-538.152);%
\path[draw=black,line width=0.478bp,line cap=round,line join=round] (331.418,-535.283) .. controls (331.873,-537.005) and (332.885,-537.818) .. (333.905,-538.153) .. controls (332.885,-538.487) and (331.873,-539.3) .. (331.418,-541.022);%
\path[draw=black,line width=0.478bp] (297.838,-548.709) -- (297.838,-586.198);%
\path[draw=black,line width=0.478bp,line cap=round,line join=round] (300.707,-583.95) .. controls (298.985,-584.404) and (298.172,-585.417) .. (297.838,-586.437) .. controls (297.503,-585.417) and (296.69,-584.404) .. (294.969,-583.95);%
}%
\TextAt{21}{372.653}{205.367}{\Character{13}}%
\NativeDiagram{%
\path[draw=black,line width=0.478bp] (132.226,-597.436) -- (144.5,-597.436);%
\path[draw=black,line width=0.478bp,line cap=round,line join=round] (142.253,-594.567) .. controls (142.707,-596.288) and (143.72,-597.101) .. (144.74,-597.436) .. controls (143.72,-597.771) and (142.707,-598.584) .. (142.253,-600.305);%
\path[draw=black,line width=0.478bp] (187.663,-597.436) -- (199.937,-597.436);%
\path[draw=black,line width=0.478bp,line cap=round,line join=round] (197.689,-594.567) .. controls (198.144,-596.288) and (199.156,-597.101) .. (200.176,-597.436) .. controls (199.156,-597.771) and (198.144,-598.584) .. (197.689,-600.305);%
\path[draw=black,line width=0.478bp] (263.157,-597.436) -- (275.432,-597.436);%
\path[draw=black,line width=0.478bp,line cap=round,line join=round] (273.184,-594.567) .. controls (273.638,-596.288) and (274.651,-597.101) .. (275.671,-597.436) .. controls (274.651,-597.771) and (273.638,-598.584) .. (273.184,-600.305);%
\path[draw=black,line width=0.478bp] (319.766,-597.436) -- (332.04,-597.436);%
\path[draw=black,line width=0.478bp,line cap=round,line join=round] (329.792,-594.567) .. controls (330.247,-596.288) and (331.259,-597.101) .. (332.279,-597.436) .. controls (331.259,-597.771) and (330.247,-598.584) .. (329.792,-600.305);%
}%
\TextAt{8}{79.226}{135.27}{For}%
\TextAt{8}{100.924}{135.27}{such}%
\TextAt{8}{128.05}{135.27}{a}%
\TextAt{8}{138.534}{135.27}{diagram}%
\TextAt{8}{184.799}{135.27}{one}%
\TextAt{8}{206.988}{135.27}{has}%
\TextAt{8}{228.59}{135.27}{in}%
\TextAt{8}{242.971}{135.27}{complete}%
\TextAt{8}{293.122}{135.27}{generality}%
\TextAt{8}{347.841}{135.27}{the}%
\TextAt{8}{368.726}{135.27}{well-known}%
\TextAt{12}{430.258}{135.27}{snake}%
\TextAt{12}{467.763}{135.27}{lemma}%
\TextAt{8}{506.605}{135.27}{:}%
\TextAt{8}{516.51}{135.27}{the}%
\TextAt{8}{79.226}{118.224}{diagram}%
\TextAt{8}{124.763}{118.224}{gives}%
\TextAt{8}{153.444}{118.224}{in}%
\TextAt{8}{167.097}{118.224}{a}%
\TextAt{8}{176.852}{118.224}{canonical}%
\TextAt{8}{228.235}{118.224}{way}%
\TextAt{8}{251.978}{118.224}{an}%
\TextAt{8}{268.237}{118.224}{exact}%
\TextAt{8}{299.13}{118.224}{sequence}%
\TextAt{9}{243.777}{94.664}{1}%
\NativeArrowAt{252.951}{94.664}{11.955}{19.26}{\longrightarrow}%
\TextAt{8}{275.533}{94.664}{ker}%
\TextAt{9}{291.469}{94.664}{\Character{40}}%
\TextAt{13}{296.021}{94.664}{\Character{11}}%
\TextAt{9}{303.543}{94.664}{\Character{41}}%
\NativeArrowAt{311.416}{94.664}{11.955}{19.26}{\longrightarrow}%
\TextAt{8}{333.998}{94.664}{ker}%
\TextAt{9}{349.934}{94.664}{\Character{40}}%
\TextAt{13}{354.487}{94.664}{\Character{12}}%
\TextAt{9}{361.758}{94.664}{\Character{41}}%
\NativeArrowAt{369.631}{94.664}{11.955}{19.26}{\longrightarrow}%
\TextAt{8}{392.213}{94.664}{ker}%
\TextAt{9}{408.149}{94.664}{\Character{40}}%
\TextAt{13}{412.702}{94.664}{\Character{13}}%
\TextAt{9}{419.424}{94.664}{\Character{41}}%
\TextAt{21}{195.627}{78.835}{\Character{14}}%
\NativeArrowAt{188.019}{72.058}{11.955}{19.26}{\longrightarrow}%
\TextAt{8}{210.601}{72.058}{coker}%
\TextAt{9}{237.596}{72.058}{\Character{40}}%
\TextAt{13}{242.148}{72.058}{\Character{11}}%
\TextAt{9}{249.67}{72.058}{\Character{41}}%
\NativeArrowAt{257.543}{72.058}{11.955}{19.26}{\longrightarrow}%
\TextAt{8}{280.125}{72.058}{coker}%
\TextAt{9}{307.12}{72.058}{\Character{40}}%
\TextAt{13}{311.672}{72.058}{\Character{12}}%
\TextAt{9}{318.944}{72.058}{\Character{41}}%
\NativeArrowAt{326.817}{72.058}{11.955}{19.26}{\longrightarrow}%
\TextAt{8}{349.399}{72.058}{coker}%
\TextAt{9}{376.394}{72.058}{\Character{40}}%
\TextAt{13}{380.946}{72.058}{\Character{13}}%
\TextAt{9}{387.668}{72.058}{\Character{41}}%
\NativeArrowAt{395.541}{72.058}{11.955}{19.26}{\longrightarrow}%
\TextAt{9}{418.123}{72.058}{1}%
\TextAt{8}{300.14}{43.033}{73}%
\Destination{78.23}{809.815}{nav.page.73}%
\Destination{79.36}{276.455}{nav.print.79}%
\Destination{79.36}{696.055}{nav.number.1.12.9}%
\Hotlink{140.714}{386.455}{6.858}{11.047}{nav.ex.1.3.4}%
\Hotlink{84.405}{386.455}{6.858}{11.047}{nav.section.1.3}%
\Hotlink{382.69}{549.455}{30.936}{11.047}{nav.number.1.12.6}%
\Hotlink{214.509}{386.455}{30.936}{11.047}{nav.number.1.12.6}%
}

% Page 86 (body page 74)
\Page{%
\TextAt{2}{79.226}{796.943}{I.}%
\TextAt{2}{89.83}{796.943}{Algebraic}%
\TextAt{2}{138.648}{796.943}{Integers}%
\Vector{q 0 G 0 g 0.299 w 7.226 18.729 m 460.769 18.729 l S Q}%
\TextAt{8}{79.226}{761.124}{relating}%
\TextAt{8}{122.908}{761.124}{the}%
\TextAt{8}{143.816}{761.124}{kernels}%
\TextAt{8}{183.648}{761.124}{and}%
\TextAt{8}{207.162}{761.124}{cokernels}%
\TextAt{8}{258.053}{761.124}{of}%
\TextAt{13}{272.134}{761.124}{\Character{11}\Character{59}}%
\TextAt{13}{284.878}{761.124}{\Character{12}}%
\TextAt{13}{292.147}{761.124}{\Character{59}}%
\TextAt{13}{297.383}{761.124}{\Character{13}}%
\TextAt{8}{308.786}{761.124}{(see}%
\TextAt{8}{333.005}{761.124}{[23],}%
\TextAt{8}{359.313}{761.124}{chap.}%
\TextAt{8}{390.956}{761.124}{III,}%
\TextAt{8}{412.162}{761.124}{§3,}%
\TextAt{8}{431.787}{761.124}{lemma}%
\TextAt{8}{470.256}{761.124}{3.3).}%
\TextAt{8}{500.481}{761.124}{In}%
\TextAt{8}{515.853}{761.124}{our}%
\TextAt{8}{79.226}{744.078}{particular}%
\TextAt{8}{133.216}{744.078}{case,}%
\TextAt{13}{161.239}{744.078}{\Character{11}\Character{59}}%
\TextAt{13}{173.971}{744.078}{\Character{12}}%
\TextAt{8}{181.276}{744.078}{,}%
\TextAt{8}{188.425}{744.078}{and}%
\TextAt{8}{211.188}{744.078}{therefore}%
\TextAt{8}{260.287}{744.078}{also}%
\TextAt{13}{283.767}{744.078}{\Character{13}}%
\TextAt{8}{290.49}{744.078}{,}%
\TextAt{8}{297.639}{744.078}{are}%
\TextAt{8}{317.15}{744.078}{surjective,}%
\TextAt{8}{372.813}{744.078}{whereas}%
\TextAt{8}{196.769}{715.836}{ker}%
\TextAt{9}{212.705}{715.836}{\Character{40}}%
\TextAt{13}{217.258}{715.836}{\Character{11}}%
\TextAt{9}{224.779}{715.836}{\Character{41}}%
\TextAt{9}{232.634}{715.836}{\Character{61}}%
\TextAt{20}{244.679}{716.371}{\Character{101}}%
\TextAt{14}{245.078}{715.836}{o}%
\TextAt{3}{250.922}{720.772}{\Character{3}}%
\TextAt{9}{255.655}{715.836}{\Character{47}}%
\TextAt{14}{261.508}{715.836}{o}%
\TextAt{3}{267.352}{720.772}{\Character{3}}%
\TextAt{8}{284.039}{715.836}{and}%
\TextAt{8}{316.852}{715.836}{ker}%
\TextAt{9}{332.788}{715.836}{\Character{40}}%
\TextAt{13}{337.341}{715.836}{\Character{12}}%
\TextAt{9}{344.612}{715.836}{\Character{41}}%
\TextAt{9}{352.467}{715.836}{\Character{61}}%
\TextAt{20}{364.911}{727.193}{\Character{77}}%
\TextAt{27}{371.949}{702.087}{p}%
\TextAt{20}{384.57}{716.371}{\Character{101}}%
\TextAt{14}{384.969}{715.836}{o}%
\TextAt{3}{390.813}{720.772}{\Character{3}}%
\TextAt{27}{390.813}{712.88}{p}%
\TextAt{9}{395.545}{715.836}{\Character{47}}%
\TextAt{14}{401.398}{715.836}{o}%
\TextAt{3}{407.243}{720.772}{\Character{3}}%
\TextAt{27}{407.243}{712.88}{p}%
\TextAt{13}{411.975}{715.836}{\Character{58}}%
\TextAt{8}{79.226}{678.896}{This}%
\TextAt{8}{105.946}{678.896}{then}%
\TextAt{8}{132.606}{678.896}{yields}%
\TextAt{8}{165.507}{678.896}{the}%
\TextAt{8}{185.663}{678.896}{exact}%
\TextAt{8}{216.555}{678.896}{sequence}%
\TextAt{9}{190.034}{651.712}{1}%
\NativeArrowAt{199.208}{651.712}{11.955}{19.26}{\longrightarrow}%
\TextAt{14}{221.79}{651.712}{o}%
\TextAt{3}{227.634}{656.648}{\Character{3}}%
\NativeArrowAt{235.687}{651.712}{11.955}{19.26}{\longrightarrow}%
\TextAt{20}{257.87}{652.248}{\Character{101}}%
\TextAt{14}{258.269}{651.712}{o}%
\TextAt{3}{264.113}{656.648}{\Character{3}}%
\NativeArrowAt{272.167}{651.712}{11.955}{19.26}{\longrightarrow}%
\TextAt{20}{294.749}{663.07}{\Character{77}}%
\TextAt{27}{301.787}{637.964}{p}%
\TextAt{20}{314.408}{652.248}{\Character{101}}%
\TextAt{14}{314.807}{651.712}{o}%
\TextAt{3}{320.651}{656.648}{\Character{3}}%
\TextAt{27}{320.651}{648.757}{p}%
\TextAt{9}{325.383}{651.712}{\Character{47}}%
\TextAt{14}{331.236}{651.712}{o}%
\TextAt{3}{337.08}{656.648}{\Character{3}}%
\TextAt{27}{337.08}{648.757}{p}%
\NativeArrowAt{272.166}{619.344}{11.955}{19.26}{\longrightarrow}%
\TextAt{13}{294.749}{619.344}{P}%
\TextAt{13}{303.907}{619.344}{ic}%
\TextAt{9}{312.951}{619.344}{\Character{40}}%
\TextAt{14}{317.503}{619.344}{o}%
\TextAt{9}{323.348}{619.344}{\Character{41}}%
\NativeArrowAt{331.221}{619.344}{11.955}{19.26}{\longrightarrow}%
\TextAt{13}{353.803}{619.344}{P}%
\TextAt{13}{362.961}{619.344}{ic}%
\TextAt{9}{372.005}{619.344}{\Character{40}}%
\TextAt{20}{376.159}{619.88}{\Character{101}}%
\TextAt{14}{376.558}{619.344}{o}%
\TextAt{9}{382.402}{619.344}{\Character{41}}%
\NativeArrowAt{390.275}{619.344}{11.955}{19.26}{\longrightarrow}%
\TextAt{9}{412.857}{619.344}{1}%
\TextAt{13}{418.71}{619.344}{\Character{58}}%
\Vector{q 0 G 0 g 0.399 w 452.334 -180.032 m 452.334 -171.962 l S Q}%
\Vector{q 0 G 0 g 0.399 w 452.534 -172.161 m 459.707 -172.161 l S Q}%
\Vector{q 0 G 0 g 0.399 w 452.534 -179.832 m 459.707 -179.832 l S Q}%
\Vector{q 0 G 0 g 0.399 w 459.906 -180.032 m 459.906 -171.962 l S Q}%
\TextAt{8}{97.159}{565.184}{A}%
\TextAt{8}{109.41}{565.184}{prime}%
\TextAt{8}{142.153}{565.184}{ideal}%
\TextAt{14}{169.696}{565.184}{p}%
\NativeMathAt{179}{565.184}{11.955}{\neq}%
\TextAt{9}{191.41}{565.184}{0}%
\TextAt{8}{200.755}{565.184}{of}%
\TextAt{14}{213.664}{565.184}{o}%
\TextAt{8}{222.985}{565.184}{is}%
\TextAt{8}{234.329}{565.184}{called}%
\TextAt{12}{267.072}{565.184}{regular}%
\TextAt{8}{312.164}{565.184}{if}%
\TextAt{14}{322.467}{565.184}{o}%
\TextAt{27}{328.312}{563.391}{p}%
\TextAt{8}{336.275}{565.184}{is}%
\TextAt{8}{347.618}{565.184}{integrally}%
\TextAt{8}{399.239}{565.184}{closed,}%
\TextAt{8}{436.68}{565.184}{and}%
\TextAt{8}{459.022}{565.184}{thus}%
\TextAt{8}{484.353}{565.184}{a}%
\TextAt{8}{493.688}{565.184}{discrete}%
\TextAt{8}{79.226}{548.138}{valuation}%
\TextAt{8}{131.206}{548.138}{ring.}%
\TextAt{8}{162.545}{548.138}{For}%
\TextAt{8}{184.421}{548.138}{the}%
\TextAt{8}{205.485}{548.138}{regular}%
\TextAt{8}{246.072}{548.138}{prime}%
\TextAt{8}{280.143}{548.138}{ideals,}%
\TextAt{8}{317.106}{548.138}{the}%
\TextAt{8}{338.17}{548.138}{summands}%
\TextAt{20}{396.685}{548.674}{\Character{101}}%
\TextAt{14}{397.084}{548.138}{o}%
\TextAt{3}{402.928}{552.477}{\Character{3}}%
\TextAt{27}{402.928}{545.183}{p}%
\TextAt{9}{407.661}{548.138}{\Character{47}}%
\TextAt{14}{413.514}{548.138}{o}%
\TextAt{3}{419.358}{552.477}{\Character{3}}%
\TextAt{27}{419.358}{545.183}{p}%
\TextAt{8}{428.895}{548.138}{in}%
\TextAt{8}{443.455}{548.138}{(12.9)}%
\TextAt{8}{478.195}{548.138}{are}%
\TextAt{8}{498.614}{548.138}{trivial.}%
\TextAt{8}{79.226}{531.092}{There}%
\TextAt{8}{114.136}{531.092}{are}%
\TextAt{8}{134.748}{531.092}{only}%
\TextAt{8}{161.54}{531.092}{finitely}%
\TextAt{8}{201.986}{531.092}{many}%
\TextAt{8}{234.959}{531.092}{non-regular}%
\TextAt{8}{298.502}{531.092}{prime}%
\TextAt{8}{332.766}{531.092}{ideals}%
\TextAt{8}{366.445}{531.092}{of}%
\TextAt{14}{380.876}{531.092}{o}%
\TextAt{8}{386.72}{531.092}{,}%
\TextAt{8}{395.245}{531.092}{namely}%
\TextAt{8}{436.994}{531.092}{the}%
\TextAt{8}{458.251}{531.092}{divisors}%
\TextAt{8}{502.08}{531.092}{of}%
\TextAt{8}{516.51}{531.092}{the}%
\TextAt{12}{79.226}{514.046}{conductor}%
\TextAt{8}{142.268}{514.046}{of}%
\TextAt{14}{156.174}{514.046}{o}%
\TextAt{8}{162.019}{514.046}{.}%
\TextAt{8}{172.196}{514.046}{This}%
\TextAt{8}{199.492}{514.046}{is}%
\TextAt{8}{211.832}{514.046}{by}%
\TextAt{8}{228.667}{514.046}{definition}%
\TextAt{8}{281.272}{514.046}{the}%
\TextAt{8}{302.005}{514.046}{biggest}%
\TextAt{8}{342.32}{514.046}{ideal}%
\TextAt{8}{370.859}{514.046}{of}%
\TextAt{20}{384.367}{514.582}{\Character{101}}%
\TextAt{14}{384.765}{514.046}{o}%
\TextAt{8}{395.083}{514.046}{which}%
\TextAt{8}{429.146}{514.046}{is}%
\TextAt{8}{441.486}{514.046}{contained}%
\TextAt{8}{495.071}{514.046}{in}%
\TextAt{14}{509.3}{514.046}{o}%
\TextAt{8}{515.145}{514.046}{,}%
\TextAt{8}{523.014}{514.046}{in}%
\TextAt{8}{79.226}{497}{other}%
\TextAt{8}{109.796}{497}{words,}%
\TextAt{14}{253.402}{479.954}{f}%
\TextAt{9}{260.624}{479.954}{\Character{61}}%
\TextAt{16}{273.049}{479.954}{\Character{102}}%
\TextAt{13}{279.027}{479.954}{a}%
\NativeMathAt{288.493}{479.954}{11.955}{\in}%
\TextAt{20}{299.385}{480.49}{\Character{101}}%
\TextAt{14}{299.784}{479.954}{o}%
\TextAt{16}{308.949}{479.954}{\Character{106}}%
\TextAt{13}{315.59}{479.954}{a}%
\TextAt{20}{321.337}{480.49}{\Character{101}}%
\TextAt{14}{321.735}{479.954}{o}%
\NativeMathAt{330.9}{479.954}{11.955}{\subseteq}%
\TextAt{14}{343.52}{479.954}{o}%
\TextAt{16}{349.364}{479.954}{\Character{103}}%
\TextAt{13}{355.342}{479.954}{\Character{58}}%
\TextAt{8}{79.226}{456.758}{Since}%
\TextAt{20}{109.385}{457.294}{\Character{101}}%
\TextAt{14}{109.784}{456.758}{o}%
\TextAt{8}{119.525}{456.758}{is}%
\TextAt{8}{131.289}{456.758}{a}%
\TextAt{8}{141.044}{456.758}{finitely}%
\TextAt{8}{180.389}{456.758}{generated}%
\TextAt{14}{233.721}{456.758}{o}%
\TextAt{8}{239.565}{456.758}{-module,}%
\TextAt{8}{288.008}{456.758}{we}%
\TextAt{8}{305.235}{456.758}{have}%
\TextAt{14}{332.23}{456.758}{f}%
\NativeMathAt{339.451}{456.758}{11.955}{\neq}%
\TextAt{9}{351.861}{456.758}{0}%
\TextAt{8}{357.729}{456.758}{.}%
\TextAt{12}{79.226}{430.963}{Proposition}%
\TextAt{12}{151.937}{430.963}{(12.10).}%
\TextAt{8}{200.715}{430.963}{For}%
\TextAt{8}{221.684}{430.963}{any}%
\TextAt{8}{243.801}{430.963}{prime}%
\TextAt{8}{276.965}{430.963}{ideal}%
\TextAt{14}{304.928}{430.963}{p}%
\NativeMathAt{314.232}{430.963}{11.955}{\neq}%
\TextAt{9}{326.642}{430.963}{0}%
\TextAt{8}{336.408}{430.963}{of}%
\TextAt{14}{349.738}{430.963}{o}%
\TextAt{8}{359.479}{430.963}{one}%
\TextAt{8}{380.939}{430.963}{has}%
\TextAt{14}{239.501}{402.721}{p}%
\TextAt{18}{248.805}{402.721}{\Character{45}}%
\TextAt{14}{255.447}{402.721}{f}%
\NativeArrowAt{274.623}{402.721}{11.955}{21.926}{\Longleftrightarrow}%
\TextAt{14}{311.817}{402.721}{p}%
\TextAt{8}{321.697}{402.721}{is}%
\TextAt{8}{333.461}{402.721}{regular}%
\TextAt{13}{369.243}{402.721}{\Character{58}}%
\TextAt{8}{79.226}{374.479}{If}%
\TextAt{8}{90.918}{374.479}{this}%
\TextAt{8}{113.741}{374.479}{is}%
\TextAt{8}{125.505}{374.479}{the}%
\TextAt{8}{145.661}{374.479}{case,}%
\TextAt{8}{173.684}{374.479}{then}%
\TextAt{20}{200.015}{375.596}{\Character{101}}%
\TextAt{14}{200.344}{374.479}{p}%
\TextAt{9}{209.648}{374.479}{\Character{61}}%
\TextAt{14}{222.073}{374.479}{p}%
\TextAt{20}{227.658}{375.014}{\Character{101}}%
\TextAt{14}{228.056}{374.479}{o}%
\TextAt{8}{237.798}{374.479}{is}%
\TextAt{8}{249.562}{374.479}{a}%
\TextAt{8}{259.317}{374.479}{prime}%
\TextAt{8}{292.481}{374.479}{ideal}%
\TextAt{8}{320.444}{374.479}{of}%
\TextAt{20}{333.375}{375.014}{\Character{101}}%
\TextAt{14}{333.774}{374.479}{o}%
\TextAt{8}{343.516}{374.479}{and}%
\TextAt{14}{366.278}{374.479}{o}%
\TextAt{27}{372.123}{372.686}{p}%
\TextAt{9}{379.93}{374.479}{\Character{61}}%
\TextAt{20}{391.957}{375.014}{\Character{101}}%
\TextAt{14}{392.356}{374.479}{o}%
\TextAt{30}{397.842}{372.841}{\Character{101}}%
\TextAt{27}{398.2}{372.096}{p}%
\TextAt{8}{402.687}{374.479}{.}%
\TextAt{19}{79.226}{339.934}{Proof.}%
\TextAt{8}{116}{339.934}{Assume}%
\TextAt{14}{160.006}{339.934}{p}%
\TextAt{18}{170.408}{339.934}{\Character{45}}%
\TextAt{14}{178.148}{339.934}{f}%
\TextAt{8}{182.049}{339.934}{,}%
\TextAt{8}{190.003}{339.934}{i.e.,}%
\TextAt{14}{212.914}{339.934}{p}%
\TextAt{18}{223.316}{339.934}{\Character{43}}%
\TextAt{14}{237.034}{339.934}{f}%
\TextAt{8}{240.934}{339.934}{,}%
\TextAt{8}{248.889}{339.934}{and}%
\TextAt{8}{272.296}{339.934}{let}%
\TextAt{13}{289.845}{339.934}{t}%
\NativeMathAt{298.492}{339.934}{11.955}{\in}%
\TextAt{14}{310.881}{339.934}{f}%
\TextAt{16}{317.877}{339.934}{\Character{110}}%
\TextAt{14}{326.951}{339.934}{p}%
\TextAt{8}{332.934}{339.934}{.}%
\TextAt{8}{343.316}{339.934}{Then}%
\TextAt{13}{374.517}{339.934}{t}%
\TextAt{20}{378.346}{340.47}{\Character{101}}%
\TextAt{14}{378.745}{339.934}{o}%
\NativeMathAt{389.008}{339.934}{11.955}{\subseteq}%
\TextAt{14}{402.726}{339.934}{o}%
\TextAt{8}{408.57}{339.934}{,}%
\TextAt{8}{416.525}{339.934}{hence}%
\TextAt{20}{449.276}{340.47}{\Character{101}}%
\TextAt{14}{449.675}{339.934}{o}%
\NativeMathAt{459.938}{339.934}{11.955}{\subseteq}%
\TextAt{17}{474.851}{344.641}{1}%
\Vector{q 0 G 0 g 0.478 w 402.851 -431.077 m 407.086 -431.077 l S Q}%
\TextAt{21}{475.44}{335.812}{t}%
\TextAt{14}{480.281}{339.934}{o}%
\NativeMathAt{490.544}{339.934}{11.955}{\subseteq}%
\TextAt{14}{504.262}{339.934}{o}%
\TextAt{27}{510.106}{338.141}{p}%
\TextAt{8}{514.593}{339.934}{.}%
\TextAt{8}{524.975}{339.934}{If}%
\TextAt{14}{79.226}{322.888}{m}%
\TextAt{9}{92.136}{322.888}{\Character{61}}%
\TextAt{14}{104.983}{322.888}{po}%
\TextAt{27}{116.81}{321.095}{p}%
\TextAt{8}{125.442}{322.888}{is}%
\TextAt{8}{137.454}{322.888}{the}%
\TextAt{8}{157.858}{322.888}{maximal}%
\TextAt{8}{205.914}{322.888}{ideal}%
\TextAt{8}{234.125}{322.888}{of}%
\TextAt{14}{247.702}{322.888}{o}%
\TextAt{27}{253.546}{321.095}{p}%
\TextAt{8}{262.178}{322.888}{then,}%
\TextAt{8}{292.399}{322.888}{putting}%
\TextAt{20}{333.946}{324.005}{\Character{101}}%
\TextAt{14}{334.275}{322.888}{p}%
\TextAt{9}{344.001}{322.888}{\Character{61}}%
\TextAt{14}{356.848}{322.888}{m}%
\TextAt{16}{368.84}{322.888}{\Character{92}}%
\TextAt{20}{379.237}{323.424}{\Character{101}}%
\TextAt{14}{379.636}{322.888}{o}%
\TextAt{8}{385.48}{322.888}{,}%
\TextAt{20}{392.609}{324.005}{\Character{101}}%
\TextAt{14}{392.939}{322.888}{p}%
\TextAt{8}{403.067}{322.888}{is}%
\TextAt{8}{415.078}{322.888}{a}%
\TextAt{8}{425.081}{322.888}{prime}%
\TextAt{8}{458.492}{322.888}{ideal}%
\TextAt{8}{486.703}{322.888}{of}%
\TextAt{20}{499.882}{323.424}{\Character{101}}%
\TextAt{14}{500.281}{322.888}{o}%
\TextAt{8}{510.27}{322.888}{such}%
\TextAt{8}{79.226}{305.842}{that}%
\TextAt{14}{104.314}{305.842}{p}%
\NativeMathAt{113.618}{305.842}{11.955}{\subseteq}%
\TextAt{20}{125.908}{306.959}{\Character{101}}%
\TextAt{14}{126.237}{305.842}{p}%
\TextAt{16}{134.302}{305.842}{\Character{92}}%
\TextAt{14}{144.353}{305.842}{o}%
\TextAt{8}{150.198}{305.842}{,}%
\TextAt{8}{157.122}{305.842}{hence}%
\TextAt{14}{189.347}{305.842}{p}%
\TextAt{9}{198.651}{305.842}{\Character{61}}%
\TextAt{20}{210.747}{306.959}{\Character{101}}%
\TextAt{14}{211.076}{305.842}{p}%
\TextAt{16}{219.141}{305.842}{\Character{92}}%
\TextAt{14}{229.192}{305.842}{o}%
\TextAt{8}{238.653}{305.842}{because}%
\TextAt{14}{281.673}{305.842}{p}%
\TextAt{8}{291.272}{305.842}{is}%
\TextAt{8}{302.755}{305.842}{maximal.}%
\TextAt{8}{355.021}{305.842}{Trivially,}%
\TextAt{14}{404.218}{305.842}{o}%
\TextAt{27}{410.063}{304.049}{p}%
\NativeMathAt{417.87}{305.842}{11.955}{\subseteq}%
\TextAt{20}{430.091}{306.378}{\Character{101}}%
\TextAt{14}{430.49}{305.842}{o}%
\TextAt{30}{435.976}{304.204}{\Character{101}}%
\TextAt{27}{436.334}{303.46}{p}%
\TextAt{8}{440.821}{305.842}{,}%
\TextAt{8}{447.745}{305.842}{and}%
\TextAt{8}{470.226}{305.842}{if}%
\TextAt{8}{480.669}{305.842}{conversely}%
\TextAt{13}{79.226}{288.796}{a}%
\TextAt{9}{85.371}{288.796}{\Character{47}}%
\TextAt{13}{91.224}{288.796}{s}%
\NativeMathAt{101.648}{288.796}{11.955}{\in}%
\TextAt{20}{114.129}{289.332}{\Character{101}}%
\TextAt{14}{114.528}{288.796}{o}%
\TextAt{30}{120.014}{287.158}{\Character{101}}%
\TextAt{27}{120.372}{286.414}{p}%
\TextAt{8}{124.859}{288.796}{,}%
\TextAt{8}{133.174}{288.796}{for}%
\TextAt{13}{151.991}{288.796}{a}%
\NativeMathAt{163.046}{288.796}{11.955}{\in}%
\TextAt{20}{175.527}{289.332}{\Character{101}}%
\TextAt{14}{175.926}{288.796}{o}%
\TextAt{8}{181.77}{288.796}{,}%
\TextAt{13}{190.085}{288.796}{s}%
\NativeMathAt{200.509}{288.796}{11.955}{\in}%
\TextAt{20}{212.99}{289.332}{\Character{101}}%
\TextAt{14}{213.388}{288.796}{o}%
\TextAt{16}{222.525}{288.796}{\Character{110}}%
\TextAt{20}{231.466}{289.913}{\Character{101}}%
\TextAt{14}{231.795}{288.796}{p}%
\TextAt{8}{237.778}{288.796}{,}%
\TextAt{8}{246.093}{288.796}{then}%
\TextAt{13}{273.685}{288.796}{ta}%
\NativeMathAt{288.967}{288.796}{11.955}{\in}%
\TextAt{14}{301.847}{288.796}{o}%
\TextAt{8}{312.521}{288.796}{and}%
\TextAt{13}{336.216}{288.796}{ts}%
\NativeMathAt{350.867}{288.796}{11.955}{\in}%
\TextAt{14}{363.747}{288.796}{o}%
\TextAt{16}{372.884}{288.796}{\Character{110}}%
\TextAt{14}{382.153}{288.796}{p}%
\TextAt{8}{388.137}{288.796}{,}%
\TextAt{8}{396.451}{288.796}{hence}%
\TextAt{13}{429.89}{288.796}{a}%
\TextAt{9}{436.035}{288.796}{\Character{47}}%
\TextAt{13}{441.888}{288.796}{s}%
\TextAt{9}{452.312}{288.796}{\Character{61}}%
\TextAt{13}{466.326}{288.796}{ta}%
\TextAt{9}{476.698}{288.796}{\Character{47}\Character{40}}%
\TextAt{13}{487.103}{288.796}{ts}%
\TextAt{9}{496.845}{288.796}{\Character{41}}%
\NativeMathAt{506.307}{288.796}{11.955}{\in}%
\TextAt{14}{519.187}{288.796}{o}%
\TextAt{27}{525.031}{287.003}{p}%
\TextAt{8}{529.518}{288.796}{.}%
\TextAt{8}{79.226}{271.75}{Therefore}%
\TextAt{14}{132.223}{271.75}{o}%
\TextAt{27}{138.068}{269.957}{p}%
\TextAt{9}{145.875}{271.75}{\Character{61}}%
\TextAt{20}{157.902}{272.286}{\Character{101}}%
\TextAt{14}{158.301}{271.75}{o}%
\TextAt{30}{163.787}{270.112}{\Character{101}}%
\TextAt{27}{164.145}{269.368}{p}%
\TextAt{8}{168.632}{271.75}{.}%
\TextAt{8}{177.08}{271.75}{Thus,}%
\TextAt{8}{209.981}{271.75}{by}%
\TextAt{8}{226.24}{271.75}{(11.5),}%
\TextAt{14}{263.325}{271.75}{o}%
\TextAt{27}{269.169}{269.957}{p}%
\TextAt{8}{277.553}{271.75}{is}%
\TextAt{8}{289.317}{271.75}{a}%
\TextAt{8}{299.073}{271.75}{valuation}%
\TextAt{8}{350.145}{271.75}{ring,}%
\TextAt{8}{377.463}{271.75}{i.e.,}%
\TextAt{14}{399.568}{271.75}{p}%
\TextAt{8}{409.448}{271.75}{is}%
\TextAt{8}{421.212}{271.75}{regular.}%
\TextAt{8}{97.159}{251.716}{One}%
\TextAt{8}{121.895}{251.716}{has}%
\TextAt{8}{142.793}{251.716}{furthermore}%
\TextAt{8}{207.532}{251.716}{that}%
\TextAt{20}{232.596}{252.832}{\Character{101}}%
\TextAt{14}{232.925}{251.716}{p}%
\TextAt{9}{242.271}{251.716}{\Character{61}}%
\TextAt{14}{254.739}{251.716}{p}%
\TextAt{20}{260.323}{252.251}{\Character{101}}%
\TextAt{14}{260.722}{251.716}{o}%
\TextAt{8}{266.566}{251.716}{.}%
\TextAt{8}{275.089}{251.716}{In}%
\TextAt{8}{289.735}{251.716}{fact,}%
\TextAt{20}{315.773}{252.832}{\Character{101}}%
\TextAt{14}{316.103}{251.716}{p}%
\TextAt{8}{326.008}{251.716}{is}%
\TextAt{8}{337.797}{251.716}{the}%
\TextAt{8}{357.978}{251.716}{only}%
\TextAt{8}{383.694}{251.716}{prime}%
\TextAt{8}{416.883}{251.716}{ideal}%
\TextAt{8}{444.87}{251.716}{of}%
\TextAt{20}{457.826}{252.251}{\Character{101}}%
\TextAt{14}{458.225}{251.716}{o}%
\TextAt{8}{467.992}{251.716}{above}%
\TextAt{14}{501.192}{251.716}{p}%
\TextAt{8}{507.175}{251.716}{.}%
\TextAt{8}{515.698}{251.716}{For}%
\TextAt{8}{79.226}{234.67}{if}%
\TextAt{20}{89.553}{235.205}{\Character{101}}%
\TextAt{14}{89.95}{234.67}{q}%
\TextAt{8}{99.695}{234.67}{is}%
\TextAt{8}{111.459}{234.67}{another}%
\TextAt{8}{154.39}{234.67}{one,}%
\TextAt{8}{179.101}{234.67}{then}%
\TextAt{20}{205.362}{235.205}{\Character{101}}%
\TextAt{14}{205.761}{234.67}{o}%
\TextAt{30}{211.247}{233.031}{\Character{101}}%
\TextAt{27}{211.605}{232.287}{p}%
\TextAt{9}{219.413}{234.67}{\Character{61}}%
\TextAt{14}{231.838}{234.67}{o}%
\TextAt{27}{237.683}{232.876}{p}%
\NativeMathAt{245.49}{234.67}{11.955}{\subseteq}%
\TextAt{20}{257.711}{235.205}{\Character{101}}%
\TextAt{14}{258.11}{234.67}{o}%
\TextAt{30}{263.551}{233.031}{\Character{101}}%
\TextAt{27}{263.954}{232.674}{q}%
\TextAt{8}{268.35}{234.67}{,}%
\TextAt{8}{275.5}{234.67}{and}%
\TextAt{8}{298.262}{234.67}{therefore}%
\TextAt{20}{239.082}{207.544}{\Character{101}}%
\TextAt{14}{239.411}{206.428}{p}%
\TextAt{9}{248.715}{206.428}{\Character{61}}%
\TextAt{20}{260.742}{206.963}{\Character{101}}%
\TextAt{14}{261.141}{206.428}{o}%
\TextAt{16}{269.642}{206.428}{\Character{92}}%
\TextAt{20}{279.939}{207.544}{\Character{101}}%
\TextAt{14}{280.269}{206.428}{p}%
\TextAt{20}{285.853}{206.963}{\Character{101}}%
\TextAt{14}{286.252}{206.428}{o}%
\TextAt{30}{291.738}{204.789}{\Character{101}}%
\TextAt{27}{292.096}{204.045}{p}%
\NativeMathAt{299.904}{206.428}{11.955}{\subseteq}%
\TextAt{20}{312.124}{206.963}{\Character{101}}%
\TextAt{14}{312.523}{206.428}{o}%
\TextAt{16}{321.024}{206.428}{\Character{92}}%
\TextAt{20}{331.253}{206.963}{\Character{101}}%
\TextAt{14}{331.651}{206.428}{q}%
\TextAt{20}{337.099}{206.963}{\Character{101}}%
\TextAt{14}{337.498}{206.428}{o}%
\TextAt{30}{342.939}{204.789}{\Character{101}}%
\TextAt{27}{343.342}{204.433}{q}%
\TextAt{9}{351.06}{206.428}{\Character{61}}%
\TextAt{20}{363.088}{206.963}{\Character{101}}%
\TextAt{14}{363.485}{206.428}{q}%
\TextAt{13}{369.332}{206.428}{\Character{59}}%
\TextAt{8}{79.226}{178.186}{hence}%
\TextAt{20}{111.941}{179.303}{\Character{101}}%
\TextAt{14}{112.27}{178.186}{p}%
\TextAt{9}{122.49}{178.186}{\Character{61}}%
\TextAt{20}{135.435}{178.721}{\Character{101}}%
\TextAt{14}{135.832}{178.186}{q}%
\TextAt{8}{141.68}{178.186}{.}%
\TextAt{8}{151.742}{178.186}{Consequently,}%
\TextAt{14}{227.265}{178.186}{p}%
\TextAt{20}{232.849}{178.721}{\Character{101}}%
\TextAt{14}{233.248}{178.186}{o}%
\TextAt{9}{243.33}{178.186}{\Character{61}}%
\TextAt{20}{256.342}{179.303}{\Character{101}}%
\TextAt{14}{256.672}{178.186}{p}%
\TextAt{21}{262.655}{182.524}{e}%
\TextAt{8}{267.081}{178.186}{,}%
\TextAt{8}{274.903}{178.186}{with}%
\TextAt{13}{302.101}{178.186}{e}%
\NativeMathAt{311.763}{178.186}{11.955}{\geq}%
\TextAt{9}{325.299}{178.186}{1}%
\TextAt{8}{331.152}{178.186}{,}%
\TextAt{8}{338.974}{178.186}{and}%
\TextAt{8}{362.274}{178.186}{furthermore}%
\TextAt{14}{427.525}{178.186}{m}%
\TextAt{9}{440.929}{178.186}{\Character{61}}%
\TextAt{14}{454.271}{178.186}{po}%
\TextAt{27}{466.098}{176.393}{p}%
\TextAt{9}{474.823}{178.186}{\Character{61}}%
\TextAt{9}{488.153}{178.186}{\Character{40}}%
\TextAt{14}{492.717}{178.186}{p}%
\TextAt{20}{498.301}{178.721}{\Character{101}}%
\TextAt{14}{498.7}{178.186}{o}%
\TextAt{9}{504.544}{178.186}{\Character{41}}%
\TextAt{14}{509.096}{178.186}{o}%
\TextAt{27}{514.941}{176.393}{p}%
\TextAt{9}{523.665}{178.186}{\Character{61}}%
\TextAt{20}{78.897}{162.257}{\Character{101}}%
\TextAt{14}{79.226}{161.14}{p}%
\TextAt{21}{85.209}{165.478}{e}%
\TextAt{14}{89.636}{161.14}{o}%
\TextAt{27}{95.48}{159.347}{p}%
\TextAt{9}{103.288}{161.14}{\Character{61}}%
\TextAt{14}{115.713}{161.14}{m}%
\TextAt{21}{124.88}{165.478}{e}%
\TextAt{8}{129.306}{161.14}{,}%
\TextAt{8}{136.456}{161.14}{i.e.,}%
\TextAt{13}{158.561}{161.14}{e}%
\TextAt{9}{167.307}{161.14}{\Character{61}}%
\TextAt{9}{179.717}{161.14}{1}%
\TextAt{8}{189.483}{161.14}{and}%
\TextAt{8}{212.245}{161.14}{thus}%
\TextAt{20}{237.668}{162.257}{\Character{101}}%
\TextAt{14}{237.997}{161.14}{p}%
\TextAt{9}{247.301}{161.14}{\Character{61}}%
\TextAt{14}{259.726}{161.14}{p}%
\TextAt{20}{265.311}{161.675}{\Character{101}}%
\TextAt{14}{265.709}{161.14}{o}%
\TextAt{8}{271.554}{161.14}{.}%
\TextAt{8}{97.159}{141.105}{Conversely,}%
\TextAt{8}{158.926}{141.105}{assume}%
\TextAt{14}{199.567}{141.105}{o}%
\TextAt{27}{205.412}{139.312}{p}%
\TextAt{8}{213.993}{141.105}{is}%
\TextAt{8}{225.953}{141.105}{a}%
\TextAt{8}{235.906}{141.105}{discrete}%
\TextAt{8}{279.081}{141.105}{valuation}%
\TextAt{8}{330.351}{141.105}{ring.}%
\TextAt{8}{359.558}{141.105}{Being}%
\TextAt{8}{392.751}{141.105}{a}%
\TextAt{8}{402.703}{141.105}{principal}%
\TextAt{8}{451.677}{141.105}{ideal}%
\TextAt{8}{479.837}{141.105}{domain,}%
\TextAt{8}{524.963}{141.105}{it}%
\TextAt{8}{79.226}{124.059}{is}%
\TextAt{8}{90.663}{124.059}{integrally}%
\TextAt{8}{142.377}{124.059}{closed,}%
\TextAt{8}{179.894}{124.059}{and}%
\TextAt{8}{202.33}{124.059}{since}%
\TextAt{20}{230.273}{124.595}{\Character{101}}%
\TextAt{14}{230.672}{124.059}{o}%
\TextAt{8}{240.087}{124.059}{is}%
\TextAt{8}{251.524}{124.059}{integral}%
\TextAt{8}{293.806}{124.059}{over}%
\TextAt{14}{318.525}{124.059}{o}%
\TextAt{8}{324.369}{124.059}{,}%
\TextAt{8}{331.257}{124.059}{hence}%
\TextAt{8}{363.437}{124.059}{a}%
\TextAt{8}{372.865}{124.059}{fortiori}%
\TextAt{8}{411.895}{124.059}{over}%
\TextAt{14}{436.614}{124.059}{o}%
\TextAt{27}{442.459}{122.266}{p}%
\TextAt{8}{446.946}{124.059}{,}%
\TextAt{8}{453.833}{124.059}{we}%
\TextAt{8}{470.734}{124.059}{have}%
\TextAt{20}{497.003}{124.595}{\Character{101}}%
\TextAt{14}{497.402}{124.059}{o}%
\NativeMathAt{506.567}{124.059}{11.955}{\subseteq}%
\TextAt{14}{519.187}{124.059}{o}%
\TextAt{27}{525.031}{122.266}{p}%
\TextAt{8}{529.518}{124.059}{.}%
\TextAt{8}{79.226}{107.013}{Let}%
\TextAt{13}{101.059}{107.013}{x}%
\TextAt{17}{107.711}{105.22}{1}%
\TextAt{13}{112.443}{107.013}{\Character{59}}%
\TextAt{13}{117.667}{107.013}{\Character{58}}%
\TextAt{13}{122.904}{107.013}{\Character{58}}%
\TextAt{13}{128.14}{107.013}{\Character{58}}%
\TextAt{13}{133.365}{107.013}{\Character{59}}%
\TextAt{13}{138.601}{107.013}{x}%
\TextAt{21}{145.316}{105.22}{n}%
\TextAt{8}{155.713}{107.013}{be}%
\TextAt{8}{172.501}{107.013}{a}%
\TextAt{8}{183.119}{107.013}{system}%
\TextAt{8}{222.801}{107.013}{of}%
\TextAt{8}{236.994}{107.013}{generators}%
\TextAt{8}{294.513}{107.013}{of}%
\TextAt{8}{308.706}{107.013}{the}%
\TextAt{14}{329.726}{107.013}{o}%
\TextAt{8}{335.57}{107.013}{-module}%
\TextAt{20}{381.225}{107.549}{\Character{101}}%
\TextAt{14}{381.624}{107.013}{o}%
\TextAt{8}{387.468}{107.013}{.}%
\TextAt{8}{398.506}{107.013}{We}%
\TextAt{8}{419.514}{107.013}{may}%
\TextAt{8}{445.746}{107.013}{write}%
\TextAt{13}{476.521}{107.013}{x}%
\TextAt{21}{483.173}{105.22}{i}%
\TextAt{9}{491.347}{107.013}{\Character{61}}%
\TextAt{13}{505.243}{107.013}{a}%
\TextAt{21}{511.388}{105.22}{i}%
\TextAt{9}{514.769}{107.013}{\Character{47}}%
\TextAt{13}{520.622}{107.013}{s}%
\TextAt{21}{526.136}{105.22}{i}%
\TextAt{8}{529.518}{107.013}{,}%
\TextAt{8}{79.226}{89.967}{with}%
\TextAt{13}{105.576}{89.967}{a}%
\TextAt{21}{111.721}{88.174}{i}%
\NativeMathAt{118.423}{89.967}{11.955}{\in}%
\TextAt{14}{129.714}{89.967}{o}%
\TextAt{8}{135.559}{89.967}{,}%
\TextAt{13}{142.46}{89.967}{s}%
\TextAt{21}{147.974}{88.174}{i}%
\NativeMathAt{154.676}{89.967}{11.955}{\in}%
\TextAt{14}{165.967}{89.967}{o}%
\TextAt{16}{173.835}{89.967}{\Character{110}}%
\TextAt{14}{181.835}{89.967}{p}%
\TextAt{8}{187.819}{89.967}{.}%
\TextAt{8}{196.164}{89.967}{Setting}%
\TextAt{13}{236.179}{89.967}{s}%
\TextAt{9}{245.013}{89.967}{\Character{61}}%
\TextAt{13}{257.439}{89.967}{s}%
\TextAt{17}{262.953}{88.174}{1}%
\TextAt{16}{269.678}{89.967}{\Character{1}}%
\TextAt{16}{274.974}{89.967}{\Character{1}}%
\TextAt{16}{280.282}{89.967}{\Character{1}}%
\TextAt{13}{285.618}{89.967}{s}%
\TextAt{21}{291.132}{88.174}{n}%
\NativeMathAt{300.089}{89.967}{11.955}{\in}%
\TextAt{14}{311.38}{89.967}{o}%
\TextAt{16}{319.248}{89.967}{\Character{110}}%
\TextAt{14}{327.248}{89.967}{p}%
\TextAt{8}{333.231}{89.967}{,}%
\TextAt{8}{340.133}{89.967}{we}%
\TextAt{8}{357.05}{89.967}{find}%
\TextAt{13}{380.149}{89.967}{sx}%
\TextAt{17}{392.315}{88.174}{1}%
\TextAt{13}{397.047}{89.967}{\Character{59}}%
\TextAt{13}{402.271}{89.967}{\Character{58}}%
\TextAt{13}{407.508}{89.967}{\Character{58}}%
\TextAt{13}{412.744}{89.967}{\Character{58}}%
\TextAt{13}{417.969}{89.967}{\Character{59}}%
\TextAt{13}{423.205}{89.967}{sx}%
\TextAt{21}{435.434}{88.174}{n}%
\NativeMathAt{444.391}{89.967}{11.955}{\in}%
\TextAt{14}{455.682}{89.967}{o}%
\TextAt{8}{465.114}{89.967}{and}%
\TextAt{8}{487.567}{89.967}{therefore}%
\TextAt{13}{79.226}{72.921}{s}%
\TextAt{20}{84.341}{73.457}{\Character{101}}%
\TextAt{14}{84.74}{72.921}{o}%
\NativeMathAt{93.905}{72.921}{11.955}{\subseteq}%
\TextAt{14}{106.525}{72.921}{o}%
\TextAt{8}{112.369}{72.921}{,}%
\TextAt{8}{119.518}{72.921}{i.e.,}%
\TextAt{13}{141.623}{72.921}{s}%
\NativeMathAt{150.458}{72.921}{11.955}{\in}%
\TextAt{14}{161.749}{72.921}{f}%
\TextAt{16}{168.306}{72.921}{\Character{110}}%
\TextAt{14}{176.94}{72.921}{p}%
\TextAt{8}{182.923}{72.921}{.}%
\TextAt{8}{191.372}{72.921}{It}%
\TextAt{8}{204.044}{72.921}{follows}%
\TextAt{8}{242.48}{72.921}{that}%
\TextAt{14}{267.849}{72.921}{p}%
\TextAt{18}{277.153}{72.921}{\Character{45}}%
\TextAt{14}{283.795}{72.921}{f}%
\TextAt{8}{287.695}{72.921}{.}%
\Vector{q 0 G 0 g 0.399 w 452.334 -701.079 m 452.334 -693.009 l S Q}%
\Vector{q 0 G 0 g 0.399 w 452.534 -693.208 m 459.707 -693.208 l S Q}%
\Vector{q 0 G 0 g 0.399 w 452.534 -700.879 m 459.707 -700.879 l S Q}%
\Vector{q 0 G 0 g 0.399 w 459.906 -701.079 m 459.906 -693.009 l S Q}%
\TextAt{8}{300.14}{43.033}{74}%
\Destination{78.23}{809.815}{nav.page.74}%
\Destination{79.36}{354.955}{nav.print.80}%
\Destination{79.36}{444.955}{nav.number.1.12.10}%
\Hotlink{390.456}{758.955}{14.294}{11.047}{nav.chapter.3}%
\Hotlink{417.341}{758.955}{6.858}{11.047}{nav.section.3.3}%
\Hotlink{332.505}{758.955}{19.22}{11.047}{nav.cite.N023}%
\Hotlink{442.955}{545.955}{30.936}{11.047}{nav.number.1.12.9}%
\Hotlink{225.74}{269.555}{30.936}{11.047}{nav.number.1.11.5}%
}

% Page 87 (body page 75)
\Page{%
\TextAt{2}{79.226}{796.943}{I.}%
\TextAt{2}{89.83}{796.943}{Algebraic}%
\TextAt{2}{138.648}{796.943}{Integers}%
\Vector{q 0 G 0 g 0.299 w 7.226 18.729 m 460.769 18.729 l S Q}%
\TextAt{8}{97.159}{761.124}{We}%
\TextAt{8}{117.303}{761.124}{now}%
\TextAt{8}{141.692}{761.124}{obtain}%
\TextAt{8}{178.119}{761.124}{the}%
\TextAt{8}{198.276}{761.124}{following}%
\TextAt{8}{247.71}{761.124}{simple}%
\TextAt{8}{284.186}{761.124}{description}%
\TextAt{8}{344.081}{761.124}{for}%
\TextAt{8}{361.966}{761.124}{the}%
\TextAt{8}{382.122}{761.124}{sum}%
\TextAt{20}{406.893}{770.09}{\Character{76}}%
\TextAt{27}{420.177}{757.637}{p}%
\TextAt{20}{426.258}{761.659}{\Character{101}}%
\TextAt{14}{426.656}{761.124}{o}%
\TextAt{3}{432.501}{765.462}{\Character{3}}%
\TextAt{27}{432.501}{758.168}{p}%
\TextAt{9}{437.233}{761.124}{\Character{47}}%
\TextAt{14}{443.086}{761.124}{o}%
\TextAt{3}{448.93}{765.462}{\Character{3}}%
\TextAt{27}{448.93}{758.168}{p}%
\TextAt{8}{457.56}{761.124}{in}%
\TextAt{8}{471.213}{761.124}{(12.9).}%
\TextAt{12}{79.226}{733.614}{Proposition}%
\TextAt{12}{151.937}{733.614}{(12.11).}%
\TextAt{20}{200.715}{744.972}{\Character{77}}%
\TextAt{27}{207.753}{719.866}{p}%
\TextAt{20}{220.374}{734.15}{\Character{101}}%
\TextAt{14}{220.773}{733.614}{o}%
\TextAt{3}{226.617}{738.55}{\Character{3}}%
\TextAt{27}{226.617}{730.659}{p}%
\TextAt{9}{231.349}{733.614}{\Character{47}}%
\TextAt{14}{237.202}{733.614}{o}%
\TextAt{3}{243.046}{738.55}{\Character{3}}%
\TextAt{27}{243.046}{730.659}{p}%
\NativeMathAt{251.1}{733.614}{11.955}{\cong}%
\TextAt{9}{263.719}{733.614}{\Character{40}}%
\TextAt{20}{267.872}{734.15}{\Character{101}}%
\TextAt{14}{268.271}{733.614}{o}%
\TextAt{9}{274.115}{733.614}{\Character{47}}%
\TextAt{14}{279.968}{733.614}{f}%
\TextAt{9}{283.869}{733.614}{\Character{41}}%
\TextAt{3}{288.421}{738.55}{\Character{3}}%
\TextAt{9}{293.154}{733.614}{\Character{47}\Character{40}}%
\TextAt{14}{303.559}{733.614}{o}%
\TextAt{9}{309.403}{733.614}{\Character{47}}%
\TextAt{14}{315.256}{733.614}{f}%
\TextAt{9}{319.157}{733.614}{\Character{41}}%
\TextAt{3}{323.709}{738.55}{\Character{3}}%
\TextAt{8}{328.441}{733.614}{.}%
\TextAt{19}{79.226}{689.268}{Proof.}%
\TextAt{8}{116}{689.268}{We}%
\TextAt{8}{136.145}{689.268}{apply}%
\TextAt{8}{168.34}{689.268}{the}%
\TextAt{8}{188.496}{689.268}{Chinese}%
\TextAt{8}{232.121}{689.268}{remainder}%
\TextAt{8}{287.402}{689.268}{theorem}%
\TextAt{8}{332.927}{689.268}{(12.3)}%
\TextAt{8}{366.76}{689.268}{repeatedly.}%
\TextAt{8}{427.56}{689.268}{We}%
\TextAt{8}{447.704}{689.268}{have}%
\TextAt{14}{263.366}{660.267}{o}%
\TextAt{9}{269.21}{660.267}{\Character{47}}%
\TextAt{14}{275.063}{660.267}{f}%
\NativeMathAt{282.285}{660.267}{11.955}{\cong}%
\TextAt{20}{294.904}{671.624}{\Character{77}}%
\TextAt{27}{301.942}{646.518}{p}%
\TextAt{14}{314.962}{660.267}{o}%
\TextAt{27}{320.806}{658.474}{p}%
\TextAt{9}{325.293}{660.267}{\Character{47}}%
\TextAt{14}{331.146}{660.267}{fo}%
\TextAt{27}{340.891}{658.474}{p}%
\TextAt{13}{345.378}{660.267}{\Character{58}}%
\TextAt{8}{517.802}{660.267}{(1)}%
\ErrataAt{79.19}{621.677}{455.81}{The only nonzero prime ideals of the integral closure $\widetilde{\mathfrak o}_{\mathfrak p}$ of $\mathfrak o_{\mathfrak p}$ are the finitely many above}%
\TextAt{14}{79.226}{604.649}{po}%
\TextAt{27}{91.053}{602.856}{p}%
\TextAt{8}{95.54}{604.649}{.}%
\TextAt{8}{105.619}{604.649}{They}%
\TextAt{8}{136.397}{604.649}{give}%
\TextAt{8}{161.006}{604.649}{the}%
\TextAt{8}{181.706}{604.649}{localizations}%
\TextAt{20}{248.585}{605.184}{\Character{101}}%
\TextAt{14}{248.984}{604.649}{o}%
\TextAt{30}{254.47}{603.011}{\Character{101}}%
\TextAt{27}{254.828}{602.266}{p}%
\TextAt{8}{259.315}{604.649}{,}%
\TextAt{8}{267.143}{604.649}{where}%
\TextAt{20}{301.167}{605.766}{\Character{101}}%
\TextAt{14}{301.496}{604.649}{p}%
\TextAt{8}{311.92}{604.649}{varies}%
\TextAt{8}{345.376}{604.649}{over}%
\TextAt{8}{370.966}{604.649}{the}%
\TextAt{8}{391.665}{604.649}{prime}%
\TextAt{8}{425.373}{604.649}{ideals}%
\TextAt{8}{458.494}{604.649}{above}%
\TextAt{14}{492.213}{604.649}{p}%
\TextAt{8}{502.637}{604.649}{of}%
\TextAt{8}{516.51}{604.649}{the}%
\TextAt{8}{79.226}{587.603}{ring}%
\TextAt{20}{103.794}{588.138}{\Character{101}}%
\TextAt{14}{104.193}{587.603}{o}%
\TextAt{8}{110.037}{587.603}{.}%
\TextAt{8}{121.189}{587.603}{At}%
\TextAt{8}{138.994}{587.603}{the}%
\TextAt{8}{160.052}{587.603}{same}%
\TextAt{8}{190.279}{587.603}{time,}%
\TextAt{20}{220.918}{588.138}{\Character{101}}%
\TextAt{14}{221.317}{587.603}{o}%
\TextAt{27}{227.161}{585.81}{p}%
\TextAt{8}{236.447}{587.603}{is}%
\TextAt{8}{249.112}{587.603}{the}%
\TextAt{8}{270.169}{587.603}{localization}%
\TextAt{8}{333.189}{587.603}{of}%
\TextAt{20}{347.021}{588.138}{\Character{101}}%
\TextAt{14}{347.42}{587.603}{o}%
\TextAt{8}{358.063}{587.603}{with}%
\TextAt{8}{385.624}{587.603}{respect}%
\TextAt{8}{426.575}{587.603}{to}%
\TextAt{8}{441.786}{587.603}{the}%
\TextAt{8}{462.844}{587.603}{multiplicative}%
\TextAt{8}{79.226}{570.557}{subset}%
\TextAt{14}{115.873}{570.557}{o}%
\TextAt{16}{124.89}{570.557}{\Character{110}}%
\TextAt{14}{134.041}{570.557}{p}%
\TextAt{8}{140.024}{570.557}{.}%
\TextAt{8}{150.744}{570.557}{Since}%
\TextAt{14}{182.059}{570.557}{f}%
\TextAt{8}{190.615}{570.557}{is}%
\TextAt{8}{203.136}{570.557}{an}%
\TextAt{8}{220.152}{570.557}{ideal}%
\TextAt{8}{248.873}{570.557}{of}%
\TextAt{20}{262.562}{571.092}{\Character{101}}%
\TextAt{14}{262.96}{570.557}{o}%
\TextAt{8}{268.805}{570.557}{,}%
\TextAt{8}{276.901}{570.557}{it}%
\TextAt{8}{289.362}{570.557}{follows}%
\TextAt{8}{328.555}{570.557}{that}%
\TextAt{14}{354.682}{570.557}{f}%
\TextAt{20}{358.183}{571.092}{\Character{101}}%
\TextAt{14}{358.582}{570.557}{o}%
\TextAt{27}{364.426}{568.764}{p}%
\TextAt{9}{373.525}{570.557}{\Character{61}}%
\TextAt{14}{387.241}{570.557}{fo}%
\TextAt{27}{396.986}{568.764}{p}%
\TextAt{8}{401.473}{570.557}{.}%
\TextAt{8}{412.193}{570.557}{The}%
\TextAt{8}{437.004}{570.557}{Chinese}%
\TextAt{8}{481.386}{570.557}{remainder}%
\TextAt{8}{79.226}{553.511}{theorem}%
\TextAt{8}{124.751}{553.511}{yields}%
\TextAt{20}{257.184}{537}{\Character{101}}%
\TextAt{14}{257.583}{536.465}{o}%
\TextAt{27}{263.427}{534.672}{p}%
\TextAt{9}{267.914}{536.465}{\Character{47}}%
\TextAt{14}{273.767}{536.465}{f}%
\TextAt{20}{277.269}{537}{\Character{101}}%
\TextAt{14}{277.668}{536.465}{o}%
\TextAt{27}{283.512}{534.672}{p}%
\NativeMathAt{291.319}{536.465}{11.955}{\cong}%
\TextAt{20}{303.939}{547.822}{\Character{77}}%
\TextAt{30}{305.332}{522.141}{\Character{101}}%
\TextAt{27}{305.69}{521.396}{p}%
\NativeMathAt{309.678}{521.396}{7.97}{\supseteq}%
\TextAt{27}{316.265}{521.396}{p}%
\TextAt{20}{323.598}{537}{\Character{101}}%
\TextAt{14}{323.997}{536.465}{o}%
\TextAt{30}{329.483}{534.827}{\Character{101}}%
\TextAt{27}{329.841}{534.082}{p}%
\TextAt{9}{334.328}{536.465}{\Character{47}}%
\TextAt{14}{340.181}{536.465}{f}%
\TextAt{20}{343.683}{537}{\Character{101}}%
\TextAt{14}{344.081}{536.465}{o}%
\TextAt{30}{349.568}{534.827}{\Character{101}}%
\TextAt{27}{349.926}{534.082}{p}%
\TextAt{8}{79.226}{502.925}{and}%
\TextAt{20}{219.731}{486.415}{\Character{101}}%
\TextAt{14}{220.13}{485.879}{o}%
\TextAt{9}{225.974}{485.879}{\Character{47}}%
\TextAt{14}{231.827}{485.879}{f}%
\NativeMathAt{239.049}{485.879}{11.955}{\cong}%
\TextAt{20}{251.668}{497.237}{\Character{77}}%
\TextAt{27}{258.706}{472.131}{p}%
\TextAt{20}{271.726}{497.237}{\Character{77}}%
\TextAt{30}{273.119}{471.555}{\Character{101}}%
\TextAt{27}{273.477}{470.811}{p}%
\NativeMathAt{277.466}{470.811}{7.97}{\supseteq}%
\TextAt{27}{284.052}{470.811}{p}%
\TextAt{20}{291.385}{486.415}{\Character{101}}%
\TextAt{14}{291.784}{485.879}{o}%
\TextAt{30}{297.271}{484.241}{\Character{101}}%
\TextAt{27}{297.628}{483.497}{p}%
\TextAt{9}{302.115}{485.879}{\Character{47}}%
\TextAt{14}{307.968}{485.879}{f}%
\TextAt{20}{311.47}{486.415}{\Character{101}}%
\TextAt{14}{311.869}{485.879}{o}%
\TextAt{30}{317.355}{484.241}{\Character{101}}%
\TextAt{27}{317.713}{483.497}{p}%
\NativeMathAt{325.521}{485.879}{11.955}{\cong}%
\TextAt{20}{338.14}{497.237}{\Character{77}}%
\TextAt{27}{345.178}{472.131}{p}%
\TextAt{20}{357.799}{486.415}{\Character{101}}%
\TextAt{14}{358.198}{485.879}{o}%
\TextAt{27}{364.042}{484.086}{p}%
\TextAt{9}{368.529}{485.879}{\Character{47}}%
\TextAt{14}{374.382}{485.879}{f}%
\TextAt{20}{377.884}{486.415}{\Character{101}}%
\TextAt{14}{378.283}{485.879}{o}%
\TextAt{27}{384.127}{484.086}{p}%
\TextAt{13}{388.614}{485.879}{\Character{58}}%
\TextAt{8}{517.802}{485.879}{(2)}%
\TextAt{8}{79.226}{451.67}{Passing}%
\TextAt{8}{121.464}{451.67}{to}%
\TextAt{8}{135.774}{451.67}{unit}%
\TextAt{8}{160.485}{451.67}{groups,}%
\TextAt{8}{201.528}{451.67}{we}%
\TextAt{8}{218.755}{451.67}{get}%
\TextAt{8}{238.266}{451.67}{from}%
\TextAt{8}{265.906}{451.67}{(1)}%
\TextAt{8}{284.771}{451.67}{and}%
\TextAt{8}{307.534}{451.67}{(2)}%
\TextAt{8}{326.399}{451.67}{that}%
\TextAt{9}{207.829}{422.669}{\Character{40}}%
\TextAt{20}{211.982}{423.204}{\Character{101}}%
\TextAt{14}{212.381}{422.669}{o}%
\TextAt{9}{218.225}{422.669}{\Character{47}}%
\TextAt{14}{224.078}{422.669}{f}%
\TextAt{9}{227.979}{422.669}{\Character{41}}%
\TextAt{3}{232.531}{427.605}{\Character{3}}%
\TextAt{9}{237.263}{422.669}{\Character{47}\Character{40}}%
\TextAt{14}{247.669}{422.669}{o}%
\TextAt{9}{253.513}{422.669}{\Character{47}}%
\TextAt{14}{259.366}{422.669}{f}%
\TextAt{9}{263.267}{422.669}{\Character{41}}%
\TextAt{3}{267.819}{427.605}{\Character{3}}%
\NativeMathAt{275.872}{422.669}{11.955}{\cong}%
\TextAt{20}{288.491}{434.026}{\Character{77}}%
\TextAt{27}{295.53}{408.92}{p}%
\TextAt{9}{306.557}{422.669}{\Character{40}}%
\TextAt{20}{310.71}{423.204}{\Character{101}}%
\TextAt{14}{311.109}{422.669}{o}%
\TextAt{27}{316.953}{420.876}{p}%
\TextAt{9}{321.44}{422.669}{\Character{47}}%
\TextAt{14}{327.293}{422.669}{f}%
\TextAt{20}{330.795}{423.204}{\Character{101}}%
\TextAt{14}{331.194}{422.669}{o}%
\TextAt{27}{337.038}{420.876}{p}%
\TextAt{9}{341.525}{422.669}{\Character{41}}%
\TextAt{3}{346.077}{427.605}{\Character{3}}%
\TextAt{9}{350.809}{422.669}{\Character{47}\Character{40}}%
\TextAt{14}{361.215}{422.669}{o}%
\TextAt{27}{367.059}{420.876}{p}%
\TextAt{9}{371.546}{422.669}{\Character{47}}%
\TextAt{14}{377.399}{422.669}{fo}%
\TextAt{27}{387.144}{420.876}{p}%
\TextAt{9}{391.63}{422.669}{\Character{41}}%
\TextAt{3}{396.183}{427.605}{\Character{3}}%
\TextAt{13}{400.915}{422.669}{\Character{58}}%
\TextAt{8}{517.802}{422.669}{(3)}%
\TextAt{8}{79.226}{384.97}{For}%
\TextAt{14}{100.196}{384.97}{f}%
\NativeMathAt{107.417}{384.97}{11.955}{\subseteq}%
\TextAt{14}{120.036}{384.97}{p}%
\TextAt{8}{129.917}{384.97}{we}%
\TextAt{8}{147.144}{384.97}{now}%
\TextAt{8}{171.533}{384.97}{consider}%
\TextAt{8}{217.118}{384.97}{the}%
\TextAt{8}{237.274}{384.97}{homomorphism}%
\TextAt{13}{230.172}{355.969}{\Character{39}}%
\TextAt{9}{241.162}{355.969}{\Character{58}}%
\TextAt{20}{247.335}{356.504}{\Character{101}}%
\TextAt{14}{247.734}{355.969}{o}%
\TextAt{3}{253.578}{360.905}{\Character{3}}%
\TextAt{27}{253.578}{353.013}{p}%
\NativeArrowAt{261.631}{355.969}{11.955}{19.26}{\longrightarrow}%
\TextAt{9}{284.213}{355.969}{\Character{40}}%
\TextAt{20}{288.367}{356.504}{\Character{101}}%
\TextAt{14}{288.766}{355.969}{o}%
\TextAt{27}{294.61}{354.175}{p}%
\TextAt{9}{299.097}{355.969}{\Character{47}}%
\TextAt{14}{304.95}{355.969}{f}%
\TextAt{20}{308.452}{356.504}{\Character{101}}%
\TextAt{14}{308.85}{355.969}{o}%
\TextAt{27}{314.695}{354.175}{p}%
\TextAt{9}{319.181}{355.969}{\Character{41}}%
\TextAt{3}{323.734}{360.905}{\Character{3}}%
\TextAt{9}{328.466}{355.969}{\Character{47}\Character{40}}%
\TextAt{14}{338.871}{355.969}{o}%
\TextAt{27}{344.716}{354.175}{p}%
\TextAt{9}{349.202}{355.969}{\Character{47}}%
\TextAt{14}{355.055}{355.969}{fo}%
\TextAt{27}{364.8}{354.175}{p}%
\TextAt{9}{369.287}{355.969}{\Character{41}}%
\TextAt{3}{373.839}{360.905}{\Character{3}}%
\TextAt{13}{378.572}{355.969}{\Character{58}}%
\TextAt{8}{79.226}{326.967}{It}%
\TextAt{8}{92.771}{326.967}{is}%
\TextAt{8}{105.406}{326.967}{surjective.}%
\TextAt{8}{164.985}{326.967}{In}%
\TextAt{8}{180.478}{326.967}{fact,}%
\TextAt{8}{207.905}{326.967}{if}%
\TextAt{13}{219.501}{326.967}{\Character{34}}%
\TextAt{8}{228.3}{326.967}{mod}%
\TextAt{14}{253.738}{326.967}{f}%
\TextAt{20}{257.24}{327.503}{\Character{101}}%
\TextAt{14}{257.638}{326.967}{o}%
\TextAt{27}{263.483}{325.174}{p}%
\TextAt{8}{272.739}{326.967}{is}%
\TextAt{8}{285.375}{326.967}{a}%
\TextAt{8}{296.002}{326.967}{unit}%
\TextAt{8}{321.585}{326.967}{in}%
\TextAt{20}{335.711}{327.503}{\Character{101}}%
\TextAt{14}{336.11}{326.967}{o}%
\TextAt{27}{341.954}{325.174}{p}%
\TextAt{9}{346.441}{326.967}{\Character{47}}%
\TextAt{14}{352.294}{326.967}{f}%
\TextAt{20}{355.796}{327.503}{\Character{101}}%
\TextAt{14}{356.195}{326.967}{o}%
\TextAt{27}{362.039}{325.174}{p}%
\TextAt{8}{366.526}{326.967}{,}%
\TextAt{8}{374.765}{326.967}{then}%
\TextAt{13}{402.297}{326.967}{\Character{34}}%
\TextAt{8}{412.544}{326.967}{is}%
\TextAt{8}{425.18}{326.967}{a}%
\TextAt{8}{435.808}{326.967}{unit}%
\TextAt{8}{461.391}{326.967}{in}%
\TextAt{20}{475.517}{327.503}{\Character{101}}%
\TextAt{14}{475.916}{326.967}{o}%
\TextAt{27}{481.76}{325.174}{p}%
\TextAt{8}{486.247}{326.967}{.}%
\TextAt{8}{497.311}{326.967}{This}%
\TextAt{8}{524.903}{326.967}{is}%
\TextAt{8}{79.226}{309.922}{so}%
\TextAt{8}{94.412}{309.922}{because}%
\TextAt{8}{138.53}{309.922}{the}%
\TextAt{8}{159.503}{309.922}{units}%
\TextAt{8}{189.645}{309.922}{in}%
\TextAt{8}{204.114}{309.922}{any}%
\TextAt{8}{227.047}{309.922}{ring}%
\TextAt{8}{251.929}{309.922}{are}%
\TextAt{8}{272.256}{309.922}{precisely}%
\TextAt{8}{320.929}{309.922}{those}%
\TextAt{8}{352.374}{309.922}{elements}%
\TextAt{8}{401.047}{309.922}{that}%
\TextAt{8}{427.232}{309.922}{are}%
\TextAt{8}{447.559}{309.922}{not}%
\TextAt{8}{469.189}{309.922}{contained}%
\TextAt{8}{523.014}{309.922}{in}%
\TextAt{8}{79.226}{292.876}{any}%
\TextAt{8}{101.974}{292.876}{maximal}%
\TextAt{8}{150.413}{292.876}{ideal,}%
\TextAt{8}{182.417}{292.876}{and}%
\TextAt{8}{205.81}{292.876}{the}%
\TextAt{8}{226.598}{292.876}{preimages}%
\TextAt{8}{281.923}{292.876}{of}%
\TextAt{8}{295.884}{292.876}{the}%
\TextAt{8}{316.671}{292.876}{maximal}%
\TextAt{8}{365.111}{292.876}{ideals}%
\TextAt{8}{398.319}{292.876}{of}%
\TextAt{20}{411.881}{293.411}{\Character{101}}%
\TextAt{14}{412.28}{292.876}{o}%
\TextAt{27}{418.124}{291.082}{p}%
\TextAt{9}{422.611}{292.876}{\Character{47}}%
\TextAt{14}{428.464}{292.876}{f}%
\TextAt{20}{431.966}{293.411}{\Character{101}}%
\TextAt{14}{432.365}{292.876}{o}%
\TextAt{27}{438.209}{291.082}{p}%
\TextAt{8}{447.224}{292.876}{give}%
\TextAt{8}{471.921}{292.876}{precisely}%
\TextAt{8}{520.408}{292.876}{all}%
\TextAt{8}{79.226}{275.83}{the}%
\TextAt{8}{99.987}{275.83}{maximal}%
\TextAt{8}{148.4}{275.83}{ideals}%
\TextAt{8}{181.583}{275.83}{of}%
\TextAt{20}{195.118}{276.365}{\Character{101}}%
\TextAt{14}{195.517}{275.83}{o}%
\TextAt{27}{201.361}{274.036}{p}%
\TextAt{8}{205.848}{275.83}{,}%
\TextAt{8}{213.753}{275.83}{since}%
\TextAt{14}{243.026}{275.83}{f}%
\TextAt{20}{246.528}{276.365}{\Character{101}}%
\TextAt{14}{246.927}{275.83}{o}%
\TextAt{27}{252.771}{274.036}{p}%
\NativeMathAt{261.609}{275.83}{11.955}{\subseteq}%
\TextAt{14}{275.258}{275.83}{p}%
\TextAt{20}{280.843}{276.365}{\Character{101}}%
\TextAt{14}{281.241}{275.83}{o}%
\TextAt{27}{287.086}{274.036}{p}%
\TextAt{8}{291.572}{275.83}{.}%
\TextAt{8}{301.834}{275.83}{The}%
\TextAt{8}{326.493}{275.83}{kernel}%
\TextAt{8}{361.564}{275.83}{of}%
\TextAt{13}{375.498}{275.83}{\Character{39}}%
\TextAt{8}{387.669}{275.83}{is}%
\TextAt{8}{400.037}{275.83}{a}%
\TextAt{8}{410.397}{275.83}{subgroup}%
\TextAt{8}{461.799}{275.83}{of}%
\TextAt{20}{475.335}{276.365}{\Character{101}}%
\TextAt{14}{475.734}{275.83}{o}%
\TextAt{3}{481.578}{280.168}{\Character{3}}%
\TextAt{27}{481.578}{272.874}{p}%
\TextAt{8}{490.812}{275.83}{which}%
\TextAt{8}{524.903}{275.83}{is}%
\TextAt{8}{79.226}{258.784}{contained}%
\TextAt{8}{131.981}{258.784}{in}%
\TextAt{14}{145.379}{258.784}{o}%
\TextAt{27}{151.224}{256.99}{p}%
\TextAt{8}{155.71}{258.784}{,}%
\TextAt{8}{162.656}{258.784}{and}%
\TextAt{8}{185.164}{258.784}{which}%
\TextAt{8}{218.396}{258.784}{contains}%
\TextAt{14}{264.062}{258.784}{o}%
\TextAt{3}{269.906}{263.122}{\Character{3}}%
\TextAt{27}{269.906}{255.828}{p}%
\TextAt{8}{274.638}{258.784}{.}%
\TextAt{8}{283.002}{258.784}{It}%
\TextAt{8}{295.42}{258.784}{is}%
\TextAt{8}{306.929}{258.784}{therefore}%
\TextAt{8}{355.775}{258.784}{equal}%
\TextAt{8}{386.413}{258.784}{to}%
\TextAt{14}{400.469}{258.784}{o}%
\TextAt{3}{406.313}{263.122}{\Character{3}}%
\TextAt{27}{406.313}{255.828}{p}%
\TextAt{8}{411.045}{258.784}{.}%
\TextAt{8}{419.409}{258.784}{We}%
\TextAt{8}{439.299}{258.784}{now}%
\TextAt{8}{463.433}{258.784}{conclude}%
\TextAt{8}{511.298}{258.784}{that}%
\TextAt{20}{235.321}{230.318}{\Character{101}}%
\TextAt{14}{235.72}{229.782}{o}%
\TextAt{3}{241.564}{234.719}{\Character{3}}%
\TextAt{27}{241.564}{226.827}{p}%
\TextAt{9}{246.296}{229.782}{\Character{47}}%
\TextAt{14}{252.149}{229.782}{o}%
\TextAt{3}{257.994}{234.719}{\Character{3}}%
\TextAt{27}{257.994}{226.827}{p}%
\NativeMathAt{266.047}{229.782}{11.955}{\cong}%
\TextAt{9}{278.666}{229.782}{\Character{40}}%
\TextAt{20}{282.82}{230.318}{\Character{101}}%
\TextAt{14}{283.218}{229.782}{o}%
\TextAt{27}{289.063}{227.989}{p}%
\TextAt{9}{293.549}{229.782}{\Character{47}}%
\TextAt{14}{299.402}{229.782}{f}%
\TextAt{20}{302.904}{230.318}{\Character{101}}%
\TextAt{14}{303.303}{229.782}{o}%
\TextAt{27}{309.147}{227.989}{p}%
\TextAt{9}{313.634}{229.782}{\Character{41}}%
\TextAt{3}{318.186}{234.719}{\Character{3}}%
\TextAt{9}{322.919}{229.782}{\Character{47}\Character{40}}%
\TextAt{14}{333.324}{229.782}{o}%
\TextAt{27}{339.168}{227.989}{p}%
\TextAt{9}{343.655}{229.782}{\Character{47}}%
\TextAt{14}{349.508}{229.782}{fo}%
\TextAt{27}{359.253}{227.989}{p}%
\TextAt{9}{363.74}{229.782}{\Character{41}}%
\TextAt{3}{368.292}{234.719}{\Character{3}}%
\TextAt{13}{373.024}{229.782}{\Character{58}}%
\TextAt{8}{79.226}{200.781}{This}%
\TextAt{8}{107.257}{200.781}{remains}%
\TextAt{8}{152.205}{200.781}{true}%
\TextAt{8}{178.228}{200.781}{also}%
\TextAt{8}{203.02}{200.781}{for}%
\TextAt{14}{222.216}{200.781}{p}%
\TextAt{18}{233.755}{200.781}{\Character{43}}%
\TextAt{14}{248.609}{200.781}{f}%
\TextAt{8}{257.718}{200.781}{because}%
\TextAt{8}{302.332}{200.781}{then}%
\TextAt{8}{330.303}{200.781}{both}%
\TextAt{8}{359.255}{200.781}{sides}%
\TextAt{8}{388.649}{200.781}{are}%
\TextAt{8}{409.471}{200.781}{equal}%
\TextAt{8}{441.675}{200.781}{to}%
\TextAt{9}{457.297}{200.781}{1}%
\TextAt{8}{468.359}{200.781}{according}%
\TextAt{8}{522.357}{200.781}{to}%
\TextAt{8}{79.226}{183.735}{(12.10).}%
\TextAt{8}{123.468}{183.735}{The}%
\TextAt{8}{147.522}{183.735}{claim}%
\TextAt{8}{178.737}{183.735}{of}%
\TextAt{8}{192.067}{183.735}{the}%
\TextAt{8}{212.223}{183.735}{proposition}%
\TextAt{8}{273.757}{183.735}{now}%
\TextAt{8}{298.145}{183.735}{follows}%
\TextAt{8}{336.581}{183.735}{from}%
\TextAt{8}{364.221}{183.735}{(3).}%
\Vector{q 0 G 0 g 0.399 w 452.334 -590.265 m 452.334 -582.195 l S Q}%
\Vector{q 0 G 0 g 0.399 w 452.534 -582.394 m 459.707 -582.394 l S Q}%
\Vector{q 0 G 0 g 0.399 w 452.534 -590.065 m 459.707 -590.065 l S Q}%
\Vector{q 0 G 0 g 0.399 w 459.906 -590.265 m 459.906 -582.195 l S Q}%
\TextAt{8}{97.159}{154.734}{Our}%
\TextAt{8}{121.011}{154.734}{study}%
\TextAt{8}{153.052}{154.734}{of}%
\TextAt{8}{166.167}{154.734}{one-dimensional}%
\TextAt{8}{251.863}{154.734}{noetherian}%
\TextAt{8}{309.859}{154.734}{integral}%
\TextAt{8}{352.253}{154.734}{domains}%
\TextAt{8}{398.281}{154.734}{was}%
\TextAt{8}{420.567}{154.734}{motivated}%
\TextAt{8}{475.323}{154.734}{by}%
\TextAt{8}{491.368}{154.734}{the}%
\TextAt{8}{511.31}{154.734}{con-}%
\TextAt{8}{79.226}{137.688}{sideration}%
\TextAt{8}{133.931}{137.688}{of}%
\TextAt{8}{147.916}{137.688}{orders.}%
\TextAt{8}{189.618}{137.688}{For}%
\TextAt{8}{211.243}{137.688}{them,}%
\TextAt{8}{245.226}{137.688}{(12.9)}%
\TextAt{8}{279.715}{137.688}{and}%
\TextAt{8}{303.133}{137.688}{(12.11)}%
\TextAt{8}{343.479}{137.688}{imply}%
\TextAt{8}{376.976}{137.688}{the}%
\TextAt{8}{397.788}{137.688}{following}%
\TextAt{8}{447.878}{137.688}{generalization}%
\TextAt{8}{523.337}{137.688}{of}%
\TextAt{8}{79.226}{120.642}{Dirichlet’s}%
\TextAt{8}{135.38}{120.642}{unit}%
\TextAt{8}{160.091}{120.642}{theorem}%
\TextAt{8}{205.616}{120.642}{and}%
\TextAt{8}{228.379}{120.642}{of}%
\TextAt{8}{241.709}{120.642}{the}%
\TextAt{8}{261.865}{120.642}{theorem}%
\TextAt{8}{307.391}{120.642}{on}%
\TextAt{8}{323.65}{120.642}{the}%
\TextAt{8}{343.806}{120.642}{finiteness}%
\TextAt{8}{394.651}{120.642}{of}%
\TextAt{8}{407.981}{120.642}{the}%
\TextAt{8}{428.138}{120.642}{class}%
\TextAt{8}{455.575}{120.642}{group.}%
\TextAt{12}{79.226}{94.63}{Theorem}%
\TextAt{12}{136.764}{94.63}{(12.12).}%
\TextAt{8}{185.996}{94.63}{Let}%
\TextAt{14}{207.421}{94.63}{o}%
\TextAt{8}{217.618}{94.63}{be}%
\TextAt{8}{233.997}{94.63}{an}%
\TextAt{8}{250.711}{94.63}{order}%
\TextAt{8}{281.736}{94.63}{in}%
\TextAt{8}{295.844}{94.63}{an}%
\TextAt{8}{312.558}{94.63}{algebraic}%
\TextAt{8}{362.448}{94.63}{number}%
\TextAt{8}{405.499}{94.63}{field}%
\TextAt{13}{431.311}{94.63}{K}%
\TextAt{8}{442.129}{94.63}{,}%
\TextAt{14}{449.847}{94.63}{o}%
\TextAt{21}{455.691}{92.837}{K}%
\TextAt{8}{468.247}{94.63}{the}%
\TextAt{8}{488.858}{94.63}{maximal}%
\TextAt{8}{79.226}{77.584}{order,}%
\TextAt{8}{113.047}{77.584}{and}%
\TextAt{14}{135.81}{77.584}{f}%
\TextAt{8}{143.608}{77.584}{the}%
\TextAt{8}{163.764}{77.584}{conductor}%
\TextAt{8}{218.399}{77.584}{of}%
\TextAt{14}{231.729}{77.584}{o}%
\TextAt{8}{237.574}{77.584}{.}%
\TextAt{8}{300.14}{43.033}{75}%
\Destination{522.36}{502.925}{nav.AMS.10}%
\Destination{522.36}{439.715}{nav.AMS.11}%
\Destination{522.36}{677.315}{nav.AMS.9}%
\Destination{78.23}{809.815}{nav.page.75}%
\Destination{79.36}{341.955}{nav.print.81}%
\Destination{79.36}{747.655}{nav.number.1.12.11}%
\Destination{79.36}{108.655}{nav.number.1.12.12}%
\Hotlink{470.713}{758.955}{30.936}{11.047}{nav.number.1.12.9}%
\Hotlink{332.427}{687.055}{30.936}{11.047}{nav.number.1.12.3}%
\Hotlink{78.726}{181.555}{36.794}{11.047}{nav.number.1.12.10}%
\Hotlink{244.726}{135.455}{30.936}{11.047}{nav.number.1.12.9}%
\Hotlink{302.633}{135.455}{36.794}{11.047}{nav.number.1.12.11}%
\Destination{79.226}{585}{errata.59}%
\Destination{79.226}{633}{errata.181}%
}

% Page 88 (body page 76)
\Page{%
\TextAt{2}{79.226}{796.943}{I.}%
\TextAt{2}{89.83}{796.943}{Algebraic}%
\TextAt{2}{138.648}{796.943}{Integers}%
\Vector{q 0 G 0 g 0.299 w 7.226 18.729 m 460.769 18.729 l S Q}%
\TextAt{8}{97.159}{761.124}{Then}%
\TextAt{8}{127.716}{761.124}{the}%
\TextAt{8}{147.873}{761.124}{groups}%
\TextAt{14}{185.663}{761.124}{o}%
\TextAt{3}{191.507}{765.462}{\Character{3}}%
\TextAt{21}{191.507}{758.103}{K}%
\TextAt{9}{199.71}{761.124}{\Character{47}}%
\TextAt{14}{205.563}{761.124}{o}%
\TextAt{3}{211.407}{765.462}{\Character{3}}%
\TextAt{8}{220.037}{761.124}{and}%
\TextAt{13}{242.799}{761.124}{P}%
\TextAt{13}{251.969}{761.124}{ic}%
\TextAt{9}{261.002}{761.124}{\Character{40}}%
\TextAt{14}{265.554}{761.124}{o}%
\TextAt{9}{271.399}{761.124}{\Character{41}}%
\TextAt{8}{279.848}{761.124}{are}%
\TextAt{8}{299.359}{761.124}{finite}%
\TextAt{8}{329.271}{761.124}{and}%
\TextAt{8}{352.034}{761.124}{one}%
\TextAt{8}{373.493}{761.124}{has}%
\TextAt{8}{227.147}{729.635}{\char35\relax }%
\TextAt{13}{236.903}{729.635}{P}%
\TextAt{13}{246.061}{729.635}{ic}%
\TextAt{9}{255.105}{729.635}{\Character{40}}%
\TextAt{14}{259.658}{729.635}{o}%
\TextAt{9}{265.502}{729.635}{\Character{41}}%
\TextAt{9}{273.357}{729.635}{\Character{61}}%
\TextAt{13}{301.336}{737.722}{h}%
\TextAt{21}{308.075}{735.929}{K}%
\Vector{q 0 G 0 g 0.478 w 214.996 -41.377 m 258.617 -41.377 l S Q}%
\TextAt{9}{286.996}{721.434}{\Character{40}}%
\TextAt{14}{291.548}{721.434}{o}%
\TextAt{3}{297.393}{725.552}{\Character{3}}%
\TextAt{21}{297.393}{718.193}{K}%
\TextAt{9}{308.916}{721.434}{\Character{58}}%
\TextAt{14}{315.489}{721.434}{o}%
\TextAt{3}{321.333}{724.888}{\Character{3}}%
\TextAt{9}{326.065}{721.434}{\Character{41}}%
\TextAt{8}{333.008}{737.722}{\char35\relax }%
\TextAt{9}{342.764}{737.722}{\Character{40}}%
\TextAt{14}{347.316}{737.722}{o}%
\TextAt{21}{353.16}{735.929}{K}%
\TextAt{9}{361.363}{737.722}{\Character{47}}%
\TextAt{14}{367.216}{737.722}{f}%
\TextAt{9}{371.116}{737.722}{\Character{41}}%
\TextAt{3}{375.669}{742.061}{\Character{3}}%
\Vector{q 0 G 0 g 0.478 w 261.008 -41.377 m 308.401 -41.377 l S Q}%
\TextAt{8}{337.11}{721.434}{\char35\relax }%
\TextAt{9}{346.865}{721.434}{\Character{40}}%
\TextAt{14}{351.417}{721.434}{o}%
\TextAt{9}{357.262}{721.434}{\Character{47}}%
\TextAt{14}{363.115}{721.434}{f}%
\TextAt{9}{367.015}{721.434}{\Character{41}}%
\TextAt{3}{371.568}{724.888}{\Character{3}}%
\TextAt{13}{381.597}{729.635}{\Character{59}}%
\TextAt{8}{79.226}{698.362}{where}%
\TextAt{13}{113.035}{698.362}{h}%
\TextAt{21}{119.774}{696.569}{K}%
\TextAt{8}{131.874}{698.362}{is}%
\TextAt{8}{143.638}{698.362}{the}%
\TextAt{8}{163.794}{698.362}{class}%
\TextAt{8}{191.231}{698.362}{number}%
\TextAt{8}{233.828}{698.362}{of}%
\TextAt{13}{247.158}{698.362}{K}%
\TextAt{8}{257.975}{698.362}{.}%
\TextAt{8}{266.424}{698.362}{In}%
\TextAt{8}{281.045}{698.362}{particular,}%
\TextAt{8}{338.286}{698.362}{one}%
\TextAt{8}{359.746}{698.362}{has}%
\TextAt{8}{380.62}{698.362}{that}%
\TextAt{8}{221.113}{670.898}{rank}%
\TextAt{9}{244.211}{670.898}{\Character{40}}%
\TextAt{14}{248.763}{670.898}{o}%
\TextAt{3}{254.607}{675.834}{\Character{3}}%
\TextAt{9}{259.34}{670.898}{\Character{41}}%
\TextAt{9}{267.195}{670.898}{\Character{61}}%
\TextAt{8}{279.638}{670.898}{rank}%
\TextAt{9}{302.736}{670.898}{\Character{40}}%
\TextAt{14}{307.288}{670.898}{o}%
\TextAt{3}{313.132}{675.834}{\Character{3}}%
\TextAt{21}{313.132}{667.943}{K}%
\TextAt{9}{321.335}{670.898}{\Character{41}}%
\TextAt{9}{329.19}{670.898}{\Character{61}}%
\TextAt{13}{341.633}{670.898}{r}%
\TextAt{9}{349.891}{670.898}{\Character{43}}%
\TextAt{13}{361.652}{670.898}{s}%
\TextAt{16}{369.823}{670.898}{\Character{0}}%
\TextAt{9}{381.778}{670.898}{1}%
\TextAt{13}{387.631}{670.898}{\Character{58}}%
\TextAt{19}{79.226}{626.38}{Proof.}%
\TextAt{8}{116}{626.38}{By}%
\TextAt{8}{134.363}{626.38}{(12.9)}%
\TextAt{8}{168.197}{626.38}{and}%
\TextAt{8}{190.959}{626.38}{(12.11),}%
\TextAt{8}{233.902}{626.38}{and}%
\TextAt{8}{256.665}{626.38}{since}%
\TextAt{13}{285.333}{626.38}{P}%
\TextAt{13}{294.503}{626.38}{ic}%
\TextAt{9}{303.536}{626.38}{\Character{40}}%
\TextAt{14}{308.088}{626.38}{o}%
\TextAt{21}{313.932}{624.586}{K}%
\TextAt{9}{322.135}{626.38}{\Character{41}}%
\TextAt{9}{329.99}{626.38}{\Character{61}}%
\TextAt{13}{342.434}{626.38}{C}%
\TextAt{13}{351.651}{626.38}{l}%
\TextAt{21}{355.154}{624.586}{K}%
\TextAt{8}{363.357}{626.38}{,}%
\TextAt{8}{370.506}{626.38}{we}%
\TextAt{8}{387.733}{626.38}{have}%
\TextAt{8}{414.728}{626.38}{the}%
\TextAt{8}{434.885}{626.38}{exact}%
\TextAt{8}{465.777}{626.38}{sequence}%
\TextAt{9}{155.024}{598.916}{1}%
\NativeArrowAt{164.198}{598.916}{11.955}{19.26}{\longrightarrow}%
\TextAt{14}{186.78}{598.916}{o}%
\TextAt{3}{192.624}{603.852}{\Character{3}}%
\TextAt{21}{192.624}{595.96}{K}%
\TextAt{9}{200.827}{598.916}{\Character{47}}%
\TextAt{14}{206.68}{598.916}{o}%
\TextAt{3}{212.524}{603.852}{\Character{3}}%
\NativeArrowAt{220.577}{598.916}{11.955}{19.26}{\longrightarrow}%
\TextAt{9}{243.159}{598.916}{\Character{40}}%
\TextAt{14}{247.711}{598.916}{o}%
\TextAt{21}{253.556}{597.122}{K}%
\TextAt{9}{261.758}{598.916}{\Character{47}}%
\TextAt{14}{267.611}{598.916}{f}%
\TextAt{9}{271.512}{598.916}{\Character{41}}%
\TextAt{3}{276.064}{603.852}{\Character{3}}%
\TextAt{9}{280.796}{598.916}{\Character{47}\Character{40}}%
\TextAt{14}{291.202}{598.916}{o}%
\TextAt{9}{297.046}{598.916}{\Character{47}}%
\TextAt{14}{302.899}{598.916}{f}%
\TextAt{9}{306.799}{598.916}{\Character{41}}%
\TextAt{3}{311.352}{603.852}{\Character{3}}%
\NativeArrowAt{319.405}{598.916}{11.955}{19.26}{\longrightarrow}%
\TextAt{13}{341.987}{598.916}{P}%
\TextAt{13}{351.145}{598.916}{ic}%
\TextAt{9}{360.189}{598.916}{\Character{40}}%
\TextAt{14}{364.742}{598.916}{o}%
\TextAt{9}{370.586}{598.916}{\Character{41}}%
\NativeArrowAt{378.459}{598.916}{11.955}{19.26}{\longrightarrow}%
\TextAt{13}{401.041}{598.916}{C}%
\TextAt{13}{410.259}{598.916}{l}%
\TextAt{21}{413.762}{597.122}{K}%
\NativeArrowAt{425.285}{598.916}{11.955}{19.26}{\longrightarrow}%
\TextAt{9}{447.867}{598.916}{1}%
\TextAt{13}{453.72}{598.916}{\Character{58}}%
\TextAt{8}{79.226}{571.452}{This}%
\TextAt{8}{105.946}{571.452}{gives}%
\TextAt{8}{134.626}{571.452}{the}%
\TextAt{8}{154.783}{571.452}{claim.}%
\Vector{q 0 G 0 g 0.399 w 452.334 -202.548 m 452.334 -194.479 l S Q}%
\Vector{q 0 G 0 g 0.399 w 452.534 -194.678 m 459.707 -194.678 l S Q}%
\Vector{q 0 G 0 g 0.399 w 452.534 -202.349 m 459.707 -202.349 l S Q}%
\Vector{q 0 G 0 g 0.399 w 459.906 -202.548 m 459.906 -194.479 l S Q}%
\TextAt{8}{97.159}{542.89}{The}%
\TextAt{8}{121.075}{542.89}{definition}%
\TextAt{8}{172.967}{542.89}{of}%
\TextAt{8}{186.16}{542.89}{the}%
\TextAt{8}{206.179}{542.89}{Picard}%
\TextAt{8}{243.27}{542.89}{group}%
\TextAt{8}{276.308}{542.89}{of}%
\TextAt{8}{289.501}{542.89}{a}%
\TextAt{8}{299.119}{542.89}{one-dimensional}%
\TextAt{8}{384.891}{542.89}{noetherian}%
\TextAt{8}{442.964}{542.89}{integral}%
\TextAt{8}{485.435}{542.89}{domain}%
\TextAt{14}{526.925}{542.89}{o}%
\TextAt{8}{79.226}{525.844}{avoids}%
\TextAt{8}{113.871}{525.844}{the}%
\TextAt{8}{133.154}{525.844}{problem}%
\TextAt{8}{177.805}{525.844}{of}%
\TextAt{8}{190.262}{525.844}{the}%
\TextAt{8}{209.544}{525.844}{uniqueness}%
\TextAt{8}{267.646}{525.844}{of}%
\TextAt{8}{280.102}{525.844}{prime}%
\TextAt{8}{312.392}{525.844}{ideal}%
\TextAt{8}{339.481}{525.844}{decomposition}%
\TextAt{8}{415.742}{525.844}{by}%
\TextAt{8}{431.128}{525.844}{restricting}%
\TextAt{8}{486.252}{525.844}{attention}%
\TextAt{8}{79.226}{508.798}{to}%
\TextAt{8}{93.587}{508.798}{the}%
\TextAt{8}{113.794}{508.798}{invertible}%
\TextAt{8}{165.551}{508.798}{ideals,}%
\TextAt{8}{201.444}{508.798}{and}%
\TextAt{8}{224.258}{508.798}{thus}%
\TextAt{8}{250.06}{508.798}{leaving}%
\TextAt{8}{289.79}{508.798}{aside}%
\TextAt{8}{319.167}{508.798}{the}%
\TextAt{8}{339.374}{508.798}{information}%
\TextAt{8}{402.847}{508.798}{carried}%
\TextAt{8}{441.919}{508.798}{by}%
\TextAt{8}{458.229}{508.798}{noninvertibles.}%
\TextAt{8}{79.226}{491.752}{But}%
\TextAt{8}{101.641}{491.752}{there}%
\TextAt{8}{130.727}{491.752}{is}%
\TextAt{8}{141.665}{491.752}{another}%
\TextAt{8}{183.771}{491.752}{important}%
\TextAt{8}{238.238}{491.752}{generalization}%
\TextAt{8}{312.215}{491.752}{of}%
\TextAt{8}{324.72}{491.752}{the}%
\TextAt{8}{344.05}{491.752}{ideal}%
\TextAt{8}{371.188}{491.752}{class}%
\TextAt{8}{397.799}{491.752}{group}%
\TextAt{8}{430.149}{491.752}{which}%
\TextAt{8}{462.81}{491.752}{does}%
\TextAt{8}{488.381}{491.752}{take}%
\TextAt{8}{512.924}{491.752}{into}%
\TextAt{8}{79.226}{474.706}{account}%
\TextAt{19}{122.106}{474.706}{all}%
\TextAt{8}{139.091}{474.706}{prime}%
\TextAt{8}{171.88}{474.706}{ideals}%
\TextAt{8}{204.084}{474.706}{of}%
\TextAt{14}{217.04}{474.706}{o}%
\TextAt{8}{222.884}{474.706}{.}%
\TextAt{8}{231.208}{474.706}{It}%
\TextAt{8}{243.506}{474.706}{is}%
\TextAt{8}{254.896}{474.706}{based}%
\TextAt{8}{287.1}{474.706}{on}%
\TextAt{8}{302.985}{474.706}{an}%
\TextAt{8}{318.87}{474.706}{artificial}%
\TextAt{8}{364.678}{474.706}{re-introduction}%
\TextAt{8}{444.954}{474.706}{of}%
\TextAt{8}{457.91}{474.706}{the}%
\TextAt{8}{477.692}{474.706}{uniqueness}%
\TextAt{8}{79.226}{457.66}{of}%
\TextAt{8}{92.391}{457.66}{prime}%
\TextAt{8}{125.39}{457.66}{decomposition.}%
\TextAt{8}{207.021}{457.66}{This}%
\TextAt{8}{233.575}{457.66}{group}%
\TextAt{8}{266.586}{457.66}{is}%
\TextAt{8}{278.185}{457.66}{called}%
\TextAt{8}{311.184}{457.66}{the}%
\TextAt{12}{331.175}{457.66}{divisor}%
\TextAt{12}{375.064}{457.66}{class}%
\TextAt{12}{406.233}{457.66}{group}%
\TextAt{8}{440.137}{457.66}{,}%
\TextAt{8}{447.155}{457.66}{or}%
\TextAt{12}{461.3}{457.66}{Chow}%
\TextAt{12}{498.865}{457.66}{group}%
\TextAt{8}{79.226}{440.614}{of}%
\TextAt{14}{92.556}{440.614}{o}%
\TextAt{8}{98.4}{440.614}{.}%
\TextAt{8}{106.849}{440.614}{Its}%
\TextAt{8}{124.136}{440.614}{definition}%
\TextAt{8}{176.165}{440.614}{starts}%
\TextAt{8}{208.814}{440.614}{from}%
\TextAt{8}{236.455}{440.614}{the}%
\TextAt{8}{256.611}{440.614}{free}%
\TextAt{8}{279.039}{440.614}{abelian}%
\TextAt{8}{319.687}{440.614}{group}%
\TextAt{13}{263.581}{413.15}{D}%
\TextAt{13}{273.623}{413.15}{iv}%
\TextAt{9}{283.708}{413.15}{\Character{40}}%
\TextAt{14}{288.26}{413.15}{o}%
\TextAt{9}{294.104}{413.15}{\Character{41}}%
\TextAt{9}{301.959}{413.15}{\Character{61}}%
\TextAt{20}{314.403}{424.507}{\Character{77}}%
\TextAt{27}{321.441}{399.401}{p}%
\TextAt{18}{334.461}{413.15}{Z}%
\TextAt{14}{342.431}{413.15}{p}%
\TextAt{8}{79.226}{376.318}{on}%
\TextAt{8}{95.301}{376.318}{the}%
\TextAt{8}{115.273}{376.318}{set}%
\TextAt{8}{133.356}{376.318}{of}%
\TextAt{8}{146.502}{376.318}{all}%
\TextAt{8}{162.577}{376.318}{maximal}%
\TextAt{8}{210.201}{376.318}{ideals}%
\TextAt{14}{242.595}{376.318}{p}%
\TextAt{8}{252.291}{376.318}{of}%
\TextAt{14}{265.437}{376.318}{o}%
\TextAt{8}{274.994}{376.318}{(i.e.,}%
\TextAt{8}{301.507}{376.318}{the}%
\TextAt{8}{321.479}{376.318}{set}%
\TextAt{8}{339.562}{376.318}{of}%
\TextAt{8}{352.708}{376.318}{all}%
\TextAt{8}{368.783}{376.318}{prime}%
\TextAt{8}{401.762}{376.318}{ideals}%
\NativeMathAt{434.156}{376.318}{11.955}{\neq}%
\TextAt{9}{446.565}{376.318}{0}%
\TextAt{8}{452.434}{376.318}{).}%
\TextAt{8}{465.376}{376.318}{This}%
\TextAt{8}{491.912}{376.318}{group}%
\TextAt{8}{524.903}{376.318}{is}%
\TextAt{8}{79.226}{359.272}{called}%
\TextAt{8}{112.39}{359.272}{the}%
\TextAt{12}{132.546}{359.272}{divisor}%
\TextAt{12}{176.625}{359.272}{group}%
\TextAt{8}{214.427}{359.272}{of}%
\TextAt{14}{227.757}{359.272}{o}%
\TextAt{8}{233.601}{359.272}{.}%
\TextAt{8}{242.05}{359.272}{Its}%
\TextAt{8}{259.337}{359.272}{elements}%
\TextAt{8}{307.193}{359.272}{are}%
\TextAt{8}{326.704}{359.272}{formal}%
\TextAt{8}{363.454}{359.272}{sums}%
\TextAt{13}{274.743}{331.808}{D}%
\TextAt{9}{288.108}{331.808}{\Character{61}}%
\TextAt{20}{300.534}{343.166}{\Character{88}}%
\TextAt{27}{307.174}{318.06}{p}%
\TextAt{13}{319.795}{331.808}{n}%
\TextAt{27}{326.783}{330.015}{p}%
\TextAt{14}{331.269}{331.808}{p}%
\TextAt{8}{79.226}{294.977}{with}%
\TextAt{13}{106.296}{294.977}{n}%
\TextAt{27}{113.284}{293.184}{p}%
\NativeMathAt{121.79}{294.977}{11.955}{\in}%
\TextAt{18}{133.78}{294.977}{Z}%
\TextAt{8}{146.057}{294.977}{and}%
\TextAt{13}{169.23}{294.977}{n}%
\TextAt{27}{176.217}{293.184}{p}%
\TextAt{9}{184.724}{294.977}{\Character{61}}%
\TextAt{9}{197.839}{294.977}{0}%
\TextAt{8}{208.008}{294.977}{for}%
\TextAt{8}{226.303}{294.977}{almost}%
\TextAt{8}{264.503}{294.977}{all}%
\TextAt{14}{281.172}{294.977}{p}%
\TextAt{8}{287.156}{294.977}{,}%
\TextAt{8}{294.817}{294.977}{called}%
\TextAt{12}{328.391}{294.977}{divisors}%
\TextAt{8}{377.602}{294.977}{(or}%
\TextAt{12}{396.877}{294.977}{0-cycles}%
\TextAt{8}{442.331}{294.977}{).}%
\TextAt{8}{456.564}{294.977}{Corollary}%
\TextAt{8}{508.692}{294.977}{(3.9)}%
\TextAt{8}{79.226}{277.931}{simply}%
\TextAt{8}{117.614}{277.931}{says}%
\TextAt{8}{143.39}{277.931}{that,}%
\TextAt{8}{173.176}{277.931}{in}%
\TextAt{8}{187.762}{277.931}{the}%
\TextAt{8}{208.851}{277.931}{case}%
\TextAt{8}{234.554}{277.931}{of}%
\TextAt{8}{248.817}{277.931}{a}%
\TextAt{8}{259.505}{277.931}{Dedekind}%
\TextAt{8}{312.622}{277.931}{domain,}%
\TextAt{8}{358.667}{277.931}{the}%
\TextAt{8}{379.756}{277.931}{divisor}%
\TextAt{8}{418.802}{277.931}{group}%
\TextAt{13}{452.91}{277.931}{D}%
\TextAt{13}{462.952}{277.931}{iv}%
\TextAt{9}{473.036}{277.931}{\Character{40}}%
\TextAt{14}{477.589}{277.931}{o}%
\TextAt{9}{483.433}{277.931}{\Character{41}}%
\TextAt{8}{492.815}{277.931}{and}%
\TextAt{8}{516.51}{277.931}{the}%
\TextAt{8}{79.226}{260.885}{group}%
\TextAt{8}{113.124}{260.885}{of}%
\TextAt{8}{127.176}{260.885}{ideals}%
\TextAt{8}{160.477}{260.885}{are}%
\TextAt{8}{180.71}{260.885}{canonically}%
\TextAt{8}{242.248}{260.885}{isomorphic.}%
\TextAt{8}{308.216}{260.885}{The}%
\TextAt{8}{332.992}{260.885}{additive}%
\TextAt{8}{378.594}{260.885}{notation}%
\TextAt{8}{426.157}{260.885}{and}%
\TextAt{8}{449.642}{260.885}{the}%
\TextAt{8}{470.521}{260.885}{name}%
\TextAt{8}{502.458}{260.885}{of}%
\TextAt{8}{516.51}{260.885}{the}%
\TextAt{8}{79.226}{243.839}{group}%
\TextAt{8}{112.321}{243.839}{stem}%
\TextAt{8}{140.263}{243.839}{from}%
\TextAt{8}{167.822}{243.839}{function}%
\TextAt{8}{213.589}{243.839}{theory}%
\TextAt{8}{250.259}{243.839}{where}%
\TextAt{8}{283.987}{243.839}{divisors}%
\TextAt{8}{326.634}{243.839}{for}%
\TextAt{8}{344.438}{243.839}{analytic}%
\TextAt{8}{388.914}{243.839}{functions}%
\TextAt{8}{439.296}{243.839}{play}%
\TextAt{8}{464.584}{243.839}{the}%
\TextAt{8}{484.659}{243.839}{same}%
\TextAt{8}{513.904}{243.839}{role}%
\TextAt{8}{79.226}{226.793}{as}%
\TextAt{8}{93.596}{226.793}{ideals}%
\TextAt{8}{126.174}{226.793}{do}%
\TextAt{8}{142.433}{226.793}{for}%
\TextAt{8}{160.318}{226.793}{algebraic}%
\TextAt{8}{209.753}{226.793}{numbers}%
\TextAt{8}{256.964}{226.793}{(see}%
\TextAt{8}{280.432}{226.793}{chap.}%
\TextAt{8}{311.324}{226.793}{III,}%
\TextAt{8}{331.779}{226.793}{§3).}%
\TextAt{8}{97.159}{206.758}{In}%
\TextAt{8}{113.056}{206.758}{order}%
\TextAt{8}{144.901}{206.758}{to}%
\TextAt{8}{160.488}{206.758}{define}%
\TextAt{8}{195.573}{206.758}{the}%
\TextAt{8}{217.005}{206.758}{divisor}%
\TextAt{8}{256.394}{206.758}{class}%
\TextAt{8}{285.107}{206.758}{group}%
\TextAt{8}{319.559}{206.758}{we}%
\TextAt{8}{338.062}{206.758}{have}%
\TextAt{8}{366.333}{206.758}{to}%
\TextAt{8}{381.919}{206.758}{associate}%
\TextAt{8}{432.427}{206.758}{to}%
\TextAt{8}{448.013}{206.758}{every}%
\TextAt{13}{480.181}{206.758}{f}%
\NativeMathAt{492.723}{206.758}{11.955}{\in}%
\TextAt{13}{506.188}{206.758}{K}%
\TextAt{3}{517.006}{211.097}{\Character{3}}%
\TextAt{8}{526.911}{206.758}{a}%
\TextAt{8}{79.226}{189.712}{“principal}%
\TextAt{8}{133.36}{189.712}{divisor”}%
\TextAt{8}{176.83}{189.712}{div}%
\TextAt{9}{192.766}{189.712}{\Character{40}}%
\TextAt{13}{197.319}{189.712}{f}%
\TextAt{9}{204.365}{189.712}{\Character{41}}%
\TextAt{8}{208.917}{189.712}{.}%
\TextAt{8}{217.31}{189.712}{We}%
\TextAt{8}{237.288}{189.712}{use}%
\TextAt{8}{257.338}{189.712}{the}%
\TextAt{8}{277.328}{189.712}{case}%
\TextAt{8}{301.932}{189.712}{of}%
\TextAt{8}{315.096}{189.712}{a}%
\TextAt{8}{324.685}{189.712}{Dedekind}%
\TextAt{8}{376.703}{189.712}{domain}%
\TextAt{8}{418.164}{189.712}{to}%
\TextAt{8}{432.308}{189.712}{guide}%
\TextAt{8}{463.357}{189.712}{us.}%
\TextAt{8}{482.868}{189.712}{There}%
\TextAt{8}{516.51}{189.712}{the}%
\TextAt{8}{79.226}{172.666}{principal}%
\TextAt{8}{128.003}{172.666}{ideal}%
\TextAt{9}{155.966}{172.666}{\Character{40}}%
\TextAt{13}{160.519}{172.666}{f}%
\TextAt{9}{167.565}{172.666}{\Character{41}}%
\TextAt{8}{176.015}{172.666}{was}%
\TextAt{8}{198.514}{172.666}{given}%
\TextAt{8}{229.084}{172.666}{by}%
\TextAt{9}{266.679}{155.62}{\Character{40}}%
\TextAt{13}{271.231}{155.62}{f}%
\TextAt{9}{278.277}{155.62}{\Character{41}}%
\TextAt{9}{286.132}{155.62}{\Character{61}}%
\TextAt{20}{298.576}{166.978}{\Character{89}}%
\TextAt{27}{304.22}{141.872}{p}%
\TextAt{14}{315.845}{155.62}{p}%
\TextAt{21}{321.828}{160.556}{v}%
\TextAt{28}{325.951}{159.56}{p}%
\TextAt{17}{330.034}{160.556}{\Character{40}}%
\TextAt{21}{333.327}{160.556}{f}%
\TextAt{17}{338.274}{160.556}{\Character{41}}%
\TextAt{13}{342.065}{155.62}{\Character{59}}%
\TextAt{8}{79.226}{124.059}{where}%
\TextAt{13}{113.703}{124.059}{v}%
\TextAt{27}{119.362}{122.266}{p}%
\TextAt{9}{128.307}{124.059}{\Character{58}}%
\TextAt{13}{136.016}{124.059}{K}%
\TextAt{3}{146.834}{128.398}{\Character{3}}%
\NativeMathAt{156.024}{124.059}{11.955}{\rightarrow}%
\TextAt{18}{172.437}{124.059}{Z}%
\TextAt{8}{184.972}{124.059}{is}%
\TextAt{8}{197.403}{124.059}{the}%
\TextAt{14}{218.227}{124.059}{p}%
\TextAt{8}{224.21}{124.059}{-adic}%
\TextAt{8}{253.486}{124.059}{exponential}%
\TextAt{8}{316.918}{124.059}{valuation}%
\TextAt{8}{368.658}{124.059}{associated}%
\TextAt{8}{425.06}{124.059}{to}%
\TextAt{8}{440.038}{124.059}{the}%
\TextAt{8}{460.861}{124.059}{valuation}%
\TextAt{8}{512.601}{124.059}{ring}%
\TextAt{14}{79.226}{107.013}{o}%
\TextAt{27}{85.07}{105.22}{p}%
\TextAt{8}{89.557}{107.013}{.}%
\TextAt{8}{101.356}{107.013}{In}%
\TextAt{8}{117.094}{107.013}{general,}%
\TextAt{14}{162.067}{107.013}{o}%
\TextAt{27}{167.911}{105.22}{p}%
\TextAt{8}{177.412}{107.013}{is}%
\TextAt{8}{190.293}{107.013}{not}%
\TextAt{8}{212.223}{107.013}{anymore}%
\TextAt{8}{260.826}{107.013}{a}%
\TextAt{8}{271.698}{107.013}{discrete}%
\TextAt{8}{315.794}{107.013}{valuation}%
\TextAt{8}{367.983}{107.013}{ring.}%
\TextAt{8}{399.95}{107.013}{Nevertheless,}%
\TextAt{14}{472.025}{107.013}{o}%
\TextAt{27}{477.87}{105.22}{p}%
\TextAt{8}{487.371}{107.013}{defines}%
\TextAt{8}{526.911}{107.013}{a}%
\TextAt{8}{79.226}{89.967}{homomorphism}%
\TextAt{8}{265.638}{72.921}{ord}%
\TextAt{27}{282.555}{71.128}{p}%
\TextAt{9}{290.362}{72.921}{\Character{58}}%
\TextAt{13}{296.935}{72.921}{K}%
\TextAt{3}{307.752}{77.858}{\Character{3}}%
\NativeArrowAt{315.806}{72.921}{11.955}{19.26}{\longrightarrow}%
\TextAt{18}{338.388}{72.921}{Z}%
\TextAt{8}{300.14}{43.033}{76}%
\Destination{78.23}{809.815}{nav.page.76}%
\Destination{79.36}{557.855}{nav.print.82}%
\Hotlink{310.824}{224.555}{14.294}{11.047}{nav.chapter.3}%
\Hotlink{336.958}{224.555}{6.858}{11.047}{nav.section.3.3}%
\Hotlink{133.863}{624.155}{30.936}{11.047}{nav.number.1.12.9}%
\Hotlink{190.459}{624.155}{36.794}{11.047}{nav.number.1.12.11}%
\Hotlink{508.192}{292.755}{25.078}{11.047}{nav.number.1.3.9}%
}

% Page 89 (body page 77)
\Page{%
\TextAt{2}{79.226}{796.943}{I.}%
\TextAt{2}{89.83}{796.943}{Algebraic}%
\TextAt{2}{138.648}{796.943}{Integers}%
\Vector{q 0 G 0 g 0.299 w 7.226 18.729 m 460.769 18.729 l S Q}%
\TextAt{8}{79.226}{761.124}{which}%
\TextAt{8}{112.713}{761.124}{generalizes}%
\TextAt{8}{171.305}{761.124}{the}%
\TextAt{8}{191.461}{761.124}{valuation}%
\TextAt{8}{242.534}{761.124}{function.}%
\TextAt{8}{292.933}{761.124}{If}%
\TextAt{13}{304.625}{761.124}{f}%
\TextAt{9}{314.992}{761.124}{\Character{61}}%
\TextAt{13}{327.418}{761.124}{a}%
\TextAt{9}{333.563}{761.124}{\Character{47}}%
\TextAt{13}{339.416}{761.124}{b}%
\NativeMathAt{347.714}{761.124}{11.955}{\in}%
\TextAt{13}{359.005}{761.124}{K}%
\TextAt{3}{369.822}{765.462}{\Character{3}}%
\TextAt{8}{374.555}{761.124}{,}%
\TextAt{8}{381.704}{761.124}{with}%
\TextAt{13}{408.364}{761.124}{a\Character{59}}%
\TextAt{13}{419.733}{761.124}{b}%
\NativeMathAt{428.051}{761.124}{11.955}{\in}%
\TextAt{14}{439.342}{761.124}{o}%
\TextAt{8}{445.186}{761.124}{,}%
\TextAt{8}{452.335}{761.124}{then}%
\TextAt{8}{478.995}{761.124}{we}%
\TextAt{8}{496.223}{761.124}{put}%
\TextAt{8}{215.846}{732.123}{ord}%
\TextAt{27}{232.762}{730.329}{p}%
\TextAt{9}{237.249}{732.123}{\Character{40}}%
\TextAt{13}{241.801}{732.123}{f}%
\TextAt{9}{248.848}{732.123}{\Character{41}}%
\TextAt{9}{256.703}{732.123}{\Character{61}}%
\TextAt{13}{269.147}{732.123}{\Character{96}}%
\TextAt{27}{274.059}{730.329}{o}%
\TextAt{28}{277.955}{729.333}{p}%
\TextAt{9}{282.536}{732.123}{\Character{40}}%
\TextAt{14}{287.088}{732.123}{o}%
\TextAt{27}{292.932}{730.329}{p}%
\TextAt{9}{297.419}{732.123}{\Character{47}}%
\TextAt{13}{303.272}{732.123}{a}%
\TextAt{14}{309.417}{732.123}{o}%
\TextAt{27}{315.261}{730.329}{p}%
\TextAt{9}{319.748}{732.123}{\Character{41}}%
\TextAt{16}{326.957}{732.123}{\Character{0}}%
\TextAt{13}{338.912}{732.123}{\Character{96}}%
\TextAt{27}{343.824}{730.329}{o}%
\TextAt{28}{347.72}{729.333}{p}%
\TextAt{9}{352.301}{732.123}{\Character{40}}%
\TextAt{14}{356.854}{732.123}{o}%
\TextAt{27}{362.698}{730.329}{p}%
\TextAt{9}{367.185}{732.123}{\Character{47}}%
\TextAt{13}{373.038}{732.123}{b}%
\TextAt{14}{378.015}{732.123}{o}%
\TextAt{27}{383.859}{730.329}{p}%
\TextAt{9}{388.346}{732.123}{\Character{41}}%
\TextAt{13}{392.898}{732.123}{\Character{59}}%
\TextAt{8}{79.226}{703.121}{where}%
\TextAt{13}{112.191}{703.121}{\Character{96}}%
\TextAt{27}{117.103}{701.328}{o}%
\TextAt{28}{120.999}{700.332}{p}%
\TextAt{9}{125.58}{703.121}{\Character{40}}%
\TextAt{13}{130.132}{703.121}{M}%
\TextAt{9}{142.706}{703.121}{\Character{41}}%
\TextAt{8}{150.311}{703.121}{denotes}%
\TextAt{8}{191.8}{703.121}{the}%
\TextAt{12}{211.112}{703.121}{length}%
\TextAt{8}{250.952}{703.121}{of}%
\TextAt{8}{263.437}{703.121}{an}%
\TextAt{14}{278.852}{703.121}{o}%
\TextAt{27}{284.696}{701.328}{p}%
\TextAt{8}{289.183}{703.121}{-module}%
\TextAt{13}{333.529}{703.121}{M}%
\TextAt{8}{346.103}{703.121}{,}%
\TextAt{8}{352.577}{703.121}{i.e.,}%
\TextAt{8}{374.006}{703.121}{the}%
\TextAt{8}{393.318}{703.121}{maximal}%
\TextAt{8}{440.283}{703.121}{length}%
\TextAt{8}{475.208}{703.121}{of}%
\TextAt{8}{487.694}{703.121}{a}%
\TextAt{8}{496.605}{703.121}{strictly}%
\TextAt{8}{79.226}{686.075}{decreasing}%
\TextAt{8}{135.87}{686.075}{chain}%
\TextAt{13}{226.456}{669.029}{M}%
\TextAt{9}{242.35}{669.029}{\Character{61}}%
\TextAt{13}{254.776}{669.029}{M}%
\TextAt{17}{266.101}{667.236}{0}%
\TextAt{18}{274.154}{669.029}{\Character{41}}%
\TextAt{13}{286.773}{669.029}{M}%
\TextAt{17}{298.098}{667.236}{1}%
\TextAt{18}{306.151}{669.029}{\Character{41}}%
\TextAt{16}{318.77}{669.029}{\Character{1}}%
\TextAt{16}{324.066}{669.029}{\Character{1}}%
\TextAt{16}{329.374}{669.029}{\Character{1}}%
\TextAt{18}{336.039}{669.029}{\Character{41}}%
\TextAt{13}{348.658}{669.029}{M}%
\TextAt{21}{359.983}{667.236}{\Character{96}}%
\TextAt{9}{367.261}{669.029}{\Character{61}}%
\TextAt{9}{379.671}{669.029}{0}%
\TextAt{8}{79.226}{645.508}{of}%
\TextAt{14}{92.682}{645.508}{o}%
\TextAt{27}{98.526}{643.715}{p}%
\TextAt{8}{103.013}{645.508}{-submodules.}%
\TextAt{8}{175.368}{645.508}{In}%
\TextAt{8}{190.114}{645.508}{the}%
\TextAt{8}{210.396}{645.508}{special}%
\TextAt{8}{248.623}{645.508}{case}%
\TextAt{8}{273.52}{645.508}{where}%
\TextAt{14}{307.454}{645.508}{o}%
\TextAt{27}{313.299}{643.715}{p}%
\TextAt{8}{321.808}{645.508}{is}%
\TextAt{8}{333.698}{645.508}{a}%
\TextAt{8}{343.579}{645.508}{discrete}%
\TextAt{8}{386.683}{645.508}{valuation}%
\TextAt{8}{437.881}{645.508}{ring}%
\TextAt{8}{462.073}{645.508}{with}%
\TextAt{8}{488.858}{645.508}{maximal}%
\TextAt{8}{79.226}{628.462}{ideal}%
\TextAt{14}{107.189}{628.462}{m}%
\TextAt{8}{116.356}{628.462}{,}%
\TextAt{8}{123.505}{628.462}{the}%
\TextAt{8}{143.662}{628.462}{value}%
\TextAt{13}{173.908}{628.462}{\Character{23}}%
\TextAt{9}{183.767}{628.462}{\Character{61}}%
\TextAt{13}{196.192}{628.462}{v}%
\TextAt{27}{201.852}{626.669}{p}%
\TextAt{9}{206.339}{628.462}{\Character{40}}%
\TextAt{13}{210.891}{628.462}{a}%
\TextAt{9}{217.036}{628.462}{\Character{41}}%
\TextAt{8}{225.486}{628.462}{of}%
\TextAt{13}{238.816}{628.462}{a}%
\NativeMathAt{248.281}{628.462}{11.955}{\in}%
\TextAt{14}{259.572}{628.462}{o}%
\TextAt{27}{265.417}{626.669}{p}%
\TextAt{8}{269.903}{628.462}{,}%
\TextAt{8}{277.053}{628.462}{for}%
\TextAt{13}{294.938}{628.462}{a}%
\NativeMathAt{304.403}{628.462}{11.955}{\neq}%
\TextAt{9}{316.812}{628.462}{0}%
\TextAt{8}{322.682}{628.462}{,}%
\TextAt{8}{329.831}{628.462}{is}%
\TextAt{8}{341.595}{628.462}{given}%
\TextAt{8}{372.164}{628.462}{by}%
\TextAt{8}{388.423}{628.462}{the}%
\TextAt{8}{408.58}{628.462}{equation}%
\TextAt{13}{281.086}{599.461}{a}%
\TextAt{14}{287.231}{599.461}{o}%
\TextAt{27}{293.076}{597.667}{p}%
\TextAt{9}{300.883}{599.461}{\Character{61}}%
\TextAt{14}{313.309}{599.461}{m}%
\TextAt{21}{322.476}{604.397}{\Character{23}}%
\TextAt{13}{327.658}{599.461}{\Character{58}}%
\TextAt{8}{79.226}{570.46}{It}%
\TextAt{8}{91.509}{570.46}{is}%
\TextAt{8}{102.884}{570.46}{equal}%
\TextAt{8}{133.386}{570.46}{to}%
\TextAt{8}{147.307}{570.46}{the}%
\TextAt{8}{167.074}{570.46}{length}%
\TextAt{8}{202.455}{570.46}{of}%
\TextAt{8}{215.395}{570.46}{the}%
\TextAt{14}{235.162}{570.46}{o}%
\TextAt{27}{241.007}{568.666}{p}%
\TextAt{8}{245.493}{570.46}{-module}%
\TextAt{14}{290.295}{570.46}{o}%
\TextAt{27}{296.139}{568.666}{p}%
\TextAt{9}{300.626}{570.46}{\Character{47}}%
\TextAt{14}{306.479}{570.46}{m}%
\TextAt{21}{315.645}{574.798}{\Character{23}}%
\TextAt{8}{320.827}{570.46}{,}%
\TextAt{8}{327.665}{570.46}{because}%
\TextAt{8}{370.577}{570.46}{the}%
\TextAt{8}{390.344}{570.46}{longest}%
\TextAt{8}{429.694}{570.46}{chain}%
\TextAt{8}{460.197}{570.46}{of}%
\TextAt{8}{473.137}{570.46}{submodules}%
\TextAt{8}{79.226}{553.414}{is}%
\TextAt{14}{210.91}{536.368}{o}%
\TextAt{27}{216.754}{534.574}{p}%
\TextAt{9}{221.241}{536.368}{\Character{47}}%
\TextAt{14}{227.094}{536.368}{m}%
\TextAt{21}{236.261}{541.304}{\Character{23}}%
\NativeMathAt{244.764}{536.368}{11.955}{\supset}%
\TextAt{14}{257.383}{536.368}{m}%
\TextAt{9}{266.55}{536.368}{\Character{47}}%
\TextAt{14}{272.403}{536.368}{m}%
\TextAt{21}{281.57}{541.304}{\Character{23}}%
\NativeMathAt{290.072}{536.368}{11.955}{\supset}%
\TextAt{16}{302.673}{536.368}{\Character{1}}%
\TextAt{16}{307.981}{536.368}{\Character{1}}%
\TextAt{16}{313.289}{536.368}{\Character{1}}%
\NativeMathAt{319.912}{536.368}{11.955}{\supset}%
\TextAt{14}{332.58}{536.368}{m}%
\TextAt{21}{341.746}{541.304}{\Character{23}}%
\TextAt{9}{346.928}{536.368}{\Character{47}}%
\TextAt{14}{352.781}{536.368}{m}%
\TextAt{21}{361.948}{541.304}{\Character{23}}%
\TextAt{9}{370.451}{536.368}{\Character{61}}%
\TextAt{9}{382.861}{536.368}{\Character{40}0\Character{41}}%
\TextAt{13}{397.834}{536.368}{\Character{58}}%
\TextAt{8}{79.226}{512.846}{Thus}%
\TextAt{8}{108.875}{512.846}{the}%
\TextAt{8}{129.031}{512.846}{function}%
\TextAt{8}{174.879}{512.846}{ord}%
\TextAt{27}{191.796}{511.053}{p}%
\TextAt{8}{200.18}{512.846}{agrees}%
\TextAt{8}{235.364}{512.846}{with}%
\TextAt{8}{262.024}{512.846}{the}%
\TextAt{8}{282.181}{512.846}{exponential}%
\TextAt{8}{344.945}{512.846}{valuation}%
\TextAt{13}{396.018}{512.846}{v}%
\TextAt{27}{401.677}{511.053}{p}%
\TextAt{8}{410.061}{512.846}{in}%
\TextAt{8}{423.714}{512.846}{this}%
\TextAt{8}{446.537}{512.846}{case.}%
\TextAt{8}{97.159}{492.811}{The}%
\TextAt{8}{120.712}{492.811}{property}%
\TextAt{8}{168.02}{492.811}{of}%
\TextAt{8}{180.849}{492.811}{the}%
\TextAt{8}{200.505}{492.811}{function}%
\TextAt{8}{245.853}{492.811}{ord}%
\TextAt{27}{262.769}{491.018}{p}%
\TextAt{8}{270.653}{492.811}{to}%
\TextAt{8}{284.462}{492.811}{be}%
\TextAt{8}{299.886}{492.811}{a}%
\TextAt{8}{309.141}{492.811}{homomorphism}%
\TextAt{8}{391.31}{492.811}{follows}%
\TextAt{8}{429.245}{492.811}{from}%
\TextAt{8}{456.385}{492.811}{the}%
\TextAt{8}{476.041}{492.811}{fact}%
\TextAt{8}{498.626}{492.811}{(which}%
\ErrataAt{79.22}{475.755}{453.64}{is easily proved) that $\ell_{\mathfrak o_{\mathfrak p}}$ is additive on short exact sequences of}%
\ErrataAt{79.22}{458.755}{453.64}{$\mathfrak o_{\mathfrak p}$-modules.}%
\TextAt{8}{97.159}{438.684}{Using}%
\TextAt{8}{130.691}{438.684}{the}%
\TextAt{8}{151.48}{438.684}{functions}%
\TextAt{8}{202.574}{438.684}{ord}%
\TextAt{27}{219.491}{436.891}{p}%
\TextAt{9}{228.375}{438.684}{\Character{58}}%
\TextAt{13}{236.024}{438.684}{K}%
\TextAt{3}{246.842}{443.023}{\Character{3}}%
\NativeMathAt{255.972}{438.684}{11.955}{\rightarrow}%
\TextAt{18}{272.325}{438.684}{Z}%
\TextAt{8}{280.295}{438.684}{,}%
\TextAt{8}{288.234}{438.684}{we}%
\TextAt{8}{306.093}{438.684}{can}%
\TextAt{8}{328.184}{438.684}{now}%
\TextAt{8}{353.205}{438.684}{associate}%
\TextAt{8}{403.068}{438.684}{to}%
\TextAt{8}{418.01}{438.684}{every}%
\TextAt{8}{449.534}{438.684}{element}%
\TextAt{13}{493.408}{438.684}{f}%
\NativeMathAt{504.852}{438.684}{11.955}{\in}%
\TextAt{13}{517.22}{438.684}{K}%
\TextAt{3}{528.037}{443.023}{\Character{3}}%
\TextAt{8}{79.226}{421.638}{the}%
\TextAt{8}{99.383}{421.638}{divisor}%
\TextAt{8}{249.056}{404.592}{div}%
\TextAt{9}{264.992}{404.592}{\Character{40}}%
\TextAt{13}{269.544}{404.592}{f}%
\TextAt{9}{276.591}{404.592}{\Character{41}}%
\TextAt{9}{284.446}{404.592}{\Character{61}}%
\TextAt{20}{296.889}{415.95}{\Character{88}}%
\TextAt{27}{303.529}{390.844}{p}%
\TextAt{8}{316.151}{404.592}{ord}%
\TextAt{27}{333.067}{402.799}{p}%
\TextAt{9}{337.554}{404.592}{\Character{40}}%
\TextAt{13}{342.106}{404.592}{f}%
\TextAt{9}{349.153}{404.592}{\Character{41}}%
\TextAt{14}{353.705}{404.592}{p}%
\TextAt{13}{359.688}{404.592}{\Character{59}}%
\TextAt{8}{79.226}{372.373}{and}%
\TextAt{8}{101.989}{372.373}{thus}%
\TextAt{8}{127.74}{372.373}{obtain}%
\TextAt{8}{164.168}{372.373}{a}%
\TextAt{8}{173.923}{372.373}{canonical}%
\TextAt{8}{225.306}{372.373}{homomorphism}%
\TextAt{8}{253.193}{343.372}{div}%
\TextAt{9}{272.45}{343.372}{\Character{58}}%
\TextAt{13}{279.023}{343.372}{K}%
\TextAt{3}{289.841}{348.308}{\Character{3}}%
\NativeArrowAt{297.894}{343.372}{11.955}{19.26}{\longrightarrow}%
\TextAt{13}{320.476}{343.372}{D}%
\TextAt{13}{330.518}{343.372}{iv}%
\TextAt{9}{340.602}{343.372}{\Character{40}}%
\TextAt{14}{345.154}{343.372}{o}%
\TextAt{9}{350.998}{343.372}{\Character{41}}%
\TextAt{13}{355.551}{343.372}{\Character{58}}%
\TextAt{8}{79.226}{314.37}{The}%
\TextAt{8}{102.098}{314.37}{elements}%
\TextAt{8}{148.774}{314.37}{div}%
\TextAt{9}{164.71}{314.371}{\Character{40}}%
\TextAt{13}{169.262}{314.37}{f}%
\TextAt{9}{176.309}{314.371}{\Character{41}}%
\TextAt{8}{183.577}{314.37}{are}%
\TextAt{8}{201.906}{314.37}{called}%
\TextAt{12}{233.888}{314.37}{principal}%
\TextAt{12}{288.671}{314.37}{divisors}%
\TextAt{8}{333.574}{314.37}{.}%
\TextAt{8}{341.629}{314.37}{They}%
\TextAt{8}{370.682}{314.37}{form}%
\TextAt{8}{397.141}{314.37}{a}%
\TextAt{8}{405.715}{314.37}{subgroup}%
\TextAt{16}{455.331}{314.37}{\Character{80}}%
\TextAt{9}{464.629}{314.371}{\Character{40}}%
\TextAt{14}{469.182}{314.37}{o}%
\TextAt{9}{475.026}{314.371}{\Character{41}}%
\TextAt{8}{482.294}{314.37}{of}%
\TextAt{13}{494.443}{314.37}{D}%
\TextAt{13}{504.485}{314.37}{iv}%
\TextAt{9}{514.569}{314.371}{\Character{40}}%
\TextAt{14}{519.121}{314.37}{o}%
\TextAt{9}{524.965}{314.371}{\Character{41}}%
\TextAt{8}{529.518}{314.37}{.}%
\TextAt{8}{79.226}{297.325}{Two}%
\TextAt{8}{106.06}{297.325}{divisors}%
\TextAt{13}{150.265}{297.325}{D}%
\TextAt{8}{165.684}{297.325}{and}%
\TextAt{13}{189.924}{297.325}{D}%
\TextAt{3}{199.969}{301.663}{\Character{48}}%
\TextAt{8}{208.138}{297.325}{which}%
\TextAt{8}{243.102}{297.325}{differ}%
\TextAt{8}{274.814}{297.325}{only}%
\TextAt{8}{301.983}{297.325}{by}%
\TextAt{8}{319.719}{297.325}{a}%
\TextAt{8}{330.951}{297.325}{principal}%
\TextAt{8}{381.206}{297.325}{divisor}%
\TextAt{8}{420.796}{297.325}{are}%
\TextAt{8}{441.784}{297.325}{called}%
\TextAt{12}{476.425}{297.325}{rationally}%
\TextAt{12}{79.226}{280.279}{equivalent}%
\TextAt{8}{138.775}{280.279}{.}%
\TextAt{12}{79.226}{254.266}{Definition}%
\TextAt{12}{142.015}{254.266}{(12.13).}%
\TextAt{8}{190.792}{254.266}{The}%
\TextAt{8}{214.846}{254.266}{quotient}%
\TextAt{8}{261.028}{254.266}{group}%
\TextAt{13}{245.768}{225.265}{C}%
\TextAt{13}{254.986}{225.265}{H}%
\TextAt{17}{265.624}{230.201}{1}%
\TextAt{9}{270.357}{225.265}{\Character{40}}%
\TextAt{14}{274.909}{225.265}{o}%
\TextAt{9}{280.753}{225.265}{\Character{41}}%
\TextAt{9}{288.608}{225.265}{\Character{61}}%
\TextAt{13}{301.052}{225.265}{D}%
\TextAt{13}{311.094}{225.265}{iv}%
\TextAt{9}{321.178}{225.265}{\Character{40}}%
\TextAt{14}{325.73}{225.265}{o}%
\TextAt{9}{331.575}{225.265}{\Character{41}\Character{47}}%
\TextAt{16}{341.98}{225.265}{\Character{80}}%
\TextAt{9}{351.278}{225.265}{\Character{40}}%
\TextAt{14}{355.831}{225.265}{o}%
\TextAt{9}{361.675}{225.265}{\Character{41}}%
\TextAt{8}{79.226}{196.264}{is}%
\TextAt{8}{90.99}{196.264}{called}%
\TextAt{8}{124.154}{196.264}{the}%
\TextAt{12}{144.31}{196.264}{divisor}%
\TextAt{12}{188.389}{196.264}{class}%
\TextAt{12}{219.747}{196.264}{group}%
\TextAt{8}{257.549}{196.264}{or}%
\TextAt{12}{271.86}{196.264}{Chow}%
\TextAt{12}{309.614}{196.264}{group}%
\TextAt{8}{347.416}{196.264}{of}%
\TextAt{14}{360.746}{196.264}{o}%
\TextAt{8}{366.591}{196.264}{.}%
\TextAt{8}{97.159}{170.252}{The}%
\TextAt{8}{121.213}{170.252}{Chow}%
\TextAt{8}{154.054}{170.252}{group}%
\TextAt{8}{187.229}{170.252}{is}%
\TextAt{8}{198.993}{170.252}{related}%
\TextAt{8}{238.015}{170.252}{to}%
\TextAt{8}{252.325}{170.252}{the}%
\TextAt{8}{272.481}{170.252}{Picard}%
\TextAt{8}{309.71}{170.252}{group}%
\TextAt{8}{342.885}{170.252}{by}%
\TextAt{8}{359.144}{170.252}{a}%
\TextAt{8}{368.9}{170.252}{canonical}%
\TextAt{8}{420.283}{170.252}{homomorphism}%
\TextAt{8}{243.787}{141.25}{div}%
\TextAt{9}{263.044}{141.25}{\Character{58}}%
\TextAt{13}{269.617}{141.25}{P}%
\TextAt{13}{278.775}{141.25}{ic}%
\TextAt{9}{287.819}{141.25}{\Character{40}}%
\TextAt{14}{292.372}{141.25}{o}%
\TextAt{9}{298.216}{141.25}{\Character{41}}%
\NativeArrowAt{306.089}{141.25}{11.955}{19.26}{\longrightarrow}%
\TextAt{13}{328.671}{141.25}{C}%
\TextAt{13}{337.889}{141.25}{H}%
\TextAt{17}{348.527}{146.187}{1}%
\TextAt{9}{353.26}{141.25}{\Character{40}}%
\TextAt{14}{357.812}{141.25}{o}%
\TextAt{9}{363.656}{141.25}{\Character{41}}%
\TextAt{8}{300.14}{43.033}{77}%
\Destination{78.23}{809.815}{nav.page.77}%
\Destination{79.36}{585.455}{nav.print.83}%
\Destination{79.36}{268.255}{nav.number.1.12.13}%
\Destination{79.226}{489.755}{errata.11}%
}

% Page 90 (body page 78)
\Page{%
\TextAt{2}{79.226}{796.943}{I.}%
\TextAt{2}{89.83}{796.943}{Algebraic}%
\TextAt{2}{138.648}{796.943}{Integers}%
\Vector{q 0 G 0 g 0.299 w 7.226 18.729 m 460.769 18.729 l S Q}%
\TextAt{8}{79.226}{761.124}{which}%
\TextAt{8}{112.819}{761.124}{is}%
\TextAt{8}{124.688}{761.124}{defined}%
\TextAt{8}{165.107}{761.124}{as}%
\TextAt{8}{179.583}{761.124}{follows.}%
\TextAt{8}{222.888}{761.124}{If}%
\TextAt{14}{234.686}{761.124}{a}%
\TextAt{8}{244.673}{761.124}{is}%
\TextAt{8}{256.542}{761.124}{an}%
\TextAt{8}{272.907}{761.124}{invertible}%
\TextAt{8}{324.719}{761.124}{ideal,}%
\TextAt{8}{356.067}{761.124}{then,}%
\TextAt{8}{386.111}{761.124}{by}%
\TextAt{8}{402.476}{761.124}{(12.4),}%
\TextAt{14}{439.693}{761.124}{ao}%
\TextAt{27}{451.521}{759.33}{p}%
\TextAt{8}{456.008}{761.124}{,}%
\TextAt{8}{463.289}{761.124}{for}%
\TextAt{8}{481.28}{761.124}{any}%
\TextAt{8}{503.503}{761.124}{prime}%
\TextAt{8}{79.226}{744.078}{ideal}%
\TextAt{14}{107.189}{744.078}{p}%
\NativeMathAt{116.493}{744.078}{11.955}{\neq}%
\TextAt{9}{128.903}{744.078}{0}%
\TextAt{8}{134.772}{744.078}{,}%
\TextAt{8}{141.921}{744.078}{is}%
\TextAt{8}{153.685}{744.078}{a}%
\TextAt{8}{163.44}{744.078}{principal}%
\TextAt{8}{212.217}{744.078}{ideal}%
\TextAt{13}{240.18}{744.078}{a}%
\TextAt{27}{246.325}{742.284}{p}%
\TextAt{14}{250.812}{744.078}{o}%
\TextAt{27}{256.656}{742.284}{p}%
\TextAt{8}{261.143}{744.078}{,}%
\TextAt{13}{268.293}{744.078}{a}%
\TextAt{27}{274.437}{742.284}{p}%
\NativeMathAt{282.245}{744.078}{11.955}{\in}%
\TextAt{13}{293.536}{744.078}{K}%
\TextAt{3}{304.354}{748.416}{\Character{3}}%
\TextAt{8}{309.086}{744.078}{,}%
\TextAt{8}{316.235}{744.078}{and}%
\TextAt{8}{338.998}{744.078}{we}%
\TextAt{8}{356.225}{744.078}{put}%
\TextAt{8}{242.149}{715.077}{div}%
\TextAt{9}{258.086}{715.077}{\Character{40}}%
\TextAt{14}{262.638}{715.077}{a}%
\TextAt{9}{268.621}{715.077}{\Character{41}}%
\TextAt{9}{276.476}{715.077}{\Character{61}}%
\TextAt{20}{288.92}{726.434}{\Character{88}}%
\TextAt{27}{295.56}{701.328}{p}%
\TextAt{16}{308.181}{715.077}{\Character{0}}%
\TextAt{8}{319.472}{715.077}{ord}%
\TextAt{27}{336.388}{713.283}{p}%
\TextAt{9}{340.875}{715.077}{\Character{40}}%
\TextAt{13}{345.427}{715.077}{a}%
\TextAt{27}{351.572}{713.283}{p}%
\TextAt{9}{356.059}{715.077}{\Character{41}}%
\TextAt{14}{360.612}{715.077}{p}%
\TextAt{13}{366.595}{715.077}{\Character{58}}%
\TextAt{8}{79.226}{677.377}{This}%
\TextAt{8}{105.946}{677.377}{gives}%
\TextAt{8}{134.626}{677.377}{us}%
\TextAt{8}{149.642}{677.377}{a}%
\TextAt{8}{159.398}{677.377}{homomorphism}%
\TextAt{8}{251.331}{648.376}{div}%
\TextAt{9}{270.588}{648.376}{\Character{58}}%
\TextAt{13}{277.161}{648.376}{J}%
\TextAt{9}{284.738}{648.376}{\Character{40}}%
\TextAt{14}{289.29}{648.376}{o}%
\TextAt{9}{295.134}{648.376}{\Character{41}}%
\NativeArrowAt{303.007}{648.376}{11.955}{19.26}{\longrightarrow}%
\TextAt{13}{325.59}{648.376}{D}%
\TextAt{13}{335.632}{648.376}{iv}%
\TextAt{9}{345.716}{648.376}{\Character{40}}%
\TextAt{14}{350.268}{648.376}{o}%
\TextAt{9}{356.112}{648.376}{\Character{41}}%
\TextAt{8}{79.226}{619.375}{of}%
\TextAt{8}{92.647}{619.375}{the}%
\TextAt{8}{112.894}{619.375}{ideal}%
\TextAt{8}{140.949}{619.375}{group}%
\TextAt{13}{174.215}{619.375}{J}%
\TextAt{9}{181.792}{619.375}{\Character{40}}%
\TextAt{14}{186.345}{619.375}{o}%
\TextAt{9}{192.189}{619.375}{\Character{41}}%
\TextAt{8}{200.729}{619.375}{which}%
\TextAt{8}{234.307}{619.375}{takes}%
\TextAt{8}{264.381}{619.375}{principal}%
\TextAt{8}{313.249}{619.375}{ideals}%
\TextAt{8}{345.918}{619.375}{into}%
\TextAt{8}{369.752}{619.375}{principal}%
\TextAt{8}{418.62}{619.375}{divisors,}%
\TextAt{8}{464.713}{619.375}{and}%
\TextAt{8}{487.567}{619.375}{therefore}%
\TextAt{8}{79.226}{602.329}{induces}%
\TextAt{8}{120.902}{602.329}{a}%
\TextAt{8}{130.657}{602.329}{homomorphism}%
\TextAt{8}{242.161}{585.283}{div}%
\TextAt{9}{261.418}{585.283}{\Character{58}}%
\TextAt{13}{267.991}{585.283}{P}%
\TextAt{13}{277.149}{585.283}{ic}%
\TextAt{9}{286.193}{585.283}{\Character{40}}%
\TextAt{14}{290.746}{585.283}{o}%
\TextAt{9}{296.59}{585.283}{\Character{41}}%
\NativeArrowAt{304.463}{585.283}{11.955}{19.26}{\longrightarrow}%
\TextAt{13}{327.045}{585.283}{C}%
\TextAt{13}{336.262}{585.283}{H}%
\TextAt{17}{346.901}{590.219}{1}%
\TextAt{9}{351.634}{585.283}{\Character{40}}%
\TextAt{14}{356.186}{585.283}{o}%
\TextAt{9}{362.03}{585.283}{\Character{41}}%
\TextAt{13}{366.583}{585.283}{\Character{58}}%
\TextAt{8}{79.226}{561.762}{In}%
\TextAt{8}{93.847}{561.762}{the}%
\TextAt{8}{114.004}{561.762}{special}%
\TextAt{8}{152.105}{561.762}{case}%
\TextAt{8}{176.876}{561.762}{of}%
\TextAt{8}{190.206}{561.762}{a}%
\TextAt{8}{199.961}{561.762}{Dedekind}%
\TextAt{8}{252.146}{561.762}{domain}%
\TextAt{8}{293.774}{561.762}{we}%
\TextAt{8}{311.001}{561.762}{obtain:}%
\TextAt{12}{79.226}{535.749}{Proposition}%
\TextAt{12}{151.937}{535.749}{(12.14).}%
\TextAt{8}{200.715}{535.749}{If}%
\TextAt{14}{212.407}{535.749}{o}%
\TextAt{8}{222.148}{535.749}{is}%
\TextAt{8}{233.912}{535.749}{a}%
\TextAt{8}{243.668}{535.749}{Dedekind}%
\TextAt{8}{295.852}{535.749}{domain,}%
\TextAt{8}{340.732}{535.749}{then}%
\TextAt{8}{243.787}{506.748}{div}%
\TextAt{9}{263.044}{506.748}{\Character{58}}%
\TextAt{13}{269.617}{506.748}{P}%
\TextAt{13}{278.775}{506.748}{ic}%
\TextAt{9}{287.819}{506.748}{\Character{40}}%
\TextAt{14}{292.372}{506.748}{o}%
\TextAt{9}{298.216}{506.748}{\Character{41}}%
\NativeArrowAt{306.089}{506.748}{11.955}{19.26}{\longrightarrow}%
\TextAt{13}{328.671}{506.748}{C}%
\TextAt{13}{337.889}{506.748}{H}%
\TextAt{17}{348.527}{511.684}{1}%
\TextAt{9}{353.26}{506.748}{\Character{40}}%
\TextAt{14}{357.812}{506.748}{o}%
\TextAt{9}{363.656}{506.748}{\Character{41}}%
\TextAt{8}{79.226}{477.747}{is}%
\TextAt{8}{90.99}{477.747}{an}%
\TextAt{8}{107.249}{477.747}{isomorphism.}%
\TextAt{12}{79.226}{445.757}{Exercise}%
\TextAt{12}{132.427}{445.757}{I.12.1.}%
\TextAt{8}{172.751}{445.757}{Show}%
\TextAt{8}{203.644}{445.757}{that}%
\TextAt{18}{184.934}{420.622}{C}%
\TextAt{9}{193.568}{420.622}{\Character{91}}%
\TextAt{13}{196.82}{420.622}{X}%
\TextAt{13}{206.815}{420.622}{\Character{59}}%
\TextAt{13}{212.039}{420.622}{Y}%
\TextAt{9}{221.462}{420.622}{\Character{93}\Character{47}\Character{40}}%
\TextAt{13}{235.119}{420.622}{X}%
\TextAt{13}{245.759}{420.622}{Y}%
\TextAt{16}{257.824}{420.622}{\Character{0}}%
\TextAt{13}{269.779}{420.622}{X}%
\TextAt{9}{280.434}{420.622}{\Character{41}}%
\TextAt{13}{284.986}{420.622}{\Character{59}}%
\TextAt{18}{298.201}{420.622}{C}%
\TextAt{9}{306.835}{420.622}{\Character{91}}%
\TextAt{13}{310.087}{420.622}{X}%
\TextAt{13}{320.082}{420.622}{\Character{59}}%
\TextAt{13}{325.306}{420.622}{Y}%
\TextAt{9}{334.729}{420.622}{\Character{93}\Character{47}\Character{40}}%
\TextAt{13}{348.386}{420.622}{X}%
\TextAt{13}{359.026}{420.622}{Y}%
\TextAt{16}{371.091}{420.622}{\Character{0}}%
\TextAt{9}{383.046}{420.622}{1\Character{41}}%
\TextAt{13}{393.451}{420.622}{\Character{59}}%
\TextAt{18}{184.934}{403.576}{C}%
\TextAt{9}{193.568}{403.576}{\Character{91}}%
\TextAt{13}{196.82}{403.576}{X}%
\TextAt{13}{206.815}{403.576}{\Character{59}}%
\TextAt{13}{212.039}{403.576}{Y}%
\TextAt{9}{221.462}{403.576}{\Character{93}\Character{47}\Character{40}}%
\TextAt{13}{235.119}{403.576}{X}%
\TextAt{17}{245.774}{407.915}{2}%
\TextAt{16}{253.163}{403.576}{\Character{0}}%
\TextAt{13}{265.118}{403.576}{Y}%
\TextAt{17}{274.511}{407.915}{3}%
\TextAt{9}{279.244}{403.576}{\Character{41}}%
\TextAt{13}{283.796}{403.576}{\Character{59}}%
\TextAt{18}{298.201}{403.576}{C}%
\TextAt{9}{306.835}{403.576}{\Character{91}}%
\TextAt{13}{310.087}{403.576}{X}%
\TextAt{13}{320.082}{403.576}{\Character{59}}%
\TextAt{13}{325.306}{403.576}{Y}%
\TextAt{9}{334.729}{403.576}{\Character{93}\Character{47}\Character{40}}%
\TextAt{13}{348.386}{403.576}{Y}%
\TextAt{17}{357.779}{407.915}{2}%
\TextAt{16}{365.168}{403.576}{\Character{0}}%
\TextAt{13}{377.123}{403.576}{X}%
\TextAt{17}{387.778}{407.915}{2}%
\TextAt{16}{395.167}{403.576}{\Character{0}}%
\TextAt{13}{407.122}{403.576}{X}%
\TextAt{17}{417.777}{407.915}{3}%
\TextAt{9}{422.51}{403.576}{\Character{41}}%
\TextAt{8}{79.226}{377.095}{are}%
\TextAt{8}{98.418}{377.095}{one-dimensional}%
\TextAt{8}{184.009}{377.095}{noetherian}%
\TextAt{8}{241.899}{377.095}{rings.}%
\TextAt{8}{275.024}{377.095}{Which}%
\TextAt{8}{311.766}{377.095}{ones}%
\TextAt{8}{337.521}{377.095}{are}%
\TextAt{8}{356.713}{377.095}{integral}%
\TextAt{8}{399.002}{377.095}{domains?}%
\TextAt{8}{451.961}{377.095}{Determine}%
\TextAt{8}{508.704}{377.095}{their}%
\TextAt{8}{79.226}{360.049}{normalizations.}%
\ErrataAt{97.1}{340.02}{437.9}{\textbf{Hint:} For instance in the last example, put $t=Y/X$ and show that the homomor-}%
\TextAt{8}{79.226}{322.968}{phism}%
\TextAt{18}{113.753}{322.968}{C}%
\TextAt{9}{122.387}{322.968}{\Character{91}}%
\TextAt{13}{125.639}{322.968}{X}%
\TextAt{13}{135.634}{322.968}{\Character{59}}%
\TextAt{13}{140.858}{322.968}{Y}%
\TextAt{9}{150.281}{322.968}{\Character{93}}%
\NativeMathAt{156.853}{322.968}{11.955}{\rightarrow}%
\TextAt{18}{172.129}{322.968}{C}%
\TextAt{9}{180.764}{322.968}{\Character{91}}%
\TextAt{13}{184.015}{322.968}{t}%
\TextAt{9}{188.242}{322.968}{\Character{93}}%
\TextAt{8}{191.494}{322.968}{,}%
\TextAt{13}{198.643}{322.968}{X}%
\NativeArrowAt{212.619}{322.968}{11.955}{11.955}{\mapsto}%
\TextAt{13}{227.895}{322.968}{t}%
\TextAt{17}{232.122}{327.306}{2}%
\TextAt{16}{239.511}{322.968}{\Character{0}}%
\TextAt{9}{251.467}{322.968}{1}%
\TextAt{8}{257.32}{322.968}{,}%
\TextAt{13}{264.469}{322.968}{Y}%
\NativeArrowAt{277.183}{322.968}{11.955}{11.955}{\mapsto}%
\TextAt{13}{292.459}{322.968}{t}%
\TextAt{9}{296.686}{322.968}{\Character{40}}%
\TextAt{13}{301.238}{322.968}{t}%
\TextAt{17}{305.465}{327.306}{2}%
\TextAt{16}{312.854}{322.968}{\Character{0}}%
\TextAt{9}{324.81}{322.968}{1\Character{41}}%
\TextAt{8}{335.215}{322.968}{,}%
\TextAt{8}{342.364}{322.968}{has}%
\TextAt{8}{363.238}{322.968}{kernel}%
\TextAt{9}{397.705}{322.968}{\Character{40}}%
\TextAt{13}{402.257}{322.968}{Y}%
\TextAt{17}{411.65}{327.306}{2}%
\TextAt{16}{419.039}{322.968}{\Character{0}}%
\TextAt{13}{430.994}{322.968}{X}%
\TextAt{17}{441.649}{327.306}{2}%
\TextAt{16}{449.038}{322.968}{\Character{0}}%
\TextAt{13}{460.993}{322.968}{X}%
\TextAt{17}{471.649}{327.306}{3}%
\TextAt{9}{476.381}{322.968}{\Character{41}}%
\TextAt{8}{480.933}{322.968}{.}%
\TextAt{12}{79.226}{290.978}{Exercise}%
\TextAt{12}{131.773}{290.978}{I.12.2.}%
\TextAt{8}{171.53}{290.978}{Let}%
\TextAt{13}{191.932}{290.978}{a}%
\TextAt{8}{201.406}{290.978}{and}%
\TextAt{13}{223.601}{290.978}{b}%
\TextAt{8}{231.908}{290.978}{be}%
\TextAt{8}{247.264}{290.978}{positive}%
\TextAt{8}{290.01}{290.978}{integers}%
\TextAt{8}{332.756}{290.978}{that}%
\TextAt{8}{357.557}{290.978}{are}%
\TextAt{8}{376.499}{290.978}{not}%
\TextAt{8}{396.746}{290.978}{perfect}%
\TextAt{8}{435.187}{290.978}{squares.}%
\TextAt{8}{480.974}{290.978}{Show}%
\TextAt{8}{511.298}{290.978}{that}%
\TextAt{8}{79.226}{273.932}{the}%
\TextAt{8}{100.048}{273.932}{fundamental}%
\TextAt{8}{168.355}{273.932}{unit}%
\TextAt{8}{193.731}{273.932}{of}%
\TextAt{8}{207.726}{273.932}{the}%
\TextAt{8}{228.548}{273.932}{order}%
\TextAt{18}{259.782}{273.932}{Z}%
\TextAt{9}{270.863}{273.932}{\Character{43}}%
\TextAt{18}{283.077}{273.932}{Z}%
\TextAt{16}{291.047}{282.543}{\Character{112}}%
\Vector{q 0 G 0 g 0.478 w 229.01 -491.218 m 235.155 -491.218 l S Q}%
\TextAt{13}{301.01}{273.932}{a}%
\TextAt{8}{311.717}{273.932}{of}%
\TextAt{8}{325.712}{273.932}{the}%
\TextAt{8}{346.534}{273.932}{field}%
\TextAt{18}{372.556}{273.932}{Q}%
\TextAt{9}{381.854}{273.932}{\Character{40}}%
\TextAt{16}{386.407}{282.543}{\Character{112}}%
\Vector{q 0 G 0 g 0.478 w 324.369 -491.218 m 330.514 -491.218 l S Q}%
\TextAt{13}{396.369}{273.932}{a}%
\TextAt{9}{402.514}{273.932}{\Character{41}}%
\TextAt{8}{411.629}{273.932}{is}%
\TextAt{8}{424.058}{273.932}{also}%
\TextAt{8}{448.203}{273.932}{the}%
\TextAt{8}{469.025}{273.932}{fundamental}%
\TextAt{8}{79.226}{256.886}{unit}%
\TextAt{8}{103.937}{256.886}{of}%
\TextAt{8}{117.267}{256.886}{the}%
\TextAt{8}{137.424}{256.886}{order}%
\TextAt{18}{167.993}{256.886}{Z}%
\TextAt{9}{178.62}{256.886}{\Character{43}}%
\TextAt{18}{190.381}{256.886}{Z}%
\TextAt{16}{198.352}{265.497}{\Character{112}}%
\Vector{q 0 G 0 g 0.478 w 136.314 -508.264 m 142.459 -508.264 l S Q}%
\TextAt{13}{208.314}{256.886}{a}%
\TextAt{9}{217.116}{256.886}{\Character{43}}%
\TextAt{18}{228.877}{256.886}{Z}%
\TextAt{16}{236.847}{266.576}{\Character{112}}%
\Vector{q 0 G 0 g 0.478 w 174.81 -507.184 m 189.086 -507.184 l S Q}%
\TextAt{16}{246.81}{256.886}{\Character{0}}%
\TextAt{13}{256.108}{256.886}{b}%
\TextAt{9}{263.742}{256.886}{\Character{43}}%
\TextAt{18}{275.504}{256.886}{Z}%
\TextAt{16}{283.474}{265.497}{\Character{112}}%
\Vector{q 0 G 0 g 0.478 w 221.436 -508.264 m 227.581 -508.264 l S Q}%
\TextAt{13}{293.436}{256.886}{a}%
\TextAt{16}{299.581}{266.576}{\Character{112}}%
\Vector{q 0 G 0 g 0.478 w 237.544 -507.184 m 251.82 -507.184 l S Q}%
\TextAt{16}{309.544}{256.886}{\Character{0}}%
\TextAt{13}{318.842}{256.886}{b}%
\TextAt{8}{327.717}{256.886}{in}%
\TextAt{8}{341.37}{256.886}{the}%
\TextAt{8}{361.526}{256.886}{field}%
\TextAt{18}{386.883}{256.886}{Q}%
\TextAt{9}{396.182}{256.886}{\Character{40}}%
\TextAt{16}{400.734}{265.497}{\Character{112}}%
\Vector{q 0 G 0 g 0.478 w 338.697 -508.264 m 344.842 -508.264 l S Q}%
\TextAt{13}{410.697}{256.886}{a\Character{59}}%
\TextAt{16}{422.086}{266.576}{\Character{112}}%
\Vector{q 0 G 0 g 0.478 w 360.048 -507.184 m 374.324 -507.184 l S Q}%
\TextAt{16}{432.048}{256.886}{\Character{0}}%
\TextAt{13}{441.347}{256.886}{b}%
\TextAt{9}{446.324}{256.886}{\Character{41}}%
\TextAt{8}{450.876}{256.886}{.}%
\TextAt{12}{79.226}{224.896}{Exercise}%
\TextAt{12}{132.444}{224.896}{I.12.3.}%
\TextAt{8}{172.784}{224.896}{Let}%
\TextAt{13}{193.768}{224.896}{K}%
\TextAt{8}{208.498}{224.896}{be}%
\TextAt{8}{224.437}{224.896}{a}%
\TextAt{8}{234.208}{224.896}{number}%
\TextAt{8}{276.819}{224.896}{field}%
\TextAt{8}{302.191}{224.896}{of}%
\TextAt{8}{315.536}{224.896}{degree}%
\TextAt{13}{351.966}{224.896}{n}%
\TextAt{9}{362.3}{224.896}{\Character{61}}%
\TextAt{9}{374.733}{224.896}{\Character{91}}%
\TextAt{13}{378.003}{224.896}{K}%
\TextAt{9}{392.167}{224.896}{\Character{58}}%
\TextAt{18}{398.765}{224.896}{Q}%
\TextAt{9}{408.064}{224.896}{\Character{93}}%
\TextAt{8}{411.315}{224.896}{.}%
\TextAt{8}{419.809}{224.896}{A}%
\TextAt{12}{432.496}{224.896}{complete}%
\TextAt{12}{489.635}{224.896}{module}%
\TextAt{8}{79.226}{207.85}{in}%
\TextAt{13}{92.879}{207.85}{K}%
\TextAt{8}{107.594}{207.85}{is}%
\TextAt{8}{119.358}{207.85}{a}%
\TextAt{8}{129.113}{207.85}{subgroup}%
\TextAt{8}{179.911}{207.85}{of}%
\TextAt{8}{193.241}{207.85}{the}%
\TextAt{8}{213.397}{207.85}{form}%
\TextAt{13}{248.22}{178.849}{M}%
\TextAt{9}{264.115}{178.849}{\Character{61}}%
\TextAt{18}{276.54}{178.849}{Z}%
\TextAt{13}{284.51}{178.849}{\Character{11}}%
\TextAt{17}{291.956}{177.056}{1}%
\TextAt{9}{299.345}{178.849}{\Character{43}}%
\TextAt{16}{311.106}{178.849}{\Character{1}}%
\TextAt{16}{316.402}{178.849}{\Character{1}}%
\TextAt{16}{321.71}{178.849}{\Character{1}}%
\TextAt{9}{327.71}{178.849}{\Character{43}}%
\TextAt{18}{339.472}{178.849}{Z}%
\TextAt{13}{347.442}{178.849}{\Character{11}}%
\TextAt{21}{354.888}{177.056}{n}%
\TextAt{13}{360.524}{178.849}{\Character{59}}%
\TextAt{8}{79.226}{149.848}{where}%
\TextAt{13}{112.057}{149.848}{\Character{11}}%
\TextAt{17}{119.502}{148.055}{1}%
\TextAt{13}{124.235}{149.848}{\Character{59}}%
\TextAt{13}{129.459}{149.848}{\Character{58}}%
\TextAt{13}{134.696}{149.848}{\Character{58}}%
\TextAt{13}{139.92}{149.848}{\Character{58}}%
\TextAt{13}{145.157}{149.848}{\Character{59}}%
\TextAt{13}{150.393}{149.848}{\Character{11}}%
\TextAt{21}{157.901}{148.055}{n}%
\TextAt{8}{166.456}{149.848}{are}%
\TextAt{8}{184.988}{149.848}{linearly}%
\TextAt{8}{225.96}{149.848}{independent}%
\TextAt{8}{291.309}{149.848}{elements}%
\TextAt{8}{338.187}{149.848}{of}%
\TextAt{13}{350.538}{149.848}{K}%
\TextAt{8}{361.356}{149.848}{.}%
\TextAt{8}{369.478}{149.848}{Show}%
\TextAt{8}{399.391}{149.848}{that}%
\TextAt{8}{423.781}{149.848}{the}%
\TextAt{8}{442.959}{149.848}{ring}%
\TextAt{8}{466.046}{149.848}{of}%
\TextAt{8}{478.397}{149.848}{multipliers}%
\TextAt{14}{243.464}{120.847}{o}%
\TextAt{9}{252.629}{120.847}{\Character{61}}%
\TextAt{16}{265.054}{120.847}{\Character{102}}%
\TextAt{13}{271.032}{120.847}{\Character{11}}%
\NativeMathAt{281.874}{120.847}{11.955}{\in}%
\TextAt{13}{293.165}{120.847}{K}%
\TextAt{16}{307.304}{120.847}{\Character{106}}%
\TextAt{13}{313.945}{120.847}{\Character{11}M}%
\NativeMathAt{337.362}{120.847}{11.955}{\subseteq}%
\TextAt{13}{349.981}{120.847}{M}%
\TextAt{16}{362.555}{120.847}{\Character{103}}%
\TextAt{8}{79.226}{91.846}{is}%
\TextAt{8}{90.99}{91.846}{an}%
\TextAt{8}{107.249}{91.846}{order}%
\TextAt{8}{137.818}{91.846}{in}%
\TextAt{13}{151.471}{91.846}{K}%
\TextAt{8}{162.289}{91.846}{,}%
\TextAt{8}{169.438}{91.846}{but}%
\TextAt{8}{190.898}{91.846}{in}%
\TextAt{8}{204.55}{91.846}{general}%
\TextAt{8}{244.875}{91.846}{not}%
\TextAt{8}{265.689}{91.846}{the}%
\TextAt{8}{285.846}{91.846}{maximal}%
\TextAt{8}{333.654}{91.846}{order.}%
\TextAt{8}{300.14}{43.033}{78}%
\Destination{78.23}{809.815}{nav.page.78}%
\Destination{79.226}{459.755}{nav.ex.1.12.1}%
\Destination{79.226}{304.955}{nav.ex.1.12.2}%
\Destination{79.226}{238.855}{nav.ex.1.12.3}%
\Destination{79.36}{550.755}{nav.print.84}%
\Destination{79.36}{549.755}{nav.number.1.12.14}%
\Hotlink{401.976}{758.955}{30.936}{11.047}{nav.number.1.12.4}%
\Destination{79.226}{357.02}{errata.52}%
}

% Page 91 (body page 79)
\Page{%
\TextAt{2}{79.226}{796.943}{I.}%
\TextAt{2}{89.83}{796.943}{Algebraic}%
\TextAt{2}{138.648}{796.943}{Integers}%
\Vector{q 0 G 0 g 0.299 w 7.226 18.729 m 460.769 18.729 l S Q}%
\TextAt{12}{79.226}{761.124}{Exercise}%
\TextAt{12}{132.791}{761.124}{I.12.4.}%
\TextAt{8}{173.432}{761.124}{Determine}%
\TextAt{8}{230.81}{761.124}{the}%
\TextAt{8}{251.283}{761.124}{ring}%
\TextAt{8}{275.665}{761.124}{of}%
\TextAt{8}{289.312}{761.124}{multipliers}%
\TextAt{14}{347.898}{761.124}{o}%
\TextAt{8}{357.956}{761.124}{of}%
\TextAt{8}{371.602}{761.124}{the}%
\TextAt{8}{392.075}{761.124}{complete}%
\TextAt{8}{441.814}{761.124}{module}%
\TextAt{13}{483.424}{761.124}{M}%
\TextAt{9}{499.858}{761.124}{\Character{61}}%
\TextAt{18}{512.822}{761.124}{Z}%
\TextAt{9}{523.665}{761.124}{\Character{43}}%
\TextAt{18}{79.226}{744.078}{Z}%
\TextAt{16}{87.196}{753.967}{\Character{112}}%
\Vector{q 0 G 0 g 0.478 w 25.159 -19.794 m 31.012 -19.794 l S Q}%
\TextAt{9}{97.159}{744.078}{2}%
\TextAt{8}{107.017}{744.078}{in}%
\TextAt{18}{120.778}{744.078}{Q}%
\TextAt{9}{130.076}{744.078}{\Character{40}}%
\TextAt{16}{134.628}{753.967}{\Character{112}}%
\Vector{q 0 G 0 g 0.478 w 72.591 -19.794 m 78.444 -19.794 l S Q}%
\TextAt{9}{144.591}{744.078}{2\Character{41}}%
\TextAt{8}{154.996}{744.078}{.}%
\TextAt{8}{163.768}{744.078}{Show}%
\TextAt{8}{194.768}{744.078}{that}%
\TextAt{13}{220.244}{744.078}{\Character{34}}%
\TextAt{9}{229.227}{744.078}{\Character{61}}%
\TextAt{9}{241.828}{744.078}{1}%
\TextAt{9}{250.4}{744.078}{\Character{43}}%
\TextAt{16}{262.254}{753.967}{\Character{112}}%
\Vector{q 0 G 0 g 0.478 w 200.216 -19.794 m 206.069 -19.794 l S Q}%
\TextAt{9}{272.216}{744.078}{2}%
\TextAt{8}{282.074}{744.078}{is}%
\TextAt{8}{293.946}{744.078}{a}%
\TextAt{8}{303.809}{744.078}{fundamental}%
\TextAt{8}{371.559}{744.078}{unit}%
\TextAt{8}{396.378}{744.078}{of}%
\TextAt{14}{409.816}{744.078}{o}%
\TextAt{8}{415.66}{744.078}{.}%
\TextAt{8}{424.432}{744.078}{Determine}%
\TextAt{8}{481.601}{744.078}{all}%
\TextAt{8}{497.968}{744.078}{integer}%
\TextAt{8}{79.226}{727.032}{solutions}%
\TextAt{8}{128.135}{727.032}{of}%
\TextAt{8}{141.465}{727.032}{“Pell’s}%
\TextAt{8}{178.095}{727.032}{equation”}%
\TextAt{13}{272.213}{709.986}{x}%
\TextAt{17}{278.865}{714.922}{2}%
\TextAt{16}{286.254}{709.986}{\Character{0}}%
\TextAt{9}{298.209}{709.986}{2}%
\TextAt{13}{304.062}{709.986}{y}%
\TextAt{17}{310.199}{714.922}{2}%
\TextAt{9}{318.252}{709.986}{\Character{61}}%
\TextAt{9}{330.661}{709.986}{7}%
\TextAt{13}{336.531}{709.986}{\Character{58}}%
\TextAt{12}{97.159}{683.883}{Hint:}%
\TextAt{13}{132.698}{683.883}{N}%
\TextAt{9}{143.32}{683.883}{\Character{40}}%
\TextAt{13}{147.873}{683.883}{x}%
\TextAt{9}{157.181}{683.883}{\Character{43}}%
\TextAt{13}{168.943}{683.883}{y}%
\TextAt{16}{175.079}{693.773}{\Character{112}}%
\Vector{q 0 G 0 g 0.478 w 113.042 -79.988 m 118.895 -79.988 l S Q}%
\TextAt{9}{185.042}{683.883}{2\Character{41}}%
\TextAt{9}{198.743}{683.883}{\Character{61}}%
\TextAt{13}{211.194}{683.883}{x}%
\TextAt{17}{217.846}{688.222}{2}%
\TextAt{16}{225.235}{683.883}{\Character{0}}%
\TextAt{9}{237.19}{683.883}{2}%
\TextAt{13}{243.043}{683.883}{y}%
\TextAt{17}{249.179}{688.222}{2}%
\TextAt{8}{253.912}{683.883}{,}%
\TextAt{13}{261.061}{683.883}{N}%
\TextAt{9}{271.684}{683.883}{\Character{40}3}%
\TextAt{9}{284.727}{683.883}{\Character{43}}%
\TextAt{16}{296.507}{693.773}{\Character{112}}%
\Vector{q 0 G 0 g 0.478 w 234.47 -79.988 m 240.323 -79.988 l S Q}%
\TextAt{9}{306.47}{683.883}{2\Character{41}}%
\TextAt{9}{320.171}{683.883}{\Character{61}}%
\TextAt{13}{332.621}{683.883}{N}%
\TextAt{9}{343.244}{683.883}{\Character{40}5}%
\TextAt{9}{356.287}{683.883}{\Character{43}}%
\TextAt{9}{368.039}{683.883}{3}%
\TextAt{16}{373.92}{693.773}{\Character{112}}%
\Vector{q 0 G 0 g 0.478 w 311.883 -79.988 m 317.736 -79.988 l S Q}%
\TextAt{9}{383.883}{683.883}{2\Character{41}}%
\TextAt{9}{397.584}{683.883}{\Character{61}}%
\TextAt{9}{410.005}{683.883}{7}%
\TextAt{8}{415.887}{683.883}{.}%
\TextAt{12}{79.226}{652.438}{Exercise}%
\TextAt{12}{134.487}{652.438}{I.12.5.}%
\TextAt{8}{176.604}{652.438}{In}%
\TextAt{8}{193.016}{652.438}{a}%
\TextAt{8}{204.563}{652.438}{one-dimensional}%
\TextAt{8}{292.264}{652.438}{noetherian}%
\TextAt{8}{352.266}{652.438}{integral}%
\TextAt{8}{396.665}{652.438}{domain}%
\TextAt{8}{440.085}{652.438}{the}%
\TextAt{8}{462.033}{652.438}{regular}%
\TextAt{8}{503.503}{652.438}{prime}%
\TextAt{8}{79.226}{635.392}{ideals}%
\NativeMathAt{111.804}{635.392}{11.955}{\neq}%
\TextAt{9}{124.214}{635.392}{0}%
\TextAt{8}{133.98}{635.392}{are}%
\TextAt{8}{153.491}{635.392}{precisely}%
\TextAt{8}{201.347}{635.392}{the}%
\TextAt{8}{221.504}{635.392}{invertible}%
\TextAt{8}{273.21}{635.392}{prime}%
\TextAt{8}{306.373}{635.392}{ideals.}%
\TextAt{15}{79.226}{582.941}{§}%
\TextAt{15}{93.785}{582.941}{13}%
\TextAt{15}{130.385}{582.941}{One-dimensional}%
\TextAt{15}{277.612}{582.941}{Schemes}%
\TextAt{8}{79.226}{551.055}{The}%
\TextAt{8}{103.21}{551.055}{first}%
\TextAt{8}{127.265}{551.055}{approach}%
\TextAt{8}{177.61}{551.055}{to}%
\TextAt{8}{191.851}{551.055}{the}%
\TextAt{8}{211.937}{551.055}{theory}%
\TextAt{8}{248.617}{551.055}{of}%
\TextAt{8}{261.877}{551.055}{algebraic}%
\TextAt{8}{311.242}{551.055}{number}%
\TextAt{8}{353.768}{551.055}{fields}%
\TextAt{8}{383.67}{551.055}{is}%
\TextAt{8}{395.364}{551.055}{dominated}%
\TextAt{8}{453.18}{551.055}{by}%
\TextAt{8}{469.369}{551.055}{the}%
\TextAt{8}{489.456}{551.055}{methods}%
\TextAt{8}{79.226}{534.009}{of}%
\TextAt{8}{93.857}{534.009}{arithmetic}%
\TextAt{8}{151.742}{534.009}{and}%
\TextAt{8}{175.806}{534.009}{algebra.}%
\TextAt{8}{225.242}{534.009}{But}%
\TextAt{8}{249.784}{534.009}{the}%
\TextAt{8}{271.241}{534.009}{theory}%
\TextAt{8}{309.292}{534.009}{may}%
\TextAt{8}{335.962}{534.009}{also}%
\TextAt{8}{360.743}{534.009}{be}%
\TextAt{8}{377.968}{534.009}{treated}%
\TextAt{8}{419.594}{534.009}{fundamentally}%
\TextAt{8}{497.97}{534.009}{from}%
\TextAt{8}{526.911}{534.009}{a}%
\TextAt{8}{79.226}{516.963}{geometric}%
\TextAt{8}{133.538}{516.963}{point}%
\TextAt{8}{165.088}{516.963}{of}%
\TextAt{8}{179.398}{516.963}{view,}%
\TextAt{8}{210.858}{516.963}{which}%
\TextAt{8}{245.325}{516.963}{will}%
\TextAt{8}{268.41}{516.963}{bring}%
\TextAt{8}{299.96}{516.963}{out}%
\TextAt{8}{321.754}{516.963}{novel}%
\TextAt{8}{352.981}{516.963}{aspects}%
\TextAt{8}{394.728}{516.963}{in}%
\TextAt{8}{409.361}{516.963}{a}%
\TextAt{8}{420.097}{516.963}{variety}%
\TextAt{8}{459.788}{516.963}{of}%
\TextAt{8}{474.098}{516.963}{ways.}%
\TextAt{8}{509.947}{516.963}{This}%
\TextAt{8}{79.226}{499.917}{geometric}%
\TextAt{8}{133.197}{499.917}{interpretation}%
\TextAt{8}{208.316}{499.917}{hinges}%
\TextAt{8}{244.784}{499.917}{on}%
\TextAt{8}{261.682}{499.917}{the}%
\TextAt{8}{282.477}{499.917}{possibility}%
\TextAt{8}{338.851}{499.917}{of}%
\TextAt{8}{352.819}{499.917}{viewing}%
\TextAt{8}{396.054}{499.917}{numbers}%
\TextAt{8}{443.904}{499.917}{as}%
\TextAt{8}{458.912}{499.917}{functions}%
\TextAt{8}{510.014}{499.917}{on}%
\TextAt{8}{526.911}{499.917}{a}%
\TextAt{8}{79.226}{482.871}{topological}%
\TextAt{8}{138.751}{482.871}{space.}%
\TextAt{8}{97.159}{462.836}{In}%
\TextAt{8}{111.78}{462.836}{order}%
\TextAt{8}{142.349}{462.836}{to}%
\TextAt{8}{156.66}{462.836}{explain}%
\TextAt{8}{197.307}{462.836}{this,}%
\TextAt{8}{223.382}{462.836}{let}%
\TextAt{8}{240.286}{462.836}{us}%
\TextAt{8}{255.302}{462.836}{start}%
\TextAt{8}{283.337}{462.836}{from}%
\TextAt{8}{310.977}{462.836}{polynomials}%
\TextAt{13}{247.858}{434.788}{f}%
\TextAt{9}{254.904}{434.788}{\Character{40}}%
\TextAt{13}{259.457}{434.788}{x}%
\TextAt{9}{266.109}{434.788}{\Character{41}}%
\TextAt{9}{273.964}{434.788}{\Character{61}}%
\TextAt{13}{286.407}{434.788}{a}%
\TextAt{21}{292.552}{432.994}{n}%
\TextAt{13}{298.189}{434.788}{x}%
\TextAt{21}{304.841}{439.724}{n}%
\TextAt{9}{313.134}{434.788}{\Character{43}}%
\TextAt{16}{324.895}{434.788}{\Character{1}}%
\TextAt{16}{330.191}{434.788}{\Character{1}}%
\TextAt{16}{335.499}{434.788}{\Character{1}}%
\TextAt{9}{341.499}{434.788}{\Character{43}}%
\TextAt{13}{353.261}{434.788}{a}%
\TextAt{17}{359.406}{432.994}{0}%
\TextAt{8}{79.226}{406.739}{with}%
\TextAt{8}{106.096}{406.739}{complex}%
\TextAt{8}{152.155}{406.739}{coefficients}%
\TextAt{13}{211.603}{406.739}{a}%
\TextAt{21}{217.748}{404.945}{i}%
\NativeMathAt{224.808}{406.739}{11.955}{\in}%
\TextAt{18}{236.457}{406.739}{C}%
\TextAt{8}{245.092}{406.739}{,}%
\TextAt{8}{252.503}{406.739}{which}%
\TextAt{8}{286.2}{406.739}{may}%
\TextAt{8}{311.779}{406.739}{be}%
\TextAt{8}{327.914}{406.739}{immediately}%
\TextAt{8}{394.786}{406.739}{interpreted}%
\TextAt{8}{455.155}{406.739}{as}%
\TextAt{8}{469.735}{406.739}{functions}%
\TextAt{8}{520.408}{406.739}{on}%
\TextAt{8}{79.226}{389.693}{the}%
\TextAt{8}{98.777}{389.693}{complex}%
\TextAt{8}{144.019}{389.693}{plane.}%
\TextAt{8}{179.583}{389.693}{This}%
\TextAt{8}{205.697}{389.693}{property}%
\TextAt{8}{252.9}{389.693}{may}%
\TextAt{8}{277.663}{389.693}{be}%
\TextAt{8}{292.981}{389.693}{formulated}%
\TextAt{8}{351.566}{389.693}{in}%
\TextAt{8}{364.613}{389.693}{a}%
\TextAt{8}{373.762}{389.693}{purely}%
\TextAt{8}{409.249}{389.693}{algebraic}%
\TextAt{8}{458.078}{389.693}{way}%
\TextAt{8}{481.215}{389.693}{as}%
\TextAt{8}{494.979}{389.693}{follows.}%
\TextAt{8}{79.226}{372.647}{Let}%
\TextAt{13}{99.469}{372.647}{a}%
\NativeMathAt{108.935}{372.647}{11.955}{\in}%
\TextAt{18}{120.226}{372.647}{C}%
\TextAt{8}{132.031}{372.647}{be}%
\TextAt{8}{147.229}{372.647}{a}%
\TextAt{8}{156.258}{372.647}{point}%
\TextAt{8}{186.101}{372.647}{in}%
\TextAt{8}{199.027}{372.647}{the}%
\TextAt{8}{218.457}{372.647}{complex}%
\TextAt{8}{263.579}{372.647}{plane.}%
\TextAt{8}{299.102}{372.647}{The}%
\TextAt{8}{322.43}{372.647}{set}%
\TextAt{8}{339.971}{372.647}{of}%
\TextAt{8}{352.574}{372.647}{all}%
\TextAt{8}{368.107}{372.647}{functions}%
\TextAt{13}{417.843}{372.647}{f}%
\TextAt{9}{424.89}{372.647}{\Character{40}}%
\TextAt{13}{429.442}{372.647}{x}%
\TextAt{9}{436.094}{372.647}{\Character{41}}%
\TextAt{8}{443.817}{372.647}{in}%
\TextAt{8}{456.744}{372.647}{the}%
\TextAt{8}{476.174}{372.647}{polynomial}%
\TextAt{8}{79.226}{355.601}{ring}%
\TextAt{18}{102.861}{355.601}{C}%
\TextAt{9}{111.495}{355.601}{\Character{91}}%
\TextAt{13}{114.747}{355.601}{x}%
\TextAt{9}{121.399}{355.601}{\Character{93}}%
\TextAt{8}{128.116}{355.601}{which}%
\TextAt{8}{161.172}{355.601}{vanish}%
\TextAt{8}{196.905}{355.601}{at}%
\TextAt{8}{210.784}{355.601}{the}%
\TextAt{8}{230.509}{355.601}{point}%
\TextAt{13}{260.647}{355.601}{a}%
\TextAt{8}{270.258}{355.601}{forms}%
\TextAt{8}{302.082}{355.601}{the}%
\TextAt{8}{321.807}{355.601}{maximal}%
\TextAt{8}{369.184}{355.601}{ideal}%
\TextAt{14}{396.716}{355.601}{p}%
\TextAt{9}{406.02}{355.601}{\Character{61}}%
\TextAt{9}{418.43}{355.601}{\Character{40}}%
\TextAt{13}{422.998}{355.601}{x}%
\TextAt{16}{431.425}{355.601}{\Character{0}}%
\TextAt{13}{442.498}{355.601}{a}%
\TextAt{9}{448.643}{355.601}{\Character{41}}%
\TextAt{8}{456.662}{355.601}{of}%
\TextAt{18}{469.56}{355.601}{C}%
\TextAt{9}{478.195}{355.601}{\Character{91}}%
\TextAt{13}{481.446}{355.601}{x}%
\TextAt{9}{488.098}{355.601}{\Character{93}}%
\TextAt{8}{491.35}{355.601}{.}%
\TextAt{8}{499.655}{355.601}{In}%
\TextAt{8}{513.844}{355.601}{this}%
\TextAt{8}{79.226}{338.555}{way}%
\TextAt{8}{103.451}{338.555}{the}%
\TextAt{8}{124.09}{338.555}{points}%
\TextAt{8}{159.756}{338.555}{of}%
\TextAt{8}{173.569}{338.555}{the}%
\TextAt{8}{194.207}{338.555}{complex}%
\TextAt{8}{240.537}{338.555}{plane}%
\TextAt{8}{272.235}{338.555}{correspond}%
\TextAt{8}{332.289}{338.555}{1–1}%
\TextAt{8}{354.243}{338.555}{to}%
\TextAt{8}{369.036}{338.555}{the}%
\TextAt{8}{389.674}{338.555}{maximal}%
\TextAt{8}{437.965}{338.555}{ideals}%
\TextAt{8}{471.025}{338.555}{of}%
\TextAt{18}{484.838}{338.555}{C}%
\TextAt{9}{493.472}{338.555}{\Character{91}}%
\TextAt{13}{496.724}{338.555}{x}%
\TextAt{9}{503.376}{338.555}{\Character{93}}%
\TextAt{8}{506.627}{338.555}{.}%
\TextAt{8}{516.522}{338.555}{We}%
\TextAt{8}{79.226}{321.509}{denote}%
\TextAt{8}{116.945}{321.509}{the}%
\TextAt{8}{137.101}{321.509}{set}%
\TextAt{8}{155.369}{321.509}{of}%
\TextAt{8}{168.699}{321.509}{all}%
\TextAt{8}{184.958}{321.509}{these}%
\TextAt{8}{214.929}{321.509}{maximal}%
\TextAt{8}{262.738}{321.509}{ideals}%
\TextAt{8}{295.316}{321.509}{by}%
\TextAt{13}{263.384}{293.46}{M}%
\TextAt{9}{279.278}{293.46}{\Character{61}}%
\TextAt{8}{291.703}{293.46}{Max}%
\TextAt{9}{314.466}{293.46}{\Character{40}}%
\TextAt{18}{319.018}{293.46}{C}%
\TextAt{9}{327.653}{293.46}{\Character{91}}%
\TextAt{13}{330.904}{293.46}{x}%
\TextAt{9}{337.556}{293.46}{\Character{93}\Character{41}}%
\TextAt{13}{345.36}{293.46}{\Character{58}}%
\TextAt{8}{79.226}{265.411}{We}%
\TextAt{8}{99.398}{265.411}{may}%
\TextAt{8}{124.795}{265.411}{view}%
\TextAt{13}{151.806}{265.411}{M}%
\TextAt{8}{168.305}{265.411}{as}%
\TextAt{8}{182.702}{265.411}{a}%
\TextAt{8}{192.486}{265.411}{new}%
\TextAt{8}{216.567}{265.411}{kind}%
\TextAt{8}{242.932}{265.411}{of}%
\TextAt{8}{256.29}{265.411}{space}%
\TextAt{8}{287.593}{265.411}{and}%
\TextAt{8}{310.383}{265.411}{may}%
\TextAt{8}{335.78}{265.411}{interpret}%
\TextAt{8}{384.262}{265.411}{the}%
\TextAt{8}{404.446}{265.411}{elements}%
\TextAt{13}{452.331}{265.411}{f}%
\TextAt{9}{459.377}{265.411}{\Character{40}}%
\TextAt{13}{463.929}{265.411}{x}%
\TextAt{9}{470.581}{265.411}{\Character{41}}%
\TextAt{8}{479.059}{265.411}{of}%
\TextAt{8}{492.417}{265.411}{the}%
\TextAt{8}{512.601}{265.411}{ring}%
\TextAt{18}{79.226}{248.365}{C}%
\TextAt{9}{87.86}{248.365}{\Character{91}}%
\TextAt{13}{91.112}{248.365}{x}%
\TextAt{9}{97.764}{248.365}{\Character{93}}%
\TextAt{8}{104.228}{248.365}{as}%
\TextAt{8}{117.914}{248.365}{functions}%
\TextAt{8}{167.692}{248.365}{on}%
\TextAt{13}{183.266}{248.365}{M}%
\TextAt{8}{199.052}{248.365}{as}%
\TextAt{8}{212.738}{248.365}{follows.}%
\TextAt{8}{255.496}{248.365}{For}%
\TextAt{8}{275.781}{248.365}{every}%
\TextAt{8}{305.988}{248.365}{point}%
\TextAt{14}{335.873}{248.365}{p}%
\TextAt{9}{345.177}{248.365}{\Character{61}}%
\TextAt{9}{357.587}{248.365}{\Character{40}}%
\TextAt{13}{362.155}{248.365}{x}%
\TextAt{16}{370.063}{248.365}{\Character{0}}%
\TextAt{13}{380.618}{248.365}{a}%
\TextAt{9}{386.763}{248.365}{\Character{41}}%
\TextAt{8}{394.528}{248.365}{of}%
\TextAt{13}{407.173}{248.365}{M}%
\TextAt{8}{422.959}{248.365}{we}%
\TextAt{8}{439.502}{248.365}{have}%
\TextAt{8}{465.812}{248.365}{the}%
\TextAt{8}{485.284}{248.365}{canonical}%
\TextAt{8}{79.226}{231.319}{isomorphism}%
\TextAt{18}{270.291}{214.273}{C}%
\TextAt{9}{278.925}{214.273}{\Character{91}}%
\TextAt{13}{282.176}{214.273}{x}%
\TextAt{9}{288.829}{214.273}{\Character{93}\Character{47}}%
\TextAt{14}{297.933}{214.273}{p}%
\NativeMathAt{313.574}{221.05}{7.97}{\sim}%
\NativeArrowAt{307.237}{214.273}{11.955}{19.26}{\longrightarrow}%
\TextAt{18}{329.819}{214.273}{C}%
\TextAt{13}{338.453}{214.273}{\Character{59}}%
\TextAt{8}{79.226}{191.16}{which}%
\TextAt{8}{112.713}{191.16}{sends}%
\TextAt{8}{144.047}{191.16}{the}%
\TextAt{8}{164.204}{191.16}{residue}%
\TextAt{8}{203.931}{191.16}{class}%
\TextAt{13}{231.368}{191.16}{f}%
\TextAt{9}{238.414}{191.16}{\Character{40}}%
\TextAt{13}{242.966}{191.16}{x}%
\TextAt{9}{249.618}{191.16}{\Character{41}}%
\TextAt{8}{257.492}{191.16}{mod}%
\TextAt{14}{282.93}{191.16}{p}%
\TextAt{8}{292.81}{191.16}{to}%
\TextAt{13}{307.12}{191.16}{f}%
\TextAt{9}{314.167}{191.16}{\Character{40}}%
\TextAt{13}{318.719}{191.16}{a}%
\TextAt{9}{324.864}{191.16}{\Character{41}}%
\TextAt{8}{329.416}{191.16}{.}%
\TextAt{8}{337.865}{191.16}{We}%
\TextAt{8}{358.009}{191.16}{may}%
\TextAt{8}{383.378}{191.16}{thus}%
\TextAt{8}{409.129}{191.16}{view}%
\TextAt{8}{436.112}{191.16}{this}%
\TextAt{8}{458.935}{191.16}{residue}%
\TextAt{8}{498.662}{191.16}{class}%
\TextAt{13}{238.44}{163.111}{f}%
\TextAt{9}{245.487}{163.111}{\Character{40}}%
\TextAt{14}{250.039}{163.111}{p}%
\TextAt{9}{256.022}{163.111}{\Character{41}}%
\TextAt{9}{263.877}{163.111}{\Character{58}\Character{61}}%
\TextAt{13}{279.572}{163.111}{f}%
\TextAt{9}{286.619}{163.111}{\Character{40}}%
\TextAt{13}{291.171}{163.111}{x}%
\TextAt{9}{297.823}{163.111}{\Character{41}}%
\TextAt{8}{305.696}{163.111}{mod}%
\TextAt{14}{331.134}{163.111}{p}%
\NativeMathAt{340.438}{163.111}{11.955}{\in}%
\TextAt{13}{351.729}{163.111}{\Character{20}}%
\TextAt{9}{358.468}{163.111}{\Character{40}}%
\TextAt{14}{363.02}{163.111}{p}%
\TextAt{9}{369.003}{163.111}{\Character{41}}%
\TextAt{8}{79.226}{135.062}{in}%
\TextAt{8}{93.938}{135.062}{the}%
\TextAt{8}{115.153}{135.062}{residue}%
\TextAt{8}{155.938}{135.062}{class}%
\TextAt{8}{184.434}{135.062}{field}%
\TextAt{13}{210.85}{135.062}{\Character{20}}%
\TextAt{9}{217.588}{135.062}{\Character{40}}%
\TextAt{14}{222.141}{135.062}{p}%
\TextAt{9}{228.124}{135.062}{\Character{41}}%
\TextAt{9}{237.784}{135.062}{\Character{61}}%
\TextAt{18}{252.031}{135.062}{C}%
\TextAt{9}{260.665}{135.062}{\Character{91}}%
\TextAt{13}{263.917}{135.062}{x}%
\TextAt{9}{270.569}{135.062}{\Character{93}\Character{47}}%
\TextAt{14}{279.674}{135.062}{p}%
\TextAt{8}{290.613}{135.062}{as}%
\TextAt{8}{306.042}{135.062}{the}%
\TextAt{8}{327.257}{135.062}{“value”}%
\TextAt{8}{369.608}{135.062}{of}%
\TextAt{13}{383.997}{135.062}{f}%
\TextAt{8}{396}{135.062}{at}%
\TextAt{8}{411.369}{135.062}{the}%
\TextAt{8}{432.584}{135.062}{point}%
\TextAt{14}{464.212}{135.062}{p}%
\NativeMathAt{475.32}{135.062}{11.955}{\in}%
\TextAt{13}{488.415}{135.062}{M}%
\TextAt{8}{500.989}{135.062}{.}%
\TextAt{8}{512.613}{135.062}{The}%
\TextAt{8}{79.226}{118.016}{topology}%
\TextAt{8}{127.444}{118.016}{on}%
\TextAt{18}{143.778}{118.016}{C}%
\TextAt{8}{156.385}{118.016}{cannot}%
\TextAt{8}{194.836}{118.016}{be}%
\TextAt{8}{210.835}{118.016}{transferred}%
\TextAt{8}{270.482}{118.016}{to}%
\TextAt{13}{284.867}{118.016}{M}%
\TextAt{8}{301.413}{118.016}{by}%
\TextAt{8}{317.747}{118.016}{algebraic}%
\TextAt{8}{367.256}{118.016}{means.}%
\TextAt{8}{407.861}{118.016}{All}%
\TextAt{8}{427.112}{118.016}{that}%
\TextAt{8}{452.556}{118.016}{can}%
\TextAt{8}{474.09}{118.016}{be}%
\TextAt{8}{490.09}{118.016}{salvaged}%
\TextAt{8}{79.226}{100.97}{algebraically}%
\TextAt{8}{147.203}{100.97}{are}%
\TextAt{8}{166.714}{100.97}{the}%
\TextAt{8}{186.87}{100.97}{point}%
\TextAt{8}{217.44}{100.97}{sets}%
\TextAt{8}{240.322}{100.97}{defined}%
\TextAt{8}{280.635}{100.97}{by}%
\TextAt{8}{296.894}{100.97}{equations}%
\TextAt{8}{349.317}{100.97}{of}%
\TextAt{8}{362.647}{100.97}{the}%
\TextAt{8}{382.804}{100.97}{form}%
\TextAt{13}{283.797}{72.921}{f}%
\TextAt{9}{290.843}{72.921}{\Character{40}}%
\TextAt{13}{295.395}{72.921}{x}%
\TextAt{9}{302.047}{72.921}{\Character{41}}%
\TextAt{9}{309.902}{72.921}{\Character{61}}%
\TextAt{9}{322.323}{72.921}{0}%
\TextAt{8}{300.14}{43.033}{79}%
\Destination{78.23}{809.815}{nav.page.79}%
\Destination{79.23}{609.355}{nav.section.1.13}%
\Destination{79.226}{775.155}{nav.ex.1.12.4}%
\Destination{79.226}{666.455}{nav.ex.1.12.5}%
\Destination{79.36}{477.855}{nav.print.85}%
\Destination{79.23}{609.355}{outline.18}%
}

% Page 92 (body page 80)
\Page{%
\TextAt{2}{79.226}{796.943}{I.}%
\TextAt{2}{89.83}{796.943}{Algebraic}%
\TextAt{2}{138.648}{796.943}{Integers}%
\Vector{q 0 G 0 g 0.299 w 7.226 18.729 m 460.769 18.729 l S Q}%
\TextAt{8}{79.226}{761.124}{(i.e.,}%
\TextAt{8}{106.563}{761.124}{only}%
\TextAt{8}{132.796}{761.124}{the}%
\TextAt{8}{153.494}{761.124}{finite}%
\TextAt{8}{183.947}{761.124}{sets}%
\TextAt{8}{207.371}{761.124}{and}%
\TextAt{13}{230.675}{761.124}{M}%
\TextAt{8}{247.688}{761.124}{itself).}%
\TextAt{8}{287.087}{761.124}{These}%
\TextAt{8}{321.497}{761.124}{sets}%
\TextAt{8}{344.921}{761.124}{are}%
\TextAt{8}{364.973}{761.124}{defined}%
\TextAt{8}{405.827}{761.124}{to}%
\TextAt{8}{420.679}{761.124}{be}%
\TextAt{8}{437.145}{761.124}{the}%
\TextAt{8}{457.843}{761.124}{closed}%
\TextAt{8}{492.911}{761.124}{subsets.}%
\TextAt{8}{79.226}{744.078}{In}%
\TextAt{8}{93.847}{744.078}{the}%
\TextAt{8}{114.004}{744.078}{new}%
\TextAt{8}{138.058}{744.078}{formulation}%
\TextAt{8}{201.157}{744.078}{they}%
\TextAt{8}{227.494}{744.078}{are}%
\TextAt{8}{247.005}{744.078}{the}%
\TextAt{8}{267.161}{744.078}{sets}%
\TextAt{13}{171.969}{718.876}{V}%
\TextAt{9}{181.395}{718.876}{\Character{40}}%
\TextAt{13}{185.947}{718.876}{f}%
\TextAt{9}{192.993}{718.876}{\Character{41}}%
\TextAt{9}{200.848}{718.876}{\Character{61}}%
\TextAt{16}{213.292}{718.876}{\Character{102}}%
\TextAt{14}{219.27}{718.876}{p}%
\NativeMathAt{228.573}{718.876}{11.955}{\in}%
\TextAt{13}{239.864}{718.876}{M}%
\TextAt{16}{255.759}{718.876}{\Character{106}}%
\TextAt{13}{262.401}{718.876}{f}%
\TextAt{9}{269.447}{718.876}{\Character{40}}%
\TextAt{14}{273.999}{718.876}{p}%
\TextAt{9}{279.982}{718.876}{\Character{41}}%
\TextAt{9}{287.837}{718.876}{\Character{61}}%
\TextAt{9}{300.258}{718.876}{0}%
\TextAt{16}{306.134}{718.876}{\Character{103}}%
\TextAt{9}{315.432}{718.876}{\Character{61}}%
\TextAt{16}{327.858}{718.876}{\Character{102}}%
\TextAt{14}{333.836}{718.876}{p}%
\NativeMathAt{343.14}{718.876}{11.955}{\in}%
\TextAt{13}{354.43}{718.876}{M}%
\TextAt{16}{370.325}{718.876}{\Character{106}}%
\TextAt{14}{376.967}{718.876}{p}%
\NativeMathAt{386.271}{718.876}{11.955}{\supseteq}%
\TextAt{9}{398.89}{718.876}{\Character{40}}%
\TextAt{13}{403.442}{718.876}{f}%
\TextAt{9}{410.489}{718.876}{\Character{40}}%
\TextAt{13}{415.041}{718.876}{x}%
\TextAt{9}{421.693}{718.876}{\Character{41}\Character{41}}%
\TextAt{16}{430.798}{718.876}{\Character{103}}%
\TextAt{13}{436.775}{718.876}{\Character{58}}%
\TextAt{8}{79.226}{693.675}{The}%
\TextAt{8}{104.678}{693.675}{algebraic}%
\TextAt{8}{155.51}{693.675}{interpretation}%
\TextAt{8}{231.388}{693.675}{of}%
\TextAt{8}{246.116}{693.675}{functions}%
\TextAt{8}{297.976}{693.675}{given}%
\TextAt{8}{329.943}{693.675}{above}%
\TextAt{8}{364.517}{693.675}{leads}%
\TextAt{8}{395.24}{693.675}{to}%
\TextAt{8}{410.948}{693.675}{the}%
\TextAt{8}{432.503}{693.675}{following}%
\TextAt{8}{483.335}{693.675}{geometric}%
\TextAt{8}{79.226}{676.629}{perception}%
\TextAt{8}{136.629}{676.629}{of}%
\TextAt{8}{149.81}{676.629}{completely}%
\TextAt{8}{208.516}{676.629}{general}%
\TextAt{8}{248.691}{676.629}{rings.}%
\TextAt{8}{281.873}{676.629}{For}%
\TextAt{8}{302.693}{676.629}{an}%
\TextAt{8}{318.803}{676.629}{arbitrary}%
\TextAt{8}{368.423}{676.629}{ring}%
\TextAt{14}{392.339}{676.629}{o}%
\TextAt{8}{398.184}{676.629}{,}%
\TextAt{8}{405.213}{676.629}{one}%
\TextAt{8}{426.524}{676.629}{introduces}%
\TextAt{8}{483.018}{676.629}{the}%
\TextAt{12}{503.025}{676.629}{spec-}%
\TextAt{12}{79.226}{659.583}{trum}%
\TextAt{13}{273.619}{642.537}{X}%
\TextAt{9}{287.595}{642.537}{\Character{61}}%
\TextAt{8}{300.02}{642.537}{Spec}%
\TextAt{9}{323.428}{642.537}{\Character{40}}%
\TextAt{14}{327.981}{642.537}{o}%
\TextAt{9}{333.825}{642.537}{\Character{41}}%
\TextAt{8}{79.226}{620.644}{as}%
\TextAt{8}{94.474}{620.644}{being}%
\TextAt{8}{126.89}{620.644}{the}%
\TextAt{8}{147.924}{620.644}{set}%
\TextAt{8}{167.07}{620.644}{of}%
\TextAt{8}{181.278}{620.644}{all}%
\TextAt{8}{198.415}{620.644}{prime}%
\TextAt{8}{232.456}{620.644}{ideals}%
\TextAt{14}{265.912}{620.644}{p}%
\TextAt{8}{276.67}{620.644}{of}%
\TextAt{14}{290.878}{620.644}{o}%
\TextAt{8}{296.723}{620.644}{.}%
\TextAt{8}{307.805}{620.644}{The}%
\TextAt{12}{332.737}{620.644}{Zariski}%
\TextAt{12}{378.387}{620.644}{topology}%
\TextAt{8}{434.008}{620.644}{on}%
\TextAt{13}{451.145}{620.644}{X}%
\TextAt{8}{466.575}{620.644}{is}%
\TextAt{8}{479.217}{620.644}{defined}%
\TextAt{8}{520.408}{620.644}{by}%
\TextAt{8}{79.226}{603.598}{stipulating}%
\TextAt{8}{137.83}{603.598}{that}%
\TextAt{8}{163.199}{603.598}{the}%
\TextAt{8}{183.356}{603.598}{sets}%
\TextAt{13}{257.964}{586.552}{V}%
\TextAt{9}{267.39}{586.552}{\Character{40}}%
\TextAt{14}{271.942}{586.552}{a}%
\TextAt{9}{277.925}{586.552}{\Character{41}}%
\TextAt{9}{285.78}{586.552}{\Character{61}}%
\TextAt{16}{298.224}{586.552}{\Character{102}}%
\TextAt{14}{304.202}{586.552}{p}%
\TextAt{16}{313.506}{586.552}{\Character{106}}%
\TextAt{14}{320.147}{586.552}{p}%
\NativeMathAt{329.451}{586.552}{11.955}{\supseteq}%
\TextAt{14}{342.071}{586.552}{a}%
\TextAt{16}{348.054}{586.552}{\Character{103}}%
\TextAt{8}{79.226}{564.659}{be}%
\TextAt{8}{94.417}{564.659}{the}%
\TextAt{8}{113.84}{564.659}{closed}%
\TextAt{8}{147.633}{564.659}{sets,}%
\TextAt{8}{173.18}{564.659}{where}%
\TextAt{14}{206.256}{564.659}{a}%
\TextAt{8}{215.403}{564.659}{varies}%
\TextAt{8}{247.582}{564.659}{over}%
\TextAt{8}{271.894}{564.659}{the}%
\TextAt{8}{291.317}{564.659}{ideals}%
\TextAt{8}{323.162}{564.659}{of}%
\TextAt{14}{335.758}{564.659}{o}%
\TextAt{8}{341.602}{564.659}{.}%
\TextAt{8}{349.806}{564.659}{This}%
\TextAt{8}{375.792}{564.659}{does}%
\TextAt{8}{401.456}{564.659}{make}%
\TextAt{13}{431.292}{564.659}{X}%
\TextAt{8}{445.111}{564.659}{into}%
\TextAt{8}{468.12}{564.659}{a}%
\TextAt{8}{477.142}{564.659}{topological}%
\TextAt{8}{79.226}{547.613}{space}%
\TextAt{8}{110.327}{547.613}{(observe}%
\TextAt{8}{156.396}{547.613}{that}%
\TextAt{13}{181.592}{547.613}{V}%
\TextAt{9}{191.017}{547.613}{\Character{40}}%
\TextAt{14}{195.57}{547.613}{a}%
\TextAt{9}{201.553}{547.613}{\Character{41}}%
\TextAt{16}{208.407}{547.613}{\Character{91}}%
\TextAt{13}{218.679}{547.613}{V}%
\TextAt{9}{228.104}{547.613}{\Character{40}}%
\TextAt{14}{232.657}{547.613}{b}%
\TextAt{9}{238.792}{547.613}{\Character{41}}%
\TextAt{9}{246.647}{547.613}{\Character{61}}%
\TextAt{13}{259.09}{547.613}{V}%
\TextAt{9}{268.516}{547.613}{\Character{40}}%
\TextAt{14}{273.068}{547.613}{ab}%
\TextAt{9}{285.186}{547.613}{\Character{41}}%
\TextAt{8}{289.738}{547.613}{)}%
\TextAt{8}{298.017}{547.613}{which,}%
\TextAt{8}{334.617}{547.613}{however,}%
\TextAt{8}{382.609}{547.613}{is}%
\TextAt{8}{394.2}{547.613}{usually}%
\TextAt{8}{434.088}{547.613}{not}%
\TextAt{8}{454.728}{547.613}{Hausdorff.}%
\TextAt{8}{512.613}{547.613}{The}%
\TextAt{8}{79.226}{530.567}{closed}%
\TextAt{8}{113.753}{530.567}{points}%
\TextAt{8}{148.937}{530.567}{correspond}%
\TextAt{8}{208.509}{530.567}{to}%
\TextAt{8}{222.82}{530.567}{the}%
\TextAt{8}{242.976}{530.567}{maximal}%
\TextAt{8}{290.785}{530.567}{ideals}%
\TextAt{8}{323.363}{530.567}{of}%
\TextAt{14}{336.693}{530.567}{o}%
\TextAt{8}{342.537}{530.567}{.}%
\TextAt{8}{97.159}{510.532}{The}%
\TextAt{8}{121.933}{510.532}{elements}%
\TextAt{13}{170.509}{510.532}{f}%
\NativeMathAt{182.103}{510.532}{11.955}{\in}%
\TextAt{14}{194.621}{510.532}{o}%
\TextAt{8}{205.083}{510.532}{now}%
\TextAt{8}{230.192}{510.532}{play}%
\TextAt{8}{256.28}{510.532}{the}%
\TextAt{8}{277.157}{510.532}{role}%
\TextAt{8}{300.639}{510.532}{of}%
\TextAt{8}{314.689}{510.532}{functions}%
\TextAt{8}{365.872}{510.532}{on}%
\TextAt{8}{382.851}{510.532}{the}%
\TextAt{8}{403.728}{510.532}{topological}%
\TextAt{8}{463.972}{510.532}{space}%
\TextAt{13}{495.967}{510.532}{X}%
\TextAt{8}{506.622}{510.532}{:}%
\TextAt{8}{516.51}{510.532}{the}%
\TextAt{8}{79.226}{493.486}{“value”}%
\TextAt{8}{120.519}{493.486}{of}%
\TextAt{13}{133.849}{493.486}{f}%
\TextAt{8}{144.793}{493.486}{at}%
\TextAt{8}{159.103}{493.486}{the}%
\TextAt{8}{179.26}{493.486}{point}%
\TextAt{14}{209.829}{493.486}{p}%
\TextAt{8}{219.71}{493.486}{is}%
\TextAt{8}{231.474}{493.486}{defined}%
\TextAt{8}{271.786}{493.486}{to}%
\TextAt{8}{286.097}{493.486}{be}%
\TextAt{13}{264.538}{468.285}{f}%
\TextAt{9}{271.584}{468.284}{\Character{40}}%
\TextAt{14}{276.136}{468.285}{p}%
\TextAt{9}{282.119}{468.284}{\Character{41}}%
\TextAt{9}{289.974}{468.284}{\Character{58}\Character{61}}%
\TextAt{13}{305.67}{468.285}{f}%
\TextAt{8}{316.037}{468.284}{mod}%
\TextAt{14}{341.475}{468.285}{p}%
\TextAt{8}{79.226}{443.083}{and}%
\TextAt{8}{102.488}{443.083}{is}%
\TextAt{8}{114.752}{443.083}{an}%
\TextAt{8}{131.511}{443.083}{element}%
\TextAt{8}{175.252}{443.083}{of}%
\TextAt{8}{189.082}{443.083}{the}%
\TextAt{8}{209.738}{443.083}{residue}%
\TextAt{8}{249.965}{443.083}{class}%
\TextAt{8}{277.901}{443.083}{field}%
\TextAt{13}{303.758}{443.083}{\Character{20}}%
\TextAt{9}{310.496}{443.083}{\Character{40}}%
\TextAt{14}{315.049}{443.083}{p}%
\TextAt{9}{321.032}{443.083}{\Character{41}}%
\TextAt{8}{325.584}{443.083}{,}%
\TextAt{8}{333.358}{443.083}{i.e.,}%
\TextAt{8}{356.088}{443.083}{in}%
\TextAt{8}{370.24}{443.083}{the}%
\TextAt{8}{390.896}{443.083}{field}%
\TextAt{8}{416.753}{443.083}{of}%
\TextAt{8}{430.582}{443.083}{fractions}%
\TextAt{8}{478.951}{443.083}{of}%
\TextAt{14}{492.78}{443.083}{o}%
\TextAt{9}{498.624}{443.083}{\Character{47}}%
\TextAt{14}{504.477}{443.083}{p}%
\TextAt{8}{510.461}{443.083}{.}%
\TextAt{8}{520.408}{443.083}{So}%
\TextAt{8}{79.226}{426.037}{the}%
\TextAt{8}{99.383}{426.037}{values}%
\TextAt{8}{134.244}{426.037}{of}%
\TextAt{13}{147.574}{426.037}{f}%
\TextAt{8}{158.518}{426.037}{do}%
\TextAt{8}{174.777}{426.037}{not}%
\TextAt{8}{195.591}{426.037}{in}%
\TextAt{8}{209.243}{426.037}{general}%
\TextAt{8}{249.568}{426.037}{lie}%
\TextAt{8}{265.17}{426.037}{in}%
\TextAt{8}{278.823}{426.037}{a}%
\TextAt{8}{288.578}{426.037}{single}%
\TextAt{8}{321.156}{426.037}{field.}%
\TextAt{8}{97.159}{406.002}{Admitting}%
\TextAt{8}{154.657}{406.002}{also}%
\TextAt{8}{179.052}{406.002}{the}%
\TextAt{8}{200.122}{406.002}{non-maximal}%
\TextAt{8}{271.608}{406.002}{prime}%
\TextAt{8}{305.686}{406.002}{ideals}%
\TextAt{8}{339.179}{406.002}{as}%
\TextAt{8}{354.463}{406.002}{non-closed}%
\TextAt{8}{412.667}{406.002}{points,}%
\TextAt{8}{452.246}{406.002}{turns}%
\TextAt{8}{483.789}{406.002}{out}%
\TextAt{8}{505.518}{406.002}{to}%
\TextAt{8}{520.743}{406.002}{be}%
\TextAt{8}{79.226}{388.956}{extremely}%
\TextAt{8}{134.326}{388.956}{useful—and}%
\TextAt{8}{199.563}{388.956}{has}%
\TextAt{8}{221.558}{388.956}{an}%
\TextAt{8}{238.939}{388.956}{intuitive}%
\TextAt{8}{286.567}{388.956}{reason}%
\TextAt{8}{324.176}{388.956}{as}%
\TextAt{8}{339.668}{388.956}{well.}%
\TextAt{8}{371.315}{388.956}{For}%
\TextAt{8}{393.406}{388.956}{instance}%
\TextAt{8}{440.113}{388.956}{in}%
\TextAt{8}{454.887}{388.956}{the}%
\TextAt{8}{476.166}{388.956}{case}%
\TextAt{8}{502.059}{388.956}{of}%
\TextAt{8}{516.51}{388.956}{the}%
\TextAt{8}{79.226}{371.91}{ring}%
\TextAt{14}{104.479}{371.91}{o}%
\TextAt{9}{115.667}{371.91}{\Character{61}}%
\TextAt{18}{130.115}{371.91}{C}%
\TextAt{9}{138.749}{371.91}{\Character{91}}%
\TextAt{13}{142.001}{371.91}{x}%
\TextAt{9}{148.653}{371.91}{\Character{93}}%
\TextAt{8}{151.905}{371.91}{,}%
\TextAt{8}{160.538}{371.91}{the}%
\TextAt{8}{181.881}{371.91}{point}%
\TextAt{14}{213.638}{371.91}{p}%
\TextAt{9}{224.964}{371.91}{\Character{61}}%
\TextAt{9}{239.394}{371.91}{\Character{40}0\Character{41}}%
\TextAt{8}{259.455}{371.91}{has}%
\TextAt{8}{281.515}{371.91}{residue}%
\TextAt{8}{322.429}{371.91}{class}%
\TextAt{8}{351.053}{371.91}{field}%
\TextAt{13}{377.597}{371.91}{\Character{20}}%
\TextAt{9}{384.336}{371.91}{\Character{40}}%
\TextAt{14}{388.888}{371.91}{p}%
\TextAt{9}{394.871}{371.91}{\Character{41}}%
\TextAt{9}{404.746}{371.91}{\Character{61}}%
\TextAt{18}{419.216}{371.91}{C}%
\TextAt{9}{427.85}{371.91}{\Character{40}}%
\TextAt{13}{432.402}{371.91}{x}%
\TextAt{9}{439.054}{371.91}{\Character{41}}%
\TextAt{8}{443.607}{371.91}{.}%
\TextAt{8}{455.616}{371.91}{The}%
\TextAt{8}{480.857}{371.91}{“value”}%
\TextAt{8}{523.337}{371.91}{of}%
\TextAt{8}{79.226}{354.864}{a}%
\TextAt{8}{89.802}{354.864}{polynomial}%
\TextAt{13}{151.117}{354.864}{f}%
\NativeMathAt{162.883}{354.864}{11.955}{\in}%
\TextAt{18}{175.573}{354.864}{C}%
\TextAt{9}{184.207}{354.864}{\Character{91}}%
\TextAt{13}{187.459}{354.864}{x}%
\TextAt{9}{194.111}{354.864}{\Character{93}}%
\TextAt{8}{202.081}{354.864}{at}%
\TextAt{8}{217.212}{354.864}{this}%
\TextAt{8}{240.855}{354.864}{point}%
\TextAt{8}{272.245}{354.864}{is}%
\TextAt{13}{284.83}{354.864}{f}%
\TextAt{9}{291.877}{354.864}{\Character{40}}%
\TextAt{13}{296.429}{354.864}{x}%
\TextAt{9}{303.081}{354.864}{\Character{41}}%
\TextAt{8}{312.352}{354.864}{itself,}%
\TextAt{8}{344.975}{354.864}{viewed}%
\TextAt{8}{384.16}{354.864}{as}%
\TextAt{8}{399.351}{354.864}{an}%
\TextAt{8}{416.431}{354.864}{element}%
\TextAt{8}{460.494}{354.864}{of}%
\TextAt{18}{474.645}{354.864}{C}%
\TextAt{9}{483.279}{354.864}{\Character{40}}%
\TextAt{13}{487.832}{354.864}{x}%
\TextAt{9}{494.484}{354.864}{\Character{41}}%
\TextAt{8}{499.036}{354.864}{.}%
\TextAt{8}{509.947}{354.864}{This}%
\TextAt{8}{79.226}{337.818}{element}%
\TextAt{8}{122.576}{337.818}{should}%
\TextAt{8}{159.816}{337.818}{be}%
\TextAt{8}{175.848}{337.818}{thought}%
\TextAt{8}{219.867}{337.818}{of}%
\TextAt{8}{233.305}{337.818}{as}%
\TextAt{8}{247.783}{337.818}{the}%
\TextAt{8}{268.047}{337.818}{value}%
\TextAt{8}{298.401}{337.818}{of}%
\TextAt{13}{311.839}{337.818}{f}%
\TextAt{8}{322.891}{337.818}{at}%
\TextAt{8}{337.309}{337.818}{the}%
\TextAt{8}{357.573}{337.818}{unknown}%
\TextAt{8}{407.761}{337.818}{place}%
\TextAt{13}{437.78}{337.818}{x}%
\TextAt{8}{444.432}{337.818}{—which}%
\TextAt{8}{489.731}{337.818}{one}%
\TextAt{8}{511.298}{337.818}{may}%
\TextAt{8}{79.226}{320.773}{imagine}%
\TextAt{8}{123.441}{320.773}{to}%
\TextAt{8}{138.39}{320.773}{be}%
\TextAt{8}{154.952}{320.773}{everywhere}%
\TextAt{8}{216.395}{320.773}{or}%
\TextAt{8}{231.343}{320.773}{nowhere}%
\TextAt{8}{277.83}{320.773}{at}%
\TextAt{8}{292.779}{320.773}{all.}%
\TextAt{8}{315.504}{320.773}{This}%
\TextAt{8}{342.862}{320.773}{intuition}%
\TextAt{8}{391.309}{320.773}{complies}%
\TextAt{8}{439.481}{320.773}{with}%
\TextAt{8}{466.779}{320.773}{the}%
\TextAt{8}{487.574}{320.773}{fact}%
\TextAt{8}{511.298}{320.773}{that}%
\TextAt{8}{79.226}{303.727}{the}%
\TextAt{8}{99.164}{303.727}{closure}%
\TextAt{8}{138.027}{303.727}{of}%
\TextAt{8}{151.138}{303.727}{the}%
\TextAt{8}{171.076}{303.727}{point}%
\TextAt{14}{201.427}{303.727}{p}%
\TextAt{9}{210.731}{303.727}{\Character{61}}%
\TextAt{9}{223.141}{303.727}{\Character{40}0\Character{41}}%
\TextAt{8}{241.793}{303.727}{in}%
\TextAt{8}{255.227}{303.727}{the}%
\TextAt{8}{275.165}{303.727}{Zariski}%
\TextAt{8}{313.705}{303.727}{topology}%
\TextAt{8}{361.63}{303.727}{of}%
\TextAt{13}{374.741}{303.727}{X}%
\TextAt{8}{389.075}{303.727}{is}%
\TextAt{8}{400.621}{303.727}{the}%
\TextAt{8}{420.559}{303.727}{total}%
\TextAt{8}{448.315}{303.727}{space}%
\TextAt{13}{479.371}{303.727}{X}%
\TextAt{8}{490.026}{303.727}{.}%
\TextAt{8}{498.402}{303.727}{This}%
\TextAt{8}{524.903}{303.727}{is}%
\TextAt{8}{79.226}{286.681}{why}%
\TextAt{14}{103.937}{286.681}{p}%
\TextAt{8}{113.818}{286.681}{is}%
\TextAt{8}{125.582}{286.681}{also}%
\TextAt{8}{149.062}{286.681}{called}%
\TextAt{8}{182.225}{286.681}{the}%
\TextAt{12}{202.382}{286.681}{generic}%
\TextAt{12}{248.553}{286.681}{point}%
\TextAt{8}{283.103}{286.681}{of}%
\TextAt{13}{296.433}{286.681}{X}%
\TextAt{8}{307.088}{286.681}{.}%
\TextAt{12}{97.159}{266.646}{Example:}%
\TextAt{8}{157.122}{266.646}{The}%
\TextAt{8}{181.176}{266.646}{space}%
\TextAt{13}{212.451}{266.646}{X}%
\TextAt{9}{226.427}{266.646}{\Character{61}}%
\TextAt{8}{238.852}{266.646}{Spec}%
\TextAt{9}{262.26}{266.646}{\Character{40}}%
\TextAt{18}{266.813}{266.646}{Z}%
\TextAt{9}{274.783}{266.646}{\Character{41}}%
\TextAt{8}{283.232}{266.646}{may}%
\TextAt{8}{308.601}{266.646}{be}%
\TextAt{8}{324.526}{266.646}{represented}%
\TextAt{8}{386.692}{266.646}{by}%
\TextAt{8}{402.951}{266.646}{a}%
\TextAt{8}{412.707}{266.646}{line.}%
\NativeDiagram{%
\path[draw=black,line width=0.399bp] (78.092,-524.918) -- (389.907,-524.918);%
\path[fill=black] (106.439,-524.918) (107.933,-524.918) .. controls (107.933,-524.093) and (107.264,-523.424) .. (106.439,-523.424) .. controls (105.613,-523.424) and (104.944,-524.093) .. (104.944,-524.918) .. controls (104.944,-525.743) and (105.613,-526.412) .. (106.439,-526.412) .. controls (107.264,-526.412) and (107.933,-525.743) .. (107.933,-524.918) -- cycle (106.439,-524.918);%
}%
\TextAt{9}{175.512}{233.792}{2}%
\NativeDiagram{%
\path[fill=black] (146.124,-524.918) (147.619,-524.918) .. controls (147.619,-524.093) and (146.95,-523.424) .. (146.124,-523.424) .. controls (145.299,-523.424) and (144.63,-524.093) .. (144.63,-524.918) .. controls (144.63,-525.743) and (145.299,-526.412) .. (146.124,-526.412) .. controls (146.95,-526.412) and (147.619,-525.743) .. (147.619,-524.918) -- cycle (146.124,-524.918);%
}%
\TextAt{9}{215.198}{233.792}{3}%
\NativeDiagram{%
\path[fill=black] (182.975,-524.918) (184.47,-524.918) .. controls (184.47,-524.093) and (183.8,-523.424) .. (182.975,-523.424) .. controls (182.15,-523.424) and (181.481,-524.093) .. (181.481,-524.918) .. controls (181.481,-525.743) and (182.15,-526.412) .. (182.975,-526.412) .. controls (183.8,-526.412) and (184.47,-525.743) .. (184.47,-524.918) -- cycle (182.975,-524.918);%
}%
\TextAt{9}{252.049}{233.792}{5}%
\NativeDiagram{%
\path[fill=black] (214.157,-524.918) (215.651,-524.918) .. controls (215.651,-524.093) and (214.982,-523.424) .. (214.157,-523.424) .. controls (213.331,-523.424) and (212.662,-524.093) .. (212.662,-524.918) .. controls (212.662,-525.743) and (213.331,-526.412) .. (214.157,-526.412) .. controls (214.982,-526.412) and (215.651,-525.743) .. (215.651,-524.918) -- cycle (214.157,-524.918);%
}%
\TextAt{9}{283.23}{233.792}{7}%
\NativeDiagram{%
\path[fill=black] (239.669,-524.918) (241.163,-524.918) .. controls (241.163,-524.093) and (240.494,-523.424) .. (239.669,-523.424) .. controls (238.843,-523.424) and (238.174,-524.093) .. (238.174,-524.918) .. controls (238.174,-525.743) and (238.843,-526.412) .. (239.669,-526.412) .. controls (240.494,-526.412) and (241.163,-525.743) .. (241.163,-524.918) -- cycle (239.669,-524.918);%
}%
\TextAt{9}{305.816}{233.792}{11}%
\NativeDiagram{%
\path[fill=black] (258.094,-524.918) (259.091,-524.918) .. controls (259.091,-524.368) and (258.645,-523.922) .. (258.094,-523.922) .. controls (257.544,-523.922) and (257.098,-524.368) .. (257.098,-524.918) .. controls (257.098,-525.468) and (257.544,-525.914) .. (258.094,-525.914) .. controls (258.645,-525.914) and (259.091,-525.468) .. (259.091,-524.918) -- cycle (258.094,-524.918);%
\path[fill=black] (266.598,-524.918) (267.594,-524.918) .. controls (267.594,-524.368) and (267.148,-523.922) .. (266.598,-523.922) .. controls (266.048,-523.922) and (265.602,-524.368) .. (265.602,-524.918) .. controls (265.602,-525.468) and (266.048,-525.914) .. (266.598,-525.914) .. controls (267.148,-525.914) and (267.594,-525.468) .. (267.594,-524.918) -- cycle (266.598,-524.918);%
\path[fill=black] (273.685,-524.918) (274.681,-524.918) .. controls (274.681,-524.368) and (274.235,-523.922) .. (273.685,-523.922) .. controls (273.135,-523.922) and (272.689,-524.368) .. (272.689,-524.918) .. controls (272.689,-525.468) and (273.135,-525.914) .. (273.685,-525.914) .. controls (274.235,-525.914) and (274.681,-525.468) .. (274.681,-524.918) -- cycle (273.685,-524.918);%
\path[fill=black] (279.355,-524.918) (280.351,-524.918) .. controls (280.351,-524.368) and (279.905,-523.922) .. (279.355,-523.922) .. controls (278.804,-523.922) and (278.358,-524.368) .. (278.358,-524.918) .. controls (278.358,-525.468) and (278.804,-525.914) .. (279.355,-525.914) .. controls (279.905,-525.914) and (280.351,-525.468) .. (280.351,-524.918) -- cycle (279.355,-524.918);%
\path[fill=black] (283.606,-524.918) (284.603,-524.918) .. controls (284.603,-524.368) and (284.157,-523.922) .. (283.606,-523.922) .. controls (283.056,-523.922) and (282.61,-524.368) .. (282.61,-524.918) .. controls (282.61,-525.468) and (283.056,-525.914) .. (283.606,-525.914) .. controls (284.157,-525.914) and (284.603,-525.468) .. (284.603,-524.918) -- cycle (283.606,-524.918);%
\path[draw=black,line width=0.399bp,pattern=dots] (323.292,-524.918) (330.095,-524.918) .. controls (330.095,-522.883) and (327.049,-521.233) .. (323.292,-521.233) .. controls (319.534,-521.233) and (316.488,-522.883) .. (316.488,-524.918) .. controls (316.488,-526.953) and (319.534,-528.603) .. (323.292,-528.603) .. controls (327.049,-528.603) and (330.095,-526.953) .. (330.095,-524.918) -- cycle (323.292,-524.918);%
}%
\TextAt{8}{362.122}{232.839}{generic}%
\TextAt{8}{401.79}{232.839}{point}%
\TextAt{8}{79.226}{199.663}{For}%
\TextAt{8}{100.747}{199.663}{every}%
\TextAt{8}{132.191}{199.663}{prime}%
\TextAt{8}{165.906}{199.663}{number}%
\TextAt{8}{209.054}{199.663}{one}%
\TextAt{8}{231.065}{199.663}{has}%
\TextAt{8}{252.491}{199.663}{a}%
\TextAt{8}{262.798}{199.663}{closed}%
\TextAt{8}{297.876}{199.663}{point,}%
\TextAt{8}{332.387}{199.663}{and}%
\TextAt{8}{355.701}{199.663}{there}%
\TextAt{8}{386.165}{199.663}{is}%
\TextAt{8}{398.48}{199.663}{also}%
\TextAt{8}{422.512}{199.663}{the}%
\TextAt{8}{443.22}{199.663}{generic}%
\TextAt{8}{483.439}{199.663}{point}%
\TextAt{9}{514.56}{199.663}{\Character{40}0\Character{41}}%
\TextAt{8}{529.518}{199.663}{,}%
\TextAt{8}{79.226}{182.618}{the}%
\TextAt{8}{99.33}{182.618}{closure}%
\TextAt{8}{138.36}{182.618}{of}%
\TextAt{8}{151.638}{182.618}{which}%
\TextAt{8}{185.072}{182.618}{is}%
\TextAt{8}{196.784}{182.618}{the}%
\TextAt{8}{216.888}{182.618}{total}%
\TextAt{8}{244.811}{182.618}{space}%
\TextAt{13}{276.034}{182.618}{X}%
\TextAt{8}{286.689}{182.618}{.}%
\TextAt{8}{295.12}{182.618}{The}%
\TextAt{8}{319.122}{182.618}{nonempty}%
\TextAt{8}{373.705}{182.618}{open}%
\TextAt{8}{401.938}{182.618}{sets}%
\TextAt{8}{424.769}{182.618}{in}%
\TextAt{13}{438.369}{182.618}{X}%
\TextAt{8}{452.87}{182.618}{are}%
\TextAt{8}{472.328}{182.618}{obtained}%
\TextAt{8}{520.408}{182.618}{by}%
\TextAt{8}{79.226}{165.572}{throwing}%
\TextAt{8}{128.734}{165.572}{out}%
\TextAt{8}{149.945}{165.572}{finitely}%
\TextAt{8}{189.685}{165.572}{many}%
\TextAt{8}{221.954}{165.572}{prime}%
\TextAt{8}{255.514}{165.572}{numbers}%
\TextAt{13}{303.121}{165.572}{p}%
\TextAt{17}{308.996}{163.778}{1}%
\TextAt{13}{313.729}{165.572}{\Character{59}}%
\TextAt{13}{318.953}{165.572}{\Character{58}}%
\TextAt{13}{324.19}{165.572}{\Character{58}}%
\TextAt{13}{329.414}{165.572}{\Character{58}}%
\TextAt{13}{334.651}{165.572}{\Character{59}}%
\TextAt{13}{339.887}{165.572}{p}%
\TextAt{21}{345.825}{163.778}{n}%
\TextAt{8}{351.461}{165.572}{.}%
\TextAt{8}{361.098}{165.572}{The}%
\TextAt{8}{385.548}{165.572}{integers}%
\TextAt{13}{429.258}{165.572}{a}%
\NativeMathAt{439.399}{165.572}{11.955}{\in}%
\TextAt{18}{451.365}{165.572}{Z}%
\TextAt{8}{463.629}{165.572}{are}%
\TextAt{8}{483.536}{165.572}{viewed}%
\TextAt{8}{522.297}{165.572}{as}%
\TextAt{8}{79.226}{148.526}{functions}%
\TextAt{8}{129.689}{148.526}{on}%
\TextAt{13}{145.948}{148.526}{X}%
\TextAt{8}{160.5}{148.526}{by}%
\TextAt{8}{176.759}{148.526}{defining}%
\TextAt{8}{220.982}{148.526}{the}%
\TextAt{8}{241.138}{148.526}{value}%
\TextAt{8}{271.385}{148.526}{of}%
\TextAt{13}{284.715}{148.526}{a}%
\TextAt{8}{294.757}{148.526}{at}%
\TextAt{8}{309.067}{148.526}{the}%
\TextAt{8}{329.224}{148.526}{point}%
\TextAt{9}{359.793}{148.526}{\Character{40}}%
\TextAt{13}{364.345}{148.526}{p}%
\TextAt{9}{370.221}{148.526}{\Character{41}}%
\TextAt{8}{378.67}{148.526}{to}%
\TextAt{8}{392.981}{148.526}{be}%
\TextAt{8}{408.905}{148.526}{the}%
\TextAt{8}{429.061}{148.526}{residue}%
\TextAt{8}{468.788}{148.526}{class}%
\TextAt{13}{244.407}{123.324}{a}%
\TextAt{9}{250.552}{123.324}{\Character{40}}%
\TextAt{13}{255.104}{123.324}{p}%
\TextAt{9}{260.979}{123.324}{\Character{41}}%
\TextAt{9}{268.834}{123.324}{\Character{61}}%
\TextAt{13}{281.278}{123.324}{a}%
\TextAt{8}{290.744}{123.324}{mod}%
\TextAt{13}{316.182}{123.324}{p}%
\NativeMathAt{325.378}{123.324}{11.955}{\in}%
\TextAt{18}{336.669}{123.324}{Z}%
\TextAt{9}{344.639}{123.324}{\Character{47}}%
\TextAt{13}{350.492}{123.324}{p}%
\TextAt{18}{356.367}{123.324}{Z}%
\TextAt{13}{364.337}{123.324}{\Character{58}}%
\TextAt{8}{79.226}{98.123}{The}%
\TextAt{8}{103.28}{98.123}{fields}%
\TextAt{8}{133.252}{98.123}{of}%
\TextAt{8}{146.582}{98.123}{values}%
\TextAt{8}{181.443}{98.123}{are}%
\TextAt{8}{200.954}{98.123}{then}%
\TextAt{18}{192.39}{72.921}{Z}%
\TextAt{9}{200.36}{72.921}{\Character{47}2}%
\TextAt{18}{212.066}{72.921}{Z}%
\TextAt{13}{220.036}{72.921}{\Character{59}}%
\TextAt{18}{229.178}{72.921}{Z}%
\TextAt{9}{237.148}{72.921}{\Character{47}3}%
\TextAt{18}{248.854}{72.921}{Z}%
\TextAt{13}{256.824}{72.921}{\Character{59}}%
\TextAt{18}{265.965}{72.921}{Z}%
\TextAt{9}{273.936}{72.921}{\Character{47}5}%
\TextAt{18}{285.642}{72.921}{Z}%
\TextAt{13}{293.612}{72.921}{\Character{59}}%
\TextAt{18}{302.753}{72.921}{Z}%
\TextAt{9}{310.723}{72.921}{\Character{47}7}%
\TextAt{18}{322.429}{72.921}{Z}%
\TextAt{13}{330.399}{72.921}{\Character{59}}%
\TextAt{18}{339.541}{72.921}{Z}%
\TextAt{9}{347.511}{72.921}{\Character{47}11}%
\TextAt{18}{365.07}{72.921}{Z}%
\TextAt{13}{373.04}{72.921}{\Character{59}}%
\TextAt{13}{382.162}{72.921}{\Character{58}}%
\TextAt{13}{387.398}{72.921}{\Character{58}}%
\TextAt{13}{392.635}{72.921}{\Character{58}}%
\TextAt{13}{397.859}{72.921}{\Character{59}}%
\TextAt{18}{407.056}{72.921}{Q}%
\TextAt{13}{416.354}{72.921}{\Character{58}}%
\TextAt{8}{300.14}{43.033}{80}%
\Destination{78.23}{809.815}{nav.page.80}%
\Destination{79.36}{525.555}{nav.print.86}%
}

% Page 93 (body page 81)
\Page{%
\TextAt{2}{79.226}{796.943}{I.}%
\TextAt{2}{89.83}{796.943}{Algebraic}%
\TextAt{2}{138.648}{796.943}{Integers}%
\Vector{q 0 G 0 g 0.299 w 7.226 18.729 m 460.769 18.729 l S Q}%
\TextAt{8}{79.226}{761.124}{Thus}%
\TextAt{8}{108.875}{761.124}{every}%
\TextAt{8}{139.767}{761.124}{prime}%
\TextAt{8}{172.931}{761.124}{field}%
\TextAt{8}{198.288}{761.124}{occurs}%
\TextAt{8}{234.44}{761.124}{exactly}%
\TextAt{8}{274.765}{761.124}{once.}%
\TextAt{8}{97.159}{741.089}{An}%
\TextAt{8}{116.592}{741.089}{important}%
\TextAt{8}{172.141}{741.089}{refinement}%
\TextAt{8}{229.95}{741.089}{of}%
\TextAt{8}{243.537}{741.089}{the}%
\TextAt{8}{263.95}{741.089}{geometric}%
\TextAt{8}{317.539}{741.089}{interpretation}%
\TextAt{8}{392.276}{741.089}{of}%
\TextAt{8}{405.863}{741.089}{elements}%
\TextAt{8}{453.976}{741.089}{of}%
\TextAt{8}{467.563}{741.089}{the}%
\TextAt{8}{487.976}{741.089}{ring}%
\TextAt{14}{512.298}{741.089}{o}%
\TextAt{8}{522.297}{741.089}{as}%
\TextAt{8}{79.226}{724.043}{functions}%
\TextAt{8}{130.799}{724.043}{on}%
\TextAt{8}{148.169}{724.043}{the}%
\TextAt{8}{169.436}{724.043}{space}%
\TextAt{13}{201.821}{724.043}{X}%
\TextAt{9}{217.689}{724.043}{\Character{61}}%
\TextAt{8}{232.007}{724.043}{Spec}%
\TextAt{9}{255.415}{724.043}{\Character{40}}%
\TextAt{14}{259.967}{724.043}{o}%
\TextAt{9}{265.812}{724.043}{\Character{41}}%
\TextAt{8}{275.372}{724.043}{is}%
\TextAt{8}{288.246}{724.043}{obtained}%
\TextAt{8}{337.488}{724.043}{by}%
\TextAt{8}{354.857}{724.043}{forming}%
\TextAt{8}{399.222}{724.043}{the}%
\TextAt{12}{420.488}{724.043}{structure}%
\TextAt{12}{480.083}{724.043}{sheaf}%
\TextAt{14}{515.588}{724.043}{o}%
\TextAt{21}{521.433}{722.25}{X}%
\TextAt{8}{529.518}{724.043}{.}%
\TextAt{8}{79.226}{706.997}{This}%
\TextAt{8}{106.116}{706.997}{means}%
\TextAt{8}{142.116}{706.997}{the}%
\TextAt{8}{162.443}{706.997}{following.}%
\TextAt{8}{216.939}{706.997}{Let}%
\TextAt{13}{238.078}{706.997}{U}%
\NativeMathAt{250.888}{706.997}{11.955}{\neq}%
\TextAt{18}{263.604}{706.997}{\Character{63}}%
\TextAt{8}{276.97}{706.997}{be}%
\TextAt{8}{293.064}{706.997}{an}%
\TextAt{8}{309.493}{706.997}{open}%
\TextAt{8}{337.949}{706.997}{subset}%
\TextAt{8}{374.009}{706.997}{of}%
\TextAt{13}{387.509}{706.997}{X}%
\TextAt{8}{398.164}{706.997}{.}%
\TextAt{8}{407.123}{706.997}{If}%
\TextAt{14}{418.986}{706.997}{o}%
\TextAt{8}{428.897}{706.997}{is}%
\TextAt{8}{440.831}{706.997}{a}%
\TextAt{8}{450.757}{706.997}{one-dimensional}%
\TextAt{8}{79.226}{689.951}{integral}%
\TextAt{8}{121.834}{689.951}{domain,}%
\TextAt{8}{166.714}{689.951}{then}%
\TextAt{8}{193.374}{689.951}{the}%
\TextAt{8}{213.531}{689.951}{ring}%
\TextAt{8}{237.596}{689.951}{of}%
\TextAt{8}{250.926}{689.951}{“regular}%
\TextAt{8}{296.129}{689.951}{functions”}%
\TextAt{8}{352.115}{689.951}{on}%
\TextAt{13}{368.374}{689.951}{U}%
\TextAt{8}{381.47}{689.951}{is}%
\TextAt{8}{393.234}{689.951}{given}%
\TextAt{8}{423.804}{689.951}{by}%
\ErrataAt{176}{658.015}{250}{\makebox[250bp][c]{$\displaystyle \mathfrak o(U)=\bigcap_{\mathfrak p\in U}\mathfrak o_{\mathfrak p}$,}}%
\ErrataAt{79.19}{626.132}{455.81}{where the intersection is in the fraction field. On the empty open set the ring is $0$.}%
\ErrataAt{79.19}{609.086}{455.81}{In the general case, $\mathfrak o(U)$ is defined to consist of all elements}%
\TextAt{13}{264.241}{579.436}{s}%
\TextAt{9}{273.076}{579.436}{\Character{61}}%
\TextAt{9}{285.486}{579.436}{\Character{40}}%
\TextAt{13}{290.054}{579.436}{s}%
\TextAt{27}{295.568}{577.643}{p}%
\TextAt{9}{300.055}{579.436}{\Character{41}}%
\NativeMathAt{307.928}{579.436}{11.955}{\in}%
\TextAt{20}{319.687}{590.794}{\Character{89}}%
\TextAt{27}{319.219}{565.422}{p}%
\NativeMathAt{323.207}{565.422}{7.97}{\in}%
\TextAt{21}{328.853}{565.422}{U}%
\TextAt{14}{337.423}{579.436}{o}%
\TextAt{27}{343.268}{577.643}{p}%
\TextAt{8}{79.226}{542.183}{which}%
\TextAt{8}{112.181}{542.183}{locally}%
\TextAt{8}{148.722}{542.183}{are}%
\TextAt{8}{167.701}{542.183}{quotients}%
\TextAt{8}{217.967}{542.183}{of}%
\TextAt{8}{230.766}{542.183}{two}%
\TextAt{8}{252.351}{542.183}{elements}%
\TextAt{8}{299.676}{542.183}{of}%
\TextAt{14}{312.474}{542.183}{o}%
\TextAt{8}{318.319}{542.183}{.}%
\TextAt{8}{326.59}{542.183}{More}%
\TextAt{8}{356.293}{542.183}{precisely,}%
\TextAt{8}{405.995}{542.183}{this}%
\TextAt{8}{428.286}{542.183}{means}%
\TextAt{8}{463.584}{542.183}{that}%
\TextAt{8}{488.421}{542.183}{for}%
\TextAt{8}{505.775}{542.183}{every}%
\TextAt{14}{79.226}{525.137}{p}%
\NativeMathAt{89.812}{525.137}{11.955}{\in}%
\TextAt{13}{102.385}{525.137}{U}%
\TextAt{8}{111.584}{525.137}{,}%
\TextAt{8}{119.674}{525.137}{there}%
\TextAt{8}{150.338}{525.137}{exists}%
\TextAt{8}{183.405}{525.137}{a}%
\TextAt{8}{193.913}{525.137}{neighbourhood}%
\TextAt{13}{274.669}{525.137}{V}%
\NativeMathAt{288.698}{525.137}{11.955}{\subseteq}%
\TextAt{13}{302.599}{525.137}{U}%
\TextAt{8}{316.448}{525.137}{of}%
\TextAt{14}{330.531}{525.137}{p}%
\TextAt{8}{336.514}{525.137}{,}%
\TextAt{8}{344.603}{525.137}{and}%
\TextAt{8}{368.118}{525.137}{elements}%
\TextAt{13}{416.727}{525.137}{f}%
\TextAt{13}{423.111}{525.137}{\Character{59}}%
\TextAt{13}{428.348}{525.137}{g}%
\NativeMathAt{439.005}{525.137}{11.955}{\in}%
\TextAt{14}{451.578}{525.137}{o}%
\TextAt{8}{462.072}{525.137}{such}%
\TextAt{8}{489.221}{525.137}{that,}%
\TextAt{8}{518.782}{525.137}{for}%
\TextAt{8}{79.226}{508.091}{each}%
\TextAt{14}{105.454}{508.091}{q}%
\NativeMathAt{114.623}{508.091}{11.955}{\in}%
\TextAt{13}{125.913}{508.091}{V}%
\TextAt{8}{135.339}{508.091}{,}%
\TextAt{8}{142.401}{508.091}{one}%
\TextAt{8}{163.751}{508.091}{has}%
\TextAt{13}{184.516}{508.091}{g}%
\TextAt{9}{190.55}{508.091}{\Character{40}}%
\TextAt{14}{195.103}{508.091}{q}%
\TextAt{9}{200.95}{508.091}{\Character{41}}%
\NativeMathAt{208.823}{508.091}{11.955}{\neq}%
\TextAt{9}{221.233}{508.091}{0}%
\TextAt{8}{230.89}{508.091}{and}%
\TextAt{13}{253.543}{508.091}{s}%
\TextAt{27}{259.057}{506.298}{q}%
\TextAt{9}{266.775}{508.091}{\Character{61}}%
\TextAt{13}{279.2}{508.091}{f}%
\TextAt{9}{286.247}{508.091}{\Character{47}}%
\TextAt{13}{292.1}{508.091}{g}%
\TextAt{8}{301.922}{508.091}{in}%
\TextAt{14}{315.466}{508.091}{o}%
\TextAt{27}{321.31}{506.298}{q}%
\TextAt{8}{325.706}{508.091}{.}%
\TextAt{8}{334.118}{508.091}{These}%
\TextAt{8}{367.878}{508.091}{quotients}%
\TextAt{8}{418.566}{508.091}{have}%
\TextAt{8}{445.452}{508.091}{to}%
\TextAt{8}{459.653}{508.091}{be}%
\TextAt{8}{475.468}{508.091}{understood}%
\TextAt{8}{79.226}{491.045}{in}%
\TextAt{8}{92.505}{491.045}{the}%
\TextAt{8}{112.287}{491.045}{more}%
\TextAt{8}{141.179}{491.045}{general}%
\TextAt{8}{181.13}{491.045}{sense}%
\TextAt{8}{210.787}{491.045}{of}%
\TextAt{8}{223.743}{491.045}{commutative}%
\TextAt{8}{293.294}{491.045}{algebra}%
\TextAt{8}{333.902}{491.045}{(see}%
\TextAt{8}{356.996}{491.045}{§11,}%
\TextAt{8}{381.241}{491.045}{exercise}%
\TextAt{8}{424.168}{491.045}{1).}%
\TextAt{8}{442.905}{491.045}{We}%
\TextAt{8}{462.675}{491.045}{leave}%
\TextAt{8}{491.244}{491.045}{it}%
\TextAt{8}{502.574}{491.045}{to}%
\TextAt{8}{516.51}{491.045}{the}%
\TextAt{8}{79.226}{473.999}{reader}%
\TextAt{8}{115.304}{473.999}{to}%
\TextAt{8}{129.922}{473.999}{check}%
\TextAt{8}{161.767}{473.999}{that}%
\TextAt{8}{187.444}{473.999}{one}%
\TextAt{8}{209.211}{473.999}{gets}%
\TextAt{8}{233.644}{473.999}{back}%
\TextAt{8}{261.269}{473.999}{the}%
\TextAt{8}{281.733}{473.999}{above}%
\TextAt{8}{315.216}{473.999}{definition}%
\TextAt{8}{367.553}{473.999}{in}%
\TextAt{8}{381.513}{473.999}{the}%
\TextAt{8}{401.978}{473.999}{case}%
\TextAt{8}{427.056}{473.999}{of}%
\TextAt{8}{440.694}{473.999}{a}%
\TextAt{8}{450.757}{473.999}{one-dimensional}%
\TextAt{8}{79.226}{456.953}{integral}%
\TextAt{8}{121.834}{456.953}{domain}%
\TextAt{14}{163.462}{456.953}{o}%
\TextAt{8}{169.307}{456.953}{.}%
\TextAt{8}{97.159}{436.918}{If}%
\TextAt{13}{108.851}{436.918}{V}%
\NativeMathAt{121.598}{436.918}{11.955}{\subseteq}%
\TextAt{13}{134.217}{436.918}{U}%
\TextAt{8}{147.313}{436.918}{are}%
\TextAt{8}{166.824}{436.918}{two}%
\TextAt{8}{188.941}{436.918}{open}%
\TextAt{8}{217.227}{436.918}{sets}%
\TextAt{8}{240.109}{436.918}{of}%
\TextAt{13}{253.439}{436.918}{X}%
\TextAt{8}{264.094}{436.918}{,}%
\TextAt{8}{271.244}{436.918}{then}%
\TextAt{8}{297.904}{436.918}{the}%
\TextAt{8}{318.06}{436.918}{projection}%
\TextAt{20}{264.861}{420.105}{\Character{89}}%
\TextAt{27}{264.393}{394.734}{p}%
\NativeMathAt{268.382}{394.734}{7.97}{\in}%
\TextAt{21}{274.027}{394.734}{U}%
\TextAt{14}{282.598}{408.748}{o}%
\TextAt{27}{288.442}{406.954}{p}%
\NativeArrowAt{296.249}{408.748}{11.955}{19.26}{\longrightarrow}%
\TextAt{20}{319.417}{420.105}{\Character{89}}%
\TextAt{27}{318.831}{394.734}{p}%
\NativeMathAt{322.82}{394.734}{7.97}{\in}%
\TextAt{21}{328.466}{394.734}{V}%
\TextAt{14}{337.272}{408.748}{o}%
\TextAt{27}{343.116}{406.954}{p}%
\TextAt{8}{79.226}{371.614}{induces}%
\TextAt{8}{120.902}{371.614}{a}%
\TextAt{8}{130.657}{371.614}{homomorphism}%
\TextAt{13}{252.249}{354.568}{\Character{26}}%
\TextAt{21}{258.287}{352.774}{U}%
\TextAt{21}{264.854}{352.774}{V}%
\TextAt{9}{275.497}{354.568}{\Character{58}}%
\TextAt{14}{282.07}{354.568}{o}%
\TextAt{9}{287.914}{354.568}{\Character{40}}%
\TextAt{13}{292.466}{354.568}{U}%
\TextAt{9}{301.665}{354.568}{\Character{41}}%
\NativeArrowAt{309.538}{354.568}{11.955}{19.26}{\longrightarrow}%
\TextAt{14}{332.12}{354.568}{o}%
\TextAt{9}{337.965}{354.568}{\Character{40}}%
\TextAt{13}{342.517}{354.568}{V}%
\TextAt{9}{351.943}{354.568}{\Character{41}}%
\TextAt{13}{356.495}{354.568}{\Character{59}}%
\TextAt{8}{79.226}{331.402}{called}%
\TextAt{8}{113.109}{331.402}{the}%
\TextAt{12}{133.985}{331.402}{restriction}%
\TextAt{8}{199.178}{331.402}{from}%
\TextAt{13}{227.537}{331.402}{U}%
\TextAt{8}{241.353}{331.402}{to}%
\TextAt{13}{256.382}{331.402}{V}%
\TextAt{8}{265.808}{331.402}{.}%
\TextAt{8}{276.414}{331.402}{The}%
\TextAt{8}{301.187}{331.402}{system}%
\TextAt{8}{340.724}{331.402}{of}%
\TextAt{8}{354.774}{331.402}{rings}%
\TextAt{14}{384.173}{331.402}{o}%
\TextAt{9}{390.017}{331.402}{\Character{40}}%
\TextAt{13}{394.57}{331.402}{U}%
\TextAt{9}{403.769}{331.402}{\Character{41}}%
\TextAt{8}{412.938}{331.402}{and}%
\TextAt{8}{436.419}{331.402}{mappings}%
\TextAt{13}{489.885}{331.402}{\Character{26}}%
\TextAt{21}{495.922}{329.609}{U}%
\TextAt{21}{502.489}{329.609}{V}%
\TextAt{8}{514.428}{331.402}{is}%
\TextAt{8}{526.911}{331.402}{a}%
\TextAt{12}{79.226}{314.356}{sheaf}%
\TextAt{8}{113.621}{314.356}{on}%
\TextAt{13}{129.88}{314.356}{X}%
\TextAt{8}{140.535}{314.356}{.}%
\TextAt{8}{148.984}{314.356}{This}%
\TextAt{8}{175.703}{314.356}{notion}%
\TextAt{8}{212.131}{314.356}{means}%
\TextAt{8}{247.961}{314.356}{the}%
\TextAt{8}{268.117}{314.356}{following.}%
\TextAt{12}{79.226}{288.581}{Definition}%
\TextAt{12}{140.987}{288.581}{(13.1).}%
\TextAt{8}{182.139}{288.581}{Let}%
\TextAt{13}{202.215}{288.581}{X}%
\TextAt{8}{215.874}{288.581}{be}%
\TextAt{8}{230.905}{288.581}{a}%
\TextAt{8}{239.766}{288.581}{topological}%
\TextAt{8}{298.398}{288.581}{space.}%
\TextAt{8}{333.925}{288.581}{A}%
\TextAt{12}{345.703}{288.581}{presheaf}%
\TextAt{16}{398.309}{288.581}{\Character{70}}%
\TextAt{8}{411.093}{288.581}{of}%
\TextAt{8}{423.529}{288.581}{abelian}%
\TextAt{8}{463.283}{288.581}{groups}%
\TextAt{8}{500.18}{288.581}{(rings,}%
\TextAt{8}{79.226}{271.535}{etc.)}%
\TextAt{8}{107.185}{271.535}{consists}%
\TextAt{8}{150.296}{271.535}{of}%
\TextAt{8}{163.626}{271.535}{the}%
\TextAt{8}{183.782}{271.535}{following}%
\TextAt{8}{233.217}{271.535}{data.}%
\TextAt{8}{86.973}{251.5}{(1)}%
\TextAt{8}{107.918}{251.5}{For}%
\TextAt{8}{128.888}{251.5}{every}%
\TextAt{8}{159.78}{251.5}{open}%
\TextAt{8}{188.066}{251.5}{set}%
\TextAt{13}{206.333}{251.5}{U}%
\TextAt{8}{215.532}{251.5}{,}%
\TextAt{8}{222.682}{251.5}{an}%
\TextAt{8}{238.941}{251.5}{abelian}%
\TextAt{8}{279.588}{251.5}{group}%
\TextAt{8}{312.764}{251.5}{(a}%
\TextAt{8}{327.074}{251.5}{ring,}%
\TextAt{8}{354.392}{251.5}{etc.)}%
\TextAt{16}{382.351}{251.5}{\Character{70}}%
\TextAt{9}{392.131}{251.5}{\Character{40}}%
\TextAt{13}{396.683}{251.5}{U}%
\TextAt{9}{405.882}{251.5}{\Character{41}}%
\TextAt{8}{414.332}{251.5}{is}%
\TextAt{8}{426.096}{251.5}{given.}%
\ErrataAt{86.86}{234.455}{446}{(2) For every inclusion $V\subseteq U$, a homomorphism $\rho_{UV}\colon\mathcal F(U)\to\mathcal F(V)$ is given, which}%
\TextAt{8}{107.918}{217.408}{is}%
\TextAt{8}{119.682}{217.408}{called}%
\TextAt{12}{152.846}{217.408}{restriction}%
\TextAt{8}{213.423}{217.408}{.}%
\TextAt{8}{300.14}{43.033}{81}%
\Destination{79.23}{265.545}{nav.Item.15}%
\Destination{79.23}{248.495}{nav.Item.16}%
\Destination{78.23}{809.815}{nav.page.81}%
\Destination{79.36}{641.155}{nav.print.87}%
\Destination{79.36}{302.555}{nav.number.1.13.1}%
\Hotlink{362.175}{488.855}{12.716}{11.047}{nav.section.1.11}%
\Hotlink{489.493}{139.755}{10.963}{9.572}{nav.section.1.11}%
\Destination{79.226}{248.455}{errata.12}%
\Destination{79.226}{675.015}{errata.53}%
}

% Page 94 (body page 82)
\Page{%
\TextAt{2}{79.226}{796.943}{I.}%
\TextAt{2}{89.83}{796.943}{Algebraic}%
\TextAt{2}{138.648}{796.943}{Integers}%
\Vector{q 0 G 0 g 0.299 w 7.226 18.729 m 460.769 18.729 l S Q}%
\TextAt{8}{79.226}{761.124}{These}%
\TextAt{8}{113.095}{761.124}{data}%
\TextAt{8}{139.767}{761.124}{are}%
\TextAt{8}{159.278}{761.124}{subject}%
\TextAt{8}{199.651}{761.124}{to}%
\TextAt{8}{213.961}{761.124}{the}%
\TextAt{8}{234.117}{761.124}{following}%
\TextAt{8}{283.552}{761.124}{conditions:}%
\TextAt{8}{86.973}{741.089}{(a)}%
\TextAt{16}{107.918}{741.089}{\Character{70}}%
\TextAt{9}{117.698}{741.089}{\Character{40}}%
\TextAt{18}{122.251}{741.089}{\Character{63}}%
\TextAt{9}{131.549}{741.089}{\Character{41}}%
\TextAt{9}{139.404}{741.089}{\Character{61}}%
\TextAt{9}{151.825}{741.089}{0}%
\TextAt{8}{157.701}{741.089}{,}%
\TextAt{8}{86.327}{724.043}{(b)}%
\TextAt{13}{107.918}{724.043}{\Character{26}}%
\TextAt{21}{113.956}{722.25}{U}%
\TextAt{21}{120.523}{722.25}{U}%
\TextAt{8}{131.508}{724.043}{is}%
\TextAt{8}{143.271}{724.043}{the}%
\TextAt{8}{163.428}{724.043}{identity}%
\TextAt{8}{206.682}{724.043}{id}%
\TextAt{9}{219.758}{724.043}{\Character{58}}%
\TextAt{16}{226.33}{724.043}{\Character{70}}%
\TextAt{9}{236.11}{724.043}{\Character{40}}%
\TextAt{13}{240.663}{724.043}{U}%
\TextAt{9}{249.862}{724.043}{\Character{41}}%
\NativeMathAt{257.735}{724.043}{11.955}{\rightarrow}%
\TextAt{16}{273.002}{724.043}{\Character{70}}%
\TextAt{9}{282.791}{724.043}{\Character{40}}%
\TextAt{13}{287.343}{724.043}{U}%
\TextAt{9}{296.542}{724.043}{\Character{41}}%
\TextAt{8}{301.095}{724.043}{,}%
\TextAt{8}{87.631}{706.997}{(c)}%
\TextAt{13}{107.918}{706.997}{\Character{26}}%
\TextAt{21}{113.956}{705.204}{U}%
\TextAt{21}{120.523}{705.204}{W}%
\TextAt{9}{133.519}{706.997}{\Character{61}}%
\TextAt{13}{145.944}{706.997}{\Character{26}}%
\TextAt{21}{151.982}{705.204}{V}%
\TextAt{21}{158.789}{705.204}{W}%
\TextAt{16}{171.116}{706.997}{\Character{14}}%
\TextAt{13}{179.75}{706.997}{\Character{26}}%
\TextAt{21}{185.788}{705.204}{U}%
\TextAt{21}{192.355}{705.204}{V}%
\TextAt{8}{199.677}{706.997}{,}%
\TextAt{8}{206.827}{706.997}{for}%
\TextAt{8}{224.712}{706.997}{open}%
\TextAt{8}{252.998}{706.997}{sets}%
\TextAt{13}{275.88}{706.997}{W}%
\NativeMathAt{291.878}{706.997}{11.955}{\subseteq}%
\TextAt{13}{304.497}{706.997}{V}%
\NativeMathAt{317.244}{706.997}{11.955}{\subseteq}%
\TextAt{13}{329.863}{706.997}{U}%
\TextAt{8}{339.062}{706.997}{.}%
\TextAt{8}{97.159}{682.197}{The}%
\TextAt{8}{121.374}{682.197}{elements}%
\TextAt{13}{169.392}{682.197}{s}%
\NativeMathAt{178.502}{682.197}{11.955}{\in}%
\TextAt{16}{190.051}{682.197}{\Character{70}}%
\TextAt{9}{199.848}{682.197}{\Character{40}}%
\TextAt{13}{204.4}{682.197}{U}%
\TextAt{9}{213.599}{682.197}{\Character{41}}%
\TextAt{8}{222.211}{682.197}{are}%
\TextAt{8}{241.883}{682.197}{called}%
\TextAt{8}{275.208}{682.197}{the}%
\TextAt{12}{295.526}{682.197}{sections}%
\TextAt{8}{345.492}{682.197}{of}%
\TextAt{8}{358.984}{682.197}{the}%
\TextAt{8}{379.302}{682.197}{presheaf}%
\TextAt{16}{425.371}{682.197}{\Character{70}}%
\TextAt{8}{439.21}{682.197}{over}%
\TextAt{13}{464.417}{682.197}{U}%
\TextAt{8}{473.616}{682.197}{.}%
\TextAt{8}{482.549}{682.197}{If}%
\TextAt{13}{494.403}{682.197}{V}%
\NativeMathAt{507.424}{682.197}{11.955}{\subseteq}%
\TextAt{13}{520.319}{682.197}{U}%
\TextAt{8}{529.518}{682.197}{,}%
\TextAt{8}{79.226}{665.151}{then}%
\TextAt{8}{106.371}{665.151}{one}%
\TextAt{8}{128.315}{665.151}{usually}%
\TextAt{8}{168.862}{665.151}{writes}%
\TextAt{13}{203.873}{665.151}{\Character{26}}%
\TextAt{21}{209.911}{663.358}{U}%
\TextAt{21}{216.478}{663.358}{V}%
\TextAt{9}{223.801}{665.151}{\Character{40}}%
\TextAt{13}{228.353}{665.151}{s}%
\TextAt{9}{233.867}{665.151}{\Character{41}}%
\TextAt{9}{242.547}{665.151}{\Character{61}}%
\TextAt{13}{255.818}{665.151}{s}%
\TextAt{16}{261.332}{665.151}{\Character{106}}%
\TextAt{21}{264.653}{663.358}{V}%
\TextAt{8}{271.965}{665.151}{.}%
\TextAt{8}{281.868}{665.151}{The}%
\TextAt{8}{306.406}{665.151}{definition}%
\TextAt{8}{358.92}{665.151}{of}%
\TextAt{8}{372.735}{665.151}{a}%
\TextAt{8}{382.975}{665.151}{presheaf}%
\TextAt{8}{429.368}{665.151}{can}%
\TextAt{8}{451.312}{665.151}{be}%
\TextAt{8}{467.721}{665.151}{reformulated}%
\TextAt{8}{79.226}{648.105}{most}%
\TextAt{8}{108.595}{648.105}{concisely}%
\TextAt{8}{158.443}{648.105}{in}%
\TextAt{8}{172.784}{648.105}{the}%
\TextAt{8}{193.628}{648.105}{language}%
\TextAt{8}{243.106}{648.105}{of}%
\TextAt{8}{257.124}{648.105}{categories.}%
\TextAt{8}{317.789}{648.105}{The}%
\TextAt{8}{342.531}{648.105}{open}%
\TextAt{8}{371.505}{648.105}{sets}%
\TextAt{8}{395.076}{648.105}{of}%
\TextAt{8}{409.094}{648.105}{the}%
\TextAt{8}{429.938}{648.105}{topological}%
\TextAt{8}{490.151}{648.105}{space}%
\TextAt{13}{522.114}{648.105}{X}%
\TextAt{8}{79.226}{631.059}{form}%
\TextAt{8}{106.9}{631.059}{a}%
\TextAt{8}{116.688}{631.059}{category}%
\TextAt{13}{163.885}{631.059}{X}%
\TextAt{7}{173.6}{629.266}{top}%
\TextAt{8}{190.255}{631.059}{in}%
\TextAt{8}{203.941}{631.059}{which}%
\TextAt{8}{237.46}{631.059}{only}%
\TextAt{8}{263.185}{631.059}{inclusions}%
\TextAt{8}{316.67}{631.059}{are}%
\TextAt{8}{336.214}{631.059}{admitted}%
\TextAt{8}{386.327}{631.059}{as}%
\TextAt{8}{400.73}{631.059}{morphisms.}%
\TextAt{8}{464.69}{631.059}{A}%
\TextAt{8}{477.396}{631.059}{presheaf}%
\TextAt{8}{523.337}{631.059}{of}%
\TextAt{8}{79.226}{614.013}{abelian}%
\TextAt{8}{119.874}{614.013}{groups}%
\TextAt{8}{157.664}{614.013}{(rings)}%
\TextAt{8}{195.454}{614.013}{is}%
\TextAt{8}{207.218}{614.013}{then}%
\TextAt{8}{233.878}{614.013}{simply}%
\TextAt{8}{271.334}{614.013}{a}%
\TextAt{8}{281.089}{614.013}{contravariant}%
\TextAt{8}{352.665}{614.013}{functor}%
\TextAt{16}{240.362}{589.257}{\Character{70}}%
\TextAt{9}{253.463}{589.257}{\Character{58}}%
\TextAt{13}{260.035}{589.257}{X}%
\TextAt{7}{269.751}{587.464}{top}%
\NativeArrowAt{285.796}{589.257}{11.955}{19.26}{\longrightarrow}%
\TextAt{9}{308.378}{589.257}{\Character{40}}%
\TextAt{13}{312.93}{589.257}{ab}%
\TextAt{9}{324.052}{589.257}{\Character{41}}%
\TextAt{13}{328.604}{589.257}{\Character{59}}%
\TextAt{9}{337.746}{589.257}{\Character{40}}%
\TextAt{8}{342.298}{589.257}{rings}%
\TextAt{9}{367.081}{589.257}{\Character{41}}%
\TextAt{8}{79.226}{564.5}{into}%
\TextAt{8}{102.969}{564.5}{the}%
\TextAt{8}{123.126}{564.5}{category}%
\TextAt{8}{170.289}{564.5}{of}%
\TextAt{8}{183.619}{564.5}{abelian}%
\TextAt{8}{224.266}{564.5}{groups}%
\TextAt{8}{262.057}{564.5}{(resp.}%
\TextAt{8}{295.934}{564.5}{rings)}%
\TextAt{8}{329.169}{564.5}{such}%
\TextAt{8}{355.566}{564.5}{that}%
\TextAt{16}{380.935}{564.5}{\Character{70}}%
\TextAt{9}{390.715}{564.5}{\Character{40}}%
\TextAt{18}{395.267}{564.5}{\Character{63}}%
\TextAt{9}{404.566}{564.5}{\Character{41}}%
\TextAt{9}{412.421}{564.5}{\Character{61}}%
\TextAt{9}{424.842}{564.5}{0}%
\TextAt{8}{430.717}{564.5}{.}%
\TextAt{12}{79.226}{539.701}{Definition}%
\TextAt{12}{142.112}{539.701}{(13.2).}%
\TextAt{8}{184.242}{539.701}{A}%
\TextAt{8}{196.999}{539.701}{presheaf}%
\TextAt{16}{242.991}{539.701}{\Character{70}}%
\TextAt{8}{256.753}{539.701}{on}%
\TextAt{8}{273.096}{539.701}{the}%
\TextAt{8}{293.337}{539.701}{topological}%
\TextAt{8}{352.946}{539.701}{space}%
\TextAt{13}{384.305}{539.701}{X}%
\TextAt{8}{398.942}{539.701}{is}%
\TextAt{8}{410.79}{539.701}{called}%
\TextAt{8}{444.038}{539.701}{a}%
\TextAt{12}{453.878}{539.701}{sheaf}%
\TextAt{8}{488.358}{539.701}{if,}%
\TextAt{8}{502.439}{539.701}{for}%
\TextAt{8}{520.408}{539.701}{all}%
\TextAt{8}{79.226}{522.655}{open}%
\TextAt{8}{107.512}{522.655}{coverings}%
\TextAt{9}{157.987}{522.655}{\Character{40}}%
\TextAt{13}{162.539}{522.655}{U}%
\TextAt{21}{170.489}{520.861}{i}%
\TextAt{9}{173.871}{522.655}{\Character{41}}%
\TextAt{8}{182.32}{522.655}{of}%
\TextAt{8}{195.65}{522.655}{the}%
\TextAt{8}{215.807}{522.655}{open}%
\TextAt{8}{244.093}{522.655}{sets}%
\TextAt{13}{266.975}{522.655}{U}%
\TextAt{8}{276.174}{522.655}{,}%
\TextAt{8}{283.323}{522.655}{one}%
\TextAt{8}{304.783}{522.655}{has:}%
\TextAt{8}{89.579}{502.62}{(i)}%
\TextAt{8}{107.918}{502.62}{If}%
\TextAt{13}{119.611}{502.62}{s\Character{59}}%
\TextAt{13}{130.347}{502.62}{s}%
\TextAt{3}{135.883}{506.958}{\Character{48}}%
\NativeMathAt{141.999}{502.62}{11.955}{\in}%
\TextAt{16}{153.273}{502.62}{\Character{70}}%
\TextAt{9}{163.07}{502.62}{\Character{40}}%
\TextAt{13}{167.622}{502.62}{U}%
\TextAt{9}{176.821}{502.62}{\Character{41}}%
\TextAt{8}{185.271}{502.62}{are}%
\TextAt{8}{204.782}{502.62}{two}%
\TextAt{8}{226.899}{502.62}{sections}%
\TextAt{8}{270.595}{502.62}{such}%
\TextAt{8}{296.992}{502.62}{that}%
\TextAt{13}{322.361}{502.62}{s}%
\TextAt{16}{327.875}{502.62}{\Character{106}}%
\TextAt{21}{331.196}{500.827}{U}%
\TextAt{22}{336.925}{499.612}{i}%
\TextAt{9}{343.905}{502.62}{\Character{61}}%
\TextAt{13}{356.331}{502.62}{s}%
\TextAt{3}{361.845}{506.958}{\Character{48}}%
\TextAt{16}{364.64}{502.62}{\Character{106}}%
\TextAt{21}{367.961}{500.827}{U}%
\TextAt{22}{373.69}{499.612}{i}%
\TextAt{8}{381.247}{502.62}{for}%
\TextAt{8}{399.132}{502.62}{all}%
\TextAt{13}{415.391}{502.62}{i}%
\TextAt{8}{419.384}{502.62}{,}%
\TextAt{8}{426.534}{502.62}{then}%
\TextAt{13}{453.194}{502.62}{s}%
\TextAt{9}{462.028}{502.62}{\Character{61}}%
\TextAt{13}{474.454}{502.62}{s}%
\TextAt{3}{479.968}{506.958}{\Character{48}}%
\TextAt{8}{482.763}{502.62}{.}%
\TextAt{8}{86.327}{485.574}{(ii)}%
\TextAt{8}{107.918}{485.574}{If}%
\TextAt{13}{119.611}{485.574}{s}%
\TextAt{21}{125.125}{483.781}{i}%
\NativeMathAt{131.827}{485.574}{11.955}{\in}%
\TextAt{16}{143.101}{485.574}{\Character{70}}%
\TextAt{9}{152.898}{485.574}{\Character{40}}%
\TextAt{13}{157.45}{485.574}{U}%
\TextAt{21}{165.4}{483.781}{i}%
\TextAt{9}{168.782}{485.574}{\Character{41}}%
\TextAt{8}{177.231}{485.574}{is}%
\TextAt{8}{188.995}{485.574}{a}%
\TextAt{8}{198.751}{485.574}{family}%
\TextAt{8}{234.52}{485.574}{of}%
\TextAt{8}{247.85}{485.574}{sections}%
\TextAt{8}{291.547}{485.574}{such}%
\TextAt{8}{317.944}{485.574}{that}%
\TextAt{13}{275.088}{460.817}{s}%
\TextAt{21}{280.602}{459.024}{i}%
\TextAt{16}{283.983}{460.817}{\Character{106}}%
\TextAt{21}{287.304}{459.024}{U}%
\TextAt{22}{293.033}{457.809}{i}%
\TextAt{3}{296.195}{459.024}{\Character{92}}%
\TextAt{21}{301.84}{459.024}{U}%
\TextAt{22}{307.569}{457.809}{j}%
\TextAt{9}{315.291}{460.817}{\Character{61}}%
\TextAt{13}{327.717}{460.817}{s}%
\TextAt{21}{333.231}{459.024}{j}%
\TextAt{16}{337.613}{460.817}{\Character{106}}%
\TextAt{21}{340.934}{459.024}{U}%
\TextAt{22}{346.663}{457.809}{i}%
\TextAt{3}{349.824}{459.024}{\Character{92}}%
\TextAt{21}{355.47}{459.024}{U}%
\TextAt{22}{361.199}{457.809}{j}%
\TextAt{8}{107.918}{436.061}{for}%
\TextAt{8}{125.803}{436.061}{all}%
\TextAt{13}{142.062}{436.061}{i\Character{59}}%
\TextAt{13}{151.28}{436.061}{j}%
\TextAt{8}{156.816}{436.061}{,}%
\TextAt{8}{163.965}{436.061}{then}%
\TextAt{8}{190.625}{436.061}{there}%
\TextAt{8}{220.537}{436.061}{exists}%
\TextAt{8}{252.851}{436.061}{a}%
\TextAt{8}{262.607}{436.061}{section}%
\TextAt{13}{301.688}{436.061}{s}%
\NativeMathAt{310.523}{436.061}{11.955}{\in}%
\TextAt{16}{321.797}{436.061}{\Character{70}}%
\TextAt{9}{331.594}{436.061}{\Character{40}}%
\TextAt{13}{336.146}{436.061}{U}%
\TextAt{9}{345.345}{436.061}{\Character{41}}%
\TextAt{8}{353.795}{436.061}{such}%
\TextAt{8}{380.192}{436.061}{that}%
\TextAt{13}{405.561}{436.061}{s}%
\TextAt{16}{411.075}{436.061}{\Character{106}}%
\TextAt{21}{414.396}{434.267}{U}%
\TextAt{22}{420.125}{433.052}{i}%
\TextAt{9}{427.106}{436.061}{\Character{61}}%
\TextAt{13}{439.531}{436.061}{s}%
\TextAt{21}{445.045}{434.267}{i}%
\TextAt{8}{452.324}{436.061}{for}%
\TextAt{8}{470.209}{436.061}{all}%
\TextAt{13}{486.468}{436.061}{i}%
\TextAt{8}{490.461}{436.061}{.}%
\TextAt{8}{97.159}{411.261}{The}%
\TextAt{12}{122.139}{411.261}{stalk}%
\TextAt{8}{154.89}{411.261}{of}%
\TextAt{8}{169.146}{411.261}{the}%
\TextAt{8}{190.229}{411.261}{sheaf}%
\TextAt{16}{220.804}{411.261}{\Character{70}}%
\TextAt{8}{235.407}{411.261}{at}%
\TextAt{8}{250.644}{411.261}{the}%
\TextAt{8}{271.726}{411.261}{point}%
\TextAt{13}{303.222}{411.261}{x}%
\NativeMathAt{314.773}{411.261}{11.955}{\in}%
\TextAt{13}{327.642}{411.261}{X}%
\TextAt{8}{343.121}{411.261}{is}%
\TextAt{8}{355.811}{411.261}{defined}%
\TextAt{8}{397.05}{411.261}{to}%
\TextAt{8}{412.287}{411.261}{be}%
\TextAt{8}{429.137}{411.261}{the}%
\TextAt{8}{450.22}{411.261}{direct}%
\TextAt{8}{484.31}{411.261}{limit}%
\TextAt{8}{513.199}{411.261}{(see}%
\TextAt{8}{79.226}{394.215}{chap.}%
\TextAt{8}{110.118}{394.215}{IV,}%
\TextAt{8}{130.263}{394.215}{§2)}%
\TextAt{16}{266.037}{377.169}{\Character{70}}%
\TextAt{21}{274.629}{375.376}{x}%
\TextAt{9}{283.215}{377.169}{\Character{61}}%
\NativeMathAt{296.006}{377.169}{11.955}{\varinjlim}%
\TextAt{21}{295.641}{364.789}{U}%
\NativeMathAt{302.219}{364.789}{7.97}{\ni}%
\TextAt{21}{307.864}{364.789}{x}%
\TextAt{16}{314.624}{377.169}{\Character{70}}%
\TextAt{9}{324.404}{377.169}{\Character{40}}%
\TextAt{13}{328.956}{377.169}{U}%
\TextAt{9}{338.155}{377.169}{\Character{41}}%
\TextAt{13}{342.707}{377.169}{\Character{59}}%
\TextAt{8}{79.226}{348.471}{where}%
\TextAt{13}{112.121}{348.471}{U}%
\TextAt{8}{124.302}{348.471}{varies}%
\TextAt{8}{156.3}{348.471}{over}%
\TextAt{8}{180.431}{348.471}{all}%
\TextAt{8}{195.776}{348.471}{open}%
\TextAt{8}{223.147}{348.471}{neighbourhoods}%
\TextAt{8}{306.851}{348.471}{of}%
\TextAt{13}{319.266}{348.471}{x}%
\TextAt{8}{325.918}{348.471}{.}%
\TextAt{8}{334.061}{348.471}{In}%
\TextAt{8}{347.768}{348.471}{other}%
\TextAt{8}{377.422}{348.471}{words,}%
\TextAt{8}{413.5}{348.471}{two}%
\TextAt{8}{434.703}{348.471}{sections}%
\TextAt{13}{477.484}{348.471}{s}%
\TextAt{21}{482.998}{346.678}{U}%
\NativeMathAt{493.395}{348.471}{11.955}{\in}%
\TextAt{16}{504.669}{348.471}{\Character{70}}%
\TextAt{9}{514.466}{348.471}{\Character{40}}%
\TextAt{13}{519.018}{348.471}{U}%
\TextAt{9}{528.217}{348.471}{\Character{41}}%
\TextAt{8}{79.226}{331.425}{and}%
\TextAt{13}{102.878}{331.425}{s}%
\TextAt{21}{108.392}{329.632}{V}%
\NativeMathAt{120.541}{331.425}{11.955}{\in}%
\TextAt{16}{133.333}{331.425}{\Character{70}}%
\TextAt{9}{143.128}{331.425}{\Character{40}}%
\TextAt{13}{147.68}{331.425}{V}%
\TextAt{9}{157.106}{331.425}{\Character{41}}%
\TextAt{8}{166.445}{331.425}{are}%
\TextAt{8}{186.846}{331.425}{called}%
\TextAt{8}{220.899}{331.425}{equivalent}%
\TextAt{8}{277.404}{331.425}{in}%
\TextAt{8}{291.947}{331.425}{the}%
\TextAt{8}{312.993}{331.425}{disjoint}%
\TextAt{8}{355.57}{331.425}{union}%
\TextAt{20}{388.978}{340.392}{\Character{83}}%
\TextAt{21}{398.94}{327.938}{U}%
\NativeMathAt{405.518}{327.938}{7.97}{\ni}%
\TextAt{21}{411.164}{327.938}{x}%
\TextAt{16}{418.421}{331.425}{\Character{70}}%
\TextAt{9}{428.201}{331.425}{\Character{40}}%
\TextAt{13}{432.754}{331.425}{U}%
\TextAt{9}{441.953}{331.425}{\Character{41}}%
\TextAt{8}{451.292}{331.425}{if}%
\TextAt{8}{462.906}{331.425}{there}%
\TextAt{8}{493.707}{331.425}{exists}%
\TextAt{8}{526.912}{331.425}{a}%
\TextAt{8}{79.226}{314.379}{neighbourhood}%
\TextAt{13}{160.181}{314.379}{W}%
\NativeMathAt{177.799}{314.379}{11.955}{\subseteq}%
\TextAt{13}{192.039}{314.379}{U}%
\TextAt{16}{204.543}{314.379}{\Character{92}}%
\TextAt{13}{215.818}{314.379}{V}%
\TextAt{8}{230.092}{314.379}{of}%
\TextAt{13}{244.373}{314.379}{x}%
\TextAt{8}{255.873}{314.379}{such}%
\TextAt{8}{283.221}{314.379}{that}%
\TextAt{13}{309.54}{314.379}{s}%
\TextAt{21}{315.054}{312.586}{U}%
\TextAt{16}{322.13}{314.379}{\Character{106}}%
\TextAt{21}{325.451}{312.586}{W}%
\TextAt{9}{340.056}{314.379}{\Character{61}}%
\TextAt{13}{354.102}{314.379}{s}%
\TextAt{21}{359.616}{312.586}{V}%
\TextAt{16}{366.928}{314.379}{\Character{106}}%
\TextAt{21}{370.249}{312.586}{W}%
\TextAt{8}{379.912}{314.379}{.}%
\TextAt{8}{391.213}{314.379}{The}%
\TextAt{8}{416.218}{314.379}{equivalence}%
\TextAt{8}{478.953}{314.379}{classes}%
\TextAt{8}{517.156}{314.379}{are}%
\TextAt{8}{79.226}{297.333}{called}%
\TextAt{12}{112.39}{297.333}{germs}%
\TextAt{12}{151.758}{297.333}{of}%
\TextAt{12}{167.085}{297.333}{sections}%
\TextAt{8}{216.89}{297.333}{at}%
\TextAt{13}{231.2}{297.333}{x}%
\TextAt{8}{237.852}{297.333}{.}%
\TextAt{8}{246.301}{297.333}{They}%
\TextAt{8}{276.535}{297.333}{are}%
\TextAt{8}{296.046}{297.333}{the}%
\TextAt{8}{316.203}{297.333}{elements}%
\TextAt{8}{364.059}{297.333}{of}%
\TextAt{16}{377.389}{297.333}{\Character{70}}%
\TextAt{21}{385.982}{295.54}{x}%
\TextAt{8}{391.247}{297.333}{.}%
\TextAt{8}{97.159}{277.298}{We}%
\TextAt{8}{117.303}{277.298}{now}%
\TextAt{8}{141.692}{277.298}{return}%
\TextAt{8}{177.462}{277.298}{to}%
\TextAt{8}{191.772}{277.298}{the}%
\TextAt{8}{211.929}{277.298}{spectrum}%
\TextAt{13}{263.037}{277.298}{X}%
\TextAt{9}{277.013}{277.298}{\Character{61}}%
\TextAt{8}{289.438}{277.298}{Spec}%
\TextAt{9}{312.847}{277.298}{\Character{40}}%
\TextAt{14}{317.399}{277.298}{o}%
\TextAt{9}{323.243}{277.298}{\Character{41}}%
\TextAt{8}{331.693}{277.298}{of}%
\TextAt{8}{345.023}{277.298}{a}%
\TextAt{8}{354.778}{277.298}{ring}%
\TextAt{14}{378.844}{277.298}{o}%
\TextAt{8}{388.586}{277.298}{and}%
\TextAt{8}{411.348}{277.298}{obtain}%
\TextAt{8}{447.776}{277.298}{the}%
\TextAt{12}{79.226}{252.499}{Proposition}%
\TextAt{12}{152.154}{252.499}{(13.3).}%
\TextAt{8}{194.388}{252.499}{The}%
\TextAt{8}{218.63}{252.499}{rings}%
\TextAt{14}{247.499}{252.499}{o}%
\TextAt{9}{253.343}{252.499}{\Character{40}}%
\TextAt{13}{257.895}{252.499}{U}%
\TextAt{9}{267.094}{252.499}{\Character{41}}%
\TextAt{8}{271.647}{252.499}{,}%
\TextAt{8}{279.031}{252.499}{together}%
\TextAt{8}{325.402}{252.499}{with}%
\TextAt{8}{352.25}{252.499}{the}%
\TextAt{8}{372.595}{252.499}{restriction}%
\TextAt{8}{428.781}{252.499}{mappings}%
\TextAt{13}{481.715}{252.499}{\Character{26}}%
\TextAt{21}{487.753}{250.706}{U}%
\TextAt{21}{494.32}{250.706}{V}%
\TextAt{8}{501.642}{252.499}{,}%
\TextAt{8}{509.027}{252.499}{form}%
\TextAt{8}{79.226}{235.453}{a}%
\TextAt{8}{89.434}{235.453}{sheaf}%
\TextAt{8}{119.536}{235.453}{on}%
\TextAt{13}{136.248}{235.453}{X}%
\TextAt{8}{146.903}{235.453}{.}%
\TextAt{8}{156.71}{235.453}{It}%
\TextAt{8}{169.835}{235.453}{is}%
\TextAt{8}{182.052}{235.453}{denoted}%
\TextAt{8}{226.727}{235.453}{by}%
\TextAt{14}{243.439}{235.453}{o}%
\TextAt{21}{249.283}{233.66}{X}%
\TextAt{8}{261.718}{235.453}{and}%
\TextAt{8}{284.934}{235.453}{called}%
\TextAt{8}{318.55}{235.453}{the}%
\TextAt{12}{339.16}{235.453}{structure}%
\TextAt{12}{397.998}{235.453}{sheaf}%
\TextAt{8}{432.846}{235.453}{on}%
\TextAt{13}{449.558}{235.453}{X}%
\TextAt{8}{460.213}{235.453}{.}%
\TextAt{8}{470.02}{235.453}{The}%
\TextAt{8}{494.526}{235.453}{stalk}%
\TextAt{8}{523.337}{235.453}{of}%
\TextAt{14}{79.226}{218.407}{o}%
\TextAt{21}{85.07}{216.614}{X}%
\TextAt{8}{97.053}{218.407}{at}%
\TextAt{8}{111.363}{218.407}{the}%
\TextAt{8}{131.52}{218.407}{point}%
\TextAt{14}{162.089}{218.407}{p}%
\NativeMathAt{171.393}{218.407}{11.955}{\in}%
\TextAt{13}{182.684}{218.407}{X}%
\TextAt{8}{197.236}{218.407}{is}%
\TextAt{8}{209}{218.407}{the}%
\TextAt{8}{229.157}{218.407}{localization}%
\TextAt{14}{291.276}{218.407}{o}%
\TextAt{27}{297.12}{216.614}{p}%
\TextAt{8}{301.607}{218.407}{,}%
\TextAt{8}{308.756}{218.407}{i.e.,}%
\TextAt{14}{330.861}{218.407}{o}%
\TextAt{21}{336.705}{216.614}{X\Character{59}}%
\TextAt{27}{346.174}{216.614}{p}%
\NativeMathAt{353.982}{218.407}{11.955}{\cong}%
\TextAt{14}{366.601}{218.407}{o}%
\TextAt{27}{372.445}{216.614}{p}%
\TextAt{8}{376.932}{218.407}{.}%
\TextAt{8}{97.159}{193.607}{The}%
\TextAt{8}{121.213}{193.607}{proof}%
\TextAt{8}{151.782}{193.607}{of}%
\TextAt{8}{165.112}{193.607}{this}%
\TextAt{8}{187.935}{193.607}{proposition}%
\TextAt{8}{249.468}{193.607}{follows}%
\TextAt{8}{287.904}{193.607}{immediately}%
\TextAt{8}{354.566}{193.607}{from}%
\TextAt{8}{382.206}{193.607}{the}%
\TextAt{8}{402.362}{193.607}{definitions.}%
\TextAt{8}{97.159}{173.573}{The}%
\TextAt{8}{120.927}{173.573}{couple}%
\TextAt{9}{157.057}{173.573}{\Character{40}}%
\TextAt{13}{161.609}{173.572}{X}%
\TextAt{13}{171.604}{173.572}{\Character{59}}%
\TextAt{14}{176.858}{173.573}{o}%
\TextAt{21}{182.702}{171.779}{X}%
\TextAt{9}{190.787}{173.572}{\Character{41}}%
\TextAt{8}{198.951}{173.573}{is}%
\TextAt{8}{210.43}{173.573}{called}%
\TextAt{8}{243.307}{173.573}{an}%
\TextAt{12}{259.281}{173.573}{affine}%
\TextAt{12}{294.794}{173.573}{scheme}%
\TextAt{8}{336.661}{173.573}{.}%
\TextAt{8}{345.014}{173.573}{Usually,}%
\TextAt{8}{389.39}{173.573}{however,}%
\TextAt{8}{437.293}{173.573}{the}%
\TextAt{8}{457.164}{173.573}{structure}%
\TextAt{8}{507.018}{173.573}{sheaf}%
\TextAt{14}{79.226}{156.527}{o}%
\TextAt{21}{85.07}{154.733}{X}%
\TextAt{8}{97.053}{156.527}{is}%
\TextAt{8}{108.817}{156.527}{dropped}%
\TextAt{8}{154.665}{156.527}{from}%
\TextAt{8}{182.305}{156.527}{the}%
\TextAt{8}{202.462}{156.527}{notation.}%
\TextAt{8}{253.853}{156.527}{Now}%
\TextAt{8}{280.513}{156.527}{let}%
\TextAt{13}{277.024}{131.77}{\Character{39}}%
\TextAt{9}{288.013}{131.77}{\Character{58}}%
\TextAt{14}{294.585}{131.77}{o}%
\NativeArrowAt{303.751}{131.77}{11.955}{19.26}{\longrightarrow}%
\TextAt{14}{326.333}{131.77}{o}%
\TextAt{3}{332.177}{136.706}{\Character{48}}%
\TextAt{8}{79.226}{107.013}{be}%
\TextAt{8}{95.184}{107.013}{a}%
\TextAt{8}{104.974}{107.013}{homomorphism}%
\TextAt{8}{187.678}{107.013}{of}%
\TextAt{8}{201.042}{107.013}{rings}%
\TextAt{8}{229.756}{107.013}{and}%
\TextAt{13}{252.552}{107.013}{X}%
\TextAt{9}{266.586}{107.013}{\Character{61}}%
\TextAt{8}{279.069}{107.013}{Spec}%
\TextAt{9}{302.478}{107.013}{\Character{40}}%
\TextAt{14}{307.03}{107.013}{o}%
\TextAt{9}{312.874}{107.013}{\Character{41}}%
\TextAt{8}{317.427}{107.013}{,}%
\TextAt{13}{324.618}{107.013}{X}%
\TextAt{3}{335.273}{111.352}{\Character{48}}%
\TextAt{9}{341.447}{107.013}{\Character{61}}%
\TextAt{8}{353.93}{107.013}{Spec}%
\TextAt{9}{377.338}{107.013}{\Character{40}}%
\TextAt{14}{381.891}{107.013}{o}%
\TextAt{3}{387.735}{111.352}{\Character{48}}%
\TextAt{9}{390.53}{107.013}{\Character{41}}%
\TextAt{8}{395.082}{107.013}{.}%
\TextAt{8}{403.632}{107.013}{Then}%
\TextAt{13}{434.224}{107.013}{\Character{39}}%
\TextAt{8}{445.823}{107.013}{induces}%
\TextAt{8}{487.533}{107.013}{a}%
\TextAt{8}{497.322}{107.013}{contin-}%
\TextAt{8}{79.226}{89.967}{uous}%
\TextAt{8}{106.603}{89.967}{map}%
\TextAt{13}{216.247}{72.921}{f}%
\TextAt{9}{226.615}{72.921}{\Character{58}}%
\TextAt{13}{233.187}{72.921}{X}%
\TextAt{3}{243.842}{77.858}{\Character{48}}%
\NativeArrowAt{249.958}{72.921}{11.955}{19.26}{\longrightarrow}%
\TextAt{13}{272.54}{72.921}{X}%
\TextAt{13}{282.535}{72.921}{\Character{59}}%
\TextAt{13}{311.669}{72.921}{f}%
\TextAt{9}{318.746}{72.921}{\Character{40}}%
\TextAt{14}{323.298}{72.921}{p}%
\TextAt{3}{329.281}{77.858}{\Character{48}}%
\TextAt{9}{332.076}{72.921}{\Character{41}}%
\TextAt{9}{339.931}{72.921}{\Character{58}\Character{61}}%
\TextAt{13}{355.627}{72.921}{\Character{39}}%
\TextAt{3}{363.295}{77.858}{\Character{0}}%
\TextAt{17}{369.882}{77.858}{1}%
\TextAt{9}{374.614}{72.921}{\Character{40}}%
\TextAt{14}{379.166}{72.921}{p}%
\TextAt{3}{385.149}{77.858}{\Character{48}}%
\TextAt{9}{387.944}{72.921}{\Character{41}}%
\TextAt{13}{392.497}{72.921}{\Character{59}}%
\TextAt{8}{300.14}{43.033}{82}%
\Destination{79.23}{755.685}{nav.Item.17}%
\Destination{79.23}{738.085}{nav.Item.18}%
\Destination{79.23}{721.045}{nav.Item.19}%
\Destination{79.23}{516.675}{nav.Item.20}%
\Destination{79.23}{499.615}{nav.Item.21}%
\Destination{78.23}{809.815}{nav.page.82}%
\Destination{79.36}{554.655}{nav.print.88}%
\Destination{79.36}{553.655}{nav.number.1.13.2}%
\Destination{79.36}{266.455}{nav.number.1.13.3}%
\Hotlink{109.618}{392.055}{13.995}{11.047}{nav.chapter.4}%
\Hotlink{135.442}{392.055}{6.858}{11.047}{nav.section.4.2}%
}

% Page 95 (body page 83)
\Page{%
\TextAt{2}{79.226}{796.943}{I.}%
\TextAt{2}{89.83}{796.943}{Algebraic}%
\TextAt{2}{138.648}{796.943}{Integers}%
\Vector{q 0 G 0 g 0.299 w 7.226 18.729 m 460.769 18.729 l S Q}%
\TextAt{8}{79.226}{761.124}{and,}%
\TextAt{8}{105.241}{761.124}{for}%
\TextAt{8}{123.126}{761.124}{every}%
\TextAt{8}{154.018}{761.124}{open}%
\TextAt{8}{182.304}{761.124}{subset}%
\TextAt{13}{218.193}{761.124}{U}%
\TextAt{8}{231.289}{761.124}{of}%
\TextAt{13}{244.619}{761.124}{X}%
\TextAt{8}{255.275}{761.124}{,}%
\TextAt{8}{262.424}{761.124}{a}%
\TextAt{8}{272.179}{761.124}{homomorphism}%
\TextAt{13}{205.753}{732.123}{f}%
\TextAt{3}{212.8}{737.059}{\Character{3}}%
\TextAt{21}{211.523}{729.167}{U}%
\TextAt{9}{221.92}{732.123}{\Character{58}}%
\TextAt{14}{228.493}{732.123}{o}%
\TextAt{9}{234.337}{732.123}{\Character{40}}%
\TextAt{13}{238.889}{732.123}{U}%
\TextAt{9}{248.088}{732.123}{\Character{41}}%
\NativeArrowAt{255.962}{732.123}{11.955}{19.26}{\longrightarrow}%
\TextAt{14}{278.544}{732.123}{o}%
\TextAt{9}{284.388}{732.123}{\Character{40}}%
\TextAt{13}{288.94}{732.123}{U}%
\TextAt{3}{298.139}{737.059}{\Character{48}}%
\TextAt{9}{300.934}{732.123}{\Character{41}}%
\TextAt{13}{305.487}{732.123}{\Character{59}}%
\TextAt{13}{334.622}{732.123}{s}%
\NativeArrowAt{343.476}{732.123}{11.955}{19.26}{\longmapsto}%
\TextAt{13}{366.058}{732.123}{s}%
\TextAt{16}{374.229}{732.123}{\Character{14}}%
\TextAt{13}{382.863}{732.123}{f}%
\TextAt{16}{389.909}{732.123}{\Character{106}}%
\TextAt{21}{393.23}{730.329}{U}%
\TextAt{23}{399.808}{732.599}{\Character{48}}%
\TextAt{13}{402.991}{732.123}{\Character{59}}%
\TextAt{8}{79.226}{703.121}{where}%
\TextAt{13}{113.035}{703.121}{U}%
\TextAt{3}{122.234}{707.46}{\Character{48}}%
\TextAt{9}{128.35}{703.121}{\Character{61}}%
\TextAt{13}{140.776}{703.121}{f}%
\TextAt{3}{147.822}{707.46}{\Character{0}}%
\TextAt{17}{154.409}{707.46}{1}%
\TextAt{9}{159.141}{703.121}{\Character{40}}%
\TextAt{13}{163.693}{703.121}{U}%
\TextAt{9}{172.892}{703.121}{\Character{41}}%
\TextAt{8}{177.445}{703.121}{.}%
\TextAt{8}{185.893}{703.121}{The}%
\TextAt{8}{209.947}{703.121}{maps}%
\TextAt{13}{240.576}{703.121}{f}%
\TextAt{3}{247.622}{707.46}{\Character{3}}%
\TextAt{21}{246.346}{700.101}{U}%
\TextAt{8}{257.319}{703.121}{have}%
\TextAt{8}{284.314}{703.121}{the}%
\TextAt{8}{304.471}{703.121}{following}%
\TextAt{8}{353.905}{703.121}{two}%
\TextAt{8}{376.022}{703.121}{properties.}%
\TextAt{8}{91.528}{683.087}{a)}%
\TextAt{8}{107.918}{683.087}{If}%
\TextAt{13}{119.611}{683.087}{V}%
\NativeMathAt{132.357}{683.087}{11.955}{\subseteq}%
\TextAt{13}{144.976}{683.087}{U}%
\TextAt{8}{158.073}{683.087}{are}%
\TextAt{8}{177.584}{683.087}{open}%
\TextAt{8}{205.87}{683.087}{sets,}%
\TextAt{8}{232.004}{683.087}{then}%
\TextAt{8}{258.664}{683.087}{the}%
\TextAt{8}{278.82}{683.087}{diagram}%
\TextAt{14}{273.2}{650.453}{o}%
\TextAt{9}{279.045}{650.453}{\Character{40}}%
\TextAt{13}{283.597}{650.453}{U}%
\TextAt{9}{292.796}{650.453}{\Character{41}}%
\TextAt{14}{336.572}{650.453}{o}%
\TextAt{9}{342.416}{650.453}{\Character{40}}%
\TextAt{13}{346.969}{650.453}{U}%
\TextAt{3}{356.168}{654.792}{\Character{48}}%
\TextAt{9}{358.963}{650.453}{\Character{41}}%
\TextAt{14}{273.087}{608.219}{o}%
\TextAt{9}{278.932}{608.219}{\Character{40}}%
\TextAt{13}{283.484}{608.219}{V}%
\TextAt{9}{292.91}{608.219}{\Character{41}}%
\TextAt{14}{336.459}{608.219}{o}%
\TextAt{9}{342.303}{608.219}{\Character{40}}%
\TextAt{13}{346.856}{608.219}{V}%
\TextAt{3}{356.281}{612.558}{\Character{48}}%
\TextAt{9}{359.076}{608.219}{\Character{41}}%
\NativeDiagram{%
\path[draw=black,line width=0.478bp] (230.7,-120.558) -- (258.743,-120.558);%
\path[draw=black,line width=0.478bp,line cap=round,line join=round] (256.495,-117.689) .. controls (256.949,-119.41) and (257.962,-120.223) .. (258.982,-120.558) .. controls (257.962,-120.892) and (256.949,-121.705) .. (256.495,-123.427);%
}%
\TextAt{21}{311.85}{658.964}{f}%
\TextAt{23}{316.797}{661.777}{\Character{3}}%
\TextAt{22}{315.919}{656.258}{U}%
\NativeDiagram{%
\path[draw=black,line width=0.478bp] (213.274,-131.114) -- (213.274,-151.757);%
\path[draw=black,line width=0.478bp,line cap=round,line join=round] (216.143,-149.51) .. controls (214.422,-149.964) and (213.609,-150.977) .. (213.274,-151.996) .. controls (212.939,-150.977) and (212.126,-149.964) .. (210.405,-149.51);%
}%
\TextAt{21}{266.058}{631.384}{\Character{26}}%
\TextAt{22}{270.395}{630.045}{U}%
\TextAt{22}{276.044}{630.045}{V}%
\NativeDiagram{%
\path[draw=black,line width=0.478bp] (278.045,-131.114) -- (278.045,-151.757);%
\path[draw=black,line width=0.478bp,line cap=round,line join=round] (280.914,-149.51) .. controls (279.192,-149.964) and (278.379,-150.977) .. (278.045,-151.996) .. controls (277.71,-150.977) and (276.897,-149.964) .. (275.175,-149.51);%
}%
\TextAt{21}{352.86}{631.755}{\Character{26}}%
\TextAt{22}{357.197}{629.464}{U}%
\TextAt{23}{362.852}{631.179}{\Character{48}}%
\TextAt{22}{365.536}{629.464}{V}%
\TextAt{23}{371.447}{631.179}{\Character{48}}%
\NativeDiagram{%
\path[draw=black,line width=0.478bp] (230.813,-162.792) -- (258.629,-162.792);%
\path[draw=black,line width=0.478bp,line cap=round,line join=round] (256.382,-159.923) .. controls (256.836,-161.645) and (257.848,-162.458) .. (258.868,-162.792) .. controls (257.848,-163.127) and (256.836,-163.94) .. (256.382,-165.661);%
}%
\TextAt{21}{311.722}{602.798}{f}%
\TextAt{23}{316.669}{605.611}{\Character{3}}%
\TextAt{22}{315.791}{600.092}{V}%
\TextAt{8}{107.918}{576.758}{is}%
\TextAt{8}{119.682}{576.758}{commutative.}%
\TextAt{8}{90.882}{559.712}{b)}%
\TextAt{8}{107.918}{559.712}{For}%
\TextAt{14}{128.888}{559.712}{p}%
\TextAt{3}{134.871}{564.05}{\Character{48}}%
\NativeMathAt{140.987}{559.712}{11.955}{\in}%
\TextAt{13}{152.278}{559.712}{U}%
\TextAt{3}{161.477}{564.05}{\Character{48}}%
\NativeMathAt{167.593}{559.712}{11.955}{\subseteq}%
\TextAt{13}{180.212}{559.712}{X}%
\TextAt{3}{190.867}{564.05}{\Character{48}}%
\TextAt{8}{197.56}{559.712}{and}%
\TextAt{13}{220.322}{559.712}{a}%
\NativeMathAt{229.788}{559.712}{11.955}{\in}%
\TextAt{14}{241.079}{559.712}{o}%
\TextAt{9}{246.923}{559.712}{\Character{40}}%
\TextAt{13}{251.476}{559.712}{U}%
\TextAt{9}{260.675}{559.712}{\Character{41}}%
\TextAt{8}{269.124}{559.712}{one}%
\TextAt{8}{290.584}{559.712}{has}%
\TextAt{13}{229.231}{530.71}{a}%
\TextAt{9}{235.375}{530.711}{\Character{40}}%
\TextAt{13}{239.928}{530.71}{f}%
\TextAt{9}{246.974}{530.711}{\Character{40}}%
\TextAt{14}{251.527}{530.71}{p}%
\TextAt{3}{257.51}{535.647}{\Character{48}}%
\TextAt{9}{260.305}{530.71}{\Character{41}\Character{41}}%
\TextAt{9}{272.702}{530.71}{\Character{61}}%
\TextAt{9}{285.124}{530.71}{0}%
\NativeArrowAt{306.249}{530.71}{11.955}{19.103}{\Longrightarrow}%
\TextAt{13}{340.628}{530.71}{f}%
\TextAt{3}{347.675}{535.647}{\Character{3}}%
\TextAt{21}{346.398}{527.755}{U}%
\TextAt{9}{353.474}{530.711}{\Character{40}}%
\TextAt{13}{358.026}{530.71}{a}%
\TextAt{9}{364.171}{530.711}{\Character{41}\Character{40}}%
\TextAt{14}{373.276}{530.71}{p}%
\TextAt{3}{379.259}{535.647}{\Character{48}}%
\TextAt{9}{382.054}{530.71}{\Character{41}}%
\TextAt{9}{389.909}{530.71}{\Character{61}}%
\TextAt{9}{402.33}{530.71}{0}%
\TextAt{13}{408.206}{530.711}{\Character{58}}%
\TextAt{8}{79.226}{501.709}{A}%
\TextAt{8}{92.426}{501.709}{continuous}%
\TextAt{8}{151.557}{501.709}{map}%
\TextAt{13}{178.099}{501.709}{f}%
\TextAt{9}{189.365}{501.709}{\Character{58}}%
\TextAt{13}{196.837}{501.709}{X}%
\TextAt{3}{207.492}{506.048}{\Character{48}}%
\NativeMathAt{214.507}{501.709}{11.955}{\rightarrow}%
\TextAt{13}{230.682}{501.709}{X}%
\TextAt{8}{245.762}{501.709}{together}%
\TextAt{8}{292.472}{501.709}{with}%
\TextAt{8}{319.66}{501.709}{a}%
\TextAt{8}{329.942}{501.709}{family}%
\TextAt{8}{366.24}{501.709}{of}%
\TextAt{8}{380.097}{501.709}{homomorphisms}%
\TextAt{13}{467.909}{501.709}{f}%
\TextAt{3}{474.956}{506.048}{\Character{3}}%
\TextAt{21}{473.679}{498.689}{U}%
\TextAt{9}{484.975}{501.709}{\Character{58}}%
\TextAt{14}{492.447}{501.709}{o}%
\TextAt{9}{498.291}{501.709}{\Character{40}}%
\TextAt{13}{502.843}{501.709}{U}%
\TextAt{9}{512.042}{501.709}{\Character{41}}%
\NativeMathAt{520.814}{501.709}{11.955}{\rightarrow}%
\TextAt{14}{79.226}{484.663}{o}%
\TextAt{9}{85.07}{484.663}{\Character{40}}%
\TextAt{13}{89.623}{484.663}{U}%
\TextAt{3}{98.822}{489.002}{\Character{48}}%
\TextAt{9}{101.617}{484.663}{\Character{41}}%
\TextAt{8}{110.853}{484.663}{which}%
\TextAt{8}{145.127}{484.663}{satisfy}%
\TextAt{8}{182.461}{484.663}{conditions}%
\TextAt{8}{239.245}{484.663}{a)}%
\TextAt{8}{254.343}{484.663}{and}%
\TextAt{8}{277.892}{484.663}{b)}%
\TextAt{8}{293.635}{484.663}{is}%
\TextAt{8}{306.186}{484.663}{called}%
\TextAt{8}{340.136}{484.663}{a}%
\TextAt{12}{350.679}{484.663}{morphism}%
\TextAt{8}{414.015}{484.663}{from}%
\TextAt{8}{442.442}{484.663}{the}%
\TextAt{8}{463.385}{484.663}{scheme}%
\TextAt{13}{504.222}{484.663}{X}%
\TextAt{3}{514.877}{489.002}{\Character{48}}%
\TextAt{8}{522.357}{484.663}{to}%
\TextAt{8}{79.226}{467.617}{the}%
\TextAt{8}{99.351}{467.617}{scheme}%
\TextAt{13}{139.369}{467.617}{X}%
\TextAt{8}{150.024}{467.617}{.}%
\TextAt{8}{158.462}{467.617}{When}%
\TextAt{8}{192.562}{467.617}{referring}%
\TextAt{8}{239.682}{467.617}{to}%
\TextAt{8}{253.96}{467.617}{such}%
\TextAt{8}{280.326}{467.617}{a}%
\TextAt{8}{290.05}{467.617}{morphism,}%
\TextAt{8}{347.971}{467.617}{the}%
\TextAt{8}{368.096}{467.617}{maps}%
\TextAt{13}{398.693}{467.617}{f}%
\TextAt{3}{405.74}{471.956}{\Character{3}}%
\TextAt{21}{404.463}{464.597}{U}%
\TextAt{8}{415.405}{467.617}{are}%
\TextAt{8}{434.884}{467.617}{usually}%
\TextAt{8}{474.914}{467.617}{not}%
\TextAt{8}{495.697}{467.617}{written}%
\TextAt{8}{79.226}{450.571}{explicitly.}%
\TextAt{8}{133.509}{450.571}{One}%
\TextAt{8}{158.18}{450.571}{can}%
\TextAt{8}{179.599}{450.571}{show}%
\TextAt{8}{208.562}{450.571}{that}%
\TextAt{8}{233.89}{450.571}{every}%
\TextAt{8}{264.742}{450.571}{morphism}%
\TextAt{8}{319.396}{450.571}{between}%
\TextAt{8}{364.546}{450.571}{two}%
\TextAt{8}{386.622}{450.571}{affine}%
\TextAt{8}{417.797}{450.571}{schemes}%
\TextAt{13}{462.421}{450.571}{X}%
\TextAt{3}{473.076}{454.91}{\Character{48}}%
\TextAt{9}{479.192}{450.571}{\Character{61}}%
\TextAt{8}{491.617}{450.571}{Spec}%
\TextAt{9}{515.026}{450.571}{\Character{40}}%
\TextAt{14}{519.578}{450.571}{o}%
\TextAt{3}{525.422}{454.91}{\Character{48}}%
\TextAt{9}{528.217}{450.571}{\Character{41}}%
\TextAt{8}{79.226}{433.525}{and}%
\TextAt{13}{101.441}{433.525}{X}%
\TextAt{9}{115.417}{433.525}{\Character{61}}%
\TextAt{8}{127.843}{433.525}{Spec}%
\TextAt{9}{151.251}{433.525}{\Character{40}}%
\TextAt{14}{155.803}{433.525}{o}%
\TextAt{9}{161.647}{433.525}{\Character{41}}%
\TextAt{8}{169.55}{433.525}{is}%
\TextAt{8}{180.766}{433.525}{induced}%
\TextAt{8}{223.783}{433.525}{in}%
\TextAt{8}{236.888}{433.525}{the}%
\TextAt{8}{256.497}{433.525}{way}%
\TextAt{8}{279.693}{433.525}{described}%
\TextAt{8}{330.899}{433.525}{above}%
\TextAt{8}{363.527}{433.525}{by}%
\TextAt{8}{379.239}{433.525}{a}%
\TextAt{8}{388.447}{433.525}{ring}%
\TextAt{8}{411.965}{433.525}{homomorphism}%
\TextAt{13}{494.088}{433.525}{\Character{39}}%
\TextAt{9}{505.077}{433.525}{\Character{58}}%
\TextAt{14}{511.649}{433.525}{o}%
\NativeMathAt{520.814}{433.525}{11.955}{\rightarrow}%
\TextAt{14}{79.226}{416.479}{o}%
\TextAt{3}{85.07}{420.818}{\Character{48}}%
\TextAt{8}{87.865}{416.48}{.}%
\TextAt{8}{96.163}{416.48}{The}%
\TextAt{8}{119.765}{416.48}{proofs}%
\TextAt{8}{154.497}{416.48}{of}%
\TextAt{8}{167.374}{416.48}{the}%
\TextAt{8}{187.079}{416.48}{above}%
\TextAt{8}{219.802}{416.48}{claims}%
\TextAt{8}{255.18}{416.48}{are}%
\TextAt{8}{274.238}{416.48}{easy,}%
\TextAt{8}{301.899}{416.48}{although}%
\TextAt{8}{350.236}{416.48}{some}%
\TextAt{8}{379.11}{416.48}{of}%
\TextAt{8}{391.988}{416.48}{them}%
\TextAt{8}{421.448}{416.48}{are}%
\TextAt{8}{440.506}{416.48}{a}%
\TextAt{8}{449.81}{416.48}{bit}%
\TextAt{8}{467.565}{416.48}{lengthy.}%
\TextAt{8}{512.613}{416.48}{The}%
\TextAt{8}{79.226}{399.434}{notion}%
\TextAt{8}{115.586}{399.434}{of}%
\TextAt{8}{128.849}{399.434}{scheme}%
\TextAt{8}{168.832}{399.434}{is}%
\TextAt{8}{180.528}{399.434}{the}%
\TextAt{8}{200.617}{399.434}{basis}%
\TextAt{8}{229.29}{399.434}{of}%
\TextAt{8}{242.553}{399.434}{a}%
\TextAt{8}{252.241}{399.434}{very}%
\TextAt{8}{277.866}{399.434}{extensive}%
\TextAt{8}{328.261}{399.434}{theory}%
\TextAt{8}{364.944}{399.434}{which}%
\TextAt{8}{398.363}{399.434}{occupies}%
\TextAt{8}{444.85}{399.434}{a}%
\TextAt{8}{454.538}{399.434}{central}%
\TextAt{8}{493.169}{399.434}{place}%
\TextAt{8}{523.014}{399.434}{in}%
\TextAt{8}{79.226}{382.388}{mathematics.}%
\TextAt{8}{152.577}{382.388}{As}%
\TextAt{8}{169.248}{382.388}{introductions}%
\TextAt{8}{240.243}{382.388}{into}%
\TextAt{8}{263.37}{382.388}{this}%
\TextAt{8}{285.576}{382.388}{important}%
\TextAt{8}{340.252}{382.388}{discipline}%
\TextAt{8}{391.067}{382.388}{let}%
\TextAt{8}{407.355}{382.388}{us}%
\TextAt{8}{421.755}{382.388}{recommend}%
\TextAt{8}{483.569}{382.388}{the}%
\TextAt{8}{503.109}{382.388}{books}%
\TextAt{8}{79.226}{365.342}{[51]}%
\TextAt{8}{101.343}{365.342}{and}%
\TextAt{8}{124.106}{365.342}{[104].}%
\TextAt{8}{97.159}{345.307}{We}%
\TextAt{8}{117.22}{345.307}{will}%
\TextAt{8}{139.241}{345.307}{now}%
\TextAt{8}{163.546}{345.307}{confine}%
\TextAt{8}{203.13}{345.307}{ourselves}%
\TextAt{8}{252.601}{345.307}{to}%
\TextAt{8}{266.828}{345.307}{considering}%
\TextAt{8}{327.943}{345.307}{noetherian}%
\TextAt{8}{386.069}{345.307}{integral}%
\TextAt{8}{428.593}{345.307}{domains}%
\TextAt{14}{474.753}{345.307}{o}%
\TextAt{8}{484.411}{345.307}{of}%
\TextAt{8}{497.657}{345.307}{dimen-}%
\TextAt{8}{79.226}{328.261}{sion}%
\NativeMathAt{103.119}{328.261}{11.955}{\leq}%
\TextAt{9}{115.738}{328.261}{1}%
\TextAt{8}{121.591}{328.261}{,}%
\TextAt{8}{128.554}{328.261}{and}%
\TextAt{8}{151.083}{328.261}{propose}%
\TextAt{8}{194.164}{328.261}{to}%
\TextAt{8}{208.241}{328.261}{illustrate}%
\TextAt{8}{257.503}{328.261}{geometrically,}%
\TextAt{8}{331.462}{328.261}{via}%
\TextAt{8}{350.417}{328.261}{the}%
\TextAt{8}{370.34}{328.261}{scheme-theoretic}%
\TextAt{8}{458.934}{328.261}{interpretation,}%
\TextAt{8}{79.226}{311.215}{some}%
\TextAt{8}{108.552}{311.215}{of}%
\TextAt{8}{121.882}{311.215}{the}%
\TextAt{8}{142.039}{311.215}{facts}%
\TextAt{8}{169.739}{311.215}{treated}%
\TextAt{8}{210.064}{311.215}{in}%
\TextAt{8}{223.716}{311.215}{previous}%
\TextAt{8}{270.282}{311.215}{sections.}%
\TextAt{12}{97.159}{291.18}{1.}%
\TextAt{12}{117.017}{291.18}{Fields.}%
\TextAt{8}{163.767}{291.18}{If}%
\TextAt{13}{176.447}{291.18}{K}%
\TextAt{8}{192.15}{291.18}{is}%
\TextAt{8}{204.901}{291.18}{a}%
\TextAt{8}{215.644}{291.18}{field,}%
\TextAt{8}{245.487}{291.18}{then}%
\TextAt{8}{273.135}{291.18}{the}%
\TextAt{8}{294.279}{291.18}{scheme}%
\TextAt{8}{335.316}{291.18}{Spec}%
\TextAt{9}{358.724}{291.18}{\Character{40}}%
\TextAt{13}{363.276}{291.18}{K}%
\TextAt{9}{374.094}{291.18}{\Character{41}}%
\TextAt{8}{383.531}{291.18}{consists}%
\TextAt{8}{427.629}{291.18}{of}%
\TextAt{8}{441.947}{291.18}{a}%
\TextAt{8}{452.69}{291.18}{single}%
\TextAt{8}{486.255}{291.18}{point}%
\TextAt{9}{517.812}{291.18}{\Character{40}0\Character{41}}%
\TextAt{8}{79.226}{274.134}{on}%
\TextAt{8}{97.102}{274.134}{top}%
\TextAt{8}{119.533}{274.134}{of}%
\TextAt{8}{134.48}{274.134}{which}%
\TextAt{8}{169.583}{274.134}{the}%
\TextAt{8}{191.356}{274.134}{field}%
\TextAt{8}{218.33}{274.134}{itself}%
\TextAt{8}{248.292}{274.134}{sits}%
\TextAt{8}{270.843}{274.134}{as}%
\TextAt{8}{286.83}{274.134}{the}%
\TextAt{8}{308.603}{274.134}{structure}%
\TextAt{8}{360.36}{274.134}{sheaf.}%
\TextAt{8}{399.41}{274.134}{One}%
\TextAt{8}{425.738}{274.134}{must}%
\TextAt{8}{456.358}{274.134}{not}%
\TextAt{8}{478.789}{274.134}{think}%
\TextAt{8}{511.298}{274.134}{that}%
\TextAt{8}{79.226}{257.088}{these}%
\TextAt{8}{109.031}{257.088}{one-point}%
\TextAt{8}{160.894}{257.088}{schemes}%
\TextAt{8}{205.392}{257.088}{are}%
\TextAt{8}{224.736}{257.088}{all}%
\TextAt{8}{240.829}{257.088}{the}%
\TextAt{8}{260.819}{257.088}{same}%
\TextAt{8}{289.978}{257.088}{because}%
\TextAt{8}{333.114}{257.088}{they}%
\TextAt{8}{359.284}{257.088}{differ}%
\TextAt{8}{389.352}{257.088}{essentially}%
\TextAt{8}{445.244}{257.088}{in}%
\TextAt{8}{458.73}{257.088}{their}%
\TextAt{8}{486.527}{257.088}{structure}%
\TextAt{8}{79.226}{240.042}{sheaves.}%
\TextAt{12}{97.159}{220.007}{2.}%
\TextAt{12}{113.769}{220.007}{Valuation}%
\TextAt{12}{174.818}{220.007}{rings.}%
\TextAt{8}{212.648}{220.007}{If}%
\TextAt{14}{224.386}{220.007}{o}%
\TextAt{8}{234.174}{220.007}{is}%
\TextAt{8}{245.985}{220.007}{a}%
\TextAt{8}{255.786}{220.007}{discrete}%
\TextAt{8}{298.812}{220.007}{valuation}%
\TextAt{8}{349.931}{220.007}{ring}%
\TextAt{8}{374.043}{220.007}{with}%
\TextAt{8}{400.749}{220.007}{maximal}%
\TextAt{8}{448.604}{220.007}{ideal}%
\TextAt{14}{476.614}{220.007}{p}%
\TextAt{8}{482.597}{220.007}{,}%
\TextAt{8}{489.804}{220.007}{then}%
\TextAt{8}{516.51}{220.007}{the}%
\TextAt{8}{79.226}{202.961}{scheme}%
\TextAt{13}{118.888}{202.961}{X}%
\TextAt{9}{132.864}{202.961}{\Character{61}}%
\TextAt{8}{145.289}{202.961}{Spec}%
\TextAt{9}{168.697}{202.961}{\Character{40}}%
\TextAt{14}{173.25}{202.961}{o}%
\TextAt{9}{179.094}{202.961}{\Character{41}}%
\TextAt{8}{187.155}{202.961}{consists}%
\TextAt{8}{229.877}{202.961}{of}%
\TextAt{8}{242.819}{202.961}{two}%
\TextAt{8}{264.548}{202.961}{points,}%
\TextAt{8}{302.673}{202.961}{the}%
\TextAt{8}{322.441}{202.961}{closed}%
\TextAt{8}{356.579}{202.961}{point}%
\TextAt{13}{386.76}{202.961}{x}%
\TextAt{9}{396.733}{202.961}{\Character{61}}%
\TextAt{14}{409.159}{202.961}{p}%
\TextAt{8}{418.651}{202.961}{with}%
\TextAt{8}{444.922}{202.961}{residue}%
\TextAt{8}{484.261}{202.961}{class}%
\TextAt{8}{511.31}{202.961}{field}%
\TextAt{13}{79.226}{185.915}{\Character{20}}%
\TextAt{9}{85.965}{185.916}{\Character{40}}%
\TextAt{14}{90.517}{185.915}{p}%
\TextAt{9}{96.5}{185.916}{\Character{41}}%
\TextAt{9}{104.881}{185.916}{\Character{61}}%
\TextAt{14}{117.846}{185.915}{o}%
\TextAt{9}{123.691}{185.916}{\Character{47}}%
\TextAt{14}{129.544}{185.915}{p}%
\TextAt{8}{135.527}{185.915}{,}%
\TextAt{8}{143.06}{185.915}{and}%
\TextAt{8}{166.13}{185.915}{the}%
\TextAt{8}{186.594}{185.915}{generic}%
\TextAt{8}{226.568}{185.915}{point}%
\TextAt{13}{257.445}{185.915}{\Character{17}}%
\TextAt{9}{267.494}{185.916}{\Character{61}}%
\TextAt{9}{280.43}{185.916}{\Character{40}0\Character{41}}%
\TextAt{8}{299.606}{185.915}{with}%
\TextAt{8}{326.573}{185.915}{residue}%
\TextAt{8}{366.608}{185.915}{class}%
\TextAt{8}{394.352}{185.915}{field}%
\TextAt{13}{420.016}{185.915}{\Character{20}}%
\TextAt{9}{426.755}{185.916}{\Character{40}}%
\TextAt{13}{431.307}{185.915}{\Character{17}}%
\TextAt{9}{437.512}{185.916}{\Character{41}}%
\TextAt{9}{445.893}{185.916}{\Character{61}}%
\TextAt{13}{458.858}{185.915}{K}%
\TextAt{8}{469.675}{185.915}{,}%
\TextAt{8}{477.209}{185.915}{the}%
\TextAt{8}{497.673}{185.915}{field}%
\TextAt{8}{523.337}{185.915}{of}%
\TextAt{8}{79.226}{168.87}{fractions}%
\TextAt{8}{127.493}{168.87}{of}%
\TextAt{14}{141.22}{168.87}{o}%
\TextAt{8}{147.065}{168.87}{.}%
\TextAt{8}{156.707}{168.87}{One}%
\TextAt{8}{181.816}{168.87}{should}%
\TextAt{8}{219.347}{168.87}{think}%
\TextAt{8}{250.637}{168.87}{of}%
\TextAt{13}{264.365}{168.87}{X}%
\TextAt{8}{279.315}{168.87}{as}%
\TextAt{8}{294.083}{168.87}{a}%
\TextAt{8}{304.236}{168.87}{point}%
\TextAt{13}{335.204}{168.87}{x}%
\TextAt{8}{346.151}{168.87}{with}%
\TextAt{8}{373.209}{168.87}{an}%
\TextAt{8}{389.866}{168.87}{infinitesimal}%
\TextAt{8}{456.663}{168.87}{neighbourhood}%
\TextAt{8}{79.226}{151.824}{described}%
\TextAt{8}{130.98}{151.824}{by}%
\TextAt{8}{147.239}{151.824}{the}%
\TextAt{8}{167.396}{151.824}{generic}%
\TextAt{8}{207.063}{151.824}{point}%
\TextAt{13}{237.632}{151.824}{\Character{17}}%
\TextAt{8}{243.837}{151.824}{:}%
\TextAt{13}{253.977}{123.271}{X}%
\TextAt{9}{267.953}{123.271}{\Character{58}}%
\NativeDiagram{%
\path[draw=black,line width=0.399bp] (204.764,-646.644) -- (289.804,-646.644);%
\path[fill=black] (238.78,-646.644) (240.773,-646.644) .. controls (240.773,-645.544) and (239.881,-644.651) .. (238.78,-644.651) .. controls (237.68,-644.651) and (236.788,-645.544) .. (236.788,-646.644) .. controls (236.788,-647.744) and (237.68,-648.637) .. (238.78,-648.637) .. controls (239.881,-648.637) and (240.773,-647.744) .. (240.773,-646.644) -- cycle (238.78,-646.644);%
\path[draw=black,line width=0.399bp,pattern=dots] (267.127,-646.644) (272.229,-646.644) .. controls (272.229,-645.078) and (269.945,-643.809) .. (267.127,-643.809) .. controls (264.309,-643.809) and (262.025,-645.078) .. (262.025,-646.644) .. controls (262.025,-648.21) and (264.309,-649.479) .. (267.127,-649.479) .. controls (269.945,-649.479) and (272.229,-648.21) .. (272.229,-646.644) -- cycle (267.127,-646.644);%
}%
\TextAt{13}{307.454}{114.623}{x}%
\TextAt{13}{336.025}{114.623}{\Character{17}}%
\TextAt{8}{79.226}{78.316}{This}%
\TextAt{8}{105.946}{78.316}{intuition}%
\TextAt{8}{153.755}{78.316}{is}%
\TextAt{8}{165.519}{78.316}{justified}%
\TextAt{8}{210.123}{78.316}{by}%
\TextAt{8}{226.382}{78.316}{the}%
\TextAt{8}{246.539}{78.316}{following}%
\TextAt{8}{295.973}{78.316}{observation.}%
\TextAt{8}{300.14}{43.033}{83}%
\Destination{79.23}{697.115}{nav.Item.22}%
\Destination{79.23}{576.635}{nav.Item.23}%
\Destination{78.23}{809.815}{nav.page.83}%
\Destination{79.36}{698.055}{nav.print.89}%
\Hotlink{78.726}{363.155}{19.22}{11.047}{nav.cite.N051}%
\Hotlink{123.606}{363.155}{25.078}{11.047}{nav.cite.N104}%
}

% Page 96 (body page 84)
\Page{%
\TextAt{2}{79.226}{796.943}{I.}%
\TextAt{2}{89.83}{796.943}{Algebraic}%
\TextAt{2}{138.648}{796.943}{Integers}%
\Vector{q 0 G 0 g 0.299 w 7.226 18.729 m 460.769 18.729 l S Q}%
\TextAt{8}{97.159}{761.124}{The}%
\TextAt{8}{121.213}{761.124}{discrete}%
\TextAt{8}{164.192}{761.124}{valuation}%
\TextAt{8}{215.264}{761.124}{rings}%
\TextAt{8}{243.944}{761.124}{arise}%
\TextAt{8}{271.322}{761.124}{as}%
\TextAt{8}{285.692}{761.124}{localizations}%
\TextAt{14}{228.387}{729.263}{o}%
\TextAt{27}{234.232}{727.47}{p}%
\TextAt{9}{242.039}{729.263}{\Character{61}}%
\TextAt{20}{254.465}{746.12}{\Character{26}}%
\TextAt{13}{264.627}{737.351}{f}%
\Vector{q 0 G 0 g 0.478 w 192.627 -41.748 m 199.673 -41.748 l S Q}%
\TextAt{13}{265.133}{721.063}{g}%
\TextAt{20}{276.766}{746.598}{\Character{12}}%
\TextAt{20}{276.766}{739.425}{\Character{12}}%
\TextAt{20}{276.766}{732.252}{\Character{12}}%
\TextAt{20}{276.766}{725.079}{\Character{12}}%
\TextAt{13}{284.648}{729.263}{f}%
\TextAt{13}{291.032}{729.263}{\Character{59}}%
\TextAt{13}{296.269}{729.263}{g}%
\NativeMathAt{305.644}{729.263}{11.955}{\in}%
\TextAt{14}{316.935}{729.263}{o}%
\TextAt{13}{322.779}{729.263}{\Character{59}}%
\TextAt{13}{331.901}{729.263}{g}%
\TextAt{9}{337.955}{729.263}{\Character{40}}%
\TextAt{14}{342.507}{729.263}{p}%
\TextAt{9}{348.49}{729.263}{\Character{41}}%
\NativeMathAt{356.363}{729.263}{11.955}{\neq}%
\TextAt{9}{368.773}{729.263}{0}%
\TextAt{20}{374.642}{746.12}{\Character{27}}%
\ErrataAt{79.19}{696.882}{455.81}{of Dedekind domains $\mathfrak o$. If $\mathfrak o$ has infinitely many nonzero prime ideals, every neighbourhood}%
\ErrataAt{79.19}{679.836}{455.81}{of $\mathfrak p$ in $X=\operatorname{Spec}(\mathfrak o)$ contains a nonzero prime $\mathfrak q\ne\mathfrak p$. By the Chinese remainder}%
\ErrataAt{79.19}{662.79}{455.81}{theorem there is $g\in\mathfrak o$ with $g\equiv0\pmod{\mathfrak q}$ and $g\equiv1\pmod{\mathfrak p}$. Thus $1/g$ is not regular at $\mathfrak q$.}%
\ErrataAt{79.19}{645.744}{455.81}{In the semilocal case all of $\mathfrak o_{\mathfrak p}$ is regular on the open $\{(0),\mathfrak p\}$. In either case,}%
\TextAt{8}{79.226}{628.681}{every}%
\TextAt{8}{110.472}{628.681}{element}%
\TextAt{13}{154.068}{628.681}{f}%
\TextAt{9}{161.114}{628.681}{\Character{47}}%
\TextAt{13}{166.967}{628.681}{g}%
\NativeMathAt{176.926}{628.681}{11.955}{\in}%
\TextAt{14}{188.82}{628.681}{o}%
\TextAt{27}{194.664}{626.888}{p}%
\TextAt{8}{203.402}{628.681}{is}%
\TextAt{8}{215.52}{628.681}{defined}%
\TextAt{8}{256.186}{628.681}{on}%
\TextAt{8}{272.799}{628.681}{a}%
\TextAt{8}{282.909}{628.681}{sufficiently}%
\TextAt{8}{341.855}{628.681}{small}%
\TextAt{8}{372.838}{628.681}{neighbourhood;}%
\TextAt{8}{456.625}{628.681}{hence}%
\TextAt{8}{489.485}{628.681}{one}%
\TextAt{8}{511.298}{628.681}{may}%
\TextAt{8}{79.226}{611.635}{say}%
\TextAt{8}{100.192}{611.635}{that}%
\TextAt{8}{126.298}{611.635}{all}%
\TextAt{8}{143.294}{611.635}{elements}%
\TextAt{13}{191.888}{611.635}{f}%
\TextAt{9}{198.935}{611.635}{\Character{47}}%
\TextAt{13}{204.788}{611.635}{g}%
\TextAt{8}{215.456}{611.635}{of}%
\TextAt{8}{229.524}{611.635}{the}%
\TextAt{8}{250.418}{611.635}{discrete}%
\TextAt{8}{294.134}{611.635}{valuation}%
\TextAt{8}{345.944}{611.635}{ring}%
\TextAt{14}{370.747}{611.635}{o}%
\TextAt{27}{376.591}{609.842}{p}%
\TextAt{8}{385.713}{611.635}{are}%
\TextAt{8}{405.961}{611.635}{like}%
\TextAt{8}{428.158}{611.635}{functions}%
\TextAt{8}{479.358}{611.635}{defined}%
\TextAt{8}{520.408}{611.635}{on}%
\TextAt{8}{79.226}{594.589}{a}%
\TextAt{8}{89.431}{594.589}{“germ”}%
\TextAt{8}{130.194}{594.589}{of}%
\TextAt{8}{143.974}{594.589}{neighbourhoods}%
\TextAt{8}{229.042}{594.589}{of}%
\TextAt{14}{242.822}{594.589}{p}%
\TextAt{8}{248.805}{594.589}{.}%
\TextAt{8}{258.603}{594.589}{Thus}%
\TextAt{8}{288.701}{594.589}{Spec}%
\TextAt{9}{312.109}{594.589}{\Character{40}}%
\TextAt{14}{316.662}{594.589}{o}%
\TextAt{27}{322.506}{592.796}{p}%
\TextAt{9}{326.993}{594.589}{\Character{41}}%
\TextAt{8}{335.892}{594.589}{may}%
\TextAt{8}{361.711}{594.589}{be}%
\TextAt{8}{378.085}{594.589}{thought}%
\TextAt{8}{422.446}{594.589}{of}%
\TextAt{8}{436.226}{594.589}{as}%
\TextAt{8}{451.046}{594.589}{such}%
\TextAt{8}{477.892}{594.589}{a}%
\TextAt{8}{488.098}{594.589}{“germ}%
\TextAt{8}{523.337}{594.589}{of}%
\TextAt{8}{79.226}{577.543}{neighbourhoods”}%
\TextAt{8}{169.368}{577.543}{of}%
\TextAt{14}{182.698}{577.543}{p}%
\TextAt{8}{188.681}{577.543}{.}%
\TextAt{8}{97.159}{557.508}{We}%
\TextAt{8}{118.387}{557.508}{want}%
\TextAt{8}{148.092}{557.508}{to}%
\TextAt{8}{163.486}{557.508}{point}%
\TextAt{8}{195.139}{557.508}{out}%
\TextAt{8}{217.037}{557.508}{a}%
\TextAt{8}{227.877}{557.508}{small}%
\TextAt{8}{259.59}{557.508}{discrepancy}%
\TextAt{8}{324.144}{557.508}{of}%
\TextAt{8}{338.558}{557.508}{intuitions.}%
\TextAt{8}{398.784}{557.508}{Considering}%
\TextAt{8}{464.318}{557.508}{the}%
\TextAt{8}{485.559}{557.508}{spectrum}%
\ErrataAt{79.19}{540.462}{455.81}{of the one-dimensional ring $\mathbb C[x]$, whose closed points constitute the complex plane, we}%
\TextAt{8}{79.226}{523.416}{will}%
\TextAt{8}{102.132}{523.416}{not}%
\TextAt{8}{123.746}{523.416}{want}%
\TextAt{8}{153.167}{523.416}{to}%
\TextAt{8}{168.277}{523.416}{visualize}%
\TextAt{8}{216.288}{523.416}{the}%
\TextAt{8}{237.245}{523.416}{infinitesimal}%
\TextAt{8}{304.444}{523.416}{neighbourhood}%
\TextAt{13}{385.249}{523.416}{X}%
\TextAt{27}{394.964}{521.623}{p}%
\TextAt{9}{404.135}{523.416}{\Character{61}}%
\TextAt{8}{417.924}{523.416}{Spec}%
\TextAt{9}{441.333}{523.416}{\Character{40}}%
\TextAt{18}{445.885}{523.416}{C}%
\TextAt{9}{454.519}{523.416}{\Character{91}}%
\TextAt{13}{457.771}{523.416}{x}%
\TextAt{9}{464.423}{523.416}{\Character{93}}%
\TextAt{27}{467.675}{521.623}{p}%
\TextAt{9}{472.162}{523.416}{\Character{41}}%
\TextAt{8}{481.412}{523.416}{of}%
\TextAt{8}{495.542}{523.416}{a}%
\TextAt{8}{506.097}{523.416}{point}%
\TextAt{14}{79.226}{506.37}{p}%
\TextAt{9}{88.53}{506.37}{\Character{61}}%
\TextAt{9}{100.94}{506.37}{\Character{40}}%
\TextAt{13}{105.508}{506.37}{x}%
\TextAt{16}{114.817}{506.37}{\Character{0}}%
\TextAt{13}{126.772}{506.37}{a}%
\TextAt{9}{132.917}{506.37}{\Character{41}}%
\TextAt{8}{141.366}{506.37}{as}%
\TextAt{8}{155.737}{506.37}{a}%
\TextAt{8}{165.492}{506.37}{small}%
\TextAt{8}{196.121}{506.37}{line}%
\TextAt{8}{218.226}{506.37}{segment,}%
\TextAt{8}{266.74}{506.37}{but}%
\TextAt{8}{288.2}{506.37}{rather}%
\TextAt{8}{323.324}{506.37}{as}%
\TextAt{8}{337.694}{506.37}{a}%
\TextAt{8}{347.45}{506.37}{little}%
\TextAt{8}{375.413}{506.37}{disc:}%
\NativeDiagram{%
\path[draw=black,line width=0.399bp] (233.998,-312.211) (276.518,-312.211) .. controls (276.518,-305.166) and (257.482,-299.455) .. (233.998,-299.455) .. controls (210.514,-299.455) and (191.478,-305.166) .. (191.478,-312.211) .. controls (191.478,-319.256) and (210.514,-324.967) .. (233.998,-324.967) .. controls (257.482,-324.967) and (276.518,-319.256) .. (276.518,-312.211) -- cycle (233.998,-312.211);%
\path[fill=black] (233.998,-312.211) (235.592,-312.211) .. controls (235.592,-311.331) and (234.878,-310.617) .. (233.998,-310.617) .. controls (233.118,-310.617) and (232.404,-311.331) .. (232.404,-312.211) .. controls (232.404,-313.091) and (233.118,-313.805) .. (233.998,-313.805) .. controls (234.878,-313.805) and (235.592,-313.091) .. (235.592,-312.211) -- cycle (233.998,-312.211);%
\path[draw=black,line width=0.399bp,pattern=dots] (256.676,-312.211) (263.195,-312.211) .. controls (263.195,-310.332) and (260.276,-308.81) .. (256.676,-308.81) .. controls (253.075,-308.81) and (250.156,-310.332) .. (250.156,-312.211) .. controls (250.156,-314.09) and (253.075,-315.612) .. (256.676,-315.612) .. controls (260.276,-315.612) and (263.195,-314.09) .. (263.195,-312.211) -- cycle (256.676,-312.211);%
}%
\TextAt{13}{302.672}{478.729}{x}%
\TextAt{13}{325.573}{481.054}{\Character{17}}%
\TextAt{8}{79.226}{419.588}{This}%
\TextAt{8}{105.75}{419.588}{two-dimensional}%
\TextAt{8}{192.121}{419.588}{nature}%
\TextAt{8}{228.998}{419.588}{is}%
\TextAt{8}{240.565}{419.588}{actually}%
\TextAt{8}{284.926}{419.588}{inherent}%
\TextAt{8}{330.578}{419.588}{in}%
\TextAt{8}{344.034}{419.588}{all}%
\TextAt{8}{360.097}{419.588}{discrete}%
\TextAt{8}{402.88}{419.588}{valuation}%
\TextAt{8}{453.756}{419.588}{rings}%
\TextAt{8}{482.24}{419.588}{with}%
\TextAt{8}{508.704}{419.588}{alge-}%
\TextAt{8}{79.226}{402.542}{braically}%
\TextAt{8}{127.367}{402.542}{closed}%
\TextAt{8}{162.226}{402.542}{residue}%
\TextAt{8}{202.286}{402.542}{field.}%
\TextAt{8}{233.191}{402.542}{But}%
\TextAt{8}{256.764}{402.542}{the}%
\TextAt{8}{277.253}{402.542}{algebraic}%
\TextAt{8}{327.02}{402.542}{justification}%
\TextAt{8}{391.48}{402.542}{of}%
\TextAt{8}{405.143}{402.542}{this}%
\TextAt{8}{428.298}{402.542}{intuition}%
\TextAt{8}{476.439}{402.542}{is}%
\TextAt{8}{488.535}{402.542}{provided}%
\TextAt{8}{79.226}{385.496}{only}%
\TextAt{8}{106.095}{385.496}{by}%
\TextAt{8}{123.53}{385.496}{the}%
\TextAt{8}{144.863}{385.496}{introduction}%
\TextAt{8}{213.037}{385.496}{of}%
\TextAt{8}{227.544}{385.496}{a}%
\TextAt{8}{238.476}{385.496}{new}%
\TextAt{8}{263.706}{385.496}{topology,}%
\TextAt{8}{315.592}{385.496}{the}%
\TextAt{12}{336.925}{385.496}{étale}%
\TextAt{12}{370.546}{385.496}{topology}%
\TextAt{8}{421.391}{385.496}{,}%
\TextAt{8}{430.011}{385.496}{which}%
\TextAt{8}{464.675}{385.496}{is}%
\TextAt{8}{477.615}{385.496}{much}%
\TextAt{8}{510.007}{385.496}{finer}%
\TextAt{8}{79.226}{368.451}{than}%
\TextAt{8}{106.544}{368.451}{the}%
\TextAt{8}{126.7}{368.451}{Zariski}%
\TextAt{8}{165.459}{368.451}{topology}%
\TextAt{8}{213.602}{368.451}{(see}%
\TextAt{8}{237.07}{368.451}{[103],}%
\TextAt{8}{268.297}{368.451}{[132]).}%
\TextAt{12}{97.159}{348.416}{3.}%
\TextAt{12}{113.225}{348.416}{Dedekind}%
\TextAt{12}{172.387}{348.416}{rings.}%
\TextAt{8}{209.744}{348.416}{The}%
\TextAt{8}{232.796}{348.416}{spectrum}%
\TextAt{13}{282.903}{348.416}{X}%
\TextAt{9}{296.879}{348.416}{\Character{61}}%
\TextAt{8}{309.305}{348.416}{Spec}%
\TextAt{9}{332.713}{348.416}{\Character{40}}%
\TextAt{14}{337.265}{348.416}{o}%
\TextAt{9}{343.109}{348.416}{\Character{41}}%
\TextAt{8}{350.558}{348.416}{of}%
\TextAt{8}{362.886}{348.416}{a}%
\TextAt{8}{371.64}{348.416}{Dedekind}%
\TextAt{8}{422.823}{348.416}{domain}%
\TextAt{14}{463.45}{348.416}{o}%
\TextAt{8}{472.19}{348.416}{is}%
\TextAt{8}{482.952}{348.416}{visualized}%
\TextAt{8}{79.226}{331.37}{as}%
\TextAt{8}{93.919}{331.37}{a}%
\TextAt{8}{103.997}{331.37}{smooth}%
\TextAt{8}{145.685}{331.37}{curve.}%
\TextAt{8}{182.419}{331.37}{At}%
\TextAt{8}{199.647}{331.37}{each}%
\TextAt{8}{226.307}{331.37}{point}%
\TextAt{14}{257.199}{331.37}{p}%
\TextAt{8}{267.402}{331.37}{one}%
\TextAt{8}{289.185}{331.37}{may}%
\TextAt{8}{314.876}{331.37}{consider}%
\TextAt{8}{360.784}{331.37}{the}%
\TextAt{8}{381.264}{331.37}{localization}%
\TextAt{14}{443.705}{331.37}{o}%
\TextAt{27}{449.55}{329.577}{p}%
\TextAt{8}{454.037}{331.37}{.}%
\TextAt{8}{463.453}{331.37}{The}%
\TextAt{8}{487.83}{331.37}{inclusion}%
\TextAt{14}{79.226}{314.324}{o}%
\NativeArrowAt{88.391}{314.324}{11.955}{13.214}{\hookrightarrow}%
\TextAt{14}{104.926}{314.324}{o}%
\TextAt{27}{110.771}{312.531}{p}%
\TextAt{8}{119.155}{314.324}{induces}%
\TextAt{8}{160.831}{314.324}{a}%
\TextAt{8}{170.586}{314.324}{morphism}%
\TextAt{13}{241.552}{286.2}{f}%
\TextAt{9}{251.919}{286.2}{\Character{58}}%
\TextAt{13}{258.492}{286.2}{X}%
\TextAt{27}{268.207}{284.406}{p}%
\TextAt{9}{276.015}{286.2}{\Character{61}}%
\TextAt{8}{288.44}{286.2}{Spec}%
\TextAt{9}{311.848}{286.2}{\Character{40}}%
\TextAt{14}{316.401}{286.2}{o}%
\TextAt{27}{322.245}{284.406}{p}%
\TextAt{9}{326.732}{286.2}{\Character{41}}%
\NativeArrowAt{334.605}{286.2}{11.955}{19.26}{\longrightarrow}%
\TextAt{13}{357.187}{286.2}{X}%
\TextAt{13}{367.182}{286.2}{\Character{59}}%
\TextAt{8}{79.226}{258.075}{which}%
\TextAt{8}{112.713}{258.075}{extracts}%
\TextAt{8}{157.329}{258.075}{the}%
\TextAt{8}{177.486}{258.075}{scheme}%
\TextAt{13}{217.536}{258.075}{X}%
\TextAt{27}{227.251}{256.282}{p}%
\TextAt{8}{235.635}{258.075}{from}%
\TextAt{13}{263.275}{258.075}{X}%
\TextAt{8}{277.828}{258.075}{as}%
\TextAt{8}{292.198}{258.075}{an}%
\TextAt{8}{308.457}{258.075}{“infinitesimal}%
\TextAt{8}{380.379}{258.075}{neighbourhood”}%
\TextAt{8}{465.907}{258.075}{of}%
\TextAt{14}{479.237}{258.075}{p}%
\TextAt{8}{485.22}{258.075}{:}%
\NativeDiagram{%
\path[draw=black,line width=0.399bp] (92.752,-587.495) .. controls (149.446,-553.479) and (189.131,-584.66) .. (223.147,-593.164) .. controls (254.329,-601.668) and (282.676,-604.503) .. (302.518,-587.495);%
\path[fill=black] (143.776,-573.038) (145.271,-573.038) .. controls (145.271,-572.213) and (144.602,-571.544) .. (143.776,-571.544) .. controls (142.951,-571.544) and (142.282,-572.213) .. (142.282,-573.038) .. controls (142.282,-573.864) and (142.951,-574.533) .. (143.776,-574.533) .. controls (144.602,-574.533) and (145.271,-573.864) .. (145.271,-573.038) -- cycle (143.776,-573.038);%
}%
\TextAt{14}{212.785}{187.963}{p}%
\NativeDiagram{%
\path[draw=black,line width=0.399bp] (126.768,-544.975) .. controls (138.107,-539.873) and (152.28,-539.873) .. (162.202,-544.975);%
\path[fill=black] (143.776,-541.29) (145.271,-541.29) .. controls (145.271,-540.464) and (144.602,-539.795) .. (143.776,-539.795) .. controls (142.951,-539.795) and (142.282,-540.464) .. (142.282,-541.29) .. controls (142.282,-542.115) and (142.951,-542.784) .. (143.776,-542.784) .. controls (144.602,-542.784) and (145.271,-542.115) .. (145.271,-541.29) -- cycle (143.776,-541.29);%
}%
\TextAt{13}{238.386}{229.425}{X}%
\TextAt{27}{248.101}{227.632}{p}%
\NativeDiagram{%
\path[draw=black,line width=0.399bp] (143.776,-549.227) -- (143.776,-564.986);%
\path[draw=black,line width=0.399bp,line cap=round,line join=round] (146.167,-563.113) .. controls (144.733,-563.492) and (144.055,-564.335) .. (143.776,-565.185) .. controls (143.497,-564.335) and (142.82,-563.492) .. (141.385,-563.113);%
}%
\TextAt{13}{219.96}{213.706}{f}%
\NativeDiagram{%
\path[draw=black,line width=0.399bp,pattern=dots] (257.164,-599.684) (263.683,-599.684) .. controls (263.683,-597.805) and (260.764,-596.282) .. (257.164,-596.282) .. controls (253.563,-596.282) and (250.644,-597.805) .. (250.644,-599.684) .. controls (250.644,-601.562) and (253.563,-603.085) .. (257.164,-603.085) .. controls (260.764,-603.085) and (263.683,-601.562) .. (263.683,-599.684) -- cycle (257.164,-599.684);%
}%
\TextAt{8}{295.994}{154.672}{generic}%
\TextAt{8}{335.661}{154.672}{point}%
\TextAt{13}{378.702}{183.516}{X}%
\TextAt{9}{392.678}{183.516}{\Character{61}}%
\TextAt{8}{405.104}{183.516}{Spec}%
\TextAt{9}{428.512}{183.516}{\Character{40}}%
\TextAt{14}{433.064}{183.516}{o}%
\TextAt{9}{438.908}{183.516}{\Character{41}}%
\TextAt{8}{300.14}{43.033}{84}%
\Destination{78.23}{809.815}{nav.page.84}%
\Destination{79.36}{776.155}{nav.print.90}%
\Hotlink{236.57}{366.255}{25.078}{11.047}{nav.cite.N103}%
\Hotlink{267.797}{366.255}{25.078}{11.047}{nav.cite.N132}%
\Destination{79.226}{713.882}{errata.54}%
\Destination{79.226}{557.462}{errata.64}%
}

% Page 97 (body page 85)
\Page{%
\TextAt{2}{79.226}{796.943}{I.}%
\TextAt{2}{89.83}{796.943}{Algebraic}%
\TextAt{2}{138.648}{796.943}{Integers}%
\Vector{q 0 G 0 g 0.299 w 7.226 18.729 m 460.769 18.729 l S Q}%
\TextAt{12}{79.226}{761.124}{4.}%
\TextAt{12}{100.751}{761.124}{Singularities.}%
\TextAt{8}{186.932}{761.124}{We}%
\TextAt{8}{208.547}{761.124}{now}%
\TextAt{8}{234.406}{761.124}{consider}%
\TextAt{8}{281.462}{761.124}{a}%
\TextAt{8}{292.688}{761.124}{one-dimensional}%
\TextAt{8}{380.068}{761.124}{noetherian}%
\TextAt{8}{439.748}{761.124}{integral}%
\TextAt{8}{483.827}{761.124}{domain}%
\TextAt{14}{526.925}{761.124}{o}%
\TextAt{8}{79.226}{744.078}{which}%
\TextAt{8}{114.033}{744.078}{is}%
\TextAt{8}{127.117}{744.078}{not}%
\TextAt{8}{149.252}{744.078}{a}%
\TextAt{8}{160.328}{744.078}{Dedekind}%
\TextAt{8}{213.832}{744.078}{domain,}%
\TextAt{8}{260.363}{744.078}{e.g.,}%
\TextAt{8}{286.725}{744.078}{an}%
\TextAt{8}{304.304}{744.078}{order}%
\TextAt{8}{336.194}{744.078}{in}%
\TextAt{8}{351.167}{744.078}{an}%
\TextAt{8}{368.747}{744.078}{algebraic}%
\TextAt{8}{419.502}{744.078}{number}%
\TextAt{8}{463.419}{744.078}{field}%
\TextAt{8}{490.096}{744.078}{which}%
\TextAt{8}{524.903}{744.078}{is}%
\TextAt{8}{79.226}{727.032}{different}%
\TextAt{8}{126.197}{727.032}{from}%
\TextAt{8}{154.638}{727.032}{the}%
\TextAt{8}{175.595}{727.032}{maximal}%
\TextAt{8}{224.204}{727.032}{order.}%
\TextAt{8}{261.726}{727.032}{Again}%
\TextAt{8}{296.671}{727.032}{we}%
\TextAt{8}{314.698}{727.032}{view}%
\TextAt{8}{342.482}{727.032}{the}%
\TextAt{8}{363.439}{727.032}{scheme}%
\TextAt{13}{404.289}{727.032}{X}%
\TextAt{9}{419.629}{727.032}{\Character{61}}%
\TextAt{8}{433.419}{727.032}{Spec}%
\TextAt{9}{456.827}{727.032}{\Character{40}}%
\TextAt{14}{461.379}{727.032}{o}%
\TextAt{9}{467.223}{727.032}{\Character{41}}%
\TextAt{8}{476.474}{727.032}{as}%
\TextAt{8}{491.644}{727.032}{a}%
\TextAt{8}{502.2}{727.032}{curve.}%
\TextAt{8}{79.226}{709.986}{But}%
\TextAt{8}{102.157}{709.986}{now}%
\TextAt{8}{126.235}{709.986}{the}%
\TextAt{8}{146.082}{709.986}{curve}%
\TextAt{8}{176.986}{709.986}{will}%
\TextAt{8}{198.781}{709.986}{not}%
\TextAt{8}{219.285}{709.986}{be}%
\TextAt{8}{234.899}{709.986}{everywhere}%
\TextAt{8}{295.393}{709.986}{smooth,}%
\TextAt{8}{339.762}{709.986}{but}%
\TextAt{8}{360.911}{709.986}{will}%
\TextAt{8}{382.706}{709.986}{have}%
\TextAt{12}{409.391}{709.986}{singularities}%
\TextAt{8}{483.645}{709.986}{at}%
\TextAt{8}{497.645}{709.986}{certain}%
\TextAt{8}{79.226}{692.94}{points.}%
\NativeDiagram{%
\path[draw=black,line width=0.399bp] (174.958,-152.799) .. controls (115.429,-148.547) and (84.248,-165.555) .. (109.76,-174.06) .. controls (126.768,-178.311) and (166.454,-158.469) .. (194.8,-138.626);%
\path[draw=black,line width=0.399bp] (174.958,-152.799) .. controls (211.809,-152.799) and (227.399,-161.303) .. (225.982,-186.816) .. controls (245.825,-168.39) and (291.179,-152.799) .. (311.023,-166.973);%
\path[fill=black] (174.958,-152.799) (176.452,-152.799) .. controls (176.452,-151.974) and (175.783,-151.305) .. (174.958,-151.305) .. controls (174.132,-151.305) and (173.463,-151.974) .. (173.463,-152.799) .. controls (173.463,-153.625) and (174.132,-154.294) .. (174.958,-154.294) .. controls (175.783,-154.294) and (176.452,-153.625) .. (176.452,-152.799) -- cycle (174.958,-152.799);%
\path[draw=black,line width=0.399bp,pattern=dots] (282.676,-161.303) (288.912,-161.303) .. controls (288.912,-159.425) and (286.12,-157.902) .. (282.676,-157.902) .. controls (279.231,-157.902) and (276.439,-159.425) .. (276.439,-161.303) .. controls (276.439,-163.182) and (279.231,-164.705) .. (282.676,-164.705) .. controls (286.12,-164.705) and (288.912,-163.182) .. (288.912,-161.303) -- cycle (282.676,-161.303);%
}%
\TextAt{8}{321.506}{592.202}{generic}%
\TextAt{8}{361.173}{592.202}{point}%
\TextAt{13}{387.206}{604.038}{X}%
\TextAt{9}{401.182}{604.038}{\Character{61}}%
\TextAt{8}{413.607}{604.038}{Spec}%
\TextAt{9}{437.015}{604.038}{\Character{40}}%
\TextAt{14}{441.568}{604.038}{o}%
\TextAt{9}{447.412}{604.038}{\Character{41}}%
\NativeDiagram{%
\path[draw=black,line width=0.399bp] (159.367,-107.444) -- (191.966,-108.862);%
\path[draw=black,line width=0.399bp] (160.784,-114.531) -- (189.131,-100.358);%
\path[fill=black] (174.958,-108.011) (176.452,-108.011) .. controls (176.452,-107.186) and (175.783,-106.517) .. (174.958,-106.517) .. controls (174.132,-106.517) and (173.463,-107.186) .. (173.463,-108.011) .. controls (173.463,-108.837) and (174.132,-109.506) .. (174.958,-109.506) .. controls (175.783,-109.506) and (176.452,-108.837) .. (176.452,-108.011) -- cycle (174.958,-108.011);%
}%
\TextAt{13}{270.984}{663.722}{X}%
\TextAt{27}{280.699}{661.929}{p}%
\TextAt{9}{288.507}{663.722}{\Character{61}}%
\TextAt{8}{300.932}{663.722}{Spec}%
\TextAt{9}{324.34}{663.722}{\Character{40}}%
\TextAt{14}{328.893}{663.722}{o}%
\TextAt{27}{334.737}{661.929}{p}%
\TextAt{9}{339.224}{663.722}{\Character{41}}%
\NativeDiagram{%
\path[draw=black,line width=0.399bp] (174.958,-118.783) -- (174.958,-141.062);%
\path[draw=black,line width=0.399bp,line cap=round,line join=round] (177.349,-139.189) .. controls (175.914,-139.568) and (175.237,-140.411) .. (174.958,-141.261) .. controls (174.679,-140.411) and (174.001,-139.568) .. (172.567,-139.189);%
}%
\TextAt{8}{79.226}{555.769}{These}%
\TextAt{8}{113.873}{555.769}{are}%
\TextAt{8}{134.162}{555.769}{precisely}%
\TextAt{8}{182.797}{555.769}{the}%
\TextAt{8}{203.732}{555.769}{nongeneric}%
\TextAt{8}{263.042}{555.769}{points}%
\TextAt{14}{299.005}{555.769}{p}%
\TextAt{8}{309.663}{555.769}{for}%
\TextAt{8}{328.326}{555.769}{which}%
\TextAt{8}{362.591}{555.769}{the}%
\TextAt{8}{383.526}{555.769}{localization}%
\TextAt{14}{446.423}{555.769}{o}%
\TextAt{27}{452.267}{553.976}{p}%
\TextAt{8}{461.43}{555.769}{is}%
\TextAt{8}{473.972}{555.769}{no}%
\TextAt{8}{491.009}{555.769}{longer}%
\TextAt{8}{526.911}{555.769}{a}%
\TextAt{8}{79.226}{538.723}{discrete}%
\TextAt{8}{122.295}{538.723}{valuation}%
\TextAt{8}{173.458}{538.723}{ring,}%
\TextAt{8}{200.888}{538.723}{that}%
\TextAt{8}{226.347}{538.723}{is}%
\TextAt{8}{238.201}{538.723}{to}%
\TextAt{8}{252.601}{538.723}{say,}%
\TextAt{8}{275.214}{538.723}{the}%
\TextAt{8}{295.46}{538.723}{maximal}%
\TextAt{8}{343.359}{538.723}{ideal}%
\TextAt{14}{371.412}{538.723}{po}%
\TextAt{27}{383.24}{536.93}{p}%
\TextAt{8}{391.714}{538.723}{is}%
\TextAt{8}{403.568}{538.723}{not}%
\TextAt{8}{424.472}{538.723}{generated}%
\TextAt{8}{477.894}{538.723}{by}%
\TextAt{8}{494.243}{538.723}{a}%
\TextAt{8}{504.089}{538.723}{single}%
\TextAt{8}{79.226}{521.677}{element.}%
\TextAt{8}{126.819}{521.677}{For}%
\TextAt{8}{147.189}{521.677}{example,}%
\TextAt{8}{195.809}{521.677}{in}%
\TextAt{8}{208.863}{521.677}{the}%
\TextAt{8}{228.42}{521.677}{one-dimensional}%
\TextAt{8}{313.73}{521.677}{ring}%
\TextAt{14}{337.197}{521.677}{o}%
\TextAt{9}{346.362}{521.678}{\Character{61}}%
\TextAt{18}{358.787}{521.677}{C}%
\TextAt{9}{367.421}{521.678}{\Character{91}}%
\TextAt{13}{370.673}{521.677}{x\Character{59}}%
\TextAt{13}{382.545}{521.677}{y}%
\TextAt{9}{388.706}{521.678}{\Character{93}\Character{47}\Character{40}}%
\TextAt{13}{402.363}{521.677}{y}%
\TextAt{17}{408.5}{526.016}{2}%
\TextAt{16}{414.663}{521.678}{\Character{0}}%
\TextAt{13}{425.392}{521.677}{x}%
\TextAt{17}{432.044}{526.016}{3}%
\TextAt{9}{436.777}{521.678}{\Character{41}}%
\TextAt{8}{441.329}{521.677}{,}%
\TextAt{8}{447.999}{521.677}{the}%
\TextAt{8}{467.556}{521.677}{closed}%
\TextAt{8}{501.483}{521.677}{points}%
\TextAt{8}{79.226}{504.632}{of}%
\TextAt{8}{92.556}{504.632}{the}%
\TextAt{8}{112.713}{504.632}{scheme}%
\TextAt{13}{152.762}{504.632}{X}%
\TextAt{8}{167.315}{504.632}{are}%
\TextAt{8}{186.826}{504.632}{given}%
\TextAt{8}{217.395}{504.632}{by}%
\TextAt{8}{233.654}{504.632}{the}%
\TextAt{8}{253.811}{504.632}{prime}%
\TextAt{8}{286.974}{504.632}{ideals}%
\TextAt{14}{224.027}{475.63}{p}%
\TextAt{9}{233.331}{475.63}{\Character{61}}%
\TextAt{9}{245.741}{475.63}{\Character{40}}%
\TextAt{13}{250.309}{475.63}{x}%
\TextAt{16}{259.618}{475.63}{\Character{0}}%
\TextAt{13}{271.573}{475.63}{a\Character{59}}%
\TextAt{13}{282.942}{475.63}{y}%
\TextAt{16}{291.755}{475.63}{\Character{0}}%
\TextAt{13}{303.71}{475.63}{b}%
\TextAt{9}{308.688}{475.63}{\Character{41}}%
\TextAt{8}{316.561}{475.63}{mod}%
\TextAt{9}{341.999}{475.63}{\Character{40}}%
\TextAt{13}{346.551}{475.63}{y}%
\TextAt{17}{352.688}{480.567}{2}%
\TextAt{16}{360.077}{475.63}{\Character{0}}%
\TextAt{13}{372.032}{475.63}{x}%
\TextAt{17}{378.684}{480.567}{3}%
\TextAt{9}{383.416}{475.63}{\Character{41}}%
\TextAt{8}{79.226}{446.629}{where}%
\TextAt{9}{113.035}{446.629}{\Character{40}}%
\TextAt{13}{117.588}{446.629}{a\Character{59}}%
\TextAt{13}{128.957}{446.629}{b}%
\TextAt{9}{133.954}{446.629}{\Character{41}}%
\TextAt{8}{142.404}{446.629}{varies}%
\TextAt{8}{175.316}{446.629}{over}%
\TextAt{8}{200.362}{446.629}{the}%
\TextAt{8}{220.519}{446.629}{points}%
\TextAt{8}{255.703}{446.629}{of}%
\TextAt{18}{269.033}{446.629}{C}%
\TextAt{17}{277.667}{450.968}{2}%
\TextAt{8}{286.297}{446.629}{which}%
\TextAt{8}{319.783}{446.629}{satisfy}%
\TextAt{8}{356.33}{446.629}{the}%
\TextAt{8}{376.487}{446.629}{equation}%
\TextAt{13}{275.973}{417.628}{b}%
\TextAt{17}{280.95}{422.564}{2}%
\TextAt{16}{288.339}{417.628}{\Character{0}}%
\TextAt{13}{300.294}{417.628}{a}%
\TextAt{17}{306.439}{422.564}{3}%
\TextAt{9}{314.492}{417.628}{\Character{61}}%
\TextAt{9}{326.901}{417.628}{0}%
\TextAt{13}{332.771}{417.628}{\Character{58}}%
\NativeDiagram{%
\path[draw=black,line width=0.399bp] (195.248,-437.98) -- (257.779,-437.98);%
\path[draw=black,line width=0.399bp,line cap=round,line join=round] (255.906,-435.589) .. controls (256.284,-437.024) and (257.128,-437.701) .. (257.978,-437.98) .. controls (257.128,-438.259) and (256.284,-438.936) .. (255.906,-440.371);%
}%
\TextAt{13}{334.361}{333.446}{x}%
\NativeDiagram{%
\path[draw=black,line width=0.399bp] (202.051,-489.004) -- (202.051,-387.354);%
\path[draw=black,line width=0.399bp,line cap=round,line join=round] (199.66,-389.227) .. controls (201.095,-388.849) and (201.772,-388.005) .. (202.051,-387.155) .. controls (202.33,-388.005) and (203.007,-388.849) .. (204.442,-389.227);%
}%
\TextAt{13}{270.983}{393.553}{y}%
\NativeDiagram{%
\path[draw=black,line width=0.399bp] (241.727,-480.83) .. controls (241.727,-480.83) and (240.448,-478.773) .. (239.965,-478.007) .. controls (239.481,-477.241) and (238.714,-476.041) .. (238.242,-475.31) .. controls (237.769,-474.579) and (237.021,-473.434) .. (236.56,-472.737) .. controls (236.098,-472.04) and (235.368,-470.95) .. (234.917,-470.286) .. controls (234.467,-469.622) and (233.754,-468.584) .. (233.315,-467.952) .. controls (232.875,-467.321) and (232.18,-466.334) .. (231.752,-465.734) .. controls (231.324,-465.134) and (230.647,-464.197) .. (230.23,-463.627) .. controls (229.813,-463.058) and (229.154,-462.17) .. (228.748,-461.631) .. controls (228.342,-461.091) and (227.7,-460.251) .. (227.306,-459.74) .. controls (226.911,-459.23) and (226.287,-458.435) .. (225.903,-457.953) .. controls (225.52,-457.471) and (224.914,-456.721) .. (224.541,-456.267) .. controls (224.169,-455.813) and (223.58,-455.106) .. (223.219,-454.679) .. controls (222.858,-454.251) and (222.287,-453.587) .. (221.937,-453.185) .. controls (221.587,-452.783) and (221.034,-452.16) .. (220.695,-451.783) .. controls (220.356,-451.406) and (219.821,-450.822) .. (219.493,-450.47) .. controls (219.165,-450.117) and (218.648,-449.572) .. (218.331,-449.243) .. controls (218.014,-448.914) and (217.515,-448.405) .. (217.209,-448.099) .. controls (216.903,-447.792) and (216.422,-447.32) .. (216.127,-447.035) .. controls (215.833,-446.751) and (215.369,-446.312) .. (215.086,-446.049) .. controls (214.802,-445.785) and (214.356,-445.38) .. (214.084,-445.137) .. controls (213.812,-444.893) and (213.383,-444.519) .. (213.122,-444.296) .. controls (212.861,-444.072) and (212.451,-443.729) .. (212.201,-443.524) .. controls (211.951,-443.319) and (211.558,-443.005) .. (211.319,-442.817) .. controls (211.08,-442.63) and (210.706,-442.344) .. (210.478,-442.174) .. controls (210.25,-442.004) and (209.893,-441.744) .. (209.676,-441.59) .. controls (209.459,-441.436) and (209.121,-441.202) .. (208.915,-441.063) .. controls (208.709,-440.924) and (208.388,-440.714) .. (208.193,-440.59) .. controls (207.999,-440.466) and (207.695,-440.278) .. (207.512,-440.168) .. controls (207.328,-440.057) and (207.043,-439.892) .. (206.87,-439.794) .. controls (206.698,-439.697) and (206.43,-439.55) .. (206.269,-439.465) .. controls (206.108,-439.38) and (205.858,-439.253) .. (205.708,-439.179) .. controls (205.558,-439.105) and (205.326,-438.995) .. (205.187,-438.932) .. controls (205.048,-438.868) and (204.834,-438.775) .. (204.706,-438.721) .. controls (204.578,-438.667) and (204.382,-438.589) .. (204.265,-438.545) .. controls (204.148,-438.5) and (203.969,-438.435) .. (203.863,-438.398) .. controls (203.758,-438.362) and (203.597,-438.309) .. (203.503,-438.279) .. controls (203.408,-438.25) and (203.265,-438.209) .. (203.181,-438.186) .. controls (203.098,-438.163) and (202.973,-438.131) .. (202.901,-438.114) .. controls (202.828,-438.097) and (202.721,-438.073) .. (202.66,-438.061) .. controls (202.599,-438.049) and (202.509,-438.033) .. (202.459,-438.025) .. controls (202.409,-438.016) and (202.337,-438.006) .. (202.298,-438.001) .. controls (202.259,-437.996) and (202.206,-437.99) .. (202.178,-437.987) .. controls (202.15,-437.985) and (202.113,-437.983) .. (202.097,-437.982) .. controls (202.08,-437.981) and (202.062,-437.98) .. (202.056,-437.98) .. controls (202.051,-437.98) and (202.05,-437.98) .. (202.056,-437.98) .. controls (202.061,-437.98) and (202.079,-437.979) .. (202.095,-437.978) .. controls (202.112,-437.977) and (202.147,-437.975) .. (202.174,-437.973) .. controls (202.202,-437.97) and (202.255,-437.965) .. (202.294,-437.96) .. controls (202.333,-437.955) and (202.403,-437.945) .. (202.453,-437.936) .. controls (202.503,-437.928) and (202.592,-437.912) .. (202.653,-437.9) .. controls (202.714,-437.888) and (202.821,-437.865) .. (202.893,-437.848) .. controls (202.965,-437.831) and (203.089,-437.799) .. (203.173,-437.777) .. controls (203.256,-437.754) and (203.398,-437.713) .. (203.492,-437.684) .. controls (203.587,-437.654) and (203.747,-437.602) .. (203.852,-437.566) .. controls (203.957,-437.529) and (204.135,-437.465) .. (204.252,-437.42) .. controls (204.368,-437.376) and (204.564,-437.298) .. (204.692,-437.245) .. controls (204.819,-437.191) and (205.033,-437.098) .. (205.171,-437.035) .. controls (205.31,-436.972) and (205.542,-436.863) .. (205.692,-436.789) .. controls (205.841,-436.716) and (206.091,-436.589) .. (206.252,-436.504) .. controls (206.413,-436.419) and (206.68,-436.274) .. (206.852,-436.177) .. controls (207.024,-436.08) and (207.309,-435.914) .. (207.492,-435.804) .. controls (207.675,-435.694) and (207.978,-435.508) .. (208.172,-435.384) .. controls (208.366,-435.26) and (208.687,-435.05) .. (208.892,-434.912) .. controls (209.098,-434.774) and (209.436,-434.54) .. (209.652,-434.387) .. controls (209.869,-434.233) and (210.225,-433.975) .. (210.453,-433.805) .. controls (210.68,-433.635) and (211.054,-433.35) .. (211.293,-433.163) .. controls (211.532,-432.976) and (211.923,-432.663) .. (212.173,-432.458) .. controls (212.423,-432.254) and (212.833,-431.911) .. (213.094,-431.688) .. controls (213.355,-431.465) and (213.782,-431.092) .. (214.054,-430.85) .. controls (214.326,-430.607) and (214.772,-430.203) .. (215.055,-429.94) .. controls (215.338,-429.677) and (215.801,-429.24) .. (216.095,-428.956) .. controls (216.39,-428.672) and (216.87,-428.2) .. (217.176,-427.894) .. controls (217.481,-427.589) and (217.98,-427.081) .. (218.297,-426.753) .. controls (218.613,-426.425) and (219.13,-425.88) .. (219.458,-425.529) .. controls (219.785,-425.177) and (220.319,-424.594) .. (220.658,-424.218) .. controls (220.997,-423.842) and (221.549,-423.22) .. (221.899,-422.819) .. controls (222.249,-422.418) and (222.819,-421.755) .. (223.18,-421.328) .. controls (223.541,-420.901) and (224.129,-420.196) .. (224.501,-419.742) .. controls (224.873,-419.289) and (225.478,-418.54) .. (225.862,-418.059) .. controls (226.245,-417.578) and (226.868,-416.785) .. (227.263,-416.275) .. controls (227.657,-415.766) and (228.298,-414.926) .. (228.704,-414.388) .. controls (229.109,-413.849) and (229.768,-412.963) .. (230.185,-412.394) .. controls (230.601,-411.826) and (231.278,-410.89) .. (231.706,-410.291) .. controls (232.134,-409.692) and (232.828,-408.707) .. (233.267,-408.077) .. controls (233.706,-407.446) and (234.418,-406.409) .. (234.868,-405.746) .. controls (235.318,-405.083) and (236.048,-403.995) .. (236.509,-403.299) .. controls (236.97,-402.602) and (237.718,-401.46) .. (238.19,-400.73) .. controls (238.663,-400) and (239.429,-398.801) .. (239.912,-398.036) .. controls (240.395,-397.272) and (241.673,-395.217) .. (241.673,-395.217);%
}%
\TextAt{8}{79.226}{254.8}{The}%
\TextAt{8}{104.113}{254.8}{only}%
\TextAt{8}{130.638}{254.8}{singular}%
\TextAt{8}{175.765}{254.8}{point}%
\TextAt{8}{207.168}{254.8}{is}%
\TextAt{8}{219.765}{254.8}{the}%
\TextAt{8}{240.755}{254.8}{origin.}%
\TextAt{8}{280.981}{254.8}{It}%
\TextAt{8}{294.486}{254.8}{corresponds}%
\TextAt{8}{359.507}{254.8}{to}%
\TextAt{8}{374.651}{254.8}{the}%
\TextAt{8}{395.64}{254.8}{maximal}%
\TextAt{8}{444.282}{254.8}{ideal}%
\TextAt{14}{473.079}{254.8}{p}%
\TextAt{17}{479.062}{253.007}{0}%
\TextAt{9}{488.535}{254.8}{\Character{61}}%
\TextAt{9}{502.367}{254.8}{\Character{40}}%
\Vector{q 0 G 0 g 0.478 w 434.932 -512.379 m 441.585 -512.379 l S Q}%
\TextAt{13}{506.932}{254.8}{x\Character{59}}%
\Vector{q 0 G 0 g 0.478 w 446.829 -512.379 m 452.965 -512.379 l S Q}%
\TextAt{13}{518.829}{254.8}{y}%
\TextAt{9}{524.965}{254.8}{\Character{41}}%
\TextAt{8}{529.518}{254.8}{,}%
\TextAt{8}{79.226}{237.754}{where}%
\Vector{q 0 G 0 g 0.478 w 41.839 -529.425 m 48.491 -529.425 l S Q}%
\TextAt{13}{113.839}{237.754}{x}%
\TextAt{9}{125.181}{237.754}{\Character{61}}%
\TextAt{13}{138.976}{237.754}{x}%
\TextAt{8}{148.949}{237.754}{mod}%
\TextAt{9}{174.387}{237.754}{\Character{40}}%
\TextAt{13}{178.939}{237.754}{y}%
\TextAt{17}{185.076}{242.093}{2}%
\TextAt{16}{193.013}{237.754}{\Character{0}}%
\TextAt{13}{205.516}{237.754}{x}%
\TextAt{17}{212.168}{242.093}{3}%
\TextAt{9}{216.9}{237.754}{\Character{41}}%
\TextAt{8}{221.452}{237.754}{,}%
\Vector{q 0 G 0 g 0.478 w 157.606 -529.425 m 163.743 -529.425 l S Q}%
\TextAt{13}{229.606}{237.754}{y}%
\TextAt{9}{240.433}{237.754}{\Character{61}}%
\TextAt{13}{254.228}{237.754}{y}%
\TextAt{8}{263.685}{237.754}{mod}%
\TextAt{9}{289.123}{237.754}{\Character{40}}%
\TextAt{13}{293.675}{237.754}{y}%
\TextAt{17}{299.812}{242.093}{2}%
\TextAt{16}{307.749}{237.754}{\Character{0}}%
\TextAt{13}{320.252}{237.754}{x}%
\TextAt{17}{326.904}{242.093}{3}%
\TextAt{9}{331.636}{237.754}{\Character{41}}%
\NativeMathAt{340.879}{237.754}{11.955}{\in}%
\TextAt{14}{353.539}{237.754}{o}%
\TextAt{8}{359.383}{237.754}{.}%
\TextAt{8}{370.242}{237.754}{The}%
\TextAt{8}{395.099}{237.754}{maximal}%
\TextAt{8}{443.712}{237.754}{ideal}%
\TextAt{14}{472.478}{237.754}{p}%
\TextAt{17}{478.462}{235.961}{0}%
\TextAt{14}{483.194}{237.754}{o}%
\TextAt{27}{489.038}{235.961}{p}%
\TextAt{24}{493.027}{234.854}{0}%
\TextAt{8}{502.377}{237.754}{of}%
\TextAt{8}{516.511}{237.754}{the}%
\TextAt{8}{79.226}{220.708}{local}%
\TextAt{8}{106.576}{220.708}{ring}%
\TextAt{8}{130.352}{220.708}{is}%
\TextAt{8}{141.825}{220.708}{generated}%
\TextAt{8}{194.867}{220.708}{by}%
\TextAt{8}{210.836}{220.708}{the}%
\TextAt{8}{230.702}{220.708}{elements}%
\Vector{q 0 G 0 g 0.478 w 206.268 -546.471 m 212.92 -546.471 l S Q}%
\TextAt{13}{278.268}{220.708}{x\Character{59}}%
\Vector{q 0 G 0 g 0.478 w 218.164 -546.471 m 224.301 -546.471 l S Q}%
\TextAt{13}{290.164}{220.708}{y}%
\TextAt{8}{296.301}{220.708}{,}%
\TextAt{8}{303.218}{220.708}{and}%
\TextAt{8}{325.69}{220.708}{cannot}%
\TextAt{8}{363.776}{220.708}{be}%
\TextAt{8}{379.41}{220.708}{generated}%
\TextAt{8}{432.452}{220.708}{by}%
\TextAt{8}{448.421}{220.708}{a}%
\TextAt{8}{457.886}{220.708}{single}%
\TextAt{8}{490.173}{220.708}{element.}%
\TextAt{12}{97.159}{200.673}{5.}%
\TextAt{12}{114.236}{200.673}{Normalization.}%
\TextAt{8}{207.405}{200.673}{Passing}%
\TextAt{8}{249.824}{200.673}{to}%
\TextAt{8}{264.316}{200.673}{the}%
\TextAt{8}{284.654}{200.673}{normalization}%
\TextAt{20}{358.595}{201.209}{\Character{101}}%
\TextAt{14}{358.994}{200.673}{o}%
\TextAt{8}{368.917}{200.673}{of}%
\TextAt{8}{382.429}{200.673}{a}%
\TextAt{8}{392.366}{200.673}{one-dimensional}%
\TextAt{8}{478.457}{200.673}{noetherian}%
\TextAt{8}{79.226}{183.627}{integral}%
\TextAt{8}{120.966}{183.627}{domain}%
\TextAt{14}{161.726}{183.627}{o}%
\TextAt{8}{170.6}{183.627}{means,}%
\TextAt{8}{208.987}{183.627}{in}%
\TextAt{8}{221.772}{183.627}{geometric}%
\TextAt{8}{274.236}{183.627}{terms,}%
\TextAt{8}{309.372}{183.627}{taking}%
\TextAt{8}{344.608}{183.627}{the}%
\TextAt{12}{363.897}{183.627}{resolution}%
\TextAt{12}{425.448}{183.627}{of}%
\TextAt{12}{439.776}{183.627}{the}%
\TextAt{12}{462.102}{183.627}{singularities}%
\TextAt{8}{79.226}{166.581}{that}%
\TextAt{8}{105.358}{166.581}{were}%
\TextAt{8}{133.104}{166.581}{just}%
\TextAt{8}{157.012}{166.581}{discussed.}%
\TextAt{8}{214.758}{166.581}{Indeed,}%
\TextAt{8}{256.993}{166.581}{if}%
\TextAt{13}{268.48}{166.581}{X}%
\TextAt{9}{283.756}{166.581}{\Character{61}}%
\TextAt{8}{297.482}{166.581}{Spec}%
\TextAt{9}{320.89}{166.581}{\Character{40}}%
\TextAt{14}{325.442}{166.581}{o}%
\TextAt{9}{331.287}{166.581}{\Character{41}}%
\TextAt{8}{340.499}{166.581}{and}%
\TextAt{20}{367.007}{169.603}{\Character{101}}%
\TextAt{13}{364.025}{166.581}{X}%
\TextAt{9}{379.301}{166.581}{\Character{61}}%
\TextAt{8}{393.027}{166.581}{Spec}%
\TextAt{9}{416.435}{166.581}{\Character{40}}%
\TextAt{20}{420.589}{167.117}{\Character{101}}%
\TextAt{14}{420.988}{166.581}{o}%
\TextAt{9}{426.832}{166.581}{\Character{41}}%
\TextAt{8}{431.384}{166.581}{,}%
\TextAt{8}{439.487}{166.581}{then}%
\TextAt{8}{466.91}{166.581}{the}%
\TextAt{8}{487.83}{166.581}{inclusion}%
\TextAt{14}{79.226}{149.535}{o}%
\NativeArrowAt{88.391}{149.535}{11.955}{13.214}{\hookrightarrow}%
\TextAt{20}{104.528}{150.071}{\Character{101}}%
\TextAt{14}{104.926}{149.535}{o}%
\TextAt{8}{114.668}{149.535}{induces}%
\TextAt{8}{156.344}{149.535}{a}%
\TextAt{8}{166.099}{149.535}{morphism}%
\TextAt{13}{220.794}{149.535}{f}%
\TextAt{9}{231.161}{149.535}{\Character{58}}%
\TextAt{20}{240.716}{152.557}{\Character{101}}%
\TextAt{13}{237.734}{149.535}{X}%
\NativeMathAt{251.71}{149.535}{11.955}{\rightarrow}%
\TextAt{13}{266.986}{149.535}{X}%
\TextAt{8}{277.641}{149.535}{.}%
\TextAt{8}{300.14}{43.033}{85}%
\Destination{78.23}{809.815}{nav.page.85}%
\Destination{79.36}{776.155}{nav.print.91}%
}

% Page 98 (body page 86)
\Page{%
\TextAt{2}{79.226}{796.943}{I.}%
\TextAt{2}{89.83}{796.943}{Algebraic}%
\TextAt{2}{138.648}{796.943}{Integers}%
\Vector{q 0 G 0 g 0.299 w 7.226 18.729 m 460.769 18.729 l S Q}%
\NativeDiagram{%
\path[draw=black,line width=0.399bp] (173.663,-25.215) .. controls (139.647,-25.215) and (114.134,-37.971) .. (136.812,-43.64) .. controls (156.655,-49.31) and (193.506,-18.128) .. (216.183,-1.12);%
\path[draw=black,line width=0.399bp] (197.757,-28.05) .. controls (238.861,-26.632) and (258.703,-39.388) .. (262.955,-59.231) .. controls (270.042,-64.9) and (275.711,-54.979) .. (268.624,-49.31);%
\path[draw=black,line width=0.399bp] (275.711,-49.31) .. controls (298.389,-35.136) and (323.901,-35.136) .. (335.24,-36.553);%
\fill (262.955,-59.231) circle[radius=1.48bp];%
}%
\TextAt{20}{414.406}{734.641}{\Character{101}}%
\TextAt{13}{411.424}{731.619}{X}%
\NativeDiagram{%
\path[fill=black] (195.802,-17.759) (197.296,-17.759) .. controls (197.296,-16.934) and (196.627,-16.265) .. (195.802,-16.265) .. controls (194.976,-16.265) and (194.307,-16.934) .. (194.307,-17.759) .. controls (194.307,-18.585) and (194.976,-19.254) .. (195.802,-19.254) .. controls (196.627,-19.254) and (197.296,-18.585) .. (197.296,-17.759) -- cycle (195.802,-17.759);%
\path[fill=black] (197.757,-28.05) (199.252,-28.05) .. controls (199.252,-27.224) and (198.583,-26.555) .. (197.757,-26.555) .. controls (196.932,-26.555) and (196.263,-27.224) .. (196.263,-28.05) .. controls (196.263,-28.875) and (196.932,-29.544) .. (197.757,-29.544) .. controls (198.583,-29.544) and (199.252,-28.875) .. (199.252,-28.05) -- cycle (197.757,-28.05);%
\path[draw=black,line width=0.399bp] (195.773,-91.83) .. controls (145.316,-90.412) and (114.134,-106.003) .. (136.812,-111.673) .. controls (156.655,-117.342) and (193.506,-86.16) .. (216.183,-69.152);%
\path[draw=black,line width=0.399bp] (195.773,-91.83) .. controls (238.861,-90.412) and (261.538,-103.168) .. (262.955,-123.011) .. controls (287.05,-104.586) and (318.231,-103.168) .. (335.24,-106.003);%
}%
\TextAt{13}{411.424}{663.912}{X}%
\NativeDiagram{%
\path[fill=black] (195.773,-91.83) (197.268,-91.83) .. controls (197.268,-91.004) and (196.599,-90.335) .. (195.773,-90.335) .. controls (194.948,-90.335) and (194.279,-91.004) .. (194.279,-91.83) .. controls (194.279,-92.655) and (194.948,-93.324) .. (195.773,-93.324) .. controls (196.599,-93.324) and (197.268,-92.655) .. (197.268,-91.83) -- cycle (195.773,-91.83);%
\path[fill=black] (262.955,-123.011) (264.45,-123.011) .. controls (264.45,-122.186) and (263.781,-121.517) .. (262.955,-121.517) .. controls (262.13,-121.517) and (261.461,-122.186) .. (261.461,-123.011) .. controls (261.461,-123.836) and (262.13,-124.505) .. (262.955,-124.505) .. controls (263.781,-124.505) and (264.45,-123.836) .. (264.45,-123.011) -- cycle (262.955,-123.011);%
\path[draw=black,line width=0.399bp,dash pattern=on 2.989bp off 2.989bp] (195.773,-35.136) -- (195.773,-82.927);%
\path[draw=black,line width=0.399bp,line cap=round,line join=round] (198.164,-81.054) .. controls (196.73,-81.433) and (196.052,-82.277) .. (195.773,-83.127) .. controls (195.494,-82.277) and (194.817,-81.433) .. (193.382,-81.054);%
\path[draw=black,line width=0.399bp] (304.058,-57.813) -- (304.058,-94.266);%
\path[draw=black,line width=0.399bp,line cap=round,line join=round] (306.449,-92.393) .. controls (305.014,-92.771) and (304.337,-93.615) .. (304.058,-94.465) .. controls (303.779,-93.615) and (303.102,-92.771) .. (301.667,-92.393);%
}%
\TextAt{13}{380.242}{694.772}{f}%
\NativeDiagram{%
\path[draw=black,line width=0.399bp] (345.161,-46.475) -- (345.161,-95.683);%
\path[draw=black,line width=0.399bp,line cap=round,line join=round] (347.552,-93.81) .. controls (346.118,-94.189) and (345.44,-95.032) .. (345.161,-95.882) .. controls (344.882,-95.032) and (344.205,-94.189) .. (342.77,-93.81);%
}%
\TextAt{8}{79.226}{619.498}{Since}%
\TextAt{20}{110.607}{620.033}{\Character{101}}%
\TextAt{14}{111.006}{619.498}{o}%
\TextAt{8}{121.969}{619.498}{is}%
\TextAt{8}{134.955}{619.498}{a}%
\TextAt{8}{145.932}{619.498}{Dedekind}%
\TextAt{8}{199.339}{619.498}{domain,}%
\TextAt{8}{245.746}{619.498}{the}%
\TextAt{8}{267.124}{619.498}{scheme}%
\TextAt{20}{311.378}{622.52}{\Character{101}}%
\TextAt{13}{308.396}{619.498}{X}%
\TextAt{8}{324.171}{619.498}{is}%
\TextAt{8}{337.156}{619.498}{to}%
\TextAt{8}{352.689}{619.498}{be}%
\TextAt{8}{369.835}{619.498}{considered}%
\TextAt{8}{428.346}{619.498}{as}%
\TextAt{8}{443.938}{619.498}{smooth.}%
\TextAt{8}{493.52}{619.498}{If}%
\TextAt{14}{506.434}{619.498}{p}%
\TextAt{20}{512.019}{620.033}{\Character{101}}%
\TextAt{14}{512.417}{619.498}{o}%
\TextAt{9}{523.665}{619.498}{\Character{61}}%
\TextAt{20}{78.897}{603.569}{\Character{101}}%
\TextAt{14}{79.226}{602.452}{p}%
\TextAt{21}{85.209}{607.677}{e}%
\TextAt{24}{89.138}{606.57}{1}%
\TextAt{17}{85.209}{599.521}{1}%
\TextAt{16}{95.779}{602.452}{\Character{1}}%
\TextAt{16}{101.075}{602.452}{\Character{1}}%
\TextAt{16}{106.383}{602.452}{\Character{1}}%
\TextAt{20}{111.39}{603.569}{\Character{101}}%
\TextAt{14}{111.719}{602.452}{p}%
\TextAt{21}{117.703}{606.79}{e}%
\TextAt{22}{121.631}{605.794}{r}%
\TextAt{21}{117.703}{599.497}{r}%
\TextAt{8}{129.971}{602.452}{is}%
\TextAt{8}{141.522}{602.452}{the}%
\TextAt{8}{161.465}{602.452}{prime}%
\TextAt{8}{194.415}{602.452}{factorization}%
\TextAt{8}{262.179}{602.452}{of}%
\TextAt{14}{275.295}{602.452}{p}%
\TextAt{8}{284.962}{602.452}{in}%
\TextAt{20}{298.003}{602.987}{\Character{101}}%
\TextAt{14}{298.402}{602.452}{o}%
\TextAt{8}{304.246}{602.452}{,}%
\TextAt{8}{311.224}{602.452}{then}%
\TextAt{20}{337.342}{603.569}{\Character{101}}%
\TextAt{14}{337.671}{602.452}{p}%
\TextAt{17}{343.654}{600.659}{1}%
\TextAt{13}{348.386}{602.452}{\Character{59}}%
\TextAt{13}{353.61}{602.452}{\Character{58}}%
\TextAt{13}{358.847}{602.452}{\Character{58}}%
\TextAt{13}{364.083}{602.452}{\Character{58}}%
\TextAt{13}{369.308}{602.452}{\Character{59}}%
\TextAt{20}{374.278}{603.569}{\Character{101}}%
\TextAt{14}{374.607}{602.452}{p}%
\TextAt{21}{380.59}{600.659}{r}%
\TextAt{8}{388.829}{602.452}{are}%
\TextAt{8}{408.127}{602.452}{the}%
\TextAt{8}{428.07}{602.452}{different}%
\TextAt{8}{474.027}{602.452}{points}%
\TextAt{8}{508.998}{602.452}{of}%
\TextAt{20}{525.097}{605.474}{\Character{101}}%
\TextAt{13}{522.114}{602.452}{X}%
\ErrataAt{79.19}{585.382}{455.81}{that are mapped to $\mathfrak p$ by $f$. If the normalization of $\mathfrak o_{\mathfrak p}$ is finite, then $\mathfrak o_{\mathfrak p}$ is a discrete valuation ring}%
\ErrataAt{79.19}{568.336}{455.81}{if and only if $r=e_1=f_1=1$, where $f_1=[\widetilde{\mathfrak o}/\widetilde{\mathfrak p}_1:\mathfrak o/\mathfrak p]$.}%
\TextAt{12}{97.159}{548.325}{6.}%
\TextAt{12}{115.706}{548.325}{Extensions.}%
\TextAt{8}{189.794}{548.325}{Let}%
\TextAt{14}{211.371}{548.325}{o}%
\TextAt{8}{221.72}{548.325}{be}%
\TextAt{8}{238.252}{548.325}{a}%
\TextAt{8}{248.615}{548.325}{Dedekind}%
\TextAt{8}{301.407}{548.325}{domain}%
\TextAt{8}{343.642}{548.325}{with}%
\TextAt{8}{370.91}{548.325}{field}%
\TextAt{8}{396.874}{548.325}{of}%
\TextAt{8}{410.812}{548.325}{fractions}%
\TextAt{13}{459.288}{548.325}{K}%
\TextAt{8}{470.106}{548.325}{.}%
\TextAt{8}{480.377}{548.325}{Let}%
\TextAt{13}{501.954}{548.325}{L}%
\TextAt{16}{514.275}{548.325}{\Character{106}}%
\TextAt{13}{521.952}{548.325}{K}%
\TextAt{8}{79.226}{531.279}{be}%
\TextAt{8}{95.745}{531.279}{a}%
\TextAt{8}{106.095}{531.279}{finite}%
\TextAt{8}{136.601}{531.279}{separable}%
\TextAt{8}{188.639}{531.279}{extension,}%
\TextAt{8}{244.399}{531.279}{and}%
\TextAt{16}{267.756}{531.279}{\Character{79}}%
\TextAt{8}{282.098}{531.279}{the}%
\TextAt{8}{302.849}{531.279}{integral}%
\TextAt{8}{346.051}{531.279}{closure}%
\TextAt{8}{385.727}{531.279}{of}%
\TextAt{14}{399.652}{531.279}{o}%
\TextAt{8}{409.988}{531.279}{in}%
\TextAt{13}{424.235}{531.279}{L}%
\TextAt{8}{432.2}{531.279}{.}%
\TextAt{8}{442.431}{531.279}{Let}%
\TextAt{13}{463.995}{531.279}{Y}%
\TextAt{9}{477.722}{531.279}{\Character{61}}%
\TextAt{8}{491.161}{531.279}{Spec}%
\TextAt{9}{514.569}{531.279}{\Character{40}}%
\TextAt{14}{519.121}{531.279}{o}%
\TextAt{9}{524.965}{531.279}{\Character{41}}%
\TextAt{8}{529.518}{531.279}{,}%
\TextAt{13}{79.226}{514.233}{X}%
\TextAt{9}{93.202}{514.233}{\Character{61}}%
\TextAt{8}{105.628}{514.233}{Spec}%
\TextAt{9}{129.036}{514.233}{\Character{40}}%
\TextAt{16}{133.588}{514.233}{\Character{79}}%
\TextAt{9}{143.438}{514.233}{\Character{41}}%
\TextAt{8}{147.99}{514.233}{,}%
\TextAt{8}{155.139}{514.233}{and}%
\TextAt{13}{274.552}{497.187}{f}%
\TextAt{9}{284.92}{497.187}{\Character{58}}%
\TextAt{13}{291.492}{497.187}{X}%
\NativeArrowAt{305.468}{497.187}{11.955}{19.26}{\longrightarrow}%
\TextAt{13}{328.05}{497.187}{Y}%
\TextAt{8}{79.226}{473.666}{the}%
\TextAt{8}{99.383}{473.666}{morphism}%
\TextAt{8}{154.077}{473.666}{induced}%
\TextAt{8}{197.642}{473.666}{by}%
\TextAt{8}{213.901}{473.666}{the}%
\TextAt{8}{234.058}{473.666}{inclusion}%
\TextAt{14}{282.894}{473.666}{o}%
\NativeArrowAt{292.06}{473.666}{11.955}{13.214}{\hookrightarrow}%
\TextAt{16}{308.586}{473.666}{\Character{79}}%
\TextAt{8}{318.444}{473.666}{.}%
\TextAt{8}{326.893}{473.666}{If}%
\TextAt{14}{338.585}{473.666}{p}%
\TextAt{8}{348.465}{473.666}{is}%
\TextAt{8}{360.229}{473.666}{a}%
\TextAt{8}{369.985}{473.666}{maximal}%
\TextAt{8}{417.793}{473.666}{ideal}%
\TextAt{8}{445.757}{473.666}{of}%
\TextAt{14}{459.087}{473.666}{o}%
\TextAt{8}{468.828}{473.666}{and}%
\TextAt{14}{262.768}{444.665}{p}%
\TextAt{16}{268.751}{444.665}{\Character{79}}%
\TextAt{9}{281.921}{444.665}{\Character{61}}%
\TextAt{14}{294.347}{444.665}{P}%
\TextAt{21}{304.24}{449.889}{e}%
\TextAt{24}{308.168}{448.782}{1}%
\TextAt{17}{304.24}{441.733}{1}%
\TextAt{16}{314.81}{444.665}{\Character{1}}%
\TextAt{16}{320.106}{444.665}{\Character{1}}%
\TextAt{16}{325.414}{444.665}{\Character{1}}%
\TextAt{14}{330.75}{444.665}{P}%
\TextAt{21}{340.644}{449.601}{e}%
\TextAt{22}{344.572}{448.605}{r}%
\TextAt{21}{340.644}{441.709}{r}%
\TextAt{8}{79.226}{415.663}{the}%
\TextAt{8}{98.457}{415.663}{prime}%
\TextAt{8}{130.694}{415.663}{decomposition}%
\TextAt{8}{206.903}{415.663}{of}%
\TextAt{14}{219.307}{415.663}{p}%
\TextAt{8}{228.262}{415.663}{in}%
\TextAt{16}{240.989}{415.663}{\Character{79}}%
\TextAt{8}{250.838}{415.663}{,}%
\TextAt{8}{257.247}{415.663}{then}%
\TextAt{14}{282.981}{415.663}{P}%
\TextAt{17}{292.874}{413.87}{1}%
\TextAt{13}{297.606}{415.663}{\Character{59}}%
\TextAt{13}{302.83}{415.663}{\Character{58}}%
\TextAt{13}{308.067}{415.663}{\Character{58}}%
\TextAt{13}{313.303}{415.663}{\Character{58}}%
\TextAt{13}{318.528}{415.663}{\Character{59}}%
\TextAt{14}{323.827}{415.663}{P}%
\TextAt{21}{333.72}{413.87}{r}%
\TextAt{8}{341.247}{415.663}{are}%
\TextAt{8}{359.831}{415.663}{the}%
\TextAt{8}{379.062}{415.663}{different}%
\TextAt{8}{424.307}{415.663}{points}%
\TextAt{8}{458.565}{415.663}{of}%
\TextAt{13}{470.969}{415.663}{X}%
\TextAt{8}{484.596}{415.663}{which}%
\TextAt{8}{517.156}{415.663}{are}%
\TextAt{8}{79.226}{398.617}{mapped}%
\TextAt{8}{123.303}{398.617}{to}%
\TextAt{14}{137.146}{398.617}{p}%
\TextAt{8}{146.559}{398.617}{by}%
\TextAt{13}{162.35}{398.617}{f}%
\TextAt{8}{169.397}{398.617}{.}%
\TextAt{8}{177.689}{398.617}{The}%
\TextAt{8}{201.275}{398.617}{morphism}%
\TextAt{13}{255.503}{398.617}{f}%
\TextAt{8}{265.979}{398.617}{is}%
\TextAt{8}{277.275}{398.617}{a}%
\TextAt{8}{286.563}{398.617}{“ramified}%
\TextAt{8}{337.144}{398.617}{covering.}%
\TextAt{8}{380.733}{398.617}{”}%
\TextAt{8}{391.296}{398.617}{It}%
\TextAt{8}{403.501}{398.617}{is}%
\TextAt{8}{414.797}{398.617}{graphically}%
\TextAt{8}{474.5}{398.617}{represented}%
\TextAt{8}{79.226}{381.572}{by}%
\TextAt{8}{95.485}{381.572}{the}%
\TextAt{8}{115.642}{381.572}{following}%
\TextAt{8}{165.076}{381.572}{picture:}%
\NativeDiagram{%
\path[draw=black,line width=0.399bp] (125.473,-413.278) .. controls (179.332,-409.026) and (210.513,-416.112) .. (238.86,-435.955) .. controls (267.207,-458.633) and (295.554,-464.302) .. (323.901,-464.302);%
\path[draw=black,line width=0.399bp] (125.473,-428.868) .. controls (139.646,-428.868) and (148.151,-441.624) .. (165.158,-441.624) .. controls (193.505,-464.302) and (210.513,-447.294) .. (238.86,-435.955) .. controls (272.876,-414.695) and (301.223,-421.782) .. (323.901,-427.451);%
\path[draw=black,line width=0.399bp] (125.473,-452.963) .. controls (139.646,-447.294) and (148.151,-444.459) .. (165.158,-441.624) .. controls (199.175,-437.373) and (216.183,-430.286) .. (238.86,-435.955) .. controls (267.207,-444.459) and (304.058,-447.294) .. (323.901,-447.294);%
\path[fill=black] (165.158,-441.624) (166.852,-441.624) .. controls (166.852,-440.689) and (166.094,-439.931) .. (165.158,-439.931) .. controls (164.223,-439.931) and (163.465,-440.689) .. (163.465,-441.624) .. controls (163.465,-442.56) and (164.223,-443.318) .. (165.158,-443.318) .. controls (166.094,-443.318) and (166.852,-442.56) .. (166.852,-441.624) -- cycle (165.158,-441.624);%
\path[fill=black] (238.86,-435.955) (240.554,-435.955) .. controls (240.554,-435.02) and (239.796,-434.261) .. (238.86,-434.261) .. controls (237.925,-434.261) and (237.167,-435.02) .. (237.167,-435.955) .. controls (237.167,-436.89) and (237.925,-437.649) .. (238.86,-437.649) .. controls (239.796,-437.649) and (240.554,-436.89) .. (240.554,-435.955) -- cycle (238.86,-435.955);%
}%
\TextAt{13}{400.085}{322.621}{X}%
\NativeDiagram{%
\path[draw=black,line width=0.399bp] (199.175,-469.971) -- (199.175,-486.581);%
\path[draw=black,line width=0.399bp,line cap=round,line join=round] (201.566,-484.708) .. controls (200.131,-485.086) and (199.454,-485.93) .. (199.175,-486.78) .. controls (198.896,-485.93) and (198.219,-485.086) .. (196.784,-484.708);%
\path[draw=black,line width=0.399bp] (125.473,-495.484) -- (323.901,-495.484);%
\path[fill=black] (165.158,-495.484) (166.852,-495.484) .. controls (166.852,-494.548) and (166.094,-493.79) .. (165.158,-493.79) .. controls (164.223,-493.79) and (163.465,-494.548) .. (163.465,-495.484) .. controls (163.465,-496.419) and (164.223,-497.177) .. (165.158,-497.177) .. controls (166.094,-497.177) and (166.852,-496.419) .. (166.852,-495.484) -- cycle (165.158,-495.484);%
\path[fill=black] (238.86,-495.484) (240.554,-495.484) .. controls (240.554,-494.548) and (239.796,-493.79) .. (238.86,-493.79) .. controls (237.925,-493.79) and (237.167,-494.548) .. (237.167,-495.484) .. controls (237.167,-496.419) and (237.925,-497.177) .. (238.86,-497.177) .. controls (239.796,-497.177) and (240.554,-496.419) .. (240.554,-495.484) -- cycle (238.86,-495.484);%
}%
\TextAt{13}{400.085}{274.432}{Y}%
\NativeDiagram{%
\path[draw=black,line width=0.399bp] (330.987,-457.215) -- (330.987,-486.581);%
\path[draw=black,line width=0.399bp,line cap=round,line join=round] (333.378,-484.708) .. controls (331.944,-485.086) and (331.266,-485.93) .. (330.987,-486.78) .. controls (330.708,-485.93) and (330.031,-485.086) .. (328.596,-484.708);%
\path[draw=black,line width=0.399bp] (165.158,-512.491) -- (165.158,-501.551);%
\path[draw=black,line width=0.399bp,line cap=round,line join=round] (162.767,-503.424) .. controls (164.202,-503.046) and (164.879,-502.202) .. (165.158,-501.352) .. controls (165.437,-502.202) and (166.115,-503.046) .. (167.549,-503.424);%
\path[draw=black,line width=0.399bp] (238.86,-512.491) -- (238.86,-501.551);%
\path[draw=black,line width=0.399bp,line cap=round,line join=round] (236.469,-503.424) .. controls (237.904,-503.046) and (238.581,-502.202) .. (238.86,-501.352) .. controls (239.139,-502.202) and (239.817,-503.046) .. (241.251,-503.424);%
}%
\TextAt{8}{235.603}{254.208}{ramified}%
\TextAt{8}{281.128}{254.208}{points}%
\TextAt{8}{79.226}{217.907}{This}%
\TextAt{8}{105.613}{217.907}{picture,}%
\TextAt{8}{148.265}{217.907}{however,}%
\TextAt{8}{196.13}{217.907}{is}%
\TextAt{8}{207.561}{217.907}{a}%
\TextAt{8}{216.983}{217.907}{fair}%
\TextAt{8}{237.786}{217.907}{rendering}%
\TextAt{8}{289.482}{217.907}{of}%
\TextAt{8}{302.479}{217.907}{the}%
\TextAt{8}{322.302}{217.907}{algebraic}%
\TextAt{8}{371.403}{217.907}{situation}%
\TextAt{8}{419.919}{217.907}{only}%
\TextAt{8}{445.277}{217.907}{in}%
\TextAt{8}{458.597}{217.907}{the}%
\TextAt{8}{478.42}{217.907}{case}%
\TextAt{8}{502.858}{217.907}{where}%
\TextAt{8}{79.226}{200.861}{the}%
\TextAt{8}{98.758}{200.861}{residue}%
\TextAt{8}{137.86}{200.861}{class}%
\TextAt{8}{164.672}{200.861}{fields}%
\TextAt{8}{194.018}{200.861}{of}%
\TextAt{14}{206.723}{200.861}{o}%
\TextAt{8}{215.84}{200.861}{are}%
\TextAt{8}{234.726}{200.861}{algebraically}%
\TextAt{8}{302.078}{200.861}{closed}%
\TextAt{8}{335.98}{200.861}{(like}%
\TextAt{8}{361.369}{200.861}{for}%
\TextAt{8}{378.629}{200.861}{the}%
\TextAt{8}{398.16}{200.861}{ring}%
\TextAt{18}{421.601}{200.861}{C}%
\TextAt{9}{430.235}{200.861}{\Character{91}}%
\TextAt{13}{433.487}{200.861}{x}%
\TextAt{9}{440.139}{200.861}{\Character{93}}%
\TextAt{8}{443.391}{200.861}{).}%
\TextAt{8}{456.186}{200.861}{Then,}%
\TextAt{8}{489.495}{200.861}{from}%
\TextAt{8}{516.51}{200.861}{the}%
\TextAt{8}{79.226}{183.815}{fundamental}%
\TextAt{8}{147.152}{183.815}{identity}%
\TextAt{20}{190.69}{192.782}{\Character{80}}%
\TextAt{21}{203.309}{180.328}{i}%
\TextAt{13}{208.683}{183.815}{e}%
\TextAt{21}{214.109}{182.022}{i}%
\TextAt{13}{217.49}{183.815}{f}%
\TextAt{21}{223.26}{182.022}{i}%
\TextAt{9}{230.446}{183.815}{\Character{61}}%
\TextAt{13}{243.355}{183.815}{n}%
\TextAt{8}{250.343}{183.815}{,}%
\TextAt{8}{257.847}{183.815}{there}%
\TextAt{8}{288.043}{183.815}{are}%
\TextAt{8}{307.837}{183.815}{exactly}%
\TextAt{13}{348.446}{183.815}{n}%
\TextAt{9}{359.238}{183.815}{\Character{61}}%
\TextAt{9}{372.138}{183.815}{\Character{91}}%
\TextAt{13}{375.399}{183.815}{L}%
\TextAt{9}{387.168}{183.815}{\Character{58}}%
\TextAt{13}{394.225}{183.815}{K}%
\TextAt{9}{405.042}{183.815}{\Character{93}}%
\TextAt{8}{412.475}{183.815}{points}%
\TextAt{14}{447.943}{183.815}{P}%
\TextAt{17}{457.837}{182.022}{1}%
\TextAt{13}{462.569}{183.815}{\Character{59}}%
\TextAt{13}{467.793}{183.815}{\Character{58}}%
\TextAt{13}{473.03}{183.815}{\Character{58}}%
\TextAt{13}{478.254}{183.815}{\Character{58}}%
\TextAt{13}{483.491}{183.815}{\Character{59}}%
\TextAt{14}{488.79}{183.815}{P}%
\TextAt{21}{498.683}{182.022}{n}%
\TextAt{8}{508.5}{183.815}{of}%
\TextAt{13}{522.114}{183.815}{X}%
\TextAt{8}{79.226}{166.769}{lying}%
\TextAt{8}{107.26}{166.769}{above}%
\TextAt{8}{139.526}{166.769}{each}%
\TextAt{8}{164.954}{166.769}{point}%
\TextAt{14}{194.613}{166.769}{p}%
\TextAt{8}{203.584}{166.769}{of}%
\TextAt{13}{216.005}{166.769}{Y}%
\TextAt{8}{225.398}{166.769}{,}%
\TextAt{8}{231.82}{166.769}{except}%
\TextAt{8}{267.648}{166.769}{when}%
\TextAt{14}{297.296}{166.769}{p}%
\TextAt{8}{306.267}{166.769}{is}%
\TextAt{8}{317.121}{166.769}{ramified}%
\TextAt{8}{361.737}{166.769}{in}%
\TextAt{16}{374.48}{166.769}{\Character{79}}%
\TextAt{8}{384.33}{166.769}{.}%
\TextAt{8}{392.475}{166.769}{At}%
\TextAt{8}{408.47}{166.769}{a}%
\TextAt{8}{417.316}{166.769}{point}%
\TextAt{14}{446.976}{166.769}{p}%
\TextAt{8}{455.947}{166.769}{of}%
\TextAt{8}{468.367}{166.769}{ramification,}%
\TextAt{8}{79.226}{149.723}{several}%
\TextAt{8}{116.819}{149.723}{of}%
\TextAt{8}{129.307}{149.723}{the}%
\TextAt{8}{148.621}{149.723}{points}%
\TextAt{14}{182.962}{149.723}{P}%
\TextAt{17}{192.855}{147.93}{1}%
\TextAt{13}{197.588}{149.723}{\Character{59}}%
\TextAt{13}{202.812}{149.723}{\Character{58}}%
\TextAt{13}{208.049}{149.723}{\Character{58}}%
\TextAt{13}{213.273}{149.723}{\Character{58}}%
\TextAt{13}{218.51}{149.723}{\Character{59}}%
\TextAt{14}{223.808}{149.723}{P}%
\TextAt{21}{233.702}{147.93}{n}%
\TextAt{8}{242.393}{149.723}{coalesce.}%
\TextAt{8}{290.945}{149.723}{This}%
\TextAt{8}{316.822}{149.723}{also}%
\TextAt{8}{339.459}{149.723}{explains}%
\TextAt{8}{383.879}{149.723}{the}%
\TextAt{8}{403.193}{149.723}{terminology}%
\TextAt{8}{467.075}{149.723}{of}%
\TextAt{8}{479.563}{149.723}{ideals}%
\TextAt{8}{511.298}{149.723}{that}%
\TextAt{8}{79.226}{132.677}{“ramify.}%
\TextAt{8}{118.583}{132.677}{”}%
\TextAt{8}{97.159}{112.643}{If}%
\TextAt{13}{108.959}{112.643}{L}%
\TextAt{16}{120.429}{112.643}{\Character{106}}%
\TextAt{13}{127.254}{112.643}{K}%
\TextAt{8}{142.078}{112.643}{is}%
\TextAt{8}{153.949}{112.643}{Galois}%
\TextAt{8}{189.971}{112.643}{with}%
\TextAt{8}{216.739}{112.643}{Galois}%
\TextAt{8}{252.76}{112.643}{group}%
\TextAt{13}{286.044}{112.643}{G}%
\TextAt{9}{298.782}{112.643}{\Character{61}}%
\TextAt{13}{311.392}{112.643}{G}%
\TextAt{9}{320.626}{112.643}{\Character{40}}%
\TextAt{13}{325.178}{112.643}{L}%
\TextAt{16}{336.647}{112.643}{\Character{106}}%
\TextAt{13}{343.473}{112.643}{K}%
\TextAt{9}{354.291}{112.643}{\Character{41}}%
\TextAt{8}{358.843}{112.643}{,}%
\TextAt{8}{366.128}{112.643}{then}%
\TextAt{8}{392.896}{112.643}{every}%
\TextAt{8}{423.896}{112.643}{automorphism}%
\TextAt{13}{501.473}{112.643}{\Character{27}}%
\NativeMathAt{512.061}{112.643}{11.955}{\in}%
\TextAt{13}{523.536}{112.643}{G}%
\TextAt{8}{79.226}{95.597}{induces}%
\TextAt{8}{121.773}{95.597}{via}%
\TextAt{13}{141.832}{95.597}{\Character{27}}%
\TextAt{9}{153.719}{95.597}{\Character{58}}%
\TextAt{16}{161.776}{95.597}{\Character{79}}%
\NativeMathAt{176.421}{95.597}{11.955}{\rightarrow}%
\TextAt{16}{193.182}{95.597}{\Character{79}}%
\TextAt{8}{207.81}{95.597}{an}%
\TextAt{8}{224.94}{95.597}{automorphism}%
\TextAt{8}{303.28}{95.597}{of}%
\TextAt{8}{317.481}{95.597}{schemes}%
\TextAt{13}{363.017}{95.597}{\Character{27}}%
\TextAt{9}{374.904}{95.597}{\Character{58}}%
\TextAt{13}{382.961}{95.597}{X}%
\NativeMathAt{398.421}{95.597}{11.955}{\rightarrow}%
\TextAt{13}{415.181}{95.597}{X}%
\TextAt{8}{425.837}{95.597}{.}%
\TextAt{8}{436.898}{95.597}{Since}%
\TextAt{8}{468.326}{95.597}{the}%
\TextAt{8}{489.354}{95.597}{ring}%
\TextAt{14}{514.29}{95.597}{o}%
\TextAt{8}{524.903}{95.597}{is}%
\TextAt{8}{300.14}{43.033}{86}%
\Destination{78.23}{809.815}{nav.page.86}%
\Destination{79.36}{634.455}{nav.print.92}%
\Destination{79.36}{215.855}{nav.print.93}%
\Destination{79.226}{600}{errata.188}%
\Destination{79.226}{750}{errata.189}%
}

% Page 99 (body page 87)
\Page{%
\TextAt{2}{79.226}{796.943}{I.}%
\TextAt{2}{89.83}{796.943}{Algebraic}%
\TextAt{2}{138.648}{796.943}{Integers}%
\Vector{q 0 G 0 g 0.299 w 7.226 18.729 m 460.769 18.729 l S Q}%
\TextAt{8}{79.226}{761.124}{fixed,}%
\TextAt{8}{110.764}{761.124}{the}%
\TextAt{8}{130.92}{761.124}{diagram}%
\TextAt{13}{251.647}{745.127}{X}%
\TextAt{13}{349.691}{745.127}{X}%
\TextAt{13}{301.3}{706.679}{Y}%
\NativeDiagram{%
\path[draw=black,line width=0.478bp] (195.655,-25.883) -- (271.863,-25.883);%
\path[draw=black,line width=0.478bp,line cap=round,line join=round] (269.615,-23.014) .. controls (270.069,-24.736) and (271.082,-25.549) .. (272.102,-25.883) .. controls (271.082,-26.218) and (270.069,-27.031) .. (269.615,-28.753);%
}%
\TextAt{21}{303.435}{750.932}{\Character{27}}%
\NativeDiagram{%
\path[draw=black,line width=0.478bp] (194.62,-33.451) -- (223.573,-56.161);%
\path[draw=black,line width=0.478bp,line cap=round,line join=round] (223.575,-52.517) .. controls (222.87,-54.152) and (223.165,-55.416) .. (223.761,-56.309) .. controls (222.752,-55.943) and (221.454,-55.958) .. (220.034,-57.032);%
}%
\TextAt{21}{273.523}{720.696}{f}%
\NativeDiagram{%
\path[draw=black,line width=0.478bp] (273.376,-33.451) -- (244.423,-56.161);%
\path[draw=black,line width=0.478bp,line cap=round,line join=round] (247.962,-57.032) .. controls (246.542,-55.958) and (245.244,-55.943) .. (244.235,-56.309) .. controls (244.831,-55.416) and (245.126,-54.152) .. (244.421,-52.517);%
}%
\TextAt{21}{333.527}{720.696}{f}%
\TextAt{8}{79.226}{686.654}{is}%
\TextAt{8}{92.257}{686.654}{commutative.}%
\TextAt{8}{170.534}{686.654}{Such}%
\TextAt{8}{200.086}{686.654}{an}%
\TextAt{8}{217.612}{686.654}{automorphism}%
\TextAt{8}{296.348}{686.654}{is}%
\TextAt{8}{309.379}{686.654}{called}%
\TextAt{8}{343.809}{686.654}{a}%
\TextAt{12}{354.831}{686.654}{covering}%
\TextAt{12}{409.417}{686.654}{transformation}%
\TextAt{8}{501.914}{686.654}{of}%
\TextAt{8}{516.51}{686.654}{the}%
\TextAt{8}{79.226}{669.608}{ramified}%
\TextAt{8}{125.109}{669.608}{covering}%
\TextAt{13}{171.326}{669.608}{X}%
\TextAt{9}{181.981}{669.608}{\Character{47}}%
\TextAt{13}{187.834}{669.608}{Y}%
\TextAt{8}{197.227}{669.608}{.}%
\TextAt{8}{206.747}{669.608}{The}%
\TextAt{8}{231.158}{669.608}{group}%
\TextAt{8}{264.691}{669.608}{of}%
\TextAt{8}{278.378}{669.608}{covering}%
\TextAt{8}{324.595}{669.608}{transformations}%
\TextAt{8}{409.32}{669.608}{is}%
\TextAt{8}{421.441}{669.608}{denoted}%
\TextAt{8}{466.02}{669.608}{by}%
\TextAt{8}{482.636}{669.608}{Aut}%
\TextAt{21}{502.47}{667.815}{Y}%
\TextAt{9}{509.758}{669.608}{\Character{40}}%
\TextAt{13}{514.31}{669.608}{X}%
\TextAt{9}{524.965}{669.608}{\Character{41}}%
\TextAt{8}{529.518}{669.608}{.}%
\TextAt{8}{79.226}{652.563}{We}%
\TextAt{8}{99.371}{652.563}{thus}%
\TextAt{8}{125.122}{652.563}{have}%
\TextAt{8}{152.117}{652.563}{a}%
\TextAt{8}{161.872}{652.563}{canonical}%
\TextAt{8}{213.256}{652.563}{isomorphism}%
\TextAt{13}{249.42}{624.13}{G}%
\TextAt{9}{258.653}{624.13}{\Character{40}}%
\TextAt{13}{263.206}{624.13}{L}%
\TextAt{16}{274.491}{624.13}{\Character{106}}%
\TextAt{13}{281.133}{624.13}{K}%
\TextAt{9}{291.95}{624.13}{\Character{41}}%
\NativeMathAt{299.824}{624.129}{11.955}{\cong}%
\TextAt{8}{312.443}{624.13}{Aut}%
\TextAt{21}{332.276}{622.336}{Y}%
\TextAt{9}{339.565}{624.13}{\Character{40}}%
\TextAt{13}{344.117}{624.13}{X}%
\TextAt{9}{354.772}{624.13}{\Character{41}}%
\TextAt{13}{359.324}{624.13}{\Character{58}}%
\TextAt{8}{79.226}{595.697}{In}%
\TextAt{8}{95.593}{595.697}{chap.}%
\TextAt{8}{128.231}{595.697}{II,}%
\TextAt{8}{145.89}{595.697}{§7,}%
\TextAt{8}{166.758}{595.697}{we}%
\TextAt{8}{185.731}{595.697}{will}%
\TextAt{8}{209.582}{595.697}{see}%
\TextAt{8}{230.241}{595.697}{that}%
\TextAt{8}{257.356}{595.697}{the}%
\TextAt{8}{279.259}{595.697}{composite}%
\TextAt{8}{336.022}{595.697}{of}%
\TextAt{8}{351.098}{595.697}{two}%
\TextAt{8}{374.961}{595.697}{unramified}%
\TextAt{8}{435.239}{595.697}{extensions}%
\TextAt{8}{493.366}{595.697}{of}%
\TextAt{13}{508.442}{595.697}{K}%
\TextAt{8}{524.903}{595.697}{is}%
\TextAt{8}{79.226}{578.651}{again}%
\TextAt{8}{110.911}{578.651}{unramified.}%
\TextAt{8}{175.368}{578.651}{The}%
\TextAt{8}{199.879}{578.651}{composite}%
\TextAt{20}{258.093}{581.673}{\Character{101}}%
\TextAt{13}{255.355}{578.651}{K}%
\TextAt{8}{266.172}{578.651}{,}%
\TextAt{8}{273.894}{578.651}{taken}%
\TextAt{8}{306.224}{578.651}{inside}%
\TextAt{8}{339.905}{578.651}{some}%
\TextAt{8}{369.689}{578.651}{algebraic}%
\TextAt{8}{419.581}{578.651}{closure}%
\Vector{q 0 G 0 g 0.478 w 387.12 -185.506 m 397.938 -185.506 l S Q}%
\TextAt{13}{459.12}{578.651}{K}%
\TextAt{8}{474.293}{578.651}{of}%
\TextAt{13}{488.081}{578.651}{K}%
\TextAt{8}{498.899}{578.651}{,}%
\TextAt{8}{506.62}{578.651}{of}%
\TextAt{8}{520.408}{578.651}{all}%
\TextAt{8}{79.226}{561.605}{unramified}%
\TextAt{8}{138.169}{561.605}{extensions}%
\TextAt{13}{194.959}{561.605}{L}%
\TextAt{16}{206.943}{561.605}{\Character{106}}%
\TextAt{13}{214.284}{561.605}{K}%
\TextAt{8}{229.409}{561.605}{is}%
\TextAt{8}{241.583}{561.605}{called}%
\TextAt{8}{275.157}{561.605}{the}%
\TextAt{12}{295.723}{561.605}{maximal}%
\TextAt{12}{350.77}{561.605}{unramified}%
\TextAt{12}{418.741}{561.605}{extension}%
\TextAt{8}{478.377}{561.605}{of}%
\TextAt{13}{492.117}{561.605}{K}%
\TextAt{8}{502.935}{561.605}{.}%
\TextAt{8}{512.613}{561.605}{The}%
\TextAt{8}{79.226}{544.559}{integral}%
\TextAt{8}{122.15}{544.559}{closure}%
\TextAt{20}{161.149}{545.094}{\Character{101}}%
\TextAt{14}{161.548}{544.559}{o}%
\TextAt{8}{171.605}{544.559}{of}%
\TextAt{14}{185.251}{544.559}{o}%
\TextAt{8}{195.309}{544.559}{in}%
\TextAt{20}{212.016}{547.581}{\Character{101}}%
\TextAt{13}{209.278}{544.559}{K}%
\TextAt{8}{224.309}{544.559}{is}%
\TextAt{8}{236.388}{544.559}{still}%
\TextAt{8}{259.527}{544.559}{a}%
\TextAt{8}{269.598}{544.559}{one-dimensional}%
\TextAt{8}{355.824}{544.559}{integral}%
\TextAt{8}{398.748}{544.559}{domain,}%
\TextAt{8}{444.023}{544.559}{but}%
\TextAt{8}{465.798}{544.559}{in}%
\TextAt{8}{479.767}{544.559}{general}%
\TextAt{8}{520.408}{544.559}{no}%
\TextAt{8}{79.226}{527.513}{longer}%
\TextAt{8}{114.179}{527.513}{noetherian,}%
\TextAt{8}{175.503}{527.513}{and,}%
\TextAt{8}{201.38}{527.513}{as}%
\TextAt{8}{215.578}{527.513}{a}%
\TextAt{8}{225.161}{527.513}{rule,}%
\TextAt{8}{251.684}{527.513}{there}%
\TextAt{8}{281.424}{527.513}{will}%
\TextAt{8}{303.357}{527.513}{be}%
\TextAt{8}{319.11}{527.513}{infinitely}%
\TextAt{8}{368.038}{527.513}{many}%
\TextAt{8}{399.738}{527.513}{prime}%
\TextAt{8}{432.73}{527.513}{ideals}%
\TextAt{8}{465.136}{527.513}{lying}%
\TextAt{8}{493.908}{527.513}{above}%
\TextAt{8}{526.911}{527.513}{a}%
\TextAt{8}{79.226}{510.467}{given}%
\TextAt{8}{109.796}{510.467}{prime}%
\TextAt{8}{142.959}{510.467}{ideal}%
\TextAt{14}{170.922}{510.467}{p}%
\NativeMathAt{180.226}{510.467}{11.955}{\neq}%
\TextAt{9}{192.636}{510.467}{0}%
\TextAt{8}{202.402}{510.467}{of}%
\TextAt{14}{215.732}{510.467}{o}%
\TextAt{8}{221.576}{510.467}{.}%
\TextAt{8}{230.025}{510.467}{The}%
\TextAt{8}{254.078}{510.467}{scheme}%
\TextAt{20}{295.504}{513.489}{\Character{101}}%
\TextAt{13}{294.128}{510.467}{Y}%
\TextAt{9}{306.842}{510.467}{\Character{61}}%
\TextAt{8}{319.268}{510.467}{Spec}%
\TextAt{9}{342.676}{510.467}{\Character{40}}%
\TextAt{20}{346.83}{511.002}{\Character{101}}%
\TextAt{14}{347.228}{510.467}{o}%
\TextAt{9}{353.073}{510.467}{\Character{41}}%
\TextAt{8}{361.522}{510.467}{with}%
\TextAt{8}{388.182}{510.467}{the}%
\TextAt{8}{408.339}{510.467}{morphism}%
\TextAt{13}{275.183}{482.034}{f}%
\TextAt{9}{285.551}{482.034}{\Character{58}}%
\TextAt{20}{293.499}{485.056}{\Character{101}}%
\TextAt{13}{292.123}{482.034}{Y}%
\NativeArrowAt{304.837}{482.034}{11.955}{19.26}{\longrightarrow}%
\TextAt{13}{327.419}{482.034}{Y}%
\TextAt{8}{79.226}{453.601}{is}%
\TextAt{8}{92.654}{453.601}{called}%
\TextAt{8}{127.482}{453.601}{the}%
\TextAt{12}{149.303}{453.601}{universal}%
\TextAt{12}{208.34}{453.601}{covering}%
\TextAt{8}{262.547}{453.601}{of}%
\TextAt{13}{277.541}{453.601}{Y}%
\TextAt{8}{286.934}{453.601}{.}%
\TextAt{8}{300.376}{453.601}{It}%
\TextAt{8}{314.713}{453.601}{plays}%
\TextAt{8}{346.361}{453.601}{the}%
\TextAt{8}{368.182}{453.601}{same}%
\TextAt{8}{399.172}{453.601}{role}%
\TextAt{8}{423.599}{453.601}{for}%
\TextAt{8}{443.148}{453.601}{schemes}%
\TextAt{8}{489.477}{453.601}{that}%
\TextAt{8}{516.51}{453.601}{the}%
\TextAt{8}{79.226}{436.555}{universal}%
\TextAt{8}{130.088}{436.555}{covering}%
\TextAt{8}{177.316}{436.555}{space}%
\TextAt{20}{212.94}{439.577}{\Character{101}}%
\TextAt{13}{209.958}{436.555}{X}%
\NativeMathAt{226.264}{436.555}{11.955}{\rightarrow}%
\TextAt{13}{243.871}{436.555}{X}%
\TextAt{8}{259.791}{436.555}{of}%
\TextAt{8}{274.488}{436.555}{a}%
\TextAt{8}{285.611}{436.555}{topological}%
\TextAt{8}{346.503}{436.555}{space}%
\TextAt{8}{379.146}{436.555}{plays}%
\TextAt{8}{410.497}{436.555}{in}%
\TextAt{8}{425.517}{436.555}{topology.}%
\TextAt{8}{481.334}{436.555}{There}%
\TextAt{8}{516.51}{436.555}{the}%
\TextAt{8}{79.226}{419.509}{group}%
\TextAt{8}{111.921}{419.509}{of}%
\TextAt{8}{124.77}{419.509}{covering}%
\TextAt{8}{170.15}{419.509}{transformations}%
\TextAt{8}{254.037}{419.509}{Aut}%
\TextAt{21}{273.87}{417.715}{X}%
\TextAt{9}{281.955}{419.509}{\Character{40}}%
\TextAt{20}{289.49}{422.531}{\Character{101}}%
\TextAt{13}{286.508}{419.509}{X}%
\TextAt{9}{297.163}{419.509}{\Character{41}}%
\TextAt{8}{305.132}{419.509}{is}%
\TextAt{8}{316.415}{419.509}{canonically}%
\TextAt{8}{376.75}{419.509}{isomorphic}%
\TextAt{8}{435.519}{419.509}{to}%
\TextAt{8}{449.349}{419.509}{the}%
\TextAt{8}{469.025}{419.509}{fundamental}%
\TextAt{8}{79.226}{402.463}{group}%
\TextAt{13}{111.992}{402.463}{\Character{25}}%
\TextAt{17}{118.632}{400.669}{1}%
\TextAt{9}{123.365}{402.463}{\Character{40}}%
\TextAt{13}{127.917}{402.463}{X}%
\TextAt{9}{138.572}{402.463}{\Character{41}}%
\TextAt{8}{143.124}{402.463}{.}%
\TextAt{8}{151.436}{402.463}{Therefore}%
\TextAt{8}{204.024}{402.463}{we}%
\TextAt{8}{220.842}{402.463}{define}%
\TextAt{8}{254.241}{402.463}{in}%
\TextAt{8}{267.484}{402.463}{our}%
\TextAt{8}{287.889}{402.463}{present}%
\TextAt{8}{328.186}{402.463}{context}%
\TextAt{8}{369.405}{402.463}{the}%
\TextAt{12}{389.151}{402.463}{fundamental}%
\TextAt{12}{466.197}{402.463}{group}%
\TextAt{8}{503.59}{402.463}{of}%
\TextAt{8}{516.51}{402.463}{the}%
\TextAt{8}{79.226}{385.417}{scheme}%
\TextAt{13}{119.276}{385.417}{Y}%
\TextAt{8}{132.567}{385.417}{by}%
\TextAt{13}{225.913}{368.371}{\Character{25}}%
\TextAt{17}{232.553}{366.578}{1}%
\TextAt{9}{237.285}{368.371}{\Character{40}}%
\TextAt{13}{241.838}{368.371}{Y}%
\TextAt{9}{251.231}{368.371}{\Character{41}}%
\TextAt{9}{259.086}{368.371}{\Character{61}}%
\TextAt{8}{271.529}{368.371}{Aut}%
\TextAt{21}{291.363}{366.578}{Y}%
\TextAt{9}{298.651}{368.371}{\Character{40}}%
\TextAt{20}{304.579}{371.393}{\Character{101}}%
\TextAt{13}{303.203}{368.371}{Y}%
\TextAt{9}{312.597}{368.371}{\Character{41}}%
\TextAt{9}{320.452}{368.371}{\Character{61}}%
\TextAt{13}{332.895}{368.371}{G}%
\TextAt{9}{342.129}{368.371}{\Character{40}}%
\TextAt{20}{349.419}{371.393}{\Character{101}}%
\TextAt{13}{346.681}{368.371}{K}%
\TextAt{16}{360.82}{368.371}{\Character{106}}%
\TextAt{13}{367.461}{368.371}{K}%
\TextAt{9}{378.279}{368.371}{\Character{41}}%
\TextAt{13}{382.831}{368.371}{\Character{58}}%
\TextAt{8}{79.226}{345.093}{This}%
\TextAt{8}{105.975}{345.093}{establishes}%
\TextAt{8}{164.071}{345.093}{a}%
\TextAt{8}{173.856}{345.093}{first}%
\TextAt{8}{198.011}{345.093}{link}%
\TextAt{8}{221.126}{345.093}{of}%
\TextAt{8}{234.485}{345.093}{Galois}%
\TextAt{8}{270.428}{345.093}{theory}%
\TextAt{8}{307.207}{345.093}{with}%
\TextAt{8}{333.897}{345.093}{classical}%
\TextAt{8}{378.925}{345.093}{topology.}%
\TextAt{8}{430.728}{345.093}{This}%
\TextAt{8}{457.477}{345.093}{link}%
\TextAt{8}{480.592}{345.093}{is}%
\TextAt{8}{492.385}{345.093}{pursued}%
\TextAt{8}{79.226}{328.047}{much}%
\TextAt{8}{110.441}{328.047}{further}%
\TextAt{8}{149.786}{328.047}{in}%
\TextAt{12}{163.438}{328.047}{étale}%
\TextAt{12}{195.705}{328.047}{topology}%
\TextAt{8}{246.551}{328.047}{.}%
\TextAt{8}{97.159}{308.012}{The}%
\TextAt{8}{122.462}{308.012}{geometric}%
\TextAt{8}{177.042}{308.012}{point}%
\TextAt{8}{208.861}{308.012}{of}%
\TextAt{8}{223.44}{308.012}{view}%
\TextAt{8}{251.671}{308.012}{of}%
\TextAt{8}{266.25}{308.012}{algebraic}%
\TextAt{8}{316.934}{308.012}{number}%
\TextAt{8}{360.779}{308.012}{fields}%
\TextAt{8}{391.999}{308.012}{explained}%
\TextAt{8}{445.6}{308.012}{in}%
\TextAt{8}{460.501}{308.012}{this}%
\TextAt{8}{484.573}{308.012}{section}%
\TextAt{8}{524.903}{308.012}{is}%
\TextAt{8}{79.226}{290.966}{corroborated}%
\TextAt{8}{148.463}{290.966}{very}%
\TextAt{8}{174.112}{290.966}{convincingly}%
\TextAt{8}{241.388}{290.966}{by}%
\TextAt{8}{257.604}{290.966}{the}%
\TextAt{8}{277.717}{290.966}{theory}%
\TextAt{8}{314.424}{290.966}{of}%
\TextAt{8}{327.711}{290.966}{function}%
\TextAt{8}{373.516}{290.966}{fields}%
\TextAt{8}{403.444}{290.966}{of}%
\TextAt{8}{416.731}{290.966}{algebraic}%
\TextAt{8}{466.122}{290.966}{curves}%
\TextAt{8}{501.909}{290.966}{over}%
\TextAt{8}{526.911}{290.966}{a}%
\TextAt{8}{79.226}{273.92}{finite}%
\TextAt{8}{109.138}{273.92}{field}%
\TextAt{18}{134.495}{273.92}{F}%
\TextAt{21}{141.801}{272.127}{p}%
\TextAt{8}{146.562}{273.92}{.}%
\TextAt{8}{155.01}{273.92}{In}%
\TextAt{8}{169.631}{273.92}{fact,}%
\TextAt{8}{195.969}{273.92}{a}%
\TextAt{8}{205.724}{273.92}{very}%
\TextAt{8}{231.416}{273.92}{close}%
\TextAt{8}{259.439}{273.92}{analogy}%
\TextAt{8}{302.704}{273.92}{exists}%
\TextAt{8}{335.019}{273.92}{between}%
\TextAt{8}{380.21}{273.92}{both}%
\TextAt{8}{407.85}{273.92}{theories.}%
\TextAt{15}{79.226}{227.118}{§}%
\TextAt{15}{93.785}{227.118}{14}%
\TextAt{15}{130.385}{227.118}{Function}%
\TextAt{15}{210.902}{227.118}{Fields}%
\TextAt{8}{79.226}{195.232}{We}%
\TextAt{8}{100.603}{195.232}{conclude}%
\TextAt{8}{149.956}{195.232}{this}%
\TextAt{8}{174.011}{195.232}{chapter}%
\TextAt{8}{217.194}{195.232}{with}%
\TextAt{8}{245.087}{195.232}{a}%
\TextAt{8}{256.075}{195.232}{brief}%
\TextAt{8}{284.291}{195.232}{sketch}%
\TextAt{8}{321.03}{195.232}{of}%
\TextAt{8}{335.593}{195.232}{the}%
\TextAt{8}{356.982}{195.232}{theory}%
\TextAt{8}{394.965}{195.232}{of}%
\TextAt{12}{410.042}{195.232}{function}%
\TextAt{12}{464.159}{195.232}{fields}%
\TextAt{8}{494.286}{195.232}{.}%
\TextAt{8}{506.432}{195.232}{They}%
\TextAt{8}{79.226}{178.186}{represent}%
\TextAt{8}{129.302}{178.186}{a}%
\TextAt{8}{138.67}{178.186}{striking}%
\TextAt{8}{180.951}{178.186}{analogy}%
\TextAt{8}{223.83}{178.186}{with}%
\TextAt{8}{250.103}{178.186}{algebraic}%
\TextAt{8}{299.151}{178.186}{number}%
\TextAt{8}{341.36}{178.186}{fields,}%
\TextAt{8}{374.274}{178.186}{and}%
\TextAt{8}{396.649}{178.186}{since}%
\TextAt{8}{424.931}{178.186}{they}%
\TextAt{8}{450.881}{178.186}{are}%
\TextAt{8}{470.005}{178.186}{immediately}%
\TextAt{8}{79.226}{161.14}{related}%
\TextAt{8}{117.689}{161.14}{to}%
\TextAt{8}{131.441}{161.14}{geometry,}%
\TextAt{8}{184.326}{161.14}{they}%
\TextAt{8}{210.104}{161.14}{actually}%
\TextAt{8}{254.102}{161.14}{serve}%
\TextAt{8}{282.869}{161.14}{as}%
\TextAt{8}{296.681}{161.14}{an}%
\TextAt{8}{312.381}{161.14}{important}%
\TextAt{8}{367.115}{161.14}{model}%
\TextAt{8}{401.346}{161.14}{for}%
\TextAt{8}{418.672}{161.14}{the}%
\TextAt{8}{438.27}{161.14}{theory}%
\TextAt{8}{474.461}{161.14}{of}%
\TextAt{8}{487.232}{161.14}{algebraic}%
\TextAt{8}{79.226}{144.094}{number}%
\TextAt{8}{121.822}{144.094}{fields.}%
\TextAt{8}{97.159}{124.059}{The}%
\TextAt{8}{121.004}{124.059}{ring}%
\TextAt{18}{144.862}{124.059}{Z}%
\TextAt{8}{156.521}{124.059}{of}%
\TextAt{8}{169.643}{124.059}{integers}%
\TextAt{8}{212.748}{124.059}{with}%
\TextAt{8}{239.199}{124.059}{its}%
\TextAt{8}{255.31}{124.059}{field}%
\TextAt{8}{280.458}{124.059}{of}%
\TextAt{8}{293.58}{124.059}{fractions}%
\TextAt{18}{341.24}{124.059}{Q}%
\TextAt{8}{354.228}{124.059}{exhibits}%
\TextAt{8}{397.978}{124.059}{obvious}%
\TextAt{8}{440.115}{124.059}{analogies}%
\TextAt{8}{490.059}{124.059}{with}%
\TextAt{8}{516.51}{124.059}{the}%
\TextAt{8}{79.226}{107.013}{polynomial}%
\TextAt{8}{139.091}{107.013}{ring}%
\TextAt{18}{162.528}{107.013}{F}%
\TextAt{21}{169.834}{105.22}{p}%
\TextAt{9}{174.595}{107.013}{\Character{91}}%
\TextAt{13}{177.847}{107.013}{t}%
\TextAt{9}{182.074}{107.013}{\Character{93}}%
\TextAt{8}{188.595}{107.013}{over}%
\TextAt{8}{213.013}{107.013}{the}%
\TextAt{8}{232.541}{107.013}{field}%
\TextAt{18}{257.269}{107.013}{F}%
\TextAt{21}{264.575}{105.22}{p}%
\TextAt{8}{272.605}{107.013}{with}%
\TextAt{13}{298.637}{107.013}{p}%
\TextAt{8}{307.781}{107.013}{elements}%
\TextAt{8}{355.01}{107.013}{and}%
\TextAt{8}{377.144}{107.013}{its}%
\TextAt{8}{392.834}{107.013}{field}%
\TextAt{8}{417.563}{107.013}{of}%
\TextAt{8}{430.265}{107.013}{fractions}%
\TextAt{18}{477.505}{107.013}{F}%
\TextAt{21}{484.811}{105.22}{p}%
\TextAt{9}{489.572}{107.013}{\Character{40}}%
\TextAt{13}{494.124}{107.013}{t}%
\TextAt{9}{498.351}{107.013}{\Character{41}}%
\TextAt{8}{502.904}{107.013}{.}%
\TextAt{8}{511.143}{107.013}{Like}%
\TextAt{18}{79.226}{89.967}{Z}%
\TextAt{8}{87.196}{89.967}{,}%
\TextAt{18}{95.482}{89.967}{F}%
\TextAt{21}{102.788}{88.174}{p}%
\TextAt{9}{107.549}{89.967}{\Character{91}}%
\TextAt{13}{110.8}{89.967}{t}%
\TextAt{9}{115.028}{89.967}{\Character{93}}%
\TextAt{8}{123.086}{89.967}{is}%
\TextAt{8}{135.759}{89.967}{also}%
\TextAt{8}{160.148}{89.967}{a}%
\TextAt{8}{170.813}{89.967}{principal}%
\TextAt{8}{220.499}{89.967}{ideal}%
\TextAt{8}{249.371}{89.967}{domain.}%
\TextAt{8}{298.277}{89.967}{The}%
\TextAt{8}{323.24}{89.967}{prime}%
\TextAt{8}{357.313}{89.967}{numbers}%
\TextAt{8}{405.433}{89.967}{correspond}%
\TextAt{8}{465.915}{89.967}{to}%
\TextAt{8}{481.135}{89.967}{the}%
\TextAt{8}{502.2}{89.967}{monic}%
\TextAt{8}{79.226}{72.921}{irreducible}%
\TextAt{8}{137.501}{72.921}{polynomials}%
\TextAt{13}{203.01}{72.921}{p}%
\TextAt{9}{208.885}{72.921}{\Character{40}}%
\TextAt{13}{213.437}{72.921}{t}%
\TextAt{9}{217.664}{72.921}{\Character{41}}%
\NativeMathAt{226.219}{72.921}{11.955}{\in}%
\TextAt{18}{238.193}{72.921}{F}%
\TextAt{21}{245.499}{71.128}{p}%
\TextAt{9}{250.259}{72.921}{\Character{91}}%
\TextAt{13}{253.511}{72.921}{t}%
\TextAt{9}{257.738}{72.921}{\Character{93}}%
\TextAt{8}{260.99}{72.921}{.}%
\TextAt{8}{270.639}{72.921}{Like}%
\TextAt{8}{296.564}{72.921}{the}%
\TextAt{8}{317.12}{72.921}{prime}%
\TextAt{8}{350.684}{72.921}{numbers}%
\TextAt{8}{398.296}{72.921}{they}%
\TextAt{8}{425.033}{72.921}{have}%
\TextAt{8}{452.428}{72.921}{finite}%
\TextAt{8}{482.74}{72.921}{fields}%
\TextAt{18}{513.112}{72.921}{F}%
\TextAt{21}{520.418}{70.619}{p}%
\TextAt{22}{524.681}{72.888}{d}%
\TextAt{8}{529.518}{72.921}{,}%
\TextAt{8}{300.14}{43.033}{87}%
\Destination{78.23}{809.815}{nav.page.87}%
\Destination{79.23}{252.625}{nav.section.1.14}%
\Destination{79.36}{305.955}{nav.print.94}%
\Destination{79.23}{252.625}{outline.19}%
\Hotlink{127.731}{593.455}{9.763}{11.047}{nav.chapter.2}%
\Hotlink{151.069}{593.455}{6.858}{11.047}{nav.section.2.7}%
}

% Page 100 (body page 88)
\Page{%
\TextAt{2}{79.226}{796.943}{I.}%
\TextAt{2}{89.83}{796.943}{Algebraic}%
\TextAt{2}{138.648}{796.943}{Integers}%
\Vector{q 0 G 0 g 0.299 w 7.226 18.729 m 460.769 18.729 l S Q}%
\TextAt{13}{79.226}{761.124}{d}%
\TextAt{9}{89.015}{761.124}{\Character{61}}%
\TextAt{8}{101.826}{761.124}{deg}%
\TextAt{9}{119.388}{761.124}{\Character{40}}%
\TextAt{13}{123.94}{761.124}{p}%
\TextAt{9}{129.815}{761.124}{\Character{40}}%
\TextAt{13}{134.368}{761.124}{t}%
\TextAt{9}{138.595}{761.124}{\Character{41}\Character{41}}%
\TextAt{8}{147.699}{761.124}{,}%
\TextAt{8}{155.131}{761.124}{as}%
\TextAt{8}{169.727}{761.124}{their}%
\TextAt{8}{197.917}{761.124}{residue}%
\TextAt{8}{237.87}{761.124}{class}%
\TextAt{8}{265.533}{761.124}{rings.}%
\TextAt{8}{299.442}{761.124}{The}%
\TextAt{8}{323.722}{761.124}{difference}%
\TextAt{8}{376.288}{761.124}{is,}%
\TextAt{8}{391.586}{761.124}{however,}%
\TextAt{8}{440.001}{761.124}{that}%
\TextAt{8}{465.596}{761.124}{now}%
\TextAt{8}{490.21}{761.124}{all}%
\TextAt{8}{506.695}{761.124}{these}%
\TextAt{8}{79.226}{744.078}{fields}%
\TextAt{8}{110.288}{744.078}{have}%
\TextAt{8}{138.373}{744.078}{the}%
\TextAt{8}{159.619}{744.078}{same}%
\TextAt{8}{190.035}{744.078}{characteristic.}%
\TextAt{8}{269.79}{744.078}{The}%
\TextAt{8}{294.934}{744.078}{geometric}%
\TextAt{8}{349.356}{744.078}{character}%
\TextAt{8}{401.506}{744.078}{of}%
\TextAt{8}{415.927}{744.078}{the}%
\TextAt{8}{437.173}{744.078}{ring}%
\TextAt{18}{462.329}{744.078}{F}%
\TextAt{21}{469.635}{742.284}{p}%
\TextAt{9}{474.396}{744.078}{\Character{91}}%
\TextAt{13}{477.647}{744.078}{t}%
\TextAt{9}{481.874}{744.078}{\Character{93}}%
\TextAt{8}{490.113}{744.078}{becomes}%
\TextAt{8}{79.226}{727.032}{much}%
\TextAt{8}{110.941}{727.032}{more}%
\TextAt{8}{140.706}{727.032}{apparent}%
\TextAt{8}{190.318}{727.032}{in}%
\TextAt{8}{204.47}{727.032}{that,}%
\TextAt{8}{233.715}{727.032}{for}%
\TextAt{8}{252.099}{727.032}{an}%
\TextAt{8}{268.858}{727.032}{element}%
\TextAt{13}{312.599}{727.032}{f}%
\TextAt{9}{323.818}{727.032}{\Character{61}}%
\TextAt{13}{337.095}{727.032}{f}%
\TextAt{9}{344.141}{727.032}{\Character{40}}%
\TextAt{13}{348.693}{727.032}{t}%
\TextAt{9}{352.92}{727.032}{\Character{41}}%
\NativeMathAt{361.645}{727.032}{11.955}{\in}%
\TextAt{18}{373.787}{727.032}{F}%
\TextAt{21}{381.093}{725.238}{p}%
\TextAt{9}{385.854}{727.032}{\Character{91}}%
\TextAt{13}{389.106}{727.032}{t}%
\TextAt{9}{393.333}{727.032}{\Character{93}}%
\TextAt{8}{396.584}{727.032}{,}%
\TextAt{8}{404.358}{727.032}{the}%
\TextAt{8}{425.014}{727.032}{value}%
\TextAt{8}{455.76}{727.032}{of}%
\TextAt{13}{469.589}{727.032}{f}%
\TextAt{8}{481.033}{727.032}{at}%
\TextAt{8}{495.843}{727.032}{a}%
\TextAt{8}{506.098}{727.032}{point}%
\TextAt{14}{79.226}{709.986}{p}%
\TextAt{9}{88.796}{709.986}{\Character{61}}%
\TextAt{9}{101.469}{709.986}{\Character{40}}%
\TextAt{13}{106.039}{709.986}{p}%
\TextAt{9}{111.915}{709.986}{\Character{40}}%
\TextAt{13}{116.467}{709.986}{t}%
\TextAt{9}{120.694}{709.986}{\Character{41}\Character{41}}%
\TextAt{8}{133.852}{709.986}{of}%
\TextAt{8}{147.338}{709.986}{the}%
\TextAt{8}{167.65}{709.986}{affine}%
\TextAt{8}{199.021}{709.986}{scheme}%
\TextAt{13}{239.227}{709.986}{X}%
\TextAt{9}{253.469}{709.986}{\Character{61}}%
\TextAt{8}{266.16}{709.986}{Spec}%
\TextAt{9}{289.568}{709.986}{\Character{40}}%
\TextAt{18}{294.121}{709.986}{F}%
\TextAt{21}{301.427}{708.193}{p}%
\TextAt{9}{306.187}{709.986}{\Character{91}}%
\TextAt{13}{309.439}{709.986}{t}%
\TextAt{9}{313.666}{709.986}{\Character{93}\Character{41}}%
\TextAt{8}{325.524}{709.986}{is}%
\TextAt{8}{337.443}{709.986}{actually}%
\TextAt{8}{382.156}{709.986}{given}%
\TextAt{8}{412.882}{709.986}{by}%
\TextAt{8}{429.297}{709.986}{the}%
\TextAt{8}{449.609}{709.986}{value}%
\TextAt{13}{480.011}{709.986}{f}%
\TextAt{9}{487.058}{709.986}{\Character{40}}%
\TextAt{13}{491.61}{709.986}{a}%
\TextAt{9}{497.755}{709.986}{\Character{41}}%
\NativeMathAt{505.894}{709.986}{11.955}{\in}%
\TextAt{18}{517.451}{709.986}{F}%
\TextAt{21}{524.757}{708.193}{p}%
\TextAt{8}{529.518}{709.986}{,}%
\TextAt{8}{79.226}{692.94}{if}%
\TextAt{13}{90.19}{692.94}{p}%
\TextAt{9}{96.065}{692.94}{\Character{40}}%
\TextAt{13}{100.618}{692.94}{t}%
\TextAt{9}{104.845}{692.94}{\Character{41}}%
\TextAt{9}{113.118}{692.94}{\Character{61}}%
\TextAt{13}{125.963}{692.94}{t}%
\TextAt{16}{133.011}{692.94}{\Character{0}}%
\TextAt{13}{145.13}{692.94}{a}%
\TextAt{8}{151.275}{692.94}{,}%
\TextAt{8}{158.724}{692.94}{or}%
\TextAt{8}{173.275}{692.94}{more}%
\TextAt{8}{202.782}{692.94}{generally}%
\TextAt{8}{252.78}{692.94}{by}%
\TextAt{13}{269.279}{692.94}{f}%
\TextAt{9}{276.325}{692.94}{\Character{40}}%
\TextAt{13}{280.878}{692.94}{\Character{11}}%
\TextAt{9}{288.4}{692.94}{\Character{41}}%
\NativeMathAt{296.682}{692.94}{11.955}{\in}%
\TextAt{18}{308.383}{692.94}{F}%
\TextAt{21}{315.689}{690.637}{p}%
\TextAt{22}{319.952}{692.907}{d}%
\TextAt{8}{324.788}{692.94}{,}%
\TextAt{8}{332.238}{692.94}{if}%
\TextAt{13}{343.202}{692.94}{\Character{11}}%
\NativeMathAt{354.455}{692.94}{11.955}{\in}%
\TextAt{18}{366.155}{692.94}{F}%
\TextAt{21}{373.461}{690.637}{p}%
\TextAt{22}{377.724}{692.907}{d}%
\TextAt{8}{386.699}{692.94}{is}%
\TextAt{8}{398.703}{692.94}{a}%
\TextAt{8}{408.699}{692.94}{zero}%
\TextAt{8}{433.65}{692.94}{of}%
\TextAt{13}{447.221}{692.94}{p}%
\TextAt{9}{453.096}{692.94}{\Character{40}}%
\TextAt{13}{457.648}{692.94}{t}%
\TextAt{9}{461.875}{692.94}{\Character{41}}%
\TextAt{8}{466.428}{692.94}{.}%
\TextAt{8}{475.597}{692.94}{This}%
\TextAt{8}{502.558}{692.94}{is}%
\TextAt{8}{514.562}{692.94}{due}%
\TextAt{8}{79.226}{675.894}{to}%
\TextAt{8}{93.536}{675.894}{the}%
\TextAt{8}{113.693}{675.894}{isomorphism}%
\TextAt{18}{265.901}{658.848}{F}%
\TextAt{21}{273.207}{657.055}{p}%
\TextAt{9}{277.968}{658.848}{\Character{91}}%
\TextAt{13}{281.22}{658.848}{t}%
\TextAt{9}{285.447}{658.848}{\Character{93}\Character{47}}%
\TextAt{14}{294.551}{658.848}{p}%
\NativeMathAt{310.193}{665.625}{7.97}{\sim}%
\NativeArrowAt{303.855}{658.848}{11.955}{19.26}{\longrightarrow}%
\TextAt{18}{326.437}{658.848}{F}%
\TextAt{21}{333.743}{656.545}{p}%
\TextAt{22}{338.006}{658.815}{d}%
\TextAt{13}{342.843}{658.848}{\Character{59}}%
\TextAt{8}{79.226}{635.326}{which}%
\TextAt{8}{113.45}{635.326}{takes}%
\TextAt{8}{144.17}{635.326}{the}%
\TextAt{8}{165.064}{635.326}{residue}%
\TextAt{8}{205.528}{635.326}{class}%
\TextAt{13}{233.702}{635.326}{f}%
\TextAt{9}{240.748}{635.326}{\Character{40}}%
\TextAt{14}{245.301}{635.326}{p}%
\TextAt{9}{251.284}{635.326}{\Character{41}}%
\TextAt{9}{260.394}{635.326}{\Character{61}}%
\TextAt{13}{274.094}{635.326}{f}%
\TextAt{8}{284.461}{635.326}{mod}%
\TextAt{14}{309.899}{635.326}{p}%
\TextAt{8}{320.517}{635.326}{to}%
\TextAt{13}{335.564}{635.326}{f}%
\TextAt{9}{342.611}{635.326}{\Character{40}}%
\TextAt{13}{347.163}{635.326}{\Character{11}}%
\TextAt{9}{354.685}{635.326}{\Character{41}}%
\TextAt{8}{359.237}{635.326}{.}%
\TextAt{8}{369.896}{635.326}{In}%
\TextAt{8}{385.255}{635.326}{the}%
\TextAt{8}{406.148}{635.326}{analogy}%
\TextAt{8}{450.151}{635.326}{between,}%
\TextAt{8}{499.514}{635.326}{on}%
\TextAt{8}{516.51}{635.326}{the}%
\TextAt{8}{79.226}{618.28}{one}%
\TextAt{8}{102.358}{618.28}{hand,}%
\TextAt{8}{136.967}{618.28}{the}%
\TextAt{8}{158.796}{618.28}{progression}%
\TextAt{8}{221.739}{618.28}{of}%
\TextAt{8}{236.742}{618.28}{the}%
\TextAt{8}{258.571}{618.28}{prime}%
\TextAt{8}{293.407}{618.28}{numbers}%
\TextAt{9}{342.291}{618.28}{2}%
\TextAt{13}{348.144}{618.28}{\Character{59}}%
\TextAt{9}{353.388}{618.28}{3}%
\TextAt{13}{359.241}{618.28}{\Character{59}}%
\TextAt{9}{364.485}{618.28}{5}%
\TextAt{13}{370.338}{618.28}{\Character{59}}%
\TextAt{9}{375.582}{618.28}{7}%
\TextAt{13}{381.435}{618.28}{\Character{59}}%
\TextAt{13}{386.659}{618.28}{\Character{58}}%
\TextAt{13}{391.896}{618.28}{\Character{58}}%
\TextAt{13}{397.132}{618.28}{\Character{58}}%
\TextAt{8}{400.419}{618.28}{,}%
\TextAt{8}{409.659}{618.28}{and}%
\TextAt{8}{434.095}{618.28}{the}%
\TextAt{8}{455.924}{618.28}{growing}%
\TextAt{8}{501.508}{618.28}{of}%
\TextAt{8}{516.51}{618.28}{the}%
\TextAt{8}{79.226}{601.234}{cardinalities}%
\TextAt{13}{144.993}{601.234}{p\Character{59}}%
\TextAt{13}{156.087}{601.234}{p}%
\TextAt{17}{161.987}{605.573}{2}%
\TextAt{13}{166.719}{601.234}{\Character{59}}%
\TextAt{13}{171.943}{601.234}{p}%
\TextAt{17}{177.839}{605.573}{3}%
\TextAt{13}{182.571}{601.234}{\Character{59}}%
\TextAt{13}{187.795}{601.234}{p}%
\TextAt{17}{193.69}{605.573}{4}%
\TextAt{13}{198.423}{601.234}{\Character{59}}%
\TextAt{13}{203.647}{601.234}{\Character{58}}%
\TextAt{13}{208.884}{601.234}{\Character{58}}%
\TextAt{13}{214.108}{601.234}{\Character{58}}%
\TextAt{8}{221.317}{601.234}{of}%
\TextAt{8}{234.66}{601.234}{the}%
\TextAt{8}{254.83}{601.234}{residue}%
\TextAt{8}{294.57}{601.234}{fields}%
\TextAt{18}{324.555}{601.234}{F}%
\TextAt{21}{331.861}{598.932}{p}%
\TextAt{22}{336.123}{601.201}{d}%
\TextAt{8}{344.871}{601.234}{on}%
\TextAt{8}{361.143}{601.234}{the}%
\TextAt{8}{381.312}{601.234}{other,}%
\TextAt{8}{415.15}{601.234}{resides}%
\TextAt{8}{453.001}{601.234}{one}%
\TextAt{8}{474.474}{601.234}{of}%
\TextAt{8}{487.817}{601.234}{the}%
\TextAt{8}{507.986}{601.234}{most}%
\TextAt{8}{79.226}{584.188}{profound}%
\TextAt{8}{128.984}{584.188}{mysteries}%
\TextAt{8}{180.486}{584.188}{of}%
\TextAt{8}{193.816}{584.188}{arithmetic.}%
\TextAt{8}{97.159}{564.153}{One}%
\TextAt{8}{123.085}{564.153}{obtains}%
\TextAt{8}{165.342}{564.153}{the}%
\TextAt{8}{186.713}{564.153}{same}%
\TextAt{8}{217.254}{564.153}{arithmetic}%
\TextAt{8}{275.053}{564.153}{theory}%
\TextAt{8}{313.018}{564.153}{for}%
\TextAt{8}{332.118}{564.153}{the}%
\TextAt{8}{353.489}{564.153}{finite}%
\TextAt{8}{384.616}{564.153}{extensions}%
\TextAt{13}{442.211}{564.153}{K}%
\TextAt{8}{458.141}{564.153}{of}%
\TextAt{18}{472.686}{564.153}{F}%
\TextAt{21}{479.992}{562.36}{p}%
\TextAt{9}{484.753}{564.153}{\Character{40}}%
\TextAt{13}{489.305}{564.153}{t}%
\TextAt{9}{493.532}{564.153}{\Character{41}}%
\TextAt{8}{503.197}{564.153}{as}%
\TextAt{8}{518.782}{564.153}{for}%
\TextAt{8}{79.226}{547.108}{algebraic}%
\TextAt{8}{129.727}{547.108}{number}%
\TextAt{8}{173.39}{547.108}{fields.}%
\TextAt{8}{211.113}{547.108}{This}%
\TextAt{8}{238.899}{547.108}{is}%
\TextAt{8}{251.73}{547.108}{clear}%
\TextAt{8}{280.76}{547.108}{from}%
\TextAt{8}{309.467}{547.108}{what}%
\TextAt{8}{339.8}{547.108}{we}%
\TextAt{8}{358.094}{547.108}{have}%
\TextAt{8}{386.155}{547.108}{developed}%
\TextAt{8}{441.522}{547.108}{for}%
\TextAt{8}{460.474}{547.108}{arbitrary}%
\TextAt{8}{511.31}{547.108}{one-}%
\TextAt{8}{79.226}{530.062}{dimensional}%
\TextAt{8}{145.219}{530.062}{noetherian}%
\TextAt{8}{204.971}{530.062}{integral}%
\TextAt{8}{249.122}{530.062}{domains.}%
\TextAt{8}{304.543}{530.062}{But}%
\TextAt{8}{329.327}{530.062}{the}%
\TextAt{8}{351.026}{530.062}{crucial}%
\TextAt{8}{390.287}{530.062}{difference}%
\TextAt{8}{444.169}{530.062}{with}%
\TextAt{8}{472.372}{530.062}{the}%
\TextAt{8}{494.071}{530.062}{number}%
\TextAt{8}{79.226}{513.016}{field}%
\TextAt{8}{104.447}{513.016}{case}%
\TextAt{8}{129.082}{513.016}{is}%
\TextAt{8}{140.709}{513.016}{seen}%
\TextAt{8}{165.989}{513.016}{in}%
\TextAt{8}{179.506}{513.016}{that}%
\TextAt{8}{204.738}{513.016}{the}%
\TextAt{8}{224.759}{513.016}{function}%
\TextAt{8}{270.47}{513.016}{field}%
\TextAt{13}{295.691}{513.016}{K}%
\TextAt{8}{310.27}{513.016}{hides}%
\TextAt{8}{340.105}{513.016}{away}%
\TextAt{8}{369.247}{513.016}{a}%
\TextAt{8}{378.866}{513.016}{finite}%
\TextAt{8}{408.641}{513.016}{number}%
\TextAt{8}{451.101}{513.016}{of}%
\TextAt{8}{464.295}{513.016}{further}%
\TextAt{8}{503.503}{513.016}{prime}%
\TextAt{8}{79.226}{495.97}{ideals,}%
\TextAt{8}{115.238}{495.97}{besides}%
\TextAt{8}{155.494}{495.97}{the}%
\TextAt{8}{175.796}{495.97}{prime}%
\TextAt{8}{209.106}{495.97}{ideals}%
\TextAt{8}{241.829}{495.97}{of}%
\TextAt{14}{255.305}{495.97}{o}%
\TextAt{8}{261.149}{495.97}{,}%
\TextAt{8}{268.481}{495.97}{which}%
\TextAt{8}{302.113}{495.97}{must}%
\TextAt{8}{331.262}{495.97}{be}%
\TextAt{8}{347.333}{495.97}{taken}%
\TextAt{8}{379.351}{495.97}{into}%
\TextAt{8}{403.24}{495.97}{account}%
\TextAt{8}{446.64}{495.97}{in}%
\TextAt{8}{460.438}{495.97}{a}%
\TextAt{8}{470.34}{495.97}{fully-fledged}%
\TextAt{8}{79.226}{478.924}{development}%
\TextAt{8}{147.191}{478.924}{of}%
\TextAt{8}{160.521}{478.924}{the}%
\TextAt{8}{180.678}{478.924}{theory.}%
\TextAt{8}{97.159}{458.889}{This}%
\TextAt{8}{123.497}{458.889}{phenomenon}%
\TextAt{8}{191.404}{458.889}{appears}%
\TextAt{8}{234.336}{458.889}{already}%
\TextAt{8}{275.259}{458.889}{for}%
\TextAt{8}{292.763}{458.889}{the}%
\TextAt{8}{312.538}{458.889}{rational}%
\TextAt{8}{355.745}{458.889}{function}%
\TextAt{8}{401.211}{458.889}{field}%
\TextAt{18}{426.186}{458.889}{F}%
\TextAt{21}{433.492}{457.096}{p}%
\TextAt{9}{438.253}{458.889}{\Character{40}}%
\TextAt{13}{442.806}{458.889}{t}%
\TextAt{9}{447.033}{458.889}{\Character{41}}%
\TextAt{8}{451.585}{458.889}{,}%
\TextAt{8}{458.429}{458.889}{where}%
\TextAt{8}{491.857}{458.889}{it}%
\TextAt{8}{503.179}{458.889}{is}%
\TextAt{8}{514.562}{458.889}{due}%
\TextAt{8}{79.226}{441.843}{to}%
\TextAt{8}{93.201}{441.843}{the}%
\TextAt{8}{113.023}{441.843}{fact}%
\TextAt{8}{135.773}{441.843}{that}%
\TextAt{8}{160.807}{441.843}{the}%
\TextAt{8}{180.628}{441.843}{choice}%
\TextAt{8}{215.083}{441.843}{of}%
\TextAt{8}{228.077}{441.843}{the}%
\TextAt{8}{247.899}{441.843}{unknown}%
\TextAt{13}{297.644}{441.843}{t}%
\TextAt{8}{305.433}{441.843}{which}%
\TextAt{8}{338.585}{441.843}{determines}%
\TextAt{8}{397.487}{441.843}{the}%
\TextAt{8}{417.309}{441.843}{ring}%
\TextAt{8}{441.039}{441.843}{of}%
\TextAt{8}{454.034}{441.843}{integrality}%
\TextAt{18}{509.972}{441.843}{F}%
\TextAt{21}{517.278}{440.05}{p}%
\TextAt{9}{522.039}{441.843}{\Character{91}}%
\TextAt{13}{525.291}{441.843}{t}%
\TextAt{9}{529.518}{441.843}{\Character{93}}%
\TextAt{8}{79.226}{424.797}{is}%
\TextAt{8}{92.051}{424.797}{totally}%
\TextAt{8}{130.521}{424.797}{arbitrary.}%
\TextAt{8}{187.045}{424.797}{A}%
\TextAt{8}{200.778}{424.797}{different}%
\TextAt{8}{248.01}{424.797}{choice,}%
\TextAt{8}{287.378}{424.797}{say}%
\TextAt{13}{308.668}{424.797}{t}%
\TextAt{3}{312.895}{429.135}{\Character{48}}%
\TextAt{9}{320.82}{424.797}{\Character{61}}%
\TextAt{9}{335.047}{424.797}{1\Character{47}}%
\TextAt{13}{346.76}{424.797}{t}%
\TextAt{8}{350.987}{424.797}{,}%
\TextAt{8}{359.463}{424.797}{determines}%
\TextAt{8}{419.762}{424.797}{a}%
\TextAt{8}{430.579}{424.797}{completely}%
\TextAt{8}{490.496}{424.797}{different}%
\TextAt{8}{79.226}{407.751}{ring}%
\TextAt{18}{103.455}{407.751}{F}%
\TextAt{21}{110.761}{405.958}{p}%
\TextAt{9}{115.522}{407.751}{\Character{91}1\Character{47}}%
\TextAt{13}{130.48}{407.751}{t}%
\TextAt{9}{134.707}{407.751}{\Character{93}}%
\TextAt{8}{137.959}{407.751}{,}%
\TextAt{8}{145.312}{407.751}{and}%
\TextAt{8}{168.239}{407.751}{thus}%
\TextAt{8}{194.154}{407.751}{completely}%
\TextAt{8}{253.172}{407.751}{different}%
\TextAt{8}{299.507}{407.751}{prime}%
\TextAt{8}{332.834}{407.751}{ideals.}%
\TextAt{8}{370.453}{407.751}{It}%
\TextAt{8}{383.289}{407.751}{is}%
\TextAt{8}{395.217}{407.751}{therefore}%
\TextAt{8}{444.48}{407.751}{crucial}%
\TextAt{8}{482.362}{407.751}{to}%
\TextAt{8}{496.836}{407.751}{build}%
\TextAt{8}{526.911}{407.751}{a}%
\TextAt{8}{79.226}{390.705}{theory}%
\TextAt{8}{116.38}{390.705}{which}%
\TextAt{8}{150.271}{390.705}{is}%
\TextAt{8}{162.438}{390.705}{independent}%
\TextAt{8}{229.17}{390.705}{of}%
\TextAt{8}{242.903}{390.705}{such}%
\TextAt{8}{269.704}{390.705}{choices.}%
\TextAt{8}{314.871}{390.705}{This}%
\TextAt{8}{341.995}{390.705}{may}%
\TextAt{8}{367.768}{390.705}{be}%
\TextAt{8}{384.096}{390.705}{done}%
\TextAt{8}{412.463}{390.705}{either}%
\TextAt{8}{446.03}{390.705}{via}%
\TextAt{8}{465.622}{390.705}{the}%
\TextAt{8}{486.183}{390.705}{theory}%
\TextAt{8}{523.337}{390.705}{of}%
\TextAt{8}{79.226}{373.659}{valuations,}%
\TextAt{8}{138.165}{373.659}{or}%
\TextAt{8}{152.475}{373.659}{scheme}%
\TextAt{8}{192.525}{373.659}{theoretically,}%
\TextAt{8}{262.116}{373.659}{i.e.,}%
\TextAt{8}{284.221}{373.659}{in}%
\TextAt{8}{297.874}{373.659}{a}%
\TextAt{8}{307.63}{373.659}{geometric}%
\TextAt{8}{360.962}{373.659}{way.}%
\TextAt{8}{97.159}{353.624}{Let}%
\TextAt{8}{118.513}{353.624}{us}%
\TextAt{8}{133.914}{353.624}{first}%
\TextAt{8}{158.424}{353.624}{sketch}%
\TextAt{8}{194.316}{353.624}{the}%
\TextAt{8}{214.857}{353.624}{more}%
\TextAt{8}{244.508}{353.624}{naive}%
\TextAt{8}{275.462}{353.624}{method,}%
\TextAt{8}{321.791}{353.624}{via}%
\TextAt{8}{341.364}{353.624}{the}%
\TextAt{8}{361.906}{353.624}{theory}%
\TextAt{8}{399.041}{353.624}{of}%
\TextAt{8}{412.756}{353.624}{valuations.}%
\TextAt{8}{474.148}{353.624}{Let}%
\TextAt{13}{495.502}{353.624}{K}%
\TextAt{8}{510.602}{353.624}{be}%
\TextAt{8}{526.911}{353.624}{a}%
\TextAt{8}{79.226}{336.578}{finite}%
\TextAt{8}{109.19}{336.578}{extension}%
\TextAt{8}{161.007}{336.578}{of}%
\TextAt{18}{174.389}{336.578}{F}%
\TextAt{21}{181.695}{334.785}{p}%
\TextAt{9}{186.456}{336.578}{\Character{40}}%
\TextAt{13}{191.008}{336.578}{t}%
\TextAt{9}{195.235}{336.578}{\Character{41}}%
\TextAt{8}{203.737}{336.578}{and}%
\TextAt{14}{226.551}{336.578}{o}%
\TextAt{8}{236.344}{336.578}{the}%
\TextAt{8}{256.553}{336.578}{integral}%
\TextAt{8}{299.212}{336.578}{closure}%
\TextAt{8}{338.346}{336.578}{of}%
\TextAt{18}{351.727}{336.578}{F}%
\TextAt{21}{359.033}{334.785}{p}%
\TextAt{9}{363.794}{336.578}{\Character{91}}%
\TextAt{13}{367.046}{336.578}{t}%
\TextAt{9}{371.273}{336.578}{\Character{93}}%
\TextAt{8}{378.474}{336.578}{in}%
\TextAt{13}{392.178}{336.578}{K}%
\TextAt{8}{402.996}{336.578}{.}%
\TextAt{8}{411.599}{336.578}{By}%
\TextAt{8}{430.014}{336.578}{§11,}%
\TextAt{8}{454.623}{336.578}{for}%
\TextAt{8}{472.559}{336.578}{every}%
\TextAt{8}{503.503}{336.578}{prime}%
\TextAt{8}{79.226}{319.532}{ideal}%
\TextAt{14}{106.959}{319.532}{p}%
\NativeMathAt{116.262}{319.532}{11.955}{\neq}%
\TextAt{9}{128.672}{319.532}{0}%
\TextAt{8}{138.208}{319.532}{of}%
\TextAt{14}{151.307}{319.532}{o}%
\TextAt{8}{160.818}{319.532}{there}%
\TextAt{8}{190.499}{319.532}{is}%
\TextAt{8}{202.032}{319.532}{an}%
\TextAt{8}{218.06}{319.532}{associated}%
\TextAt{12}{273.564}{319.532}{normalized}%
\TextAt{12}{342.304}{319.532}{discrete}%
\TextAt{12}{392.023}{319.532}{valuation}%
\TextAt{8}{445.846}{319.532}{,}%
\TextAt{8}{452.81}{319.532}{i.e.,}%
\TextAt{8}{474.731}{319.532}{a}%
\TextAt{8}{484.255}{319.532}{surjective}%
\TextAt{8}{79.226}{302.486}{function}%
\TextAt{13}{255.036}{285.44}{v}%
\TextAt{27}{260.695}{283.647}{p}%
\TextAt{9}{268.503}{285.44}{\Character{58}}%
\TextAt{13}{275.075}{285.44}{K}%
\NativeArrowAt{289.214}{285.44}{11.955}{19.26}{\longrightarrow}%
\TextAt{18}{311.796}{285.44}{Z}%
\TextAt{16}{322.423}{285.44}{\Character{91}}%
\TextAt{16}{333.039}{285.44}{\Character{102}}%
\NativeMathAt{339.017}{285.44}{11.955}{\infty}%
\TextAt{16}{350.972}{285.44}{\Character{103}}%
\TextAt{8}{79.226}{261.919}{satisfying}%
\TextAt{8}{131.387}{261.919}{the}%
\TextAt{8}{151.543}{261.919}{properties}%
\TextAt{8}{89.579}{241.884}{(i)}%
\TextAt{13}{107.918}{241.884}{v}%
\TextAt{27}{113.578}{240.091}{p}%
\TextAt{9}{118.065}{241.884}{\Character{40}0\Character{41}}%
\TextAt{9}{136.309}{241.884}{\Character{61}}%
\NativeMathAt{148.769}{241.884}{11.955}{\infty}%
\TextAt{8}{160.724}{241.884}{,}%
\TextAt{8}{86.327}{224.838}{(ii)}%
\TextAt{13}{107.918}{224.838}{v}%
\TextAt{27}{113.578}{223.045}{p}%
\TextAt{9}{118.065}{224.838}{\Character{40}}%
\TextAt{13}{122.617}{224.838}{ab}%
\TextAt{9}{133.739}{224.838}{\Character{41}}%
\TextAt{9}{141.594}{224.838}{\Character{61}}%
\TextAt{13}{154.038}{224.838}{v}%
\TextAt{27}{159.697}{223.045}{p}%
\TextAt{9}{164.184}{224.838}{\Character{40}}%
\TextAt{13}{168.736}{224.838}{a}%
\TextAt{9}{174.881}{224.838}{\Character{41}}%
\TextAt{9}{182.078}{224.838}{\Character{43}}%
\TextAt{13}{193.851}{224.838}{v}%
\TextAt{27}{199.511}{223.045}{p}%
\TextAt{9}{203.997}{224.838}{\Character{40}}%
\TextAt{13}{208.55}{224.838}{b}%
\TextAt{9}{213.527}{224.838}{\Character{41}}%
\TextAt{8}{218.079}{224.838}{,}%
\TextAt{8}{83.076}{207.792}{(iii)}%
\TextAt{13}{107.918}{207.792}{v}%
\TextAt{27}{113.578}{205.999}{p}%
\TextAt{9}{118.065}{207.792}{\Character{40}}%
\TextAt{13}{122.617}{207.792}{a}%
\TextAt{9}{131.419}{207.792}{\Character{43}}%
\TextAt{13}{143.18}{207.792}{b}%
\TextAt{9}{148.157}{207.792}{\Character{41}}%
\NativeMathAt{156.03}{207.792}{11.955}{\geq}%
\TextAt{8}{168.649}{207.792}{min}%
\TextAt{16}{188.16}{207.792}{\Character{102}}%
\TextAt{13}{194.138}{207.792}{v}%
\TextAt{27}{199.797}{205.999}{p}%
\TextAt{9}{204.284}{207.792}{\Character{40}}%
\TextAt{13}{208.836}{207.792}{a}%
\TextAt{9}{214.981}{207.792}{\Character{41}}%
\TextAt{13}{219.534}{207.792}{\Character{59}}%
\TextAt{13}{224.758}{207.792}{v}%
\TextAt{27}{230.437}{205.999}{p}%
\TextAt{9}{234.924}{207.792}{\Character{40}}%
\TextAt{13}{239.476}{207.792}{b}%
\TextAt{9}{244.453}{207.792}{\Character{41}}%
\TextAt{16}{249.006}{207.792}{\Character{103}}%
\TextAt{8}{254.983}{207.792}{.}%
\TextAt{8}{79.226}{187.757}{The}%
\TextAt{8}{102.329}{187.757}{relation}%
\TextAt{8}{144.308}{187.757}{between}%
\TextAt{8}{188.547}{187.757}{the}%
\TextAt{8}{207.752}{187.757}{valuations}%
\TextAt{8}{262.488}{187.757}{and}%
\TextAt{8}{284.299}{187.757}{the}%
\TextAt{8}{303.504}{187.757}{prime}%
\TextAt{8}{335.716}{187.757}{decomposition}%
\TextAt{8}{411.9}{187.757}{in}%
\TextAt{8}{424.601}{187.757}{the}%
\TextAt{8}{443.806}{187.757}{Dedekind}%
\TextAt{8}{495.039}{187.757}{domain}%
\TextAt{14}{79.226}{170.711}{o}%
\TextAt{8}{88.968}{170.711}{is}%
\TextAt{8}{100.732}{170.711}{given}%
\TextAt{8}{131.301}{170.711}{by}%
\TextAt{9}{267.354}{153.665}{\Character{40}}%
\TextAt{13}{271.906}{153.665}{a}%
\TextAt{9}{278.051}{153.665}{\Character{41}}%
\TextAt{9}{285.906}{153.665}{\Character{61}}%
\TextAt{20}{298.35}{165.023}{\Character{89}}%
\TextAt{27}{303.993}{139.917}{p}%
\TextAt{14}{315.618}{153.665}{p}%
\TextAt{21}{321.601}{158.601}{v}%
\TextAt{28}{325.725}{157.605}{p}%
\TextAt{17}{329.808}{158.601}{\Character{40}}%
\TextAt{21}{333.101}{158.601}{a}%
\TextAt{17}{337.599}{158.601}{\Character{41}}%
\TextAt{13}{341.39}{153.665}{\Character{58}}%
\TextAt{8}{79.226}{121.326}{The}%
\TextAt{8}{104.357}{121.326}{definition}%
\TextAt{8}{157.463}{121.326}{of}%
\TextAt{8}{171.87}{121.326}{a}%
\TextAt{8}{182.702}{121.326}{discrete}%
\TextAt{8}{226.758}{121.326}{valuation}%
\TextAt{8}{278.907}{121.326}{of}%
\TextAt{13}{293.314}{121.326}{K}%
\TextAt{8}{309.106}{121.326}{does}%
\TextAt{8}{336.58}{121.326}{not}%
\TextAt{8}{358.471}{121.326}{require}%
\TextAt{8}{398.892}{121.326}{the}%
\TextAt{8}{420.126}{121.326}{subring}%
\TextAt{14}{462.89}{121.326}{o}%
\TextAt{8}{473.709}{121.326}{to}%
\TextAt{8}{489.096}{121.326}{be}%
\TextAt{8}{506.097}{121.326}{given}%
\TextAt{8}{79.226}{104.28}{in}%
\TextAt{8}{93.895}{104.28}{advance,}%
\TextAt{8}{142.973}{104.28}{and}%
\TextAt{8}{166.751}{104.28}{in}%
\TextAt{8}{181.42}{104.28}{fact,}%
\TextAt{8}{209.027}{104.28}{aside}%
\TextAt{8}{239.369}{104.28}{from}%
\TextAt{8}{268.025}{104.28}{those}%
\TextAt{8}{299.67}{104.28}{arising}%
\TextAt{8}{338.476}{104.28}{from}%
\TextAt{14}{367.132}{104.28}{o}%
\TextAt{8}{372.976}{104.28}{,}%
\TextAt{8}{381.395}{104.28}{there}%
\TextAt{8}{412.322}{104.28}{are}%
\TextAt{8}{432.849}{104.28}{finitely}%
\TextAt{8}{473.209}{104.28}{many}%
\TextAt{8}{506.097}{104.28}{other}%
\TextAt{8}{79.226}{87.234}{discrete}%
\TextAt{8}{121.654}{87.234}{valuations}%
\TextAt{8}{176.789}{87.234}{of}%
\TextAt{13}{189.568}{87.234}{K}%
\TextAt{8}{200.386}{87.234}{.}%
\TextAt{8}{208.65}{87.234}{In}%
\TextAt{8}{222.72}{87.234}{the}%
\TextAt{8}{242.325}{87.234}{case}%
\TextAt{8}{266.545}{87.234}{of}%
\TextAt{8}{279.323}{87.234}{the}%
\TextAt{8}{298.928}{87.234}{field}%
\TextAt{18}{323.734}{87.234}{F}%
\TextAt{21}{331.04}{85.441}{p}%
\TextAt{9}{335.801}{87.234}{\Character{40}}%
\TextAt{13}{340.353}{87.234}{t}%
\TextAt{9}{344.58}{87.234}{\Character{41}}%
\TextAt{8}{352.478}{87.234}{there}%
\TextAt{8}{381.839}{87.234}{is}%
\TextAt{8}{393.051}{87.234}{one}%
\TextAt{8}{413.959}{87.234}{more}%
\TextAt{8}{442.674}{87.234}{valuation,}%
\TextAt{8}{496.557}{87.234}{besides}%
\TextAt{8}{300.14}{43.033}{88}%
\Destination{79.23}{256.475}{nav.Item.24}%
\Destination{79.23}{238.585}{nav.Item.25}%
\Destination{79.23}{221.535}{nav.Item.26}%
\Destination{78.23}{809.815}{nav.page.88}%
\Destination{79.36}{334.555}{nav.print.95}%
\Hotlink{435.193}{334.355}{12.716}{11.047}{nav.section.1.11}%
}

% Page 101 (body page 89)
\Page{%
\TextAt{2}{79.226}{796.943}{I.}%
\TextAt{2}{89.83}{796.943}{Algebraic}%
\TextAt{2}{138.648}{796.943}{Integers}%
\Vector{q 0 G 0 g 0.299 w 7.226 18.729 m 460.769 18.729 l S Q}%
\TextAt{8}{79.226}{761.124}{the}%
\TextAt{8}{98.965}{761.124}{ones}%
\TextAt{8}{124.621}{761.124}{associated}%
\TextAt{8}{179.938}{761.124}{to}%
\TextAt{8}{193.831}{761.124}{the}%
\TextAt{8}{213.569}{761.124}{prime}%
\TextAt{8}{246.315}{761.124}{ideals}%
\TextAt{14}{278.475}{761.124}{p}%
\TextAt{9}{287.779}{761.124}{\Character{61}}%
\TextAt{9}{300.189}{761.124}{\Character{40}}%
\TextAt{13}{304.757}{761.124}{p}%
\TextAt{9}{310.632}{761.124}{\Character{40}}%
\TextAt{13}{315.184}{761.124}{t}%
\TextAt{9}{319.412}{761.124}{\Character{41}\Character{41}}%
\TextAt{8}{331.996}{761.124}{of}%
\TextAt{18}{344.908}{761.124}{F}%
\TextAt{21}{352.214}{759.33}{p}%
\TextAt{9}{356.975}{761.124}{\Character{91}}%
\TextAt{13}{360.226}{761.124}{t}%
\TextAt{9}{364.454}{761.124}{\Character{93}}%
\TextAt{8}{367.705}{761.124}{,}%
\TextAt{8}{374.52}{761.124}{namely,}%
\TextAt{8}{417.105}{761.124}{the}%
\TextAt{12}{436.844}{761.124}{degree}%
\TextAt{12}{478.947}{761.124}{valuation}%
\TextAt{13}{79.226}{744.078}{v}%
\NativeMathAt{84.885}{742.284}{7.97}{\infty}%
\TextAt{8}{93.852}{744.078}{.}%
\TextAt{8}{102.3}{744.078}{For}%
\TextAt{13}{123.27}{744.078}{f}%
\TextAt{9}{130.316}{744.078}{\Character{47}}%
\TextAt{13}{136.169}{744.078}{g}%
\NativeMathAt{145.524}{744.078}{11.955}{\in}%
\TextAt{18}{156.815}{744.078}{F}%
\TextAt{21}{164.121}{742.284}{p}%
\TextAt{9}{168.882}{744.078}{\Character{40}}%
\TextAt{13}{173.434}{744.078}{t}%
\TextAt{9}{177.661}{744.078}{\Character{41}}%
\TextAt{8}{182.214}{744.078}{,}%
\TextAt{13}{189.363}{744.078}{f}%
\TextAt{13}{195.747}{744.078}{\Character{59}}%
\TextAt{13}{200.984}{744.078}{g}%
\NativeMathAt{210.358}{744.078}{11.955}{\in}%
\TextAt{18}{221.649}{744.078}{F}%
\TextAt{21}{228.955}{742.284}{p}%
\TextAt{9}{233.716}{744.078}{\Character{91}}%
\TextAt{13}{236.968}{744.078}{t}%
\TextAt{9}{241.195}{744.078}{\Character{93}}%
\TextAt{8}{244.447}{744.078}{,}%
\TextAt{8}{251.596}{744.078}{it}%
\TextAt{8}{263.3}{744.078}{is}%
\TextAt{8}{275.064}{744.078}{defined}%
\TextAt{8}{315.377}{744.078}{by}%
\TextAt{13}{235.154}{710.48}{v}%
\NativeMathAt{240.813}{708.687}{7.97}{\infty}%
\TextAt{20}{249.78}{727.337}{\Character{18}}%
\TextAt{13}{259.775}{718.568}{f}%
\Vector{q 0 G 0 g 0.478 w 187.775 -60.531 m 194.822 -60.531 l S Q}%
\TextAt{13}{260.282}{702.28}{g}%
\TextAt{20}{268.017}{727.337}{\Character{19}}%
\TextAt{9}{280.139}{710.48}{\Character{61}}%
\TextAt{8}{292.564}{710.48}{deg}%
\TextAt{9}{310.126}{710.48}{\Character{40}}%
\TextAt{13}{314.679}{710.48}{g}%
\TextAt{9}{320.713}{710.48}{\Character{41}}%
\TextAt{16}{327.922}{710.48}{\Character{0}}%
\TextAt{8}{339.877}{710.48}{deg}%
\TextAt{9}{357.439}{710.48}{\Character{40}}%
\TextAt{13}{361.991}{710.48}{f}%
\TextAt{9}{369.038}{710.48}{\Character{41}}%
\TextAt{13}{373.59}{710.48}{\Character{58}}%
\TextAt{8}{79.226}{677.237}{It}%
\TextAt{8}{92.296}{677.237}{is}%
\TextAt{8}{104.456}{677.237}{associated}%
\TextAt{8}{160.588}{677.237}{to}%
\TextAt{8}{175.296}{677.237}{the}%
\TextAt{8}{195.849}{677.237}{prime}%
\TextAt{8}{229.41}{677.237}{ideal}%
\TextAt{14}{257.77}{677.237}{p}%
\TextAt{9}{267.75}{677.237}{\Character{61}}%
\TextAt{13}{280.852}{677.237}{y}%
\TextAt{18}{286.989}{677.237}{F}%
\TextAt{21}{294.294}{675.444}{p}%
\TextAt{9}{299.055}{677.237}{\Character{91}}%
\TextAt{13}{302.307}{677.237}{y}%
\TextAt{9}{308.444}{677.237}{\Character{93}}%
\TextAt{8}{315.99}{677.237}{of}%
\TextAt{8}{329.717}{677.237}{the}%
\TextAt{8}{350.27}{677.237}{ring}%
\TextAt{18}{374.733}{677.237}{F}%
\TextAt{21}{382.039}{675.444}{p}%
\TextAt{9}{386.8}{677.237}{\Character{91}}%
\TextAt{13}{390.051}{677.237}{y}%
\TextAt{9}{396.188}{677.237}{\Character{93}}%
\TextAt{8}{399.44}{677.237}{,}%
\TextAt{8}{407.085}{677.237}{where}%
\TextAt{13}{441.291}{677.237}{y}%
\TextAt{9}{451.425}{677.237}{\Character{61}}%
\TextAt{9}{464.516}{677.237}{1\Character{47}}%
\TextAt{13}{476.233}{677.237}{t}%
\TextAt{8}{480.46}{677.237}{.}%
\TextAt{8}{490.099}{677.237}{One}%
\TextAt{8}{515.207}{677.237}{can}%
\TextAt{8}{79.226}{660.192}{show}%
\TextAt{8}{108.229}{660.192}{that}%
\TextAt{8}{133.598}{660.192}{this}%
\TextAt{8}{156.421}{660.192}{exhausts}%
\TextAt{8}{204.349}{660.192}{all}%
\TextAt{8}{220.608}{660.192}{normalized}%
\TextAt{8}{280.444}{660.192}{valuations}%
\TextAt{8}{336.131}{660.192}{of}%
\TextAt{8}{349.461}{660.192}{the}%
\TextAt{8}{369.617}{660.192}{field}%
\TextAt{18}{394.974}{660.192}{F}%
\TextAt{21}{402.28}{658.398}{p}%
\TextAt{9}{407.041}{660.192}{\Character{40}}%
\TextAt{13}{411.593}{660.192}{t}%
\TextAt{9}{415.82}{660.192}{\Character{41}}%
\TextAt{8}{420.373}{660.192}{.}%
\TextAt{8}{97.159}{640.157}{For}%
\TextAt{8}{118.327}{640.157}{an}%
\TextAt{8}{134.785}{640.157}{arbitrary}%
\TextAt{8}{184.754}{640.157}{finite}%
\TextAt{8}{214.865}{640.157}{extension}%
\TextAt{13}{266.83}{640.157}{K}%
\TextAt{8}{281.744}{640.157}{of}%
\TextAt{18}{295.273}{640.157}{F}%
\TextAt{21}{302.579}{638.363}{p}%
\TextAt{9}{307.34}{640.157}{\Character{40}}%
\TextAt{13}{311.892}{640.157}{t}%
\TextAt{9}{316.119}{640.157}{\Character{41}}%
\TextAt{8}{320.672}{640.157}{,}%
\TextAt{8}{328.07}{640.157}{instead}%
\TextAt{8}{368.654}{640.157}{of}%
\TextAt{8}{382.183}{640.157}{restricting}%
\TextAt{8}{438.38}{640.157}{attention}%
\TextAt{8}{488.994}{640.157}{to}%
\TextAt{8}{503.503}{640.157}{prime}%
\TextAt{8}{79.226}{623.111}{ideals,}%
\TextAt{8}{115.797}{623.111}{one}%
\TextAt{8}{137.849}{623.111}{now}%
\TextAt{8}{162.83}{623.111}{considers}%
\TextAt{8}{213.623}{623.111}{all}%
\TextAt{8}{230.474}{623.111}{normalized}%
\TextAt{8}{290.903}{623.111}{discrete}%
\TextAt{8}{334.474}{623.111}{valuations}%
\TextAt{13}{390.754}{623.111}{v}%
\TextAt{27}{396.413}{621.318}{p}%
\TextAt{8}{405.39}{623.111}{of}%
\TextAt{13}{419.313}{623.111}{K}%
\TextAt{8}{434.621}{623.111}{in}%
\TextAt{8}{448.866}{623.111}{the}%
\TextAt{8}{469.615}{623.111}{above}%
\TextAt{8}{503.384}{623.111}{sense,}%
\TextAt{8}{79.226}{606.065}{where}%
\TextAt{8}{113.152}{606.065}{the}%
\TextAt{8}{133.426}{606.065}{index}%
\TextAt{14}{165.08}{606.065}{p}%
\TextAt{8}{175.078}{606.065}{has}%
\TextAt{8}{196.069}{606.065}{kept}%
\TextAt{8}{222.2}{606.065}{only}%
\TextAt{8}{248.009}{606.065}{a}%
\TextAt{8}{257.881}{606.065}{symbolic}%
\TextAt{8}{306.512}{606.065}{value.}%
\TextAt{8}{341.66}{606.065}{As}%
\TextAt{8}{359.064}{606.065}{an}%
\TextAt{8}{375.44}{606.065}{analogue}%
\TextAt{8}{424.346}{606.065}{of}%
\TextAt{8}{437.793}{606.065}{the}%
\TextAt{8}{458.067}{606.065}{ideal}%
\TextAt{8}{486.147}{606.065}{group}%
\TextAt{8}{519.439}{606.065}{we}%
\TextAt{8}{79.226}{589.019}{form}%
\TextAt{8}{106.866}{589.019}{the}%
\TextAt{8}{127.023}{589.019}{“divisor}%
\TextAt{8}{170.659}{589.019}{group”,}%
\TextAt{8}{212.61}{589.019}{i.e.,}%
\TextAt{8}{234.715}{589.019}{the}%
\TextAt{8}{254.872}{589.019}{free}%
\TextAt{8}{277.299}{589.019}{abelian}%
\TextAt{8}{317.947}{589.019}{group}%
\TextAt{8}{351.123}{589.019}{generated}%
\TextAt{8}{404.455}{589.019}{by}%
\TextAt{8}{420.714}{589.019}{these}%
\TextAt{8}{450.685}{589.019}{symbols,}%
\TextAt{13}{162.555}{552.727}{D}%
\TextAt{13}{172.597}{552.727}{iv}%
\TextAt{9}{182.681}{552.727}{\Character{40}}%
\TextAt{13}{187.234}{552.727}{K}%
\TextAt{9}{198.051}{552.727}{\Character{41}}%
\TextAt{9}{205.906}{552.727}{\Character{61}}%
\TextAt{20}{218.35}{573.171}{\Character{40}}%
\TextAt{20}{227.981}{564.085}{\Character{88}}%
\TextAt{27}{234.621}{538.979}{p}%
\TextAt{13}{247.242}{552.727}{n}%
\TextAt{27}{254.229}{550.934}{p}%
\TextAt{14}{258.716}{552.727}{p}%
\TextAt{20}{268.597}{573.649}{\Character{12}}%
\TextAt{20}{268.597}{566.476}{\Character{12}}%
\TextAt{20}{268.597}{559.303}{\Character{12}}%
\TextAt{20}{268.597}{552.129}{\Character{12}}%
\TextAt{20}{268.597}{544.956}{\Character{12}}%
\TextAt{13}{276.479}{552.727}{n}%
\TextAt{27}{283.467}{550.934}{p}%
\NativeMathAt{291.274}{552.727}{11.955}{\in}%
\TextAt{18}{302.565}{552.727}{Z}%
\TextAt{13}{310.536}{552.727}{\Character{59}}%
\TextAt{13}{319.658}{552.727}{n}%
\TextAt{27}{326.665}{550.934}{p}%
\TextAt{9}{334.472}{552.727}{\Character{61}}%
\TextAt{9}{346.882}{552.727}{0}%
\TextAt{8}{356.648}{552.727}{for}%
\TextAt{8}{374.533}{552.727}{almost}%
\TextAt{8}{412.323}{552.727}{all}%
\TextAt{14}{428.583}{552.727}{p}%
\TextAt{20}{434.566}{573.171}{\Character{41}}%
\TextAt{13}{446.189}{552.727}{\Character{58}}%
\TextAt{8}{79.226}{515.06}{We}%
\TextAt{8}{99.371}{515.06}{consider}%
\TextAt{8}{144.956}{515.06}{the}%
\TextAt{8}{165.112}{515.06}{mapping}%
\TextAt{8}{186.442}{486.092}{div}%
\TextAt{9}{205.699}{486.092}{\Character{58}}%
\TextAt{13}{212.271}{486.092}{K}%
\TextAt{3}{223.089}{491.028}{\Character{3}}%
\NativeArrowAt{231.142}{486.092}{11.955}{19.26}{\longrightarrow}%
\TextAt{13}{253.724}{486.092}{D}%
\TextAt{13}{263.766}{486.092}{iv}%
\TextAt{9}{273.85}{486.092}{\Character{40}}%
\TextAt{13}{278.403}{486.092}{K}%
\TextAt{9}{289.22}{486.092}{\Character{41}}%
\TextAt{13}{293.773}{486.092}{\Character{59}}%
\TextAt{8}{322.927}{486.092}{div}%
\TextAt{9}{338.863}{486.092}{\Character{40}}%
\TextAt{13}{343.416}{486.092}{f}%
\TextAt{9}{350.462}{486.092}{\Character{41}}%
\TextAt{9}{358.317}{486.092}{\Character{61}}%
\TextAt{20}{370.761}{497.449}{\Character{88}}%
\TextAt{27}{377.401}{472.343}{p}%
\TextAt{13}{390.022}{486.092}{v}%
\TextAt{27}{395.681}{484.298}{p}%
\TextAt{9}{400.168}{486.092}{\Character{40}}%
\TextAt{13}{404.72}{486.092}{f}%
\TextAt{9}{411.767}{486.092}{\Character{41}}%
\TextAt{14}{416.319}{486.092}{p}%
\TextAt{13}{422.302}{486.092}{\Character{59}}%
\TextAt{8}{79.226}{447.755}{the}%
\TextAt{8}{99.383}{447.755}{image}%
\TextAt{8}{133.204}{447.755}{of}%
\TextAt{8}{146.534}{447.755}{which}%
\TextAt{8}{180.02}{447.755}{is}%
\TextAt{8}{191.784}{447.755}{written}%
\TextAt{13}{232.754}{447.755}{P}%
\TextAt{9}{241.926}{447.755}{\Character{40}}%
\TextAt{13}{246.478}{447.755}{K}%
\TextAt{9}{257.296}{447.755}{\Character{41}}%
\TextAt{8}{261.848}{447.755}{,}%
\TextAt{8}{268.997}{447.755}{and}%
\TextAt{8}{291.76}{447.755}{we}%
\TextAt{8}{308.987}{447.755}{define}%
\TextAt{8}{342.796}{447.755}{the}%
\TextAt{12}{362.953}{447.755}{divisor}%
\TextAt{12}{407.031}{447.755}{class}%
\TextAt{12}{438.39}{447.755}{group}%
\TextAt{8}{476.192}{447.755}{of}%
\TextAt{13}{489.522}{447.755}{K}%
\TextAt{8}{504.237}{447.755}{by}%
\TextAt{13}{242.549}{418.787}{C}%
\TextAt{13}{251.767}{418.787}{l}%
\TextAt{9}{255.532}{418.787}{\Character{40}}%
\TextAt{13}{260.084}{418.787}{K}%
\TextAt{9}{270.902}{418.787}{\Character{41}}%
\TextAt{9}{278.757}{418.787}{\Character{61}}%
\TextAt{13}{291.201}{418.787}{D}%
\TextAt{13}{301.243}{418.787}{iv}%
\TextAt{9}{311.327}{418.787}{\Character{40}}%
\TextAt{13}{315.879}{418.787}{K}%
\TextAt{9}{326.697}{418.787}{\Character{41}\Character{47}}%
\TextAt{13}{337.102}{418.787}{P}%
\TextAt{9}{346.273}{418.787}{\Character{40}}%
\TextAt{13}{350.825}{418.787}{K}%
\TextAt{9}{361.643}{418.787}{\Character{41}}%
\TextAt{13}{366.195}{418.787}{\Character{58}}%
\TextAt{8}{79.226}{389.818}{Unlike}%
\TextAt{8}{115.7}{389.818}{the}%
\TextAt{8}{135.592}{389.818}{ideal}%
\TextAt{8}{163.291}{389.818}{class}%
\TextAt{8}{190.463}{389.818}{group}%
\TextAt{8}{223.374}{389.818}{of}%
\TextAt{8}{236.44}{389.818}{an}%
\TextAt{8}{252.434}{389.818}{algebraic}%
\TextAt{8}{301.605}{389.818}{number}%
\TextAt{8}{343.936}{389.818}{field,}%
\TextAt{8}{372.333}{389.818}{this}%
\TextAt{8}{394.891}{389.818}{group}%
\TextAt{8}{427.803}{389.818}{is}%
\TextAt{8}{439.302}{389.818}{not}%
\TextAt{8}{459.851}{389.818}{finite.}%
\TextAt{8}{494.226}{389.818}{Rather,}%
\TextAt{8}{79.226}{372.772}{one}%
\TextAt{8}{100.686}{372.772}{has}%
\TextAt{8}{121.559}{372.772}{the}%
\TextAt{8}{141.716}{372.772}{canonical}%
\TextAt{8}{193.099}{372.772}{homomorphism}%
\TextAt{8}{257.255}{343.803}{deg}%
\TextAt{9}{278.138}{343.803}{\Character{58}}%
\TextAt{13}{284.71}{343.803}{C}%
\TextAt{13}{293.928}{343.803}{l}%
\TextAt{9}{297.694}{343.803}{\Character{40}}%
\TextAt{13}{302.246}{343.803}{K}%
\TextAt{9}{313.064}{343.803}{\Character{41}}%
\NativeArrowAt{320.937}{343.803}{11.955}{19.26}{\longrightarrow}%
\TextAt{18}{343.519}{343.803}{Z}%
\TextAt{13}{351.489}{343.803}{\Character{59}}%
\TextAt{8}{79.226}{314.834}{which}%
\TextAt{8}{113.984}{314.834}{associates}%
\TextAt{8}{169.101}{314.834}{to}%
\TextAt{8}{184.682}{314.834}{the}%
\TextAt{8}{206.11}{314.834}{class}%
\TextAt{8}{234.818}{314.834}{of}%
\TextAt{14}{249.419}{314.834}{p}%
\TextAt{8}{260.57}{314.834}{the}%
\TextAt{8}{281.998}{314.834}{degree}%
\TextAt{8}{319.684}{314.834}{deg}%
\TextAt{9}{337.247}{314.834}{\Character{40}}%
\TextAt{14}{341.799}{314.834}{p}%
\TextAt{9}{347.782}{314.834}{\Character{41}}%
\TextAt{9}{357.8}{314.834}{\Character{61}}%
\TextAt{9}{372.386}{314.834}{\Character{91}}%
\TextAt{13}{375.664}{314.834}{\Character{20}}%
\TextAt{9}{382.403}{314.834}{\Character{40}}%
\TextAt{14}{386.955}{314.834}{p}%
\TextAt{9}{392.938}{314.834}{\Character{41}}%
\TextAt{9}{402.957}{314.834}{\Character{58}}%
\TextAt{18}{411.716}{314.834}{F}%
\TextAt{21}{419.022}{313.041}{p}%
\TextAt{9}{423.783}{314.834}{\Character{93}}%
\TextAt{8}{432.203}{314.834}{of}%
\TextAt{8}{446.804}{314.834}{the}%
\TextAt{8}{468.232}{314.834}{residue}%
\TextAt{8}{509.23}{314.834}{class}%
\TextAt{8}{79.226}{297.788}{field}%
\TextAt{8}{105.046}{297.788}{of}%
\TextAt{8}{118.839}{297.788}{the}%
\TextAt{8}{139.459}{297.788}{valuation}%
\TextAt{8}{190.995}{297.788}{ring}%
\TextAt{8}{215.524}{297.788}{of}%
\TextAt{14}{229.317}{297.788}{p}%
\TextAt{8}{235.3}{297.788}{,}%
\TextAt{8}{243.028}{297.788}{and}%
\TextAt{8}{266.254}{297.788}{which}%
\TextAt{8}{300.203}{297.788}{associates}%
\TextAt{8}{354.513}{297.788}{to}%
\TextAt{8}{369.286}{297.788}{the}%
\TextAt{8}{389.906}{297.788}{class}%
\TextAt{8}{417.806}{297.788}{of}%
\TextAt{8}{431.599}{297.788}{an}%
\TextAt{8}{448.321}{297.788}{arbitrary}%
\TextAt{8}{498.554}{297.788}{divisor}%
\TextAt{14}{79.226}{280.742}{a}%
\TextAt{9}{88.53}{280.742}{\Character{61}}%
\TextAt{20}{100.956}{289.709}{\Character{80}}%
\TextAt{27}{113.575}{277.255}{p}%
\TextAt{13}{120.054}{280.742}{n}%
\TextAt{27}{127.042}{278.949}{p}%
\TextAt{14}{131.529}{280.742}{p}%
\TextAt{8}{141.409}{280.742}{the}%
\TextAt{8}{161.566}{280.742}{sum}%
\TextAt{8}{247.485}{262.199}{deg}%
\TextAt{9}{265.047}{262.199}{\Character{40}}%
\TextAt{14}{269.599}{262.199}{a}%
\TextAt{9}{275.582}{262.199}{\Character{41}}%
\TextAt{9}{283.437}{262.199}{\Character{61}}%
\TextAt{20}{295.881}{273.557}{\Character{88}}%
\TextAt{27}{302.521}{248.451}{p}%
\TextAt{13}{315.142}{262.199}{n}%
\TextAt{27}{322.13}{260.406}{p}%
\TextAt{8}{328.609}{262.199}{deg}%
\TextAt{9}{346.171}{262.199}{\Character{40}}%
\TextAt{14}{350.724}{262.199}{p}%
\TextAt{9}{356.707}{262.199}{\Character{41}}%
\TextAt{13}{361.259}{262.199}{\Character{58}}%
\TextAt{8}{79.226}{229.324}{For}%
\TextAt{8}{99.026}{229.324}{a}%
\TextAt{8}{107.612}{229.324}{principal}%
\TextAt{8}{155.22}{229.324}{divisor}%
\TextAt{8}{192.163}{229.324}{div}%
\TextAt{9}{208.1}{229.324}{\Character{40}}%
\TextAt{13}{212.652}{229.324}{f}%
\TextAt{9}{219.698}{229.324}{\Character{41}}%
\TextAt{8}{224.251}{229.324}{,}%
\TextAt{13}{230.464}{229.324}{f}%
\NativeMathAt{240.832}{229.324}{11.955}{\in}%
\TextAt{13}{252.123}{229.324}{K}%
\TextAt{3}{262.94}{233.662}{\Character{3}}%
\TextAt{8}{267.673}{229.324}{,}%
\TextAt{8}{273.886}{229.324}{we}%
\TextAt{8}{289.944}{229.324}{find}%
\TextAt{8}{312.183}{229.324}{by}%
\TextAt{8}{327.273}{229.324}{an}%
\TextAt{8}{342.362}{229.324}{easy}%
\TextAt{8}{366.944}{229.324}{calculation}%
\TextAt{8}{424.965}{229.324}{that}%
\TextAt{8}{449.164}{229.324}{deg}%
\TextAt{9}{466.726}{229.324}{\Character{40}}%
\TextAt{8}{471.279}{229.324}{div}%
\TextAt{9}{487.215}{229.324}{\Character{40}}%
\TextAt{13}{491.767}{229.324}{f}%
\TextAt{9}{498.814}{229.324}{\Character{41}\Character{41}}%
\TextAt{9}{511.212}{229.324}{\Character{61}}%
\TextAt{9}{523.633}{229.324}{0}%
\TextAt{8}{529.518}{229.324}{,}%
\TextAt{8}{79.226}{212.278}{so}%
\TextAt{8}{94.207}{212.278}{that}%
\TextAt{8}{120.186}{212.278}{the}%
\TextAt{8}{140.953}{212.278}{mapping}%
\TextAt{8}{189.696}{212.278}{deg}%
\TextAt{8}{211.766}{212.278}{is}%
\TextAt{8}{224.14}{212.278}{indeed}%
\TextAt{8}{261.812}{212.278}{well-defined.}%
\TextAt{8}{332.239}{212.278}{As}%
\TextAt{8}{350.137}{212.278}{an}%
\TextAt{8}{367.006}{212.278}{analogue}%
\TextAt{8}{416.406}{212.278}{of}%
\TextAt{8}{430.347}{212.278}{the}%
\TextAt{8}{451.114}{212.278}{finiteness}%
\TextAt{8}{502.57}{212.278}{of}%
\TextAt{8}{516.51}{212.278}{the}%
\TextAt{8}{79.226}{195.232}{class}%
\TextAt{8}{107.251}{195.232}{number}%
\TextAt{8}{150.436}{195.232}{of}%
\TextAt{8}{164.354}{195.232}{an}%
\TextAt{8}{181.201}{195.232}{algebraic}%
\TextAt{8}{231.224}{195.232}{number}%
\TextAt{8}{274.409}{195.232}{field,}%
\TextAt{8}{303.752}{195.232}{one}%
\TextAt{8}{325.8}{195.232}{obtains}%
\TextAt{8}{367.43}{195.232}{here}%
\TextAt{8}{393.376}{195.232}{the}%
\TextAt{8}{414.12}{195.232}{fact}%
\TextAt{8}{437.794}{195.232}{that,}%
\TextAt{8}{467.15}{195.232}{if}%
\TextAt{8}{478.462}{195.232}{not}%
\TextAt{13}{499.864}{195.232}{C}%
\TextAt{13}{509.082}{195.232}{l}%
\TextAt{9}{512.847}{195.232}{\Character{40}}%
\TextAt{13}{517.399}{195.232}{K}%
\TextAt{9}{528.217}{195.232}{\Character{41}}%
\TextAt{8}{79.226}{178.186}{itself,}%
\TextAt{8}{112.176}{178.186}{the}%
\TextAt{8}{133.414}{178.186}{kernel}%
\TextAt{13}{168.963}{178.186}{C}%
\TextAt{13}{178.181}{178.186}{l}%
\TextAt{17}{181.946}{182.524}{0}%
\TextAt{9}{186.679}{178.186}{\Character{40}}%
\TextAt{13}{191.231}{178.186}{K}%
\TextAt{9}{202.049}{178.186}{\Character{41}}%
\TextAt{8}{211.58}{178.186}{of}%
\TextAt{8}{225.992}{178.186}{deg}%
\TextAt{8}{248.534}{178.186}{is}%
\TextAt{8}{261.38}{178.186}{finite.}%
\TextAt{8}{299.088}{178.186}{The}%
\TextAt{8}{324.224}{178.186}{infinitude}%
\TextAt{8}{377.98}{178.186}{of}%
\TextAt{8}{392.392}{178.186}{the}%
\TextAt{8}{413.63}{178.186}{class}%
\TextAt{8}{442.149}{178.186}{group}%
\TextAt{8}{476.407}{178.186}{of}%
\TextAt{8}{490.819}{178.186}{function}%
\TextAt{8}{79.226}{161.14}{fields}%
\TextAt{8}{108.688}{161.14}{must}%
\TextAt{8}{137.182}{161.14}{not}%
\TextAt{8}{157.487}{161.14}{be}%
\TextAt{8}{172.902}{161.14}{considered}%
\TextAt{8}{229.682}{161.14}{as}%
\TextAt{8}{243.543}{161.14}{strange.}%
\TextAt{8}{288.966}{161.14}{On}%
\TextAt{8}{307.968}{161.14}{the}%
\TextAt{8}{327.615}{161.14}{contrary,}%
\TextAt{8}{376.319}{161.14}{it}%
\TextAt{8}{387.514}{161.14}{is}%
\TextAt{8}{398.769}{161.14}{rather}%
\TextAt{8}{433.384}{161.14}{the}%
\TextAt{8}{453.031}{161.14}{finiteness}%
\TextAt{8}{503.367}{161.14}{in}%
\TextAt{8}{516.51}{161.14}{the}%
\TextAt{8}{79.226}{144.094}{number}%
\TextAt{8}{121.8}{144.094}{field}%
\TextAt{8}{147.135}{144.094}{case}%
\TextAt{8}{171.884}{144.094}{that}%
\TextAt{8}{197.231}{144.094}{should}%
\TextAt{8}{234.342}{144.094}{be}%
\TextAt{8}{250.244}{144.094}{regarded}%
\TextAt{8}{298.353}{144.094}{as}%
\TextAt{8}{312.701}{144.094}{a}%
\TextAt{8}{322.434}{144.094}{deficiency}%
\TextAt{8}{376.055}{144.094}{which}%
\TextAt{8}{409.519}{144.094}{calls}%
\TextAt{8}{435.572}{144.094}{for}%
\TextAt{8}{453.434}{144.094}{correction.}%
\TextAt{8}{512.613}{144.094}{The}%
\TextAt{8}{79.226}{127.048}{adequate}%
\TextAt{8}{128.815}{127.048}{appreciation}%
\TextAt{8}{195.644}{127.048}{of}%
\TextAt{8}{208.805}{127.048}{this}%
\TextAt{8}{231.46}{127.048}{situation}%
\TextAt{8}{280.14}{127.048}{and}%
\TextAt{8}{302.734}{127.048}{its}%
\TextAt{8}{318.885}{127.048}{amendment}%
\TextAt{8}{382.127}{127.048}{will}%
\TextAt{8}{404.064}{127.048}{be}%
\TextAt{8}{419.82}{127.048}{explained}%
\TextAt{8}{472.003}{127.048}{in}%
\TextAt{8}{485.488}{127.048}{chap.}%
\TextAt{8}{516.212}{127.048}{III,}%
\TextAt{8}{79.226}{110.002}{§1.}%
\TextAt{8}{97.159}{89.967}{The}%
\TextAt{8}{121.607}{89.967}{ideal,}%
\TextAt{8}{153.315}{89.967}{completely}%
\TextAt{8}{212.564}{89.967}{satisfactory}%
\TextAt{8}{275.532}{89.967}{framework}%
\TextAt{8}{333.49}{89.967}{for}%
\TextAt{8}{351.77}{89.967}{the}%
\TextAt{8}{372.32}{89.967}{theory}%
\TextAt{8}{409.465}{89.967}{of}%
\TextAt{8}{423.189}{89.967}{function}%
\TextAt{8}{469.431}{89.967}{fields}%
\TextAt{8}{499.797}{89.967}{is}%
\TextAt{8}{511.956}{89.967}{pro-}%
\TextAt{8}{79.226}{72.921}{vided}%
\TextAt{8}{110.963}{72.921}{by}%
\TextAt{8}{127.421}{72.921}{the}%
\TextAt{8}{147.777}{72.921}{notion}%
\TextAt{8}{184.404}{72.921}{of}%
\TextAt{8}{197.933}{72.921}{scheme.}%
\TextAt{8}{243.132}{72.921}{In}%
\TextAt{8}{257.952}{72.921}{the}%
\TextAt{8}{278.308}{72.921}{last}%
\TextAt{8}{300.684}{72.921}{section}%
\TextAt{8}{339.964}{72.921}{we}%
\TextAt{8}{357.391}{72.921}{introduced}%
\TextAt{19}{416.123}{72.921}{affine}%
\TextAt{8}{448.553}{72.921}{schemes}%
\TextAt{8}{493.417}{72.921}{as}%
\TextAt{8}{507.986}{72.921}{pairs}%
\TextAt{8}{300.14}{43.033}{89}%
\Destination{78.23}{809.815}{nav.page.89}%
\Destination{79.36}{404.855}{nav.print.96}%
\Hotlink{515.712}{124.855}{14.294}{11.047}{nav.chapter.3}%
\Hotlink{84.405}{107.755}{6.858}{11.047}{nav.section.3.1}%
}

% Page 102 (body page 90)
\Page{%
\TextAt{2}{79.226}{796.943}{I.}%
\TextAt{2}{89.83}{796.943}{Algebraic}%
\TextAt{2}{138.648}{796.943}{Integers}%
\Vector{q 0 G 0 g 0.299 w 7.226 18.729 m 460.769 18.729 l S Q}%
\TextAt{9}{79.226}{761.124}{\Character{40}}%
\TextAt{13}{83.778}{761.124}{X}%
\TextAt{13}{93.773}{761.124}{\Character{59}}%
\TextAt{14}{99.027}{761.124}{o}%
\TextAt{21}{104.872}{759.33}{X}%
\TextAt{9}{112.957}{761.124}{\Character{41}}%
\TextAt{8}{120.779}{761.124}{consisting}%
\TextAt{8}{174.261}{761.124}{of}%
\TextAt{8}{186.964}{761.124}{a}%
\TextAt{8}{196.092}{761.124}{topological}%
\TextAt{8}{254.989}{761.124}{space}%
\TextAt{13}{285.637}{761.124}{X}%
\TextAt{9}{299.613}{761.124}{\Character{61}}%
\TextAt{8}{312.038}{761.124}{Spec}%
\TextAt{9}{335.447}{761.124}{\Character{40}}%
\TextAt{14}{339.999}{761.124}{o}%
\TextAt{9}{345.843}{761.124}{\Character{41}}%
\TextAt{8}{353.666}{761.124}{and}%
\TextAt{8}{375.801}{761.124}{a}%
\TextAt{8}{384.929}{761.124}{sheaf}%
\TextAt{8}{413.951}{761.124}{of}%
\TextAt{8}{426.654}{761.124}{rings}%
\TextAt{14}{454.707}{761.124}{o}%
\TextAt{21}{460.551}{759.33}{X}%
\TextAt{8}{471.906}{761.124}{on}%
\TextAt{13}{487.538}{761.124}{X}%
\TextAt{8}{498.193}{761.124}{.}%
\TextAt{8}{506.432}{761.124}{More}%
\TextAt{8}{79.226}{744.078}{generally,}%
\TextAt{8}{130.933}{744.078}{a}%
\TextAt{12}{140.287}{744.078}{scheme}%
\TextAt{8}{185.649}{744.078}{is}%
\TextAt{8}{197.011}{744.078}{a}%
\TextAt{8}{206.364}{744.078}{topological}%
\TextAt{8}{265.487}{744.078}{space}%
\TextAt{13}{296.359}{744.078}{X}%
\TextAt{8}{310.509}{744.078}{with}%
\TextAt{8}{336.767}{744.078}{a}%
\TextAt{8}{346.121}{744.078}{sheaf}%
\TextAt{8}{375.367}{744.078}{of}%
\TextAt{8}{388.295}{744.078}{rings}%
\TextAt{14}{416.574}{744.078}{o}%
\TextAt{21}{422.418}{742.284}{X}%
\TextAt{8}{433.998}{744.078}{such}%
\TextAt{8}{459.993}{744.078}{that,}%
\TextAt{8}{488.292}{744.078}{for}%
\TextAt{8}{505.775}{744.078}{every}%
\TextAt{8}{79.226}{727.032}{point}%
\TextAt{8}{109.427}{727.032}{of}%
\TextAt{13}{122.388}{727.032}{X}%
\TextAt{8}{133.043}{727.032}{,}%
\TextAt{8}{139.897}{727.032}{there}%
\TextAt{8}{169.44}{727.032}{exists}%
\TextAt{8}{201.386}{727.032}{a}%
\TextAt{8}{210.773}{727.032}{neighbourhood}%
\TextAt{13}{290.408}{727.032}{U}%
\TextAt{8}{303.135}{727.032}{which,}%
\TextAt{8}{339.579}{727.032}{together}%
\TextAt{8}{385.393}{727.032}{with}%
\TextAt{8}{411.684}{727.032}{the}%
\TextAt{8}{431.471}{727.032}{restriction}%
\TextAt{14}{487.1}{727.032}{o}%
\TextAt{21}{492.945}{725.238}{U}%
\TextAt{8}{503.549}{727.032}{of}%
\TextAt{8}{516.51}{727.032}{the}%
\TextAt{8}{79.226}{709.986}{sheaf}%
\TextAt{14}{108.98}{709.986}{o}%
\TextAt{21}{114.825}{708.193}{X}%
\TextAt{8}{126.912}{709.986}{to}%
\TextAt{13}{141.328}{709.986}{U}%
\TextAt{8}{150.527}{709.986}{,}%
\TextAt{8}{157.808}{709.986}{is}%
\TextAt{8}{169.677}{709.986}{isomorphic}%
\TextAt{8}{229.032}{709.986}{to}%
\TextAt{8}{243.448}{709.986}{an}%
\TextAt{8}{259.812}{709.986}{affine}%
\TextAt{8}{291.132}{709.986}{scheme}%
\TextAt{8}{331.287}{709.986}{in}%
\TextAt{8}{345.046}{709.986}{the}%
\TextAt{8}{365.307}{709.986}{sense}%
\TextAt{8}{395.444}{709.986}{of}%
\TextAt{8}{408.879}{709.986}{§13.}%
\TextAt{8}{435.038}{709.986}{This}%
\TextAt{8}{461.863}{709.986}{generalization}%
\TextAt{8}{79.226}{692.94}{of}%
\TextAt{8}{92.564}{692.94}{affine}%
\TextAt{8}{123.787}{692.94}{schemes}%
\TextAt{8}{168.459}{692.94}{is}%
\TextAt{8}{180.231}{692.94}{the}%
\TextAt{8}{200.395}{692.94}{correct}%
\TextAt{8}{239.425}{692.94}{notion}%
\TextAt{8}{275.86}{692.94}{for}%
\TextAt{8}{293.753}{692.94}{a}%
\TextAt{8}{303.516}{692.94}{function}%
\TextAt{8}{349.372}{692.94}{field}%
\TextAt{13}{374.737}{692.94}{K}%
\TextAt{8}{385.554}{692.94}{.}%
\TextAt{8}{394.026}{692.94}{It}%
\TextAt{8}{406.707}{692.94}{shows}%
\TextAt{8}{440.332}{692.94}{all}%
\TextAt{8}{456.599}{692.94}{prime}%
\TextAt{8}{489.771}{692.94}{ideals}%
\TextAt{8}{522.357}{692.94}{at}%
\TextAt{8}{79.226}{675.894}{once,}%
\TextAt{8}{109.138}{675.894}{and}%
\TextAt{8}{131.901}{675.894}{misses}%
\TextAt{8}{167.85}{675.894}{none.}%
\TextAt{8}{97.159}{655.859}{In}%
\TextAt{8}{112.013}{655.859}{the}%
\TextAt{8}{132.403}{655.859}{case}%
\TextAt{13}{157.407}{655.859}{K}%
\TextAt{9}{171.943}{655.859}{\Character{61}}%
\TextAt{18}{184.766}{655.859}{F}%
\TextAt{21}{192.072}{654.066}{p}%
\TextAt{9}{196.832}{655.859}{\Character{40}}%
\TextAt{13}{201.385}{655.859}{t}%
\TextAt{9}{205.612}{655.859}{\Character{41}}%
\TextAt{8}{214.295}{655.859}{for}%
\TextAt{8}{232.413}{655.859}{instance,}%
\TextAt{8}{281.541}{655.859}{the}%
\TextAt{8}{301.931}{655.859}{corresponding}%
\TextAt{8}{377.35}{655.859}{scheme}%
\TextAt{9}{417.633}{655.859}{\Character{40}}%
\TextAt{13}{422.185}{655.859}{X}%
\TextAt{13}{432.18}{655.859}{\Character{59}}%
\TextAt{14}{437.434}{655.859}{o}%
\TextAt{21}{443.278}{654.066}{X}%
\TextAt{9}{451.363}{655.859}{\Character{41}}%
\TextAt{8}{460.046}{655.859}{is}%
\TextAt{8}{472.043}{655.859}{obtained}%
\TextAt{8}{520.408}{655.859}{by}%
\TextAt{8}{79.226}{638.813}{gluing}%
\TextAt{8}{114.765}{638.813}{the}%
\TextAt{8}{135.335}{638.813}{two}%
\TextAt{8}{157.867}{638.813}{rings}%
\TextAt{13}{186.961}{638.813}{A}%
\TextAt{9}{199.763}{638.813}{\Character{61}}%
\TextAt{18}{212.895}{638.813}{F}%
\TextAt{21}{220.201}{637.02}{p}%
\TextAt{9}{224.962}{638.813}{\Character{91}}%
\TextAt{13}{228.213}{638.813}{u}%
\TextAt{9}{234.876}{638.813}{\Character{93}}%
\TextAt{8}{242.439}{638.813}{and}%
\TextAt{13}{265.616}{638.813}{B}%
\TextAt{9}{279.139}{638.813}{\Character{61}}%
\TextAt{18}{292.27}{638.813}{F}%
\TextAt{21}{299.576}{637.02}{p}%
\TextAt{9}{304.337}{638.813}{\Character{91}}%
\TextAt{13}{307.589}{638.813}{v}%
\TextAt{9}{313.677}{638.813}{\Character{93}}%
\TextAt{8}{316.929}{638.813}{,}%
\TextAt{8}{324.596}{638.813}{or}%
\TextAt{8}{339.32}{638.813}{more}%
\TextAt{8}{369.001}{638.813}{precisely}%
\TextAt{8}{417.271}{638.813}{the}%
\TextAt{8}{437.842}{638.813}{two}%
\TextAt{8}{460.373}{638.813}{affine}%
\TextAt{8}{492.002}{638.813}{schemes}%
\TextAt{13}{79.226}{621.767}{U}%
\TextAt{9}{93.182}{621.767}{\Character{61}}%
\TextAt{8}{107.044}{621.767}{Spec}%
\TextAt{9}{130.452}{621.767}{\Character{40}}%
\TextAt{13}{135.005}{621.767}{A}%
\TextAt{9}{143.78}{621.767}{\Character{41}}%
\TextAt{8}{153.073}{621.767}{and}%
\TextAt{13}{176.678}{621.767}{V}%
\TextAt{9}{190.861}{621.767}{\Character{61}}%
\TextAt{8}{204.723}{621.767}{Spec}%
\TextAt{9}{228.131}{621.767}{\Character{40}}%
\TextAt{13}{232.683}{621.767}{B}%
\TextAt{9}{242.18}{621.767}{\Character{41}}%
\TextAt{8}{246.732}{621.767}{.}%
\TextAt{8}{257.709}{621.767}{Removing}%
\TextAt{8}{313.354}{621.767}{from}%
\TextAt{13}{341.838}{621.767}{U}%
\TextAt{8}{355.777}{621.767}{the}%
\TextAt{8}{376.776}{621.767}{point}%
\TextAt{14}{408.188}{621.767}{p}%
\TextAt{17}{414.172}{619.974}{0}%
\TextAt{9}{423.661}{621.767}{\Character{61}}%
\TextAt{9}{437.505}{621.767}{\Character{40}}%
\TextAt{13}{442.075}{621.767}{u}%
\TextAt{9}{448.738}{621.767}{\Character{41}}%
\TextAt{8}{453.29}{621.767}{,}%
\TextAt{8}{461.493}{621.767}{and}%
\TextAt{8}{485.098}{621.767}{the}%
\TextAt{8}{506.098}{621.767}{point}%
\TextAt{14}{79.226}{604.721}{p}%
\NativeMathAt{85.209}{602.928}{7.97}{\infty}%
\TextAt{9}{97.577}{604.721}{\Character{61}}%
\TextAt{9}{110.07}{604.721}{\Character{40}}%
\TextAt{13}{114.635}{604.721}{v}%
\TextAt{9}{120.723}{604.721}{\Character{41}}%
\TextAt{8}{129.22}{604.721}{from}%
\TextAt{13}{156.908}{604.721}{V}%
\TextAt{8}{166.334}{604.721}{,}%
\TextAt{8}{173.542}{604.721}{one}%
\TextAt{8}{195.049}{604.721}{has}%
\TextAt{13}{215.969}{604.721}{U}%
\TextAt{16}{227.857}{604.721}{\Character{110}}%
\TextAt{16}{236.513}{604.721}{\Character{102}}%
\TextAt{14}{242.501}{604.721}{p}%
\TextAt{17}{248.485}{602.928}{0}%
\TextAt{16}{253.217}{604.721}{\Character{103}}%
\TextAt{9}{262.596}{604.721}{\Character{61}}%
\TextAt{8}{275.102}{604.721}{Spec}%
\TextAt{9}{298.51}{604.721}{\Character{40}}%
\TextAt{18}{303.062}{604.721}{F}%
\TextAt{21}{310.368}{602.928}{p}%
\TextAt{9}{315.129}{604.721}{\Character{91}}%
\TextAt{13}{318.381}{604.721}{u\Character{59}}%
\TextAt{13}{330.264}{604.721}{u}%
\TextAt{3}{336.95}{609.06}{\Character{0}}%
\TextAt{17}{343.536}{609.06}{1}%
\TextAt{9}{348.269}{604.721}{\Character{93}\Character{41}}%
\TextAt{8}{356.073}{604.721}{,}%
\TextAt{13}{363.281}{604.721}{V}%
\TextAt{16}{375.395}{604.721}{\Character{110}}%
\TextAt{16}{384.051}{604.721}{\Character{102}}%
\TextAt{14}{390.039}{604.721}{p}%
\NativeMathAt{396.022}{602.928}{7.97}{\infty}%
\TextAt{16}{404.989}{604.721}{\Character{103}}%
\TextAt{9}{414.368}{604.721}{\Character{61}}%
\TextAt{8}{426.874}{604.721}{Spec}%
\TextAt{9}{450.282}{604.721}{\Character{40}}%
\TextAt{18}{454.834}{604.721}{F}%
\TextAt{21}{462.14}{602.928}{p}%
\TextAt{9}{466.901}{604.721}{\Character{91}}%
\TextAt{13}{470.153}{604.721}{v}%
\TextAt{13}{476.226}{604.721}{\Character{59}}%
\TextAt{13}{481.463}{604.721}{v}%
\TextAt{3}{487.573}{609.06}{\Character{0}}%
\TextAt{17}{494.16}{609.06}{1}%
\TextAt{9}{498.892}{604.721}{\Character{93}\Character{41}}%
\TextAt{8}{506.696}{604.721}{,}%
\TextAt{8}{513.904}{604.721}{and}%
\TextAt{8}{79.226}{587.675}{the}%
\TextAt{8}{99.383}{587.675}{isomorphism}%
\TextAt{13}{167.802}{587.675}{f}%
\TextAt{9}{178.169}{587.675}{\Character{58}}%
\TextAt{18}{184.742}{587.675}{F}%
\TextAt{21}{192.048}{585.882}{p}%
\TextAt{9}{196.809}{587.675}{\Character{91}}%
\TextAt{13}{200.06}{587.675}{u\Character{59}}%
\TextAt{13}{211.944}{587.675}{u}%
\TextAt{3}{218.629}{592.014}{\Character{0}}%
\TextAt{17}{225.216}{592.014}{1}%
\TextAt{9}{229.948}{587.675}{\Character{93}}%
\NativeMathAt{236.521}{587.675}{11.955}{\rightarrow}%
\TextAt{18}{251.797}{587.675}{F}%
\TextAt{21}{259.103}{585.882}{p}%
\TextAt{9}{263.863}{587.675}{\Character{91}}%
\TextAt{13}{267.115}{587.675}{v}%
\TextAt{13}{273.188}{587.675}{\Character{59}}%
\TextAt{13}{278.425}{587.675}{v}%
\TextAt{3}{284.536}{592.014}{\Character{0}}%
\TextAt{17}{291.122}{592.014}{1}%
\TextAt{9}{295.855}{587.675}{\Character{93}}%
\TextAt{8}{299.106}{587.675}{,}%
\TextAt{13}{306.255}{587.675}{u}%
\NativeArrowAt{316.239}{587.675}{11.955}{11.955}{\mapsto}%
\TextAt{13}{331.515}{587.675}{v}%
\TextAt{3}{337.603}{592.014}{\Character{0}}%
\TextAt{17}{344.19}{592.014}{1}%
\TextAt{8}{348.922}{587.675}{,}%
\TextAt{8}{356.071}{587.675}{yields}%
\TextAt{8}{388.972}{587.675}{a}%
\TextAt{8}{398.727}{587.675}{bijection}%
\TextAt{13}{190.002}{558.674}{\Character{39}}%
\TextAt{9}{200.991}{558.674}{\Character{58}}%
\TextAt{13}{207.564}{558.674}{V}%
\TextAt{16}{219.646}{558.674}{\Character{110}}%
\TextAt{16}{228.278}{558.674}{\Character{102}}%
\TextAt{14}{234.258}{558.674}{p}%
\NativeMathAt{240.241}{556.881}{7.97}{\infty}%
\TextAt{16}{249.208}{558.674}{\Character{103}}%
\NativeArrowAt{258.497}{558.674}{11.955}{19.272}{\longrightarrow}%
\TextAt{13}{281.088}{558.674}{U}%
\TextAt{16}{292.944}{558.674}{\Character{110}}%
\TextAt{16}{301.576}{558.674}{\Character{102}}%
\TextAt{14}{307.556}{558.674}{p}%
\TextAt{17}{313.539}{556.881}{0}%
\TextAt{16}{318.271}{558.674}{\Character{103}}%
\TextAt{13}{324.249}{558.674}{\Character{59}}%
\TextAt{14}{353.403}{558.674}{p}%
\NativeArrowAt{362.707}{558.674}{11.955}{19.26}{\longmapsto}%
\TextAt{13}{385.289}{558.674}{f}%
\TextAt{3}{392.336}{563.61}{\Character{0}}%
\TextAt{17}{398.922}{563.61}{1}%
\TextAt{9}{403.654}{558.674}{\Character{40}}%
\TextAt{14}{408.207}{558.674}{p}%
\TextAt{9}{414.19}{558.674}{\Character{41}}%
\TextAt{13}{418.742}{558.674}{\Character{58}}%
\TextAt{8}{79.226}{529.673}{We}%
\TextAt{8}{100.107}{529.673}{now}%
\TextAt{8}{125.231}{529.673}{identify}%
\TextAt{8}{168.564}{529.673}{in}%
\TextAt{8}{182.953}{529.673}{the}%
\TextAt{8}{203.845}{529.673}{union}%
\TextAt{13}{237.1}{529.673}{U}%
\TextAt{16}{249.457}{529.673}{\Character{91}}%
\TextAt{13}{260.586}{529.673}{V}%
\TextAt{8}{274.645}{529.673}{the}%
\TextAt{8}{295.538}{529.673}{points}%
\TextAt{8}{331.458}{529.673}{of}%
\TextAt{13}{345.524}{529.673}{V}%
\TextAt{16}{358.108}{529.673}{\Character{110}}%
\TextAt{16}{367.242}{529.673}{\Character{102}}%
\TextAt{14}{373.222}{529.673}{p}%
\NativeMathAt{379.205}{527.88}{7.97}{\infty}%
\TextAt{16}{388.171}{529.673}{\Character{103}}%
\TextAt{8}{398.783}{529.673}{with}%
\TextAt{8}{426.179}{529.673}{those}%
\TextAt{8}{457.544}{529.673}{of}%
\TextAt{13}{471.61}{529.673}{U}%
\TextAt{16}{483.968}{529.673}{\Character{110}}%
\TextAt{16}{493.102}{529.673}{\Character{102}}%
\TextAt{14}{499.081}{529.673}{p}%
\TextAt{17}{505.065}{527.88}{0}%
\TextAt{16}{509.797}{529.673}{\Character{103}}%
\TextAt{8}{520.408}{529.673}{by}%
\TextAt{8}{79.226}{512.627}{means}%
\TextAt{8}{114.811}{512.627}{of}%
\TextAt{13}{127.895}{512.627}{\Character{39}}%
\TextAt{8}{135.564}{512.627}{,}%
\TextAt{8}{142.517}{512.627}{and}%
\TextAt{8}{165.034}{512.627}{obtain}%
\TextAt{8}{201.217}{512.627}{a}%
\TextAt{8}{210.727}{512.627}{topological}%
\TextAt{8}{270.006}{512.627}{space}%
\TextAt{13}{301.036}{512.627}{X}%
\TextAt{8}{311.691}{512.627}{.}%
\TextAt{8}{320.058}{512.627}{It}%
\TextAt{8}{332.485}{512.627}{is}%
\TextAt{8}{344.004}{512.627}{immediately}%
\TextAt{8}{410.421}{512.627}{obvious}%
\TextAt{8}{452.521}{512.627}{how}%
\TextAt{8}{476.664}{512.627}{to}%
\TextAt{8}{490.729}{512.627}{obtain}%
\TextAt{8}{526.911}{512.627}{a}%
\TextAt{8}{79.226}{495.581}{sheaf}%
\TextAt{8}{108.575}{495.581}{of}%
\TextAt{8}{121.605}{495.581}{rings}%
\TextAt{14}{149.986}{495.581}{o}%
\TextAt{21}{155.83}{493.788}{X}%
\TextAt{8}{167.512}{495.581}{on}%
\TextAt{13}{183.471}{495.581}{X}%
\TextAt{8}{197.724}{495.581}{from}%
\TextAt{8}{225.064}{495.581}{the}%
\TextAt{8}{244.921}{495.581}{two}%
\TextAt{8}{266.738}{495.581}{sheaves}%
\TextAt{14}{307.862}{495.581}{o}%
\TextAt{21}{313.707}{493.788}{U}%
\TextAt{8}{324.38}{495.581}{and}%
\TextAt{14}{346.843}{495.581}{o}%
\TextAt{21}{352.687}{493.788}{V}%
\TextAt{8}{359.999}{495.581}{.}%
\TextAt{8}{368.347}{495.581}{Removing}%
\TextAt{8}{422.849}{495.581}{from}%
\TextAt{13}{450.19}{495.581}{X}%
\TextAt{8}{464.442}{495.581}{the}%
\TextAt{8}{484.299}{495.581}{point}%
\TextAt{14}{514.568}{495.581}{p}%
\NativeMathAt{520.551}{493.788}{7.97}{\infty}%
\TextAt{8}{529.518}{495.581}{,}%
\TextAt{8}{79.226}{478.535}{resp.}%
\TextAt{14}{108.548}{478.535}{p}%
\TextAt{17}{114.531}{476.742}{0}%
\TextAt{8}{119.264}{478.535}{,}%
\TextAt{8}{126.413}{478.535}{one}%
\TextAt{8}{147.872}{478.535}{gets}%
\TextAt{8}{171.998}{478.535}{canonical}%
\TextAt{8}{223.381}{478.535}{isomorphisms}%
\TextAt{9}{228.006}{451.201}{\Character{40}}%
\TextAt{13}{232.558}{451.201}{X}%
\TextAt{16}{245.87}{451.201}{\Character{110}}%
\TextAt{16}{254.502}{451.201}{\Character{102}}%
\TextAt{14}{260.482}{451.201}{p}%
\NativeMathAt{266.465}{449.407}{7.97}{\infty}%
\TextAt{16}{275.432}{451.201}{\Character{103}}%
\TextAt{13}{281.409}{451.201}{\Character{59}}%
\TextAt{14}{286.653}{451.201}{o}%
\TextAt{21}{292.498}{449.341}{X}%
\TextAt{3}{300.085}{449.341}{\Character{110}\Character{102}}%
\TextAt{27}{308.553}{449.341}{p}%
\NativeMathAt{312.542}{448.345}{5.978}{\infty}%
\TextAt{3}{320.346}{449.341}{\Character{103}}%
\TextAt{9}{325.078}{451.201}{\Character{41}}%
\NativeMathAt{332.951}{451.201}{11.955}{\cong}%
\TextAt{9}{345.57}{451.201}{\Character{40}}%
\TextAt{13}{350.123}{451.201}{U\Character{59}}%
\TextAt{14}{363.265}{451.201}{o}%
\TextAt{21}{369.11}{449.407}{U}%
\TextAt{9}{376.186}{451.201}{\Character{41}}%
\TextAt{13}{380.738}{451.201}{\Character{59}}%
\TextAt{9}{235.893}{431.166}{\Character{40}}%
\TextAt{13}{240.446}{431.166}{X}%
\TextAt{16}{253.757}{431.166}{\Character{110}}%
\TextAt{16}{262.389}{431.166}{\Character{102}}%
\TextAt{14}{268.369}{431.166}{p}%
\TextAt{17}{274.352}{429.373}{0}%
\TextAt{16}{279.085}{431.166}{\Character{103}}%
\TextAt{13}{285.062}{431.166}{\Character{59}}%
\TextAt{14}{290.306}{431.166}{o}%
\TextAt{21}{296.151}{429.306}{X}%
\TextAt{3}{303.738}{429.306}{\Character{110}\Character{102}}%
\TextAt{27}{312.206}{429.306}{p}%
\TextAt{24}{316.195}{428.199}{0}%
\TextAt{3}{320.346}{429.306}{\Character{103}}%
\TextAt{9}{325.078}{431.166}{\Character{41}}%
\NativeMathAt{332.951}{431.166}{11.955}{\cong}%
\TextAt{9}{345.57}{431.166}{\Character{40}}%
\TextAt{13}{350.123}{431.166}{V}%
\TextAt{13}{357.583}{431.166}{\Character{59}}%
\TextAt{14}{362.842}{431.166}{o}%
\TextAt{21}{368.686}{429.373}{V}%
\TextAt{9}{375.997}{431.166}{\Character{41}}%
\TextAt{13}{380.55}{431.166}{\Character{58}}%
\TextAt{8}{79.226}{404.134}{The}%
\TextAt{8}{102.34}{404.134}{pair}%
\TextAt{9}{125.465}{404.134}{\Character{40}}%
\TextAt{13}{130.017}{404.134}{X}%
\TextAt{13}{140.012}{404.134}{\Character{59}}%
\TextAt{14}{145.266}{404.134}{o}%
\TextAt{21}{151.111}{402.341}{X}%
\TextAt{9}{159.196}{404.134}{\Character{41}}%
\TextAt{8}{166.705}{404.134}{is}%
\TextAt{8}{177.529}{404.134}{the}%
\TextAt{8}{196.745}{404.134}{scheme}%
\TextAt{8}{235.854}{404.134}{corresponding}%
\TextAt{8}{310.1}{404.134}{to}%
\TextAt{8}{323.47}{404.134}{the}%
\TextAt{8}{342.686}{404.134}{field}%
\TextAt{18}{367.103}{404.134}{F}%
\TextAt{21}{374.409}{402.341}{p}%
\TextAt{9}{379.17}{404.134}{\Character{40}}%
\TextAt{13}{383.722}{404.134}{t}%
\TextAt{9}{387.949}{404.134}{\Character{41}}%
\TextAt{8}{392.501}{404.134}{.}%
\TextAt{8}{400.636}{404.134}{It}%
\TextAt{8}{412.368}{404.134}{is}%
\TextAt{8}{423.192}{404.134}{called}%
\TextAt{8}{455.415}{404.134}{the}%
\TextAt{12}{474.632}{404.134}{projective}%
\TextAt{12}{79.226}{387.088}{line}%
\TextAt{8}{104.212}{387.088}{over}%
\TextAt{18}{129.259}{387.088}{F}%
\TextAt{21}{136.564}{385.295}{p}%
\TextAt{8}{145.223}{387.088}{and}%
\TextAt{8}{167.985}{387.088}{denoted}%
\TextAt{18}{212.208}{387.088}{P}%
\TextAt{17}{219.514}{391.427}{1}%
\TextAt{25}{219.514}{384.023}{F}%
\TextAt{22}{224.384}{383.027}{p}%
\TextAt{8}{229.221}{387.088}{.}%
\NativeDiagram{%
\path[draw=black,line width=0.399bp] (168.072,-407.95) .. controls (147.143,-410.891) and (132.561,-430.242) .. (135.502,-451.171) .. controls (138.444,-472.1) and (157.795,-486.682) .. (178.724,-483.741) .. controls (199.653,-480.8) and (214.235,-461.448) .. (211.294,-440.519) .. controls (208.919,-423.615) and (195.629,-410.325) .. (178.724,-407.95);%
\path[draw=black,line width=0.399bp] (177.658,-476.161) .. controls (194.402,-473.808) and (206.067,-458.327) .. (203.714,-441.584) .. controls (201.361,-424.841) and (185.88,-413.175) .. (169.137,-415.528) .. controls (152.394,-417.881) and (140.728,-433.362) .. (143.081,-450.105) .. controls (144.981,-463.629) and (155.613,-474.261) .. (169.137,-476.161);%
\path[fill=black] (173.398,-484.113) (174.892,-484.113) .. controls (174.892,-483.288) and (174.223,-482.619) .. (173.398,-482.619) .. controls (172.573,-482.619) and (171.904,-483.288) .. (171.904,-484.113) .. controls (171.904,-484.939) and (172.573,-485.608) .. (173.398,-485.608) .. controls (174.223,-485.608) and (174.892,-484.939) .. (174.892,-484.113) -- cycle (173.398,-484.113);%
}%
\TextAt{14}{240.04}{278.022}{p}%
\TextAt{17}{246.023}{276.229}{0}%
\NativeDiagram{%
\path[fill=black] (173.398,-415.23) (174.892,-415.23) .. controls (174.892,-414.405) and (174.223,-413.736) .. (173.398,-413.736) .. controls (172.573,-413.736) and (171.904,-414.405) .. (171.904,-415.23) .. controls (171.904,-416.056) and (172.573,-416.725) .. (173.398,-416.725) .. controls (174.223,-416.725) and (174.892,-416.056) .. (174.892,-415.23) -- cycle (173.398,-415.23);%
}%
\TextAt{14}{237.923}{347.187}{p}%
\NativeMathAt{243.906}{345.394}{7.97}{\infty}%
\TextAt{13}{192.33}{324.07}{U}%
\TextAt{13}{219.818}{324.07}{V}%
\NativeDiagram{%
\path[draw=black,line width=0.399bp] (224.422,-414.663) .. controls (218.753,-423.167) and (207.414,-421.75) .. (200.726,-421.75);%
\path[draw=black,line width=0.399bp,line cap=round,line join=round] (202.599,-424.141) .. controls (202.22,-422.706) and (201.377,-422.029) .. (200.527,-421.75) .. controls (201.377,-421.471) and (202.22,-420.794) .. (202.599,-419.359);%
}%
\TextAt{13}{300.606}{357.925}{\Character{39}}%
\NativeDiagram{%
\path[draw=black,line width=0.399bp] (224.422,-445.845) -- (242.449,-445.845);%
\path[draw=black,line width=0.399bp] (224.422,-448.004) -- (224.422,-443.686);%
\path[draw=black,line width=0.399bp,line cap=round,line join=round] (240.576,-443.454) .. controls (240.955,-444.889) and (241.798,-445.566) .. (242.648,-445.845) .. controls (241.798,-446.124) and (240.955,-446.801) .. (240.576,-448.236);%
\path[draw=black,line width=0.399bp] (286.785,-445.845) (325.054,-445.845) .. controls (325.054,-424.71) and (307.921,-407.577) .. (286.785,-407.577) .. controls (265.65,-407.577) and (248.517,-424.71) .. (248.517,-445.845) .. controls (248.517,-466.98) and (265.65,-484.113) .. (286.785,-484.113) .. controls (307.921,-484.113) and (325.054,-466.98) .. (325.054,-445.845) -- cycle (286.785,-445.845);%
\path[fill=black] (286.785,-407.577) (288.28,-407.577) .. controls (288.28,-406.751) and (287.611,-406.082) .. (286.785,-406.082) .. controls (285.96,-406.082) and (285.291,-406.751) .. (285.291,-407.577) .. controls (285.291,-408.402) and (285.96,-409.071) .. (286.785,-409.071) .. controls (287.611,-409.071) and (288.28,-408.402) .. (288.28,-407.577) -- cycle (286.785,-407.577);%
\path[fill=black] (286.785,-484.113) (288.28,-484.113) .. controls (288.28,-483.288) and (287.611,-482.619) .. (286.785,-482.619) .. controls (285.96,-482.619) and (285.291,-483.288) .. (285.291,-484.113) .. controls (285.291,-484.939) and (285.96,-485.608) .. (286.785,-485.608) .. controls (287.611,-485.608) and (288.28,-484.939) .. (288.28,-484.113) -- cycle (286.785,-484.113);%
}%
\TextAt{14}{351.31}{354.558}{p}%
\NativeMathAt{357.293}{352.765}{7.97}{\infty}%
\TextAt{14}{353.427}{278.022}{p}%
\TextAt{17}{359.41}{276.229}{0}%
\TextAt{18}{402.655}{326.029}{P}%
\TextAt{17}{409.961}{330.367}{1}%
\TextAt{25}{409.961}{322.964}{F}%
\TextAt{22}{414.831}{321.968}{p}%
\TextAt{8}{79.226}{241.78}{More}%
\TextAt{8}{110.045}{241.78}{generally,}%
\TextAt{8}{162.805}{241.78}{one}%
\TextAt{8}{184.848}{241.78}{may}%
\TextAt{8}{210.802}{241.78}{similarly}%
\TextAt{8}{259.255}{241.78}{associate}%
\TextAt{8}{309.07}{241.78}{a}%
\TextAt{8}{319.41}{241.78}{scheme}%
\TextAt{9}{360.044}{241.78}{\Character{40}}%
\TextAt{13}{364.597}{241.78}{X}%
\TextAt{13}{374.591}{241.78}{\Character{59}}%
\TextAt{14}{379.846}{241.78}{o}%
\TextAt{21}{385.69}{239.987}{X}%
\TextAt{9}{393.775}{241.78}{\Character{41}}%
\TextAt{8}{402.809}{241.78}{to}%
\TextAt{8}{417.704}{241.78}{an}%
\TextAt{8}{434.547}{241.78}{arbitrary}%
\TextAt{8}{484.901}{241.78}{extension}%
\TextAt{13}{79.226}{224.734}{K}%
\TextAt{16}{93.365}{224.734}{\Character{106}}%
\TextAt{18}{100.006}{224.734}{F}%
\TextAt{21}{107.312}{222.941}{p}%
\TextAt{9}{112.073}{224.734}{\Character{40}}%
\TextAt{13}{116.626}{224.734}{t}%
\TextAt{9}{120.853}{224.734}{\Character{41}}%
\TextAt{8}{125.405}{224.734}{.}%
\TextAt{8}{133.853}{224.734}{For}%
\TextAt{8}{154.823}{224.734}{the}%
\TextAt{8}{174.979}{224.734}{precise}%
\TextAt{8}{213.403}{224.734}{description}%
\TextAt{8}{273.298}{224.734}{of}%
\TextAt{8}{286.629}{224.734}{this}%
\TextAt{8}{309.451}{224.734}{procedure}%
\TextAt{8}{363.751}{224.734}{we}%
\TextAt{8}{380.979}{224.734}{refer}%
\TextAt{8}{407.962}{224.734}{the}%
\TextAt{8}{428.118}{224.734}{reader}%
\TextAt{8}{463.888}{224.734}{to}%
\TextAt{8}{478.198}{224.734}{[51].}%
\TextAt{8}{300.14}{43.033}{90}%
\Destination{78.23}{809.815}{nav.page.90}%
\Destination{79.36}{419.155}{nav.print.97}%
\Hotlink{414.058}{707.755}{12.716}{11.047}{nav.section.1.13}%
\Hotlink{477.698}{222.555}{19.22}{11.047}{nav.cite.N051}%
\Destination{79.226}{390}{errata.184}%
}


% Chapter 2
% Page 103 (body page 91)
\Page{%
\TextAt{1}{79.226}{673.054}{Chapter}%
\TextAt{1}{188.19}{673.054}{II}%
\TextAt{1}{79.226}{614.873}{The}%
\TextAt{1}{136.112}{614.873}{Theory}%
\TextAt{1}{233.079}{614.873}{of}%
\TextAt{1}{264.856}{614.873}{Valuations}%
\TextAt{15}{79.226}{546.17}{§}%
\TextAt{15}{93.785}{546.17}{1}%
\TextAt{15}{120.693}{546.17}{The}%
\TextAt{31}{160.202}{546.17}{p}%
\TextAt{15}{168.663}{546.17}{-adic}%
\TextAt{15}{215.747}{546.17}{Numbers}%
\TextAt{8}{79.226}{514.284}{The}%
\TextAt{13}{102.666}{514.284}{p}%
\TextAt{8}{108.542}{514.284}{-adic}%
\TextAt{8}{136.537}{514.284}{numbers}%
\TextAt{8}{183.134}{514.284}{were}%
\TextAt{8}{209.504}{514.284}{invented}%
\TextAt{8}{255.718}{514.284}{at}%
\TextAt{8}{269.415}{514.284}{the}%
\TextAt{8}{288.958}{514.284}{beginning}%
\TextAt{8}{342}{514.284}{of}%
\TextAt{8}{354.716}{514.284}{the}%
\TextAt{8}{374.259}{514.284}{twentieth}%
\TextAt{8}{425.352}{514.284}{century}%
\TextAt{8}{467.012}{514.284}{by}%
\TextAt{8}{482.657}{514.284}{the}%
\TextAt{8}{502.2}{514.284}{math-}%
\TextAt{8}{79.226}{497.238}{ematician}%
\TextAt{26}{132.895}{497.238}{Kurt}%
\TextAt{26}{168.674}{497.238}{Hensel}%
\TextAt{8}{214.919}{497.238}{(1861–1941)}%
\TextAt{8}{280.986}{497.238}{with}%
\TextAt{8}{307.983}{497.238}{a}%
\TextAt{8}{318.076}{497.238}{view}%
\TextAt{8}{345.396}{497.238}{to}%
\TextAt{8}{360.043}{497.238}{introduce}%
\TextAt{8}{412.409}{497.238}{into}%
\TextAt{8}{436.489}{497.238}{number}%
\TextAt{8}{479.423}{497.238}{theory}%
\TextAt{8}{516.51}{497.238}{the}%
\TextAt{8}{79.226}{480.192}{powerful}%
\TextAt{8}{126.696}{480.192}{method}%
\TextAt{8}{169.289}{480.192}{of}%
\TextAt{8}{182.615}{480.192}{power}%
\TextAt{8}{216.756}{480.192}{series}%
\TextAt{8}{248.087}{480.192}{expansion}%
\TextAt{8}{302.455}{480.192}{which}%
\TextAt{8}{335.938}{480.192}{plays}%
\TextAt{8}{365.918}{480.192}{such}%
\TextAt{8}{392.311}{480.192}{a}%
\TextAt{8}{402.063}{480.192}{predominant}%
\TextAt{8}{470.682}{480.192}{role}%
\TextAt{8}{493.441}{480.192}{in}%
\TextAt{8}{507.09}{480.192}{func-}%
\TextAt{8}{79.226}{463.147}{tion}%
\TextAt{8}{103.563}{463.147}{theory.}%
\TextAt{8}{144.696}{463.147}{The}%
\TextAt{8}{169.02}{463.147}{idea}%
\TextAt{8}{194.002}{463.147}{originated}%
\TextAt{8}{249.566}{463.147}{from}%
\TextAt{8}{277.477}{463.147}{the}%
\TextAt{8}{297.904}{463.147}{observation}%
\TextAt{8}{360.365}{463.147}{made}%
\TextAt{8}{391.851}{463.147}{in}%
\TextAt{8}{405.775}{463.147}{the}%
\TextAt{8}{426.202}{463.147}{last}%
\TextAt{8}{448.649}{463.147}{chapter}%
\TextAt{8}{490.871}{463.147}{that}%
\TextAt{8}{516.51}{463.147}{the}%
\TextAt{8}{79.226}{446.101}{numbers}%
\TextAt{13}{126.524}{446.101}{f}%
\NativeMathAt{137.039}{446.101}{11.955}{\in}%
\TextAt{18}{148.478}{446.101}{Z}%
\TextAt{8}{160.432}{446.101}{may}%
\TextAt{8}{185.887}{446.101}{be}%
\TextAt{8}{201.899}{446.101}{viewed}%
\TextAt{8}{240.349}{446.101}{in}%
\TextAt{8}{254.089}{446.101}{analogy}%
\TextAt{8}{297.441}{446.101}{with}%
\TextAt{8}{324.188}{446.101}{the}%
\TextAt{8}{344.431}{446.101}{polynomials}%
\TextAt{13}{409.626}{446.101}{f}%
\TextAt{9}{416.672}{446.101}{\Character{40}}%
\TextAt{13}{421.225}{446.101}{z}%
\TextAt{9}{427.195}{446.101}{\Character{41}}%
\NativeMathAt{435.216}{446.101}{11.955}{\in}%
\TextAt{18}{446.655}{446.101}{C}%
\TextAt{9}{455.289}{446.101}{\Character{91}}%
\TextAt{13}{458.541}{446.101}{z}%
\TextAt{9}{464.512}{446.101}{\Character{93}}%
\TextAt{8}{471.747}{446.101}{as}%
\TextAt{8}{486.204}{446.101}{functions}%
\TextAt{8}{79.226}{429.055}{on}%
\TextAt{8}{96.381}{429.055}{the}%
\TextAt{8}{117.434}{429.055}{space}%
\TextAt{13}{149.605}{429.055}{X}%
\TextAt{8}{165.053}{429.055}{of}%
\TextAt{8}{179.28}{429.055}{prime}%
\TextAt{8}{213.339}{429.055}{numbers}%
\TextAt{8}{261.446}{429.055}{in}%
\TextAt{18}{275.995}{429.055}{Z}%
\TextAt{8}{283.966}{429.055}{,}%
\TextAt{8}{292.235}{429.055}{associating}%
\TextAt{8}{352.775}{429.055}{to}%
\TextAt{8}{367.982}{429.055}{them}%
\TextAt{8}{398.79}{429.055}{their}%
\TextAt{8}{427.649}{429.055}{“value”}%
\TextAt{8}{469.838}{429.055}{at}%
\TextAt{8}{485.045}{429.055}{the}%
\TextAt{8}{506.098}{429.055}{point}%
\TextAt{13}{79.226}{412.009}{p}%
\NativeMathAt{88.422}{412.009}{11.955}{\in}%
\TextAt{13}{99.713}{412.009}{X}%
\TextAt{8}{110.368}{412.009}{,}%
\TextAt{8}{117.517}{412.009}{i.e.,}%
\TextAt{8}{139.623}{412.009}{the}%
\TextAt{8}{159.779}{412.009}{element}%
\TextAt{13}{264.646}{394.963}{f}%
\TextAt{9}{271.692}{394.963}{\Character{40}}%
\TextAt{13}{276.244}{394.963}{p}%
\TextAt{9}{282.119}{394.963}{\Character{41}}%
\TextAt{9}{289.974}{394.963}{\Character{58}\Character{61}}%
\TextAt{13}{305.67}{394.963}{f}%
\TextAt{8}{316.037}{394.963}{mod}%
\TextAt{13}{341.475}{394.963}{p}%
\TextAt{8}{79.226}{373.153}{in}%
\TextAt{8}{93.145}{373.153}{the}%
\TextAt{8}{113.567}{373.153}{residue}%
\TextAt{8}{153.56}{373.153}{class}%
\TextAt{8}{181.263}{373.153}{field}%
\TextAt{13}{206.885}{373.153}{\Character{20}}%
\TextAt{9}{213.624}{373.153}{\Character{40}}%
\TextAt{13}{218.176}{373.153}{p}%
\TextAt{9}{224.051}{373.153}{\Character{41}}%
\TextAt{9}{232.36}{373.153}{\Character{61}}%
\TextAt{18}{245.256}{373.153}{Z}%
\TextAt{9}{253.226}{373.153}{\Character{47}}%
\TextAt{13}{259.079}{373.153}{p}%
\TextAt{18}{264.954}{373.153}{Z}%
\TextAt{8}{272.924}{373.153}{.}%
\TextAt{8}{282.17}{373.153}{This}%
\TextAt{8}{309.155}{373.153}{point}%
\TextAt{8}{339.991}{373.153}{of}%
\TextAt{8}{353.586}{373.153}{view}%
\TextAt{8}{380.835}{373.153}{suggests}%
\TextAt{8}{426.817}{373.153}{the}%
\TextAt{8}{447.239}{373.153}{further}%
\TextAt{8}{486.85}{373.153}{question:}%
\TextAt{8}{79.226}{356.107}{whether}%
\TextAt{8}{124.76}{356.107}{not}%
\TextAt{8}{146.24}{356.107}{only}%
\TextAt{8}{172.597}{356.107}{the}%
\TextAt{8}{193.42}{356.107}{“value”}%
\TextAt{8}{235.379}{356.107}{of}%
\TextAt{8}{249.375}{356.107}{the}%
\TextAt{8}{270.197}{356.107}{integer}%
\TextAt{13}{309.562}{356.107}{f}%
\NativeMathAt{321.064}{356.107}{11.955}{\in}%
\TextAt{18}{333.49}{356.107}{Z}%
\TextAt{8}{346.023}{356.107}{at}%
\TextAt{13}{360.999}{356.107}{p}%
\TextAt{8}{366.874}{356.107}{,}%
\TextAt{8}{374.856}{356.107}{but}%
\TextAt{8}{396.981}{356.107}{also}%
\TextAt{8}{421.127}{356.107}{the}%
\TextAt{8}{441.95}{356.107}{higher}%
\TextAt{8}{478.385}{356.107}{derivatives}%
\TextAt{8}{79.226}{339.061}{of}%
\TextAt{13}{93.247}{339.061}{f}%
\TextAt{8}{104.881}{339.061}{can}%
\TextAt{8}{127.032}{339.061}{be}%
\TextAt{8}{143.647}{339.061}{reasonably}%
\TextAt{8}{202.619}{339.061}{defined.}%
\TextAt{8}{249.554}{339.061}{In}%
\TextAt{8}{264.866}{339.061}{the}%
\TextAt{8}{285.713}{339.061}{case}%
\TextAt{8}{311.175}{339.061}{of}%
\TextAt{8}{325.196}{339.061}{the}%
\TextAt{8}{346.043}{339.061}{polynomials}%
\TextAt{13}{411.842}{339.061}{f}%
\TextAt{9}{418.888}{339.061}{\Character{40}}%
\TextAt{13}{423.44}{339.061}{z}%
\TextAt{9}{429.411}{339.061}{\Character{41}}%
\NativeMathAt{438.461}{339.061}{11.955}{\in}%
\TextAt{18}{450.929}{339.061}{C}%
\TextAt{9}{459.563}{339.061}{\Character{91}}%
\TextAt{13}{462.815}{339.061}{z}%
\TextAt{9}{468.786}{339.061}{\Character{93}}%
\TextAt{8}{472.037}{339.061}{,}%
\TextAt{8}{480.05}{339.061}{the}%
\TextAt{8}{500.897}{339.061}{higher}%
\TextAt{8}{79.226}{322.015}{derivatives}%
\TextAt{8}{137.508}{322.015}{at}%
\TextAt{8}{151.818}{322.015}{the}%
\TextAt{8}{171.974}{322.015}{point}%
\TextAt{13}{202.544}{322.015}{z}%
\TextAt{9}{211.835}{322.015}{\Character{61}}%
\TextAt{13}{224.261}{322.015}{a}%
\TextAt{8}{234.303}{322.015}{are}%
\TextAt{8}{253.814}{322.015}{given}%
\TextAt{8}{284.383}{322.015}{by}%
\TextAt{8}{300.642}{322.015}{the}%
\TextAt{8}{320.799}{322.015}{coefficients}%
\TextAt{8}{380.036}{322.015}{of}%
\TextAt{8}{393.367}{322.015}{the}%
\TextAt{8}{413.523}{322.015}{expansion}%
\TextAt{13}{201.419}{297.008}{f}%
\TextAt{9}{208.466}{297.008}{\Character{40}}%
\TextAt{13}{213.018}{297.008}{z}%
\TextAt{9}{218.988}{297.008}{\Character{41}}%
\TextAt{9}{226.843}{297.008}{\Character{61}}%
\TextAt{13}{239.287}{297.008}{a}%
\TextAt{17}{245.432}{295.214}{0}%
\TextAt{9}{252.821}{297.008}{\Character{43}}%
\TextAt{13}{264.582}{297.008}{a}%
\TextAt{17}{270.727}{295.214}{1}%
\TextAt{9}{275.46}{297.008}{\Character{40}}%
\TextAt{13}{280.012}{297.008}{z}%
\TextAt{16}{288.639}{297.008}{\Character{0}}%
\TextAt{13}{300.594}{297.008}{a}%
\TextAt{9}{306.739}{297.008}{\Character{41}}%
\TextAt{9}{313.936}{297.008}{\Character{43}}%
\TextAt{16}{325.71}{297.008}{\Character{1}}%
\TextAt{16}{331.006}{297.008}{\Character{1}}%
\TextAt{16}{336.314}{297.008}{\Character{1}}%
\TextAt{9}{342.314}{297.008}{\Character{43}}%
\TextAt{13}{354.075}{297.008}{a}%
\TextAt{21}{360.22}{295.214}{n}%
\TextAt{9}{365.857}{297.008}{\Character{40}}%
\TextAt{13}{370.409}{297.008}{z}%
\TextAt{16}{379.036}{297.008}{\Character{0}}%
\TextAt{13}{390.991}{297.008}{a}%
\TextAt{9}{397.136}{297.008}{\Character{41}}%
\TextAt{21}{401.689}{301.944}{n}%
\TextAt{13}{407.325}{297.008}{\Character{59}}%
\TextAt{8}{79.226}{272.001}{and}%
\TextAt{8}{102.303}{272.001}{more}%
\TextAt{8}{131.883}{272.001}{generally,}%
\TextAt{8}{184.304}{272.001}{for}%
\TextAt{8}{202.502}{272.001}{rational}%
\TextAt{8}{246.405}{272.001}{functions}%
\TextAt{13}{297.181}{272.001}{f}%
\TextAt{9}{304.228}{272.001}{\Character{40}}%
\TextAt{13}{308.78}{272.001}{z}%
\TextAt{9}{314.75}{272.001}{\Character{41}}%
\TextAt{9}{323.143}{272.001}{\Character{61}}%
\TextAt{21}{337.596}{277.699}{g}%
\TextAt{17}{341.912}{277.699}{\Character{40}}%
\TextAt{21}{345.205}{277.699}{z}%
\TextAt{17}{349.481}{277.699}{\Character{41}}%
\Vector{q 0 G 0 g 0.478 w 265.314 -499.011 m 281.056 -499.011 l S Q}%
\TextAt{21}{337.314}{267.878}{h}%
\TextAt{17}{342.194}{267.878}{\Character{40}}%
\TextAt{21}{345.487}{267.878}{z}%
\TextAt{17}{349.763}{267.878}{\Character{41}}%
\NativeMathAt{358.107}{272.001}{11.955}{\in}%
\TextAt{18}{369.933}{272.001}{C}%
\TextAt{9}{378.567}{272.001}{\Character{40}}%
\TextAt{13}{383.119}{272.001}{z}%
\TextAt{9}{389.09}{272.001}{\Character{41}}%
\TextAt{8}{393.642}{272.001}{,}%
\TextAt{8}{401.184}{272.001}{with}%
\TextAt{13}{428.157}{272.001}{g}%
\TextAt{13}{434.171}{272.001}{\Character{59}}%
\TextAt{13}{439.407}{272.001}{h}%
\NativeMathAt{450.03}{272.001}{11.955}{\in}%
\TextAt{18}{461.855}{272.001}{C}%
\TextAt{9}{470.49}{272.001}{\Character{91}}%
\TextAt{13}{473.741}{272.001}{z}%
\TextAt{9}{479.712}{272.001}{\Character{93}}%
\TextAt{8}{482.964}{272.001}{,}%
\TextAt{8}{490.505}{272.001}{they}%
\TextAt{8}{517.156}{272.001}{are}%
\TextAt{8}{79.226}{254.955}{defined}%
\TextAt{8}{119.539}{254.955}{by}%
\TextAt{8}{135.798}{254.955}{the}%
\TextAt{8}{155.954}{254.955}{Taylor}%
\TextAt{8}{192.705}{254.955}{expansion}%
\TextAt{13}{249.637}{223.987}{f}%
\TextAt{9}{256.683}{223.987}{\Character{40}}%
\TextAt{13}{261.236}{223.987}{z}%
\TextAt{9}{267.206}{223.987}{\Character{41}}%
\TextAt{9}{275.061}{223.987}{\Character{61}}%
\NativeMathAt{291.905}{238.931}{7.97}{\infty}%
\TextAt{20}{287.505}{235.345}{\Character{88}}%
\TextAt{21}{288.387}{210.239}{\Character{23}}%
\TextAt{17}{293.071}{210.239}{\Character{61}0}%
\TextAt{13}{306.766}{223.987}{a}%
\TextAt{21}{312.911}{222.194}{\Character{23}}%
\TextAt{9}{318.093}{223.987}{\Character{40}}%
\TextAt{13}{322.645}{223.987}{z}%
\TextAt{16}{331.273}{223.987}{\Character{0}}%
\TextAt{13}{343.228}{223.987}{a}%
\TextAt{9}{349.373}{223.987}{\Character{41}}%
\TextAt{21}{353.925}{228.923}{\Character{23}}%
\TextAt{13}{359.107}{223.987}{\Character{59}}%
\TextAt{8}{79.226}{191.119}{provided}%
\TextAt{8}{127.717}{191.119}{there}%
\TextAt{8}{157.988}{191.119}{is}%
\TextAt{8}{170.111}{191.119}{no}%
\TextAt{8}{186.729}{191.119}{pole}%
\TextAt{8}{212.123}{191.119}{at}%
\TextAt{13}{226.792}{191.119}{z}%
\TextAt{9}{236.696}{191.119}{\Character{61}}%
\TextAt{13}{249.733}{191.119}{a}%
\TextAt{8}{255.878}{191.119}{,}%
\TextAt{8}{263.477}{191.119}{i.e.,}%
\TextAt{8}{286.031}{191.119}{as}%
\TextAt{8}{300.76}{191.119}{long}%
\TextAt{8}{326.488}{191.119}{as}%
\TextAt{9}{341.217}{191.119}{\Character{40}}%
\TextAt{13}{345.77}{191.119}{z}%
\TextAt{16}{354.642}{191.119}{\Character{0}}%
\TextAt{13}{366.842}{191.119}{a}%
\TextAt{9}{372.987}{191.119}{\Character{41}}%
\TextAt{18}{381.472}{191.119}{\Character{45}}%
\TextAt{13}{388.726}{191.119}{h}%
\TextAt{9}{395.465}{191.119}{\Character{40}}%
\TextAt{13}{400.017}{191.119}{z}%
\TextAt{9}{405.988}{191.119}{\Character{41}}%
\TextAt{8}{410.54}{191.119}{.}%
\TextAt{8}{420.066}{191.119}{The}%
\TextAt{8}{444.479}{191.119}{fact}%
\TextAt{8}{467.924}{191.119}{that}%
\TextAt{8}{493.652}{191.119}{such}%
\TextAt{8}{520.408}{191.119}{an}%
\TextAt{8}{79.226}{174.073}{expansion}%
\TextAt{8}{133.81}{174.073}{can}%
\TextAt{8}{155.482}{174.073}{also}%
\TextAt{8}{179.174}{174.073}{be}%
\TextAt{8}{195.311}{174.073}{written}%
\TextAt{8}{236.494}{174.073}{down,}%
\TextAt{8}{270.903}{174.073}{relative}%
\TextAt{8}{312.743}{174.073}{to}%
\TextAt{8}{327.265}{174.073}{a}%
\TextAt{8}{337.233}{174.073}{prime}%
\TextAt{8}{370.609}{174.073}{number}%
\TextAt{13}{413.417}{174.073}{p}%
\TextAt{8}{423.402}{174.073}{in}%
\TextAt{18}{437.267}{174.073}{Z}%
\TextAt{8}{445.237}{174.073}{,}%
\TextAt{8}{452.652}{174.073}{for}%
\TextAt{8}{470.749}{174.073}{any}%
\TextAt{8}{493.078}{174.073}{rational}%
\TextAt{8}{79.226}{157.028}{number}%
\TextAt{13}{121.822}{157.028}{f}%
\NativeMathAt{132.19}{157.028}{11.955}{\in}%
\TextAt{18}{143.481}{157.028}{Q}%
\TextAt{8}{156.677}{157.028}{as}%
\TextAt{8}{171.047}{157.028}{long}%
\TextAt{8}{196.416}{157.028}{as}%
\TextAt{8}{210.786}{157.028}{it}%
\TextAt{8}{222.49}{157.028}{lies}%
\TextAt{8}{242.706}{157.028}{in}%
\TextAt{8}{256.359}{157.028}{the}%
\TextAt{8}{276.515}{157.028}{local}%
\TextAt{8}{304.156}{157.028}{ring}%
\TextAt{18}{231.837}{132.02}{Z}%
\TextAt{17}{239.807}{130.161}{\Character{40}}%
\TextAt{21}{243.101}{130.161}{p}%
\TextAt{17}{247.363}{130.161}{\Character{41}}%
\TextAt{9}{254.476}{132.02}{\Character{61}}%
\TextAt{20}{266.901}{145.291}{\Character{110}}%
\TextAt{13}{276.419}{140.108}{g}%
\Vector{q 0 G 0 g 0.478 w 204.067 -638.991 m 210.805 -638.991 l S Q}%
\TextAt{13}{276.067}{123.82}{h}%
\TextAt{20}{287.322}{145.769}{\Character{12}}%
\TextAt{20}{287.322}{138.596}{\Character{12}}%
\TextAt{20}{287.322}{131.423}{\Character{12}}%
\TextAt{13}{294.627}{132.02}{g}%
\TextAt{13}{300.641}{132.02}{\Character{59}}%
\TextAt{13}{305.877}{132.02}{h}%
\NativeMathAt{315.965}{132.02}{11.955}{\in}%
\TextAt{18}{327.256}{132.02}{Z}%
\TextAt{13}{335.226}{132.02}{\Character{59}}%
\TextAt{13}{344.348}{132.02}{p}%
\TextAt{18}{353.564}{132.02}{\Character{45}}%
\TextAt{13}{360.206}{132.02}{h}%
\TextAt{20}{366.944}{145.291}{\Character{111}}%
\TextAt{13}{376.907}{132.02}{\Character{59}}%
\TextAt{8}{79.226}{107.013}{leads}%
\TextAt{8}{109.579}{107.013}{us}%
\TextAt{8}{125.622}{107.013}{to}%
\TextAt{8}{140.96}{107.013}{the}%
\TextAt{8}{162.143}{107.013}{notion}%
\TextAt{8}{199.598}{107.013}{of}%
\TextAt{13}{213.955}{107.013}{p}%
\TextAt{8}{219.83}{107.013}{-adic}%
\TextAt{8}{249.466}{107.013}{number.}%
\TextAt{8}{299.695}{107.013}{First,}%
\TextAt{8}{332.744}{107.013}{every}%
\TextAt{8}{364.663}{107.013}{positive}%
\TextAt{8}{409.004}{107.013}{integer}%
\TextAt{13}{448.73}{107.013}{f}%
\NativeMathAt{460.848}{107.013}{11.955}{\in}%
\TextAt{18}{473.889}{107.013}{N}%
\TextAt{8}{487.448}{107.013}{admits}%
\TextAt{8}{526.911}{107.013}{a}%
\TextAt{13}{79.226}{89.967}{p}%
\TextAt{12}{85.101}{89.967}{-adic}%
\TextAt{12}{117.799}{89.967}{expansion}%
\TextAt{13}{238.914}{72.921}{f}%
\TextAt{9}{249.281}{72.921}{\Character{61}}%
\TextAt{13}{261.706}{72.921}{a}%
\TextAt{17}{267.851}{71.128}{0}%
\TextAt{9}{275.24}{72.921}{\Character{43}}%
\TextAt{13}{287.002}{72.921}{a}%
\TextAt{17}{293.147}{71.128}{1}%
\TextAt{13}{297.879}{72.921}{p}%
\TextAt{9}{306.411}{72.921}{\Character{43}}%
\TextAt{16}{318.172}{72.921}{\Character{1}}%
\TextAt{16}{323.468}{72.921}{\Character{1}}%
\TextAt{16}{328.776}{72.921}{\Character{1}}%
\TextAt{9}{334.776}{72.921}{\Character{43}}%
\TextAt{13}{346.538}{72.921}{a}%
\TextAt{21}{352.683}{71.128}{n}%
\TextAt{13}{358.319}{72.921}{p}%
\TextAt{21}{364.194}{77.858}{n}%
\TextAt{13}{369.83}{72.921}{\Character{59}}%
\TextAt{8}{300.14}{43.033}{91}%
\Destination{79.23}{773.075}{nav.chapter.2}%
\Destination{78.23}{809.815}{nav.page.91}%
\Destination{79.23}{570.065}{nav.section.2.1}%
\Destination{79.36}{629.855}{nav.print.99}%
\Destination{79.23}{773.075}{outline.20}%
\Destination{79.23}{570.065}{outline.21}%
}

% Page 104 (body page 92)
\Page{%
\TextAt{2}{79.226}{796.943}{II.}%
\TextAt{2}{94.073}{796.943}{The}%
\TextAt{2}{116.492}{796.943}{Theory}%
\TextAt{2}{154.401}{796.943}{of}%
\TextAt{2}{166.826}{796.943}{Valuations}%
\Vector{q 0 G 0 g 0.299 w 7.226 18.729 m 460.769 18.729 l S Q}%
\TextAt{8}{79.226}{761.124}{with}%
\TextAt{8}{105.08}{761.124}{coefficients}%
\TextAt{13}{163.511}{761.124}{a}%
\TextAt{21}{169.656}{759.33}{i}%
\TextAt{8}{176.129}{761.124}{in}%
\TextAt{16}{188.975}{761.124}{\Character{102}}%
\TextAt{9}{194.953}{761.124}{0}%
\TextAt{13}{200.806}{761.124}{\Character{59}}%
\TextAt{9}{206.05}{761.124}{1}%
\TextAt{13}{211.903}{761.124}{\Character{59}}%
\TextAt{13}{217.127}{761.124}{\Character{58}}%
\TextAt{13}{222.364}{761.124}{\Character{58}}%
\TextAt{13}{227.588}{761.124}{\Character{58}}%
\TextAt{13}{232.825}{761.124}{\Character{59}}%
\TextAt{13}{238.061}{761.124}{p}%
\TextAt{16}{245.007}{761.124}{\Character{0}}%
\TextAt{9}{255.313}{761.124}{1}%
\TextAt{16}{261.166}{761.124}{\Character{103}}%
\TextAt{8}{267.143}{761.124}{,}%
\TextAt{8}{273.648}{761.124}{i.e.,}%
\TextAt{8}{295.108}{761.124}{in}%
\TextAt{8}{307.954}{761.124}{a}%
\TextAt{8}{316.903}{761.124}{fixed}%
\TextAt{8}{344.383}{761.124}{system}%
\TextAt{8}{382.395}{761.124}{of}%
\TextAt{8}{394.919}{761.124}{representatives}%
\TextAt{8}{473.913}{761.124}{of}%
\TextAt{8}{486.437}{761.124}{the}%
\TextAt{8}{505.787}{761.124}{“field}%
\TextAt{8}{79.226}{744.078}{of}%
\TextAt{8}{92.834}{744.078}{values”}%
\TextAt{13}{133.496}{744.078}{\Character{20}}%
\TextAt{9}{140.234}{744.078}{\Character{40}}%
\TextAt{13}{144.787}{744.078}{p}%
\TextAt{9}{150.662}{744.078}{\Character{41}}%
\TextAt{9}{158.995}{744.078}{\Character{61}}%
\TextAt{18}{171.906}{744.078}{F}%
\TextAt{21}{179.212}{742.284}{p}%
\TextAt{8}{183.973}{744.078}{.}%
\TextAt{8}{193.254}{744.078}{This}%
\TextAt{8}{220.252}{744.078}{representation}%
\TextAt{8}{297.018}{744.078}{is}%
\TextAt{8}{309.06}{744.078}{clearly}%
\TextAt{8}{346.733}{744.078}{unique.}%
\TextAt{8}{390.158}{744.078}{It}%
\TextAt{8}{403.108}{744.078}{is}%
\TextAt{8}{415.149}{744.078}{computed}%
\TextAt{8}{469.405}{744.078}{explicitly}%
\TextAt{8}{520.408}{744.078}{by}%
\TextAt{8}{79.226}{727.032}{successively}%
\TextAt{8}{142.816}{727.032}{dividing}%
\TextAt{8}{188.018}{727.032}{by}%
\TextAt{13}{204.277}{727.032}{p}%
\TextAt{8}{210.152}{727.032}{,}%
\TextAt{8}{217.302}{727.032}{forming}%
\TextAt{8}{260.555}{727.032}{the}%
\TextAt{8}{280.712}{727.032}{following}%
\TextAt{8}{330.146}{727.032}{system}%
\TextAt{8}{368.965}{727.032}{of}%
\TextAt{8}{382.295}{727.032}{equations:}%
\TextAt{13}{273.415}{699.697}{f}%
\TextAt{9}{283.782}{699.697}{\Character{61}}%
\TextAt{13}{296.208}{699.697}{a}%
\TextAt{17}{302.353}{697.904}{0}%
\TextAt{9}{309.742}{699.697}{\Character{43}}%
\TextAt{13}{321.503}{699.697}{pf}%
\TextAt{17}{333.148}{697.904}{1}%
\TextAt{13}{337.881}{699.697}{\Character{59}}%
\TextAt{13}{269.959}{679.663}{f}%
\TextAt{17}{275.729}{677.869}{1}%
\TextAt{9}{283.782}{679.663}{\Character{61}}%
\TextAt{13}{296.208}{679.663}{a}%
\TextAt{17}{302.353}{677.869}{1}%
\TextAt{9}{309.742}{679.663}{\Character{43}}%
\TextAt{13}{321.503}{679.663}{pf}%
\TextAt{17}{333.148}{677.869}{2}%
\TextAt{13}{337.881}{679.663}{\Character{59}}%
\TextAt{8}{284.359}{663.426}{.}%
\TextAt{8}{284.359}{659.441}{.}%
\TextAt{8}{284.359}{655.456}{.}%
\TextAt{13}{258.235}{635.422}{f}%
\TextAt{21}{264.005}{633.628}{n}%
\TextAt{3}{269.143}{633.628}{\Character{0}}%
\TextAt{17}{275.729}{633.628}{1}%
\TextAt{9}{283.782}{635.422}{\Character{61}}%
\TextAt{13}{296.208}{635.422}{a}%
\TextAt{21}{302.353}{633.628}{n}%
\TextAt{3}{307.491}{633.628}{\Character{0}}%
\TextAt{17}{314.078}{633.628}{1}%
\TextAt{9}{321.467}{635.422}{\Character{43}}%
\TextAt{13}{333.228}{635.422}{pf}%
\TextAt{21}{344.873}{633.628}{n}%
\TextAt{13}{350.509}{635.422}{\Character{59}}%
\TextAt{13}{269.055}{615.387}{f}%
\TextAt{21}{274.825}{613.594}{n}%
\TextAt{9}{283.782}{615.387}{\Character{61}}%
\TextAt{13}{296.208}{615.387}{a}%
\TextAt{21}{302.353}{613.594}{n}%
\TextAt{13}{307.989}{615.387}{\Character{58}}%
\TextAt{8}{79.226}{588.355}{Here}%
\TextAt{13}{107.625}{588.355}{a}%
\TextAt{21}{113.77}{586.562}{i}%
\NativeMathAt{121.784}{588.355}{11.955}{\in}%
\TextAt{16}{134.373}{588.355}{\Character{102}}%
\TextAt{9}{140.365}{588.355}{0}%
\TextAt{13}{146.218}{588.355}{\Character{59}}%
\TextAt{9}{151.462}{588.355}{1}%
\TextAt{13}{157.315}{588.355}{\Character{59}}%
\TextAt{13}{162.539}{588.355}{\Character{58}}%
\TextAt{13}{167.776}{588.355}{\Character{58}}%
\TextAt{13}{173}{588.355}{\Character{58}}%
\TextAt{13}{178.237}{588.355}{\Character{59}}%
\TextAt{13}{183.473}{588.355}{p}%
\TextAt{16}{192.592}{588.355}{\Character{0}}%
\TextAt{9}{205.072}{588.355}{1}%
\TextAt{16}{210.925}{588.355}{\Character{103}}%
\TextAt{8}{221.571}{588.355}{denotes}%
\TextAt{8}{264.674}{588.355}{the}%
\TextAt{8}{285.6}{588.355}{representative}%
\TextAt{8}{361.556}{588.355}{of}%
\TextAt{13}{375.656}{588.355}{f}%
\TextAt{21}{381.426}{586.562}{i}%
\TextAt{8}{388.129}{588.355}{mod}%
\TextAt{13}{413.566}{588.355}{p}%
\NativeMathAt{424.075}{588.355}{11.955}{\in}%
\TextAt{18}{436.678}{588.355}{Z}%
\TextAt{9}{444.648}{588.355}{\Character{47}}%
\TextAt{13}{450.501}{588.355}{p}%
\TextAt{18}{456.376}{588.355}{Z}%
\TextAt{8}{464.346}{588.355}{.}%
\TextAt{8}{475.105}{588.355}{In}%
\TextAt{8}{490.496}{588.355}{concrete}%
\TextAt{8}{79.226}{571.309}{cases,}%
\TextAt{8}{111.734}{571.309}{one}%
\TextAt{8}{133.032}{571.309}{sometimes}%
\TextAt{8}{189.574}{571.309}{writes}%
\TextAt{8}{223.939}{571.309}{the}%
\TextAt{8}{243.934}{571.309}{number}%
\TextAt{13}{286.369}{571.309}{f}%
\TextAt{8}{297.151}{571.309}{simply}%
\TextAt{8}{334.445}{571.309}{as}%
\TextAt{8}{348.653}{571.309}{the}%
\TextAt{8}{368.648}{571.309}{sequence}%
\TextAt{8}{416.989}{571.309}{of}%
\TextAt{8}{430.157}{571.309}{digits}%
\TextAt{13}{461.928}{571.309}{a}%
\TextAt{17}{468.073}{569.516}{0}%
\TextAt{13}{472.805}{571.309}{\Character{59}}%
\TextAt{13}{478.029}{571.309}{a}%
\TextAt{17}{484.194}{569.516}{1}%
\TextAt{13}{488.926}{571.309}{a}%
\TextAt{17}{495.071}{569.516}{2}%
\TextAt{16}{501.796}{571.309}{\Character{1}}%
\TextAt{16}{507.092}{571.309}{\Character{1}}%
\TextAt{16}{512.4}{571.309}{\Character{1}}%
\TextAt{13}{517.736}{571.309}{a}%
\TextAt{21}{523.881}{569.516}{n}%
\TextAt{8}{529.518}{571.309}{,}%
\TextAt{8}{79.226}{554.263}{for}%
\TextAt{8}{97.111}{554.263}{instance}%
\TextAt{9}{236.864}{541.215}{216}%
\TextAt{9}{257.714}{541.215}{\Character{61}}%
\TextAt{9}{270.135}{541.215}{0}%
\TextAt{13}{276.023}{541.215}{\Character{59}}%
\TextAt{9}{281.267}{541.215}{0011011}%
\TextAt{8}{332.2}{541.215}{(2-adic)}%
\TextAt{13}{371.88}{541.215}{\Character{59}}%
\TextAt{9}{236.864}{521.181}{216}%
\TextAt{9}{257.714}{521.181}{\Character{61}}%
\TextAt{9}{270.135}{521.181}{0}%
\TextAt{13}{276.023}{521.181}{\Character{59}}%
\TextAt{9}{281.267}{521.181}{0022}%
\TextAt{8}{332.2}{521.181}{(3-adic)}%
\TextAt{13}{371.88}{521.181}{\Character{59}}%
\TextAt{9}{236.864}{501.146}{216}%
\TextAt{9}{257.714}{501.146}{\Character{61}}%
\TextAt{9}{270.135}{501.146}{1}%
\TextAt{13}{276.023}{501.146}{\Character{59}}%
\TextAt{9}{281.267}{501.146}{331}%
\TextAt{8}{332.2}{501.146}{(5-adic)}%
\TextAt{13}{371.88}{501.146}{\Character{58}}%
\TextAt{8}{79.226}{480.024}{As}%
\TextAt{8}{96.854}{480.024}{soon}%
\TextAt{8}{124.249}{480.024}{as}%
\TextAt{8}{138.96}{480.024}{one}%
\TextAt{8}{160.76}{480.024}{tries}%
\TextAt{8}{187.175}{480.024}{to}%
\TextAt{8}{201.826}{480.024}{write}%
\TextAt{8}{232.079}{480.024}{down}%
\TextAt{8}{263.312}{480.024}{such}%
\TextAt{13}{290.049}{480.024}{p}%
\TextAt{8}{295.924}{480.024}{-adic}%
\TextAt{8}{324.874}{480.024}{expansions}%
\TextAt{8}{384.201}{480.024}{also}%
\TextAt{8}{408.022}{480.024}{for}%
\TextAt{8}{426.248}{480.024}{negative}%
\TextAt{8}{472.771}{480.024}{integers,}%
\TextAt{8}{519.762}{480.024}{let}%
\TextAt{8}{79.226}{462.978}{alone}%
\TextAt{8}{109.796}{462.978}{for}%
\TextAt{8}{127.68}{462.978}{fractions,}%
\TextAt{8}{178.801}{462.978}{one}%
\TextAt{8}{200.26}{462.978}{is}%
\TextAt{8}{212.024}{462.978}{forced}%
\TextAt{8}{246.814}{462.978}{to}%
\TextAt{8}{261.124}{462.978}{allow}%
\TextAt{8}{291.371}{462.978}{infinite}%
\TextAt{8}{331.038}{462.978}{series}%
\NativeMathAt{224.926}{443.092}{7.97}{\infty}%
\TextAt{20}{220.526}{439.506}{\Character{88}}%
\TextAt{21}{221.408}{414.4}{\Character{23}}%
\TextAt{17}{226.092}{414.4}{\Character{61}0}%
\TextAt{13}{239.787}{428.148}{a}%
\TextAt{21}{245.932}{426.355}{\Character{23}}%
\TextAt{13}{251.114}{428.148}{p}%
\TextAt{21}{256.989}{433.084}{\Character{23}}%
\TextAt{9}{265.492}{428.148}{\Character{61}}%
\TextAt{13}{277.917}{428.148}{a}%
\TextAt{17}{284.062}{426.355}{0}%
\TextAt{9}{291.451}{428.148}{\Character{43}}%
\TextAt{13}{303.213}{428.148}{a}%
\TextAt{17}{309.357}{426.355}{1}%
\TextAt{13}{314.09}{428.148}{p}%
\TextAt{9}{322.622}{428.148}{\Character{43}}%
\TextAt{13}{334.383}{428.148}{a}%
\TextAt{17}{340.528}{426.355}{2}%
\TextAt{13}{345.26}{428.148}{p}%
\TextAt{17}{351.135}{433.084}{2}%
\TextAt{9}{358.524}{428.148}{\Character{43}}%
\TextAt{16}{370.286}{428.148}{\Character{1}}%
\TextAt{16}{375.582}{428.148}{\Character{1}}%
\TextAt{16}{380.89}{428.148}{\Character{1}}%
\TextAt{13}{388.218}{428.148}{\Character{58}}%
\TextAt{8}{79.226}{391.01}{This}%
\TextAt{8}{104.858}{391.01}{notation}%
\TextAt{8}{150.61}{391.01}{should}%
\TextAt{8}{186.654}{391.01}{at}%
\TextAt{8}{199.876}{391.01}{first}%
\TextAt{8}{222.913}{391.01}{be}%
\TextAt{8}{237.749}{391.01}{understood}%
\TextAt{8}{297.859}{391.01}{in}%
\TextAt{8}{310.424}{391.01}{a}%
\TextAt{8}{319.091}{391.01}{purely}%
\TextAt{8}{354.095}{391.01}{formal}%
\TextAt{8}{389.757}{391.01}{sense,}%
\TextAt{8}{422.169}{391.01}{i.e.,}%
\TextAt{20}{443.404}{399.976}{\Character{80}}%
\NativeMathAt{456.023}{396.821}{7.97}{\infty}%
\TextAt{21}{456.023}{387.523}{\Character{23}}%
\TextAt{17}{460.707}{387.523}{\Character{61}0}%
\TextAt{13}{474.018}{391.01}{a}%
\TextAt{21}{480.163}{389.216}{\Character{23}}%
\TextAt{13}{485.345}{391.01}{p}%
\TextAt{21}{491.22}{395.348}{\Character{23}}%
\TextAt{8}{499.211}{391.01}{simply}%
\TextAt{8}{79.226}{373.964}{stands}%
\TextAt{8}{115.773}{373.964}{for}%
\TextAt{8}{133.658}{373.964}{the}%
\TextAt{8}{153.814}{373.964}{sequence}%
\TextAt{8}{202.317}{373.964}{of}%
\TextAt{8}{215.647}{373.964}{partial}%
\TextAt{8}{253.377}{373.964}{sums}%
\TextAt{13}{225.193}{337.417}{s}%
\TextAt{21}{230.707}{335.624}{n}%
\TextAt{9}{239.665}{337.417}{\Character{61}}%
\TextAt{21}{252.745}{352.361}{n}%
\TextAt{3}{257.883}{352.361}{\Character{0}}%
\TextAt{17}{264.47}{352.361}{1}%
\TextAt{20}{252.09}{348.774}{\Character{88}}%
\TextAt{21}{252.972}{323.668}{\Character{23}}%
\TextAt{17}{257.656}{323.668}{\Character{61}0}%
\TextAt{13}{271.351}{337.417}{a}%
\TextAt{21}{277.496}{335.624}{\Character{23}}%
\TextAt{13}{282.678}{337.417}{p}%
\TextAt{21}{288.553}{342.353}{\Character{23}}%
\TextAt{13}{293.735}{337.417}{\Character{59}}%
\TextAt{13}{322.87}{337.417}{n}%
\TextAt{9}{333.198}{337.417}{\Character{61}}%
\TextAt{9}{345.608}{337.417}{1}%
\TextAt{13}{351.477}{337.417}{\Character{59}}%
\TextAt{9}{356.721}{337.417}{2}%
\TextAt{13}{362.574}{337.417}{\Character{59}}%
\TextAt{13}{367.798}{337.417}{\Character{58}}%
\TextAt{13}{373.035}{337.417}{\Character{58}}%
\TextAt{13}{378.259}{337.417}{\Character{58}}%
\TextAt{13}{383.496}{337.417}{\Character{58}}%
\TextAt{12}{79.226}{291.589}{Definition}%
\TextAt{12}{142.015}{291.589}{(1.1).}%
\TextAt{8}{177.33}{291.589}{Fix}%
\TextAt{8}{198.3}{291.589}{a}%
\TextAt{8}{208.055}{291.589}{prime}%
\TextAt{8}{241.219}{291.589}{number}%
\TextAt{13}{283.815}{291.589}{p}%
\TextAt{8}{289.69}{291.589}{.}%
\TextAt{8}{298.138}{291.589}{A}%
\TextAt{13}{310.811}{291.589}{p}%
\TextAt{12}{316.686}{291.589}{-adic}%
\TextAt{12}{349.383}{291.589}{integer}%
\TextAt{8}{393.857}{291.589}{is}%
\TextAt{8}{405.621}{291.589}{a}%
\TextAt{8}{415.376}{291.589}{formal}%
\TextAt{8}{452.126}{291.589}{infinite}%
\TextAt{8}{491.793}{291.589}{series}%
\TextAt{13}{249.222}{262.588}{a}%
\TextAt{17}{255.366}{260.794}{0}%
\TextAt{9}{262.755}{262.588}{\Character{43}}%
\TextAt{13}{274.517}{262.588}{a}%
\TextAt{17}{280.662}{260.794}{1}%
\TextAt{13}{285.394}{262.588}{p}%
\TextAt{9}{293.926}{262.588}{\Character{43}}%
\TextAt{13}{305.687}{262.588}{a}%
\TextAt{17}{311.832}{260.794}{2}%
\TextAt{13}{316.564}{262.588}{p}%
\TextAt{17}{322.44}{267.524}{2}%
\TextAt{9}{329.828}{262.588}{\Character{43}}%
\TextAt{16}{341.59}{262.588}{\Character{1}}%
\TextAt{16}{346.886}{262.588}{\Character{1}}%
\TextAt{16}{352.194}{262.588}{\Character{1}}%
\TextAt{13}{359.522}{262.588}{\Character{59}}%
\TextAt{8}{79.226}{233.587}{where}%
\TextAt{9}{113.035}{233.587}{0}%
\NativeMathAt{122.209}{233.587}{11.955}{\leq}%
\TextAt{13}{134.829}{233.587}{a}%
\TextAt{21}{140.973}{231.793}{i}%
\TextAt{13}{147.676}{233.587}{\Character{60}}%
\TextAt{13}{160.086}{233.587}{p}%
\TextAt{8}{165.976}{233.587}{,}%
\TextAt{8}{173.125}{233.587}{for}%
\TextAt{8}{191.01}{233.587}{all}%
\TextAt{13}{207.269}{233.587}{i}%
\TextAt{9}{214.584}{233.587}{\Character{61}}%
\TextAt{9}{226.994}{233.587}{0}%
\TextAt{13}{232.862}{233.587}{\Character{59}}%
\TextAt{9}{238.106}{233.587}{1}%
\TextAt{13}{243.959}{233.587}{\Character{59}}%
\TextAt{9}{249.203}{233.587}{2}%
\TextAt{13}{255.056}{233.587}{\Character{59}}%
\TextAt{13}{260.28}{233.587}{\Character{58}}%
\TextAt{13}{265.517}{233.587}{\Character{58}}%
\TextAt{13}{270.753}{233.587}{\Character{58}}%
\TextAt{8}{274.041}{233.587}{.}%
\TextAt{8}{282.489}{233.587}{The}%
\TextAt{8}{306.543}{233.587}{set}%
\TextAt{8}{324.81}{233.587}{of}%
\TextAt{8}{338.14}{233.587}{all}%
\TextAt{13}{354.399}{233.587}{p}%
\TextAt{8}{360.274}{233.587}{-adic}%
\TextAt{8}{388.883}{233.587}{integers}%
\TextAt{8}{432.197}{233.587}{is}%
\TextAt{8}{443.961}{233.587}{denoted}%
\TextAt{8}{488.183}{233.587}{by}%
\TextAt{18}{504.442}{233.587}{Z}%
\TextAt{21}{512.412}{231.793}{p}%
\TextAt{8}{517.173}{233.587}{.}%
\TextAt{8}{97.159}{207.574}{The}%
\TextAt{13}{122.972}{207.574}{p}%
\TextAt{8}{128.847}{207.574}{-adic}%
\TextAt{8}{159.215}{207.574}{expansion}%
\TextAt{8}{215.346}{207.574}{of}%
\TextAt{8}{230.435}{207.574}{an}%
\TextAt{8}{248.453}{207.574}{arbitrary}%
\TextAt{8}{299.982}{207.574}{number}%
\TextAt{13}{344.337}{207.574}{f}%
\NativeMathAt{357.703}{207.574}{11.955}{\in}%
\TextAt{18}{371.991}{207.574}{Z}%
\TextAt{17}{379.962}{205.714}{\Character{40}}%
\TextAt{21}{383.255}{205.714}{p}%
\TextAt{17}{387.518}{205.714}{\Character{41}}%
\TextAt{8}{396.965}{207.574}{results}%
\TextAt{8}{435.917}{207.574}{from}%
\TextAt{8}{465.317}{207.574}{the}%
\TextAt{8}{487.232}{207.574}{following}%
\TextAt{8}{79.226}{190.528}{proposition}%
\TextAt{8}{140.759}{190.528}{about}%
\TextAt{8}{174.258}{190.528}{the}%
\TextAt{8}{194.414}{190.528}{residue}%
\TextAt{8}{234.141}{190.528}{classes}%
\TextAt{8}{271.394}{190.528}{in}%
\TextAt{18}{285.046}{190.528}{Z}%
\TextAt{9}{293.017}{190.528}{\Character{47}}%
\TextAt{13}{298.87}{190.528}{p}%
\TextAt{21}{304.745}{194.867}{n}%
\TextAt{18}{310.381}{190.528}{Z}%
\TextAt{8}{318.351}{190.528}{.}%
\TextAt{12}{79.226}{164.516}{Proposition}%
\TextAt{12}{152.224}{164.516}{(1.2).}%
\TextAt{8}{187.788}{164.516}{The}%
\TextAt{8}{212.091}{164.516}{residue}%
\TextAt{8}{252.067}{164.516}{classes}%
\TextAt{13}{289.569}{164.516}{a}%
\TextAt{8}{299.034}{164.516}{mod}%
\TextAt{13}{324.472}{164.516}{p}%
\TextAt{21}{330.347}{168.854}{n}%
\NativeMathAt{339.729}{164.516}{11.955}{\in}%
\TextAt{18}{351.444}{164.516}{Z}%
\TextAt{9}{359.414}{164.516}{\Character{47}}%
\TextAt{13}{365.267}{164.516}{p}%
\TextAt{21}{371.142}{168.854}{n}%
\TextAt{18}{376.779}{164.516}{Z}%
\TextAt{8}{388.895}{164.516}{can}%
\TextAt{8}{410.604}{164.516}{be}%
\TextAt{8}{426.777}{164.516}{uniquely}%
\TextAt{8}{474.5}{164.516}{represented}%
\TextAt{8}{79.226}{147.47}{in}%
\TextAt{8}{92.879}{147.47}{the}%
\TextAt{8}{113.035}{147.47}{form}%
\TextAt{13}{182.785}{130.424}{a}%
\NativeMathAt{192.25}{130.424}{11.955}{\equiv}%
\TextAt{13}{204.87}{130.424}{a}%
\TextAt{17}{211.015}{128.631}{0}%
\TextAt{9}{218.404}{130.424}{\Character{43}}%
\TextAt{13}{230.165}{130.424}{a}%
\TextAt{17}{236.31}{128.631}{1}%
\TextAt{13}{241.042}{130.424}{p}%
\TextAt{9}{249.574}{130.424}{\Character{43}}%
\TextAt{13}{261.335}{130.424}{a}%
\TextAt{17}{267.48}{128.631}{2}%
\TextAt{13}{272.213}{130.424}{p}%
\TextAt{17}{278.088}{135.36}{2}%
\TextAt{9}{285.477}{130.424}{\Character{43}}%
\TextAt{16}{297.238}{130.424}{\Character{1}}%
\TextAt{16}{302.534}{130.424}{\Character{1}}%
\TextAt{16}{307.842}{130.424}{\Character{1}}%
\TextAt{9}{313.842}{130.424}{\Character{43}}%
\TextAt{13}{325.604}{130.424}{a}%
\TextAt{21}{331.749}{128.631}{n}%
\TextAt{3}{336.887}{128.631}{\Character{0}}%
\TextAt{17}{343.473}{128.631}{1}%
\TextAt{13}{348.206}{130.424}{p}%
\TextAt{21}{354.081}{135.36}{n}%
\TextAt{3}{359.219}{135.36}{\Character{0}}%
\TextAt{17}{365.806}{135.36}{1}%
\TextAt{9}{382.493}{130.424}{\Character{40}}%
\TextAt{8}{387.045}{130.424}{mod}%
\TextAt{13}{413.147}{130.424}{p}%
\TextAt{21}{419.022}{135.36}{n}%
\TextAt{9}{424.659}{130.424}{\Character{41}}%
\TextAt{8}{79.226}{106.902}{where}%
\TextAt{9}{113.035}{106.902}{0}%
\NativeMathAt{122.209}{106.902}{11.955}{\leq}%
\TextAt{13}{134.829}{106.902}{a}%
\TextAt{21}{140.973}{105.109}{i}%
\TextAt{13}{147.676}{106.902}{\Character{60}}%
\TextAt{13}{160.086}{106.902}{p}%
\TextAt{8}{169.874}{106.902}{for}%
\TextAt{13}{187.758}{106.902}{i}%
\TextAt{9}{195.073}{106.902}{\Character{61}}%
\TextAt{9}{207.483}{106.902}{0}%
\TextAt{13}{213.351}{106.902}{\Character{59}}%
\TextAt{13}{218.575}{106.902}{\Character{58}}%
\TextAt{13}{223.812}{106.902}{\Character{58}}%
\TextAt{13}{229.048}{106.902}{\Character{58}}%
\TextAt{13}{234.273}{106.902}{\Character{59}}%
\TextAt{13}{239.509}{106.902}{n}%
\TextAt{16}{249.216}{106.902}{\Character{0}}%
\TextAt{9}{261.171}{106.902}{1}%
\TextAt{8}{267.024}{106.902}{.}%
\TextAt{8}{300.14}{43.033}{92}%
\Destination{78.23}{809.815}{nav.page.92}%
\Destination{79.36}{776.155}{nav.print.100}%
\Destination{79.36}{305.555}{nav.number.2.1.1}%
\Destination{79.36}{178.555}{nav.number.2.1.2}%
}

% Page 105 (body page 93)
\Page{%
\TextAt{2}{79.226}{796.943}{II.}%
\TextAt{2}{94.073}{796.943}{The}%
\TextAt{2}{116.492}{796.943}{Theory}%
\TextAt{2}{154.401}{796.943}{of}%
\TextAt{2}{166.826}{796.943}{Valuations}%
\Vector{q 0 G 0 g 0.299 w 7.226 18.729 m 460.769 18.729 l S Q}%
\TextAt{19}{79.226}{761.124}{Proof}%
\TextAt{19}{110.62}{761.124}{(induction}%
\TextAt{19}{167.407}{761.124}{on}%
\TextAt{13}{184.145}{761.124}{n}%
\TextAt{19}{191.132}{761.124}{).}%
\TextAt{8}{205.478}{761.124}{This}%
\TextAt{8}{231.479}{761.124}{is}%
\TextAt{8}{242.524}{761.124}{clear}%
\TextAt{8}{269.768}{761.124}{for}%
\TextAt{13}{286.934}{761.124}{n}%
\TextAt{9}{297.242}{761.124}{\Character{61}}%
\TextAt{9}{309.651}{761.124}{1}%
\TextAt{8}{315.521}{761.124}{.}%
\TextAt{8}{323.729}{761.124}{Assume}%
\TextAt{8}{366.372}{761.124}{the}%
\TextAt{8}{385.809}{761.124}{statement}%
\TextAt{8}{439.462}{761.124}{is}%
\TextAt{8}{450.506}{761.124}{proved}%
\TextAt{8}{487.841}{761.124}{for}%
\TextAt{13}{505.006}{761.124}{n}%
\TextAt{16}{513.18}{761.124}{\Character{0}}%
\TextAt{9}{523.665}{761.124}{1}%
\TextAt{8}{529.518}{761.124}{.}%
\TextAt{8}{79.226}{744.078}{Then}%
\TextAt{8}{109.784}{744.078}{we}%
\TextAt{8}{127.011}{744.078}{have}%
\TextAt{8}{154.006}{744.078}{a}%
\TextAt{8}{163.761}{744.078}{unique}%
\TextAt{8}{201.803}{744.078}{representation}%
\TextAt{13}{189.2}{716.51}{a}%
\TextAt{9}{198.666}{716.51}{\Character{61}}%
\TextAt{13}{211.091}{716.51}{a}%
\TextAt{17}{217.236}{714.717}{0}%
\TextAt{9}{224.625}{716.51}{\Character{43}}%
\TextAt{13}{236.387}{716.51}{a}%
\TextAt{17}{242.532}{714.717}{1}%
\TextAt{13}{247.264}{716.51}{p}%
\TextAt{9}{255.796}{716.51}{\Character{43}}%
\TextAt{13}{267.557}{716.51}{a}%
\TextAt{17}{273.702}{714.717}{2}%
\TextAt{13}{278.434}{716.51}{p}%
\TextAt{17}{284.309}{721.446}{2}%
\TextAt{9}{291.698}{716.51}{\Character{43}}%
\TextAt{16}{303.46}{716.51}{\Character{1}}%
\TextAt{16}{308.756}{716.51}{\Character{1}}%
\TextAt{16}{314.064}{716.51}{\Character{1}}%
\TextAt{9}{320.064}{716.51}{\Character{43}}%
\TextAt{13}{331.825}{716.51}{a}%
\TextAt{21}{337.97}{714.717}{n}%
\TextAt{3}{343.108}{714.717}{\Character{0}}%
\TextAt{17}{349.695}{714.717}{2}%
\TextAt{13}{354.427}{716.51}{p}%
\TextAt{21}{360.302}{721.446}{n}%
\TextAt{3}{365.441}{721.446}{\Character{0}}%
\TextAt{17}{372.027}{721.446}{2}%
\TextAt{9}{379.416}{716.51}{\Character{43}}%
\TextAt{13}{391.177}{716.51}{g}%
\TextAt{13}{397.19}{716.51}{p}%
\TextAt{21}{403.087}{721.446}{n}%
\TextAt{3}{408.225}{721.446}{\Character{0}}%
\TextAt{17}{414.812}{721.446}{1}%
\TextAt{13}{419.544}{716.51}{\Character{59}}%
\TextAt{8}{79.226}{688.942}{for}%
\TextAt{8}{98.129}{688.942}{some}%
\TextAt{8}{128.473}{688.942}{integer}%
\TextAt{13}{168.191}{688.942}{g}%
\TextAt{8}{174.225}{688.942}{.}%
\TextAt{8}{185.728}{688.942}{If}%
\TextAt{13}{198.438}{688.942}{g}%
\NativeMathAt{209.528}{688.942}{11.955}{\equiv}%
\TextAt{13}{223.883}{688.942}{a}%
\TextAt{21}{230.028}{687.149}{n}%
\TextAt{3}{235.166}{687.149}{\Character{0}}%
\TextAt{17}{241.752}{687.149}{1}%
\TextAt{8}{249.805}{688.942}{mod}%
\TextAt{13}{275.243}{688.942}{p}%
\TextAt{8}{286.034}{688.942}{such}%
\TextAt{8}{313.449}{688.942}{that}%
\TextAt{9}{339.836}{688.942}{0}%
\NativeMathAt{350.745}{688.942}{11.955}{\leq}%
\TextAt{13}{365.1}{688.942}{a}%
\TextAt{21}{371.245}{687.149}{n}%
\TextAt{3}{376.383}{687.149}{\Character{0}}%
\TextAt{17}{382.969}{687.149}{1}%
\TextAt{13}{392.757}{688.942}{\Character{60}}%
\TextAt{13}{406.9}{688.942}{p}%
\TextAt{8}{412.793}{688.942}{,}%
\TextAt{8}{421.215}{688.942}{then}%
\TextAt{13}{448.893}{688.942}{a}%
\TextAt{21}{455.038}{687.149}{n}%
\TextAt{3}{460.176}{687.149}{\Character{0}}%
\TextAt{17}{466.763}{687.149}{1}%
\TextAt{8}{476.411}{688.942}{is}%
\TextAt{8}{489.193}{688.942}{uniquely}%
\TextAt{8}{79.226}{671.896}{determined}%
\TextAt{8}{140.353}{671.896}{by}%
\TextAt{13}{156.612}{671.896}{a}%
\TextAt{8}{162.757}{671.896}{,}%
\TextAt{8}{169.906}{671.896}{and}%
\TextAt{8}{192.669}{671.896}{the}%
\TextAt{8}{212.825}{671.896}{congruence}%
\TextAt{8}{273.306}{671.896}{of}%
\TextAt{8}{286.636}{671.896}{the}%
\TextAt{8}{306.793}{671.896}{proposition}%
\TextAt{8}{368.326}{671.896}{holds.}%
\Vector{q 0 G 0 g 0.399 w 452.334 -102.104 m 452.334 -94.034 l S Q}%
\Vector{q 0 G 0 g 0.399 w 452.534 -94.233 m 459.707 -94.233 l S Q}%
\Vector{q 0 G 0 g 0.399 w 452.534 -101.905 m 459.707 -101.905 l S Q}%
\Vector{q 0 G 0 g 0.399 w 459.906 -102.104 m 459.906 -94.034 l S Q}%
\TextAt{8}{97.159}{643.305}{Every}%
\TextAt{8}{130.671}{643.305}{integer}%
\TextAt{13}{169.228}{643.305}{f}%
\TextAt{8}{180.03}{643.305}{and,}%
\TextAt{8}{205.93}{643.305}{more}%
\TextAt{8}{235.055}{643.305}{generally,}%
\TextAt{8}{286.97}{643.305}{every}%
\TextAt{8}{317.72}{643.305}{rational}%
\TextAt{8}{361.167}{643.305}{number}%
\TextAt{13}{403.621}{643.305}{f}%
\NativeMathAt{413.989}{643.305}{11.955}{\in}%
\TextAt{18}{425.28}{643.305}{Z}%
\TextAt{17}{433.25}{641.445}{\Character{40}}%
\TextAt{21}{436.543}{641.445}{p}%
\TextAt{17}{440.806}{641.445}{\Character{41}}%
\TextAt{8}{448.353}{643.305}{the}%
\TextAt{8}{468.367}{643.305}{denominator}%
\TextAt{8}{79.226}{626.259}{of}%
\TextAt{8}{92.556}{626.259}{which}%
\TextAt{8}{126.043}{626.259}{is}%
\TextAt{8}{137.806}{626.259}{not}%
\TextAt{8}{158.62}{626.259}{divisible}%
\TextAt{8}{204.528}{626.259}{by}%
\TextAt{13}{220.787}{626.259}{p}%
\TextAt{8}{226.662}{626.259}{,}%
\TextAt{8}{233.812}{626.259}{defines}%
\TextAt{8}{272.236}{626.259}{a}%
\TextAt{8}{281.991}{626.259}{sequence}%
\TextAt{8}{330.493}{626.259}{of}%
\TextAt{8}{343.823}{626.259}{residue}%
\TextAt{8}{383.55}{626.259}{classes}%
\TextAt{9}{198.88}{598.691}{\Character{22}}%
\TextAt{13}{198.399}{598.691}{s}%
\TextAt{21}{203.913}{596.898}{n}%
\TextAt{9}{212.871}{598.691}{\Character{61}}%
\TextAt{13}{225.296}{598.691}{f}%
\TextAt{8}{235.663}{598.691}{mod}%
\TextAt{13}{261.101}{598.691}{p}%
\TextAt{21}{266.976}{603.627}{n}%
\NativeMathAt{275.934}{598.691}{11.955}{\in}%
\TextAt{18}{287.225}{598.691}{Z}%
\TextAt{9}{295.195}{598.691}{\Character{47}}%
\TextAt{13}{301.048}{598.691}{p}%
\TextAt{21}{306.923}{603.627}{n}%
\TextAt{18}{312.559}{598.691}{Z}%
\TextAt{13}{320.529}{598.691}{\Character{59}}%
\TextAt{13}{349.664}{598.691}{n}%
\TextAt{9}{359.992}{598.691}{\Character{61}}%
\TextAt{9}{372.401}{598.691}{1}%
\TextAt{13}{378.271}{598.691}{\Character{59}}%
\TextAt{9}{383.515}{598.691}{2}%
\TextAt{13}{389.368}{598.691}{\Character{59}}%
\TextAt{13}{394.592}{598.691}{\Character{58}}%
\TextAt{13}{399.829}{598.691}{\Character{58}}%
\TextAt{13}{405.053}{598.691}{\Character{58}}%
\TextAt{13}{410.29}{598.691}{\Character{59}}%
\TextAt{8}{79.226}{571.123}{for}%
\TextAt{8}{97.111}{571.123}{which}%
\TextAt{8}{130.597}{571.123}{we}%
\TextAt{8}{147.825}{571.123}{find,}%
\TextAt{8}{174.485}{571.123}{by}%
\TextAt{8}{190.744}{571.123}{the}%
\TextAt{8}{210.9}{571.123}{preceding}%
\TextAt{8}{263.575}{571.123}{proposition,}%
\TextAt{9}{209.517}{545.222}{\Character{22}}%
\TextAt{13}{209.036}{545.222}{s}%
\TextAt{17}{214.55}{543.429}{1}%
\TextAt{9}{222.603}{545.222}{\Character{61}}%
\TextAt{13}{235.029}{545.222}{a}%
\TextAt{17}{241.174}{543.429}{0}%
\TextAt{8}{249.227}{545.222}{mod}%
\TextAt{13}{274.665}{545.222}{p\Character{59}}%
\TextAt{9}{209.517}{525.187}{\Character{22}}%
\TextAt{13}{209.036}{525.187}{s}%
\TextAt{17}{214.55}{523.394}{2}%
\TextAt{9}{222.603}{525.187}{\Character{61}}%
\TextAt{13}{235.029}{525.187}{a}%
\TextAt{17}{241.174}{523.394}{0}%
\TextAt{9}{248.563}{525.187}{\Character{43}}%
\TextAt{13}{260.324}{525.187}{a}%
\TextAt{17}{266.469}{523.394}{1}%
\TextAt{13}{271.201}{525.187}{p}%
\TextAt{8}{280.397}{525.187}{mod}%
\TextAt{13}{305.835}{525.187}{p}%
\TextAt{17}{311.71}{530.124}{2}%
\TextAt{13}{316.443}{525.187}{\Character{59}}%
\TextAt{9}{209.517}{505.153}{\Character{22}}%
\TextAt{13}{209.036}{505.153}{s}%
\TextAt{17}{214.55}{503.359}{3}%
\TextAt{9}{222.603}{505.153}{\Character{61}}%
\TextAt{13}{235.029}{505.153}{a}%
\TextAt{17}{241.174}{503.359}{0}%
\TextAt{9}{248.563}{505.153}{\Character{43}}%
\TextAt{13}{260.324}{505.153}{a}%
\TextAt{17}{266.469}{503.359}{1}%
\TextAt{13}{271.201}{505.153}{p}%
\TextAt{9}{279.733}{505.153}{\Character{43}}%
\TextAt{13}{291.494}{505.153}{a}%
\TextAt{17}{297.639}{503.359}{2}%
\TextAt{13}{302.372}{505.153}{p}%
\TextAt{17}{308.247}{510.089}{2}%
\TextAt{8}{316.3}{505.153}{mod}%
\TextAt{13}{341.738}{505.153}{p}%
\TextAt{17}{347.613}{510.089}{3}%
\TextAt{13}{352.345}{505.153}{\Character{59}}%
\TextAt{8}{381.5}{505.153}{etc.,}%
\TextAt{8}{79.226}{479.554}{with}%
\TextAt{8}{106.97}{479.554}{uniquely}%
\TextAt{8}{155.527}{479.554}{determined}%
\TextAt{8}{217.737}{479.554}{coefficients}%
\TextAt{13}{278.058}{479.554}{a}%
\TextAt{17}{284.203}{477.761}{0}%
\TextAt{13}{288.936}{479.554}{\Character{59}}%
\TextAt{13}{294.16}{479.554}{a}%
\TextAt{17}{300.325}{477.761}{1}%
\TextAt{13}{305.057}{479.554}{\Character{59}}%
\TextAt{13}{310.281}{479.554}{a}%
\TextAt{17}{316.446}{477.761}{2}%
\TextAt{13}{321.178}{479.554}{\Character{59}}%
\TextAt{13}{326.402}{479.554}{\Character{58}}%
\TextAt{13}{331.639}{479.554}{\Character{58}}%
\TextAt{13}{336.875}{479.554}{\Character{58}}%
\NativeMathAt{345.33}{479.554}{11.955}{\in}%
\TextAt{16}{358.457}{479.554}{\Character{102}}%
\TextAt{9}{364.444}{479.554}{0}%
\TextAt{13}{370.297}{479.554}{\Character{59}}%
\TextAt{9}{375.542}{479.554}{1}%
\TextAt{13}{381.395}{479.554}{\Character{59}}%
\TextAt{13}{386.619}{479.554}{\Character{58}}%
\TextAt{13}{391.856}{479.554}{\Character{58}}%
\TextAt{13}{397.08}{479.554}{\Character{58}}%
\TextAt{13}{402.317}{479.554}{\Character{59}}%
\TextAt{13}{407.553}{479.554}{p}%
\TextAt{16}{416.886}{479.554}{\Character{0}}%
\TextAt{9}{429.579}{479.554}{1}%
\TextAt{16}{435.432}{479.554}{\Character{103}}%
\TextAt{8}{446.391}{479.554}{which}%
\TextAt{8}{480.96}{479.554}{keep}%
\TextAt{8}{508.704}{479.554}{their}%
\TextAt{8}{79.226}{462.508}{meaning}%
\TextAt{8}{126.055}{462.508}{from}%
\TextAt{8}{153.695}{462.508}{one}%
\TextAt{8}{175.154}{462.508}{line}%
\TextAt{8}{197.26}{462.508}{to}%
\TextAt{8}{211.57}{462.508}{the}%
\TextAt{8}{231.726}{462.508}{next.}%
\TextAt{8}{262.614}{462.508}{The}%
\TextAt{8}{286.668}{462.508}{sequence}%
\TextAt{8}{335.17}{462.508}{of}%
\TextAt{8}{348.5}{462.508}{numbers}%
\TextAt{13}{163.182}{434.941}{s}%
\TextAt{21}{168.696}{433.147}{n}%
\TextAt{9}{177.653}{434.941}{\Character{61}}%
\TextAt{13}{190.079}{434.941}{a}%
\TextAt{17}{196.224}{433.147}{0}%
\TextAt{9}{203.613}{434.941}{\Character{43}}%
\TextAt{13}{215.374}{434.941}{a}%
\TextAt{17}{221.519}{433.147}{1}%
\TextAt{13}{226.251}{434.941}{p}%
\TextAt{9}{234.783}{434.941}{\Character{43}}%
\TextAt{13}{246.544}{434.941}{a}%
\TextAt{17}{252.689}{433.147}{2}%
\TextAt{13}{257.422}{434.941}{p}%
\TextAt{17}{263.297}{439.877}{2}%
\TextAt{9}{270.686}{434.941}{\Character{43}}%
\TextAt{16}{282.447}{434.941}{\Character{1}}%
\TextAt{16}{287.743}{434.941}{\Character{1}}%
\TextAt{16}{293.051}{434.941}{\Character{1}}%
\TextAt{9}{299.051}{434.941}{\Character{43}}%
\TextAt{13}{310.813}{434.941}{a}%
\TextAt{21}{316.958}{433.147}{n}%
\TextAt{3}{322.096}{433.147}{\Character{0}}%
\TextAt{17}{328.682}{433.147}{1}%
\TextAt{13}{333.415}{434.941}{p}%
\TextAt{21}{339.29}{439.877}{n}%
\TextAt{3}{344.428}{439.877}{\Character{0}}%
\TextAt{17}{351.014}{439.877}{1}%
\TextAt{13}{355.747}{434.941}{\Character{59}}%
\TextAt{13}{384.882}{434.941}{n}%
\TextAt{9}{395.21}{434.941}{\Character{61}}%
\TextAt{9}{407.62}{434.941}{1}%
\TextAt{13}{413.488}{434.941}{\Character{59}}%
\TextAt{9}{418.732}{434.941}{2}%
\TextAt{13}{424.585}{434.941}{\Character{59}}%
\TextAt{13}{429.809}{434.941}{\Character{58}}%
\TextAt{13}{435.046}{434.941}{\Character{58}}%
\TextAt{13}{440.282}{434.941}{\Character{58}}%
\TextAt{13}{445.507}{434.941}{\Character{59}}%
\TextAt{8}{79.226}{407.373}{defines}%
\TextAt{8}{117.65}{407.373}{a}%
\TextAt{13}{127.405}{407.373}{p}%
\TextAt{8}{133.281}{407.373}{-adic}%
\TextAt{8}{161.889}{407.373}{integer}%
\NativeMathAt{274.278}{399.299}{7.97}{\infty}%
\TextAt{20}{269.878}{395.712}{\Character{88}}%
\TextAt{21}{270.76}{370.606}{\Character{23}}%
\TextAt{17}{275.444}{370.606}{\Character{61}0}%
\TextAt{13}{289.139}{384.355}{a}%
\TextAt{21}{295.284}{382.562}{\Character{23}}%
\TextAt{13}{300.466}{384.355}{p}%
\TextAt{21}{306.341}{389.291}{\Character{23}}%
\NativeMathAt{314.844}{384.355}{11.955}{\in}%
\TextAt{18}{326.135}{384.355}{Z}%
\TextAt{21}{334.105}{382.562}{p}%
\TextAt{13}{338.866}{384.355}{\Character{58}}%
\TextAt{8}{79.226}{353.587}{We}%
\TextAt{8}{101.396}{353.587}{call}%
\TextAt{8}{124.881}{353.587}{it}%
\TextAt{8}{138.61}{353.587}{the}%
\TextAt{13}{160.792}{353.587}{p}%
\TextAt{12}{166.667}{353.587}{-adic}%
\TextAt{12}{201.694}{353.587}{expansion}%
\TextAt{8}{265.588}{353.587}{of}%
\TextAt{13}{280.943}{353.587}{f}%
\TextAt{8}{287.99}{353.587}{.}%
\TextAt{8}{302.514}{353.587}{In}%
\TextAt{8}{319.161}{353.587}{analogy}%
\TextAt{8}{364.452}{353.587}{with}%
\TextAt{8}{393.137}{353.587}{the}%
\TextAt{8}{415.319}{353.587}{Laurent}%
\TextAt{8}{461.411}{353.587}{series}%
\TextAt{13}{494.771}{353.587}{f}%
\TextAt{9}{501.817}{353.587}{\Character{40}}%
\TextAt{13}{506.37}{353.587}{z}%
\TextAt{9}{512.34}{353.587}{\Character{41}}%
\TextAt{9}{523.65}{353.587}{\Character{61}}%
\TextAt{20}{79.226}{345.507}{\Character{80}}%
\NativeMathAt{91.846}{342.352}{7.97}{\infty}%
\TextAt{21}{91.846}{333.054}{\Character{23}}%
\TextAt{17}{96.529}{333.054}{\Character{61}}%
\TextAt{3}{103.116}{333.054}{\Character{0}}%
\TextAt{21}{109.702}{333.054}{m}%
\TextAt{13}{119.684}{336.541}{a}%
\TextAt{21}{125.829}{334.747}{\Character{23}}%
\TextAt{9}{131.011}{336.541}{\Character{40}}%
\TextAt{13}{135.563}{336.541}{z}%
\TextAt{16}{145.148}{336.541}{\Character{0}}%
\TextAt{13}{158.062}{336.541}{a}%
\TextAt{9}{164.207}{336.541}{\Character{41}}%
\TextAt{21}{168.759}{340.879}{\Character{23}}%
\TextAt{8}{173.941}{336.541}{,}%
\TextAt{8}{182.848}{336.541}{we}%
\TextAt{8}{201.481}{336.541}{now}%
\TextAt{8}{227.275}{336.541}{extend}%
\TextAt{8}{266.722}{336.541}{the}%
\TextAt{8}{288.285}{336.541}{domain}%
\TextAt{8}{331.318}{336.541}{of}%
\TextAt{13}{346.054}{336.541}{p}%
\TextAt{8}{351.929}{336.541}{-adic}%
\TextAt{8}{381.944}{336.541}{integers}%
\TextAt{8}{426.663}{336.541}{into}%
\TextAt{8}{451.812}{336.541}{that}%
\TextAt{8}{478.587}{336.541}{of}%
\TextAt{8}{493.322}{336.541}{the}%
\TextAt{8}{514.885}{336.541}{for-}%
\TextAt{8}{79.226}{319.495}{mal}%
\TextAt{8}{101.989}{319.495}{series}%
\NativeMathAt{176.362}{313.752}{7.97}{\infty}%
\TextAt{20}{171.962}{310.165}{\Character{88}}%
\TextAt{21}{167.923}{285.059}{\Character{23}}%
\TextAt{17}{172.607}{285.059}{\Character{61}}%
\TextAt{3}{179.193}{285.059}{\Character{0}}%
\TextAt{21}{185.78}{285.059}{m}%
\TextAt{13}{195.263}{298.808}{a}%
\TextAt{21}{201.408}{297.015}{\Character{23}}%
\TextAt{13}{206.59}{298.808}{p}%
\TextAt{21}{212.465}{303.744}{\Character{23}}%
\TextAt{9}{220.968}{298.808}{\Character{61}}%
\TextAt{13}{233.393}{298.808}{a}%
\TextAt{3}{239.538}{297.015}{\Character{0}}%
\TextAt{21}{246.125}{297.015}{m}%
\TextAt{13}{254.113}{298.808}{p}%
\TextAt{3}{259.988}{303.744}{\Character{0}}%
\TextAt{21}{266.575}{303.744}{m}%
\TextAt{9}{277.22}{298.808}{\Character{43}}%
\TextAt{16}{288.982}{298.808}{\Character{1}}%
\TextAt{16}{294.278}{298.808}{\Character{1}}%
\TextAt{16}{299.586}{298.808}{\Character{1}}%
\TextAt{9}{305.586}{298.808}{\Character{43}}%
\TextAt{13}{317.347}{298.808}{a}%
\TextAt{3}{323.492}{297.015}{\Character{0}}%
\TextAt{17}{330.079}{297.015}{1}%
\TextAt{13}{334.811}{298.808}{p}%
\TextAt{3}{340.686}{303.744}{\Character{0}}%
\TextAt{17}{347.273}{303.744}{1}%
\TextAt{9}{354.662}{298.808}{\Character{43}}%
\TextAt{13}{366.423}{298.808}{a}%
\TextAt{17}{372.568}{297.015}{0}%
\TextAt{9}{379.957}{298.808}{\Character{43}}%
\TextAt{13}{391.718}{298.808}{a}%
\TextAt{17}{397.863}{297.015}{1}%
\TextAt{13}{402.595}{298.808}{p}%
\TextAt{9}{411.127}{298.808}{\Character{43}}%
\TextAt{16}{422.889}{298.808}{\Character{1}}%
\TextAt{16}{428.185}{298.808}{\Character{1}}%
\TextAt{16}{433.493}{298.808}{\Character{1}}%
\TextAt{13}{440.821}{298.808}{\Character{59}}%
\TextAt{8}{79.226}{267.759}{where}%
\TextAt{13}{113.139}{267.759}{m}%
\NativeMathAt{126.875}{267.759}{11.955}{\in}%
\TextAt{18}{138.343}{267.759}{Z}%
\TextAt{8}{150.314}{267.759}{and}%
\TextAt{9}{173.18}{267.759}{0}%
\NativeMathAt{182.53}{267.759}{11.955}{\leq}%
\TextAt{13}{195.326}{267.759}{a}%
\TextAt{21}{201.471}{265.966}{\Character{23}}%
\TextAt{13}{210.15}{267.759}{\Character{60}}%
\TextAt{13}{222.739}{267.759}{p}%
\TextAt{8}{228.627}{267.759}{.}%
\TextAt{8}{237.386}{267.759}{Such}%
\TextAt{8}{265.775}{267.759}{series}%
\TextAt{8}{297.213}{267.759}{we}%
\TextAt{8}{314.544}{267.759}{call}%
\TextAt{8}{336.107}{267.759}{simply}%
\TextAt{13}{373.666}{267.759}{p}%
\TextAt{12}{379.541}{267.759}{-adic}%
\TextAt{12}{412.357}{267.759}{numbers}%
\TextAt{8}{466.558}{267.759}{and}%
\TextAt{8}{489.424}{267.759}{we}%
\TextAt{8}{506.755}{267.759}{write}%
\TextAt{18}{79.226}{250.713}{Q}%
\TextAt{21}{88.525}{248.92}{p}%
\TextAt{8}{96.793}{250.713}{for}%
\TextAt{8}{114.288}{250.713}{the}%
\TextAt{8}{134.054}{250.713}{set}%
\TextAt{8}{151.931}{250.713}{of}%
\TextAt{8}{164.871}{250.713}{all}%
\TextAt{8}{180.74}{250.713}{these}%
\TextAt{13}{210.321}{250.713}{p}%
\TextAt{8}{216.197}{250.713}{-adic}%
\TextAt{8}{244.415}{250.713}{numbers.}%
\TextAt{8}{296.047}{250.713}{If}%
\TextAt{13}{307.349}{250.713}{f}%
\NativeMathAt{317.716}{250.713}{11.955}{\in}%
\TextAt{18}{329.007}{250.713}{Q}%
\TextAt{8}{341.813}{250.713}{is}%
\TextAt{8}{353.187}{250.713}{any}%
\TextAt{8}{374.914}{250.713}{rational}%
\TextAt{8}{418.112}{250.713}{number,}%
\TextAt{8}{463.648}{250.713}{then}%
\TextAt{8}{489.918}{250.713}{we}%
\TextAt{8}{506.755}{250.713}{write}%
\TextAt{13}{193.129}{223.146}{f}%
\TextAt{9}{203.496}{223.146}{\Character{61}}%
\TextAt{13}{217.469}{231.233}{g}%
\Vector{q 0 G 0 g 0.478 w 145.117 -547.865 m 151.856 -547.865 l S Q}%
\TextAt{13}{217.117}{214.945}{h}%
\TextAt{13}{225.051}{223.146}{p}%
\TextAt{3}{230.926}{228.082}{\Character{0}}%
\TextAt{21}{237.513}{228.082}{m}%
\TextAt{8}{269.412}{223.146}{where}%
\TextAt{13}{303.221}{223.146}{g}%
\TextAt{13}{309.235}{223.146}{\Character{59}}%
\TextAt{13}{314.471}{223.146}{h}%
\NativeMathAt{324.559}{223.146}{11.955}{\in}%
\TextAt{18}{335.85}{223.146}{Z}%
\TextAt{13}{343.82}{223.146}{\Character{59}}%
\TextAt{9}{361.019}{223.146}{\Character{40}}%
\TextAt{13}{365.572}{223.146}{g}%
\TextAt{13}{371.586}{223.146}{h\Character{59}}%
\TextAt{13}{383.553}{223.146}{p}%
\TextAt{9}{389.464}{223.146}{\Character{41}}%
\TextAt{9}{397.319}{223.146}{\Character{61}}%
\TextAt{9}{409.74}{223.146}{1}%
\TextAt{13}{415.615}{223.146}{\Character{59}}%
\TextAt{8}{79.226}{195.578}{and}%
\TextAt{8}{101.989}{195.578}{if}%
\TextAt{13}{252.84}{178.532}{a}%
\TextAt{17}{258.985}{176.739}{0}%
\TextAt{9}{266.374}{178.532}{\Character{43}}%
\TextAt{13}{278.135}{178.532}{a}%
\TextAt{17}{284.28}{176.739}{1}%
\TextAt{13}{289.012}{178.532}{p}%
\TextAt{9}{297.544}{178.532}{\Character{43}}%
\TextAt{13}{309.305}{178.532}{a}%
\TextAt{17}{315.45}{176.739}{2}%
\TextAt{13}{320.183}{178.532}{p}%
\TextAt{17}{326.058}{183.468}{2}%
\TextAt{9}{333.447}{178.532}{\Character{43}}%
\TextAt{16}{345.208}{178.532}{\Character{1}}%
\TextAt{16}{350.504}{178.532}{\Character{1}}%
\TextAt{16}{355.812}{178.532}{\Character{1}}%
\TextAt{8}{79.226}{155.625}{is}%
\TextAt{8}{90.99}{155.625}{the}%
\TextAt{13}{111.146}{155.625}{p}%
\TextAt{8}{117.022}{155.625}{-adic}%
\TextAt{8}{145.63}{155.625}{expansion}%
\TextAt{8}{200.002}{155.625}{of}%
\TextAt{13}{213.332}{155.625}{g}%
\TextAt{9}{219.367}{155.625}{\Character{47}}%
\TextAt{13}{225.22}{155.625}{h}%
\TextAt{8}{231.958}{155.625}{,}%
\TextAt{8}{239.107}{155.625}{then}%
\TextAt{8}{265.767}{155.625}{we}%
\TextAt{8}{282.995}{155.625}{attach}%
\TextAt{8}{319.099}{155.625}{to}%
\TextAt{13}{333.41}{155.625}{f}%
\TextAt{8}{344.354}{155.625}{the}%
\TextAt{13}{364.51}{155.625}{p}%
\TextAt{8}{370.385}{155.625}{-adic}%
\TextAt{8}{398.994}{155.625}{number}%
\TextAt{13}{182.45}{128.057}{a}%
\TextAt{17}{188.595}{126.264}{0}%
\TextAt{13}{193.327}{128.057}{p}%
\TextAt{3}{199.203}{132.993}{\Character{0}}%
\TextAt{21}{205.789}{132.993}{m}%
\TextAt{9}{216.434}{128.057}{\Character{43}}%
\TextAt{13}{228.196}{128.057}{a}%
\TextAt{17}{234.341}{126.264}{1}%
\TextAt{13}{239.073}{128.057}{p}%
\TextAt{3}{244.948}{132.993}{\Character{0}}%
\TextAt{21}{251.535}{132.993}{m}%
\TextAt{17}{259.025}{132.993}{\Character{43}1}%
\TextAt{9}{273.001}{128.057}{\Character{43}}%
\TextAt{16}{284.762}{128.057}{\Character{1}}%
\TextAt{16}{290.058}{128.057}{\Character{1}}%
\TextAt{16}{295.366}{128.057}{\Character{1}}%
\TextAt{9}{301.366}{128.057}{\Character{43}}%
\TextAt{13}{313.128}{128.057}{a}%
\TextAt{21}{319.273}{126.264}{m}%
\TextAt{9}{329.918}{128.057}{\Character{43}}%
\TextAt{13}{341.679}{128.057}{a}%
\TextAt{21}{347.824}{126.264}{m}%
\TextAt{17}{355.315}{126.264}{\Character{43}1}%
\TextAt{13}{366.633}{128.057}{p}%
\TextAt{9}{375.165}{128.057}{\Character{43}}%
\TextAt{16}{386.927}{128.057}{\Character{1}}%
\TextAt{16}{392.223}{128.057}{\Character{1}}%
\TextAt{16}{397.531}{128.057}{\Character{1}}%
\NativeMathAt{404.166}{128.057}{11.955}{\in}%
\TextAt{18}{415.486}{128.057}{Q}%
\TextAt{21}{424.785}{126.264}{p}%
\TextAt{8}{79.226}{100.489}{as}%
\TextAt{8}{93.596}{100.489}{its}%
\TextAt{13}{109.915}{100.489}{p}%
\TextAt{8}{115.79}{100.489}{-adic}%
\TextAt{8}{144.399}{100.489}{expansion.}%
\TextAt{8}{203.322}{100.489}{In}%
\TextAt{8}{217.943}{100.489}{this}%
\TextAt{8}{240.766}{100.489}{way}%
\TextAt{8}{264.509}{100.489}{we}%
\TextAt{8}{281.736}{100.489}{obtain}%
\TextAt{8}{318.163}{100.489}{a}%
\TextAt{8}{327.919}{100.489}{canonical}%
\TextAt{8}{379.302}{100.489}{mapping}%
\TextAt{18}{279.742}{72.921}{Q}%
\NativeArrowAt{292.361}{72.921}{11.955}{19.26}{\longrightarrow}%
\TextAt{18}{314.943}{72.921}{Q}%
\TextAt{21}{324.241}{71.128}{p}%
\TextAt{13}{329.002}{72.921}{\Character{59}}%
\TextAt{8}{300.14}{43.033}{93}%
\Destination{78.23}{809.815}{nav.page.93}%
\Destination{79.36}{776.155}{nav.print.101}%
}

% Page 106 (body page 94)
\Page{%
\TextAt{2}{79.226}{796.943}{II.}%
\TextAt{2}{94.073}{796.943}{The}%
\TextAt{2}{116.492}{796.943}{Theory}%
\TextAt{2}{154.401}{796.943}{of}%
\TextAt{2}{166.826}{796.943}{Valuations}%
\Vector{q 0 G 0 g 0.299 w 7.226 18.729 m 460.769 18.729 l S Q}%
\TextAt{8}{79.226}{761.124}{which}%
\TextAt{8}{113.351}{761.124}{takes}%
\TextAt{18}{143.974}{761.124}{Z}%
\TextAt{8}{156.48}{761.124}{into}%
\TextAt{18}{180.862}{761.124}{Z}%
\TextAt{21}{188.832}{759.33}{p}%
\TextAt{8}{198.129}{761.124}{and}%
\TextAt{8}{221.531}{761.124}{is}%
\TextAt{8}{233.933}{761.124}{injective.}%
\TextAt{8}{286.894}{761.124}{For}%
\TextAt{8}{308.502}{761.124}{if}%
\TextAt{13}{319.865}{761.124}{a\Character{59}}%
\TextAt{13}{331.234}{761.124}{b}%
\NativeMathAt{340.641}{761.124}{11.955}{\in}%
\TextAt{18}{353.02}{761.124}{Z}%
\TextAt{8}{365.527}{761.124}{have}%
\TextAt{8}{393.16}{761.124}{the}%
\TextAt{8}{413.955}{761.124}{same}%
\TextAt{13}{443.92}{761.124}{p}%
\TextAt{8}{449.795}{761.124}{-adic}%
\TextAt{8}{479.043}{761.124}{expansion,}%
\TextAt{8}{79.226}{744.078}{then}%
\TextAt{13}{106.511}{744.078}{a}%
\TextAt{16}{115.738}{744.078}{\Character{0}}%
\TextAt{13}{128.119}{744.078}{b}%
\TextAt{8}{137.619}{744.078}{is}%
\TextAt{8}{150.007}{744.078}{divisible}%
\TextAt{8}{196.54}{744.078}{by}%
\TextAt{13}{213.423}{744.078}{p}%
\TextAt{21}{219.299}{748.416}{n}%
\TextAt{8}{229.457}{744.078}{for}%
\TextAt{8}{247.967}{744.078}{every}%
\TextAt{13}{279.483}{744.078}{n}%
\TextAt{8}{286.471}{744.078}{,}%
\TextAt{8}{294.401}{744.078}{and}%
\TextAt{8}{317.789}{744.078}{hence}%
\TextAt{13}{350.919}{744.078}{a}%
\TextAt{9}{361.45}{744.078}{\Character{61}}%
\TextAt{13}{374.94}{744.078}{b}%
\TextAt{8}{379.917}{744.078}{.}%
\TextAt{8}{390.239}{744.078}{We}%
\TextAt{8}{411.009}{744.078}{now}%
\TextAt{8}{436.022}{744.078}{identify}%
\TextAt{18}{479.243}{744.078}{Q}%
\TextAt{8}{493.063}{744.078}{with}%
\TextAt{8}{520.348}{744.078}{its}%
\TextAt{8}{79.226}{727.032}{image}%
\TextAt{8}{112.875}{727.032}{in}%
\TextAt{18}{126.356}{727.032}{Q}%
\TextAt{21}{135.654}{725.238}{p}%
\TextAt{8}{140.415}{727.032}{,}%
\TextAt{8}{147.427}{727.032}{so}%
\TextAt{8}{161.625}{727.032}{that}%
\TextAt{8}{186.821}{727.032}{we}%
\TextAt{8}{203.876}{727.032}{may}%
\TextAt{8}{229.073}{727.032}{write}%
\TextAt{18}{258.813}{727.032}{Q}%
\NativeMathAt{271.432}{727.032}{11.955}{\subseteq}%
\TextAt{18}{284.051}{727.032}{Q}%
\TextAt{21}{293.35}{725.238}{p}%
\TextAt{8}{301.836}{727.032}{and}%
\TextAt{18}{324.426}{727.032}{Z}%
\NativeMathAt{335.717}{727.032}{11.955}{\subseteq}%
\TextAt{18}{348.337}{727.032}{Z}%
\TextAt{21}{356.307}{725.238}{p}%
\TextAt{8}{361.068}{727.032}{.}%
\TextAt{8}{369.459}{727.032}{Thus,}%
\TextAt{8}{402.222}{727.032}{for}%
\TextAt{8}{419.934}{727.032}{every}%
\TextAt{8}{450.654}{727.032}{rational}%
\TextAt{8}{494.071}{727.032}{number}%
\TextAt{13}{79.226}{709.986}{f}%
\NativeMathAt{89.593}{709.986}{11.955}{\in}%
\TextAt{18}{100.884}{709.986}{Q}%
\TextAt{8}{110.183}{709.986}{,}%
\TextAt{8}{117.332}{709.986}{we}%
\TextAt{8}{134.559}{709.986}{obtain}%
\TextAt{8}{170.987}{709.986}{an}%
\TextAt{8}{187.246}{709.986}{identity}%
\TextAt{13}{268.114}{686.98}{f}%
\TextAt{9}{278.481}{686.98}{\Character{61}}%
\NativeMathAt{299.346}{701.924}{7.97}{\infty}%
\TextAt{20}{294.946}{698.337}{\Character{88}}%
\TextAt{21}{290.906}{673.231}{\Character{23}}%
\TextAt{17}{295.59}{673.231}{\Character{61}}%
\TextAt{3}{302.177}{673.231}{\Character{0}}%
\TextAt{21}{308.763}{673.231}{m}%
\TextAt{13}{318.246}{686.98}{a}%
\TextAt{21}{324.391}{685.186}{\Character{23}}%
\TextAt{13}{329.573}{686.98}{p}%
\TextAt{21}{335.448}{691.916}{\Character{23}}%
\TextAt{13}{340.63}{686.98}{\Character{58}}%
\TextAt{8}{79.226}{655.317}{This}%
\TextAt{8}{105.862}{655.317}{establishes}%
\TextAt{8}{163.845}{655.317}{the}%
\TextAt{8}{183.918}{655.317}{arithmetic}%
\TextAt{8}{240.418}{655.317}{analogue}%
\TextAt{8}{289.123}{655.317}{of}%
\TextAt{8}{302.369}{655.317}{the}%
\TextAt{8}{322.442}{655.317}{function-theoretic}%
\TextAt{8}{416.984}{655.317}{power}%
\TextAt{8}{451.044}{655.317}{series}%
\TextAt{8}{482.295}{655.317}{expansion}%
\TextAt{8}{79.226}{638.271}{for}%
\TextAt{8}{97.111}{638.271}{which}%
\TextAt{8}{130.597}{638.271}{we}%
\TextAt{8}{147.825}{638.271}{were}%
\TextAt{8}{174.808}{638.271}{looking.}%
\TextAt{12}{97.159}{618.236}{Examples:}%
\TextAt{8}{162.43}{618.236}{a)}%
\TextAt{16}{176.74}{618.236}{\Character{0}}%
\TextAt{9}{186.039}{618.236}{1}%
\TextAt{9}{195.197}{618.236}{\Character{61}}%
\TextAt{9}{207.618}{618.236}{\Character{40}}%
\TextAt{13}{212.191}{618.236}{p}%
\TextAt{16}{220.722}{618.236}{\Character{0}}%
\TextAt{9}{232.678}{618.236}{1\Character{41}}%
\TextAt{9}{245.721}{618.236}{\Character{43}}%
\TextAt{9}{257.473}{618.236}{\Character{40}}%
\TextAt{13}{262.053}{618.236}{p}%
\TextAt{16}{270.585}{618.236}{\Character{0}}%
\TextAt{9}{282.54}{618.236}{1\Character{41}}%
\TextAt{13}{292.945}{618.236}{p}%
\TextAt{9}{301.477}{618.236}{\Character{43}}%
\TextAt{9}{313.229}{618.236}{\Character{40}}%
\TextAt{13}{317.791}{618.236}{p}%
\TextAt{16}{326.323}{618.236}{\Character{0}}%
\TextAt{9}{338.278}{618.236}{1\Character{41}}%
\TextAt{13}{348.683}{618.236}{p}%
\TextAt{17}{354.558}{622.575}{2}%
\TextAt{9}{361.947}{618.236}{\Character{43}}%
\TextAt{16}{373.709}{618.236}{\Character{1}}%
\TextAt{16}{379.005}{618.236}{\Character{1}}%
\TextAt{16}{384.313}{618.236}{\Character{1}}%
\TextAt{8}{389.649}{618.236}{.}%
\TextAt{8}{398.097}{618.236}{In}%
\TextAt{8}{412.718}{618.236}{fact,}%
\TextAt{8}{439.056}{618.236}{we}%
\TextAt{8}{456.283}{618.236}{have}%
\TextAt{16}{179.99}{589.235}{\Character{0}}%
\TextAt{9}{189.289}{589.235}{1}%
\TextAt{9}{198.447}{589.235}{\Character{61}}%
\TextAt{9}{210.868}{589.235}{\Character{40}}%
\TextAt{13}{215.44}{589.235}{p}%
\TextAt{16}{223.972}{589.235}{\Character{0}}%
\TextAt{9}{235.927}{589.235}{1\Character{41}}%
\TextAt{9}{248.97}{589.235}{\Character{43}}%
\TextAt{9}{260.722}{589.235}{\Character{40}}%
\TextAt{13}{265.303}{589.235}{p}%
\TextAt{16}{273.835}{589.235}{\Character{0}}%
\TextAt{9}{285.79}{589.235}{1\Character{41}}%
\TextAt{13}{296.195}{589.235}{p}%
\TextAt{9}{304.727}{589.235}{\Character{43}}%
\TextAt{16}{316.488}{589.235}{\Character{1}}%
\TextAt{16}{321.784}{589.235}{\Character{1}}%
\TextAt{16}{327.092}{589.235}{\Character{1}}%
\TextAt{9}{333.093}{589.235}{\Character{43}}%
\TextAt{9}{344.845}{589.235}{\Character{40}}%
\TextAt{13}{349.406}{589.235}{p}%
\TextAt{16}{357.938}{589.235}{\Character{0}}%
\TextAt{9}{369.893}{589.235}{1\Character{41}}%
\TextAt{13}{380.298}{589.235}{p}%
\TextAt{21}{386.174}{594.171}{n}%
\TextAt{3}{391.312}{594.171}{\Character{0}}%
\TextAt{17}{397.898}{594.171}{1}%
\TextAt{16}{405.287}{589.235}{\Character{0}}%
\TextAt{13}{417.242}{589.235}{p}%
\TextAt{21}{423.118}{594.171}{n}%
\TextAt{13}{428.754}{589.235}{\Character{59}}%
\TextAt{8}{79.226}{560.234}{hence}%
\TextAt{16}{111.732}{560.234}{\Character{0}}%
\TextAt{9}{121.031}{560.234}{1}%
\NativeMathAt{130.205}{560.234}{11.955}{\equiv}%
\TextAt{9}{142.824}{560.234}{\Character{40}}%
\TextAt{13}{147.376}{560.234}{p}%
\TextAt{16}{155.908}{560.234}{\Character{0}}%
\TextAt{9}{167.863}{560.234}{1\Character{41}}%
\TextAt{9}{180.906}{560.234}{\Character{43}}%
\TextAt{9}{192.658}{560.234}{\Character{40}}%
\TextAt{13}{197.239}{560.234}{p}%
\TextAt{16}{205.771}{560.234}{\Character{0}}%
\TextAt{9}{217.726}{560.234}{1\Character{41}}%
\TextAt{13}{228.131}{560.234}{p}%
\TextAt{9}{236.663}{560.234}{\Character{43}}%
\TextAt{16}{248.424}{560.234}{\Character{1}}%
\TextAt{16}{253.72}{560.234}{\Character{1}}%
\TextAt{16}{259.028}{560.234}{\Character{1}}%
\TextAt{9}{265.029}{560.234}{\Character{43}}%
\TextAt{9}{276.781}{560.234}{\Character{40}}%
\TextAt{13}{281.342}{560.234}{p}%
\TextAt{16}{289.874}{560.234}{\Character{0}}%
\TextAt{9}{301.829}{560.234}{1\Character{41}}%
\TextAt{13}{312.234}{560.234}{p}%
\TextAt{21}{318.11}{564.572}{n}%
\TextAt{3}{323.248}{564.572}{\Character{0}}%
\TextAt{17}{329.834}{564.572}{1}%
\TextAt{8}{337.887}{560.234}{mod}%
\TextAt{13}{363.325}{560.234}{p}%
\TextAt{21}{369.201}{564.572}{n}%
\TextAt{8}{374.837}{560.234}{.}%
\TextAt{8}{97.159}{537.468}{b)}%
\TextAt{9}{123.554}{545.556}{1}%
\Vector{q 0 G 0 g 0.478 w 41.31 -233.543 m 67.65 -233.543 l S Q}%
\TextAt{9}{113.31}{529.267}{1}%
\TextAt{16}{121.82}{529.267}{\Character{0}}%
\TextAt{13}{133.775}{529.267}{p}%
\TextAt{9}{144.167}{537.468}{\Character{61}}%
\TextAt{9}{156.577}{537.468}{1}%
\TextAt{9}{165.077}{537.468}{\Character{43}}%
\TextAt{13}{176.863}{537.468}{p}%
\TextAt{9}{185.395}{537.468}{\Character{43}}%
\TextAt{13}{197.156}{537.468}{p}%
\TextAt{17}{203.031}{542.404}{2}%
\TextAt{9}{210.42}{537.468}{\Character{43}}%
\TextAt{16}{222.182}{537.468}{\Character{1}}%
\TextAt{16}{227.478}{537.468}{\Character{1}}%
\TextAt{16}{232.786}{537.468}{\Character{1}}%
\TextAt{8}{238.122}{537.468}{.}%
\TextAt{8}{246.57}{537.468}{In}%
\TextAt{8}{261.191}{537.468}{fact,}%
\TextAt{9}{212.702}{503.919}{1}%
\TextAt{9}{221.86}{503.919}{\Character{61}}%
\TextAt{9}{234.281}{503.919}{\Character{40}1}%
\TextAt{9}{247.324}{503.919}{\Character{43}}%
\TextAt{13}{259.125}{503.919}{p}%
\TextAt{9}{267.656}{503.919}{\Character{43}}%
\TextAt{16}{279.418}{503.919}{\Character{1}}%
\TextAt{16}{284.714}{503.919}{\Character{1}}%
\TextAt{16}{290.022}{503.919}{\Character{1}}%
\TextAt{9}{296.022}{503.919}{\Character{43}}%
\TextAt{13}{307.783}{503.919}{p}%
\TextAt{21}{313.659}{508.855}{n}%
\TextAt{3}{318.797}{508.855}{\Character{0}}%
\TextAt{17}{325.383}{508.855}{1}%
\TextAt{9}{330.116}{503.919}{\Character{41}\Character{40}1}%
\TextAt{16}{347.73}{503.919}{\Character{0}}%
\TextAt{13}{359.685}{503.919}{p}%
\TextAt{9}{365.56}{503.919}{\Character{41}}%
\TextAt{9}{372.757}{503.919}{\Character{43}}%
\TextAt{13}{384.53}{503.919}{p}%
\TextAt{21}{390.406}{508.855}{n}%
\TextAt{13}{396.042}{503.919}{\Character{59}}%
\TextAt{8}{79.226}{474.917}{hence}%
\TextAt{9}{218.508}{465.959}{1}%
\Vector{q 0 G 0 g 0.478 w 136.264 -313.14 m 162.604 -313.14 l S Q}%
\TextAt{9}{208.264}{449.671}{1}%
\TextAt{16}{216.774}{449.671}{\Character{0}}%
\TextAt{13}{228.729}{449.671}{p}%
\NativeMathAt{239.121}{457.872}{11.955}{\equiv}%
\TextAt{9}{251.74}{457.872}{1}%
\TextAt{9}{260.24}{457.872}{\Character{43}}%
\TextAt{13}{272.011}{457.872}{p}%
\TextAt{9}{280.543}{457.872}{\Character{43}}%
\TextAt{16}{292.304}{457.872}{\Character{1}}%
\TextAt{16}{297.6}{457.872}{\Character{1}}%
\TextAt{16}{302.908}{457.872}{\Character{1}}%
\TextAt{9}{308.908}{457.872}{\Character{43}}%
\TextAt{13}{320.67}{457.872}{p}%
\TextAt{21}{326.545}{462.808}{n}%
\TextAt{3}{331.683}{462.808}{\Character{0}}%
\TextAt{17}{338.27}{462.808}{1}%
\TextAt{9}{354.957}{457.872}{\Character{40}}%
\TextAt{8}{359.509}{457.872}{mod}%
\TextAt{13}{385.611}{457.872}{p}%
\TextAt{21}{391.486}{462.808}{n}%
\TextAt{9}{397.123}{457.872}{\Character{41}}%
\TextAt{13}{401.675}{457.872}{\Character{58}}%
\TextAt{8}{79.226}{431.458}{One}%
\TextAt{8}{104.039}{431.458}{can}%
\TextAt{8}{125.6}{431.458}{define}%
\TextAt{8}{159.51}{431.458}{addition}%
\TextAt{8}{205.794}{431.458}{and}%
\TextAt{8}{228.658}{431.458}{multiplication}%
\TextAt{8}{303.886}{431.458}{of}%
\TextAt{13}{317.317}{431.458}{p}%
\TextAt{8}{323.193}{431.458}{-adic}%
\TextAt{8}{351.903}{431.458}{numbers}%
\TextAt{8}{399.215}{431.458}{which}%
\TextAt{8}{432.803}{431.458}{turn}%
\TextAt{18}{458.918}{431.458}{Z}%
\TextAt{21}{466.889}{429.665}{p}%
\TextAt{8}{475.648}{431.458}{into}%
\TextAt{8}{499.493}{431.458}{a}%
\TextAt{8}{509.349}{431.458}{ring,}%
\TextAt{8}{79.226}{414.412}{and}%
\TextAt{18}{101.41}{414.412}{Q}%
\TextAt{21}{110.708}{412.619}{p}%
\TextAt{8}{118.787}{414.412}{into}%
\TextAt{8}{141.951}{414.412}{its}%
\TextAt{8}{157.69}{414.412}{field}%
\TextAt{8}{182.468}{414.412}{of}%
\TextAt{8}{195.219}{414.412}{fractions.}%
\TextAt{8}{247.445}{414.412}{However,}%
\TextAt{8}{297.385}{414.412}{the}%
\TextAt{8}{316.962}{414.412}{direct}%
\TextAt{8}{349.546}{414.412}{approach,}%
\TextAt{8}{402.749}{414.412}{defining}%
\TextAt{8}{446.392}{414.412}{sum}%
\TextAt{8}{470.584}{414.412}{and}%
\TextAt{8}{492.767}{414.412}{product}%
\TextAt{8}{79.226}{397.366}{via}%
\TextAt{8}{98.115}{397.366}{the}%
\TextAt{8}{117.971}{397.366}{usual}%
\TextAt{8}{148.301}{397.366}{carry-over}%
\TextAt{8}{203.294}{397.366}{rules}%
\TextAt{8}{231.017}{397.366}{for}%
\TextAt{8}{248.602}{397.366}{digits,}%
\TextAt{8}{283.547}{397.366}{as}%
\TextAt{8}{297.617}{397.366}{one}%
\TextAt{8}{318.777}{397.366}{does}%
\TextAt{8}{344.875}{397.366}{it}%
\TextAt{8}{356.279}{397.366}{when}%
\TextAt{8}{386.537}{397.366}{dealing}%
\TextAt{8}{426.562}{397.366}{with}%
\TextAt{8}{452.922}{397.366}{real}%
\TextAt{8}{475.385}{397.366}{numbers}%
\TextAt{8}{522.297}{397.366}{as}%
\TextAt{8}{79.226}{380.32}{decimal}%
\TextAt{8}{122.755}{380.32}{fractions,}%
\TextAt{8}{174.637}{380.32}{leads}%
\TextAt{8}{204.572}{380.32}{into}%
\TextAt{8}{228.924}{380.32}{complications.}%
\TextAt{8}{308.863}{380.32}{They}%
\TextAt{8}{339.707}{380.32}{disappear}%
\TextAt{8}{393.386}{380.32}{once}%
\TextAt{8}{420.655}{380.32}{we}%
\TextAt{8}{438.492}{380.32}{use}%
\TextAt{8}{459.317}{380.32}{another}%
\TextAt{8}{502.858}{380.32}{repre-}%
\TextAt{8}{79.226}{363.274}{sentation}%
\TextAt{8}{130.041}{363.274}{of}%
\TextAt{8}{143.71}{363.274}{the}%
\TextAt{13}{164.206}{363.274}{p}%
\TextAt{8}{170.082}{363.274}{-adic}%
\TextAt{8}{199.03}{363.274}{numbers}%
\TextAt{13}{246.581}{363.274}{f}%
\TextAt{9}{257.527}{363.274}{\Character{61}}%
\TextAt{20}{270.531}{372.24}{\Character{80}}%
\NativeMathAt{283.151}{369.086}{7.97}{\infty}%
\TextAt{21}{283.151}{359.787}{\Character{23}}%
\TextAt{17}{287.835}{359.787}{\Character{61}0}%
\TextAt{13}{301.146}{363.274}{a}%
\TextAt{21}{307.291}{361.481}{\Character{23}}%
\TextAt{13}{312.473}{363.274}{p}%
\TextAt{21}{318.348}{367.612}{\Character{23}}%
\TextAt{8}{323.53}{363.274}{,}%
\TextAt{8}{331.104}{363.274}{viewing}%
\TextAt{8}{374.04}{363.274}{them}%
\TextAt{8}{404.291}{363.274}{not}%
\TextAt{8}{425.445}{363.274}{as}%
\TextAt{8}{440.155}{363.274}{sequences}%
\TextAt{8}{493.611}{363.274}{of}%
\TextAt{8}{507.281}{363.274}{sums}%
\TextAt{8}{79.226}{346.228}{of}%
\TextAt{8}{92.556}{346.228}{integers}%
\TextAt{13}{258.81}{321.505}{s}%
\TextAt{21}{264.324}{319.712}{n}%
\TextAt{9}{273.281}{321.505}{\Character{61}}%
\TextAt{21}{286.362}{336.449}{n}%
\TextAt{3}{291.5}{336.449}{\Character{0}}%
\TextAt{17}{298.086}{336.449}{1}%
\TextAt{20}{285.707}{332.863}{\Character{88}}%
\TextAt{21}{286.589}{307.757}{\Character{23}}%
\TextAt{17}{291.273}{307.757}{\Character{61}0}%
\TextAt{13}{304.968}{321.505}{a}%
\TextAt{21}{311.113}{319.712}{\Character{23}}%
\TextAt{13}{316.295}{321.505}{p}%
\TextAt{21}{322.17}{326.441}{\Character{23}}%
\NativeMathAt{330.673}{321.505}{11.955}{\in}%
\TextAt{18}{341.964}{321.505}{Z}%
\TextAt{13}{349.934}{321.505}{\Character{59}}%
\TextAt{8}{79.226}{290.673}{but}%
\TextAt{8}{100.686}{290.673}{rather}%
\TextAt{8}{135.81}{290.673}{as}%
\TextAt{8}{150.18}{290.673}{sequences}%
\TextAt{8}{203.297}{290.673}{of}%
\TextAt{8}{216.627}{290.673}{residue}%
\TextAt{8}{256.354}{290.673}{classes}%
\TextAt{9}{241.736}{261.671}{\Character{22}}%
\TextAt{13}{241.255}{261.671}{s}%
\TextAt{21}{246.769}{259.878}{n}%
\TextAt{9}{255.726}{261.671}{\Character{61}}%
\TextAt{13}{268.152}{261.671}{s}%
\TextAt{21}{273.666}{259.878}{n}%
\TextAt{8}{282.623}{261.671}{mod}%
\TextAt{13}{308.061}{261.671}{p}%
\TextAt{21}{313.936}{266.608}{n}%
\NativeMathAt{322.893}{261.671}{11.955}{\in}%
\TextAt{18}{334.184}{261.671}{Z}%
\TextAt{9}{342.154}{261.671}{\Character{47}}%
\TextAt{13}{348.007}{261.671}{p}%
\TextAt{21}{353.882}{266.608}{n}%
\TextAt{18}{359.519}{261.671}{Z}%
\TextAt{13}{367.489}{261.671}{\Character{58}}%
\TextAt{8}{79.226}{232.67}{The}%
\TextAt{8}{104.011}{232.67}{terms}%
\TextAt{8}{137.321}{232.67}{of}%
\TextAt{8}{151.382}{232.67}{such}%
\TextAt{8}{178.511}{232.67}{a}%
\TextAt{8}{188.997}{232.67}{sequence}%
\TextAt{8}{238.231}{232.67}{lie}%
\TextAt{8}{254.564}{232.67}{in}%
\TextAt{8}{268.948}{232.67}{different}%
\TextAt{8}{315.85}{232.67}{rings}%
\TextAt{18}{345.262}{232.67}{Z}%
\TextAt{9}{353.232}{232.67}{\Character{47}}%
\TextAt{13}{359.085}{232.67}{p}%
\TextAt{21}{364.961}{237.009}{n}%
\TextAt{18}{370.597}{232.67}{Z}%
\TextAt{8}{378.567}{232.67}{,}%
\TextAt{8}{386.631}{232.67}{but}%
\TextAt{8}{408.822}{232.67}{these}%
\TextAt{8}{439.525}{232.67}{are}%
\TextAt{8}{459.767}{232.67}{related}%
\TextAt{8}{499.52}{232.67}{by}%
\TextAt{8}{516.51}{232.67}{the}%
\TextAt{8}{79.226}{215.624}{canonical}%
\TextAt{8}{130.609}{215.624}{projections}%
\TextAt{18}{219.31}{198.578}{Z}%
\TextAt{9}{227.281}{198.578}{\Character{47}}%
\TextAt{13}{233.134}{198.578}{p}%
\TextAt{18}{239.009}{198.578}{Z}%
\TextAt{21}{254.299}{205.4}{\Character{21}}%
\TextAt{24}{259.239}{204.293}{1}%
\NativeArrowAt{250.3}{198.578}{11.955}{15.673}{\longleftarrow}%
\TextAt{18}{269.298}{198.578}{Z}%
\TextAt{9}{277.268}{198.578}{\Character{47}}%
\TextAt{13}{283.121}{198.578}{p}%
\TextAt{17}{288.996}{203.515}{2}%
\TextAt{18}{293.729}{198.578}{Z}%
\TextAt{21}{309.018}{205.4}{\Character{21}}%
\TextAt{24}{313.958}{204.293}{2}%
\NativeArrowAt{305.019}{198.578}{11.955}{15.673}{\longleftarrow}%
\TextAt{18}{324.018}{198.578}{Z}%
\TextAt{9}{331.988}{198.578}{\Character{47}}%
\TextAt{13}{337.841}{198.578}{p}%
\TextAt{17}{343.716}{203.515}{3}%
\TextAt{18}{348.448}{198.578}{Z}%
\TextAt{21}{363.738}{205.4}{\Character{21}}%
\TextAt{24}{368.678}{204.293}{3}%
\NativeArrowAt{359.739}{198.578}{11.955}{15.673}{\longleftarrow}%
\TextAt{16}{378.724}{198.578}{\Character{1}}%
\TextAt{16}{384.032}{198.578}{\Character{1}}%
\TextAt{16}{389.328}{198.578}{\Character{1}}%
\TextAt{8}{79.226}{175.057}{and}%
\TextAt{8}{101.989}{175.057}{we}%
\TextAt{8}{119.216}{175.057}{find}%
\TextAt{13}{269.153}{158.011}{\Character{21}}%
\TextAt{21}{275.982}{156.217}{n}%
\TextAt{9}{281.618}{158.011}{\Character{40}}%
\TextAt{9}{286.639}{158.011}{\Character{22}}%
\TextAt{13}{286.171}{158.011}{s}%
\TextAt{21}{291.685}{156.217}{n}%
\TextAt{17}{296.823}{156.217}{\Character{43}1}%
\TextAt{9}{308.142}{158.011}{\Character{41}}%
\TextAt{9}{315.997}{158.011}{\Character{61}}%
\TextAt{9}{328.896}{158.011}{\Character{22}}%
\TextAt{13}{328.44}{158.011}{s}%
\TextAt{21}{333.954}{156.217}{n}%
\TextAt{13}{339.591}{158.011}{\Character{58}}%
\TextAt{8}{79.226}{134.489}{In}%
\TextAt{8}{93.847}{134.489}{the}%
\TextAt{8}{114.004}{134.489}{direct}%
\TextAt{8}{147.167}{134.489}{product}%
\NativeMathAt{214.339}{114.591}{7.97}{\infty}%
\TextAt{20}{210.935}{111.005}{\Character{89}}%
\TextAt{21}{210.594}{85.899}{n}%
\TextAt{17}{215.732}{85.899}{\Character{61}1}%
\TextAt{18}{228.545}{99.647}{Z}%
\TextAt{9}{236.515}{99.647}{\Character{47}}%
\TextAt{13}{242.368}{99.647}{p}%
\TextAt{21}{248.243}{104.583}{n}%
\TextAt{18}{253.88}{99.647}{Z}%
\TextAt{9}{265.171}{99.647}{\Character{61}}%
\TextAt{16}{277.596}{99.647}{\Character{102}}%
\TextAt{9}{283.574}{99.647}{\Character{40}}%
\TextAt{13}{288.126}{99.647}{x}%
\TextAt{21}{294.778}{97.854}{n}%
\TextAt{9}{300.415}{99.647}{\Character{41}}%
\TextAt{21}{304.967}{97.854}{n}%
\NativeMathAt{310.105}{97.854}{7.97}{\in}%
\TextAt{25}{315.751}{97.854}{N}%
\TextAt{16}{325.326}{99.647}{\Character{106}}%
\TextAt{13}{331.968}{99.647}{x}%
\TextAt{21}{338.62}{97.854}{n}%
\NativeMathAt{347.577}{99.647}{11.955}{\in}%
\TextAt{18}{358.868}{99.647}{Z}%
\TextAt{9}{366.838}{99.647}{\Character{47}}%
\TextAt{13}{372.691}{99.647}{p}%
\TextAt{21}{378.566}{104.583}{n}%
\TextAt{18}{384.202}{99.647}{Z}%
\TextAt{16}{392.173}{99.647}{\Character{103}}%
\TextAt{13}{398.15}{99.647}{\Character{59}}%
\TextAt{8}{300.14}{43.033}{94}%
\Destination{78.23}{809.815}{nav.page.94}%
\Destination{79.36}{776.155}{nav.print.102}%
}

% Page 107 (body page 95)
\Page{%
\TextAt{2}{79.226}{796.943}{II.}%
\TextAt{2}{94.073}{796.943}{The}%
\TextAt{2}{116.492}{796.943}{Theory}%
\TextAt{2}{154.401}{796.943}{of}%
\TextAt{2}{166.826}{796.943}{Valuations}%
\Vector{q 0 G 0 g 0.299 w 7.226 18.729 m 460.769 18.729 l S Q}%
\TextAt{8}{79.226}{761.124}{we}%
\TextAt{8}{96.454}{761.124}{now}%
\TextAt{8}{120.842}{761.124}{consider}%
\TextAt{8}{166.427}{761.124}{all}%
\TextAt{8}{182.686}{761.124}{elements}%
\TextAt{9}{230.543}{761.124}{\Character{40}}%
\TextAt{13}{235.095}{761.124}{x}%
\TextAt{21}{241.747}{759.33}{n}%
\TextAt{9}{247.384}{761.124}{\Character{41}}%
\TextAt{21}{251.936}{759.33}{n}%
\NativeMathAt{257.074}{759.33}{7.97}{\in}%
\TextAt{25}{262.72}{759.33}{N}%
\TextAt{8}{272.871}{761.124}{with}%
\TextAt{8}{299.531}{761.124}{the}%
\TextAt{8}{319.688}{761.124}{property}%
\TextAt{8}{367.497}{761.124}{that}%
\TextAt{13}{208.658}{732.123}{\Character{21}}%
\TextAt{21}{215.486}{730.329}{n}%
\TextAt{9}{221.123}{732.123}{\Character{40}}%
\TextAt{13}{225.675}{732.123}{x}%
\TextAt{21}{232.327}{730.329}{n}%
\TextAt{17}{237.465}{730.329}{\Character{43}1}%
\TextAt{9}{248.784}{732.123}{\Character{41}}%
\TextAt{9}{256.639}{732.123}{\Character{61}}%
\TextAt{13}{269.083}{732.123}{x}%
\TextAt{21}{275.735}{730.329}{n}%
\TextAt{8}{305.281}{732.123}{for}%
\TextAt{8}{323.166}{732.123}{all}%
\TextAt{13}{339.425}{732.123}{n}%
\TextAt{9}{349.734}{732.123}{\Character{61}}%
\TextAt{9}{362.144}{732.123}{1}%
\TextAt{13}{368.012}{732.123}{\Character{59}}%
\TextAt{9}{373.257}{732.123}{2}%
\TextAt{13}{379.11}{732.123}{\Character{59}}%
\TextAt{13}{384.334}{732.123}{\Character{58}}%
\TextAt{13}{389.571}{732.123}{\Character{58}}%
\TextAt{13}{394.795}{732.123}{\Character{58}}%
\TextAt{13}{400.032}{732.123}{\Character{58}}%
\TextAt{8}{79.226}{703.121}{This}%
\TextAt{8}{105.58}{703.121}{set}%
\TextAt{8}{123.482}{703.121}{is}%
\TextAt{8}{134.88}{703.121}{called}%
\TextAt{8}{167.677}{703.121}{the}%
\TextAt{12}{187.468}{703.121}{projective}%
\TextAt{12}{249.668}{703.121}{limit}%
\TextAt{8}{280.876}{703.121}{of}%
\TextAt{8}{293.84}{703.121}{the}%
\TextAt{8}{313.63}{703.121}{rings}%
\TextAt{18}{341.945}{703.121}{Z}%
\TextAt{9}{349.915}{703.121}{\Character{47}}%
\TextAt{13}{355.768}{703.121}{p}%
\TextAt{21}{361.643}{707.46}{n}%
\TextAt{18}{367.279}{703.121}{Z}%
\TextAt{8}{378.781}{703.121}{and}%
\TextAt{8}{401.178}{703.121}{is}%
\TextAt{8}{412.576}{703.121}{denoted}%
\TextAt{8}{456.432}{703.121}{by}%
\NativeMathAt{472.325}{703.121}{11.955}{\varprojlim}%
\TextAt{21}{488.584}{697.682}{n}%
\TextAt{18}{496.213}{703.121}{Z}%
\TextAt{9}{504.183}{703.121}{\Character{47}}%
\TextAt{13}{510.036}{703.121}{p}%
\TextAt{21}{515.911}{707.46}{n}%
\TextAt{18}{521.548}{703.121}{Z}%
\TextAt{8}{529.518}{703.121}{.}%
\TextAt{8}{79.226}{686.075}{In}%
\TextAt{8}{93.847}{686.075}{other}%
\TextAt{8}{124.417}{686.075}{words,}%
\TextAt{8}{161.227}{686.075}{we}%
\TextAt{8}{178.454}{686.075}{have}%
\NativeMathAt{132.369}{649.895}{11.955}{\varprojlim}%
\TextAt{21}{137.93}{637.781}{n}%
\TextAt{18}{150.621}{649.895}{Z}%
\TextAt{9}{158.591}{649.895}{\Character{47}}%
\TextAt{13}{164.444}{649.895}{p}%
\TextAt{21}{170.319}{654.831}{n}%
\TextAt{18}{175.956}{649.895}{Z}%
\TextAt{9}{187.247}{649.895}{\Character{61}}%
\TextAt{20}{199.672}{670.338}{\Character{40}}%
\TextAt{9}{209.303}{649.895}{\Character{40}}%
\TextAt{13}{213.855}{649.895}{x}%
\TextAt{21}{220.507}{648.102}{n}%
\TextAt{9}{226.143}{649.895}{\Character{41}}%
\TextAt{21}{230.696}{648.102}{n}%
\NativeMathAt{235.834}{648.102}{7.97}{\in}%
\TextAt{25}{241.479}{648.102}{N}%
\NativeMathAt{251.055}{649.895}{11.955}{\in}%
\NativeMathAt{266.091}{664.839}{7.97}{\infty}%
\TextAt{20}{262.687}{661.252}{\Character{89}}%
\TextAt{21}{262.346}{636.146}{n}%
\TextAt{17}{267.484}{636.146}{\Character{61}1}%
\TextAt{18}{280.297}{649.895}{Z}%
\TextAt{9}{288.267}{649.895}{\Character{47}}%
\TextAt{13}{294.12}{649.895}{p}%
\TextAt{21}{299.995}{654.831}{n}%
\TextAt{18}{305.632}{649.895}{Z}%
\TextAt{20}{316.923}{670.817}{\Character{12}}%
\TextAt{20}{316.923}{663.644}{\Character{12}}%
\TextAt{20}{316.923}{656.47}{\Character{12}}%
\TextAt{20}{316.923}{649.297}{\Character{12}}%
\TextAt{20}{316.923}{642.124}{\Character{12}}%
\TextAt{13}{324.228}{649.895}{\Character{21}}%
\TextAt{21}{331.057}{648.102}{n}%
\TextAt{9}{336.693}{649.895}{\Character{40}}%
\TextAt{13}{341.246}{649.895}{x}%
\TextAt{21}{347.898}{648.102}{n}%
\TextAt{17}{353.036}{648.102}{\Character{43}1}%
\TextAt{9}{364.355}{649.895}{\Character{41}}%
\TextAt{9}{372.21}{649.895}{\Character{61}}%
\TextAt{13}{384.653}{649.895}{x}%
\TextAt{21}{391.305}{648.102}{n}%
\TextAt{13}{396.942}{649.895}{\Character{59}}%
\TextAt{13}{406.064}{649.895}{n}%
\TextAt{9}{416.392}{649.895}{\Character{61}}%
\TextAt{9}{428.802}{649.895}{1}%
\TextAt{13}{434.67}{649.895}{\Character{59}}%
\TextAt{9}{439.914}{649.895}{2}%
\TextAt{13}{445.767}{649.895}{\Character{59}}%
\TextAt{13}{450.991}{649.895}{\Character{58}}%
\TextAt{13}{456.228}{649.895}{\Character{58}}%
\TextAt{13}{461.464}{649.895}{\Character{58}}%
\TextAt{20}{464.751}{670.338}{\Character{41}}%
\TextAt{13}{476.375}{649.895}{\Character{58}}%
\TextAt{8}{79.226}{613.583}{The}%
\TextAt{8}{102.883}{613.583}{modified}%
\TextAt{8}{150.283}{613.583}{representation}%
\TextAt{8}{226.376}{613.583}{of}%
\TextAt{8}{239.309}{613.583}{the}%
\TextAt{13}{259.069}{613.583}{p}%
\TextAt{8}{264.944}{613.583}{-adic}%
\TextAt{8}{293.156}{613.583}{numbers}%
\TextAt{8}{339.97}{613.583}{alluded}%
\TextAt{8}{380.543}{613.583}{to}%
\TextAt{8}{394.457}{613.583}{above}%
\TextAt{8}{427.236}{613.583}{now}%
\TextAt{8}{451.228}{613.583}{follows}%
\TextAt{8}{489.267}{613.583}{from}%
\TextAt{8}{516.51}{613.583}{the}%
\TextAt{12}{79.226}{587.571}{Proposition}%
\TextAt{12}{151.937}{587.571}{(1.3).}%
\TextAt{8}{187.253}{587.571}{Associating}%
\TextAt{8}{249.814}{587.571}{to}%
\TextAt{8}{264.125}{587.571}{every}%
\TextAt{13}{295.017}{587.571}{p}%
\TextAt{8}{300.892}{587.571}{-adic}%
\TextAt{8}{329.501}{587.571}{integer}%
\TextAt{13}{273.779}{552.047}{f}%
\TextAt{9}{284.146}{552.047}{\Character{61}}%
\NativeMathAt{300.972}{566.991}{7.97}{\infty}%
\TextAt{20}{296.572}{563.405}{\Character{88}}%
\TextAt{21}{297.454}{538.299}{\Character{23}}%
\TextAt{17}{302.137}{538.299}{\Character{61}0}%
\TextAt{13}{315.833}{552.047}{a}%
\TextAt{21}{321.978}{550.254}{\Character{23}}%
\TextAt{13}{327.16}{552.047}{p}%
\TextAt{21}{333.035}{556.983}{\Character{23}}%
\TextAt{8}{79.226}{515.185}{the}%
\TextAt{8}{99.383}{515.185}{sequence}%
\TextAt{9}{147.885}{515.185}{\Character{40}}%
\TextAt{9}{152.906}{515.185}{\Character{22}}%
\TextAt{13}{152.437}{515.185}{s}%
\TextAt{21}{157.951}{513.392}{n}%
\TextAt{9}{163.587}{515.185}{\Character{41}}%
\TextAt{21}{168.14}{513.392}{n}%
\NativeMathAt{173.278}{513.392}{7.97}{\in}%
\TextAt{25}{178.923}{513.392}{N}%
\TextAt{8}{189.075}{515.185}{of}%
\TextAt{8}{202.405}{515.185}{residue}%
\TextAt{8}{242.132}{515.185}{classes}%
\TextAt{9}{226.489}{477.969}{\Character{22}}%
\TextAt{13}{226.008}{477.969}{s}%
\TextAt{21}{231.522}{476.176}{n}%
\TextAt{9}{240.479}{477.969}{\Character{61}}%
\TextAt{21}{253.559}{492.913}{n}%
\TextAt{3}{258.697}{492.913}{\Character{0}}%
\TextAt{17}{265.284}{492.913}{1}%
\TextAt{20}{252.904}{489.327}{\Character{88}}%
\TextAt{21}{253.786}{464.221}{\Character{23}}%
\TextAt{17}{258.47}{464.221}{\Character{61}0}%
\TextAt{13}{272.165}{477.969}{a}%
\TextAt{21}{278.31}{476.176}{\Character{23}}%
\TextAt{13}{283.492}{477.969}{p}%
\TextAt{21}{289.368}{482.906}{\Character{23}}%
\TextAt{8}{297.87}{477.969}{mod}%
\TextAt{13}{323.308}{477.969}{p}%
\TextAt{21}{329.183}{482.906}{n}%
\NativeMathAt{338.141}{477.969}{11.955}{\in}%
\TextAt{18}{349.432}{477.969}{Z}%
\TextAt{9}{357.402}{477.969}{\Character{47}}%
\TextAt{13}{363.255}{477.969}{p}%
\TextAt{21}{369.13}{482.906}{n}%
\TextAt{18}{374.766}{477.969}{Z}%
\TextAt{13}{382.736}{477.969}{\Character{59}}%
\TextAt{8}{79.226}{441.777}{yields}%
\TextAt{8}{112.127}{441.777}{a}%
\TextAt{8}{121.882}{441.777}{bijection}%
\TextAt{18}{262.321}{424.731}{Z}%
\TextAt{21}{270.291}{422.938}{p}%
\NativeMathAt{280.96}{431.508}{7.97}{\sim}%
\NativeArrowAt{278.373}{424.731}{11.955}{13.175}{\longrightarrow}%
\NativeMathAt{294.867}{424.731}{11.955}{\varprojlim}%
\TextAt{21}{300.427}{412.617}{n}%
\TextAt{18}{313.118}{424.731}{Z}%
\TextAt{9}{321.088}{424.731}{\Character{47}}%
\TextAt{13}{326.941}{424.731}{p}%
\TextAt{21}{332.817}{429.667}{n}%
\TextAt{18}{338.453}{424.731}{Z}%
\TextAt{13}{346.423}{424.731}{\Character{58}}%
\TextAt{8}{97.159}{386.016}{The}%
\TextAt{8}{122.765}{386.016}{proof}%
\TextAt{8}{154.887}{386.016}{is}%
\TextAt{8}{168.204}{386.016}{an}%
\TextAt{8}{186.016}{386.016}{immediate}%
\TextAt{8}{244.798}{386.016}{consequence}%
\TextAt{8}{312.415}{386.016}{of}%
\TextAt{8}{327.297}{386.016}{proposition}%
\TextAt{8}{390.383}{386.016}{(1.2).}%
\TextAt{8}{427.567}{386.016}{The}%
\TextAt{8}{453.173}{386.016}{projective}%
\TextAt{8}{508.704}{386.016}{limit}%
\NativeMathAt{79.226}{368.97}{11.955}{\varprojlim}%
\TextAt{21}{95.485}{363.531}{n}%
\TextAt{18}{103.114}{368.97}{Z}%
\TextAt{9}{111.084}{368.97}{\Character{47}}%
\TextAt{13}{116.937}{368.97}{p}%
\TextAt{21}{122.812}{373.309}{n}%
\TextAt{18}{128.449}{368.97}{Z}%
\TextAt{8}{139.497}{368.97}{offers}%
\TextAt{8}{169.63}{368.97}{the}%
\TextAt{8}{188.967}{368.97}{advantage}%
\TextAt{8}{243.452}{368.97}{of}%
\TextAt{8}{255.963}{368.97}{being}%
\TextAt{8}{286.681}{368.97}{clearly}%
\TextAt{8}{323.258}{368.97}{a}%
\TextAt{8}{332.194}{368.97}{ring.}%
\TextAt{8}{360.538}{368.97}{In}%
\TextAt{8}{374.339}{368.97}{fact,}%
\TextAt{8}{400.021}{368.97}{it}%
\TextAt{8}{410.906}{368.97}{is}%
\TextAt{8}{421.851}{368.97}{a}%
\TextAt{8}{430.787}{368.97}{subring}%
\TextAt{8}{471.655}{368.97}{of}%
\TextAt{8}{484.166}{368.97}{the}%
\TextAt{8}{503.503}{368.97}{direct}%
\TextAt{8}{79.226}{351.924}{product}%
\TextAt{20}{123.413}{360.891}{\Character{81}}%
\NativeMathAt{134.704}{357.736}{7.97}{\infty}%
\TextAt{21}{134.704}{348.437}{n}%
\TextAt{17}{139.843}{348.437}{\Character{61}1}%
\TextAt{18}{153.154}{351.924}{Z}%
\TextAt{9}{161.124}{351.924}{\Character{47}}%
\TextAt{13}{166.977}{351.924}{p}%
\TextAt{21}{172.852}{356.263}{n}%
\TextAt{18}{178.489}{351.924}{Z}%
\TextAt{8}{190.644}{351.924}{where}%
\TextAt{8}{224.741}{351.924}{addition}%
\TextAt{8}{271.212}{351.924}{and}%
\TextAt{8}{294.262}{351.924}{multiplication}%
\TextAt{8}{369.676}{351.924}{are}%
\TextAt{8}{389.475}{351.924}{defined}%
\TextAt{8}{430.076}{351.924}{componentwise.}%
\TextAt{8}{516.522}{351.924}{We}%
\TextAt{8}{79.226}{334.879}{identify}%
\TextAt{18}{122.915}{334.879}{Z}%
\TextAt{21}{130.885}{333.085}{p}%
\TextAt{8}{140.636}{334.879}{with}%
\NativeMathAt{168.389}{334.879}{11.955}{\varprojlim}%
\TextAt{21}{184.648}{329.439}{n}%
\TextAt{18}{192.277}{334.879}{Z}%
\TextAt{9}{200.247}{334.879}{\Character{47}}%
\TextAt{13}{206.1}{334.879}{p}%
\TextAt{21}{211.975}{339.217}{n}%
\TextAt{18}{217.611}{334.879}{Z}%
\TextAt{8}{230.572}{334.879}{and}%
\TextAt{8}{254.427}{334.879}{obtain}%
\TextAt{8}{291.947}{334.879}{the}%
\TextAt{12}{313.196}{334.879}{ring}%
\TextAt{12}{342.38}{334.879}{of}%
\TextAt{13}{358.964}{334.879}{p}%
\TextAt{12}{364.839}{334.879}{-adic}%
\TextAt{12}{398.793}{334.879}{integers}%
\TextAt{18}{449.667}{334.879}{Z}%
\TextAt{21}{457.637}{333.085}{p}%
\TextAt{8}{462.398}{334.879}{.}%
\TextAt{8}{474.125}{334.879}{Since}%
\TextAt{8}{505.775}{334.879}{every}%
\TextAt{8}{79.226}{317.833}{element}%
\TextAt{13}{122.468}{317.833}{f}%
\NativeMathAt{132.835}{317.833}{11.955}{\in}%
\TextAt{18}{144.126}{317.833}{Q}%
\TextAt{21}{153.425}{316.039}{p}%
\TextAt{8}{162.083}{317.833}{admits}%
\TextAt{8}{200.519}{317.833}{a}%
\TextAt{8}{210.274}{317.833}{representation}%
\TextAt{13}{237.76}{288.831}{f}%
\TextAt{9}{248.127}{288.831}{\Character{61}}%
\TextAt{13}{260.552}{288.831}{p}%
\TextAt{3}{266.427}{293.768}{\Character{0}}%
\TextAt{21}{273.014}{293.768}{m}%
\TextAt{13}{281.003}{288.831}{g}%
\TextAt{8}{310.947}{288.831}{with}%
\TextAt{13}{337.607}{288.831}{g}%
\NativeMathAt{346.962}{288.831}{11.955}{\in}%
\TextAt{18}{358.253}{288.831}{Z}%
\TextAt{21}{366.223}{287.038}{p}%
\TextAt{13}{370.984}{288.831}{\Character{59}}%
\TextAt{8}{79.226}{259.83}{addition}%
\TextAt{8}{124.841}{259.83}{and}%
\TextAt{8}{147.035}{259.83}{multiplication}%
\TextAt{8}{221.593}{259.83}{extend}%
\TextAt{8}{259.066}{259.83}{from}%
\TextAt{18}{286.139}{259.83}{Z}%
\TextAt{21}{294.109}{258.037}{p}%
\TextAt{8}{302.199}{259.83}{to}%
\TextAt{18}{315.941}{259.83}{Q}%
\TextAt{21}{325.239}{258.037}{p}%
\TextAt{8}{333.33}{259.83}{and}%
\TextAt{18}{355.524}{259.83}{Q}%
\TextAt{21}{364.823}{258.037}{p}%
\TextAt{8}{372.913}{259.83}{becomes}%
\TextAt{8}{418.898}{259.83}{the}%
\TextAt{8}{438.486}{259.83}{field}%
\TextAt{8}{463.275}{259.83}{of}%
\TextAt{8}{476.037}{259.83}{fractions}%
\TextAt{8}{523.337}{259.83}{of}%
\TextAt{18}{79.226}{242.784}{Z}%
\TextAt{21}{87.196}{240.991}{p}%
\TextAt{8}{91.957}{242.784}{.}%
\TextAt{8}{100.124}{242.784}{In}%
\TextAt{18}{113.9}{242.784}{Z}%
\TextAt{21}{121.87}{240.991}{p}%
\TextAt{8}{126.631}{242.784}{,}%
\TextAt{8}{133.104}{242.784}{we}%
\TextAt{8}{149.487}{242.784}{found}%
\TextAt{8}{181.482}{242.784}{the}%
\TextAt{8}{200.794}{242.784}{rational}%
\TextAt{8}{243.537}{242.784}{integers}%
\TextAt{13}{286.006}{242.784}{a}%
\NativeMathAt{295.471}{242.784}{11.955}{\in}%
\TextAt{18}{306.762}{242.784}{Z}%
\TextAt{8}{317.785}{242.784}{which}%
\TextAt{8}{350.426}{242.784}{were}%
\TextAt{8}{376.564}{242.784}{determined}%
\TextAt{8}{436.846}{242.784}{by}%
\TextAt{8}{452.26}{242.784}{the}%
\TextAt{8}{471.571}{242.784}{congruences}%
\TextAt{13}{199.11}{213.783}{a}%
\NativeMathAt{208.576}{213.783}{11.955}{\equiv}%
\TextAt{13}{221.195}{213.783}{a}%
\TextAt{17}{227.34}{211.99}{0}%
\TextAt{9}{234.729}{213.783}{\Character{43}}%
\TextAt{13}{246.49}{213.783}{a}%
\TextAt{17}{252.635}{211.99}{1}%
\TextAt{13}{257.368}{213.783}{p}%
\TextAt{9}{265.9}{213.783}{\Character{43}}%
\TextAt{16}{277.661}{213.783}{\Character{1}}%
\TextAt{16}{282.957}{213.783}{\Character{1}}%
\TextAt{16}{288.265}{213.783}{\Character{1}}%
\TextAt{9}{294.265}{213.783}{\Character{43}}%
\TextAt{13}{306.027}{213.783}{a}%
\TextAt{21}{312.171}{211.99}{n}%
\TextAt{3}{317.31}{211.99}{\Character{0}}%
\TextAt{17}{323.896}{211.99}{1}%
\TextAt{13}{328.628}{213.783}{p}%
\TextAt{21}{334.504}{218.719}{n}%
\TextAt{3}{339.642}{218.719}{\Character{0}}%
\TextAt{17}{346.228}{218.719}{1}%
\TextAt{9}{362.916}{213.783}{\Character{40}}%
\TextAt{8}{367.468}{213.783}{mod}%
\TextAt{13}{393.57}{213.783}{p}%
\TextAt{21}{399.445}{218.719}{n}%
\TextAt{9}{405.082}{213.783}{\Character{41}}%
\TextAt{13}{409.634}{213.783}{\Character{59}}%
\TextAt{9}{79.226}{184.782}{0}%
\NativeMathAt{88.4}{184.782}{11.955}{\leq}%
\TextAt{13}{101.019}{184.782}{a}%
\TextAt{21}{107.164}{182.989}{i}%
\TextAt{13}{113.866}{184.782}{\Character{60}}%
\TextAt{13}{126.276}{184.782}{p}%
\TextAt{8}{132.167}{184.782}{.}%
\TextAt{8}{140.615}{184.782}{Making}%
\TextAt{8}{182.889}{184.782}{the}%
\TextAt{8}{203.045}{184.782}{identification}%
\TextAt{18}{265.981}{155.781}{Z}%
\TextAt{21}{273.951}{153.988}{p}%
\TextAt{9}{282.033}{155.781}{\Character{61}}%
\NativeMathAt{294.458}{155.781}{11.955}{\varprojlim}%
\TextAt{21}{300.019}{143.666}{n}%
\TextAt{18}{312.71}{155.781}{Z}%
\TextAt{9}{320.68}{155.781}{\Character{47}}%
\TextAt{13}{326.533}{155.781}{p}%
\TextAt{21}{332.408}{160.717}{n}%
\TextAt{18}{338.045}{155.781}{Z}%
\TextAt{8}{79.226}{121.103}{the}%
\TextAt{8}{99.383}{121.103}{subset}%
\TextAt{18}{135.272}{121.103}{Z}%
\TextAt{8}{147.14}{121.103}{is}%
\TextAt{8}{158.903}{121.103}{taken}%
\TextAt{8}{190.776}{121.103}{to}%
\TextAt{8}{205.086}{121.103}{the}%
\TextAt{8}{225.243}{121.103}{set}%
\TextAt{8}{243.51}{121.103}{of}%
\TextAt{8}{256.84}{121.103}{tuples}%
\TextAt{9}{187.875}{86.261}{\Character{40}}%
\TextAt{13}{192.427}{86.261}{a}%
\TextAt{8}{201.893}{86.261}{mod}%
\TextAt{13}{227.331}{86.261}{p\Character{59}}%
\TextAt{13}{238.425}{86.261}{a}%
\TextAt{8}{247.916}{86.261}{mod}%
\TextAt{13}{273.354}{86.261}{p}%
\TextAt{17}{279.229}{91.197}{2}%
\TextAt{13}{283.961}{86.261}{\Character{59}}%
\TextAt{13}{289.185}{86.261}{a}%
\TextAt{8}{298.671}{86.261}{mod}%
\TextAt{13}{324.109}{86.261}{p}%
\TextAt{17}{329.984}{91.197}{3}%
\TextAt{13}{334.716}{86.261}{\Character{59}}%
\TextAt{13}{339.94}{86.261}{\Character{58}}%
\TextAt{13}{345.177}{86.261}{\Character{58}}%
\TextAt{13}{350.413}{86.261}{\Character{58}}%
\TextAt{9}{353.701}{86.261}{\Character{41}}%
\NativeMathAt{361.574}{86.261}{11.955}{\in}%
\NativeMathAt{376.61}{101.205}{7.97}{\infty}%
\TextAt{20}{373.206}{97.619}{\Character{89}}%
\TextAt{21}{372.865}{72.513}{n}%
\TextAt{17}{378.003}{72.513}{\Character{61}1}%
\TextAt{18}{390.816}{86.261}{Z}%
\TextAt{9}{398.786}{86.261}{\Character{47}}%
\TextAt{13}{404.639}{86.261}{p}%
\TextAt{21}{410.514}{91.197}{n}%
\TextAt{18}{416.151}{86.261}{Z}%
\TextAt{8}{300.14}{43.033}{95}%
\Destination{78.23}{809.815}{nav.page.95}%
\Destination{79.36}{776.155}{nav.print.103}%
\Destination{79.36}{401.055}{nav.print.104}%
\Destination{79.36}{601.555}{nav.number.2.1.3}%
\Hotlink{389.883}{383.855}{25.078}{11.047}{nav.number.2.1.2}%
}

% Page 108 (body page 96)
\Page{%
\TextAt{2}{79.226}{796.943}{II.}%
\TextAt{2}{94.073}{796.943}{The}%
\TextAt{2}{116.492}{796.943}{Theory}%
\TextAt{2}{154.401}{796.943}{of}%
\TextAt{2}{166.826}{796.943}{Valuations}%
\Vector{q 0 G 0 g 0.299 w 7.226 18.729 m 460.769 18.729 l S Q}%
\TextAt{8}{79.226}{761.124}{and}%
\TextAt{8}{103.175}{761.124}{thereby}%
\TextAt{8}{146.635}{761.124}{is}%
\TextAt{8}{159.585}{761.124}{realized}%
\TextAt{8}{203.691}{761.124}{as}%
\TextAt{8}{219.247}{761.124}{a}%
\TextAt{8}{230.189}{761.124}{subring}%
\TextAt{8}{273.063}{761.124}{of}%
\TextAt{18}{287.579}{761.124}{Z}%
\TextAt{21}{295.549}{759.33}{p}%
\TextAt{8}{300.31}{761.124}{.}%
\TextAt{8}{312.318}{761.124}{We}%
\TextAt{8}{333.648}{761.124}{obtain}%
\TextAt{18}{371.262}{761.124}{Q}%
\TextAt{8}{385.644}{761.124}{as}%
\TextAt{8}{401.201}{761.124}{a}%
\TextAt{8}{412.143}{761.124}{subfield}%
\TextAt{8}{456.308}{761.124}{of}%
\TextAt{8}{470.824}{761.124}{the}%
\TextAt{8}{492.167}{761.124}{field}%
\TextAt{18}{518.71}{761.124}{Q}%
\TextAt{21}{528.009}{759.33}{p}%
\TextAt{8}{79.226}{744.078}{of}%
\TextAt{13}{93.058}{744.078}{p}%
\TextAt{8}{98.933}{744.078}{-adic}%
\TextAt{8}{128.043}{744.078}{numbers}%
\TextAt{8}{175.756}{744.078}{in}%
\TextAt{8}{189.91}{744.078}{the}%
\TextAt{8}{210.568}{744.078}{same}%
\TextAt{8}{240.396}{744.078}{way.}%
\TextAt{8}{269.215}{744.078}{Despite}%
\TextAt{8}{311.882}{744.078}{their}%
\TextAt{8}{340.347}{744.078}{origin}%
\TextAt{8}{374.024}{744.078}{in}%
\TextAt{8}{388.179}{744.078}{function-theoretic}%
\TextAt{8}{483.306}{744.078}{ideas,}%
\TextAt{8}{516.51}{744.078}{the}%
\TextAt{13}{79.226}{727.032}{p}%
\TextAt{8}{85.101}{727.032}{-adic}%
\TextAt{8}{114.348}{727.032}{numbers}%
\TextAt{8}{162.196}{727.032}{live}%
\TextAt{8}{184.293}{727.032}{up}%
\TextAt{8}{201.835}{727.032}{to}%
\TextAt{8}{216.783}{727.032}{their}%
\TextAt{8}{245.384}{727.032}{destiny}%
\TextAt{8}{286.406}{727.032}{entirely}%
\TextAt{8}{329.317}{727.032}{within}%
\TextAt{8}{366.37}{727.032}{arithmetic,}%
\TextAt{8}{427.002}{727.032}{more}%
\TextAt{8}{456.906}{727.032}{precisely}%
\TextAt{8}{505.4}{727.032}{at}%
\TextAt{8}{520.348}{727.032}{its}%
\TextAt{8}{79.226}{709.986}{classical}%
\TextAt{8}{124.225}{709.986}{heart,}%
\TextAt{8}{158.047}{709.986}{the}%
\TextAt{8}{178.203}{709.986}{Diophantine}%
\TextAt{8}{244.71}{709.986}{equations.}%
\TextAt{8}{301.684}{709.986}{Such}%
\TextAt{8}{329.97}{709.986}{an}%
\TextAt{8}{346.229}{709.986}{equation}%
\TextAt{13}{261.097}{682.035}{F}%
\TextAt{9}{270.301}{682.035}{\Character{40}}%
\TextAt{13}{274.853}{682.035}{x}%
\TextAt{17}{281.505}{680.242}{1}%
\TextAt{13}{286.238}{682.035}{\Character{59}}%
\TextAt{13}{291.462}{682.035}{\Character{58}}%
\TextAt{13}{296.699}{682.035}{\Character{58}}%
\TextAt{13}{301.923}{682.035}{\Character{58}}%
\TextAt{13}{307.16}{682.035}{\Character{59}}%
\TextAt{13}{312.396}{682.035}{x}%
\TextAt{21}{319.11}{680.242}{n}%
\TextAt{9}{324.747}{682.035}{\Character{41}}%
\TextAt{9}{332.602}{682.035}{\Character{61}}%
\TextAt{9}{345.023}{682.035}{0}%
\TextAt{8}{79.226}{654.084}{is}%
\TextAt{8}{90.061}{654.084}{given}%
\TextAt{8}{119.701}{654.084}{by}%
\TextAt{8}{135.03}{654.084}{a}%
\TextAt{8}{143.856}{654.084}{polynomial}%
\TextAt{13}{203.42}{654.084}{F}%
\NativeMathAt{215.944}{654.084}{11.955}{\in}%
\TextAt{18}{227.235}{654.084}{Z}%
\TextAt{9}{235.205}{654.084}{\Character{91}}%
\TextAt{13}{238.457}{654.084}{x}%
\TextAt{17}{245.109}{652.291}{1}%
\TextAt{13}{249.842}{654.084}{\Character{59}}%
\TextAt{13}{255.066}{654.084}{\Character{58}}%
\TextAt{13}{260.303}{654.084}{\Character{58}}%
\TextAt{13}{265.527}{654.084}{\Character{58}}%
\TextAt{13}{270.764}{654.084}{\Character{59}}%
\TextAt{13}{276}{654.084}{x}%
\TextAt{21}{282.714}{652.291}{n}%
\TextAt{9}{288.351}{654.084}{\Character{93}}%
\TextAt{8}{291.602}{654.084}{,}%
\TextAt{8}{298.008}{654.084}{and}%
\TextAt{8}{319.841}{654.084}{the}%
\TextAt{8}{339.068}{654.084}{question}%
\TextAt{8}{384.704}{654.084}{is}%
\TextAt{8}{395.539}{654.084}{whether}%
\TextAt{8}{439.477}{654.084}{it}%
\TextAt{8}{450.252}{654.084}{admits}%
\TextAt{8}{487.758}{654.084}{solutions}%
\TextAt{8}{79.226}{637.038}{in}%
\TextAt{8}{91.674}{637.038}{integers.}%
\TextAt{8}{139.137}{637.038}{This}%
\TextAt{8}{164.652}{637.038}{difficult}%
\TextAt{8}{206.367}{637.038}{problem}%
\TextAt{8}{250.687}{637.038}{can}%
\TextAt{8}{270.942}{637.038}{be}%
\TextAt{8}{285.662}{637.038}{weakened}%
\TextAt{8}{336.809}{637.038}{by}%
\TextAt{8}{351.863}{637.038}{considering,}%
\TextAt{8}{415.349}{637.038}{instead}%
\TextAt{8}{454.529}{637.038}{of}%
\TextAt{8}{466.655}{637.038}{the}%
\TextAt{8}{485.606}{637.038}{equation,}%
\TextAt{8}{79.226}{619.992}{all}%
\TextAt{8}{95.485}{619.992}{the}%
\TextAt{8}{115.642}{619.992}{congruences}%
\TextAt{13}{230.674}{602.946}{F}%
\TextAt{9}{239.878}{602.946}{\Character{40}}%
\TextAt{13}{244.43}{602.946}{x}%
\TextAt{17}{251.082}{601.153}{1}%
\TextAt{13}{255.814}{602.946}{\Character{59}}%
\TextAt{13}{261.038}{602.946}{\Character{58}}%
\TextAt{13}{266.275}{602.946}{\Character{58}}%
\TextAt{13}{271.511}{602.946}{\Character{58}}%
\TextAt{13}{276.736}{602.946}{\Character{59}}%
\TextAt{13}{281.972}{602.946}{x}%
\TextAt{21}{288.687}{601.153}{n}%
\TextAt{9}{294.324}{602.946}{\Character{41}}%
\NativeMathAt{302.197}{602.946}{11.955}{\equiv}%
\TextAt{9}{314.816}{602.946}{0}%
\TextAt{9}{332.605}{602.946}{\Character{40}}%
\TextAt{8}{337.176}{602.946}{mod}%
\TextAt{13}{363.278}{602.946}{m}%
\TextAt{9}{373.518}{602.946}{\Character{41}}%
\TextAt{13}{378.07}{602.946}{\Character{58}}%
\TextAt{8}{79.226}{579.875}{By}%
\TextAt{8}{97.589}{579.875}{the}%
\TextAt{8}{117.746}{579.875}{Chinese}%
\TextAt{8}{161.37}{579.875}{remainder}%
\TextAt{8}{216.651}{579.875}{theorem,}%
\TextAt{8}{265.428}{579.875}{this}%
\TextAt{8}{288.25}{579.875}{amounts}%
\TextAt{8}{335.473}{579.875}{to}%
\TextAt{8}{349.784}{579.875}{considering}%
\TextAt{8}{410.982}{579.875}{the}%
\TextAt{8}{431.139}{579.875}{congruences}%
\TextAt{13}{231.891}{551.924}{F}%
\TextAt{9}{241.094}{551.924}{\Character{40}}%
\TextAt{13}{245.647}{551.924}{x}%
\TextAt{17}{252.299}{550.13}{1}%
\TextAt{13}{257.031}{551.924}{\Character{59}}%
\TextAt{13}{262.255}{551.924}{\Character{58}}%
\TextAt{13}{267.492}{551.924}{\Character{58}}%
\TextAt{13}{272.728}{551.924}{\Character{58}}%
\TextAt{13}{277.953}{551.924}{\Character{59}}%
\TextAt{13}{283.189}{551.924}{x}%
\TextAt{21}{289.904}{550.13}{n}%
\TextAt{9}{295.54}{551.924}{\Character{41}}%
\NativeMathAt{303.414}{551.924}{11.955}{\equiv}%
\TextAt{9}{316.033}{551.924}{0}%
\TextAt{9}{333.822}{551.924}{\Character{40}}%
\TextAt{8}{338.393}{551.924}{mod}%
\TextAt{13}{364.495}{551.924}{p}%
\TextAt{21}{370.37}{556.86}{\Character{23}}%
\TextAt{9}{375.552}{551.924}{\Character{41}}%
\TextAt{8}{79.226}{523.973}{modulo}%
\TextAt{8}{120.363}{523.973}{all}%
\TextAt{8}{135.808}{523.973}{prime}%
\TextAt{8}{168.158}{523.973}{powers.}%
\TextAt{8}{211.196}{523.973}{The}%
\TextAt{8}{234.437}{523.973}{hope}%
\TextAt{8}{261.909}{523.973}{is}%
\TextAt{8}{272.859}{523.973}{to}%
\TextAt{8}{286.355}{523.973}{obtain}%
\TextAt{8}{321.969}{523.973}{in}%
\TextAt{8}{334.808}{523.973}{this}%
\TextAt{8}{356.817}{523.973}{way}%
\TextAt{8}{379.746}{523.973}{information}%
\TextAt{8}{442.354}{523.973}{about}%
\TextAt{8}{475.039}{523.973}{the}%
\TextAt{8}{494.381}{523.973}{original}%
\TextAt{8}{79.226}{506.927}{equation.}%
\TextAt{8}{131.432}{506.927}{This}%
\TextAt{8}{157.692}{506.927}{plethora}%
\TextAt{8}{203.415}{506.927}{of}%
\TextAt{8}{216.285}{506.927}{congruences}%
\TextAt{8}{280.921}{506.927}{is}%
\TextAt{8}{292.224}{506.927}{now}%
\TextAt{8}{316.153}{506.927}{synthesized}%
\TextAt{8}{377.597}{506.927}{again}%
\TextAt{8}{408.363}{506.927}{into}%
\TextAt{8}{431.646}{506.927}{a}%
\TextAt{8}{440.941}{506.927}{single}%
\TextAt{8}{473.059}{506.927}{equation}%
\TextAt{8}{520.408}{506.927}{by}%
\TextAt{8}{79.226}{489.881}{means}%
\TextAt{8}{115.056}{489.881}{of}%
\TextAt{8}{128.386}{489.881}{the}%
\TextAt{13}{148.542}{489.881}{p}%
\TextAt{8}{154.417}{489.881}{-adic}%
\TextAt{8}{183.026}{489.881}{numbers.}%
\TextAt{8}{234.788}{489.881}{In}%
\TextAt{8}{249.409}{489.881}{fact,}%
\TextAt{8}{275.746}{489.881}{one}%
\TextAt{8}{297.206}{489.881}{has}%
\TextAt{8}{318.08}{489.881}{the}%
\TextAt{12}{79.226}{464.169}{Proposition}%
\TextAt{12}{152.238}{464.169}{(1.4).}%
\TextAt{8}{187.815}{464.169}{Let}%
\TextAt{13}{209.046}{464.169}{F}%
\TextAt{9}{218.249}{464.169}{\Character{40}}%
\TextAt{13}{222.802}{464.169}{x}%
\TextAt{17}{229.454}{462.375}{1}%
\TextAt{13}{234.186}{464.169}{\Character{59}}%
\TextAt{13}{239.41}{464.169}{\Character{58}}%
\TextAt{13}{244.647}{464.169}{\Character{58}}%
\TextAt{13}{249.871}{464.169}{\Character{58}}%
\TextAt{13}{255.108}{464.169}{\Character{59}}%
\TextAt{13}{260.344}{464.169}{x}%
\TextAt{21}{267.059}{462.375}{n}%
\TextAt{9}{272.695}{464.169}{\Character{41}}%
\TextAt{8}{281.406}{464.169}{be}%
\TextAt{8}{297.592}{464.169}{a}%
\TextAt{8}{307.609}{464.169}{polynomial}%
\TextAt{8}{368.363}{464.169}{with}%
\TextAt{8}{395.285}{464.169}{integer}%
\TextAt{8}{434.245}{464.169}{coefficients,}%
\TextAt{8}{497.061}{464.169}{and}%
\TextAt{8}{520.085}{464.169}{fix}%
\TextAt{8}{79.226}{447.123}{a}%
\TextAt{8}{88.982}{447.123}{prime}%
\TextAt{8}{122.145}{447.123}{number}%
\TextAt{13}{164.741}{447.123}{p}%
\TextAt{8}{170.617}{447.123}{.}%
\TextAt{8}{179.065}{447.123}{The}%
\TextAt{8}{203.119}{447.123}{congruence}%
\TextAt{13}{231.891}{419.172}{F}%
\TextAt{9}{241.094}{419.172}{\Character{40}}%
\TextAt{13}{245.647}{419.172}{x}%
\TextAt{17}{252.299}{417.379}{1}%
\TextAt{13}{257.031}{419.172}{\Character{59}}%
\TextAt{13}{262.255}{419.172}{\Character{58}}%
\TextAt{13}{267.492}{419.172}{\Character{58}}%
\TextAt{13}{272.728}{419.172}{\Character{58}}%
\TextAt{13}{277.953}{419.172}{\Character{59}}%
\TextAt{13}{283.189}{419.172}{x}%
\TextAt{21}{289.904}{417.379}{n}%
\TextAt{9}{295.54}{419.172}{\Character{41}}%
\NativeMathAt{303.414}{419.172}{11.955}{\equiv}%
\TextAt{9}{316.033}{419.172}{0}%
\TextAt{9}{333.822}{419.172}{\Character{40}}%
\TextAt{8}{338.393}{419.172}{mod}%
\TextAt{13}{364.495}{419.172}{p}%
\TextAt{21}{370.37}{424.108}{\Character{23}}%
\TextAt{9}{375.552}{419.172}{\Character{41}}%
\TextAt{8}{79.226}{391.221}{is}%
\TextAt{8}{90.99}{391.221}{solvable}%
\TextAt{8}{134.961}{391.221}{for}%
\TextAt{8}{152.846}{391.221}{arbitrary}%
\TextAt{13}{202.615}{391.221}{\Character{23}}%
\NativeMathAt{212.474}{391.221}{11.955}{\geq}%
\TextAt{9}{225.094}{391.221}{1}%
\TextAt{8}{234.844}{391.221}{if}%
\TextAt{8}{245.568}{391.221}{and}%
\TextAt{8}{268.33}{391.221}{only}%
\TextAt{8}{294.022}{391.221}{if}%
\TextAt{8}{304.746}{391.221}{the}%
\TextAt{8}{324.902}{391.221}{equation}%
\TextAt{13}{261.097}{363.27}{F}%
\TextAt{9}{270.301}{363.27}{\Character{40}}%
\TextAt{13}{274.853}{363.27}{x}%
\TextAt{17}{281.505}{361.477}{1}%
\TextAt{13}{286.238}{363.27}{\Character{59}}%
\TextAt{13}{291.462}{363.27}{\Character{58}}%
\TextAt{13}{296.699}{363.27}{\Character{58}}%
\TextAt{13}{301.923}{363.27}{\Character{58}}%
\TextAt{13}{307.16}{363.27}{\Character{59}}%
\TextAt{13}{312.396}{363.27}{x}%
\TextAt{21}{319.11}{361.477}{n}%
\TextAt{9}{324.747}{363.27}{\Character{41}}%
\TextAt{9}{332.602}{363.27}{\Character{61}}%
\TextAt{9}{345.023}{363.27}{0}%
\TextAt{8}{79.226}{335.319}{is}%
\TextAt{8}{90.99}{335.319}{solvable}%
\TextAt{8}{134.961}{335.319}{in}%
\TextAt{13}{148.614}{335.319}{p}%
\TextAt{8}{154.489}{335.319}{-adic}%
\TextAt{8}{183.098}{335.319}{integers.}%
\TextAt{19}{79.226}{300.941}{Proof.}%
\TextAt{8}{116}{300.941}{As}%
\TextAt{8}{133.287}{300.941}{established}%
\TextAt{8}{193.243}{300.941}{above,}%
\TextAt{8}{229.67}{300.941}{we}%
\TextAt{8}{246.897}{300.941}{view}%
\TextAt{8}{273.88}{300.941}{the}%
\TextAt{8}{294.037}{300.941}{ring}%
\TextAt{18}{318.102}{300.941}{Z}%
\TextAt{21}{326.073}{299.147}{p}%
\TextAt{8}{334.731}{300.941}{as}%
\TextAt{8}{349.101}{300.941}{the}%
\TextAt{8}{369.257}{300.941}{projective}%
\TextAt{8}{423.235}{300.941}{limit}%
\TextAt{18}{231.439}{266.125}{Z}%
\TextAt{21}{239.409}{264.332}{p}%
\TextAt{9}{247.49}{266.125}{\Character{61}}%
\NativeMathAt{259.916}{266.125}{11.955}{\varprojlim}%
\TextAt{21}{265.703}{254.011}{\Character{23}}%
\TextAt{18}{278.167}{266.125}{Z}%
\TextAt{9}{286.138}{266.125}{\Character{47}}%
\TextAt{13}{291.991}{266.125}{p}%
\TextAt{21}{297.866}{271.062}{\Character{23}}%
\TextAt{18}{303.048}{266.125}{Z}%
\NativeMathAt{314.339}{266.125}{11.955}{\subseteq}%
\NativeMathAt{330.476}{281.069}{7.97}{\infty}%
\TextAt{20}{327.072}{277.483}{\Character{89}}%
\TextAt{21}{326.958}{252.377}{\Character{23}}%
\TextAt{17}{331.642}{252.377}{\Character{61}1}%
\TextAt{18}{344.455}{266.125}{Z}%
\TextAt{9}{352.425}{266.125}{\Character{47}}%
\TextAt{13}{358.278}{266.125}{p}%
\TextAt{21}{364.153}{271.062}{\Character{23}}%
\TextAt{18}{369.335}{266.125}{Z}%
\TextAt{13}{377.305}{266.125}{\Character{58}}%
\TextAt{8}{79.226}{230.983}{Viewed}%
\TextAt{8}{119.887}{230.983}{over}%
\TextAt{8}{144.635}{230.983}{the}%
\TextAt{8}{164.494}{230.983}{ring}%
\TextAt{8}{188.262}{230.983}{on}%
\TextAt{8}{204.223}{230.983}{the}%
\TextAt{8}{224.082}{230.983}{right,}%
\TextAt{8}{255.394}{230.983}{the}%
\TextAt{8}{275.252}{230.983}{equation}%
\TextAt{13}{322.763}{230.983}{F}%
\TextAt{9}{335.288}{230.983}{\Character{61}}%
\TextAt{9}{347.697}{230.983}{0}%
\TextAt{8}{357.166}{230.983}{splits}%
\TextAt{8}{387.557}{230.983}{up}%
\TextAt{8}{404.164}{230.983}{into}%
\TextAt{8}{427.609}{230.983}{components}%
\TextAt{8}{491.762}{230.983}{over}%
\TextAt{8}{516.51}{230.983}{the}%
\TextAt{8}{79.226}{213.937}{individual}%
\TextAt{8}{134.184}{213.937}{rings}%
\TextAt{18}{162.865}{213.937}{Z}%
\TextAt{9}{170.835}{213.937}{\Character{47}}%
\TextAt{13}{176.688}{213.937}{p}%
\TextAt{21}{182.563}{218.276}{\Character{23}}%
\TextAt{18}{187.745}{213.937}{Z}%
\TextAt{8}{195.715}{213.937}{,}%
\TextAt{8}{202.864}{213.937}{namely,}%
\TextAt{8}{245.783}{213.937}{the}%
\TextAt{8}{265.94}{213.937}{congruences}%
\TextAt{13}{230.265}{185.987}{F}%
\TextAt{9}{239.469}{185.987}{\Character{40}}%
\TextAt{13}{244.021}{185.987}{x}%
\TextAt{17}{250.673}{184.193}{1}%
\TextAt{13}{255.405}{185.987}{\Character{59}}%
\TextAt{13}{260.629}{185.987}{\Character{58}}%
\TextAt{13}{265.866}{185.987}{\Character{58}}%
\TextAt{13}{271.102}{185.987}{\Character{58}}%
\TextAt{13}{276.327}{185.987}{\Character{59}}%
\TextAt{13}{281.563}{185.987}{x}%
\TextAt{21}{288.278}{184.193}{n}%
\TextAt{9}{293.915}{185.987}{\Character{41}}%
\NativeMathAt{301.788}{185.987}{11.955}{\equiv}%
\TextAt{9}{314.407}{185.987}{0}%
\TextAt{9}{332.196}{185.987}{\Character{40}}%
\TextAt{8}{336.767}{185.987}{mod}%
\TextAt{13}{362.869}{185.987}{p}%
\TextAt{21}{368.745}{190.923}{\Character{23}}%
\TextAt{9}{373.927}{185.987}{\Character{41}}%
\TextAt{13}{378.479}{185.987}{\Character{58}}%
\TextAt{8}{79.226}{158.036}{If}%
\TextAt{8}{90.918}{158.036}{now}%
\TextAt{9}{208.515}{140.99}{\Character{40}}%
\TextAt{13}{213.068}{140.99}{x}%
\TextAt{17}{219.72}{139.196}{1}%
\TextAt{13}{224.452}{140.99}{\Character{59}}%
\TextAt{13}{229.676}{140.99}{\Character{58}}%
\TextAt{13}{234.913}{140.99}{\Character{58}}%
\TextAt{13}{240.149}{140.99}{\Character{58}}%
\TextAt{13}{245.374}{140.99}{\Character{59}}%
\TextAt{13}{250.61}{140.99}{x}%
\TextAt{21}{257.325}{139.196}{n}%
\TextAt{9}{262.961}{140.99}{\Character{41}}%
\TextAt{9}{270.816}{140.99}{\Character{61}}%
\TextAt{9}{283.237}{140.99}{\Character{40}}%
\TextAt{13}{287.812}{140.99}{x}%
\TextAt{17}{294.464}{147.1}{\Character{40}}%
\TextAt{21}{297.758}{147.1}{\Character{23}}%
\TextAt{17}{302.441}{147.1}{\Character{41}}%
\TextAt{17}{294.464}{138.059}{1}%
\TextAt{13}{306.233}{140.99}{\Character{59}}%
\TextAt{13}{311.457}{140.99}{\Character{58}}%
\TextAt{13}{316.694}{140.99}{\Character{58}}%
\TextAt{13}{321.918}{140.99}{\Character{58}}%
\TextAt{13}{327.155}{140.99}{\Character{59}}%
\TextAt{13}{332.391}{140.99}{x}%
\TextAt{17}{339.106}{145.926}{\Character{40}}%
\TextAt{21}{342.399}{145.926}{\Character{23}}%
\TextAt{17}{347.083}{145.926}{\Character{41}}%
\TextAt{21}{339.106}{138.034}{n}%
\TextAt{9}{350.874}{140.99}{\Character{41}}%
\TextAt{21}{355.427}{139.196}{\Character{23}}%
\NativeMathAt{360.11}{139.196}{7.97}{\in}%
\TextAt{25}{365.756}{139.196}{N}%
\NativeMathAt{375.331}{140.99}{11.955}{\in}%
\TextAt{18}{386.622}{140.99}{Z}%
\TextAt{21}{394.592}{145.926}{n}%
\TextAt{21}{394.592}{138.034}{p}%
\TextAt{13}{400.229}{140.99}{\Character{59}}%
\TextAt{8}{79.226}{117.918}{with}%
\TextAt{9}{105.969}{117.918}{\Character{40}}%
\TextAt{13}{110.521}{117.918}{x}%
\TextAt{17}{117.173}{124.029}{\Character{40}}%
\TextAt{21}{120.466}{124.029}{\Character{23}}%
\TextAt{17}{125.15}{124.029}{\Character{41}}%
\TextAt{21}{117.173}{114.854}{i}%
\TextAt{9}{128.942}{117.918}{\Character{41}}%
\TextAt{21}{133.494}{116.125}{\Character{23}}%
\NativeMathAt{138.178}{116.125}{7.97}{\in}%
\TextAt{25}{143.823}{116.125}{N}%
\NativeMathAt{153.539}{117.918}{11.955}{\in}%
\TextAt{18}{164.971}{117.918}{Z}%
\TextAt{21}{172.941}{116.125}{p}%
\TextAt{9}{181.163}{117.918}{\Character{61}}%
\NativeMathAt{193.729}{117.918}{11.955}{\varprojlim}%
\TextAt{21}{209.988}{112.479}{\Character{23}}%
\TextAt{18}{217.163}{117.918}{Z}%
\TextAt{9}{225.133}{117.918}{\Character{47}}%
\TextAt{13}{230.986}{117.918}{p}%
\TextAt{21}{236.861}{122.257}{\Character{23}}%
\TextAt{18}{242.043}{117.918}{Z}%
\TextAt{8}{250.013}{117.918}{,}%
\TextAt{8}{257.266}{117.918}{is}%
\TextAt{8}{269.112}{117.918}{a}%
\TextAt{13}{278.95}{117.918}{p}%
\TextAt{8}{284.825}{117.918}{-adic}%
\TextAt{8}{313.516}{117.918}{solution}%
\TextAt{8}{357.893}{117.918}{of}%
\TextAt{8}{371.305}{117.918}{the}%
\TextAt{8}{391.544}{117.918}{equation}%
\TextAt{13}{439.435}{117.918}{F}%
\TextAt{9}{448.639}{117.918}{\Character{40}}%
\TextAt{13}{453.191}{117.918}{x}%
\TextAt{17}{459.843}{116.125}{1}%
\TextAt{13}{464.576}{117.918}{\Character{59}}%
\TextAt{13}{469.8}{117.918}{\Character{58}}%
\TextAt{13}{475.037}{117.918}{\Character{58}}%
\TextAt{13}{480.261}{117.918}{\Character{58}}%
\TextAt{13}{485.498}{117.918}{\Character{59}}%
\TextAt{13}{490.734}{117.918}{x}%
\TextAt{21}{497.449}{116.125}{n}%
\TextAt{9}{503.085}{117.918}{\Character{41}}%
\TextAt{9}{511.083}{117.918}{\Character{61}}%
\TextAt{9}{523.648}{117.918}{0}%
\TextAt{8}{529.518}{117.918}{,}%
\TextAt{8}{79.226}{100.872}{then}%
\TextAt{8}{105.886}{100.872}{the}%
\TextAt{8}{126.043}{100.872}{congruences}%
\TextAt{8}{191.139}{100.872}{are}%
\TextAt{8}{210.649}{100.872}{solved}%
\TextAt{8}{245.833}{100.872}{by}%
\TextAt{13}{178.998}{72.921}{F}%
\TextAt{9}{188.202}{72.921}{\Character{40}}%
\TextAt{13}{192.754}{72.921}{x}%
\TextAt{17}{199.406}{79.032}{\Character{40}}%
\TextAt{21}{202.699}{79.032}{\Character{23}}%
\TextAt{17}{207.383}{79.032}{\Character{41}}%
\TextAt{17}{199.406}{69.99}{1}%
\TextAt{13}{211.175}{72.921}{\Character{59}}%
\TextAt{13}{216.399}{72.921}{\Character{58}}%
\TextAt{13}{221.636}{72.921}{\Character{58}}%
\TextAt{13}{226.86}{72.921}{\Character{58}}%
\TextAt{13}{232.097}{72.921}{\Character{59}}%
\TextAt{13}{237.333}{72.921}{x}%
\TextAt{17}{244.048}{77.858}{\Character{40}}%
\TextAt{21}{247.341}{77.858}{\Character{23}}%
\TextAt{17}{252.025}{77.858}{\Character{41}}%
\TextAt{21}{244.048}{69.966}{n}%
\TextAt{9}{255.816}{72.921}{\Character{41}}%
\NativeMathAt{263.689}{72.921}{11.955}{\equiv}%
\TextAt{9}{276.309}{72.921}{0}%
\TextAt{9}{294.098}{72.921}{\Character{40}}%
\TextAt{8}{298.669}{72.921}{mod}%
\TextAt{13}{324.771}{72.921}{p}%
\TextAt{21}{330.646}{77.858}{\Character{23}}%
\TextAt{9}{335.828}{72.921}{\Character{41}}%
\TextAt{13}{340.38}{72.921}{\Character{59}}%
\TextAt{13}{369.515}{72.921}{\Character{23}}%
\TextAt{9}{379.394}{72.921}{\Character{61}}%
\TextAt{9}{391.803}{72.921}{1}%
\TextAt{13}{397.672}{72.921}{\Character{59}}%
\TextAt{9}{402.916}{72.921}{2}%
\TextAt{13}{408.769}{72.921}{\Character{59}}%
\TextAt{13}{413.993}{72.921}{\Character{58}}%
\TextAt{13}{419.23}{72.921}{\Character{58}}%
\TextAt{13}{424.466}{72.921}{\Character{58}}%
\TextAt{13}{429.691}{72.921}{\Character{58}}%
\TextAt{8}{300.14}{43.033}{96}%
\Destination{78.23}{809.815}{nav.page.96}%
\Destination{79.36}{479.155}{nav.print.105}%
\Destination{79.36}{478.155}{nav.number.2.1.4}%
}

% Page 109 (body page 97)
\Page{%
\TextAt{2}{79.226}{796.943}{II.}%
\TextAt{2}{94.073}{796.943}{The}%
\TextAt{2}{116.492}{796.943}{Theory}%
\TextAt{2}{154.401}{796.943}{of}%
\TextAt{2}{166.826}{796.943}{Valuations}%
\Vector{q 0 G 0 g 0.299 w 7.226 18.729 m 460.769 18.729 l S Q}%
\TextAt{8}{79.226}{760.991}{Conversely,}%
\TextAt{8}{140.747}{760.991}{let}%
\TextAt{8}{157.652}{760.991}{a}%
\TextAt{8}{167.407}{760.991}{solution}%
\TextAt{9}{211.701}{760.991}{\Character{40}}%
\TextAt{13}{216.254}{760.991}{x}%
\TextAt{17}{222.906}{767.101}{\Character{40}}%
\TextAt{21}{226.199}{767.101}{\Character{23}}%
\TextAt{17}{230.883}{767.101}{\Character{41}}%
\TextAt{17}{222.906}{758.06}{1}%
\TextAt{13}{234.674}{760.991}{\Character{59}}%
\TextAt{13}{239.898}{760.991}{\Character{58}}%
\TextAt{13}{245.135}{760.991}{\Character{58}}%
\TextAt{13}{250.371}{760.991}{\Character{58}}%
\TextAt{13}{255.596}{760.991}{\Character{59}}%
\TextAt{13}{260.832}{760.991}{x}%
\TextAt{17}{267.547}{767.101}{\Character{40}}%
\TextAt{21}{270.84}{767.101}{\Character{23}}%
\TextAt{17}{275.524}{767.101}{\Character{41}}%
\TextAt{21}{267.547}{759.764}{n}%
\TextAt{9}{279.316}{760.991}{\Character{41}}%
\TextAt{8}{287.765}{760.991}{of}%
\TextAt{8}{301.095}{760.991}{the}%
\TextAt{8}{321.252}{760.991}{congruence}%
\TextAt{13}{231.891}{737.583}{F}%
\TextAt{9}{241.094}{737.583}{\Character{40}}%
\TextAt{13}{245.647}{737.583}{x}%
\TextAt{17}{252.299}{735.789}{1}%
\TextAt{13}{257.031}{737.583}{\Character{59}}%
\TextAt{13}{262.255}{737.583}{\Character{58}}%
\TextAt{13}{267.492}{737.583}{\Character{58}}%
\TextAt{13}{272.728}{737.583}{\Character{58}}%
\TextAt{13}{277.953}{737.583}{\Character{59}}%
\TextAt{13}{283.189}{737.583}{x}%
\TextAt{21}{289.904}{735.789}{n}%
\TextAt{9}{295.54}{737.583}{\Character{41}}%
\NativeMathAt{303.414}{737.583}{11.955}{\equiv}%
\TextAt{9}{316.033}{737.583}{0}%
\TextAt{9}{333.822}{737.583}{\Character{40}}%
\TextAt{8}{338.393}{737.583}{mod}%
\TextAt{13}{364.495}{737.583}{p}%
\TextAt{21}{370.37}{742.519}{\Character{23}}%
\TextAt{9}{375.552}{737.583}{\Character{41}}%
\TextAt{8}{79.226}{714.174}{be}%
\TextAt{8}{97.495}{714.174}{given}%
\TextAt{8}{130.409}{714.174}{for}%
\TextAt{8}{150.639}{714.174}{every}%
\TextAt{13}{183.875}{714.174}{\Character{23}}%
\NativeMathAt{197.73}{714.174}{11.955}{\geq}%
\TextAt{9}{214.344}{714.174}{1}%
\TextAt{8}{220.197}{714.174}{.}%
\TextAt{8}{235.68}{714.174}{If}%
\TextAt{8}{249.716}{714.174}{the}%
\TextAt{8}{272.217}{714.174}{elements}%
\TextAt{9}{322.418}{714.174}{\Character{40}}%
\TextAt{13}{326.971}{714.174}{x}%
\TextAt{17}{333.623}{720.285}{\Character{40}}%
\TextAt{21}{336.916}{720.285}{\Character{23}}%
\TextAt{17}{341.6}{720.285}{\Character{41}}%
\TextAt{21}{333.623}{711.11}{i}%
\TextAt{9}{345.391}{714.174}{\Character{41}}%
\TextAt{21}{349.944}{712.381}{\Character{23}}%
\NativeMathAt{354.628}{712.381}{7.97}{\in}%
\TextAt{25}{360.273}{712.381}{N}%
\NativeMathAt{373.844}{714.174}{11.955}{\in}%
\TextAt{20}{389.13}{723.141}{\Character{81}}%
\NativeMathAt{400.421}{719.986}{7.97}{\infty}%
\TextAt{21}{400.421}{710.687}{\Character{23}}%
\TextAt{17}{405.105}{710.687}{\Character{61}1}%
\TextAt{18}{418.417}{714.174}{Z}%
\TextAt{9}{426.387}{714.174}{\Character{47}}%
\TextAt{13}{432.24}{714.174}{p}%
\TextAt{21}{438.115}{718.513}{\Character{23}}%
\TextAt{18}{443.297}{714.174}{Z}%
\TextAt{8}{457.509}{714.174}{are}%
\TextAt{8}{479.364}{714.174}{already}%
\TextAt{8}{523.014}{714.174}{in}%
\NativeMathAt{79.226}{697.128}{11.955}{\varprojlim}%
\TextAt{21}{95.485}{691.689}{\Character{23}}%
\TextAt{18}{102.66}{697.128}{Z}%
\TextAt{9}{110.63}{697.128}{\Character{47}}%
\TextAt{13}{116.483}{697.128}{p}%
\TextAt{21}{122.358}{701.467}{\Character{23}}%
\TextAt{18}{127.54}{697.128}{Z}%
\TextAt{8}{135.51}{697.128}{,}%
\TextAt{8}{143.598}{697.128}{for}%
\TextAt{8}{162.233}{697.128}{all}%
\TextAt{13}{179.243}{697.128}{i}%
\TextAt{9}{187.836}{697.128}{\Character{61}}%
\TextAt{9}{201.525}{697.128}{1}%
\TextAt{13}{207.394}{697.128}{\Character{59}}%
\TextAt{13}{212.618}{697.128}{\Character{58}}%
\TextAt{13}{217.855}{697.128}{\Character{58}}%
\TextAt{13}{223.079}{697.128}{\Character{58}}%
\TextAt{13}{228.316}{697.128}{\Character{59}}%
\TextAt{13}{233.552}{697.128}{n}%
\TextAt{8}{240.602}{697.128}{,}%
\TextAt{8}{248.69}{697.128}{then}%
\TextAt{8}{276.101}{697.128}{we}%
\TextAt{8}{294.079}{697.128}{have}%
\TextAt{8}{321.824}{697.128}{a}%
\TextAt{13}{332.33}{697.128}{p}%
\TextAt{8}{338.205}{697.128}{-adic}%
\TextAt{8}{367.565}{697.128}{solution}%
\TextAt{8}{412.609}{697.128}{of}%
\TextAt{8}{426.69}{697.128}{the}%
\TextAt{8}{447.597}{697.128}{equation}%
\TextAt{13}{496.156}{697.128}{F}%
\TextAt{9}{509.96}{697.128}{\Character{61}}%
\TextAt{9}{523.649}{697.128}{0}%
\TextAt{8}{529.518}{697.128}{.}%
\TextAt{8}{79.226}{680.082}{But}%
\TextAt{8}{102.496}{680.082}{this}%
\TextAt{8}{125.347}{680.082}{is}%
\TextAt{8}{137.139}{680.082}{not}%
\TextAt{8}{157.982}{680.082}{automatically}%
\TextAt{8}{231.846}{680.082}{the}%
\TextAt{8}{252.031}{680.082}{case.}%
\TextAt{8}{281.439}{680.082}{We}%
\TextAt{8}{301.612}{680.082}{will}%
\TextAt{8}{323.745}{680.082}{therefore}%
\TextAt{8}{372.874}{680.082}{extract}%
\TextAt{8}{412.905}{680.082}{a}%
\TextAt{8}{422.689}{680.082}{subsequence}%
\TextAt{8}{488.841}{680.082}{from}%
\TextAt{8}{516.51}{680.082}{the}%
\TextAt{8}{79.226}{663.036}{sequence}%
\TextAt{9}{128.788}{663.036}{\Character{40}}%
\TextAt{13}{133.34}{663.036}{x}%
\TextAt{17}{139.992}{669.147}{\Character{40}}%
\TextAt{21}{143.286}{669.147}{\Character{23}}%
\TextAt{17}{147.969}{669.147}{\Character{41}}%
\TextAt{17}{139.992}{660.105}{1}%
\TextAt{13}{151.761}{663.036}{\Character{59}}%
\TextAt{13}{156.985}{663.036}{\Character{58}}%
\TextAt{13}{162.222}{663.036}{\Character{58}}%
\TextAt{13}{167.446}{663.036}{\Character{58}}%
\TextAt{13}{172.683}{663.036}{\Character{59}}%
\TextAt{13}{177.919}{663.036}{x}%
\TextAt{17}{184.634}{669.147}{\Character{40}}%
\TextAt{21}{187.927}{669.147}{\Character{23}}%
\TextAt{17}{192.611}{669.147}{\Character{41}}%
\TextAt{21}{184.634}{661.81}{n}%
\TextAt{9}{196.402}{663.036}{\Character{41}}%
\TextAt{8}{205.912}{663.036}{which}%
\TextAt{8}{240.458}{663.036}{fits}%
\TextAt{8}{261.088}{663.036}{our}%
\TextAt{8}{282.962}{663.036}{needs.}%
\TextAt{8}{322.612}{663.036}{For}%
\TextAt{8}{344.641}{663.036}{simplicity}%
\TextAt{8}{399.092}{663.036}{of}%
\TextAt{8}{413.482}{663.036}{notation}%
\TextAt{8}{461.382}{663.036}{we}%
\TextAt{8}{479.669}{663.036}{only}%
\TextAt{8}{506.42}{663.036}{carry}%
\TextAt{8}{79.226}{645.99}{this}%
\TextAt{8}{102.366}{645.99}{out}%
\TextAt{8}{123.497}{645.99}{in}%
\TextAt{8}{137.467}{645.99}{the}%
\TextAt{8}{157.941}{645.99}{case}%
\TextAt{13}{183.03}{645.99}{n}%
\TextAt{9}{193.879}{645.99}{\Character{61}}%
\TextAt{9}{206.838}{645.99}{1}%
\TextAt{8}{212.698}{645.99}{,}%
\TextAt{8}{220.244}{645.99}{writing}%
\TextAt{13}{260.887}{645.99}{x}%
\TextAt{21}{267.539}{644.197}{\Character{23}}%
\TextAt{9}{276.582}{645.99}{\Character{61}}%
\TextAt{13}{289.549}{645.99}{x}%
\TextAt{17}{296.201}{652.101}{\Character{40}}%
\TextAt{21}{299.494}{652.101}{\Character{23}}%
\TextAt{17}{304.178}{652.101}{\Character{41}}%
\TextAt{17}{296.201}{643.059}{1}%
\TextAt{8}{307.969}{645.99}{.}%
\TextAt{8}{317.37}{645.99}{The}%
\TextAt{8}{341.741}{645.99}{general}%
\TextAt{8}{382.383}{645.99}{case}%
\TextAt{8}{407.472}{645.99}{follows}%
\TextAt{8}{446.225}{645.99}{exactly}%
\TextAt{8}{486.867}{645.99}{the}%
\TextAt{8}{507.341}{645.99}{same}%
\TextAt{8}{79.226}{628.944}{pattern.}%
\TextAt{8}{127.89}{628.944}{In}%
\TextAt{8}{143.339}{628.944}{what}%
\TextAt{8}{173.434}{628.944}{follows,}%
\TextAt{8}{216.157}{628.944}{we}%
\TextAt{8}{234.212}{628.944}{view}%
\TextAt{9}{262.023}{628.944}{\Character{40}}%
\TextAt{13}{266.576}{628.944}{x}%
\TextAt{21}{273.228}{627.151}{\Character{23}}%
\TextAt{9}{278.41}{628.944}{\Character{41}}%
\TextAt{8}{287.688}{628.944}{as}%
\TextAt{8}{302.886}{628.944}{a}%
\TextAt{8}{313.47}{628.944}{sequence}%
\TextAt{8}{362.8}{628.944}{in}%
\TextAt{18}{377.281}{628.944}{Z}%
\TextAt{8}{385.251}{628.944}{.}%
\TextAt{8}{396.184}{628.944}{Since}%
\TextAt{18}{427.57}{628.944}{Z}%
\TextAt{9}{435.54}{628.944}{\Character{47}}%
\TextAt{13}{441.393}{628.944}{p}%
\TextAt{18}{447.268}{628.944}{Z}%
\TextAt{8}{459.964}{628.944}{is}%
\TextAt{8}{472.556}{628.944}{finite,}%
\TextAt{8}{506.755}{628.944}{there}%
\TextAt{8}{79.226}{611.898}{are}%
\TextAt{8}{98.617}{611.898}{infinitely}%
\TextAt{8}{147.598}{611.898}{many}%
\TextAt{8}{179.35}{611.898}{terms}%
\TextAt{13}{211.809}{611.898}{x}%
\TextAt{21}{218.461}{610.105}{\Character{23}}%
\TextAt{8}{227.42}{611.898}{which}%
\TextAt{8}{260.787}{611.898}{mod}%
\TextAt{13}{287.005}{611.898}{p}%
\TextAt{8}{296.658}{611.898}{are}%
\TextAt{8}{316.049}{611.898}{congruent}%
\TextAt{8}{370.241}{611.898}{to}%
\TextAt{8}{384.432}{611.898}{the}%
\TextAt{8}{404.469}{611.898}{same}%
\TextAt{8}{433.675}{611.898}{element}%
\TextAt{13}{476.797}{611.898}{y}%
\TextAt{17}{482.505}{610.105}{1}%
\NativeMathAt{490.558}{611.898}{11.955}{\in}%
\TextAt{18}{501.849}{611.898}{Z}%
\TextAt{9}{509.819}{611.898}{\Character{47}}%
\TextAt{13}{515.672}{611.898}{p}%
\TextAt{18}{521.548}{611.898}{Z}%
\TextAt{8}{529.518}{611.898}{.}%
\TextAt{8}{79.226}{594.852}{Hence}%
\TextAt{8}{114.004}{594.852}{we}%
\TextAt{8}{131.231}{594.852}{may}%
\TextAt{8}{156.6}{594.852}{choose}%
\TextAt{8}{193.733}{594.852}{a}%
\TextAt{8}{203.488}{594.852}{subsequence}%
\TextAt{16}{269.612}{594.852}{\Character{102}}%
\TextAt{13}{275.59}{594.852}{x}%
\TextAt{17}{282.242}{600.963}{\Character{40}1\Character{41}}%
\TextAt{21}{282.242}{593.626}{\Character{23}}%
\TextAt{16}{293.561}{594.852}{\Character{103}}%
\TextAt{8}{303.436}{594.852}{of}%
\TextAt{16}{316.766}{594.852}{\Character{102}}%
\TextAt{13}{322.743}{594.852}{x}%
\TextAt{21}{329.395}{593.059}{\Character{23}}%
\TextAt{16}{334.577}{594.852}{\Character{103}}%
\TextAt{8}{344.452}{594.852}{such}%
\TextAt{8}{370.85}{594.852}{that}%
\TextAt{13}{166.78}{571.444}{x}%
\TextAt{17}{173.433}{576.38}{\Character{40}1\Character{41}}%
\TextAt{21}{173.433}{568.489}{\Character{23}}%
\NativeMathAt{188.072}{571.444}{11.955}{\equiv}%
\TextAt{13}{200.692}{571.444}{y}%
\TextAt{17}{206.399}{569.651}{1}%
\TextAt{9}{223.087}{571.444}{\Character{40}}%
\TextAt{8}{227.639}{571.444}{mod}%
\TextAt{13}{253.741}{571.444}{p}%
\TextAt{9}{259.616}{571.444}{\Character{41}}%
\TextAt{8}{288.079}{571.444}{and}%
\TextAt{13}{330.854}{571.444}{F}%
\TextAt{9}{340.058}{571.444}{\Character{40}}%
\TextAt{13}{344.61}{571.444}{x}%
\TextAt{17}{351.262}{576.38}{\Character{40}1\Character{41}}%
\TextAt{21}{351.262}{568.489}{\Character{23}}%
\TextAt{9}{362.581}{571.444}{\Character{41}}%
\NativeMathAt{370.454}{571.444}{11.955}{\equiv}%
\TextAt{9}{383.074}{571.444}{0}%
\TextAt{9}{400.863}{571.444}{\Character{40}}%
\TextAt{8}{405.434}{571.444}{mod}%
\TextAt{13}{431.536}{571.444}{p}%
\TextAt{9}{437.411}{571.444}{\Character{41}}%
\TextAt{13}{441.964}{571.444}{\Character{58}}%
\TextAt{8}{79.226}{548.036}{Likewise,}%
\TextAt{8}{129.522}{548.036}{we}%
\TextAt{8}{146.749}{548.036}{may}%
\TextAt{8}{172.118}{548.036}{extract}%
\TextAt{8}{212.12}{548.036}{from}%
\TextAt{16}{239.76}{548.036}{\Character{102}}%
\TextAt{13}{245.738}{548.036}{x}%
\TextAt{17}{252.39}{554.146}{\Character{40}1\Character{41}}%
\TextAt{21}{252.39}{546.809}{\Character{23}}%
\TextAt{16}{263.709}{548.036}{\Character{103}}%
\TextAt{8}{273.584}{548.036}{a}%
\TextAt{8}{283.339}{548.036}{subsequence}%
\TextAt{16}{349.463}{548.036}{\Character{102}}%
\TextAt{13}{355.441}{548.036}{x}%
\TextAt{17}{362.093}{554.146}{\Character{40}2\Character{41}}%
\TextAt{21}{362.093}{546.809}{\Character{23}}%
\TextAt{16}{373.412}{548.036}{\Character{103}}%
\TextAt{8}{383.287}{548.036}{such}%
\TextAt{8}{409.684}{548.036}{that}%
\TextAt{13}{162.048}{524.628}{x}%
\TextAt{17}{168.7}{529.564}{\Character{40}2\Character{41}}%
\TextAt{21}{168.7}{521.672}{\Character{23}}%
\NativeMathAt{183.34}{524.628}{11.955}{\equiv}%
\TextAt{13}{195.959}{524.628}{y}%
\TextAt{17}{201.667}{522.834}{2}%
\TextAt{9}{218.354}{524.628}{\Character{40}}%
\TextAt{8}{222.907}{524.628}{mod}%
\TextAt{13}{249.009}{524.628}{p}%
\TextAt{17}{254.884}{529.564}{2}%
\TextAt{9}{259.616}{524.628}{\Character{41}}%
\TextAt{8}{288.079}{524.628}{and}%
\TextAt{13}{330.854}{524.628}{F}%
\TextAt{9}{340.058}{524.628}{\Character{40}}%
\TextAt{13}{344.61}{524.628}{x}%
\TextAt{17}{351.262}{529.564}{\Character{40}2\Character{41}}%
\TextAt{21}{351.262}{521.672}{\Character{23}}%
\TextAt{9}{362.581}{524.628}{\Character{41}}%
\NativeMathAt{370.454}{524.628}{11.955}{\equiv}%
\TextAt{9}{383.074}{524.628}{0}%
\TextAt{9}{400.863}{524.628}{\Character{40}}%
\TextAt{8}{405.434}{524.628}{mod}%
\TextAt{13}{431.536}{524.628}{p}%
\TextAt{17}{437.411}{529.564}{2}%
\TextAt{9}{442.144}{524.628}{\Character{41}}%
\TextAt{13}{446.696}{524.628}{\Character{59}}%
\TextAt{8}{79.226}{501.219}{where}%
\TextAt{13}{112.477}{501.219}{y}%
\TextAt{17}{118.185}{499.426}{2}%
\NativeMathAt{126.238}{501.219}{11.955}{\in}%
\TextAt{18}{137.529}{501.219}{Z}%
\TextAt{9}{145.499}{501.219}{\Character{47}}%
\TextAt{13}{151.352}{501.219}{p}%
\TextAt{17}{157.227}{505.558}{2}%
\TextAt{18}{161.96}{501.219}{Z}%
\TextAt{8}{173.269}{501.219}{evidently}%
\TextAt{8}{223.114}{501.219}{satisfies}%
\TextAt{13}{265.666}{501.219}{y}%
\TextAt{17}{271.373}{499.426}{2}%
\NativeMathAt{279.427}{501.219}{11.955}{\equiv}%
\TextAt{13}{292.046}{501.219}{y}%
\TextAt{17}{297.754}{499.426}{1}%
\TextAt{8}{305.807}{501.219}{mod}%
\TextAt{13}{331.245}{501.219}{p}%
\TextAt{8}{337.12}{501.219}{.}%
\TextAt{8}{345.382}{501.219}{Continuing}%
\TextAt{8}{405.317}{501.219}{in}%
\TextAt{8}{418.412}{501.219}{this}%
\TextAt{8}{440.676}{501.219}{way,}%
\TextAt{8}{466.244}{501.219}{we}%
\TextAt{8}{482.913}{501.219}{obtain}%
\TextAt{8}{518.782}{501.219}{for}%
\TextAt{8}{79.226}{484.173}{each}%
\TextAt{13}{105.563}{484.173}{k}%
\NativeMathAt{115.374}{484.173}{11.955}{\geq}%
\TextAt{9}{127.993}{484.173}{1}%
\TextAt{8}{137.743}{484.173}{a}%
\TextAt{8}{147.499}{484.173}{subsequence}%
\TextAt{16}{213.623}{484.173}{\Character{102}}%
\TextAt{13}{219.601}{484.173}{x}%
\TextAt{17}{226.253}{490.284}{\Character{40}}%
\TextAt{21}{229.546}{490.284}{k}%
\TextAt{17}{234.167}{490.284}{\Character{41}}%
\TextAt{21}{226.253}{482.947}{\Character{23}}%
\TextAt{16}{237.959}{484.173}{\Character{103}}%
\TextAt{8}{247.834}{484.173}{of}%
\TextAt{16}{261.164}{484.173}{\Character{102}}%
\TextAt{13}{267.141}{484.173}{x}%
\TextAt{17}{273.794}{490.284}{\Character{40}}%
\TextAt{21}{277.087}{490.284}{k}%
\TextAt{3}{281.708}{490.284}{\Character{0}}%
\TextAt{17}{288.295}{490.284}{1\Character{41}}%
\TextAt{21}{273.794}{482.947}{\Character{23}}%
\TextAt{16}{296.321}{484.173}{\Character{103}}%
\TextAt{8}{306.196}{484.173}{the}%
\TextAt{8}{326.352}{484.173}{terms}%
\TextAt{8}{358.93}{484.173}{of}%
\TextAt{8}{372.26}{484.173}{which}%
\TextAt{8}{405.746}{484.173}{satisfy}%
\TextAt{8}{442.293}{484.173}{the}%
\TextAt{8}{462.45}{484.173}{congruences}%
\TextAt{13}{162.705}{460.765}{x}%
\TextAt{17}{169.357}{465.701}{\Character{40}}%
\TextAt{21}{172.651}{465.701}{k}%
\TextAt{17}{177.272}{465.701}{\Character{41}}%
\TextAt{21}{169.357}{457.809}{\Character{23}}%
\NativeMathAt{184.385}{460.765}{11.955}{\equiv}%
\TextAt{13}{197.004}{460.765}{y}%
\TextAt{21}{202.712}{458.972}{k}%
\TextAt{9}{219.786}{460.765}{\Character{40}}%
\TextAt{8}{224.339}{460.765}{mod}%
\TextAt{13}{250.441}{460.765}{p}%
\TextAt{21}{256.316}{465.701}{k}%
\TextAt{9}{261.436}{460.765}{\Character{41}}%
\TextAt{8}{289.898}{460.765}{and}%
\TextAt{13}{332.674}{460.765}{F}%
\TextAt{9}{341.878}{460.765}{\Character{40}}%
\TextAt{13}{346.43}{460.765}{x}%
\TextAt{17}{353.082}{465.701}{\Character{40}}%
\TextAt{21}{356.375}{465.701}{k}%
\TextAt{17}{360.997}{465.701}{\Character{41}}%
\TextAt{21}{353.082}{457.809}{\Character{23}}%
\TextAt{9}{364.788}{460.765}{\Character{41}}%
\NativeMathAt{372.661}{460.765}{11.955}{\equiv}%
\TextAt{9}{385.281}{460.765}{0}%
\TextAt{9}{403.07}{460.765}{\Character{40}}%
\TextAt{8}{407.641}{460.765}{mod}%
\TextAt{13}{433.743}{460.765}{p}%
\TextAt{21}{439.618}{465.701}{k}%
\TextAt{9}{444.738}{460.765}{\Character{41}}%
\TextAt{8}{79.226}{437.357}{for}%
\TextAt{8}{97.111}{437.357}{some}%
\TextAt{13}{126.437}{437.357}{y}%
\TextAt{21}{132.145}{435.563}{k}%
\NativeMathAt{140.585}{437.357}{11.955}{\in}%
\TextAt{18}{151.876}{437.357}{Z}%
\TextAt{9}{159.847}{437.357}{\Character{47}}%
\TextAt{13}{165.7}{437.357}{p}%
\TextAt{21}{171.575}{441.695}{k}%
\TextAt{18}{176.694}{437.357}{Z}%
\TextAt{8}{188.562}{437.357}{such}%
\TextAt{8}{214.959}{437.357}{that}%
\TextAt{13}{245.675}{413.948}{y}%
\TextAt{21}{251.383}{412.155}{k}%
\NativeMathAt{259.824}{413.948}{11.955}{\equiv}%
\TextAt{13}{272.443}{413.948}{y}%
\TextAt{21}{278.151}{412.155}{k}%
\TextAt{3}{282.772}{412.155}{\Character{0}}%
\TextAt{17}{289.359}{412.155}{1}%
\TextAt{9}{306.046}{413.948}{\Character{40}}%
\TextAt{8}{310.599}{413.948}{mod}%
\TextAt{13}{336.701}{413.948}{p}%
\TextAt{21}{342.576}{418.885}{k}%
\TextAt{3}{347.197}{418.885}{\Character{0}}%
\TextAt{17}{353.784}{418.885}{1}%
\TextAt{9}{358.516}{413.948}{\Character{41}}%
\TextAt{13}{363.069}{413.948}{\Character{58}}%
\TextAt{8}{79.226}{390.54}{The}%
\TextAt{13}{103.28}{390.54}{y}%
\TextAt{21}{108.988}{388.747}{k}%
\TextAt{8}{118.005}{390.54}{define}%
\TextAt{8}{151.814}{390.54}{a}%
\TextAt{13}{161.569}{390.54}{p}%
\TextAt{8}{167.445}{390.54}{-adic}%
\TextAt{8}{196.053}{390.54}{integer}%
\TextAt{13}{234.752}{390.54}{y}%
\TextAt{9}{244.21}{390.54}{\Character{61}}%
\TextAt{9}{256.62}{390.54}{\Character{40}}%
\TextAt{13}{261.187}{390.54}{y}%
\TextAt{21}{266.895}{388.747}{k}%
\TextAt{9}{272.015}{390.54}{\Character{41}}%
\TextAt{21}{276.567}{388.747}{k}%
\NativeMathAt{281.189}{388.747}{7.97}{\in}%
\TextAt{25}{286.834}{388.747}{N}%
\NativeMathAt{296.41}{390.54}{11.955}{\in}%
\NativeMathAt{307.701}{390.54}{11.955}{\varprojlim}%
\TextAt{21}{323.96}{385.101}{k}%
\TextAt{18}{331.072}{390.54}{Z}%
\TextAt{9}{339.042}{390.54}{\Character{47}}%
\TextAt{13}{344.895}{390.54}{p}%
\TextAt{21}{350.77}{394.878}{k}%
\TextAt{18}{355.89}{390.54}{Z}%
\TextAt{9}{367.181}{390.54}{\Character{61}}%
\TextAt{18}{379.606}{390.54}{Z}%
\TextAt{21}{387.576}{388.747}{p}%
\TextAt{8}{396.235}{390.54}{satisfying}%
\TextAt{13}{251.455}{367.132}{F}%
\TextAt{9}{260.659}{367.132}{\Character{40}}%
\TextAt{13}{265.211}{367.132}{y}%
\TextAt{21}{270.919}{365.338}{k}%
\TextAt{9}{276.038}{367.132}{\Character{41}}%
\NativeMathAt{283.912}{367.132}{11.955}{\equiv}%
\TextAt{9}{296.531}{367.132}{0}%
\TextAt{9}{314.32}{367.132}{\Character{40}}%
\TextAt{8}{318.891}{367.132}{mod}%
\TextAt{13}{344.993}{367.132}{p}%
\TextAt{21}{350.868}{372.068}{k}%
\TextAt{9}{355.988}{367.132}{\Character{41}}%
\TextAt{8}{79.226}{343.723}{for}%
\TextAt{8}{97.111}{343.723}{all}%
\TextAt{13}{113.37}{343.723}{k}%
\NativeMathAt{123.18}{343.723}{11.955}{\geq}%
\TextAt{9}{135.8}{343.723}{1}%
\TextAt{8}{141.653}{343.723}{.}%
\TextAt{8}{150.101}{343.723}{In}%
\TextAt{8}{164.722}{343.723}{other}%
\TextAt{8}{195.292}{343.723}{words,}%
\TextAt{13}{232.102}{343.723}{F}%
\TextAt{9}{241.305}{343.723}{\Character{40}}%
\TextAt{13}{245.858}{343.723}{y}%
\TextAt{9}{251.994}{343.723}{\Character{41}}%
\TextAt{9}{259.849}{343.723}{\Character{61}}%
\TextAt{9}{272.27}{343.723}{0}%
\TextAt{8}{278.146}{343.723}{.}%
\Vector{q 0 G 0 g 0.399 w 452.334 -430.277 m 452.334 -422.207 l S Q}%
\Vector{q 0 G 0 g 0.399 w 452.534 -422.406 m 459.707 -422.406 l S Q}%
\Vector{q 0 G 0 g 0.399 w 452.534 -430.077 m 459.707 -430.077 l S Q}%
\Vector{q 0 G 0 g 0.399 w 459.906 -430.277 m 459.906 -422.207 l S Q}%
\TextAt{12}{79.226}{311.941}{Exercise}%
\TextAt{12}{131.606}{311.941}{II.1.1.}%
\TextAt{8}{169.866}{311.941}{A}%
\TextAt{13}{181.825}{311.941}{p}%
\TextAt{8}{187.7}{311.941}{-adic}%
\TextAt{8}{215.595}{311.941}{number}%
\TextAt{13}{257.478}{311.941}{a}%
\TextAt{9}{266.943}{311.941}{\Character{61}}%
\TextAt{20}{279.369}{320.907}{\Character{80}}%
\NativeMathAt{291.988}{317.752}{7.97}{\infty}%
\TextAt{21}{291.988}{308.454}{\Character{23}}%
\TextAt{17}{296.672}{308.454}{\Character{61}}%
\TextAt{3}{303.259}{308.454}{\Character{0}}%
\TextAt{21}{309.845}{308.454}{m}%
\TextAt{13}{319.826}{311.941}{a}%
\TextAt{21}{325.971}{310.147}{\Character{23}}%
\TextAt{13}{331.153}{311.941}{p}%
\TextAt{21}{337.028}{316.279}{\Character{23}}%
\NativeMathAt{345.531}{311.941}{11.955}{\in}%
\TextAt{18}{356.822}{311.941}{Q}%
\TextAt{21}{366.121}{310.147}{p}%
\TextAt{8}{374.066}{311.941}{is}%
\TextAt{8}{385.116}{311.941}{a}%
\TextAt{8}{394.158}{311.941}{rational}%
\TextAt{8}{437.033}{311.941}{number}%
\TextAt{8}{478.916}{311.941}{if}%
\TextAt{8}{488.926}{311.941}{and}%
\TextAt{8}{510.975}{311.941}{only}%
\TextAt{8}{79.226}{294.895}{if}%
\TextAt{8}{89.95}{294.895}{the}%
\TextAt{8}{110.106}{294.895}{sequence}%
\TextAt{8}{158.608}{294.895}{of}%
\TextAt{8}{171.939}{294.895}{digits}%
\TextAt{8}{203.871}{294.895}{is}%
\TextAt{8}{215.635}{294.895}{periodic}%
\TextAt{8}{260.502}{294.895}{(possibly}%
\TextAt{8}{310.057}{294.895}{with}%
\TextAt{8}{336.717}{294.895}{a}%
\TextAt{8}{346.472}{294.895}{finite}%
\TextAt{8}{376.384}{294.895}{string}%
\TextAt{8}{409.619}{294.895}{before}%
\TextAt{8}{444.732}{294.895}{the}%
\TextAt{8}{464.888}{294.895}{first}%
\TextAt{8}{489.014}{294.895}{period).}%
\TextAt{12}{97.159}{274.86}{Hint:}%
\TextAt{8}{132.698}{274.86}{Write}%
\TextAt{13}{165.204}{274.86}{p}%
\TextAt{21}{171.079}{279.198}{m}%
\TextAt{13}{179.068}{274.86}{a}%
\TextAt{9}{188.533}{274.86}{\Character{61}}%
\TextAt{13}{200.959}{274.86}{b}%
\TextAt{9}{208.593}{274.86}{\Character{43}}%
\TextAt{13}{220.354}{274.86}{c}%
\TextAt{21}{232.748}{280.116}{p}%
\TextAt{22}{237.011}{282.928}{\Character{96}}%
\Vector{q 0 G 0 g 0.478 w 154.587 -496.151 m 174.77 -496.151 l S Q}%
\TextAt{17}{226.587}{270.737}{1}%
\TextAt{3}{230.822}{270.737}{\Character{0}}%
\TextAt{21}{237.408}{270.737}{p}%
\TextAt{22}{241.671}{273.006}{n}%
\TextAt{8}{247.965}{274.86}{,}%
\TextAt{9}{255.114}{274.86}{0}%
\NativeMathAt{264.288}{274.86}{11.955}{\leq}%
\TextAt{13}{276.908}{274.86}{b}%
\TextAt{13}{285.193}{274.86}{\Character{60}}%
\TextAt{13}{297.614}{274.86}{p}%
\TextAt{21}{303.506}{279.198}{\Character{96}}%
\TextAt{8}{307.464}{274.86}{,}%
\TextAt{9}{314.613}{274.86}{0}%
\NativeMathAt{323.787}{274.86}{11.955}{\leq}%
\TextAt{13}{336.406}{274.86}{c}%
\TextAt{13}{344.751}{274.86}{\Character{60}}%
\TextAt{13}{357.172}{274.86}{p}%
\TextAt{21}{363.065}{279.198}{n}%
\TextAt{8}{368.702}{274.86}{.}%
\TextAt{12}{79.226}{246.066}{Exercise}%
\TextAt{12}{132.039}{246.066}{II.1.2.}%
\TextAt{8}{170.676}{246.066}{A}%
\TextAt{13}{183.011}{246.066}{p}%
\TextAt{8}{188.886}{246.066}{-adic}%
\TextAt{8}{217.157}{246.066}{integer}%
\TextAt{13}{255.519}{246.066}{a}%
\TextAt{9}{264.985}{246.066}{\Character{61}}%
\TextAt{13}{277.41}{246.066}{a}%
\TextAt{17}{283.555}{244.273}{0}%
\TextAt{9}{290.255}{246.066}{\Character{43}}%
\TextAt{13}{301.327}{246.066}{a}%
\TextAt{17}{307.472}{244.273}{1}%
\TextAt{13}{312.204}{246.066}{p}%
\TextAt{9}{320.047}{246.066}{\Character{43}}%
\TextAt{13}{331.119}{246.066}{a}%
\TextAt{17}{337.264}{244.273}{2}%
\TextAt{13}{341.996}{246.066}{p}%
\TextAt{17}{347.871}{250.404}{2}%
\TextAt{9}{354.571}{246.066}{\Character{43}}%
\TextAt{16}{365.643}{246.066}{\Character{1}}%
\TextAt{16}{370.939}{246.066}{\Character{1}}%
\TextAt{16}{376.247}{246.066}{\Character{1}}%
\TextAt{8}{385.143}{246.066}{is}%
\TextAt{8}{396.57}{246.066}{a}%
\TextAt{8}{405.988}{246.066}{unit}%
\TextAt{8}{430.362}{246.066}{in}%
\TextAt{8}{443.678}{246.066}{the}%
\TextAt{8}{463.498}{246.066}{ring}%
\TextAt{18}{487.226}{246.066}{Z}%
\TextAt{21}{495.196}{244.273}{p}%
\TextAt{8}{503.518}{246.066}{if}%
\TextAt{8}{513.904}{246.066}{and}%
\TextAt{8}{79.226}{229.02}{only}%
\TextAt{8}{104.918}{229.02}{if}%
\TextAt{13}{115.642}{229.02}{a}%
\TextAt{17}{121.787}{227.227}{0}%
\NativeMathAt{129.84}{229.02}{11.955}{\neq}%
\TextAt{9}{142.25}{229.02}{0}%
\TextAt{8}{148.118}{229.02}{.}%
\TextAt{12}{79.226}{200.226}{Exercise}%
\TextAt{12}{132.427}{200.226}{II.1.3.}%
\TextAt{8}{171.401}{200.226}{Show}%
\TextAt{8}{202.293}{200.226}{that}%
\TextAt{8}{227.662}{200.226}{the}%
\TextAt{8}{247.818}{200.226}{equation}%
\TextAt{13}{295.627}{200.226}{x}%
\TextAt{17}{302.279}{204.564}{2}%
\TextAt{9}{310.332}{200.226}{\Character{61}}%
\TextAt{9}{322.742}{200.226}{2}%
\TextAt{8}{332.508}{200.226}{has}%
\TextAt{8}{353.382}{200.226}{a}%
\TextAt{8}{363.137}{200.226}{solution}%
\TextAt{8}{407.431}{200.226}{in}%
\TextAt{18}{421.084}{200.226}{Z}%
\TextAt{17}{429.054}{198.433}{7}%
\TextAt{8}{433.786}{200.226}{.}%
\TextAt{12}{79.226}{171.432}{Exercise}%
\TextAt{12}{132.427}{171.432}{II.1.4.}%
\TextAt{8}{171.401}{171.432}{Write}%
\TextAt{8}{203.907}{171.432}{the}%
\TextAt{8}{224.063}{171.432}{numbers}%
\TextAt{17}{272.47}{176.139}{2}%
\Vector{q 0 G 0 g 0.478 w 200.47 -599.579 m 204.704 -599.579 l S Q}%
\TextAt{17}{272.47}{167.309}{3}%
\TextAt{8}{281.797}{171.432}{and}%
\TextAt{16}{304.559}{171.432}{\Character{0}}%
\TextAt{17}{315.053}{176.139}{2}%
\Vector{q 0 G 0 g 0.478 w 243.053 -599.579 m 247.287 -599.579 l S Q}%
\TextAt{17}{315.053}{167.309}{3}%
\TextAt{8}{324.38}{171.432}{as}%
\TextAt{8}{338.75}{171.432}{5-adic}%
\TextAt{8}{373.217}{171.432}{numbers.}%
\TextAt{12}{79.226}{142.638}{Exercise}%
\TextAt{12}{132.209}{142.638}{II.1.5.}%
\TextAt{8}{170.993}{142.638}{The}%
\TextAt{8}{194.857}{142.638}{field}%
\TextAt{18}{220.025}{142.638}{Q}%
\TextAt{21}{229.323}{140.845}{p}%
\TextAt{8}{237.792}{142.638}{of}%
\TextAt{13}{250.933}{142.638}{p}%
\TextAt{8}{256.808}{142.638}{-adic}%
\TextAt{8}{285.227}{142.638}{numbers}%
\TextAt{8}{332.248}{142.638}{has}%
\TextAt{8}{352.933}{142.638}{no}%
\TextAt{8}{369.002}{142.638}{automorphisms}%
\TextAt{8}{450.897}{142.638}{except}%
\TextAt{8}{487.446}{142.638}{the}%
\TextAt{8}{507.413}{142.638}{iden-}%
\TextAt{8}{79.226}{125.592}{tity.}%
\TextAt{12}{79.226}{96.798}{Exercise}%
\TextAt{12}{131.787}{96.798}{II.1.6.}%
\TextAt{8}{170.205}{96.798}{How}%
\TextAt{8}{196.309}{96.798}{is}%
\TextAt{8}{207.517}{96.798}{the}%
\TextAt{8}{227.117}{96.798}{addition,}%
\TextAt{8}{276.107}{96.798}{subtraction,}%
\TextAt{8}{340.77}{96.798}{multiplication}%
\TextAt{8}{415.34}{96.798}{and}%
\TextAt{8}{437.547}{96.798}{division}%
\TextAt{8}{480.304}{96.798}{of}%
\TextAt{8}{493.078}{96.798}{rational}%
\TextAt{8}{79.226}{79.752}{numbers}%
\TextAt{8}{126.437}{79.752}{reflected}%
\TextAt{8}{173.254}{79.752}{in}%
\TextAt{8}{186.906}{79.752}{the}%
\TextAt{8}{207.063}{79.752}{representation}%
\TextAt{8}{283.552}{79.752}{by}%
\TextAt{13}{299.811}{79.752}{p}%
\TextAt{8}{305.686}{79.752}{-adic}%
\TextAt{8}{334.295}{79.752}{digits?}%
\TextAt{8}{300.14}{43.033}{97}%
\Destination{78.23}{809.815}{nav.page.97}%
\Destination{79.226}{325.955}{nav.ex.2.1.1}%
\Destination{79.226}{260.055}{nav.ex.2.1.2}%
\Destination{79.226}{214.255}{nav.ex.2.1.3}%
\Destination{79.226}{185.455}{nav.ex.2.1.4}%
\Destination{79.226}{156.655}{nav.ex.2.1.5}%
\Destination{79.226}{110.755}{nav.ex.2.1.6}%
\Destination{79.36}{586.455}{nav.print.106}%
}

% Page 110 (body page 98)
\Page{%
\TextAt{2}{79.226}{796.943}{II.}%
\TextAt{2}{94.073}{796.943}{The}%
\TextAt{2}{116.492}{796.943}{Theory}%
\TextAt{2}{154.401}{796.943}{of}%
\TextAt{2}{166.826}{796.943}{Valuations}%
\Vector{q 0 G 0 g 0.299 w 7.226 18.729 m 460.769 18.729 l S Q}%
\TextAt{15}{79.226}{760.804}{§}%
\TextAt{15}{93.785}{760.804}{2}%
\TextAt{15}{120.693}{760.804}{The}%
\TextAt{31}{160.202}{760.804}{p}%
\TextAt{15}{168.663}{760.804}{-adic}%
\TextAt{15}{215.747}{760.804}{Absolute}%
\TextAt{15}{297.451}{760.804}{Value}%
\TextAt{8}{79.226}{728.918}{The}%
\TextAt{8}{103.28}{728.918}{representation}%
\TextAt{8}{179.769}{728.918}{of}%
\TextAt{8}{193.099}{728.918}{a}%
\TextAt{13}{202.855}{728.918}{p}%
\TextAt{8}{208.73}{728.918}{-adic}%
\TextAt{8}{237.338}{728.918}{integer}%
\TextAt{13}{208.174}{700.387}{a}%
\TextAt{17}{214.319}{698.594}{0}%
\TextAt{9}{221.708}{700.387}{\Character{43}}%
\TextAt{13}{233.469}{700.387}{a}%
\TextAt{17}{239.614}{698.594}{1}%
\TextAt{13}{244.346}{700.387}{p}%
\TextAt{9}{252.878}{700.387}{\Character{43}}%
\TextAt{13}{264.639}{700.387}{a}%
\TextAt{17}{270.784}{698.594}{2}%
\TextAt{13}{275.517}{700.387}{p}%
\TextAt{17}{281.392}{705.323}{2}%
\TextAt{9}{288.781}{700.387}{\Character{43}}%
\TextAt{16}{300.542}{700.387}{\Character{1}}%
\TextAt{16}{305.838}{700.387}{\Character{1}}%
\TextAt{16}{311.146}{700.387}{\Character{1}}%
\TextAt{13}{318.475}{700.387}{\Character{59}}%
\TextAt{9}{347.629}{700.387}{0}%
\NativeMathAt{356.803}{700.387}{11.955}{\leq}%
\TextAt{13}{369.422}{700.387}{a}%
\TextAt{21}{375.567}{698.594}{i}%
\TextAt{13}{382.27}{700.387}{\Character{60}}%
\TextAt{13}{394.68}{700.387}{p\Character{59}}%
\TextAt{8}{517.802}{700.387}{(1)}%
\TextAt{8}{79.226}{671.856}{resembles}%
\TextAt{8}{131.697}{671.856}{very}%
\TextAt{8}{157.389}{671.856}{much}%
\TextAt{8}{188.604}{671.856}{the}%
\TextAt{8}{208.76}{671.856}{decimal}%
\TextAt{8}{251.679}{671.856}{fraction}%
\TextAt{8}{294.933}{671.856}{representation}%
\TextAt{13}{177.443}{637.619}{a}%
\TextAt{17}{183.588}{635.825}{0}%
\TextAt{9}{190.977}{637.619}{\Character{43}}%
\TextAt{13}{202.739}{637.619}{a}%
\TextAt{17}{208.883}{635.825}{1}%
\TextAt{20}{215.608}{654.476}{\Character{18}}%
\TextAt{9}{228.531}{645.706}{1}%
\Vector{q 0 G 0 g 0.478 w 153.604 -133.392 m 165.31 -133.392 l S Q}%
\TextAt{9}{225.604}{629.418}{10}%
\TextAt{20}{238.506}{654.476}{\Character{19}}%
\TextAt{9}{249.963}{637.619}{\Character{43}}%
\TextAt{13}{261.724}{637.619}{a}%
\TextAt{17}{267.869}{635.825}{2}%
\TextAt{20}{274.594}{654.476}{\Character{18}}%
\TextAt{9}{287.516}{645.706}{1}%
\Vector{q 0 G 0 g 0.478 w 212.59 -133.392 m 224.296 -133.392 l S Q}%
\TextAt{9}{284.59}{629.418}{10}%
\TextAt{20}{297.491}{654.476}{\Character{19}}%
\TextAt{17}{306.292}{651.799}{2}%
\TextAt{9}{313.681}{637.619}{\Character{43}}%
\TextAt{16}{325.442}{637.619}{\Character{1}}%
\TextAt{16}{330.738}{637.619}{\Character{1}}%
\TextAt{16}{336.046}{637.619}{\Character{1}}%
\TextAt{13}{343.375}{637.619}{\Character{59}}%
\TextAt{9}{372.529}{637.619}{0}%
\NativeMathAt{381.703}{637.619}{11.955}{\leq}%
\TextAt{13}{394.322}{637.619}{a}%
\TextAt{21}{400.467}{635.825}{i}%
\TextAt{13}{407.169}{637.619}{\Character{60}}%
\TextAt{9}{419.595}{637.619}{10}%
\TextAt{13}{431.301}{637.619}{\Character{59}}%
\TextAt{8}{79.226}{605.363}{of}%
\TextAt{8}{93.693}{605.363}{a}%
\TextAt{8}{104.585}{605.363}{real}%
\TextAt{8}{128.484}{605.363}{number}%
\TextAt{8}{172.217}{605.363}{between}%
\TextAt{8}{218.544}{605.363}{0}%
\TextAt{8}{229.436}{605.363}{and}%
\TextAt{8}{253.336}{605.363}{10.}%
\TextAt{8}{276.91}{605.363}{But}%
\TextAt{8}{301.287}{605.363}{it}%
\TextAt{8}{314.128}{605.363}{does}%
\TextAt{8}{341.662}{605.363}{not}%
\TextAt{8}{363.612}{605.363}{converge}%
\TextAt{8}{412.558}{605.363}{as}%
\TextAt{8}{428.064}{605.363}{the}%
\TextAt{8}{449.357}{605.363}{decimal}%
\TextAt{8}{493.413}{605.363}{fraction}%
\TextAt{8}{79.226}{588.317}{does.}%
\TextAt{8}{112.701}{588.317}{Nonetheless,}%
\TextAt{8}{181.182}{588.317}{the}%
\TextAt{8}{202.181}{588.317}{field}%
\TextAt{18}{228.38}{588.317}{Q}%
\TextAt{21}{237.679}{586.524}{p}%
\TextAt{8}{247.179}{588.317}{of}%
\TextAt{13}{261.352}{588.317}{p}%
\TextAt{8}{267.227}{588.317}{-adic}%
\TextAt{8}{296.678}{588.317}{numbers}%
\TextAt{8}{344.731}{588.317}{can}%
\TextAt{8}{367.033}{588.317}{be}%
\TextAt{8}{383.8}{588.317}{constructed}%
\TextAt{8}{447.79}{588.317}{from}%
\TextAt{8}{476.273}{588.317}{the}%
\TextAt{8}{497.272}{588.317}{field}%
\TextAt{18}{523.471}{588.317}{Q}%
\TextAt{8}{79.226}{571.271}{in}%
\TextAt{8}{94.167}{571.271}{the}%
\TextAt{8}{115.611}{571.271}{same}%
\TextAt{8}{146.224}{571.271}{fashion}%
\TextAt{8}{187.574}{571.271}{as}%
\TextAt{8}{203.231}{571.271}{the}%
\TextAt{8}{224.676}{571.271}{field}%
\TextAt{8}{251.32}{571.271}{of}%
\TextAt{8}{265.938}{571.271}{real}%
\TextAt{8}{289.988}{571.271}{numbers}%
\TextAt{18}{338.487}{571.271}{R}%
\TextAt{8}{347.121}{571.271}{.}%
\TextAt{8}{359.432}{571.271}{The}%
\TextAt{8}{384.774}{571.271}{key}%
\TextAt{8}{407.198}{571.271}{to}%
\TextAt{8}{422.796}{571.271}{this}%
\TextAt{8}{446.906}{571.271}{is}%
\TextAt{8}{459.958}{571.271}{to}%
\TextAt{8}{475.556}{571.271}{replace}%
\TextAt{8}{516.51}{571.271}{the}%
\TextAt{8}{79.226}{554.225}{ordinary}%
\TextAt{8}{126.912}{554.225}{absolute}%
\TextAt{8}{173.677}{554.225}{value}%
\TextAt{8}{204.446}{554.225}{by}%
\TextAt{8}{221.228}{554.225}{a}%
\TextAt{8}{231.506}{554.225}{new}%
\TextAt{8}{256.083}{554.225}{“}%
\TextAt{13}{261.606}{554.225}{p}%
\TextAt{8}{267.481}{554.225}{-adic”}%
\TextAt{8}{302.136}{554.225}{absolute}%
\TextAt{8}{348.901}{554.225}{value}%
\TextAt{16}{379.67}{554.225}{\Character{106}}%
\TextAt{16}{387.393}{554.225}{\Character{106}}%
\TextAt{21}{390.732}{552.432}{p}%
\TextAt{8}{399.913}{554.225}{with}%
\TextAt{8}{427.096}{554.225}{respect}%
\TextAt{8}{467.668}{554.225}{to}%
\TextAt{8}{482.501}{554.225}{which}%
\TextAt{8}{516.51}{554.225}{the}%
\TextAt{8}{79.226}{537.179}{series}%
\TextAt{8}{111.371}{537.179}{(1)}%
\TextAt{8}{131.046}{537.179}{converge}%
\TextAt{8}{179.665}{537.179}{so}%
\TextAt{8}{194.846}{537.179}{that}%
\TextAt{8}{221.025}{537.179}{the}%
\TextAt{13}{241.991}{537.179}{p}%
\TextAt{8}{247.866}{537.179}{-adic}%
\TextAt{8}{277.285}{537.179}{numbers}%
\TextAt{8}{325.307}{537.179}{appear}%
\TextAt{8}{364.816}{537.179}{in}%
\TextAt{8}{379.279}{537.179}{the}%
\TextAt{8}{400.245}{537.179}{usual}%
\TextAt{8}{431.685}{537.179}{manner}%
\TextAt{8}{474.768}{537.179}{as}%
\TextAt{8}{489.949}{537.179}{limits}%
\TextAt{8}{523.337}{537.179}{of}%
\TextAt{8}{79.226}{520.133}{Cauchy}%
\TextAt{8}{121.966}{520.133}{sequences}%
\TextAt{8}{175.872}{520.133}{of}%
\TextAt{8}{189.992}{520.133}{rational}%
\TextAt{8}{234.37}{520.133}{numbers.}%
\TextAt{8}{288.5}{520.133}{This}%
\TextAt{8}{316.009}{520.133}{approach}%
\TextAt{8}{367.213}{520.133}{was}%
\TextAt{8}{390.502}{520.133}{proposed}%
\TextAt{8}{441.109}{520.133}{by}%
\TextAt{8}{458.157}{520.133}{the}%
\TextAt{8}{479.103}{520.133}{Hungarian}%
\TextAt{8}{79.226}{503.087}{mathematician}%
\TextAt{26}{159.804}{503.087}{J.}%
\TextAt{26}{175.58}{503.087}{Kurschák}%
\TextAt{8}{234.986}{503.087}{.}%
\TextAt{8}{245.155}{503.087}{The}%
\TextAt{13}{269.783}{503.087}{p}%
\TextAt{8}{275.658}{503.087}{-adic}%
\TextAt{8}{304.84}{503.087}{absolute}%
\TextAt{8}{351.656}{503.087}{value}%
\TextAt{16}{382.477}{503.087}{\Character{106}}%
\TextAt{16}{390.248}{503.087}{\Character{106}}%
\TextAt{21}{393.59}{501.294}{p}%
\TextAt{8}{402.822}{503.087}{is}%
\TextAt{8}{415.159}{503.087}{defined}%
\TextAt{8}{456.046}{503.087}{as}%
\TextAt{8}{470.99}{503.087}{follows.}%
\TextAt{8}{515.697}{503.087}{Let}%
\TextAt{13}{79.226}{486.041}{a}%
\TextAt{9}{89.676}{486.041}{\Character{61}}%
\TextAt{21}{104.303}{490.748}{b}%
\Vector{q 0 G 0 g 0.478 w 32.281 -284.97 m 35.948 -284.97 l S Q}%
\TextAt{21}{104.281}{481.919}{c}%
\TextAt{8}{109.144}{486.041}{,}%
\TextAt{13}{117.015}{486.041}{b\Character{59}}%
\TextAt{13}{127.213}{486.041}{c}%
\NativeMathAt{136.579}{486.041}{11.955}{\in}%
\TextAt{18}{148.853}{486.041}{Z}%
\TextAt{8}{161.298}{486.041}{be}%
\TextAt{8}{177.8}{486.041}{a}%
\TextAt{8}{188.133}{486.041}{nonzero}%
\TextAt{8}{232.286}{486.041}{rational}%
\TextAt{8}{276.452}{486.041}{number.}%
\TextAt{8}{325.331}{486.041}{We}%
\TextAt{8}{346.053}{486.041}{extract}%
\TextAt{8}{386.633}{486.041}{from}%
\TextAt{13}{414.85}{486.041}{b}%
\TextAt{8}{424.302}{486.041}{and}%
\TextAt{8}{447.642}{486.041}{from}%
\TextAt{13}{475.86}{486.041}{c}%
\TextAt{8}{485.372}{486.041}{as}%
\TextAt{8}{500.32}{486.041}{high}%
\TextAt{8}{526.912}{486.041}{a}%
\TextAt{8}{79.226}{468.995}{power}%
\TextAt{8}{113.37}{468.995}{of}%
\TextAt{8}{126.7}{468.995}{the}%
\TextAt{8}{146.857}{468.995}{prime}%
\TextAt{8}{180.02}{468.995}{number}%
\TextAt{13}{222.616}{468.995}{p}%
\TextAt{8}{232.389}{468.995}{as}%
\TextAt{8}{246.759}{468.995}{possible,}%
\TextAt{13}{238.091}{437.335}{a}%
\TextAt{9}{247.557}{437.335}{\Character{61}}%
\TextAt{13}{259.982}{437.335}{p}%
\TextAt{21}{265.857}{442.271}{m}%
\TextAt{13}{275.072}{445.423}{b}%
\TextAt{3}{280.049}{449.761}{\Character{48}}%
\Vector{q 0 G 0 g 0.478 w 203.041 -333.676 m 210.875 -333.676 l S Q}%
\TextAt{13}{275.041}{429.135}{c}%
\TextAt{3}{280.079}{432.588}{\Character{48}}%
\TextAt{13}{284.07}{437.335}{\Character{59}}%
\TextAt{9}{313.225}{437.335}{\Character{40}}%
\TextAt{13}{317.777}{437.335}{b}%
\TextAt{3}{322.754}{442.271}{\Character{48}}%
\TextAt{13}{325.549}{437.335}{c}%
\TextAt{3}{330.587}{442.271}{\Character{48}}%
\TextAt{13}{333.382}{437.335}{\Character{59}}%
\TextAt{13}{338.606}{437.335}{p}%
\TextAt{9}{344.501}{437.335}{\Character{41}}%
\TextAt{9}{352.356}{437.335}{\Character{61}}%
\TextAt{9}{364.777}{437.335}{1}%
\TextAt{13}{370.653}{437.335}{\Character{59}}%
\TextAt{8}{517.802}{437.335}{(2)}%
\TextAt{8}{79.226}{408.804}{and}%
\TextAt{8}{101.989}{408.804}{we}%
\TextAt{8}{119.216}{408.804}{put}%
\TextAt{16}{279.598}{389.696}{\Character{106}}%
\TextAt{13}{282.918}{389.696}{a}%
\TextAt{16}{289.063}{389.696}{\Character{106}}%
\TextAt{21}{292.384}{387.903}{p}%
\TextAt{9}{300.466}{389.696}{\Character{61}}%
\TextAt{9}{318.092}{397.784}{1}%
\Vector{q 0 G 0 g 0.478 w 242.087 -381.315 m 255.951 -381.315 l S Q}%
\TextAt{13}{314.087}{381.495}{p}%
\TextAt{21}{319.962}{384.949}{m}%
\TextAt{13}{329.146}{389.696}{\Character{58}}%
\TextAt{8}{79.226}{363.484}{Thus}%
\TextAt{8}{108.587}{363.484}{the}%
\TextAt{13}{128.455}{363.484}{p}%
\TextAt{8}{134.33}{363.484}{-adic}%
\TextAt{8}{162.651}{363.484}{value}%
\TextAt{8}{192.609}{363.484}{no}%
\TextAt{8}{208.58}{363.484}{longer}%
\TextAt{8}{243.416}{363.484}{measures}%
\TextAt{8}{293.327}{363.484}{the}%
\TextAt{8}{313.196}{363.484}{size}%
\TextAt{8}{335.072}{363.484}{of}%
\TextAt{8}{348.114}{363.484}{a}%
\TextAt{8}{357.581}{363.484}{number}%
\TextAt{13}{399.889}{363.484}{a}%
\NativeMathAt{409.355}{363.484}{11.955}{\in}%
\TextAt{18}{420.646}{363.484}{N}%
\TextAt{8}{429.281}{363.484}{.}%
\TextAt{8}{437.633}{363.484}{Instead}%
\TextAt{8}{478.698}{363.484}{it}%
\TextAt{8}{490.113}{363.484}{becomes}%
\TextAt{8}{79.226}{346.438}{small}%
\TextAt{8}{108.968}{346.438}{if}%
\TextAt{8}{118.805}{346.438}{the}%
\TextAt{8}{138.074}{346.438}{number}%
\TextAt{8}{179.783}{346.438}{is}%
\TextAt{8}{190.66}{346.438}{divisible}%
\TextAt{8}{235.68}{346.438}{by}%
\TextAt{8}{251.052}{346.438}{a}%
\TextAt{8}{259.92}{346.438}{high}%
\TextAt{8}{285.048}{346.438}{power}%
\TextAt{8}{318.304}{346.438}{of}%
\TextAt{13}{330.747}{346.438}{p}%
\TextAt{8}{336.622}{346.438}{.}%
\TextAt{8}{344.775}{346.438}{This}%
\TextAt{8}{370.608}{346.438}{elaborates}%
\TextAt{8}{425.396}{346.438}{on}%
\TextAt{8}{440.768}{346.438}{the}%
\TextAt{8}{460.037}{346.438}{idea}%
\TextAt{8}{483.861}{346.438}{suggested}%
\TextAt{8}{79.226}{329.392}{in}%
\TextAt{8}{94.006}{329.392}{(1.4)}%
\TextAt{8}{123.108}{329.392}{that}%
\TextAt{8}{149.604}{329.392}{an}%
\TextAt{8}{166.99}{329.392}{integer}%
\TextAt{8}{206.816}{329.392}{has}%
\TextAt{8}{228.816}{329.392}{to}%
\TextAt{8}{244.254}{329.392}{be}%
\TextAt{8}{261.305}{329.392}{0}%
\TextAt{8}{272.187}{329.392}{if}%
\TextAt{8}{284.038}{329.392}{it}%
\TextAt{8}{296.869}{329.392}{is}%
\TextAt{8}{309.76}{329.392}{infinitely}%
\TextAt{8}{359.987}{329.392}{divisible}%
\TextAt{8}{407.022}{329.392}{by}%
\TextAt{13}{424.408}{329.392}{p}%
\TextAt{8}{430.283}{329.392}{.}%
\TextAt{8}{442.112}{329.392}{In}%
\TextAt{8}{457.86}{329.392}{particular,}%
\TextAt{8}{516.51}{329.392}{the}%
\TextAt{8}{79.226}{312.346}{summands}%
\TextAt{8}{137.449}{312.346}{of}%
\TextAt{8}{150.995}{312.346}{a}%
\TextAt{13}{160.967}{312.346}{p}%
\TextAt{8}{166.842}{312.346}{-adic}%
\TextAt{8}{195.667}{312.346}{series}%
\TextAt{13}{227.218}{312.346}{a}%
\TextAt{17}{233.363}{310.553}{0}%
\TextAt{9}{240.899}{312.346}{\Character{43}}%
\TextAt{13}{252.808}{312.346}{a}%
\TextAt{17}{258.953}{310.553}{1}%
\TextAt{13}{263.685}{312.346}{p}%
\TextAt{9}{272.364}{312.346}{\Character{43}}%
\TextAt{13}{284.273}{312.346}{a}%
\TextAt{17}{290.418}{310.553}{2}%
\TextAt{13}{295.15}{312.346}{p}%
\TextAt{17}{301.025}{316.684}{2}%
\TextAt{9}{308.562}{312.346}{\Character{43}}%
\TextAt{16}{320.47}{312.346}{\Character{1}}%
\TextAt{16}{325.766}{312.346}{\Character{1}}%
\TextAt{16}{331.074}{312.346}{\Character{1}}%
\TextAt{8}{340.524}{312.346}{form}%
\TextAt{8}{368.381}{312.346}{a}%
\TextAt{8}{378.352}{312.346}{sequence}%
\TextAt{8}{427.071}{312.346}{converging}%
\TextAt{8}{485.509}{312.346}{to}%
\TextAt{8}{500.035}{312.346}{0}%
\TextAt{8}{510.007}{312.346}{with}%
\TextAt{8}{79.226}{295.3}{respect}%
\TextAt{8}{119.445}{295.3}{to}%
\TextAt{16}{133.924}{295.3}{\Character{106}}%
\TextAt{16}{141.3}{295.3}{\Character{106}}%
\TextAt{21}{144.631}{293.507}{p}%
\TextAt{8}{149.392}{295.3}{.}%
\TextAt{8}{158.347}{295.3}{The}%
\TextAt{8}{182.569}{295.3}{exponent}%
\TextAt{13}{233.141}{295.3}{m}%
\TextAt{8}{247.446}{295.3}{in}%
\TextAt{8}{261.268}{295.3}{the}%
\TextAt{8}{281.593}{295.3}{representation}%
\TextAt{8}{358.251}{295.3}{(2)}%
\TextAt{8}{377.285}{295.3}{of}%
\TextAt{8}{390.783}{295.3}{the}%
\TextAt{8}{411.108}{295.3}{number}%
\TextAt{13}{453.873}{295.3}{a}%
\TextAt{8}{464.084}{295.3}{is}%
\TextAt{8}{476.017}{295.3}{denoted}%
\TextAt{8}{520.408}{295.3}{by}%
\TextAt{13}{79.226}{278.254}{v}%
\TextAt{21}{84.885}{276.461}{p}%
\TextAt{9}{89.646}{278.254}{\Character{40}}%
\TextAt{13}{94.199}{278.254}{a}%
\TextAt{9}{100.344}{278.254}{\Character{41}}%
\TextAt{8}{104.896}{278.254}{,}%
\TextAt{8}{112.045}{278.254}{and}%
\TextAt{8}{134.808}{278.254}{one}%
\TextAt{8}{156.267}{278.254}{puts}%
\TextAt{8}{182.342}{278.254}{formally}%
\TextAt{13}{228.524}{278.254}{v}%
\TextAt{21}{234.184}{276.461}{p}%
\TextAt{9}{238.945}{278.254}{\Character{40}0\Character{41}}%
\TextAt{9}{257.189}{278.254}{\Character{61}}%
\NativeMathAt{269.648}{278.254}{11.955}{\infty}%
\TextAt{8}{281.604}{278.254}{.}%
\TextAt{8}{290.052}{278.254}{This}%
\TextAt{8}{316.772}{278.254}{gives}%
\TextAt{8}{345.452}{278.254}{the}%
\TextAt{8}{365.609}{278.254}{function}%
\TextAt{13}{254.033}{249.723}{v}%
\TextAt{21}{259.692}{247.929}{p}%
\TextAt{9}{267.774}{249.723}{\Character{58}}%
\TextAt{18}{274.346}{249.723}{Q}%
\NativeArrowAt{286.965}{249.723}{11.955}{19.26}{\longrightarrow}%
\TextAt{18}{309.547}{249.723}{Z}%
\TextAt{16}{320.174}{249.723}{\Character{91}}%
\TextAt{16}{330.79}{249.723}{\Character{102}}%
\NativeMathAt{336.768}{249.723}{11.955}{\infty}%
\TextAt{16}{348.723}{249.723}{\Character{103}}%
\TextAt{13}{354.711}{249.723}{\Character{59}}%
\TextAt{8}{79.226}{221.191}{which}%
\TextAt{8}{112.713}{221.191}{is}%
\TextAt{8}{124.476}{221.191}{easily}%
\TextAt{8}{156.732}{221.191}{checked}%
\TextAt{8}{199.651}{221.191}{to}%
\TextAt{8}{213.961}{221.191}{satisfy}%
\TextAt{8}{250.508}{221.191}{the}%
\TextAt{8}{270.664}{221.191}{properties}%
\TextAt{8}{91.528}{201.157}{1)}%
\TextAt{13}{107.918}{201.157}{v}%
\TextAt{21}{113.578}{199.363}{p}%
\TextAt{9}{118.339}{201.157}{\Character{40}}%
\TextAt{13}{122.891}{201.157}{a}%
\TextAt{9}{129.036}{201.157}{\Character{41}}%
\TextAt{9}{136.891}{201.157}{\Character{61}}%
\NativeMathAt{149.335}{201.157}{11.955}{\infty}%
\NativeArrowAt{164.602}{201.157}{11.955}{21.938}{\Longleftrightarrow}%
\TextAt{13}{189.849}{201.157}{a}%
\TextAt{9}{199.315}{201.157}{\Character{61}}%
\TextAt{9}{211.725}{201.157}{0}%
\TextAt{8}{217.594}{201.157}{,}%
\TextAt{8}{91.528}{184.111}{2)}%
\TextAt{13}{107.918}{184.111}{v}%
\TextAt{21}{113.578}{182.317}{p}%
\TextAt{9}{118.339}{184.111}{\Character{40}}%
\TextAt{13}{122.891}{184.111}{ab}%
\TextAt{9}{134.013}{184.111}{\Character{41}}%
\TextAt{9}{141.868}{184.111}{\Character{61}}%
\TextAt{13}{154.312}{184.111}{v}%
\TextAt{21}{159.971}{182.317}{p}%
\TextAt{9}{164.732}{184.111}{\Character{40}}%
\TextAt{13}{169.284}{184.111}{a}%
\TextAt{9}{175.429}{184.111}{\Character{41}}%
\TextAt{9}{182.626}{184.111}{\Character{43}}%
\TextAt{13}{194.399}{184.111}{v}%
\TextAt{21}{200.059}{182.317}{p}%
\TextAt{9}{204.82}{184.111}{\Character{40}}%
\TextAt{13}{209.372}{184.111}{b}%
\TextAt{9}{214.349}{184.111}{\Character{41}}%
\TextAt{8}{218.901}{184.111}{,}%
\TextAt{8}{91.528}{167.065}{3)}%
\TextAt{13}{107.918}{167.065}{v}%
\TextAt{21}{113.578}{165.271}{p}%
\TextAt{9}{118.339}{167.065}{\Character{40}}%
\TextAt{13}{122.891}{167.065}{a}%
\TextAt{9}{131.693}{167.065}{\Character{43}}%
\TextAt{13}{143.454}{167.065}{b}%
\TextAt{9}{148.431}{167.065}{\Character{41}}%
\NativeMathAt{156.304}{167.065}{11.955}{\geq}%
\TextAt{8}{168.923}{167.065}{min}%
\TextAt{16}{188.434}{167.065}{\Character{102}}%
\TextAt{13}{194.412}{167.065}{v}%
\TextAt{21}{200.071}{165.271}{p}%
\TextAt{9}{204.832}{167.065}{\Character{40}}%
\TextAt{13}{209.384}{167.065}{a}%
\TextAt{9}{215.529}{167.065}{\Character{41}}%
\TextAt{13}{220.082}{167.065}{\Character{59}}%
\TextAt{13}{225.306}{167.065}{v}%
\TextAt{21}{230.985}{165.271}{p}%
\TextAt{9}{235.746}{167.065}{\Character{40}}%
\TextAt{13}{240.298}{167.065}{b}%
\TextAt{9}{245.275}{167.065}{\Character{41}}%
\TextAt{16}{249.828}{167.065}{\Character{103}}%
\TextAt{8}{255.805}{167.065}{,}%
\TextAt{8}{79.226}{147.03}{where}%
\TextAt{13}{113.049}{147.03}{x}%
\TextAt{9}{122.367}{147.03}{\Character{43}}%
\NativeMathAt{134.138}{147.03}{11.955}{\infty}%
\TextAt{9}{149.437}{147.03}{\Character{61}}%
\NativeMathAt{161.885}{147.03}{11.955}{\infty}%
\TextAt{8}{173.841}{147.03}{,}%
\NativeMathAt{181.007}{147.03}{11.955}{\infty}%
\TextAt{9}{195.628}{147.03}{\Character{43}}%
\NativeMathAt{207.399}{147.03}{11.955}{\infty}%
\TextAt{9}{222.698}{147.03}{\Character{61}}%
\NativeMathAt{235.146}{147.03}{11.955}{\infty}%
\TextAt{8}{251.013}{147.03}{and}%
\NativeMathAt{273.789}{147.03}{11.955}{\infty}%
\TextAt{13}{289.088}{147.03}{\Character{62}}%
\TextAt{13}{301.521}{147.03}{x}%
\TextAt{8}{308.189}{147.03}{,}%
\TextAt{8}{315.355}{147.03}{for}%
\TextAt{8}{333.253}{147.03}{all}%
\TextAt{13}{349.526}{147.03}{x}%
\NativeMathAt{359.522}{147.03}{11.955}{\in}%
\TextAt{18}{370.836}{147.03}{Z}%
\TextAt{8}{378.806}{147.03}{.}%
\TextAt{8}{387.296}{147.03}{The}%
\TextAt{8}{411.363}{147.03}{function}%
\TextAt{13}{457.225}{147.03}{v}%
\TextAt{21}{462.884}{145.237}{p}%
\TextAt{8}{471.556}{147.03}{is}%
\TextAt{8}{483.333}{147.03}{called}%
\TextAt{8}{516.51}{147.03}{the}%
\TextAt{13}{79.226}{129.984}{p}%
\TextAt{12}{85.101}{129.984}{-adic}%
\TextAt{12}{117.799}{129.984}{exponential}%
\TextAt{12}{190.056}{129.984}{valuation}%
\TextAt{8}{247.775}{129.984}{of}%
\TextAt{18}{261.105}{129.984}{Q}%
\TextAt{8}{270.404}{129.984}{.}%
\TextAt{8}{278.852}{129.984}{The}%
\TextAt{13}{302.906}{129.984}{p}%
\TextAt{12}{308.781}{129.984}{-adic}%
\TextAt{12}{341.478}{129.984}{absolute}%
\TextAt{12}{394.595}{129.984}{value}%
\TextAt{8}{428.727}{129.984}{is}%
\TextAt{8}{440.491}{129.984}{given}%
\TextAt{8}{471.061}{129.984}{by}%
\TextAt{16}{206.356}{101.453}{\Character{106}}%
\TextAt{16}{213.565}{101.453}{\Character{106}}%
\TextAt{21}{216.895}{99.659}{p}%
\TextAt{9}{224.977}{101.453}{\Character{58}}%
\TextAt{18}{231.55}{101.453}{Q}%
\NativeArrowAt{244.169}{101.453}{11.955}{19.26}{\longrightarrow}%
\TextAt{18}{266.751}{101.453}{R}%
\TextAt{13}{275.385}{101.453}{\Character{59}}%
\TextAt{13}{304.52}{101.453}{a}%
\NativeArrowAt{314.005}{101.453}{11.955}{19.26}{\longmapsto}%
\TextAt{16}{336.576}{101.453}{\Character{106}}%
\TextAt{13}{339.908}{101.453}{a}%
\TextAt{16}{346.053}{101.453}{\Character{106}}%
\TextAt{21}{349.374}{99.659}{p}%
\TextAt{9}{357.456}{101.453}{\Character{61}}%
\TextAt{13}{369.881}{101.453}{p}%
\TextAt{3}{375.757}{106.389}{\Character{0}}%
\TextAt{21}{382.343}{106.389}{v}%
\TextAt{22}{386.467}{105.393}{p}%
\TextAt{17}{390.805}{106.389}{\Character{40}}%
\TextAt{21}{394.098}{106.389}{a}%
\TextAt{17}{398.596}{106.389}{\Character{41}}%
\TextAt{13}{402.388}{101.453}{\Character{58}}%
\TextAt{8}{79.226}{72.921}{In}%
\TextAt{8}{93.847}{72.921}{view}%
\TextAt{8}{120.83}{72.921}{of}%
\TextAt{8}{134.16}{72.921}{1),}%
\TextAt{8}{151.722}{72.921}{2),}%
\TextAt{8}{169.284}{72.921}{3),}%
\TextAt{8}{186.847}{72.921}{it}%
\TextAt{8}{198.551}{72.921}{satisfies}%
\TextAt{8}{241.661}{72.921}{the}%
\TextAt{8}{261.817}{72.921}{conditions}%
\TextAt{8}{317.815}{72.921}{of}%
\TextAt{8}{331.146}{72.921}{a}%
\TextAt{8}{340.901}{72.921}{norm}%
\TextAt{8}{371.47}{72.921}{on}%
\TextAt{18}{387.729}{72.921}{Q}%
\TextAt{8}{397.028}{72.921}{:}%
\TextAt{8}{300.14}{43.033}{98}%
\Destination{522.36}{717.435}{nav.AMS.12}%
\Destination{522.36}{454.385}{nav.AMS.13}%
\Destination{79.23}{215.765}{nav.Item.27}%
\Destination{79.23}{197.815}{nav.Item.28}%
\Destination{79.23}{180.765}{nav.Item.29}%
\Destination{208.17}{717.435}{nav.equation.2.2.1}%
\Destination{78.23}{809.815}{nav.page.98}%
\Destination{79.23}{773.075}{nav.section.2.2}%
\Destination{79.36}{686.855}{nav.print.107}%
\Destination{79.23}{773.075}{outline.22}%
\Hotlink{93.506}{327.155}{25.078}{11.047}{nav.number.2.1.4}%
}

% Page 111 (body page 99)
\Page{%
\TextAt{2}{79.226}{796.943}{II.}%
\TextAt{2}{94.073}{796.943}{The}%
\TextAt{2}{116.492}{796.943}{Theory}%
\TextAt{2}{154.401}{796.943}{of}%
\TextAt{2}{166.826}{796.943}{Valuations}%
\Vector{q 0 G 0 g 0.299 w 7.226 18.729 m 460.769 18.729 l S Q}%
\TextAt{8}{91.528}{761.124}{1)}%
\TextAt{16}{107.918}{761.124}{\Character{106}}%
\TextAt{13}{111.239}{761.124}{a}%
\TextAt{16}{117.384}{761.124}{\Character{106}}%
\TextAt{21}{120.705}{759.33}{p}%
\TextAt{9}{128.787}{761.124}{\Character{61}}%
\TextAt{9}{141.197}{761.124}{0}%
\NativeArrowAt{150.386}{761.124}{11.955}{21.926}{\Longleftrightarrow}%
\TextAt{13}{175.625}{761.124}{a}%
\TextAt{9}{185.091}{761.124}{\Character{61}}%
\TextAt{9}{197.501}{761.124}{0}%
\TextAt{8}{203.369}{761.124}{,}%
\TextAt{8}{91.528}{744.078}{2)}%
\TextAt{16}{107.918}{744.078}{\Character{106}}%
\TextAt{13}{111.239}{744.078}{ab}%
\TextAt{16}{122.361}{744.078}{\Character{106}}%
\TextAt{21}{125.682}{742.284}{p}%
\TextAt{9}{133.764}{744.078}{\Character{61}}%
\TextAt{16}{146.19}{744.078}{\Character{106}}%
\TextAt{13}{149.51}{744.078}{a}%
\TextAt{16}{155.655}{744.078}{\Character{106}}%
\TextAt{21}{158.976}{742.284}{p}%
\TextAt{16}{163.737}{744.078}{\Character{106}}%
\TextAt{13}{167.058}{744.078}{b}%
\TextAt{16}{172.035}{744.078}{\Character{106}}%
\TextAt{21}{175.356}{742.284}{p}%
\TextAt{8}{180.117}{744.078}{,}%
\TextAt{8}{91.528}{727.032}{3)}%
\TextAt{16}{107.918}{727.032}{\Character{106}}%
\TextAt{13}{111.239}{727.032}{a}%
\TextAt{9}{120.041}{727.032}{\Character{43}}%
\TextAt{13}{131.802}{727.032}{b}%
\TextAt{16}{136.779}{727.032}{\Character{106}}%
\TextAt{21}{140.1}{725.238}{p}%
\NativeMathAt{148.182}{727.032}{11.955}{\leq}%
\TextAt{8}{160.801}{727.032}{max}%
\TextAt{16}{182.596}{727.032}{\Character{102}\Character{106}}%
\TextAt{13}{191.894}{727.032}{a}%
\TextAt{16}{198.039}{727.032}{\Character{106}}%
\TextAt{21}{201.36}{725.238}{p}%
\TextAt{13}{206.121}{727.032}{\Character{59}}%
\TextAt{16}{211.365}{727.032}{\Character{106}}%
\TextAt{13}{214.686}{727.032}{b}%
\TextAt{16}{219.663}{727.032}{\Character{106}}%
\TextAt{21}{222.984}{725.238}{p}%
\TextAt{16}{227.745}{727.032}{\Character{103}}%
\NativeMathAt{237.034}{727.032}{11.955}{\leq}%
\TextAt{16}{249.647}{727.032}{\Character{106}}%
\TextAt{13}{252.984}{727.032}{a}%
\TextAt{16}{259.128}{727.032}{\Character{106}}%
\TextAt{21}{262.449}{725.238}{p}%
\TextAt{9}{269.867}{727.032}{\Character{43}}%
\TextAt{16}{281.628}{727.032}{\Character{106}}%
\TextAt{13}{284.949}{727.032}{b}%
\TextAt{16}{289.926}{727.032}{\Character{106}}%
\TextAt{21}{293.247}{725.238}{p}%
\TextAt{8}{298.008}{727.032}{.}%
\TextAt{8}{79.226}{706.997}{One}%
\TextAt{8}{104.862}{706.997}{can}%
\TextAt{8}{127.246}{706.997}{show}%
\TextAt{8}{157.174}{706.997}{that}%
\TextAt{8}{183.468}{706.997}{the}%
\TextAt{8}{204.549}{706.997}{absolute}%
\TextAt{8}{251.716}{706.997}{values}%
\TextAt{16}{287.502}{706.997}{\Character{106}}%
\TextAt{16}{295.632}{706.997}{\Character{106}}%
\TextAt{21}{298.966}{705.204}{p}%
\TextAt{8}{308.549}{706.997}{and}%
\TextAt{16}{332.236}{706.997}{\Character{106}}%
\TextAt{16}{340.365}{706.997}{\Character{106}}%
\TextAt{8}{348.522}{706.997}{essentially}%
\TextAt{8}{405.505}{706.997}{exhaust}%
\TextAt{8}{449.743}{706.997}{all}%
\TextAt{8}{466.927}{706.997}{norms}%
\TextAt{8}{503.035}{706.997}{on}%
\TextAt{18}{520.219}{706.997}{Q}%
\TextAt{8}{529.518}{706.997}{:}%
\TextAt{8}{79.226}{689.951}{any}%
\TextAt{8}{102.44}{689.951}{further}%
\TextAt{8}{142.882}{689.951}{norm}%
\TextAt{8}{174.548}{689.951}{is}%
\TextAt{8}{187.409}{689.951}{a}%
\TextAt{8}{198.261}{689.951}{power}%
\TextAt{16}{233.502}{689.951}{\Character{106}}%
\TextAt{16}{241.799}{689.951}{\Character{106}}%
\TextAt{21}{245.138}{694.289}{s}%
\TextAt{21}{245.138}{686.996}{p}%
\TextAt{8}{254.894}{689.951}{or}%
\TextAt{16}{270.301}{689.951}{\Character{106}}%
\TextAt{16}{278.598}{689.951}{\Character{106}}%
\TextAt{21}{281.937}{694.289}{s}%
\TextAt{8}{286.351}{689.951}{,}%
\TextAt{8}{294.872}{689.951}{for}%
\TextAt{8}{313.854}{689.951}{some}%
\TextAt{8}{344.277}{689.951}{real}%
\TextAt{8}{368.136}{689.951}{number}%
\TextAt{13}{411.83}{689.951}{s}%
\TextAt{13}{422.53}{689.951}{\Character{62}}%
\TextAt{9}{436.829}{689.951}{0}%
\TextAt{8}{447.676}{689.951}{(see}%
\TextAt{8}{472.241}{689.951}{(3.7)).}%
\TextAt{8}{512.613}{689.951}{The}%
\TextAt{8}{79.226}{672.905}{usual}%
\TextAt{8}{110.136}{672.905}{absolute}%
\TextAt{8}{156.659}{672.905}{value}%
\TextAt{16}{187.186}{672.905}{\Character{106}}%
\TextAt{16}{194.67}{672.905}{\Character{106}}%
\TextAt{8}{202.183}{672.905}{is}%
\TextAt{8}{214.228}{672.905}{denoted}%
\TextAt{8}{258.73}{672.905}{in}%
\TextAt{8}{272.664}{672.905}{this}%
\TextAt{8}{295.767}{672.905}{context}%
\TextAt{8}{337.675}{672.905}{by}%
\TextAt{16}{354.214}{672.905}{\Character{106}}%
\TextAt{16}{361.698}{672.905}{\Character{106}}%
\NativeMathAt{365.034}{671.112}{7.97}{\infty}%
\TextAt{8}{374.001}{672.905}{.}%
\TextAt{8}{383.29}{672.905}{The}%
\TextAt{8}{407.625}{672.905}{good}%
\TextAt{8}{436.526}{672.905}{reason}%
\TextAt{8}{473.293}{672.905}{for}%
\TextAt{8}{491.459}{672.905}{this}%
\TextAt{8}{514.562}{672.905}{will}%
\TextAt{8}{79.226}{655.859}{be}%
\TextAt{8}{95.71}{655.859}{explained}%
\TextAt{8}{148.621}{655.859}{in}%
\TextAt{8}{162.834}{655.859}{due}%
\TextAt{8}{185.499}{655.859}{course.}%
\TextAt{8}{227.558}{655.859}{In}%
\TextAt{8}{242.739}{655.859}{conjunction}%
\TextAt{8}{306.709}{655.859}{with}%
\TextAt{8}{333.929}{655.859}{the}%
\TextAt{8}{354.645}{655.859}{absolute}%
\TextAt{8}{401.447}{655.859}{values}%
\TextAt{16}{436.868}{655.859}{\Character{106}}%
\TextAt{16}{444.627}{655.859}{\Character{106}}%
\TextAt{21}{447.967}{654.066}{p}%
\TextAt{8}{452.728}{655.859}{,}%
\TextAt{8}{460.577}{655.859}{it}%
\TextAt{8}{472.84}{655.859}{satisfies}%
\TextAt{8}{516.51}{655.859}{the}%
\TextAt{8}{79.226}{638.813}{following}%
\TextAt{8}{128.661}{638.813}{important}%
\TextAt{12}{183.953}{638.813}{product}%
\TextAt{12}{234.667}{638.813}{formula}%
\TextAt{8}{279.607}{638.813}{:}%
\TextAt{12}{79.226}{613.241}{Proposition}%
\TextAt{12}{151.937}{613.241}{(2.1).}%
\TextAt{8}{187.253}{613.241}{For}%
\TextAt{8}{208.222}{613.241}{every}%
\TextAt{8}{239.115}{613.241}{rational}%
\TextAt{8}{282.703}{613.241}{number}%
\TextAt{13}{325.299}{613.241}{a}%
\NativeMathAt{334.765}{613.241}{11.955}{\neq}%
\TextAt{9}{347.175}{613.241}{0}%
\TextAt{8}{353.044}{613.241}{,}%
\TextAt{8}{360.193}{613.241}{one}%
\TextAt{8}{381.652}{613.241}{has}%
\TextAt{20}{276.164}{597.138}{\Character{89}}%
\TextAt{21}{281.671}{572.032}{p}%
\TextAt{16}{293.433}{585.781}{\Character{106}}%
\TextAt{13}{296.754}{585.781}{a}%
\TextAt{16}{302.899}{585.781}{\Character{106}}%
\TextAt{21}{306.22}{583.988}{p}%
\TextAt{9}{314.301}{585.781}{\Character{61}}%
\TextAt{9}{326.711}{585.781}{1}%
\TextAt{13}{332.58}{585.781}{\Character{59}}%
\TextAt{8}{79.226}{549.58}{where}%
\TextAt{13}{113.035}{549.58}{p}%
\TextAt{8}{122.808}{549.58}{varies}%
\TextAt{8}{155.72}{549.58}{over}%
\TextAt{8}{180.767}{549.58}{all}%
\TextAt{8}{197.026}{549.58}{prime}%
\TextAt{8}{230.189}{549.58}{numbers}%
\TextAt{8}{277.4}{549.58}{as}%
\TextAt{8}{291.77}{549.58}{well}%
\TextAt{8}{315.501}{549.58}{as}%
\TextAt{8}{329.871}{549.58}{the}%
\TextAt{8}{350.028}{549.58}{symbol}%
\NativeMathAt{390.09}{549.58}{11.955}{\infty}%
\TextAt{8}{402.045}{549.58}{.}%
\TextAt{19}{79.226}{515.482}{Proof.}%
\TextAt{8}{116}{515.482}{In}%
\TextAt{8}{130.621}{515.482}{the}%
\TextAt{8}{150.778}{515.482}{prime}%
\TextAt{8}{183.941}{515.482}{factorization}%
\TextAt{13}{271.334}{488.022}{a}%
\TextAt{9}{280.8}{488.022}{\Character{61}}%
\NativeMathAt{293.225}{488.022}{11.955}{\pm}%
\TextAt{20}{306.537}{499.379}{\Character{89}}%
\TextAt{21}{304.516}{473.919}{p}%
\NativeMathAt{308.779}{473.919}{7.97}{\neq}%
\NativeMathAt{315.365}{473.919}{7.97}{\infty}%
\TextAt{13}{325.826}{488.022}{p}%
\TextAt{21}{331.701}{492.958}{v}%
\TextAt{22}{335.825}{491.962}{p}%
\TextAt{8}{79.226}{450.797}{of}%
\TextAt{13}{92.373}{450.797}{a}%
\TextAt{8}{98.518}{450.797}{,}%
\TextAt{8}{105.52}{450.797}{the}%
\TextAt{8}{125.494}{450.797}{exponent}%
\TextAt{13}{175.713}{450.797}{v}%
\TextAt{21}{181.373}{449.004}{p}%
\TextAt{8}{189.848}{450.797}{of}%
\TextAt{13}{202.994}{450.797}{p}%
\TextAt{8}{212.584}{450.797}{is}%
\TextAt{8}{224.164}{450.797}{precisely}%
\TextAt{8}{271.838}{450.797}{the}%
\TextAt{8}{291.811}{450.797}{exponential}%
\TextAt{8}{354.392}{450.797}{valuation}%
\TextAt{13}{405.281}{450.797}{v}%
\TextAt{21}{410.941}{449.004}{p}%
\TextAt{9}{415.701}{450.797}{\Character{40}}%
\TextAt{13}{420.254}{450.797}{a}%
\TextAt{9}{426.399}{450.797}{\Character{41}}%
\TextAt{8}{434.665}{450.797}{and}%
\TextAt{8}{457.245}{450.797}{the}%
\TextAt{8}{477.218}{450.797}{sign}%
\TextAt{8}{501.16}{450.797}{equals}%
\TextAt{21}{86.676}{438.458}{a}%
\Vector{q 0 G 0 g 0.478 w 8.422 -337.26 m 25.428 -337.26 l S Q}%
\TextAt{3}{80.422}{429.628}{\Character{106}}%
\TextAt{21}{82.774}{429.628}{a}%
\TextAt{3}{87.272}{429.628}{\Character{106}}%
\NativeMathAt{89.624}{428.632}{5.978}{\infty}%
\TextAt{8}{98.624}{433.751}{.}%
\TextAt{8}{107.072}{433.751}{The}%
\TextAt{8}{131.126}{433.751}{equation}%
\TextAt{8}{178.935}{433.751}{therefore}%
\TextAt{8}{228.035}{433.751}{reads}%
\TextAt{13}{259.734}{400.433}{a}%
\TextAt{9}{269.199}{400.433}{\Character{61}}%
\TextAt{13}{290.624}{408.521}{a}%
\Vector{q 0 G 0 g 0.478 w 210.82 -370.578 m 232.574 -370.578 l S Q}%
\TextAt{16}{282.82}{392.233}{\Character{106}}%
\TextAt{13}{286.141}{392.233}{a}%
\TextAt{16}{292.286}{392.233}{\Character{106}}%
\NativeMathAt{295.607}{390.44}{7.97}{\infty}%
\TextAt{20}{309.782}{411.791}{\Character{89}}%
\TextAt{21}{307.762}{386.331}{p}%
\NativeMathAt{312.024}{386.331}{7.97}{\neq}%
\NativeMathAt{318.611}{386.331}{7.97}{\infty}%
\TextAt{9}{336.115}{408.521}{1}%
\Vector{q 0 G 0 g 0.478 w 258.267 -370.578 m 275.815 -370.578 l S Q}%
\TextAt{16}{330.267}{392.233}{\Character{106}}%
\TextAt{13}{333.588}{392.233}{a}%
\TextAt{16}{339.733}{392.233}{\Character{106}}%
\TextAt{21}{343.054}{390.44}{p}%
\TextAt{13}{349.01}{400.433}{\Character{59}}%
\TextAt{8}{79.226}{363.209}{so}%
\TextAt{8}{93.596}{363.209}{that}%
\TextAt{8}{118.965}{363.209}{one}%
\TextAt{8}{140.425}{363.209}{has}%
\TextAt{8}{161.298}{363.209}{indeed}%
\TextAt{20}{198.359}{372.175}{\Character{81}}%
\TextAt{21}{209.65}{359.722}{p}%
\TextAt{16}{216.404}{363.209}{\Character{106}}%
\TextAt{13}{219.725}{363.209}{a}%
\TextAt{16}{225.87}{363.209}{\Character{106}}%
\TextAt{21}{229.191}{361.415}{p}%
\TextAt{9}{237.272}{363.209}{\Character{61}}%
\TextAt{9}{249.681}{363.209}{1}%
\TextAt{8}{255.551}{363.209}{.}%
\Vector{q 0 G 0 g 0.399 w 452.334 -410.791 m 452.334 -402.722 l S Q}%
\Vector{q 0 G 0 g 0.399 w 452.534 -402.921 m 459.707 -402.921 l S Q}%
\Vector{q 0 G 0 g 0.399 w 452.534 -410.592 m 459.707 -410.592 l S Q}%
\Vector{q 0 G 0 g 0.399 w 459.906 -410.791 m 459.906 -402.722 l S Q}%
\TextAt{8}{97.159}{334.648}{The}%
\TextAt{8}{121.454}{334.648}{notation}%
\TextAt{16}{168.535}{334.648}{\Character{106}}%
\TextAt{16}{175.983}{334.648}{\Character{106}}%
\NativeMathAt{179.315}{332.855}{7.97}{\infty}%
\TextAt{8}{192.42}{334.648}{for}%
\TextAt{8}{210.546}{334.648}{the}%
\TextAt{8}{230.944}{334.648}{ordinary}%
\TextAt{8}{278.348}{334.648}{absolute}%
\TextAt{8}{324.832}{334.648}{value}%
\TextAt{8}{355.319}{334.648}{is}%
\TextAt{8}{367.324}{334.648}{motivated}%
\TextAt{8}{422.535}{334.648}{by}%
\TextAt{8}{439.035}{334.648}{the}%
\TextAt{8}{459.433}{334.648}{analogy}%
\TextAt{8}{502.939}{334.648}{of}%
\TextAt{8}{516.51}{334.648}{the}%
\TextAt{8}{79.226}{317.602}{field}%
\TextAt{8}{104.524}{317.602}{of}%
\TextAt{8}{117.794}{317.602}{rational}%
\TextAt{8}{161.323}{317.602}{numbers}%
\TextAt{18}{208.475}{317.602}{Q}%
\TextAt{8}{221.611}{317.602}{with}%
\TextAt{8}{248.212}{317.602}{the}%
\TextAt{8}{268.309}{317.602}{rational}%
\TextAt{8}{311.838}{317.602}{function}%
\TextAt{8}{357.627}{317.602}{field}%
\TextAt{13}{382.924}{317.602}{k}%
\TextAt{9}{389.413}{317.602}{\Character{40}}%
\TextAt{13}{393.966}{317.602}{t}%
\TextAt{9}{398.193}{317.602}{\Character{41}}%
\TextAt{8}{406.583}{317.602}{over}%
\TextAt{8}{431.57}{317.602}{a}%
\TextAt{8}{441.266}{317.602}{finite}%
\TextAt{8}{471.118}{317.602}{field}%
\TextAt{13}{496.416}{317.602}{k}%
\TextAt{8}{502.905}{317.602}{,}%
\TextAt{8}{510.007}{317.602}{with}%
\TextAt{8}{79.226}{300.556}{which}%
\TextAt{8}{113.616}{300.556}{we}%
\TextAt{8}{131.747}{300.556}{started}%
\TextAt{8}{172.39}{300.556}{our}%
\TextAt{8}{194.108}{300.556}{considerations.}%
\TextAt{8}{277.596}{300.556}{Instead}%
\TextAt{8}{319.852}{300.556}{of}%
\TextAt{18}{334.086}{300.556}{Z}%
\TextAt{8}{342.056}{300.556}{,}%
\TextAt{8}{350.335}{300.556}{we}%
\TextAt{8}{368.466}{300.556}{have}%
\TextAt{8}{396.364}{300.556}{inside}%
\TextAt{13}{430.491}{300.556}{k}%
\TextAt{9}{436.981}{300.556}{\Character{40}}%
\TextAt{13}{441.533}{300.556}{t}%
\TextAt{9}{445.76}{300.556}{\Character{41}}%
\TextAt{8}{455.114}{300.556}{the}%
\TextAt{8}{476.174}{300.556}{polynomial}%
\TextAt{8}{79.226}{283.51}{ring}%
\TextAt{13}{103.341}{283.51}{k}%
\TextAt{9}{109.83}{283.51}{\Character{91}}%
\TextAt{13}{113.082}{283.51}{t}%
\TextAt{9}{117.309}{283.51}{\Character{93}}%
\TextAt{8}{120.561}{283.51}{,}%
\TextAt{8}{127.771}{283.51}{the}%
\TextAt{8}{147.977}{283.51}{prime}%
\TextAt{8}{181.19}{283.51}{ideals}%
\TextAt{14}{213.817}{283.51}{p}%
\NativeMathAt{223.204}{283.51}{11.955}{\neq}%
\TextAt{9}{235.697}{283.51}{0}%
\TextAt{8}{245.513}{283.51}{of}%
\TextAt{8}{258.892}{283.51}{which}%
\TextAt{8}{292.427}{283.51}{are}%
\TextAt{8}{311.987}{283.51}{given}%
\TextAt{8}{342.606}{283.51}{by}%
\TextAt{8}{358.914}{283.51}{the}%
\TextAt{8}{379.119}{283.51}{monic}%
\TextAt{8}{413.635}{283.51}{irreducible}%
\TextAt{8}{471.559}{283.51}{polynomials}%
\TextAt{13}{79.226}{266.464}{p}%
\TextAt{9}{85.101}{266.464}{\Character{40}}%
\TextAt{13}{89.654}{266.464}{t}%
\TextAt{9}{93.881}{266.464}{\Character{41}}%
\NativeMathAt{101.754}{266.464}{11.955}{\in}%
\TextAt{13}{113.045}{266.464}{k}%
\TextAt{9}{119.534}{266.464}{\Character{91}}%
\TextAt{13}{122.786}{266.464}{t}%
\TextAt{9}{127.013}{266.464}{\Character{93}}%
\TextAt{8}{130.265}{266.464}{.}%
\TextAt{8}{138.713}{266.464}{For}%
\TextAt{8}{159.683}{266.464}{every}%
\TextAt{8}{190.575}{266.464}{such}%
\TextAt{14}{216.972}{266.464}{p}%
\TextAt{8}{222.955}{266.464}{,}%
\TextAt{8}{230.104}{266.464}{one}%
\TextAt{8}{251.564}{266.464}{defines}%
\TextAt{8}{289.988}{266.464}{an}%
\TextAt{8}{306.247}{266.464}{absolute}%
\TextAt{8}{352.489}{266.464}{value}%
\TextAt{16}{266.359}{239.004}{\Character{106}}%
\TextAt{16}{273.556}{239.004}{\Character{106}}%
\TextAt{27}{276.898}{237.211}{p}%
\TextAt{9}{284.706}{239.004}{\Character{58}}%
\TextAt{13}{291.278}{239.004}{k}%
\TextAt{9}{297.768}{239.004}{\Character{40}}%
\TextAt{13}{302.32}{239.004}{t}%
\TextAt{9}{306.547}{239.004}{\Character{41}}%
\NativeArrowAt{314.42}{239.004}{11.955}{19.26}{\longrightarrow}%
\TextAt{18}{337.002}{239.004}{R}%
\TextAt{8}{79.226}{211.544}{as}%
\TextAt{8}{94.282}{211.544}{follows.}%
\TextAt{8}{139.325}{211.544}{Let}%
\TextAt{13}{160.98}{211.544}{f}%
\TextAt{9}{168.026}{211.544}{\Character{40}}%
\TextAt{13}{172.578}{211.544}{t}%
\TextAt{9}{176.806}{211.544}{\Character{41}}%
\TextAt{9}{185.832}{211.544}{\Character{61}}%
\TextAt{21}{200.918}{217.242}{g}%
\TextAt{17}{205.234}{217.242}{\Character{40}}%
\TextAt{21}{208.527}{217.242}{t}%
\TextAt{17}{211.585}{217.242}{\Character{41}}%
\Vector{q 0 G 0 g 0.478 w 128.636 -559.467 m 143.16 -559.467 l S Q}%
\TextAt{21}{200.636}{207.421}{h}%
\TextAt{17}{205.516}{207.421}{\Character{40}}%
\TextAt{21}{208.809}{207.421}{t}%
\TextAt{17}{211.867}{207.421}{\Character{41}}%
\TextAt{8}{216.356}{211.544}{,}%
\TextAt{13}{224.362}{211.544}{g}%
\TextAt{9}{230.396}{211.544}{\Character{40}}%
\TextAt{13}{234.948}{211.544}{t}%
\TextAt{9}{239.175}{211.544}{\Character{41}}%
\TextAt{13}{243.728}{211.544}{\Character{59}}%
\TextAt{13}{248.952}{211.544}{h}%
\TextAt{9}{255.711}{211.544}{\Character{40}}%
\TextAt{13}{260.263}{211.544}{t}%
\TextAt{9}{264.49}{211.544}{\Character{41}}%
\NativeMathAt{273.531}{211.544}{11.955}{\in}%
\TextAt{13}{285.99}{211.544}{k}%
\TextAt{9}{292.48}{211.544}{\Character{91}}%
\TextAt{13}{295.731}{211.544}{t}%
\TextAt{9}{299.959}{211.544}{\Character{93}}%
\TextAt{8}{307.793}{211.544}{be}%
\TextAt{8}{324.403}{211.544}{a}%
\TextAt{8}{334.844}{211.544}{nonzero}%
\TextAt{8}{379.106}{211.544}{rational}%
\TextAt{8}{423.38}{211.544}{function.}%
\TextAt{8}{475.835}{211.544}{We}%
\TextAt{8}{496.665}{211.544}{extract}%
\TextAt{8}{79.226}{194.498}{from}%
\TextAt{13}{106.866}{194.498}{g}%
\TextAt{9}{112.901}{194.498}{\Character{40}}%
\TextAt{13}{117.453}{194.498}{t}%
\TextAt{9}{121.68}{194.498}{\Character{41}}%
\TextAt{8}{130.13}{194.498}{and}%
\TextAt{13}{152.893}{194.498}{h}%
\TextAt{9}{159.631}{194.498}{\Character{40}}%
\TextAt{13}{164.183}{194.498}{t}%
\TextAt{9}{168.411}{194.498}{\Character{41}}%
\TextAt{8}{176.86}{194.498}{the}%
\TextAt{8}{197.017}{194.498}{highest}%
\TextAt{8}{237.401}{194.498}{possible}%
\TextAt{8}{281.42}{194.498}{power}%
\TextAt{8}{315.564}{194.498}{of}%
\TextAt{8}{328.894}{194.498}{the}%
\TextAt{8}{349.051}{194.498}{irreducible}%
\TextAt{8}{406.926}{194.498}{polynomial}%
\TextAt{13}{467.419}{194.498}{p}%
\TextAt{9}{473.294}{194.498}{\Character{40}}%
\TextAt{13}{477.846}{194.498}{t}%
\TextAt{9}{482.073}{194.498}{\Character{41}}%
\TextAt{8}{486.626}{194.498}{,}%
\TextAt{13}{219.606}{163.045}{f}%
\TextAt{9}{226.652}{163.045}{\Character{40}}%
\TextAt{13}{231.205}{163.045}{t}%
\TextAt{9}{235.432}{163.045}{\Character{41}}%
\TextAt{9}{243.287}{163.045}{\Character{61}}%
\TextAt{13}{255.73}{163.045}{p}%
\TextAt{9}{261.606}{163.045}{\Character{40}}%
\TextAt{13}{266.158}{163.045}{t}%
\TextAt{9}{270.385}{163.045}{\Character{41}}%
\TextAt{21}{274.937}{167.981}{m}%
\TextAt{20}{284.495}{171.132}{\Character{101}}%
\TextAt{13}{284.474}{171.132}{g}%
\TextAt{9}{290.508}{171.132}{\Character{40}}%
\TextAt{13}{295.06}{171.132}{t}%
\TextAt{9}{299.287}{171.132}{\Character{41}}%
\Vector{q 0 G 0 g 0.478 w 212.122 -607.967 m 232.192 -607.967 l S Q}%
\TextAt{20}{283.845}{155.725}{\Character{101}}%
\TextAt{13}{284.122}{152.571}{h}%
\TextAt{9}{290.86}{152.571}{\Character{40}}%
\TextAt{13}{295.412}{152.571}{t}%
\TextAt{9}{299.64}{152.571}{\Character{41}}%
\TextAt{13}{305.387}{163.045}{\Character{59}}%
\TextAt{9}{334.542}{163.045}{\Character{40}}%
\TextAt{20}{339.116}{163.045}{\Character{101}}%
\TextAt{13}{339.094}{163.045}{g}%
\TextAt{20}{344.852}{166.199}{\Character{101}}%
\TextAt{13}{345.129}{163.045}{h\Character{59}}%
\TextAt{13}{357.084}{163.045}{p}%
\TextAt{9}{362.986}{163.045}{\Character{41}}%
\TextAt{9}{370.841}{163.045}{\Character{61}}%
\TextAt{9}{383.262}{163.045}{1}%
\TextAt{13}{389.138}{163.045}{\Character{59}}%
\TextAt{8}{79.226}{129.875}{and}%
\TextAt{8}{101.989}{129.875}{put}%
\TextAt{13}{230.471}{112.829}{v}%
\TextAt{27}{236.13}{111.035}{p}%
\TextAt{9}{240.617}{112.829}{\Character{40}}%
\TextAt{13}{245.169}{112.829}{f}%
\TextAt{9}{252.216}{112.829}{\Character{41}}%
\TextAt{9}{260.071}{112.829}{\Character{61}}%
\TextAt{13}{272.515}{112.829}{m\Character{59}}%
\TextAt{16}{311.908}{112.829}{\Character{106}}%
\TextAt{13}{315.229}{112.829}{f}%
\TextAt{16}{322.276}{112.829}{\Character{106}}%
\TextAt{27}{325.597}{111.035}{p}%
\TextAt{9}{333.404}{112.829}{\Character{61}}%
\TextAt{13}{345.83}{112.829}{q}%
\TextAt{3}{351.449}{119.074}{\Character{0}}%
\TextAt{21}{358.035}{119.074}{v}%
\TextAt{28}{362.159}{118.077}{p}%
\TextAt{17}{366.242}{119.074}{\Character{40}}%
\TextAt{21}{369.535}{119.074}{f}%
\TextAt{17}{374.482}{119.074}{\Character{41}}%
\TextAt{27}{351.02}{110.858}{p}%
\TextAt{13}{378.273}{112.829}{\Character{59}}%
\TextAt{8}{79.226}{89.967}{where}%
\TextAt{13}{112.833}{89.967}{q}%
\TextAt{27}{118.023}{88.174}{p}%
\TextAt{9}{125.831}{89.967}{\Character{61}}%
\TextAt{13}{138.256}{89.967}{q}%
\TextAt{21}{143.875}{94.306}{d}%
\TextAt{28}{148.233}{93.31}{p}%
\TextAt{8}{152.813}{89.967}{,}%
\TextAt{13}{159.801}{89.967}{d}%
\TextAt{27}{165.883}{88.174}{p}%
\TextAt{8}{174.065}{89.967}{being}%
\TextAt{8}{205.4}{89.967}{the}%
\TextAt{8}{225.354}{89.967}{degree}%
\TextAt{8}{261.567}{89.967}{of}%
\TextAt{8}{274.694}{89.967}{the}%
\TextAt{8}{294.648}{89.967}{residue}%
\TextAt{8}{334.172}{89.967}{class}%
\TextAt{8}{361.407}{89.967}{field}%
\TextAt{8}{386.561}{89.967}{of}%
\TextAt{14}{399.689}{89.967}{p}%
\TextAt{8}{409.366}{89.967}{over}%
\TextAt{13}{434.21}{89.967}{k}%
\TextAt{8}{444.394}{89.967}{and}%
\TextAt{13}{466.954}{89.967}{q}%
\TextAt{8}{476.268}{89.967}{a}%
\TextAt{8}{485.821}{89.967}{fixed}%
\TextAt{8}{513.904}{89.967}{real}%
\TextAt{8}{79.226}{72.921}{number}%
\TextAt{13}{121.957}{72.921}{\Character{62}}%
\TextAt{9}{134.613}{72.921}{1}%
\TextAt{8}{140.466}{72.921}{.}%
\TextAt{8}{149.319}{72.921}{Furthermore}%
\TextAt{8}{217.251}{72.921}{we}%
\TextAt{8}{234.614}{72.921}{put}%
\TextAt{13}{256.208}{72.921}{v}%
\TextAt{27}{261.867}{71.128}{p}%
\TextAt{9}{266.354}{72.921}{\Character{40}0\Character{41}}%
\TextAt{9}{284.837}{72.921}{\Character{61}}%
\NativeMathAt{297.518}{72.921}{11.955}{\infty}%
\TextAt{8}{313.505}{72.921}{and}%
\TextAt{16}{336.403}{72.921}{\Character{106}}%
\TextAt{9}{339.724}{72.921}{0}%
\TextAt{16}{345.577}{72.921}{\Character{106}}%
\TextAt{27}{348.898}{71.128}{p}%
\TextAt{9}{356.935}{72.921}{\Character{61}}%
\TextAt{9}{369.584}{72.921}{0}%
\TextAt{8}{375.444}{72.921}{,}%
\TextAt{8}{382.762}{72.921}{and}%
\TextAt{8}{405.659}{72.921}{obtain}%
\TextAt{8}{442.221}{72.921}{for}%
\TextAt{13}{460.241}{72.921}{v}%
\TextAt{27}{465.901}{71.128}{p}%
\TextAt{8}{474.42}{72.921}{and}%
\TextAt{16}{497.317}{72.921}{\Character{106}}%
\TextAt{16}{504.658}{72.921}{\Character{106}}%
\TextAt{27}{507.991}{71.128}{p}%
\TextAt{8}{516.51}{72.921}{the}%
\TextAt{8}{300.14}{43.033}{99}%
\Destination{79.23}{773.075}{nav.Item.30}%
\Destination{79.23}{757.785}{nav.Item.31}%
\Destination{79.23}{740.735}{nav.Item.32}%
\Destination{78.23}{809.815}{nav.page.99}%
\Destination{79.36}{721.955}{nav.print.108}%
\Destination{79.36}{627.255}{nav.number.2.2.1}%
\Hotlink{471.741}{687.755}{25.078}{11.047}{nav.number.2.3.7}%
}

% Page 112 (body page 100)
\Page{%
\TextAt{2}{79.226}{796.943}{II.}%
\TextAt{2}{94.073}{796.943}{The}%
\TextAt{2}{116.492}{796.943}{Theory}%
\TextAt{2}{154.401}{796.943}{of}%
\TextAt{2}{166.826}{796.943}{Valuations}%
\Vector{q 0 G 0 g 0.299 w 7.226 18.729 m 460.769 18.729 l S Q}%
\TextAt{8}{79.226}{761.124}{same}%
\TextAt{8}{108.715}{761.124}{conditions}%
\TextAt{8}{164.877}{761.124}{1),}%
\TextAt{8}{182.643}{761.124}{2),}%
\TextAt{8}{200.409}{761.124}{3)}%
\TextAt{8}{214.883}{761.124}{as}%
\TextAt{8}{229.416}{761.124}{for}%
\TextAt{13}{247.464}{761.124}{v}%
\TextAt{21}{253.123}{759.33}{p}%
\TextAt{8}{261.945}{761.124}{and}%
\TextAt{16}{284.871}{761.124}{\Character{106}}%
\TextAt{16}{292.235}{761.124}{\Character{106}}%
\TextAt{21}{295.573}{759.33}{p}%
\TextAt{8}{304.395}{761.124}{above.}%
\TextAt{8}{342.611}{761.124}{In}%
\TextAt{8}{357.395}{761.124}{the}%
\TextAt{8}{377.715}{761.124}{case}%
\TextAt{14}{402.649}{761.124}{p}%
\TextAt{9}{412.231}{761.124}{\Character{61}}%
\TextAt{9}{424.927}{761.124}{\Character{40}}%
\TextAt{13}{429.487}{761.124}{t}%
\TextAt{16}{436.482}{761.124}{\Character{0}}%
\TextAt{13}{448.549}{761.124}{a}%
\TextAt{9}{454.694}{761.124}{\Character{41}}%
\TextAt{8}{463.307}{761.124}{for}%
\TextAt{13}{481.355}{761.124}{a}%
\NativeMathAt{491.099}{761.124}{11.955}{\in}%
\TextAt{13}{502.668}{761.124}{k}%
\TextAt{8}{509.157}{761.124}{,}%
\TextAt{8}{516.51}{761.124}{the}%
\TextAt{8}{79.226}{744.078}{valuation}%
\TextAt{13}{131.089}{744.078}{v}%
\TextAt{27}{136.748}{742.284}{p}%
\TextAt{9}{141.235}{744.078}{\Character{40}}%
\TextAt{13}{145.787}{744.078}{f}%
\TextAt{9}{152.833}{744.078}{\Character{41}}%
\TextAt{8}{162.073}{744.078}{is}%
\TextAt{8}{174.627}{744.078}{clearly}%
\TextAt{8}{212.813}{744.078}{the}%
\TextAt{8}{233.759}{744.078}{order}%
\TextAt{8}{265.119}{744.078}{of}%
\TextAt{8}{279.239}{744.078}{the}%
\TextAt{8}{300.185}{744.078}{zero,}%
\TextAt{8}{329.136}{744.078}{resp.}%
\TextAt{8}{360.828}{744.078}{pole,}%
\TextAt{8}{390.101}{744.078}{of}%
\TextAt{8}{404.221}{744.078}{the}%
\TextAt{8}{425.168}{744.078}{function}%
\TextAt{13}{471.806}{744.078}{f}%
\TextAt{9}{483.519}{744.078}{\Character{61}}%
\TextAt{13}{497.291}{744.078}{f}%
\TextAt{9}{504.337}{744.078}{\Character{40}}%
\TextAt{13}{508.89}{744.078}{t}%
\TextAt{9}{513.117}{744.078}{\Character{41}}%
\TextAt{8}{522.357}{744.078}{at}%
\TextAt{13}{79.226}{727.032}{t}%
\TextAt{9}{86.774}{727.032}{\Character{61}}%
\TextAt{13}{99.2}{727.032}{a}%
\TextAt{8}{105.345}{727.032}{.}%
\TextAt{8}{113.793}{727.032}{But}%
\TextAt{8}{137.034}{727.032}{for}%
\TextAt{8}{154.919}{727.032}{the}%
\TextAt{8}{175.075}{727.032}{function}%
\TextAt{8}{220.923}{727.032}{field}%
\TextAt{13}{246.28}{727.032}{k}%
\TextAt{9}{252.77}{727.032}{\Character{40}}%
\TextAt{13}{257.322}{727.032}{t}%
\TextAt{9}{261.549}{727.032}{\Character{41}}%
\TextAt{8}{266.101}{727.032}{,}%
\TextAt{8}{273.251}{727.032}{there}%
\TextAt{8}{303.162}{727.032}{is}%
\TextAt{8}{314.926}{727.032}{one}%
\TextAt{8}{336.386}{727.032}{more}%
\TextAt{8}{365.652}{727.032}{exponential}%
\TextAt{8}{428.417}{727.032}{valuation}%
\TextAt{13}{246.668}{699.821}{v}%
\NativeMathAt{252.328}{698.028}{7.97}{\infty}%
\TextAt{9}{264.615}{699.821}{\Character{58}}%
\TextAt{13}{271.187}{699.821}{k}%
\TextAt{9}{277.677}{699.821}{\Character{40}}%
\TextAt{13}{282.229}{699.821}{t}%
\TextAt{9}{286.456}{699.821}{\Character{41}}%
\NativeArrowAt{294.33}{699.821}{11.955}{19.26}{\longrightarrow}%
\TextAt{18}{316.912}{699.821}{Z}%
\TextAt{16}{327.538}{699.821}{\Character{91}}%
\TextAt{16}{338.154}{699.821}{\Character{102}}%
\NativeMathAt{344.132}{699.821}{11.955}{\infty}%
\TextAt{16}{356.087}{699.821}{\Character{103}}%
\TextAt{13}{362.076}{699.821}{\Character{59}}%
\TextAt{8}{79.226}{672.611}{namely}%
\TextAt{13}{240.751}{655.565}{v}%
\NativeMathAt{246.411}{653.771}{7.97}{\infty}%
\TextAt{9}{255.377}{655.565}{\Character{40}}%
\TextAt{13}{259.929}{655.565}{f}%
\TextAt{9}{266.976}{655.565}{\Character{41}}%
\TextAt{9}{274.831}{655.565}{\Character{61}}%
\TextAt{8}{287.274}{655.565}{deg}%
\TextAt{9}{304.837}{655.565}{\Character{40}}%
\TextAt{13}{309.389}{655.565}{h}%
\TextAt{9}{316.127}{655.565}{\Character{41}}%
\TextAt{16}{323.336}{655.565}{\Character{0}}%
\TextAt{8}{335.292}{655.565}{deg}%
\TextAt{9}{352.854}{655.565}{\Character{40}}%
\TextAt{13}{357.406}{655.565}{g}%
\TextAt{9}{363.44}{655.565}{\Character{41}}%
\TextAt{13}{367.993}{655.565}{\Character{59}}%
\TextAt{8}{79.226}{632.81}{where}%
\TextAt{13}{113.804}{632.81}{f}%
\TextAt{9}{125.482}{632.81}{\Character{61}}%
\TextAt{21}{140.696}{638.066}{g}%
\Vector{q 0 G 0 g 0.478 w 68.414 -138.201 m 73.294 -138.201 l S Q}%
\TextAt{21}{140.414}{628.688}{h}%
\NativeMathAt{151.121}{632.81}{11.955}{\neq}%
\TextAt{9}{164.846}{632.81}{0}%
\TextAt{8}{170.71}{632.81}{,}%
\TextAt{13}{178.821}{632.81}{g}%
\TextAt{13}{184.835}{632.81}{\Character{59}}%
\TextAt{13}{190.071}{632.81}{h}%
\NativeMathAt{201.469}{632.81}{11.955}{\in}%
\TextAt{13}{214.071}{632.81}{k}%
\TextAt{9}{220.561}{632.81}{\Character{91}}%
\TextAt{13}{223.812}{632.81}{t}%
\TextAt{9}{228.039}{632.81}{\Character{93}}%
\TextAt{8}{231.291}{632.81}{.}%
\TextAt{8}{242.047}{632.81}{It}%
\TextAt{8}{255.488}{632.81}{describes}%
\TextAt{8}{306.123}{632.81}{the}%
\TextAt{8}{327.048}{632.81}{order}%
\TextAt{8}{358.387}{632.81}{of}%
\TextAt{8}{372.486}{632.81}{zero,}%
\TextAt{8}{401.41}{632.81}{resp.}%
\TextAt{8}{433.04}{632.81}{pole,}%
\TextAt{8}{462.287}{632.81}{of}%
\TextAt{13}{476.386}{632.81}{f}%
\TextAt{9}{483.433}{632.81}{\Character{40}}%
\TextAt{13}{487.985}{632.81}{t}%
\TextAt{9}{492.212}{632.81}{\Character{41}}%
\TextAt{8}{501.431}{632.81}{at}%
\TextAt{8}{516.51}{632.81}{the}%
\TextAt{8}{79.226}{615.764}{point}%
\TextAt{8}{109.848}{615.764}{at}%
\TextAt{8}{124.211}{615.764}{infinity}%
\NativeMathAt{164.588}{615.764}{11.955}{\infty}%
\TextAt{8}{176.543}{615.764}{,}%
\TextAt{8}{183.758}{615.764}{i.e.,}%
\TextAt{8}{205.928}{615.764}{the}%
\TextAt{8}{226.137}{615.764}{order}%
\TextAt{8}{256.759}{615.764}{of}%
\TextAt{8}{270.141}{615.764}{zero,}%
\TextAt{8}{298.17}{615.764}{resp.}%
\TextAt{8}{327.649}{615.764}{pole,}%
\TextAt{8}{356}{615.764}{of}%
\TextAt{8}{369.383}{615.764}{the}%
\TextAt{8}{389.592}{615.764}{function}%
\TextAt{13}{435.492}{615.764}{f}%
\TextAt{9}{442.538}{615.764}{\Character{40}1\Character{47}}%
\TextAt{13}{458.797}{615.764}{t}%
\TextAt{9}{463.024}{615.764}{\Character{41}}%
\TextAt{8}{471.526}{615.764}{at}%
\TextAt{8}{485.889}{615.764}{the}%
\TextAt{8}{506.097}{615.764}{point}%
\TextAt{13}{79.226}{598.718}{t}%
\TextAt{9}{87.241}{598.718}{\Character{61}}%
\TextAt{9}{100.117}{598.718}{0}%
\TextAt{8}{105.987}{598.718}{.}%
\TextAt{8}{115.258}{598.718}{It}%
\TextAt{8}{128.205}{598.718}{is}%
\TextAt{8}{140.243}{598.718}{associated}%
\TextAt{8}{196.252}{598.718}{to}%
\TextAt{8}{210.836}{598.718}{the}%
\TextAt{8}{231.267}{598.718}{prime}%
\TextAt{8}{264.705}{598.718}{ideal}%
\TextAt{14}{292.942}{598.718}{p}%
\TextAt{9}{302.713}{598.718}{\Character{61}}%
\TextAt{9}{315.589}{598.718}{\Character{40}1\Character{47}}%
\TextAt{13}{331.864}{598.718}{t}%
\TextAt{9}{336.091}{598.718}{\Character{41}}%
\TextAt{8}{344.815}{598.718}{of}%
\TextAt{8}{358.419}{598.718}{the}%
\TextAt{8}{378.85}{598.718}{ring}%
\TextAt{13}{403.19}{598.718}{k}%
\TextAt{9}{409.679}{598.718}{\Character{91}1\Character{47}}%
\TextAt{13}{424.637}{598.718}{t}%
\TextAt{9}{428.864}{598.718}{\Character{93}}%
\NativeMathAt{435.904}{598.718}{11.955}{\subseteq}%
\TextAt{13}{448.99}{598.718}{k}%
\TextAt{9}{455.48}{598.718}{\Character{40}}%
\TextAt{13}{460.032}{598.718}{t}%
\TextAt{9}{464.259}{598.718}{\Character{41}}%
\TextAt{8}{472.983}{598.718}{in}%
\TextAt{8}{486.91}{598.718}{the}%
\TextAt{8}{507.341}{598.718}{same}%
\TextAt{8}{79.226}{581.673}{way}%
\TextAt{8}{102.917}{581.673}{as}%
\TextAt{8}{117.235}{581.673}{the}%
\TextAt{8}{137.34}{581.673}{exponential}%
\TextAt{8}{200.053}{581.673}{valuations}%
\TextAt{13}{255.688}{581.673}{v}%
\TextAt{27}{261.347}{579.879}{p}%
\TextAt{8}{269.68}{581.673}{are}%
\TextAt{8}{289.139}{581.673}{associated}%
\TextAt{8}{344.822}{581.673}{to}%
\TextAt{8}{359.08}{581.673}{the}%
\TextAt{8}{379.185}{581.673}{prime}%
\TextAt{8}{412.297}{581.673}{ideals}%
\TextAt{14}{444.823}{581.673}{p}%
\TextAt{8}{454.651}{581.673}{of}%
\TextAt{13}{467.929}{581.673}{k}%
\TextAt{9}{474.419}{581.673}{\Character{91}}%
\TextAt{13}{477.671}{581.673}{t}%
\TextAt{9}{481.898}{581.673}{\Character{93}}%
\TextAt{8}{485.149}{581.673}{.}%
\TextAt{8}{493.58}{581.673}{Putting}%
\TextAt{16}{267.089}{554.462}{\Character{106}}%
\TextAt{13}{270.41}{554.462}{f}%
\TextAt{16}{277.457}{554.462}{\Character{106}}%
\NativeMathAt{280.777}{552.669}{7.97}{\infty}%
\TextAt{9}{293.065}{554.462}{\Character{61}}%
\TextAt{13}{305.49}{554.462}{q}%
\TextAt{3}{311.109}{559.398}{\Character{0}}%
\TextAt{21}{317.696}{559.398}{v}%
\NativeMathAt{321.819}{558.402}{5.978}{\infty}%
\TextAt{17}{329.623}{559.398}{\Character{40}}%
\TextAt{21}{332.917}{559.398}{f}%
\TextAt{17}{337.863}{559.398}{\Character{41}}%
\TextAt{13}{341.655}{554.462}{\Character{59}}%
\TextAt{8}{79.226}{527.251}{the}%
\TextAt{8}{99.383}{527.251}{unique}%
\TextAt{8}{137.424}{527.251}{factorization}%
\TextAt{8}{205.401}{527.251}{in}%
\TextAt{13}{219.054}{527.251}{k}%
\TextAt{9}{225.543}{527.251}{\Character{40}}%
\TextAt{13}{230.096}{527.251}{t}%
\TextAt{9}{234.323}{527.251}{\Character{41}}%
\TextAt{8}{242.773}{527.251}{yields,}%
\TextAt{8}{278.925}{527.251}{as}%
\TextAt{8}{293.295}{527.251}{in}%
\TextAt{8}{306.948}{527.251}{(2.1)}%
\TextAt{8}{334.923}{527.251}{above,}%
\TextAt{8}{371.35}{527.251}{the}%
\TextAt{8}{391.507}{527.251}{formula}%
\TextAt{20}{275.85}{511.398}{\Character{89}}%
\TextAt{27}{281.494}{486.292}{p}%
\TextAt{16}{293.119}{500.041}{\Character{106}}%
\TextAt{13}{296.44}{500.041}{f}%
\TextAt{16}{303.486}{500.041}{\Character{106}}%
\TextAt{27}{306.807}{498.248}{p}%
\TextAt{9}{314.615}{500.041}{\Character{61}}%
\TextAt{9}{327.024}{500.041}{1}%
\TextAt{13}{332.893}{500.041}{\Character{59}}%
\TextAt{8}{79.226}{463.463}{where}%
\TextAt{14}{113.235}{463.463}{p}%
\TextAt{8}{123.315}{463.463}{varies}%
\TextAt{8}{156.428}{463.463}{over}%
\TextAt{8}{181.673}{463.463}{the}%
\TextAt{8}{202.029}{463.463}{prime}%
\TextAt{8}{235.393}{463.463}{ideals}%
\TextAt{8}{268.17}{463.463}{of}%
\TextAt{13}{281.7}{463.463}{k}%
\TextAt{9}{288.19}{463.463}{\Character{91}}%
\TextAt{13}{291.441}{463.463}{t}%
\TextAt{9}{295.668}{463.463}{\Character{93}}%
\TextAt{8}{303.017}{463.463}{as}%
\TextAt{8}{317.587}{463.463}{well}%
\TextAt{8}{341.518}{463.463}{as}%
\TextAt{8}{356.088}{463.463}{the}%
\TextAt{8}{376.444}{463.463}{symbol}%
\NativeMathAt{416.705}{463.463}{11.955}{\infty}%
\TextAt{8}{428.66}{463.463}{,}%
\TextAt{8}{436.059}{463.463}{which}%
\TextAt{8}{469.745}{463.463}{now}%
\TextAt{8}{494.334}{463.463}{denotes}%
\TextAt{8}{79.226}{446.417}{the}%
\TextAt{8}{99.902}{446.417}{point}%
\TextAt{8}{130.991}{446.417}{at}%
\TextAt{8}{145.821}{446.417}{infinity}%
\TextAt{8}{186.665}{446.417}{(see}%
\TextAt{8}{210.652}{446.417}{chap.}%
\TextAt{8}{242.064}{446.417}{I,}%
\TextAt{8}{253.953}{446.417}{§14,}%
\TextAt{8}{279.146}{446.417}{p.}%
\TextAt{8}{293.319}{446.417}{95).}%
\TextAt{8}{319.596}{446.417}{In}%
\TextAt{8}{334.737}{446.417}{view}%
\TextAt{8}{362.239}{446.417}{of}%
\TextAt{8}{376.089}{446.417}{the}%
\TextAt{8}{396.765}{446.417}{product}%
\TextAt{8}{441.184}{446.417}{formula}%
\TextAt{8}{484.634}{446.417}{(2.1),}%
\TextAt{8}{516.51}{446.417}{the}%
\TextAt{8}{79.226}{429.371}{above}%
\TextAt{8}{111.416}{429.371}{consideration}%
\TextAt{8}{182.042}{429.371}{shows}%
\TextAt{8}{214.674}{429.371}{that}%
\TextAt{8}{239.057}{429.371}{the}%
\TextAt{8}{258.227}{429.371}{ordinary}%
\TextAt{8}{304.405}{429.371}{absolute}%
\TextAt{8}{349.662}{429.371}{value}%
\TextAt{16}{378.922}{429.371}{\Character{106}}%
\TextAt{16}{385.139}{429.371}{\Character{106}}%
\TextAt{8}{391.387}{429.371}{of}%
\TextAt{18}{403.731}{429.371}{Q}%
\TextAt{8}{415.942}{429.371}{should}%
\TextAt{8}{452.089}{429.371}{be}%
\TextAt{8}{467.027}{429.371}{thought}%
\TextAt{8}{509.953}{429.371}{of}%
\TextAt{8}{522.297}{429.371}{as}%
\TextAt{8}{79.226}{412.325}{being}%
\TextAt{8}{109.805}{412.325}{associated}%
\TextAt{8}{164.582}{412.325}{to}%
\TextAt{8}{177.934}{412.325}{a}%
\TextAt{8}{186.731}{412.325}{virtual}%
\TextAt{8}{223.826}{412.325}{point}%
\TextAt{8}{253.437}{412.325}{at}%
\TextAt{8}{266.789}{412.325}{infinity.}%
\TextAt{8}{310.365}{412.325}{This}%
\TextAt{8}{336.126}{412.325}{point}%
\TextAt{8}{365.737}{412.325}{of}%
\TextAt{8}{378.109}{412.325}{view}%
\TextAt{8}{404.133}{412.325}{justifies}%
\TextAt{8}{445.89}{412.325}{the}%
\TextAt{8}{465.088}{412.325}{notation}%
\TextAt{16}{510.97}{412.325}{\Character{106}}%
\TextAt{16}{517.211}{412.325}{\Character{106}}%
\NativeMathAt{520.551}{410.532}{7.97}{\infty}%
\TextAt{8}{529.518}{412.325}{,}%
\TextAt{8}{79.226}{395.279}{obeys}%
\TextAt{8}{111.347}{395.279}{our}%
\TextAt{8}{131.703}{395.279}{constant}%
\TextAt{19}{178.469}{395.279}{leitmotiv}%
\TextAt{8}{226.55}{395.279}{to}%
\TextAt{8}{240.403}{395.279}{study}%
\TextAt{8}{272.201}{395.279}{numbers}%
\TextAt{8}{318.954}{395.279}{as}%
\TextAt{8}{332.867}{395.279}{functions}%
\TextAt{8}{382.873}{395.279}{from}%
\TextAt{8}{410.056}{395.279}{a}%
\TextAt{8}{419.354}{395.279}{geometric}%
\TextAt{8}{472.228}{395.279}{perspective,}%
\TextAt{8}{79.226}{378.233}{and}%
\TextAt{8}{102.713}{378.233}{it}%
\TextAt{8}{115.141}{378.233}{will}%
\TextAt{8}{137.97}{378.233}{fulfill}%
\TextAt{8}{168.929}{378.233}{the}%
\TextAt{8}{189.81}{378.233}{expectations}%
\TextAt{8}{258.236}{378.233}{thus}%
\TextAt{8}{284.712}{378.233}{raised}%
\TextAt{8}{319.317}{378.233}{in}%
\TextAt{8}{333.694}{378.233}{an}%
\TextAt{8}{350.677}{378.233}{ever}%
\TextAt{8}{376.112}{378.233}{growing}%
\TextAt{8}{420.748}{378.233}{and}%
\TextAt{8}{444.235}{378.233}{amazing}%
\TextAt{8}{491.142}{378.233}{manner.}%
\TextAt{8}{79.226}{361.187}{The}%
\TextAt{8}{102.701}{361.187}{decisive}%
\TextAt{8}{145.101}{361.187}{difference}%
\TextAt{8}{196.862}{361.187}{between}%
\TextAt{8}{241.474}{361.187}{the}%
\TextAt{8}{261.052}{361.187}{absolute}%
\TextAt{8}{306.716}{361.187}{value}%
\TextAt{16}{336.384}{361.187}{\Character{106}}%
\TextAt{16}{343.007}{361.187}{\Character{106}}%
\NativeMathAt{346.344}{359.394}{7.97}{\infty}%
\TextAt{8}{358.629}{361.187}{of}%
\TextAt{18}{371.381}{361.187}{Q}%
\TextAt{8}{383.998}{361.187}{and}%
\TextAt{8}{406.182}{361.187}{the}%
\TextAt{8}{425.759}{361.187}{absolute}%
\TextAt{8}{471.423}{361.187}{value}%
\TextAt{16}{501.091}{361.187}{\Character{106}}%
\TextAt{16}{507.714}{361.187}{\Character{106}}%
\NativeMathAt{511.052}{359.394}{7.97}{\infty}%
\TextAt{8}{523.337}{361.187}{of}%
\TextAt{13}{79.226}{344.141}{k}%
\TextAt{9}{85.716}{344.141}{\Character{40}}%
\TextAt{13}{90.268}{344.141}{t}%
\TextAt{9}{94.495}{344.141}{\Character{41}}%
\TextAt{8}{102.111}{344.141}{is,}%
\TextAt{8}{116.46}{344.141}{however,}%
\TextAt{8}{163.924}{344.141}{that}%
\TextAt{8}{188.459}{344.141}{the}%
\TextAt{8}{207.781}{344.141}{former}%
\TextAt{8}{244.343}{344.141}{is}%
\TextAt{8}{255.273}{344.141}{not}%
\TextAt{8}{275.253}{344.141}{derived}%
\TextAt{8}{315.39}{344.141}{from}%
\TextAt{8}{342.196}{344.141}{any}%
\TextAt{8}{363.479}{344.141}{exponential}%
\TextAt{8}{425.41}{344.141}{valuation}%
\TextAt{13}{475.649}{344.141}{v}%
\TextAt{27}{481.308}{342.348}{p}%
\TextAt{8}{488.858}{344.141}{attached}%
\TextAt{8}{79.226}{327.095}{to}%
\TextAt{8}{93.988}{327.095}{a}%
\TextAt{8}{104.195}{327.095}{prime}%
\TextAt{8}{137.811}{327.095}{ideal.}%
\TextAt{8}{171.68}{327.095}{Having}%
\TextAt{8}{212.134}{327.095}{introduced}%
\TextAt{8}{271.118}{327.095}{the}%
\TextAt{13}{291.726}{327.095}{p}%
\TextAt{8}{297.601}{327.095}{-adic}%
\TextAt{8}{326.662}{327.095}{absolute}%
\TextAt{8}{373.356}{327.095}{value}%
\TextAt{16}{404.054}{327.095}{\Character{106}}%
\TextAt{16}{411.705}{327.095}{\Character{106}}%
\TextAt{21}{415.045}{325.302}{p}%
\TextAt{8}{424.155}{327.095}{on}%
\TextAt{8}{440.866}{327.095}{the}%
\TextAt{8}{461.474}{327.095}{field}%
\TextAt{18}{487.283}{327.095}{Q}%
\TextAt{8}{496.581}{327.095}{,}%
\TextAt{8}{504.295}{327.095}{let}%
\TextAt{8}{521.651}{327.095}{us}%
\TextAt{8}{79.226}{310.049}{now}%
\TextAt{8}{103.398}{310.049}{give}%
\TextAt{8}{127.247}{310.049}{a}%
\TextAt{8}{136.785}{310.049}{new}%
\TextAt{8}{160.622}{310.049}{definition}%
\TextAt{8}{212.434}{310.049}{of}%
\TextAt{8}{225.548}{310.049}{the}%
\TextAt{8}{245.487}{310.049}{field}%
\TextAt{18}{270.627}{310.049}{Q}%
\TextAt{21}{279.926}{308.256}{p}%
\TextAt{8}{288.367}{310.049}{of}%
\TextAt{13}{301.48}{310.049}{p}%
\TextAt{8}{307.355}{310.049}{-adic}%
\TextAt{8}{335.747}{310.049}{numbers,}%
\TextAt{8}{386.037}{310.049}{imitating}%
\TextAt{8}{436.558}{310.049}{the}%
\TextAt{8}{456.497}{310.049}{construction}%
\TextAt{8}{523.337}{310.049}{of}%
\TextAt{8}{79.226}{293.003}{the}%
\TextAt{8}{99.703}{293.003}{field}%
\TextAt{8}{125.381}{293.003}{of}%
\TextAt{8}{139.032}{293.003}{real}%
\TextAt{8}{162.116}{293.003}{numbers.}%
\TextAt{8}{214.84}{293.003}{We}%
\TextAt{8}{235.306}{293.003}{will}%
\TextAt{8}{257.732}{293.003}{verify}%
\TextAt{8}{290.571}{293.003}{afterwards}%
\TextAt{8}{348.193}{293.003}{that}%
\TextAt{8}{373.882}{293.003}{this}%
\TextAt{8}{397.026}{293.003}{new,}%
\TextAt{8}{424.733}{293.003}{analytic}%
\TextAt{8}{469.61}{293.003}{construction}%
\TextAt{8}{79.226}{275.957}{does}%
\TextAt{8}{105.599}{275.957}{agree}%
\TextAt{8}{136.144}{275.957}{with}%
\TextAt{8}{162.78}{275.957}{Hensel’s}%
\TextAt{8}{208.065}{275.957}{definition,}%
\TextAt{8}{263.327}{275.957}{which}%
\TextAt{8}{296.789}{275.957}{was}%
\TextAt{8}{319.264}{275.957}{motivated}%
\TextAt{8}{374.209}{275.957}{by}%
\TextAt{8}{390.444}{275.957}{function}%
\TextAt{8}{436.268}{275.957}{theory.}%
\TextAt{8}{476.581}{275.957}{A}%
\TextAt{12}{489.229}{275.957}{Cauchy}%
\TextAt{12}{79.226}{258.911}{sequence}%
\TextAt{8}{135.161}{258.911}{with}%
\TextAt{8}{162.129}{258.911}{respect}%
\TextAt{8}{202.487}{258.911}{to}%
\TextAt{16}{217.105}{258.911}{\Character{106}}%
\TextAt{16}{224.613}{258.911}{\Character{106}}%
\TextAt{21}{227.952}{257.118}{p}%
\TextAt{8}{236.919}{258.911}{is}%
\TextAt{8}{248.99}{258.911}{by}%
\TextAt{8}{265.557}{258.911}{definition}%
\TextAt{8}{317.894}{258.911}{a}%
\TextAt{8}{327.958}{258.911}{sequence}%
\TextAt{16}{376.768}{258.911}{\Character{102}}%
\TextAt{13}{382.745}{258.911}{x}%
\TextAt{21}{389.397}{257.118}{n}%
\TextAt{16}{395.034}{258.911}{\Character{103}}%
\TextAt{8}{405.217}{258.911}{of}%
\TextAt{8}{418.855}{258.911}{rational}%
\TextAt{8}{462.751}{258.911}{numbers}%
\TextAt{8}{510.27}{258.911}{such}%
\TextAt{8}{79.226}{241.865}{that}%
\TextAt{8}{104.595}{241.865}{for}%
\TextAt{8}{122.48}{241.865}{every}%
\TextAt{13}{153.372}{241.865}{\Character{34}}%
\TextAt{13}{162.159}{241.865}{\Character{62}}%
\TextAt{9}{174.596}{241.865}{0}%
\TextAt{8}{180.449}{241.865}{,}%
\TextAt{8}{187.599}{241.865}{there}%
\TextAt{8}{217.51}{241.865}{exists}%
\TextAt{8}{249.825}{241.865}{a}%
\TextAt{8}{259.581}{241.865}{positive}%
\TextAt{8}{302.894}{241.865}{integer}%
\TextAt{13}{341.593}{241.865}{n}%
\TextAt{17}{348.581}{240.072}{0}%
\TextAt{8}{357.21}{241.865}{satisfying}%
\TextAt{16}{213.195}{214.655}{\Character{106}}%
\TextAt{13}{216.516}{214.655}{x}%
\TextAt{21}{223.168}{212.861}{n}%
\TextAt{16}{231.461}{214.655}{\Character{0}}%
\TextAt{13}{243.416}{214.655}{x}%
\TextAt{21}{250.069}{212.861}{m}%
\TextAt{16}{258.057}{214.655}{\Character{106}}%
\TextAt{21}{261.378}{212.861}{p}%
\TextAt{13}{269.46}{214.655}{\Character{60}}%
\TextAt{13}{281.87}{214.655}{\Character{34}}%
\TextAt{8}{311.274}{214.655}{for}%
\TextAt{8}{329.159}{214.655}{all}%
\TextAt{13}{345.418}{214.655}{n\Character{59}}%
\TextAt{13}{357.624}{214.655}{m}%
\NativeMathAt{371.209}{214.655}{11.955}{\geq}%
\TextAt{13}{383.829}{214.655}{n}%
\TextAt{17}{390.816}{212.861}{0}%
\TextAt{13}{395.549}{214.655}{\Character{58}}%
\TextAt{12}{97.159}{184.455}{Example:}%
\TextAt{8}{157.122}{184.455}{Every}%
\TextAt{8}{190.776}{184.455}{formal}%
\TextAt{8}{227.526}{184.455}{series}%
\NativeMathAt{246.001}{166.229}{7.97}{\infty}%
\TextAt{20}{241.601}{162.642}{\Character{88}}%
\TextAt{21}{242.483}{137.536}{\Character{23}}%
\TextAt{17}{247.167}{137.536}{\Character{61}0}%
\TextAt{13}{260.862}{151.285}{a}%
\TextAt{21}{267.007}{149.491}{\Character{23}}%
\TextAt{13}{272.189}{151.285}{p}%
\TextAt{21}{278.065}{156.221}{\Character{23}}%
\TextAt{13}{283.247}{151.285}{\Character{59}}%
\TextAt{9}{312.401}{151.285}{0}%
\NativeMathAt{321.575}{151.285}{11.955}{\leq}%
\TextAt{13}{334.194}{151.285}{a}%
\TextAt{21}{340.339}{149.491}{\Character{23}}%
\TextAt{13}{348.842}{151.285}{\Character{60}}%
\TextAt{13}{361.252}{151.285}{p\Character{59}}%
\TextAt{8}{79.226}{116.763}{provides}%
\TextAt{8}{125.469}{116.763}{a}%
\TextAt{8}{135.224}{116.763}{Cauchy}%
\TextAt{8}{177.175}{116.763}{sequence}%
\TextAt{8}{225.677}{116.763}{via}%
\TextAt{8}{244.865}{116.763}{its}%
\TextAt{8}{261.184}{116.763}{partial}%
\TextAt{8}{298.914}{116.763}{sums}%
\TextAt{13}{269.532}{81.888}{x}%
\TextAt{21}{276.184}{80.095}{n}%
\TextAt{9}{285.141}{81.888}{\Character{61}}%
\TextAt{21}{298.222}{96.832}{n}%
\TextAt{3}{303.36}{96.832}{\Character{0}}%
\TextAt{17}{309.946}{96.832}{1}%
\TextAt{20}{297.567}{93.245}{\Character{88}}%
\TextAt{21}{298.449}{68.139}{\Character{23}}%
\TextAt{17}{303.133}{68.139}{\Character{61}0}%
\TextAt{13}{316.828}{81.888}{a}%
\TextAt{21}{322.973}{80.095}{\Character{23}}%
\TextAt{13}{328.155}{81.888}{p}%
\TextAt{21}{334.03}{86.824}{\Character{23}}%
\TextAt{13}{339.212}{81.888}{\Character{59}}%
\TextAt{8}{297.211}{43.033}{100}%
\Destination{78.23}{809.815}{nav.page.100}%
\Destination{79.36}{776.155}{nav.print.109}%
\Hotlink{241.564}{444.255}{5.22}{11.047}{nav.chapter.1}%
\Hotlink{259.132}{444.255}{12.716}{11.047}{nav.section.1.14}%
\Hotlink{292.819}{444.255}{12.716}{11.047}{nav.print.95}%
\Hotlink{306.448}{525.055}{25.078}{11.047}{nav.number.2.2.1}%
\Hotlink{484.134}{444.255}{25.078}{11.047}{nav.number.2.2.1}%
}

% Page 113 (body page 101)
\Page{%
\TextAt{2}{79.226}{796.943}{II.}%
\TextAt{2}{94.073}{796.943}{The}%
\TextAt{2}{116.492}{796.943}{Theory}%
\TextAt{2}{154.401}{796.943}{of}%
\TextAt{2}{166.826}{796.943}{Valuations}%
\Vector{q 0 G 0 g 0.299 w 7.226 18.729 m 460.769 18.729 l S Q}%
\TextAt{8}{79.226}{761.124}{because}%
\TextAt{8}{122.528}{761.124}{for}%
\TextAt{13}{140.413}{761.124}{n}%
\TextAt{13}{150.706}{761.124}{\Character{62}}%
\TextAt{13}{163.128}{761.124}{m}%
\TextAt{8}{177.283}{761.124}{one}%
\TextAt{8}{198.743}{761.124}{has}%
\TextAt{16}{179.911}{726.578}{\Character{106}}%
\TextAt{13}{183.232}{726.578}{x}%
\TextAt{21}{189.884}{724.784}{n}%
\TextAt{16}{198.177}{726.578}{\Character{0}}%
\TextAt{13}{210.133}{726.578}{x}%
\TextAt{21}{216.785}{724.784}{m}%
\TextAt{16}{224.773}{726.578}{\Character{106}}%
\TextAt{21}{228.094}{724.784}{p}%
\TextAt{9}{236.176}{726.578}{\Character{61}}%
\TextAt{20}{248.601}{747.499}{\Character{12}}%
\TextAt{20}{248.601}{740.326}{\Character{12}}%
\TextAt{20}{248.601}{733.153}{\Character{12}}%
\TextAt{20}{248.601}{725.98}{\Character{12}}%
\TextAt{20}{248.601}{718.807}{\Character{12}}%
\TextAt{21}{253.987}{741.522}{n}%
\TextAt{3}{259.126}{741.522}{\Character{0}}%
\TextAt{17}{265.712}{741.522}{1}%
\TextAt{20}{253.333}{737.935}{\Character{88}}%
\TextAt{21}{252.586}{712.829}{\Character{23}}%
\TextAt{17}{257.27}{712.829}{\Character{61}}%
\TextAt{21}{263.857}{712.829}{m}%
\TextAt{13}{273.34}{726.578}{a}%
\TextAt{21}{279.485}{724.784}{\Character{23}}%
\TextAt{13}{284.667}{726.578}{p}%
\TextAt{21}{290.542}{731.514}{\Character{23}}%
\TextAt{20}{295.724}{747.499}{\Character{12}}%
\TextAt{20}{295.724}{740.326}{\Character{12}}%
\TextAt{20}{295.724}{733.153}{\Character{12}}%
\TextAt{20}{295.724}{725.98}{\Character{12}}%
\TextAt{20}{295.724}{718.807}{\Character{12}}%
\TextAt{21}{299.709}{711.135}{p}%
\NativeMathAt{307.791}{726.578}{11.955}{\leq}%
\TextAt{8}{324.756}{726.578}{max}%
\TextAt{21}{320.41}{719.137}{m}%
\NativeMathAt{327.901}{719.137}{7.97}{\leq}%
\TextAt{21}{334.487}{719.137}{\Character{23}}%
\TextAt{21}{339.165}{719.137}{\Character{60}n}%
\TextAt{16}{350.896}{726.578}{\Character{102}\Character{106}}%
\TextAt{13}{360.194}{726.578}{a}%
\TextAt{21}{366.339}{724.784}{\Character{23}}%
\TextAt{13}{371.521}{726.578}{p}%
\TextAt{21}{377.396}{731.514}{\Character{23}}%
\TextAt{16}{382.578}{726.578}{\Character{106}}%
\TextAt{21}{385.899}{724.784}{p}%
\TextAt{16}{390.66}{726.578}{\Character{103}}%
\NativeMathAt{399.949}{726.578}{11.955}{\leq}%
\TextAt{9}{417.779}{734.665}{1}%
\Vector{q 0 G 0 g 0.478 w 341.773 -44.434 m 355.637 -44.434 l S Q}%
\TextAt{13}{413.773}{718.377}{p}%
\TextAt{21}{419.649}{721.831}{m}%
\TextAt{13}{428.833}{726.578}{\Character{58}}%
\TextAt{8}{79.226}{687.668}{A}%
\TextAt{8}{91.918}{687.668}{sequence}%
\TextAt{16}{140.44}{687.668}{\Character{102}}%
\TextAt{13}{146.418}{687.668}{x}%
\TextAt{21}{153.07}{685.875}{n}%
\TextAt{16}{158.706}{687.668}{\Character{103}}%
\TextAt{8}{168.601}{687.668}{in}%
\TextAt{18}{182.274}{687.668}{Q}%
\TextAt{8}{195.49}{687.668}{is}%
\TextAt{8}{207.273}{687.668}{called}%
\TextAt{8}{240.457}{687.668}{a}%
\TextAt{12}{250.232}{687.668}{nullsequence}%
\TextAt{8}{327.937}{687.668}{with}%
\TextAt{8}{354.617}{687.668}{respect}%
\TextAt{8}{394.686}{687.668}{to}%
\TextAt{16}{409.017}{687.668}{\Character{106}}%
\TextAt{16}{416.238}{687.668}{\Character{106}}%
\TextAt{21}{419.576}{685.875}{p}%
\TextAt{8}{428.254}{687.668}{if}%
\TextAt{16}{438.997}{687.668}{\Character{106}}%
\TextAt{13}{442.318}{687.668}{x}%
\TextAt{21}{448.97}{685.875}{n}%
\TextAt{16}{454.607}{687.668}{\Character{106}}%
\TextAt{21}{457.928}{685.875}{p}%
\TextAt{8}{466.606}{687.668}{is}%
\TextAt{8}{478.389}{687.668}{a}%
\TextAt{8}{488.165}{687.668}{sequence}%
\TextAt{8}{79.226}{670.622}{converging}%
\TextAt{8}{137.448}{670.622}{to}%
\TextAt{8}{151.758}{670.622}{0}%
\TextAt{8}{161.514}{670.622}{in}%
\TextAt{8}{175.166}{670.622}{the}%
\TextAt{8}{195.323}{670.622}{usual}%
\TextAt{8}{225.952}{670.622}{sense.}%
\TextAt{12}{97.159}{650.587}{Example:}%
\TextAt{9}{159.271}{650.587}{1}%
\TextAt{13}{165.124}{650.587}{\Character{59}}%
\TextAt{13}{170.348}{650.587}{p\Character{59}}%
\TextAt{13}{181.455}{650.587}{p}%
\TextAt{17}{187.362}{654.925}{2}%
\TextAt{13}{192.095}{650.587}{\Character{59}}%
\TextAt{13}{197.319}{650.587}{p}%
\TextAt{17}{203.214}{654.925}{3}%
\TextAt{13}{207.946}{650.587}{\Character{59}}%
\TextAt{13}{213.17}{650.587}{\Character{58}}%
\TextAt{13}{218.407}{650.587}{\Character{58}}%
\TextAt{13}{223.643}{650.587}{\Character{58}}%
\TextAt{8}{226.93}{650.587}{.}%
\TextAt{8}{238.602}{650.587}{The}%
\TextAt{8}{263.73}{650.587}{Cauchy}%
\TextAt{8}{306.755}{650.587}{sequences}%
\TextAt{8}{360.946}{650.587}{form}%
\TextAt{8}{389.661}{650.587}{a}%
\TextAt{8}{400.49}{650.587}{ring}%
\TextAt{13}{425.631}{650.587}{R}%
\TextAt{8}{434.639}{650.587}{,}%
\TextAt{8}{443.131}{650.587}{the}%
\TextAt{8}{464.362}{650.587}{nullsequences}%
\TextAt{8}{79.226}{633.541}{form}%
\TextAt{8}{106.229}{633.541}{a}%
\TextAt{8}{115.348}{633.541}{maximal}%
\TextAt{8}{162.52}{633.541}{ideal}%
\TextAt{14}{189.846}{633.541}{m}%
\TextAt{8}{199.012}{633.541}{,}%
\TextAt{8}{205.652}{633.541}{and}%
\TextAt{8}{227.778}{633.541}{we}%
\TextAt{8}{244.368}{633.541}{define}%
\TextAt{8}{277.54}{633.541}{afresh}%
\TextAt{8}{311.107}{633.541}{the}%
\TextAt{8}{330.626}{633.541}{field}%
\TextAt{8}{355.346}{633.541}{of}%
\TextAt{13}{368.039}{633.541}{p}%
\TextAt{8}{373.914}{633.541}{-adic}%
\TextAt{8}{401.886}{633.541}{numbers}%
\TextAt{8}{448.46}{633.541}{to}%
\TextAt{8}{462.133}{633.541}{be}%
\TextAt{8}{477.42}{633.541}{the}%
\TextAt{8}{496.94}{633.541}{residue}%
\TextAt{8}{79.226}{616.495}{class}%
\TextAt{8}{106.663}{616.495}{field}%
\TextAt{18}{275.829}{599.449}{Q}%
\TextAt{21}{285.128}{597.656}{p}%
\TextAt{9}{293.209}{599.449}{\Character{58}\Character{61}}%
\TextAt{13}{308.887}{599.449}{R}%
\TextAt{9}{317.895}{599.449}{\Character{47}}%
\TextAt{14}{323.748}{599.449}{m}%
\TextAt{13}{332.915}{599.449}{\Character{58}}%
\TextAt{8}{79.226}{575.927}{We}%
\TextAt{8}{98.689}{575.927}{embed}%
\TextAt{18}{135.069}{575.927}{Q}%
\TextAt{8}{147.584}{575.927}{in}%
\TextAt{18}{160.555}{575.927}{Q}%
\TextAt{21}{169.854}{574.134}{p}%
\TextAt{8}{177.831}{575.927}{by}%
\TextAt{8}{193.408}{575.927}{associating}%
\TextAt{8}{252.372}{575.927}{to}%
\TextAt{8}{266.001}{575.927}{every}%
\TextAt{8}{296.211}{575.927}{element}%
\TextAt{13}{338.772}{575.927}{a}%
\NativeMathAt{348.238}{575.927}{11.955}{\in}%
\TextAt{18}{359.529}{575.927}{Q}%
\TextAt{8}{372.043}{575.927}{the}%
\TextAt{8}{391.519}{575.927}{residue}%
\TextAt{8}{430.564}{575.927}{class}%
\TextAt{8}{457.32}{575.927}{of}%
\TextAt{8}{469.969}{575.927}{the}%
\TextAt{8}{489.444}{575.927}{constant}%
\TextAt{8}{79.226}{558.881}{sequence}%
\TextAt{9}{127.39}{558.881}{\Character{40}}%
\TextAt{13}{131.942}{558.881}{a\Character{59}}%
\TextAt{13}{143.311}{558.881}{a\Character{59}}%
\TextAt{13}{154.693}{558.881}{a\Character{59}}%
\TextAt{13}{166.062}{558.881}{\Character{58}}%
\TextAt{13}{171.299}{558.881}{\Character{58}}%
\TextAt{13}{176.535}{558.881}{\Character{58}}%
\TextAt{9}{179.849}{558.881}{\Character{41}}%
\TextAt{8}{184.401}{558.881}{.}%
\TextAt{8}{192.737}{558.881}{The}%
\TextAt{13}{216.452}{558.881}{p}%
\TextAt{8}{222.327}{558.881}{-adic}%
\TextAt{8}{250.597}{558.881}{absolute}%
\TextAt{8}{296.501}{558.881}{value}%
\TextAt{16}{326.409}{558.881}{\Character{106}}%
\TextAt{16}{333.271}{558.881}{\Character{106}}%
\TextAt{21}{336.609}{557.088}{p}%
\TextAt{8}{344.929}{558.881}{on}%
\TextAt{18}{360.849}{558.881}{Q}%
\TextAt{8}{373.706}{558.881}{is}%
\TextAt{8}{385.131}{558.881}{extended}%
\TextAt{8}{434.538}{558.881}{to}%
\TextAt{18}{448.509}{558.881}{Q}%
\TextAt{21}{457.808}{557.088}{p}%
\TextAt{8}{466.127}{558.881}{by}%
\TextAt{8}{482.048}{558.881}{giving}%
\TextAt{8}{516.51}{558.881}{the}%
\TextAt{8}{79.226}{541.835}{element}%
\TextAt{13}{122.468}{541.835}{x}%
\TextAt{9}{132.441}{541.836}{\Character{61}}%
\TextAt{16}{144.866}{541.835}{\Character{102}}%
\TextAt{13}{150.844}{541.836}{x}%
\TextAt{21}{157.496}{540.042}{n}%
\TextAt{16}{163.132}{541.835}{\Character{103}}%
\TextAt{8}{172.431}{541.835}{mod}%
\TextAt{14}{197.869}{541.835}{m}%
\NativeMathAt{210.356}{541.835}{11.955}{\in}%
\TextAt{13}{221.647}{541.835}{R}%
\TextAt{9}{230.656}{541.835}{\Character{47}}%
\TextAt{14}{236.509}{541.835}{m}%
\TextAt{8}{249.573}{541.835}{the}%
\TextAt{8}{269.729}{541.835}{absolute}%
\TextAt{8}{315.972}{541.835}{value}%
\TextAt{16}{250.343}{512.834}{\Character{106}}%
\TextAt{13}{253.664}{512.834}{x}%
\TextAt{16}{260.316}{512.834}{\Character{106}}%
\TextAt{21}{263.637}{511.041}{p}%
\TextAt{9}{271.719}{512.834}{\Character{58}\Character{61}}%
\TextAt{8}{290.304}{512.834}{lim}%
\TextAt{21}{287.396}{505.661}{n}%
\NativeMathAt{292.534}{505.661}{7.97}{\rightarrow}%
\NativeMathAt{300.998}{505.661}{7.97}{\infty}%
\TextAt{16}{311.463}{512.834}{\Character{106}}%
\TextAt{13}{314.784}{512.834}{x}%
\TextAt{21}{321.436}{511.041}{n}%
\TextAt{16}{327.073}{512.834}{\Character{106}}%
\TextAt{21}{330.394}{511.041}{p}%
\NativeMathAt{338.475}{512.834}{11.955}{\in}%
\TextAt{18}{349.766}{512.834}{R}%
\TextAt{13}{358.401}{512.834}{\Character{58}}%
\TextAt{8}{79.226}{482.548}{This}%
\TextAt{8}{107.277}{482.548}{limit}%
\TextAt{8}{136.572}{482.548}{exists}%
\TextAt{8}{170.218}{482.548}{because}%
\TextAt{16}{214.85}{482.548}{\Character{102}\Character{106}}%
\TextAt{13}{224.149}{482.548}{x}%
\TextAt{21}{230.801}{480.755}{n}%
\TextAt{16}{236.437}{482.548}{\Character{106}}%
\TextAt{21}{239.758}{480.755}{p}%
\TextAt{16}{244.519}{482.548}{\Character{103}}%
\TextAt{8}{255.725}{482.548}{is}%
\TextAt{8}{268.821}{482.548}{a}%
\TextAt{8}{279.907}{482.548}{Cauchy}%
\TextAt{8}{323.189}{482.548}{sequence}%
\TextAt{8}{373.022}{482.548}{in}%
\TextAt{18}{388.006}{482.548}{R}%
\TextAt{8}{396.641}{482.548}{,}%
\TextAt{8}{405.454}{482.548}{and}%
\TextAt{8}{429.548}{482.548}{it}%
\TextAt{8}{442.583}{482.548}{is}%
\TextAt{8}{455.678}{482.548}{independent}%
\TextAt{8}{523.337}{482.548}{of}%
\TextAt{8}{79.226}{465.502}{the}%
\TextAt{8}{99.754}{465.502}{choice}%
\TextAt{8}{134.915}{465.502}{of}%
\TextAt{8}{148.616}{465.502}{the}%
\TextAt{8}{169.144}{465.502}{sequence}%
\TextAt{16}{218.018}{465.502}{\Character{102}}%
\TextAt{13}{223.995}{465.502}{x}%
\TextAt{21}{230.647}{463.709}{n}%
\TextAt{16}{236.284}{465.502}{\Character{103}}%
\TextAt{8}{246.53}{465.502}{within}%
\TextAt{8}{283.317}{465.502}{its}%
\TextAt{8}{300.007}{465.502}{class}%
\TextAt{8}{327.816}{465.502}{mod}%
\TextAt{14}{354.524}{465.502}{m}%
\TextAt{8}{367.96}{465.502}{because}%
\TextAt{8}{411.633}{465.502}{any}%
\TextAt{13}{434.121}{465.502}{p}%
\TextAt{8}{439.997}{465.502}{-adic}%
\TextAt{8}{468.977}{465.502}{nullsequence}%
\TextAt{16}{79.226}{448.456}{\Character{102}}%
\TextAt{13}{85.204}{448.456}{y}%
\TextAt{21}{90.911}{446.663}{n}%
\TextAt{16}{96.548}{448.456}{\Character{103}}%
\NativeMathAt{106.22}{448.456}{11.955}{\in}%
\TextAt{14}{117.899}{448.456}{m}%
\TextAt{8}{131.187}{448.456}{satisfies}%
\TextAt{8}{174.521}{448.456}{of}%
\TextAt{8}{188.074}{448.456}{course}%
\TextAt{8}{224.127}{448.456}{lim}%
\TextAt{21}{240.386}{446.663}{n}%
\NativeMathAt{245.524}{446.663}{7.97}{\rightarrow}%
\NativeMathAt{253.988}{446.663}{7.97}{\infty}%
\TextAt{16}{264.952}{448.456}{\Character{106}}%
\TextAt{13}{268.273}{448.456}{y}%
\TextAt{21}{273.98}{446.663}{n}%
\TextAt{16}{279.617}{448.456}{\Character{106}}%
\TextAt{21}{282.938}{446.663}{p}%
\TextAt{9}{291.4}{448.456}{\Character{61}}%
\TextAt{9}{304.192}{448.456}{0}%
\TextAt{8}{310.06}{448.456}{.}%
\TextAt{8}{319.179}{448.456}{The}%
\TextAt{13}{343.456}{448.456}{p}%
\TextAt{8}{349.331}{448.456}{-adic}%
\TextAt{8}{378.163}{448.456}{exponential}%
\TextAt{8}{441.151}{448.456}{valuation}%
\TextAt{13}{492.447}{448.456}{v}%
\TextAt{21}{498.107}{446.663}{p}%
\TextAt{8}{506.989}{448.456}{on}%
\TextAt{18}{523.471}{448.456}{Q}%
\TextAt{8}{79.226}{431.41}{extends}%
\TextAt{8}{121.882}{431.41}{to}%
\TextAt{8}{136.193}{431.41}{an}%
\TextAt{8}{152.452}{431.41}{exponential}%
\TextAt{8}{215.216}{431.41}{valuation}%
\TextAt{13}{251.652}{402.409}{v}%
\TextAt{21}{257.311}{400.616}{p}%
\TextAt{9}{265.393}{402.409}{\Character{58}}%
\TextAt{18}{271.966}{402.409}{Q}%
\TextAt{21}{281.264}{400.616}{p}%
\NativeArrowAt{289.346}{402.409}{11.955}{19.26}{\longrightarrow}%
\TextAt{18}{311.928}{402.409}{Z}%
\TextAt{16}{322.555}{402.409}{\Character{91}}%
\TextAt{16}{333.171}{402.409}{\Character{102}}%
\NativeMathAt{339.149}{402.409}{11.955}{\infty}%
\TextAt{16}{351.104}{402.409}{\Character{103}}%
\TextAt{13}{357.092}{402.409}{\Character{58}}%
\ErrataAt{79.19}{373.408}{455.81}{Let $x$ be the class of the Cauchy sequence $(x_n)$ and use $v_p(0)=\infty$.}%
\ErrataAt{238}{344.407}{172}{$\displaystyle v_p(x):=\lim_{n\to\infty}v_p(x_n).$}%
\ErrataAt{79.19}{315.405}{455.81}{If $x=0$, then $|x_n|_p\to0$ and the valuation limit is $\infty$. If $x\ne0$, then $|x_n|_p\to|x|_p>0$.}%
\ErrataAt{79.19}{298.359}{455.81}{The set of possible absolute values has no positive accumulation point. Thus}%
\ErrataAt{192}{269.358}{240}{$v_p(x)=v_p(x_n)\quad\text{for sufficiently large $n$, if $x\ne0$}.$}%
\TextAt{8}{98.319}{238.256}{Again}%
\TextAt{8}{132.463}{238.256}{we}%
\TextAt{8}{149.69}{238.256}{find}%
\TextAt{8}{173.098}{238.256}{for}%
\TextAt{8}{190.983}{238.256}{all}%
\TextAt{13}{207.242}{238.256}{x}%
\NativeMathAt{217.215}{238.256}{11.955}{\in}%
\TextAt{18}{228.506}{238.256}{Q}%
\TextAt{21}{237.805}{236.463}{p}%
\TextAt{8}{246.463}{238.256}{that}%
\TextAt{16}{271.084}{209.255}{\Character{106}}%
\TextAt{13}{274.405}{209.255}{x}%
\TextAt{16}{281.057}{209.255}{\Character{106}}%
\TextAt{21}{284.378}{207.462}{p}%
\TextAt{9}{292.459}{209.255}{\Character{61}}%
\TextAt{13}{304.885}{209.255}{p}%
\TextAt{3}{310.76}{214.191}{\Character{0}}%
\TextAt{21}{317.347}{214.191}{v}%
\TextAt{22}{321.47}{213.195}{p}%
\TextAt{17}{325.809}{214.191}{\Character{40}}%
\TextAt{21}{329.102}{214.191}{x}%
\TextAt{17}{333.869}{214.191}{\Character{41}}%
\TextAt{13}{337.66}{209.255}{\Character{58}}%
\TextAt{8}{79.226}{180.254}{As}%
\TextAt{8}{96.513}{180.254}{for}%
\TextAt{8}{114.398}{180.254}{the}%
\TextAt{8}{134.555}{180.254}{field}%
\TextAt{8}{159.912}{180.254}{of}%
\TextAt{8}{173.242}{180.254}{real}%
\TextAt{8}{196.004}{180.254}{numbers}%
\TextAt{8}{243.215}{180.254}{one}%
\TextAt{8}{264.675}{180.254}{proves}%
\TextAt{8}{300.839}{180.254}{the}%
\TextAt{12}{79.226}{154.242}{Proposition}%
\TextAt{12}{154.094}{154.242}{(2.2).}%
\TextAt{8}{191.285}{154.242}{The}%
\TextAt{8}{217.214}{154.242}{field}%
\TextAt{18}{244.446}{154.242}{Q}%
\TextAt{21}{253.745}{152.448}{p}%
\TextAt{8}{264.278}{154.242}{of}%
\TextAt{13}{279.483}{154.242}{p}%
\TextAt{8}{285.358}{154.242}{-adic}%
\TextAt{8}{315.842}{154.242}{numbers}%
\TextAt{8}{364.928}{154.242}{is}%
\TextAt{8}{378.567}{154.242}{complete}%
\TextAt{8}{429.865}{154.242}{with}%
\TextAt{8}{458.4}{154.242}{respect}%
\TextAt{8}{500.325}{154.242}{to}%
\TextAt{8}{516.51}{154.242}{the}%
\TextAt{8}{79.226}{137.196}{absolute}%
\TextAt{8}{125.469}{137.196}{value}%
\TextAt{16}{155.715}{137.196}{\Character{106}}%
\TextAt{16}{162.924}{137.196}{\Character{106}}%
\TextAt{21}{166.254}{135.402}{p}%
\TextAt{8}{171.015}{137.196}{,}%
\TextAt{8}{178.165}{137.196}{i.e.,}%
\TextAt{8}{200.27}{137.196}{every}%
\TextAt{8}{231.162}{137.196}{Cauchy}%
\TextAt{8}{273.113}{137.196}{sequence}%
\TextAt{8}{321.615}{137.196}{in}%
\TextAt{18}{335.267}{137.196}{Q}%
\TextAt{21}{344.566}{135.402}{p}%
\TextAt{8}{353.224}{137.196}{converges}%
\TextAt{8}{405.648}{137.196}{with}%
\TextAt{8}{432.308}{137.196}{respect}%
\TextAt{8}{472.358}{137.196}{to}%
\TextAt{16}{486.668}{137.196}{\Character{106}}%
\TextAt{16}{493.865}{137.196}{\Character{106}}%
\TextAt{21}{497.207}{135.402}{p}%
\TextAt{8}{501.968}{137.196}{.}%
\TextAt{8}{297.211}{43.033}{101}%
\Destination{78.23}{809.815}{nav.page.101}%
\Destination{79.36}{776.155}{nav.print.110}%
\Destination{79.36}{388.455}{nav.print.111}%
\Destination{79.36}{168.255}{nav.number.2.2.2}%
\Hotlink{478.432}{98.455}{21.483}{9.572}{nav.number.2.2.3}%
\Destination{79.226}{390.408}{errata.72}%
}

% Page 114 (body page 102)
\Page{%
\TextAt{2}{79.226}{796.943}{II.}%
\TextAt{2}{94.073}{796.943}{The}%
\TextAt{2}{116.492}{796.943}{Theory}%
\TextAt{2}{154.401}{796.943}{of}%
\TextAt{2}{166.826}{796.943}{Valuations}%
\Vector{q 0 G 0 g 0.299 w 7.226 18.729 m 460.769 18.729 l S Q}%
\TextAt{8}{97.159}{761.124}{As}%
\TextAt{8}{114.987}{761.124}{well}%
\TextAt{8}{139.26}{761.124}{as}%
\TextAt{8}{154.171}{761.124}{the}%
\TextAt{8}{174.869}{761.124}{field}%
\TextAt{18}{200.767}{761.124}{R}%
\TextAt{8}{209.401}{761.124}{,}%
\TextAt{8}{217.227}{761.124}{we}%
\TextAt{8}{234.996}{761.124}{thus}%
\TextAt{8}{261.289}{761.124}{obtain}%
\TextAt{8}{298.257}{761.124}{for}%
\TextAt{8}{316.684}{761.124}{each}%
\TextAt{8}{343.562}{761.124}{prime}%
\TextAt{8}{377.267}{761.124}{number}%
\TextAt{13}{420.405}{761.124}{p}%
\TextAt{8}{430.719}{761.124}{a}%
\TextAt{8}{441.015}{761.124}{new}%
\TextAt{8}{465.61}{761.124}{field}%
\TextAt{18}{491.509}{761.124}{Q}%
\TextAt{21}{500.807}{759.33}{p}%
\TextAt{8}{510.007}{761.124}{with}%
\TextAt{8}{79.226}{744.078}{equal}%
\TextAt{8}{110.118}{744.078}{rights}%
\TextAt{8}{143.031}{744.078}{and}%
\TextAt{8}{165.794}{744.078}{standing,}%
\TextAt{8}{216.591}{744.078}{so}%
\TextAt{8}{230.961}{744.078}{that}%
\TextAt{18}{256.33}{744.078}{Q}%
\TextAt{8}{269.526}{744.078}{has}%
\TextAt{8}{290.4}{744.078}{given}%
\TextAt{8}{320.969}{744.078}{rise}%
\TextAt{8}{342.488}{744.078}{to}%
\TextAt{8}{356.799}{744.078}{the}%
\TextAt{8}{376.955}{744.078}{infinite}%
\TextAt{8}{416.622}{744.078}{family}%
\TextAt{8}{452.392}{744.078}{of}%
\TextAt{8}{465.722}{744.078}{fields}%
\TextAt{18}{222.256}{716.32}{Q}%
\TextAt{17}{231.555}{714.527}{2}%
\TextAt{13}{236.287}{716.32}{\Character{59}}%
\TextAt{18}{241.531}{716.32}{Q}%
\TextAt{17}{250.83}{714.527}{3}%
\TextAt{13}{255.562}{716.32}{\Character{59}}%
\TextAt{18}{260.806}{716.32}{Q}%
\TextAt{17}{270.105}{714.527}{5}%
\TextAt{13}{274.837}{716.32}{\Character{59}}%
\TextAt{18}{280.081}{716.32}{Q}%
\TextAt{17}{289.38}{714.527}{7}%
\TextAt{13}{294.112}{716.32}{\Character{59}}%
\TextAt{18}{299.356}{716.32}{Q}%
\TextAt{17}{308.655}{714.527}{11}%
\TextAt{13}{317.621}{716.32}{\Character{59}}%
\TextAt{13}{322.845}{716.32}{\Character{58}}%
\TextAt{13}{328.082}{716.32}{\Character{58}}%
\TextAt{13}{333.318}{716.32}{\Character{58}}%
\TextAt{13}{338.543}{716.32}{\Character{59}}%
\TextAt{18}{343.842}{716.32}{Q}%
\NativeMathAt{353.141}{714.527}{7.97}{\infty}%
\TextAt{9}{365.428}{716.32}{\Character{61}}%
\TextAt{18}{377.853}{716.32}{R}%
\TextAt{13}{386.488}{716.32}{\Character{58}}%
\TextAt{8}{79.226}{688.563}{An}%
\TextAt{8}{98.401}{688.563}{important}%
\TextAt{8}{153.692}{688.563}{special}%
\TextAt{8}{191.791}{688.563}{property}%
\TextAt{8}{239.598}{688.563}{of}%
\TextAt{8}{252.926}{688.563}{the}%
\TextAt{13}{273.081}{688.563}{p}%
\TextAt{8}{278.956}{688.563}{-adic}%
\TextAt{8}{307.563}{688.563}{absolute}%
\TextAt{8}{353.804}{688.563}{values}%
\TextAt{16}{388.664}{688.563}{\Character{106}}%
\TextAt{16}{395.861}{688.563}{\Character{106}}%
\TextAt{21}{399.202}{686.769}{p}%
\TextAt{8}{407.858}{688.563}{lies}%
\TextAt{8}{428.073}{688.563}{in}%
\TextAt{8}{441.724}{688.563}{the}%
\TextAt{8}{461.879}{688.563}{fact}%
\TextAt{8}{484.962}{688.563}{that}%
\TextAt{8}{510.33}{688.563}{they}%
\TextAt{8}{79.226}{671.517}{do}%
\TextAt{8}{95.485}{671.517}{not}%
\TextAt{8}{116.299}{671.517}{only}%
\TextAt{8}{141.991}{671.517}{satisfy}%
\TextAt{8}{178.538}{671.517}{the}%
\TextAt{8}{198.694}{671.517}{usual}%
\TextAt{8}{229.323}{671.517}{triangle}%
\TextAt{8}{272.254}{671.517}{inequality,}%
\TextAt{8}{328.838}{671.517}{but}%
\TextAt{8}{350.298}{671.517}{also}%
\TextAt{8}{373.778}{671.517}{the}%
\TextAt{8}{393.934}{671.517}{stronger}%
\TextAt{8}{439.531}{671.517}{version}%
\TextAt{16}{239.803}{643.759}{\Character{106}}%
\TextAt{13}{243.124}{643.759}{x}%
\TextAt{9}{252.433}{643.759}{\Character{43}}%
\TextAt{13}{264.194}{643.759}{y}%
\TextAt{16}{270.331}{643.759}{\Character{106}}%
\TextAt{21}{273.652}{641.966}{p}%
\NativeMathAt{281.734}{643.759}{11.955}{\leq}%
\TextAt{8}{294.353}{643.759}{max}%
\TextAt{16}{316.147}{643.759}{\Character{102}\Character{106}}%
\TextAt{13}{325.446}{643.759}{x}%
\TextAt{16}{332.098}{643.759}{\Character{106}}%
\TextAt{21}{335.419}{641.966}{p}%
\TextAt{13}{340.18}{643.759}{\Character{59}}%
\TextAt{16}{345.424}{643.759}{\Character{106}}%
\TextAt{13}{348.745}{643.759}{y}%
\TextAt{16}{354.881}{643.759}{\Character{106}}%
\TextAt{21}{358.202}{641.966}{p}%
\TextAt{16}{362.963}{643.759}{\Character{103}}%
\TextAt{13}{368.941}{643.759}{\Character{58}}%
\TextAt{8}{79.226}{616.001}{This}%
\TextAt{8}{106.408}{616.001}{yields}%
\TextAt{8}{139.771}{616.001}{the}%
\TextAt{8}{160.389}{616.001}{following}%
\TextAt{8}{210.286}{616.001}{remarkable}%
\TextAt{8}{271.563}{616.001}{proposition,}%
\TextAt{8}{336.926}{616.001}{which}%
\TextAt{8}{370.874}{616.001}{gives}%
\TextAt{8}{400.017}{616.001}{us}%
\TextAt{8}{415.494}{616.001}{a}%
\TextAt{8}{425.712}{616.001}{new}%
\TextAt{8}{450.228}{616.001}{definition}%
\TextAt{8}{502.718}{616.001}{of}%
\TextAt{8}{516.51}{616.001}{the}%
\TextAt{13}{79.226}{598.955}{p}%
\TextAt{8}{85.101}{598.955}{-adic}%
\TextAt{8}{113.71}{598.955}{integers.}%
\TextAt{12}{79.226}{573.298}{Proposition}%
\TextAt{12}{151.937}{573.298}{(2.3).}%
\TextAt{8}{187.253}{573.298}{The}%
\TextAt{8}{211.307}{573.298}{set}%
\TextAt{18}{241.589}{545.541}{Z}%
\TextAt{21}{249.559}{543.747}{p}%
\TextAt{9}{257.641}{545.541}{\Character{58}\Character{61}}%
\TextAt{16}{273.318}{545.541}{\Character{102}}%
\TextAt{13}{279.295}{545.541}{x}%
\NativeMathAt{289.268}{545.541}{11.955}{\in}%
\TextAt{18}{300.559}{545.541}{Q}%
\TextAt{21}{309.858}{543.747}{p}%
\TextAt{16}{317.94}{545.541}{\Character{106}}%
\TextAt{16}{324.563}{545.541}{\Character{106}}%
\TextAt{13}{327.902}{545.541}{x}%
\TextAt{16}{334.554}{545.541}{\Character{106}}%
\TextAt{21}{337.875}{543.747}{p}%
\NativeMathAt{345.957}{545.541}{11.955}{\leq}%
\TextAt{9}{358.576}{545.541}{1}%
\TextAt{16}{364.429}{545.541}{\Character{103}}%
\TextAt{8}{79.226}{517.783}{is}%
\TextAt{8}{90.99}{517.783}{a}%
\TextAt{8}{100.745}{517.783}{subring}%
\TextAt{8}{142.433}{517.783}{of}%
\TextAt{18}{155.763}{517.783}{Q}%
\TextAt{21}{165.062}{515.99}{p}%
\TextAt{8}{169.823}{517.783}{.}%
\TextAt{8}{178.271}{517.783}{It}%
\TextAt{8}{190.943}{517.783}{is}%
\TextAt{8}{202.707}{517.783}{the}%
\TextAt{8}{222.864}{517.783}{closure}%
\TextAt{8}{261.945}{517.783}{with}%
\TextAt{8}{288.605}{517.783}{respect}%
\TextAt{8}{328.655}{517.783}{to}%
\TextAt{16}{342.965}{517.783}{\Character{106}}%
\TextAt{16}{350.174}{517.783}{\Character{106}}%
\TextAt{21}{353.504}{515.99}{p}%
\TextAt{8}{362.163}{517.783}{of}%
\TextAt{8}{375.493}{517.783}{the}%
\TextAt{8}{395.649}{517.783}{ring}%
\TextAt{18}{419.715}{517.783}{Z}%
\TextAt{8}{431.583}{517.783}{in}%
\TextAt{8}{445.235}{517.783}{the}%
\TextAt{8}{465.392}{517.783}{field}%
\TextAt{18}{490.749}{517.783}{Q}%
\TextAt{21}{500.047}{515.99}{p}%
\TextAt{8}{504.808}{517.783}{.}%
\TextAt{19}{79.226}{483.515}{Proof.}%
\TextAt{8}{116}{483.515}{That}%
\TextAt{18}{145.267}{483.515}{Z}%
\TextAt{21}{153.237}{481.722}{p}%
\TextAt{8}{161.895}{483.515}{is}%
\TextAt{8}{173.659}{483.515}{closed}%
\TextAt{8}{208.185}{483.515}{under}%
\TextAt{8}{241.349}{483.515}{addition}%
\TextAt{8}{287.532}{483.515}{and}%
\TextAt{8}{310.294}{483.515}{multiplication}%
\TextAt{8}{385.421}{483.515}{follows}%
\TextAt{8}{423.857}{483.515}{from}%
\TextAt{16}{168.694}{455.757}{\Character{106}}%
\TextAt{13}{172.015}{455.757}{x}%
\TextAt{9}{181.324}{455.757}{\Character{43}}%
\TextAt{13}{193.085}{455.757}{y}%
\TextAt{16}{199.222}{455.757}{\Character{106}}%
\TextAt{21}{202.543}{453.964}{p}%
\NativeMathAt{210.625}{455.757}{11.955}{\leq}%
\TextAt{8}{223.244}{455.757}{max}%
\TextAt{16}{245.038}{455.757}{\Character{102}\Character{106}}%
\TextAt{13}{254.337}{455.757}{x}%
\TextAt{16}{260.989}{455.757}{\Character{106}}%
\TextAt{21}{264.31}{453.964}{p}%
\TextAt{13}{269.071}{455.757}{\Character{59}}%
\TextAt{16}{274.315}{455.757}{\Character{106}}%
\TextAt{13}{277.636}{455.757}{y}%
\TextAt{16}{283.772}{455.757}{\Character{106}}%
\TextAt{21}{287.093}{453.964}{p}%
\TextAt{16}{291.854}{455.757}{\Character{103}}%
\TextAt{8}{321.742}{455.757}{and}%
\TextAt{16}{364.518}{455.757}{\Character{106}}%
\TextAt{13}{367.839}{455.757}{xy}%
\TextAt{16}{380.627}{455.757}{\Character{106}}%
\TextAt{21}{383.948}{453.964}{p}%
\TextAt{9}{392.03}{455.757}{\Character{61}}%
\TextAt{16}{404.456}{455.757}{\Character{106}}%
\TextAt{13}{407.776}{455.757}{x}%
\TextAt{16}{414.428}{455.757}{\Character{106}}%
\TextAt{21}{417.749}{453.964}{p}%
\TextAt{16}{422.51}{455.757}{\Character{106}}%
\TextAt{13}{425.831}{455.757}{y}%
\TextAt{16}{431.968}{455.757}{\Character{106}}%
\TextAt{21}{435.289}{453.964}{p}%
\TextAt{13}{440.05}{455.757}{\Character{58}}%
\TextAt{8}{79.226}{428}{If}%
\TextAt{16}{92.685}{428}{\Character{102}}%
\TextAt{13}{98.662}{428}{x}%
\TextAt{21}{105.314}{426.206}{n}%
\TextAt{16}{110.951}{428}{\Character{103}}%
\TextAt{8}{122.592}{428}{is}%
\TextAt{8}{136.122}{428}{a}%
\TextAt{8}{147.644}{428}{Cauchy}%
\TextAt{8}{191.361}{428}{sequence}%
\TextAt{8}{241.629}{428}{in}%
\TextAt{18}{257.048}{428}{Z}%
\TextAt{8}{270.682}{428}{and}%
\TextAt{13}{295.211}{428}{x}%
\TextAt{9}{308.194}{428}{\Character{61}}%
\TextAt{8}{323.629}{428}{lim}%
\TextAt{21}{339.888}{426.206}{n}%
\NativeMathAt{345.027}{426.206}{7.97}{\rightarrow}%
\NativeMathAt{353.491}{426.206}{7.97}{\infty}%
\TextAt{13}{364.454}{428}{x}%
\TextAt{21}{371.106}{426.206}{n}%
\TextAt{8}{376.742}{428}{,}%
\TextAt{8}{386.099}{428}{then}%
\TextAt{16}{414.526}{428}{\Character{106}}%
\TextAt{13}{417.847}{428}{x}%
\TextAt{21}{424.499}{426.206}{n}%
\TextAt{16}{430.135}{428}{\Character{106}}%
\TextAt{21}{433.456}{426.206}{p}%
\NativeMathAt{444.548}{428}{11.955}{\leq}%
\TextAt{9}{460.177}{428}{1}%
\TextAt{8}{471.694}{428}{implies}%
\TextAt{8}{513.187}{428}{also}%
\TextAt{16}{79.226}{410.954}{\Character{106}}%
\TextAt{13}{82.547}{410.954}{x}%
\TextAt{16}{89.199}{410.954}{\Character{106}}%
\TextAt{21}{92.52}{409.161}{p}%
\NativeMathAt{102.547}{410.954}{11.955}{\leq}%
\TextAt{9}{117.111}{410.954}{1}%
\TextAt{8}{122.964}{410.954}{,}%
\TextAt{8}{131.54}{410.954}{hence}%
\TextAt{13}{165.188}{410.954}{x}%
\NativeMathAt{177.106}{410.954}{11.955}{\in}%
\TextAt{18}{190.342}{410.954}{Z}%
\TextAt{21}{198.312}{409.161}{p}%
\TextAt{8}{203.073}{410.954}{.}%
\TextAt{8}{214.945}{410.954}{Conversely,}%
\TextAt{8}{277.893}{410.954}{let}%
\TextAt{13}{295.939}{410.954}{x}%
\TextAt{9}{307.857}{410.954}{\Character{61}}%
\TextAt{8}{322.228}{410.954}{lim}%
\TextAt{21}{338.487}{409.161}{n}%
\NativeMathAt{343.625}{409.161}{7.97}{\rightarrow}%
\NativeMathAt{352.089}{409.161}{7.97}{\infty}%
\TextAt{13}{363.052}{410.954}{x}%
\TextAt{21}{369.704}{409.161}{n}%
\NativeMathAt{380.607}{410.954}{11.955}{\in}%
\TextAt{18}{393.843}{410.954}{Z}%
\TextAt{21}{401.813}{409.161}{p}%
\TextAt{8}{406.574}{410.954}{,}%
\TextAt{8}{415.15}{410.954}{for}%
\TextAt{8}{434.176}{410.954}{a}%
\TextAt{8}{445.073}{410.954}{Cauchy}%
\TextAt{8}{488.165}{410.954}{sequence}%
\ErrataAt{79.19}{393.908}{455.81}{$(x_n)$ in $\mathbb Q$. For $n\ge n_0$, we have $|x_n-x|_p\le1$, so $|x_n|_p\le1$ and $x_n=a_n/b_n$, with}%
\TextAt{13}{79.226}{376.862}{a}%
\TextAt{21}{85.371}{375.069}{n}%
\TextAt{13}{91.007}{376.862}{\Character{59}}%
\TextAt{13}{96.231}{376.862}{b}%
\TextAt{21}{101.229}{375.069}{n}%
\NativeMathAt{111.419}{376.862}{11.955}{\in}%
\TextAt{18}{123.944}{376.862}{Z}%
\TextAt{8}{131.914}{376.862}{,}%
\TextAt{9}{139.968}{376.862}{\Character{40}}%
\TextAt{13}{144.52}{376.862}{b}%
\TextAt{21}{149.497}{375.069}{n}%
\TextAt{13}{155.134}{376.862}{\Character{59}}%
\TextAt{13}{160.358}{376.862}{p}%
\TextAt{9}{166.253}{376.862}{\Character{41}}%
\TextAt{9}{175.339}{376.862}{\Character{61}}%
\TextAt{9}{188.992}{376.862}{1}%
\TextAt{8}{194.871}{376.862}{.}%
\TextAt{8}{205.491}{376.862}{Choosing}%
\TextAt{8}{257.335}{376.862}{for}%
\TextAt{8}{275.944}{376.862}{each}%
\TextAt{13}{303.005}{376.862}{n}%
\NativeMathAt{314.547}{376.862}{11.955}{\geq}%
\TextAt{13}{328.4}{376.862}{n}%
\TextAt{17}{335.387}{375.069}{0}%
\TextAt{8}{344.741}{376.862}{a}%
\TextAt{8}{355.22}{376.862}{solution}%
\TextAt{13}{400.238}{376.862}{y}%
\TextAt{21}{405.945}{375.069}{n}%
\NativeMathAt{416.136}{376.862}{11.955}{\in}%
\TextAt{18}{428.66}{376.862}{Z}%
\TextAt{8}{441.252}{376.862}{of}%
\TextAt{8}{455.306}{376.862}{the}%
\TextAt{8}{476.186}{376.862}{congruence}%
\TextAt{13}{79.226}{359.816}{b}%
\TextAt{21}{84.203}{358.023}{n}%
\TextAt{13}{89.84}{359.816}{y}%
\TextAt{21}{95.547}{358.023}{n}%
\NativeMathAt{105.217}{359.816}{11.955}{\equiv}%
\TextAt{13}{118.548}{359.816}{a}%
\TextAt{21}{124.693}{358.023}{n}%
\TextAt{8}{133.65}{359.816}{mod}%
\TextAt{13}{159.088}{359.816}{p}%
\TextAt{21}{164.963}{364.154}{n}%
\TextAt{8}{174.915}{359.816}{yields}%
\TextAt{16}{208.233}{359.816}{\Character{106}}%
\TextAt{13}{211.554}{359.816}{x}%
\TextAt{21}{218.206}{358.023}{n}%
\TextAt{16}{226.784}{359.816}{\Character{0}}%
\TextAt{13}{239.024}{359.816}{y}%
\TextAt{21}{244.732}{358.023}{n}%
\TextAt{16}{250.368}{359.816}{\Character{106}}%
\TextAt{21}{253.689}{358.023}{p}%
\NativeMathAt{262.483}{359.816}{11.955}{\leq}%
\TextAt{9}{275.814}{359.816}{1\Character{47}}%
\TextAt{13}{287.52}{359.816}{p}%
\TextAt{21}{293.395}{364.154}{n}%
\TextAt{8}{303.347}{359.816}{and}%
\TextAt{8}{326.528}{359.816}{hence}%
\TextAt{13}{359.452}{359.816}{x}%
\TextAt{9}{370.137}{359.816}{\Character{61}}%
\TextAt{8}{383.274}{359.816}{lim}%
\TextAt{21}{399.533}{358.023}{n}%
\NativeMathAt{404.671}{358.023}{7.97}{\rightarrow}%
\NativeMathAt{413.135}{358.023}{7.97}{\infty}%
\TextAt{13}{424.099}{359.816}{y}%
\TextAt{21}{429.807}{358.023}{n}%
\TextAt{8}{435.443}{359.816}{,}%
\TextAt{8}{443.114}{359.816}{so}%
\TextAt{8}{457.902}{359.816}{that}%
\TextAt{13}{483.689}{359.816}{x}%
\TextAt{8}{494.657}{359.816}{belongs}%
\TextAt{8}{79.226}{342.77}{to}%
\TextAt{8}{93.536}{342.77}{the}%
\TextAt{8}{113.693}{342.77}{closure}%
\TextAt{8}{152.774}{342.77}{of}%
\TextAt{18}{166.104}{342.77}{Z}%
\TextAt{8}{174.075}{342.77}{.}%
\Vector{q 0 G 0 g 0.399 w 452.334 -431.23 m 452.334 -423.16 l S Q}%
\Vector{q 0 G 0 g 0.399 w 452.534 -423.36 m 459.707 -423.36 l S Q}%
\Vector{q 0 G 0 g 0.399 w 452.534 -431.031 m 459.707 -431.031 l S Q}%
\Vector{q 0 G 0 g 0.399 w 459.906 -431.23 m 459.906 -423.16 l S Q}%
\TextAt{8}{97.159}{314.124}{The}%
\TextAt{8}{121.213}{314.124}{group}%
\TextAt{8}{154.388}{314.124}{of}%
\TextAt{8}{167.718}{314.124}{units}%
\TextAt{8}{197.044}{314.124}{of}%
\TextAt{18}{210.374}{314.124}{Z}%
\TextAt{21}{218.344}{312.331}{p}%
\TextAt{8}{227.003}{314.124}{is}%
\TextAt{8}{238.767}{314.124}{obviously}%
\TextAt{18}{242.35}{286.366}{Z}%
\TextAt{3}{250.32}{291.303}{\Character{3}}%
\TextAt{21}{250.32}{283.411}{p}%
\TextAt{9}{258.402}{286.366}{\Character{61}}%
\TextAt{16}{270.827}{286.366}{\Character{102}}%
\TextAt{13}{276.805}{286.366}{x}%
\NativeMathAt{286.778}{286.366}{11.955}{\in}%
\TextAt{18}{298.069}{286.366}{Z}%
\TextAt{21}{306.039}{284.573}{p}%
\TextAt{16}{314.121}{286.366}{\Character{106}}%
\TextAt{16}{320.744}{286.366}{\Character{106}}%
\TextAt{13}{324.083}{286.366}{x}%
\TextAt{16}{330.735}{286.366}{\Character{106}}%
\TextAt{21}{334.056}{284.573}{p}%
\TextAt{9}{342.138}{286.366}{\Character{61}}%
\TextAt{9}{354.548}{286.366}{1}%
\TextAt{16}{360.417}{286.366}{\Character{103}}%
\TextAt{13}{366.394}{286.366}{\Character{58}}%
\TextAt{8}{79.226}{258.609}{Every}%
\TextAt{8}{112.88}{258.609}{element}%
\TextAt{13}{156.122}{258.609}{x}%
\NativeMathAt{166.095}{258.609}{11.955}{\in}%
\TextAt{18}{177.386}{258.609}{Q}%
\TextAt{3}{186.684}{262.947}{\Character{3}}%
\TextAt{21}{186.684}{255.653}{p}%
\TextAt{8}{195.342}{258.609}{admits}%
\TextAt{8}{233.778}{258.609}{a}%
\TextAt{8}{243.534}{258.609}{unique}%
\TextAt{8}{281.575}{258.609}{representation}%
\TextAt{13}{210.881}{230.851}{x}%
\TextAt{9}{220.854}{230.851}{\Character{61}}%
\TextAt{13}{233.28}{230.851}{p}%
\TextAt{21}{239.155}{235.787}{m}%
\TextAt{13}{247.143}{230.851}{u}%
\TextAt{8}{277.716}{230.851}{with}%
\TextAt{13}{304.376}{230.851}{m}%
\NativeMathAt{317.936}{230.851}{11.955}{\in}%
\TextAt{18}{329.227}{230.851}{Z}%
\TextAt{8}{341.095}{230.851}{and}%
\TextAt{13}{363.857}{230.851}{u}%
\NativeMathAt{373.841}{230.851}{11.955}{\in}%
\TextAt{18}{385.132}{230.851}{Z}%
\TextAt{3}{393.102}{235.787}{\Character{3}}%
\TextAt{21}{393.102}{227.896}{p}%
\TextAt{13}{397.863}{230.851}{\Character{58}}%
\TextAt{8}{79.226}{203.094}{For}%
\TextAt{8}{101.416}{203.094}{if}%
\TextAt{13}{113.361}{203.094}{v}%
\TextAt{21}{119.02}{201.3}{p}%
\TextAt{9}{123.781}{203.094}{\Character{40}}%
\TextAt{13}{128.333}{203.094}{x}%
\TextAt{9}{134.985}{203.094}{\Character{41}}%
\TextAt{9}{144.92}{203.094}{\Character{61}}%
\TextAt{13}{159.445}{203.094}{m}%
\NativeMathAt{175.085}{203.094}{11.955}{\in}%
\TextAt{18}{188.456}{203.094}{Z}%
\TextAt{8}{196.426}{203.094}{,}%
\TextAt{8}{205.102}{203.094}{then}%
\TextAt{13}{232.982}{203.094}{v}%
\TextAt{21}{238.642}{201.3}{p}%
\TextAt{9}{243.402}{203.094}{\Character{40}}%
\TextAt{13}{247.955}{203.094}{xp}%
\TextAt{3}{260.482}{207.432}{\Character{0}}%
\TextAt{21}{267.069}{207.432}{m}%
\TextAt{9}{275.057}{203.094}{\Character{41}}%
\TextAt{9}{284.992}{203.094}{\Character{61}}%
\TextAt{9}{299.493}{203.094}{0}%
\TextAt{8}{305.369}{203.094}{,}%
\TextAt{8}{314.045}{203.094}{hence}%
\TextAt{16}{347.771}{203.094}{\Character{106}}%
\TextAt{13}{351.092}{203.094}{xp}%
\TextAt{3}{363.62}{207.432}{\Character{0}}%
\TextAt{21}{370.206}{207.432}{m}%
\TextAt{16}{378.195}{203.094}{\Character{106}}%
\TextAt{21}{381.516}{201.3}{p}%
\TextAt{9}{391.678}{203.094}{\Character{61}}%
\TextAt{9}{406.168}{203.094}{1}%
\TextAt{8}{412.036}{203.094}{,}%
\TextAt{8}{420.712}{203.094}{i.e.,}%
\TextAt{13}{444.343}{203.094}{u}%
\TextAt{9}{456.406}{203.094}{\Character{61}}%
\TextAt{13}{470.912}{203.094}{xp}%
\TextAt{3}{483.439}{207.432}{\Character{0}}%
\TextAt{21}{490.026}{207.432}{m}%
\NativeMathAt{503.415}{203.094}{11.955}{\in}%
\TextAt{18}{516.787}{203.094}{Z}%
\TextAt{3}{524.757}{207.432}{\Character{3}}%
\TextAt{21}{524.757}{200.138}{p}%
\TextAt{8}{529.518}{203.094}{.}%
\TextAt{8}{79.226}{186.048}{Furthermore}%
\TextAt{8}{147.024}{186.048}{we}%
\TextAt{8}{164.251}{186.048}{have}%
\TextAt{8}{191.246}{186.048}{the}%
\TextAt{12}{79.226}{160.391}{Proposition}%
\TextAt{12}{151.937}{160.391}{(2.4).}%
\TextAt{8}{187.253}{160.391}{The}%
\TextAt{8}{211.307}{160.391}{nonzero}%
\TextAt{8}{254.883}{160.391}{ideals}%
\TextAt{8}{287.461}{160.391}{of}%
\TextAt{8}{300.791}{160.391}{the}%
\TextAt{8}{320.948}{160.391}{ring}%
\TextAt{18}{345.013}{160.391}{Z}%
\TextAt{21}{352.984}{158.597}{p}%
\TextAt{8}{361.642}{160.391}{are}%
\TextAt{8}{381.153}{160.391}{the}%
\TextAt{8}{401.309}{160.391}{principal}%
\TextAt{8}{450.086}{160.391}{ideals}%
\TextAt{13}{231.205}{132.633}{p}%
\TextAt{21}{237.08}{137.569}{n}%
\TextAt{18}{242.716}{132.633}{Z}%
\TextAt{21}{250.686}{130.84}{p}%
\TextAt{9}{258.768}{132.633}{\Character{61}}%
\TextAt{16}{271.194}{132.633}{\Character{102}}%
\TextAt{13}{277.171}{132.633}{x}%
\NativeMathAt{287.144}{132.633}{11.955}{\in}%
\TextAt{18}{298.435}{132.633}{Q}%
\TextAt{21}{307.734}{130.84}{p}%
\TextAt{16}{315.815}{132.633}{\Character{106}}%
\TextAt{13}{322.457}{132.633}{v}%
\TextAt{21}{328.116}{130.84}{p}%
\TextAt{9}{332.877}{132.633}{\Character{40}}%
\TextAt{13}{337.43}{132.633}{x}%
\TextAt{9}{344.082}{132.633}{\Character{41}}%
\NativeMathAt{351.955}{132.633}{11.955}{\geq}%
\TextAt{13}{364.574}{132.633}{n}%
\TextAt{16}{371.562}{132.633}{\Character{103}}%
\TextAt{13}{377.539}{132.633}{\Character{59}}%
\TextAt{8}{79.226}{104.875}{with}%
\TextAt{13}{105.886}{104.875}{n}%
\NativeMathAt{116.195}{104.875}{11.955}{\geq}%
\TextAt{9}{128.814}{104.875}{0}%
\TextAt{8}{134.667}{104.875}{,}%
\TextAt{8}{141.816}{104.875}{and}%
\TextAt{8}{164.579}{104.875}{one}%
\TextAt{8}{186.038}{104.875}{has}%
\TextAt{18}{258.336}{87.829}{Z}%
\TextAt{21}{266.306}{86.036}{p}%
\TextAt{9}{271.067}{87.829}{\Character{47}}%
\TextAt{13}{276.92}{87.829}{p}%
\TextAt{21}{282.795}{92.766}{n}%
\TextAt{18}{288.432}{87.829}{Z}%
\TextAt{21}{296.402}{86.036}{p}%
\NativeMathAt{304.484}{87.829}{11.955}{\cong}%
\TextAt{18}{317.103}{87.829}{Z}%
\TextAt{9}{325.073}{87.829}{\Character{47}}%
\TextAt{13}{330.926}{87.829}{p}%
\TextAt{21}{336.801}{92.766}{n}%
\TextAt{18}{342.438}{87.829}{Z}%
\TextAt{13}{350.408}{87.829}{\Character{58}}%
\TextAt{8}{297.211}{43.033}{102}%
\Destination{78.23}{809.815}{nav.page.102}%
\Destination{79.36}{498.555}{nav.print.112}%
\Destination{79.36}{587.255}{nav.number.2.2.3}%
\Destination{79.36}{174.355}{nav.number.2.2.4}%
\Destination{79.226}{410.908}{errata.73}%
}

% Page 115 (body page 103)
\Page{%
\TextAt{2}{79.226}{796.943}{II.}%
\TextAt{2}{94.073}{796.943}{The}%
\TextAt{2}{116.492}{796.943}{Theory}%
\TextAt{2}{154.401}{796.943}{of}%
\TextAt{2}{166.826}{796.943}{Valuations}%
\Vector{q 0 G 0 g 0.299 w 7.226 18.729 m 460.769 18.729 l S Q}%
\TextAt{19}{79.226}{761.124}{Proof.}%
\TextAt{8}{116}{761.124}{Let}%
\TextAt{14}{136.995}{761.124}{a}%
\NativeMathAt{146.341}{761.124}{11.955}{\neq}%
\TextAt{9}{158.798}{761.124}{\Character{40}0\Character{41}}%
\TextAt{8}{177.69}{761.124}{be}%
\TextAt{8}{193.639}{761.124}{an}%
\TextAt{8}{209.923}{761.124}{ideal}%
\TextAt{8}{237.912}{761.124}{of}%
\TextAt{18}{251.267}{761.124}{Z}%
\TextAt{21}{259.237}{759.33}{p}%
\TextAt{8}{267.92}{761.124}{and}%
\TextAt{13}{290.708}{761.124}{x}%
\TextAt{9}{300.724}{761.124}{\Character{61}}%
\TextAt{13}{313.192}{761.124}{p}%
\TextAt{21}{319.067}{765.462}{m}%
\TextAt{13}{327.056}{761.124}{u}%
\TextAt{8}{333.719}{761.124}{,}%
\TextAt{13}{340.899}{761.124}{u}%
\NativeMathAt{350.925}{761.124}{11.955}{\in}%
\TextAt{18}{362.259}{761.124}{Z}%
\TextAt{3}{370.229}{765.462}{\Character{3}}%
\TextAt{21}{370.229}{758.168}{p}%
\TextAt{8}{374.99}{761.124}{,}%
\TextAt{8}{382.171}{761.124}{an}%
\TextAt{8}{398.455}{761.124}{element}%
\TextAt{8}{441.722}{761.124}{of}%
\TextAt{14}{455.077}{761.124}{a}%
\TextAt{8}{464.983}{761.124}{with}%
\TextAt{8}{491.668}{761.124}{smallest}%
\TextAt{8}{79.226}{744.078}{possible}%
\TextAt{13}{122.673}{744.078}{m}%
\TextAt{8}{136.238}{744.078}{(since}%
\TextAt{16}{168.89}{744.078}{\Character{106}}%
\TextAt{13}{172.211}{744.078}{x}%
\TextAt{16}{178.863}{744.078}{\Character{106}}%
\TextAt{21}{182.184}{742.284}{p}%
\NativeMathAt{190.266}{744.078}{11.955}{\leq}%
\TextAt{9}{202.885}{744.078}{1}%
\TextAt{8}{208.738}{744.078}{,}%
\TextAt{8}{215.43}{744.078}{one}%
\TextAt{8}{236.318}{744.078}{has}%
\TextAt{13}{256.62}{744.078}{m}%
\NativeMathAt{270.18}{744.078}{11.955}{\geq}%
\TextAt{9}{282.799}{744.078}{0}%
\TextAt{8}{288.652}{744.078}{).}%
\TextAt{8}{301.465}{744.078}{Then}%
\TextAt{14}{331.451}{744.078}{a}%
\TextAt{9}{340.755}{744.078}{\Character{61}}%
\TextAt{13}{353.18}{744.078}{p}%
\TextAt{21}{359.055}{748.416}{m}%
\TextAt{18}{367.044}{744.078}{Z}%
\TextAt{21}{375.014}{742.284}{p}%
\TextAt{8}{383.101}{744.078}{because}%
\TextAt{13}{425.831}{744.078}{y}%
\TextAt{9}{435.288}{744.078}{\Character{61}}%
\TextAt{13}{447.714}{744.078}{p}%
\TextAt{21}{453.589}{748.416}{n}%
\TextAt{13}{459.225}{744.078}{u}%
\TextAt{3}{465.887}{748.416}{\Character{48}}%
\NativeMathAt{472.003}{744.078}{11.955}{\in}%
\TextAt{14}{483.294}{744.078}{a}%
\TextAt{8}{489.277}{744.078}{,}%
\TextAt{13}{495.969}{744.078}{u}%
\TextAt{3}{502.632}{748.416}{\Character{48}}%
\NativeMathAt{508.748}{744.078}{11.955}{\in}%
\TextAt{18}{520.039}{744.078}{Z}%
\TextAt{3}{528.009}{748.416}{\Character{3}}%
\TextAt{21}{528.009}{741.122}{p}%
\TextAt{8}{79.226}{727.032}{implies}%
\TextAt{13}{118.953}{727.032}{n}%
\NativeMathAt{129.262}{727.032}{11.955}{\geq}%
\TextAt{13}{141.881}{727.032}{m}%
\TextAt{8}{152.12}{727.032}{,}%
\TextAt{8}{159.269}{727.032}{hence}%
\TextAt{13}{191.776}{727.032}{y}%
\TextAt{9}{201.233}{727.032}{\Character{61}}%
\TextAt{9}{213.643}{727.032}{\Character{40}}%
\TextAt{13}{218.211}{727.032}{p}%
\TextAt{21}{224.086}{731.37}{n}%
\TextAt{3}{229.224}{731.37}{\Character{0}}%
\TextAt{21}{235.811}{731.37}{m}%
\TextAt{13}{243.799}{727.032}{u}%
\TextAt{3}{250.462}{731.37}{\Character{48}}%
\TextAt{9}{253.257}{727.032}{\Character{41}}%
\TextAt{13}{257.809}{727.032}{p}%
\TextAt{21}{263.684}{731.37}{m}%
\NativeMathAt{274.994}{727.032}{11.955}{\in}%
\TextAt{13}{286.285}{727.032}{p}%
\TextAt{21}{292.16}{731.37}{m}%
\TextAt{18}{300.149}{727.032}{Z}%
\TextAt{21}{308.119}{725.238}{p}%
\TextAt{8}{312.88}{727.032}{.}%
\TextAt{8}{321.328}{727.032}{The}%
\TextAt{8}{345.382}{727.032}{homomorphism}%
\TextAt{18}{205.848}{699.572}{Z}%
\NativeArrowAt{217.139}{699.572}{11.955}{19.26}{\longrightarrow}%
\TextAt{18}{239.721}{699.572}{Z}%
\TextAt{21}{247.691}{697.778}{p}%
\TextAt{9}{252.452}{699.572}{\Character{47}}%
\TextAt{13}{258.305}{699.572}{p}%
\TextAt{21}{264.18}{704.508}{n}%
\TextAt{18}{269.816}{699.572}{Z}%
\TextAt{21}{277.787}{697.778}{p}%
\TextAt{13}{282.548}{699.572}{\Character{59}}%
\TextAt{13}{311.683}{699.572}{a}%
\NativeArrowAt{321.168}{699.572}{11.955}{19.26}{\longmapsto}%
\TextAt{13}{343.75}{699.572}{a}%
\TextAt{8}{353.216}{699.572}{mod}%
\TextAt{13}{378.653}{699.572}{p}%
\TextAt{21}{384.529}{704.508}{n}%
\TextAt{18}{390.165}{699.572}{Z}%
\TextAt{21}{398.135}{697.778}{p}%
\TextAt{13}{402.896}{699.572}{\Character{59}}%
\TextAt{8}{79.226}{672.111}{has}%
\TextAt{8}{100.128}{672.111}{kernel}%
\TextAt{13}{134.623}{672.111}{p}%
\TextAt{21}{140.498}{676.45}{n}%
\TextAt{18}{146.135}{672.111}{Z}%
\TextAt{8}{158.03}{672.111}{and}%
\TextAt{8}{180.821}{672.111}{is}%
\TextAt{8}{192.613}{672.111}{surjective.}%
\TextAt{8}{249.66}{672.111}{Indeed,}%
\TextAt{8}{290.977}{672.111}{for}%
\TextAt{8}{308.89}{672.111}{every}%
\TextAt{13}{339.81}{672.111}{x}%
\NativeMathAt{349.831}{672.111}{11.955}{\in}%
\TextAt{18}{361.17}{672.111}{Z}%
\TextAt{21}{369.14}{670.318}{p}%
\TextAt{8}{373.901}{672.111}{,}%
\TextAt{8}{381.086}{672.111}{there}%
\TextAt{8}{411.026}{672.111}{exists}%
\TextAt{8}{443.369}{672.111}{by}%
\TextAt{8}{459.656}{672.111}{(2.3)}%
\TextAt{8}{487.659}{672.111}{an}%
\TextAt{13}{503.947}{672.111}{a}%
\NativeMathAt{513.46}{672.111}{11.955}{\in}%
\TextAt{18}{524.799}{672.111}{Z}%
\TextAt{8}{79.226}{655.066}{such}%
\TextAt{8}{105.623}{655.066}{that}%
\TextAt{16}{270.045}{638.02}{\Character{106}}%
\TextAt{13}{273.366}{638.02}{x}%
\TextAt{16}{282.675}{638.02}{\Character{0}}%
\TextAt{13}{294.63}{638.02}{a}%
\TextAt{16}{300.775}{638.02}{\Character{106}}%
\TextAt{21}{304.096}{636.226}{p}%
\NativeMathAt{312.177}{638.02}{11.955}{\leq}%
\TextAt{9}{328.821}{646.107}{1}%
\Vector{q 0 G 0 g 0.478 w 253.992 -132.992 m 265.504 -132.992 l S Q}%
\TextAt{13}{325.992}{629.819}{p}%
\TextAt{21}{331.867}{633.273}{n}%
\TextAt{13}{338.699}{638.02}{\Character{59}}%
\TextAt{8}{79.226}{611.716}{i.e.,}%
\TextAt{13}{101.696}{611.716}{v}%
\TextAt{21}{107.356}{609.923}{p}%
\TextAt{9}{112.117}{611.716}{\Character{40}}%
\TextAt{13}{116.669}{611.716}{x}%
\TextAt{16}{126.177}{611.716}{\Character{0}}%
\TextAt{13}{138.331}{611.716}{a}%
\TextAt{9}{144.476}{611.716}{\Character{41}}%
\NativeMathAt{152.847}{611.716}{11.955}{\geq}%
\TextAt{13}{165.964}{611.716}{n}%
\TextAt{8}{172.952}{611.716}{,}%
\TextAt{8}{180.466}{611.716}{therefore}%
\TextAt{13}{229.859}{611.716}{x}%
\TextAt{16}{239.366}{611.716}{\Character{0}}%
\TextAt{13}{251.521}{611.716}{a}%
\NativeMathAt{261.484}{611.716}{11.955}{\in}%
\TextAt{13}{273.273}{611.716}{p}%
\TextAt{21}{279.148}{616.055}{n}%
\TextAt{18}{284.785}{611.716}{Z}%
\TextAt{21}{292.755}{609.923}{p}%
\TextAt{8}{301.705}{611.716}{and}%
\TextAt{8}{324.76}{611.716}{hence}%
\TextAt{13}{357.558}{611.716}{x}%
\NativeMathAt{368.029}{611.716}{11.955}{\equiv}%
\TextAt{13}{381.147}{611.716}{a}%
\TextAt{8}{390.612}{611.716}{mod}%
\TextAt{13}{416.05}{611.716}{p}%
\TextAt{21}{421.925}{616.055}{n}%
\TextAt{18}{427.562}{611.716}{Z}%
\TextAt{21}{435.532}{609.923}{p}%
\TextAt{8}{440.293}{611.716}{.}%
\TextAt{8}{449.618}{611.716}{So}%
\TextAt{8}{466.169}{611.716}{we}%
\TextAt{8}{483.688}{611.716}{obtain}%
\TextAt{8}{520.408}{611.716}{an}%
\TextAt{8}{79.226}{594.67}{isomorphism}%
\TextAt{18}{258.336}{577.624}{Z}%
\TextAt{21}{266.306}{575.831}{p}%
\TextAt{9}{271.067}{577.624}{\Character{47}}%
\TextAt{13}{276.92}{577.624}{p}%
\TextAt{21}{282.795}{582.561}{n}%
\TextAt{18}{288.432}{577.624}{Z}%
\TextAt{21}{296.402}{575.831}{p}%
\NativeMathAt{304.484}{577.624}{11.955}{\cong}%
\TextAt{18}{317.103}{577.624}{Z}%
\TextAt{9}{325.073}{577.624}{\Character{47}}%
\TextAt{13}{330.926}{577.624}{p}%
\TextAt{21}{336.801}{582.561}{n}%
\TextAt{18}{342.438}{577.624}{Z}%
\TextAt{13}{350.408}{577.624}{\Character{58}}%
\Vector{q 0 G 0 g 0.399 w 452.334 -219.237 m 452.334 -211.167 l S Q}%
\Vector{q 0 G 0 g 0.399 w 452.534 -211.366 m 459.707 -211.366 l S Q}%
\Vector{q 0 G 0 g 0.399 w 452.534 -219.038 m 459.707 -219.038 l S Q}%
\Vector{q 0 G 0 g 0.399 w 459.906 -219.237 m 459.906 -211.167 l S Q}%
\TextAt{8}{97.159}{526.202}{We}%
\TextAt{8}{116.934}{526.202}{now}%
\TextAt{8}{140.954}{526.202}{want}%
\TextAt{8}{169.205}{526.202}{to}%
\TextAt{8}{183.147}{526.202}{establish}%
\TextAt{8}{231.029}{526.202}{the}%
\TextAt{8}{250.816}{526.202}{link}%
\TextAt{8}{273.533}{526.202}{with}%
\TextAt{8}{299.824}{526.202}{Hensel’s}%
\TextAt{8}{344.765}{526.202}{definition}%
\TextAt{8}{396.424}{526.202}{of}%
\TextAt{8}{409.385}{526.202}{the}%
\TextAt{8}{429.173}{526.202}{ring}%
\TextAt{18}{452.87}{526.202}{Z}%
\TextAt{21}{460.84}{524.409}{p}%
\TextAt{8}{469.129}{526.202}{and}%
\TextAt{8}{491.523}{526.202}{the}%
\TextAt{8}{511.31}{526.202}{field}%
\TextAt{18}{79.226}{509.156}{Q}%
\TextAt{21}{88.525}{507.363}{p}%
\TextAt{8}{97.183}{509.156}{which}%
\TextAt{8}{130.669}{509.156}{was}%
\TextAt{8}{153.169}{509.156}{given}%
\TextAt{8}{183.738}{509.156}{in}%
\TextAt{8}{197.391}{509.156}{§1.}%
\TextAt{8}{217.376}{509.156}{There}%
\TextAt{8}{251.185}{509.156}{we}%
\TextAt{8}{268.413}{509.156}{defined}%
\TextAt{8}{308.726}{509.156}{the}%
\TextAt{13}{328.882}{509.156}{p}%
\TextAt{8}{334.757}{509.156}{-adic}%
\TextAt{8}{363.366}{509.156}{integers}%
\TextAt{8}{406.68}{509.156}{as}%
\TextAt{8}{421.05}{509.156}{formal}%
\TextAt{8}{457.8}{509.156}{series}%
\NativeMathAt{246.001}{489.776}{7.97}{\infty}%
\TextAt{20}{241.601}{486.189}{\Character{88}}%
\TextAt{21}{242.483}{461.083}{\Character{23}}%
\TextAt{17}{247.167}{461.083}{\Character{61}0}%
\TextAt{13}{260.862}{474.832}{a}%
\TextAt{21}{267.007}{473.039}{\Character{23}}%
\TextAt{13}{272.189}{474.832}{p}%
\TextAt{21}{278.065}{479.768}{\Character{23}}%
\TextAt{13}{283.247}{474.832}{\Character{59}}%
\TextAt{9}{312.401}{474.832}{0}%
\NativeMathAt{321.575}{474.832}{11.955}{\leq}%
\TextAt{13}{334.194}{474.832}{a}%
\TextAt{21}{340.339}{473.039}{\Character{23}}%
\TextAt{13}{348.842}{474.832}{\Character{60}}%
\TextAt{13}{361.252}{474.832}{p\Character{59}}%
\TextAt{8}{79.226}{440.061}{which}%
\TextAt{8}{112.713}{440.061}{we}%
\TextAt{8}{129.94}{440.061}{identified}%
\TextAt{8}{180.989}{440.061}{with}%
\TextAt{8}{207.649}{440.061}{sequences}%
\TextAt{9}{196.828}{412.601}{\Character{22}}%
\TextAt{13}{196.348}{412.601}{s}%
\TextAt{21}{201.862}{410.807}{n}%
\TextAt{9}{210.819}{412.601}{\Character{61}}%
\TextAt{13}{223.244}{412.601}{s}%
\TextAt{21}{228.758}{410.807}{n}%
\TextAt{8}{237.715}{412.601}{mod}%
\TextAt{13}{263.153}{412.601}{p}%
\TextAt{21}{269.028}{417.537}{n}%
\NativeMathAt{277.986}{412.601}{11.955}{\in}%
\TextAt{18}{289.277}{412.601}{Z}%
\TextAt{9}{297.247}{412.601}{\Character{47}}%
\TextAt{13}{303.1}{412.601}{p}%
\TextAt{21}{308.975}{417.537}{n}%
\TextAt{18}{314.611}{412.601}{Z}%
\TextAt{13}{322.581}{412.601}{\Character{59}}%
\TextAt{13}{351.716}{412.601}{n}%
\TextAt{9}{362.044}{412.601}{\Character{61}}%
\TextAt{9}{374.454}{412.601}{1}%
\TextAt{13}{380.323}{412.601}{\Character{59}}%
\TextAt{9}{385.567}{412.601}{2}%
\TextAt{13}{391.42}{412.601}{\Character{59}}%
\TextAt{13}{396.644}{412.601}{\Character{58}}%
\TextAt{13}{401.881}{412.601}{\Character{58}}%
\TextAt{13}{407.105}{412.601}{\Character{58}}%
\TextAt{13}{412.342}{412.601}{\Character{59}}%
\TextAt{8}{79.226}{385.141}{where}%
\TextAt{13}{113.035}{385.141}{s}%
\TextAt{21}{118.549}{383.347}{n}%
\TextAt{8}{128.083}{385.141}{was}%
\TextAt{8}{150.583}{385.141}{the}%
\TextAt{8}{170.739}{385.141}{partial}%
\TextAt{8}{208.47}{385.141}{sum}%
\TextAt{13}{270.101}{360.549}{s}%
\TextAt{21}{275.615}{358.756}{n}%
\TextAt{9}{284.572}{360.549}{\Character{61}}%
\TextAt{21}{297.653}{375.493}{n}%
\TextAt{3}{302.791}{375.493}{\Character{0}}%
\TextAt{17}{309.377}{375.493}{1}%
\TextAt{20}{296.998}{371.907}{\Character{88}}%
\TextAt{21}{297.88}{346.801}{\Character{23}}%
\TextAt{17}{302.564}{346.801}{\Character{61}0}%
\TextAt{13}{316.259}{360.549}{a}%
\TextAt{21}{322.404}{358.756}{\Character{23}}%
\TextAt{13}{327.586}{360.549}{p}%
\TextAt{21}{333.461}{365.485}{\Character{23}}%
\TextAt{13}{338.643}{360.549}{\Character{58}}%
\TextAt{8}{79.226}{330.497}{These}%
\TextAt{8}{113.095}{330.497}{sequences}%
\TextAt{8}{166.212}{330.497}{constituted}%
\TextAt{8}{227.41}{330.497}{the}%
\TextAt{8}{247.567}{330.497}{projective}%
\TextAt{8}{301.544}{330.497}{limit}%
\NativeMathAt{180.5}{295.726}{11.955}{\varprojlim}%
\TextAt{21}{186.06}{283.611}{n}%
\TextAt{18}{198.751}{295.726}{Z}%
\TextAt{9}{206.722}{295.726}{\Character{47}}%
\TextAt{13}{212.575}{295.726}{p}%
\TextAt{21}{218.45}{300.662}{n}%
\TextAt{18}{224.086}{295.726}{Z}%
\TextAt{9}{235.377}{295.726}{\Character{61}}%
\TextAt{20}{247.802}{316.169}{\Character{40}}%
\TextAt{9}{257.433}{295.726}{\Character{40}}%
\TextAt{13}{261.985}{295.726}{x}%
\TextAt{21}{268.637}{293.932}{n}%
\TextAt{9}{274.274}{295.726}{\Character{41}}%
\NativeMathAt{282.147}{295.726}{11.955}{\in}%
\NativeMathAt{297.183}{310.67}{7.97}{\infty}%
\TextAt{20}{293.779}{307.083}{\Character{89}}%
\TextAt{21}{293.438}{281.977}{n}%
\TextAt{17}{298.576}{281.977}{\Character{61}1}%
\TextAt{18}{311.389}{295.726}{Z}%
\TextAt{9}{319.359}{295.726}{\Character{47}}%
\TextAt{13}{325.212}{295.726}{p}%
\TextAt{21}{331.088}{300.662}{n}%
\TextAt{18}{336.724}{295.726}{Z}%
\TextAt{20}{348.015}{316.647}{\Character{12}}%
\TextAt{20}{348.015}{309.474}{\Character{12}}%
\TextAt{20}{348.015}{302.301}{\Character{12}}%
\TextAt{20}{348.015}{295.128}{\Character{12}}%
\TextAt{20}{348.015}{287.955}{\Character{12}}%
\TextAt{13}{355.321}{295.726}{x}%
\TextAt{21}{361.973}{293.933}{n}%
\TextAt{17}{367.111}{293.933}{\Character{43}1}%
\NativeArrowAt{381.751}{295.726}{11.955}{19.26}{\longmapsto}%
\TextAt{13}{404.333}{295.726}{x}%
\TextAt{21}{410.985}{293.933}{n}%
\TextAt{20}{416.621}{316.169}{\Character{41}}%
\TextAt{13}{428.244}{295.726}{\Character{58}}%
\TextAt{8}{79.226}{260.955}{We}%
\TextAt{8}{99.371}{260.955}{viewed}%
\TextAt{8}{137.735}{260.955}{the}%
\TextAt{13}{157.891}{260.955}{p}%
\TextAt{8}{163.766}{260.955}{-adic}%
\TextAt{8}{192.375}{260.955}{integers}%
\TextAt{8}{235.689}{260.955}{as}%
\TextAt{8}{250.059}{260.955}{elements}%
\TextAt{8}{297.915}{260.955}{of}%
\TextAt{8}{311.245}{260.955}{this}%
\TextAt{8}{334.068}{260.955}{ring.}%
\TextAt{8}{362.684}{260.955}{Since}%
\TextAt{18}{258.336}{233.495}{Z}%
\TextAt{21}{266.306}{231.701}{p}%
\TextAt{9}{271.067}{233.495}{\Character{47}}%
\TextAt{13}{276.92}{233.495}{p}%
\TextAt{21}{282.795}{238.431}{n}%
\TextAt{18}{288.432}{233.495}{Z}%
\TextAt{21}{296.402}{231.701}{p}%
\NativeMathAt{304.484}{233.494}{11.955}{\cong}%
\TextAt{18}{317.103}{233.495}{Z}%
\TextAt{9}{325.073}{233.495}{\Character{47}}%
\TextAt{13}{330.926}{233.495}{p}%
\TextAt{21}{336.801}{238.431}{n}%
\TextAt{18}{342.438}{233.495}{Z}%
\TextAt{13}{350.408}{233.495}{\Character{59}}%
\TextAt{8}{79.226}{206.034}{we}%
\TextAt{8}{96.454}{206.034}{obtain,}%
\TextAt{8}{136.133}{206.034}{for}%
\TextAt{8}{154.018}{206.034}{every}%
\TextAt{13}{184.91}{206.034}{n}%
\NativeMathAt{195.218}{206.034}{11.955}{\geq}%
\TextAt{9}{207.838}{206.034}{1}%
\TextAt{8}{213.691}{206.034}{,}%
\TextAt{8}{220.84}{206.034}{a}%
\TextAt{8}{230.595}{206.034}{surjective}%
\TextAt{8}{283.007}{206.034}{homomorphism}%
\TextAt{18}{268.403}{178.574}{Z}%
\TextAt{21}{276.373}{176.781}{p}%
\NativeArrowAt{284.455}{178.574}{11.955}{19.26}{\longrightarrow}%
\TextAt{18}{307.037}{178.574}{Z}%
\TextAt{9}{315.007}{178.574}{\Character{47}}%
\TextAt{13}{320.86}{178.574}{p}%
\TextAt{21}{326.735}{183.511}{n}%
\TextAt{18}{332.371}{178.574}{Z}%
\TextAt{13}{340.341}{178.574}{\Character{58}}%
\TextAt{8}{79.226}{151.114}{It}%
\TextAt{8}{91.899}{151.114}{is}%
\TextAt{8}{103.663}{151.114}{clear}%
\TextAt{8}{131.626}{151.114}{that}%
\TextAt{8}{156.995}{151.114}{the}%
\TextAt{8}{177.151}{151.114}{family}%
\TextAt{8}{212.921}{151.114}{of}%
\TextAt{8}{226.251}{151.114}{these}%
\TextAt{8}{256.222}{151.114}{homomorphisms}%
\TextAt{8}{343.507}{151.114}{yields}%
\TextAt{8}{376.408}{151.114}{a}%
\TextAt{8}{386.163}{151.114}{homomorphism}%
\TextAt{18}{259.277}{123.654}{Z}%
\TextAt{21}{267.247}{121.861}{p}%
\NativeArrowAt{275.329}{123.654}{11.955}{19.26}{\longrightarrow}%
\NativeMathAt{297.911}{123.654}{11.955}{\varprojlim}%
\TextAt{21}{303.471}{111.54}{n}%
\TextAt{18}{316.162}{123.654}{Z}%
\TextAt{9}{324.132}{123.654}{\Character{47}}%
\TextAt{13}{329.985}{123.654}{p}%
\TextAt{21}{335.861}{128.59}{n}%
\TextAt{18}{341.497}{123.654}{Z}%
\TextAt{13}{349.467}{123.654}{\Character{58}}%
\TextAt{8}{79.226}{89.967}{It}%
\TextAt{8}{91.826}{89.967}{is}%
\TextAt{8}{103.518}{89.967}{now}%
\TextAt{8}{127.834}{89.967}{possible}%
\TextAt{8}{171.78}{89.967}{to}%
\TextAt{8}{186.018}{89.967}{identify}%
\TextAt{8}{228.542}{89.967}{both}%
\TextAt{8}{256.11}{89.967}{definitions}%
\TextAt{8}{312.681}{89.967}{given}%
\TextAt{8}{343.178}{89.967}{for}%
\TextAt{18}{360.991}{89.967}{Z}%
\TextAt{21}{368.961}{88.174}{p}%
\TextAt{8}{377.547}{89.967}{(and}%
\TextAt{8}{404.792}{89.967}{therefore}%
\TextAt{8}{453.819}{89.967}{also}%
\TextAt{8}{477.227}{89.967}{for}%
\TextAt{18}{495.04}{89.967}{Q}%
\TextAt{21}{504.338}{88.174}{p}%
\TextAt{8}{509.099}{89.967}{)}%
\TextAt{8}{517.479}{89.967}{via}%
\TextAt{8}{79.226}{72.921}{the}%
\TextAt{8}{297.211}{43.033}{103}%
\Destination{78.23}{809.815}{nav.page.103}%
\Destination{79.36}{541.155}{nav.print.113}%
\Hotlink{202.57}{506.955}{6.858}{11.047}{nav.section.2.1}%
\Hotlink{459.156}{669.955}{25.078}{11.047}{nav.number.2.2.3}%
}

% Page 116 (body page 104)
\Page{%
\TextAt{2}{79.226}{796.943}{II.}%
\TextAt{2}{94.073}{796.943}{The}%
\TextAt{2}{116.492}{796.943}{Theory}%
\TextAt{2}{154.401}{796.943}{of}%
\TextAt{2}{166.826}{796.943}{Valuations}%
\Vector{q 0 G 0 g 0.299 w 7.226 18.729 m 460.769 18.729 l S Q}%
\TextAt{12}{79.226}{761.124}{Proposition}%
\TextAt{12}{151.937}{761.124}{(2.5).}%
\TextAt{8}{187.253}{761.124}{The}%
\TextAt{8}{211.307}{761.124}{homomorphism}%
\TextAt{18}{260.903}{733.767}{Z}%
\TextAt{21}{268.873}{731.974}{p}%
\NativeArrowAt{276.955}{733.767}{11.955}{19.26}{\longrightarrow}%
\NativeMathAt{299.537}{733.767}{11.955}{\varprojlim}%
\TextAt{21}{305.097}{721.653}{n}%
\TextAt{18}{317.788}{733.767}{Z}%
\TextAt{9}{325.758}{733.767}{\Character{47}}%
\TextAt{13}{331.611}{733.767}{p}%
\TextAt{21}{337.486}{738.703}{n}%
\TextAt{18}{343.123}{733.767}{Z}%
\TextAt{8}{79.226}{700.853}{is}%
\TextAt{8}{90.99}{700.853}{an}%
\TextAt{8}{107.249}{700.853}{isomorphism.}%
\TextAt{19}{79.226}{666.814}{Proof.}%
\TextAt{8}{116}{666.814}{If}%
\TextAt{13}{126.632}{666.814}{x}%
\NativeMathAt{136.605}{666.814}{11.955}{\in}%
\TextAt{18}{147.896}{666.814}{Z}%
\TextAt{21}{155.866}{665.021}{p}%
\TextAt{8}{163.463}{666.814}{is}%
\TextAt{8}{174.166}{666.814}{mapped}%
\TextAt{8}{217.651}{666.814}{to}%
\TextAt{8}{230.9}{666.814}{zero,}%
\TextAt{8}{258.015}{666.814}{this}%
\TextAt{8}{279.777}{666.814}{means}%
\TextAt{8}{314.545}{666.814}{that}%
\TextAt{13}{338.854}{666.814}{x}%
\NativeMathAt{348.827}{666.814}{11.955}{\in}%
\TextAt{13}{360.118}{666.814}{p}%
\TextAt{21}{365.993}{671.153}{n}%
\TextAt{18}{371.629}{666.814}{Z}%
\TextAt{21}{379.599}{665.021}{p}%
\TextAt{8}{387.197}{666.814}{for}%
\TextAt{8}{404.021}{666.814}{all}%
\TextAt{13}{419.219}{666.814}{n}%
\NativeMathAt{429.528}{666.814}{11.955}{\geq}%
\TextAt{9}{442.147}{666.814}{1}%
\TextAt{8}{448}{666.814}{,}%
\TextAt{8}{454.301}{666.814}{i.e.,}%
\TextAt{16}{475.557}{666.814}{\Character{106}}%
\TextAt{13}{478.878}{666.814}{x}%
\TextAt{16}{485.53}{666.814}{\Character{106}}%
\TextAt{21}{488.851}{665.021}{p}%
\NativeMathAt{496.933}{666.814}{11.955}{\leq}%
\TextAt{9}{509.552}{666.814}{1\Character{47}}%
\TextAt{13}{521.258}{666.814}{p}%
\TextAt{21}{527.133}{671.153}{n}%
\TextAt{8}{79.226}{649.768}{for}%
\TextAt{8}{97.951}{649.768}{all}%
\TextAt{13}{115.051}{649.768}{n}%
\NativeMathAt{126.791}{649.768}{11.955}{\geq}%
\TextAt{9}{140.842}{649.768}{1}%
\TextAt{8}{146.695}{649.768}{,}%
\TextAt{8}{154.895}{649.768}{so}%
\TextAt{8}{170.105}{649.768}{that}%
\TextAt{16}{196.314}{649.768}{\Character{106}}%
\TextAt{13}{199.635}{649.768}{x}%
\TextAt{16}{206.287}{649.768}{\Character{106}}%
\TextAt{21}{209.608}{647.975}{p}%
\TextAt{9}{219.122}{649.768}{\Character{61}}%
\TextAt{9}{232.966}{649.768}{0}%
\TextAt{8}{243.569}{649.768}{and}%
\TextAt{8}{267.172}{649.768}{thus}%
\TextAt{13}{293.764}{649.768}{x}%
\TextAt{9}{305.169}{649.768}{\Character{61}}%
\TextAt{9}{319.013}{649.768}{0}%
\TextAt{8}{324.879}{649.768}{.}%
\TextAt{8}{335.848}{649.768}{This}%
\TextAt{8}{363.408}{649.768}{shows}%
\TextAt{8}{397.866}{649.768}{injectivity.}%
\TextAt{8}{459.238}{649.768}{An}%
\TextAt{8}{479.255}{649.768}{element}%
\TextAt{8}{523.337}{649.768}{of}%
\NativeMathAt{79.226}{632.722}{11.955}{\varprojlim}%
\TextAt{21}{95.485}{627.283}{n}%
\TextAt{18}{103.114}{632.722}{Z}%
\TextAt{9}{111.084}{632.722}{\Character{47}}%
\TextAt{13}{116.937}{632.722}{p}%
\TextAt{21}{122.812}{637.061}{n}%
\TextAt{18}{128.449}{632.722}{Z}%
\TextAt{8}{140.316}{632.722}{is}%
\TextAt{8}{152.08}{632.722}{given}%
\TextAt{8}{182.649}{632.722}{by}%
\TextAt{8}{198.908}{632.722}{a}%
\TextAt{8}{208.664}{632.722}{sequence}%
\TextAt{8}{257.166}{632.722}{of}%
\TextAt{8}{270.496}{632.722}{partial}%
\TextAt{8}{308.227}{632.722}{sums}%
\TextAt{13}{228.153}{594.7}{s}%
\TextAt{21}{233.667}{592.907}{n}%
\TextAt{9}{242.624}{594.7}{\Character{61}}%
\TextAt{21}{255.705}{609.644}{n}%
\TextAt{3}{260.843}{609.644}{\Character{0}}%
\TextAt{17}{267.429}{609.644}{1}%
\TextAt{20}{255.05}{606.058}{\Character{88}}%
\TextAt{21}{255.932}{580.952}{\Character{23}}%
\TextAt{17}{260.616}{580.952}{\Character{61}0}%
\TextAt{13}{274.311}{594.7}{a}%
\TextAt{21}{280.456}{592.907}{\Character{23}}%
\TextAt{13}{285.638}{594.7}{p}%
\TextAt{21}{291.513}{599.636}{\Character{23}}%
\TextAt{13}{296.695}{594.7}{\Character{59}}%
\TextAt{9}{325.849}{594.7}{0}%
\NativeMathAt{335.023}{594.7}{11.955}{\leq}%
\TextAt{13}{347.643}{594.7}{a}%
\TextAt{21}{353.787}{592.907}{\Character{23}}%
\TextAt{13}{362.29}{594.7}{\Character{60}}%
\TextAt{13}{374.7}{594.7}{p\Character{58}}%
\TextAt{8}{79.226}{560.033}{We}%
\TextAt{8}{99.858}{560.033}{saw}%
\TextAt{8}{122.844}{560.033}{above}%
\TextAt{8}{156.507}{560.033}{that}%
\TextAt{8}{182.363}{560.033}{this}%
\TextAt{8}{205.673}{560.033}{sequence}%
\TextAt{8}{254.662}{560.033}{is}%
\TextAt{8}{266.913}{560.033}{a}%
\TextAt{8}{277.155}{560.033}{Cauchy}%
\TextAt{8}{319.593}{560.033}{sequence}%
\TextAt{8}{368.582}{560.033}{in}%
\TextAt{18}{382.722}{560.033}{Z}%
\TextAt{21}{390.693}{558.24}{p}%
\TextAt{8}{395.453}{560.033}{,}%
\TextAt{8}{403.212}{560.033}{and}%
\TextAt{8}{426.461}{560.033}{thus}%
\TextAt{8}{452.7}{560.033}{converges}%
\TextAt{8}{505.61}{560.033}{to}%
\TextAt{8}{520.408}{560.033}{an}%
\TextAt{8}{79.226}{542.987}{element}%
\TextAt{13}{258.679}{522.3}{x}%
\TextAt{9}{268.652}{522.3}{\Character{61}}%
\NativeMathAt{285.477}{537.244}{7.97}{\infty}%
\TextAt{20}{281.077}{533.657}{\Character{88}}%
\TextAt{21}{281.959}{508.551}{\Character{23}}%
\TextAt{17}{286.643}{508.551}{\Character{61}0}%
\TextAt{13}{300.338}{522.3}{a}%
\TextAt{21}{306.483}{520.507}{\Character{23}}%
\TextAt{13}{311.665}{522.3}{p}%
\TextAt{21}{317.54}{527.236}{\Character{23}}%
\NativeMathAt{326.043}{522.3}{11.955}{\in}%
\TextAt{18}{337.334}{522.3}{Z}%
\TextAt{21}{345.304}{520.507}{p}%
\TextAt{13}{350.065}{522.3}{\Character{58}}%
\TextAt{8}{79.226}{492.172}{Since}%
\TextAt{13}{240.042}{471.354}{x}%
\TextAt{16}{249.351}{471.354}{\Character{0}}%
\TextAt{13}{261.306}{471.354}{s}%
\TextAt{21}{266.82}{469.561}{n}%
\TextAt{9}{275.777}{471.354}{\Character{61}}%
\NativeMathAt{292.603}{486.298}{7.97}{\infty}%
\TextAt{20}{288.203}{482.711}{\Character{88}}%
\TextAt{21}{288.633}{457.605}{\Character{23}}%
\TextAt{17}{293.316}{457.605}{\Character{61}}%
\TextAt{21}{299.903}{457.605}{n}%
\TextAt{13}{307.464}{471.354}{a}%
\TextAt{21}{313.609}{469.561}{\Character{23}}%
\TextAt{13}{318.791}{471.354}{p}%
\TextAt{21}{324.666}{476.29}{\Character{23}}%
\NativeMathAt{333.169}{471.354}{11.955}{\in}%
\TextAt{13}{344.46}{471.354}{p}%
\TextAt{21}{350.335}{476.29}{n}%
\TextAt{18}{355.971}{471.354}{Z}%
\TextAt{21}{363.941}{469.561}{p}%
\TextAt{13}{368.702}{471.354}{\Character{59}}%
\TextAt{8}{79.226}{440.676}{one}%
\TextAt{8}{100.815}{440.676}{has}%
\TextAt{13}{121.818}{440.676}{x}%
\NativeMathAt{132.011}{440.676}{11.955}{\equiv}%
\TextAt{13}{144.851}{440.676}{s}%
\TextAt{21}{150.365}{438.883}{n}%
\TextAt{8}{159.322}{440.676}{mod}%
\TextAt{13}{184.76}{440.676}{p}%
\TextAt{21}{190.635}{445.015}{n}%
\TextAt{8}{200.298}{440.676}{for}%
\TextAt{8}{218.312}{440.676}{all}%
\TextAt{13}{234.701}{440.676}{n}%
\TextAt{8}{241.688}{440.676}{,}%
\TextAt{8}{248.999}{440.676}{i.e.,}%
\TextAt{13}{271.266}{440.676}{x}%
\TextAt{8}{281.945}{440.676}{is}%
\TextAt{8}{293.838}{440.676}{mapped}%
\TextAt{8}{338.512}{440.676}{to}%
\TextAt{8}{352.952}{440.676}{the}%
\TextAt{8}{373.238}{440.676}{element}%
\TextAt{8}{416.609}{440.676}{of}%
\NativeMathAt{430.068}{440.676}{11.955}{\varprojlim}%
\TextAt{21}{446.327}{435.237}{n}%
\TextAt{18}{453.956}{440.676}{Z}%
\TextAt{9}{461.926}{440.676}{\Character{47}}%
\TextAt{13}{467.779}{440.676}{p}%
\TextAt{21}{473.654}{445.015}{n}%
\TextAt{18}{479.29}{440.676}{Z}%
\TextAt{8}{491.287}{440.676}{which}%
\TextAt{8}{524.903}{440.676}{is}%
\TextAt{8}{79.226}{423.63}{defined}%
\TextAt{8}{119.539}{423.63}{by}%
\TextAt{8}{135.798}{423.63}{the}%
\TextAt{8}{155.954}{423.63}{given}%
\TextAt{8}{186.524}{423.63}{sequence}%
\TextAt{9}{235.026}{423.63}{\Character{40}}%
\TextAt{13}{239.578}{423.63}{s}%
\TextAt{21}{245.092}{421.837}{n}%
\TextAt{9}{250.729}{423.63}{\Character{41}}%
\TextAt{21}{255.281}{421.837}{n}%
\NativeMathAt{260.419}{421.837}{7.97}{\in}%
\TextAt{25}{266.065}{421.837}{N}%
\TextAt{8}{272.319}{423.63}{.}%
\TextAt{8}{280.767}{423.63}{This}%
\TextAt{8}{307.487}{423.63}{shows}%
\TextAt{8}{341.105}{423.63}{surjectivity.}%
\Vector{q 0 G 0 g 0.399 w 452.334 -350.37 m 452.334 -342.3 l S Q}%
\Vector{q 0 G 0 g 0.399 w 452.534 -342.499 m 459.707 -342.499 l S Q}%
\Vector{q 0 G 0 g 0.399 w 452.534 -350.17 m 459.707 -350.17 l S Q}%
\Vector{q 0 G 0 g 0.399 w 459.906 -350.37 m 459.906 -342.3 l S Q}%
\TextAt{8}{97.159}{395.099}{We}%
\TextAt{8}{117.303}{395.099}{emphasize}%
\TextAt{8}{173.289}{395.099}{that}%
\TextAt{8}{198.658}{395.099}{the}%
\TextAt{8}{218.815}{395.099}{elements}%
\TextAt{8}{266.671}{395.099}{on}%
\TextAt{8}{282.93}{395.099}{the}%
\TextAt{8}{303.087}{395.099}{right}%
\TextAt{8}{331.385}{395.099}{hand}%
\TextAt{8}{360.651}{395.099}{side}%
\TextAt{8}{384.119}{395.099}{of}%
\TextAt{8}{397.449}{395.099}{the}%
\TextAt{8}{417.605}{395.099}{isomorphism}%
\TextAt{18}{263.947}{367.742}{Z}%
\TextAt{21}{271.917}{365.949}{p}%
\NativeMathAt{282.586}{374.519}{7.97}{\sim}%
\NativeArrowAt{279.999}{367.742}{11.955}{13.175}{\longrightarrow}%
\NativeMathAt{296.493}{367.742}{11.955}{\varprojlim}%
\TextAt{21}{302.053}{355.628}{n}%
\TextAt{18}{314.744}{367.742}{Z}%
\TextAt{9}{322.714}{367.742}{\Character{47}}%
\TextAt{13}{328.567}{367.742}{p}%
\TextAt{21}{334.442}{372.679}{n}%
\TextAt{18}{340.079}{367.742}{Z}%
\TextAt{8}{79.226}{334.709}{are}%
\TextAt{8}{98.737}{334.709}{given}%
\TextAt{8}{129.306}{334.709}{formally}%
\TextAt{8}{175.489}{334.709}{by}%
\TextAt{8}{191.748}{334.709}{sequences}%
\TextAt{8}{244.865}{334.709}{of}%
\TextAt{8}{258.195}{334.709}{partial}%
\TextAt{8}{295.926}{334.709}{sums}%
\TextAt{13}{225.193}{299.676}{s}%
\TextAt{21}{230.707}{297.882}{n}%
\TextAt{9}{239.665}{299.676}{\Character{61}}%
\TextAt{21}{252.745}{314.62}{n}%
\TextAt{3}{257.883}{314.62}{\Character{0}}%
\TextAt{17}{264.47}{314.62}{1}%
\TextAt{20}{252.09}{311.033}{\Character{88}}%
\TextAt{21}{252.972}{285.927}{\Character{23}}%
\TextAt{17}{257.656}{285.927}{\Character{61}0}%
\TextAt{13}{271.351}{299.676}{a}%
\TextAt{21}{277.496}{297.882}{\Character{23}}%
\TextAt{13}{282.678}{299.676}{p}%
\TextAt{21}{288.553}{304.612}{\Character{23}}%
\TextAt{13}{293.735}{299.676}{\Character{59}}%
\TextAt{13}{322.87}{299.676}{n}%
\TextAt{9}{333.198}{299.676}{\Character{61}}%
\TextAt{9}{345.608}{299.676}{1}%
\TextAt{13}{351.477}{299.676}{\Character{59}}%
\TextAt{9}{356.721}{299.676}{2}%
\TextAt{13}{362.574}{299.676}{\Character{59}}%
\TextAt{13}{367.798}{299.676}{\Character{58}}%
\TextAt{13}{373.035}{299.676}{\Character{58}}%
\TextAt{13}{378.259}{299.676}{\Character{58}}%
\TextAt{13}{383.496}{299.676}{\Character{58}}%
\TextAt{8}{79.226}{265.008}{On}%
\TextAt{8}{99.817}{265.008}{the}%
\TextAt{8}{121.054}{265.008}{left,}%
\TextAt{8}{146.136}{265.008}{however,}%
\TextAt{8}{195.618}{265.008}{these}%
\TextAt{8}{226.67}{265.008}{sequences}%
\TextAt{8}{280.867}{265.008}{converge}%
\TextAt{8}{329.756}{265.008}{with}%
\TextAt{8}{357.496}{265.008}{respect}%
\TextAt{8}{398.627}{265.008}{to}%
\TextAt{8}{414.017}{265.008}{the}%
\TextAt{8}{435.254}{265.008}{absolute}%
\TextAt{8}{482.577}{265.008}{value}%
\TextAt{8}{513.904}{265.008}{and}%
\TextAt{8}{79.226}{247.962}{yield}%
\TextAt{8}{107.512}{247.962}{the}%
\TextAt{8}{127.668}{247.962}{elements}%
\TextAt{8}{175.525}{247.962}{of}%
\TextAt{18}{188.855}{247.962}{Z}%
\TextAt{21}{196.825}{246.169}{p}%
\TextAt{8}{205.483}{247.962}{in}%
\TextAt{8}{219.136}{247.962}{the}%
\TextAt{8}{239.293}{247.962}{familiar}%
\TextAt{8}{282.547}{247.962}{way,}%
\TextAt{8}{308.561}{247.962}{as}%
\TextAt{8}{322.931}{247.962}{convergent}%
\TextAt{8}{381.476}{247.962}{infinite}%
\TextAt{8}{421.143}{247.962}{series}%
\TextAt{13}{272.35}{213.741}{x}%
\TextAt{9}{282.323}{213.741}{\Character{61}}%
\NativeMathAt{299.149}{228.685}{7.97}{\infty}%
\TextAt{20}{294.749}{225.099}{\Character{88}}%
\TextAt{21}{295.631}{199.993}{\Character{23}}%
\TextAt{17}{300.314}{199.993}{\Character{61}0}%
\TextAt{13}{314.01}{213.741}{a}%
\TextAt{21}{320.155}{211.948}{\Character{23}}%
\TextAt{13}{325.337}{213.741}{p}%
\TextAt{21}{331.212}{218.677}{\Character{23}}%
\TextAt{13}{336.394}{213.741}{\Character{58}}%
\TextAt{8}{79.226}{179.074}{Yet}%
\TextAt{8}{99.834}{179.074}{another,}%
\TextAt{8}{145.344}{179.074}{very}%
\TextAt{8}{170.196}{179.074}{elegant}%
\TextAt{8}{209.358}{179.074}{method}%
\TextAt{8}{251.114}{179.074}{to}%
\TextAt{8}{264.584}{179.074}{introduce}%
\TextAt{8}{315.773}{179.074}{the}%
\TextAt{13}{335.09}{179.074}{p}%
\TextAt{8}{340.965}{179.074}{-adic}%
\TextAt{8}{368.733}{179.074}{numbers}%
\TextAt{8}{415.104}{179.074}{comes}%
\TextAt{8}{448.791}{179.074}{about}%
\TextAt{8}{481.449}{179.074}{as}%
\TextAt{8}{494.979}{179.074}{follows.}%
\TextAt{8}{79.226}{162.028}{Let}%
\TextAt{18}{100.431}{162.028}{Z}%
\TextAt{9}{108.401}{162.028}{\Character{91}\Character{91}}%
\TextAt{13}{114.904}{162.028}{X}%
\TextAt{9}{125.56}{162.028}{\Character{93}\Character{93}}%
\TextAt{8}{136.196}{162.028}{denote}%
\TextAt{8}{174.15}{162.028}{the}%
\TextAt{8}{194.542}{162.028}{ring}%
\TextAt{8}{218.843}{162.028}{of}%
\TextAt{8}{232.408}{162.028}{all}%
\TextAt{8}{248.903}{162.028}{formal}%
\TextAt{8}{285.889}{162.028}{power}%
\TextAt{8}{320.268}{162.028}{series}%
\TextAt{20}{351.838}{170.994}{\Character{80}}%
\NativeMathAt{364.457}{167.839}{7.97}{\infty}%
\TextAt{21}{364.457}{158.541}{i}%
\TextAt{17}{367.341}{158.541}{\Character{61}0}%
\TextAt{13}{380.652}{162.028}{a}%
\TextAt{21}{386.797}{160.235}{i}%
\TextAt{13}{390.178}{162.028}{X}%
\TextAt{21}{400.833}{166.366}{i}%
\TextAt{8}{408.347}{162.028}{with}%
\TextAt{8}{435.243}{162.028}{integer}%
\TextAt{8}{474.177}{162.028}{coefficients.}%
\TextAt{8}{79.226}{144.982}{Then}%
\TextAt{8}{109.784}{144.982}{one}%
\TextAt{8}{131.243}{144.982}{has}%
\TextAt{8}{152.117}{144.982}{the}%
\TextAt{12}{79.226}{119.439}{Proposition}%
\TextAt{12}{151.937}{119.439}{(2.6).}%
\TextAt{8}{187.253}{119.439}{There}%
\TextAt{8}{221.062}{119.439}{is}%
\TextAt{8}{232.826}{119.439}{a}%
\TextAt{8}{242.582}{119.439}{canonical}%
\TextAt{8}{293.965}{119.439}{isomorphism}%
\TextAt{18}{251.171}{92.083}{Z}%
\TextAt{21}{259.141}{90.289}{p}%
\NativeMathAt{267.222}{92.083}{11.955}{\cong}%
\TextAt{18}{279.842}{92.083}{Z}%
\TextAt{9}{287.812}{92.083}{\Character{91}\Character{91}}%
\TextAt{13}{294.315}{92.083}{X}%
\TextAt{9}{304.97}{92.083}{\Character{93}\Character{93}\Character{47}\Character{40}}%
\TextAt{13}{321.879}{92.083}{X}%
\TextAt{16}{335.191}{92.083}{\Character{0}}%
\TextAt{13}{347.146}{92.083}{p}%
\TextAt{9}{353.021}{92.083}{\Character{41}}%
\TextAt{13}{357.573}{92.083}{\Character{58}}%
\TextAt{8}{297.211}{43.033}{104}%
\Destination{78.23}{809.815}{nav.page.104}%
\Destination{79.36}{681.855}{nav.print.114}%
\Destination{79.36}{775.155}{nav.number.2.2.5}%
\Destination{79.36}{133.455}{nav.number.2.2.6}%
}

% Page 117 (body page 105)
\Page{%
\TextAt{2}{79.226}{796.943}{II.}%
\TextAt{2}{94.073}{796.943}{The}%
\TextAt{2}{116.492}{796.943}{Theory}%
\TextAt{2}{154.401}{796.943}{of}%
\TextAt{2}{166.826}{796.943}{Valuations}%
\Vector{q 0 G 0 g 0.299 w 7.226 18.729 m 460.769 18.729 l S Q}%
\TextAt{19}{79.226}{761.124}{Proof.}%
\TextAt{8}{116}{761.124}{Consider}%
\TextAt{8}{163.71}{761.124}{the}%
\TextAt{8}{182.74}{761.124}{visibly}%
\TextAt{8}{218.745}{761.124}{surjective}%
\TextAt{8}{270.03}{761.124}{homomorphism}%
\TextAt{18}{351.573}{761.124}{Z}%
\TextAt{9}{359.543}{761.124}{\Character{91}\Character{91}}%
\TextAt{13}{366.046}{761.124}{X}%
\TextAt{9}{376.701}{761.124}{\Character{93}\Character{93}}%
\NativeArrowAt{386.525}{761.124}{11.955}{19.26}{\longrightarrow}%
\TextAt{18}{409.107}{761.124}{Z}%
\TextAt{21}{417.078}{759.33}{p}%
\TextAt{8}{424.609}{761.124}{which}%
\TextAt{8}{456.968}{761.124}{to}%
\TextAt{8}{470.152}{761.124}{every}%
\TextAt{8}{499.917}{761.124}{formal}%
\TextAt{8}{79.226}{744.078}{power}%
\TextAt{8}{113.717}{744.078}{series}%
\TextAt{20}{145.399}{753.044}{\Character{80}}%
\NativeMathAt{158.018}{749.889}{7.97}{\infty}%
\TextAt{21}{158.018}{740.591}{\Character{23}}%
\TextAt{17}{162.702}{740.591}{\Character{61}0}%
\TextAt{13}{176.013}{744.078}{a}%
\TextAt{21}{182.158}{742.284}{\Character{23}}%
\TextAt{13}{187.34}{744.078}{X}%
\TextAt{21}{197.995}{748.416}{\Character{23}}%
\TextAt{8}{207.422}{744.078}{associates}%
\TextAt{8}{261.615}{744.078}{the}%
\TextAt{8}{282.118}{744.078}{convergent}%
\TextAt{8}{341.009}{744.078}{series}%
\TextAt{20}{372.691}{753.044}{\Character{80}}%
\NativeMathAt{385.31}{749.889}{7.97}{\infty}%
\TextAt{21}{385.31}{740.591}{\Character{23}}%
\TextAt{17}{389.994}{740.591}{\Character{61}0}%
\TextAt{13}{403.306}{744.078}{a}%
\TextAt{21}{409.45}{742.284}{\Character{23}}%
\TextAt{13}{414.632}{744.078}{p}%
\TextAt{21}{420.508}{748.416}{\Character{23}}%
\TextAt{8}{425.69}{744.078}{.}%
\TextAt{8}{435.179}{744.078}{The}%
\TextAt{8}{459.58}{744.078}{principal}%
\TextAt{8}{508.704}{744.078}{ideal}%
\TextAt{9}{79.226}{727.032}{\Character{40}}%
\TextAt{13}{83.778}{727.032}{X}%
\TextAt{16}{97.746}{727.032}{\Character{0}}%
\TextAt{13}{110.357}{727.032}{p}%
\TextAt{9}{116.232}{727.032}{\Character{41}}%
\TextAt{8}{125.643}{727.032}{clearly}%
\TextAt{8}{164.001}{727.032}{belongs}%
\TextAt{8}{206.973}{727.032}{to}%
\TextAt{8}{222.246}{727.032}{the}%
\TextAt{8}{243.364}{727.032}{kernel}%
\TextAt{8}{278.792}{727.032}{of}%
\TextAt{8}{293.084}{727.032}{this}%
\TextAt{8}{316.869}{727.032}{mapping.}%
\TextAt{8}{372.437}{727.032}{In}%
\TextAt{8}{388.02}{727.032}{order}%
\TextAt{8}{419.551}{727.032}{to}%
\TextAt{8}{434.823}{727.032}{show}%
\TextAt{8}{464.788}{727.032}{that}%
\TextAt{8}{491.119}{727.032}{it}%
\TextAt{8}{503.785}{727.032}{is}%
\TextAt{8}{516.51}{727.032}{the}%
\TextAt{8}{79.226}{709.986}{whole}%
\TextAt{8}{112.535}{709.986}{kernel,}%
\TextAt{8}{150.436}{709.986}{let}%
\TextAt{13}{167.486}{709.986}{f}%
\TextAt{9}{174.533}{709.986}{\Character{40}}%
\TextAt{13}{179.085}{709.986}{X}%
\TextAt{9}{189.74}{709.986}{\Character{41}}%
\TextAt{9}{197.846}{709.986}{\Character{61}}%
\TextAt{20}{210.535}{718.952}{\Character{80}}%
\NativeMathAt{223.154}{715.797}{7.97}{\infty}%
\TextAt{21}{223.154}{706.499}{\Character{23}}%
\TextAt{17}{227.838}{706.499}{\Character{61}0}%
\TextAt{13}{241.15}{709.986}{a}%
\TextAt{21}{247.295}{708.193}{\Character{23}}%
\TextAt{13}{252.477}{709.986}{X}%
\TextAt{21}{263.132}{714.324}{\Character{23}}%
\TextAt{8}{272.357}{709.986}{be}%
\TextAt{8}{288.427}{709.986}{a}%
\TextAt{8}{298.328}{709.986}{power}%
\TextAt{8}{332.617}{709.986}{series}%
\TextAt{8}{364.097}{709.986}{such}%
\TextAt{8}{390.64}{709.986}{that}%
\TextAt{13}{416.155}{709.986}{f}%
\TextAt{9}{423.201}{709.986}{\Character{40}}%
\TextAt{13}{427.753}{709.986}{p}%
\TextAt{9}{433.628}{709.986}{\Character{41}}%
\TextAt{9}{441.734}{709.986}{\Character{61}}%
\TextAt{20}{454.423}{718.952}{\Character{80}}%
\NativeMathAt{467.043}{715.797}{7.97}{\infty}%
\TextAt{21}{467.043}{706.499}{\Character{23}}%
\TextAt{17}{471.727}{706.499}{\Character{61}0}%
\TextAt{13}{485.038}{709.986}{a}%
\TextAt{21}{491.183}{708.193}{\Character{23}}%
\TextAt{13}{496.365}{709.986}{p}%
\TextAt{21}{502.24}{714.324}{\Character{23}}%
\TextAt{9}{510.991}{709.986}{\Character{61}}%
\TextAt{9}{523.652}{709.986}{0}%
\TextAt{8}{529.518}{709.986}{.}%
\TextAt{8}{79.226}{692.94}{Since}%
\TextAt{18}{109.784}{692.94}{Z}%
\TextAt{21}{117.754}{691.147}{p}%
\TextAt{9}{122.515}{692.94}{\Character{47}}%
\TextAt{13}{128.368}{692.94}{p}%
\TextAt{21}{134.243}{697.278}{n}%
\TextAt{18}{139.879}{692.94}{Z}%
\TextAt{21}{147.849}{691.147}{p}%
\NativeMathAt{155.931}{692.94}{11.955}{\cong}%
\TextAt{18}{168.55}{692.94}{Z}%
\TextAt{9}{176.52}{692.94}{\Character{47}}%
\TextAt{13}{182.373}{692.94}{p}%
\TextAt{21}{188.249}{697.278}{n}%
\TextAt{18}{193.885}{692.94}{Z}%
\TextAt{8}{201.855}{692.94}{,}%
\TextAt{8}{209.004}{692.94}{this}%
\TextAt{8}{231.827}{692.94}{means}%
\TextAt{8}{267.656}{692.94}{that}%
\TextAt{13}{200.882}{666.049}{a}%
\TextAt{17}{207.027}{664.256}{0}%
\TextAt{9}{214.416}{666.049}{\Character{43}}%
\TextAt{13}{226.177}{666.049}{a}%
\TextAt{17}{232.322}{664.256}{1}%
\TextAt{13}{237.054}{666.049}{p}%
\TextAt{9}{245.586}{666.049}{\Character{43}}%
\TextAt{16}{257.348}{666.049}{\Character{1}}%
\TextAt{16}{262.644}{666.049}{\Character{1}}%
\TextAt{16}{267.952}{666.049}{\Character{1}}%
\TextAt{9}{273.952}{666.049}{\Character{43}}%
\TextAt{13}{285.713}{666.049}{a}%
\TextAt{21}{291.858}{664.256}{n}%
\TextAt{3}{296.996}{664.256}{\Character{0}}%
\TextAt{17}{303.583}{664.256}{1}%
\TextAt{13}{308.315}{666.049}{p}%
\TextAt{21}{314.19}{670.985}{n}%
\TextAt{3}{319.329}{670.985}{\Character{0}}%
\TextAt{17}{325.915}{670.985}{1}%
\NativeMathAt{333.968}{666.049}{11.955}{\equiv}%
\TextAt{9}{346.588}{666.049}{0}%
\TextAt{9}{364.377}{666.049}{\Character{40}}%
\TextAt{8}{368.948}{666.049}{mod}%
\TextAt{13}{395.05}{666.049}{p}%
\TextAt{21}{400.925}{670.985}{n}%
\TextAt{9}{406.561}{666.049}{\Character{41}}%
\TextAt{8}{79.226}{639.159}{for}%
\TextAt{8}{97.111}{639.159}{all}%
\TextAt{13}{113.37}{639.159}{n}%
\TextAt{8}{120.358}{639.159}{.}%
\TextAt{8}{128.806}{639.159}{We}%
\TextAt{8}{148.951}{639.159}{put,}%
\TextAt{8}{173.662}{639.159}{for}%
\TextAt{13}{191.547}{639.159}{n}%
\NativeMathAt{201.855}{639.159}{11.955}{\geq}%
\TextAt{9}{214.475}{639.159}{1}%
\TextAt{8}{220.328}{639.159}{,}%
\TextAt{13}{204.746}{610.206}{b}%
\TextAt{21}{209.723}{608.413}{n}%
\TextAt{3}{214.862}{608.413}{\Character{0}}%
\TextAt{17}{221.448}{608.413}{1}%
\TextAt{9}{229.501}{610.206}{\Character{61}}%
\TextAt{16}{241.927}{610.206}{\Character{0}}%
\TextAt{9}{255.25}{618.294}{1}%
\Vector{q 0 G 0 g 0.478 w 180.421 -160.805 m 191.932 -160.805 l S Q}%
\TextAt{13}{252.421}{602.006}{p}%
\TextAt{21}{258.296}{605.459}{n}%
\TextAt{9}{265.128}{610.206}{\Character{40}}%
\TextAt{13}{269.68}{610.206}{a}%
\TextAt{17}{275.825}{608.413}{0}%
\TextAt{9}{283.214}{610.206}{\Character{43}}%
\TextAt{13}{294.975}{610.206}{a}%
\TextAt{17}{301.12}{608.413}{1}%
\TextAt{13}{305.853}{610.206}{p}%
\TextAt{9}{314.384}{610.206}{\Character{43}}%
\TextAt{16}{326.146}{610.206}{\Character{1}}%
\TextAt{16}{331.442}{610.206}{\Character{1}}%
\TextAt{16}{336.75}{610.206}{\Character{1}}%
\TextAt{9}{342.75}{610.206}{\Character{43}}%
\TextAt{13}{354.511}{610.206}{a}%
\TextAt{21}{360.656}{608.413}{n}%
\TextAt{3}{365.794}{608.413}{\Character{0}}%
\TextAt{17}{372.381}{608.413}{1}%
\TextAt{13}{377.113}{610.206}{p}%
\TextAt{21}{382.988}{615.143}{n}%
\TextAt{3}{388.127}{615.143}{\Character{0}}%
\TextAt{17}{394.713}{615.143}{1}%
\TextAt{9}{399.445}{610.206}{\Character{41}}%
\TextAt{13}{403.998}{610.206}{\Character{58}}%
\TextAt{8}{79.226}{580.543}{Then}%
\TextAt{8}{109.784}{580.543}{we}%
\TextAt{8}{127.011}{580.543}{obtain}%
\TextAt{8}{163.438}{580.543}{successively}%
\TextAt{13}{249.052}{565.176}{a}%
\TextAt{17}{255.197}{563.383}{0}%
\TextAt{9}{263.25}{565.176}{\Character{61}}%
\TextAt{16}{275.676}{565.176}{\Character{0}}%
\TextAt{13}{284.974}{565.176}{pb}%
\TextAt{17}{295.826}{563.383}{0}%
\TextAt{13}{300.559}{565.176}{\Character{59}}%
\TextAt{13}{249.052}{545.141}{a}%
\TextAt{17}{255.197}{543.348}{1}%
\TextAt{9}{263.25}{545.141}{\Character{61}}%
\TextAt{13}{275.676}{545.141}{b}%
\TextAt{17}{280.653}{543.348}{0}%
\TextAt{16}{288.042}{545.141}{\Character{0}}%
\TextAt{13}{299.997}{545.141}{pb}%
\TextAt{17}{310.849}{543.348}{1}%
\TextAt{13}{315.581}{545.141}{\Character{59}}%
\TextAt{13}{249.052}{525.107}{a}%
\TextAt{17}{255.197}{523.313}{2}%
\TextAt{9}{263.25}{525.107}{\Character{61}}%
\TextAt{13}{275.676}{525.107}{b}%
\TextAt{17}{280.653}{523.313}{1}%
\TextAt{16}{288.042}{525.107}{\Character{0}}%
\TextAt{13}{299.997}{525.107}{pb}%
\TextAt{17}{310.849}{523.313}{2}%
\TextAt{13}{315.581}{525.107}{\Character{59}}%
\TextAt{8}{344.736}{525.107}{etc.}%
\TextAt{8}{79.226}{505.129}{But}%
\TextAt{8}{102.467}{505.129}{this}%
\TextAt{8}{125.289}{505.129}{amounts}%
\TextAt{8}{172.512}{505.129}{to}%
\TextAt{8}{186.823}{505.129}{the}%
\TextAt{8}{206.979}{505.129}{equality}%
\TextAt{13}{156.703}{478.238}{a}%
\TextAt{17}{162.848}{476.445}{0}%
\TextAt{9}{170.237}{478.238}{\Character{43}}%
\TextAt{13}{181.998}{478.238}{a}%
\TextAt{17}{188.143}{476.445}{1}%
\TextAt{13}{192.875}{478.238}{X}%
\TextAt{9}{206.187}{478.238}{\Character{43}}%
\TextAt{13}{217.948}{478.238}{a}%
\TextAt{17}{224.093}{476.445}{2}%
\TextAt{13}{228.826}{478.238}{X}%
\TextAt{17}{239.481}{483.174}{2}%
\TextAt{9}{246.87}{478.238}{\Character{43}}%
\TextAt{16}{258.631}{478.238}{\Character{1}}%
\TextAt{16}{263.927}{478.238}{\Character{1}}%
\TextAt{16}{269.235}{478.238}{\Character{1}}%
\TextAt{9}{275.9}{478.238}{\Character{61}}%
\TextAt{9}{288.31}{478.238}{\Character{40}}%
\TextAt{13}{292.877}{478.238}{X}%
\TextAt{16}{306.189}{478.238}{\Character{0}}%
\TextAt{13}{318.144}{478.238}{p}%
\TextAt{9}{324.019}{478.238}{\Character{41}\Character{40}}%
\TextAt{13}{333.124}{478.238}{b}%
\TextAt{17}{338.101}{476.445}{0}%
\TextAt{9}{345.49}{478.238}{\Character{43}}%
\TextAt{13}{357.251}{478.238}{b}%
\TextAt{17}{362.229}{476.445}{1}%
\TextAt{13}{366.961}{478.238}{X}%
\TextAt{9}{380.273}{478.238}{\Character{43}}%
\TextAt{13}{392.034}{478.238}{b}%
\TextAt{17}{397.011}{476.445}{2}%
\TextAt{13}{401.743}{478.238}{X}%
\TextAt{17}{412.399}{483.174}{2}%
\TextAt{9}{419.788}{478.238}{\Character{43}}%
\TextAt{16}{431.549}{478.238}{\Character{1}}%
\TextAt{16}{436.845}{478.238}{\Character{1}}%
\TextAt{16}{442.153}{478.238}{\Character{1}}%
\TextAt{9}{447.489}{478.238}{\Character{41}}%
\TextAt{13}{452.041}{478.238}{\Character{59}}%
\TextAt{8}{79.226}{451.347}{i.e.,}%
\TextAt{13}{101.331}{451.347}{f}%
\TextAt{9}{108.378}{451.347}{\Character{40}}%
\TextAt{13}{112.93}{451.347}{X}%
\TextAt{9}{123.585}{451.347}{\Character{41}}%
\TextAt{8}{132.035}{451.347}{belongs}%
\TextAt{8}{174.045}{451.347}{to}%
\TextAt{8}{188.356}{451.347}{the}%
\TextAt{8}{208.512}{451.347}{principal}%
\TextAt{8}{257.289}{451.347}{ideal}%
\TextAt{9}{285.252}{451.347}{\Character{40}}%
\TextAt{13}{289.805}{451.347}{X}%
\TextAt{16}{303.116}{451.347}{\Character{0}}%
\TextAt{13}{315.072}{451.347}{p}%
\TextAt{9}{320.947}{451.347}{\Character{41}}%
\TextAt{8}{325.499}{451.347}{.}%
\Vector{q 0 G 0 g 0.399 w 452.334 -322.653 m 452.334 -314.583 l S Q}%
\Vector{q 0 G 0 g 0.399 w 452.534 -314.782 m 459.707 -314.782 l S Q}%
\Vector{q 0 G 0 g 0.399 w 452.534 -322.453 m 459.707 -322.453 l S Q}%
\Vector{q 0 G 0 g 0.399 w 459.906 -322.653 m 459.906 -314.583 l S Q}%
\TextAt{12}{79.226}{417.575}{Exercise}%
\TextAt{12}{132.427}{417.575}{II.2.1.}%
\TextAt{16}{171.401}{417.575}{\Character{106}}%
\TextAt{13}{174.721}{417.575}{x}%
\TextAt{16}{184.03}{417.575}{\Character{0}}%
\TextAt{13}{195.985}{417.575}{y}%
\TextAt{16}{202.122}{417.575}{\Character{106}}%
\TextAt{21}{205.443}{415.782}{p}%
\NativeMathAt{213.525}{417.575}{11.955}{\geq}%
\TextAt{20}{226.144}{427.737}{\Character{12}}%
\TextAt{20}{226.144}{420.564}{\Character{12}}%
\TextAt{16}{230.129}{417.575}{\Character{106}}%
\TextAt{13}{233.45}{417.575}{x}%
\TextAt{16}{240.102}{417.575}{\Character{106}}%
\TextAt{21}{243.423}{415.782}{p}%
\TextAt{16}{250.84}{417.575}{\Character{0}}%
\TextAt{16}{262.783}{417.575}{\Character{106}}%
\TextAt{13}{266.117}{417.575}{y}%
\TextAt{16}{272.253}{417.575}{\Character{106}}%
\TextAt{21}{275.574}{415.782}{p}%
\TextAt{20}{280.335}{427.737}{\Character{12}}%
\TextAt{20}{280.335}{420.564}{\Character{12}}%
\TextAt{8}{284.32}{417.575}{.}%
\TextAt{12}{79.226}{386.791}{Exercise}%
\TextAt{12}{133.446}{386.791}{II.2.2.}%
\TextAt{8}{173.307}{386.791}{Let}%
\TextAt{13}{195.162}{386.791}{n}%
\TextAt{8}{206.934}{386.791}{be}%
\TextAt{8}{223.745}{386.791}{a}%
\TextAt{8}{234.386}{386.791}{natural}%
\TextAt{8}{276.255}{386.791}{number,}%
\TextAt{13}{323.211}{386.791}{n}%
\TextAt{9}{335.03}{386.791}{\Character{61}}%
\TextAt{13}{348.966}{386.791}{a}%
\TextAt{17}{355.111}{384.998}{0}%
\TextAt{9}{363.104}{386.791}{\Character{43}}%
\TextAt{13}{375.47}{386.791}{a}%
\TextAt{17}{381.615}{384.998}{1}%
\TextAt{13}{386.347}{386.791}{p}%
\TextAt{9}{395.483}{386.791}{\Character{43}}%
\TextAt{16}{407.849}{386.791}{\Character{1}}%
\TextAt{16}{413.145}{386.791}{\Character{1}}%
\TextAt{16}{418.453}{386.791}{\Character{1}}%
\TextAt{9}{425.057}{386.791}{\Character{43}}%
\TextAt{13}{437.423}{386.791}{a}%
\TextAt{21}{443.568}{384.998}{r}%
\TextAt{3}{447.625}{384.998}{\Character{0}}%
\TextAt{17}{454.211}{384.998}{1}%
\TextAt{13}{458.943}{386.791}{p}%
\TextAt{21}{464.819}{391.129}{r}%
\TextAt{3}{468.875}{391.129}{\Character{0}}%
\TextAt{17}{475.462}{391.129}{1}%
\TextAt{8}{484.978}{386.791}{its}%
\TextAt{13}{502.183}{386.791}{p}%
\TextAt{8}{508.058}{386.791}{-adic}%
\TextAt{8}{79.226}{369.745}{expansion,}%
\TextAt{8}{136.85}{369.745}{with}%
\TextAt{9}{163.51}{369.745}{0}%
\NativeMathAt{172.684}{369.745}{11.955}{\leq}%
\TextAt{13}{185.303}{369.745}{a}%
\TextAt{21}{191.448}{367.952}{i}%
\TextAt{13}{198.15}{369.745}{\Character{60}}%
\TextAt{13}{210.56}{369.745}{p}%
\TextAt{8}{216.451}{369.745}{,}%
\TextAt{8}{223.6}{369.745}{and}%
\TextAt{13}{246.363}{369.745}{s}%
\TextAt{9}{255.198}{369.745}{\Character{61}}%
\TextAt{13}{267.623}{369.745}{a}%
\TextAt{17}{273.768}{367.952}{0}%
\TextAt{9}{281.157}{369.745}{\Character{43}}%
\TextAt{13}{292.918}{369.745}{a}%
\TextAt{17}{299.063}{367.952}{1}%
\TextAt{9}{306.452}{369.745}{\Character{43}}%
\TextAt{16}{318.214}{369.745}{\Character{1}}%
\TextAt{16}{323.51}{369.745}{\Character{1}}%
\TextAt{16}{328.818}{369.745}{\Character{1}}%
\TextAt{9}{334.818}{369.745}{\Character{43}}%
\TextAt{13}{346.579}{369.745}{a}%
\TextAt{21}{352.724}{367.952}{r}%
\TextAt{3}{356.781}{367.952}{\Character{0}}%
\TextAt{17}{363.367}{367.952}{1}%
\TextAt{8}{368.1}{369.745}{.}%
\TextAt{8}{376.548}{369.745}{Show}%
\TextAt{8}{407.44}{369.745}{that}%
\TextAt{13}{432.809}{369.745}{v}%
\TextAt{21}{438.468}{367.952}{p}%
\TextAt{9}{443.229}{369.745}{\Character{40}}%
\TextAt{13}{447.781}{369.745}{n}%
\TextAt{9}{454.769}{369.745}{\Character{33}\Character{41}}%
\TextAt{9}{465.863}{369.745}{\Character{61}}%
\TextAt{21}{479.515}{374.452}{n}%
\TextAt{3}{484.653}{374.452}{\Character{0}}%
\TextAt{21}{491.24}{374.452}{s}%
\Vector{q 0 G 0 g 0.478 w 407.515 -401.266 m 423.155 -401.266 l S Q}%
\TextAt{21}{479.793}{365.622}{p}%
\TextAt{3}{484.056}{365.622}{\Character{0}}%
\TextAt{17}{490.643}{365.622}{1}%
\TextAt{8}{496.351}{369.745}{.}%
\TextAt{12}{79.226}{338.961}{Exercise}%
\TextAt{12}{132.427}{338.961}{II.2.3.}%
\TextAt{8}{171.401}{338.961}{The}%
\TextAt{8}{195.454}{338.961}{sequence}%
\TextAt{9}{243.956}{338.961}{1}%
\TextAt{13}{249.809}{338.961}{\Character{59}}%
\TextAt{17}{258.366}{343.668}{1}%
\Vector{q 0 G 0 g 0.478 w 184.249 -432.05 m 192.717 -432.05 l S Q}%
\TextAt{17}{256.249}{334.838}{10}%
\TextAt{13}{265.913}{338.961}{\Character{59}}%
\TextAt{17}{276.545}{343.668}{1}%
\Vector{q 0 G 0 g 0.478 w 200.353 -432.05 m 212.972 -432.05 l S Q}%
\TextAt{17}{272.353}{334.838}{10}%
\TextAt{24}{280.821}{337.108}{2}%
\TextAt{13}{286.168}{338.961}{\Character{59}}%
\TextAt{17}{296.8}{343.668}{1}%
\Vector{q 0 G 0 g 0.478 w 220.607 -432.05 m 233.227 -432.05 l S Q}%
\TextAt{17}{292.607}{334.838}{10}%
\TextAt{24}{301.076}{337.108}{3}%
\TextAt{13}{306.422}{338.961}{\Character{59}}%
\TextAt{13}{311.646}{338.961}{\Character{58}}%
\TextAt{13}{316.883}{338.961}{\Character{58}}%
\TextAt{13}{322.119}{338.961}{\Character{58}}%
\TextAt{8}{329.304}{338.961}{does}%
\TextAt{8}{355.701}{338.961}{not}%
\TextAt{8}{376.515}{338.961}{converge}%
\TextAt{8}{424.323}{338.961}{in}%
\TextAt{18}{437.976}{338.961}{Q}%
\TextAt{21}{447.275}{337.168}{p}%
\TextAt{8}{452.036}{338.961}{,}%
\TextAt{8}{459.185}{338.961}{for}%
\TextAt{8}{477.07}{338.961}{any}%
\TextAt{13}{499.187}{338.961}{p}%
\TextAt{8}{505.062}{338.961}{.}%
\TextAt{12}{79.226}{308.177}{Exercise}%
\TextAt{12}{132.474}{308.177}{II.2.4.}%
\TextAt{8}{171.488}{308.177}{Let}%
\TextAt{13}{192.498}{308.177}{\Character{34}}%
\NativeMathAt{201.367}{308.177}{11.955}{\in}%
\TextAt{9}{212.727}{308.177}{1}%
\TextAt{9}{221.251}{308.177}{\Character{43}}%
\TextAt{13}{233.054}{308.177}{p}%
\TextAt{18}{238.929}{308.177}{Z}%
\TextAt{21}{246.899}{306.384}{p}%
\TextAt{8}{251.66}{308.177}{,}%
\TextAt{8}{258.86}{308.177}{and}%
\TextAt{8}{281.664}{308.177}{let}%
\TextAt{13}{298.609}{308.177}{\Character{11}}%
\TextAt{9}{309.521}{308.177}{\Character{61}}%
\TextAt{13}{322.017}{308.177}{a}%
\TextAt{17}{328.161}{306.384}{0}%
\TextAt{9}{335.578}{308.177}{\Character{43}}%
\TextAt{13}{347.367}{308.177}{a}%
\TextAt{17}{353.512}{306.384}{1}%
\TextAt{13}{358.245}{308.177}{p}%
\TextAt{9}{366.804}{308.177}{\Character{43}}%
\TextAt{13}{378.593}{308.177}{a}%
\TextAt{17}{384.738}{306.384}{2}%
\TextAt{13}{389.471}{308.177}{p}%
\TextAt{17}{395.346}{312.516}{2}%
\TextAt{9}{402.763}{308.177}{\Character{43}}%
\TextAt{16}{414.552}{308.177}{\Character{1}}%
\TextAt{16}{419.848}{308.177}{\Character{1}}%
\TextAt{16}{425.156}{308.177}{\Character{1}}%
\TextAt{8}{434.43}{308.177}{be}%
\TextAt{8}{450.395}{308.177}{a}%
\TextAt{13}{460.192}{308.177}{p}%
\TextAt{8}{466.067}{308.177}{-adic}%
\TextAt{8}{494.716}{308.177}{integer,}%
\TextAt{8}{79.226}{291.131}{and}%
\TextAt{8}{102.982}{291.131}{write}%
\TextAt{13}{133.886}{291.131}{s}%
\TextAt{21}{139.4}{289.338}{n}%
\TextAt{9}{150.049}{291.131}{\Character{61}}%
\TextAt{13}{164.167}{291.131}{a}%
\TextAt{17}{170.312}{289.338}{0}%
\TextAt{9}{178.378}{291.131}{\Character{43}}%
\TextAt{13}{190.816}{291.131}{a}%
\TextAt{17}{196.961}{289.338}{1}%
\TextAt{13}{201.693}{291.131}{p}%
\TextAt{9}{210.901}{291.131}{\Character{43}}%
\TextAt{16}{223.34}{291.131}{\Character{1}}%
\TextAt{16}{228.636}{291.131}{\Character{1}}%
\TextAt{16}{233.944}{291.131}{\Character{1}}%
\TextAt{9}{240.621}{291.131}{\Character{43}}%
\TextAt{13}{253.059}{291.131}{a}%
\TextAt{21}{259.204}{289.338}{n}%
\TextAt{3}{264.342}{289.338}{\Character{0}}%
\TextAt{17}{270.928}{289.338}{1}%
\TextAt{13}{275.661}{291.131}{p}%
\TextAt{21}{281.536}{295.47}{n}%
\TextAt{3}{286.674}{295.47}{\Character{0}}%
\TextAt{17}{293.261}{295.47}{1}%
\TextAt{8}{297.993}{291.131}{.}%
\TextAt{8}{309.42}{291.131}{Show}%
\TextAt{8}{341.305}{291.131}{that}%
\TextAt{8}{367.667}{291.131}{the}%
\TextAt{8}{388.816}{291.131}{sequence}%
\TextAt{13}{438.311}{291.131}{\Character{34}}%
\TextAt{21}{443.789}{295.47}{s}%
\TextAt{22}{447.705}{294.473}{n}%
\TextAt{8}{458.192}{291.131}{converges}%
\TextAt{8}{511.608}{291.131}{to}%
\TextAt{8}{526.911}{291.131}{a}%
\TextAt{8}{79.226}{274.085}{number}%
\TextAt{13}{121.445}{274.085}{\Character{34}}%
\TextAt{21}{126.923}{278.424}{\Character{11}}%
\TextAt{8}{136.379}{274.085}{in}%
\TextAt{9}{149.654}{274.085}{1}%
\TextAt{9}{157.377}{274.085}{\Character{43}}%
\TextAt{13}{168.38}{274.085}{p}%
\TextAt{18}{174.255}{274.085}{Z}%
\TextAt{21}{182.225}{272.292}{p}%
\TextAt{8}{186.986}{274.085}{.}%
\TextAt{8}{195.308}{274.085}{Show}%
\TextAt{8}{225.823}{274.085}{furthermore}%
\TextAt{8}{290.158}{274.085}{that}%
\TextAt{9}{315.149}{274.085}{1}%
\TextAt{9}{322.872}{274.085}{\Character{43}}%
\TextAt{13}{333.875}{274.085}{p}%
\TextAt{18}{339.75}{274.085}{Z}%
\TextAt{21}{347.72}{272.292}{p}%
\TextAt{8}{356.001}{274.085}{is}%
\TextAt{8}{367.387}{274.085}{thus}%
\TextAt{8}{392.76}{274.085}{turned}%
\TextAt{8}{430.101}{274.085}{into}%
\TextAt{8}{453.466}{274.085}{a}%
\TextAt{8}{462.844}{274.085}{multiplicative}%
\TextAt{18}{79.226}{257.039}{Z}%
\TextAt{21}{87.196}{255.246}{p}%
\TextAt{8}{91.957}{257.039}{-module.}%
\TextAt{12}{79.226}{226.255}{Exercise}%
\TextAt{12}{132.427}{226.255}{II.2.5.}%
\TextAt{8}{171.401}{226.255}{For}%
\TextAt{8}{192.37}{226.255}{every}%
\TextAt{13}{223.262}{226.255}{a}%
\NativeMathAt{232.728}{226.255}{11.955}{\in}%
\TextAt{18}{244.019}{226.255}{Z}%
\TextAt{8}{251.989}{226.255}{,}%
\TextAt{9}{259.138}{226.255}{\Character{40}}%
\TextAt{13}{263.69}{226.255}{a\Character{59}}%
\TextAt{13}{275.059}{226.255}{p}%
\TextAt{9}{280.955}{226.255}{\Character{41}}%
\TextAt{9}{288.81}{226.255}{\Character{61}}%
\TextAt{9}{301.231}{226.255}{1}%
\TextAt{8}{307.106}{226.255}{,}%
\TextAt{8}{314.256}{226.255}{the}%
\TextAt{8}{334.412}{226.255}{sequence}%
\TextAt{16}{382.914}{226.255}{\Character{102}}%
\TextAt{13}{388.892}{226.255}{a}%
\TextAt{21}{395.037}{230.594}{p}%
\TextAt{22}{399.299}{233.407}{n}%
\TextAt{16}{404.896}{226.255}{\Character{103}}%
\TextAt{21}{410.874}{224.462}{n}%
\NativeMathAt{416.012}{224.462}{7.97}{\in}%
\TextAt{25}{421.658}{224.462}{N}%
\TextAt{8}{431.81}{226.255}{converges}%
\TextAt{8}{484.233}{226.255}{in}%
\TextAt{18}{497.886}{226.255}{Q}%
\TextAt{21}{507.184}{224.462}{p}%
\TextAt{8}{511.945}{226.255}{.}%
\TextAt{12}{79.226}{195.472}{Exercise}%
\TextAt{12}{132.427}{195.472}{II.2.6.}%
\TextAt{8}{171.401}{195.472}{The}%
\TextAt{8}{195.454}{195.472}{fields}%
\TextAt{18}{225.426}{195.472}{Q}%
\TextAt{21}{234.724}{193.678}{p}%
\TextAt{8}{243.383}{195.472}{and}%
\TextAt{18}{266.145}{195.472}{Q}%
\TextAt{21}{275.444}{193.678}{q}%
\TextAt{8}{283.906}{195.472}{are}%
\TextAt{8}{303.417}{195.472}{not}%
\TextAt{8}{324.231}{195.472}{isomorphic,}%
\TextAt{8}{386.733}{195.472}{unless}%
\TextAt{13}{421.319}{195.472}{p}%
\TextAt{9}{430.515}{195.472}{\Character{61}}%
\TextAt{13}{442.94}{195.472}{q}%
\TextAt{8}{448.559}{195.472}{.}%
\TextAt{12}{79.226}{164.688}{Exercise}%
\TextAt{12}{132.427}{164.688}{II.2.7.}%
\TextAt{8}{171.401}{164.688}{The}%
\TextAt{8}{195.454}{164.688}{algebraic}%
\TextAt{8}{244.889}{164.688}{closure}%
\TextAt{8}{283.97}{164.688}{of}%
\TextAt{18}{297.3}{164.688}{Q}%
\TextAt{21}{306.599}{162.894}{p}%
\TextAt{8}{315.257}{164.688}{has}%
\TextAt{8}{336.131}{164.688}{infinite}%
\TextAt{8}{375.798}{164.688}{degree.}%
\TextAt{12}{79.226}{133.904}{Exercise}%
\TextAt{12}{132.558}{133.904}{II.2.8.}%
\TextAt{8}{171.645}{133.904}{In}%
\TextAt{8}{186.38}{133.904}{the}%
\TextAt{8}{206.65}{133.904}{ring}%
\TextAt{18}{230.83}{133.904}{Z}%
\TextAt{21}{238.8}{132.111}{p}%
\TextAt{9}{243.561}{133.904}{\Character{91}\Character{91}}%
\TextAt{13}{250.064}{133.904}{X}%
\TextAt{9}{260.719}{133.904}{\Character{93}\Character{93}}%
\TextAt{8}{271.234}{133.904}{of}%
\TextAt{8}{284.678}{133.904}{formal}%
\TextAt{8}{321.542}{133.904}{power}%
\TextAt{8}{355.799}{133.904}{series}%
\TextAt{20}{387.248}{142.87}{\Character{80}}%
\NativeMathAt{399.867}{139.715}{7.97}{\infty}%
\TextAt{21}{399.867}{130.417}{\Character{23}}%
\TextAt{17}{404.551}{130.417}{\Character{61}0}%
\TextAt{13}{417.862}{133.904}{a}%
\TextAt{21}{424.007}{132.111}{\Character{23}}%
\TextAt{13}{429.189}{133.904}{X}%
\TextAt{21}{439.844}{138.242}{\Character{23}}%
\TextAt{8}{449.038}{133.904}{over}%
\TextAt{18}{474.197}{133.904}{Z}%
\TextAt{21}{482.168}{132.111}{p}%
\TextAt{8}{486.928}{133.904}{,}%
\TextAt{8}{494.22}{133.904}{one}%
\TextAt{8}{515.793}{133.904}{has}%
\TextAt{8}{79.226}{116.858}{the}%
\TextAt{8}{98.336}{116.858}{following}%
\TextAt{12}{146.724}{116.858}{division}%
\TextAt{12}{195.313}{116.858}{with}%
\TextAt{12}{224.762}{116.858}{remainder}%
\TextAt{8}{284.466}{116.858}{.}%
\TextAt{8}{292.566}{116.858}{Let}%
\TextAt{13}{312.489}{116.858}{f}%
\TextAt{13}{318.873}{116.858}{\Character{59}}%
\TextAt{13}{324.11}{116.858}{g}%
\NativeMathAt{333.484}{116.858}{11.955}{\in}%
\TextAt{18}{344.775}{116.858}{Z}%
\TextAt{21}{352.745}{115.065}{p}%
\TextAt{9}{357.506}{116.858}{\Character{91}\Character{91}}%
\TextAt{13}{364.009}{116.858}{X}%
\TextAt{9}{374.664}{116.858}{\Character{93}\Character{93}}%
\TextAt{8}{384.019}{116.858}{and}%
\TextAt{8}{405.735}{116.858}{let}%
\TextAt{13}{421.593}{116.858}{f}%
\TextAt{9}{428.639}{116.858}{\Character{40}}%
\TextAt{13}{433.191}{116.858}{X}%
\TextAt{9}{443.846}{116.858}{\Character{41}}%
\TextAt{9}{451.701}{116.858}{\Character{61}}%
\TextAt{13}{464.145}{116.858}{a}%
\TextAt{17}{470.29}{115.065}{0}%
\TextAt{9}{475.539}{116.858}{\Character{43}}%
\TextAt{13}{485.16}{116.858}{a}%
\TextAt{17}{491.305}{115.065}{1}%
\TextAt{13}{496.037}{116.858}{X}%
\TextAt{9}{507.208}{116.858}{\Character{43}}%
\TextAt{16}{516.829}{116.858}{\Character{1}}%
\TextAt{16}{522.125}{116.858}{\Character{1}}%
\TextAt{16}{527.433}{116.858}{\Character{1}}%
\TextAt{8}{79.226}{99.812}{such}%
\TextAt{8}{105.623}{99.812}{that}%
\TextAt{13}{130.992}{99.812}{p}%
\TextAt{16}{140.188}{99.812}{\Character{106}}%
\TextAt{13}{146.83}{99.812}{a}%
\TextAt{21}{152.975}{98.019}{\Character{23}}%
\TextAt{8}{162.054}{99.812}{for}%
\TextAt{13}{179.939}{99.812}{\Character{23}}%
\TextAt{9}{189.798}{99.812}{\Character{61}}%
\TextAt{9}{202.208}{99.812}{0}%
\TextAt{13}{208.076}{99.812}{\Character{59}}%
\TextAt{13}{213.3}{99.812}{\Character{58}}%
\TextAt{13}{218.537}{99.812}{\Character{58}}%
\TextAt{13}{223.773}{99.812}{\Character{58}}%
\TextAt{13}{228.998}{99.812}{\Character{59}}%
\TextAt{13}{234.234}{99.812}{n}%
\TextAt{16}{243.941}{99.812}{\Character{0}}%
\TextAt{9}{255.896}{99.812}{1}%
\TextAt{8}{261.749}{99.812}{,}%
\TextAt{8}{268.899}{99.812}{but}%
\TextAt{13}{290.358}{99.812}{p}%
\TextAt{18}{299.554}{99.812}{\Character{45}}%
\TextAt{13}{306.196}{99.812}{a}%
\TextAt{21}{312.341}{98.019}{n}%
\TextAt{8}{317.977}{99.812}{.}%
\TextAt{8}{326.425}{99.812}{Then}%
\TextAt{8}{356.983}{99.812}{one}%
\TextAt{8}{378.442}{99.812}{may}%
\TextAt{8}{403.811}{99.812}{write}%
\TextAt{8}{433.723}{99.812}{in}%
\TextAt{8}{447.376}{99.812}{a}%
\TextAt{8}{457.131}{99.812}{unique}%
\TextAt{8}{495.173}{99.812}{way}%
\TextAt{13}{277.465}{72.921}{g}%
\TextAt{9}{286.82}{72.921}{\Character{61}}%
\TextAt{13}{299.245}{72.921}{q}%
\TextAt{13}{304.852}{72.921}{f}%
\TextAt{9}{314.568}{72.921}{\Character{43}}%
\TextAt{13}{326.329}{72.921}{r}%
\TextAt{13}{331.279}{72.921}{\Character{59}}%
\TextAt{8}{297.211}{43.033}{105}%
\Destination{78.23}{809.815}{nav.page.105}%
\Destination{79.226}{431.555}{nav.ex.2.2.1}%
\Destination{79.226}{400.755}{nav.ex.2.2.2}%
\Destination{79.226}{352.955}{nav.ex.2.2.3}%
\Destination{79.226}{322.155}{nav.ex.2.2.4}%
\Destination{79.226}{240.255}{nav.ex.2.2.5}%
\Destination{79.226}{209.455}{nav.ex.2.2.6}%
\Destination{79.226}{178.655}{nav.ex.2.2.7}%
\Destination{79.226}{147.855}{nav.ex.2.2.8}%
\Destination{79.36}{776.155}{nav.print.115}%
}

% Page 118 (body page 106)
\Page{%
\TextAt{2}{79.226}{796.943}{II.}%
\TextAt{2}{94.073}{796.943}{The}%
\TextAt{2}{116.492}{796.943}{Theory}%
\TextAt{2}{154.401}{796.943}{of}%
\TextAt{2}{166.826}{796.943}{Valuations}%
\Vector{q 0 G 0 g 0.299 w 7.226 18.729 m 460.769 18.729 l S Q}%
\TextAt{8}{79.226}{761.124}{where}%
\TextAt{13}{113.035}{761.124}{q}%
\NativeMathAt{121.975}{761.124}{11.955}{\in}%
\TextAt{18}{133.266}{761.124}{Z}%
\TextAt{21}{141.236}{759.33}{p}%
\TextAt{9}{145.997}{761.124}{\Character{91}\Character{91}}%
\TextAt{13}{152.501}{761.124}{X}%
\TextAt{9}{163.156}{761.124}{\Character{93}\Character{93}}%
\TextAt{8}{169.659}{761.124}{,}%
\TextAt{8}{176.808}{761.124}{and}%
\TextAt{13}{199.571}{761.124}{r}%
\NativeMathAt{208.492}{761.124}{11.955}{\in}%
\TextAt{18}{219.783}{761.124}{Z}%
\TextAt{21}{227.753}{759.33}{p}%
\TextAt{9}{232.514}{761.124}{\Character{91}}%
\TextAt{13}{235.766}{761.124}{X}%
\TextAt{9}{246.421}{761.124}{\Character{93}}%
\TextAt{8}{253.57}{761.124}{is}%
\TextAt{8}{265.334}{761.124}{a}%
\TextAt{8}{275.089}{761.124}{polynomial}%
\TextAt{8}{335.583}{761.124}{of}%
\TextAt{8}{348.913}{761.124}{degree}%
\NativeMathAt{385.328}{761.124}{11.955}{\leq}%
\TextAt{13}{397.947}{761.124}{n}%
\TextAt{16}{407.592}{761.124}{\Character{0}}%
\TextAt{9}{419.547}{761.124}{1}%
\TextAt{8}{425.4}{761.124}{.}%
\TextAt{12}{97.159}{741.089}{Hint:}%
\TextAt{8}{135.307}{741.089}{Let}%
\TextAt{13}{157.581}{741.089}{\Character{28}}%
\TextAt{8}{169.203}{741.089}{be}%
\TextAt{8}{186.432}{741.089}{the}%
\TextAt{8}{207.893}{741.089}{operator}%
\TextAt{13}{256.361}{741.089}{\Character{28}}%
\TextAt{9}{262.78}{741.089}{\Character{40}}%
\TextAt{20}{267.332}{750.055}{\Character{80}}%
\NativeMathAt{279.951}{746.901}{7.97}{\infty}%
\TextAt{21}{279.951}{737.602}{\Character{23}}%
\TextAt{17}{284.635}{737.602}{\Character{61}0}%
\TextAt{13}{297.947}{741.089}{b}%
\TextAt{21}{302.924}{739.296}{\Character{23}}%
\TextAt{13}{308.106}{741.089}{X}%
\TextAt{21}{318.761}{745.427}{\Character{23}}%
\TextAt{9}{323.943}{741.089}{\Character{41}}%
\TextAt{9}{334.021}{741.089}{\Character{61}}%
\TextAt{20}{348.689}{750.055}{\Character{80}}%
\NativeMathAt{361.308}{746.901}{7.97}{\infty}%
\TextAt{21}{361.308}{737.602}{\Character{23}}%
\TextAt{17}{365.992}{737.602}{\Character{61}}%
\TextAt{21}{372.578}{737.602}{n}%
\TextAt{13}{380.207}{741.089}{b}%
\TextAt{21}{385.184}{739.296}{\Character{23}}%
\TextAt{13}{390.366}{741.089}{X}%
\TextAt{21}{401.022}{745.427}{\Character{23}}%
\TextAt{3}{405.705}{745.427}{\Character{0}}%
\TextAt{21}{412.292}{745.427}{n}%
\TextAt{8}{417.928}{741.089}{.}%
\TextAt{8}{430.291}{741.089}{Show}%
\TextAt{8}{462.488}{741.089}{that}%
\TextAt{13}{489.162}{741.089}{U}%
\TextAt{9}{498.361}{741.089}{\Character{40}}%
\TextAt{13}{502.913}{741.089}{X}%
\TextAt{9}{513.568}{741.089}{\Character{41}}%
\TextAt{9}{523.646}{741.089}{\Character{61}}%
\TextAt{13}{79.226}{724.043}{a}%
\TextAt{21}{85.371}{722.25}{n}%
\TextAt{9}{94.164}{724.043}{\Character{43}}%
\TextAt{13}{106.426}{724.043}{a}%
\TextAt{21}{112.571}{722.25}{n}%
\TextAt{17}{117.709}{722.25}{\Character{43}1}%
\TextAt{13}{129.028}{724.043}{X}%
\TextAt{9}{142.84}{724.043}{\Character{43}}%
\TextAt{16}{155.101}{724.043}{\Character{1}}%
\TextAt{16}{160.397}{724.043}{\Character{1}}%
\TextAt{16}{165.705}{724.043}{\Character{1}}%
\TextAt{9}{173.62}{724.043}{\Character{61}}%
\TextAt{13}{187.296}{724.043}{\Character{28}}%
\TextAt{9}{193.715}{724.043}{\Character{40}}%
\TextAt{13}{198.267}{724.043}{f}%
\TextAt{9}{205.314}{724.043}{\Character{40}}%
\TextAt{13}{209.866}{724.043}{X}%
\TextAt{9}{220.521}{724.043}{\Character{41}\Character{41}}%
\TextAt{8}{234.257}{724.043}{is}%
\TextAt{8}{246.755}{724.043}{a}%
\TextAt{8}{257.244}{724.043}{unit}%
\TextAt{8}{282.689}{724.043}{in}%
\TextAt{18}{297.076}{724.043}{Z}%
\TextAt{21}{305.046}{722.25}{p}%
\TextAt{9}{309.807}{724.043}{\Character{91}\Character{91}}%
\TextAt{13}{316.31}{724.043}{X}%
\TextAt{9}{326.965}{724.043}{\Character{93}\Character{93}}%
\TextAt{8}{338.1}{724.043}{and}%
\TextAt{8}{361.596}{724.043}{write}%
\TextAt{13}{392.242}{724.043}{f}%
\TextAt{9}{399.288}{724.043}{\Character{40}}%
\TextAt{13}{403.84}{724.043}{X}%
\TextAt{9}{414.495}{724.043}{\Character{41}}%
\TextAt{9}{423.605}{724.043}{\Character{61}}%
\TextAt{13}{437.295}{724.043}{pP}%
\TextAt{9}{452.341}{724.043}{\Character{40}}%
\TextAt{13}{456.894}{724.043}{X}%
\TextAt{9}{467.549}{724.043}{\Character{41}}%
\TextAt{9}{475.248}{724.043}{\Character{43}}%
\TextAt{13}{487.519}{724.043}{X}%
\TextAt{21}{498.174}{728.381}{n}%
\TextAt{13}{503.811}{724.043}{U}%
\TextAt{9}{513.01}{724.043}{\Character{40}}%
\TextAt{13}{517.562}{724.043}{X}%
\TextAt{9}{528.217}{724.043}{\Character{41}}%
\TextAt{8}{79.226}{706.997}{with}%
\TextAt{8}{105.886}{706.997}{a}%
\TextAt{8}{115.642}{706.997}{polynomial}%
\TextAt{13}{176.135}{706.997}{P}%
\TextAt{9}{185.306}{706.997}{\Character{40}}%
\TextAt{13}{189.858}{706.997}{X}%
\TextAt{9}{200.513}{706.997}{\Character{41}}%
\TextAt{8}{208.963}{706.997}{of}%
\TextAt{8}{222.293}{706.997}{degree}%
\NativeMathAt{258.708}{706.997}{11.955}{\leq}%
\TextAt{13}{271.328}{706.997}{n}%
\TextAt{16}{280.972}{706.997}{\Character{0}}%
\TextAt{9}{292.927}{706.997}{1}%
\TextAt{8}{298.78}{706.997}{.}%
\TextAt{8}{307.229}{706.997}{Show}%
\TextAt{8}{338.121}{706.997}{that}%
\TextAt{13}{199.121}{672.941}{q}%
\TextAt{9}{204.74}{672.941}{\Character{40}}%
\TextAt{13}{209.293}{672.941}{X}%
\TextAt{9}{219.948}{672.941}{\Character{41}}%
\TextAt{9}{227.803}{672.941}{\Character{61}}%
\TextAt{9}{252.995}{681.028}{1}%
\Vector{q 0 G 0 g 0.478 w 169.442 -98.071 m 198.401 -98.071 l S Q}%
\TextAt{13}{241.442}{664.74}{U}%
\TextAt{9}{250.641}{664.74}{\Character{40}}%
\TextAt{13}{255.193}{664.74}{X}%
\TextAt{9}{265.848}{664.74}{\Character{41}}%
\NativeMathAt{277.989}{687.885}{7.97}{\infty}%
\TextAt{20}{273.589}{684.298}{\Character{88}}%
\TextAt{21}{275.371}{659.104}{i}%
\TextAt{17}{278.254}{659.104}{\Character{61}0}%
\TextAt{9}{290.857}{672.941}{\Character{40}}%
\TextAt{16}{295.41}{672.941}{\Character{0}}%
\TextAt{9}{304.708}{672.941}{1\Character{41}}%
\TextAt{21}{315.114}{677.877}{i}%
\TextAt{13}{318.495}{672.941}{p}%
\TextAt{21}{324.37}{677.877}{i}%
\TextAt{20}{329.744}{689.797}{\Character{18}}%
\TextAt{13}{338.544}{672.941}{\Character{28}}%
\TextAt{16}{347.62}{672.941}{\Character{14}}%
\TextAt{13}{357.463}{681.028}{P}%
\Vector{q 0 G 0 g 0.478 w 285.449 -98.071 m 294.649 -98.071 l S Q}%
\TextAt{13}{357.449}{664.74}{U}%
\TextAt{20}{367.844}{689.797}{\Character{19}}%
\TextAt{21}{376.644}{687.121}{i}%
\TextAt{16}{382.682}{672.941}{\Character{14}}%
\TextAt{13}{391.317}{672.941}{\Character{28}}%
\TextAt{9}{397.736}{672.941}{\Character{40}}%
\TextAt{13}{402.288}{672.941}{g}%
\TextAt{9}{408.322}{672.941}{\Character{41}}%
\TextAt{8}{79.226}{637.445}{is}%
\TextAt{8}{90.99}{637.445}{a}%
\TextAt{8}{100.745}{637.445}{well-defined}%
\TextAt{8}{164.789}{637.445}{power}%
\TextAt{8}{198.933}{637.445}{series}%
\TextAt{8}{230.268}{637.445}{in}%
\TextAt{18}{243.921}{637.445}{Z}%
\TextAt{21}{251.891}{635.652}{p}%
\TextAt{9}{256.652}{637.445}{\Character{91}\Character{91}}%
\TextAt{13}{263.155}{637.445}{X}%
\TextAt{9}{273.81}{637.445}{\Character{93}\Character{93}}%
\TextAt{8}{284.211}{637.445}{such}%
\TextAt{8}{310.608}{637.445}{that}%
\TextAt{13}{335.977}{637.445}{\Character{28}}%
\TextAt{9}{342.396}{637.445}{\Character{40}}%
\TextAt{13}{346.948}{637.445}{q}%
\TextAt{13}{352.555}{637.445}{f}%
\TextAt{9}{359.614}{637.445}{\Character{41}}%
\TextAt{9}{367.469}{637.445}{\Character{61}}%
\TextAt{13}{379.912}{637.445}{\Character{28}}%
\TextAt{9}{386.331}{637.445}{\Character{40}}%
\TextAt{13}{390.883}{637.445}{g}%
\TextAt{9}{396.918}{637.445}{\Character{41}}%
\TextAt{8}{401.47}{637.445}{.}%
\TextAt{12}{79.226}{606.287}{Exercise}%
\TextAt{12}{131.974}{606.287}{II.2.9.}%
\TextAt{12}{170.555}{606.287}{(}%
\TextAt{13}{175.791}{606.287}{p}%
\TextAt{12}{181.666}{606.287}{-adic}%
\TextAt{12}{213.911}{606.287}{Weierstrass}%
\TextAt{12}{285.417}{606.287}{Preparation}%
\TextAt{12}{359.505}{606.287}{Theorem).}%
\TextAt{8}{426.079}{606.287}{Every}%
\TextAt{8}{459.34}{606.287}{nonzero}%
\TextAt{8}{502.523}{606.287}{power}%
\TextAt{8}{79.226}{589.241}{series}%
\TextAt{13}{236.007}{568.554}{f}%
\TextAt{9}{243.053}{568.554}{\Character{40}}%
\TextAt{13}{247.605}{568.554}{X}%
\TextAt{9}{258.261}{568.554}{\Character{41}}%
\TextAt{9}{266.116}{568.554}{\Character{61}}%
\NativeMathAt{282.959}{583.498}{7.97}{\infty}%
\TextAt{20}{278.559}{579.911}{\Character{88}}%
\TextAt{21}{279.441}{554.805}{\Character{23}}%
\TextAt{17}{284.125}{554.805}{\Character{61}0}%
\TextAt{13}{297.82}{568.554}{a}%
\TextAt{21}{303.965}{566.76}{\Character{23}}%
\TextAt{13}{309.147}{568.554}{X}%
\TextAt{21}{319.802}{573.49}{\Character{23}}%
\NativeMathAt{328.305}{568.554}{11.955}{\in}%
\TextAt{18}{339.596}{568.554}{Z}%
\TextAt{21}{347.566}{566.76}{p}%
\TextAt{9}{352.327}{568.554}{\Character{91}\Character{91}}%
\TextAt{13}{358.831}{568.554}{X}%
\TextAt{9}{369.486}{568.554}{\Character{93}\Character{93}}%
\TextAt{8}{79.226}{538.464}{admits}%
\TextAt{8}{117.662}{538.464}{a}%
\TextAt{8}{127.417}{538.464}{unique}%
\TextAt{8}{165.459}{538.464}{representation}%
\TextAt{13}{248.417}{510.918}{f}%
\TextAt{9}{255.464}{510.918}{\Character{40}}%
\TextAt{13}{260.016}{510.918}{X}%
\TextAt{9}{270.671}{510.918}{\Character{41}}%
\TextAt{9}{278.526}{510.918}{\Character{61}}%
\TextAt{13}{290.97}{510.918}{p}%
\TextAt{21}{296.845}{515.854}{\Character{22}}%
\TextAt{13}{302.437}{510.918}{P}%
\TextAt{9}{311.608}{510.918}{\Character{40}}%
\TextAt{13}{316.161}{510.918}{X}%
\TextAt{9}{326.816}{510.918}{\Character{41}}%
\TextAt{13}{331.368}{510.918}{U}%
\TextAt{9}{340.567}{510.918}{\Character{40}}%
\TextAt{13}{345.119}{510.918}{X}%
\TextAt{9}{355.775}{510.918}{\Character{41}}%
\TextAt{13}{360.327}{510.918}{\Character{59}}%
\TextAt{8}{79.226}{483.372}{where}%
\TextAt{13}{114.368}{483.372}{U}%
\TextAt{9}{123.567}{483.372}{\Character{40}}%
\TextAt{13}{128.12}{483.372}{X}%
\TextAt{9}{138.775}{483.372}{\Character{41}}%
\TextAt{8}{148.557}{483.372}{is}%
\TextAt{8}{161.654}{483.372}{a}%
\TextAt{8}{172.742}{483.372}{unit}%
\TextAt{8}{198.787}{483.372}{in}%
\TextAt{18}{213.772}{483.372}{Z}%
\TextAt{21}{221.742}{481.579}{p}%
\TextAt{9}{226.503}{483.372}{\Character{91}\Character{91}}%
\TextAt{13}{233.007}{483.372}{X}%
\TextAt{9}{243.662}{483.372}{\Character{93}\Character{93}}%
\TextAt{8}{255.395}{483.372}{and}%
\TextAt{13}{279.491}{483.372}{P}%
\TextAt{9}{288.662}{483.372}{\Character{40}}%
\TextAt{13}{293.214}{483.372}{X}%
\TextAt{9}{303.869}{483.372}{\Character{41}}%
\NativeMathAt{314.014}{483.372}{11.955}{\in}%
\TextAt{18}{327.576}{483.372}{Z}%
\TextAt{21}{335.546}{481.579}{p}%
\TextAt{9}{340.307}{483.372}{\Character{91}}%
\TextAt{13}{343.559}{483.372}{X}%
\TextAt{9}{354.214}{483.372}{\Character{93}}%
\TextAt{8}{362.696}{483.372}{is}%
\TextAt{8}{375.793}{483.372}{a}%
\TextAt{8}{386.881}{483.372}{monic}%
\TextAt{8}{422.681}{483.372}{polynomial}%
\TextAt{8}{484.507}{483.372}{satisfying}%
\TextAt{13}{79.226}{466.326}{P}%
\TextAt{9}{88.397}{466.326}{\Character{40}}%
\TextAt{13}{92.95}{466.326}{X}%
\TextAt{9}{103.605}{466.326}{\Character{41}}%
\NativeMathAt{111.478}{466.326}{11.955}{\equiv}%
\TextAt{13}{124.097}{466.326}{X}%
\TextAt{21}{134.752}{470.664}{n}%
\TextAt{8}{143.709}{466.326}{mod}%
\TextAt{13}{169.147}{466.326}{p}%
\TextAt{8}{175.022}{466.326}{.}%
\TextAt{15}{79.226}{414.093}{§}%
\TextAt{15}{93.785}{414.093}{3}%
\TextAt{15}{120.693}{414.093}{Valuations}%
\TextAt{8}{79.226}{382.207}{The}%
\TextAt{8}{102.404}{382.207}{procedure}%
\TextAt{8}{155.829}{382.207}{we}%
\TextAt{8}{172.181}{382.207}{performed}%
\TextAt{8}{227.232}{382.207}{in}%
\TextAt{8}{240.009}{382.207}{the}%
\TextAt{8}{259.29}{382.207}{previous}%
\TextAt{8}{304.98}{382.207}{section}%
\TextAt{8}{343.186}{382.207}{with}%
\TextAt{8}{368.97}{382.207}{the}%
\TextAt{8}{388.251}{382.207}{field}%
\TextAt{18}{412.732}{382.207}{Q}%
\TextAt{8}{425.053}{382.207}{in}%
\TextAt{8}{437.83}{382.207}{order}%
\TextAt{8}{467.524}{382.207}{to}%
\TextAt{8}{480.959}{382.207}{obtain}%
\TextAt{8}{516.51}{382.207}{the}%
\TextAt{13}{79.226}{365.161}{p}%
\TextAt{8}{85.101}{365.161}{-adic}%
\TextAt{8}{113.198}{365.161}{numbers}%
\TextAt{8}{159.898}{365.161}{can}%
\TextAt{8}{180.846}{365.161}{be}%
\TextAt{8}{196.259}{365.161}{generalized}%
\TextAt{8}{256.228}{365.161}{to}%
\TextAt{8}{270.027}{365.161}{arbitrary}%
\TextAt{8}{319.285}{365.161}{fields}%
\TextAt{8}{348.745}{365.161}{using}%
\TextAt{8}{378.863}{365.161}{the}%
\TextAt{8}{398.508}{365.161}{concept}%
\TextAt{8}{440.915}{365.161}{of}%
\TextAt{8}{453.734}{365.161}{(multiplicative)}%
\TextAt{8}{79.226}{348.115}{valuation.}%
\TextAt{12}{79.226}{322.518}{Definition}%
\TextAt{12}{142.015}{322.518}{(3.1).}%
\TextAt{8}{177.33}{322.518}{A}%
\TextAt{12}{190.003}{322.518}{valuation}%
\TextAt{8}{247.722}{322.518}{of}%
\TextAt{8}{261.052}{322.518}{a}%
\TextAt{8}{270.808}{322.518}{field}%
\TextAt{13}{296.165}{322.518}{K}%
\TextAt{8}{310.88}{322.518}{is}%
\TextAt{8}{322.644}{322.518}{a}%
\TextAt{8}{332.399}{322.518}{function}%
\TextAt{16}{273.104}{294.972}{\Character{106}}%
\TextAt{16}{280.313}{294.972}{\Character{106}}%
\TextAt{9}{286.964}{294.972}{\Character{58}}%
\TextAt{13}{293.537}{294.972}{K}%
\NativeArrowAt{307.675}{294.972}{11.955}{19.26}{\longrightarrow}%
\TextAt{18}{330.257}{294.972}{R}%
\TextAt{8}{79.226}{267.426}{enjoying}%
\TextAt{8}{125.732}{267.426}{the}%
\TextAt{8}{145.888}{267.426}{properties}%
\TextAt{8}{89.579}{247.391}{(i)}%
\TextAt{16}{107.918}{247.391}{\Character{106}}%
\TextAt{13}{111.239}{247.391}{x}%
\TextAt{16}{117.891}{247.391}{\Character{106}}%
\NativeMathAt{124.514}{247.391}{11.955}{\geq}%
\TextAt{9}{137.153}{247.391}{0}%
\TextAt{8}{143.005}{247.391}{,}%
\TextAt{8}{150.155}{247.391}{and}%
\TextAt{16}{172.917}{247.391}{\Character{106}}%
\TextAt{13}{176.238}{247.391}{x}%
\TextAt{16}{182.89}{247.391}{\Character{106}}%
\TextAt{9}{189.532}{247.391}{\Character{61}}%
\TextAt{9}{201.942}{247.391}{0}%
\NativeArrowAt{211.131}{247.391}{11.955}{21.926}{\Longleftrightarrow}%
\TextAt{13}{236.37}{247.391}{x}%
\TextAt{9}{246.343}{247.391}{\Character{61}}%
\TextAt{9}{258.753}{247.391}{0}%
\TextAt{8}{264.621}{247.391}{,}%
\TextAt{8}{86.327}{230.345}{(ii)}%
\TextAt{16}{107.918}{230.345}{\Character{106}}%
\TextAt{13}{111.239}{230.345}{xy}%
\TextAt{16}{124.028}{230.345}{\Character{106}}%
\TextAt{9}{130.67}{230.345}{\Character{61}}%
\TextAt{16}{143.095}{230.345}{\Character{106}}%
\TextAt{13}{146.416}{230.345}{x}%
\TextAt{16}{153.068}{230.345}{\Character{106}\Character{106}}%
\TextAt{13}{159.71}{230.345}{y}%
\TextAt{16}{165.847}{230.345}{\Character{106}}%
\TextAt{8}{169.168}{230.345}{,}%
\TextAt{8}{83.076}{213.299}{(iii)}%
\TextAt{16}{107.918}{213.299}{\Character{106}}%
\TextAt{13}{111.239}{213.299}{x}%
\TextAt{9}{120.548}{213.299}{\Character{43}}%
\TextAt{13}{132.309}{213.299}{y}%
\TextAt{16}{138.446}{213.299}{\Character{106}}%
\NativeMathAt{145.069}{213.299}{11.955}{\leq}%
\TextAt{16}{157.682}{213.299}{\Character{106}}%
\TextAt{13}{161.028}{213.299}{x}%
\TextAt{16}{167.68}{213.299}{\Character{106}}%
\TextAt{9}{173.658}{213.299}{\Character{43}}%
\TextAt{16}{185.419}{213.299}{\Character{106}}%
\TextAt{13}{188.74}{213.299}{y}%
\TextAt{16}{194.877}{213.299}{\Character{106}}%
\TextAt{8}{214.05}{213.299}{“triangle}%
\TextAt{8}{262.504}{213.299}{inequality”.}%
\TextAt{8}{97.159}{187.702}{We}%
\TextAt{8}{117.275}{187.702}{tacitly}%
\TextAt{8}{153.997}{187.702}{exclude}%
\TextAt{8}{195.907}{187.702}{in}%
\TextAt{8}{209.531}{187.702}{the}%
\TextAt{8}{229.659}{187.702}{sequel}%
\TextAt{8}{264.48}{187.702}{the}%
\TextAt{8}{284.608}{187.702}{case}%
\TextAt{8}{309.351}{187.702}{where}%
\TextAt{16}{343.132}{187.702}{\Character{106}}%
\TextAt{16}{350.305}{187.702}{\Character{106}}%
\TextAt{8}{357.512}{187.702}{is}%
\TextAt{8}{369.247}{187.702}{the}%
\TextAt{8}{389.375}{187.702}{trivial}%
\TextAt{8}{424.148}{187.702}{valuation}%
\TextAt{8}{475.192}{187.702}{of}%
\TextAt{13}{488.494}{187.702}{K}%
\TextAt{8}{503.18}{187.702}{which}%
\TextAt{8}{79.226}{170.657}{satisfies}%
\TextAt{16}{122.336}{170.657}{\Character{106}}%
\TextAt{13}{125.657}{170.657}{x}%
\TextAt{16}{132.309}{170.657}{\Character{106}}%
\TextAt{9}{138.951}{170.657}{\Character{61}}%
\TextAt{9}{151.361}{170.657}{1}%
\TextAt{8}{161.127}{170.657}{for}%
\TextAt{8}{179.012}{170.657}{all}%
\TextAt{13}{195.271}{170.657}{x}%
\NativeMathAt{205.244}{170.657}{11.955}{\neq}%
\TextAt{9}{217.654}{170.657}{0}%
\TextAt{8}{223.522}{170.657}{.}%
\TextAt{8}{231.971}{170.657}{Defining}%
\TextAt{8}{278.632}{170.657}{the}%
\TextAt{8}{298.788}{170.657}{distance}%
\TextAt{8}{344.373}{170.657}{between}%
\TextAt{8}{389.564}{170.657}{two}%
\TextAt{8}{411.681}{170.657}{points}%
\TextAt{13}{446.865}{170.657}{x\Character{59}}%
\TextAt{13}{458.737}{170.657}{y}%
\NativeMathAt{468.219}{170.657}{11.955}{\in}%
\TextAt{13}{479.51}{170.657}{K}%
\TextAt{8}{494.225}{170.657}{by}%
\TextAt{13}{264.493}{143.11}{d}%
\TextAt{9}{270.576}{143.11}{\Character{40}}%
\TextAt{13}{275.128}{143.11}{x\Character{59}}%
\TextAt{13}{287}{143.11}{y}%
\TextAt{9}{293.161}{143.11}{\Character{41}}%
\TextAt{9}{301.016}{143.11}{\Character{61}}%
\TextAt{16}{313.46}{143.11}{\Character{106}}%
\TextAt{13}{316.781}{143.11}{x}%
\TextAt{16}{326.09}{143.11}{\Character{0}}%
\TextAt{13}{338.045}{143.11}{y}%
\TextAt{16}{344.181}{143.11}{\Character{106}}%
\TextAt{8}{79.226}{115.564}{makes}%
\TextAt{13}{114.41}{115.564}{K}%
\TextAt{8}{129.125}{115.564}{into}%
\TextAt{8}{152.868}{115.564}{a}%
\TextAt{8}{162.624}{115.564}{metric}%
\TextAt{8}{199.039}{115.564}{space,}%
\TextAt{8}{233.566}{115.564}{and}%
\TextAt{8}{256.328}{115.564}{hence}%
\TextAt{8}{288.834}{115.564}{in}%
\TextAt{8}{302.487}{115.564}{particular}%
\TextAt{8}{356.477}{115.564}{a}%
\TextAt{8}{366.232}{115.564}{topological}%
\TextAt{8}{425.757}{115.564}{space.}%
\TextAt{12}{79.226}{89.967}{Definition}%
\TextAt{12}{143.22}{89.967}{(3.2).}%
\TextAt{8}{179.583}{89.967}{Two}%
\TextAt{8}{205.987}{89.967}{valuations}%
\TextAt{8}{262.722}{89.967}{of}%
\TextAt{13}{277.099}{89.967}{K}%
\TextAt{8}{292.862}{89.967}{are}%
\TextAt{8}{313.42}{89.967}{called}%
\TextAt{12}{347.631}{89.967}{equivalent}%
\TextAt{8}{412.124}{89.967}{if}%
\TextAt{8}{423.896}{89.967}{they}%
\TextAt{8}{451.28}{89.967}{define}%
\TextAt{8}{486.137}{89.967}{the}%
\TextAt{8}{507.341}{89.967}{same}%
\TextAt{8}{79.226}{72.921}{topology}%
\TextAt{8}{127.37}{72.921}{on}%
\TextAt{13}{143.629}{72.921}{K}%
\TextAt{8}{154.446}{72.921}{.}%
\TextAt{8}{297.211}{43.033}{106}%
\Destination{79.23}{261.985}{nav.Item.33}%
\Destination{79.23}{244.395}{nav.Item.34}%
\Destination{79.23}{227.345}{nav.Item.35}%
\Destination{78.23}{809.815}{nav.page.106}%
\Destination{79.23}{439.955}{nav.section.2.3}%
\Destination{79.226}{620.255}{nav.ex.2.2.9}%
\Destination{79.36}{756.055}{nav.print.116}%
\Destination{79.36}{336.555}{nav.number.2.3.1}%
\Destination{79.36}{103.955}{nav.number.2.3.2}%
\Destination{79.23}{439.955}{outline.23}%
}

% Page 119 (body page 107)
\Page{%
\TextAt{2}{79.226}{796.943}{II.}%
\TextAt{2}{94.073}{796.943}{The}%
\TextAt{2}{116.492}{796.943}{Theory}%
\TextAt{2}{154.401}{796.943}{of}%
\TextAt{2}{166.826}{796.943}{Valuations}%
\Vector{q 0 G 0 g 0.299 w 7.226 18.729 m 460.769 18.729 l S Q}%
\TextAt{12}{79.226}{761.124}{Proposition}%
\TextAt{12}{152.028}{761.124}{(3.3).}%
\TextAt{8}{187.423}{761.124}{Two}%
\TextAt{8}{212.859}{761.124}{valuations}%
\TextAt{16}{268.625}{761.124}{\Character{106}}%
\TextAt{16}{275.906}{761.124}{\Character{106}}%
\TextAt{17}{279.243}{759.33}{1}%
\TextAt{8}{287.952}{761.124}{and}%
\TextAt{16}{310.794}{761.124}{\Character{106}}%
\TextAt{16}{318.075}{761.124}{\Character{106}}%
\TextAt{17}{321.412}{759.33}{2}%
\TextAt{8}{330.121}{761.124}{on}%
\TextAt{13}{346.459}{761.124}{K}%
\TextAt{8}{361.253}{761.124}{are}%
\TextAt{8}{380.843}{761.124}{equivalent}%
\TextAt{8}{436.537}{761.124}{if}%
\TextAt{8}{447.34}{761.124}{and}%
\TextAt{8}{470.182}{761.124}{only}%
\TextAt{8}{495.952}{761.124}{if}%
\TextAt{8}{506.755}{761.124}{there}%
\TextAt{8}{79.226}{744.078}{exists}%
\TextAt{8}{111.541}{744.078}{a}%
\TextAt{8}{121.296}{744.078}{real}%
\TextAt{8}{144.059}{744.078}{number}%
\TextAt{13}{186.655}{744.078}{s}%
\TextAt{13}{195.478}{744.078}{\Character{62}}%
\TextAt{9}{207.916}{744.078}{0}%
\TextAt{8}{217.666}{744.078}{such}%
\TextAt{8}{244.063}{744.078}{that}%
\TextAt{8}{269.432}{744.078}{one}%
\TextAt{8}{290.891}{744.078}{has}%
\TextAt{16}{280.098}{715.357}{\Character{106}}%
\TextAt{13}{283.419}{715.357}{x}%
\TextAt{16}{290.071}{715.357}{\Character{106}}%
\TextAt{17}{293.392}{713.564}{1}%
\TextAt{9}{301.445}{715.357}{\Character{61}}%
\TextAt{16}{313.871}{715.357}{\Character{106}}%
\TextAt{13}{317.192}{715.357}{x}%
\TextAt{16}{323.844}{715.357}{\Character{106}}%
\TextAt{21}{327.165}{720.293}{s}%
\TextAt{17}{327.165}{712.402}{2}%
\TextAt{8}{79.226}{686.637}{for}%
\TextAt{8}{97.111}{686.637}{all}%
\TextAt{13}{113.37}{686.637}{x}%
\NativeMathAt{123.343}{686.637}{11.955}{\in}%
\TextAt{13}{134.634}{686.637}{K}%
\TextAt{8}{145.452}{686.637}{.}%
\TextAt{19}{79.226}{651.818}{Proof.}%
\TextAt{8}{116}{651.818}{If}%
\TextAt{16}{126.964}{651.818}{\Character{106}}%
\TextAt{16}{133.444}{651.818}{\Character{106}}%
\TextAt{17}{136.776}{650.025}{1}%
\TextAt{9}{144.829}{651.818}{\Character{61}}%
\TextAt{16}{157.254}{651.818}{\Character{106}}%
\TextAt{16}{163.734}{651.818}{\Character{106}}%
\TextAt{21}{167.065}{656.157}{s}%
\TextAt{17}{167.065}{648.863}{2}%
\TextAt{8}{171.798}{651.818}{,}%
\TextAt{8}{178.364}{651.818}{with}%
\TextAt{13}{204.296}{651.818}{s}%
\TextAt{13}{213.119}{651.818}{\Character{62}}%
\TextAt{9}{225.557}{651.818}{0}%
\TextAt{8}{231.41}{651.818}{,}%
\TextAt{8}{237.977}{651.818}{then}%
\TextAt{16}{263.909}{651.818}{\Character{106}}%
\TextAt{16}{270.389}{651.818}{\Character{106}}%
\TextAt{17}{273.72}{650.025}{1}%
\TextAt{8}{281.621}{651.818}{and}%
\TextAt{16}{303.656}{651.818}{\Character{106}}%
\TextAt{16}{310.136}{651.818}{\Character{106}}%
\TextAt{17}{313.467}{650.025}{2}%
\TextAt{8}{321.369}{651.818}{are}%
\TextAt{8}{340.152}{651.818}{obviously}%
\TextAt{8}{391.201}{651.818}{equivalent.}%
\TextAt{8}{451.125}{651.818}{For}%
\TextAt{8}{471.366}{651.818}{an}%
\TextAt{8}{486.897}{651.818}{arbitrary}%
\TextAt{8}{79.226}{634.772}{valuation}%
\TextAt{16}{131.008}{634.772}{\Character{106}}%
\TextAt{16}{138.922}{634.772}{\Character{106}}%
\TextAt{8}{146.864}{634.772}{on}%
\TextAt{13}{163.832}{634.772}{K}%
\TextAt{8}{174.65}{634.772}{,}%
\TextAt{8}{182.686}{634.772}{the}%
\TextAt{8}{203.552}{634.772}{inequality}%
\TextAt{16}{258.573}{634.772}{\Character{106}}%
\TextAt{13}{261.894}{634.772}{x}%
\TextAt{16}{268.546}{634.772}{\Character{106}}%
\TextAt{13}{276.397}{634.772}{\Character{60}}%
\TextAt{9}{290.031}{634.772}{1}%
\TextAt{8}{300.491}{634.772}{is}%
\TextAt{8}{312.965}{634.772}{tantamount}%
\TextAt{8}{377.431}{634.772}{to}%
\TextAt{8}{392.451}{634.772}{the}%
\TextAt{8}{413.317}{634.772}{condition}%
\TextAt{8}{465.409}{634.772}{that}%
\TextAt{16}{491.488}{634.772}{\Character{102}}%
\TextAt{13}{497.465}{634.772}{x}%
\TextAt{21}{504.117}{639.111}{n}%
\TextAt{16}{509.754}{634.772}{\Character{103}}%
\TextAt{21}{515.731}{632.979}{n}%
\NativeMathAt{520.87}{632.979}{7.97}{\in}%
\TextAt{25}{526.515}{632.979}{N}%
\TextAt{8}{79.226}{617.727}{converges}%
\TextAt{8}{131.914}{617.727}{to}%
\TextAt{8}{146.488}{617.727}{zero}%
\TextAt{8}{171.464}{617.727}{in}%
\TextAt{8}{185.381}{617.727}{the}%
\TextAt{8}{205.801}{617.727}{topology}%
\TextAt{8}{254.209}{617.727}{defined}%
\TextAt{8}{294.786}{617.727}{by}%
\TextAt{16}{311.309}{617.727}{\Character{106}}%
\TextAt{16}{318.781}{617.727}{\Character{106}}%
\TextAt{8}{322.113}{617.727}{.}%
\TextAt{8}{331.354}{617.727}{Therefore}%
\TextAt{8}{384.615}{617.727}{if}%
\TextAt{16}{395.603}{617.727}{\Character{106}}%
\TextAt{16}{403.075}{617.727}{\Character{106}}%
\TextAt{17}{406.407}{615.933}{1}%
\TextAt{8}{415.3}{617.727}{and}%
\TextAt{16}{438.327}{617.727}{\Character{106}}%
\TextAt{16}{445.799}{617.727}{\Character{106}}%
\TextAt{17}{449.131}{615.933}{2}%
\TextAt{8}{458.025}{617.727}{are}%
\TextAt{8}{477.8}{617.727}{equivalent,}%
\TextAt{8}{79.226}{600.681}{one}%
\TextAt{8}{100.686}{600.681}{has}%
\TextAt{8}{121.559}{600.681}{the}%
\TextAt{8}{141.716}{600.681}{implication}%
\TextAt{16}{251.892}{583.635}{\Character{106}}%
\TextAt{13}{255.213}{583.635}{x}%
\TextAt{16}{261.865}{583.635}{\Character{106}}%
\TextAt{17}{265.186}{581.841}{1}%
\TextAt{13}{273.239}{583.635}{\Character{60}}%
\TextAt{9}{285.664}{583.635}{1}%
\NativeArrowAt{294.822}{583.635}{11.955}{19.083}{\Longrightarrow}%
\TextAt{16}{317.217}{583.635}{\Character{106}}%
\TextAt{13}{320.547}{583.635}{x}%
\TextAt{16}{327.199}{583.635}{\Character{106}}%
\TextAt{17}{330.52}{581.841}{2}%
\TextAt{13}{338.574}{583.635}{\Character{60}}%
\TextAt{9}{350.999}{583.635}{1}%
\TextAt{13}{356.852}{583.635}{\Character{58}}%
\TextAt{8}{517.802}{583.635}{(*)}%
\TextAt{8}{79.226}{560.233}{Now}%
\TextAt{8}{105.681}{560.233}{let}%
\TextAt{13}{122.38}{560.233}{y}%
\NativeMathAt{131.837}{560.233}{11.955}{\in}%
\TextAt{13}{143.128}{560.233}{K}%
\TextAt{8}{157.638}{560.233}{be}%
\TextAt{8}{173.357}{560.233}{a}%
\TextAt{8}{182.907}{560.233}{fixed}%
\TextAt{8}{210.987}{560.233}{element}%
\TextAt{8}{254.023}{560.233}{satisfying}%
\TextAt{16}{305.978}{560.233}{\Character{106}}%
\TextAt{13}{309.299}{560.233}{y}%
\TextAt{16}{315.436}{560.233}{\Character{106}}%
\TextAt{17}{318.757}{558.44}{1}%
\TextAt{13}{326.81}{560.233}{\Character{62}}%
\TextAt{9}{339.235}{560.233}{1}%
\TextAt{8}{345.088}{560.233}{.}%
\TextAt{8}{353.468}{560.233}{Let}%
\TextAt{13}{374.232}{560.233}{x}%
\NativeMathAt{384.205}{560.233}{11.955}{\in}%
\TextAt{13}{395.496}{560.233}{K}%
\TextAt{8}{406.314}{560.233}{,}%
\TextAt{13}{413.298}{560.233}{x}%
\NativeMathAt{423.271}{560.233}{11.955}{\neq}%
\TextAt{9}{435.68}{560.233}{0}%
\TextAt{8}{441.55}{560.233}{.}%
\TextAt{8}{449.93}{560.233}{Then}%
\TextAt{16}{480.282}{560.233}{\Character{106}}%
\TextAt{13}{483.602}{560.233}{x}%
\TextAt{16}{490.255}{560.233}{\Character{106}}%
\TextAt{17}{493.575}{558.44}{1}%
\TextAt{9}{501.629}{560.233}{\Character{61}}%
\TextAt{16}{514.054}{560.233}{\Character{106}}%
\TextAt{13}{517.375}{560.233}{y}%
\TextAt{16}{523.512}{560.233}{\Character{106}}%
\TextAt{21}{526.832}{564.572}{\Character{11}}%
\TextAt{17}{526.833}{557.278}{1}%
\TextAt{8}{79.226}{543.187}{for}%
\TextAt{8}{98.618}{543.187}{some}%
\TextAt{13}{129.451}{543.187}{\Character{11}}%
\NativeMathAt{142.861}{543.187}{11.955}{\in}%
\TextAt{18}{156.719}{543.187}{R}%
\TextAt{8}{165.354}{543.187}{.}%
\TextAt{8}{178.322}{543.187}{Let}%
\TextAt{13}{200.798}{543.187}{m}%
\TextAt{21}{211.037}{541.394}{i}%
\TextAt{9}{214.419}{543.187}{\Character{47}}%
\TextAt{13}{220.272}{543.187}{n}%
\TextAt{21}{227.259}{541.394}{i}%
\TextAt{8}{236.045}{543.187}{be}%
\TextAt{8}{253.476}{543.187}{a}%
\TextAt{8}{264.738}{543.187}{sequence}%
\TextAt{8}{314.747}{543.187}{of}%
\TextAt{8}{329.583}{543.187}{rational}%
\TextAt{8}{374.679}{543.187}{numbers}%
\TextAt{8}{423.396}{543.187}{(with}%
\TextAt{13}{456.118}{543.187}{n}%
\TextAt{21}{463.106}{541.394}{i}%
\TextAt{13}{472.375}{543.187}{\Character{62}}%
\TextAt{9}{487.368}{543.187}{0}%
\TextAt{8}{493.221}{543.187}{)}%
\TextAt{8}{503.18}{543.187}{which}%
\TextAt{8}{79.226}{526.141}{converges}%
\TextAt{8}{131.65}{526.141}{to}%
\TextAt{13}{145.96}{526.141}{\Character{11}}%
\TextAt{8}{157.379}{526.141}{from}%
\TextAt{8}{185.019}{526.141}{above.}%
\TextAt{8}{222.746}{526.141}{Then}%
\TextAt{8}{253.303}{526.141}{we}%
\TextAt{8}{270.531}{526.141}{have}%
\TextAt{16}{297.525}{526.141}{\Character{106}}%
\TextAt{13}{300.846}{526.141}{x}%
\TextAt{16}{307.498}{526.141}{\Character{106}}%
\TextAt{17}{310.819}{524.348}{1}%
\TextAt{9}{318.873}{526.141}{\Character{61}}%
\TextAt{16}{331.298}{526.141}{\Character{106}}%
\TextAt{13}{334.619}{526.141}{y}%
\TextAt{16}{340.756}{526.141}{\Character{106}}%
\TextAt{21}{344.076}{530.48}{\Character{11}}%
\TextAt{17}{344.076}{523.186}{1}%
\TextAt{13}{353.334}{526.141}{\Character{60}}%
\TextAt{16}{365.76}{526.141}{\Character{106}}%
\TextAt{13}{369.081}{526.141}{y}%
\TextAt{16}{375.217}{526.141}{\Character{106}}%
\TextAt{21}{378.538}{532.252}{m}%
\TextAt{22}{386.029}{531.037}{i}%
\TextAt{17}{389.191}{532.252}{\Character{47}}%
\TextAt{21}{393.425}{532.252}{n}%
\TextAt{22}{398.563}{531.037}{i}%
\TextAt{17}{378.538}{523.21}{1}%
\TextAt{8}{402.223}{526.141}{,}%
\TextAt{8}{409.372}{526.141}{hence}%
\TextAt{20}{237.538}{510.482}{\Character{12}}%
\TextAt{20}{237.538}{503.308}{\Character{12}}%
\TextAt{20}{237.538}{496.135}{\Character{12}}%
\TextAt{20}{237.538}{488.962}{\Character{12}}%
\TextAt{13}{243.637}{501.234}{x}%
\TextAt{21}{250.289}{505.573}{n}%
\TextAt{22}{255.427}{504.358}{i}%
\Vector{q 0 G 0 g 0.478 w 170.718 -277.865 m 188.005 -277.865 l S Q}%
\TextAt{13}{242.718}{484.946}{y}%
\TextAt{21}{248.855}{488.4}{m}%
\TextAt{22}{256.345}{487.184}{i}%
\TextAt{20}{261.201}{510.482}{\Character{12}}%
\TextAt{20}{261.201}{503.309}{\Character{12}}%
\TextAt{20}{261.201}{496.135}{\Character{12}}%
\TextAt{20}{261.201}{488.962}{\Character{12}}%
\TextAt{17}{265.186}{481.291}{1}%
\TextAt{13}{273.239}{493.147}{\Character{60}}%
\TextAt{9}{285.664}{493.147}{1}%
\NativeArrowAt{294.822}{493.147}{11.955}{19.083}{\Longrightarrow}%
\TextAt{20}{317.226}{510.482}{\Character{12}}%
\TextAt{20}{317.226}{503.309}{\Character{12}}%
\TextAt{20}{317.226}{496.135}{\Character{12}}%
\TextAt{20}{317.226}{488.962}{\Character{12}}%
\TextAt{13}{323.326}{501.234}{x}%
\TextAt{21}{329.978}{505.573}{n}%
\TextAt{22}{335.116}{504.358}{i}%
\Vector{q 0 G 0 g 0.478 w 250.407 -277.865 m 267.694 -277.865 l S Q}%
\TextAt{13}{322.407}{484.946}{y}%
\TextAt{21}{328.544}{488.4}{m}%
\TextAt{22}{336.034}{487.184}{i}%
\TextAt{20}{340.89}{510.482}{\Character{12}}%
\TextAt{20}{340.89}{503.309}{\Character{12}}%
\TextAt{20}{340.89}{496.135}{\Character{12}}%
\TextAt{20}{340.89}{488.962}{\Character{12}}%
\TextAt{17}{344.875}{481.291}{2}%
\TextAt{13}{352.928}{493.147}{\Character{60}}%
\TextAt{9}{365.353}{493.147}{1}%
\TextAt{13}{371.206}{493.147}{\Character{59}}%
\TextAt{8}{79.226}{456.532}{so}%
\TextAt{8}{93.851}{456.532}{that}%
\TextAt{16}{119.474}{456.532}{\Character{106}}%
\TextAt{13}{122.795}{456.532}{x}%
\TextAt{16}{129.447}{456.532}{\Character{106}}%
\TextAt{17}{132.768}{454.739}{2}%
\NativeMathAt{141.255}{456.532}{11.955}{\leq}%
\TextAt{16}{154.298}{456.532}{\Character{106}}%
\TextAt{13}{157.629}{456.532}{y}%
\TextAt{16}{163.765}{456.532}{\Character{106}}%
\TextAt{21}{167.086}{462.642}{m}%
\TextAt{22}{174.577}{461.427}{i}%
\TextAt{17}{177.738}{462.642}{\Character{47}}%
\TextAt{21}{181.973}{462.642}{n}%
\TextAt{22}{187.111}{461.427}{i}%
\TextAt{17}{167.086}{453.601}{2}%
\TextAt{8}{190.771}{456.532}{,}%
\TextAt{8}{198.238}{456.532}{and}%
\TextAt{8}{221.255}{456.532}{thus}%
\TextAt{16}{247.261}{456.532}{\Character{106}}%
\TextAt{13}{250.582}{456.532}{x}%
\TextAt{16}{257.234}{456.532}{\Character{106}}%
\TextAt{17}{260.555}{454.739}{2}%
\NativeMathAt{269.042}{456.532}{11.955}{\leq}%
\TextAt{16}{282.085}{456.532}{\Character{106}}%
\TextAt{13}{285.416}{456.532}{y}%
\TextAt{16}{291.552}{456.532}{\Character{106}}%
\TextAt{21}{294.873}{460.871}{\Character{11}}%
\TextAt{17}{294.873}{453.577}{2}%
\TextAt{8}{300.81}{456.532}{.}%
\TextAt{8}{310.022}{456.532}{Using}%
\TextAt{8}{343.177}{456.532}{a}%
\TextAt{8}{353.187}{456.532}{sequence}%
\TextAt{13}{401.944}{456.532}{m}%
\TextAt{21}{412.183}{454.739}{i}%
\TextAt{9}{415.564}{456.532}{\Character{47}}%
\TextAt{13}{421.417}{456.532}{n}%
\TextAt{21}{428.405}{454.739}{i}%
\TextAt{8}{435.938}{456.532}{which}%
\TextAt{8}{469.679}{456.532}{converges}%
\TextAt{8}{522.357}{456.532}{to}%
\TextAt{13}{79.226}{439.486}{\Character{11}}%
\TextAt{8}{90.852}{439.486}{from}%
\TextAt{8}{118.699}{439.486}{below}%
\TextAt{8}{152.069}{439.486}{(*)}%
\TextAt{8}{171.14}{439.486}{tells}%
\TextAt{8}{196.118}{439.486}{us}%
\TextAt{8}{211.34}{439.486}{that}%
\TextAt{16}{236.915}{439.486}{\Character{106}}%
\TextAt{13}{240.236}{439.486}{x}%
\TextAt{16}{246.888}{439.486}{\Character{106}}%
\TextAt{17}{250.209}{437.693}{2}%
\NativeMathAt{258.614}{439.486}{11.955}{\geq}%
\TextAt{16}{271.573}{439.486}{\Character{106}}%
\TextAt{13}{274.906}{439.486}{y}%
\TextAt{16}{281.043}{439.486}{\Character{106}}%
\TextAt{21}{284.364}{443.825}{\Character{11}}%
\TextAt{17}{284.364}{436.531}{2}%
\TextAt{8}{290.301}{439.486}{.}%
\TextAt{8}{299.369}{439.486}{So}%
\TextAt{8}{315.834}{439.486}{we}%
\TextAt{8}{333.268}{439.486}{have}%
\TextAt{16}{360.469}{439.486}{\Character{106}}%
\TextAt{13}{363.79}{439.486}{x}%
\TextAt{16}{370.442}{439.486}{\Character{106}}%
\TextAt{17}{373.763}{437.693}{2}%
\TextAt{9}{382.168}{439.486}{\Character{61}}%
\TextAt{16}{394.946}{439.486}{\Character{106}}%
\TextAt{13}{398.266}{439.486}{y}%
\TextAt{16}{404.403}{439.486}{\Character{106}}%
\TextAt{21}{407.724}{443.825}{\Character{11}}%
\TextAt{17}{407.724}{436.531}{2}%
\TextAt{8}{413.661}{439.486}{.}%
\TextAt{8}{422.729}{439.486}{For}%
\TextAt{8}{443.905}{439.486}{all}%
\TextAt{13}{460.37}{439.486}{x}%
\NativeMathAt{470.695}{439.486}{11.955}{\in}%
\TextAt{13}{482.338}{439.486}{K}%
\TextAt{8}{493.155}{439.486}{,}%
\TextAt{13}{500.563}{439.486}{x}%
\NativeMathAt{510.887}{439.486}{11.955}{\neq}%
\TextAt{9}{523.655}{439.486}{0}%
\TextAt{8}{529.518}{439.486}{,}%
\TextAt{8}{79.226}{422.44}{we}%
\TextAt{8}{96.454}{422.44}{therefore}%
\TextAt{8}{145.553}{422.44}{get}%
\ErrataAt{179}{401.939}{276}{$\displaystyle \log|x|_1=s\log|x|_2,\qquad s:=\frac{\log|y|_1}{\log|y|_2},$}%
\TextAt{8}{79.226}{374.431}{hence}%
\TextAt{16}{111.732}{374.431}{\Character{106}}%
\TextAt{13}{115.053}{374.431}{x}%
\TextAt{16}{121.705}{374.431}{\Character{106}}%
\TextAt{17}{125.026}{372.638}{1}%
\TextAt{9}{133.079}{374.431}{\Character{61}}%
\TextAt{16}{145.505}{374.431}{\Character{106}}%
\TextAt{13}{148.826}{374.431}{x}%
\TextAt{16}{155.478}{374.431}{\Character{106}}%
\TextAt{21}{158.799}{378.77}{s}%
\TextAt{17}{158.799}{371.476}{2}%
\TextAt{8}{163.531}{374.431}{.}%
\TextAt{8}{171.979}{374.431}{But}%
\TextAt{16}{195.22}{374.431}{\Character{106}}%
\TextAt{13}{198.541}{374.431}{y}%
\TextAt{16}{204.678}{374.431}{\Character{106}}%
\TextAt{17}{207.999}{372.638}{1}%
\TextAt{13}{216.052}{374.431}{\Character{62}}%
\TextAt{9}{228.477}{374.431}{1}%
\TextAt{8}{238.228}{374.431}{implies}%
\TextAt{16}{277.955}{374.431}{\Character{106}}%
\TextAt{13}{281.275}{374.431}{y}%
\TextAt{16}{287.412}{374.431}{\Character{106}}%
\TextAt{17}{290.733}{372.638}{2}%
\TextAt{13}{298.786}{374.431}{\Character{62}}%
\TextAt{9}{311.212}{374.431}{1}%
\TextAt{8}{317.065}{374.431}{,}%
\TextAt{8}{324.214}{374.431}{hence}%
\TextAt{13}{356.72}{374.431}{s}%
\TextAt{13}{365.543}{374.431}{\Character{62}}%
\TextAt{9}{377.98}{374.431}{0}%
\TextAt{8}{383.833}{374.431}{.}%
\Vector{q 0 G 0 g 0.399 w 452.334 -399.569 m 452.334 -391.499 l S Q}%
\Vector{q 0 G 0 g 0.399 w 452.534 -391.698 m 459.707 -391.698 l S Q}%
\Vector{q 0 G 0 g 0.399 w 452.534 -399.369 m 459.707 -399.369 l S Q}%
\Vector{q 0 G 0 g 0.399 w 459.906 -399.569 m 459.906 -391.499 l S Q}%
\TextAt{8}{97.159}{345.51}{The}%
\TextAt{8}{120.79}{345.51}{proof}%
\TextAt{8}{150.937}{345.51}{shows}%
\TextAt{8}{184.132}{345.51}{that}%
\TextAt{8}{209.079}{345.51}{the}%
\TextAt{8}{228.813}{345.51}{equivalence}%
\TextAt{8}{290.175}{345.51}{of}%
\TextAt{16}{303.082}{345.51}{\Character{106}}%
\TextAt{16}{309.861}{345.51}{\Character{106}}%
\TextAt{17}{313.199}{343.717}{1}%
\TextAt{8}{321.406}{345.51}{and}%
\TextAt{16}{343.746}{345.51}{\Character{106}}%
\TextAt{16}{350.525}{345.51}{\Character{106}}%
\TextAt{17}{353.863}{343.717}{2}%
\TextAt{8}{362.07}{345.51}{is}%
\TextAt{8}{373.411}{345.51}{also}%
\TextAt{8}{396.469}{345.51}{equivalent}%
\TextAt{8}{451.662}{345.51}{to}%
\TextAt{8}{465.55}{345.51}{the}%
\TextAt{8}{485.284}{345.51}{condition}%
\TextAt{16}{251.892}{316.79}{\Character{106}}%
\TextAt{13}{255.213}{316.79}{x}%
\TextAt{16}{261.865}{316.79}{\Character{106}}%
\TextAt{17}{265.186}{314.997}{1}%
\TextAt{13}{273.239}{316.79}{\Character{60}}%
\TextAt{9}{285.664}{316.79}{1}%
\NativeArrowAt{294.822}{316.79}{11.955}{19.083}{\Longrightarrow}%
\TextAt{16}{317.217}{316.79}{\Character{106}}%
\TextAt{13}{320.547}{316.79}{x}%
\TextAt{16}{327.199}{316.79}{\Character{106}}%
\TextAt{17}{330.52}{314.997}{2}%
\TextAt{13}{338.574}{316.79}{\Character{60}}%
\TextAt{9}{350.999}{316.79}{1}%
\TextAt{13}{356.852}{316.79}{\Character{58}}%
\TextAt{8}{79.226}{288.07}{We}%
\TextAt{8}{98.198}{288.07}{use}%
\TextAt{8}{117.241}{288.07}{this}%
\TextAt{8}{138.89}{288.07}{for}%
\TextAt{8}{155.602}{288.07}{the}%
\TextAt{8}{174.586}{288.07}{proof}%
\TextAt{8}{203.982}{288.07}{of}%
\TextAt{8}{216.139}{288.07}{the}%
\TextAt{8}{235.122}{288.07}{following}%
\TextAt{8}{283.384}{288.07}{approximation}%
\TextAt{8}{360.278}{288.07}{theorem,}%
\TextAt{8}{408.117}{288.07}{which}%
\TextAt{8}{440.431}{288.07}{may}%
\TextAt{8}{464.626}{288.07}{be}%
\TextAt{8}{479.378}{288.07}{considered}%
\TextAt{8}{79.226}{271.024}{a}%
\TextAt{8}{88.982}{271.024}{variant}%
\TextAt{8}{128.673}{271.024}{of}%
\TextAt{8}{142.003}{271.024}{the}%
\TextAt{8}{162.159}{271.024}{Chinese}%
\TextAt{8}{205.784}{271.024}{remainder}%
\TextAt{8}{261.064}{271.024}{theorem.}%
\TextAt{12}{79.226}{245.091}{Approximation}%
\TextAt{12}{173.21}{245.091}{Theorem}%
\TextAt{12}{230.982}{245.091}{(3.4).}%
\TextAt{8}{266.956}{245.091}{Let}%
\TextAt{16}{288.583}{245.091}{\Character{106}}%
\TextAt{16}{296.45}{245.091}{\Character{106}}%
\TextAt{17}{299.781}{243.298}{1}%
\TextAt{13}{304.513}{245.091}{\Character{59}}%
\TextAt{13}{309.737}{245.091}{\Character{58}}%
\TextAt{13}{314.974}{245.091}{\Character{58}}%
\TextAt{13}{320.21}{245.091}{\Character{58}}%
\TextAt{13}{325.435}{245.091}{\Character{59}}%
\TextAt{16}{330.734}{245.091}{\Character{106}}%
\TextAt{16}{338.601}{245.091}{\Character{106}}%
\TextAt{21}{341.932}{243.298}{n}%
\TextAt{8}{352.124}{245.091}{be}%
\TextAt{8}{368.707}{245.091}{pairwise}%
\TextAt{8}{414.95}{245.091}{inequivalent}%
\TextAt{8}{480.98}{245.091}{valuations}%
\TextAt{8}{79.226}{228.045}{of}%
\TextAt{8}{92.177}{228.045}{the}%
\TextAt{8}{111.955}{228.045}{field}%
\TextAt{13}{136.933}{228.045}{K}%
\TextAt{8}{151.27}{228.045}{and}%
\TextAt{8}{173.654}{228.045}{let}%
\TextAt{13}{190.18}{228.045}{a}%
\TextAt{17}{196.325}{226.252}{1}%
\TextAt{13}{201.057}{228.045}{\Character{59}}%
\TextAt{13}{206.281}{228.045}{\Character{58}}%
\TextAt{13}{211.518}{228.045}{\Character{58}}%
\TextAt{13}{216.742}{228.045}{\Character{58}}%
\TextAt{13}{221.979}{228.045}{\Character{59}}%
\TextAt{13}{227.215}{228.045}{a}%
\TextAt{21}{233.423}{226.252}{n}%
\NativeMathAt{242.38}{228.045}{11.955}{\in}%
\TextAt{13}{253.671}{228.045}{K}%
\TextAt{8}{268.007}{228.045}{be}%
\TextAt{8}{283.553}{228.045}{given}%
\TextAt{8}{313.743}{228.045}{elements.}%
\TextAt{8}{366.025}{228.045}{Then}%
\TextAt{8}{396.203}{228.045}{for}%
\TextAt{8}{413.709}{228.045}{every}%
\TextAt{13}{444.223}{228.045}{\Character{34}}%
\TextAt{13}{453.01}{228.045}{\Character{62}}%
\TextAt{9}{465.447}{228.045}{0}%
\TextAt{8}{474.819}{228.045}{there}%
\TextAt{8}{504.352}{228.045}{exists}%
\TextAt{8}{79.226}{210.999}{an}%
\TextAt{13}{95.485}{210.999}{x}%
\NativeMathAt{105.458}{210.999}{11.955}{\in}%
\TextAt{13}{116.749}{210.999}{K}%
\TextAt{8}{131.464}{210.999}{such}%
\TextAt{8}{157.861}{210.999}{that}%
\TextAt{16}{214.926}{193.954}{\Character{106}}%
\TextAt{13}{218.246}{193.954}{x}%
\TextAt{16}{227.555}{193.954}{\Character{0}}%
\TextAt{13}{239.51}{193.954}{a}%
\TextAt{21}{245.655}{192.16}{i}%
\TextAt{16}{249.037}{193.954}{\Character{106}}%
\TextAt{21}{252.357}{192.16}{i}%
\TextAt{13}{259.06}{193.954}{\Character{60}}%
\TextAt{13}{271.47}{193.954}{\Character{34}}%
\TextAt{8}{300.873}{193.954}{for}%
\TextAt{8}{318.758}{193.954}{all}%
\TextAt{13}{335.017}{193.954}{i}%
\TextAt{9}{342.332}{193.954}{\Character{61}}%
\TextAt{9}{354.742}{193.954}{1}%
\TextAt{13}{360.61}{193.954}{\Character{59}}%
\TextAt{13}{365.834}{193.954}{\Character{58}}%
\TextAt{13}{371.071}{193.954}{\Character{58}}%
\TextAt{13}{376.295}{193.954}{\Character{58}}%
\TextAt{13}{381.532}{193.954}{\Character{59}}%
\TextAt{13}{386.768}{193.954}{n\Character{58}}%
\TextAt{19}{79.226}{152.78}{Proof.}%
\TextAt{8}{116}{152.78}{By}%
\TextAt{8}{134.439}{152.78}{the}%
\TextAt{8}{154.672}{152.78}{above}%
\TextAt{8}{187.924}{152.78}{remark,}%
\TextAt{8}{231.273}{152.78}{since}%
\TextAt{16}{260.017}{152.78}{\Character{106}}%
\TextAt{16}{267.298}{152.78}{\Character{106}}%
\TextAt{17}{270.633}{150.986}{1}%
\TextAt{8}{279.338}{152.78}{and}%
\TextAt{16}{302.177}{152.78}{\Character{106}}%
\TextAt{16}{309.458}{152.78}{\Character{106}}%
\TextAt{21}{312.792}{150.986}{n}%
\TextAt{8}{322.402}{152.78}{are}%
\TextAt{8}{341.989}{152.78}{inequivalent,}%
\TextAt{8}{410.707}{152.78}{there}%
\TextAt{8}{440.695}{152.78}{exists}%
\TextAt{13}{473.086}{152.78}{\Character{11}}%
\NativeMathAt{484.058}{152.78}{11.955}{\in}%
\TextAt{13}{495.479}{152.78}{K}%
\TextAt{8}{510.27}{152.78}{such}%
\TextAt{8}{79.226}{135.734}{that}%
\TextAt{16}{104.101}{135.734}{\Character{106}}%
\TextAt{13}{107.421}{135.734}{\Character{11}}%
\TextAt{16}{114.943}{135.734}{\Character{106}}%
\TextAt{17}{118.264}{133.941}{1}%
\TextAt{13}{126.317}{135.734}{\Character{60}}%
\TextAt{9}{138.743}{135.734}{1}%
\TextAt{8}{147.999}{135.734}{and}%
\TextAt{16}{170.267}{135.734}{\Character{106}}%
\TextAt{13}{173.588}{135.734}{\Character{11}}%
\TextAt{16}{181.109}{135.734}{\Character{106}}%
\TextAt{21}{184.43}{133.941}{n}%
\NativeMathAt{193.387}{135.734}{11.955}{\geq}%
\TextAt{9}{206.007}{135.734}{1}%
\TextAt{8}{211.86}{135.734}{.}%
\TextAt{8}{220.143}{135.734}{By}%
\TextAt{8}{238.012}{135.734}{the}%
\TextAt{8}{257.674}{135.734}{same}%
\TextAt{8}{286.506}{135.734}{token,}%
\TextAt{8}{321.234}{135.734}{there}%
\TextAt{8}{350.652}{135.734}{exists}%
\TextAt{13}{382.472}{135.734}{\Character{12}}%
\NativeMathAt{393.064}{135.734}{11.955}{\in}%
\TextAt{13}{404.355}{135.734}{K}%
\TextAt{8}{418.576}{135.734}{such}%
\TextAt{8}{444.478}{135.734}{that}%
\TextAt{16}{469.353}{135.734}{\Character{106}}%
\TextAt{13}{472.674}{135.734}{\Character{12}}%
\TextAt{16}{479.945}{135.734}{\Character{106}}%
\TextAt{21}{483.266}{133.941}{n}%
\TextAt{13}{492.223}{135.734}{\Character{60}}%
\TextAt{9}{504.648}{135.734}{1}%
\TextAt{8}{513.904}{135.734}{and}%
\TextAt{16}{79.226}{118.688}{\Character{106}}%
\TextAt{13}{82.547}{118.688}{\Character{12}}%
\TextAt{16}{89.818}{118.688}{\Character{106}}%
\TextAt{17}{93.139}{116.895}{1}%
\NativeMathAt{101.632}{118.688}{11.955}{\geq}%
\TextAt{9}{114.69}{118.688}{1}%
\TextAt{8}{120.543}{118.688}{.}%
\TextAt{8}{129.765}{118.688}{Putting}%
\TextAt{13}{173.109}{118.688}{y}%
\TextAt{9}{183.006}{118.688}{\Character{61}}%
\TextAt{13}{195.871}{118.688}{\Character{12}}%
\TextAt{9}{203.142}{118.688}{\Character{47}}%
\TextAt{13}{208.995}{118.688}{\Character{11}}%
\TextAt{8}{216.517}{118.688}{,}%
\TextAt{8}{223.988}{118.688}{one}%
\TextAt{8}{245.706}{118.688}{finds}%
\TextAt{16}{273.986}{118.688}{\Character{106}}%
\TextAt{13}{277.307}{118.688}{y}%
\TextAt{16}{283.444}{118.688}{\Character{106}}%
\TextAt{17}{286.765}{116.895}{1}%
\TextAt{13}{295.257}{118.688}{\Character{62}}%
\TextAt{9}{308.122}{118.688}{1}%
\TextAt{8}{318.13}{118.688}{and}%
\TextAt{16}{341.151}{118.688}{\Character{106}}%
\TextAt{13}{344.472}{118.688}{y}%
\TextAt{16}{350.608}{118.688}{\Character{106}}%
\TextAt{21}{353.929}{116.895}{n}%
\TextAt{13}{363.326}{118.688}{\Character{60}}%
\TextAt{9}{376.19}{118.688}{1}%
\TextAt{8}{382.043}{118.688}{.}%
\TextAt{8}{391.265}{118.688}{We}%
\TextAt{8}{411.667}{118.688}{now}%
\TextAt{8}{436.314}{118.688}{prove}%
\TextAt{8}{468.121}{118.688}{by}%
\TextAt{8}{484.638}{118.688}{induction}%
\TextAt{8}{79.226}{101.642}{on}%
\TextAt{13}{95.485}{101.642}{n}%
\TextAt{8}{106.37}{101.642}{that}%
\TextAt{8}{131.739}{101.642}{there}%
\TextAt{8}{161.651}{101.642}{exists}%
\TextAt{13}{193.966}{101.642}{z}%
\NativeMathAt{203.257}{101.642}{11.955}{\in}%
\TextAt{13}{214.548}{101.642}{K}%
\TextAt{8}{229.263}{101.642}{such}%
\TextAt{8}{255.66}{101.642}{that}%
\TextAt{16}{187.178}{72.921}{\Character{106}}%
\TextAt{13}{190.499}{72.921}{z}%
\TextAt{16}{196.47}{72.921}{\Character{106}}%
\TextAt{17}{199.791}{71.128}{1}%
\TextAt{13}{207.844}{72.921}{\Character{62}}%
\TextAt{9}{220.269}{72.921}{1}%
\TextAt{8}{250.033}{72.921}{and}%
\TextAt{16}{292.808}{72.921}{\Character{106}}%
\TextAt{13}{296.129}{72.921}{z}%
\TextAt{16}{302.1}{72.921}{\Character{106}}%
\TextAt{21}{305.421}{71.128}{j}%
\TextAt{13}{313.124}{72.921}{\Character{60}}%
\TextAt{9}{325.549}{72.921}{1}%
\TextAt{8}{343.357}{72.921}{for}%
\TextAt{13}{361.242}{72.921}{j}%
\TextAt{9}{370.079}{72.921}{\Character{61}}%
\TextAt{9}{382.488}{72.921}{2}%
\TextAt{13}{388.357}{72.921}{\Character{59}}%
\TextAt{13}{393.581}{72.921}{\Character{58}}%
\TextAt{13}{398.818}{72.921}{\Character{58}}%
\TextAt{13}{404.054}{72.921}{\Character{58}}%
\TextAt{13}{409.279}{72.921}{\Character{59}}%
\TextAt{13}{414.515}{72.921}{n\Character{58}}%
\TextAt{8}{297.211}{43.033}{107}%
\Destination{522.36}{600.685}{nav.AMS.14}%
\Destination{251.89}{600.685}{nav.equation.2.3.1}%
\Destination{78.23}{809.815}{nav.page.107}%
\Destination{79.36}{776.155}{nav.print.117}%
\Destination{79.36}{167.755}{nav.print.118}%
\Destination{79.36}{775.155}{nav.number.2.3.3}%
\Destination{79.226}{418.939}{errata.74}%
\Destination{79.226}{259.5}{nav.number.2.3.4}%
}

% Page 120 (body page 108)
\Page{%
\TextAt{2}{79.226}{796.943}{II.}%
\TextAt{2}{94.073}{796.943}{The}%
\TextAt{2}{116.492}{796.943}{Theory}%
\TextAt{2}{154.401}{796.943}{of}%
\TextAt{2}{166.826}{796.943}{Valuations}%
\Vector{q 0 G 0 g 0.299 w 7.226 18.729 m 460.769 18.729 l S Q}%
\TextAt{8}{79.226}{761.124}{We}%
\TextAt{8}{99.371}{761.124}{have}%
\TextAt{8}{126.365}{761.124}{just}%
\TextAt{8}{149.511}{761.124}{done}%
\TextAt{8}{177.474}{761.124}{this}%
\TextAt{8}{200.296}{761.124}{for}%
\TextAt{13}{218.181}{761.124}{n}%
\TextAt{9}{228.49}{761.124}{\Character{61}}%
\TextAt{9}{240.9}{761.124}{2}%
\TextAt{8}{246.768}{761.124}{.}%
\TextAt{8}{255.216}{761.124}{Assume}%
\TextAt{8}{298.578}{761.124}{we}%
\TextAt{8}{315.805}{761.124}{have}%
\TextAt{8}{342.8}{761.124}{found}%
\TextAt{13}{375.641}{761.124}{z}%
\NativeMathAt{384.932}{761.124}{11.955}{\in}%
\TextAt{13}{396.223}{761.124}{K}%
\TextAt{8}{410.938}{761.124}{satisfying}%
\TextAt{16}{176.946}{732.259}{\Character{106}}%
\TextAt{13}{180.267}{732.259}{z}%
\TextAt{16}{186.238}{732.259}{\Character{106}}%
\TextAt{17}{189.558}{730.466}{1}%
\TextAt{13}{197.612}{732.259}{\Character{62}}%
\TextAt{9}{210.037}{732.259}{1}%
\TextAt{8}{239.8}{732.259}{and}%
\TextAt{16}{282.576}{732.259}{\Character{106}}%
\TextAt{13}{285.897}{732.259}{z}%
\TextAt{16}{291.867}{732.259}{\Character{106}}%
\TextAt{21}{295.188}{730.466}{j}%
\TextAt{13}{302.891}{732.259}{\Character{60}}%
\TextAt{9}{315.317}{732.259}{1}%
\TextAt{8}{333.125}{732.259}{for}%
\TextAt{13}{351.01}{732.259}{j}%
\TextAt{9}{359.846}{732.259}{\Character{61}}%
\TextAt{9}{372.256}{732.259}{2}%
\TextAt{13}{378.125}{732.259}{\Character{59}}%
\TextAt{13}{383.349}{732.259}{\Character{58}}%
\TextAt{13}{388.586}{732.259}{\Character{58}}%
\TextAt{13}{393.81}{732.259}{\Character{58}}%
\TextAt{13}{399.047}{732.259}{\Character{59}}%
\TextAt{13}{404.283}{732.259}{n}%
\TextAt{16}{413.99}{732.259}{\Character{0}}%
\TextAt{9}{425.945}{732.259}{1}%
\TextAt{13}{431.798}{732.259}{\Character{58}}%
\TextAt{8}{79.226}{703.394}{If}%
\TextAt{16}{92.514}{703.394}{\Character{106}}%
\TextAt{13}{95.835}{703.394}{z}%
\TextAt{16}{101.806}{703.394}{\Character{106}}%
\TextAt{21}{105.127}{701.601}{n}%
\NativeMathAt{116.803}{703.394}{11.955}{\leq}%
\TextAt{9}{132.142}{703.394}{1}%
\TextAt{8}{137.995}{703.394}{,}%
\TextAt{8}{147.139}{703.394}{then}%
\TextAt{13}{175.395}{703.394}{z}%
\TextAt{21}{181.366}{707.733}{m}%
\TextAt{13}{189.354}{703.394}{y}%
\TextAt{8}{200.984}{703.394}{will}%
\TextAt{8}{224.685}{703.394}{do,}%
\TextAt{8}{246.191}{703.394}{for}%
\TextAt{13}{265.672}{703.394}{m}%
\TextAt{8}{281.404}{703.394}{large.}%
\TextAt{8}{319.364}{703.394}{If}%
\TextAt{8}{332.652}{703.394}{however}%
\TextAt{16}{379.127}{703.394}{\Character{106}}%
\TextAt{13}{382.448}{703.394}{z}%
\TextAt{16}{388.419}{703.394}{\Character{106}}%
\TextAt{21}{391.74}{701.601}{n}%
\TextAt{13}{403.416}{703.394}{\Character{62}}%
\TextAt{9}{418.561}{703.394}{1}%
\TextAt{8}{424.414}{703.394}{,}%
\TextAt{8}{433.558}{703.394}{the}%
\TextAt{8}{455.311}{703.394}{sequence}%
\TextAt{13}{505.409}{703.394}{t}%
\TextAt{21}{509.636}{701.601}{m}%
\TextAt{9}{523.665}{703.394}{\Character{61}}%
\TextAt{13}{79.226}{686.348}{z}%
\TextAt{21}{85.197}{690.687}{m}%
\TextAt{9}{93.185}{686.348}{\Character{47}\Character{40}1}%
\TextAt{9}{112.851}{686.348}{\Character{43}}%
\TextAt{13}{125.411}{686.348}{z}%
\TextAt{21}{131.382}{690.687}{m}%
\TextAt{9}{139.37}{686.348}{\Character{41}}%
\TextAt{8}{148.957}{686.348}{will}%
\TextAt{8}{172.198}{686.348}{converge}%
\TextAt{8}{221.144}{686.348}{to}%
\TextAt{8}{236.591}{686.348}{1}%
\TextAt{8}{247.483}{686.348}{with}%
\TextAt{8}{275.279}{686.348}{respect}%
\TextAt{8}{316.465}{686.348}{to}%
\TextAt{16}{331.912}{686.348}{\Character{106}}%
\TextAt{16}{340.257}{686.348}{\Character{106}}%
\TextAt{17}{343.588}{684.555}{1}%
\TextAt{8}{353.354}{686.348}{and}%
\TextAt{16}{377.254}{686.348}{\Character{106}}%
\TextAt{16}{385.599}{686.348}{\Character{106}}%
\TextAt{21}{388.929}{684.555}{n}%
\TextAt{8}{394.566}{686.348}{,}%
\TextAt{8}{403.135}{686.348}{and}%
\TextAt{8}{427.035}{686.348}{to}%
\TextAt{8}{442.482}{686.348}{0}%
\TextAt{8}{453.374}{686.348}{with}%
\TextAt{8}{481.17}{686.348}{respect}%
\TextAt{8}{522.357}{686.348}{to}%
\TextAt{16}{79.226}{669.302}{\Character{106}}%
\TextAt{16}{86.567}{669.302}{\Character{106}}%
\TextAt{17}{89.897}{667.509}{2}%
\TextAt{13}{94.629}{669.302}{\Character{59}}%
\TextAt{13}{99.853}{669.302}{\Character{58}}%
\TextAt{13}{105.09}{669.302}{\Character{58}}%
\TextAt{13}{110.314}{669.302}{\Character{58}}%
\TextAt{13}{115.551}{669.302}{\Character{59}}%
\TextAt{16}{120.85}{669.302}{\Character{106}}%
\TextAt{16}{128.179}{669.302}{\Character{106}}%
\TextAt{21}{131.52}{667.509}{n}%
\TextAt{3}{136.659}{667.509}{\Character{0}}%
\TextAt{17}{143.245}{667.509}{1}%
\TextAt{8}{147.977}{669.302}{.}%
\TextAt{8}{156.82}{669.302}{Hence,}%
\TextAt{8}{195.014}{669.302}{for}%
\TextAt{13}{213.03}{669.302}{m}%
\TextAt{8}{227.298}{669.302}{large,}%
\TextAt{13}{259.335}{669.302}{t}%
\TextAt{21}{263.562}{667.509}{m}%
\TextAt{13}{271.551}{669.302}{y}%
\TextAt{8}{281.716}{669.302}{will}%
\TextAt{8}{303.952}{669.302}{suffice.}%
\TextAt{8}{344.07}{669.302}{The}%
\TextAt{8}{368.255}{669.302}{sequence}%
\TextAt{13}{416.888}{669.302}{z}%
\TextAt{21}{422.859}{673.641}{m}%
\TextAt{9}{430.848}{669.302}{\Character{47}\Character{40}1}%
\TextAt{9}{449.821}{669.302}{\Character{43}}%
\TextAt{13}{461.703}{669.302}{z}%
\TextAt{21}{467.674}{673.641}{m}%
\TextAt{9}{475.662}{669.302}{\Character{41}}%
\TextAt{8}{484.243}{669.302}{converges}%
\TextAt{8}{79.226}{652.256}{to}%
\TextAt{8}{92.82}{652.256}{1}%
\TextAt{8}{101.859}{652.256}{with}%
\TextAt{8}{127.802}{652.256}{respect}%
\TextAt{8}{167.135}{652.256}{to}%
\TextAt{16}{180.729}{652.256}{\Character{106}}%
\TextAt{16}{187.221}{652.256}{\Character{106}}%
\TextAt{17}{190.552}{650.463}{1}%
\TextAt{8}{198.465}{652.256}{and}%
\TextAt{8}{220.511}{652.256}{to}%
\TextAt{8}{234.105}{652.256}{0}%
\TextAt{8}{243.144}{652.256}{with}%
\TextAt{8}{269.087}{652.256}{respect}%
\TextAt{8}{308.42}{652.256}{to}%
\TextAt{16}{322.014}{652.256}{\Character{106}}%
\TextAt{16}{328.506}{652.256}{\Character{106}}%
\TextAt{17}{331.837}{650.463}{2}%
\TextAt{13}{336.569}{652.256}{\Character{59}}%
\TextAt{13}{341.793}{652.256}{\Character{58}}%
\TextAt{13}{347.03}{652.256}{\Character{58}}%
\TextAt{13}{352.254}{652.256}{\Character{58}}%
\TextAt{13}{357.491}{652.256}{\Character{59}}%
\TextAt{16}{362.79}{652.256}{\Character{106}}%
\TextAt{16}{369.282}{652.256}{\Character{106}}%
\TextAt{21}{372.612}{650.463}{n}%
\TextAt{8}{378.249}{652.256}{.}%
\TextAt{8}{386.458}{652.256}{For}%
\TextAt{8}{406.711}{652.256}{every}%
\TextAt{13}{436.886}{652.256}{i}%
\TextAt{8}{444.061}{652.256}{we}%
\TextAt{8}{460.571}{652.256}{may}%
\TextAt{8}{485.224}{652.256}{construct}%
\TextAt{8}{79.226}{635.21}{in}%
\TextAt{8}{92.071}{635.21}{this}%
\TextAt{8}{114.085}{635.21}{way}%
\TextAt{8}{137.019}{635.21}{a}%
\TextAt{13}{145.966}{635.21}{z}%
\TextAt{21}{151.404}{633.417}{i}%
\TextAt{8}{157.874}{635.21}{which}%
\TextAt{8}{190.552}{635.21}{is}%
\TextAt{8}{201.508}{635.21}{very}%
\TextAt{8}{226.391}{635.21}{close}%
\TextAt{8}{253.606}{635.21}{to}%
\TextAt{8}{267.108}{635.21}{1}%
\TextAt{8}{276.055}{635.21}{with}%
\TextAt{8}{301.906}{635.21}{respect}%
\TextAt{8}{341.148}{635.21}{to}%
\TextAt{16}{354.65}{635.21}{\Character{106}}%
\TextAt{16}{361.046}{635.21}{\Character{106}}%
\TextAt{21}{364.381}{633.417}{i}%
\TextAt{8}{367.762}{635.21}{,}%
\TextAt{8}{374.264}{635.21}{and}%
\TextAt{8}{396.219}{635.21}{very}%
\TextAt{8}{421.102}{635.21}{close}%
\TextAt{8}{448.316}{635.21}{to}%
\TextAt{8}{461.818}{635.21}{0}%
\TextAt{8}{470.765}{635.21}{with}%
\TextAt{8}{496.617}{635.21}{respect}%
\TextAt{8}{79.226}{618.164}{to}%
\TextAt{16}{93.536}{618.164}{\Character{106}}%
\TextAt{16}{100.745}{618.164}{\Character{106}}%
\TextAt{21}{104.076}{616.371}{j}%
\TextAt{8}{112.355}{618.164}{for}%
\TextAt{13}{130.24}{618.164}{j}%
\NativeMathAt{139.076}{618.164}{11.955}{\neq}%
\TextAt{13}{151.502}{618.164}{i}%
\TextAt{8}{155.495}{618.164}{.}%
\TextAt{8}{163.944}{618.164}{The}%
\TextAt{8}{187.997}{618.164}{element}%
\TextAt{13}{251.455}{601.118}{x}%
\TextAt{9}{261.428}{601.118}{\Character{61}}%
\TextAt{13}{273.854}{601.118}{a}%
\TextAt{17}{279.999}{599.325}{1}%
\TextAt{13}{284.731}{601.118}{z}%
\TextAt{17}{290.169}{599.325}{1}%
\TextAt{9}{297.558}{601.118}{\Character{43}}%
\TextAt{16}{309.319}{601.118}{\Character{1}}%
\TextAt{16}{314.615}{601.118}{\Character{1}}%
\TextAt{16}{319.923}{601.118}{\Character{1}}%
\TextAt{9}{325.924}{601.118}{\Character{43}}%
\TextAt{13}{337.685}{601.118}{a}%
\TextAt{21}{343.83}{599.325}{n}%
\TextAt{13}{349.466}{601.118}{z}%
\TextAt{21}{354.904}{599.325}{n}%
\TextAt{8}{79.226}{577.655}{then}%
\TextAt{8}{105.886}{577.655}{satisfies}%
\TextAt{8}{148.997}{577.655}{the}%
\TextAt{8}{169.153}{577.655}{statement}%
\TextAt{8}{223.525}{577.655}{of}%
\TextAt{8}{236.855}{577.655}{the}%
\TextAt{8}{257.011}{577.655}{approximation}%
\TextAt{8}{335.079}{577.655}{theorem.}%
\Vector{q 0 G 0 g 0.399 w 452.334 -196.345 m 452.334 -188.275 l S Q}%
\Vector{q 0 G 0 g 0.399 w 452.534 -188.474 m 459.707 -188.474 l S Q}%
\Vector{q 0 G 0 g 0.399 w 452.534 -196.146 m 459.707 -196.146 l S Q}%
\Vector{q 0 G 0 g 0.399 w 459.906 -196.345 m 459.906 -188.275 l S Q}%
\TextAt{12}{79.226}{542.754}{Definition}%
\TextAt{12}{141.914}{542.754}{(3.5).}%
\TextAt{8}{177.142}{542.754}{The}%
\TextAt{8}{201.109}{542.754}{valuation}%
\TextAt{16}{252.094}{542.754}{\Character{106}}%
\TextAt{16}{259.207}{542.754}{\Character{106}}%
\TextAt{8}{266.356}{542.754}{is}%
\TextAt{8}{278.033}{542.754}{called}%
\TextAt{12}{311.109}{542.754}{nonarchimedean}%
\TextAt{8}{410.417}{542.754}{if}%
\TextAt{16}{421.053}{542.754}{\Character{106}}%
\TextAt{13}{424.374}{542.754}{n}%
\TextAt{16}{431.362}{542.754}{\Character{106}}%
\TextAt{8}{438.493}{542.754}{stays}%
\TextAt{8}{467.803}{542.754}{bounded,}%
\TextAt{8}{518.782}{542.754}{for}%
\TextAt{8}{79.226}{525.709}{all}%
\TextAt{13}{95.485}{525.709}{n}%
\NativeMathAt{105.794}{525.708}{11.955}{\in}%
\TextAt{18}{117.085}{525.709}{N}%
\TextAt{8}{125.719}{525.709}{.}%
\TextAt{8}{134.167}{525.709}{Otherwise}%
\TextAt{8}{189.508}{525.709}{it}%
\TextAt{8}{201.212}{525.709}{is}%
\TextAt{8}{212.976}{525.709}{called}%
\TextAt{12}{246.139}{525.709}{archimedean}%
\TextAt{8}{319.963}{525.709}{.}%
\TextAt{12}{79.226}{493.797}{Proposition}%
\TextAt{12}{153.269}{493.797}{(3.6).}%
\TextAt{8}{189.743}{493.797}{The}%
\TextAt{8}{214.954}{493.797}{valuation}%
\TextAt{16}{267.185}{493.797}{\Character{106}}%
\TextAt{16}{275.542}{493.797}{\Character{106}}%
\TextAt{8}{283.937}{493.797}{is}%
\TextAt{8}{296.859}{493.797}{nonarchimedean}%
\TextAt{8}{384.847}{493.797}{if}%
\TextAt{8}{396.729}{493.797}{and}%
\TextAt{8}{420.649}{493.797}{only}%
\TextAt{8}{447.499}{493.797}{if}%
\TextAt{8}{459.38}{493.797}{it}%
\TextAt{8}{472.242}{493.797}{satisfies}%
\TextAt{8}{516.51}{493.797}{the}%
\TextAt{12}{79.226}{476.751}{strong}%
\TextAt{12}{120.687}{476.751}{triangle}%
\TextAt{12}{170.265}{476.751}{inequality}%
\TextAt{16}{246.945}{447.886}{\Character{106}}%
\TextAt{13}{250.266}{447.886}{x}%
\TextAt{9}{259.574}{447.886}{\Character{43}}%
\TextAt{13}{271.336}{447.886}{y}%
\TextAt{16}{277.472}{447.886}{\Character{106}}%
\NativeMathAt{284.095}{447.886}{11.955}{\leq}%
\TextAt{8}{296.733}{447.886}{max}%
\TextAt{16}{318.528}{447.886}{\Character{102}\Character{106}}%
\TextAt{13}{327.826}{447.886}{x}%
\TextAt{16}{334.478}{447.886}{\Character{106}}%
\TextAt{13}{337.799}{447.886}{\Character{59}}%
\TextAt{16}{343.043}{447.886}{\Character{106}}%
\TextAt{13}{346.364}{447.886}{y}%
\TextAt{16}{352.501}{447.886}{\Character{106}\Character{103}}%
\TextAt{13}{361.799}{447.886}{\Character{58}}%
\TextAt{19}{79.226}{401.166}{Proof.}%
\TextAt{8}{116}{401.166}{If}%
\TextAt{8}{127.692}{401.166}{the}%
\TextAt{8}{147.849}{401.166}{strong}%
\TextAt{8}{183.69}{401.166}{triangle}%
\TextAt{8}{226.621}{401.166}{inequality}%
\TextAt{8}{280.934}{401.166}{holds,}%
\TextAt{8}{314.815}{401.166}{then}%
\TextAt{8}{341.475}{401.166}{one}%
\TextAt{8}{362.934}{401.166}{has}%
\TextAt{16}{248.222}{372.302}{\Character{106}}%
\TextAt{13}{251.543}{372.302}{n}%
\TextAt{16}{258.53}{372.302}{\Character{106}}%
\TextAt{9}{265.172}{372.302}{\Character{61}}%
\TextAt{16}{277.598}{372.302}{\Character{106}}%
\TextAt{9}{280.918}{372.302}{1}%
\TextAt{9}{289.418}{372.302}{\Character{43}}%
\TextAt{16}{301.189}{372.302}{\Character{1}}%
\TextAt{16}{306.485}{372.302}{\Character{1}}%
\TextAt{16}{311.793}{372.302}{\Character{1}}%
\TextAt{9}{317.794}{372.302}{\Character{43}}%
\TextAt{9}{329.546}{372.302}{1}%
\TextAt{16}{335.408}{372.302}{\Character{106}}%
\NativeMathAt{342.031}{372.302}{11.955}{\leq}%
\TextAt{9}{354.669}{372.302}{1}%
\TextAt{13}{360.522}{372.302}{\Character{58}}%
\TextAt{8}{79.226}{343.437}{Conversely,}%
\TextAt{8}{142.737}{343.437}{let}%
\TextAt{16}{161.233}{343.437}{\Character{106}}%
\TextAt{13}{164.554}{343.437}{n}%
\TextAt{16}{171.542}{343.437}{\Character{106}}%
\NativeMathAt{180.879}{343.437}{11.955}{\leq}%
\TextAt{13}{196.227}{343.437}{N}%
\TextAt{8}{212.339}{343.437}{for}%
\TextAt{8}{231.815}{343.437}{all}%
\TextAt{13}{249.666}{343.437}{n}%
\NativeMathAt{262.687}{343.437}{11.955}{\in}%
\TextAt{18}{276.69}{343.437}{N}%
\TextAt{8}{285.324}{343.437}{.}%
\TextAt{8}{298.548}{343.437}{Let}%
\TextAt{13}{321.109}{343.437}{x\Character{59}}%
\TextAt{13}{332.981}{343.437}{y}%
\NativeMathAt{345.175}{343.437}{11.955}{\in}%
\TextAt{13}{359.178}{343.437}{K}%
\TextAt{8}{375.485}{343.437}{and}%
\TextAt{8}{399.839}{343.437}{suppose}%
\TextAt{16}{445.449}{343.437}{\Character{106}}%
\TextAt{13}{448.77}{343.437}{x}%
\TextAt{16}{455.422}{343.437}{\Character{106}}%
\NativeMathAt{464.759}{343.437}{11.955}{\geq}%
\TextAt{16}{480.086}{343.437}{\Character{106}}%
\TextAt{13}{483.429}{343.437}{y}%
\TextAt{16}{489.566}{343.437}{\Character{106}}%
\TextAt{8}{492.886}{343.437}{.}%
\TextAt{8}{506.11}{343.437}{Then}%
\TextAt{16}{79.226}{326.391}{\Character{106}}%
\TextAt{13}{82.547}{326.391}{x}%
\TextAt{16}{89.199}{326.391}{\Character{106}}%
\TextAt{21}{92.52}{330.729}{\Character{23}}%
\TextAt{16}{97.702}{326.391}{\Character{106}}%
\TextAt{13}{101.023}{326.391}{y}%
\TextAt{16}{107.16}{326.391}{\Character{106}}%
\TextAt{21}{110.48}{330.729}{n}%
\TextAt{3}{115.619}{330.729}{\Character{0}}%
\TextAt{21}{122.205}{330.729}{\Character{23}}%
\NativeMathAt{130.708}{326.391}{11.955}{\leq}%
\TextAt{16}{143.309}{326.391}{\Character{106}}%
\TextAt{13}{146.648}{326.391}{x}%
\TextAt{16}{153.3}{326.391}{\Character{106}}%
\TextAt{21}{156.621}{330.729}{n}%
\TextAt{8}{166.155}{326.391}{for}%
\TextAt{13}{184.04}{326.391}{\Character{23}}%
\NativeMathAt{193.899}{326.391}{11.955}{\geq}%
\TextAt{9}{206.518}{326.391}{0}%
\TextAt{8}{216.268}{326.391}{and}%
\TextAt{8}{239.031}{326.391}{one}%
\TextAt{8}{260.491}{326.391}{gets}%
\TextAt{16}{184.745}{291.016}{\Character{106}}%
\TextAt{13}{188.066}{291.016}{x}%
\TextAt{9}{197.374}{291.016}{\Character{43}}%
\TextAt{13}{209.136}{291.016}{y}%
\TextAt{16}{215.272}{291.016}{\Character{106}}%
\TextAt{21}{218.593}{295.952}{n}%
\NativeMathAt{227.551}{291.016}{11.955}{\leq}%
\TextAt{21}{246.235}{305.96}{n}%
\TextAt{20}{240.17}{302.374}{\Character{88}}%
\TextAt{21}{241.052}{277.268}{\Character{23}}%
\TextAt{17}{245.736}{277.268}{\Character{61}0}%
\TextAt{20}{259.431}{308.351}{\Character{12}}%
\TextAt{20}{259.431}{301.178}{\Character{12}}%
\TextAt{20}{259.431}{294.005}{\Character{12}}%
\TextAt{20}{259.431}{286.832}{\Character{12}}%
\TextAt{20}{263.416}{307.873}{\Character{18}}%
\TextAt{13}{272.216}{299.104}{n}%
\TextAt{13}{272.441}{282.815}{\Character{23}}%
\TextAt{20}{279.204}{307.873}{\Character{19}}%
\TextAt{20}{288.004}{308.351}{\Character{12}}%
\TextAt{20}{288.004}{301.178}{\Character{12}}%
\TextAt{20}{288.004}{294.005}{\Character{12}}%
\TextAt{20}{288.004}{286.832}{\Character{12}}%
\TextAt{16}{293.982}{291.016}{\Character{106}}%
\TextAt{13}{297.303}{291.016}{x}%
\TextAt{16}{303.955}{291.016}{\Character{106}}%
\TextAt{21}{307.276}{295.952}{\Character{23}}%
\TextAt{16}{312.458}{291.016}{\Character{106}}%
\TextAt{13}{315.779}{291.016}{y}%
\TextAt{16}{321.915}{291.016}{\Character{106}}%
\TextAt{21}{325.236}{295.952}{n}%
\TextAt{3}{330.374}{295.952}{\Character{0}}%
\TextAt{21}{336.961}{295.952}{\Character{23}}%
\NativeMathAt{345.464}{291.016}{11.955}{\leq}%
\TextAt{13}{358.083}{291.016}{N}%
\TextAt{9}{368.706}{291.016}{\Character{40}}%
\TextAt{13}{373.258}{291.016}{n}%
\TextAt{9}{382.902}{291.016}{\Character{43}}%
\TextAt{9}{394.654}{291.016}{1\Character{41}}%
\TextAt{16}{405.069}{291.016}{\Character{106}}%
\TextAt{13}{408.39}{291.016}{x}%
\TextAt{16}{415.042}{291.016}{\Character{106}}%
\TextAt{21}{418.363}{295.952}{n}%
\TextAt{13}{423.999}{291.016}{\Character{59}}%
\TextAt{8}{79.226}{254.96}{hence}%
\TextAt{16}{156.233}{237.914}{\Character{106}}%
\TextAt{13}{159.554}{237.914}{x}%
\TextAt{9}{168.863}{237.914}{\Character{43}}%
\TextAt{13}{180.624}{237.914}{y}%
\TextAt{16}{186.761}{237.914}{\Character{106}}%
\NativeMathAt{193.384}{237.914}{11.955}{\leq}%
\TextAt{13}{206.022}{237.914}{N}%
\TextAt{17}{216.644}{242.85}{1\Character{47}}%
\TextAt{21}{225.113}{242.85}{n}%
\TextAt{9}{230.749}{237.914}{\Character{40}1}%
\TextAt{9}{243.792}{237.914}{\Character{43}}%
\TextAt{13}{255.572}{237.914}{n}%
\TextAt{9}{262.56}{237.914}{\Character{41}}%
\TextAt{17}{267.112}{242.85}{1\Character{47}}%
\TextAt{21}{275.581}{242.85}{n}%
\TextAt{16}{281.217}{237.914}{\Character{106}}%
\TextAt{13}{284.538}{237.914}{x}%
\TextAt{16}{291.19}{237.914}{\Character{106}}%
\TextAt{9}{297.832}{237.914}{\Character{61}}%
\TextAt{13}{310.257}{237.914}{N}%
\TextAt{17}{320.88}{242.85}{1\Character{47}}%
\TextAt{21}{329.348}{242.85}{n}%
\TextAt{9}{334.985}{237.914}{\Character{40}1}%
\TextAt{9}{348.028}{237.914}{\Character{43}}%
\TextAt{13}{359.808}{237.914}{n}%
\TextAt{9}{366.795}{237.914}{\Character{41}}%
\TextAt{17}{371.348}{242.85}{1\Character{47}}%
\TextAt{21}{379.816}{242.85}{n}%
\TextAt{8}{387.445}{237.914}{max}%
\TextAt{16}{409.239}{237.914}{\Character{102}\Character{106}}%
\TextAt{13}{418.538}{237.914}{x}%
\TextAt{16}{425.19}{237.914}{\Character{106}}%
\TextAt{13}{428.511}{237.914}{\Character{59}}%
\TextAt{16}{433.755}{237.914}{\Character{106}}%
\TextAt{13}{437.076}{237.914}{y}%
\TextAt{16}{443.212}{237.914}{\Character{106}\Character{103}}%
\TextAt{13}{452.511}{237.914}{\Character{59}}%
\TextAt{8}{79.226}{214.451}{and}%
\TextAt{8}{101.989}{214.451}{thus}%
\TextAt{16}{127.74}{214.451}{\Character{106}}%
\TextAt{13}{131.061}{214.451}{x}%
\TextAt{9}{140.37}{214.451}{\Character{43}}%
\TextAt{13}{152.131}{214.451}{y}%
\TextAt{16}{158.268}{214.451}{\Character{106}}%
\NativeMathAt{164.891}{214.451}{11.955}{\leq}%
\TextAt{8}{177.529}{214.451}{max}%
\TextAt{16}{199.323}{214.451}{\Character{102}\Character{106}}%
\TextAt{13}{208.622}{214.451}{x}%
\TextAt{16}{215.274}{214.451}{\Character{106}}%
\TextAt{13}{218.595}{214.451}{\Character{59}}%
\TextAt{16}{223.839}{214.451}{\Character{106}}%
\TextAt{13}{227.16}{214.451}{y}%
\TextAt{16}{233.296}{214.451}{\Character{106}\Character{103}}%
\TextAt{8}{246.492}{214.451}{by}%
\TextAt{8}{262.751}{214.451}{letting}%
\TextAt{13}{299.824}{214.451}{n}%
\NativeMathAt{310.133}{214.451}{11.955}{\rightarrow}%
\NativeMathAt{325.4}{214.451}{11.955}{\infty}%
\TextAt{8}{337.364}{214.451}{.}%
\Vector{q 0 G 0 g 0.399 w 452.334 -559.549 m 452.334 -551.479 l S Q}%
\Vector{q 0 G 0 g 0.399 w 452.534 -551.678 m 459.707 -551.678 l S Q}%
\Vector{q 0 G 0 g 0.399 w 452.534 -559.35 m 459.707 -559.35 l S Q}%
\Vector{q 0 G 0 g 0.399 w 459.906 -559.549 m 459.906 -551.479 l S Q}%
\TextAt{12}{79.226}{182.52}{Remark.}%
\TextAt{8}{133.383}{182.52}{The}%
\TextAt{8}{157.437}{182.52}{strong}%
\TextAt{8}{193.278}{182.52}{triangle}%
\TextAt{8}{236.209}{182.52}{inequality}%
\TextAt{8}{290.522}{182.52}{immediately}%
\TextAt{8}{357.184}{182.52}{implies}%
\TextAt{8}{396.911}{182.52}{that}%
\TextAt{16}{213.278}{153.655}{\Character{106}}%
\TextAt{13}{216.599}{153.655}{x}%
\TextAt{16}{223.251}{153.655}{\Character{106}}%
\NativeMathAt{229.874}{153.655}{11.955}{\neq}%
\TextAt{16}{242.318}{153.655}{\Character{106}}%
\TextAt{13}{245.639}{153.655}{y}%
\TextAt{16}{251.776}{153.655}{\Character{106}}%
\NativeArrowAt{258.417}{153.655}{11.955}{19.067}{\Longrightarrow}%
\TextAt{16}{280.796}{153.655}{\Character{106}}%
\TextAt{13}{284.126}{153.655}{x}%
\TextAt{9}{293.435}{153.655}{\Character{43}}%
\TextAt{13}{305.196}{153.655}{y}%
\TextAt{16}{311.333}{153.655}{\Character{106}}%
\TextAt{9}{317.975}{153.655}{\Character{61}}%
\TextAt{8}{330.4}{153.655}{max}%
\TextAt{16}{352.195}{153.655}{\Character{102}\Character{106}}%
\TextAt{13}{361.493}{153.655}{x}%
\TextAt{16}{368.145}{153.655}{\Character{106}}%
\TextAt{13}{371.466}{153.655}{\Character{59}}%
\TextAt{16}{376.71}{153.655}{\Character{106}}%
\TextAt{13}{380.031}{153.655}{y}%
\TextAt{16}{386.168}{153.655}{\Character{106}\Character{103}}%
\TextAt{13}{395.466}{153.655}{\Character{58}}%
\TextAt{8}{97.159}{118.832}{One}%
\TextAt{8}{121.628}{118.832}{may}%
\TextAt{8}{146.755}{118.832}{extend}%
\TextAt{8}{184.555}{118.832}{the}%
\TextAt{8}{204.47}{118.832}{nonarchimedean}%
\TextAt{8}{291.058}{118.832}{valuation}%
\TextAt{16}{341.889}{118.832}{\Character{106}}%
\TextAt{16}{348.847}{118.832}{\Character{106}}%
\TextAt{8}{355.842}{118.832}{of}%
\TextAt{13}{368.93}{118.832}{K}%
\TextAt{8}{383.403}{118.832}{to}%
\TextAt{8}{397.472}{118.832}{a}%
\TextAt{8}{406.985}{118.832}{valuation}%
\TextAt{8}{457.816}{118.832}{of}%
\TextAt{8}{470.904}{118.832}{the}%
\TextAt{8}{490.819}{118.832}{function}%
\TextAt{8}{79.226}{101.786}{field}%
\TextAt{13}{104.583}{101.786}{K}%
\TextAt{9}{115.401}{101.786}{\Character{40}}%
\TextAt{13}{119.953}{101.786}{t}%
\TextAt{9}{124.18}{101.786}{\Character{41}}%
\TextAt{8}{132.63}{101.786}{in}%
\TextAt{8}{146.283}{101.786}{a}%
\TextAt{8}{156.038}{101.786}{canonical}%
\TextAt{8}{207.421}{101.786}{way}%
\TextAt{8}{231.164}{101.786}{by}%
\TextAt{8}{247.423}{101.786}{setting,}%
\TextAt{8}{289.111}{101.786}{for}%
\TextAt{8}{306.996}{101.786}{a}%
\TextAt{8}{316.752}{101.786}{polynomial}%
\TextAt{13}{377.245}{101.786}{f}%
\TextAt{9}{384.291}{101.786}{\Character{40}}%
\TextAt{13}{388.843}{101.786}{t}%
\TextAt{9}{393.071}{101.786}{\Character{41}}%
\TextAt{9}{400.926}{101.786}{\Character{61}}%
\TextAt{13}{413.369}{101.786}{a}%
\TextAt{17}{419.514}{99.993}{0}%
\TextAt{9}{426.903}{101.786}{\Character{43}}%
\TextAt{13}{438.664}{101.786}{a}%
\TextAt{17}{444.809}{99.993}{1}%
\TextAt{13}{449.542}{101.786}{t}%
\TextAt{9}{456.426}{101.786}{\Character{43}}%
\TextAt{16}{468.187}{101.786}{\Character{1}}%
\TextAt{16}{473.483}{101.786}{\Character{1}}%
\TextAt{16}{478.791}{101.786}{\Character{1}}%
\TextAt{9}{484.791}{101.786}{\Character{43}}%
\TextAt{13}{496.553}{101.786}{a}%
\TextAt{21}{502.697}{99.993}{n}%
\TextAt{13}{508.334}{101.786}{t}%
\TextAt{21}{512.561}{106.124}{n}%
\TextAt{8}{518.197}{101.786}{,}%
\TextAt{16}{241.699}{72.921}{\Character{106}}%
\TextAt{13}{245.019}{72.921}{f}%
\TextAt{16}{252.066}{72.921}{\Character{106}}%
\TextAt{9}{258.708}{72.921}{\Character{61}}%
\TextAt{8}{271.133}{72.921}{max}%
\TextAt{16}{292.927}{72.921}{\Character{102}\Character{106}}%
\TextAt{13}{302.226}{72.921}{a}%
\TextAt{17}{308.371}{71.128}{0}%
\TextAt{16}{313.103}{72.921}{\Character{106}}%
\TextAt{13}{316.424}{72.921}{\Character{59}}%
\TextAt{13}{321.648}{72.921}{\Character{58}}%
\TextAt{13}{326.885}{72.921}{\Character{58}}%
\TextAt{13}{332.109}{72.921}{\Character{58}}%
\TextAt{13}{337.346}{72.921}{\Character{59}}%
\TextAt{16}{342.645}{72.921}{\Character{106}}%
\TextAt{13}{345.966}{72.921}{a}%
\TextAt{21}{352.111}{71.128}{n}%
\TextAt{16}{357.747}{72.921}{\Character{106}\Character{103}}%
\TextAt{13}{367.045}{72.921}{\Character{58}}%
\TextAt{8}{297.211}{43.033}{108}%
\Destination{78.23}{809.815}{nav.page.108}%
\Destination{79.36}{556.755}{nav.number.2.3.5}%
\Destination{79.36}{507.755}{nav.number.2.3.6}%
}

% Page 121 (body page 109)
\Page{%
\TextAt{2}{79.226}{796.943}{II.}%
\TextAt{2}{94.073}{796.943}{The}%
\TextAt{2}{116.492}{796.943}{Theory}%
\TextAt{2}{154.401}{796.943}{of}%
\TextAt{2}{166.826}{796.943}{Valuations}%
\Vector{q 0 G 0 g 0.299 w 7.226 18.729 m 460.769 18.729 l S Q}%
\TextAt{8}{79.226}{761.124}{The}%
\TextAt{8}{103.082}{761.124}{triangle}%
\TextAt{8}{145.815}{761.124}{inequality}%
\TextAt{16}{199.929}{761.124}{\Character{106}}%
\TextAt{13}{203.25}{761.124}{f}%
\TextAt{9}{212.549}{761.124}{\Character{43}}%
\TextAt{13}{223.905}{761.124}{g}%
\TextAt{16}{229.94}{761.124}{\Character{106}}%
\NativeMathAt{236.563}{761.124}{11.955}{\leq}%
\TextAt{8}{249.201}{761.124}{max}%
\TextAt{16}{270.995}{761.124}{\Character{102}\Character{106}}%
\TextAt{13}{280.293}{761.124}{f}%
\TextAt{16}{287.34}{761.124}{\Character{106}}%
\TextAt{13}{290.661}{761.124}{\Character{59}}%
\TextAt{16}{295.905}{761.124}{\Character{106}}%
\TextAt{13}{299.226}{761.124}{g}%
\TextAt{16}{305.26}{761.124}{\Character{106}\Character{103}}%
\TextAt{8}{318.258}{761.124}{is}%
\TextAt{8}{329.824}{761.124}{immediate.}%
\TextAt{8}{391.538}{761.124}{The}%
\TextAt{8}{415.394}{761.124}{proof}%
\TextAt{8}{445.766}{761.124}{that}%
\TextAt{16}{470.936}{761.124}{\Character{106}}%
\TextAt{13}{474.257}{761.124}{f}%
\TextAt{13}{481.287}{761.124}{g}%
\TextAt{16}{487.338}{761.124}{\Character{106}}%
\TextAt{9}{493.98}{761.124}{\Character{61}}%
\TextAt{16}{506.405}{761.124}{\Character{106}}%
\TextAt{13}{509.726}{761.124}{f}%
\TextAt{16}{516.773}{761.124}{\Character{106}\Character{106}}%
\TextAt{13}{523.414}{761.124}{g}%
\TextAt{16}{529.449}{761.124}{\Character{106}}%
\TextAt{8}{79.226}{744.078}{is}%
\TextAt{8}{91.726}{744.078}{the}%
\TextAt{8}{112.619}{744.078}{same}%
\TextAt{8}{142.682}{744.078}{as}%
\TextAt{8}{157.788}{744.078}{the}%
\TextAt{8}{178.681}{744.078}{proof}%
\TextAt{8}{209.987}{744.078}{of}%
\TextAt{8}{224.053}{744.078}{Gauss’s}%
\TextAt{8}{267.326}{744.078}{lemma}%
\TextAt{8}{305.781}{744.078}{for}%
\TextAt{8}{324.403}{744.078}{polynomials}%
\TextAt{8}{390.247}{744.078}{over}%
\TextAt{8}{416.03}{744.078}{factorial}%
\TextAt{8}{462.626}{744.078}{rings}%
\TextAt{8}{492.043}{744.078}{once}%
\TextAt{8}{519.439}{744.078}{we}%
\TextAt{8}{79.226}{727.032}{replace}%
\TextAt{8}{118.621}{727.032}{the}%
\TextAt{8}{138.504}{727.032}{content}%
\TextAt{8}{179.86}{727.032}{of}%
\TextAt{13}{192.917}{727.032}{f}%
\TextAt{8}{203.588}{727.032}{in}%
\TextAt{8}{216.968}{727.032}{this}%
\TextAt{8}{239.518}{727.032}{lemma}%
\TextAt{8}{276.964}{727.032}{by}%
\TextAt{8}{292.95}{727.032}{the}%
\TextAt{8}{312.834}{727.032}{absolute}%
\TextAt{8}{358.804}{727.032}{value}%
\TextAt{16}{388.778}{727.032}{\Character{106}}%
\TextAt{13}{392.098}{727.032}{f}%
\TextAt{16}{399.145}{727.032}{\Character{106}}%
\TextAt{8}{402.466}{727.032}{.}%
\TextAt{8}{410.823}{727.032}{For}%
\TextAt{8}{431.52}{727.032}{the}%
\TextAt{8}{451.404}{727.032}{field}%
\TextAt{18}{476.488}{727.032}{Q}%
\TextAt{8}{485.786}{727.032}{,}%
\TextAt{8}{492.717}{727.032}{we}%
\TextAt{8}{509.672}{727.032}{have}%
\TextAt{8}{79.226}{709.986}{the}%
\TextAt{8}{100.133}{709.986}{usual}%
\TextAt{8}{131.512}{709.986}{absolute}%
\TextAt{8}{178.505}{709.986}{value}%
\TextAt{16}{209.502}{709.986}{\Character{106}}%
\TextAt{16}{217.452}{709.986}{\Character{106}}%
\NativeMathAt{220.791}{708.193}{7.97}{\infty}%
\TextAt{9}{234.357}{709.986}{\Character{61}}%
\TextAt{16}{248.061}{709.986}{\Character{106}}%
\TextAt{16}{256.011}{709.986}{\Character{106}}%
\TextAt{8}{259.35}{709.986}{,}%
\TextAt{8}{267.437}{709.986}{this}%
\TextAt{8}{291.01}{709.986}{being}%
\TextAt{8}{323.298}{709.986}{the}%
\TextAt{8}{344.204}{709.986}{archimedean}%
\TextAt{8}{412.92}{709.986}{valuation,}%
\TextAt{8}{468.182}{709.986}{and}%
\TextAt{8}{491.694}{709.986}{for}%
\TextAt{8}{510.33}{709.986}{each}%
\TextAt{8}{79.226}{692.94}{prime}%
\TextAt{8}{112.39}{692.94}{number}%
\TextAt{13}{154.986}{692.94}{p}%
\TextAt{8}{164.759}{692.94}{the}%
\TextAt{8}{184.915}{692.94}{nonarchimedean}%
\TextAt{8}{271.745}{692.94}{valuation}%
\TextAt{16}{322.818}{692.94}{\Character{106}}%
\TextAt{16}{330.015}{692.94}{\Character{106}}%
\TextAt{21}{333.357}{691.147}{p}%
\TextAt{8}{338.118}{692.94}{.}%
\TextAt{8}{346.566}{692.94}{As}%
\TextAt{8}{363.853}{692.94}{a}%
\TextAt{8}{373.609}{692.94}{matter}%
\TextAt{8}{411.985}{692.94}{of}%
\TextAt{8}{425.315}{692.94}{fact:}%
\TextAt{12}{79.226}{667.546}{Proposition}%
\TextAt{12}{152.874}{667.546}{(3.7).}%
\TextAt{8}{189.003}{667.546}{Every}%
\TextAt{8}{223.471}{667.546}{valuation}%
\TextAt{8}{275.357}{667.546}{of}%
\TextAt{18}{289.501}{667.546}{Q}%
\TextAt{8}{303.511}{667.546}{is}%
\TextAt{8}{316.089}{667.546}{equivalent}%
\TextAt{8}{372.518}{667.546}{to}%
\TextAt{8}{387.642}{667.546}{one}%
\TextAt{8}{409.916}{667.546}{of}%
\TextAt{8}{424.06}{667.546}{the}%
\TextAt{8}{445.03}{667.546}{valuations}%
\TextAt{16}{501.531}{667.546}{\Character{106}}%
\TextAt{16}{509.553}{667.546}{\Character{106}}%
\TextAt{21}{512.884}{665.753}{p}%
\TextAt{8}{522.357}{667.546}{or}%
\TextAt{16}{79.226}{650.5}{\Character{106}}%
\TextAt{16}{86.435}{650.5}{\Character{106}}%
\NativeMathAt{89.765}{648.707}{7.97}{\infty}%
\TextAt{8}{98.732}{650.5}{.}%
\TextAt{19}{79.226}{616.759}{Proof.}%
\TextAt{8}{116}{616.759}{Let}%
\TextAt{16}{137.053}{616.759}{\Character{107}}%
\TextAt{16}{147}{616.759}{\Character{107}}%
\TextAt{8}{156.969}{616.759}{be}%
\TextAt{8}{172.976}{616.759}{a}%
\TextAt{8}{182.814}{616.759}{nonarchimedean}%
\TextAt{8}{269.728}{616.759}{valuation}%
\TextAt{8}{320.883}{616.759}{of}%
\TextAt{18}{334.296}{616.759}{Q}%
\TextAt{8}{343.595}{616.759}{.}%
\TextAt{8}{352.292}{616.759}{Then}%
\TextAt{16}{382.932}{616.759}{\Character{107}}%
\TextAt{13}{388.91}{616.759}{n}%
\TextAt{16}{395.898}{616.759}{\Character{107}}%
\TextAt{9}{405.337}{616.759}{\Character{61}}%
\TextAt{16}{417.904}{616.759}{\Character{107}}%
\TextAt{9}{423.882}{616.759}{1}%
\TextAt{9}{432.43}{616.759}{\Character{43}}%
\TextAt{16}{444.266}{616.759}{\Character{1}}%
\TextAt{16}{449.562}{616.759}{\Character{1}}%
\TextAt{16}{454.87}{616.759}{\Character{1}}%
\TextAt{9}{460.927}{616.759}{\Character{43}}%
\TextAt{9}{472.727}{616.759}{1}%
\TextAt{16}{478.598}{616.759}{\Character{107}}%
\NativeMathAt{488.031}{616.759}{11.955}{\leq}%
\TextAt{9}{500.798}{616.759}{1}%
\TextAt{8}{506.651}{616.759}{,}%
\TextAt{8}{513.904}{616.759}{and}%
\TextAt{8}{79.226}{599.713}{there}%
\TextAt{8}{109.409}{599.713}{is}%
\TextAt{8}{121.445}{599.713}{a}%
\TextAt{8}{131.471}{599.713}{prime}%
\TextAt{8}{164.906}{599.713}{number}%
\TextAt{13}{207.774}{599.713}{p}%
\TextAt{8}{217.818}{599.713}{such}%
\TextAt{8}{244.486}{599.713}{that}%
\TextAt{16}{270.126}{599.713}{\Character{107}}%
\TextAt{13}{276.104}{599.713}{p}%
\TextAt{16}{281.979}{599.713}{\Character{107}}%
\TextAt{13}{291.74}{599.713}{\Character{60}}%
\TextAt{9}{304.628}{599.713}{1}%
\TextAt{8}{314.649}{599.713}{because,}%
\TextAt{8}{361.542}{599.713}{if}%
\TextAt{8}{372.537}{599.713}{not,}%
\TextAt{8}{396.942}{599.713}{unique}%
\TextAt{8}{435.255}{599.713}{prime}%
\TextAt{8}{468.69}{599.713}{factorization}%
\TextAt{8}{79.226}{582.667}{would}%
\TextAt{8}{113.37}{582.667}{imply}%
\TextAt{16}{146.211}{582.667}{\Character{107}}%
\TextAt{13}{152.189}{582.667}{x}%
\TextAt{16}{158.841}{582.667}{\Character{107}}%
\TextAt{9}{168.139}{582.667}{\Character{61}}%
\TextAt{9}{180.549}{582.667}{1}%
\TextAt{8}{190.315}{582.667}{for}%
\TextAt{8}{208.2}{582.667}{all}%
\TextAt{13}{224.459}{582.667}{x}%
\NativeMathAt{234.432}{582.667}{11.955}{\in}%
\TextAt{18}{245.723}{582.667}{Q}%
\TextAt{3}{255.021}{587.006}{\Character{3}}%
\TextAt{8}{259.754}{582.667}{.}%
\TextAt{8}{268.202}{582.667}{The}%
\TextAt{8}{292.256}{582.667}{set}%
\TextAt{14}{249.961}{555.832}{a}%
\TextAt{9}{259.265}{555.832}{\Character{61}}%
\TextAt{16}{271.691}{555.832}{\Character{102}}%
\TextAt{13}{277.668}{555.832}{a}%
\NativeMathAt{287.134}{555.832}{11.955}{\in}%
\TextAt{18}{298.425}{555.832}{Z}%
\TextAt{16}{309.716}{555.832}{\Character{106}}%
\TextAt{16}{316.339}{555.832}{\Character{107}}%
\TextAt{13}{322.335}{555.832}{a}%
\TextAt{16}{328.48}{555.832}{\Character{107}}%
\TextAt{13}{337.779}{555.832}{\Character{60}}%
\TextAt{9}{350.204}{555.832}{1}%
\TextAt{16}{356.057}{555.832}{\Character{103}}%
\TextAt{8}{79.226}{528.997}{is}%
\TextAt{8}{91.02}{528.997}{an}%
\TextAt{8}{107.309}{528.997}{ideal}%
\TextAt{8}{135.302}{528.997}{of}%
\TextAt{18}{148.662}{528.997}{Z}%
\TextAt{8}{160.559}{528.997}{satisfying}%
\TextAt{13}{212.75}{528.997}{p}%
\TextAt{18}{218.625}{528.997}{Z}%
\NativeMathAt{229.967}{528.997}{11.955}{\subseteq}%
\TextAt{14}{242.637}{528.997}{a}%
\NativeMathAt{251.992}{528.997}{11.955}{\neq}%
\TextAt{18}{264.469}{528.997}{Z}%
\TextAt{8}{272.439}{528.997}{,}%
\TextAt{8}{279.625}{528.997}{and}%
\TextAt{8}{302.418}{528.997}{since}%
\TextAt{13}{331.116}{528.997}{p}%
\TextAt{18}{336.991}{528.997}{Z}%
\TextAt{8}{348.889}{528.997}{is}%
\TextAt{8}{360.683}{528.997}{a}%
\TextAt{8}{370.468}{528.997}{maximal}%
\TextAt{8}{418.307}{528.997}{ideal,}%
\TextAt{8}{449.559}{528.997}{we}%
\TextAt{8}{466.816}{528.997}{have}%
\TextAt{14}{493.841}{528.997}{a}%
\TextAt{9}{503.196}{528.997}{\Character{61}}%
\TextAt{13}{515.672}{528.997}{p}%
\TextAt{18}{521.548}{528.997}{Z}%
\TextAt{8}{529.518}{528.997}{.}%
\TextAt{8}{79.226}{511.951}{If}%
\TextAt{8}{90.918}{511.951}{now}%
\TextAt{13}{115.307}{511.951}{a}%
\NativeMathAt{124.773}{511.951}{11.955}{\in}%
\TextAt{18}{136.064}{511.951}{Z}%
\TextAt{8}{147.931}{511.951}{and}%
\TextAt{13}{170.694}{511.951}{a}%
\TextAt{9}{180.16}{511.951}{\Character{61}}%
\TextAt{13}{192.585}{511.951}{bp}%
\TextAt{21}{203.437}{516.289}{m}%
\TextAt{8}{215.323}{511.951}{with}%
\TextAt{13}{241.983}{511.951}{p}%
\TextAt{18}{251.179}{511.951}{\Character{45}}%
\TextAt{13}{257.821}{511.951}{b}%
\TextAt{8}{262.798}{511.951}{,}%
\TextAt{8}{269.947}{511.951}{so}%
\TextAt{8}{284.317}{511.951}{that}%
\TextAt{13}{309.686}{511.951}{b}%
\NativeMathAt{317.984}{511.951}{11.955}{\notin}%
\TextAt{14}{329.275}{511.951}{a}%
\TextAt{8}{335.258}{511.951}{,}%
\TextAt{8}{342.408}{511.951}{then}%
\TextAt{16}{369.068}{511.951}{\Character{107}}%
\TextAt{13}{375.045}{511.951}{b}%
\TextAt{16}{380.022}{511.951}{\Character{107}}%
\TextAt{9}{389.321}{511.951}{\Character{61}}%
\TextAt{9}{401.731}{511.951}{1}%
\TextAt{8}{411.497}{511.951}{and}%
\TextAt{8}{434.259}{511.951}{hence}%
\TextAt{16}{259.518}{485.115}{\Character{107}}%
\TextAt{13}{265.496}{485.115}{a}%
\TextAt{16}{271.641}{485.115}{\Character{107}}%
\TextAt{9}{280.939}{485.115}{\Character{61}}%
\TextAt{16}{293.365}{485.115}{\Character{107}}%
\TextAt{13}{299.342}{485.115}{p}%
\TextAt{16}{305.217}{485.115}{\Character{107}}%
\TextAt{21}{311.195}{490.051}{m}%
\TextAt{9}{322.504}{485.115}{\Character{61}}%
\TextAt{16}{334.93}{485.115}{\Character{106}}%
\TextAt{13}{338.251}{485.115}{a}%
\TextAt{16}{344.396}{485.115}{\Character{106}}%
\TextAt{21}{347.717}{490.051}{s}%
\TextAt{21}{347.717}{482.16}{p}%
\TextAt{8}{79.226}{458.28}{where}%
\TextAt{13}{114.911}{458.28}{s}%
\TextAt{9}{126.942}{458.28}{\Character{61}}%
\TextAt{16}{142.564}{458.28}{\Character{0}}%
\TextAt{8}{153.855}{458.28}{log}%
\TextAt{16}{170.815}{458.28}{\Character{107}}%
\TextAt{13}{176.793}{458.28}{p}%
\TextAt{16}{182.668}{458.28}{\Character{107}}%
\TextAt{9}{188.646}{458.28}{\Character{47}}%
\TextAt{8}{196.491}{458.28}{log}%
\TextAt{13}{213.451}{458.28}{p}%
\TextAt{8}{219.327}{458.28}{.}%
\TextAt{8}{233.402}{458.28}{Consequently}%
\TextAt{16}{307.857}{458.28}{\Character{107}}%
\TextAt{16}{319.597}{458.28}{\Character{107}}%
\TextAt{8}{331.358}{458.28}{is}%
\TextAt{8}{344.998}{458.28}{equivalent}%
\TextAt{8}{402.489}{458.28}{to}%
\TextAt{16}{418.675}{458.28}{\Character{106}}%
\TextAt{16}{427.749}{458.28}{\Character{106}}%
\TextAt{21}{431.09}{456.487}{p}%
\TextAt{8}{435.851}{458.28}{.}%
\TextAt{8}{449.926}{458.28}{Now}%
\TextAt{8}{478.461}{458.28}{let}%
\TextAt{16}{497.241}{458.28}{\Character{107}}%
\TextAt{16}{508.981}{458.28}{\Character{107}}%
\TextAt{8}{520.743}{458.28}{be}%
\TextAt{8}{79.226}{441.234}{archimedean.}%
\TextAt{8}{151.742}{441.234}{Then}%
\TextAt{8}{182.3}{441.234}{one}%
\TextAt{8}{203.759}{441.234}{has,}%
\TextAt{8}{227.885}{441.234}{for}%
\TextAt{8}{245.77}{441.234}{every}%
\TextAt{8}{276.662}{441.234}{two}%
\TextAt{8}{298.779}{441.234}{natural}%
\TextAt{8}{339.761}{441.234}{numbers}%
\TextAt{13}{386.972}{441.234}{n\Character{59}}%
\TextAt{13}{399.178}{441.234}{m}%
\TextAt{13}{412.736}{441.234}{\Character{62}}%
\TextAt{9}{425.189}{441.234}{1}%
\TextAt{8}{431.042}{441.234}{,}%
\TextAt{16}{247.011}{414.399}{\Character{107}}%
\TextAt{13}{252.989}{414.399}{m}%
\TextAt{16}{263.228}{414.399}{\Character{107}}%
\TextAt{17}{269.206}{419.335}{1\Character{47}}%
\TextAt{7}{279.085}{419.335}{log}%
\TextAt{21}{291.312}{419.335}{m}%
\TextAt{9}{302.622}{414.399}{\Character{61}}%
\TextAt{16}{315.047}{414.399}{\Character{107}}%
\TextAt{13}{321.025}{414.399}{n}%
\TextAt{16}{328.012}{414.399}{\Character{107}}%
\TextAt{17}{333.99}{419.335}{1\Character{47}}%
\TextAt{7}{343.87}{419.335}{log}%
\TextAt{21}{356.097}{419.335}{n}%
\TextAt{13}{361.733}{414.399}{\Character{58}}%
\TextAt{8}{517.802}{414.399}{(*)}%
\TextAt{8}{79.226}{387.563}{In}%
\TextAt{8}{93.847}{387.563}{fact,}%
\TextAt{8}{120.185}{387.563}{we}%
\TextAt{8}{137.412}{387.563}{may}%
\TextAt{8}{162.781}{387.563}{write}%
\TextAt{13}{238.912}{370.517}{m}%
\TextAt{9}{252.472}{370.517}{\Character{61}}%
\TextAt{13}{264.898}{370.517}{a}%
\TextAt{17}{271.043}{368.724}{0}%
\TextAt{9}{278.432}{370.517}{\Character{43}}%
\TextAt{13}{290.193}{370.517}{a}%
\TextAt{17}{296.338}{368.724}{1}%
\TextAt{13}{301.07}{370.517}{n}%
\TextAt{9}{310.714}{370.517}{\Character{43}}%
\TextAt{16}{322.476}{370.517}{\Character{1}}%
\TextAt{16}{327.772}{370.517}{\Character{1}}%
\TextAt{16}{333.08}{370.517}{\Character{1}}%
\TextAt{9}{339.08}{370.517}{\Character{43}}%
\TextAt{13}{350.841}{370.517}{a}%
\TextAt{21}{356.986}{368.724}{r}%
\TextAt{13}{361.541}{370.517}{n}%
\TextAt{21}{368.529}{375.453}{r}%
\TextAt{8}{79.226}{347.924}{where}%
\TextAt{13}{114}{347.924}{a}%
\TextAt{21}{120.145}{346.13}{i}%
\NativeMathAt{128.491}{347.924}{11.955}{\in}%
\TextAt{16}{141.415}{347.924}{\Character{102}}%
\TextAt{9}{147.404}{347.924}{0}%
\TextAt{13}{153.257}{347.924}{\Character{59}}%
\TextAt{9}{158.501}{347.924}{1}%
\TextAt{13}{164.354}{347.924}{\Character{59}}%
\TextAt{13}{169.578}{347.924}{\Character{58}}%
\TextAt{13}{174.815}{347.924}{\Character{58}}%
\TextAt{13}{180.051}{347.924}{\Character{58}}%
\TextAt{13}{185.276}{347.924}{\Character{59}}%
\TextAt{13}{190.512}{347.924}{n}%
\TextAt{16}{200.877}{347.924}{\Character{0}}%
\TextAt{9}{213.49}{347.924}{1}%
\TextAt{16}{219.343}{347.924}{\Character{103}}%
\TextAt{8}{230.183}{347.924}{and}%
\TextAt{13}{253.91}{347.924}{n}%
\TextAt{21}{260.898}{352.262}{r}%
\NativeMathAt{270.417}{347.924}{11.955}{\leq}%
\TextAt{13}{284.681}{347.924}{m}%
\TextAt{8}{294.92}{347.924}{.}%
\TextAt{8}{306.263}{347.924}{Hence,}%
\TextAt{8}{345.498}{347.924}{observing}%
\TextAt{8}{398.886}{347.924}{that}%
\TextAt{13}{425.22}{347.924}{r}%
\NativeMathAt{435.785}{347.924}{11.955}{\leq}%
\TextAt{8}{450.049}{347.924}{log}%
\TextAt{13}{467.009}{347.924}{m}%
\TextAt{9}{477.249}{347.924}{\Character{47}}%
\TextAt{8}{485.094}{347.924}{log}%
\TextAt{13}{502.055}{347.924}{n}%
\TextAt{8}{513.904}{347.924}{and}%
\TextAt{16}{79.226}{330.878}{\Character{107}}%
\TextAt{13}{85.204}{330.878}{a}%
\TextAt{21}{91.349}{329.084}{i}%
\TextAt{16}{94.73}{330.878}{\Character{107}}%
\TextAt{9}{104.028}{330.878}{\Character{61}}%
\TextAt{16}{116.454}{330.878}{\Character{107}}%
\TextAt{9}{122.432}{330.878}{1}%
\TextAt{9}{130.932}{330.878}{\Character{43}}%
\TextAt{16}{142.702}{330.878}{\Character{1}}%
\TextAt{16}{147.998}{330.878}{\Character{1}}%
\TextAt{16}{153.306}{330.878}{\Character{1}}%
\TextAt{9}{159.307}{330.878}{\Character{43}}%
\TextAt{9}{171.059}{330.878}{1}%
\TextAt{16}{176.921}{330.878}{\Character{107}}%
\NativeMathAt{186.21}{330.878}{11.955}{\leq}%
\TextAt{13}{198.839}{330.878}{a}%
\TextAt{21}{204.984}{329.084}{i}%
\TextAt{16}{208.365}{330.878}{\Character{107}}%
\TextAt{9}{214.343}{330.878}{1}%
\TextAt{16}{220.196}{330.878}{\Character{107}}%
\NativeMathAt{229.485}{330.878}{11.955}{\leq}%
\TextAt{13}{242.114}{330.878}{n}%
\TextAt{8}{249.101}{330.878}{,}%
\TextAt{8}{256.25}{330.878}{one}%
\TextAt{8}{277.71}{330.878}{gets}%
\TextAt{8}{301.835}{330.878}{the}%
\TextAt{8}{321.992}{330.878}{inequality}%
\ErrataAt{100}{299.768}{435}{$\displaystyle N=\max(1,\|n\|),\qquad \|m\|\le\sum_i\|a_i\|\|n\|^i\le n(r+1)N^r\le n(1+\log m/\log n)N^{\log m/\log n}.$}%
\TextAt{8}{98.003}{264.077}{Substituting}%
\TextAt{8}{164.744}{264.077}{here}%
\TextAt{13}{189.786}{264.077}{m}%
\TextAt{21}{200.025}{268.415}{k}%
\TextAt{8}{208.727}{264.077}{for}%
\TextAt{13}{226.297}{264.077}{m}%
\TextAt{8}{236.536}{264.077}{,}%
\TextAt{8}{243.433}{264.077}{taking}%
\TextAt{13}{279.222}{264.077}{k}%
\TextAt{8}{285.712}{264.077}{-th}%
\TextAt{8}{304.25}{264.077}{roots}%
\TextAt{8}{333.595}{264.077}{on}%
\TextAt{8}{349.539}{264.077}{both}%
\TextAt{8}{376.864}{264.077}{sides,}%
\TextAt{8}{407.946}{264.077}{and}%
\TextAt{8}{430.393}{264.077}{letting}%
\TextAt{13}{467.151}{264.077}{k}%
\TextAt{8}{477.223}{264.077}{tend}%
\TextAt{8}{503.567}{264.077}{to}%
\NativeMathAt{517.562}{264.077}{11.955}{\infty}%
\TextAt{8}{529.518}{264.077}{,}%
\TextAt{8}{79.226}{247.031}{one}%
\TextAt{8}{100.686}{247.031}{finally}%
\TextAt{8}{136.133}{247.031}{obtains}%
\ErrataAt{234}{220.196}{246}{$\displaystyle\|m\|\le N^{\log m/\log n}.$}%
\ErrataAt{79.19}{193.36}{455.81}{As $\|m\|>1$ for some $m>1$, it follows that $\|n\|>1$ for every $n>1$. Swapping $m,n$ gives (*).}%
\ErrataAt{79.19}{176.314}{455.81}{Put $c=\|n\|^{1/\log n}=e^s$ with $s>0$. For each positive rational $x=a/b$ we have}%
\TextAt{16}{256.346}{149.479}{\Character{107}}%
\TextAt{13}{262.323}{149.479}{x}%
\TextAt{16}{268.975}{149.479}{\Character{107}}%
\TextAt{9}{278.274}{149.479}{\Character{61}}%
\TextAt{13}{290.699}{149.479}{e}%
\TextAt{21}{296.125}{154.415}{s}%
\TextAt{7}{301.452}{154.415}{log}%
\TextAt{21}{313.679}{154.415}{x}%
\TextAt{9}{322.265}{149.479}{\Character{61}}%
\TextAt{16}{334.69}{149.479}{\Character{106}}%
\TextAt{13}{338.011}{149.479}{x}%
\TextAt{16}{344.663}{149.479}{\Character{106}}%
\TextAt{21}{347.984}{154.415}{s}%
\TextAt{13}{352.398}{149.479}{\Character{58}}%
\TextAt{8}{79.226}{122.644}{Therefore}%
\TextAt{16}{132.223}{122.644}{\Character{107}}%
\TextAt{16}{142.098}{122.644}{\Character{107}}%
\TextAt{8}{151.973}{122.644}{is}%
\TextAt{8}{163.737}{122.644}{equivalent}%
\TextAt{8}{219.353}{122.644}{to}%
\TextAt{8}{233.663}{122.644}{the}%
\TextAt{8}{253.819}{122.644}{usual}%
\TextAt{8}{284.449}{122.644}{absolute}%
\TextAt{8}{330.691}{122.644}{value}%
\TextAt{16}{360.938}{122.644}{\Character{106}}%
\TextAt{16}{368.135}{122.644}{\Character{106}}%
\TextAt{8}{375.374}{122.644}{on}%
\TextAt{18}{391.633}{122.644}{Q}%
\TextAt{8}{400.932}{122.644}{.}%
\Vector{q 0 G 0 g 0.399 w 452.334 -651.356 m 452.334 -643.287 l S Q}%
\Vector{q 0 G 0 g 0.399 w 452.534 -643.486 m 459.707 -643.486 l S Q}%
\Vector{q 0 G 0 g 0.399 w 459.906 -651.356 m 459.906 -643.287 l S Q}%
\TextAt{8}{297.211}{43.033}{109}%
\Destination{522.36}{431.445}{nav.AMS.15}%
\Destination{78.23}{809.815}{nav.page.109}%
\Destination{79.36}{776.155}{nav.print.119}%
\Destination{79.36}{724.955}{nav.print.120}%
\Destination{79.36}{681.555}{nav.number.2.3.7}%
\Destination{79.226}{316.768}{errata.75}%
}

% Page 122 (body page 110)
\Page{%
\TextAt{2}{79.226}{796.943}{II.}%
\TextAt{2}{94.073}{796.943}{The}%
\TextAt{2}{116.492}{796.943}{Theory}%
\TextAt{2}{154.401}{796.943}{of}%
\TextAt{2}{166.826}{796.943}{Valuations}%
\Vector{q 0 G 0 g 0.299 w 7.226 18.729 m 460.769 18.729 l S Q}%
\TextAt{8}{97.159}{761.124}{Let}%
\TextAt{16}{118.128}{761.124}{\Character{106}}%
\TextAt{16}{125.337}{761.124}{\Character{106}}%
\TextAt{8}{132.565}{761.124}{be}%
\TextAt{8}{148.489}{761.124}{a}%
\TextAt{8}{158.245}{761.124}{nonarchimedean}%
\TextAt{8}{245.075}{761.124}{valuation}%
\TextAt{8}{296.147}{761.124}{of}%
\TextAt{8}{309.477}{761.124}{the}%
\TextAt{8}{329.634}{761.124}{field}%
\TextAt{13}{354.991}{761.124}{K}%
\TextAt{8}{365.808}{761.124}{.}%
\TextAt{8}{374.257}{761.124}{Putting}%
\TextAt{13}{179.448}{732.123}{v}%
\TextAt{9}{185.536}{732.123}{\Character{40}}%
\TextAt{13}{190.089}{732.123}{x}%
\TextAt{9}{196.741}{732.123}{\Character{41}}%
\TextAt{9}{204.596}{732.123}{\Character{61}}%
\TextAt{16}{217.039}{732.123}{\Character{0}}%
\TextAt{8}{228.33}{732.123}{log}%
\TextAt{16}{245.291}{732.123}{\Character{106}}%
\TextAt{13}{248.612}{732.123}{x}%
\TextAt{16}{255.264}{732.123}{\Character{106}}%
\TextAt{8}{282.495}{732.123}{for}%
\TextAt{13}{300.38}{732.123}{x}%
\NativeMathAt{310.353}{732.123}{11.955}{\neq}%
\TextAt{9}{322.762}{732.123}{0}%
\TextAt{13}{328.631}{732.123}{\Character{59}}%
\TextAt{8}{357.786}{732.123}{and}%
\TextAt{13}{380.548}{732.123}{v}%
\TextAt{9}{386.637}{732.123}{\Character{40}0\Character{41}}%
\TextAt{9}{404.881}{732.123}{\Character{61}}%
\NativeMathAt{417.341}{732.123}{11.955}{\infty}%
\TextAt{13}{429.296}{732.123}{\Character{59}}%
\TextAt{8}{79.226}{703.121}{we}%
\TextAt{8}{96.454}{703.121}{obtain}%
\TextAt{8}{132.881}{703.121}{a}%
\TextAt{8}{142.636}{703.121}{function}%
\TextAt{13}{256.733}{686.075}{v}%
\TextAt{9}{266.142}{686.075}{\Character{58}}%
\TextAt{13}{272.714}{686.075}{K}%
\NativeArrowAt{286.853}{686.075}{11.955}{19.26}{\longrightarrow}%
\TextAt{18}{309.435}{686.075}{R}%
\TextAt{16}{320.726}{686.075}{\Character{91}}%
\TextAt{16}{331.342}{686.075}{\Character{102}}%
\NativeMathAt{337.32}{686.075}{11.955}{\infty}%
\TextAt{16}{349.275}{686.075}{\Character{103}}%
\TextAt{8}{79.226}{662.554}{verifying}%
\TextAt{8}{127.358}{662.554}{the}%
\TextAt{8}{147.514}{662.554}{properties}%
\TextAt{8}{89.579}{642.519}{(i)}%
\TextAt{13}{107.918}{642.519}{v}%
\TextAt{9}{114.007}{642.519}{\Character{40}}%
\TextAt{13}{118.559}{642.519}{x}%
\TextAt{9}{125.211}{642.519}{\Character{41}}%
\TextAt{9}{133.066}{642.519}{\Character{61}}%
\NativeMathAt{145.51}{642.519}{11.955}{\infty}%
\NativeArrowAt{160.777}{642.519}{11.955}{21.938}{\Longleftrightarrow}%
\TextAt{13}{186.024}{642.519}{x}%
\TextAt{9}{195.997}{642.519}{\Character{61}}%
\TextAt{9}{208.407}{642.519}{0}%
\TextAt{8}{214.276}{642.519}{,}%
\TextAt{8}{86.327}{625.473}{(ii)}%
\TextAt{13}{107.918}{625.473}{v}%
\TextAt{9}{114.007}{625.473}{\Character{40}}%
\TextAt{13}{118.559}{625.473}{xy}%
\TextAt{9}{131.348}{625.473}{\Character{41}}%
\TextAt{9}{139.203}{625.473}{\Character{61}}%
\TextAt{13}{151.646}{625.473}{v}%
\TextAt{9}{157.735}{625.473}{\Character{40}}%
\TextAt{13}{162.287}{625.473}{x}%
\TextAt{9}{168.939}{625.473}{\Character{41}}%
\TextAt{9}{176.136}{625.473}{\Character{43}}%
\TextAt{13}{187.909}{625.473}{v}%
\TextAt{9}{193.998}{625.473}{\Character{40}}%
\TextAt{13}{198.55}{625.473}{y}%
\TextAt{9}{204.687}{625.473}{\Character{41}}%
\TextAt{8}{209.239}{625.473}{,}%
\TextAt{8}{83.076}{608.427}{(iii)}%
\TextAt{13}{107.918}{608.427}{v}%
\TextAt{9}{114.007}{608.427}{\Character{40}}%
\TextAt{13}{118.559}{608.427}{x}%
\TextAt{9}{127.868}{608.427}{\Character{43}}%
\TextAt{13}{139.629}{608.427}{y}%
\TextAt{9}{145.766}{608.427}{\Character{41}}%
\NativeMathAt{153.639}{608.427}{11.955}{\geq}%
\TextAt{8}{166.258}{608.427}{min}%
\TextAt{16}{185.769}{608.427}{\Character{102}}%
\TextAt{13}{191.747}{608.427}{v}%
\TextAt{9}{197.835}{608.427}{\Character{40}}%
\TextAt{13}{202.387}{608.427}{x}%
\TextAt{9}{209.039}{608.427}{\Character{41}}%
\TextAt{13}{213.592}{608.427}{\Character{59}}%
\TextAt{13}{218.816}{608.427}{v}%
\TextAt{9}{224.924}{608.427}{\Character{40}}%
\TextAt{13}{229.476}{608.427}{y}%
\TextAt{9}{235.613}{608.427}{\Character{41}}%
\TextAt{16}{240.165}{608.427}{\Character{103}}%
\TextAt{8}{246.143}{608.427}{,}%
\TextAt{8}{79.226}{588.392}{where}%
\TextAt{8}{114.45}{588.392}{we}%
\TextAt{8}{133.093}{588.392}{fix}%
\TextAt{8}{151.089}{588.392}{the}%
\TextAt{8}{172.661}{588.392}{following}%
\TextAt{8}{223.51}{588.392}{conventions}%
\TextAt{8}{288.084}{588.392}{regarding}%
\TextAt{8}{341.54}{588.392}{elements}%
\TextAt{13}{390.811}{588.392}{a}%
\NativeMathAt{402.688}{588.392}{11.955}{\in}%
\TextAt{18}{416.39}{588.392}{R}%
\TextAt{8}{430.337}{588.392}{and}%
\TextAt{8}{454.514}{588.392}{the}%
\TextAt{8}{476.086}{588.392}{symbol}%
\NativeMathAt{517.562}{588.392}{11.955}{\infty}%
\TextAt{8}{529.518}{588.392}{:}%
\TextAt{13}{79.226}{571.346}{a}%
\TextAt{13}{89.842}{571.346}{\Character{60}}%
\NativeMathAt{103.419}{571.346}{11.955}{\infty}%
\TextAt{8}{115.375}{571.346}{,}%
\TextAt{13}{123.368}{571.346}{a}%
\TextAt{9}{132.63}{571.346}{\Character{43}}%
\NativeMathAt{144.852}{571.346}{11.955}{\infty}%
\TextAt{9}{161.279}{571.346}{\Character{61}}%
\NativeMathAt{174.855}{571.346}{11.955}{\infty}%
\TextAt{8}{186.811}{571.346}{,}%
\NativeMathAt{194.804}{571.346}{11.955}{\infty}%
\TextAt{9}{209.876}{571.346}{\Character{43}}%
\NativeMathAt{222.098}{571.346}{11.955}{\infty}%
\TextAt{9}{238.525}{571.346}{\Character{61}}%
\NativeMathAt{252.102}{571.346}{11.955}{\infty}%
\TextAt{8}{264.057}{571.346}{.}%
\TextAt{8}{274.532}{571.346}{A}%
\TextAt{8}{287.879}{571.346}{function}%
\TextAt{13}{334.402}{571.346}{v}%
\TextAt{8}{345.063}{571.346}{on}%
\TextAt{13}{361.998}{571.346}{K}%
\TextAt{8}{377.388}{571.346}{with}%
\TextAt{8}{404.724}{571.346}{these}%
\TextAt{8}{435.371}{571.346}{properties}%
\TextAt{8}{491.064}{571.346}{is}%
\TextAt{8}{503.503}{571.346}{called}%
\TextAt{8}{79.226}{554.3}{an}%
\TextAt{12}{96.031}{554.3}{exponential}%
\TextAt{12}{168.915}{554.3}{valuation}%
\TextAt{8}{227.18}{554.3}{of}%
\TextAt{13}{241.055}{554.3}{K}%
\TextAt{8}{251.873}{554.3}{.}%
\TextAt{8}{261.957}{554.3}{We}%
\TextAt{8}{282.647}{554.3}{exclude}%
\TextAt{8}{325.131}{554.3}{the}%
\TextAt{8}{345.833}{554.3}{case}%
\TextAt{8}{371.149}{554.3}{of}%
\TextAt{8}{385.025}{554.3}{the}%
\TextAt{8}{405.726}{554.3}{trivial}%
\TextAt{8}{441.073}{554.3}{function}%
\TextAt{13}{487.467}{554.3}{v}%
\TextAt{9}{493.555}{554.3}{\Character{40}}%
\TextAt{13}{498.107}{554.3}{x}%
\TextAt{9}{504.759}{554.3}{\Character{41}}%
\TextAt{9}{513.546}{554.3}{\Character{61}}%
\TextAt{9}{526.9}{554.3}{0}%
\TextAt{8}{79.226}{537.254}{for}%
\TextAt{13}{97.458}{537.254}{x}%
\NativeMathAt{108.021}{537.254}{11.955}{\neq}%
\TextAt{9}{121.028}{537.254}{0}%
\TextAt{8}{126.89}{537.254}{,}%
\TextAt{13}{134.472}{537.254}{v}%
\TextAt{9}{140.56}{537.254}{\Character{40}0\Character{41}}%
\TextAt{9}{159.401}{537.254}{\Character{61}}%
\NativeMathAt{172.445}{537.254}{11.955}{\infty}%
\TextAt{8}{184.4}{537.254}{.}%
\TextAt{8}{193.888}{537.254}{Two}%
\TextAt{8}{219.591}{537.254}{exponential}%
\TextAt{8}{282.702}{537.254}{valuations}%
\TextAt{13}{338.736}{537.254}{v}%
\TextAt{17}{344.395}{535.461}{1}%
\TextAt{8}{353.371}{537.254}{and}%
\TextAt{13}{376.48}{537.254}{v}%
\TextAt{17}{382.139}{535.461}{2}%
\TextAt{8}{391.116}{537.254}{of}%
\TextAt{13}{404.792}{537.254}{K}%
\TextAt{8}{419.853}{537.254}{are}%
\TextAt{8}{439.711}{537.254}{called}%
\TextAt{12}{473.221}{537.254}{equivalent}%
\TextAt{8}{79.226}{520.208}{if}%
\TextAt{13}{90.396}{520.208}{v}%
\TextAt{17}{96.056}{518.415}{1}%
\TextAt{9}{104.87}{520.208}{\Character{61}}%
\TextAt{13}{118.056}{520.208}{sv}%
\TextAt{17}{129.229}{518.415}{2}%
\TextAt{8}{133.962}{520.208}{,}%
\TextAt{8}{141.669}{520.208}{for}%
\TextAt{8}{160}{520.208}{some}%
\TextAt{8}{189.773}{520.208}{real}%
\TextAt{8}{212.982}{520.208}{number}%
\TextAt{13}{256.025}{520.208}{s}%
\TextAt{13}{265.613}{520.208}{\Character{62}}%
\TextAt{9}{278.806}{520.208}{0}%
\TextAt{8}{284.659}{520.208}{.}%
\TextAt{8}{294.447}{520.208}{For}%
\TextAt{8}{315.863}{520.208}{every}%
\TextAt{8}{347.202}{520.208}{exponential}%
\TextAt{8}{410.413}{520.208}{valuation}%
\TextAt{13}{461.932}{520.208}{v}%
\TextAt{8}{472.364}{520.208}{we}%
\TextAt{8}{490.038}{520.208}{obtain}%
\TextAt{8}{526.911}{520.208}{a}%
\TextAt{8}{79.226}{503.162}{valuation}%
\TextAt{8}{130.299}{503.162}{in}%
\TextAt{8}{143.951}{503.162}{the}%
\TextAt{8}{164.108}{503.162}{sense}%
\TextAt{8}{194.139}{503.162}{of}%
\TextAt{8}{207.469}{503.162}{(3.1)}%
\TextAt{8}{235.444}{503.162}{by}%
\TextAt{8}{251.703}{503.162}{putting}%
\TextAt{16}{275.619}{474.161}{\Character{106}}%
\TextAt{13}{278.939}{474.161}{x}%
\TextAt{16}{285.592}{474.161}{\Character{106}}%
\TextAt{9}{292.233}{474.161}{\Character{61}}%
\TextAt{13}{304.659}{474.161}{q}%
\TextAt{3}{310.278}{479.097}{\Character{0}}%
\TextAt{21}{316.864}{479.097}{v}%
\TextAt{17}{321.274}{479.097}{\Character{40}}%
\TextAt{21}{324.567}{479.097}{x}%
\TextAt{17}{329.334}{479.097}{\Character{41}}%
\TextAt{13}{333.125}{474.161}{\Character{59}}%
\TextAt{8}{79.226}{445.16}{for}%
\TextAt{8}{98.343}{445.16}{some}%
\TextAt{8}{128.901}{445.16}{fixed}%
\TextAt{8}{158.419}{445.16}{real}%
\TextAt{8}{182.413}{445.16}{number}%
\TextAt{13}{226.242}{445.16}{q}%
\TextAt{13}{237.277}{445.16}{\Character{62}}%
\TextAt{9}{251.806}{445.16}{1}%
\TextAt{8}{257.659}{445.16}{.}%
\TextAt{8}{269.803}{445.16}{To}%
\TextAt{8}{288.262}{445.16}{distinguish}%
\TextAt{8}{348.804}{445.16}{it}%
\TextAt{8}{361.74}{445.16}{from}%
\TextAt{13}{390.612}{445.16}{v}%
\TextAt{8}{396.701}{445.16}{,}%
\TextAt{8}{405.39}{445.16}{we}%
\TextAt{8}{423.849}{445.16}{call}%
\TextAt{16}{446.54}{445.16}{\Character{106}}%
\TextAt{16}{454.98}{445.16}{\Character{106}}%
\TextAt{8}{463.441}{445.16}{an}%
\TextAt{8}{480.932}{445.16}{associated}%
\TextAt{12}{79.226}{428.114}{multiplicative}%
\TextAt{12}{164.261}{428.114}{valuation}%
\TextAt{8}{218.083}{428.114}{,}%
\TextAt{8}{225.45}{428.114}{or}%
\TextAt{12}{239.935}{428.114}{absolute}%
\TextAt{12}{293.252}{428.114}{value}%
\TextAt{8}{323.487}{428.114}{.}%
\TextAt{8}{332.458}{428.114}{Replacing}%
\TextAt{13}{386.778}{428.114}{v}%
\TextAt{8}{396.938}{428.114}{by}%
\TextAt{8}{413.371}{428.114}{an}%
\TextAt{8}{429.805}{428.114}{equivalent}%
\TextAt{8}{485.594}{428.114}{valuation}%
\TextAt{13}{79.226}{411.068}{sv}%
\TextAt{8}{94.22}{411.068}{(i.e.,}%
\TextAt{8}{120.476}{411.068}{replacing}%
\TextAt{13}{170.05}{411.068}{q}%
\TextAt{8}{179.061}{411.068}{by}%
\TextAt{13}{194.814}{411.068}{q}%
\TextAt{3}{200.434}{415.407}{\Character{48}}%
\TextAt{9}{206.549}{411.068}{\Character{61}}%
\TextAt{13}{218.975}{411.068}{q}%
\TextAt{21}{224.594}{415.407}{s}%
\TextAt{8}{229.008}{411.068}{)}%
\TextAt{8}{236.955}{411.068}{changes}%
\TextAt{16}{279.763}{411.068}{\Character{106}}%
\TextAt{16}{286.458}{411.068}{\Character{106}}%
\TextAt{8}{293.188}{411.068}{into}%
\TextAt{8}{316.425}{411.068}{the}%
\TextAt{8}{336.076}{411.068}{equivalent}%
\TextAt{8}{391.186}{411.068}{multiplicative}%
\TextAt{8}{464.503}{411.068}{valuation}%
\TextAt{16}{515.07}{411.068}{\Character{106}}%
\TextAt{16}{521.765}{411.068}{\Character{106}}%
\TextAt{21}{525.104}{415.407}{s}%
\TextAt{8}{529.518}{411.068}{.}%
\TextAt{8}{79.226}{394.022}{The}%
\TextAt{8}{103.28}{394.022}{conditions}%
\TextAt{8}{159.278}{394.022}{(i),}%
\TextAt{8}{178.789}{394.022}{(ii),}%
\TextAt{8}{201.551}{394.022}{(iii)}%
\TextAt{8}{224.314}{394.022}{immediately}%
\TextAt{8}{290.976}{394.022}{imply}%
\TextAt{8}{323.817}{394.022}{the}%
\TextAt{12}{79.226}{368.01}{Proposition}%
\TextAt{12}{151.937}{368.01}{(3.8).}%
\TextAt{8}{187.253}{368.01}{The}%
\TextAt{8}{211.307}{368.01}{subset}%
\TextAt{14}{193.968}{339.009}{o}%
\TextAt{9}{203.133}{339.009}{\Character{61}}%
\TextAt{16}{215.558}{339.009}{\Character{102}}%
\TextAt{13}{221.536}{339.009}{x}%
\NativeMathAt{231.509}{339.009}{11.955}{\in}%
\TextAt{13}{242.8}{339.009}{K}%
\TextAt{16}{256.938}{339.009}{\Character{106}}%
\TextAt{13}{263.58}{339.009}{v}%
\TextAt{9}{269.668}{339.009}{\Character{40}}%
\TextAt{13}{274.22}{339.009}{x}%
\TextAt{9}{280.872}{339.009}{\Character{41}}%
\NativeMathAt{288.746}{339.009}{11.955}{\geq}%
\TextAt{9}{301.365}{339.009}{0}%
\TextAt{16}{307.218}{339.009}{\Character{103}}%
\TextAt{9}{316.516}{339.009}{\Character{61}}%
\TextAt{16}{328.942}{339.009}{\Character{102}}%
\TextAt{13}{334.919}{339.009}{x}%
\NativeMathAt{344.892}{339.009}{11.955}{\in}%
\TextAt{13}{356.183}{339.009}{K}%
\TextAt{16}{370.322}{339.009}{\Character{106}}%
\TextAt{16}{376.945}{339.009}{\Character{106}}%
\TextAt{13}{380.284}{339.009}{x}%
\TextAt{16}{386.937}{339.009}{\Character{106}}%
\NativeMathAt{393.56}{339.009}{11.955}{\leq}%
\TextAt{9}{406.198}{339.009}{1}%
\TextAt{16}{412.051}{339.009}{\Character{103}}%
\TextAt{8}{79.226}{310.008}{is}%
\TextAt{8}{90.99}{310.008}{a}%
\TextAt{8}{100.745}{310.008}{ring}%
\TextAt{8}{124.811}{310.008}{with}%
\TextAt{8}{151.471}{310.008}{group}%
\TextAt{8}{184.647}{310.008}{of}%
\TextAt{8}{197.977}{310.008}{units}%
\TextAt{14}{191.795}{281.006}{o}%
\TextAt{3}{197.639}{285.943}{\Character{3}}%
\TextAt{9}{205.693}{281.006}{\Character{61}}%
\TextAt{16}{218.118}{281.006}{\Character{102}}%
\TextAt{13}{224.096}{281.006}{x}%
\NativeMathAt{234.069}{281.006}{11.955}{\in}%
\TextAt{13}{245.36}{281.006}{K}%
\TextAt{16}{259.498}{281.006}{\Character{106}}%
\TextAt{13}{266.14}{281.006}{v}%
\TextAt{9}{272.228}{281.006}{\Character{40}}%
\TextAt{13}{276.78}{281.006}{x}%
\TextAt{9}{283.432}{281.006}{\Character{41}}%
\TextAt{9}{291.287}{281.006}{\Character{61}}%
\TextAt{9}{303.708}{281.006}{0}%
\TextAt{16}{309.584}{281.006}{\Character{103}}%
\TextAt{9}{318.883}{281.006}{\Character{61}}%
\TextAt{16}{331.308}{281.006}{\Character{102}}%
\TextAt{13}{337.286}{281.006}{x}%
\NativeMathAt{347.259}{281.006}{11.955}{\in}%
\TextAt{13}{358.55}{281.006}{K}%
\TextAt{16}{372.688}{281.006}{\Character{106}}%
\TextAt{16}{379.311}{281.006}{\Character{106}}%
\TextAt{13}{382.651}{281.006}{x}%
\TextAt{16}{389.303}{281.006}{\Character{106}}%
\TextAt{9}{395.944}{281.006}{\Character{61}}%
\TextAt{9}{408.353}{281.006}{1}%
\TextAt{16}{414.223}{281.006}{\Character{103}}%
\TextAt{8}{79.226}{252.005}{and}%
\TextAt{8}{101.989}{252.005}{the}%
\TextAt{8}{122.145}{252.005}{unique}%
\TextAt{8}{160.187}{252.005}{maximal}%
\TextAt{8}{207.995}{252.005}{ideal}%
\TextAt{14}{192.466}{223.004}{p}%
\TextAt{9}{201.77}{223.004}{\Character{61}}%
\TextAt{16}{214.196}{223.004}{\Character{102}}%
\TextAt{13}{220.173}{223.004}{x}%
\NativeMathAt{230.146}{223.004}{11.955}{\in}%
\TextAt{13}{241.437}{223.004}{K}%
\TextAt{16}{255.576}{223.004}{\Character{106}}%
\TextAt{13}{262.217}{223.004}{v}%
\TextAt{9}{268.305}{223.004}{\Character{40}}%
\TextAt{13}{272.858}{223.004}{x}%
\TextAt{9}{279.51}{223.004}{\Character{41}}%
\TextAt{13}{287.383}{223.004}{\Character{62}}%
\TextAt{9}{299.809}{223.004}{0}%
\TextAt{16}{305.662}{223.004}{\Character{103}}%
\TextAt{9}{314.96}{223.004}{\Character{61}}%
\TextAt{16}{327.385}{223.004}{\Character{102}}%
\TextAt{13}{333.363}{223.004}{x}%
\NativeMathAt{343.336}{223.004}{11.955}{\in}%
\TextAt{13}{354.627}{223.004}{K}%
\TextAt{16}{368.765}{223.004}{\Character{106}}%
\TextAt{16}{375.388}{223.004}{\Character{106}}%
\TextAt{13}{378.728}{223.004}{x}%
\TextAt{16}{385.38}{223.004}{\Character{106}}%
\TextAt{13}{392.022}{223.004}{\Character{60}}%
\TextAt{9}{404.447}{223.004}{1}%
\TextAt{16}{410.3}{223.004}{\Character{103}}%
\TextAt{13}{416.278}{223.004}{\Character{58}}%
\TextAt{14}{97.159}{185.037}{o}%
\TextAt{8}{106.936}{185.037}{is}%
\TextAt{8}{118.735}{185.037}{an}%
\TextAt{8}{135.03}{185.037}{integral}%
\TextAt{8}{177.674}{185.037}{domain}%
\TextAt{8}{219.337}{185.037}{with}%
\TextAt{8}{246.033}{185.037}{field}%
\TextAt{8}{271.425}{185.037}{of}%
\TextAt{8}{284.791}{185.037}{fractions}%
\TextAt{13}{332.695}{185.037}{K}%
\TextAt{8}{347.445}{185.037}{and}%
\TextAt{8}{370.244}{185.037}{has}%
\TextAt{8}{391.153}{185.037}{the}%
\TextAt{8}{411.345}{185.037}{property}%
\TextAt{8}{459.189}{185.037}{that,}%
\TextAt{8}{487.854}{185.037}{for}%
\TextAt{8}{505.775}{185.037}{every}%
\TextAt{13}{79.226}{167.991}{x}%
\NativeMathAt{89.199}{167.991}{11.955}{\in}%
\TextAt{13}{100.49}{167.991}{K}%
\TextAt{8}{111.308}{167.991}{,}%
\TextAt{8}{118.168}{167.991}{either}%
\TextAt{13}{150.97}{167.991}{x}%
\NativeMathAt{160.943}{167.991}{11.955}{\in}%
\TextAt{14}{172.234}{167.991}{o}%
\TextAt{8}{181.614}{167.991}{or}%
\TextAt{13}{195.563}{167.991}{x}%
\TextAt{3}{202.215}{172.329}{\Character{0}}%
\TextAt{17}{208.802}{172.329}{1}%
\NativeMathAt{216.855}{167.991}{11.955}{\in}%
\TextAt{14}{228.146}{167.991}{o}%
\TextAt{8}{233.99}{167.991}{.}%
\TextAt{8}{242.318}{167.991}{Such}%
\TextAt{8}{270.243}{167.991}{a}%
\TextAt{8}{279.637}{167.991}{ring}%
\TextAt{8}{303.341}{167.991}{is}%
\TextAt{8}{314.744}{167.991}{called}%
\TextAt{8}{347.546}{167.991}{a}%
\TextAt{12}{356.94}{167.991}{valuation}%
\TextAt{12}{414.83}{167.991}{ring}%
\TextAt{8}{438.274}{167.991}{.}%
\TextAt{8}{446.602}{167.991}{Its}%
\TextAt{8}{463.528}{167.991}{only}%
\TextAt{8}{488.858}{167.991}{maximal}%
\TextAt{8}{79.226}{150.945}{ideal}%
\TextAt{8}{107.734}{150.945}{is}%
\TextAt{14}{120.043}{150.945}{p}%
\TextAt{9}{130.275}{150.945}{\Character{61}}%
\TextAt{16}{143.629}{150.945}{\Character{102}}%
\TextAt{13}{149.607}{150.945}{x}%
\NativeMathAt{160.508}{150.945}{11.955}{\in}%
\TextAt{14}{172.728}{150.945}{o}%
\TextAt{16}{182.821}{150.945}{\Character{106}}%
\TextAt{13}{190.391}{150.945}{x}%
\TextAt{3}{197.043}{155.283}{\Character{0}}%
\TextAt{17}{203.63}{155.283}{1}%
\NativeMathAt{212.612}{150.945}{11.955}{\notin}%
\TextAt{14}{224.831}{150.945}{o}%
\TextAt{16}{230.675}{150.945}{\Character{103}}%
\TextAt{8}{236.653}{150.945}{.}%
\TextAt{8}{246.736}{150.945}{The}%
\TextAt{8}{271.334}{150.945}{field}%
\TextAt{14}{297.236}{150.945}{o}%
\TextAt{9}{303.08}{150.945}{\Character{47}}%
\TextAt{14}{308.933}{150.945}{p}%
\TextAt{8}{319.359}{150.945}{is}%
\TextAt{8}{331.667}{150.945}{called}%
\TextAt{8}{365.376}{150.945}{the}%
\TextAt{12}{386.077}{150.945}{residue}%
\TextAt{12}{432.946}{150.945}{class}%
\TextAt{12}{464.931}{150.945}{field}%
\TextAt{8}{494.193}{150.945}{of}%
\TextAt{14}{508.067}{150.945}{o}%
\TextAt{8}{513.912}{150.945}{.}%
\TextAt{8}{523.994}{150.945}{A}%
\TextAt{8}{79.226}{133.899}{valuation}%
\TextAt{8}{130.119}{133.899}{ring}%
\TextAt{8}{154.004}{133.899}{is}%
\TextAt{8}{165.588}{133.899}{always}%
\TextAt{8}{202.875}{133.899}{integrally}%
\TextAt{8}{254.736}{133.899}{closed.}%
\TextAt{8}{293.753}{133.899}{For}%
\TextAt{8}{314.543}{133.899}{if}%
\TextAt{13}{325.086}{133.899}{x}%
\NativeMathAt{335.059}{133.899}{11.955}{\in}%
\TextAt{13}{346.35}{133.899}{K}%
\TextAt{8}{360.885}{133.899}{is}%
\TextAt{8}{372.469}{133.899}{integral}%
\TextAt{8}{414.897}{133.899}{over}%
\TextAt{14}{439.763}{133.899}{o}%
\TextAt{8}{445.607}{133.899}{,}%
\TextAt{8}{452.612}{133.899}{then}%
\TextAt{8}{479.092}{133.899}{there}%
\TextAt{8}{508.824}{133.899}{is}%
\TextAt{8}{520.408}{133.899}{an}%
\TextAt{8}{79.226}{116.853}{equation}%
\TextAt{13}{237.569}{99.807}{x}%
\TextAt{21}{244.221}{104.743}{n}%
\TextAt{9}{252.514}{99.807}{\Character{43}}%
\TextAt{13}{264.276}{99.807}{a}%
\TextAt{17}{270.421}{98.014}{1}%
\TextAt{13}{275.153}{99.807}{x}%
\TextAt{21}{281.805}{104.743}{n}%
\TextAt{3}{286.943}{104.743}{\Character{0}}%
\TextAt{17}{293.53}{104.743}{1}%
\TextAt{9}{300.919}{99.807}{\Character{43}}%
\TextAt{16}{312.68}{99.807}{\Character{1}}%
\TextAt{16}{317.976}{99.807}{\Character{1}}%
\TextAt{16}{323.284}{99.807}{\Character{1}}%
\TextAt{9}{329.284}{99.807}{\Character{43}}%
\TextAt{13}{341.046}{99.807}{a}%
\TextAt{21}{347.191}{98.014}{n}%
\TextAt{9}{356.148}{99.807}{\Character{61}}%
\TextAt{9}{368.558}{99.807}{0}%
\TextAt{8}{297.211}{43.033}{110}%
\Destination{79.23}{657.115}{nav.Item.36}%
\Destination{79.23}{639.515}{nav.Item.37}%
\Destination{79.23}{622.475}{nav.Item.38}%
\Destination{78.23}{809.815}{nav.page.110}%
\Destination{79.36}{200.055}{nav.print.121}%
\Destination{79.36}{382.055}{nav.number.2.3.8}%
\Hotlink{206.969}{500.955}{25.078}{11.047}{nav.number.2.3.1}%
}

% Page 123 (body page 111)
\Page{%
\TextAt{2}{79.226}{796.943}{II.}%
\TextAt{2}{94.073}{796.943}{The}%
\TextAt{2}{116.492}{796.943}{Theory}%
\TextAt{2}{154.401}{796.943}{of}%
\TextAt{2}{166.826}{796.943}{Valuations}%
\Vector{q 0 G 0 g 0.299 w 7.226 18.729 m 460.769 18.729 l S Q}%
\TextAt{8}{79.226}{761.124}{with}%
\TextAt{13}{106.694}{761.124}{a}%
\TextAt{21}{112.839}{759.33}{i}%
\NativeMathAt{120.917}{761.124}{11.955}{\in}%
\TextAt{14}{133.585}{761.124}{o}%
\TextAt{8}{144.134}{761.124}{and}%
\TextAt{8}{167.705}{761.124}{the}%
\TextAt{8}{188.669}{761.124}{hypothesis}%
\TextAt{13}{247.16}{761.124}{x}%
\NativeMathAt{258.51}{761.124}{11.955}{\notin}%
\TextAt{14}{271.177}{761.124}{o}%
\TextAt{8}{277.021}{761.124}{,}%
\TextAt{8}{285.18}{761.124}{so}%
\TextAt{8}{300.358}{761.124}{that}%
\TextAt{13}{326.535}{761.124}{x}%
\TextAt{3}{333.187}{765.462}{\Character{0}}%
\TextAt{17}{339.773}{765.462}{1}%
\NativeMathAt{349.203}{761.124}{11.955}{\in}%
\TextAt{14}{361.87}{761.124}{o}%
\TextAt{8}{367.715}{761.124}{,}%
\TextAt{8}{375.874}{761.124}{would}%
\TextAt{8}{410.825}{761.124}{imply}%
\TextAt{8}{444.474}{761.124}{the}%
\TextAt{8}{465.438}{761.124}{contradiction}%
\TextAt{13}{79.226}{744.078}{x}%
\TextAt{9}{89.199}{744.078}{\Character{61}}%
\TextAt{16}{101.625}{744.078}{\Character{0}}%
\TextAt{13}{110.923}{744.078}{a}%
\TextAt{17}{117.068}{742.284}{1}%
\TextAt{16}{124.036}{744.078}{\Character{0}}%
\TextAt{13}{135.571}{744.078}{a}%
\TextAt{17}{141.716}{742.284}{2}%
\TextAt{13}{146.448}{744.078}{x}%
\TextAt{3}{153.1}{748.416}{\Character{0}}%
\TextAt{17}{159.687}{748.416}{1}%
\TextAt{16}{166.655}{744.078}{\Character{0}}%
\TextAt{16}{178.18}{744.078}{\Character{1}}%
\TextAt{16}{183.476}{744.078}{\Character{1}}%
\TextAt{16}{188.784}{744.078}{\Character{1}}%
\TextAt{16}{194.331}{744.078}{\Character{0}}%
\TextAt{13}{205.908}{744.078}{a}%
\TextAt{21}{212.053}{742.284}{n}%
\TextAt{9}{217.689}{744.078}{\Character{40}}%
\TextAt{13}{222.242}{744.078}{x}%
\TextAt{3}{228.894}{748.416}{\Character{0}}%
\TextAt{17}{235.48}{748.416}{1}%
\TextAt{9}{240.213}{744.078}{\Character{41}}%
\TextAt{21}{244.765}{748.416}{n}%
\TextAt{3}{249.903}{748.416}{\Character{0}}%
\TextAt{17}{256.49}{748.416}{1}%
\NativeMathAt{264.543}{744.078}{11.955}{\in}%
\TextAt{14}{275.834}{744.078}{o}%
\TextAt{8}{281.678}{744.078}{.}%
\TextAt{8}{290.058}{744.078}{An}%
\TextAt{8}{309.028}{744.078}{exponential}%
\TextAt{8}{371.587}{744.078}{valuation}%
\TextAt{13}{422.454}{744.078}{v}%
\TextAt{8}{432.234}{744.078}{is}%
\TextAt{8}{443.792}{744.078}{called}%
\TextAt{12}{476.75}{744.078}{discrete}%
\TextAt{8}{525.943}{744.078}{if}%
\TextAt{8}{79.226}{727.032}{it}%
\TextAt{8}{90.93}{727.032}{admits}%
\TextAt{8}{129.366}{727.032}{a}%
\TextAt{8}{139.122}{727.032}{smallest}%
\TextAt{8}{184.121}{727.032}{positive}%
\TextAt{8}{227.434}{727.032}{value}%
\TextAt{13}{257.681}{727.032}{s}%
\TextAt{8}{263.195}{727.032}{.}%
\TextAt{8}{271.643}{727.032}{In}%
\TextAt{8}{286.264}{727.032}{this}%
\TextAt{8}{309.087}{727.032}{case,}%
\TextAt{8}{337.11}{727.032}{one}%
\TextAt{8}{358.569}{727.032}{finds}%
\TextAt{13}{274.385}{699.305}{v}%
\TextAt{9}{280.474}{699.305}{\Character{40}}%
\TextAt{13}{285.026}{699.305}{K}%
\TextAt{3}{295.844}{704.241}{\Character{3}}%
\TextAt{9}{300.576}{699.305}{\Character{41}}%
\TextAt{9}{308.431}{699.305}{\Character{61}}%
\TextAt{13}{320.875}{699.305}{s}%
\TextAt{18}{326.389}{699.305}{Z}%
\TextAt{13}{334.359}{699.305}{\Character{58}}%
\TextAt{8}{79.226}{671.578}{It}%
\TextAt{8}{92.995}{671.578}{is}%
\TextAt{8}{105.855}{671.578}{called}%
\TextAt{12}{140.115}{671.578}{normalized}%
\TextAt{8}{209.63}{671.578}{if}%
\TextAt{13}{221.45}{671.578}{s}%
\TextAt{9}{232.153}{671.578}{\Character{61}}%
\TextAt{9}{246.44}{671.578}{1}%
\TextAt{8}{252.299}{671.578}{.}%
\TextAt{8}{264.036}{671.578}{Dividing}%
\TextAt{8}{312.774}{671.578}{by}%
\TextAt{13}{330.129}{671.578}{s}%
\TextAt{8}{340.636}{671.578}{we}%
\TextAt{8}{358.96}{671.578}{may}%
\TextAt{8}{385.425}{671.578}{always}%
\TextAt{8}{423.989}{671.578}{pass}%
\TextAt{8}{450.573}{671.578}{to}%
\TextAt{8}{465.98}{671.578}{a}%
\TextAt{8}{476.831}{671.578}{normalized}%
\TextAt{8}{79.226}{654.532}{valuation}%
\TextAt{8}{130.299}{654.532}{without}%
\TextAt{8}{173.875}{654.532}{changing}%
\TextAt{8}{222.987}{654.532}{the}%
\TextAt{8}{243.143}{654.532}{invariants}%
\TextAt{14}{296.882}{654.532}{o}%
\TextAt{8}{302.726}{654.532}{,}%
\TextAt{14}{309.875}{654.532}{o}%
\TextAt{3}{315.72}{658.87}{\Character{3}}%
\TextAt{8}{320.452}{654.532}{,}%
\TextAt{14}{327.601}{654.532}{p}%
\TextAt{8}{333.584}{654.532}{.}%
\TextAt{8}{342.033}{654.532}{Having}%
\TextAt{8}{382.035}{654.532}{done}%
\TextAt{8}{409.998}{654.532}{so,}%
\TextAt{8}{427.62}{654.532}{an}%
\TextAt{8}{443.879}{654.532}{element}%
\TextAt{13}{222.46}{626.805}{\Character{25}}%
\NativeMathAt{232.85}{626.805}{11.955}{\in}%
\TextAt{14}{244.141}{626.805}{o}%
\TextAt{8}{273.896}{626.805}{such}%
\TextAt{8}{300.293}{626.805}{that}%
\TextAt{13}{345.674}{626.805}{v}%
\TextAt{9}{351.763}{626.805}{\Character{40}}%
\TextAt{13}{356.315}{626.805}{\Character{25}}%
\TextAt{9}{363.384}{626.805}{\Character{41}}%
\TextAt{9}{371.239}{626.805}{\Character{61}}%
\TextAt{9}{383.66}{626.805}{1}%
\TextAt{8}{79.226}{599.078}{is}%
\TextAt{8}{90.99}{599.078}{a}%
\TextAt{12}{100.745}{599.078}{prime}%
\TextAt{12}{139.289}{599.078}{element}%
\TextAt{8}{184.982}{599.078}{,}%
\TextAt{8}{192.131}{599.078}{and}%
\TextAt{8}{214.893}{599.078}{every}%
\TextAt{8}{245.786}{599.078}{element}%
\TextAt{13}{289.027}{599.078}{x}%
\NativeMathAt{299}{599.078}{11.955}{\in}%
\TextAt{13}{310.291}{599.078}{K}%
\TextAt{3}{321.109}{603.416}{\Character{3}}%
\TextAt{8}{329.739}{599.078}{admits}%
\TextAt{8}{368.175}{599.078}{a}%
\TextAt{8}{377.93}{599.078}{unique}%
\TextAt{8}{415.971}{599.078}{representation}%
\TextAt{13}{283.938}{571.351}{x}%
\TextAt{9}{293.911}{571.351}{\Character{61}}%
\TextAt{13}{306.337}{571.351}{u\Character{25}}%
\TextAt{21}{320.069}{576.287}{m}%
\TextAt{8}{79.226}{543.624}{with}%
\TextAt{13}{105.886}{543.624}{m}%
\NativeMathAt{119.446}{543.624}{11.955}{\in}%
\TextAt{18}{130.737}{543.624}{Z}%
\TextAt{8}{142.605}{543.624}{and}%
\TextAt{13}{165.367}{543.624}{u}%
\NativeMathAt{175.351}{543.624}{11.955}{\in}%
\TextAt{14}{186.642}{543.624}{o}%
\TextAt{3}{192.486}{547.962}{\Character{3}}%
\TextAt{8}{197.218}{543.624}{.}%
\TextAt{8}{205.667}{543.624}{For}%
\TextAt{8}{226.636}{543.624}{if}%
\TextAt{13}{237.36}{543.624}{v}%
\TextAt{9}{243.448}{543.624}{\Character{40}}%
\TextAt{13}{248}{543.624}{x}%
\TextAt{9}{254.652}{543.624}{\Character{41}}%
\TextAt{9}{262.507}{543.624}{\Character{61}}%
\TextAt{13}{274.951}{543.624}{m}%
\TextAt{8}{285.19}{543.624}{,}%
\TextAt{8}{292.339}{543.624}{then}%
\TextAt{13}{318.999}{543.624}{v}%
\TextAt{9}{325.088}{543.624}{\Character{40}}%
\TextAt{13}{329.64}{543.624}{x\Character{25}}%
\TextAt{3}{343.361}{547.962}{\Character{0}}%
\TextAt{21}{349.948}{547.962}{m}%
\TextAt{9}{357.937}{543.624}{\Character{41}}%
\TextAt{9}{365.792}{543.624}{\Character{61}}%
\TextAt{9}{378.213}{543.624}{0}%
\TextAt{8}{384.088}{543.624}{,}%
\TextAt{8}{391.237}{543.624}{hence}%
\TextAt{13}{423.743}{543.624}{u}%
\TextAt{9}{433.727}{543.624}{\Character{61}}%
\TextAt{13}{446.152}{543.624}{x\Character{25}}%
\TextAt{3}{459.874}{547.962}{\Character{0}}%
\TextAt{21}{466.46}{547.962}{m}%
\NativeMathAt{477.77}{543.624}{11.955}{\in}%
\TextAt{14}{489.061}{543.624}{o}%
\TextAt{3}{494.905}{547.962}{\Character{3}}%
\TextAt{8}{499.637}{543.624}{.}%
\TextAt{12}{79.226}{517.975}{Proposition}%
\TextAt{12}{151.937}{517.975}{(3.9).}%
\TextAt{8}{187.253}{517.975}{If}%
\TextAt{13}{198.945}{517.975}{v}%
\TextAt{8}{208.931}{517.975}{is}%
\TextAt{8}{220.695}{517.975}{a}%
\TextAt{8}{230.45}{517.975}{discrete}%
\TextAt{8}{273.429}{517.975}{exponential}%
\TextAt{8}{336.194}{517.975}{valuation}%
\TextAt{8}{387.266}{517.975}{of}%
\TextAt{13}{400.596}{517.975}{K}%
\TextAt{8}{411.414}{517.975}{,}%
\TextAt{8}{418.563}{517.975}{then}%
\TextAt{14}{246.384}{490.248}{o}%
\TextAt{9}{255.549}{490.248}{\Character{61}}%
\TextAt{16}{267.974}{490.248}{\Character{102}}%
\TextAt{13}{273.952}{490.248}{x}%
\NativeMathAt{283.925}{490.248}{11.955}{\in}%
\TextAt{13}{295.216}{490.248}{K}%
\TextAt{16}{309.354}{490.248}{\Character{106}}%
\TextAt{13}{315.996}{490.248}{v}%
\TextAt{9}{322.084}{490.248}{\Character{40}}%
\TextAt{13}{326.637}{490.248}{x}%
\TextAt{9}{333.289}{490.248}{\Character{41}}%
\NativeMathAt{341.162}{490.248}{11.955}{\geq}%
\TextAt{9}{353.781}{490.248}{0}%
\TextAt{16}{359.634}{490.248}{\Character{103}}%
\TextAt{8}{79.226}{462.521}{is}%
\TextAt{8}{91.156}{462.521}{a}%
\TextAt{8}{101.076}{462.521}{principal}%
\TextAt{8}{150.019}{462.521}{ideal}%
\TextAt{8}{178.148}{462.521}{domain,}%
\TextAt{8}{223.234}{462.521}{hence}%
\TextAt{8}{255.906}{462.521}{a}%
\TextAt{12}{265.827}{462.521}{discrete}%
\TextAt{12}{316.002}{462.521}{valuation}%
\TextAt{12}{374.497}{462.521}{ring}%
\TextAt{8}{402.004}{462.521}{(see}%
\TextAt{8}{425.638}{462.521}{I,}%
\TextAt{8}{437.173}{462.521}{(11.3)).}%
\TextAt{8}{480.608}{462.521}{Suppose}%
\TextAt{13}{526.681}{462.521}{v}%
\TextAt{8}{79.226}{445.475}{is}%
\TextAt{8}{90.99}{445.475}{normalized.}%
\TextAt{8}{155.377}{445.475}{Then}%
\TextAt{8}{185.934}{445.475}{the}%
\TextAt{8}{206.09}{445.475}{nonzero}%
\TextAt{8}{249.667}{445.475}{ideals}%
\TextAt{8}{282.245}{445.475}{of}%
\TextAt{14}{295.575}{445.475}{o}%
\TextAt{8}{305.316}{445.475}{are}%
\TextAt{8}{324.827}{445.475}{given}%
\TextAt{8}{355.397}{445.475}{by}%
\TextAt{14}{195.187}{417.748}{p}%
\TextAt{21}{201.171}{422.684}{n}%
\TextAt{9}{210.128}{417.748}{\Character{61}}%
\TextAt{13}{222.553}{417.748}{\Character{25}}%
\TextAt{21}{229.622}{422.684}{n}%
\TextAt{14}{235.259}{417.748}{o}%
\TextAt{9}{244.424}{417.748}{\Character{61}}%
\TextAt{16}{256.849}{417.748}{\Character{102}}%
\TextAt{13}{262.827}{417.748}{x}%
\NativeMathAt{272.8}{417.748}{11.955}{\in}%
\TextAt{13}{284.091}{417.748}{K}%
\TextAt{16}{298.229}{417.748}{\Character{106}}%
\TextAt{13}{304.871}{417.748}{v}%
\TextAt{9}{310.959}{417.748}{\Character{40}}%
\TextAt{13}{315.512}{417.748}{x}%
\TextAt{9}{322.164}{417.748}{\Character{41}}%
\NativeMathAt{330.037}{417.748}{11.955}{\geq}%
\TextAt{13}{342.656}{417.748}{n}%
\TextAt{16}{349.644}{417.748}{\Character{103}}%
\TextAt{13}{355.621}{417.748}{\Character{59}}%
\TextAt{13}{384.756}{417.748}{n}%
\NativeMathAt{395.084}{417.748}{11.955}{\geq}%
\TextAt{9}{407.704}{417.748}{0}%
\TextAt{13}{413.557}{417.748}{\Character{59}}%
\TextAt{8}{79.226}{390.021}{where}%
\TextAt{13}{113.035}{390.021}{\Character{25}}%
\TextAt{8}{124.002}{390.021}{is}%
\TextAt{8}{135.766}{390.021}{a}%
\TextAt{8}{145.521}{390.021}{prime}%
\TextAt{8}{178.685}{390.021}{element,}%
\TextAt{8}{225.179}{390.021}{i.e.,}%
\TextAt{13}{247.284}{390.021}{v}%
\TextAt{9}{253.372}{390.021}{\Character{40}}%
\TextAt{13}{257.924}{390.021}{\Character{25}}%
\TextAt{9}{264.994}{390.021}{\Character{41}}%
\TextAt{9}{272.849}{390.021}{\Character{61}}%
\TextAt{9}{285.27}{390.021}{1}%
\TextAt{8}{291.145}{390.021}{.}%
\TextAt{8}{299.594}{390.021}{One}%
\TextAt{8}{324.305}{390.021}{has}%
\TextAt{14}{267.606}{362.294}{p}%
\TextAt{21}{273.589}{367.23}{n}%
\TextAt{9}{279.225}{362.294}{\Character{47}}%
\TextAt{14}{285.078}{362.294}{p}%
\TextAt{21}{291.061}{367.23}{n}%
\TextAt{17}{296.199}{367.23}{\Character{43}1}%
\NativeMathAt{310.839}{362.294}{11.955}{\cong}%
\TextAt{14}{323.458}{362.294}{o}%
\TextAt{9}{329.302}{362.294}{\Character{47}}%
\TextAt{14}{335.155}{362.294}{p}%
\TextAt{13}{341.138}{362.294}{\Character{58}}%
\TextAt{19}{79.226}{317.362}{Proof.}%
\TextAt{8}{116}{317.362}{Let}%
\TextAt{14}{137.019}{317.362}{a}%
\NativeMathAt{146.407}{317.362}{11.955}{\neq}%
\TextAt{9}{158.9}{317.362}{0}%
\TextAt{8}{168.717}{317.362}{be}%
\TextAt{8}{184.691}{317.362}{an}%
\TextAt{8}{200.999}{317.362}{ideal}%
\TextAt{8}{229.012}{317.362}{of}%
\TextAt{14}{242.391}{317.362}{o}%
\TextAt{8}{252.183}{317.362}{and}%
\TextAt{13}{274.995}{317.362}{x}%
\NativeMathAt{285.052}{317.362}{11.955}{\neq}%
\TextAt{9}{297.545}{317.362}{0}%
\TextAt{8}{307.362}{317.362}{an}%
\TextAt{8}{323.67}{317.362}{element}%
\TextAt{8}{366.961}{317.362}{in}%
\TextAt{14}{380.664}{317.362}{a}%
\TextAt{8}{390.594}{317.362}{with}%
\TextAt{8}{417.303}{317.362}{smallest}%
\TextAt{8}{462.352}{317.362}{possible}%
\TextAt{8}{506.42}{317.362}{value}%
\TextAt{13}{79.226}{300.316}{v}%
\TextAt{9}{85.314}{300.316}{\Character{40}}%
\TextAt{13}{89.867}{300.316}{x}%
\TextAt{9}{96.519}{300.316}{\Character{41}}%
\TextAt{9}{104.374}{300.316}{\Character{61}}%
\TextAt{13}{116.817}{300.316}{n}%
\TextAt{8}{123.805}{300.316}{.}%
\TextAt{8}{132.015}{300.316}{Then}%
\TextAt{13}{161.859}{300.316}{x}%
\TextAt{9}{171.831}{300.316}{\Character{61}}%
\TextAt{13}{184.257}{300.316}{u\Character{25}}%
\TextAt{21}{197.989}{304.655}{n}%
\TextAt{8}{203.625}{300.316}{,}%
\TextAt{13}{210.203}{300.316}{u}%
\NativeMathAt{220.186}{300.316}{11.955}{\in}%
\TextAt{14}{231.477}{300.316}{o}%
\TextAt{3}{237.321}{304.655}{\Character{3}}%
\TextAt{8}{242.054}{300.316}{,}%
\TextAt{8}{248.632}{300.316}{so}%
\TextAt{8}{262.287}{300.316}{that}%
\TextAt{13}{286.942}{300.316}{\Character{25}}%
\TextAt{21}{294.011}{304.655}{n}%
\TextAt{14}{299.648}{300.316}{o}%
\NativeMathAt{308.813}{300.316}{11.955}{\subseteq}%
\TextAt{14}{321.432}{300.316}{a}%
\TextAt{8}{327.415}{300.316}{.}%
\TextAt{8}{335.626}{300.316}{If}%
\TextAt{13}{346.604}{300.316}{y}%
\TextAt{9}{356.061}{300.316}{\Character{61}}%
\TextAt{13}{368.487}{300.316}{\Character{34}\Character{25}}%
\TextAt{21}{381.034}{304.655}{m}%
\NativeMathAt{392.343}{300.316}{11.955}{\in}%
\TextAt{14}{403.634}{300.316}{a}%
\TextAt{8}{412.801}{300.316}{is}%
\TextAt{8}{423.85}{300.316}{arbitrary}%
\TextAt{8}{472.906}{300.316}{with}%
\TextAt{13}{498.851}{300.316}{\Character{34}}%
\NativeMathAt{507.65}{300.316}{11.955}{\in}%
\TextAt{14}{518.941}{300.316}{o}%
\TextAt{3}{524.785}{304.655}{\Character{3}}%
\TextAt{8}{529.518}{300.316}{,}%
\TextAt{8}{79.226}{283.27}{then}%
\TextAt{13}{105.886}{283.27}{m}%
\TextAt{9}{119.446}{283.27}{\Character{61}}%
\TextAt{13}{131.872}{283.27}{v}%
\TextAt{9}{137.96}{283.27}{\Character{40}}%
\TextAt{13}{142.512}{283.27}{y}%
\TextAt{9}{148.649}{283.27}{\Character{41}}%
\NativeMathAt{156.522}{283.27}{11.955}{\geq}%
\TextAt{13}{169.141}{283.27}{n}%
\TextAt{8}{176.129}{283.27}{,}%
\TextAt{8}{183.278}{283.27}{hence}%
\TextAt{13}{215.784}{283.27}{y}%
\TextAt{9}{225.242}{283.27}{\Character{61}}%
\TextAt{9}{237.652}{283.27}{\Character{40}}%
\TextAt{13}{242.22}{283.27}{\Character{34}\Character{25}}%
\TextAt{21}{254.767}{287.609}{m}%
\TextAt{3}{262.257}{287.609}{\Character{0}}%
\TextAt{21}{268.844}{287.609}{n}%
\TextAt{9}{274.48}{283.27}{\Character{41}}%
\TextAt{13}{279.033}{283.27}{\Character{25}}%
\TextAt{21}{286.102}{287.609}{n}%
\NativeMathAt{295.059}{283.27}{11.955}{\in}%
\TextAt{13}{306.35}{283.27}{\Character{25}}%
\TextAt{21}{313.419}{287.609}{n}%
\TextAt{14}{319.056}{283.27}{o}%
\TextAt{8}{324.9}{283.27}{,}%
\TextAt{8}{332.049}{283.27}{so}%
\TextAt{8}{346.419}{283.27}{that}%
\TextAt{14}{371.788}{283.27}{a}%
\TextAt{9}{381.092}{283.27}{\Character{61}}%
\TextAt{13}{393.517}{283.27}{\Character{25}}%
\TextAt{21}{400.587}{287.609}{n}%
\TextAt{14}{406.223}{283.27}{o}%
\TextAt{8}{412.067}{283.27}{.}%
\TextAt{8}{420.516}{283.27}{The}%
\TextAt{8}{444.569}{283.27}{isomorphism}%
\TextAt{14}{269.231}{255.543}{p}%
\TextAt{21}{275.214}{260.479}{n}%
\TextAt{9}{280.851}{255.543}{\Character{47}}%
\TextAt{14}{286.704}{255.543}{p}%
\TextAt{21}{292.687}{260.479}{n}%
\TextAt{17}{297.825}{260.479}{\Character{43}1}%
\NativeMathAt{312.465}{255.543}{11.955}{\cong}%
\TextAt{14}{325.084}{255.543}{o}%
\TextAt{9}{330.928}{255.543}{\Character{47}}%
\TextAt{14}{336.781}{255.543}{p}%
\TextAt{8}{79.226}{227.816}{results}%
\TextAt{8}{116.419}{227.816}{from}%
\TextAt{8}{144.059}{227.816}{the}%
\TextAt{8}{164.215}{227.816}{correspondence}%
\TextAt{13}{245.893}{227.816}{a\Character{25}}%
\TextAt{21}{259.107}{232.155}{n}%
\NativeArrowAt{268.065}{227.816}{11.955}{19.26}{\longmapsto}%
\TextAt{13}{290.647}{227.816}{a}%
\TextAt{8}{300.112}{227.816}{mod}%
\TextAt{14}{325.55}{227.816}{p}%
\TextAt{8}{331.533}{227.816}{.}%
\Vector{q 0 G 0 g 0.399 w 452.334 -546.184 m 452.334 -538.114 l S Q}%
\Vector{q 0 G 0 g 0.399 w 452.534 -538.313 m 459.707 -538.313 l S Q}%
\Vector{q 0 G 0 g 0.399 w 452.534 -545.985 m 459.707 -545.985 l S Q}%
\Vector{q 0 G 0 g 0.399 w 459.906 -546.184 m 459.906 -538.114 l S Q}%
\TextAt{8}{97.159}{199.179}{In}%
\TextAt{8}{111.78}{199.179}{a}%
\TextAt{8}{121.536}{199.179}{discretely}%
\TextAt{8}{173.947}{199.179}{valued}%
\TextAt{8}{210.697}{199.179}{field}%
\TextAt{13}{236.054}{199.179}{K}%
\TextAt{8}{250.769}{199.179}{the}%
\TextAt{8}{270.926}{199.179}{chain}%
\TextAt{14}{250.515}{171.452}{o}%
\NativeMathAt{259.68}{171.452}{11.955}{\supseteq}%
\TextAt{14}{272.299}{171.452}{p}%
\NativeMathAt{281.603}{171.452}{11.955}{\supseteq}%
\TextAt{14}{294.222}{171.452}{p}%
\TextAt{17}{300.205}{176.388}{2}%
\NativeMathAt{308.259}{171.452}{11.955}{\supseteq}%
\TextAt{14}{320.878}{171.452}{p}%
\TextAt{17}{326.861}{176.388}{3}%
\NativeMathAt{334.914}{171.452}{11.955}{\supseteq}%
\TextAt{16}{347.515}{171.452}{\Character{1}}%
\TextAt{16}{352.823}{171.452}{\Character{1}}%
\TextAt{16}{358.131}{171.452}{\Character{1}}%
\TextAt{8}{79.226}{143.725}{consisting}%
\TextAt{8}{134.45}{143.725}{of}%
\TextAt{8}{148.894}{143.725}{the}%
\TextAt{8}{170.165}{143.725}{ideals}%
\TextAt{8}{203.857}{143.725}{of}%
\TextAt{8}{218.301}{143.725}{the}%
\TextAt{8}{239.572}{143.725}{valuation}%
\TextAt{8}{291.759}{143.725}{ring}%
\TextAt{14}{316.939}{143.725}{o}%
\TextAt{8}{327.795}{143.725}{forms}%
\TextAt{8}{361.164}{143.725}{a}%
\TextAt{8}{372.034}{143.725}{basis}%
\TextAt{8}{401.889}{143.725}{of}%
\TextAt{8}{416.333}{143.725}{neighbourhoods}%
\TextAt{8}{502.066}{143.725}{of}%
\TextAt{8}{516.51}{143.725}{the}%
\TextAt{8}{79.226}{126.679}{zero}%
\TextAt{8}{103.99}{126.679}{element.}%
\TextAt{8}{151.94}{126.679}{Indeed,}%
\TextAt{8}{193.287}{126.679}{if}%
\TextAt{13}{204.063}{126.679}{v}%
\TextAt{8}{214.101}{126.679}{is}%
\TextAt{8}{225.918}{126.679}{a}%
\TextAt{8}{235.726}{126.679}{normalized}%
\TextAt{8}{295.614}{126.679}{exponential}%
\TextAt{8}{358.431}{126.679}{valuation}%
\TextAt{8}{409.556}{126.679}{and}%
\TextAt{16}{432.371}{126.679}{\Character{106}}%
\TextAt{16}{439.628}{126.679}{\Character{106}}%
\TextAt{9}{446.373}{126.679}{\Character{61}}%
\TextAt{13}{458.888}{126.679}{q}%
\TextAt{3}{464.507}{131.018}{\Character{0}}%
\TextAt{21}{471.093}{131.018}{v}%
\TextAt{8}{479.951}{126.679}{(}%
\TextAt{13}{484.506}{126.679}{q}%
\TextAt{13}{493.532}{126.679}{\Character{62}}%
\TextAt{9}{506.05}{126.679}{1}%
\TextAt{8}{511.903}{126.679}{)}%
\TextAt{8}{520.408}{126.679}{an}%
\TextAt{8}{79.226}{109.633}{associated}%
\TextAt{8}{134.961}{109.633}{multiplicative}%
\TextAt{8}{208.784}{109.633}{valuation,}%
\TextAt{8}{263.109}{109.633}{then}%
\TextAt{14}{232.619}{78.301}{p}%
\TextAt{21}{238.602}{83.238}{n}%
\TextAt{9}{247.559}{78.301}{\Character{61}}%
\TextAt{20}{259.984}{95.158}{\Character{26}}%
\TextAt{13}{268.951}{78.301}{x}%
\NativeMathAt{278.924}{78.301}{11.955}{\in}%
\TextAt{13}{290.215}{78.301}{K}%
\TextAt{20}{304.353}{95.636}{\Character{12}}%
\TextAt{20}{304.353}{88.463}{\Character{12}}%
\TextAt{20}{304.353}{81.29}{\Character{12}}%
\TextAt{20}{304.353}{74.117}{\Character{12}}%
\TextAt{16}{311.659}{78.301}{\Character{106}}%
\TextAt{13}{314.98}{78.301}{x}%
\TextAt{16}{321.632}{78.301}{\Character{106}}%
\TextAt{13}{328.274}{78.301}{\Character{60}}%
\TextAt{9}{350.006}{86.389}{1}%
\Vector{q 0 G 0 g 0.478 w 269.895 -692.71 m 291.971 -692.71 l S Q}%
\TextAt{13}{341.895}{70.101}{q}%
\TextAt{21}{347.514}{73.554}{n}%
\TextAt{3}{352.652}{73.554}{\Character{0}}%
\TextAt{17}{359.239}{73.554}{1}%
\TextAt{20}{365.166}{95.158}{\Character{27}}%
\TextAt{13}{376.125}{78.301}{\Character{58}}%
\TextAt{8}{297.211}{43.033}{111}%
\Destination{78.23}{809.815}{nav.page.111}%
\Destination{79.36}{531.955}{nav.number.2.3.9}%
\Hotlink{436.7}{460.3}{29.21}{11.047}{nav.number.1.11.3}%
}

% Page 124 (body page 112)
\Page{%
\TextAt{2}{79.226}{796.943}{II.}%
\TextAt{2}{94.073}{796.943}{The}%
\TextAt{2}{116.492}{796.943}{Theory}%
\TextAt{2}{154.401}{796.943}{of}%
\TextAt{2}{166.826}{796.943}{Valuations}%
\Vector{q 0 G 0 g 0.299 w 7.226 18.729 m 460.769 18.729 l S Q}%
\TextAt{8}{79.226}{761.124}{As}%
\TextAt{8}{97.702}{761.124}{a}%
\TextAt{8}{108.646}{761.124}{basis}%
\TextAt{8}{138.575}{761.124}{of}%
\TextAt{8}{153.094}{761.124}{neighbourhoods}%
\TextAt{8}{238.902}{761.124}{of}%
\TextAt{8}{253.42}{761.124}{the}%
\TextAt{8}{274.766}{761.124}{element}%
\TextAt{8}{319.196}{761.124}{1}%
\TextAt{8}{330.141}{761.124}{of}%
\TextAt{13}{344.659}{761.124}{K}%
\TextAt{3}{355.477}{765.462}{\Character{3}}%
\TextAt{8}{360.209}{761.124}{,}%
\TextAt{8}{368.845}{761.124}{we}%
\TextAt{8}{387.261}{761.124}{obtain}%
\TextAt{8}{424.877}{761.124}{in}%
\TextAt{8}{439.719}{761.124}{the}%
\TextAt{8}{461.064}{761.124}{same}%
\TextAt{8}{491.579}{761.124}{way}%
\TextAt{8}{516.51}{761.124}{the}%
\TextAt{8}{79.226}{744.078}{descending}%
\TextAt{8}{138.464}{744.078}{chain}%
\TextAt{14}{231.176}{727.032}{o}%
\TextAt{3}{237.02}{731.968}{\Character{3}}%
\TextAt{9}{245.073}{727.032}{\Character{61}}%
\TextAt{13}{257.498}{727.032}{U}%
\TextAt{17}{266.697}{731.968}{\Character{40}0\Character{41}}%
\NativeMathAt{281.337}{727.032}{11.955}{\supseteq}%
\TextAt{13}{293.956}{727.032}{U}%
\TextAt{17}{303.155}{731.968}{\Character{40}1\Character{41}}%
\NativeMathAt{317.795}{727.032}{11.955}{\supseteq}%
\TextAt{13}{330.414}{727.032}{U}%
\TextAt{17}{339.614}{731.968}{\Character{40}2\Character{41}}%
\NativeMathAt{354.253}{727.032}{11.955}{\supseteq}%
\TextAt{16}{366.854}{727.032}{\Character{1}}%
\TextAt{16}{372.162}{727.032}{\Character{1}}%
\TextAt{16}{377.47}{727.032}{\Character{1}}%
\TextAt{8}{79.226}{705.772}{of}%
\TextAt{8}{92.556}{705.772}{subgroups}%
\TextAt{13}{162.43}{678.312}{U}%
\TextAt{17}{171.629}{683.248}{\Character{40}}%
\TextAt{21}{174.922}{683.248}{n}%
\TextAt{17}{180.06}{683.248}{\Character{41}}%
\TextAt{9}{187.173}{678.312}{\Character{61}}%
\TextAt{9}{199.583}{678.312}{1}%
\TextAt{9}{208.083}{678.312}{\Character{43}}%
\TextAt{14}{219.869}{678.312}{p}%
\TextAt{21}{225.852}{683.248}{n}%
\TextAt{9}{234.809}{678.312}{\Character{61}}%
\TextAt{20}{247.235}{695.169}{\Character{26}}%
\TextAt{13}{256.201}{678.312}{x}%
\NativeMathAt{266.174}{678.312}{11.955}{\in}%
\TextAt{13}{277.465}{678.312}{K}%
\TextAt{3}{288.283}{683.248}{\Character{3}}%
\TextAt{20}{296.336}{695.647}{\Character{12}}%
\TextAt{20}{296.336}{688.474}{\Character{12}}%
\TextAt{20}{296.336}{681.301}{\Character{12}}%
\TextAt{20}{296.336}{674.128}{\Character{12}}%
\TextAt{16}{303.642}{678.312}{\Character{106}}%
\TextAt{9}{306.963}{678.312}{1}%
\TextAt{16}{315.472}{678.312}{\Character{0}}%
\TextAt{13}{327.427}{678.312}{x}%
\TextAt{16}{334.08}{678.312}{\Character{106}}%
\TextAt{13}{340.721}{678.312}{\Character{60}}%
\TextAt{9}{362.454}{686.4}{1}%
\Vector{q 0 G 0 g 0.478 w 282.342 -92.699 m 304.418 -92.699 l S Q}%
\TextAt{13}{354.342}{670.112}{q}%
\TextAt{21}{359.961}{673.565}{n}%
\TextAt{3}{365.1}{673.565}{\Character{0}}%
\TextAt{17}{371.686}{673.565}{1}%
\TextAt{20}{377.614}{695.169}{\Character{27}}%
\TextAt{13}{388.573}{678.312}{\Character{59}}%
\TextAt{13}{417.708}{678.312}{n}%
\TextAt{13}{428.013}{678.312}{\Character{62}}%
\TextAt{9}{440.461}{678.312}{0}%
\TextAt{13}{446.314}{678.312}{\Character{59}}%
\TextAt{8}{79.226}{648.965}{of}%
\TextAt{14}{92.285}{648.965}{o}%
\TextAt{3}{98.13}{653.303}{\Character{3}}%
\TextAt{8}{102.862}{648.965}{.}%
\TextAt{8}{111.22}{648.965}{(Observe}%
\TextAt{8}{160.443}{648.965}{that}%
\TextAt{9}{185.541}{648.965}{1}%
\TextAt{9}{193.479}{648.965}{\Character{43}}%
\TextAt{14}{204.705}{648.965}{p}%
\TextAt{21}{210.688}{653.303}{n}%
\TextAt{8}{219.951}{648.965}{is}%
\TextAt{8}{231.444}{648.965}{closed}%
\TextAt{8}{265.699}{648.965}{under}%
\TextAt{8}{298.592}{648.965}{multiplication}%
\TextAt{8}{373.448}{648.965}{and}%
\TextAt{8}{395.939}{648.965}{that,}%
\TextAt{8}{424.343}{648.965}{if}%
\TextAt{13}{434.796}{648.965}{x}%
\NativeMathAt{444.769}{648.965}{11.955}{\in}%
\TextAt{13}{456.06}{648.965}{U}%
\TextAt{17}{465.259}{653.303}{\Character{40}}%
\TextAt{21}{468.552}{653.303}{n}%
\TextAt{17}{473.691}{653.303}{\Character{41}}%
\TextAt{8}{477.482}{648.965}{,}%
\TextAt{8}{484.415}{648.965}{then}%
\TextAt{8}{510.804}{648.965}{so}%
\TextAt{8}{524.903}{648.965}{is}%
\TextAt{13}{79.226}{631.919}{x}%
\TextAt{3}{85.878}{636.257}{\Character{0}}%
\TextAt{17}{92.465}{636.257}{1}%
\TextAt{8}{101.136}{631.919}{because}%
\TextAt{16}{144.48}{631.919}{\Character{106}}%
\TextAt{9}{147.801}{631.919}{1}%
\TextAt{16}{156.339}{631.919}{\Character{0}}%
\TextAt{13}{168.323}{631.919}{x}%
\TextAt{3}{174.975}{636.257}{\Character{0}}%
\TextAt{17}{181.562}{636.257}{1}%
\TextAt{16}{186.294}{631.919}{\Character{106}}%
\TextAt{9}{193.007}{631.919}{\Character{61}}%
\TextAt{16}{205.504}{631.919}{\Character{106}}%
\TextAt{13}{208.825}{631.919}{x}%
\TextAt{3}{215.477}{636.257}{\Character{0}}%
\TextAt{17}{222.064}{636.257}{1}%
\TextAt{16}{226.796}{631.919}{\Character{106}\Character{106}}%
\TextAt{13}{233.438}{631.919}{x}%
\TextAt{16}{242.775}{631.919}{\Character{0}}%
\TextAt{9}{254.759}{631.919}{1}%
\TextAt{16}{260.612}{631.919}{\Character{106}}%
\TextAt{9}{267.325}{631.919}{\Character{61}}%
\TextAt{16}{279.822}{631.919}{\Character{106}}%
\TextAt{9}{283.143}{631.919}{1}%
\TextAt{16}{291.681}{631.919}{\Character{0}}%
\TextAt{13}{303.665}{631.919}{x}%
\TextAt{16}{310.317}{631.919}{\Character{106}}%
\TextAt{13}{317.03}{631.919}{\Character{60}}%
\TextAt{9}{329.527}{631.919}{1\Character{47}}%
\TextAt{13}{341.233}{631.919}{q}%
\TextAt{21}{346.853}{636.257}{n}%
\TextAt{3}{351.991}{636.257}{\Character{0}}%
\TextAt{17}{358.577}{636.257}{1}%
\TextAt{8}{363.31}{631.919}{.)}%
\TextAt{13}{376.439}{631.919}{U}%
\TextAt{17}{385.638}{636.257}{\Character{40}}%
\TextAt{21}{388.931}{636.257}{n}%
\TextAt{17}{394.069}{636.257}{\Character{41}}%
\TextAt{8}{401.8}{631.919}{is}%
\TextAt{8}{413.606}{631.919}{called}%
\TextAt{8}{446.811}{631.919}{the}%
\TextAt{13}{467.01}{631.919}{n}%
\TextAt{12}{473.997}{631.919}{-th}%
\TextAt{12}{495.72}{631.919}{higher}%
\TextAt{12}{79.226}{614.873}{unit}%
\TextAt{12}{107.78}{614.873}{group}%
\TextAt{8}{145.711}{614.873}{and}%
\TextAt{13}{168.602}{614.873}{U}%
\TextAt{17}{177.801}{619.211}{\Character{40}1\Character{41}}%
\TextAt{8}{193.146}{614.873}{the}%
\TextAt{8}{213.432}{614.873}{group}%
\TextAt{8}{246.736}{614.873}{of}%
\TextAt{12}{260.709}{614.873}{principal}%
\TextAt{12}{316.999}{614.873}{units}%
\TextAt{8}{346.229}{614.873}{.}%
\TextAt{8}{355.064}{614.873}{Regarding}%
\TextAt{8}{411.298}{614.873}{the}%
\TextAt{8}{431.584}{614.873}{successive}%
\TextAt{8}{485.869}{614.873}{quotients}%
\TextAt{8}{79.226}{597.827}{of}%
\TextAt{8}{92.556}{597.827}{the}%
\TextAt{8}{112.713}{597.827}{chain}%
\TextAt{8}{143.605}{597.827}{of}%
\TextAt{8}{156.935}{597.827}{higher}%
\TextAt{8}{192.705}{597.827}{unit}%
\TextAt{8}{217.416}{597.827}{groups,}%
\TextAt{8}{258.458}{597.827}{we}%
\TextAt{8}{275.685}{597.827}{have}%
\TextAt{8}{302.68}{597.827}{the}%
\TextAt{12}{79.226}{573.323}{Proposition}%
\TextAt{12}{151.937}{573.323}{(3.10).}%
\TextAt{14}{193.984}{573.323}{o}%
\TextAt{3}{199.828}{577.661}{\Character{3}}%
\TextAt{9}{204.56}{573.323}{\Character{47}}%
\TextAt{13}{210.413}{573.323}{U}%
\TextAt{17}{219.612}{577.661}{\Character{40}}%
\TextAt{21}{222.906}{577.661}{n}%
\TextAt{17}{228.044}{577.661}{\Character{41}}%
\NativeMathAt{235.156}{573.323}{11.955}{\cong}%
\TextAt{9}{247.775}{573.323}{\Character{40}}%
\TextAt{14}{252.328}{573.323}{o}%
\TextAt{9}{258.172}{573.323}{\Character{47}}%
\TextAt{14}{264.025}{573.323}{p}%
\TextAt{21}{270.008}{577.661}{n}%
\TextAt{9}{275.644}{573.323}{\Character{41}}%
\TextAt{3}{280.197}{577.661}{\Character{3}}%
\TextAt{8}{288.826}{573.323}{and}%
\TextAt{13}{311.589}{573.323}{U}%
\TextAt{17}{320.788}{577.661}{\Character{40}}%
\TextAt{21}{324.081}{577.661}{n}%
\TextAt{17}{329.22}{577.661}{\Character{41}}%
\TextAt{9}{333.011}{573.323}{\Character{47}}%
\TextAt{13}{338.864}{573.323}{U}%
\TextAt{17}{348.063}{577.661}{\Character{40}}%
\TextAt{21}{351.356}{577.661}{n}%
\TextAt{17}{356.494}{577.661}{\Character{43}1\Character{41}}%
\NativeMathAt{374.427}{573.323}{11.955}{\cong}%
\TextAt{14}{387.047}{573.323}{o}%
\TextAt{9}{392.891}{573.323}{\Character{47}}%
\TextAt{14}{398.744}{573.323}{p}%
\TextAt{8}{404.727}{573.323}{,}%
\TextAt{8}{411.876}{573.323}{for}%
\TextAt{13}{429.761}{573.323}{n}%
\NativeMathAt{440.07}{573.323}{11.955}{\geq}%
\TextAt{9}{452.689}{573.323}{1}%
\TextAt{8}{458.542}{573.323}{.}%
\TextAt{19}{79.226}{541.36}{Proof.}%
\TextAt{8}{116}{541.36}{The}%
\TextAt{8}{139.892}{541.36}{first}%
\TextAt{8}{163.856}{541.36}{isomorphism}%
\TextAt{8}{232.113}{541.36}{is}%
\TextAt{8}{243.715}{541.36}{induced}%
\TextAt{8}{287.118}{541.36}{by}%
\TextAt{8}{303.215}{541.36}{the}%
\TextAt{8}{323.209}{541.36}{canonical}%
\TextAt{8}{374.431}{541.36}{and}%
\TextAt{8}{397.032}{541.36}{obviously}%
\TextAt{8}{448.647}{541.36}{surjective}%
\TextAt{8}{500.897}{541.36}{homo-}%
\TextAt{8}{79.226}{524.314}{morphism}%
\TextAt{14}{213.175}{507.268}{o}%
\TextAt{3}{219.02}{512.204}{\Character{3}}%
\NativeArrowAt{227.073}{507.268}{11.955}{19.26}{\longrightarrow}%
\TextAt{9}{249.655}{507.268}{\Character{40}}%
\TextAt{14}{254.207}{507.268}{o}%
\TextAt{9}{260.051}{507.268}{\Character{47}}%
\TextAt{14}{265.904}{507.268}{p}%
\TextAt{21}{271.887}{512.204}{n}%
\TextAt{9}{277.524}{507.268}{\Character{41}}%
\TextAt{3}{282.076}{512.204}{\Character{3}}%
\TextAt{13}{286.808}{507.268}{\Character{59}}%
\TextAt{13}{315.943}{507.268}{u}%
\NativeArrowAt{325.946}{507.268}{11.955}{19.26}{\longmapsto}%
\TextAt{13}{348.528}{507.268}{u}%
\TextAt{8}{358.511}{507.268}{mod}%
\TextAt{14}{383.949}{507.268}{p}%
\TextAt{21}{389.932}{512.204}{n}%
\TextAt{13}{395.569}{507.268}{\Character{59}}%
\TextAt{8}{79.226}{486.008}{the}%
\TextAt{8}{100.228}{486.008}{kernel}%
\TextAt{8}{135.541}{486.008}{of}%
\TextAt{8}{149.717}{486.008}{which}%
\TextAt{8}{184.049}{486.008}{is}%
\TextAt{13}{196.659}{486.008}{U}%
\TextAt{17}{205.858}{490.346}{\Character{40}}%
\TextAt{21}{209.151}{490.346}{n}%
\TextAt{17}{214.289}{490.346}{\Character{41}}%
\TextAt{8}{218.081}{486.008}{.}%
\TextAt{8}{229.067}{486.008}{The}%
\TextAt{8}{253.966}{486.008}{second}%
\TextAt{8}{292.591}{486.008}{isomorphism}%
\TextAt{8}{361.856}{486.008}{is}%
\TextAt{8}{374.466}{486.008}{given,}%
\TextAt{8}{409.344}{486.008}{once}%
\TextAt{8}{436.85}{486.008}{we}%
\TextAt{8}{454.923}{486.008}{choose}%
\TextAt{8}{492.902}{486.008}{a}%
\TextAt{8}{503.503}{486.008}{prime}%
\TextAt{8}{79.226}{468.962}{element}%
\TextAt{13}{122.468}{468.962}{\Character{25}}%
\TextAt{8}{129.537}{468.962}{,}%
\TextAt{8}{136.686}{468.962}{by}%
\TextAt{8}{152.946}{468.962}{the}%
\TextAt{8}{173.102}{468.962}{surjective}%
\TextAt{8}{225.513}{468.962}{homomorphism}%
\TextAt{13}{177.053}{445.239}{U}%
\TextAt{17}{186.252}{450.175}{\Character{40}}%
\TextAt{21}{189.545}{450.175}{n}%
\TextAt{17}{194.684}{450.175}{\Character{41}}%
\TextAt{9}{201.796}{445.239}{\Character{61}}%
\TextAt{9}{214.206}{445.239}{1}%
\TextAt{9}{222.706}{445.239}{\Character{43}}%
\TextAt{13}{234.492}{445.239}{\Character{25}}%
\TextAt{21}{241.562}{450.175}{n}%
\TextAt{14}{247.198}{445.239}{o}%
\NativeArrowAt{256.363}{445.239}{11.955}{19.26}{\longrightarrow}%
\TextAt{14}{278.945}{445.239}{o}%
\TextAt{9}{284.789}{445.239}{\Character{47}}%
\TextAt{14}{290.642}{445.239}{p}%
\TextAt{13}{296.625}{445.239}{\Character{59}}%
\TextAt{9}{325.78}{445.239}{1}%
\TextAt{9}{334.28}{445.239}{\Character{43}}%
\TextAt{13}{346.051}{445.239}{\Character{25}}%
\TextAt{21}{353.12}{450.175}{n}%
\TextAt{13}{358.756}{445.239}{a}%
\NativeArrowAt{368.222}{445.239}{11.955}{19.26}{\longmapsto}%
\TextAt{13}{390.804}{445.239}{a}%
\TextAt{8}{400.27}{445.239}{mod}%
\TextAt{14}{425.708}{445.239}{p}%
\TextAt{13}{431.691}{445.239}{\Character{59}}%
\TextAt{8}{79.226}{421.515}{which}%
\TextAt{8}{112.713}{421.515}{has}%
\TextAt{8}{133.586}{421.515}{kernel}%
\TextAt{13}{168.053}{421.515}{U}%
\TextAt{17}{177.252}{425.854}{\Character{40}}%
\TextAt{21}{180.545}{425.854}{n}%
\TextAt{17}{185.684}{425.854}{\Character{43}1\Character{41}}%
\TextAt{8}{200.296}{421.515}{.}%
\Vector{q 0 G 0 g 0.399 w 452.334 -352.485 m 452.334 -344.415 l S Q}%
\Vector{q 0 G 0 g 0.399 w 452.534 -344.614 m 459.707 -344.614 l S Q}%
\Vector{q 0 G 0 g 0.399 w 452.534 -352.286 m 459.707 -352.286 l S Q}%
\Vector{q 0 G 0 g 0.399 w 459.906 -352.485 m 459.906 -344.415 l S Q}%
\TextAt{12}{79.226}{389.552}{Exercise}%
\TextAt{12}{132.419}{389.552}{II.3.1.}%
\TextAt{8}{171.385}{389.552}{Show}%
\TextAt{8}{202.27}{389.552}{that}%
\TextAt{16}{227.632}{389.552}{\Character{106}}%
\TextAt{13}{230.953}{389.552}{z}%
\TextAt{16}{236.924}{389.552}{\Character{106}}%
\TextAt{9}{243.565}{389.552}{\Character{61}}%
\TextAt{9}{255.975}{389.552}{\Character{40}}%
\TextAt{13}{260.543}{389.552}{z}%
\TextAt{9}{267.223}{389.552}{\Character{22}}%
\TextAt{13}{266.514}{389.552}{z}%
\TextAt{9}{272.484}{389.552}{\Character{41}}%
\TextAt{17}{277.037}{393.891}{1\Character{47}2}%
\TextAt{9}{293.558}{389.552}{\Character{61}}%
\TextAt{20}{305.984}{399.342}{\Character{112}}%
\Vector{q 0 G 0 g 0.478 w 245.939 -374.419 m 292.786 -374.419 l S Q}%
\TextAt{16}{317.939}{389.552}{\Character{106}}%
\TextAt{8}{323.252}{389.552}{N}%
\TextAt{25}{332.027}{387.693}{C}%
\TextAt{3}{337.784}{387.693}{\Character{106}}%
\TextAt{25}{340.136}{387.693}{R}%
\TextAt{9}{346.39}{389.552}{\Character{40}}%
\TextAt{13}{350.943}{389.552}{z}%
\TextAt{9}{356.913}{389.552}{\Character{41}}%
\TextAt{16}{361.466}{389.552}{\Character{106}}%
\TextAt{8}{368.677}{389.552}{is}%
\TextAt{8}{380.434}{389.552}{the}%
\TextAt{8}{400.583}{389.552}{only}%
\TextAt{8}{426.267}{389.552}{valuation}%
\TextAt{8}{477.333}{389.552}{of}%
\TextAt{18}{490.656}{389.552}{C}%
\TextAt{8}{503.18}{389.552}{which}%
\TextAt{8}{79.226}{372.506}{extends}%
\TextAt{8}{121.882}{372.506}{the}%
\TextAt{8}{142.039}{372.506}{absolute}%
\TextAt{8}{188.281}{372.506}{value}%
\TextAt{16}{218.528}{372.506}{\Character{106}}%
\TextAt{16}{225.725}{372.506}{\Character{106}}%
\TextAt{8}{232.964}{372.506}{of}%
\TextAt{18}{246.294}{372.506}{R}%
\TextAt{8}{254.929}{372.506}{.}%
\TextAt{12}{79.226}{343.532}{Exercise}%
\TextAt{12}{132.888}{343.532}{II.3.2.}%
\TextAt{8}{172.264}{343.532}{What}%
\TextAt{8}{205.506}{343.532}{is}%
\TextAt{8}{217.671}{343.532}{the}%
\TextAt{8}{238.229}{343.532}{relation}%
\TextAt{8}{281.562}{343.532}{between}%
\TextAt{8}{327.154}{343.532}{the}%
\TextAt{8}{347.712}{343.532}{Chinese}%
\TextAt{8}{391.737}{343.532}{remainder}%
\TextAt{8}{447.42}{343.532}{theorem}%
\TextAt{8}{493.346}{343.532}{and}%
\TextAt{8}{516.51}{343.532}{the}%
\TextAt{8}{79.226}{326.486}{approximation}%
\TextAt{8}{157.293}{326.486}{theorem}%
\TextAt{8}{202.819}{326.486}{(3.4)?}%
\TextAt{12}{79.226}{297.512}{Exercise}%
\TextAt{12}{132.65}{297.512}{II.3.3.}%
\TextAt{8}{171.819}{297.512}{Let}%
\TextAt{13}{192.983}{297.512}{k}%
\TextAt{8}{203.564}{297.512}{be}%
\TextAt{8}{219.683}{297.512}{a}%
\TextAt{8}{229.633}{297.512}{field}%
\TextAt{8}{255.185}{297.512}{and}%
\TextAt{13}{278.142}{297.512}{K}%
\TextAt{9}{292.612}{297.512}{\Character{61}}%
\TextAt{13}{305.369}{297.512}{k}%
\TextAt{9}{311.858}{297.512}{\Character{40}}%
\TextAt{13}{316.411}{297.512}{t}%
\TextAt{9}{320.638}{297.512}{\Character{41}}%
\TextAt{8}{329.282}{297.512}{the}%
\TextAt{8}{349.633}{297.512}{function}%
\TextAt{8}{395.676}{297.512}{field}%
\TextAt{8}{421.227}{297.512}{in}%
\TextAt{8}{435.075}{297.512}{one}%
\TextAt{8}{456.729}{297.512}{variable.}%
\TextAt{8}{505.775}{297.512}{Show}%
\TextAt{8}{79.226}{280.466}{that}%
\TextAt{8}{105.901}{280.466}{the}%
\TextAt{8}{127.364}{280.466}{valuations}%
\TextAt{13}{184.358}{280.466}{v}%
\TextAt{27}{190.017}{278.673}{p}%
\TextAt{8}{199.707}{280.466}{associated}%
\TextAt{8}{256.749}{280.466}{to}%
\TextAt{8}{272.365}{280.466}{the}%
\TextAt{8}{293.828}{280.466}{prime}%
\TextAt{8}{328.298}{280.466}{ideals}%
\TextAt{14}{362.182}{280.466}{p}%
\TextAt{9}{373.712}{280.466}{\Character{61}}%
\TextAt{9}{388.357}{280.466}{\Character{40}}%
\TextAt{13}{392.916}{280.466}{p}%
\TextAt{9}{398.791}{280.466}{\Character{40}}%
\TextAt{13}{403.344}{280.466}{t}%
\TextAt{9}{407.571}{280.466}{\Character{41}\Character{41}}%
\TextAt{8}{421.879}{280.466}{of}%
\TextAt{13}{436.516}{280.466}{k}%
\TextAt{9}{443.005}{280.466}{\Character{91}}%
\TextAt{13}{446.257}{280.466}{t}%
\TextAt{9}{450.484}{280.466}{\Character{93}}%
\TextAt{8}{453.736}{280.466}{,}%
\TextAt{8}{462.518}{280.466}{together}%
\TextAt{8}{510.007}{280.466}{with}%
\ErrataAt{79.22}{263.455}{453.64}{the degree valuation $v_\infty$, are, up to equivalence, all nontrivial valuations on $K$ trivial}%
\ErrataAt{79.22}{246.355}{453.64}{on $k$. What are the residue class fields?}%
\TextAt{12}{79.226}{217.4}{Exercise}%
\TextAt{12}{132.921}{217.4}{II.3.4.}%
\TextAt{8}{172.324}{217.4}{Let}%
\TextAt{14}{193.723}{217.4}{o}%
\TextAt{8}{203.894}{217.4}{be}%
\TextAt{8}{220.248}{217.4}{an}%
\TextAt{8}{236.937}{217.4}{arbitrary}%
\TextAt{8}{287.135}{217.4}{valuation}%
\TextAt{8}{338.637}{217.4}{ring}%
\TextAt{8}{363.133}{217.4}{with}%
\TextAt{8}{390.222}{217.4}{field}%
\TextAt{8}{416.009}{217.4}{of}%
\TextAt{8}{429.768}{217.4}{fractions}%
\TextAt{13}{478.066}{217.4}{K}%
\TextAt{8}{488.884}{217.4}{,}%
\TextAt{8}{496.57}{217.4}{and}%
\TextAt{8}{519.762}{217.4}{let}%
\TextAt{9}{79.226}{200.354}{\Character{0}}%
\TextAt{9}{90.703}{200.354}{\Character{61}}%
\TextAt{13}{104}{200.354}{K}%
\TextAt{3}{114.817}{204.693}{\Character{3}}%
\TextAt{9}{119.55}{200.354}{\Character{47}}%
\TextAt{14}{125.403}{200.354}{o}%
\TextAt{3}{131.247}{204.693}{\Character{3}}%
\TextAt{8}{135.979}{200.354}{.}%
\TextAt{8}{145.936}{200.354}{Then}%
\TextAt{9}{176.996}{200.354}{\Character{0}}%
\TextAt{8}{188.71}{200.354}{becomes}%
\TextAt{8}{235.766}{200.354}{a}%
\TextAt{8}{246.024}{200.354}{totally}%
\TextAt{8}{283.934}{200.354}{ordered}%
\TextAt{8}{326.71}{200.354}{group}%
\TextAt{8}{360.389}{200.354}{if}%
\TextAt{8}{371.615}{200.354}{we}%
\TextAt{8}{389.345}{200.354}{define}%
\TextAt{13}{423.657}{200.354}{x}%
\TextAt{8}{433.63}{200.354}{mod}%
\TextAt{14}{459.068}{200.354}{o}%
\TextAt{3}{464.912}{204.693}{\Character{3}}%
\NativeMathAt{473.822}{200.354}{11.955}{\geq}%
\TextAt{13}{487.298}{200.354}{y}%
\TextAt{8}{496.755}{200.354}{mod}%
\TextAt{14}{522.193}{200.354}{o}%
\TextAt{3}{528.037}{204.693}{\Character{3}}%
\TextAt{8}{79.226}{183.308}{to}%
\TextAt{8}{93.536}{183.308}{mean}%
\TextAt{13}{124.751}{183.308}{x}%
\TextAt{9}{131.404}{183.308}{\Character{47}}%
\TextAt{13}{137.257}{183.308}{y}%
\NativeMathAt{146.714}{183.308}{11.955}{\in}%
\TextAt{14}{158.005}{183.308}{o}%
\TextAt{8}{163.849}{183.308}{.}%
\TextAt{8}{172.298}{183.308}{Write}%
\TextAt{9}{204.804}{183.308}{\Character{0}}%
\TextAt{8}{216.015}{183.308}{additively}%
\TextAt{8}{270.328}{183.308}{and}%
\TextAt{8}{293.09}{183.308}{show}%
\TextAt{8}{322.093}{183.308}{that}%
\TextAt{8}{347.462}{183.308}{the}%
\TextAt{8}{367.619}{183.308}{function}%
\TextAt{13}{255.767}{159.585}{v}%
\TextAt{9}{265.176}{159.585}{\Character{58}}%
\TextAt{13}{271.748}{159.585}{K}%
\NativeArrowAt{285.887}{159.585}{11.955}{19.26}{\longrightarrow}%
\TextAt{9}{308.469}{159.585}{\Character{0}}%
\TextAt{16}{318.44}{159.585}{\Character{91}}%
\TextAt{16}{329.056}{159.585}{\Character{102}}%
\NativeMathAt{335.034}{159.585}{11.955}{\infty}%
\TextAt{16}{346.989}{159.585}{\Character{103}}%
\TextAt{13}{352.977}{159.585}{\Character{59}}%
\TextAt{13}{79.226}{135.861}{v}%
\TextAt{9}{85.314}{135.861}{\Character{40}0\Character{41}}%
\TextAt{9}{103.558}{135.861}{\Character{61}}%
\NativeMathAt{116.018}{135.861}{11.955}{\infty}%
\TextAt{8}{127.974}{135.861}{,}%
\TextAt{13}{135.123}{135.861}{v}%
\TextAt{9}{141.211}{135.861}{\Character{40}}%
\TextAt{13}{145.763}{135.861}{x}%
\TextAt{9}{152.415}{135.861}{\Character{41}}%
\TextAt{9}{160.27}{135.861}{\Character{61}}%
\TextAt{13}{172.714}{135.861}{x}%
\TextAt{8}{182.687}{135.861}{mod}%
\TextAt{14}{208.125}{135.861}{o}%
\TextAt{3}{213.969}{140.2}{\Character{3}}%
\TextAt{8}{222.599}{135.861}{for}%
\TextAt{13}{240.484}{135.861}{x}%
\NativeMathAt{250.457}{135.861}{11.955}{\in}%
\TextAt{13}{261.748}{135.861}{K}%
\TextAt{3}{272.565}{140.2}{\Character{3}}%
\TextAt{8}{277.298}{135.861}{,}%
\TextAt{8}{284.447}{135.861}{satisfies}%
\TextAt{8}{327.557}{135.861}{the}%
\TextAt{8}{347.714}{135.861}{conditions}%
\TextAt{8}{297.211}{43.033}{112}%
\Destination{78.23}{809.815}{nav.page.112}%
\Destination{79.226}{403.555}{nav.ex.2.3.1}%
\Destination{79.226}{357.555}{nav.ex.2.3.2}%
\Destination{79.226}{311.555}{nav.ex.2.3.3}%
\Destination{79.226}{231.355}{nav.ex.2.3.4}%
\Destination{79.36}{776.155}{nav.print.122}%
\Destination{79.36}{404.555}{nav.print.123}%
\Destination{79.36}{587.355}{nav.number.2.3.10}%
\Destination{79.226}{277.455}{errata.13}%
\Hotlink{202.36}{324.3}{23.35}{11.047}{nav.number.2.3.4}%
}

% Page 125 (body page 113)
\Page{%
\TextAt{2}{79.226}{796.943}{II.}%
\TextAt{2}{94.073}{796.943}{The}%
\TextAt{2}{116.492}{796.943}{Theory}%
\TextAt{2}{154.401}{796.943}{of}%
\TextAt{2}{166.826}{796.943}{Valuations}%
\Vector{q 0 G 0 g 0.299 w 7.226 18.729 m 460.769 18.729 l S Q}%
\TextAt{8}{91.528}{761.124}{1)}%
\TextAt{13}{107.918}{761.124}{v}%
\TextAt{9}{114.007}{761.124}{\Character{40}}%
\TextAt{13}{118.559}{761.124}{x}%
\TextAt{9}{125.211}{761.124}{\Character{41}}%
\TextAt{9}{133.066}{761.124}{\Character{61}}%
\NativeMathAt{145.51}{761.124}{11.955}{\infty}%
\NativeArrowAt{160.777}{761.124}{11.955}{21.938}{\Longleftrightarrow}%
\TextAt{13}{186.024}{761.124}{x}%
\TextAt{9}{195.997}{761.124}{\Character{61}}%
\TextAt{9}{208.407}{761.124}{0}%
\TextAt{8}{214.276}{761.124}{,}%
\TextAt{8}{91.528}{744.078}{2)}%
\TextAt{13}{107.918}{744.078}{v}%
\TextAt{9}{114.007}{744.078}{\Character{40}}%
\TextAt{13}{118.559}{744.078}{xy}%
\TextAt{9}{131.348}{744.078}{\Character{41}}%
\TextAt{9}{139.203}{744.078}{\Character{61}}%
\TextAt{13}{151.646}{744.078}{v}%
\TextAt{9}{157.735}{744.078}{\Character{40}}%
\TextAt{13}{162.287}{744.078}{x}%
\TextAt{9}{168.939}{744.078}{\Character{41}}%
\TextAt{9}{176.136}{744.078}{\Character{43}}%
\TextAt{13}{187.909}{744.078}{v}%
\TextAt{9}{193.998}{744.078}{\Character{40}}%
\TextAt{13}{198.55}{744.078}{y}%
\TextAt{9}{204.687}{744.078}{\Character{41}}%
\TextAt{8}{209.239}{744.078}{,}%
\TextAt{8}{91.528}{727.032}{3)}%
\TextAt{13}{107.918}{727.032}{v}%
\TextAt{9}{114.007}{727.032}{\Character{40}}%
\TextAt{13}{118.559}{727.032}{x}%
\TextAt{9}{127.868}{727.032}{\Character{43}}%
\TextAt{13}{139.629}{727.032}{y}%
\TextAt{9}{145.766}{727.032}{\Character{41}}%
\NativeMathAt{153.639}{727.032}{11.955}{\geq}%
\TextAt{8}{166.258}{727.032}{min}%
\TextAt{16}{185.769}{727.032}{\Character{102}}%
\TextAt{13}{191.747}{727.032}{v}%
\TextAt{9}{197.835}{727.032}{\Character{40}}%
\TextAt{13}{202.387}{727.032}{x}%
\TextAt{9}{209.039}{727.032}{\Character{41}}%
\TextAt{13}{213.592}{727.032}{\Character{59}}%
\TextAt{13}{218.816}{727.032}{v}%
\TextAt{9}{224.924}{727.032}{\Character{40}}%
\TextAt{13}{229.476}{727.032}{y}%
\TextAt{9}{235.613}{727.032}{\Character{41}}%
\TextAt{16}{240.165}{727.032}{\Character{103}}%
\TextAt{8}{246.143}{727.032}{.}%
\TextAt{13}{79.226}{706.997}{v}%
\TextAt{8}{89.212}{706.997}{is}%
\TextAt{8}{100.976}{706.997}{called}%
\TextAt{8}{134.139}{706.997}{a}%
\TextAt{12}{143.895}{706.997}{Krull}%
\TextAt{12}{179.366}{706.997}{valuation}%
\TextAt{8}{233.188}{706.997}{.}%
\TextAt{15}{79.226}{654.133}{§}%
\TextAt{15}{93.785}{654.133}{4}%
\TextAt{15}{120.693}{654.133}{Completions}%
\TextAt{12}{79.226}{616.27}{Definition}%
\TextAt{12}{143.071}{616.27}{(4.1).}%
\TextAt{8}{179.304}{616.27}{A}%
\TextAt{8}{192.894}{616.27}{valued}%
\TextAt{8}{230.562}{616.27}{field}%
\TextAt{9}{256.837}{616.27}{\Character{40}}%
\TextAt{13}{261.39}{616.27}{K}%
\TextAt{13}{271.552}{616.27}{\Character{59}}%
\TextAt{16}{276.801}{616.27}{\Character{106}}%
\TextAt{16}{284.919}{616.27}{\Character{106}}%
\TextAt{9}{288.258}{616.27}{\Character{41}}%
\TextAt{8}{297.626}{616.27}{is}%
\TextAt{8}{310.308}{616.27}{called}%
\TextAt{12}{344.39}{616.27}{complete}%
\TextAt{8}{401.844}{616.27}{if}%
\TextAt{8}{413.486}{616.27}{every}%
\TextAt{8}{445.296}{616.27}{Cauchy}%
\TextAt{8}{488.165}{616.27}{sequence}%
\TextAt{16}{79.226}{599.224}{\Character{102}}%
\TextAt{13}{85.204}{599.224}{a}%
\TextAt{21}{91.349}{597.431}{n}%
\TextAt{16}{96.985}{599.224}{\Character{103}}%
\TextAt{21}{102.963}{597.431}{n}%
\NativeMathAt{108.101}{597.431}{7.97}{\in}%
\TextAt{25}{113.746}{597.431}{N}%
\TextAt{8}{123.898}{599.224}{in}%
\TextAt{13}{137.551}{599.224}{K}%
\TextAt{8}{152.266}{599.224}{converges}%
\TextAt{8}{204.689}{599.224}{to}%
\TextAt{8}{219}{599.224}{an}%
\TextAt{8}{235.259}{599.224}{element}%
\TextAt{13}{278.501}{599.224}{a}%
\NativeMathAt{287.966}{599.224}{11.955}{\in}%
\TextAt{13}{299.257}{599.224}{K}%
\TextAt{8}{310.075}{599.224}{,}%
\TextAt{8}{317.224}{599.224}{i.e.,}%
\TextAt{8}{264.857}{570.223}{lim}%
\TextAt{21}{261.949}{563.05}{n}%
\NativeMathAt{267.087}{563.05}{7.97}{\rightarrow}%
\NativeMathAt{275.551}{563.05}{7.97}{\infty}%
\TextAt{16}{286.016}{570.223}{\Character{106}}%
\TextAt{13}{289.337}{570.223}{a}%
\TextAt{21}{295.482}{568.43}{n}%
\TextAt{16}{303.775}{570.223}{\Character{0}}%
\TextAt{13}{315.73}{570.223}{a}%
\TextAt{16}{321.875}{570.223}{\Character{106}}%
\TextAt{9}{328.517}{570.223}{\Character{61}}%
\TextAt{9}{340.927}{570.223}{0}%
\TextAt{13}{346.795}{570.223}{\Character{58}}%
\TextAt{8}{97.159}{530.97}{Here,}%
\TextAt{8}{128.794}{530.97}{as}%
\TextAt{8}{143.768}{530.97}{usual,}%
\TextAt{8}{178.404}{530.97}{we}%
\TextAt{8}{196.236}{530.97}{call}%
\TextAt{16}{218.3}{530.97}{\Character{102}}%
\TextAt{13}{224.277}{530.97}{a}%
\TextAt{21}{230.422}{529.177}{n}%
\TextAt{16}{236.058}{530.97}{\Character{103}}%
\TextAt{21}{242.036}{529.177}{n}%
\NativeMathAt{247.174}{529.177}{7.97}{\in}%
\TextAt{25}{252.82}{529.177}{N}%
\TextAt{8}{263.576}{530.97}{a}%
\TextAt{12}{273.935}{530.97}{Cauchy}%
\TextAt{12}{322.654}{530.97}{sequence}%
\TextAt{8}{378.885}{530.97}{if}%
\TextAt{8}{390.213}{530.97}{for}%
\TextAt{8}{408.702}{530.97}{every}%
\TextAt{13}{440.199}{530.97}{\Character{34}}%
\TextAt{13}{450.014}{530.97}{\Character{62}}%
\TextAt{9}{463.482}{530.97}{0}%
\TextAt{8}{473.836}{530.97}{there}%
\TextAt{8}{504.352}{530.97}{exists}%
\TextAt{13}{79.226}{513.924}{N}%
\NativeMathAt{93.17}{513.924}{11.955}{\in}%
\TextAt{18}{104.461}{513.924}{N}%
\TextAt{8}{116.992}{513.924}{such}%
\TextAt{8}{143.389}{513.924}{that}%
\TextAt{16}{216.957}{496.878}{\Character{106}}%
\TextAt{13}{220.278}{496.878}{a}%
\TextAt{21}{226.423}{495.085}{n}%
\TextAt{16}{234.716}{496.878}{\Character{0}}%
\TextAt{13}{246.671}{496.878}{a}%
\TextAt{21}{252.816}{495.085}{m}%
\TextAt{16}{260.804}{496.878}{\Character{106}}%
\TextAt{13}{267.446}{496.878}{\Character{60}}%
\TextAt{13}{279.856}{496.878}{\Character{34}}%
\TextAt{8}{309.26}{496.878}{for}%
\TextAt{8}{327.145}{496.878}{all}%
\TextAt{13}{343.404}{496.878}{n\Character{59}}%
\TextAt{13}{355.61}{496.878}{m}%
\NativeMathAt{369.196}{496.878}{11.955}{\geq}%
\TextAt{13}{381.815}{496.878}{N}%
\TextAt{13}{391.786}{496.878}{\Character{58}}%
\TextAt{8}{79.226}{473.356}{From}%
\TextAt{8}{110.857}{473.356}{any}%
\TextAt{8}{133.88}{473.356}{valued}%
\TextAt{8}{171.536}{473.356}{field}%
\TextAt{9}{197.799}{473.356}{\Character{40}}%
\TextAt{13}{202.351}{473.356}{K}%
\TextAt{13}{212.513}{473.356}{\Character{59}}%
\TextAt{16}{217.762}{473.356}{\Character{106}}%
\TextAt{16}{225.868}{473.356}{\Character{106}}%
\TextAt{9}{229.207}{473.356}{\Character{41}}%
\TextAt{8}{238.563}{473.356}{we}%
\TextAt{8}{256.696}{473.356}{get}%
\TextAt{8}{277.113}{473.356}{a}%
\TextAt{8}{287.774}{473.356}{complete}%
\TextAt{8}{338.103}{473.356}{valued}%
\TextAt{8}{375.759}{473.356}{field}%
\TextAt{9}{402.022}{473.356}{\Character{40}}%
\TextAt{20}{409.312}{476.378}{\Character{98}}%
\TextAt{13}{406.574}{473.356}{K}%
\TextAt{13}{417.382}{473.356}{\Character{59}}%
\TextAt{16}{422.636}{473.356}{\Character{106}}%
\TextAt{16}{430.742}{473.356}{\Character{106}}%
\TextAt{9}{434.081}{473.356}{\Character{41}}%
\TextAt{8}{443.437}{473.356}{by}%
\TextAt{8}{460.602}{473.356}{the}%
\TextAt{8}{481.664}{473.356}{process}%
\TextAt{8}{523.337}{473.356}{of}%
\TextAt{12}{79.226}{456.31}{completion}%
\TextAt{8}{143.676}{456.31}{.}%
\TextAt{8}{152.038}{456.31}{This}%
\TextAt{8}{178.497}{456.31}{completion}%
\TextAt{8}{238.072}{456.31}{is}%
\TextAt{8}{249.575}{456.31}{obtained}%
\TextAt{8}{297.446}{456.31}{in}%
\TextAt{8}{310.838}{456.31}{the}%
\TextAt{8}{330.733}{456.31}{same}%
\TextAt{8}{359.798}{456.31}{way}%
\TextAt{8}{383.281}{456.31}{as}%
\TextAt{8}{397.39}{456.31}{the}%
\TextAt{8}{417.286}{456.31}{field}%
\TextAt{8}{442.382}{456.31}{of}%
\TextAt{8}{455.451}{456.31}{real}%
\TextAt{8}{477.953}{456.31}{numbers}%
\TextAt{8}{524.903}{456.31}{is}%
\TextAt{8}{79.226}{439.264}{constructed}%
\TextAt{8}{142.523}{439.264}{from}%
\TextAt{8}{170.312}{439.264}{the}%
\TextAt{8}{190.618}{439.264}{field}%
\TextAt{8}{216.124}{439.264}{of}%
\TextAt{8}{229.604}{439.264}{rational}%
\TextAt{8}{273.341}{439.264}{numbers.}%
\TextAt{8}{325.551}{439.264}{Take}%
\TextAt{8}{353.986}{439.264}{the}%
\TextAt{8}{374.292}{439.264}{ring}%
\TextAt{13}{398.507}{439.264}{R}%
\TextAt{8}{411.562}{439.264}{of}%
\TextAt{8}{425.042}{439.264}{all}%
\TextAt{8}{441.45}{439.264}{Cauchy}%
\TextAt{8}{483.55}{439.264}{sequences}%
\TextAt{8}{79.226}{422.218}{of}%
\TextAt{9}{92.991}{422.218}{\Character{40}}%
\TextAt{13}{97.543}{422.218}{K}%
\TextAt{13}{107.705}{422.218}{\Character{59}}%
\TextAt{16}{112.954}{422.218}{\Character{106}}%
\TextAt{16}{120.593}{422.218}{\Character{106}}%
\TextAt{9}{123.928}{422.218}{\Character{41}}%
\TextAt{8}{128.48}{422.218}{,}%
\TextAt{8}{136.172}{422.218}{consider}%
\TextAt{8}{182.192}{422.218}{therein}%
\TextAt{8}{222.294}{422.218}{the}%
\TextAt{8}{242.884}{422.218}{maximal}%
\TextAt{8}{291.128}{422.218}{ideal}%
\TextAt{14}{319.525}{422.218}{m}%
\TextAt{8}{333.024}{422.218}{of}%
\TextAt{8}{346.788}{422.218}{all}%
\TextAt{8}{363.481}{422.218}{nullsequences}%
\TextAt{8}{436.221}{422.218}{with}%
\TextAt{8}{463.315}{422.218}{respect}%
\TextAt{8}{503.799}{422.218}{to}%
\TextAt{16}{518.544}{422.218}{\Character{106}}%
\TextAt{16}{526.183}{422.218}{\Character{106}}%
\TextAt{8}{529.518}{422.218}{,}%
\TextAt{8}{79.226}{405.172}{and}%
\TextAt{8}{101.989}{405.172}{define}%
\TextAt{20}{281.814}{391.148}{\Character{98}}%
\TextAt{13}{279.076}{388.126}{K}%
\TextAt{9}{293.214}{388.126}{\Character{61}}%
\TextAt{13}{305.64}{388.126}{R}%
\TextAt{9}{314.648}{388.126}{\Character{47}}%
\TextAt{14}{320.501}{388.126}{m}%
\TextAt{13}{329.668}{388.126}{\Character{58}}%
\TextAt{8}{79.226}{364.605}{One}%
\TextAt{8}{105.215}{364.605}{embeds}%
\TextAt{8}{148.169}{364.605}{the}%
\TextAt{8}{169.603}{364.605}{field}%
\TextAt{13}{196.237}{364.605}{K}%
\TextAt{8}{212.23}{364.605}{into}%
\TextAt{20}{239.989}{367.627}{\Character{98}}%
\TextAt{13}{237.251}{364.605}{K}%
\TextAt{8}{253.243}{364.605}{by}%
\TextAt{8}{270.78}{364.605}{sending}%
\TextAt{8}{314.391}{364.605}{every}%
\TextAt{13}{346.561}{364.605}{a}%
\NativeMathAt{358.204}{364.605}{11.955}{\in}%
\TextAt{13}{371.672}{364.605}{K}%
\TextAt{8}{387.665}{364.605}{to}%
\TextAt{8}{403.253}{364.605}{the}%
\TextAt{8}{424.687}{364.605}{class}%
\TextAt{8}{453.402}{364.605}{of}%
\TextAt{8}{468.01}{364.605}{the}%
\TextAt{8}{489.444}{364.605}{constant}%
\TextAt{8}{79.226}{347.559}{Cauchy}%
\TextAt{8}{121.639}{347.559}{sequence}%
\TextAt{9}{170.603}{347.559}{\Character{40}}%
\TextAt{13}{175.156}{347.559}{a\Character{59}}%
\TextAt{13}{186.525}{347.559}{a\Character{59}}%
\TextAt{13}{197.907}{347.559}{a\Character{59}}%
\TextAt{13}{209.276}{347.559}{\Character{58}}%
\TextAt{13}{214.513}{347.559}{\Character{58}}%
\TextAt{13}{219.749}{347.559}{\Character{58}}%
\TextAt{9}{223.063}{347.559}{\Character{41}}%
\TextAt{8}{227.615}{347.559}{.}%
\TextAt{8}{237.45}{347.559}{The}%
\TextAt{8}{261.966}{347.559}{valuation}%
\TextAt{16}{313.501}{347.559}{\Character{106}}%
\TextAt{16}{321.164}{347.559}{\Character{106}}%
\TextAt{8}{328.862}{347.559}{is}%
\TextAt{8}{341.088}{347.559}{extended}%
\TextAt{8}{391.296}{347.559}{from}%
\TextAt{13}{419.398}{347.559}{K}%
\TextAt{8}{434.576}{347.559}{to}%
\TextAt{20}{452.086}{350.581}{\Character{98}}%
\TextAt{13}{449.348}{347.559}{K}%
\TextAt{8}{464.525}{347.559}{by}%
\TextAt{8}{481.247}{347.559}{giving}%
\TextAt{8}{516.51}{347.559}{the}%
\TextAt{8}{79.226}{330.513}{element}%
\TextAt{13}{122.468}{330.513}{a}%
\NativeMathAt{131.934}{330.513}{11.955}{\in}%
\TextAt{20}{145.963}{333.535}{\Character{98}}%
\TextAt{13}{143.225}{330.513}{K}%
\TextAt{8}{157.94}{330.513}{which}%
\TextAt{8}{191.426}{330.513}{is}%
\TextAt{8}{203.19}{330.513}{represented}%
\TextAt{8}{265.357}{330.513}{by}%
\TextAt{8}{281.616}{330.513}{the}%
\TextAt{8}{301.772}{330.513}{Cauchy}%
\TextAt{8}{343.723}{330.513}{sequence}%
\TextAt{16}{392.225}{330.513}{\Character{102}}%
\TextAt{13}{398.203}{330.513}{a}%
\TextAt{21}{404.348}{328.72}{n}%
\TextAt{16}{409.984}{330.513}{\Character{103}}%
\TextAt{21}{415.962}{328.72}{n}%
\NativeMathAt{421.1}{328.72}{7.97}{\in}%
\TextAt{25}{426.746}{328.72}{N}%
\TextAt{8}{436.897}{330.513}{the}%
\TextAt{8}{457.054}{330.513}{absolute}%
\TextAt{8}{503.296}{330.513}{value}%
\TextAt{16}{268.86}{301.512}{\Character{106}}%
\TextAt{13}{272.181}{301.512}{a}%
\TextAt{16}{278.326}{301.512}{\Character{106}}%
\TextAt{9}{284.968}{301.512}{\Character{61}}%
\TextAt{8}{300.301}{301.512}{lim}%
\TextAt{21}{297.393}{294.339}{n}%
\NativeMathAt{302.531}{294.339}{7.97}{\rightarrow}%
\NativeMathAt{310.995}{294.339}{7.97}{\infty}%
\TextAt{16}{321.461}{301.512}{\Character{106}}%
\TextAt{13}{324.782}{301.512}{a}%
\TextAt{21}{330.926}{299.718}{n}%
\TextAt{16}{336.563}{301.512}{\Character{106}}%
\TextAt{13}{339.884}{301.512}{\Character{58}}%
\TextAt{8}{79.226}{270.03}{This}%
\TextAt{8}{105.892}{270.03}{limit}%
\TextAt{8}{133.801}{270.03}{exists}%
\TextAt{8}{166.062}{270.03}{because}%
\TextAt{20}{209.31}{280.192}{\Character{12}}%
\TextAt{20}{209.31}{273.019}{\Character{12}}%
\TextAt{16}{213.295}{270.03}{\Character{106}}%
\TextAt{13}{216.616}{270.03}{a}%
\TextAt{21}{222.761}{268.236}{n}%
\TextAt{16}{228.397}{270.03}{\Character{106}}%
\TextAt{16}{234.255}{270.03}{\Character{0}}%
\TextAt{16}{246.091}{270.03}{\Character{106}}%
\TextAt{13}{249.43}{270.03}{a}%
\TextAt{21}{255.575}{268.236}{m}%
\TextAt{16}{263.564}{270.03}{\Character{106}}%
\TextAt{20}{266.885}{280.192}{\Character{12}}%
\TextAt{20}{266.885}{273.019}{\Character{12}}%
\NativeMathAt{274.191}{270.03}{11.955}{\leq}%
\TextAt{16}{286.792}{270.03}{\Character{106}}%
\TextAt{13}{290.131}{270.03}{a}%
\TextAt{21}{296.276}{268.236}{n}%
\TextAt{16}{304.459}{270.03}{\Character{0}}%
\TextAt{13}{316.304}{270.03}{a}%
\TextAt{21}{322.449}{268.236}{m}%
\TextAt{16}{330.437}{270.03}{\Character{106}}%
\TextAt{8}{337.602}{270.03}{implies}%
\TextAt{8}{377.275}{270.03}{that}%
\TextAt{16}{402.59}{270.03}{\Character{106}}%
\TextAt{13}{405.911}{270.03}{a}%
\TextAt{21}{412.056}{268.236}{n}%
\TextAt{16}{417.692}{270.03}{\Character{106}}%
\TextAt{8}{424.856}{270.03}{is}%
\TextAt{8}{436.566}{270.03}{a}%
\TextAt{8}{446.268}{270.03}{Cauchy}%
\TextAt{8}{488.165}{270.03}{sequence}%
\TextAt{8}{79.226}{252.984}{of}%
\TextAt{8}{91.939}{252.984}{real}%
\TextAt{8}{114.084}{252.984}{numbers.}%
\TextAt{8}{165.641}{252.984}{As}%
\TextAt{8}{182.311}{252.984}{in}%
\TextAt{8}{195.346}{252.984}{the}%
\TextAt{8}{214.885}{252.984}{case}%
\TextAt{8}{239.039}{252.984}{of}%
\TextAt{8}{251.752}{252.984}{the}%
\TextAt{8}{271.291}{252.984}{field}%
\TextAt{8}{296.031}{252.984}{of}%
\TextAt{8}{308.744}{252.984}{real}%
\TextAt{8}{330.889}{252.984}{numbers,}%
\TextAt{8}{380.858}{252.984}{one}%
\TextAt{8}{401.701}{252.984}{proves}%
\TextAt{8}{437.248}{252.984}{that}%
\TextAt{20}{464.738}{256.006}{\Character{98}}%
\TextAt{13}{462}{252.984}{K}%
\TextAt{8}{476.097}{252.984}{is}%
\TextAt{8}{487.244}{252.984}{complete}%
\TextAt{8}{79.226}{235.938}{with}%
\TextAt{8}{105.81}{235.938}{respect}%
\TextAt{8}{145.783}{235.938}{to}%
\TextAt{8}{160.018}{235.938}{the}%
\TextAt{8}{180.098}{235.938}{extended}%
\TextAt{16}{229.767}{235.938}{\Character{106}}%
\TextAt{16}{236.892}{235.938}{\Character{106}}%
\TextAt{8}{240.23}{235.938}{,}%
\TextAt{8}{247.318}{235.938}{and}%
\TextAt{8}{270.004}{235.938}{that}%
\TextAt{8}{295.297}{235.938}{each}%
\TextAt{13}{321.558}{235.938}{a}%
\NativeMathAt{331.024}{235.938}{11.955}{\in}%
\TextAt{20}{345.053}{238.96}{\Character{98}}%
\TextAt{13}{342.315}{235.938}{K}%
\TextAt{8}{356.953}{235.938}{is}%
\TextAt{8}{368.641}{235.938}{a}%
\TextAt{8}{378.32}{235.938}{limit}%
\TextAt{8}{406.207}{235.938}{of}%
\TextAt{8}{419.461}{235.938}{a}%
\TextAt{8}{429.14}{235.938}{sequence}%
\TextAt{16}{477.566}{235.938}{\Character{102}}%
\TextAt{13}{483.543}{235.938}{a}%
\TextAt{21}{489.688}{234.145}{n}%
\TextAt{16}{495.325}{235.938}{\Character{103}}%
\TextAt{8}{505.123}{235.938}{in}%
\TextAt{13}{518.7}{235.938}{K}%
\TextAt{8}{529.518}{235.938}{.}%
\TextAt{8}{79.226}{218.892}{Finally}%
\TextAt{8}{119.706}{218.892}{one}%
\TextAt{8}{141.81}{218.892}{proves}%
\TextAt{8}{178.619}{218.892}{the}%
\TextAt{12}{199.421}{218.892}{uniqueness}%
\TextAt{12}{268.259}{218.892}{of}%
\TextAt{12}{284.328}{218.892}{the}%
\TextAt{12}{308.394}{218.892}{completion}%
\TextAt{9}{377.387}{218.892}{\Character{40}}%
\TextAt{20}{384.677}{221.914}{\Character{98}}%
\TextAt{13}{381.939}{218.892}{K}%
\TextAt{13}{392.747}{218.892}{\Character{59}}%
\TextAt{16}{398.001}{218.892}{\Character{106}}%
\TextAt{16}{405.844}{218.892}{\Character{106}}%
\TextAt{9}{409.185}{218.892}{\Character{41}}%
\TextAt{8}{413.737}{218.892}{:}%
\TextAt{8}{423.475}{218.892}{if}%
\TextAt{9}{434.844}{218.892}{\Character{40}}%
\TextAt{20}{442.135}{221.914}{\Character{98}}%
\TextAt{13}{439.396}{218.892}{K}%
\TextAt{3}{450.214}{223.23}{\Character{48}}%
\TextAt{13}{453.009}{218.892}{\Character{59}}%
\TextAt{16}{458.253}{218.892}{\Character{106}}%
\TextAt{16}{466.096}{218.892}{\Character{106}}%
\TextAt{3}{469.437}{223.23}{\Character{48}}%
\TextAt{9}{472.232}{218.892}{\Character{41}}%
\TextAt{8}{481.327}{218.892}{is}%
\TextAt{8}{493.736}{218.892}{another}%
\TextAt{8}{79.226}{201.846}{complete}%
\TextAt{8}{128.649}{201.846}{valued}%
\TextAt{8}{165.399}{201.846}{field}%
\TextAt{8}{190.756}{201.846}{that}%
\TextAt{8}{216.125}{201.846}{contains}%
\TextAt{9}{262.045}{201.846}{\Character{40}}%
\TextAt{13}{266.597}{201.846}{K}%
\TextAt{13}{276.759}{201.846}{\Character{59}}%
\TextAt{16}{282.008}{201.846}{\Character{106}}%
\TextAt{16}{289.217}{201.846}{\Character{106}}%
\TextAt{9}{292.548}{201.846}{\Character{41}}%
\TextAt{8}{300.997}{201.846}{as}%
\TextAt{8}{315.367}{201.846}{a}%
\TextAt{8}{325.123}{201.846}{dense}%
\TextAt{8}{357.043}{201.846}{subfield,}%
\TextAt{8}{403.274}{201.846}{then}%
\TextAt{8}{429.934}{201.846}{mapping}%
\TextAt{16}{239.371}{172.845}{\Character{106}}%
\TextAt{16}{246.58}{172.845}{\Character{106}}%
\TextAt{8}{249.91}{172.845}{-}%
\TextAt{8}{258.708}{172.845}{lim}%
\TextAt{21}{255.8}{165.672}{n}%
\NativeMathAt{260.938}{165.672}{7.97}{\rightarrow}%
\NativeMathAt{269.402}{165.672}{7.97}{\infty}%
\TextAt{13}{279.868}{172.845}{a}%
\TextAt{21}{286.012}{171.051}{n}%
\NativeArrowAt{294.97}{172.845}{11.955}{19.26}{\longmapsto}%
\TextAt{16}{317.541}{172.845}{\Character{106}}%
\TextAt{16}{324.75}{172.845}{\Character{106}}%
\TextAt{3}{328.091}{177.781}{\Character{48}}%
\TextAt{8}{330.886}{172.845}{-}%
\TextAt{8}{339.684}{172.845}{lim}%
\TextAt{21}{336.776}{165.672}{n}%
\NativeMathAt{341.914}{165.672}{7.97}{\rightarrow}%
\NativeMathAt{350.378}{165.672}{7.97}{\infty}%
\TextAt{13}{360.843}{172.845}{a}%
\TextAt{21}{366.988}{171.051}{n}%
\TextAt{8}{79.226}{139.868}{gives}%
\TextAt{8}{106.636}{139.868}{a}%
\TextAt{13}{115.12}{139.868}{K}%
\TextAt{8}{125.938}{139.868}{-isomorphism}%
\TextAt{13}{196.984}{139.868}{\Character{27}}%
\TextAt{9}{207.387}{139.868}{\Character{58}}%
\TextAt{20}{216.698}{142.89}{\Character{98}}%
\TextAt{13}{213.959}{139.868}{K}%
\NativeArrowAt{228.098}{139.868}{11.955}{19.26}{\longrightarrow}%
\TextAt{20}{253.418}{142.89}{\Character{98}}%
\TextAt{13}{250.68}{139.868}{K}%
\TextAt{3}{261.497}{144.207}{\Character{48}}%
\TextAt{8}{266.919}{139.868}{such}%
\TextAt{8}{292.045}{139.868}{that}%
\TextAt{16}{316.143}{139.868}{\Character{106}}%
\TextAt{13}{319.464}{139.868}{a}%
\TextAt{16}{325.609}{139.868}{\Character{106}}%
\TextAt{9}{332.25}{139.868}{\Character{61}}%
\TextAt{16}{344.676}{139.868}{\Character{106}}%
\TextAt{13}{347.997}{139.868}{\Character{27}}%
\TextAt{13}{355.062}{139.868}{a}%
\TextAt{16}{361.224}{139.868}{\Character{106}}%
\TextAt{3}{364.545}{144.207}{\Character{48}}%
\TextAt{8}{367.34}{139.868}{.}%
\TextAt{8}{375.365}{139.868}{The}%
\TextAt{8}{398.147}{139.868}{fields}%
\TextAt{18}{426.848}{139.868}{R}%
\TextAt{8}{438.109}{139.868}{and}%
\TextAt{18}{459.6}{139.868}{C}%
\TextAt{8}{470.861}{139.868}{are}%
\TextAt{8}{489.101}{139.868}{the}%
\TextAt{8}{507.986}{139.868}{most}%
\TextAt{8}{79.226}{122.822}{familiar}%
\TextAt{8}{122.625}{122.822}{examples}%
\TextAt{8}{173.234}{122.822}{of}%
\TextAt{8}{186.709}{122.822}{complete}%
\TextAt{8}{236.277}{122.822}{fields.}%
\TextAt{8}{271.236}{122.822}{They}%
\TextAt{8}{301.616}{122.822}{are}%
\TextAt{8}{321.273}{122.822}{complete}%
\TextAt{8}{370.841}{122.822}{with}%
\TextAt{8}{397.646}{122.822}{respect}%
\TextAt{8}{437.841}{122.822}{to}%
\TextAt{8}{452.297}{122.822}{an}%
\TextAt{8}{468.702}{122.822}{archimedean}%
\TextAt{8}{79.226}{105.776}{valuation.}%
\TextAt{8}{136.5}{105.776}{Amazingly}%
\TextAt{8}{195.582}{105.776}{enough,}%
\TextAt{8}{239.846}{105.776}{there}%
\TextAt{8}{270.308}{105.776}{are}%
\TextAt{8}{290.369}{105.776}{no}%
\TextAt{8}{307.178}{105.776}{others}%
\TextAt{8}{342.912}{105.776}{of}%
\TextAt{8}{356.793}{105.776}{this}%
\TextAt{8}{380.165}{105.776}{type.}%
\TextAt{8}{412.703}{105.776}{More}%
\TextAt{8}{443.488}{105.776}{precisely}%
\TextAt{8}{491.895}{105.776}{we}%
\TextAt{8}{509.672}{105.776}{have}%
\TextAt{8}{79.226}{88.73}{the}%
\TextAt{8}{297.211}{43.033}{113}%
\Destination{79.23}{773.075}{nav.Item.39}%
\Destination{79.23}{758.125}{nav.Item.40}%
\Destination{79.23}{741.075}{nav.Item.41}%
\Destination{78.23}{809.815}{nav.page.113}%
\Destination{79.23}{682.865}{nav.section.2.4}%
\Destination{79.36}{471.355}{nav.print.124}%
\Destination{79.36}{630.255}{nav.number.2.4.1}%
\Destination{79.23}{682.865}{outline.24}%
}

% Page 126 (body page 114)
\Page{%
\TextAt{2}{79.226}{796.943}{II.}%
\TextAt{2}{94.073}{796.943}{The}%
\TextAt{2}{116.492}{796.943}{Theory}%
\TextAt{2}{154.401}{796.943}{of}%
\TextAt{2}{166.826}{796.943}{Valuations}%
\Vector{q 0 G 0 g 0.299 w 7.226 18.729 m 460.769 18.729 l S Q}%
\TextAt{12}{79.226}{761.124}{Theorem}%
\TextAt{12}{136.875}{761.124}{(4.2)}%
\TextAt{12}{169.669}{761.124}{(}%
\TextAt{26}{174.906}{761.124}{Ostrowski}%
\TextAt{12}{238.495}{761.124}{).}%
\TextAt{8}{251.923}{761.124}{Let}%
\TextAt{13}{273.444}{761.124}{K}%
\TextAt{8}{288.711}{761.124}{be}%
\TextAt{8}{305.187}{761.124}{a}%
\TextAt{8}{315.495}{761.124}{field}%
\TextAt{8}{341.403}{761.124}{which}%
\TextAt{8}{375.442}{761.124}{is}%
\TextAt{8}{387.757}{761.124}{complete}%
\TextAt{8}{437.732}{761.124}{with}%
\TextAt{8}{464.944}{761.124}{respect}%
\TextAt{8}{505.546}{761.124}{to}%
\TextAt{8}{520.408}{761.124}{an}%
\TextAt{8}{79.226}{744.078}{archimedean}%
\TextAt{8}{146.623}{744.078}{valuation}%
\TextAt{16}{197.127}{744.078}{\Character{106}}%
\TextAt{16}{203.762}{744.078}{\Character{106}}%
\TextAt{8}{207.098}{744.078}{.}%
\TextAt{8}{215.357}{744.078}{Then}%
\TextAt{8}{245.346}{744.078}{there}%
\TextAt{8}{274.689}{744.078}{is}%
\TextAt{8}{285.885}{744.078}{an}%
\TextAt{8}{301.576}{744.078}{isomorphism}%
\TextAt{13}{369.427}{744.078}{\Character{27}}%
\TextAt{8}{379.838}{744.078}{from}%
\TextAt{13}{406.91}{744.078}{K}%
\TextAt{8}{421.057}{744.078}{onto}%
\TextAt{18}{446.838}{744.078}{R}%
\TextAt{8}{458.801}{744.078}{or}%
\TextAt{18}{472.543}{744.078}{C}%
\TextAt{8}{484.506}{744.078}{satisfying}%
\TextAt{16}{233.475}{715.428}{\Character{106}}%
\TextAt{13}{236.796}{715.428}{a}%
\TextAt{16}{242.941}{715.428}{\Character{106}}%
\TextAt{9}{249.582}{715.428}{\Character{61}}%
\TextAt{16}{262.008}{715.428}{\Character{106}}%
\TextAt{13}{265.329}{715.428}{\Character{27}}%
\TextAt{13}{272.394}{715.428}{a}%
\TextAt{16}{278.556}{715.428}{\Character{106}}%
\TextAt{21}{281.877}{720.365}{s}%
\TextAt{8}{310.201}{715.428}{for}%
\TextAt{8}{328.086}{715.428}{all}%
\TextAt{13}{344.345}{715.428}{a}%
\NativeMathAt{353.811}{715.428}{11.955}{\in}%
\TextAt{13}{365.102}{715.428}{K}%
\TextAt{13}{375.264}{715.428}{\Character{59}}%
\TextAt{8}{79.226}{686.779}{for}%
\TextAt{8}{97.111}{686.779}{some}%
\TextAt{8}{126.437}{686.779}{fixed}%
\TextAt{13}{154.723}{686.779}{s}%
\NativeMathAt{163.558}{686.779}{11.955}{\in}%
\TextAt{9}{174.849}{686.779}{\Character{40}0}%
\TextAt{13}{185.254}{686.779}{\Character{59}}%
\TextAt{9}{190.498}{686.779}{1\Character{93}}%
\TextAt{8}{199.603}{686.779}{.}%
\TextAt{19}{79.226}{652.002}{Proof.}%
\TextAt{8}{116}{652.002}{We}%
\TextAt{8}{136.324}{652.002}{may}%
\TextAt{8}{161.872}{652.002}{assume}%
\TextAt{8}{202.496}{652.002}{without}%
\TextAt{8}{246.252}{652.002}{loss}%
\TextAt{8}{268.667}{652.002}{of}%
\TextAt{8}{282.177}{652.002}{generality}%
\TextAt{8}{336.346}{652.002}{that}%
\TextAt{18}{361.894}{652.002}{R}%
\NativeMathAt{374.154}{652.002}{11.955}{\subseteq}%
\TextAt{13}{387.079}{652.002}{K}%
\TextAt{8}{401.974}{652.002}{and}%
\TextAt{8}{424.915}{652.002}{that}%
\TextAt{8}{450.464}{652.002}{the}%
\TextAt{8}{470.799}{652.002}{valuation}%
\TextAt{16}{522.051}{652.002}{\Character{106}}%
\TextAt{16}{529.427}{652.002}{\Character{106}}%
\TextAt{8}{79.226}{634.956}{of}%
\TextAt{13}{92.576}{634.956}{K}%
\TextAt{8}{107.311}{634.956}{is}%
\TextAt{8}{119.094}{634.956}{an}%
\TextAt{8}{135.373}{634.956}{extension}%
\TextAt{8}{187.158}{634.956}{of}%
\TextAt{8}{200.508}{634.956}{the}%
\TextAt{8}{220.684}{634.956}{usual}%
\TextAt{8}{251.333}{634.956}{absolute}%
\TextAt{8}{297.595}{634.956}{value}%
\TextAt{8}{327.861}{634.956}{of}%
\TextAt{18}{341.211}{634.956}{R}%
\TextAt{8}{349.845}{634.956}{.}%
\TextAt{8}{358.353}{634.956}{In}%
\TextAt{8}{372.993}{634.956}{fact,}%
\TextAt{8}{399.355}{634.956}{replacing}%
\TextAt{16}{449.455}{634.956}{\Character{106}}%
\TextAt{16}{456.676}{634.956}{\Character{106}}%
\TextAt{8}{463.931}{634.956}{by}%
\TextAt{16}{480.21}{634.956}{\Character{106}}%
\TextAt{16}{487.431}{634.956}{\Character{106}}%
\TextAt{21}{490.769}{639.294}{s}%
\TextAt{23}{494.684}{642.107}{\Character{0}}%
\TextAt{24}{500.441}{642.107}{1}%
\TextAt{8}{509.007}{634.956}{for}%
\TextAt{8}{526.911}{634.956}{a}%
\TextAt{8}{79.226}{617.91}{suitable}%
\TextAt{13}{122.209}{617.91}{s}%
\TextAt{13}{131.032}{617.91}{\Character{62}}%
\TextAt{9}{143.47}{617.91}{0}%
\TextAt{8}{149.323}{617.91}{,}%
\TextAt{8}{155.949}{617.91}{we}%
\TextAt{8}{172.523}{617.91}{may}%
\TextAt{8}{197.239}{617.91}{assume}%
\TextAt{8}{237.03}{617.91}{by}%
\TextAt{8}{252.636}{617.91}{(3.7)}%
\TextAt{8}{279.958}{617.91}{that}%
\TextAt{8}{304.673}{617.91}{the}%
\TextAt{8}{324.177}{617.91}{restriction}%
\TextAt{8}{379.521}{617.91}{of}%
\TextAt{16}{392.198}{617.91}{\Character{106}}%
\TextAt{16}{398.749}{617.91}{\Character{106}}%
\TextAt{8}{405.328}{617.91}{to}%
\TextAt{18}{418.985}{617.91}{Q}%
\TextAt{8}{431.528}{617.91}{is}%
\TextAt{8}{442.639}{617.91}{equal}%
\TextAt{8}{472.877}{617.91}{to}%
\TextAt{8}{486.535}{617.91}{the}%
\TextAt{8}{506.038}{617.91}{usual}%
\TextAt{8}{79.226}{600.864}{absolute}%
\TextAt{8}{125.571}{600.864}{value.}%
\TextAt{8}{160.675}{600.864}{Then}%
\TextAt{8}{191.335}{600.864}{taking}%
\TextAt{8}{227.542}{600.864}{the}%
\TextAt{8}{247.8}{600.864}{closure}%
\TextAt{20}{288.312}{603.952}{\Character{98}}%
\TextAt{18}{286.984}{600.864}{Q}%
\TextAt{8}{300.282}{600.864}{in}%
\TextAt{13}{314.037}{600.864}{K}%
\TextAt{8}{328.854}{600.864}{we}%
\TextAt{8}{346.184}{600.864}{find}%
\TextAt{8}{369.695}{600.864}{that}%
\TextAt{20}{396.494}{603.952}{\Character{98}}%
\TextAt{18}{395.166}{600.864}{Q}%
\TextAt{8}{408.464}{600.864}{is}%
\TextAt{8}{420.33}{600.864}{complete}%
\TextAt{8}{469.855}{600.864}{with}%
\TextAt{8}{496.617}{600.864}{respect}%
\TextAt{8}{79.226}{583.818}{to}%
\TextAt{8}{94.579}{583.818}{the}%
\TextAt{8}{115.778}{583.818}{restriction}%
\TextAt{8}{172.819}{583.818}{of}%
\TextAt{16}{187.192}{583.818}{\Character{106}}%
\TextAt{16}{195.441}{583.818}{\Character{106}}%
\TextAt{8}{203.714}{583.818}{to}%
\TextAt{20}{220.395}{586.906}{\Character{98}}%
\TextAt{18}{219.067}{583.818}{Q}%
\TextAt{8}{228.365}{583.818}{,}%
\TextAt{8}{236.818}{583.818}{in}%
\TextAt{8}{251.514}{583.818}{other}%
\TextAt{8}{283.126}{583.818}{words,}%
\TextAt{8}{321.239}{583.818}{it}%
\TextAt{8}{333.986}{583.818}{is}%
\TextAt{8}{346.792}{583.818}{a}%
\TextAt{8}{357.591}{583.818}{completion}%
\TextAt{8}{418.469}{583.818}{of}%
\TextAt{9}{432.842}{583.818}{\Character{40}}%
\TextAt{18}{437.394}{583.818}{Q}%
\TextAt{13}{446.693}{583.818}{\Character{59}}%
\TextAt{16}{451.937}{583.818}{\Character{106}}%
\TextAt{16}{460.186}{583.818}{\Character{106}}%
\TextAt{9}{463.519}{583.818}{\Character{41}}%
\TextAt{8}{468.071}{583.818}{.}%
\TextAt{8}{479.647}{583.818}{In}%
\TextAt{8}{495.311}{583.818}{view}%
\TextAt{8}{523.337}{583.818}{of}%
\TextAt{8}{79.226}{566.772}{the}%
\TextAt{8}{99.483}{566.772}{uniqueness}%
\TextAt{8}{158.558}{566.772}{of}%
\TextAt{8}{171.988}{566.772}{completions,}%
\TextAt{8}{239.815}{566.772}{there}%
\TextAt{8}{269.827}{566.772}{is}%
\TextAt{8}{281.691}{566.772}{an}%
\TextAt{8}{298.05}{566.772}{isomorphism}%
\TextAt{13}{366.569}{566.772}{\Character{27}}%
\TextAt{9}{377.143}{566.772}{\Character{58}}%
\TextAt{18}{383.886}{566.772}{R}%
\NativeArrowAt{396.012}{566.772}{11.955}{19.26}{\longrightarrow}%
\TextAt{20}{420.093}{569.86}{\Character{98}}%
\TextAt{18}{418.764}{566.772}{Q}%
\TextAt{8}{432.06}{566.772}{such}%
\TextAt{8}{458.557}{566.772}{that}%
\TextAt{16}{484.026}{566.772}{\Character{106}}%
\TextAt{13}{487.347}{566.772}{a}%
\TextAt{16}{493.492}{566.772}{\Character{106}}%
\TextAt{9}{500.304}{566.772}{\Character{61}}%
\TextAt{16}{512.9}{566.772}{\Character{106}}%
\TextAt{13}{516.221}{566.772}{\Character{27}}%
\TextAt{13}{523.287}{566.772}{a}%
\TextAt{16}{529.449}{566.772}{\Character{106}}%
\TextAt{8}{79.226}{549.726}{as}%
\TextAt{8}{94.051}{549.726}{required.}%
\TextAt{8}{145.816}{549.726}{In}%
\TextAt{8}{160.892}{549.726}{order}%
\TextAt{8}{191.916}{549.726}{to}%
\TextAt{8}{206.682}{549.726}{prove}%
\TextAt{8}{238.687}{549.726}{that}%
\TextAt{13}{264.51}{549.726}{K}%
\TextAt{9}{279.424}{549.726}{\Character{61}}%
\TextAt{18}{292.626}{549.726}{R}%
\TextAt{8}{305.612}{549.726}{or}%
\TextAt{9}{320.378}{549.726}{\Character{61}}%
\TextAt{18}{333.579}{549.726}{C}%
\TextAt{8}{346.565}{549.726}{we}%
\TextAt{8}{364.248}{549.726}{show}%
\TextAt{8}{393.706}{549.726}{that}%
\TextAt{8}{419.53}{549.726}{each}%
\TextAt{13}{446.323}{549.726}{\Character{24}}%
\NativeMathAt{456.109}{549.726}{11.955}{\in}%
\TextAt{13}{468.176}{549.726}{K}%
\TextAt{8}{483.346}{549.726}{satisfies}%
\TextAt{8}{526.911}{549.726}{a}%
\TextAt{8}{79.226}{532.68}{quadratic}%
\TextAt{8}{130.956}{532.68}{equation}%
\TextAt{8}{178.132}{532.68}{over}%
\TextAt{18}{202.544}{532.68}{R}%
\TextAt{8}{211.178}{532.68}{.}%
\TextAt{8}{219.416}{532.68}{For}%
\TextAt{8}{239.752}{532.68}{this,}%
\TextAt{8}{265.319}{532.68}{consider}%
\TextAt{8}{310.271}{532.68}{the}%
\TextAt{8}{329.794}{532.68}{continuous}%
\TextAt{8}{387.764}{532.68}{function}%
\TextAt{13}{432.979}{532.68}{f}%
\TextAt{9}{443.346}{532.68}{\Character{58}}%
\TextAt{18}{449.919}{532.68}{C}%
\NativeArrowAt{461.874}{532.68}{11.955}{19.26}{\longrightarrow}%
\TextAt{18}{484.456}{532.68}{R}%
\TextAt{8}{496.354}{532.68}{defined}%
\TextAt{8}{79.226}{515.634}{by}%
\TextAt{13}{235.179}{498.588}{f}%
\TextAt{9}{242.225}{498.588}{\Character{40}}%
\TextAt{13}{246.778}{498.588}{z}%
\TextAt{9}{252.748}{498.588}{\Character{41}}%
\TextAt{9}{260.603}{498.588}{\Character{61}}%
\TextAt{20}{273.047}{508.75}{\Character{12}}%
\TextAt{20}{273.047}{501.577}{\Character{12}}%
\TextAt{13}{277.032}{498.588}{\Character{24}}%
\TextAt{17}{282.722}{503.524}{2}%
\TextAt{16}{290.111}{498.588}{\Character{0}}%
\TextAt{9}{302.067}{498.588}{\Character{40}}%
\TextAt{13}{306.619}{498.588}{z}%
\TextAt{9}{315.246}{498.588}{\Character{43}}%
\TextAt{9}{327.703}{498.588}{\Character{22}}%
\TextAt{13}{327.007}{498.588}{z}%
\TextAt{9}{332.978}{498.588}{\Character{41}}%
\TextAt{13}{337.53}{498.588}{\Character{24}}%
\TextAt{9}{345.878}{498.588}{\Character{43}}%
\TextAt{13}{357.639}{498.588}{z}%
\TextAt{9}{364.319}{498.588}{\Character{22}}%
\TextAt{13}{363.609}{498.588}{z}%
\TextAt{20}{369.58}{508.75}{\Character{12}}%
\TextAt{20}{369.58}{501.577}{\Character{12}}%
\TextAt{13}{373.565}{498.588}{\Character{58}}%
\TextAt{8}{79.226}{475.217}{Note}%
\TextAt{8}{107.305}{475.217}{here}%
\TextAt{8}{132.456}{475.217}{that}%
\TextAt{13}{157.618}{475.217}{z}%
\TextAt{9}{165.823}{475.217}{\Character{43}}%
\TextAt{9}{177.862}{475.217}{\Character{22}}%
\TextAt{13}{177.162}{475.217}{z}%
\TextAt{13}{183.116}{475.217}{\Character{59}}%
\TextAt{13}{188.352}{475.217}{z}%
\TextAt{9}{195.056}{475.217}{\Character{22}}%
\TextAt{13}{194.347}{475.217}{z}%
\NativeMathAt{203.638}{475.217}{11.955}{\in}%
\TextAt{18}{214.929}{475.217}{R}%
\NativeMathAt{226.884}{475.217}{11.955}{\subseteq}%
\TextAt{13}{239.504}{475.217}{K}%
\TextAt{8}{250.321}{475.217}{.}%
\TextAt{8}{258.701}{475.217}{Since}%
\TextAt{8}{289.052}{475.217}{lim}%
\TextAt{21}{305.311}{473.424}{z}%
\NativeMathAt{309.586}{473.424}{7.97}{\rightarrow}%
\NativeMathAt{318.05}{473.424}{7.97}{\infty}%
\TextAt{13}{329.014}{475.217}{f}%
\TextAt{9}{336.06}{475.217}{\Character{40}}%
\TextAt{13}{340.613}{475.217}{z}%
\TextAt{9}{346.583}{475.217}{\Character{41}}%
\TextAt{9}{354.438}{475.217}{\Character{61}}%
\NativeMathAt{366.882}{475.217}{11.955}{\infty}%
\TextAt{8}{378.837}{475.217}{,}%
\TextAt{13}{385.821}{475.217}{f}%
\TextAt{9}{392.867}{475.217}{\Character{40}}%
\TextAt{13}{397.42}{475.217}{z}%
\TextAt{9}{403.39}{475.217}{\Character{41}}%
\TextAt{8}{411.633}{475.217}{has}%
\TextAt{8}{432.3}{475.217}{a}%
\TextAt{8}{441.849}{475.217}{minimum}%
\TextAt{13}{493.994}{475.217}{m}%
\TextAt{8}{504.234}{475.217}{.}%
\TextAt{8}{512.613}{475.217}{The}%
\TextAt{8}{79.226}{458.171}{set}%
\TextAt{13}{244.556}{441.125}{S}%
\TextAt{9}{255.772}{441.125}{\Character{61}}%
\TextAt{16}{268.198}{441.125}{\Character{102}}%
\TextAt{13}{274.175}{441.125}{z}%
\NativeMathAt{283.467}{441.125}{11.955}{\in}%
\TextAt{18}{294.758}{441.125}{C}%
\TextAt{16}{306.713}{441.125}{\Character{106}}%
\TextAt{13}{313.355}{441.125}{f}%
\TextAt{9}{320.401}{441.125}{\Character{40}}%
\TextAt{13}{324.953}{441.125}{z}%
\TextAt{9}{330.924}{441.125}{\Character{41}}%
\TextAt{9}{338.779}{441.125}{\Character{61}}%
\TextAt{13}{351.223}{441.125}{m}%
\TextAt{16}{361.462}{441.125}{\Character{103}}%
\TextAt{8}{79.226}{417.754}{is}%
\TextAt{8}{92.698}{417.754}{therefore}%
\TextAt{8}{143.506}{417.754}{nonempty,}%
\TextAt{8}{202.548}{417.754}{bounded,}%
\TextAt{8}{255.731}{417.754}{and}%
\TextAt{8}{280.202}{417.754}{closed,}%
\TextAt{8}{320.115}{417.754}{and}%
\TextAt{8}{344.586}{417.754}{there}%
\TextAt{8}{376.206}{417.754}{is}%
\TextAt{8}{389.678}{417.754}{a}%
\TextAt{13}{401.141}{417.754}{z}%
\TextAt{17}{406.579}{415.961}{0}%
\NativeMathAt{417.543}{417.754}{11.955}{\in}%
\TextAt{13}{431.745}{417.754}{S}%
\TextAt{8}{445.245}{417.754}{such}%
\TextAt{8}{473.351}{417.754}{that}%
\TextAt{16}{500.427}{417.754}{\Character{106}}%
\TextAt{13}{503.748}{417.754}{z}%
\TextAt{17}{509.186}{415.961}{0}%
\TextAt{16}{513.919}{417.754}{\Character{106}}%
\NativeMathAt{523.459}{417.754}{11.955}{\geq}%
\TextAt{16}{79.226}{400.708}{\Character{106}}%
\TextAt{13}{82.547}{400.708}{z}%
\TextAt{16}{88.518}{400.708}{\Character{106}}%
\TextAt{8}{96.96}{400.708}{for}%
\TextAt{8}{116.07}{400.708}{all}%
\TextAt{13}{133.554}{400.708}{z}%
\NativeMathAt{144.932}{400.708}{11.955}{\in}%
\TextAt{13}{158.31}{400.708}{S}%
\TextAt{8}{166.205}{400.708}{.}%
\TextAt{8}{178.327}{400.708}{It}%
\TextAt{8}{192.224}{400.708}{suffices}%
\TextAt{8}{233.236}{400.708}{to}%
\TextAt{8}{248.77}{400.708}{show}%
\TextAt{8}{278.998}{400.708}{that}%
\TextAt{13}{305.592}{400.708}{m}%
\TextAt{9}{321.239}{400.708}{\Character{61}}%
\TextAt{9}{335.741}{400.708}{0}%
\TextAt{8}{341.604}{400.708}{,}%
\TextAt{8}{350.284}{400.708}{because}%
\TextAt{8}{394.81}{400.708}{then}%
\TextAt{8}{422.695}{400.708}{one}%
\TextAt{8}{445.379}{400.708}{has}%
\TextAt{8}{467.477}{400.708}{the}%
\TextAt{8}{488.858}{400.708}{equation}%
\TextAt{13}{79.226}{383.662}{\Character{24}}%
\TextAt{17}{84.917}{388.001}{2}%
\TextAt{16}{92.306}{383.662}{\Character{0}}%
\TextAt{9}{104.261}{383.662}{\Character{40}}%
\TextAt{13}{108.813}{383.662}{z}%
\TextAt{17}{114.251}{381.869}{0}%
\TextAt{9}{121.64}{383.662}{\Character{43}}%
\TextAt{9}{134.097}{383.662}{\Character{22}}%
\TextAt{13}{133.401}{383.662}{z}%
\TextAt{17}{138.839}{381.869}{0}%
\TextAt{9}{143.571}{383.662}{\Character{41}}%
\TextAt{13}{148.124}{383.662}{\Character{24}}%
\TextAt{9}{156.471}{383.662}{\Character{43}}%
\TextAt{13}{168.232}{383.662}{z}%
\TextAt{17}{173.67}{381.869}{0}%
\TextAt{9}{179.112}{383.662}{\Character{22}}%
\TextAt{13}{178.402}{383.662}{z}%
\TextAt{17}{183.84}{381.869}{0}%
\TextAt{9}{191.893}{383.662}{\Character{61}}%
\TextAt{9}{204.303}{383.662}{0}%
\TextAt{8}{210.172}{383.662}{.}%
\TextAt{8}{218.62}{383.662}{Assume}%
\TextAt{13}{261.982}{383.662}{m}%
\TextAt{13}{275.527}{383.662}{\Character{62}}%
\TextAt{9}{287.967}{383.662}{0}%
\TextAt{8}{293.82}{383.662}{.}%
\TextAt{8}{302.269}{383.662}{Consider}%
\TextAt{8}{351.105}{383.662}{the}%
\TextAt{8}{371.262}{383.662}{real}%
\TextAt{8}{394.024}{383.662}{polynomial}%
\TextAt{13}{220.02}{355.013}{g}%
\TextAt{9}{226.055}{355.013}{\Character{40}}%
\TextAt{13}{230.607}{355.013}{x}%
\TextAt{9}{237.259}{355.013}{\Character{41}}%
\TextAt{9}{245.114}{355.013}{\Character{61}}%
\TextAt{13}{257.558}{355.013}{x}%
\TextAt{17}{264.21}{359.949}{2}%
\TextAt{16}{271.599}{355.013}{\Character{0}}%
\TextAt{9}{283.554}{355.013}{\Character{40}}%
\TextAt{13}{288.106}{355.013}{z}%
\TextAt{17}{293.544}{353.22}{0}%
\TextAt{9}{300.933}{355.013}{\Character{43}}%
\TextAt{9}{313.39}{355.013}{\Character{22}}%
\TextAt{13}{312.695}{355.013}{z}%
\TextAt{17}{318.132}{353.22}{0}%
\TextAt{9}{322.865}{355.013}{\Character{41}}%
\TextAt{13}{327.417}{355.013}{x}%
\TextAt{9}{336.726}{355.013}{\Character{43}}%
\TextAt{13}{348.487}{355.013}{z}%
\TextAt{17}{353.925}{353.22}{0}%
\TextAt{9}{359.366}{355.013}{\Character{22}}%
\TextAt{13}{358.657}{355.013}{z}%
\TextAt{17}{364.095}{353.22}{0}%
\TextAt{9}{371.484}{355.013}{\Character{43}}%
\TextAt{13}{383.246}{355.013}{\Character{34}\Character{59}}%
\TextAt{8}{79.226}{326.364}{where}%
\TextAt{9}{113.465}{326.364}{0}%
\TextAt{13}{123.371}{326.364}{\Character{60}}%
\TextAt{13}{136.522}{326.364}{\Character{34}}%
\TextAt{13}{146.05}{326.364}{\Character{60}}%
\TextAt{13}{159.201}{326.364}{m}%
\TextAt{8}{169.457}{326.364}{,}%
\TextAt{8}{177.143}{326.364}{with}%
\TextAt{8}{204.232}{326.364}{the}%
\TextAt{8}{224.819}{326.364}{roots}%
\TextAt{13}{254.909}{326.364}{z}%
\TextAt{17}{260.347}{324.57}{1}%
\TextAt{13}{265.079}{326.364}{\Character{59}}%
\TextAt{9}{271.032}{326.364}{\Character{22}}%
\TextAt{13}{270.323}{326.364}{z}%
\TextAt{17}{275.761}{324.57}{1}%
\NativeMathAt{284.547}{326.364}{11.955}{\in}%
\TextAt{18}{296.57}{326.364}{C}%
\TextAt{8}{305.204}{326.364}{.}%
\TextAt{8}{314.941}{326.364}{We}%
\TextAt{8}{335.515}{326.364}{have}%
\TextAt{13}{362.94}{326.364}{z}%
\TextAt{17}{368.378}{324.57}{1}%
\TextAt{9}{373.819}{326.364}{\Character{22}}%
\TextAt{13}{373.11}{326.364}{z}%
\TextAt{17}{378.548}{324.57}{1}%
\TextAt{9}{387.333}{326.364}{\Character{61}}%
\TextAt{13}{400.491}{326.364}{z}%
\TextAt{17}{405.929}{324.57}{0}%
\TextAt{9}{411.37}{326.364}{\Character{22}}%
\TextAt{13}{410.661}{326.364}{z}%
\TextAt{17}{416.099}{324.57}{0}%
\TextAt{9}{423.781}{326.364}{\Character{43}}%
\TextAt{13}{435.835}{326.364}{\Character{34}}%
\TextAt{8}{441.313}{326.364}{,}%
\TextAt{8}{448.999}{326.364}{hence}%
\TextAt{16}{481.935}{326.364}{\Character{106}}%
\TextAt{13}{485.256}{326.364}{z}%
\TextAt{17}{490.694}{324.57}{1}%
\TextAt{16}{495.426}{326.364}{\Character{106}}%
\TextAt{13}{502.8}{326.364}{\Character{62}}%
\TextAt{16}{515.958}{326.364}{\Character{106}}%
\TextAt{13}{519.278}{326.364}{z}%
\TextAt{17}{524.716}{324.57}{0}%
\TextAt{16}{529.449}{326.364}{\Character{106}}%
\TextAt{8}{79.226}{309.318}{and}%
\TextAt{8}{101.989}{309.318}{thus}%
\TextAt{13}{278.219}{292.272}{f}%
\TextAt{9}{285.265}{292.272}{\Character{40}}%
\TextAt{13}{289.817}{292.272}{z}%
\TextAt{17}{295.255}{290.478}{1}%
\TextAt{9}{299.988}{292.272}{\Character{41}}%
\TextAt{13}{307.861}{292.272}{\Character{62}}%
\TextAt{13}{320.271}{292.272}{m\Character{58}}%
\TextAt{8}{79.226}{268.901}{For}%
\TextAt{8}{100.196}{268.901}{fixed}%
\TextAt{13}{128.481}{268.901}{n}%
\NativeMathAt{138.79}{268.901}{11.955}{\in}%
\TextAt{18}{150.081}{268.901}{N}%
\TextAt{8}{158.715}{268.901}{,}%
\TextAt{8}{165.864}{268.901}{consider}%
\TextAt{8}{211.449}{268.901}{on}%
\TextAt{8}{227.708}{268.901}{the}%
\TextAt{8}{247.865}{268.901}{other}%
\TextAt{8}{278.434}{268.901}{hand}%
\TextAt{8}{307.701}{268.901}{the}%
\TextAt{8}{327.857}{268.901}{real}%
\TextAt{8}{350.62}{268.901}{polynomial}%
\TextAt{13}{164.337}{232.587}{G}%
\TextAt{9}{173.57}{232.587}{\Character{40}}%
\TextAt{13}{178.123}{232.587}{x}%
\TextAt{9}{184.775}{232.587}{\Character{41}}%
\TextAt{9}{192.63}{232.587}{\Character{61}}%
\TextAt{9}{205.051}{232.587}{\Character{91}}%
\TextAt{13}{208.325}{232.587}{g}%
\TextAt{9}{214.359}{232.587}{\Character{40}}%
\TextAt{13}{218.912}{232.587}{x}%
\TextAt{9}{225.564}{232.587}{\Character{41}}%
\TextAt{16}{232.773}{232.587}{\Character{0}}%
\TextAt{13}{244.728}{232.587}{\Character{34}}%
\TextAt{9}{250.206}{232.587}{\Character{93}}%
\TextAt{21}{253.458}{237.523}{n}%
\TextAt{16}{261.751}{232.587}{\Character{0}}%
\TextAt{9}{273.706}{232.587}{\Character{40}}%
\TextAt{16}{278.258}{232.587}{\Character{0}}%
\TextAt{13}{287.557}{232.587}{\Character{34}}%
\TextAt{9}{293.035}{232.587}{\Character{41}}%
\TextAt{21}{297.587}{237.523}{n}%
\TextAt{9}{306.544}{232.587}{\Character{61}}%
\TextAt{17}{321.921}{247.531}{2}%
\TextAt{21}{326.156}{247.531}{n}%
\TextAt{20}{318.97}{243.944}{\Character{89}}%
\TextAt{21}{319.756}{218.75}{i}%
\TextAt{17}{322.639}{218.75}{\Character{61}1}%
\TextAt{9}{334.246}{232.587}{\Character{40}}%
\TextAt{13}{338.798}{232.587}{x}%
\TextAt{16}{348.107}{232.587}{\Character{0}}%
\TextAt{13}{360.062}{232.587}{\Character{11}}%
\TextAt{21}{367.508}{230.793}{i}%
\TextAt{9}{370.889}{232.587}{\Character{41}}%
\TextAt{9}{378.744}{232.587}{\Character{61}}%
\TextAt{17}{394.139}{247.531}{2}%
\TextAt{21}{398.373}{247.531}{n}%
\TextAt{20}{391.187}{243.944}{\Character{89}}%
\TextAt{21}{391.974}{218.75}{i}%
\TextAt{17}{394.857}{218.75}{\Character{61}1}%
\TextAt{9}{406.464}{232.587}{\Character{40}}%
\TextAt{13}{411.016}{232.587}{x}%
\TextAt{16}{420.325}{232.587}{\Character{0}}%
\TextAt{9}{433.439}{232.587}{\Character{22}}%
\TextAt{13}{432.28}{232.587}{\Character{11}}%
\TextAt{21}{439.725}{230.793}{i}%
\TextAt{9}{443.107}{232.587}{\Character{41}}%
\TextAt{8}{79.226}{195.988}{with}%
\TextAt{8}{106.297}{195.988}{roots}%
\TextAt{13}{136.368}{195.988}{\Character{11}}%
\TextAt{17}{143.814}{194.195}{1}%
\TextAt{13}{148.546}{195.988}{\Character{59}}%
\TextAt{13}{153.77}{195.988}{\Character{58}}%
\TextAt{13}{159.007}{195.988}{\Character{58}}%
\TextAt{13}{164.243}{195.988}{\Character{58}}%
\TextAt{13}{169.468}{195.988}{\Character{59}}%
\TextAt{13}{174.704}{195.988}{\Character{11}}%
\TextAt{17}{182.213}{194.195}{2}%
\TextAt{21}{186.447}{194.195}{n}%
\NativeMathAt{196.104}{195.988}{11.955}{\in}%
\TextAt{18}{208.094}{195.988}{C}%
\TextAt{8}{216.729}{195.988}{.}%
\TextAt{8}{226.409}{195.988}{It}%
\TextAt{8}{239.492}{195.988}{follows}%
\TextAt{8}{278.339}{195.988}{that}%
\TextAt{13}{304.118}{195.988}{G}%
\TextAt{9}{313.352}{195.988}{\Character{40}}%
\TextAt{13}{317.904}{195.988}{z}%
\TextAt{17}{323.342}{194.195}{1}%
\TextAt{9}{328.074}{195.988}{\Character{41}}%
\TextAt{9}{336.634}{195.988}{\Character{61}}%
\TextAt{9}{349.749}{195.988}{0}%
\TextAt{8}{355.625}{195.988}{;}%
\TextAt{8}{363.39}{195.988}{say,}%
\TextAt{13}{386.403}{195.988}{z}%
\TextAt{17}{391.841}{194.195}{1}%
\TextAt{9}{400.594}{195.988}{\Character{61}}%
\TextAt{13}{413.719}{195.988}{\Character{11}}%
\TextAt{17}{421.165}{194.195}{1}%
\TextAt{8}{425.897}{195.988}{.}%
\TextAt{8}{435.578}{195.988}{We}%
\TextAt{8}{456.133}{195.988}{may}%
\TextAt{8}{481.912}{195.988}{substitute}%
\TextAt{13}{79.226}{178.942}{\Character{24}}%
\NativeMathAt{88.237}{178.942}{11.955}{\in}%
\TextAt{13}{99.528}{178.942}{K}%
\TextAt{8}{114.243}{178.942}{into}%
\TextAt{8}{137.986}{178.942}{the}%
\TextAt{8}{158.143}{178.942}{polynomial}%
\TextAt{13}{213.198}{142.628}{G}%
\TextAt{9}{222.431}{142.628}{\Character{40}}%
\TextAt{13}{226.984}{142.628}{x}%
\TextAt{9}{233.636}{142.628}{\Character{41}}%
\TextAt{17}{238.188}{147.564}{2}%
\TextAt{9}{246.241}{142.628}{\Character{61}}%
\TextAt{17}{261.618}{157.572}{2}%
\TextAt{21}{265.853}{157.572}{n}%
\TextAt{20}{258.667}{153.985}{\Character{89}}%
\TextAt{21}{259.453}{128.791}{i}%
\TextAt{17}{262.336}{128.791}{\Character{61}1}%
\TextAt{20}{273.943}{152.312}{\Character{0}}%
\TextAt{13}{279.422}{142.628}{x}%
\TextAt{17}{286.074}{147.564}{2}%
\TextAt{16}{293.463}{142.628}{\Character{0}}%
\TextAt{9}{305.418}{142.628}{\Character{40}}%
\TextAt{13}{309.971}{142.628}{\Character{11}}%
\TextAt{21}{317.416}{140.835}{i}%
\TextAt{9}{323.454}{142.628}{\Character{43}}%
\TextAt{9}{336.366}{142.628}{\Character{22}}%
\TextAt{13}{335.216}{142.628}{\Character{11}}%
\TextAt{21}{342.661}{140.835}{i}%
\TextAt{9}{346.042}{142.628}{\Character{41}}%
\TextAt{13}{350.595}{142.628}{x}%
\TextAt{9}{359.904}{142.628}{\Character{43}}%
\TextAt{13}{371.665}{142.628}{\Character{11}}%
\TextAt{21}{379.11}{140.835}{i}%
\TextAt{9}{383.651}{142.628}{\Character{22}}%
\TextAt{13}{382.492}{142.628}{\Character{11}}%
\TextAt{21}{389.937}{140.835}{i}%
\TextAt{20}{393.319}{152.312}{\Character{1}}%
\TextAt{8}{79.226}{106.699}{and}%
\TextAt{8}{101.989}{106.699}{get}%
\TextAt{16}{219.074}{81.976}{\Character{106}}%
\TextAt{13}{222.395}{81.976}{G}%
\TextAt{9}{231.629}{81.976}{\Character{40}}%
\TextAt{13}{236.181}{81.976}{\Character{24}}%
\TextAt{9}{241.872}{81.976}{\Character{41}}%
\TextAt{16}{246.424}{81.976}{\Character{106}}%
\TextAt{17}{249.745}{86.912}{2}%
\TextAt{9}{257.798}{81.976}{\Character{61}}%
\TextAt{17}{273.175}{96.92}{2}%
\TextAt{21}{277.409}{96.92}{n}%
\TextAt{20}{270.223}{93.334}{\Character{89}}%
\TextAt{21}{271.01}{68.139}{i}%
\TextAt{17}{273.893}{68.139}{\Character{61}1}%
\TextAt{13}{287.492}{81.976}{f}%
\TextAt{9}{294.538}{81.976}{\Character{40}}%
\TextAt{13}{299.091}{81.976}{\Character{11}}%
\TextAt{21}{306.536}{80.183}{i}%
\TextAt{9}{309.918}{81.976}{\Character{41}}%
\NativeMathAt{317.791}{81.976}{11.955}{\geq}%
\TextAt{13}{330.41}{81.976}{f}%
\TextAt{9}{337.457}{81.976}{\Character{40}}%
\TextAt{13}{342.009}{81.976}{\Character{11}}%
\TextAt{17}{349.455}{80.183}{1}%
\TextAt{9}{354.187}{81.976}{\Character{41}}%
\TextAt{13}{358.739}{81.976}{m}%
\TextAt{17}{368.978}{86.912}{2}%
\TextAt{21}{373.213}{86.912}{n}%
\TextAt{3}{378.351}{86.912}{\Character{0}}%
\TextAt{17}{384.937}{86.912}{1}%
\TextAt{13}{389.67}{81.976}{\Character{58}}%
\TextAt{8}{297.211}{43.033}{114}%
\Destination{78.23}{809.815}{nav.page.114}%
\Destination{79.36}{490.255}{nav.print.125}%
\Destination{79.36}{775.155}{nav.number.2.4.2}%
\Hotlink{252.136}{615.755}{25.078}{11.047}{nav.number.2.3.7}%
}

% Page 127 (body page 115)
\Page{%
\TextAt{2}{79.226}{796.943}{II.}%
\TextAt{2}{94.073}{796.943}{The}%
\TextAt{2}{116.492}{796.943}{Theory}%
\TextAt{2}{154.401}{796.943}{of}%
\TextAt{2}{166.826}{796.943}{Valuations}%
\Vector{q 0 G 0 g 0.299 w 7.226 18.729 m 460.769 18.729 l S Q}%
\TextAt{8}{79.226}{761.124}{From}%
\TextAt{8}{109.951}{761.124}{this}%
\TextAt{8}{132.773}{761.124}{and}%
\TextAt{8}{155.536}{761.124}{the}%
\TextAt{8}{175.692}{761.124}{inequality}%
\TextAt{16}{138.308}{733.032}{\Character{106}}%
\TextAt{13}{141.629}{733.032}{G}%
\TextAt{9}{150.862}{733.032}{\Character{40}}%
\TextAt{13}{155.415}{733.032}{\Character{24}}%
\TextAt{9}{161.105}{733.032}{\Character{41}}%
\TextAt{16}{165.658}{733.032}{\Character{106}}%
\NativeMathAt{172.281}{733.032}{11.955}{\leq}%
\TextAt{20}{184.919}{743.194}{\Character{12}}%
\TextAt{20}{184.919}{736.021}{\Character{12}}%
\TextAt{13}{188.904}{733.032}{\Character{24}}%
\TextAt{17}{194.594}{737.969}{2}%
\TextAt{16}{201.983}{733.032}{\Character{0}}%
\TextAt{9}{213.938}{733.032}{\Character{40}}%
\TextAt{13}{218.491}{733.032}{z}%
\TextAt{17}{223.928}{731.239}{0}%
\TextAt{9}{231.317}{733.032}{\Character{43}}%
\TextAt{9}{243.774}{733.032}{\Character{22}}%
\TextAt{13}{243.079}{733.032}{z}%
\TextAt{17}{248.517}{731.239}{0}%
\TextAt{9}{253.249}{733.032}{\Character{41}}%
\TextAt{13}{257.801}{733.032}{\Character{24}}%
\TextAt{9}{266.148}{733.032}{\Character{43}}%
\TextAt{13}{277.91}{733.032}{z}%
\TextAt{17}{283.348}{731.239}{0}%
\TextAt{9}{288.789}{733.032}{\Character{22}}%
\TextAt{13}{288.08}{733.032}{z}%
\TextAt{17}{293.518}{731.239}{0}%
\TextAt{20}{298.25}{743.194}{\Character{12}}%
\TextAt{20}{298.25}{736.021}{\Character{12}}%
\TextAt{21}{302.235}{740.04}{n}%
\TextAt{9}{310.528}{733.032}{\Character{43}}%
\TextAt{16}{322.29}{733.032}{\Character{106}}%
\TextAt{13}{325.61}{733.032}{\Character{34}}%
\TextAt{16}{331.088}{733.032}{\Character{106}}%
\TextAt{21}{334.409}{737.969}{n}%
\TextAt{9}{343.366}{733.032}{\Character{61}}%
\TextAt{13}{355.792}{733.032}{f}%
\TextAt{9}{362.838}{733.032}{\Character{40}}%
\TextAt{13}{367.391}{733.032}{z}%
\TextAt{17}{372.829}{731.239}{0}%
\TextAt{9}{377.561}{733.032}{\Character{41}}%
\TextAt{21}{382.113}{737.969}{n}%
\TextAt{9}{390.406}{733.032}{\Character{43}}%
\TextAt{13}{402.168}{733.032}{\Character{34}}%
\TextAt{21}{407.646}{737.969}{n}%
\TextAt{9}{416.603}{733.032}{\Character{61}}%
\TextAt{13}{429.028}{733.032}{m}%
\TextAt{21}{439.268}{737.969}{n}%
\TextAt{9}{447.561}{733.032}{\Character{43}}%
\TextAt{13}{459.322}{733.032}{\Character{34}}%
\TextAt{21}{464.8}{737.969}{n}%
\TextAt{13}{470.436}{733.032}{\Character{59}}%
\TextAt{8}{79.226}{704.941}{it}%
\TextAt{8}{90.93}{704.941}{follows}%
\TextAt{8}{129.366}{704.941}{that}%
\TextAt{13}{154.735}{704.941}{f}%
\TextAt{9}{161.781}{704.941}{\Character{40}}%
\TextAt{13}{166.334}{704.941}{\Character{11}}%
\TextAt{17}{173.779}{703.148}{1}%
\TextAt{9}{178.512}{704.941}{\Character{41}}%
\TextAt{13}{183.064}{704.941}{m}%
\TextAt{17}{193.303}{709.28}{2}%
\TextAt{21}{197.537}{709.28}{n}%
\TextAt{3}{202.676}{709.28}{\Character{0}}%
\TextAt{17}{209.262}{709.28}{1}%
\NativeMathAt{217.315}{704.941}{11.955}{\leq}%
\TextAt{9}{229.935}{704.941}{\Character{40}}%
\TextAt{13}{234.487}{704.941}{m}%
\TextAt{21}{244.726}{709.28}{n}%
\TextAt{9}{253.019}{704.941}{\Character{43}}%
\TextAt{13}{264.781}{704.941}{\Character{34}}%
\TextAt{21}{270.259}{709.28}{n}%
\TextAt{9}{275.895}{704.941}{\Character{41}}%
\TextAt{17}{280.447}{709.28}{2}%
\TextAt{8}{289.077}{704.941}{and}%
\TextAt{8}{311.84}{704.941}{hence}%
\TextAt{13}{245.326}{680.944}{f}%
\TextAt{9}{252.373}{680.944}{\Character{40}}%
\TextAt{13}{256.925}{680.944}{\Character{11}}%
\TextAt{17}{264.371}{679.151}{1}%
\TextAt{9}{269.103}{680.944}{\Character{41}}%
\Vector{q 0 G 0 g 0.478 w 173.326 -98.155 m 201.655 -98.155 l S Q}%
\TextAt{13}{254.371}{664.656}{m}%
\NativeMathAt{278.172}{672.857}{11.955}{\leq}%
\TextAt{20}{290.791}{686.127}{\Character{16}}%
\TextAt{9}{297.931}{672.857}{1}%
\TextAt{9}{306.431}{672.857}{\Character{43}}%
\TextAt{20}{318.202}{686.127}{\Character{16}}%
\TextAt{13}{328.918}{680.944}{\Character{34}}%
\Vector{q 0 G 0 g 0.478 w 254.537 -98.155 m 264.777 -98.155 l S Q}%
\TextAt{13}{326.537}{664.656}{m}%
\TextAt{20}{337.972}{686.127}{\Character{17}}%
\TextAt{21}{345.112}{683.45}{n}%
\TextAt{20}{350.748}{686.127}{\Character{17}}%
\TextAt{17}{357.888}{683.727}{2}%
\TextAt{13}{364.613}{672.857}{\Character{58}}%
\TextAt{8}{79.226}{643.648}{For}%
\TextAt{13}{101.402}{643.648}{n}%
\NativeMathAt{113.767}{643.648}{11.955}{\rightarrow}%
\NativeMathAt{131.09}{643.648}{11.955}{\infty}%
\TextAt{8}{148.158}{643.648}{we}%
\TextAt{8}{166.592}{643.648}{have}%
\TextAt{13}{194.793}{643.648}{f}%
\TextAt{9}{201.84}{643.648}{\Character{40}}%
\TextAt{13}{206.392}{643.648}{\Character{11}}%
\TextAt{17}{213.838}{641.855}{1}%
\TextAt{9}{218.57}{643.648}{\Character{41}}%
\NativeMathAt{228.499}{643.648}{11.955}{\leq}%
\TextAt{13}{243.175}{643.648}{m}%
\TextAt{8}{253.414}{643.648}{,}%
\TextAt{8}{262.071}{643.648}{which}%
\TextAt{8}{296.764}{643.648}{contradicts}%
\TextAt{8}{358.201}{643.648}{the}%
\TextAt{8}{379.564}{643.648}{inequality}%
\TextAt{13}{435.083}{643.648}{f}%
\TextAt{9}{442.129}{643.648}{\Character{40}}%
\TextAt{13}{446.682}{643.648}{\Character{11}}%
\TextAt{17}{454.127}{641.855}{1}%
\TextAt{9}{458.859}{643.648}{\Character{41}}%
\TextAt{13}{468.789}{643.648}{\Character{62}}%
\TextAt{13}{483.255}{643.648}{m}%
\TextAt{8}{498.614}{643.648}{proved}%
\TextAt{8}{79.226}{626.602}{before.}%
\Vector{q 0 G 0 g 0.399 w 452.334 -147.398 m 452.334 -139.328 l S Q}%
\Vector{q 0 G 0 g 0.399 w 452.534 -139.527 m 459.707 -139.527 l S Q}%
\Vector{q 0 G 0 g 0.399 w 452.534 -147.199 m 459.707 -147.199 l S Q}%
\Vector{q 0 G 0 g 0.399 w 459.906 -147.398 m 459.906 -139.328 l S Q}%
\TextAt{8}{97.159}{597.861}{In}%
\TextAt{8}{111.755}{597.861}{view}%
\TextAt{8}{138.712}{597.861}{of}%
\TextAt{26}{152.531}{597.861}{Ostrowski}%
\TextAt{8}{216.121}{597.861}{’s}%
\TextAt{8}{227.86}{597.861}{theorem,}%
\TextAt{8}{276.616}{597.861}{we}%
\TextAt{8}{293.819}{597.861}{will}%
\TextAt{8}{315.899}{597.861}{henceforth}%
\TextAt{8}{373.426}{597.861}{restrict}%
\TextAt{8}{413.785}{597.861}{attention}%
\TextAt{8}{464.175}{597.861}{to}%
\TextAt{8}{478.46}{597.861}{the}%
\TextAt{8}{498.591}{597.861}{case}%
\TextAt{8}{523.337}{597.861}{of}%
\TextAt{8}{79.226}{580.815}{nonarchimedean}%
\TextAt{8}{165.936}{580.815}{valuations.}%
\TextAt{8}{226.133}{580.815}{In}%
\TextAt{8}{240.634}{580.815}{this}%
\TextAt{8}{263.335}{580.815}{case}%
\TextAt{8}{287.985}{580.815}{it}%
\TextAt{8}{299.568}{580.815}{is}%
\TextAt{8}{311.211}{580.815}{usually}%
\TextAt{8}{351.152}{580.815}{expedient—both}%
\TextAt{8}{439.476}{580.815}{with}%
\TextAt{8}{466.015}{580.815}{regard}%
\TextAt{8}{502.321}{580.815}{to}%
\TextAt{8}{516.51}{580.815}{the}%
\TextAt{8}{79.226}{563.769}{substance}%
\TextAt{8}{133.132}{563.769}{and}%
\TextAt{8}{156.348}{563.769}{to}%
\TextAt{8}{171.113}{563.769}{practical}%
\TextAt{8}{219.698}{563.769}{technique—to}%
\TextAt{8}{294.944}{563.769}{work}%
\TextAt{8}{324.019}{563.769}{with}%
\TextAt{8}{351.133}{563.769}{the}%
\TextAt{8}{371.743}{563.769}{exponential}%
\TextAt{8}{434.962}{563.769}{valuations}%
\TextAt{13}{491.103}{563.769}{v}%
\TextAt{8}{501.543}{563.769}{rather}%
\TextAt{8}{79.226}{546.723}{than}%
\TextAt{8}{106.248}{546.723}{the}%
\TextAt{8}{126.108}{546.723}{multiplicative}%
\TextAt{8}{199.635}{546.723}{valuations.}%
\TextAt{8}{259.775}{546.723}{So}%
\TextAt{8}{275.738}{546.723}{let}%
\TextAt{13}{292.346}{546.723}{v}%
\TextAt{8}{302.036}{546.723}{be}%
\TextAt{8}{317.664}{546.723}{an}%
\TextAt{8}{333.627}{546.723}{exponential}%
\TextAt{8}{396.095}{546.723}{valuation}%
\TextAt{8}{446.872}{546.723}{of}%
\TextAt{8}{459.906}{546.723}{the}%
\TextAt{8}{479.766}{546.723}{field}%
\TextAt{13}{504.827}{546.723}{K}%
\TextAt{8}{515.645}{546.723}{.}%
\TextAt{8}{523.994}{546.723}{It}%
\TextAt{8}{79.226}{529.677}{is}%
\TextAt{8}{90.99}{529.677}{canonically}%
\TextAt{8}{151.806}{529.677}{continued}%
\TextAt{8}{205.138}{529.677}{to}%
\TextAt{8}{219.448}{529.677}{an}%
\TextAt{8}{235.707}{529.677}{exponential}%
\TextAt{8}{298.472}{529.677}{valuation}%
\TextAt{9}{349.987}{529.677}{\Character{94}}%
\TextAt{13}{349.544}{529.677}{v}%
\TextAt{8}{359.53}{529.677}{of}%
\TextAt{8}{372.86}{529.677}{the}%
\TextAt{8}{393.017}{529.677}{completion}%
\TextAt{20}{455.59}{532.699}{\Character{98}}%
\TextAt{13}{452.852}{529.677}{K}%
\TextAt{8}{467.567}{529.677}{by}%
\TextAt{8}{483.826}{529.677}{setting}%
\TextAt{9}{260.752}{501.586}{\Character{94}}%
\TextAt{13}{260.309}{501.586}{v}%
\TextAt{9}{266.397}{501.586}{\Character{40}}%
\TextAt{13}{270.95}{501.586}{a}%
\TextAt{9}{277.095}{501.586}{\Character{41}}%
\TextAt{9}{284.95}{501.586}{\Character{61}}%
\TextAt{8}{300.301}{501.586}{lim}%
\TextAt{21}{297.393}{494.413}{n}%
\NativeMathAt{302.531}{494.413}{7.97}{\rightarrow}%
\NativeMathAt{310.995}{494.413}{7.97}{\infty}%
\TextAt{13}{321.461}{501.586}{v}%
\TextAt{9}{327.549}{501.586}{\Character{40}}%
\TextAt{13}{332.101}{501.586}{a}%
\TextAt{21}{338.246}{499.792}{n}%
\TextAt{9}{343.882}{501.586}{\Character{41}}%
\TextAt{13}{348.435}{501.586}{\Character{59}}%
\TextAt{8}{79.226}{469.519}{where}%
\TextAt{13}{112.895}{469.519}{a}%
\TextAt{9}{122.361}{469.519}{\Character{61}}%
\TextAt{8}{134.786}{469.519}{lim}%
\TextAt{21}{151.045}{467.726}{n}%
\NativeMathAt{156.184}{467.726}{7.97}{\rightarrow}%
\NativeMathAt{164.648}{467.726}{7.97}{\infty}%
\TextAt{13}{175.611}{469.519}{a}%
\TextAt{21}{181.756}{467.726}{n}%
\NativeMathAt{190.713}{469.519}{11.955}{\in}%
\TextAt{20}{204.742}{472.541}{\Character{98}}%
\TextAt{13}{202.004}{469.519}{K}%
\TextAt{8}{212.822}{469.519}{,}%
\TextAt{13}{219.859}{469.519}{a}%
\TextAt{21}{226.004}{467.726}{n}%
\NativeMathAt{234.961}{469.519}{11.955}{\in}%
\TextAt{13}{246.252}{469.519}{K}%
\TextAt{8}{257.069}{469.519}{.}%
\TextAt{8}{265.471}{469.519}{Observe}%
\TextAt{8}{310.27}{469.519}{here}%
\TextAt{8}{335.487}{469.519}{that}%
\TextAt{8}{360.715}{469.519}{the}%
\TextAt{8}{380.732}{469.519}{sequence}%
\TextAt{13}{429.093}{469.519}{v}%
\TextAt{9}{435.182}{469.519}{\Character{40}}%
\TextAt{13}{439.734}{469.519}{a}%
\TextAt{21}{445.879}{467.726}{n}%
\TextAt{9}{451.515}{469.519}{\Character{41}}%
\TextAt{8}{459.825}{469.519}{has}%
\TextAt{8}{480.558}{469.519}{to}%
\TextAt{8}{494.728}{469.519}{become}%
\TextAt{8}{79.226}{452.473}{stationary}%
\TextAt{8}{135.809}{452.473}{(provided}%
\TextAt{13}{189.39}{452.473}{a}%
\NativeMathAt{200.382}{452.473}{11.955}{\neq}%
\TextAt{9}{214.322}{452.473}{0}%
\TextAt{8}{220.186}{452.473}{)}%
\TextAt{8}{229.534}{452.473}{because,}%
\TextAt{8}{277.206}{452.473}{for}%
\TextAt{13}{295.986}{452.473}{n}%
\NativeMathAt{307.821}{452.473}{11.955}{\geq}%
\TextAt{13}{321.966}{452.473}{n}%
\TextAt{17}{328.953}{450.68}{0}%
\TextAt{8}{333.686}{452.473}{,}%
\TextAt{8}{341.954}{452.473}{one}%
\TextAt{8}{364.309}{452.473}{has}%
\TextAt{9}{386.521}{452.473}{\Character{94}}%
\TextAt{13}{386.078}{452.473}{v}%
\TextAt{9}{392.166}{452.473}{\Character{40}}%
\TextAt{13}{396.718}{452.473}{a}%
\TextAt{16}{406.13}{452.473}{\Character{0}}%
\TextAt{13}{418.696}{452.473}{a}%
\TextAt{21}{424.841}{450.68}{n}%
\TextAt{9}{430.477}{452.473}{\Character{41}}%
\TextAt{13}{439.876}{452.473}{\Character{62}}%
\TextAt{9}{454.27}{452.473}{\Character{94}}%
\TextAt{13}{453.827}{452.473}{v}%
\TextAt{9}{459.915}{452.473}{\Character{40}}%
\TextAt{13}{464.468}{452.473}{a}%
\TextAt{9}{470.613}{452.473}{\Character{41}}%
\TextAt{8}{475.165}{452.473}{,}%
\TextAt{8}{483.433}{452.473}{so}%
\TextAt{8}{498.699}{452.473}{that}%
\TextAt{8}{524.963}{452.473}{it}%
\TextAt{8}{79.226}{435.427}{follows}%
\TextAt{8}{117.662}{435.427}{from}%
\TextAt{8}{145.302}{435.427}{the}%
\TextAt{8}{165.459}{435.427}{remark}%
\TextAt{8}{205.461}{435.427}{on}%
\TextAt{8}{221.72}{435.427}{p.}%
\TextAt{8}{235.373}{435.427}{119}%
\TextAt{13}{169.56}{407.336}{v}%
\TextAt{9}{175.648}{407.336}{\Character{40}}%
\TextAt{13}{180.201}{407.336}{a}%
\TextAt{21}{186.346}{405.543}{n}%
\TextAt{9}{191.982}{407.336}{\Character{41}}%
\TextAt{9}{199.837}{407.336}{\Character{61}}%
\TextAt{9}{212.7}{407.336}{\Character{94}}%
\TextAt{13}{212.281}{407.336}{v}%
\TextAt{9}{218.369}{407.336}{\Character{40}}%
\TextAt{13}{222.921}{407.336}{a}%
\TextAt{21}{229.066}{405.543}{n}%
\TextAt{16}{237.359}{407.336}{\Character{0}}%
\TextAt{13}{249.314}{407.336}{a}%
\TextAt{9}{258.116}{407.336}{\Character{43}}%
\TextAt{13}{269.877}{407.336}{a}%
\TextAt{9}{276.022}{407.336}{\Character{41}}%
\TextAt{9}{283.877}{407.336}{\Character{61}}%
\TextAt{8}{296.321}{407.336}{min}%
\TextAt{16}{315.832}{407.336}{\Character{102}}%
\TextAt{9}{322.252}{407.336}{\Character{94}}%
\TextAt{13}{321.809}{407.336}{v}%
\TextAt{9}{327.897}{407.336}{\Character{40}}%
\TextAt{13}{332.45}{407.336}{a}%
\TextAt{21}{338.595}{405.543}{n}%
\TextAt{16}{346.888}{407.336}{\Character{0}}%
\TextAt{13}{358.843}{407.336}{a}%
\TextAt{9}{364.988}{407.336}{\Character{41}}%
\TextAt{13}{369.54}{407.336}{\Character{59}}%
\TextAt{9}{375.227}{407.336}{\Character{94}}%
\TextAt{13}{374.784}{407.336}{v}%
\TextAt{9}{380.873}{407.336}{\Character{40}}%
\TextAt{13}{385.425}{407.336}{a}%
\TextAt{9}{391.57}{407.336}{\Character{41}}%
\TextAt{16}{396.122}{407.336}{\Character{103}}%
\TextAt{9}{405.421}{407.336}{\Character{61}}%
\TextAt{9}{418.273}{407.336}{\Character{94}}%
\TextAt{13}{417.846}{407.336}{v}%
\TextAt{9}{423.934}{407.336}{\Character{40}}%
\TextAt{13}{428.487}{407.336}{a}%
\TextAt{9}{434.632}{407.336}{\Character{41}}%
\TextAt{13}{439.184}{407.336}{\Character{58}}%
\TextAt{8}{79.226}{379.245}{Therefore}%
\TextAt{8}{132.223}{379.245}{it}%
\TextAt{8}{143.928}{379.245}{follows}%
\TextAt{8}{182.363}{379.245}{that}%
\TextAt{13}{265.756}{362.199}{v}%
\TextAt{9}{271.844}{362.199}{\Character{40}}%
\TextAt{13}{276.397}{362.199}{K}%
\TextAt{3}{287.214}{367.135}{\Character{3}}%
\TextAt{9}{291.947}{362.199}{\Character{41}}%
\TextAt{9}{299.802}{362.199}{\Character{61}}%
\TextAt{9}{312.665}{362.199}{\Character{94}}%
\TextAt{13}{312.245}{362.199}{v}%
\TextAt{9}{318.333}{362.199}{\Character{40}}%
\TextAt{20}{325.624}{365.221}{\Character{98}}%
\TextAt{13}{322.886}{362.199}{K}%
\TextAt{3}{333.703}{367.135}{\Character{3}}%
\TextAt{9}{338.436}{362.199}{\Character{41}}%
\TextAt{13}{342.988}{362.199}{\Character{59}}%
\TextAt{8}{79.226}{339.067}{and}%
\TextAt{8}{102.714}{339.067}{if}%
\TextAt{13}{114.164}{339.067}{v}%
\TextAt{8}{124.875}{339.067}{is}%
\TextAt{8}{137.364}{339.067}{discrete}%
\TextAt{8}{181.068}{339.067}{and}%
\TextAt{8}{204.557}{339.067}{normalized,}%
\TextAt{8}{268.551}{339.067}{then}%
\TextAt{8}{295.936}{339.067}{so}%
\TextAt{8}{311.032}{339.067}{is}%
\TextAt{8}{323.521}{339.067}{the}%
\TextAt{8}{344.403}{339.067}{extension}%
\TextAt{9}{397.338}{339.067}{\Character{94}}%
\TextAt{13}{396.895}{339.067}{v}%
\TextAt{8}{402.983}{339.067}{.}%
\TextAt{8}{413.608}{339.067}{In}%
\TextAt{8}{428.955}{339.067}{the}%
\TextAt{8}{449.836}{339.067}{nonarchimedean}%
\TextAt{8}{79.226}{322.021}{case,}%
\TextAt{8}{108.693}{322.021}{for}%
\TextAt{8}{127.733}{322.021}{a}%
\TextAt{8}{138.643}{322.021}{sequence}%
\TextAt{16}{188.3}{322.021}{\Character{102}}%
\TextAt{13}{194.278}{322.021}{a}%
\TextAt{21}{200.423}{320.228}{n}%
\TextAt{16}{206.059}{322.021}{\Character{103}}%
\TextAt{21}{212.037}{320.228}{n}%
\NativeMathAt{217.175}{320.228}{7.97}{\in}%
\TextAt{25}{222.82}{320.228}{N}%
\TextAt{8}{234.127}{322.021}{to}%
\TextAt{8}{249.592}{322.021}{be}%
\TextAt{8}{266.672}{322.021}{a}%
\TextAt{8}{277.582}{322.021}{Cauchy}%
\TextAt{8}{320.688}{322.021}{sequence,}%
\TextAt{8}{373.885}{322.021}{it}%
\TextAt{8}{386.745}{322.021}{suffices}%
\TextAt{8}{427.686}{322.021}{that}%
\TextAt{13}{454.21}{322.021}{a}%
\TextAt{21}{460.355}{320.228}{n}%
\TextAt{17}{465.493}{320.228}{\Character{43}1}%
\TextAt{16}{480.256}{322.021}{\Character{0}}%
\TextAt{13}{492.999}{322.021}{a}%
\TextAt{21}{499.143}{320.228}{n}%
\TextAt{8}{509.832}{322.021}{be}%
\TextAt{8}{526.911}{322.021}{a}%
\TextAt{8}{79.226}{304.975}{nullsequence.}%
\TextAt{8}{156.008}{304.975}{In}%
\TextAt{8}{172.143}{304.975}{fact,}%
\TextAt{13}{200.373}{304.975}{v}%
\TextAt{9}{206.461}{304.975}{\Character{40}}%
\TextAt{13}{211.013}{304.975}{a}%
\TextAt{21}{217.158}{303.182}{n}%
\TextAt{16}{226.483}{304.975}{\Character{0}}%
\TextAt{13}{239.47}{304.975}{a}%
\TextAt{21}{245.615}{303.182}{m}%
\TextAt{9}{253.603}{304.975}{\Character{41}}%
\NativeMathAt{264.056}{304.975}{11.955}{\geq}%
\TextAt{8}{279.255}{304.975}{min}%
\TextAt{21}{298.766}{303.182}{m}%
\NativeMathAt{306.256}{303.182}{7.97}{\leq}%
\TextAt{21}{312.843}{303.182}{i\Character{60}n}%
\TextAt{16}{327.949}{304.975}{\Character{102}}%
\TextAt{13}{333.927}{304.975}{v}%
\TextAt{9}{340.015}{304.975}{\Character{40}}%
\TextAt{13}{344.567}{304.975}{a}%
\TextAt{21}{350.712}{303.182}{i}%
\TextAt{17}{353.595}{303.182}{\Character{43}1}%
\TextAt{16}{368.602}{304.975}{\Character{0}}%
\TextAt{13}{381.589}{304.975}{a}%
\TextAt{21}{387.734}{303.182}{i}%
\TextAt{9}{391.116}{304.975}{\Character{41}}%
\TextAt{16}{395.668}{304.975}{\Character{103}}%
\TextAt{8}{401.646}{304.975}{.}%
\TextAt{8}{414.635}{304.975}{By}%
\TextAt{8}{434.512}{304.975}{the}%
\TextAt{8}{456.182}{304.975}{same}%
\TextAt{8}{487.022}{304.975}{token}%
\TextAt{8}{520.408}{304.975}{an}%
\TextAt{8}{79.226}{287.929}{infinite}%
\TextAt{8}{120.204}{287.929}{series}%
\TextAt{20}{152.849}{296.896}{\Character{80}}%
\NativeMathAt{165.468}{293.741}{7.97}{\infty}%
\TextAt{21}{165.468}{284.442}{\Character{23}}%
\TextAt{17}{170.152}{284.442}{\Character{61}0}%
\TextAt{13}{183.463}{287.929}{a}%
\TextAt{21}{189.608}{286.136}{\Character{23}}%
\TextAt{8}{199.998}{287.929}{converges}%
\TextAt{8}{253.732}{287.929}{in}%
\TextAt{20}{271.434}{290.951}{\Character{98}}%
\TextAt{13}{268.695}{287.929}{K}%
\TextAt{8}{284.721}{287.929}{if}%
\TextAt{8}{296.755}{287.929}{and}%
\TextAt{8}{320.828}{287.929}{only}%
\TextAt{8}{347.83}{287.929}{if}%
\TextAt{8}{359.865}{287.929}{the}%
\TextAt{8}{381.332}{287.929}{sequence}%
\TextAt{8}{431.144}{287.929}{of}%
\TextAt{8}{445.785}{287.929}{its}%
\TextAt{8}{463.414}{287.929}{terms}%
\TextAt{13}{497.302}{287.929}{a}%
\TextAt{21}{503.447}{286.136}{\Character{23}}%
\TextAt{8}{513.837}{287.929}{is}%
\TextAt{8}{526.911}{287.929}{a}%
\TextAt{8}{79.226}{270.883}{nullsequence.}%
\TextAt{8}{155.33}{270.883}{The}%
\TextAt{8}{180.671}{270.883}{following}%
\TextAt{8}{231.394}{270.883}{proposition}%
\TextAt{8}{294.214}{270.883}{is}%
\TextAt{8}{307.266}{270.883}{proved}%
\TextAt{8}{346.607}{270.883}{exactly}%
\TextAt{8}{388.219}{270.883}{as}%
\TextAt{8}{403.877}{270.883}{its}%
\TextAt{8}{421.483}{270.883}{analogue,}%
\TextAt{8}{475.134}{270.883}{proposition}%
\TextAt{8}{79.226}{253.837}{(2.4),}%
\TextAt{8}{110.453}{253.837}{in}%
\TextAt{8}{124.106}{253.837}{the}%
\TextAt{8}{144.262}{253.837}{special}%
\TextAt{8}{182.363}{253.837}{case}%
\TextAt{9}{207.135}{253.837}{\Character{40}}%
\TextAt{18}{211.687}{253.837}{Q}%
\TextAt{13}{220.985}{253.837}{\Character{59}}%
\TextAt{13}{226.209}{253.837}{v}%
\TextAt{21}{231.889}{252.044}{p}%
\TextAt{9}{236.65}{253.837}{\Character{41}}%
\TextAt{8}{241.202}{253.837}{.}%
\TextAt{12}{79.226}{228.085}{Proposition}%
\TextAt{12}{152.081}{228.085}{(4.3).}%
\TextAt{8}{187.521}{228.085}{If}%
\TextAt{14}{199.338}{228.085}{o}%
\NativeMathAt{208.715}{228.085}{11.955}{\subseteq}%
\TextAt{13}{221.547}{228.085}{K}%
\TextAt{8}{232.365}{228.085}{,}%
\TextAt{8}{239.67}{228.085}{resp.}%
\TextAt{20}{268.967}{228.62}{\Character{98}}%
\TextAt{14}{269.366}{228.085}{o}%
\NativeMathAt{278.744}{228.085}{11.955}{\subseteq}%
\TextAt{20}{294.314}{231.107}{\Character{98}}%
\TextAt{13}{291.575}{228.085}{K}%
\TextAt{8}{302.393}{228.085}{,}%
\TextAt{8}{309.698}{228.085}{is}%
\TextAt{8}{321.587}{228.085}{the}%
\TextAt{8}{341.868}{228.085}{valuation}%
\TextAt{8}{393.065}{228.085}{ring}%
\TextAt{8}{417.255}{228.085}{of}%
\TextAt{13}{430.71}{228.085}{v}%
\TextAt{8}{436.798}{228.085}{,}%
\TextAt{8}{444.103}{228.085}{resp.}%
\TextAt{8}{473.799}{228.085}{of}%
\TextAt{9}{487.697}{228.085}{\Character{94}}%
\TextAt{13}{487.254}{228.085}{v}%
\TextAt{8}{493.342}{228.085}{,}%
\TextAt{8}{500.647}{228.085}{and}%
\TextAt{14}{523.535}{228.085}{p}%
\TextAt{8}{529.518}{228.085}{,}%
\TextAt{8}{79.226}{211.039}{resp.}%
\TextAt{20}{108.219}{212.155}{\Character{98}}%
\TextAt{14}{108.548}{211.039}{p}%
\TextAt{8}{114.531}{211.039}{,}%
\TextAt{8}{121.68}{211.039}{is}%
\TextAt{8}{133.444}{211.039}{the}%
\TextAt{8}{153.601}{211.039}{maximal}%
\TextAt{8}{201.409}{211.039}{ideal,}%
\TextAt{8}{232.624}{211.039}{then}%
\TextAt{8}{259.284}{211.039}{one}%
\TextAt{8}{280.744}{211.039}{has}%
\TextAt{20}{279.949}{183.483}{\Character{98}}%
\TextAt{14}{280.347}{182.948}{o}%
\TextAt{9}{286.192}{182.948}{\Character{47}}%
\TextAt{20}{291.715}{184.064}{\Character{98}}%
\TextAt{14}{292.045}{182.948}{p}%
\NativeMathAt{301.349}{182.947}{11.955}{\cong}%
\TextAt{14}{313.968}{182.948}{o}%
\TextAt{9}{319.812}{182.948}{\Character{47}}%
\TextAt{14}{325.665}{182.948}{p}%
\TextAt{8}{79.226}{154.856}{and,}%
\TextAt{8}{105.241}{154.856}{if}%
\TextAt{13}{115.964}{154.856}{v}%
\TextAt{8}{125.95}{154.856}{is}%
\TextAt{8}{137.714}{154.856}{discrete,}%
\TextAt{8}{183.945}{154.856}{one}%
\TextAt{8}{205.404}{154.856}{has}%
\TextAt{8}{226.278}{154.856}{furthermore}%
\TextAt{20}{236.693}{127.3}{\Character{98}}%
\TextAt{14}{237.092}{126.765}{o}%
\TextAt{9}{242.936}{126.765}{\Character{47}}%
\TextAt{20}{248.459}{127.882}{\Character{98}}%
\TextAt{14}{248.789}{126.765}{p}%
\TextAt{21}{256.183}{131.701}{n}%
\NativeMathAt{265.14}{126.765}{11.955}{\cong}%
\TextAt{14}{277.76}{126.765}{o}%
\TextAt{9}{283.604}{126.765}{\Character{47}}%
\TextAt{14}{289.457}{126.765}{p}%
\TextAt{21}{295.44}{131.701}{n}%
\TextAt{8}{324.987}{126.765}{for}%
\TextAt{13}{342.872}{126.765}{n}%
\NativeMathAt{353.18}{126.765}{11.955}{\geq}%
\TextAt{9}{365.799}{126.765}{1}%
\TextAt{13}{371.652}{126.765}{\Character{58}}%
\TextAt{8}{97.159}{89.967}{Generalizing}%
\TextAt{8}{165.397}{89.967}{the}%
\TextAt{13}{186.078}{89.967}{p}%
\TextAt{8}{191.954}{89.967}{-adic}%
\TextAt{8}{221.087}{89.967}{expansion}%
\TextAt{8}{275.983}{89.967}{to}%
\TextAt{8}{290.818}{89.967}{the}%
\TextAt{8}{311.499}{89.967}{case}%
\TextAt{8}{336.795}{89.967}{of}%
\TextAt{8}{350.649}{89.967}{an}%
\TextAt{8}{367.433}{89.967}{arbitrary}%
\TextAt{8}{417.726}{89.967}{discrete}%
\TextAt{8}{461.23}{89.967}{valuation}%
\TextAt{13}{512.827}{89.967}{v}%
\TextAt{8}{523.337}{89.967}{of}%
\TextAt{8}{79.226}{72.921}{the}%
\TextAt{8}{99.383}{72.921}{field}%
\TextAt{13}{124.739}{72.921}{K}%
\TextAt{8}{135.557}{72.921}{,}%
\TextAt{8}{142.706}{72.921}{we}%
\TextAt{8}{159.934}{72.921}{have}%
\TextAt{8}{186.929}{72.921}{the}%
\TextAt{8}{297.211}{43.033}{115}%
\Destination{78.23}{809.815}{nav.page.115}%
\Destination{79.36}{561.755}{nav.print.126}%
\Destination{79.36}{242.055}{nav.number.2.4.3}%
\Hotlink{234.873}{433.255}{18.574}{11.047}{nav.print.119}%
\Hotlink{78.726}{251.655}{25.078}{11.047}{nav.number.2.2.4}%
}

% Page 128 (body page 116)
\Page{%
\TextAt{2}{79.226}{796.943}{II.}%
\TextAt{2}{94.073}{796.943}{The}%
\TextAt{2}{116.492}{796.943}{Theory}%
\TextAt{2}{154.401}{796.943}{of}%
\TextAt{2}{166.826}{796.943}{Valuations}%
\Vector{q 0 G 0 g 0.299 w 7.226 18.729 m 460.769 18.729 l S Q}%
\TextAt{12}{79.226}{761.124}{Proposition}%
\TextAt{12}{153.297}{761.124}{(4.4).}%
\TextAt{8}{189.794}{761.124}{Let}%
\TextAt{13}{211.945}{761.124}{R}%
\NativeMathAt{226.288}{761.124}{11.955}{\subseteq}%
\TextAt{14}{240.921}{761.124}{o}%
\TextAt{8}{251.844}{761.124}{be}%
\TextAt{8}{268.95}{761.124}{a}%
\TextAt{8}{279.887}{761.124}{system}%
\TextAt{8}{319.887}{761.124}{of}%
\TextAt{8}{334.399}{761.124}{representatives}%
\TextAt{8}{415.381}{761.124}{for}%
\TextAt{13}{434.448}{761.124}{\Character{20}}%
\TextAt{9}{446.521}{761.124}{\Character{61}}%
\TextAt{14}{460.96}{761.124}{o}%
\TextAt{9}{466.804}{761.124}{\Character{47}}%
\TextAt{14}{472.657}{761.124}{p}%
\TextAt{8}{483.719}{761.124}{such}%
\TextAt{8}{511.298}{761.124}{that}%
\TextAt{9}{79.226}{744.078}{0}%
\NativeMathAt{90.564}{744.078}{11.955}{\in}%
\TextAt{13}{104.019}{744.078}{R}%
\TextAt{8}{113.028}{744.078}{,}%
\TextAt{8}{121.765}{744.078}{and}%
\TextAt{8}{145.797}{744.078}{let}%
\TextAt{13}{163.972}{744.078}{\Character{25}}%
\NativeMathAt{176.526}{744.078}{11.955}{\in}%
\TextAt{14}{189.982}{744.078}{o}%
\TextAt{8}{200.993}{744.078}{be}%
\TextAt{8}{218.188}{744.078}{a}%
\TextAt{8}{229.213}{744.078}{prime}%
\TextAt{8}{263.647}{744.078}{element.}%
\TextAt{8}{315.249}{744.078}{Then}%
\TextAt{8}{347.077}{744.078}{every}%
\TextAt{13}{379.239}{744.078}{x}%
\NativeMathAt{391.376}{744.078}{11.955}{\neq}%
\TextAt{9}{405.949}{744.078}{0}%
\TextAt{8}{416.986}{744.078}{in}%
\TextAt{20}{434.647}{747.1}{\Character{98}}%
\TextAt{13}{431.909}{744.078}{K}%
\TextAt{8}{447.894}{744.078}{admits}%
\TextAt{8}{487.6}{744.078}{a}%
\TextAt{8}{498.626}{744.078}{unique}%
\TextAt{8}{79.226}{727.032}{representation}%
\TextAt{8}{155.715}{727.032}{as}%
\TextAt{8}{170.085}{727.032}{a}%
\TextAt{8}{179.841}{727.032}{convergent}%
\TextAt{8}{238.385}{727.032}{series}%
\TextAt{13}{227.369}{698.648}{x}%
\TextAt{9}{237.342}{698.648}{\Character{61}}%
\TextAt{13}{249.767}{698.648}{\Character{25}}%
\TextAt{21}{256.837}{703.585}{m}%
\TextAt{9}{264.825}{698.648}{\Character{40}}%
\TextAt{13}{269.378}{698.648}{a}%
\TextAt{17}{275.523}{696.855}{0}%
\TextAt{9}{282.912}{698.648}{\Character{43}}%
\TextAt{13}{294.673}{698.648}{a}%
\TextAt{17}{300.818}{696.855}{1}%
\TextAt{13}{305.55}{698.648}{\Character{25}}%
\TextAt{9}{315.276}{698.648}{\Character{43}}%
\TextAt{13}{327.037}{698.648}{a}%
\TextAt{17}{333.182}{696.855}{2}%
\TextAt{13}{337.915}{698.648}{\Character{25}}%
\TextAt{17}{344.984}{703.585}{2}%
\TextAt{9}{352.373}{698.648}{\Character{43}}%
\TextAt{16}{364.134}{698.648}{\Character{1}}%
\TextAt{16}{369.43}{698.648}{\Character{1}}%
\TextAt{16}{374.738}{698.648}{\Character{1}}%
\TextAt{9}{380.074}{698.648}{\Character{41}}%
\TextAt{8}{79.226}{670.265}{where}%
\TextAt{13}{113.035}{670.265}{a}%
\TextAt{21}{119.18}{668.472}{i}%
\NativeMathAt{125.882}{670.265}{11.955}{\in}%
\TextAt{13}{137.173}{670.265}{R}%
\TextAt{8}{146.182}{670.265}{,}%
\TextAt{13}{153.331}{670.265}{a}%
\TextAt{17}{159.476}{668.472}{0}%
\NativeMathAt{167.529}{670.265}{11.955}{\neq}%
\TextAt{9}{179.939}{670.265}{0}%
\TextAt{8}{185.808}{670.265}{,}%
\TextAt{13}{192.957}{670.265}{m}%
\NativeMathAt{206.517}{670.265}{11.955}{\in}%
\TextAt{18}{217.808}{670.265}{Z}%
\TextAt{8}{225.778}{670.265}{.}%
\TextAt{19}{79.226}{635.64}{Proof.}%
\TextAt{8}{116}{635.64}{Let}%
\TextAt{13}{138.364}{635.64}{x}%
\TextAt{9}{150.714}{635.64}{\Character{61}}%
\TextAt{13}{165.516}{635.64}{\Character{25}}%
\TextAt{21}{172.585}{639.978}{m}%
\TextAt{13}{180.574}{635.64}{u}%
\TextAt{8}{192.528}{635.64}{with}%
\TextAt{13}{220.583}{635.64}{u}%
\NativeMathAt{232.942}{635.64}{11.955}{\in}%
\TextAt{20}{246.211}{636.175}{\Character{98}}%
\TextAt{14}{246.61}{635.64}{o}%
\TextAt{3}{253.866}{639.978}{\Character{3}}%
\TextAt{8}{258.598}{635.64}{.}%
\TextAt{8}{271.23}{635.64}{Since}%
\TextAt{20}{302.783}{636.175}{\Character{98}}%
\TextAt{14}{303.182}{635.64}{o}%
\TextAt{9}{309.026}{635.64}{\Character{47}}%
\TextAt{20}{314.55}{636.756}{\Character{98}}%
\TextAt{14}{314.879}{635.64}{p}%
\NativeMathAt{326.56}{635.639}{11.955}{\cong}%
\TextAt{14}{341.556}{635.64}{o}%
\TextAt{9}{347.4}{635.64}{\Character{47}}%
\TextAt{14}{353.253}{635.64}{p}%
\TextAt{8}{359.236}{635.64}{,}%
\TextAt{8}{368.128}{635.64}{the}%
\TextAt{8}{389.679}{635.64}{class}%
\TextAt{13}{418.511}{635.64}{u}%
\TextAt{8}{428.494}{635.64}{mod}%
\TextAt{20}{453.603}{636.756}{\Character{98}}%
\TextAt{14}{453.932}{635.64}{p}%
\TextAt{8}{465.207}{635.64}{has}%
\TextAt{8}{487.476}{635.64}{a}%
\TextAt{8}{498.626}{635.64}{unique}%
\TextAt{8}{79.226}{618.594}{representative}%
\TextAt{13}{154.464}{618.594}{a}%
\TextAt{17}{160.609}{616.8}{0}%
\NativeMathAt{168.751}{618.594}{11.955}{\in}%
\TextAt{13}{180.13}{618.594}{R}%
\TextAt{8}{189.139}{618.594}{,}%
\TextAt{13}{196.353}{618.594}{a}%
\TextAt{17}{202.498}{616.8}{0}%
\NativeMathAt{210.64}{618.594}{11.955}{\neq}%
\TextAt{9}{223.145}{618.594}{0}%
\TextAt{8}{229.007}{618.594}{.}%
\TextAt{8}{237.611}{618.594}{We}%
\TextAt{8}{257.808}{618.594}{thus}%
\TextAt{8}{283.611}{618.594}{have}%
\TextAt{13}{310.658}{618.594}{u}%
\TextAt{9}{320.73}{618.594}{\Character{61}}%
\TextAt{13}{333.244}{618.594}{a}%
\TextAt{17}{339.389}{616.8}{0}%
\TextAt{9}{346.813}{618.594}{\Character{43}}%
\TextAt{13}{358.61}{618.594}{\Character{25}}%
\TextAt{13}{365.664}{618.594}{b}%
\TextAt{17}{370.656}{616.8}{1}%
\TextAt{8}{375.389}{618.594}{,}%
\TextAt{8}{382.603}{618.594}{for}%
\TextAt{8}{400.54}{618.594}{some}%
\TextAt{13}{429.918}{618.594}{b}%
\TextAt{17}{434.895}{616.8}{1}%
\NativeMathAt{443.037}{618.594}{11.955}{\in}%
\TextAt{20}{454.018}{619.129}{\Character{98}}%
\TextAt{14}{454.416}{618.594}{o}%
\TextAt{8}{460.261}{618.594}{.}%
\TextAt{8}{468.865}{618.594}{Assume}%
\TextAt{8}{512.278}{618.594}{now}%
\TextAt{8}{79.226}{601.548}{that}%
\TextAt{13}{104.595}{601.548}{a}%
\TextAt{17}{110.74}{599.754}{0}%
\TextAt{13}{115.472}{601.548}{\Character{59}}%
\TextAt{13}{120.696}{601.548}{\Character{58}}%
\TextAt{13}{125.933}{601.548}{\Character{58}}%
\TextAt{13}{131.169}{601.548}{\Character{58}}%
\TextAt{13}{136.394}{601.548}{\Character{59}}%
\TextAt{13}{141.63}{601.548}{a}%
\TextAt{21}{147.838}{599.754}{n}%
\TextAt{3}{152.976}{599.754}{\Character{0}}%
\TextAt{17}{159.563}{599.754}{1}%
\NativeMathAt{167.616}{601.548}{11.955}{\in}%
\TextAt{13}{178.907}{601.548}{R}%
\TextAt{8}{191.813}{601.548}{have}%
\TextAt{8}{218.808}{601.548}{been}%
\TextAt{8}{246.436}{601.548}{found,}%
\TextAt{8}{282.529}{601.548}{satisfying}%
\TextAt{13}{209.848}{573.164}{u}%
\TextAt{9}{219.831}{573.164}{\Character{61}}%
\TextAt{13}{232.257}{573.164}{a}%
\TextAt{17}{238.402}{571.371}{0}%
\TextAt{9}{245.791}{573.164}{\Character{43}}%
\TextAt{13}{257.552}{573.164}{a}%
\TextAt{17}{263.697}{571.371}{1}%
\TextAt{13}{268.429}{573.164}{\Character{25}}%
\TextAt{9}{278.155}{573.164}{\Character{43}}%
\TextAt{16}{289.917}{573.164}{\Character{1}}%
\TextAt{16}{295.213}{573.164}{\Character{1}}%
\TextAt{16}{300.521}{573.164}{\Character{1}}%
\TextAt{9}{306.521}{573.164}{\Character{43}}%
\TextAt{13}{318.282}{573.164}{a}%
\TextAt{21}{324.427}{571.371}{n}%
\TextAt{3}{329.565}{571.371}{\Character{0}}%
\TextAt{17}{336.152}{571.371}{1}%
\TextAt{13}{340.884}{573.164}{\Character{25}}%
\TextAt{21}{347.954}{578.101}{n}%
\TextAt{3}{353.092}{578.101}{\Character{0}}%
\TextAt{17}{359.678}{578.101}{1}%
\TextAt{9}{367.067}{573.164}{\Character{43}}%
\TextAt{13}{378.829}{573.164}{\Character{25}}%
\TextAt{21}{385.898}{578.101}{n}%
\TextAt{13}{391.534}{573.164}{b}%
\TextAt{21}{396.511}{571.371}{n}%
\TextAt{8}{79.226}{544.781}{for}%
\TextAt{8}{98.13}{544.781}{some}%
\TextAt{13}{128.475}{544.781}{b}%
\TextAt{21}{133.452}{542.988}{n}%
\NativeMathAt{144.146}{544.781}{11.955}{\in}%
\TextAt{20}{156.775}{545.317}{\Character{98}}%
\TextAt{14}{157.174}{544.781}{o}%
\TextAt{8}{163.018}{544.781}{,}%
\TextAt{8}{171.441}{544.781}{and}%
\TextAt{8}{195.223}{544.781}{that}%
\TextAt{8}{221.611}{544.781}{the}%
\TextAt{13}{242.786}{544.781}{a}%
\TextAt{21}{248.931}{542.988}{i}%
\TextAt{8}{257.229}{544.781}{are}%
\TextAt{8}{277.759}{544.781}{uniquely}%
\TextAt{8}{326.252}{544.781}{determined}%
\TextAt{8}{388.398}{544.781}{by}%
\TextAt{8}{405.676}{544.781}{this}%
\TextAt{8}{429.517}{544.781}{equation.}%
\TextAt{8}{484.934}{544.781}{Then}%
\TextAt{8}{516.51}{544.781}{the}%
\TextAt{8}{79.226}{527.735}{representative}%
\TextAt{13}{154.128}{527.735}{a}%
\TextAt{21}{160.273}{525.942}{n}%
\NativeMathAt{169.23}{527.735}{11.955}{\in}%
\TextAt{13}{180.521}{527.735}{R}%
\TextAt{8}{193.143}{527.735}{of}%
\TextAt{13}{206.189}{527.735}{b}%
\TextAt{21}{211.166}{525.942}{n}%
\TextAt{8}{220.124}{527.735}{mod}%
\TextAt{13}{245.562}{527.735}{\Character{25}}%
\TextAt{20}{252.232}{528.271}{\Character{98}}%
\TextAt{14}{252.631}{527.735}{o}%
\NativeMathAt{261.796}{527.735}{11.955}{\in}%
\TextAt{20}{272.688}{528.271}{\Character{98}}%
\TextAt{14}{273.087}{527.735}{o}%
\TextAt{9}{278.931}{527.735}{\Character{47}}%
\TextAt{20}{284.455}{528.852}{\Character{98}}%
\TextAt{14}{284.784}{527.735}{p}%
\NativeMathAt{294.088}{527.735}{11.955}{\cong}%
\TextAt{14}{306.707}{527.735}{o}%
\TextAt{9}{312.552}{527.735}{\Character{47}}%
\TextAt{14}{318.405}{527.735}{p}%
\TextAt{8}{328.001}{527.735}{is}%
\TextAt{8}{339.481}{527.735}{also}%
\TextAt{8}{362.677}{527.735}{uniquely}%
\TextAt{8}{409.867}{527.735}{determined}%
\TextAt{8}{470.71}{527.735}{by}%
\TextAt{13}{486.685}{527.735}{u}%
\TextAt{8}{496.961}{527.735}{and}%
\TextAt{8}{519.44}{527.735}{we}%
\TextAt{8}{79.226}{510.689}{have}%
\TextAt{13}{106.221}{510.689}{b}%
\TextAt{21}{111.198}{508.896}{n}%
\TextAt{9}{120.155}{510.689}{\Character{61}}%
\TextAt{13}{132.581}{510.689}{a}%
\TextAt{21}{138.726}{508.896}{n}%
\TextAt{9}{147.019}{510.689}{\Character{43}}%
\TextAt{13}{158.78}{510.689}{\Character{25}}%
\TextAt{13}{165.834}{510.689}{b}%
\TextAt{21}{170.826}{508.896}{n}%
\TextAt{17}{175.965}{508.896}{\Character{43}1}%
\TextAt{8}{187.283}{510.689}{,}%
\TextAt{8}{194.433}{510.689}{for}%
\TextAt{8}{212.317}{510.689}{some}%
\TextAt{13}{241.643}{510.689}{b}%
\TextAt{21}{246.621}{508.896}{n}%
\TextAt{17}{251.759}{508.896}{\Character{43}1}%
\NativeMathAt{266.398}{510.689}{11.955}{\in}%
\TextAt{20}{277.291}{511.225}{\Character{98}}%
\TextAt{14}{277.689}{510.689}{o}%
\TextAt{8}{283.534}{510.689}{.}%
\TextAt{8}{291.982}{510.689}{Hence}%
\TextAt{13}{177.949}{482.306}{u}%
\TextAt{9}{187.932}{482.306}{\Character{61}}%
\TextAt{13}{200.358}{482.306}{a}%
\TextAt{17}{206.503}{480.513}{0}%
\TextAt{9}{213.892}{482.306}{\Character{43}}%
\TextAt{13}{225.653}{482.306}{a}%
\TextAt{17}{231.798}{480.513}{1}%
\TextAt{13}{236.53}{482.306}{\Character{25}}%
\TextAt{9}{246.256}{482.306}{\Character{43}}%
\TextAt{16}{258.018}{482.306}{\Character{1}}%
\TextAt{16}{263.314}{482.306}{\Character{1}}%
\TextAt{16}{268.622}{482.306}{\Character{1}}%
\TextAt{9}{274.622}{482.306}{\Character{43}}%
\TextAt{13}{286.383}{482.306}{a}%
\TextAt{21}{292.528}{480.513}{n}%
\TextAt{3}{297.666}{480.513}{\Character{0}}%
\TextAt{17}{304.253}{480.513}{1}%
\TextAt{13}{308.985}{482.306}{\Character{25}}%
\TextAt{21}{316.055}{487.242}{n}%
\TextAt{3}{321.193}{487.242}{\Character{0}}%
\TextAt{17}{327.779}{487.242}{1}%
\TextAt{9}{335.168}{482.306}{\Character{43}}%
\TextAt{13}{346.93}{482.306}{a}%
\TextAt{21}{353.075}{480.513}{n}%
\TextAt{13}{358.711}{482.306}{\Character{25}}%
\TextAt{21}{365.78}{487.242}{n}%
\TextAt{9}{374.073}{482.306}{\Character{43}}%
\TextAt{13}{385.834}{482.306}{\Character{25}}%
\TextAt{21}{392.904}{487.242}{n}%
\TextAt{17}{398.042}{487.242}{\Character{43}1}%
\TextAt{13}{409.361}{482.306}{b}%
\TextAt{21}{414.338}{480.513}{n}%
\TextAt{17}{419.476}{480.513}{\Character{43}1}%
\TextAt{13}{430.795}{482.306}{\Character{58}}%
\TextAt{8}{79.226}{453.923}{In}%
\TextAt{8}{94.575}{453.923}{this}%
\TextAt{8}{118.126}{453.923}{way}%
\TextAt{8}{142.596}{453.923}{we}%
\TextAt{8}{160.552}{453.923}{find}%
\TextAt{8}{184.688}{453.923}{an}%
\TextAt{8}{201.675}{453.923}{infinite}%
\TextAt{8}{242.07}{453.923}{series}%
\TextAt{20}{274.132}{462.889}{\Character{80}}%
\NativeMathAt{286.752}{459.734}{7.97}{\infty}%
\TextAt{21}{286.752}{450.436}{\Character{23}}%
\TextAt{17}{291.436}{450.436}{\Character{61}0}%
\TextAt{13}{304.747}{453.923}{a}%
\TextAt{21}{310.892}{452.129}{\Character{23}}%
\TextAt{13}{316.074}{453.923}{\Character{25}}%
\TextAt{21}{323.143}{458.261}{\Character{23}}%
\TextAt{8}{332.95}{453.923}{which}%
\TextAt{8}{367.165}{453.923}{is}%
\TextAt{8}{379.656}{453.923}{uniquely}%
\TextAt{8}{427.858}{453.923}{determined}%
\TextAt{8}{489.713}{453.923}{by}%
\TextAt{13}{506.7}{453.923}{u}%
\TextAt{8}{513.362}{453.923}{.}%
\TextAt{8}{523.994}{453.923}{It}%
\TextAt{8}{79.226}{436.877}{converges}%
\TextAt{8}{131.65}{436.877}{to}%
\TextAt{13}{145.96}{436.877}{u}%
\TextAt{8}{156.52}{436.877}{because}%
\TextAt{8}{199.821}{436.877}{the}%
\TextAt{8}{219.978}{436.877}{remainder}%
\TextAt{8}{275.258}{436.877}{term}%
\TextAt{13}{303.222}{436.877}{\Character{25}}%
\TextAt{21}{310.291}{441.215}{n}%
\TextAt{17}{315.429}{441.215}{\Character{43}1}%
\TextAt{13}{326.748}{436.877}{b}%
\TextAt{21}{331.725}{435.083}{n}%
\TextAt{17}{336.863}{435.083}{\Character{43}1}%
\TextAt{8}{352.079}{436.877}{tends}%
\TextAt{8}{383.354}{436.877}{to}%
\TextAt{8}{397.664}{436.877}{zero.}%
\Vector{q 0 G 0 g 0.399 w 452.334 -337.123 m 452.334 -329.054 l S Q}%
\Vector{q 0 G 0 g 0.399 w 452.534 -329.253 m 459.707 -329.253 l S Q}%
\Vector{q 0 G 0 g 0.399 w 452.534 -336.924 m 459.707 -336.924 l S Q}%
\Vector{q 0 G 0 g 0.399 w 459.906 -337.123 m 459.906 -329.054 l S Q}%
\TextAt{8}{97.159}{408.052}{In}%
\TextAt{8}{112.906}{408.052}{the}%
\TextAt{8}{134.189}{408.052}{case}%
\TextAt{8}{160.086}{408.052}{of}%
\TextAt{8}{174.542}{408.052}{the}%
\TextAt{8}{195.825}{408.052}{field}%
\TextAt{8}{222.308}{408.052}{of}%
\TextAt{8}{236.764}{408.052}{rational}%
\TextAt{8}{281.479}{408.052}{numbers}%
\TextAt{18}{329.816}{408.052}{Q}%
\TextAt{8}{344.138}{408.052}{and}%
\TextAt{8}{368.027}{408.052}{the}%
\TextAt{13}{389.309}{408.052}{p}%
\TextAt{8}{395.185}{408.052}{-adic}%
\TextAt{8}{424.919}{408.052}{valuation}%
\TextAt{13}{477.118}{408.052}{v}%
\TextAt{21}{482.777}{406.259}{p}%
\TextAt{8}{492.562}{408.052}{with}%
\TextAt{8}{520.348}{408.052}{its}%
\TextAt{8}{79.226}{391.006}{completion}%
\TextAt{18}{140.188}{391.006}{Q}%
\TextAt{21}{149.486}{389.213}{p}%
\TextAt{8}{154.247}{391.006}{,}%
\TextAt{8}{162.804}{391.006}{the}%
\TextAt{8}{184.086}{391.006}{numbers}%
\TextAt{9}{232.423}{391.006}{0}%
\TextAt{13}{238.276}{391.006}{\Character{59}}%
\TextAt{9}{243.52}{391.006}{1}%
\TextAt{13}{249.373}{391.006}{\Character{59}}%
\TextAt{13}{254.597}{391.006}{\Character{58}}%
\TextAt{13}{259.834}{391.006}{\Character{58}}%
\TextAt{13}{265.07}{391.006}{\Character{58}}%
\TextAt{13}{270.295}{391.006}{\Character{59}}%
\TextAt{13}{275.531}{391.006}{p}%
\TextAt{16}{284.894}{391.006}{\Character{0}}%
\TextAt{9}{297.616}{391.006}{1}%
\TextAt{8}{308.493}{391.006}{form}%
\TextAt{8}{337.259}{391.006}{a}%
\TextAt{8}{348.14}{391.006}{system}%
\TextAt{8}{388.085}{391.006}{of}%
\TextAt{8}{402.541}{391.006}{representatives}%
\TextAt{13}{483.468}{391.006}{R}%
\TextAt{8}{497.5}{391.006}{for}%
\TextAt{8}{516.51}{391.006}{the}%
\TextAt{8}{79.226}{373.96}{residue}%
\TextAt{8}{119.954}{373.96}{class}%
\TextAt{8}{148.392}{373.96}{field}%
\TextAt{18}{174.75}{373.96}{Z}%
\TextAt{9}{182.721}{373.96}{\Character{47}}%
\TextAt{13}{188.574}{373.96}{p}%
\TextAt{18}{194.449}{373.96}{Z}%
\TextAt{8}{207.317}{373.96}{of}%
\TextAt{8}{221.648}{373.96}{the}%
\TextAt{8}{242.806}{373.96}{valuation,}%
\TextAt{8}{298.382}{373.96}{and}%
\TextAt{8}{322.145}{373.96}{we}%
\TextAt{8}{340.374}{373.96}{get}%
\TextAt{8}{360.886}{373.96}{back}%
\TextAt{8}{389.204}{373.96}{the}%
\TextAt{8}{410.362}{373.96}{representation}%
\TextAt{8}{487.852}{373.96}{of}%
\TextAt{13}{502.183}{373.96}{p}%
\TextAt{8}{508.058}{373.96}{-adic}%
\TextAt{8}{79.226}{356.914}{numbers}%
\TextAt{8}{126.437}{356.914}{which}%
\TextAt{8}{159.924}{356.914}{has}%
\TextAt{8}{180.797}{356.914}{already}%
\TextAt{8}{222.102}{356.914}{been}%
\TextAt{8}{249.731}{356.914}{discussed}%
\TextAt{8}{300.636}{356.914}{in}%
\TextAt{8}{314.289}{356.914}{§2:}%
\TextAt{13}{227.534}{328.531}{x}%
\TextAt{9}{237.507}{328.531}{\Character{61}}%
\TextAt{13}{249.933}{328.531}{p}%
\TextAt{21}{255.808}{333.467}{m}%
\TextAt{9}{263.797}{328.531}{\Character{40}}%
\TextAt{13}{268.349}{328.531}{a}%
\TextAt{17}{274.494}{326.738}{0}%
\TextAt{9}{281.883}{328.531}{\Character{43}}%
\TextAt{13}{293.644}{328.531}{a}%
\TextAt{17}{299.789}{326.738}{1}%
\TextAt{13}{304.521}{328.531}{p}%
\TextAt{9}{313.053}{328.531}{\Character{43}}%
\TextAt{13}{324.814}{328.531}{a}%
\TextAt{17}{330.959}{326.738}{2}%
\TextAt{13}{335.692}{328.531}{p}%
\TextAt{17}{341.567}{333.467}{2}%
\TextAt{9}{348.956}{328.531}{\Character{43}}%
\TextAt{16}{360.717}{328.531}{\Character{1}}%
\TextAt{16}{366.013}{328.531}{\Character{1}}%
\TextAt{16}{371.321}{328.531}{\Character{1}}%
\TextAt{9}{376.657}{328.531}{\Character{41}}%
\TextAt{13}{381.21}{328.531}{\Character{59}}%
\TextAt{8}{79.226}{300.148}{where}%
\TextAt{9}{114.114}{300.148}{0}%
\NativeMathAt{125.127}{300.148}{11.955}{\leq}%
\TextAt{13}{139.585}{300.148}{a}%
\TextAt{21}{145.73}{298.354}{i}%
\TextAt{13}{154.27}{300.148}{\Character{60}}%
\TextAt{13}{168.521}{300.148}{p}%
\TextAt{8}{179.386}{300.148}{and}%
\TextAt{13}{203.228}{300.148}{m}%
\NativeMathAt{218.626}{300.148}{11.955}{\in}%
\TextAt{18}{231.756}{300.148}{Z}%
\TextAt{8}{239.726}{300.148}{.}%
\TextAt{8}{251.411}{300.148}{In}%
\TextAt{8}{267.111}{300.148}{the}%
\TextAt{8}{288.347}{300.148}{case}%
\TextAt{8}{314.197}{300.148}{of}%
\TextAt{8}{328.606}{300.148}{the}%
\TextAt{8}{349.841}{300.148}{rational}%
\TextAt{8}{394.508}{300.148}{function}%
\TextAt{8}{441.435}{300.148}{field}%
\TextAt{13}{467.871}{300.148}{k}%
\TextAt{9}{474.361}{300.148}{\Character{40}}%
\TextAt{13}{478.913}{300.148}{t}%
\TextAt{9}{483.14}{300.148}{\Character{41}}%
\TextAt{8}{492.669}{300.148}{and}%
\TextAt{8}{516.51}{300.148}{the}%
\TextAt{8}{79.226}{283.102}{valuation}%
\TextAt{13}{129.702}{283.102}{v}%
\TextAt{27}{135.362}{281.308}{p}%
\TextAt{8}{143.149}{283.102}{attached}%
\TextAt{8}{190.362}{283.102}{to}%
\TextAt{8}{204.076}{283.102}{a}%
\TextAt{8}{213.235}{283.102}{prime}%
\TextAt{8}{245.802}{283.102}{ideal}%
\TextAt{14}{273.169}{283.102}{p}%
\TextAt{9}{282.473}{283.102}{\Character{61}}%
\TextAt{9}{294.882}{283.102}{\Character{40}}%
\TextAt{13}{299.451}{283.102}{t}%
\TextAt{16}{305.115}{283.102}{\Character{0}}%
\TextAt{13}{315.851}{283.102}{a}%
\TextAt{9}{321.996}{283.102}{\Character{41}}%
\TextAt{8}{329.849}{283.102}{of}%
\TextAt{13}{342.583}{283.102}{k}%
\TextAt{9}{349.072}{283.102}{\Character{91}}%
\TextAt{13}{352.324}{283.102}{t}%
\TextAt{9}{356.551}{283.102}{\Character{93}}%
\TextAt{8}{363.104}{283.102}{(see}%
\ErrataAt{385.97}{283.102}{19.53}{\hyperlink{nav.section.2.3}{\S3}),}%
\TextAt{8}{408.739}{283.102}{we}%
\TextAt{8}{425.371}{283.102}{may}%
\TextAt{8}{450.143}{283.102}{take}%
\TextAt{8}{474.916}{283.102}{as}%
\TextAt{8}{488.689}{283.102}{a}%
\TextAt{8}{497.848}{283.102}{system}%
\TextAt{8}{79.226}{266.056}{of}%
\TextAt{8}{92.404}{266.056}{representatives}%
\TextAt{13}{172.052}{266.056}{R}%
\TextAt{8}{184.805}{266.056}{the}%
\TextAt{8}{204.809}{266.056}{field}%
\TextAt{8}{230.013}{266.056}{of}%
\TextAt{8}{243.19}{266.056}{coefficients}%
\TextAt{13}{302.276}{266.056}{k}%
\TextAt{8}{312.51}{266.056}{itself.}%
\TextAt{8}{345.356}{266.056}{The}%
\TextAt{8}{369.257}{266.056}{completion}%
\TextAt{8}{428.94}{266.056}{then}%
\TextAt{8}{455.447}{266.056}{turns}%
\TextAt{8}{485.924}{266.056}{out}%
\TextAt{8}{506.585}{266.056}{to}%
\TextAt{8}{520.743}{266.056}{be}%
\ErrataAt{79.19}{249.01}{455.81}{the formal Laurent series field $k((x))$, with $x=t-a$, whose elements have the form}%
\TextAt{13}{177.424}{220.626}{f}%
\TextAt{9}{184.47}{220.626}{\Character{40}}%
\TextAt{13}{189.023}{220.626}{t}%
\TextAt{9}{193.25}{220.626}{\Character{41}}%
\TextAt{9}{201.105}{220.626}{\Character{61}}%
\TextAt{9}{213.526}{220.626}{\Character{40}}%
\TextAt{13}{218.101}{220.626}{t}%
\TextAt{16}{224.985}{220.626}{\Character{0}}%
\TextAt{13}{236.94}{220.626}{a}%
\TextAt{9}{243.085}{220.626}{\Character{41}}%
\TextAt{21}{247.637}{225.563}{m}%
\TextAt{20}{255.626}{230.31}{\Character{0}}%
\TextAt{13}{261.105}{220.626}{a}%
\TextAt{17}{267.25}{218.833}{0}%
\TextAt{9}{274.639}{220.626}{\Character{43}}%
\TextAt{13}{286.4}{220.626}{a}%
\TextAt{17}{292.545}{218.833}{1}%
\TextAt{9}{297.278}{220.626}{\Character{40}}%
\TextAt{13}{301.83}{220.626}{t}%
\TextAt{16}{308.714}{220.626}{\Character{0}}%
\TextAt{13}{320.669}{220.626}{a}%
\TextAt{9}{326.814}{220.626}{\Character{41}}%
\TextAt{9}{334.011}{220.626}{\Character{43}}%
\TextAt{13}{345.784}{220.626}{a}%
\TextAt{17}{351.929}{218.833}{2}%
\TextAt{9}{356.662}{220.626}{\Character{40}}%
\TextAt{13}{361.214}{220.626}{t}%
\TextAt{16}{368.098}{220.626}{\Character{0}}%
\TextAt{13}{380.053}{220.626}{a}%
\TextAt{9}{386.198}{220.626}{\Character{41}}%
\TextAt{17}{390.75}{225.563}{2}%
\TextAt{9}{398.139}{220.626}{\Character{43}}%
\TextAt{16}{409.9}{220.626}{\Character{1}}%
\TextAt{16}{415.196}{220.626}{\Character{1}}%
\TextAt{16}{420.504}{220.626}{\Character{1}}%
\TextAt{20}{425.841}{230.31}{\Character{1}}%
\TextAt{13}{431.32}{220.626}{\Character{59}}%
\TextAt{8}{79.226}{192.243}{with}%
\TextAt{13}{107.689}{192.243}{a}%
\TextAt{21}{113.834}{190.45}{i}%
\NativeMathAt{123.608}{192.243}{11.955}{\in}%
\TextAt{13}{137.972}{192.243}{k}%
\TextAt{8}{150.161}{192.243}{and}%
\TextAt{13}{174.727}{192.243}{m}%
\NativeMathAt{191.359}{192.243}{11.955}{\in}%
\TextAt{18}{205.722}{192.243}{Z}%
\TextAt{8}{213.693}{192.243}{.}%
\TextAt{8}{227.549}{192.243}{The}%
\TextAt{8}{253.406}{192.243}{motivating}%
\TextAt{8}{314.088}{192.243}{analogy}%
\TextAt{8}{359.157}{192.243}{of}%
\TextAt{8}{374.289}{192.243}{the}%
\TextAt{8}{396.249}{192.243}{beginning}%
\TextAt{8}{451.706}{192.243}{of}%
\TextAt{8}{466.839}{192.243}{this}%
\TextAt{8}{491.464}{192.243}{chapter,}%
\TextAt{8}{79.226}{175.197}{between}%
\TextAt{8}{125.301}{175.197}{power}%
\TextAt{8}{160.33}{175.197}{series}%
\TextAt{8}{192.549}{175.197}{and}%
\TextAt{13}{216.197}{175.197}{p}%
\TextAt{8}{222.072}{175.197}{-adic}%
\TextAt{8}{251.565}{175.197}{numbers,}%
\TextAt{8}{303.134}{175.197}{thus}%
\TextAt{8}{329.77}{175.197}{appears}%
\TextAt{8}{373.969}{175.197}{as}%
\TextAt{8}{389.223}{175.197}{two}%
\TextAt{8}{412.225}{175.197}{special}%
\TextAt{8}{451.211}{175.197}{instances}%
\TextAt{8}{502.296}{175.197}{of}%
\TextAt{8}{516.51}{175.197}{the}%
\TextAt{8}{79.226}{158.151}{same}%
\TextAt{8}{107.829}{158.151}{concrete}%
\TextAt{8}{153.278}{158.151}{mathematical}%
\TextAt{8}{226.055}{158.151}{situation.}%
\TextAt{8}{279.214}{158.151}{In}%
\TextAt{8}{293.113}{158.151}{§1}%
\TextAt{8}{307.824}{158.151}{we}%
\TextAt{8}{324.329}{158.151}{identified}%
\TextAt{8}{374.655}{158.151}{the}%
\TextAt{8}{394.089}{158.151}{ring}%
\TextAt{18}{417.432}{158.151}{Z}%
\TextAt{21}{425.402}{156.358}{p}%
\TextAt{8}{433.337}{158.151}{of}%
\TextAt{13}{445.945}{158.151}{p}%
\TextAt{8}{451.82}{158.151}{-adic}%
\TextAt{8}{479.706}{158.151}{integers}%
\TextAt{8}{522.297}{158.151}{as}%
\TextAt{8}{79.226}{141.105}{being}%
\TextAt{8}{110.383}{141.105}{the}%
\TextAt{8}{130.159}{141.105}{projective}%
\TextAt{8}{183.756}{141.105}{limit}%
\NativeMathAt{211.338}{141.105}{11.955}{\varprojlim}%
\TextAt{21}{227.597}{135.666}{n}%
\TextAt{18}{235.226}{141.105}{Z}%
\TextAt{9}{243.196}{141.105}{\Character{47}}%
\TextAt{13}{249.049}{141.105}{p}%
\TextAt{21}{254.924}{145.444}{n}%
\TextAt{18}{260.561}{141.105}{Z}%
\TextAt{8}{268.531}{141.105}{.}%
\TextAt{8}{276.852}{141.105}{We}%
\TextAt{8}{296.616}{141.105}{obtain}%
\TextAt{8}{332.663}{141.105}{a}%
\TextAt{8}{342.037}{141.105}{similar}%
\TextAt{8}{380.093}{141.105}{result}%
\TextAt{8}{412.29}{141.105}{in}%
\TextAt{8}{425.562}{141.105}{the}%
\TextAt{8}{445.338}{141.105}{general}%
\TextAt{8}{485.282}{141.105}{setting}%
\TextAt{8}{523.337}{141.105}{of}%
\TextAt{8}{79.226}{124.059}{valuation}%
\TextAt{8}{130.074}{124.059}{theory.}%
\TextAt{8}{170.32}{124.059}{To}%
\TextAt{8}{187.323}{124.059}{explain}%
\TextAt{8}{227.746}{124.059}{this,}%
\TextAt{8}{253.641}{124.059}{let}%
\TextAt{13}{270.321}{124.059}{K}%
\TextAt{8}{284.812}{124.059}{be}%
\TextAt{8}{300.512}{124.059}{complete}%
\TextAt{8}{349.71}{124.059}{with}%
\TextAt{8}{376.146}{124.059}{respect}%
\TextAt{8}{415.971}{124.059}{to}%
\TextAt{8}{430.057}{124.059}{a}%
\TextAt{8}{439.588}{124.059}{discrete}%
\TextAt{8}{482.343}{124.059}{valuation.}%
\TextAt{8}{79.226}{107.013}{Let}%
\TextAt{14}{100.333}{107.013}{o}%
\TextAt{8}{110.212}{107.013}{be}%
\TextAt{8}{126.274}{107.013}{the}%
\TextAt{8}{146.568}{107.013}{valuation}%
\TextAt{8}{197.778}{107.013}{ring}%
\TextAt{8}{221.981}{107.013}{with}%
\TextAt{8}{248.778}{107.013}{the}%
\TextAt{8}{269.072}{107.013}{maximal}%
\TextAt{8}{317.018}{107.013}{ideal}%
\TextAt{14}{345.119}{107.013}{p}%
\TextAt{8}{351.102}{107.013}{.}%
\TextAt{8}{359.963}{107.013}{We}%
\TextAt{8}{380.245}{107.013}{then}%
\TextAt{8}{407.042}{107.013}{have}%
\TextAt{8}{434.174}{107.013}{for}%
\TextAt{8}{452.197}{107.013}{every}%
\TextAt{13}{483.226}{107.013}{n}%
\NativeMathAt{493.769}{107.013}{11.955}{\geq}%
\TextAt{9}{506.623}{107.013}{1}%
\TextAt{8}{516.51}{107.013}{the}%
\TextAt{8}{79.226}{89.967}{canonical}%
\TextAt{8}{130.609}{89.967}{homomorphisms}%
\TextAt{14}{278.466}{72.921}{o}%
\NativeArrowAt{287.631}{72.921}{11.955}{19.26}{\longrightarrow}%
\TextAt{14}{310.213}{72.921}{o}%
\TextAt{9}{316.057}{72.921}{\Character{47}}%
\TextAt{14}{321.91}{72.921}{p}%
\TextAt{21}{327.893}{77.858}{n}%
\TextAt{8}{297.211}{43.033}{116}%
\Destination{78.23}{809.815}{nav.page.116}%
\Destination{79.36}{650.655}{nav.print.127}%
\Destination{79.36}{775.155}{nav.number.2.4.4}%
\Hotlink{319.468}{354.755}{6.858}{11.047}{nav.section.2.2}%
\Hotlink{298.292}{155.955}{6.858}{11.047}{nav.section.2.1}%
\Destination{79.226}{300.102}{errata.82}%
}

% Page 129 (body page 117)
\Page{%
\TextAt{2}{79.226}{796.943}{II.}%
\TextAt{2}{94.073}{796.943}{The}%
\TextAt{2}{116.492}{796.943}{Theory}%
\TextAt{2}{154.401}{796.943}{of}%
\TextAt{2}{166.826}{796.943}{Valuations}%
\Vector{q 0 G 0 g 0.299 w 7.226 18.729 m 460.769 18.729 l S Q}%
\TextAt{8}{79.226}{761.124}{and}%
\TextAt{14}{230.674}{744.078}{o}%
\TextAt{9}{236.518}{744.078}{\Character{47}}%
\TextAt{14}{242.371}{744.078}{p}%
\TextAt{21}{255.674}{750.899}{\Character{21}}%
\TextAt{24}{260.614}{749.792}{1}%
\NativeArrowAt{251.675}{744.078}{11.955}{15.673}{\longleftarrow}%
\TextAt{14}{270.674}{744.078}{o}%
\TextAt{9}{276.518}{744.078}{\Character{47}}%
\TextAt{14}{282.371}{744.078}{p}%
\TextAt{17}{288.354}{749.014}{2}%
\TextAt{21}{300.406}{750.899}{\Character{21}}%
\TextAt{24}{305.346}{749.792}{2}%
\NativeArrowAt{296.407}{744.078}{11.955}{15.673}{\longleftarrow}%
\TextAt{14}{315.405}{744.078}{o}%
\TextAt{9}{321.25}{744.078}{\Character{47}}%
\TextAt{14}{327.103}{744.078}{p}%
\TextAt{17}{333.086}{749.014}{3}%
\TextAt{21}{345.138}{750.899}{\Character{21}}%
\TextAt{24}{350.078}{749.792}{3}%
\NativeArrowAt{341.139}{744.078}{11.955}{15.673}{\longleftarrow}%
\TextAt{16}{360.124}{744.078}{\Character{1}}%
\TextAt{16}{365.42}{744.078}{\Character{1}}%
\TextAt{16}{370.728}{744.078}{\Character{1}}%
\TextAt{13}{378.07}{744.078}{\Character{58}}%
\TextAt{8}{79.226}{722.875}{This}%
\TextAt{8}{105.946}{722.875}{gives}%
\TextAt{8}{134.626}{722.875}{us}%
\TextAt{8}{149.642}{722.875}{a}%
\TextAt{8}{159.398}{722.875}{homomorphism}%
\TextAt{14}{269.34}{705.829}{o}%
\NativeArrowAt{278.505}{705.829}{11.955}{19.26}{\longrightarrow}%
\NativeMathAt{301.087}{705.829}{11.955}{\varprojlim}%
\TextAt{21}{306.648}{693.715}{n}%
\TextAt{14}{319.339}{705.829}{o}%
\TextAt{9}{325.183}{705.829}{\Character{47}}%
\TextAt{14}{331.036}{705.829}{p}%
\TextAt{21}{337.019}{710.766}{n}%
\TextAt{8}{79.226}{679.07}{into}%
\TextAt{8}{102.969}{679.07}{the}%
\TextAt{8}{123.126}{679.07}{projective}%
\TextAt{8}{177.103}{679.07}{limit}%
\NativeMathAt{184.781}{648.17}{11.955}{\varprojlim}%
\TextAt{21}{190.342}{636.056}{n}%
\TextAt{14}{203.033}{648.17}{o}%
\TextAt{9}{208.877}{648.17}{\Character{47}}%
\TextAt{14}{214.73}{648.17}{p}%
\TextAt{21}{220.713}{653.106}{n}%
\TextAt{9}{229.67}{648.17}{\Character{61}}%
\TextAt{20}{242.096}{668.614}{\Character{40}}%
\TextAt{9}{251.727}{648.17}{\Character{40}}%
\TextAt{13}{256.279}{648.17}{x}%
\TextAt{21}{262.931}{646.377}{n}%
\TextAt{9}{268.567}{648.17}{\Character{41}}%
\NativeMathAt{276.44}{648.17}{11.955}{\in}%
\NativeMathAt{291.477}{663.114}{7.97}{\infty}%
\TextAt{20}{288.073}{659.528}{\Character{89}}%
\TextAt{21}{287.731}{634.422}{n}%
\TextAt{17}{292.87}{634.422}{\Character{61}1}%
\TextAt{14}{305.683}{648.17}{o}%
\TextAt{9}{311.527}{648.17}{\Character{47}}%
\TextAt{14}{317.38}{648.17}{p}%
\TextAt{21}{323.363}{653.106}{n}%
\TextAt{20}{332.32}{669.092}{\Character{12}}%
\TextAt{20}{332.32}{661.919}{\Character{12}}%
\TextAt{20}{332.32}{654.745}{\Character{12}}%
\TextAt{20}{332.32}{647.572}{\Character{12}}%
\TextAt{20}{332.32}{640.399}{\Character{12}}%
\TextAt{13}{339.626}{648.17}{\Character{21}}%
\TextAt{21}{346.455}{646.377}{n}%
\TextAt{9}{352.091}{648.17}{\Character{40}}%
\TextAt{13}{356.643}{648.17}{x}%
\TextAt{21}{363.295}{646.377}{n}%
\TextAt{17}{368.434}{646.377}{\Character{43}1}%
\TextAt{9}{379.752}{648.17}{\Character{41}}%
\TextAt{9}{387.607}{648.17}{\Character{61}}%
\TextAt{13}{400.051}{648.17}{x}%
\TextAt{21}{406.703}{646.377}{n}%
\TextAt{20}{412.339}{668.614}{\Character{41}}%
\TextAt{13}{423.963}{648.17}{\Character{58}}%
\TextAt{8}{79.226}{616.72}{Considering}%
\TextAt{8}{144.849}{616.72}{the}%
\TextAt{8}{166.177}{616.72}{rings}%
\TextAt{14}{196.03}{616.72}{o}%
\TextAt{9}{201.874}{616.72}{\Character{47}}%
\TextAt{14}{207.727}{616.72}{p}%
\TextAt{21}{213.71}{621.059}{n}%
\TextAt{8}{224.416}{616.72}{as}%
\TextAt{8}{239.958}{616.72}{topological}%
\TextAt{8}{300.655}{616.72}{rings,}%
\TextAt{8}{334.053}{616.72}{for}%
\TextAt{8}{353.11}{616.72}{the}%
\TextAt{8}{374.439}{616.72}{discrete}%
\TextAt{8}{418.59}{616.72}{topology,}%
\TextAt{8}{470.47}{616.72}{gives}%
\TextAt{8}{500.323}{616.72}{us}%
\TextAt{8}{516.51}{616.72}{the}%
\TextAt{8}{79.226}{599.674}{product}%
\TextAt{8}{122.86}{599.674}{topology}%
\TextAt{8}{170.739}{599.674}{on}%
\TextAt{20}{186.733}{608.641}{\Character{81}}%
\NativeMathAt{198.024}{605.486}{7.97}{\infty}%
\TextAt{21}{198.024}{596.187}{n}%
\TextAt{17}{203.162}{596.187}{\Character{61}1}%
\TextAt{14}{216.473}{599.674}{o}%
\TextAt{9}{222.318}{599.674}{\Character{47}}%
\TextAt{14}{228.171}{599.674}{p}%
\TextAt{21}{234.154}{604.013}{n}%
\TextAt{8}{239.79}{599.674}{,}%
\TextAt{8}{246.727}{599.674}{and}%
\TextAt{8}{269.225}{599.674}{the}%
\TextAt{8}{289.116}{599.674}{projective}%
\TextAt{8}{342.829}{599.674}{limit}%
\NativeMathAt{370.527}{599.674}{11.955}{\varprojlim}%
\TextAt{21}{386.786}{594.235}{n}%
\TextAt{14}{394.414}{599.674}{o}%
\TextAt{9}{400.259}{599.674}{\Character{47}}%
\TextAt{14}{406.112}{599.674}{p}%
\TextAt{21}{412.095}{604.013}{n}%
\TextAt{8}{421.363}{599.674}{becomes}%
\TextAt{8}{467.652}{599.674}{a}%
\TextAt{8}{477.142}{599.674}{topological}%
\TextAt{8}{79.226}{582.628}{ring}%
\TextAt{8}{103.292}{582.628}{in}%
\TextAt{8}{116.945}{582.628}{a}%
\TextAt{8}{126.7}{582.628}{canonical}%
\TextAt{8}{178.083}{582.628}{way,}%
\TextAt{8}{204.098}{582.628}{being}%
\TextAt{8}{235.636}{582.628}{a}%
\TextAt{8}{245.391}{582.628}{closed}%
\TextAt{8}{279.918}{582.628}{subset}%
\TextAt{8}{315.807}{582.628}{of}%
\TextAt{8}{329.137}{582.628}{the}%
\TextAt{8}{349.293}{582.628}{product}%
\TextAt{8}{393.193}{582.628}{(see}%
\TextAt{8}{416.661}{582.628}{chap.}%
\TextAt{8}{447.553}{582.628}{IV,}%
\TextAt{8}{467.697}{582.628}{§2).}%
\TextAt{12}{79.226}{558.162}{Proposition}%
\TextAt{12}{151.937}{558.162}{(4.5).}%
\TextAt{8}{187.253}{558.162}{The}%
\TextAt{8}{211.307}{558.162}{canonical}%
\TextAt{8}{262.69}{558.162}{mapping}%
\TextAt{14}{269.34}{534.573}{o}%
\NativeArrowAt{278.505}{534.573}{11.955}{19.26}{\longrightarrow}%
\NativeMathAt{301.087}{534.573}{11.955}{\varprojlim}%
\TextAt{21}{306.648}{522.459}{n}%
\TextAt{14}{319.339}{534.573}{o}%
\TextAt{9}{325.183}{534.573}{\Character{47}}%
\TextAt{14}{331.036}{534.573}{p}%
\TextAt{21}{337.019}{539.509}{n}%
\TextAt{8}{79.226}{505.307}{is}%
\TextAt{8}{90.99}{505.307}{an}%
\TextAt{8}{107.249}{505.307}{isomorphism}%
\TextAt{8}{175.669}{505.307}{and}%
\TextAt{8}{198.431}{505.307}{a}%
\TextAt{8}{208.187}{505.307}{homeomorphism.}%
\TextAt{8}{300.608}{505.307}{The}%
\TextAt{8}{324.662}{505.307}{same}%
\TextAt{8}{353.988}{505.307}{is}%
\TextAt{8}{365.752}{505.307}{true}%
\TextAt{8}{390.463}{505.307}{for}%
\TextAt{8}{408.348}{505.307}{the}%
\TextAt{8}{428.504}{505.307}{mapping}%
\TextAt{14}{258.081}{481.718}{o}%
\TextAt{3}{263.925}{486.654}{\Character{3}}%
\NativeArrowAt{271.978}{481.718}{11.955}{19.26}{\longrightarrow}%
\NativeMathAt{294.56}{481.718}{11.955}{\varprojlim}%
\TextAt{21}{300.121}{469.604}{n}%
\TextAt{14}{312.812}{481.718}{o}%
\TextAt{3}{318.656}{486.654}{\Character{3}}%
\TextAt{9}{323.388}{481.718}{\Character{47}}%
\TextAt{13}{329.241}{481.718}{U}%
\TextAt{17}{338.44}{486.654}{\Character{40}}%
\TextAt{21}{341.734}{486.654}{n}%
\TextAt{17}{346.872}{486.654}{\Character{41}}%
\TextAt{13}{350.663}{481.718}{\Character{58}}%
\TextAt{19}{79.226}{436.785}{Proof.}%
\TextAt{8}{116}{436.785}{The}%
\TextAt{8}{140.617}{436.785}{map}%
\TextAt{8}{167.194}{436.785}{is}%
\TextAt{8}{179.521}{436.785}{injective}%
\TextAt{8}{226.577}{436.785}{since}%
\TextAt{8}{255.808}{436.785}{its}%
\TextAt{8}{272.69}{436.785}{kernel}%
\TextAt{8}{307.719}{436.785}{is}%
\TextAt{20}{320.046}{445.752}{\Character{84}}%
\NativeMathAt{330.008}{442.597}{7.97}{\infty}%
\TextAt{21}{330.008}{433.298}{n}%
\TextAt{17}{335.147}{433.298}{\Character{61}1}%
\TextAt{14}{348.458}{436.785}{p}%
\TextAt{21}{354.441}{441.124}{n}%
\TextAt{9}{364.357}{436.785}{\Character{61}}%
\TextAt{16}{377.742}{436.785}{\Character{102}}%
\TextAt{9}{383.719}{436.785}{0}%
\TextAt{16}{389.572}{436.785}{\Character{103}}%
\TextAt{8}{395.55}{436.785}{.}%
\TextAt{8}{405.686}{436.785}{To}%
\TextAt{8}{423.476}{436.785}{prove}%
\TextAt{8}{455.589}{436.785}{surjectivity,}%
\TextAt{8}{519.762}{436.785}{let}%
\TextAt{14}{79.226}{419.739}{p}%
\TextAt{9}{88.563}{419.739}{\Character{61}}%
\TextAt{13}{101.022}{419.739}{\Character{25}}%
\TextAt{14}{108.091}{419.739}{o}%
\TextAt{8}{117.852}{419.739}{and}%
\TextAt{8}{140.634}{419.739}{let}%
\TextAt{13}{157.558}{419.739}{R}%
\NativeMathAt{169.921}{419.739}{11.955}{\subseteq}%
\TextAt{14}{182.573}{419.739}{o}%
\TextAt{8}{188.417}{419.739}{,}%
\TextAt{13}{195.591}{419.739}{R}%
\NativeMathAt{207.953}{419.739}{11.955}{\ni}%
\TextAt{9}{219.277}{419.739}{0}%
\TextAt{8}{225.13}{419.739}{,}%
\TextAt{8}{232.304}{419.739}{be}%
\TextAt{8}{248.248}{419.739}{a}%
\TextAt{8}{258.023}{419.739}{system}%
\TextAt{8}{296.86}{419.739}{of}%
\TextAt{8}{310.21}{419.739}{representatives}%
\TextAt{8}{390.03}{419.739}{of}%
\TextAt{14}{403.38}{419.739}{o}%
\TextAt{9}{409.224}{419.739}{\Character{47}}%
\TextAt{14}{415.077}{419.739}{p}%
\TextAt{8}{421.06}{419.739}{.}%
\TextAt{8}{429.566}{419.739}{We}%
\TextAt{8}{449.73}{419.739}{saw}%
\TextAt{8}{472.249}{419.739}{in}%
\TextAt{8}{485.922}{419.739}{the}%
\TextAt{8}{506.098}{419.739}{proof}%
\TextAt{8}{79.226}{402.693}{of}%
\TextAt{8}{93.413}{402.693}{(4.4)}%
\TextAt{8}{122.244}{402.693}{(and}%
\TextAt{8}{150.418}{402.693}{in}%
\TextAt{8}{164.927}{402.693}{fact}%
\TextAt{8}{188.869}{402.693}{already}%
\TextAt{8}{231.03}{402.693}{in}%
\TextAt{8}{245.54}{402.693}{(1.2))}%
\TextAt{8}{278.926}{402.693}{that}%
\TextAt{8}{305.151}{402.693}{the}%
\TextAt{8}{326.164}{402.693}{elements}%
\TextAt{13}{374.877}{402.693}{a}%
\TextAt{8}{384.343}{402.693}{mod}%
\TextAt{14}{409.781}{402.693}{p}%
\TextAt{21}{415.764}{407.032}{n}%
\NativeMathAt{426.18}{402.693}{11.955}{\in}%
\TextAt{14}{438.931}{402.693}{o}%
\TextAt{9}{444.775}{402.693}{\Character{47}}%
\TextAt{14}{450.628}{402.693}{p}%
\TextAt{21}{456.611}{407.032}{n}%
\TextAt{8}{467.001}{402.693}{can}%
\TextAt{8}{489.317}{402.693}{be}%
\TextAt{8}{506.098}{402.693}{given}%
\TextAt{8}{79.226}{385.647}{uniquely}%
\TextAt{8}{126.7}{385.647}{in}%
\TextAt{8}{140.353}{385.647}{the}%
\TextAt{8}{160.509}{385.647}{form}%
\TextAt{13}{197.862}{362.058}{a}%
\NativeMathAt{207.328}{362.058}{11.955}{\equiv}%
\TextAt{13}{219.947}{362.058}{a}%
\TextAt{17}{226.092}{360.265}{0}%
\TextAt{9}{233.481}{362.058}{\Character{43}}%
\TextAt{13}{245.242}{362.058}{a}%
\TextAt{17}{251.387}{360.265}{1}%
\TextAt{13}{256.12}{362.058}{\Character{25}}%
\TextAt{9}{265.846}{362.058}{\Character{43}}%
\TextAt{16}{277.607}{362.058}{\Character{1}}%
\TextAt{16}{282.903}{362.058}{\Character{1}}%
\TextAt{16}{288.211}{362.058}{\Character{1}}%
\TextAt{9}{294.211}{362.058}{\Character{43}}%
\TextAt{13}{305.973}{362.058}{a}%
\TextAt{21}{312.118}{360.265}{n}%
\TextAt{3}{317.256}{360.265}{\Character{0}}%
\TextAt{17}{323.842}{360.265}{1}%
\TextAt{13}{328.575}{362.058}{\Character{25}}%
\TextAt{21}{335.644}{366.995}{n}%
\TextAt{3}{340.782}{366.994}{\Character{0}}%
\TextAt{17}{347.369}{366.995}{1}%
\TextAt{9}{364.056}{362.058}{\Character{40}}%
\TextAt{8}{368.608}{362.058}{mod}%
\TextAt{14}{394.71}{362.058}{p}%
\TextAt{21}{400.693}{366.995}{n}%
\TextAt{9}{406.33}{362.058}{\Character{41}}%
\TextAt{13}{410.882}{362.058}{\Character{59}}%
\TextAt{8}{79.226}{338.469}{where}%
\TextAt{13}{113.035}{338.469}{a}%
\TextAt{21}{119.18}{336.676}{i}%
\NativeMathAt{125.882}{338.469}{11.955}{\in}%
\TextAt{13}{137.173}{338.469}{R}%
\TextAt{8}{146.182}{338.469}{.}%
\TextAt{8}{154.63}{338.469}{Each}%
\TextAt{8}{183.729}{338.469}{element}%
\TextAt{13}{226.971}{338.469}{s}%
\NativeMathAt{235.806}{338.469}{11.955}{\in}%
\NativeMathAt{247.097}{338.469}{11.955}{\varprojlim}%
\TextAt{21}{263.356}{333.03}{n}%
\TextAt{14}{270.985}{338.469}{o}%
\TextAt{9}{276.829}{338.469}{\Character{47}}%
\TextAt{14}{282.682}{338.469}{p}%
\TextAt{21}{288.665}{342.808}{n}%
\TextAt{8}{298.199}{338.469}{is}%
\TextAt{8}{309.963}{338.469}{therefore}%
\TextAt{8}{359.062}{338.469}{given}%
\TextAt{8}{389.632}{338.469}{by}%
\TextAt{8}{405.891}{338.469}{a}%
\TextAt{8}{415.646}{338.469}{sequence}%
\TextAt{8}{464.148}{338.469}{of}%
\TextAt{8}{477.478}{338.469}{sums}%
\TextAt{13}{179.939}{314.88}{s}%
\TextAt{21}{185.453}{313.087}{n}%
\TextAt{9}{194.41}{314.88}{\Character{61}}%
\TextAt{13}{206.836}{314.88}{a}%
\TextAt{17}{212.981}{313.087}{0}%
\TextAt{9}{220.37}{314.88}{\Character{43}}%
\TextAt{13}{232.131}{314.88}{a}%
\TextAt{17}{238.276}{313.087}{1}%
\TextAt{13}{243.008}{314.88}{\Character{25}}%
\TextAt{9}{252.734}{314.88}{\Character{43}}%
\TextAt{16}{264.496}{314.88}{\Character{1}}%
\TextAt{16}{269.792}{314.88}{\Character{1}}%
\TextAt{16}{275.1}{314.88}{\Character{1}}%
\TextAt{9}{281.1}{314.88}{\Character{43}}%
\TextAt{13}{292.861}{314.88}{a}%
\TextAt{21}{299.006}{313.087}{n}%
\TextAt{3}{304.144}{313.087}{\Character{0}}%
\TextAt{17}{310.731}{313.087}{1}%
\TextAt{13}{315.463}{314.88}{\Character{25}}%
\TextAt{21}{322.533}{319.816}{n}%
\TextAt{3}{327.671}{319.816}{\Character{0}}%
\TextAt{17}{334.257}{319.816}{1}%
\TextAt{13}{338.99}{314.88}{\Character{59}}%
\TextAt{13}{368.125}{314.88}{n}%
\TextAt{9}{378.452}{314.88}{\Character{61}}%
\TextAt{9}{390.862}{314.88}{1}%
\TextAt{13}{396.731}{314.88}{\Character{59}}%
\TextAt{9}{401.975}{314.88}{2}%
\TextAt{13}{407.828}{314.88}{\Character{59}}%
\TextAt{13}{413.052}{314.88}{\Character{58}}%
\TextAt{13}{418.289}{314.88}{\Character{58}}%
\TextAt{13}{423.525}{314.88}{\Character{58}}%
\TextAt{13}{428.75}{314.88}{\Character{59}}%
\TextAt{8}{79.226}{291.291}{with}%
\TextAt{8}{105.886}{291.291}{fixed}%
\TextAt{8}{134.172}{291.291}{coefficients}%
\TextAt{13}{193.41}{291.291}{a}%
\TextAt{21}{199.555}{289.498}{i}%
\NativeMathAt{206.257}{291.291}{11.955}{\in}%
\TextAt{13}{217.548}{291.291}{R}%
\TextAt{8}{226.556}{291.291}{,}%
\TextAt{8}{233.706}{291.291}{and}%
\TextAt{8}{256.468}{291.291}{it}%
\TextAt{8}{268.172}{291.291}{is}%
\TextAt{8}{279.936}{291.291}{thus}%
\TextAt{8}{305.688}{291.291}{the}%
\TextAt{8}{325.844}{291.291}{image}%
\TextAt{8}{359.665}{291.291}{of}%
\TextAt{8}{372.995}{291.291}{the}%
\TextAt{8}{393.152}{291.291}{element}%
\TextAt{13}{236.043}{261.73}{x}%
\TextAt{9}{246.016}{261.73}{\Character{61}}%
\TextAt{8}{261.349}{261.73}{lim}%
\TextAt{21}{258.441}{254.556}{n}%
\NativeMathAt{263.58}{254.556}{7.97}{\rightarrow}%
\NativeMathAt{272.044}{254.556}{7.97}{\infty}%
\TextAt{13}{282.509}{261.73}{s}%
\TextAt{21}{288.023}{259.936}{n}%
\TextAt{9}{296.98}{261.73}{\Character{61}}%
\NativeMathAt{313.806}{276.674}{7.97}{\infty}%
\TextAt{20}{309.406}{273.087}{\Character{88}}%
\TextAt{21}{310.288}{247.981}{\Character{23}}%
\TextAt{17}{314.971}{247.981}{\Character{61}0}%
\TextAt{13}{328.667}{261.73}{a}%
\TextAt{21}{334.812}{259.936}{\Character{23}}%
\TextAt{13}{339.994}{261.73}{\Character{25}}%
\TextAt{21}{347.063}{266.666}{\Character{23}}%
\NativeMathAt{355.566}{261.73}{11.955}{\in}%
\TextAt{14}{366.857}{261.73}{o}%
\TextAt{13}{372.701}{261.73}{\Character{58}}%
\TextAt{8}{79.226}{230.003}{The}%
\TextAt{8}{102.24}{230.003}{sets}%
\TextAt{13}{124.081}{230.003}{P}%
\TextAt{21}{131.627}{228.21}{n}%
\TextAt{9}{140.584}{230.003}{\Character{61}}%
\TextAt{20}{153.009}{238.969}{\Character{81}}%
\TextAt{21}{164.3}{226.516}{\Character{23}}%
\TextAt{21}{168.978}{226.516}{\Character{62}n}%
\TextAt{14}{183.199}{230.003}{o}%
\TextAt{9}{189.044}{230.003}{\Character{47}}%
\TextAt{14}{194.897}{230.003}{p}%
\TextAt{21}{200.88}{234.341}{\Character{23}}%
\TextAt{8}{208.919}{230.003}{form}%
\TextAt{8}{235.519}{230.003}{a}%
\TextAt{8}{244.234}{230.003}{basis}%
\TextAt{8}{271.934}{230.003}{of}%
\TextAt{8}{284.223}{230.003}{neighbourhoods}%
\TextAt{8}{367.802}{230.003}{of}%
\TextAt{8}{380.091}{230.003}{the}%
\TextAt{8}{399.207}{230.003}{zero}%
\TextAt{8}{422.878}{230.003}{element}%
\TextAt{8}{465.08}{230.003}{of}%
\TextAt{20}{477.369}{238.969}{\Character{81}}%
\NativeMathAt{488.66}{235.815}{7.97}{\infty}%
\TextAt{21}{488.66}{226.516}{\Character{23}}%
\TextAt{17}{493.344}{226.516}{\Character{61}1}%
\TextAt{14}{506.655}{230.003}{o}%
\TextAt{9}{512.5}{230.003}{\Character{47}}%
\TextAt{14}{518.353}{230.003}{p}%
\TextAt{21}{524.336}{234.341}{\Character{23}}%
\TextAt{8}{529.518}{230.003}{.}%
\TextAt{8}{79.226}{212.957}{Under}%
\TextAt{8}{114.661}{212.957}{the}%
\TextAt{8}{134.818}{212.957}{bijection}%
\TextAt{14}{269.567}{195.911}{o}%
\NativeArrowAt{278.732}{195.911}{11.955}{19.26}{\longrightarrow}%
\NativeMathAt{301.314}{195.911}{11.955}{\varprojlim}%
\TextAt{21}{307.102}{183.797}{\Character{23}}%
\TextAt{14}{319.566}{195.911}{o}%
\TextAt{9}{325.41}{195.911}{\Character{47}}%
\TextAt{14}{331.263}{195.911}{p}%
\TextAt{21}{337.246}{200.847}{\Character{23}}%
\TextAt{8}{79.226}{169.032}{the}%
\TextAt{8}{99.659}{169.032}{basis}%
\TextAt{8}{128.675}{169.032}{of}%
\TextAt{8}{142.281}{169.032}{neighbourhoods}%
\TextAt{14}{227.175}{169.032}{p}%
\TextAt{21}{233.158}{173.371}{n}%
\TextAt{8}{242.968}{169.032}{of}%
\TextAt{8}{256.574}{169.032}{zero}%
\TextAt{8}{281.561}{169.032}{in}%
\TextAt{14}{295.49}{169.032}{o}%
\TextAt{8}{305.508}{169.032}{is}%
\TextAt{8}{317.548}{169.032}{mapped}%
\TextAt{8}{362.368}{169.032}{onto}%
\TextAt{8}{388.994}{169.032}{the}%
\TextAt{8}{409.426}{169.032}{basis}%
\TextAt{8}{438.442}{169.032}{of}%
\TextAt{8}{452.048}{169.032}{neighbourhoods}%
\TextAt{13}{79.226}{151.986}{P}%
\TextAt{21}{86.771}{150.193}{n}%
\TextAt{16}{95.396}{151.986}{\Character{92}}%
\NativeMathAt{106.355}{151.986}{11.955}{\varprojlim}%
\TextAt{21}{122.614}{146.547}{\Character{23}}%
\TextAt{14}{129.788}{151.986}{o}%
\TextAt{9}{135.633}{151.986}{\Character{47}}%
\TextAt{14}{141.486}{151.986}{p}%
\TextAt{21}{147.469}{156.325}{\Character{23}}%
\TextAt{8}{157.035}{151.986}{of}%
\TextAt{8}{170.852}{151.986}{zero}%
\TextAt{8}{196.05}{151.986}{in}%
\NativeMathAt{210.19}{151.986}{11.955}{\varprojlim}%
\TextAt{21}{226.449}{146.547}{\Character{23}}%
\TextAt{14}{233.623}{151.986}{o}%
\TextAt{9}{239.468}{151.986}{\Character{47}}%
\TextAt{14}{245.321}{151.986}{p}%
\TextAt{21}{251.304}{156.325}{\Character{23}}%
\TextAt{8}{256.486}{151.986}{.}%
\TextAt{8}{266.394}{151.986}{Thus}%
\TextAt{8}{296.53}{151.986}{the}%
\TextAt{8}{317.173}{151.986}{bijection}%
\TextAt{8}{365.457}{151.986}{is}%
\TextAt{8}{377.708}{151.986}{a}%
\TextAt{8}{387.95}{151.986}{homeomorphism.}%
\TextAt{8}{481.832}{151.986}{It}%
\TextAt{8}{494.991}{151.986}{induces}%
\TextAt{8}{79.226}{134.94}{an}%
\TextAt{8}{95.485}{134.94}{isomorphism}%
\TextAt{8}{163.905}{134.94}{and}%
\TextAt{8}{186.667}{134.94}{homeomorphism}%
\TextAt{8}{274.538}{134.94}{on}%
\TextAt{8}{290.797}{134.94}{the}%
\TextAt{8}{310.953}{134.94}{group}%
\TextAt{8}{344.129}{134.94}{of}%
\TextAt{8}{357.459}{134.94}{units}%
\TextAt{14}{187.801}{103.475}{o}%
\TextAt{3}{193.645}{108.411}{\Character{3}}%
\NativeMathAt{201.698}{103.475}{11.955}{\cong}%
\TextAt{20}{214.318}{123.918}{\Character{32}}%
\NativeMathAt{223.782}{103.475}{11.955}{\varprojlim}%
\TextAt{21}{229.343}{91.36}{n}%
\TextAt{14}{242.034}{103.475}{o}%
\TextAt{9}{247.878}{103.475}{\Character{47}}%
\TextAt{14}{253.731}{103.475}{p}%
\TextAt{21}{259.714}{108.411}{n}%
\TextAt{20}{265.35}{123.918}{\Character{33}}%
\TextAt{3}{274.815}{121.242}{\Character{3}}%
\NativeMathAt{282.868}{103.475}{11.955}{\cong}%
\NativeMathAt{295.487}{103.475}{11.955}{\varprojlim}%
\TextAt{21}{301.048}{91.36}{n}%
\TextAt{9}{311.746}{103.475}{\Character{40}}%
\TextAt{14}{316.299}{103.475}{o}%
\TextAt{9}{322.143}{103.475}{\Character{47}}%
\TextAt{14}{327.996}{103.475}{p}%
\TextAt{21}{333.979}{108.411}{n}%
\TextAt{9}{339.615}{103.475}{\Character{41}}%
\TextAt{3}{344.168}{108.411}{\Character{3}}%
\NativeMathAt{352.221}{103.475}{11.955}{\cong}%
\NativeMathAt{364.84}{103.475}{11.955}{\varprojlim}%
\TextAt{21}{370.401}{91.36}{n}%
\TextAt{14}{383.092}{103.475}{o}%
\TextAt{3}{388.936}{108.411}{\Character{3}}%
\TextAt{9}{393.668}{103.475}{\Character{47}}%
\TextAt{13}{399.521}{103.475}{U}%
\TextAt{17}{408.72}{108.411}{\Character{40}}%
\TextAt{21}{412.014}{108.411}{n}%
\TextAt{17}{417.152}{108.411}{\Character{41}}%
\TextAt{13}{420.943}{103.475}{\Character{58}}%
\Vector{q 0 G 0 g 0.399 w 452.334 -701.079 m 452.334 -693.009 l S Q}%
\Vector{q 0 G 0 g 0.399 w 452.534 -693.208 m 459.707 -693.208 l S Q}%
\Vector{q 0 G 0 g 0.399 w 452.534 -700.879 m 459.707 -700.879 l S Q}%
\Vector{q 0 G 0 g 0.399 w 459.906 -701.079 m 459.906 -693.009 l S Q}%
\TextAt{8}{297.211}{43.033}{117}%
\Destination{78.23}{809.815}{nav.page.117}%
\Destination{79.36}{737.855}{nav.print.128}%
\Destination{79.36}{572.155}{nav.number.2.4.5}%
\Hotlink{447.053}{580.455}{13.995}{11.047}{nav.chapter.4}%
\Hotlink{472.876}{580.455}{6.858}{11.047}{nav.section.4.2}%
\Hotlink{92.913}{400.455}{25.078}{11.047}{nav.number.2.4.4}%
\Hotlink{245.04}{400.455}{25.078}{11.047}{nav.number.2.1.2}%
}

% Page 130 (body page 118)
\Page{%
\TextAt{2}{79.226}{796.943}{II.}%
\TextAt{2}{94.073}{796.943}{The}%
\TextAt{2}{116.492}{796.943}{Theory}%
\TextAt{2}{154.401}{796.943}{of}%
\TextAt{2}{166.826}{796.943}{Valuations}%
\Vector{q 0 G 0 g 0.299 w 7.226 18.729 m 460.769 18.729 l S Q}%
\TextAt{8}{97.159}{761.124}{One}%
\TextAt{8}{122.393}{761.124}{of}%
\TextAt{8}{136.245}{761.124}{our}%
\TextAt{8}{157.582}{761.124}{chief}%
\TextAt{8}{185.41}{761.124}{concerns}%
\TextAt{8}{233.466}{761.124}{will}%
\TextAt{8}{256.094}{761.124}{be}%
\TextAt{8}{272.541}{761.124}{to}%
\TextAt{8}{287.374}{761.124}{study}%
\TextAt{8}{320.151}{761.124}{the}%
\TextAt{8}{340.83}{761.124}{finite}%
\TextAt{8}{371.265}{761.124}{extensions}%
\TextAt{13}{428.168}{761.124}{L}%
\TextAt{16}{440.344}{761.124}{\Character{106}}%
\TextAt{13}{447.876}{761.124}{K}%
\TextAt{8}{463.114}{761.124}{of}%
\TextAt{8}{476.966}{761.124}{a}%
\TextAt{8}{487.244}{761.124}{complete}%
\TextAt{8}{79.226}{744.078}{valued}%
\TextAt{8}{116.605}{744.078}{field}%
\TextAt{13}{142.591}{744.078}{K}%
\TextAt{8}{153.409}{744.078}{.}%
\TextAt{8}{163.744}{744.078}{This}%
\TextAt{8}{191.093}{744.078}{means}%
\TextAt{8}{227.552}{744.078}{that}%
\TextAt{8}{253.55}{744.078}{we}%
\TextAt{8}{271.406}{744.078}{have}%
\TextAt{8}{299.03}{744.078}{to}%
\TextAt{8}{313.97}{744.078}{turn}%
\TextAt{8}{340.613}{744.078}{to}%
\TextAt{8}{355.552}{744.078}{the}%
\TextAt{8}{376.338}{744.078}{question}%
\TextAt{8}{423.532}{744.078}{of}%
\TextAt{8}{437.491}{744.078}{factoring}%
\TextAt{8}{487.232}{744.078}{algebraic}%
\TextAt{8}{79.226}{727.032}{equations}%
\TextAt{13}{206.994}{709.986}{f}%
\TextAt{9}{214.04}{709.986}{\Character{40}}%
\TextAt{13}{218.592}{709.986}{x}%
\TextAt{9}{225.244}{709.986}{\Character{41}}%
\TextAt{9}{233.099}{709.986}{\Character{61}}%
\TextAt{13}{245.543}{709.986}{a}%
\TextAt{21}{251.688}{708.193}{n}%
\TextAt{13}{257.324}{709.986}{x}%
\TextAt{21}{263.976}{714.922}{n}%
\TextAt{9}{272.269}{709.986}{\Character{43}}%
\TextAt{13}{284.031}{709.986}{a}%
\TextAt{21}{290.176}{708.193}{n}%
\TextAt{3}{295.314}{708.193}{\Character{0}}%
\TextAt{17}{301.9}{708.193}{1}%
\TextAt{13}{306.633}{709.986}{x}%
\TextAt{21}{313.285}{714.922}{n}%
\TextAt{3}{318.423}{714.922}{\Character{0}}%
\TextAt{17}{325.01}{714.922}{1}%
\TextAt{9}{332.399}{709.986}{\Character{43}}%
\TextAt{16}{344.16}{709.986}{\Character{1}}%
\TextAt{16}{349.456}{709.986}{\Character{1}}%
\TextAt{16}{354.764}{709.986}{\Character{1}}%
\TextAt{9}{360.764}{709.986}{\Character{43}}%
\TextAt{13}{372.525}{709.986}{a}%
\TextAt{17}{378.67}{708.193}{0}%
\TextAt{9}{386.724}{709.986}{\Character{61}}%
\TextAt{9}{399.134}{709.986}{0}%
\TextAt{8}{79.226}{686.813}{over}%
\TextAt{8}{105.304}{686.813}{complete}%
\TextAt{8}{155.759}{686.813}{valued}%
\TextAt{8}{193.542}{686.813}{fields.}%
\TextAt{8}{231.161}{686.813}{For}%
\TextAt{8}{253.162}{686.813}{this,}%
\TextAt{8}{280.527}{686.813}{Hensel’s}%
\TextAt{8}{326.869}{686.813}{seminal}%
\TextAt{8}{370.235}{686.813}{“lemma”}%
\TextAt{8}{420.032}{686.813}{is}%
\TextAt{8}{432.828}{686.813}{of}%
\TextAt{8}{447.19}{686.813}{fundamental}%
\TextAt{8}{515.865}{686.813}{im-}%
\TextAt{8}{79.226}{669.767}{portance.}%
\TextAt{8}{134.895}{669.767}{Let}%
\TextAt{13}{156.753}{669.767}{K}%
\TextAt{8}{172.356}{669.767}{again}%
\TextAt{8}{204.471}{669.767}{be}%
\TextAt{8}{221.283}{669.767}{a}%
\TextAt{8}{231.927}{669.767}{field}%
\TextAt{8}{258.172}{669.767}{which}%
\TextAt{8}{292.546}{669.767}{is}%
\TextAt{8}{305.198}{669.767}{complete}%
\TextAt{8}{355.509}{669.767}{with}%
\TextAt{8}{383.057}{669.767}{respect}%
\TextAt{8}{423.995}{669.767}{to}%
\TextAt{8}{439.193}{669.767}{a}%
\TextAt{8}{449.836}{669.767}{nonarchimedean}%
\TextAt{8}{79.226}{652.721}{valuation}%
\TextAt{16}{129.982}{652.721}{\Character{106}}%
\TextAt{16}{136.868}{652.721}{\Character{106}}%
\TextAt{8}{140.205}{652.721}{.}%
\TextAt{8}{148.547}{652.721}{Let}%
\TextAt{14}{169.2}{652.721}{o}%
\TextAt{8}{178.625}{652.721}{be}%
\TextAt{8}{194.233}{652.721}{the}%
\TextAt{8}{214.073}{652.721}{corresponding}%
\TextAt{8}{288.942}{652.721}{valuation}%
\TextAt{8}{339.698}{652.721}{ring}%
\TextAt{8}{363.448}{652.721}{with}%
\TextAt{8}{389.791}{652.721}{maximal}%
\TextAt{8}{437.283}{652.721}{ideal}%
\TextAt{14}{464.93}{652.721}{p}%
\TextAt{8}{474.494}{652.721}{and}%
\TextAt{8}{496.94}{652.721}{residue}%
\TextAt{8}{79.226}{635.675}{class}%
\TextAt{8}{106.706}{635.675}{field}%
\TextAt{13}{132.105}{635.675}{\Character{20}}%
\TextAt{9}{142.237}{635.675}{\Character{61}}%
\TextAt{14}{154.735}{635.675}{o}%
\TextAt{9}{160.58}{635.675}{\Character{47}}%
\TextAt{14}{166.433}{635.675}{p}%
\TextAt{8}{172.416}{635.675}{.}%
\TextAt{8}{180.992}{635.675}{We}%
\TextAt{8}{201.179}{635.675}{call}%
\TextAt{8}{222.681}{635.675}{a}%
\TextAt{8}{232.479}{635.675}{polynomial}%
\TextAt{13}{293.015}{635.675}{f}%
\TextAt{9}{300.061}{635.675}{\Character{40}}%
\TextAt{13}{304.614}{635.675}{x}%
\TextAt{9}{311.266}{635.675}{\Character{41}}%
\TextAt{9}{319.192}{635.675}{\Character{61}}%
\TextAt{13}{331.709}{635.675}{a}%
\TextAt{17}{337.854}{633.882}{0}%
\TextAt{9}{345.272}{635.675}{\Character{43}}%
\TextAt{13}{357.063}{635.675}{a}%
\TextAt{17}{363.208}{633.882}{1}%
\TextAt{13}{367.94}{635.675}{x}%
\TextAt{9}{377.278}{635.675}{\Character{43}}%
\TextAt{16}{389.068}{635.675}{\Character{1}}%
\TextAt{16}{394.364}{635.675}{\Character{1}}%
\TextAt{16}{399.672}{635.675}{\Character{1}}%
\TextAt{9}{405.702}{635.675}{\Character{43}}%
\TextAt{13}{417.492}{635.675}{a}%
\TextAt{21}{423.637}{633.882}{n}%
\TextAt{13}{429.273}{635.675}{x}%
\TextAt{21}{435.925}{640.014}{n}%
\NativeMathAt{444.955}{635.675}{11.955}{\in}%
\TextAt{14}{456.319}{635.675}{o}%
\TextAt{9}{462.163}{635.675}{\Character{91}}%
\TextAt{13}{465.414}{635.675}{x}%
\TextAt{9}{472.067}{635.675}{\Character{93}}%
\TextAt{12}{479.258}{635.675}{primitive}%
\TextAt{8}{79.226}{618.629}{if}%
\TextAt{13}{89.95}{618.629}{f}%
\TextAt{9}{96.996}{618.629}{\Character{40}}%
\TextAt{13}{101.549}{618.629}{x}%
\TextAt{9}{108.201}{618.629}{\Character{41}}%
\NativeMathAt{116.074}{618.629}{11.955}{\not\equiv}%
\TextAt{9}{128.693}{618.629}{0}%
\TextAt{8}{137.867}{618.629}{mod}%
\TextAt{14}{163.305}{618.629}{p}%
\TextAt{8}{169.288}{618.629}{,}%
\TextAt{8}{176.437}{618.629}{i.e.,}%
\TextAt{8}{198.542}{618.629}{if}%
\TextAt{16}{230.899}{601.583}{\Character{106}}%
\TextAt{13}{234.22}{601.583}{f}%
\TextAt{16}{241.266}{601.583}{\Character{106}}%
\TextAt{9}{247.908}{601.583}{\Character{61}}%
\TextAt{8}{260.333}{601.583}{max}%
\TextAt{16}{282.128}{601.583}{\Character{102}\Character{106}}%
\TextAt{13}{291.426}{601.583}{a}%
\TextAt{17}{297.571}{599.79}{0}%
\TextAt{16}{302.303}{601.583}{\Character{106}}%
\TextAt{13}{305.624}{601.583}{\Character{59}}%
\TextAt{13}{310.848}{601.583}{\Character{58}}%
\TextAt{13}{316.085}{601.583}{\Character{58}}%
\TextAt{13}{321.321}{601.583}{\Character{58}}%
\TextAt{13}{326.546}{601.583}{\Character{59}}%
\TextAt{16}{331.845}{601.583}{\Character{106}}%
\TextAt{13}{335.166}{601.583}{a}%
\TextAt{21}{341.311}{599.79}{n}%
\TextAt{16}{346.947}{601.583}{\Character{106}\Character{103}}%
\TextAt{9}{359.567}{601.583}{\Character{61}}%
\TextAt{9}{371.976}{601.583}{1}%
\TextAt{13}{377.845}{601.583}{\Character{58}}%
\TextAt{12}{79.226}{569.677}{Hensel’s}%
\TextAt{12}{132.568}{569.677}{Lemma}%
\TextAt{12}{180.734}{569.677}{(4.6).}%
\TextAt{8}{216.485}{569.677}{If}%
\TextAt{8}{228.612}{569.677}{a}%
\TextAt{8}{238.803}{569.677}{primitive}%
\TextAt{8}{289.318}{569.677}{polynomial}%
\TextAt{13}{350.246}{569.677}{f}%
\TextAt{9}{357.293}{569.677}{\Character{40}}%
\TextAt{13}{361.845}{569.677}{x}%
\TextAt{9}{368.497}{569.677}{\Character{41}}%
\NativeMathAt{377.112}{569.677}{11.955}{\in}%
\TextAt{14}{389.144}{569.677}{o}%
\TextAt{9}{394.988}{569.677}{\Character{91}}%
\TextAt{13}{398.24}{569.677}{x}%
\TextAt{9}{404.892}{569.677}{\Character{93}}%
\TextAt{8}{412.476}{569.677}{admits}%
\TextAt{8}{451.347}{569.677}{modulo}%
\TextAt{14}{493.733}{569.677}{p}%
\TextAt{8}{504.048}{569.677}{a}%
\TextAt{8}{514.239}{569.677}{fac-}%
\TextAt{8}{79.226}{552.631}{torization}%
\TextAt{13}{237.911}{535.585}{f}%
\TextAt{9}{244.957}{535.585}{\Character{40}}%
\TextAt{13}{249.509}{535.585}{x}%
\TextAt{9}{256.161}{535.585}{\Character{41}}%
\NativeMathAt{264.035}{535.585}{11.955}{\equiv}%
\TextAt{9}{277.07}{535.585}{\Character{22}}%
\TextAt{13}{276.654}{535.585}{g}%
\TextAt{9}{282.688}{535.585}{\Character{40}}%
\TextAt{13}{287.241}{535.585}{x}%
\TextAt{9}{293.893}{535.585}{\Character{41}}%
\TextAt{9}{298.563}{538.74}{\Character{22}}%
\TextAt{13}{298.445}{535.585}{h}%
\TextAt{9}{305.184}{535.585}{\Character{40}}%
\TextAt{13}{309.736}{535.585}{x}%
\TextAt{9}{316.388}{535.585}{\Character{41}}%
\TextAt{9}{332.874}{535.585}{\Character{40}}%
\TextAt{8}{337.448}{535.585}{mod}%
\TextAt{14}{363.55}{535.585}{p}%
\TextAt{9}{369.533}{535.585}{\Character{41}}%
\TextAt{8}{79.226}{512.413}{into}%
\TextAt{8}{102.969}{512.413}{relatively}%
\TextAt{8}{154.03}{512.413}{prime}%
\TextAt{8}{187.193}{512.413}{polynomials}%
\TextAt{9}{252.717}{512.413}{\Character{22}}%
\TextAt{13}{252.301}{512.413}{g}%
\TextAt{13}{258.315}{512.413}{\Character{59}}%
\TextAt{9}{263.697}{515.567}{\Character{22}}%
\TextAt{13}{263.58}{512.413}{h}%
\NativeMathAt{273.639}{512.413}{11.955}{\in}%
\TextAt{13}{284.93}{512.413}{\Character{20}}%
\TextAt{9}{291.668}{512.413}{\Character{91}}%
\TextAt{13}{294.92}{512.413}{x}%
\TextAt{9}{301.572}{512.413}{\Character{93}}%
\TextAt{8}{304.824}{512.413}{,}%
\TextAt{8}{311.973}{512.413}{then}%
\TextAt{13}{338.633}{512.413}{f}%
\TextAt{9}{345.68}{512.413}{\Character{40}}%
\TextAt{13}{350.232}{512.413}{x}%
\TextAt{9}{356.884}{512.413}{\Character{41}}%
\TextAt{8}{365.334}{512.413}{admits}%
\TextAt{8}{403.77}{512.413}{a}%
\TextAt{8}{413.525}{512.413}{factorization}%
\TextAt{13}{264.58}{484.226}{f}%
\TextAt{9}{271.626}{484.226}{\Character{40}}%
\TextAt{13}{276.179}{484.226}{x}%
\TextAt{9}{282.831}{484.226}{\Character{41}}%
\TextAt{9}{290.686}{484.226}{\Character{61}}%
\TextAt{13}{303.129}{484.226}{g}%
\TextAt{9}{309.164}{484.226}{\Character{40}}%
\TextAt{13}{313.716}{484.226}{x}%
\TextAt{9}{320.368}{484.226}{\Character{41}}%
\TextAt{13}{324.92}{484.226}{h}%
\TextAt{9}{331.659}{484.226}{\Character{40}}%
\TextAt{13}{336.211}{484.226}{x}%
\TextAt{9}{342.863}{484.226}{\Character{41}}%
\TextAt{8}{79.226}{456.04}{into}%
\TextAt{8}{102.969}{456.04}{polynomials}%
\TextAt{13}{168.077}{456.04}{g}%
\TextAt{13}{174.091}{456.04}{\Character{59}}%
\TextAt{13}{179.327}{456.04}{h}%
\NativeMathAt{189.415}{456.04}{11.955}{\in}%
\TextAt{14}{200.706}{456.04}{o}%
\TextAt{9}{206.55}{456.04}{\Character{91}}%
\TextAt{13}{209.802}{456.04}{x}%
\TextAt{9}{216.454}{456.04}{\Character{93}}%
\TextAt{8}{223.603}{456.04}{such}%
\TextAt{8}{250}{456.04}{that}%
\TextAt{8}{275.369}{456.04}{deg}%
\TextAt{9}{292.931}{456.04}{\Character{40}}%
\TextAt{13}{297.483}{456.04}{g}%
\TextAt{9}{303.517}{456.04}{\Character{41}}%
\TextAt{9}{311.372}{456.04}{\Character{61}}%
\TextAt{8}{323.816}{456.04}{deg}%
\TextAt{9}{341.378}{456.04}{\Character{40}}%
\TextAt{9}{346.327}{456.04}{\Character{22}}%
\TextAt{13}{345.93}{456.04}{g}%
\TextAt{9}{351.965}{456.04}{\Character{41}}%
\TextAt{8}{360.414}{456.04}{and}%
\TextAt{13}{157.658}{427.853}{g}%
\TextAt{9}{163.692}{427.853}{\Character{40}}%
\TextAt{13}{168.244}{427.853}{x}%
\TextAt{9}{174.896}{427.853}{\Character{41}}%
\NativeMathAt{182.77}{427.853}{11.955}{\equiv}%
\TextAt{9}{195.805}{427.853}{\Character{22}}%
\TextAt{13}{195.389}{427.853}{g}%
\TextAt{9}{201.423}{427.853}{\Character{40}}%
\TextAt{13}{205.976}{427.853}{x}%
\TextAt{9}{212.628}{427.853}{\Character{41}}%
\TextAt{9}{229.114}{427.853}{\Character{40}}%
\TextAt{8}{233.687}{427.853}{mod}%
\TextAt{14}{259.789}{427.853}{p}%
\TextAt{9}{265.772}{427.853}{\Character{41}}%
\TextAt{8}{294.235}{427.853}{and}%
\TextAt{13}{337.011}{427.853}{h}%
\TextAt{9}{343.749}{427.853}{\Character{40}}%
\TextAt{13}{348.302}{427.853}{x}%
\TextAt{9}{354.954}{427.853}{\Character{41}}%
\NativeMathAt{362.827}{427.853}{11.955}{\equiv}%
\TextAt{9}{375.564}{431.008}{\Character{22}}%
\TextAt{13}{375.446}{427.853}{h}%
\TextAt{9}{382.185}{427.853}{\Character{40}}%
\TextAt{13}{386.737}{427.853}{x}%
\TextAt{9}{393.389}{427.853}{\Character{41}}%
\TextAt{9}{409.875}{427.853}{\Character{40}}%
\TextAt{8}{414.449}{427.853}{mod}%
\TextAt{14}{440.551}{427.853}{p}%
\TextAt{9}{446.534}{427.853}{\Character{41}}%
\TextAt{13}{451.086}{427.853}{\Character{58}}%
\TextAt{19}{79.226}{382.199}{Proof.}%
\TextAt{8}{116}{382.199}{Let}%
\TextAt{13}{137.698}{382.199}{d}%
\TextAt{9}{148.343}{382.199}{\Character{61}}%
\TextAt{8}{162.009}{382.199}{deg}%
\TextAt{9}{179.571}{382.199}{\Character{40}}%
\TextAt{13}{184.124}{382.199}{f}%
\TextAt{9}{191.17}{382.199}{\Character{41}}%
\TextAt{8}{195.722}{382.199}{,}%
\TextAt{13}{203.782}{382.199}{m}%
\TextAt{9}{218.583}{382.199}{\Character{61}}%
\TextAt{8}{232.25}{382.199}{deg}%
\TextAt{9}{249.812}{382.199}{\Character{40}}%
\TextAt{9}{254.762}{382.199}{\Character{22}}%
\TextAt{13}{254.364}{382.199}{g}%
\TextAt{9}{260.399}{382.199}{\Character{41}}%
\TextAt{8}{264.951}{382.199}{,}%
\TextAt{8}{273.011}{382.199}{hence}%
\TextAt{13}{306.245}{382.199}{d}%
\TextAt{16}{315.481}{382.199}{\Character{0}}%
\TextAt{13}{327.932}{382.199}{m}%
\NativeMathAt{342.734}{382.199}{11.955}{\geq}%
\TextAt{8}{356.594}{382.199}{deg}%
\TextAt{9}{374.156}{382.199}{\Character{40}}%
\TextAt{9}{378.826}{385.354}{\Character{22}}%
\TextAt{13}{378.709}{382.199}{h}%
\TextAt{9}{385.447}{382.199}{\Character{41}}%
\TextAt{8}{389.999}{382.199}{.}%
\TextAt{8}{400.633}{382.199}{Let}%
\TextAt{13}{422.33}{382.199}{g}%
\TextAt{17}{427.936}{380.406}{0}%
\TextAt{13}{432.668}{382.199}{\Character{59}}%
\TextAt{13}{437.892}{382.199}{h}%
\TextAt{17}{444.651}{380.406}{0}%
\NativeMathAt{453.945}{382.199}{11.955}{\in}%
\TextAt{14}{466.477}{382.199}{o}%
\TextAt{9}{472.321}{382.199}{\Character{91}}%
\TextAt{13}{475.573}{382.199}{x}%
\TextAt{9}{482.225}{382.199}{\Character{93}}%
\TextAt{8}{490.102}{382.199}{be}%
\TextAt{8}{506.755}{382.199}{poly-}%
\TextAt{8}{79.226}{365.153}{nomials}%
\TextAt{8}{123.578}{365.153}{such}%
\TextAt{8}{151.337}{365.153}{that}%
\TextAt{13}{178.067}{365.153}{g}%
\TextAt{17}{183.673}{363.36}{0}%
\NativeMathAt{194.046}{365.153}{11.955}{\equiv}%
\TextAt{9}{209.401}{365.153}{\Character{22}}%
\TextAt{13}{208.986}{365.153}{g}%
\TextAt{8}{218.341}{365.153}{mod}%
\TextAt{14}{243.779}{365.153}{p}%
\TextAt{8}{249.762}{365.153}{,}%
\TextAt{13}{258.613}{365.153}{h}%
\TextAt{17}{265.351}{363.36}{0}%
\NativeMathAt{275.725}{365.153}{11.955}{\equiv}%
\TextAt{9}{290.782}{368.308}{\Character{22}}%
\TextAt{13}{290.664}{365.153}{h}%
\TextAt{8}{300.724}{365.153}{mod}%
\TextAt{14}{326.162}{365.153}{p}%
\TextAt{8}{337.404}{365.153}{and}%
\TextAt{8}{361.528}{365.153}{deg}%
\TextAt{9}{379.09}{365.153}{\Character{40}}%
\TextAt{13}{383.642}{365.153}{g}%
\TextAt{17}{389.248}{363.36}{0}%
\TextAt{9}{393.98}{365.153}{\Character{41}}%
\TextAt{9}{404.154}{365.153}{\Character{61}}%
\TextAt{13}{418.919}{365.153}{m}%
\TextAt{8}{429.158}{365.153}{,}%
\TextAt{8}{438.009}{365.153}{deg}%
\TextAt{9}{455.572}{365.153}{\Character{40}}%
\TextAt{13}{460.124}{365.153}{h}%
\TextAt{17}{466.862}{363.36}{0}%
\TextAt{9}{471.595}{365.153}{\Character{41}}%
\NativeMathAt{481.788}{365.153}{11.955}{\leq}%
\TextAt{13}{496.728}{365.153}{d}%
\TextAt{16}{506.395}{365.153}{\Character{0}}%
\TextAt{13}{519.278}{365.153}{m}%
\TextAt{8}{529.518}{365.153}{.}%
\TextAt{8}{79.226}{348.107}{Since}%
\TextAt{9}{109.757}{348.107}{\Character{40}}%
\TextAt{9}{114.707}{348.107}{\Character{22}}%
\TextAt{13}{114.309}{348.107}{g}%
\TextAt{13}{120.323}{348.107}{\Character{59}}%
\TextAt{9}{125.705}{351.262}{\Character{22}}%
\TextAt{13}{125.588}{348.107}{h}%
\TextAt{9}{132.326}{348.107}{\Character{41}}%
\TextAt{9}{140.181}{348.107}{\Character{61}}%
\TextAt{9}{152.602}{348.107}{1}%
\TextAt{8}{158.478}{348.107}{,}%
\TextAt{8}{165.606}{348.107}{there}%
\TextAt{8}{195.491}{348.107}{exist}%
\TextAt{8}{223.165}{348.107}{polynomials}%
\TextAt{13}{288.246}{348.107}{a}%
\TextAt{9}{294.391}{348.107}{\Character{40}}%
\TextAt{13}{298.944}{348.107}{x}%
\TextAt{9}{305.596}{348.107}{\Character{41}}%
\TextAt{13}{310.148}{348.107}{\Character{59}}%
\TextAt{13}{315.372}{348.107}{b}%
\TextAt{9}{320.369}{348.107}{\Character{40}}%
\TextAt{13}{324.922}{348.107}{x}%
\TextAt{9}{331.574}{348.107}{\Character{41}}%
\NativeMathAt{339.447}{348.107}{11.955}{\in}%
\TextAt{14}{350.738}{348.107}{o}%
\TextAt{9}{356.582}{348.107}{\Character{91}}%
\TextAt{13}{359.834}{348.107}{x}%
\TextAt{9}{366.486}{348.107}{\Character{93}}%
\TextAt{8}{373.609}{348.107}{satisfying}%
\TextAt{13}{425.742}{348.107}{ag}%
\TextAt{17}{437.493}{346.314}{0}%
\TextAt{9}{444.828}{348.107}{\Character{43}}%
\TextAt{13}{456.535}{348.107}{bh}%
\TextAt{17}{468.25}{346.314}{0}%
\NativeMathAt{476.304}{348.107}{11.955}{\equiv}%
\TextAt{9}{488.923}{348.107}{1}%
\TextAt{8}{498.097}{348.107}{mod}%
\TextAt{14}{523.535}{348.107}{p}%
\TextAt{8}{529.518}{348.107}{.}%
\ErrataAt{79.19}{331.061}{455.81}{If $f=g_0h_0$, we are done. Otherwise choose $\pi$ of maximal absolute value among the}%
\ErrataAt{79.19}{314.015}{455.81}{coefficients of $f-g_0h_0$ and $ag_0+bh_0-1$. Then $0<|\pi|<1$. We seek $g,h$ in the}%
\TextAt{8}{79.226}{296.969}{following}%
\TextAt{8}{128.661}{296.969}{form:}%
\TextAt{13}{238.029}{281.59}{g}%
\TextAt{9}{247.384}{281.59}{\Character{61}}%
\TextAt{13}{259.809}{281.59}{g}%
\TextAt{17}{265.415}{279.797}{0}%
\TextAt{9}{272.804}{281.59}{\Character{43}}%
\TextAt{13}{284.565}{281.59}{p}%
\TextAt{17}{290.44}{279.797}{1}%
\TextAt{13}{295.172}{281.59}{\Character{25}}%
\TextAt{9}{304.898}{281.59}{\Character{43}}%
\TextAt{13}{316.66}{281.59}{p}%
\TextAt{17}{322.535}{279.797}{2}%
\TextAt{13}{327.267}{281.59}{\Character{25}}%
\TextAt{17}{334.336}{286.526}{2}%
\TextAt{9}{341.725}{281.59}{\Character{43}}%
\TextAt{16}{353.487}{281.59}{\Character{1}}%
\TextAt{16}{358.783}{281.59}{\Character{1}}%
\TextAt{16}{364.091}{281.59}{\Character{1}}%
\TextAt{13}{371.419}{281.59}{\Character{59}}%
\TextAt{13}{237.325}{261.555}{h}%
\TextAt{9}{247.384}{261.555}{\Character{61}}%
\TextAt{13}{259.809}{261.555}{h}%
\TextAt{17}{266.548}{259.762}{0}%
\TextAt{9}{273.937}{261.555}{\Character{43}}%
\TextAt{13}{285.698}{261.555}{q}%
\TextAt{17}{290.889}{259.762}{1}%
\TextAt{13}{295.621}{261.555}{\Character{25}}%
\TextAt{9}{305.347}{261.555}{\Character{43}}%
\TextAt{13}{317.108}{261.555}{q}%
\TextAt{17}{322.298}{259.762}{2}%
\TextAt{13}{327.031}{261.555}{\Character{25}}%
\TextAt{17}{334.1}{266.492}{2}%
\TextAt{9}{341.489}{261.555}{\Character{43}}%
\TextAt{16}{353.25}{261.555}{\Character{1}}%
\TextAt{16}{358.546}{261.555}{\Character{1}}%
\TextAt{16}{363.854}{261.555}{\Character{1}}%
\TextAt{13}{371.183}{261.555}{\Character{59}}%
\TextAt{8}{79.226}{240.352}{where}%
\TextAt{13}{113.723}{240.352}{p}%
\TextAt{21}{119.599}{238.559}{i}%
\TextAt{13}{122.98}{240.352}{\Character{59}}%
\TextAt{13}{128.204}{240.352}{q}%
\TextAt{21}{133.414}{238.559}{i}%
\NativeMathAt{141.289}{240.352}{11.955}{\in}%
\TextAt{14}{153.752}{240.352}{o}%
\TextAt{9}{159.597}{240.352}{\Character{91}}%
\TextAt{13}{162.848}{240.352}{x}%
\TextAt{9}{169.5}{240.352}{\Character{93}}%
\TextAt{8}{177.338}{240.352}{are}%
\TextAt{8}{197.536}{240.352}{polynomials}%
\TextAt{8}{263.332}{240.352}{of}%
\TextAt{8}{277.35}{240.352}{degree}%
\TextAt{13}{314.454}{240.352}{\Character{60}}%
\TextAt{13}{328.035}{240.352}{m}%
\TextAt{8}{338.291}{240.352}{,}%
\TextAt{8}{346.301}{240.352}{resp.}%
\NativeMathAt{377.687}{240.352}{11.955}{\leq}%
\TextAt{13}{391.479}{240.352}{d}%
\TextAt{16}{400.687}{240.352}{\Character{0}}%
\TextAt{13}{413.111}{240.352}{m}%
\TextAt{8}{423.351}{240.352}{.}%
\TextAt{8}{433.863}{240.352}{We}%
\TextAt{8}{454.696}{240.352}{then}%
\TextAt{8}{482.044}{240.352}{determine}%
\TextAt{8}{79.226}{223.306}{successively}%
\TextAt{8}{142.816}{223.306}{the}%
\TextAt{8}{162.972}{223.306}{polynomials}%
\TextAt{13}{220.497}{196.798}{g}%
\TextAt{21}{226.102}{195.005}{n}%
\TextAt{3}{231.241}{195.005}{\Character{0}}%
\TextAt{17}{237.827}{195.005}{1}%
\TextAt{9}{245.88}{196.798}{\Character{61}}%
\TextAt{13}{258.306}{196.798}{g}%
\TextAt{17}{263.911}{195.005}{0}%
\TextAt{9}{271.3}{196.798}{\Character{43}}%
\TextAt{13}{283.061}{196.798}{p}%
\TextAt{17}{288.936}{195.005}{1}%
\TextAt{13}{293.669}{196.798}{\Character{25}}%
\TextAt{9}{303.395}{196.798}{\Character{43}}%
\TextAt{16}{315.156}{196.798}{\Character{1}}%
\TextAt{16}{320.452}{196.798}{\Character{1}}%
\TextAt{16}{325.76}{196.798}{\Character{1}}%
\TextAt{9}{331.76}{196.798}{\Character{43}}%
\TextAt{13}{343.522}{196.798}{p}%
\TextAt{21}{349.397}{195.005}{n}%
\TextAt{3}{354.535}{195.005}{\Character{0}}%
\TextAt{17}{361.122}{195.005}{1}%
\TextAt{13}{365.854}{196.798}{\Character{25}}%
\TextAt{21}{372.923}{201.735}{n}%
\TextAt{3}{378.061}{201.735}{\Character{0}}%
\TextAt{17}{384.648}{201.735}{1}%
\TextAt{13}{389.38}{196.798}{\Character{59}}%
\TextAt{13}{219.364}{176.764}{h}%
\TextAt{21}{226.102}{174.97}{n}%
\TextAt{3}{231.241}{174.97}{\Character{0}}%
\TextAt{17}{237.827}{174.97}{1}%
\TextAt{9}{245.88}{176.764}{\Character{61}}%
\TextAt{13}{258.306}{176.764}{h}%
\TextAt{17}{265.044}{174.97}{0}%
\TextAt{9}{272.433}{176.764}{\Character{43}}%
\TextAt{13}{284.195}{176.764}{q}%
\TextAt{17}{289.385}{174.97}{1}%
\TextAt{13}{294.117}{176.764}{\Character{25}}%
\TextAt{9}{303.843}{176.764}{\Character{43}}%
\TextAt{16}{315.604}{176.764}{\Character{1}}%
\TextAt{16}{320.9}{176.764}{\Character{1}}%
\TextAt{16}{326.208}{176.764}{\Character{1}}%
\TextAt{9}{332.209}{176.764}{\Character{43}}%
\TextAt{13}{343.97}{176.764}{q}%
\TextAt{21}{349.16}{174.97}{n}%
\TextAt{3}{354.298}{174.97}{\Character{0}}%
\TextAt{17}{360.885}{174.97}{1}%
\TextAt{13}{365.617}{176.764}{\Character{25}}%
\TextAt{21}{372.687}{181.7}{n}%
\TextAt{3}{377.825}{181.7}{\Character{0}}%
\TextAt{17}{384.411}{181.7}{1}%
\TextAt{13}{389.144}{176.764}{\Character{59}}%
\TextAt{8}{79.226}{151.216}{in}%
\TextAt{8}{92.879}{151.216}{such}%
\TextAt{8}{119.276}{151.216}{a}%
\TextAt{8}{129.031}{151.216}{way}%
\TextAt{8}{152.774}{151.216}{that}%
\TextAt{8}{178.143}{151.216}{one}%
\TextAt{8}{199.603}{151.216}{has}%
\TextAt{13}{240.316}{123.03}{f}%
\NativeMathAt{250.683}{123.03}{11.955}{\equiv}%
\TextAt{13}{263.303}{123.03}{g}%
\TextAt{21}{268.908}{121.236}{n}%
\TextAt{3}{274.046}{121.236}{\Character{0}}%
\TextAt{17}{280.633}{121.236}{1}%
\TextAt{13}{285.365}{123.03}{h}%
\TextAt{21}{292.104}{121.236}{n}%
\TextAt{3}{297.242}{121.236}{\Character{0}}%
\TextAt{17}{303.828}{121.236}{1}%
\TextAt{9}{320.516}{123.03}{\Character{40}}%
\TextAt{8}{325.068}{123.03}{mod}%
\TextAt{13}{351.17}{123.03}{\Character{25}}%
\TextAt{21}{358.239}{127.966}{n}%
\TextAt{9}{363.876}{123.03}{\Character{41}}%
\TextAt{13}{368.428}{123.03}{\Character{58}}%
\TextAt{8}{297.211}{43.033}{118}%
\Destination{78.23}{809.815}{nav.page.118}%
\Destination{79.36}{776.155}{nav.print.129}%
\Destination{79.226}{348.061}{errata.76}%
\Destination{79.226}{577}{nav.number.2.4.6}%
}

% Page 131 (body page 119)
\Page{%
\TextAt{2}{79.226}{796.943}{II.}%
\TextAt{2}{94.073}{796.943}{The}%
\TextAt{2}{116.492}{796.943}{Theory}%
\TextAt{2}{154.401}{796.943}{of}%
\TextAt{2}{166.826}{796.943}{Valuations}%
\Vector{q 0 G 0 g 0.299 w 7.226 18.729 m 460.769 18.729 l S Q}%
\TextAt{8}{79.226}{761.124}{Passing}%
\TextAt{8}{122.281}{761.124}{to}%
\TextAt{8}{137.408}{761.124}{the}%
\TextAt{8}{158.381}{761.124}{limit}%
\TextAt{8}{187.162}{761.124}{as}%
\TextAt{13}{202.349}{761.124}{n}%
\NativeMathAt{214.049}{761.124}{11.955}{\rightarrow}%
\NativeMathAt{230.715}{761.124}{11.955}{\infty}%
\TextAt{8}{242.673}{761.124}{,}%
\TextAt{8}{250.843}{761.124}{we}%
\TextAt{8}{268.888}{761.124}{will}%
\TextAt{8}{291.81}{761.124}{finally}%
\TextAt{8}{328.074}{761.124}{obtain}%
\TextAt{8}{365.318}{761.124}{the}%
\TextAt{8}{386.292}{761.124}{identity}%
\TextAt{13}{430.362}{761.124}{f}%
\TextAt{9}{442.122}{761.124}{\Character{61}}%
\TextAt{13}{455.94}{761.124}{g}%
\TextAt{13}{461.954}{761.124}{h}%
\TextAt{8}{468.712}{761.124}{.}%
\TextAt{8}{479.612}{761.124}{For}%
\TextAt{13}{501.398}{761.124}{n}%
\TextAt{9}{513.099}{761.124}{\Character{61}}%
\TextAt{9}{526.907}{761.124}{1}%
\TextAt{8}{79.226}{744.078}{the}%
\TextAt{8}{100.175}{744.078}{congruence}%
\TextAt{8}{161.448}{744.078}{is}%
\TextAt{8}{174.004}{744.078}{satisfied}%
\TextAt{8}{219.796}{744.078}{in}%
\TextAt{8}{234.241}{744.078}{view}%
\TextAt{8}{262.016}{744.078}{of}%
\TextAt{8}{276.138}{744.078}{our}%
\TextAt{8}{297.744}{744.078}{choice}%
\TextAt{8}{333.326}{744.078}{of}%
\TextAt{13}{347.448}{744.078}{\Character{25}}%
\TextAt{8}{354.517}{744.078}{.}%
\TextAt{8}{365.342}{744.078}{Let}%
\TextAt{8}{387.104}{744.078}{us}%
\TextAt{8}{402.912}{744.078}{assume}%
\TextAt{8}{444.148}{744.078}{that}%
\TextAt{8}{470.309}{744.078}{it}%
\TextAt{8}{482.806}{744.078}{is}%
\TextAt{8}{495.362}{744.078}{already}%
\TextAt{8}{79.226}{727.032}{established}%
\TextAt{8}{139.181}{727.032}{for}%
\TextAt{8}{157.066}{727.032}{some}%
\TextAt{13}{186.392}{727.032}{n}%
\NativeMathAt{196.701}{727.032}{11.955}{\geq}%
\TextAt{9}{209.32}{727.032}{1}%
\TextAt{8}{215.173}{727.032}{.}%
\TextAt{8}{223.621}{727.032}{Then,}%
\TextAt{8}{257.431}{727.032}{in}%
\TextAt{8}{271.083}{727.032}{view}%
\TextAt{8}{298.066}{727.032}{of}%
\TextAt{8}{311.396}{727.032}{the}%
\TextAt{8}{331.553}{727.032}{relation}%
\TextAt{13}{201.319}{698.031}{g}%
\TextAt{21}{206.924}{696.237}{n}%
\TextAt{9}{215.881}{698.031}{\Character{61}}%
\TextAt{13}{228.307}{698.031}{g}%
\TextAt{21}{233.912}{696.237}{n}%
\TextAt{3}{239.05}{696.237}{\Character{0}}%
\TextAt{17}{245.637}{696.237}{1}%
\TextAt{9}{253.026}{698.031}{\Character{43}}%
\TextAt{13}{264.787}{698.031}{p}%
\TextAt{21}{270.662}{696.237}{n}%
\TextAt{13}{276.298}{698.031}{\Character{25}}%
\TextAt{21}{283.368}{702.967}{n}%
\TextAt{13}{289.004}{698.031}{\Character{59}}%
\TextAt{13}{318.139}{698.031}{h}%
\TextAt{21}{324.897}{696.237}{n}%
\TextAt{9}{333.854}{698.031}{\Character{61}}%
\TextAt{13}{346.28}{698.031}{h}%
\TextAt{21}{353.018}{696.237}{n}%
\TextAt{3}{358.156}{696.237}{\Character{0}}%
\TextAt{17}{364.743}{696.237}{1}%
\TextAt{9}{372.132}{698.031}{\Character{43}}%
\TextAt{13}{383.893}{698.031}{q}%
\TextAt{21}{389.083}{696.237}{n}%
\TextAt{13}{394.72}{698.031}{\Character{25}}%
\TextAt{21}{401.789}{702.967}{n}%
\TextAt{13}{407.425}{698.031}{\Character{59}}%
\TextAt{8}{79.226}{669.029}{the}%
\TextAt{8}{99.383}{669.029}{condition}%
\TextAt{8}{150.766}{669.029}{on}%
\TextAt{13}{167.025}{669.029}{g}%
\TextAt{21}{172.63}{667.236}{n}%
\TextAt{13}{178.267}{669.029}{\Character{59}}%
\TextAt{13}{183.491}{669.029}{h}%
\TextAt{21}{190.249}{667.236}{n}%
\TextAt{8}{199.783}{669.029}{reduces}%
\TextAt{8}{241.459}{669.029}{to}%
\TextAt{13}{175.688}{640.028}{f}%
\TextAt{16}{185.391}{640.028}{\Character{0}}%
\TextAt{13}{197.346}{640.028}{g}%
\TextAt{21}{202.951}{638.235}{n}%
\TextAt{3}{208.09}{638.235}{\Character{0}}%
\TextAt{17}{214.676}{638.235}{1}%
\TextAt{13}{219.408}{640.028}{h}%
\TextAt{21}{226.147}{638.235}{n}%
\TextAt{3}{231.285}{638.235}{\Character{0}}%
\TextAt{17}{237.872}{638.235}{1}%
\NativeMathAt{245.925}{640.028}{11.955}{\equiv}%
\TextAt{9}{258.544}{640.028}{\Character{40}}%
\TextAt{13}{263.096}{640.028}{g}%
\TextAt{21}{268.702}{638.235}{n}%
\TextAt{3}{273.84}{638.235}{\Character{0}}%
\TextAt{17}{280.426}{638.235}{1}%
\TextAt{13}{285.159}{640.028}{q}%
\TextAt{21}{290.349}{638.235}{n}%
\TextAt{9}{298.642}{640.028}{\Character{43}}%
\TextAt{13}{310.403}{640.028}{h}%
\TextAt{21}{317.142}{638.235}{n}%
\TextAt{3}{322.28}{638.235}{\Character{0}}%
\TextAt{17}{328.867}{638.235}{1}%
\TextAt{13}{333.599}{640.028}{p}%
\TextAt{21}{339.474}{638.235}{n}%
\TextAt{9}{345.11}{640.028}{\Character{41}}%
\TextAt{13}{349.663}{640.028}{\Character{25}}%
\TextAt{21}{356.732}{644.965}{n}%
\TextAt{9}{374.323}{640.028}{\Character{40}}%
\TextAt{8}{378.876}{640.028}{mod}%
\TextAt{13}{404.978}{640.028}{\Character{25}}%
\TextAt{21}{412.047}{644.965}{n}%
\TextAt{17}{417.185}{644.964}{\Character{43}1}%
\TextAt{9}{428.504}{640.028}{\Character{41}}%
\TextAt{13}{433.056}{640.028}{\Character{58}}%
\TextAt{8}{79.226}{611.027}{Dividing}%
\TextAt{8}{126.867}{611.027}{by}%
\TextAt{13}{143.127}{611.027}{\Character{25}}%
\TextAt{21}{150.196}{615.366}{n}%
\TextAt{8}{155.832}{611.027}{,}%
\TextAt{8}{162.981}{611.027}{this}%
\TextAt{8}{185.804}{611.027}{means}%
\TextAt{13}{185.324}{582.026}{g}%
\TextAt{21}{190.929}{580.233}{n}%
\TextAt{3}{196.067}{580.233}{\Character{0}}%
\TextAt{17}{202.654}{580.233}{1}%
\TextAt{13}{207.386}{582.026}{q}%
\TextAt{21}{212.577}{580.233}{n}%
\TextAt{9}{220.87}{582.026}{\Character{43}}%
\TextAt{13}{232.631}{582.026}{h}%
\TextAt{21}{239.369}{580.233}{n}%
\TextAt{3}{244.508}{580.233}{\Character{0}}%
\TextAt{17}{251.094}{580.233}{1}%
\TextAt{13}{255.826}{582.026}{p}%
\TextAt{21}{261.702}{580.233}{n}%
\NativeMathAt{270.659}{582.026}{11.955}{\equiv}%
\TextAt{13}{283.278}{582.026}{g}%
\TextAt{17}{288.883}{580.233}{0}%
\TextAt{13}{293.616}{582.026}{q}%
\TextAt{21}{298.806}{580.233}{n}%
\TextAt{9}{307.099}{582.026}{\Character{43}}%
\TextAt{13}{318.86}{582.026}{h}%
\TextAt{17}{325.599}{580.233}{0}%
\TextAt{13}{330.331}{582.026}{p}%
\TextAt{21}{336.206}{580.233}{n}%
\NativeMathAt{345.163}{582.026}{11.955}{\equiv}%
\TextAt{13}{357.783}{582.026}{f}%
\TextAt{21}{363.553}{580.233}{n}%
\TextAt{9}{381.144}{582.026}{\Character{40}}%
\TextAt{8}{385.696}{582.026}{mod}%
\TextAt{13}{411.798}{582.026}{\Character{25}}%
\TextAt{9}{418.868}{582.026}{\Character{41}}%
\TextAt{13}{423.42}{582.026}{\Character{59}}%
\TextAt{8}{79.226}{553.025}{where}%
\TextAt{13}{113.035}{553.025}{f}%
\TextAt{21}{118.805}{551.232}{n}%
\TextAt{9}{127.763}{553.025}{\Character{61}}%
\TextAt{13}{140.188}{553.025}{\Character{25}}%
\TextAt{3}{147.257}{557.363}{\Character{0}}%
\TextAt{21}{153.844}{557.363}{n}%
\TextAt{9}{159.48}{553.025}{\Character{40}}%
\TextAt{13}{164.032}{553.025}{f}%
\TextAt{16}{173.736}{553.025}{\Character{0}}%
\TextAt{13}{185.691}{553.025}{g}%
\TextAt{21}{191.296}{551.232}{n}%
\TextAt{3}{196.434}{551.232}{\Character{0}}%
\TextAt{17}{203.021}{551.232}{1}%
\TextAt{13}{207.753}{553.025}{h}%
\TextAt{21}{214.492}{551.232}{n}%
\TextAt{3}{219.63}{551.232}{\Character{0}}%
\TextAt{17}{226.216}{551.232}{1}%
\TextAt{9}{230.949}{553.025}{\Character{41}}%
\NativeMathAt{238.822}{553.025}{11.955}{\in}%
\TextAt{14}{250.113}{553.025}{o}%
\TextAt{9}{255.957}{553.025}{\Character{91}}%
\TextAt{13}{259.209}{553.025}{x}%
\TextAt{9}{265.861}{553.025}{\Character{93}}%
\TextAt{8}{269.112}{553.025}{.}%
\TextAt{8}{277.561}{553.025}{Since}%
\TextAt{13}{308.118}{553.025}{g}%
\TextAt{17}{313.723}{551.232}{0}%
\TextAt{13}{318.456}{553.025}{a}%
\TextAt{9}{327.257}{553.025}{\Character{43}}%
\TextAt{13}{339.019}{553.025}{h}%
\TextAt{17}{345.757}{551.232}{0}%
\TextAt{13}{350.49}{553.025}{b}%
\NativeMathAt{358.787}{553.025}{11.955}{\equiv}%
\TextAt{9}{371.407}{553.025}{1}%
\TextAt{8}{380.581}{553.025}{mod}%
\TextAt{13}{406.019}{553.025}{\Character{25}}%
\TextAt{8}{413.088}{553.025}{,}%
\TextAt{8}{420.237}{553.025}{one}%
\TextAt{8}{441.697}{553.025}{has}%
\TextAt{13}{228.503}{524.024}{g}%
\TextAt{17}{234.108}{522.231}{0}%
\TextAt{13}{238.84}{524.024}{af}%
\TextAt{21}{250.755}{522.231}{n}%
\TextAt{9}{259.048}{524.024}{\Character{43}}%
\TextAt{13}{270.81}{524.024}{h}%
\TextAt{17}{277.548}{522.231}{0}%
\TextAt{13}{282.28}{524.024}{bf}%
\TextAt{21}{293.028}{522.231}{n}%
\NativeMathAt{301.985}{524.024}{11.955}{\equiv}%
\TextAt{13}{314.604}{524.024}{f}%
\TextAt{21}{320.374}{522.231}{n}%
\TextAt{9}{337.965}{524.024}{\Character{40}}%
\TextAt{8}{342.518}{524.024}{mod}%
\TextAt{13}{368.62}{524.024}{\Character{25}}%
\TextAt{9}{375.689}{524.024}{\Character{41}}%
\TextAt{13}{380.241}{524.024}{\Character{58}}%
\TextAt{8}{79.226}{495.023}{At}%
\TextAt{8}{96.322}{495.023}{this}%
\TextAt{8}{119.335}{495.023}{point}%
\TextAt{8}{150.096}{495.023}{we}%
\TextAt{8}{167.515}{495.023}{would}%
\TextAt{8}{201.85}{495.023}{like}%
\TextAt{8}{223.5}{495.023}{to}%
\TextAt{8}{238.002}{495.023}{put}%
\TextAt{13}{259.652}{495.023}{q}%
\TextAt{21}{264.843}{493.229}{n}%
\TextAt{9}{274.126}{495.023}{\Character{61}}%
\TextAt{13}{286.877}{495.023}{af}%
\TextAt{21}{298.792}{493.229}{n}%
\TextAt{8}{308.517}{495.023}{and}%
\TextAt{13}{331.47}{495.023}{p}%
\TextAt{21}{337.346}{493.229}{n}%
\TextAt{9}{346.628}{495.023}{\Character{61}}%
\TextAt{13}{359.38}{495.023}{bf}%
\TextAt{21}{370.127}{493.229}{n}%
\TextAt{8}{375.763}{495.023}{,}%
\TextAt{8}{383.151}{495.023}{but}%
\TextAt{8}{404.802}{495.023}{the}%
\TextAt{8}{425.149}{495.023}{degrees}%
\TextAt{8}{466.371}{495.023}{might}%
\TextAt{8}{500.06}{495.023}{be}%
\TextAt{8}{516.176}{495.023}{too}%
\TextAt{8}{79.226}{477.977}{big.}%
\TextAt{8}{103.288}{477.977}{For}%
\TextAt{8}{124.257}{477.977}{this}%
\TextAt{8}{147.08}{477.977}{reason,}%
\TextAt{8}{186.819}{477.977}{we}%
\TextAt{8}{204.046}{477.977}{write}%
\TextAt{13}{227.972}{448.976}{b}%
\TextAt{9}{232.949}{448.976}{\Character{40}}%
\TextAt{13}{237.502}{448.976}{x}%
\TextAt{9}{244.154}{448.976}{\Character{41}}%
\TextAt{13}{248.706}{448.976}{f}%
\TextAt{21}{254.476}{447.182}{n}%
\TextAt{9}{260.112}{448.976}{\Character{40}}%
\TextAt{13}{264.665}{448.976}{x}%
\TextAt{9}{271.317}{448.976}{\Character{41}}%
\TextAt{9}{279.172}{448.976}{\Character{61}}%
\TextAt{13}{291.615}{448.976}{q}%
\TextAt{9}{297.235}{448.976}{\Character{40}}%
\TextAt{13}{301.787}{448.976}{x}%
\TextAt{9}{308.439}{448.976}{\Character{41}}%
\TextAt{13}{312.991}{448.976}{g}%
\TextAt{17}{318.597}{447.182}{0}%
\TextAt{9}{323.329}{448.976}{\Character{40}}%
\TextAt{13}{327.881}{448.976}{x}%
\TextAt{9}{334.533}{448.976}{\Character{41}}%
\TextAt{9}{341.73}{448.976}{\Character{43}}%
\TextAt{13}{353.504}{448.976}{p}%
\TextAt{21}{359.379}{447.182}{n}%
\TextAt{9}{365.015}{448.976}{\Character{40}}%
\TextAt{13}{369.567}{448.976}{x}%
\TextAt{9}{376.219}{448.976}{\Character{41}}%
\TextAt{13}{380.772}{448.976}{\Character{59}}%
\TextAt{8}{79.226}{419.974}{where}%
\TextAt{8}{113.947}{419.974}{deg}%
\TextAt{9}{131.509}{419.974}{\Character{40}}%
\TextAt{13}{136.061}{419.974}{p}%
\TextAt{21}{141.936}{418.181}{n}%
\TextAt{9}{147.573}{419.974}{\Character{41}}%
\TextAt{13}{156.999}{419.974}{\Character{60}}%
\TextAt{8}{170.978}{419.974}{deg}%
\TextAt{9}{188.54}{419.974}{\Character{40}}%
\TextAt{13}{193.092}{419.974}{g}%
\TextAt{17}{198.697}{418.181}{0}%
\TextAt{9}{203.43}{419.974}{\Character{41}}%
\TextAt{9}{212.839}{419.974}{\Character{61}}%
\TextAt{13}{226.835}{419.974}{m}%
\TextAt{8}{237.074}{419.974}{.}%
\TextAt{8}{248.256}{419.974}{Since}%
\TextAt{13}{279.725}{419.974}{g}%
\TextAt{17}{285.33}{418.181}{0}%
\NativeMathAt{294.937}{419.974}{11.955}{\equiv}%
\TextAt{9}{309.525}{419.974}{\Character{22}}%
\TextAt{13}{309.109}{419.974}{g}%
\TextAt{8}{318.464}{419.974}{mod}%
\TextAt{14}{343.902}{419.974}{p}%
\TextAt{8}{354.694}{419.974}{and}%
\TextAt{8}{378.368}{419.974}{deg}%
\TextAt{9}{395.93}{419.974}{\Character{40}}%
\TextAt{13}{400.483}{419.974}{g}%
\TextAt{17}{406.088}{418.181}{0}%
\TextAt{9}{410.82}{419.974}{\Character{41}}%
\TextAt{9}{420.229}{419.974}{\Character{61}}%
\TextAt{8}{434.225}{419.974}{deg}%
\TextAt{9}{451.787}{419.974}{\Character{40}}%
\TextAt{9}{456.737}{419.974}{\Character{22}}%
\TextAt{13}{456.34}{419.974}{g}%
\TextAt{9}{462.374}{419.974}{\Character{41}}%
\TextAt{8}{466.926}{419.974}{,}%
\TextAt{8}{475.215}{419.974}{the}%
\TextAt{8}{496.282}{419.974}{highest}%
\TextAt{8}{79.226}{402.928}{coefficient}%
\TextAt{8}{133.849}{402.928}{of}%
\TextAt{13}{147.179}{402.928}{g}%
\TextAt{17}{152.785}{401.135}{0}%
\TextAt{8}{161.414}{402.928}{is}%
\TextAt{8}{173.178}{402.928}{a}%
\TextAt{8}{182.934}{402.928}{unit;}%
\TextAt{8}{210.897}{402.928}{hence}%
\TextAt{13}{243.403}{402.928}{q}%
\TextAt{9}{249.022}{402.928}{\Character{40}}%
\TextAt{13}{253.574}{402.928}{x}%
\TextAt{9}{260.226}{402.928}{\Character{41}}%
\NativeMathAt{268.1}{402.928}{11.955}{\in}%
\TextAt{14}{279.391}{402.928}{o}%
\TextAt{9}{285.235}{402.928}{\Character{91}}%
\TextAt{13}{288.486}{402.928}{x}%
\TextAt{9}{295.139}{402.928}{\Character{93}}%
\TextAt{8}{302.288}{402.928}{and}%
\TextAt{8}{325.05}{402.928}{we}%
\TextAt{8}{342.278}{402.928}{obtain}%
\TextAt{8}{378.705}{402.928}{the}%
\TextAt{8}{398.862}{402.928}{congruence}%
\TextAt{13}{210.632}{373.927}{g}%
\TextAt{17}{216.238}{372.134}{0}%
\TextAt{9}{220.97}{373.927}{\Character{40}}%
\TextAt{13}{225.522}{373.927}{af}%
\TextAt{21}{237.437}{372.134}{n}%
\TextAt{9}{245.73}{373.927}{\Character{43}}%
\TextAt{13}{257.492}{373.927}{h}%
\TextAt{17}{264.23}{372.134}{0}%
\TextAt{13}{268.962}{373.927}{q}%
\TextAt{9}{274.582}{373.927}{\Character{41}}%
\TextAt{9}{281.779}{373.927}{\Character{43}}%
\TextAt{13}{293.552}{373.927}{h}%
\TextAt{17}{300.29}{372.134}{0}%
\TextAt{13}{305.023}{373.927}{p}%
\TextAt{21}{310.898}{372.134}{n}%
\NativeMathAt{319.855}{373.927}{11.955}{\equiv}%
\TextAt{13}{332.474}{373.927}{f}%
\TextAt{21}{338.244}{372.134}{n}%
\TextAt{9}{355.836}{373.927}{\Character{40}}%
\TextAt{8}{360.388}{373.927}{mod}%
\TextAt{13}{386.49}{373.927}{\Character{25}}%
\TextAt{9}{393.559}{373.927}{\Character{41}}%
\TextAt{13}{398.112}{373.927}{\Character{58}}%
\TextAt{8}{79.226}{344.926}{Omitting}%
\TextAt{8}{131.517}{344.926}{now}%
\TextAt{8}{157.458}{344.926}{from}%
\TextAt{8}{186.651}{344.926}{the}%
\TextAt{8}{208.361}{344.926}{polynomial}%
\TextAt{13}{270.406}{344.926}{af}%
\TextAt{21}{282.321}{343.133}{n}%
\TextAt{9}{291.673}{344.926}{\Character{43}}%
\TextAt{13}{304.493}{344.926}{h}%
\TextAt{17}{311.231}{343.133}{0}%
\TextAt{13}{315.964}{344.926}{q}%
\TextAt{8}{327.033}{344.926}{all}%
\TextAt{8}{344.845}{344.926}{coefficients}%
\TextAt{8}{405.636}{344.926}{divisible}%
\TextAt{8}{453.096}{344.926}{by}%
\TextAt{13}{470.908}{344.926}{\Character{25}}%
\TextAt{8}{477.977}{344.926}{,}%
\TextAt{8}{487.068}{344.926}{we}%
\TextAt{8}{505.848}{344.926}{get}%
\TextAt{8}{526.911}{344.926}{a}%
\TextAt{8}{79.226}{327.88}{polynomial}%
\TextAt{13}{141.027}{327.88}{q}%
\TextAt{21}{146.218}{326.087}{n}%
\TextAt{8}{157.06}{327.88}{such}%
\TextAt{8}{184.765}{327.88}{that}%
\TextAt{13}{211.442}{327.88}{g}%
\TextAt{17}{217.047}{326.087}{0}%
\TextAt{13}{221.779}{327.88}{q}%
\TextAt{21}{226.97}{326.087}{n}%
\TextAt{9}{236.154}{327.88}{\Character{43}}%
\TextAt{13}{248.807}{327.88}{h}%
\TextAt{17}{255.546}{326.087}{0}%
\TextAt{13}{260.278}{327.88}{p}%
\TextAt{21}{266.153}{326.087}{n}%
\NativeMathAt{277.34}{327.88}{11.955}{\equiv}%
\TextAt{13}{292.189}{327.88}{f}%
\TextAt{21}{297.959}{326.087}{n}%
\TextAt{8}{306.916}{327.88}{mod}%
\TextAt{13}{332.354}{327.88}{\Character{25}}%
\TextAt{8}{344.629}{327.88}{and}%
\TextAt{8}{368.699}{327.88}{which,}%
\TextAt{8}{407.073}{327.88}{in}%
\TextAt{8}{422.034}{327.88}{view}%
\TextAt{8}{450.325}{327.88}{of}%
\TextAt{8}{464.963}{327.88}{deg}%
\TextAt{9}{482.525}{327.88}{\Character{40}}%
\TextAt{13}{487.078}{327.88}{f}%
\TextAt{21}{492.848}{326.087}{n}%
\TextAt{9}{498.484}{327.88}{\Character{41}}%
\NativeMathAt{508.586}{327.88}{11.955}{\leq}%
\TextAt{13}{523.435}{327.88}{d}%
\TextAt{8}{529.518}{327.88}{,}%
\TextAt{8}{79.226}{310.834}{deg}%
\TextAt{9}{96.788}{310.834}{\Character{40}}%
\TextAt{13}{101.341}{310.834}{g}%
\TextAt{17}{106.946}{309.041}{0}%
\TextAt{9}{111.678}{310.834}{\Character{41}}%
\TextAt{9}{119.533}{310.834}{\Character{61}}%
\TextAt{13}{131.977}{310.834}{m}%
\TextAt{8}{146.114}{310.834}{and}%
\TextAt{8}{168.876}{310.834}{deg}%
\TextAt{9}{186.438}{310.834}{\Character{40}}%
\TextAt{13}{190.991}{310.834}{h}%
\TextAt{17}{197.729}{309.041}{0}%
\TextAt{13}{202.462}{310.834}{p}%
\TextAt{21}{208.337}{309.041}{n}%
\TextAt{9}{213.973}{310.834}{\Character{41}}%
\TextAt{13}{221.846}{310.834}{\Character{60}}%
\TextAt{9}{234.272}{310.834}{\Character{40}}%
\TextAt{13}{238.824}{310.834}{d}%
\TextAt{16}{247.563}{310.834}{\Character{0}}%
\TextAt{13}{259.518}{310.834}{m}%
\TextAt{9}{269.758}{310.834}{\Character{41}}%
\TextAt{9}{276.955}{310.834}{\Character{43}}%
\TextAt{13}{288.728}{310.834}{m}%
\TextAt{9}{302.288}{310.834}{\Character{61}}%
\TextAt{13}{314.714}{310.834}{d}%
\TextAt{8}{320.796}{310.834}{,}%
\TextAt{8}{327.946}{310.834}{has}%
\TextAt{8}{348.819}{310.834}{degree}%
\NativeMathAt{385.235}{310.834}{11.955}{\leq}%
\TextAt{13}{397.854}{310.834}{d}%
\TextAt{16}{406.593}{310.834}{\Character{0}}%
\TextAt{13}{418.549}{310.834}{m}%
\TextAt{8}{432.685}{310.834}{as}%
\TextAt{8}{447.055}{310.834}{required.}%
\Vector{q 0 G 0 g 0.399 w 452.334 -463.166 m 452.334 -455.096 l S Q}%
\Vector{q 0 G 0 g 0.399 w 452.534 -455.295 m 459.707 -455.295 l S Q}%
\Vector{q 0 G 0 g 0.399 w 452.534 -462.967 m 459.707 -462.967 l S Q}%
\Vector{q 0 G 0 g 0.399 w 459.906 -463.166 m 459.906 -455.096 l S Q}%
\TextAt{12}{97.159}{281.833}{Example:}%
\TextAt{8}{156.606}{281.833}{The}%
\TextAt{8}{179.628}{281.833}{polynomial}%
\TextAt{13}{239.09}{281.833}{x}%
\TextAt{21}{245.742}{286.171}{p}%
\TextAt{3}{250.004}{286.171}{\Character{0}}%
\TextAt{17}{256.591}{286.171}{1}%
\TextAt{16}{261.87}{281.833}{\Character{0}}%
\TextAt{9}{271.715}{281.833}{1}%
\NativeMathAt{280.889}{281.833}{11.955}{\in}%
\TextAt{18}{292.18}{281.833}{Z}%
\TextAt{21}{300.15}{280.04}{p}%
\TextAt{9}{304.911}{281.833}{\Character{91}}%
\TextAt{13}{308.163}{281.833}{x}%
\TextAt{9}{314.815}{281.833}{\Character{93}}%
\TextAt{8}{320.932}{281.833}{splits}%
\TextAt{8}{350.589}{281.833}{over}%
\TextAt{8}{374.603}{281.833}{the}%
\TextAt{8}{393.728}{281.833}{residue}%
\TextAt{8}{432.423}{281.833}{class}%
\TextAt{8}{458.829}{281.833}{field}%
\TextAt{18}{483.154}{281.833}{Z}%
\TextAt{21}{491.124}{280.04}{p}%
\TextAt{9}{495.885}{281.833}{\Character{47}}%
\TextAt{13}{501.738}{281.833}{p}%
\TextAt{18}{507.613}{281.833}{Z}%
\TextAt{21}{515.583}{280.04}{p}%
\TextAt{9}{523.665}{281.833}{\Character{61}}%
\TextAt{18}{79.226}{264.787}{F}%
\TextAt{21}{86.532}{262.994}{p}%
\TextAt{8}{95.202}{264.787}{into}%
\TextAt{8}{118.957}{264.787}{distinct}%
\TextAt{8}{161.302}{264.787}{linear}%
\TextAt{8}{193.832}{264.787}{factors.}%
\TextAt{8}{236.531}{264.787}{Applying}%
\TextAt{8}{287.269}{264.787}{(repeatedly)}%
\TextAt{8}{353.62}{264.787}{Hensel’s}%
\TextAt{8}{398.941}{264.787}{lemma,}%
\TextAt{8}{439.926}{264.787}{we}%
\TextAt{8}{457.166}{264.787}{see}%
\TextAt{8}{476.09}{264.787}{that}%
\TextAt{8}{501.471}{264.787}{it}%
\TextAt{8}{513.187}{264.787}{also}%
\TextAt{8}{79.226}{247.741}{splits}%
\TextAt{8}{109.11}{247.741}{into}%
\TextAt{8}{132.048}{247.741}{linear}%
\TextAt{8}{163.761}{247.741}{factors}%
\TextAt{8}{201.069}{247.741}{over}%
\TextAt{18}{225.31}{247.741}{Z}%
\TextAt{21}{233.28}{245.948}{p}%
\TextAt{8}{238.041}{247.741}{.}%
\TextAt{8}{246.221}{247.741}{We}%
\TextAt{8}{265.56}{247.741}{thus}%
\TextAt{8}{290.507}{247.741}{obtain}%
\TextAt{8}{326.129}{247.741}{the}%
\TextAt{8}{345.48}{247.741}{astonishing}%
\TextAt{8}{405.945}{247.741}{result}%
\TextAt{8}{437.718}{247.741}{that}%
\TextAt{8}{462.282}{247.741}{the}%
\TextAt{8}{481.633}{247.741}{field}%
\TextAt{18}{506.185}{247.741}{Q}%
\TextAt{21}{515.484}{245.948}{p}%
\TextAt{8}{523.337}{247.741}{of}%
\TextAt{13}{79.226}{230.695}{p}%
\TextAt{8}{85.101}{230.695}{-adic}%
\TextAt{8}{113.194}{230.695}{numbers}%
\TextAt{8}{159.889}{230.695}{contains}%
\TextAt{8}{205.292}{230.695}{the}%
\TextAt{9}{224.933}{230.695}{\Character{40}}%
\TextAt{13}{229.485}{230.695}{p}%
\TextAt{16}{236.962}{230.695}{\Character{0}}%
\TextAt{9}{247.861}{230.695}{1\Character{41}}%
\TextAt{8}{258.267}{230.695}{-th}%
\TextAt{8}{276.604}{230.695}{roots}%
\TextAt{8}{305.748}{230.695}{of}%
\TextAt{8}{318.562}{230.695}{unity.}%
\TextAt{8}{352.53}{230.695}{These,}%
\TextAt{8}{389.238}{230.695}{together}%
\TextAt{8}{434.905}{230.695}{with}%
\TextAt{8}{461.049}{230.695}{0,}%
\TextAt{8}{473.643}{230.695}{even}%
\TextAt{8}{499.787}{230.695}{form}%
\TextAt{8}{526.911}{230.695}{a}%
\TextAt{8}{79.226}{213.649}{system}%
\TextAt{8}{118.045}{213.649}{of}%
\TextAt{8}{131.375}{213.649}{representatives}%
\TextAt{8}{211.175}{213.649}{for}%
\TextAt{8}{229.06}{213.649}{the}%
\TextAt{8}{249.217}{213.649}{residue}%
\TextAt{8}{288.944}{213.649}{class}%
\TextAt{8}{316.381}{213.649}{field,}%
\TextAt{8}{344.99}{213.649}{which}%
\TextAt{8}{378.476}{213.649}{is}%
\TextAt{8}{390.24}{213.649}{closed}%
\TextAt{8}{424.766}{213.649}{under}%
\TextAt{8}{457.93}{213.649}{multiplication.}%
\TextAt{12}{79.226}{187.637}{Corollary}%
\TextAt{12}{139.482}{187.637}{(4.7).}%
\TextAt{8}{175.204}{187.637}{Let}%
\TextAt{8}{196.58}{187.637}{the}%
\TextAt{8}{217.142}{187.637}{field}%
\TextAt{13}{242.906}{187.637}{K}%
\TextAt{8}{258.027}{187.637}{be}%
\TextAt{8}{274.358}{187.637}{complete}%
\TextAt{8}{324.187}{187.637}{with}%
\TextAt{8}{351.254}{187.637}{respect}%
\TextAt{8}{391.71}{187.637}{to}%
\TextAt{8}{406.427}{187.637}{the}%
\TextAt{8}{426.99}{187.637}{nonarchimedean}%
\TextAt{8}{514.227}{187.637}{val-}%
\TextAt{8}{79.226}{170.591}{uation}%
\TextAt{16}{116.143}{170.591}{\Character{106}}%
\TextAt{16}{123.83}{170.591}{\Character{106}}%
\TextAt{8}{127.172}{170.591}{.}%
\TextAt{8}{137.09}{170.591}{Then,}%
\TextAt{8}{171.511}{170.591}{for}%
\TextAt{8}{189.886}{170.591}{every}%
\TextAt{8}{221.268}{170.591}{irreducible}%
\TextAt{8}{279.633}{170.591}{polynomial}%
\TextAt{13}{340.616}{170.591}{f}%
\TextAt{9}{347.662}{170.591}{\Character{40}}%
\TextAt{13}{352.214}{170.591}{x}%
\TextAt{9}{358.866}{170.591}{\Character{41}}%
\TextAt{9}{367.557}{170.591}{\Character{61}}%
\TextAt{13}{380.834}{170.591}{a}%
\TextAt{17}{386.979}{168.798}{0}%
\TextAt{9}{394.702}{170.591}{\Character{43}}%
\TextAt{13}{406.797}{170.591}{a}%
\TextAt{17}{412.942}{168.798}{1}%
\TextAt{13}{417.675}{170.591}{x}%
\TextAt{9}{427.317}{170.591}{\Character{43}}%
\TextAt{16}{439.412}{170.591}{\Character{1}}%
\TextAt{16}{444.708}{170.591}{\Character{1}}%
\TextAt{16}{450.016}{170.591}{\Character{1}}%
\TextAt{9}{456.35}{170.591}{\Character{43}}%
\TextAt{13}{468.446}{170.591}{a}%
\TextAt{21}{474.591}{168.798}{n}%
\TextAt{13}{480.227}{170.591}{x}%
\TextAt{21}{486.879}{174.929}{n}%
\NativeMathAt{496.671}{170.591}{11.955}{\in}%
\TextAt{13}{508.796}{170.591}{K}%
\TextAt{9}{519.614}{170.591}{\Character{91}}%
\TextAt{13}{522.866}{170.591}{x}%
\TextAt{9}{529.518}{170.591}{\Character{93}}%
\TextAt{8}{79.226}{153.545}{such}%
\TextAt{8}{105.623}{153.545}{that}%
\TextAt{13}{130.992}{153.545}{a}%
\TextAt{17}{137.137}{151.752}{0}%
\TextAt{13}{141.869}{153.545}{a}%
\TextAt{21}{148.014}{151.752}{n}%
\NativeMathAt{156.971}{153.545}{11.955}{\neq}%
\TextAt{9}{169.381}{153.545}{0}%
\TextAt{8}{175.25}{153.545}{,}%
\TextAt{8}{182.399}{153.545}{one}%
\TextAt{8}{203.859}{153.545}{has}%
\TextAt{16}{252.187}{136.499}{\Character{106}}%
\TextAt{13}{255.508}{136.499}{f}%
\TextAt{16}{262.554}{136.499}{\Character{106}}%
\TextAt{9}{269.196}{136.499}{\Character{61}}%
\TextAt{8}{281.621}{136.499}{max}%
\TextAt{16}{303.416}{136.499}{\Character{102}\Character{106}}%
\TextAt{13}{312.714}{136.499}{a}%
\TextAt{17}{318.859}{134.706}{0}%
\TextAt{16}{323.591}{136.499}{\Character{106}}%
\TextAt{13}{326.912}{136.499}{\Character{59}}%
\TextAt{16}{332.157}{136.499}{\Character{106}}%
\TextAt{13}{335.477}{136.499}{a}%
\TextAt{21}{341.622}{134.706}{n}%
\TextAt{16}{347.259}{136.499}{\Character{106}\Character{103}}%
\TextAt{13}{356.557}{136.499}{\Character{58}}%
\TextAt{8}{79.226}{112.977}{In}%
\TextAt{8}{93.847}{112.977}{particular,}%
\TextAt{13}{151.089}{112.977}{a}%
\TextAt{21}{157.234}{111.184}{n}%
\TextAt{9}{166.191}{112.977}{\Character{61}}%
\TextAt{9}{178.601}{112.977}{1}%
\TextAt{8}{188.367}{112.977}{and}%
\TextAt{13}{211.129}{112.977}{a}%
\TextAt{17}{217.274}{111.184}{0}%
\NativeMathAt{225.327}{112.977}{11.955}{\in}%
\TextAt{14}{236.618}{112.977}{o}%
\TextAt{8}{246.36}{112.977}{imply}%
\TextAt{8}{279.201}{112.977}{that}%
\TextAt{13}{304.57}{112.977}{f}%
\NativeMathAt{314.937}{112.977}{11.955}{\in}%
\TextAt{14}{326.228}{112.977}{o}%
\TextAt{9}{332.072}{112.977}{\Character{91}}%
\TextAt{13}{335.324}{112.977}{x}%
\TextAt{9}{341.976}{112.977}{\Character{93}}%
\TextAt{8}{345.228}{112.977}{.}%
\TextAt{8}{297.211}{43.033}{119}%
\Destination{78.23}{809.815}{nav.page.119}%
\Destination{79.36}{776.155}{nav.print.130}%
\Destination{79.36}{279.755}{nav.print.131}%
\Destination{79.36}{201.655}{nav.number.2.4.7}%
}

% Page 132 (body page 120)
\Page{%
\TextAt{2}{79.226}{796.943}{II.}%
\TextAt{2}{94.073}{796.943}{The}%
\TextAt{2}{116.492}{796.943}{Theory}%
\TextAt{2}{154.401}{796.943}{of}%
\TextAt{2}{166.826}{796.943}{Valuations}%
\Vector{q 0 G 0 g 0.299 w 7.226 18.729 m 460.769 18.729 l S Q}%
\TextAt{19}{79.226}{761.124}{Proof.}%
\TextAt{8}{116}{761.124}{After}%
\TextAt{8}{147.228}{761.124}{multiplying}%
\TextAt{8}{210.341}{761.124}{by}%
\TextAt{8}{227.27}{761.124}{a}%
\TextAt{8}{237.696}{761.124}{suitable}%
\TextAt{8}{282.003}{761.124}{element}%
\TextAt{8}{325.915}{761.124}{of}%
\TextAt{13}{339.916}{761.124}{K}%
\TextAt{8}{355.301}{761.124}{we}%
\TextAt{8}{373.199}{761.124}{may}%
\TextAt{8}{399.239}{761.124}{assume}%
\TextAt{8}{440.354}{761.124}{that}%
\TextAt{13}{466.393}{761.124}{f}%
\NativeMathAt{477.903}{761.124}{11.955}{\in}%
\TextAt{14}{490.337}{761.124}{o}%
\TextAt{9}{496.181}{761.124}{\Character{91}}%
\TextAt{13}{499.433}{761.124}{x}%
\TextAt{9}{506.085}{761.124}{\Character{93}}%
\TextAt{8}{513.904}{761.124}{and}%
\TextAt{16}{79.226}{744.078}{\Character{106}}%
\TextAt{13}{82.547}{744.078}{f}%
\TextAt{16}{89.593}{744.078}{\Character{106}}%
\TextAt{9}{97.114}{744.078}{\Character{61}}%
\TextAt{9}{110.408}{744.078}{1}%
\TextAt{8}{116.271}{744.078}{.}%
\TextAt{8}{126.266}{744.078}{Let}%
\TextAt{13}{147.751}{744.078}{a}%
\TextAt{21}{153.896}{742.284}{r}%
\TextAt{8}{162.864}{744.078}{be}%
\TextAt{8}{179.303}{744.078}{the}%
\TextAt{8}{199.975}{744.078}{first}%
\TextAt{8}{224.617}{744.078}{one}%
\TextAt{8}{246.592}{744.078}{among}%
\TextAt{8}{284.838}{744.078}{the}%
\TextAt{8}{305.51}{744.078}{coefficients}%
\TextAt{13}{365.263}{744.078}{a}%
\TextAt{17}{371.408}{742.284}{0}%
\TextAt{13}{376.14}{744.078}{\Character{59}}%
\TextAt{13}{381.364}{744.078}{\Character{58}}%
\TextAt{13}{386.601}{744.078}{\Character{58}}%
\TextAt{13}{391.837}{744.078}{\Character{58}}%
\TextAt{13}{397.062}{744.078}{\Character{59}}%
\TextAt{13}{402.298}{744.078}{a}%
\TextAt{21}{408.506}{742.284}{n}%
\TextAt{8}{418.556}{744.078}{such}%
\TextAt{8}{445.468}{744.078}{that}%
\TextAt{16}{471.353}{744.078}{\Character{106}}%
\TextAt{13}{474.673}{744.078}{a}%
\TextAt{21}{480.818}{742.284}{r}%
\TextAt{16}{485.373}{744.078}{\Character{106}}%
\TextAt{9}{492.894}{744.078}{\Character{61}}%
\TextAt{9}{506.188}{744.078}{1}%
\TextAt{8}{512.051}{744.078}{.}%
\TextAt{8}{522.046}{744.078}{In}%
\TextAt{8}{79.226}{727.032}{other}%
\TextAt{8}{109.796}{727.032}{words,}%
\TextAt{8}{146.605}{727.032}{we}%
\TextAt{8}{163.833}{727.032}{have}%
\TextAt{13}{182.879}{700.317}{f}%
\TextAt{9}{189.925}{700.317}{\Character{40}}%
\TextAt{13}{194.478}{700.317}{x}%
\TextAt{9}{201.13}{700.317}{\Character{41}}%
\NativeMathAt{209.003}{700.317}{11.955}{\equiv}%
\TextAt{13}{221.622}{700.317}{x}%
\TextAt{21}{228.274}{705.253}{r}%
\TextAt{9}{232.829}{700.317}{\Character{40}}%
\TextAt{13}{237.381}{700.317}{a}%
\TextAt{21}{243.526}{698.524}{r}%
\TextAt{9}{250.738}{700.317}{\Character{43}}%
\TextAt{13}{262.499}{700.317}{a}%
\TextAt{21}{268.644}{698.524}{r}%
\TextAt{17}{272.701}{698.524}{\Character{43}1}%
\TextAt{13}{284.019}{700.317}{x}%
\TextAt{9}{293.328}{700.317}{\Character{43}}%
\TextAt{16}{305.089}{700.317}{\Character{1}}%
\TextAt{16}{310.385}{700.317}{\Character{1}}%
\TextAt{16}{315.693}{700.317}{\Character{1}}%
\TextAt{9}{321.694}{700.317}{\Character{43}}%
\TextAt{13}{333.455}{700.317}{a}%
\TextAt{21}{339.6}{698.524}{n}%
\TextAt{13}{345.236}{700.317}{x}%
\TextAt{21}{351.889}{705.253}{n}%
\TextAt{3}{357.027}{705.253}{\Character{0}}%
\TextAt{21}{363.613}{705.253}{r}%
\TextAt{9}{368.168}{700.317}{\Character{41}}%
\TextAt{9}{384.654}{700.317}{\Character{40}}%
\TextAt{8}{389.228}{700.317}{mod}%
\TextAt{14}{415.33}{700.317}{p}%
\TextAt{9}{421.313}{700.317}{\Character{41}}%
\TextAt{13}{425.865}{700.317}{\Character{58}}%
\TextAt{8}{79.226}{673.602}{If}%
\TextAt{8}{92.1}{673.602}{one}%
\TextAt{8}{114.741}{673.602}{had}%
\TextAt{8}{138.685}{673.602}{max}%
\TextAt{16}{160.479}{673.602}{\Character{102}\Character{106}}%
\TextAt{13}{169.778}{673.602}{a}%
\TextAt{17}{175.923}{671.809}{0}%
\TextAt{16}{180.655}{673.602}{\Character{106}}%
\TextAt{13}{183.976}{673.602}{\Character{59}}%
\TextAt{16}{189.22}{673.602}{\Character{106}}%
\TextAt{13}{192.541}{673.602}{a}%
\TextAt{21}{198.686}{671.809}{n}%
\TextAt{16}{204.322}{673.602}{\Character{106}\Character{103}}%
\TextAt{13}{218.955}{673.602}{\Character{60}}%
\TextAt{9}{233.394}{673.602}{1}%
\TextAt{8}{239.247}{673.602}{,}%
\TextAt{8}{247.873}{673.602}{then}%
\TextAt{9}{275.715}{673.602}{0}%
\TextAt{13}{286.902}{673.602}{\Character{60}}%
\TextAt{13}{301.332}{673.602}{r}%
\TextAt{13}{312.259}{673.602}{\Character{60}}%
\TextAt{13}{326.689}{673.602}{n}%
\TextAt{8}{338.781}{673.602}{and}%
\TextAt{8}{362.725}{673.602}{the}%
\TextAt{8}{384.063}{673.602}{congruence}%
\TextAt{8}{445.726}{673.602}{would}%
\TextAt{8}{481.051}{673.602}{contradict}%
\TextAt{8}{79.226}{656.556}{Hensel’s}%
\TextAt{8}{124.536}{656.556}{lemma.}%
\Vector{q 0 G 0 g 0.399 w 452.334 -117.444 m 452.334 -109.374 l S Q}%
\Vector{q 0 G 0 g 0.399 w 452.534 -109.573 m 459.707 -109.573 l S Q}%
\Vector{q 0 G 0 g 0.399 w 452.534 -117.244 m 459.707 -117.244 l S Q}%
\Vector{q 0 G 0 g 0.399 w 459.906 -117.444 m 459.906 -109.374 l S Q}%
\TextAt{8}{97.159}{628.208}{From}%
\TextAt{8}{127.847}{628.208}{this}%
\TextAt{8}{150.633}{628.208}{corollary}%
\TextAt{8}{199.062}{628.208}{we}%
\TextAt{8}{216.253}{628.208}{can}%
\TextAt{8}{237.675}{628.208}{now}%
\TextAt{8}{262.027}{628.208}{deduce}%
\TextAt{8}{301}{628.208}{the}%
\TextAt{8}{321.12}{628.208}{following}%
\TextAt{8}{370.518}{628.208}{theorem}%
\TextAt{8}{416.006}{628.208}{on}%
\TextAt{8}{432.228}{628.208}{extensions}%
\TextAt{8}{488.572}{628.208}{of}%
\TextAt{8}{501.865}{628.208}{valua-}%
\TextAt{8}{79.226}{611.162}{tions.}%
\TextAt{12}{79.226}{585.803}{Theorem}%
\TextAt{12}{136.634}{585.803}{(4.8).}%
\TextAt{8}{172.292}{585.803}{Let}%
\TextAt{13}{193.604}{585.803}{K}%
\TextAt{8}{208.661}{585.803}{be}%
\TextAt{8}{224.928}{585.803}{complete}%
\TextAt{8}{274.693}{585.803}{with}%
\TextAt{8}{301.695}{585.803}{respect}%
\TextAt{8}{342.087}{585.803}{to}%
\TextAt{8}{356.74}{585.803}{the}%
\TextAt{8}{377.239}{585.803}{valuation}%
\TextAt{16}{428.654}{585.803}{\Character{106}}%
\TextAt{16}{436.198}{585.803}{\Character{106}}%
\TextAt{8}{439.535}{585.803}{.}%
\TextAt{8}{449.01}{585.803}{Then}%
\TextAt{16}{479.91}{585.803}{\Character{106}}%
\TextAt{16}{487.454}{585.803}{\Character{106}}%
\TextAt{8}{495.031}{585.803}{may}%
\TextAt{8}{520.743}{585.803}{be}%
\TextAt{8}{79.226}{568.757}{extended}%
\TextAt{8}{129.808}{568.757}{in}%
\TextAt{8}{144.297}{568.757}{a}%
\TextAt{8}{154.889}{568.757}{unique}%
\TextAt{8}{193.767}{568.757}{way}%
\TextAt{8}{218.347}{568.757}{to}%
\TextAt{8}{233.493}{568.757}{a}%
\TextAt{8}{244.085}{568.757}{valuation}%
\TextAt{8}{295.994}{568.757}{of}%
\TextAt{8}{310.161}{568.757}{any}%
\TextAt{8}{333.114}{568.757}{given}%
\TextAt{8}{364.52}{568.757}{algebraic}%
\TextAt{8}{414.791}{568.757}{extension}%
\TextAt{13}{467.394}{568.757}{L}%
\TextAt{16}{480.104}{568.757}{\Character{106}}%
\TextAt{13}{488.172}{568.757}{K}%
\TextAt{8}{498.989}{568.757}{.}%
\TextAt{8}{509.947}{568.757}{This}%
\TextAt{8}{79.226}{551.711}{extension}%
\TextAt{8}{130.992}{551.711}{is}%
\TextAt{8}{142.756}{551.711}{given}%
\TextAt{8}{173.325}{551.711}{by}%
\TextAt{8}{189.584}{551.711}{the}%
\TextAt{8}{209.741}{551.711}{formula}%
\TextAt{16}{257.625}{523.679}{\Character{106}}%
\TextAt{13}{260.945}{523.679}{\Character{11}}%
\TextAt{16}{268.467}{523.679}{\Character{106}}%
\TextAt{9}{275.109}{523.679}{\Character{61}}%
\TextAt{22}{290.855}{528.38}{n}%
\TextAt{20}{288.814}{537.639}{\Character{113}}%
\Vector{q 0 G 0 g 0.478 w 228.77 -236.122 m 279.119 -236.122 l S Q}%
\TextAt{16}{300.77}{523.679}{\Character{106}}%
\TextAt{8}{306.083}{523.679}{N}%
\TextAt{21}{314.858}{521.819}{L}%
\TextAt{3}{320.617}{521.819}{\Character{106}}%
\TextAt{21}{322.969}{521.819}{K}%
\TextAt{9}{331.172}{523.679}{\Character{40}}%
\TextAt{13}{335.724}{523.679}{\Character{11}}%
\TextAt{9}{343.246}{523.679}{\Character{41}}%
\TextAt{16}{347.798}{523.679}{\Character{106}}%
\TextAt{13}{351.119}{523.679}{\Character{59}}%
\TextAt{8}{79.226}{496.964}{when}%
\TextAt{13}{109.784}{496.964}{L}%
\TextAt{16}{121.069}{496.964}{\Character{106}}%
\TextAt{13}{127.711}{496.964}{K}%
\TextAt{8}{142.426}{496.964}{has}%
\TextAt{8}{163.299}{496.964}{finite}%
\TextAt{8}{193.211}{496.964}{degree}%
\TextAt{13}{229.627}{496.964}{n}%
\TextAt{8}{236.614}{496.964}{.}%
\TextAt{8}{245.063}{496.964}{In}%
\TextAt{8}{259.684}{496.964}{this}%
\TextAt{8}{282.506}{496.964}{case}%
\TextAt{13}{307.277}{496.964}{L}%
\TextAt{8}{319.139}{496.964}{is}%
\TextAt{8}{330.903}{496.964}{again}%
\TextAt{8}{362.13}{496.964}{complete.}%
\TextAt{19}{79.226}{463.292}{Proof.}%
\TextAt{8}{116}{463.292}{If}%
\TextAt{8}{128.687}{463.292}{the}%
\TextAt{8}{149.837}{463.292}{valuation}%
\TextAt{16}{201.904}{463.292}{\Character{106}}%
\TextAt{16}{210.105}{463.292}{\Character{106}}%
\TextAt{8}{218.329}{463.292}{is}%
\TextAt{8}{231.087}{463.292}{archimedean,}%
\TextAt{8}{303.546}{463.292}{then}%
\TextAt{8}{331.2}{463.292}{by}%
\TextAt{8}{348.454}{463.292}{Ostrowski’s}%
\TextAt{8}{412.081}{463.292}{theorem,}%
\TextAt{13}{462.101}{463.292}{K}%
\TextAt{9}{477.933}{463.292}{\Character{61}}%
\TextAt{18}{492.053}{463.292}{R}%
\TextAt{8}{505.579}{463.292}{or}%
\TextAt{18}{520.883}{463.292}{C}%
\TextAt{8}{529.518}{463.292}{.}%
\TextAt{8}{79.226}{446.246}{We}%
\TextAt{8}{100.568}{446.246}{have}%
\TextAt{8}{128.76}{446.246}{N}%
\TextAt{25}{137.536}{444.386}{C}%
\TextAt{3}{143.292}{444.386}{\Character{106}}%
\TextAt{25}{145.644}{444.386}{R}%
\TextAt{9}{151.898}{446.246}{\Character{40}}%
\TextAt{13}{156.451}{446.246}{z}%
\TextAt{9}{162.421}{446.246}{\Character{41}}%
\TextAt{9}{172.32}{446.246}{\Character{61}}%
\TextAt{13}{186.802}{446.246}{z}%
\TextAt{9}{193.481}{446.246}{\Character{22}}%
\TextAt{13}{192.772}{446.246}{z}%
\TextAt{9}{204.104}{446.246}{\Character{61}}%
\TextAt{16}{218.571}{446.246}{\Character{106}}%
\TextAt{13}{221.892}{446.246}{z}%
\TextAt{16}{227.862}{446.246}{\Character{106}}%
\TextAt{17}{231.183}{450.584}{2}%
\TextAt{8}{241.01}{446.246}{and}%
\TextAt{8}{264.971}{446.246}{the}%
\TextAt{8}{286.324}{446.246}{theorem}%
\TextAt{8}{333.047}{446.246}{is}%
\TextAt{8}{346.009}{446.246}{part}%
\TextAt{8}{372.575}{446.246}{of}%
\TextAt{8}{387.103}{446.246}{classical}%
\TextAt{8}{433.3}{446.246}{analysis.}%
\TextAt{8}{485.474}{446.246}{So}%
\TextAt{8}{502.931}{446.246}{let}%
\TextAt{16}{521.033}{446.246}{\Character{106}}%
\TextAt{16}{529.438}{446.246}{\Character{106}}%
\TextAt{8}{79.226}{429.2}{be}%
\TextAt{8}{96.922}{429.2}{nonarchimedean.}%
\TextAt{8}{193.619}{429.2}{Since}%
\TextAt{8}{225.948}{429.2}{every}%
\TextAt{8}{258.612}{429.2}{algebraic}%
\TextAt{8}{309.819}{429.2}{extension}%
\TextAt{13}{363.356}{429.2}{L}%
\TextAt{16}{377.661}{429.2}{\Character{106}}%
\TextAt{13}{387.322}{429.2}{K}%
\TextAt{8}{403.809}{429.2}{is}%
\TextAt{8}{417.345}{429.2}{the}%
\TextAt{8}{439.273}{429.2}{union}%
\TextAt{8}{473.563}{429.2}{of}%
\TextAt{8}{488.664}{429.2}{its}%
\TextAt{8}{506.755}{429.2}{finite}%
\TextAt{8}{79.226}{412.154}{subextensions,}%
\TextAt{8}{156.803}{412.154}{we}%
\TextAt{8}{174.031}{412.154}{may}%
\TextAt{8}{199.4}{412.154}{assume}%
\TextAt{8}{239.844}{412.154}{that}%
\TextAt{8}{265.213}{412.154}{the}%
\TextAt{8}{285.369}{412.154}{degree}%
\TextAt{13}{321.785}{412.154}{n}%
\TextAt{9}{332.093}{412.154}{\Character{61}}%
\TextAt{9}{344.502}{412.154}{\Character{91}}%
\TextAt{13}{347.77}{412.154}{L}%
\TextAt{9}{359.056}{412.154}{\Character{58}}%
\TextAt{13}{365.628}{412.154}{K}%
\TextAt{9}{376.446}{412.154}{\Character{93}}%
\TextAt{8}{383.595}{412.154}{is}%
\TextAt{8}{395.359}{412.154}{finite.}%
\TextAt{12}{97.159}{392.119}{Existence}%
\TextAt{12}{157.659}{392.119}{of}%
\TextAt{12}{173.077}{392.119}{the}%
\TextAt{12}{196.492}{392.119}{extended}%
\TextAt{12}{254.219}{392.119}{valuation:}%
\TextAt{8}{317.138}{392.119}{let}%
\TextAt{14}{334.121}{392.119}{o}%
\TextAt{8}{343.942}{392.119}{be}%
\TextAt{8}{359.945}{392.119}{the}%
\TextAt{8}{380.181}{392.119}{valuation}%
\TextAt{8}{431.332}{392.119}{ring}%
\TextAt{8}{455.477}{392.119}{of}%
\TextAt{13}{468.886}{392.119}{K}%
\TextAt{8}{483.68}{392.119}{and}%
\TextAt{16}{506.522}{392.119}{\Character{79}}%
\TextAt{8}{520.348}{392.119}{its}%
\TextAt{8}{79.226}{375.073}{integral}%
\TextAt{8}{121.834}{375.073}{closure}%
\TextAt{8}{160.916}{375.073}{in}%
\TextAt{13}{174.569}{375.073}{L}%
\TextAt{8}{182.533}{375.073}{.}%
\TextAt{8}{190.981}{375.073}{Then}%
\TextAt{8}{221.539}{375.073}{one}%
\TextAt{8}{242.998}{375.073}{has}%
\TextAt{16}{234.48}{348.359}{\Character{79}}%
\TextAt{9}{247.651}{348.359}{\Character{61}}%
\TextAt{16}{260.076}{348.359}{\Character{102}}%
\TextAt{13}{266.054}{348.359}{\Character{11}}%
\NativeMathAt{276.897}{348.359}{11.955}{\in}%
\TextAt{13}{288.188}{348.359}{L}%
\TextAt{16}{299.473}{348.359}{\Character{106}}%
\TextAt{8}{306.115}{348.359}{N}%
\TextAt{21}{314.89}{346.499}{L}%
\TextAt{3}{320.649}{346.499}{\Character{106}}%
\TextAt{21}{323.001}{346.499}{K}%
\TextAt{9}{331.204}{348.359}{\Character{40}}%
\TextAt{13}{335.756}{348.359}{\Character{11}}%
\TextAt{9}{343.278}{348.359}{\Character{41}}%
\NativeMathAt{351.151}{348.359}{11.955}{\in}%
\TextAt{14}{362.442}{348.359}{o}%
\TextAt{16}{368.286}{348.359}{\Character{103}}%
\TextAt{13}{374.264}{348.359}{\Character{58}}%
\TextAt{8}{517.802}{348.359}{(*)}%
\TextAt{8}{79.226}{321.644}{The}%
\TextAt{8}{103.767}{321.644}{implication}%
\TextAt{13}{165.394}{321.644}{\Character{11}}%
\NativeMathAt{177.067}{321.644}{11.955}{\in}%
\TextAt{16}{189.178}{321.644}{\Character{79}}%
\NativeArrowAt{203.19}{321.644}{11.955}{19.067}{\Longrightarrow}%
\TextAt{8}{226.409}{321.644}{N}%
\TextAt{21}{235.184}{319.784}{L}%
\TextAt{3}{240.943}{319.784}{\Character{106}}%
\TextAt{21}{243.295}{319.784}{K}%
\TextAt{9}{251.498}{321.644}{\Character{40}}%
\TextAt{13}{256.05}{321.644}{\Character{11}}%
\TextAt{9}{263.572}{321.644}{\Character{41}}%
\NativeMathAt{272.276}{321.644}{11.955}{\in}%
\TextAt{14}{284.398}{321.644}{o}%
\TextAt{8}{294.627}{321.644}{is}%
\TextAt{8}{306.878}{321.644}{evident}%
\TextAt{8}{348.336}{321.644}{(see}%
\TextAt{8}{372.292}{321.644}{chap.}%
\TextAt{8}{403.671}{321.644}{I,}%
\TextAt{8}{415.528}{321.644}{§2,}%
\TextAt{8}{434.824}{321.644}{p.}%
\TextAt{8}{448.964}{321.644}{12).}%
\TextAt{8}{475.146}{321.644}{Conversely,}%
\TextAt{8}{79.226}{304.598}{let}%
\TextAt{13}{96.131}{304.598}{\Character{11}}%
\NativeMathAt{106.973}{304.598}{11.955}{\in}%
\TextAt{13}{118.264}{304.598}{L}%
\TextAt{3}{126.229}{308.936}{\Character{3}}%
\TextAt{8}{134.858}{304.598}{and}%
\TextAt{8}{157.621}{304.598}{N}%
\TextAt{21}{166.396}{302.738}{L}%
\TextAt{3}{172.155}{302.738}{\Character{106}}%
\TextAt{21}{174.508}{302.738}{K}%
\TextAt{9}{182.71}{304.598}{\Character{40}}%
\TextAt{13}{187.263}{304.598}{\Character{11}}%
\TextAt{9}{194.784}{304.598}{\Character{41}}%
\NativeMathAt{202.657}{304.598}{11.955}{\in}%
\TextAt{14}{213.948}{304.598}{o}%
\TextAt{8}{219.793}{304.598}{.}%
\TextAt{8}{228.241}{304.598}{Let}%
\TextAt{13}{205.563}{277.883}{f}%
\TextAt{9}{212.609}{277.883}{\Character{40}}%
\TextAt{13}{217.162}{277.883}{x}%
\TextAt{9}{223.814}{277.883}{\Character{41}}%
\TextAt{9}{231.669}{277.883}{\Character{61}}%
\TextAt{13}{244.112}{277.883}{x}%
\TextAt{21}{250.764}{282.819}{d}%
\TextAt{9}{258.276}{277.883}{\Character{43}}%
\TextAt{13}{270.038}{277.883}{a}%
\TextAt{21}{276.183}{276.09}{d}%
\TextAt{3}{280.54}{276.09}{\Character{0}}%
\TextAt{17}{287.127}{276.09}{1}%
\TextAt{13}{291.859}{277.883}{x}%
\TextAt{21}{298.511}{282.819}{d}%
\TextAt{3}{302.868}{282.819}{\Character{0}}%
\TextAt{17}{309.455}{282.819}{1}%
\TextAt{9}{316.844}{277.883}{\Character{43}}%
\TextAt{16}{328.605}{277.883}{\Character{1}}%
\TextAt{16}{333.901}{277.883}{\Character{1}}%
\TextAt{16}{339.209}{277.883}{\Character{1}}%
\TextAt{9}{345.209}{277.883}{\Character{43}}%
\TextAt{13}{356.971}{277.883}{a}%
\TextAt{17}{363.116}{276.09}{0}%
\NativeMathAt{371.169}{277.883}{11.955}{\in}%
\TextAt{13}{382.46}{277.883}{K}%
\TextAt{9}{393.277}{277.883}{\Character{91}}%
\TextAt{13}{396.529}{277.883}{x}%
\TextAt{9}{403.181}{277.883}{\Character{93}}%
\TextAt{8}{79.226}{251.168}{be}%
\TextAt{8}{95.06}{251.168}{the}%
\TextAt{8}{115.126}{251.168}{minimal}%
\TextAt{8}{160.56}{251.168}{polynomial}%
\TextAt{8}{220.963}{251.168}{of}%
\TextAt{13}{234.202}{251.168}{\Character{11}}%
\TextAt{8}{245.531}{251.168}{over}%
\TextAt{13}{270.486}{251.168}{K}%
\TextAt{8}{281.304}{251.168}{.}%
\TextAt{8}{289.722}{251.168}{Then}%
\TextAt{8}{320.189}{251.168}{N}%
\TextAt{21}{328.964}{249.309}{L}%
\TextAt{3}{334.723}{249.309}{\Character{106}}%
\TextAt{21}{337.075}{249.309}{K}%
\TextAt{9}{345.278}{251.168}{\Character{40}}%
\TextAt{13}{349.83}{251.168}{\Character{11}}%
\TextAt{9}{357.352}{251.168}{\Character{41}}%
\TextAt{9}{365.207}{251.168}{\Character{61}}%
\NativeMathAt{377.65}{251.168}{11.955}{\pm}%
\TextAt{13}{386.949}{251.168}{a}%
\TextAt{21}{393.094}{255.507}{m}%
\TextAt{17}{393.094}{248.213}{0}%
\NativeMathAt{404.403}{251.168}{11.955}{\in}%
\TextAt{14}{415.694}{251.168}{o}%
\TextAt{8}{421.539}{251.168}{,}%
\TextAt{8}{428.615}{251.168}{so}%
\TextAt{8}{442.895}{251.168}{that}%
\TextAt{16}{468.173}{251.168}{\Character{106}}%
\TextAt{13}{471.494}{251.168}{a}%
\TextAt{17}{477.639}{249.375}{0}%
\TextAt{16}{482.371}{251.168}{\Character{106}}%
\NativeMathAt{488.994}{251.168}{11.955}{\leq}%
\TextAt{9}{501.632}{251.168}{1}%
\TextAt{8}{507.485}{251.168}{,}%
\TextAt{8}{514.562}{251.168}{i.e.,}%
\TextAt{13}{79.226}{234.122}{a}%
\TextAt{17}{85.371}{232.329}{0}%
\NativeMathAt{93.781}{234.123}{11.955}{\in}%
\TextAt{14}{105.429}{234.123}{o}%
\TextAt{8}{111.273}{234.122}{.}%
\TextAt{8}{120.349}{234.122}{By}%
\TextAt{8}{138.922}{234.122}{(4.7)}%
\TextAt{8}{167.106}{234.122}{this}%
\TextAt{8}{190.138}{234.122}{gives}%
\TextAt{13}{219.028}{234.122}{f}%
\TextAt{9}{226.075}{234.123}{\Character{40}}%
\TextAt{13}{230.627}{234.122}{x}%
\TextAt{9}{237.279}{234.123}{\Character{41}}%
\NativeMathAt{245.509}{234.122}{11.955}{\in}%
\TextAt{14}{257.157}{234.123}{o}%
\TextAt{9}{263.001}{234.122}{\Character{91}}%
\TextAt{13}{266.253}{234.123}{x}%
\TextAt{9}{272.905}{234.122}{\Character{93}}%
\TextAt{8}{276.156}{234.122}{,}%
\TextAt{8}{283.567}{234.122}{i.e.,}%
\TextAt{13}{305.934}{234.122}{\Character{11}}%
\NativeMathAt{317.133}{234.123}{11.955}{\in}%
\TextAt{16}{328.765}{234.123}{\Character{79}}%
\TextAt{8}{338.631}{234.122}{.}%
\TextAt{8}{347.707}{234.122}{For}%
\TextAt{8}{368.886}{234.122}{the}%
\TextAt{8}{389.252}{234.122}{function}%
\TextAt{16}{435.309}{234.122}{\Character{106}}%
\TextAt{13}{438.63}{234.123}{\Character{11}}%
\TextAt{16}{446.152}{234.122}{\Character{106}}%
\TextAt{9}{453.151}{234.123}{\Character{61}}%
\TextAt{22}{469.254}{238.123}{n}%
\TextAt{20}{467.213}{243.913}{\Character{112}}%
\Vector{q 0 G 0 g 0.478 w 407.168 -529.848 m 457.518 -529.848 l S Q}%
\TextAt{16}{479.168}{234.122}{\Character{106}}%
\TextAt{8}{484.481}{234.123}{N}%
\TextAt{21}{493.256}{232.263}{L}%
\TextAt{3}{499.015}{232.263}{\Character{106}}%
\TextAt{21}{501.368}{232.263}{K}%
\TextAt{9}{509.57}{234.122}{\Character{40}}%
\TextAt{13}{514.123}{234.123}{\Character{11}}%
\TextAt{9}{521.645}{234.122}{\Character{41}}%
\TextAt{16}{526.197}{234.123}{\Character{106}}%
\TextAt{8}{529.518}{234.122}{,}%
\TextAt{8}{79.226}{217.077}{the}%
\TextAt{8}{100.308}{217.077}{conditions}%
\TextAt{16}{157.231}{217.077}{\Character{106}}%
\TextAt{13}{160.552}{217.077}{\Character{11}}%
\TextAt{16}{168.074}{217.077}{\Character{106}}%
\TextAt{9}{176.293}{217.077}{\Character{61}}%
\TextAt{9}{190.281}{217.077}{0}%
\NativeArrowAt{201.046}{217.077}{11.955}{21.926}{\Longleftrightarrow}%
\TextAt{13}{227.861}{217.077}{\Character{11}}%
\TextAt{9}{240.281}{217.077}{\Character{61}}%
\TextAt{9}{254.269}{217.077}{0}%
\TextAt{8}{264.959}{217.077}{and}%
\TextAt{16}{288.647}{217.077}{\Character{106}}%
\TextAt{13}{291.968}{217.077}{\Character{11}\Character{12}}%
\TextAt{16}{306.761}{217.077}{\Character{106}}%
\TextAt{9}{314.98}{217.077}{\Character{61}}%
\TextAt{16}{328.982}{217.077}{\Character{106}}%
\TextAt{13}{332.303}{217.077}{\Character{11}}%
\TextAt{16}{339.825}{217.077}{\Character{106}\Character{106}}%
\TextAt{13}{346.466}{217.077}{\Character{12}}%
\TextAt{16}{353.738}{217.077}{\Character{106}}%
\TextAt{8}{361.881}{217.077}{are}%
\TextAt{8}{382.318}{217.077}{obvious.}%
\TextAt{8}{431.99}{217.077}{The}%
\TextAt{8}{456.969}{217.077}{strong}%
\TextAt{8}{493.736}{217.077}{triangle}%
\TextAt{8}{79.226}{200.031}{inequality}%
\TextAt{16}{246.566}{182.985}{\Character{106}}%
\TextAt{13}{249.887}{182.985}{\Character{11}}%
\TextAt{9}{260.066}{182.985}{\Character{43}}%
\TextAt{13}{271.827}{182.985}{\Character{12}}%
\TextAt{16}{279.098}{182.985}{\Character{106}}%
\NativeMathAt{285.721}{182.985}{11.955}{\leq}%
\TextAt{8}{298.359}{182.985}{max}%
\TextAt{16}{320.153}{182.985}{\Character{102}\Character{106}}%
\TextAt{13}{329.452}{182.985}{\Character{11}}%
\TextAt{16}{336.974}{182.985}{\Character{106}}%
\TextAt{13}{340.295}{182.985}{\Character{59}}%
\TextAt{16}{345.539}{182.985}{\Character{106}}%
\TextAt{13}{348.86}{182.985}{\Character{12}}%
\TextAt{16}{356.131}{182.985}{\Character{106}\Character{103}}%
\TextAt{8}{79.226}{160.443}{reduces,}%
\TextAt{8}{124.154}{160.443}{after}%
\TextAt{8}{151.794}{160.443}{dividing}%
\TextAt{8}{196.997}{160.443}{by}%
\TextAt{13}{213.256}{160.443}{\Character{11}}%
\TextAt{8}{224.675}{160.443}{or}%
\TextAt{13}{238.985}{160.443}{\Character{12}}%
\TextAt{8}{246.256}{160.443}{,}%
\TextAt{8}{253.405}{160.443}{to}%
\TextAt{8}{267.716}{160.443}{the}%
\TextAt{8}{287.872}{160.443}{implication}%
\TextAt{16}{245.425}{133.728}{\Character{106}}%
\TextAt{13}{248.746}{133.728}{\Character{11}}%
\TextAt{16}{256.268}{133.728}{\Character{106}}%
\NativeMathAt{262.891}{133.728}{11.955}{\leq}%
\TextAt{9}{275.529}{133.728}{1}%
\NativeArrowAt{284.687}{133.728}{11.955}{19.083}{\Longrightarrow}%
\TextAt{16}{307.082}{133.728}{\Character{106}}%
\TextAt{13}{310.412}{133.728}{\Character{11}}%
\TextAt{9}{320.59}{133.728}{\Character{43}}%
\TextAt{9}{332.342}{133.728}{1}%
\TextAt{16}{338.205}{133.728}{\Character{106}}%
\NativeMathAt{344.828}{133.728}{11.955}{\leq}%
\TextAt{9}{357.466}{133.728}{1}%
\TextAt{13}{363.319}{133.728}{\Character{59}}%
\TextAt{8}{79.226}{107.013}{and}%
\TextAt{8}{102.881}{107.013}{then,}%
\TextAt{8}{133.908}{107.013}{by}%
\TextAt{8}{151.06}{107.013}{(*),}%
\TextAt{8}{174.292}{107.013}{to}%
\TextAt{13}{189.495}{107.013}{\Character{11}}%
\NativeMathAt{201.858}{107.013}{11.955}{\in}%
\TextAt{16}{214.65}{107.013}{\Character{79}}%
\NativeArrowAt{229.36}{107.013}{11.955}{19.068}{\Longrightarrow}%
\TextAt{13}{253.269}{107.013}{\Character{11}}%
\TextAt{9}{264.056}{107.013}{\Character{43}}%
\TextAt{9}{276.418}{107.013}{1}%
\NativeMathAt{287.12}{107.013}{11.955}{\in}%
\TextAt{16}{299.912}{107.013}{\Character{79}}%
\TextAt{8}{309.781}{107.013}{,}%
\TextAt{8}{318.046}{107.013}{which}%
\TextAt{8}{352.424}{107.013}{is}%
\TextAt{8}{365.081}{107.013}{trivially}%
\TextAt{8}{410.207}{107.013}{true.}%
\TextAt{8}{442.146}{107.013}{Thus}%
\TextAt{8}{472.687}{107.013}{the}%
\TextAt{8}{493.736}{107.013}{formula}%
\TextAt{16}{79.226}{89.967}{\Character{106}}%
\TextAt{13}{82.547}{89.967}{\Character{11}}%
\TextAt{16}{90.069}{89.967}{\Character{106}}%
\TextAt{9}{96.71}{89.967}{\Character{61}}%
\TextAt{22}{112.457}{93.968}{n}%
\TextAt{20}{110.416}{99.757}{\Character{112}}%
\Vector{q 0 G 0 g 0.478 w 50.371 -674.004 m 100.721 -674.004 l S Q}%
\TextAt{16}{122.371}{89.967}{\Character{106}}%
\TextAt{8}{127.684}{89.967}{N}%
\TextAt{21}{136.46}{88.108}{L}%
\TextAt{3}{142.219}{88.108}{\Character{106}}%
\TextAt{21}{144.571}{88.108}{K}%
\TextAt{9}{152.774}{89.967}{\Character{40}}%
\TextAt{13}{157.326}{89.967}{\Character{11}}%
\TextAt{9}{164.848}{89.967}{\Character{41}}%
\TextAt{16}{169.4}{89.967}{\Character{106}}%
\TextAt{8}{176.497}{89.967}{does}%
\TextAt{8}{202.772}{89.967}{define}%
\TextAt{8}{236.459}{89.967}{a}%
\TextAt{8}{246.093}{89.967}{valuation}%
\TextAt{8}{297.044}{89.967}{of}%
\TextAt{13}{310.252}{89.967}{L}%
\TextAt{8}{321.992}{89.967}{and,}%
\TextAt{8}{347.909}{89.967}{restricted}%
\TextAt{8}{399.876}{89.967}{to}%
\TextAt{13}{414.065}{89.967}{K}%
\TextAt{8}{424.882}{89.967}{,}%
\TextAt{8}{431.934}{89.967}{it}%
\TextAt{8}{443.517}{89.967}{clearly}%
\TextAt{8}{480.791}{89.967}{gives}%
\TextAt{8}{509.349}{89.967}{back}%
\TextAt{8}{79.226}{72.921}{the}%
\TextAt{8}{99.383}{72.921}{given}%
\TextAt{8}{129.952}{72.921}{valuation.}%
\TextAt{8}{185.575}{72.921}{Equally}%
\TextAt{8}{228.662}{72.921}{obviously}%
\TextAt{8}{280.44}{72.921}{it}%
\TextAt{8}{292.144}{72.921}{has}%
\TextAt{16}{313.017}{72.921}{\Character{79}}%
\TextAt{8}{326.765}{72.921}{as}%
\TextAt{8}{341.135}{72.921}{its}%
\TextAt{8}{357.454}{72.921}{valuation}%
\TextAt{8}{408.526}{72.921}{ring.}%
\TextAt{8}{297.211}{43.033}{120}%
\Destination{522.36}{365.405}{nav.AMS.16}%
\Destination{234.48}{365.405}{nav.equation.2.4.1}%
\Destination{78.23}{809.815}{nav.page.120}%
\Destination{79.36}{407.155}{nav.print.132}%
\Destination{79.36}{599.755}{nav.number.2.4.8}%
\Hotlink{403.171}{319.455}{5.22}{11.047}{nav.chapter.1}%
\Hotlink{420.707}{319.455}{6.858}{11.047}{nav.section.1.2}%
\Hotlink{448.464}{319.455}{12.716}{11.047}{nav.print.12}%
\Hotlink{138.422}{231.955}{25.078}{11.047}{nav.number.2.4.7}%
}

% Page 133 (body page 121)
\Page{%
\TextAt{2}{79.226}{796.943}{II.}%
\TextAt{2}{94.073}{796.943}{The}%
\TextAt{2}{116.492}{796.943}{Theory}%
\TextAt{2}{154.401}{796.943}{of}%
\TextAt{2}{166.826}{796.943}{Valuations}%
\Vector{q 0 G 0 g 0.299 w 7.226 18.729 m 460.769 18.729 l S Q}%
\TextAt{12}{97.159}{761.124}{Uniqueness}%
\TextAt{12}{166.889}{761.124}{of}%
\TextAt{12}{180.981}{761.124}{the}%
\TextAt{12}{203.07}{761.124}{extended}%
\TextAt{12}{259.47}{761.124}{valuation:}%
\TextAt{8}{321.694}{761.124}{let}%
\TextAt{16}{337.525}{761.124}{\Character{106}}%
\TextAt{16}{343.658}{761.124}{\Character{106}}%
\TextAt{3}{346.99}{765.462}{\Character{48}}%
\TextAt{8}{352.609}{761.124}{be}%
\TextAt{8}{367.459}{761.124}{another}%
\TextAt{8}{409.316}{761.124}{extension}%
\TextAt{8}{460.008}{761.124}{with}%
\TextAt{8}{485.594}{761.124}{valuation}%
\TextAt{8}{79.226}{744.078}{ring}%
\TextAt{16}{103.718}{744.078}{\Character{79}}%
\TextAt{3}{113.567}{748.416}{\Character{48}}%
\TextAt{8}{116.362}{744.078}{.}%
\TextAt{8}{126.088}{744.078}{Let}%
\TextAt{14}{147.483}{744.078}{P}%
\TextAt{8}{157.376}{744.078}{,}%
\TextAt{8}{165.058}{744.078}{resp.}%
\TextAt{14}{195.657}{744.078}{P}%
\TextAt{3}{205.55}{748.416}{\Character{48}}%
\TextAt{8}{208.345}{744.078}{,}%
\TextAt{8}{216.026}{744.078}{be}%
\TextAt{8}{232.376}{744.078}{the}%
\TextAt{8}{252.959}{744.078}{maximal}%
\TextAt{8}{301.193}{744.078}{ideal}%
\TextAt{8}{329.582}{744.078}{of}%
\TextAt{16}{343.338}{744.078}{\Character{79}}%
\TextAt{8}{353.187}{744.078}{,}%
\TextAt{8}{360.869}{744.078}{resp.}%
\TextAt{16}{391.468}{744.078}{\Character{79}}%
\TextAt{3}{401.318}{748.416}{\Character{48}}%
\TextAt{8}{404.113}{744.078}{.}%
\TextAt{8}{413.838}{744.078}{We}%
\TextAt{8}{434.408}{744.078}{show}%
\TextAt{8}{463.837}{744.078}{that}%
\TextAt{16}{489.632}{744.078}{\Character{79}}%
\NativeMathAt{503.524}{744.078}{11.955}{\subseteq}%
\TextAt{16}{516.854}{744.078}{\Character{79}}%
\TextAt{3}{526.723}{748.416}{\Character{48}}%
\TextAt{8}{529.518}{744.078}{.}%
\TextAt{8}{79.226}{727.032}{Let}%
\TextAt{13}{100.196}{727.032}{\Character{11}}%
\NativeMathAt{111.038}{727.032}{11.955}{\in}%
\TextAt{16}{122.312}{727.032}{\Character{79}}%
\TextAt{16}{134.817}{727.032}{\Character{110}}%
\TextAt{16}{143.449}{727.032}{\Character{79}}%
\TextAt{3}{153.32}{731.37}{\Character{48}}%
\TextAt{8}{160.012}{727.032}{and}%
\TextAt{8}{182.775}{727.032}{let}%
\TextAt{13}{230.266}{709.986}{f}%
\TextAt{9}{237.312}{709.986}{\Character{40}}%
\TextAt{13}{241.864}{709.986}{x}%
\TextAt{9}{248.516}{709.986}{\Character{41}}%
\TextAt{9}{256.371}{709.986}{\Character{61}}%
\TextAt{13}{268.815}{709.986}{x}%
\TextAt{21}{275.467}{714.922}{d}%
\TextAt{9}{282.979}{709.986}{\Character{43}}%
\TextAt{13}{294.741}{709.986}{a}%
\TextAt{17}{300.885}{708.193}{1}%
\TextAt{13}{305.618}{709.986}{x}%
\TextAt{21}{312.27}{714.922}{d}%
\TextAt{3}{316.627}{714.922}{\Character{0}}%
\TextAt{17}{323.214}{714.922}{1}%
\TextAt{9}{330.603}{709.986}{\Character{43}}%
\TextAt{16}{342.364}{709.986}{\Character{1}}%
\TextAt{16}{347.66}{709.986}{\Character{1}}%
\TextAt{16}{352.968}{709.986}{\Character{1}}%
\TextAt{9}{358.968}{709.986}{\Character{43}}%
\TextAt{13}{370.73}{709.986}{a}%
\TextAt{21}{376.875}{708.193}{d}%
\TextAt{8}{79.226}{688.903}{be}%
\TextAt{8}{94.57}{688.903}{the}%
\TextAt{8}{114.146}{688.903}{minimal}%
\TextAt{8}{159.091}{688.903}{polynomial}%
\TextAt{8}{219.004}{688.903}{of}%
\TextAt{13}{231.754}{688.903}{\Character{11}}%
\TextAt{8}{242.593}{688.903}{over}%
\TextAt{13}{267.059}{688.903}{K}%
\TextAt{8}{277.876}{688.903}{.}%
\TextAt{8}{286.131}{688.903}{Then}%
\TextAt{8}{316.108}{688.903}{one}%
\TextAt{8}{336.988}{688.903}{has}%
\TextAt{13}{357.281}{688.903}{a}%
\TextAt{17}{363.426}{687.109}{1}%
\TextAt{13}{368.158}{688.903}{\Character{59}}%
\TextAt{13}{373.382}{688.903}{\Character{58}}%
\TextAt{13}{378.619}{688.903}{\Character{58}}%
\TextAt{13}{383.855}{688.903}{\Character{58}}%
\TextAt{13}{389.08}{688.903}{\Character{59}}%
\TextAt{13}{394.316}{688.903}{a}%
\TextAt{21}{400.524}{687.109}{d}%
\NativeMathAt{408.7}{688.903}{11.955}{\in}%
\TextAt{14}{419.991}{688.903}{o}%
\TextAt{8}{429.153}{688.903}{and}%
\TextAt{13}{451.335}{688.903}{\Character{11}}%
\TextAt{3}{458.857}{693.241}{\Character{0}}%
\TextAt{17}{465.443}{693.241}{1}%
\NativeMathAt{473.497}{688.903}{11.955}{\in}%
\TextAt{14}{484.788}{688.903}{P}%
\TextAt{3}{494.681}{693.241}{\Character{48}}%
\TextAt{8}{497.476}{688.903}{,}%
\TextAt{8}{504.161}{688.903}{hence}%
\TextAt{9}{79.226}{671.857}{1}%
\TextAt{9}{88.384}{671.857}{\Character{61}}%
\TextAt{16}{100.825}{671.857}{\Character{0}}%
\TextAt{13}{110.124}{671.857}{a}%
\TextAt{17}{116.269}{670.063}{1}%
\TextAt{13}{121.001}{671.857}{\Character{11}}%
\TextAt{3}{128.523}{676.195}{\Character{0}}%
\TextAt{17}{135.109}{676.195}{1}%
\TextAt{16}{142.363}{671.857}{\Character{0}}%
\TextAt{16}{154.163}{671.857}{\Character{1}}%
\TextAt{16}{159.471}{671.857}{\Character{1}}%
\TextAt{16}{164.779}{671.857}{\Character{1}}%
\TextAt{16}{170.613}{671.857}{\Character{0}}%
\TextAt{13}{182.473}{671.857}{a}%
\TextAt{21}{188.618}{670.063}{d}%
\TextAt{9}{193.473}{671.857}{\Character{40}}%
\TextAt{13}{198.026}{671.857}{\Character{11}}%
\TextAt{3}{205.547}{676.195}{\Character{0}}%
\TextAt{17}{212.134}{676.195}{1}%
\TextAt{9}{216.866}{671.857}{\Character{41}}%
\TextAt{21}{221.419}{676.195}{d}%
\NativeMathAt{229.595}{671.857}{11.955}{\in}%
\TextAt{14}{240.886}{671.857}{P}%
\TextAt{3}{250.779}{676.195}{\Character{48}}%
\TextAt{8}{253.574}{671.857}{,}%
\TextAt{8}{260.671}{671.857}{a}%
\TextAt{8}{270.36}{671.857}{contradiction.}%
\TextAt{8}{346.118}{671.857}{This}%
\TextAt{8}{372.772}{671.857}{shows}%
\TextAt{8}{406.323}{671.857}{the}%
\TextAt{8}{426.414}{671.857}{inclusion}%
\TextAt{16}{475.185}{671.857}{\Character{79}}%
\NativeMathAt{488.348}{671.857}{11.955}{\subseteq}%
\TextAt{16}{500.96}{671.857}{\Character{79}}%
\TextAt{3}{510.824}{676.195}{\Character{48}}%
\TextAt{8}{513.619}{671.857}{.}%
\TextAt{8}{522.046}{671.857}{In}%
\TextAt{8}{79.226}{654.811}{other}%
\TextAt{8}{110.443}{654.811}{words,}%
\TextAt{8}{148.062}{654.811}{we}%
\TextAt{8}{165.936}{654.811}{have}%
\TextAt{8}{193.579}{654.811}{that}%
\TextAt{16}{219.595}{654.811}{\Character{106}}%
\TextAt{13}{222.916}{654.811}{\Character{11}}%
\TextAt{16}{230.437}{654.811}{\Character{106}}%
\NativeMathAt{238.172}{654.811}{11.955}{\leq}%
\TextAt{9}{251.904}{654.811}{1}%
\NativeArrowAt{262.173}{654.811}{11.955}{19.075}{\Longrightarrow}%
\TextAt{16}{285.672}{654.811}{\Character{106}}%
\TextAt{13}{288.993}{654.811}{\Character{11}}%
\TextAt{16}{296.515}{654.811}{\Character{106}}%
\TextAt{3}{299.836}{659.149}{\Character{48}}%
\NativeMathAt{307.055}{654.811}{11.955}{\leq}%
\TextAt{9}{320.777}{654.811}{1}%
\TextAt{8}{331.175}{654.811}{and}%
\TextAt{8}{354.585}{654.811}{this}%
\TextAt{8}{378.054}{654.811}{implies}%
\TextAt{8}{418.429}{654.811}{that}%
\TextAt{8}{444.445}{654.811}{the}%
\TextAt{8}{465.249}{654.811}{valuations}%
\TextAt{16}{521.583}{654.811}{\Character{106}}%
\TextAt{16}{529.438}{654.811}{\Character{106}}%
\TextAt{8}{79.226}{637.765}{and}%
\TextAt{16}{101.717}{637.765}{\Character{106}}%
\TextAt{16}{108.651}{637.765}{\Character{106}}%
\TextAt{3}{111.985}{642.103}{\Character{48}}%
\TextAt{8}{118.406}{637.765}{are}%
\TextAt{8}{137.645}{637.765}{equivalent.}%
\TextAt{8}{197.721}{637.765}{For}%
\TextAt{8}{218.419}{637.765}{if}%
\TextAt{8}{228.871}{637.765}{they}%
\TextAt{8}{254.937}{637.765}{were}%
\TextAt{8}{281.648}{637.765}{not,}%
\TextAt{8}{305.497}{637.765}{then}%
\TextAt{8}{331.885}{637.765}{the}%
\TextAt{8}{351.77}{637.765}{approximation}%
\TextAt{8}{429.566}{637.765}{theorem}%
\TextAt{8}{474.819}{637.765}{(3.4)}%
\TextAt{8}{502.523}{637.765}{would}%
\TextAt{8}{79.226}{620.719}{allow}%
\TextAt{8}{110.417}{620.719}{us}%
\TextAt{8}{126.376}{620.719}{to}%
\TextAt{8}{141.63}{620.719}{find}%
\TextAt{8}{165.982}{620.719}{an}%
\TextAt{13}{183.185}{620.719}{\Character{11}}%
\NativeMathAt{195.636}{620.719}{11.955}{\in}%
\TextAt{13}{208.536}{620.719}{L}%
\TextAt{8}{221.341}{620.719}{such}%
\TextAt{8}{248.682}{620.719}{that}%
\TextAt{16}{274.995}{620.719}{\Character{106}}%
\TextAt{13}{278.316}{620.719}{\Character{11}}%
\TextAt{16}{285.838}{620.719}{\Character{106}}%
\TextAt{13}{294.088}{620.719}{\Character{60}}%
\TextAt{9}{308.122}{620.719}{1}%
\TextAt{8}{318.816}{620.719}{and}%
\TextAt{16}{342.523}{620.719}{\Character{106}}%
\TextAt{13}{345.843}{620.719}{\Character{11}}%
\TextAt{16}{353.365}{620.719}{\Character{106}}%
\TextAt{3}{356.686}{625.057}{\Character{48}}%
\TextAt{13}{364.41}{620.719}{\Character{62}}%
\TextAt{9}{378.444}{620.719}{1}%
\TextAt{8}{384.297}{620.719}{.}%
\TextAt{8}{395.577}{620.719}{Thus}%
\TextAt{16}{426.17}{620.719}{\Character{106}}%
\TextAt{16}{434.312}{620.719}{\Character{106}}%
\TextAt{8}{442.494}{620.719}{and}%
\TextAt{16}{466.201}{620.719}{\Character{106}}%
\TextAt{16}{474.343}{620.719}{\Character{106}}%
\TextAt{3}{477.684}{625.057}{\Character{48}}%
\TextAt{8}{485.32}{620.719}{are}%
\TextAt{8}{505.775}{620.719}{equal}%
\TextAt{8}{79.226}{603.673}{because}%
\TextAt{8}{122.137}{603.673}{they}%
\TextAt{8}{148.083}{603.673}{agree}%
\TextAt{8}{178.261}{603.673}{on}%
\TextAt{13}{194.129}{603.673}{K}%
\TextAt{8}{204.946}{603.673}{.}%
\TextAt{8}{213.264}{603.673}{The}%
\TextAt{8}{236.927}{603.673}{fact}%
\TextAt{8}{259.621}{603.673}{that}%
\TextAt{13}{284.599}{603.673}{L}%
\TextAt{8}{296.07}{603.673}{is}%
\TextAt{8}{307.442}{603.673}{again}%
\TextAt{8}{338.278}{603.673}{complete}%
\TextAt{8}{387.31}{603.673}{with}%
\TextAt{8}{413.578}{603.673}{respect}%
\TextAt{8}{453.237}{603.673}{to}%
\TextAt{8}{467.156}{603.673}{the}%
\TextAt{8}{486.921}{603.673}{extended}%
\TextAt{8}{79.226}{586.627}{valuation}%
\TextAt{8}{130.299}{586.627}{is}%
\TextAt{8}{142.063}{586.627}{deduced}%
\TextAt{8}{187.576}{586.627}{from}%
\TextAt{8}{215.216}{586.627}{the}%
\TextAt{8}{235.373}{586.627}{following}%
\TextAt{8}{284.807}{586.627}{general}%
\TextAt{8}{325.132}{586.627}{result.}%
\Vector{q 0 G 0 g 0.399 w 452.334 -187.373 m 452.334 -179.303 l S Q}%
\Vector{q 0 G 0 g 0.399 w 452.534 -179.503 m 459.707 -179.503 l S Q}%
\Vector{q 0 G 0 g 0.399 w 452.534 -187.174 m 459.707 -187.174 l S Q}%
\Vector{q 0 G 0 g 0.399 w 459.906 -187.373 m 459.906 -179.303 l S Q}%
\TextAt{12}{79.226}{554.899}{Proposition}%
\TextAt{12}{152.848}{554.899}{(4.9).}%
\TextAt{8}{188.955}{554.899}{Let}%
\TextAt{13}{210.716}{554.899}{K}%
\TextAt{8}{226.222}{554.899}{be}%
\TextAt{8}{242.938}{554.899}{complete}%
\TextAt{8}{293.152}{554.899}{with}%
\TextAt{8}{320.604}{554.899}{respect}%
\TextAt{8}{361.445}{554.899}{to}%
\TextAt{8}{376.547}{554.899}{the}%
\TextAt{8}{397.495}{554.899}{valuation}%
\TextAt{16}{449.358}{554.899}{\Character{106}}%
\TextAt{16}{457.356}{554.899}{\Character{106}}%
\TextAt{8}{465.378}{554.899}{and}%
\TextAt{8}{488.932}{554.899}{let}%
\TextAt{13}{506.628}{554.899}{V}%
\TextAt{8}{520.743}{554.899}{be}%
\TextAt{8}{79.226}{537.853}{an}%
\TextAt{13}{96.205}{537.853}{n}%
\TextAt{8}{103.193}{537.853}{-dimensional}%
\TextAt{8}{172.26}{537.853}{normed}%
\TextAt{8}{215.253}{537.853}{vector}%
\TextAt{8}{251.097}{537.853}{space}%
\TextAt{8}{283.092}{537.853}{over}%
\TextAt{13}{308.857}{537.853}{K}%
\TextAt{8}{319.675}{537.853}{.}%
\TextAt{8}{330.283}{537.853}{Then,}%
\TextAt{8}{364.991}{537.853}{for}%
\TextAt{8}{383.596}{537.853}{any}%
\TextAt{8}{406.433}{537.853}{basis}%
\TextAt{13}{435.893}{537.853}{v}%
\TextAt{17}{441.552}{536.06}{1}%
\TextAt{13}{446.284}{537.853}{\Character{59}}%
\TextAt{13}{451.508}{537.853}{\Character{58}}%
\TextAt{13}{456.745}{537.853}{\Character{58}}%
\TextAt{13}{461.981}{537.853}{\Character{58}}%
\TextAt{13}{467.206}{537.853}{\Character{59}}%
\TextAt{13}{472.442}{537.853}{v}%
\TextAt{21}{478.165}{536.06}{n}%
\TextAt{8}{488.418}{537.853}{of}%
\TextAt{13}{502.468}{537.853}{V}%
\TextAt{8}{516.51}{537.853}{the}%
\TextAt{8}{79.226}{520.807}{maximum}%
\TextAt{8}{133.861}{520.807}{norm}%
\TextAt{16}{199.612}{503.761}{\Character{107}}%
\TextAt{13}{205.59}{503.761}{x}%
\TextAt{17}{212.242}{501.968}{1}%
\TextAt{13}{216.974}{503.761}{v}%
\TextAt{17}{222.633}{501.968}{1}%
\TextAt{9}{230.022}{503.761}{\Character{43}}%
\TextAt{16}{241.784}{503.761}{\Character{1}}%
\TextAt{16}{247.08}{503.761}{\Character{1}}%
\TextAt{16}{252.388}{503.761}{\Character{1}}%
\TextAt{9}{258.388}{503.761}{\Character{43}}%
\TextAt{13}{270.149}{503.761}{x}%
\TextAt{21}{276.801}{501.968}{n}%
\TextAt{13}{282.438}{503.761}{v}%
\TextAt{21}{288.097}{501.968}{n}%
\TextAt{16}{293.733}{503.761}{\Character{107}}%
\TextAt{9}{303.032}{503.761}{\Character{61}}%
\TextAt{8}{315.457}{503.761}{max}%
\TextAt{16}{337.251}{503.761}{\Character{102}\Character{106}}%
\TextAt{13}{346.55}{503.761}{x}%
\TextAt{17}{353.202}{501.968}{1}%
\TextAt{16}{357.934}{503.761}{\Character{106}}%
\TextAt{13}{361.255}{503.761}{\Character{59}}%
\TextAt{13}{366.479}{503.761}{\Character{58}}%
\TextAt{13}{371.716}{503.761}{\Character{58}}%
\TextAt{13}{376.952}{503.761}{\Character{58}}%
\TextAt{13}{382.177}{503.761}{\Character{59}}%
\TextAt{16}{387.476}{503.761}{\Character{106}}%
\TextAt{13}{390.797}{503.761}{x}%
\TextAt{21}{397.449}{501.968}{n}%
\TextAt{16}{403.085}{503.761}{\Character{106}\Character{103}}%
\TextAt{8}{79.226}{482.678}{is}%
\TextAt{8}{90.99}{482.678}{equivalent}%
\TextAt{8}{146.605}{482.678}{to}%
\TextAt{8}{160.916}{482.678}{the}%
\TextAt{8}{181.072}{482.678}{given}%
\TextAt{8}{211.642}{482.678}{norm}%
\TextAt{8}{242.211}{482.678}{on}%
\TextAt{13}{258.47}{482.678}{V}%
\TextAt{8}{267.896}{482.678}{.}%
\TextAt{8}{276.344}{482.678}{In}%
\TextAt{8}{290.965}{482.678}{particular,}%
\TextAt{13}{348.207}{482.678}{V}%
\TextAt{8}{361.53}{482.678}{is}%
\TextAt{8}{373.294}{482.678}{complete}%
\TextAt{8}{422.716}{482.678}{and}%
\TextAt{8}{445.479}{482.678}{the}%
\TextAt{8}{465.635}{482.678}{isomorphism}%
\TextAt{13}{178.357}{459.367}{K}%
\TextAt{21}{189.174}{464.303}{n}%
\NativeArrowAt{198.131}{459.367}{11.955}{19.26}{\longrightarrow}%
\TextAt{13}{220.713}{459.367}{V}%
\TextAt{13}{228.173}{459.367}{\Character{59}}%
\TextAt{9}{257.343}{459.367}{\Character{40}}%
\TextAt{13}{261.895}{459.367}{x}%
\TextAt{17}{268.547}{457.574}{1}%
\TextAt{13}{273.279}{459.367}{\Character{59}}%
\TextAt{13}{278.503}{459.367}{\Character{58}}%
\TextAt{13}{283.74}{459.367}{\Character{58}}%
\TextAt{13}{288.976}{459.367}{\Character{58}}%
\TextAt{13}{294.201}{459.367}{\Character{59}}%
\TextAt{13}{299.437}{459.367}{x}%
\TextAt{21}{306.152}{457.574}{n}%
\TextAt{9}{311.789}{459.367}{\Character{41}}%
\NativeArrowAt{319.662}{459.367}{11.955}{19.26}{\longmapsto}%
\TextAt{13}{342.244}{459.367}{x}%
\TextAt{17}{348.896}{457.574}{1}%
\TextAt{13}{353.628}{459.367}{v}%
\TextAt{17}{359.287}{457.574}{1}%
\TextAt{9}{366.676}{459.367}{\Character{43}}%
\TextAt{16}{378.438}{459.367}{\Character{1}}%
\TextAt{16}{383.734}{459.367}{\Character{1}}%
\TextAt{16}{389.042}{459.367}{\Character{1}}%
\TextAt{9}{395.042}{459.367}{\Character{43}}%
\TextAt{13}{406.803}{459.367}{x}%
\TextAt{21}{413.455}{457.574}{n}%
\TextAt{13}{419.092}{459.367}{v}%
\TextAt{21}{424.751}{457.574}{n}%
\TextAt{13}{430.387}{459.367}{\Character{59}}%
\TextAt{8}{79.226}{436.056}{is}%
\TextAt{8}{90.99}{436.056}{a}%
\TextAt{8}{100.745}{436.056}{homeomorphism.}%
\TextAt{19}{79.226}{404.328}{Proof.}%
\TextAt{8}{116}{404.328}{Let}%
\TextAt{13}{137.422}{404.328}{v}%
\TextAt{17}{143.081}{402.535}{1}%
\TextAt{13}{147.813}{404.328}{\Character{59}}%
\TextAt{13}{153.037}{404.328}{\Character{58}}%
\TextAt{13}{158.274}{404.328}{\Character{58}}%
\TextAt{13}{163.51}{404.328}{\Character{58}}%
\TextAt{13}{168.735}{404.328}{\Character{59}}%
\TextAt{13}{173.971}{404.328}{v}%
\TextAt{21}{179.693}{402.535}{n}%
\TextAt{8}{189.679}{404.328}{be}%
\TextAt{8}{206.055}{404.328}{a}%
\TextAt{8}{216.262}{404.328}{basis}%
\TextAt{8}{245.455}{404.328}{and}%
\TextAt{16}{268.669}{404.328}{\Character{107}}%
\TextAt{16}{278.986}{404.328}{\Character{107}}%
\TextAt{8}{289.323}{404.328}{be}%
\TextAt{8}{305.699}{404.328}{the}%
\TextAt{8}{326.307}{404.328}{corresponding}%
\TextAt{8}{401.945}{404.328}{maximum}%
\TextAt{8}{457.032}{404.328}{norm}%
\TextAt{8}{488.054}{404.328}{on}%
\TextAt{13}{504.765}{404.328}{V}%
\TextAt{8}{514.19}{404.328}{.}%
\TextAt{8}{523.994}{404.328}{It}%
\TextAt{8}{79.226}{387.282}{suffices}%
\TextAt{8}{119.013}{387.282}{to}%
\TextAt{8}{133.323}{387.282}{show}%
\TextAt{8}{162.327}{387.282}{that,}%
\TextAt{8}{190.947}{387.282}{for}%
\TextAt{8}{208.832}{387.282}{every}%
\TextAt{8}{239.724}{387.282}{norm}%
\TextAt{16}{270.294}{387.282}{\Character{106}}%
\TextAt{16}{277.491}{387.282}{\Character{106}}%
\TextAt{8}{284.73}{387.282}{on}%
\TextAt{13}{300.989}{387.282}{V}%
\TextAt{8}{310.415}{387.282}{,}%
\TextAt{8}{317.564}{387.282}{there}%
\TextAt{8}{347.476}{387.282}{exist}%
\TextAt{8}{375.176}{387.282}{constants}%
\TextAt{13}{427.014}{387.282}{\Character{26}\Character{59}}%
\TextAt{13}{438.276}{387.282}{\Character{26}}%
\TextAt{3}{444.333}{391.621}{\Character{48}}%
\TextAt{13}{450.449}{387.282}{\Character{62}}%
\TextAt{9}{462.875}{387.282}{0}%
\TextAt{8}{472.625}{387.282}{such}%
\TextAt{8}{499.022}{387.282}{that}%
\TextAt{13}{212.346}{363.971}{\Character{26}}%
\TextAt{16}{218.384}{363.971}{\Character{107}}%
\TextAt{13}{224.361}{363.971}{x}%
\TextAt{16}{231.013}{363.971}{\Character{107}}%
\NativeMathAt{240.302}{363.971}{11.955}{\leq}%
\TextAt{16}{252.915}{363.971}{\Character{106}}%
\TextAt{13}{256.252}{363.971}{x}%
\TextAt{16}{262.904}{363.971}{\Character{106}}%
\NativeMathAt{269.527}{363.971}{11.955}{\leq}%
\TextAt{13}{282.165}{363.971}{\Character{26}}%
\TextAt{3}{288.203}{368.907}{\Character{48}}%
\TextAt{16}{290.998}{363.971}{\Character{107}}%
\TextAt{13}{296.976}{363.971}{x}%
\TextAt{16}{303.628}{363.971}{\Character{107}}%
\TextAt{8}{333.516}{363.971}{for}%
\TextAt{8}{351.401}{363.971}{all}%
\TextAt{13}{367.66}{363.971}{x}%
\NativeMathAt{377.632}{363.971}{11.955}{\in}%
\TextAt{13}{388.923}{363.971}{V}%
\TextAt{13}{396.383}{363.971}{\Character{58}}%
\TextAt{8}{79.226}{340.66}{Then}%
\TextAt{8}{110.377}{340.66}{the}%
\TextAt{8}{131.126}{340.66}{norm}%
\TextAt{16}{162.289}{340.66}{\Character{106}}%
\TextAt{16}{170.084}{340.66}{\Character{106}}%
\TextAt{8}{177.911}{340.66}{defines}%
\TextAt{8}{216.929}{340.66}{the}%
\TextAt{8}{237.678}{340.66}{same}%
\TextAt{8}{267.597}{340.66}{topology}%
\TextAt{8}{316.334}{340.66}{on}%
\TextAt{13}{333.186}{340.66}{V}%
\TextAt{8}{347.102}{340.66}{as}%
\TextAt{8}{362.065}{340.66}{the}%
\TextAt{8}{382.815}{340.66}{norm}%
\TextAt{16}{413.977}{340.66}{\Character{107}}%
\TextAt{16}{424.438}{340.66}{\Character{107}}%
\TextAt{8}{430.423}{340.66}{,}%
\TextAt{8}{438.314}{340.66}{and}%
\TextAt{8}{461.669}{340.66}{we}%
\TextAt{8}{479.49}{340.66}{obtain}%
\TextAt{8}{516.51}{340.66}{the}%
\TextAt{8}{79.226}{323.614}{topological}%
\TextAt{8}{139.579}{323.614}{isomorphism}%
\TextAt{13}{208.827}{323.614}{K}%
\TextAt{21}{219.644}{327.952}{n}%
\NativeArrowAt{230.013}{323.614}{11.955}{19.26}{\longrightarrow}%
\TextAt{13}{254.006}{323.614}{V}%
\TextAt{8}{263.432}{323.614}{,}%
\TextAt{9}{271.616}{323.614}{\Character{40}}%
\TextAt{13}{276.169}{323.614}{x}%
\TextAt{17}{282.821}{321.82}{1}%
\TextAt{13}{287.553}{323.614}{\Character{59}}%
\TextAt{13}{292.777}{323.614}{\Character{58}}%
\TextAt{13}{298.014}{323.614}{\Character{58}}%
\TextAt{13}{303.25}{323.614}{\Character{58}}%
\TextAt{13}{308.475}{323.614}{\Character{59}}%
\TextAt{13}{313.711}{323.614}{x}%
\TextAt{21}{320.426}{321.82}{n}%
\TextAt{9}{326.062}{323.614}{\Character{41}}%
\NativeArrowAt{335.347}{323.614}{11.955}{19.26}{\longmapsto}%
\TextAt{13}{359.34}{323.614}{x}%
\TextAt{17}{365.992}{321.82}{1}%
\TextAt{13}{370.725}{323.614}{v}%
\TextAt{17}{376.384}{321.82}{1}%
\TextAt{9}{384.337}{323.614}{\Character{43}}%
\TextAt{16}{396.663}{323.614}{\Character{1}}%
\TextAt{16}{401.959}{323.614}{\Character{1}}%
\TextAt{16}{407.267}{323.614}{\Character{1}}%
\TextAt{9}{413.832}{323.614}{\Character{43}}%
\TextAt{13}{426.158}{323.614}{x}%
\TextAt{21}{432.81}{321.82}{n}%
\TextAt{13}{438.446}{323.614}{v}%
\TextAt{21}{444.106}{321.82}{n}%
\TextAt{8}{449.742}{323.614}{.}%
\TextAt{8}{460.675}{323.614}{In}%
\TextAt{8}{476.124}{323.614}{fact,}%
\TextAt{16}{503.497}{323.614}{\Character{107}}%
\TextAt{16}{514.197}{323.614}{\Character{107}}%
\TextAt{8}{524.903}{323.614}{is}%
\TextAt{8}{79.226}{306.568}{transformed}%
\TextAt{8}{143.847}{306.568}{into}%
\TextAt{8}{166.781}{306.568}{the}%
\TextAt{8}{186.128}{306.568}{maximum}%
\TextAt{8}{239.954}{306.568}{norm}%
\TextAt{8}{269.714}{306.568}{on}%
\TextAt{13}{285.164}{306.568}{K}%
\TextAt{21}{295.982}{310.906}{n}%
\TextAt{8}{301.618}{306.568}{.}%
\TextAt{8}{309.797}{306.568}{For}%
\TextAt{13}{329.957}{306.568}{\Character{26}}%
\TextAt{3}{335.994}{310.906}{\Character{48}}%
\TextAt{8}{341.877}{306.568}{we}%
\TextAt{8}{358.296}{306.568}{may}%
\TextAt{8}{382.855}{306.568}{obviously}%
\TextAt{8}{433.824}{306.568}{take}%
\TextAt{16}{458.383}{306.568}{\Character{106}}%
\TextAt{13}{461.704}{306.568}{v}%
\TextAt{17}{467.363}{304.774}{1}%
\TextAt{16}{472.096}{306.568}{\Character{106}}%
\TextAt{9}{476.418}{306.568}{\Character{43}}%
\TextAt{16}{486.525}{306.568}{\Character{1}}%
\TextAt{16}{491.821}{306.568}{\Character{1}}%
\TextAt{16}{497.129}{306.568}{\Character{1}}%
\TextAt{9}{501.474}{306.568}{\Character{43}}%
\TextAt{16}{511.58}{306.568}{\Character{106}}%
\TextAt{13}{514.901}{306.568}{v}%
\TextAt{21}{520.56}{304.774}{n}%
\TextAt{16}{526.197}{306.568}{\Character{106}}%
\TextAt{8}{529.518}{306.568}{.}%
\TextAt{8}{79.226}{289.522}{The}%
\TextAt{8}{103.002}{289.522}{existence}%
\TextAt{8}{152.529}{289.522}{of}%
\TextAt{13}{165.581}{289.522}{\Character{26}}%
\TextAt{8}{175.239}{289.522}{is}%
\TextAt{8}{186.725}{289.522}{proved}%
\TextAt{8}{224.5}{289.522}{by}%
\TextAt{8}{240.481}{289.522}{induction}%
\TextAt{8}{292.232}{289.522}{on}%
\TextAt{13}{308.213}{289.522}{n}%
\TextAt{8}{315.201}{289.522}{.}%
\TextAt{8}{323.557}{289.522}{For}%
\TextAt{13}{344.248}{289.522}{n}%
\TextAt{9}{354.557}{289.522}{\Character{61}}%
\TextAt{9}{366.966}{289.522}{1}%
\TextAt{8}{376.454}{289.522}{we}%
\TextAt{8}{393.404}{289.522}{may}%
\TextAt{8}{418.495}{289.522}{take}%
\TextAt{13}{443.586}{289.522}{\Character{26}}%
\TextAt{9}{452.944}{289.522}{\Character{61}}%
\TextAt{16}{465.37}{289.522}{\Character{106}}%
\TextAt{13}{468.691}{289.522}{v}%
\TextAt{17}{474.35}{287.728}{1}%
\TextAt{16}{479.082}{289.522}{\Character{106}}%
\TextAt{8}{482.403}{289.522}{.}%
\TextAt{8}{490.759}{289.522}{Suppose}%
\TextAt{8}{79.226}{272.476}{that}%
\TextAt{8}{104.595}{272.476}{everything}%
\TextAt{8}{162.159}{272.476}{is}%
\TextAt{8}{173.923}{272.476}{proved}%
\TextAt{8}{211.976}{272.476}{for}%
\TextAt{9}{229.861}{272.476}{\Character{40}}%
\TextAt{13}{234.414}{272.476}{n}%
\TextAt{16}{244.058}{272.476}{\Character{0}}%
\TextAt{9}{256.013}{272.476}{1\Character{41}}%
\TextAt{8}{266.418}{272.476}{-dimensional}%
\TextAt{8}{334.766}{272.476}{vector}%
\TextAt{8}{369.89}{272.476}{spaces.}%
\TextAt{8}{410.331}{272.476}{Let}%
\TextAt{13}{189.063}{249.164}{V}%
\TextAt{21}{195.888}{247.371}{i}%
\TextAt{9}{202.59}{249.164}{\Character{61}}%
\TextAt{13}{215.015}{249.164}{K}%
\TextAt{13}{225.823}{249.164}{v}%
\TextAt{17}{231.492}{247.371}{1}%
\TextAt{9}{238.881}{249.164}{\Character{43}}%
\TextAt{16}{250.642}{249.164}{\Character{1}}%
\TextAt{16}{255.938}{249.164}{\Character{1}}%
\TextAt{16}{261.246}{249.164}{\Character{1}}%
\TextAt{9}{267.247}{249.164}{\Character{43}}%
\TextAt{13}{279.008}{249.164}{K}%
\TextAt{13}{289.815}{249.164}{v}%
\TextAt{21}{295.485}{247.371}{i}%
\TextAt{3}{298.368}{247.371}{\Character{0}}%
\TextAt{17}{304.955}{247.371}{1}%
\TextAt{9}{312.344}{249.164}{\Character{43}}%
\TextAt{13}{324.105}{249.164}{K}%
\TextAt{13}{334.913}{249.164}{v}%
\TextAt{21}{340.582}{247.371}{i}%
\TextAt{17}{343.465}{247.371}{\Character{43}1}%
\TextAt{9}{357.441}{249.164}{\Character{43}}%
\TextAt{16}{369.202}{249.164}{\Character{1}}%
\TextAt{16}{374.498}{249.164}{\Character{1}}%
\TextAt{16}{379.806}{249.164}{\Character{1}}%
\TextAt{9}{385.806}{249.164}{\Character{43}}%
\TextAt{13}{397.568}{249.164}{K}%
\TextAt{13}{408.375}{249.164}{v}%
\TextAt{21}{414.044}{247.371}{n}%
\TextAt{13}{419.681}{249.164}{\Character{59}}%
\TextAt{8}{79.226}{225.853}{so}%
\TextAt{8}{95.091}{225.853}{that}%
\TextAt{13}{121.955}{225.853}{V}%
\TextAt{9}{137.25}{225.853}{\Character{61}}%
\TextAt{13}{152.223}{225.853}{V}%
\TextAt{21}{159.048}{224.06}{i}%
\TextAt{9}{166.105}{225.853}{\Character{43}}%
\TextAt{13}{178.885}{225.853}{K}%
\TextAt{13}{189.693}{225.853}{v}%
\TextAt{21}{195.362}{224.06}{i}%
\TextAt{8}{198.744}{225.853}{.}%
\TextAt{8}{211.678}{225.853}{Then}%
\TextAt{13}{243.73}{225.853}{V}%
\TextAt{21}{250.555}{224.06}{i}%
\TextAt{8}{259.328}{225.853}{is}%
\TextAt{8}{272.587}{225.853}{complete}%
\TextAt{8}{323.505}{225.853}{with}%
\TextAt{8}{351.66}{225.853}{respect}%
\TextAt{8}{393.205}{225.853}{to}%
\TextAt{8}{409.011}{225.853}{the}%
\TextAt{8}{430.663}{225.853}{restriction}%
\TextAt{8}{488.156}{225.853}{of}%
\TextAt{16}{502.981}{225.853}{\Character{106}}%
\TextAt{16}{511.684}{225.853}{\Character{106}}%
\TextAt{8}{520.408}{225.853}{by}%
\TextAt{8}{79.226}{208.807}{induction,}%
\TextAt{8}{135.228}{208.807}{hence}%
\TextAt{8}{168.312}{208.807}{it}%
\TextAt{8}{180.593}{208.807}{is}%
\TextAt{8}{192.934}{208.807}{closed}%
\TextAt{8}{228.037}{208.807}{in}%
\TextAt{13}{242.267}{208.807}{V}%
\TextAt{8}{251.693}{208.807}{.}%
\TextAt{8}{261.873}{208.807}{Thus}%
\TextAt{13}{292.099}{208.807}{V}%
\TextAt{21}{298.923}{207.014}{i}%
\TextAt{9}{305.354}{208.807}{\Character{43}}%
\TextAt{13}{317.509}{208.807}{v}%
\TextAt{21}{323.168}{207.014}{i}%
\TextAt{8}{331.024}{208.807}{is}%
\TextAt{8}{343.365}{208.807}{also}%
\TextAt{8}{367.422}{208.807}{closed.}%
\TextAt{8}{408.231}{208.807}{Since}%
\TextAt{9}{439.366}{208.807}{0}%
\NativeMathAt{449.523}{208.807}{11.955}{\notin}%
\TextAt{20}{461.797}{217.774}{\Character{83}}%
\TextAt{21}{471.76}{214.619}{n}%
\TextAt{21}{471.76}{205.32}{i}%
\TextAt{17}{474.643}{205.32}{\Character{61}1}%
\TextAt{9}{485.962}{208.807}{\Character{40}}%
\TextAt{13}{490.514}{208.807}{V}%
\TextAt{21}{497.339}{207.014}{i}%
\TextAt{9}{503.77}{208.807}{\Character{43}}%
\TextAt{13}{515.925}{208.807}{v}%
\TextAt{21}{521.584}{207.014}{i}%
\TextAt{9}{524.965}{208.807}{\Character{41}}%
\TextAt{8}{529.518}{208.807}{,}%
\TextAt{8}{79.226}{191.761}{there}%
\TextAt{8}{109.769}{191.761}{exists}%
\TextAt{8}{142.715}{191.761}{a}%
\TextAt{8}{153.101}{191.761}{neighbourhood}%
\TextAt{8}{233.736}{191.761}{of}%
\TextAt{8}{247.696}{191.761}{0}%
\TextAt{8}{258.083}{191.761}{which}%
\TextAt{8}{292.2}{191.761}{is}%
\TextAt{8}{304.595}{191.761}{disjoint}%
\TextAt{8}{346.913}{191.761}{from}%
\TextAt{20}{375.184}{200.728}{\Character{83}}%
\TextAt{21}{385.147}{197.573}{n}%
\TextAt{21}{385.147}{188.274}{i}%
\TextAt{17}{388.03}{188.274}{\Character{61}1}%
\TextAt{9}{399.349}{191.761}{\Character{40}}%
\TextAt{13}{403.901}{191.761}{V}%
\TextAt{21}{410.726}{189.968}{i}%
\TextAt{9}{417.194}{191.761}{\Character{43}}%
\TextAt{13}{429.385}{191.761}{v}%
\TextAt{21}{435.044}{189.968}{i}%
\TextAt{9}{438.426}{191.761}{\Character{41}}%
\TextAt{8}{442.978}{191.761}{,}%
\TextAt{8}{450.916}{191.761}{i.e.,}%
\TextAt{8}{473.809}{191.761}{there}%
\TextAt{8}{504.352}{191.761}{exists}%
\TextAt{13}{79.226}{174.715}{\Character{26}}%
\TextAt{13}{88.575}{174.715}{\Character{62}}%
\TextAt{9}{101.01}{174.715}{0}%
\TextAt{8}{110.761}{174.715}{such}%
\TextAt{8}{137.158}{174.715}{that}%
\TextAt{16}{174.288}{151.404}{\Character{106}}%
\TextAt{13}{177.609}{151.404}{w}%
\TextAt{21}{185.977}{149.611}{i}%
\TextAt{9}{192.015}{151.404}{\Character{43}}%
\TextAt{13}{203.776}{151.404}{v}%
\TextAt{21}{209.435}{149.611}{i}%
\TextAt{16}{212.816}{151.404}{\Character{106}}%
\NativeMathAt{219.439}{151.404}{11.955}{\geq}%
\TextAt{13}{232.077}{151.404}{\Character{26}}%
\TextAt{8}{262.025}{151.404}{for}%
\TextAt{8}{279.91}{151.404}{all}%
\TextAt{13}{296.169}{151.404}{w}%
\TextAt{21}{304.537}{149.611}{i}%
\NativeMathAt{311.239}{151.404}{11.955}{\in}%
\TextAt{13}{322.53}{151.404}{V}%
\TextAt{21}{329.355}{149.611}{i}%
\TextAt{8}{336.633}{151.404}{and}%
\TextAt{8}{359.396}{151.404}{all}%
\TextAt{13}{375.655}{151.404}{i}%
\TextAt{9}{382.969}{151.404}{\Character{61}}%
\TextAt{9}{395.379}{151.404}{1}%
\TextAt{13}{401.248}{151.404}{\Character{59}}%
\TextAt{13}{406.472}{151.404}{\Character{58}}%
\TextAt{13}{411.709}{151.404}{\Character{58}}%
\TextAt{13}{416.933}{151.404}{\Character{58}}%
\TextAt{13}{422.17}{151.404}{\Character{59}}%
\TextAt{13}{427.406}{151.404}{n\Character{58}}%
\TextAt{8}{79.226}{128.092}{For}%
\TextAt{13}{100.196}{128.092}{x}%
\TextAt{9}{110.168}{128.092}{\Character{61}}%
\TextAt{13}{122.594}{128.092}{x}%
\TextAt{17}{129.246}{126.299}{1}%
\TextAt{13}{133.978}{128.092}{v}%
\TextAt{17}{139.638}{126.299}{1}%
\TextAt{9}{147.027}{128.092}{\Character{43}}%
\TextAt{16}{158.788}{128.092}{\Character{1}}%
\TextAt{16}{164.084}{128.092}{\Character{1}}%
\TextAt{16}{169.392}{128.092}{\Character{1}}%
\TextAt{9}{175.392}{128.092}{\Character{43}}%
\TextAt{13}{187.154}{128.092}{x}%
\TextAt{21}{193.806}{126.299}{n}%
\TextAt{13}{199.442}{128.092}{v}%
\TextAt{21}{205.101}{126.299}{n}%
\NativeMathAt{214.058}{128.092}{11.955}{\neq}%
\TextAt{9}{226.468}{128.092}{0}%
\TextAt{8}{236.234}{128.092}{and}%
\TextAt{16}{258.997}{128.092}{\Character{106}}%
\TextAt{13}{262.318}{128.092}{x}%
\TextAt{21}{268.97}{126.299}{r}%
\TextAt{16}{273.525}{128.092}{\Character{106}}%
\TextAt{9}{280.166}{128.092}{\Character{61}}%
\TextAt{8}{292.592}{128.092}{max}%
\TextAt{16}{314.386}{128.092}{\Character{102}\Character{106}}%
\TextAt{13}{323.685}{128.092}{x}%
\TextAt{21}{330.337}{126.299}{i}%
\TextAt{16}{333.718}{128.092}{\Character{106}\Character{103}}%
\TextAt{8}{343.016}{128.092}{,}%
\TextAt{8}{350.166}{128.092}{one}%
\TextAt{8}{371.625}{128.092}{finds}%
\TextAt{16}{192.931}{100.507}{\Character{106}}%
\TextAt{13}{196.252}{100.507}{x}%
\TextAt{3}{202.904}{105.443}{\Character{0}}%
\TextAt{17}{209.49}{105.443}{1}%
\TextAt{21}{202.904}{97.551}{r}%
\TextAt{13}{214.223}{100.507}{x}%
\TextAt{16}{220.875}{100.507}{\Character{106}}%
\TextAt{9}{227.516}{100.507}{\Character{61}}%
\TextAt{20}{239.942}{117.842}{\Character{12}}%
\TextAt{20}{239.942}{110.669}{\Character{12}}%
\TextAt{20}{239.942}{103.496}{\Character{12}}%
\TextAt{20}{239.942}{96.323}{\Character{12}}%
\TextAt{13}{245.122}{108.595}{x}%
\TextAt{17}{251.775}{106.801}{1}%
\Vector{q 0 G 0 g 0.478 w 173.122 -670.504 m 184.507 -670.504 l S Q}%
\TextAt{13}{245.211}{92.306}{x}%
\TextAt{21}{251.863}{90.513}{r}%
\TextAt{13}{257.702}{100.507}{v}%
\TextAt{17}{263.362}{98.714}{1}%
\TextAt{9}{270.751}{100.507}{\Character{43}}%
\TextAt{16}{282.512}{100.507}{\Character{1}}%
\TextAt{16}{287.808}{100.507}{\Character{1}}%
\TextAt{16}{293.116}{100.507}{\Character{1}}%
\TextAt{9}{299.116}{100.507}{\Character{43}}%
\TextAt{13}{310.878}{100.507}{v}%
\TextAt{21}{316.537}{98.714}{r}%
\TextAt{9}{323.748}{100.507}{\Character{43}}%
\TextAt{16}{335.51}{100.507}{\Character{1}}%
\TextAt{16}{340.806}{100.507}{\Character{1}}%
\TextAt{16}{346.114}{100.507}{\Character{1}}%
\TextAt{9}{352.114}{100.507}{\Character{43}}%
\TextAt{13}{365.071}{108.595}{x}%
\TextAt{21}{371.723}{106.801}{n}%
\Vector{q 0 G 0 g 0.478 w 293.071 -670.504 m 305.359 -670.504 l S Q}%
\TextAt{13}{365.612}{92.306}{x}%
\TextAt{21}{372.264}{90.513}{r}%
\TextAt{13}{378.555}{100.507}{v}%
\TextAt{21}{384.214}{98.714}{n}%
\TextAt{20}{389.85}{117.842}{\Character{12}}%
\TextAt{20}{389.85}{110.669}{\Character{12}}%
\TextAt{20}{389.85}{103.496}{\Character{12}}%
\TextAt{20}{389.85}{96.323}{\Character{12}}%
\NativeMathAt{397.156}{100.507}{11.955}{\geq}%
\TextAt{13}{409.775}{100.507}{\Character{26}\Character{59}}%
\TextAt{8}{79.226}{72.921}{so}%
\TextAt{8}{93.596}{72.921}{that}%
\TextAt{8}{118.965}{72.921}{one}%
\TextAt{8}{140.425}{72.921}{has}%
\TextAt{16}{161.298}{72.921}{\Character{106}}%
\TextAt{13}{164.619}{72.921}{x}%
\TextAt{16}{171.271}{72.921}{\Character{106}}%
\NativeMathAt{177.894}{72.921}{11.955}{\geq}%
\TextAt{13}{190.532}{72.921}{\Character{26}}%
\TextAt{16}{196.57}{72.921}{\Character{106}}%
\TextAt{13}{199.891}{72.921}{x}%
\TextAt{21}{206.543}{71.128}{r}%
\TextAt{16}{211.098}{72.921}{\Character{106}}%
\TextAt{9}{217.74}{72.921}{\Character{61}}%
\TextAt{13}{230.165}{72.921}{\Character{26}}%
\TextAt{16}{236.203}{72.921}{\Character{107}}%
\TextAt{13}{242.18}{72.921}{x}%
\TextAt{16}{248.833}{72.921}{\Character{107}}%
\TextAt{8}{254.81}{72.921}{.}%
\Vector{q 0 G 0 g 0.399 w 452.334 -701.079 m 452.334 -693.009 l S Q}%
\Vector{q 0 G 0 g 0.399 w 452.534 -693.208 m 459.707 -693.208 l S Q}%
\Vector{q 0 G 0 g 0.399 w 452.534 -700.879 m 459.707 -700.879 l S Q}%
\Vector{q 0 G 0 g 0.399 w 459.906 -701.079 m 459.906 -693.009 l S Q}%
\TextAt{8}{297.211}{43.033}{121}%
\Destination{78.23}{809.815}{nav.page.121}%
\Destination{79.36}{497.655}{nav.print.133}%
\Destination{79.36}{568.855}{nav.number.2.4.9}%
\Hotlink{474.35}{635.6}{23.35}{11.047}{nav.number.2.3.4}%
}

% Page 134 (body page 122)
\Page{%
\TextAt{2}{79.226}{796.943}{II.}%
\TextAt{2}{94.073}{796.943}{The}%
\TextAt{2}{116.492}{796.943}{Theory}%
\TextAt{2}{154.401}{796.943}{of}%
\TextAt{2}{166.826}{796.943}{Valuations}%
\Vector{q 0 G 0 g 0.299 w 7.226 18.729 m 460.769 18.729 l S Q}%
\TextAt{8}{97.159}{761.124}{The}%
\TextAt{8}{120.858}{761.124}{fact}%
\TextAt{8}{143.588}{761.124}{that}%
\TextAt{8}{168.601}{761.124}{an}%
\TextAt{8}{184.505}{761.124}{exponential}%
\TextAt{8}{246.915}{761.124}{valuation}%
\TextAt{13}{297.632}{761.124}{v}%
\TextAt{8}{307.262}{761.124}{on}%
\TextAt{13}{323.166}{761.124}{K}%
\TextAt{8}{337.526}{761.124}{associated}%
\TextAt{8}{392.906}{761.124}{with}%
\TextAt{16}{419.211}{761.124}{\Character{106}}%
\TextAt{16}{426.061}{761.124}{\Character{106}}%
\TextAt{8}{432.937}{761.124}{extends}%
\TextAt{8}{475.238}{761.124}{uniquely}%
\TextAt{8}{522.357}{761.124}{to}%
\TextAt{13}{79.226}{744.078}{L}%
\TextAt{8}{91.088}{744.078}{is}%
\TextAt{8}{102.852}{744.078}{a}%
\TextAt{8}{112.607}{744.078}{trivial}%
\TextAt{8}{147.409}{744.078}{consequence}%
\TextAt{8}{213.473}{744.078}{of}%
\TextAt{8}{226.803}{744.078}{theorem}%
\TextAt{8}{272.328}{744.078}{(4.8).}%
\TextAt{8}{304.854}{744.078}{The}%
\TextAt{8}{328.908}{744.078}{extension}%
\TextAt{13}{380.674}{744.078}{w}%
\TextAt{8}{393.261}{744.078}{is}%
\TextAt{8}{405.025}{744.078}{given}%
\TextAt{8}{435.594}{744.078}{by}%
\TextAt{8}{451.853}{744.078}{the}%
\TextAt{8}{472.01}{744.078}{formula}%
\TextAt{13}{252.323}{713.267}{w}%
\TextAt{9}{261.013}{713.267}{\Character{40}}%
\TextAt{13}{265.565}{713.267}{\Character{11}}%
\TextAt{9}{273.087}{713.267}{\Character{41}}%
\TextAt{9}{280.942}{713.267}{\Character{61}}%
\TextAt{9}{295.148}{721.355}{1}%
\Vector{q 0 G 0 g 0.478 w 222.581 -57.744 m 229.569 -57.744 l S Q}%
\TextAt{13}{294.581}{705.066}{n}%
\TextAt{13}{302.764}{713.267}{v}%
\TextAt{9}{308.852}{713.267}{\Character{40}}%
\TextAt{8}{313.405}{713.267}{N}%
\TextAt{21}{322.18}{711.407}{L}%
\TextAt{3}{327.939}{711.407}{\Character{106}}%
\TextAt{21}{330.291}{711.407}{K}%
\TextAt{9}{338.494}{713.267}{\Character{40}}%
\TextAt{13}{343.046}{713.267}{\Character{11}}%
\TextAt{9}{350.568}{713.267}{\Character{41}\Character{41}}%
\TextAt{8}{79.226}{684.082}{if}%
\TextAt{13}{89.95}{684.082}{n}%
\TextAt{9}{100.258}{684.082}{\Character{61}}%
\TextAt{9}{112.668}{684.082}{\Character{91}}%
\TextAt{13}{115.936}{684.082}{L}%
\TextAt{9}{127.221}{684.082}{\Character{58}}%
\TextAt{13}{133.793}{684.082}{K}%
\TextAt{9}{144.611}{684.082}{\Character{93}}%
\TextAt{13}{151.184}{684.082}{\Character{60}}%
\NativeMathAt{163.609}{684.082}{11.955}{\infty}%
\TextAt{8}{175.564}{684.082}{.}%
\ErrataAt{79.22}{658.355}{453.64}{\textbf{Exercise II.4.1.} An infinite separable algebraic extension of a nontrivially valued}%
\ErrataAt{79.22}{642.555}{453.64}{complete field $K$ is never complete.}%
\TextAt{12}{79.226}{609.834}{Exercise}%
\TextAt{12}{133.686}{609.834}{II.4.2.}%
\TextAt{8}{173.754}{609.834}{Let}%
\TextAt{13}{195.819}{609.834}{X}%
\TextAt{17}{205.534}{608.04}{0}%
\TextAt{13}{210.266}{609.834}{\Character{59}}%
\TextAt{13}{215.49}{609.834}{X}%
\TextAt{17}{225.225}{608.04}{1}%
\TextAt{13}{229.958}{609.834}{\Character{59}}%
\TextAt{13}{235.182}{609.834}{\Character{58}}%
\TextAt{13}{240.419}{609.834}{\Character{58}}%
\TextAt{13}{245.643}{609.834}{\Character{58}}%
\TextAt{8}{253.934}{609.834}{be}%
\TextAt{8}{270.953}{609.834}{an}%
\TextAt{8}{288.307}{609.834}{infinite}%
\TextAt{8}{329.069}{609.834}{sequence}%
\TextAt{8}{378.665}{609.834}{of}%
\TextAt{8}{393.09}{609.834}{unknowns,}%
\TextAt{13}{452.405}{609.834}{p}%
\TextAt{8}{463.272}{609.834}{a}%
\TextAt{8}{474.123}{609.834}{fixed}%
\TextAt{8}{503.503}{609.834}{prime}%
\TextAt{8}{79.226}{592.788}{number}%
\TextAt{8}{121.612}{592.788}{and}%
\TextAt{13}{144.163}{592.788}{W}%
\TextAt{21}{155.215}{590.994}{n}%
\TextAt{9}{164.172}{592.788}{\Character{61}}%
\TextAt{13}{176.598}{592.788}{X}%
\TextAt{21}{187.253}{598.455}{p}%
\TextAt{22}{191.516}{601.268}{n}%
\TextAt{17}{186.313}{589.857}{0}%
\TextAt{9}{199.338}{592.788}{\Character{43}}%
\TextAt{13}{210.668}{592.788}{pX}%
\TextAt{21}{227.199}{598.455}{p}%
\TextAt{22}{231.461}{601.268}{n}%
\TextAt{23}{236.062}{601.268}{\Character{0}}%
\TextAt{24}{241.818}{601.268}{1}%
\TextAt{17}{226.259}{589.857}{1}%
\TextAt{9}{248.693}{592.788}{\Character{43}}%
\TextAt{16}{260.023}{592.788}{\Character{1}}%
\TextAt{16}{265.319}{592.788}{\Character{1}}%
\TextAt{16}{270.627}{592.788}{\Character{1}}%
\TextAt{9}{276.197}{592.788}{\Character{43}}%
\TextAt{13}{287.527}{592.788}{p}%
\TextAt{21}{293.402}{597.126}{n}%
\TextAt{13}{299.038}{592.788}{X}%
\TextAt{21}{308.754}{590.995}{n}%
\TextAt{8}{314.39}{592.788}{,}%
\TextAt{13}{321.37}{592.788}{n}%
\NativeMathAt{331.679}{592.788}{11.955}{\geq}%
\TextAt{9}{344.298}{592.788}{0}%
\TextAt{8}{350.151}{592.788}{.}%
\TextAt{8}{358.529}{592.788}{Show}%
\TextAt{8}{389.211}{592.788}{that}%
\TextAt{8}{414.369}{592.788}{there}%
\TextAt{8}{444.07}{592.788}{exist}%
\TextAt{8}{471.559}{592.788}{polynomials}%
\TextAt{13}{79.226}{575.742}{S}%
\TextAt{17}{86.425}{573.949}{0}%
\TextAt{13}{91.158}{575.742}{\Character{59}}%
\TextAt{13}{96.382}{575.742}{S}%
\TextAt{17}{103.601}{573.949}{1}%
\TextAt{13}{108.334}{575.742}{\Character{59}}%
\TextAt{13}{113.558}{575.742}{\Character{58}}%
\TextAt{13}{118.795}{575.742}{\Character{58}}%
\TextAt{13}{124.019}{575.742}{\Character{58}}%
\TextAt{8}{127.318}{575.742}{;}%
\TextAt{13}{134.467}{575.742}{P}%
\TextAt{17}{142.012}{573.949}{0}%
\TextAt{13}{146.745}{575.742}{\Character{59}}%
\TextAt{13}{151.969}{575.742}{P}%
\TextAt{17}{159.534}{573.949}{1}%
\TextAt{13}{164.266}{575.742}{\Character{59}}%
\TextAt{13}{169.49}{575.742}{\Character{58}}%
\TextAt{13}{174.727}{575.742}{\Character{58}}%
\TextAt{13}{179.963}{575.742}{\Character{58}}%
\NativeMathAt{186.571}{575.742}{11.955}{\in}%
\TextAt{18}{197.862}{575.742}{Z}%
\TextAt{9}{205.832}{575.742}{\Character{91}}%
\TextAt{13}{209.084}{575.742}{X}%
\TextAt{17}{218.799}{573.949}{0}%
\TextAt{13}{223.532}{575.742}{\Character{59}}%
\TextAt{13}{228.756}{575.742}{X}%
\TextAt{17}{238.491}{573.949}{1}%
\TextAt{13}{243.223}{575.742}{\Character{59}}%
\TextAt{13}{248.447}{575.742}{\Character{58}}%
\TextAt{13}{253.684}{575.742}{\Character{58}}%
\TextAt{13}{258.92}{575.742}{\Character{58}}%
\TextAt{9}{264.2}{575.742}{\Character{59}}%
\TextAt{13}{269.444}{575.742}{Y}%
\TextAt{17}{276.236}{573.949}{0}%
\TextAt{13}{280.968}{575.742}{\Character{59}}%
\TextAt{13}{286.192}{575.742}{Y}%
\TextAt{17}{293.004}{573.949}{1}%
\TextAt{13}{297.737}{575.742}{\Character{59}}%
\TextAt{13}{302.961}{575.742}{\Character{58}}%
\TextAt{13}{308.198}{575.742}{\Character{58}}%
\TextAt{13}{313.422}{575.742}{\Character{58}}%
\TextAt{9}{316.721}{575.742}{\Character{93}}%
\TextAt{8}{323.87}{575.742}{such}%
\TextAt{8}{350.267}{575.742}{that}%
\TextAt{13}{176.355}{547.674}{W}%
\TextAt{21}{187.407}{545.881}{n}%
\TextAt{9}{193.043}{547.674}{\Character{40}}%
\TextAt{13}{197.596}{547.674}{S}%
\TextAt{17}{204.795}{545.881}{0}%
\TextAt{13}{209.527}{547.674}{\Character{59}}%
\TextAt{13}{214.751}{547.674}{S}%
\TextAt{17}{221.971}{545.881}{1}%
\TextAt{13}{226.703}{547.674}{\Character{59}}%
\TextAt{13}{231.927}{547.674}{\Character{58}}%
\TextAt{13}{237.164}{547.674}{\Character{58}}%
\TextAt{13}{242.388}{547.674}{\Character{58}}%
\TextAt{9}{245.687}{547.674}{\Character{41}}%
\TextAt{9}{253.542}{547.674}{\Character{61}}%
\TextAt{13}{265.986}{547.674}{W}%
\TextAt{21}{277.037}{545.881}{n}%
\TextAt{9}{282.674}{547.674}{\Character{40}}%
\TextAt{13}{287.226}{547.674}{X}%
\TextAt{17}{296.941}{545.881}{0}%
\TextAt{13}{301.673}{547.674}{\Character{59}}%
\TextAt{13}{306.897}{547.674}{X}%
\TextAt{17}{316.633}{545.881}{1}%
\TextAt{13}{321.365}{547.674}{\Character{59}}%
\TextAt{13}{326.589}{547.674}{\Character{58}}%
\TextAt{13}{331.826}{547.674}{\Character{58}}%
\TextAt{13}{337.062}{547.674}{\Character{58}}%
\TextAt{9}{340.349}{547.674}{\Character{41}}%
\TextAt{9}{347.546}{547.674}{\Character{43}}%
\TextAt{13}{359.32}{547.674}{W}%
\TextAt{21}{370.371}{545.881}{n}%
\TextAt{9}{376.007}{547.674}{\Character{40}}%
\TextAt{13}{380.56}{547.674}{Y}%
\TextAt{17}{387.352}{545.881}{0}%
\TextAt{13}{392.084}{547.674}{\Character{59}}%
\TextAt{13}{397.308}{547.674}{Y}%
\TextAt{17}{404.12}{545.881}{1}%
\TextAt{13}{408.852}{547.674}{\Character{59}}%
\TextAt{13}{414.076}{547.674}{\Character{58}}%
\TextAt{13}{419.313}{547.674}{\Character{58}}%
\TextAt{13}{424.549}{547.674}{\Character{58}}%
\TextAt{9}{427.836}{547.674}{\Character{41}}%
\TextAt{13}{432.389}{547.674}{\Character{59}}%
\TextAt{13}{178.901}{519.606}{W}%
\TextAt{21}{189.953}{517.813}{n}%
\TextAt{9}{195.589}{519.606}{\Character{40}}%
\TextAt{13}{200.142}{519.606}{P}%
\TextAt{17}{207.687}{517.813}{0}%
\TextAt{13}{212.419}{519.606}{\Character{59}}%
\TextAt{13}{217.643}{519.606}{P}%
\TextAt{17}{225.209}{517.813}{1}%
\TextAt{13}{229.941}{519.606}{\Character{59}}%
\TextAt{13}{235.165}{519.606}{\Character{58}}%
\TextAt{13}{240.402}{519.606}{\Character{58}}%
\TextAt{13}{245.626}{519.606}{\Character{58}}%
\TextAt{9}{248.925}{519.606}{\Character{41}}%
\TextAt{9}{256.78}{519.606}{\Character{61}}%
\TextAt{13}{269.224}{519.606}{W}%
\TextAt{21}{280.275}{517.813}{n}%
\TextAt{9}{285.911}{519.606}{\Character{40}}%
\TextAt{13}{290.464}{519.606}{X}%
\TextAt{17}{300.179}{517.813}{0}%
\TextAt{13}{304.911}{519.606}{\Character{59}}%
\TextAt{13}{310.135}{519.606}{X}%
\TextAt{17}{319.871}{517.813}{1}%
\TextAt{13}{324.603}{519.606}{\Character{59}}%
\TextAt{13}{329.827}{519.606}{\Character{58}}%
\TextAt{13}{335.064}{519.606}{\Character{58}}%
\TextAt{13}{340.288}{519.606}{\Character{58}}%
\TextAt{9}{343.587}{519.606}{\Character{41}}%
\TextAt{16}{350.796}{519.606}{\Character{1}}%
\TextAt{13}{356.774}{519.606}{W}%
\TextAt{21}{367.825}{517.813}{n}%
\TextAt{9}{373.462}{519.606}{\Character{40}}%
\TextAt{13}{378.014}{519.606}{Y}%
\TextAt{17}{384.806}{517.813}{0}%
\TextAt{13}{389.538}{519.606}{\Character{59}}%
\TextAt{13}{394.762}{519.606}{Y}%
\TextAt{17}{401.574}{517.813}{1}%
\TextAt{13}{406.306}{519.606}{\Character{59}}%
\TextAt{13}{411.53}{519.606}{\Character{58}}%
\TextAt{13}{416.767}{519.606}{\Character{58}}%
\TextAt{13}{422.003}{519.606}{\Character{58}}%
\TextAt{9}{425.291}{519.606}{\Character{41}}%
\TextAt{13}{429.843}{519.606}{\Character{58}}%
\TextAt{12}{79.226}{482.074}{Exercise}%
\TextAt{12}{133.724}{482.074}{II.4.3.}%
\TextAt{8}{173.826}{482.074}{Let}%
\TextAt{13}{195.924}{482.074}{A}%
\TextAt{8}{209.725}{482.074}{be}%
\TextAt{8}{226.777}{482.074}{a}%
\TextAt{8}{237.661}{482.074}{commutative}%
\TextAt{8}{308.715}{482.074}{ring.}%
\TextAt{8}{340.716}{482.074}{For}%
\TextAt{13}{362.814}{482.074}{a}%
\TextAt{9}{374.202}{482.074}{\Character{61}}%
\TextAt{9}{388.536}{482.074}{\Character{40}}%
\TextAt{13}{393.103}{482.074}{a}%
\TextAt{17}{399.248}{480.281}{0}%
\TextAt{13}{403.98}{482.074}{\Character{59}}%
\TextAt{13}{409.204}{482.074}{a}%
\TextAt{17}{415.369}{480.281}{1}%
\TextAt{13}{420.101}{482.074}{\Character{59}}%
\TextAt{13}{425.325}{482.074}{\Character{58}}%
\TextAt{13}{430.562}{482.074}{\Character{58}}%
\TextAt{13}{435.798}{482.074}{\Character{58}}%
\TextAt{9}{439.086}{482.074}{\Character{41}}%
\TextAt{8}{443.638}{482.074}{,}%
\TextAt{13}{452.197}{482.074}{b}%
\TextAt{9}{462.418}{482.074}{\Character{61}}%
\TextAt{9}{476.752}{482.074}{\Character{40}}%
\TextAt{13}{481.318}{482.074}{b}%
\TextAt{17}{486.295}{480.281}{0}%
\TextAt{13}{491.028}{482.074}{\Character{59}}%
\TextAt{13}{496.252}{482.074}{b}%
\TextAt{17}{501.249}{480.281}{1}%
\TextAt{13}{505.981}{482.074}{\Character{59}}%
\TextAt{13}{511.205}{482.074}{\Character{58}}%
\TextAt{13}{516.442}{482.074}{\Character{58}}%
\TextAt{13}{521.678}{482.074}{\Character{58}}%
\TextAt{9}{524.965}{482.074}{\Character{41}}%
\TextAt{8}{529.518}{482.074}{,}%
\TextAt{13}{79.226}{465.028}{a}%
\TextAt{21}{85.371}{463.235}{i}%
\TextAt{13}{88.752}{465.028}{\Character{59}}%
\TextAt{13}{93.976}{465.028}{b}%
\TextAt{21}{98.974}{463.235}{i}%
\NativeMathAt{105.676}{465.028}{11.955}{\in}%
\TextAt{13}{116.967}{465.028}{A}%
\TextAt{8}{125.742}{465.028}{,}%
\TextAt{8}{132.891}{465.028}{put}%
\TextAt{13}{142.916}{436.961}{a}%
\TextAt{9}{151.718}{436.961}{\Character{43}}%
\TextAt{13}{163.479}{436.961}{b}%
\TextAt{9}{171.777}{436.961}{\Character{61}}%
\TextAt{9}{184.187}{436.961}{\Character{40}}%
\TextAt{13}{188.755}{436.961}{S}%
\TextAt{17}{195.954}{435.167}{0}%
\TextAt{9}{200.687}{436.961}{\Character{40}}%
\TextAt{13}{205.239}{436.961}{a\Character{59}}%
\TextAt{13}{216.608}{436.961}{b}%
\TextAt{9}{221.605}{436.961}{\Character{41}}%
\TextAt{13}{226.158}{436.961}{\Character{59}}%
\TextAt{13}{231.382}{436.961}{S}%
\TextAt{17}{238.601}{435.167}{1}%
\TextAt{9}{243.333}{436.961}{\Character{40}}%
\TextAt{13}{247.886}{436.961}{a\Character{59}}%
\TextAt{13}{259.255}{436.961}{b}%
\TextAt{9}{264.252}{436.961}{\Character{41}}%
\TextAt{13}{268.804}{436.961}{\Character{59}}%
\TextAt{13}{274.028}{436.961}{\Character{58}}%
\TextAt{13}{279.265}{436.961}{\Character{58}}%
\TextAt{13}{284.501}{436.961}{\Character{58}}%
\TextAt{9}{287.788}{436.961}{\Character{41}}%
\TextAt{13}{292.341}{436.961}{\Character{59}}%
\TextAt{13}{321.476}{436.961}{a}%
\TextAt{16}{330.297}{436.961}{\Character{1}}%
\TextAt{13}{336.274}{436.961}{b}%
\TextAt{9}{344.572}{436.961}{\Character{61}}%
\TextAt{9}{356.982}{436.961}{\Character{40}}%
\TextAt{13}{361.55}{436.961}{P}%
\TextAt{17}{369.095}{435.167}{0}%
\TextAt{9}{373.828}{436.961}{\Character{40}}%
\TextAt{13}{378.38}{436.961}{a\Character{59}}%
\TextAt{13}{389.749}{436.961}{b}%
\TextAt{9}{394.746}{436.961}{\Character{41}}%
\TextAt{13}{399.299}{436.961}{\Character{59}}%
\TextAt{13}{404.523}{436.961}{P}%
\TextAt{17}{412.088}{435.167}{1}%
\TextAt{9}{416.82}{436.961}{\Character{40}}%
\TextAt{13}{421.373}{436.961}{a\Character{59}}%
\TextAt{13}{432.742}{436.961}{b}%
\TextAt{9}{437.739}{436.961}{\Character{41}}%
\TextAt{13}{442.291}{436.961}{\Character{59}}%
\TextAt{13}{447.515}{436.961}{\Character{58}}%
\TextAt{13}{452.752}{436.961}{\Character{58}}%
\TextAt{13}{457.988}{436.961}{\Character{58}}%
\TextAt{9}{461.275}{436.961}{\Character{41}}%
\TextAt{13}{465.828}{436.961}{\Character{58}}%
\TextAt{8}{79.226}{408.893}{Show}%
\TextAt{8}{111.188}{408.893}{that}%
\TextAt{8}{137.626}{408.893}{with}%
\TextAt{8}{165.356}{408.893}{these}%
\TextAt{8}{196.397}{408.893}{operations}%
\TextAt{8}{254.445}{408.893}{the}%
\TextAt{8}{275.671}{408.893}{vectors}%
\TextAt{13}{316.48}{408.893}{a}%
\TextAt{9}{327.768}{408.893}{\Character{61}}%
\TextAt{9}{342.007}{408.893}{\Character{40}}%
\TextAt{13}{346.569}{408.893}{a}%
\TextAt{17}{352.714}{407.1}{0}%
\TextAt{13}{357.446}{408.893}{\Character{59}}%
\TextAt{13}{362.67}{408.893}{a}%
\TextAt{17}{368.835}{407.1}{1}%
\TextAt{13}{373.567}{408.893}{\Character{59}}%
\TextAt{13}{378.791}{408.893}{\Character{58}}%
\TextAt{13}{384.028}{408.893}{\Character{58}}%
\TextAt{13}{389.264}{408.893}{\Character{58}}%
\TextAt{9}{392.551}{408.893}{\Character{41}}%
\TextAt{8}{402.071}{408.893}{form}%
\TextAt{8}{430.781}{408.893}{a}%
\TextAt{8}{441.606}{408.893}{commutative}%
\TextAt{8}{512.601}{408.893}{ring}%
\TextAt{13}{79.226}{391.847}{W}%
\TextAt{9}{91.903}{391.847}{\Character{40}}%
\TextAt{13}{96.456}{391.847}{A}%
\TextAt{9}{105.231}{391.847}{\Character{41}}%
\TextAt{8}{113.681}{391.847}{with}%
\TextAt{8}{140.341}{391.847}{1.}%
\TextAt{8}{154.647}{391.847}{It}%
\TextAt{8}{167.32}{391.847}{is}%
\TextAt{8}{179.084}{391.847}{called}%
\TextAt{8}{212.247}{391.847}{the}%
\TextAt{12}{232.404}{391.847}{ring}%
\TextAt{12}{260.331}{391.847}{of}%
\TextAt{12}{275.657}{391.847}{Witt}%
\TextAt{12}{308.247}{391.847}{vectors}%
\TextAt{8}{353.761}{391.847}{over}%
\TextAt{13}{378.807}{391.847}{A}%
\TextAt{8}{387.582}{391.847}{.}%
\TextAt{12}{79.226}{360.39}{Exercise}%
\TextAt{12}{131.112}{360.39}{II.4.4.}%
\TextAt{8}{168.943}{360.39}{Assume}%
\TextAt{13}{211.162}{360.39}{pA}%
\TextAt{9}{229.133}{360.39}{\Character{61}}%
\TextAt{9}{241.543}{360.39}{0}%
\TextAt{8}{247.412}{360.39}{.}%
\TextAt{8}{255.479}{360.39}{For}%
\TextAt{8}{275.305}{360.39}{every}%
\TextAt{8}{305.055}{360.39}{Witt}%
\TextAt{8}{332.198}{360.39}{vector}%
\TextAt{13}{366.179}{360.39}{a}%
\TextAt{9}{375.645}{360.39}{\Character{61}}%
\TextAt{9}{388.055}{360.39}{\Character{40}}%
\TextAt{13}{392.623}{360.39}{a}%
\TextAt{17}{398.768}{358.597}{0}%
\TextAt{13}{403.5}{360.39}{\Character{59}}%
\TextAt{13}{408.724}{360.39}{a}%
\TextAt{17}{414.889}{358.597}{1}%
\TextAt{13}{419.622}{360.39}{\Character{59}}%
\TextAt{13}{424.846}{360.39}{\Character{58}}%
\TextAt{13}{430.083}{360.39}{\Character{58}}%
\TextAt{13}{435.307}{360.39}{\Character{58}}%
\TextAt{9}{438.606}{360.39}{\Character{41}}%
\NativeMathAt{446.479}{360.39}{11.955}{\in}%
\TextAt{13}{457.77}{360.39}{W}%
\TextAt{9}{470.447}{360.39}{\Character{40}}%
\TextAt{13}{475}{360.39}{A}%
\TextAt{9}{483.775}{360.39}{\Character{41}}%
\TextAt{8}{491.082}{360.39}{consider}%
\TextAt{8}{79.226}{343.345}{the}%
\TextAt{8}{99.383}{343.345}{“ghost}%
\TextAt{8}{136.193}{343.345}{components”}%
\TextAt{13}{201.205}{315.277}{a}%
\TextAt{17}{207.35}{320.213}{\Character{40}}%
\TextAt{21}{210.643}{320.213}{n}%
\TextAt{17}{215.781}{320.213}{\Character{41}}%
\TextAt{9}{222.893}{315.277}{\Character{61}}%
\TextAt{13}{235.319}{315.277}{W}%
\TextAt{21}{246.37}{313.484}{n}%
\TextAt{9}{252.007}{315.277}{\Character{40}}%
\TextAt{13}{256.559}{315.277}{a}%
\TextAt{9}{262.704}{315.277}{\Character{41}}%
\TextAt{9}{270.559}{315.277}{\Character{61}}%
\TextAt{13}{283.003}{315.277}{a}%
\TextAt{21}{289.148}{320.944}{p}%
\TextAt{22}{293.41}{323.757}{n}%
\TextAt{17}{289.148}{312.346}{0}%
\TextAt{9}{301.664}{315.277}{\Character{43}}%
\TextAt{13}{313.425}{315.277}{pa}%
\TextAt{21}{325.446}{320.944}{p}%
\TextAt{22}{329.708}{323.757}{n}%
\TextAt{23}{334.309}{323.757}{\Character{0}}%
\TextAt{24}{340.065}{323.757}{1}%
\TextAt{17}{325.445}{312.346}{1}%
\TextAt{9}{347.371}{315.277}{\Character{43}}%
\TextAt{16}{359.132}{315.277}{\Character{1}}%
\TextAt{16}{364.428}{315.277}{\Character{1}}%
\TextAt{16}{369.736}{315.277}{\Character{1}}%
\TextAt{9}{375.737}{315.277}{\Character{43}}%
\TextAt{13}{387.498}{315.277}{p}%
\TextAt{21}{393.373}{320.213}{n}%
\TextAt{13}{399.009}{315.277}{a}%
\TextAt{21}{405.154}{313.484}{n}%
\TextAt{8}{79.226}{287.209}{as}%
\TextAt{8}{93.596}{287.209}{well}%
\TextAt{8}{117.327}{287.209}{as}%
\TextAt{8}{131.697}{287.209}{the}%
\TextAt{8}{151.854}{287.209}{mappings}%
\TextAt{13}{204.6}{287.209}{V}%
\TextAt{13}{212.06}{287.209}{\Character{59}}%
\TextAt{13}{217.296}{287.209}{F}%
\TextAt{9}{229.843}{287.209}{\Character{58}}%
\TextAt{13}{236.416}{287.209}{W}%
\TextAt{9}{249.093}{287.209}{\Character{40}}%
\TextAt{13}{253.645}{287.209}{A}%
\TextAt{9}{262.421}{287.209}{\Character{41}}%
\NativeArrowAt{270.294}{287.209}{11.955}{19.26}{\longrightarrow}%
\TextAt{13}{292.876}{287.209}{W}%
\TextAt{9}{305.553}{287.209}{\Character{40}}%
\TextAt{13}{310.106}{287.209}{A}%
\TextAt{9}{318.881}{287.209}{\Character{41}}%
\TextAt{8}{327.331}{287.209}{defined}%
\TextAt{8}{367.644}{287.209}{by}%
\TextAt{13}{179.158}{259.142}{V}%
\TextAt{13}{188.567}{259.142}{a}%
\TextAt{9}{198.05}{259.142}{\Character{61}}%
\TextAt{9}{210.46}{259.142}{\Character{40}0}%
\TextAt{13}{220.881}{259.142}{\Character{59}}%
\TextAt{13}{226.105}{259.142}{a}%
\TextAt{17}{232.27}{257.348}{0}%
\TextAt{13}{237.002}{259.142}{\Character{59}}%
\TextAt{13}{242.226}{259.142}{a}%
\TextAt{17}{248.391}{257.348}{1}%
\TextAt{13}{253.124}{259.142}{\Character{59}}%
\TextAt{13}{258.348}{259.142}{\Character{58}}%
\TextAt{13}{263.585}{259.142}{\Character{58}}%
\TextAt{13}{268.809}{259.142}{\Character{58}}%
\TextAt{9}{272.108}{259.142}{\Character{41}}%
\TextAt{8}{300.57}{259.142}{and}%
\TextAt{13}{343.346}{259.142}{F}%
\TextAt{13}{352.528}{259.142}{a}%
\TextAt{9}{362.015}{259.142}{\Character{61}}%
\TextAt{9}{374.425}{259.142}{\Character{40}}%
\TextAt{13}{378.993}{259.142}{a}%
\TextAt{21}{385.138}{264.809}{p}%
\TextAt{17}{385.138}{256.21}{0}%
\TextAt{13}{389.899}{259.142}{\Character{59}}%
\TextAt{13}{395.123}{259.142}{a}%
\TextAt{21}{401.288}{264.809}{p}%
\TextAt{17}{401.288}{256.21}{1}%
\TextAt{13}{406.049}{259.142}{\Character{59}}%
\TextAt{13}{411.273}{259.142}{\Character{58}}%
\TextAt{13}{416.51}{259.142}{\Character{58}}%
\TextAt{13}{421.734}{259.142}{\Character{58}}%
\TextAt{9}{425.033}{259.142}{\Character{41}}%
\TextAt{13}{429.585}{259.142}{\Character{59}}%
\TextAt{8}{79.226}{231.074}{called}%
\TextAt{8}{112.39}{231.074}{respectively}%
\TextAt{8}{176.183}{231.074}{“transfer”}%
\TextAt{8}{230.543}{231.074}{(“Verschiebung”}%
\TextAt{8}{317.409}{231.074}{in}%
\TextAt{8}{331.062}{231.074}{German)}%
\TextAt{8}{380.568}{231.074}{and}%
\TextAt{8}{403.331}{231.074}{“Frobenius”.}%
\TextAt{8}{471.172}{231.074}{Show}%
\TextAt{8}{502.065}{231.074}{that}%
\ErrataAt{152.36}{203.055}{372}{$(Va)^{(n)}=p\,a^{(n-1)}\quad\text{and}\quad a^{(n)}=(Fa)^{(n-1)}+p^n a_n\quad(n\ge1).$}%
\TextAt{8}{297.211}{43.033}{122}%
\Destination{78.23}{809.815}{nav.page.122}%
\Destination{79.226}{672.355}{nav.ex.2.4.1}%
\Destination{79.226}{623.855}{nav.ex.2.4.2}%
\Destination{79.226}{496.055}{nav.ex.2.4.3}%
\Destination{79.226}{374.355}{nav.ex.2.4.4}%
\Destination{79.36}{673.355}{nav.print.134}%
\Hotlink{271.828}{741.855}{25.078}{11.047}{nav.number.2.4.8}%
\Destination{79.226}{672.355}{errata.14}%
\Destination{79.226}{217.055}{errata.15}%
}

% Page 135 (body page 123)
\Page{%
\TextAt{2}{79.226}{796.943}{II.}%
\TextAt{2}{94.073}{796.943}{The}%
\TextAt{2}{116.492}{796.943}{Theory}%
\TextAt{2}{154.401}{796.943}{of}%
\TextAt{2}{166.826}{796.943}{Valuations}%
\Vector{q 0 G 0 g 0.299 w 7.226 18.729 m 460.769 18.729 l S Q}%
\TextAt{12}{79.226}{761.124}{Exercise}%
\TextAt{12}{132.536}{761.124}{II.4.5.}%
\TextAt{8}{171.605}{761.124}{Let}%
\TextAt{13}{192.67}{761.124}{k}%
\TextAt{8}{203.152}{761.124}{be}%
\TextAt{8}{219.172}{761.124}{a}%
\TextAt{8}{229.023}{761.124}{field}%
\TextAt{8}{254.475}{761.124}{of}%
\TextAt{8}{267.9}{761.124}{characteristic}%
\TextAt{13}{339.93}{761.124}{p}%
\TextAt{8}{345.805}{761.124}{.}%
\TextAt{8}{354.539}{761.124}{Then}%
\TextAt{13}{385.192}{761.124}{V}%
\TextAt{8}{398.61}{761.124}{is}%
\TextAt{8}{410.469}{761.124}{a}%
\TextAt{8}{420.32}{761.124}{homomorphism}%
\TextAt{8}{503.085}{761.124}{of}%
\TextAt{8}{516.51}{761.124}{the}%
\TextAt{8}{79.226}{744.078}{additive}%
\TextAt{8}{124.106}{744.078}{group}%
\TextAt{8}{157.281}{744.078}{of}%
\TextAt{13}{170.611}{744.078}{W}%
\TextAt{9}{183.289}{744.078}{\Character{40}}%
\TextAt{13}{187.841}{744.078}{k}%
\TextAt{9}{194.331}{744.078}{\Character{41}}%
\TextAt{8}{202.78}{744.078}{and}%
\TextAt{13}{225.543}{744.078}{F}%
\TextAt{8}{238.644}{744.078}{is}%
\TextAt{8}{250.408}{744.078}{a}%
\TextAt{8}{260.163}{744.078}{ring}%
\TextAt{8}{284.229}{744.078}{homomorphism,}%
\TextAt{8}{370.151}{744.078}{and}%
\TextAt{8}{392.914}{744.078}{one}%
\TextAt{8}{414.373}{744.078}{has}%
\TextAt{13}{257.841}{715.706}{V}%
\TextAt{13}{267.25}{715.706}{F}%
\TextAt{13}{276.443}{715.706}{a}%
\TextAt{9}{285.936}{715.706}{\Character{61}}%
\TextAt{13}{298.362}{715.706}{F}%
\TextAt{13}{307.544}{715.706}{V}%
\TextAt{13}{316.964}{715.706}{a}%
\TextAt{9}{326.457}{715.706}{\Character{61}}%
\TextAt{13}{338.882}{715.706}{pa\Character{58}}%
\TextAt{12}{79.226}{672.751}{Exercise}%
\TextAt{12}{131.436}{672.751}{II.4.6.}%
\TextAt{8}{169.549}{672.751}{If}%
\TextAt{13}{180.38}{672.751}{k}%
\TextAt{8}{189.905}{672.751}{is}%
\TextAt{8}{200.808}{672.751}{a}%
\TextAt{8}{209.702}{672.751}{perfect}%
\TextAt{8}{247.85}{672.751}{field}%
\TextAt{8}{272.346}{672.751}{of}%
\TextAt{8}{284.815}{672.751}{characteristic}%
\TextAt{13}{355.888}{672.751}{p}%
\TextAt{8}{361.763}{672.751}{,}%
\TextAt{8}{368.223}{672.751}{then}%
\TextAt{13}{394.022}{672.751}{W}%
\TextAt{9}{406.699}{672.751}{\Character{40}}%
\TextAt{13}{411.252}{672.751}{k}%
\TextAt{9}{417.741}{672.751}{\Character{41}}%
\TextAt{8}{425.33}{672.751}{is}%
\TextAt{8}{436.232}{672.751}{a}%
\TextAt{8}{445.127}{672.751}{complete}%
\TextAt{8}{493.688}{672.751}{discrete}%
\TextAt{8}{79.226}{655.705}{valuation}%
\TextAt{8}{130.299}{655.705}{ring}%
\TextAt{8}{154.364}{655.705}{with}%
\TextAt{8}{181.024}{655.705}{residue}%
\TextAt{8}{220.751}{655.705}{class}%
\TextAt{8}{248.189}{655.705}{field}%
\TextAt{13}{273.545}{655.705}{k}%
\TextAt{8}{280.035}{655.705}{.}%
\TextAt{15}{79.226}{603.114}{§}%
\TextAt{15}{93.785}{603.114}{5}%
\TextAt{15}{120.693}{603.114}{Local}%
\TextAt{15}{172.425}{603.114}{Fields}%
\TextAt{8}{79.226}{571.228}{Among}%
\TextAt{8}{121.196}{571.228}{all}%
\TextAt{8}{138.777}{571.228}{complete}%
\TextAt{8}{189.521}{571.228}{(nonarchimedean)}%
\TextAt{8}{286.784}{571.228}{valued}%
\TextAt{8}{324.856}{571.228}{fields,}%
\TextAt{8}{359.732}{571.228}{those}%
\TextAt{8}{391.683}{571.228}{arising}%
\TextAt{8}{430.795}{571.228}{as}%
\TextAt{8}{446.487}{571.228}{completions}%
\TextAt{8}{512.259}{571.228}{of}%
\TextAt{8}{526.911}{571.228}{a}%
\TextAt{12}{79.226}{554.182}{global}%
\TextAt{12}{117.323}{554.182}{field}%
\TextAt{8}{142.142}{554.182}{,}%
\TextAt{8}{148.356}{554.182}{i.e.,}%
\TextAt{8}{169.527}{554.182}{of}%
\TextAt{8}{181.689}{554.182}{a}%
\TextAt{8}{190.276}{554.182}{finite}%
\TextAt{8}{219.02}{554.182}{extension}%
\TextAt{8}{269.618}{554.182}{of}%
\TextAt{8}{281.78}{554.182}{either}%
\TextAt{18}{313.776}{554.182}{Q}%
\TextAt{8}{325.804}{554.182}{or}%
\TextAt{18}{338.946}{554.182}{F}%
\TextAt{21}{346.252}{552.389}{p}%
\TextAt{9}{351.013}{554.182}{\Character{40}}%
\TextAt{13}{355.565}{554.182}{t}%
\TextAt{9}{359.792}{554.182}{\Character{41}}%
\TextAt{8}{364.345}{554.182}{,}%
\TextAt{8}{370.559}{554.182}{have}%
\TextAt{8}{396.386}{554.182}{the}%
\TextAt{8}{415.375}{554.182}{most}%
\TextAt{8}{442.887}{554.182}{eminent}%
\TextAt{8}{486.264}{554.182}{relevance}%
\TextAt{8}{79.226}{537.136}{for}%
\TextAt{8}{96.517}{537.136}{number}%
\TextAt{8}{138.52}{537.136}{theory.}%
\TextAt{8}{178.643}{537.136}{The}%
\TextAt{8}{202.103}{537.136}{valuation}%
\TextAt{8}{252.582}{537.136}{on}%
\TextAt{8}{268.248}{537.136}{such}%
\TextAt{8}{294.051}{537.136}{a}%
\TextAt{8}{303.213}{537.136}{completion}%
\TextAt{8}{362.455}{537.136}{is}%
\TextAt{8}{373.625}{537.136}{discrete}%
\TextAt{8}{416.011}{537.136}{and}%
\TextAt{8}{438.18}{537.136}{has}%
\TextAt{8}{458.46}{537.136}{a}%
\TextAt{8}{467.622}{537.136}{finite}%
\TextAt{8}{496.94}{537.136}{residue}%
\TextAt{8}{79.226}{520.09}{class}%
\TextAt{8}{107.703}{520.09}{field,}%
\TextAt{8}{137.611}{520.09}{as}%
\TextAt{8}{153.02}{520.09}{we}%
\TextAt{8}{171.287}{520.09}{shall}%
\TextAt{8}{199.703}{520.09}{see}%
\TextAt{8}{219.656}{520.09}{shortly.}%
\TextAt{8}{265.761}{520.09}{In}%
\TextAt{8}{281.421}{520.09}{contrast}%
\TextAt{8}{327.735}{520.09}{to}%
\TextAt{8}{343.084}{520.09}{the}%
\TextAt{8}{364.28}{520.09}{global}%
\TextAt{8}{399.798}{520.09}{fields,}%
\TextAt{8}{434.321}{520.09}{all}%
\TextAt{8}{451.619}{520.09}{fields}%
\TextAt{8}{482.63}{520.09}{which}%
\TextAt{8}{517.156}{520.09}{are}%
\TextAt{8}{79.226}{503.044}{complete}%
\TextAt{8}{127.887}{503.044}{with}%
\TextAt{8}{153.785}{503.044}{respect}%
\TextAt{8}{193.073}{503.044}{to}%
\TextAt{8}{206.621}{503.044}{a}%
\TextAt{8}{215.615}{503.044}{discrete}%
\TextAt{8}{257.832}{503.044}{valuation}%
\TextAt{8}{308.142}{503.044}{and}%
\TextAt{8}{330.143}{503.044}{have}%
\TextAt{8}{356.376}{503.044}{a}%
\TextAt{8}{365.369}{503.044}{finite}%
\TextAt{8}{394.519}{503.044}{residue}%
\TextAt{8}{433.484}{503.044}{class}%
\TextAt{8}{460.159}{503.044}{field}%
\TextAt{8}{484.754}{503.044}{are}%
\TextAt{8}{503.503}{503.044}{called}%
\TextAt{12}{79.226}{485.998}{local}%
\TextAt{12}{111.397}{485.998}{fields}%
\TextAt{8}{141.524}{485.998}{.}%
\TextAt{8}{151.498}{485.998}{For}%
\TextAt{8}{172.977}{485.998}{such}%
\TextAt{8}{199.882}{485.998}{a}%
\TextAt{8}{210.146}{485.998}{local}%
\TextAt{8}{238.296}{485.998}{field,}%
\TextAt{8}{267.54}{485.998}{the}%
\TextAt{8}{288.205}{485.998}{normalized}%
\TextAt{8}{348.55}{485.998}{exponential}%
\TextAt{8}{411.823}{485.998}{valuation}%
\TextAt{8}{463.404}{485.998}{is}%
\TextAt{8}{475.677}{485.998}{denoted}%
\TextAt{8}{520.408}{485.998}{by}%
\TextAt{13}{79.226}{468.952}{v}%
\TextAt{27}{84.885}{467.159}{p}%
\TextAt{8}{89.372}{468.952}{,}%
\TextAt{8}{96.521}{468.952}{and}%
\TextAt{16}{119.284}{468.952}{\Character{106}}%
\TextAt{16}{126.493}{468.952}{\Character{106}}%
\TextAt{27}{129.823}{465.465}{p}%
\TextAt{8}{138.208}{468.952}{denotes}%
\TextAt{8}{180.541}{468.952}{the}%
\TextAt{8}{200.697}{468.952}{absolute}%
\TextAt{8}{246.94}{468.952}{value}%
\TextAt{8}{277.186}{468.952}{normalized}%
\TextAt{8}{337.022}{468.952}{by}%
\TextAt{16}{271.477}{440.581}{\Character{106}}%
\TextAt{13}{274.798}{440.581}{x}%
\TextAt{16}{281.45}{440.581}{\Character{106}}%
\TextAt{27}{284.771}{437.094}{p}%
\TextAt{9}{292.578}{440.581}{\Character{61}}%
\TextAt{13}{305.004}{440.581}{q}%
\TextAt{3}{310.623}{445.517}{\Character{0}}%
\TextAt{21}{317.209}{445.517}{v}%
\TextAt{28}{321.333}{444.521}{p}%
\TextAt{17}{325.416}{445.517}{\Character{40}}%
\TextAt{21}{328.709}{445.517}{x}%
\TextAt{17}{333.476}{445.517}{\Character{41}}%
\TextAt{13}{337.267}{440.581}{\Character{59}}%
\TextAt{8}{79.226}{412.209}{where}%
\TextAt{13}{113.035}{412.209}{q}%
\TextAt{8}{122.552}{412.209}{is}%
\TextAt{8}{134.316}{412.209}{the}%
\TextAt{8}{154.472}{412.209}{cardinality}%
\TextAt{8}{213.017}{412.209}{of}%
\TextAt{8}{226.347}{412.209}{the}%
\TextAt{8}{246.503}{412.209}{residue}%
\TextAt{8}{286.23}{412.209}{class}%
\TextAt{8}{313.667}{412.209}{field.}%
\TextAt{12}{79.226}{386.377}{Proposition}%
\TextAt{12}{151.937}{386.377}{(5.1).}%
\TextAt{8}{187.253}{386.377}{A}%
\TextAt{8}{199.926}{386.377}{local}%
\TextAt{8}{227.566}{386.377}{field}%
\TextAt{13}{252.923}{386.377}{K}%
\TextAt{8}{267.638}{386.377}{is}%
\TextAt{8}{279.402}{386.377}{locally}%
\TextAt{8}{316.475}{386.377}{compact.}%
\TextAt{8}{367.854}{386.377}{Its}%
\TextAt{8}{385.141}{386.377}{valuation}%
\TextAt{8}{436.214}{386.377}{ring}%
\TextAt{14}{460.279}{386.377}{o}%
\TextAt{8}{470.021}{386.377}{is}%
\TextAt{8}{481.785}{386.377}{compact.}%
\TextAt{19}{79.226}{351.758}{Proof.}%
\TextAt{8}{116}{351.758}{By}%
\TextAt{8}{136.965}{351.758}{(4.5)}%
\TextAt{8}{167.542}{351.758}{we}%
\TextAt{8}{187.371}{351.758}{have}%
\TextAt{14}{216.968}{351.758}{o}%
\NativeMathAt{230.567}{351.758}{11.955}{\cong}%
\NativeMathAt{247.62}{351.758}{11.955}{\varprojlim}%
\TextAt{14}{265.872}{351.758}{o}%
\TextAt{9}{271.716}{351.758}{\Character{47}}%
\TextAt{14}{277.569}{351.758}{p}%
\TextAt{21}{283.552}{356.096}{n}%
\TextAt{8}{289.189}{351.758}{,}%
\TextAt{8}{299.59}{351.758}{both}%
\TextAt{8}{329.832}{351.758}{algebraically}%
\TextAt{8}{400.411}{351.758}{and}%
\TextAt{8}{425.776}{351.758}{topologically.}%
\TextAt{8}{506.109}{351.758}{Since}%
\TextAt{14}{79.226}{334.712}{p}%
\TextAt{21}{85.209}{339.05}{\Character{23}}%
\TextAt{9}{90.391}{334.712}{\Character{47}}%
\TextAt{14}{96.244}{334.712}{p}%
\TextAt{21}{102.227}{339.05}{\Character{23}}%
\TextAt{17}{106.911}{339.05}{\Character{43}1}%
\NativeMathAt{121.551}{334.712}{11.955}{\cong}%
\TextAt{14}{134.17}{334.712}{o}%
\TextAt{9}{140.014}{334.712}{\Character{47}}%
\TextAt{14}{145.867}{334.712}{p}%
\TextAt{8}{155.404}{334.712}{(see}%
\TextAt{8}{178.527}{334.712}{(3.9)),}%
\TextAt{8}{214.033}{334.712}{the}%
\TextAt{8}{233.846}{334.712}{rings}%
\TextAt{14}{262.182}{334.712}{o}%
\TextAt{9}{268.026}{334.712}{\Character{47}}%
\TextAt{14}{273.879}{334.712}{p}%
\TextAt{21}{279.862}{339.05}{n}%
\TextAt{8}{289.051}{334.712}{are}%
\TextAt{8}{308.218}{334.712}{finite,}%
\TextAt{8}{341.106}{334.712}{hence}%
\TextAt{8}{373.268}{334.712}{compact.}%
\TextAt{8}{424.532}{334.712}{Being}%
\TextAt{8}{457.184}{334.712}{a}%
\TextAt{8}{466.595}{334.712}{closed}%
\TextAt{8}{500.777}{334.712}{subset}%
\TextAt{8}{79.226}{317.666}{of}%
\TextAt{8}{92.412}{317.666}{the}%
\TextAt{8}{112.424}{317.666}{compact}%
\TextAt{8}{159.108}{317.666}{product}%
\TextAt{20}{202.863}{326.632}{\Character{81}}%
\NativeMathAt{214.154}{323.477}{7.97}{\infty}%
\TextAt{21}{214.154}{314.179}{n}%
\TextAt{17}{219.293}{314.179}{\Character{61}1}%
\TextAt{14}{232.604}{317.666}{o}%
\TextAt{9}{238.448}{317.666}{\Character{47}}%
\TextAt{14}{244.301}{317.666}{p}%
\TextAt{21}{250.284}{322.004}{n}%
\TextAt{8}{255.921}{317.666}{,}%
\TextAt{8}{262.954}{317.666}{it}%
\TextAt{8}{274.514}{317.666}{follows}%
\TextAt{8}{312.806}{317.666}{that}%
\TextAt{8}{338.031}{317.666}{the}%
\TextAt{8}{358.043}{317.666}{projective}%
\TextAt{8}{411.876}{317.666}{limit}%
\NativeMathAt{439.695}{317.666}{11.955}{\varprojlim}%
\TextAt{14}{457.947}{317.666}{o}%
\TextAt{9}{463.791}{317.666}{\Character{47}}%
\TextAt{14}{469.644}{317.666}{p}%
\TextAt{21}{475.627}{322.004}{n}%
\TextAt{8}{481.263}{317.666}{,}%
\TextAt{8}{488.297}{317.666}{and}%
\TextAt{8}{510.915}{317.666}{thus}%
\TextAt{14}{79.226}{300.62}{o}%
\TextAt{8}{85.07}{300.62}{,}%
\TextAt{8}{93.531}{300.62}{is}%
\TextAt{8}{106.345}{300.62}{also}%
\TextAt{8}{130.874}{300.62}{compact.}%
\TextAt{8}{185.402}{300.62}{For}%
\TextAt{8}{207.421}{300.62}{every}%
\TextAt{13}{239.362}{300.62}{a}%
\NativeMathAt{250.617}{300.62}{11.955}{\in}%
\TextAt{13}{263.696}{300.62}{K}%
\TextAt{8}{274.514}{300.62}{,}%
\TextAt{8}{282.975}{300.62}{the}%
\TextAt{8}{304.18}{300.62}{set}%
\TextAt{13}{323.497}{300.62}{a}%
\TextAt{9}{333.014}{300.62}{\Character{43}}%
\TextAt{14}{345.491}{300.62}{o}%
\TextAt{8}{356.282}{300.62}{is}%
\TextAt{8}{369.096}{300.62}{an}%
\TextAt{8}{386.404}{300.62}{open,}%
\TextAt{8}{419.254}{300.62}{and}%
\TextAt{8}{443.066}{300.62}{at}%
\TextAt{8}{458.425}{300.62}{the}%
\TextAt{8}{479.631}{300.62}{same}%
\TextAt{8}{510.007}{300.62}{time}%
\TextAt{8}{79.226}{283.574}{compact}%
\TextAt{8}{126.055}{283.574}{neighbourhood,}%
\TextAt{8}{209.31}{283.574}{so}%
\TextAt{8}{223.68}{283.574}{that}%
\TextAt{13}{249.049}{283.574}{K}%
\TextAt{8}{263.764}{283.574}{is}%
\TextAt{8}{275.528}{283.574}{locally}%
\TextAt{8}{312.601}{283.574}{compact.}%
\Vector{q 0 G 0 g 0.399 w 452.334 -490.426 m 452.334 -482.356 l S Q}%
\Vector{q 0 G 0 g 0.399 w 452.534 -482.556 m 459.707 -482.556 l S Q}%
\Vector{q 0 G 0 g 0.399 w 452.534 -490.227 m 459.707 -490.227 l S Q}%
\Vector{q 0 G 0 g 0.399 w 459.906 -490.426 m 459.906 -482.356 l S Q}%
\TextAt{8}{97.159}{254.753}{In}%
\TextAt{8}{111.744}{254.753}{happy}%
\TextAt{8}{146.831}{254.753}{concord}%
\TextAt{8}{190.371}{254.753}{with}%
\TextAt{8}{216.995}{254.753}{the}%
\TextAt{8}{237.115}{254.753}{definition}%
\TextAt{8}{289.107}{254.753}{of}%
\TextAt{8}{302.401}{254.753}{global}%
\TextAt{8}{336.843}{254.753}{fields}%
\TextAt{8}{366.778}{254.753}{as}%
\TextAt{8}{381.112}{254.753}{the}%
\TextAt{8}{401.232}{254.753}{finite}%
\TextAt{8}{431.107}{254.753}{extensions}%
\TextAt{8}{487.451}{254.753}{of}%
\TextAt{18}{500.745}{254.753}{Q}%
\TextAt{8}{513.904}{254.753}{and}%
\TextAt{18}{79.226}{237.707}{F}%
\TextAt{21}{86.532}{235.913}{p}%
\TextAt{9}{91.293}{237.707}{\Character{40}}%
\TextAt{13}{95.845}{237.707}{t}%
\TextAt{9}{100.072}{237.707}{\Character{41}}%
\TextAt{8}{104.625}{237.707}{,}%
\TextAt{8}{111.774}{237.707}{we}%
\TextAt{8}{129.001}{237.707}{now}%
\TextAt{8}{153.39}{237.707}{obtain}%
\TextAt{8}{189.817}{237.707}{the}%
\TextAt{8}{209.974}{237.707}{following}%
\TextAt{8}{259.408}{237.707}{characterization}%
\TextAt{8}{344.948}{237.707}{of}%
\TextAt{8}{358.278}{237.707}{local}%
\TextAt{8}{385.918}{237.707}{fields.}%
\TextAt{12}{79.226}{211.874}{Proposition}%
\TextAt{12}{153.032}{211.874}{(5.2).}%
\TextAt{8}{189.3}{211.874}{The}%
\TextAt{8}{214.305}{211.874}{local}%
\TextAt{8}{242.898}{211.874}{fields}%
\TextAt{8}{273.821}{211.874}{are}%
\TextAt{8}{294.284}{211.874}{precisely}%
\TextAt{8}{343.092}{211.874}{the}%
\TextAt{8}{364.2}{211.874}{finite}%
\TextAt{8}{395.064}{211.874}{extensions}%
\TextAt{8}{452.397}{211.874}{of}%
\TextAt{8}{466.678}{211.874}{the}%
\TextAt{8}{487.787}{211.874}{fields}%
\TextAt{18}{518.71}{211.874}{Q}%
\TextAt{21}{528.009}{210.081}{p}%
\TextAt{8}{79.226}{194.828}{and}%
\TextAt{18}{101.989}{194.828}{F}%
\TextAt{21}{109.295}{193.035}{p}%
\TextAt{9}{114.056}{194.828}{\Character{40}\Character{40}}%
\TextAt{13}{123.16}{194.828}{t}%
\TextAt{9}{127.387}{194.828}{\Character{41}\Character{41}}%
\TextAt{8}{136.492}{194.828}{.}%
\TextAt{19}{79.226}{160.209}{Proof.}%
\TextAt{8}{116}{160.209}{A}%
\TextAt{8}{129.328}{160.209}{finite}%
\TextAt{8}{159.895}{160.209}{extension}%
\TextAt{13}{212.317}{160.209}{K}%
\TextAt{8}{227.688}{160.209}{of}%
\TextAt{13}{241.673}{160.209}{k}%
\TextAt{9}{252.601}{160.209}{\Character{61}}%
\TextAt{18}{266.144}{160.209}{Q}%
\TextAt{21}{275.442}{158.416}{p}%
\TextAt{8}{284.756}{160.209}{or}%
\TextAt{13}{299.722}{160.209}{k}%
\TextAt{9}{310.65}{160.209}{\Character{61}}%
\TextAt{18}{324.193}{160.209}{F}%
\TextAt{21}{331.499}{158.416}{p}%
\TextAt{9}{336.26}{160.209}{\Character{40}\Character{40}}%
\TextAt{13}{345.364}{160.209}{t}%
\TextAt{9}{349.592}{160.209}{\Character{41}\Character{41}}%
\TextAt{8}{363.249}{160.209}{is}%
\TextAt{8}{375.669}{160.209}{again}%
\TextAt{8}{407.552}{160.209}{complete,}%
\TextAt{8}{461.046}{160.209}{by}%
\TextAt{8}{477.96}{160.209}{(4.8),}%
\TextAt{8}{510.007}{160.209}{with}%
\TextAt{8}{79.226}{141.105}{respect}%
\TextAt{8}{121.184}{141.105}{to}%
\TextAt{8}{137.403}{141.105}{the}%
\TextAt{8}{159.468}{141.105}{extended}%
\TextAt{8}{211.121}{141.105}{valuation}%
\TextAt{16}{264.102}{141.105}{\Character{106}}%
\TextAt{13}{267.423}{141.105}{\Character{11}}%
\TextAt{16}{274.945}{141.105}{\Character{106}}%
\TextAt{9}{284.839}{141.105}{\Character{61}}%
\TextAt{22}{303.837}{145.624}{n}%
\TextAt{20}{301.796}{154.914}{\Character{113}}%
\Vector{q 0 G 0 g 0.478 w 241.751 -618.847 m 290.898 -618.847 l S Q}%
\TextAt{20}{313.751}{151.267}{\Character{12}}%
\TextAt{20}{313.751}{144.094}{\Character{12}}%
\TextAt{13}{317.736}{141.105}{N}%
\TextAt{21}{327.11}{139.246}{K}%
\TextAt{3}{334.815}{139.246}{\Character{106}}%
\TextAt{21}{337.167}{139.246}{k}%
\TextAt{9}{342.287}{141.105}{\Character{40}}%
\TextAt{13}{346.839}{141.105}{\Character{11}}%
\TextAt{9}{354.361}{141.105}{\Character{41}}%
\TextAt{20}{358.913}{151.267}{\Character{12}}%
\TextAt{20}{358.913}{144.094}{\Character{12}}%
\TextAt{8}{362.898}{141.105}{,}%
\TextAt{8}{372.433}{141.105}{which}%
\TextAt{8}{407.828}{141.105}{itself}%
\TextAt{8}{438.082}{141.105}{is}%
\TextAt{8}{451.754}{141.105}{obviously}%
\TextAt{8}{505.44}{141.105}{again}%
\TextAt{8}{79.226}{124.059}{discrete.}%
\TextAt{8}{127.883}{124.059}{Since}%
\TextAt{13}{158.816}{124.059}{K}%
\TextAt{16}{173.594}{124.059}{\Character{106}}%
\TextAt{13}{180.876}{124.059}{k}%
\TextAt{8}{191.639}{124.059}{is}%
\TextAt{8}{203.778}{124.059}{of}%
\TextAt{8}{217.484}{124.059}{finite}%
\TextAt{8}{247.771}{124.059}{degree,}%
\TextAt{8}{287.908}{124.059}{so}%
\TextAt{8}{302.654}{124.059}{is}%
\TextAt{8}{314.793}{124.059}{the}%
\TextAt{8}{335.326}{124.059}{residue}%
\TextAt{8}{375.428}{124.059}{class}%
\TextAt{8}{403.241}{124.059}{field}%
\TextAt{8}{428.973}{124.059}{extension}%
\TextAt{13}{481.115}{124.059}{\Character{20}}%
\TextAt{16}{491.815}{124.059}{\Character{106}}%
\TextAt{18}{499.096}{124.059}{F}%
\TextAt{21}{506.402}{122.266}{p}%
\TextAt{8}{511.163}{124.059}{,}%
\TextAt{8}{518.782}{124.059}{for}%
\TextAt{8}{79.226}{107.013}{if}%
\TextAt{9}{91.244}{107.013}{\Character{22}}%
\TextAt{13}{90.52}{107.013}{x}%
\TextAt{17}{97.172}{105.22}{1}%
\TextAt{13}{101.904}{107.013}{\Character{59}}%
\TextAt{13}{107.128}{107.013}{\Character{58}}%
\TextAt{13}{112.365}{107.013}{\Character{58}}%
\TextAt{13}{117.589}{107.013}{\Character{58}}%
\TextAt{13}{122.826}{107.013}{\Character{59}}%
\TextAt{9}{128.85}{107.013}{\Character{22}}%
\TextAt{13}{128.125}{107.013}{x}%
\TextAt{21}{134.777}{105.22}{n}%
\NativeMathAt{144.705}{107.013}{11.955}{\in}%
\TextAt{13}{156.967}{107.013}{\Character{20}}%
\TextAt{8}{168.172}{107.013}{are}%
\TextAt{8}{188.253}{107.013}{linearly}%
\TextAt{8}{230.773}{107.013}{independent,}%
\TextAt{8}{301.065}{107.013}{then}%
\TextAt{8}{328.294}{107.013}{any}%
\TextAt{8}{350.981}{107.013}{choice}%
\TextAt{8}{386.34}{107.013}{of}%
\TextAt{8}{400.24}{107.013}{preimages}%
\TextAt{13}{455.505}{107.013}{x}%
\TextAt{17}{462.157}{105.22}{1}%
\TextAt{13}{466.889}{107.013}{\Character{59}}%
\TextAt{13}{472.113}{107.013}{\Character{58}}%
\TextAt{13}{477.35}{107.013}{\Character{58}}%
\TextAt{13}{482.586}{107.013}{\Character{58}}%
\TextAt{13}{487.811}{107.013}{\Character{59}}%
\TextAt{13}{493.047}{107.013}{x}%
\TextAt{21}{499.762}{105.22}{n}%
\NativeMathAt{509.69}{107.013}{11.955}{\in}%
\TextAt{13}{521.952}{107.013}{K}%
\TextAt{8}{79.226}{89.967}{is}%
\TextAt{8}{91.782}{89.967}{linearly}%
\TextAt{8}{134.524}{89.967}{independent}%
\TextAt{8}{201.643}{89.967}{over}%
\TextAt{13}{227.481}{89.967}{k}%
\TextAt{8}{233.971}{89.967}{.}%
\TextAt{8}{244.794}{89.967}{Indeed,}%
\TextAt{8}{287.065}{89.967}{dividing}%
\TextAt{8}{333.059}{89.967}{any}%
\TextAt{8}{355.968}{89.967}{nontrivial}%
\TextAt{13}{410.104}{89.967}{k}%
\TextAt{8}{416.593}{89.967}{-linear}%
\TextAt{8}{453.8}{89.967}{relation}%
\TextAt{13}{497.523}{89.967}{\Character{21}}%
\TextAt{17}{504.352}{88.174}{1}%
\TextAt{13}{509.084}{89.967}{x}%
\TextAt{17}{515.736}{88.174}{1}%
\TextAt{9}{523.665}{89.967}{\Character{43}}%
\TextAt{16}{79.226}{72.921}{\Character{1}}%
\TextAt{16}{84.522}{72.921}{\Character{1}}%
\TextAt{16}{89.83}{72.921}{\Character{1}}%
\TextAt{9}{96.177}{72.921}{\Character{43}}%
\TextAt{13}{108.286}{72.921}{\Character{21}}%
\TextAt{21}{115.114}{71.128}{n}%
\TextAt{13}{120.751}{72.921}{x}%
\TextAt{21}{127.403}{71.128}{n}%
\TextAt{9}{137.227}{72.921}{\Character{61}}%
\TextAt{9}{150.509}{72.921}{0}%
\TextAt{8}{156.373}{72.921}{,}%
\TextAt{13}{164.158}{72.921}{\Character{21}}%
\TextAt{21}{170.987}{71.128}{i}%
\NativeMathAt{178.556}{72.921}{11.955}{\in}%
\TextAt{13}{190.715}{72.921}{k}%
\TextAt{8}{197.204}{72.921}{,}%
\TextAt{8}{204.99}{72.921}{by}%
\TextAt{8}{221.758}{72.921}{the}%
\TextAt{8}{242.423}{72.921}{coefficient}%
\TextAt{13}{297.555}{72.921}{\Character{21}}%
\TextAt{21}{304.384}{71.128}{i}%
\TextAt{8}{312.172}{72.921}{with}%
\TextAt{8}{339.341}{72.921}{biggest}%
\TextAt{8}{379.589}{72.921}{absolute}%
\TextAt{8}{426.34}{72.921}{value,}%
\TextAt{8}{460.475}{72.921}{yields}%
\TextAt{8}{493.884}{72.921}{a}%
\TextAt{8}{504.149}{72.921}{linear}%
\TextAt{8}{297.211}{43.033}{123}%
\Destination{78.23}{809.815}{nav.page.123}%
\Destination{79.23}{629.515}{nav.section.2.5}%
\Destination{79.226}{775.155}{nav.ex.2.4.5}%
\Destination{79.226}{686.755}{nav.ex.2.4.6}%
\Destination{79.36}{401.355}{nav.print.135}%
\Destination{79.36}{400.355}{nav.number.2.5.1}%
\Destination{79.36}{225.855}{nav.number.2.5.2}%
\Destination{79.23}{629.515}{outline.25}%
\Hotlink{136.465}{349.555}{25.078}{11.047}{nav.number.2.4.5}%
\Hotlink{178.027}{332.555}{25.078}{11.047}{nav.number.2.3.9}%
\Hotlink{477.46}{158.055}{25.078}{11.047}{nav.number.2.4.8}%
}

% Page 136 (body page 124)
\Page{%
\TextAt{2}{79.226}{796.943}{II.}%
\TextAt{2}{94.073}{796.943}{The}%
\TextAt{2}{116.492}{796.943}{Theory}%
\TextAt{2}{154.401}{796.943}{of}%
\TextAt{2}{166.826}{796.943}{Valuations}%
\Vector{q 0 G 0 g 0.299 w 7.226 18.729 m 460.769 18.729 l S Q}%
\TextAt{8}{79.226}{761.124}{combination}%
\TextAt{8}{146.816}{761.124}{with}%
\TextAt{8}{174.393}{761.124}{coefficients}%
\TextAt{8}{234.547}{761.124}{in}%
\TextAt{8}{249.116}{761.124}{the}%
\TextAt{8}{270.189}{761.124}{valuation}%
\TextAt{8}{322.177}{761.124}{ring}%
\TextAt{8}{347.159}{761.124}{of}%
\TextAt{13}{361.406}{761.124}{k}%
\TextAt{8}{372.709}{761.124}{with}%
\TextAt{8}{400.285}{761.124}{1}%
\TextAt{8}{410.957}{761.124}{as}%
\TextAt{13}{426.243}{761.124}{i}%
\TextAt{8}{430.237}{761.124}{-th}%
\TextAt{8}{450.006}{761.124}{coefficient,}%
\TextAt{8}{509.027}{761.124}{from}%
\TextAt{8}{79.226}{744.078}{which}%
\TextAt{8}{112.82}{744.078}{we}%
\TextAt{8}{130.156}{744.078}{obtain}%
\TextAt{8}{166.691}{744.078}{a}%
\TextAt{8}{176.554}{744.078}{nontrivial}%
\TextAt{8}{230.006}{744.078}{relation}%
\TextAt{9}{273.532}{747.233}{\Character{22}}%
\TextAt{13}{273.045}{744.078}{\Character{21}}%
\TextAt{17}{279.873}{742.284}{1}%
\TextAt{9}{285.33}{744.078}{\Character{22}}%
\TextAt{13}{284.605}{744.078}{x}%
\TextAt{17}{291.258}{742.284}{1}%
\TextAt{9}{298.72}{744.078}{\Character{43}}%
\TextAt{16}{310.555}{744.078}{\Character{1}}%
\TextAt{16}{315.851}{744.078}{\Character{1}}%
\TextAt{16}{321.159}{744.078}{\Character{1}}%
\TextAt{9}{327.233}{744.078}{\Character{43}}%
\TextAt{9}{339.555}{747.233}{\Character{22}}%
\TextAt{13}{339.067}{744.078}{\Character{21}}%
\TextAt{21}{345.896}{742.284}{n}%
\TextAt{9}{352.257}{744.078}{\Character{22}}%
\TextAt{13}{351.532}{744.078}{x}%
\TextAt{21}{358.184}{742.284}{n}%
\TextAt{9}{367.325}{744.078}{\Character{61}}%
\TextAt{9}{379.926}{744.078}{0}%
\TextAt{8}{389.793}{744.078}{by}%
\TextAt{8}{406.159}{744.078}{reducing}%
\TextAt{8}{453.741}{744.078}{to}%
\TextAt{13}{468.159}{744.078}{\Character{20}}%
\TextAt{8}{474.898}{744.078}{.}%
\TextAt{8}{483.67}{744.078}{Therefore}%
\TextAt{13}{79.226}{727.032}{K}%
\TextAt{8}{93.941}{727.032}{is}%
\TextAt{8}{105.705}{727.032}{a}%
\TextAt{8}{115.461}{727.032}{local}%
\TextAt{8}{143.101}{727.032}{field.}%
\TextAt{8}{97.159}{706.997}{Conversely,}%
\TextAt{8}{160.831}{706.997}{let}%
\TextAt{13}{179.455}{706.997}{K}%
\TextAt{8}{195.891}{706.997}{be}%
\TextAt{8}{213.535}{706.997}{a}%
\TextAt{8}{225.011}{706.997}{local}%
\TextAt{8}{254.372}{706.997}{field,}%
\TextAt{13}{285.131}{706.997}{v}%
\TextAt{8}{296.837}{706.997}{its}%
\TextAt{8}{314.876}{706.997}{discrete}%
\TextAt{8}{359.575}{706.997}{exponential}%
\TextAt{8}{424.06}{706.997}{valuation,}%
\TextAt{8}{480.535}{706.997}{and}%
\TextAt{13}{505.018}{706.997}{p}%
\TextAt{8}{516.51}{706.997}{the}%
\TextAt{8}{79.226}{689.951}{characteristic}%
\TextAt{8}{151.967}{689.951}{of}%
\TextAt{8}{166.104}{689.951}{its}%
\TextAt{8}{183.23}{689.951}{residue}%
\TextAt{8}{223.764}{689.951}{class}%
\TextAt{8}{252.008}{689.951}{field}%
\TextAt{13}{278.172}{689.951}{\Character{20}}%
\TextAt{8}{284.91}{689.951}{.}%
\TextAt{8}{295.78}{689.951}{If}%
\TextAt{13}{308.279}{689.951}{K}%
\TextAt{8}{323.801}{689.951}{has}%
\TextAt{8}{345.481}{689.951}{characteristic}%
\TextAt{8}{418.223}{689.951}{0,}%
\TextAt{8}{432.238}{689.951}{then}%
\TextAt{8}{459.705}{689.951}{the}%
\TextAt{8}{480.669}{689.951}{restriction}%
\TextAt{8}{79.226}{672.905}{of}%
\TextAt{13}{92.935}{672.905}{v}%
\TextAt{8}{103.299}{672.905}{to}%
\TextAt{18}{117.988}{672.905}{Q}%
\TextAt{8}{131.563}{672.905}{is}%
\TextAt{8}{143.706}{672.905}{equivalent}%
\TextAt{8}{199.7}{672.905}{to}%
\TextAt{8}{214.389}{672.905}{the}%
\TextAt{13}{234.924}{672.905}{p}%
\TextAt{8}{240.799}{672.905}{-adic}%
\TextAt{8}{269.786}{672.905}{valuation}%
\TextAt{13}{321.238}{672.905}{v}%
\TextAt{21}{326.897}{671.112}{p}%
\TextAt{8}{335.934}{672.905}{of}%
\TextAt{18}{349.643}{672.905}{Q}%
\TextAt{8}{363.217}{672.905}{because}%
\TextAt{13}{406.898}{672.905}{v}%
\TextAt{9}{412.986}{672.905}{\Character{40}}%
\TextAt{13}{417.538}{672.905}{p}%
\TextAt{9}{423.413}{672.905}{\Character{41}}%
\TextAt{13}{431.932}{672.905}{\Character{62}}%
\TextAt{9}{445.003}{672.905}{0}%
\TextAt{8}{450.856}{672.905}{.}%
\TextAt{8}{460.44}{672.905}{In}%
\TextAt{8}{475.44}{672.905}{view}%
\TextAt{8}{502.802}{672.905}{of}%
\TextAt{8}{516.511}{672.905}{the}%
\TextAt{8}{79.226}{655.859}{completeness}%
\TextAt{8}{150.184}{655.859}{of}%
\TextAt{13}{164.116}{655.859}{K}%
\TextAt{8}{174.933}{655.859}{,}%
\TextAt{8}{182.835}{655.859}{the}%
\TextAt{8}{203.593}{655.859}{closure}%
\TextAt{8}{243.276}{655.859}{of}%
\TextAt{18}{257.208}{655.859}{Q}%
\TextAt{8}{271.005}{655.859}{in}%
\TextAt{13}{285.26}{655.859}{K}%
\TextAt{8}{300.577}{655.859}{is}%
\TextAt{8}{312.942}{655.859}{the}%
\TextAt{8}{333.701}{655.859}{completion}%
\TextAt{8}{394.138}{655.859}{of}%
\TextAt{18}{408.07}{655.859}{Q}%
\TextAt{8}{421.867}{655.859}{with}%
\TextAt{8}{449.129}{655.859}{respect}%
\TextAt{8}{489.78}{655.859}{to}%
\TextAt{13}{504.693}{655.859}{v}%
\TextAt{21}{510.352}{654.066}{p}%
\TextAt{8}{515.113}{655.859}{,}%
\TextAt{8}{523.014}{655.859}{in}%
\TextAt{8}{79.226}{638.813}{other}%
\TextAt{8}{111.056}{638.813}{words}%
\TextAt{18}{145.874}{638.813}{Q}%
\TextAt{21}{155.172}{637.02}{p}%
\NativeMathAt{165.402}{638.813}{11.955}{\subseteq}%
\TextAt{13}{180.168}{638.813}{K}%
\TextAt{8}{190.986}{638.813}{.}%
\TextAt{8}{203.215}{638.813}{The}%
\TextAt{8}{228.529}{638.813}{fact}%
\TextAt{8}{252.874}{638.813}{that}%
\TextAt{13}{279.503}{638.813}{K}%
\TextAt{16}{295.789}{638.813}{\Character{106}}%
\TextAt{18}{304.578}{638.813}{Q}%
\TextAt{21}{313.877}{637.02}{p}%
\TextAt{8}{323.795}{638.813}{is}%
\TextAt{8}{336.819}{638.813}{of}%
\TextAt{8}{351.409}{638.813}{finite}%
\TextAt{8}{382.581}{638.813}{degree}%
\TextAt{8}{420.257}{638.813}{results}%
\TextAt{8}{458.71}{638.813}{from}%
\TextAt{8}{487.61}{638.813}{the}%
\TextAt{8}{509.027}{638.813}{local}%
\TextAt{8}{79.226}{621.767}{compactness}%
\TextAt{8}{148.015}{621.767}{of}%
\TextAt{8}{162.372}{621.767}{the}%
\TextAt{8}{183.555}{621.767}{vector}%
\TextAt{8}{219.706}{621.767}{space}%
\TextAt{13}{252.008}{621.767}{K}%
\TextAt{8}{262.826}{621.767}{,}%
\TextAt{8}{271.259}{621.767}{by}%
\TextAt{8}{288.545}{621.767}{a}%
\TextAt{8}{299.327}{621.767}{general}%
\TextAt{8}{340.679}{621.767}{theorem}%
\TextAt{8}{387.231}{621.767}{of}%
\TextAt{8}{401.588}{621.767}{topological}%
\TextAt{8}{462.14}{621.767}{linear}%
\TextAt{8}{495.685}{621.767}{algebra}%
\TextAt{8}{79.226}{604.721}{(see}%
\TextAt{8}{103.451}{604.721}{[18],}%
\TextAt{8}{129.767}{604.721}{chap.}%
\TextAt{8}{161.417}{604.721}{I,}%
\TextAt{8}{173.543}{604.721}{§2,}%
\TextAt{8}{193.176}{604.721}{n}%
\TextAt{7}{199.679}{609.06}{o}%
\TextAt{8}{209.064}{604.721}{4,}%
\TextAt{8}{223.018}{604.721}{th.}%
\TextAt{8}{241.983}{604.721}{3),}%
\TextAt{8}{260.492}{604.721}{but}%
\TextAt{8}{282.709}{604.721}{it}%
\TextAt{8}{295.171}{604.721}{also}%
\TextAt{8}{319.408}{604.721}{follows}%
\TextAt{8}{358.601}{604.721}{from}%
\TextAt{8}{386.999}{604.721}{(6.8)}%
\TextAt{8}{415.731}{604.721}{below.}%
\TextAt{8}{455.718}{604.721}{If}%
\TextAt{8}{468.167}{604.721}{on}%
\TextAt{8}{485.184}{604.721}{the}%
\TextAt{8}{506.097}{604.721}{other}%
\TextAt{8}{79.226}{587.675}{hand}%
\TextAt{8}{108.685}{587.675}{the}%
\TextAt{8}{129.034}{587.675}{characteristic}%
\TextAt{8}{201.16}{587.675}{of}%
\TextAt{13}{214.683}{587.675}{K}%
\TextAt{8}{229.59}{587.675}{is}%
\TextAt{8}{241.547}{587.675}{not}%
\TextAt{8}{262.553}{587.675}{equal}%
\TextAt{8}{293.638}{587.675}{to}%
\TextAt{8}{308.14}{587.675}{zero,}%
\TextAt{8}{336.344}{587.675}{then}%
\TextAt{8}{363.197}{587.675}{it}%
\TextAt{8}{375.093}{587.675}{has}%
\TextAt{8}{396.159}{587.675}{to}%
\TextAt{8}{410.662}{587.675}{equal}%
\TextAt{13}{441.747}{587.675}{p}%
\TextAt{8}{447.622}{587.675}{.}%
\TextAt{8}{456.647}{587.675}{In}%
\TextAt{8}{471.461}{587.675}{this}%
\TextAt{8}{494.476}{587.675}{case}%
\TextAt{8}{519.439}{587.675}{we}%
\TextAt{8}{79.226}{570.629}{find}%
\TextAt{13}{103.05}{570.629}{K}%
\TextAt{9}{117.897}{570.629}{\Character{61}}%
\TextAt{13}{131.032}{570.629}{\Character{20}}%
\TextAt{9}{137.77}{570.629}{\Character{40}\Character{40}}%
\TextAt{13}{146.875}{570.629}{t}%
\TextAt{9}{151.102}{570.629}{\Character{41}\Character{41}}%
\TextAt{8}{160.207}{570.629}{,}%
\TextAt{8}{167.876}{570.629}{for}%
\TextAt{8}{186.176}{570.629}{a}%
\TextAt{8}{196.348}{570.629}{prime}%
\TextAt{8}{229.927}{570.629}{element}%
\TextAt{13}{273.585}{570.629}{t}%
\TextAt{8}{282.125}{570.629}{of}%
\TextAt{13}{295.871}{570.629}{K}%
\TextAt{8}{311.002}{570.629}{(see}%
\TextAt{8}{334.886}{570.629}{p.}%
\TextAt{8}{348.955}{570.629}{127),}%
\TextAt{8}{378.753}{570.629}{hence}%
\TextAt{18}{411.675}{570.629}{F}%
\TextAt{21}{418.981}{568.836}{p}%
\TextAt{9}{423.741}{570.629}{\Character{40}\Character{40}}%
\TextAt{13}{432.846}{570.629}{t}%
\TextAt{9}{437.073}{570.629}{\Character{41}\Character{41}}%
\NativeMathAt{450.207}{570.629}{11.955}{\subseteq}%
\TextAt{13}{463.535}{570.629}{K}%
\TextAt{8}{474.353}{570.629}{.}%
\TextAt{8}{484.049}{570.629}{In}%
\TextAt{8}{499.086}{570.629}{fact,}%
\TextAt{8}{525.943}{570.629}{if}%
\TextAt{13}{79.226}{553.583}{\Character{20}}%
\TextAt{9}{89.975}{553.583}{\Character{61}}%
\TextAt{18}{103.09}{553.583}{F}%
\TextAt{21}{110.396}{551.79}{p}%
\TextAt{9}{115.157}{553.583}{\Character{40}}%
\TextAt{13}{119.709}{553.583}{\Character{11}}%
\TextAt{9}{127.231}{553.583}{\Character{41}}%
\TextAt{8}{136.086}{553.583}{and}%
\TextAt{13}{159.253}{553.583}{p}%
\TextAt{9}{165.128}{553.583}{\Character{40}}%
\TextAt{13}{169.68}{553.583}{X}%
\TextAt{9}{180.335}{553.583}{\Character{41}}%
\NativeMathAt{188.898}{553.583}{11.955}{\in}%
\TextAt{18}{200.879}{553.583}{F}%
\TextAt{21}{208.185}{551.79}{p}%
\TextAt{9}{212.946}{553.583}{\Character{91}}%
\TextAt{13}{216.197}{553.583}{X}%
\TextAt{9}{226.852}{553.583}{\Character{93}}%
\NativeMathAt{234.115}{553.583}{11.955}{\subseteq}%
\TextAt{13}{247.424}{553.583}{K}%
\TextAt{9}{258.241}{553.583}{\Character{91}}%
\TextAt{13}{261.493}{553.583}{X}%
\TextAt{9}{272.148}{553.583}{\Character{93}}%
\TextAt{8}{279.702}{553.583}{is}%
\TextAt{8}{291.87}{553.583}{the}%
\TextAt{8}{312.431}{553.583}{minimal}%
\TextAt{8}{358.361}{553.583}{polynomial}%
\TextAt{8}{419.259}{553.583}{of}%
\TextAt{13}{432.994}{553.583}{\Character{11}}%
\TextAt{8}{444.818}{553.583}{over}%
\TextAt{18}{470.268}{553.583}{F}%
\TextAt{21}{477.574}{551.79}{p}%
\TextAt{8}{482.335}{553.583}{,}%
\TextAt{8}{489.99}{553.583}{then,}%
\TextAt{8}{520.408}{553.583}{by}%
\TextAt{8}{79.226}{536.537}{Hensel’s}%
\TextAt{8}{125.458}{536.537}{lemma,}%
\TextAt{13}{167.58}{536.537}{p}%
\TextAt{9}{173.455}{536.537}{\Character{40}}%
\TextAt{13}{178.007}{536.537}{X}%
\TextAt{9}{188.662}{536.537}{\Character{41}}%
\TextAt{8}{198.034}{536.537}{splits}%
\TextAt{8}{229.644}{536.537}{over}%
\TextAt{13}{255.612}{536.537}{K}%
\TextAt{8}{271.248}{536.537}{into}%
\TextAt{8}{295.913}{536.537}{linear}%
\TextAt{8}{329.352}{536.537}{factors.}%
\TextAt{8}{374.781}{536.537}{We}%
\TextAt{8}{395.846}{536.537}{may}%
\TextAt{8}{422.137}{536.537}{therefore}%
\TextAt{8}{472.158}{536.537}{view}%
\TextAt{13}{500.062}{536.537}{\Character{20}}%
\TextAt{8}{511.62}{536.537}{as}%
\TextAt{8}{526.911}{536.537}{a}%
\TextAt{8}{79.226}{519.491}{subfield}%
\TextAt{8}{122.241}{519.491}{of}%
\TextAt{13}{135.607}{519.491}{K}%
\TextAt{8}{146.424}{519.491}{,}%
\TextAt{8}{153.619}{519.491}{and}%
\TextAt{8}{176.417}{519.491}{then}%
\TextAt{8}{203.113}{519.491}{the}%
\TextAt{8}{223.305}{519.491}{elements}%
\TextAt{8}{271.198}{519.491}{of}%
\TextAt{13}{284.564}{519.491}{K}%
\TextAt{8}{299.315}{519.491}{turn}%
\TextAt{8}{325.365}{519.491}{out}%
\TextAt{8}{346.215}{519.491}{to}%
\TextAt{8}{360.561}{519.491}{be,}%
\TextAt{8}{379.782}{519.491}{by}%
\TextAt{8}{396.077}{519.491}{(4.4),}%
\TextAt{8}{427.349}{519.491}{the}%
\TextAt{8}{447.541}{519.491}{Laurent}%
\TextAt{8}{491.644}{519.491}{series}%
\TextAt{8}{523.014}{519.491}{in}%
\TextAt{13}{79.226}{502.445}{t}%
\TextAt{8}{87.351}{502.445}{with}%
\TextAt{8}{114.011}{502.445}{coefficients}%
\TextAt{8}{173.249}{502.445}{in}%
\TextAt{13}{186.901}{502.445}{\Character{20}}%
\TextAt{8}{193.64}{502.445}{.}%
\Vector{q 0 G 0 g 0.399 w 452.334 -271.555 m 452.334 -263.485 l S Q}%
\Vector{q 0 G 0 g 0.399 w 452.534 -263.684 m 459.707 -263.684 l S Q}%
\Vector{q 0 G 0 g 0.399 w 452.534 -271.355 m 459.707 -271.355 l S Q}%
\Vector{q 0 G 0 g 0.399 w 459.906 -271.555 m 459.906 -263.485 l S Q}%
\TextAt{12}{79.226}{470.455}{Remark.}%
\TextAt{8}{132.88}{470.455}{One}%
\TextAt{8}{157.089}{470.455}{can}%
\TextAt{8}{178.046}{470.455}{show}%
\TextAt{8}{206.546}{470.455}{that}%
\TextAt{8}{231.412}{470.455}{a}%
\TextAt{8}{240.665}{470.455}{field}%
\TextAt{13}{265.519}{470.455}{K}%
\TextAt{8}{279.731}{470.455}{which}%
\TextAt{8}{312.715}{470.455}{is}%
\TextAt{8}{323.976}{470.455}{locally}%
\TextAt{8}{360.546}{470.455}{compact}%
\TextAt{8}{406.872}{470.455}{with}%
\TextAt{8}{433.029}{470.455}{respect}%
\TextAt{8}{472.576}{470.455}{to}%
\TextAt{8}{486.384}{470.455}{a}%
\TextAt{8}{495.637}{470.455}{nondis-}%
\TextAt{8}{79.226}{453.409}{crete}%
\TextAt{8}{108.546}{453.409}{topology}%
\TextAt{8}{157.401}{453.409}{is}%
\TextAt{8}{169.876}{453.409}{isomorphic}%
\TextAt{8}{229.838}{453.409}{either}%
\TextAt{8}{263.713}{453.409}{to}%
\TextAt{18}{278.734}{453.409}{R}%
\TextAt{8}{291.978}{453.409}{or}%
\TextAt{18}{306.999}{453.409}{C}%
\TextAt{8}{315.634}{453.409}{,}%
\TextAt{8}{323.672}{453.409}{or}%
\TextAt{8}{338.694}{453.409}{to}%
\TextAt{8}{353.716}{453.409}{a}%
\TextAt{8}{364.182}{453.409}{finite}%
\TextAt{8}{394.806}{453.409}{extension}%
\TextAt{8}{447.283}{453.409}{of}%
\TextAt{18}{461.324}{453.409}{Q}%
\TextAt{21}{470.623}{451.616}{p}%
\TextAt{8}{479.993}{453.409}{or}%
\TextAt{18}{495.014}{453.409}{F}%
\TextAt{21}{502.32}{451.616}{p}%
\TextAt{9}{507.081}{453.409}{\Character{40}\Character{40}}%
\TextAt{13}{516.186}{453.409}{t}%
\TextAt{9}{520.413}{453.409}{\Character{41}\Character{41}}%
\TextAt{8}{529.518}{453.409}{,}%
\TextAt{8}{79.226}{436.363}{i.e.,}%
\TextAt{8}{101.331}{436.363}{to}%
\TextAt{8}{115.642}{436.363}{a}%
\TextAt{8}{125.397}{436.363}{local}%
\TextAt{8}{153.037}{436.363}{field}%
\TextAt{8}{178.394}{436.363}{(see}%
\TextAt{8}{201.862}{436.363}{[137],}%
\TextAt{8}{233.089}{436.363}{chap.}%
\TextAt{8}{263.981}{436.363}{I,}%
\TextAt{8}{275.351}{436.363}{§3).}%
\ErrataAt{97.14}{413.31}{437.86}{We have just seen that the local fields of characteristic $p$ are the Laurent series fields}%
\TextAt{18}{79.226}{396.294}{F}%
\TextAt{21}{86.532}{394.501}{q}%
\TextAt{9}{91.097}{396.294}{\Character{40}\Character{40}}%
\TextAt{13}{100.202}{396.294}{t}%
\TextAt{9}{104.429}{396.294}{\Character{41}\Character{41}}%
\TextAt{8}{113.533}{396.294}{,}%
\TextAt{8}{120.736}{396.294}{with}%
\TextAt{13}{147.439}{396.294}{q}%
\TextAt{9}{156.451}{396.294}{\Character{61}}%
\TextAt{13}{168.949}{396.294}{p}%
\TextAt{21}{174.825}{400.632}{f}%
\TextAt{8}{180.269}{396.294}{.}%
\TextAt{8}{188.846}{396.294}{The}%
\TextAt{8}{212.942}{396.294}{local}%
\TextAt{8}{240.625}{396.294}{fields}%
\TextAt{8}{270.639}{396.294}{of}%
\TextAt{8}{284.012}{396.294}{characteristic}%
\TextAt{8}{355.989}{396.294}{0,}%
\TextAt{8}{369.049}{396.294}{i.e.,}%
\TextAt{8}{391.208}{396.294}{the}%
\TextAt{8}{411.407}{396.294}{finite}%
\TextAt{8}{441.361}{396.294}{extensions}%
\TextAt{13}{497.785}{396.294}{K}%
\TextAt{16}{511.996}{396.294}{\Character{106}}%
\TextAt{18}{518.71}{396.294}{Q}%
\TextAt{21}{528.009}{394.501}{p}%
\TextAt{8}{79.226}{379.248}{of}%
\TextAt{8}{92.86}{379.248}{the}%
\TextAt{8}{113.321}{379.248}{fields}%
\TextAt{8}{143.597}{379.248}{of}%
\TextAt{13}{157.231}{379.248}{p}%
\TextAt{8}{163.106}{379.248}{-adic}%
\TextAt{8}{192.019}{379.248}{numbers}%
\TextAt{18}{239.534}{379.248}{Q}%
\TextAt{21}{248.832}{377.455}{p}%
\TextAt{8}{253.593}{379.248}{,}%
\TextAt{8}{261.123}{379.248}{are}%
\TextAt{8}{280.937}{379.248}{called}%
\TextAt{13}{314.405}{379.248}{p}%
\TextAt{12}{320.28}{379.248}{-adic}%
\TextAt{12}{353.328}{379.248}{number}%
\TextAt{12}{403.052}{379.248}{fields}%
\TextAt{8}{433.179}{379.248}{.}%
\TextAt{8}{442.54}{379.248}{For}%
\TextAt{8}{463.814}{379.248}{them}%
\TextAt{8}{494.029}{379.248}{one}%
\TextAt{8}{515.793}{379.248}{has}%
\TextAt{8}{79.226}{362.202}{an}%
\TextAt{19}{96.376}{362.202}{exponential}%
\TextAt{19}{158.602}{362.202}{function}%
\TextAt{8}{205.508}{362.202}{and}%
\TextAt{8}{229.161}{362.202}{a}%
\TextAt{19}{239.807}{362.202}{logarithm}%
\TextAt{19}{292.625}{362.202}{function}%
\TextAt{8}{334.468}{362.202}{.}%
\TextAt{8}{345.588}{362.202}{In}%
\TextAt{8}{361.1}{362.202}{contrast}%
\TextAt{8}{407.265}{362.202}{to}%
\TextAt{8}{422.465}{362.202}{the}%
\TextAt{8}{443.512}{362.202}{real}%
\TextAt{8}{467.166}{362.202}{and}%
\TextAt{8}{490.819}{362.202}{complex}%
\TextAt{8}{79.226}{345.156}{case,}%
\TextAt{8}{108.092}{345.156}{however,}%
\TextAt{8}{157.066}{345.156}{the}%
\TextAt{8}{177.897}{345.156}{former}%
\TextAt{8}{215.967}{345.156}{is}%
\TextAt{8}{228.405}{345.156}{not}%
\TextAt{8}{249.894}{345.156}{defined}%
\TextAt{8}{290.881}{345.156}{on}%
\TextAt{8}{307.814}{345.156}{all}%
\TextAt{8}{324.747}{345.156}{of}%
\TextAt{13}{338.752}{345.156}{K}%
\TextAt{8}{349.569}{345.156}{,}%
\TextAt{8}{357.561}{345.156}{whereas}%
\TextAt{8}{402.518}{345.156}{the}%
\TextAt{8}{423.348}{345.156}{latter}%
\TextAt{8}{455.895}{345.156}{is}%
\TextAt{8}{468.333}{345.156}{given}%
\TextAt{8}{499.577}{345.156}{on}%
\TextAt{8}{516.51}{345.156}{the}%
\TextAt{8}{79.226}{328.11}{whole}%
\TextAt{8}{113.127}{328.11}{multiplicative}%
\TextAt{8}{187.687}{328.11}{group}%
\TextAt{13}{221.6}{328.11}{K}%
\TextAt{3}{232.418}{332.448}{\Character{3}}%
\TextAt{8}{237.15}{328.11}{.}%
\TextAt{8}{247.81}{328.11}{For}%
\TextAt{8}{269.517}{328.11}{the}%
\TextAt{8}{290.41}{328.11}{definition}%
\TextAt{8}{343.177}{328.11}{of}%
\TextAt{8}{357.244}{328.11}{the}%
\TextAt{8}{378.137}{328.11}{logarithm}%
\TextAt{8}{432.219}{328.11}{we}%
\TextAt{8}{450.183}{328.11}{make}%
\TextAt{8}{481.49}{328.11}{use}%
\TextAt{8}{502.443}{328.11}{of}%
\TextAt{8}{516.51}{328.11}{the}%
\TextAt{8}{79.226}{311.064}{following}%
\TextAt{8}{128.661}{311.064}{fact.}%
\TextAt{12}{79.226}{285.052}{Proposition}%
\TextAt{12}{151.165}{285.052}{(5.3).}%
\TextAt{8}{185.808}{285.052}{The}%
\TextAt{8}{209.19}{285.052}{multiplicative}%
\TextAt{8}{282.341}{285.052}{group}%
\TextAt{8}{314.845}{285.052}{of}%
\TextAt{8}{327.503}{285.052}{a}%
\TextAt{8}{336.587}{285.052}{local}%
\TextAt{8}{363.555}{285.052}{field}%
\TextAt{13}{388.24}{285.052}{K}%
\TextAt{8}{402.284}{285.052}{admits}%
\TextAt{8}{440.048}{285.052}{the}%
\TextAt{8}{459.532}{285.052}{decomposition}%
\TextAt{13}{244.552}{256.051}{K}%
\TextAt{3}{255.37}{260.987}{\Character{3}}%
\TextAt{9}{263.423}{256.051}{\Character{61}}%
\TextAt{9}{275.833}{256.051}{\Character{40}}%
\TextAt{13}{280.4}{256.051}{\Character{25}}%
\TextAt{9}{287.47}{256.051}{\Character{41}}%
\TextAt{16}{294.679}{256.051}{\Character{2}}%
\TextAt{13}{306.634}{256.051}{\Character{22}}%
\TextAt{21}{313.677}{254.257}{q}%
\TextAt{3}{317.744}{254.257}{\Character{0}}%
\TextAt{17}{324.33}{254.257}{1}%
\TextAt{16}{331.719}{256.051}{\Character{2}}%
\TextAt{13}{343.674}{256.051}{U}%
\TextAt{17}{352.873}{260.987}{\Character{40}1\Character{41}}%
\TextAt{13}{364.192}{256.051}{\Character{58}}%
\TextAt{8}{79.226}{227.049}{Here}%
\TextAt{13}{107.002}{227.049}{\Character{25}}%
\TextAt{8}{118.116}{227.049}{is}%
\TextAt{8}{130.028}{227.049}{a}%
\TextAt{8}{139.931}{227.049}{prime}%
\TextAt{8}{173.242}{227.049}{element,}%
\TextAt{9}{219.92}{227.049}{\Character{40}}%
\TextAt{13}{224.473}{227.049}{\Character{25}}%
\TextAt{9}{231.542}{227.049}{\Character{41}}%
\TextAt{9}{239.648}{227.049}{\Character{61}}%
\TextAt{16}{252.344}{227.049}{\Character{102}}%
\TextAt{13}{258.321}{227.049}{\Character{25}}%
\TextAt{21}{265.39}{231.388}{k}%
\TextAt{16}{274.083}{227.049}{\Character{106}}%
\TextAt{13}{280.976}{227.049}{k}%
\NativeMathAt{291.038}{227.049}{11.955}{\in}%
\TextAt{18}{302.58}{227.049}{Z}%
\TextAt{16}{310.55}{227.049}{\Character{103}}%
\TextAt{8}{316.528}{227.049}{,}%
\TextAt{13}{323.862}{227.049}{q}%
\TextAt{9}{333.053}{227.049}{\Character{61}}%
\TextAt{8}{345.73}{227.049}{\char35\relax }%
\TextAt{13}{355.485}{227.049}{\Character{20}}%
\TextAt{8}{366.269}{227.049}{is}%
\TextAt{8}{378.18}{227.049}{the}%
\TextAt{8}{398.484}{227.049}{number}%
\TextAt{8}{441.228}{227.049}{of}%
\TextAt{8}{454.706}{227.049}{elements}%
\TextAt{8}{502.71}{227.049}{in}%
\TextAt{8}{516.51}{227.049}{the}%
\TextAt{8}{79.226}{210.003}{residue}%
\TextAt{8}{118.953}{210.003}{class}%
\TextAt{8}{146.39}{210.003}{field}%
\TextAt{13}{171.747}{210.003}{\Character{20}}%
\TextAt{9}{181.807}{210.003}{\Character{61}}%
\TextAt{14}{194.232}{210.003}{o}%
\TextAt{9}{200.076}{210.003}{\Character{47}}%
\TextAt{14}{205.929}{210.003}{p}%
\TextAt{8}{211.912}{210.003}{,}%
\TextAt{8}{219.062}{210.003}{and}%
\TextAt{13}{241.824}{210.003}{U}%
\TextAt{17}{251.023}{214.342}{\Character{40}1\Character{41}}%
\TextAt{9}{265.663}{210.003}{\Character{61}}%
\TextAt{9}{278.072}{210.003}{1}%
\TextAt{9}{286.573}{210.003}{\Character{43}}%
\TextAt{14}{298.359}{210.003}{p}%
\TextAt{8}{308.24}{210.003}{is}%
\TextAt{8}{320.004}{210.003}{the}%
\TextAt{8}{340.16}{210.003}{group}%
\TextAt{8}{373.336}{210.003}{of}%
\TextAt{12}{387.18}{210.003}{principal}%
\TextAt{12}{443.321}{210.003}{units}%
\TextAt{8}{472.552}{210.003}{.}%
\TextAt{19}{79.226}{175.025}{Proof.}%
\TextAt{8}{116}{175.025}{For}%
\TextAt{8}{136.739}{175.025}{every}%
\TextAt{13}{167.4}{175.025}{\Character{11}}%
\NativeMathAt{178.243}{175.025}{11.955}{\in}%
\TextAt{13}{189.534}{175.025}{K}%
\TextAt{3}{200.351}{179.363}{\Character{3}}%
\TextAt{8}{205.084}{175.025}{,}%
\TextAt{8}{212.048}{175.025}{one}%
\TextAt{8}{233.277}{175.025}{has}%
\TextAt{8}{253.92}{175.025}{a}%
\TextAt{8}{263.445}{175.025}{unique}%
\TextAt{8}{301.255}{175.025}{representation}%
\TextAt{13}{377.513}{175.025}{\Character{11}}%
\TextAt{9}{388.356}{175.025}{\Character{61}}%
\TextAt{13}{400.781}{175.025}{\Character{25}}%
\TextAt{21}{407.851}{179.363}{n}%
\TextAt{13}{413.487}{175.025}{u}%
\TextAt{8}{423.816}{175.025}{with}%
\TextAt{13}{450.245}{175.025}{n}%
\NativeMathAt{460.554}{175.025}{11.955}{\in}%
\TextAt{18}{471.845}{175.025}{Z}%
\TextAt{8}{479.815}{175.025}{,}%
\TextAt{13}{486.779}{175.025}{u}%
\NativeMathAt{496.763}{175.025}{11.955}{\in}%
\TextAt{14}{508.054}{175.025}{o}%
\TextAt{3}{513.898}{179.363}{\Character{3}}%
\TextAt{8}{522.297}{175.025}{so}%
\TextAt{8}{79.226}{157.979}{that}%
\TextAt{13}{105.006}{157.979}{K}%
\TextAt{3}{115.824}{162.317}{\Character{3}}%
\TextAt{9}{124.577}{157.979}{\Character{61}}%
\TextAt{9}{137.692}{157.979}{\Character{40}}%
\TextAt{13}{142.255}{157.979}{\Character{25}}%
\TextAt{9}{149.324}{157.979}{\Character{41}}%
\TextAt{16}{156.813}{157.979}{\Character{2}}%
\TextAt{14}{169.048}{157.979}{o}%
\TextAt{3}{174.892}{162.317}{\Character{3}}%
\TextAt{8}{179.625}{157.979}{.}%
\TextAt{8}{189.306}{157.979}{Since}%
\TextAt{8}{220.274}{157.979}{the}%
\TextAt{8}{240.841}{157.979}{polynomial}%
\TextAt{13}{301.745}{157.979}{X}%
\TextAt{21}{312.4}{162.317}{q}%
\TextAt{3}{316.467}{162.317}{\Character{0}}%
\TextAt{17}{323.053}{162.317}{1}%
\TextAt{16}{330.722}{157.979}{\Character{0}}%
\TextAt{9}{342.958}{157.979}{1}%
\TextAt{8}{353.119}{157.979}{splits}%
\TextAt{8}{384.219}{157.979}{into}%
\TextAt{8}{408.372}{157.979}{linear}%
\TextAt{8}{441.301}{157.979}{factors}%
\TextAt{8}{479.825}{157.979}{over}%
\TextAt{13}{505.282}{157.979}{K}%
\TextAt{8}{520.408}{157.979}{by}%
\TextAt{8}{79.226}{140.933}{Hensel’s}%
\TextAt{8}{123.296}{140.933}{lemma,}%
\TextAt{14}{163.274}{140.933}{o}%
\TextAt{3}{169.119}{145.271}{\Character{3}}%
\TextAt{8}{176.508}{140.933}{contains}%
\TextAt{8}{221.188}{140.933}{the}%
\TextAt{8}{240.104}{140.933}{group}%
\TextAt{13}{272.04}{140.933}{\Character{22}}%
\TextAt{21}{279.083}{139.14}{q}%
\TextAt{3}{283.15}{139.14}{\Character{0}}%
\TextAt{17}{289.736}{139.14}{1}%
\TextAt{8}{297.126}{140.933}{of}%
\TextAt{9}{309.216}{140.933}{\Character{40}}%
\TextAt{13}{313.768}{140.933}{q}%
\TextAt{16}{319.508}{140.933}{\Character{0}}%
\TextAt{9}{328.927}{140.933}{1\Character{41}}%
\TextAt{8}{339.333}{140.933}{-th}%
\TextAt{8}{356.946}{140.933}{roots}%
\TextAt{8}{385.367}{140.933}{of}%
\TextAt{8}{397.457}{140.933}{unity.}%
\TextAt{8}{431.183}{140.933}{The}%
\TextAt{8}{453.997}{140.933}{homomorphism}%
\TextAt{14}{79.226}{123.887}{o}%
\TextAt{3}{85.07}{128.225}{\Character{3}}%
\NativeMathAt{94.343}{123.887}{11.955}{\rightarrow}%
\TextAt{13}{110.839}{123.887}{\Character{20}}%
\TextAt{3}{117.577}{128.225}{\Character{3}}%
\TextAt{8}{122.31}{123.887}{,}%
\TextAt{13}{130.354}{123.887}{u}%
\NativeArrowAt{141.557}{123.887}{11.955}{11.955}{\mapsto}%
\TextAt{9}{158.782}{123.887}{\Character{22}}%
\TextAt{13}{158.052}{123.887}{u}%
\TextAt{9}{169.255}{123.887}{\Character{61}}%
\TextAt{13}{182.9}{123.887}{u}%
\TextAt{8}{192.884}{123.887}{mod}%
\TextAt{14}{218.321}{123.887}{p}%
\TextAt{8}{224.305}{123.887}{,}%
\TextAt{8}{232.348}{123.887}{has}%
\TextAt{8}{253.938}{123.887}{kernel}%
\TextAt{13}{289.12}{123.887}{U}%
\TextAt{17}{298.319}{128.225}{\Character{40}1\Character{41}}%
\TextAt{8}{314.251}{123.887}{and}%
\TextAt{8}{337.73}{123.887}{maps}%
\TextAt{13}{369.074}{123.887}{\Character{22}}%
\TextAt{21}{376.117}{122.094}{q}%
\TextAt{3}{380.184}{122.094}{\Character{0}}%
\TextAt{17}{386.771}{122.094}{1}%
\TextAt{8}{396.116}{123.887}{bijectively}%
\TextAt{8}{452.758}{123.887}{onto}%
\TextAt{13}{479.823}{123.887}{\Character{20}}%
\TextAt{3}{486.562}{128.225}{\Character{3}}%
\TextAt{8}{491.294}{123.887}{.}%
\TextAt{8}{501.889}{123.887}{Hence}%
\TextAt{14}{79.226}{106.841}{o}%
\TextAt{3}{85.07}{111.179}{\Character{3}}%
\TextAt{9}{93.124}{106.841}{\Character{61}}%
\TextAt{13}{105.549}{106.841}{\Character{22}}%
\TextAt{21}{112.592}{105.048}{q}%
\TextAt{3}{116.659}{105.048}{\Character{0}}%
\TextAt{17}{123.245}{105.048}{1}%
\TextAt{16}{130.634}{106.841}{\Character{2}}%
\TextAt{13}{142.589}{106.841}{U}%
\TextAt{17}{151.788}{111.179}{\Character{40}1\Character{41}}%
\TextAt{8}{163.107}{106.841}{.}%
\Vector{q 0 G 0 g 0.399 w 452.334 -667.159 m 452.334 -659.089 l S Q}%
\Vector{q 0 G 0 g 0.399 w 452.534 -659.289 m 459.707 -659.289 l S Q}%
\Vector{q 0 G 0 g 0.399 w 452.534 -666.96 m 459.707 -666.96 l S Q}%
\Vector{q 0 G 0 g 0.399 w 459.906 -667.159 m 459.906 -659.089 l S Q}%
\TextAt{8}{297.211}{43.033}{124}%
\Destination{78.23}{809.815}{nav.page.124}%
\Destination{79.36}{485.455}{nav.print.136}%
\Destination{79.36}{299.055}{nav.number.2.5.3}%
%
%
\Hotlink{102.951}{602.555}{19.22}{11.047}{nav.cite.N018}%
\Hotlink{348.455}{568.455}{18.574}{11.047}{nav.print.127}%
%
%
\Hotlink{201.362}{434.155}{25.078}{11.047}{nav.cite.N137}%
\Hotlink{395.577}{517.255}{25.078}{11.047}{nav.number.2.4.4}%
\Destination{79.226}{427}{errata.192}%
\Hotlink{388.16}{602.5}{21.722}{11.047}{nav.number.2.6.8}%
}

% Page 137 (body page 125)
\Page{%
\TextAt{2}{79.226}{796.943}{II.}%
\TextAt{2}{94.073}{796.943}{The}%
\TextAt{2}{116.492}{796.943}{Theory}%
\TextAt{2}{154.401}{796.943}{of}%
\TextAt{2}{166.826}{796.943}{Valuations}%
\Vector{q 0 G 0 g 0.299 w 7.226 18.729 m 460.769 18.729 l S Q}%
\TextAt{12}{79.226}{761.124}{Proposition}%
\TextAt{12}{152.126}{761.124}{(5.4).}%
\TextAt{8}{187.606}{761.124}{For}%
\TextAt{8}{208.74}{761.124}{a}%
\TextAt{13}{218.66}{761.124}{p}%
\TextAt{8}{224.535}{761.124}{-adic}%
\TextAt{8}{253.308}{761.124}{number}%
\TextAt{8}{296.068}{761.124}{field}%
\TextAt{13}{321.59}{761.124}{K}%
\TextAt{8}{336.469}{761.124}{there}%
\TextAt{8}{366.545}{761.124}{is}%
\TextAt{8}{378.473}{761.124}{a}%
\TextAt{8}{388.393}{761.124}{uniquely}%
\TextAt{8}{436.031}{761.124}{determined}%
\TextAt{8}{497.322}{761.124}{contin-}%
\TextAt{8}{79.226}{744.078}{uous}%
\TextAt{8}{106.603}{744.078}{homomorphism}%
\TextAt{8}{267.432}{727.032}{log}%
\TextAt{9}{285.721}{727.032}{\Character{58}}%
\TextAt{13}{292.293}{727.032}{K}%
\TextAt{3}{303.111}{731.968}{\Character{3}}%
\NativeArrowAt{311.164}{727.032}{11.955}{19.26}{\longrightarrow}%
\TextAt{13}{333.746}{727.032}{K}%
\TextAt{8}{79.226}{704.029}{such}%
\TextAt{8}{105.623}{704.029}{that}%
\TextAt{8}{130.992}{704.029}{log}%
\TextAt{13}{147.952}{704.029}{p}%
\TextAt{9}{157.148}{704.029}{\Character{61}}%
\TextAt{9}{169.558}{704.029}{0}%
\TextAt{8}{175.427}{704.029}{,}%
\TextAt{8}{182.576}{704.029}{which}%
\TextAt{8}{216.062}{704.029}{on}%
\TextAt{8}{232.322}{704.029}{principal}%
\TextAt{8}{281.099}{704.029}{units}%
\TextAt{9}{310.425}{704.029}{\Character{40}1}%
\TextAt{9}{323.468}{704.029}{\Character{43}}%
\TextAt{13}{335.248}{704.029}{x}%
\TextAt{9}{341.9}{704.029}{\Character{41}}%
\NativeMathAt{349.773}{704.029}{11.955}{\in}%
\TextAt{13}{361.064}{704.029}{U}%
\TextAt{17}{370.263}{708.367}{\Character{40}1\Character{41}}%
\TextAt{8}{385.479}{704.029}{is}%
\TextAt{8}{397.243}{704.029}{given}%
\TextAt{8}{427.813}{704.029}{by}%
\TextAt{8}{444.072}{704.029}{the}%
\TextAt{8}{464.228}{704.029}{series}%
\TextAt{8}{223.112}{671.737}{log}%
\TextAt{9}{238.08}{671.737}{\Character{40}1}%
\TextAt{9}{251.123}{671.737}{\Character{43}}%
\TextAt{13}{262.904}{671.737}{x}%
\TextAt{9}{269.556}{671.737}{\Character{41}}%
\TextAt{9}{277.411}{671.737}{\Character{61}}%
\TextAt{13}{289.854}{671.737}{x}%
\TextAt{16}{299.163}{671.737}{\Character{0}}%
\TextAt{13}{312.314}{679.825}{x}%
\TextAt{17}{318.966}{684.163}{2}%
\Vector{q 0 G 0 g 0.478 w 240.314 -99.274 m 251.698 -99.274 l S Q}%
\TextAt{9}{315.079}{663.536}{2}%
\TextAt{9}{327.55}{671.737}{\Character{43}}%
\TextAt{13}{340.507}{679.825}{x}%
\TextAt{17}{347.159}{684.163}{3}%
\Vector{q 0 G 0 g 0.478 w 268.507 -99.274 m 279.892 -99.274 l S Q}%
\TextAt{9}{343.273}{663.536}{3}%
\TextAt{16}{355.744}{671.737}{\Character{0}}%
\TextAt{16}{367.687}{671.737}{\Character{1}}%
\TextAt{16}{372.983}{671.737}{\Character{1}}%
\TextAt{16}{378.292}{671.737}{\Character{1}}%
\TextAt{13}{385.632}{671.737}{\Character{58}}%
\TextAt{19}{79.226}{626.706}{Proof.}%
\TextAt{8}{116}{626.706}{By}%
\TextAt{8}{134.101}{626.706}{§4,}%
\TextAt{8}{152.576}{626.706}{we}%
\TextAt{8}{169.541}{626.706}{can}%
\TextAt{8}{190.738}{626.706}{think}%
\TextAt{8}{221.367}{626.706}{of}%
\TextAt{8}{234.435}{626.706}{the}%
\TextAt{13}{254.328}{626.706}{p}%
\TextAt{8}{260.204}{626.706}{-adic}%
\TextAt{8}{288.549}{626.706}{valuation}%
\TextAt{13}{339.359}{626.706}{v}%
\TextAt{21}{345.019}{624.912}{p}%
\TextAt{8}{353.414}{626.706}{of}%
\TextAt{18}{366.481}{626.706}{Q}%
\TextAt{21}{375.78}{624.912}{p}%
\TextAt{8}{384.175}{626.706}{as}%
\TextAt{8}{398.283}{626.706}{extended}%
\TextAt{8}{447.766}{626.706}{to}%
\TextAt{13}{461.813}{626.706}{K}%
\TextAt{8}{472.631}{626.706}{.}%
\TextAt{8}{480.992}{626.706}{Observing}%
\TextAt{8}{79.226}{609.66}{that}%
\TextAt{13}{104.86}{609.66}{v}%
\TextAt{21}{110.519}{607.867}{p}%
\TextAt{9}{115.28}{609.66}{\Character{40}}%
\TextAt{13}{119.832}{609.66}{x}%
\TextAt{9}{126.485}{609.66}{\Character{41}}%
\TextAt{13}{134.809}{609.66}{\Character{62}}%
\TextAt{9}{147.686}{609.66}{0}%
\TextAt{8}{153.539}{609.66}{,}%
\TextAt{8}{161.019}{609.66}{so}%
\TextAt{8}{175.655}{609.66}{that}%
\TextAt{13}{201.288}{609.66}{c}%
\TextAt{9}{210.099}{609.66}{\Character{61}}%
\TextAt{13}{222.976}{609.66}{p}%
\TextAt{21}{228.851}{613.998}{v}%
\TextAt{22}{232.974}{613.002}{p}%
\TextAt{17}{237.313}{613.998}{\Character{40}}%
\TextAt{21}{240.606}{613.998}{x}%
\TextAt{17}{245.373}{613.998}{\Character{41}}%
\TextAt{13}{252.937}{609.66}{\Character{62}}%
\TextAt{9}{265.814}{609.66}{1}%
\TextAt{8}{271.667}{609.66}{,}%
\TextAt{8}{279.147}{609.66}{and}%
\TextAt{13}{302.174}{609.66}{p}%
\TextAt{21}{308.05}{613.998}{v}%
\TextAt{22}{312.173}{613.002}{p}%
\TextAt{17}{316.512}{613.998}{\Character{40}}%
\TextAt{21}{319.805}{613.998}{\Character{23}}%
\TextAt{17}{324.489}{613.998}{\Character{41}}%
\NativeMathAt{332.053}{609.66}{11.955}{\leq}%
\TextAt{13}{345.123}{609.66}{\Character{23}}%
\TextAt{8}{351.661}{609.66}{,}%
\TextAt{8}{359.142}{609.66}{giving}%
\TextAt{13}{394.208}{609.66}{v}%
\TextAt{21}{399.867}{607.867}{p}%
\TextAt{9}{404.628}{609.66}{\Character{40}}%
\TextAt{13}{409.181}{609.66}{\Character{23}}%
\TextAt{9}{415.719}{609.66}{\Character{41}}%
\NativeMathAt{424.043}{609.66}{11.955}{\leq}%
\TextAt{8}{437.114}{609.66}{ln}%
\TextAt{13}{448.862}{609.66}{\Character{23}}%
\TextAt{9}{455.4}{609.66}{\Character{47}}%
\TextAt{8}{463.245}{609.66}{ln}%
\TextAt{13}{474.993}{609.66}{p}%
\TextAt{8}{485.031}{609.66}{(with}%
\TextAt{8}{516.511}{609.66}{the}%
\TextAt{8}{79.226}{592.614}{usual}%
\TextAt{8}{109.855}{592.614}{logarithm),}%
\TextAt{8}{171.006}{592.614}{we}%
\TextAt{8}{188.233}{592.614}{compute}%
\TextAt{8}{235.707}{592.614}{the}%
\TextAt{8}{255.864}{592.614}{valuation}%
\TextAt{8}{306.936}{592.614}{of}%
\TextAt{8}{320.266}{592.614}{the}%
\TextAt{8}{340.423}{592.614}{terms}%
\TextAt{13}{373.001}{592.614}{x}%
\TextAt{21}{379.653}{596.952}{\Character{23}}%
\TextAt{9}{384.835}{592.614}{\Character{47}}%
\TextAt{13}{390.688}{592.614}{\Character{23}}%
\TextAt{8}{401.123}{592.614}{of}%
\TextAt{8}{414.453}{592.614}{the}%
\TextAt{8}{434.609}{592.614}{series,}%
\TextAt{13}{168.99}{560.537}{v}%
\TextAt{21}{174.65}{558.744}{p}%
\TextAt{20}{181.403}{577.394}{\Character{18}}%
\TextAt{13}{191.399}{568.625}{x}%
\TextAt{21}{198.051}{572.963}{\Character{23}}%
\Vector{q 0 G 0 g 0.478 w 119.399 -210.474 m 131.233 -210.474 l S Q}%
\TextAt{13}{194.047}{552.337}{\Character{23}}%
\TextAt{20}{204.428}{577.394}{\Character{19}}%
\TextAt{9}{216.55}{560.537}{\Character{61}}%
\TextAt{13}{228.975}{560.537}{\Character{23}}%
\TextAt{13}{235.503}{560.537}{v}%
\TextAt{21}{241.172}{558.744}{p}%
\TextAt{9}{245.933}{560.537}{\Character{40}}%
\TextAt{13}{250.486}{560.537}{x}%
\TextAt{9}{257.138}{560.537}{\Character{41}}%
\TextAt{16}{264.347}{560.537}{\Character{0}}%
\TextAt{13}{276.302}{560.537}{v}%
\TextAt{21}{281.961}{558.744}{p}%
\TextAt{9}{286.722}{560.537}{\Character{40}}%
\TextAt{13}{291.274}{560.537}{\Character{23}}%
\TextAt{9}{297.812}{560.537}{\Character{41}}%
\NativeMathAt{305.685}{560.537}{11.955}{\geq}%
\TextAt{13}{318.305}{560.537}{\Character{23}}%
\TextAt{8}{326.457}{568.625}{ln}%
\TextAt{13}{338.205}{568.625}{c}%
\Vector{q 0 G 0 g 0.478 w 254.038 -210.474 m 271.661 -210.474 l S Q}%
\TextAt{8}{326.038}{552.337}{ln}%
\TextAt{13}{337.786}{552.337}{p}%
\TextAt{16}{347.513}{560.537}{\Character{0}}%
\TextAt{8}{360.664}{568.625}{ln}%
\TextAt{13}{372.412}{568.625}{\Character{23}}%
\Vector{q 0 G 0 g 0.478 w 288.664 -210.474 m 306.95 -210.474 l S Q}%
\TextAt{8}{360.996}{552.337}{ln}%
\TextAt{13}{372.743}{552.337}{p}%
\TextAt{9}{383.466}{560.537}{\Character{61}}%
\TextAt{8}{397.087}{568.625}{ln}%
\TextAt{9}{406.843}{568.625}{\Character{40}}%
\TextAt{13}{411.395}{568.625}{c}%
\TextAt{21}{416.433}{572.963}{\Character{23}}%
\TextAt{9}{421.615}{568.625}{\Character{47}}%
\TextAt{13}{427.468}{568.625}{\Character{23}}%
\TextAt{9}{434.006}{568.625}{\Character{41}}%
\Vector{q 0 G 0 g 0.478 w 325.087 -210.474 m 366.558 -210.474 l S Q}%
\TextAt{8}{409.011}{552.337}{ln}%
\TextAt{13}{420.759}{552.337}{p}%
\TextAt{13}{439.754}{560.537}{\Character{58}}%
\TextAt{8}{79.226}{528.473}{This}%
\TextAt{8}{106.333}{528.473}{shows}%
\TextAt{8}{140.338}{528.473}{that}%
\TextAt{13}{166.094}{528.473}{x}%
\TextAt{21}{172.746}{532.811}{\Character{23}}%
\TextAt{9}{177.928}{528.473}{\Character{47}}%
\TextAt{13}{183.781}{528.473}{\Character{23}}%
\TextAt{8}{194.603}{528.473}{is}%
\TextAt{8}{206.754}{528.473}{a}%
\TextAt{8}{216.896}{528.473}{nullsequence,}%
\TextAt{8}{288.322}{528.473}{i.e.,}%
\TextAt{8}{310.91}{528.473}{the}%
\TextAt{8}{331.454}{528.473}{logarithm}%
\TextAt{8}{385.185}{528.473}{series}%
\TextAt{8}{416.906}{528.473}{converges.}%
\TextAt{8}{475.041}{528.473}{It}%
\TextAt{8}{488.101}{528.473}{defines}%
\TextAt{8}{526.911}{528.473}{a}%
\TextAt{8}{79.226}{511.427}{homomorphism}%
\TextAt{8}{161.896}{511.427}{because}%
\TextAt{8}{192.372}{483.636}{log}%
\TextAt{9}{207.339}{483.636}{\Character{40}\Character{40}1}%
\TextAt{9}{224.925}{483.636}{\Character{43}}%
\TextAt{13}{236.715}{483.636}{x}%
\TextAt{9}{243.367}{483.636}{\Character{41}\Character{40}1}%
\TextAt{9}{260.953}{483.636}{\Character{43}}%
\TextAt{13}{272.743}{483.636}{y}%
\TextAt{9}{278.879}{483.636}{\Character{41}\Character{41}}%
\TextAt{9}{291.276}{483.636}{\Character{61}}%
\TextAt{8}{303.73}{483.636}{log}%
\TextAt{9}{318.698}{483.636}{\Character{40}1}%
\TextAt{9}{331.741}{483.636}{\Character{43}}%
\TextAt{13}{343.522}{483.636}{x}%
\TextAt{9}{350.174}{483.636}{\Character{41}}%
\TextAt{9}{357.371}{483.636}{\Character{43}}%
\TextAt{8}{369.144}{483.636}{log}%
\TextAt{9}{384.112}{483.636}{\Character{40}1}%
\TextAt{9}{397.155}{483.636}{\Character{43}}%
\TextAt{13}{408.935}{483.636}{y}%
\TextAt{9}{415.072}{483.636}{\Character{41}}%
\TextAt{8}{79.226}{455.846}{is}%
\TextAt{8}{90.003}{455.846}{an}%
\TextAt{8}{105.274}{455.846}{identity}%
\TextAt{8}{147.54}{455.846}{of}%
\TextAt{8}{159.883}{455.846}{formal}%
\TextAt{8}{195.646}{455.846}{power}%
\TextAt{8}{228.802}{455.846}{series}%
\TextAt{8}{259.149}{455.846}{and}%
\TextAt{8}{280.925}{455.846}{all}%
\TextAt{8}{296.196}{455.846}{series}%
\TextAt{8}{326.543}{455.846}{in}%
\TextAt{8}{339.208}{455.846}{it}%
\TextAt{8}{349.925}{455.846}{converge}%
\TextAt{8}{396.746}{455.846}{provided}%
\TextAt{9}{443.89}{455.846}{1}%
\TextAt{9}{450.37}{455.846}{\Character{43}}%
\TextAt{13}{460.123}{455.846}{x\Character{59}}%
\TextAt{9}{472.019}{455.846}{1}%
\TextAt{9}{478.499}{455.846}{\Character{43}}%
\TextAt{13}{488.251}{455.846}{y}%
\NativeMathAt{497.709}{455.846}{11.955}{\in}%
\TextAt{13}{509}{455.846}{U}%
\TextAt{17}{518.199}{460.184}{\Character{40}1\Character{41}}%
\TextAt{8}{529.518}{455.846}{.}%
\TextAt{8}{97.159}{435.811}{For}%
\TextAt{8}{118.128}{435.811}{every}%
\TextAt{13}{149.02}{435.811}{\Character{11}}%
\NativeMathAt{159.863}{435.811}{11.955}{\in}%
\TextAt{13}{171.154}{435.811}{K}%
\TextAt{3}{181.972}{440.15}{\Character{3}}%
\TextAt{8}{186.704}{435.811}{,}%
\TextAt{8}{193.853}{435.811}{choosing}%
\TextAt{8}{241.399}{435.811}{a}%
\TextAt{8}{251.154}{435.811}{prime}%
\TextAt{8}{284.318}{435.811}{element}%
\TextAt{13}{327.56}{435.811}{\Character{25}}%
\TextAt{8}{334.629}{435.811}{,}%
\TextAt{8}{341.778}{435.811}{we}%
\TextAt{8}{359.006}{435.811}{have}%
\TextAt{8}{386}{435.811}{a}%
\TextAt{8}{395.756}{435.811}{unique}%
\TextAt{8}{433.797}{435.811}{representation}%
\TextAt{13}{258.256}{408.021}{\Character{11}}%
\TextAt{9}{269.099}{408.021}{\Character{61}}%
\TextAt{13}{281.524}{408.021}{\Character{25}}%
\TextAt{21}{288.593}{412.957}{v}%
\TextAt{28}{292.717}{411.961}{p}%
\TextAt{17}{296.8}{412.957}{\Character{40}}%
\TextAt{21}{300.093}{412.957}{\Character{11}}%
\TextAt{17}{305.532}{412.957}{\Character{41}}%
\TextAt{13}{309.323}{408.021}{\Character{33}}%
\TextAt{9}{317.041}{408.021}{\Character{40}}%
\TextAt{13}{321.594}{408.021}{\Character{11}}%
\TextAt{9}{329.115}{408.021}{\Character{41}}%
\TextAt{16}{333.668}{408.021}{\Character{104}}%
\TextAt{13}{338.317}{408.021}{\Character{11}}%
\TextAt{16}{345.839}{408.021}{\Character{105}}%
\TextAt{13}{350.488}{408.021}{\Character{59}}%
\TextAt{8}{79.226}{380.231}{where}%
\TextAt{13}{114.421}{380.231}{v}%
\TextAt{27}{120.081}{378.437}{p}%
\TextAt{9}{130.25}{380.23}{\Character{61}}%
\TextAt{13}{145.037}{380.23}{ev}%
\TextAt{21}{156.122}{378.437}{p}%
\TextAt{8}{166.166}{380.231}{is}%
\TextAt{8}{179.316}{380.231}{the}%
\TextAt{8}{200.858}{380.231}{normalized}%
\TextAt{8}{262.08}{380.231}{valuation}%
\TextAt{8}{314.538}{380.231}{of}%
\TextAt{13}{329.254}{380.231}{K}%
\TextAt{8}{340.072}{380.231}{,}%
\TextAt{13}{348.953}{380.231}{\Character{33}}%
\TextAt{9}{356.672}{380.23}{\Character{40}}%
\TextAt{13}{361.224}{380.231}{\Character{11}}%
\TextAt{9}{368.746}{380.23}{\Character{41}}%
\NativeMathAt{378.98}{380.231}{11.955}{\in}%
\TextAt{13}{392.633}{380.23}{\Character{22}}%
\TextAt{21}{399.676}{378.437}{q}%
\TextAt{3}{403.743}{378.437}{\Character{0}}%
\TextAt{17}{410.329}{378.437}{1}%
\TextAt{8}{415.062}{380.231}{,}%
\TextAt{16}{423.943}{380.231}{\Character{104}}%
\TextAt{13}{428.593}{380.23}{\Character{11}}%
\TextAt{16}{436.114}{380.231}{\Character{105}}%
\NativeMathAt{446.431}{380.231}{11.955}{\in}%
\TextAt{13}{460.099}{380.23}{U}%
\TextAt{17}{469.298}{384.569}{\Character{40}1\Character{41}}%
\TextAt{8}{480.617}{380.231}{.}%
\TextAt{8}{493.223}{380.231}{As}%
\TextAt{8}{511.896}{380.231}{sug-}%
\TextAt{8}{79.226}{363.185}{gested}%
\TextAt{8}{115.686}{363.185}{by}%
\TextAt{8}{132.575}{363.185}{the}%
\TextAt{8}{153.361}{363.185}{equation}%
\TextAt{13}{201.8}{363.185}{p}%
\TextAt{9}{212.07}{363.185}{\Character{61}}%
\TextAt{13}{225.569}{363.185}{\Character{25}}%
\TextAt{21}{232.638}{367.523}{e}%
\TextAt{13}{237.065}{363.185}{\Character{33}}%
\TextAt{9}{244.783}{363.185}{\Character{40}}%
\TextAt{13}{249.335}{363.185}{p}%
\TextAt{9}{255.211}{363.185}{\Character{41}}%
\TextAt{16}{259.763}{363.185}{\Character{104}}%
\TextAt{13}{264.412}{363.185}{p}%
\TextAt{16}{270.287}{363.185}{\Character{105}}%
\TextAt{8}{274.937}{363.185}{,}%
\TextAt{8}{282.873}{363.185}{we}%
\TextAt{8}{300.731}{363.185}{define}%
\TextAt{8}{335.17}{363.185}{log}%
\TextAt{13}{352.131}{363.185}{\Character{25}}%
\TextAt{9}{363.594}{363.185}{\Character{61}}%
\TextAt{16}{377.094}{363.185}{\Character{0}}%
\TextAt{17}{387.588}{367.892}{1}%
\Vector{q 0 G 0 g 0.478 w 315.588 -407.827 m 319.822 -407.827 l S Q}%
\TextAt{21}{387.741}{359.062}{e}%
\TextAt{8}{395.01}{363.185}{log}%
\TextAt{16}{409.978}{363.185}{\Character{104}}%
\TextAt{13}{414.627}{363.185}{p}%
\TextAt{16}{420.502}{363.185}{\Character{105}}%
\TextAt{8}{429.679}{363.185}{and}%
\TextAt{8}{453.072}{363.185}{thus}%
\TextAt{8}{479.453}{363.185}{obtain}%
\TextAt{8}{516.511}{363.185}{the}%
\TextAt{8}{79.226}{346.139}{homomorphism}%
\TextAt{8}{161.896}{346.139}{log}%
\TextAt{9}{180.185}{346.139}{\Character{58}}%
\TextAt{13}{186.757}{346.139}{K}%
\TextAt{3}{197.575}{350.477}{\Character{3}}%
\NativeMathAt{205.628}{346.139}{11.955}{\rightarrow}%
\TextAt{13}{220.904}{346.139}{K}%
\TextAt{8}{235.619}{346.139}{by}%
\TextAt{8}{234.757}{318.348}{log}%
\TextAt{13}{251.718}{318.348}{\Character{11}}%
\TextAt{9}{262.56}{318.348}{\Character{61}}%
\TextAt{13}{274.986}{318.348}{v}%
\TextAt{27}{280.645}{316.555}{p}%
\TextAt{9}{285.132}{318.348}{\Character{40}}%
\TextAt{13}{289.684}{318.348}{\Character{11}}%
\TextAt{9}{297.206}{318.348}{\Character{41}}%
\TextAt{8}{303.751}{318.348}{log}%
\TextAt{13}{320.711}{318.348}{\Character{25}}%
\TextAt{9}{330.437}{318.348}{\Character{43}}%
\TextAt{8}{342.198}{318.348}{log}%
\TextAt{16}{357.166}{318.348}{\Character{104}}%
\TextAt{13}{361.816}{318.348}{\Character{11}}%
\TextAt{16}{369.337}{318.348}{\Character{105}}%
\TextAt{13}{373.987}{318.348}{\Character{58}}%
\TextAt{8}{79.226}{290.558}{It}%
\TextAt{8}{92.755}{290.558}{is}%
\TextAt{8}{105.376}{290.558}{obviously}%
\TextAt{8}{158.011}{290.558}{continuous}%
\TextAt{8}{217.472}{290.558}{and}%
\TextAt{8}{241.091}{290.558}{has}%
\TextAt{8}{262.822}{290.558}{the}%
\TextAt{8}{283.835}{290.558}{property}%
\TextAt{8}{332.5}{290.558}{that}%
\TextAt{8}{358.726}{290.558}{log}%
\TextAt{13}{375.686}{290.558}{p}%
\TextAt{9}{386.342}{290.558}{\Character{61}}%
\TextAt{9}{400.21}{290.558}{0}%
\TextAt{8}{406.081}{290.558}{.}%
\TextAt{8}{417.099}{290.558}{If}%
\TextAt{13}{429.648}{290.558}{\Character{21}}%
\TextAt{9}{441.258}{290.558}{\Character{58}}%
\TextAt{13}{449.29}{290.558}{K}%
\TextAt{3}{460.108}{294.896}{\Character{3}}%
\NativeMathAt{469.621}{290.558}{11.955}{\rightarrow}%
\TextAt{13}{486.357}{290.558}{K}%
\TextAt{8}{501.929}{290.558}{is}%
\TextAt{8}{514.55}{290.558}{any}%
\TextAt{8}{79.226}{273.512}{continuation}%
\TextAt{8}{146.35}{273.512}{of}%
\TextAt{8}{159.149}{273.512}{log}%
\TextAt{9}{177.438}{273.512}{\Character{58}}%
\TextAt{13}{184.01}{273.512}{U}%
\TextAt{17}{193.209}{277.85}{\Character{40}1\Character{41}}%
\NativeMathAt{207.849}{273.512}{11.955}{\rightarrow}%
\TextAt{13}{223.125}{273.512}{K}%
\TextAt{8}{237.309}{273.512}{such}%
\TextAt{8}{263.175}{273.512}{that}%
\TextAt{13}{288.013}{273.512}{\Character{21}}%
\TextAt{9}{294.842}{273.512}{\Character{40}}%
\TextAt{13}{299.394}{273.512}{p}%
\TextAt{9}{305.269}{273.512}{\Character{41}}%
\TextAt{9}{313.124}{273.512}{\Character{61}}%
\TextAt{9}{325.545}{273.512}{0}%
\TextAt{8}{331.421}{273.512}{,}%
\TextAt{8}{338.145}{273.512}{then}%
\TextAt{8}{364.275}{273.512}{we}%
\TextAt{8}{380.971}{273.512}{find}%
\TextAt{8}{403.848}{273.512}{that}%
\TextAt{13}{428.686}{273.512}{\Character{21}}%
\TextAt{9}{435.515}{273.512}{\Character{40}}%
\TextAt{13}{440.067}{273.512}{\Character{24}}%
\TextAt{9}{445.758}{273.512}{\Character{41}}%
\TextAt{9}{453.613}{273.512}{\Character{61}}%
\TextAt{17}{472.579}{278.219}{1}%
\Vector{q 0 G 0 g 0.478 w 395.252 -497.499 m 410.139 -497.499 l S Q}%
\TextAt{21}{467.252}{269.389}{q}%
\TextAt{3}{471.319}{269.389}{\Character{0}}%
\TextAt{17}{477.905}{269.389}{1}%
\TextAt{13}{483.335}{273.512}{\Character{21}}%
\TextAt{9}{490.163}{273.512}{\Character{40}}%
\TextAt{13}{494.716}{273.512}{\Character{24}}%
\TextAt{21}{500.406}{277.85}{q}%
\TextAt{3}{504.473}{277.85}{\Character{0}}%
\TextAt{17}{511.059}{277.85}{1}%
\TextAt{9}{515.792}{273.512}{\Character{41}}%
\TextAt{9}{523.647}{273.512}{\Character{61}}%
\TextAt{9}{79.226}{256.466}{0}%
\TextAt{8}{90.821}{256.466}{for}%
\TextAt{8}{110.55}{256.466}{each}%
\TextAt{13}{138.732}{256.466}{\Character{24}}%
\NativeMathAt{150.886}{256.466}{11.955}{\in}%
\TextAt{13}{165.321}{256.466}{\Character{22}}%
\TextAt{21}{172.364}{254.673}{q}%
\TextAt{3}{176.43}{254.673}{\Character{0}}%
\TextAt{17}{183.017}{254.673}{1}%
\TextAt{8}{187.749}{256.466}{.}%
\TextAt{8}{201.731}{256.466}{It}%
\TextAt{8}{216.248}{256.466}{follows}%
\TextAt{8}{256.528}{256.466}{that}%
\TextAt{9}{283.741}{256.466}{0}%
\TextAt{9}{296.043}{256.466}{\Character{61}}%
\TextAt{13}{311.627}{256.466}{e\Character{21}}%
\TextAt{9}{323.881}{256.466}{\Character{40}}%
\TextAt{13}{328.433}{256.466}{\Character{25}}%
\TextAt{9}{335.503}{256.466}{\Character{41}}%
\TextAt{9}{343.955}{256.466}{\Character{43}}%
\TextAt{13}{356.987}{256.466}{\Character{21}}%
\TextAt{9}{363.816}{256.466}{\Character{40}}%
\TextAt{16}{368.368}{256.466}{\Character{104}}%
\TextAt{13}{373.017}{256.466}{p}%
\TextAt{16}{378.893}{256.466}{\Character{105}}%
\TextAt{9}{383.542}{256.466}{\Character{41}}%
\TextAt{9}{394.541}{256.466}{\Character{61}}%
\TextAt{13}{410.127}{256.466}{e\Character{21}}%
\TextAt{9}{422.381}{256.466}{\Character{40}}%
\TextAt{13}{426.933}{256.466}{\Character{25}}%
\TextAt{9}{434.002}{256.466}{\Character{41}}%
\TextAt{9}{442.454}{256.466}{\Character{43}}%
\TextAt{8}{455.487}{256.466}{log}%
\TextAt{16}{470.455}{256.466}{\Character{104}}%
\TextAt{13}{475.104}{256.466}{p}%
\TextAt{16}{480.98}{256.466}{\Character{105}}%
\TextAt{8}{485.629}{256.466}{,}%
\TextAt{8}{495.084}{256.466}{so}%
\TextAt{8}{511.298}{256.466}{that}%
\TextAt{13}{79.226}{239.42}{\Character{21}}%
\TextAt{9}{86.055}{239.42}{\Character{40}}%
\TextAt{13}{90.607}{239.42}{\Character{25}}%
\TextAt{9}{97.676}{239.42}{\Character{41}}%
\TextAt{9}{106.045}{239.42}{\Character{61}}%
\TextAt{8}{119.001}{239.42}{log}%
\TextAt{13}{135.962}{239.42}{\Character{25}}%
\TextAt{8}{143.031}{239.42}{,}%
\TextAt{8}{150.557}{239.42}{and}%
\TextAt{8}{173.62}{239.42}{thus}%
\TextAt{13}{199.673}{239.42}{\Character{21}}%
\TextAt{9}{206.501}{239.42}{\Character{40}}%
\TextAt{13}{211.054}{239.42}{\Character{11}}%
\TextAt{9}{218.575}{239.42}{\Character{41}}%
\TextAt{9}{226.944}{239.42}{\Character{61}}%
\TextAt{13}{239.901}{239.42}{v}%
\TextAt{27}{245.56}{237.627}{p}%
\TextAt{9}{250.047}{239.42}{\Character{40}}%
\TextAt{13}{254.599}{239.42}{\Character{11}}%
\TextAt{9}{262.121}{239.42}{\Character{41}}%
\TextAt{13}{266.673}{239.42}{\Character{21}}%
\TextAt{9}{273.502}{239.42}{\Character{40}}%
\TextAt{13}{278.054}{239.42}{\Character{25}}%
\TextAt{9}{285.123}{239.42}{\Character{41}}%
\TextAt{9}{292.523}{239.42}{\Character{43}}%
\TextAt{13}{304.504}{239.42}{\Character{21}}%
\TextAt{9}{311.332}{239.42}{\Character{40}}%
\TextAt{16}{315.885}{239.42}{\Character{104}}%
\TextAt{13}{320.534}{239.42}{\Character{11}}%
\TextAt{16}{328.056}{239.42}{\Character{105}}%
\TextAt{9}{332.705}{239.42}{\Character{41}}%
\TextAt{9}{341.074}{239.42}{\Character{61}}%
\TextAt{13}{354.03}{239.42}{v}%
\TextAt{27}{359.689}{237.627}{p}%
\TextAt{9}{364.176}{239.42}{\Character{40}}%
\TextAt{13}{368.729}{239.42}{\Character{11}}%
\TextAt{9}{376.25}{239.42}{\Character{41}}%
\TextAt{8}{382.795}{239.42}{log}%
\TextAt{13}{399.755}{239.42}{\Character{25}}%
\TextAt{9}{409.687}{239.42}{\Character{43}}%
\TextAt{8}{421.653}{239.42}{log}%
\TextAt{16}{436.621}{239.42}{\Character{104}}%
\TextAt{13}{441.27}{239.42}{\Character{11}}%
\TextAt{16}{448.792}{239.42}{\Character{105}}%
\TextAt{9}{457.275}{239.42}{\Character{61}}%
\TextAt{8}{470.214}{239.42}{log}%
\TextAt{13}{487.174}{239.42}{\Character{11}}%
\TextAt{8}{494.696}{239.42}{,}%
\TextAt{8}{502.222}{239.42}{for}%
\TextAt{8}{520.408}{239.42}{all}%
\TextAt{13}{79.226}{222.374}{\Character{11}}%
\NativeMathAt{90.069}{222.374}{11.955}{\in}%
\TextAt{13}{101.36}{222.374}{K}%
\TextAt{3}{112.177}{226.712}{\Character{3}}%
\TextAt{8}{116.91}{222.374}{.}%
\TextAt{8}{125.358}{222.374}{log}%
\TextAt{8}{144.223}{222.374}{is}%
\TextAt{8}{155.987}{222.374}{therefore}%
\TextAt{8}{205.087}{222.374}{uniquely}%
\TextAt{8}{252.561}{222.374}{determined}%
\TextAt{8}{313.688}{222.374}{and}%
\TextAt{8}{336.45}{222.374}{independent}%
\TextAt{8}{402.778}{222.374}{of}%
\TextAt{8}{416.108}{222.374}{the}%
\TextAt{8}{436.264}{222.374}{choice}%
\TextAt{8}{471.054}{222.374}{of}%
\TextAt{13}{484.384}{222.374}{\Character{25}}%
\TextAt{8}{491.453}{222.374}{.}%
\Vector{q 0 G 0 g 0.399 w 452.334 -551.626 m 452.334 -543.556 l S Q}%
\Vector{q 0 G 0 g 0.399 w 452.534 -543.756 m 459.707 -543.756 l S Q}%
\Vector{q 0 G 0 g 0.399 w 452.534 -551.427 m 459.707 -551.427 l S Q}%
\Vector{q 0 G 0 g 0.399 w 459.906 -551.626 m 459.906 -543.556 l S Q}%
\TextAt{12}{79.226}{188.087}{Proposition}%
\TextAt{12}{153.833}{188.087}{(5.5).}%
\TextAt{8}{190.797}{188.087}{Let}%
\TextAt{13}{213.414}{188.087}{K}%
\TextAt{16}{230.361}{188.087}{\Character{106}}%
\TextAt{18}{239.811}{188.087}{Q}%
\TextAt{21}{249.11}{186.294}{p}%
\TextAt{8}{259.416}{188.087}{be}%
\TextAt{8}{276.988}{188.087}{a}%
\TextAt{13}{288.391}{188.087}{p}%
\TextAt{8}{294.267}{188.087}{-adic}%
\TextAt{8}{324.523}{188.087}{number}%
\TextAt{8}{368.768}{188.087}{field}%
\TextAt{8}{395.772}{188.087}{with}%
\TextAt{8}{424.08}{188.087}{valuation}%
\TextAt{8}{476.801}{188.087}{ring}%
\TextAt{14}{502.515}{188.087}{o}%
\TextAt{8}{513.904}{188.087}{and}%
\TextAt{8}{79.226}{171.041}{maximal}%
\TextAt{8}{127.035}{171.041}{ideal}%
\TextAt{14}{154.998}{171.041}{p}%
\TextAt{8}{160.981}{171.041}{,}%
\TextAt{8}{168.13}{171.041}{and}%
\TextAt{8}{190.893}{171.041}{let}%
\TextAt{13}{207.798}{171.041}{p}%
\TextAt{14}{213.673}{171.041}{o}%
\TextAt{9}{222.838}{171.041}{\Character{61}}%
\TextAt{14}{235.263}{171.041}{p}%
\TextAt{21}{241.246}{175.38}{e}%
\TextAt{8}{245.673}{171.041}{.}%
\TextAt{8}{254.121}{171.041}{Then}%
\TextAt{8}{284.679}{171.041}{the}%
\TextAt{8}{304.835}{171.041}{power}%
\TextAt{8}{338.979}{171.041}{series}%
\TextAt{8}{108.608}{139.413}{exp}%
\TextAt{9}{126.493}{139.413}{\Character{40}}%
\TextAt{13}{131.045}{139.413}{x}%
\TextAt{9}{137.697}{139.413}{\Character{41}}%
\TextAt{9}{145.552}{139.413}{\Character{61}}%
\TextAt{9}{157.973}{139.413}{1}%
\TextAt{9}{166.473}{139.413}{\Character{43}}%
\TextAt{13}{178.267}{139.413}{x}%
\TextAt{9}{187.576}{139.413}{\Character{43}}%
\TextAt{13}{200.533}{147.501}{x}%
\TextAt{17}{207.185}{151.84}{2}%
\Vector{q 0 G 0 g 0.478 w 128.533 -631.598 m 139.917 -631.598 l S Q}%
\TextAt{9}{201.673}{131.213}{2\Character{33}}%
\TextAt{9}{215.769}{139.413}{\Character{43}}%
\TextAt{13}{228.726}{147.501}{x}%
\TextAt{17}{235.378}{151.84}{3}%
\Vector{q 0 G 0 g 0.478 w 156.726 -631.598 m 168.11 -631.598 l S Q}%
\TextAt{9}{229.866}{131.213}{3\Character{33}}%
\TextAt{9}{243.963}{139.413}{\Character{43}}%
\TextAt{16}{255.724}{139.413}{\Character{1}}%
\TextAt{16}{261.02}{139.413}{\Character{1}}%
\TextAt{16}{266.328}{139.413}{\Character{1}}%
\TextAt{8}{295.574}{139.413}{and}%
\TextAt{8}{340.343}{139.413}{log}%
\TextAt{9}{355.31}{139.413}{\Character{40}1}%
\TextAt{9}{368.353}{139.413}{\Character{43}}%
\TextAt{13}{380.134}{139.413}{z}%
\TextAt{9}{386.104}{139.413}{\Character{41}}%
\TextAt{9}{393.959}{139.413}{\Character{61}}%
\TextAt{13}{406.403}{139.413}{z}%
\TextAt{16}{415.03}{139.413}{\Character{0}}%
\TextAt{13}{428.181}{147.501}{z}%
\TextAt{17}{434.152}{151.84}{2}%
\Vector{q 0 G 0 g 0.478 w 356.181 -631.598 m 366.884 -631.598 l S Q}%
\TextAt{9}{430.606}{131.213}{2}%
\TextAt{9}{442.736}{139.413}{\Character{43}}%
\TextAt{13}{455.693}{147.501}{z}%
\TextAt{17}{461.664}{151.84}{3}%
\Vector{q 0 G 0 g 0.478 w 383.693 -631.598 m 394.396 -631.598 l S Q}%
\TextAt{9}{458.118}{131.213}{3}%
\TextAt{16}{470.248}{139.413}{\Character{0}}%
\TextAt{16}{482.191}{139.413}{\Character{1}}%
\TextAt{16}{487.487}{139.413}{\Character{1}}%
\TextAt{16}{492.796}{139.413}{\Character{1}}%
\TextAt{13}{500.136}{139.413}{\Character{59}}%
\TextAt{8}{79.226}{110.506}{yield,}%
\TextAt{8}{110.764}{110.506}{for}%
\TextAt{13}{128.649}{110.506}{n}%
\TextAt{13}{138.942}{110.506}{\Character{62}}%
\TextAt{13}{151.364}{110.506}{e}%
\TextAt{9}{156.808}{110.506}{\Character{47}\Character{40}}%
\TextAt{13}{167.213}{110.506}{p}%
\TextAt{16}{175.745}{110.506}{\Character{0}}%
\TextAt{9}{187.7}{110.506}{1\Character{41}}%
\TextAt{8}{198.106}{110.506}{,}%
\TextAt{8}{205.255}{110.506}{two}%
\TextAt{8}{227.372}{110.506}{mutually}%
\TextAt{8}{276.807}{110.506}{inverse}%
\TextAt{8}{315.565}{110.506}{isomorphisms}%
\TextAt{8}{388.599}{110.506}{(and}%
\TextAt{8}{415.917}{110.506}{homeomorphisms)}%
\TextAt{14}{261.178}{79.397}{p}%
\TextAt{21}{267.161}{83.736}{n}%
\TextAt{13}{326.142}{79.397}{U}%
\TextAt{17}{335.341}{83.736}{\Character{40}}%
\TextAt{21}{338.634}{83.736}{n}%
\TextAt{17}{343.773}{83.736}{\Character{41}}%
\NativeDiagram{%
\path[draw=black,line width=0.478bp] (206.15,-688.729) -- (248.313,-688.729);%
\path[draw=black,line width=0.478bp,line cap=round,line join=round] (246.065,-685.859) .. controls (246.52,-687.581) and (247.532,-688.394) .. (248.552,-688.729) .. controls (247.532,-689.063) and (246.52,-689.876) .. (246.065,-691.598);%
}%
\TextAt{7}{293.003}{89.633}{exp}%
\NativeDiagram{%
\path[draw=black,line width=0.478bp] (248.791,-694.5) -- (206.629,-694.5);%
\path[draw=black,line width=0.478bp,line cap=round,line join=round] (208.876,-697.369) .. controls (208.422,-695.647) and (207.409,-694.834) .. (206.389,-694.5) .. controls (207.409,-694.165) and (208.422,-693.352) .. (208.876,-691.63);%
}%
\TextAt{7}{294.063}{71.154}{log}%
\TextAt{13}{352.717}{79.397}{\Character{58}}%
\TextAt{8}{297.211}{43.033}{125}%
\Destination{78.23}{809.815}{nav.page.125}%
\Destination{79.36}{607.655}{nav.print.137}%
\Destination{79.36}{775.155}{nav.number.2.5.4}%
\Destination{79.36}{202.055}{nav.number.2.5.5}%
\Hotlink{139.28}{624.555}{6.858}{11.047}{nav.section.2.4}%
}

% Page 138 (body page 126)
\Page{%
\TextAt{2}{79.226}{796.943}{II.}%
\TextAt{2}{94.073}{796.943}{The}%
\TextAt{2}{116.492}{796.943}{Theory}%
\TextAt{2}{154.401}{796.943}{of}%
\TextAt{2}{166.826}{796.943}{Valuations}%
\Vector{q 0 G 0 g 0.299 w 7.226 18.729 m 460.769 18.729 l S Q}%
\TextAt{8}{97.159}{761.124}{We}%
\TextAt{8}{117.303}{761.124}{prepare}%
\TextAt{8}{159.577}{761.124}{the}%
\TextAt{8}{179.733}{761.124}{proof}%
\TextAt{8}{210.303}{761.124}{by}%
\TextAt{8}{226.562}{761.124}{the}%
\TextAt{8}{246.718}{761.124}{following}%
\TextAt{8}{296.153}{761.124}{elementary}%
\TextAt{8}{355.988}{761.124}{lemma.}%
\TextAt{12}{79.226}{735.965}{Lemma}%
\TextAt{12}{127.784}{735.965}{(5.6).}%
\TextAt{8}{163.875}{735.965}{Let}%
\TextAt{13}{185.62}{735.965}{\Character{23}}%
\TextAt{9}{196.8}{735.965}{\Character{61}}%
\TextAt{20}{210.547}{744.932}{\Character{80}}%
\TextAt{21}{223.166}{741.777}{r}%
\TextAt{21}{223.166}{732.478}{i}%
\TextAt{17}{226.05}{732.478}{\Character{61}0}%
\TextAt{13}{239.361}{735.965}{a}%
\TextAt{21}{245.506}{734.172}{i}%
\TextAt{13}{248.887}{735.965}{p}%
\TextAt{21}{254.762}{740.304}{i}%
\TextAt{8}{258.144}{735.965}{,}%
\TextAt{9}{266.262}{735.965}{0}%
\NativeMathAt{276.758}{735.965}{11.955}{\leq}%
\TextAt{13}{290.699}{735.965}{a}%
\TextAt{21}{296.844}{734.172}{i}%
\TextAt{13}{304.867}{735.965}{\Character{60}}%
\TextAt{13}{318.603}{735.965}{p}%
\TextAt{8}{324.49}{735.965}{,}%
\TextAt{8}{332.608}{735.965}{be}%
\TextAt{8}{349.308}{735.965}{the}%
\TextAt{13}{370.24}{735.965}{p}%
\TextAt{8}{376.115}{735.965}{-adic}%
\TextAt{8}{405.499}{735.965}{expansion}%
\TextAt{8}{460.647}{735.965}{of}%
\TextAt{8}{474.753}{735.965}{the}%
\TextAt{8}{495.685}{735.965}{natural}%
\TextAt{8}{79.226}{718.919}{number}%
\TextAt{13}{121.822}{718.919}{\Character{23}}%
\NativeMathAt{131.681}{718.919}{11.955}{\in}%
\TextAt{18}{142.972}{718.919}{N}%
\TextAt{8}{151.606}{718.919}{.}%
\TextAt{8}{160.055}{718.919}{Then}%
\TextAt{13}{232.673}{697.884}{v}%
\TextAt{21}{238.333}{696.091}{p}%
\TextAt{9}{243.093}{697.884}{\Character{40}}%
\TextAt{13}{247.646}{697.884}{\Character{23}}%
\TextAt{9}{254.184}{697.884}{\Character{33}\Character{41}}%
\TextAt{9}{265.278}{697.884}{\Character{61}}%
\TextAt{9}{289.173}{705.972}{1}%
\Vector{q 0 G 0 g 0.478 w 206.93 -73.127 m 233.269 -73.127 l S Q}%
\TextAt{13}{278.93}{689.684}{p}%
\TextAt{16}{287.461}{689.684}{\Character{0}}%
\TextAt{9}{299.417}{689.684}{1}%
\TextAt{21}{315.064}{712.829}{r}%
\TextAt{20}{308.458}{709.242}{\Character{88}}%
\TextAt{21}{310.24}{684.048}{i}%
\TextAt{17}{313.123}{684.048}{\Character{61}0}%
\TextAt{13}{327.719}{697.884}{a}%
\TextAt{21}{333.864}{696.091}{i}%
\TextAt{9}{337.245}{697.884}{\Character{40}}%
\TextAt{13}{341.797}{697.884}{p}%
\TextAt{21}{347.672}{702.821}{i}%
\TextAt{16}{353.71}{697.884}{\Character{0}}%
\TextAt{9}{365.665}{697.884}{1\Character{41}}%
\TextAt{13}{376.071}{697.884}{\Character{58}}%
\TextAt{19}{79.226}{651.47}{Proof.}%
\TextAt{8}{116}{651.47}{Let}%
\TextAt{9}{136.97}{651.47}{\Character{91}}%
\TextAt{13}{140.221}{651.47}{c}%
\TextAt{9}{145.259}{651.47}{\Character{93}}%
\TextAt{8}{152.408}{651.47}{signify}%
\TextAt{8}{189.541}{651.47}{the}%
\TextAt{8}{209.697}{651.47}{biggest}%
\TextAt{8}{249.436}{651.47}{integer}%
\NativeMathAt{288.135}{651.47}{11.955}{\leq}%
\TextAt{13}{300.755}{651.47}{c}%
\TextAt{8}{305.793}{651.47}{.}%
\TextAt{8}{314.241}{651.47}{Then}%
\TextAt{8}{344.798}{651.47}{we}%
\TextAt{8}{362.026}{651.47}{have}%
\TextAt{9}{228.09}{626.586}{\Character{91}}%
\TextAt{13}{231.341}{626.586}{\Character{23}}%
\TextAt{9}{237.879}{626.586}{\Character{47}}%
\TextAt{13}{243.732}{626.586}{p}%
\TextAt{9}{249.607}{626.586}{\Character{93}}%
\TextAt{9}{256.158}{626.586}{\Character{61}}%
\TextAt{13}{268.605}{626.586}{a}%
\TextAt{17}{274.75}{624.793}{1}%
\TextAt{9}{282.139}{626.586}{\Character{43}}%
\TextAt{13}{293.9}{626.586}{a}%
\TextAt{17}{300.045}{624.793}{2}%
\TextAt{13}{304.778}{626.586}{p}%
\TextAt{9}{313.31}{626.586}{\Character{43}}%
\TextAt{16}{325.071}{626.586}{\Character{1}}%
\TextAt{16}{330.367}{626.586}{\Character{1}}%
\TextAt{16}{335.675}{626.586}{\Character{1}}%
\TextAt{9}{341.675}{626.586}{\Character{43}}%
\TextAt{13}{353.436}{626.586}{a}%
\TextAt{21}{359.581}{624.793}{r}%
\TextAt{13}{364.136}{626.586}{p}%
\TextAt{21}{370.011}{631.522}{r}%
\TextAt{3}{374.068}{631.522}{\Character{0}}%
\TextAt{17}{380.654}{631.522}{1}%
\TextAt{13}{385.387}{626.586}{\Character{59}}%
\TextAt{9}{223.357}{606.551}{\Character{91}}%
\TextAt{13}{226.609}{606.551}{\Character{23}}%
\TextAt{9}{233.147}{606.551}{\Character{47}}%
\TextAt{13}{239}{606.551}{p}%
\TextAt{17}{244.875}{611.487}{2}%
\TextAt{9}{249.607}{606.551}{\Character{93}}%
\TextAt{9}{256.158}{606.551}{\Character{61}}%
\TextAt{13}{268.605}{606.551}{a}%
\TextAt{17}{274.75}{604.758}{2}%
\TextAt{9}{282.139}{606.551}{\Character{43}}%
\TextAt{16}{293.9}{606.551}{\Character{1}}%
\TextAt{16}{299.196}{606.551}{\Character{1}}%
\TextAt{16}{304.504}{606.551}{\Character{1}}%
\TextAt{9}{310.505}{606.551}{\Character{43}}%
\TextAt{13}{322.266}{606.551}{a}%
\TextAt{21}{328.411}{604.758}{r}%
\TextAt{13}{332.966}{606.551}{p}%
\TextAt{21}{338.841}{611.487}{r}%
\TextAt{3}{342.898}{611.487}{\Character{0}}%
\TextAt{17}{349.484}{611.487}{2}%
\TextAt{13}{354.216}{606.551}{\Character{59}}%
\TextAt{8}{256.756}{590.315}{.}%
\TextAt{8}{256.756}{586.33}{.}%
\TextAt{8}{256.756}{582.345}{.}%
\TextAt{9}{223.535}{562.31}{\Character{91}}%
\TextAt{13}{226.786}{562.31}{\Character{23}}%
\TextAt{9}{233.324}{562.31}{\Character{47}}%
\TextAt{13}{239.177}{562.31}{p}%
\TextAt{21}{245.053}{567.246}{r}%
\TextAt{9}{249.607}{562.31}{\Character{93}}%
\TextAt{9}{256.158}{562.31}{\Character{61}}%
\TextAt{13}{268.605}{562.31}{a}%
\TextAt{21}{274.75}{560.517}{r}%
\TextAt{13}{279.305}{562.31}{\Character{58}}%
\TextAt{8}{79.226}{537.759}{Now}%
\TextAt{8}{106.026}{537.759}{we}%
\TextAt{8}{123.394}{537.759}{count}%
\TextAt{8}{155.729}{537.759}{how}%
\TextAt{8}{180.258}{537.759}{many}%
\TextAt{8}{212.27}{537.759}{numbers}%
\TextAt{9}{259.621}{537.759}{1}%
\TextAt{13}{265.474}{537.759}{\Character{59}}%
\TextAt{9}{270.718}{537.759}{2}%
\TextAt{13}{276.571}{537.759}{\Character{59}}%
\TextAt{13}{281.795}{537.759}{\Character{58}}%
\TextAt{13}{287.032}{537.759}{\Character{58}}%
\TextAt{13}{292.268}{537.759}{\Character{58}}%
\TextAt{13}{297.493}{537.759}{\Character{59}}%
\TextAt{13}{302.729}{537.759}{\Character{23}}%
\TextAt{8}{313.367}{537.759}{are}%
\TextAt{8}{333.018}{537.759}{divisible}%
\TextAt{8}{379.066}{537.759}{by}%
\TextAt{13}{395.465}{537.759}{p}%
\TextAt{8}{401.34}{537.759}{,}%
\TextAt{8}{408.665}{537.759}{and}%
\TextAt{8}{431.567}{537.759}{then}%
\TextAt{8}{458.367}{537.759}{by}%
\TextAt{13}{474.766}{537.759}{p}%
\TextAt{17}{480.642}{542.097}{2}%
\TextAt{8}{485.374}{537.759}{,}%
\TextAt{8}{492.698}{537.759}{etc.}%
\TextAt{8}{516.522}{537.759}{We}%
\TextAt{8}{79.226}{520.713}{find}%
\TextAt{13}{125.207}{503.667}{v}%
\TextAt{21}{130.866}{501.874}{p}%
\TextAt{9}{135.627}{503.667}{\Character{40}}%
\TextAt{13}{140.179}{503.667}{\Character{23}}%
\TextAt{9}{146.717}{503.667}{\Character{33}\Character{41}}%
\TextAt{9}{157.811}{503.667}{\Character{61}}%
\TextAt{9}{170.233}{503.667}{\Character{91}}%
\TextAt{13}{173.519}{503.667}{\Character{23}}%
\TextAt{9}{180.057}{503.667}{\Character{47}}%
\TextAt{13}{185.91}{503.667}{p}%
\TextAt{9}{191.785}{503.667}{\Character{93}}%
\TextAt{9}{197.679}{503.667}{\Character{43}}%
\TextAt{16}{209.455}{503.667}{\Character{1}}%
\TextAt{16}{214.751}{503.667}{\Character{1}}%
\TextAt{16}{220.059}{503.667}{\Character{1}}%
\TextAt{9}{226.059}{503.667}{\Character{43}}%
\TextAt{9}{237.811}{503.667}{\Character{91}}%
\TextAt{13}{241.072}{503.667}{\Character{23}}%
\TextAt{9}{247.61}{503.667}{\Character{47}}%
\TextAt{13}{253.463}{503.667}{p}%
\TextAt{21}{259.338}{508.603}{r}%
\TextAt{9}{263.893}{503.667}{\Character{93}}%
\TextAt{9}{270.444}{503.667}{\Character{61}}%
\TextAt{13}{282.891}{503.667}{a}%
\TextAt{17}{289.036}{501.874}{1}%
\TextAt{9}{296.425}{503.667}{\Character{43}}%
\TextAt{9}{308.177}{503.667}{\Character{40}}%
\TextAt{13}{312.738}{503.667}{p}%
\TextAt{9}{321.27}{503.667}{\Character{43}}%
\TextAt{9}{333.022}{503.667}{1\Character{41}}%
\TextAt{13}{343.437}{503.667}{a}%
\TextAt{17}{349.582}{501.874}{2}%
\TextAt{9}{356.971}{503.667}{\Character{43}}%
\TextAt{16}{368.732}{503.667}{\Character{1}}%
\TextAt{16}{374.028}{503.667}{\Character{1}}%
\TextAt{16}{379.336}{503.667}{\Character{1}}%
\TextAt{9}{385.336}{503.667}{\Character{43}}%
\TextAt{9}{397.088}{503.667}{\Character{40}}%
\TextAt{13}{401.65}{503.667}{p}%
\TextAt{21}{407.525}{508.603}{r}%
\TextAt{3}{411.582}{508.603}{\Character{0}}%
\TextAt{17}{418.168}{508.603}{1}%
\TextAt{9}{425.557}{503.667}{\Character{43}}%
\TextAt{16}{437.318}{503.667}{\Character{1}}%
\TextAt{16}{442.614}{503.667}{\Character{1}}%
\TextAt{16}{447.922}{503.667}{\Character{1}}%
\TextAt{9}{453.923}{503.667}{\Character{43}}%
\TextAt{9}{465.675}{503.667}{1\Character{41}}%
\TextAt{13}{476.089}{503.667}{a}%
\TextAt{21}{482.234}{501.874}{r}%
\TextAt{8}{79.226}{481.426}{and}%
\TextAt{8}{101.989}{481.426}{hence}%
\TextAt{9}{119.801}{451.773}{\Character{40}}%
\TextAt{13}{124.354}{451.773}{p}%
\TextAt{16}{132.885}{451.773}{\Character{0}}%
\TextAt{9}{144.841}{451.773}{1\Character{41}}%
\TextAt{13}{155.246}{451.773}{v}%
\TextAt{21}{160.905}{449.979}{p}%
\TextAt{9}{165.666}{451.773}{\Character{40}}%
\TextAt{13}{170.218}{451.773}{\Character{23}}%
\TextAt{9}{176.756}{451.773}{\Character{33}\Character{41}}%
\TextAt{9}{187.85}{451.773}{\Character{61}}%
\TextAt{9}{200.272}{451.773}{\Character{40}}%
\TextAt{13}{204.859}{451.773}{p}%
\TextAt{16}{213.391}{451.773}{\Character{0}}%
\TextAt{9}{225.346}{451.773}{1\Character{41}}%
\TextAt{13}{235.751}{451.773}{a}%
\TextAt{17}{241.896}{449.979}{1}%
\TextAt{9}{249.285}{451.773}{\Character{43}}%
\TextAt{9}{261.037}{451.773}{\Character{40}}%
\TextAt{13}{265.599}{451.773}{p}%
\TextAt{17}{271.474}{456.709}{2}%
\TextAt{16}{278.863}{451.773}{\Character{0}}%
\TextAt{9}{290.818}{451.773}{1\Character{41}}%
\TextAt{13}{301.223}{451.773}{a}%
\TextAt{17}{307.368}{449.979}{2}%
\TextAt{9}{314.757}{451.773}{\Character{43}}%
\TextAt{16}{326.519}{451.773}{\Character{1}}%
\TextAt{16}{331.815}{451.773}{\Character{1}}%
\TextAt{16}{337.123}{451.773}{\Character{1}}%
\TextAt{9}{343.123}{451.773}{\Character{43}}%
\TextAt{9}{354.875}{451.773}{\Character{40}}%
\TextAt{13}{359.437}{451.773}{p}%
\TextAt{21}{365.312}{456.709}{r}%
\TextAt{16}{372.523}{451.773}{\Character{0}}%
\TextAt{9}{384.478}{451.773}{1\Character{41}}%
\TextAt{13}{394.884}{451.773}{a}%
\TextAt{21}{401.029}{449.979}{r}%
\TextAt{9}{408.904}{451.773}{\Character{61}}%
\TextAt{21}{427.936}{466.717}{r}%
\TextAt{20}{421.33}{463.13}{\Character{88}}%
\TextAt{21}{423.112}{437.936}{i}%
\TextAt{17}{425.995}{437.936}{\Character{61}0}%
\TextAt{13}{440.591}{451.773}{a}%
\TextAt{21}{446.736}{449.979}{i}%
\TextAt{9}{450.117}{451.773}{\Character{40}}%
\TextAt{13}{454.669}{451.773}{p}%
\TextAt{21}{460.544}{456.709}{i}%
\TextAt{16}{466.582}{451.773}{\Character{0}}%
\TextAt{9}{478.538}{451.773}{1\Character{41}}%
\TextAt{13}{488.943}{451.773}{\Character{58}}%
\Vector{q 0 G 0 g 0.399 w 452.334 -355.292 m 452.334 -347.223 l S Q}%
\Vector{q 0 G 0 g 0.399 w 452.534 -347.422 m 459.707 -347.422 l S Q}%
\Vector{q 0 G 0 g 0.399 w 452.534 -355.093 m 459.707 -355.093 l S Q}%
\Vector{q 0 G 0 g 0.399 w 459.906 -355.292 m 459.906 -347.223 l S Q}%
\TextAt{19}{79.226}{390.56}{Proof}%
\TextAt{19}{110.62}{390.56}{of}%
\TextAt{19}{124.369}{390.56}{(5.5).}%
\TextAt{8}{159.039}{390.56}{We}%
\TextAt{8}{179.382}{390.56}{again}%
\TextAt{8}{210.808}{390.56}{think}%
\TextAt{8}{241.9}{390.56}{of}%
\TextAt{8}{255.429}{390.56}{the}%
\TextAt{13}{275.784}{390.56}{p}%
\TextAt{8}{281.659}{390.56}{-adic}%
\TextAt{8}{310.467}{390.56}{valuation}%
\TextAt{13}{361.739}{390.56}{v}%
\TextAt{21}{367.398}{388.767}{p}%
\TextAt{8}{376.255}{390.56}{of}%
\TextAt{18}{389.784}{390.56}{Q}%
\TextAt{21}{399.083}{388.767}{p}%
\TextAt{8}{407.94}{390.56}{as}%
\TextAt{8}{422.509}{390.56}{being}%
\TextAt{8}{454.246}{390.56}{extended}%
\TextAt{8}{504.191}{390.56}{to}%
\TextAt{13}{518.7}{390.56}{K}%
\TextAt{8}{529.518}{390.56}{.}%
\TextAt{8}{79.226}{373.514}{Then}%
\TextAt{13}{110.402}{373.514}{v}%
\TextAt{27}{116.061}{371.721}{p}%
\TextAt{9}{124.922}{373.514}{\Character{61}}%
\TextAt{13}{138.4}{373.514}{ev}%
\TextAt{21}{149.485}{371.721}{p}%
\TextAt{8}{158.761}{373.514}{is}%
\TextAt{8}{171.143}{373.514}{the}%
\TextAt{8}{191.918}{373.514}{normalized}%
\TextAt{8}{252.371}{373.514}{valuation}%
\TextAt{8}{304.062}{373.514}{of}%
\TextAt{13}{318.01}{373.514}{K}%
\TextAt{8}{328.828}{373.514}{.}%
\TextAt{8}{339.13}{373.514}{For}%
\TextAt{8}{360.717}{373.514}{every}%
\TextAt{8}{392.227}{373.514}{natural}%
\TextAt{8}{433.828}{373.514}{number}%
\TextAt{13}{477.042}{373.514}{\Character{23}}%
\TextAt{13}{487.933}{373.514}{\Character{62}}%
\TextAt{9}{501.433}{373.514}{1}%
\TextAt{8}{507.286}{373.514}{,}%
\TextAt{8}{515.207}{373.514}{one}%
\TextAt{8}{79.226}{356.468}{has}%
\TextAt{8}{100.1}{356.468}{the}%
\TextAt{8}{120.256}{356.468}{estimate}%
\TextAt{13}{269.005}{346.386}{v}%
\TextAt{21}{274.664}{344.593}{p}%
\TextAt{9}{279.425}{346.386}{\Character{40}}%
\TextAt{13}{283.978}{346.386}{\Character{23}}%
\TextAt{9}{290.515}{346.386}{\Character{41}}%
\Vector{q 0 G 0 g 0.478 w 196.535 -432.713 m 223.538 -432.713 l S Q}%
\TextAt{13}{268.535}{330.098}{\Character{23}}%
\TextAt{16}{277.73}{330.098}{\Character{0}}%
\TextAt{9}{289.685}{330.098}{1}%
\NativeMathAt{300.054}{338.299}{11.955}{\leq}%
\TextAt{9}{324.112}{346.386}{1}%
\Vector{q 0 G 0 g 0.478 w 241.869 -432.713 m 268.209 -432.713 l S Q}%
\TextAt{13}{313.869}{330.098}{p}%
\TextAt{16}{322.401}{330.098}{\Character{0}}%
\TextAt{9}{334.356}{330.098}{1}%
\TextAt{13}{341.404}{338.299}{\Character{59}}%
\TextAt{8}{79.226}{312.616}{for}%
\TextAt{8}{97.111}{312.616}{if}%
\TextAt{13}{107.835}{312.616}{\Character{23}}%
\TextAt{9}{117.694}{312.616}{\Character{61}}%
\TextAt{13}{130.119}{312.616}{p}%
\TextAt{21}{135.994}{316.954}{a}%
\TextAt{13}{140.99}{312.616}{\Character{23}}%
\TextAt{17}{146.774}{310.822}{0}%
\TextAt{8}{151.507}{312.616}{,}%
\TextAt{8}{158.656}{312.616}{with}%
\TextAt{9}{185.316}{312.616}{\Character{40}}%
\TextAt{13}{189.868}{312.616}{\Character{23}}%
\TextAt{17}{195.652}{310.822}{0}%
\TextAt{13}{200.384}{312.616}{\Character{59}}%
\TextAt{13}{205.608}{312.616}{p}%
\TextAt{9}{211.504}{312.616}{\Character{41}}%
\TextAt{9}{219.359}{312.616}{\Character{61}}%
\TextAt{9}{231.78}{312.616}{1}%
\TextAt{8}{241.553}{312.616}{and}%
\TextAt{13}{264.315}{312.616}{a}%
\TextAt{13}{273.772}{312.616}{\Character{62}}%
\TextAt{9}{286.206}{312.616}{0}%
\TextAt{8}{292.059}{312.616}{,}%
\TextAt{8}{299.209}{312.616}{then}%
\TextAt{13}{145.432}{290.698}{v}%
\TextAt{21}{151.091}{288.905}{p}%
\TextAt{9}{155.852}{290.698}{\Character{40}}%
\TextAt{13}{160.404}{290.698}{\Character{23}}%
\TextAt{9}{166.942}{290.698}{\Character{41}}%
\Vector{q 0 G 0 g 0.478 w 72.962 -488.401 m 99.964 -488.401 l S Q}%
\TextAt{13}{144.962}{274.409}{\Character{23}}%
\TextAt{16}{154.156}{274.409}{\Character{0}}%
\TextAt{9}{166.111}{274.409}{1}%
\TextAt{9}{176.481}{282.61}{\Character{61}}%
\TextAt{13}{207.955}{290.698}{a}%
\Vector{q 0 G 0 g 0.478 w 118.102 -488.401 m 159.954 -488.401 l S Q}%
\TextAt{13}{190.102}{274.409}{p}%
\TextAt{21}{195.977}{277.863}{a}%
\TextAt{13}{200.973}{274.409}{\Character{23}}%
\TextAt{17}{206.757}{272.616}{0}%
\TextAt{16}{214.146}{274.409}{\Character{0}}%
\TextAt{9}{226.101}{274.409}{1}%
\NativeMathAt{236.47}{282.61}{11.955}{\leq}%
\TextAt{13}{262.881}{290.698}{a}%
\Vector{q 0 G 0 g 0.478 w 178.285 -488.401 m 209.621 -488.401 l S Q}%
\TextAt{13}{250.285}{274.409}{p}%
\TextAt{21}{256.16}{277.863}{a}%
\TextAt{16}{263.813}{274.409}{\Character{0}}%
\TextAt{9}{275.768}{274.409}{1}%
\TextAt{9}{286.138}{282.61}{\Character{61}}%
\TextAt{9}{310.002}{290.698}{1}%
\Vector{q 0 G 0 g 0.478 w 227.759 -488.401 m 254.098 -488.401 l S Q}%
\TextAt{13}{299.759}{274.409}{p}%
\TextAt{16}{308.29}{274.409}{\Character{0}}%
\TextAt{9}{320.246}{274.409}{1}%
\TextAt{13}{370.728}{290.698}{a}%
\Vector{q 0 G 0 g 0.478 w 256.49 -488.401 m 347.111 -488.401 l S Q}%
\TextAt{13}{328.49}{274.409}{p}%
\TextAt{21}{334.365}{277.863}{a}%
\TextAt{3}{338.863}{277.863}{\Character{0}}%
\TextAt{17}{345.449}{277.863}{1}%
\TextAt{9}{352.838}{274.409}{\Character{43}}%
\TextAt{16}{364.599}{274.409}{\Character{1}}%
\TextAt{16}{369.895}{274.409}{\Character{1}}%
\TextAt{16}{375.203}{274.409}{\Character{1}}%
\TextAt{9}{381.204}{274.409}{\Character{43}}%
\TextAt{13}{392.965}{274.409}{p}%
\TextAt{9}{401.497}{274.409}{\Character{43}}%
\TextAt{9}{413.249}{274.409}{1}%
\NativeMathAt{423.628}{282.61}{11.955}{\leq}%
\TextAt{9}{447.686}{290.698}{1}%
\Vector{q 0 G 0 g 0.478 w 365.442 -488.401 m 391.782 -488.401 l S Q}%
\TextAt{13}{437.442}{274.409}{p}%
\TextAt{16}{445.974}{274.409}{\Character{0}}%
\TextAt{9}{457.929}{274.409}{1}%
\TextAt{13}{464.978}{282.61}{\Character{58}}%
\TextAt{8}{79.226}{253.156}{For}%
\TextAt{13}{100.196}{253.156}{v}%
\TextAt{21}{105.855}{251.362}{p}%
\TextAt{9}{110.616}{253.156}{\Character{40}}%
\TextAt{13}{115.168}{253.156}{z}%
\TextAt{9}{121.139}{253.156}{\Character{41}}%
\TextAt{13}{129.012}{253.156}{\Character{62}}%
\TextAt{9}{141.437}{253.156}{1\Character{47}\Character{40}}%
\TextAt{13}{157.696}{253.156}{p}%
\TextAt{16}{166.227}{253.156}{\Character{0}}%
\TextAt{9}{178.183}{253.156}{1\Character{41}}%
\TextAt{8}{188.588}{253.156}{,}%
\TextAt{13}{195.737}{253.156}{z}%
\NativeMathAt{205.028}{253.156}{11.955}{\neq}%
\TextAt{9}{217.439}{253.156}{0}%
\TextAt{8}{223.307}{253.156}{,}%
\TextAt{8}{230.456}{253.156}{i.e.,}%
\TextAt{13}{252.561}{253.156}{v}%
\TextAt{27}{258.221}{251.362}{p}%
\TextAt{9}{262.707}{253.156}{\Character{40}}%
\TextAt{13}{267.26}{253.156}{z}%
\TextAt{9}{273.23}{253.156}{\Character{41}}%
\TextAt{13}{281.103}{253.156}{\Character{62}}%
\TextAt{13}{293.512}{253.156}{e}%
\TextAt{9}{298.954}{253.156}{\Character{47}\Character{40}}%
\TextAt{13}{309.36}{253.156}{p}%
\TextAt{16}{317.892}{253.156}{\Character{0}}%
\TextAt{9}{329.847}{253.156}{1\Character{41}}%
\TextAt{8}{340.252}{253.156}{,}%
\TextAt{8}{347.401}{253.156}{this}%
\TextAt{8}{370.224}{253.156}{yields}%
\TextAt{13}{120.54}{222.515}{v}%
\TextAt{21}{126.199}{220.721}{p}%
\TextAt{20}{132.952}{239.372}{\Character{18}}%
\TextAt{13}{142.948}{230.602}{z}%
\TextAt{21}{148.919}{234.941}{\Character{23}}%
\Vector{q 0 G 0 g 0.478 w 70.948 -548.497 m 82.101 -548.497 l S Q}%
\TextAt{13}{145.256}{214.314}{\Character{23}}%
\TextAt{20}{155.296}{239.372}{\Character{19}}%
\TextAt{16}{166.753}{222.515}{\Character{0}}%
\TextAt{13}{178.709}{222.515}{v}%
\TextAt{21}{184.368}{220.721}{p}%
\TextAt{9}{189.129}{222.515}{\Character{40}}%
\TextAt{13}{193.681}{222.515}{z}%
\TextAt{9}{199.652}{222.515}{\Character{41}}%
\TextAt{9}{207.507}{222.515}{\Character{61}}%
\TextAt{9}{219.928}{222.515}{\Character{40}}%
\TextAt{13}{224.503}{222.515}{\Character{23}}%
\TextAt{16}{233.697}{222.515}{\Character{0}}%
\TextAt{9}{245.652}{222.515}{1\Character{41}}%
\TextAt{13}{256.058}{222.515}{v}%
\TextAt{21}{261.717}{220.721}{p}%
\TextAt{9}{266.478}{222.515}{\Character{40}}%
\TextAt{13}{271.03}{222.515}{z}%
\TextAt{9}{277.001}{222.515}{\Character{41}}%
\TextAt{16}{284.21}{222.515}{\Character{0}}%
\TextAt{13}{296.165}{222.515}{v}%
\TextAt{21}{301.824}{220.721}{p}%
\TextAt{9}{306.585}{222.515}{\Character{40}}%
\TextAt{13}{311.137}{222.515}{\Character{23}}%
\TextAt{9}{317.675}{222.515}{\Character{41}}%
\TextAt{13}{325.549}{222.515}{\Character{62}}%
\TextAt{9}{337.974}{222.515}{\Character{40}}%
\TextAt{13}{342.526}{222.515}{\Character{23}}%
\TextAt{16}{351.721}{222.515}{\Character{0}}%
\TextAt{9}{363.676}{222.515}{1\Character{41}}%
\TextAt{20}{376.074}{239.372}{\Character{18}}%
\TextAt{9}{396.313}{230.602}{1}%
\Vector{q 0 G 0 g 0.478 w 314.07 -548.497 m 340.41 -548.497 l S Q}%
\TextAt{13}{386.07}{214.314}{p}%
\TextAt{16}{394.602}{214.314}{\Character{0}}%
\TextAt{9}{406.557}{214.314}{1}%
\TextAt{16}{416.262}{222.515}{\Character{0}}%
\TextAt{13}{429.883}{230.602}{v}%
\TextAt{21}{435.542}{228.809}{p}%
\TextAt{9}{440.303}{230.602}{\Character{40}}%
\TextAt{13}{444.855}{230.602}{\Character{23}}%
\TextAt{9}{451.393}{230.602}{\Character{41}}%
\Vector{q 0 G 0 g 0.478 w 357.413 -548.497 m 384.415 -548.497 l S Q}%
\TextAt{13}{429.413}{214.314}{\Character{23}}%
\TextAt{16}{438.607}{214.314}{\Character{0}}%
\TextAt{9}{450.562}{214.314}{1}%
\TextAt{20}{457.611}{239.372}{\Character{19}}%
\NativeMathAt{469.732}{222.515}{11.955}{\geq}%
\TextAt{9}{482.351}{222.515}{0}%
\TextAt{13}{488.204}{222.515}{\Character{59}}%
\TextAt{8}{79.226}{190.878}{and}%
\TextAt{8}{101.989}{190.878}{thus}%
\TextAt{13}{127.74}{190.878}{v}%
\TextAt{27}{133.4}{189.085}{p}%
\TextAt{9}{137.886}{190.878}{\Character{40}}%
\TextAt{8}{142.439}{190.878}{log}%
\TextAt{9}{157.407}{190.878}{\Character{40}1}%
\TextAt{9}{170.45}{190.878}{\Character{43}}%
\TextAt{13}{182.23}{190.878}{z}%
\TextAt{9}{188.2}{190.878}{\Character{41}\Character{41}}%
\TextAt{9}{200.598}{190.878}{\Character{61}}%
\TextAt{13}{213.051}{190.878}{v}%
\TextAt{27}{218.711}{189.085}{p}%
\TextAt{9}{223.198}{190.878}{\Character{40}}%
\TextAt{13}{227.75}{190.878}{z}%
\TextAt{9}{233.72}{190.878}{\Character{41}}%
\TextAt{8}{238.273}{190.878}{.}%
\TextAt{8}{246.721}{190.878}{For}%
\TextAt{13}{267.691}{190.878}{n}%
\TextAt{13}{277.984}{190.878}{\Character{62}}%
\TextAt{13}{290.406}{190.878}{e}%
\TextAt{9}{295.85}{190.878}{\Character{47}\Character{40}}%
\TextAt{13}{306.255}{190.878}{p}%
\TextAt{16}{314.787}{190.878}{\Character{0}}%
\TextAt{9}{326.742}{190.878}{1\Character{41}}%
\TextAt{8}{337.147}{190.878}{,}%
\TextAt{8}{344.297}{190.878}{log}%
\TextAt{8}{363.162}{190.878}{therefore}%
\TextAt{8}{412.262}{190.878}{maps}%
\TextAt{13}{442.891}{190.878}{U}%
\TextAt{17}{452.09}{195.217}{\Character{40}}%
\TextAt{21}{455.383}{195.217}{n}%
\TextAt{17}{460.522}{195.217}{\Character{41}}%
\TextAt{8}{468.21}{190.878}{into}%
\TextAt{14}{491.953}{190.878}{p}%
\TextAt{21}{497.936}{195.217}{n}%
\TextAt{8}{503.573}{190.878}{.}%
\TextAt{8}{97.159}{170.844}{For}%
\TextAt{8}{117.345}{170.844}{the}%
\TextAt{8}{136.718}{170.844}{exponential}%
\TextAt{8}{198.7}{170.844}{series}%
\TextAt{20}{229.251}{179.81}{\Character{80}}%
\NativeMathAt{241.87}{176.655}{7.97}{\infty}%
\TextAt{21}{241.87}{167.357}{\Character{23}}%
\TextAt{17}{246.554}{167.357}{\Character{61}0}%
\TextAt{13}{259.865}{170.844}{x}%
\TextAt{21}{266.517}{175.182}{\Character{23}}%
\TextAt{9}{271.7}{170.844}{\Character{47}}%
\TextAt{13}{277.553}{170.844}{\Character{23}}%
\TextAt{9}{284.09}{170.844}{\Character{33}}%
\TextAt{8}{287.342}{170.844}{,}%
\TextAt{8}{293.865}{170.844}{we}%
\TextAt{8}{310.309}{170.844}{compute}%
\TextAt{8}{357}{170.844}{the}%
\TextAt{8}{376.373}{170.844}{valuations}%
\TextAt{13}{431.277}{170.844}{v}%
\TextAt{21}{436.936}{169.05}{p}%
\TextAt{9}{441.697}{170.844}{\Character{40}}%
\TextAt{13}{446.249}{170.844}{x}%
\TextAt{21}{452.901}{175.182}{\Character{23}}%
\TextAt{9}{458.083}{170.844}{\Character{47}}%
\TextAt{13}{463.936}{170.844}{\Character{23}}%
\TextAt{9}{470.474}{170.844}{\Character{33}\Character{41}}%
\TextAt{8}{481.392}{170.844}{as}%
\TextAt{8}{494.979}{170.844}{follows.}%
\TextAt{8}{79.226}{153.798}{Writing,}%
\TextAt{8}{125.397}{153.798}{for}%
\TextAt{13}{143.282}{153.798}{\Character{23}}%
\TextAt{13}{153.121}{153.798}{\Character{62}}%
\TextAt{9}{165.566}{153.798}{0}%
\TextAt{8}{171.419}{153.798}{,}%
\TextAt{13}{199.202}{136.752}{\Character{23}}%
\TextAt{9}{209.061}{136.752}{\Character{61}}%
\TextAt{13}{221.486}{136.752}{a}%
\TextAt{17}{227.631}{134.958}{0}%
\TextAt{9}{235.02}{136.752}{\Character{43}}%
\TextAt{13}{246.781}{136.752}{a}%
\TextAt{17}{252.926}{134.958}{1}%
\TextAt{13}{257.659}{136.752}{p}%
\TextAt{9}{266.19}{136.752}{\Character{43}}%
\TextAt{16}{277.952}{136.752}{\Character{1}}%
\TextAt{16}{283.248}{136.752}{\Character{1}}%
\TextAt{16}{288.556}{136.752}{\Character{1}}%
\TextAt{9}{294.556}{136.752}{\Character{43}}%
\TextAt{13}{306.317}{136.752}{a}%
\TextAt{21}{312.462}{134.958}{r}%
\TextAt{13}{317.017}{136.752}{p}%
\TextAt{21}{322.892}{141.688}{r}%
\TextAt{13}{327.447}{136.752}{\Character{59}}%
\TextAt{9}{356.601}{136.752}{0}%
\NativeMathAt{365.775}{136.752}{11.955}{\leq}%
\TextAt{13}{378.395}{136.752}{a}%
\TextAt{21}{384.539}{134.958}{i}%
\TextAt{13}{391.242}{136.752}{\Character{60}}%
\TextAt{13}{403.651}{136.752}{p\Character{59}}%
\TextAt{8}{79.226}{114.511}{we}%
\TextAt{8}{96.454}{114.511}{get}%
\TextAt{8}{115.964}{114.511}{from}%
\TextAt{8}{143.605}{114.511}{(5.6)}%
\TextAt{8}{171.58}{114.511}{that}%
\TextAt{13}{145}{81.976}{v}%
\TextAt{21}{150.659}{80.183}{p}%
\TextAt{9}{155.42}{81.976}{\Character{40}}%
\TextAt{13}{159.973}{81.976}{\Character{23}}%
\TextAt{9}{166.511}{81.976}{\Character{33}\Character{41}}%
\TextAt{9}{177.605}{81.976}{\Character{61}}%
\TextAt{9}{201.5}{90.064}{1}%
\Vector{q 0 G 0 g 0.478 w 119.256 -689.035 m 145.596 -689.035 l S Q}%
\TextAt{13}{191.256}{73.776}{p}%
\TextAt{16}{199.788}{73.776}{\Character{0}}%
\TextAt{9}{211.743}{73.776}{1}%
\TextAt{21}{227.39}{96.92}{r}%
\TextAt{20}{220.784}{93.334}{\Character{88}}%
\TextAt{21}{222.567}{68.139}{i}%
\TextAt{17}{225.45}{68.139}{\Character{61}0}%
\TextAt{13}{240.045}{81.976}{a}%
\TextAt{21}{246.19}{80.183}{i}%
\TextAt{9}{249.572}{81.976}{\Character{40}}%
\TextAt{13}{254.124}{81.976}{p}%
\TextAt{21}{259.999}{86.912}{i}%
\TextAt{16}{266.037}{81.976}{\Character{0}}%
\TextAt{9}{277.992}{81.976}{1\Character{41}}%
\TextAt{9}{291.693}{81.976}{\Character{61}}%
\TextAt{9}{315.583}{90.064}{1}%
\Vector{q 0 G 0 g 0.478 w 233.339 -689.035 m 259.679 -689.035 l S Q}%
\TextAt{13}{305.339}{73.776}{p}%
\TextAt{16}{313.871}{73.776}{\Character{0}}%
\TextAt{9}{325.826}{73.776}{1}%
\TextAt{20}{332.875}{91.66}{\Character{0}}%
\TextAt{13}{338.354}{81.976}{\Character{23}}%
\TextAt{16}{347.549}{81.976}{\Character{0}}%
\TextAt{9}{359.504}{81.976}{\Character{40}}%
\TextAt{13}{364.056}{81.976}{a}%
\TextAt{17}{370.201}{80.183}{0}%
\TextAt{9}{377.59}{81.976}{\Character{43}}%
\TextAt{13}{389.352}{81.976}{a}%
\TextAt{17}{395.497}{80.183}{1}%
\TextAt{9}{402.886}{81.976}{\Character{43}}%
\TextAt{16}{414.647}{81.976}{\Character{1}}%
\TextAt{16}{419.943}{81.976}{\Character{1}}%
\TextAt{16}{425.251}{81.976}{\Character{1}}%
\TextAt{9}{431.251}{81.976}{\Character{43}}%
\TextAt{13}{443.012}{81.976}{a}%
\TextAt{21}{449.157}{80.183}{r}%
\TextAt{9}{453.712}{81.976}{\Character{41}}%
\TextAt{20}{458.264}{91.66}{\Character{1}}%
\TextAt{13}{463.744}{81.976}{\Character{58}}%
\TextAt{8}{297.211}{43.033}{126}%
\Destination{78.23}{809.815}{nav.page.126}%
\Destination{79.36}{750.955}{nav.print.138}%
\Destination{79.36}{749.955}{nav.number.2.5.6}%
\Hotlink{123.869}{388.355}{26.106}{11.047}{nav.number.2.5.5}%
\Hotlink{143.105}{112.355}{25.078}{11.047}{nav.number.2.5.6}%
}

% Page 139 (body page 127)
\Page{%
\TextAt{2}{79.226}{796.943}{II.}%
\TextAt{2}{94.073}{796.943}{The}%
\TextAt{2}{116.492}{796.943}{Theory}%
\TextAt{2}{154.401}{796.943}{of}%
\TextAt{2}{166.826}{796.943}{Valuations}%
\Vector{q 0 G 0 g 0.299 w 7.226 18.729 m 460.769 18.729 l S Q}%
\TextAt{8}{79.226}{761.124}{Putting}%
\TextAt{13}{122.313}{761.124}{s}%
\TextAt{21}{127.827}{759.33}{\Character{23}}%
\TextAt{9}{136.329}{761.124}{\Character{61}}%
\TextAt{13}{148.755}{761.124}{a}%
\TextAt{17}{154.9}{759.33}{0}%
\TextAt{9}{162.289}{761.124}{\Character{43}}%
\TextAt{16}{174.05}{761.124}{\Character{1}}%
\TextAt{16}{179.346}{761.124}{\Character{1}}%
\TextAt{16}{184.654}{761.124}{\Character{1}}%
\TextAt{9}{190.654}{761.124}{\Character{43}}%
\TextAt{13}{202.416}{761.124}{a}%
\TextAt{21}{208.561}{759.33}{r}%
\TextAt{8}{217.013}{761.124}{this}%
\TextAt{8}{239.835}{761.124}{becomes}%
\TextAt{13}{156.325}{730.212}{v}%
\TextAt{21}{161.984}{728.418}{p}%
\TextAt{20}{168.737}{747.069}{\Character{18}}%
\TextAt{13}{178.733}{738.299}{x}%
\TextAt{21}{185.385}{742.638}{\Character{23}}%
\Vector{q 0 G 0 g 0.478 w 106.733 -40.8 m 118.567 -40.8 l S Q}%
\TextAt{13}{179.756}{722.011}{\Character{23}}%
\TextAt{9}{186.294}{722.011}{\Character{33}}%
\TextAt{20}{191.763}{747.069}{\Character{19}}%
\TextAt{9}{203.884}{730.212}{\Character{61}}%
\TextAt{13}{216.31}{730.212}{\Character{23}}%
\TextAt{13}{222.838}{730.212}{v}%
\TextAt{21}{228.507}{728.418}{p}%
\TextAt{9}{233.268}{730.212}{\Character{40}}%
\TextAt{13}{237.82}{730.212}{x}%
\TextAt{9}{244.472}{730.212}{\Character{41}}%
\TextAt{16}{251.681}{730.212}{\Character{0}}%
\TextAt{13}{264.832}{738.299}{\Character{23}}%
\TextAt{16}{274.026}{738.299}{\Character{0}}%
\TextAt{13}{285.982}{738.299}{s}%
\TextAt{21}{291.496}{736.506}{\Character{23}}%
\Vector{q 0 G 0 g 0.478 w 192.832 -40.8 m 224.678 -40.8 l S Q}%
\TextAt{13}{267.585}{722.011}{p}%
\TextAt{16}{276.117}{722.011}{\Character{0}}%
\TextAt{9}{288.072}{722.011}{1}%
\TextAt{9}{301.194}{730.212}{\Character{61}}%
\TextAt{13}{313.619}{730.212}{\Character{23}}%
\TextAt{20}{322.15}{747.069}{\Character{18}}%
\TextAt{13}{330.95}{730.212}{v}%
\TextAt{21}{336.61}{728.418}{p}%
\TextAt{9}{341.37}{730.212}{\Character{40}}%
\TextAt{13}{345.923}{730.212}{x}%
\TextAt{9}{352.575}{730.212}{\Character{41}}%
\TextAt{16}{359.784}{730.212}{\Character{0}}%
\TextAt{9}{383.178}{738.299}{1}%
\Vector{q 0 G 0 g 0.478 w 300.935 -40.8 m 327.274 -40.8 l S Q}%
\TextAt{13}{372.935}{722.011}{p}%
\TextAt{16}{381.466}{722.011}{\Character{0}}%
\TextAt{9}{393.421}{722.011}{1}%
\TextAt{20}{400.47}{747.069}{\Character{19}}%
\TextAt{9}{411.927}{730.212}{\Character{43}}%
\TextAt{13}{432.706}{738.299}{s}%
\TextAt{21}{438.22}{736.506}{\Character{23}}%
\Vector{q 0 G 0 g 0.478 w 352.884 -40.8 m 379.224 -40.8 l S Q}%
\TextAt{13}{424.884}{722.011}{p}%
\TextAt{16}{433.416}{722.011}{\Character{0}}%
\TextAt{9}{445.371}{722.011}{1}%
\TextAt{13}{452.419}{730.212}{\Character{58}}%
\TextAt{8}{79.226}{698.761}{For}%
\TextAt{13}{98.998}{698.761}{v}%
\TextAt{27}{104.657}{696.968}{p}%
\TextAt{9}{109.144}{698.761}{\Character{40}}%
\TextAt{13}{113.697}{698.761}{x}%
\TextAt{9}{120.349}{698.761}{\Character{41}}%
\TextAt{13}{128.222}{698.761}{\Character{62}}%
\TextAt{13}{140.632}{698.761}{e}%
\TextAt{9}{146.073}{698.761}{\Character{47}\Character{40}}%
\TextAt{13}{156.478}{698.761}{p}%
\TextAt{16}{162.562}{698.761}{\Character{0}}%
\TextAt{9}{172.068}{698.761}{1\Character{41}}%
\TextAt{8}{182.474}{698.761}{,}%
\TextAt{8}{188.665}{698.761}{i.e.,}%
\TextAt{13}{209.812}{698.761}{v}%
\TextAt{21}{215.471}{696.968}{p}%
\TextAt{9}{220.232}{698.761}{\Character{40}}%
\TextAt{13}{224.785}{698.761}{x}%
\TextAt{9}{231.437}{698.761}{\Character{41}}%
\TextAt{13}{239.31}{698.761}{\Character{62}}%
\TextAt{9}{251.735}{698.761}{1\Character{47}\Character{40}}%
\TextAt{13}{267.994}{698.761}{p}%
\TextAt{16}{274.077}{698.761}{\Character{0}}%
\TextAt{9}{283.584}{698.761}{1\Character{41}}%
\TextAt{8}{293.989}{698.761}{,}%
\TextAt{8}{300.18}{698.761}{this}%
\TextAt{8}{321.806}{698.761}{implies}%
\TextAt{8}{360.335}{698.761}{the}%
\TextAt{8}{379.294}{698.761}{convergence}%
\TextAt{8}{442.81}{698.761}{of}%
\TextAt{8}{454.943}{698.761}{the}%
\TextAt{8}{473.902}{698.761}{exponential}%
\TextAt{8}{79.226}{681.716}{series.}%
\TextAt{8}{115.112}{681.716}{If}%
\TextAt{8}{126.804}{681.716}{furthermore}%
\TextAt{13}{191.517}{681.716}{x}%
\NativeMathAt{201.49}{681.715}{11.955}{\neq}%
\TextAt{9}{213.9}{681.716}{0}%
\TextAt{8}{223.666}{681.716}{and}%
\TextAt{13}{246.428}{681.716}{\Character{23}}%
\TextAt{13}{256.267}{681.716}{\Character{62}}%
\TextAt{9}{268.713}{681.715}{1}%
\TextAt{8}{274.566}{681.716}{,}%
\TextAt{8}{281.715}{681.716}{then}%
\TextAt{8}{308.375}{681.716}{one}%
\TextAt{8}{329.835}{681.716}{has}%
\TextAt{13}{141.31}{650.93}{v}%
\TextAt{21}{146.969}{649.136}{p}%
\TextAt{20}{153.722}{667.787}{\Character{18}}%
\TextAt{13}{163.718}{659.017}{x}%
\TextAt{21}{170.37}{663.356}{\Character{23}}%
\Vector{q 0 G 0 g 0.478 w 91.718 -120.082 m 103.552 -120.082 l S Q}%
\TextAt{13}{164.741}{642.729}{\Character{23}}%
\TextAt{9}{171.278}{642.729}{\Character{33}}%
\TextAt{20}{176.748}{667.787}{\Character{19}}%
\TextAt{16}{188.205}{650.93}{\Character{0}}%
\TextAt{13}{200.16}{650.93}{v}%
\TextAt{21}{205.819}{649.136}{p}%
\TextAt{9}{210.58}{650.93}{\Character{40}}%
\TextAt{13}{215.133}{650.93}{x}%
\TextAt{9}{221.785}{650.93}{\Character{41}}%
\TextAt{9}{229.64}{650.93}{\Character{61}}%
\TextAt{9}{242.061}{650.93}{\Character{40}}%
\TextAt{13}{246.636}{650.93}{\Character{23}}%
\TextAt{16}{255.83}{650.93}{\Character{0}}%
\TextAt{9}{267.785}{650.93}{1\Character{41}}%
\TextAt{13}{278.191}{650.93}{v}%
\TextAt{21}{283.85}{649.136}{p}%
\TextAt{9}{288.611}{650.93}{\Character{40}}%
\TextAt{13}{293.163}{650.93}{x}%
\TextAt{9}{299.815}{650.93}{\Character{41}}%
\TextAt{16}{307.024}{650.93}{\Character{0}}%
\TextAt{13}{320.175}{659.017}{\Character{23}}%
\TextAt{16}{329.37}{659.017}{\Character{0}}%
\TextAt{9}{341.325}{659.017}{1}%
\Vector{q 0 G 0 g 0.478 w 248.175 -120.082 m 275.178 -120.082 l S Q}%
\TextAt{13}{320.506}{642.729}{p}%
\TextAt{16}{329.038}{642.729}{\Character{0}}%
\TextAt{9}{340.993}{642.729}{1}%
\TextAt{9}{351.03}{650.93}{\Character{43}}%
\TextAt{13}{363.987}{659.017}{s}%
\TextAt{21}{369.501}{657.224}{\Character{23}}%
\TextAt{16}{377.339}{659.017}{\Character{0}}%
\TextAt{9}{389.295}{659.017}{1}%
\Vector{q 0 G 0 g 0.478 w 291.987 -120.082 m 323.148 -120.082 l S Q}%
\TextAt{13}{366.397}{642.729}{p}%
\TextAt{16}{374.929}{642.729}{\Character{0}}%
\TextAt{9}{386.884}{642.729}{1}%
\TextAt{13}{399.664}{650.93}{\Character{62}}%
\TextAt{13}{413.285}{659.017}{s}%
\TextAt{21}{418.799}{657.224}{\Character{23}}%
\TextAt{16}{426.638}{659.017}{\Character{0}}%
\TextAt{9}{438.593}{659.017}{1}%
\Vector{q 0 G 0 g 0.478 w 341.285 -120.082 m 372.446 -120.082 l S Q}%
\TextAt{13}{415.695}{642.729}{p}%
\TextAt{16}{424.227}{642.729}{\Character{0}}%
\TextAt{9}{436.182}{642.729}{1}%
\NativeMathAt{448.962}{650.93}{11.955}{\geq}%
\TextAt{9}{461.581}{650.93}{0}%
\TextAt{13}{467.434}{650.93}{\Character{58}}%
\TextAt{8}{79.226}{618.13}{Therefore}%
\TextAt{13}{131.799}{618.13}{v}%
\TextAt{27}{137.458}{616.337}{p}%
\TextAt{9}{141.945}{618.13}{\Character{40}}%
\TextAt{8}{146.498}{618.13}{exp}%
\TextAt{9}{164.383}{618.13}{\Character{40}}%
\TextAt{13}{168.935}{618.13}{x}%
\TextAt{9}{175.587}{618.13}{\Character{41}}%
\TextAt{16}{181.928}{618.13}{\Character{0}}%
\TextAt{9}{193.016}{618.13}{1\Character{41}}%
\TextAt{9}{206.717}{618.13}{\Character{61}}%
\TextAt{13}{219.168}{618.13}{v}%
\TextAt{27}{224.827}{616.337}{p}%
\TextAt{9}{229.314}{618.13}{\Character{40}}%
\TextAt{13}{233.866}{618.13}{x}%
\TextAt{9}{240.518}{618.13}{\Character{41}}%
\TextAt{8}{245.07}{618.13}{,}%
\TextAt{8}{251.88}{618.13}{i.e.,}%
\TextAt{8}{273.646}{618.13}{for}%
\TextAt{13}{291.107}{618.13}{n}%
\TextAt{13}{301.4}{618.13}{\Character{62}}%
\TextAt{13}{313.822}{618.13}{e}%
\TextAt{9}{319.266}{618.13}{\Character{47}\Character{40}}%
\TextAt{13}{329.671}{618.13}{p}%
\TextAt{16}{337.335}{618.13}{\Character{0}}%
\TextAt{9}{348.423}{618.13}{1\Character{41}}%
\TextAt{8}{358.828}{618.13}{,}%
\TextAt{8}{365.638}{618.13}{exp}%
\TextAt{8}{386.996}{618.13}{maps}%
\TextAt{8}{417.201}{618.13}{the}%
\TextAt{8}{436.933}{618.13}{group}%
\TextAt{14}{469.685}{618.13}{p}%
\TextAt{21}{475.668}{622.468}{n}%
\TextAt{8}{484.777}{618.13}{into}%
\TextAt{13}{508.096}{618.13}{U}%
\TextAt{17}{517.295}{622.468}{\Character{40}}%
\TextAt{21}{520.588}{622.468}{n}%
\TextAt{17}{525.726}{622.468}{\Character{41}}%
\TextAt{8}{529.518}{618.13}{.}%
\TextAt{8}{79.226}{601.084}{Furthermore,}%
\TextAt{8}{150.276}{601.084}{one}%
\TextAt{8}{171.735}{601.084}{has}%
\TextAt{8}{192.609}{601.084}{for}%
\TextAt{13}{210.494}{601.084}{v}%
\TextAt{27}{216.153}{599.291}{p}%
\TextAt{9}{220.64}{601.084}{\Character{40}}%
\TextAt{13}{225.192}{601.084}{x}%
\TextAt{9}{231.844}{601.084}{\Character{41}}%
\TextAt{13}{236.397}{601.084}{\Character{59}}%
\TextAt{13}{241.621}{601.084}{v}%
\TextAt{27}{247.3}{599.291}{p}%
\TextAt{9}{251.787}{601.084}{\Character{40}}%
\TextAt{13}{256.339}{601.084}{z}%
\TextAt{9}{262.31}{601.084}{\Character{41}}%
\TextAt{13}{270.183}{601.084}{\Character{62}}%
\TextAt{13}{282.593}{601.084}{e}%
\TextAt{9}{288.034}{601.084}{\Character{47}\Character{40}}%
\TextAt{13}{298.439}{601.084}{p}%
\TextAt{16}{306.971}{601.084}{\Character{0}}%
\TextAt{9}{318.926}{601.084}{1\Character{41}}%
\TextAt{8}{333.229}{601.084}{that}%
\TextAt{8}{180.999}{573.908}{exp}%
\TextAt{8}{200.869}{573.908}{log}%
\TextAt{9}{215.844}{573.908}{\Character{40}1}%
\TextAt{9}{228.887}{573.908}{\Character{43}}%
\TextAt{13}{240.668}{573.908}{z}%
\TextAt{9}{246.638}{573.908}{\Character{41}}%
\TextAt{9}{254.493}{573.908}{\Character{61}}%
\TextAt{9}{266.914}{573.908}{1}%
\TextAt{9}{275.414}{573.908}{\Character{43}}%
\TextAt{13}{287.208}{573.908}{z}%
\TextAt{8}{317.089}{573.908}{and}%
\TextAt{8}{361.857}{573.908}{log}%
\TextAt{8}{378.81}{573.908}{exp}%
\TextAt{13}{398.695}{573.908}{x}%
\TextAt{9}{408.667}{573.908}{\Character{61}}%
\TextAt{13}{421.093}{573.908}{x\Character{59}}%
\TextAt{8}{79.226}{546.732}{for}%
\TextAt{8}{97.403}{546.732}{these}%
\TextAt{8}{127.667}{546.732}{are}%
\TextAt{8}{147.47}{546.732}{identities}%
\TextAt{8}{198.225}{546.732}{of}%
\TextAt{8}{211.847}{546.732}{formal}%
\TextAt{8}{248.889}{546.732}{power}%
\TextAt{8}{283.325}{546.732}{series}%
\TextAt{8}{314.952}{546.732}{and}%
\TextAt{8}{338.006}{546.732}{all}%
\TextAt{8}{354.557}{546.732}{of}%
\TextAt{8}{368.18}{546.732}{the}%
\TextAt{8}{388.628}{546.732}{series}%
\TextAt{8}{420.255}{546.732}{converge.}%
\TextAt{8}{473.491}{546.732}{This}%
\TextAt{8}{500.502}{546.732}{proves}%
\TextAt{8}{79.226}{529.686}{the}%
\TextAt{8}{99.383}{529.686}{proposition.}%
\Vector{q 0 G 0 g 0.399 w 452.334 -244.314 m 452.334 -236.244 l S Q}%
\Vector{q 0 G 0 g 0.399 w 452.534 -236.443 m 459.707 -236.443 l S Q}%
\Vector{q 0 G 0 g 0.399 w 452.534 -244.115 m 459.707 -244.115 l S Q}%
\Vector{q 0 G 0 g 0.399 w 459.906 -244.314 m 459.906 -236.244 l S Q}%
\TextAt{8}{97.159}{501.207}{For}%
\TextAt{8}{117.761}{501.207}{an}%
\TextAt{8}{133.652}{501.207}{arbitrary}%
\TextAt{8}{183.054}{501.207}{local}%
\TextAt{8}{210.327}{501.207}{field}%
\TextAt{13}{235.316}{501.207}{K}%
\TextAt{8}{246.134}{501.207}{,}%
\TextAt{8}{252.989}{501.207}{the}%
\TextAt{8}{272.778}{501.207}{group}%
\TextAt{8}{305.586}{501.207}{of}%
\TextAt{8}{318.549}{501.207}{principal}%
\TextAt{8}{366.958}{501.207}{units}%
\TextAt{13}{395.917}{501.207}{U}%
\TextAt{17}{405.116}{505.545}{\Character{40}1\Character{41}}%
\TextAt{8}{419.964}{501.207}{is}%
\TextAt{8}{431.361}{501.207}{a}%
\TextAt{18}{440.749}{501.207}{Z}%
\TextAt{21}{448.719}{499.413}{p}%
\TextAt{8}{453.48}{501.207}{-module}%
\TextAt{8}{498.303}{501.207}{(where}%
\TextAt{13}{79.226}{484.161}{p}%
\TextAt{9}{88.422}{484.161}{\Character{61}}%
\TextAt{8}{100.848}{484.161}{char}%
\TextAt{9}{122.965}{484.161}{\Character{40}}%
\TextAt{13}{127.517}{484.161}{\Character{20}}%
\TextAt{9}{134.256}{484.161}{\Character{41}}%
\TextAt{8}{138.808}{484.161}{)}%
\TextAt{8}{146.984}{484.161}{in}%
\TextAt{8}{160.36}{484.161}{a}%
\TextAt{8}{169.838}{484.161}{canonical}%
\TextAt{8}{220.945}{484.161}{way,}%
\TextAt{8}{246.738}{484.161}{i.e.,}%
\TextAt{8}{268.622}{484.161}{for}%
\TextAt{8}{286.23}{484.161}{every}%
\TextAt{9}{316.846}{484.161}{1}%
\TextAt{9}{324.772}{484.161}{\Character{43}}%
\TextAt{13}{335.985}{484.161}{x}%
\NativeMathAt{345.958}{484.161}{11.955}{\in}%
\TextAt{13}{357.249}{484.161}{U}%
\TextAt{17}{366.448}{488.499}{\Character{40}1\Character{41}}%
\TextAt{8}{381.388}{484.161}{and}%
\TextAt{8}{403.874}{484.161}{every}%
\TextAt{13}{434.489}{484.161}{z}%
\NativeMathAt{443.781}{484.161}{11.955}{\in}%
\TextAt{18}{455.072}{484.161}{Z}%
\TextAt{21}{463.042}{482.367}{p}%
\TextAt{8}{467.803}{484.161}{,}%
\TextAt{8}{474.731}{484.161}{one}%
\TextAt{8}{495.913}{484.161}{has}%
\TextAt{8}{516.51}{484.161}{the}%
\TextAt{8}{79.226}{467.115}{power}%
\TextAt{9}{113.122}{467.115}{\Character{40}1}%
\TextAt{9}{125.651}{467.115}{\Character{43}}%
\TextAt{13}{136.932}{467.115}{x}%
\TextAt{9}{143.584}{467.115}{\Character{41}}%
\TextAt{21}{148.136}{471.453}{z}%
\NativeMathAt{156.231}{467.115}{11.955}{\in}%
\TextAt{13}{167.522}{467.115}{U}%
\TextAt{17}{176.721}{471.453}{\Character{40}1\Character{41}}%
\TextAt{8}{188.04}{467.115}{.}%
\TextAt{8}{196.405}{467.115}{This}%
\TextAt{8}{222.877}{467.115}{is}%
\TextAt{8}{234.393}{467.115}{a}%
\TextAt{8}{243.901}{467.115}{consequence}%
\TextAt{8}{309.717}{467.115}{of}%
\TextAt{8}{322.799}{467.115}{the}%
\TextAt{8}{342.708}{467.115}{fact}%
\TextAt{8}{365.546}{467.115}{that}%
\TextAt{13}{390.667}{467.115}{U}%
\TextAt{17}{399.866}{471.453}{\Character{40}1\Character{41}}%
\TextAt{9}{411.185}{467.115}{\Character{47}}%
\TextAt{13}{417.037}{467.115}{U}%
\TextAt{17}{426.237}{471.453}{\Character{40}}%
\TextAt{21}{429.53}{471.453}{n}%
\TextAt{17}{434.668}{471.453}{\Character{43}1\Character{41}}%
\TextAt{8}{452.93}{467.115}{has}%
\TextAt{8}{473.555}{467.115}{order}%
\TextAt{13}{503.877}{467.115}{q}%
\TextAt{21}{509.496}{471.453}{n}%
\TextAt{8}{518.782}{467.115}{for}%
\TextAt{8}{79.226}{450.069}{all}%
\TextAt{13}{95.883}{450.069}{n}%
\TextAt{8}{107.166}{450.069}{(where}%
\TextAt{13}{145.928}{450.069}{q}%
\TextAt{9}{155.546}{450.069}{\Character{61}}%
\TextAt{8}{168.65}{450.069}{\char35\relax }%
\TextAt{14}{178.405}{450.069}{o}%
\TextAt{9}{184.249}{450.069}{\Character{47}}%
\TextAt{14}{190.102}{450.069}{p}%
\TextAt{8}{196.085}{450.069}{—the}%
\TextAt{8}{228.344}{450.069}{reason}%
\TextAt{8}{265.229}{450.069}{for}%
\TextAt{8}{283.512}{450.069}{this}%
\TextAt{8}{306.732}{450.069}{is}%
\TextAt{8}{318.894}{450.069}{that}%
\TextAt{13}{344.66}{450.069}{U}%
\TextAt{17}{353.859}{454.407}{\Character{40}}%
\TextAt{21}{357.153}{454.407}{i}%
\TextAt{17}{360.036}{454.407}{\Character{41}}%
\TextAt{9}{363.827}{450.069}{\Character{47}}%
\TextAt{13}{369.68}{450.069}{U}%
\TextAt{17}{378.879}{454.407}{\Character{40}}%
\TextAt{21}{382.173}{454.407}{i}%
\TextAt{17}{385.056}{454.407}{\Character{43}1\Character{41}}%
\NativeMathAt{403.667}{450.068}{11.955}{\cong}%
\TextAt{14}{416.964}{450.069}{o}%
\TextAt{9}{422.808}{450.069}{\Character{47}}%
\TextAt{14}{428.661}{450.069}{p}%
\TextAt{8}{434.644}{450.069}{,}%
\TextAt{8}{442.291}{450.069}{by}%
\TextAt{8}{458.948}{450.069}{(3.10),}%
\TextAt{8}{496.53}{450.069}{so}%
\TextAt{8}{511.298}{450.069}{that}%
\TextAt{13}{79.226}{433.023}{U}%
\TextAt{17}{88.425}{437.361}{\Character{40}1\Character{41}}%
\TextAt{9}{99.744}{433.023}{\Character{47}}%
\TextAt{13}{105.597}{433.023}{U}%
\TextAt{17}{114.796}{437.361}{\Character{40}}%
\TextAt{21}{118.089}{437.361}{n}%
\TextAt{17}{123.228}{437.361}{\Character{43}1\Character{41}}%
\TextAt{8}{141.737}{433.023}{is}%
\TextAt{8}{153.501}{433.023}{a}%
\TextAt{18}{163.256}{433.023}{Z}%
\TextAt{9}{171.226}{433.023}{\Character{47}}%
\TextAt{13}{177.079}{433.023}{q}%
\TextAt{21}{182.699}{437.361}{n}%
\TextAt{18}{188.335}{433.023}{Z}%
\TextAt{8}{196.305}{433.023}{-module)}%
\TextAt{8}{246.051}{433.023}{and}%
\TextAt{8}{268.813}{433.023}{of}%
\TextAt{8}{282.143}{433.023}{the}%
\TextAt{8}{302.3}{433.023}{formulas}%
\TextAt{13}{174.576}{405.847}{U}%
\TextAt{17}{183.775}{410.783}{\Character{40}1\Character{41}}%
\TextAt{9}{198.414}{405.847}{\Character{61}}%
\NativeMathAt{210.84}{405.847}{11.955}{\varprojlim}%
\TextAt{21}{216.4}{393.732}{n}%
\TextAt{13}{229.091}{405.847}{U}%
\TextAt{17}{238.29}{410.783}{\Character{40}1\Character{41}}%
\TextAt{9}{249.609}{405.847}{\Character{47}}%
\TextAt{13}{255.462}{405.847}{U}%
\TextAt{17}{264.661}{410.783}{\Character{40}}%
\TextAt{21}{267.955}{410.783}{n}%
\TextAt{17}{273.093}{410.783}{\Character{43}1\Character{41}}%
\TextAt{8}{311.615}{405.847}{and}%
\TextAt{18}{354.391}{405.847}{Z}%
\TextAt{21}{362.361}{404.053}{p}%
\TextAt{9}{370.443}{405.847}{\Character{61}}%
\NativeMathAt{382.868}{405.847}{11.955}{\varprojlim}%
\TextAt{21}{388.429}{393.732}{n}%
\TextAt{18}{401.12}{405.847}{Z}%
\TextAt{9}{409.09}{405.847}{\Character{47}}%
\TextAt{13}{414.943}{405.847}{q}%
\TextAt{21}{420.562}{410.783}{n}%
\TextAt{18}{426.198}{405.847}{Z}%
\TextAt{13}{434.168}{405.847}{\Character{58}}%
\TextAt{8}{79.226}{371.095}{This}%
\TextAt{8}{105.946}{371.095}{obviously}%
\TextAt{8}{157.724}{371.095}{extends}%
\TextAt{8}{200.38}{371.095}{the}%
\TextAt{18}{220.536}{371.095}{Z}%
\TextAt{8}{228.506}{371.095}{-module}%
\TextAt{8}{273.697}{371.095}{structure}%
\TextAt{8}{323.837}{371.095}{of}%
\TextAt{13}{337.167}{371.095}{U}%
\TextAt{17}{346.366}{375.433}{\Character{40}1\Character{41}}%
\TextAt{8}{357.685}{371.095}{.}%
\TextAt{8}{366.133}{371.095}{The}%
\TextAt{8}{390.187}{371.095}{function}%
\TextAt{13}{266.663}{343.919}{f}%
\TextAt{9}{273.709}{343.919}{\Character{40}}%
\TextAt{13}{278.262}{343.919}{z}%
\TextAt{9}{284.232}{343.919}{\Character{41}}%
\TextAt{9}{292.087}{343.919}{\Character{61}}%
\TextAt{9}{304.508}{343.919}{\Character{40}1}%
\TextAt{9}{317.551}{343.919}{\Character{43}}%
\TextAt{13}{329.354}{343.919}{x}%
\TextAt{9}{336.006}{343.919}{\Character{41}}%
\TextAt{21}{340.559}{348.855}{z}%
\TextAt{8}{79.226}{316.743}{is}%
\TextAt{8}{91.384}{316.743}{continuous}%
\TextAt{8}{150.383}{316.743}{because}%
\TextAt{8}{194.079}{316.743}{the}%
\TextAt{8}{214.63}{316.743}{congruence}%
\TextAt{13}{275.505}{316.743}{z}%
\NativeMathAt{285.469}{316.743}{11.955}{\equiv}%
\TextAt{13}{298.76}{316.743}{z}%
\TextAt{3}{304.731}{321.081}{\Character{48}}%
\TextAt{8}{310.847}{316.743}{mod}%
\TextAt{13}{336.285}{316.743}{q}%
\TextAt{21}{341.904}{321.081}{n}%
\TextAt{18}{347.54}{316.743}{Z}%
\TextAt{21}{355.51}{314.95}{p}%
\TextAt{8}{364.563}{316.743}{implies}%
\TextAt{9}{404.685}{316.743}{\Character{40}1}%
\TextAt{9}{417.991}{316.743}{\Character{43}}%
\TextAt{13}{430.045}{316.743}{x}%
\TextAt{9}{436.698}{316.743}{\Character{41}}%
\TextAt{21}{441.25}{321.081}{z}%
\NativeMathAt{450.017}{316.743}{11.955}{\equiv}%
\TextAt{9}{463.308}{316.743}{\Character{40}1}%
\TextAt{9}{476.614}{316.743}{\Character{43}}%
\TextAt{13}{488.669}{316.743}{x}%
\TextAt{9}{495.321}{316.743}{\Character{41}}%
\TextAt{21}{499.873}{321.081}{z}%
\TextAt{23}{504.149}{323.894}{\Character{48}}%
\TextAt{8}{510.652}{316.743}{mod}%
\TextAt{13}{79.226}{299.697}{U}%
\TextAt{17}{88.425}{304.035}{\Character{40}}%
\TextAt{21}{91.718}{304.035}{n}%
\TextAt{17}{96.857}{304.035}{\Character{43}1\Character{41}}%
\TextAt{8}{111.469}{299.697}{,}%
\TextAt{8}{121.048}{299.697}{so}%
\TextAt{8}{137.362}{299.697}{that}%
\TextAt{8}{164.676}{299.697}{the}%
\TextAt{8}{186.776}{299.697}{neighbourhood}%
\TextAt{13}{268.724}{299.697}{z}%
\TextAt{9}{278.677}{299.697}{\Character{43}}%
\TextAt{13}{291.763}{299.697}{q}%
\TextAt{21}{297.383}{304.035}{n}%
\TextAt{18}{303.019}{299.697}{Z}%
\TextAt{21}{310.989}{297.904}{p}%
\TextAt{8}{321.592}{299.697}{of}%
\TextAt{13}{336.866}{299.697}{z}%
\TextAt{8}{348.678}{299.697}{is}%
\TextAt{8}{362.386}{299.697}{mapped}%
\TextAt{8}{408.875}{299.697}{into}%
\TextAt{8}{434.562}{299.697}{the}%
\TextAt{8}{456.663}{299.697}{neighbourhood}%
\TextAt{9}{79.226}{282.651}{\Character{40}1}%
\TextAt{9}{92.269}{282.651}{\Character{43}}%
\TextAt{13}{104.049}{282.651}{x}%
\TextAt{9}{110.702}{282.651}{\Character{41}}%
\TextAt{21}{115.254}{286.989}{z}%
\TextAt{13}{120.028}{282.651}{U}%
\TextAt{17}{129.227}{286.989}{\Character{40}}%
\TextAt{21}{132.52}{286.989}{n}%
\TextAt{17}{137.658}{286.989}{\Character{43}1\Character{41}}%
\TextAt{8}{156.168}{282.651}{of}%
\TextAt{13}{169.498}{282.651}{f}%
\TextAt{9}{176.544}{282.651}{\Character{40}}%
\TextAt{13}{181.096}{282.651}{z}%
\TextAt{9}{187.067}{282.651}{\Character{41}}%
\TextAt{8}{191.619}{282.651}{.}%
\TextAt{8}{200.068}{282.651}{In}%
\TextAt{8}{214.689}{282.651}{particular,}%
\TextAt{9}{271.93}{282.651}{\Character{40}1}%
\TextAt{9}{284.973}{282.651}{\Character{43}}%
\TextAt{13}{296.754}{282.651}{x}%
\TextAt{9}{303.406}{282.651}{\Character{41}}%
\TextAt{21}{307.958}{286.989}{z}%
\TextAt{8}{316.629}{282.651}{may}%
\TextAt{8}{341.998}{282.651}{be}%
\TextAt{8}{357.922}{282.651}{expressed}%
\TextAt{8}{410.394}{282.651}{as}%
\TextAt{8}{424.764}{282.651}{the}%
\TextAt{8}{444.92}{282.651}{limit}%
\TextAt{9}{246.001}{255.475}{\Character{40}1}%
\TextAt{9}{259.044}{255.475}{\Character{43}}%
\TextAt{13}{270.824}{255.475}{x}%
\TextAt{9}{277.476}{255.475}{\Character{41}}%
\TextAt{21}{282.028}{260.411}{z}%
\TextAt{9}{290.123}{255.475}{\Character{61}}%
\TextAt{8}{304.329}{255.475}{lim}%
\TextAt{21}{302.548}{248.214}{i}%
\NativeMathAt{305.432}{248.214}{7.97}{\rightarrow}%
\NativeMathAt{313.896}{248.214}{7.97}{\infty}%
\TextAt{9}{322.368}{255.475}{\Character{40}1}%
\TextAt{9}{335.411}{255.475}{\Character{43}}%
\TextAt{13}{347.192}{255.475}{x}%
\TextAt{9}{353.844}{255.475}{\Character{41}}%
\TextAt{21}{358.396}{260.411}{z}%
\TextAt{22}{362.335}{259.196}{i}%
\TextAt{8}{79.226}{226.925}{of}%
\TextAt{8}{92.556}{226.925}{ordinary}%
\TextAt{8}{139.719}{226.925}{powers}%
\TextAt{9}{178.478}{226.925}{\Character{40}1}%
\TextAt{9}{191.521}{226.925}{\Character{43}}%
\TextAt{13}{203.301}{226.925}{x}%
\TextAt{9}{209.953}{226.925}{\Character{41}}%
\TextAt{21}{214.506}{231.264}{z}%
\TextAt{22}{218.445}{230.049}{i}%
\TextAt{8}{222.105}{226.925}{,}%
\TextAt{13}{229.254}{226.925}{z}%
\TextAt{21}{234.692}{225.132}{i}%
\NativeMathAt{241.394}{226.925}{11.955}{\in}%
\TextAt{18}{252.685}{226.925}{Z}%
\TextAt{8}{260.655}{226.925}{,}%
\TextAt{8}{267.804}{226.925}{if}%
\TextAt{13}{278.528}{226.925}{z}%
\TextAt{9}{287.819}{226.925}{\Character{61}}%
\TextAt{8}{300.245}{226.925}{lim}%
\TextAt{21}{316.504}{225.132}{i}%
\NativeMathAt{319.387}{225.132}{7.97}{\rightarrow}%
\NativeMathAt{327.851}{225.132}{7.97}{\infty}%
\TextAt{13}{338.814}{226.925}{z}%
\TextAt{21}{344.252}{225.132}{i}%
\TextAt{8}{347.633}{226.925}{.}%
\TextAt{8}{97.159}{206.891}{After}%
\TextAt{8}{129.552}{206.891}{this}%
\TextAt{8}{154.211}{206.891}{discussion}%
\TextAt{8}{210.862}{206.891}{we}%
\TextAt{8}{229.925}{206.891}{can}%
\TextAt{8}{253.221}{206.891}{now}%
\TextAt{8}{279.445}{206.891}{determine}%
\TextAt{8}{335.905}{206.891}{explicitly}%
\TextAt{8}{388.467}{206.891}{the}%
\TextAt{8}{410.459}{206.891}{structure}%
\TextAt{8}{462.435}{206.891}{of}%
\TextAt{8}{477.601}{206.891}{the}%
\TextAt{8}{499.594}{206.891}{locally}%
\TextAt{8}{79.226}{189.845}{compact}%
\TextAt{8}{126.055}{189.845}{multiplicative}%
\TextAt{8}{199.878}{189.845}{group}%
\TextAt{13}{233.053}{189.845}{K}%
\TextAt{3}{243.871}{194.183}{\Character{3}}%
\TextAt{8}{252.501}{189.845}{of}%
\TextAt{8}{265.831}{189.845}{a}%
\TextAt{8}{275.586}{189.845}{local}%
\TextAt{8}{303.226}{189.845}{field}%
\TextAt{13}{328.583}{189.845}{K}%
\TextAt{8}{339.401}{189.845}{.}%
\TextAt{12}{79.226}{164.354}{Proposition}%
\TextAt{12}{153.051}{164.354}{(5.7).}%
\TextAt{8}{189.334}{164.354}{Let}%
\TextAt{13}{211.271}{164.354}{K}%
\TextAt{8}{226.954}{164.354}{be}%
\TextAt{8}{243.846}{164.354}{a}%
\TextAt{8}{254.569}{164.354}{local}%
\TextAt{8}{283.177}{164.354}{field}%
\TextAt{8}{309.501}{164.354}{and}%
\TextAt{13}{333.231}{164.354}{q}%
\TextAt{9}{343.821}{164.354}{\Character{61}}%
\TextAt{13}{357.895}{164.354}{p}%
\TextAt{21}{363.77}{168.692}{f}%
\TextAt{8}{374.08}{164.354}{the}%
\TextAt{8}{395.204}{164.354}{number}%
\TextAt{8}{438.768}{164.354}{of}%
\TextAt{8}{453.066}{164.354}{elements}%
\TextAt{8}{501.89}{164.354}{in}%
\TextAt{8}{516.51}{164.354}{the}%
\TextAt{8}{79.226}{147.308}{residue}%
\TextAt{8}{118.953}{147.308}{class}%
\TextAt{8}{146.39}{147.308}{field.}%
\TextAt{8}{176.298}{147.308}{Then}%
\TextAt{8}{206.856}{147.308}{the}%
\TextAt{8}{227.012}{147.308}{following}%
\TextAt{8}{276.447}{147.308}{hold.}%
\TextAt{8}{89.579}{127.273}{(i)}%
\TextAt{8}{107.918}{127.273}{If}%
\TextAt{13}{119.611}{127.273}{K}%
\TextAt{8}{134.326}{127.273}{has}%
\TextAt{8}{155.199}{127.273}{characteristic}%
\TextAt{8}{227.134}{127.273}{0,}%
\TextAt{8}{240.141}{127.273}{then}%
\TextAt{8}{266.801}{127.273}{one}%
\TextAt{8}{288.26}{127.273}{has}%
\TextAt{8}{309.134}{127.273}{(both}%
\TextAt{8}{341.329}{127.273}{algebraically}%
\TextAt{8}{409.307}{127.273}{and}%
\TextAt{8}{432.069}{127.273}{topologically)}%
\TextAt{13}{225.834}{100.097}{K}%
\TextAt{3}{236.652}{105.033}{\Character{3}}%
\NativeMathAt{244.705}{100.097}{11.955}{\cong}%
\TextAt{18}{257.324}{100.097}{Z}%
\TextAt{16}{267.951}{100.097}{\Character{8}}%
\TextAt{18}{279.906}{100.097}{Z}%
\TextAt{9}{287.877}{100.097}{\Character{47}\Character{40}}%
\TextAt{13}{298.282}{100.097}{q}%
\TextAt{16}{306.558}{100.097}{\Character{0}}%
\TextAt{9}{318.513}{100.097}{1\Character{41}}%
\TextAt{18}{328.918}{100.097}{Z}%
\TextAt{16}{339.545}{100.097}{\Character{8}}%
\TextAt{18}{351.5}{100.097}{Z}%
\TextAt{9}{359.47}{100.097}{\Character{47}}%
\TextAt{13}{365.323}{100.097}{p}%
\TextAt{21}{371.198}{105.033}{a}%
\TextAt{18}{376.195}{100.097}{Z}%
\TextAt{16}{386.821}{100.097}{\Character{8}}%
\TextAt{18}{398.777}{100.097}{Z}%
\TextAt{21}{406.747}{105.033}{d}%
\TextAt{21}{406.747}{97.142}{p}%
\TextAt{13}{411.602}{100.097}{\Character{59}}%
\TextAt{8}{107.918}{72.921}{where}%
\TextAt{13}{141.728}{72.921}{a}%
\NativeMathAt{151.193}{72.921}{11.955}{\geq}%
\TextAt{9}{163.813}{72.921}{0}%
\TextAt{8}{173.563}{72.921}{and}%
\TextAt{13}{196.326}{72.921}{d}%
\TextAt{9}{205.729}{72.921}{\Character{61}}%
\TextAt{9}{218.139}{72.921}{\Character{91}}%
\TextAt{13}{221.406}{72.921}{K}%
\TextAt{9}{235.545}{72.921}{\Character{58}}%
\TextAt{18}{242.117}{72.921}{Q}%
\TextAt{21}{251.416}{71.128}{p}%
\TextAt{9}{256.177}{72.921}{\Character{93}}%
\TextAt{8}{259.429}{72.921}{.}%
\TextAt{8}{297.211}{43.033}{127}%
\Destination{79.23}{141.865}{nav.Item.42}%
\Destination{78.23}{809.815}{nav.page.127}%
\Destination{79.36}{713.755}{nav.print.139}%
\Destination{79.36}{178.355}{nav.number.2.5.7}%
\Hotlink{458.448}{447.855}{30.936}{11.047}{nav.number.2.3.10}%
}

% Page 140 (body page 128)
\Page{%
\TextAt{2}{79.226}{796.943}{II.}%
\TextAt{2}{94.073}{796.943}{The}%
\TextAt{2}{116.492}{796.943}{Theory}%
\TextAt{2}{154.401}{796.943}{of}%
\TextAt{2}{166.826}{796.943}{Valuations}%
\Vector{q 0 G 0 g 0.299 w 7.226 18.729 m 460.769 18.729 l S Q}%
\TextAt{8}{86.327}{761.124}{(ii)}%
\TextAt{8}{107.918}{761.124}{If}%
\TextAt{13}{119.611}{761.124}{K}%
\TextAt{8}{134.326}{761.124}{has}%
\TextAt{8}{155.199}{761.124}{characteristic}%
\TextAt{13}{227.134}{761.124}{p}%
\TextAt{8}{233.009}{761.124}{,}%
\TextAt{8}{240.158}{761.124}{then}%
\TextAt{8}{266.818}{761.124}{one}%
\TextAt{8}{288.278}{761.124}{has}%
\TextAt{8}{309.151}{761.124}{(both}%
\TextAt{8}{341.347}{761.124}{algebraically}%
\TextAt{8}{409.324}{761.124}{and}%
\TextAt{8}{432.086}{761.124}{topologically)}%
\TextAt{13}{248.773}{734.039}{K}%
\TextAt{3}{259.591}{738.975}{\Character{3}}%
\NativeMathAt{267.644}{734.039}{11.955}{\cong}%
\TextAt{18}{280.263}{734.039}{Z}%
\TextAt{16}{290.89}{734.039}{\Character{8}}%
\TextAt{18}{302.845}{734.039}{Z}%
\TextAt{9}{310.815}{734.039}{\Character{47}\Character{40}}%
\TextAt{13}{321.221}{734.039}{q}%
\TextAt{16}{329.496}{734.039}{\Character{0}}%
\TextAt{9}{341.452}{734.039}{1\Character{41}}%
\TextAt{18}{351.857}{734.039}{Z}%
\TextAt{16}{362.484}{734.039}{\Character{8}}%
\TextAt{18}{374.439}{734.039}{Z}%
\TextAt{25}{382.409}{738.975}{N}%
\TextAt{21}{382.409}{731.083}{p}%
\TextAt{13}{388.663}{734.039}{\Character{58}}%
\TextAt{19}{79.226}{693.105}{Proof.}%
\TextAt{8}{116}{693.105}{By}%
\TextAt{8}{134.363}{693.105}{(5.3)}%
\TextAt{8}{162.338}{693.105}{we}%
\TextAt{8}{179.566}{693.105}{have}%
\TextAt{8}{206.561}{693.105}{(both}%
\TextAt{8}{238.756}{693.105}{algebraically}%
\TextAt{8}{306.733}{693.105}{and}%
\TextAt{8}{329.496}{693.105}{topologically)}%
\TextAt{13}{179.235}{666.02}{K}%
\TextAt{3}{190.053}{670.956}{\Character{3}}%
\TextAt{9}{198.106}{666.02}{\Character{61}}%
\TextAt{9}{210.516}{666.02}{\Character{40}}%
\TextAt{13}{215.084}{666.02}{\Character{25}}%
\TextAt{9}{222.153}{666.02}{\Character{41}}%
\TextAt{16}{229.362}{666.02}{\Character{2}}%
\TextAt{13}{241.317}{666.02}{\Character{22}}%
\TextAt{21}{248.36}{664.227}{q}%
\TextAt{3}{252.427}{664.227}{\Character{0}}%
\TextAt{17}{259.013}{664.227}{1}%
\TextAt{16}{266.402}{666.02}{\Character{2}}%
\TextAt{13}{278.357}{666.02}{U}%
\TextAt{17}{287.557}{670.956}{\Character{40}1\Character{41}}%
\NativeMathAt{302.196}{666.02}{11.955}{\cong}%
\TextAt{18}{314.815}{666.02}{Z}%
\TextAt{16}{325.442}{666.02}{\Character{8}}%
\TextAt{18}{337.397}{666.02}{Z}%
\TextAt{9}{345.368}{666.02}{\Character{47}\Character{40}}%
\TextAt{13}{355.773}{666.02}{q}%
\TextAt{16}{364.049}{666.02}{\Character{0}}%
\TextAt{9}{376.004}{666.02}{1\Character{41}}%
\TextAt{18}{386.409}{666.02}{Z}%
\TextAt{16}{397.036}{666.02}{\Character{8}}%
\TextAt{13}{408.991}{666.02}{U}%
\TextAt{17}{418.19}{670.956}{\Character{40}1\Character{41}}%
\TextAt{13}{429.509}{666.02}{\Character{58}}%
\TextAt{8}{79.226}{638.935}{This}%
\TextAt{8}{105.946}{638.935}{reduces}%
\TextAt{8}{147.622}{638.935}{us}%
\TextAt{8}{162.637}{638.935}{to}%
\TextAt{8}{176.948}{638.935}{the}%
\TextAt{8}{197.104}{638.935}{computation}%
\TextAt{8}{265.404}{638.935}{of}%
\TextAt{8}{278.734}{638.935}{the}%
\TextAt{18}{298.89}{638.935}{Z}%
\TextAt{21}{306.861}{637.142}{p}%
\TextAt{8}{311.621}{638.935}{-module}%
\TextAt{13}{356.812}{638.935}{U}%
\TextAt{17}{366.011}{643.273}{\Character{40}1\Character{41}}%
\TextAt{8}{377.33}{638.935}{.}%
\TextAt{8}{97.159}{618.9}{(i)}%
\TextAt{8}{113.418}{618.9}{Assume}%
\TextAt{8}{156.779}{618.9}{char}%
\TextAt{9}{178.896}{618.9}{\Character{40}}%
\TextAt{13}{183.449}{618.9}{K}%
\TextAt{9}{194.266}{618.9}{\Character{41}}%
\TextAt{9}{202.121}{618.9}{\Character{61}}%
\TextAt{9}{214.542}{618.9}{0}%
\TextAt{8}{220.418}{618.9}{.}%
\TextAt{8}{228.866}{618.9}{For}%
\TextAt{13}{249.836}{618.9}{n}%
\TextAt{8}{260.721}{618.9}{sufficiently}%
\TextAt{8}{319.313}{618.9}{big,}%
\TextAt{8}{342.076}{618.9}{(5.5)}%
\TextAt{8}{370.051}{618.9}{gives}%
\TextAt{8}{398.731}{618.9}{us}%
\TextAt{8}{413.747}{618.9}{the}%
\TextAt{8}{433.903}{618.9}{isomorphism}%
\TextAt{8}{234.429}{591.815}{log}%
\TextAt{9}{252.718}{591.815}{\Character{58}}%
\TextAt{13}{259.29}{591.815}{U}%
\TextAt{17}{268.489}{596.752}{\Character{40}}%
\TextAt{21}{271.783}{596.752}{n}%
\TextAt{17}{276.921}{596.752}{\Character{41}}%
\NativeArrowAt{284.033}{591.815}{11.955}{19.26}{\longrightarrow}%
\TextAt{14}{306.615}{591.815}{p}%
\TextAt{21}{312.598}{596.752}{n}%
\TextAt{9}{321.555}{591.815}{\Character{61}}%
\TextAt{13}{333.981}{591.815}{\Character{25}}%
\TextAt{21}{341.05}{596.752}{n}%
\TextAt{14}{346.686}{591.815}{o}%
\NativeMathAt{355.851}{591.815}{11.955}{\cong}%
\TextAt{14}{368.471}{591.815}{o}%
\TextAt{13}{374.315}{591.815}{\Character{58}}%
\TextAt{8}{79.226}{564.73}{Since}%
\TextAt{8}{111.039}{564.73}{log,}%
\TextAt{8}{134.726}{564.73}{exp,}%
\TextAt{8}{161.33}{564.73}{and}%
\TextAt{13}{185.349}{564.73}{f}%
\TextAt{9}{192.395}{564.73}{\Character{40}}%
\TextAt{13}{196.948}{564.73}{z}%
\TextAt{9}{202.918}{564.73}{\Character{41}}%
\TextAt{9}{212.913}{564.73}{\Character{61}}%
\TextAt{9}{227.474}{564.73}{\Character{40}1}%
\TextAt{9}{241.378}{564.73}{\Character{43}}%
\TextAt{13}{254.033}{564.73}{x}%
\TextAt{9}{260.685}{564.73}{\Character{41}}%
\TextAt{21}{265.237}{569.069}{z}%
\TextAt{8}{275.165}{564.73}{are}%
\TextAt{8}{295.931}{564.73}{continuous,}%
\TextAt{8}{359.357}{564.73}{this}%
\TextAt{8}{383.436}{564.73}{is}%
\TextAt{8}{396.455}{564.73}{a}%
\TextAt{8}{407.467}{564.73}{topological}%
\TextAt{8}{468.247}{564.73}{isomorphism}%
\TextAt{8}{79.226}{547.684}{of}%
\TextAt{18}{94.238}{547.684}{Z}%
\TextAt{21}{102.209}{545.891}{p}%
\TextAt{8}{106.97}{547.684}{-modules.}%
\TextAt{8}{166.373}{547.684}{By}%
\TextAt{8}{186.418}{547.684}{chap.}%
\TextAt{8}{218.993}{547.684}{I,}%
\TextAt{8}{232.044}{547.684}{(2.9),}%
\TextAt{14}{265.374}{547.684}{o}%
\TextAt{8}{276.798}{547.684}{admits}%
\TextAt{8}{316.916}{547.684}{an}%
\TextAt{8}{334.858}{547.684}{integral}%
\TextAt{8}{379.148}{547.684}{basis}%
\TextAt{13}{409.571}{547.684}{\Character{11}}%
\TextAt{17}{417.016}{545.891}{1}%
\TextAt{13}{421.749}{547.684}{\Character{59}}%
\TextAt{13}{426.973}{547.684}{\Character{58}}%
\TextAt{13}{432.21}{547.684}{\Character{58}}%
\TextAt{13}{437.434}{547.684}{\Character{58}}%
\TextAt{13}{442.671}{547.684}{\Character{59}}%
\TextAt{13}{447.907}{547.684}{\Character{11}}%
\TextAt{21}{455.415}{545.891}{d}%
\TextAt{8}{465.85}{547.684}{over}%
\TextAt{18}{492.579}{547.684}{Z}%
\TextAt{21}{500.549}{545.891}{p}%
\TextAt{8}{505.31}{547.684}{,}%
\TextAt{8}{514.562}{547.684}{i.e.,}%
\TextAt{14}{79.226}{530.638}{o}%
\TextAt{9}{88.699}{530.638}{\Character{61}}%
\TextAt{18}{101.433}{530.638}{Z}%
\TextAt{21}{109.403}{528.845}{p}%
\TextAt{13}{114.164}{530.638}{\Character{11}}%
\TextAt{17}{121.61}{528.845}{1}%
\TextAt{16}{129.122}{530.638}{\Character{8}}%
\TextAt{16}{141.185}{530.638}{\Character{1}}%
\TextAt{16}{146.493}{530.638}{\Character{1}}%
\TextAt{16}{151.789}{530.638}{\Character{1}}%
\TextAt{16}{157.886}{530.638}{\Character{8}}%
\TextAt{18}{170.006}{530.638}{Z}%
\TextAt{21}{177.976}{528.845}{p}%
\TextAt{13}{182.737}{530.638}{\Character{11}}%
\TextAt{21}{190.183}{528.845}{d}%
\NativeMathAt{198.667}{530.638}{11.955}{\cong}%
\TextAt{18}{211.595}{530.638}{Z}%
\TextAt{21}{219.565}{534.977}{d}%
\TextAt{21}{219.565}{527.683}{p}%
\TextAt{8}{224.42}{530.638}{.}%
\TextAt{8}{233.411}{530.638}{Therefore}%
\TextAt{13}{286.589}{530.638}{U}%
\TextAt{17}{295.788}{534.977}{\Character{40}}%
\TextAt{21}{299.082}{534.977}{n}%
\TextAt{17}{304.22}{534.977}{\Character{41}}%
\NativeMathAt{311.64}{530.638}{11.955}{\cong}%
\TextAt{18}{324.568}{530.638}{Z}%
\TextAt{21}{332.538}{534.977}{d}%
\TextAt{21}{332.538}{527.683}{p}%
\TextAt{8}{337.393}{530.638}{.}%
\TextAt{8}{346.384}{530.638}{Since}%
\TextAt{8}{377.122}{530.638}{the}%
\TextAt{8}{397.46}{530.638}{index}%
\TextAt{9}{429.178}{530.638}{\Character{40}}%
\TextAt{13}{433.731}{530.638}{U}%
\TextAt{17}{442.93}{534.977}{\Character{40}1\Character{41}}%
\TextAt{9}{457.877}{530.638}{\Character{58}}%
\TextAt{13}{464.758}{530.638}{U}%
\TextAt{17}{473.957}{534.977}{\Character{40}}%
\TextAt{21}{477.25}{534.977}{n}%
\TextAt{17}{482.389}{534.977}{\Character{41}}%
\TextAt{9}{486.18}{530.638}{\Character{41}}%
\TextAt{8}{494.81}{530.638}{is}%
\TextAt{8}{506.755}{530.638}{finite}%
\TextAt{8}{79.226}{513.593}{and}%
\TextAt{13}{102.061}{513.593}{U}%
\TextAt{17}{111.26}{517.931}{\Character{40}}%
\TextAt{21}{114.553}{517.931}{n}%
\TextAt{17}{119.691}{517.931}{\Character{41}}%
\TextAt{8}{127.452}{513.593}{is}%
\TextAt{8}{139.287}{513.593}{a}%
\TextAt{8}{149.114}{513.593}{finitely}%
\TextAt{8}{188.531}{513.593}{generated}%
\TextAt{18}{241.934}{513.593}{Z}%
\TextAt{21}{249.905}{511.799}{p}%
\TextAt{8}{254.666}{513.593}{-module}%
\TextAt{8}{299.928}{513.593}{of}%
\TextAt{8}{313.33}{513.593}{rank}%
\TextAt{13}{340.396}{513.593}{d}%
\TextAt{8}{346.479}{513.593}{,}%
\TextAt{8}{353.718}{513.593}{so}%
\TextAt{8}{368.16}{513.593}{is}%
\TextAt{13}{379.995}{513.593}{U}%
\TextAt{17}{389.194}{517.931}{\Character{40}1\Character{41}}%
\TextAt{8}{400.513}{513.593}{.}%
\TextAt{8}{409.177}{513.593}{The}%
\TextAt{8}{433.302}{513.593}{torsion}%
\TextAt{8}{472.468}{513.593}{subgroup}%
\TextAt{8}{523.337}{513.593}{of}%
\TextAt{13}{79.226}{496.547}{U}%
\TextAt{17}{88.425}{500.885}{\Character{40}1\Character{41}}%
\TextAt{8}{104.166}{496.547}{is}%
\TextAt{8}{116.455}{496.547}{the}%
\TextAt{8}{137.137}{496.547}{group}%
\TextAt{13}{170.838}{496.547}{\Character{22}}%
\TextAt{21}{177.88}{494.753}{p}%
\TextAt{22}{182.143}{497.023}{a}%
\TextAt{8}{191.596}{496.547}{of}%
\TextAt{8}{205.451}{496.547}{roots}%
\TextAt{8}{235.637}{496.547}{of}%
\TextAt{8}{249.492}{496.547}{unity}%
\TextAt{8}{280.587}{496.547}{in}%
\TextAt{13}{294.764}{496.547}{K}%
\TextAt{8}{310.004}{496.547}{of}%
\TextAt{13}{323.86}{496.547}{p}%
\TextAt{8}{329.735}{496.547}{-power}%
\TextAt{8}{368.301}{496.547}{order.}%
\TextAt{8}{404.997}{496.547}{By}%
\TextAt{8}{423.885}{496.547}{the}%
\TextAt{8}{444.566}{496.547}{main}%
\TextAt{8}{474.358}{496.547}{theorem}%
\TextAt{8}{520.408}{496.547}{on}%
\TextAt{8}{79.226}{479.501}{modules}%
\TextAt{8}{125.807}{479.501}{over}%
\TextAt{8}{151.526}{479.501}{principal}%
\TextAt{8}{200.976}{479.501}{ideal}%
\TextAt{8}{229.612}{479.501}{domains,}%
\TextAt{8}{279.947}{479.501}{there}%
\TextAt{8}{310.532}{479.501}{exists}%
\TextAt{8}{343.519}{479.501}{in}%
\TextAt{13}{357.845}{479.501}{U}%
\TextAt{17}{367.044}{483.839}{\Character{40}1\Character{41}}%
\TextAt{8}{382.933}{479.501}{a}%
\TextAt{8}{393.361}{479.501}{free,}%
\TextAt{8}{419.882}{479.501}{finitely}%
\TextAt{8}{459.899}{479.501}{generated}%
\TextAt{8}{513.904}{479.501}{and}%
\TextAt{8}{79.226}{462.455}{therefore}%
\TextAt{8}{128.326}{462.455}{closed,}%
\TextAt{18}{166.104}{462.455}{Z}%
\TextAt{21}{174.075}{460.661}{p}%
\TextAt{8}{178.835}{462.455}{-submodule}%
\TextAt{13}{241.648}{462.455}{V}%
\TextAt{8}{254.971}{462.455}{of}%
\TextAt{8}{268.301}{462.455}{rank}%
\TextAt{13}{295.296}{462.455}{d}%
\TextAt{8}{305.276}{462.455}{such}%
\TextAt{8}{331.673}{462.455}{that}%
\TextAt{13}{228.032}{435.37}{U}%
\TextAt{17}{237.231}{440.306}{\Character{40}1\Character{41}}%
\TextAt{9}{251.871}{435.37}{\Character{61}}%
\TextAt{13}{264.296}{435.37}{\Character{22}}%
\TextAt{21}{271.339}{433.576}{p}%
\TextAt{22}{275.602}{435.846}{a}%
\TextAt{16}{283.289}{435.37}{\Character{2}}%
\TextAt{13}{295.244}{435.37}{V}%
\NativeMathAt{307.991}{435.37}{11.955}{\cong}%
\TextAt{18}{320.61}{435.37}{Z}%
\TextAt{9}{328.58}{435.37}{\Character{47}}%
\TextAt{13}{334.433}{435.37}{p}%
\TextAt{21}{340.308}{440.306}{a}%
\TextAt{18}{345.304}{435.37}{Z}%
\TextAt{16}{355.931}{435.37}{\Character{8}}%
\TextAt{18}{367.886}{435.37}{Z}%
\TextAt{21}{375.857}{440.306}{d}%
\TextAt{21}{375.857}{432.414}{p}%
\TextAt{13}{380.712}{435.37}{\Character{59}}%
\TextAt{8}{79.226}{408.285}{both}%
\TextAt{8}{106.866}{408.285}{algebraically}%
\TextAt{8}{174.844}{408.285}{and}%
\TextAt{8}{197.606}{408.285}{topologically.}%
\TextAt{8}{97.159}{388.25}{(ii)}%
\TextAt{8}{116.67}{388.25}{If}%
\TextAt{8}{128.362}{388.25}{char}%
\TextAt{9}{150.479}{388.25}{\Character{40}}%
\TextAt{13}{155.031}{388.25}{K}%
\TextAt{9}{165.849}{388.25}{\Character{41}}%
\TextAt{9}{173.704}{388.25}{\Character{61}}%
\TextAt{13}{186.148}{388.25}{p}%
\TextAt{8}{192.023}{388.25}{,}%
\TextAt{8}{199.172}{388.25}{we}%
\TextAt{8}{216.399}{388.25}{have}%
\TextAt{13}{243.394}{388.25}{K}%
\NativeMathAt{257.533}{388.25}{11.955}{\cong}%
\TextAt{18}{270.152}{388.25}{F}%
\TextAt{21}{277.458}{386.457}{q}%
\TextAt{9}{282.023}{388.25}{\Character{40}\Character{40}}%
\TextAt{13}{291.127}{388.25}{t}%
\TextAt{9}{295.355}{388.25}{\Character{41}\Character{41}}%
\TextAt{8}{308.357}{388.25}{(see}%
\TextAt{8}{331.825}{388.25}{p.}%
\TextAt{8}{345.477}{388.25}{127)}%
\TextAt{8}{371.504}{388.25}{and}%
\TextAt{13}{238.438}{361.165}{U}%
\TextAt{17}{247.637}{366.101}{\Character{40}1\Character{41}}%
\TextAt{9}{262.277}{361.165}{\Character{61}}%
\TextAt{9}{274.687}{361.165}{1}%
\TextAt{9}{283.187}{361.165}{\Character{43}}%
\TextAt{14}{294.973}{361.165}{p}%
\TextAt{9}{304.277}{361.165}{\Character{61}}%
\TextAt{9}{316.687}{361.165}{1}%
\TextAt{9}{325.199}{361.165}{\Character{43}}%
\TextAt{13}{336.974}{361.165}{t}%
\TextAt{18}{341.201}{361.165}{F}%
\TextAt{21}{348.507}{359.372}{q}%
\TextAt{9}{353.072}{361.165}{\Character{91}\Character{91}}%
\TextAt{13}{359.575}{361.165}{t}%
\TextAt{9}{363.802}{361.165}{\Character{93}\Character{93}}%
\TextAt{13}{370.306}{361.165}{\Character{58}}%
\TextAt{8}{79.226}{334.08}{The}%
\TextAt{8}{103.28}{334.08}{following}%
\TextAt{8}{152.715}{334.08}{argument}%
\TextAt{8}{205.078}{334.08}{is}%
\TextAt{8}{216.842}{334.08}{taken}%
\TextAt{8}{248.715}{334.08}{from}%
\TextAt{8}{276.355}{334.08}{the}%
\TextAt{8}{296.511}{334.08}{book}%
\TextAt{8}{325.455}{334.08}{[79]}%
\TextAt{8}{347.572}{334.08}{of}%
\TextAt{26}{361.416}{334.08}{K.}%
\TextAt{26}{379.827}{334.08}{Iwasawa}%
\TextAt{8}{428.939}{334.08}{.}%
\TextAt{8}{97.159}{314.045}{Let}%
\TextAt{13}{118.054}{314.045}{\Character{33}}%
\TextAt{17}{125.343}{312.252}{1}%
\TextAt{13}{130.075}{314.045}{\Character{59}}%
\TextAt{13}{135.299}{314.045}{\Character{58}}%
\TextAt{13}{140.536}{314.045}{\Character{58}}%
\TextAt{13}{145.772}{314.045}{\Character{58}}%
\TextAt{13}{150.997}{314.045}{\Character{59}}%
\TextAt{13}{156.233}{314.045}{\Character{33}}%
\TextAt{21}{163.586}{312.252}{f}%
\TextAt{8}{172.853}{314.045}{be}%
\TextAt{8}{188.703}{314.045}{a}%
\TextAt{8}{198.384}{314.045}{basis}%
\TextAt{8}{227.05}{314.045}{of}%
\TextAt{18}{240.306}{314.045}{F}%
\TextAt{21}{247.612}{312.252}{q}%
\TextAt{16}{255.498}{314.045}{\Character{106}}%
\TextAt{18}{262.139}{314.045}{F}%
\TextAt{21}{269.445}{312.252}{p}%
\TextAt{8}{274.206}{314.045}{.}%
\TextAt{8}{282.63}{314.045}{For}%
\TextAt{8}{303.525}{314.045}{every}%
\TextAt{8}{334.343}{314.045}{natural}%
\TextAt{8}{375.25}{314.045}{number}%
\TextAt{13}{417.772}{314.045}{n}%
\TextAt{8}{428.583}{314.045}{relatively}%
\TextAt{8}{479.569}{314.045}{prime}%
\TextAt{8}{512.658}{314.045}{to}%
\TextAt{13}{526.894}{314.045}{p}%
\TextAt{8}{79.226}{296.999}{we}%
\TextAt{8}{96.454}{296.999}{consider}%
\TextAt{8}{142.039}{296.999}{the}%
\TextAt{8}{162.195}{296.999}{continuous}%
\TextAt{8}{220.799}{296.999}{homomorphism}%
\TextAt{13}{169.795}{261.484}{g}%
\TextAt{21}{175.401}{259.69}{n}%
\TextAt{9}{184.358}{261.484}{\Character{58}}%
\TextAt{18}{190.93}{261.484}{Z}%
\TextAt{21}{198.9}{266.42}{f}%
\TextAt{21}{198.9}{258.528}{p}%
\NativeArrowAt{207.666}{261.484}{11.955}{19.26}{\longrightarrow}%
\TextAt{13}{230.248}{261.484}{U}%
\TextAt{17}{239.447}{266.42}{\Character{40}}%
\TextAt{21}{242.74}{266.42}{n}%
\TextAt{17}{247.879}{266.42}{\Character{41}}%
\TextAt{13}{251.67}{261.484}{\Character{59}}%
\TextAt{13}{280.805}{261.484}{g}%
\TextAt{21}{286.43}{259.69}{n}%
\TextAt{9}{292.066}{261.484}{\Character{40}}%
\TextAt{13}{296.619}{261.484}{a}%
\TextAt{17}{302.763}{259.69}{1}%
\TextAt{13}{307.496}{261.484}{\Character{59}}%
\TextAt{13}{312.72}{261.484}{\Character{58}}%
\TextAt{13}{317.957}{261.484}{\Character{58}}%
\TextAt{13}{323.181}{261.484}{\Character{58}}%
\TextAt{13}{328.418}{261.484}{\Character{59}}%
\TextAt{13}{333.654}{261.484}{a}%
\TextAt{21}{339.862}{259.69}{f}%
\TextAt{9}{345.306}{261.484}{\Character{41}}%
\TextAt{9}{353.161}{261.484}{\Character{61}}%
\TextAt{21}{370.77}{276.915}{f}%
\TextAt{20}{365.605}{272.841}{\Character{89}}%
\TextAt{21}{366.391}{247.647}{i}%
\TextAt{17}{369.274}{247.647}{\Character{61}1}%
\TextAt{9}{380.881}{261.484}{\Character{40}1}%
\TextAt{9}{393.924}{261.484}{\Character{43}}%
\TextAt{13}{405.704}{261.484}{\Character{33}}%
\TextAt{21}{412.994}{259.69}{i}%
\TextAt{13}{416.375}{261.484}{t}%
\TextAt{21}{420.602}{266.42}{n}%
\TextAt{9}{426.238}{261.484}{\Character{41}}%
\TextAt{21}{430.791}{266.42}{a}%
\TextAt{22}{435.289}{265.205}{i}%
\TextAt{13}{438.949}{261.484}{\Character{58}}%
\TextAt{8}{79.226}{226.999}{This}%
\TextAt{8}{105.946}{226.999}{function}%
\TextAt{8}{151.794}{226.999}{has}%
\TextAt{8}{172.668}{226.999}{the}%
\TextAt{8}{192.824}{226.999}{following}%
\TextAt{8}{242.259}{226.999}{properties.}%
\TextAt{8}{301.827}{226.999}{If}%
\TextAt{13}{313.52}{226.999}{m}%
\TextAt{9}{327.08}{226.999}{\Character{61}}%
\TextAt{13}{339.505}{226.999}{np}%
\TextAt{21}{352.368}{231.338}{s}%
\TextAt{8}{356.782}{226.999}{,}%
\TextAt{13}{363.931}{226.999}{s}%
\NativeMathAt{372.766}{226.999}{11.955}{\geq}%
\TextAt{9}{385.385}{226.999}{0}%
\TextAt{8}{391.238}{226.999}{,}%
\TextAt{8}{398.388}{226.999}{then}%
\TextAt{13}{246.915}{199.915}{U}%
\TextAt{17}{256.114}{204.851}{\Character{40}}%
\TextAt{21}{259.407}{204.851}{m}%
\TextAt{17}{266.898}{204.851}{\Character{41}}%
\TextAt{9}{274.01}{199.915}{\Character{61}}%
\TextAt{13}{286.435}{199.915}{g}%
\TextAt{21}{292.041}{198.121}{n}%
\TextAt{9}{297.677}{199.915}{\Character{40}}%
\TextAt{13}{302.229}{199.915}{p}%
\TextAt{21}{308.104}{204.851}{s}%
\TextAt{18}{312.519}{199.915}{Z}%
\TextAt{21}{320.489}{204.851}{f}%
\TextAt{21}{320.489}{196.959}{p}%
\TextAt{9}{325.934}{199.915}{\Character{41}}%
\TextAt{13}{330.486}{199.915}{U}%
\TextAt{17}{339.685}{204.851}{\Character{40}}%
\TextAt{21}{342.978}{204.851}{m}%
\TextAt{17}{350.469}{204.851}{\Character{43}1\Character{41}}%
\TextAt{8}{517.802}{199.915}{(1)}%
\TextAt{8}{79.226}{172.83}{and}%
\TextAt{8}{101.989}{172.83}{for}%
\TextAt{13}{119.874}{172.83}{\Character{11}}%
\TextAt{9}{130.716}{172.83}{\Character{61}}%
\TextAt{9}{143.126}{172.83}{\Character{40}}%
\TextAt{13}{147.694}{172.83}{a}%
\TextAt{17}{153.839}{171.036}{1}%
\TextAt{13}{158.571}{172.83}{\Character{59}}%
\TextAt{13}{163.795}{172.83}{\Character{58}}%
\TextAt{13}{169.032}{172.83}{\Character{58}}%
\TextAt{13}{174.268}{172.83}{\Character{58}}%
\TextAt{13}{179.493}{172.83}{\Character{59}}%
\TextAt{13}{184.729}{172.83}{a}%
\TextAt{21}{190.937}{171.036}{f}%
\TextAt{9}{196.382}{172.83}{\Character{41}}%
\NativeMathAt{204.255}{172.83}{11.955}{\in}%
\TextAt{18}{215.546}{172.83}{Z}%
\TextAt{21}{223.516}{177.168}{f}%
\TextAt{21}{223.516}{169.874}{p}%
\TextAt{8}{228.961}{172.83}{,}%
\TextAt{13}{213.743}{145.745}{\Character{11}}%
\NativeMathAt{224.586}{145.745}{11.955}{\notin}%
\TextAt{13}{235.877}{145.745}{p}%
\TextAt{18}{241.752}{145.745}{Z}%
\TextAt{21}{249.722}{150.681}{f}%
\TextAt{21}{249.722}{142.789}{p}%
\NativeArrowAt{270.443}{145.745}{11.955}{21.926}{\Longleftrightarrow}%
\TextAt{13}{307.637}{145.745}{g}%
\TextAt{21}{313.242}{143.951}{n}%
\TextAt{9}{318.878}{145.745}{\Character{40}}%
\TextAt{13}{323.431}{145.745}{p}%
\TextAt{21}{329.306}{150.681}{s}%
\TextAt{13}{333.72}{145.745}{\Character{11}}%
\TextAt{9}{341.242}{145.745}{\Character{41}}%
\NativeMathAt{349.115}{145.745}{11.955}{\notin}%
\TextAt{13}{360.406}{145.745}{U}%
\TextAt{17}{369.605}{150.681}{\Character{40}}%
\TextAt{21}{372.898}{150.681}{m}%
\TextAt{17}{380.389}{150.681}{\Character{43}1\Character{41}}%
\TextAt{13}{395.001}{145.745}{\Character{58}}%
\TextAt{8}{517.802}{145.745}{(2)}%
\TextAt{8}{79.226}{118.66}{Indeed,}%
\TextAt{8}{120.507}{118.66}{for}%
\TextAt{13}{138.392}{118.66}{\Character{33}}%
\TextAt{9}{149.431}{118.66}{\Character{61}}%
\TextAt{20}{161.857}{127.626}{\Character{80}}%
\TextAt{21}{174.476}{124.471}{f}%
\TextAt{21}{174.476}{115.173}{i}%
\TextAt{17}{177.359}{115.173}{\Character{61}1}%
\TextAt{13}{190.671}{118.66}{b}%
\TextAt{21}{195.648}{116.867}{i}%
\TextAt{13}{199.029}{118.66}{\Character{33}}%
\TextAt{21}{206.318}{116.867}{i}%
\NativeMathAt{213.02}{118.66}{11.955}{\in}%
\TextAt{18}{224.311}{118.66}{F}%
\TextAt{21}{231.617}{116.867}{q}%
\TextAt{8}{236.182}{118.66}{,}%
\TextAt{13}{243.331}{118.66}{b}%
\TextAt{21}{248.308}{116.867}{i}%
\NativeMathAt{255.011}{118.66}{11.955}{\in}%
\TextAt{18}{266.302}{118.66}{Z}%
\TextAt{8}{274.272}{118.66}{,}%
\TextAt{13}{281.421}{118.66}{b}%
\TextAt{21}{286.398}{116.867}{i}%
\NativeMathAt{293.1}{118.66}{11.955}{\equiv}%
\TextAt{13}{305.719}{118.66}{a}%
\TextAt{21}{311.864}{116.867}{i}%
\TextAt{8}{318.566}{118.66}{mod}%
\TextAt{13}{344.004}{118.66}{p}%
\TextAt{8}{349.88}{118.66}{,}%
\TextAt{8}{357.029}{118.66}{we}%
\TextAt{8}{374.256}{118.66}{have}%
\TextAt{13}{195.364}{81.976}{g}%
\TextAt{21}{200.969}{80.183}{n}%
\TextAt{9}{206.605}{81.976}{\Character{40}}%
\TextAt{13}{211.158}{81.976}{\Character{11}}%
\TextAt{9}{218.68}{81.976}{\Character{41}}%
\NativeMathAt{226.553}{81.976}{11.955}{\equiv}%
\TextAt{21}{244.337}{97.407}{f}%
\TextAt{20}{239.172}{93.334}{\Character{89}}%
\TextAt{21}{239.958}{68.139}{i}%
\TextAt{17}{242.841}{68.139}{\Character{61}1}%
\TextAt{9}{254.448}{81.976}{\Character{40}1}%
\TextAt{9}{267.491}{81.976}{\Character{43}}%
\TextAt{13}{279.271}{81.976}{\Character{33}}%
\TextAt{21}{286.561}{80.183}{i}%
\TextAt{13}{289.942}{81.976}{t}%
\TextAt{21}{294.169}{86.912}{n}%
\TextAt{9}{299.805}{81.976}{\Character{41}}%
\TextAt{21}{304.358}{86.912}{b}%
\TextAt{22}{307.98}{85.697}{i}%
\NativeMathAt{314.961}{81.976}{11.955}{\equiv}%
\TextAt{9}{327.58}{81.976}{1}%
\TextAt{9}{336.08}{81.976}{\Character{43}}%
\TextAt{13}{347.851}{81.976}{\Character{33}}%
\TextAt{13}{355.562}{81.976}{t}%
\TextAt{21}{359.797}{86.912}{n}%
\TextAt{8}{368.754}{81.976}{mod}%
\TextAt{14}{394.192}{81.976}{p}%
\TextAt{21}{400.175}{86.912}{n}%
\TextAt{17}{405.313}{86.912}{\Character{43}1}%
\TextAt{8}{297.211}{43.033}{128}%
\Destination{522.36}{216.965}{nav.AMS.17}%
\Destination{522.36}{162.795}{nav.AMS.18}%
\Destination{79.23}{773.075}{nav.Item.43}%
\Destination{78.23}{809.815}{nav.page.128}%
\Destination{79.36}{776.155}{nav.print.140}%
\Hotlink{218.493}{545.455}{5.22}{11.047}{nav.chapter.1}%
\Hotlink{344.977}{386.055}{18.574}{11.047}{nav.print.127}%
\Hotlink{324.955}{331.855}{19.22}{11.047}{nav.cite.N079}%
\Hotlink{133.863}{690.855}{25.078}{11.047}{nav.number.2.5.3}%
\Hotlink{341.576}{616.655}{25.078}{11.047}{nav.number.2.5.5}%
\Hotlink{231.544}{545.455}{25.078}{11.047}{nav.number.1.2.9}%
}

% Page 141 (body page 129)
\Page{%
\TextAt{2}{79.226}{796.943}{II.}%
\TextAt{2}{94.073}{796.943}{The}%
\TextAt{2}{116.492}{796.943}{Theory}%
\TextAt{2}{154.401}{796.943}{of}%
\TextAt{2}{166.826}{796.943}{Valuations}%
\Vector{q 0 G 0 g 0.299 w 7.226 18.729 m 460.769 18.729 l S Q}%
\TextAt{8}{79.226}{761.124}{and}%
\TextAt{8}{101.989}{761.124}{hence,}%
\TextAt{8}{137.747}{761.124}{since}%
\TextAt{8}{166.415}{761.124}{we}%
\TextAt{8}{183.643}{761.124}{are}%
\TextAt{8}{203.153}{761.124}{in}%
\TextAt{8}{216.806}{761.124}{characteristic}%
\TextAt{13}{288.74}{761.124}{p}%
\TextAt{8}{294.616}{761.124}{,}%
\TextAt{13}{199.885}{732.417}{g}%
\TextAt{21}{205.491}{730.623}{n}%
\TextAt{9}{211.127}{732.417}{\Character{40}}%
\TextAt{13}{215.679}{732.417}{p}%
\TextAt{21}{221.554}{737.353}{s}%
\TextAt{13}{225.968}{732.417}{\Character{11}}%
\TextAt{9}{233.49}{732.417}{\Character{41}}%
\TextAt{9}{241.345}{732.417}{\Character{61}}%
\TextAt{13}{253.789}{732.417}{g}%
\TextAt{21}{259.394}{730.623}{n}%
\TextAt{9}{265.03}{732.417}{\Character{40}}%
\TextAt{13}{269.583}{732.417}{\Character{11}}%
\TextAt{9}{277.104}{732.417}{\Character{41}}%
\TextAt{21}{281.657}{737.353}{p}%
\TextAt{22}{285.919}{740.166}{s}%
\NativeMathAt{293.73}{732.417}{11.955}{\equiv}%
\TextAt{9}{306.35}{732.417}{1}%
\TextAt{9}{314.85}{732.417}{\Character{43}}%
\TextAt{13}{326.621}{732.417}{\Character{33}}%
\TextAt{21}{334.339}{737.353}{p}%
\TextAt{22}{338.602}{740.166}{s}%
\TextAt{13}{343.092}{732.417}{t}%
\TextAt{21}{347.319}{737.353}{m}%
\TextAt{8}{358.628}{732.417}{mod}%
\TextAt{14}{384.066}{732.417}{p}%
\TextAt{21}{390.049}{737.353}{m}%
\TextAt{17}{397.54}{737.353}{\Character{43}1}%
\TextAt{13}{408.859}{732.417}{\Character{58}}%
\TextAt{8}{79.226}{703.71}{As}%
\TextAt{13}{96.517}{703.71}{\Character{11}}%
\TextAt{8}{107.939}{703.71}{varies}%
\TextAt{8}{140.855}{703.71}{over}%
\TextAt{8}{165.905}{703.71}{the}%
\TextAt{8}{186.065}{703.71}{elements}%
\TextAt{8}{233.925}{703.71}{of}%
\TextAt{18}{247.258}{703.71}{Z}%
\TextAt{21}{255.229}{708.048}{f}%
\TextAt{21}{255.229}{700.754}{p}%
\TextAt{8}{260.673}{703.71}{,}%
\TextAt{13}{267.827}{703.71}{\Character{33}}%
\TextAt{8}{275.545}{703.71}{,}%
\TextAt{8}{282.699}{703.71}{and}%
\TextAt{8}{305.465}{703.71}{thus}%
\TextAt{8}{331.22}{703.71}{also}%
\TextAt{13}{354.703}{703.71}{\Character{33}}%
\TextAt{21}{362.421}{708.048}{p}%
\TextAt{22}{366.684}{710.861}{s}%
\TextAt{8}{371.174}{703.71}{,}%
\TextAt{8}{378.328}{703.71}{varies}%
\TextAt{8}{411.244}{703.71}{over}%
\TextAt{8}{436.293}{703.71}{the}%
\TextAt{8}{456.453}{703.71}{elements}%
\TextAt{8}{504.313}{703.71}{of}%
\TextAt{18}{517.647}{703.71}{F}%
\TextAt{21}{524.953}{701.917}{q}%
\TextAt{8}{529.518}{703.71}{,}%
\TextAt{8}{79.226}{686.664}{and}%
\TextAt{8}{100.808}{686.664}{we}%
\TextAt{8}{116.854}{686.664}{get}%
\TextAt{8}{135.184}{686.664}{(1).}%
\TextAt{8}{158.206}{686.664}{Furthermore}%
\TextAt{8}{224.823}{686.664}{one}%
\TextAt{8}{245.102}{686.664}{has}%
\TextAt{13}{264.794}{686.664}{g}%
\TextAt{21}{270.399}{684.871}{n}%
\TextAt{9}{276.036}{686.664}{\Character{40}}%
\TextAt{13}{280.588}{686.664}{p}%
\TextAt{21}{286.463}{691.002}{s}%
\TextAt{13}{290.877}{686.664}{\Character{11}}%
\TextAt{9}{298.399}{686.664}{\Character{41}}%
\NativeMathAt{306.272}{686.664}{11.955}{\equiv}%
\TextAt{9}{318.892}{686.664}{1}%
\TextAt{8}{328.065}{686.664}{mod}%
\TextAt{14}{353.503}{686.664}{p}%
\TextAt{21}{359.486}{691.002}{m}%
\TextAt{17}{366.977}{691.002}{\Character{43}1}%
\NativeArrowAt{381.617}{686.664}{11.955}{21.926}{\Longleftrightarrow}%
\TextAt{13}{406.855}{686.664}{\Character{33}}%
\TextAt{9}{417.894}{686.664}{\Character{61}}%
\TextAt{9}{430.303}{686.664}{0}%
\NativeArrowAt{439.494}{686.664}{11.955}{21.926}{\Longleftrightarrow}%
\TextAt{13}{464.732}{686.664}{b}%
\TextAt{21}{469.709}{684.871}{i}%
\NativeMathAt{476.412}{686.664}{11.955}{\equiv}%
\TextAt{9}{489.031}{686.664}{0}%
\TextAt{8}{498.205}{686.664}{mod}%
\TextAt{13}{523.643}{686.664}{p}%
\TextAt{8}{529.518}{686.664}{,}%
\TextAt{8}{79.226}{669.618}{for}%
\TextAt{13}{97.495}{669.618}{i}%
\TextAt{9}{105.463}{669.618}{\Character{61}}%
\TextAt{9}{118.53}{669.618}{1}%
\TextAt{13}{124.396}{669.618}{\Character{59}}%
\TextAt{13}{129.62}{669.618}{\Character{58}}%
\TextAt{13}{134.857}{669.618}{\Character{58}}%
\TextAt{13}{140.093}{669.618}{\Character{58}}%
\TextAt{13}{145.318}{669.618}{\Character{59}}%
\TextAt{13}{150.554}{669.618}{f}%
\TextAt{8}{157.663}{669.618}{,}%
\NativeArrowAt{165.293}{669.618}{11.955}{21.926}{\Longleftrightarrow}%
\TextAt{13}{191.185}{669.618}{a}%
\TextAt{21}{197.33}{667.825}{i}%
\NativeMathAt{204.687}{669.618}{11.955}{\equiv}%
\TextAt{9}{217.96}{669.618}{0}%
\TextAt{8}{227.134}{669.618}{mod}%
\TextAt{13}{252.572}{669.618}{p}%
\TextAt{8}{258.447}{669.618}{,}%
\TextAt{8}{266.076}{669.618}{for}%
\TextAt{13}{284.345}{669.618}{i}%
\TextAt{9}{292.314}{669.618}{\Character{61}}%
\TextAt{9}{305.381}{669.618}{1}%
\TextAt{13}{311.246}{669.618}{\Character{59}}%
\TextAt{13}{316.47}{669.618}{\Character{58}}%
\TextAt{13}{321.707}{669.618}{\Character{58}}%
\TextAt{13}{326.943}{669.618}{\Character{58}}%
\TextAt{13}{332.168}{669.618}{\Character{59}}%
\TextAt{13}{337.404}{669.618}{f}%
\TextAt{8}{344.514}{669.618}{,}%
\NativeArrowAt{352.143}{669.618}{11.955}{21.926}{\Longleftrightarrow}%
\TextAt{13}{378.036}{669.618}{\Character{11}}%
\NativeMathAt{389.532}{669.618}{11.955}{\in}%
\TextAt{13}{401.478}{669.618}{p}%
\TextAt{18}{407.353}{669.618}{Z}%
\TextAt{21}{415.323}{673.956}{f}%
\TextAt{21}{415.323}{666.662}{p}%
\TextAt{8}{420.768}{669.618}{,}%
\TextAt{8}{428.397}{669.618}{and}%
\TextAt{8}{451.543}{669.618}{this}%
\TextAt{8}{474.75}{669.618}{amounts}%
\TextAt{8}{522.357}{669.618}{to}%
\TextAt{8}{79.226}{652.572}{(2).}%
\TextAt{8}{97.159}{632.537}{We}%
\TextAt{8}{117.303}{632.537}{now}%
\TextAt{8}{141.692}{632.537}{consider}%
\TextAt{8}{187.277}{632.537}{the}%
\TextAt{8}{207.433}{632.537}{continuous}%
\TextAt{8}{266.038}{632.537}{homomorphism}%
\TextAt{8}{348.708}{632.537}{of}%
\TextAt{18}{362.038}{632.537}{Z}%
\TextAt{21}{370.008}{630.744}{p}%
\TextAt{8}{374.769}{632.537}{-modules}%
\TextAt{13}{209.583}{603.83}{g}%
\TextAt{9}{218.938}{603.83}{\Character{61}}%
\TextAt{20}{238.305}{615.188}{\Character{89}}%
\TextAt{17}{231.363}{589.285}{\Character{40}}%
\TextAt{21}{234.656}{589.285}{n\Character{59}p}%
\TextAt{17}{246.41}{589.285}{\Character{41}\Character{61}1}%
\TextAt{13}{262.516}{603.83}{g}%
\TextAt{21}{268.121}{602.037}{n}%
\TextAt{9}{277.079}{603.83}{\Character{58}}%
\TextAt{13}{283.651}{603.83}{A}%
\TextAt{9}{295.747}{603.83}{\Character{61}}%
\TextAt{20}{315.115}{615.188}{\Character{89}}%
\TextAt{17}{308.173}{589.285}{\Character{40}}%
\TextAt{21}{311.466}{589.285}{n\Character{59}p}%
\TextAt{17}{323.219}{589.285}{\Character{41}\Character{61}1}%
\TextAt{18}{339.326}{603.83}{Z}%
\TextAt{21}{347.296}{608.766}{f}%
\TextAt{21}{347.296}{600.875}{p}%
\NativeArrowAt{356.062}{603.83}{11.955}{19.26}{\longrightarrow}%
\TextAt{13}{378.644}{603.83}{U}%
\TextAt{17}{387.843}{608.766}{\Character{40}1\Character{41}}%
\TextAt{13}{399.161}{603.83}{\Character{59}}%
\TextAt{8}{79.226}{563.566}{where}%
\TextAt{8}{111.775}{563.566}{the}%
\TextAt{8}{130.67}{563.566}{product}%
\TextAt{20}{173.309}{572.533}{\Character{81}}%
\TextAt{17}{184.6}{560.079}{\Character{40}}%
\TextAt{21}{187.893}{560.079}{n\Character{59}p}%
\TextAt{17}{199.647}{560.079}{\Character{41}\Character{61}1}%
\TextAt{18}{216.251}{563.566}{Z}%
\TextAt{21}{224.221}{567.905}{f}%
\TextAt{21}{224.221}{560.611}{p}%
\TextAt{8}{232.303}{563.566}{is}%
\TextAt{8}{242.806}{563.566}{taken}%
\TextAt{8}{273.418}{563.566}{over}%
\TextAt{8}{297.203}{563.566}{all}%
\TextAt{13}{312.202}{563.566}{n}%
\NativeMathAt{322.51}{563.566}{11.955}{\geq}%
\TextAt{9}{335.129}{563.566}{1}%
\TextAt{8}{343.619}{563.566}{such}%
\TextAt{8}{368.755}{563.566}{that}%
\TextAt{9}{392.863}{563.566}{\Character{40}}%
\TextAt{13}{397.416}{563.566}{n\Character{59}}%
\TextAt{13}{409.622}{563.566}{p}%
\TextAt{9}{415.523}{563.566}{\Character{41}}%
\TextAt{9}{423.378}{563.566}{\Character{61}}%
\TextAt{9}{435.799}{563.566}{1}%
\TextAt{8}{441.674}{563.566}{,}%
\TextAt{8}{447.815}{563.566}{each}%
\TextAt{8}{472.891}{563.566}{factor}%
\TextAt{8}{505.129}{563.566}{being}%
\TextAt{8}{79.226}{546.52}{a}%
\TextAt{8}{88.139}{546.52}{copy}%
\TextAt{8}{114.613}{546.52}{of}%
\TextAt{18}{127.1}{546.52}{Z}%
\TextAt{21}{135.07}{550.859}{f}%
\TextAt{21}{135.07}{543.565}{p}%
\TextAt{8}{140.515}{546.52}{.}%
\TextAt{8}{148.683}{546.52}{Observe}%
\TextAt{8}{192.779}{546.52}{that}%
\TextAt{8}{217.305}{546.52}{the}%
\TextAt{8}{236.618}{546.52}{product}%
\TextAt{13}{279.675}{546.52}{g}%
\TextAt{9}{285.709}{546.52}{\Character{40}}%
\TextAt{13}{290.261}{546.52}{\Character{24}}%
\TextAt{9}{295.952}{546.52}{\Character{41}}%
\TextAt{9}{303.807}{546.52}{\Character{61}}%
\TextAt{20}{316.251}{555.487}{\Character{81}}%
\TextAt{13}{329.534}{546.52}{g}%
\TextAt{21}{335.139}{544.727}{n}%
\TextAt{9}{340.776}{546.52}{\Character{40}}%
\TextAt{13}{345.328}{546.52}{\Character{11}}%
\TextAt{21}{352.774}{544.727}{n}%
\TextAt{9}{358.41}{546.52}{\Character{41}}%
\TextAt{8}{366.017}{546.52}{converges}%
\TextAt{8}{417.597}{546.52}{because}%
\TextAt{13}{460.056}{546.52}{g}%
\TextAt{21}{465.661}{544.727}{n}%
\TextAt{9}{471.297}{546.52}{\Character{40}}%
\TextAt{13}{475.85}{546.52}{\Character{11}}%
\TextAt{21}{483.295}{544.727}{n}%
\TextAt{9}{488.932}{546.52}{\Character{41}}%
\NativeMathAt{496.805}{546.52}{11.955}{\in}%
\TextAt{13}{508.096}{546.52}{U}%
\TextAt{17}{517.295}{550.859}{\Character{40}}%
\TextAt{21}{520.588}{550.859}{n}%
\TextAt{17}{525.726}{550.859}{\Character{41}}%
\TextAt{8}{529.518}{546.52}{.}%
\TextAt{8}{79.226}{529.474}{Let}%
\TextAt{13}{100.253}{529.474}{m}%
\TextAt{9}{113.911}{529.474}{\Character{61}}%
\TextAt{13}{126.435}{529.474}{np}%
\TextAt{21}{139.297}{533.813}{s}%
\TextAt{8}{143.711}{529.474}{,}%
\TextAt{8}{150.933}{529.474}{with}%
\TextAt{9}{177.65}{529.474}{\Character{40}}%
\TextAt{13}{182.202}{529.474}{n\Character{59}}%
\TextAt{13}{194.408}{529.474}{p}%
\TextAt{9}{200.309}{529.474}{\Character{41}}%
\TextAt{9}{208.259}{529.474}{\Character{61}}%
\TextAt{9}{220.776}{529.474}{1}%
\TextAt{8}{226.657}{529.474}{,}%
\TextAt{8}{233.878}{529.474}{be}%
\TextAt{8}{249.86}{529.474}{any}%
\TextAt{8}{272.034}{529.474}{natural}%
\TextAt{8}{313.074}{529.474}{number.}%
\TextAt{8}{360.394}{529.474}{As}%
\TextAt{13}{377.739}{529.474}{g}%
\TextAt{21}{383.344}{527.681}{n}%
\TextAt{9}{388.98}{529.474}{\Character{40}}%
\TextAt{18}{393.533}{529.474}{Z}%
\TextAt{21}{401.503}{533.813}{f}%
\TextAt{21}{401.503}{526.519}{p}%
\TextAt{9}{406.948}{529.474}{\Character{41}}%
\NativeMathAt{414.919}{529.474}{11.955}{\subseteq}%
\TextAt{13}{427.636}{529.474}{g}%
\TextAt{9}{433.67}{529.474}{\Character{40}}%
\TextAt{13}{438.223}{529.474}{A}%
\TextAt{9}{446.998}{529.474}{\Character{41}}%
\TextAt{8}{451.55}{529.474}{,}%
\TextAt{8}{458.772}{529.474}{it}%
\TextAt{8}{470.533}{529.474}{follows}%
\TextAt{8}{509.027}{529.474}{from}%
\TextAt{8}{79.226}{512.428}{(1)}%
\TextAt{8}{97.589}{512.428}{that}%
\TextAt{8}{122.455}{512.428}{each}%
\TextAt{8}{148.289}{512.428}{coset}%
\TextAt{8}{177.112}{512.428}{of}%
\TextAt{13}{189.939}{512.428}{U}%
\TextAt{17}{199.138}{516.767}{\Character{40}}%
\TextAt{21}{202.432}{516.767}{m}%
\TextAt{17}{209.922}{516.767}{\Character{41}}%
\TextAt{9}{213.713}{512.428}{\Character{47}}%
\TextAt{13}{219.566}{512.428}{U}%
\TextAt{17}{228.765}{516.767}{\Character{40}}%
\TextAt{21}{232.059}{516.767}{m}%
\TextAt{17}{239.549}{516.767}{\Character{43}1\Character{41}}%
\TextAt{8}{257.556}{512.428}{is}%
\TextAt{8}{268.817}{512.428}{represented}%
\TextAt{8}{330.481}{512.428}{by}%
\TextAt{8}{346.237}{512.428}{an}%
\TextAt{8}{361.993}{512.428}{element}%
\TextAt{8}{404.732}{512.428}{of}%
\TextAt{13}{417.559}{512.428}{g}%
\TextAt{9}{423.594}{512.428}{\Character{40}}%
\TextAt{13}{428.146}{512.428}{A}%
\TextAt{9}{436.921}{512.428}{\Character{41}}%
\TextAt{8}{441.474}{512.428}{.}%
\TextAt{8}{449.754}{512.428}{This}%
\TextAt{8}{475.971}{512.428}{means}%
\TextAt{8}{511.298}{512.428}{that}%
\TextAt{13}{79.226}{495.382}{g}%
\TextAt{9}{85.26}{495.382}{\Character{40}}%
\TextAt{13}{89.813}{495.382}{A}%
\TextAt{9}{98.588}{495.382}{\Character{41}}%
\TextAt{8}{107.038}{495.382}{is}%
\TextAt{8}{118.802}{495.382}{dense}%
\TextAt{8}{150.722}{495.382}{in}%
\TextAt{13}{164.375}{495.382}{U}%
\TextAt{17}{173.574}{499.721}{\Character{40}1\Character{41}}%
\TextAt{8}{184.893}{495.382}{.}%
\TextAt{8}{193.341}{495.382}{Since}%
\TextAt{13}{223.898}{495.382}{A}%
\TextAt{8}{236.571}{495.382}{is}%
\TextAt{8}{248.335}{495.382}{compact}%
\TextAt{8}{295.163}{495.382}{and}%
\TextAt{13}{317.926}{495.382}{g}%
\TextAt{8}{327.858}{495.382}{is}%
\TextAt{8}{339.622}{495.382}{continuous,}%
\TextAt{13}{401.478}{495.382}{g}%
\TextAt{8}{411.409}{495.382}{is}%
\TextAt{8}{423.173}{495.382}{actually}%
\TextAt{8}{467.73}{495.382}{surjective.}%
\TextAt{8}{97.159}{475.347}{On}%
\TextAt{8}{116.785}{475.347}{the}%
\TextAt{8}{137.057}{475.347}{other}%
\TextAt{8}{167.741}{475.347}{hand,}%
\TextAt{8}{200.403}{475.347}{let}%
\TextAt{13}{217.423}{475.347}{\Character{24}}%
\TextAt{9}{226.631}{475.347}{\Character{61}}%
\TextAt{9}{239.244}{475.347}{\Character{40}}%
\TextAt{13}{243.805}{475.347}{\Character{58}}%
\TextAt{13}{249.029}{475.347}{\Character{58}}%
\TextAt{13}{254.266}{475.347}{\Character{58}}%
\TextAt{13}{259.49}{475.347}{\Character{59}}%
\TextAt{13}{264.727}{475.347}{\Character{11}}%
\TextAt{21}{272.227}{473.554}{n}%
\TextAt{13}{277.863}{475.347}{\Character{59}}%
\TextAt{13}{283.087}{475.347}{\Character{58}}%
\TextAt{13}{288.324}{475.347}{\Character{58}}%
\TextAt{13}{293.56}{475.347}{\Character{58}}%
\TextAt{9}{296.847}{475.347}{\Character{41}}%
\NativeMathAt{304.917}{475.347}{11.955}{\in}%
\TextAt{13}{316.404}{475.347}{A}%
\TextAt{8}{325.18}{475.347}{,}%
\TextAt{13}{332.473}{475.347}{\Character{24}}%
\NativeMathAt{341.68}{475.347}{11.955}{\neq}%
\TextAt{9}{354.293}{475.347}{0}%
\TextAt{8}{360.155}{475.347}{,}%
\TextAt{8}{367.448}{475.347}{i.e.,}%
\TextAt{13}{389.697}{475.347}{\Character{11}}%
\TextAt{21}{397.143}{473.554}{n}%
\NativeMathAt{406.296}{475.347}{11.955}{\neq}%
\TextAt{9}{418.909}{475.347}{0}%
\TextAt{8}{428.784}{475.347}{for}%
\TextAt{8}{446.784}{475.347}{some}%
\TextAt{13}{476.225}{475.347}{n}%
\TextAt{8}{483.213}{475.347}{.}%
\TextAt{8}{492.007}{475.347}{Such}%
\TextAt{8}{520.408}{475.347}{an}%
\TextAt{13}{79.226}{458.301}{\Character{11}}%
\TextAt{21}{86.672}{456.508}{n}%
\TextAt{8}{96.158}{458.301}{is}%
\TextAt{8}{107.874}{458.301}{of}%
\TextAt{8}{121.156}{458.301}{the}%
\TextAt{8}{141.265}{458.301}{form}%
\TextAt{13}{168.858}{458.301}{\Character{11}}%
\TextAt{21}{176.303}{456.508}{n}%
\TextAt{9}{185.26}{458.301}{\Character{61}}%
\TextAt{13}{197.686}{458.301}{p}%
\TextAt{21}{203.561}{462.64}{s}%
\TextAt{13}{207.975}{458.301}{\Character{12}}%
\TextAt{21}{214.584}{456.508}{n}%
\TextAt{8}{224.07}{458.301}{with}%
\TextAt{13}{250.682}{458.301}{s}%
\TextAt{9}{259.517}{458.301}{\Character{61}}%
\TextAt{13}{271.942}{458.301}{s}%
\TextAt{9}{277.456}{458.301}{\Character{40}}%
\TextAt{13}{282.008}{458.301}{\Character{11}}%
\TextAt{21}{289.454}{456.508}{n}%
\TextAt{9}{295.09}{458.301}{\Character{41}}%
\NativeMathAt{302.964}{458.301}{11.955}{\geq}%
\TextAt{9}{315.583}{458.301}{0}%
\TextAt{8}{321.436}{458.301}{,}%
\TextAt{8}{328.547}{458.301}{and}%
\TextAt{13}{351.262}{458.301}{\Character{12}}%
\TextAt{21}{357.87}{456.508}{n}%
\NativeMathAt{366.828}{458.301}{11.955}{\in}%
\TextAt{18}{378.118}{458.301}{Z}%
\TextAt{21}{386.089}{462.64}{f}%
\TextAt{21}{386.089}{455.346}{p}%
\TextAt{18}{394.093}{458.301}{r}%
\TextAt{13}{405.95}{458.301}{p}%
\TextAt{18}{411.825}{458.301}{Z}%
\TextAt{21}{419.796}{462.64}{f}%
\TextAt{21}{419.796}{455.346}{p}%
\TextAt{8}{425.24}{458.301}{.}%
\TextAt{8}{433.673}{458.301}{It}%
\TextAt{8}{446.298}{458.301}{now}%
\TextAt{8}{470.638}{458.301}{follows}%
\TextAt{8}{509.027}{458.301}{from}%
\TextAt{8}{79.226}{441.255}{(2)}%
\TextAt{8}{98.091}{441.255}{that}%
\TextAt{13}{151.933}{424.21}{g}%
\TextAt{21}{157.538}{422.416}{n}%
\TextAt{9}{163.174}{424.21}{\Character{40}}%
\TextAt{13}{167.727}{424.21}{\Character{11}}%
\TextAt{21}{175.172}{422.416}{n}%
\TextAt{9}{180.809}{424.21}{\Character{41}}%
\NativeMathAt{188.682}{424.21}{11.955}{\in}%
\TextAt{13}{199.973}{424.21}{U}%
\TextAt{17}{209.172}{429.146}{\Character{40}}%
\TextAt{21}{212.465}{429.146}{m}%
\TextAt{17}{219.956}{429.146}{\Character{41}}%
\TextAt{13}{223.747}{424.21}{\Character{59}}%
\TextAt{13}{252.882}{424.21}{g}%
\TextAt{21}{258.507}{422.416}{n}%
\TextAt{9}{264.143}{424.21}{\Character{40}}%
\TextAt{13}{268.696}{424.21}{\Character{11}}%
\TextAt{21}{276.141}{422.416}{n}%
\TextAt{9}{281.778}{424.21}{\Character{41}}%
\NativeMathAt{289.651}{424.21}{11.955}{\notin}%
\TextAt{13}{300.942}{424.21}{U}%
\TextAt{17}{310.141}{429.146}{\Character{40}}%
\TextAt{21}{313.434}{429.146}{m}%
\TextAt{17}{320.924}{429.146}{\Character{43}1\Character{41}}%
\TextAt{8}{347.492}{424.21}{for}%
\TextAt{13}{365.377}{424.21}{m}%
\TextAt{9}{378.937}{424.21}{\Character{61}}%
\TextAt{13}{391.362}{424.21}{m}%
\TextAt{9}{401.601}{424.21}{\Character{40}}%
\TextAt{13}{406.154}{424.21}{\Character{11}}%
\TextAt{21}{413.599}{422.416}{n}%
\TextAt{9}{419.236}{424.21}{\Character{41}}%
\TextAt{9}{427.091}{424.21}{\Character{61}}%
\TextAt{13}{439.534}{424.21}{np}%
\TextAt{21}{452.397}{429.146}{s}%
\TextAt{13}{456.811}{424.21}{\Character{58}}%
\TextAt{8}{79.226}{400.814}{Since}%
\TextAt{8}{109.523}{400.814}{the}%
\TextAt{13}{129.419}{400.814}{n}%
\TextAt{8}{140.044}{400.814}{are}%
\TextAt{8}{159.294}{400.814}{prime}%
\TextAt{8}{192.198}{400.814}{to}%
\TextAt{13}{206.248}{400.814}{p}%
\TextAt{8}{212.123}{400.814}{,}%
\TextAt{8}{219.064}{400.814}{all}%
\TextAt{8}{235.063}{400.814}{the}%
\TextAt{13}{254.959}{400.814}{m}%
\TextAt{9}{265.198}{400.814}{\Character{40}}%
\TextAt{13}{269.75}{400.814}{\Character{11}}%
\TextAt{21}{277.196}{399.021}{n}%
\TextAt{9}{282.832}{400.814}{\Character{41}}%
\TextAt{8}{291.022}{400.814}{have}%
\TextAt{8}{317.756}{400.814}{to}%
\TextAt{8}{331.806}{400.814}{be}%
\TextAt{8}{347.47}{400.814}{distinct,}%
\TextAt{8}{392.847}{400.814}{for}%
\TextAt{8}{410.471}{400.814}{all}%
\TextAt{13}{426.47}{400.814}{\Character{11}}%
\TextAt{21}{433.916}{399.021}{n}%
\NativeMathAt{442.873}{400.814}{11.955}{\neq}%
\TextAt{9}{455.283}{400.814}{0}%
\TextAt{8}{461.151}{400.814}{.}%
\TextAt{8}{469.513}{400.814}{Let}%
\TextAt{13}{490.222}{400.814}{n}%
\TextAt{8}{500.846}{400.814}{be}%
\TextAt{8}{516.51}{400.814}{the}%
\TextAt{8}{79.226}{383.768}{natural}%
\TextAt{8}{119.872}{383.768}{number,}%
\TextAt{8}{165.452}{383.768}{prime}%
\TextAt{8}{198.279}{383.768}{to}%
\TextAt{13}{212.253}{383.768}{p}%
\TextAt{8}{221.69}{383.768}{and}%
\TextAt{8}{244.116}{383.768}{such}%
\TextAt{8}{270.177}{383.768}{that}%
\TextAt{13}{295.21}{383.768}{\Character{11}}%
\TextAt{21}{302.656}{381.975}{n}%
\NativeMathAt{311.613}{383.768}{11.955}{\neq}%
\TextAt{9}{324.023}{383.768}{0}%
\TextAt{8}{329.891}{383.768}{,}%
\TextAt{8}{336.772}{383.768}{which}%
\TextAt{8}{369.922}{383.768}{satisfies}%
\TextAt{13}{412.696}{383.768}{m}%
\TextAt{9}{422.936}{383.768}{\Character{40}}%
\TextAt{13}{427.488}{383.768}{\Character{11}}%
\TextAt{21}{434.933}{381.975}{n}%
\TextAt{9}{440.57}{383.768}{\Character{41}}%
\TextAt{13}{448.443}{383.768}{\Character{60}}%
\TextAt{13}{460.853}{383.768}{m}%
\TextAt{9}{471.108}{383.768}{\Character{40}}%
\TextAt{13}{475.66}{383.768}{\Character{11}}%
\TextAt{21}{483.106}{381.975}{n}%
\TextAt{23}{488.244}{384.244}{\Character{48}}%
\TextAt{9}{491.426}{383.768}{\Character{41}}%
\TextAt{8}{495.979}{383.768}{,}%
\TextAt{8}{502.859}{383.768}{for}%
\TextAt{8}{520.408}{383.768}{all}%
\TextAt{13}{79.226}{366.722}{n}%
\TextAt{3}{86.214}{371.06}{\Character{48}}%
\NativeMathAt{92.33}{366.722}{11.955}{\neq}%
\TextAt{13}{104.755}{366.722}{n}%
\TextAt{8}{115.64}{366.722}{such}%
\TextAt{8}{142.037}{366.722}{that}%
\TextAt{13}{167.406}{366.722}{\Character{11}}%
\TextAt{21}{174.852}{364.929}{n}%
\TextAt{23}{179.99}{367.198}{\Character{48}}%
\NativeMathAt{186.493}{366.722}{11.955}{\neq}%
\TextAt{9}{198.903}{366.722}{0}%
\TextAt{8}{204.772}{366.722}{.}%
\TextAt{8}{213.22}{366.722}{Then}%
\TextAt{8}{243.777}{366.722}{one}%
\TextAt{8}{265.237}{366.722}{has,}%
\TextAt{8}{289.362}{366.722}{for}%
\TextAt{8}{307.247}{366.722}{all}%
\TextAt{13}{323.506}{366.722}{n}%
\TextAt{3}{330.494}{371.06}{\Character{48}}%
\NativeMathAt{336.61}{366.722}{11.955}{\neq}%
\TextAt{13}{349.035}{366.722}{n}%
\TextAt{8}{356.023}{366.722}{,}%
\TextAt{8}{363.172}{366.722}{that}%
\TextAt{13}{182.854}{338.015}{g}%
\TextAt{21}{188.459}{336.222}{n}%
\TextAt{23}{193.597}{338.491}{\Character{48}}%
\TextAt{9}{196.78}{338.015}{\Character{40}}%
\TextAt{13}{201.332}{338.015}{\Character{11}}%
\TextAt{21}{208.778}{336.222}{n}%
\TextAt{23}{213.916}{338.491}{\Character{48}}%
\TextAt{9}{217.099}{338.015}{\Character{41}}%
\NativeMathAt{224.972}{338.015}{11.955}{\in}%
\TextAt{13}{236.263}{338.015}{U}%
\TextAt{17}{245.462}{342.951}{\Character{40}}%
\TextAt{21}{248.755}{342.951}{m}%
\TextAt{17}{256.246}{342.951}{\Character{43}1\Character{41}}%
\TextAt{8}{282.813}{338.015}{where}%
\TextAt{13}{316.622}{338.015}{m}%
\TextAt{9}{330.182}{338.015}{\Character{61}}%
\TextAt{13}{342.608}{338.015}{m}%
\TextAt{9}{352.847}{338.015}{\Character{40}}%
\TextAt{13}{357.399}{338.015}{\Character{11}}%
\TextAt{21}{364.845}{336.222}{n}%
\TextAt{9}{370.481}{338.015}{\Character{41}}%
\TextAt{13}{378.354}{338.015}{\Character{60}}%
\TextAt{13}{390.764}{338.015}{m}%
\TextAt{9}{401.019}{338.015}{\Character{40}}%
\TextAt{13}{405.571}{338.015}{\Character{11}}%
\TextAt{21}{413.017}{336.222}{n}%
\TextAt{23}{418.155}{338.491}{\Character{48}}%
\TextAt{9}{421.338}{338.015}{\Character{41}}%
\TextAt{13}{425.89}{338.015}{\Character{58}}%
\TextAt{8}{79.226}{309.308}{Consequently}%
\TextAt{13}{226.7}{292.262}{g}%
\TextAt{9}{232.734}{292.262}{\Character{40}}%
\TextAt{13}{237.286}{292.262}{\Character{24}}%
\TextAt{9}{242.977}{292.262}{\Character{41}}%
\NativeMathAt{250.85}{292.262}{11.955}{\equiv}%
\TextAt{13}{263.469}{292.262}{g}%
\TextAt{21}{269.075}{290.469}{n}%
\TextAt{9}{274.711}{292.262}{\Character{40}}%
\TextAt{13}{279.263}{292.262}{\Character{11}}%
\TextAt{21}{286.709}{290.469}{n}%
\TextAt{9}{292.345}{292.262}{\Character{41}}%
\NativeMathAt{300.218}{292.262}{11.955}{\not\equiv}%
\TextAt{9}{312.838}{292.262}{1}%
\TextAt{8}{322.011}{292.262}{mod}%
\TextAt{13}{347.449}{292.262}{U}%
\TextAt{17}{356.648}{297.198}{\Character{40}}%
\TextAt{21}{359.942}{297.198}{m}%
\TextAt{17}{367.432}{297.198}{\Character{43}1\Character{41}}%
\TextAt{13}{382.044}{292.262}{\Character{59}}%
\TextAt{8}{79.226}{268.866}{and}%
\TextAt{8}{102.748}{268.866}{so}%
\TextAt{13}{117.877}{268.866}{g}%
\TextAt{9}{123.912}{268.866}{\Character{40}}%
\TextAt{13}{128.464}{268.866}{\Character{24}}%
\TextAt{9}{134.154}{268.866}{\Character{41}}%
\NativeMathAt{143.322}{268.866}{11.955}{\neq}%
\TextAt{9}{157.023}{268.866}{1}%
\TextAt{8}{162.894}{268.866}{.}%
\TextAt{8}{173.62}{268.866}{This}%
\TextAt{8}{201.099}{268.866}{shows}%
\TextAt{8}{235.476}{268.866}{the}%
\TextAt{8}{256.392}{268.866}{injectivity}%
\TextAt{8}{312.432}{268.866}{of}%
\TextAt{13}{326.521}{268.866}{g}%
\TextAt{8}{332.556}{268.866}{.}%
\TextAt{8}{343.282}{268.866}{Since}%
\TextAt{13}{374.599}{268.866}{A}%
\TextAt{9}{387.989}{268.866}{\Character{61}}%
\TextAt{18}{401.708}{268.866}{Z}%
\TextAt{25}{409.678}{273.205}{N}%
\TextAt{21}{409.678}{265.911}{p}%
\TextAt{8}{415.933}{268.866}{,}%
\TextAt{8}{424.031}{268.866}{this}%
\TextAt{8}{447.613}{268.866}{proves}%
\TextAt{8}{484.536}{268.866}{the}%
\TextAt{8}{505.452}{268.866}{claim}%
\TextAt{8}{79.226}{251.821}{(ii).}%
\Vector{q 0 G 0 g 0.399 w 452.334 -522.179 m 452.334 -514.11 l S Q}%
\Vector{q 0 G 0 g 0.399 w 452.534 -514.309 m 459.707 -514.309 l S Q}%
\Vector{q 0 G 0 g 0.399 w 452.534 -521.98 m 459.707 -521.98 l S Q}%
\Vector{q 0 G 0 g 0.399 w 459.906 -522.179 m 459.906 -514.11 l S Q}%
\TextAt{12}{79.226}{217.01}{Corollary}%
\TextAt{12}{140.469}{217.01}{(5.8).}%
\TextAt{8}{177.049}{217.01}{If}%
\TextAt{8}{190.005}{217.01}{the}%
\TextAt{8}{211.426}{217.01}{natural}%
\TextAt{8}{253.673}{217.01}{number}%
\TextAt{13}{297.534}{217.01}{n}%
\TextAt{8}{309.683}{217.01}{is}%
\TextAt{8}{322.712}{217.01}{not}%
\TextAt{8}{344.79}{217.01}{divisible}%
\TextAt{8}{391.962}{217.01}{by}%
\TextAt{8}{409.486}{217.01}{the}%
\TextAt{8}{430.907}{217.01}{characteristic}%
\TextAt{8}{504.105}{217.01}{of}%
\TextAt{13}{518.7}{217.01}{K}%
\TextAt{8}{529.518}{217.01}{,}%
\TextAt{8}{79.226}{199.964}{then}%
\TextAt{8}{106.012}{199.964}{one}%
\TextAt{8}{127.597}{199.964}{finds}%
\TextAt{8}{155.745}{199.964}{the}%
\TextAt{8}{176.028}{199.964}{following}%
\TextAt{8}{225.588}{199.964}{indices}%
\TextAt{8}{264.137}{199.964}{for}%
\TextAt{8}{282.148}{199.964}{the}%
\TextAt{8}{302.43}{199.964}{subgroups}%
\TextAt{8}{357.968}{199.964}{of}%
\TextAt{13}{371.423}{199.964}{n}%
\TextAt{8}{378.411}{199.964}{-th}%
\TextAt{8}{397.39}{199.964}{powers}%
\TextAt{13}{436.274}{199.964}{K}%
\TextAt{3}{447.092}{204.302}{\Character{3}}%
\TextAt{21}{451.326}{204.302}{n}%
\TextAt{8}{460.985}{199.964}{and}%
\TextAt{13}{483.874}{199.964}{U}%
\TextAt{21}{493.073}{204.302}{n}%
\TextAt{8}{502.732}{199.964}{in}%
\TextAt{8}{516.51}{199.964}{the}%
\TextAt{8}{79.226}{182.918}{multiplicative}%
\TextAt{8}{153.049}{182.918}{group}%
\TextAt{13}{186.225}{182.918}{K}%
\TextAt{3}{197.043}{187.256}{\Character{3}}%
\TextAt{8}{205.672}{182.918}{and}%
\TextAt{8}{228.435}{182.918}{in}%
\TextAt{8}{242.088}{182.918}{the}%
\TextAt{8}{262.244}{182.918}{unit}%
\TextAt{8}{286.956}{182.918}{group}%
\TextAt{13}{320.131}{182.918}{U}%
\TextAt{8}{329.33}{182.918}{:}%
\TextAt{9}{204.565}{154.211}{\Character{40}}%
\TextAt{13}{209.118}{154.211}{K}%
\TextAt{3}{219.935}{159.147}{\Character{3}}%
\TextAt{9}{227.989}{154.211}{\Character{58}}%
\TextAt{13}{234.561}{154.211}{K}%
\TextAt{3}{245.379}{159.147}{\Character{3}}%
\TextAt{21}{249.613}{159.147}{n}%
\TextAt{9}{255.249}{154.211}{\Character{41}}%
\TextAt{9}{263.104}{154.211}{\Character{61}}%
\TextAt{13}{275.548}{154.211}{n}%
\TextAt{9}{282.536}{154.211}{\Character{40}}%
\TextAt{13}{287.088}{154.211}{U}%
\TextAt{9}{299.608}{154.211}{\Character{58}}%
\TextAt{13}{306.18}{154.211}{U}%
\TextAt{21}{315.379}{159.147}{n}%
\TextAt{9}{321.016}{154.211}{\Character{41}}%
\TextAt{9}{328.871}{154.211}{\Character{61}}%
\TextAt{13}{348.074}{162.299}{n}%
\Vector{q 0 G 0 g 0.478 w 270.51 -616.8 m 288.626 -616.8 l S Q}%
\TextAt{16}{342.51}{146.01}{\Character{106}}%
\TextAt{13}{345.831}{146.01}{n}%
\TextAt{16}{352.818}{146.01}{\Character{106}}%
\TextAt{27}{356.139}{142.523}{p}%
\TextAt{8}{361.822}{154.211}{\char35\relax }%
\TextAt{13}{371.577}{154.211}{\Character{22}}%
\TextAt{21}{378.62}{152.418}{n}%
\TextAt{9}{384.256}{154.211}{\Character{40}}%
\TextAt{13}{388.809}{154.211}{K}%
\TextAt{9}{399.626}{154.211}{\Character{41}}%
\TextAt{13}{404.179}{154.211}{\Character{58}}%
\TextAt{19}{79.226}{101.628}{Proof.}%
\TextAt{8}{116}{101.628}{The}%
\TextAt{8}{140.054}{101.628}{first}%
\TextAt{8}{164.18}{101.628}{equality}%
\TextAt{8}{208.736}{101.628}{is}%
\TextAt{8}{220.5}{101.628}{a}%
\TextAt{8}{230.256}{101.628}{consequence}%
\TextAt{8}{296.32}{101.628}{of}%
\TextAt{13}{309.65}{101.628}{K}%
\TextAt{3}{320.468}{105.967}{\Character{3}}%
\TextAt{9}{328.521}{101.628}{\Character{61}}%
\TextAt{9}{340.93}{101.628}{\Character{40}}%
\TextAt{13}{345.499}{101.628}{\Character{25}}%
\TextAt{9}{352.568}{101.628}{\Character{41}}%
\TextAt{16}{359.777}{101.628}{\Character{2}}%
\TextAt{13}{371.732}{101.628}{U}%
\TextAt{8}{380.931}{101.628}{.}%
\TextAt{8}{389.38}{101.628}{By}%
\TextAt{8}{407.743}{101.628}{(5.7),}%
\TextAt{8}{438.97}{101.628}{we}%
\TextAt{8}{456.197}{101.628}{have}%
\TextAt{13}{197.491}{72.921}{U}%
\NativeMathAt{210.011}{72.921}{11.955}{\cong}%
\TextAt{13}{222.63}{72.921}{\Character{22}}%
\TextAt{9}{229.673}{72.921}{\Character{40}}%
\TextAt{13}{234.225}{72.921}{K}%
\TextAt{9}{245.043}{72.921}{\Character{41}}%
\TextAt{16}{252.252}{72.921}{\Character{2}}%
\TextAt{18}{264.207}{72.921}{Z}%
\TextAt{21}{272.177}{77.858}{d}%
\TextAt{21}{272.177}{69.966}{p}%
\TextAt{13}{277.033}{72.921}{\Character{59}}%
\TextAt{8}{294.232}{72.921}{resp.}%
\TextAt{13}{330.313}{72.921}{U}%
\NativeMathAt{342.832}{72.921}{11.955}{\cong}%
\TextAt{13}{355.452}{72.921}{\Character{22}}%
\TextAt{9}{362.495}{72.921}{\Character{40}}%
\TextAt{13}{367.047}{72.921}{K}%
\TextAt{9}{377.865}{72.921}{\Character{41}}%
\TextAt{16}{385.074}{72.921}{\Character{2}}%
\TextAt{18}{397.029}{72.921}{Z}%
\TextAt{25}{404.999}{77.858}{N}%
\TextAt{21}{404.999}{69.966}{p}%
\TextAt{13}{411.253}{72.921}{\Character{59}}%
\TextAt{8}{297.211}{43.033}{129}%
\Destination{78.23}{809.815}{nav.page.129}%
\Destination{79.36}{776.155}{nav.print.141}%
\Destination{79.36}{232.055}{nav.print.142}%
\Destination{79.36}{231.055}{nav.number.2.5.8}%
\Hotlink{407.243}{99.455}{25.078}{11.047}{nav.number.2.5.7}%
}

% Page 142 (body page 130)
\Page{%
\TextAt{2}{79.226}{796.943}{II.}%
\TextAt{2}{94.073}{796.943}{The}%
\TextAt{2}{116.492}{796.943}{Theory}%
\TextAt{2}{154.401}{796.943}{of}%
\TextAt{2}{166.826}{796.943}{Valuations}%
\Vector{q 0 G 0 g 0.299 w 7.226 18.729 m 460.769 18.729 l S Q}%
\TextAt{8}{79.226}{761.124}{when}%
\TextAt{8}{109.784}{761.124}{char}%
\TextAt{9}{131.901}{761.124}{\Character{40}}%
\TextAt{13}{136.453}{761.124}{K}%
\TextAt{9}{147.271}{761.124}{\Character{41}}%
\TextAt{9}{155.126}{761.124}{\Character{61}}%
\TextAt{9}{167.547}{761.124}{0}%
\TextAt{8}{173.422}{761.124}{,}%
\TextAt{8}{180.571}{761.124}{resp.}%
\TextAt{13}{208.594}{761.124}{p}%
\TextAt{13}{217.776}{761.124}{\Character{62}}%
\TextAt{9}{230.216}{761.124}{0}%
\TextAt{8}{236.069}{761.124}{.}%
\TextAt{8}{244.517}{761.124}{From}%
\TextAt{8}{275.242}{761.124}{the}%
\TextAt{8}{295.398}{761.124}{exact}%
\TextAt{8}{326.291}{761.124}{sequence}%
\TextAt{9}{161.439}{732.123}{1}%
\NativeArrowAt{170.613}{732.123}{11.955}{19.26}{\longrightarrow}%
\TextAt{13}{193.195}{732.123}{\Character{22}}%
\TextAt{21}{200.237}{730.329}{n}%
\TextAt{9}{205.874}{732.123}{\Character{40}}%
\TextAt{13}{210.426}{732.123}{K}%
\TextAt{9}{221.244}{732.123}{\Character{41}}%
\NativeArrowAt{229.117}{732.123}{11.955}{19.26}{\longrightarrow}%
\TextAt{13}{251.699}{732.123}{\Character{22}}%
\TextAt{9}{258.742}{732.123}{\Character{40}}%
\TextAt{13}{263.294}{732.123}{K}%
\TextAt{9}{274.112}{732.123}{\Character{41}}%
\TextAt{21}{284.688}{738.9}{n}%
\NativeArrowAt{281.985}{732.123}{11.955}{11.955}{\longrightarrow}%
\TextAt{13}{297.261}{732.123}{\Character{22}}%
\TextAt{9}{304.304}{732.123}{\Character{40}}%
\TextAt{13}{308.857}{732.123}{K}%
\TextAt{9}{319.674}{732.123}{\Character{41}}%
\NativeArrowAt{327.547}{732.123}{11.955}{19.26}{\longrightarrow}%
\TextAt{13}{350.129}{732.123}{\Character{22}}%
\TextAt{9}{357.172}{732.123}{\Character{40}}%
\TextAt{13}{361.725}{732.123}{K}%
\TextAt{9}{372.542}{732.123}{\Character{41}\Character{47}}%
\TextAt{13}{382.948}{732.123}{\Character{22}}%
\TextAt{9}{389.991}{732.123}{\Character{40}}%
\TextAt{13}{394.543}{732.123}{K}%
\TextAt{9}{405.361}{732.123}{\Character{41}}%
\TextAt{21}{409.913}{737.059}{n}%
\NativeArrowAt{418.87}{732.123}{11.955}{19.26}{\longrightarrow}%
\TextAt{9}{441.452}{732.123}{1}%
\TextAt{13}{447.305}{732.123}{\Character{59}}%
\TextAt{8}{79.226}{703.121}{one}%
\TextAt{8}{100.686}{703.121}{has}%
\TextAt{8}{121.559}{703.121}{\char35\relax }%
\TextAt{13}{131.315}{703.121}{\Character{22}}%
\TextAt{21}{138.358}{701.328}{n}%
\TextAt{9}{143.994}{703.121}{\Character{40}}%
\TextAt{13}{148.546}{703.121}{K}%
\TextAt{9}{159.364}{703.121}{\Character{41}}%
\TextAt{9}{167.219}{703.121}{\Character{61}}%
\TextAt{8}{179.663}{703.121}{\char35\relax }%
\TextAt{13}{189.418}{703.121}{\Character{22}}%
\TextAt{9}{196.461}{703.121}{\Character{40}}%
\TextAt{13}{201.013}{703.121}{K}%
\TextAt{9}{211.831}{703.121}{\Character{41}\Character{47}}%
\TextAt{13}{222.236}{703.121}{\Character{22}}%
\TextAt{9}{229.279}{703.121}{\Character{40}}%
\TextAt{13}{233.832}{703.121}{K}%
\TextAt{9}{244.649}{703.121}{\Character{41}}%
\TextAt{21}{249.202}{707.46}{n}%
\TextAt{8}{254.838}{703.121}{.}%
\TextAt{8}{263.286}{703.121}{When}%
\TextAt{8}{297.418}{703.121}{char}%
\TextAt{9}{319.536}{703.121}{\Character{40}}%
\TextAt{13}{324.088}{703.121}{K}%
\TextAt{9}{334.906}{703.121}{\Character{41}}%
\TextAt{9}{342.761}{703.121}{\Character{61}}%
\TextAt{9}{355.182}{703.121}{0}%
\TextAt{8}{361.057}{703.121}{,}%
\TextAt{8}{368.206}{703.121}{this}%
\TextAt{8}{391.029}{703.121}{gives:}%
\TextAt{9}{135.256}{674.12}{\Character{40}}%
\TextAt{13}{139.808}{674.12}{U}%
\TextAt{9}{152.328}{674.12}{\Character{58}}%
\TextAt{13}{158.901}{674.12}{U}%
\TextAt{21}{168.1}{679.056}{n}%
\TextAt{9}{173.736}{674.12}{\Character{41}}%
\TextAt{9}{181.591}{674.12}{\Character{61}}%
\TextAt{8}{194.035}{674.12}{\char35\relax }%
\TextAt{13}{203.79}{674.12}{\Character{22}}%
\TextAt{21}{210.833}{672.327}{n}%
\TextAt{9}{216.469}{674.12}{\Character{40}}%
\TextAt{13}{221.022}{674.12}{K}%
\TextAt{9}{231.839}{674.12}{\Character{41}}%
\TextAt{8}{236.392}{674.12}{\char35\relax }%
\TextAt{9}{246.147}{674.12}{\Character{40}}%
\TextAt{18}{250.699}{674.12}{Z}%
\TextAt{21}{258.67}{672.327}{p}%
\TextAt{9}{263.43}{674.12}{\Character{47}}%
\TextAt{13}{269.283}{674.12}{n}%
\TextAt{18}{276.271}{674.12}{Z}%
\TextAt{21}{284.241}{672.327}{p}%
\TextAt{9}{289.002}{674.12}{\Character{41}}%
\TextAt{21}{293.554}{679.056}{d}%
\TextAt{9}{301.731}{674.12}{\Character{61}}%
\TextAt{8}{314.156}{674.12}{\char35\relax }%
\TextAt{13}{323.912}{674.12}{\Character{22}}%
\TextAt{21}{330.955}{672.327}{n}%
\TextAt{9}{336.591}{674.12}{\Character{40}}%
\TextAt{13}{341.143}{674.12}{K}%
\TextAt{9}{351.961}{674.12}{\Character{41}}%
\TextAt{13}{356.513}{674.12}{p}%
\TextAt{21}{362.388}{679.056}{dv}%
\TextAt{22}{370.869}{678.06}{p}%
\TextAt{17}{375.208}{679.056}{\Character{40}}%
\TextAt{21}{378.501}{679.056}{n}%
\TextAt{17}{383.639}{679.056}{\Character{41}}%
\TextAt{9}{390.751}{674.12}{\Character{61}}%
\TextAt{8}{403.177}{674.12}{\char35\relax }%
\TextAt{13}{412.932}{674.12}{\Character{22}}%
\TextAt{21}{419.975}{672.327}{n}%
\TextAt{9}{425.612}{674.12}{\Character{40}}%
\TextAt{13}{430.164}{674.12}{K}%
\TextAt{9}{440.982}{674.12}{\Character{41}\Character{47}}%
\TextAt{16}{453.379}{674.12}{\Character{106}}%
\TextAt{13}{456.7}{674.12}{n}%
\TextAt{16}{463.688}{674.12}{\Character{106}}%
\TextAt{27}{467.009}{670.633}{p}%
\TextAt{13}{473.488}{674.12}{\Character{59}}%
\TextAt{8}{79.226}{645.119}{and}%
\TextAt{8}{102.681}{645.119}{when}%
\TextAt{8}{133.93}{645.119}{char}%
\TextAt{9}{156.047}{645.119}{\Character{40}}%
\TextAt{13}{160.599}{645.119}{K}%
\TextAt{9}{171.417}{645.119}{\Character{41}}%
\TextAt{9}{180.455}{645.119}{\Character{61}}%
\TextAt{13}{194.073}{645.119}{p}%
\TextAt{8}{199.948}{645.119}{,}%
\TextAt{8}{207.962}{645.119}{one}%
\TextAt{8}{230.113}{645.119}{gets}%
\TextAt{8}{254.93}{645.119}{simply}%
\TextAt{9}{293.078}{645.119}{\Character{40}}%
\TextAt{13}{297.63}{645.119}{U}%
\TextAt{9}{311.329}{645.119}{\Character{58}}%
\TextAt{13}{319.08}{645.119}{U}%
\TextAt{21}{328.279}{649.458}{n}%
\TextAt{9}{333.915}{645.119}{\Character{41}}%
\TextAt{9}{342.953}{645.119}{\Character{61}}%
\TextAt{8}{356.572}{645.119}{\char35\relax }%
\TextAt{13}{366.327}{645.119}{\Character{22}}%
\TextAt{21}{373.37}{643.326}{n}%
\TextAt{9}{379.006}{645.119}{\Character{40}}%
\TextAt{13}{383.559}{645.119}{K}%
\TextAt{9}{394.376}{645.119}{\Character{41}}%
\TextAt{9}{403.414}{645.119}{\Character{61}}%
\TextAt{8}{417.033}{645.119}{\char35\relax }%
\TextAt{13}{426.788}{645.119}{\Character{22}}%
\TextAt{21}{433.831}{643.326}{n}%
\TextAt{9}{439.467}{645.119}{\Character{40}}%
\TextAt{13}{444.02}{645.119}{K}%
\TextAt{9}{454.837}{645.119}{\Character{41}\Character{47}}%
\TextAt{16}{467.235}{645.119}{\Character{106}}%
\TextAt{13}{470.556}{645.119}{n}%
\TextAt{16}{477.544}{645.119}{\Character{106}}%
\TextAt{27}{480.865}{641.632}{p}%
\TextAt{8}{485.351}{645.119}{,}%
\TextAt{8}{493.365}{645.119}{because}%
\TextAt{9}{79.226}{628.073}{\Character{40}}%
\TextAt{13}{83.778}{628.073}{n\Character{59}}%
\TextAt{13}{95.984}{628.073}{p}%
\TextAt{9}{101.885}{628.073}{\Character{41}}%
\TextAt{9}{109.74}{628.073}{\Character{61}}%
\TextAt{9}{122.161}{628.073}{1}%
\TextAt{8}{128.037}{628.073}{,}%
\TextAt{8}{135.186}{628.073}{i.e.,}%
\TextAt{13}{157.291}{628.073}{n}%
\TextAt{18}{164.279}{628.073}{Z}%
\TextAt{21}{172.249}{626.28}{p}%
\TextAt{9}{180.331}{628.073}{\Character{61}}%
\TextAt{18}{192.756}{628.073}{Z}%
\TextAt{21}{200.726}{626.28}{p}%
\TextAt{8}{205.487}{628.073}{.}%
\Vector{q 0 G 0 g 0.399 w 452.334 -145.927 m 452.334 -137.857 l S Q}%
\Vector{q 0 G 0 g 0.399 w 452.534 -138.056 m 459.707 -138.056 l S Q}%
\Vector{q 0 G 0 g 0.399 w 452.534 -145.728 m 459.707 -145.728 l S Q}%
\Vector{q 0 G 0 g 0.399 w 459.906 -145.927 m 459.906 -137.857 l S Q}%
\TextAt{12}{79.226}{593.094}{Exercise}%
\TextAt{12}{130.979}{593.094}{II.5.1.}%
\TextAt{8}{168.694}{593.094}{The}%
\TextAt{8}{191.489}{593.094}{logarithm}%
\TextAt{8}{243.574}{593.094}{function}%
\TextAt{8}{288.163}{593.094}{can}%
\TextAt{8}{308.364}{593.094}{be}%
\TextAt{8}{323.03}{593.094}{continued}%
\TextAt{8}{375.103}{593.094}{to}%
\TextAt{8}{388.155}{593.094}{a}%
\TextAt{8}{396.651}{593.094}{continuous}%
\TextAt{8}{453.997}{593.094}{homomorphism}%
\ErrataAt{79.19}{575.937}{455.81}{$\log\colon\overline{\mathbb Q}_p^{\,*}\to\overline{\mathbb Q}_p$, and the exponential to a continuous homomorphism}%
\ErrataAt{79.19}{557.842}{455.81}{$\exp\colon\{x\in\overline{\mathbb Q}_p:v_p(x)>1/(p-1)\}\to\overline{\mathbb Q}_p^{\,*}$, where $v_p$ is the unique extension}%
\ErrataAt{79.19}{540.796}{455.81}{of the normalized valuation on $\mathbb Q_p$.}%
\TextAt{12}{79.226}{508.826}{Exercise}%
\TextAt{12}{132.622}{508.826}{II.5.2.}%
\TextAt{8}{171.766}{508.826}{Let}%
\TextAt{13}{192.905}{508.826}{K}%
\TextAt{16}{207.333}{508.826}{\Character{106}}%
\TextAt{18}{214.264}{508.826}{Q}%
\TextAt{21}{223.562}{507.032}{p}%
\TextAt{8}{232.39}{508.826}{be}%
\TextAt{8}{248.484}{508.826}{a}%
\TextAt{13}{258.41}{508.826}{p}%
\TextAt{8}{264.285}{508.826}{-adic}%
\TextAt{8}{293.063}{508.826}{number}%
\TextAt{8}{335.829}{508.826}{field.}%
\TextAt{8}{366.246}{508.826}{For}%
\TextAt{9}{387.386}{508.826}{1}%
\TextAt{9}{395.994}{508.826}{\Character{43}}%
\TextAt{13}{407.888}{508.826}{x}%
\NativeMathAt{418.15}{508.826}{11.955}{\in}%
\TextAt{13}{429.731}{508.826}{U}%
\TextAt{17}{438.93}{513.164}{\Character{40}1\Character{41}}%
\TextAt{8}{454.316}{508.826}{and}%
\TextAt{13}{477.248}{508.826}{z}%
\NativeMathAt{486.829}{508.826}{11.955}{\in}%
\TextAt{18}{498.409}{508.826}{Z}%
\TextAt{21}{506.379}{507.032}{p}%
\TextAt{8}{515.207}{508.826}{one}%
\TextAt{8}{79.226}{491.78}{has}%
\TextAt{9}{248.481}{471.093}{\Character{40}1}%
\TextAt{9}{261.524}{471.093}{\Character{43}}%
\TextAt{13}{273.304}{471.093}{x}%
\TextAt{9}{279.956}{471.093}{\Character{41}}%
\TextAt{21}{284.509}{476.029}{z}%
\TextAt{9}{292.603}{471.093}{\Character{61}}%
\NativeMathAt{309.429}{486.037}{7.97}{\infty}%
\TextAt{20}{305.029}{482.45}{\Character{88}}%
\TextAt{21}{305.911}{457.344}{\Character{23}}%
\TextAt{17}{310.595}{457.344}{\Character{61}0}%
\TextAt{20}{324.29}{487.95}{\Character{18}}%
\TextAt{13}{333.374}{479.181}{z}%
\TextAt{13}{333.09}{462.892}{\Character{23}}%
\TextAt{20}{339.628}{487.95}{\Character{19}}%
\TextAt{13}{348.429}{471.093}{x}%
\TextAt{21}{355.081}{476.029}{\Character{23}}%
\TextAt{13}{360.263}{471.093}{\Character{58}}%
\ErrataAt{79.19}{439.7}{455.81}{The series converges for $v_{\mathfrak p}(x)>0$ and $z\in\mathbb Z_p$. The stricter bound $e/(p-1)$ enters Exercise \hyperlink{nav.ex.2.5.3}{3}.}%
\TextAt{12}{79.226}{407.72}{Exercise}%
\TextAt{12}{132.427}{407.72}{II.5.3.}%
\ErrataAt{171}{407.72}{364}{Under the hypotheses of Exercise \hyperlink{nav.ex.2.5.2}{2}, the exponential identity holds when}%
\ErrataAt{79.19}{395}{455.81}{$v_{\mathfrak p}(x)>e/(p-1)$. The logarithmic identity holds for every $1+x\in U^{(1)}$:}%
\ErrataAt{130}{378.7}{356}{$(1+x)^z=\exp\bigl(z\log(1+x)\bigr)\qquad\text{and}\qquad\log\bigl((1+x)^z\bigr)=z\log(1+x).$}%
\TextAt{12}{79.226}{334.774}{Exercise}%
\TextAt{12}{133.383}{334.774}{II.5.4.}%
\TextAt{8}{173.188}{334.774}{For}%
\TextAt{8}{194.989}{334.774}{a}%
\TextAt{13}{205.576}{334.774}{p}%
\TextAt{8}{211.451}{334.774}{-adic}%
\TextAt{8}{240.891}{334.774}{number}%
\TextAt{8}{284.319}{334.774}{field}%
\TextAt{13}{310.507}{334.774}{K}%
\TextAt{8}{321.325}{334.774}{,}%
\TextAt{8}{329.514}{334.774}{every}%
\TextAt{8}{361.237}{334.774}{subgroup}%
\TextAt{8}{412.866}{334.774}{of}%
\TextAt{8}{427.027}{334.774}{finite}%
\TextAt{8}{457.771}{334.774}{index}%
\TextAt{8}{490.14}{334.774}{in}%
\TextAt{13}{504.624}{334.774}{K}%
\TextAt{3}{515.442}{339.113}{\Character{3}}%
\TextAt{8}{524.903}{334.774}{is}%
\TextAt{8}{79.226}{317.728}{both}%
\TextAt{8}{106.866}{317.728}{open}%
\TextAt{8}{135.152}{317.728}{and}%
\TextAt{8}{157.915}{317.728}{closed.}%
\ErrataAt{79.22}{285.755}{453.64}{\textbf{Exercise II.5.5.} If $K$ is a $p$-adic number field, the subgroups $U^n$ of $n$-th powers}%
\ErrataAt{79.22}{268.655}{453.64}{in $U=\mathfrak o^*$, for $n\ge1$, form a basis of neighbourhoods of $1$ in $K^*$.}%
\TextAt{8}{297.211}{43.033}{130}%
\Destination{78.23}{809.815}{nav.page.130}%
\Destination{79.226}{607.055}{nav.ex.2.5.1}%
\Destination{79.226}{522.855}{nav.ex.2.5.2}%
\Destination{79.226}{421.755}{nav.ex.2.5.3}%
\Destination{79.226}{348.755}{nav.ex.2.5.4}%
\Destination{79.226}{299.755}{nav.ex.2.5.5}%
\Destination{79.226}{590.055}{errata.16}%
\Destination{79.226}{299.755}{errata.17}%
\Destination{79.226}{592.937}{errata.77}%
\Destination{79.226}{424.72}{errata.78}%
\Destination{79.226}{456.7}{errata.83}%
}

% Page 143 (body page 131)
\Page{%
\TextAt{2}{79.226}{796.943}{II.}%
\TextAt{2}{94.073}{796.943}{The}%
\TextAt{2}{116.492}{796.943}{Theory}%
\TextAt{2}{154.401}{796.943}{of}%
\TextAt{2}{166.826}{796.943}{Valuations}%
\Vector{q 0 G 0 g 0.299 w 7.226 18.729 m 460.769 18.729 l S Q}%
\TextAt{12}{79.226}{761.124}{Exercise}%
\TextAt{12}{131.569}{761.124}{II.5.6.}%
\TextAt{8}{169.797}{761.124}{Let}%
\TextAt{13}{190.02}{761.124}{K}%
\TextAt{8}{203.989}{761.124}{be}%
\TextAt{8}{219.168}{761.124}{a}%
\TextAt{13}{228.177}{761.124}{p}%
\TextAt{8}{234.052}{761.124}{-adic}%
\TextAt{8}{261.915}{761.124}{number}%
\TextAt{8}{303.766}{761.124}{field,}%
\TextAt{13}{331.778}{761.124}{v}%
\TextAt{27}{337.437}{759.33}{p}%
\TextAt{8}{345.075}{761.124}{the}%
\TextAt{8}{364.486}{761.124}{normalized}%
\TextAt{8}{423.576}{761.124}{exponential}%
\TextAt{8}{485.594}{761.124}{valuation}%
\TextAt{8}{79.226}{744.078}{of}%
\TextAt{13}{92.169}{744.078}{K}%
\TextAt{8}{102.986}{744.078}{,}%
\TextAt{8}{109.825}{744.078}{and}%
\TextAt{8}{132.201}{744.078}{d}%
\TextAt{13}{138.704}{744.078}{x}%
\TextAt{8}{148.866}{744.078}{the}%
\TextAt{8}{168.635}{744.078}{Haar}%
\TextAt{8}{197.191}{744.078}{measure}%
\TextAt{8}{242.388}{744.078}{on}%
\TextAt{8}{258.26}{744.078}{the}%
\TextAt{8}{278.029}{744.078}{locally}%
\TextAt{8}{314.714}{744.078}{compact}%
\TextAt{8}{361.155}{744.078}{additive}%
\TextAt{8}{405.647}{744.078}{group}%
\TextAt{13}{438.435}{744.078}{K}%
\TextAt{17}{449.253}{748.416}{\Character{43}}%
\TextAt{8}{456.337}{744.078}{,}%
\TextAt{8}{463.176}{744.078}{scaled}%
\TextAt{8}{497.315}{744.078}{so}%
\TextAt{8}{511.298}{744.078}{that}%
\ErrataAt{79.22}{727.055}{453.64}{$\int_{\mathfrak o}dx=1$. Then $\int_{a\mathfrak o}dx=|a|_{\mathfrak p}=q^{-v_{\mathfrak p}(a)}$ for $a\ne0$, where $q=|\mathfrak o/\mathfrak p|$. Furthermore,}%
\TextAt{13}{247.641}{693.888}{I}%
\TextAt{9}{253.744}{693.888}{\Character{40}}%
\TextAt{13}{258.296}{693.888}{f}%
\TextAt{9}{265.343}{693.888}{\Character{41}}%
\TextAt{9}{273.198}{693.888}{\Character{61}}%
\TextAt{20}{285.641}{710.16}{\Character{90}}%
\TextAt{21}{292.283}{683.095}{K}%
\TextAt{25}{299.988}{683.095}{r}%
\TextAt{3}{306.187}{683.095}{\Character{102}}%
\TextAt{17}{310.421}{683.095}{0}%
\TextAt{3}{314.655}{683.095}{\Character{103}}%
\TextAt{13}{321.38}{693.888}{f}%
\TextAt{9}{328.426}{693.888}{\Character{40}}%
\TextAt{13}{332.978}{693.888}{x}%
\TextAt{9}{339.631}{693.888}{\Character{41}}%
\TextAt{8}{347.691}{701.976}{d}%
\TextAt{13}{354.195}{701.976}{x}%
\Vector{q 0 G 0 g 0.478 w 273.378 -77.123 m 291.159 -77.123 l S Q}%
\TextAt{16}{345.378}{685.687}{\Character{106}}%
\TextAt{13}{348.699}{685.687}{x}%
\TextAt{16}{355.351}{685.687}{\Character{106}}%
\TextAt{27}{358.672}{682.2}{p}%
\TextAt{8}{79.226}{658.304}{is}%
\TextAt{8}{90.99}{658.304}{a}%
\TextAt{8}{100.745}{658.304}{Haar}%
\TextAt{8}{129.689}{658.304}{measure}%
\TextAt{8}{175.274}{658.304}{on}%
\TextAt{8}{191.533}{658.304}{the}%
\TextAt{8}{211.689}{658.304}{locally}%
\TextAt{8}{248.762}{658.304}{compact}%
\TextAt{8}{295.591}{658.304}{group}%
\TextAt{13}{328.766}{658.304}{K}%
\TextAt{3}{339.584}{662.642}{\Character{3}}%
\TextAt{8}{344.316}{658.304}{.}%
\TextAt{15}{79.226}{605.589}{§}%
\TextAt{15}{93.785}{605.589}{6}%
\TextAt{15}{120.693}{605.589}{Henselian}%
\TextAt{15}{209.318}{605.589}{Fields}%
\TextAt{8}{79.226}{573.703}{Most}%
\TextAt{8}{109.237}{573.703}{results}%
\TextAt{8}{146.791}{573.703}{on}%
\TextAt{8}{163.412}{573.703}{complete}%
\TextAt{8}{213.196}{573.703}{valued}%
\TextAt{8}{250.308}{573.703}{fields}%
\TextAt{8}{280.642}{573.703}{can}%
\TextAt{8}{302.463}{573.703}{be}%
\TextAt{8}{318.749}{573.703}{derived}%
\TextAt{8}{360.082}{573.703}{from}%
\TextAt{8}{388.084}{573.703}{Hensel’s}%
\TextAt{8}{433.756}{573.703}{lemma}%
\TextAt{8}{471.836}{573.703}{alone,}%
\TextAt{8}{506.109}{573.703}{with-}%
\TextAt{8}{79.226}{556.657}{out}%
\TextAt{8}{101.599}{556.657}{the}%
\TextAt{8}{123.314}{556.657}{full}%
\TextAt{8}{145.352}{556.657}{strength}%
\TextAt{8}{193.153}{556.657}{of}%
\TextAt{8}{208.041}{556.657}{completeness.}%
\TextAt{8}{287.624}{556.657}{This}%
\TextAt{8}{315.903}{556.657}{lemma}%
\TextAt{8}{355.18}{556.657}{is}%
\TextAt{8}{368.503}{556.657}{valid}%
\TextAt{8}{398.359}{556.657}{in}%
\TextAt{8}{413.571}{556.657}{a}%
\TextAt{8}{424.885}{556.657}{much}%
\TextAt{8}{457.658}{556.657}{bigger}%
\TextAt{8}{494.341}{556.657}{class}%
\TextAt{8}{523.337}{556.657}{of}%
\TextAt{8}{79.226}{539.611}{nonarchimedean}%
\TextAt{8}{165.409}{539.611}{valued}%
\TextAt{8}{201.511}{539.611}{fields}%
\TextAt{8}{230.835}{539.611}{than}%
\TextAt{8}{257.504}{539.611}{the}%
\TextAt{8}{277.013}{539.611}{complete}%
\TextAt{8}{325.787}{539.611}{ones.}%
\TextAt{8}{356.197}{539.611}{For}%
\TextAt{8}{376.518}{539.611}{example,}%
\TextAt{8}{425.1}{539.611}{let}%
\TextAt{9}{441.356}{539.611}{\Character{40}}%
\TextAt{13}{445.909}{539.611}{K}%
\TextAt{13}{456.071}{539.611}{\Character{59}}%
\TextAt{13}{461.307}{539.611}{v}%
\TextAt{9}{467.408}{539.611}{\Character{41}}%
\TextAt{8}{475.21}{539.611}{be}%
\TextAt{8}{490.486}{539.611}{a}%
\TextAt{8}{499.594}{539.611}{nonar-}%
\TextAt{8}{79.226}{522.565}{chimedean}%
\TextAt{8}{136.421}{522.565}{valued}%
\TextAt{8}{172.813}{522.565}{field}%
\TextAt{8}{197.813}{522.565}{and}%
\TextAt{9}{220.218}{522.565}{\Character{40}}%
\TextAt{20}{227.509}{525.587}{\Character{98}}%
\TextAt{13}{224.77}{522.565}{K}%
\TextAt{13}{235.578}{522.565}{\Character{59}}%
\TextAt{9}{241.275}{522.565}{\Character{94}}%
\TextAt{13}{240.832}{522.565}{v}%
\TextAt{9}{246.92}{522.565}{\Character{41}}%
\TextAt{8}{255.013}{522.565}{its}%
\TextAt{8}{270.974}{522.565}{completion.}%
\TextAt{8}{335.241}{522.565}{Let}%
\TextAt{14}{355.853}{522.565}{o}%
\TextAt{8}{361.697}{522.565}{,}%
\TextAt{8}{368.56}{522.565}{resp.}%
\TextAt{20}{395.827}{523.1}{\Character{98}}%
\TextAt{14}{396.226}{522.565}{o}%
\TextAt{8}{402.07}{522.565}{,}%
\TextAt{8}{408.933}{522.565}{be}%
\TextAt{8}{424.5}{522.565}{the}%
\TextAt{8}{444.299}{522.565}{valuation}%
\TextAt{8}{495.014}{522.565}{rings}%
\TextAt{8}{523.337}{522.565}{of}%
\TextAt{13}{79.226}{505.519}{K}%
\TextAt{8}{90.044}{505.519}{,}%
\TextAt{8}{96.989}{505.519}{resp.}%
\TextAt{20}{127.495}{508.541}{\Character{98}}%
\TextAt{13}{124.756}{505.519}{K}%
\TextAt{8}{135.574}{505.519}{.}%
\TextAt{8}{143.937}{505.519}{We}%
\TextAt{8}{163.826}{505.519}{then}%
\TextAt{8}{190.231}{505.519}{consider}%
\TextAt{8}{235.561}{505.519}{the}%
\TextAt{8}{255.462}{505.519}{separable}%
\TextAt{8}{306.649}{505.519}{closure}%
\TextAt{13}{345.475}{505.519}{K}%
\TextAt{21}{355.435}{503.726}{v}%
\TextAt{8}{363.984}{505.519}{of}%
\TextAt{13}{377.059}{505.519}{K}%
\TextAt{8}{391.519}{505.519}{in}%
\TextAt{20}{407.654}{508.541}{\Character{98}}%
\TextAt{13}{404.916}{505.519}{K}%
\TextAt{8}{419.376}{505.519}{and}%
\TextAt{8}{441.883}{505.519}{the}%
\TextAt{8}{461.784}{505.519}{valuation}%
\TextAt{8}{512.601}{505.519}{ring}%
\TextAt{14}{79.226}{488.473}{o}%
\TextAt{21}{85.07}{486.68}{v}%
\NativeMathAt{93.299}{488.473}{11.955}{\subseteq}%
\TextAt{13}{105.918}{488.473}{K}%
\TextAt{21}{115.877}{486.68}{v}%
\TextAt{8}{124.682}{488.473}{with}%
\TextAt{8}{151.342}{488.473}{maximal}%
\TextAt{8}{199.151}{488.473}{ideal}%
\TextAt{14}{227.114}{488.473}{p}%
\TextAt{21}{233.097}{486.68}{v}%
\TextAt{8}{241.902}{488.473}{which}%
\TextAt{8}{275.389}{488.473}{is}%
\TextAt{8}{287.152}{488.473}{associated}%
\TextAt{8}{342.887}{488.473}{to}%
\TextAt{8}{357.198}{488.473}{the}%
\TextAt{8}{377.354}{488.473}{restriction}%
\TextAt{8}{433.352}{488.473}{of}%
\TextAt{9}{447.125}{488.473}{\Character{94}}%
\TextAt{13}{446.682}{488.473}{v}%
\TextAt{8}{456.668}{488.473}{to}%
\TextAt{13}{470.978}{488.473}{K}%
\TextAt{21}{480.937}{486.68}{v}%
\TextAt{8}{485.845}{488.473}{,}%
\TextAt{13}{228.443}{459.815}{K}%
\NativeMathAt{242.582}{459.815}{11.955}{\subseteq}%
\TextAt{13}{255.201}{459.815}{K}%
\TextAt{21}{265.16}{458.022}{v}%
\NativeMathAt{273.389}{459.815}{11.955}{\subseteq}%
\TextAt{20}{288.746}{462.837}{\Character{98}}%
\TextAt{13}{286.008}{459.815}{K}%
\TextAt{13}{296.815}{459.815}{\Character{59}}%
\TextAt{14}{325.98}{459.815}{o}%
\NativeMathAt{335.145}{459.815}{11.955}{\subseteq}%
\TextAt{14}{347.765}{459.815}{o}%
\TextAt{21}{353.609}{458.022}{v}%
\NativeMathAt{361.837}{459.815}{11.955}{\subseteq}%
\TextAt{20}{374.058}{460.351}{\Character{98}}%
\TextAt{14}{374.456}{459.815}{o}%
\TextAt{13}{380.301}{459.815}{\Character{58}}%
\TextAt{8}{79.226}{431.158}{Then}%
\TextAt{8}{110.528}{431.158}{Hensel’s}%
\TextAt{8}{156.583}{431.158}{lemma}%
\TextAt{8}{195.046}{431.158}{holds}%
\TextAt{8}{226.42}{431.158}{in}%
\TextAt{8}{240.818}{431.158}{the}%
\TextAt{8}{261.719}{431.158}{ring}%
\TextAt{14}{286.529}{431.158}{o}%
\TextAt{21}{292.374}{429.364}{v}%
\TextAt{8}{301.923}{431.158}{as}%
\TextAt{8}{317.038}{431.158}{well}%
\TextAt{8}{341.514}{431.158}{as}%
\TextAt{8}{356.629}{431.158}{in}%
\TextAt{8}{371.026}{431.158}{the}%
\TextAt{8}{391.927}{431.158}{ring}%
\TextAt{20}{416.339}{431.693}{\Character{98}}%
\TextAt{14}{416.738}{431.158}{o}%
\TextAt{8}{427.224}{431.158}{even}%
\TextAt{8}{454.629}{431.158}{though}%
\TextAt{13}{495.053}{431.158}{K}%
\TextAt{21}{505.012}{429.364}{v}%
\TextAt{8}{514.562}{431.158}{will}%
\TextAt{8}{79.226}{414.112}{not,}%
\TextAt{8}{103.989}{414.112}{as}%
\TextAt{8}{118.917}{414.112}{a}%
\TextAt{8}{129.231}{414.112}{rule,}%
\TextAt{8}{156.588}{414.112}{be}%
\TextAt{8}{173.07}{414.112}{complete.}%
\TextAt{8}{228.717}{414.112}{When}%
\TextAt{13}{263.407}{414.112}{K}%
\TextAt{21}{273.366}{412.318}{v}%
\TextAt{8}{282.729}{414.112}{is}%
\TextAt{8}{295.051}{414.112}{algebraically}%
\TextAt{8}{363.586}{414.112}{closed}%
\TextAt{8}{398.67}{414.112}{in}%
\TextAt{20}{415.619}{417.134}{\Character{98}}%
\TextAt{13}{412.881}{414.112}{K}%
\TextAt{8}{423.699}{414.112}{—hence}%
\TextAt{8}{468.467}{414.112}{in}%
\TextAt{8}{482.677}{414.112}{particular}%
\TextAt{8}{79.226}{397.066}{if}%
\TextAt{8}{90.839}{397.066}{char}%
\TextAt{9}{112.956}{397.066}{\Character{40}}%
\TextAt{13}{117.509}{397.066}{K}%
\TextAt{9}{128.326}{397.066}{\Character{41}}%
\TextAt{9}{137.699}{397.066}{\Character{61}}%
\TextAt{9}{151.639}{397.066}{0}%
\TextAt{8}{157.509}{397.066}{—this}%
\TextAt{8}{192.925}{397.066}{is}%
\TextAt{8}{205.579}{397.066}{immediately}%
\TextAt{8}{273.13}{397.066}{obvious}%
\TextAt{8}{316.365}{397.066}{(otherwise}%
\TextAt{8}{373.898}{397.066}{this}%
\TextAt{8}{397.609}{397.066}{follows}%
\TextAt{8}{436.935}{397.066}{from}%
\TextAt{8}{465.464}{397.066}{(6.6)}%
\TextAt{8}{494.329}{397.066}{and}%
\TextAt{8}{517.981}{397.066}{§6,}%
\TextAt{8}{79.226}{380.02}{exercise}%
\TextAt{8}{122.528}{380.02}{3}%
\TextAt{8}{132.283}{380.02}{below).}%
\TextAt{8}{174.553}{380.02}{Indeed,}%
\TextAt{8}{215.834}{380.02}{by}%
\TextAt{8}{232.093}{380.02}{(4.3)}%
\TextAt{8}{260.068}{380.02}{we}%
\TextAt{8}{277.295}{380.02}{have}%
\TextAt{14}{257.198}{351.362}{o}%
\TextAt{9}{263.042}{351.362}{\Character{47}}%
\TextAt{14}{268.895}{351.362}{p}%
\TextAt{9}{278.199}{351.362}{\Character{61}}%
\TextAt{14}{290.624}{351.362}{o}%
\TextAt{21}{296.469}{349.569}{v}%
\TextAt{9}{301.376}{351.362}{\Character{47}}%
\TextAt{14}{307.229}{351.362}{p}%
\TextAt{21}{313.212}{349.569}{v}%
\TextAt{9}{321.441}{351.362}{\Character{61}}%
\TextAt{20}{333.467}{351.897}{\Character{98}}%
\TextAt{14}{333.866}{351.362}{o}%
\TextAt{9}{339.71}{351.362}{\Character{47}}%
\TextAt{20}{345.234}{352.479}{\Character{98}}%
\TextAt{14}{345.563}{351.362}{p}%
\TextAt{13}{351.546}{351.362}{\Character{59}}%
\TextAt{8}{79.226}{322.704}{and}%
\TextAt{8}{102.326}{322.704}{if}%
\TextAt{8}{113.387}{322.704}{a}%
\TextAt{8}{123.48}{322.704}{primitive}%
\TextAt{8}{173.898}{322.704}{polynomial}%
\TextAt{13}{234.728}{322.704}{f}%
\TextAt{9}{241.775}{322.704}{\Character{40}}%
\TextAt{13}{246.327}{322.704}{x}%
\TextAt{9}{252.979}{322.704}{\Character{41}}%
\NativeMathAt{261.427}{322.704}{11.955}{\in}%
\TextAt{14}{273.293}{322.704}{o}%
\TextAt{21}{279.137}{320.911}{v}%
\TextAt{9}{284.045}{322.704}{\Character{91}}%
\TextAt{13}{287.296}{322.704}{x}%
\TextAt{9}{293.948}{322.704}{\Character{93}}%
\TextAt{8}{301.435}{322.704}{splits}%
\TextAt{8}{332.461}{322.704}{over}%
\TextAt{14}{357.844}{322.704}{o}%
\TextAt{21}{363.689}{320.911}{v}%
\TextAt{9}{368.596}{322.704}{\Character{47}}%
\TextAt{14}{374.449}{322.704}{p}%
\TextAt{21}{380.432}{320.911}{v}%
\TextAt{8}{389.575}{322.704}{into}%
\TextAt{8}{413.655}{322.704}{relatively}%
\TextAt{8}{465.053}{322.704}{prime}%
\TextAt{8}{498.554}{322.704}{factors}%
\TextAt{9}{79.642}{305.658}{\Character{22}}%
\TextAt{13}{79.226}{305.658}{g}%
\TextAt{9}{85.26}{305.658}{\Character{40}}%
\TextAt{13}{89.813}{305.658}{x}%
\TextAt{9}{96.465}{305.658}{\Character{41}}%
\TextAt{13}{101.017}{305.658}{\Character{59}}%
\TextAt{9}{106.379}{308.813}{\Character{22}}%
\TextAt{13}{106.261}{305.658}{h}%
\TextAt{9}{113}{305.658}{\Character{40}}%
\TextAt{13}{117.552}{305.658}{x}%
\TextAt{9}{124.204}{305.658}{\Character{41}}%
\TextAt{8}{128.757}{305.658}{,}%
\TextAt{8}{135.906}{305.658}{then}%
\TextAt{8}{162.566}{305.658}{we}%
\TextAt{8}{179.793}{305.658}{have}%
\TextAt{8}{206.788}{305.658}{by}%
\TextAt{8}{223.047}{305.658}{Hensel’s}%
\TextAt{8}{268.357}{305.658}{lemma}%
\TextAt{8}{306.076}{305.658}{(4.6)}%
\TextAt{8}{334.051}{305.658}{a}%
\TextAt{8}{343.806}{305.658}{factorization}%
\TextAt{8}{411.783}{305.658}{in}%
\TextAt{20}{425.037}{306.194}{\Character{98}}%
\TextAt{14}{425.436}{305.658}{o}%
\TextAt{13}{264.58}{277.001}{f}%
\TextAt{9}{271.626}{277.001}{\Character{40}}%
\TextAt{13}{276.179}{277.001}{x}%
\TextAt{9}{282.831}{277.001}{\Character{41}}%
\TextAt{9}{290.686}{277.001}{\Character{61}}%
\TextAt{13}{303.129}{277.001}{g}%
\TextAt{9}{309.164}{277.001}{\Character{40}}%
\TextAt{13}{313.716}{277.001}{x}%
\TextAt{9}{320.368}{277.001}{\Character{41}}%
\TextAt{13}{324.92}{277.001}{h}%
\TextAt{9}{331.659}{277.001}{\Character{40}}%
\TextAt{13}{336.211}{277.001}{x}%
\TextAt{9}{342.863}{277.001}{\Character{41}}%
\TextAt{8}{79.226}{248.343}{such}%
\TextAt{8}{106.565}{248.343}{that}%
\TextAt{13}{132.876}{248.343}{g}%
\NativeMathAt{143.837}{248.343}{11.955}{\equiv}%
\TextAt{9}{158.477}{248.343}{\Character{22}}%
\TextAt{13}{158.062}{248.343}{g}%
\TextAt{8}{167.417}{248.343}{mod}%
\TextAt{20}{192.525}{249.46}{\Character{98}}%
\TextAt{14}{192.855}{248.343}{p}%
\TextAt{8}{198.838}{248.343}{,}%
\TextAt{13}{207.164}{248.343}{h}%
\NativeMathAt{218.829}{248.343}{11.955}{\equiv}%
\TextAt{9}{233.172}{251.498}{\Character{22}}%
\TextAt{13}{233.054}{248.343}{h}%
\TextAt{8}{243.113}{248.343}{mod}%
\TextAt{20}{268.222}{249.46}{\Character{98}}%
\TextAt{14}{268.551}{248.343}{p}%
\TextAt{8}{274.534}{248.343}{,}%
\TextAt{8}{282.861}{248.343}{deg}%
\TextAt{9}{300.423}{248.343}{\Character{40}}%
\TextAt{13}{304.976}{248.343}{g}%
\TextAt{9}{311.01}{248.343}{\Character{41}}%
\TextAt{9}{320.478}{248.343}{\Character{61}}%
\TextAt{8}{334.519}{248.343}{deg}%
\TextAt{9}{352.082}{248.343}{\Character{40}}%
\TextAt{9}{357.031}{248.343}{\Character{22}}%
\TextAt{13}{356.634}{248.343}{g}%
\TextAt{9}{362.668}{248.343}{\Character{41}}%
\TextAt{8}{367.221}{248.343}{.}%
\TextAt{8}{378.495}{248.343}{But}%
\TextAt{8}{402.678}{248.343}{this}%
\TextAt{8}{426.443}{248.343}{factorization}%
\TextAt{8}{495.362}{248.343}{already}%
\TextAt{8}{79.226}{231.297}{takes}%
\TextAt{8}{110.27}{231.297}{place}%
\TextAt{8}{141.243}{231.297}{over}%
\TextAt{14}{167.35}{231.297}{o}%
\TextAt{21}{173.194}{229.504}{v}%
\TextAt{8}{183.06}{231.297}{once}%
\TextAt{8}{210.781}{231.297}{the}%
\TextAt{8}{231.998}{231.297}{highest}%
\TextAt{8}{273.443}{231.297}{coefficient}%
\TextAt{8}{329.127}{231.297}{of}%
\TextAt{13}{343.518}{231.297}{g}%
\TextAt{8}{354.51}{231.297}{is}%
\TextAt{8}{367.335}{231.297}{chosen}%
\TextAt{8}{405.851}{231.297}{to}%
\TextAt{8}{421.222}{231.297}{be}%
\TextAt{8}{438.207}{231.297}{in}%
\TextAt{14}{452.921}{231.297}{o}%
\TextAt{3}{458.765}{235.635}{\Character{3}}%
\TextAt{21}{458.765}{228.342}{v}%
\TextAt{8}{463.673}{231.297}{,}%
\TextAt{8}{472.148}{231.297}{because}%
\TextAt{8}{516.51}{231.297}{the}%
\TextAt{8}{79.226}{214.251}{coefficients}%
\TextAt{8}{138.464}{214.251}{of}%
\TextAt{13}{151.794}{214.251}{f}%
\TextAt{8}{158.84}{214.251}{,}%
\TextAt{8}{165.99}{214.251}{and}%
\TextAt{8}{188.752}{214.251}{therefore}%
\TextAt{8}{237.852}{214.251}{also}%
\TextAt{8}{261.332}{214.251}{those}%
\TextAt{8}{291.961}{214.251}{of}%
\TextAt{13}{305.291}{214.251}{g}%
\TextAt{8}{315.223}{214.251}{and}%
\TextAt{13}{337.986}{214.251}{h}%
\TextAt{8}{348.622}{214.251}{are}%
\TextAt{8}{368.132}{214.251}{algebraic}%
\TextAt{8}{417.567}{214.251}{over}%
\TextAt{13}{442.613}{214.251}{K}%
\TextAt{8}{453.431}{214.251}{.}%
\TextAt{8}{97.159}{194.216}{The}%
\TextAt{8}{121.663}{194.216}{valued}%
\TextAt{8}{158.864}{194.216}{field}%
\TextAt{13}{184.672}{194.216}{K}%
\TextAt{21}{194.631}{192.423}{v}%
\TextAt{8}{203.887}{194.216}{is}%
\TextAt{8}{216.102}{194.216}{called}%
\TextAt{8}{249.716}{194.216}{the}%
\TextAt{12}{270.323}{194.216}{henselization}%
\TextAt{8}{350.372}{194.216}{of}%
\TextAt{8}{364.152}{194.216}{the}%
\TextAt{8}{384.76}{194.216}{field}%
\TextAt{13}{410.567}{194.216}{K}%
\TextAt{8}{425.733}{194.216}{with}%
\TextAt{8}{452.844}{194.216}{respect}%
\TextAt{8}{493.344}{194.216}{to}%
\TextAt{13}{508.106}{194.216}{v}%
\TextAt{8}{514.194}{194.216}{.}%
\TextAt{8}{523.994}{194.216}{It}%
\TextAt{8}{79.226}{177.17}{enjoys}%
\TextAt{8}{114.551}{177.17}{all}%
\TextAt{8}{130.628}{177.17}{the}%
\TextAt{8}{150.602}{177.17}{relevant}%
\TextAt{8}{194.655}{177.17}{algebraic}%
\TextAt{8}{243.907}{177.17}{properties}%
\TextAt{8}{298.743}{177.17}{of}%
\TextAt{8}{311.891}{177.17}{the}%
\TextAt{8}{331.865}{177.17}{completion}%
\TextAt{20}{394.257}{180.192}{\Character{98}}%
\TextAt{13}{391.519}{177.17}{K}%
\TextAt{8}{402.337}{177.17}{,}%
\TextAt{8}{409.34}{177.17}{but}%
\TextAt{8}{430.618}{177.17}{offers}%
\TextAt{8}{461.388}{177.17}{the}%
\TextAt{8}{481.362}{177.17}{advantage}%
\TextAt{8}{79.226}{160.124}{of}%
\TextAt{8}{93.105}{160.124}{being}%
\TextAt{8}{125.191}{160.124}{itself}%
\TextAt{8}{154.085}{160.124}{an}%
\TextAt{8}{170.893}{160.124}{algebraic}%
\TextAt{8}{220.876}{160.124}{extension}%
\TextAt{8}{273.19}{160.124}{of}%
\TextAt{13}{287.069}{160.124}{K}%
\TextAt{8}{302.333}{160.124}{which}%
\TextAt{8}{336.368}{160.124}{can}%
\TextAt{8}{358.376}{160.124}{also}%
\TextAt{8}{382.404}{160.124}{be}%
\TextAt{8}{398.877}{160.124}{obtained}%
\TextAt{8}{447.557}{160.124}{in}%
\TextAt{8}{461.759}{160.124}{a}%
\TextAt{8}{472.062}{160.124}{purely}%
\TextAt{8}{508.704}{160.124}{alge-}%
\TextAt{8}{79.226}{143.078}{braic}%
\TextAt{8}{108.886}{143.078}{manner,}%
\TextAt{8}{154.903}{143.078}{without}%
\TextAt{8}{198.874}{143.078}{the}%
\TextAt{8}{219.424}{143.078}{analytic}%
\TextAt{8}{264.374}{143.078}{recourse}%
\TextAt{8}{310.353}{143.078}{to}%
\TextAt{8}{325.057}{143.078}{the}%
\TextAt{8}{345.607}{143.078}{completion}%
\TextAt{8}{405.836}{143.078}{(see}%
\TextAt{8}{429.698}{143.078}{§9,}%
\TextAt{8}{448.876}{143.078}{exercise}%
\TextAt{8}{492.571}{143.078}{4).}%
\TextAt{8}{512.613}{143.078}{The}%
\TextAt{8}{79.226}{126.032}{consequence}%
\TextAt{8}{146.165}{126.032}{is}%
\TextAt{8}{158.803}{126.032}{that}%
\TextAt{8}{185.046}{126.032}{taking}%
\TextAt{8}{222.024}{126.032}{the}%
\TextAt{8}{243.055}{126.032}{henselization}%
\TextAt{8}{313.58}{126.032}{of}%
\TextAt{8}{327.784}{126.032}{an}%
\TextAt{8}{344.917}{126.032}{infinite}%
\TextAt{8}{385.459}{126.032}{algebraic}%
\TextAt{8}{435.767}{126.032}{extension}%
\TextAt{13}{488.407}{126.032}{L}%
\TextAt{16}{501.182}{126.032}{\Character{106}}%
\TextAt{13}{509.314}{126.032}{K}%
\TextAt{8}{524.903}{126.032}{is}%
\TextAt{8}{79.226}{108.986}{possible}%
\TextAt{8}{123.245}{108.986}{within}%
\TextAt{8}{159.661}{108.986}{the}%
\TextAt{8}{179.817}{108.986}{category}%
\TextAt{8}{226.98}{108.986}{of}%
\TextAt{8}{240.31}{108.986}{algebraic}%
\TextAt{8}{289.745}{108.986}{extensions.}%
\TextAt{8}{350.676}{108.986}{Let}%
\TextAt{8}{371.646}{108.986}{us}%
\TextAt{8}{386.661}{108.986}{define}%
\TextAt{8}{420.471}{108.986}{in}%
\TextAt{8}{434.123}{108.986}{general:}%
\TextAt{8}{297.211}{43.033}{131}%
\Destination{78.23}{809.815}{nav.page.131}%
\Destination{79.23}{631.985}{nav.section.2.6}%
\Destination{79.226}{775.155}{nav.ex.2.5.6}%
\Destination{79.36}{588.655}{nav.print.143}%
\Destination{79.23}{631.985}{outline.26}%
\Hotlink{523.16}{394.855}{6.858}{11.047}{nav.section.2.6}%
\Hotlink{122.028}{377.855}{6.858}{11.047}{nav.ex.2.6.3}%
\Hotlink{492.071}{140.855}{6.858}{11.047}{nav.ex.2.9.4}%
\Hotlink{434.877}{140.855}{6.858}{11.047}{nav.section.2.9}%
\Hotlink{464.964}{394.855}{25.078}{11.047}{nav.number.2.6.6}%
\Hotlink{231.593}{377.855}{25.078}{11.047}{nav.number.2.4.3}%
\Destination{79.226}{741.055}{errata.18}%
\Hotlink{307.243}{303.5}{21.722}{11.047}{nav.number.2.4.6}%
}

% Page 144 (body page 132)
\Page{%
\TextAt{2}{79.226}{796.943}{II.}%
\TextAt{2}{94.073}{796.943}{The}%
\TextAt{2}{116.492}{796.943}{Theory}%
\TextAt{2}{154.401}{796.943}{of}%
\TextAt{2}{166.826}{796.943}{Valuations}%
\Vector{q 0 G 0 g 0.299 w 7.226 18.729 m 460.769 18.729 l S Q}%
\TextAt{12}{79.226}{761.124}{Definition}%
\TextAt{12}{141.474}{761.124}{(6.1).}%
\TextAt{8}{176.319}{761.124}{A}%
\TextAt{12}{188.521}{761.124}{henselian}%
\TextAt{12}{246.476}{761.124}{field}%
\TextAt{8}{274.722}{761.124}{is}%
\TextAt{8}{286.016}{761.124}{a}%
\TextAt{8}{295.301}{761.124}{field}%
\TextAt{8}{320.188}{761.124}{with}%
\TextAt{8}{346.378}{761.124}{a}%
\TextAt{8}{355.663}{761.124}{nonarchimedean}%
\TextAt{8}{442.023}{761.124}{valuation}%
\TextAt{13}{492.625}{761.124}{v}%
\TextAt{8}{502.14}{761.124}{whose}%
\TextAt{8}{79.226}{744.078}{valuation}%
\TextAt{8}{129.687}{744.078}{ring}%
\TextAt{14}{153.141}{744.078}{o}%
\TextAt{8}{162.271}{744.078}{satisfies}%
\TextAt{8}{204.769}{744.078}{Hensel’s}%
\TextAt{8}{249.467}{744.078}{lemma}%
\TextAt{8}{286.574}{744.078}{in}%
\TextAt{8}{299.615}{744.078}{the}%
\TextAt{8}{319.16}{744.078}{sense}%
\TextAt{8}{348.579}{744.078}{of}%
\TextAt{8}{361.298}{744.078}{(4.6).}%
\TextAt{8}{393.62}{744.078}{One}%
\TextAt{8}{417.719}{744.078}{also}%
\TextAt{8}{440.587}{744.078}{calls}%
\TextAt{8}{466.05}{744.078}{the}%
\TextAt{8}{485.594}{744.078}{valuation}%
\TextAt{13}{79.226}{727.032}{v}%
\TextAt{8}{85.314}{727.032}{,}%
\TextAt{8}{92.464}{727.032}{or}%
\TextAt{8}{106.774}{727.032}{the}%
\TextAt{8}{126.93}{727.032}{ring}%
\TextAt{14}{150.996}{727.032}{o}%
\TextAt{8}{156.84}{727.032}{,}%
\TextAt{8}{163.99}{727.032}{henselian.}%
\TextAt{12}{79.226}{697.994}{Theorem}%
\TextAt{12}{137.003}{697.994}{(6.2).}%
\TextAt{8}{172.982}{697.994}{Let}%
\TextAt{13}{194.614}{697.994}{K}%
\TextAt{8}{209.993}{697.994}{be}%
\TextAt{8}{226.58}{697.994}{a}%
\TextAt{8}{236.999}{697.994}{henselian}%
\TextAt{8}{288.448}{697.994}{field}%
\TextAt{8}{314.468}{697.994}{with}%
\TextAt{8}{341.791}{697.994}{respect}%
\TextAt{8}{382.504}{697.994}{to}%
\TextAt{8}{397.477}{697.994}{the}%
\TextAt{8}{418.297}{697.994}{valuation}%
\TextAt{16}{470.033}{697.994}{\Character{106}}%
\TextAt{16}{477.23}{697.994}{\Character{106}}%
\TextAt{8}{480.572}{697.994}{.}%
\TextAt{8}{491.01}{697.994}{Then}%
\TextAt{16}{522.23}{697.994}{\Character{106}}%
\TextAt{16}{529.439}{697.994}{\Character{106}}%
\TextAt{8}{79.226}{680.948}{admits}%
\TextAt{8}{117.607}{680.948}{one}%
\TextAt{8}{139.011}{680.948}{and}%
\TextAt{8}{161.719}{680.948}{only}%
\TextAt{8}{187.356}{680.948}{one}%
\TextAt{8}{208.76}{680.948}{extension}%
\TextAt{8}{260.471}{680.948}{to}%
\TextAt{8}{274.726}{680.948}{any}%
\TextAt{8}{296.788}{680.948}{given}%
\TextAt{8}{327.302}{680.948}{algebraic}%
\TextAt{8}{376.682}{680.948}{extension}%
\TextAt{13}{428.393}{680.948}{L}%
\TextAt{16}{439.678}{680.948}{\Character{106}}%
\TextAt{13}{446.32}{680.948}{K}%
\TextAt{8}{457.137}{680.948}{.}%
\TextAt{8}{465.567}{680.948}{It}%
\TextAt{8}{478.185}{680.948}{is}%
\TextAt{8}{489.894}{680.948}{given}%
\TextAt{8}{520.408}{680.948}{by}%
\TextAt{16}{257.657}{654.826}{\Character{106}}%
\TextAt{13}{260.978}{654.826}{\Character{11}}%
\TextAt{16}{268.5}{654.826}{\Character{106}}%
\TextAt{9}{275.142}{654.826}{\Character{61}}%
\TextAt{22}{290.888}{660.045}{n}%
\TextAt{20}{288.847}{669.218}{\Character{113}}%
\Vector{q 0 G 0 g 0.478 w 228.802 -104.543 m 279.087 -104.543 l S Q}%
\TextAt{20}{300.802}{664.988}{\Character{12}}%
\TextAt{20}{300.802}{657.815}{\Character{12}}%
\TextAt{13}{304.787}{654.826}{N}%
\TextAt{21}{314.161}{652.966}{L}%
\TextAt{3}{319.92}{652.966}{\Character{106}}%
\TextAt{21}{322.273}{652.966}{K}%
\TextAt{9}{330.475}{654.826}{\Character{40}}%
\TextAt{13}{335.028}{654.826}{\Character{11}}%
\TextAt{9}{342.549}{654.826}{\Character{41}}%
\TextAt{20}{347.102}{664.988}{\Character{12}}%
\TextAt{20}{347.102}{657.815}{\Character{12}}%
\TextAt{13}{351.087}{654.826}{\Character{59}}%
\TextAt{8}{79.226}{630.991}{if}%
\TextAt{13}{90.195}{630.991}{L}%
\TextAt{16}{101.899}{630.991}{\Character{106}}%
\TextAt{13}{108.959}{630.991}{K}%
\TextAt{8}{123.919}{630.991}{has}%
\TextAt{8}{145.038}{630.991}{finite}%
\TextAt{8}{175.195}{630.991}{degree}%
\TextAt{13}{211.856}{630.991}{n}%
\TextAt{8}{218.844}{630.991}{.}%
\TextAt{8}{228.028}{630.991}{In}%
\TextAt{8}{242.894}{630.991}{any}%
\TextAt{8}{265.257}{630.991}{case,}%
\TextAt{8}{293.586}{630.991}{the}%
\TextAt{8}{313.988}{630.991}{valuation}%
\TextAt{8}{365.306}{630.991}{ring}%
\TextAt{8}{389.617}{630.991}{of}%
\TextAt{8}{403.193}{630.991}{the}%
\TextAt{8}{423.594}{630.991}{extended}%
\TextAt{8}{473.585}{630.991}{valuation}%
\TextAt{8}{524.903}{630.991}{is}%
\TextAt{8}{79.226}{613.945}{the}%
\TextAt{8}{99.383}{613.945}{integral}%
\TextAt{8}{141.991}{613.945}{closure}%
\TextAt{8}{181.072}{613.945}{of}%
\TextAt{8}{194.402}{613.945}{the}%
\TextAt{8}{214.559}{613.945}{valuation}%
\TextAt{8}{265.631}{613.945}{ring}%
\TextAt{8}{289.697}{613.945}{of}%
\TextAt{13}{303.027}{613.945}{K}%
\TextAt{8}{317.742}{613.945}{in}%
\TextAt{13}{331.395}{613.945}{L}%
\TextAt{8}{339.359}{613.945}{.}%
\TextAt{8}{97.159}{589.409}{The}%
\TextAt{8}{122.68}{589.409}{proof}%
\TextAt{8}{154.717}{589.409}{of}%
\TextAt{8}{169.515}{589.409}{this}%
\TextAt{8}{193.805}{589.409}{theorem}%
\TextAt{8}{240.797}{589.409}{is}%
\TextAt{19}{254.029}{589.409}{verbatim}%
\TextAt{8}{303.161}{589.409}{the}%
\TextAt{8}{324.785}{589.409}{same}%
\TextAt{8}{355.579}{589.409}{as}%
\TextAt{8}{371.416}{589.409}{in}%
\TextAt{8}{386.537}{589.409}{the}%
\TextAt{8}{408.161}{589.409}{case}%
\TextAt{8}{434.399}{589.409}{of}%
\TextAt{8}{449.197}{589.409}{a}%
\TextAt{8}{460.42}{589.409}{complete}%
\TextAt{8}{511.31}{589.409}{field}%
\TextAt{8}{79.226}{572.363}{(see}%
\TextAt{8}{102.773}{572.363}{(4.8)).}%
\TextAt{8}{140.091}{572.363}{What}%
\TextAt{8}{173.011}{572.363}{is}%
\TextAt{8}{184.853}{572.363}{remarkable}%
\TextAt{8}{245.748}{572.363}{about}%
\TextAt{8}{279.326}{572.363}{our}%
\TextAt{8}{300.219}{572.363}{current}%
\TextAt{8}{340.945}{572.363}{setting}%
\TextAt{8}{379.46}{572.363}{is}%
\TextAt{8}{391.303}{572.363}{that,}%
\TextAt{8}{420.022}{572.363}{conversely,}%
\TextAt{8}{478.39}{572.363}{the}%
\TextAt{8}{498.626}{572.363}{unique}%
\TextAt{8}{79.226}{555.317}{extendability}%
\TextAt{8}{150.646}{555.317}{also}%
\TextAt{8}{174.974}{555.317}{characterizes}%
\TextAt{8}{245.151}{555.317}{henselian}%
\TextAt{8}{296.785}{555.317}{fields.}%
\TextAt{8}{333.853}{555.317}{In}%
\TextAt{8}{349.323}{555.317}{order}%
\TextAt{8}{380.741}{555.317}{to}%
\TextAt{8}{395.899}{555.317}{prove}%
\TextAt{8}{428.298}{555.317}{this,}%
\TextAt{8}{455.432}{555.317}{we}%
\TextAt{8}{473.508}{555.317}{appeal}%
\TextAt{8}{511.753}{555.317}{to}%
\TextAt{8}{526.911}{555.317}{a}%
\TextAt{8}{79.226}{538.271}{method}%
\TextAt{8}{122.329}{538.271}{which}%
\TextAt{8}{156.323}{538.271}{allows}%
\TextAt{8}{191.691}{538.271}{us}%
\TextAt{8}{207.213}{538.271}{to}%
\TextAt{8}{222.031}{538.271}{express}%
\TextAt{8}{263.305}{538.271}{the}%
\TextAt{8}{283.968}{538.271}{valuations}%
\TextAt{8}{340.162}{538.271}{of}%
\TextAt{8}{353.999}{538.271}{the}%
\TextAt{8}{374.662}{538.271}{roots}%
\TextAt{8}{404.83}{538.271}{of}%
\TextAt{8}{418.667}{538.271}{a}%
\TextAt{8}{428.929}{538.271}{polynomial}%
\TextAt{8}{489.929}{538.271}{in}%
\TextAt{8}{504.089}{538.271}{terms}%
\TextAt{8}{79.226}{521.225}{of}%
\TextAt{8}{92.823}{521.225}{the}%
\TextAt{8}{113.247}{521.225}{valuations}%
\TextAt{8}{169.201}{521.225}{of}%
\TextAt{8}{182.798}{521.225}{the}%
\TextAt{8}{203.221}{521.225}{coefficients.}%
\TextAt{8}{267.811}{521.225}{It}%
\TextAt{8}{280.751}{521.225}{relies}%
\TextAt{8}{310.99}{521.225}{on}%
\TextAt{8}{327.516}{521.225}{the}%
\TextAt{8}{347.939}{521.225}{notion}%
\TextAt{8}{384.634}{521.225}{of}%
\TextAt{12}{398.745}{521.225}{Newton}%
\TextAt{12}{449.347}{521.225}{polygon}%
\TextAt{8}{495.697}{521.225}{,}%
\TextAt{8}{503.18}{521.225}{which}%
\TextAt{8}{79.226}{504.179}{arises}%
\TextAt{8}{111.218}{504.179}{as}%
\TextAt{8}{125.588}{504.179}{follows.}%
\TextAt{8}{97.159}{484.144}{Let}%
\TextAt{13}{118.128}{484.144}{v}%
\TextAt{8}{128.114}{484.144}{be}%
\TextAt{8}{144.038}{484.144}{an}%
\TextAt{8}{160.297}{484.144}{arbitrary}%
\TextAt{8}{210.067}{484.144}{exponential}%
\TextAt{8}{272.831}{484.144}{valuation}%
\TextAt{8}{323.904}{484.144}{of}%
\TextAt{8}{337.234}{484.144}{the}%
\TextAt{8}{357.39}{484.144}{field}%
\TextAt{13}{382.747}{484.144}{K}%
\TextAt{8}{397.462}{484.144}{and}%
\TextAt{8}{420.225}{484.144}{let}%
\TextAt{13}{212.592}{460.309}{f}%
\TextAt{9}{219.638}{460.309}{\Character{40}}%
\TextAt{13}{224.19}{460.309}{x}%
\TextAt{9}{230.843}{460.309}{\Character{41}}%
\TextAt{9}{238.698}{460.309}{\Character{61}}%
\TextAt{13}{251.141}{460.309}{a}%
\TextAt{17}{257.286}{458.516}{0}%
\TextAt{9}{264.675}{460.309}{\Character{43}}%
\TextAt{13}{276.436}{460.309}{a}%
\TextAt{17}{282.581}{458.516}{1}%
\TextAt{13}{287.314}{460.309}{x}%
\TextAt{9}{296.622}{460.309}{\Character{43}}%
\TextAt{16}{308.384}{460.309}{\Character{1}}%
\TextAt{16}{313.68}{460.309}{\Character{1}}%
\TextAt{16}{318.988}{460.309}{\Character{1}}%
\TextAt{9}{324.988}{460.309}{\Character{43}}%
\TextAt{13}{336.749}{460.309}{a}%
\TextAt{21}{342.894}{458.516}{n}%
\TextAt{13}{348.531}{460.309}{x}%
\TextAt{21}{355.183}{465.246}{n}%
\NativeMathAt{364.14}{460.309}{11.955}{\in}%
\TextAt{13}{375.431}{460.309}{K}%
\TextAt{9}{386.249}{460.309}{\Character{91}}%
\TextAt{13}{389.5}{460.309}{x}%
\TextAt{9}{396.152}{460.309}{\Character{93}}%
\TextAt{8}{79.226}{436.475}{be}%
\TextAt{8}{94.051}{436.475}{a}%
\TextAt{8}{102.708}{436.475}{polynomial}%
\TextAt{8}{162.102}{436.475}{satisfying}%
\TextAt{13}{213.163}{436.475}{a}%
\TextAt{17}{219.308}{434.681}{0}%
\TextAt{13}{224.04}{436.475}{a}%
\TextAt{21}{230.185}{434.681}{n}%
\NativeMathAt{239.142}{436.475}{11.955}{\neq}%
\TextAt{9}{251.552}{436.475}{0}%
\TextAt{8}{257.421}{436.475}{.}%
\TextAt{8}{265.503}{436.475}{To}%
\TextAt{8}{281.631}{436.475}{each}%
\TextAt{8}{306.869}{436.475}{term}%
\TextAt{13}{333.733}{436.475}{a}%
\TextAt{21}{339.878}{434.681}{i}%
\TextAt{13}{343.259}{436.475}{x}%
\TextAt{21}{349.911}{440.813}{i}%
\TextAt{8}{356.091}{436.475}{we}%
\TextAt{8}{372.219}{436.475}{associate}%
\TextAt{8}{420.351}{436.475}{a}%
\TextAt{8}{429.008}{436.475}{point}%
\TextAt{9}{458.478}{436.475}{\Character{40}}%
\TextAt{13}{463.03}{436.475}{i\Character{59}}%
\TextAt{13}{472.248}{436.475}{v}%
\TextAt{9}{478.356}{436.475}{\Character{40}}%
\TextAt{13}{482.908}{436.475}{a}%
\TextAt{21}{489.053}{434.681}{i}%
\TextAt{9}{492.435}{436.475}{\Character{41}\Character{41}}%
\NativeMathAt{504.86}{436.475}{11.955}{\in}%
\TextAt{18}{516.151}{436.475}{R}%
\TextAt{17}{524.785}{440.813}{2}%
\TextAt{8}{529.518}{436.475}{,}%
\TextAt{8}{79.226}{419.429}{ignoring}%
\TextAt{8}{125.382}{419.429}{however}%
\TextAt{8}{170.879}{419.429}{the}%
\TextAt{8}{191.654}{419.429}{point}%
\TextAt{9}{222.842}{419.429}{\Character{40}}%
\TextAt{13}{227.394}{419.429}{i\Character{59}}%
\NativeMathAt{236.631}{419.429}{11.955}{\infty}%
\TextAt{9}{248.587}{419.429}{\Character{41}}%
\TextAt{8}{257.655}{419.429}{if}%
\TextAt{13}{268.997}{419.429}{a}%
\TextAt{21}{275.141}{417.635}{i}%
\TextAt{9}{282.897}{419.429}{\Character{61}}%
\TextAt{9}{296.359}{419.429}{0}%
\TextAt{8}{302.229}{419.429}{.}%
\TextAt{8}{312.532}{419.429}{We}%
\TextAt{8}{333.294}{419.429}{now}%
\TextAt{8}{358.301}{419.429}{take}%
\TextAt{8}{384.288}{419.429}{the}%
\TextAt{8}{405.063}{419.429}{lower}%
\TextAt{8}{436.25}{419.429}{convex}%
\TextAt{8}{475.245}{419.429}{envelope}%
\TextAt{8}{523.337}{419.429}{of}%
\TextAt{8}{79.226}{402.383}{the}%
\TextAt{8}{99.383}{402.383}{set}%
\TextAt{8}{117.65}{402.383}{of}%
\TextAt{8}{130.98}{402.383}{points}%
\TextAt{16}{212.235}{385.337}{\Character{102}}%
\TextAt{9}{218.212}{385.337}{\Character{40}0}%
\TextAt{13}{228.618}{385.337}{\Character{59}}%
\TextAt{13}{233.842}{385.337}{v}%
\TextAt{9}{239.95}{385.337}{\Character{40}}%
\TextAt{13}{244.502}{385.337}{a}%
\TextAt{17}{250.647}{383.543}{0}%
\TextAt{9}{255.38}{385.337}{\Character{41}\Character{41}}%
\TextAt{13}{264.484}{385.337}{\Character{59}}%
\TextAt{9}{269.728}{385.337}{\Character{40}1}%
\TextAt{13}{280.134}{385.337}{\Character{59}}%
\TextAt{13}{285.358}{385.337}{v}%
\TextAt{9}{291.466}{385.337}{\Character{40}}%
\TextAt{13}{296.018}{385.337}{a}%
\TextAt{17}{302.163}{383.543}{1}%
\TextAt{9}{306.896}{385.337}{\Character{41}\Character{41}}%
\TextAt{13}{316}{385.337}{\Character{59}}%
\TextAt{13}{321.224}{385.337}{\Character{58}}%
\TextAt{13}{326.461}{385.337}{\Character{58}}%
\TextAt{13}{331.697}{385.337}{\Character{58}}%
\TextAt{13}{336.922}{385.337}{\Character{59}}%
\TextAt{9}{342.221}{385.337}{\Character{40}}%
\TextAt{13}{346.773}{385.337}{n\Character{59}}%
\TextAt{13}{358.979}{385.337}{v}%
\TextAt{9}{365.093}{385.337}{\Character{40}}%
\TextAt{13}{369.646}{385.337}{a}%
\TextAt{21}{375.791}{383.543}{n}%
\TextAt{9}{381.427}{385.337}{\Character{41}\Character{41}}%
\TextAt{16}{390.532}{385.337}{\Character{103}}%
\TextAt{13}{396.509}{385.337}{\Character{58}}%
\TextAt{8}{79.226}{364.029}{This}%
\TextAt{8}{105.946}{364.029}{produces}%
\TextAt{8}{155.106}{364.029}{a}%
\TextAt{8}{164.861}{364.029}{polygonal}%
\TextAt{8}{218.205}{364.029}{chain}%
\TextAt{8}{249.097}{364.029}{which}%
\TextAt{8}{282.584}{364.029}{is}%
\TextAt{8}{294.347}{364.029}{called}%
\TextAt{8}{327.511}{364.029}{the}%
\TextAt{12}{347.668}{364.029}{Newton}%
\TextAt{12}{397.963}{364.029}{polygon}%
\TextAt{8}{448.211}{364.029}{of}%
\TextAt{13}{461.541}{364.029}{f}%
\TextAt{9}{468.587}{364.029}{\Character{40}}%
\TextAt{13}{473.139}{364.029}{x}%
\TextAt{9}{479.791}{364.029}{\Character{41}}%
\TextAt{8}{484.344}{364.029}{.}%
\NativeDiagram{%
\path[draw=black,line width=0.399bp] (141.561,-497.072) -- (338.399,-497.072);%
\path[draw=black,line width=0.399bp,fill=black] (340.895,-497.072) -- (337.324,-495.726) -- (338.498,-497.072) -- (337.324,-498.418) -- cycle;%
\path[draw=black,line width=0.399bp] (153.75,-526.553) -- (153.75,-428.702);%
\path[draw=black,line width=0.399bp,fill=black] (153.75,-426.205) -- (152.404,-429.777) -- (153.75,-428.602) -- (155.096,-429.777) -- cycle;%
\path[draw=black,line width=0.797bp] (153.75,-474.394) -- (219.571,-521.564) -- (290.268,-512.039) -- (312.209,-432.441);%
\path[fill=black] (153.75,-474.394) (155.344,-474.394) .. controls (155.344,-473.514) and (154.63,-472.8) .. (153.75,-472.8) .. controls (152.87,-472.8) and (152.156,-473.514) .. (152.156,-474.394) .. controls (152.156,-475.275) and (152.87,-475.989) .. (153.75,-475.989) .. controls (154.63,-475.989) and (155.344,-475.275) .. (155.344,-474.394) -- cycle (153.75,-474.394);%
\path[fill=black] (175.69,-441.512) (177.284,-441.512) .. controls (177.284,-440.632) and (176.571,-439.918) .. (175.69,-439.918) .. controls (174.81,-439.918) and (174.096,-440.632) .. (174.096,-441.512) .. controls (174.096,-442.392) and (174.81,-443.106) .. (175.69,-443.106) .. controls (176.571,-443.106) and (177.284,-442.392) .. (177.284,-441.512) -- cycle (175.69,-441.512);%
\path[fill=black] (198.85,-474.394) (200.444,-474.394) .. controls (200.444,-473.514) and (199.73,-472.8) .. (198.85,-472.8) .. controls (197.97,-472.8) and (197.256,-473.514) .. (197.256,-474.394) .. controls (197.256,-475.275) and (197.97,-475.989) .. (198.85,-475.989) .. controls (199.73,-475.989) and (200.444,-475.275) .. (200.444,-474.394) -- cycle (198.85,-474.394);%
\path[fill=black] (219.571,-521.564) (221.165,-521.564) .. controls (221.165,-520.683) and (220.452,-519.97) .. (219.571,-519.97) .. controls (218.691,-519.97) and (217.977,-520.683) .. (217.977,-521.564) .. controls (217.977,-522.444) and (218.691,-523.158) .. (219.571,-523.158) .. controls (220.452,-523.158) and (221.165,-522.444) .. (221.165,-521.564) -- cycle (219.571,-521.564);%
\path[fill=black] (243.95,-507.957) (245.544,-507.957) .. controls (245.544,-507.077) and (244.83,-506.363) .. (243.95,-506.363) .. controls (243.069,-506.363) and (242.355,-507.077) .. (242.355,-507.957) .. controls (242.355,-508.837) and (243.069,-509.551) .. (243.95,-509.551) .. controls (244.83,-509.551) and (245.544,-508.837) .. (245.544,-507.957) -- cycle (243.95,-507.957);%
\path[fill=black] (262.233,-466.684) (263.827,-466.684) .. controls (263.827,-465.804) and (263.114,-465.09) .. (262.233,-465.09) .. controls (261.353,-465.09) and (260.639,-465.804) .. (260.639,-466.684) .. controls (260.639,-467.565) and (261.353,-468.278) .. (262.233,-468.278) .. controls (263.114,-468.278) and (263.827,-467.565) .. (263.827,-466.684) -- cycle (262.233,-466.684);%
\path[fill=black] (290.268,-512.039) (291.863,-512.039) .. controls (291.863,-511.159) and (291.149,-510.445) .. (290.268,-510.445) .. controls (289.388,-510.445) and (288.674,-511.159) .. (288.674,-512.039) .. controls (288.674,-512.92) and (289.388,-513.633) .. (290.268,-513.633) .. controls (291.149,-513.633) and (291.863,-512.92) .. (291.863,-512.039) -- cycle (290.268,-512.039);%
\path[fill=black] (297.582,-439.244) (299.176,-439.244) .. controls (299.176,-438.364) and (298.462,-437.65) .. (297.582,-437.65) .. controls (296.701,-437.65) and (295.988,-438.364) .. (295.988,-439.244) .. controls (295.988,-440.125) and (296.701,-440.838) .. (297.582,-440.838) .. controls (298.462,-440.838) and (299.176,-440.125) .. (299.176,-439.244) -- cycle (297.582,-439.244);%
\path[fill=black] (312.209,-432.441) (313.803,-432.441) .. controls (313.803,-431.561) and (313.089,-430.847) .. (312.209,-430.847) .. controls (311.328,-430.847) and (310.615,-431.561) .. (310.615,-432.441) .. controls (310.615,-433.321) and (311.328,-434.035) .. (312.209,-434.035) .. controls (313.089,-434.035) and (313.803,-433.321) .. (313.803,-432.441) -- cycle (312.209,-432.441);%
}%
\TextAt{9}{175.294}{296.617}{\Character{40}0}%
\TextAt{13}{185.699}{296.617}{\Character{59}}%
\TextAt{13}{190.923}{296.617}{v}%
\TextAt{9}{197.032}{296.617}{\Character{40}}%
\TextAt{13}{201.584}{296.617}{a}%
\TextAt{17}{207.729}{294.824}{0}%
\TextAt{9}{212.461}{296.617}{\Character{41}\Character{41}}%
\TextAt{9}{270.04}{239.286}{\Character{40}}%
\TextAt{13}{274.592}{239.286}{i\Character{59}}%
\TextAt{13}{283.81}{239.286}{v}%
\TextAt{9}{289.918}{239.286}{\Character{40}}%
\TextAt{13}{294.471}{239.286}{a}%
\TextAt{21}{300.616}{237.493}{i}%
\TextAt{9}{303.997}{239.286}{\Character{41}\Character{41}}%
\TextAt{9}{366.452}{248.81}{\Character{40}}%
\TextAt{13}{371.004}{248.81}{j\Character{59}}%
\TextAt{13}{381.094}{248.81}{v}%
\TextAt{9}{387.202}{248.81}{\Character{40}}%
\TextAt{13}{391.754}{248.81}{a}%
\TextAt{21}{397.899}{247.017}{j}%
\TextAt{9}{402.281}{248.81}{\Character{41}\Character{41}}%
\TextAt{9}{388.393}{338.57}{\Character{40}}%
\TextAt{13}{392.945}{338.57}{n\Character{59}}%
\TextAt{13}{405.151}{338.57}{v}%
\TextAt{9}{411.265}{338.57}{\Character{40}}%
\TextAt{13}{415.818}{338.57}{a}%
\TextAt{21}{421.963}{336.777}{n}%
\TextAt{9}{427.599}{338.57}{\Character{41}\Character{41}}%
\TextAt{8}{79.226}{206.744}{The}%
\TextAt{8}{103.956}{206.744}{polygon}%
\TextAt{8}{148.867}{206.744}{consists}%
\TextAt{8}{192.653}{206.744}{of}%
\TextAt{8}{206.659}{206.744}{a}%
\TextAt{8}{217.091}{206.744}{sequence}%
\TextAt{8}{266.269}{206.744}{of}%
\TextAt{8}{280.276}{206.744}{line}%
\TextAt{8}{303.057}{206.744}{segments}%
\TextAt{13}{353.61}{206.744}{S}%
\TextAt{17}{360.81}{204.951}{1}%
\TextAt{13}{365.542}{206.744}{\Character{59}}%
\TextAt{13}{370.766}{206.744}{S}%
\TextAt{17}{377.985}{204.951}{2}%
\TextAt{13}{382.718}{206.744}{\Character{59}}%
\TextAt{13}{387.942}{206.744}{\Character{58}}%
\TextAt{13}{393.179}{206.744}{\Character{58}}%
\TextAt{13}{398.403}{206.744}{\Character{58}}%
\TextAt{8}{406.275}{206.744}{whose}%
\TextAt{8}{441.478}{206.744}{slopes}%
\TextAt{8}{476.418}{206.744}{are}%
\TextAt{8}{496.605}{206.744}{strictly}%
\TextAt{8}{79.226}{189.698}{increasing,}%
\TextAt{8}{137.173}{189.698}{and}%
\TextAt{8}{159.935}{189.698}{which}%
\TextAt{8}{193.422}{189.698}{are}%
\TextAt{8}{212.933}{189.698}{subject}%
\TextAt{8}{253.305}{189.698}{to}%
\TextAt{8}{267.616}{189.698}{the}%
\TextAt{8}{287.772}{189.698}{following}%
\TextAt{12}{79.226}{165.162}{Proposition}%
\TextAt{12}{151.406}{165.162}{(6.3).}%
\TextAt{8}{186.259}{165.162}{Let}%
\TextAt{13}{206.766}{165.162}{f}%
\TextAt{9}{213.813}{165.162}{\Character{40}}%
\TextAt{13}{218.365}{165.162}{x}%
\TextAt{9}{225.017}{165.162}{\Character{41}}%
\TextAt{9}{232.872}{165.162}{\Character{61}}%
\TextAt{13}{245.316}{165.162}{a}%
\TextAt{17}{251.461}{163.369}{0}%
\TextAt{9}{257.905}{165.162}{\Character{43}}%
\TextAt{13}{268.721}{165.162}{a}%
\TextAt{17}{274.866}{163.369}{1}%
\TextAt{13}{279.598}{165.162}{x}%
\TextAt{9}{287.962}{165.162}{\Character{43}}%
\TextAt{16}{298.778}{165.162}{\Character{1}}%
\TextAt{16}{304.074}{165.162}{\Character{1}}%
\TextAt{16}{309.382}{165.162}{\Character{1}}%
\TextAt{9}{314.437}{165.162}{\Character{43}}%
\TextAt{13}{325.254}{165.162}{a}%
\TextAt{21}{331.399}{163.369}{n}%
\TextAt{13}{337.035}{165.162}{x}%
\TextAt{21}{343.687}{169.5}{n}%
\TextAt{8}{349.323}{165.162}{,}%
\TextAt{13}{356.103}{165.162}{a}%
\TextAt{17}{362.248}{163.369}{0}%
\TextAt{13}{366.98}{165.162}{a}%
\TextAt{21}{373.125}{163.369}{n}%
\NativeMathAt{382.082}{165.162}{11.955}{\neq}%
\TextAt{9}{394.492}{165.162}{0}%
\TextAt{8}{400.361}{165.162}{,}%
\TextAt{8}{407.14}{165.162}{be}%
\TextAt{8}{422.602}{165.162}{a}%
\TextAt{8}{431.896}{165.162}{polynomial}%
\TextAt{8}{491.927}{165.162}{over}%
\TextAt{8}{516.51}{165.162}{the}%
\TextAt{8}{79.226}{148.116}{field}%
\TextAt{13}{104.237}{148.116}{K}%
\TextAt{8}{115.055}{148.116}{,}%
\TextAt{13}{121.927}{148.116}{v}%
\TextAt{8}{131.567}{148.116}{an}%
\TextAt{8}{147.48}{148.116}{exponential}%
\TextAt{8}{209.898}{148.116}{valuation}%
\TextAt{8}{260.625}{148.116}{of}%
\TextAt{13}{273.609}{148.116}{K}%
\TextAt{8}{284.426}{148.116}{,}%
\TextAt{8}{291.299}{148.116}{and}%
\TextAt{13}{313.715}{148.116}{w}%
\TextAt{8}{325.956}{148.116}{an}%
\TextAt{8}{341.869}{148.116}{extension}%
\TextAt{8}{393.289}{148.116}{to}%
\TextAt{8}{407.254}{148.116}{the}%
\TextAt{8}{427.064}{148.116}{splitting}%
\TextAt{8}{472.96}{148.116}{field}%
\TextAt{13}{497.971}{148.116}{L}%
\TextAt{8}{509.487}{148.116}{of}%
\TextAt{13}{522.471}{148.116}{f}%
\TextAt{8}{529.518}{148.116}{.}%
\TextAt{8}{97.159}{128.081}{If}%
\TextAt{9}{107.791}{128.081}{\Character{40}}%
\TextAt{13}{112.344}{128.081}{r}%
\TextAt{13}{117.294}{128.081}{\Character{59}}%
\TextAt{13}{122.518}{128.081}{v}%
\TextAt{9}{128.626}{128.081}{\Character{40}}%
\TextAt{13}{133.179}{128.081}{a}%
\TextAt{21}{139.323}{126.288}{r}%
\TextAt{9}{143.878}{128.081}{\Character{41}\Character{41}}%
\NativeMathAt{156.304}{128.081}{11.955}{\leftrightarrow}%
\TextAt{9}{171.58}{128.081}{\Character{40}}%
\TextAt{13}{176.132}{128.081}{s\Character{59}}%
\TextAt{13}{186.868}{128.081}{v}%
\TextAt{9}{192.978}{128.081}{\Character{40}}%
\TextAt{13}{197.531}{128.081}{a}%
\TextAt{21}{203.676}{126.288}{s}%
\TextAt{9}{208.09}{128.081}{\Character{41}\Character{41}}%
\TextAt{8}{220.032}{128.081}{is}%
\TextAt{8}{230.736}{128.081}{a}%
\TextAt{8}{239.432}{128.081}{line}%
\TextAt{8}{260.477}{128.081}{segment}%
\TextAt{8}{304.68}{128.081}{of}%
\TextAt{8}{316.95}{128.081}{slope}%
\TextAt{16}{345.539}{128.081}{\Character{0}}%
\TextAt{13}{354.838}{128.081}{m}%
\TextAt{8}{367.915}{128.081}{occurring}%
\TextAt{8}{418.561}{128.081}{in}%
\TextAt{8}{431.154}{128.081}{the}%
\TextAt{8}{450.251}{128.081}{Newton}%
\TextAt{8}{492.433}{128.081}{polygon}%
\TextAt{8}{79.226}{111.035}{of}%
\TextAt{13}{92.556}{111.035}{f}%
\TextAt{8}{99.603}{111.035}{,}%
\TextAt{8}{106.752}{111.035}{then}%
\TextAt{13}{133.412}{111.035}{f}%
\TextAt{9}{140.458}{111.035}{\Character{40}}%
\TextAt{13}{145.011}{111.035}{x}%
\TextAt{9}{151.663}{111.035}{\Character{41}}%
\TextAt{8}{160.112}{111.035}{has}%
\TextAt{8}{180.986}{111.035}{precisely}%
\TextAt{13}{228.843}{111.035}{s}%
\TextAt{16}{237.013}{111.035}{\Character{0}}%
\TextAt{13}{248.969}{111.035}{r}%
\TextAt{8}{258.466}{111.035}{roots}%
\TextAt{13}{288.127}{111.035}{\Character{11}}%
\TextAt{17}{295.573}{109.242}{1}%
\TextAt{13}{300.305}{111.035}{\Character{59}}%
\TextAt{13}{305.529}{111.035}{\Character{58}}%
\TextAt{13}{310.766}{111.035}{\Character{58}}%
\TextAt{13}{315.99}{111.035}{\Character{58}}%
\TextAt{13}{321.227}{111.035}{\Character{59}}%
\TextAt{13}{326.463}{111.035}{\Character{11}}%
\TextAt{21}{333.971}{109.242}{s}%
\TextAt{3}{337.887}{109.242}{\Character{0}}%
\TextAt{21}{344.474}{109.242}{r}%
\TextAt{8}{352.926}{111.035}{of}%
\TextAt{8}{366.256}{111.035}{value}%
\TextAt{13}{233.525}{87.2}{w}%
\TextAt{9}{242.214}{87.2}{\Character{40}}%
\TextAt{13}{246.766}{87.2}{\Character{11}}%
\TextAt{17}{254.212}{85.407}{1}%
\TextAt{9}{258.944}{87.2}{\Character{41}}%
\TextAt{9}{266.799}{87.2}{\Character{61}}%
\TextAt{16}{279.243}{87.2}{\Character{1}}%
\TextAt{16}{284.539}{87.2}{\Character{1}}%
\TextAt{16}{289.847}{87.2}{\Character{1}}%
\TextAt{9}{296.511}{87.2}{\Character{61}}%
\TextAt{13}{308.937}{87.2}{w}%
\TextAt{9}{317.626}{87.2}{\Character{40}}%
\TextAt{13}{322.179}{87.2}{\Character{11}}%
\TextAt{21}{329.624}{85.407}{s}%
\TextAt{3}{333.54}{85.407}{\Character{0}}%
\TextAt{21}{340.127}{85.407}{r}%
\TextAt{9}{344.682}{87.2}{\Character{41}}%
\TextAt{9}{352.537}{87.2}{\Character{61}}%
\TextAt{13}{364.98}{87.2}{m\Character{58}}%
\TextAt{8}{297.211}{43.033}{132}%
\Destination{78.23}{809.815}{nav.page.132}%
\Destination{79.36}{712.955}{nav.print.144}%
\Destination{79.36}{775.155}{nav.number.2.6.1}%
\Destination{79.36}{711.955}{nav.number.2.6.2}%
\Destination{79.36}{179.155}{nav.number.2.6.3}%
\Hotlink{102.273}{570.155}{25.078}{11.047}{nav.number.2.4.8}%
\Hotlink{362.456}{741.9}{21.722}{11.047}{nav.number.2.4.6}%
}

% Page 145 (body page 133)
\Page{%
\TextAt{2}{79.226}{796.943}{II.}%
\TextAt{2}{94.073}{796.943}{The}%
\TextAt{2}{116.492}{796.943}{Theory}%
\TextAt{2}{154.401}{796.943}{of}%
\TextAt{2}{166.826}{796.943}{Valuations}%
\Vector{q 0 G 0 g 0.299 w 7.226 18.729 m 460.769 18.729 l S Q}%
\TextAt{19}{79.226}{761.124}{Proof.}%
\TextAt{8}{116}{761.124}{Dividing}%
\TextAt{8}{164.596}{761.124}{by}%
\TextAt{13}{181.81}{761.124}{a}%
\TextAt{21}{187.955}{759.33}{n}%
\TextAt{8}{198.444}{761.124}{only}%
\TextAt{8}{225.09}{761.124}{shifts}%
\TextAt{8}{257.057}{761.124}{the}%
\TextAt{8}{278.168}{761.124}{polygon}%
\TextAt{8}{323.357}{761.124}{up}%
\TextAt{8}{341.216}{761.124}{or}%
\TextAt{8}{356.482}{761.124}{down.}%
\TextAt{8}{394.789}{761.124}{Thus}%
\TextAt{8}{425.393}{761.124}{we}%
\TextAt{8}{443.575}{761.124}{may}%
\TextAt{8}{469.899}{761.124}{assume}%
\TextAt{8}{511.298}{761.124}{that}%
\TextAt{13}{79.226}{744.078}{a}%
\TextAt{21}{85.371}{742.284}{n}%
\TextAt{9}{94.328}{744.078}{\Character{61}}%
\TextAt{9}{106.738}{744.078}{1}%
\TextAt{8}{112.607}{744.078}{.}%
\TextAt{8}{121.055}{744.078}{We}%
\TextAt{8}{141.2}{744.078}{number}%
\TextAt{8}{183.796}{744.078}{the}%
\TextAt{8}{203.952}{744.078}{roots}%
\TextAt{13}{233.613}{744.078}{\Character{11}}%
\TextAt{17}{241.059}{742.284}{1}%
\TextAt{13}{245.791}{744.078}{\Character{59}}%
\TextAt{13}{251.015}{744.078}{\Character{58}}%
\TextAt{13}{256.252}{744.078}{\Character{58}}%
\TextAt{13}{261.476}{744.078}{\Character{58}}%
\TextAt{13}{266.713}{744.078}{\Character{59}}%
\TextAt{13}{271.949}{744.078}{\Character{11}}%
\TextAt{21}{279.457}{742.284}{n}%
\NativeMathAt{288.414}{744.078}{11.955}{\in}%
\TextAt{13}{299.705}{744.078}{L}%
\TextAt{8}{311.567}{744.078}{of}%
\TextAt{13}{324.897}{744.078}{f}%
\TextAt{8}{335.841}{744.078}{in}%
\TextAt{8}{349.494}{744.078}{such}%
\TextAt{8}{375.891}{744.078}{a}%
\TextAt{8}{385.646}{744.078}{way}%
\TextAt{8}{409.389}{744.078}{that}%
\TextAt{13}{238.373}{721.198}{w}%
\TextAt{9}{247.063}{721.198}{\Character{40}}%
\TextAt{13}{251.615}{721.198}{\Character{11}}%
\TextAt{17}{259.061}{719.404}{1}%
\TextAt{9}{263.793}{721.198}{\Character{41}}%
\TextAt{9}{271.648}{721.198}{\Character{61}}%
\TextAt{16}{284.092}{721.198}{\Character{1}}%
\TextAt{16}{289.388}{721.198}{\Character{1}}%
\TextAt{16}{294.696}{721.198}{\Character{1}}%
\TextAt{9}{301.36}{721.198}{\Character{61}}%
\TextAt{13}{313.786}{721.198}{w}%
\TextAt{9}{322.475}{721.198}{\Character{40}}%
\TextAt{13}{327.028}{721.198}{\Character{11}}%
\TextAt{21}{334.473}{719.404}{s}%
\TextAt{24}{338.389}{718.297}{1}%
\TextAt{9}{343.038}{721.198}{\Character{41}}%
\TextAt{9}{350.893}{721.198}{\Character{61}}%
\TextAt{13}{363.337}{721.198}{m}%
\TextAt{17}{373.576}{719.404}{1}%
\TextAt{13}{378.308}{721.198}{\Character{59}}%
\TextAt{13}{223.72}{701.163}{w}%
\TextAt{9}{232.409}{701.163}{\Character{40}}%
\TextAt{13}{236.962}{701.163}{\Character{11}}%
\TextAt{21}{244.407}{699.37}{s}%
\TextAt{24}{248.323}{698.263}{1}%
\TextAt{17}{252.474}{699.37}{\Character{43}1}%
\TextAt{9}{263.793}{701.163}{\Character{41}}%
\TextAt{9}{271.648}{701.163}{\Character{61}}%
\TextAt{16}{284.092}{701.163}{\Character{1}}%
\TextAt{16}{289.388}{701.163}{\Character{1}}%
\TextAt{16}{294.696}{701.163}{\Character{1}}%
\TextAt{9}{301.36}{701.163}{\Character{61}}%
\TextAt{13}{313.786}{701.163}{w}%
\TextAt{9}{322.475}{701.163}{\Character{40}}%
\TextAt{13}{327.028}{701.163}{\Character{11}}%
\TextAt{21}{334.473}{699.37}{s}%
\TextAt{24}{338.389}{698.263}{2}%
\TextAt{9}{343.038}{701.163}{\Character{41}}%
\TextAt{9}{350.893}{701.163}{\Character{61}}%
\TextAt{13}{363.337}{701.163}{m}%
\TextAt{17}{373.576}{699.37}{2}%
\TextAt{13}{378.308}{701.163}{\Character{59}}%
\TextAt{16}{274.235}{681.128}{\Character{1}}%
\TextAt{16}{279.531}{681.128}{\Character{1}}%
\TextAt{16}{284.839}{681.128}{\Character{1}}%
\TextAt{16}{314.058}{681.128}{\Character{1}}%
\TextAt{16}{319.354}{681.128}{\Character{1}}%
\TextAt{16}{324.662}{681.128}{\Character{1}}%
\TextAt{16}{353.881}{681.128}{\Character{1}}%
\TextAt{16}{359.177}{681.128}{\Character{1}}%
\TextAt{16}{364.485}{681.128}{\Character{1}}%
\TextAt{13}{224.522}{661.093}{w}%
\TextAt{9}{233.212}{661.093}{\Character{40}}%
\TextAt{13}{237.764}{661.093}{\Character{11}}%
\TextAt{21}{245.21}{659.3}{s}%
\TextAt{22}{249.126}{658.304}{t}%
\TextAt{17}{252.474}{659.3}{\Character{43}1}%
\TextAt{9}{263.793}{661.093}{\Character{41}}%
\TextAt{9}{271.648}{661.093}{\Character{61}}%
\TextAt{16}{284.092}{661.093}{\Character{1}}%
\TextAt{16}{289.388}{661.093}{\Character{1}}%
\TextAt{16}{294.696}{661.093}{\Character{1}}%
\TextAt{9}{301.36}{661.093}{\Character{61}}%
\TextAt{13}{313.786}{661.093}{w}%
\TextAt{9}{322.475}{661.093}{\Character{40}}%
\TextAt{13}{327.028}{661.093}{\Character{11}}%
\TextAt{21}{334.473}{659.3}{n}%
\TextAt{9}{340.11}{661.093}{\Character{41}}%
\TextAt{9}{347.965}{661.093}{\Character{61}}%
\TextAt{13}{360.408}{661.093}{m}%
\TextAt{21}{370.647}{659.3}{t}%
\TextAt{17}{373.705}{659.3}{\Character{43}1}%
\TextAt{13}{385.024}{661.093}{\Character{59}}%
\TextAt{8}{79.226}{639.054}{where}%
\TextAt{13}{114.531}{639.054}{m}%
\TextAt{17}{124.771}{637.261}{1}%
\TextAt{13}{135.373}{639.054}{\Character{60}}%
\TextAt{13}{150.341}{639.054}{m}%
\TextAt{17}{160.587}{637.261}{2}%
\TextAt{13}{171.19}{639.054}{\Character{60}}%
\TextAt{16}{186.165}{639.054}{\Character{1}}%
\TextAt{16}{191.461}{639.054}{\Character{1}}%
\TextAt{16}{196.769}{639.054}{\Character{1}}%
\TextAt{13}{205.983}{639.054}{\Character{60}}%
\TextAt{13}{220.951}{639.054}{m}%
\TextAt{21}{231.197}{637.261}{t}%
\TextAt{17}{234.255}{637.261}{\Character{43}1}%
\TextAt{8}{245.574}{639.054}{.}%
\TextAt{8}{258.51}{639.054}{Viewing}%
\TextAt{8}{305.197}{639.054}{the}%
\TextAt{8}{326.849}{639.054}{coefficients}%
\TextAt{13}{387.583}{639.054}{a}%
\TextAt{21}{393.728}{637.261}{i}%
\TextAt{8}{402.503}{639.054}{as}%
\TextAt{8}{418.369}{639.054}{elementary}%
\TextAt{8}{479.7}{639.054}{symmetric}%
\TextAt{8}{79.226}{622.008}{functions}%
\TextAt{8}{129.689}{622.008}{of}%
\TextAt{8}{143.019}{622.008}{the}%
\TextAt{8}{163.175}{622.008}{roots}%
\TextAt{13}{192.836}{622.008}{\Character{11}}%
\TextAt{21}{200.282}{620.215}{j}%
\TextAt{8}{204.664}{622.008}{,}%
\TextAt{8}{211.813}{622.008}{we}%
\TextAt{8}{229.041}{622.008}{immediately}%
\TextAt{8}{295.703}{622.008}{find}%
\TextAt{13}{214.13}{598.224}{v}%
\TextAt{9}{220.219}{598.224}{\Character{40}}%
\TextAt{13}{224.771}{598.224}{a}%
\TextAt{21}{230.916}{596.43}{n}%
\TextAt{9}{236.552}{598.224}{\Character{41}}%
\TextAt{9}{244.407}{598.224}{\Character{61}}%
\TextAt{13}{256.851}{598.224}{v}%
\TextAt{9}{262.939}{598.224}{\Character{40}1\Character{41}}%
\TextAt{9}{281.183}{598.224}{\Character{61}}%
\TextAt{9}{293.604}{598.224}{0}%
\TextAt{13}{299.496}{598.224}{\Character{59}}%
\TextAt{13}{203.31}{578.189}{v}%
\TextAt{9}{209.398}{578.189}{\Character{40}}%
\TextAt{13}{213.95}{578.189}{a}%
\TextAt{21}{220.095}{576.396}{n}%
\TextAt{3}{225.233}{576.396}{\Character{0}}%
\TextAt{17}{231.82}{576.396}{1}%
\TextAt{9}{236.552}{578.189}{\Character{41}}%
\NativeMathAt{244.425}{578.189}{11.955}{\geq}%
\TextAt{8}{257.045}{578.189}{min}%
\TextAt{21}{265.359}{570.927}{i}%
\TextAt{16}{276.556}{578.189}{\Character{102}}%
\TextAt{13}{282.533}{578.189}{w}%
\TextAt{9}{291.223}{578.189}{\Character{40}}%
\TextAt{13}{295.775}{578.189}{\Character{11}}%
\TextAt{21}{303.221}{576.396}{i}%
\TextAt{9}{306.602}{578.189}{\Character{41}}%
\TextAt{16}{311.154}{578.189}{\Character{103}}%
\TextAt{9}{320.453}{578.189}{\Character{61}}%
\TextAt{13}{332.878}{578.189}{m}%
\TextAt{17}{343.117}{576.396}{1}%
\TextAt{13}{347.85}{578.189}{\Character{59}}%
\TextAt{13}{203.31}{553.815}{v}%
\TextAt{9}{209.398}{553.815}{\Character{40}}%
\TextAt{13}{213.95}{553.815}{a}%
\TextAt{21}{220.095}{552.021}{n}%
\TextAt{3}{225.233}{552.021}{\Character{0}}%
\TextAt{17}{231.82}{552.021}{2}%
\TextAt{9}{236.552}{553.815}{\Character{41}}%
\NativeMathAt{244.425}{553.815}{11.955}{\geq}%
\TextAt{8}{257.045}{553.815}{min}%
\TextAt{21}{262.24}{546.553}{i\Character{59}j}%
\TextAt{16}{276.556}{553.815}{\Character{102}}%
\TextAt{13}{282.533}{553.815}{w}%
\TextAt{9}{291.223}{553.815}{\Character{40}}%
\TextAt{13}{295.775}{553.815}{\Character{11}}%
\TextAt{21}{303.221}{552.021}{i}%
\TextAt{13}{306.602}{553.815}{\Character{11}}%
\TextAt{21}{314.048}{552.021}{j}%
\TextAt{9}{318.43}{553.815}{\Character{41}}%
\TextAt{16}{322.982}{553.815}{\Character{103}}%
\TextAt{9}{332.28}{553.815}{\Character{61}}%
\TextAt{9}{344.69}{553.815}{2}%
\TextAt{13}{350.559}{553.815}{m}%
\TextAt{17}{360.798}{552.021}{1}%
\TextAt{13}{365.531}{553.815}{\Character{59}}%
\TextAt{16}{246.994}{527.891}{\Character{1}}%
\TextAt{16}{252.29}{527.891}{\Character{1}}%
\TextAt{16}{257.598}{527.891}{\Character{1}}%
\TextAt{16}{286.817}{527.891}{\Character{1}}%
\TextAt{16}{292.113}{527.891}{\Character{1}}%
\TextAt{16}{297.421}{527.891}{\Character{1}}%
\TextAt{16}{326.639}{527.891}{\Character{1}}%
\TextAt{16}{331.936}{527.891}{\Character{1}}%
\TextAt{16}{337.244}{527.891}{\Character{1}}%
\TextAt{13}{199.477}{507.856}{v}%
\TextAt{9}{205.565}{507.856}{\Character{40}}%
\TextAt{13}{210.118}{507.856}{a}%
\TextAt{21}{216.262}{506.063}{n}%
\TextAt{3}{221.401}{506.063}{\Character{0}}%
\TextAt{21}{227.987}{506.063}{s}%
\TextAt{24}{231.903}{504.956}{1}%
\TextAt{9}{236.552}{507.856}{\Character{41}}%
\TextAt{9}{244.407}{507.856}{\Character{61}}%
\TextAt{8}{262.006}{507.856}{min}%
\TextAt{21}{256.851}{500.595}{i}%
\TextAt{24}{259.734}{499.488}{1}%
\TextAt{21}{263.885}{500.594}{\Character{59}\Character{58}\Character{58}\Character{58}\Character{59}i}%
\TextAt{22}{278.53}{499.598}{s}%
\TextAt{24}{282.024}{497.805}{1}%
\TextAt{16}{286.673}{507.856}{\Character{102}}%
\TextAt{13}{292.65}{507.856}{w}%
\TextAt{9}{301.34}{507.856}{\Character{40}}%
\TextAt{13}{305.892}{507.856}{\Character{11}}%
\TextAt{21}{313.338}{506.063}{i}%
\TextAt{24}{316.221}{504.956}{1}%
\TextAt{16}{322.863}{507.856}{\Character{1}}%
\TextAt{16}{328.159}{507.856}{\Character{1}}%
\TextAt{16}{333.467}{507.856}{\Character{1}}%
\TextAt{13}{338.803}{507.856}{\Character{11}}%
\TextAt{21}{346.249}{506.063}{i}%
\TextAt{22}{349.132}{505.066}{s}%
\TextAt{24}{352.626}{503.273}{1}%
\TextAt{9}{357.773}{507.856}{\Character{41}}%
\TextAt{16}{362.325}{507.856}{\Character{103}}%
\TextAt{9}{371.624}{507.856}{\Character{61}}%
\TextAt{13}{384.049}{507.856}{s}%
\TextAt{17}{389.563}{506.063}{1}%
\TextAt{13}{394.295}{507.856}{m}%
\TextAt{17}{404.535}{506.063}{1}%
\TextAt{13}{409.267}{507.856}{\Character{59}}%
\TextAt{8}{79.226}{479.684}{the}%
\TextAt{8}{99.383}{479.684}{latter}%
\TextAt{8}{131.255}{479.684}{because}%
\TextAt{8}{174.557}{479.684}{the}%
\TextAt{8}{194.713}{479.684}{value}%
\TextAt{8}{224.96}{479.684}{of}%
\TextAt{8}{238.29}{479.684}{the}%
\TextAt{8}{258.446}{479.684}{term}%
\TextAt{13}{286.409}{479.684}{\Character{11}}%
\TextAt{17}{293.855}{477.891}{1}%
\TextAt{16}{300.58}{479.684}{\Character{1}}%
\TextAt{16}{305.876}{479.684}{\Character{1}}%
\TextAt{16}{311.184}{479.684}{\Character{1}}%
\TextAt{13}{316.52}{479.684}{\Character{11}}%
\TextAt{21}{323.965}{477.891}{s}%
\TextAt{24}{327.881}{476.784}{1}%
\TextAt{8}{336.428}{479.684}{is}%
\TextAt{8}{348.192}{479.684}{smaller}%
\TextAt{8}{388.576}{479.684}{than}%
\TextAt{8}{415.894}{479.684}{that}%
\TextAt{8}{441.263}{479.684}{of}%
\TextAt{8}{454.593}{479.684}{all}%
\TextAt{8}{470.852}{479.684}{the}%
\TextAt{8}{491.008}{479.684}{others,}%
\TextAt{13}{157.267}{456.342}{v}%
\TextAt{9}{163.356}{456.342}{\Character{40}}%
\TextAt{13}{167.908}{456.342}{a}%
\TextAt{21}{174.053}{454.549}{n}%
\TextAt{3}{179.191}{454.549}{\Character{0}}%
\TextAt{21}{185.778}{454.549}{s}%
\TextAt{24}{189.693}{453.442}{1}%
\TextAt{3}{193.844}{454.549}{\Character{0}}%
\TextAt{17}{200.431}{454.549}{1}%
\TextAt{9}{205.163}{456.342}{\Character{41}}%
\NativeMathAt{213.036}{456.342}{11.955}{\geq}%
\TextAt{8}{235.433}{456.342}{min}%
\TextAt{21}{225.656}{449.081}{i}%
\TextAt{24}{228.539}{447.974}{1}%
\TextAt{21}{232.69}{449.081}{\Character{59}\Character{58}\Character{58}\Character{58}\Character{59}i}%
\TextAt{22}{247.335}{447.974}{s}%
\TextAt{24}{250.829}{446.181}{1}%
\TextAt{24}{254.98}{447.974}{\Character{43}1}%
\TextAt{16}{264.721}{456.342}{\Character{102}}%
\TextAt{13}{270.698}{456.342}{w}%
\TextAt{9}{279.388}{456.342}{\Character{40}}%
\TextAt{13}{283.94}{456.342}{\Character{11}}%
\TextAt{21}{291.386}{454.549}{i}%
\TextAt{24}{294.269}{453.442}{1}%
\TextAt{16}{300.911}{456.342}{\Character{1}}%
\TextAt{16}{306.207}{456.342}{\Character{1}}%
\TextAt{16}{311.515}{456.342}{\Character{1}}%
\TextAt{13}{316.851}{456.342}{\Character{11}}%
\TextAt{21}{324.296}{454.549}{i}%
\TextAt{22}{327.18}{453.442}{s}%
\TextAt{24}{330.673}{451.649}{1}%
\TextAt{24}{334.824}{453.442}{\Character{43}1}%
\TextAt{9}{345.064}{456.342}{\Character{41}}%
\TextAt{16}{349.616}{456.342}{\Character{103}}%
\TextAt{9}{358.914}{456.342}{\Character{61}}%
\TextAt{13}{371.34}{456.342}{s}%
\TextAt{17}{376.854}{454.549}{1}%
\TextAt{13}{381.586}{456.342}{m}%
\TextAt{17}{391.826}{454.549}{1}%
\TextAt{9}{399.215}{456.342}{\Character{43}}%
\TextAt{13}{410.976}{456.342}{m}%
\TextAt{17}{421.215}{454.549}{2}%
\TextAt{13}{425.947}{456.342}{\Character{59}}%
\TextAt{13}{157.267}{429.068}{v}%
\TextAt{9}{163.356}{429.068}{\Character{40}}%
\TextAt{13}{167.908}{429.068}{a}%
\TextAt{21}{174.053}{427.275}{n}%
\TextAt{3}{179.191}{427.275}{\Character{0}}%
\TextAt{21}{185.778}{427.275}{s}%
\TextAt{24}{189.693}{426.168}{1}%
\TextAt{3}{193.844}{427.275}{\Character{0}}%
\TextAt{17}{200.431}{427.275}{2}%
\TextAt{9}{205.163}{429.068}{\Character{41}}%
\NativeMathAt{213.036}{429.068}{11.955}{\geq}%
\TextAt{8}{235.433}{429.068}{min}%
\TextAt{21}{225.656}{421.806}{i}%
\TextAt{24}{228.539}{420.699}{1}%
\TextAt{21}{232.69}{421.806}{\Character{59}\Character{58}\Character{58}\Character{58}\Character{59}i}%
\TextAt{22}{247.335}{420.699}{s}%
\TextAt{24}{250.829}{418.906}{1}%
\TextAt{24}{254.98}{420.699}{\Character{43}2}%
\TextAt{16}{264.721}{429.068}{\Character{102}}%
\TextAt{13}{270.698}{429.068}{w}%
\TextAt{9}{279.388}{429.068}{\Character{40}}%
\TextAt{13}{283.94}{429.068}{\Character{11}}%
\TextAt{21}{291.386}{427.275}{i}%
\TextAt{24}{294.269}{426.168}{1}%
\TextAt{16}{300.911}{429.068}{\Character{1}}%
\TextAt{16}{306.207}{429.068}{\Character{1}}%
\TextAt{16}{311.515}{429.068}{\Character{1}}%
\TextAt{13}{316.851}{429.068}{\Character{11}}%
\TextAt{21}{324.296}{427.275}{i}%
\TextAt{22}{327.18}{426.168}{s}%
\TextAt{24}{330.673}{424.374}{1}%
\TextAt{24}{334.824}{426.168}{\Character{43}2}%
\TextAt{9}{345.064}{429.068}{\Character{41}}%
\TextAt{16}{349.616}{429.068}{\Character{103}}%
\TextAt{9}{358.914}{429.068}{\Character{61}}%
\TextAt{13}{371.34}{429.068}{s}%
\TextAt{17}{376.854}{427.275}{1}%
\TextAt{13}{381.586}{429.068}{m}%
\TextAt{17}{391.826}{427.275}{1}%
\TextAt{9}{399.215}{429.068}{\Character{43}}%
\TextAt{9}{410.967}{429.068}{2}%
\TextAt{13}{416.829}{429.068}{m}%
\TextAt{17}{427.068}{427.275}{2}%
\TextAt{13}{431.8}{429.068}{\Character{59}}%
\TextAt{16}{215.606}{401.794}{\Character{1}}%
\TextAt{16}{220.902}{401.794}{\Character{1}}%
\TextAt{16}{226.21}{401.794}{\Character{1}}%
\TextAt{16}{255.417}{401.794}{\Character{1}}%
\TextAt{16}{260.725}{401.794}{\Character{1}}%
\TextAt{16}{266.033}{401.794}{\Character{1}}%
\TextAt{16}{295.24}{401.794}{\Character{1}}%
\TextAt{16}{300.548}{401.794}{\Character{1}}%
\TextAt{16}{305.856}{401.794}{\Character{1}}%
\TextAt{13}{168.088}{381.759}{v}%
\TextAt{9}{174.176}{381.759}{\Character{40}}%
\TextAt{13}{178.729}{381.759}{a}%
\TextAt{21}{184.874}{379.965}{n}%
\TextAt{3}{190.012}{379.965}{\Character{0}}%
\TextAt{21}{196.598}{379.965}{s}%
\TextAt{24}{200.514}{378.859}{2}%
\TextAt{9}{205.163}{381.759}{\Character{41}}%
\TextAt{9}{213.018}{381.759}{\Character{61}}%
\TextAt{8}{230.618}{381.759}{min}%
\TextAt{21}{225.462}{374.497}{i}%
\TextAt{24}{228.345}{373.39}{1}%
\TextAt{21}{232.496}{374.497}{\Character{59}\Character{58}\Character{58}\Character{58}\Character{59}i}%
\TextAt{22}{247.141}{373.501}{s}%
\TextAt{24}{250.635}{371.708}{2}%
\TextAt{16}{255.284}{381.759}{\Character{102}}%
\TextAt{13}{261.262}{381.759}{w}%
\TextAt{9}{269.951}{381.759}{\Character{40}}%
\TextAt{13}{274.503}{381.759}{\Character{11}}%
\TextAt{21}{281.949}{379.966}{i}%
\TextAt{24}{284.832}{378.859}{1}%
\TextAt{16}{291.474}{381.759}{\Character{1}}%
\TextAt{16}{296.77}{381.759}{\Character{1}}%
\TextAt{16}{302.078}{381.759}{\Character{1}}%
\TextAt{13}{307.414}{381.759}{\Character{11}}%
\TextAt{21}{314.86}{379.966}{i}%
\TextAt{22}{317.743}{378.969}{s}%
\TextAt{24}{321.237}{377.176}{2}%
\TextAt{9}{326.384}{381.759}{\Character{41}}%
\TextAt{16}{330.936}{381.759}{\Character{103}}%
\TextAt{9}{340.235}{381.759}{\Character{61}}%
\TextAt{13}{352.66}{381.759}{s}%
\TextAt{17}{358.174}{379.966}{1}%
\TextAt{13}{362.906}{381.759}{m}%
\TextAt{17}{373.146}{379.966}{1}%
\TextAt{9}{380.535}{381.759}{\Character{43}}%
\TextAt{9}{392.287}{381.759}{\Character{40}}%
\TextAt{13}{396.848}{381.759}{s}%
\TextAt{17}{402.362}{379.966}{2}%
\TextAt{16}{409.751}{381.759}{\Character{0}}%
\TextAt{13}{421.707}{381.759}{s}%
\TextAt{17}{427.221}{379.966}{1}%
\TextAt{9}{431.953}{381.759}{\Character{41}}%
\TextAt{13}{436.505}{381.759}{m}%
\TextAt{17}{446.744}{379.966}{2}%
\TextAt{13}{451.477}{381.759}{\Character{59}}%
\TextAt{8}{79.226}{353.587}{and}%
\TextAt{8}{101.951}{353.587}{so}%
\TextAt{8}{116.283}{353.587}{on.}%
\TextAt{8}{137.081}{353.587}{From}%
\TextAt{8}{167.768}{353.587}{this}%
\TextAt{8}{190.552}{353.587}{result}%
\TextAt{8}{223.092}{353.587}{one}%
\TextAt{8}{244.514}{353.587}{concludes}%
\TextAt{8}{297.211}{353.587}{that}%
\TextAt{8}{322.542}{353.587}{the}%
\TextAt{8}{342.66}{353.587}{vertices}%
\TextAt{8}{384.956}{353.587}{of}%
\TextAt{8}{398.248}{353.587}{the}%
\TextAt{8}{418.367}{353.587}{Newton}%
\TextAt{8}{461.571}{353.587}{polygon,}%
\TextAt{8}{509.027}{353.587}{from}%
\TextAt{8}{79.226}{336.541}{right}%
\TextAt{8}{107.524}{336.541}{to}%
\TextAt{8}{121.834}{336.541}{left,}%
\TextAt{8}{145.565}{336.541}{are}%
\TextAt{8}{165.076}{336.541}{given}%
\TextAt{8}{195.646}{336.541}{by}%
\TextAt{9}{148.9}{311.982}{\Character{40}}%
\TextAt{13}{153.453}{311.982}{n\Character{59}}%
\TextAt{9}{165.684}{311.982}{0\Character{41}}%
\TextAt{13}{176.09}{311.982}{\Character{59}}%
\TextAt{9}{193.289}{311.982}{\Character{40}}%
\TextAt{13}{197.841}{311.982}{n}%
\TextAt{16}{207.486}{311.982}{\Character{0}}%
\TextAt{13}{219.441}{311.982}{s}%
\TextAt{17}{224.955}{310.189}{1}%
\TextAt{13}{229.687}{311.982}{\Character{59}}%
\TextAt{13}{234.911}{311.982}{s}%
\TextAt{17}{240.445}{310.189}{1}%
\TextAt{13}{245.178}{311.982}{m}%
\TextAt{17}{255.417}{310.189}{1}%
\TextAt{9}{260.149}{311.982}{\Character{41}}%
\TextAt{13}{264.701}{311.982}{\Character{59}}%
\TextAt{9}{281.901}{311.982}{\Character{40}}%
\TextAt{13}{286.453}{311.982}{n}%
\TextAt{16}{296.097}{311.982}{\Character{0}}%
\TextAt{13}{308.053}{311.982}{s}%
\TextAt{17}{313.567}{310.189}{2}%
\TextAt{13}{318.299}{311.982}{\Character{59}}%
\TextAt{13}{323.523}{311.982}{s}%
\TextAt{17}{329.057}{310.189}{1}%
\TextAt{13}{333.789}{311.982}{m}%
\TextAt{17}{344.029}{310.189}{1}%
\TextAt{9}{351.418}{311.982}{\Character{43}}%
\TextAt{9}{363.17}{311.982}{\Character{40}}%
\TextAt{13}{367.731}{311.982}{s}%
\TextAt{17}{373.245}{310.189}{2}%
\TextAt{16}{380.634}{311.982}{\Character{0}}%
\TextAt{13}{392.589}{311.982}{s}%
\TextAt{17}{398.103}{310.189}{1}%
\TextAt{9}{402.836}{311.982}{\Character{41}}%
\TextAt{13}{407.388}{311.982}{m}%
\TextAt{17}{417.627}{310.189}{2}%
\TextAt{9}{422.36}{311.982}{\Character{41}}%
\TextAt{13}{426.912}{311.982}{\Character{59}}%
\TextAt{13}{444.092}{311.982}{\Character{58}}%
\TextAt{13}{449.328}{311.982}{\Character{58}}%
\TextAt{13}{454.552}{311.982}{\Character{58}}%
\TextAt{13}{459.789}{311.982}{\Character{58}}%
\TextAt{8}{79.226}{287.423}{The}%
\TextAt{8}{103.28}{287.423}{slope}%
\TextAt{8}{132.929}{287.423}{of}%
\TextAt{8}{146.259}{287.423}{the}%
\TextAt{8}{166.415}{287.423}{extreme}%
\TextAt{8}{210.96}{287.423}{right-hand}%
\TextAt{8}{268.524}{287.423}{line}%
\TextAt{8}{290.629}{287.423}{segment}%
\TextAt{8}{335.892}{287.423}{is}%
\TextAt{9}{261.522}{268.758}{0}%
\TextAt{16}{270.032}{268.758}{\Character{0}}%
\TextAt{13}{281.987}{268.758}{s}%
\TextAt{17}{287.501}{266.965}{1}%
\TextAt{13}{292.234}{268.758}{m}%
\TextAt{17}{302.473}{266.965}{1}%
\Vector{q 0 G 0 g 0.478 w 181.089 -510.341 m 243.639 -510.341 l S Q}%
\TextAt{13}{253.089}{252.47}{n}%
\TextAt{16}{262.733}{252.47}{\Character{0}}%
\TextAt{9}{274.688}{252.47}{\Character{40}}%
\TextAt{13}{279.241}{252.47}{n}%
\TextAt{16}{288.885}{252.47}{\Character{0}}%
\TextAt{13}{300.84}{252.47}{s}%
\TextAt{17}{306.354}{250.677}{1}%
\TextAt{9}{311.086}{252.47}{\Character{41}}%
\TextAt{9}{320.155}{260.671}{\Character{61}}%
\TextAt{16}{332.581}{260.671}{\Character{0}}%
\TextAt{13}{341.879}{260.671}{m}%
\TextAt{17}{352.118}{258.877}{1}%
\TextAt{13}{356.851}{260.671}{\Character{59}}%
\TextAt{8}{79.226}{232.556}{and,}%
\TextAt{8}{105.241}{232.556}{proceeding}%
\TextAt{8}{164.096}{232.556}{further}%
\TextAt{8}{203.44}{232.556}{to}%
\TextAt{8}{217.751}{232.556}{the}%
\TextAt{8}{237.907}{232.556}{left,}%
\TextAt{9}{120.355}{212.629}{\Character{40}}%
\TextAt{13}{124.908}{212.629}{s}%
\TextAt{17}{130.422}{210.836}{1}%
\TextAt{13}{135.154}{212.629}{m}%
\TextAt{17}{145.393}{210.836}{1}%
\TextAt{9}{152.782}{212.629}{\Character{43}}%
\TextAt{16}{164.543}{212.629}{\Character{1}}%
\TextAt{16}{169.839}{212.629}{\Character{1}}%
\TextAt{16}{175.147}{212.629}{\Character{1}}%
\TextAt{9}{181.148}{212.629}{\Character{43}}%
\TextAt{9}{192.9}{212.629}{\Character{40}}%
\TextAt{13}{197.461}{212.629}{s}%
\TextAt{21}{202.975}{210.836}{j}%
\TextAt{16}{210.014}{212.629}{\Character{0}}%
\TextAt{13}{221.969}{212.629}{s}%
\TextAt{21}{227.483}{210.836}{j}%
\TextAt{3}{231.367}{210.836}{\Character{0}}%
\TextAt{17}{237.954}{210.836}{1}%
\TextAt{9}{242.686}{212.629}{\Character{41}}%
\TextAt{13}{247.239}{212.629}{m}%
\TextAt{21}{257.478}{210.836}{j}%
\TextAt{9}{261.86}{212.629}{\Character{41}}%
\TextAt{16}{269.069}{212.629}{\Character{0}}%
\TextAt{9}{281.024}{212.629}{\Character{40}}%
\TextAt{13}{285.576}{212.629}{s}%
\TextAt{17}{291.09}{210.836}{1}%
\TextAt{13}{295.823}{212.629}{m}%
\TextAt{17}{306.062}{210.836}{1}%
\TextAt{9}{313.451}{212.629}{\Character{43}}%
\TextAt{16}{325.212}{212.629}{\Character{1}}%
\TextAt{16}{330.508}{212.629}{\Character{1}}%
\TextAt{16}{335.816}{212.629}{\Character{1}}%
\TextAt{9}{341.817}{212.629}{\Character{43}}%
\TextAt{9}{353.569}{212.629}{\Character{40}}%
\TextAt{13}{358.13}{212.629}{s}%
\TextAt{21}{363.644}{210.836}{j}%
\TextAt{17}{367.528}{210.836}{\Character{43}1}%
\TextAt{16}{381.504}{212.629}{\Character{0}}%
\TextAt{13}{393.459}{212.629}{s}%
\TextAt{21}{398.973}{210.836}{j}%
\TextAt{9}{403.355}{212.629}{\Character{41}}%
\TextAt{13}{407.907}{212.629}{m}%
\TextAt{21}{418.147}{210.836}{j}%
\TextAt{17}{422.031}{210.836}{\Character{43}1}%
\TextAt{9}{433.35}{212.629}{\Character{41}}%
\Vector{q 0 G 0 g 0.478 w 48.355 -566.47 m 365.902 -566.47 l S Q}%
\TextAt{9}{225.812}{196.341}{\Character{40}}%
\TextAt{13}{230.364}{196.341}{n}%
\TextAt{16}{240.009}{196.341}{\Character{0}}%
\TextAt{13}{251.964}{196.341}{s}%
\TextAt{21}{257.478}{194.548}{j}%
\TextAt{9}{261.86}{196.341}{\Character{41}}%
\TextAt{16}{269.069}{196.341}{\Character{0}}%
\TextAt{9}{281.024}{196.341}{\Character{40}}%
\TextAt{13}{285.576}{196.341}{n}%
\TextAt{16}{295.221}{196.341}{\Character{0}}%
\TextAt{13}{307.176}{196.341}{s}%
\TextAt{21}{312.69}{194.548}{j}%
\TextAt{17}{316.574}{194.548}{\Character{43}1}%
\TextAt{9}{327.893}{196.341}{\Character{41}}%
\TextAt{9}{442.418}{204.542}{\Character{61}}%
\TextAt{16}{454.844}{204.542}{\Character{0}}%
\TextAt{13}{464.142}{204.542}{m}%
\TextAt{21}{474.381}{202.748}{j}%
\TextAt{17}{478.265}{202.748}{\Character{43}1}%
\TextAt{13}{489.584}{204.542}{\Character{58}}%
\Vector{q 0 G 0 g 0.399 w 452.334 -597.581 m 452.334 -589.511 l S Q}%
\Vector{q 0 G 0 g 0.399 w 452.534 -589.711 m 459.707 -589.711 l S Q}%
\Vector{q 0 G 0 g 0.399 w 452.534 -597.382 m 459.707 -597.382 l S Q}%
\Vector{q 0 G 0 g 0.399 w 459.906 -597.581 m 459.906 -589.511 l S Q}%
\TextAt{8}{97.159}{148.687}{We}%
\TextAt{8}{117.035}{148.687}{emphasize}%
\TextAt{8}{172.753}{148.687}{that,}%
\TextAt{8}{201.159}{148.687}{according}%
\TextAt{8}{253.576}{148.687}{to}%
\TextAt{8}{267.618}{148.687}{the}%
\TextAt{8}{287.506}{148.687}{preceding}%
\TextAt{8}{339.912}{148.687}{proposition,}%
\TextAt{8}{404.483}{148.687}{the}%
\TextAt{8}{424.371}{148.687}{Newton}%
\TextAt{8}{467.344}{148.687}{polygon}%
\TextAt{8}{511.31}{148.687}{con-}%
\TextAt{8}{79.226}{131.641}{sists}%
\TextAt{8}{105.473}{131.641}{of}%
\TextAt{8}{119.502}{131.641}{precisely}%
\TextAt{8}{168.057}{131.641}{one}%
\TextAt{8}{190.215}{131.641}{segment}%
\TextAt{8}{236.176}{131.641}{if}%
\TextAt{8}{247.598}{131.641}{and}%
\TextAt{8}{271.06}{131.641}{only}%
\TextAt{8}{297.45}{131.641}{if}%
\TextAt{8}{308.872}{131.641}{the}%
\TextAt{8}{329.727}{131.641}{roots}%
\TextAt{13}{360.087}{131.641}{\Character{11}}%
\TextAt{17}{367.532}{129.848}{1}%
\TextAt{13}{372.265}{131.641}{\Character{59}}%
\TextAt{13}{377.489}{131.641}{\Character{58}}%
\TextAt{13}{382.726}{131.641}{\Character{58}}%
\TextAt{13}{387.95}{131.641}{\Character{58}}%
\TextAt{13}{393.187}{131.641}{\Character{59}}%
\TextAt{13}{398.423}{131.641}{\Character{11}}%
\TextAt{21}{405.931}{129.848}{n}%
\TextAt{8}{416.164}{131.641}{of}%
\TextAt{13}{430.192}{131.641}{f}%
\TextAt{8}{441.835}{131.641}{all}%
\TextAt{8}{458.792}{131.641}{have}%
\TextAt{8}{486.486}{131.641}{the}%
\TextAt{8}{507.341}{131.641}{same}%
\TextAt{8}{79.226}{114.595}{value.}%
\TextAt{8}{114.02}{114.595}{In}%
\TextAt{8}{128.629}{114.595}{general,}%
\TextAt{13}{172.196}{114.595}{f}%
\TextAt{9}{179.242}{114.595}{\Character{40}}%
\TextAt{13}{183.794}{114.595}{x}%
\TextAt{9}{190.446}{114.595}{\Character{41}}%
\TextAt{8}{198.884}{114.595}{factors}%
\TextAt{8}{236.985}{114.595}{into}%
\TextAt{8}{260.716}{114.595}{a}%
\TextAt{8}{270.459}{114.595}{product}%
\TextAt{8}{314.347}{114.595}{according}%
\TextAt{8}{367.021}{114.595}{to}%
\TextAt{8}{381.319}{114.595}{the}%
\TextAt{8}{401.463}{114.595}{slopes}%
\TextAt{16}{435.715}{114.595}{\Character{0}}%
\TextAt{13}{445.013}{114.595}{m}%
\TextAt{21}{455.253}{112.802}{r}%
\TextAt{13}{463.128}{114.595}{\Character{60}}%
\TextAt{16}{475.554}{114.595}{\Character{1}}%
\TextAt{16}{480.85}{114.595}{\Character{1}}%
\TextAt{16}{486.158}{114.595}{\Character{1}}%
\TextAt{13}{492.822}{114.595}{\Character{60}}%
\TextAt{16}{505.248}{114.595}{\Character{0}}%
\TextAt{13}{514.546}{114.595}{m}%
\TextAt{17}{524.785}{112.802}{1}%
\TextAt{8}{529.518}{114.595}{,}%
\TextAt{13}{256.622}{83.526}{f}%
\TextAt{9}{263.668}{83.526}{\Character{40}}%
\TextAt{13}{268.22}{83.526}{x}%
\TextAt{9}{274.872}{83.526}{\Character{41}}%
\TextAt{9}{282.727}{83.526}{\Character{61}}%
\TextAt{13}{295.171}{83.526}{a}%
\TextAt{21}{301.316}{81.733}{n}%
\TextAt{21}{314.555}{98.47}{r}%
\TextAt{20}{308.945}{94.883}{\Character{89}}%
\TextAt{21}{309.231}{69.689}{j}%
\TextAt{17}{313.115}{69.689}{\Character{61}1}%
\TextAt{13}{326.213}{83.526}{f}%
\TextAt{21}{331.983}{81.733}{j}%
\TextAt{9}{336.366}{83.526}{\Character{40}}%
\TextAt{13}{340.918}{83.526}{x}%
\TextAt{9}{347.57}{83.526}{\Character{41}}%
\TextAt{13}{352.122}{83.526}{\Character{59}}%
\TextAt{8}{297.211}{43.033}{133}%
\Destination{78.23}{809.815}{nav.page.133}%
\Destination{79.36}{776.155}{nav.print.145}%
\Destination{79.36}{302.455}{nav.print.146}%
}

% Page 146 (body page 134)
\Page{%
\TextAt{2}{79.226}{796.943}{II.}%
\TextAt{2}{94.073}{796.943}{The}%
\TextAt{2}{116.492}{796.943}{Theory}%
\TextAt{2}{154.401}{796.943}{of}%
\TextAt{2}{166.826}{796.943}{Valuations}%
\Vector{q 0 G 0 g 0.299 w 7.226 18.729 m 460.769 18.729 l S Q}%
\TextAt{8}{79.226}{761.124}{where}%
\TextAt{13}{243.218}{744.078}{f}%
\TextAt{21}{248.988}{742.284}{j}%
\TextAt{9}{253.37}{744.078}{\Character{40}}%
\TextAt{13}{257.922}{744.078}{x}%
\TextAt{9}{264.574}{744.078}{\Character{41}}%
\TextAt{9}{272.429}{744.078}{\Character{61}}%
\TextAt{20}{296.964}{755.435}{\Character{89}}%
\TextAt{21}{284.873}{729.532}{w}%
\TextAt{17}{291.164}{729.532}{\Character{40}}%
\TextAt{21}{294.457}{729.532}{\Character{11}}%
\TextAt{22}{299.896}{728.317}{i}%
\TextAt{17}{303.058}{729.532}{\Character{41}\Character{61}}%
\TextAt{21}{312.938}{729.532}{m}%
\TextAt{22}{320.428}{728.317}{j}%
\TextAt{9}{324.331}{744.078}{\Character{40}}%
\TextAt{13}{328.883}{744.078}{x}%
\TextAt{16}{338.192}{744.078}{\Character{0}}%
\TextAt{13}{350.147}{744.078}{\Character{11}}%
\TextAt{21}{357.593}{742.284}{i}%
\TextAt{9}{360.974}{744.078}{\Character{41}}%
\TextAt{13}{365.526}{744.078}{\Character{58}}%
\TextAt{8}{79.226}{709.521}{Here}%
\TextAt{8}{106.163}{709.521}{the}%
\TextAt{8}{125.628}{709.521}{factor}%
\TextAt{13}{158.435}{709.521}{f}%
\TextAt{21}{164.205}{707.728}{j}%
\TextAt{8}{171.793}{709.521}{corresponds}%
\TextAt{8}{235.289}{709.521}{to}%
\TextAt{8}{248.908}{709.521}{the}%
\TextAt{9}{268.373}{709.521}{\Character{40}}%
\TextAt{13}{272.925}{709.521}{r}%
\TextAt{16}{279.769}{709.521}{\Character{0}}%
\TextAt{13}{290.31}{709.521}{j}%
\TextAt{9}{297.068}{709.521}{\Character{43}}%
\TextAt{9}{307.397}{709.521}{1\Character{41}}%
\TextAt{8}{317.821}{709.521}{-th}%
\TextAt{8}{335.983}{709.521}{segment}%
\TextAt{8}{380.554}{709.521}{of}%
\TextAt{8}{393.192}{709.521}{the}%
\TextAt{8}{412.657}{709.521}{Newton}%
\TextAt{8}{455.208}{709.521}{polygon,}%
\TextAt{8}{502.14}{709.521}{whose}%
\TextAt{8}{79.226}{692.475}{slope}%
\TextAt{8}{108.875}{692.475}{equals}%
\TextAt{8}{144.382}{692.475}{minus}%
\TextAt{8}{178.586}{692.475}{the}%
\TextAt{8}{198.742}{692.475}{value}%
\TextAt{8}{228.989}{692.475}{of}%
\TextAt{8}{242.319}{692.475}{the}%
\TextAt{8}{262.475}{692.475}{roots}%
\TextAt{8}{292.136}{692.475}{of}%
\TextAt{13}{305.466}{692.475}{f}%
\TextAt{21}{311.236}{690.682}{j}%
\TextAt{8}{315.618}{692.475}{.}%
\TextAt{12}{79.226}{666.463}{Proposition}%
\TextAt{12}{151.926}{666.463}{(6.4).}%
\TextAt{8}{187.231}{666.463}{If}%
\TextAt{8}{198.913}{666.463}{the}%
\TextAt{8}{219.059}{666.463}{valuation}%
\TextAt{13}{270.121}{666.463}{v}%
\TextAt{8}{280.097}{666.463}{admits}%
\TextAt{8}{318.523}{666.463}{a}%
\TextAt{8}{328.268}{666.463}{unique}%
\TextAt{8}{366.299}{666.463}{extension}%
\TextAt{13}{418.055}{666.463}{w}%
\TextAt{8}{430.631}{666.463}{to}%
\TextAt{8}{444.931}{666.463}{the}%
\TextAt{8}{465.078}{666.463}{splitting}%
\TextAt{8}{511.31}{666.463}{field}%
\TextAt{13}{79.226}{649.417}{L}%
\TextAt{8}{91.088}{649.417}{of}%
\TextAt{13}{104.418}{649.417}{f}%
\TextAt{8}{111.464}{649.417}{,}%
\TextAt{8}{118.614}{649.417}{then}%
\TextAt{8}{145.274}{649.417}{the}%
\TextAt{8}{165.43}{649.417}{factorization}%
\TextAt{13}{258.247}{626.525}{f}%
\TextAt{9}{265.294}{626.525}{\Character{40}}%
\TextAt{13}{269.846}{626.525}{x}%
\TextAt{9}{276.498}{626.525}{\Character{41}}%
\TextAt{9}{284.353}{626.525}{\Character{61}}%
\TextAt{13}{296.797}{626.525}{a}%
\TextAt{21}{302.942}{624.731}{n}%
\TextAt{21}{316.18}{641.469}{r}%
\TextAt{20}{310.571}{637.882}{\Character{89}}%
\TextAt{21}{310.856}{612.688}{j}%
\TextAt{17}{314.74}{612.688}{\Character{61}1}%
\TextAt{13}{327.839}{626.525}{f}%
\TextAt{21}{333.609}{624.731}{j}%
\TextAt{9}{337.991}{626.525}{\Character{40}}%
\TextAt{13}{342.544}{626.525}{x}%
\TextAt{9}{349.196}{626.525}{\Character{41}}%
\TextAt{8}{79.226}{593.504}{is}%
\TextAt{8}{90.99}{593.504}{defined}%
\TextAt{8}{131.303}{593.504}{already}%
\TextAt{8}{172.608}{593.504}{over}%
\TextAt{13}{197.654}{593.504}{K}%
\TextAt{8}{208.472}{593.504}{,}%
\TextAt{8}{215.621}{593.504}{i.e.,}%
\TextAt{13}{237.726}{593.504}{f}%
\TextAt{21}{243.496}{591.711}{j}%
\TextAt{9}{247.878}{593.504}{\Character{40}}%
\TextAt{13}{252.431}{593.504}{x}%
\TextAt{9}{259.083}{593.504}{\Character{41}}%
\TextAt{9}{266.938}{593.504}{\Character{61}}%
\TextAt{20}{279.381}{602.471}{\Character{81}}%
\TextAt{21}{290.672}{590.017}{w}%
\TextAt{17}{296.964}{590.017}{\Character{40}}%
\TextAt{21}{300.257}{590.017}{\Character{11}}%
\TextAt{22}{305.696}{588.802}{i}%
\TextAt{17}{308.858}{590.017}{\Character{41}\Character{61}}%
\TextAt{21}{318.737}{590.017}{m}%
\TextAt{22}{326.228}{588.802}{j}%
\TextAt{9}{330.629}{593.504}{\Character{40}}%
\TextAt{13}{335.181}{593.504}{x}%
\TextAt{16}{344.49}{593.504}{\Character{0}}%
\TextAt{13}{356.445}{593.504}{\Character{11}}%
\TextAt{21}{363.891}{591.711}{i}%
\TextAt{9}{367.272}{593.504}{\Character{41}}%
\NativeMathAt{375.145}{593.504}{11.955}{\in}%
\TextAt{13}{386.436}{593.504}{K}%
\TextAt{9}{397.254}{593.504}{\Character{91}}%
\TextAt{13}{400.505}{593.504}{x}%
\TextAt{9}{407.157}{593.504}{\Character{93}}%
\TextAt{8}{410.409}{593.504}{.}%
\TextAt{19}{79.226}{558.525}{Proof.}%
\TextAt{8}{116}{558.525}{We}%
\TextAt{8}{137.785}{558.525}{may}%
\TextAt{8}{164.795}{558.525}{clearly}%
\TextAt{8}{203.831}{558.525}{assume}%
\TextAt{8}{245.916}{558.525}{that}%
\TextAt{13}{272.925}{558.525}{a}%
\TextAt{21}{279.07}{556.732}{n}%
\TextAt{9}{290.823}{558.525}{\Character{61}}%
\TextAt{9}{306.03}{558.525}{1}%
\TextAt{8}{311.897}{558.525}{.}%
\TextAt{8}{325.267}{558.525}{The}%
\TextAt{8}{350.961}{558.525}{statement}%
\TextAt{8}{406.974}{558.525}{is}%
\TextAt{8}{420.378}{558.525}{obvious}%
\TextAt{8}{464.364}{558.525}{when}%
\TextAt{13}{496.562}{558.525}{f}%
\TextAt{9}{503.608}{558.525}{\Character{40}}%
\TextAt{13}{508.161}{558.525}{x}%
\TextAt{9}{514.813}{558.525}{\Character{41}}%
\TextAt{8}{524.903}{558.525}{is}%
\TextAt{8}{79.226}{541.479}{irreducible}%
\TextAt{8}{137.709}{541.479}{because}%
\TextAt{8}{181.618}{541.479}{then}%
\TextAt{8}{208.886}{541.479}{one}%
\TextAt{8}{230.954}{541.479}{has}%
\TextAt{13}{252.435}{541.479}{\Character{11}}%
\TextAt{21}{259.881}{539.686}{i}%
\TextAt{9}{267.619}{541.479}{\Character{61}}%
\TextAt{13}{281.08}{541.479}{\Character{27}}%
\TextAt{21}{287.734}{539.686}{i}%
\TextAt{13}{291.115}{541.479}{\Character{11}}%
\TextAt{17}{298.56}{539.686}{1}%
\TextAt{8}{307.798}{541.479}{for}%
\TextAt{8}{326.291}{541.479}{some}%
\TextAt{13}{356.225}{541.479}{\Character{27}}%
\TextAt{21}{362.878}{539.686}{i}%
\NativeMathAt{370.616}{541.479}{11.955}{\in}%
\TextAt{13}{382.943}{541.479}{G}%
\TextAt{9}{392.176}{541.479}{\Character{40}}%
\TextAt{13}{396.729}{541.479}{L}%
\TextAt{16}{409.05}{541.479}{\Character{106}}%
\TextAt{13}{416.727}{541.479}{K}%
\TextAt{9}{427.545}{541.479}{\Character{41}}%
\TextAt{8}{432.097}{541.479}{,}%
\TextAt{8}{440.006}{541.479}{and}%
\TextAt{8}{463.377}{541.479}{since,}%
\TextAt{8}{496.057}{541.479}{for}%
\TextAt{8}{514.55}{541.479}{any}%
\TextAt{8}{79.226}{524.433}{extension}%
\TextAt{13}{130.146}{524.433}{w}%
\TextAt{8}{141.888}{524.433}{of}%
\TextAt{13}{154.372}{524.433}{v}%
\TextAt{8}{160.461}{524.433}{,}%
\TextAt{13}{166.933}{524.433}{w}%
\TextAt{16}{176.55}{524.433}{\Character{14}}%
\TextAt{13}{183.455}{524.433}{\Character{27}}%
\TextAt{21}{190.109}{522.64}{i}%
\TextAt{8}{196.542}{524.433}{is}%
\TextAt{8}{207.46}{524.433}{another}%
\TextAt{8}{249.546}{524.433}{one,}%
\TextAt{8}{273.581}{524.433}{the}%
\TextAt{8}{292.892}{524.433}{uniqueness}%
\TextAt{8}{351.021}{524.433}{implies}%
\TextAt{8}{389.902}{524.433}{that}%
\TextAt{13}{414.426}{524.433}{w}%
\TextAt{9}{423.115}{524.433}{\Character{40}}%
\TextAt{13}{427.667}{524.433}{\Character{11}}%
\TextAt{21}{435.113}{522.64}{i}%
\TextAt{9}{438.494}{524.433}{\Character{41}}%
\TextAt{9}{446.349}{524.433}{\Character{61}}%
\TextAt{13}{458.793}{524.433}{w}%
\TextAt{9}{467.482}{524.433}{\Character{40}}%
\TextAt{13}{472.035}{524.433}{\Character{27}}%
\TextAt{21}{478.688}{522.64}{i}%
\TextAt{13}{482.07}{524.433}{\Character{11}}%
\TextAt{17}{489.515}{522.64}{1}%
\TextAt{9}{494.247}{524.433}{\Character{41}}%
\TextAt{9}{502.102}{524.433}{\Character{61}}%
\TextAt{13}{514.546}{524.433}{m}%
\TextAt{17}{524.785}{522.64}{1}%
\TextAt{8}{529.518}{524.433}{,}%
\TextAt{8}{79.226}{507.387}{hence}%
\TextAt{13}{111.732}{507.387}{f}%
\TextAt{17}{117.502}{505.594}{1}%
\TextAt{9}{122.235}{507.387}{\Character{40}}%
\TextAt{13}{126.787}{507.387}{x}%
\TextAt{9}{133.439}{507.387}{\Character{41}}%
\TextAt{9}{141.294}{507.387}{\Character{61}}%
\TextAt{13}{153.738}{507.387}{f}%
\TextAt{9}{160.784}{507.387}{\Character{40}}%
\TextAt{13}{165.336}{507.387}{x}%
\TextAt{9}{171.988}{507.387}{\Character{41}}%
\TextAt{8}{176.541}{507.387}{.}%
\TextAt{8}{97.159}{487.353}{The}%
\TextAt{8}{121.438}{487.353}{general}%
\TextAt{8}{161.988}{487.353}{case}%
\TextAt{8}{186.984}{487.353}{follows}%
\TextAt{8}{225.645}{487.353}{by}%
\TextAt{8}{242.129}{487.353}{induction}%
\TextAt{8}{294.383}{487.353}{on}%
\TextAt{13}{310.867}{487.353}{n}%
\TextAt{8}{317.855}{487.353}{.}%
\TextAt{8}{326.979}{487.353}{For}%
\TextAt{13}{348.173}{487.353}{n}%
\TextAt{9}{358.865}{487.353}{\Character{61}}%
\TextAt{9}{371.657}{487.353}{1}%
\TextAt{8}{381.65}{487.353}{there}%
\TextAt{8}{411.787}{487.353}{is}%
\TextAt{8}{423.776}{487.353}{nothing}%
\TextAt{8}{466.932}{487.353}{to}%
\TextAt{8}{481.468}{487.353}{show.}%
\TextAt{8}{515.697}{487.353}{Let}%
\TextAt{13}{79.226}{470.307}{p}%
\TextAt{9}{85.101}{470.307}{\Character{40}}%
\TextAt{13}{89.654}{470.307}{x}%
\TextAt{9}{96.306}{470.307}{\Character{41}}%
\TextAt{8}{105.379}{470.307}{be}%
\TextAt{8}{121.927}{470.307}{the}%
\TextAt{8}{142.707}{470.307}{minimal}%
\TextAt{8}{188.855}{470.307}{polynomial}%
\TextAt{8}{249.972}{470.307}{of}%
\TextAt{13}{263.925}{470.307}{\Character{11}}%
\TextAt{17}{271.371}{468.513}{1}%
\TextAt{8}{280.624}{470.307}{and}%
\TextAt{13}{304.01}{470.307}{g}%
\TextAt{9}{310.045}{470.307}{\Character{40}}%
\TextAt{13}{314.597}{470.307}{x}%
\TextAt{9}{321.249}{470.307}{\Character{41}}%
\TextAt{9}{330.168}{470.307}{\Character{61}}%
\TextAt{13}{343.673}{470.307}{f}%
\TextAt{9}{350.719}{470.307}{\Character{40}}%
\TextAt{13}{355.271}{470.307}{x}%
\TextAt{9}{361.923}{470.307}{\Character{41}\Character{47}}%
\TextAt{13}{372.329}{470.307}{p}%
\TextAt{9}{378.204}{470.307}{\Character{40}}%
\TextAt{13}{382.756}{470.307}{x}%
\TextAt{9}{389.408}{470.307}{\Character{41}}%
\NativeMathAt{398.344}{470.307}{11.955}{\in}%
\TextAt{13}{410.697}{470.307}{K}%
\TextAt{9}{421.515}{470.307}{\Character{91}}%
\TextAt{13}{424.767}{470.307}{x}%
\TextAt{9}{431.419}{470.307}{\Character{93}}%
\TextAt{8}{434.67}{470.307}{.}%
\TextAt{8}{444.989}{470.307}{Since}%
\TextAt{8}{476.17}{470.307}{all}%
\TextAt{8}{493.053}{470.307}{roots}%
\TextAt{8}{523.337}{470.307}{of}%
\TextAt{13}{79.226}{453.261}{p}%
\TextAt{9}{85.101}{453.261}{\Character{40}}%
\TextAt{13}{89.654}{453.261}{x}%
\TextAt{9}{96.306}{453.261}{\Character{41}}%
\TextAt{8}{105.716}{453.261}{have}%
\TextAt{8}{133.672}{453.261}{the}%
\TextAt{8}{154.789}{453.261}{same}%
\TextAt{8}{185.076}{453.261}{value}%
\TextAt{13}{216.284}{453.261}{m}%
\TextAt{17}{226.523}{451.467}{1}%
\TextAt{8}{231.255}{453.261}{,}%
\TextAt{13}{239.606}{453.261}{p}%
\TextAt{9}{245.481}{453.261}{\Character{40}}%
\TextAt{13}{250.033}{453.261}{x}%
\TextAt{9}{256.685}{453.261}{\Character{41}}%
\TextAt{8}{266.096}{453.261}{is}%
\TextAt{8}{278.821}{453.261}{a}%
\TextAt{8}{289.537}{453.261}{divisor}%
\TextAt{8}{328.611}{453.261}{of}%
\TextAt{13}{342.902}{453.261}{f}%
\TextAt{17}{348.672}{451.467}{1}%
\TextAt{9}{353.404}{453.261}{\Character{40}}%
\TextAt{13}{357.957}{453.261}{x}%
\TextAt{9}{364.609}{453.261}{\Character{41}}%
\TextAt{8}{369.161}{453.261}{.}%
\TextAt{8}{380.492}{453.261}{Let}%
\TextAt{13}{402.422}{453.261}{g}%
\TextAt{17}{408.028}{451.467}{1}%
\TextAt{9}{412.76}{453.261}{\Character{40}}%
\TextAt{13}{417.312}{453.261}{x}%
\TextAt{9}{423.964}{453.261}{\Character{41}}%
\TextAt{9}{433.456}{453.261}{\Character{61}}%
\TextAt{13}{447.538}{453.261}{f}%
\TextAt{17}{453.308}{451.467}{1}%
\TextAt{9}{458.04}{453.261}{\Character{40}}%
\TextAt{13}{462.593}{453.261}{x}%
\TextAt{9}{469.245}{453.261}{\Character{41}\Character{47}}%
\TextAt{13}{479.65}{453.261}{p}%
\TextAt{9}{485.525}{453.261}{\Character{40}}%
\TextAt{13}{490.078}{453.261}{x}%
\TextAt{9}{496.73}{453.261}{\Character{41}}%
\TextAt{8}{501.282}{453.261}{.}%
\TextAt{8}{512.613}{453.261}{The}%
\TextAt{8}{79.226}{436.215}{factorization}%
\TextAt{8}{147.203}{436.215}{of}%
\TextAt{13}{160.533}{436.215}{g}%
\TextAt{9}{166.568}{436.215}{\Character{40}}%
\TextAt{13}{171.12}{436.215}{x}%
\TextAt{9}{177.772}{436.215}{\Character{41}}%
\TextAt{8}{186.222}{436.215}{according}%
\TextAt{8}{238.908}{436.215}{to}%
\TextAt{8}{253.218}{436.215}{the}%
\TextAt{8}{273.375}{436.215}{slopes}%
\TextAt{8}{307.638}{436.215}{is}%
\TextAt{13}{249.971}{400.703}{g}%
\TextAt{9}{256.005}{400.703}{\Character{40}}%
\TextAt{13}{260.558}{400.703}{x}%
\TextAt{9}{267.21}{400.703}{\Character{41}}%
\TextAt{9}{275.065}{400.703}{\Character{61}}%
\TextAt{13}{287.508}{400.703}{g}%
\TextAt{17}{293.114}{398.91}{1}%
\TextAt{9}{297.846}{400.703}{\Character{40}}%
\TextAt{13}{302.398}{400.703}{x}%
\TextAt{9}{309.051}{400.703}{\Character{41}}%
\TextAt{21}{321.205}{415.647}{r}%
\TextAt{20}{315.595}{412.061}{\Character{89}}%
\TextAt{21}{315.881}{386.867}{j}%
\TextAt{17}{319.765}{386.867}{\Character{61}2}%
\TextAt{13}{332.864}{400.703}{f}%
\TextAt{21}{338.634}{398.91}{j}%
\TextAt{9}{343.016}{400.703}{\Character{40}}%
\TextAt{13}{347.568}{400.703}{x}%
\TextAt{9}{354.22}{400.703}{\Character{41}}%
\TextAt{13}{358.773}{400.703}{\Character{58}}%
\TextAt{8}{79.226}{362.203}{Since}%
\TextAt{8}{109.784}{362.203}{deg}%
\TextAt{9}{127.346}{362.203}{\Character{40}}%
\TextAt{13}{131.898}{362.203}{g}%
\TextAt{9}{137.932}{362.203}{\Character{41}}%
\TextAt{13}{145.805}{362.203}{\Character{60}}%
\TextAt{8}{158.231}{362.203}{deg}%
\TextAt{9}{175.793}{362.203}{\Character{40}}%
\TextAt{13}{180.345}{362.203}{f}%
\TextAt{9}{187.392}{362.203}{\Character{41}}%
\TextAt{8}{191.944}{362.203}{,}%
\TextAt{8}{199.093}{362.203}{it}%
\TextAt{8}{210.797}{362.203}{follows}%
\TextAt{8}{249.233}{362.203}{that}%
\TextAt{13}{274.602}{362.203}{f}%
\TextAt{21}{280.372}{360.41}{j}%
\TextAt{9}{284.754}{362.203}{\Character{40}}%
\TextAt{13}{289.307}{362.203}{x}%
\TextAt{9}{295.959}{362.203}{\Character{41}}%
\NativeMathAt{303.832}{362.203}{11.955}{\in}%
\TextAt{13}{315.123}{362.203}{K}%
\TextAt{9}{325.941}{362.203}{\Character{91}}%
\TextAt{13}{329.192}{362.203}{x}%
\TextAt{9}{335.844}{362.203}{\Character{93}}%
\TextAt{8}{342.993}{362.203}{for}%
\TextAt{8}{360.878}{362.203}{all}%
\TextAt{13}{377.137}{362.203}{j}%
\TextAt{9}{385.974}{362.203}{\Character{61}}%
\TextAt{9}{398.384}{362.203}{1}%
\TextAt{13}{404.252}{362.203}{\Character{59}}%
\TextAt{13}{409.476}{362.203}{\Character{58}}%
\TextAt{13}{414.713}{362.203}{\Character{58}}%
\TextAt{13}{419.949}{362.203}{\Character{58}}%
\TextAt{13}{425.174}{362.203}{\Character{59}}%
\TextAt{13}{430.41}{362.203}{r}%
\TextAt{8}{436.073}{362.203}{.}%
\Vector{q 0 G 0 g 0.399 w 452.334 -411.797 m 452.334 -403.727 l S Q}%
\Vector{q 0 G 0 g 0.399 w 452.534 -403.926 m 459.707 -403.926 l S Q}%
\Vector{q 0 G 0 g 0.399 w 452.534 -411.597 m 459.707 -411.597 l S Q}%
\Vector{q 0 G 0 g 0.399 w 459.906 -411.797 m 459.906 -403.727 l S Q}%
\TextAt{8}{97.159}{333.202}{If}%
\TextAt{8}{110.185}{333.202}{the}%
\TextAt{8}{131.676}{333.202}{polynomial}%
\TextAt{13}{193.503}{333.202}{f}%
\TextAt{8}{205.782}{333.202}{is}%
\TextAt{8}{218.88}{333.202}{irreducible,}%
\TextAt{8}{281.674}{333.202}{then,}%
\TextAt{8}{313.254}{333.202}{by}%
\TextAt{8}{330.847}{333.202}{the}%
\TextAt{8}{352.338}{333.202}{above}%
\TextAt{8}{386.848}{333.202}{factorization}%
\TextAt{8}{456.159}{333.202}{result,}%
\TextAt{8}{493.657}{333.202}{there}%
\TextAt{8}{524.903}{333.202}{is}%
\TextAt{8}{79.226}{316.156}{only}%
\TextAt{8}{105.521}{316.156}{one}%
\TextAt{8}{127.584}{316.156}{slope,}%
\TextAt{8}{161.239}{316.156}{i.e.,}%
\TextAt{8}{184.099}{316.156}{the}%
\TextAt{8}{204.858}{316.156}{Newton}%
\TextAt{8}{248.704}{316.156}{polygon}%
\TextAt{8}{293.541}{316.156}{consists}%
\TextAt{8}{337.255}{316.156}{of}%
\TextAt{8}{351.189}{316.156}{a}%
\TextAt{8}{361.548}{316.156}{single}%
\TextAt{8}{394.729}{316.156}{segment.}%
\TextAt{8}{446.352}{316.156}{The}%
\TextAt{8}{471.01}{316.156}{values}%
\TextAt{8}{506.474}{316.156}{of}%
\TextAt{8}{520.408}{316.156}{all}%
\TextAt{8}{79.226}{299.11}{coefficients}%
\TextAt{8}{138.464}{299.11}{lie}%
\TextAt{8}{154.066}{299.11}{on}%
\TextAt{8}{170.325}{299.11}{or}%
\TextAt{8}{184.635}{299.11}{above}%
\TextAt{8}{217.81}{299.11}{this}%
\TextAt{8}{240.633}{299.11}{line}%
\TextAt{8}{262.738}{299.11}{segment}%
\TextAt{8}{308}{299.11}{and}%
\TextAt{8}{330.763}{299.11}{we}%
\TextAt{8}{347.99}{299.11}{get}%
\TextAt{8}{367.501}{299.11}{the}%
\TextAt{12}{79.226}{273.098}{Corollary}%
\TextAt{12}{139.021}{273.098}{(6.5).}%
\TextAt{8}{174.342}{273.098}{Let}%
\TextAt{13}{195.318}{273.098}{f}%
\TextAt{9}{202.364}{273.098}{\Character{40}}%
\TextAt{13}{206.916}{273.098}{x}%
\TextAt{9}{213.568}{273.098}{\Character{41}}%
\TextAt{9}{221.435}{273.098}{\Character{61}}%
\TextAt{13}{233.887}{273.098}{a}%
\TextAt{17}{240.032}{271.305}{0}%
\TextAt{9}{247.425}{273.098}{\Character{43}}%
\TextAt{13}{259.19}{273.098}{a}%
\TextAt{17}{265.335}{271.305}{1}%
\TextAt{13}{270.068}{273.098}{x}%
\TextAt{9}{279.381}{273.098}{\Character{43}}%
\TextAt{16}{291.146}{273.098}{\Character{1}}%
\TextAt{16}{296.442}{273.098}{\Character{1}}%
\TextAt{16}{301.75}{273.098}{\Character{1}}%
\TextAt{9}{307.754}{273.098}{\Character{43}}%
\TextAt{13}{319.52}{273.098}{a}%
\TextAt{21}{325.665}{271.305}{n}%
\TextAt{13}{331.301}{273.098}{x}%
\TextAt{21}{337.953}{277.436}{n}%
\NativeMathAt{346.92}{273.098}{11.955}{\in}%
\TextAt{13}{358.221}{273.098}{K}%
\TextAt{9}{369.039}{273.098}{\Character{91}}%
\TextAt{13}{372.291}{273.098}{x}%
\TextAt{9}{378.943}{273.098}{\Character{93}}%
\TextAt{8}{386.098}{273.098}{be}%
\TextAt{8}{402.028}{273.098}{an}%
\TextAt{8}{418.293}{273.098}{irreducible}%
\TextAt{8}{476.174}{273.098}{polynomial}%
\TextAt{8}{79.226}{256.052}{with}%
\TextAt{13}{105.982}{256.052}{a}%
\TextAt{21}{112.127}{254.259}{n}%
\NativeMathAt{121.248}{256.052}{11.955}{\neq}%
\TextAt{9}{133.825}{256.052}{0}%
\TextAt{8}{139.69}{256.052}{.}%
\TextAt{8}{148.426}{256.052}{Then,}%
\TextAt{8}{182.355}{256.052}{if}%
\TextAt{16}{193.175}{256.052}{\Character{106}}%
\TextAt{16}{200.372}{256.052}{\Character{106}}%
\TextAt{8}{207.707}{256.052}{is}%
\TextAt{8}{219.567}{256.052}{a}%
\TextAt{8}{229.418}{256.052}{nonarchimedean}%
\TextAt{8}{316.345}{256.052}{valuation}%
\TextAt{8}{367.513}{256.052}{of}%
\TextAt{13}{380.939}{256.052}{K}%
\TextAt{8}{395.75}{256.052}{with}%
\TextAt{8}{422.506}{256.052}{a}%
\TextAt{8}{432.357}{256.052}{unique}%
\TextAt{8}{470.495}{256.052}{extension}%
\TextAt{8}{522.357}{256.052}{to}%
\TextAt{8}{79.226}{239.006}{the}%
\TextAt{8}{99.383}{239.006}{splitting}%
\TextAt{8}{145.625}{239.006}{field,}%
\TextAt{8}{174.234}{239.006}{one}%
\TextAt{8}{195.693}{239.006}{has}%
\TextAt{16}{251.191}{221.96}{\Character{106}}%
\TextAt{13}{254.511}{221.96}{f}%
\TextAt{16}{261.558}{221.96}{\Character{106}}%
\TextAt{9}{268.2}{221.96}{\Character{61}}%
\TextAt{8}{280.625}{221.96}{max}%
\TextAt{16}{302.419}{221.96}{\Character{102}\Character{106}}%
\TextAt{13}{311.718}{221.96}{a}%
\TextAt{17}{317.863}{220.167}{0}%
\TextAt{16}{322.595}{221.96}{\Character{106}}%
\TextAt{13}{327.909}{221.96}{\Character{59}}%
\TextAt{16}{333.153}{221.96}{\Character{106}}%
\TextAt{13}{336.474}{221.96}{a}%
\TextAt{21}{342.619}{220.167}{n}%
\TextAt{16}{348.255}{221.96}{\Character{106}\Character{103}}%
\TextAt{13}{357.553}{221.96}{\Character{58}}%
\TextAt{8}{97.159}{189.472}{In}%
\TextAt{8}{113.496}{189.472}{(4.7)}%
\TextAt{8}{143.187}{189.472}{we}%
\TextAt{8}{162.131}{189.472}{deduced}%
\TextAt{8}{209.36}{189.472}{this}%
\TextAt{8}{233.898}{189.472}{result}%
\TextAt{8}{268.192}{189.472}{for}%
\TextAt{8}{287.793}{189.472}{complete}%
\TextAt{8}{338.932}{189.472}{fields}%
\TextAt{8}{370.62}{189.472}{from}%
\TextAt{8}{399.976}{189.472}{Hensel’s}%
\TextAt{8}{447.002}{189.472}{lemma}%
\TextAt{8}{486.437}{189.472}{and}%
\TextAt{8}{510.915}{189.472}{thus}%
\TextAt{8}{79.226}{172.426}{obtained}%
\TextAt{8}{128.385}{172.426}{the}%
\TextAt{8}{149.569}{172.426}{uniqueness}%
\TextAt{8}{209.572}{172.426}{of}%
\TextAt{8}{223.929}{172.426}{the}%
\TextAt{8}{245.114}{172.426}{extended}%
\TextAt{8}{295.887}{172.426}{valuation.}%
\TextAt{8}{354.593}{172.426}{Here}%
\TextAt{8}{383.249}{172.426}{we}%
\TextAt{8}{401.504}{172.426}{obtain}%
\TextAt{8}{438.959}{172.426}{it,}%
\TextAt{8}{455.199}{172.426}{by}%
\TextAt{8}{472.486}{172.426}{contrast,}%
\TextAt{8}{522.297}{172.426}{as}%
\TextAt{8}{79.226}{155.38}{a}%
\TextAt{8}{89.596}{155.38}{consequence}%
\TextAt{8}{156.274}{155.38}{of}%
\TextAt{8}{170.218}{155.38}{the}%
\TextAt{8}{190.988}{155.38}{uniqueness}%
\TextAt{8}{250.577}{155.38}{of}%
\TextAt{8}{264.521}{155.38}{the}%
\TextAt{8}{285.291}{155.38}{extended}%
\TextAt{8}{335.651}{155.38}{valuation.}%
\TextAt{8}{393.116}{155.38}{We}%
\TextAt{8}{413.875}{155.38}{now}%
\TextAt{8}{438.877}{155.38}{proceed}%
\TextAt{8}{482.733}{155.38}{to}%
\TextAt{8}{497.657}{155.38}{deduce}%
\TextAt{8}{79.226}{138.334}{Hensel’s}%
\TextAt{8}{124.536}{138.334}{lemma}%
\TextAt{8}{162.255}{138.334}{from}%
\TextAt{8}{189.895}{138.334}{the}%
\TextAt{8}{210.052}{138.334}{unique}%
\TextAt{8}{248.093}{138.334}{extendability.}%
\TextAt{12}{79.226}{112.322}{Theorem}%
\TextAt{12}{137.285}{112.322}{(6.6).}%
\TextAt{8}{173.509}{112.322}{A}%
\TextAt{8}{187.088}{112.322}{nonarchimedean}%
\TextAt{8}{274.827}{112.322}{valued}%
\TextAt{8}{312.485}{112.322}{field}%
\TextAt{9}{338.75}{112.322}{\Character{40}}%
\TextAt{13}{343.302}{112.322}{K}%
\TextAt{13}{353.464}{112.322}{\Character{59}}%
\TextAt{16}{358.714}{112.322}{\Character{106}}%
\TextAt{16}{365.911}{112.322}{\Character{106}}%
\TextAt{9}{369.253}{112.322}{\Character{41}}%
\TextAt{8}{378.611}{112.322}{is}%
\TextAt{8}{391.283}{112.322}{henselian}%
\TextAt{8}{442.976}{112.322}{if}%
\TextAt{8}{454.608}{112.322}{and}%
\TextAt{8}{478.279}{112.322}{only}%
\TextAt{8}{504.879}{112.322}{if}%
\TextAt{8}{516.51}{112.322}{the}%
\TextAt{8}{79.226}{95.276}{valuation}%
\TextAt{16}{130.299}{95.276}{\Character{106}}%
\TextAt{16}{137.496}{95.276}{\Character{106}}%
\TextAt{8}{144.735}{95.276}{can}%
\TextAt{8}{166.195}{95.276}{be}%
\TextAt{8}{182.119}{95.276}{uniquely}%
\TextAt{8}{229.593}{95.276}{extended}%
\TextAt{8}{279.338}{95.276}{to}%
\TextAt{8}{293.649}{95.276}{any}%
\TextAt{8}{315.766}{95.276}{algebraic}%
\TextAt{8}{365.2}{95.276}{extension.}%
\TextAt{8}{297.211}{43.033}{134}%
\Destination{78.23}{809.815}{nav.page.134}%
\Destination{79.36}{681.455}{nav.print.147}%
\Destination{79.36}{680.455}{nav.number.2.6.4}%
\Destination{79.36}{287.055}{nav.number.2.6.5}%
\Destination{79.36}{126.355}{nav.number.2.6.6}%
\Hotlink{112.996}{187.255}{25.078}{11.047}{nav.number.2.4.7}%
}

% Page 147 (body page 135)
\Page{%
\TextAt{2}{79.226}{796.943}{II.}%
\TextAt{2}{94.073}{796.943}{The}%
\TextAt{2}{116.492}{796.943}{Theory}%
\TextAt{2}{154.401}{796.943}{of}%
\TextAt{2}{166.826}{796.943}{Valuations}%
\Vector{q 0 G 0 g 0.299 w 7.226 18.729 m 460.769 18.729 l S Q}%
\TextAt{19}{79.226}{761.124}{Proof.}%
\TextAt{8}{116}{761.124}{The}%
\TextAt{8}{140.721}{761.124}{fact}%
\TextAt{8}{164.474}{761.124}{that}%
\TextAt{8}{190.51}{761.124}{a}%
\TextAt{8}{200.932}{761.124}{henselian}%
\TextAt{8}{252.385}{761.124}{valuation}%
\TextAt{16}{304.125}{761.124}{\Character{106}}%
\TextAt{16}{311.322}{761.124}{\Character{106}}%
\TextAt{8}{319.228}{761.124}{extends}%
\TextAt{8}{362.552}{761.124}{uniquely}%
\TextAt{8}{410.693}{761.124}{was}%
\TextAt{8}{433.859}{761.124}{dealt}%
\TextAt{8}{463.793}{761.124}{with}%
\TextAt{8}{491.12}{761.124}{in}%
\TextAt{8}{505.44}{761.124}{(6.2).}%
\TextAt{8}{79.226}{744.078}{Let}%
\TextAt{8}{99.342}{744.078}{us}%
\TextAt{8}{113.505}{744.078}{assume}%
\TextAt{8}{153.096}{744.078}{conversely}%
\TextAt{8}{208.241}{744.078}{that}%
\TextAt{16}{232.756}{744.078}{\Character{106}}%
\TextAt{16}{239.965}{744.078}{\Character{106}}%
\TextAt{8}{246.34}{744.078}{admits}%
\TextAt{8}{283.922}{744.078}{one}%
\TextAt{8}{304.528}{744.078}{and}%
\TextAt{8}{326.438}{744.078}{only}%
\TextAt{8}{351.276}{744.078}{one}%
\TextAt{8}{371.883}{744.078}{extension}%
\TextAt{8}{422.795}{744.078}{to}%
\TextAt{8}{436.252}{744.078}{any}%
\TextAt{8}{457.516}{744.078}{given}%
\TextAt{8}{487.232}{744.078}{algebraic}%
\TextAt{8}{79.226}{727.032}{extension.}%
\TextAt{8}{135.543}{727.032}{We}%
\TextAt{8}{155.687}{727.032}{first}%
\TextAt{8}{179.813}{727.032}{show:}%
\TextAt{8}{97.159}{706.997}{Let}%
\TextAt{13}{118.625}{706.997}{f}%
\TextAt{9}{125.672}{706.997}{\Character{40}}%
\TextAt{13}{130.224}{706.997}{x}%
\TextAt{9}{136.876}{706.997}{\Character{41}}%
\TextAt{9}{145.579}{706.997}{\Character{61}}%
\TextAt{13}{158.869}{706.997}{a}%
\TextAt{17}{165.014}{705.204}{0}%
\TextAt{9}{172.742}{706.997}{\Character{43}}%
\TextAt{13}{184.842}{706.997}{a}%
\TextAt{17}{190.987}{705.204}{1}%
\TextAt{13}{195.719}{706.997}{x}%
\TextAt{9}{205.367}{706.997}{\Character{43}}%
\TextAt{16}{217.467}{706.997}{\Character{1}}%
\TextAt{16}{222.763}{706.997}{\Character{1}}%
\TextAt{16}{228.071}{706.997}{\Character{1}}%
\TextAt{9}{234.41}{706.997}{\Character{43}}%
\TextAt{13}{246.51}{706.997}{a}%
\TextAt{21}{252.655}{705.204}{n}%
\TextAt{13}{258.292}{706.997}{x}%
\TextAt{21}{264.944}{711.335}{n}%
\NativeMathAt{274.748}{706.997}{11.955}{\in}%
\TextAt{14}{286.886}{706.997}{o}%
\TextAt{9}{292.73}{706.997}{\Character{91}}%
\TextAt{13}{295.982}{706.997}{x}%
\TextAt{9}{302.634}{706.997}{\Character{93}}%
\TextAt{8}{310.28}{706.997}{be}%
\TextAt{8}{326.702}{706.997}{a}%
\TextAt{8}{336.954}{706.997}{primitive,}%
\TextAt{8}{390.908}{706.997}{irreducible}%
\TextAt{8}{449.28}{706.997}{polynomial}%
\TextAt{8}{510.27}{706.997}{such}%
\TextAt{8}{79.226}{689.951}{that}%
\TextAt{13}{104.726}{689.951}{a}%
\TextAt{17}{110.871}{688.158}{0}%
\TextAt{13}{115.603}{689.951}{a}%
\TextAt{21}{121.748}{688.158}{n}%
\NativeMathAt{130.929}{689.951}{11.955}{\neq}%
\TextAt{9}{143.566}{689.951}{0}%
\TextAt{8}{149.43}{689.951}{,}%
\TextAt{8}{156.743}{689.951}{and}%
\TextAt{8}{179.637}{689.951}{let}%
\TextAt{9}{199.22}{693.106}{\Character{22}}%
\TextAt{13}{196.672}{689.951}{f}%
\TextAt{9}{203.719}{689.951}{\Character{40}}%
\TextAt{13}{208.271}{689.951}{x}%
\TextAt{9}{214.923}{689.951}{\Character{41}}%
\TextAt{9}{223.005}{689.951}{\Character{61}}%
\TextAt{13}{235.668}{689.951}{f}%
\TextAt{9}{242.715}{689.951}{\Character{40}}%
\TextAt{13}{247.267}{689.951}{x}%
\TextAt{9}{253.919}{689.951}{\Character{41}}%
\TextAt{8}{261.792}{689.951}{mod}%
\TextAt{14}{287.23}{689.951}{p}%
\NativeMathAt{296.757}{689.951}{11.955}{\in}%
\TextAt{13}{308.271}{689.951}{\Character{20}}%
\TextAt{9}{315.01}{689.951}{\Character{91}}%
\TextAt{13}{318.262}{689.951}{x}%
\TextAt{9}{324.914}{689.951}{\Character{93}}%
\TextAt{8}{328.165}{689.951}{.}%
\TextAt{8}{337.007}{689.951}{Then}%
\TextAt{8}{367.695}{689.951}{we}%
\TextAt{8}{385.053}{689.951}{have}%
\TextAt{8}{412.179}{689.951}{deg}%
\TextAt{9}{429.741}{689.951}{\Character{40}}%
\TextAt{9}{436.841}{693.106}{\Character{22}}%
\TextAt{13}{434.294}{689.951}{f}%
\TextAt{9}{441.34}{689.951}{\Character{41}}%
\TextAt{9}{449.422}{689.951}{\Character{61}}%
\TextAt{9}{462.058}{689.951}{0}%
\TextAt{8}{471.966}{689.951}{or}%
\TextAt{8}{486.408}{689.951}{deg}%
\TextAt{9}{503.97}{689.951}{\Character{40}}%
\TextAt{9}{511.07}{693.106}{\Character{22}}%
\TextAt{13}{508.522}{689.951}{f}%
\TextAt{9}{515.569}{689.951}{\Character{41}}%
\TextAt{9}{523.651}{689.951}{\Character{61}}%
\TextAt{8}{79.226}{672.905}{deg}%
\TextAt{9}{96.788}{672.905}{\Character{40}}%
\TextAt{13}{101.341}{672.905}{f}%
\TextAt{9}{108.387}{672.905}{\Character{41}}%
\TextAt{8}{112.939}{672.905}{,}%
\TextAt{8}{120.089}{672.905}{and}%
\TextAt{8}{142.851}{672.905}{we}%
\TextAt{8}{160.079}{672.905}{find}%
\TextAt{9}{268.866}{659.014}{\Character{22}}%
\TextAt{13}{266.318}{655.859}{f}%
\TextAt{9}{273.364}{655.859}{\Character{40}}%
\TextAt{13}{277.917}{655.859}{x}%
\TextAt{9}{284.569}{655.859}{\Character{41}}%
\TextAt{9}{292.424}{655.859}{\Character{61}}%
\TextAt{9}{304.988}{655.859}{\Character{22}}%
\TextAt{13}{304.867}{655.859}{a}%
\TextAt{9}{312.896}{655.859}{\Character{22}}%
\TextAt{13}{311.012}{655.859}{\Character{39}}%
\TextAt{9}{318.681}{655.859}{\Character{40}}%
\TextAt{13}{323.233}{655.859}{x}%
\TextAt{9}{329.885}{655.859}{\Character{41}}%
\TextAt{21}{334.437}{660.795}{m}%
\TextAt{13}{342.426}{655.859}{\Character{59}}%
\TextAt{8}{79.226}{632.337}{for}%
\TextAt{8}{97.111}{632.337}{some}%
\TextAt{8}{126.437}{632.337}{irreducible}%
\TextAt{8}{184.312}{632.337}{polynomial}%
\TextAt{9}{246.688}{632.337}{\Character{22}}%
\TextAt{13}{244.805}{632.337}{\Character{39}}%
\TextAt{9}{252.474}{632.337}{\Character{40}}%
\TextAt{13}{257.026}{632.337}{x}%
\TextAt{9}{263.678}{632.337}{\Character{41}}%
\NativeMathAt{271.551}{632.337}{11.955}{\in}%
\TextAt{13}{282.842}{632.337}{\Character{20}}%
\TextAt{9}{289.581}{632.337}{\Character{91}}%
\TextAt{13}{292.832}{632.337}{x}%
\TextAt{9}{299.484}{632.337}{\Character{93}}%
\TextAt{8}{306.633}{632.337}{and}%
\TextAt{8}{329.396}{632.337}{a}%
\TextAt{8}{339.152}{632.337}{constant}%
\TextAt{9}{386.52}{632.337}{\Character{22}}%
\TextAt{13}{386.374}{632.337}{a}%
\TextAt{8}{392.519}{632.337}{.}%
\TextAt{8}{97.159}{612.303}{As}%
\TextAt{13}{116.374}{612.303}{f}%
\TextAt{8}{129.246}{612.303}{is}%
\TextAt{8}{142.939}{612.303}{irreducible,}%
\TextAt{8}{206.476}{612.303}{the}%
\TextAt{8}{228.561}{612.303}{Newton}%
\TextAt{8}{273.731}{612.303}{polygon}%
\TextAt{8}{319.893}{612.303}{is}%
\TextAt{8}{333.585}{612.303}{a}%
\TextAt{8}{345.269}{612.303}{single}%
\TextAt{8}{379.775}{612.303}{line}%
\TextAt{8}{403.808}{612.303}{segment}%
\TextAt{8}{450.999}{612.303}{and}%
\TextAt{8}{475.69}{612.303}{thus}%
\TextAt{16}{503.37}{612.303}{\Character{106}}%
\TextAt{13}{506.691}{612.303}{f}%
\TextAt{16}{513.737}{612.303}{\Character{106}}%
\TextAt{9}{523.665}{612.303}{\Character{61}}%
\TextAt{8}{79.226}{595.257}{max}%
\TextAt{16}{101.02}{595.257}{\Character{102}\Character{106}}%
\TextAt{13}{110.319}{595.257}{a}%
\TextAt{17}{116.464}{593.463}{0}%
\TextAt{16}{121.196}{595.257}{\Character{106}}%
\TextAt{13}{126.51}{595.257}{\Character{59}}%
\TextAt{16}{131.754}{595.257}{\Character{106}}%
\TextAt{13}{135.075}{595.257}{a}%
\TextAt{21}{141.22}{593.463}{n}%
\TextAt{16}{146.856}{595.257}{\Character{106}\Character{103}}%
\TextAt{8}{156.154}{595.257}{.}%
\TextAt{8}{164.398}{595.257}{We}%
\TextAt{8}{183.929}{595.257}{may}%
\TextAt{8}{208.684}{595.257}{assume}%
\TextAt{8}{248.515}{595.257}{that}%
\TextAt{13}{273.27}{595.257}{a}%
\TextAt{21}{279.415}{593.463}{n}%
\TextAt{8}{288.335}{595.257}{is}%
\TextAt{8}{299.485}{595.257}{a}%
\TextAt{8}{308.627}{595.257}{unit,}%
\TextAt{8}{336.099}{595.257}{because}%
\TextAt{8}{378.787}{595.257}{otherwise}%
\TextAt{8}{430.262}{595.257}{the}%
\TextAt{8}{449.805}{595.257}{Newton}%
\TextAt{8}{492.433}{595.257}{polygon}%
\TextAt{8}{79.226}{578.211}{is}%
\TextAt{8}{90.99}{578.211}{a}%
\TextAt{8}{100.745}{578.211}{segment}%
\TextAt{8}{146.008}{578.211}{which}%
\TextAt{8}{179.494}{578.211}{does}%
\TextAt{8}{205.891}{578.211}{not}%
\TextAt{8}{226.705}{578.211}{lie}%
\TextAt{8}{242.307}{578.211}{on}%
\TextAt{8}{258.566}{578.211}{the}%
\TextAt{13}{278.722}{578.211}{x}%
\TextAt{8}{285.374}{578.211}{-axis}%
\TextAt{8}{313.074}{578.211}{and}%
\TextAt{8}{335.837}{578.211}{this}%
\TextAt{8}{358.659}{578.211}{means}%
\TextAt{8}{394.489}{578.211}{that}%
\TextAt{9}{422.406}{581.365}{\Character{22}}%
\TextAt{13}{419.858}{578.211}{f}%
\TextAt{9}{426.904}{578.211}{\Character{40}}%
\TextAt{13}{431.457}{578.211}{x}%
\TextAt{9}{438.109}{578.211}{\Character{41}}%
\TextAt{9}{445.964}{578.211}{\Character{61}}%
\TextAt{9}{458.529}{578.211}{\Character{22}}%
\TextAt{13}{458.407}{578.211}{a}%
\TextAt{17}{464.552}{576.417}{0}%
\TextAt{8}{469.285}{578.211}{.}%
\TextAt{8}{97.159}{558.176}{Let}%
\TextAt{13}{119.835}{558.176}{L}%
\TextAt{16}{134.03}{558.176}{\Character{106}}%
\TextAt{13}{143.581}{558.176}{K}%
\TextAt{8}{160.003}{558.176}{be}%
\TextAt{8}{177.635}{558.176}{the}%
\TextAt{8}{199.498}{558.176}{splitting}%
\TextAt{8}{247.448}{558.176}{field}%
\TextAt{8}{274.512}{558.176}{of}%
\TextAt{13}{289.55}{558.176}{f}%
\TextAt{9}{296.596}{558.176}{\Character{40}}%
\TextAt{13}{301.148}{558.176}{x}%
\TextAt{9}{307.8}{558.176}{\Character{41}}%
\TextAt{8}{317.957}{558.176}{over}%
\TextAt{13}{344.711}{558.176}{K}%
\TextAt{8}{361.133}{558.176}{and}%
\TextAt{16}{385.603}{558.176}{\Character{79}}%
\TextAt{8}{401.057}{558.176}{the}%
\TextAt{8}{422.921}{558.176}{valuation}%
\TextAt{8}{475.7}{558.176}{ring}%
\TextAt{8}{501.473}{558.176}{of}%
\TextAt{8}{516.51}{558.176}{the}%
\TextAt{8}{79.226}{541.13}{unique}%
\TextAt{8}{118.939}{541.13}{extension}%
\TextAt{16}{172.376}{541.13}{\Character{106}}%
\TextAt{16}{179.585}{541.13}{\Character{106}}%
\TextAt{8}{188.484}{541.13}{to}%
\TextAt{13}{204.466}{541.13}{L}%
\TextAt{8}{212.43}{541.13}{,}%
\TextAt{8}{221.669}{541.13}{with}%
\TextAt{8}{250}{541.13}{maximal}%
\TextAt{8}{299.48}{541.13}{ideal}%
\TextAt{14}{329.115}{541.13}{P}%
\TextAt{8}{339.008}{541.13}{.}%
\TextAt{8}{352.47}{541.13}{For}%
\TextAt{8}{375.111}{541.13}{an}%
\TextAt{8}{393.042}{541.13}{arbitrary}%
\TextAt{13}{444.482}{541.13}{K}%
\TextAt{8}{455.3}{541.13}{-automorphism}%
\TextAt{13}{79.226}{524.084}{\Character{27}}%
\NativeMathAt{91.432}{524.084}{11.955}{\in}%
\TextAt{13}{104.526}{524.084}{G}%
\TextAt{9}{118.883}{524.084}{\Character{61}}%
\TextAt{13}{133.111}{524.084}{G}%
\TextAt{9}{142.345}{524.084}{\Character{40}}%
\TextAt{13}{146.897}{524.084}{L}%
\TextAt{16}{159.985}{524.084}{\Character{106}}%
\TextAt{13}{168.43}{524.084}{K}%
\TextAt{9}{179.247}{524.084}{\Character{41}}%
\TextAt{8}{183.8}{524.084}{,}%
\TextAt{8}{192.271}{524.084}{we}%
\TextAt{8}{210.557}{524.084}{have}%
\TextAt{16}{238.609}{524.084}{\Character{106}}%
\TextAt{13}{241.93}{524.084}{\Character{27}}%
\TextAt{13}{248.996}{524.084}{\Character{11}}%
\TextAt{16}{256.534}{524.084}{\Character{106}}%
\TextAt{9}{264.979}{524.084}{\Character{61}}%
\TextAt{16}{279.207}{524.084}{\Character{106}}%
\TextAt{13}{282.528}{524.084}{\Character{11}}%
\TextAt{16}{290.05}{524.084}{\Character{106}}%
\TextAt{8}{298.326}{524.084}{for}%
\TextAt{8}{317.269}{524.084}{all}%
\TextAt{13}{334.585}{524.084}{\Character{11}}%
\NativeMathAt{347.231}{524.084}{11.955}{\in}%
\TextAt{13}{360.324}{524.084}{L}%
\TextAt{8}{368.289}{524.084}{,}%
\TextAt{8}{376.76}{524.084}{because}%
\TextAt{16}{421.12}{524.084}{\Character{106}}%
\TextAt{16}{428.317}{524.084}{\Character{106}}%
\TextAt{8}{436.614}{524.084}{and}%
\TextAt{8}{460.435}{524.084}{the}%
\TextAt{8}{481.649}{524.084}{composite}%
\TextAt{16}{79.226}{507.038}{\Character{106}}%
\TextAt{16}{86.435}{507.038}{\Character{106}}%
\TextAt{16}{93.058}{507.038}{\Character{14}}%
\TextAt{13}{102.382}{507.038}{\Character{27}}%
\TextAt{8}{114.334}{507.038}{extend}%
\TextAt{8}{153.347}{507.038}{the}%
\TextAt{8}{174.476}{507.038}{same}%
\TextAt{8}{204.774}{507.038}{valuation.}%
\TextAt{8}{263.314}{507.038}{This}%
\TextAt{8}{291.006}{507.038}{shows}%
\TextAt{8}{325.596}{507.038}{that}%
\TextAt{13}{351.937}{507.038}{\Character{27}}%
\TextAt{16}{359.019}{507.038}{\Character{79}}%
\TextAt{9}{373.847}{507.038}{\Character{61}}%
\TextAt{16}{387.929}{507.038}{\Character{79}}%
\TextAt{8}{397.779}{507.038}{,}%
\TextAt{13}{406.143}{507.038}{\Character{27}}%
\TextAt{14}{413.225}{507.038}{P}%
\TextAt{9}{428.096}{507.038}{\Character{61}}%
\TextAt{14}{442.178}{507.038}{P}%
\TextAt{8}{452.072}{507.038}{.}%
\TextAt{8}{463.436}{507.038}{If}%
\TextAt{13}{476.101}{507.038}{\Character{11}}%
\TextAt{8}{488.492}{507.038}{is}%
\TextAt{8}{501.228}{507.038}{a}%
\TextAt{8}{511.956}{507.038}{zero}%
\TextAt{8}{79.226}{489.992}{of}%
\TextAt{13}{94.091}{489.992}{f}%
\TextAt{9}{101.138}{489.992}{\Character{40}}%
\TextAt{13}{105.69}{489.992}{x}%
\TextAt{9}{112.342}{489.992}{\Character{41}}%
\TextAt{8}{122.327}{489.992}{and}%
\TextAt{13}{146.625}{489.992}{\Character{22}}%
\TextAt{8}{159.1}{489.992}{its}%
\TextAt{8}{176.954}{489.992}{multiplicity,}%
\TextAt{8}{243.91}{489.992}{then}%
\TextAt{13}{272.105}{489.992}{\Character{27}}%
\TextAt{13}{279.171}{489.992}{\Character{11}}%
\NativeMathAt{292.646}{489.992}{11.955}{\in}%
\TextAt{16}{306.538}{489.992}{\Character{79}}%
\TextAt{8}{321.835}{489.992}{for}%
\TextAt{8}{341.256}{489.992}{all}%
\TextAt{13}{359.05}{489.992}{\Character{27}}%
\NativeMathAt{372.069}{489.992}{11.955}{\in}%
\TextAt{13}{385.976}{489.992}{G}%
\TextAt{8}{395.21}{489.992}{.}%
\TextAt{8}{408.264}{489.992}{Indeed,}%
\TextAt{8}{451.464}{489.992}{if}%
\TextAt{13}{463.723}{489.992}{\Character{11}}%
\NativeMathAt{477.182}{489.992}{11.955}{\notin}%
\TextAt{16}{491.074}{489.992}{\Character{79}}%
\TextAt{8}{500.939}{489.992}{,}%
\TextAt{8}{510.007}{489.992}{then}%
\TextAt{20}{79.226}{481.913}{\Character{81}}%
\TextAt{21}{90.517}{469.459}{\Character{27}}%
\TextAt{16}{98.133}{472.946}{\Character{106}}%
\TextAt{13}{101.454}{472.946}{\Character{11}}%
\TextAt{16}{108.976}{472.946}{\Character{106}}%
\TextAt{21}{112.297}{478.758}{\Character{22}}%
\TextAt{9}{122.551}{472.946}{\Character{61}}%
\TextAt{20}{136.318}{481.912}{\Character{81}}%
\TextAt{21}{147.609}{469.459}{\Character{27}}%
\TextAt{16}{155.225}{472.946}{\Character{106}}%
\TextAt{13}{158.545}{472.946}{\Character{27}}%
\TextAt{13}{165.611}{472.946}{\Character{11}}%
\TextAt{16}{173.15}{472.946}{\Character{106}}%
\TextAt{21}{176.47}{478.758}{\Character{22}}%
\TextAt{13}{186.725}{472.946}{\Character{62}}%
\TextAt{9}{200.491}{472.946}{1}%
\TextAt{8}{211.029}{472.946}{would}%
\TextAt{8}{245.96}{472.946}{imply}%
\TextAt{8}{279.587}{472.946}{that}%
\TextAt{8}{305.743}{472.946}{the}%
\TextAt{8}{326.687}{472.946}{constant}%
\TextAt{8}{374.697}{472.946}{coefficient}%
\TextAt{13}{430.107}{472.946}{a}%
\TextAt{17}{436.252}{471.153}{0}%
\TextAt{8}{445.668}{472.946}{could}%
\TextAt{8}{477.67}{472.946}{not}%
\TextAt{8}{499.271}{472.946}{belong}%
\TextAt{8}{79.226}{455.9}{to}%
\TextAt{14}{93.981}{455.9}{o}%
\TextAt{8}{99.826}{455.9}{.}%
\TextAt{8}{109.608}{455.9}{Thus}%
\TextAt{8}{139.702}{455.9}{every}%
\TextAt{13}{171.039}{455.9}{\Character{27}}%
\NativeMathAt{182.2}{455.9}{11.955}{\in}%
\TextAt{13}{194.249}{455.9}{G}%
\TextAt{8}{207.825}{455.9}{induces}%
\TextAt{8}{249.945}{455.9}{a}%
\TextAt{13}{260.145}{455.9}{\Character{20}}%
\TextAt{8}{266.884}{455.9}{-automorphism}%
\TextAt{9}{349.31}{455.9}{\Character{22}}%
\TextAt{13}{348.696}{455.9}{\Character{27}}%
\TextAt{8}{360.12}{455.9}{of}%
\TextAt{16}{373.895}{455.9}{\Character{79}}%
\TextAt{9}{383.745}{455.9}{\Character{47}}%
\TextAt{14}{389.598}{455.9}{P}%
\TextAt{8}{399.491}{455.9}{,}%
\TextAt{8}{407.196}{455.9}{and}%
\TextAt{8}{430.404}{455.9}{the}%
\TextAt{8}{451.005}{455.9}{zeroes}%
\Vector{q 0 G 0 g 0.478 w 414.299 -311.279 m 428.903 -311.279 l S Q}%
\TextAt{13}{486.299}{455.9}{\Character{27}}%
\TextAt{13}{493.365}{455.9}{\Character{11}}%
\TextAt{9}{504.982}{455.9}{\Character{61}}%
\TextAt{9}{518.766}{455.9}{\Character{22}}%
\TextAt{13}{518.165}{455.9}{\Character{27}}%
\TextAt{9}{526.407}{455.9}{\Character{22}}%
\TextAt{13}{525.248}{455.9}{\Character{11}}%
\TextAt{8}{79.226}{438.854}{of}%
\TextAt{9}{95.741}{442.009}{\Character{22}}%
\TextAt{13}{93.194}{438.854}{f}%
\TextAt{9}{100.24}{438.854}{\Character{40}}%
\TextAt{13}{104.793}{438.854}{x}%
\TextAt{9}{111.445}{438.854}{\Character{41}}%
\TextAt{8}{120.532}{438.854}{are}%
\TextAt{8}{140.68}{438.854}{all}%
\TextAt{8}{157.577}{438.854}{conjugate}%
\TextAt{8}{211.224}{438.854}{over}%
\TextAt{13}{236.907}{438.854}{\Character{20}}%
\TextAt{8}{243.646}{438.854}{.}%
\TextAt{8}{254.007}{438.854}{It}%
\TextAt{8}{267.317}{438.854}{follows}%
\TextAt{8}{306.391}{438.854}{that}%
\TextAt{9}{334.945}{442.009}{\Character{22}}%
\TextAt{13}{332.397}{438.854}{f}%
\TextAt{9}{339.444}{438.854}{\Character{40}}%
\TextAt{13}{343.996}{438.854}{x}%
\TextAt{9}{350.648}{438.854}{\Character{41}}%
\TextAt{9}{359.591}{438.854}{\Character{61}}%
\TextAt{9}{373.243}{438.854}{\Character{22}}%
\TextAt{13}{373.12}{438.854}{a}%
\TextAt{9}{381.148}{438.854}{\Character{22}}%
\TextAt{13}{379.265}{438.854}{\Character{39}}%
\TextAt{9}{386.933}{438.854}{\Character{40}}%
\TextAt{13}{391.485}{438.854}{x}%
\TextAt{9}{398.137}{438.854}{\Character{41}}%
\TextAt{21}{402.69}{443.193}{m}%
\TextAt{8}{410.678}{438.854}{,}%
\TextAt{8}{418.625}{438.854}{if}%
\TextAt{9}{431.869}{438.854}{\Character{22}}%
\TextAt{13}{429.986}{438.854}{\Character{39}}%
\TextAt{9}{437.654}{438.854}{\Character{40}}%
\TextAt{13}{442.207}{438.854}{x}%
\TextAt{9}{448.859}{438.854}{\Character{41}}%
\TextAt{8}{457.946}{438.854}{is}%
\TextAt{8}{470.348}{438.854}{the}%
\TextAt{8}{491.142}{438.854}{minimal}%
\TextAt{8}{79.226}{421.808}{polynomial}%
\TextAt{8}{139.719}{421.808}{of}%
\TextAt{9}{154.209}{421.808}{\Character{22}}%
\TextAt{13}{153.049}{421.808}{\Character{11}}%
\TextAt{8}{164.468}{421.808}{over}%
\TextAt{13}{189.514}{421.808}{\Character{20}}%
\TextAt{8}{196.253}{421.808}{.}%
\TextAt{8}{204.701}{421.808}{Since}%
\TextAt{13}{235.259}{421.808}{a}%
\TextAt{21}{241.404}{420.015}{n}%
\NativeMathAt{250.361}{421.808}{11.955}{\in}%
\TextAt{14}{261.652}{421.808}{o}%
\TextAt{3}{267.496}{426.147}{\Character{3}}%
\TextAt{8}{272.228}{421.808}{,}%
\TextAt{8}{279.378}{421.808}{we}%
\TextAt{8}{296.605}{421.808}{have}%
\TextAt{8}{323.6}{421.808}{furthermore}%
\TextAt{8}{388.313}{421.808}{that}%
\TextAt{8}{413.682}{421.808}{deg}%
\TextAt{9}{431.244}{421.808}{\Character{40}}%
\TextAt{9}{438.344}{424.963}{\Character{22}}%
\TextAt{13}{435.796}{421.808}{f}%
\TextAt{9}{442.843}{421.808}{\Character{41}}%
\TextAt{9}{450.698}{421.808}{\Character{61}}%
\TextAt{8}{463.142}{421.808}{deg}%
\TextAt{9}{480.704}{421.808}{\Character{40}}%
\TextAt{13}{485.256}{421.808}{f}%
\TextAt{9}{492.302}{421.808}{\Character{41}}%
\TextAt{8}{496.855}{421.808}{.}%
\TextAt{8}{97.159}{401.773}{Let}%
\TextAt{8}{118.128}{401.773}{now}%
\TextAt{13}{142.517}{401.773}{f}%
\TextAt{9}{149.563}{401.773}{\Character{40}}%
\TextAt{13}{154.116}{401.773}{x}%
\TextAt{9}{160.768}{401.773}{\Character{41}}%
\NativeMathAt{168.641}{401.773}{11.955}{\in}%
\TextAt{14}{179.932}{401.773}{o}%
\TextAt{9}{185.776}{401.773}{\Character{91}}%
\TextAt{13}{189.028}{401.773}{x}%
\TextAt{9}{195.68}{401.773}{\Character{93}}%
\TextAt{8}{202.829}{401.773}{be}%
\TextAt{8}{218.753}{401.773}{an}%
\TextAt{8}{235.012}{401.773}{arbitrary}%
\TextAt{8}{284.782}{401.773}{primitive}%
\TextAt{8}{334.862}{401.773}{polynomial,}%
\TextAt{8}{398.607}{401.773}{and}%
\TextAt{8}{421.369}{401.773}{let}%
\TextAt{13}{251.586}{372.772}{f}%
\TextAt{9}{258.633}{372.772}{\Character{40}}%
\TextAt{13}{263.185}{372.772}{x}%
\TextAt{9}{269.837}{372.772}{\Character{41}}%
\TextAt{9}{277.692}{372.772}{\Character{61}}%
\TextAt{13}{290.136}{372.772}{f}%
\TextAt{17}{295.906}{370.979}{1}%
\TextAt{9}{300.638}{372.772}{\Character{40}}%
\TextAt{13}{305.191}{372.772}{x}%
\TextAt{9}{311.843}{372.772}{\Character{41}}%
\TextAt{16}{318.388}{372.772}{\Character{1}}%
\TextAt{16}{323.684}{372.772}{\Character{1}}%
\TextAt{16}{328.992}{372.772}{\Character{1}}%
\TextAt{13}{334.328}{372.772}{f}%
\TextAt{21}{340.098}{370.979}{r}%
\TextAt{9}{344.652}{372.772}{\Character{40}}%
\TextAt{13}{349.205}{372.772}{x}%
\TextAt{9}{355.857}{372.772}{\Character{41}}%
\TextAt{8}{79.226}{343.771}{be}%
\TextAt{8}{96.528}{343.771}{its}%
\TextAt{8}{114.224}{343.771}{factorization}%
\TextAt{8}{183.578}{343.771}{into}%
\TextAt{8}{208.699}{343.771}{irreducibles}%
\TextAt{8}{272.566}{343.771}{over}%
\TextAt{13}{298.989}{343.771}{K}%
\TextAt{8}{309.807}{343.771}{.}%
\TextAt{8}{322.387}{343.771}{Since}%
\TextAt{9}{354.322}{343.771}{1}%
\TextAt{9}{365.835}{343.771}{\Character{61}}%
\TextAt{16}{380.615}{343.771}{\Character{106}}%
\TextAt{13}{383.936}{343.771}{f}%
\TextAt{16}{390.982}{343.771}{\Character{106}}%
\TextAt{9}{399.971}{343.771}{\Character{61}}%
\TextAt{20}{414.744}{352.738}{\Character{81}}%
\TextAt{16}{428.027}{343.771}{\Character{106}}%
\TextAt{13}{431.348}{343.771}{f}%
\TextAt{21}{437.118}{341.978}{i}%
\TextAt{16}{440.5}{343.771}{\Character{106}}%
\TextAt{8}{443.82}{343.771}{,}%
\TextAt{8}{452.691}{343.771}{multiplying}%
\TextAt{8}{516.51}{343.771}{the}%
\TextAt{13}{79.226}{326.725}{f}%
\TextAt{21}{84.996}{324.932}{i}%
\TextAt{8}{93.55}{326.725}{by}%
\TextAt{8}{111.084}{326.725}{suitable}%
\TextAt{8}{155.995}{326.725}{constants}%
\TextAt{8}{209.108}{326.725}{yields}%
\TextAt{16}{243.284}{326.725}{\Character{106}}%
\TextAt{13}{246.605}{326.725}{f}%
\TextAt{21}{252.375}{324.932}{i}%
\TextAt{16}{255.756}{326.725}{\Character{106}}%
\TextAt{9}{264.571}{326.725}{\Character{61}}%
\TextAt{9}{279.156}{326.725}{1}%
\TextAt{8}{285.022}{326.725}{.}%
\TextAt{8}{297.296}{326.725}{The}%
\TextAt{13}{322.625}{326.725}{f}%
\TextAt{21}{328.395}{324.932}{i}%
\TextAt{9}{331.776}{326.725}{\Character{40}}%
\TextAt{13}{336.328}{326.725}{x}%
\TextAt{9}{342.98}{326.725}{\Character{41}}%
\TextAt{8}{352.705}{326.725}{are}%
\TextAt{8}{373.491}{326.725}{therefore}%
\TextAt{8}{423.866}{326.725}{primitive,}%
\TextAt{8}{478.792}{326.725}{irreducible}%
\TextAt{8}{79.226}{309.679}{polynomials}%
\TextAt{8}{144.334}{309.679}{in}%
\TextAt{14}{157.987}{309.679}{o}%
\TextAt{9}{163.831}{309.679}{\Character{91}}%
\TextAt{13}{167.083}{309.679}{x}%
\TextAt{9}{173.735}{309.679}{\Character{93}}%
\TextAt{8}{176.986}{309.679}{.}%
\TextAt{8}{185.435}{309.679}{It}%
\TextAt{8}{198.107}{309.679}{follows}%
\TextAt{8}{236.543}{309.679}{that}%
\TextAt{9}{252.508}{283.833}{\Character{22}}%
\TextAt{13}{249.961}{280.678}{f}%
\TextAt{9}{257.007}{280.678}{\Character{40}}%
\TextAt{13}{261.559}{280.678}{x}%
\TextAt{9}{268.212}{280.678}{\Character{41}}%
\TextAt{9}{276.067}{280.678}{\Character{61}}%
\TextAt{9}{291.058}{283.833}{\Character{22}}%
\TextAt{13}{288.51}{280.678}{f}%
\TextAt{17}{294.28}{278.885}{1}%
\TextAt{9}{299.012}{280.678}{\Character{40}}%
\TextAt{13}{303.565}{280.678}{x}%
\TextAt{9}{310.217}{280.678}{\Character{41}}%
\TextAt{16}{316.762}{280.678}{\Character{1}}%
\TextAt{16}{322.058}{280.678}{\Character{1}}%
\TextAt{16}{327.366}{280.678}{\Character{1}}%
\TextAt{9}{335.25}{283.833}{\Character{22}}%
\TextAt{13}{332.702}{280.678}{f}%
\TextAt{21}{338.472}{278.885}{r}%
\TextAt{9}{343.027}{280.678}{\Character{40}}%
\TextAt{13}{347.579}{280.678}{x}%
\TextAt{9}{354.231}{280.678}{\Character{41}}%
\TextAt{13}{358.783}{280.678}{\Character{59}}%
\TextAt{8}{79.226}{251.677}{where}%
\TextAt{8}{113.502}{251.677}{deg}%
\TextAt{9}{131.065}{251.677}{\Character{40}}%
\TextAt{9}{138.165}{254.832}{\Character{22}}%
\TextAt{13}{135.617}{251.677}{f}%
\TextAt{21}{141.387}{249.884}{i}%
\TextAt{9}{144.768}{251.677}{\Character{41}}%
\TextAt{9}{153.424}{251.677}{\Character{61}}%
\TextAt{9}{166.634}{251.677}{0}%
\TextAt{8}{176.876}{251.677}{or}%
\TextAt{8}{191.654}{251.677}{deg}%
\TextAt{9}{209.216}{251.677}{\Character{40}}%
\TextAt{9}{216.316}{254.832}{\Character{22}}%
\TextAt{13}{213.768}{251.677}{f}%
\TextAt{21}{219.538}{249.884}{i}%
\TextAt{9}{222.92}{251.677}{\Character{41}}%
\TextAt{9}{231.576}{251.677}{\Character{61}}%
\TextAt{8}{244.81}{251.677}{deg}%
\TextAt{9}{262.372}{251.677}{\Character{40}}%
\TextAt{13}{266.925}{251.677}{f}%
\TextAt{21}{272.695}{249.884}{i}%
\TextAt{9}{276.076}{251.677}{\Character{41}}%
\TextAt{8}{280.628}{251.677}{,}%
\TextAt{8}{288.361}{251.677}{and}%
\TextAt{9}{314.139}{254.832}{\Character{22}}%
\TextAt{13}{311.591}{251.677}{f}%
\TextAt{21}{317.361}{249.884}{i}%
\TextAt{8}{325.107}{251.677}{is,}%
\TextAt{8}{340.706}{251.677}{up}%
\TextAt{8}{358.078}{251.677}{to}%
\TextAt{8}{372.856}{251.677}{a}%
\TextAt{8}{383.078}{251.677}{constant}%
\TextAt{8}{430.768}{251.677}{factor,}%
\TextAt{8}{468.102}{251.677}{the}%
\TextAt{8}{488.726}{251.677}{power}%
\TextAt{8}{523.337}{251.677}{of}%
\TextAt{8}{79.226}{234.631}{an}%
\TextAt{8}{95.762}{234.631}{irreducible}%
\TextAt{8}{153.913}{234.631}{polynomial.}%
\TextAt{8}{219.787}{234.631}{If}%
\TextAt{9}{234.303}{237.786}{\Character{22}}%
\TextAt{13}{231.755}{234.631}{f}%
\TextAt{9}{242.594}{234.631}{\Character{61}}%
\TextAt{9}{255.888}{234.631}{\Character{22}}%
\TextAt{13}{255.49}{234.631}{g}%
\TextAt{9}{261.642}{237.786}{\Character{22}}%
\TextAt{13}{261.524}{234.631}{h}%
\TextAt{8}{272.437}{234.631}{is}%
\TextAt{8}{284.477}{234.631}{a}%
\TextAt{8}{294.509}{234.631}{factorization}%
\TextAt{8}{362.763}{234.631}{into}%
\TextAt{8}{386.782}{234.631}{relatively}%
\TextAt{8}{438.119}{234.631}{prime}%
\TextAt{8}{471.559}{234.631}{polynomials}%
\TextAt{9}{79.642}{217.585}{\Character{22}}%
\TextAt{13}{79.226}{217.585}{g}%
\TextAt{13}{85.24}{217.585}{\Character{59}}%
\TextAt{9}{90.622}{220.74}{\Character{22}}%
\TextAt{13}{90.505}{217.585}{h}%
\NativeMathAt{100.564}{217.585}{11.955}{\in}%
\TextAt{13}{111.855}{217.585}{\Character{20}}%
\TextAt{9}{118.593}{217.585}{\Character{91}}%
\TextAt{13}{121.845}{217.585}{x}%
\TextAt{9}{128.497}{217.585}{\Character{93}}%
\TextAt{8}{131.749}{217.585}{,}%
\TextAt{8}{138.898}{217.585}{then}%
\TextAt{8}{165.558}{217.585}{we}%
\TextAt{8}{182.786}{217.585}{must}%
\TextAt{8}{211.789}{217.585}{have}%
\TextAt{9}{235.23}{188.584}{\Character{22}}%
\TextAt{13}{234.814}{188.584}{g}%
\TextAt{9}{244.169}{188.584}{\Character{61}}%
\TextAt{9}{256.722}{188.584}{\Character{22}}%
\TextAt{13}{256.595}{188.584}{a}%
\TextAt{20}{264.732}{199.941}{\Character{89}}%
\TextAt{21}{265.959}{174.57}{i}%
\NativeMathAt{268.842}{174.57}{7.97}{\in}%
\TextAt{21}{274.488}{174.57}{I}%
\TextAt{9}{284.548}{191.739}{\Character{22}}%
\TextAt{13}{282.001}{188.584}{f}%
\TextAt{21}{287.771}{186.791}{i}%
\TextAt{13}{291.152}{188.584}{\Character{59}}%
\TextAt{9}{320.424}{191.739}{\Character{22}}%
\TextAt{13}{320.306}{188.584}{h}%
\TextAt{9}{330.366}{188.584}{\Character{61}}%
\TextAt{9}{342.353}{191.739}{\Character{22}}%
\TextAt{13}{342.791}{188.584}{b}%
\TextAt{20}{349.761}{199.941}{\Character{89}}%
\TextAt{21}{349.947}{174.57}{j}%
\NativeMathAt{353.831}{174.57}{7.97}{\in}%
\TextAt{21}{359.477}{174.57}{J}%
\TextAt{9}{369.577}{191.739}{\Character{22}}%
\TextAt{13}{367.029}{188.584}{f}%
\TextAt{21}{372.799}{186.791}{j}%
\TextAt{8}{79.226}{148.966}{where}%
\TextAt{9}{113.181}{148.966}{\Character{22}}%
\TextAt{13}{113.035}{148.966}{a\Character{59}}%
\TextAt{9}{123.987}{152.12}{\Character{22}}%
\TextAt{13}{124.424}{148.966}{b}%
\NativeMathAt{132.722}{148.966}{11.955}{\in}%
\TextAt{13}{144.013}{148.966}{\Character{20}}%
\TextAt{8}{154.649}{148.966}{and}%
\TextAt{16}{177.412}{148.966}{\Character{102}}%
\TextAt{9}{183.39}{148.966}{1}%
\TextAt{13}{189.243}{148.966}{\Character{59}}%
\TextAt{13}{194.467}{148.966}{\Character{58}}%
\TextAt{13}{199.704}{148.966}{\Character{58}}%
\TextAt{13}{204.928}{148.966}{\Character{58}}%
\TextAt{13}{210.165}{148.966}{\Character{59}}%
\TextAt{13}{215.401}{148.966}{r}%
\TextAt{16}{221.064}{148.966}{\Character{103}}%
\TextAt{9}{230.362}{148.966}{\Character{61}}%
\TextAt{13}{242.788}{148.966}{I}%
\TextAt{9}{253.906}{150.46}{\Character{95}}%
\TextAt{16}{251.547}{148.966}{\Character{91}}%
\TextAt{13}{262.174}{148.966}{J}%
\TextAt{8}{273.648}{148.966}{and}%
\TextAt{8}{296.411}{148.966}{deg}%
\TextAt{9}{313.973}{148.966}{\Character{40}}%
\TextAt{9}{321.073}{152.12}{\Character{22}}%
\TextAt{13}{318.526}{148.966}{f}%
\TextAt{21}{324.296}{147.172}{i}%
\TextAt{9}{327.677}{148.966}{\Character{41}}%
\TextAt{9}{335.532}{148.966}{\Character{61}}%
\TextAt{8}{347.976}{148.966}{deg}%
\TextAt{9}{365.538}{148.966}{\Character{40}}%
\TextAt{13}{370.09}{148.966}{f}%
\TextAt{21}{375.86}{147.172}{i}%
\TextAt{9}{379.241}{148.966}{\Character{41}}%
\TextAt{8}{387.691}{148.966}{for}%
\TextAt{13}{405.576}{148.966}{i}%
\NativeMathAt{412.89}{148.966}{11.955}{\in}%
\TextAt{13}{424.181}{148.966}{I}%
\TextAt{8}{430.284}{148.966}{.}%
\TextAt{8}{438.732}{148.966}{We}%
\TextAt{8}{458.877}{148.966}{now}%
\TextAt{8}{483.265}{148.966}{put}%
\TextAt{13}{233.188}{119.965}{g}%
\TextAt{9}{242.543}{119.965}{\Character{61}}%
\TextAt{13}{254.969}{119.965}{a}%
\TextAt{20}{263.106}{131.322}{\Character{89}}%
\TextAt{21}{264.333}{105.95}{i}%
\NativeMathAt{267.216}{105.95}{7.97}{\in}%
\TextAt{21}{272.862}{105.95}{I}%
\TextAt{13}{280.375}{119.965}{f}%
\TextAt{21}{286.145}{118.171}{i}%
\TextAt{13}{289.526}{119.965}{\Character{59}}%
\TextAt{13}{318.661}{119.965}{h}%
\TextAt{9}{328.74}{119.965}{\Character{61}}%
\TextAt{13}{341.165}{119.965}{b}%
\TextAt{20}{348.135}{131.322}{\Character{89}}%
\TextAt{21}{348.321}{105.95}{j}%
\NativeMathAt{352.205}{105.95}{7.97}{\in}%
\TextAt{21}{357.851}{105.95}{J}%
\TextAt{13}{365.404}{119.965}{f}%
\TextAt{21}{371.174}{118.171}{j}%
\TextAt{13}{375.556}{119.965}{\Character{59}}%
\TextAt{8}{79.226}{80.346}{for}%
\TextAt{13}{97.111}{80.346}{a\Character{59}}%
\TextAt{13}{108.48}{80.346}{b}%
\NativeMathAt{116.798}{80.346}{11.955}{\in}%
\TextAt{14}{128.089}{80.346}{o}%
\TextAt{3}{133.933}{84.685}{\Character{3}}%
\TextAt{8}{142.563}{80.346}{such}%
\TextAt{8}{168.96}{80.346}{that}%
\TextAt{13}{194.329}{80.346}{a}%
\NativeMathAt{203.795}{80.346}{11.955}{\equiv}%
\TextAt{9}{216.56}{80.346}{\Character{22}}%
\TextAt{13}{216.414}{80.346}{a}%
\TextAt{8}{222.559}{80.346}{,}%
\TextAt{13}{229.708}{80.346}{b}%
\NativeMathAt{238.006}{80.346}{11.955}{\equiv}%
\TextAt{9}{250.187}{83.501}{\Character{22}}%
\TextAt{13}{250.625}{80.346}{b}%
\TextAt{8}{258.923}{80.346}{mod}%
\TextAt{14}{284.361}{80.346}{p}%
\TextAt{8}{294.242}{80.346}{and}%
\TextAt{13}{317.004}{80.346}{f}%
\TextAt{9}{327.372}{80.346}{\Character{61}}%
\TextAt{13}{339.797}{80.346}{g}%
\TextAt{13}{345.81}{80.346}{h}%
\TextAt{8}{352.57}{80.346}{.}%
\Vector{q 0 G 0 g 0.399 w 452.334 -693.654 m 452.334 -685.584 l S Q}%
\Vector{q 0 G 0 g 0.399 w 452.534 -685.783 m 459.707 -685.783 l S Q}%
\Vector{q 0 G 0 g 0.399 w 452.534 -693.454 m 459.707 -693.454 l S Q}%
\Vector{q 0 G 0 g 0.399 w 459.906 -693.654 m 459.906 -685.584 l S Q}%
\TextAt{8}{297.211}{43.033}{135}%
\Destination{78.23}{809.815}{nav.page.135}%
\Destination{79.36}{687.955}{nav.print.148}%
\Hotlink{504.94}{758.955}{25.078}{11.047}{nav.number.2.6.2}%
}

% Page 148 (body page 136)
\Page{%
\TextAt{2}{79.226}{796.943}{II.}%
\TextAt{2}{94.073}{796.943}{The}%
\TextAt{2}{116.492}{796.943}{Theory}%
\TextAt{2}{154.401}{796.943}{of}%
\TextAt{2}{166.826}{796.943}{Valuations}%
\Vector{q 0 G 0 g 0.299 w 7.226 18.729 m 460.769 18.729 l S Q}%
\TextAt{8}{97.159}{761.124}{We}%
\TextAt{8}{116.498}{761.124}{have}%
\TextAt{8}{142.688}{761.124}{introduced}%
\TextAt{8}{200.415}{761.124}{henselian}%
\TextAt{8}{250.395}{761.124}{fields}%
\TextAt{8}{279.561}{761.124}{by}%
\TextAt{8}{295.015}{761.124}{a}%
\TextAt{8}{303.965}{761.124}{condition}%
\TextAt{8}{354.543}{761.124}{of}%
\TextAt{8}{367.068}{761.124}{which}%
\TextAt{8}{399.749}{761.124}{the}%
\TextAt{8}{419.101}{761.124}{reader}%
\TextAt{8}{454.065}{761.124}{will}%
\TextAt{8}{475.365}{761.124}{find}%
\TextAt{8}{497.968}{761.124}{weaker}%
\TextAt{8}{79.226}{744.078}{versions}%
\TextAt{8}{123.448}{744.078}{in}%
\TextAt{8}{136.969}{744.078}{the}%
\TextAt{8}{156.994}{744.078}{literature,}%
\TextAt{8}{211.524}{744.078}{restricted}%
\TextAt{8}{263.481}{744.078}{to}%
\TextAt{19}{277.659}{744.078}{monic}%
\TextAt{8}{312.64}{744.078}{polynomials}%
\TextAt{8}{377.616}{744.078}{only.}%
\TextAt{8}{406.834}{744.078}{Both}%
\TextAt{8}{435.802}{744.078}{are}%
\TextAt{8}{455.181}{744.078}{equivalent}%
\TextAt{8}{510.665}{744.078}{as}%
\TextAt{8}{524.903}{744.078}{is}%
\TextAt{8}{79.226}{727.032}{shown}%
\TextAt{8}{114.733}{727.032}{by}%
\TextAt{8}{130.992}{727.032}{the}%
\TextAt{8}{151.148}{727.032}{following}%
\TextAt{12}{79.226}{701.019}{Proposition}%
\TextAt{12}{151.907}{701.019}{(6.7).}%
\TextAt{8}{187.196}{701.019}{A}%
\TextAt{8}{199.841}{701.019}{nonarchimedean}%
\TextAt{8}{286.645}{701.019}{field}%
\TextAt{9}{311.975}{701.019}{\Character{40}}%
\TextAt{13}{316.528}{701.019}{K}%
\TextAt{13}{326.69}{701.019}{\Character{59}}%
\TextAt{13}{331.926}{701.019}{v}%
\TextAt{9}{338.027}{701.019}{\Character{41}}%
\TextAt{8}{346.45}{701.019}{is}%
\TextAt{8}{358.187}{701.019}{henselian}%
\TextAt{8}{408.946}{701.019}{if}%
\TextAt{8}{419.643}{701.019}{any}%
\TextAt{8}{441.734}{701.019}{monic}%
\TextAt{8}{476.174}{701.019}{polynomial}%
\TextAt{13}{79.226}{683.973}{f}%
\TextAt{9}{86.273}{683.973}{\Character{40}}%
\TextAt{13}{90.825}{683.973}{x}%
\TextAt{9}{97.477}{683.973}{\Character{41}}%
\NativeMathAt{105.35}{683.973}{11.955}{\in}%
\TextAt{14}{116.641}{683.973}{o}%
\TextAt{9}{122.485}{683.973}{\Character{91}}%
\TextAt{13}{125.737}{683.973}{x}%
\TextAt{9}{132.389}{683.973}{\Character{93}}%
\TextAt{8}{139.538}{683.973}{which}%
\TextAt{8}{173.025}{683.973}{splits}%
\TextAt{8}{203.713}{683.973}{over}%
\TextAt{8}{228.76}{683.973}{the}%
\TextAt{8}{248.916}{683.973}{residue}%
\TextAt{8}{288.643}{683.973}{class}%
\TextAt{8}{316.08}{683.973}{field}%
\TextAt{13}{341.437}{683.973}{\Character{20}}%
\TextAt{9}{351.496}{683.973}{\Character{61}}%
\TextAt{14}{363.922}{683.973}{o}%
\TextAt{9}{369.766}{683.973}{\Character{47}}%
\TextAt{14}{375.619}{683.973}{p}%
\TextAt{8}{385.5}{683.973}{as}%
\TextAt{13}{247.112}{654.972}{f}%
\TextAt{9}{254.159}{654.972}{\Character{40}}%
\TextAt{13}{258.711}{654.972}{x}%
\TextAt{9}{265.363}{654.972}{\Character{41}}%
\NativeMathAt{273.236}{654.972}{11.955}{\equiv}%
\TextAt{9}{286.271}{654.972}{\Character{22}}%
\TextAt{13}{285.855}{654.972}{g}%
\TextAt{9}{291.89}{654.972}{\Character{40}}%
\TextAt{13}{296.442}{654.972}{x}%
\TextAt{9}{303.094}{654.972}{\Character{41}}%
\TextAt{9}{307.764}{658.127}{\Character{22}}%
\TextAt{13}{307.646}{654.972}{h}%
\TextAt{9}{314.385}{654.972}{\Character{40}}%
\TextAt{13}{318.937}{654.972}{x}%
\TextAt{9}{325.589}{654.972}{\Character{41}}%
\TextAt{8}{333.463}{654.972}{mod}%
\TextAt{14}{358.9}{654.972}{p}%
\TextAt{8}{79.226}{625.971}{with}%
\TextAt{8}{105.886}{625.971}{relatively}%
\TextAt{8}{156.947}{625.971}{prime}%
\TextAt{8}{190.11}{625.971}{monic}%
\TextAt{8}{224.577}{625.971}{factors}%
\TextAt{9}{263.106}{625.971}{\Character{22}}%
\TextAt{13}{262.69}{625.971}{g}%
\TextAt{9}{268.724}{625.971}{\Character{40}}%
\TextAt{13}{273.277}{625.971}{x}%
\TextAt{9}{279.929}{625.971}{\Character{41}}%
\TextAt{13}{284.481}{625.971}{\Character{59}}%
\TextAt{9}{289.843}{629.126}{\Character{22}}%
\TextAt{13}{289.725}{625.971}{h}%
\TextAt{9}{296.464}{625.971}{\Character{40}}%
\TextAt{13}{301.016}{625.971}{x}%
\TextAt{9}{307.668}{625.971}{\Character{41}}%
\NativeMathAt{315.541}{625.971}{11.955}{\in}%
\TextAt{13}{326.832}{625.971}{\Character{20}}%
\TextAt{9}{333.571}{625.971}{\Character{91}}%
\TextAt{13}{336.823}{625.971}{x}%
\TextAt{9}{343.475}{625.971}{\Character{93}}%
\TextAt{8}{346.726}{625.971}{,}%
\TextAt{8}{353.876}{625.971}{admits}%
\TextAt{8}{392.311}{625.971}{itself}%
\TextAt{8}{420.657}{625.971}{a}%
\TextAt{8}{430.413}{625.971}{splitting}%
\TextAt{13}{264.58}{596.97}{f}%
\TextAt{9}{271.626}{596.97}{\Character{40}}%
\TextAt{13}{276.179}{596.97}{x}%
\TextAt{9}{282.831}{596.97}{\Character{41}}%
\TextAt{9}{290.686}{596.97}{\Character{61}}%
\TextAt{13}{303.129}{596.97}{g}%
\TextAt{9}{309.164}{596.97}{\Character{40}}%
\TextAt{13}{313.716}{596.97}{x}%
\TextAt{9}{320.368}{596.97}{\Character{41}}%
\TextAt{13}{324.92}{596.97}{h}%
\TextAt{9}{331.659}{596.97}{\Character{40}}%
\TextAt{13}{336.211}{596.97}{x}%
\TextAt{9}{342.863}{596.97}{\Character{41}}%
\TextAt{8}{79.226}{567.969}{into}%
\TextAt{8}{102.969}{567.969}{monic}%
\TextAt{8}{137.436}{567.969}{factors}%
\TextAt{13}{175.549}{567.969}{g}%
\TextAt{9}{181.583}{567.969}{\Character{40}}%
\TextAt{13}{186.136}{567.969}{x}%
\TextAt{9}{192.788}{567.969}{\Character{41}}%
\TextAt{13}{197.34}{567.969}{\Character{59}}%
\TextAt{13}{202.564}{567.969}{h}%
\TextAt{9}{209.323}{567.969}{\Character{40}}%
\TextAt{13}{213.875}{567.969}{x}%
\TextAt{9}{220.527}{567.969}{\Character{41}}%
\NativeMathAt{228.4}{567.969}{11.955}{\in}%
\TextAt{14}{239.691}{567.969}{o}%
\TextAt{9}{245.535}{567.969}{\Character{91}}%
\TextAt{13}{248.787}{567.969}{x}%
\TextAt{9}{255.439}{567.969}{\Character{93}}%
\TextAt{8}{262.588}{567.969}{such}%
\TextAt{8}{288.985}{567.969}{that}%
\TextAt{13}{176.061}{538.968}{g}%
\TextAt{9}{182.095}{538.968}{\Character{40}}%
\TextAt{13}{186.647}{538.968}{x}%
\TextAt{9}{193.299}{538.968}{\Character{41}}%
\NativeMathAt{201.173}{538.968}{11.955}{\equiv}%
\TextAt{9}{214.208}{538.968}{\Character{22}}%
\TextAt{13}{213.792}{538.968}{g}%
\TextAt{9}{219.826}{538.968}{\Character{40}}%
\TextAt{13}{224.379}{538.968}{x}%
\TextAt{9}{231.031}{538.968}{\Character{41}}%
\TextAt{8}{238.904}{538.968}{mod}%
\TextAt{14}{264.342}{538.968}{p}%
\TextAt{8}{294.235}{538.968}{and}%
\TextAt{13}{337.011}{538.968}{h}%
\TextAt{9}{343.749}{538.968}{\Character{40}}%
\TextAt{13}{348.302}{538.968}{x}%
\TextAt{9}{354.954}{538.968}{\Character{41}}%
\NativeMathAt{362.827}{538.968}{11.955}{\equiv}%
\TextAt{9}{375.564}{542.123}{\Character{22}}%
\TextAt{13}{375.446}{538.968}{h}%
\TextAt{9}{382.185}{538.968}{\Character{40}}%
\TextAt{13}{386.737}{538.968}{x}%
\TextAt{9}{393.389}{538.968}{\Character{41}}%
\TextAt{8}{401.262}{538.968}{mod}%
\TextAt{14}{426.7}{538.968}{p}%
\TextAt{13}{432.683}{538.968}{\Character{58}}%
\TextAt{19}{79.226}{492.034}{Proof}%
\TextAt{19}{110.62}{492.034}{(}%
\TextAt{26}{116.706}{492.034}{E.}%
\TextAt{26}{133.909}{492.034}{Nart}%
\TextAt{19}{164.431}{492.034}{).}%
\TextAt{8}{178.777}{492.034}{We}%
\TextAt{8}{200.701}{492.034}{have}%
\TextAt{8}{229.475}{492.034}{just}%
\TextAt{8}{254.4}{492.034}{seen}%
\TextAt{8}{281.596}{492.034}{that}%
\TextAt{8}{308.744}{492.034}{the}%
\TextAt{8}{330.68}{492.034}{property}%
\TextAt{8}{380.268}{492.034}{of}%
\TextAt{13}{395.378}{492.034}{K}%
\TextAt{8}{411.872}{492.034}{to}%
\TextAt{8}{427.962}{492.034}{be}%
\TextAt{8}{445.666}{492.034}{henselian}%
\TextAt{8}{498.231}{492.034}{follows}%
\TextAt{8}{79.226}{474.988}{from}%
\TextAt{8}{108.72}{474.988}{the}%
\TextAt{8}{130.729}{474.988}{condition}%
\TextAt{8}{183.966}{474.988}{that}%
\TextAt{8}{211.188}{474.988}{the}%
\TextAt{8}{233.197}{474.988}{Newton}%
\TextAt{8}{278.292}{474.988}{polygon}%
\TextAt{8}{324.38}{474.988}{of}%
\TextAt{8}{339.563}{474.988}{every}%
\TextAt{8}{372.308}{474.988}{irreducible}%
\TextAt{8}{432.036}{474.988}{polynomial}%
\TextAt{13}{494.383}{474.988}{f}%
\TextAt{9}{501.429}{474.988}{\Character{40}}%
\TextAt{13}{505.982}{474.988}{x}%
\TextAt{9}{512.634}{474.988}{\Character{41}}%
\TextAt{9}{523.645}{474.988}{\Character{61}}%
\TextAt{13}{79.226}{457.942}{a}%
\TextAt{17}{85.371}{456.149}{0}%
\TextAt{9}{91.649}{457.942}{\Character{43}}%
\TextAt{13}{102.3}{457.942}{a}%
\TextAt{17}{108.445}{456.149}{1}%
\TextAt{13}{113.177}{457.942}{x}%
\TextAt{9}{121.376}{457.942}{\Character{43}}%
\TextAt{16}{132.026}{457.942}{\Character{1}}%
\TextAt{16}{137.322}{457.942}{\Character{1}}%
\TextAt{16}{142.63}{457.942}{\Character{1}}%
\TextAt{9}{147.52}{457.942}{\Character{43}}%
\TextAt{13}{158.171}{457.942}{a}%
\TextAt{21}{164.316}{456.149}{n}%
\TextAt{13}{169.952}{457.942}{x}%
\TextAt{21}{176.604}{462.28}{n}%
\NativeMathAt{185.561}{457.942}{11.955}{\in}%
\TextAt{13}{196.852}{457.942}{K}%
\TextAt{9}{207.67}{457.942}{\Character{91}}%
\TextAt{13}{210.922}{457.942}{x}%
\TextAt{9}{217.574}{457.942}{\Character{93}}%
\TextAt{8}{224.18}{457.942}{is}%
\TextAt{8}{235.4}{457.942}{a}%
\TextAt{8}{244.613}{457.942}{single}%
\TextAt{8}{276.648}{457.942}{line}%
\TextAt{8}{298.21}{457.942}{segment.}%
\TextAt{8}{347.842}{457.942}{It}%
\TextAt{8}{359.971}{457.942}{is}%
\TextAt{8}{371.192}{457.942}{therefore}%
\TextAt{8}{419.749}{457.942}{sufficient}%
\TextAt{8}{468.365}{457.942}{to}%
\TextAt{8}{482.132}{457.942}{show}%
\TextAt{8}{510.593}{457.942}{this.}%
\TextAt{8}{79.226}{440.896}{We}%
\TextAt{8}{99.537}{440.896}{may}%
\TextAt{8}{125.073}{440.896}{assume}%
\TextAt{8}{165.684}{440.896}{that}%
\TextAt{13}{191.219}{440.896}{a}%
\TextAt{21}{197.364}{439.103}{n}%
\TextAt{9}{206.605}{440.896}{\Character{61}}%
\TextAt{9}{219.301}{440.896}{1}%
\TextAt{8}{225.168}{440.896}{.}%
\TextAt{8}{234.116}{440.896}{Let}%
\TextAt{13}{255.252}{440.896}{L}%
\TextAt{16}{266.821}{440.896}{\Character{106}}%
\TextAt{13}{273.747}{440.896}{K}%
\TextAt{8}{288.628}{440.896}{be}%
\TextAt{8}{304.719}{440.896}{the}%
\TextAt{8}{325.042}{440.896}{splitting}%
\TextAt{8}{371.452}{440.896}{field}%
\TextAt{8}{396.975}{440.896}{of}%
\TextAt{13}{410.472}{440.896}{f}%
\TextAt{8}{417.518}{440.896}{.}%
\TextAt{8}{426.466}{440.896}{Then}%
\TextAt{8}{457.19}{440.896}{there}%
\TextAt{8}{487.269}{440.896}{is}%
\TextAt{8}{499.199}{440.896}{always}%
\TextAt{8}{79.226}{423.85}{an}%
\TextAt{8}{95.891}{423.85}{extension}%
\TextAt{13}{148.063}{423.85}{w}%
\TextAt{8}{161.056}{423.85}{of}%
\TextAt{13}{174.792}{423.85}{v}%
\TextAt{8}{185.184}{423.85}{to}%
\TextAt{13}{199.9}{423.85}{L}%
\TextAt{8}{207.865}{423.85}{.}%
\TextAt{8}{217.531}{423.85}{It}%
\TextAt{8}{230.61}{423.85}{is}%
\TextAt{8}{242.78}{423.85}{obtained}%
\TextAt{8}{291.318}{423.85}{for}%
\TextAt{8}{309.609}{423.85}{example}%
\TextAt{8}{355.863}{423.85}{by}%
\TextAt{8}{372.528}{423.85}{taking}%
\TextAt{8}{409.039}{423.85}{the}%
\TextAt{8}{429.601}{423.85}{completion}%
\TextAt{20}{492.581}{426.872}{\Character{98}}%
\TextAt{13}{489.843}{423.85}{K}%
\TextAt{8}{504.964}{423.85}{of}%
\TextAt{13}{518.7}{423.85}{K}%
\TextAt{8}{529.518}{423.85}{,}%
\TextAt{8}{79.226}{406.804}{extending}%
\TextAt{8}{132.733}{406.804}{the}%
\TextAt{8}{152.741}{406.804}{valuation}%
\TextAt{8}{203.665}{406.804}{of}%
\TextAt{20}{219.585}{409.826}{\Character{98}}%
\TextAt{13}{216.846}{406.804}{K}%
\TextAt{8}{231.413}{406.804}{in}%
\TextAt{8}{244.918}{406.804}{a}%
\TextAt{8}{254.525}{406.804}{unique}%
\TextAt{8}{292.418}{406.804}{way}%
\TextAt{8}{316.012}{406.804}{to}%
\TextAt{8}{330.174}{406.804}{a}%
\TextAt{8}{339.781}{406.804}{valuation}%
\TextAt{9}{391.148}{406.804}{\Character{22}}%
\TextAt{13}{390.705}{406.804}{v}%
\TextAt{8}{400.543}{406.804}{of}%
\TextAt{8}{413.724}{406.804}{the}%
\TextAt{8}{433.732}{406.804}{algebraic}%
\TextAt{8}{483.019}{406.804}{closure}%
\Vector{q 0 G 0 g 0.478 w 449.952 -353.866 m 460.769 -353.866 l S Q}%
\TextAt{20}{524.69}{409.826}{\Character{98}}%
\TextAt{13}{521.952}{406.804}{K}%
\TextAt{8}{79.226}{389.758}{of}%
\TextAt{20}{95.073}{392.78}{\Character{98}}%
\TextAt{13}{92.334}{389.758}{K}%
\TextAt{8}{103.152}{389.758}{,}%
\TextAt{8}{110.124}{389.758}{embedding}%
\TextAt{13}{169.08}{389.758}{L}%
\TextAt{8}{180.721}{389.758}{into}%
\Vector{q 0 G 0 g 0.478 w 132.242 -370.912 m 143.059 -370.912 l S Q}%
\TextAt{20}{206.98}{392.78}{\Character{98}}%
\TextAt{13}{204.242}{389.758}{K}%
\TextAt{8}{215.059}{389.758}{,}%
\TextAt{8}{222.031}{389.758}{and}%
\TextAt{8}{244.572}{389.758}{restricting}%
\TextAt{9}{300.791}{389.758}{\Character{22}}%
\TextAt{13}{300.349}{389.758}{v}%
\TextAt{8}{310.112}{389.758}{to}%
\TextAt{13}{324.201}{389.758}{L}%
\TextAt{8}{332.166}{389.758}{.}%
\TextAt{8}{340.54}{389.758}{It}%
\TextAt{8}{352.991}{389.758}{is}%
\TextAt{8}{364.533}{389.758}{also}%
\TextAt{8}{387.791}{389.758}{possible}%
\TextAt{8}{431.588}{389.758}{to}%
\TextAt{8}{445.677}{389.758}{get}%
\TextAt{8}{464.966}{389.758}{the}%
\TextAt{8}{484.901}{389.758}{extension}%
\TextAt{13}{79.226}{372.712}{w}%
\TextAt{8}{92.011}{372.712}{directly,}%
\TextAt{8}{137.127}{372.712}{without}%
\TextAt{8}{180.901}{372.712}{passing}%
\TextAt{8}{222.201}{372.712}{through}%
\TextAt{8}{266.634}{372.712}{the}%
\TextAt{8}{286.988}{372.712}{completion.}%
\TextAt{8}{351.969}{372.712}{For}%
\TextAt{8}{373.137}{372.712}{this}%
\TextAt{8}{396.158}{372.712}{we}%
\TextAt{8}{413.583}{372.712}{refer}%
\TextAt{8}{440.764}{372.712}{to}%
\TextAt{8}{455.273}{372.712}{[93],}%
\TextAt{8}{480.889}{372.712}{chap.}%
\TextAt{8}{511.979}{372.712}{XII,}%
\TextAt{8}{79.226}{355.666}{§4,}%
\TextAt{8}{97.912}{355.666}{th.}%
\TextAt{8}{116.12}{355.666}{1.}%
\TextAt{8}{97.159}{335.631}{Assume}%
\TextAt{8}{140.52}{335.631}{now}%
\TextAt{8}{164.909}{335.631}{that}%
\TextAt{8}{190.278}{335.631}{the}%
\TextAt{8}{210.434}{335.631}{Newton}%
\TextAt{8}{253.676}{335.631}{polygon}%
\TextAt{8}{297.91}{335.631}{of}%
\TextAt{13}{311.24}{335.631}{f}%
\TextAt{8}{322.184}{335.631}{consists}%
\TextAt{8}{365.294}{335.631}{of}%
\TextAt{8}{378.624}{335.631}{more}%
\TextAt{8}{407.891}{335.631}{than}%
\TextAt{8}{435.208}{335.631}{one}%
\TextAt{8}{456.668}{335.631}{segment:}%
\NativeDiagram{%
\path[draw=black,line width=0.399bp] (139.605,-494.352) -- (325.329,-494.352);%
\path[draw=black,line width=0.399bp,fill=black] (327.826,-494.352) -- (324.255,-493.006) -- (325.429,-494.352) -- (324.255,-495.698) -- cycle;%
\path[draw=black,line width=0.399bp] (154.912,-525.676) -- (154.912,-454.045);%
\path[draw=black,line width=0.399bp,fill=black] (154.912,-451.548) -- (153.566,-455.12) -- (154.912,-453.945) -- (156.258,-455.12) -- cycle;%
\path[draw=black,line width=0.797bp] (182.975,-504.713) -- (233.999,-521.338) -- (273.542,-518.929) -- (299.054,-494.352);%
\path[draw=black,line width=0.399bp,dash pattern=on 2.989bp off 2.989bp] (273.542,-494.352) -- (273.542,-518.929);%
\path[fill=black] (170.219,-481.1) (171.813,-481.1) .. controls (171.813,-480.219) and (171.1,-479.506) .. (170.219,-479.506) .. controls (169.339,-479.506) and (168.625,-480.219) .. (168.625,-481.1) .. controls (168.625,-481.98) and (169.339,-482.694) .. (170.219,-482.694) .. controls (171.1,-482.694) and (171.813,-481.98) .. (171.813,-481.1) -- cycle (170.219,-481.1);%
\path[fill=black] (174.301,-488.328) (175.895,-488.328) .. controls (175.895,-487.448) and (175.181,-486.734) .. (174.301,-486.734) .. controls (173.421,-486.734) and (172.707,-487.448) .. (172.707,-488.328) .. controls (172.707,-489.209) and (173.421,-489.922) .. (174.301,-489.922) .. controls (175.181,-489.922) and (175.895,-489.209) .. (175.895,-488.328) -- cycle (174.301,-488.328);%
\path[fill=black] (178.638,-496.28) (180.232,-496.28) .. controls (180.232,-495.399) and (179.518,-494.686) .. (178.638,-494.686) .. controls (177.758,-494.686) and (177.044,-495.399) .. (177.044,-496.28) .. controls (177.044,-497.16) and (177.758,-497.874) .. (178.638,-497.874) .. controls (179.518,-497.874) and (180.232,-497.16) .. (180.232,-496.28) -- cycle (178.638,-496.28);%
\path[fill=black] (182.975,-504.713) (184.569,-504.713) .. controls (184.569,-503.832) and (183.856,-503.119) .. (182.975,-503.119) .. controls (182.095,-503.119) and (181.381,-503.832) .. (181.381,-504.713) .. controls (181.381,-505.593) and (182.095,-506.307) .. (182.975,-506.307) .. controls (183.856,-506.307) and (184.569,-505.593) .. (184.569,-504.713) -- cycle (182.975,-504.713);%
\path[fill=black] (233.999,-521.338) (235.593,-521.338) .. controls (235.593,-520.458) and (234.88,-519.744) .. (233.999,-519.744) .. controls (233.119,-519.744) and (232.405,-520.458) .. (232.405,-521.338) .. controls (232.405,-522.219) and (233.119,-522.932) .. (233.999,-522.932) .. controls (234.88,-522.932) and (235.593,-522.219) .. (235.593,-521.338) -- cycle (233.999,-521.338);%
\path[fill=black] (273.542,-518.929) (275.136,-518.929) .. controls (275.136,-518.049) and (274.423,-517.335) .. (273.542,-517.335) .. controls (272.662,-517.335) and (271.948,-518.049) .. (271.948,-518.929) .. controls (271.948,-519.809) and (272.662,-520.523) .. (273.542,-520.523) .. controls (274.423,-520.523) and (275.136,-519.809) .. (275.136,-518.929) -- cycle (273.542,-518.929);%
}%
\TextAt{13}{340.423}{283.832}{m}%
\TextAt{13}{335.933}{264.545}{e}%
\TextAt{8}{79.226}{222.113}{Let}%
\TextAt{8}{100.617}{222.113}{the}%
\TextAt{8}{121.196}{222.113}{last}%
\TextAt{8}{143.794}{222.113}{segment}%
\TextAt{8}{189.479}{222.113}{be}%
\TextAt{8}{205.825}{222.113}{given}%
\TextAt{8}{236.816}{222.113}{by}%
\TextAt{8}{253.497}{222.113}{the}%
\TextAt{8}{274.075}{222.113}{points}%
\TextAt{9}{309.681}{222.113}{\Character{40}}%
\TextAt{13}{314.234}{222.113}{m\Character{59}}%
\TextAt{13}{329.692}{222.113}{e}%
\TextAt{9}{335.143}{222.113}{\Character{41}}%
\TextAt{8}{344.014}{222.113}{and}%
\TextAt{9}{367.199}{222.113}{\Character{40}}%
\TextAt{13}{371.751}{222.113}{n\Character{59}}%
\TextAt{9}{383.983}{222.113}{0\Character{41}}%
\TextAt{8}{394.388}{222.113}{.}%
\TextAt{8}{404.102}{222.113}{If}%
\TextAt{13}{416.216}{222.113}{e}%
\TextAt{9}{425.681}{222.113}{\Character{61}}%
\TextAt{9}{438.808}{222.113}{0}%
\TextAt{8}{444.679}{222.113}{,}%
\TextAt{8}{452.356}{222.113}{we}%
\TextAt{8}{470.005}{222.113}{immediately}%
\TextAt{8}{79.226}{205.067}{have}%
\TextAt{8}{107.183}{205.067}{a}%
\TextAt{8}{117.9}{205.067}{contradiction.}%
\TextAt{8}{196.567}{205.067}{Because}%
\TextAt{8}{242.289}{205.067}{then}%
\TextAt{8}{269.911}{205.067}{we}%
\TextAt{8}{288.1}{205.067}{have}%
\TextAt{13}{316.057}{205.067}{v}%
\TextAt{9}{322.145}{205.067}{\Character{40}}%
\TextAt{13}{326.698}{205.067}{a}%
\TextAt{21}{332.843}{203.274}{i}%
\TextAt{9}{336.224}{205.067}{\Character{41}}%
\NativeMathAt{345.737}{205.067}{11.955}{\geq}%
\TextAt{9}{359.995}{205.067}{0}%
\TextAt{8}{365.848}{205.067}{,}%
\TextAt{8}{374.2}{205.067}{so}%
\TextAt{8}{389.532}{205.067}{that}%
\TextAt{13}{415.863}{205.067}{f}%
\TextAt{9}{422.91}{205.067}{\Character{40}}%
\TextAt{13}{427.462}{205.067}{x}%
\TextAt{9}{434.114}{205.067}{\Character{41}}%
\NativeMathAt{443.627}{205.067}{11.955}{\in}%
\TextAt{14}{456.557}{205.067}{o}%
\TextAt{9}{462.402}{205.067}{\Character{91}}%
\TextAt{13}{465.653}{205.067}{x}%
\TextAt{9}{472.305}{205.067}{\Character{93}}%
\TextAt{8}{475.557}{205.067}{,}%
\TextAt{8}{483.909}{205.067}{and}%
\TextAt{13}{507.633}{205.067}{a}%
\TextAt{17}{513.778}{203.274}{0}%
\NativeMathAt{523.471}{205.067}{11.955}{\equiv}%
\TextAt{16}{79.226}{188.021}{\Character{1}}%
\TextAt{16}{84.522}{188.021}{\Character{1}}%
\TextAt{16}{89.83}{188.021}{\Character{1}}%
\NativeMathAt{97.338}{188.021}{11.955}{\equiv}%
\TextAt{13}{110.87}{188.022}{a}%
\TextAt{21}{117.015}{186.228}{m}%
\TextAt{3}{124.505}{186.228}{\Character{0}}%
\TextAt{17}{131.092}{186.228}{1}%
\NativeMathAt{140.023}{188.021}{11.955}{\equiv}%
\TextAt{9}{153.52}{188.021}{0}%
\TextAt{8}{162.694}{188.021}{mod}%
\TextAt{14}{188.132}{188.021}{p}%
\TextAt{8}{194.115}{188.021}{,}%
\TextAt{13}{201.908}{188.021}{a}%
\TextAt{21}{208.053}{186.228}{m}%
\NativeMathAt{220.24}{188.021}{11.955}{\not\equiv}%
\TextAt{9}{233.737}{188.021}{0}%
\TextAt{8}{242.911}{188.021}{mod}%
\TextAt{14}{268.349}{188.021}{p}%
\TextAt{8}{274.332}{188.021}{.}%
\TextAt{8}{284.326}{188.021}{Therefore}%
\TextAt{13}{337.838}{188.021}{f}%
\TextAt{9}{344.885}{188.021}{\Character{40}}%
\TextAt{13}{349.437}{188.021}{x}%
\TextAt{9}{356.089}{188.021}{\Character{41}}%
\NativeMathAt{364.84}{188.021}{11.955}{\equiv}%
\TextAt{9}{378.337}{188.021}{\Character{40}}%
\TextAt{13}{382.89}{188.021}{X}%
\TextAt{21}{393.545}{192.36}{n}%
\TextAt{3}{398.683}{192.36}{\Character{0}}%
\TextAt{21}{405.269}{192.36}{m}%
\TextAt{9}{416.266}{188.021}{\Character{43}}%
\TextAt{16}{428.378}{188.021}{\Character{1}}%
\TextAt{16}{433.674}{188.021}{\Character{1}}%
\TextAt{16}{438.982}{188.021}{\Character{1}}%
\TextAt{9}{445.334}{188.021}{\Character{43}}%
\TextAt{13}{457.446}{188.021}{a}%
\TextAt{21}{463.591}{186.228}{m}%
\TextAt{9}{471.58}{188.021}{\Character{41}}%
\TextAt{13}{476.132}{188.021}{X}%
\TextAt{21}{486.787}{192.36}{m}%
\TextAt{8}{498.097}{188.021}{mod}%
\TextAt{14}{523.535}{188.021}{p}%
\TextAt{8}{529.518}{188.021}{,}%
\TextAt{8}{79.226}{170.976}{with}%
\TextAt{13}{106.976}{170.976}{m}%
\TextAt{13}{122.386}{170.976}{\Character{62}}%
\TextAt{9}{136.678}{170.976}{0}%
\TextAt{8}{147.519}{170.976}{because}%
\TextAt{8}{191.911}{170.976}{there}%
\TextAt{8}{222.913}{170.976}{is}%
\TextAt{8}{235.767}{170.976}{more}%
\TextAt{8}{266.123}{170.976}{than}%
\TextAt{8}{294.531}{170.976}{one}%
\TextAt{8}{317.081}{170.976}{segment.}%
\TextAt{8}{370.165}{170.976}{In}%
\TextAt{8}{385.876}{170.976}{view}%
\TextAt{8}{413.95}{170.976}{of}%
\TextAt{8}{428.37}{170.976}{the}%
\TextAt{8}{449.617}{170.976}{condition}%
\TextAt{8}{502.09}{170.976}{of}%
\TextAt{8}{516.51}{170.976}{the}%
\TextAt{8}{79.226}{153.93}{proposition}%
\TextAt{8}{140.759}{153.93}{this}%
\TextAt{8}{163.582}{153.93}{contradicts}%
\TextAt{8}{223.812}{153.93}{the}%
\TextAt{8}{243.968}{153.93}{irreducibility}%
\TextAt{8}{313.559}{153.93}{of}%
\TextAt{13}{326.889}{153.93}{f}%
\TextAt{8}{333.936}{153.93}{.}%
\TextAt{8}{97.159}{133.895}{We}%
\TextAt{8}{118.158}{133.895}{will}%
\TextAt{8}{141.118}{133.895}{now}%
\TextAt{8}{166.361}{133.895}{reduce}%
\TextAt{8}{204.276}{133.895}{to}%
\TextAt{13}{219.441}{133.895}{e}%
\TextAt{9}{229.644}{133.895}{\Character{61}}%
\TextAt{9}{243.512}{133.895}{0}%
\TextAt{8}{254.13}{133.895}{by}%
\TextAt{8}{271.244}{133.895}{a}%
\TextAt{8}{281.854}{133.895}{transformation.}%
\TextAt{8}{368.721}{133.895}{Let}%
\TextAt{13}{390.545}{133.895}{\Character{11}}%
\NativeMathAt{402.844}{133.895}{11.955}{\in}%
\TextAt{13}{415.591}{133.895}{L}%
\TextAt{8}{428.308}{133.895}{be}%
\TextAt{8}{445.087}{133.895}{a}%
\TextAt{8}{455.697}{133.895}{root}%
\TextAt{8}{481.597}{133.895}{of}%
\TextAt{13}{495.782}{133.895}{f}%
\TextAt{9}{502.828}{133.895}{\Character{40}}%
\TextAt{13}{507.381}{133.895}{x}%
\TextAt{9}{514.033}{133.895}{\Character{41}}%
\TextAt{8}{523.337}{133.895}{of}%
\TextAt{8}{79.226}{116.849}{minimum}%
\TextAt{8}{132.056}{116.849}{value}%
\TextAt{13}{162.78}{116.849}{w}%
\TextAt{9}{171.47}{116.849}{\Character{40}}%
\TextAt{13}{176.022}{116.849}{\Character{11}}%
\TextAt{9}{183.544}{116.849}{\Character{41}}%
\TextAt{8}{192.471}{116.849}{and}%
\TextAt{8}{215.712}{116.849}{let}%
\TextAt{13}{233.094}{116.849}{a}%
\NativeMathAt{243.374}{116.849}{11.955}{\in}%
\TextAt{13}{255.479}{116.849}{K}%
\TextAt{8}{270.672}{116.849}{such}%
\TextAt{8}{297.547}{116.849}{that}%
\TextAt{13}{323.394}{116.849}{v}%
\TextAt{9}{329.482}{116.849}{\Character{40}}%
\TextAt{13}{334.034}{116.849}{a}%
\TextAt{9}{340.179}{116.849}{\Character{41}}%
\TextAt{9}{348.846}{116.849}{\Character{61}}%
\TextAt{13}{362.106}{116.849}{e}%
\TextAt{8}{367.532}{116.849}{.}%
\TextAt{8}{377.413}{116.849}{We}%
\TextAt{8}{398.036}{116.849}{consider}%
\TextAt{8}{444.098}{116.849}{the}%
\TextAt{8}{464.733}{116.849}{characteristic}%
\TextAt{8}{79.226}{99.803}{polynomial}%
\TextAt{13}{140.86}{99.803}{g}%
\TextAt{9}{146.894}{99.803}{\Character{40}}%
\TextAt{13}{151.447}{99.803}{x}%
\TextAt{9}{158.099}{99.803}{\Character{41}}%
\TextAt{8}{167.689}{99.803}{of}%
\TextAt{13}{182.16}{99.803}{a}%
\TextAt{3}{188.305}{104.141}{\Character{0}}%
\TextAt{17}{194.891}{104.141}{1}%
\TextAt{13}{199.624}{99.803}{\Character{11}}%
\TextAt{21}{207.146}{104.141}{r}%
\NativeMathAt{216.965}{99.803}{11.955}{\in}%
\TextAt{13}{230.2}{99.803}{K}%
\TextAt{9}{241.018}{99.803}{\Character{40}}%
\TextAt{13}{245.57}{99.803}{\Character{11}}%
\TextAt{9}{253.092}{99.803}{\Character{41}}%
\TextAt{8}{257.644}{99.803}{,}%
\TextAt{13}{266.219}{99.803}{r}%
\TextAt{9}{277.085}{99.803}{\Character{61}}%
\TextAt{13}{291.454}{99.803}{n}%
\TextAt{16}{301.876}{99.803}{\Character{0}}%
\TextAt{13}{314.609}{99.803}{m}%
\TextAt{8}{324.848}{99.803}{.}%
\TextAt{8}{336.719}{99.803}{If}%
\TextAt{13}{349.552}{99.803}{f}%
\TextAt{9}{356.598}{99.803}{\Character{40}}%
\TextAt{13}{361.151}{99.803}{x}%
\TextAt{9}{367.803}{99.803}{\Character{41}}%
\TextAt{9}{377.606}{99.803}{\Character{61}}%
\TextAt{20}{391.989}{108.769}{\Character{81}}%
\TextAt{21}{403.28}{105.614}{n}%
\TextAt{21}{403.28}{96.316}{i}%
\TextAt{17}{406.164}{96.316}{\Character{61}1}%
\TextAt{9}{417.482}{99.803}{\Character{40}}%
\TextAt{13}{422.035}{99.803}{x}%
\TextAt{16}{432.121}{99.803}{\Character{0}}%
\TextAt{13}{444.854}{99.803}{\Character{11}}%
\TextAt{21}{452.299}{98.01}{i}%
\TextAt{9}{455.681}{99.803}{\Character{41}}%
\TextAt{8}{460.233}{99.803}{,}%
\TextAt{8}{468.808}{99.803}{then}%
\TextAt{13}{496.609}{99.803}{g}%
\TextAt{9}{502.643}{99.803}{\Character{40}}%
\TextAt{13}{507.196}{99.803}{x}%
\TextAt{9}{513.848}{99.803}{\Character{41}}%
\TextAt{9}{523.651}{99.803}{\Character{61}}%
\TextAt{8}{297.211}{43.033}{136}%
\Destination{78.23}{809.815}{nav.page.136}%
\Destination{79.36}{507.055}{nav.print.149}%
\Destination{79.36}{715.055}{nav.number.2.6.7}%
\Hotlink{454.773}{370.555}{19.22}{11.047}{nav.cite.N093}%
%
}

% Page 149 (body page 137)
\Page{%
\TextAt{2}{79.226}{796.943}{II.}%
\TextAt{2}{94.073}{796.943}{The}%
\TextAt{2}{116.492}{796.943}{Theory}%
\TextAt{2}{154.401}{796.943}{of}%
\TextAt{2}{166.826}{796.943}{Valuations}%
\Vector{q 0 G 0 g 0.299 w 7.226 18.729 m 460.769 18.729 l S Q}%
\TextAt{20}{79.226}{770.09}{\Character{81}}%
\TextAt{21}{90.517}{766.935}{n}%
\TextAt{21}{90.517}{757.637}{i}%
\TextAt{17}{93.4}{757.637}{\Character{61}1}%
\TextAt{9}{104.719}{761.124}{\Character{40}}%
\TextAt{13}{109.271}{761.124}{x}%
\TextAt{16}{117.701}{761.124}{\Character{0}}%
\TextAt{13}{128.776}{761.124}{\Character{11}}%
\TextAt{21}{136.298}{765.462}{r}%
\TextAt{21}{136.222}{758.168}{i}%
\TextAt{13}{140.853}{761.124}{a}%
\TextAt{3}{146.998}{765.462}{\Character{0}}%
\TextAt{17}{153.584}{765.462}{1}%
\TextAt{9}{158.316}{761.124}{\Character{41}}%
\TextAt{8}{162.869}{761.124}{.}%
\TextAt{8}{171.174}{761.124}{Proposition}%
\TextAt{8}{233.735}{761.124}{(6.3)}%
\TextAt{8}{261.28}{761.124}{shows}%
\TextAt{8}{294.468}{761.124}{that}%
\TextAt{8}{319.407}{761.124}{the}%
\TextAt{8}{339.133}{761.124}{Newton}%
\TextAt{8}{381.945}{761.124}{polygon}%
\TextAt{8}{425.749}{761.124}{of}%
\TextAt{13}{438.649}{761.124}{g}%
\TextAt{9}{444.683}{761.124}{\Character{40}}%
\TextAt{13}{449.235}{761.124}{x}%
\TextAt{9}{455.888}{761.124}{\Character{41}}%
\TextAt{8}{463.907}{761.124}{also}%
\TextAt{8}{486.957}{761.124}{has}%
\TextAt{8}{507.401}{761.124}{more}%
\TextAt{8}{79.226}{744.078}{than}%
\TextAt{8}{106.544}{744.078}{one}%
\TextAt{8}{128.003}{744.078}{segment,}%
\TextAt{8}{176.517}{744.078}{the}%
\TextAt{8}{196.674}{744.078}{last}%
\TextAt{8}{218.851}{744.078}{one}%
\TextAt{8}{240.31}{744.078}{of}%
\TextAt{8}{253.64}{744.078}{slope}%
\TextAt{16}{199.262}{715.077}{\Character{0}}%
\TextAt{13}{208.56}{715.077}{w}%
\TextAt{9}{217.25}{715.077}{\Character{40}}%
\TextAt{13}{221.802}{715.077}{a}%
\TextAt{3}{227.947}{720.013}{\Character{0}}%
\TextAt{17}{234.534}{720.013}{1}%
\TextAt{13}{239.266}{715.077}{\Character{11}}%
\TextAt{21}{246.788}{720.013}{r}%
\TextAt{9}{251.343}{715.077}{\Character{41}}%
\TextAt{9}{259.198}{715.077}{\Character{61}}%
\TextAt{13}{271.641}{715.077}{v}%
\TextAt{9}{277.729}{715.077}{\Character{40}}%
\TextAt{13}{282.282}{715.077}{a}%
\TextAt{9}{288.427}{715.077}{\Character{41}}%
\TextAt{16}{295.636}{715.077}{\Character{0}}%
\TextAt{13}{307.591}{715.077}{r}%
\TextAt{13}{313.186}{715.077}{w}%
\TextAt{9}{321.881}{715.077}{\Character{40}}%
\TextAt{13}{326.433}{715.077}{\Character{11}}%
\TextAt{9}{333.955}{715.077}{\Character{41}}%
\TextAt{9}{341.81}{715.077}{\Character{61}}%
\TextAt{13}{354.254}{715.077}{e}%
\TextAt{16}{362.336}{715.077}{\Character{0}}%
\TextAt{13}{374.291}{715.077}{r}%
\TextAt{13}{381.174}{723.164}{e}%
\Vector{q 0 G 0 g 0.478 w 309.087 -55.935 m 314.687 -55.935 l S Q}%
\TextAt{13}{381.087}{706.876}{r}%
\TextAt{9}{391.204}{715.077}{\Character{61}}%
\TextAt{9}{403.613}{715.077}{0}%
\TextAt{13}{409.482}{715.077}{\Character{58}}%
\TextAt{8}{79.226}{684.45}{Since}%
\TextAt{13}{110.185}{684.45}{g}%
\TextAt{9}{116.219}{684.45}{\Character{40}}%
\TextAt{13}{120.772}{684.45}{x}%
\TextAt{9}{127.424}{684.45}{\Character{41}}%
\TextAt{8}{136.275}{684.45}{is}%
\TextAt{8}{148.441}{684.45}{a}%
\TextAt{8}{158.598}{684.45}{power}%
\TextAt{8}{193.144}{684.45}{of}%
\TextAt{8}{206.875}{684.45}{the}%
\TextAt{8}{227.433}{684.45}{minimal}%
\TextAt{8}{273.36}{684.45}{polynomial}%
\TextAt{8}{334.255}{684.45}{of}%
\TextAt{13}{347.987}{684.45}{a}%
\TextAt{3}{354.132}{688.788}{\Character{0}}%
\TextAt{17}{360.718}{688.788}{1}%
\TextAt{13}{365.45}{684.45}{\Character{11}}%
\TextAt{21}{372.972}{688.788}{r}%
\TextAt{8}{377.527}{684.45}{,}%
\TextAt{8}{385.178}{684.45}{hence}%
\TextAt{8}{418.086}{684.45}{of}%
\TextAt{8}{431.818}{684.45}{an}%
\TextAt{8}{448.478}{684.45}{irreducible}%
\TextAt{8}{506.755}{684.45}{poly-}%
\TextAt{8}{79.226}{667.404}{nomial,}%
\TextAt{8}{120.854}{667.404}{this}%
\TextAt{8}{143.676}{667.404}{produces}%
\TextAt{8}{192.836}{667.404}{the}%
\TextAt{8}{212.993}{667.404}{same}%
\TextAt{8}{242.319}{667.404}{contradiction}%
\TextAt{8}{313.547}{667.404}{as}%
\TextAt{8}{327.918}{667.404}{before.}%
\Vector{q 0 G 0 g 0.399 w 452.334 -106.596 m 452.334 -98.526 l S Q}%
\Vector{q 0 G 0 g 0.399 w 452.534 -98.726 m 459.707 -98.726 l S Q}%
\Vector{q 0 G 0 g 0.399 w 452.534 -106.397 m 459.707 -106.397 l S Q}%
\Vector{q 0 G 0 g 0.399 w 459.906 -106.596 m 459.906 -98.526 l S Q}%
\TextAt{8}{97.159}{638.403}{Let}%
\TextAt{13}{119.041}{638.403}{K}%
\TextAt{8}{134.668}{638.403}{be}%
\TextAt{8}{151.504}{638.403}{a}%
\TextAt{8}{162.172}{638.403}{field}%
\TextAt{8}{188.441}{638.403}{which}%
\TextAt{8}{222.84}{638.403}{is}%
\TextAt{8}{235.516}{638.403}{henselian}%
\TextAt{8}{287.214}{638.403}{with}%
\TextAt{8}{314.786}{638.403}{respect}%
\TextAt{8}{355.748}{638.403}{to}%
\TextAt{8}{370.971}{638.403}{the}%
\TextAt{8}{392.04}{638.403}{exponential}%
\TextAt{8}{455.717}{638.403}{valuation}%
\TextAt{13}{507.701}{638.403}{v}%
\TextAt{8}{513.79}{638.403}{.}%
\TextAt{8}{524.975}{638.403}{If}%
\TextAt{13}{79.226}{621.357}{L}%
\TextAt{16}{90.511}{621.357}{\Character{106}}%
\TextAt{13}{97.153}{621.357}{K}%
\TextAt{8}{110.959}{621.357}{is}%
\TextAt{8}{121.813}{621.357}{a}%
\TextAt{8}{130.659}{621.357}{finite}%
\TextAt{8}{159.661}{621.357}{extension}%
\TextAt{8}{210.517}{621.357}{of}%
\TextAt{8}{222.937}{621.357}{degree}%
\TextAt{13}{258.443}{621.357}{n}%
\TextAt{8}{265.431}{621.357}{,}%
\TextAt{8}{271.852}{621.357}{then}%
\TextAt{13}{297.603}{621.357}{v}%
\TextAt{8}{306.679}{621.357}{extends}%
\TextAt{8}{348.425}{621.357}{uniquely}%
\TextAt{8}{394.989}{621.357}{to}%
\TextAt{8}{408.39}{621.357}{an}%
\TextAt{8}{423.739}{621.357}{exponential}%
\TextAt{8}{485.594}{621.357}{valuation}%
\TextAt{13}{79.226}{604.311}{w}%
\TextAt{8}{91.813}{604.311}{of}%
\TextAt{13}{105.143}{604.311}{L}%
\TextAt{8}{113.108}{604.311}{,}%
\TextAt{8}{120.257}{604.311}{namely}%
\TextAt{13}{250.398}{585.083}{w}%
\TextAt{9}{259.088}{585.083}{\Character{40}}%
\TextAt{13}{263.64}{585.083}{\Character{11}}%
\TextAt{9}{271.162}{585.083}{\Character{41}}%
\TextAt{9}{279.017}{585.083}{\Character{61}}%
\TextAt{9}{293.223}{593.171}{1}%
\Vector{q 0 G 0 g 0.478 w 220.656 -185.928 m 227.643 -185.928 l S Q}%
\TextAt{13}{292.656}{576.883}{n}%
\TextAt{13}{300.839}{585.083}{v}%
\TextAt{9}{306.927}{585.083}{\Character{40}}%
\TextAt{13}{311.479}{585.083}{N}%
\TextAt{21}{320.853}{583.224}{L}%
\TextAt{3}{326.612}{583.224}{\Character{106}}%
\TextAt{21}{328.965}{583.224}{K}%
\TextAt{9}{337.167}{585.083}{\Character{40}}%
\TextAt{13}{341.72}{585.083}{\Character{11}}%
\TextAt{9}{349.241}{585.083}{\Character{41}\Character{41}}%
\TextAt{13}{358.346}{585.083}{\Character{58}}%
\TextAt{8}{79.226}{560.444}{This}%
\TextAt{8}{106.394}{560.444}{follows}%
\TextAt{8}{145.277}{560.444}{from}%
\TextAt{8}{173.366}{560.444}{(6.2)}%
\TextAt{8}{201.789}{560.444}{by}%
\TextAt{8}{218.495}{560.444}{taking}%
\TextAt{8}{255.048}{560.444}{the}%
\TextAt{8}{275.652}{560.444}{logarithm.}%
\TextAt{8}{334.89}{560.444}{For}%
\TextAt{8}{356.308}{560.444}{the}%
\TextAt{8}{376.912}{560.444}{value}%
\TextAt{8}{407.606}{560.444}{groups}%
\TextAt{8}{445.844}{560.444}{and}%
\TextAt{8}{469.055}{560.444}{residue}%
\TextAt{8}{509.23}{560.444}{class}%
\TextAt{8}{79.226}{543.398}{fields}%
\TextAt{8}{109.198}{543.398}{of}%
\TextAt{13}{122.528}{543.398}{v}%
\TextAt{8}{132.513}{543.398}{and}%
\TextAt{13}{155.276}{543.398}{w}%
\TextAt{8}{163.966}{543.398}{,}%
\TextAt{8}{171.115}{543.398}{one}%
\TextAt{8}{192.574}{543.398}{gets}%
\TextAt{8}{216.7}{543.398}{the}%
\TextAt{8}{236.856}{543.398}{inclusions}%
\TextAt{13}{217.688}{514.397}{v}%
\TextAt{9}{223.777}{514.397}{\Character{40}}%
\TextAt{13}{228.329}{514.397}{K}%
\TextAt{3}{239.147}{519.333}{\Character{3}}%
\TextAt{9}{243.879}{514.397}{\Character{41}}%
\NativeMathAt{251.752}{514.397}{11.955}{\subseteq}%
\TextAt{13}{264.371}{514.397}{w}%
\TextAt{9}{273.061}{514.397}{\Character{40}}%
\TextAt{13}{277.613}{514.397}{L}%
\TextAt{3}{285.578}{519.333}{\Character{3}}%
\TextAt{9}{290.31}{514.397}{\Character{41}}%
\TextAt{8}{318.773}{514.397}{and}%
\TextAt{13}{361.548}{514.397}{\Character{20}}%
\NativeMathAt{371.608}{514.397}{11.955}{\subseteq}%
\TextAt{13}{384.227}{514.397}{\Character{21}\Character{58}}%
\TextAt{8}{79.226}{485.396}{The}%
\TextAt{8}{103.28}{485.396}{index}%
\TextAt{13}{227.788}{468.35}{e}%
\TextAt{9}{236.534}{468.35}{\Character{61}}%
\TextAt{13}{248.959}{468.35}{e}%
\TextAt{9}{254.385}{468.35}{\Character{40}}%
\TextAt{13}{258.937}{468.35}{w}%
\TextAt{16}{270.948}{468.35}{\Character{106}}%
\TextAt{13}{277.589}{468.35}{v}%
\TextAt{9}{283.677}{468.35}{\Character{41}}%
\TextAt{9}{291.532}{468.35}{\Character{61}}%
\TextAt{9}{303.953}{468.35}{\Character{40}}%
\TextAt{13}{308.528}{468.35}{w}%
\TextAt{9}{317.218}{468.35}{\Character{40}}%
\TextAt{13}{321.77}{468.35}{L}%
\TextAt{3}{329.735}{473.286}{\Character{3}}%
\TextAt{9}{334.467}{468.35}{\Character{41}}%
\TextAt{9}{342.322}{468.35}{\Character{58}}%
\TextAt{13}{348.913}{468.35}{v}%
\TextAt{9}{355.001}{468.35}{\Character{40}}%
\TextAt{13}{359.553}{468.35}{K}%
\TextAt{3}{370.371}{473.286}{\Character{3}}%
\TextAt{9}{375.103}{468.35}{\Character{41}\Character{41}}%
\TextAt{8}{79.226}{444.828}{is}%
\TextAt{8}{90.99}{444.828}{called}%
\TextAt{8}{124.154}{444.828}{the}%
\TextAt{12}{144.31}{444.828}{ramification}%
\TextAt{12}{218.958}{444.828}{index}%
\TextAt{8}{254.776}{444.828}{of}%
\TextAt{8}{268.106}{444.828}{the}%
\TextAt{8}{288.262}{444.828}{extension}%
\TextAt{13}{340.028}{444.828}{L}%
\TextAt{16}{351.314}{444.828}{\Character{106}}%
\TextAt{13}{357.955}{444.828}{K}%
\TextAt{8}{372.67}{444.828}{and}%
\TextAt{8}{395.433}{444.828}{the}%
\TextAt{8}{415.589}{444.828}{degree}%
\TextAt{13}{251.301}{415.827}{f}%
\TextAt{9}{261.668}{415.827}{\Character{61}}%
\TextAt{13}{274.093}{415.827}{f}%
\TextAt{9}{281.14}{415.827}{\Character{40}}%
\TextAt{13}{285.692}{415.827}{w}%
\TextAt{16}{297.703}{415.827}{\Character{106}}%
\TextAt{13}{304.344}{415.827}{v}%
\TextAt{9}{310.433}{415.827}{\Character{41}}%
\TextAt{9}{318.288}{415.827}{\Character{61}}%
\TextAt{9}{330.709}{415.827}{\Character{91}}%
\TextAt{13}{333.983}{415.827}{\Character{21}}%
\TextAt{9}{344.132}{415.827}{\Character{58}}%
\TextAt{13}{350.705}{415.827}{\Character{20}}%
\TextAt{9}{357.443}{415.827}{\Character{93}}%
\TextAt{8}{79.226}{386.826}{is}%
\TextAt{8}{91.728}{386.826}{called}%
\TextAt{8}{125.63}{386.826}{the}%
\TextAt{12}{146.525}{386.826}{inertia}%
\TextAt{12}{190.222}{386.826}{degree}%
\TextAt{8}{228.323}{386.826}{.}%
\TextAt{8}{238.986}{386.826}{If}%
\TextAt{13}{251.417}{386.826}{v}%
\TextAt{8}{257.505}{386.826}{,}%
\TextAt{8}{265.577}{386.826}{and}%
\TextAt{8}{289.078}{386.826}{hence}%
\TextAt{13}{322.323}{386.826}{w}%
\TextAt{9}{335.591}{386.826}{\Character{61}}%
\TextAt{17}{350.923}{391.533}{1}%
\Vector{q 0 G 0 g 0.478 w 278.471 -384.185 m 283.609 -384.185 l S Q}%
\TextAt{21}{350.471}{382.704}{n}%
\TextAt{13}{356.804}{386.826}{v}%
\TextAt{16}{366.053}{386.826}{\Character{14}}%
\TextAt{13}{375.19}{386.826}{N}%
\TextAt{21}{384.564}{384.966}{L}%
\TextAt{3}{390.323}{384.966}{\Character{106}}%
\TextAt{21}{392.675}{384.966}{K}%
\TextAt{8}{400.878}{386.826}{,}%
\TextAt{8}{408.95}{386.826}{is}%
\TextAt{8}{421.452}{386.826}{discrete}%
\TextAt{8}{465.17}{386.826}{and}%
\TextAt{8}{488.671}{386.826}{if}%
\TextAt{14}{500.133}{386.826}{o}%
\TextAt{13}{505.977}{386.826}{\Character{59}}%
\TextAt{14}{511.221}{386.826}{p}%
\TextAt{13}{517.204}{386.826}{\Character{59}}%
\TextAt{13}{522.428}{386.826}{\Character{25}}%
\TextAt{8}{529.518}{386.826}{,}%
\TextAt{8}{79.226}{369.78}{resp.}%
\TextAt{16}{106.335}{369.78}{\Character{79}}%
\TextAt{13}{116.185}{369.78}{\Character{59}}%
\TextAt{14}{121.429}{369.78}{P}%
\TextAt{13}{131.322}{369.78}{\Character{59}}%
\TextAt{9}{136.566}{369.78}{\Character{5}}%
\TextAt{8}{145.341}{369.78}{,}%
\TextAt{8}{151.759}{369.78}{are}%
\TextAt{8}{170.356}{369.78}{the}%
\TextAt{8}{189.598}{369.78}{valuation}%
\TextAt{8}{239.756}{369.78}{ring,}%
\TextAt{8}{266.342}{369.78}{the}%
\TextAt{8}{285.584}{369.78}{maximal}%
\TextAt{8}{332.478}{369.78}{ideal}%
\TextAt{8}{359.527}{369.78}{and}%
\TextAt{8}{381.376}{369.78}{a}%
\TextAt{8}{390.217}{369.78}{prime}%
\TextAt{8}{422.466}{369.78}{element}%
\TextAt{8}{464.794}{369.78}{of}%
\TextAt{13}{477.209}{369.78}{K}%
\TextAt{8}{488.027}{369.78}{,}%
\TextAt{8}{494.445}{369.78}{resp.}%
\TextAt{13}{521.553}{369.78}{L}%
\TextAt{8}{529.518}{369.78}{,}%
\TextAt{8}{79.226}{352.734}{then}%
\TextAt{8}{105.886}{352.734}{one}%
\TextAt{8}{127.346}{352.734}{has}%
\TextAt{13}{251.901}{335.688}{e}%
\TextAt{9}{260.647}{335.688}{\Character{61}}%
\TextAt{9}{273.057}{335.688}{\Character{40}}%
\TextAt{13}{277.625}{335.688}{w}%
\TextAt{9}{286.315}{335.688}{\Character{40}\Character{5}\Character{41}}%
\TextAt{18}{304.195}{335.688}{Z}%
\TextAt{9}{315.486}{335.688}{\Character{58}}%
\TextAt{13}{322.058}{335.688}{v}%
\TextAt{9}{328.146}{335.688}{\Character{40}}%
\TextAt{13}{332.699}{335.688}{\Character{25}}%
\TextAt{9}{339.768}{335.688}{\Character{41}}%
\TextAt{18}{344.32}{335.688}{Z}%
\TextAt{9}{352.291}{335.688}{\Character{41}}%
\TextAt{13}{356.843}{335.688}{\Character{59}}%
\TextAt{8}{79.226}{312.167}{so}%
\TextAt{8}{93.596}{312.167}{that}%
\TextAt{13}{118.965}{312.167}{v}%
\TextAt{9}{125.053}{312.167}{\Character{40}}%
\TextAt{13}{129.606}{312.167}{\Character{25}}%
\TextAt{9}{136.675}{312.167}{\Character{41}}%
\TextAt{9}{144.53}{312.167}{\Character{61}}%
\TextAt{13}{156.974}{312.167}{ew}%
\TextAt{9}{171.089}{312.167}{\Character{40}\Character{5}\Character{41}}%
\TextAt{8}{188.969}{312.167}{,}%
\TextAt{8}{196.118}{312.167}{and}%
\TextAt{8}{218.88}{312.167}{we}%
\TextAt{8}{236.108}{312.167}{find}%
\TextAt{13}{283.624}{295.121}{\Character{25}}%
\TextAt{9}{294.014}{295.121}{\Character{61}}%
\TextAt{13}{306.44}{295.121}{\Character{34}}%
\TextAt{9}{311.918}{295.121}{\Character{5}}%
\TextAt{21}{320.693}{300.057}{e}%
\TextAt{13}{325.12}{295.121}{\Character{59}}%
\TextAt{8}{79.226}{271.599}{for}%
\TextAt{8}{97.395}{271.599}{some}%
\TextAt{8}{127.005}{271.599}{unit}%
\TextAt{13}{152.001}{271.599}{\Character{34}}%
\NativeMathAt{161.284}{271.599}{11.955}{\in}%
\TextAt{16}{173.048}{271.599}{\Character{79}}%
\TextAt{3}{182.909}{275.937}{\Character{3}}%
\TextAt{8}{187.641}{271.599}{.}%
\TextAt{8}{196.942}{271.599}{From}%
\TextAt{8}{227.95}{271.599}{this}%
\TextAt{8}{251.057}{271.599}{one}%
\TextAt{8}{272.801}{271.599}{deduces}%
\TextAt{8}{316.709}{271.599}{the}%
\TextAt{8}{337.15}{271.599}{familiar}%
\TextAt{8}{380.687}{271.599}{(see}%
\TextAt{8}{404.44}{271.599}{chap.}%
\TextAt{8}{435.616}{271.599}{I)}%
\TextAt{8}{448.572}{271.599}{interpretation}%
\TextAt{8}{523.337}{271.599}{of}%
\TextAt{8}{79.226}{254.553}{the}%
\TextAt{8}{99.383}{254.553}{ramification}%
\TextAt{8}{164.431}{254.553}{index:}%
\TextAt{14}{200.519}{254.553}{p}%
\TextAt{16}{206.502}{254.553}{\Character{79}}%
\TextAt{9}{219.673}{254.553}{\Character{61}}%
\TextAt{13}{232.098}{254.553}{\Character{25}}%
\TextAt{16}{239.168}{254.553}{\Character{79}}%
\TextAt{9}{252.338}{254.553}{\Character{61}}%
\TextAt{9}{264.748}{254.553}{\Character{5}}%
\TextAt{21}{273.539}{258.891}{e}%
\TextAt{16}{277.966}{254.553}{\Character{79}}%
\TextAt{9}{291.136}{254.553}{\Character{61}}%
\TextAt{14}{303.562}{254.553}{P}%
\TextAt{21}{313.455}{258.891}{e}%
\TextAt{8}{317.882}{254.553}{,}%
\TextAt{8}{325.031}{254.553}{or}%
\TextAt{14}{286.347}{225.552}{p}%
\TextAt{9}{295.651}{225.552}{\Character{61}}%
\TextAt{14}{308.077}{225.552}{P}%
\TextAt{21}{317.97}{230.488}{e}%
\TextAt{13}{322.397}{225.552}{\Character{58}}%
\TextAt{12}{79.226}{187.584}{Proposition}%
\TextAt{12}{151.937}{187.584}{(6.8).}%
\TextAt{8}{187.253}{187.584}{One}%
\TextAt{8}{211.964}{187.584}{has}%
\TextAt{9}{232.838}{187.584}{\Character{91}}%
\TextAt{13}{236.09}{187.584}{L}%
\TextAt{9}{247.375}{187.584}{\Character{58}}%
\TextAt{13}{253.948}{187.584}{K}%
\TextAt{9}{264.765}{187.584}{\Character{93}}%
\NativeMathAt{271.338}{187.584}{11.955}{\geq}%
\TextAt{13}{283.957}{187.584}{ef}%
\TextAt{8}{300.326}{187.584}{and}%
\TextAt{8}{323.089}{187.584}{the}%
\TextAt{12}{343.245}{187.584}{fundamental}%
\TextAt{12}{420.763}{187.584}{identity}%
\TextAt{9}{272.999}{158.583}{\Character{91}}%
\TextAt{13}{276.25}{158.583}{L}%
\TextAt{9}{287.536}{158.583}{\Character{58}}%
\TextAt{13}{294.108}{158.583}{K}%
\TextAt{9}{304.926}{158.583}{\Character{93}}%
\TextAt{9}{311.477}{158.583}{\Character{61}}%
\TextAt{13}{323.924}{158.583}{ef}%
\TextAt{13}{335.724}{158.583}{\Character{59}}%
\TextAt{8}{79.226}{129.582}{if}%
\TextAt{13}{89.95}{129.582}{v}%
\TextAt{8}{99.936}{129.582}{is}%
\TextAt{8}{111.699}{129.582}{discrete}%
\TextAt{8}{154.678}{129.582}{and}%
\TextAt{13}{177.441}{129.582}{L}%
\TextAt{16}{188.726}{129.582}{\Character{106}}%
\TextAt{13}{195.368}{129.582}{K}%
\TextAt{8}{210.083}{129.582}{is}%
\TextAt{8}{221.847}{129.582}{separable.}%
\TextAt{8}{297.211}{43.033}{137}%
\Destination{78.23}{809.815}{nav.page.137}%
\Destination{79.36}{575.455}{nav.print.150}%
\Destination{79.36}{201.555}{nav.number.2.6.8}%
\Hotlink{435.116}{269.355}{5.22}{11.047}{nav.chapter.1}%
\Hotlink{233.235}{758.955}{25.078}{11.047}{nav.number.2.6.3}%
\Hotlink{172.866}{558.255}{25.078}{11.047}{nav.number.2.6.2}%
}

% Page 150 (body page 138)
\Page{%
\TextAt{2}{79.226}{796.943}{II.}%
\TextAt{2}{94.073}{796.943}{The}%
\TextAt{2}{116.492}{796.943}{Theory}%
\TextAt{2}{154.401}{796.943}{of}%
\TextAt{2}{166.826}{796.943}{Valuations}%
\Vector{q 0 G 0 g 0.299 w 7.226 18.729 m 460.769 18.729 l S Q}%
\TextAt{19}{79.226}{761.124}{Proof.}%
\TextAt{8}{116}{761.124}{Let}%
\TextAt{13}{137.514}{761.124}{\Character{33}}%
\TextAt{17}{144.803}{759.33}{1}%
\TextAt{13}{149.535}{761.124}{\Character{59}}%
\TextAt{13}{154.759}{761.124}{\Character{58}}%
\TextAt{13}{159.996}{761.124}{\Character{58}}%
\TextAt{13}{165.232}{761.124}{\Character{58}}%
\TextAt{13}{170.457}{761.124}{\Character{59}}%
\TextAt{13}{175.693}{761.124}{\Character{33}}%
\TextAt{21}{183.045}{759.33}{f}%
\TextAt{8}{192.932}{761.124}{be}%
\TextAt{8}{209.4}{761.124}{representatives}%
\TextAt{8}{289.745}{761.124}{of}%
\TextAt{8}{303.619}{761.124}{a}%
\TextAt{8}{313.918}{761.124}{basis}%
\TextAt{8}{343.203}{761.124}{of}%
\TextAt{13}{357.077}{761.124}{\Character{21}}%
\TextAt{16}{368.153}{761.124}{\Character{106}}%
\TextAt{13}{375.722}{761.124}{\Character{20}}%
\TextAt{8}{386.902}{761.124}{and}%
\TextAt{8}{410.209}{761.124}{let}%
\TextAt{13}{427.657}{761.124}{\Character{25}}%
\TextAt{17}{434.298}{759.33}{0}%
\TextAt{13}{439.03}{761.124}{\Character{59}}%
\TextAt{13}{444.254}{761.124}{\Character{58}}%
\TextAt{13}{449.491}{761.124}{\Character{58}}%
\TextAt{13}{454.715}{761.124}{\Character{58}}%
\TextAt{13}{459.952}{761.124}{\Character{59}}%
\TextAt{13}{465.188}{761.124}{\Character{25}}%
\TextAt{21}{471.891}{759.33}{e}%
\TextAt{3}{475.819}{759.33}{\Character{0}}%
\TextAt{17}{482.406}{759.33}{1}%
\NativeMathAt{491.386}{761.124}{11.955}{\in}%
\TextAt{13}{503.604}{761.124}{L}%
\TextAt{3}{511.569}{765.462}{\Character{3}}%
\TextAt{8}{520.743}{761.124}{be}%
\TextAt{8}{79.226}{744.078}{elements}%
\TextAt{8}{127.326}{744.078}{the}%
\TextAt{8}{147.726}{744.078}{values}%
\TextAt{8}{182.831}{744.078}{of}%
\TextAt{8}{196.405}{744.078}{which}%
\TextAt{8}{230.135}{744.078}{represent}%
\TextAt{8}{280.841}{744.078}{the}%
\TextAt{8}{301.241}{744.078}{various}%
\TextAt{8}{341.558}{744.078}{cosets}%
\TextAt{8}{375.742}{744.078}{in}%
\TextAt{13}{389.639}{744.078}{w}%
\TextAt{9}{398.328}{744.078}{\Character{40}}%
\TextAt{13}{402.881}{744.078}{L}%
\TextAt{3}{410.845}{748.416}{\Character{3}}%
\TextAt{9}{415.578}{744.078}{\Character{41}\Character{47}}%
\TextAt{13}{425.983}{744.078}{v}%
\TextAt{9}{432.071}{744.078}{\Character{40}}%
\TextAt{13}{436.623}{744.078}{K}%
\TextAt{3}{447.441}{748.416}{\Character{3}}%
\TextAt{9}{452.173}{744.078}{\Character{41}}%
\TextAt{8}{460.867}{744.078}{(the}%
\TextAt{8}{485.822}{744.078}{finiteness}%
\TextAt{8}{79.226}{727.032}{of}%
\TextAt{13}{92.968}{727.032}{e}%
\TextAt{8}{102.703}{727.032}{will}%
\TextAt{8}{125.22}{727.032}{be}%
\TextAt{8}{141.556}{727.032}{a}%
\TextAt{8}{151.723}{727.032}{consequence}%
\TextAt{8}{218.2}{727.032}{of}%
\TextAt{8}{231.942}{727.032}{what}%
\TextAt{8}{261.62}{727.032}{follows).}%
\TextAt{8}{310.397}{727.032}{If}%
\TextAt{13}{322.501}{727.032}{v}%
\TextAt{8}{332.899}{727.032}{is}%
\TextAt{8}{345.075}{727.032}{discrete,}%
\TextAt{8}{391.82}{727.032}{we}%
\TextAt{8}{409.459}{727.032}{may}%
\TextAt{8}{435.24}{727.032}{choose}%
\TextAt{8}{472.785}{727.032}{for}%
\TextAt{8}{491.082}{727.032}{instance}%
\TextAt{13}{79.226}{709.986}{\Character{25}}%
\TextAt{21}{85.866}{708.193}{i}%
\TextAt{9}{92.569}{709.986}{\Character{61}}%
\TextAt{9}{104.979}{709.986}{\Character{5}}%
\TextAt{21}{113.769}{714.324}{i}%
\TextAt{8}{117.151}{709.986}{.}%
\TextAt{8}{125.599}{709.986}{We}%
\TextAt{8}{145.743}{709.986}{show}%
\TextAt{8}{174.747}{709.986}{that}%
\TextAt{8}{200.116}{709.986}{the}%
\TextAt{8}{220.272}{709.986}{elements}%
\TextAt{13}{207.608}{683.979}{\Character{33}}%
\TextAt{21}{214.898}{682.185}{j}%
\TextAt{13}{219.28}{683.979}{\Character{25}}%
\TextAt{21}{225.92}{682.185}{i}%
\TextAt{13}{229.301}{683.979}{\Character{59}}%
\TextAt{13}{246.481}{683.979}{j}%
\TextAt{9}{255.337}{683.979}{\Character{61}}%
\TextAt{9}{267.747}{683.979}{1}%
\TextAt{13}{273.616}{683.979}{\Character{59}}%
\TextAt{13}{278.84}{683.979}{\Character{58}}%
\TextAt{13}{284.077}{683.979}{\Character{58}}%
\TextAt{13}{289.301}{683.979}{\Character{58}}%
\TextAt{13}{294.538}{683.979}{\Character{59}}%
\TextAt{13}{299.774}{683.979}{f}%
\TextAt{13}{306.158}{683.979}{\Character{59}}%
\TextAt{13}{323.35}{683.979}{i}%
\TextAt{9}{330.746}{683.979}{\Character{61}}%
\TextAt{9}{343.156}{683.979}{0}%
\TextAt{13}{349.025}{683.979}{\Character{59}}%
\TextAt{13}{354.249}{683.979}{\Character{58}}%
\TextAt{13}{359.486}{683.979}{\Character{58}}%
\TextAt{13}{364.71}{683.979}{\Character{58}}%
\TextAt{13}{369.947}{683.979}{\Character{59}}%
\TextAt{13}{375.183}{683.979}{e}%
\TextAt{16}{383.327}{683.979}{\Character{0}}%
\TextAt{9}{395.283}{683.979}{1}%
\TextAt{13}{401.136}{683.979}{\Character{59}}%
\TextAt{8}{79.226}{657.972}{are}%
\TextAt{8}{98.607}{657.972}{linearly}%
\TextAt{8}{140.427}{657.972}{independent}%
\TextAt{8}{206.624}{657.972}{over}%
\TextAt{13}{231.54}{657.972}{K}%
\TextAt{8}{242.358}{657.972}{,}%
\TextAt{8}{249.403}{657.972}{and}%
\TextAt{8}{272.035}{657.972}{in}%
\TextAt{8}{285.558}{657.972}{the}%
\TextAt{8}{305.584}{657.972}{discrete}%
\TextAt{8}{348.432}{657.972}{case}%
\TextAt{8}{373.073}{657.972}{form}%
\TextAt{8}{400.583}{657.972}{even}%
\TextAt{8}{427.113}{657.972}{a}%
\TextAt{8}{436.738}{657.972}{basis}%
\TextAt{8}{465.348}{657.972}{of}%
\TextAt{13}{478.548}{657.972}{L}%
\TextAt{16}{489.833}{657.972}{\Character{106}}%
\TextAt{13}{496.475}{657.972}{K}%
\TextAt{8}{507.293}{657.972}{.}%
\TextAt{8}{515.697}{657.972}{Let}%
\TextAt{21}{259.645}{637.808}{e}%
\TextAt{3}{263.574}{637.808}{\Character{0}}%
\TextAt{17}{270.16}{637.808}{1}%
\TextAt{20}{258.385}{634.222}{\Character{88}}%
\TextAt{21}{260.168}{609.027}{i}%
\TextAt{17}{263.051}{609.027}{\Character{61}0}%
\TextAt{21}{283.807}{638.295}{f}%
\TextAt{20}{277.647}{634.222}{\Character{88}}%
\TextAt{21}{278.928}{609.027}{j}%
\TextAt{17}{282.813}{609.027}{\Character{61}1}%
\TextAt{13}{296.908}{622.864}{a}%
\TextAt{21}{303.053}{621.071}{ij}%
\TextAt{13}{310.318}{622.864}{\Character{33}}%
\TextAt{21}{317.607}{621.071}{j}%
\TextAt{13}{321.989}{622.864}{\Character{25}}%
\TextAt{21}{328.63}{621.071}{i}%
\TextAt{9}{335.332}{622.864}{\Character{61}}%
\TextAt{9}{347.742}{622.864}{0}%
\TextAt{8}{79.226}{584.978}{with}%
\TextAt{13}{104.612}{584.978}{a}%
\TextAt{21}{110.757}{583.185}{ij}%
\NativeMathAt{121.343}{584.978}{11.955}{\in}%
\TextAt{13}{132.634}{584.978}{K}%
\TextAt{8}{143.452}{584.978}{.}%
\TextAt{8}{151.476}{584.978}{Assume}%
\TextAt{8}{193.563}{584.978}{that}%
\TextAt{8}{217.658}{584.978}{not}%
\TextAt{8}{237.199}{584.978}{all}%
\TextAt{13}{252.184}{584.978}{a}%
\TextAt{21}{258.329}{583.185}{ij}%
\TextAt{9}{268.915}{584.978}{\Character{61}}%
\TextAt{9}{281.325}{584.978}{0}%
\TextAt{8}{287.193}{584.978}{.}%
\TextAt{8}{295.217}{584.978}{Then}%
\TextAt{8}{324.501}{584.978}{there}%
\TextAt{8}{353.139}{584.978}{exist}%
\TextAt{8}{379.565}{584.978}{nonzero}%
\TextAt{8}{421.868}{584.978}{sums}%
\TextAt{13}{449.98}{584.978}{s}%
\TextAt{21}{455.494}{583.185}{i}%
\TextAt{9}{462.196}{584.978}{\Character{61}}%
\TextAt{20}{474.621}{593.945}{\Character{80}}%
\TextAt{21}{487.241}{590.79}{f}%
\TextAt{21}{487.241}{581.491}{j}%
\TextAt{17}{491.125}{581.491}{\Character{61}1}%
\TextAt{13}{504.436}{584.978}{a}%
\TextAt{21}{510.581}{583.185}{ij}%
\TextAt{13}{517.846}{584.978}{\Character{33}}%
\TextAt{21}{525.136}{583.185}{j}%
\TextAt{8}{529.518}{584.978}{,}%
\TextAt{8}{79.226}{567.932}{and}%
\TextAt{8}{102.457}{567.932}{each}%
\TextAt{8}{129.262}{567.932}{time}%
\TextAt{8}{156.391}{567.932}{that}%
\TextAt{13}{182.228}{567.932}{s}%
\TextAt{21}{187.742}{566.139}{i}%
\NativeMathAt{195.242}{567.932}{11.955}{\neq}%
\TextAt{9}{208.453}{567.932}{0}%
\TextAt{8}{218.684}{567.932}{we}%
\TextAt{8}{236.379}{567.932}{find}%
\TextAt{13}{260.256}{567.932}{w}%
\TextAt{9}{268.945}{567.932}{\Character{40}}%
\TextAt{13}{273.498}{567.932}{s}%
\TextAt{21}{279.012}{566.139}{i}%
\TextAt{9}{282.393}{567.932}{\Character{41}}%
\NativeMathAt{291.064}{567.932}{11.955}{\in}%
\TextAt{13}{303.153}{567.932}{v}%
\TextAt{9}{309.241}{567.932}{\Character{40}}%
\TextAt{13}{313.793}{567.932}{K}%
\TextAt{3}{324.611}{572.271}{\Character{3}}%
\TextAt{9}{329.343}{567.932}{\Character{41}}%
\TextAt{8}{333.896}{567.932}{.}%
\TextAt{8}{343.748}{567.932}{In}%
\TextAt{8}{358.838}{567.932}{fact,}%
\TextAt{8}{385.76}{567.932}{dividing}%
\TextAt{13}{431.431}{567.932}{s}%
\TextAt{21}{436.945}{566.139}{i}%
\TextAt{8}{444.692}{567.932}{by}%
\TextAt{8}{461.419}{567.932}{the}%
\TextAt{8}{482.044}{567.932}{coefficient}%
\TextAt{13}{79.226}{550.886}{a}%
\TextAt{21}{85.371}{549.093}{i\Character{23}}%
\TextAt{8}{97.794}{550.886}{of}%
\TextAt{8}{111.584}{550.886}{minimum}%
\TextAt{8}{164.396}{550.886}{value,}%
\TextAt{8}{198.469}{550.886}{we}%
\TextAt{8}{216.157}{550.886}{get}%
\TextAt{8}{236.128}{550.886}{a}%
\TextAt{8}{246.343}{550.886}{linear}%
\TextAt{8}{279.321}{550.886}{combination}%
\TextAt{8}{346.455}{550.886}{of}%
\TextAt{8}{360.245}{550.886}{the}%
\TextAt{13}{380.862}{550.886}{\Character{33}}%
\TextAt{17}{388.151}{549.093}{1}%
\TextAt{13}{392.884}{550.886}{\Character{59}}%
\TextAt{13}{398.108}{550.886}{\Character{58}}%
\TextAt{13}{403.345}{550.886}{\Character{58}}%
\TextAt{13}{408.569}{550.886}{\Character{58}}%
\TextAt{13}{413.806}{550.886}{\Character{59}}%
\TextAt{13}{419.042}{550.886}{\Character{33}}%
\TextAt{21}{426.394}{549.093}{f}%
\TextAt{8}{436.196}{550.886}{with}%
\TextAt{8}{463.316}{550.886}{coefficients}%
\TextAt{8}{523.014}{550.886}{in}%
\TextAt{8}{79.226}{533.84}{the}%
\TextAt{8}{99.661}{533.84}{valuation}%
\TextAt{8}{151.013}{533.84}{ring}%
\TextAt{14}{175.357}{533.84}{o}%
\NativeMathAt{184.998}{533.84}{11.955}{\subseteq}%
\TextAt{13}{198.092}{533.84}{K}%
\TextAt{8}{213.086}{533.84}{one}%
\TextAt{8}{234.824}{533.84}{of}%
\TextAt{8}{248.433}{533.84}{which}%
\TextAt{8}{282.199}{533.84}{equals}%
\TextAt{8}{317.984}{533.84}{1.}%
\TextAt{8}{333.127}{533.84}{This}%
\TextAt{8}{360.126}{533.84}{linear}%
\TextAt{8}{392.923}{533.84}{combination}%
\TextAt{8}{459.875}{533.84}{is}%
\NativeMathAt{471.918}{533.84}{11.955}{\not\equiv}%
\TextAt{9}{485.013}{533.84}{0}%
\TextAt{8}{494.187}{533.84}{mod}%
\TextAt{14}{519.624}{533.84}{P}%
\TextAt{8}{529.518}{533.84}{,}%
\TextAt{8}{79.226}{516.794}{hence}%
\TextAt{8}{111.732}{516.794}{a}%
\TextAt{8}{121.488}{516.794}{unit,}%
\TextAt{8}{149.451}{516.794}{so}%
\TextAt{8}{163.821}{516.794}{that}%
\TextAt{13}{189.19}{516.794}{w}%
\TextAt{9}{197.879}{516.794}{\Character{40}}%
\TextAt{13}{202.432}{516.794}{s}%
\TextAt{21}{207.946}{515.001}{i}%
\TextAt{9}{211.327}{516.794}{\Character{41}}%
\TextAt{9}{219.182}{516.794}{\Character{61}}%
\TextAt{13}{231.626}{516.794}{w}%
\TextAt{9}{240.315}{516.794}{\Character{40}}%
\TextAt{13}{244.867}{516.794}{a}%
\TextAt{21}{251.012}{515.001}{i\Character{23}}%
\TextAt{9}{259.078}{516.794}{\Character{41}}%
\NativeMathAt{266.951}{516.794}{11.955}{\in}%
\TextAt{13}{278.242}{516.794}{v}%
\TextAt{9}{284.33}{516.794}{\Character{40}}%
\TextAt{13}{288.882}{516.794}{K}%
\TextAt{3}{299.7}{521.133}{\Character{3}}%
\TextAt{9}{304.432}{516.794}{\Character{41}}%
\TextAt{8}{308.985}{516.794}{.}%
\TextAt{8}{97.159}{496.759}{In}%
\TextAt{8}{110.835}{496.759}{the}%
\TextAt{8}{130.047}{496.759}{sum}%
\TextAt{20}{153.873}{505.726}{\Character{80}}%
\TextAt{21}{166.492}{502.571}{e}%
\TextAt{3}{170.42}{502.571}{\Character{0}}%
\TextAt{17}{177.007}{502.571}{1}%
\TextAt{21}{166.492}{493.272}{i}%
\TextAt{17}{169.375}{493.272}{\Character{61}0}%
\TextAt{13}{183.732}{496.759}{s}%
\TextAt{21}{189.246}{494.966}{i}%
\TextAt{13}{192.627}{496.759}{\Character{25}}%
\TextAt{21}{199.267}{494.966}{i}%
\TextAt{8}{202.649}{496.759}{,}%
\TextAt{8}{209.042}{496.759}{two}%
\TextAt{8}{230.214}{496.759}{nonzero}%
\TextAt{8}{272.845}{496.759}{summands}%
\TextAt{8}{329.907}{496.759}{must}%
\TextAt{8}{357.965}{496.759}{have}%
\TextAt{8}{384.015}{496.759}{the}%
\TextAt{8}{403.226}{496.759}{same}%
\TextAt{8}{431.607}{496.759}{value,}%
\TextAt{8}{464.35}{496.759}{say}%
\TextAt{13}{483.633}{496.759}{w}%
\TextAt{9}{492.322}{496.759}{\Character{40}}%
\TextAt{13}{496.875}{496.759}{s}%
\TextAt{21}{502.389}{494.966}{i}%
\TextAt{13}{505.77}{496.759}{\Character{25}}%
\TextAt{21}{512.41}{494.966}{i}%
\TextAt{9}{515.792}{496.759}{\Character{41}}%
\TextAt{9}{523.647}{496.759}{\Character{61}}%
\TextAt{13}{79.226}{479.714}{w}%
\TextAt{9}{87.916}{479.714}{\Character{40}}%
\TextAt{13}{92.468}{479.714}{s}%
\TextAt{21}{97.982}{477.92}{j}%
\TextAt{13}{102.364}{479.714}{\Character{25}}%
\TextAt{21}{109.004}{477.92}{j}%
\TextAt{9}{113.387}{479.714}{\Character{41}}%
\TextAt{8}{117.939}{479.714}{,}%
\TextAt{13}{126.174}{479.714}{i}%
\NativeMathAt{134.969}{479.714}{11.955}{\neq}%
\TextAt{13}{148.875}{479.714}{j}%
\TextAt{8}{154.391}{479.714}{,}%
\TextAt{8}{162.626}{479.714}{because}%
\TextAt{8}{206.796}{479.714}{otherwise}%
\TextAt{8}{259.754}{479.714}{it}%
\TextAt{8}{272.327}{479.714}{could}%
\TextAt{8}{304.411}{479.714}{not}%
\TextAt{8}{326.094}{479.714}{be}%
\TextAt{8}{342.887}{479.714}{zero}%
\TextAt{8}{368.467}{479.714}{(observe}%
\TextAt{8}{415.578}{479.714}{that}%
\TextAt{13}{441.816}{479.714}{w}%
\TextAt{9}{450.506}{479.714}{\Character{40}}%
\TextAt{13}{455.058}{479.714}{x}%
\TextAt{9}{461.71}{479.714}{\Character{41}}%
\NativeMathAt{471.064}{479.714}{11.955}{\neq}%
\TextAt{13}{484.97}{479.714}{w}%
\TextAt{9}{493.659}{479.714}{\Character{40}}%
\TextAt{13}{498.212}{479.714}{y}%
\TextAt{9}{504.348}{479.714}{\Character{41}}%
\NativeArrowAt{513.685}{479.714}{11.955}{19.084}{\Longrightarrow}%
\TextAt{13}{79.226}{462.668}{w}%
\TextAt{9}{87.916}{462.668}{\Character{40}}%
\TextAt{13}{92.468}{462.668}{x}%
\TextAt{9}{101.777}{462.668}{\Character{43}}%
\TextAt{13}{113.538}{462.668}{y}%
\TextAt{9}{119.675}{462.668}{\Character{41}}%
\TextAt{9}{127.53}{462.668}{\Character{61}}%
\TextAt{8}{139.973}{462.668}{min}%
\TextAt{16}{159.484}{462.668}{\Character{102}}%
\TextAt{13}{165.462}{462.668}{w}%
\TextAt{9}{174.151}{462.668}{\Character{40}}%
\TextAt{13}{178.704}{462.668}{x}%
\TextAt{9}{185.356}{462.668}{\Character{41}}%
\TextAt{13}{189.908}{462.668}{\Character{59}}%
\TextAt{13}{195.132}{462.668}{w}%
\TextAt{9}{203.842}{462.668}{\Character{40}}%
\TextAt{13}{208.394}{462.668}{y}%
\TextAt{9}{214.531}{462.668}{\Character{41}}%
\TextAt{16}{219.083}{462.668}{\Character{103}}%
\TextAt{8}{225.061}{462.668}{).}%
\TextAt{8}{238.064}{462.668}{It}%
\TextAt{8}{250.736}{462.668}{follows}%
\TextAt{8}{289.172}{462.668}{that}%
\TextAt{13}{174.349}{436.66}{w}%
\TextAt{9}{183.038}{436.66}{\Character{40}}%
\TextAt{13}{187.59}{436.66}{\Character{25}}%
\TextAt{21}{194.231}{434.867}{i}%
\TextAt{9}{197.612}{436.66}{\Character{41}}%
\TextAt{9}{205.467}{436.66}{\Character{61}}%
\TextAt{13}{217.911}{436.66}{w}%
\TextAt{9}{226.6}{436.66}{\Character{40}}%
\TextAt{13}{231.153}{436.66}{\Character{25}}%
\TextAt{21}{237.793}{434.867}{j}%
\TextAt{9}{242.175}{436.66}{\Character{41}}%
\TextAt{9}{249.372}{436.66}{\Character{43}}%
\TextAt{13}{261.145}{436.66}{w}%
\TextAt{9}{269.835}{436.66}{\Character{40}}%
\TextAt{13}{274.387}{436.66}{s}%
\TextAt{21}{279.901}{434.867}{j}%
\TextAt{9}{284.283}{436.66}{\Character{41}}%
\TextAt{16}{291.492}{436.66}{\Character{0}}%
\TextAt{13}{303.447}{436.66}{w}%
\TextAt{9}{312.137}{436.66}{\Character{40}}%
\TextAt{13}{316.689}{436.66}{s}%
\TextAt{21}{322.203}{434.867}{i}%
\TextAt{9}{325.585}{436.66}{\Character{41}}%
\NativeMathAt{333.458}{436.66}{11.955}{\equiv}%
\TextAt{13}{346.077}{436.66}{w}%
\TextAt{9}{354.767}{436.66}{\Character{40}}%
\TextAt{13}{359.319}{436.66}{\Character{25}}%
\TextAt{21}{365.959}{434.867}{j}%
\TextAt{9}{370.341}{436.66}{\Character{41}}%
\TextAt{8}{378.215}{436.66}{mod}%
\TextAt{13}{403.652}{436.66}{v}%
\TextAt{9}{409.741}{436.66}{\Character{40}}%
\TextAt{13}{414.293}{436.66}{K}%
\TextAt{3}{425.111}{441.597}{\Character{3}}%
\TextAt{9}{429.843}{436.66}{\Character{41}}%
\TextAt{13}{434.395}{436.66}{\Character{59}}%
\TextAt{8}{79.226}{410.653}{a}%
\TextAt{8}{89.369}{410.653}{contradiction.}%
\TextAt{8}{166.311}{410.653}{This}%
\TextAt{8}{193.419}{410.653}{shows}%
\TextAt{8}{227.424}{410.653}{the}%
\TextAt{8}{247.968}{410.653}{linear}%
\TextAt{8}{280.874}{410.653}{independence}%
\TextAt{8}{353.757}{410.653}{of}%
\TextAt{8}{367.475}{410.653}{the}%
\TextAt{13}{388.019}{410.653}{\Character{33}}%
\TextAt{21}{395.308}{408.86}{j}%
\TextAt{13}{399.69}{410.653}{\Character{25}}%
\TextAt{21}{406.331}{408.86}{i}%
\TextAt{8}{409.712}{410.653}{.}%
\TextAt{8}{419.323}{410.653}{In}%
\TextAt{8}{434.331}{410.653}{particular,}%
\TextAt{8}{492.057}{410.653}{we}%
\TextAt{8}{509.672}{410.653}{have}%
\TextAt{13}{79.226}{393.607}{ef}%
\NativeMathAt{95.019}{393.607}{11.955}{\leq}%
\TextAt{9}{107.638}{393.607}{\Character{91}}%
\TextAt{13}{110.89}{393.607}{L}%
\TextAt{9}{122.175}{393.607}{\Character{58}}%
\TextAt{13}{128.748}{393.607}{K}%
\TextAt{9}{139.565}{393.607}{\Character{93}}%
\TextAt{8}{142.817}{393.607}{.}%
\TextAt{8}{97.159}{373.573}{Assume}%
\TextAt{8}{140.647}{373.573}{now}%
\TextAt{8}{165.162}{373.573}{that}%
\TextAt{13}{190.657}{373.573}{v}%
\TextAt{8}{196.746}{373.573}{,}%
\TextAt{8}{204.053}{373.573}{and}%
\TextAt{8}{226.942}{373.573}{thus}%
\TextAt{8}{252.82}{373.573}{also}%
\TextAt{13}{276.427}{373.573}{w}%
\TextAt{8}{285.116}{373.573}{,}%
\TextAt{8}{292.424}{373.573}{is}%
\TextAt{8}{304.314}{373.573}{discrete}%
\TextAt{8}{347.42}{373.573}{and}%
\TextAt{8}{370.309}{373.573}{let}%
\TextAt{9}{387.34}{373.573}{\Character{5}}%
\TextAt{8}{400.14}{373.573}{be}%
\TextAt{8}{416.19}{373.573}{a}%
\TextAt{8}{426.072}{373.573}{prime}%
\TextAt{8}{459.363}{373.573}{element}%
\TextAt{8}{502.731}{373.573}{in}%
\TextAt{8}{516.51}{373.573}{the}%
\TextAt{8}{79.226}{356.527}{valuation}%
\TextAt{8}{130.299}{356.527}{ring}%
\TextAt{16}{154.364}{356.527}{\Character{79}}%
\TextAt{8}{168.112}{356.527}{of}%
\TextAt{13}{181.442}{356.527}{w}%
\TextAt{8}{190.131}{356.527}{.}%
\TextAt{8}{198.579}{356.527}{We}%
\TextAt{8}{218.724}{356.527}{consider}%
\TextAt{8}{264.309}{356.527}{the}%
\TextAt{14}{284.465}{356.527}{o}%
\TextAt{8}{290.31}{356.527}{-module}%
\TextAt{13}{258.808}{321.957}{M}%
\TextAt{9}{274.703}{321.957}{\Character{61}}%
\TextAt{21}{288.388}{336.901}{e}%
\TextAt{3}{292.316}{336.901}{\Character{0}}%
\TextAt{17}{298.903}{336.901}{1}%
\TextAt{20}{287.128}{333.315}{\Character{88}}%
\TextAt{21}{288.91}{308.12}{i}%
\TextAt{17}{291.794}{308.12}{\Character{61}0}%
\TextAt{21}{312.55}{337.388}{f}%
\TextAt{20}{306.389}{333.315}{\Character{88}}%
\TextAt{21}{307.671}{308.12}{j}%
\TextAt{17}{311.555}{308.12}{\Character{61}1}%
\TextAt{14}{325.65}{321.957}{o}%
\TextAt{13}{331.495}{321.957}{\Character{33}}%
\TextAt{21}{338.784}{320.164}{j}%
\TextAt{13}{343.166}{321.957}{\Character{25}}%
\TextAt{21}{349.806}{320.164}{i}%
\TextAt{8}{79.226}{285.81}{where}%
\TextAt{13}{113.402}{285.81}{\Character{25}}%
\TextAt{21}{120.043}{284.017}{i}%
\TextAt{9}{127.37}{285.81}{\Character{61}}%
\TextAt{9}{140.413}{285.81}{\Character{5}}%
\TextAt{21}{149.196}{290.149}{i}%
\TextAt{8}{156.842}{285.81}{and}%
\TextAt{8}{179.972}{285.81}{show}%
\TextAt{8}{209.342}{285.81}{that}%
\TextAt{13}{235.078}{285.81}{M}%
\TextAt{9}{251.598}{285.81}{\Character{61}}%
\TextAt{16}{264.649}{285.81}{\Character{79}}%
\TextAt{8}{274.498}{285.81}{,}%
\TextAt{8}{282.106}{285.81}{i.e.,}%
\TextAt{16}{304.67}{285.81}{\Character{102}}%
\TextAt{13}{310.648}{285.81}{\Character{33}}%
\TextAt{21}{317.937}{284.017}{j}%
\TextAt{13}{322.319}{285.81}{\Character{25}}%
\TextAt{21}{328.96}{284.017}{i}%
\TextAt{16}{332.341}{285.81}{\Character{103}}%
\TextAt{8}{342.583}{285.81}{is}%
\TextAt{8}{354.714}{285.81}{even}%
\TextAt{8}{381.741}{285.81}{an}%
\TextAt{8}{398.367}{285.81}{integral}%
\TextAt{8}{441.342}{285.81}{basis}%
\TextAt{8}{470.449}{285.81}{of}%
\TextAt{16}{484.146}{285.81}{\Character{79}}%
\TextAt{8}{498.26}{285.81}{over}%
\TextAt{14}{523.674}{285.81}{o}%
\TextAt{8}{529.518}{285.81}{.}%
\TextAt{8}{79.226}{268.764}{We}%
\TextAt{8}{99.371}{268.764}{put}%
\TextAt{13}{272.799}{243.287}{N}%
\TextAt{9}{286.743}{243.287}{\Character{61}}%
\TextAt{21}{305.329}{258.718}{f}%
\TextAt{20}{299.168}{254.645}{\Character{88}}%
\TextAt{21}{300.45}{229.45}{j}%
\TextAt{17}{304.334}{229.45}{\Character{61}1}%
\TextAt{14}{318.429}{243.287}{o}%
\TextAt{13}{324.273}{243.287}{\Character{33}}%
\TextAt{21}{331.563}{241.494}{j}%
\TextAt{13}{335.945}{243.287}{\Character{59}}%
\TextAt{8}{79.226}{211.042}{so}%
\TextAt{8}{93.596}{211.042}{that}%
\TextAt{13}{118.965}{211.042}{M}%
\TextAt{9}{134.86}{211.042}{\Character{61}}%
\TextAt{13}{147.285}{211.042}{N}%
\TextAt{9}{160.564}{211.042}{\Character{43}}%
\TextAt{9}{172.316}{211.042}{\Character{5}}%
\TextAt{13}{181.101}{211.042}{N}%
\TextAt{9}{194.38}{211.042}{\Character{43}}%
\TextAt{16}{206.142}{211.042}{\Character{1}}%
\TextAt{16}{211.438}{211.042}{\Character{1}}%
\TextAt{16}{216.746}{211.042}{\Character{1}}%
\TextAt{9}{222.746}{211.042}{\Character{43}}%
\TextAt{9}{234.498}{211.042}{\Character{5}}%
\TextAt{21}{243.283}{215.38}{e}%
\TextAt{3}{247.211}{215.38}{\Character{0}}%
\TextAt{17}{253.797}{215.38}{1}%
\TextAt{13}{258.53}{211.042}{N}%
\TextAt{8}{269.152}{211.042}{.}%
\TextAt{8}{277.601}{211.042}{We}%
\TextAt{8}{297.745}{211.042}{find}%
\TextAt{8}{321.153}{211.042}{that}%
\TextAt{16}{269.741}{185.035}{\Character{79}}%
\TextAt{9}{282.912}{185.035}{\Character{61}}%
\TextAt{13}{295.337}{185.035}{N}%
\TextAt{9}{308.616}{185.035}{\Character{43}}%
\TextAt{9}{320.368}{185.035}{\Character{5}}%
\TextAt{16}{329.153}{185.035}{\Character{79}}%
\TextAt{13}{339.003}{185.035}{\Character{59}}%
\TextAt{8}{79.226}{159.027}{because,}%
\TextAt{8}{125.78}{159.027}{for}%
\TextAt{13}{143.665}{159.027}{\Character{11}}%
\NativeMathAt{154.507}{159.027}{11.955}{\in}%
\TextAt{16}{165.781}{159.027}{\Character{79}}%
\TextAt{8}{175.648}{159.027}{,}%
\TextAt{8}{182.797}{159.027}{we}%
\TextAt{8}{200.024}{159.027}{have}%
\TextAt{13}{227.019}{159.027}{\Character{11}}%
\NativeMathAt{237.862}{159.027}{11.955}{\equiv}%
\TextAt{13}{250.481}{159.027}{a}%
\TextAt{17}{256.626}{157.234}{1}%
\TextAt{13}{261.358}{159.027}{\Character{33}}%
\TextAt{17}{268.648}{157.234}{1}%
\TextAt{9}{276.037}{159.027}{\Character{43}}%
\TextAt{16}{287.798}{159.027}{\Character{1}}%
\TextAt{16}{293.094}{159.027}{\Character{1}}%
\TextAt{16}{298.402}{159.027}{\Character{1}}%
\TextAt{9}{304.402}{159.027}{\Character{43}}%
\TextAt{13}{316.164}{159.027}{a}%
\TextAt{21}{322.308}{157.234}{f}%
\TextAt{13}{327.753}{159.027}{\Character{33}}%
\TextAt{21}{335.043}{157.234}{f}%
\TextAt{8}{343.808}{159.027}{mod}%
\TextAt{9}{369.246}{159.027}{\Character{5}}%
\TextAt{16}{378.022}{159.027}{\Character{79}}%
\TextAt{8}{387.871}{159.027}{,}%
\TextAt{13}{395.021}{159.027}{a}%
\TextAt{21}{401.166}{157.234}{i}%
\NativeMathAt{407.868}{159.027}{11.955}{\in}%
\TextAt{14}{419.159}{159.027}{o}%
\TextAt{8}{425.003}{159.027}{.}%
\TextAt{8}{433.451}{159.027}{This}%
\TextAt{8}{460.171}{159.027}{implies}%
\TextAt{16}{145.892}{133.02}{\Character{79}}%
\TextAt{9}{159.063}{133.02}{\Character{61}}%
\TextAt{13}{171.488}{133.02}{N}%
\TextAt{9}{184.768}{133.02}{\Character{43}}%
\TextAt{9}{196.52}{133.02}{\Character{5}\Character{40}}%
\TextAt{13}{209.857}{133.02}{N}%
\TextAt{9}{223.136}{133.02}{\Character{43}}%
\TextAt{9}{234.888}{133.02}{\Character{5}}%
\TextAt{16}{243.672}{133.02}{\Character{79}}%
\TextAt{9}{253.522}{133.02}{\Character{41}}%
\TextAt{9}{261.377}{133.02}{\Character{61}}%
\TextAt{16}{273.821}{133.02}{\Character{1}}%
\TextAt{16}{279.117}{133.02}{\Character{1}}%
\TextAt{16}{284.425}{133.02}{\Character{1}}%
\TextAt{9}{291.089}{133.02}{\Character{61}}%
\TextAt{13}{303.515}{133.02}{N}%
\TextAt{9}{316.794}{133.02}{\Character{43}}%
\TextAt{9}{328.546}{133.02}{\Character{5}}%
\TextAt{13}{337.331}{133.02}{N}%
\TextAt{9}{350.61}{133.02}{\Character{43}}%
\TextAt{16}{362.371}{133.02}{\Character{1}}%
\TextAt{16}{367.667}{133.02}{\Character{1}}%
\TextAt{16}{372.975}{133.02}{\Character{1}}%
\TextAt{9}{378.976}{133.02}{\Character{43}}%
\TextAt{9}{390.728}{133.02}{\Character{5}}%
\TextAt{21}{399.512}{137.957}{e}%
\TextAt{3}{403.441}{137.957}{\Character{0}}%
\TextAt{17}{410.027}{137.957}{1}%
\TextAt{13}{414.76}{133.02}{N}%
\TextAt{9}{428.039}{133.02}{\Character{43}}%
\TextAt{9}{439.791}{133.02}{\Character{5}}%
\TextAt{21}{448.576}{137.957}{e}%
\TextAt{16}{453.002}{133.02}{\Character{79}}%
\TextAt{13}{462.852}{133.02}{\Character{59}}%
\TextAt{8}{79.226}{107.013}{so}%
\TextAt{8}{94.057}{107.013}{that}%
\TextAt{16}{119.887}{107.013}{\Character{79}}%
\TextAt{9}{133.843}{107.013}{\Character{61}}%
\TextAt{13}{147.053}{107.013}{M}%
\TextAt{9}{162.598}{107.013}{\Character{43}}%
\TextAt{14}{174.673}{107.013}{P}%
\TextAt{21}{184.567}{111.352}{e}%
\TextAt{9}{193.099}{107.013}{\Character{61}}%
\TextAt{13}{206.31}{107.013}{M}%
\TextAt{9}{221.854}{107.013}{\Character{43}}%
\TextAt{14}{233.93}{107.013}{p}%
\TextAt{16}{239.913}{107.013}{\Character{79}}%
\TextAt{8}{249.763}{107.013}{.}%
\TextAt{8}{259.593}{107.013}{Since}%
\TextAt{13}{290.612}{107.013}{L}%
\TextAt{16}{302.682}{107.013}{\Character{106}}%
\TextAt{13}{310.109}{107.013}{K}%
\TextAt{8}{325.285}{107.013}{is}%
\TextAt{8}{337.51}{107.013}{separable,}%
\TextAt{16}{392.781}{107.013}{\Character{79}}%
\TextAt{8}{406.989}{107.013}{is}%
\TextAt{8}{419.214}{107.013}{a}%
\TextAt{8}{429.43}{107.013}{finitely}%
\TextAt{8}{469.235}{107.013}{generated}%
\TextAt{14}{523.028}{107.013}{o}%
\TextAt{8}{528.872}{107.013}{-}%
\TextAt{8}{79.226}{89.967}{module}%
\TextAt{8}{120.578}{89.967}{(see}%
\TextAt{8}{144.105}{89.967}{chap.}%
\TextAt{8}{175.057}{89.967}{I,}%
\ErrataAt{186.48}{89.967}{38.52}{\hyperlink{nav.number.1.2.9}{(2.9)}),}%
\TextAt{8}{228.199}{89.967}{and}%
\TextAt{8}{251.02}{89.967}{we}%
\TextAt{8}{268.307}{89.967}{conclude}%
\TextAt{16}{316.485}{89.967}{\Character{79}}%
\TextAt{9}{329.757}{89.967}{\Character{61}}%
\TextAt{13}{342.283}{89.967}{M}%
\TextAt{8}{358.813}{89.967}{from}%
\TextAt{8}{386.512}{89.967}{Nakayama’s}%
\TextAt{8}{452.014}{89.967}{lemma}%
\TextAt{8}{489.791}{89.967}{(chap.}%
\TextAt{8}{525.298}{89.967}{I,}%
\TextAt{8}{79.226}{72.921}{§11,}%
\TextAt{8}{103.77}{72.921}{exercise}%
\TextAt{8}{147.072}{72.921}{7).}%
\Vector{q 0 G 0 g 0.399 w 452.334 -701.079 m 452.334 -693.009 l S Q}%
\Vector{q 0 G 0 g 0.399 w 452.534 -693.208 m 459.707 -693.208 l S Q}%
\Vector{q 0 G 0 g 0.399 w 452.534 -700.879 m 459.707 -700.879 l S Q}%
\Vector{q 0 G 0 g 0.399 w 459.906 -701.079 m 459.906 -693.009 l S Q}%
\TextAt{8}{297.211}{43.033}{138}%
\Destination{78.23}{809.815}{nav.page.138}%
\Destination{79.36}{565.855}{nav.print.151}%
\Hotlink{174.557}{87.755}{5.22}{11.047}{nav.chapter.1}%
\Hotlink{524.798}{87.755}{5.22}{11.047}{nav.chapter.1}%
\Destination{79.226}{106.967}{errata.84}%
\Hotlink{79.655}{70.7}{16.008}{11.047}{nav.section.1.11}%
\Hotlink{147.715}{70.7}{5.033}{11.047}{nav.ex.1.11.7}%
}

% Page 151 (body page 139)
\Page{%
\TextAt{2}{79.226}{796.943}{II.}%
\TextAt{2}{94.073}{796.943}{The}%
\TextAt{2}{116.492}{796.943}{Theory}%
\TextAt{2}{154.401}{796.943}{of}%
\TextAt{2}{166.826}{796.943}{Valuations}%
\Vector{q 0 G 0 g 0.299 w 7.226 18.729 m 460.769 18.729 l S Q}%
\TextAt{12}{79.226}{761.124}{Remark.}%
\TextAt{8}{132.51}{761.124}{We}%
\TextAt{8}{151.78}{761.124}{had}%
\TextAt{8}{173.67}{761.124}{already}%
\TextAt{8}{214.101}{761.124}{proved}%
\TextAt{8}{251.281}{761.124}{the}%
\TextAt{8}{270.564}{761.124}{identity}%
\TextAt{9}{312.944}{761.124}{\Character{91}}%
\TextAt{13}{316.196}{761.124}{L}%
\TextAt{9}{327.481}{761.124}{\Character{58}}%
\TextAt{13}{334.053}{761.124}{K}%
\TextAt{9}{344.871}{761.124}{\Character{93}}%
\TextAt{9}{351.422}{761.124}{\Character{61}}%
\TextAt{13}{363.869}{761.124}{ef}%
\TextAt{8}{379.365}{761.124}{in}%
\TextAt{8}{392.144}{761.124}{a}%
\TextAt{8}{401.026}{761.124}{somewhat}%
\TextAt{8}{454.847}{761.124}{different}%
\TextAt{8}{500.145}{761.124}{way}%
\TextAt{8}{523.014}{761.124}{in}%
\TextAt{8}{79.226}{744.078}{chap.}%
\TextAt{8}{109.587}{744.078}{I,}%
\TextAt{8}{120.425}{744.078}{(8.2),}%
\TextAt{8}{151.226}{744.078}{also}%
\TextAt{8}{174.175}{744.078}{in}%
\TextAt{8}{187.296}{744.078}{the}%
\TextAt{8}{206.921}{744.078}{case}%
\TextAt{8}{231.161}{744.078}{where}%
\TextAt{13}{264.439}{744.078}{v}%
\TextAt{8}{273.893}{744.078}{was}%
\TextAt{8}{295.861}{744.078}{discrete}%
\TextAt{8}{338.308}{744.078}{and}%
\TextAt{13}{360.54}{744.078}{L}%
\TextAt{16}{371.825}{744.078}{\Character{106}}%
\TextAt{13}{378.467}{744.078}{K}%
\TextAt{8}{392.65}{744.078}{separable.}%
\TextAt{8}{448.467}{744.078}{Both}%
\TextAt{8}{477.034}{744.078}{hypotheses}%
\TextAt{8}{79.226}{727.032}{are}%
\TextAt{8}{99.577}{727.032}{actually}%
\TextAt{8}{144.974}{727.032}{needed.}%
\TextAt{8}{191.054}{727.032}{But,}%
\TextAt{8}{218.596}{727.032}{strangely}%
\TextAt{8}{269.911}{727.032}{enough,}%
\TextAt{8}{314.537}{727.032}{the}%
\TextAt{8}{335.533}{727.032}{separability}%
\TextAt{8}{399.532}{727.032}{condition}%
\TextAt{8}{451.755}{727.032}{can}%
\TextAt{8}{474.055}{727.032}{be}%
\TextAt{8}{490.819}{727.032}{dropped}%
\TextAt{8}{79.226}{709.986}{once}%
\TextAt{13}{106.141}{709.986}{K}%
\TextAt{8}{121.111}{709.986}{is}%
\TextAt{8}{133.129}{709.986}{complete}%
\TextAt{8}{182.807}{709.986}{with}%
\TextAt{8}{209.721}{709.986}{respect}%
\TextAt{8}{250.026}{709.986}{to}%
\TextAt{8}{264.591}{709.986}{the}%
\TextAt{8}{285.002}{709.986}{discrete}%
\TextAt{8}{328.235}{709.986}{valuation.}%
\TextAt{8}{384.623}{709.986}{In}%
\TextAt{8}{399.499}{709.986}{this}%
\TextAt{8}{422.576}{709.986}{case,}%
\TextAt{8}{450.917}{709.986}{one}%
\TextAt{8}{472.631}{709.986}{deduces}%
\TextAt{8}{516.51}{709.986}{the}%
\TextAt{8}{79.226}{692.94}{equality}%
\TextAt{16}{124.683}{692.94}{\Character{79}}%
\TextAt{9}{139.386}{692.94}{\Character{61}}%
\TextAt{13}{153.345}{692.94}{M}%
\TextAt{8}{170.715}{692.94}{in}%
\TextAt{8}{185.267}{692.94}{the}%
\TextAt{8}{206.323}{692.94}{above}%
\TextAt{8}{240.398}{692.94}{proof}%
\TextAt{8}{271.867}{692.94}{from}%
\TextAt{16}{300.407}{692.94}{\Character{79}}%
\TextAt{9}{315.111}{692.94}{\Character{61}}%
\TextAt{13}{329.069}{692.94}{M}%
\TextAt{9}{344.913}{692.94}{\Character{43}}%
\TextAt{14}{357.287}{692.94}{p}%
\TextAt{16}{363.27}{692.94}{\Character{79}}%
\TextAt{8}{373.12}{692.94}{,}%
\TextAt{8}{381.394}{692.94}{not}%
\TextAt{8}{403.107}{692.94}{by}%
\TextAt{8}{420.266}{692.94}{means}%
\TextAt{8}{456.995}{692.94}{of}%
\TextAt{8}{471.224}{692.94}{Nakayama’s}%
\TextAt{8}{79.226}{675.894}{lemma,}%
\TextAt{8}{120.197}{675.894}{but}%
\TextAt{8}{141.656}{675.894}{rather}%
\TextAt{8}{176.78}{675.894}{like}%
\TextAt{8}{198.24}{675.894}{this:}%
\TextAt{8}{225.613}{675.894}{as}%
\TextAt{14}{239.983}{675.894}{p}%
\TextAt{21}{245.966}{680.232}{i}%
\TextAt{13}{249.348}{675.894}{M}%
\NativeMathAt{265.242}{675.894}{11.955}{\subseteq}%
\TextAt{13}{277.861}{675.894}{M}%
\TextAt{8}{290.435}{675.894}{,}%
\TextAt{8}{297.584}{675.894}{we}%
\TextAt{8}{314.812}{675.894}{get}%
\TextAt{8}{334.322}{675.894}{successively}%
\TextAt{16}{172.373}{651.467}{\Character{79}}%
\TextAt{9}{185.544}{651.467}{\Character{61}}%
\TextAt{13}{197.969}{651.467}{M}%
\TextAt{9}{213.199}{651.467}{\Character{43}}%
\TextAt{14}{224.961}{651.467}{p}%
\TextAt{9}{230.944}{651.467}{\Character{40}}%
\TextAt{13}{235.496}{651.467}{M}%
\TextAt{9}{250.726}{651.467}{\Character{43}}%
\TextAt{14}{262.488}{651.467}{p}%
\TextAt{16}{268.471}{651.467}{\Character{79}}%
\TextAt{9}{278.321}{651.467}{\Character{41}}%
\TextAt{9}{286.176}{651.467}{\Character{61}}%
\TextAt{13}{298.619}{651.467}{M}%
\TextAt{9}{313.849}{651.467}{\Character{43}}%
\TextAt{14}{325.611}{651.467}{p}%
\TextAt{17}{331.594}{656.403}{2}%
\TextAt{16}{336.326}{651.467}{\Character{79}}%
\TextAt{9}{349.497}{651.467}{\Character{61}}%
\TextAt{16}{361.922}{651.467}{\Character{1}}%
\TextAt{16}{367.218}{651.467}{\Character{1}}%
\TextAt{16}{372.526}{651.467}{\Character{1}}%
\TextAt{9}{379.191}{651.467}{\Character{61}}%
\TextAt{13}{391.616}{651.467}{M}%
\TextAt{9}{406.847}{651.467}{\Character{43}}%
\TextAt{14}{418.608}{651.467}{p}%
\TextAt{21}{424.591}{656.403}{\Character{23}}%
\TextAt{16}{429.773}{651.467}{\Character{79}}%
\TextAt{8}{79.226}{627.04}{for}%
\TextAt{8}{97.392}{627.04}{all}%
\TextAt{13}{113.932}{627.04}{\Character{23}}%
\NativeMathAt{124.27}{627.04}{11.955}{\geq}%
\TextAt{9}{137.369}{627.04}{1}%
\TextAt{8}{143.222}{627.04}{,}%
\TextAt{8}{150.722}{627.04}{and}%
\TextAt{8}{173.766}{627.04}{since}%
\TextAt{16}{202.715}{627.04}{\Character{102}}%
\TextAt{14}{208.693}{627.04}{p}%
\TextAt{21}{214.676}{631.379}{\Character{23}}%
\TextAt{16}{219.858}{627.04}{\Character{79}}%
\TextAt{16}{229.697}{627.04}{\Character{103}}%
\TextAt{21}{235.686}{625.247}{\Character{23}}%
\NativeMathAt{240.369}{625.247}{7.97}{\in}%
\TextAt{25}{246.015}{625.247}{N}%
\TextAt{8}{256.448}{627.04}{is}%
\TextAt{8}{268.493}{627.04}{a}%
\TextAt{8}{278.529}{627.04}{basis}%
\TextAt{8}{307.551}{627.04}{of}%
\TextAt{8}{321.162}{627.04}{neighbourhoods}%
\TextAt{8}{406.062}{627.04}{of}%
\TextAt{8}{419.673}{627.04}{zero}%
\TextAt{8}{444.665}{627.04}{in}%
\TextAt{16}{458.599}{627.04}{\Character{79}}%
\TextAt{8}{468.449}{627.04}{,}%
\TextAt{13}{475.95}{627.04}{M}%
\TextAt{8}{492.702}{627.04}{is}%
\TextAt{8}{504.747}{627.04}{dense}%
\TextAt{8}{79.226}{609.994}{in}%
\TextAt{16}{92.879}{609.994}{\Character{79}}%
\TextAt{8}{102.729}{609.994}{.}%
\TextAt{8}{111.177}{609.994}{Since}%
\TextAt{14}{141.734}{609.994}{o}%
\TextAt{8}{151.476}{609.994}{is}%
\TextAt{8}{163.24}{609.994}{closed}%
\TextAt{8}{197.767}{609.994}{in}%
\TextAt{13}{211.419}{609.994}{K}%
\TextAt{8}{222.237}{609.994}{,}%
\TextAt{8}{229.386}{609.994}{(4.9)}%
\TextAt{8}{257.361}{609.994}{implies}%
\TextAt{8}{297.088}{609.994}{that}%
\TextAt{13}{322.457}{609.994}{M}%
\TextAt{8}{338.928}{609.994}{is}%
\TextAt{8}{350.692}{609.994}{closed}%
\TextAt{8}{385.219}{609.994}{in}%
\TextAt{16}{398.871}{609.994}{\Character{79}}%
\TextAt{8}{408.721}{609.994}{,}%
\TextAt{8}{415.87}{609.994}{so}%
\TextAt{8}{430.241}{609.994}{that}%
\TextAt{13}{455.609}{609.994}{M}%
\TextAt{9}{471.504}{609.994}{\Character{61}}%
\TextAt{16}{483.929}{609.994}{\Character{79}}%
\TextAt{8}{493.779}{609.994}{.}%
\TextAt{12}{79.226}{582.954}{Exercise}%
\TextAt{12}{132.215}{582.954}{II.6.1.}%
\TextAt{8}{171.006}{582.954}{In}%
\TextAt{8}{185.443}{582.954}{a}%
\TextAt{8}{195.015}{582.954}{henselian}%
\TextAt{8}{245.617}{582.954}{field}%
\TextAt{8}{270.79}{582.954}{the}%
\TextAt{8}{290.763}{582.954}{zeroes}%
\TextAt{8}{325.429}{582.954}{of}%
\TextAt{8}{338.575}{582.954}{a}%
\TextAt{8}{348.147}{582.954}{polynomial}%
\TextAt{8}{408.456}{582.954}{are}%
\TextAt{8}{427.783}{582.954}{continuous}%
\TextAt{8}{486.204}{582.954}{functions}%
\TextAt{8}{79.226}{565.908}{of}%
\TextAt{8}{93.604}{565.908}{its}%
\TextAt{8}{110.97}{565.908}{coefficients.}%
\TextAt{8}{177.902}{565.908}{More}%
\TextAt{8}{209.185}{565.908}{precisely,}%
\TextAt{8}{260.622}{565.908}{one}%
\TextAt{8}{283.13}{565.908}{has:}%
\TextAt{8}{310.65}{565.908}{let}%
\TextAt{13}{328.602}{565.908}{f}%
\TextAt{9}{335.649}{565.908}{\Character{40}}%
\TextAt{13}{340.201}{565.908}{x}%
\TextAt{9}{346.853}{565.908}{\Character{41}}%
\NativeMathAt{356.512}{565.908}{11.955}{\in}%
\TextAt{13}{369.588}{565.908}{K}%
\TextAt{9}{380.406}{565.908}{\Character{91}}%
\TextAt{13}{383.657}{565.908}{x}%
\TextAt{9}{390.31}{565.908}{\Character{93}}%
\TextAt{8}{398.506}{565.908}{be}%
\TextAt{8}{415.478}{565.908}{a}%
\TextAt{8}{426.281}{565.908}{monic}%
\TextAt{8}{461.796}{565.908}{polynomial}%
\TextAt{8}{523.337}{565.908}{of}%
\TextAt{8}{79.226}{548.862}{degree}%
\TextAt{13}{115.642}{548.862}{n}%
\TextAt{8}{126.527}{548.862}{and}%
\TextAt{13}{252.912}{525.844}{f}%
\TextAt{9}{259.959}{525.844}{\Character{40}}%
\TextAt{13}{264.511}{525.844}{x}%
\TextAt{9}{271.163}{525.844}{\Character{41}}%
\TextAt{9}{279.018}{525.844}{\Character{61}}%
\TextAt{21}{297.071}{540.788}{r}%
\TextAt{20}{291.462}{537.201}{\Character{89}}%
\TextAt{21}{292.248}{512.007}{i}%
\TextAt{17}{295.131}{512.007}{\Character{61}1}%
\TextAt{9}{306.738}{525.844}{\Character{40}}%
\TextAt{13}{311.29}{525.844}{x}%
\TextAt{16}{320.599}{525.844}{\Character{0}}%
\TextAt{13}{332.554}{525.844}{\Character{11}}%
\TextAt{21}{340}{524.05}{i}%
\TextAt{9}{343.381}{525.844}{\Character{41}}%
\TextAt{21}{347.933}{530.78}{m}%
\TextAt{22}{355.424}{529.565}{i}%
\TextAt{8}{79.226}{496.883}{its}%
\TextAt{8}{97.039}{496.883}{decomposition}%
\TextAt{8}{175.667}{496.883}{into}%
\TextAt{8}{200.904}{496.883}{linear}%
\TextAt{8}{234.916}{496.883}{factors,}%
\TextAt{8}{278.149}{496.883}{with}%
\TextAt{13}{306.302}{496.883}{m}%
\TextAt{21}{316.542}{495.09}{i}%
\NativeMathAt{325.79}{496.883}{11.955}{\geq}%
\TextAt{9}{340.955}{496.883}{1}%
\TextAt{8}{346.808}{496.883}{,}%
\TextAt{13}{355.824}{496.883}{\Character{11}}%
\TextAt{21}{363.27}{495.09}{i}%
\NativeMathAt{372.518}{496.883}{11.955}{\neq}%
\TextAt{13}{387.489}{496.883}{\Character{11}}%
\TextAt{21}{394.935}{495.09}{j}%
\TextAt{8}{404.708}{496.883}{for}%
\TextAt{13}{424.087}{496.883}{i}%
\NativeMathAt{433.947}{496.883}{11.955}{\neq}%
\TextAt{13}{448.918}{496.883}{j}%
\TextAt{8}{454.434}{496.883}{.}%
\TextAt{8}{467.364}{496.883}{If}%
\TextAt{8}{480.55}{496.883}{the}%
\TextAt{8}{502.2}{496.883}{monic}%
\TextAt{8}{79.226}{479.837}{polynomial}%
\TextAt{13}{139.789}{479.837}{g}%
\TextAt{9}{145.823}{479.837}{\Character{40}}%
\TextAt{13}{150.375}{479.837}{x}%
\TextAt{9}{157.027}{479.837}{\Character{41}}%
\TextAt{8}{165.546}{479.837}{of}%
\TextAt{8}{178.946}{479.837}{degree}%
\TextAt{13}{215.43}{479.837}{n}%
\TextAt{8}{226.385}{479.837}{has}%
\TextAt{8}{247.328}{479.837}{all}%
\TextAt{8}{263.656}{479.837}{coefficients}%
\TextAt{8}{322.963}{479.837}{sufficiently}%
\TextAt{8}{381.625}{479.837}{close}%
\TextAt{8}{409.717}{479.837}{to}%
\TextAt{8}{424.097}{479.837}{those}%
\TextAt{8}{454.795}{479.837}{of}%
\TextAt{13}{468.194}{479.837}{f}%
\TextAt{9}{475.241}{479.837}{\Character{40}}%
\TextAt{13}{479.793}{479.837}{x}%
\TextAt{9}{486.445}{479.837}{\Character{41}}%
\TextAt{8}{490.998}{479.837}{,}%
\TextAt{8}{498.233}{479.837}{then}%
\TextAt{8}{524.963}{479.837}{it}%
\TextAt{8}{79.226}{462.791}{has}%
\TextAt{13}{99.609}{462.791}{r}%
\TextAt{8}{108.616}{462.791}{roots}%
\TextAt{13}{137.786}{462.791}{\Character{12}}%
\TextAt{17}{144.394}{460.998}{1}%
\TextAt{13}{149.127}{462.791}{\Character{59}}%
\TextAt{13}{154.351}{462.791}{\Character{58}}%
\TextAt{13}{159.588}{462.791}{\Character{58}}%
\TextAt{13}{164.812}{462.791}{\Character{58}}%
\TextAt{13}{170.049}{462.791}{\Character{59}}%
\TextAt{13}{175.285}{462.791}{\Character{12}}%
\TextAt{21}{181.956}{460.998}{r}%
\TextAt{8}{189.917}{462.791}{which}%
\TextAt{8}{222.913}{462.791}{approximate}%
\TextAt{8}{290.076}{462.791}{the}%
\TextAt{13}{309.742}{462.791}{\Character{11}}%
\TextAt{17}{317.187}{460.998}{1}%
\TextAt{13}{321.919}{462.791}{\Character{59}}%
\TextAt{13}{327.143}{462.791}{\Character{58}}%
\TextAt{13}{332.38}{462.791}{\Character{58}}%
\TextAt{13}{337.616}{462.791}{\Character{58}}%
\TextAt{13}{342.841}{462.791}{\Character{59}}%
\TextAt{13}{348.077}{462.791}{\Character{11}}%
\TextAt{21}{355.586}{460.998}{r}%
\TextAt{8}{363.547}{462.791}{to}%
\TextAt{8}{377.366}{462.791}{any}%
\TextAt{8}{398.993}{462.791}{previously}%
\TextAt{8}{454.5}{462.791}{given}%
\TextAt{8}{484.578}{462.791}{precision.}%
\TextAt{12}{79.226}{433.415}{Exercise}%
\TextAt{12}{133.238}{433.415}{II.6.2.}%
\TextAt{12}{172.918}{433.415}{(Krasner’s}%
\TextAt{12}{239.483}{433.415}{Lemma).}%
\TextAt{8}{298.957}{433.415}{Let}%
\TextAt{13}{320.631}{433.415}{\Character{11}}%
\NativeMathAt{332.676}{433.415}{11.955}{\in}%
\Vector{q 0 G 0 g 0.478 w 273.17 -330.742 m 283.988 -330.742 l S Q}%
\TextAt{13}{345.17}{433.415}{K}%
\TextAt{8}{360.591}{433.415}{be}%
\TextAt{8}{377.22}{433.415}{separable}%
\TextAt{8}{429.369}{433.415}{over}%
\TextAt{13}{455.121}{433.415}{K}%
\TextAt{8}{470.541}{433.415}{and}%
\TextAt{8}{494.01}{433.415}{let}%
\TextAt{13}{511.62}{433.415}{\Character{11}}%
\TextAt{9}{523.665}{433.415}{\Character{61}}%
\TextAt{13}{79.226}{416.369}{\Character{11}}%
\TextAt{17}{86.672}{414.576}{1}%
\TextAt{13}{91.404}{416.369}{\Character{59}}%
\TextAt{13}{96.628}{416.369}{\Character{58}}%
\TextAt{13}{101.865}{416.369}{\Character{58}}%
\TextAt{13}{107.089}{416.369}{\Character{58}}%
\TextAt{13}{112.326}{416.369}{\Character{59}}%
\TextAt{13}{117.562}{416.369}{\Character{11}}%
\TextAt{21}{125.07}{414.576}{n}%
\TextAt{8}{134.604}{416.369}{be}%
\TextAt{8}{150.528}{416.369}{its}%
\TextAt{8}{166.847}{416.369}{conjugates}%
\TextAt{8}{224.471}{416.369}{over}%
\TextAt{13}{249.517}{416.369}{K}%
\TextAt{8}{260.335}{416.369}{.}%
\TextAt{8}{268.783}{416.369}{If}%
\TextAt{13}{280.475}{416.369}{\Character{12}}%
\NativeMathAt{291.068}{416.369}{11.955}{\in}%
\Vector{q 0 G 0 g 0.478 w 230.359 -347.788 m 241.176 -347.788 l S Q}%
\TextAt{13}{302.359}{416.369}{K}%
\TextAt{8}{317.074}{416.369}{is}%
\TextAt{8}{328.837}{416.369}{such}%
\TextAt{8}{355.234}{416.369}{that}%
\TextAt{16}{207.38}{391.943}{\Character{106}}%
\TextAt{13}{210.701}{391.943}{\Character{11}}%
\TextAt{16}{220.879}{391.943}{\Character{0}}%
\TextAt{13}{232.834}{391.943}{\Character{12}}%
\TextAt{16}{240.106}{391.943}{\Character{106}}%
\TextAt{13}{246.747}{391.943}{\Character{60}}%
\TextAt{16}{259.173}{391.943}{\Character{106}}%
\TextAt{13}{262.494}{391.943}{\Character{11}}%
\TextAt{16}{272.672}{391.943}{\Character{0}}%
\TextAt{13}{284.627}{391.943}{\Character{11}}%
\TextAt{21}{292.073}{390.149}{i}%
\TextAt{16}{295.454}{391.943}{\Character{106}}%
\TextAt{8}{324.678}{391.943}{for}%
\TextAt{13}{342.563}{391.943}{i}%
\TextAt{9}{349.877}{391.943}{\Character{61}}%
\TextAt{9}{362.286}{391.943}{2}%
\TextAt{13}{368.156}{391.943}{\Character{59}}%
\TextAt{13}{373.38}{391.943}{\Character{58}}%
\TextAt{13}{378.617}{391.943}{\Character{58}}%
\TextAt{13}{383.841}{391.943}{\Character{58}}%
\TextAt{13}{389.078}{391.943}{\Character{59}}%
\TextAt{13}{394.314}{391.943}{n\Character{59}}%
\TextAt{8}{79.226}{367.516}{then}%
\TextAt{8}{105.886}{367.516}{one}%
\TextAt{8}{127.346}{367.516}{has}%
\TextAt{13}{148.219}{367.516}{K}%
\TextAt{9}{159.037}{367.516}{\Character{40}}%
\TextAt{13}{163.589}{367.516}{\Character{11}}%
\TextAt{9}{171.111}{367.516}{\Character{41}}%
\NativeMathAt{178.984}{367.516}{11.955}{\subseteq}%
\TextAt{13}{191.604}{367.516}{K}%
\TextAt{9}{202.421}{367.516}{\Character{40}}%
\TextAt{13}{206.974}{367.516}{\Character{12}}%
\TextAt{9}{214.245}{367.516}{\Character{41}}%
\TextAt{8}{218.797}{367.516}{.}%
\TextAt{12}{79.226}{338.14}{Exercise}%
\TextAt{12}{132.01}{338.14}{II.6.3.}%
\TextAt{8}{170.623}{338.14}{A}%
\TextAt{8}{182.933}{338.14}{field}%
\TextAt{8}{207.928}{338.14}{which}%
\TextAt{8}{241.053}{338.14}{is}%
\TextAt{8}{252.455}{338.14}{henselian}%
\TextAt{8}{302.878}{338.14}{with}%
\TextAt{8}{329.177}{338.14}{respect}%
\TextAt{8}{368.865}{338.14}{to}%
\TextAt{8}{382.813}{338.14}{two}%
\TextAt{8}{404.569}{338.14}{inequivalent}%
\TextAt{8}{469.578}{338.14}{valuations}%
\TextAt{8}{524.903}{338.14}{is}%
\TextAt{8}{79.226}{321.094}{separably}%
\TextAt{8}{131.65}{321.094}{closed}%
\TextAt{8}{166.176}{321.094}{(Theorem}%
\TextAt{8}{220.154}{321.094}{of}%
\TextAt{26}{233.998}{321.094}{F.K.}%
\TextAt{26}{264.746}{321.094}{Schmidt}%
\TextAt{8}{313.512}{321.094}{).}%
\TextAt{12}{79.226}{291.718}{Exercise}%
\TextAt{12}{134.422}{291.718}{II.6.4.}%
\TextAt{8}{175.131}{291.718}{A}%
\TextAt{8}{189.537}{291.718}{separably}%
\TextAt{8}{243.695}{291.718}{closed}%
\TextAt{8}{279.956}{291.718}{field}%
\TextAt{13}{307.048}{291.718}{K}%
\TextAt{8}{323.498}{291.718}{is}%
\TextAt{8}{336.997}{291.718}{henselian}%
\TextAt{8}{389.517}{291.718}{with}%
\TextAt{8}{417.912}{291.718}{respect}%
\TextAt{8}{459.697}{291.718}{to}%
\TextAt{8}{475.742}{291.718}{any}%
\TextAt{8}{499.594}{291.718}{nonar-}%
\TextAt{8}{79.226}{274.672}{chimedean}%
\TextAt{8}{136.778}{274.672}{valuation.}%
\TextAt{8}{97.159}{254.637}{More}%
\TextAt{8}{127.608}{254.637}{generally,}%
\TextAt{8}{179.906}{254.637}{every}%
\TextAt{8}{211.013}{254.637}{valuation}%
\TextAt{8}{262.3}{254.637}{of}%
\TextAt{13}{275.845}{254.637}{K}%
\TextAt{8}{290.775}{254.637}{admits}%
\TextAt{8}{329.426}{254.637}{a}%
\TextAt{8}{339.396}{254.637}{unique}%
\TextAt{8}{377.652}{254.637}{extension}%
\TextAt{8}{429.633}{254.637}{to}%
\TextAt{8}{444.158}{254.637}{any}%
\TextAt{8}{466.49}{254.637}{purely}%
\TextAt{8}{502.798}{254.637}{insep-}%
\TextAt{8}{79.226}{237.591}{arable}%
\TextAt{8}{114.35}{237.591}{extension}%
\TextAt{13}{166.116}{237.591}{L}%
\TextAt{16}{177.402}{237.591}{\Character{106}}%
\TextAt{13}{184.043}{237.591}{K}%
\TextAt{8}{194.861}{237.591}{.}%
\TextAt{12}{97.159}{217.557}{Hint:}%
\TextAt{8}{132.698}{217.557}{If}%
\TextAt{13}{144.39}{217.557}{\Character{11}}%
\TextAt{21}{151.912}{221.895}{p}%
\TextAt{9}{159.993}{217.557}{\Character{61}}%
\TextAt{13}{172.419}{217.557}{a}%
\NativeMathAt{181.885}{217.557}{11.955}{\in}%
\TextAt{13}{193.175}{217.557}{K}%
\TextAt{8}{203.993}{217.557}{,}%
\TextAt{8}{211.142}{217.557}{one}%
\TextAt{8}{232.602}{217.557}{is}%
\TextAt{8}{244.366}{217.557}{forced}%
\TextAt{8}{279.155}{217.557}{to}%
\TextAt{8}{293.466}{217.557}{put}%
\TextAt{13}{314.925}{217.557}{w}%
\TextAt{9}{323.615}{217.557}{\Character{40}}%
\TextAt{13}{328.167}{217.557}{\Character{11}}%
\TextAt{9}{335.689}{217.557}{\Character{41}}%
\TextAt{9}{343.544}{217.557}{\Character{61}}%
\TextAt{17}{357.197}{222.264}{1}%
\Vector{q 0 G 0 g 0.478 w 285.183 -553.455 m 289.446 -553.455 l S Q}%
\TextAt{21}{357.183}{213.434}{p}%
\TextAt{13}{362.641}{217.557}{v}%
\TextAt{9}{368.729}{217.557}{\Character{40}}%
\TextAt{13}{373.282}{217.557}{a}%
\TextAt{9}{379.427}{217.557}{\Character{41}}%
\TextAt{8}{383.979}{217.557}{.}%
\TextAt{12}{79.226}{188.181}{Exercise}%
\TextAt{12}{133.273}{188.181}{II.6.5.}%
\TextAt{8}{172.983}{188.181}{Let}%
\TextAt{13}{194.688}{188.181}{K}%
\TextAt{8}{210.139}{188.181}{be}%
\TextAt{8}{226.8}{188.181}{a}%
\TextAt{8}{237.291}{188.181}{nonarchimedean}%
\TextAt{8}{324.857}{188.181}{valued}%
\TextAt{8}{362.343}{188.181}{field,}%
\TextAt{14}{391.872}{188.181}{o}%
\TextAt{8}{402.35}{188.181}{the}%
\TextAt{8}{423.242}{188.181}{valuation}%
\TextAt{8}{475.05}{188.181}{ring,}%
\TextAt{8}{503.288}{188.181}{and}%
\TextAt{14}{526.786}{188.181}{p}%
\TextAt{8}{79.226}{171.135}{the}%
\TextAt{8}{98.912}{171.135}{maximal}%
\TextAt{8}{146.251}{171.135}{ideal.}%
\TextAt{13}{178.608}{171.135}{K}%
\TextAt{8}{192.853}{171.135}{is}%
\TextAt{8}{204.146}{171.135}{henselian}%
\TextAt{8}{254.461}{171.135}{if}%
\TextAt{8}{264.715}{171.135}{and}%
\TextAt{8}{287.007}{171.135}{only}%
\TextAt{8}{312.229}{171.135}{if}%
\TextAt{8}{322.482}{171.135}{every}%
\TextAt{8}{352.904}{171.135}{polynomial}%
\TextAt{13}{412.927}{171.135}{f}%
\TextAt{9}{419.973}{171.135}{\Character{40}}%
\TextAt{13}{424.525}{171.135}{x}%
\TextAt{9}{431.177}{171.135}{\Character{41}}%
\TextAt{9}{439.032}{171.135}{\Character{61}}%
\TextAt{13}{451.476}{171.135}{x}%
\TextAt{21}{458.128}{175.473}{n}%
\TextAt{9}{465.459}{171.135}{\Character{43}}%
\TextAt{13}{476.259}{171.135}{a}%
\TextAt{21}{482.404}{169.341}{n}%
\TextAt{3}{487.542}{169.341}{\Character{0}}%
\TextAt{17}{494.129}{169.341}{1}%
\TextAt{13}{498.861}{171.135}{x}%
\TextAt{21}{505.513}{175.473}{n}%
\TextAt{3}{510.651}{175.473}{\Character{0}}%
\TextAt{17}{517.238}{175.473}{1}%
\TextAt{9}{523.665}{171.135}{\Character{43}}%
\TextAt{16}{79.226}{154.089}{\Character{1}}%
\TextAt{16}{84.522}{154.089}{\Character{1}}%
\TextAt{16}{89.83}{154.089}{\Character{1}}%
\TextAt{9}{95.83}{154.089}{\Character{43}}%
\TextAt{13}{107.592}{154.089}{a}%
\TextAt{17}{113.737}{152.295}{0}%
\NativeMathAt{121.79}{154.089}{11.955}{\in}%
\TextAt{14}{133.081}{154.089}{o}%
\TextAt{9}{138.925}{154.089}{\Character{91}}%
\TextAt{13}{142.177}{154.089}{x}%
\TextAt{9}{148.829}{154.089}{\Character{93}}%
\TextAt{8}{155.978}{154.089}{such}%
\TextAt{8}{182.375}{154.089}{that}%
\TextAt{13}{207.744}{154.089}{a}%
\TextAt{17}{213.889}{152.295}{0}%
\NativeMathAt{221.942}{154.089}{11.955}{\in}%
\TextAt{14}{233.233}{154.089}{p}%
\TextAt{8}{243.113}{154.089}{and}%
\TextAt{13}{265.876}{154.089}{a}%
\TextAt{17}{272.021}{152.295}{1}%
\NativeMathAt{280.074}{154.089}{11.955}{\notin}%
\TextAt{14}{291.365}{154.089}{p}%
\TextAt{8}{301.245}{154.089}{has}%
\TextAt{8}{322.119}{154.089}{a}%
\TextAt{8}{331.875}{154.089}{zero}%
\TextAt{13}{356.586}{154.089}{\Character{11}}%
\NativeMathAt{367.429}{154.089}{11.955}{\in}%
\TextAt{14}{378.719}{154.089}{p}%
\TextAt{8}{384.703}{154.089}{.}%
\TextAt{12}{97.159}{134.054}{Hint:}%
\TextAt{8}{132.698}{134.054}{The}%
\TextAt{8}{156.751}{134.054}{Newton}%
\TextAt{8}{199.993}{134.054}{polygon.}%
\TextAt{12}{79.226}{107.013}{Remark.}%
\TextAt{8}{133.824}{107.013}{A}%
\TextAt{8}{146.936}{107.013}{local}%
\TextAt{8}{175.017}{107.013}{ring}%
\TextAt{14}{199.523}{107.013}{o}%
\TextAt{8}{209.706}{107.013}{with}%
\TextAt{8}{236.806}{107.013}{maximal}%
\TextAt{8}{285.055}{107.013}{ideal}%
\TextAt{14}{313.459}{107.013}{p}%
\TextAt{8}{323.78}{107.013}{is}%
\TextAt{8}{335.984}{107.013}{called}%
\TextAt{12}{369.589}{107.013}{henselian}%
\TextAt{8}{427.94}{107.013}{if}%
\TextAt{8}{439.104}{107.013}{Hensel’s}%
\TextAt{8}{484.855}{107.013}{lemma}%
\TextAt{8}{523.014}{107.013}{in}%
\TextAt{8}{79.226}{89.967}{the}%
\TextAt{8}{100.37}{89.967}{sense}%
\TextAt{8}{131.389}{89.967}{of}%
\TextAt{8}{145.707}{89.967}{(6.7)}%
\TextAt{8}{174.67}{89.967}{holds}%
\TextAt{8}{206.286}{89.967}{for}%
\TextAt{8}{225.159}{89.967}{it.}%
\TextAt{8}{244.377}{89.967}{A}%
\TextAt{8}{258.036}{89.967}{characterization}%
\TextAt{8}{344.563}{89.967}{of}%
\TextAt{8}{358.881}{89.967}{these}%
\TextAt{8}{389.84}{89.967}{rings}%
\TextAt{8}{419.508}{89.967}{which}%
\TextAt{8}{453.982}{89.967}{is}%
\TextAt{8}{466.734}{89.967}{important}%
\TextAt{8}{523.014}{89.967}{in}%
\TextAt{8}{79.226}{72.921}{algebraic}%
\TextAt{8}{128.661}{72.921}{geometry}%
\TextAt{8}{179.721}{72.921}{is}%
\TextAt{8}{191.485}{72.921}{the}%
\TextAt{8}{211.642}{72.921}{following:}%
\TextAt{8}{297.211}{43.033}{139}%
\Destination{78.23}{809.815}{nav.page.139}%
\Destination{79.226}{596.955}{nav.ex.2.6.1}%
\Destination{79.226}{447.455}{nav.ex.2.6.2}%
\Destination{79.226}{352.155}{nav.ex.2.6.3}%
\Destination{79.226}{305.755}{nav.ex.2.6.4}%
\Destination{79.226}{202.155}{nav.ex.2.6.5}%
\Destination{79.36}{597.955}{nav.print.152}%
\Hotlink{109.087}{741.855}{5.22}{11.047}{nav.chapter.1}%
\Hotlink{119.925}{741.855}{25.078}{11.047}{nav.number.1.8.2}%
\Hotlink{228.886}{607.755}{25.078}{11.047}{nav.number.2.4.9}%
\Hotlink{145.207}{87.755}{25.078}{11.047}{nav.number.2.6.7}%
}

% Page 152 (body page 140)
\Page{%
\TextAt{2}{79.226}{796.943}{II.}%
\TextAt{2}{94.073}{796.943}{The}%
\TextAt{2}{116.492}{796.943}{Theory}%
\TextAt{2}{154.401}{796.943}{of}%
\TextAt{2}{166.826}{796.943}{Valuations}%
\Vector{q 0 G 0 g 0.299 w 7.226 18.729 m 460.769 18.729 l S Q}%
\TextAt{8}{97.159}{761.124}{A}%
\TextAt{8}{110.372}{761.124}{local}%
\TextAt{8}{138.554}{761.124}{ring}%
\TextAt{14}{163.161}{761.124}{o}%
\TextAt{8}{173.444}{761.124}{is}%
\TextAt{8}{185.75}{761.124}{henselian}%
\TextAt{8}{237.077}{761.124}{if}%
\TextAt{8}{248.342}{761.124}{and}%
\TextAt{8}{271.646}{761.124}{only}%
\TextAt{8}{297.879}{761.124}{if}%
\TextAt{8}{309.144}{761.124}{every}%
\TextAt{8}{340.578}{761.124}{finite}%
\TextAt{8}{371.031}{761.124}{commutative}%
\TextAt{14}{441.498}{761.124}{o}%
\TextAt{8}{447.343}{761.124}{-algebra}%
\TextAt{13}{492.764}{761.124}{A}%
\TextAt{8}{505.978}{761.124}{splits}%
\TextAt{8}{79.226}{744.078}{into}%
\TextAt{8}{102.969}{744.078}{a}%
\TextAt{8}{112.725}{744.078}{direct}%
\TextAt{8}{145.888}{744.078}{product}%
\TextAt{13}{189.788}{744.078}{A}%
\TextAt{9}{201.884}{744.078}{\Character{61}}%
\TextAt{20}{214.309}{753.044}{\Character{81}}%
\TextAt{21}{225.6}{749.889}{r}%
\TextAt{21}{225.6}{740.591}{i}%
\TextAt{17}{228.483}{740.591}{\Character{61}1}%
\TextAt{13}{241.795}{744.078}{A}%
\TextAt{21}{250.57}{742.284}{i}%
\TextAt{8}{257.849}{744.078}{of}%
\TextAt{8}{271.179}{744.078}{local}%
\TextAt{8}{298.819}{744.078}{rings}%
\TextAt{13}{327.5}{744.078}{A}%
\TextAt{21}{336.275}{742.284}{i}%
\TextAt{8}{339.656}{744.078}{.}%
\TextAt{8}{97.159}{724.043}{The}%
\TextAt{8}{121.213}{724.043}{proof}%
\TextAt{8}{151.782}{724.043}{is}%
\TextAt{8}{163.546}{724.043}{not}%
\TextAt{8}{184.36}{724.043}{straightforward,}%
\TextAt{8}{269.971}{724.043}{we}%
\TextAt{8}{287.198}{724.043}{refer}%
\TextAt{8}{314.181}{724.043}{to}%
\TextAt{8}{328.491}{724.043}{[103],}%
\TextAt{8}{359.718}{724.043}{chap.}%
\TextAt{8}{390.61}{724.043}{I,}%
\TextAt{8}{401.98}{724.043}{§4,}%
\TextAt{8}{420.666}{724.043}{th.}%
\TextAt{8}{438.873}{724.043}{4.2.}%
\TextAt{15}{79.226}{675.084}{§}%
\TextAt{15}{93.785}{675.084}{7}%
\TextAt{15}{120.693}{675.084}{Unramified}%
\TextAt{15}{222.023}{675.084}{and}%
\TextAt{15}{259.415}{675.084}{Tamely}%
\TextAt{15}{327.726}{675.084}{Ramified}%
\TextAt{15}{409.93}{675.084}{Extensions}%
\TextAt{8}{79.226}{643.198}{In}%
\TextAt{8}{93.695}{643.198}{this}%
\TextAt{8}{116.365}{643.198}{section}%
\TextAt{8}{155.294}{643.198}{we}%
\TextAt{8}{172.369}{643.198}{fix}%
\TextAt{8}{188.799}{643.198}{a}%
\TextAt{8}{198.402}{643.198}{base}%
\TextAt{8}{224.324}{643.198}{field}%
\TextAt{13}{249.528}{643.198}{K}%
\TextAt{8}{264.091}{643.198}{which}%
\TextAt{8}{297.425}{643.198}{is}%
\TextAt{8}{309.037}{643.198}{henselian}%
\TextAt{8}{359.67}{643.198}{with}%
\TextAt{8}{386.178}{643.198}{respect}%
\TextAt{8}{426.075}{643.198}{to}%
\TextAt{8}{440.233}{643.198}{a}%
\TextAt{8}{449.836}{643.198}{nonarchimedean}%
\TextAt{8}{79.226}{626.152}{valuation}%
\TextAt{13}{131.161}{626.152}{v}%
\TextAt{8}{142.009}{626.152}{or}%
\TextAt{16}{157.182}{626.152}{\Character{106}}%
\TextAt{16}{164.391}{626.152}{\Character{106}}%
\TextAt{8}{167.721}{626.152}{.}%
\TextAt{8}{178.757}{626.152}{As}%
\TextAt{8}{196.907}{626.152}{before,}%
\TextAt{8}{236.349}{626.152}{we}%
\TextAt{8}{254.439}{626.152}{denote}%
\TextAt{8}{293.02}{626.152}{the}%
\TextAt{8}{314.039}{626.152}{valuation}%
\TextAt{8}{365.974}{626.152}{ring,}%
\TextAt{8}{394.37}{626.152}{the}%
\TextAt{8}{415.388}{626.152}{maximal}%
\TextAt{8}{464.06}{626.152}{ideal}%
\TextAt{8}{492.885}{626.152}{and}%
\TextAt{8}{516.51}{626.152}{the}%
\TextAt{8}{79.226}{609.106}{residue}%
\TextAt{8}{120.089}{609.106}{class}%
\TextAt{8}{148.661}{609.106}{field}%
\TextAt{8}{175.153}{609.106}{by}%
\TextAt{14}{192.548}{609.106}{o}%
\TextAt{13}{198.392}{609.106}{\Character{59}}%
\TextAt{14}{203.636}{609.106}{p}%
\TextAt{13}{209.619}{609.106}{\Character{59}}%
\TextAt{13}{214.843}{609.106}{\Character{20}}%
\TextAt{8}{221.602}{609.106}{,}%
\TextAt{8}{230.17}{609.106}{respectively.}%
\TextAt{8}{300.94}{609.106}{If}%
\TextAt{13}{313.767}{609.106}{L}%
\TextAt{16}{326.988}{609.106}{\Character{106}}%
\TextAt{13}{335.564}{609.106}{K}%
\TextAt{8}{351.414}{609.106}{is}%
\TextAt{8}{364.314}{609.106}{an}%
\TextAt{8}{381.708}{609.106}{algebraic}%
\TextAt{8}{432.278}{609.106}{extension,}%
\TextAt{8}{488.715}{609.106}{then}%
\TextAt{8}{516.51}{609.106}{the}%
\TextAt{8}{79.226}{592.06}{corresponding}%
\TextAt{8}{155.686}{592.06}{invariants}%
\TextAt{8}{210.698}{592.06}{are}%
\TextAt{8}{231.482}{592.06}{labelled}%
\TextAt{13}{275.998}{592.06}{w}%
\TextAt{13}{284.666}{592.06}{\Character{59}}%
\TextAt{16}{289.931}{592.06}{\Character{79}}%
\TextAt{13}{299.781}{592.06}{\Character{59}}%
\TextAt{14}{305.025}{592.06}{P}%
\TextAt{13}{314.919}{592.06}{\Character{59}}%
\TextAt{13}{320.143}{592.06}{\Character{21}}%
\TextAt{8}{326.991}{592.06}{,}%
\TextAt{8}{335.732}{592.06}{respectively.}%
\TextAt{8}{406.917}{592.06}{An}%
\TextAt{8}{427.366}{592.06}{especially}%
\TextAt{8}{481.374}{592.06}{important}%
\TextAt{8}{79.226}{575.014}{role}%
\TextAt{8}{101.975}{575.014}{among}%
\TextAt{8}{139.691}{575.014}{these}%
\TextAt{8}{169.649}{575.014}{extensions}%
\TextAt{8}{226.015}{575.014}{is}%
\TextAt{8}{237.765}{575.014}{played}%
\TextAt{8}{274.502}{575.014}{by}%
\TextAt{8}{290.747}{575.014}{the}%
\TextAt{8}{310.889}{575.014}{unramified}%
\TextAt{8}{369.407}{575.014}{extensions,}%
\TextAt{8}{429.029}{575.014}{which}%
\TextAt{8}{462.501}{575.014}{are}%
\TextAt{8}{481.998}{575.014}{defined}%
\TextAt{8}{522.297}{575.014}{as}%
\TextAt{8}{79.226}{557.968}{follows.}%
\TextAt{12}{79.226}{532.855}{Definition}%
\TextAt{12}{141.761}{532.855}{(7.1).}%
\TextAt{8}{176.856}{532.855}{A}%
\TextAt{8}{189.308}{532.855}{finite}%
\TextAt{8}{218.999}{532.855}{extension}%
\TextAt{13}{270.545}{532.855}{L}%
\TextAt{16}{281.83}{532.855}{\Character{106}}%
\TextAt{13}{288.472}{532.855}{K}%
\TextAt{8}{302.966}{532.855}{is}%
\TextAt{8}{314.51}{532.855}{called}%
\TextAt{12}{347.453}{532.855}{unramified}%
\TextAt{8}{414.146}{532.855}{if}%
\TextAt{8}{424.649}{532.855}{the}%
\TextAt{8}{444.585}{532.855}{extension}%
\TextAt{13}{496.13}{532.855}{\Character{21}}%
\TextAt{16}{506.28}{532.855}{\Character{106}}%
\TextAt{13}{512.921}{532.855}{\Character{20}}%
\TextAt{8}{523.337}{532.855}{of}%
\TextAt{8}{79.226}{515.809}{the}%
\TextAt{8}{99.383}{515.809}{residue}%
\TextAt{8}{139.11}{515.809}{class}%
\TextAt{8}{166.547}{515.809}{field}%
\TextAt{8}{191.904}{515.809}{is}%
\TextAt{8}{203.668}{515.809}{separable}%
\TextAt{8}{255.111}{515.809}{and}%
\TextAt{8}{277.873}{515.809}{one}%
\TextAt{8}{299.333}{515.809}{has}%
\TextAt{9}{263.928}{489.958}{\Character{91}}%
\TextAt{13}{267.179}{489.958}{L}%
\TextAt{9}{278.465}{489.958}{\Character{58}}%
\TextAt{13}{285.037}{489.958}{K}%
\TextAt{9}{295.855}{489.958}{\Character{93}}%
\TextAt{9}{302.406}{489.958}{\Character{61}}%
\TextAt{9}{314.828}{489.958}{\Character{91}}%
\TextAt{13}{318.104}{489.958}{\Character{21}}%
\TextAt{9}{328.254}{489.958}{\Character{58}}%
\TextAt{13}{334.826}{489.958}{\Character{20}}%
\TextAt{9}{341.565}{489.958}{\Character{93}}%
\TextAt{13}{344.816}{489.958}{\Character{58}}%
\TextAt{8}{79.226}{464.107}{An}%
\TextAt{8}{100.211}{464.107}{arbitrary}%
\TextAt{8}{151.789}{464.107}{algebraic}%
\TextAt{8}{203.032}{464.107}{extension}%
\TextAt{13}{256.607}{464.107}{L}%
\TextAt{16}{270.974}{464.107}{\Character{106}}%
\TextAt{13}{280.698}{464.107}{K}%
\TextAt{8}{297.222}{464.107}{is}%
\TextAt{8}{310.794}{464.107}{called}%
\TextAt{8}{345.767}{464.107}{unramified}%
\TextAt{8}{406.108}{464.107}{if}%
\TextAt{8}{418.64}{464.107}{it}%
\TextAt{8}{432.153}{464.107}{is}%
\TextAt{8}{445.726}{464.107}{a}%
\TextAt{8}{457.29}{464.107}{union}%
\TextAt{8}{491.616}{464.107}{of}%
\TextAt{8}{506.755}{464.107}{finite}%
\TextAt{8}{79.226}{447.061}{unramified}%
\TextAt{8}{137.759}{447.061}{subextensions.}%
\TextAt{12}{79.226}{419.41}{Remark.}%
\TextAt{8}{133.049}{419.41}{This}%
\TextAt{8}{159.435}{419.41}{definition}%
\TextAt{8}{211.13}{419.41}{does}%
\TextAt{8}{237.192}{419.41}{not}%
\TextAt{8}{257.672}{419.41}{require}%
\TextAt{13}{296.683}{419.41}{K}%
\TextAt{8}{311.064}{419.41}{to}%
\TextAt{8}{325.04}{419.41}{be}%
\TextAt{8}{340.63}{419.41}{henselian;}%
\TextAt{8}{394.445}{419.41}{it}%
\TextAt{8}{405.815}{419.41}{applies}%
\TextAt{8}{444.562}{419.41}{in}%
\TextAt{8}{457.881}{419.41}{all}%
\TextAt{8}{473.806}{419.41}{cases}%
\TextAt{8}{502.858}{419.41}{where}%
\TextAt{13}{79.226}{402.364}{v}%
\TextAt{8}{89.212}{402.364}{extends}%
\TextAt{8}{131.868}{402.364}{uniquely}%
\TextAt{8}{179.342}{402.364}{to}%
\TextAt{13}{193.652}{402.364}{L}%
\TextAt{8}{201.617}{402.364}{.}%
\TextAt{12}{79.226}{374.713}{Proposition}%
\TextAt{12}{152.079}{374.713}{(7.2).}%
\TextAt{8}{187.517}{374.713}{Let}%
\TextAt{13}{208.609}{374.713}{L}%
\TextAt{16}{220.104}{374.713}{\Character{106}}%
\TextAt{13}{226.955}{374.713}{K}%
\TextAt{8}{241.793}{374.713}{and}%
\TextAt{13}{264.678}{374.713}{K}%
\TextAt{3}{275.496}{379.051}{\Character{48}}%
\TextAt{16}{281.821}{374.713}{\Character{106}}%
\TextAt{13}{288.672}{374.713}{K}%
\TextAt{8}{303.51}{374.713}{be}%
\TextAt{8}{319.557}{374.713}{two}%
\TextAt{8}{341.797}{374.713}{extensions}%
\TextAt{8}{398.3}{374.713}{inside}%
\TextAt{8}{431.646}{374.713}{an}%
\TextAt{8}{448.028}{374.713}{algebraic}%
\TextAt{8}{497.585}{374.713}{closure}%
\Vector{q 0 G 0 g 0.478 w 7.226 -406.49 m 18.044 -406.49 l S Q}%
\TextAt{13}{79.226}{357.667}{K}%
\TextAt{16}{93.365}{357.667}{\Character{106}}%
\TextAt{13}{100.006}{357.667}{K}%
\TextAt{8}{114.721}{357.667}{and}%
\TextAt{8}{137.484}{357.667}{let}%
\TextAt{13}{154.389}{357.667}{L}%
\TextAt{3}{162.353}{362.005}{\Character{48}}%
\TextAt{9}{168.469}{357.667}{\Character{61}}%
\TextAt{13}{180.895}{357.667}{LK}%
\TextAt{3}{199.677}{362.005}{\Character{48}}%
\TextAt{8}{202.472}{357.667}{.}%
\TextAt{8}{210.92}{357.667}{Then}%
\TextAt{8}{241.478}{357.667}{one}%
\TextAt{8}{262.937}{357.667}{has}%
\TextAt{13}{189.49}{331.816}{L}%
\TextAt{16}{200.775}{331.816}{\Character{106}}%
\TextAt{13}{207.417}{331.816}{K}%
\TextAt{8}{222.132}{331.816}{unramified}%
\NativeArrowAt{292.043}{331.816}{11.955}{19.067}{\Longrightarrow}%
\TextAt{13}{326.387}{331.816}{L}%
\TextAt{3}{334.351}{336.752}{\Character{48}}%
\TextAt{16}{340.467}{331.816}{\Character{106}}%
\TextAt{13}{347.109}{331.816}{K}%
\TextAt{3}{357.926}{336.752}{\Character{48}}%
\TextAt{8}{364.619}{331.816}{unramified}%
\TextAt{13}{419.254}{331.816}{\Character{58}}%
\TextAt{8}{79.226}{305.964}{Each}%
\TextAt{8}{108.325}{305.964}{subextension}%
\TextAt{8}{178.036}{305.964}{of}%
\TextAt{8}{191.366}{305.964}{an}%
\TextAt{8}{207.625}{305.964}{unramified}%
\TextAt{8}{266.157}{305.964}{extension}%
\TextAt{8}{317.923}{305.964}{is}%
\TextAt{8}{329.687}{305.964}{unramified.}%
\TextAt{19}{79.226}{272.786}{Proof.}%
\TextAt{8}{116}{272.786}{The}%
\TextAt{8}{141.436}{272.786}{notations}%
\TextAt{14}{194.274}{272.786}{o}%
\TextAt{13}{200.118}{272.786}{\Character{59}}%
\TextAt{14}{205.362}{272.786}{p}%
\TextAt{13}{211.345}{272.786}{\Character{59}}%
\TextAt{13}{216.569}{272.786}{\Character{20}}%
\TextAt{8}{223.328}{272.786}{;}%
\TextAt{14}{232.55}{272.786}{o}%
\TextAt{3}{238.395}{277.124}{\Character{48}}%
\TextAt{13}{241.19}{272.786}{\Character{59}}%
\TextAt{14}{246.434}{272.786}{p}%
\TextAt{3}{252.417}{277.124}{\Character{48}}%
\TextAt{13}{255.212}{272.786}{\Character{59}}%
\TextAt{13}{260.436}{272.786}{\Character{20}}%
\TextAt{3}{267.195}{277.124}{\Character{48}}%
\TextAt{8}{269.99}{272.786}{;}%
\TextAt{16}{279.212}{272.786}{\Character{79}}%
\TextAt{13}{289.062}{272.786}{\Character{59}}%
\TextAt{14}{294.306}{272.786}{P}%
\TextAt{13}{304.2}{272.786}{\Character{59}}%
\TextAt{13}{309.424}{272.786}{\Character{21}}%
\TextAt{8}{316.272}{272.786}{;}%
\TextAt{16}{325.495}{272.786}{\Character{79}}%
\TextAt{3}{335.345}{277.124}{\Character{48}}%
\TextAt{13}{338.14}{272.786}{\Character{59}}%
\TextAt{14}{343.384}{272.786}{P}%
\TextAt{3}{353.277}{277.124}{\Character{48}}%
\TextAt{13}{356.072}{272.786}{\Character{59}}%
\TextAt{13}{361.296}{272.786}{\Character{21}}%
\TextAt{3}{368.145}{277.124}{\Character{48}}%
\TextAt{8}{376.219}{272.786}{are}%
\TextAt{8}{397.113}{272.786}{self-explanatory.}%
\TextAt{8}{489.771}{272.786}{We}%
\TextAt{8}{511.298}{272.786}{may}%
\TextAt{8}{79.226}{255.74}{assume}%
\TextAt{8}{120.378}{255.74}{that}%
\TextAt{13}{146.455}{255.74}{L}%
\TextAt{16}{158.947}{255.74}{\Character{106}}%
\TextAt{13}{166.795}{255.74}{K}%
\TextAt{8}{182.218}{255.74}{is}%
\TextAt{8}{194.69}{255.74}{finite.}%
\TextAt{8}{231.277}{255.74}{Then}%
\TextAt{13}{262.542}{255.74}{\Character{21}}%
\TextAt{16}{273.898}{255.74}{\Character{106}}%
\TextAt{13}{281.746}{255.74}{\Character{20}}%
\TextAt{8}{293.09}{255.74}{is}%
\TextAt{8}{305.562}{255.74}{also}%
\TextAt{8}{329.75}{255.74}{finite}%
\TextAt{8}{360.37}{255.74}{and,}%
\TextAt{8}{387.269}{255.74}{being}%
\TextAt{8}{419.515}{255.74}{separable,}%
\TextAt{8}{475.095}{255.74}{is}%
\TextAt{8}{487.567}{255.74}{therefore}%
\TextAt{8}{79.226}{238.694}{generated}%
\TextAt{8}{133.154}{238.694}{by}%
\TextAt{8}{150.008}{238.694}{a}%
\TextAt{8}{160.359}{238.694}{primitive}%
\TextAt{8}{211.035}{238.694}{element}%
\TextAt{9}{256.032}{238.694}{\Character{22}}%
\TextAt{13}{254.872}{238.694}{\Character{11}}%
\TextAt{8}{262.394}{238.694}{,}%
\TextAt{13}{270.287}{238.694}{\Character{21}}%
\TextAt{9}{281.452}{238.694}{\Character{61}}%
\TextAt{13}{294.892}{238.694}{\Character{20}}%
\TextAt{9}{301.631}{238.694}{\Character{40}}%
\TextAt{9}{307.322}{238.694}{\Character{22}}%
\TextAt{13}{306.183}{238.694}{\Character{11}}%
\TextAt{9}{313.705}{238.694}{\Character{41}}%
\TextAt{8}{318.257}{238.694}{.}%
\TextAt{8}{328.492}{238.694}{Let}%
\TextAt{13}{350.057}{238.694}{\Character{11}}%
\NativeMathAt{361.914}{238.694}{11.955}{\in}%
\TextAt{16}{374.204}{238.694}{\Character{79}}%
\TextAt{8}{388.562}{238.694}{be}%
\TextAt{8}{405.082}{238.694}{a}%
\TextAt{8}{415.433}{238.694}{lifting,}%
\TextAt{13}{453.573}{238.694}{f}%
\TextAt{9}{460.62}{238.694}{\Character{40}}%
\TextAt{13}{465.172}{238.694}{x}%
\TextAt{9}{471.824}{238.694}{\Character{41}}%
\NativeMathAt{480.712}{238.694}{11.955}{\in}%
\TextAt{14}{493.018}{238.694}{o}%
\TextAt{9}{498.862}{238.694}{\Character{91}}%
\TextAt{13}{502.114}{238.694}{x}%
\TextAt{9}{508.766}{238.694}{\Character{93}}%
\TextAt{8}{516.51}{238.694}{the}%
\TextAt{8}{79.226}{221.648}{minimal}%
\TextAt{8}{124.751}{221.648}{polynomial}%
\TextAt{8}{185.245}{221.648}{of}%
\TextAt{13}{198.575}{221.648}{\Character{11}}%
\TextAt{8}{209.994}{221.648}{and}%
\TextAt{9}{235.304}{224.802}{\Character{22}}%
\TextAt{13}{232.756}{221.648}{f}%
\TextAt{9}{239.803}{221.648}{\Character{40}}%
\TextAt{13}{244.355}{221.648}{x}%
\TextAt{9}{251.007}{221.648}{\Character{41}}%
\TextAt{9}{258.862}{221.648}{\Character{61}}%
\TextAt{13}{271.306}{221.648}{f}%
\TextAt{9}{278.352}{221.648}{\Character{40}}%
\TextAt{13}{282.905}{221.648}{x}%
\TextAt{9}{289.557}{221.648}{\Character{41}}%
\TextAt{8}{297.43}{221.648}{mod}%
\TextAt{14}{322.868}{221.648}{p}%
\NativeMathAt{332.172}{221.648}{11.955}{\in}%
\TextAt{13}{343.463}{221.648}{\Character{20}}%
\TextAt{9}{350.201}{221.648}{\Character{91}}%
\TextAt{13}{353.453}{221.648}{x}%
\TextAt{9}{360.105}{221.648}{\Character{93}}%
\TextAt{8}{363.357}{221.648}{.}%
\TextAt{8}{371.805}{221.648}{Since}%
\TextAt{9}{156.217}{195.796}{\Character{91}}%
\TextAt{13}{159.469}{195.796}{\Character{21}}%
\TextAt{9}{169.618}{195.796}{\Character{58}}%
\TextAt{13}{176.19}{195.796}{\Character{20}}%
\TextAt{9}{182.929}{195.796}{\Character{93}}%
\NativeMathAt{189.501}{195.796}{11.955}{\leq}%
\TextAt{8}{202.121}{195.796}{deg}%
\TextAt{9}{219.683}{195.796}{\Character{40}}%
\TextAt{9}{226.783}{198.951}{\Character{22}}%
\TextAt{13}{224.235}{195.796}{f}%
\TextAt{9}{231.282}{195.796}{\Character{41}}%
\TextAt{9}{239.137}{195.796}{\Character{61}}%
\TextAt{8}{251.58}{195.796}{deg}%
\TextAt{9}{269.142}{195.796}{\Character{40}}%
\TextAt{13}{273.695}{195.796}{f}%
\TextAt{9}{280.741}{195.796}{\Character{41}}%
\TextAt{9}{288.596}{195.796}{\Character{61}}%
\TextAt{9}{301.017}{195.796}{\Character{91}}%
\TextAt{13}{304.291}{195.796}{K}%
\TextAt{9}{315.109}{195.796}{\Character{40}}%
\TextAt{13}{319.661}{195.796}{\Character{11}}%
\TextAt{9}{327.183}{195.796}{\Character{41}}%
\TextAt{9}{335.038}{195.796}{\Character{58}}%
\TextAt{13}{341.629}{195.796}{K}%
\TextAt{9}{352.447}{195.796}{\Character{93}}%
\NativeMathAt{359.019}{195.796}{11.955}{\leq}%
\TextAt{9}{371.638}{195.796}{\Character{91}}%
\TextAt{13}{374.89}{195.796}{L}%
\TextAt{9}{386.175}{195.796}{\Character{58}}%
\TextAt{13}{392.748}{195.796}{K}%
\TextAt{9}{403.566}{195.796}{\Character{93}}%
\TextAt{9}{410.117}{195.796}{\Character{61}}%
\TextAt{9}{422.539}{195.796}{\Character{91}}%
\TextAt{13}{425.815}{195.796}{\Character{21}}%
\TextAt{9}{435.964}{195.796}{\Character{58}}%
\TextAt{13}{442.537}{195.796}{\Character{20}}%
\TextAt{9}{449.276}{195.796}{\Character{93}}%
\TextAt{13}{452.527}{195.796}{\Character{59}}%
\TextAt{8}{79.226}{169.945}{one}%
\TextAt{8}{100.686}{169.945}{has}%
\TextAt{13}{121.559}{169.945}{L}%
\TextAt{9}{132.845}{169.945}{\Character{61}}%
\TextAt{13}{145.27}{169.945}{K}%
\TextAt{9}{156.088}{169.945}{\Character{40}}%
\TextAt{13}{160.64}{169.945}{\Character{11}}%
\TextAt{9}{168.162}{169.945}{\Character{41}}%
\TextAt{8}{176.612}{169.945}{and}%
\TextAt{9}{201.922}{173.1}{\Character{22}}%
\TextAt{13}{199.374}{169.945}{f}%
\TextAt{9}{206.421}{169.945}{\Character{40}}%
\TextAt{13}{210.973}{169.945}{x}%
\TextAt{9}{217.625}{169.945}{\Character{41}}%
\TextAt{8}{226.075}{169.945}{is}%
\TextAt{8}{237.839}{169.945}{the}%
\TextAt{8}{257.995}{169.945}{minimal}%
\TextAt{8}{303.52}{169.945}{polynomial}%
\TextAt{8}{364.014}{169.945}{of}%
\TextAt{9}{378.503}{169.945}{\Character{22}}%
\TextAt{13}{377.344}{169.945}{\Character{11}}%
\TextAt{8}{388.763}{169.945}{over}%
\TextAt{13}{413.809}{169.945}{\Character{20}}%
\TextAt{8}{420.547}{169.945}{.}%
\TextAt{8}{97.159}{149.91}{We}%
\TextAt{8}{116.827}{149.91}{thus}%
\TextAt{8}{142.103}{149.91}{have}%
\TextAt{13}{168.622}{149.91}{L}%
\TextAt{3}{176.586}{154.249}{\Character{48}}%
\TextAt{9}{182.702}{149.91}{\Character{61}}%
\TextAt{13}{195.127}{149.91}{K}%
\TextAt{3}{205.945}{154.249}{\Character{48}}%
\TextAt{9}{208.74}{149.91}{\Character{40}}%
\TextAt{13}{213.293}{149.91}{\Character{11}}%
\TextAt{9}{220.814}{149.91}{\Character{41}}%
\TextAt{8}{225.367}{149.91}{.}%
\TextAt{8}{233.656}{149.91}{In}%
\TextAt{8}{247.801}{149.91}{order}%
\TextAt{8}{277.895}{149.91}{to}%
\TextAt{8}{291.729}{149.91}{prove}%
\TextAt{8}{322.803}{149.91}{that}%
\TextAt{13}{347.696}{149.91}{L}%
\TextAt{3}{355.66}{154.249}{\Character{48}}%
\TextAt{16}{361.776}{149.91}{\Character{106}}%
\TextAt{13}{368.418}{149.91}{K}%
\TextAt{3}{379.235}{154.249}{\Character{48}}%
\TextAt{8}{385.452}{149.91}{is}%
\TextAt{8}{396.74}{149.91}{unramified,}%
\TextAt{8}{458.143}{149.91}{let}%
\TextAt{13}{474.572}{149.91}{g}%
\TextAt{9}{480.606}{149.91}{\Character{40}}%
\TextAt{13}{485.159}{149.91}{x}%
\TextAt{9}{491.811}{149.91}{\Character{41}}%
\NativeMathAt{499.684}{149.91}{11.955}{\in}%
\TextAt{14}{510.975}{149.91}{o}%
\TextAt{3}{516.819}{154.249}{\Character{48}}%
\TextAt{9}{519.614}{149.91}{\Character{91}}%
\TextAt{13}{522.866}{149.91}{x}%
\TextAt{9}{529.518}{149.91}{\Character{93}}%
\TextAt{8}{79.226}{132.865}{be}%
\TextAt{8}{94.589}{132.865}{the}%
\TextAt{8}{114.185}{132.865}{minimal}%
\TextAt{8}{159.149}{132.865}{polynomial}%
\TextAt{8}{219.081}{132.865}{of}%
\TextAt{13}{231.85}{132.865}{\Character{11}}%
\TextAt{8}{242.708}{132.865}{over}%
\TextAt{13}{267.192}{132.865}{K}%
\TextAt{3}{278.01}{137.203}{\Character{48}}%
\TextAt{8}{284.141}{132.865}{and}%
\TextAt{9}{306.759}{132.865}{\Character{22}}%
\TextAt{13}{306.343}{132.865}{g}%
\TextAt{9}{312.377}{132.865}{\Character{40}}%
\TextAt{13}{316.929}{132.865}{x}%
\TextAt{9}{323.582}{132.865}{\Character{41}}%
\TextAt{9}{331.437}{132.865}{\Character{61}}%
\TextAt{13}{343.88}{132.865}{g}%
\TextAt{9}{349.914}{132.865}{\Character{40}}%
\TextAt{13}{354.467}{132.865}{x}%
\TextAt{9}{361.119}{132.865}{\Character{41}}%
\TextAt{8}{368.992}{132.865}{mod}%
\TextAt{14}{394.43}{132.865}{p}%
\TextAt{3}{400.413}{137.203}{\Character{48}}%
\NativeMathAt{406.529}{132.865}{11.955}{\in}%
\TextAt{13}{417.82}{132.865}{\Character{20}}%
\TextAt{3}{424.558}{137.203}{\Character{48}}%
\TextAt{9}{427.354}{132.865}{\Character{91}}%
\TextAt{13}{430.605}{132.865}{x}%
\TextAt{9}{437.257}{132.865}{\Character{93}}%
\TextAt{8}{440.509}{132.865}{.}%
\TextAt{8}{448.77}{132.865}{Being}%
\TextAt{8}{481.205}{132.865}{a}%
\TextAt{8}{490.4}{132.865}{factor}%
\TextAt{8}{523.337}{132.865}{of}%
\TextAt{9}{81.774}{118.973}{\Character{22}}%
\TextAt{13}{79.226}{115.819}{f}%
\TextAt{9}{86.273}{115.819}{\Character{40}}%
\TextAt{13}{90.825}{115.819}{x}%
\TextAt{9}{97.477}{115.819}{\Character{41}}%
\TextAt{8}{102.029}{115.819}{,}%
\TextAt{9}{109.669}{115.819}{\Character{22}}%
\TextAt{13}{109.253}{115.819}{g}%
\TextAt{9}{115.288}{115.819}{\Character{40}}%
\TextAt{13}{119.84}{115.819}{x}%
\TextAt{9}{126.492}{115.819}{\Character{41}}%
\TextAt{8}{135.002}{115.819}{is}%
\TextAt{8}{146.826}{115.819}{separable}%
\TextAt{8}{198.329}{115.819}{and}%
\TextAt{8}{221.151}{115.819}{hence}%
\TextAt{8}{253.717}{115.819}{irreducible}%
\TextAt{8}{311.652}{115.819}{over}%
\TextAt{13}{336.758}{115.819}{\Character{20}}%
\TextAt{3}{343.497}{120.157}{\Character{48}}%
\TextAt{8}{346.292}{115.819}{,}%
\TextAt{8}{353.516}{115.819}{because}%
\TextAt{8}{396.878}{115.819}{otherwise}%
\TextAt{13}{449.026}{115.819}{g}%
\TextAt{9}{455.061}{115.819}{\Character{40}}%
\TextAt{13}{459.613}{115.819}{x}%
\TextAt{9}{466.265}{115.819}{\Character{41}}%
\TextAt{8}{474.775}{115.819}{is}%
\TextAt{8}{486.599}{115.819}{reducible}%
\TextAt{8}{79.226}{98.773}{by}%
\TextAt{8}{95.485}{98.773}{Hensel’s}%
\TextAt{8}{140.795}{98.773}{lemma.}%
\TextAt{8}{183.065}{98.773}{We}%
\TextAt{8}{203.209}{98.773}{obtain}%
\TextAt{9}{150.128}{72.921}{\Character{91}}%
\TextAt{13}{153.38}{72.921}{\Character{21}}%
\TextAt{3}{160.208}{77.858}{\Character{48}}%
\TextAt{9}{166.324}{72.921}{\Character{58}}%
\TextAt{13}{172.896}{72.921}{\Character{20}}%
\TextAt{3}{179.635}{77.858}{\Character{48}}%
\TextAt{9}{182.43}{72.921}{\Character{93}}%
\NativeMathAt{189.003}{72.921}{11.955}{\leq}%
\TextAt{9}{201.622}{72.921}{\Character{91}}%
\TextAt{13}{204.874}{72.921}{L}%
\TextAt{3}{212.838}{77.858}{\Character{48}}%
\TextAt{9}{218.954}{72.921}{\Character{58}}%
\TextAt{13}{225.526}{72.921}{K}%
\TextAt{3}{236.344}{77.858}{\Character{48}}%
\TextAt{9}{239.139}{72.921}{\Character{93}}%
\TextAt{9}{245.69}{72.921}{\Character{61}}%
\TextAt{8}{258.137}{72.921}{deg}%
\TextAt{9}{275.699}{72.921}{\Character{40}}%
\TextAt{13}{280.252}{72.921}{g}%
\TextAt{9}{286.286}{72.921}{\Character{41}}%
\TextAt{9}{294.141}{72.921}{\Character{61}}%
\TextAt{8}{306.585}{72.921}{deg}%
\TextAt{9}{324.147}{72.921}{\Character{40}}%
\TextAt{9}{329.096}{72.921}{\Character{22}}%
\TextAt{13}{328.699}{72.921}{g}%
\TextAt{9}{334.733}{72.921}{\Character{41}}%
\TextAt{9}{342.588}{72.921}{\Character{61}}%
\TextAt{9}{355.009}{72.921}{\Character{91}}%
\TextAt{13}{358.284}{72.921}{\Character{20}}%
\TextAt{3}{365.022}{77.858}{\Character{48}}%
\TextAt{9}{367.817}{72.921}{\Character{40}}%
\TextAt{9}{373.52}{72.921}{\Character{22}}%
\TextAt{13}{372.37}{72.921}{\Character{11}}%
\TextAt{9}{379.891}{72.921}{\Character{41}}%
\TextAt{9}{387.746}{72.921}{\Character{58}}%
\TextAt{13}{394.337}{72.921}{\Character{20}}%
\TextAt{3}{401.075}{77.858}{\Character{48}}%
\TextAt{9}{403.871}{72.921}{\Character{93}}%
\NativeMathAt{410.443}{72.921}{11.955}{\leq}%
\TextAt{9}{423.062}{72.921}{\Character{91}}%
\TextAt{13}{426.314}{72.921}{\Character{21}}%
\TextAt{3}{433.142}{77.858}{\Character{48}}%
\TextAt{9}{439.258}{72.921}{\Character{58}}%
\TextAt{13}{445.831}{72.921}{\Character{20}}%
\TextAt{3}{452.569}{77.858}{\Character{48}}%
\TextAt{9}{455.365}{72.921}{\Character{93}}%
\TextAt{13}{458.616}{72.921}{\Character{58}}%
\TextAt{8}{297.211}{43.033}{140}%
\Destination{78.23}{809.815}{nav.page.140}%
\Destination{79.23}{700.945}{nav.section.2.7}%
\Destination{79.36}{547.855}{nav.print.153}%
\Destination{79.36}{546.855}{nav.number.2.7.1}%
\Destination{79.36}{388.755}{nav.number.2.7.2}%
\Destination{79.23}{700.945}{outline.27}%
%
%
\Hotlink{327.991}{721.855}{25.078}{11.047}{nav.cite.N103}%
}

% Page 153 (body page 141)
\Page{%
\TextAt{2}{79.226}{796.943}{II.}%
\TextAt{2}{94.073}{796.943}{The}%
\TextAt{2}{116.492}{796.943}{Theory}%
\TextAt{2}{154.401}{796.943}{of}%
\TextAt{2}{166.826}{796.943}{Valuations}%
\Vector{q 0 G 0 g 0.299 w 7.226 18.729 m 460.769 18.729 l S Q}%
\TextAt{8}{79.226}{761.124}{This}%
\TextAt{8}{105.946}{761.124}{implies}%
\TextAt{9}{145.673}{761.124}{\Character{91}}%
\TextAt{13}{148.925}{761.124}{L}%
\TextAt{3}{156.889}{765.462}{\Character{48}}%
\TextAt{9}{163.005}{761.124}{\Character{58}}%
\TextAt{13}{169.578}{761.124}{K}%
\TextAt{3}{180.395}{765.462}{\Character{48}}%
\TextAt{9}{183.19}{761.124}{\Character{93}}%
\TextAt{9}{189.741}{761.124}{\Character{61}}%
\TextAt{9}{202.163}{761.124}{\Character{91}}%
\TextAt{13}{205.44}{761.124}{\Character{21}}%
\TextAt{3}{212.268}{765.462}{\Character{48}}%
\TextAt{9}{218.384}{761.124}{\Character{58}}%
\TextAt{13}{224.957}{761.124}{\Character{20}}%
\TextAt{3}{231.695}{765.462}{\Character{48}}%
\TextAt{9}{234.49}{761.124}{\Character{93}}%
\TextAt{8}{237.742}{761.124}{,}%
\TextAt{8}{244.891}{761.124}{i.e.,}%
\TextAt{13}{266.996}{761.124}{L}%
\TextAt{3}{274.961}{765.462}{\Character{48}}%
\TextAt{16}{281.077}{761.124}{\Character{106}}%
\TextAt{13}{287.719}{761.124}{K}%
\TextAt{3}{298.536}{765.462}{\Character{48}}%
\TextAt{8}{305.229}{761.124}{is}%
\TextAt{8}{316.993}{761.124}{unramified.}%
\TextAt{8}{97.159}{741.089}{If}%
\TextAt{13}{109.7}{741.089}{L}%
\TextAt{16}{122.431}{741.089}{\Character{106}}%
\TextAt{13}{130.519}{741.089}{K}%
\TextAt{8}{146.083}{741.089}{is}%
\TextAt{8}{158.695}{741.089}{a}%
\TextAt{8}{169.3}{741.089}{subextension}%
\TextAt{8}{239.859}{741.089}{of}%
\TextAt{8}{254.037}{741.089}{the}%
\TextAt{8}{275.043}{741.089}{unramified}%
\TextAt{8}{334.424}{741.089}{extension}%
\TextAt{13}{387.038}{741.089}{L}%
\TextAt{3}{395.003}{745.427}{\Character{48}}%
\TextAt{16}{402.565}{741.089}{\Character{106}}%
\TextAt{13}{410.653}{741.089}{K}%
\TextAt{8}{421.47}{741.089}{,}%
\TextAt{8}{429.681}{741.089}{then}%
\TextAt{8}{457.189}{741.089}{it}%
\TextAt{8}{469.742}{741.089}{follows}%
\TextAt{8}{509.027}{741.089}{from}%
\TextAt{8}{79.226}{724.043}{what}%
\TextAt{8}{108.223}{724.043}{we}%
\TextAt{8}{125.18}{724.043}{have}%
\TextAt{8}{151.905}{724.043}{just}%
\TextAt{8}{174.78}{724.043}{proved}%
\TextAt{8}{212.564}{724.043}{that}%
\TextAt{13}{237.663}{724.043}{L}%
\TextAt{3}{245.628}{728.381}{\Character{48}}%
\TextAt{16}{251.743}{724.043}{\Character{106}}%
\TextAt{13}{258.385}{724.043}{L}%
\TextAt{8}{269.977}{724.043}{is}%
\TextAt{8}{281.471}{724.043}{unramified.}%
\TextAt{8}{344.465}{724.043}{Hence}%
\TextAt{8}{378.973}{724.043}{so}%
\TextAt{8}{393.073}{724.043}{is}%
\TextAt{13}{404.567}{724.043}{L}%
\TextAt{16}{415.852}{724.043}{\Character{106}}%
\TextAt{13}{422.494}{724.043}{K}%
\TextAt{8}{433.312}{724.043}{,}%
\TextAt{8}{440.245}{724.043}{by}%
\TextAt{8}{456.234}{724.043}{the}%
\TextAt{8}{476.121}{724.043}{formula}%
\TextAt{8}{518.782}{724.043}{for}%
\TextAt{8}{79.226}{706.997}{the}%
\TextAt{8}{99.383}{706.997}{degree.}%
\Vector{q 0 G 0 g 0.399 w 452.334 -67.003 m 452.334 -58.933 l S Q}%
\Vector{q 0 G 0 g 0.399 w 452.534 -59.132 m 459.707 -59.132 l S Q}%
\Vector{q 0 G 0 g 0.399 w 452.534 -66.804 m 459.707 -66.804 l S Q}%
\Vector{q 0 G 0 g 0.399 w 459.906 -67.003 m 459.906 -58.933 l S Q}%
\TextAt{12}{79.226}{672.018}{Corollary}%
\TextAt{12}{139.014}{672.018}{(7.3).}%
\TextAt{8}{174.33}{672.018}{The}%
\TextAt{8}{198.383}{672.018}{composite}%
\TextAt{8}{253.401}{672.018}{of}%
\TextAt{8}{266.731}{672.018}{two}%
\TextAt{8}{288.848}{672.018}{unramified}%
\TextAt{8}{347.381}{672.018}{extensions}%
\TextAt{8}{403.761}{672.018}{of}%
\TextAt{13}{417.091}{672.018}{K}%
\TextAt{8}{431.806}{672.018}{is}%
\TextAt{8}{443.57}{672.018}{again}%
\TextAt{8}{474.797}{672.018}{unramified.}%
\TextAt{19}{79.226}{637.04}{Proof.}%
\TextAt{8}{116}{637.04}{It}%
\TextAt{8}{130.099}{637.04}{suffices}%
\TextAt{8}{171.312}{637.04}{to}%
\TextAt{8}{187.049}{637.04}{show}%
\TextAt{8}{217.478}{637.04}{this}%
\TextAt{8}{241.727}{637.04}{for}%
\TextAt{8}{261.038}{637.04}{two}%
\TextAt{8}{284.582}{637.04}{finite}%
\TextAt{8}{315.92}{637.04}{extensions}%
\TextAt{13}{373.727}{637.04}{L}%
\TextAt{16}{387.443}{637.04}{\Character{106}}%
\TextAt{13}{396.515}{637.04}{K}%
\TextAt{8}{412.656}{637.04}{and}%
\TextAt{13}{436.845}{637.04}{L}%
\TextAt{3}{444.81}{641.378}{\Character{48}}%
\TextAt{16}{453.356}{637.04}{\Character{106}}%
\TextAt{13}{462.428}{637.04}{K}%
\TextAt{8}{473.246}{637.04}{.}%
\TextAt{13}{485.973}{637.04}{L}%
\TextAt{16}{499.689}{637.04}{\Character{106}}%
\TextAt{13}{508.762}{637.04}{K}%
\TextAt{8}{524.903}{637.04}{is}%
\TextAt{8}{79.226}{619.994}{unramified,}%
\TextAt{8}{140.888}{619.994}{hence}%
\TextAt{8}{173.241}{619.994}{so}%
\TextAt{8}{187.458}{619.994}{is}%
\TextAt{13}{199.069}{619.994}{LL}%
\TextAt{3}{214.998}{624.332}{\Character{48}}%
\TextAt{16}{221.114}{619.994}{\Character{106}}%
\TextAt{13}{227.756}{619.994}{L}%
\TextAt{3}{235.72}{624.332}{\Character{48}}%
\TextAt{8}{238.515}{619.994}{,}%
\TextAt{8}{245.542}{619.994}{by}%
\TextAt{8}{261.648}{619.994}{(7.2).}%
\TextAt{8}{294.123}{619.994}{This}%
\TextAt{8}{320.69}{619.994}{implies}%
\TextAt{8}{360.264}{619.994}{that}%
\TextAt{13}{385.48}{619.994}{LL}%
\TextAt{3}{401.409}{624.332}{\Character{48}}%
\TextAt{16}{407.525}{619.994}{\Character{106}}%
\TextAt{13}{414.166}{619.994}{K}%
\TextAt{8}{428.728}{619.994}{is}%
\TextAt{8}{440.339}{619.994}{unramified}%
\TextAt{8}{498.719}{619.994}{as}%
\TextAt{8}{512.936}{619.994}{well}%
\TextAt{8}{79.226}{602.948}{because}%
\TextAt{8}{123.456}{602.948}{separability}%
\TextAt{8}{187.544}{602.948}{is}%
\TextAt{8}{200.236}{602.948}{transitive}%
\TextAt{8}{253.266}{602.948}{and}%
\TextAt{8}{276.957}{602.948}{the}%
\TextAt{8}{298.042}{602.948}{degrees}%
\TextAt{8}{340}{602.948}{of}%
\TextAt{8}{354.259}{602.948}{field}%
\TextAt{8}{380.544}{602.948}{(and}%
\TextAt{8}{408.79}{602.948}{residue}%
\TextAt{8}{449.446}{602.948}{field)}%
\TextAt{8}{480.286}{602.948}{extensions}%
\TextAt{8}{79.226}{585.902}{are}%
\TextAt{8}{98.737}{585.902}{multiplicative.}%
\Vector{q 0 G 0 g 0.399 w 452.334 -188.098 m 452.334 -180.029 l S Q}%
\Vector{q 0 G 0 g 0.399 w 452.534 -180.228 m 459.707 -180.228 l S Q}%
\Vector{q 0 G 0 g 0.399 w 452.534 -187.899 m 459.707 -187.899 l S Q}%
\Vector{q 0 G 0 g 0.399 w 459.906 -188.098 m 459.906 -180.029 l S Q}%
\TextAt{12}{79.226}{550.923}{Definition}%
\TextAt{12}{141.951}{550.923}{(7.4).}%
\TextAt{8}{177.211}{550.923}{Let}%
\TextAt{13}{198.124}{550.923}{L}%
\TextAt{16}{209.41}{550.923}{\Character{106}}%
\TextAt{13}{216.051}{550.923}{K}%
\TextAt{8}{230.711}{550.923}{be}%
\TextAt{8}{246.579}{550.923}{an}%
\TextAt{8}{262.783}{550.923}{algebraic}%
\TextAt{8}{312.162}{550.923}{extension.}%
\TextAt{8}{368.46}{550.923}{Then}%
\TextAt{8}{398.962}{550.923}{the}%
\TextAt{8}{419.063}{550.923}{composite}%
\TextAt{13}{474.025}{550.923}{T}%
\TextAt{16}{485.832}{550.923}{\Character{106}}%
\TextAt{13}{492.474}{550.923}{K}%
\TextAt{8}{507.133}{550.923}{of}%
\TextAt{8}{520.408}{550.923}{all}%
\TextAt{8}{79.226}{533.877}{unramified}%
\TextAt{8}{137.759}{533.877}{subextensions}%
\TextAt{8}{212.084}{533.877}{is}%
\TextAt{8}{223.848}{533.877}{called}%
\TextAt{8}{257.011}{533.877}{the}%
\TextAt{12}{277.168}{533.877}{maximal}%
\TextAt{12}{331.743}{533.877}{unramified}%
\TextAt{12}{399.242}{533.877}{subextension}%
\TextAt{8}{479.091}{533.877}{of}%
\TextAt{13}{492.421}{533.877}{L}%
\TextAt{16}{503.706}{533.877}{\Character{106}}%
\TextAt{13}{510.348}{533.877}{K}%
\TextAt{8}{521.165}{533.877}{.}%
\TextAt{12}{79.226}{501.887}{Proposition}%
\TextAt{12}{152.775}{501.887}{(7.5).}%
\TextAt{8}{188.819}{501.887}{The}%
\TextAt{8}{213.602}{501.887}{residue}%
\TextAt{8}{254.057}{501.887}{class}%
\TextAt{8}{282.223}{501.887}{field}%
\TextAt{8}{308.308}{501.887}{of}%
\TextAt{13}{322.367}{501.887}{T}%
\TextAt{8}{335.479}{501.887}{is}%
\TextAt{8}{347.972}{501.887}{the}%
\TextAt{8}{368.856}{501.887}{separable}%
\TextAt{8}{421.028}{501.887}{closure}%
\TextAt{13}{460.838}{501.887}{\Character{21}}%
\TextAt{21}{467.666}{500.094}{s}%
\TextAt{8}{476.706}{501.887}{of}%
\TextAt{13}{490.765}{501.887}{\Character{20}}%
\TextAt{8}{502.129}{501.887}{in}%
\TextAt{8}{516.51}{501.887}{the}%
\TextAt{8}{79.226}{484.841}{residue}%
\TextAt{8}{119.348}{484.841}{class}%
\TextAt{8}{147.181}{484.841}{field}%
\TextAt{8}{172.933}{484.841}{extension}%
\TextAt{13}{225.094}{484.841}{\Character{21}}%
\TextAt{16}{235.917}{484.841}{\Character{106}}%
\TextAt{13}{243.232}{484.841}{\Character{20}}%
\TextAt{8}{254.264}{484.841}{of}%
\TextAt{13}{267.989}{484.841}{L}%
\TextAt{16}{279.948}{484.841}{\Character{106}}%
\TextAt{13}{287.263}{484.841}{K}%
\TextAt{8}{298.081}{484.841}{,}%
\TextAt{8}{305.724}{484.841}{whereas}%
\TextAt{8}{350.401}{484.841}{the}%
\TextAt{8}{370.953}{484.841}{value}%
\TextAt{8}{401.595}{484.841}{group}%
\TextAt{8}{435.166}{484.841}{of}%
\TextAt{13}{448.891}{484.841}{T}%
\TextAt{8}{461.671}{484.841}{equals}%
\TextAt{8}{497.573}{484.841}{that}%
\TextAt{8}{523.337}{484.841}{of}%
\TextAt{13}{79.226}{467.795}{K}%
\TextAt{8}{90.044}{467.795}{.}%
\TextAt{19}{79.226}{432.816}{Proof.}%
\TextAt{8}{116}{432.816}{Let}%
\TextAt{13}{137.449}{432.816}{\Character{21}}%
\TextAt{17}{144.278}{431.023}{0}%
\TextAt{8}{153.387}{432.816}{be}%
\TextAt{8}{169.792}{432.816}{the}%
\TextAt{8}{190.428}{432.816}{residue}%
\TextAt{8}{230.635}{432.816}{class}%
\TextAt{8}{258.552}{432.816}{field}%
\TextAt{8}{284.388}{432.816}{of}%
\TextAt{13}{298.198}{432.816}{T}%
\TextAt{8}{311.062}{432.816}{and}%
\TextAt{8}{334.305}{432.816}{assume}%
\TextAt{9}{376.388}{432.816}{\Character{22}}%
\TextAt{13}{375.229}{432.816}{\Character{11}}%
\NativeMathAt{386.889}{432.816}{11.955}{\in}%
\TextAt{13}{398.998}{432.816}{\Character{21}}%
\TextAt{8}{410.204}{432.816}{is}%
\TextAt{8}{422.447}{432.816}{separable}%
\TextAt{8}{474.37}{432.816}{over}%
\TextAt{13}{499.896}{432.816}{\Character{20}}%
\TextAt{8}{506.635}{432.816}{.}%
\TextAt{8}{516.522}{432.816}{We}%
\TextAt{8}{79.226}{415.77}{have}%
\TextAt{8}{105.285}{415.77}{to}%
\TextAt{8}{118.66}{415.77}{show}%
\TextAt{8}{146.727}{415.77}{that}%
\TextAt{9}{172.32}{415.77}{\Character{22}}%
\TextAt{13}{171.16}{415.77}{\Character{11}}%
\NativeMathAt{182.003}{415.77}{11.955}{\in}%
\TextAt{13}{193.294}{415.77}{\Character{21}}%
\TextAt{17}{200.122}{413.977}{0}%
\TextAt{8}{204.854}{415.77}{.}%
\TextAt{8}{212.991}{415.77}{Let}%
\TextAt{9}{235.572}{418.925}{\Character{22}}%
\TextAt{13}{233.024}{415.77}{f}%
\TextAt{9}{240.071}{415.77}{\Character{40}}%
\TextAt{13}{244.623}{415.77}{x}%
\TextAt{9}{251.275}{415.77}{\Character{41}}%
\NativeMathAt{259.148}{415.77}{11.955}{\in}%
\TextAt{13}{270.439}{415.77}{\Character{20}}%
\TextAt{9}{277.178}{415.77}{\Character{91}}%
\TextAt{13}{280.429}{415.77}{x}%
\TextAt{9}{287.082}{415.77}{\Character{93}}%
\TextAt{8}{293.295}{415.77}{be}%
\TextAt{8}{308.283}{415.77}{the}%
\TextAt{8}{327.504}{415.77}{minimal}%
\TextAt{8}{372.093}{415.77}{polynomial}%
\TextAt{8}{431.651}{415.77}{of}%
\TextAt{9}{445.204}{415.77}{\Character{22}}%
\TextAt{13}{444.045}{415.77}{\Character{11}}%
\TextAt{8}{454.528}{415.77}{and}%
\TextAt{13}{476.355}{415.77}{f}%
\TextAt{9}{483.401}{415.77}{\Character{40}}%
\TextAt{13}{487.954}{415.77}{x}%
\TextAt{9}{494.606}{415.77}{\Character{41}}%
\NativeMathAt{502.479}{415.77}{11.955}{\in}%
\TextAt{14}{513.77}{415.77}{o}%
\TextAt{9}{519.614}{415.77}{\Character{91}}%
\TextAt{13}{522.866}{415.77}{x}%
\TextAt{9}{529.518}{415.77}{\Character{93}}%
\TextAt{8}{79.226}{398.724}{a}%
\TextAt{8}{89.941}{398.724}{monic}%
\TextAt{8}{125.367}{398.724}{polynomial}%
\TextAt{8}{186.82}{398.724}{such}%
\TextAt{8}{214.176}{398.724}{that}%
\TextAt{9}{243.052}{401.879}{\Character{22}}%
\TextAt{13}{240.504}{398.724}{f}%
\TextAt{9}{252.507}{398.724}{\Character{61}}%
\TextAt{13}{266.567}{398.724}{f}%
\TextAt{8}{276.934}{398.724}{mod}%
\TextAt{14}{302.372}{398.724}{p}%
\TextAt{8}{308.355}{398.724}{.}%
\TextAt{8}{319.682}{398.724}{Then}%
\TextAt{13}{351.199}{398.724}{f}%
\TextAt{9}{358.245}{398.724}{\Character{40}}%
\TextAt{13}{362.797}{398.724}{x}%
\TextAt{9}{369.449}{398.724}{\Character{41}}%
\TextAt{8}{378.859}{398.724}{is}%
\TextAt{8}{391.582}{398.724}{irreducible}%
\TextAt{8}{450.416}{398.724}{and}%
\TextAt{8}{474.138}{398.724}{by}%
\TextAt{8}{491.357}{398.724}{Hensel’s}%
\TextAt{8}{79.226}{381.678}{lemma}%
\TextAt{8}{117.818}{381.678}{has}%
\TextAt{8}{139.566}{381.678}{a}%
\TextAt{8}{150.195}{381.678}{root}%
\TextAt{13}{176.115}{381.678}{\Character{11}}%
\TextAt{8}{188.407}{381.678}{in}%
\TextAt{13}{202.934}{381.678}{L}%
\TextAt{8}{215.669}{381.678}{such}%
\TextAt{8}{242.94}{381.678}{that}%
\TextAt{9}{270.342}{381.678}{\Character{22}}%
\TextAt{13}{269.182}{381.678}{\Character{11}}%
\TextAt{9}{281.514}{381.678}{\Character{61}}%
\TextAt{13}{295.428}{381.678}{\Character{11}}%
\TextAt{8}{306.271}{381.678}{mod}%
\TextAt{14}{331.709}{381.678}{P}%
\TextAt{8}{341.602}{381.678}{,}%
\TextAt{8}{349.843}{381.678}{i.e.,}%
\TextAt{9}{373.04}{381.678}{\Character{91}}%
\TextAt{13}{376.292}{381.678}{K}%
\TextAt{9}{387.109}{381.678}{\Character{40}}%
\TextAt{13}{391.662}{381.678}{\Character{11}}%
\TextAt{9}{399.184}{381.678}{\Character{41}}%
\TextAt{9}{408.533}{381.678}{\Character{58}}%
\TextAt{13}{416.607}{381.678}{K}%
\TextAt{9}{427.424}{381.678}{\Character{93}}%
\TextAt{9}{435.47}{381.678}{\Character{61}}%
\TextAt{9}{449.374}{381.678}{\Character{91}}%
\TextAt{13}{452.652}{381.678}{\Character{20}}%
\TextAt{9}{459.39}{381.678}{\Character{40}}%
\TextAt{9}{465.093}{381.678}{\Character{22}}%
\TextAt{13}{463.943}{381.678}{\Character{11}}%
\TextAt{9}{471.464}{381.678}{\Character{41}}%
\TextAt{9}{480.813}{381.678}{\Character{58}}%
\TextAt{13}{488.888}{381.678}{\Character{20}}%
\TextAt{9}{495.626}{381.678}{\Character{93}}%
\TextAt{8}{498.878}{381.678}{.}%
\TextAt{8}{509.947}{381.678}{This}%
\TextAt{8}{79.226}{364.632}{implies}%
\TextAt{8}{118.953}{364.632}{that}%
\TextAt{13}{144.322}{364.632}{K}%
\TextAt{9}{155.14}{364.632}{\Character{40}}%
\TextAt{13}{159.692}{364.632}{\Character{11}}%
\TextAt{9}{167.214}{364.632}{\Character{41}}%
\TextAt{16}{175.087}{364.632}{\Character{106}}%
\TextAt{13}{181.729}{364.632}{K}%
\TextAt{8}{196.444}{364.632}{is}%
\TextAt{8}{208.208}{364.632}{unramified,}%
\TextAt{8}{269.992}{364.632}{so}%
\TextAt{8}{284.362}{364.632}{that}%
\TextAt{13}{309.731}{364.632}{K}%
\TextAt{9}{320.549}{364.632}{\Character{40}}%
\TextAt{13}{325.101}{364.632}{\Character{11}}%
\TextAt{9}{332.623}{364.632}{\Character{41}}%
\NativeMathAt{340.496}{364.632}{11.955}{\subseteq}%
\TextAt{13}{353.115}{364.632}{T}%
\TextAt{8}{361.602}{364.632}{,}%
\TextAt{8}{368.751}{364.632}{and}%
\TextAt{8}{391.514}{364.632}{thus}%
\TextAt{9}{418.425}{364.632}{\Character{22}}%
\TextAt{13}{417.265}{364.632}{\Character{11}}%
\NativeMathAt{428.108}{364.632}{11.955}{\in}%
\TextAt{13}{439.399}{364.632}{\Character{21}}%
\TextAt{17}{446.227}{362.839}{0}%
\TextAt{8}{450.96}{364.632}{.}%
\TextAt{8}{97.159}{344.598}{In}%
\TextAt{8}{111.832}{344.598}{order}%
\TextAt{8}{142.452}{344.598}{to}%
\TextAt{8}{156.814}{344.598}{prove}%
\TextAt{13}{188.415}{344.598}{w}%
\TextAt{9}{197.105}{344.598}{\Character{40}}%
\TextAt{13}{201.657}{344.598}{T}%
\TextAt{3}{210.144}{348.936}{\Character{3}}%
\TextAt{9}{214.876}{344.598}{\Character{41}}%
\TextAt{9}{222.826}{344.598}{\Character{61}}%
\TextAt{13}{235.35}{344.598}{v}%
\TextAt{9}{241.439}{344.598}{\Character{40}}%
\TextAt{13}{245.991}{344.598}{K}%
\TextAt{3}{256.809}{348.936}{\Character{3}}%
\TextAt{9}{261.541}{344.598}{\Character{41}}%
\TextAt{8}{270.042}{344.598}{we}%
\TextAt{8}{287.321}{344.598}{may}%
\TextAt{8}{312.741}{344.598}{suppose}%
\TextAt{13}{356.812}{344.598}{L}%
\TextAt{16}{368.185}{344.598}{\Character{106}}%
\TextAt{13}{374.914}{344.598}{K}%
\TextAt{8}{389.681}{344.598}{to}%
\TextAt{8}{404.042}{344.598}{be}%
\TextAt{8}{420.018}{344.598}{finite.}%
\TextAt{8}{454.635}{344.598}{The}%
\TextAt{8}{478.74}{344.598}{claim}%
\TextAt{8}{510.007}{344.598}{then}%
\TextAt{8}{79.226}{327.552}{follows}%
\TextAt{8}{117.662}{327.552}{from}%
\TextAt{9}{154.725}{298.551}{\Character{91}}%
\TextAt{13}{157.977}{298.551}{T}%
\TextAt{9}{169.785}{298.551}{\Character{58}}%
\TextAt{13}{176.357}{298.551}{K}%
\TextAt{9}{187.175}{298.551}{\Character{93}}%
\NativeMathAt{193.747}{298.551}{11.955}{\geq}%
\TextAt{9}{206.367}{298.551}{\Character{40}}%
\TextAt{13}{210.919}{298.551}{w}%
\TextAt{9}{219.609}{298.551}{\Character{40}}%
\TextAt{13}{224.161}{298.551}{T}%
\TextAt{3}{232.648}{303.487}{\Character{3}}%
\TextAt{9}{237.38}{298.551}{\Character{41}}%
\TextAt{9}{245.235}{298.551}{\Character{58}}%
\TextAt{13}{251.826}{298.551}{v}%
\TextAt{9}{257.914}{298.551}{\Character{40}}%
\TextAt{13}{262.466}{298.551}{K}%
\TextAt{3}{273.284}{303.487}{\Character{3}}%
\TextAt{9}{278.016}{298.551}{\Character{41}\Character{41}\Character{91}}%
\TextAt{13}{290.373}{298.551}{\Character{21}}%
\TextAt{17}{297.201}{296.757}{0}%
\TextAt{9}{305.254}{298.551}{\Character{58}}%
\TextAt{13}{311.827}{298.551}{\Character{20}}%
\TextAt{9}{318.565}{298.551}{\Character{93}}%
\TextAt{9}{325.116}{298.551}{\Character{61}}%
\TextAt{9}{337.538}{298.551}{\Character{40}}%
\TextAt{13}{342.116}{298.551}{w}%
\TextAt{9}{350.805}{298.551}{\Character{40}}%
\TextAt{13}{355.357}{298.551}{T}%
\TextAt{3}{363.844}{303.487}{\Character{3}}%
\TextAt{9}{368.577}{298.551}{\Character{41}}%
\TextAt{9}{376.432}{298.551}{\Character{58}}%
\TextAt{13}{383.022}{298.551}{v}%
\TextAt{9}{389.11}{298.551}{\Character{40}}%
\TextAt{13}{393.663}{298.551}{K}%
\TextAt{3}{404.48}{303.487}{\Character{3}}%
\TextAt{9}{409.213}{298.551}{\Character{41}\Character{41}\Character{91}}%
\TextAt{13}{421.569}{298.551}{T}%
\TextAt{9}{433.377}{298.551}{\Character{58}}%
\TextAt{13}{439.949}{298.551}{K}%
\TextAt{9}{450.767}{298.551}{\Character{93}}%
\TextAt{13}{454.019}{298.551}{\Character{58}}%
\Vector{q 0 G 0 g 0.399 w 452.334 -504.451 m 452.334 -496.381 l S Q}%
\Vector{q 0 G 0 g 0.399 w 452.534 -496.58 m 459.707 -496.58 l S Q}%
\Vector{q 0 G 0 g 0.399 w 452.534 -504.251 m 459.707 -504.251 l S Q}%
\Vector{q 0 G 0 g 0.399 w 459.906 -504.451 m 459.906 -496.381 l S Q}%
\TextAt{8}{97.159}{240.548}{The}%
\TextAt{8}{122.714}{240.548}{composite}%
\TextAt{8}{179.233}{240.548}{of}%
\TextAt{8}{194.064}{240.548}{all}%
\TextAt{8}{211.824}{240.548}{unramified}%
\TextAt{8}{271.857}{240.548}{extensions}%
\TextAt{8}{329.739}{240.548}{inside}%
\TextAt{8}{364.464}{240.548}{the}%
\TextAt{8}{386.121}{240.548}{algebraic}%
\TextAt{8}{437.057}{240.548}{closure}%
\Vector{q 0 G 0 g 0.478 w 405.64 -523.609 m 416.457 -523.609 l S Q}%
\TextAt{13}{477.64}{240.548}{K}%
\TextAt{8}{493.856}{240.548}{of}%
\TextAt{13}{508.687}{240.548}{K}%
\TextAt{8}{524.903}{240.548}{is}%
\TextAt{8}{79.226}{223.502}{simply}%
\TextAt{8}{116.566}{223.502}{called}%
\TextAt{8}{149.613}{223.502}{the}%
\TextAt{12}{169.654}{223.502}{maximal}%
\TextAt{12}{224.096}{223.502}{unramified}%
\TextAt{12}{291.461}{223.502}{extension}%
\TextAt{13}{350.571}{223.502}{K}%
\TextAt{21}{360.531}{221.709}{nr}%
\TextAt{16}{373.544}{223.502}{\Character{106}}%
\TextAt{13}{380.186}{223.502}{K}%
\TextAt{8}{394.785}{223.502}{of}%
\TextAt{13}{407.999}{223.502}{K}%
\TextAt{8}{422.598}{223.502}{(}%
\TextAt{13}{427.153}{223.502}{nr}%
\TextAt{9}{443.062}{223.502}{\Character{61}}%
\TextAt{8}{455.488}{223.502}{‘non}%
\TextAt{8}{481.386}{223.502}{ramifiée’).}%
\TextAt{8}{79.226}{206.456}{Its}%
\TextAt{8}{97.31}{206.456}{residue}%
\TextAt{8}{137.835}{206.456}{class}%
\TextAt{8}{166.069}{206.456}{field}%
\TextAt{8}{192.223}{206.456}{is}%
\TextAt{8}{204.784}{206.456}{the}%
\TextAt{8}{225.737}{206.456}{separable}%
\TextAt{8}{277.977}{206.456}{closure}%
\TextAt{9}{318.299}{206.456}{\Character{22}}%
\TextAt{13}{317.856}{206.456}{\Character{20}}%
\TextAt{21}{324.594}{204.663}{s}%
\TextAt{16}{333.688}{206.456}{\Character{106}}%
\TextAt{13}{341.688}{206.456}{\Character{20}}%
\TextAt{8}{348.426}{206.456}{.}%
\TextAt{13}{359.266}{206.456}{K}%
\TextAt{21}{369.225}{204.663}{nr}%
\TextAt{8}{383.612}{206.456}{contains}%
\TextAt{8}{430.329}{206.456}{all}%
\TextAt{8}{447.385}{206.456}{roots}%
\TextAt{8}{477.843}{206.456}{of}%
\TextAt{8}{491.97}{206.456}{unity}%
\TextAt{8}{523.337}{206.456}{of}%
\TextAt{8}{79.226}{189.41}{order}%
\TextAt{13}{109.755}{189.41}{m}%
\TextAt{8}{123.85}{189.41}{not}%
\TextAt{8}{144.623}{189.41}{divisible}%
\TextAt{8}{190.49}{189.41}{by}%
\TextAt{8}{206.708}{189.41}{the}%
\TextAt{8}{226.824}{189.41}{characteristic}%
\TextAt{8}{298.717}{189.41}{of}%
\TextAt{13}{312.006}{189.41}{\Character{20}}%
\TextAt{8}{322.601}{189.41}{because}%
\TextAt{8}{365.862}{189.41}{the}%
\TextAt{8}{385.977}{189.41}{separable}%
\TextAt{8}{437.379}{189.41}{polynomial}%
\TextAt{13}{497.831}{189.41}{x}%
\TextAt{21}{504.484}{193.749}{m}%
\TextAt{16}{515.045}{189.41}{\Character{0}}%
\TextAt{9}{526.916}{189.41}{1}%
\TextAt{8}{79.226}{172.364}{splits}%
\TextAt{8}{110.063}{172.364}{over}%
\TextAt{9}{135.7}{172.364}{\Character{22}}%
\TextAt{13}{135.257}{172.364}{\Character{20}}%
\TextAt{21}{141.996}{170.571}{s}%
\TextAt{8}{150.456}{172.364}{and}%
\TextAt{8}{173.366}{172.364}{hence}%
\TextAt{8}{206.021}{172.364}{also}%
\TextAt{8}{229.649}{172.364}{over}%
\TextAt{13}{254.843}{172.364}{K}%
\TextAt{21}{264.802}{170.571}{nr}%
\TextAt{8}{274.495}{172.364}{,}%
\TextAt{8}{281.829}{172.364}{by}%
\TextAt{8}{298.236}{172.364}{Hensel’s}%
\TextAt{8}{343.695}{172.364}{lemma.}%
\TextAt{8}{386.409}{172.364}{If}%
\TextAt{13}{398.249}{172.364}{\Character{20}}%
\TextAt{8}{409.033}{172.364}{is}%
\TextAt{8}{420.945}{172.364}{a}%
\TextAt{8}{430.848}{172.364}{finite}%
\TextAt{8}{460.908}{172.364}{field,}%
\TextAt{8}{489.702}{172.364}{then}%
\TextAt{8}{516.51}{172.364}{the}%
\TextAt{8}{79.226}{155.318}{extension}%
\TextAt{13}{130.863}{155.318}{K}%
\TextAt{21}{140.823}{153.525}{nr}%
\TextAt{16}{153.836}{155.318}{\Character{106}}%
\TextAt{13}{160.478}{155.318}{K}%
\TextAt{8}{175.065}{155.318}{is}%
\TextAt{8}{186.7}{155.318}{even}%
\TextAt{8}{213.231}{155.318}{generated}%
\TextAt{8}{266.435}{155.318}{by}%
\TextAt{8}{282.565}{155.318}{these}%
\TextAt{8}{312.408}{155.318}{roots}%
\TextAt{8}{341.94}{155.318}{of}%
\TextAt{8}{355.142}{155.318}{unity}%
\TextAt{8}{385.582}{155.318}{because}%
\TextAt{8}{428.755}{155.318}{they}%
\TextAt{8}{454.964}{155.318}{generate}%
\TextAt{9}{502.107}{155.318}{\Character{22}}%
\TextAt{13}{501.664}{155.318}{\Character{20}}%
\TextAt{21}{508.402}{153.525}{s}%
\TextAt{16}{516.137}{155.318}{\Character{106}}%
\TextAt{13}{522.779}{155.318}{\Character{20}}%
\TextAt{8}{529.518}{155.318}{.}%
\TextAt{8}{97.159}{135.284}{If}%
\TextAt{8}{108.761}{135.284}{the}%
\TextAt{8}{128.828}{135.284}{characteristic}%
\TextAt{13}{200.672}{135.284}{p}%
\TextAt{9}{209.868}{135.284}{\Character{61}}%
\TextAt{8}{222.294}{135.284}{char}%
\TextAt{9}{244.411}{135.284}{\Character{40}}%
\TextAt{13}{248.963}{135.284}{\Character{20}}%
\TextAt{9}{255.701}{135.284}{\Character{41}}%
\TextAt{8}{264.061}{135.284}{of}%
\TextAt{8}{277.301}{135.284}{the}%
\TextAt{8}{297.368}{135.284}{residue}%
\TextAt{8}{337.005}{135.284}{class}%
\TextAt{8}{364.352}{135.284}{field}%
\TextAt{8}{389.619}{135.284}{is}%
\TextAt{8}{401.293}{135.284}{positive,}%
\TextAt{8}{447.787}{135.284}{then}%
\TextAt{8}{474.357}{135.284}{one}%
\TextAt{8}{495.727}{135.284}{has}%
\TextAt{8}{516.51}{135.284}{the}%
\TextAt{8}{79.226}{118.238}{following}%
\TextAt{8}{128.661}{118.238}{weaker}%
\TextAt{8}{167.36}{118.238}{notion}%
\TextAt{8}{203.787}{118.238}{accompanying}%
\TextAt{8}{279.894}{118.238}{that}%
\TextAt{8}{305.263}{118.238}{of}%
\TextAt{8}{318.593}{118.238}{an}%
\TextAt{8}{334.852}{118.238}{unramified}%
\TextAt{8}{393.384}{118.238}{extension.}%
\TextAt{12}{79.226}{92.225}{Definition}%
\TextAt{12}{141.889}{92.225}{(7.6).}%
\TextAt{8}{177.095}{92.225}{An}%
\TextAt{8}{196.162}{92.225}{algebraic}%
\TextAt{8}{245.488}{92.225}{extension}%
\TextAt{13}{297.144}{92.225}{L}%
\TextAt{16}{308.43}{92.225}{\Character{106}}%
\TextAt{13}{315.071}{92.225}{K}%
\TextAt{8}{329.677}{92.225}{is}%
\TextAt{8}{341.332}{92.225}{called}%
\TextAt{12}{374.386}{92.225}{tamely}%
\TextAt{12}{418.71}{92.225}{ramified}%
\TextAt{8}{470.57}{92.225}{if}%
\TextAt{8}{481.185}{92.225}{the}%
\TextAt{8}{501.232}{92.225}{exten-}%
\TextAt{8}{79.226}{75.179}{sion}%
\TextAt{13}{102.612}{75.179}{\Character{21}}%
\TextAt{16}{112.761}{75.179}{\Character{106}}%
\TextAt{13}{119.403}{75.179}{\Character{20}}%
\TextAt{8}{129.299}{75.179}{of}%
\TextAt{8}{141.89}{75.179}{the}%
\TextAt{8}{161.307}{75.179}{residue}%
\TextAt{8}{200.294}{75.179}{class}%
\TextAt{8}{226.991}{75.179}{fields}%
\TextAt{8}{256.223}{75.179}{is}%
\TextAt{8}{267.248}{75.179}{separable}%
\TextAt{8}{317.951}{75.179}{and}%
\TextAt{8}{339.974}{75.179}{one}%
\TextAt{8}{360.694}{75.179}{has}%
\TextAt{9}{380.828}{75.179}{\Character{40}\Character{91}}%
\TextAt{13}{388.632}{75.179}{L}%
\TextAt{9}{399.918}{75.179}{\Character{58}}%
\TextAt{13}{406.49}{75.179}{T}%
\TextAt{9}{414.977}{75.179}{\Character{93}}%
\TextAt{13}{418.229}{75.179}{\Character{59}}%
\TextAt{13}{423.453}{75.179}{p}%
\TextAt{9}{429.348}{75.179}{\Character{41}}%
\TextAt{9}{437.203}{75.179}{\Character{61}}%
\TextAt{9}{449.624}{75.179}{1}%
\TextAt{8}{455.499}{75.179}{.}%
\TextAt{8}{463.701}{75.179}{In}%
\TextAt{8}{477.583}{75.179}{the}%
\TextAt{8}{497}{75.179}{infinite}%
\TextAt{8}{297.211}{43.033}{141}%
\Destination{78.23}{809.815}{nav.page.141}%
\Destination{79.36}{652.055}{nav.print.154}%
\Destination{79.36}{686.055}{nav.number.2.7.3}%
\Destination{79.36}{564.955}{nav.number.2.7.4}%
\Destination{79.36}{515.855}{nav.number.2.7.5}%
\Destination{79.36}{106.255}{nav.number.2.7.6}%
\Hotlink{261.148}{617.755}{25.078}{11.047}{nav.number.2.7.2}%
}

% Page 154 (body page 142)
\Page{%
\TextAt{2}{79.226}{796.943}{II.}%
\TextAt{2}{94.073}{796.943}{The}%
\TextAt{2}{116.492}{796.943}{Theory}%
\TextAt{2}{154.401}{796.943}{of}%
\TextAt{2}{166.826}{796.943}{Valuations}%
\Vector{q 0 G 0 g 0.299 w 7.226 18.729 m 460.769 18.729 l S Q}%
\TextAt{8}{79.226}{761.124}{case}%
\TextAt{8}{104.189}{761.124}{this}%
\TextAt{8}{127.202}{761.124}{latter}%
\TextAt{8}{159.266}{761.124}{condition}%
\TextAt{8}{210.841}{761.124}{is}%
\TextAt{8}{222.796}{761.124}{taken}%
\TextAt{8}{254.86}{761.124}{to}%
\TextAt{8}{269.361}{761.124}{mean}%
\TextAt{8}{300.767}{761.124}{that}%
\TextAt{8}{326.328}{761.124}{the}%
\TextAt{8}{346.675}{761.124}{degree}%
\TextAt{8}{383.282}{761.124}{of}%
\TextAt{8}{396.803}{761.124}{each}%
\TextAt{8}{423.332}{761.124}{finite}%
\TextAt{8}{453.435}{761.124}{subextension}%
\TextAt{8}{523.337}{761.124}{of}%
\TextAt{13}{79.226}{744.078}{L}%
\TextAt{16}{90.511}{744.078}{\Character{106}}%
\TextAt{13}{97.153}{744.078}{T}%
\TextAt{8}{109.537}{744.078}{is}%
\TextAt{8}{121.301}{744.078}{prime}%
\TextAt{8}{154.465}{744.078}{to}%
\TextAt{13}{168.775}{744.078}{p}%
\TextAt{8}{174.65}{744.078}{.}%
\TextAt{8}{97.159}{718.745}{As}%
\TextAt{8}{116.027}{718.745}{before,}%
\TextAt{8}{156.367}{718.745}{in}%
\TextAt{8}{171.6}{718.745}{this}%
\TextAt{8}{196.003}{718.745}{definition}%
\TextAt{13}{249.613}{718.745}{K}%
\TextAt{8}{265.908}{718.745}{need}%
\TextAt{8}{294.795}{718.745}{not}%
\TextAt{8}{317.189}{718.745}{be}%
\TextAt{8}{334.694}{718.745}{henselian.}%
\TextAt{8}{394.773}{718.745}{We}%
\TextAt{8}{416.498}{718.745}{apply}%
\TextAt{8}{450.274}{718.745}{it}%
\TextAt{8}{463.558}{718.745}{whenever}%
\TextAt{8}{516.51}{718.745}{the}%
\TextAt{8}{79.226}{701.699}{valuation}%
\TextAt{13}{129.151}{701.699}{v}%
\TextAt{8}{137.989}{701.699}{of}%
\TextAt{13}{150.172}{701.699}{K}%
\TextAt{8}{163.739}{701.699}{has}%
\TextAt{8}{183.465}{701.699}{a}%
\TextAt{8}{192.073}{701.699}{unique}%
\TextAt{8}{228.967}{701.699}{extension}%
\TextAt{8}{279.585}{701.699}{to}%
\TextAt{13}{292.748}{701.699}{L}%
\TextAt{8}{300.712}{701.699}{.}%
\TextAt{8}{308.778}{701.699}{When}%
\TextAt{8}{341.763}{701.699}{the}%
\TextAt{8}{360.771}{701.699}{fundamental}%
\TextAt{8}{427.266}{701.699}{identity}%
\TextAt{13}{469.372}{701.699}{ef}%
\TextAt{9}{485.165}{701.699}{\Character{61}}%
\TextAt{9}{497.575}{701.699}{\Character{91}}%
\TextAt{13}{500.842}{701.699}{L}%
\TextAt{9}{512.128}{701.699}{\Character{58}}%
\TextAt{13}{518.7}{701.699}{K}%
\TextAt{9}{529.518}{701.699}{\Character{93}}%
\TextAt{8}{79.226}{684.653}{holds}%
\TextAt{8}{109.068}{684.653}{and}%
\TextAt{13}{131.042}{684.653}{\Character{21}}%
\TextAt{16}{141.192}{684.653}{\Character{106}}%
\TextAt{13}{147.833}{684.653}{\Character{20}}%
\TextAt{8}{157.682}{684.653}{is}%
\TextAt{8}{168.658}{684.653}{separable,}%
\TextAt{8}{222.722}{684.653}{to}%
\TextAt{8}{236.245}{684.653}{say}%
\TextAt{8}{255.685}{684.653}{that}%
\TextAt{8}{280.266}{684.653}{the}%
\TextAt{8}{299.635}{684.653}{extension}%
\TextAt{8}{350.613}{684.653}{is}%
\TextAt{8}{361.589}{684.653}{unramified,}%
\TextAt{8}{422.743}{684.653}{resp.}%
\TextAt{8}{449.979}{684.653}{tamely}%
\TextAt{8}{487.89}{684.653}{ramified,}%
\TextAt{8}{79.226}{667.607}{simply}%
\TextAt{8}{116.682}{667.607}{amounts}%
\TextAt{8}{163.905}{667.607}{to}%
\TextAt{8}{178.215}{667.607}{saying}%
\TextAt{8}{214.057}{667.607}{that}%
\TextAt{13}{239.425}{667.607}{e}%
\TextAt{9}{248.172}{667.607}{\Character{61}}%
\TextAt{9}{260.582}{667.607}{1}%
\TextAt{8}{266.45}{667.607}{,}%
\TextAt{8}{273.599}{667.607}{resp.}%
\TextAt{9}{301.622}{667.607}{\Character{40}}%
\TextAt{13}{306.175}{667.607}{e\Character{59}}%
\TextAt{13}{316.815}{667.607}{p}%
\TextAt{9}{322.719}{667.607}{\Character{41}}%
\TextAt{9}{330.574}{667.607}{\Character{61}}%
\TextAt{9}{342.995}{667.607}{1}%
\TextAt{8}{348.871}{667.607}{.}%
\TextAt{12}{79.226}{642.274}{Proposition}%
\TextAt{12}{150.686}{642.274}{(7.7).}%
\TextAt{8}{184.914}{642.274}{A}%
\TextAt{8}{196.498}{642.274}{finite}%
\TextAt{8}{225.322}{642.274}{extension}%
\TextAt{13}{276.001}{642.274}{L}%
\TextAt{16}{287.286}{642.274}{\Character{106}}%
\TextAt{13}{293.928}{642.274}{K}%
\TextAt{8}{307.555}{642.274}{is}%
\TextAt{8}{318.232}{642.274}{tamely}%
\TextAt{8}{355.843}{642.274}{ramified}%
\TextAt{8}{400.281}{642.274}{if}%
\TextAt{8}{409.917}{642.274}{and}%
\TextAt{8}{431.592}{642.274}{only}%
\TextAt{8}{456.196}{642.274}{if}%
\TextAt{8}{465.832}{642.274}{the}%
\TextAt{8}{484.901}{642.274}{extension}%
\TextAt{13}{79.226}{625.228}{L}%
\TextAt{16}{90.511}{625.228}{\Character{106}}%
\TextAt{13}{97.153}{625.228}{T}%
\TextAt{8}{109.537}{625.228}{is}%
\TextAt{8}{121.301}{625.228}{generated}%
\TextAt{8}{174.633}{625.228}{by}%
\TextAt{8}{190.892}{625.228}{radicals}%
\TextAt{13}{243.186}{598.604}{L}%
\TextAt{9}{254.472}{598.604}{\Character{61}}%
\TextAt{13}{266.897}{598.604}{T}%
\TextAt{20}{275.384}{608.288}{\Character{0}}%
\TextAt{22}{284.184}{602.249}{m}%
\TextAt{24}{290.722}{600.456}{1}%
\TextAt{16}{288.232}{606.902}{\Character{112}}%
\Vector{q 0 G 0 g 0.478 w 226.194 -166.859 m 237.072 -166.859 l S Q}%
\TextAt{13}{298.194}{598.604}{a}%
\TextAt{17}{304.339}{596.81}{1}%
\TextAt{13}{309.072}{598.604}{\Character{59}}%
\TextAt{13}{314.296}{598.604}{\Character{58}}%
\TextAt{13}{319.533}{598.604}{\Character{58}}%
\TextAt{13}{324.757}{598.604}{\Character{58}}%
\TextAt{13}{329.994}{598.604}{\Character{59}}%
\TextAt{22}{338.613}{602.249}{m}%
\TextAt{22}{345.151}{601.253}{r}%
\TextAt{16}{342.667}{606.902}{\Character{112}}%
\Vector{q 0 G 0 g 0.478 w 280.63 -166.859 m 291.33 -166.859 l S Q}%
\TextAt{13}{352.63}{598.604}{a}%
\TextAt{21}{358.775}{596.81}{r}%
\TextAt{20}{363.33}{608.288}{\Character{1}}%
\TextAt{8}{79.226}{571.98}{such}%
\TextAt{8}{105.623}{571.98}{that}%
\TextAt{9}{130.992}{571.98}{\Character{40}}%
\TextAt{13}{135.544}{571.98}{m}%
\TextAt{21}{145.784}{570.187}{i}%
\TextAt{13}{149.165}{571.98}{\Character{59}}%
\TextAt{13}{154.389}{571.98}{p}%
\TextAt{9}{160.284}{571.98}{\Character{41}}%
\TextAt{9}{168.139}{571.98}{\Character{61}}%
\TextAt{9}{180.56}{571.98}{1}%
\TextAt{8}{186.436}{571.98}{.}%
\TextAt{8}{194.884}{571.98}{In}%
\TextAt{8}{209.505}{571.98}{this}%
\TextAt{8}{232.328}{571.98}{case}%
\TextAt{8}{257.099}{571.98}{the}%
\TextAt{8}{277.255}{571.98}{fundamental}%
\TextAt{8}{344.898}{571.98}{identity}%
\TextAt{8}{388.151}{571.98}{always}%
\TextAt{8}{425.619}{571.98}{holds:}%
\TextAt{9}{272.999}{545.356}{\Character{91}}%
\TextAt{13}{276.25}{545.356}{L}%
\TextAt{9}{287.536}{545.356}{\Character{58}}%
\TextAt{13}{294.108}{545.356}{K}%
\TextAt{9}{304.926}{545.356}{\Character{93}}%
\TextAt{9}{311.477}{545.356}{\Character{61}}%
\TextAt{13}{323.924}{545.356}{ef}%
\TextAt{13}{335.724}{545.356}{\Character{58}}%
\TextAt{19}{79.226}{502.158}{Proof.}%
\TextAt{8}{116}{502.158}{We}%
\TextAt{8}{136.846}{502.158}{may}%
\TextAt{8}{162.916}{502.158}{assume}%
\TextAt{8}{204.062}{502.158}{that}%
\TextAt{13}{230.132}{502.158}{K}%
\TextAt{9}{245.466}{502.158}{\Character{61}}%
\TextAt{13}{259.086}{502.158}{T}%
\TextAt{8}{272.172}{502.158}{because}%
\TextAt{13}{316.175}{502.158}{L}%
\TextAt{16}{328.655}{502.158}{\Character{106}}%
\TextAt{13}{336.492}{502.158}{K}%
\TextAt{8}{351.908}{502.158}{is}%
\TextAt{8}{364.373}{502.158}{obviously}%
\TextAt{8}{416.852}{502.158}{tamely}%
\TextAt{8}{456.253}{502.158}{ramified}%
\TextAt{8}{502.479}{502.158}{if}%
\TextAt{8}{513.904}{502.158}{and}%
\TextAt{8}{79.226}{485.112}{only}%
\TextAt{8}{106.136}{485.112}{if}%
\TextAt{13}{118.078}{485.112}{L}%
\TextAt{16}{131.44}{485.112}{\Character{106}}%
\TextAt{13}{140.158}{485.112}{T}%
\TextAt{8}{153.76}{485.112}{is}%
\TextAt{8}{166.742}{485.112}{tamely}%
\TextAt{8}{206.66}{485.112}{ramified,}%
\TextAt{8}{256.96}{485.112}{and}%
\TextAt{8}{280.94}{485.112}{if}%
\TextAt{8}{292.883}{485.112}{this}%
\TextAt{8}{316.923}{485.112}{is}%
\TextAt{8}{329.906}{485.112}{the}%
\TextAt{8}{351.28}{485.112}{case,}%
\TextAt{8}{380.826}{485.112}{then}%
\TextAt{9}{408.704}{485.112}{\Character{91}}%
\TextAt{13}{411.956}{485.112}{T}%
\TextAt{9}{425.84}{485.112}{\Character{58}}%
\TextAt{13}{434.489}{485.112}{K}%
\TextAt{9}{445.306}{485.112}{\Character{93}}%
\TextAt{9}{453.938}{485.112}{\Character{61}}%
\TextAt{9}{468.427}{485.112}{\Character{91}}%
\TextAt{13}{471.708}{485.112}{\Character{21}}%
\TextAt{9}{483.934}{485.112}{\Character{58}}%
\TextAt{13}{492.582}{485.112}{\Character{20}}%
\TextAt{9}{499.321}{485.112}{\Character{93}}%
\TextAt{9}{507.953}{485.112}{\Character{61}}%
\TextAt{13}{522.471}{485.112}{f}%
\TextAt{8}{529.518}{485.112}{.}%
\TextAt{8}{79.226}{468.066}{Let}%
\TextAt{13}{100.91}{468.066}{L}%
\TextAt{16}{113.412}{468.066}{\Character{106}}%
\TextAt{13}{121.27}{468.066}{K}%
\TextAt{8}{136.699}{468.066}{be}%
\TextAt{8}{153.338}{468.066}{tamely}%
\TextAt{8}{192.751}{468.066}{ramified,}%
\TextAt{8}{242.42}{468.066}{so}%
\TextAt{8}{257.505}{468.066}{that}%
\TextAt{13}{283.587}{468.066}{\Character{20}}%
\TextAt{9}{294.864}{468.066}{\Character{61}}%
\TextAt{13}{308.506}{468.066}{\Character{21}}%
\TextAt{8}{319.946}{468.066}{and}%
\TextAt{9}{343.423}{468.066}{\Character{40}\Character{91}}%
\TextAt{13}{351.227}{468.066}{L}%
\TextAt{9}{363.729}{468.066}{\Character{58}}%
\TextAt{13}{371.518}{468.066}{K}%
\TextAt{9}{382.336}{468.066}{\Character{93}}%
\TextAt{13}{385.587}{468.066}{\Character{59}}%
\TextAt{13}{390.811}{468.066}{p}%
\TextAt{9}{396.707}{468.066}{\Character{41}}%
\TextAt{9}{405.781}{468.066}{\Character{61}}%
\TextAt{9}{419.422}{468.066}{1}%
\TextAt{8}{425.292}{468.066}{.}%
\TextAt{8}{435.883}{468.066}{We}%
\TextAt{8}{456.741}{468.066}{first}%
\TextAt{8}{481.581}{468.066}{show}%
\TextAt{8}{511.298}{468.066}{that}%
\TextAt{13}{79.226}{451.02}{e}%
\TextAt{9}{88.944}{451.02}{\Character{61}}%
\TextAt{9}{102.322}{451.02}{1}%
\TextAt{8}{112.66}{451.02}{implies}%
\TextAt{13}{152.957}{451.02}{L}%
\TextAt{9}{165.214}{451.02}{\Character{61}}%
\TextAt{13}{178.61}{451.02}{K}%
\TextAt{8}{189.428}{451.02}{.}%
\TextAt{8}{199.586}{451.02}{Let}%
\TextAt{13}{221.125}{451.02}{\Character{11}}%
\NativeMathAt{232.939}{451.02}{11.955}{\in}%
\TextAt{13}{245.201}{451.02}{L}%
\TextAt{18}{256.211}{451.02}{r}%
\TextAt{13}{268.554}{451.02}{K}%
\TextAt{8}{279.372}{451.02}{.}%
\TextAt{8}{289.53}{451.02}{Writing}%
\TextAt{13}{333.019}{451.02}{\Character{11}}%
\TextAt{9}{344.833}{451.02}{\Character{61}}%
\TextAt{13}{358.229}{451.02}{\Character{11}}%
\TextAt{17}{365.675}{449.227}{1}%
\TextAt{13}{370.407}{451.02}{\Character{59}}%
\TextAt{13}{375.631}{451.02}{\Character{58}}%
\TextAt{13}{380.868}{451.02}{\Character{58}}%
\TextAt{13}{386.092}{451.02}{\Character{58}}%
\TextAt{13}{391.329}{451.02}{\Character{59}}%
\TextAt{13}{396.565}{451.02}{\Character{11}}%
\TextAt{21}{404.073}{449.227}{m}%
\TextAt{8}{416.529}{451.02}{for}%
\TextAt{8}{434.984}{451.02}{the}%
\TextAt{8}{455.71}{451.02}{conjugates}%
\TextAt{8}{513.904}{451.02}{and}%
\TextAt{13}{79.226}{433.974}{a}%
\TextAt{9}{88.928}{433.974}{\Character{61}}%
\TextAt{8}{101.59}{433.974}{Tr}%
\TextAt{9}{114.597}{433.974}{\Character{40}}%
\TextAt{13}{119.149}{433.974}{\Character{11}}%
\TextAt{9}{126.671}{433.974}{\Character{41}}%
\TextAt{9}{134.765}{433.974}{\Character{61}}%
\TextAt{20}{147.442}{442.941}{\Character{80}}%
\TextAt{21}{160.062}{439.786}{m}%
\TextAt{21}{160.062}{430.487}{i}%
\TextAt{17}{162.945}{430.487}{\Character{61}1}%
\TextAt{13}{176.256}{433.974}{\Character{11}}%
\TextAt{21}{183.702}{432.181}{i}%
\TextAt{8}{187.083}{433.974}{,}%
\TextAt{8}{194.405}{433.974}{the}%
\TextAt{8}{214.7}{433.974}{element}%
\TextAt{13}{258.081}{433.974}{\Character{12}}%
\TextAt{9}{268.909}{433.974}{\Character{61}}%
\TextAt{13}{281.571}{433.974}{\Character{11}}%
\TextAt{16}{291.844}{433.974}{\Character{0}}%
\TextAt{17}{306.717}{438.681}{1}%
\Vector{q 0 G 0 g 0.478 w 233.089 -337.037 m 240.58 -337.037 l S Q}%
\TextAt{21}{305.089}{429.851}{m}%
\TextAt{13}{313.775}{433.974}{a}%
\NativeMathAt{323.477}{433.974}{11.955}{\in}%
\TextAt{13}{335.004}{433.974}{L}%
\TextAt{18}{345.72}{433.974}{r}%
\TextAt{13}{357.77}{433.974}{K}%
\TextAt{8}{372.623}{433.974}{has}%
\TextAt{8}{393.636}{433.974}{trace}%
\TextAt{8}{423.04}{433.974}{Tr}%
\TextAt{9}{436.048}{433.974}{\Character{40}}%
\TextAt{13}{440.6}{433.974}{\Character{12}}%
\TextAt{9}{447.871}{433.974}{\Character{41}}%
\TextAt{9}{455.965}{433.974}{\Character{61}}%
\TextAt{20}{468.642}{442.941}{\Character{80}}%
\TextAt{21}{481.262}{439.786}{m}%
\TextAt{21}{481.262}{430.487}{i}%
\TextAt{17}{484.145}{430.487}{\Character{61}1}%
\TextAt{13}{497.456}{433.974}{\Character{12}}%
\TextAt{21}{504.065}{432.181}{i}%
\TextAt{9}{511.003}{433.974}{\Character{61}}%
\TextAt{9}{523.652}{433.974}{0}%
\TextAt{8}{529.518}{433.974}{.}%
\TextAt{8}{79.226}{416.928}{Since}%
\TextAt{13}{109.951}{416.928}{v}%
\TextAt{9}{116.039}{416.928}{\Character{40}}%
\TextAt{13}{120.592}{416.928}{K}%
\TextAt{3}{131.409}{421.267}{\Character{3}}%
\TextAt{9}{136.142}{416.928}{\Character{41}}%
\TextAt{9}{144.284}{416.928}{\Character{61}}%
\TextAt{13}{157.011}{416.928}{w}%
\TextAt{9}{165.701}{416.928}{\Character{40}}%
\TextAt{13}{170.253}{416.928}{L}%
\TextAt{3}{178.217}{421.267}{\Character{3}}%
\TextAt{9}{182.95}{416.928}{\Character{41}}%
\TextAt{8}{187.502}{416.928}{,}%
\TextAt{8}{194.861}{416.928}{we}%
\TextAt{8}{212.255}{416.928}{may}%
\TextAt{8}{237.792}{416.928}{choose}%
\TextAt{8}{275.092}{416.928}{a}%
\TextAt{13}{285.015}{416.928}{b}%
\NativeMathAt{293.598}{416.928}{11.955}{\in}%
\TextAt{13}{305.175}{416.928}{K}%
\TextAt{3}{315.992}{421.267}{\Character{3}}%
\TextAt{8}{324.789}{416.928}{such}%
\TextAt{8}{351.354}{416.928}{that}%
\TextAt{13}{376.89}{416.928}{v}%
\TextAt{9}{382.979}{416.928}{\Character{40}}%
\TextAt{13}{387.531}{416.928}{b}%
\TextAt{9}{392.508}{416.928}{\Character{41}}%
\TextAt{9}{400.649}{416.928}{\Character{61}}%
\TextAt{13}{413.377}{416.928}{w}%
\TextAt{9}{422.067}{416.928}{\Character{40}}%
\TextAt{13}{426.619}{416.928}{\Character{12}}%
\TextAt{9}{433.89}{416.928}{\Character{41}}%
\TextAt{8}{442.508}{416.928}{and}%
\TextAt{8}{465.438}{416.928}{obtain}%
\TextAt{8}{502.033}{416.928}{a}%
\TextAt{8}{511.956}{416.928}{unit}%
\TextAt{13}{79.226}{399.882}{\Character{34}}%
\TextAt{9}{88.897}{399.882}{\Character{61}}%
\TextAt{13}{102.195}{399.882}{\Character{12}}%
\TextAt{9}{109.466}{399.882}{\Character{47}}%
\TextAt{13}{115.319}{399.882}{b}%
\NativeMathAt{124.489}{399.882}{11.955}{\in}%
\TextAt{13}{136.652}{399.882}{L}%
\TextAt{18}{147.622}{399.882}{r}%
\TextAt{13}{159.926}{399.882}{K}%
\TextAt{8}{175.153}{399.882}{with}%
\TextAt{8}{202.324}{399.882}{trace}%
\TextAt{20}{232.102}{408.849}{\Character{80}}%
\TextAt{21}{244.722}{405.694}{m}%
\TextAt{21}{244.722}{396.395}{i}%
\TextAt{17}{247.605}{396.395}{\Character{61}1}%
\TextAt{13}{260.916}{399.882}{\Character{34}}%
\TextAt{21}{266.394}{398.089}{i}%
\TextAt{9}{273.969}{399.882}{\Character{61}}%
\TextAt{9}{287.251}{399.882}{0}%
\TextAt{8}{293.119}{399.882}{.}%
\TextAt{8}{303.103}{399.882}{But}%
\TextAt{8}{326.855}{399.882}{the}%
\TextAt{8}{347.523}{399.882}{conjugates}%
\TextAt{13}{405.659}{399.882}{\Character{34}}%
\TextAt{21}{411.137}{398.089}{i}%
\TextAt{8}{418.927}{399.882}{have}%
\TextAt{8}{446.434}{399.882}{the}%
\TextAt{8}{467.102}{399.882}{same}%
\TextAt{8}{496.94}{399.882}{residue}%
\TextAt{8}{79.226}{382.836}{classes}%
\TextAt{9}{117.644}{382.836}{\Character{22}}%
\TextAt{13}{116.856}{382.836}{\Character{34}}%
\TextAt{21}{122.334}{381.043}{i}%
\TextAt{8}{129.991}{382.836}{in}%
\TextAt{13}{144.022}{382.836}{\Character{20}}%
\TextAt{8}{150.76}{382.836}{,}%
\TextAt{8}{158.382}{382.836}{because}%
\TextAt{13}{202.061}{382.836}{\Character{21}}%
\TextAt{9}{212.855}{382.836}{\Character{61}}%
\TextAt{13}{225.924}{382.836}{\Character{20}}%
\TextAt{8}{232.663}{382.836}{.}%
\TextAt{8}{242.245}{382.836}{Hence}%
\TextAt{9}{277.401}{382.836}{0}%
\TextAt{9}{287.204}{382.836}{\Character{61}}%
\TextAt{20}{300.288}{391.803}{\Character{80}}%
\TextAt{21}{312.907}{388.648}{m}%
\TextAt{21}{312.907}{379.349}{i}%
\TextAt{17}{315.791}{379.349}{\Character{61}1}%
\TextAt{9}{329.89}{382.836}{\Character{22}}%
\TextAt{13}{329.102}{382.836}{\Character{34}}%
\TextAt{21}{334.58}{381.043}{i}%
\TextAt{9}{341.926}{382.836}{\Character{61}}%
\TextAt{13}{354.996}{382.836}{m}%
\TextAt{9}{366.023}{382.836}{\Character{22}}%
\TextAt{13}{365.235}{382.836}{\Character{34}}%
\TextAt{8}{370.713}{382.836}{,}%
\TextAt{8}{378.334}{382.836}{and}%
\TextAt{8}{401.475}{382.836}{thus}%
\TextAt{13}{427.604}{382.836}{m}%
\NativeMathAt{441.809}{382.836}{11.955}{\equiv}%
\TextAt{9}{455.072}{382.836}{0}%
\TextAt{8}{464.246}{382.836}{mod}%
\TextAt{13}{489.684}{382.836}{p}%
\TextAt{8}{495.559}{382.836}{,}%
\TextAt{8}{503.18}{382.836}{which}%
\TextAt{8}{79.226}{365.79}{contradicts}%
\TextAt{13}{139.456}{365.79}{p}%
\TextAt{18}{148.652}{365.79}{\Character{45}}%
\TextAt{9}{155.294}{365.79}{\Character{91}}%
\TextAt{13}{158.546}{365.79}{L}%
\TextAt{9}{169.831}{365.79}{\Character{58}}%
\TextAt{13}{176.403}{365.79}{K}%
\TextAt{9}{187.221}{365.79}{\Character{93}}%
\TextAt{8}{194.37}{365.79}{and}%
\TextAt{13}{217.133}{365.79}{m}%
\TextAt{16}{230.693}{365.79}{\Character{106}}%
\TextAt{9}{237.335}{365.79}{\Character{91}}%
\TextAt{13}{240.586}{365.79}{L}%
\TextAt{9}{251.872}{365.79}{\Character{58}}%
\TextAt{13}{258.444}{365.79}{K}%
\TextAt{9}{269.262}{365.79}{\Character{93}}%
\TextAt{8}{272.513}{365.79}{.}%
\ErrataAt{97.15}{345.798}{437.85}{Choose $\omega_1,\ldots,\omega_r\in w(L^*)$ lifting generators of the cyclic direct factors of the quotient}%
\ErrataAt{79.226}{328.752}{458.774}{$w(L^*)/v(K^*)$. Let $m_i$ be the order of $\omega_i$ modulo $v(K^*)$. Since $w(L^*)=\frac1n v(N_{L\mid K}(L^*))\subseteq\frac1n v(K^*)$,}%
\TextAt{8}{79.226}{311.663}{where}%
\TextAt{13}{113.329}{311.663}{n}%
\TextAt{9}{124.137}{311.663}{\Character{61}}%
\TextAt{9}{137.049}{311.663}{\Character{91}}%
\TextAt{13}{140.314}{311.663}{L}%
\TextAt{9}{152.099}{311.663}{\Character{58}}%
\TextAt{13}{159.172}{311.663}{K}%
\TextAt{9}{169.989}{311.663}{\Character{93}}%
\TextAt{8}{173.241}{311.663}{,}%
\TextAt{8}{180.757}{311.663}{we}%
\TextAt{8}{198.278}{311.663}{have}%
\TextAt{13}{225.566}{311.663}{m}%
\TextAt{21}{235.805}{309.87}{i}%
\TextAt{16}{243.007}{311.663}{\Character{106}}%
\TextAt{13}{250.149}{311.663}{n}%
\TextAt{8}{257.136}{311.663}{,}%
\TextAt{8}{264.652}{311.663}{so}%
\TextAt{8}{279.316}{311.663}{that}%
\TextAt{9}{304.978}{311.663}{\Character{40}}%
\TextAt{13}{309.53}{311.663}{m}%
\TextAt{21}{319.769}{309.87}{i}%
\TextAt{13}{323.151}{311.663}{\Character{59}}%
\TextAt{13}{328.375}{311.663}{p}%
\TextAt{9}{334.27}{311.663}{\Character{41}}%
\TextAt{9}{342.627}{311.663}{\Character{61}}%
\TextAt{9}{355.55}{311.663}{1}%
\TextAt{8}{361.421}{311.663}{.}%
\TextAt{8}{370.75}{311.663}{Let}%
\TextAt{13}{392.012}{311.663}{\Character{13}}%
\TextAt{21}{398.084}{309.87}{i}%
\NativeMathAt{405.286}{311.663}{11.955}{\in}%
\TextAt{13}{417.077}{311.663}{L}%
\TextAt{3}{425.042}{316.002}{\Character{3}}%
\TextAt{8}{433.965}{311.663}{be}%
\TextAt{8}{450.182}{311.663}{an}%
\TextAt{8}{466.735}{311.663}{element}%
\TextAt{8}{510.27}{311.663}{such}%
\TextAt{8}{79.226}{294.618}{that}%
\TextAt{13}{104.963}{294.618}{w}%
\TextAt{9}{113.653}{294.618}{\Character{40}}%
\TextAt{13}{118.205}{294.618}{\Character{13}}%
\TextAt{21}{124.277}{292.824}{i}%
\TextAt{9}{127.658}{294.618}{\Character{41}}%
\TextAt{9}{136.146}{294.618}{\Character{61}}%
\TextAt{13}{149.211}{294.618}{\Character{33}}%
\TextAt{21}{156.501}{292.824}{i}%
\TextAt{8}{159.882}{294.618}{.}%
\TextAt{8}{169.435}{294.618}{Then}%
\TextAt{13}{200.36}{294.618}{w}%
\TextAt{9}{209.05}{294.618}{\Character{40}}%
\TextAt{13}{213.602}{294.618}{\Character{13}}%
\TextAt{21}{220.324}{299.95}{m}%
\TextAt{22}{227.815}{298.735}{i}%
\TextAt{21}{219.674}{291.554}{i}%
\TextAt{9}{231.475}{294.618}{\Character{41}}%
\TextAt{9}{239.963}{294.618}{\Character{61}}%
\TextAt{13}{253.028}{294.618}{v}%
\TextAt{9}{259.116}{294.618}{\Character{40}}%
\TextAt{13}{263.669}{294.618}{c}%
\TextAt{21}{268.707}{292.824}{i}%
\TextAt{9}{272.088}{294.618}{\Character{41}}%
\TextAt{8}{276.64}{294.618}{,}%
\TextAt{8}{284.25}{294.618}{with}%
\TextAt{13}{311.278}{294.618}{c}%
\TextAt{21}{316.316}{292.824}{i}%
\NativeMathAt{323.645}{294.618}{11.955}{\in}%
\TextAt{13}{335.564}{294.618}{K}%
\TextAt{8}{346.381}{294.618}{,}%
\TextAt{8}{353.991}{294.618}{so}%
\TextAt{8}{368.729}{294.618}{that}%
\TextAt{13}{394.466}{294.618}{\Character{13}}%
\TextAt{21}{401.188}{299.95}{m}%
\TextAt{22}{408.679}{298.735}{i}%
\TextAt{21}{400.538}{291.554}{i}%
\TextAt{9}{416.287}{294.618}{\Character{61}}%
\TextAt{13}{429.339}{294.618}{c}%
\TextAt{21}{434.377}{292.824}{i}%
\TextAt{13}{437.759}{294.618}{\Character{34}}%
\TextAt{21}{443.237}{292.824}{i}%
\TextAt{8}{450.884}{294.618}{for}%
\TextAt{8}{469.137}{294.618}{some}%
\TextAt{8}{498.831}{294.618}{unit}%
\TextAt{13}{523.91}{294.618}{\Character{34}}%
\TextAt{21}{529.388}{292.824}{i}%
\TextAt{8}{79.226}{277.572}{in}%
\TextAt{13}{93.191}{277.572}{L}%
\TextAt{8}{101.155}{277.572}{.}%
\TextAt{8}{110.54}{277.572}{As}%
\TextAt{13}{128.139}{277.572}{\Character{21}}%
\TextAt{9}{138.82}{277.572}{\Character{61}}%
\TextAt{13}{151.777}{277.572}{\Character{20}}%
\TextAt{8}{162.725}{277.572}{we}%
\TextAt{8}{180.264}{277.572}{may}%
\TextAt{8}{205.945}{277.572}{write}%
\TextAt{13}{236.169}{277.572}{\Character{34}}%
\TextAt{21}{241.647}{275.778}{i}%
\TextAt{9}{248.881}{277.572}{\Character{61}}%
\TextAt{13}{261.838}{277.572}{b}%
\TextAt{21}{266.815}{275.778}{i}%
\TextAt{13}{270.196}{277.572}{u}%
\TextAt{21}{276.859}{275.778}{i}%
\TextAt{8}{280.24}{277.572}{,}%
\TextAt{8}{287.779}{277.572}{where}%
\TextAt{13}{321.9}{277.572}{b}%
\TextAt{21}{326.877}{275.778}{i}%
\NativeMathAt{334.111}{277.572}{11.955}{\in}%
\TextAt{13}{345.934}{277.572}{K}%
\TextAt{8}{360.961}{277.572}{and}%
\TextAt{13}{384.035}{277.572}{u}%
\TextAt{21}{390.698}{275.778}{i}%
\TextAt{8}{398.289}{277.572}{is}%
\TextAt{8}{410.364}{277.572}{a}%
\TextAt{8}{420.432}{277.572}{unit}%
\TextAt{8}{445.455}{277.572}{in}%
\TextAt{13}{459.42}{277.572}{L}%
\TextAt{8}{471.594}{277.572}{which}%
\TextAt{8}{505.392}{277.572}{tends}%
\TextAt{8}{79.226}{260.526}{to}%
\TextAt{8}{93.846}{260.526}{1}%
\TextAt{8}{103.911}{260.526}{in}%
\TextAt{13}{117.874}{260.526}{\Character{21}}%
\TextAt{8}{124.703}{260.526}{.}%
\TextAt{8}{134.08}{260.526}{By}%
\TextAt{8}{152.753}{260.526}{Hensel’s}%
\TextAt{8}{198.373}{260.526}{lemma}%
\TextAt{8}{236.401}{260.526}{the}%
\TextAt{8}{256.867}{260.526}{equation}%
\TextAt{13}{304.986}{260.526}{x}%
\TextAt{21}{311.638}{264.864}{m}%
\TextAt{22}{319.129}{263.649}{i}%
\TextAt{16}{325.656}{260.526}{\Character{0}}%
\TextAt{13}{337.823}{260.526}{u}%
\TextAt{21}{344.485}{258.732}{i}%
\TextAt{9}{351.715}{260.526}{\Character{61}}%
\TextAt{9}{364.651}{260.526}{0}%
\TextAt{8}{374.729}{260.526}{has}%
\TextAt{8}{395.912}{260.526}{a}%
\TextAt{8}{405.977}{260.526}{solution}%
\TextAt{13}{450.581}{260.526}{\Character{12}}%
\TextAt{21}{457.189}{258.732}{i}%
\NativeMathAt{464.419}{260.526}{11.955}{\in}%
\TextAt{13}{476.238}{260.526}{L}%
\TextAt{8}{484.203}{260.526}{.}%
\TextAt{8}{493.58}{260.526}{Putting}%
\TextAt{13}{79.226}{243.48}{\Character{11}}%
\TextAt{21}{86.672}{241.686}{i}%
\TextAt{9}{93.374}{243.48}{\Character{61}}%
\TextAt{13}{105.799}{243.48}{\Character{13}}%
\TextAt{21}{111.871}{241.686}{i}%
\TextAt{13}{115.253}{243.48}{\Character{12}}%
\TextAt{3}{122.524}{248.428}{\Character{0}}%
\TextAt{17}{129.11}{248.428}{1}%
\TextAt{21}{121.861}{240.416}{i}%
\NativeMathAt{137.163}{243.48}{11.955}{\in}%
\TextAt{13}{148.454}{243.48}{L}%
\TextAt{8}{156.419}{243.48}{,}%
\TextAt{8}{163.568}{243.48}{we}%
\TextAt{8}{180.796}{243.48}{find}%
\TextAt{13}{204.204}{243.48}{w}%
\TextAt{9}{212.893}{243.48}{\Character{40}}%
\TextAt{13}{217.446}{243.48}{\Character{11}}%
\TextAt{21}{224.891}{241.686}{i}%
\TextAt{9}{228.272}{243.48}{\Character{41}}%
\TextAt{9}{236.127}{243.48}{\Character{61}}%
\TextAt{13}{248.571}{243.48}{\Character{33}}%
\TextAt{21}{255.86}{241.686}{i}%
\TextAt{8}{263.139}{243.48}{and}%
\TextAt{13}{239.441}{216.856}{\Character{11}}%
\TextAt{21}{246.962}{222.189}{m}%
\TextAt{22}{254.453}{220.974}{i}%
\TextAt{21}{246.886}{213.792}{i}%
\TextAt{9}{261.434}{216.856}{\Character{61}}%
\TextAt{13}{273.859}{216.856}{a}%
\TextAt{21}{280.004}{215.063}{i}%
\TextAt{13}{283.385}{216.856}{\Character{59}}%
\TextAt{13}{312.52}{216.856}{i}%
\TextAt{9}{319.854}{216.856}{\Character{61}}%
\TextAt{9}{332.264}{216.856}{1}%
\TextAt{13}{338.132}{216.856}{\Character{59}}%
\TextAt{13}{343.356}{216.856}{\Character{58}}%
\TextAt{13}{348.593}{216.856}{\Character{58}}%
\TextAt{13}{353.829}{216.856}{\Character{58}}%
\TextAt{13}{359.054}{216.856}{\Character{59}}%
\TextAt{13}{364.29}{216.856}{r}%
\TextAt{13}{369.251}{216.856}{\Character{59}}%
\TextAt{8}{79.226}{190.232}{where}%
\TextAt{13}{113.11}{190.232}{a}%
\TextAt{21}{119.255}{188.439}{i}%
\TextAt{9}{126.085}{190.232}{\Character{61}}%
\TextAt{13}{138.639}{190.232}{c}%
\TextAt{21}{143.677}{188.439}{i}%
\TextAt{13}{147.058}{190.232}{b}%
\TextAt{21}{152.035}{188.439}{i}%
\NativeMathAt{158.865}{190.232}{11.955}{\in}%
\TextAt{13}{170.284}{190.232}{K}%
\TextAt{8}{181.102}{190.232}{,}%
\TextAt{8}{188.345}{190.232}{i.e.,}%
\TextAt{8}{210.544}{190.232}{we}%
\TextAt{8}{227.847}{190.232}{have}%
\TextAt{13}{254.916}{190.232}{K}%
\TextAt{9}{265.734}{190.232}{\Character{40}}%
\TextAt{22}{273.607}{193.177}{m}%
\TextAt{24}{280.145}{191.384}{1}%
\TextAt{16}{277.655}{197.946}{\Character{112}}%
\Vector{q 0 G 0 g 0.478 w 215.617 -575.815 m 226.495 -575.815 l S Q}%
\TextAt{13}{287.617}{190.232}{a}%
\TextAt{17}{293.762}{188.439}{1}%
\TextAt{13}{298.495}{190.232}{\Character{59}}%
\TextAt{13}{303.719}{190.232}{\Character{58}}%
\TextAt{13}{308.956}{190.232}{\Character{58}}%
\TextAt{13}{314.18}{190.232}{\Character{58}}%
\TextAt{13}{319.417}{190.232}{\Character{59}}%
\TextAt{22}{328.036}{193.177}{m}%
\TextAt{22}{334.574}{192.181}{r}%
\TextAt{16}{332.09}{197.946}{\Character{112}}%
\Vector{q 0 G 0 g 0.478 w 270.053 -575.815 m 280.753 -575.815 l S Q}%
\TextAt{13}{342.053}{190.232}{a}%
\TextAt{21}{348.198}{188.439}{r}%
\TextAt{9}{352.753}{190.232}{\Character{41}}%
\NativeMathAt{360.754}{190.232}{11.955}{\subseteq}%
\TextAt{13}{373.501}{190.232}{L}%
\TextAt{8}{381.466}{190.232}{.}%
\TextAt{8}{390.139}{190.232}{By}%
\TextAt{8}{408.578}{190.232}{construction,}%
\TextAt{8}{478.98}{190.232}{both}%
\TextAt{8}{506.695}{190.232}{fields}%
\TextAt{8}{79.226}{173.186}{have}%
\TextAt{8}{106.25}{173.186}{the}%
\TextAt{8}{126.435}{173.186}{same}%
\TextAt{8}{155.79}{173.186}{value}%
\TextAt{8}{186.066}{173.186}{group}%
\TextAt{8}{219.271}{173.186}{and}%
\TextAt{8}{242.062}{173.186}{the}%
\TextAt{8}{262.248}{173.186}{same}%
\TextAt{8}{291.603}{173.186}{residue}%
\TextAt{8}{331.359}{173.186}{class}%
\TextAt{8}{358.825}{173.186}{field.}%
\TextAt{8}{388.82}{173.186}{So,}%
\TextAt{8}{408.367}{173.186}{by}%
\TextAt{8}{424.655}{173.186}{what}%
\TextAt{8}{453.951}{173.186}{we}%
\TextAt{8}{471.207}{173.186}{proved}%
\TextAt{8}{509.29}{173.186}{first,}%
\TextAt{8}{79.226}{156.14}{we}%
\TextAt{8}{96.454}{156.14}{have}%
\TextAt{13}{240.395}{139.094}{L}%
\TextAt{9}{251.68}{139.094}{\Character{61}}%
\TextAt{13}{264.106}{139.094}{K}%
\TextAt{20}{274.924}{148.778}{\Character{0}}%
\TextAt{22}{283.724}{142.739}{m}%
\TextAt{24}{290.262}{140.946}{1}%
\TextAt{16}{287.771}{147.392}{\Character{112}}%
\Vector{q 0 G 0 g 0.478 w 225.734 -626.369 m 236.611 -626.369 l S Q}%
\TextAt{13}{297.734}{139.094}{a}%
\TextAt{17}{303.879}{137.301}{1}%
\TextAt{13}{308.611}{139.094}{\Character{59}}%
\TextAt{13}{313.835}{139.094}{\Character{58}}%
\TextAt{13}{319.072}{139.094}{\Character{58}}%
\TextAt{13}{324.308}{139.094}{\Character{58}}%
\TextAt{13}{329.533}{139.094}{\Character{59}}%
\TextAt{22}{338.153}{142.739}{m}%
\TextAt{22}{344.691}{141.743}{r}%
\TextAt{16}{342.207}{147.392}{\Character{112}}%
\Vector{q 0 G 0 g 0.478 w 280.17 -626.369 m 290.869 -626.369 l S Q}%
\TextAt{13}{352.17}{139.094}{a}%
\TextAt{21}{358.315}{137.301}{r}%
\TextAt{20}{362.869}{148.778}{\Character{1}}%
\TextAt{13}{368.349}{139.094}{\Character{58}}%
\TextAt{8}{79.226}{116.591}{The}%
\TextAt{8}{104.382}{116.591}{inequality}%
\TextAt{9}{159.796}{116.591}{\Character{91}}%
\TextAt{13}{163.048}{116.591}{L}%
\TextAt{9}{176.211}{116.591}{\Character{58}}%
\TextAt{13}{184.661}{116.591}{K}%
\TextAt{9}{195.479}{116.591}{\Character{93}}%
\NativeMathAt{203.929}{116.591}{11.955}{\leq}%
\TextAt{13}{218.427}{116.591}{e}%
\TextAt{8}{228.851}{116.591}{and}%
\TextAt{8}{252.716}{116.591}{thus,}%
\TextAt{8}{283.097}{116.591}{in}%
\TextAt{8}{297.851}{116.591}{view}%
\TextAt{8}{325.936}{116.591}{of}%
\TextAt{8}{340.368}{116.591}{(6.8),}%
\TextAt{8}{372.972}{116.591}{the}%
\TextAt{8}{394.231}{116.591}{equality}%
\TextAt{9}{439.89}{116.591}{\Character{91}}%
\TextAt{13}{443.141}{116.591}{L}%
\TextAt{9}{456.305}{116.591}{\Character{58}}%
\TextAt{13}{464.755}{116.591}{K}%
\TextAt{9}{475.573}{116.591}{\Character{93}}%
\TextAt{9}{484.001}{116.591}{\Character{61}}%
\TextAt{13}{498.326}{116.591}{e}%
\TextAt{8}{503.752}{116.591}{,}%
\TextAt{8}{512.278}{116.591}{now}%
\TextAt{8}{79.226}{99.545}{follows}%
\TextAt{8}{117.662}{99.545}{by}%
\TextAt{8}{133.921}{99.545}{induction}%
\TextAt{8}{185.95}{99.545}{on}%
\TextAt{13}{202.209}{99.545}{r}%
\TextAt{8}{207.809}{99.545}{.}%
\TextAt{8}{216.258}{99.545}{If}%
\TextAt{13}{227.95}{99.545}{L}%
\TextAt{17}{235.914}{97.752}{1}%
\TextAt{9}{243.968}{99.545}{\Character{61}}%
\TextAt{13}{256.393}{99.545}{K}%
\TextAt{9}{267.211}{99.545}{\Character{40}}%
\TextAt{22}{275.084}{102.49}{m}%
\TextAt{24}{281.622}{100.697}{1}%
\TextAt{16}{279.131}{107.26}{\Character{112}}%
\Vector{q 0 G 0 g 0.478 w 217.094 -666.501 m 227.971 -666.501 l S Q}%
\TextAt{13}{289.094}{99.545}{a}%
\TextAt{17}{295.239}{97.752}{1}%
\TextAt{9}{299.971}{99.545}{\Character{41}}%
\TextAt{8}{304.523}{99.545}{,}%
\TextAt{8}{311.673}{99.545}{then}%
\TextAt{13}{338.333}{99.545}{\Character{33}}%
\TextAt{17}{345.622}{97.752}{1}%
\NativeMathAt{353.675}{99.545}{11.955}{\in}%
\TextAt{13}{364.966}{99.545}{w}%
\TextAt{9}{373.656}{99.545}{\Character{40}}%
\TextAt{13}{378.208}{99.545}{L}%
\TextAt{3}{386.172}{103.884}{\Character{3}}%
\TextAt{17}{386.172}{96.59}{1}%
\TextAt{9}{390.905}{99.545}{\Character{41}}%
\TextAt{8}{399.354}{99.545}{yields}%
\TextAt{13}{188.998}{72.921}{e}%
\TextAt{9}{194.423}{72.921}{\Character{40}}%
\TextAt{13}{198.976}{72.921}{L}%
\TextAt{17}{206.94}{71.128}{1}%
\TextAt{16}{214.993}{72.921}{\Character{106}}%
\TextAt{13}{221.635}{72.921}{K}%
\TextAt{9}{232.453}{72.921}{\Character{41}}%
\TextAt{9}{240.308}{72.921}{\Character{61}}%
\TextAt{9}{252.729}{72.921}{\Character{40}}%
\TextAt{13}{257.304}{72.921}{w}%
\TextAt{9}{265.993}{72.921}{\Character{40}}%
\TextAt{13}{270.545}{72.921}{L}%
\TextAt{3}{278.51}{77.858}{\Character{3}}%
\TextAt{17}{278.51}{69.966}{1}%
\TextAt{9}{283.242}{72.921}{\Character{41}}%
\TextAt{9}{291.097}{72.921}{\Character{58}}%
\TextAt{13}{297.688}{72.921}{v}%
\TextAt{9}{303.776}{72.921}{\Character{40}}%
\TextAt{13}{308.328}{72.921}{K}%
\TextAt{3}{319.146}{77.858}{\Character{3}}%
\TextAt{9}{323.878}{72.921}{\Character{41}\Character{41}}%
\NativeMathAt{336.304}{72.921}{11.955}{\geq}%
\TextAt{13}{348.923}{72.921}{m}%
\TextAt{17}{359.163}{71.128}{1}%
\NativeMathAt{367.216}{72.921}{11.955}{\geq}%
\TextAt{9}{379.835}{72.921}{\Character{91}}%
\TextAt{13}{383.087}{72.921}{L}%
\TextAt{17}{391.051}{71.128}{1}%
\TextAt{9}{399.104}{72.921}{\Character{58}}%
\TextAt{13}{405.677}{72.921}{K}%
\TextAt{9}{416.494}{72.921}{\Character{93}}%
\TextAt{13}{419.746}{72.921}{\Character{58}}%
\TextAt{8}{297.211}{43.033}{142}%
\Destination{78.23}{809.815}{nav.page.142}%
\Destination{79.36}{733.755}{nav.print.155}%
\Destination{79.36}{656.255}{nav.number.2.7.7}%
\Hotlink{339.868}{114.355}{25.078}{11.047}{nav.number.2.6.8}%
\Destination{79.226}{361}{errata.217}%
}

% Page 155 (body page 143)
\Page{%
\TextAt{2}{79.226}{796.943}{II.}%
\TextAt{2}{94.073}{796.943}{The}%
\TextAt{2}{116.492}{796.943}{Theory}%
\TextAt{2}{154.401}{796.943}{of}%
\TextAt{2}{166.826}{796.943}{Valuations}%
\Vector{q 0 G 0 g 0.299 w 7.226 18.729 m 460.769 18.729 l S Q}%
\TextAt{8}{79.226}{761.124}{Also}%
\TextAt{13}{106.594}{761.124}{e}%
\TextAt{9}{112.02}{761.124}{\Character{40}}%
\TextAt{13}{116.572}{761.124}{L}%
\TextAt{16}{129.513}{761.124}{\Character{106}}%
\TextAt{13}{137.81}{761.124}{L}%
\TextAt{17}{145.774}{759.33}{1}%
\TextAt{9}{150.507}{761.124}{\Character{41}}%
\NativeMathAt{160.035}{761.124}{11.955}{\geq}%
\TextAt{9}{174.31}{761.124}{\Character{91}}%
\TextAt{13}{177.561}{761.124}{L}%
\TextAt{9}{190.502}{761.124}{\Character{58}}%
\TextAt{13}{198.73}{761.124}{L}%
\TextAt{17}{206.694}{759.33}{1}%
\TextAt{9}{211.427}{761.124}{\Character{93}}%
\TextAt{8}{214.678}{761.124}{,}%
\TextAt{8}{223.042}{761.124}{because}%
\TextAt{13}{267.315}{761.124}{w}%
\TextAt{9}{276.004}{761.124}{\Character{40}}%
\TextAt{13}{280.556}{761.124}{L}%
\TextAt{3}{288.521}{765.462}{\Character{3}}%
\TextAt{9}{293.253}{761.124}{\Character{41}\Character{47}}%
\TextAt{13}{303.658}{761.124}{w}%
\TextAt{9}{312.348}{761.124}{\Character{40}}%
\TextAt{13}{316.9}{761.124}{L}%
\TextAt{3}{324.865}{765.462}{\Character{3}}%
\TextAt{17}{324.865}{758.168}{1}%
\TextAt{9}{329.597}{761.124}{\Character{41}}%
\TextAt{8}{339.018}{761.124}{is}%
\TextAt{8}{351.753}{761.124}{generated}%
\TextAt{8}{406.057}{761.124}{by}%
\TextAt{8}{423.287}{761.124}{the}%
\TextAt{8}{444.415}{761.124}{residue}%
\TextAt{8}{485.113}{761.124}{classes}%
\TextAt{8}{523.337}{761.124}{of}%
\TextAt{13}{79.226}{744.078}{\Character{33}}%
\TextAt{17}{86.515}{742.284}{2}%
\TextAt{13}{91.248}{744.078}{\Character{59}}%
\TextAt{13}{96.472}{744.078}{\Character{58}}%
\TextAt{13}{101.709}{744.078}{\Character{58}}%
\TextAt{13}{106.933}{744.078}{\Character{58}}%
\TextAt{13}{112.17}{744.078}{\Character{59}}%
\TextAt{13}{117.406}{744.078}{\Character{33}}%
\TextAt{21}{124.758}{742.284}{r}%
\TextAt{8}{129.312}{744.078}{.}%
\TextAt{8}{137.761}{744.078}{Thus}%
\TextAt{13}{175.288}{715.077}{e}%
\TextAt{9}{184.035}{715.077}{\Character{61}}%
\TextAt{13}{196.46}{715.077}{e}%
\TextAt{9}{201.886}{715.077}{\Character{40}}%
\TextAt{13}{206.438}{715.077}{L}%
\TextAt{16}{217.723}{715.077}{\Character{106}}%
\TextAt{13}{224.365}{715.077}{L}%
\TextAt{17}{232.329}{713.283}{1}%
\TextAt{9}{237.062}{715.077}{\Character{41}}%
\TextAt{13}{241.614}{715.077}{e}%
\TextAt{9}{247.04}{715.077}{\Character{40}}%
\TextAt{13}{251.592}{715.077}{L}%
\TextAt{17}{259.556}{713.283}{1}%
\TextAt{16}{267.609}{715.077}{\Character{106}}%
\TextAt{13}{274.251}{715.077}{K}%
\TextAt{9}{285.069}{715.077}{\Character{41}}%
\NativeMathAt{292.942}{715.077}{11.955}{\geq}%
\TextAt{9}{305.561}{715.077}{\Character{91}}%
\TextAt{13}{308.813}{715.077}{L}%
\TextAt{9}{320.098}{715.077}{\Character{58}}%
\TextAt{13}{326.671}{715.077}{L}%
\TextAt{17}{334.635}{713.283}{1}%
\TextAt{9}{339.368}{715.077}{\Character{93}\Character{91}}%
\TextAt{13}{345.871}{715.077}{L}%
\TextAt{17}{353.835}{713.283}{1}%
\TextAt{9}{361.889}{715.077}{\Character{58}}%
\TextAt{13}{368.461}{715.077}{K}%
\TextAt{9}{379.279}{715.077}{\Character{93}}%
\TextAt{9}{385.83}{715.077}{\Character{61}}%
\TextAt{9}{398.252}{715.077}{\Character{91}}%
\TextAt{13}{401.528}{715.077}{L}%
\TextAt{9}{412.814}{715.077}{\Character{58}}%
\TextAt{13}{419.386}{715.077}{K}%
\TextAt{9}{430.204}{715.077}{\Character{93}}%
\TextAt{13}{433.456}{715.077}{\Character{58}}%
\TextAt{8}{79.226}{686.075}{In}%
\TextAt{8}{93.192}{686.075}{order}%
\TextAt{8}{123.106}{686.075}{to}%
\TextAt{8}{136.761}{686.075}{prove}%
\TextAt{8}{167.656}{686.075}{that}%
\TextAt{8}{192.37}{686.075}{an}%
\TextAt{8}{207.973}{686.075}{extension}%
\TextAt{13}{259.084}{686.075}{L}%
\TextAt{9}{270.369}{686.075}{\Character{61}}%
\TextAt{13}{282.795}{686.075}{K}%
\TextAt{9}{293.613}{686.075}{\Character{40}}%
\TextAt{22}{301.486}{689.02}{m}%
\TextAt{24}{308.024}{687.227}{1}%
\TextAt{16}{305.533}{693.79}{\Character{112}}%
\Vector{q 0 G 0 g 0.478 w 243.496 -79.971 m 254.373 -79.971 l S Q}%
\TextAt{13}{315.496}{686.075}{a}%
\TextAt{17}{321.641}{684.282}{1}%
\TextAt{13}{326.373}{686.075}{\Character{59}}%
\TextAt{13}{331.597}{686.075}{\Character{58}}%
\TextAt{13}{336.834}{686.075}{\Character{58}}%
\TextAt{13}{342.058}{686.075}{\Character{58}}%
\TextAt{13}{347.295}{686.075}{\Character{59}}%
\TextAt{22}{355.915}{689.02}{m}%
\TextAt{22}{362.453}{688.024}{r}%
\TextAt{16}{359.969}{693.79}{\Character{112}}%
\Vector{q 0 G 0 g 0.478 w 297.932 -79.971 m 308.631 -79.971 l S Q}%
\TextAt{13}{369.932}{686.075}{a}%
\TextAt{21}{376.077}{684.282}{r}%
\TextAt{9}{380.631}{686.075}{\Character{41}}%
\TextAt{8}{388.426}{686.075}{is}%
\TextAt{8}{399.535}{686.075}{tamely}%
\TextAt{8}{437.578}{686.075}{ramified,}%
\TextAt{8}{485.831}{686.075}{it}%
\TextAt{8}{496.88}{686.075}{suffices}%
\TextAt{8}{79.226}{669.029}{to}%
\TextAt{8}{94.377}{669.029}{look}%
\TextAt{8}{120.586}{669.029}{at}%
\TextAt{8}{135.737}{669.029}{the}%
\TextAt{8}{156.734}{669.029}{case}%
\TextAt{13}{182.346}{669.029}{r}%
\TextAt{9}{192.7}{669.029}{\Character{61}}%
\TextAt{9}{206.544}{669.029}{1}%
\TextAt{8}{212.41}{669.029}{,}%
\TextAt{8}{220.61}{669.029}{i.e.,}%
\TextAt{13}{243.766}{669.029}{L}%
\TextAt{9}{256.484}{669.029}{\Character{61}}%
\TextAt{13}{270.342}{669.029}{K}%
\TextAt{9}{281.159}{669.029}{\Character{40}}%
\TextAt{22}{289.032}{673.05}{m}%
\TextAt{16}{288.929}{677.641}{\Character{112}}%
\Vector{q 0 G 0 g 0.478 w 226.891 -96.12 m 233.036 -96.12 l S Q}%
\TextAt{13}{298.891}{669.029}{a}%
\TextAt{9}{305.036}{669.029}{\Character{41}}%
\TextAt{8}{309.589}{669.029}{,}%
\TextAt{8}{317.789}{669.029}{where}%
\TextAt{9}{352.438}{669.029}{\Character{40}}%
\TextAt{13}{356.991}{669.029}{m\Character{59}}%
\TextAt{13}{372.449}{669.029}{p}%
\TextAt{9}{378.349}{669.029}{\Character{41}}%
\TextAt{9}{387.638}{669.029}{\Character{61}}%
\TextAt{9}{401.494}{669.029}{1}%
\TextAt{8}{407.366}{669.029}{.}%
\TextAt{8}{418.336}{669.029}{The}%
\TextAt{8}{443.23}{669.029}{general}%
\TextAt{8}{484.395}{669.029}{case}%
\TextAt{8}{510.007}{669.029}{then}%
\TextAt{8}{79.226}{651.983}{follows}%
\TextAt{8}{116.845}{651.983}{by}%
\TextAt{8}{132.287}{651.983}{induction.}%
\TextAt{8}{188.594}{651.983}{We}%
\TextAt{8}{207.922}{651.983}{may}%
\TextAt{8}{232.473}{651.983}{assume}%
\TextAt{8}{272.101}{651.983}{without}%
\TextAt{8}{314.86}{651.983}{loss}%
\TextAt{8}{336.28}{651.983}{of}%
\TextAt{8}{348.792}{651.983}{generality}%
\TextAt{8}{401.965}{651.983}{that}%
\TextAt{13}{426.517}{651.983}{\Character{20}}%
\TextAt{8}{436.335}{651.983}{is}%
\TextAt{8}{447.282}{651.983}{separably}%
\TextAt{8}{498.889}{651.983}{closed.}%
\TextAt{8}{79.226}{634.938}{This}%
\TextAt{8}{106.657}{634.938}{is}%
\TextAt{8}{119.133}{634.938}{seen}%
\TextAt{8}{145.261}{634.938}{by}%
\TextAt{8}{162.231}{634.938}{passing}%
\TextAt{8}{204.044}{634.938}{to}%
\TextAt{8}{219.066}{634.938}{the}%
\TextAt{8}{239.934}{634.938}{maximal}%
\TextAt{8}{288.454}{634.938}{unramified}%
\TextAt{8}{347.698}{634.938}{extension}%
\TextAt{13}{400.175}{634.938}{K}%
\TextAt{17}{410.134}{633.144}{1}%
\TextAt{9}{419.399}{634.938}{\Character{61}}%
\TextAt{13}{433.037}{634.938}{K}%
\TextAt{21}{442.996}{633.144}{nr}%
\TextAt{8}{452.689}{634.938}{,}%
\TextAt{8}{460.728}{634.938}{which}%
\TextAt{8}{494.925}{634.938}{has}%
\TextAt{8}{516.51}{634.938}{the}%
\TextAt{8}{79.226}{617.892}{separable}%
\TextAt{8}{130.669}{617.892}{closure}%
\TextAt{13}{169.751}{617.892}{\Character{20}}%
\TextAt{17}{176.489}{616.098}{1}%
\TextAt{9}{184.542}{617.892}{\Character{61}}%
\TextAt{9}{197.394}{617.892}{\Character{22}}%
\TextAt{13}{196.968}{617.892}{\Character{20}}%
\TextAt{21}{203.706}{616.098}{s}%
\TextAt{8}{212.018}{617.892}{of}%
\TextAt{13}{225.348}{617.892}{\Character{20}}%
\TextAt{8}{235.984}{617.892}{as}%
\TextAt{8}{250.354}{617.892}{its}%
\TextAt{8}{266.673}{617.892}{residue}%
\TextAt{8}{306.4}{617.892}{class}%
\TextAt{8}{333.837}{617.892}{field.}%
\TextAt{8}{363.745}{617.892}{We}%
\TextAt{8}{383.889}{617.892}{obtain}%
\TextAt{8}{420.317}{617.892}{the}%
\TextAt{8}{440.473}{617.892}{following}%
\TextAt{8}{489.908}{617.892}{diagram}%
\TextAt{13}{266.371}{589.94}{L}%
\TextAt{13}{330.104}{589.94}{L}%
\TextAt{17}{338.069}{588.147}{1}%
\TextAt{13}{264.944}{538.939}{K}%
\TextAt{13}{329.106}{538.939}{K}%
\TextAt{17}{339.065}{537.145}{1}%
\NativeDiagram{%
\path[draw=black,line width=0.478bp] (207.688,-181.07) -- (252.753,-181.07);%
\path[draw=black,line width=0.478bp] (198.354,-222.312) -- (198.354,-188.638);%
\path[draw=black,line width=0.478bp] (209.115,-232.072) -- (251.755,-232.072);%
\path[draw=black,line width=0.478bp] (264.453,-222.312) -- (264.453,-190.432);%
}%
\TextAt{13}{348.951}{567.628}{\Character{59}}%
\TextAt{8}{79.226}{510.247}{where}%
\TextAt{13}{114.026}{510.247}{L}%
\TextAt{16}{125.323}{510.247}{\Character{92}}%
\TextAt{13}{136.625}{510.247}{K}%
\TextAt{17}{146.584}{508.454}{1}%
\TextAt{9}{156.325}{510.247}{\Character{61}}%
\TextAt{13}{170.439}{510.247}{T}%
\TextAt{9}{183.935}{510.247}{\Character{61}}%
\TextAt{13}{198.049}{510.247}{K}%
\TextAt{8}{213.755}{510.247}{and}%
\TextAt{13}{237.508}{510.247}{L}%
\TextAt{17}{245.473}{508.454}{1}%
\TextAt{9}{255.214}{510.247}{\Character{61}}%
\TextAt{13}{269.328}{510.247}{K}%
\TextAt{17}{279.287}{508.454}{1}%
\TextAt{9}{284.019}{510.247}{\Character{40}}%
\TextAt{22}{291.893}{514.268}{m}%
\TextAt{16}{291.789}{518.858}{\Character{112}}%
\Vector{q 0 G 0 g 0.478 w 229.752 -254.903 m 235.896 -254.903 l S Q}%
\TextAt{13}{301.752}{510.247}{a}%
\TextAt{9}{307.896}{510.247}{\Character{41}}%
\TextAt{8}{312.449}{510.247}{.}%
\TextAt{8}{323.869}{510.247}{If}%
\TextAt{8}{336.552}{510.247}{now}%
\TextAt{13}{361.932}{510.247}{L}%
\TextAt{17}{369.896}{508.454}{1}%
\TextAt{16}{379.637}{510.247}{\Character{106}}%
\TextAt{13}{387.968}{510.247}{K}%
\TextAt{17}{397.927}{508.454}{1}%
\TextAt{8}{407.547}{510.247}{is}%
\TextAt{8}{420.302}{510.247}{tamely}%
\TextAt{8}{459.991}{510.247}{ramified,}%
\TextAt{8}{510.007}{510.247}{then}%
\TextAt{13}{79.226}{493.201}{\Character{21}}%
\TextAt{17}{86.055}{491.408}{1}%
\TextAt{16}{94.108}{493.201}{\Character{106}}%
\TextAt{13}{100.749}{493.201}{\Character{20}}%
\TextAt{17}{107.488}{491.408}{1}%
\TextAt{8}{116.103}{493.201}{is}%
\TextAt{8}{127.852}{493.201}{separable;}%
\TextAt{8}{182.537}{493.201}{hence}%
\TextAt{13}{215.029}{493.201}{\Character{21}}%
\TextAt{17}{221.857}{491.408}{1}%
\TextAt{9}{229.91}{493.201}{\Character{61}}%
\TextAt{13}{242.336}{493.201}{\Character{20}}%
\TextAt{17}{249.074}{491.408}{1}%
\TextAt{8}{257.689}{493.201}{and}%
\TextAt{13}{280.437}{493.201}{p}%
\TextAt{18}{289.633}{493.201}{\Character{45}}%
\TextAt{9}{296.275}{493.201}{\Character{91}}%
\TextAt{13}{299.527}{493.201}{L}%
\TextAt{17}{307.491}{491.408}{1}%
\TextAt{9}{315.544}{493.201}{\Character{58}}%
\TextAt{13}{322.117}{493.201}{K}%
\TextAt{17}{332.076}{491.408}{1}%
\TextAt{9}{336.808}{493.201}{\Character{93}}%
\TextAt{9}{343.359}{493.201}{\Character{61}}%
\TextAt{9}{355.781}{493.201}{\Character{91}}%
\TextAt{13}{359.058}{493.201}{L}%
\TextAt{9}{370.343}{493.201}{\Character{58}}%
\TextAt{13}{376.916}{493.201}{K}%
\TextAt{9}{387.733}{493.201}{\Character{93}}%
\TextAt{9}{394.284}{493.201}{\Character{61}}%
\TextAt{9}{406.706}{493.201}{\Character{91}}%
\TextAt{13}{409.983}{493.201}{L}%
\TextAt{9}{421.268}{493.201}{\Character{58}}%
\TextAt{13}{427.841}{493.201}{T}%
\TextAt{9}{436.328}{493.201}{\Character{93}}%
\TextAt{8}{439.579}{493.201}{,}%
\TextAt{8}{446.717}{493.201}{i.e.,}%
\TextAt{13}{468.81}{493.201}{L}%
\TextAt{16}{480.096}{493.201}{\Character{106}}%
\TextAt{13}{486.737}{493.201}{K}%
\TextAt{8}{501.438}{493.201}{is}%
\TextAt{8}{513.187}{493.201}{also}%
\TextAt{8}{79.226}{476.155}{tamely}%
\TextAt{8}{117.925}{476.155}{ramified.}%
\ErrataAt{78.9}{456.142}{460.1}{Put $\alpha=\sqrt[m]{a}$ and $d=[K(\alpha):K]$. Since $\mu_m\subseteq K$, the map $\sigma\mapsto\sigma(\alpha)/\alpha$}%
\ErrataAt{79.226}{439.096}{455.774}{identifies the Galois group with $\mu_d$. Thus $d\mid m$, and $\alpha^j\in K$ exactly when $d\mid j$.}%
\ErrataAt{79.226}{422.05}{455.774}{Put $n=\operatorname{ord}(w(\alpha)\bmod v(K^*))$, so $d=hn$ for some $h$. Choose $b\in K^*$ with $v(b)=w(\alpha^n)$.}%
\ErrataAt{79.226}{405.004}{455.774}{The unit $\alpha^d/b^h$ has an $h$th root $c\in K$ by Hensel\textquoteright{}s lemma, since $(h,p)=1$}%
\ErrataAt{79.226}{387.958}{455.774}{and the residue field is separably closed. Then $\alpha^n/(bc)\in\mu_h\subseteq K$, so $d\mid n$ and $n=d$.}%
\ErrataAt{79.226}{370.912}{455.774}{The ramification index $e$ is divisible by $n$, while the fundamental inequality gives $[L:K]\ge ef$.}%
\ErrataAt{79.226}{353.866}{455.774}{Consequently, recalling that $[L:K]=d=n$, we have}%
\TextAt{13}{239.844}{307.797}{e}%
\NativeMathAt{248.59}{307.797}{11.955}{\geq}%
\TextAt{13}{261.21}{307.797}{n}%
\TextAt{9}{271.518}{307.797}{\Character{61}}%
\TextAt{9}{283.928}{307.797}{\Character{91}}%
\TextAt{13}{287.195}{307.797}{L}%
\TextAt{9}{298.481}{307.797}{\Character{58}}%
\TextAt{13}{305.053}{307.797}{K}%
\TextAt{9}{315.871}{307.797}{\Character{93}}%
\NativeMathAt{322.443}{307.797}{11.955}{\geq}%
\TextAt{13}{335.063}{307.797}{ef}%
\NativeMathAt{350.855}{307.797}{11.955}{\geq}%
\TextAt{13}{363.475}{307.797}{e\Character{59}}%
\TextAt{8}{79.226}{278.796}{in}%
\TextAt{8}{93.765}{278.796}{other}%
\TextAt{8}{125.221}{278.796}{words}%
\TextAt{13}{159.665}{278.796}{f}%
\TextAt{9}{171.543}{278.796}{\Character{61}}%
\TextAt{9}{185.471}{278.796}{1}%
\TextAt{8}{191.331}{278.796}{,}%
\TextAt{8}{199.588}{278.796}{and}%
\TextAt{8}{223.237}{278.796}{so}%
\TextAt{13}{238.493}{278.796}{\Character{21}}%
\TextAt{9}{250.153}{278.796}{\Character{61}}%
\TextAt{13}{264.088}{278.796}{\Character{20}}%
\TextAt{8}{275.611}{278.796}{and}%
\TextAt{13}{299.259}{278.796}{p}%
\TextAt{18}{309.966}{278.796}{\Character{45}}%
\TextAt{13}{318.117}{278.796}{n}%
\TextAt{9}{329.936}{278.796}{\Character{61}}%
\TextAt{13}{343.872}{278.796}{e}%
\TextAt{8}{349.297}{278.796}{.}%
\TextAt{8}{360.404}{278.796}{This}%
\TextAt{8}{388.01}{278.796}{shows}%
\TextAt{8}{422.514}{278.796}{that}%
\TextAt{13}{448.769}{278.796}{L}%
\TextAt{16}{461.565}{278.796}{\Character{106}}%
\TextAt{13}{469.717}{278.796}{K}%
\TextAt{8}{485.318}{278.796}{is}%
\TextAt{8}{497.968}{278.796}{tamely}%
\TextAt{8}{79.226}{261.75}{ramified.}%
\Vector{q 0 G 0 g 0.399 w 452.334 -512.25 m 452.334 -504.18 l S Q}%
\Vector{q 0 G 0 g 0.399 w 452.534 -504.379 m 459.707 -504.379 l S Q}%
\Vector{q 0 G 0 g 0.399 w 452.534 -512.05 m 459.707 -512.05 l S Q}%
\Vector{q 0 G 0 g 0.399 w 459.906 -512.25 m 459.906 -504.18 l S Q}%
\TextAt{12}{79.226}{226.772}{Corollary}%
\TextAt{12}{139.703}{226.772}{(7.8).}%
\TextAt{8}{175.617}{226.772}{Let}%
\TextAt{13}{197.185}{226.772}{L}%
\TextAt{16}{209.49}{226.772}{\Character{106}}%
\TextAt{13}{217.152}{226.772}{K}%
\TextAt{8}{232.466}{226.772}{and}%
\TextAt{13}{255.827}{226.772}{K}%
\TextAt{3}{266.645}{231.11}{\Character{48}}%
\TextAt{16}{273.781}{226.772}{\Character{106}}%
\TextAt{13}{281.443}{226.772}{K}%
\TextAt{8}{296.757}{226.772}{be}%
\TextAt{8}{313.28}{226.772}{two}%
\TextAt{8}{335.996}{226.772}{extensions}%
\TextAt{8}{392.975}{226.772}{inside}%
\TextAt{8}{426.797}{226.772}{the}%
\TextAt{8}{447.552}{226.772}{algebraic}%
\TextAt{8}{497.585}{226.772}{closure}%
\Vector{q 0 G 0 g 0.478 w 7.226 -554.431 m 18.044 -554.431 l S Q}%
\TextAt{13}{79.226}{209.726}{K}%
\TextAt{16}{93.365}{209.726}{\Character{106}}%
\TextAt{13}{100.006}{209.726}{K}%
\TextAt{8}{110.824}{209.726}{,}%
\TextAt{8}{117.973}{209.726}{and}%
\TextAt{13}{140.736}{209.726}{L}%
\TextAt{3}{148.7}{214.064}{\Character{48}}%
\TextAt{9}{154.816}{209.726}{\Character{61}}%
\TextAt{13}{167.242}{209.726}{LK}%
\TextAt{3}{186.024}{214.064}{\Character{48}}%
\TextAt{8}{188.819}{209.726}{.}%
\TextAt{8}{197.267}{209.726}{Then}%
\TextAt{8}{227.825}{209.726}{we}%
\TextAt{8}{245.052}{209.726}{have:}%
\TextAt{13}{163.798}{180.725}{L}%
\TextAt{16}{175.084}{180.725}{\Character{106}}%
\TextAt{13}{181.725}{180.725}{K}%
\TextAt{8}{196.44}{180.725}{tamely}%
\TextAt{8}{235.139}{180.725}{ramified}%
\NativeArrowAt{292.043}{180.725}{11.955}{19.067}{\Longrightarrow}%
\TextAt{13}{326.387}{180.725}{L}%
\TextAt{3}{334.351}{185.661}{\Character{48}}%
\TextAt{16}{340.467}{180.725}{\Character{106}}%
\TextAt{13}{347.109}{180.725}{K}%
\TextAt{3}{357.926}{185.661}{\Character{48}}%
\TextAt{8}{364.619}{180.725}{tamely}%
\TextAt{8}{403.318}{180.725}{ramified}%
\TextAt{13}{444.946}{180.725}{\Character{58}}%
\TextAt{8}{79.226}{151.723}{Every}%
\TextAt{8}{112.88}{151.723}{subextension}%
\TextAt{8}{182.591}{151.723}{of}%
\TextAt{8}{195.921}{151.723}{a}%
\TextAt{8}{205.676}{151.723}{tamely}%
\TextAt{8}{244.375}{151.723}{ramified}%
\TextAt{8}{289.9}{151.723}{extension}%
\TextAt{8}{341.666}{151.723}{is}%
\TextAt{8}{353.43}{151.723}{tamely}%
\TextAt{8}{392.129}{151.723}{ramified.}%
\TextAt{19}{79.226}{116.745}{Proof.}%
\TextAt{8}{116}{116.745}{We}%
\TextAt{8}{136.96}{116.745}{may}%
\TextAt{8}{163.145}{116.745}{assume}%
\TextAt{8}{204.405}{116.745}{without}%
\TextAt{8}{248.798}{116.745}{loss}%
\TextAt{8}{271.85}{116.745}{of}%
\TextAt{8}{285.996}{116.745}{generality}%
\TextAt{8}{340.801}{116.745}{that}%
\TextAt{13}{366.986}{116.745}{L}%
\TextAt{16}{379.661}{116.745}{\Character{106}}%
\TextAt{13}{387.693}{116.745}{K}%
\TextAt{8}{403.224}{116.745}{is}%
\TextAt{8}{415.804}{116.745}{finite}%
\TextAt{8}{446.531}{116.745}{and}%
\TextAt{8}{470.11}{116.745}{consider}%
\TextAt{8}{516.51}{116.745}{the}%
\TextAt{8}{297.211}{43.033}{143}%
\Destination{78.23}{809.815}{nav.page.143}%
\Destination{79.36}{776.155}{nav.print.156}%
\Destination{79.36}{240.755}{nav.number.2.7.8}%
\Destination{79.226}{468}{errata.193}%
}

% Page 156 (body page 144)
\Page{%
\TextAt{2}{79.226}{796.943}{II.}%
\TextAt{2}{94.073}{796.943}{The}%
\TextAt{2}{116.492}{796.943}{Theory}%
\TextAt{2}{154.401}{796.943}{of}%
\TextAt{2}{166.826}{796.943}{Valuations}%
\Vector{q 0 G 0 g 0.299 w 7.226 18.729 m 460.769 18.729 l S Q}%
\TextAt{8}{79.226}{761.124}{diagram}%
\TextAt{13}{266.911}{744.529}{L}%
\TextAt{13}{331.073}{744.529}{L}%
\TextAt{3}{339.037}{748.868}{\Character{48}}%
\TextAt{13}{266.65}{694.725}{T}%
\TextAt{13}{330.812}{694.725}{T}%
\TextAt{3}{339.299}{699.063}{\Character{48}}%
\TextAt{13}{265.484}{644.92}{K}%
\TextAt{13}{329.647}{644.92}{K}%
\TextAt{3}{340.464}{649.258}{\Character{48}}%
\NativeDiagram{%
\path[draw=black,line width=0.478bp] (208.227,-26.48) -- (253.721,-26.48);%
\path[draw=black,line width=0.478bp] (198.893,-66.526) -- (198.893,-34.048);%
\path[draw=black,line width=0.478bp] (208.489,-76.285) -- (253.46,-76.285);%
\path[draw=black,line width=0.478bp] (264.453,-65.929) -- (264.453,-34.048);%
\path[draw=black,line width=0.478bp] (198.893,-116.331) -- (198.893,-83.853);%
\path[draw=black,line width=0.478bp] (209.654,-126.091) -- (252.295,-126.091);%
\path[draw=black,line width=0.478bp] (264.453,-115.734) -- (264.453,-83.853);%
}%
\TextAt{13}{348.412}{699.109}{\Character{58}}%
\TextAt{8}{79.226}{626.714}{The}%
\TextAt{8}{105.556}{626.714}{inclusion}%
\TextAt{13}{156.668}{626.714}{T}%
\NativeMathAt{172.354}{626.714}{11.955}{\subseteq}%
\TextAt{13}{188.852}{626.714}{T}%
\TextAt{3}{197.338}{631.052}{\Character{48}}%
\TextAt{8}{206.307}{626.714}{follows}%
\TextAt{8}{247.018}{626.714}{from}%
\TextAt{8}{276.934}{626.714}{(7.2).}%
\TextAt{8}{316.288}{626.714}{If}%
\TextAt{13}{330.255}{626.714}{L}%
\TextAt{16}{345.419}{626.714}{\Character{106}}%
\TextAt{13}{355.939}{626.714}{K}%
\TextAt{8}{372.93}{626.714}{is}%
\TextAt{8}{386.969}{626.714}{tamely}%
\TextAt{8}{427.944}{626.714}{ramified,}%
\TextAt{8}{479.566}{626.714}{then}%
\TextAt{13}{508.501}{626.714}{L}%
\TextAt{9}{523.665}{626.714}{\Character{61}}%
\TextAt{13}{79.226}{609.668}{T}%
\TextAt{9}{87.713}{609.668}{\Character{40}}%
\TextAt{22}{95.586}{612.613}{m}%
\TextAt{24}{102.124}{610.819}{1}%
\TextAt{16}{99.633}{617.382}{\Character{112}}%
\Vector{q 0 G 0 g 0.478 w 37.596 -156.379 m 48.473 -156.379 l S Q}%
\TextAt{13}{109.596}{609.668}{a}%
\TextAt{17}{115.741}{607.874}{1}%
\TextAt{13}{120.473}{609.668}{\Character{59}}%
\TextAt{13}{125.697}{609.668}{\Character{58}}%
\TextAt{13}{130.934}{609.668}{\Character{58}}%
\TextAt{13}{136.17}{609.668}{\Character{58}}%
\TextAt{13}{141.395}{609.668}{\Character{59}}%
\TextAt{22}{150.015}{612.613}{m}%
\TextAt{22}{156.553}{611.616}{r}%
\TextAt{16}{154.069}{617.382}{\Character{112}}%
\Vector{q 0 G 0 g 0.478 w 92.032 -156.379 m 102.732 -156.379 l S Q}%
\TextAt{13}{164.032}{609.668}{a}%
\TextAt{21}{170.177}{607.874}{r}%
\TextAt{9}{174.732}{609.668}{\Character{41}}%
\TextAt{8}{179.284}{609.668}{,}%
\TextAt{9}{188.576}{609.668}{\Character{40}}%
\TextAt{13}{193.128}{609.668}{m}%
\TextAt{21}{203.367}{607.874}{i}%
\TextAt{13}{206.749}{609.668}{\Character{59}}%
\TextAt{13}{211.973}{609.668}{p}%
\TextAt{9}{217.868}{609.668}{\Character{41}}%
\TextAt{9}{228.652}{609.668}{\Character{61}}%
\TextAt{9}{243.99}{609.668}{1}%
\TextAt{8}{249.862}{609.668}{;}%
\TextAt{8}{259.582}{609.668}{hence}%
\TextAt{13}{293.803}{609.668}{L}%
\TextAt{3}{301.767}{614.006}{\Character{48}}%
\TextAt{9}{310.804}{609.668}{\Character{61}}%
\TextAt{13}{326.151}{609.668}{LK}%
\TextAt{3}{344.933}{614.006}{\Character{48}}%
\TextAt{9}{353.97}{609.668}{\Character{61}}%
\TextAt{13}{369.317}{609.668}{LT}%
\TextAt{3}{385.768}{614.006}{\Character{48}}%
\TextAt{9}{394.805}{609.668}{\Character{61}}%
\TextAt{13}{410.152}{609.668}{T}%
\TextAt{3}{418.639}{614.006}{\Character{48}}%
\TextAt{9}{421.434}{609.668}{\Character{40}}%
\TextAt{22}{429.307}{612.613}{m}%
\TextAt{24}{435.845}{610.819}{1}%
\TextAt{16}{433.354}{617.382}{\Character{112}}%
\Vector{q 0 G 0 g 0.478 w 371.317 -156.379 m 382.194 -156.379 l S Q}%
\TextAt{13}{443.317}{609.668}{a}%
\TextAt{17}{449.462}{607.874}{1}%
\TextAt{13}{454.194}{609.668}{\Character{59}}%
\TextAt{13}{459.418}{609.668}{\Character{58}}%
\TextAt{13}{464.655}{609.668}{\Character{58}}%
\TextAt{13}{469.891}{609.668}{\Character{58}}%
\TextAt{13}{475.116}{609.668}{\Character{59}}%
\TextAt{22}{483.736}{612.613}{m}%
\TextAt{22}{490.274}{611.616}{r}%
\TextAt{16}{487.79}{617.382}{\Character{112}}%
\Vector{q 0 G 0 g 0.478 w 425.753 -156.379 m 436.453 -156.379 l S Q}%
\TextAt{13}{497.753}{609.668}{a}%
\TextAt{21}{503.898}{607.874}{r}%
\TextAt{9}{508.453}{609.668}{\Character{41}}%
\TextAt{8}{513.005}{609.668}{,}%
\TextAt{8}{522.297}{609.668}{so}%
\TextAt{8}{79.226}{592.622}{that}%
\TextAt{13}{104.595}{592.622}{L}%
\TextAt{3}{112.56}{596.96}{\Character{48}}%
\TextAt{16}{118.675}{592.622}{\Character{106}}%
\TextAt{13}{125.317}{592.622}{K}%
\TextAt{3}{136.135}{596.96}{\Character{48}}%
\TextAt{8}{142.827}{592.622}{is}%
\TextAt{8}{154.591}{592.622}{also}%
\TextAt{8}{178.071}{592.622}{tamely}%
\TextAt{8}{216.77}{592.622}{ramified,}%
\TextAt{8}{265.547}{592.622}{by}%
\TextAt{8}{281.806}{592.622}{(7.7).}%
\TextAt{8}{97.159}{572.587}{The}%
\TextAt{8}{120.892}{572.587}{claim}%
\TextAt{8}{151.787}{572.587}{concerning}%
\TextAt{8}{209.999}{572.587}{the}%
\TextAt{8}{229.835}{572.587}{subextensions}%
\TextAt{8}{303.84}{572.587}{follows}%
\TextAt{8}{341.955}{572.587}{exactly}%
\TextAt{8}{381.96}{572.587}{as}%
\TextAt{8}{396.01}{572.587}{in}%
\TextAt{8}{409.342}{572.587}{the}%
\TextAt{8}{429.178}{572.587}{unramified}%
\TextAt{8}{487.39}{572.587}{case.}%
\Vector{q 0 G 0 g 0.399 w 452.334 -201.413 m 452.334 -193.343 l S Q}%
\Vector{q 0 G 0 g 0.399 w 452.534 -193.543 m 459.707 -193.543 l S Q}%
\Vector{q 0 G 0 g 0.399 w 452.534 -201.214 m 459.707 -201.214 l S Q}%
\Vector{q 0 G 0 g 0.399 w 459.906 -201.413 m 459.906 -193.343 l S Q}%
\TextAt{12}{79.226}{541.089}{Corollary}%
\TextAt{12}{139.014}{541.089}{(7.9).}%
\TextAt{8}{174.33}{541.089}{The}%
\TextAt{8}{198.383}{541.089}{composite}%
\TextAt{8}{253.401}{541.089}{of}%
\TextAt{8}{266.731}{541.089}{tamely}%
\TextAt{8}{305.43}{541.089}{ramified}%
\TextAt{8}{350.955}{541.089}{extensions}%
\TextAt{8}{407.336}{541.089}{is}%
\TextAt{8}{419.1}{541.089}{tamely}%
\TextAt{8}{457.799}{541.089}{ramified.}%
\TextAt{19}{79.226}{509.592}{Proof.}%
\TextAt{8}{116}{509.592}{This}%
\TextAt{8}{144.419}{509.592}{follows}%
\TextAt{8}{184.554}{509.592}{from}%
\TextAt{8}{213.894}{509.592}{(7.8),}%
\TextAt{8}{247.244}{509.592}{exactly}%
\TextAt{8}{289.268}{509.592}{as}%
\TextAt{8}{305.337}{509.592}{(7.3)}%
\TextAt{8}{335.012}{509.592}{followed}%
\TextAt{8}{381.913}{509.592}{from}%
\TextAt{8}{411.253}{509.592}{(7.2)}%
\TextAt{8}{440.927}{509.592}{in}%
\TextAt{8}{456.279}{509.592}{the}%
\TextAt{8}{478.134}{509.592}{unramified}%
\TextAt{8}{79.226}{492.546}{case.}%
\Vector{q 0 G 0 g 0.399 w 452.334 -281.454 m 452.334 -273.384 l S Q}%
\Vector{q 0 G 0 g 0.399 w 452.534 -273.584 m 459.707 -273.584 l S Q}%
\Vector{q 0 G 0 g 0.399 w 452.534 -281.255 m 459.707 -281.255 l S Q}%
\Vector{q 0 G 0 g 0.399 w 459.906 -281.454 m 459.906 -273.384 l S Q}%
\TextAt{12}{79.226}{461.048}{Definition}%
\TextAt{12}{141.244}{461.048}{(7.10).}%
\TextAt{8}{182.62}{461.048}{Let}%
\TextAt{13}{202.92}{461.048}{L}%
\TextAt{16}{214.205}{461.048}{\Character{106}}%
\TextAt{13}{220.847}{461.048}{K}%
\TextAt{8}{234.892}{461.048}{be}%
\TextAt{8}{250.146}{461.048}{an}%
\TextAt{8}{265.735}{461.048}{algebraic}%
\TextAt{8}{314.5}{461.048}{extension.}%
\TextAt{8}{370.593}{461.048}{Then}%
\TextAt{8}{400.48}{461.048}{the}%
\TextAt{8}{419.967}{461.048}{composite}%
\TextAt{13}{474.315}{461.048}{V}%
\TextAt{16}{487.061}{461.048}{\Character{106}}%
\TextAt{13}{493.703}{461.048}{K}%
\TextAt{8}{507.748}{461.048}{of}%
\TextAt{8}{520.408}{461.048}{all}%
\TextAt{8}{79.226}{444.002}{tamely}%
\TextAt{8}{118.177}{444.002}{ramified}%
\TextAt{8}{163.954}{444.002}{subextensions}%
\TextAt{8}{238.531}{444.002}{is}%
\TextAt{8}{250.546}{444.002}{called}%
\TextAt{8}{283.962}{444.002}{the}%
\TextAt{12}{304.37}{444.002}{maximal}%
\TextAt{12}{359.235}{444.002}{tamely}%
\TextAt{12}{403.974}{444.002}{ramified}%
\TextAt{12}{456.818}{444.002}{subextension}%
\TextAt{8}{79.226}{426.956}{of}%
\TextAt{13}{92.556}{426.956}{L}%
\TextAt{16}{103.841}{426.956}{\Character{106}}%
\TextAt{13}{110.483}{426.956}{K}%
\TextAt{8}{121.301}{426.956}{.}%
\TextAt{8}{97.159}{402.685}{Let}%
\TextAt{13}{118.745}{402.685}{w}%
\TextAt{9}{127.435}{402.685}{\Character{40}}%
\TextAt{13}{131.987}{402.685}{L}%
\TextAt{3}{139.952}{407.023}{\Character{3}}%
\TextAt{9}{144.684}{402.685}{\Character{41}}%
\TextAt{17}{149.236}{407.023}{\Character{40}}%
\TextAt{21}{152.529}{407.023}{p}%
\TextAt{17}{156.792}{407.023}{\Character{41}}%
\TextAt{8}{165.098}{402.685}{denote}%
\TextAt{8}{203.433}{402.685}{the}%
\TextAt{8}{224.207}{402.685}{subgroup}%
\TextAt{8}{275.621}{402.685}{of}%
\TextAt{8}{289.568}{402.685}{all}%
\TextAt{8}{306.444}{402.685}{elements}%
\TextAt{13}{354.917}{402.685}{\Character{33}}%
\NativeMathAt{367.008}{402.685}{11.955}{\in}%
\TextAt{13}{379.35}{402.685}{w}%
\TextAt{9}{388.039}{402.685}{\Character{40}}%
\TextAt{13}{392.592}{402.685}{L}%
\TextAt{3}{400.556}{407.023}{\Character{3}}%
\TextAt{9}{405.289}{402.685}{\Character{41}}%
\TextAt{8}{414.355}{402.685}{such}%
\TextAt{8}{441.369}{402.685}{that}%
\TextAt{13}{467.355}{402.685}{m\Character{33}}%
\NativeMathAt{489.684}{402.685}{11.955}{\in}%
\TextAt{13}{502.027}{402.685}{v}%
\TextAt{9}{508.115}{402.685}{\Character{40}}%
\TextAt{13}{512.667}{402.685}{K}%
\TextAt{3}{523.485}{407.023}{\Character{3}}%
\TextAt{9}{528.217}{402.685}{\Character{41}}%
\TextAt{8}{79.226}{385.639}{for}%
\TextAt{8}{96.963}{385.639}{some}%
\TextAt{13}{126.14}{385.639}{m}%
\TextAt{8}{140.129}{385.639}{satisfying}%
\TextAt{9}{192.141}{385.639}{\Character{40}}%
\TextAt{13}{196.693}{385.639}{m\Character{59}}%
\TextAt{13}{212.151}{385.639}{p}%
\TextAt{9}{218.052}{385.639}{\Character{41}}%
\TextAt{9}{225.907}{385.639}{\Character{61}}%
\TextAt{9}{238.328}{385.639}{1}%
\TextAt{8}{244.203}{385.639}{.}%
\TextAt{8}{252.602}{385.639}{The}%
\TextAt{8}{276.508}{385.639}{quotient}%
\TextAt{8}{322.542}{385.639}{group}%
\TextAt{13}{355.569}{385.639}{w}%
\TextAt{9}{364.259}{385.639}{\Character{40}}%
\TextAt{13}{368.811}{385.639}{L}%
\TextAt{3}{376.776}{389.977}{\Character{3}}%
\TextAt{9}{381.508}{385.639}{\Character{41}}%
\TextAt{17}{386.06}{389.977}{\Character{40}}%
\TextAt{21}{389.353}{389.977}{p}%
\TextAt{17}{393.616}{389.977}{\Character{41}}%
\TextAt{9}{397.408}{385.639}{\Character{47}}%
\TextAt{13}{403.261}{385.639}{v}%
\TextAt{9}{409.349}{385.639}{\Character{40}}%
\TextAt{13}{413.901}{385.639}{K}%
\TextAt{3}{424.719}{389.977}{\Character{3}}%
\TextAt{9}{429.451}{385.639}{\Character{41}}%
\TextAt{8}{437.753}{385.639}{then}%
\TextAt{8}{464.264}{385.639}{consists}%
\TextAt{8}{507.226}{385.639}{of}%
\TextAt{8}{520.408}{385.639}{all}%
\TextAt{8}{79.226}{368.593}{elements}%
\TextAt{8}{127.083}{368.593}{of}%
\TextAt{13}{140.413}{368.593}{w}%
\TextAt{9}{149.102}{368.593}{\Character{40}}%
\TextAt{13}{153.655}{368.593}{L}%
\TextAt{3}{161.619}{372.931}{\Character{3}}%
\TextAt{9}{166.351}{368.593}{\Character{41}\Character{47}}%
\TextAt{13}{176.757}{368.593}{v}%
\TextAt{9}{182.845}{368.593}{\Character{40}}%
\TextAt{13}{187.397}{368.593}{K}%
\TextAt{3}{198.215}{372.931}{\Character{3}}%
\TextAt{9}{202.947}{368.593}{\Character{41}}%
\TextAt{8}{211.397}{368.593}{whose}%
\TextAt{8}{245.923}{368.593}{order}%
\TextAt{8}{276.493}{368.593}{is}%
\TextAt{8}{288.257}{368.593}{prime}%
\TextAt{8}{321.42}{368.593}{to}%
\TextAt{13}{335.731}{368.593}{p}%
\TextAt{8}{341.606}{368.593}{.}%
\TextAt{12}{79.226}{344.321}{Proposition}%
\TextAt{12}{153.018}{344.321}{(7.11).}%
\TextAt{8}{196.004}{344.321}{The}%
\TextAt{8}{220.997}{344.321}{maximal}%
\TextAt{8}{269.745}{344.321}{tamely}%
\TextAt{8}{309.383}{344.321}{ramified}%
\TextAt{8}{355.847}{344.321}{subextension}%
\TextAt{13}{426.497}{344.321}{V}%
\TextAt{16}{440.844}{344.321}{\Character{106}}%
\TextAt{13}{449.087}{344.321}{K}%
\TextAt{8}{464.741}{344.321}{of}%
\TextAt{13}{479.01}{344.321}{L}%
\TextAt{16}{491.896}{344.321}{\Character{106}}%
\TextAt{13}{500.139}{344.321}{K}%
\TextAt{8}{515.793}{344.321}{has}%
\TextAt{8}{79.226}{327.275}{value}%
\TextAt{8}{109.53}{327.275}{group}%
\TextAt{13}{142.764}{327.275}{w}%
\TextAt{9}{151.453}{327.275}{\Character{40}}%
\TextAt{13}{156.006}{327.275}{V}%
\TextAt{3}{165.431}{331.613}{\Character{3}}%
\TextAt{9}{170.164}{327.275}{\Character{41}}%
\TextAt{9}{178.114}{327.275}{\Character{61}}%
\TextAt{13}{190.659}{327.275}{w}%
\TextAt{9}{199.349}{327.275}{\Character{40}}%
\TextAt{13}{203.901}{327.275}{L}%
\TextAt{3}{211.865}{331.613}{\Character{3}}%
\TextAt{9}{216.598}{327.275}{\Character{41}}%
\TextAt{17}{221.15}{331.613}{\Character{40}}%
\TextAt{21}{224.443}{331.613}{p}%
\TextAt{17}{228.706}{331.613}{\Character{41}}%
\TextAt{8}{236.453}{327.275}{and}%
\TextAt{8}{259.273}{327.275}{residue}%
\TextAt{8}{299.058}{327.275}{class}%
\TextAt{8}{326.552}{327.275}{field}%
\TextAt{8}{351.967}{327.275}{equal}%
\TextAt{8}{382.917}{327.275}{to}%
\TextAt{8}{397.285}{327.275}{the}%
\TextAt{8}{417.499}{327.275}{separable}%
\TextAt{8}{469}{327.275}{closure}%
\TextAt{13}{508.139}{327.275}{\Character{21}}%
\TextAt{21}{514.968}{325.482}{s}%
\TextAt{8}{523.337}{327.275}{of}%
\TextAt{13}{79.226}{310.229}{\Character{20}}%
\TextAt{8}{89.862}{310.229}{in}%
\TextAt{13}{103.515}{310.229}{\Character{21}}%
\TextAt{16}{113.664}{310.229}{\Character{106}}%
\TextAt{13}{120.306}{310.229}{\Character{20}}%
\TextAt{8}{127.044}{310.229}{.}%
\TextAt{19}{79.226}{278.731}{Proof.}%
\TextAt{8}{116}{278.731}{We}%
\TextAt{8}{137.781}{278.731}{may}%
\TextAt{8}{164.786}{278.731}{restrict}%
\TextAt{8}{206.806}{278.731}{to}%
\TextAt{8}{222.753}{278.731}{the}%
\TextAt{8}{244.545}{278.731}{case}%
\TextAt{8}{270.953}{278.731}{of}%
\TextAt{8}{285.919}{278.731}{a}%
\TextAt{8}{297.31}{278.731}{finite}%
\TextAt{8}{328.858}{278.731}{extension}%
\TextAt{13}{382.26}{278.731}{L}%
\TextAt{16}{396.334}{278.731}{\Character{106}}%
\TextAt{13}{405.764}{278.731}{K}%
\TextAt{8}{416.581}{278.731}{.}%
\TextAt{8}{429.938}{278.731}{By}%
\TextAt{8}{449.937}{278.731}{passing}%
\TextAt{8}{492.675}{278.731}{from}%
\TextAt{13}{521.952}{278.731}{K}%
\TextAt{8}{79.226}{261.686}{to}%
\TextAt{8}{95.324}{261.686}{the}%
\TextAt{8}{117.268}{261.686}{maximal}%
\TextAt{8}{166.864}{261.686}{unramified}%
\TextAt{8}{227.184}{261.686}{subextension,}%
\TextAt{8}{302.381}{261.686}{we}%
\TextAt{8}{321.396}{261.686}{may}%
\TextAt{8}{348.552}{261.686}{assume}%
\TextAt{8}{390.784}{261.686}{by}%
\TextAt{8}{408.83}{261.686}{(7.5)}%
\TextAt{8}{438.593}{261.686}{that}%
\TextAt{13}{465.749}{261.686}{\Character{21}}%
\TextAt{21}{472.578}{259.892}{s}%
\TextAt{9}{483.359}{261.686}{\Character{61}}%
\TextAt{13}{498.83}{261.686}{\Character{20}}%
\TextAt{8}{505.569}{261.686}{.}%
\TextAt{8}{519.38}{261.686}{As}%
\TextAt{13}{79.226}{244.64}{p}%
\TextAt{18}{88.547}{244.64}{\Character{45}}%
\TextAt{13}{95.313}{244.64}{e}%
\TextAt{9}{100.739}{244.64}{\Character{40}}%
\TextAt{13}{105.291}{244.64}{V}%
\TextAt{16}{118.162}{244.64}{\Character{106}}%
\TextAt{13}{124.929}{244.64}{K}%
\TextAt{9}{135.747}{244.64}{\Character{41}}%
\TextAt{9}{143.733}{244.64}{\Character{61}}%
\TextAt{8}{156.295}{244.64}{\char35\relax }%
\TextAt{13}{166.05}{244.64}{w}%
\TextAt{9}{174.74}{244.64}{\Character{40}}%
\TextAt{13}{179.292}{244.64}{V}%
\TextAt{3}{188.718}{248.978}{\Character{3}}%
\TextAt{9}{193.45}{244.64}{\Character{41}\Character{47}}%
\TextAt{13}{203.855}{244.64}{v}%
\TextAt{9}{209.944}{244.64}{\Character{40}}%
\TextAt{13}{214.496}{244.64}{K}%
\TextAt{3}{225.314}{248.978}{\Character{3}}%
\TextAt{9}{230.046}{244.64}{\Character{41}}%
\TextAt{8}{234.598}{244.64}{,}%
\TextAt{8}{241.839}{244.64}{we}%
\TextAt{8}{259.14}{244.64}{certainly}%
\TextAt{8}{307.667}{244.64}{have}%
\TextAt{13}{334.735}{244.64}{w}%
\TextAt{9}{343.425}{244.64}{\Character{40}}%
\TextAt{13}{347.977}{244.64}{V}%
\TextAt{3}{357.403}{248.978}{\Character{3}}%
\TextAt{9}{362.135}{244.64}{\Character{41}}%
\NativeMathAt{370.133}{244.64}{11.955}{\subseteq}%
\TextAt{13}{382.877}{244.64}{w}%
\TextAt{9}{391.566}{244.64}{\Character{40}}%
\TextAt{13}{396.119}{244.64}{L}%
\TextAt{3}{404.083}{248.978}{\Character{3}}%
\TextAt{9}{408.816}{244.64}{\Character{41}}%
\TextAt{17}{413.368}{248.978}{\Character{40}}%
\TextAt{21}{416.661}{248.978}{p}%
\TextAt{17}{420.924}{248.978}{\Character{41}}%
\TextAt{8}{424.715}{244.64}{.}%
\TextAt{8}{433.383}{244.64}{Conversely}%
\TextAt{8}{492.706}{244.64}{we}%
\TextAt{8}{510.007}{244.64}{find,}%
\TextAt{8}{79.226}{227.594}{as}%
\TextAt{8}{93.941}{227.594}{in}%
\TextAt{8}{107.938}{227.594}{the}%
\TextAt{8}{128.439}{227.594}{proof}%
\TextAt{8}{159.353}{227.594}{of}%
\TextAt{8}{173.027}{227.594}{(7.7),}%
\TextAt{8}{204.685}{227.594}{for}%
\TextAt{8}{222.914}{227.594}{every}%
\TextAt{13}{254.151}{227.594}{\Character{33}}%
\NativeMathAt{265.777}{227.594}{11.955}{\in}%
\TextAt{13}{277.655}{227.594}{w}%
\TextAt{9}{286.344}{227.594}{\Character{40}}%
\TextAt{13}{290.897}{227.594}{L}%
\TextAt{3}{298.861}{231.932}{\Character{3}}%
\TextAt{9}{303.593}{227.594}{\Character{41}}%
\TextAt{17}{308.146}{231.932}{\Character{40}}%
\TextAt{21}{311.439}{231.932}{p}%
\TextAt{17}{315.702}{231.932}{\Character{41}}%
\TextAt{8}{323.735}{227.594}{a}%
\TextAt{8}{333.835}{227.594}{radical}%
\TextAt{13}{372.555}{227.594}{\Character{11}}%
\TextAt{9}{383.985}{227.594}{\Character{61}}%
\TextAt{22}{400.318}{231.614}{m}%
\TextAt{16}{400.215}{236.205}{\Character{112}}%
\Vector{q 0 G 0 g 0.478 w 338.177 -537.556 m 344.322 -537.556 l S Q}%
\TextAt{13}{410.177}{227.594}{a}%
\NativeMathAt{420.23}{227.594}{11.955}{\in}%
\TextAt{13}{432.108}{227.594}{L}%
\TextAt{8}{444.314}{227.594}{such}%
\TextAt{8}{471.056}{227.594}{that}%
\TextAt{13}{496.769}{227.594}{a}%
\NativeMathAt{506.822}{227.594}{11.955}{\in}%
\TextAt{13}{518.7}{227.594}{K}%
\TextAt{8}{529.518}{227.594}{,}%
\TextAt{9}{79.226}{210.548}{\Character{40}}%
\TextAt{13}{83.778}{210.548}{m\Character{59}}%
\TextAt{13}{99.236}{210.548}{p}%
\TextAt{9}{105.137}{210.548}{\Character{41}}%
\TextAt{9}{112.992}{210.548}{\Character{61}}%
\TextAt{9}{125.413}{210.548}{1}%
\TextAt{8}{135.186}{210.548}{and}%
\TextAt{13}{157.949}{210.548}{w}%
\TextAt{9}{166.638}{210.548}{\Character{40}}%
\TextAt{13}{171.191}{210.548}{\Character{11}}%
\TextAt{9}{178.712}{210.548}{\Character{41}}%
\TextAt{9}{186.567}{210.548}{\Character{61}}%
\TextAt{13}{199.011}{210.548}{\Character{33}}%
\TextAt{8}{206.729}{210.548}{,}%
\TextAt{8}{213.878}{210.548}{so}%
\TextAt{8}{228.248}{210.548}{that}%
\TextAt{8}{253.617}{210.548}{one}%
\TextAt{8}{275.077}{210.548}{has}%
\TextAt{13}{295.951}{210.548}{\Character{11}}%
\NativeMathAt{306.793}{210.548}{11.955}{\in}%
\TextAt{13}{318.084}{210.548}{V}%
\TextAt{8}{327.51}{210.548}{,}%
\TextAt{8}{334.659}{210.548}{and}%
\TextAt{13}{357.422}{210.548}{\Character{33}}%
\NativeMathAt{368.461}{210.548}{11.955}{\in}%
\TextAt{13}{379.752}{210.548}{w}%
\TextAt{9}{388.441}{210.548}{\Character{40}}%
\TextAt{13}{392.993}{210.548}{V}%
\TextAt{3}{402.419}{214.886}{\Character{3}}%
\TextAt{9}{407.151}{210.548}{\Character{41}}%
\TextAt{8}{411.704}{210.548}{.}%
\Vector{q 0 G 0 g 0.399 w 452.334 -563.452 m 452.334 -555.383 l S Q}%
\Vector{q 0 G 0 g 0.399 w 452.534 -555.582 m 459.707 -555.582 l S Q}%
\Vector{q 0 G 0 g 0.399 w 452.534 -563.253 m 459.707 -563.253 l S Q}%
\Vector{q 0 G 0 g 0.399 w 459.906 -563.452 m 459.906 -555.383 l S Q}%
\TextAt{8}{97.159}{183.287}{The}%
\TextAt{8}{121.213}{183.287}{results}%
\TextAt{8}{158.405}{183.287}{obtained}%
\TextAt{8}{206.537}{183.287}{in}%
\TextAt{8}{220.19}{183.287}{this}%
\TextAt{8}{243.012}{183.287}{section}%
\TextAt{8}{282.093}{183.287}{may}%
\TextAt{8}{307.462}{183.287}{be}%
\TextAt{8}{323.387}{183.287}{summarized}%
\TextAt{8}{388.482}{183.287}{in}%
\TextAt{8}{402.135}{183.287}{the}%
\TextAt{8}{422.292}{183.287}{following}%
\TextAt{8}{471.726}{183.287}{picture:}%
\TextAt{13}{203.456}{162.045}{K}%
\NativeMathAt{234.199}{162.045}{11.955}{\subseteq}%
\TextAt{13}{264.723}{162.045}{T}%
\NativeMathAt{294.436}{162.045}{11.955}{\subseteq}%
\TextAt{13}{329.904}{162.045}{V}%
\NativeMathAt{365.498}{162.045}{11.955}{\subseteq}%
\TextAt{13}{396.023}{162.045}{L}%
\TextAt{13}{205.496}{144.999}{\Character{20}}%
\NativeMathAt{234.199}{144.999}{11.955}{\subseteq}%
\TextAt{13}{263.346}{144.999}{\Character{21}}%
\TextAt{21}{270.174}{143.206}{s}%
\TextAt{9}{294.533}{144.999}{\Character{61}}%
\TextAt{13}{328.995}{144.999}{\Character{21}}%
\TextAt{21}{335.824}{143.206}{s}%
\NativeMathAt{365.498}{144.999}{11.955}{\subseteq}%
\TextAt{13}{396.591}{144.999}{\Character{21}}%
\TextAt{13}{193.493}{127.953}{v}%
\TextAt{9}{199.582}{127.953}{\Character{40}}%
\TextAt{13}{204.134}{127.953}{K}%
\TextAt{3}{214.952}{132.291}{\Character{3}}%
\TextAt{9}{219.684}{127.953}{\Character{41}}%
\TextAt{9}{234.281}{127.953}{\Character{61}}%
\TextAt{13}{253.46}{127.953}{w}%
\TextAt{9}{262.15}{127.953}{\Character{40}}%
\TextAt{13}{266.702}{127.953}{T}%
\TextAt{3}{275.189}{132.291}{\Character{3}}%
\TextAt{9}{279.921}{127.953}{\Character{41}}%
\NativeMathAt{294.436}{127.953}{11.955}{\subseteq}%
\TextAt{13}{313.697}{127.953}{w}%
\TextAt{9}{322.387}{127.953}{\Character{40}}%
\TextAt{13}{326.939}{127.953}{L}%
\TextAt{3}{334.904}{132.291}{\Character{3}}%
\TextAt{9}{339.636}{127.953}{\Character{41}}%
\TextAt{17}{344.188}{132.291}{\Character{40}}%
\TextAt{21}{347.482}{132.291}{p}%
\TextAt{17}{351.744}{132.291}{\Character{41}}%
\NativeMathAt{365.498}{127.953}{11.955}{\subseteq}%
\TextAt{13}{384.759}{127.953}{w}%
\TextAt{9}{393.449}{127.953}{\Character{40}}%
\TextAt{13}{398.001}{127.953}{L}%
\TextAt{3}{405.966}{132.291}{\Character{3}}%
\TextAt{9}{410.698}{127.953}{\Character{41}}%
\TextAt{13}{420.232}{145.419}{\Character{58}}%
\TextAt{8}{79.226}{107.013}{If}%
\TextAt{13}{90.91}{107.013}{L}%
\TextAt{16}{102.195}{107.013}{\Character{106}}%
\TextAt{13}{108.837}{107.013}{K}%
\TextAt{8}{123.543}{107.013}{is}%
\TextAt{8}{135.298}{107.013}{finite}%
\TextAt{8}{165.201}{107.013}{and}%
\TextAt{13}{187.955}{107.013}{e}%
\TextAt{9}{196.702}{107.013}{\Character{61}}%
\TextAt{13}{209.127}{107.013}{e}%
\TextAt{3}{214.553}{111.352}{\Character{48}}%
\TextAt{13}{217.348}{107.013}{p}%
\TextAt{21}{223.223}{111.352}{a}%
\TextAt{8}{232.108}{107.013}{where}%
\TextAt{9}{265.908}{107.013}{\Character{40}}%
\TextAt{13}{270.461}{107.013}{e}%
\TextAt{3}{275.886}{111.352}{\Character{48}}%
\TextAt{13}{278.681}{107.013}{\Character{59}}%
\TextAt{13}{283.905}{107.013}{p}%
\TextAt{9}{289.8}{107.013}{\Character{41}}%
\TextAt{9}{297.655}{107.013}{\Character{61}}%
\TextAt{9}{310.076}{107.013}{1}%
\TextAt{8}{315.952}{107.013}{,}%
\TextAt{8}{323.094}{107.013}{then}%
\TextAt{9}{349.746}{107.013}{\Character{91}}%
\TextAt{13}{352.997}{107.013}{V}%
\TextAt{9}{365.744}{107.013}{\Character{58}}%
\TextAt{13}{372.316}{107.013}{T}%
\TextAt{9}{380.803}{107.013}{\Character{93}}%
\TextAt{9}{387.354}{107.013}{\Character{61}}%
\TextAt{13}{399.801}{107.013}{e}%
\TextAt{3}{405.227}{111.352}{\Character{48}}%
\TextAt{8}{408.022}{107.013}{.}%
\TextAt{8}{416.467}{107.013}{The}%
\TextAt{8}{440.512}{107.013}{extension}%
\TextAt{13}{492.27}{107.013}{L}%
\TextAt{16}{503.555}{107.013}{\Character{106}}%
\TextAt{13}{510.197}{107.013}{K}%
\TextAt{8}{524.903}{107.013}{is}%
\TextAt{8}{79.226}{89.967}{called}%
\TextAt{12}{112.323}{89.967}{totally}%
\TextAt{12}{155.059}{89.967}{(or}%
\TextAt{12}{176.932}{89.967}{purely)}%
\TextAt{12}{223.995}{89.967}{ramified}%
\TextAt{8}{275.898}{89.967}{if}%
\TextAt{13}{286.555}{89.967}{T}%
\TextAt{9}{298.363}{89.967}{\Character{61}}%
\TextAt{13}{310.789}{89.967}{K}%
\TextAt{8}{321.606}{89.967}{,}%
\TextAt{8}{328.702}{89.967}{and}%
\TextAt{12}{351.399}{89.967}{wildly}%
\TextAt{12}{391.325}{89.967}{ramified}%
\TextAt{8}{443.227}{89.967}{if}%
\TextAt{8}{453.885}{89.967}{it}%
\TextAt{8}{465.523}{89.967}{is}%
\TextAt{8}{477.22}{89.967}{not}%
\TextAt{8}{497.968}{89.967}{tamely}%
\TextAt{8}{79.226}{72.921}{ramified,}%
\TextAt{8}{128.003}{72.921}{i.e.,}%
\TextAt{8}{150.108}{72.921}{if}%
\TextAt{13}{160.832}{72.921}{V}%
\NativeMathAt{173.579}{72.921}{11.955}{\neq}%
\TextAt{13}{186.004}{72.921}{L}%
\TextAt{8}{193.969}{72.921}{.}%
\TextAt{8}{297.211}{43.033}{144}%
\Destination{78.23}{809.815}{nav.page.144}%
\Destination{79.36}{641.755}{nav.print.157}%
\Destination{79.36}{555.055}{nav.number.2.7.9}%
\Destination{79.36}{475.055}{nav.number.2.7.10}%
\Destination{79.36}{358.355}{nav.number.2.7.11}%
\Hotlink{276.434}{624.555}{25.078}{11.047}{nav.number.2.7.2}%
\Hotlink{281.306}{590.455}{25.078}{11.047}{nav.number.2.7.7}%
\Hotlink{213.394}{507.355}{25.078}{11.047}{nav.number.2.7.8}%
\Hotlink{304.837}{507.355}{25.078}{11.047}{nav.number.2.7.3}%
\Hotlink{410.753}{507.355}{25.078}{11.047}{nav.number.2.7.2}%
\Hotlink{408.33}{259.455}{25.078}{11.047}{nav.number.2.7.5}%
\Hotlink{172.527}{225.355}{25.078}{11.047}{nav.number.2.7.7}%
}

% Page 157 (body page 145)
\Page{%
\TextAt{2}{79.226}{796.943}{II.}%
\TextAt{2}{94.073}{796.943}{The}%
\TextAt{2}{116.492}{796.943}{Theory}%
\TextAt{2}{154.401}{796.943}{of}%
\TextAt{2}{166.826}{796.943}{Valuations}%
\Vector{q 0 G 0 g 0.299 w 7.226 18.729 m 460.769 18.729 l S Q}%
\TextAt{12}{79.226}{761.124}{Important}%
\TextAt{12}{145.892}{761.124}{Example:}%
\TextAt{8}{212.614}{761.124}{Consider}%
\TextAt{8}{262.992}{761.124}{the}%
\TextAt{8}{284.69}{761.124}{extension}%
\TextAt{18}{337.998}{761.124}{Q}%
\TextAt{21}{347.296}{759.33}{p}%
\TextAt{9}{352.057}{761.124}{\Character{40}}%
\TextAt{13}{356.61}{761.124}{\Character{16}}%
\TextAt{9}{362.625}{761.124}{\Character{41}}%
\TextAt{16}{373.126}{761.124}{\Character{106}}%
\TextAt{18}{382.394}{761.124}{Q}%
\TextAt{21}{391.693}{759.33}{p}%
\TextAt{8}{401.893}{761.124}{for}%
\TextAt{8}{421.319}{761.124}{a}%
\TextAt{8}{432.616}{761.124}{primitive}%
\TextAt{13}{484.238}{761.124}{n}%
\TextAt{8}{491.226}{761.124}{-th}%
\TextAt{8}{511.621}{761.124}{root}%
\TextAt{8}{79.226}{744.078}{of}%
\TextAt{8}{93.466}{744.078}{unity}%
\TextAt{13}{124.945}{744.078}{\Character{16}}%
\TextAt{8}{130.961}{744.078}{.}%
\TextAt{8}{142.138}{744.078}{In}%
\TextAt{8}{157.669}{744.078}{the}%
\TextAt{8}{178.735}{744.078}{two}%
\TextAt{8}{201.762}{744.078}{cases}%
\TextAt{9}{232.057}{744.078}{\Character{40}}%
\TextAt{13}{236.609}{744.078}{n\Character{59}}%
\TextAt{13}{248.815}{744.078}{p}%
\TextAt{9}{254.716}{744.078}{\Character{41}}%
\TextAt{9}{264.125}{744.078}{\Character{61}}%
\TextAt{9}{278.088}{744.078}{1}%
\TextAt{8}{288.776}{744.078}{and}%
\TextAt{13}{312.448}{744.078}{n}%
\TextAt{9}{324.307}{744.078}{\Character{61}}%
\TextAt{13}{338.282}{744.078}{p}%
\TextAt{21}{344.157}{748.416}{s}%
\TextAt{8}{348.571}{744.078}{,}%
\TextAt{8}{356.858}{744.078}{this}%
\TextAt{8}{380.59}{744.078}{extension}%
\TextAt{8}{433.266}{744.078}{behaves}%
\TextAt{8}{477.812}{744.078}{completely}%
\TextAt{8}{79.226}{727.032}{differently.}%
\TextAt{8}{139.953}{727.032}{Let}%
\TextAt{8}{161.44}{727.032}{us}%
\TextAt{8}{176.974}{727.032}{first}%
\TextAt{8}{201.617}{727.032}{look}%
\TextAt{8}{227.503}{727.032}{at}%
\TextAt{8}{242.331}{727.032}{the}%
\TextAt{8}{263.005}{727.032}{case}%
\TextAt{9}{288.294}{727.032}{\Character{40}}%
\TextAt{13}{292.847}{727.032}{n\Character{59}}%
\TextAt{13}{305.053}{727.032}{p}%
\TextAt{9}{310.953}{727.032}{\Character{41}}%
\TextAt{9}{319.692}{727.032}{\Character{61}}%
\TextAt{9}{332.998}{727.032}{1}%
\TextAt{8}{343.284}{727.032}{and}%
\TextAt{8}{366.565}{727.032}{choose}%
\TextAt{8}{404.215}{727.032}{as}%
\TextAt{8}{419.103}{727.032}{our}%
\TextAt{8}{440.435}{727.032}{base}%
\TextAt{8}{467.026}{727.032}{field,}%
\TextAt{8}{496.282}{727.032}{instead}%
\TextAt{8}{79.226}{709.986}{of}%
\TextAt{18}{92.828}{709.986}{Q}%
\TextAt{21}{102.127}{708.193}{p}%
\TextAt{8}{106.888}{709.986}{,}%
\TextAt{8}{114.377}{709.986}{any}%
\TextAt{8}{136.766}{709.986}{discretely}%
\TextAt{8}{189.449}{709.986}{valued}%
\TextAt{8}{226.471}{709.986}{complete}%
\TextAt{8}{276.166}{709.986}{field}%
\TextAt{13}{301.795}{709.986}{K}%
\TextAt{8}{316.782}{709.986}{with}%
\TextAt{8}{343.714}{709.986}{finite}%
\TextAt{8}{373.898}{709.986}{residue}%
\TextAt{8}{413.897}{709.986}{class}%
\TextAt{8}{441.606}{709.986}{field}%
\TextAt{13}{467.235}{709.986}{\Character{20}}%
\TextAt{9}{477.758}{709.986}{\Character{61}}%
\TextAt{18}{490.647}{709.986}{F}%
\TextAt{21}{497.953}{708.193}{q}%
\TextAt{8}{502.518}{709.986}{,}%
\TextAt{8}{510.007}{709.986}{with}%
\TextAt{13}{79.226}{692.94}{q}%
\TextAt{9}{88.166}{692.94}{\Character{61}}%
\TextAt{13}{100.592}{692.94}{p}%
\TextAt{21}{106.467}{697.278}{r}%
\TextAt{8}{111.021}{692.94}{.}%
\TextAt{12}{79.226}{666.927}{Proposition}%
\TextAt{12}{153.557}{666.927}{(7.12).}%
\TextAt{8}{197.012}{666.927}{Let}%
\TextAt{13}{219.389}{666.927}{L}%
\TextAt{9}{233.074}{666.927}{\Character{61}}%
\TextAt{13}{247.9}{666.927}{K}%
\TextAt{9}{258.717}{666.927}{\Character{40}}%
\TextAt{13}{263.27}{666.927}{\Character{16}}%
\TextAt{9}{269.285}{666.927}{\Character{41}}%
\TextAt{8}{273.838}{666.927}{,}%
\TextAt{8}{282.747}{666.927}{and}%
\TextAt{8}{306.918}{666.927}{let}%
\TextAt{16}{325.231}{666.927}{\Character{79}}%
\TextAt{16}{340.797}{666.927}{\Character{106}}%
\TextAt{14}{349.842}{666.927}{o}%
\TextAt{8}{355.687}{666.927}{,}%
\TextAt{8}{364.596}{666.927}{resp.}%
\TextAt{13}{394.027}{666.927}{\Character{21}}%
\TextAt{16}{406.576}{666.927}{\Character{106}}%
\TextAt{13}{415.618}{666.927}{\Character{20}}%
\TextAt{8}{422.356}{666.927}{,}%
\TextAt{8}{431.266}{666.927}{be}%
\TextAt{8}{448.598}{666.927}{the}%
\TextAt{8}{470.163}{666.927}{extension}%
\TextAt{8}{523.337}{666.927}{of}%
\TextAt{8}{79.226}{649.881}{valuation}%
\TextAt{8}{130.94}{649.881}{rings,}%
\TextAt{8}{163.673}{649.881}{resp.}%
\TextAt{8}{192.338}{649.881}{residue}%
\TextAt{8}{232.706}{649.881}{class}%
\TextAt{8}{260.784}{649.881}{fields,}%
\TextAt{8}{294.809}{649.881}{of}%
\TextAt{13}{308.78}{649.881}{L}%
\TextAt{16}{321.158}{649.881}{\Character{106}}%
\TextAt{13}{328.892}{649.881}{K}%
\TextAt{8}{339.71}{649.881}{.}%
\TextAt{8}{350.082}{649.881}{Suppose}%
\TextAt{8}{396.631}{649.881}{that}%
\TextAt{9}{422.641}{649.881}{\Character{40}}%
\TextAt{13}{427.193}{649.881}{n\Character{59}}%
\TextAt{13}{439.399}{649.881}{p}%
\TextAt{9}{445.3}{649.881}{\Character{41}}%
\TextAt{9}{454.254}{649.881}{\Character{61}}%
\TextAt{9}{467.764}{649.881}{1}%
\TextAt{8}{473.637}{649.881}{.}%
\TextAt{8}{484.009}{649.881}{Then}%
\TextAt{8}{515.207}{649.881}{one}%
\TextAt{8}{79.226}{632.836}{has:}%
\TextAt{8}{89.579}{612.801}{(i)}%
\TextAt{8}{107.918}{612.801}{The}%
\TextAt{8}{133.377}{612.801}{extension}%
\TextAt{13}{186.547}{612.801}{L}%
\TextAt{16}{200.225}{612.801}{\Character{106}}%
\TextAt{13}{209.26}{612.801}{K}%
\TextAt{8}{225.38}{612.801}{is}%
\TextAt{8}{238.548}{612.801}{unramified}%
\TextAt{8}{298.485}{612.801}{of}%
\TextAt{8}{313.219}{612.801}{degree}%
\TextAt{13}{351.039}{612.801}{f}%
\TextAt{8}{358.085}{612.801}{,}%
\TextAt{8}{366.99}{612.801}{where}%
\TextAt{13}{402.204}{612.801}{f}%
\TextAt{8}{414.552}{612.801}{is}%
\TextAt{8}{427.72}{612.801}{the}%
\TextAt{8}{449.281}{612.801}{smallest}%
\TextAt{8}{495.685}{612.801}{natural}%
\TextAt{8}{107.918}{595.755}{number}%
\TextAt{8}{150.515}{595.755}{such}%
\TextAt{8}{176.912}{595.755}{that}%
\TextAt{13}{202.281}{595.755}{q}%
\TextAt{21}{207.9}{600.093}{f}%
\NativeMathAt{216.666}{595.755}{11.955}{\equiv}%
\TextAt{9}{229.285}{595.755}{1}%
\TextAt{8}{238.459}{595.755}{mod}%
\TextAt{13}{263.897}{595.755}{n}%
\TextAt{8}{270.884}{595.755}{.}%
\TextAt{8}{86.327}{578.709}{(ii)}%
\TextAt{8}{107.918}{578.709}{The}%
\TextAt{8}{132.245}{578.709}{Galois}%
\TextAt{8}{168.431}{578.709}{group}%
\TextAt{13}{201.88}{578.709}{G}%
\TextAt{9}{211.113}{578.709}{\Character{40}}%
\TextAt{13}{215.666}{578.709}{L}%
\TextAt{16}{227.416}{578.709}{\Character{106}}%
\TextAt{13}{234.523}{578.709}{K}%
\TextAt{9}{245.34}{578.709}{\Character{41}}%
\TextAt{8}{254.063}{578.709}{is}%
\TextAt{8}{266.1}{578.709}{canonically}%
\TextAt{8}{327.189}{578.709}{isomorphic}%
\TextAt{8}{386.711}{578.709}{to}%
\TextAt{13}{401.294}{578.709}{G}%
\TextAt{9}{410.528}{578.709}{\Character{40}}%
\TextAt{13}{415.08}{578.709}{\Character{21}}%
\TextAt{16}{425.695}{578.709}{\Character{106}}%
\TextAt{13}{432.801}{578.709}{\Character{20}}%
\TextAt{9}{439.54}{578.709}{\Character{41}}%
\TextAt{8}{448.263}{578.709}{and}%
\TextAt{8}{471.298}{578.709}{is}%
\TextAt{8}{483.335}{578.709}{generated}%
\TextAt{8}{107.918}{561.663}{by}%
\TextAt{8}{124.178}{561.663}{the}%
\TextAt{8}{144.334}{561.663}{automorphism}%
\TextAt{13}{221.803}{561.663}{\Character{39}}%
\TextAt{9}{232.793}{561.663}{\Character{58}}%
\TextAt{13}{239.365}{561.663}{\Character{16}}%
\NativeArrowAt{248.702}{561.663}{11.955}{11.955}{\mapsto}%
\TextAt{13}{263.978}{561.663}{\Character{16}}%
\TextAt{21}{269.993}{566.001}{q}%
\TextAt{8}{274.558}{561.663}{.}%
\TextAt{8}{83.076}{544.617}{(iii)}%
\TextAt{16}{107.918}{544.617}{\Character{79}}%
\TextAt{9}{121.089}{544.617}{\Character{61}}%
\TextAt{14}{133.515}{544.617}{o}%
\TextAt{9}{139.359}{544.617}{\Character{91}}%
\TextAt{13}{142.61}{544.617}{\Character{16}}%
\TextAt{9}{148.626}{544.617}{\Character{93}}%
\TextAt{8}{151.878}{544.617}{.}%
\TextAt{19}{79.226}{509.638}{Proof.}%
\TextAt{8}{116}{509.638}{(i)}%
\TextAt{8}{132.043}{509.638}{If}%
\TextAt{13}{143.519}{509.638}{\Character{30}}%
\TextAt{9}{150.447}{509.638}{\Character{40}}%
\TextAt{13}{155}{509.638}{X}%
\TextAt{9}{165.655}{509.638}{\Character{41}}%
\TextAt{8}{173.888}{509.638}{is}%
\TextAt{8}{185.436}{509.638}{the}%
\TextAt{8}{205.376}{509.638}{minimal}%
\TextAt{8}{250.685}{509.638}{polynomial}%
\TextAt{8}{310.962}{509.638}{of}%
\TextAt{13}{324.076}{509.638}{\Character{16}}%
\TextAt{8}{333.773}{509.638}{over}%
\TextAt{13}{358.603}{509.638}{K}%
\TextAt{8}{369.421}{509.638}{,}%
\TextAt{8}{376.397}{509.638}{then}%
\TextAt{8}{402.841}{509.638}{the}%
\TextAt{8}{422.781}{509.638}{reduction}%
\TextAt{9}{476.107}{512.793}{\Character{22}}%
\TextAt{13}{474.594}{509.638}{\Character{30}}%
\TextAt{9}{481.522}{509.638}{\Character{40}}%
\TextAt{13}{486.074}{509.638}{X}%
\TextAt{9}{496.729}{509.638}{\Character{41}}%
\TextAt{8}{504.963}{509.638}{is}%
\TextAt{8}{516.51}{509.638}{the}%
\TextAt{8}{79.226}{492.592}{minimal}%
\TextAt{8}{125.739}{492.592}{polynomial}%
\TextAt{8}{187.219}{492.592}{of}%
\TextAt{9}{202.593}{495.747}{\Character{22}}%
\TextAt{13}{201.536}{492.592}{\Character{16}}%
\TextAt{9}{212.555}{492.592}{\Character{61}}%
\TextAt{13}{226.663}{492.592}{\Character{16}}%
\TextAt{8}{235.999}{492.592}{mod}%
\TextAt{14}{261.437}{492.592}{P}%
\TextAt{8}{276.215}{492.592}{over}%
\TextAt{13}{302.248}{492.592}{\Character{20}}%
\TextAt{8}{308.986}{492.592}{.}%
\TextAt{8}{320.396}{492.592}{Indeed,}%
\TextAt{8}{362.911}{492.592}{being}%
\TextAt{8}{395.436}{492.592}{a}%
\TextAt{8}{406.179}{492.592}{divisor}%
\TextAt{8}{445.279}{492.592}{of}%
\TextAt{13}{459.596}{492.592}{X}%
\TextAt{21}{470.252}{496.931}{n}%
\TextAt{16}{479.217}{492.592}{\Character{0}}%
\TextAt{9}{491.846}{492.592}{1}%
\TextAt{8}{497.699}{492.592}{,}%
\TextAt{9}{507.595}{495.747}{\Character{22}}%
\TextAt{13}{506.082}{492.592}{\Character{30}}%
\TextAt{9}{513.01}{492.592}{\Character{40}}%
\TextAt{13}{517.562}{492.592}{X}%
\TextAt{9}{528.217}{492.592}{\Character{41}}%
\TextAt{8}{79.226}{475.546}{is}%
\TextAt{8}{92.192}{475.546}{separable}%
\TextAt{8}{144.837}{475.546}{and}%
\TextAt{8}{168.802}{475.546}{by}%
\TextAt{8}{186.263}{475.546}{Hensel’s}%
\TextAt{8}{232.775}{475.546}{lemma}%
\TextAt{8}{271.696}{475.546}{cannot}%
\TextAt{8}{311.274}{475.546}{split}%
\TextAt{8}{338.55}{475.546}{into}%
\TextAt{8}{363.495}{475.546}{factors.}%
\TextAt{13}{409.766}{475.546}{\Character{30}}%
\TextAt{8}{421.793}{475.546}{and}%
\TextAt{9}{447.271}{478.701}{\Character{22}}%
\TextAt{13}{445.758}{475.546}{\Character{30}}%
\TextAt{8}{457.785}{475.546}{have}%
\TextAt{8}{485.982}{475.546}{the}%
\TextAt{8}{507.341}{475.546}{same}%
\TextAt{8}{79.226}{458.5}{degree,}%
\TextAt{8}{119.63}{458.5}{so}%
\TextAt{8}{134.589}{458.5}{that}%
\TextAt{9}{160.547}{458.5}{\Character{91}}%
\TextAt{13}{163.798}{458.5}{L}%
\TextAt{9}{176.088}{458.5}{\Character{58}}%
\TextAt{13}{183.664}{458.5}{K}%
\TextAt{9}{194.482}{458.5}{\Character{93}}%
\TextAt{9}{202.038}{458.5}{\Character{61}}%
\TextAt{9}{215.463}{458.5}{\Character{91}}%
\TextAt{13}{218.739}{458.5}{\Character{20}}%
\TextAt{9}{225.477}{458.5}{\Character{40}}%
\TextAt{9}{231.087}{461.655}{\Character{22}}%
\TextAt{13}{230.03}{458.5}{\Character{16}}%
\TextAt{9}{236.045}{458.5}{\Character{41}}%
\TextAt{9}{244.904}{458.5}{\Character{58}}%
\TextAt{13}{252.499}{458.5}{\Character{20}}%
\TextAt{9}{259.237}{458.5}{\Character{93}}%
\TextAt{9}{266.793}{458.5}{\Character{61}}%
\TextAt{9}{280.218}{458.5}{\Character{91}}%
\TextAt{13}{283.495}{458.5}{\Character{21}}%
\TextAt{9}{294.648}{458.5}{\Character{58}}%
\TextAt{13}{302.224}{458.5}{\Character{20}}%
\TextAt{9}{308.963}{458.5}{\Character{93}}%
\TextAt{9}{316.519}{458.5}{\Character{61}\Character{58}}%
\TextAt{13}{333.22}{458.5}{f}%
\TextAt{8}{340.266}{458.5}{.}%
\TextAt{13}{350.482}{458.5}{L}%
\TextAt{16}{362.771}{458.5}{\Character{106}}%
\TextAt{13}{370.416}{458.5}{K}%
\TextAt{8}{385.721}{458.5}{is}%
\TextAt{8}{398.074}{458.5}{therefore}%
\TextAt{8}{447.762}{458.5}{unramified.}%
\TextAt{8}{512.613}{458.5}{The}%
\TextAt{8}{79.226}{441.454}{polynomial}%
\TextAt{13}{139.65}{441.454}{X}%
\TextAt{21}{150.305}{445.793}{n}%
\TextAt{16}{158.456}{441.454}{\Character{0}}%
\TextAt{9}{170.269}{441.454}{1}%
\TextAt{8}{179.95}{441.454}{splits}%
\TextAt{8}{210.569}{441.454}{over}%
\TextAt{16}{235.545}{441.454}{\Character{79}}%
\TextAt{8}{249.223}{441.454}{and}%
\TextAt{8}{271.916}{441.454}{thus}%
\TextAt{8}{297.598}{441.454}{(because}%
\TextAt{9}{345.385}{441.454}{\Character{40}}%
\TextAt{13}{349.937}{441.454}{n\Character{59}}%
\TextAt{13}{362.143}{441.454}{p}%
\TextAt{9}{368.044}{441.454}{\Character{41}}%
\TextAt{9}{375.899}{441.454}{\Character{61}}%
\TextAt{9}{388.32}{441.454}{1}%
\TextAt{8}{394.196}{441.454}{)}%
\TextAt{8}{402.579}{441.454}{over}%
\TextAt{13}{427.555}{441.454}{\Character{21}}%
\TextAt{8}{438.212}{441.454}{into}%
\TextAt{8}{461.885}{441.454}{distinct}%
\TextAt{8}{504.149}{441.454}{linear}%
\TextAt{8}{79.226}{424.408}{factors,}%
\TextAt{8}{120.626}{424.408}{so}%
\TextAt{8}{135.023}{424.408}{that}%
\TextAt{13}{160.42}{424.408}{\Character{21}}%
\TextAt{9}{170.616}{424.408}{\Character{61}}%
\TextAt{18}{183.089}{424.408}{F}%
\TextAt{21}{190.395}{422.106}{q}%
\TextAt{22}{194.461}{424.375}{f}%
\TextAt{8}{203.658}{424.408}{contains}%
\TextAt{8}{249.606}{424.408}{the}%
\TextAt{8}{269.79}{424.408}{group}%
\TextAt{13}{302.993}{424.408}{\Character{22}}%
\TextAt{21}{310.036}{422.615}{n}%
\TextAt{8}{319.597}{424.408}{of}%
\TextAt{13}{332.955}{424.408}{n}%
\TextAt{8}{339.943}{424.408}{-th}%
\TextAt{8}{358.823}{424.408}{roots}%
\TextAt{8}{388.512}{424.408}{of}%
\TextAt{8}{401.869}{424.408}{unity}%
\TextAt{8}{432.466}{424.408}{and}%
\TextAt{8}{455.257}{424.408}{is}%
\TextAt{8}{467.048}{424.408}{generated}%
\TextAt{8}{520.408}{424.408}{by}%
\TextAt{8}{79.226}{407.362}{it.}%
\TextAt{8}{96.047}{407.362}{Consequently}%
\TextAt{13}{168.816}{407.362}{f}%
\TextAt{8}{179.948}{407.362}{is}%
\TextAt{8}{191.901}{407.362}{the}%
\TextAt{8}{212.246}{407.362}{smallest}%
\TextAt{8}{257.434}{407.362}{number}%
\TextAt{8}{300.219}{407.362}{such}%
\TextAt{8}{326.804}{407.362}{that}%
\TextAt{13}{352.362}{407.362}{\Character{22}}%
\TextAt{21}{359.405}{405.569}{n}%
\NativeMathAt{368.684}{407.362}{11.955}{\subseteq}%
\TextAt{18}{381.625}{407.362}{F}%
\TextAt{3}{388.93}{411.701}{\Character{3}}%
\TextAt{21}{388.93}{403.368}{q}%
\TextAt{22}{392.997}{405.637}{f}%
\TextAt{8}{398.269}{407.362}{,}%
\TextAt{8}{405.654}{407.362}{i.e.,}%
\TextAt{8}{427.995}{407.362}{such}%
\TextAt{8}{454.581}{407.362}{that}%
\TextAt{13}{480.138}{407.362}{n}%
\TextAt{16}{490.768}{407.362}{\Character{106}}%
\TextAt{13}{497.732}{407.362}{q}%
\TextAt{21}{503.351}{411.701}{f}%
\TextAt{16}{511.581}{407.362}{\Character{0}}%
\TextAt{9}{523.665}{407.362}{1}%
\TextAt{8}{529.518}{407.362}{.}%
\TextAt{8}{79.226}{390.316}{This}%
\TextAt{8}{105.946}{390.316}{shows}%
\TextAt{8}{139.564}{390.316}{(i).}%
\TextAt{8}{160.374}{390.316}{(ii)}%
\TextAt{8}{179.885}{390.316}{results}%
\TextAt{8}{217.077}{390.316}{trivially}%
\TextAt{8}{261.311}{390.316}{from}%
\TextAt{8}{288.952}{390.316}{this.}%
\TextAt{8}{97.159}{370.282}{(iii)}%
\TextAt{8}{119.993}{370.282}{Since}%
\TextAt{13}{150.622}{370.282}{L}%
\TextAt{16}{162.028}{370.282}{\Character{106}}%
\TextAt{13}{168.792}{370.282}{K}%
\TextAt{8}{183.578}{370.282}{is}%
\TextAt{8}{195.413}{370.282}{unramified,}%
\TextAt{8}{257.287}{370.282}{we}%
\TextAt{8}{274.586}{370.282}{have}%
\TextAt{14}{301.652}{370.282}{p}%
\TextAt{16}{307.635}{370.282}{\Character{79}}%
\TextAt{9}{320.927}{370.282}{\Character{61}}%
\TextAt{14}{333.474}{370.282}{P}%
\TextAt{8}{343.367}{370.282}{,}%
\TextAt{8}{350.605}{370.282}{and}%
\TextAt{8}{373.439}{370.282}{since}%
\TextAt{9}{402.179}{370.282}{1}%
\TextAt{13}{408.032}{370.282}{\Character{59}}%
\TextAt{13}{413.256}{370.282}{\Character{16}}%
\TextAt{13}{419.27}{370.282}{\Character{59}}%
\TextAt{13}{424.494}{370.282}{\Character{58}}%
\TextAt{13}{429.731}{370.282}{\Character{58}}%
\TextAt{13}{434.967}{370.282}{\Character{58}}%
\TextAt{13}{440.192}{370.282}{\Character{59}}%
\TextAt{13}{445.428}{370.282}{\Character{16}}%
\TextAt{21}{451.528}{374.62}{f}%
\TextAt{3}{456.475}{374.62}{\Character{0}}%
\TextAt{17}{463.062}{374.62}{1}%
\TextAt{8}{471.763}{370.282}{represents}%
\TextAt{8}{526.911}{370.282}{a}%
\TextAt{8}{79.226}{353.236}{basis}%
\TextAt{8}{107.966}{353.236}{of}%
\TextAt{13}{121.296}{353.236}{\Character{21}}%
\TextAt{16}{131.446}{353.236}{\Character{106}}%
\TextAt{13}{138.087}{353.236}{\Character{20}}%
\TextAt{8}{144.826}{353.236}{,}%
\TextAt{8}{151.975}{353.236}{we}%
\TextAt{8}{169.203}{353.236}{have}%
\TextAt{16}{196.197}{353.236}{\Character{79}}%
\TextAt{9}{209.368}{353.236}{\Character{61}}%
\TextAt{14}{221.793}{353.236}{o}%
\TextAt{9}{227.638}{353.236}{\Character{91}}%
\TextAt{13}{230.889}{353.236}{\Character{16}}%
\TextAt{9}{236.905}{353.236}{\Character{93}}%
\TextAt{9}{242.799}{353.236}{\Character{43}}%
\TextAt{14}{254.575}{353.236}{p}%
\TextAt{16}{260.558}{353.236}{\Character{79}}%
\TextAt{8}{270.407}{353.236}{,}%
\TextAt{8}{277.557}{353.236}{and}%
\TextAt{16}{300.319}{353.236}{\Character{79}}%
\TextAt{9}{313.49}{353.236}{\Character{61}}%
\TextAt{14}{325.915}{353.236}{o}%
\TextAt{9}{331.76}{353.236}{\Character{91}}%
\TextAt{13}{335.011}{353.236}{\Character{16}}%
\TextAt{9}{341.027}{353.236}{\Character{93}}%
\TextAt{8}{348.176}{353.236}{by}%
\TextAt{8}{364.435}{353.236}{Nakayama’s}%
\TextAt{8}{429.878}{353.236}{lemma.}%
\Vector{q 0 G 0 g 0.399 w 452.334 -420.764 m 452.334 -412.695 l S Q}%
\Vector{q 0 G 0 g 0.399 w 452.534 -412.894 m 459.707 -412.894 l S Q}%
\Vector{q 0 G 0 g 0.399 w 452.534 -420.565 m 459.707 -420.565 l S Q}%
\Vector{q 0 G 0 g 0.399 w 459.906 -420.764 m 459.906 -412.695 l S Q}%
\TextAt{12}{79.226}{318.257}{Proposition}%
\TextAt{12}{151.937}{318.257}{(7.13).}%
\ErrataAt{193.95}{318.257}{341.05}{Let $\zeta$ be a primitive $p^m$-th root of unity, with $m\ge1$. Then one has:}%
\TextAt{8}{89.579}{298.222}{(i)}%
\TextAt{18}{107.918}{298.222}{Q}%
\TextAt{21}{117.217}{296.429}{p}%
\TextAt{9}{121.978}{298.222}{\Character{40}}%
\TextAt{13}{126.53}{298.222}{\Character{16}}%
\TextAt{9}{132.546}{298.222}{\Character{41}}%
\TextAt{16}{140.419}{298.222}{\Character{106}}%
\TextAt{18}{147.061}{298.222}{Q}%
\TextAt{21}{156.359}{296.429}{p}%
\TextAt{8}{165.017}{298.222}{is}%
\TextAt{8}{176.781}{298.222}{totally}%
\TextAt{8}{214.189}{298.222}{ramified}%
\TextAt{8}{259.714}{298.222}{of}%
\TextAt{8}{273.044}{298.222}{degree}%
\TextAt{13}{309.46}{298.222}{\Character{39}}%
\TextAt{9}{317.128}{298.222}{\Character{40}}%
\TextAt{13}{321.681}{298.222}{p}%
\TextAt{21}{327.556}{302.561}{m}%
\TextAt{9}{335.544}{298.222}{\Character{41}}%
\TextAt{9}{343.399}{298.222}{\Character{61}}%
\TextAt{9}{355.82}{298.222}{\Character{40}}%
\TextAt{13}{360.395}{298.222}{p}%
\TextAt{16}{368.927}{298.222}{\Character{0}}%
\TextAt{9}{380.882}{298.222}{1\Character{41}}%
\TextAt{13}{391.288}{298.222}{p}%
\TextAt{21}{397.163}{302.561}{m}%
\TextAt{3}{404.653}{302.561}{\Character{0}}%
\TextAt{17}{411.24}{302.561}{1}%
\TextAt{8}{415.972}{298.222}{.}%
\TextAt{8}{86.327}{281.176}{(ii)}%
\TextAt{13}{107.918}{281.176}{G}%
\TextAt{9}{117.152}{281.176}{\Character{40}}%
\TextAt{18}{121.704}{281.176}{Q}%
\TextAt{21}{131.003}{279.383}{p}%
\TextAt{9}{135.764}{281.176}{\Character{40}}%
\TextAt{13}{140.316}{281.176}{\Character{16}}%
\TextAt{9}{146.332}{281.176}{\Character{41}}%
\TextAt{16}{154.205}{281.176}{\Character{106}}%
\TextAt{18}{160.847}{281.176}{Q}%
\TextAt{21}{170.145}{279.383}{p}%
\TextAt{9}{174.906}{281.176}{\Character{41}}%
\NativeMathAt{182.779}{281.176}{11.955}{\cong}%
\TextAt{9}{195.398}{281.176}{\Character{40}}%
\TextAt{18}{199.951}{281.176}{Z}%
\TextAt{9}{207.921}{281.176}{\Character{47}}%
\TextAt{13}{213.774}{281.176}{p}%
\TextAt{21}{219.649}{285.515}{m}%
\TextAt{18}{227.638}{281.176}{Z}%
\TextAt{9}{235.608}{281.176}{\Character{41}}%
\TextAt{3}{240.16}{285.515}{\Character{3}}%
\TextAt{8}{244.893}{281.176}{.}%
\TextAt{8}{83.076}{264.13}{(iii)}%
\TextAt{18}{107.918}{264.13}{Z}%
\TextAt{21}{115.889}{262.337}{p}%
\TextAt{9}{120.65}{264.13}{\Character{91}}%
\TextAt{13}{123.901}{264.13}{\Character{16}}%
\TextAt{9}{129.917}{264.13}{\Character{93}}%
\TextAt{8}{137.066}{264.13}{is}%
\TextAt{8}{148.83}{264.13}{the}%
\TextAt{8}{168.986}{264.13}{valuation}%
\TextAt{8}{220.059}{264.13}{ring}%
\TextAt{8}{244.124}{264.13}{of}%
\TextAt{18}{257.454}{264.13}{Q}%
\TextAt{21}{266.753}{262.337}{p}%
\TextAt{9}{271.514}{264.13}{\Character{40}}%
\TextAt{13}{276.066}{264.13}{\Character{16}}%
\TextAt{9}{282.082}{264.13}{\Character{41}}%
\TextAt{8}{286.634}{264.13}{.}%
\TextAt{8}{83.398}{247.084}{(iv)}%
\TextAt{9}{107.918}{247.084}{1}%
\TextAt{16}{116.428}{247.084}{\Character{0}}%
\TextAt{13}{128.383}{247.084}{\Character{16}}%
\TextAt{8}{138.296}{247.084}{is}%
\TextAt{8}{150.06}{247.084}{a}%
\TextAt{8}{159.816}{247.084}{prime}%
\TextAt{8}{192.979}{247.084}{element}%
\TextAt{8}{236.221}{247.084}{of}%
\TextAt{18}{249.551}{247.084}{Z}%
\TextAt{21}{257.521}{245.291}{p}%
\TextAt{9}{262.282}{247.084}{\Character{91}}%
\TextAt{13}{265.534}{247.084}{\Character{16}}%
\TextAt{9}{271.549}{247.084}{\Character{93}}%
\TextAt{8}{278.698}{247.084}{with}%
\TextAt{8}{305.359}{247.084}{norm}%
\TextAt{13}{335.928}{247.084}{p}%
\TextAt{8}{341.803}{247.084}{.}%
\TextAt{19}{79.226}{212.105}{Proof.}%
\TextAt{13}{116}{212.105}{\Character{24}}%
\TextAt{9}{125.012}{212.105}{\Character{61}}%
\TextAt{13}{137.437}{212.105}{\Character{16}}%
\TextAt{21}{143.453}{216.444}{p}%
\TextAt{22}{147.715}{219.257}{m}%
\TextAt{23}{154.253}{219.257}{\Character{0}}%
\TextAt{24}{160.009}{219.257}{1}%
\TextAt{8}{168.556}{212.105}{is}%
\TextAt{8}{180.32}{212.105}{a}%
\TextAt{8}{190.075}{212.105}{primitive}%
\TextAt{13}{240.156}{212.105}{p}%
\TextAt{8}{246.031}{212.105}{-th}%
\TextAt{8}{264.884}{212.105}{root}%
\TextAt{8}{289.93}{212.105}{of}%
\TextAt{8}{303.26}{212.105}{unity,}%
\TextAt{8}{336.101}{212.105}{i.e.,}%
\TextAt{13}{120.673}{183.104}{\Character{24}}%
\TextAt{21}{126.363}{188.041}{p}%
\TextAt{3}{130.626}{188.041}{\Character{0}}%
\TextAt{17}{137.213}{188.041}{1}%
\TextAt{9}{144.602}{183.104}{\Character{43}}%
\TextAt{13}{156.363}{183.104}{\Character{24}}%
\TextAt{21}{162.053}{188.041}{p}%
\TextAt{3}{166.316}{188.041}{\Character{0}}%
\TextAt{17}{172.903}{188.041}{2}%
\TextAt{9}{180.292}{183.104}{\Character{43}}%
\TextAt{16}{192.053}{183.104}{\Character{1}}%
\TextAt{16}{197.349}{183.104}{\Character{1}}%
\TextAt{16}{202.657}{183.104}{\Character{1}}%
\TextAt{9}{208.657}{183.104}{\Character{43}}%
\TextAt{9}{220.409}{183.104}{1}%
\TextAt{9}{229.579}{183.104}{\Character{61}}%
\TextAt{9}{241.988}{183.104}{0}%
\TextAt{13}{247.871}{183.104}{\Character{59}}%
\TextAt{8}{265.07}{183.104}{hence}%
\TextAt{13}{305.634}{183.104}{\Character{16}}%
\TextAt{17}{311.65}{188.041}{\Character{40}}%
\TextAt{21}{314.943}{188.041}{p}%
\TextAt{3}{319.206}{188.041}{\Character{0}}%
\TextAt{17}{325.792}{188.041}{1\Character{41}}%
\TextAt{21}{333.32}{188.041}{p}%
\TextAt{22}{337.582}{190.853}{m}%
\TextAt{23}{344.12}{190.853}{\Character{0}}%
\TextAt{24}{349.876}{190.853}{1}%
\TextAt{9}{357.182}{183.104}{\Character{43}}%
\TextAt{13}{368.944}{183.104}{\Character{16}}%
\TextAt{17}{374.959}{188.041}{\Character{40}}%
\TextAt{21}{378.252}{188.041}{p}%
\TextAt{3}{382.515}{188.041}{\Character{0}}%
\TextAt{17}{389.102}{188.041}{2\Character{41}}%
\TextAt{21}{396.629}{188.041}{p}%
\TextAt{22}{400.892}{190.853}{m}%
\TextAt{23}{407.43}{190.853}{\Character{0}}%
\TextAt{24}{413.186}{190.853}{1}%
\TextAt{9}{420.492}{183.104}{\Character{43}}%
\TextAt{16}{432.253}{183.104}{\Character{1}}%
\TextAt{16}{437.549}{183.104}{\Character{1}}%
\TextAt{16}{442.857}{183.104}{\Character{1}}%
\TextAt{9}{448.858}{183.104}{\Character{43}}%
\TextAt{9}{460.61}{183.104}{1}%
\TextAt{9}{469.768}{183.104}{\Character{61}}%
\TextAt{9}{482.189}{183.104}{0}%
\TextAt{13}{488.071}{183.104}{\Character{58}}%
\TextAt{8}{79.226}{154.103}{Denoting}%
\TextAt{8}{130.032}{154.103}{by}%
\TextAt{13}{146.527}{154.103}{\Character{30}}%
\TextAt{8}{157.589}{154.103}{the}%
\TextAt{8}{177.981}{154.103}{polynomial}%
\TextAt{8}{238.71}{154.103}{on}%
\TextAt{8}{255.205}{154.103}{the}%
\TextAt{8}{275.597}{154.103}{left,}%
\TextAt{13}{299.623}{154.103}{\Character{16}}%
\TextAt{16}{308.456}{154.103}{\Character{0}}%
\TextAt{9}{320.572}{154.103}{1}%
\TextAt{8}{330.558}{154.103}{is}%
\TextAt{8}{342.558}{154.103}{a}%
\TextAt{8}{352.549}{154.103}{root}%
\TextAt{8}{377.831}{154.103}{of}%
\TextAt{8}{391.397}{154.103}{the}%
\TextAt{8}{411.789}{154.103}{equation}%
\TextAt{13}{459.834}{154.103}{\Character{30}}%
\TextAt{9}{466.762}{154.103}{\Character{40}}%
\TextAt{13}{471.315}{154.103}{X}%
\TextAt{9}{484.787}{154.103}{\Character{43}}%
\TextAt{9}{496.694}{154.103}{1\Character{41}}%
\TextAt{9}{510.814}{154.103}{\Character{61}}%
\TextAt{9}{523.629}{154.103}{0}%
\TextAt{8}{529.518}{154.103}{.}%
\TextAt{8}{79.226}{137.057}{But}%
\TextAt{8}{102.467}{137.057}{this}%
\TextAt{8}{125.289}{137.057}{is}%
\TextAt{8}{137.053}{137.057}{irreducible}%
\TextAt{8}{194.928}{137.057}{because}%
\TextAt{8}{238.23}{137.057}{it}%
\TextAt{8}{249.934}{137.057}{satisfies}%
\TextAt{8}{293.044}{137.057}{Eisenstein’s}%
\TextAt{8}{356.467}{137.057}{criterion:}%
\TextAt{13}{407.846}{137.057}{\Character{30}}%
\TextAt{9}{414.774}{137.057}{\Character{40}1\Character{41}}%
\TextAt{9}{433.018}{137.057}{\Character{61}}%
\TextAt{13}{445.478}{137.057}{p}%
\TextAt{8}{455.25}{137.057}{and}%
\TextAt{13}{188.959}{103.304}{\Character{30}}%
\TextAt{9}{195.887}{103.304}{\Character{40}}%
\TextAt{13}{200.44}{103.304}{X}%
\TextAt{9}{211.095}{103.304}{\Character{41}}%
\TextAt{9}{218.95}{103.304}{\Character{61}}%
\TextAt{13}{237.293}{111.392}{X}%
\TextAt{21}{247.948}{115.731}{p}%
\TextAt{22}{252.211}{118.543}{m}%
\TextAt{16}{262.402}{111.392}{\Character{0}}%
\TextAt{9}{274.357}{111.392}{1}%
\Vector{q 0 G 0 g 0.478 w 160.589 -667.707 m 212.915 -667.707 l S Q}%
\TextAt{13}{232.589}{95.044}{X}%
\TextAt{21}{243.244}{98.498}{p}%
\TextAt{22}{247.507}{100.767}{m}%
\TextAt{23}{254.045}{100.767}{\Character{0}}%
\TextAt{24}{259.801}{100.767}{1}%
\TextAt{16}{267.107}{95.044}{\Character{0}}%
\TextAt{9}{279.062}{95.044}{1}%
\NativeMathAt{289.431}{103.304}{11.955}{\equiv}%
\TextAt{9}{302.05}{103.304}{\Character{40}}%
\TextAt{13}{306.603}{103.304}{X}%
\TextAt{16}{319.915}{103.304}{\Character{0}}%
\TextAt{9}{331.87}{103.304}{1\Character{41}}%
\TextAt{21}{342.275}{108.241}{p}%
\TextAt{22}{346.538}{111.053}{m}%
\TextAt{23}{353.076}{111.053}{\Character{0}}%
\TextAt{24}{358.832}{111.053}{1}%
\TextAt{17}{362.983}{108.241}{\Character{40}}%
\TextAt{21}{366.276}{108.241}{p}%
\TextAt{3}{370.539}{108.241}{\Character{0}}%
\TextAt{17}{377.125}{108.241}{1\Character{41}}%
\TextAt{8}{388.472}{103.304}{mod}%
\TextAt{13}{413.91}{103.304}{p\Character{58}}%
\TextAt{8}{297.211}{43.033}{145}%
\Destination{79.23}{629.725}{nav.Item.44}%
\Destination{79.23}{593.425}{nav.Item.45}%
\Destination{79.23}{559.225}{nav.Item.46}%
\Destination{79.23}{312.265}{nav.Item.47}%
\Destination{79.23}{294.875}{nav.Item.48}%
\Destination{79.23}{277.835}{nav.Item.49}%
\Destination{79.23}{260.785}{nav.Item.50}%
\Destination{78.23}{809.815}{nav.page.145}%
\Destination{79.23}{778.165}{nav.sectionz2az.19}%
\Destination{79.36}{776.155}{nav.print.158}%
\Destination{79.36}{333.255}{nav.print.159}%
\Destination{79.36}{680.955}{nav.number.2.7.12}%
\Destination{79.36}{332.255}{nav.number.2.7.13}%
\Destination{79.226}{335.257}{errata.79}%
}

% Page 158 (body page 146)
\Page{%
\TextAt{2}{79.226}{796.943}{II.}%
\TextAt{2}{94.073}{796.943}{The}%
\TextAt{2}{116.492}{796.943}{Theory}%
\TextAt{2}{154.401}{796.943}{of}%
\TextAt{2}{166.826}{796.943}{Valuations}%
\Vector{q 0 G 0 g 0.299 w 7.226 18.729 m 460.769 18.729 l S Q}%
\TextAt{8}{79.226}{761.124}{It}%
\TextAt{8}{91.144}{761.124}{follows}%
\TextAt{8}{128.826}{761.124}{that}%
\TextAt{9}{153.441}{761.124}{\Character{91}}%
\TextAt{18}{156.692}{761.124}{Q}%
\TextAt{21}{165.991}{759.33}{p}%
\TextAt{9}{170.752}{761.124}{\Character{40}}%
\TextAt{13}{175.304}{761.124}{\Character{16}}%
\TextAt{9}{181.32}{761.124}{\Character{41}}%
\TextAt{9}{189.175}{761.124}{\Character{58}}%
\TextAt{18}{195.765}{761.124}{Q}%
\TextAt{21}{205.064}{759.33}{p}%
\TextAt{9}{209.825}{761.124}{\Character{93}}%
\TextAt{9}{216.376}{761.124}{\Character{61}}%
\TextAt{13}{228.823}{761.124}{\Character{39}}%
\TextAt{9}{236.491}{761.124}{\Character{40}}%
\TextAt{13}{241.043}{761.124}{p}%
\TextAt{21}{246.919}{765.462}{m}%
\TextAt{9}{254.907}{761.124}{\Character{41}}%
\TextAt{8}{259.46}{761.124}{.}%
\TextAt{8}{267.657}{761.124}{The}%
\TextAt{8}{290.956}{761.124}{canonical}%
\TextAt{8}{341.585}{761.124}{injection}%
\TextAt{13}{388.628}{761.124}{G}%
\TextAt{9}{397.861}{761.124}{\Character{40}}%
\TextAt{18}{402.414}{761.124}{Q}%
\TextAt{21}{411.712}{759.33}{p}%
\TextAt{9}{416.473}{761.124}{\Character{40}}%
\TextAt{13}{421.025}{761.124}{\Character{16}}%
\TextAt{9}{427.041}{761.124}{\Character{41}}%
\TextAt{16}{434.914}{761.124}{\Character{106}}%
\TextAt{18}{441.556}{761.124}{Q}%
\TextAt{21}{450.854}{759.33}{p}%
\TextAt{9}{455.615}{761.124}{\Character{41}}%
\NativeArrowAt{463.489}{761.124}{11.955}{13.214}{\hookrightarrow}%
\TextAt{9}{480.024}{761.124}{\Character{40}}%
\TextAt{18}{484.576}{761.124}{Z}%
\TextAt{9}{492.546}{761.124}{\Character{47}}%
\TextAt{13}{498.399}{761.124}{p}%
\TextAt{21}{504.274}{765.462}{m}%
\TextAt{18}{512.263}{761.124}{Z}%
\TextAt{9}{520.233}{761.124}{\Character{41}}%
\TextAt{3}{524.785}{765.462}{\Character{3}}%
\TextAt{8}{529.518}{761.124}{,}%
\TextAt{13}{79.226}{744.078}{\Character{27}}%
\NativeArrowAt{90.383}{744.078}{11.955}{11.955}{\mapsto}%
\TextAt{13}{106.413}{744.078}{n}%
\TextAt{9}{113.401}{744.078}{\Character{40}}%
\TextAt{13}{117.953}{744.078}{\Character{27}}%
\TextAt{9}{125.035}{744.078}{\Character{41}}%
\TextAt{8}{129.588}{744.078}{,}%
\TextAt{8}{137.29}{744.078}{where}%
\TextAt{13}{171.541}{744.078}{\Character{27}}%
\TextAt{13}{178.607}{744.078}{\Character{16}}%
\TextAt{9}{188.714}{744.078}{\Character{61}}%
\TextAt{13}{201.893}{744.078}{\Character{16}}%
\TextAt{21}{207.909}{748.416}{n}%
\TextAt{17}{213.047}{748.416}{\Character{40}}%
\TextAt{21}{216.34}{748.416}{\Character{27}}%
\TextAt{17}{221.466}{748.416}{\Character{41}}%
\TextAt{8}{225.257}{744.078}{,}%
\TextAt{8}{232.959}{744.078}{is}%
\TextAt{8}{245.165}{744.078}{therefore}%
\TextAt{8}{294.708}{744.078}{bijective,}%
\TextAt{8}{345.006}{744.078}{since}%
\TextAt{8}{374.117}{744.078}{both}%
\TextAt{8}{402.199}{744.078}{groups}%
\TextAt{8}{440.432}{744.078}{have}%
\TextAt{8}{467.869}{744.078}{order}%
\TextAt{13}{498.881}{744.078}{\Character{39}}%
\TextAt{9}{506.549}{744.078}{\Character{40}}%
\TextAt{13}{511.102}{744.078}{p}%
\TextAt{21}{516.977}{748.416}{m}%
\TextAt{9}{524.965}{744.078}{\Character{41}}%
\TextAt{8}{529.518}{744.078}{.}%
\TextAt{8}{79.226}{727.032}{Thus}%
\TextAt{13}{198.03}{709.986}{N}%
\TextAt{25}{207.404}{708.126}{Q}%
\TextAt{22}{213.603}{707.13}{p}%
\TextAt{17}{217.942}{708.126}{\Character{40}}%
\TextAt{21}{221.235}{708.126}{\Character{16}}%
\TextAt{17}{225.504}{708.126}{\Character{41}}%
\TextAt{3}{228.797}{708.126}{\Character{106}}%
\TextAt{25}{231.149}{708.126}{Q}%
\TextAt{22}{237.348}{707.13}{p}%
\TextAt{9}{242.185}{709.986}{\Character{40}1}%
\TextAt{16}{255.247}{709.986}{\Character{0}}%
\TextAt{13}{267.202}{709.986}{\Character{16}}%
\TextAt{9}{273.218}{709.986}{\Character{41}}%
\TextAt{9}{281.073}{709.986}{\Character{61}}%
\TextAt{20}{293.517}{721.343}{\Character{89}}%
\TextAt{21}{298.592}{696.237}{\Character{27}}%
\TextAt{9}{308.793}{709.986}{\Character{40}1}%
\TextAt{16}{321.855}{709.986}{\Character{0}}%
\TextAt{13}{333.81}{709.986}{\Character{27}}%
\TextAt{13}{340.875}{709.986}{\Character{16}}%
\TextAt{9}{346.908}{709.986}{\Character{41}}%
\TextAt{9}{354.763}{709.986}{\Character{61}}%
\TextAt{13}{367.206}{709.986}{\Character{30}}%
\TextAt{9}{374.135}{709.986}{\Character{40}1\Character{41}}%
\TextAt{9}{392.379}{709.986}{\Character{61}}%
\TextAt{13}{404.838}{709.986}{p\Character{58}}%
\TextAt{8}{79.226}{679.153}{Writing}%
\TextAt{13}{123.095}{679.153}{w}%
\TextAt{8}{136.631}{679.153}{for}%
\TextAt{8}{155.465}{679.153}{the}%
\TextAt{8}{176.571}{679.153}{extension}%
\TextAt{8}{229.286}{679.153}{of}%
\TextAt{8}{243.565}{679.153}{the}%
\TextAt{8}{264.671}{679.153}{normalized}%
\TextAt{8}{325.456}{679.153}{valuation}%
\TextAt{13}{377.478}{679.153}{v}%
\TextAt{21}{383.137}{677.36}{p}%
\TextAt{8}{392.745}{679.153}{of}%
\TextAt{18}{407.024}{679.153}{Q}%
\TextAt{21}{416.323}{677.36}{p}%
\TextAt{8}{421.083}{679.153}{,}%
\TextAt{8}{429.419}{679.153}{we}%
\TextAt{8}{447.596}{679.153}{find}%
\TextAt{8}{471.954}{679.153}{furthermore}%
\TextAt{8}{79.226}{662.107}{that}%
\TextAt{13}{104.822}{662.107}{\Character{39}}%
\TextAt{9}{112.49}{662.107}{\Character{40}}%
\TextAt{13}{117.042}{662.107}{p}%
\TextAt{21}{122.917}{666.446}{m}%
\TextAt{9}{130.906}{662.107}{\Character{41}}%
\TextAt{13}{135.458}{662.107}{w}%
\TextAt{9}{144.148}{662.107}{\Character{40}}%
\TextAt{13}{148.7}{662.107}{\Character{16}}%
\TextAt{16}{157.527}{662.107}{\Character{0}}%
\TextAt{9}{169.637}{662.107}{1\Character{41}}%
\TextAt{9}{183.732}{662.107}{\Character{61}}%
\TextAt{13}{196.561}{662.107}{v}%
\TextAt{21}{202.22}{660.314}{p}%
\TextAt{9}{206.981}{662.107}{\Character{40}}%
\TextAt{13}{211.533}{662.107}{p}%
\TextAt{9}{217.408}{662.107}{\Character{41}}%
\TextAt{9}{225.657}{662.107}{\Character{61}}%
\TextAt{9}{238.461}{662.107}{1}%
\TextAt{8}{244.332}{662.107}{,}%
\TextAt{8}{251.765}{662.107}{i.e.,}%
\TextAt{18}{274.153}{662.107}{Q}%
\TextAt{21}{283.452}{660.314}{p}%
\TextAt{9}{288.213}{662.107}{\Character{40}}%
\TextAt{13}{292.765}{662.107}{\Character{16}}%
\TextAt{9}{298.781}{662.107}{\Character{41}}%
\TextAt{16}{307.04}{662.107}{\Character{106}}%
\TextAt{18}{314.068}{662.107}{Q}%
\TextAt{21}{323.366}{660.314}{p}%
\TextAt{8}{332.251}{662.107}{is}%
\TextAt{8}{344.242}{662.107}{totally}%
\TextAt{8}{381.876}{662.107}{ramified}%
\TextAt{8}{427.628}{662.107}{and}%
\TextAt{13}{450.617}{662.107}{\Character{16}}%
\TextAt{16}{459.444}{662.107}{\Character{0}}%
\TextAt{9}{471.554}{662.107}{1}%
\TextAt{8}{481.531}{662.107}{is}%
\TextAt{8}{493.521}{662.107}{a}%
\TextAt{8}{503.503}{662.107}{prime}%
\TextAt{8}{79.226}{645.061}{element}%
\TextAt{8}{121.804}{645.061}{of}%
\TextAt{18}{134.471}{645.061}{Q}%
\TextAt{21}{143.769}{643.268}{p}%
\TextAt{9}{148.53}{645.061}{\Character{40}}%
\TextAt{13}{153.083}{645.061}{\Character{16}}%
\TextAt{9}{159.098}{645.061}{\Character{41}}%
\TextAt{8}{163.651}{645.061}{.}%
\TextAt{8}{171.878}{645.061}{As}%
\TextAt{8}{188.501}{645.061}{in}%
\TextAt{8}{201.491}{645.061}{the}%
\TextAt{8}{220.984}{645.061}{proof}%
\TextAt{8}{250.889}{645.061}{of}%
\TextAt{8}{263.556}{645.061}{(6.8),}%
\TextAt{8}{294.252}{645.061}{it}%
\TextAt{8}{305.292}{645.061}{follows}%
\TextAt{8}{343.065}{645.061}{that}%
\TextAt{18}{367.77}{645.061}{Z}%
\TextAt{21}{375.74}{643.268}{p}%
\TextAt{9}{380.501}{645.061}{\Character{91}}%
\TextAt{13}{383.753}{645.061}{\Character{16}}%
\TextAt{16}{391.068}{645.061}{\Character{0}}%
\TextAt{9}{401.666}{645.061}{1\Character{93}}%
\TextAt{9}{414.064}{645.061}{\Character{61}}%
\TextAt{18}{426.517}{645.061}{Z}%
\TextAt{21}{434.488}{643.268}{p}%
\TextAt{9}{439.248}{645.061}{\Character{91}}%
\TextAt{13}{442.5}{645.061}{\Character{16}}%
\TextAt{9}{448.516}{645.061}{\Character{93}}%
\TextAt{8}{455.001}{645.061}{is}%
\TextAt{8}{466.102}{645.061}{the}%
\TextAt{8}{485.594}{645.061}{valuation}%
\TextAt{8}{79.226}{628.015}{ring}%
\TextAt{8}{103.292}{628.015}{of}%
\TextAt{18}{116.622}{628.015}{Q}%
\TextAt{21}{125.92}{626.222}{p}%
\TextAt{9}{130.681}{628.015}{\Character{40}}%
\TextAt{13}{135.234}{628.015}{\Character{16}}%
\TextAt{9}{141.249}{628.015}{\Character{41}}%
\TextAt{8}{145.802}{628.015}{.}%
\TextAt{8}{154.25}{628.015}{This}%
\TextAt{8}{180.97}{628.015}{concludes}%
\TextAt{8}{233.704}{628.015}{the}%
\TextAt{8}{253.86}{628.015}{proof.}%
\Vector{q 0 G 0 g 0.399 w 452.334 -145.985 m 452.334 -137.915 l S Q}%
\Vector{q 0 G 0 g 0.399 w 452.534 -138.114 m 459.707 -138.114 l S Q}%
\Vector{q 0 G 0 g 0.399 w 452.534 -145.785 m 459.707 -145.785 l S Q}%
\Vector{q 0 G 0 g 0.399 w 459.906 -145.985 m 459.906 -137.915 l S Q}%
\TextAt{8}{97.159}{599.014}{If}%
\TextAt{13}{108.283}{599.014}{\Character{16}}%
\TextAt{21}{113.395}{597.221}{n}%
\TextAt{8}{122.36}{599.014}{is}%
\TextAt{8}{133.556}{599.014}{a}%
\TextAt{8}{142.743}{599.014}{primitive}%
\TextAt{13}{192.254}{599.014}{n}%
\TextAt{8}{199.242}{599.014}{-th}%
\TextAt{8}{217.527}{599.014}{root}%
\TextAt{8}{242.004}{599.014}{of}%
\TextAt{8}{254.766}{599.014}{unity}%
\TextAt{8}{284.767}{599.014}{and}%
\TextAt{13}{306.961}{599.014}{n}%
\TextAt{9}{317.269}{599.014}{\Character{61}}%
\TextAt{13}{329.695}{599.014}{n}%
\TextAt{3}{336.682}{603.353}{\Character{48}}%
\TextAt{13}{339.477}{599.014}{p}%
\TextAt{21}{345.353}{603.353}{m}%
\TextAt{8}{353.341}{599.014}{,}%
\TextAt{8}{360.036}{599.014}{with}%
\TextAt{9}{386.127}{599.014}{\Character{40}}%
\TextAt{13}{390.679}{599.014}{n}%
\TextAt{3}{397.667}{603.353}{\Character{48}}%
\TextAt{13}{400.462}{599.014}{\Character{59}}%
\TextAt{13}{405.686}{599.014}{p}%
\TextAt{9}{411.581}{599.014}{\Character{41}}%
\TextAt{9}{419.436}{599.014}{\Character{61}}%
\TextAt{9}{431.857}{599.014}{1}%
\TextAt{8}{437.733}{599.014}{,}%
\TextAt{8}{444.427}{599.014}{then}%
\TextAt{8}{470.519}{599.014}{propositions}%
\TextAt{8}{79.226}{581.968}{(7.12)}%
\TextAt{8}{113.171}{581.968}{and}%
\TextAt{8}{136.044}{581.968}{(7.13)}%
\TextAt{8}{169.989}{581.968}{yield}%
\TextAt{8}{198.386}{581.968}{the}%
\TextAt{8}{218.654}{581.968}{following}%
\TextAt{8}{268.2}{581.968}{result}%
\TextAt{8}{300.889}{581.968}{for}%
\TextAt{8}{318.885}{581.968}{the}%
\TextAt{8}{339.153}{581.968}{maximal}%
\TextAt{8}{387.073}{581.968}{unramified}%
\TextAt{8}{445.717}{581.968}{and}%
\TextAt{8}{468.59}{581.968}{the}%
\TextAt{8}{488.858}{581.968}{maximal}%
\TextAt{8}{79.226}{564.922}{tamely}%
\TextAt{8}{117.925}{564.922}{ramified}%
\TextAt{8}{163.45}{564.922}{extension:}%
\ErrataAt{199}{535.921}{257}{$\mathbb Q_p\subseteq T=\mathbb Q_p(\zeta_{n^{\prime}})\subseteq V\subseteq\mathbb Q_p(\zeta_n)$,}%
\ErrataAt{79.19}{516.921}{455.81}{where $V=T$ if $m=0$ and $V=T(\zeta_p)$ if $m\ge1$.}%
\TextAt{12}{79.226}{497.954}{Exercise}%
\TextAt{12}{133.217}{497.954}{II.7.1.}%
\TextAt{8}{172.877}{497.954}{The}%
\TextAt{8}{197.618}{497.954}{maximal}%
\TextAt{8}{246.113}{497.954}{unramified}%
\TextAt{8}{305.333}{497.954}{extension}%
\TextAt{8}{357.785}{497.954}{of}%
\TextAt{18}{371.802}{497.954}{Q}%
\TextAt{21}{381.101}{496.16}{p}%
\TextAt{8}{390.446}{497.954}{is}%
\TextAt{8}{402.896}{497.954}{obtained}%
\TextAt{8}{451.715}{497.954}{by}%
\TextAt{8}{468.661}{497.954}{adjoining}%
\TextAt{8}{520.408}{497.954}{all}%
\TextAt{8}{79.226}{480.908}{roots}%
\TextAt{8}{108.887}{480.908}{of}%
\TextAt{8}{122.217}{480.908}{unity}%
\TextAt{8}{152.786}{480.908}{of}%
\TextAt{8}{166.116}{480.908}{order}%
\TextAt{8}{196.686}{480.908}{prime}%
\TextAt{8}{229.849}{480.908}{to}%
\TextAt{13}{244.16}{480.908}{p}%
\TextAt{8}{250.035}{480.908}{.}%
\TextAt{12}{79.226}{448.918}{Exercise}%
\TextAt{12}{133.554}{448.918}{II.7.2.}%
\TextAt{8}{173.507}{448.918}{Let}%
\TextAt{13}{195.457}{448.918}{K}%
\TextAt{8}{211.151}{448.918}{be}%
\TextAt{8}{228.055}{448.918}{henselian}%
\TextAt{8}{279.821}{448.918}{and}%
\TextAt{13}{303.563}{448.918}{K}%
\TextAt{21}{313.522}{447.125}{nr}%
\TextAt{16}{328.206}{448.918}{\Character{106}}%
\TextAt{13}{336.517}{448.918}{K}%
\TextAt{8}{352.212}{448.918}{the}%
\TextAt{8}{373.348}{448.918}{maximal}%
\TextAt{8}{422.137}{448.918}{unramified}%
\TextAt{8}{481.649}{448.918}{extension.}%
\TextAt{8}{79.226}{431.872}{Show}%
\TextAt{8}{111.787}{431.872}{that}%
\TextAt{8}{138.825}{431.872}{the}%
\TextAt{8}{160.65}{431.872}{subextensions}%
\TextAt{8}{236.645}{431.872}{of}%
\TextAt{13}{251.644}{431.872}{K}%
\TextAt{21}{261.603}{430.079}{nr}%
\TextAt{16}{277.461}{431.872}{\Character{106}}%
\TextAt{13}{286.947}{431.872}{K}%
\TextAt{8}{303.331}{431.872}{correspond}%
\TextAt{8}{364.572}{431.872}{1–1}%
\TextAt{8}{387.713}{431.872}{to}%
\TextAt{8}{403.692}{431.872}{the}%
\TextAt{8}{425.517}{431.872}{subextensions}%
\TextAt{8}{501.511}{431.872}{of}%
\TextAt{8}{516.51}{431.872}{the}%
\TextAt{8}{79.226}{414.826}{separable}%
\TextAt{8}{130.669}{414.826}{closure}%
\TextAt{9}{170.193}{414.826}{\Character{22}}%
\TextAt{13}{169.751}{414.826}{\Character{20}}%
\TextAt{21}{176.489}{413.033}{s}%
\TextAt{16}{184.224}{414.826}{\Character{106}}%
\TextAt{13}{190.866}{414.826}{\Character{20}}%
\TextAt{8}{197.604}{414.826}{.}%
\TextAt{12}{79.226}{382.836}{Exercise}%
\TextAt{12}{133.198}{382.836}{II.7.3.}%
\TextAt{8}{172.842}{382.836}{Let}%
\TextAt{13}{194.482}{382.836}{L}%
\TextAt{16}{206.91}{382.836}{\Character{106}}%
\TextAt{13}{214.694}{382.836}{K}%
\TextAt{8}{230.08}{382.836}{be}%
\TextAt{8}{246.675}{382.836}{totally}%
\TextAt{8}{284.753}{382.836}{and}%
\TextAt{8}{308.186}{382.836}{tamely}%
\TextAt{8}{347.555}{382.836}{ramified,}%
\TextAt{8}{397.17}{382.836}{and}%
\TextAt{8}{420.603}{382.836}{let}%
\TextAt{9}{438.179}{382.836}{\Character{1}}%
\TextAt{8}{447.934}{382.836}{,}%
\TextAt{8}{455.921}{382.836}{resp.}%
\TextAt{9}{484.614}{382.836}{\Character{0}}%
\TextAt{8}{491.928}{382.836}{,}%
\TextAt{8}{499.916}{382.836}{be}%
\TextAt{8}{516.51}{382.836}{the}%
\TextAt{8}{79.226}{365.79}{value}%
\TextAt{8}{109.599}{365.79}{group}%
\TextAt{8}{142.901}{365.79}{of}%
\TextAt{13}{156.357}{365.79}{L}%
\TextAt{8}{164.321}{365.79}{,}%
\TextAt{8}{171.628}{365.79}{resp.}%
\TextAt{13}{199.777}{365.79}{K}%
\TextAt{8}{210.595}{365.79}{.}%
\TextAt{8}{219.421}{365.79}{Show}%
\TextAt{8}{250.44}{365.79}{that}%
\TextAt{8}{275.935}{365.79}{the}%
\TextAt{8}{296.217}{365.79}{intermediate}%
\TextAt{8}{364.308}{365.79}{fields}%
\TextAt{8}{394.406}{365.79}{of}%
\TextAt{13}{407.862}{365.79}{L}%
\TextAt{16}{419.362}{365.79}{\Character{106}}%
\TextAt{13}{426.219}{365.79}{K}%
\TextAt{8}{441.06}{365.79}{correspond}%
\TextAt{8}{500.759}{365.79}{1–1}%
\TextAt{8}{522.357}{365.79}{to}%
\TextAt{8}{79.226}{348.744}{the}%
\TextAt{8}{99.383}{348.744}{subgroups}%
\TextAt{8}{154.795}{348.744}{of}%
\TextAt{9}{168.125}{348.744}{\Character{1}\Character{47}\Character{0}}%
\TextAt{8}{191.047}{348.744}{.}%
\TextAt{15}{79.226}{295.88}{§}%
\TextAt{15}{93.785}{295.88}{8}%
\TextAt{15}{120.693}{295.88}{Extensions}%
\TextAt{15}{218.339}{295.88}{of}%
\TextAt{15}{240.409}{295.88}{Valuations}%
\TextAt{8}{79.226}{263.994}{Having}%
\TextAt{8}{118.816}{263.994}{seen}%
\TextAt{8}{143.82}{263.994}{that}%
\TextAt{8}{168.776}{263.994}{the}%
\TextAt{8}{188.52}{263.994}{henselian}%
\TextAt{8}{238.893}{263.994}{valuations}%
\TextAt{8}{294.168}{263.994}{extend}%
\TextAt{8}{331.797}{263.994}{uniquely}%
\TextAt{8}{378.859}{263.994}{to}%
\TextAt{8}{392.756}{263.994}{algebraic}%
\TextAt{8}{441.779}{263.994}{extensions}%
\TextAt{8}{497.747}{263.994}{we}%
\TextAt{8}{514.562}{263.994}{will}%
\TextAt{8}{79.226}{246.948}{now}%
\TextAt{8}{102.906}{246.948}{study}%
\TextAt{8}{134.453}{246.948}{the}%
\TextAt{8}{153.901}{246.948}{question}%
\TextAt{8}{199.759}{246.948}{of}%
\TextAt{8}{212.38}{246.948}{how}%
\TextAt{8}{236.061}{246.948}{a}%
\TextAt{8}{245.108}{246.948}{valuation}%
\TextAt{13}{295.472}{246.948}{v}%
\TextAt{8}{304.749}{246.948}{of}%
\TextAt{8}{317.371}{246.948}{a}%
\TextAt{8}{326.418}{246.948}{field}%
\TextAt{13}{351.067}{246.948}{K}%
\TextAt{8}{365.074}{246.948}{extends}%
\TextAt{8}{407.022}{246.948}{to}%
\TextAt{8}{420.624}{246.948}{an}%
\TextAt{8}{436.175}{246.948}{algebraic}%
\TextAt{8}{484.901}{246.948}{extension}%
\TextAt{8}{79.226}{229.902}{in}%
\TextAt{8}{92.386}{229.902}{general.}%
\TextAt{8}{137.097}{229.902}{So}%
\TextAt{8}{152.863}{229.902}{let}%
\TextAt{13}{169.274}{229.902}{v}%
\TextAt{8}{178.767}{229.902}{be}%
\TextAt{8}{194.198}{229.902}{an}%
\TextAt{8}{209.963}{229.902}{arbitrary}%
\TextAt{8}{259.24}{229.902}{archimedean}%
\TextAt{8}{326.712}{229.902}{or}%
\TextAt{8}{340.529}{229.902}{nonarchimedean}%
\TextAt{8}{426.866}{229.902}{valuation.}%
\TextAt{8}{482.325}{229.902}{There}%
\TextAt{8}{515.641}{229.902}{is}%
\TextAt{8}{526.911}{229.902}{a}%
\TextAt{8}{79.226}{212.857}{little}%
\TextAt{8}{107.631}{212.857}{discrepancy}%
\TextAt{8}{171.543}{212.857}{in}%
\TextAt{8}{185.638}{212.857}{notation}%
\TextAt{8}{232.921}{212.857}{here,}%
\TextAt{8}{262.082}{212.857}{because}%
\TextAt{8}{305.825}{212.857}{archimedean}%
\TextAt{8}{374.233}{212.857}{valuations}%
\TextAt{8}{430.362}{212.857}{manifest}%
\TextAt{8}{478.015}{212.857}{themselves}%
\TextAt{8}{79.226}{195.811}{only}%
\TextAt{8}{106.421}{195.811}{as}%
\TextAt{8}{122.295}{195.811}{absolute}%
\TextAt{8}{170.041}{195.811}{values}%
\TextAt{8}{206.405}{195.811}{while}%
\TextAt{8}{238.466}{195.811}{the}%
\TextAt{8}{260.126}{195.811}{letter}%
\TextAt{13}{292.844}{195.811}{v}%
\TextAt{8}{304.333}{195.811}{has}%
\TextAt{8}{326.71}{195.811}{hitherto}%
\TextAt{8}{373.093}{195.811}{been}%
\TextAt{8}{402.225}{195.811}{used}%
\TextAt{8}{430.448}{195.811}{for}%
\TextAt{8}{449.836}{195.811}{nonarchimedean}%
\TextAt{8}{79.226}{178.765}{exponential}%
\TextAt{8}{140.838}{178.765}{valuations.}%
\TextAt{8}{200.692}{178.765}{In}%
\TextAt{8}{214.16}{178.765}{spite}%
\TextAt{8}{241.03}{178.765}{of}%
\TextAt{8}{253.208}{178.765}{this,}%
\TextAt{8}{278.36}{178.765}{it}%
\TextAt{8}{288.911}{178.765}{will}%
\TextAt{8}{309.863}{178.765}{prove}%
\TextAt{8}{340.26}{178.765}{advantageous,}%
\TextAt{8}{414.871}{178.765}{and}%
\TextAt{8}{436.481}{178.765}{agrees}%
\TextAt{8}{470.512}{178.765}{with}%
\TextAt{8}{496.019}{178.765}{current}%
\TextAt{8}{79.226}{161.719}{usage,}%
\TextAt{8}{114.915}{161.719}{to}%
\TextAt{8}{129.629}{161.719}{employ}%
\TextAt{8}{170.358}{161.719}{the}%
\TextAt{8}{190.918}{161.719}{letter}%
\TextAt{13}{222.537}{161.719}{v}%
\TextAt{8}{232.926}{161.719}{simultaneously}%
\TextAt{8}{312.808}{161.719}{for}%
\TextAt{8}{331.097}{161.719}{both}%
\TextAt{8}{359.141}{161.719}{types}%
\TextAt{8}{390.496}{161.719}{of}%
\TextAt{8}{404.23}{161.719}{valuations,}%
\TextAt{8}{463.674}{161.719}{to}%
\TextAt{8}{478.388}{161.719}{denote}%
\TextAt{8}{516.51}{161.719}{the}%
\TextAt{8}{79.226}{144.673}{corresponding}%
\TextAt{8}{155.249}{144.673}{multiplicative}%
\TextAt{8}{229.91}{144.673}{valuation}%
\TextAt{8}{281.82}{144.673}{in}%
\TextAt{8}{296.31}{144.673}{both}%
\TextAt{8}{324.787}{144.673}{cases}%
\TextAt{8}{355.01}{144.673}{by}%
\TextAt{16}{372.107}{144.673}{\Character{106}}%
\TextAt{16}{379.304}{144.673}{\Character{106}}%
\TextAt{21}{382.646}{141.186}{v}%
\TextAt{8}{392.288}{144.673}{and}%
\TextAt{8}{415.888}{144.673}{the}%
\TextAt{8}{436.882}{144.673}{completion}%
\TextAt{8}{497.555}{144.673}{by}%
\TextAt{13}{514.651}{144.673}{K}%
\TextAt{21}{524.61}{142.879}{v}%
\TextAt{8}{529.518}{144.673}{.}%
\TextAt{8}{79.226}{127.627}{Where}%
\TextAt{8}{116.61}{127.627}{confusion}%
\TextAt{8}{168.376}{127.627}{lurks,}%
\TextAt{8}{200.631}{127.627}{we}%
\TextAt{8}{217.858}{127.627}{will}%
\TextAt{8}{239.963}{127.627}{supply}%
\TextAt{8}{277.419}{127.627}{clarifying}%
\TextAt{8}{328.802}{127.627}{remarks.}%
\TextAt{8}{297.211}{43.033}{146}%
\Destination{78.23}{809.815}{nav.page.146}%
\Destination{79.23}{321.745}{nav.section.2.8}%
\Destination{79.226}{511.955}{nav.ex.2.7.1}%
\Destination{79.226}{462.955}{nav.ex.2.7.2}%
\Destination{79.226}{396.855}{nav.ex.2.7.3}%
\Destination{79.36}{278.955}{nav.print.160}%
\Destination{79.23}{321.745}{outline.28}%
\Hotlink{263.056}{642.855}{25.078}{11.047}{nav.number.2.6.8}%
\Hotlink{78.726}{579.755}{30.936}{11.047}{nav.number.2.7.12}%
\Hotlink{135.544}{579.755}{30.936}{11.047}{nav.number.2.7.13}%
\Destination{79.226}{552.921}{errata.80}%
}

% Page 159 (body page 147)
\Page{%
\TextAt{2}{79.226}{796.943}{II.}%
\TextAt{2}{94.073}{796.943}{The}%
\TextAt{2}{116.492}{796.943}{Theory}%
\TextAt{2}{154.401}{796.943}{of}%
\TextAt{2}{166.826}{796.943}{Valuations}%
\Vector{q 0 G 0 g 0.299 w 7.226 18.729 m 460.769 18.729 l S Q}%
\TextAt{8}{97.159}{761.124}{For}%
\TextAt{8}{117.449}{761.124}{every}%
\TextAt{8}{147.662}{761.124}{valuation}%
\TextAt{13}{198.056}{761.124}{v}%
\TextAt{8}{207.362}{761.124}{of}%
\TextAt{13}{220.013}{761.124}{K}%
\TextAt{8}{234.049}{761.124}{we}%
\TextAt{8}{250.598}{761.124}{consider}%
\TextAt{8}{295.504}{761.124}{the}%
\TextAt{8}{314.981}{761.124}{completion}%
\TextAt{13}{374.138}{761.124}{K}%
\TextAt{21}{384.097}{759.33}{v}%
\TextAt{8}{392.223}{761.124}{and}%
\TextAt{8}{414.306}{761.124}{an}%
\TextAt{8}{429.886}{761.124}{algebraic}%
\TextAt{8}{478.642}{761.124}{closure}%
\Vector{q 0 G 0 g 0.478 w 445.044 -3.033 m 455.862 -3.033 l S Q}%
\TextAt{13}{517.044}{761.124}{K}%
\TextAt{21}{527.862}{759.33}{v}%
\TextAt{8}{79.226}{744.078}{of}%
\TextAt{13}{92.006}{744.078}{K}%
\TextAt{21}{101.965}{742.284}{v}%
\TextAt{8}{106.873}{744.078}{.}%
\TextAt{8}{115.138}{744.078}{The}%
\TextAt{8}{138.642}{744.078}{canonical}%
\TextAt{8}{189.475}{744.078}{extension}%
\TextAt{8}{240.691}{744.078}{of}%
\TextAt{13}{253.472}{744.078}{v}%
\TextAt{8}{262.907}{744.078}{to}%
\TextAt{13}{276.668}{744.078}{K}%
\TextAt{21}{286.627}{742.284}{v}%
\TextAt{8}{294.882}{744.078}{is}%
\TextAt{8}{306.096}{744.078}{again}%
\TextAt{8}{336.773}{744.078}{denoted}%
\TextAt{8}{380.445}{744.078}{by}%
\TextAt{13}{396.154}{744.078}{v}%
\TextAt{8}{405.59}{744.078}{and}%
\TextAt{8}{427.803}{744.078}{the}%
\TextAt{8}{447.41}{744.078}{unique}%
\TextAt{8}{484.901}{744.078}{extension}%
\TextAt{8}{79.226}{727.032}{of}%
\TextAt{8}{92.556}{727.032}{this}%
\TextAt{8}{115.379}{727.032}{latter}%
\TextAt{8}{147.251}{727.032}{valuation}%
\TextAt{8}{198.324}{727.032}{to}%
\Vector{q 0 G 0 g 0.478 w 140.634 -37.125 m 151.452 -37.125 l S Q}%
\TextAt{13}{212.634}{727.032}{K}%
\TextAt{21}{223.452}{725.238}{v}%
\TextAt{8}{232.257}{727.032}{by}%
\TextAt{9}{248.958}{727.032}{\Character{22}}%
\TextAt{13}{248.516}{727.032}{v}%
\TextAt{8}{254.604}{727.032}{.}%
\TextAt{8}{97.159}{706.997}{Let}%
\TextAt{13}{118.128}{706.997}{L}%
\TextAt{16}{129.414}{706.997}{\Character{106}}%
\TextAt{13}{136.055}{706.997}{K}%
\TextAt{8}{150.77}{706.997}{be}%
\TextAt{8}{166.695}{706.997}{an}%
\TextAt{8}{182.954}{706.997}{algebraic}%
\TextAt{8}{232.388}{706.997}{extension.}%
\TextAt{8}{288.705}{706.997}{Choosing}%
\TextAt{8}{339.825}{706.997}{a}%
\TextAt{13}{349.581}{706.997}{K}%
\TextAt{8}{360.399}{706.997}{-embedding}%
\TextAt{13}{271.42}{680.705}{\Character{28}}%
\TextAt{9}{281.159}{680.706}{\Character{58}}%
\TextAt{13}{287.732}{680.705}{L}%
\NativeArrowAt{299.017}{680.706}{11.955}{19.26}{\longrightarrow}%
\Vector{q 0 G 0 g 0.478 w 249.599 -83.451 m 260.417 -83.451 l S Q}%
\TextAt{13}{321.599}{680.705}{K}%
\TextAt{21}{332.417}{678.912}{v}%
\TextAt{13}{337.324}{680.706}{\Character{59}}%
\TextAt{8}{79.226}{654.414}{we}%
\TextAt{8}{96.454}{654.414}{obtain}%
\TextAt{8}{132.881}{654.414}{by}%
\TextAt{8}{149.14}{654.414}{restriction}%
\TextAt{8}{205.138}{654.414}{of}%
\TextAt{9}{218.911}{654.414}{\Character{22}}%
\TextAt{13}{218.468}{654.414}{v}%
\TextAt{8}{228.454}{654.414}{to}%
\TextAt{13}{242.764}{654.414}{\Character{28}}%
\TextAt{13}{249.172}{654.414}{L}%
\TextAt{8}{261.045}{654.414}{an}%
\TextAt{8}{277.304}{654.414}{extension}%
\TextAt{13}{281.881}{628.123}{w}%
\TextAt{9}{293.891}{628.123}{\Character{61}}%
\TextAt{9}{306.743}{628.123}{\Character{22}}%
\TextAt{13}{306.317}{628.123}{v}%
\TextAt{16}{315.062}{628.123}{\Character{14}}%
\TextAt{13}{323.696}{628.123}{\Character{28}}%
\TextAt{8}{79.226}{601.831}{of}%
\TextAt{8}{92.26}{601.831}{the}%
\TextAt{8}{112.121}{601.831}{valuation}%
\TextAt{13}{162.898}{601.831}{v}%
\TextAt{8}{172.588}{601.831}{to}%
\TextAt{13}{186.602}{601.831}{L}%
\TextAt{8}{194.567}{601.831}{.}%
\TextAt{8}{202.917}{601.831}{In}%
\TextAt{8}{217.242}{601.831}{other}%
\TextAt{8}{247.516}{601.831}{words,}%
\TextAt{8}{284.089}{601.831}{if}%
\TextAt{13}{294.517}{601.831}{v}%
\TextAt{8}{300.605}{601.831}{,}%
\TextAt{8}{307.518}{601.831}{resp.}%
\TextAt{9}{335.688}{601.831}{\Character{22}}%
\TextAt{13}{335.245}{601.831}{v}%
\TextAt{8}{341.333}{601.831}{,}%
\TextAt{8}{348.246}{601.831}{are}%
\TextAt{8}{367.461}{601.831}{given}%
\TextAt{8}{397.735}{601.831}{by}%
\TextAt{8}{413.698}{601.831}{the}%
\TextAt{8}{433.559}{601.831}{absolute}%
\TextAt{8}{479.505}{601.831}{values}%
\TextAt{16}{514.071}{601.831}{\Character{106}}%
\TextAt{16}{521.268}{601.831}{\Character{106}}%
\TextAt{21}{524.61}{598.344}{v}%
\TextAt{8}{529.518}{601.831}{,}%
\TextAt{8}{79.226}{584.785}{resp.}%
\TextAt{16}{107.709}{584.785}{\Character{106}}%
\TextAt{16}{114.906}{584.785}{\Character{106}}%
\TextAt{17}{118.571}{581.298}{\Character{22}}%
\TextAt{21}{118.248}{581.298}{v}%
\TextAt{8}{123.155}{584.785}{,}%
\TextAt{8}{130.879}{584.785}{on}%
\TextAt{13}{147.597}{584.785}{K}%
\TextAt{13}{157.759}{584.785}{\Character{59}}%
\TextAt{13}{162.995}{584.785}{K}%
\TextAt{21}{172.968}{582.992}{v}%
\TextAt{8}{177.875}{584.785}{,}%
\TextAt{8}{185.599}{584.785}{resp.}%
\Vector{q 0 G 0 g 0.478 w 142.081 -179.372 m 152.899 -179.372 l S Q}%
\TextAt{13}{214.081}{584.785}{K}%
\TextAt{21}{224.899}{582.992}{v}%
\TextAt{8}{229.807}{584.785}{,}%
\TextAt{8}{237.53}{584.785}{where}%
\TextAt{16}{271.799}{584.785}{\Character{106}}%
\TextAt{16}{278.996}{584.785}{\Character{106}}%
\TextAt{17}{282.661}{581.298}{\Character{22}}%
\TextAt{21}{282.338}{581.298}{v}%
\TextAt{8}{291.602}{584.785}{extends}%
\TextAt{8}{334.718}{584.785}{precisely}%
\TextAt{8}{383.034}{584.785}{the}%
\TextAt{8}{403.65}{584.785}{absolute}%
\TextAt{8}{450.352}{584.785}{value}%
\TextAt{16}{481.058}{584.785}{\Character{106}}%
\TextAt{16}{488.255}{584.785}{\Character{106}}%
\TextAt{21}{491.597}{581.298}{v}%
\TextAt{8}{500.861}{584.785}{of}%
\TextAt{13}{514.651}{584.785}{K}%
\TextAt{21}{524.61}{582.992}{v}%
\TextAt{8}{529.518}{584.785}{,}%
\TextAt{8}{79.226}{567.739}{then}%
\TextAt{8}{105.886}{567.739}{we}%
\TextAt{8}{123.114}{567.739}{obtain}%
\TextAt{8}{159.541}{567.739}{on}%
\TextAt{13}{175.8}{567.739}{L}%
\TextAt{8}{187.662}{567.739}{the}%
\TextAt{8}{207.818}{567.739}{multiplicative}%
\TextAt{8}{281.642}{567.739}{valuation}%
\TextAt{16}{273.151}{541.448}{\Character{106}}%
\TextAt{13}{276.472}{541.448}{x}%
\TextAt{16}{283.124}{541.448}{\Character{106}}%
\TextAt{21}{286.445}{537.961}{w}%
\TextAt{9}{296.555}{541.448}{\Character{61}}%
\TextAt{16}{308.98}{541.448}{\Character{106}}%
\TextAt{13}{312.301}{541.448}{\Character{28}}%
\TextAt{13}{318.709}{541.448}{x}%
\TextAt{16}{325.372}{541.448}{\Character{106}}%
\TextAt{17}{329.016}{537.961}{\Character{22}}%
\TextAt{21}{328.693}{537.961}{v}%
\TextAt{13}{335.593}{541.448}{\Character{58}}%
\TextAt{8}{79.226}{515.156}{The}%
\TextAt{8}{104.717}{515.156}{mapping}%
\TextAt{13}{154.285}{515.156}{\Character{28}}%
\TextAt{9}{166.474}{515.156}{\Character{58}}%
\TextAt{13}{175.495}{515.156}{L}%
\NativeMathAt{189.229}{515.156}{11.955}{\rightarrow}%
\Vector{q 0 G 0 g 0.478 w 134.954 -249.001 m 145.771 -249.001 l S Q}%
\TextAt{13}{206.954}{515.156}{K}%
\TextAt{21}{217.771}{513.363}{v}%
\TextAt{8}{228.013}{515.156}{is}%
\TextAt{8}{241.214}{515.156}{obviously}%
\TextAt{8}{294.429}{515.156}{continuous}%
\TextAt{8}{354.47}{515.156}{with}%
\TextAt{8}{382.567}{515.156}{respect}%
\TextAt{8}{424.054}{515.156}{to}%
\TextAt{8}{439.801}{515.156}{this}%
\TextAt{8}{464.06}{515.156}{valuation.}%
\TextAt{8}{523.994}{515.156}{It}%
\TextAt{8}{79.226}{498.11}{extends}%
\TextAt{8}{121.882}{498.11}{in}%
\TextAt{8}{135.535}{498.11}{a}%
\TextAt{8}{145.29}{498.11}{unique}%
\TextAt{8}{183.332}{498.11}{way}%
\TextAt{8}{207.075}{498.11}{to}%
\TextAt{8}{221.385}{498.11}{a}%
\TextAt{8}{231.141}{498.11}{continuous}%
\TextAt{13}{289.745}{498.11}{K}%
\TextAt{8}{300.562}{498.11}{-embedding}%
\TextAt{13}{268.025}{471.819}{\Character{28}}%
\TextAt{9}{277.765}{471.819}{\Character{58}}%
\TextAt{13}{284.337}{471.819}{L}%
\TextAt{21}{292.302}{470.025}{w}%
\NativeArrowAt{302.412}{471.819}{11.955}{19.26}{\longrightarrow}%
\Vector{q 0 G 0 g 0.478 w 252.994 -292.338 m 263.812 -292.338 l S Q}%
\TextAt{13}{324.994}{471.819}{K}%
\TextAt{21}{335.812}{470.025}{v}%
\TextAt{13}{340.719}{471.819}{\Character{59}}%
\TextAt{8}{79.226}{445.527}{where,}%
\TextAt{8}{117.244}{445.527}{in}%
\TextAt{8}{131.662}{445.527}{the}%
\TextAt{8}{152.584}{445.527}{case}%
\TextAt{8}{178.12}{445.527}{of}%
\TextAt{8}{192.216}{445.527}{an}%
\TextAt{8}{209.24}{445.527}{infinite}%
\TextAt{8}{249.673}{445.527}{extension}%
\TextAt{13}{302.204}{445.527}{L}%
\TextAt{16}{314.794}{445.527}{\Character{106}}%
\TextAt{13}{322.74}{445.527}{K}%
\TextAt{8}{333.557}{445.527}{,}%
\TextAt{13}{341.663}{445.527}{L}%
\TextAt{21}{349.628}{443.734}{w}%
\TextAt{8}{361.08}{445.527}{does}%
\TextAt{19}{388.242}{445.527}{not}%
\TextAt{8}{409.834}{445.527}{mean}%
\TextAt{8}{441.814}{445.527}{the}%
\TextAt{8}{462.736}{445.527}{completion}%
\TextAt{8}{523.337}{445.527}{of}%
\TextAt{13}{79.226}{428.481}{L}%
\TextAt{8}{92.602}{428.481}{with}%
\TextAt{8}{120.775}{428.481}{respect}%
\TextAt{8}{162.339}{428.481}{to}%
\TextAt{13}{178.163}{428.481}{w}%
\TextAt{8}{186.852}{428.481}{,}%
\TextAt{8}{195.894}{428.481}{but}%
\TextAt{8}{218.867}{428.481}{the}%
\TextAt{8}{240.537}{428.481}{union}%
\TextAt{13}{274.569}{428.481}{L}%
\TextAt{21}{282.533}{426.688}{w}%
\TextAt{9}{295.223}{428.481}{\Character{61}}%
\TextAt{20}{310.228}{437.448}{\Character{83}}%
\TextAt{21}{320.191}{424.994}{i}%
\TextAt{13}{325.565}{428.481}{L}%
\TextAt{21}{333.529}{426.688}{iw}%
\TextAt{8}{348.613}{428.481}{of}%
\TextAt{8}{363.456}{428.481}{the}%
\TextAt{8}{385.126}{428.481}{completions}%
\TextAt{13}{451.09}{428.481}{L}%
\TextAt{21}{459.055}{426.688}{iw}%
\TextAt{8}{474.139}{428.481}{of}%
\TextAt{8}{488.982}{428.481}{all}%
\TextAt{8}{506.755}{428.481}{finite}%
\TextAt{8}{79.226}{411.435}{subextensions}%
\TextAt{13}{153.775}{411.435}{L}%
\TextAt{21}{161.74}{409.642}{i}%
\TextAt{16}{168.823}{411.435}{\Character{106}}%
\TextAt{13}{175.846}{411.435}{K}%
\TextAt{8}{190.785}{411.435}{of}%
\TextAt{13}{204.339}{411.435}{L}%
\TextAt{16}{216.006}{411.435}{\Character{106}}%
\TextAt{13}{223.029}{411.435}{K}%
\TextAt{8}{233.847}{411.435}{.}%
\TextAt{8}{242.967}{411.435}{This}%
\TextAt{8}{269.91}{411.435}{union}%
\TextAt{8}{302.652}{411.435}{will}%
\TextAt{8}{324.981}{411.435}{be}%
\TextAt{8}{341.129}{411.435}{henceforth}%
\TextAt{8}{398.905}{411.435}{called}%
\TextAt{8}{432.293}{411.435}{the}%
\TextAt{12}{452.673}{411.435}{localization}%
\TextAt{8}{523.337}{411.435}{of}%
\TextAt{13}{79.226}{394.389}{L}%
\TextAt{8}{91.088}{394.389}{with}%
\TextAt{8}{117.748}{394.389}{respect}%
\TextAt{8}{157.798}{394.389}{to}%
\TextAt{13}{172.108}{394.389}{w}%
\TextAt{8}{180.798}{394.389}{.}%
\TextAt{8}{189.246}{394.389}{When}%
\TextAt{9}{223.378}{394.389}{\Character{91}}%
\TextAt{13}{226.63}{394.389}{L}%
\TextAt{9}{237.915}{394.389}{\Character{58}}%
\TextAt{13}{244.488}{394.389}{K}%
\TextAt{9}{255.305}{394.389}{\Character{93}}%
\TextAt{13}{261.878}{394.389}{\Character{60}}%
\NativeMathAt{274.303}{394.389}{11.955}{\infty}%
\TextAt{8}{286.258}{394.389}{,}%
\TextAt{13}{293.408}{394.389}{\Character{28}}%
\TextAt{8}{303.724}{394.389}{is}%
\TextAt{8}{315.488}{394.389}{given}%
\TextAt{8}{346.057}{394.389}{by}%
\TextAt{8}{362.316}{394.389}{the}%
\TextAt{8}{382.473}{394.389}{rule}%
\TextAt{13}{199.388}{368.098}{x}%
\TextAt{9}{209.361}{368.098}{\Character{61}}%
\TextAt{13}{221.786}{368.098}{w}%
\TextAt{8}{230.476}{368.098}{-}%
\TextAt{8}{239.273}{368.098}{lim}%
\TextAt{21}{236.365}{360.925}{n}%
\NativeMathAt{241.504}{360.925}{7.97}{\rightarrow}%
\NativeMathAt{249.968}{360.925}{7.97}{\infty}%
\TextAt{13}{260.433}{368.098}{x}%
\TextAt{21}{267.085}{366.305}{n}%
\NativeArrowAt{287.997}{368.098}{11.955}{19.26}{\longmapsto}%
\TextAt{13}{322.534}{368.098}{\Character{28}}%
\TextAt{13}{328.942}{368.098}{x}%
\TextAt{9}{338.926}{368.098}{\Character{58}\Character{61}}%
\TextAt{9}{355.018}{368.098}{\Character{22}}%
\TextAt{13}{354.603}{368.098}{v}%
\TextAt{8}{360.692}{368.098}{-}%
\TextAt{8}{369.49}{368.098}{lim}%
\TextAt{21}{366.582}{360.925}{n}%
\NativeMathAt{371.72}{360.925}{7.97}{\rightarrow}%
\NativeMathAt{380.184}{360.925}{7.97}{\infty}%
\TextAt{13}{390.649}{368.098}{\Character{28}}%
\TextAt{13}{397.057}{368.098}{x}%
\TextAt{21}{403.72}{366.305}{n}%
\TextAt{13}{409.356}{368.098}{\Character{59}}%
\TextAt{8}{79.226}{340.521}{where}%
\TextAt{16}{112.631}{340.521}{\Character{102}}%
\TextAt{13}{118.608}{340.521}{x}%
\TextAt{21}{125.26}{338.728}{n}%
\TextAt{16}{130.897}{340.521}{\Character{103}}%
\TextAt{21}{136.874}{338.728}{n}%
\NativeMathAt{142.012}{338.728}{7.97}{\in}%
\TextAt{25}{147.658}{338.728}{N}%
\TextAt{8}{157.405}{340.521}{is}%
\TextAt{8}{168.764}{340.521}{a}%
\TextAt{13}{178.115}{340.521}{w}%
\TextAt{8}{186.804}{340.521}{-Cauchy}%
\TextAt{8}{232.247}{340.521}{sequence}%
\TextAt{8}{280.345}{340.521}{in}%
\TextAt{13}{293.593}{340.521}{L}%
\TextAt{8}{301.557}{340.521}{,}%
\TextAt{8}{308.382}{340.521}{and}%
\TextAt{8}{330.74}{340.521}{hence}%
\TextAt{16}{362.841}{340.521}{\Character{102}}%
\TextAt{13}{368.819}{340.521}{\Character{28}}%
\TextAt{13}{375.227}{340.521}{x}%
\TextAt{21}{381.89}{338.728}{n}%
\TextAt{16}{387.526}{340.521}{\Character{103}}%
\TextAt{21}{393.504}{338.728}{n}%
\NativeMathAt{398.642}{338.728}{7.97}{\in}%
\TextAt{25}{404.288}{338.728}{N}%
\TextAt{8}{414.035}{340.521}{a}%
\TextAt{9}{423.828}{340.521}{\Character{22}}%
\TextAt{13}{423.385}{340.521}{v}%
\TextAt{8}{429.474}{340.521}{-Cauchy}%
\TextAt{8}{474.917}{340.521}{sequence}%
\TextAt{8}{523.014}{340.521}{in}%
\Vector{q 0 G 0 g 0.478 w 7.226 -440.682 m 18.044 -440.682 l S Q}%
\TextAt{13}{79.226}{323.475}{K}%
\TextAt{21}{90.044}{321.682}{v}%
\TextAt{8}{94.951}{323.475}{.}%
\TextAt{8}{104.111}{323.475}{Note}%
\TextAt{8}{132.633}{323.475}{here}%
\TextAt{8}{158.227}{323.475}{that}%
\TextAt{8}{183.833}{323.475}{the}%
\TextAt{8}{204.227}{323.475}{sequence}%
\TextAt{13}{252.966}{323.475}{\Character{28}}%
\TextAt{13}{259.374}{323.475}{x}%
\TextAt{21}{266.037}{321.682}{n}%
\TextAt{8}{275.807}{323.475}{converges}%
\TextAt{8}{328.468}{323.475}{in}%
\TextAt{8}{342.357}{323.475}{the}%
\TextAt{8}{362.751}{323.475}{finite}%
\TextAt{8}{392.9}{323.475}{complete}%
\TextAt{8}{442.559}{323.475}{extension}%
\TextAt{13}{494.562}{323.475}{\Character{28}}%
\TextAt{13}{500.97}{323.475}{L}%
\TextAt{16}{511.764}{323.475}{\Character{1}}%
\TextAt{13}{517.903}{323.475}{K}%
\TextAt{21}{527.862}{321.682}{v}%
\TextAt{8}{79.226}{306.429}{of}%
\TextAt{13}{92.556}{306.429}{K}%
\TextAt{21}{102.515}{304.636}{v}%
\TextAt{8}{107.423}{306.429}{.}%
\TextAt{8}{115.871}{306.429}{We}%
\TextAt{8}{136.016}{306.429}{consider}%
\TextAt{8}{181.601}{306.429}{the}%
\TextAt{8}{201.757}{306.429}{diagram}%
\TextAt{8}{247.294}{306.429}{of}%
\TextAt{8}{260.624}{306.429}{fields}%
\TextAt{13}{267.91}{281.187}{L}%
\TextAt{13}{330.702}{281.187}{L}%
\TextAt{21}{338.666}{279.394}{w}%
\TextAt{13}{266.483}{230.186}{K}%
\TextAt{13}{330.646}{230.186}{K}%
\TextAt{21}{340.605}{228.393}{v}%
\NativeDiagram{%
\path[draw=black,line width=0.478bp] (209.226,-489.823) -- (253.35,-489.823);%
\path[draw=black,line width=0.478bp] (199.892,-531.065) -- (199.892,-497.391);%
\path[draw=black,line width=0.478bp] (210.653,-540.825) -- (253.294,-540.825);%
\path[draw=black,line width=0.478bp] (266.079,-531.065) -- (266.079,-499.185);%
}%
\TextAt{8}{517.802}{258.875}{(*)}%
\TextAt{8}{79.226}{205.355}{The}%
\TextAt{8}{102.683}{205.355}{canonical}%
\TextAt{8}{153.469}{205.355}{extension}%
\TextAt{8}{204.638}{205.355}{of}%
\TextAt{8}{217.371}{205.355}{the}%
\TextAt{8}{236.93}{205.355}{valuation}%
\TextAt{13}{287.406}{205.355}{w}%
\TextAt{8}{299.395}{205.355}{from}%
\TextAt{13}{326.439}{205.355}{L}%
\TextAt{8}{337.704}{205.355}{to}%
\TextAt{13}{351.417}{205.355}{L}%
\TextAt{21}{359.381}{203.562}{w}%
\TextAt{8}{369.471}{205.355}{is}%
\TextAt{8}{380.638}{205.355}{precisely}%
\TextAt{8}{427.897}{205.355}{the}%
\TextAt{8}{447.457}{205.355}{unique}%
\TextAt{8}{484.901}{205.355}{extension}%
\TextAt{8}{79.226}{188.309}{of}%
\TextAt{8}{92.556}{188.309}{the}%
\TextAt{8}{112.713}{188.309}{valuation}%
\TextAt{13}{163.785}{188.309}{v}%
\TextAt{8}{173.771}{188.309}{from}%
\TextAt{13}{201.411}{188.309}{K}%
\TextAt{21}{211.37}{186.516}{v}%
\TextAt{8}{220.175}{188.309}{to}%
\TextAt{8}{234.485}{188.309}{the}%
\TextAt{8}{254.642}{188.309}{extension}%
\TextAt{13}{306.408}{188.309}{L}%
\TextAt{21}{314.372}{186.516}{w}%
\TextAt{16}{324.483}{188.309}{\Character{106}}%
\TextAt{13}{331.124}{188.309}{K}%
\TextAt{21}{341.083}{186.516}{v}%
\TextAt{8}{345.991}{188.309}{.}%
\TextAt{8}{354.439}{188.309}{We}%
\TextAt{8}{374.584}{188.309}{have}%
\TextAt{13}{277.706}{162.017}{L}%
\TextAt{21}{285.671}{160.224}{w}%
\TextAt{9}{295.781}{162.017}{\Character{61}}%
\TextAt{13}{308.206}{162.017}{LK}%
\TextAt{21}{326.13}{160.224}{v}%
\TextAt{13}{331.038}{162.017}{\Character{59}}%
\TextAt{8}{79.226}{135.726}{because}%
\TextAt{8}{122.264}{135.726}{if}%
\TextAt{13}{132.723}{135.726}{L}%
\TextAt{16}{144.009}{135.726}{\Character{106}}%
\TextAt{13}{150.65}{135.726}{K}%
\TextAt{8}{165.101}{135.726}{is}%
\TextAt{8}{176.601}{135.726}{finite,}%
\TextAt{8}{209.553}{135.726}{then}%
\TextAt{8}{235.949}{135.726}{the}%
\TextAt{8}{255.842}{135.726}{field}%
\TextAt{13}{280.935}{135.726}{LK}%
\TextAt{21}{298.858}{133.933}{v}%
\NativeMathAt{307.087}{135.726}{11.955}{\subseteq}%
\TextAt{13}{319.706}{135.726}{L}%
\TextAt{21}{327.67}{133.933}{w}%
\TextAt{8}{338.093}{135.726}{is}%
\TextAt{8}{349.593}{135.726}{complete}%
\TextAt{8}{398.751}{135.726}{by}%
\TextAt{8}{414.746}{135.726}{(4.8),}%
\TextAt{8}{445.762}{135.726}{contains}%
\TextAt{8}{491.418}{135.726}{the}%
\TextAt{8}{511.31}{135.726}{field}%
\TextAt{13}{79.226}{118.68}{L}%
\TextAt{8}{91.508}{118.68}{and}%
\TextAt{8}{114.69}{118.68}{therefore}%
\TextAt{8}{164.21}{118.68}{has}%
\TextAt{8}{185.503}{118.68}{to}%
\TextAt{8}{200.233}{118.68}{be}%
\TextAt{8}{216.577}{118.68}{its}%
\TextAt{8}{233.315}{118.68}{completion.}%
\TextAt{8}{298.961}{118.68}{If}%
\TextAt{13}{311.073}{118.68}{L}%
\TextAt{21}{319.037}{116.887}{w}%
\TextAt{16}{329.863}{118.68}{\Character{106}}%
\TextAt{13}{337.219}{118.68}{K}%
\TextAt{21}{347.179}{116.887}{v}%
\TextAt{8}{356.403}{118.68}{has}%
\TextAt{8}{377.697}{118.68}{degree}%
\TextAt{13}{414.532}{118.68}{n}%
\TextAt{13}{425.543}{118.68}{\Character{60}}%
\NativeMathAt{438.696}{118.68}{11.955}{\infty}%
\TextAt{8}{450.651}{118.68}{,}%
\TextAt{8}{458.325}{118.68}{then,}%
\TextAt{8}{488.761}{118.68}{by}%
\TextAt{8}{505.44}{118.68}{(4.8),}%
\TextAt{8}{79.226}{101.634}{the}%
\TextAt{8}{99.383}{101.634}{absolute}%
\TextAt{8}{145.625}{101.634}{values}%
\TextAt{8}{180.486}{101.634}{corresponding}%
\TextAt{8}{255.672}{101.634}{to}%
\TextAt{13}{269.983}{101.634}{v}%
\TextAt{8}{279.968}{101.634}{and}%
\TextAt{13}{302.731}{101.634}{w}%
\TextAt{8}{315.318}{101.634}{satisfy}%
\TextAt{8}{351.865}{101.634}{the}%
\TextAt{8}{372.021}{101.634}{relation}%
\TextAt{16}{247.737}{73.842}{\Character{106}}%
\TextAt{13}{251.058}{73.842}{x}%
\TextAt{16}{257.71}{73.842}{\Character{106}}%
\TextAt{21}{261.03}{70.355}{w}%
\TextAt{9}{271.141}{73.842}{\Character{61}}%
\TextAt{22}{286.887}{78.763}{n}%
\TextAt{20}{284.846}{87.985}{\Character{113}}%
\Vector{q 0 G 0 g 0.478 w 224.801 -685.776 m 289.007 -685.776 l S Q}%
\TextAt{20}{296.801}{84.004}{\Character{12}}%
\TextAt{20}{296.801}{76.831}{\Character{12}}%
\TextAt{13}{300.786}{73.842}{N}%
\TextAt{21}{310.16}{71.983}{L}%
\TextAt{22}{315.919}{70.986}{w}%
\TextAt{3}{321.918}{71.982}{\Character{106}}%
\TextAt{21}{324.27}{71.982}{K}%
\TextAt{22}{331.411}{70.986}{v}%
\TextAt{9}{336.358}{73.842}{\Character{40}}%
\TextAt{13}{340.91}{73.842}{x}%
\TextAt{9}{347.562}{73.842}{\Character{41}}%
\TextAt{20}{352.115}{84.004}{\Character{12}}%
\TextAt{20}{352.115}{76.831}{\Character{12}}%
\TextAt{21}{356.1}{69.16}{v}%
\TextAt{13}{361.007}{73.842}{\Character{58}}%
\TextAt{8}{297.211}{43.033}{147}%
\Destination{522.36}{275.925}{nav.AMS.20}%
\Destination{78.23}{809.815}{nav.page.147}%
\Destination{79.36}{355.555}{nav.print.161}%
\Hotlink{414.246}{133.555}{25.078}{11.047}{nav.number.2.4.8}%
\Hotlink{504.94}{116.455}{25.078}{11.047}{nav.number.2.4.8}%
\Destination{79.226}{320}{nav.local.2.8.diagram}%
}

% Page 160 (body page 148)
\Page{%
\TextAt{2}{79.226}{796.943}{II.}%
\TextAt{2}{94.073}{796.943}{The}%
\TextAt{2}{116.492}{796.943}{Theory}%
\TextAt{2}{154.401}{796.943}{of}%
\TextAt{2}{166.826}{796.943}{Valuations}%
\Vector{q 0 G 0 g 0.299 w 7.226 18.729 m 460.769 18.729 l S Q}%
\TextAt{8}{79.226}{761.124}{The}%
\TextAt{8}{103.386}{761.124}{field}%
\TextAt{8}{128.848}{761.124}{diagram}%
\TextAt{8}{174.491}{761.124}{(*)}%
\TextAt{8}{193.462}{761.124}{is}%
\TextAt{8}{205.331}{761.124}{of}%
\TextAt{8}{218.767}{761.124}{central}%
\TextAt{8}{257.571}{761.124}{importance}%
\TextAt{8}{319.139}{761.124}{for}%
\TextAt{8}{337.129}{761.124}{algebraic}%
\TextAt{8}{386.669}{761.124}{number}%
\TextAt{8}{429.371}{761.124}{theory.}%
\TextAt{8}{470.009}{761.124}{It}%
\TextAt{8}{482.787}{761.124}{shows}%
\TextAt{8}{516.51}{761.124}{the}%
\TextAt{8}{79.226}{744.078}{passage}%
\TextAt{8}{122.715}{744.078}{from}%
\TextAt{8}{151.439}{744.078}{the}%
\TextAt{8}{172.679}{744.078}{“global}%
\TextAt{8}{213.765}{744.078}{extension”}%
\TextAt{13}{272.138}{744.078}{L}%
\TextAt{16}{285.27}{744.078}{\Character{106}}%
\TextAt{13}{293.758}{744.078}{K}%
\TextAt{8}{309.557}{744.078}{to}%
\TextAt{8}{324.951}{744.078}{the}%
\TextAt{8}{346.191}{744.078}{“local}%
\TextAt{8}{380.438}{744.078}{extension”}%
\TextAt{13}{438.811}{744.078}{L}%
\TextAt{21}{446.776}{742.284}{w}%
\TextAt{16}{458.733}{744.078}{\Character{106}}%
\TextAt{13}{467.221}{744.078}{K}%
\TextAt{21}{477.18}{742.284}{v}%
\TextAt{8}{487.069}{744.078}{and}%
\TextAt{8}{510.915}{744.078}{thus}%
\TextAt{8}{79.226}{727.032}{represents}%
\TextAt{8}{134.643}{727.032}{one}%
\TextAt{8}{156.442}{727.032}{of}%
\TextAt{8}{170.112}{727.032}{the}%
\TextAt{8}{190.608}{727.032}{most}%
\TextAt{8}{219.628}{727.032}{important}%
\TextAt{8}{275.26}{727.032}{methods}%
\TextAt{8}{322.811}{727.032}{of}%
\TextAt{8}{336.481}{727.032}{algebraic}%
\TextAt{8}{386.255}{727.032}{number}%
\TextAt{8}{429.191}{727.032}{theory,}%
\TextAt{8}{468.637}{727.032}{the}%
\TextAt{8}{489.133}{727.032}{so-called}%
\TextAt{12}{79.226}{709.986}{local-to-global}%
\TextAt{12}{166.865}{709.986}{principle}%
\TextAt{8}{218.117}{709.986}{.}%
\TextAt{8}{226.99}{709.986}{This}%
\TextAt{8}{253.851}{709.986}{terminology}%
\TextAt{8}{318.718}{709.986}{arises}%
\TextAt{8}{350.851}{709.986}{from}%
\TextAt{8}{378.633}{709.986}{the}%
\TextAt{8}{398.931}{709.986}{case}%
\TextAt{8}{423.844}{709.986}{of}%
\TextAt{8}{437.315}{709.986}{a}%
\TextAt{8}{447.212}{709.986}{function}%
\TextAt{8}{493.202}{709.986}{field}%
\TextAt{13}{518.7}{709.986}{K}%
\TextAt{8}{529.518}{709.986}{,}%
\TextAt{8}{79.226}{692.94}{for}%
\TextAt{8}{98.368}{692.94}{example}%
\TextAt{13}{145.473}{692.94}{K}%
\TextAt{9}{161.754}{692.94}{\Character{61}}%
\TextAt{18}{176.322}{692.94}{C}%
\TextAt{9}{184.956}{692.94}{\Character{40}}%
\TextAt{13}{189.508}{692.94}{t}%
\TextAt{9}{193.736}{692.94}{\Character{41}}%
\TextAt{8}{198.288}{692.94}{,}%
\TextAt{8}{207.008}{692.94}{where}%
\TextAt{8}{242.075}{692.94}{the}%
\TextAt{8}{263.488}{692.94}{elements}%
\TextAt{8}{312.602}{692.94}{of}%
\TextAt{8}{327.189}{692.94}{the}%
\TextAt{8}{348.603}{692.94}{extension}%
\TextAt{13}{401.625}{692.94}{L}%
\TextAt{8}{414.744}{692.94}{are}%
\TextAt{8}{435.512}{692.94}{algebraic}%
\TextAt{8}{486.204}{692.94}{functions}%
\TextAt{8}{79.226}{675.894}{on}%
\TextAt{8}{95.759}{675.894}{a}%
\TextAt{8}{105.789}{675.894}{Riemann}%
\TextAt{8}{155.653}{675.894}{surface,}%
\TextAt{8}{198.652}{675.894}{hence}%
\TextAt{8}{231.432}{675.894}{on}%
\TextAt{8}{247.965}{675.894}{a}%
\TextAt{8}{257.994}{675.894}{global}%
\TextAt{8}{292.747}{675.894}{object,}%
\TextAt{8}{331.454}{675.894}{whereas}%
\TextAt{8}{376.01}{675.894}{passing}%
\TextAt{8}{417.386}{675.894}{to}%
\TextAt{13}{431.97}{675.894}{K}%
\TextAt{21}{441.93}{674.101}{v}%
\TextAt{8}{451.009}{675.894}{and}%
\TextAt{13}{474.045}{675.894}{L}%
\TextAt{21}{482.01}{674.101}{w}%
\TextAt{8}{492.971}{675.894}{signifies}%
\TextAt{8}{79.226}{658.848}{looking}%
\TextAt{8}{121.207}{658.848}{at}%
\TextAt{8}{136.515}{658.848}{power}%
\TextAt{8}{171.658}{658.848}{series}%
\TextAt{8}{203.991}{658.848}{expansions,}%
\TextAt{8}{267.477}{658.848}{i.e.,}%
\TextAt{8}{290.83}{658.848}{the}%
\TextAt{8}{311.985}{658.848}{local}%
\TextAt{8}{340.624}{658.848}{study}%
\TextAt{8}{373.877}{658.848}{of}%
\TextAt{8}{388.205}{658.848}{functions.}%
\TextAt{8}{446.214}{658.848}{The}%
\TextAt{8}{471.266}{658.848}{diagram}%
\TextAt{8}{517.802}{658.848}{(*)}%
\TextAt{8}{79.226}{641.802}{thus}%
\TextAt{8}{105.726}{641.802}{expresses}%
\TextAt{8}{157.058}{641.802}{in}%
\TextAt{8}{171.459}{641.802}{an}%
\TextAt{8}{188.467}{641.802}{abstract}%
\TextAt{8}{234.813}{641.802}{manner}%
\TextAt{8}{277.836}{641.802}{our}%
\TextAt{8}{299.398}{641.802}{original}%
\TextAt{8}{342.433}{641.802}{goal,}%
\TextAt{8}{371.344}{641.802}{to}%
\TextAt{8}{386.403}{641.802}{provide}%
\TextAt{8}{428.78}{641.802}{methods}%
\TextAt{8}{476.74}{641.802}{of}%
\TextAt{8}{490.819}{641.802}{function}%
\TextAt{8}{79.226}{624.756}{theory}%
\TextAt{8}{115.976}{624.756}{for}%
\TextAt{8}{133.861}{624.756}{use}%
\TextAt{8}{154.077}{624.756}{in}%
\TextAt{8}{167.73}{624.756}{the}%
\TextAt{8}{187.887}{624.756}{theory}%
\TextAt{8}{224.637}{624.756}{of}%
\TextAt{8}{237.967}{624.756}{numbers}%
\TextAt{8}{285.178}{624.756}{by}%
\TextAt{8}{301.437}{624.756}{means}%
\TextAt{8}{337.267}{624.756}{of}%
\TextAt{8}{350.597}{624.756}{valuations.}%
\TextAt{8}{97.159}{604.721}{We}%
\TextAt{8}{116.929}{604.721}{saw}%
\TextAt{8}{139.055}{604.721}{that}%
\TextAt{8}{164.05}{604.721}{every}%
\TextAt{13}{194.568}{604.721}{K}%
\TextAt{8}{205.385}{604.721}{-embedding}%
\TextAt{13}{268.087}{604.721}{\Character{28}}%
\TextAt{9}{277.827}{604.721}{\Character{58}}%
\TextAt{13}{284.399}{604.721}{L}%
\NativeMathAt{295.684}{604.721}{11.955}{\rightarrow}%
\Vector{q 0 G 0 g 0.478 w 238.96 -159.436 m 249.778 -159.436 l S Q}%
\TextAt{13}{310.96}{604.721}{K}%
\TextAt{21}{321.778}{602.928}{v}%
\TextAt{8}{330.209}{604.721}{gave}%
\TextAt{8}{356.184}{604.721}{us}%
\TextAt{8}{370.826}{604.721}{an}%
\TextAt{8}{386.711}{604.721}{extension}%
\TextAt{13}{438.102}{604.721}{w}%
\TextAt{9}{450.113}{604.721}{\Character{61}}%
\TextAt{9}{462.965}{604.721}{\Character{22}}%
\TextAt{13}{462.538}{604.721}{v}%
\TextAt{16}{470.518}{604.721}{\Character{14}}%
\TextAt{13}{478.387}{604.721}{\Character{28}}%
\TextAt{8}{488.33}{604.721}{of}%
\TextAt{13}{501.286}{604.721}{v}%
\TextAt{8}{507.374}{604.721}{.}%
\TextAt{8}{515.698}{604.721}{For}%
\TextAt{8}{79.226}{587.675}{every}%
\TextAt{8}{110.118}{587.675}{automorphism}%
\TextAt{13}{187.588}{587.675}{\Character{27}}%
\NativeMathAt{197.991}{587.675}{11.955}{\in}%
\TextAt{13}{209.282}{587.675}{G}%
\TextAt{9}{218.516}{587.675}{\Character{40}}%
\Vector{q 0 G 0 g 0.478 w 151.068 -176.482 m 161.886 -176.482 l S Q}%
\TextAt{13}{223.068}{587.675}{K}%
\TextAt{21}{233.886}{585.882}{v}%
\TextAt{16}{242.114}{587.675}{\Character{106}}%
\TextAt{13}{248.756}{587.675}{K}%
\TextAt{21}{258.715}{585.882}{v}%
\TextAt{9}{263.622}{587.675}{\Character{41}}%
\TextAt{8}{272.072}{587.675}{of}%
\Vector{q 0 G 0 g 0.478 w 213.402 -176.482 m 224.22 -176.482 l S Q}%
\TextAt{13}{285.402}{587.675}{K}%
\TextAt{21}{296.22}{585.882}{v}%
\TextAt{8}{305.025}{587.675}{over}%
\TextAt{13}{330.071}{587.675}{K}%
\TextAt{21}{340.03}{585.882}{v}%
\TextAt{8}{344.938}{587.675}{,}%
\TextAt{8}{352.087}{587.675}{we}%
\TextAt{8}{369.314}{587.675}{obtain}%
\TextAt{8}{405.742}{587.675}{with}%
\TextAt{8}{432.402}{587.675}{the}%
\TextAt{8}{452.558}{587.675}{composite}%
\TextAt{13}{267.693}{558.674}{L}%
\TextAt{21}{281.911}{565.451}{\Character{28}}%
\NativeArrowAt{278.979}{558.674}{11.955}{11.955}{\longrightarrow}%
\Vector{q 0 G 0 g 0.478 w 222.255 -205.483 m 233.072 -205.483 l S Q}%
\TextAt{13}{294.255}{558.674}{K}%
\TextAt{21}{305.072}{556.881}{v}%
\TextAt{21}{316.01}{565.451}{\Character{27}}%
\NativeArrowAt{313.301}{558.674}{11.955}{11.955}{\longrightarrow}%
\Vector{q 0 G 0 g 0.478 w 256.577 -205.483 m 267.395 -205.483 l S Q}%
\TextAt{13}{328.577}{558.674}{K}%
\TextAt{21}{339.395}{556.881}{v}%
\TextAt{8}{79.226}{529.673}{a}%
\TextAt{8}{89.419}{529.673}{new}%
\TextAt{13}{113.91}{529.673}{K}%
\TextAt{8}{124.728}{529.673}{-embedding}%
\TextAt{13}{188.241}{529.673}{\Character{28}}%
\TextAt{3}{194.66}{534.011}{\Character{48}}%
\TextAt{9}{201.521}{529.673}{\Character{61}}%
\TextAt{13}{214.692}{529.673}{\Character{27}}%
\TextAt{16}{224.729}{529.673}{\Character{14}}%
\TextAt{13}{233.661}{529.673}{\Character{28}}%
\TextAt{8}{244.415}{529.673}{of}%
\TextAt{13}{258.182}{529.673}{L}%
\TextAt{8}{266.147}{529.673}{.}%
\TextAt{8}{275.907}{529.673}{It}%
\TextAt{8}{289.017}{529.673}{will}%
\TextAt{8}{311.559}{529.673}{be}%
\TextAt{8}{327.921}{529.673}{said}%
\TextAt{8}{352.484}{529.673}{to}%
\TextAt{8}{367.232}{529.673}{be}%
\TextAt{19}{383.593}{529.673}{conjugate}%
\TextAt{19}{436.068}{529.673}{to}%
\TextAt{13}{450.585}{529.673}{\Character{28}}%
\TextAt{19}{461.657}{529.673}{over}%
\TextAt{13}{487.986}{529.673}{K}%
\TextAt{21}{497.945}{527.88}{v}%
\TextAt{8}{502.853}{529.673}{.}%
\TextAt{8}{512.613}{529.673}{The}%
\TextAt{8}{79.226}{512.627}{following}%
\TextAt{8}{128.661}{512.627}{result}%
\TextAt{8}{161.239}{512.627}{gives}%
\TextAt{8}{189.919}{512.627}{us}%
\TextAt{8}{204.935}{512.627}{a}%
\TextAt{8}{214.69}{512.627}{complete}%
\TextAt{8}{264.113}{512.627}{description}%
\TextAt{8}{324.008}{512.627}{of}%
\TextAt{8}{337.338}{512.627}{the}%
\TextAt{8}{357.495}{512.627}{possible}%
\TextAt{8}{401.514}{512.627}{extensions}%
\TextAt{8}{457.894}{512.627}{of}%
\TextAt{13}{471.224}{512.627}{v}%
\TextAt{8}{481.21}{512.627}{to}%
\TextAt{13}{495.52}{512.627}{L}%
\TextAt{8}{503.485}{512.627}{.}%
\TextAt{12}{79.226}{486.615}{Extension}%
\TextAt{12}{141.441}{486.615}{Theorem}%
\TextAt{12}{198.168}{486.615}{(8.1).}%
\TextAt{8}{233.234}{486.615}{Let}%
\TextAt{13}{253.954}{486.615}{L}%
\TextAt{16}{265.239}{486.615}{\Character{106}}%
\TextAt{13}{271.881}{486.615}{K}%
\TextAt{8}{286.346}{486.615}{be}%
\TextAt{8}{302.021}{486.615}{an}%
\TextAt{8}{318.031}{486.615}{algebraic}%
\TextAt{8}{367.216}{486.615}{field}%
\TextAt{8}{392.323}{486.615}{extension}%
\TextAt{8}{443.839}{486.615}{and}%
\TextAt{13}{466.352}{486.615}{v}%
\TextAt{8}{476.088}{486.615}{a}%
\TextAt{8}{485.594}{486.615}{valuation}%
\TextAt{8}{79.226}{469.569}{of}%
\TextAt{13}{92.556}{469.569}{K}%
\TextAt{8}{103.374}{469.569}{.}%
\TextAt{8}{111.822}{469.569}{Then}%
\TextAt{8}{142.38}{469.569}{one}%
\TextAt{8}{163.839}{469.569}{has:}%
\TextAt{8}{89.579}{449.534}{(i)}%
\TextAt{8}{107.918}{449.534}{Every}%
\TextAt{8}{142.434}{449.534}{extension}%
\TextAt{13}{195.062}{449.534}{w}%
\TextAt{8}{208.511}{449.534}{of}%
\TextAt{8}{222.703}{449.534}{the}%
\TextAt{8}{243.721}{449.534}{valuation}%
\TextAt{13}{295.656}{449.534}{v}%
\TextAt{8}{306.503}{449.534}{arises}%
\TextAt{8}{339.357}{449.534}{as}%
\TextAt{8}{354.59}{449.534}{the}%
\TextAt{8}{375.608}{449.534}{composite}%
\TextAt{13}{431.488}{449.534}{w}%
\TextAt{9}{444.967}{449.534}{\Character{61}}%
\TextAt{9}{459.289}{449.534}{\Character{22}}%
\TextAt{13}{458.861}{449.534}{v}%
\TextAt{16}{468.194}{449.534}{\Character{14}}%
\TextAt{13}{477.416}{449.534}{\Character{28}}%
\TextAt{8}{488.594}{449.534}{for}%
\TextAt{8}{507.341}{449.534}{some}%
\TextAt{13}{107.918}{432.488}{K}%
\TextAt{8}{118.736}{432.488}{-embedding}%
\TextAt{13}{181.812}{432.488}{\Character{28}}%
\TextAt{9}{191.551}{432.488}{\Character{58}}%
\TextAt{13}{198.124}{432.488}{L}%
\NativeMathAt{209.409}{432.488}{11.955}{\rightarrow}%
\Vector{q 0 G 0 g 0.478 w 152.685 -331.669 m 163.503 -331.669 l S Q}%
\TextAt{13}{224.685}{432.488}{K}%
\TextAt{21}{235.503}{430.695}{v}%
\TextAt{8}{240.411}{432.488}{.}%
\TextAt{8}{86.327}{415.442}{(ii)}%
\TextAt{8}{107.918}{415.442}{Two}%
\TextAt{8}{133.382}{415.442}{extensions}%
\TextAt{9}{190.312}{415.442}{\Character{22}}%
\TextAt{13}{189.87}{415.442}{v}%
\TextAt{16}{198.687}{415.442}{\Character{14}}%
\TextAt{13}{207.394}{415.442}{\Character{28}}%
\TextAt{8}{217.818}{415.442}{and}%
\TextAt{9}{241.13}{415.442}{\Character{22}}%
\TextAt{13}{240.687}{415.442}{v}%
\TextAt{16}{249.505}{415.442}{\Character{14}}%
\TextAt{13}{258.212}{415.442}{\Character{28}}%
\TextAt{3}{264.631}{419.78}{\Character{48}}%
\TextAt{8}{271.43}{415.442}{are}%
\TextAt{8}{291.048}{415.442}{equal}%
\TextAt{8}{322.047}{415.442}{if}%
\TextAt{8}{332.878}{415.442}{and}%
\TextAt{8}{355.747}{415.442}{only}%
\TextAt{8}{381.546}{415.442}{if}%
\TextAt{13}{392.376}{415.442}{\Character{28}}%
\TextAt{8}{402.799}{415.442}{and}%
\TextAt{13}{425.669}{415.442}{\Character{28}}%
\TextAt{3}{432.088}{419.78}{\Character{48}}%
\TextAt{8}{438.887}{415.442}{are}%
\TextAt{8}{458.505}{415.442}{conjugate}%
\TextAt{8}{511.621}{415.442}{over}%
\TextAt{13}{107.918}{398.396}{K}%
\TextAt{21}{117.878}{396.603}{v}%
\TextAt{8}{122.785}{398.396}{.}%
\TextAt{19}{79.226}{363.417}{Proof.}%
\TextAt{8}{116}{363.417}{(i)}%
\TextAt{8}{133.835}{363.417}{Let}%
\TextAt{13}{156.38}{363.417}{w}%
\TextAt{8}{170.542}{363.417}{be}%
\TextAt{8}{188.042}{363.417}{an}%
\TextAt{8}{205.877}{363.417}{extension}%
\TextAt{8}{259.218}{363.417}{of}%
\TextAt{13}{274.124}{363.417}{v}%
\TextAt{8}{285.685}{363.417}{to}%
\TextAt{13}{301.571}{363.417}{L}%
\TextAt{8}{309.535}{363.417}{,}%
\TextAt{8}{318.654}{363.417}{and}%
\TextAt{13}{342.992}{363.417}{L}%
\TextAt{21}{350.957}{361.624}{w}%
\TextAt{8}{363.219}{363.417}{the}%
\TextAt{8}{384.951}{363.417}{localization}%
\TextAt{8}{448.646}{363.417}{of}%
\TextAt{8}{463.552}{363.417}{the}%
\TextAt{8}{485.284}{363.417}{canonical}%
\TextAt{8}{79.226}{346.371}{valuation,}%
\TextAt{8}{133.987}{346.371}{which}%
\TextAt{8}{167.823}{346.371}{is}%
\TextAt{8}{179.936}{346.371}{again}%
\TextAt{8}{211.513}{346.371}{denoted}%
\TextAt{8}{256.085}{346.371}{by}%
\TextAt{13}{272.693}{346.371}{w}%
\TextAt{8}{281.383}{346.371}{.}%
\TextAt{8}{290.879}{346.371}{This}%
\TextAt{8}{317.949}{346.371}{is}%
\TextAt{8}{330.062}{346.371}{the}%
\TextAt{8}{350.568}{346.371}{unique}%
\TextAt{8}{388.959}{346.371}{extension}%
\TextAt{8}{441.074}{346.371}{of}%
\TextAt{8}{454.753}{346.371}{the}%
\TextAt{8}{475.259}{346.371}{valuation}%
\TextAt{13}{526.681}{346.371}{v}%
\TextAt{8}{79.226}{329.325}{from}%
\TextAt{13}{107.566}{329.325}{K}%
\TextAt{21}{117.525}{327.532}{v}%
\TextAt{8}{127.03}{329.325}{to}%
\TextAt{13}{142.041}{329.325}{L}%
\TextAt{21}{150.005}{327.532}{w}%
\TextAt{8}{156.794}{329.325}{.}%
\TextAt{8}{167.342}{329.325}{Choosing}%
\TextAt{8}{219.163}{329.325}{any}%
\TextAt{13}{241.98}{329.325}{K}%
\TextAt{21}{251.939}{327.532}{v}%
\TextAt{8}{256.846}{329.325}{-embedding}%
\TextAt{13}{320.622}{329.325}{\Character{28}}%
\TextAt{9}{331.554}{329.325}{\Character{58}}%
\TextAt{13}{339.319}{329.325}{L}%
\TextAt{21}{347.284}{327.532}{w}%
\NativeMathAt{358.587}{329.325}{11.955}{\rightarrow}%
\Vector{q 0 G 0 g 0.478 w 303.055 -434.832 m 313.873 -434.832 l S Q}%
\TextAt{13}{375.055}{329.325}{K}%
\TextAt{21}{385.873}{327.532}{v}%
\TextAt{8}{390.781}{329.325}{,}%
\TextAt{8}{398.805}{329.325}{the}%
\TextAt{8}{419.661}{329.325}{valuation}%
\TextAt{9}{471.876}{329.325}{\Character{22}}%
\TextAt{13}{471.433}{329.325}{v}%
\TextAt{16}{480.655}{329.325}{\Character{14}}%
\TextAt{13}{489.767}{329.325}{\Character{28}}%
\TextAt{8}{500.783}{329.325}{has}%
\TextAt{8}{522.357}{329.325}{to}%
\TextAt{8}{79.226}{312.279}{coincide}%
\TextAt{8}{124.335}{312.279}{with}%
\TextAt{13}{151.236}{312.279}{w}%
\TextAt{8}{159.925}{312.279}{.}%
\TextAt{8}{169.097}{312.279}{The}%
\TextAt{8}{193.392}{312.279}{restriction}%
\TextAt{8}{249.631}{312.279}{of}%
\TextAt{13}{263.202}{312.279}{\Character{28}}%
\TextAt{8}{273.759}{312.279}{to}%
\TextAt{13}{288.31}{312.279}{L}%
\TextAt{8}{300.413}{312.279}{is}%
\TextAt{8}{312.418}{312.279}{therefore}%
\TextAt{8}{361.759}{312.279}{a}%
\TextAt{13}{371.756}{312.279}{K}%
\TextAt{8}{382.573}{312.279}{-embedding}%
\TextAt{13}{445.89}{312.279}{\Character{28}}%
\TextAt{9}{456.04}{312.279}{\Character{58}}%
\TextAt{13}{463.023}{312.279}{L}%
\NativeMathAt{474.719}{312.279}{11.955}{\rightarrow}%
\Vector{q 0 G 0 g 0.478 w 418.406 -451.878 m 429.224 -451.878 l S Q}%
\TextAt{13}{490.406}{312.279}{K}%
\TextAt{21}{501.224}{310.486}{v}%
\TextAt{8}{510.27}{312.279}{such}%
\TextAt{8}{79.226}{295.233}{that}%
\TextAt{13}{104.595}{295.233}{w}%
\TextAt{9}{116.605}{295.233}{\Character{61}}%
\TextAt{9}{129.457}{295.233}{\Character{22}}%
\TextAt{13}{129.031}{295.233}{v}%
\TextAt{16}{137.776}{295.233}{\Character{14}}%
\TextAt{13}{146.41}{295.233}{\Character{28}}%
\TextAt{8}{152.829}{295.233}{.}%
\TextAt{8}{97.159}{275.199}{(ii)}%
\TextAt{8}{116.848}{275.199}{Let}%
\TextAt{13}{137.996}{275.199}{\Character{28}}%
\TextAt{8}{148.491}{275.199}{and}%
\TextAt{13}{171.432}{275.199}{\Character{27}}%
\TextAt{13}{178.498}{275.199}{\Character{28}}%
\TextAt{8}{184.933}{275.199}{,}%
\TextAt{8}{192.306}{275.199}{with}%
\TextAt{13}{219.144}{275.199}{\Character{27}}%
\NativeMathAt{229.852}{275.199}{11.955}{\in}%
\TextAt{13}{241.447}{275.199}{G}%
\TextAt{9}{250.68}{275.199}{\Character{40}}%
\Vector{q 0 G 0 g 0.478 w 183.233 -488.958 m 194.05 -488.958 l S Q}%
\TextAt{13}{255.233}{275.199}{K}%
\TextAt{21}{266.05}{273.405}{v}%
\TextAt{16}{274.583}{275.199}{\Character{106}}%
\TextAt{13}{281.529}{275.199}{K}%
\TextAt{21}{291.488}{273.405}{v}%
\TextAt{9}{296.396}{275.199}{\Character{41}}%
\TextAt{8}{300.948}{275.199}{,}%
\TextAt{8}{308.32}{275.199}{be}%
\TextAt{8}{324.423}{275.199}{two}%
\TextAt{8}{346.719}{275.199}{embeddings}%
\TextAt{8}{410.69}{275.199}{of}%
\TextAt{13}{424.198}{275.199}{L}%
\TextAt{8}{436.239}{275.199}{conjugate}%
\TextAt{8}{489.426}{275.199}{over}%
\TextAt{13}{514.651}{275.199}{K}%
\TextAt{21}{524.61}{273.405}{v}%
\TextAt{8}{529.518}{275.199}{.}%
\TextAt{8}{79.226}{258.153}{Since}%
\TextAt{9}{110.713}{258.153}{\Character{22}}%
\TextAt{13}{110.27}{258.153}{v}%
\TextAt{8}{120.743}{258.153}{is}%
\TextAt{8}{132.994}{258.153}{the}%
\TextAt{8}{153.637}{258.153}{only}%
\TextAt{8}{179.816}{258.153}{extension}%
\TextAt{8}{232.069}{258.153}{of}%
\TextAt{8}{245.886}{258.153}{the}%
\TextAt{8}{266.529}{258.153}{valuation}%
\TextAt{13}{318.088}{258.153}{v}%
\TextAt{8}{328.561}{258.153}{from}%
\TextAt{13}{356.688}{258.153}{K}%
\TextAt{21}{366.647}{256.359}{v}%
\TextAt{8}{375.939}{258.153}{to}%
\Vector{q 0 G 0 g 0.478 w 318.737 -506.004 m 329.554 -506.004 l S Q}%
\TextAt{13}{390.737}{258.153}{K}%
\TextAt{21}{401.554}{256.359}{v}%
\TextAt{8}{406.462}{258.153}{,}%
\TextAt{8}{414.22}{258.153}{one}%
\TextAt{8}{436.166}{258.153}{has}%
\TextAt{9}{457.97}{258.153}{\Character{22}}%
\TextAt{13}{457.527}{258.153}{v}%
\TextAt{9}{467.766}{258.153}{\Character{61}}%
\TextAt{9}{481.455}{258.153}{\Character{22}}%
\TextAt{13}{481.021}{258.153}{v}%
\TextAt{16}{490.098}{258.153}{\Character{14}}%
\TextAt{13}{499.064}{258.153}{\Character{27}}%
\TextAt{8}{506.146}{258.153}{,}%
\TextAt{8}{513.904}{258.153}{and}%
\TextAt{8}{79.226}{241.107}{thus}%
\TextAt{9}{105.867}{241.107}{\Character{22}}%
\TextAt{13}{105.424}{241.107}{v}%
\TextAt{16}{114.473}{241.107}{\Character{14}}%
\TextAt{13}{123.411}{241.107}{\Character{28}}%
\TextAt{9}{133.912}{241.107}{\Character{61}}%
\TextAt{9}{147.529}{241.107}{\Character{22}}%
\TextAt{13}{147.098}{241.107}{v}%
\TextAt{16}{156.147}{241.107}{\Character{14}}%
\TextAt{9}{165.085}{241.107}{\Character{40}}%
\TextAt{13}{169.638}{241.107}{\Character{27}}%
\TextAt{16}{179.681}{241.107}{\Character{14}}%
\TextAt{13}{188.619}{241.107}{\Character{28}}%
\TextAt{9}{195.038}{241.107}{\Character{41}}%
\TextAt{8}{199.59}{241.107}{.}%
\TextAt{8}{209.378}{241.107}{The}%
\TextAt{8}{233.878}{241.107}{extensions}%
\TextAt{8}{290.704}{241.107}{induced}%
\TextAt{8}{334.715}{241.107}{to}%
\TextAt{13}{349.472}{241.107}{L}%
\TextAt{8}{361.78}{241.107}{by}%
\TextAt{13}{378.485}{241.107}{\Character{28}}%
\TextAt{8}{389.248}{241.107}{and}%
\TextAt{8}{412.457}{241.107}{by}%
\TextAt{13}{429.162}{241.107}{\Character{27}}%
\TextAt{13}{436.228}{241.107}{\Character{28}}%
\TextAt{8}{447.007}{241.107}{are}%
\TextAt{8}{466.964}{241.107}{therefore}%
\TextAt{8}{516.51}{241.107}{the}%
\TextAt{8}{79.226}{224.061}{same.}%
\TextAt{8}{97.159}{204.026}{Conversely,}%
\TextAt{8}{159.76}{204.026}{let}%
\TextAt{13}{177.529}{204.026}{\Character{28}}%
\TextAt{13}{183.291}{204.026}{\Character{59}}%
\TextAt{13}{188.516}{204.026}{\Character{28}}%
\TextAt{3}{194.96}{208.364}{\Character{48}}%
\TextAt{9}{202.548}{204.026}{\Character{58}}%
\TextAt{13}{210.593}{204.026}{L}%
\NativeMathAt{223.351}{204.026}{11.955}{\rightarrow}%
\Vector{q 0 G 0 g 0.478 w 168.099 -560.131 m 178.917 -560.131 l S Q}%
\TextAt{13}{240.099}{204.026}{K}%
\TextAt{21}{250.917}{202.233}{v}%
\TextAt{8}{260.585}{204.026}{be}%
\TextAt{8}{277.374}{204.026}{two}%
\TextAt{13}{300.355}{204.026}{K}%
\TextAt{8}{311.172}{204.026}{-embeddings}%
\TextAt{8}{379.726}{204.026}{such}%
\TextAt{8}{406.987}{204.026}{that}%
\TextAt{9}{433.663}{204.026}{\Character{22}}%
\TextAt{13}{433.22}{204.026}{v}%
\TextAt{16}{442.554}{204.026}{\Character{14}}%
\TextAt{13}{451.777}{204.026}{\Character{28}}%
\TextAt{9}{462.989}{204.026}{\Character{61}}%
\TextAt{9}{477.311}{204.026}{\Character{22}}%
\TextAt{13}{476.887}{204.026}{v}%
\TextAt{16}{486.22}{204.026}{\Character{14}}%
\TextAt{13}{495.443}{204.026}{\Character{28}}%
\TextAt{3}{501.862}{208.364}{\Character{48}}%
\TextAt{8}{504.657}{204.026}{.}%
\TextAt{8}{515.698}{204.026}{Let}%
\TextAt{13}{79.226}{186.98}{\Character{27}}%
\TextAt{9}{89.629}{186.98}{\Character{58}}%
\TextAt{13}{96.202}{186.98}{\Character{28}}%
\TextAt{13}{102.61}{186.98}{L}%
\NativeMathAt{113.906}{186.98}{11.955}{\rightarrow}%
\TextAt{13}{129.182}{186.98}{\Character{28}}%
\TextAt{3}{135.601}{191.318}{\Character{48}}%
\TextAt{13}{138.396}{186.98}{L}%
\TextAt{8}{149.967}{186.98}{be}%
\TextAt{8}{165.6}{186.98}{the}%
\TextAt{13}{185.465}{186.98}{K}%
\TextAt{8}{196.283}{186.98}{-isomorphism}%
\TextAt{13}{268.308}{186.98}{\Character{27}}%
\TextAt{9}{278.712}{186.98}{\Character{61}}%
\TextAt{13}{291.137}{186.98}{\Character{28}}%
\TextAt{3}{297.556}{191.318}{\Character{48}}%
\TextAt{16}{302.412}{186.98}{\Character{14}}%
\TextAt{13}{310.451}{186.98}{\Character{28}}%
\TextAt{3}{316.87}{191.318}{\Character{0}}%
\TextAt{17}{323.456}{191.318}{1}%
\TextAt{8}{328.189}{186.98}{.}%
\TextAt{8}{336.54}{186.98}{We}%
\TextAt{8}{356.393}{186.98}{can}%
\TextAt{8}{377.561}{186.98}{extend}%
\TextAt{13}{415.312}{186.98}{\Character{27}}%
\TextAt{8}{426}{186.98}{to}%
\TextAt{8}{440.019}{186.98}{a}%
\TextAt{13}{449.483}{186.98}{K}%
\TextAt{21}{459.442}{185.187}{v}%
\TextAt{8}{464.35}{186.98}{-isomorphism}%
\TextAt{13}{243.651}{157.979}{\Character{27}}%
\TextAt{9}{254.054}{157.979}{\Character{58}}%
\TextAt{13}{260.627}{157.979}{\Character{28}}%
\TextAt{13}{267.035}{157.979}{L}%
\TextAt{16}{277.667}{157.979}{\Character{1}}%
\TextAt{13}{283.644}{157.979}{K}%
\TextAt{21}{293.603}{156.186}{v}%
\NativeArrowAt{301.832}{157.979}{11.955}{19.26}{\longrightarrow}%
\TextAt{13}{324.414}{157.979}{\Character{28}}%
\TextAt{3}{330.833}{162.915}{\Character{48}}%
\TextAt{13}{333.628}{157.979}{L}%
\TextAt{16}{344.249}{157.979}{\Character{1}}%
\TextAt{13}{350.226}{157.979}{K}%
\TextAt{21}{360.186}{156.186}{v}%
\TextAt{13}{365.093}{157.979}{\Character{58}}%
\TextAt{8}{79.226}{128.978}{Indeed,}%
\TextAt{13}{120.507}{128.978}{\Character{28}}%
\TextAt{13}{126.915}{128.978}{L}%
\TextAt{8}{138.788}{128.978}{is}%
\TextAt{8}{150.552}{128.978}{dense}%
\TextAt{8}{182.472}{128.978}{in}%
\TextAt{13}{196.125}{128.978}{\Character{28}}%
\TextAt{13}{202.533}{128.978}{L}%
\TextAt{16}{213.165}{128.978}{\Character{1}}%
\TextAt{13}{219.143}{128.978}{K}%
\TextAt{21}{229.102}{127.184}{v}%
\TextAt{8}{234.01}{128.978}{,}%
\TextAt{8}{241.159}{128.978}{so}%
\TextAt{8}{255.529}{128.978}{every}%
\TextAt{8}{286.421}{128.978}{element}%
\TextAt{13}{329.663}{128.978}{x}%
\NativeMathAt{339.636}{128.978}{11.955}{\in}%
\TextAt{13}{350.927}{128.978}{\Character{28}}%
\TextAt{13}{357.335}{128.978}{L}%
\TextAt{16}{367.967}{128.978}{\Character{1}}%
\TextAt{13}{373.944}{128.978}{K}%
\TextAt{21}{383.904}{127.184}{v}%
\TextAt{8}{392.708}{128.978}{can}%
\TextAt{8}{414.168}{128.978}{be}%
\TextAt{8}{430.092}{128.978}{written}%
\TextAt{8}{471.063}{128.978}{as}%
\TextAt{8}{485.433}{128.978}{a}%
\TextAt{8}{495.188}{128.978}{limit}%
\TextAt{13}{273.411}{99.976}{x}%
\TextAt{9}{283.384}{99.976}{\Character{61}}%
\TextAt{8}{298.718}{99.976}{lim}%
\TextAt{21}{295.81}{92.803}{n}%
\NativeMathAt{300.948}{92.803}{7.97}{\rightarrow}%
\NativeMathAt{309.412}{92.803}{7.97}{\infty}%
\TextAt{13}{319.877}{99.976}{\Character{28}}%
\TextAt{13}{326.285}{99.976}{x}%
\TextAt{21}{332.948}{98.183}{n}%
\TextAt{8}{297.211}{43.033}{148}%
\Destination{79.23}{466.465}{nav.Item.51}%
\Destination{79.23}{430.035}{nav.Item.52}%
\Destination{78.23}{809.815}{nav.page.148}%
\Destination{79.36}{464.555}{nav.print.162}%
\Destination{79.226}{490}{nav.number.2.8.1}%
\Hotlink{174.491}{758.924}{19}{11.047}{nav.local.2.8.diagram}%
\Hotlink{517.802}{656.648}{18}{11.047}{nav.local.2.8.diagram}%
}

% Page 161 (body page 149)
\Page{%
\TextAt{2}{79.226}{796.943}{II.}%
\TextAt{2}{94.073}{796.943}{The}%
\TextAt{2}{116.492}{796.943}{Theory}%
\TextAt{2}{154.401}{796.943}{of}%
\TextAt{2}{166.826}{796.943}{Valuations}%
\Vector{q 0 G 0 g 0.299 w 7.226 18.729 m 460.769 18.729 l S Q}%
\TextAt{8}{79.226}{761.124}{for}%
\TextAt{8}{97.374}{761.124}{some}%
\TextAt{8}{126.963}{761.124}{sequence}%
\TextAt{13}{175.728}{761.124}{x}%
\TextAt{21}{182.38}{759.33}{n}%
\TextAt{8}{192.176}{761.124}{which}%
\TextAt{8}{225.926}{761.124}{belongs}%
\TextAt{8}{268.199}{761.124}{to}%
\TextAt{8}{282.772}{761.124}{a}%
\TextAt{8}{292.79}{761.124}{finite}%
\TextAt{8}{322.965}{761.124}{subextension}%
\TextAt{8}{392.938}{761.124}{of}%
\TextAt{13}{406.531}{761.124}{L}%
\TextAt{8}{414.496}{761.124}{.}%
\TextAt{8}{423.733}{761.124}{As}%
\TextAt{9}{441.725}{761.124}{\Character{22}}%
\TextAt{13}{441.283}{761.124}{v}%
\TextAt{16}{450.207}{761.124}{\Character{14}}%
\TextAt{13}{459.02}{761.124}{\Character{28}}%
\TextAt{9}{469.208}{761.124}{\Character{61}}%
\TextAt{9}{482.514}{761.124}{\Character{22}}%
\TextAt{13}{482.081}{761.124}{v}%
\TextAt{16}{491.005}{761.124}{\Character{14}}%
\TextAt{13}{499.819}{761.124}{\Character{28}}%
\TextAt{3}{506.238}{765.462}{\Character{48}}%
\TextAt{8}{509.033}{761.124}{,}%
\TextAt{8}{516.51}{761.124}{the}%
\TextAt{8}{79.226}{744.078}{sequence}%
\TextAt{13}{127.728}{744.078}{\Character{28}}%
\TextAt{3}{134.147}{748.416}{\Character{48}}%
\TextAt{13}{136.942}{744.078}{x}%
\TextAt{21}{143.594}{742.284}{n}%
\TextAt{9}{152.552}{744.078}{\Character{61}}%
\TextAt{13}{164.977}{744.078}{\Character{27}}%
\TextAt{13}{172.043}{744.078}{\Character{28}}%
\TextAt{13}{178.463}{744.078}{x}%
\TextAt{21}{185.13}{742.284}{n}%
\TextAt{8}{194.664}{744.078}{converges}%
\TextAt{8}{247.088}{744.078}{to}%
\TextAt{8}{261.398}{744.078}{an}%
\TextAt{8}{277.657}{744.078}{element}%
\TextAt{13}{264.703}{715.36}{\Character{27}}%
\TextAt{13}{271.769}{715.36}{x}%
\TextAt{9}{281.758}{715.36}{\Character{58}\Character{61}}%
\TextAt{8}{300.343}{715.36}{lim}%
\TextAt{21}{297.435}{708.187}{n}%
\NativeMathAt{302.574}{708.187}{7.97}{\rightarrow}%
\NativeMathAt{311.038}{708.187}{7.97}{\infty}%
\TextAt{13}{321.503}{715.36}{\Character{27}}%
\TextAt{13}{328.569}{715.36}{\Character{28}}%
\TextAt{13}{334.988}{715.36}{x}%
\TextAt{21}{341.656}{713.567}{n}%
\TextAt{8}{79.226}{685.557}{in}%
\TextAt{13}{94.253}{685.557}{\Character{28}}%
\TextAt{3}{100.672}{689.896}{\Character{48}}%
\TextAt{13}{103.467}{685.557}{L}%
\TextAt{16}{115.024}{685.557}{\Character{1}}%
\TextAt{13}{121.938}{685.557}{K}%
\TextAt{21}{131.897}{683.764}{v}%
\TextAt{8}{136.805}{685.557}{.}%
\TextAt{8}{149.374}{685.557}{Clearly}%
\TextAt{8}{191.396}{685.557}{the}%
\TextAt{8}{212.926}{685.557}{correspondence}%
\TextAt{13}{295.977}{685.557}{x}%
\NativeArrowAt{308.291}{685.557}{11.955}{11.955}{\mapsto}%
\TextAt{13}{325.908}{685.557}{\Character{27}}%
\TextAt{13}{332.974}{685.557}{x}%
\TextAt{8}{344.914}{685.557}{does}%
\TextAt{8}{372.685}{685.557}{not}%
\TextAt{8}{394.872}{685.557}{depend}%
\TextAt{8}{436.882}{685.557}{on}%
\TextAt{8}{454.514}{685.557}{the}%
\TextAt{8}{476.044}{685.557}{choice}%
\TextAt{8}{512.208}{685.557}{of}%
\TextAt{8}{526.911}{685.557}{a}%
\TextAt{8}{79.226}{668.512}{sequence}%
\TextAt{16}{128.876}{668.512}{\Character{102}}%
\TextAt{13}{134.853}{668.512}{x}%
\TextAt{21}{141.505}{666.718}{n}%
\TextAt{16}{147.142}{668.512}{\Character{103}}%
\TextAt{8}{153.119}{668.512}{,}%
\TextAt{8}{161.703}{668.512}{and}%
\TextAt{8}{185.613}{668.512}{yields}%
\TextAt{8}{219.661}{668.512}{an}%
\TextAt{8}{237.067}{668.512}{isomorphism}%
\TextAt{13}{306.634}{668.512}{\Character{28}}%
\TextAt{13}{313.042}{668.512}{L}%
\TextAt{16}{324.456}{668.512}{\Character{1}}%
\TextAt{13}{331.216}{668.512}{K}%
\TextAt{21}{341.175}{666.718}{v}%
\NativeMathAt{351.359}{668.512}{11.955}{\rightarrow}%
\TextAt{13}{368.59}{668.512}{\Character{28}}%
\TextAt{3}{375.009}{672.85}{\Character{48}}%
\TextAt{13}{377.804}{668.512}{L}%
\TextAt{16}{389.207}{668.512}{\Character{1}}%
\TextAt{13}{395.967}{668.512}{K}%
\TextAt{21}{405.926}{666.718}{v}%
\TextAt{8}{415.878}{668.512}{which}%
\TextAt{8}{450.512}{668.512}{leaves}%
\TextAt{13}{485.218}{668.512}{K}%
\TextAt{21}{495.177}{666.718}{v}%
\TextAt{8}{505.129}{668.512}{fixed.}%
\TextAt{8}{79.226}{651.466}{Extending}%
\TextAt{13}{135.301}{651.466}{\Character{27}}%
\TextAt{8}{145.94}{651.466}{to}%
\TextAt{8}{159.909}{651.466}{a}%
\TextAt{13}{169.324}{651.466}{K}%
\TextAt{21}{179.283}{649.672}{v}%
\TextAt{8}{184.19}{651.466}{-automorphism}%
\TextAt{9}{265.831}{651.466}{\Character{126}}%
\TextAt{13}{265.216}{651.466}{\Character{27}}%
\NativeMathAt{275.619}{651.466}{11.955}{\in}%
\TextAt{13}{286.91}{651.466}{G}%
\TextAt{9}{296.144}{651.466}{\Character{40}}%
\Vector{q 0 G 0 g 0.478 w 228.696 -112.691 m 239.514 -112.691 l S Q}%
\TextAt{13}{300.696}{651.466}{K}%
\TextAt{21}{311.514}{649.672}{v}%
\TextAt{16}{319.742}{651.466}{\Character{106}}%
\TextAt{13}{326.384}{651.466}{K}%
\TextAt{21}{336.343}{649.672}{v}%
\TextAt{9}{341.251}{651.466}{\Character{41}}%
\TextAt{8}{349.359}{651.466}{gives}%
\TextAt{13}{377.699}{651.466}{\Character{28}}%
\TextAt{3}{384.118}{655.804}{\Character{48}}%
\TextAt{9}{390.233}{651.466}{\Character{61}}%
\TextAt{9}{403.264}{651.466}{\Character{126}}%
\TextAt{13}{402.659}{651.466}{\Character{27}}%
\TextAt{16}{411.7}{651.466}{\Character{14}}%
\TextAt{13}{419.637}{651.466}{\Character{28}}%
\TextAt{8}{426.056}{651.466}{,}%
\TextAt{8}{432.932}{651.466}{so}%
\TextAt{8}{446.961}{651.466}{that}%
\TextAt{13}{471.989}{651.466}{\Character{28}}%
\TextAt{8}{481.964}{651.466}{and}%
\TextAt{13}{504.386}{651.466}{\Character{28}}%
\TextAt{3}{510.805}{655.804}{\Character{48}}%
\TextAt{8}{517.156}{651.466}{are}%
\TextAt{8}{79.226}{634.42}{indeed}%
\TextAt{8}{116.287}{634.42}{conjugate}%
\TextAt{8}{169.296}{634.42}{over}%
\TextAt{13}{194.342}{634.42}{K}%
\TextAt{21}{204.302}{632.626}{v}%
\TextAt{8}{209.209}{634.42}{.}%
\Vector{q 0 G 0 g 0.399 w 452.334 -139.58 m 452.334 -131.511 l S Q}%
\Vector{q 0 G 0 g 0.399 w 452.534 -131.71 m 459.707 -131.71 l S Q}%
\Vector{q 0 G 0 g 0.399 w 452.534 -139.381 m 459.707 -139.381 l S Q}%
\Vector{q 0 G 0 g 0.399 w 459.906 -139.58 m 459.906 -131.511 l S Q}%
\TextAt{8}{97.159}{605.499}{Those}%
\TextAt{8}{131.849}{605.499}{who}%
\TextAt{8}{156.725}{605.499}{prefer}%
\TextAt{8}{190.375}{605.499}{to}%
\TextAt{8}{204.849}{605.499}{be}%
\TextAt{8}{220.937}{605.499}{given}%
\TextAt{8}{251.671}{605.499}{an}%
\TextAt{8}{268.094}{605.499}{extension}%
\TextAt{13}{320.023}{605.499}{L}%
\TextAt{16}{331.588}{605.499}{\Character{106}}%
\TextAt{13}{338.509}{605.499}{K}%
\TextAt{8}{353.388}{605.499}{by}%
\TextAt{8}{369.811}{605.499}{an}%
\TextAt{8}{386.234}{605.499}{algebraic}%
\TextAt{8}{435.833}{605.499}{equation}%
\TextAt{13}{483.805}{605.499}{f}%
\TextAt{9}{490.852}{605.499}{\Character{40}}%
\TextAt{13}{495.404}{605.499}{X}%
\TextAt{9}{506.059}{605.499}{\Character{41}}%
\TextAt{9}{514.201}{605.499}{\Character{61}}%
\TextAt{9}{526.897}{605.499}{0}%
\TextAt{8}{79.226}{588.453}{will}%
\TextAt{8}{101.331}{588.453}{appreciate}%
\TextAt{8}{157.915}{588.453}{the}%
\TextAt{8}{178.071}{588.453}{following}%
\TextAt{8}{227.506}{588.453}{concrete}%
\TextAt{8}{273.677}{588.453}{variant}%
\TextAt{8}{313.368}{588.453}{of}%
\TextAt{8}{326.698}{588.453}{the}%
\TextAt{8}{346.855}{588.453}{above}%
\TextAt{8}{380.03}{588.453}{extension}%
\TextAt{8}{431.796}{588.453}{theorem.}%
\TextAt{8}{97.159}{568.419}{Let}%
\TextAt{13}{118.151}{568.419}{L}%
\TextAt{9}{129.476}{568.419}{\Character{61}}%
\TextAt{13}{141.941}{568.419}{K}%
\TextAt{9}{152.759}{568.419}{\Character{40}}%
\TextAt{13}{157.311}{568.419}{\Character{11}}%
\TextAt{9}{164.833}{568.419}{\Character{41}}%
\TextAt{8}{173.306}{568.419}{be}%
\TextAt{8}{189.253}{568.419}{generated}%
\TextAt{8}{242.608}{568.419}{by}%
\TextAt{8}{258.89}{568.419}{the}%
\TextAt{8}{279.07}{568.419}{zero}%
\TextAt{13}{303.804}{568.419}{\Character{11}}%
\TextAt{8}{315.247}{568.419}{of}%
\TextAt{8}{328.6}{568.419}{an}%
\TextAt{8}{344.882}{568.419}{irreducible}%
\TextAt{8}{402.78}{568.419}{polynomial}%
\TextAt{13}{463.297}{568.419}{f}%
\TextAt{9}{470.343}{568.419}{\Character{40}}%
\TextAt{13}{474.895}{568.419}{X}%
\TextAt{9}{485.55}{568.419}{\Character{41}}%
\NativeMathAt{493.463}{568.419}{11.955}{\in}%
\TextAt{13}{504.793}{568.419}{K}%
\TextAt{9}{515.611}{568.419}{\Character{91}}%
\TextAt{13}{518.863}{568.419}{X}%
\TextAt{9}{529.518}{568.419}{\Character{93}}%
\TextAt{8}{79.226}{551.373}{and}%
\TextAt{8}{101.989}{551.373}{let}%
\TextAt{13}{233.439}{534.327}{f}%
\TextAt{9}{240.485}{534.327}{\Character{40}}%
\TextAt{13}{245.038}{534.327}{X}%
\TextAt{9}{255.693}{534.327}{\Character{41}}%
\TextAt{9}{263.548}{534.327}{\Character{61}}%
\TextAt{13}{275.991}{534.327}{f}%
\TextAt{17}{281.761}{532.534}{1}%
\TextAt{9}{286.494}{534.327}{\Character{40}}%
\TextAt{13}{291.046}{534.327}{X}%
\TextAt{9}{301.701}{534.327}{\Character{41}}%
\TextAt{21}{306.253}{539.263}{m}%
\TextAt{24}{313.744}{538.156}{1}%
\TextAt{16}{320.386}{534.327}{\Character{1}}%
\TextAt{16}{325.682}{534.327}{\Character{1}}%
\TextAt{16}{330.99}{534.327}{\Character{1}}%
\TextAt{13}{336.326}{534.327}{f}%
\TextAt{21}{342.096}{532.534}{r}%
\TextAt{9}{346.65}{534.327}{\Character{40}}%
\TextAt{13}{351.203}{534.327}{X}%
\TextAt{9}{361.858}{534.327}{\Character{41}}%
\TextAt{21}{366.41}{539.263}{m}%
\TextAt{22}{373.901}{538.267}{r}%
\TextAt{8}{79.226}{510.927}{be}%
\TextAt{8}{95.262}{510.927}{the}%
\TextAt{8}{115.53}{510.927}{decomposition}%
\TextAt{8}{192.777}{510.927}{of}%
\TextAt{13}{206.218}{510.927}{f}%
\TextAt{9}{213.265}{510.927}{\Character{40}}%
\TextAt{13}{217.817}{510.927}{X}%
\TextAt{9}{228.472}{510.927}{\Character{41}}%
\TextAt{8}{237.034}{510.927}{into}%
\TextAt{8}{260.888}{510.927}{irreducible}%
\TextAt{8}{318.875}{510.927}{factors}%
\TextAt{13}{357.1}{510.927}{f}%
\TextAt{17}{362.87}{509.133}{1}%
\TextAt{9}{367.602}{510.927}{\Character{40}}%
\TextAt{13}{372.155}{510.927}{X}%
\TextAt{9}{382.81}{510.927}{\Character{41}}%
\TextAt{13}{387.362}{510.927}{\Character{59}}%
\TextAt{13}{392.586}{510.927}{\Character{58}}%
\TextAt{13}{397.823}{510.927}{\Character{58}}%
\TextAt{13}{403.047}{510.927}{\Character{58}}%
\TextAt{13}{408.284}{510.927}{\Character{59}}%
\TextAt{13}{413.52}{510.927}{f}%
\TextAt{21}{419.353}{509.133}{r}%
\TextAt{9}{423.907}{510.927}{\Character{40}}%
\TextAt{13}{428.46}{510.927}{X}%
\TextAt{9}{439.115}{510.927}{\Character{41}}%
\TextAt{8}{447.676}{510.927}{over}%
\TextAt{8}{472.834}{510.927}{the}%
\TextAt{8}{493.102}{510.927}{comple-}%
\TextAt{8}{79.226}{493.881}{tion}%
\TextAt{13}{103.359}{493.881}{K}%
\TextAt{21}{113.318}{492.087}{v}%
\TextAt{8}{118.226}{493.881}{.}%
\TextAt{8}{126.877}{493.881}{Of}%
\TextAt{8}{143.526}{493.881}{course,}%
\TextAt{8}{182.692}{493.881}{the}%
\TextAt{13}{202.915}{493.881}{m}%
\TextAt{21}{213.155}{492.087}{i}%
\TextAt{8}{220.501}{493.881}{are}%
\TextAt{8}{240.079}{493.881}{one}%
\TextAt{8}{261.606}{493.881}{if}%
\TextAt{13}{272.397}{493.881}{f}%
\TextAt{8}{283.409}{493.881}{is}%
\TextAt{8}{295.24}{493.881}{separable.}%
\TextAt{8}{351.436}{493.881}{The}%
\TextAt{13}{375.557}{493.881}{K}%
\TextAt{8}{386.375}{493.881}{-embeddings}%
\TextAt{13}{454.133}{493.881}{\Character{28}}%
\TextAt{9}{463.987}{493.881}{\Character{58}}%
\TextAt{13}{470.675}{493.881}{L}%
\NativeMathAt{482.075}{493.881}{11.955}{\rightarrow}%
\Vector{q 0 G 0 g 0.478 w 425.466 -270.276 m 436.284 -270.276 l S Q}%
\TextAt{13}{497.466}{493.881}{K}%
\TextAt{21}{508.284}{492.087}{v}%
\TextAt{8}{517.156}{493.881}{are}%
\TextAt{8}{79.226}{476.835}{then}%
\TextAt{8}{105.886}{476.835}{given}%
\TextAt{8}{136.456}{476.835}{by}%
\TextAt{8}{152.715}{476.835}{the}%
\TextAt{8}{172.871}{476.835}{zeroes}%
\TextAt{13}{207.72}{476.835}{\Character{12}}%
\TextAt{8}{218.889}{476.835}{of}%
\TextAt{13}{232.219}{476.835}{f}%
\TextAt{9}{239.265}{476.835}{\Character{40}}%
\TextAt{13}{243.818}{476.835}{X}%
\TextAt{9}{254.473}{476.835}{\Character{41}}%
\TextAt{8}{262.923}{476.835}{which}%
\TextAt{8}{296.409}{476.835}{lie}%
\TextAt{8}{312.011}{476.835}{in}%
\Vector{q 0 G 0 g 0.478 w 253.663 -287.322 m 264.481 -287.322 l S Q}%
\TextAt{13}{325.663}{476.835}{K}%
\TextAt{21}{336.481}{475.041}{v}%
\TextAt{8}{341.389}{476.835}{:}%
\TextAt{13}{233.811}{448.117}{\Character{28}}%
\TextAt{9}{243.551}{448.117}{\Character{58}}%
\TextAt{13}{250.123}{448.117}{L}%
\NativeArrowAt{261.408}{448.117}{11.955}{19.26}{\longrightarrow}%
\Vector{q 0 G 0 g 0.478 w 211.99 -316.04 m 222.808 -316.04 l S Q}%
\TextAt{13}{283.99}{448.117}{K}%
\TextAt{21}{294.808}{446.324}{v}%
\TextAt{13}{299.716}{448.117}{\Character{59}}%
\TextAt{13}{328.851}{448.117}{\Character{28}}%
\TextAt{9}{335.289}{448.117}{\Character{40}}%
\TextAt{13}{339.841}{448.117}{\Character{11}}%
\TextAt{9}{347.363}{448.117}{\Character{41}}%
\TextAt{9}{355.218}{448.117}{\Character{61}}%
\TextAt{13}{367.662}{448.117}{\Character{12}}%
\TextAt{13}{374.931}{448.117}{\Character{58}}%
\TextAt{8}{79.226}{419.4}{Two}%
\TextAt{8}{104.667}{419.4}{embeddings}%
\TextAt{13}{168.543}{419.4}{\Character{28}}%
\TextAt{8}{178.943}{419.4}{and}%
\TextAt{13}{201.789}{419.4}{\Character{28}}%
\TextAt{3}{208.208}{423.738}{\Character{48}}%
\TextAt{8}{214.984}{419.4}{are}%
\TextAt{8}{234.578}{419.4}{conjugate}%
\TextAt{8}{287.671}{419.4}{over}%
\TextAt{13}{312.8}{419.4}{K}%
\TextAt{21}{322.76}{417.606}{v}%
\TextAt{8}{331.648}{419.4}{if}%
\TextAt{8}{342.455}{419.4}{and}%
\TextAt{8}{365.301}{419.4}{only}%
\TextAt{8}{391.077}{419.4}{if}%
\TextAt{8}{401.884}{419.4}{the}%
\TextAt{8}{422.124}{419.4}{zeroes}%
\TextAt{13}{457.057}{419.4}{\Character{28}}%
\TextAt{9}{463.476}{419.4}{\Character{40}}%
\TextAt{13}{468.028}{419.4}{\Character{11}}%
\TextAt{9}{475.55}{419.4}{\Character{41}}%
\TextAt{8}{484.083}{419.4}{and}%
\TextAt{13}{506.929}{419.4}{\Character{28}}%
\TextAt{3}{513.348}{423.738}{\Character{48}}%
\TextAt{9}{516.143}{419.4}{\Character{40}}%
\TextAt{13}{520.695}{419.4}{\Character{11}}%
\TextAt{9}{528.217}{419.4}{\Character{41}}%
\TextAt{8}{79.226}{402.354}{are}%
\TextAt{8}{98.14}{402.354}{conjugate}%
\TextAt{8}{150.552}{402.354}{over}%
\TextAt{13}{175.001}{402.354}{K}%
\TextAt{21}{184.961}{400.56}{v}%
\TextAt{8}{189.868}{402.354}{,}%
\TextAt{8}{196.54}{402.354}{i.e.,}%
\TextAt{8}{218.167}{402.354}{if}%
\TextAt{8}{228.294}{402.354}{they}%
\TextAt{8}{254.034}{402.354}{are}%
\TextAt{8}{272.948}{402.354}{zeroes}%
\TextAt{8}{307.201}{402.354}{of}%
\TextAt{8}{319.934}{402.354}{the}%
\TextAt{8}{339.493}{402.354}{same}%
\TextAt{8}{368.222}{402.354}{irreducible}%
\TextAt{8}{425.5}{402.354}{factor}%
\TextAt{13}{458.402}{402.354}{f}%
\TextAt{21}{464.172}{400.56}{i}%
\TextAt{8}{467.553}{402.354}{.}%
\TextAt{8}{475.802}{402.354}{With}%
\TextAt{8}{505.44}{402.354}{(8.1),}%
\TextAt{8}{79.226}{385.308}{this}%
\TextAt{8}{102.049}{385.308}{gives}%
\TextAt{8}{130.729}{385.308}{the}%
\TextAt{12}{79.226}{359.376}{Proposition}%
\TextAt{12}{153.87}{359.376}{(8.2).}%
\TextAt{8}{190.865}{359.376}{Suppose}%
\TextAt{8}{238.452}{359.376}{the}%
\TextAt{8}{260.289}{359.376}{extension}%
\TextAt{13}{313.734}{359.376}{L}%
\TextAt{16}{327.882}{359.376}{\Character{106}}%
\TextAt{13}{337.386}{359.376}{K}%
\TextAt{8}{353.781}{359.376}{is}%
\TextAt{8}{367.224}{359.376}{generated}%
\TextAt{8}{422.236}{359.376}{by}%
\TextAt{8}{440.175}{359.376}{the}%
\TextAt{8}{462.011}{359.376}{zero}%
\TextAt{13}{488.402}{359.376}{\Character{11}}%
\TextAt{8}{501.501}{359.376}{of}%
\TextAt{8}{516.51}{359.376}{the}%
\TextAt{8}{79.226}{342.33}{irreducible}%
\TextAt{8}{137.101}{342.33}{polynomial}%
\TextAt{13}{197.594}{342.33}{f}%
\TextAt{9}{204.641}{342.33}{\Character{40}}%
\TextAt{13}{209.193}{342.33}{X}%
\TextAt{9}{219.848}{342.33}{\Character{41}}%
\NativeMathAt{227.721}{342.33}{11.955}{\in}%
\TextAt{13}{239.012}{342.33}{K}%
\TextAt{9}{249.83}{342.33}{\Character{91}}%
\TextAt{13}{253.082}{342.33}{X}%
\TextAt{9}{263.737}{342.33}{\Character{93}}%
\TextAt{8}{266.988}{342.33}{.}%
\TextAt{8}{97.159}{322.296}{Then}%
\TextAt{8}{128.864}{322.296}{the}%
\TextAt{8}{150.167}{322.296}{valuations}%
\TextAt{13}{207.002}{322.296}{w}%
\TextAt{17}{215.37}{320.502}{1}%
\TextAt{13}{220.102}{322.296}{\Character{59}}%
\TextAt{13}{225.326}{322.296}{\Character{58}}%
\TextAt{13}{230.563}{322.296}{\Character{58}}%
\TextAt{13}{235.787}{322.296}{\Character{58}}%
\TextAt{13}{241.024}{322.296}{\Character{59}}%
\TextAt{13}{246.26}{322.296}{w}%
\TextAt{21}{254.691}{320.502}{r}%
\TextAt{8}{264.29}{322.296}{extending}%
\TextAt{13}{319.092}{322.296}{v}%
\TextAt{8}{330.225}{322.296}{to}%
\TextAt{13}{345.683}{322.296}{L}%
\TextAt{8}{358.692}{322.296}{correspond}%
\TextAt{8}{419.412}{322.296}{1–1}%
\TextAt{8}{442.031}{322.296}{to}%
\TextAt{8}{457.488}{322.296}{the}%
\TextAt{8}{478.792}{322.296}{irreducible}%
\TextAt{8}{79.226}{305.25}{factors}%
\TextAt{13}{117.339}{305.25}{f}%
\TextAt{17}{123.109}{303.456}{1}%
\TextAt{13}{127.842}{305.25}{\Character{59}}%
\TextAt{13}{133.066}{305.25}{\Character{58}}%
\TextAt{13}{138.303}{305.25}{\Character{58}}%
\TextAt{13}{143.527}{305.25}{\Character{58}}%
\TextAt{13}{148.764}{305.25}{\Character{59}}%
\TextAt{13}{154}{305.25}{f}%
\TextAt{21}{159.832}{303.456}{r}%
\TextAt{8}{168.284}{305.25}{in}%
\TextAt{8}{181.937}{305.25}{the}%
\TextAt{8}{202.094}{305.25}{decomposition}%
\TextAt{13}{233.439}{276.532}{f}%
\TextAt{9}{240.485}{276.532}{\Character{40}}%
\TextAt{13}{245.038}{276.532}{X}%
\TextAt{9}{255.693}{276.532}{\Character{41}}%
\TextAt{9}{263.548}{276.532}{\Character{61}}%
\TextAt{13}{275.991}{276.532}{f}%
\TextAt{17}{281.761}{274.739}{1}%
\TextAt{9}{286.494}{276.532}{\Character{40}}%
\TextAt{13}{291.046}{276.532}{X}%
\TextAt{9}{301.701}{276.532}{\Character{41}}%
\TextAt{21}{306.253}{281.468}{m}%
\TextAt{24}{313.744}{280.361}{1}%
\TextAt{16}{320.386}{276.532}{\Character{1}}%
\TextAt{16}{325.682}{276.532}{\Character{1}}%
\TextAt{16}{330.99}{276.532}{\Character{1}}%
\TextAt{13}{336.326}{276.532}{f}%
\TextAt{21}{342.096}{274.739}{r}%
\TextAt{9}{346.65}{276.532}{\Character{40}}%
\TextAt{13}{351.203}{276.532}{X}%
\TextAt{9}{361.858}{276.532}{\Character{41}}%
\TextAt{21}{366.41}{281.468}{m}%
\TextAt{22}{373.901}{280.472}{r}%
\TextAt{8}{79.226}{247.815}{of}%
\TextAt{13}{92.556}{247.815}{f}%
\TextAt{8}{103.5}{247.815}{over}%
\TextAt{8}{128.546}{247.815}{the}%
\TextAt{8}{148.702}{247.815}{completion}%
\TextAt{13}{208.538}{247.815}{K}%
\TextAt{21}{218.497}{246.021}{v}%
\TextAt{8}{223.405}{247.815}{.}%
\TextAt{8}{97.159}{221.883}{The}%
\TextAt{8}{122.624}{221.883}{extended}%
\TextAt{8}{173.781}{221.883}{valuation}%
\TextAt{13}{226.266}{221.883}{w}%
\TextAt{21}{234.633}{220.09}{i}%
\TextAt{8}{243.324}{221.883}{is}%
\TextAt{8}{256.499}{221.883}{explicitly}%
\TextAt{8}{308.637}{221.883}{obtained}%
\TextAt{8}{358.18}{221.883}{from}%
\TextAt{8}{387.232}{221.883}{the}%
\TextAt{8}{408.8}{221.883}{factor}%
\TextAt{13}{443.71}{221.883}{f}%
\TextAt{21}{449.48}{220.09}{i}%
\TextAt{8}{458.17}{221.883}{as}%
\TextAt{8}{473.952}{221.883}{follows:}%
\TextAt{8}{519.762}{221.883}{let}%
\TextAt{13}{79.226}{204.837}{\Character{11}}%
\TextAt{21}{86.672}{203.044}{i}%
\NativeMathAt{93.374}{204.837}{11.955}{\in}%
\Vector{q 0 G 0 g 0.478 w 32.665 -559.32 m 43.483 -559.32 l S Q}%
\TextAt{13}{104.665}{204.837}{K}%
\TextAt{21}{115.482}{203.044}{v}%
\TextAt{8}{124.287}{204.837}{be}%
\TextAt{8}{140.212}{204.837}{a}%
\TextAt{8}{149.967}{204.837}{zero}%
\TextAt{8}{174.679}{204.837}{of}%
\TextAt{13}{188.009}{204.837}{f}%
\TextAt{21}{193.779}{203.044}{i}%
\TextAt{8}{201.057}{204.837}{and}%
\TextAt{8}{223.82}{204.837}{let}%
\TextAt{13}{233.694}{176.12}{\Character{28}}%
\TextAt{21}{238.777}{174.327}{i}%
\TextAt{9}{245.479}{176.12}{\Character{58}}%
\TextAt{13}{252.052}{176.12}{L}%
\NativeArrowAt{263.337}{176.12}{11.955}{19.26}{\longrightarrow}%
\Vector{q 0 G 0 g 0.478 w 213.919 -588.037 m 224.737 -588.037 l S Q}%
\TextAt{13}{285.919}{176.12}{K}%
\TextAt{21}{296.737}{174.327}{v}%
\TextAt{13}{301.644}{176.12}{\Character{59}}%
\TextAt{13}{330.779}{176.12}{\Character{11}}%
\NativeArrowAt{341.641}{176.12}{11.955}{19.26}{\longmapsto}%
\TextAt{13}{364.224}{176.12}{\Character{11}}%
\TextAt{21}{371.669}{174.327}{i}%
\TextAt{13}{375.05}{176.12}{\Character{59}}%
\TextAt{8}{79.226}{147.402}{be}%
\TextAt{8}{95.15}{147.402}{the}%
\TextAt{8}{115.307}{147.402}{corresponding}%
\TextAt{13}{190.493}{147.402}{K}%
\TextAt{8}{201.311}{147.402}{-embedding}%
\TextAt{8}{264.386}{147.402}{of}%
\TextAt{13}{277.716}{147.402}{L}%
\TextAt{8}{289.578}{147.402}{into}%
\Vector{q 0 G 0 g 0.478 w 241.321 -616.755 m 252.139 -616.755 l S Q}%
\TextAt{13}{313.321}{147.402}{K}%
\TextAt{21}{324.139}{145.609}{v}%
\TextAt{8}{329.046}{147.402}{.}%
\TextAt{8}{337.495}{147.402}{Then}%
\TextAt{8}{368.052}{147.402}{one}%
\TextAt{8}{389.511}{147.402}{has}%
\TextAt{13}{277.702}{118.685}{w}%
\TextAt{21}{286.07}{116.892}{i}%
\TextAt{9}{292.772}{118.685}{\Character{61}}%
\TextAt{9}{305.624}{118.685}{\Character{22}}%
\TextAt{13}{305.198}{118.685}{v}%
\TextAt{16}{313.943}{118.685}{\Character{14}}%
\TextAt{13}{322.577}{118.685}{\Character{28}}%
\TextAt{21}{327.66}{116.892}{i}%
\TextAt{13}{331.042}{118.685}{\Character{58}}%
\TextAt{13}{79.226}{89.967}{\Character{28}}%
\TextAt{21}{84.31}{88.174}{i}%
\TextAt{8}{91.588}{89.967}{extends}%
\TextAt{8}{134.245}{89.967}{to}%
\TextAt{8}{148.555}{89.967}{an}%
\TextAt{8}{164.814}{89.967}{isomorphism}%
\TextAt{13}{260.662}{72.921}{\Character{28}}%
\TextAt{21}{265.745}{71.128}{i}%
\TextAt{9}{272.447}{72.921}{\Character{58}}%
\TextAt{13}{279.02}{72.921}{L}%
\TextAt{21}{286.984}{71.128}{w}%
\TextAt{22}{293.061}{69.913}{i}%
\NativeMathAt{302.629}{79.698}{7.97}{\sim}%
\NativeArrowAt{300.042}{72.921}{11.955}{13.175}{\longrightarrow}%
\TextAt{13}{316.536}{72.921}{K}%
\TextAt{21}{326.495}{71.128}{v}%
\TextAt{9}{331.402}{72.921}{\Character{40}}%
\TextAt{13}{335.955}{72.921}{\Character{11}}%
\TextAt{21}{343.4}{71.128}{i}%
\TextAt{9}{346.782}{72.921}{\Character{41}}%
\TextAt{8}{297.211}{43.033}{149}%
\Destination{78.23}{809.815}{nav.page.149}%
\Destination{79.36}{525.955}{nav.print.163}%
\Destination{79.36}{373.355}{nav.number.2.8.2}%
\Hotlink{505.44}{400.154}{28.2}{11.047}{nav.number.2.8.1}%
}

% Page 162 (body page 150)
\Page{%
\TextAt{2}{79.226}{796.943}{II.}%
\TextAt{2}{94.073}{796.943}{The}%
\TextAt{2}{116.492}{796.943}{Theory}%
\TextAt{2}{154.401}{796.943}{of}%
\TextAt{2}{166.826}{796.943}{Valuations}%
\Vector{q 0 G 0 g 0.299 w 7.226 18.729 m 460.769 18.729 l S Q}%
\TextAt{8}{79.226}{761.124}{on}%
\TextAt{8}{95.485}{761.124}{the}%
\TextAt{8}{115.642}{761.124}{completion}%
\TextAt{13}{175.477}{761.124}{L}%
\TextAt{21}{183.442}{759.33}{w}%
\TextAt{22}{189.519}{758.115}{i}%
\TextAt{8}{197.076}{761.124}{of}%
\TextAt{13}{210.406}{761.124}{L}%
\TextAt{8}{222.268}{761.124}{with}%
\TextAt{8}{248.928}{761.124}{respect}%
\TextAt{8}{288.978}{761.124}{to}%
\TextAt{13}{303.288}{761.124}{w}%
\TextAt{21}{311.656}{759.33}{i}%
\TextAt{8}{315.037}{761.124}{.}%
\TextAt{8}{97.159}{741.089}{Let}%
\TextAt{13}{119.373}{741.089}{L}%
\TextAt{16}{132.78}{741.089}{\Character{106}}%
\TextAt{13}{141.543}{741.089}{K}%
\TextAt{8}{157.502}{741.089}{be}%
\TextAt{8}{174.672}{741.089}{again}%
\TextAt{8}{207.143}{741.089}{an}%
\TextAt{8}{224.647}{741.089}{arbitrary}%
\TextAt{8}{275.661}{741.089}{finite}%
\TextAt{8}{306.818}{741.089}{extension.}%
\TextAt{8}{366.869}{741.089}{We}%
\TextAt{8}{388.258}{741.089}{will}%
\TextAt{8}{411.608}{741.089}{write}%
\TextAt{13}{442.765}{741.089}{w}%
\TextAt{16}{456.896}{741.089}{\Character{106}}%
\TextAt{13}{465.659}{741.089}{v}%
\TextAt{8}{476.89}{741.089}{to}%
\TextAt{8}{492.445}{741.089}{indicate}%
\TextAt{8}{79.226}{724.043}{that}%
\TextAt{13}{105.779}{724.043}{w}%
\TextAt{8}{119.55}{724.043}{is}%
\TextAt{8}{132.498}{724.043}{an}%
\TextAt{8}{149.94}{724.043}{extension}%
\TextAt{8}{202.89}{724.043}{of}%
\TextAt{8}{217.404}{724.043}{the}%
\TextAt{8}{238.744}{724.043}{valuation}%
\TextAt{13}{291.001}{724.043}{v}%
\TextAt{8}{302.17}{724.043}{of}%
\TextAt{13}{316.684}{724.043}{K}%
\TextAt{8}{332.583}{724.043}{to}%
\TextAt{13}{348.077}{724.043}{L}%
\TextAt{8}{356.042}{724.043}{.}%
\TextAt{8}{368.042}{724.043}{The}%
\TextAt{8}{393.28}{724.043}{inclusions}%
\TextAt{13}{447.915}{724.043}{L}%
\NativeArrowAt{461.209}{724.043}{11.955}{13.223}{\hookrightarrow}%
\TextAt{13}{479.771}{724.043}{L}%
\TextAt{21}{487.735}{722.25}{w}%
\TextAt{8}{499.606}{724.043}{induce}%
\TextAt{8}{79.226}{706.997}{homomorphisms}%
\TextAt{13}{166.511}{706.997}{L}%
\TextAt{16}{177.132}{706.997}{\Character{10}}%
\TextAt{21}{186.431}{705.204}{K}%
\TextAt{13}{197.29}{706.997}{K}%
\TextAt{21}{207.249}{705.204}{v}%
\NativeMathAt{215.477}{706.997}{11.955}{\rightarrow}%
\TextAt{13}{230.753}{706.997}{L}%
\TextAt{21}{238.718}{705.204}{w}%
\TextAt{8}{245.507}{706.997}{,}%
\TextAt{8}{252.657}{706.997}{via}%
\TextAt{13}{271.845}{706.997}{a}%
\TextAt{16}{280.646}{706.997}{\Character{10}}%
\TextAt{13}{292.601}{706.997}{b}%
\NativeArrowAt{300.899}{706.997}{11.955}{11.955}{\mapsto}%
\TextAt{13}{316.175}{706.997}{ab}%
\TextAt{8}{327.297}{706.997}{,}%
\TextAt{8}{334.447}{706.997}{and}%
\TextAt{8}{357.209}{706.997}{hence}%
\TextAt{8}{389.715}{706.997}{a}%
\TextAt{8}{399.471}{706.997}{canonical}%
\TextAt{8}{450.854}{706.997}{homomorphism}%
\TextAt{13}{243.806}{677.996}{\Character{39}}%
\TextAt{9}{254.795}{677.996}{\Character{58}}%
\TextAt{13}{261.367}{677.996}{L}%
\TextAt{16}{271.988}{677.996}{\Character{10}}%
\TextAt{21}{281.287}{676.203}{K}%
\TextAt{13}{292.146}{677.996}{K}%
\TextAt{21}{302.105}{676.203}{v}%
\NativeArrowAt{310.334}{677.996}{11.955}{19.26}{\longrightarrow}%
\TextAt{20}{332.916}{689.353}{\Character{89}}%
\TextAt{21}{334.027}{663.45}{w}%
\TextAt{3}{340.319}{663.45}{\Character{106}}%
\TextAt{21}{342.671}{663.45}{v}%
\TextAt{13}{350.184}{677.996}{L}%
\TextAt{21}{358.149}{676.203}{w}%
\TextAt{13}{364.938}{677.996}{\Character{58}}%
\TextAt{8}{79.226}{638.894}{To}%
\TextAt{8}{95.957}{638.894}{begin}%
\TextAt{8}{126.998}{638.894}{with,}%
\TextAt{8}{156.512}{638.894}{the}%
\TextAt{8}{176.172}{638.894}{tensor}%
\TextAt{8}{210.859}{638.894}{product}%
\TextAt{8}{254.262}{638.894}{is}%
\TextAt{8}{265.529}{638.894}{taken}%
\TextAt{8}{296.905}{638.894}{in}%
\TextAt{8}{310.061}{638.894}{the}%
\TextAt{8}{329.721}{638.894}{sense}%
\TextAt{8}{359.255}{638.894}{of}%
\TextAt{8}{372.089}{638.894}{vector}%
\TextAt{8}{406.716}{638.894}{spaces,}%
\TextAt{8}{445.46}{638.894}{i.e.,}%
\TextAt{8}{467.168}{638.894}{the}%
\TextAt{13}{486.828}{638.894}{K}%
\TextAt{8}{497.645}{638.894}{-vector}%
\TextAt{8}{79.226}{621.848}{space}%
\TextAt{13}{111.411}{621.848}{L}%
\TextAt{8}{124.184}{621.848}{is}%
\TextAt{8}{136.859}{621.848}{lifted}%
\TextAt{8}{168.004}{621.848}{to}%
\TextAt{8}{183.225}{621.848}{a}%
\TextAt{13}{193.891}{621.848}{K}%
\TextAt{21}{203.85}{620.055}{v}%
\TextAt{8}{208.757}{621.848}{-vector}%
\TextAt{8}{248.69}{621.848}{space}%
\TextAt{13}{280.875}{621.848}{L}%
\TextAt{16}{292.117}{621.848}{\Character{10}}%
\TextAt{21}{301.416}{620.055}{K}%
\TextAt{13}{312.896}{621.848}{K}%
\TextAt{21}{322.855}{620.055}{v}%
\TextAt{8}{327.762}{621.848}{.}%
\TextAt{8}{338.943}{621.848}{This}%
\TextAt{8}{366.573}{621.848}{latter,}%
\TextAt{8}{402.836}{621.848}{however,}%
\TextAt{8}{452.105}{621.848}{is}%
\TextAt{8}{464.78}{621.848}{in}%
\TextAt{8}{479.343}{621.848}{fact}%
\TextAt{8}{503.339}{621.848}{a}%
\TextAt{13}{514.005}{621.848}{K}%
\TextAt{21}{523.965}{620.055}{v}%
\TextAt{8}{528.872}{621.848}{-}%
\TextAt{8}{79.226}{604.802}{algebra,}%
\TextAt{8}{123.417}{604.802}{with}%
\TextAt{8}{150.022}{604.802}{the}%
\TextAt{8}{170.124}{604.802}{multiplication}%
\TextAt{9}{245.196}{604.802}{\Character{40}}%
\TextAt{13}{249.748}{604.802}{a}%
\TextAt{16}{258.438}{604.802}{\Character{10}}%
\TextAt{13}{270.282}{604.802}{b}%
\TextAt{9}{275.259}{604.802}{\Character{41}\Character{40}}%
\TextAt{13}{284.364}{604.802}{a}%
\TextAt{3}{290.509}{609.141}{\Character{48}}%
\TextAt{16}{295.849}{604.802}{\Character{10}}%
\TextAt{13}{307.692}{604.802}{b}%
\TextAt{3}{312.67}{609.141}{\Character{48}}%
\TextAt{9}{315.465}{604.802}{\Character{41}}%
\TextAt{9}{323.32}{604.802}{\Character{61}}%
\TextAt{13}{335.763}{604.802}{aa}%
\TextAt{3}{348.053}{609.141}{\Character{48}}%
\TextAt{16}{353.393}{604.802}{\Character{10}}%
\TextAt{13}{365.237}{604.802}{bb}%
\TextAt{3}{375.191}{609.141}{\Character{48}}%
\TextAt{8}{377.986}{604.802}{,}%
\TextAt{8}{385.092}{604.802}{and}%
\TextAt{13}{407.8}{604.802}{\Character{39}}%
\TextAt{8}{419.311}{604.802}{is}%
\TextAt{8}{431.021}{604.802}{a}%
\TextAt{8}{440.721}{604.802}{homomorphism}%
\TextAt{8}{523.337}{604.802}{of}%
\TextAt{13}{79.226}{587.756}{K}%
\TextAt{21}{89.185}{585.963}{v}%
\TextAt{8}{94.093}{587.756}{-algebras.}%
\TextAt{8}{148.138}{587.756}{This}%
\TextAt{8}{174.858}{587.756}{homomorphism}%
\TextAt{8}{257.528}{587.756}{is}%
\TextAt{8}{269.292}{587.756}{the}%
\TextAt{8}{289.448}{587.756}{subject}%
\TextAt{8}{329.821}{587.756}{of}%
\TextAt{8}{343.151}{587.756}{the}%
\TextAt{12}{79.226}{561.744}{Proposition}%
\TextAt{12}{151.937}{561.744}{(8.3).}%
\TextAt{8}{187.253}{561.744}{If}%
\TextAt{13}{198.945}{561.744}{L}%
\TextAt{16}{210.231}{561.744}{\Character{106}}%
\TextAt{13}{216.872}{561.744}{K}%
\TextAt{8}{231.587}{561.744}{is}%
\TextAt{8}{243.351}{561.744}{separable,}%
\TextAt{8}{298.046}{561.744}{then}%
\TextAt{13}{324.706}{561.744}{L}%
\TextAt{16}{335.327}{561.744}{\Character{10}}%
\TextAt{21}{344.626}{559.951}{K}%
\TextAt{13}{355.485}{561.744}{K}%
\TextAt{21}{365.444}{559.951}{v}%
\NativeMathAt{373.673}{561.744}{11.955}{\cong}%
\TextAt{20}{386.292}{570.711}{\Character{81}}%
\TextAt{21}{397.583}{558.257}{w}%
\TextAt{3}{403.874}{558.257}{\Character{106}}%
\TextAt{21}{406.227}{558.257}{v}%
\TextAt{13}{413.127}{561.744}{L}%
\TextAt{21}{421.091}{559.951}{w}%
\TextAt{8}{427.881}{561.744}{.}%
\TextAt{19}{79.226}{526.765}{Proof.}%
\TextAt{8}{116}{526.765}{Let}%
\TextAt{13}{136.129}{526.765}{\Character{11}}%
\TextAt{8}{146.708}{526.765}{be}%
\TextAt{8}{161.792}{526.765}{a}%
\TextAt{8}{170.707}{526.765}{primitive}%
\TextAt{8}{219.947}{526.765}{element}%
\TextAt{8}{262.348}{526.765}{for}%
\TextAt{13}{279.393}{526.765}{L}%
\TextAt{16}{290.678}{526.765}{\Character{106}}%
\TextAt{13}{297.32}{526.765}{K}%
\TextAt{8}{308.138}{526.765}{,}%
\TextAt{8}{314.615}{526.765}{so}%
\TextAt{8}{328.145}{526.765}{that}%
\TextAt{13}{352.673}{526.765}{L}%
\TextAt{9}{363.958}{526.765}{\Character{61}}%
\TextAt{13}{376.384}{526.765}{K}%
\TextAt{9}{387.202}{526.765}{\Character{40}}%
\TextAt{13}{391.754}{526.765}{\Character{11}}%
\TextAt{9}{399.276}{526.765}{\Character{41}}%
\TextAt{8}{403.828}{526.765}{,}%
\TextAt{8}{410.305}{526.765}{and}%
\TextAt{8}{432.227}{526.765}{let}%
\TextAt{13}{448.291}{526.765}{f}%
\TextAt{9}{455.338}{526.765}{\Character{40}}%
\TextAt{13}{459.89}{526.765}{X}%
\TextAt{9}{470.545}{526.765}{\Character{41}}%
\NativeMathAt{478.418}{526.765}{11.955}{\in}%
\TextAt{13}{489.709}{526.765}{K}%
\TextAt{9}{500.527}{526.765}{\Character{91}}%
\TextAt{13}{503.779}{526.765}{X}%
\TextAt{9}{514.434}{526.765}{\Character{93}}%
\TextAt{8}{520.743}{526.765}{be}%
\TextAt{8}{79.226}{509.719}{its}%
\TextAt{8}{95.616}{509.719}{minimal}%
\TextAt{8}{141.212}{509.719}{polynomial.}%
\TextAt{8}{206.47}{509.719}{To}%
\TextAt{8}{223.768}{509.719}{every}%
\TextAt{13}{254.731}{509.719}{w}%
\TextAt{16}{266.863}{509.719}{\Character{106}}%
\TextAt{13}{273.625}{509.719}{v}%
\TextAt{8}{279.714}{509.719}{,}%
\TextAt{8}{286.952}{509.719}{there}%
\TextAt{8}{316.934}{509.719}{corresponds}%
\TextAt{8}{381.193}{509.719}{an}%
\TextAt{8}{397.523}{509.719}{irreducible}%
\TextAt{8}{455.469}{509.719}{factor}%
\TextAt{13}{489.038}{509.719}{f}%
\TextAt{21}{494.808}{507.926}{w}%
\TextAt{9}{501.598}{509.719}{\Character{40}}%
\TextAt{13}{506.15}{509.719}{X}%
\TextAt{9}{516.805}{509.719}{\Character{41}}%
\NativeMathAt{524.799}{509.719}{11.955}{\in}%
\TextAt{13}{79.226}{492.673}{K}%
\TextAt{21}{89.185}{490.88}{v}%
\TextAt{9}{94.093}{492.673}{\Character{91}}%
\TextAt{13}{97.344}{492.673}{X}%
\TextAt{9}{108}{492.673}{\Character{93}}%
\TextAt{8}{115.635}{492.673}{of}%
\TextAt{13}{129.451}{492.673}{f}%
\TextAt{9}{136.497}{492.673}{\Character{40}}%
\TextAt{13}{141.049}{492.673}{X}%
\TextAt{9}{151.705}{492.673}{\Character{41}}%
\TextAt{8}{156.257}{492.673}{,}%
\TextAt{8}{164.014}{492.673}{and}%
\TextAt{8}{187.262}{492.673}{in}%
\TextAt{8}{201.401}{492.673}{view}%
\TextAt{8}{228.87}{492.673}{of}%
\TextAt{8}{242.686}{492.673}{the}%
\TextAt{8}{263.328}{492.673}{separability,}%
\TextAt{8}{329.367}{492.673}{we}%
\TextAt{8}{347.08}{492.673}{have}%
\TextAt{13}{374.561}{492.673}{f}%
\TextAt{9}{381.607}{492.673}{\Character{40}}%
\TextAt{13}{386.16}{492.673}{X}%
\TextAt{9}{396.815}{492.673}{\Character{41}}%
\TextAt{9}{405.506}{492.673}{\Character{61}}%
\TextAt{20}{418.77}{501.64}{\Character{81}}%
\TextAt{21}{430.061}{489.186}{w}%
\TextAt{3}{436.352}{489.186}{\Character{106}}%
\TextAt{21}{438.704}{489.186}{v}%
\TextAt{13}{445.604}{492.673}{f}%
\TextAt{21}{451.374}{490.88}{w}%
\TextAt{9}{458.164}{492.673}{\Character{40}}%
\TextAt{13}{462.716}{492.673}{X}%
\TextAt{9}{473.371}{492.673}{\Character{41}}%
\TextAt{8}{477.924}{492.673}{.}%
\TextAt{8}{487.83}{492.673}{Consider}%
\TextAt{8}{79.226}{475.627}{all}%
\TextAt{8}{95.7}{475.627}{the}%
\TextAt{13}{116.071}{475.627}{L}%
\TextAt{21}{124.036}{473.834}{w}%
\TextAt{8}{134.937}{475.627}{as}%
\TextAt{8}{149.522}{475.627}{embedded}%
\TextAt{8}{205.005}{475.627}{into}%
\TextAt{8}{228.963}{475.627}{an}%
\TextAt{8}{245.437}{475.627}{algebraic}%
\TextAt{8}{295.086}{475.627}{closure}%
\Vector{q 0 G 0 g 0.478 w 262.382 -288.529 m 273.2 -288.529 l S Q}%
\TextAt{13}{334.382}{475.627}{K}%
\TextAt{21}{345.2}{473.834}{v}%
\TextAt{8}{354.22}{475.627}{of}%
\TextAt{13}{367.764}{475.627}{K}%
\TextAt{21}{377.723}{473.834}{v}%
\TextAt{8}{386.743}{475.627}{and}%
\TextAt{8}{409.72}{475.627}{denote}%
\TextAt{8}{447.654}{475.627}{by}%
\TextAt{13}{464.127}{475.627}{\Character{11}}%
\TextAt{21}{471.573}{473.834}{w}%
\TextAt{8}{482.475}{475.627}{the}%
\TextAt{8}{502.846}{475.627}{image}%
\TextAt{8}{79.226}{458.582}{of}%
\TextAt{13}{92.913}{458.582}{\Character{11}}%
\TextAt{8}{104.689}{458.582}{under}%
\TextAt{13}{138.21}{458.582}{L}%
\NativeArrowAt{150.094}{458.582}{11.955}{13.225}{\hookrightarrow}%
\TextAt{13}{167.248}{458.582}{L}%
\TextAt{21}{175.212}{456.788}{w}%
\TextAt{8}{182.002}{458.582}{.}%
\TextAt{8}{191.521}{458.582}{Then}%
\TextAt{8}{222.436}{458.582}{we}%
\TextAt{8}{240.02}{458.582}{find}%
\TextAt{13}{263.785}{458.582}{L}%
\TextAt{21}{271.75}{456.788}{w}%
\TextAt{9}{282.469}{458.582}{\Character{61}}%
\TextAt{13}{295.503}{458.582}{K}%
\TextAt{21}{305.462}{456.788}{v}%
\TextAt{9}{310.369}{458.582}{\Character{40}}%
\TextAt{13}{314.922}{458.582}{\Character{11}}%
\TextAt{21}{322.367}{456.788}{w}%
\TextAt{9}{329.157}{458.582}{\Character{41}}%
\TextAt{8}{337.964}{458.582}{and}%
\TextAt{13}{361.083}{458.582}{f}%
\TextAt{21}{366.853}{456.788}{w}%
\TextAt{9}{373.643}{458.582}{\Character{40}}%
\TextAt{13}{378.195}{458.582}{X}%
\TextAt{9}{388.85}{458.582}{\Character{41}}%
\TextAt{8}{397.657}{458.582}{is}%
\TextAt{8}{409.778}{458.582}{the}%
\TextAt{8}{430.291}{458.582}{minimal}%
\TextAt{8}{476.174}{458.582}{polynomial}%
\TextAt{8}{79.226}{441.536}{of}%
\TextAt{13}{92.556}{441.536}{\Character{11}}%
\TextAt{21}{100.002}{439.742}{w}%
\TextAt{8}{110.689}{441.536}{over}%
\TextAt{13}{135.735}{441.536}{K}%
\TextAt{21}{145.694}{439.742}{v}%
\TextAt{8}{150.601}{441.536}{.}%
\TextAt{8}{159.05}{441.536}{We}%
\TextAt{8}{179.194}{441.536}{now}%
\TextAt{8}{203.583}{441.536}{get}%
\TextAt{8}{223.094}{441.536}{a}%
\TextAt{8}{232.849}{441.536}{commutative}%
\TextAt{8}{302.775}{441.536}{diagram}%
\TextAt{13}{212.278}{409.2}{K}%
\TextAt{21}{222.237}{407.407}{v}%
\TextAt{9}{227.145}{409.2}{\Character{91}}%
\TextAt{13}{230.397}{409.2}{X}%
\TextAt{9}{241.052}{409.2}{\Character{93}\Character{47}\Character{40}}%
\TextAt{13}{254.709}{409.2}{f}%
\TextAt{9}{261.755}{409.2}{\Character{41}}%
\TextAt{20}{319.651}{420.557}{\Character{89}}%
\TextAt{21}{320.763}{394.654}{w}%
\TextAt{3}{327.054}{394.654}{\Character{106}}%
\TextAt{21}{329.406}{394.654}{v}%
\TextAt{13}{336.92}{409.2}{K}%
\TextAt{21}{346.879}{407.407}{v}%
\TextAt{9}{351.787}{409.2}{\Character{91}}%
\TextAt{13}{355.038}{409.2}{X}%
\TextAt{9}{365.693}{409.2}{\Character{93}\Character{47}\Character{40}}%
\TextAt{13}{379.35}{409.2}{f}%
\TextAt{21}{385.12}{407.407}{w}%
\TextAt{9}{391.91}{409.2}{\Character{41}}%
\TextAt{13}{216.47}{337.875}{L}%
\TextAt{16}{227.091}{337.875}{\Character{10}}%
\TextAt{21}{236.39}{336.081}{K}%
\TextAt{13}{247.249}{337.875}{K}%
\TextAt{21}{257.208}{336.081}{v}%
\TextAt{20}{342.046}{349.232}{\Character{89}}%
\TextAt{21}{343.158}{323.329}{w}%
\TextAt{3}{349.449}{323.329}{\Character{106}}%
\TextAt{21}{351.801}{323.329}{v}%
\TextAt{13}{359.315}{337.875}{L}%
\TextAt{21}{367.279}{336.081}{w}%
\NativeDiagram{%
\path[draw=black,line width=0.478bp] (199.661,-361.81) -- (241.823,-361.81);%
\path[draw=black,line width=0.478bp,line cap=round,line join=round] (239.576,-358.941) .. controls (240.03,-360.662) and (241.042,-361.475) .. (242.062,-361.81) .. controls (241.042,-362.145) and (240.03,-362.957) .. (239.576,-364.679);%
\path[draw=black,line width=0.478bp] (167.294,-372.367) -- (167.294,-422.898);%
\path[draw=black,line width=0.478bp,line cap=round,line join=round] (170.163,-420.65) .. controls (168.441,-421.105) and (167.629,-422.117) .. (167.294,-423.137) .. controls (166.959,-422.117) and (166.146,-421.105) .. (164.425,-420.65);%
\path[draw=black,line width=0.478bp] (286.059,-387.112) -- (286.059,-418.514);%
\path[draw=black,line width=0.478bp,line cap=round,line join=round] (288.928,-416.267) .. controls (287.207,-416.721) and (286.394,-417.733) .. (286.059,-418.753) .. controls (285.725,-417.733) and (284.912,-416.721) .. (283.19,-416.267);%
\path[draw=black,line width=0.478bp] (195.469,-433.136) -- (264.218,-433.136);%
\path[draw=black,line width=0.478bp,line cap=round,line join=round] (261.97,-430.267) .. controls (262.424,-431.988) and (263.437,-432.801) .. (264.457,-433.136) .. controls (263.437,-433.471) and (262.424,-434.284) .. (261.97,-436.005);%
}%
\TextAt{13}{401.617}{370.948}{\Character{59}}%
\TextAt{8}{79.226}{294.394}{where}%
\TextAt{8}{112.646}{294.394}{the}%
\TextAt{8}{132.412}{294.394}{top}%
\TextAt{8}{152.837}{294.394}{arrow}%
\TextAt{8}{185.3}{294.394}{is}%
\TextAt{8}{196.674}{294.394}{an}%
\TextAt{8}{212.543}{294.394}{isomorphism}%
\TextAt{8}{280.573}{294.394}{by}%
\TextAt{8}{296.442}{294.394}{the}%
\TextAt{8}{316.209}{294.394}{Chinese}%
\TextAt{8}{359.443}{294.394}{remainder}%
\TextAt{8}{414.334}{294.394}{theorem.}%
\TextAt{8}{464.281}{294.394}{The}%
\TextAt{8}{487.945}{294.394}{arrow}%
\TextAt{8}{520.408}{294.394}{on}%
\TextAt{8}{79.226}{277.348}{the}%
\TextAt{8}{99.952}{277.348}{left}%
\TextAt{8}{121}{277.348}{is}%
\TextAt{8}{133.333}{277.348}{induced}%
\TextAt{8}{177.467}{277.348}{by}%
\TextAt{13}{194.295}{277.348}{X}%
\NativeArrowAt{209.241}{277.348}{11.955}{11.955}{\mapsto}%
\TextAt{13}{225.488}{277.348}{\Character{11}}%
\TextAt{16}{236.054}{277.348}{\Character{10}}%
\TextAt{9}{248.397}{277.348}{1}%
\TextAt{8}{258.717}{277.348}{and}%
\TextAt{8}{282.049}{277.348}{is}%
\TextAt{8}{294.382}{277.348}{an}%
\TextAt{8}{311.21}{277.348}{isomorphism}%
\TextAt{8}{380.199}{277.348}{because}%
\TextAt{13}{424.07}{277.348}{K}%
\TextAt{9}{434.887}{277.348}{\Character{91}}%
\TextAt{13}{438.139}{277.348}{X}%
\TextAt{9}{448.794}{277.348}{\Character{93}\Character{47}\Character{40}}%
\TextAt{13}{462.451}{277.348}{f}%
\TextAt{9}{469.497}{277.348}{\Character{41}}%
\NativeMathAt{478.341}{277.348}{11.955}{\cong}%
\TextAt{13}{491.93}{277.348}{K}%
\TextAt{9}{502.748}{277.348}{\Character{40}}%
\TextAt{13}{507.3}{277.348}{\Character{11}}%
\TextAt{9}{514.822}{277.348}{\Character{41}}%
\TextAt{9}{523.645}{277.348}{\Character{61}}%
\TextAt{13}{79.226}{260.302}{L}%
\TextAt{8}{87.191}{260.302}{.}%
\TextAt{8}{101.502}{260.302}{The}%
\TextAt{8}{127.51}{260.302}{arrow}%
\TextAt{8}{162.317}{260.302}{on}%
\TextAt{8}{180.53}{260.302}{the}%
\TextAt{8}{202.641}{260.302}{right}%
\TextAt{8}{232.893}{260.302}{is}%
\TextAt{8}{246.611}{260.302}{induced}%
\TextAt{8}{292.13}{260.302}{by}%
\TextAt{13}{310.344}{260.302}{X}%
\NativeArrowAt{327.65}{260.302}{11.955}{11.955}{\mapsto}%
\TextAt{13}{346.256}{260.302}{\Character{11}}%
\TextAt{21}{353.702}{258.509}{w}%
\TextAt{8}{366.343}{260.302}{and}%
\TextAt{8}{391.06}{260.302}{is}%
\TextAt{8}{404.778}{260.302}{an}%
\TextAt{8}{422.992}{260.302}{isomorphism}%
\TextAt{8}{493.365}{260.302}{because}%
\TextAt{13}{79.226}{243.256}{K}%
\TextAt{21}{89.185}{241.463}{v}%
\TextAt{9}{94.093}{243.256}{\Character{91}}%
\TextAt{13}{97.344}{243.256}{X}%
\TextAt{9}{108}{243.256}{\Character{93}\Character{47}\Character{40}}%
\TextAt{13}{121.657}{243.256}{f}%
\TextAt{21}{127.427}{241.463}{w}%
\TextAt{9}{134.216}{243.256}{\Character{41}}%
\NativeMathAt{142.089}{243.256}{11.955}{\cong}%
\TextAt{13}{154.708}{243.256}{K}%
\TextAt{21}{164.668}{241.463}{v}%
\TextAt{9}{169.575}{243.256}{\Character{40}}%
\TextAt{13}{174.128}{243.256}{\Character{11}}%
\TextAt{21}{181.573}{241.463}{w}%
\TextAt{9}{188.363}{243.256}{\Character{41}}%
\TextAt{9}{196.218}{243.256}{\Character{61}}%
\TextAt{13}{208.661}{243.256}{L}%
\TextAt{21}{216.626}{241.463}{w}%
\TextAt{8}{223.415}{243.256}{.}%
\TextAt{8}{231.863}{243.256}{Hence}%
\TextAt{8}{266.641}{243.256}{the}%
\TextAt{8}{286.797}{243.256}{bottom}%
\TextAt{8}{328.103}{243.256}{arrow}%
\TextAt{8}{360.955}{243.256}{is}%
\TextAt{8}{372.719}{243.256}{an}%
\TextAt{8}{388.978}{243.256}{isomorphism}%
\TextAt{8}{457.398}{243.256}{as}%
\TextAt{8}{471.768}{243.256}{well.}%
\Vector{q 0 G 0 g 0.399 w 452.334 -530.744 m 452.334 -522.674 l S Q}%
\Vector{q 0 G 0 g 0.399 w 452.534 -522.873 m 459.707 -522.873 l S Q}%
\Vector{q 0 G 0 g 0.399 w 452.534 -530.544 m 459.707 -530.544 l S Q}%
\Vector{q 0 G 0 g 0.399 w 459.906 -530.744 m 459.906 -522.674 l S Q}%
\TextAt{12}{79.226}{208.278}{Corollary}%
\TextAt{12}{139.014}{208.278}{(8.4).}%
\TextAt{8}{174.33}{208.278}{If}%
\TextAt{13}{186.022}{208.278}{L}%
\TextAt{16}{197.307}{208.278}{\Character{106}}%
\TextAt{13}{203.949}{208.278}{K}%
\TextAt{8}{218.664}{208.278}{is}%
\TextAt{8}{230.428}{208.278}{separable,}%
\TextAt{8}{285.123}{208.278}{then}%
\TextAt{8}{311.783}{208.278}{one}%
\TextAt{8}{333.242}{208.278}{has}%
\TextAt{9}{248.892}{179.276}{\Character{91}}%
\TextAt{13}{252.144}{179.276}{L}%
\TextAt{9}{263.429}{179.276}{\Character{58}}%
\TextAt{13}{270.002}{179.276}{K}%
\TextAt{9}{280.819}{179.276}{\Character{93}}%
\TextAt{9}{287.37}{179.276}{\Character{61}}%
\TextAt{20}{299.817}{190.634}{\Character{88}}%
\TextAt{21}{301.925}{164.731}{w}%
\TextAt{3}{308.217}{164.731}{\Character{106}}%
\TextAt{21}{310.569}{164.731}{v}%
\TextAt{9}{317.086}{179.276}{\Character{91}}%
\TextAt{13}{320.338}{179.276}{L}%
\TextAt{21}{328.302}{177.483}{w}%
\TextAt{9}{338.412}{179.276}{\Character{58}}%
\TextAt{13}{344.985}{179.276}{K}%
\TextAt{21}{354.944}{177.483}{v}%
\TextAt{9}{359.852}{179.276}{\Character{93}}%
\TextAt{8}{79.226}{140.295}{and}%
\TextAt{13}{157.638}{123.249}{N}%
\TextAt{21}{167.012}{121.389}{L}%
\TextAt{3}{172.771}{121.389}{\Character{106}}%
\TextAt{21}{175.123}{121.389}{K}%
\TextAt{9}{183.326}{123.249}{\Character{40}}%
\TextAt{13}{187.878}{123.249}{\Character{11}}%
\TextAt{9}{195.4}{123.249}{\Character{41}}%
\TextAt{9}{203.255}{123.249}{\Character{61}}%
\TextAt{20}{215.699}{134.606}{\Character{89}}%
\TextAt{21}{216.81}{108.703}{w}%
\TextAt{3}{223.101}{108.703}{\Character{106}}%
\TextAt{21}{225.454}{108.703}{v}%
\TextAt{13}{232.967}{123.249}{N}%
\TextAt{21}{242.341}{121.389}{L}%
\TextAt{22}{248.1}{120.393}{w}%
\TextAt{3}{254.098}{121.389}{\Character{106}}%
\TextAt{21}{256.451}{121.389}{K}%
\TextAt{22}{263.592}{120.393}{v}%
\TextAt{9}{268.539}{123.249}{\Character{40}}%
\TextAt{13}{273.091}{123.249}{\Character{11}}%
\TextAt{9}{280.613}{123.249}{\Character{41}}%
\TextAt{13}{285.165}{123.249}{\Character{59}}%
\TextAt{8}{314.32}{123.249}{Tr}%
\TextAt{21}{327.327}{121.389}{L}%
\TextAt{3}{333.086}{121.389}{\Character{106}}%
\TextAt{21}{335.438}{121.389}{K}%
\TextAt{9}{343.641}{123.249}{\Character{40}}%
\TextAt{13}{348.193}{123.249}{\Character{11}}%
\TextAt{9}{355.715}{123.249}{\Character{41}}%
\TextAt{9}{363.57}{123.249}{\Character{61}}%
\TextAt{20}{376.013}{134.606}{\Character{88}}%
\TextAt{21}{378.121}{108.703}{w}%
\TextAt{3}{384.413}{108.703}{\Character{106}}%
\TextAt{21}{386.765}{108.703}{v}%
\TextAt{8}{395.275}{123.249}{Tr}%
\TextAt{21}{408.282}{121.389}{L}%
\TextAt{22}{414.041}{120.393}{w}%
\TextAt{3}{420.039}{121.389}{\Character{106}}%
\TextAt{21}{422.391}{121.389}{K}%
\TextAt{22}{429.533}{120.393}{v}%
\TextAt{9}{434.48}{123.249}{\Character{40}}%
\TextAt{13}{439.032}{123.249}{\Character{11}}%
\TextAt{9}{446.554}{123.249}{\Character{41}}%
\TextAt{13}{451.106}{123.249}{\Character{58}}%
\TextAt{8}{297.211}{43.033}{150}%
\Destination{78.23}{809.815}{nav.page.150}%
\Destination{79.36}{541.755}{nav.print.164}%
\Destination{79.36}{575.755}{nav.number.2.8.3}%
\Destination{79.36}{222.255}{nav.number.2.8.4}%
}

% Page 163 (body page 151)
\Page{%
\TextAt{2}{79.226}{796.943}{II.}%
\TextAt{2}{94.073}{796.943}{The}%
\TextAt{2}{116.492}{796.943}{Theory}%
\TextAt{2}{154.401}{796.943}{of}%
\TextAt{2}{166.826}{796.943}{Valuations}%
\Vector{q 0 G 0 g 0.299 w 7.226 18.729 m 460.769 18.729 l S Q}%
\TextAt{19}{79.226}{761.124}{Proof.}%
\TextAt{8}{116}{761.124}{The}%
\TextAt{8}{139.896}{761.124}{first}%
\TextAt{8}{163.864}{761.124}{equation}%
\TextAt{8}{211.515}{761.124}{results}%
\TextAt{8}{248.549}{761.124}{from}%
\TextAt{8}{276.032}{761.124}{(8.3)}%
\TextAt{8}{303.849}{761.124}{since}%
\TextAt{9}{332.36}{761.124}{\Character{91}}%
\TextAt{13}{335.612}{761.124}{L}%
\TextAt{9}{346.897}{761.124}{\Character{58}}%
\TextAt{13}{353.469}{761.124}{K}%
\TextAt{9}{364.287}{761.124}{\Character{93}}%
\TextAt{9}{370.838}{761.124}{\Character{61}}%
\TextAt{8}{383.285}{761.124}{dim}%
\TextAt{21}{402.796}{759.33}{K}%
\TextAt{9}{410.999}{761.124}{\Character{40}}%
\TextAt{13}{415.551}{761.124}{L}%
\TextAt{9}{423.515}{761.124}{\Character{41}}%
\TextAt{9}{431.37}{761.124}{\Character{61}}%
\TextAt{8}{443.814}{761.124}{dim}%
\TextAt{21}{463.325}{759.33}{K}%
\TextAt{22}{470.466}{758.334}{v}%
\TextAt{9}{475.413}{761.124}{\Character{40}}%
\TextAt{13}{479.965}{761.124}{L}%
\TextAt{16}{490.264}{761.124}{\Character{10}}%
\TextAt{21}{499.562}{759.33}{K}%
\TextAt{13}{510.099}{761.124}{K}%
\TextAt{21}{520.058}{759.33}{v}%
\TextAt{9}{524.965}{761.124}{\Character{41}}%
\TextAt{8}{529.518}{761.124}{.}%
\TextAt{8}{79.226}{744.078}{On}%
\TextAt{8}{98.737}{744.078}{both}%
\TextAt{8}{126.377}{744.078}{sides}%
\TextAt{8}{154.46}{744.078}{of}%
\TextAt{8}{167.79}{744.078}{the}%
\TextAt{8}{187.946}{744.078}{isomorphism}%
\TextAt{13}{259.194}{720.513}{L}%
\TextAt{16}{269.815}{720.513}{\Character{10}}%
\TextAt{21}{279.113}{718.719}{K}%
\TextAt{13}{289.973}{720.513}{K}%
\TextAt{21}{299.932}{718.719}{v}%
\NativeMathAt{308.16}{720.512}{11.955}{\cong}%
\TextAt{20}{320.779}{731.87}{\Character{89}}%
\TextAt{21}{321.891}{705.967}{w}%
\TextAt{3}{328.182}{705.967}{\Character{106}}%
\TextAt{21}{330.535}{705.967}{v}%
\TextAt{13}{338.048}{720.513}{L}%
\TextAt{21}{346.013}{718.719}{w}%
\TextAt{8}{79.226}{686.847}{let}%
\TextAt{8}{96.06}{686.847}{us}%
\TextAt{8}{111.004}{686.847}{consider}%
\TextAt{8}{156.518}{686.847}{the}%
\TextAt{8}{176.604}{686.847}{endomorphism:}%
\TextAt{8}{259.88}{686.847}{multiplication}%
\TextAt{8}{334.935}{686.847}{by}%
\TextAt{13}{351.123}{686.847}{\Character{11}}%
\TextAt{8}{358.644}{686.847}{.}%
\TextAt{8}{367.069}{686.847}{The}%
\TextAt{8}{391.052}{686.847}{characteristic}%
\TextAt{8}{462.915}{686.847}{polynomial}%
\TextAt{8}{523.337}{686.847}{of}%
\TextAt{13}{79.226}{669.801}{\Character{11}}%
\TextAt{8}{89.975}{669.801}{on}%
\TextAt{8}{105.563}{669.801}{the}%
\TextAt{13}{125.049}{669.801}{K}%
\TextAt{21}{135.008}{668.008}{v}%
\TextAt{8}{139.916}{669.801}{-vector}%
\TextAt{8}{178.267}{669.801}{space}%
\TextAt{13}{208.871}{669.801}{L}%
\TextAt{16}{218.121}{669.801}{\Character{10}}%
\TextAt{21}{227.419}{668.008}{K}%
\TextAt{13}{236.907}{669.801}{K}%
\TextAt{21}{246.866}{668.008}{v}%
\TextAt{8}{255.001}{669.801}{is}%
\TextAt{8}{266.094}{669.801}{the}%
\TextAt{8}{285.58}{669.801}{same}%
\TextAt{8}{314.235}{669.801}{as}%
\TextAt{8}{327.935}{669.801}{that}%
\TextAt{8}{352.633}{669.801}{on}%
\TextAt{8}{368.222}{669.801}{the}%
\TextAt{13}{387.707}{669.801}{K}%
\TextAt{8}{398.525}{669.801}{-vector}%
\TextAt{8}{436.876}{669.801}{space}%
\TextAt{13}{467.48}{669.801}{L}%
\TextAt{8}{475.445}{669.801}{.}%
\TextAt{8}{483.67}{669.801}{Therefore}%
\TextAt{8}{164.443}{646.236}{char.}%
\TextAt{8}{193.387}{646.236}{polynomial}%
\TextAt{21}{249.982}{643.299}{L}%
\TextAt{3}{255.741}{643.299}{\Character{106}}%
\TextAt{21}{258.094}{643.299}{K}%
\TextAt{9}{266.296}{646.236}{\Character{40}}%
\TextAt{13}{270.849}{646.236}{\Character{11}}%
\TextAt{9}{278.37}{646.236}{\Character{41}}%
\TextAt{9}{286.225}{646.236}{\Character{61}}%
\TextAt{20}{298.669}{657.594}{\Character{89}}%
\TextAt{21}{299.781}{631.691}{w}%
\TextAt{3}{306.072}{631.691}{\Character{106}}%
\TextAt{21}{308.424}{631.691}{v}%
\TextAt{8}{315.938}{646.236}{char.}%
\TextAt{8}{344.881}{646.236}{polynomial}%
\TextAt{21}{401.477}{643.299}{L}%
\TextAt{22}{407.236}{642.303}{w}%
\TextAt{3}{413.234}{643.299}{\Character{106}}%
\TextAt{21}{415.586}{643.299}{K}%
\TextAt{22}{422.728}{642.303}{v}%
\TextAt{9}{427.675}{646.236}{\Character{40}}%
\TextAt{13}{432.227}{646.236}{\Character{11}}%
\TextAt{9}{439.749}{646.236}{\Character{41}}%
\TextAt{13}{444.301}{646.236}{\Character{58}}%
\TextAt{8}{79.226}{612.571}{This}%
\TextAt{8}{105.946}{612.571}{implies}%
\TextAt{8}{145.673}{612.571}{immediately}%
\TextAt{8}{212.335}{612.571}{the}%
\TextAt{8}{232.491}{612.571}{identities}%
\TextAt{8}{282.954}{612.571}{for}%
\TextAt{8}{300.839}{612.571}{the}%
\TextAt{8}{320.996}{612.571}{norm}%
\TextAt{8}{351.565}{612.571}{and}%
\TextAt{8}{374.328}{612.571}{the}%
\TextAt{8}{394.484}{612.571}{trace.}%
\Vector{q 0 G 0 g 0.399 w 452.334 -161.429 m 452.334 -153.359 l S Q}%
\Vector{q 0 G 0 g 0.399 w 452.534 -153.559 m 459.707 -153.559 l S Q}%
\Vector{q 0 G 0 g 0.399 w 452.534 -161.23 m 459.707 -161.23 l S Q}%
\Vector{q 0 G 0 g 0.399 w 459.906 -161.429 m 459.906 -153.359 l S Q}%
\TextAt{8}{97.159}{585.123}{If}%
\TextAt{13}{110.686}{585.123}{v}%
\TextAt{8}{122.506}{585.123}{is}%
\TextAt{8}{136.105}{585.123}{a}%
\TextAt{8}{147.695}{585.123}{nonarchimedean}%
\TextAt{8}{236.36}{585.123}{valuation,}%
\TextAt{8}{292.978}{585.123}{then}%
\TextAt{8}{321.473}{585.123}{we}%
\TextAt{8}{340.535}{585.123}{define,}%
\TextAt{8}{379.89}{585.123}{as}%
\TextAt{8}{396.095}{585.123}{in}%
\TextAt{8}{411.582}{585.123}{the}%
\TextAt{8}{433.574}{585.123}{henselian}%
\TextAt{8}{486.194}{585.123}{case,}%
\TextAt{8}{516.51}{585.123}{the}%
\TextAt{12}{79.226}{568.077}{ramification}%
\TextAt{12}{153.874}{568.077}{index}%
\TextAt{8}{189.692}{568.077}{of}%
\TextAt{8}{203.022}{568.077}{an}%
\TextAt{8}{219.281}{568.077}{extension}%
\TextAt{13}{271.047}{568.077}{w}%
\TextAt{16}{283.057}{568.077}{\Character{106}}%
\TextAt{13}{289.699}{568.077}{v}%
\TextAt{8}{299.685}{568.077}{by}%
\TextAt{13}{251.901}{544.512}{e}%
\TextAt{21}{257.327}{542.718}{w}%
\TextAt{9}{267.437}{544.512}{\Character{61}}%
\TextAt{9}{279.846}{544.512}{\Character{40}}%
\TextAt{13}{284.415}{544.512}{w}%
\TextAt{9}{293.104}{544.512}{\Character{40}}%
\TextAt{13}{297.657}{544.512}{L}%
\TextAt{3}{305.621}{549.448}{\Character{3}}%
\TextAt{9}{310.353}{544.512}{\Character{41}}%
\TextAt{9}{318.208}{544.512}{\Character{58}}%
\TextAt{13}{324.799}{544.512}{v}%
\TextAt{9}{330.887}{544.512}{\Character{40}}%
\TextAt{13}{335.44}{544.512}{K}%
\TextAt{3}{346.257}{549.448}{\Character{3}}%
\TextAt{9}{350.99}{544.512}{\Character{41}\Character{41}}%
\TextAt{8}{79.226}{520.947}{and}%
\TextAt{8}{101.989}{520.947}{the}%
\TextAt{12}{122.145}{520.947}{inertia}%
\TextAt{12}{164.993}{520.947}{degree}%
\TextAt{8}{206.991}{520.947}{by}%
\TextAt{13}{271.843}{503.901}{f}%
\TextAt{21}{277.613}{502.107}{w}%
\TextAt{9}{287.723}{503.901}{\Character{61}}%
\TextAt{9}{300.132}{503.901}{\Character{91}}%
\TextAt{13}{303.4}{503.901}{\Character{21}}%
\TextAt{21}{310.228}{502.107}{w}%
\TextAt{9}{320.339}{503.901}{\Character{58}}%
\TextAt{13}{326.911}{503.901}{\Character{20}}%
\TextAt{9}{333.65}{503.901}{\Character{93}}%
\TextAt{13}{336.901}{503.901}{\Character{59}}%
\TextAt{8}{79.226}{482.709}{where}%
\TextAt{13}{112.941}{482.709}{\Character{21}}%
\TextAt{21}{119.769}{480.915}{w}%
\TextAt{8}{126.559}{482.709}{,}%
\TextAt{8}{133.633}{482.709}{resp.}%
\TextAt{13}{161.561}{482.709}{\Character{20}}%
\TextAt{8}{168.3}{482.709}{,}%
\TextAt{8}{175.373}{482.709}{is}%
\TextAt{8}{187.043}{482.709}{the}%
\TextAt{8}{207.105}{482.709}{residue}%
\TextAt{8}{246.738}{482.709}{class}%
\TextAt{8}{274.08}{482.709}{field}%
\TextAt{8}{299.343}{482.709}{of}%
\TextAt{13}{312.579}{482.709}{w}%
\TextAt{8}{321.268}{482.709}{,}%
\TextAt{8}{328.342}{482.709}{resp.}%
\TextAt{13}{356.27}{482.709}{v}%
\TextAt{8}{362.359}{482.709}{.}%
\TextAt{8}{370.776}{482.709}{From}%
\TextAt{8}{401.406}{482.709}{(8.4)}%
\TextAt{8}{429.287}{482.709}{and}%
\TextAt{8}{451.955}{482.709}{(6.8),}%
\TextAt{8}{483.106}{482.709}{we}%
\TextAt{8}{500.239}{482.709}{obtain}%
\TextAt{8}{79.226}{465.663}{the}%
\TextAt{12}{99.383}{465.663}{fundamental}%
\TextAt{12}{176.9}{465.663}{identity}%
\TextAt{12}{226.777}{465.663}{of}%
\TextAt{12}{242.103}{465.663}{valuation}%
\TextAt{12}{300.409}{465.663}{theory}%
\TextAt{8}{338.582}{465.663}{:}%
\TextAt{12}{79.226}{441.204}{Proposition}%
\TextAt{12}{151.937}{441.204}{(8.5).}%
\TextAt{8}{187.253}{441.204}{If}%
\TextAt{13}{198.945}{441.204}{v}%
\TextAt{8}{208.931}{441.204}{is}%
\TextAt{8}{220.695}{441.204}{discrete}%
\TextAt{8}{263.674}{441.204}{and}%
\TextAt{13}{286.436}{441.204}{L}%
\TextAt{16}{297.722}{441.204}{\Character{106}}%
\TextAt{13}{304.363}{441.204}{K}%
\TextAt{8}{319.078}{441.204}{separable,}%
\TextAt{8}{373.773}{441.204}{then}%
\TextAt{20}{256.892}{428.996}{\Character{88}}%
\TextAt{21}{258.999}{403.093}{w}%
\TextAt{3}{265.291}{403.093}{\Character{106}}%
\TextAt{21}{267.643}{403.093}{v}%
\TextAt{13}{276.153}{417.638}{e}%
\TextAt{21}{281.578}{415.845}{w}%
\TextAt{13}{288.368}{417.638}{f}%
\TextAt{21}{294.138}{415.845}{w}%
\TextAt{9}{304.248}{417.638}{\Character{61}}%
\TextAt{9}{316.658}{417.638}{\Character{91}}%
\TextAt{13}{319.925}{417.638}{L}%
\TextAt{9}{331.21}{417.638}{\Character{58}}%
\TextAt{13}{337.783}{417.638}{K}%
\TextAt{9}{348.601}{417.638}{\Character{93}}%
\TextAt{13}{351.852}{417.638}{\Character{58}}%
\TextAt{8}{97.159}{376.01}{This}%
\TextAt{8}{124.807}{376.01}{proposition}%
\TextAt{8}{187.269}{376.01}{repeats}%
\TextAt{8}{228.905}{376.01}{what}%
\TextAt{8}{259.099}{376.01}{we}%
\TextAt{8}{277.255}{376.01}{have}%
\TextAt{8}{305.178}{376.01}{already}%
\TextAt{8}{347.412}{376.01}{seen}%
\TextAt{8}{373.757}{376.01}{in}%
\TextAt{8}{388.338}{376.01}{chap.}%
\TextAt{8}{420.159}{376.01}{I,}%
\TextAt{8}{432.457}{376.01}{(8.2),}%
\TextAt{8}{464.844}{376.01}{working}%
\TextAt{8}{510.007}{376.01}{with}%
\TextAt{8}{79.226}{358.964}{the}%
\TextAt{8}{99.749}{358.964}{prime}%
\TextAt{8}{133.279}{358.964}{decomposition.}%
\TextAt{8}{216.065}{358.964}{If}%
\TextAt{13}{228.124}{358.964}{K}%
\TextAt{8}{243.205}{358.964}{is}%
\TextAt{8}{255.336}{358.964}{the}%
\TextAt{8}{275.859}{358.964}{field}%
\TextAt{8}{301.582}{358.964}{of}%
\TextAt{8}{315.279}{358.964}{fractions}%
\TextAt{8}{363.514}{358.964}{of}%
\TextAt{8}{377.211}{358.964}{a}%
\TextAt{8}{387.333}{358.964}{Dedekind}%
\TextAt{8}{439.884}{358.964}{domain}%
\TextAt{14}{481.878}{358.964}{o}%
\TextAt{8}{487.722}{358.964}{,}%
\TextAt{8}{495.33}{358.964}{then}%
\TextAt{8}{522.357}{358.964}{to}%
\TextAt{8}{79.226}{341.918}{every}%
\TextAt{8}{110.297}{341.918}{nonzero}%
\TextAt{8}{154.052}{341.918}{prime}%
\TextAt{8}{187.394}{341.918}{ideal}%
\TextAt{14}{215.535}{341.918}{p}%
\TextAt{8}{225.594}{341.918}{of}%
\TextAt{14}{239.102}{341.918}{o}%
\TextAt{8}{249.022}{341.918}{is}%
\TextAt{8}{260.965}{341.918}{associated}%
\TextAt{8}{316.878}{341.918}{the}%
\TextAt{14}{337.213}{341.918}{p}%
\TextAt{12}{343.196}{341.918}{-adic}%
\TextAt{12}{376.099}{341.918}{valuation}%
\TextAt{13}{433.997}{341.918}{v}%
\TextAt{27}{439.656}{340.125}{p}%
\TextAt{8}{448.218}{341.918}{of}%
\TextAt{13}{461.727}{341.918}{K}%
\TextAt{8}{472.544}{341.918}{,}%
\TextAt{8}{479.917}{341.918}{defined}%
\TextAt{8}{520.408}{341.918}{by}%
\TextAt{13}{79.226}{324.872}{v}%
\TextAt{27}{84.885}{323.079}{p}%
\TextAt{9}{89.372}{324.872}{\Character{40}}%
\TextAt{13}{93.925}{324.872}{a}%
\TextAt{9}{100.07}{324.872}{\Character{41}}%
\TextAt{9}{108.666}{324.872}{\Character{61}}%
\TextAt{13}{121.847}{324.872}{\Character{23}}%
\TextAt{27}{127.63}{323.079}{p}%
\TextAt{8}{132.117}{324.872}{,}%
\TextAt{8}{139.809}{324.872}{where}%
\TextAt{9}{174.052}{324.872}{\Character{40}}%
\TextAt{13}{178.604}{324.872}{a}%
\TextAt{9}{184.749}{324.872}{\Character{41}}%
\TextAt{9}{193.345}{324.872}{\Character{61}}%
\TextAt{20}{206.526}{333.838}{\Character{81}}%
\TextAt{27}{217.817}{321.385}{p}%
\TextAt{14}{224.296}{324.872}{p}%
\TextAt{21}{230.279}{329.21}{\Character{23}}%
\TextAt{28}{234.442}{328.214}{p}%
\TextAt{8}{243.354}{324.872}{(see}%
\TextAt{8}{267.256}{324.872}{chap.}%
\TextAt{8}{298.582}{324.872}{I,}%
\TextAt{8}{310.385}{324.872}{§11,}%
\TextAt{8}{335.471}{324.872}{p.}%
\TextAt{8}{349.558}{324.872}{67).}%
\TextAt{8}{375.578}{324.872}{The}%
\TextAt{8}{400.066}{324.872}{valuation}%
\TextAt{8}{451.572}{324.872}{ring}%
\TextAt{8}{476.072}{324.872}{of}%
\TextAt{13}{489.836}{324.872}{v}%
\TextAt{27}{495.495}{323.079}{p}%
\TextAt{8}{504.313}{324.872}{is}%
\TextAt{8}{516.51}{324.872}{the}%
\TextAt{8}{79.226}{307.826}{localization}%
\TextAt{14}{141.345}{307.826}{o}%
\TextAt{27}{147.189}{306.033}{p}%
\TextAt{8}{151.676}{307.826}{.}%
\TextAt{8}{160.125}{307.826}{If}%
\TextAt{16}{171.817}{307.826}{\Character{79}}%
\TextAt{8}{185.564}{307.826}{is}%
\TextAt{8}{197.328}{307.826}{the}%
\TextAt{8}{217.484}{307.826}{integral}%
\TextAt{8}{260.092}{307.826}{closure}%
\TextAt{8}{299.174}{307.826}{of}%
\TextAt{14}{312.504}{307.826}{o}%
\TextAt{8}{322.246}{307.826}{in}%
\TextAt{13}{335.898}{307.826}{L}%
\TextAt{8}{347.76}{307.826}{and}%
\TextAt{8}{370.523}{307.826}{if}%
\TextAt{14}{262.768}{284.261}{p}%
\TextAt{16}{268.751}{284.261}{\Character{79}}%
\TextAt{9}{281.921}{284.261}{\Character{61}}%
\TextAt{14}{294.347}{284.261}{P}%
\TextAt{21}{304.24}{289.486}{e}%
\TextAt{24}{308.168}{288.379}{1}%
\TextAt{17}{304.24}{281.33}{1}%
\TextAt{16}{314.81}{284.261}{\Character{1}}%
\TextAt{16}{320.106}{284.261}{\Character{1}}%
\TextAt{16}{325.414}{284.261}{\Character{1}}%
\TextAt{14}{330.75}{284.261}{P}%
\TextAt{21}{340.644}{289.197}{e}%
\TextAt{22}{344.572}{288.201}{r}%
\TextAt{21}{340.644}{281.305}{r}%
\TextAt{8}{79.226}{260.696}{is}%
\TextAt{8}{91.798}{260.696}{the}%
\TextAt{8}{112.763}{260.696}{prime}%
\TextAt{8}{146.735}{260.696}{decomposition}%
\TextAt{8}{224.679}{260.696}{of}%
\TextAt{14}{238.817}{260.696}{p}%
\TextAt{8}{249.506}{260.696}{in}%
\TextAt{13}{263.967}{260.696}{L}%
\TextAt{8}{271.932}{260.696}{,}%
\TextAt{8}{280.092}{260.696}{then}%
\TextAt{8}{307.56}{260.696}{the}%
\TextAt{8}{328.525}{260.696}{valuations}%
\TextAt{13}{385.021}{260.696}{w}%
\TextAt{21}{393.388}{258.903}{i}%
\TextAt{9}{401.468}{260.696}{\Character{61}}%
\TextAt{17}{417.895}{265.403}{1}%
\Vector{q 0 G 0 g 0.478 w 344.467 -510.315 m 351.557 -510.315 l S Q}%
\TextAt{21}{416.467}{256.573}{e}%
\TextAt{22}{420.395}{255.358}{i}%
\TextAt{13}{424.752}{260.696}{v}%
\TextAt{27}{430.412}{258.903}{P}%
\TextAt{22}{437.007}{257.688}{i}%
\TextAt{8}{440.667}{260.696}{,}%
\TextAt{13}{448.827}{260.696}{i}%
\TextAt{9}{457.519}{260.696}{\Character{61}}%
\TextAt{9}{471.315}{260.696}{1}%
\TextAt{13}{477.175}{260.696}{\Character{59}}%
\TextAt{13}{482.399}{260.696}{\Character{58}}%
\TextAt{13}{487.636}{260.696}{\Character{58}}%
\TextAt{13}{492.86}{260.696}{\Character{58}}%
\TextAt{13}{498.097}{260.696}{\Character{59}}%
\TextAt{13}{503.333}{260.696}{r}%
\TextAt{8}{508.996}{260.696}{,}%
\TextAt{8}{517.156}{260.696}{are}%
\TextAt{8}{79.226}{243.65}{precisely}%
\TextAt{8}{127.12}{243.65}{the}%
\TextAt{8}{147.314}{243.65}{extensions}%
\TextAt{8}{203.732}{243.65}{of}%
\TextAt{13}{217.099}{243.65}{v}%
\TextAt{9}{226.572}{243.65}{\Character{61}}%
\TextAt{13}{239.061}{243.65}{v}%
\TextAt{27}{244.721}{241.857}{p}%
\TextAt{8}{253.142}{243.65}{to}%
\TextAt{13}{267.49}{243.65}{L}%
\TextAt{8}{275.454}{243.65}{,}%
\TextAt{13}{282.65}{243.65}{e}%
\TextAt{21}{288.076}{241.857}{i}%
\TextAt{8}{295.392}{243.65}{are}%
\TextAt{8}{314.94}{243.65}{the}%
\TextAt{8}{335.134}{243.65}{corresponding}%
\TextAt{8}{410.357}{243.65}{ramification}%
\TextAt{8}{475.443}{243.65}{indices}%
\TextAt{8}{513.904}{243.65}{and}%
\TextAt{13}{79.226}{226.604}{f}%
\TextAt{21}{84.996}{224.811}{i}%
\TextAt{9}{91.698}{226.604}{\Character{61}}%
\TextAt{9}{104.108}{226.604}{\Character{91}}%
\TextAt{16}{107.375}{226.604}{\Character{79}}%
\TextAt{9}{117.225}{226.604}{\Character{47}}%
\TextAt{14}{123.078}{226.604}{P}%
\TextAt{21}{132.971}{224.811}{i}%
\TextAt{9}{139.673}{226.604}{\Character{58}}%
\TextAt{14}{146.246}{226.604}{o}%
\TextAt{9}{152.09}{226.604}{\Character{47}}%
\TextAt{14}{157.943}{226.604}{p}%
\TextAt{9}{163.926}{226.604}{\Character{93}}%
\TextAt{8}{171.075}{226.604}{the}%
\TextAt{8}{191.232}{226.604}{inertia}%
\TextAt{8}{228.305}{226.604}{degrees.}%
\TextAt{8}{273.886}{226.604}{The}%
\TextAt{8}{297.94}{226.604}{fundamental}%
\TextAt{8}{365.582}{226.604}{identity}%
\TextAt{21}{268.532}{211.473}{r}%
\TextAt{20}{261.926}{207.886}{\Character{88}}%
\TextAt{21}{263.708}{182.692}{i}%
\TextAt{17}{266.591}{182.692}{\Character{61}1}%
\TextAt{13}{281.187}{196.529}{e}%
\TextAt{21}{286.612}{194.735}{i}%
\TextAt{13}{289.994}{196.529}{f}%
\TextAt{21}{295.764}{194.735}{i}%
\TextAt{9}{302.466}{196.529}{\Character{61}}%
\TextAt{9}{314.875}{196.529}{\Character{91}}%
\TextAt{13}{318.143}{196.529}{L}%
\TextAt{9}{329.428}{196.529}{\Character{58}}%
\TextAt{13}{336.001}{196.529}{K}%
\TextAt{9}{346.818}{196.529}{\Character{93}}%
\TextAt{8}{79.226}{165.564}{has}%
\TextAt{8}{100.959}{165.564}{thus}%
\TextAt{8}{127.57}{165.564}{been}%
\TextAt{8}{156.057}{165.564}{established}%
\TextAt{8}{216.872}{165.564}{in}%
\TextAt{8}{231.384}{165.564}{two}%
\TextAt{8}{254.36}{165.564}{different}%
\TextAt{8}{301.39}{165.564}{ways.}%
\TextAt{8}{336.876}{165.564}{The}%
\TextAt{19}{361.789}{165.564}{raison}%
\TextAt{19}{398.135}{165.564}{d’être}%
\TextAt{8}{431.524}{165.564}{of}%
\TextAt{8}{445.713}{165.564}{valuation}%
\TextAt{8}{497.645}{165.564}{theory,}%
\TextAt{8}{79.226}{148.518}{however,}%
\TextAt{8}{127.143}{148.518}{is}%
\TextAt{8}{138.638}{148.518}{not}%
\TextAt{8}{159.184}{148.518}{to}%
\TextAt{8}{173.226}{148.518}{reformulate}%
\TextAt{8}{235.399}{148.518}{ideal-theoretic}%
\TextAt{8}{311.871}{148.518}{knowledge,}%
\TextAt{8}{371.492}{148.518}{but}%
\TextAt{8}{392.683}{148.518}{rather,}%
\TextAt{8}{430.844}{148.518}{as}%
\TextAt{8}{444.946}{148.518}{has}%
\TextAt{8}{465.551}{148.518}{been}%
\TextAt{8}{492.911}{148.518}{stressed}%
\TextAt{8}{79.226}{131.473}{earlier,}%
\TextAt{8}{120.155}{131.473}{to}%
\TextAt{8}{135.991}{131.473}{provide}%
\TextAt{8}{179.144}{131.473}{the}%
\TextAt{8}{200.826}{131.473}{possibility}%
\TextAt{8}{258.086}{131.473}{of}%
\TextAt{8}{272.942}{131.473}{passing}%
\TextAt{8}{315.569}{131.473}{from}%
\TextAt{8}{344.735}{131.473}{the}%
\TextAt{8}{366.417}{131.473}{extension}%
\TextAt{13}{419.708}{131.473}{L}%
\TextAt{16}{433.593}{131.472}{\Character{106}}%
\TextAt{13}{442.835}{131.473}{K}%
\TextAt{8}{459.075}{131.473}{to}%
\TextAt{8}{474.911}{131.473}{the}%
\TextAt{8}{496.593}{131.473}{various}%
\TextAt{8}{79.226}{114.427}{completions}%
\TextAt{13}{144.28}{114.427}{L}%
\TextAt{21}{152.245}{112.633}{w}%
\TextAt{16}{163.383}{114.427}{\Character{106}}%
\TextAt{13}{171.054}{114.427}{K}%
\TextAt{21}{181.013}{112.633}{v}%
\TextAt{8}{190.421}{114.427}{where}%
\TextAt{8}{224.834}{114.427}{much}%
\TextAt{8}{256.652}{114.427}{simpler}%
\TextAt{8}{298.286}{114.427}{arithmetic}%
\TextAt{8}{355.474}{114.427}{laws}%
\TextAt{8}{381.829}{114.427}{apply.}%
\TextAt{8}{419.405}{114.427}{Let}%
\TextAt{8}{440.978}{114.427}{us}%
\TextAt{8}{456.597}{114.427}{also}%
\TextAt{8}{480.681}{114.427}{emphasize}%
\TextAt{8}{79.226}{97.381}{once}%
\TextAt{8}{105.886}{97.381}{more}%
\TextAt{8}{135.152}{97.381}{that}%
\TextAt{8}{160.521}{97.381}{completions}%
\TextAt{8}{224.972}{97.381}{may}%
\TextAt{8}{250.34}{97.381}{always}%
\TextAt{8}{287.808}{97.381}{be}%
\TextAt{8}{303.732}{97.381}{replaced}%
\TextAt{8}{349.903}{97.381}{with}%
\TextAt{8}{376.563}{97.381}{henselizations.}%
\TextAt{12}{79.226}{72.921}{Exercise}%
\TextAt{12}{131.421}{72.921}{II.8.1.}%
\TextAt{8}{169.52}{72.921}{Up}%
\TextAt{8}{187.822}{72.921}{to}%
\TextAt{8}{201.258}{72.921}{equivalence,}%
\TextAt{8}{265.595}{72.921}{the}%
\TextAt{8}{284.877}{72.921}{valuations}%
\TextAt{8}{339.69}{72.921}{of}%
\TextAt{8}{352.145}{72.921}{the}%
\TextAt{8}{371.427}{72.921}{field}%
\TextAt{18}{395.91}{72.921}{Q}%
\TextAt{9}{405.209}{72.921}{\Character{40}}%
\TextAt{16}{409.761}{82.811}{\Character{112}}%
\Vector{q 0 G 0 g 0.478 w 347.724 -690.95 m 353.577 -690.95 l S Q}%
\TextAt{9}{419.724}{72.921}{5\Character{41}}%
\TextAt{8}{433.152}{72.921}{are}%
\TextAt{8}{451.788}{72.921}{given}%
\TextAt{8}{481.483}{72.921}{as}%
\TextAt{8}{494.979}{72.921}{follows.}%
\TextAt{8}{297.211}{43.033}{151}%
\Destination{78.23}{809.815}{nav.page.151}%
\Destination{79.226}{86.955}{nav.ex.2.8.1}%
\Destination{79.36}{600.155}{nav.print.165}%
\Destination{79.36}{455.155}{nav.number.2.8.5}%
\Hotlink{419.659}{373.855}{5.22}{11.047}{nav.chapter.1}%
\Hotlink{298.082}{322.655}{5.22}{11.047}{nav.chapter.1}%
\Hotlink{315.564}{322.655}{12.716}{11.047}{nav.section.1.11}%
\Hotlink{349.058}{322.655}{12.716}{11.047}{nav.print.67}%
\Hotlink{275.532}{758.955}{25.078}{11.047}{nav.number.2.8.3}%
\Hotlink{400.906}{480.555}{25.078}{11.047}{nav.number.2.8.4}%
\Hotlink{451.455}{480.555}{25.078}{11.047}{nav.number.2.6.8}%
\Hotlink{431.957}{373.855}{25.078}{11.047}{nav.number.1.8.2}%
}

% Page 164 (body page 152)
\Page{%
\TextAt{2}{79.226}{796.943}{II.}%
\TextAt{2}{94.073}{796.943}{The}%
\TextAt{2}{116.492}{796.943}{Theory}%
\TextAt{2}{154.401}{796.943}{of}%
\TextAt{2}{166.826}{796.943}{Valuations}%
\Vector{q 0 G 0 g 0.299 w 7.226 18.729 m 460.769 18.729 l S Q}%
\TextAt{8}{91.528}{761.124}{1)}%
\TextAt{20}{107.918}{771.286}{\Character{12}}%
\TextAt{20}{107.918}{764.112}{\Character{12}}%
\TextAt{13}{111.904}{761.124}{a}%
\TextAt{9}{120.705}{761.124}{\Character{43}}%
\TextAt{13}{132.466}{761.124}{b}%
\TextAt{16}{137.444}{771.013}{\Character{112}}%
\Vector{q 0 G 0 g 0.478 w 75.406 -2.748 m 81.259 -2.748 l S Q}%
\TextAt{9}{147.406}{761.124}{5}%
\TextAt{20}{153.259}{771.286}{\Character{12}}%
\TextAt{20}{153.259}{764.112}{\Character{12}}%
\TextAt{17}{157.244}{756.441}{1}%
\TextAt{9}{165.297}{761.124}{\Character{61}}%
\TextAt{20}{177.723}{771.286}{\Character{12}}%
\TextAt{20}{177.723}{764.112}{\Character{12}}%
\TextAt{13}{181.708}{761.124}{a}%
\TextAt{9}{190.51}{761.124}{\Character{43}}%
\TextAt{13}{202.271}{761.124}{b}%
\TextAt{16}{207.248}{771.013}{\Character{112}}%
\Vector{q 0 G 0 g 0.478 w 145.211 -2.748 m 151.064 -2.748 l S Q}%
\TextAt{9}{217.211}{761.124}{5}%
\TextAt{20}{223.064}{771.286}{\Character{12}}%
\TextAt{20}{223.064}{764.112}{\Character{12}}%
\TextAt{8}{229.807}{761.124}{and}%
\TextAt{20}{251.43}{771.286}{\Character{12}}%
\TextAt{20}{251.43}{764.112}{\Character{12}}%
\TextAt{13}{255.416}{761.124}{a}%
\TextAt{9}{264.217}{761.124}{\Character{43}}%
\TextAt{13}{275.978}{761.124}{b}%
\TextAt{16}{280.956}{771.013}{\Character{112}}%
\Vector{q 0 G 0 g 0.478 w 218.918 -2.748 m 224.771 -2.748 l S Q}%
\TextAt{9}{290.918}{761.124}{5}%
\TextAt{20}{296.771}{771.286}{\Character{12}}%
\TextAt{20}{296.771}{764.112}{\Character{12}}%
\TextAt{17}{300.756}{756.441}{2}%
\TextAt{9}{308.809}{761.124}{\Character{61}}%
\TextAt{20}{321.235}{771.286}{\Character{12}}%
\TextAt{20}{321.235}{764.112}{\Character{12}}%
\TextAt{13}{325.22}{761.124}{a}%
\TextAt{16}{334.022}{761.124}{\Character{0}}%
\TextAt{13}{345.977}{761.124}{b}%
\TextAt{16}{350.954}{771.013}{\Character{112}}%
\Vector{q 0 G 0 g 0.478 w 288.917 -2.748 m 294.77 -2.748 l S Q}%
\TextAt{9}{360.917}{761.124}{5}%
\TextAt{20}{366.77}{771.286}{\Character{12}}%
\TextAt{20}{366.769}{764.112}{\Character{12}}%
\TextAt{8}{373.513}{761.124}{are}%
\TextAt{8}{391.885}{761.124}{the}%
\TextAt{8}{410.902}{761.124}{archimedean}%
\TextAt{8}{477.728}{761.124}{valuations.}%
\TextAt{8}{91.528}{744.078}{2)}%
\TextAt{8}{107.918}{744.078}{If}%
\TextAt{13}{120.134}{744.078}{p}%
\TextAt{9}{130.221}{744.078}{\Character{61}}%
\TextAt{9}{143.527}{744.078}{2}%
\TextAt{8}{153.811}{744.078}{or}%
\TextAt{8}{168.645}{744.078}{5}%
\TextAt{8}{178.923}{744.078}{or}%
\TextAt{8}{193.757}{744.078}{a}%
\TextAt{8}{204.035}{744.078}{prime}%
\TextAt{8}{237.722}{744.078}{number}%
\NativeMathAt{280.841}{744.078}{11.955}{\neq}%
\TextAt{9}{294.147}{744.078}{2}%
\TextAt{13}{300.011}{744.078}{\Character{59}}%
\TextAt{9}{305.255}{744.078}{5}%
\TextAt{8}{315.529}{744.078}{such}%
\TextAt{8}{342.449}{744.078}{that}%
\TextAt{20}{368.341}{753.761}{\Character{0}}%
\TextAt{21}{375.016}{749.334}{p}%
\Vector{q 0 G 0 g 0.478 w 303.016 -26.933 m 307.279 -26.933 l S Q}%
\TextAt{17}{375.03}{739.955}{5}%
\TextAt{20}{380.474}{753.761}{\Character{1}}%
\TextAt{9}{390.166}{744.078}{\Character{61}}%
\TextAt{16}{403.483}{744.078}{\Character{0}}%
\TextAt{9}{412.781}{744.078}{1}%
\TextAt{8}{418.634}{744.078}{,}%
\TextAt{8}{426.437}{744.078}{then}%
\TextAt{8}{453.62}{744.078}{there}%
\TextAt{8}{484.055}{744.078}{is}%
\TextAt{8}{496.342}{744.078}{exactly}%
\TextAt{8}{107.918}{727.032}{one}%
\TextAt{8}{129.378}{727.032}{extension}%
\TextAt{8}{181.144}{727.032}{of}%
\TextAt{16}{194.474}{727.032}{\Character{106}}%
\TextAt{16}{201.671}{727.032}{\Character{106}}%
\TextAt{21}{205.013}{723.545}{p}%
\TextAt{8}{213.671}{727.032}{to}%
\TextAt{18}{227.982}{727.032}{Q}%
\TextAt{9}{237.28}{727.032}{\Character{40}}%
\TextAt{16}{241.833}{736.921}{\Character{112}}%
\Vector{q 0 G 0 g 0.478 w 179.795 -36.84 m 185.648 -36.84 l S Q}%
\TextAt{9}{251.795}{727.032}{5\Character{41}}%
\TextAt{8}{262.201}{727.032}{,}%
\TextAt{8}{269.35}{727.032}{namely}%
\TextAt{20}{251.831}{710.212}{\Character{12}}%
\TextAt{20}{251.831}{703.038}{\Character{12}}%
\TextAt{20}{251.831}{695.865}{\Character{12}}%
\TextAt{13}{255.816}{696.463}{a}%
\TextAt{9}{264.618}{696.463}{\Character{43}}%
\TextAt{13}{276.379}{696.463}{b}%
\TextAt{16}{281.356}{706.936}{\Character{112}}%
\Vector{q 0 G 0 g 0.478 w 219.319 -66.825 m 225.172 -66.825 l S Q}%
\TextAt{9}{291.319}{696.463}{5}%
\TextAt{20}{297.172}{710.212}{\Character{12}}%
\TextAt{20}{297.172}{703.038}{\Character{12}}%
\TextAt{20}{297.172}{695.865}{\Character{12}}%
\TextAt{27}{301.157}{688.194}{p}%
\TextAt{9}{308.965}{696.463}{\Character{61}}%
\TextAt{20}{321.39}{706.625}{\Character{12}}%
\TextAt{20}{321.39}{699.452}{\Character{12}}%
\TextAt{13}{325.375}{696.463}{a}%
\TextAt{17}{331.52}{701.399}{2}%
\TextAt{16}{338.909}{696.463}{\Character{0}}%
\TextAt{9}{350.864}{696.463}{5}%
\TextAt{13}{356.717}{696.463}{b}%
\TextAt{17}{361.694}{701.399}{2}%
\TextAt{20}{366.427}{706.625}{\Character{12}}%
\TextAt{20}{366.427}{699.452}{\Character{12}}%
\TextAt{17}{370.412}{703.47}{1\Character{47}2}%
\TextAt{21}{370.412}{691.78}{p}%
\TextAt{13}{385.605}{696.463}{\Character{58}}%
\TextAt{8}{91.528}{664.742}{3)}%
\TextAt{8}{107.918}{664.742}{If}%
\TextAt{13}{120.136}{664.742}{p}%
\TextAt{8}{130.434}{664.742}{is}%
\TextAt{8}{142.723}{664.742}{a}%
\TextAt{8}{153.003}{664.742}{prime}%
\TextAt{8}{186.692}{664.742}{number}%
\NativeMathAt{229.813}{664.742}{11.955}{\neq}%
\TextAt{9}{243.119}{664.742}{2}%
\TextAt{13}{248.987}{664.742}{\Character{59}}%
\TextAt{9}{254.231}{664.742}{5}%
\TextAt{8}{264.507}{664.742}{such}%
\TextAt{8}{291.429}{664.742}{that}%
\TextAt{20}{317.323}{674.425}{\Character{0}}%
\TextAt{21}{323.998}{669.997}{p}%
\Vector{q 0 G 0 g 0.478 w 251.998 -106.27 m 256.261 -106.27 l S Q}%
\TextAt{17}{324.012}{660.619}{5}%
\TextAt{20}{329.456}{674.425}{\Character{1}}%
\TextAt{9}{339.151}{664.742}{\Character{61}}%
\TextAt{9}{352.457}{664.742}{1}%
\TextAt{8}{358.325}{664.742}{,}%
\TextAt{8}{366.131}{664.742}{then}%
\TextAt{8}{393.316}{664.742}{there}%
\TextAt{8}{423.753}{664.742}{are}%
\TextAt{8}{443.789}{664.742}{two}%
\TextAt{8}{466.431}{664.742}{extensions}%
\TextAt{8}{523.337}{664.742}{of}%
\TextAt{16}{107.918}{647.696}{\Character{106}}%
\TextAt{16}{115.127}{647.696}{\Character{106}}%
\TextAt{21}{118.458}{644.209}{p}%
\TextAt{8}{127.116}{647.696}{to}%
\TextAt{18}{141.426}{647.696}{Q}%
\TextAt{9}{150.725}{647.696}{\Character{40}}%
\TextAt{16}{155.277}{657.585}{\Character{112}}%
\Vector{q 0 G 0 g 0.478 w 93.24 -116.176 m 99.093 -116.176 l S Q}%
\TextAt{9}{165.24}{647.696}{5\Character{41}}%
\TextAt{8}{175.645}{647.696}{,}%
\TextAt{8}{182.794}{647.696}{namely}%
\TextAt{20}{165.64}{630.875}{\Character{12}}%
\TextAt{20}{165.64}{623.702}{\Character{12}}%
\TextAt{20}{165.64}{616.529}{\Character{12}}%
\TextAt{13}{169.625}{617.127}{a}%
\TextAt{9}{178.427}{617.127}{\Character{43}}%
\TextAt{13}{190.188}{617.127}{b}%
\TextAt{16}{195.165}{627.6}{\Character{112}}%
\Vector{q 0 G 0 g 0.478 w 133.128 -146.161 m 138.981 -146.161 l S Q}%
\TextAt{9}{205.128}{617.127}{5}%
\TextAt{20}{210.981}{630.875}{\Character{12}}%
\TextAt{20}{210.981}{623.702}{\Character{12}}%
\TextAt{20}{210.981}{616.529}{\Character{12}}%
\TextAt{27}{214.966}{608.858}{p}%
\TextAt{24}{218.954}{607.751}{1}%
\TextAt{9}{226.924}{617.127}{\Character{61}}%
\TextAt{16}{239.35}{617.127}{\Character{106}}%
\TextAt{13}{242.671}{617.127}{a}%
\TextAt{9}{251.472}{617.127}{\Character{43}}%
\TextAt{13}{263.234}{617.127}{b\Character{13}}%
\TextAt{16}{274.933}{617.127}{\Character{106}}%
\TextAt{21}{278.254}{613.64}{p}%
\TextAt{13}{285.007}{617.127}{\Character{59}}%
\TextAt{8}{314.162}{617.127}{resp.}%
\TextAt{20}{352.235}{630.875}{\Character{12}}%
\TextAt{20}{352.235}{623.702}{\Character{12}}%
\TextAt{20}{352.235}{616.529}{\Character{12}}%
\TextAt{13}{356.22}{617.127}{a}%
\TextAt{9}{365.022}{617.127}{\Character{43}}%
\TextAt{13}{376.783}{617.127}{b}%
\TextAt{16}{381.76}{627.6}{\Character{112}}%
\Vector{q 0 G 0 g 0.478 w 319.723 -146.161 m 325.576 -146.161 l S Q}%
\TextAt{9}{391.723}{617.127}{5}%
\TextAt{20}{397.576}{630.875}{\Character{12}}%
\TextAt{20}{397.576}{623.702}{\Character{12}}%
\TextAt{20}{397.576}{616.529}{\Character{12}}%
\TextAt{27}{401.561}{608.858}{p}%
\TextAt{24}{405.55}{607.751}{2}%
\TextAt{9}{413.52}{617.127}{\Character{61}}%
\TextAt{16}{425.945}{617.127}{\Character{106}}%
\TextAt{13}{429.266}{617.127}{a}%
\TextAt{16}{438.068}{617.127}{\Character{0}}%
\TextAt{13}{450.023}{617.127}{b\Character{13}}%
\TextAt{16}{461.722}{617.127}{\Character{106}}%
\TextAt{21}{465.043}{613.64}{p}%
\TextAt{13}{471.796}{617.127}{\Character{59}}%
\TextAt{8}{107.918}{586.093}{where}%
\TextAt{13}{141.728}{586.093}{\Character{13}}%
\TextAt{8}{152.347}{586.093}{is}%
\TextAt{8}{164.111}{586.093}{a}%
\TextAt{8}{173.867}{586.093}{solution}%
\TextAt{8}{218.161}{586.093}{of}%
\TextAt{13}{231.491}{586.093}{x}%
\TextAt{17}{238.143}{590.431}{2}%
\TextAt{16}{245.532}{586.093}{\Character{0}}%
\TextAt{9}{257.487}{586.093}{5}%
\TextAt{9}{266.645}{586.093}{\Character{61}}%
\TextAt{9}{279.066}{586.093}{0}%
\TextAt{8}{288.836}{586.093}{in}%
\TextAt{18}{302.489}{586.093}{Q}%
\TextAt{21}{311.788}{584.299}{p}%
\TextAt{8}{316.549}{586.093}{.}%
\TextAt{12}{79.226}{554.77}{Exercise}%
\TextAt{12}{132.427}{554.77}{II.8.2.}%
\TextAt{8}{171.401}{554.77}{Determine}%
\TextAt{8}{228.463}{554.77}{the}%
\TextAt{8}{248.619}{554.77}{valuations}%
\TextAt{8}{304.306}{554.77}{of}%
\TextAt{8}{317.636}{554.77}{the}%
\TextAt{8}{337.793}{554.77}{field}%
\TextAt{18}{363.149}{554.77}{Q}%
\TextAt{9}{372.448}{554.77}{\Character{40}}%
\TextAt{13}{377}{554.77}{i}%
\TextAt{9}{380.994}{554.77}{\Character{41}}%
\TextAt{8}{389.443}{554.77}{of}%
\TextAt{8}{402.773}{554.77}{the}%
\TextAt{8}{422.93}{554.77}{Gaussian}%
\TextAt{8}{473.213}{554.77}{numbers.}%
\TextAt{12}{79.226}{523.447}{Exercise}%
\TextAt{12}{133.008}{523.447}{II.8.3.}%
\TextAt{8}{172.488}{523.447}{How}%
\TextAt{8}{199.654}{523.447}{many}%
\TextAt{8}{232.032}{523.447}{extensions}%
\TextAt{8}{288.918}{523.447}{to}%
\TextAt{18}{303.734}{523.447}{Q}%
\TextAt{9}{313.033}{523.447}{\Character{40}}%
\TextAt{22}{320.906}{529.003}{n}%
\TextAt{16}{318.865}{533.337}{\Character{112}}%
\Vector{q 0 G 0 g 0.478 w 256.828 -240.424 m 262.681 -240.424 l S Q}%
\TextAt{9}{328.828}{523.447}{2\Character{41}}%
\TextAt{8}{343.636}{523.447}{does}%
\TextAt{8}{370.539}{523.447}{the}%
\TextAt{8}{391.201}{523.447}{archimedean}%
\TextAt{8}{459.672}{523.447}{absolute}%
\TextAt{8}{506.42}{523.447}{value}%
\TextAt{16}{79.226}{506.401}{\Character{106}}%
\TextAt{16}{86.435}{506.401}{\Character{106}}%
\TextAt{8}{93.663}{506.401}{of}%
\TextAt{18}{106.993}{506.401}{Q}%
\TextAt{8}{120.189}{506.401}{admit?}%
\TextAt{12}{79.226}{475.079}{Exercise}%
\TextAt{12}{133.814}{475.079}{II.8.4.}%
\TextAt{8}{173.994}{475.079}{Let}%
\TextAt{13}{196.17}{475.079}{L}%
\TextAt{16}{209.511}{475.079}{\Character{106}}%
\TextAt{13}{218.208}{475.079}{K}%
\TextAt{8}{234.129}{475.079}{be}%
\TextAt{8}{251.26}{475.079}{a}%
\TextAt{8}{262.221}{475.079}{finite}%
\TextAt{8}{293.339}{475.079}{separable}%
\TextAt{8}{345.989}{475.079}{extension,}%
\TextAt{14}{402.514}{475.079}{o}%
\TextAt{8}{413.462}{475.079}{the}%
\TextAt{8}{434.825}{475.079}{valuation}%
\TextAt{8}{487.103}{475.079}{ring}%
\TextAt{8}{512.375}{475.079}{of}%
\TextAt{8}{526.911}{475.079}{a}%
\TextAt{8}{79.226}{458.033}{discrete}%
\TextAt{8}{122.891}{458.033}{valuation}%
\TextAt{13}{174.649}{458.033}{v}%
\TextAt{8}{185.32}{458.033}{and}%
\TextAt{16}{208.769}{458.033}{\Character{79}}%
\TextAt{8}{223.201}{458.033}{its}%
\TextAt{8}{240.206}{458.033}{integral}%
\TextAt{8}{283.5}{458.033}{closure}%
\TextAt{8}{323.267}{458.033}{in}%
\TextAt{13}{337.606}{458.033}{L}%
\TextAt{8}{345.57}{458.033}{.}%
\TextAt{8}{356.076}{458.033}{If}%
\TextAt{13}{368.454}{458.033}{w}%
\TextAt{16}{381.632}{458.033}{\Character{106}}%
\TextAt{13}{389.443}{458.033}{v}%
\TextAt{8}{400.114}{458.033}{varies}%
\TextAt{8}{433.712}{458.033}{over}%
\TextAt{8}{459.444}{458.033}{the}%
\TextAt{8}{480.286}{458.033}{extensions}%
\TextAt{8}{79.226}{440.987}{of}%
\TextAt{13}{92.547}{440.987}{v}%
\TextAt{8}{102.524}{440.987}{to}%
\TextAt{13}{116.826}{440.987}{L}%
\TextAt{8}{128.679}{440.987}{and}%
\TextAt{20}{151.035}{441.522}{\Character{98}}%
\TextAt{14}{151.433}{440.987}{o}%
\TextAt{21}{157.278}{439.194}{v}%
\TextAt{8}{162.185}{440.987}{,}%
\TextAt{8}{169.327}{440.987}{resp.}%
\TextAt{20}{200.274}{444.009}{\Character{98}}%
\TextAt{16}{197.342}{440.987}{\Character{79}}%
\TextAt{21}{206.859}{439.194}{w}%
\TextAt{8}{213.649}{440.987}{,}%
\TextAt{8}{220.791}{440.987}{are}%
\TextAt{8}{240.293}{440.987}{the}%
\TextAt{8}{260.441}{440.987}{valuation}%
\TextAt{8}{311.505}{440.987}{rings}%
\TextAt{8}{340.177}{440.987}{of}%
\TextAt{8}{353.498}{440.987}{the}%
\TextAt{8}{373.646}{440.987}{completions}%
\TextAt{13}{438.087}{440.987}{K}%
\TextAt{21}{448.047}{439.194}{v}%
\TextAt{8}{452.954}{440.987}{,}%
\TextAt{8}{460.096}{440.987}{resp.}%
\TextAt{13}{488.111}{440.987}{L}%
\TextAt{21}{496.075}{439.194}{w}%
\TextAt{8}{502.865}{440.987}{,}%
\TextAt{8}{510.007}{440.987}{then}%
\TextAt{8}{79.226}{423.941}{one}%
\TextAt{8}{100.686}{423.941}{has}%
\TextAt{16}{259.81}{406.895}{\Character{79}}%
\TextAt{16}{272.315}{406.895}{\Character{10}}%
\TextAt{27}{281.615}{405.102}{o}%
\TextAt{20}{288.267}{407.43}{\Character{98}}%
\TextAt{14}{288.666}{406.895}{o}%
\TextAt{21}{294.51}{405.102}{v}%
\NativeMathAt{302.739}{406.895}{11.955}{\cong}%
\TextAt{20}{315.358}{418.252}{\Character{89}}%
\TextAt{21}{316.47}{392.349}{w}%
\TextAt{3}{322.761}{392.349}{\Character{106}}%
\TextAt{21}{325.113}{392.349}{v}%
\TextAt{20}{335.559}{409.917}{\Character{98}}%
\TextAt{16}{332.627}{406.895}{\Character{79}}%
\TextAt{21}{342.144}{405.102}{w}%
\TextAt{13}{348.934}{406.895}{\Character{58}}%
\TextAt{12}{79.226}{358.947}{Exercise}%
\TextAt{12}{133.312}{358.947}{II.8.5.}%
\TextAt{8}{173.056}{358.947}{How}%
\TextAt{8}{200.485}{358.947}{does}%
\TextAt{8}{227.652}{358.947}{proposition}%
\TextAt{8}{289.955}{358.947}{(8.2)}%
\TextAt{8}{318.7}{358.947}{relate}%
\TextAt{8}{351.988}{358.947}{to}%
\TextAt{8}{367.068}{358.947}{Dedekind’s}%
\TextAt{8}{427.888}{358.947}{proposition,}%
\TextAt{8}{493.636}{358.947}{chap.}%
\TextAt{8}{525.297}{358.947}{I,}%
\TextAt{8}{79.226}{341.901}{(8.3)?}%
\TextAt{12}{79.226}{310.578}{Exercise}%
\TextAt{12}{132.889}{310.578}{II.8.6.}%
\TextAt{8}{172.265}{310.578}{Let}%
\TextAt{13}{193.636}{310.578}{L}%
\TextAt{16}{205.606}{310.578}{\Character{106}}%
\TextAt{13}{212.932}{310.578}{K}%
\TextAt{8}{228.049}{310.578}{be}%
\TextAt{8}{244.376}{310.578}{a}%
\TextAt{8}{254.533}{310.578}{finite}%
\TextAt{8}{284.847}{310.578}{field}%
\TextAt{8}{310.605}{310.578}{extension,}%
\TextAt{13}{366.125}{310.578}{v}%
\TextAt{8}{376.513}{310.578}{a}%
\TextAt{8}{386.67}{310.578}{nonarchimedean}%
\TextAt{8}{473.902}{310.578}{exponential}%
\TextAt{8}{79.226}{293.532}{valuation,}%
\TextAt{8}{133.47}{293.532}{and}%
\TextAt{13}{156.132}{293.532}{w}%
\TextAt{8}{168.618}{293.532}{an}%
\TextAt{8}{184.776}{293.532}{extension}%
\TextAt{8}{236.441}{293.532}{to}%
\TextAt{13}{250.651}{293.532}{L}%
\TextAt{8}{258.615}{293.532}{.}%
\TextAt{8}{267.03}{293.532}{If}%
\TextAt{16}{278.621}{293.532}{\Character{79}}%
\TextAt{8}{292.268}{293.532}{is}%
\TextAt{8}{303.931}{293.532}{the}%
\TextAt{8}{323.986}{293.532}{integral}%
\TextAt{8}{366.494}{293.532}{closure}%
\TextAt{8}{405.474}{293.532}{of}%
\TextAt{8}{418.704}{293.532}{the}%
\TextAt{8}{438.759}{293.532}{valuation}%
\TextAt{8}{489.731}{293.532}{ring}%
\TextAt{14}{513.696}{293.532}{o}%
\TextAt{8}{523.337}{293.532}{of}%
\TextAt{13}{79.226}{276.486}{v}%
\TextAt{8}{89.361}{276.486}{in}%
\TextAt{13}{103.162}{276.486}{L}%
\TextAt{8}{111.127}{276.486}{,}%
\TextAt{8}{118.462}{276.486}{then}%
\TextAt{8}{145.271}{276.486}{the}%
\TextAt{8}{165.576}{276.486}{localization}%
\TextAt{16}{227.844}{276.486}{\Character{79}}%
\TextAt{27}{237.362}{274.693}{P}%
\TextAt{8}{248.502}{276.486}{of}%
\TextAt{16}{261.981}{276.486}{\Character{79}}%
\TextAt{8}{275.877}{276.486}{at}%
\TextAt{8}{290.336}{276.486}{the}%
\TextAt{8}{310.641}{276.486}{prime}%
\TextAt{8}{343.954}{276.486}{ideal}%
\TextAt{14}{372.066}{276.486}{P}%
\TextAt{9}{385.534}{276.486}{\Character{61}}%
\TextAt{16}{398.213}{276.486}{\Character{102}}%
\TextAt{13}{404.191}{276.486}{\Character{11}}%
\NativeMathAt{415.287}{276.486}{11.955}{\in}%
\TextAt{16}{426.812}{276.486}{\Character{79}}%
\TextAt{16}{440.238}{276.486}{\Character{106}}%
\TextAt{13}{447.151}{276.486}{w}%
\TextAt{9}{455.841}{276.486}{\Character{40}}%
\TextAt{13}{460.393}{276.486}{\Character{11}}%
\TextAt{9}{467.915}{276.486}{\Character{41}}%
\TextAt{13}{476.042}{276.486}{\Character{62}}%
\TextAt{9}{488.721}{276.486}{0}%
\TextAt{16}{494.574}{276.486}{\Character{103}}%
\TextAt{8}{504.598}{276.486}{is}%
\TextAt{8}{516.511}{276.486}{the}%
\TextAt{8}{79.226}{259.44}{valuation}%
\TextAt{8}{130.299}{259.44}{ring}%
\TextAt{8}{154.364}{259.44}{of}%
\TextAt{13}{167.694}{259.44}{w}%
\TextAt{8}{176.384}{259.44}{.}%
\TextAt{15}{79.226}{207.083}{§}%
\TextAt{15}{93.785}{207.083}{9}%
\TextAt{15}{120.693}{207.083}{Galois}%
\TextAt{15}{179.897}{207.083}{Theory}%
\TextAt{15}{247.244}{207.083}{of}%
\TextAt{15}{269.314}{207.083}{Valuations}%
\TextAt{8}{79.226}{175.197}{We}%
\TextAt{8}{99.063}{175.197}{now}%
\TextAt{8}{123.145}{175.197}{consider}%
\TextAt{8}{168.423}{175.197}{Galois}%
\TextAt{8}{204.029}{175.197}{extensions}%
\TextAt{13}{260.102}{175.197}{L}%
\TextAt{16}{271.388}{175.197}{\Character{106}}%
\TextAt{13}{278.029}{175.197}{K}%
\TextAt{8}{292.437}{175.197}{and}%
\TextAt{8}{314.893}{175.197}{study}%
\TextAt{8}{346.841}{175.197}{the}%
\TextAt{8}{366.69}{175.197}{effect}%
\TextAt{8}{397.263}{175.197}{of}%
\TextAt{8}{410.286}{175.197}{the}%
\TextAt{8}{430.135}{175.197}{Galois}%
\TextAt{8}{465.741}{175.197}{action}%
\TextAt{8}{500.559}{175.197}{on}%
\TextAt{8}{516.51}{175.197}{the}%
\TextAt{8}{79.226}{158.151}{extended}%
\TextAt{8}{128.842}{158.151}{valuations}%
\TextAt{13}{184.399}{158.151}{w}%
\TextAt{16}{196.41}{158.151}{\Character{106}}%
\TextAt{13}{203.052}{158.151}{v}%
\TextAt{8}{209.14}{158.151}{.}%
\TextAt{8}{217.545}{158.151}{This}%
\TextAt{8}{244.135}{158.151}{leads}%
\TextAt{8}{273.331}{158.151}{to}%
\TextAt{8}{287.512}{158.151}{a}%
\TextAt{8}{297.138}{158.151}{direct}%
\TextAt{8}{330.172}{158.151}{generalization}%
\TextAt{8}{404.846}{158.151}{of}%
\TextAt{8}{418.046}{158.151}{“Hilbert’s}%
\TextAt{8}{471.619}{158.151}{ramification}%
\TextAt{8}{79.226}{141.105}{theory”—see}%
\TextAt{8}{149.153}{141.105}{chap.}%
\TextAt{8}{180.978}{141.105}{I,}%
\TextAt{8}{193.28}{141.105}{§9,}%
\TextAt{8}{213.133}{141.105}{where}%
\TextAt{8}{247.875}{141.105}{we}%
\TextAt{8}{266.036}{141.105}{studied,}%
\TextAt{8}{311.484}{141.105}{instead}%
\TextAt{8}{352.802}{141.105}{of}%
\TextAt{8}{367.065}{141.105}{valuations}%
\TextAt{13}{423.685}{141.105}{v}%
\TextAt{8}{429.773}{141.105}{,}%
\TextAt{8}{438.089}{141.105}{the}%
\TextAt{8}{459.179}{141.105}{prime}%
\TextAt{8}{493.275}{141.105}{ideals}%
\TextAt{14}{526.786}{141.105}{p}%
\TextAt{8}{79.226}{124.059}{and}%
\TextAt{8}{102.327}{124.059}{their}%
\TextAt{8}{130.629}{124.059}{decomposition}%
\TextAt{14}{208.103}{124.059}{p}%
\TextAt{9}{217.984}{124.059}{\Character{61}}%
\TextAt{14}{230.986}{124.059}{P}%
\TextAt{21}{240.88}{129.284}{e}%
\TextAt{24}{244.808}{128.177}{1}%
\TextAt{17}{240.88}{121.128}{1}%
\TextAt{16}{251.45}{124.059}{\Character{1}}%
\TextAt{16}{256.746}{124.059}{\Character{1}}%
\TextAt{16}{262.054}{124.059}{\Character{1}}%
\TextAt{14}{267.39}{124.059}{P}%
\TextAt{21}{277.283}{128.398}{e}%
\TextAt{22}{281.212}{127.401}{r}%
\TextAt{21}{277.283}{121.104}{r}%
\TextAt{8}{290.104}{124.059}{in}%
\TextAt{8}{304.095}{124.059}{Galois}%
\TextAt{8}{340.347}{124.059}{extensions}%
\TextAt{8}{397.067}{124.059}{of}%
\TextAt{8}{410.735}{124.059}{algebraic}%
\TextAt{8}{460.509}{124.059}{number}%
\TextAt{8}{503.443}{124.059}{fields.}%
\TextAt{8}{79.226}{107.013}{The}%
\TextAt{8}{104.108}{107.013}{arguments}%
\TextAt{8}{161.915}{107.013}{stay}%
\TextAt{8}{187.527}{107.013}{the}%
\TextAt{8}{208.511}{107.013}{same,}%
\TextAt{8}{242.125}{107.013}{so}%
\TextAt{8}{257.323}{107.013}{we}%
\TextAt{8}{275.379}{107.013}{may}%
\TextAt{8}{301.576}{107.013}{be}%
\TextAt{8}{318.329}{107.013}{rather}%
\TextAt{8}{354.282}{107.013}{brief}%
\TextAt{8}{382.093}{107.013}{here.}%
\TextAt{8}{414.486}{107.013}{However,}%
\TextAt{8}{465.925}{107.013}{we}%
\TextAt{8}{483.98}{107.013}{formulate}%
\TextAt{8}{79.226}{89.967}{and}%
\TextAt{8}{102.896}{89.967}{prove}%
\TextAt{8}{135.352}{89.967}{all}%
\TextAt{8}{152.518}{89.967}{results}%
\TextAt{8}{190.617}{89.967}{for}%
\TextAt{8}{209.409}{89.967}{extensions}%
\TextAt{8}{266.696}{89.967}{that}%
\TextAt{8}{292.972}{89.967}{are}%
\TextAt{8}{313.39}{89.967}{not}%
\TextAt{8}{335.11}{89.967}{necessarily}%
\TextAt{8}{394.346}{89.967}{finite,}%
\TextAt{8}{428.644}{89.967}{using}%
\TextAt{8}{460.179}{89.967}{infinite}%
\TextAt{8}{500.754}{89.967}{Galois}%
\TextAt{8}{79.226}{72.921}{theory.}%
\TextAt{8}{121.261}{72.921}{The}%
\TextAt{8}{145.886}{72.921}{reader}%
\TextAt{8}{182.227}{72.921}{who}%
\TextAt{8}{207.51}{72.921}{happens}%
\TextAt{8}{253.989}{72.921}{not}%
\TextAt{8}{275.374}{72.921}{to}%
\TextAt{8}{290.256}{72.921}{know}%
\TextAt{8}{321.396}{72.921}{this}%
\TextAt{8}{344.79}{72.921}{theory}%
\TextAt{8}{382.111}{72.921}{should}%
\TextAt{8}{419.815}{72.921}{feel}%
\TextAt{8}{441.511}{72.921}{free}%
\TextAt{8}{464.511}{72.921}{to}%
\TextAt{8}{479.392}{72.921}{assume}%
\TextAt{8}{520.408}{72.921}{all}%
\TextAt{8}{297.211}{43.033}{152}%
\Destination{79.23}{773.075}{nav.Item.53}%
\Destination{79.23}{756.445}{nav.Item.54}%
\Destination{79.23}{675.895}{nav.Item.55}%
\Destination{78.23}{809.815}{nav.page.152}%
\Destination{79.23}{233.485}{nav.section.2.9}%
\Destination{79.226}{568.755}{nav.ex.2.8.2}%
\Destination{79.226}{537.455}{nav.ex.2.8.3}%
\Destination{79.226}{489.055}{nav.ex.2.8.4}%
\Destination{79.226}{372.955}{nav.ex.2.8.5}%
\Destination{79.226}{324.555}{nav.ex.2.8.6}%
\Destination{79.36}{759.055}{nav.print.166}%
\Destination{79.23}{233.485}{outline.29}%
\Hotlink{524.797}{356.755}{5.22}{11.047}{nav.chapter.1}%
\Hotlink{180.478}{138.955}{5.22}{11.047}{nav.chapter.1}%
\Hotlink{198.459}{138.955}{6.858}{11.047}{nav.section.1.9}%
\Hotlink{289.455}{356.755}{25.078}{11.047}{nav.number.2.8.2}%
\Hotlink{78.726}{339.655}{25.078}{11.047}{nav.number.2.8.3}%
}

% Page 165 (body page 153)
\Page{%
\TextAt{2}{79.226}{796.943}{II.}%
\TextAt{2}{94.073}{796.943}{The}%
\TextAt{2}{116.492}{796.943}{Theory}%
\TextAt{2}{154.401}{796.943}{of}%
\TextAt{2}{166.826}{796.943}{Valuations}%
\Vector{q 0 G 0 g 0.299 w 7.226 18.729 m 460.769 18.729 l S Q}%
\TextAt{8}{79.226}{761.124}{extensions}%
\TextAt{8}{135.598}{761.124}{in}%
\TextAt{8}{149.241}{761.124}{this}%
\TextAt{8}{172.055}{761.124}{section}%
\TextAt{8}{211.127}{761.124}{to}%
\TextAt{8}{225.429}{761.124}{be}%
\TextAt{8}{241.344}{761.124}{finite.}%
\TextAt{8}{275.804}{761.124}{On}%
\TextAt{8}{295.305}{761.124}{the}%
\TextAt{8}{315.453}{761.124}{other}%
\TextAt{8}{346.013}{761.124}{hand,}%
\TextAt{8}{378.524}{761.124}{we}%
\TextAt{8}{395.742}{761.124}{treat}%
\TextAt{8}{424.354}{761.124}{infinite}%
\TextAt{8}{464.012}{761.124}{Galois}%
\TextAt{8}{499.917}{761.124}{theory}%
\TextAt{8}{79.226}{744.078}{also}%
\TextAt{8}{102.706}{744.078}{in}%
\TextAt{8}{116.359}{744.078}{chap.}%
\TextAt{8}{147.251}{744.078}{IV,}%
\TextAt{8}{167.396}{744.078}{§1}%
\TextAt{8}{182.83}{744.078}{below.}%
\TextAt{8}{220.544}{744.078}{Its}%
\TextAt{8}{237.831}{744.078}{main}%
\TextAt{8}{267.098}{744.078}{result}%
\TextAt{8}{299.676}{744.078}{can}%
\TextAt{8}{321.135}{744.078}{be}%
\TextAt{8}{337.059}{744.078}{put}%
\TextAt{8}{358.519}{744.078}{in}%
\TextAt{8}{372.172}{744.078}{a}%
\TextAt{8}{381.927}{744.078}{nutshell}%
\TextAt{8}{425.886}{744.078}{like}%
\TextAt{8}{447.346}{744.078}{this:}%
\TextAt{8}{97.159}{724.043}{In}%
\TextAt{8}{110.999}{724.043}{the}%
\TextAt{8}{130.375}{724.043}{case}%
\TextAt{8}{154.366}{724.043}{of}%
\TextAt{8}{166.915}{724.043}{a}%
\TextAt{8}{175.89}{724.043}{Galois}%
\TextAt{8}{211.023}{724.043}{extension}%
\TextAt{13}{262.008}{724.043}{L}%
\TextAt{16}{273.293}{724.043}{\Character{106}}%
\TextAt{13}{279.935}{724.043}{K}%
\TextAt{8}{293.869}{724.043}{of}%
\TextAt{8}{306.419}{724.043}{infinite}%
\TextAt{8}{345.305}{724.043}{degree,}%
\TextAt{8}{384.348}{724.043}{the}%
\TextAt{8}{403.724}{724.043}{main}%
\TextAt{8}{432.21}{724.043}{theorem}%
\TextAt{8}{476.954}{724.043}{of}%
\TextAt{8}{489.504}{724.043}{ordinary}%
\TextAt{8}{79.226}{706.997}{Galois}%
\TextAt{8}{114.344}{706.997}{theory,}%
\TextAt{8}{152.73}{706.997}{concerning}%
\TextAt{8}{210.467}{706.997}{the}%
\TextAt{8}{229.828}{706.997}{1–1}%
\TextAt{8}{250.505}{706.997}{correspondence}%
\TextAt{8}{331.387}{706.997}{between}%
\TextAt{8}{375.782}{706.997}{the}%
\TextAt{8}{395.144}{706.997}{intermediate}%
\TextAt{8}{462.314}{706.997}{fields}%
\TextAt{8}{491.49}{706.997}{of}%
\TextAt{13}{504.025}{706.997}{L}%
\TextAt{16}{515.31}{706.997}{\Character{106}}%
\TextAt{13}{521.952}{706.997}{K}%
\TextAt{8}{79.226}{689.951}{and}%
\TextAt{8}{101.498}{689.951}{the}%
\TextAt{8}{121.163}{689.951}{subgroups}%
\TextAt{8}{176.084}{689.951}{of}%
\TextAt{8}{188.923}{689.951}{the}%
\TextAt{8}{208.589}{689.951}{Galois}%
\TextAt{8}{244.011}{689.951}{group}%
\TextAt{13}{276.696}{689.951}{G}%
\TextAt{9}{285.929}{689.951}{\Character{40}}%
\TextAt{13}{290.482}{689.951}{L}%
\TextAt{16}{301.767}{689.951}{\Character{106}}%
\TextAt{13}{308.409}{689.951}{K}%
\TextAt{9}{319.226}{689.951}{\Character{41}}%
\TextAt{8}{327.185}{689.951}{ceases}%
\TextAt{8}{361.28}{689.951}{to}%
\TextAt{8}{375.1}{689.951}{hold;}%
\TextAt{8}{404.039}{689.951}{there}%
\TextAt{8}{433.46}{689.951}{are}%
\TextAt{8}{452.479}{689.951}{more}%
\TextAt{8}{481.255}{689.951}{subgroups}%
\TextAt{8}{79.226}{672.905}{than}%
\TextAt{8}{106.478}{672.905}{intermediate}%
\TextAt{8}{174.377}{672.905}{fields.}%
\TextAt{8}{208.878}{672.905}{The}%
\TextAt{8}{232.866}{672.905}{correspondence}%
\TextAt{8}{314.478}{672.905}{can}%
\TextAt{8}{335.872}{672.905}{be}%
\TextAt{8}{351.73}{672.905}{salvaged,}%
\TextAt{8}{401.507}{672.905}{however,}%
\TextAt{8}{449.585}{672.905}{by}%
\TextAt{8}{465.779}{672.905}{considering}%
\TextAt{8}{526.911}{672.905}{a}%
\TextAt{8}{79.226}{655.859}{canonical}%
\TextAt{8}{130.632}{655.859}{topology}%
\TextAt{8}{178.799}{655.859}{on}%
\TextAt{8}{195.081}{655.859}{the}%
\TextAt{8}{215.26}{655.859}{group}%
\TextAt{13}{248.459}{655.859}{G}%
\TextAt{9}{257.692}{655.859}{\Character{40}}%
\TextAt{13}{262.244}{655.859}{L}%
\TextAt{16}{273.569}{655.859}{\Character{106}}%
\TextAt{13}{280.25}{655.859}{K}%
\TextAt{9}{291.067}{655.859}{\Character{41}}%
\TextAt{8}{295.62}{655.859}{,}%
\TextAt{8}{302.797}{655.859}{the}%
\TextAt{12}{322.977}{655.859}{Krull}%
\TextAt{12}{358.474}{655.859}{topology}%
\TextAt{8}{409.319}{655.859}{.}%
\TextAt{8}{417.836}{655.859}{It}%
\TextAt{8}{430.532}{655.859}{is}%
\TextAt{8}{442.319}{655.859}{given}%
\TextAt{8}{472.911}{655.859}{by}%
\TextAt{8}{489.193}{655.859}{defining,}%
\TextAt{8}{79.226}{638.813}{for}%
\TextAt{8}{98.536}{638.813}{every}%
\TextAt{13}{130.852}{638.813}{\Character{27}}%
\NativeMathAt{143.683}{638.813}{11.955}{\in}%
\TextAt{13}{157.401}{638.813}{G}%
\TextAt{9}{166.635}{638.813}{\Character{40}}%
\TextAt{13}{171.187}{638.813}{L}%
\TextAt{16}{184.9}{638.813}{\Character{106}}%
\TextAt{13}{193.969}{638.813}{K}%
\TextAt{9}{204.787}{638.813}{\Character{41}}%
\TextAt{8}{209.339}{638.813}{,}%
\TextAt{8}{218.269}{638.813}{as}%
\TextAt{8}{234.063}{638.813}{a}%
\TextAt{8}{245.243}{638.813}{basis}%
\TextAt{8}{275.408}{638.813}{of}%
\TextAt{8}{290.162}{638.813}{neighbourhoods}%
\TextAt{8}{376.206}{638.813}{the}%
\TextAt{8}{397.786}{638.813}{cosets}%
\TextAt{13}{433.152}{638.813}{\Character{27}}%
\TextAt{13}{440.218}{638.813}{G}%
\TextAt{9}{449.468}{638.813}{\Character{40}}%
\TextAt{13}{454.02}{638.813}{L}%
\TextAt{16}{467.733}{638.813}{\Character{106}}%
\TextAt{13}{476.802}{638.813}{M}%
\TextAt{9}{489.376}{638.813}{\Character{41}}%
\TextAt{8}{493.928}{638.813}{,}%
\TextAt{8}{502.858}{638.813}{where}%
\TextAt{13}{79.226}{621.767}{M}%
\TextAt{16}{96.892}{621.767}{\Character{106}}%
\TextAt{13}{105.306}{621.767}{K}%
\TextAt{8}{121.061}{621.767}{varies}%
\TextAt{8}{155.013}{621.767}{over}%
\TextAt{8}{181.099}{621.767}{the}%
\TextAt{8}{202.295}{621.767}{finite}%
\TextAt{8}{233.246}{621.767}{Galois}%
\TextAt{8}{270.199}{621.767}{subextensions}%
\TextAt{8}{345.564}{621.767}{of}%
\TextAt{13}{359.934}{621.767}{L}%
\TextAt{16}{372.991}{621.767}{\Character{106}}%
\TextAt{13}{381.405}{621.767}{K}%
\TextAt{8}{392.222}{621.767}{.}%
\TextAt{13}{403.79}{621.767}{G}%
\TextAt{9}{413.024}{621.767}{\Character{40}}%
\TextAt{13}{417.576}{621.767}{L}%
\TextAt{16}{430.633}{621.767}{\Character{106}}%
\TextAt{13}{439.046}{621.767}{K}%
\TextAt{9}{449.864}{621.767}{\Character{41}}%
\TextAt{8}{459.354}{621.767}{is}%
\TextAt{8}{472.157}{621.767}{thus}%
\TextAt{8}{498.948}{621.767}{turned}%
\TextAt{8}{79.226}{604.721}{into}%
\TextAt{8}{103.792}{604.721}{a}%
\TextAt{8}{114.369}{604.721}{compact,}%
\TextAt{8}{165.478}{604.721}{Hausdorff}%
\TextAt{8}{219.692}{604.721}{topological}%
\TextAt{8}{280.039}{604.721}{group.}%
\TextAt{8}{320.233}{604.721}{The}%
\TextAt{8}{345.109}{604.721}{main}%
\TextAt{8}{375.198}{604.721}{theorem}%
\TextAt{8}{421.546}{604.721}{of}%
\TextAt{8}{435.698}{604.721}{Galois}%
\TextAt{8}{472.434}{604.721}{theory}%
\TextAt{8}{510.007}{604.721}{then}%
\TextAt{8}{79.226}{587.675}{has}%
\TextAt{8}{101.146}{587.675}{to}%
\TextAt{8}{116.503}{587.675}{be}%
\TextAt{8}{133.474}{587.675}{modified}%
\TextAt{8}{182.317}{587.675}{in}%
\TextAt{8}{197.016}{587.675}{the}%
\TextAt{8}{218.219}{587.675}{infinite}%
\TextAt{8}{258.933}{587.675}{case}%
\TextAt{8}{284.751}{587.675}{by}%
\TextAt{8}{302.056}{587.675}{the}%
\TextAt{8}{323.259}{587.675}{condition}%
\TextAt{8}{375.689}{587.675}{that}%
\TextAt{8}{402.104}{587.675}{the}%
\TextAt{8}{423.307}{587.675}{intermediate}%
\TextAt{8}{492.319}{587.675}{fields}%
\TextAt{8}{523.337}{587.675}{of}%
\TextAt{13}{79.226}{570.629}{L}%
\TextAt{16}{91.252}{570.629}{\Character{106}}%
\TextAt{13}{98.635}{570.629}{K}%
\TextAt{8}{113.785}{570.629}{correspond}%
\TextAt{8}{173.792}{570.629}{1–1}%
\TextAt{8}{195.698}{570.629}{to}%
\TextAt{8}{210.443}{570.629}{the}%
\TextAt{19}{231.034}{570.629}{closed}%
\TextAt{8}{265.852}{570.629}{subgroups}%
\TextAt{8}{321.699}{570.629}{of}%
\TextAt{13}{335.464}{570.629}{G}%
\TextAt{9}{344.697}{570.629}{\Character{40}}%
\TextAt{13}{349.25}{570.629}{L}%
\TextAt{16}{361.276}{570.629}{\Character{106}}%
\TextAt{13}{368.658}{570.629}{K}%
\TextAt{9}{379.476}{570.629}{\Character{41}}%
\TextAt{8}{384.028}{570.629}{.}%
\TextAt{8}{393.781}{570.629}{Otherwise,}%
\TextAt{8}{452.917}{570.629}{everything}%
\TextAt{8}{510.915}{570.629}{goes}%
\TextAt{8}{79.226}{553.583}{through}%
\TextAt{8}{123.808}{553.583}{as}%
\TextAt{8}{138.525}{553.583}{in}%
\TextAt{8}{152.526}{553.583}{the}%
\TextAt{8}{173.03}{553.583}{finite}%
\TextAt{8}{203.289}{553.583}{case.}%
\TextAt{8}{233.653}{553.583}{So}%
\TextAt{8}{250.26}{553.583}{one}%
\TextAt{8}{272.067}{553.583}{tacitly}%
\TextAt{8}{309.165}{553.583}{restricts}%
\TextAt{8}{354.512}{553.583}{attention}%
\TextAt{8}{405.274}{553.583}{to}%
\TextAt{8}{419.932}{553.583}{closed}%
\TextAt{8}{454.806}{553.583}{subgroups,}%
\TextAt{8}{513.904}{553.583}{and}%
\TextAt{8}{79.226}{536.537}{accordingly}%
\TextAt{8}{141.345}{536.537}{to}%
\TextAt{19}{155.656}{536.537}{continuous}%
\TextAt{8}{214.379}{536.537}{homomorphisms}%
\TextAt{8}{301.664}{536.537}{of}%
\TextAt{13}{314.994}{536.537}{G}%
\TextAt{9}{324.228}{536.537}{\Character{40}}%
\TextAt{13}{328.78}{536.537}{L}%
\TextAt{16}{340.065}{536.537}{\Character{106}}%
\TextAt{13}{346.707}{536.537}{K}%
\TextAt{9}{357.525}{536.537}{\Character{41}}%
\TextAt{8}{362.077}{536.537}{.}%
\TextAt{8}{97.159}{516.503}{So}%
\TextAt{8}{114.536}{516.503}{let}%
\TextAt{13}{132.559}{516.503}{L}%
\TextAt{16}{145.75}{516.502}{\Character{106}}%
\TextAt{13}{154.298}{516.503}{K}%
\TextAt{8}{170.131}{516.503}{be}%
\TextAt{8}{187.174}{516.503}{an}%
\TextAt{8}{204.551}{516.503}{arbitrary,}%
\TextAt{8}{257.99}{516.503}{finite}%
\TextAt{8}{289.02}{516.503}{or}%
\TextAt{8}{304.449}{516.503}{infinite,}%
\TextAt{8}{348.765}{516.503}{Galois}%
\TextAt{8}{385.797}{516.503}{extension}%
\TextAt{8}{438.681}{516.503}{with}%
\TextAt{8}{466.46}{516.503}{Galois}%
\TextAt{8}{503.491}{516.503}{group}%
\TextAt{13}{79.226}{499.457}{G}%
\TextAt{9}{93.209}{499.457}{\Character{61}}%
\TextAt{13}{107.063}{499.457}{G}%
\TextAt{9}{116.297}{499.457}{\Character{40}}%
\TextAt{13}{120.849}{499.457}{L}%
\TextAt{16}{133.563}{499.457}{\Character{106}}%
\TextAt{13}{141.633}{499.457}{K}%
\TextAt{9}{152.45}{499.457}{\Character{41}}%
\TextAt{8}{157.003}{499.457}{.}%
\TextAt{8}{167.966}{499.457}{If}%
\TextAt{13}{180.496}{499.457}{v}%
\TextAt{8}{191.32}{499.457}{is}%
\TextAt{8}{203.922}{499.457}{an}%
\TextAt{8}{221.019}{499.457}{(archimedean}%
\TextAt{8}{294.378}{499.457}{or}%
\TextAt{8}{309.526}{499.457}{nonarchimedean)}%
\TextAt{8}{401.75}{499.457}{valuation}%
\TextAt{8}{453.66}{499.457}{of}%
\TextAt{13}{467.829}{499.457}{K}%
\TextAt{8}{483.382}{499.457}{and}%
\TextAt{13}{506.983}{499.457}{w}%
\TextAt{8}{520.408}{499.457}{an}%
\TextAt{8}{79.226}{482.411}{extension}%
\TextAt{8}{131.105}{482.411}{to}%
\TextAt{13}{145.528}{482.411}{L}%
\TextAt{8}{153.492}{482.411}{,}%
\TextAt{8}{160.782}{482.411}{then,}%
\TextAt{8}{190.835}{482.411}{for}%
\TextAt{8}{208.832}{482.411}{every}%
\TextAt{13}{239.837}{482.411}{\Character{27}}%
\NativeMathAt{250.432}{482.411}{11.955}{\in}%
\TextAt{13}{261.915}{482.411}{G}%
\TextAt{8}{271.149}{482.411}{,}%
\TextAt{13}{278.439}{482.411}{w}%
\TextAt{16}{289.862}{482.411}{\Character{14}}%
\TextAt{13}{298.573}{482.411}{\Character{27}}%
\TextAt{8}{309.665}{482.411}{also}%
\TextAt{8}{333.258}{482.411}{extends}%
\TextAt{13}{376.026}{482.411}{v}%
\TextAt{8}{382.115}{482.411}{,}%
\TextAt{8}{389.405}{482.411}{so}%
\TextAt{8}{403.887}{482.411}{that}%
\TextAt{8}{429.369}{482.411}{the}%
\TextAt{8}{449.638}{482.411}{group}%
\TextAt{13}{482.926}{482.411}{G}%
\TextAt{8}{496.17}{482.411}{acts}%
\TextAt{8}{520.408}{482.411}{on}%
\TextAt{8}{79.226}{465.365}{the}%
\TextAt{8}{99.383}{465.365}{set}%
\TextAt{8}{117.65}{465.365}{of}%
\TextAt{8}{130.98}{465.365}{extensions}%
\TextAt{13}{187.361}{465.365}{w}%
\TextAt{16}{199.371}{465.365}{\Character{106}}%
\TextAt{13}{206.013}{465.365}{v}%
\TextAt{8}{212.101}{465.365}{.}%
\TextAt{12}{79.226}{439.793}{Proposition}%
\TextAt{12}{152.964}{439.793}{(9.1).}%
\TextAt{8}{189.173}{439.793}{The}%
\TextAt{8}{214.119}{439.793}{group}%
\TextAt{13}{248.188}{439.793}{G}%
\TextAt{8}{262.212}{439.793}{acts}%
\TextAt{8}{287.23}{439.793}{transitively}%
\TextAt{8}{349.656}{439.793}{on}%
\TextAt{8}{366.808}{439.793}{the}%
\TextAt{8}{387.857}{439.793}{set}%
\TextAt{8}{407.017}{439.793}{of}%
\TextAt{8}{421.24}{439.793}{extensions}%
\TextAt{13}{478.513}{439.793}{w}%
\TextAt{16}{492.045}{439.793}{\Character{106}}%
\TextAt{13}{500.208}{439.793}{v}%
\TextAt{8}{506.297}{439.793}{,}%
\TextAt{8}{514.562}{439.793}{i.e.,}%
\TextAt{8}{79.226}{422.747}{every}%
\TextAt{8}{110.118}{422.747}{two}%
\TextAt{8}{132.235}{422.747}{extensions}%
\TextAt{8}{188.616}{422.747}{are}%
\TextAt{8}{208.127}{422.747}{conjugate.}%
\TextAt{19}{79.226}{388.651}{Proof.}%
\TextAt{8}{116}{388.651}{Let}%
\TextAt{13}{136.618}{388.651}{w}%
\TextAt{8}{148.854}{388.651}{and}%
\TextAt{13}{171.265}{388.651}{w}%
\TextAt{3}{179.955}{392.989}{\Character{48}}%
\TextAt{8}{186.296}{388.651}{be}%
\TextAt{8}{201.869}{388.651}{two}%
\TextAt{8}{223.635}{388.651}{extensions}%
\TextAt{8}{279.664}{388.651}{of}%
\TextAt{13}{292.643}{388.651}{v}%
\TextAt{8}{302.277}{388.651}{to}%
\TextAt{13}{316.236}{388.651}{L}%
\TextAt{8}{324.201}{388.651}{.}%
\TextAt{8}{332.532}{388.651}{Suppose}%
\TextAt{13}{378.088}{388.651}{L}%
\TextAt{16}{389.374}{388.651}{\Character{106}}%
\TextAt{13}{396.015}{388.651}{K}%
\TextAt{8}{410.379}{388.651}{is}%
\TextAt{8}{421.792}{388.651}{finite.}%
\TextAt{8}{456.137}{388.651}{If}%
\TextAt{13}{467.478}{388.651}{w}%
\TextAt{8}{479.714}{388.651}{and}%
\TextAt{13}{502.125}{388.651}{w}%
\TextAt{3}{510.815}{392.989}{\Character{48}}%
\TextAt{8}{517.156}{388.651}{are}%
\TextAt{8}{79.226}{371.605}{not}%
\TextAt{8}{100.04}{371.605}{conjugate,}%
\TextAt{8}{156.301}{371.605}{then}%
\TextAt{8}{182.961}{371.605}{the}%
\TextAt{8}{203.118}{371.605}{sets}%
\TextAt{16}{191.349}{344.147}{\Character{102}}%
\TextAt{13}{197.326}{344.147}{w}%
\TextAt{16}{208.673}{344.147}{\Character{14}}%
\TextAt{13}{217.307}{344.147}{\Character{27}}%
\TextAt{16}{227.71}{344.147}{\Character{106}}%
\TextAt{13}{234.352}{344.147}{\Character{27}}%
\NativeMathAt{244.755}{344.147}{11.955}{\in}%
\TextAt{13}{256.046}{344.147}{G}%
\TextAt{16}{265.28}{344.147}{\Character{103}}%
\TextAt{8}{295.168}{344.147}{and}%
\TextAt{16}{337.943}{344.147}{\Character{102}}%
\TextAt{13}{343.921}{344.147}{w}%
\TextAt{3}{352.61}{349.083}{\Character{48}}%
\TextAt{16}{358.062}{344.147}{\Character{14}}%
\TextAt{13}{366.696}{344.147}{\Character{27}}%
\TextAt{16}{377.1}{344.147}{\Character{106}}%
\TextAt{13}{383.741}{344.147}{\Character{27}}%
\NativeMathAt{394.145}{344.147}{11.955}{\in}%
\TextAt{13}{405.436}{344.147}{G}%
\TextAt{16}{414.669}{344.147}{\Character{103}}%
\TextAt{8}{79.226}{316.689}{would}%
\TextAt{8}{112.569}{316.689}{be}%
\TextAt{8}{127.693}{316.689}{disjoint.}%
\TextAt{8}{173.665}{316.689}{By}%
\TextAt{8}{191.227}{316.689}{the}%
\TextAt{8}{210.583}{316.689}{approximation}%
\TextAt{8}{287.85}{316.689}{theorem}%
\TextAt{8}{332.574}{316.689}{(3.4),}%
\TextAt{8}{363.161}{316.689}{we}%
\TextAt{8}{379.588}{316.689}{would}%
\TextAt{8}{412.931}{316.689}{be}%
\TextAt{8}{428.055}{316.689}{able}%
\TextAt{8}{451.965}{316.689}{to}%
\TextAt{8}{465.475}{316.689}{find}%
\TextAt{8}{488.083}{316.689}{an}%
\TextAt{13}{503.541}{316.689}{x}%
\NativeMathAt{513.514}{316.689}{11.955}{\in}%
\TextAt{13}{524.805}{316.689}{L}%
\TextAt{8}{79.226}{299.643}{such}%
\TextAt{8}{105.623}{299.643}{that}%
\TextAt{16}{221.551}{282.597}{\Character{106}}%
\TextAt{13}{224.872}{282.597}{\Character{27}}%
\TextAt{13}{231.938}{282.597}{x}%
\TextAt{16}{238.607}{282.597}{\Character{106}}%
\TextAt{21}{241.928}{279.11}{w}%
\TextAt{13}{252.038}{282.597}{\Character{60}}%
\TextAt{9}{264.463}{282.597}{1}%
\TextAt{8}{294.227}{282.597}{and}%
\TextAt{16}{338.995}{282.597}{\Character{106}}%
\TextAt{13}{342.316}{282.597}{\Character{27}}%
\TextAt{13}{349.382}{282.597}{x}%
\TextAt{16}{356.05}{282.597}{\Character{106}}%
\TextAt{21}{359.371}{279.11}{w}%
\TextAt{23}{365.662}{281.38}{\Character{48}}%
\TextAt{13}{372.166}{282.597}{\Character{62}}%
\TextAt{9}{384.591}{282.597}{1}%
\TextAt{8}{79.226}{259.737}{for}%
\TextAt{8}{97.58}{259.737}{all}%
\TextAt{13}{114.307}{259.737}{\Character{27}}%
\NativeMathAt{125.509}{259.737}{11.955}{\in}%
\TextAt{13}{137.598}{259.737}{G}%
\TextAt{8}{146.832}{259.737}{.}%
\TextAt{8}{156.686}{259.737}{Then}%
\TextAt{8}{187.712}{259.737}{one}%
\TextAt{8}{209.64}{259.737}{would}%
\TextAt{8}{244.253}{259.737}{have}%
\TextAt{8}{271.716}{259.737}{for}%
\TextAt{8}{290.07}{259.737}{the}%
\TextAt{8}{310.695}{259.737}{norm}%
\TextAt{13}{341.732}{259.737}{\Character{11}}%
\TextAt{9}{353.373}{259.737}{\Character{61}}%
\TextAt{13}{366.597}{259.737}{N}%
\TextAt{21}{375.971}{257.877}{L}%
\TextAt{3}{381.73}{257.877}{\Character{106}}%
\TextAt{21}{384.083}{257.877}{K}%
\TextAt{9}{392.285}{259.737}{\Character{40}}%
\TextAt{13}{396.838}{259.737}{x}%
\TextAt{9}{403.49}{259.737}{\Character{41}}%
\TextAt{9}{412.146}{259.737}{\Character{61}}%
\TextAt{20}{425.385}{268.704}{\Character{81}}%
\TextAt{21}{436.676}{256.25}{\Character{27}}%
\NativeMathAt{441.802}{256.25}{7.97}{\in}%
\TextAt{21}{447.447}{256.25}{G}%
\TextAt{13}{456.537}{259.737}{\Character{27}}%
\TextAt{13}{463.603}{259.737}{x}%
\TextAt{8}{474.637}{259.737}{that}%
\TextAt{16}{500.474}{259.737}{\Character{106}}%
\TextAt{13}{503.795}{259.737}{\Character{11}}%
\TextAt{16}{511.317}{259.737}{\Character{106}}%
\TextAt{21}{514.638}{256.25}{v}%
\TextAt{9}{523.665}{259.737}{\Character{61}}%
\TextAt{20}{79.226}{251.658}{\Character{81}}%
\TextAt{21}{90.517}{239.204}{\Character{27}}%
\TextAt{16}{98.133}{242.691}{\Character{106}}%
\TextAt{13}{101.454}{242.691}{\Character{27}}%
\TextAt{13}{108.52}{242.691}{x}%
\TextAt{16}{115.188}{242.691}{\Character{106}}%
\TextAt{21}{118.509}{239.204}{w}%
\TextAt{13}{128.62}{242.691}{\Character{60}}%
\TextAt{9}{141.045}{242.691}{1}%
\TextAt{8}{150.795}{242.691}{and}%
\TextAt{8}{173.558}{242.691}{likewise}%
\TextAt{16}{216.537}{242.691}{\Character{106}}%
\TextAt{13}{219.858}{242.691}{\Character{11}}%
\TextAt{16}{227.38}{242.691}{\Character{106}}%
\TextAt{21}{230.7}{239.204}{v}%
\TextAt{13}{238.929}{242.691}{\Character{62}}%
\TextAt{9}{251.354}{242.691}{1}%
\TextAt{8}{257.207}{242.691}{,}%
\TextAt{8}{264.357}{242.691}{a}%
\TextAt{8}{274.112}{242.691}{contradiction.}%
\TextAt{8}{97.159}{222.656}{If}%
\TextAt{13}{109.43}{222.656}{L}%
\TextAt{16}{121.703}{222.656}{\Character{106}}%
\TextAt{13}{129.332}{222.656}{K}%
\TextAt{8}{144.626}{222.656}{is}%
\TextAt{8}{156.969}{222.656}{infinite,}%
\TextAt{8}{200.612}{222.656}{then}%
\TextAt{8}{227.851}{222.656}{we}%
\TextAt{8}{245.658}{222.656}{let}%
\TextAt{13}{263.142}{222.656}{M}%
\TextAt{16}{280.024}{222.656}{\Character{106}}%
\TextAt{13}{287.652}{222.656}{K}%
\TextAt{8}{302.947}{222.656}{vary}%
\TextAt{8}{329.552}{222.656}{over}%
\TextAt{8}{355.178}{222.656}{all}%
\TextAt{8}{372.016}{222.656}{finite}%
\TextAt{8}{402.507}{222.656}{Galois}%
\TextAt{8}{439}{222.656}{subextensions}%
\TextAt{8}{513.904}{222.656}{and}%
\TextAt{8}{79.226}{205.61}{consider}%
\TextAt{8}{125.857}{205.61}{the}%
\TextAt{8}{147.059}{205.61}{sets}%
\TextAt{13}{170.987}{205.61}{X}%
\TextAt{21}{180.702}{203.817}{M}%
\TextAt{9}{195.284}{205.61}{\Character{61}}%
\TextAt{16}{209.491}{205.61}{\Character{102}}%
\TextAt{13}{215.469}{205.61}{\Character{27}}%
\NativeMathAt{227.654}{205.61}{11.955}{\in}%
\TextAt{13}{240.727}{205.61}{G}%
\TextAt{16}{255.063}{205.61}{\Character{106}}%
\TextAt{13}{263.487}{205.61}{w}%
\TextAt{16}{275.545}{205.61}{\Character{14}}%
\TextAt{13}{284.892}{205.61}{\Character{27}}%
\TextAt{16}{291.975}{205.61}{\Character{106}}%
\TextAt{21}{295.296}{203.817}{M}%
\TextAt{9}{309.878}{205.61}{\Character{61}}%
\TextAt{13}{324.085}{205.61}{w}%
\TextAt{3}{332.775}{209.949}{\Character{48}}%
\TextAt{16}{335.57}{205.61}{\Character{106}}%
\TextAt{21}{338.891}{203.817}{M}%
\TextAt{16}{348.371}{205.61}{\Character{103}}%
\TextAt{8}{354.348}{205.61}{.}%
\TextAt{8}{365.933}{205.61}{They}%
\TextAt{8}{397.214}{205.61}{are}%
\TextAt{8}{417.77}{205.61}{nonempty,}%
\TextAt{8}{475.984}{205.61}{as}%
\TextAt{8}{491.399}{205.61}{we}%
\TextAt{8}{509.672}{205.61}{have}%
\TextAt{8}{79.226}{188.564}{just}%
\TextAt{8}{104.001}{188.564}{seen,}%
\TextAt{8}{134.707}{188.564}{and}%
\TextAt{8}{159.099}{188.564}{also}%
\TextAt{8}{184.209}{188.564}{closed}%
\TextAt{8}{220.365}{188.564}{because,}%
\TextAt{8}{268.956}{188.564}{for}%
\TextAt{13}{288.47}{188.564}{\Character{27}}%
\NativeMathAt{301.651}{188.564}{11.955}{\in}%
\TextAt{13}{315.719}{188.564}{G}%
\TextAt{18}{328.72}{188.564}{r}%
\TextAt{13}{341.786}{188.564}{X}%
\TextAt{21}{351.502}{186.771}{M}%
\TextAt{8}{360.981}{188.564}{,}%
\TextAt{8}{370.168}{188.564}{the}%
\TextAt{8}{391.954}{188.564}{whole}%
\TextAt{8}{426.747}{188.564}{open}%
\TextAt{8}{456.663}{188.564}{neighbourhood}%
\TextAt{13}{79.226}{171.518}{\Character{27}}%
\TextAt{13}{86.292}{171.518}{G}%
\TextAt{9}{95.542}{171.519}{\Character{40}}%
\TextAt{13}{100.095}{171.518}{L}%
\TextAt{16}{111.522}{171.519}{\Character{106}}%
\TextAt{13}{118.305}{171.518}{M}%
\TextAt{9}{130.879}{171.519}{\Character{41}}%
\TextAt{8}{139.412}{171.518}{lies}%
\TextAt{8}{159.711}{171.518}{in}%
\TextAt{8}{173.448}{171.518}{the}%
\TextAt{8}{193.687}{171.518}{complement}%
\TextAt{8}{259.129}{171.518}{of}%
\TextAt{13}{272.543}{171.518}{X}%
\TextAt{21}{282.258}{169.725}{M}%
\TextAt{8}{291.738}{171.518}{.}%
\TextAt{8}{300.436}{171.518}{We}%
\TextAt{8}{320.664}{171.518}{have}%
\TextAt{20}{347.742}{180.485}{\Character{84}}%
\TextAt{21}{357.704}{168.032}{M}%
\TextAt{13}{369.177}{171.518}{X}%
\TextAt{21}{378.892}{169.725}{M}%
\NativeMathAt{391.834}{171.519}{11.955}{\neq}%
\TextAt{18}{404.402}{171.519}{\Character{63}}%
\TextAt{8}{413.7}{171.518}{,}%
\TextAt{8}{420.954}{171.518}{because}%
\TextAt{8}{464.339}{171.518}{otherwise}%
\TextAt{8}{516.51}{171.518}{the}%
\TextAt{8}{79.226}{154.473}{compactness}%
\TextAt{8}{147.359}{154.473}{of}%
\TextAt{13}{161.06}{154.473}{G}%
\TextAt{8}{174.561}{154.473}{would}%
\TextAt{8}{209.076}{154.473}{yield}%
\TextAt{8}{237.733}{154.473}{a}%
\TextAt{8}{247.859}{154.473}{relation}%
\TextAt{20}{291.161}{163.439}{\Character{84}}%
\TextAt{21}{301.124}{160.284}{r}%
\TextAt{21}{301.124}{150.986}{i}%
\TextAt{17}{304.007}{150.986}{\Character{61}1}%
\TextAt{13}{317.318}{154.473}{X}%
\TextAt{21}{327.033}{152.679}{M}%
\TextAt{22}{335.166}{151.464}{i}%
\TextAt{9}{342.779}{154.473}{\Character{61}}%
\TextAt{18}{355.836}{154.473}{\Character{63}}%
\TextAt{8}{369.403}{154.473}{with}%
\TextAt{8}{396.434}{154.473}{finitely}%
\TextAt{8}{436.149}{154.473}{many}%
\TextAt{13}{468.392}{154.473}{M}%
\TextAt{21}{479.717}{152.679}{i}%
\TextAt{8}{483.098}{154.473}{,}%
\TextAt{8}{490.711}{154.473}{and}%
\TextAt{8}{513.844}{154.473}{this}%
\ErrataAt{79.22}{137.455}{453.64}{is a contradiction because for $M=M_1\cdots M_r$ one has $X_M\subseteq\bigcap_{i=1}^r X_{M_i}$.}%
\Vector{q 0 G 0 g 0.399 w 452.334 -636.573 m 452.334 -628.504 l S Q}%
\Vector{q 0 G 0 g 0.399 w 452.534 -628.703 m 459.707 -628.703 l S Q}%
\Vector{q 0 G 0 g 0.399 w 459.906 -636.573 m 459.906 -628.504 l S Q}%
\TextAt{8}{297.211}{43.033}{153}%
\Destination{78.23}{809.815}{nav.page.153}%
\Destination{79.36}{704.955}{nav.print.167}%
\Destination{79.36}{453.755}{nav.number.2.9.1}%
\Hotlink{146.751}{741.855}{13.995}{11.047}{nav.chapter.4}%
\Hotlink{172.575}{741.855}{6.858}{11.047}{nav.section.4.1}%
\Destination{79.226}{151.455}{errata.19}%
\Hotlink{333.735}{313.696}{21.722}{11.047}{nav.number.2.3.4}%
}

% Page 166 (body page 154)
\Page{%
\TextAt{2}{79.226}{796.943}{II.}%
\TextAt{2}{94.073}{796.943}{The}%
\TextAt{2}{116.492}{796.943}{Theory}%
\TextAt{2}{154.401}{796.943}{of}%
\TextAt{2}{166.826}{796.943}{Valuations}%
\Vector{q 0 G 0 g 0.299 w 7.226 18.729 m 460.769 18.729 l S Q}%
\TextAt{12}{79.226}{761.124}{Definition}%
\TextAt{12}{141.782}{761.124}{(9.2).}%
\TextAt{8}{176.895}{761.124}{The}%
\TextAt{12}{200.746}{761.124}{decomposition}%
\TextAt{12}{289.328}{761.124}{group}%
\TextAt{8}{326.927}{761.124}{of}%
\TextAt{8}{340.055}{761.124}{an}%
\TextAt{8}{356.111}{761.124}{extension}%
\TextAt{13}{407.675}{761.124}{w}%
\TextAt{8}{420.059}{761.124}{of}%
\TextAt{13}{433.186}{761.124}{v}%
\TextAt{8}{442.969}{761.124}{to}%
\TextAt{13}{457.077}{761.124}{L}%
\TextAt{8}{468.736}{761.124}{is}%
\TextAt{8}{480.298}{761.124}{defined}%
\TextAt{8}{520.408}{761.124}{by}%
\TextAt{13}{182.581}{732.664}{G}%
\TextAt{21}{191.815}{730.871}{w}%
\TextAt{9}{201.925}{732.664}{\Character{61}}%
\TextAt{13}{214.35}{732.664}{G}%
\TextAt{21}{223.584}{730.871}{w}%
\TextAt{9}{230.374}{732.664}{\Character{40}}%
\TextAt{13}{234.926}{732.664}{L}%
\TextAt{16}{246.211}{732.664}{\Character{106}}%
\TextAt{13}{252.853}{732.664}{K}%
\TextAt{9}{263.671}{732.664}{\Character{41}}%
\TextAt{9}{271.526}{732.664}{\Character{61}}%
\TextAt{16}{283.969}{732.664}{\Character{102}}%
\TextAt{13}{289.947}{732.664}{\Character{27}}%
\NativeMathAt{300.35}{732.664}{11.955}{\in}%
\TextAt{13}{311.641}{732.664}{G}%
\TextAt{9}{320.875}{732.664}{\Character{40}}%
\TextAt{13}{325.427}{732.664}{L}%
\TextAt{16}{336.712}{732.664}{\Character{106}}%
\TextAt{13}{343.354}{732.664}{K}%
\TextAt{9}{354.172}{732.664}{\Character{41}}%
\TextAt{16}{362.045}{732.664}{\Character{106}}%
\TextAt{13}{368.687}{732.664}{w}%
\TextAt{16}{380.033}{732.664}{\Character{14}}%
\TextAt{13}{388.667}{732.664}{\Character{27}}%
\TextAt{9}{399.07}{732.664}{\Character{61}}%
\TextAt{13}{411.496}{732.664}{w}%
\TextAt{16}{420.185}{732.664}{\Character{103}}%
\TextAt{13}{426.163}{732.664}{\Character{58}}%
\TextAt{8}{97.159}{695.394}{If}%
\TextAt{13}{108.059}{695.394}{v}%
\TextAt{8}{117.252}{695.394}{is}%
\TextAt{8}{128.223}{695.394}{a}%
\TextAt{8}{137.186}{695.394}{nonarchimedean}%
\TextAt{8}{223.224}{695.394}{valuation,}%
\TextAt{8}{276.914}{695.394}{then}%
\TextAt{8}{302.781}{695.394}{the}%
\TextAt{8}{322.145}{695.394}{decomposition}%
\TextAt{8}{398.488}{695.394}{group}%
\TextAt{8}{430.871}{695.394}{contains}%
\TextAt{8}{475.998}{695.394}{two}%
\TextAt{8}{497.322}{695.394}{further}%
\TextAt{8}{79.226}{678.348}{canonical}%
\TextAt{8}{130.609}{678.348}{subgroups}%
\TextAt{13}{266.594}{661.302}{G}%
\TextAt{21}{275.828}{659.509}{w}%
\NativeMathAt{285.938}{661.302}{11.955}{\supseteq}%
\TextAt{13}{298.558}{661.302}{I}%
\TextAt{21}{303.72}{659.509}{w}%
\NativeMathAt{313.831}{661.302}{11.955}{\supseteq}%
\TextAt{13}{326.45}{661.302}{R}%
\TextAt{21}{335.36}{659.509}{w}%
\TextAt{13}{342.15}{661.302}{\Character{59}}%
\TextAt{8}{79.226}{638.012}{which}%
\TextAt{8}{111.875}{638.012}{are}%
\TextAt{8}{130.548}{638.012}{defined}%
\TextAt{8}{170.023}{638.012}{as}%
\TextAt{8}{183.556}{638.012}{follows.}%
\TextAt{8}{226.263}{638.012}{Let}%
\TextAt{14}{246.395}{638.012}{o}%
\TextAt{8}{252.239}{638.012}{,}%
\TextAt{8}{258.718}{638.012}{resp.}%
\TextAt{16}{285.903}{638.012}{\Character{79}}%
\TextAt{8}{295.753}{638.012}{,}%
\TextAt{8}{302.232}{638.012}{be}%
\TextAt{8}{317.319}{638.012}{the}%
\TextAt{8}{336.638}{638.012}{valuation}%
\TextAt{8}{386.872}{638.012}{ring,}%
\TextAt{14}{413.52}{638.012}{p}%
\TextAt{8}{419.503}{638.012}{,}%
\TextAt{8}{425.982}{638.012}{resp.}%
\TextAt{14}{453.167}{638.012}{P}%
\TextAt{8}{463.06}{638.012}{,}%
\TextAt{8}{469.539}{638.012}{the}%
\TextAt{8}{488.858}{638.012}{maximal}%
\TextAt{8}{79.226}{620.966}{ideal,}%
\TextAt{8}{110.441}{620.966}{and}%
\TextAt{8}{133.204}{620.966}{let}%
\TextAt{13}{150.108}{620.966}{\Character{20}}%
\TextAt{9}{160.168}{620.966}{\Character{61}}%
\TextAt{14}{172.593}{620.966}{o}%
\TextAt{9}{178.437}{620.966}{\Character{47}}%
\TextAt{14}{184.29}{620.966}{p}%
\TextAt{8}{190.274}{620.966}{,}%
\TextAt{8}{197.423}{620.966}{resp.}%
\TextAt{13}{225.446}{620.966}{\Character{21}}%
\TextAt{9}{235.595}{620.966}{\Character{61}}%
\TextAt{16}{248.02}{620.966}{\Character{79}}%
\TextAt{9}{257.87}{620.966}{\Character{47}}%
\TextAt{14}{263.723}{620.966}{P}%
\TextAt{8}{273.616}{620.966}{,}%
\TextAt{8}{280.766}{620.966}{be}%
\TextAt{8}{296.69}{620.966}{the}%
\TextAt{8}{316.846}{620.966}{residue}%
\TextAt{8}{356.573}{620.966}{class}%
\TextAt{8}{384.011}{620.966}{field}%
\TextAt{8}{409.367}{620.966}{of}%
\TextAt{13}{422.697}{620.966}{v}%
\TextAt{8}{428.786}{620.966}{,}%
\TextAt{8}{435.935}{620.966}{resp.}%
\TextAt{13}{463.958}{620.966}{w}%
\TextAt{8}{472.647}{620.966}{.}%
\TextAt{12}{79.226}{595.109}{Definition}%
\TextAt{12}{142.015}{595.109}{(9.3).}%
\TextAt{8}{177.33}{595.109}{The}%
\TextAt{12}{201.384}{595.109}{inertia}%
\TextAt{12}{244.231}{595.109}{group}%
\TextAt{8}{282.034}{595.109}{of}%
\TextAt{13}{295.364}{595.109}{w}%
\TextAt{16}{307.374}{595.109}{\Character{106}}%
\TextAt{13}{314.016}{595.109}{v}%
\TextAt{8}{324.001}{595.109}{is}%
\TextAt{8}{335.765}{595.109}{defined}%
\TextAt{8}{376.078}{595.109}{by}%
\TextAt{13}{157.49}{566.65}{I}%
\TextAt{21}{162.653}{564.856}{w}%
\TextAt{9}{172.763}{566.65}{\Character{61}}%
\TextAt{13}{185.189}{566.65}{I}%
\TextAt{21}{190.352}{564.856}{w}%
\TextAt{9}{197.141}{566.65}{\Character{40}}%
\TextAt{13}{201.693}{566.65}{L}%
\TextAt{16}{212.979}{566.65}{\Character{106}}%
\TextAt{13}{219.621}{566.65}{K}%
\TextAt{9}{230.438}{566.65}{\Character{41}}%
\TextAt{9}{238.293}{566.65}{\Character{61}}%
\TextAt{16}{250.737}{566.65}{\Character{102}}%
\TextAt{13}{256.714}{566.65}{\Character{27}}%
\NativeMathAt{267.118}{566.65}{11.955}{\in}%
\TextAt{13}{278.409}{566.65}{G}%
\TextAt{21}{287.642}{564.856}{w}%
\TextAt{16}{297.753}{566.65}{\Character{106}}%
\TextAt{13}{304.394}{566.65}{\Character{27}}%
\TextAt{13}{311.459}{566.65}{x}%
\NativeMathAt{321.45}{566.65}{11.955}{\equiv}%
\TextAt{13}{334.069}{566.65}{x}%
\TextAt{8}{344.042}{566.65}{mod}%
\TextAt{14}{369.48}{566.65}{P}%
\TextAt{8}{383.27}{566.65}{for}%
\TextAt{8}{401.155}{566.65}{all}%
\TextAt{13}{417.414}{566.65}{x}%
\NativeMathAt{427.387}{566.65}{11.955}{\in}%
\TextAt{16}{438.661}{566.65}{\Character{79}}%
\TextAt{16}{448.512}{566.65}{\Character{103}}%
\TextAt{8}{79.226}{538.191}{and}%
\TextAt{8}{101.989}{538.191}{the}%
\TextAt{12}{122.145}{538.191}{ramification}%
\TextAt{12}{196.793}{538.191}{group}%
\TextAt{8}{234.596}{538.191}{by}%
\TextAt{13}{146}{509.731}{R}%
\TextAt{21}{154.91}{507.938}{w}%
\TextAt{9}{165.021}{509.731}{\Character{61}}%
\TextAt{13}{177.446}{509.731}{R}%
\TextAt{21}{186.356}{507.938}{w}%
\TextAt{9}{193.146}{509.731}{\Character{40}}%
\TextAt{13}{197.698}{509.731}{L}%
\TextAt{16}{208.983}{509.731}{\Character{106}}%
\TextAt{13}{215.625}{509.731}{K}%
\TextAt{9}{226.443}{509.731}{\Character{41}}%
\TextAt{9}{234.298}{509.731}{\Character{61}}%
\TextAt{20}{246.741}{523.002}{\Character{110}}%
\TextAt{13}{254.711}{509.731}{\Character{27}}%
\NativeMathAt{265.115}{509.731}{11.955}{\in}%
\TextAt{13}{276.406}{509.731}{G}%
\TextAt{21}{285.639}{507.938}{w}%
\TextAt{20}{296.326}{523.48}{\Character{12}}%
\TextAt{20}{296.326}{516.307}{\Character{12}}%
\TextAt{20}{296.326}{509.134}{\Character{12}}%
\TextAt{13}{305.404}{517.819}{\Character{27}}%
\TextAt{13}{312.47}{517.819}{x}%
\Vector{q 0 G 0 g 0.478 w 233.404 -261.28 m 247.139 -261.28 l S Q}%
\TextAt{13}{308.945}{501.531}{x}%
\NativeMathAt{323.655}{509.731}{11.955}{\equiv}%
\TextAt{9}{336.274}{509.731}{1}%
\TextAt{8}{345.448}{509.731}{mod}%
\TextAt{14}{370.886}{509.731}{P}%
\TextAt{8}{384.677}{509.731}{for}%
\TextAt{8}{402.561}{509.731}{all}%
\TextAt{13}{418.82}{509.731}{x}%
\NativeMathAt{428.793}{509.731}{11.955}{\in}%
\TextAt{13}{440.084}{509.731}{L}%
\TextAt{3}{448.049}{514.668}{\Character{3}}%
\TextAt{20}{452.781}{523.002}{\Character{111}}%
\TextAt{13}{462.744}{509.731}{\Character{58}}%
\TextAt{8}{97.159}{472.461}{Observe}%
\TextAt{8}{142.55}{472.461}{in}%
\TextAt{8}{156.654}{472.461}{this}%
\TextAt{8}{179.927}{472.461}{definition}%
\TextAt{8}{232.407}{472.461}{that,}%
\TextAt{8}{261.592}{472.461}{for}%
\TextAt{13}{279.928}{472.461}{\Character{27}}%
\NativeMathAt{291.1}{472.461}{11.955}{\in}%
\TextAt{13}{303.16}{472.461}{G}%
\TextAt{21}{312.394}{470.667}{w}%
\TextAt{8}{319.183}{472.461}{,}%
\TextAt{8}{326.896}{472.461}{the}%
\TextAt{8}{347.504}{472.461}{identity}%
\TextAt{13}{391.209}{472.461}{w}%
\TextAt{16}{402.863}{472.461}{\Character{14}}%
\TextAt{13}{411.804}{472.461}{\Character{27}}%
\TextAt{9}{422.977}{472.461}{\Character{61}}%
\TextAt{13}{436.171}{472.461}{w}%
\TextAt{8}{449.209}{472.461}{implies}%
\TextAt{8}{489.387}{472.461}{that}%
\TextAt{8}{515.207}{472.461}{one}%
\TextAt{8}{79.226}{455.415}{always}%
\TextAt{8}{116.694}{455.415}{has}%
\TextAt{13}{137.567}{455.415}{\Character{27}}%
\TextAt{16}{144.65}{455.415}{\Character{79}}%
\TextAt{9}{157.82}{455.415}{\Character{61}}%
\TextAt{16}{170.246}{455.415}{\Character{79}}%
\TextAt{8}{183.993}{455.415}{and}%
\TextAt{13}{206.756}{455.415}{\Character{27}}%
\TextAt{13}{213.822}{455.415}{x}%
\TextAt{9}{220.49}{455.415}{\Character{47}}%
\TextAt{13}{226.343}{455.415}{x}%
\NativeMathAt{236.316}{455.415}{11.955}{\in}%
\TextAt{16}{247.59}{455.415}{\Character{79}}%
\TextAt{8}{257.457}{455.415}{,}%
\TextAt{8}{264.606}{455.415}{for}%
\TextAt{8}{282.491}{455.415}{all}%
\TextAt{13}{298.75}{455.415}{x}%
\NativeMathAt{308.723}{455.415}{11.955}{\in}%
\TextAt{13}{320.014}{455.415}{L}%
\TextAt{3}{327.978}{459.753}{\Character{3}}%
\TextAt{8}{332.711}{455.415}{.}%
\TextAt{8}{97.159}{435.38}{The}%
\TextAt{8}{121.747}{435.38}{subgroups}%
\TextAt{13}{177.694}{435.38}{G}%
\TextAt{21}{186.928}{433.587}{w}%
\TextAt{13}{193.717}{435.38}{\Character{59}}%
\TextAt{13}{198.941}{435.38}{I}%
\TextAt{21}{204.124}{433.587}{w}%
\TextAt{13}{210.913}{435.38}{\Character{59}}%
\TextAt{13}{216.137}{435.38}{R}%
\TextAt{21}{225.068}{433.587}{w}%
\TextAt{8}{236.289}{435.38}{of}%
\TextAt{13}{250.154}{435.38}{G}%
\TextAt{9}{263.619}{435.38}{\Character{61}}%
\TextAt{13}{276.955}{435.38}{G}%
\TextAt{9}{286.189}{435.38}{\Character{40}}%
\TextAt{13}{290.741}{435.38}{L}%
\TextAt{16}{302.937}{435.38}{\Character{106}}%
\TextAt{13}{310.49}{435.38}{K}%
\TextAt{9}{321.307}{435.38}{\Character{41}}%
\TextAt{8}{325.86}{435.38}{,}%
\TextAt{8}{333.677}{435.38}{and}%
\TextAt{8}{356.974}{435.38}{in}%
\TextAt{8}{371.162}{435.38}{fact}%
\TextAt{8}{394.781}{435.38}{all}%
\TextAt{8}{411.575}{435.38}{canonical}%
\TextAt{8}{463.493}{435.38}{subgroups}%
\TextAt{8}{519.439}{435.38}{we}%
\TextAt{8}{79.226}{418.334}{will}%
\TextAt{8}{102.482}{418.334}{encounter}%
\TextAt{8}{157.287}{418.334}{in}%
\TextAt{8}{172.09}{418.334}{the}%
\TextAt{8}{193.397}{418.334}{sequel,}%
\TextAt{8}{232.937}{418.334}{are}%
\TextAt{8}{253.598}{418.334}{all}%
\TextAt{19}{271.008}{418.334}{closed}%
\TextAt{8}{306.541}{418.334}{in}%
\TextAt{8}{321.345}{418.334}{the}%
\TextAt{8}{342.652}{418.334}{Krull}%
\TextAt{8}{374.36}{418.334}{topology.}%
\TextAt{8}{429.525}{418.334}{The}%
\TextAt{8}{454.73}{418.334}{proof}%
\TextAt{8}{486.449}{418.334}{of}%
\TextAt{8}{500.93}{418.334}{this}%
\TextAt{8}{524.903}{418.334}{is}%
\TextAt{8}{79.226}{401.288}{routine}%
\TextAt{8}{119.919}{401.288}{in}%
\TextAt{8}{133.94}{401.288}{all}%
\TextAt{8}{150.567}{401.288}{cases.}%
\TextAt{8}{185.609}{401.288}{Let}%
\TextAt{8}{206.946}{401.288}{us}%
\TextAt{8}{222.33}{401.288}{just}%
\TextAt{8}{245.843}{401.288}{illustrate}%
\TextAt{8}{295.706}{401.288}{the}%
\TextAt{8}{316.231}{401.288}{model}%
\TextAt{8}{351.388}{401.288}{of}%
\TextAt{8}{365.086}{401.288}{the}%
\TextAt{8}{385.611}{401.288}{argument}%
\TextAt{8}{438.343}{401.288}{for}%
\TextAt{8}{456.596}{401.288}{the}%
\TextAt{8}{477.121}{401.288}{example}%
\TextAt{8}{523.337}{401.288}{of}%
\TextAt{8}{79.226}{384.242}{the}%
\TextAt{8}{99.383}{384.242}{decomposition}%
\TextAt{8}{176.517}{384.242}{group.}%
\TextAt{8}{97.159}{364.207}{Let}%
\TextAt{13}{117.809}{364.207}{\Character{27}}%
\NativeMathAt{128.212}{364.207}{11.955}{\in}%
\TextAt{13}{139.503}{364.207}{G}%
\TextAt{9}{152.058}{364.207}{\Character{61}}%
\TextAt{13}{164.483}{364.207}{G}%
\TextAt{9}{173.717}{364.207}{\Character{40}}%
\TextAt{13}{178.269}{364.207}{L}%
\TextAt{16}{189.554}{364.207}{\Character{106}}%
\TextAt{13}{196.196}{364.207}{K}%
\TextAt{9}{207.014}{364.207}{\Character{41}}%
\TextAt{8}{215.144}{364.207}{be}%
\TextAt{8}{230.749}{364.207}{an}%
\TextAt{8}{246.689}{364.207}{element}%
\TextAt{8}{289.612}{364.207}{which}%
\TextAt{8}{322.779}{364.207}{belongs}%
\TextAt{8}{364.47}{364.207}{to}%
\TextAt{8}{378.461}{364.207}{the}%
\TextAt{8}{398.299}{364.207}{closure}%
\TextAt{8}{437.061}{364.207}{of}%
\TextAt{13}{450.072}{364.207}{G}%
\TextAt{21}{459.305}{362.414}{w}%
\TextAt{8}{466.095}{364.207}{.}%
\TextAt{8}{474.437}{364.207}{This}%
\TextAt{8}{500.837}{364.207}{means}%
\TextAt{8}{79.226}{347.161}{that,}%
\TextAt{8}{108.085}{347.161}{in}%
\TextAt{8}{121.928}{347.161}{every}%
\TextAt{8}{153.01}{347.161}{neighbourhood}%
\TextAt{13}{233.204}{347.161}{\Character{27}}%
\TextAt{13}{240.27}{347.161}{G}%
\TextAt{9}{249.52}{347.161}{\Character{40}}%
\TextAt{13}{254.073}{347.161}{L}%
\TextAt{16}{265.682}{347.161}{\Character{106}}%
\TextAt{13}{272.648}{347.161}{M}%
\TextAt{9}{285.222}{347.161}{\Character{41}}%
\TextAt{8}{289.774}{347.161}{,}%
\TextAt{8}{297.161}{347.161}{there}%
\TextAt{8}{327.263}{347.161}{is}%
\TextAt{8}{339.217}{347.161}{some}%
\TextAt{8}{368.733}{347.161}{element}%
\TextAt{13}{412.165}{347.161}{\Character{27}}%
\TextAt{21}{418.819}{345.368}{M}%
\TextAt{8}{432.386}{347.161}{of}%
\TextAt{13}{445.907}{347.161}{G}%
\TextAt{21}{455.14}{345.368}{w}%
\TextAt{8}{461.93}{347.161}{.}%
\TextAt{8}{470.949}{347.161}{Here}%
\TextAt{13}{498.767}{347.161}{M}%
\TextAt{16}{514.986}{347.161}{\Character{106}}%
\TextAt{13}{521.952}{347.161}{K}%
\TextAt{8}{79.226}{330.115}{varies}%
\TextAt{8}{113.159}{330.115}{over}%
\TextAt{8}{139.226}{330.115}{all}%
\TextAt{8}{156.505}{330.115}{finite}%
\TextAt{8}{187.437}{330.115}{Galois}%
\TextAt{8}{224.371}{330.115}{subextensions}%
\TextAt{8}{299.717}{330.115}{of}%
\TextAt{13}{314.067}{330.115}{L}%
\TextAt{16}{327.092}{330.115}{\Character{106}}%
\TextAt{13}{335.472}{330.115}{K}%
\TextAt{8}{346.29}{330.115}{.}%
\TextAt{8}{357.8}{330.115}{Since}%
\TextAt{13}{389.377}{330.115}{\Character{27}}%
\TextAt{21}{396.031}{328.322}{M}%
\NativeMathAt{410.57}{330.115}{11.955}{\in}%
\TextAt{13}{423.6}{330.115}{\Character{27}}%
\TextAt{13}{430.666}{330.115}{G}%
\TextAt{9}{439.916}{330.115}{\Character{40}}%
\TextAt{13}{444.469}{330.115}{L}%
\TextAt{16}{457.493}{330.115}{\Character{106}}%
\TextAt{13}{465.874}{330.115}{M}%
\TextAt{9}{478.447}{330.115}{\Character{41}}%
\TextAt{8}{483}{330.115}{,}%
\TextAt{8}{491.424}{330.115}{we}%
\TextAt{8}{509.672}{330.115}{have}%
\TextAt{13}{79.226}{313.069}{\Character{27}}%
\TextAt{21}{85.88}{311.276}{M}%
\TextAt{16}{95.359}{313.069}{\Character{106}}%
\TextAt{21}{98.68}{311.276}{M}%
\TextAt{9}{112.783}{313.069}{\Character{61}}%
\TextAt{13}{126.511}{313.069}{\Character{27}}%
\TextAt{16}{133.594}{313.069}{\Character{106}}%
\TextAt{21}{136.915}{311.276}{M}%
\TextAt{8}{146.395}{313.069}{,}%
\TextAt{8}{154.499}{313.069}{and}%
\TextAt{13}{178.026}{313.069}{w}%
\TextAt{16}{189.893}{313.069}{\Character{14}}%
\TextAt{13}{199.048}{313.069}{\Character{27}}%
\TextAt{21}{205.702}{311.276}{M}%
\TextAt{9}{219.805}{313.069}{\Character{61}}%
\TextAt{13}{233.533}{313.069}{w}%
\TextAt{8}{246.884}{313.069}{implies}%
\TextAt{8}{287.376}{313.069}{that}%
\TextAt{13}{313.509}{313.069}{w}%
\TextAt{16}{325.376}{313.069}{\Character{14}}%
\TextAt{13}{334.531}{313.069}{\Character{27}}%
\TextAt{16}{341.614}{313.069}{\Character{106}}%
\TextAt{21}{344.934}{311.276}{M}%
\TextAt{9}{359.038}{313.069}{\Character{61}}%
\TextAt{13}{372.765}{313.069}{w}%
\TextAt{16}{384.633}{313.069}{\Character{14}}%
\TextAt{13}{393.788}{313.069}{\Character{27}}%
\TextAt{21}{400.441}{311.276}{M}%
\TextAt{16}{409.921}{313.069}{\Character{106}}%
\TextAt{21}{413.242}{311.276}{M}%
\TextAt{9}{427.345}{313.069}{\Character{61}}%
\TextAt{13}{441.073}{313.069}{w}%
\TextAt{16}{449.763}{313.069}{\Character{106}}%
\TextAt{21}{453.084}{311.276}{M}%
\TextAt{8}{462.563}{313.069}{.}%
\TextAt{8}{473.305}{313.069}{As}%
\TextAt{13}{491.356}{313.069}{L}%
\TextAt{8}{503.982}{313.069}{is}%
\TextAt{8}{516.51}{313.069}{the}%
\TextAt{8}{79.226}{296.023}{union}%
\TextAt{8}{111.617}{296.023}{of}%
\TextAt{8}{124.82}{296.023}{all}%
\TextAt{8}{140.952}{296.023}{the}%
\TextAt{13}{160.981}{296.023}{M}%
\TextAt{8}{173.555}{296.023}{,}%
\TextAt{8}{180.603}{296.023}{we}%
\TextAt{8}{197.703}{296.023}{get}%
\TextAt{13}{217.087}{296.023}{w}%
\TextAt{16}{228.173}{296.023}{\Character{14}}%
\TextAt{13}{236.547}{296.023}{\Character{27}}%
\TextAt{9}{246.951}{296.023}{\Character{61}}%
\TextAt{13}{259.376}{296.023}{w}%
\TextAt{8}{268.066}{296.023}{,}%
\TextAt{8}{275.113}{296.023}{so}%
\TextAt{8}{289.356}{296.023}{that}%
\TextAt{13}{314.598}{296.023}{\Character{27}}%
\NativeMathAt{325.001}{296.023}{11.955}{\in}%
\TextAt{13}{336.292}{296.023}{G}%
\TextAt{21}{345.526}{294.23}{w}%
\TextAt{8}{352.315}{296.023}{.}%
\TextAt{8}{360.721}{296.023}{This}%
\TextAt{8}{387.314}{296.023}{shows}%
\TextAt{8}{420.805}{296.023}{that}%
\TextAt{8}{446.047}{296.023}{the}%
\TextAt{8}{466.076}{296.023}{subgroup}%
\TextAt{13}{516.746}{296.023}{G}%
\TextAt{21}{525.98}{294.23}{w}%
\TextAt{8}{79.226}{278.977}{is}%
\TextAt{8}{90.99}{278.977}{closed}%
\TextAt{8}{125.517}{278.977}{in}%
\TextAt{13}{139.169}{278.977}{G}%
\TextAt{8}{148.403}{278.977}{.}%
\TextAt{8}{97.159}{258.943}{The}%
\TextAt{8}{121.496}{258.943}{groups}%
\TextAt{13}{159.569}{258.943}{G}%
\TextAt{21}{168.803}{257.149}{w}%
\TextAt{13}{175.592}{258.943}{\Character{59}}%
\TextAt{13}{180.816}{258.943}{I}%
\TextAt{21}{185.999}{257.149}{w}%
\TextAt{13}{192.789}{258.943}{\Character{59}}%
\TextAt{13}{198.013}{258.943}{R}%
\TextAt{21}{206.943}{257.149}{w}%
\TextAt{8}{217.913}{258.943}{carry}%
\TextAt{8}{248.443}{258.943}{very}%
\TextAt{8}{274.417}{258.943}{significant}%
\TextAt{8}{330.376}{258.943}{information}%
\TextAt{8}{394.081}{258.943}{about}%
\TextAt{8}{427.862}{258.943}{the}%
\TextAt{8}{448.302}{258.943}{behaviour}%
\TextAt{8}{502.897}{258.943}{of}%
\TextAt{8}{516.51}{258.943}{the}%
\TextAt{8}{79.226}{241.897}{valuation}%
\TextAt{13}{131.411}{241.897}{v}%
\TextAt{8}{142.509}{241.897}{of}%
\TextAt{13}{156.951}{241.897}{K}%
\TextAt{8}{172.778}{241.897}{as}%
\TextAt{8}{188.261}{241.897}{it}%
\TextAt{8}{201.077}{241.897}{is}%
\TextAt{8}{213.953}{241.897}{extended}%
\TextAt{8}{264.811}{241.897}{to}%
\TextAt{13}{280.234}{241.897}{L}%
\TextAt{8}{288.198}{241.897}{.}%
\TextAt{8}{299.984}{241.897}{But}%
\TextAt{8}{324.337}{241.897}{before}%
\TextAt{8}{360.561}{241.897}{going}%
\TextAt{8}{392.901}{241.897}{into}%
\TextAt{8}{417.756}{241.897}{this,}%
\TextAt{8}{445.22}{241.897}{we}%
\TextAt{8}{463.56}{241.897}{will}%
\TextAt{8}{486.777}{241.897}{treat}%
\TextAt{8}{516.51}{241.897}{the}%
\TextAt{8}{79.226}{224.851}{functorial}%
\TextAt{8}{132.235}{224.851}{properties}%
\TextAt{8}{187.253}{224.851}{of}%
\TextAt{8}{200.583}{224.851}{the}%
\TextAt{8}{220.739}{224.851}{groups}%
\TextAt{13}{258.53}{224.851}{G}%
\TextAt{21}{267.763}{223.057}{w}%
\TextAt{13}{274.553}{224.851}{\Character{59}}%
\TextAt{13}{279.777}{224.851}{I}%
\TextAt{21}{284.96}{223.057}{w}%
\TextAt{13}{291.749}{224.851}{\Character{59}}%
\TextAt{13}{296.973}{224.851}{R}%
\TextAt{21}{305.904}{223.057}{w}%
\TextAt{8}{312.693}{224.851}{.}%
\TextAt{8}{97.159}{204.816}{Consider}%
\TextAt{8}{145.996}{204.816}{two}%
\TextAt{8}{168.113}{204.816}{Galois}%
\TextAt{8}{204.026}{204.816}{extensions}%
\TextAt{13}{260.407}{204.816}{L}%
\TextAt{16}{271.692}{204.816}{\Character{106}}%
\TextAt{13}{278.334}{204.816}{K}%
\TextAt{8}{293.049}{204.816}{and}%
\TextAt{13}{315.811}{204.816}{L}%
\TextAt{3}{323.776}{209.154}{\Character{48}}%
\TextAt{16}{329.892}{204.816}{\Character{106}}%
\TextAt{13}{336.534}{204.816}{K}%
\TextAt{3}{347.351}{209.154}{\Character{48}}%
\TextAt{8}{354.044}{204.816}{and}%
\TextAt{8}{376.806}{204.816}{a}%
\TextAt{8}{386.562}{204.816}{commutative}%
\TextAt{8}{456.488}{204.816}{diagram}%
\TextAt{13}{268.537}{176.271}{L}%
\TextAt{13}{332.699}{176.271}{L}%
\TextAt{3}{340.663}{180.61}{\Character{48}}%
\TextAt{13}{267.11}{126.466}{K}%
\TextAt{13}{331.273}{126.466}{K}%
\TextAt{3}{342.09}{130.805}{\Character{48}}%
\NativeDiagram{%
\path[draw=black,line width=0.478bp] (209.853,-594.74) -- (254.869,-594.74);%
\path[draw=black,line width=0.478bp,line cap=round,line join=round] (252.622,-591.87) .. controls (253.076,-593.592) and (254.088,-594.405) .. (255.108,-594.74) .. controls (254.088,-595.074) and (253.076,-595.887) .. (252.622,-597.609);%
}%
\TextAt{21}{302.261}{182.076}{\Character{28}}%
\NativeDiagram{%
\path[draw=black,line width=0.478bp] (200.519,-633.064) -- (200.519,-602.786);%
\path[draw=black,line width=0.478bp,line cap=round] (198.176,-633.064) .. controls (198.176,-633.883) and (198.701,-634.546) .. (199.348,-634.546) .. controls (199.995,-634.546) and (200.519,-633.883) .. (200.519,-633.064);%
\path[draw=black,line width=0.478bp,line cap=round,line join=round] (197.65,-605.033) .. controls (199.372,-604.579) and (200.184,-603.567) .. (200.519,-602.547) .. controls (200.854,-603.567) and (201.667,-604.579) .. (203.388,-605.033);%
\path[draw=black,line width=0.478bp] (211.28,-644.545) -- (253.443,-644.545);%
\path[draw=black,line width=0.478bp,line cap=round,line join=round] (251.195,-641.676) .. controls (251.649,-643.397) and (252.662,-644.21) .. (253.682,-644.545) .. controls (252.662,-644.88) and (251.649,-645.693) .. (251.195,-647.414);%
}%
\TextAt{21}{302.261}{132.27}{\Character{28}}%
\NativeDiagram{%
\path[draw=black,line width=0.478bp] (266.079,-632.467) -- (266.079,-602.786);%
\path[draw=black,line width=0.478bp,line cap=round] (263.736,-632.467) .. controls (263.736,-633.286) and (264.261,-633.949) .. (264.908,-633.949) .. controls (265.555,-633.949) and (266.079,-633.286) .. (266.079,-632.467);%
\path[draw=black,line width=0.478bp,line cap=round,line join=round] (263.21,-605.033) .. controls (264.932,-604.579) and (265.745,-603.567) .. (266.079,-602.547) .. controls (266.414,-603.567) and (267.227,-604.579) .. (268.948,-605.033);%
}%
\TextAt{8}{79.226}{101.381}{with}%
\TextAt{8}{105.512}{101.381}{homomorphisms}%
\TextAt{13}{192.422}{101.381}{\Character{28}}%
\TextAt{8}{202.364}{101.381}{which}%
\TextAt{8}{235.476}{101.381}{will}%
\TextAt{8}{257.207}{101.381}{typically}%
\TextAt{8}{304.641}{101.381}{be}%
\TextAt{8}{320.191}{101.381}{inclusions.}%
\TextAt{8}{378.069}{101.381}{They}%
\TextAt{8}{407.929}{101.381}{induce}%
\TextAt{8}{444.616}{101.381}{a}%
\TextAt{8}{453.997}{101.381}{homomorphism}%
\TextAt{13}{176.811}{72.921}{\Character{28}}%
\TextAt{3}{183.23}{77.858}{\Character{3}}%
\TextAt{9}{191.283}{72.921}{\Character{58}}%
\TextAt{13}{197.856}{72.921}{G}%
\TextAt{9}{207.089}{72.921}{\Character{40}}%
\TextAt{13}{211.642}{72.921}{L}%
\TextAt{3}{219.606}{77.858}{\Character{48}}%
\TextAt{16}{225.722}{72.921}{\Character{106}}%
\TextAt{13}{232.364}{72.921}{K}%
\TextAt{3}{243.181}{77.858}{\Character{48}}%
\TextAt{9}{245.977}{72.921}{\Character{41}}%
\NativeArrowAt{253.85}{72.921}{11.955}{19.26}{\longrightarrow}%
\TextAt{13}{276.432}{72.921}{G}%
\TextAt{9}{285.665}{72.921}{\Character{40}}%
\TextAt{13}{290.218}{72.921}{L}%
\TextAt{16}{301.503}{72.921}{\Character{106}}%
\TextAt{13}{308.145}{72.921}{K}%
\TextAt{9}{318.962}{72.921}{\Character{41}}%
\TextAt{13}{323.515}{72.921}{\Character{59}}%
\TextAt{13}{352.65}{72.921}{\Character{28}}%
\TextAt{3}{359.088}{77.858}{\Character{3}}%
\TextAt{9}{363.821}{72.921}{\Character{40}}%
\TextAt{13}{368.373}{72.921}{\Character{27}}%
\TextAt{3}{375.455}{77.858}{\Character{48}}%
\TextAt{9}{378.25}{72.921}{\Character{41}}%
\TextAt{9}{386.105}{72.921}{\Character{61}}%
\TextAt{13}{398.549}{72.921}{\Character{28}}%
\TextAt{3}{404.968}{77.858}{\Character{0}}%
\TextAt{17}{411.554}{77.858}{1}%
\TextAt{13}{416.287}{72.921}{\Character{27}}%
\TextAt{3}{423.369}{77.858}{\Character{48}}%
\TextAt{13}{426.164}{72.921}{\Character{28}}%
\TextAt{13}{431.926}{72.921}{\Character{58}}%
\TextAt{8}{297.211}{43.033}{154}%
\Destination{78.23}{809.815}{nav.page.154}%
\Destination{79.36}{710.355}{nav.print.168}%
\Destination{79.36}{775.155}{nav.number.2.9.2}%
\Destination{79.36}{609.155}{nav.number.2.9.3}%
}

% Page 167 (body page 155)
\Page{%
\TextAt{2}{79.226}{796.943}{II.}%
\TextAt{2}{94.073}{796.943}{The}%
\TextAt{2}{116.492}{796.943}{Theory}%
\TextAt{2}{154.401}{796.943}{of}%
\TextAt{2}{166.826}{796.943}{Valuations}%
\Vector{q 0 G 0 g 0.299 w 7.226 18.729 m 460.769 18.729 l S Q}%
\TextAt{8}{79.226}{761.124}{Observe}%
\TextAt{8}{124.722}{761.124}{here}%
\TextAt{8}{150.634}{761.124}{that,}%
\TextAt{13}{179.95}{761.124}{L}%
\TextAt{16}{192.183}{761.124}{\Character{106}}%
\TextAt{13}{199.772}{761.124}{K}%
\TextAt{8}{215.043}{761.124}{being}%
\TextAt{8}{247.137}{761.124}{normal,}%
\TextAt{8}{290.763}{761.124}{the}%
\TextAt{8}{311.475}{761.124}{same}%
\TextAt{8}{341.357}{761.124}{is}%
\TextAt{8}{353.677}{761.124}{true}%
\TextAt{8}{378.944}{761.124}{of}%
\TextAt{13}{392.83}{761.124}{\Character{28}}%
\TextAt{13}{399.238}{761.124}{L}%
\TextAt{16}{411.482}{761.124}{\Character{106}}%
\TextAt{13}{419.071}{761.124}{\Character{28}}%
\TextAt{13}{425.479}{761.124}{K}%
\TextAt{8}{436.307}{761.124}{,}%
\TextAt{8}{444.152}{761.124}{and}%
\TextAt{8}{467.47}{761.124}{thus}%
\TextAt{8}{493.778}{761.124}{one}%
\TextAt{8}{515.793}{761.124}{has}%
\TextAt{13}{79.226}{744.078}{\Character{27}}%
\TextAt{3}{86.309}{748.416}{\Character{48}}%
\TextAt{13}{89.104}{744.078}{\Character{28}}%
\TextAt{13}{95.512}{744.078}{L}%
\NativeMathAt{106.808}{744.078}{11.955}{\subseteq}%
\TextAt{13}{119.427}{744.078}{\Character{28}}%
\TextAt{13}{125.835}{744.078}{L}%
\TextAt{8}{133.811}{744.078}{,}%
\TextAt{8}{140.96}{744.078}{so}%
\TextAt{8}{155.33}{744.078}{that}%
\TextAt{8}{180.699}{744.078}{composing}%
\TextAt{8}{238.323}{744.078}{with}%
\TextAt{13}{264.983}{744.078}{\Character{28}}%
\TextAt{3}{271.402}{748.416}{\Character{0}}%
\TextAt{17}{277.988}{748.416}{1}%
\TextAt{8}{286.618}{744.078}{makes}%
\TextAt{8}{321.802}{744.078}{sense.}%
\TextAt{8}{97.159}{724.043}{Now}%
\TextAt{8}{124.041}{724.043}{let}%
\TextAt{13}{141.167}{724.043}{w}%
\TextAt{3}{149.856}{728.381}{\Character{48}}%
\TextAt{8}{156.77}{724.043}{be}%
\TextAt{8}{172.916}{724.043}{a}%
\TextAt{8}{182.894}{724.043}{valuation}%
\TextAt{8}{234.188}{724.043}{of}%
\TextAt{13}{247.739}{724.043}{L}%
\TextAt{3}{255.704}{728.381}{\Character{48}}%
\TextAt{8}{258.499}{724.043}{,}%
\TextAt{13}{265.925}{724.043}{v}%
\TextAt{3}{272.013}{728.381}{\Character{48}}%
\TextAt{9}{278.507}{724.043}{\Character{61}}%
\TextAt{13}{291.31}{724.043}{w}%
\TextAt{3}{300}{728.381}{\Character{48}}%
\TextAt{16}{302.795}{724.043}{\Character{106}}%
\TextAt{21}{306.116}{722.25}{K}%
\TextAt{23}{313.82}{724.519}{\Character{48}}%
\TextAt{8}{317.003}{724.043}{,}%
\TextAt{8}{324.429}{724.043}{and}%
\TextAt{13}{347.413}{724.043}{w}%
\TextAt{9}{359.801}{724.043}{\Character{61}}%
\TextAt{13}{372.605}{724.043}{w}%
\TextAt{3}{381.294}{728.381}{\Character{48}}%
\TextAt{16}{386.897}{724.043}{\Character{14}}%
\TextAt{13}{395.682}{724.043}{\Character{28}}%
\TextAt{8}{402.101}{724.043}{,}%
\TextAt{13}{409.528}{724.043}{v}%
\TextAt{9}{419.314}{724.043}{\Character{61}}%
\TextAt{13}{432.118}{724.043}{w}%
\TextAt{16}{440.807}{724.043}{\Character{106}}%
\TextAt{21}{444.128}{722.25}{K}%
\TextAt{8}{452.331}{724.043}{.}%
\TextAt{8}{461.444}{724.043}{Then}%
\TextAt{8}{492.223}{724.043}{we}%
\TextAt{8}{509.672}{724.043}{have}%
\TextAt{8}{79.226}{706.997}{the}%
\TextAt{12}{79.226}{682.154}{Proposition}%
\TextAt{12}{151.937}{682.154}{(9.4).}%
\TextAt{13}{187.253}{682.154}{\Character{28}}%
\TextAt{3}{193.672}{686.493}{\Character{3}}%
\TextAt{9}{201.725}{682.154}{\Character{58}}%
\TextAt{13}{208.298}{682.154}{G}%
\TextAt{9}{217.531}{682.154}{\Character{40}}%
\TextAt{13}{222.084}{682.154}{L}%
\TextAt{3}{230.048}{686.493}{\Character{48}}%
\TextAt{16}{236.164}{682.154}{\Character{106}}%
\TextAt{13}{242.806}{682.154}{K}%
\TextAt{3}{253.623}{686.493}{\Character{48}}%
\TextAt{9}{256.418}{682.154}{\Character{41}}%
\NativeMathAt{264.292}{682.154}{11.955}{\rightarrow}%
\TextAt{13}{279.568}{682.154}{G}%
\TextAt{9}{288.801}{682.154}{\Character{40}}%
\TextAt{13}{293.354}{682.154}{L}%
\TextAt{16}{304.639}{682.154}{\Character{106}}%
\TextAt{13}{311.281}{682.154}{K}%
\TextAt{9}{322.098}{682.154}{\Character{41}}%
\TextAt{8}{330.548}{682.154}{induces}%
\TextAt{8}{372.224}{682.154}{homomorphisms}%
\TextAt{13}{233.411}{658.363}{G}%
\TextAt{21}{242.644}{656.57}{w}%
\TextAt{23}{248.936}{658.839}{\Character{48}}%
\TextAt{9}{252.118}{658.363}{\Character{40}}%
\TextAt{13}{256.671}{658.363}{L}%
\TextAt{3}{264.635}{663.299}{\Character{48}}%
\TextAt{16}{270.751}{658.363}{\Character{106}}%
\TextAt{13}{277.393}{658.363}{K}%
\TextAt{3}{288.21}{663.299}{\Character{48}}%
\TextAt{9}{291.006}{658.363}{\Character{41}}%
\NativeArrowAt{298.879}{658.363}{11.955}{19.26}{\longrightarrow}%
\TextAt{13}{321.461}{658.363}{G}%
\TextAt{21}{330.694}{656.57}{w}%
\TextAt{9}{337.484}{658.363}{\Character{40}}%
\TextAt{13}{342.036}{658.363}{L}%
\TextAt{16}{353.321}{658.363}{\Character{106}}%
\TextAt{13}{359.963}{658.363}{K}%
\TextAt{9}{370.781}{658.363}{\Character{41}}%
\TextAt{13}{375.333}{658.363}{\Character{59}}%
\TextAt{13}{237.482}{638.328}{I}%
\TextAt{21}{242.644}{636.535}{w}%
\TextAt{23}{248.936}{638.804}{\Character{48}}%
\TextAt{9}{252.118}{638.328}{\Character{40}}%
\TextAt{13}{256.671}{638.328}{L}%
\TextAt{3}{264.635}{643.265}{\Character{48}}%
\TextAt{16}{270.751}{638.328}{\Character{106}}%
\TextAt{13}{277.393}{638.328}{K}%
\TextAt{3}{288.21}{643.265}{\Character{48}}%
\TextAt{9}{291.006}{638.328}{\Character{41}}%
\NativeArrowAt{298.879}{638.328}{11.955}{19.26}{\longrightarrow}%
\TextAt{13}{321.461}{638.328}{I}%
\TextAt{21}{326.624}{636.535}{w}%
\TextAt{9}{333.413}{638.328}{\Character{40}}%
\TextAt{13}{337.965}{638.328}{L}%
\TextAt{16}{349.251}{638.328}{\Character{106}}%
\TextAt{13}{355.892}{638.328}{K}%
\TextAt{9}{366.71}{638.328}{\Character{41}}%
\TextAt{13}{371.262}{638.328}{\Character{59}}%
\TextAt{13}{233.734}{618.294}{R}%
\TextAt{21}{242.644}{616.5}{w}%
\TextAt{23}{248.936}{618.77}{\Character{48}}%
\TextAt{9}{252.118}{618.294}{\Character{40}}%
\TextAt{13}{256.671}{618.294}{L}%
\TextAt{3}{264.635}{623.23}{\Character{48}}%
\TextAt{16}{270.751}{618.294}{\Character{106}}%
\TextAt{13}{277.393}{618.294}{K}%
\TextAt{3}{288.21}{623.23}{\Character{48}}%
\TextAt{9}{291.006}{618.294}{\Character{41}}%
\NativeArrowAt{298.879}{618.294}{11.955}{19.26}{\longrightarrow}%
\TextAt{13}{321.461}{618.294}{R}%
\TextAt{21}{330.371}{616.5}{w}%
\TextAt{9}{337.16}{618.294}{\Character{40}}%
\TextAt{13}{341.713}{618.294}{L}%
\TextAt{16}{352.998}{618.294}{\Character{106}}%
\TextAt{13}{359.64}{618.294}{K}%
\TextAt{9}{370.457}{618.294}{\Character{41}}%
\TextAt{13}{375.01}{618.294}{\Character{58}}%
\TextAt{8}{79.226}{596.025}{In}%
\TextAt{8}{93.847}{596.025}{the}%
\TextAt{8}{114.004}{596.025}{latter}%
\TextAt{8}{145.876}{596.025}{two}%
\TextAt{8}{167.993}{596.025}{cases,}%
\TextAt{13}{200.631}{596.025}{v}%
\TextAt{8}{210.617}{596.025}{is}%
\TextAt{8}{222.38}{596.025}{assumed}%
\TextAt{8}{269.328}{596.025}{to}%
\TextAt{8}{283.639}{596.025}{be}%
\TextAt{8}{299.563}{596.025}{nonarchimedean.}%
\TextAt{19}{79.226}{563.385}{Proof.}%
\TextAt{8}{116}{563.385}{Let}%
\TextAt{13}{136.97}{563.385}{\Character{27}}%
\TextAt{3}{144.052}{567.723}{\Character{48}}%
\NativeMathAt{150.168}{563.385}{11.955}{\in}%
\TextAt{13}{161.459}{563.385}{G}%
\TextAt{21}{170.693}{561.592}{w}%
\TextAt{23}{176.984}{563.861}{\Character{48}}%
\TextAt{9}{180.166}{563.385}{\Character{40}}%
\TextAt{13}{184.719}{563.385}{L}%
\TextAt{3}{192.683}{567.723}{\Character{48}}%
\TextAt{16}{198.799}{563.385}{\Character{106}}%
\TextAt{13}{205.441}{563.385}{K}%
\TextAt{3}{216.258}{567.723}{\Character{48}}%
\TextAt{9}{219.054}{563.385}{\Character{41}}%
\TextAt{8}{227.503}{563.385}{and}%
\TextAt{13}{250.266}{563.385}{\Character{27}}%
\TextAt{9}{260.669}{563.385}{\Character{61}}%
\TextAt{13}{273.095}{563.385}{\Character{28}}%
\TextAt{3}{279.514}{567.723}{\Character{3}}%
\TextAt{9}{284.246}{563.385}{\Character{40}}%
\TextAt{13}{288.798}{563.385}{\Character{27}}%
\TextAt{3}{295.881}{567.723}{\Character{48}}%
\TextAt{9}{298.676}{563.385}{\Character{41}}%
\TextAt{8}{303.228}{563.385}{.}%
\TextAt{8}{311.676}{563.385}{If}%
\TextAt{13}{323.368}{563.385}{x}%
\NativeMathAt{333.341}{563.385}{11.955}{\in}%
\TextAt{13}{344.632}{563.385}{L}%
\TextAt{8}{352.597}{563.385}{,}%
\TextAt{8}{359.746}{563.385}{then}%
\TextAt{8}{386.406}{563.385}{one}%
\TextAt{8}{407.866}{563.385}{has}%
\TextAt{16}{163.816}{538.477}{\Character{106}}%
\TextAt{13}{167.137}{538.477}{x}%
\TextAt{16}{173.789}{538.477}{\Character{106}}%
\TextAt{21}{177.11}{534.99}{w}%
\TextAt{3}{183.401}{534.99}{\Character{14}}%
\TextAt{21}{187.635}{534.99}{\Character{27}}%
\TextAt{9}{196.579}{538.477}{\Character{61}}%
\TextAt{16}{209.005}{538.477}{\Character{106}}%
\TextAt{13}{212.326}{538.477}{\Character{27}}%
\TextAt{13}{219.392}{538.477}{x}%
\TextAt{16}{226.06}{538.477}{\Character{106}}%
\TextAt{21}{229.381}{534.99}{w}%
\TextAt{9}{239.491}{538.477}{\Character{61}}%
\TextAt{20}{251.917}{548.639}{\Character{12}}%
\TextAt{20}{251.917}{541.466}{\Character{12}}%
\TextAt{13}{255.902}{538.477}{\Character{28}}%
\TextAt{3}{262.321}{543.413}{\Character{0}}%
\TextAt{17}{268.907}{543.413}{1}%
\TextAt{13}{273.64}{538.477}{\Character{27}}%
\TextAt{3}{280.722}{543.413}{\Character{48}}%
\TextAt{13}{283.517}{538.477}{\Character{28}}%
\TextAt{13}{289.925}{538.477}{x}%
\TextAt{20}{296.588}{548.639}{\Character{12}}%
\TextAt{20}{296.588}{541.466}{\Character{12}}%
\TextAt{21}{300.573}{533.795}{w}%
\TextAt{9}{310.684}{538.477}{\Character{61}}%
\TextAt{16}{323.109}{538.477}{\Character{106}}%
\TextAt{13}{326.43}{538.477}{\Character{27}}%
\TextAt{3}{333.512}{543.413}{\Character{48}}%
\TextAt{13}{336.307}{538.477}{\Character{28}}%
\TextAt{13}{342.715}{538.477}{x}%
\TextAt{16}{349.378}{538.477}{\Character{106}}%
\TextAt{21}{352.699}{534.99}{w}%
\TextAt{23}{358.991}{537.26}{\Character{48}}%
\TextAt{9}{365.494}{538.477}{\Character{61}}%
\TextAt{16}{377.919}{538.477}{\Character{106}}%
\TextAt{13}{381.24}{538.477}{\Character{28}}%
\TextAt{13}{387.648}{538.477}{x}%
\TextAt{16}{394.311}{538.477}{\Character{106}}%
\TextAt{21}{397.632}{534.99}{w}%
\TextAt{23}{403.924}{537.26}{\Character{48}}%
\TextAt{9}{410.427}{538.477}{\Character{61}}%
\TextAt{16}{422.852}{538.477}{\Character{106}}%
\TextAt{13}{426.173}{538.477}{x}%
\TextAt{16}{432.825}{538.477}{\Character{106}}%
\TextAt{21}{436.146}{534.99}{w}%
\TextAt{13}{444.928}{538.477}{\Character{59}}%
\TextAt{8}{79.226}{513.569}{so}%
\TextAt{8}{93.596}{513.569}{that}%
\TextAt{13}{118.965}{513.569}{\Character{27}}%
\NativeMathAt{129.368}{513.569}{11.955}{\in}%
\TextAt{13}{140.659}{513.569}{G}%
\TextAt{21}{149.893}{511.776}{w}%
\TextAt{9}{156.682}{513.569}{\Character{40}}%
\TextAt{13}{161.235}{513.569}{L}%
\TextAt{16}{172.52}{513.569}{\Character{106}}%
\TextAt{13}{179.162}{513.569}{K}%
\TextAt{9}{189.979}{513.569}{\Character{41}}%
\TextAt{8}{194.532}{513.569}{.}%
\TextAt{8}{202.98}{513.569}{If}%
\TextAt{13}{214.672}{513.569}{\Character{27}}%
\TextAt{3}{221.755}{517.908}{\Character{48}}%
\NativeMathAt{227.871}{513.569}{11.955}{\in}%
\TextAt{13}{239.162}{513.569}{I}%
\TextAt{21}{244.324}{511.776}{w}%
\TextAt{23}{250.616}{514.045}{\Character{48}}%
\TextAt{9}{253.798}{513.569}{\Character{40}}%
\TextAt{13}{258.351}{513.569}{L}%
\TextAt{3}{266.315}{517.908}{\Character{48}}%
\TextAt{16}{272.431}{513.569}{\Character{106}}%
\TextAt{13}{279.073}{513.569}{K}%
\TextAt{3}{289.89}{517.908}{\Character{48}}%
\TextAt{9}{292.685}{513.569}{\Character{41}}%
\TextAt{8}{301.135}{513.569}{and}%
\TextAt{13}{323.898}{513.569}{x}%
\NativeMathAt{333.871}{513.569}{11.955}{\in}%
\TextAt{16}{345.145}{513.569}{\Character{79}}%
\TextAt{8}{355.011}{513.569}{,}%
\TextAt{8}{362.161}{513.569}{then}%
\TextAt{13}{157.226}{488.661}{w}%
\TextAt{9}{165.916}{488.661}{\Character{40}}%
\TextAt{13}{170.468}{488.661}{\Character{27}}%
\TextAt{13}{177.534}{488.661}{x}%
\TextAt{16}{186.86}{488.661}{\Character{0}}%
\TextAt{13}{198.815}{488.661}{x}%
\TextAt{9}{205.467}{488.661}{\Character{41}}%
\TextAt{9}{213.322}{488.661}{\Character{61}}%
\TextAt{13}{225.765}{488.661}{w}%
\TextAt{20}{234.455}{498.345}{\Character{0}}%
\TextAt{13}{239.934}{488.661}{\Character{28}}%
\TextAt{3}{246.353}{493.598}{\Character{0}}%
\TextAt{17}{252.94}{493.598}{1}%
\TextAt{9}{257.672}{488.661}{\Character{40}}%
\TextAt{13}{262.225}{488.661}{\Character{27}}%
\TextAt{3}{269.307}{493.598}{\Character{48}}%
\TextAt{13}{272.102}{488.661}{\Character{28}}%
\TextAt{13}{278.51}{488.661}{x}%
\TextAt{16}{287.83}{488.661}{\Character{0}}%
\TextAt{13}{299.785}{488.661}{\Character{28}}%
\TextAt{13}{306.193}{488.661}{x}%
\TextAt{9}{312.856}{488.661}{\Character{41}}%
\TextAt{20}{317.408}{498.345}{\Character{1}}%
\TextAt{9}{326.209}{488.661}{\Character{61}}%
\TextAt{13}{338.634}{488.661}{w}%
\TextAt{3}{347.324}{493.598}{\Character{48}}%
\TextAt{20}{350.119}{498.345}{\Character{0}}%
\TextAt{13}{355.598}{488.661}{\Character{27}}%
\TextAt{3}{362.681}{493.598}{\Character{48}}%
\TextAt{9}{365.476}{488.661}{\Character{40}}%
\TextAt{13}{370.028}{488.661}{\Character{28}}%
\TextAt{13}{376.436}{488.661}{x}%
\TextAt{9}{383.099}{488.661}{\Character{41}}%
\TextAt{16}{390.308}{488.661}{\Character{0}}%
\TextAt{9}{402.263}{488.661}{\Character{40}}%
\TextAt{13}{406.815}{488.661}{\Character{28}}%
\TextAt{13}{413.223}{488.661}{x}%
\TextAt{9}{419.886}{488.661}{\Character{41}}%
\TextAt{20}{424.439}{498.345}{\Character{1}}%
\TextAt{13}{433.239}{488.661}{\Character{62}}%
\TextAt{9}{445.665}{488.661}{0}%
\TextAt{13}{451.518}{488.661}{\Character{59}}%
\TextAt{8}{79.226}{463.754}{and}%
\TextAt{13}{101.989}{463.754}{\Character{27}}%
\NativeMathAt{112.392}{463.754}{11.955}{\in}%
\TextAt{13}{123.683}{463.754}{I}%
\TextAt{21}{128.846}{461.96}{w}%
\TextAt{9}{135.635}{463.754}{\Character{40}}%
\TextAt{13}{140.188}{463.754}{L}%
\TextAt{16}{151.473}{463.754}{\Character{106}}%
\TextAt{13}{158.115}{463.754}{K}%
\TextAt{9}{168.932}{463.754}{\Character{41}}%
\TextAt{8}{173.485}{463.754}{.}%
\TextAt{8}{181.933}{463.754}{If}%
\TextAt{13}{193.625}{463.754}{\Character{27}}%
\TextAt{3}{200.708}{468.092}{\Character{48}}%
\NativeMathAt{206.824}{463.754}{11.955}{\in}%
\TextAt{13}{218.115}{463.754}{R}%
\TextAt{21}{227.025}{461.96}{w}%
\TextAt{23}{233.316}{464.23}{\Character{48}}%
\TextAt{9}{236.498}{463.754}{\Character{40}}%
\TextAt{13}{241.051}{463.753}{L}%
\TextAt{3}{249.015}{468.092}{\Character{48}}%
\TextAt{16}{255.131}{463.754}{\Character{106}}%
\TextAt{13}{261.773}{463.753}{K}%
\TextAt{3}{272.591}{468.092}{\Character{48}}%
\TextAt{9}{275.386}{463.754}{\Character{41}}%
\TextAt{8}{283.835}{463.754}{and}%
\TextAt{13}{306.598}{463.754}{x}%
\NativeMathAt{316.571}{463.754}{11.955}{\in}%
\TextAt{13}{327.862}{463.754}{L}%
\TextAt{3}{335.826}{468.092}{\Character{3}}%
\TextAt{8}{340.559}{463.754}{,}%
\TextAt{8}{347.708}{463.754}{then}%
\TextAt{13}{152.9}{434.571}{w}%
\TextAt{20}{163.582}{447.842}{\Character{16}}%
\TextAt{13}{171.917}{442.659}{\Character{27}}%
\TextAt{13}{178.983}{442.659}{x}%
\Vector{q 0 G 0 g 0.478 w 99.917 -336.44 m 113.652 -336.44 l S Q}%
\TextAt{13}{175.458}{426.371}{x}%
\TextAt{16}{189.504}{434.571}{\Character{0}}%
\TextAt{9}{201.459}{434.571}{1}%
\TextAt{20}{207.312}{447.842}{\Character{17}}%
\TextAt{9}{217.773}{434.571}{\Character{61}}%
\TextAt{13}{230.198}{434.571}{w}%
\TextAt{20}{240.88}{451.428}{\Character{18}}%
\TextAt{13}{249.681}{434.571}{\Character{28}}%
\TextAt{3}{256.1}{439.508}{\Character{0}}%
\TextAt{17}{262.686}{439.508}{1}%
\TextAt{20}{269.411}{451.428}{\Character{18}}%
\TextAt{13}{279.407}{442.659}{\Character{27}}%
\TextAt{3}{286.489}{446.998}{\Character{48}}%
\TextAt{13}{289.284}{442.659}{\Character{28}}%
\TextAt{13}{295.692}{442.659}{x}%
\Vector{q 0 G 0 g 0.478 w 207.407 -336.44 m 230.355 -336.44 l S Q}%
\TextAt{13}{284.346}{426.371}{\Character{28}}%
\TextAt{13}{290.754}{426.371}{x}%
\TextAt{16}{306.208}{434.571}{\Character{0}}%
\TextAt{9}{318.163}{434.571}{1}%
\TextAt{20}{324.016}{451.428}{\Character{19}\Character{19}}%
\TextAt{9}{344.937}{434.571}{\Character{61}}%
\TextAt{13}{357.363}{434.571}{w}%
\TextAt{3}{366.052}{439.508}{\Character{48}}%
\TextAt{20}{370.84}{451.428}{\Character{18}}%
\TextAt{13}{380.836}{442.659}{\Character{27}}%
\TextAt{3}{387.918}{446.998}{\Character{48}}%
\TextAt{13}{390.713}{442.659}{\Character{28}}%
\TextAt{13}{397.121}{442.659}{x}%
\Vector{q 0 G 0 g 0.478 w 308.836 -336.44 m 331.784 -336.44 l S Q}%
\TextAt{13}{385.774}{426.371}{\Character{28}}%
\TextAt{13}{392.182}{426.371}{x}%
\TextAt{16}{407.636}{434.571}{\Character{0}}%
\TextAt{9}{419.592}{434.571}{1}%
\TextAt{20}{425.445}{451.428}{\Character{19}}%
\TextAt{13}{437.566}{434.571}{\Character{62}}%
\TextAt{9}{449.991}{434.571}{0}%
\TextAt{13}{455.844}{434.571}{\Character{59}}%
\TextAt{8}{79.226}{405.389}{so}%
\TextAt{8}{93.596}{405.389}{that}%
\TextAt{13}{118.965}{405.389}{\Character{27}}%
\NativeMathAt{129.368}{405.389}{11.955}{\in}%
\TextAt{13}{140.659}{405.389}{R}%
\TextAt{21}{149.569}{403.596}{w}%
\TextAt{9}{156.359}{405.389}{\Character{40}}%
\TextAt{13}{160.911}{405.389}{L}%
\TextAt{16}{172.197}{405.389}{\Character{106}}%
\TextAt{13}{178.838}{405.389}{K}%
\TextAt{9}{189.656}{405.389}{\Character{41}}%
\TextAt{8}{194.208}{405.389}{.}%
\Vector{q 0 G 0 g 0.399 w 452.334 -368.611 m 452.334 -360.541 l S Q}%
\Vector{q 0 G 0 g 0.399 w 452.534 -360.74 m 459.707 -360.74 l S Q}%
\Vector{q 0 G 0 g 0.399 w 452.534 -368.411 m 459.707 -368.411 l S Q}%
\Vector{q 0 G 0 g 0.399 w 459.906 -368.611 m 459.906 -360.541 l S Q}%
\TextAt{8}{97.159}{377.558}{If}%
\TextAt{8}{109.358}{377.558}{the}%
\TextAt{8}{130.022}{377.558}{two}%
\TextAt{8}{152.646}{377.558}{homomorphisms}%
\TextAt{13}{240.438}{377.558}{\Character{28}}%
\TextAt{9}{251.042}{377.558}{\Character{58}}%
\TextAt{13}{258.479}{377.558}{L}%
\NativeMathAt{270.629}{377.558}{11.955}{\rightarrow}%
\TextAt{13}{286.769}{377.558}{L}%
\TextAt{3}{294.734}{381.896}{\Character{48}}%
\TextAt{8}{301.934}{377.558}{and}%
\TextAt{13}{325.203}{377.558}{\Character{28}}%
\TextAt{9}{335.808}{377.558}{\Character{58}}%
\TextAt{13}{343.244}{377.558}{K}%
\NativeMathAt{358.247}{377.558}{11.955}{\rightarrow}%
\TextAt{13}{374.388}{377.558}{K}%
\TextAt{3}{385.205}{381.896}{\Character{48}}%
\TextAt{8}{392.405}{377.558}{are}%
\TextAt{8}{412.423}{377.558}{isomorphisms,}%
\TextAt{8}{489.343}{377.558}{then}%
\TextAt{8}{516.51}{377.558}{the}%
\TextAt{8}{79.226}{360.512}{homomorphisms}%
\TextAt{8}{167.601}{360.512}{(9.4)}%
\TextAt{8}{196.666}{360.512}{are}%
\TextAt{8}{217.267}{360.512}{of}%
\TextAt{8}{231.687}{360.512}{course}%
\TextAt{8}{268.606}{360.512}{isomorphisms.}%
\TextAt{8}{349.461}{360.512}{In}%
\TextAt{8}{365.172}{360.512}{particular,}%
\TextAt{8}{423.776}{360.512}{in}%
\TextAt{8}{438.519}{360.512}{the}%
\TextAt{8}{459.765}{360.512}{case}%
\TextAt{13}{485.626}{360.512}{K}%
\TextAt{9}{501.622}{360.512}{\Character{61}}%
\TextAt{13}{515.905}{360.512}{K}%
\TextAt{3}{526.723}{364.85}{\Character{48}}%
\TextAt{8}{529.518}{360.512}{,}%
\TextAt{13}{79.226}{343.466}{L}%
\TextAt{9}{90.511}{343.466}{\Character{61}}%
\TextAt{13}{102.937}{343.466}{L}%
\TextAt{3}{110.901}{347.804}{\Character{48}}%
\TextAt{8}{113.697}{343.466}{,}%
\TextAt{8}{120.846}{343.466}{we}%
\TextAt{8}{138.073}{343.466}{find}%
\TextAt{8}{161.481}{343.466}{for}%
\TextAt{8}{179.366}{343.466}{each}%
\TextAt{13}{205.704}{343.466}{\Character{28}}%
\NativeMathAt{215.443}{343.466}{11.955}{\in}%
\TextAt{13}{226.734}{343.466}{G}%
\TextAt{9}{235.968}{343.466}{\Character{40}}%
\TextAt{13}{240.52}{343.466}{L}%
\TextAt{16}{251.806}{343.466}{\Character{106}}%
\TextAt{13}{258.447}{343.466}{K}%
\TextAt{9}{269.265}{343.466}{\Character{41}}%
\TextAt{8}{273.817}{343.466}{:}%
\TextAt{13}{159.295}{318.558}{G}%
\TextAt{21}{168.528}{316.765}{w}%
\TextAt{3}{174.82}{316.765}{\Character{14}}%
\TextAt{21}{179.054}{316.765}{\Character{28}}%
\TextAt{9}{187.551}{318.558}{\Character{61}}%
\TextAt{13}{199.977}{318.558}{\Character{28}}%
\TextAt{3}{206.396}{323.494}{\Character{0}}%
\TextAt{17}{212.982}{323.494}{1}%
\TextAt{13}{217.714}{318.558}{G}%
\TextAt{21}{226.948}{316.765}{w}%
\TextAt{13}{233.737}{318.558}{\Character{28}}%
\TextAt{13}{239.499}{318.558}{\Character{59}}%
\TextAt{13}{268.634}{318.558}{I}%
\TextAt{21}{273.823}{316.765}{w}%
\TextAt{3}{280.115}{316.765}{\Character{14}}%
\TextAt{21}{284.349}{316.765}{\Character{28}}%
\TextAt{9}{292.846}{318.558}{\Character{61}}%
\TextAt{13}{305.272}{318.558}{\Character{28}}%
\TextAt{3}{311.691}{323.494}{\Character{0}}%
\TextAt{17}{318.277}{323.494}{1}%
\TextAt{13}{323.009}{318.558}{I}%
\TextAt{21}{328.172}{316.765}{w}%
\TextAt{13}{334.962}{318.558}{\Character{28}}%
\TextAt{13}{340.724}{318.558}{\Character{59}}%
\TextAt{13}{369.859}{318.558}{R}%
\TextAt{21}{378.795}{316.765}{w}%
\TextAt{3}{385.086}{316.765}{\Character{14}}%
\TextAt{21}{389.321}{316.765}{\Character{28}}%
\TextAt{9}{397.818}{318.558}{\Character{61}}%
\TextAt{13}{410.243}{318.558}{\Character{28}}%
\TextAt{3}{416.662}{323.494}{\Character{0}}%
\TextAt{17}{423.249}{323.494}{1}%
\TextAt{13}{427.981}{318.558}{R}%
\TextAt{21}{436.891}{316.765}{w}%
\TextAt{13}{443.681}{318.558}{\Character{28}}%
\TextAt{13}{449.443}{318.558}{\Character{59}}%
\TextAt{8}{79.226}{293.65}{i.e.,}%
\TextAt{8}{101.541}{293.65}{the}%
\TextAt{8}{121.865}{293.65}{decomposition,}%
\TextAt{8}{202.462}{293.65}{inertia,}%
\TextAt{8}{242.996}{293.65}{and}%
\TextAt{8}{265.927}{293.65}{ramification}%
\TextAt{8}{331.143}{293.65}{groups}%
\TextAt{8}{369.101}{293.65}{of}%
\TextAt{8}{382.599}{293.65}{conjugate}%
\TextAt{8}{435.776}{293.65}{valuations}%
\TextAt{8}{491.631}{293.65}{are}%
\TextAt{8}{511.31}{293.65}{con-}%
\TextAt{8}{79.226}{276.604}{jugate.}%
\TextAt{8}{97.159}{256.569}{Another}%
\TextAt{8}{143.007}{256.569}{special}%
\TextAt{8}{181.108}{256.569}{case}%
\TextAt{8}{205.879}{256.569}{arises}%
\TextAt{8}{237.871}{256.569}{from}%
\TextAt{8}{265.512}{256.569}{an}%
\TextAt{8}{281.771}{256.569}{intermediate}%
\TextAt{8}{349.736}{256.569}{field}%
\TextAt{13}{375.093}{256.569}{M}%
\TextAt{8}{391.564}{256.569}{of}%
\TextAt{13}{404.894}{256.569}{L}%
\TextAt{16}{416.179}{256.569}{\Character{106}}%
\TextAt{13}{422.821}{256.569}{K}%
\TextAt{8}{437.536}{256.569}{by}%
\TextAt{8}{453.795}{256.569}{the}%
\TextAt{8}{473.951}{256.569}{diagram}%
\TextAt{13}{267.43}{232.173}{L}%
\TextAt{13}{332.47}{232.173}{L}%
\TextAt{13}{266.003}{182.965}{K}%
\TextAt{13}{330.165}{182.965}{M}%
\NativeDiagram{%
\path[draw=black,line width=2.797bp] (208.747,-538.838) -- (255.119,-538.838);%
\path[draw=black!0,line width=1.841bp] (208.747,-538.838) -- (255.119,-538.838);%
\path[draw=black,line width=0.478bp] (199.413,-578.286) -- (199.413,-546.406);%
\path[draw=black,line width=0.478bp] (211.895,-588.046) -- (252.336,-588.046);%
\path[draw=black,line width=0.478bp,line cap=round] (211.895,-585.703) .. controls (211.077,-585.703) and (210.413,-586.227) .. (210.413,-586.875) .. controls (210.413,-587.522) and (211.077,-588.046) .. (211.895,-588.046);%
\path[draw=black,line width=0.478bp,line cap=round,line join=round] (250.089,-585.177) .. controls (250.543,-586.898) and (251.555,-587.711) .. (252.575,-588.046) .. controls (251.555,-588.381) and (250.543,-589.194) .. (250.089,-590.915);%
\path[draw=black,line width=0.478bp] (264.453,-578.286) -- (264.453,-546.406);%
}%
\TextAt{13}{347.892}{211.654}{\Character{58}}%
\TextAt{13}{79.226}{160.761}{\Character{28}}%
\TextAt{3}{85.645}{165.099}{\Character{3}}%
\TextAt{8}{94.275}{160.761}{then}%
\TextAt{8}{120.935}{160.761}{becomes}%
\TextAt{8}{167.488}{160.761}{the}%
\TextAt{8}{187.645}{160.761}{inclusion}%
\TextAt{13}{236.482}{160.761}{G}%
\TextAt{9}{245.715}{160.761}{\Character{40}}%
\TextAt{13}{250.267}{160.761}{L}%
\TextAt{16}{261.553}{160.761}{\Character{106}}%
\TextAt{13}{268.195}{160.761}{M}%
\TextAt{9}{280.768}{160.761}{\Character{41}}%
\NativeArrowAt{288.641}{160.761}{11.955}{13.214}{\hookrightarrow}%
\TextAt{13}{305.176}{160.761}{G}%
\TextAt{9}{314.41}{160.761}{\Character{40}}%
\TextAt{13}{318.962}{160.761}{L}%
\TextAt{16}{330.248}{160.761}{\Character{106}}%
\TextAt{13}{336.889}{160.761}{K}%
\TextAt{9}{347.707}{160.761}{\Character{41}}%
\TextAt{8}{352.26}{160.761}{,}%
\TextAt{8}{359.409}{160.761}{and}%
\TextAt{8}{382.171}{160.761}{we}%
\TextAt{8}{399.399}{160.761}{trivially}%
\TextAt{8}{443.633}{160.761}{get}%
\TextAt{8}{463.144}{160.761}{the}%
\TextAt{12}{79.226}{135.918}{Proposition}%
\TextAt{12}{151.937}{135.918}{(9.5).}%
\TextAt{8}{187.253}{135.918}{For}%
\TextAt{8}{208.222}{135.918}{the}%
\TextAt{8}{228.379}{135.918}{extensions}%
\TextAt{13}{284.759}{135.918}{K}%
\NativeMathAt{298.898}{135.918}{11.955}{\subseteq}%
\TextAt{13}{311.517}{135.918}{M}%
\NativeMathAt{327.412}{135.918}{11.955}{\subseteq}%
\TextAt{13}{340.031}{135.918}{L}%
\TextAt{8}{347.996}{135.918}{,}%
\TextAt{8}{355.145}{135.918}{one}%
\TextAt{8}{376.604}{135.918}{has}%
\TextAt{13}{210.687}{112.127}{G}%
\TextAt{21}{219.921}{110.334}{w}%
\TextAt{9}{226.71}{112.127}{\Character{40}}%
\TextAt{13}{231.263}{112.127}{L}%
\TextAt{16}{242.548}{112.127}{\Character{106}}%
\TextAt{13}{249.19}{112.127}{M}%
\TextAt{9}{261.763}{112.127}{\Character{41}}%
\TextAt{9}{269.618}{112.127}{\Character{61}}%
\TextAt{13}{282.062}{112.127}{G}%
\TextAt{21}{291.296}{110.334}{w}%
\TextAt{9}{298.085}{112.127}{\Character{40}}%
\TextAt{13}{302.637}{112.127}{L}%
\TextAt{16}{313.923}{112.127}{\Character{106}}%
\TextAt{13}{320.564}{112.127}{K}%
\TextAt{9}{331.382}{112.127}{\Character{41}}%
\TextAt{16}{338.591}{112.127}{\Character{92}}%
\TextAt{13}{349.218}{112.127}{G}%
\TextAt{9}{358.451}{112.127}{\Character{40}}%
\TextAt{13}{363.004}{112.127}{L}%
\TextAt{16}{374.289}{112.127}{\Character{106}}%
\TextAt{13}{380.931}{112.127}{M}%
\TextAt{9}{393.504}{112.127}{\Character{41}}%
\TextAt{13}{398.057}{112.127}{\Character{59}}%
\TextAt{13}{214.758}{92.092}{I}%
\TextAt{21}{219.921}{90.299}{w}%
\TextAt{9}{226.71}{92.092}{\Character{40}}%
\TextAt{13}{231.263}{92.092}{L}%
\TextAt{16}{242.548}{92.092}{\Character{106}}%
\TextAt{13}{249.19}{92.092}{M}%
\TextAt{9}{261.763}{92.092}{\Character{41}}%
\TextAt{9}{269.618}{92.092}{\Character{61}}%
\TextAt{13}{282.062}{92.092}{I}%
\TextAt{21}{287.225}{90.299}{w}%
\TextAt{9}{294.014}{92.092}{\Character{40}}%
\TextAt{13}{298.567}{92.092}{L}%
\TextAt{16}{309.852}{92.092}{\Character{106}}%
\TextAt{13}{316.494}{92.092}{K}%
\TextAt{9}{327.311}{92.092}{\Character{41}}%
\TextAt{16}{334.52}{92.092}{\Character{92}}%
\TextAt{13}{345.147}{92.092}{G}%
\TextAt{9}{354.381}{92.092}{\Character{40}}%
\TextAt{13}{358.933}{92.092}{L}%
\TextAt{16}{370.218}{92.092}{\Character{106}}%
\TextAt{13}{376.86}{92.092}{M}%
\TextAt{9}{389.434}{92.092}{\Character{41}}%
\TextAt{13}{393.986}{92.092}{\Character{59}}%
\TextAt{13}{211.011}{72.058}{R}%
\TextAt{21}{219.921}{70.264}{w}%
\TextAt{9}{226.71}{72.058}{\Character{40}}%
\TextAt{13}{231.263}{72.058}{L}%
\TextAt{16}{242.548}{72.058}{\Character{106}}%
\TextAt{13}{249.19}{72.058}{M}%
\TextAt{9}{261.763}{72.058}{\Character{41}}%
\TextAt{9}{269.618}{72.058}{\Character{61}}%
\TextAt{13}{282.062}{72.058}{R}%
\TextAt{21}{290.972}{70.264}{w}%
\TextAt{9}{297.762}{72.058}{\Character{40}}%
\TextAt{13}{302.314}{72.058}{L}%
\TextAt{16}{313.599}{72.058}{\Character{106}}%
\TextAt{13}{320.241}{72.058}{K}%
\TextAt{9}{331.059}{72.058}{\Character{41}}%
\TextAt{16}{338.268}{72.058}{\Character{92}}%
\TextAt{13}{348.894}{72.058}{G}%
\TextAt{9}{358.128}{72.058}{\Character{40}}%
\TextAt{13}{362.68}{72.058}{L}%
\TextAt{16}{373.966}{72.058}{\Character{106}}%
\TextAt{13}{380.607}{72.058}{M}%
\TextAt{9}{393.181}{72.058}{\Character{41}}%
\TextAt{13}{397.733}{72.058}{\Character{58}}%
\TextAt{8}{297.211}{43.033}{155}%
\Destination{78.23}{809.815}{nav.page.155}%
\Destination{79.36}{776.155}{nav.print.169}%
\Destination{79.36}{696.155}{nav.number.2.9.4}%
\Destination{79.36}{149.955}{nav.number.2.9.5}%
\Hotlink{167.101}{358.355}{25.078}{11.047}{nav.number.2.9.4}%
}

% Page 168 (body page 156)
\Page{%
\TextAt{2}{79.226}{796.943}{II.}%
\TextAt{2}{94.073}{796.943}{The}%
\TextAt{2}{116.492}{796.943}{Theory}%
\TextAt{2}{154.401}{796.943}{of}%
\TextAt{2}{166.826}{796.943}{Valuations}%
\Vector{q 0 G 0 g 0.299 w 7.226 18.729 m 460.769 18.729 l S Q}%
\TextAt{8}{97.159}{761.124}{A}%
\TextAt{8}{109.831}{761.124}{particularly}%
\TextAt{8}{173.254}{761.124}{important}%
\TextAt{8}{228.546}{761.124}{special}%
\TextAt{8}{266.647}{761.124}{case}%
\TextAt{8}{291.418}{761.124}{of}%
\TextAt{8}{304.748}{761.124}{(9.4)}%
\TextAt{8}{332.724}{761.124}{occurs}%
\TextAt{8}{368.876}{761.124}{with}%
\TextAt{8}{395.536}{761.124}{the}%
\TextAt{8}{415.692}{761.124}{diagram}%
\TextAt{13}{267.91}{732.895}{L}%
\TextAt{13}{330.702}{732.895}{L}%
\TextAt{21}{338.666}{731.102}{w}%
\TextAt{13}{266.483}{681.894}{K}%
\TextAt{13}{330.646}{681.894}{K}%
\TextAt{21}{340.605}{680.101}{v}%
\NativeDiagram{%
\path[draw=black,line width=0.478bp] (209.226,-38.115) -- (253.35,-38.115);%
\path[draw=black,line width=0.478bp] (199.892,-79.357) -- (199.892,-45.683);%
\path[draw=black,line width=0.478bp] (210.653,-89.117) -- (253.294,-89.117);%
\path[draw=black,line width=0.478bp] (266.079,-79.357) -- (266.079,-47.477);%
}%
\TextAt{8}{79.226}{654.077}{which}%
\TextAt{8}{112.395}{654.077}{can}%
\TextAt{8}{133.537}{654.077}{be}%
\TextAt{8}{149.143}{654.077}{associated}%
\TextAt{8}{204.561}{654.077}{to}%
\TextAt{8}{218.553}{654.077}{any}%
\TextAt{8}{240.353}{654.077}{extension}%
\TextAt{8}{291.801}{654.077}{of}%
\TextAt{8}{304.813}{654.077}{valuations}%
\TextAt{13}{360.183}{654.077}{w}%
\TextAt{16}{372.193}{654.077}{\Character{106}}%
\TextAt{13}{378.835}{654.077}{v}%
\TextAt{8}{388.503}{654.077}{of}%
\TextAt{13}{401.515}{654.077}{L}%
\TextAt{16}{412.801}{654.077}{\Character{106}}%
\TextAt{13}{419.442}{654.077}{K}%
\TextAt{8}{430.26}{654.077}{.}%
\TextAt{8}{438.603}{654.077}{If}%
\TextAt{13}{449.977}{654.077}{L}%
\TextAt{16}{461.262}{654.077}{\Character{106}}%
\TextAt{13}{467.904}{654.077}{K}%
\TextAt{8}{482.302}{654.077}{is}%
\TextAt{8}{493.748}{654.077}{infinite,}%
\TextAt{8}{79.226}{637.031}{then}%
\TextAt{13}{106.593}{637.031}{L}%
\TextAt{21}{114.558}{635.237}{w}%
\TextAt{8}{125.951}{637.031}{has}%
\TextAt{8}{147.532}{637.031}{to}%
\TextAt{8}{162.549}{637.031}{be}%
\TextAt{8}{179.18}{637.031}{read}%
\TextAt{8}{205.902}{637.031}{as}%
\TextAt{8}{220.979}{637.031}{the}%
\TextAt{8}{241.842}{637.031}{localization}%
\TextAt{8}{304.668}{637.031}{in}%
\TextAt{8}{319.028}{637.031}{the}%
\TextAt{8}{339.891}{637.031}{sense}%
\TextAt{8}{370.629}{637.031}{of}%
\TextAt{8}{384.666}{637.031}{§8,}%
\TextAt{8}{404.236}{637.031}{p.}%
\TextAt{8}{418.595}{637.031}{160.}%
\TextAt{8}{446.739}{637.031}{(This}%
\TextAt{8}{478.72}{637.031}{distinction}%
\TextAt{8}{79.226}{619.985}{is}%
\TextAt{8}{91.5}{619.985}{rendered}%
\TextAt{8}{140.13}{619.985}{superfluous}%
\TextAt{8}{202.222}{619.985}{if}%
\TextAt{8}{213.456}{619.985}{we}%
\TextAt{8}{231.194}{619.985}{consider,}%
\TextAt{8}{280.669}{619.985}{as}%
\TextAt{8}{295.549}{619.985}{we}%
\TextAt{8}{313.287}{619.985}{may}%
\TextAt{8}{339.166}{619.985}{perfectly}%
\TextAt{8}{388.119}{619.985}{well}%
\TextAt{8}{412.36}{619.985}{do,}%
\TextAt{8}{432.509}{619.985}{the}%
\TextAt{8}{453.176}{619.985}{henselization}%
\TextAt{8}{523.337}{619.985}{of}%
\TextAt{13}{79.226}{602.939}{L}%
\TextAt{16}{91.988}{602.939}{\Character{106}}%
\TextAt{13}{100.106}{602.939}{K}%
\TextAt{8}{110.924}{602.939}{.)}%
\TextAt{8}{126.526}{602.939}{Since}%
\TextAt{8}{157.95}{602.939}{in}%
\TextAt{8}{172.469}{602.939}{the}%
\TextAt{8}{193.492}{602.939}{local}%
\TextAt{8}{221.998}{602.939}{extension}%
\TextAt{13}{274.631}{602.939}{L}%
\TextAt{21}{282.595}{601.145}{w}%
\TextAt{16}{294.182}{602.939}{\Character{106}}%
\TextAt{13}{302.3}{602.939}{K}%
\TextAt{21}{312.259}{601.145}{v}%
\TextAt{8}{321.93}{602.939}{the}%
\TextAt{8}{342.953}{602.939}{extension}%
\TextAt{8}{395.585}{602.939}{of}%
\TextAt{8}{409.782}{602.939}{the}%
\TextAt{8}{430.805}{602.939}{valuation}%
\TextAt{8}{482.743}{602.939}{is}%
\TextAt{8}{495.374}{602.939}{unique,}%
\TextAt{8}{79.226}{585.893}{we}%
\TextAt{8}{97.693}{585.893}{denote}%
\TextAt{8}{136.651}{585.893}{the}%
\TextAt{8}{158.047}{585.893}{decomposition,}%
\TextAt{8}{239.983}{585.893}{inertia,}%
\TextAt{8}{281.858}{585.893}{and}%
\TextAt{8}{305.86}{585.893}{ramification}%
\TextAt{8}{372.148}{585.893}{groups}%
\TextAt{8}{411.178}{585.893}{simply}%
\TextAt{8}{449.873}{585.893}{by}%
\TextAt{13}{467.371}{585.893}{G}%
\TextAt{9}{476.605}{585.893}{\Character{40}}%
\TextAt{13}{481.157}{585.893}{L}%
\TextAt{21}{489.122}{584.099}{w}%
\TextAt{16}{501.344}{585.893}{\Character{106}}%
\TextAt{13}{510.099}{585.893}{K}%
\TextAt{21}{520.058}{584.099}{v}%
\TextAt{9}{524.965}{585.893}{\Character{41}}%
\TextAt{8}{529.518}{585.893}{,}%
\TextAt{13}{79.226}{568.847}{I}%
\TextAt{9}{85.329}{568.847}{\Character{40}}%
\TextAt{13}{89.881}{568.847}{L}%
\TextAt{21}{97.846}{567.053}{w}%
\TextAt{16}{107.956}{568.847}{\Character{106}}%
\TextAt{13}{114.598}{568.847}{K}%
\TextAt{21}{124.557}{567.053}{v}%
\TextAt{9}{129.464}{568.847}{\Character{41}}%
\TextAt{8}{134.017}{568.847}{,}%
\TextAt{13}{141.166}{568.847}{R}%
\TextAt{9}{150.174}{568.847}{\Character{40}}%
\TextAt{13}{154.727}{568.847}{L}%
\TextAt{21}{162.691}{567.053}{w}%
\TextAt{16}{172.802}{568.847}{\Character{106}}%
\TextAt{13}{179.443}{568.847}{K}%
\TextAt{21}{189.402}{567.053}{v}%
\TextAt{9}{194.31}{568.847}{\Character{41}}%
\TextAt{8}{198.862}{568.847}{.}%
\TextAt{8}{207.311}{568.847}{In}%
\TextAt{8}{221.932}{568.847}{this}%
\TextAt{8}{244.754}{568.847}{case,}%
\TextAt{8}{272.777}{568.847}{the}%
\TextAt{8}{292.934}{568.847}{homomorphism}%
\TextAt{13}{375.604}{568.847}{\Character{28}}%
\TextAt{3}{382.023}{573.185}{\Character{3}}%
\TextAt{8}{390.652}{568.847}{is}%
\TextAt{8}{402.416}{568.847}{the}%
\TextAt{8}{422.573}{568.847}{restriction}%
\TextAt{8}{478.571}{568.847}{map}%
\TextAt{13}{199.518}{540.119}{G}%
\TextAt{9}{208.752}{540.119}{\Character{40}}%
\TextAt{13}{213.304}{540.119}{L}%
\TextAt{21}{221.269}{538.325}{w}%
\TextAt{16}{231.379}{540.119}{\Character{106}}%
\TextAt{13}{238.021}{540.119}{K}%
\TextAt{21}{247.98}{538.325}{v}%
\TextAt{9}{252.887}{540.119}{\Character{41}}%
\NativeArrowAt{260.76}{540.119}{11.955}{19.26}{\longrightarrow}%
\TextAt{13}{283.342}{540.119}{G}%
\TextAt{9}{292.576}{540.119}{\Character{40}}%
\TextAt{13}{297.128}{540.119}{L}%
\TextAt{16}{308.414}{540.119}{\Character{106}}%
\TextAt{13}{315.055}{540.119}{K}%
\TextAt{9}{325.873}{540.119}{\Character{41}}%
\TextAt{13}{330.425}{540.119}{\Character{59}}%
\TextAt{13}{359.56}{540.119}{\Character{27}}%
\NativeArrowAt{369.983}{540.119}{11.955}{19.26}{\longmapsto}%
\TextAt{13}{392.565}{540.119}{\Character{27}}%
\TextAt{16}{399.648}{540.119}{\Character{106}}%
\TextAt{21}{402.969}{538.325}{L}%
\TextAt{13}{409.226}{540.119}{\Character{59}}%
\TextAt{8}{79.226}{511.391}{and}%
\TextAt{8}{101.989}{511.391}{we}%
\TextAt{8}{119.216}{511.391}{have}%
\TextAt{8}{146.211}{511.391}{the}%
\TextAt{12}{79.226}{485.456}{Proposition}%
\TextAt{12}{151.937}{485.456}{(9.6).}%
\TextAt{13}{240.505}{469.527}{G}%
\TextAt{21}{249.739}{467.734}{w}%
\TextAt{9}{256.528}{469.527}{\Character{40}}%
\TextAt{13}{261.08}{469.527}{L}%
\TextAt{16}{272.366}{469.527}{\Character{106}}%
\TextAt{13}{279.007}{469.527}{K}%
\TextAt{9}{289.825}{469.527}{\Character{41}}%
\NativeMathAt{297.698}{469.527}{11.955}{\cong}%
\TextAt{13}{310.318}{469.527}{G}%
\TextAt{9}{319.551}{469.527}{\Character{40}}%
\TextAt{13}{324.103}{469.527}{L}%
\TextAt{21}{332.068}{467.734}{w}%
\TextAt{16}{342.178}{469.527}{\Character{106}}%
\TextAt{13}{348.82}{469.527}{K}%
\TextAt{21}{358.779}{467.734}{v}%
\TextAt{9}{363.687}{469.527}{\Character{41}}%
\TextAt{13}{368.239}{469.527}{\Character{59}}%
\TextAt{13}{244.576}{449.492}{I}%
\TextAt{21}{249.739}{447.699}{w}%
\TextAt{9}{256.528}{449.492}{\Character{40}}%
\TextAt{13}{261.08}{449.492}{L}%
\TextAt{16}{272.366}{449.492}{\Character{106}}%
\TextAt{13}{279.007}{449.492}{K}%
\TextAt{9}{289.825}{449.492}{\Character{41}}%
\NativeMathAt{297.698}{449.492}{11.955}{\cong}%
\TextAt{13}{310.317}{449.492}{I}%
\TextAt{9}{316.42}{449.492}{\Character{40}}%
\TextAt{13}{320.973}{449.492}{L}%
\TextAt{21}{328.937}{447.699}{w}%
\TextAt{16}{339.047}{449.492}{\Character{106}}%
\TextAt{13}{345.689}{449.492}{K}%
\TextAt{21}{355.648}{447.699}{v}%
\TextAt{9}{360.556}{449.492}{\Character{41}}%
\TextAt{13}{365.108}{449.492}{\Character{59}}%
\TextAt{13}{240.828}{429.458}{R}%
\TextAt{21}{249.739}{427.664}{w}%
\TextAt{9}{256.528}{429.458}{\Character{40}}%
\TextAt{13}{261.08}{429.458}{L}%
\TextAt{16}{272.366}{429.458}{\Character{106}}%
\TextAt{13}{279.007}{429.458}{K}%
\TextAt{9}{289.825}{429.458}{\Character{41}}%
\NativeMathAt{297.698}{429.457}{11.955}{\cong}%
\TextAt{13}{310.317}{429.458}{R}%
\TextAt{9}{319.326}{429.458}{\Character{40}}%
\TextAt{13}{323.878}{429.458}{L}%
\TextAt{21}{331.843}{427.664}{w}%
\TextAt{16}{341.953}{429.458}{\Character{106}}%
\TextAt{13}{348.595}{429.458}{K}%
\TextAt{21}{358.554}{427.664}{v}%
\TextAt{9}{363.461}{429.458}{\Character{41}}%
\TextAt{13}{368.014}{429.458}{\Character{58}}%
\TextAt{19}{79.226}{390.246}{Proof.}%
\TextAt{8}{116}{390.246}{The}%
\TextAt{8}{140.912}{390.246}{proposition}%
\TextAt{8}{203.304}{390.246}{derives}%
\TextAt{8}{243.244}{390.246}{from}%
\TextAt{8}{271.743}{390.246}{the}%
\TextAt{8}{292.758}{390.246}{fact}%
\TextAt{8}{316.702}{390.246}{that}%
\TextAt{8}{342.929}{390.246}{the}%
\TextAt{8}{363.944}{390.246}{decomposition}%
\TextAt{8}{441.937}{390.246}{group}%
\TextAt{13}{475.971}{390.246}{G}%
\TextAt{21}{485.205}{388.452}{w}%
\TextAt{9}{491.994}{390.246}{\Character{40}}%
\TextAt{13}{496.547}{390.246}{L}%
\TextAt{16}{509.295}{390.246}{\Character{106}}%
\TextAt{13}{517.399}{390.246}{K}%
\TextAt{9}{528.217}{390.246}{\Character{41}}%
\TextAt{8}{79.226}{373.2}{consists}%
\TextAt{8}{121.269}{373.2}{precisely}%
\TextAt{8}{168.058}{373.2}{of}%
\TextAt{8}{180.321}{373.2}{those}%
\TextAt{8}{209.882}{373.2}{automorphisms}%
\TextAt{13}{290.899}{373.2}{\Character{27}}%
\NativeMathAt{301.302}{373.2}{11.955}{\in}%
\TextAt{13}{312.593}{373.2}{G}%
\TextAt{9}{321.827}{373.2}{\Character{40}}%
\TextAt{13}{326.379}{373.2}{L}%
\TextAt{16}{337.664}{373.2}{\Character{106}}%
\TextAt{13}{344.306}{373.2}{K}%
\TextAt{9}{355.124}{373.2}{\Character{41}}%
\TextAt{8}{362.506}{373.2}{which}%
\TextAt{8}{394.925}{373.2}{are}%
\TextAt{19}{413.368}{373.2}{continuous}%
\TextAt{8}{471.025}{373.2}{with}%
\TextAt{8}{496.617}{373.2}{respect}%
\TextAt{8}{79.226}{356.154}{to}%
\TextAt{8}{93.04}{356.154}{the}%
\TextAt{8}{112.699}{356.154}{valuation}%
\TextAt{13}{163.275}{356.154}{w}%
\TextAt{8}{171.965}{356.154}{.}%
\TextAt{8}{180.247}{356.154}{Indeed,}%
\TextAt{8}{221.131}{356.154}{the}%
\TextAt{8}{240.791}{356.154}{continuity}%
\TextAt{8}{295.587}{356.154}{of}%
\TextAt{8}{308.42}{356.154}{the}%
\TextAt{13}{328.08}{356.154}{\Character{27}}%
\NativeMathAt{338.483}{356.154}{11.955}{\in}%
\TextAt{13}{349.774}{356.154}{G}%
\TextAt{21}{359.008}{354.361}{w}%
\TextAt{9}{365.797}{356.154}{\Character{40}}%
\TextAt{13}{370.349}{356.154}{L}%
\TextAt{16}{381.635}{356.154}{\Character{106}}%
\TextAt{13}{388.276}{356.154}{K}%
\TextAt{9}{399.094}{356.154}{\Character{41}}%
\TextAt{8}{407.047}{356.154}{is}%
\TextAt{8}{418.314}{356.154}{clear.}%
\TextAt{8}{450.663}{356.154}{For}%
\TextAt{8}{471.135}{356.154}{an}%
\TextAt{8}{486.897}{356.154}{arbitrary}%
\TextAt{8}{79.226}{339.108}{continuous}%
\TextAt{8}{137.83}{339.108}{automorphism}%
\TextAt{13}{215.3}{339.108}{\Character{27}}%
\TextAt{8}{222.382}{339.108}{,}%
\TextAt{8}{229.531}{339.108}{one}%
\TextAt{8}{250.991}{339.108}{has}%
\TextAt{16}{211.744}{310.38}{\Character{106}}%
\TextAt{13}{215.065}{310.38}{x}%
\TextAt{16}{221.717}{310.38}{\Character{106}}%
\TextAt{21}{225.038}{306.893}{w}%
\TextAt{13}{235.148}{310.38}{\Character{60}}%
\TextAt{9}{247.574}{310.38}{1}%
\NativeArrowAt{268.687}{310.38}{11.955}{19.083}{\Longrightarrow}%
\TextAt{16}{303.037}{310.38}{\Character{106}}%
\TextAt{13}{306.367}{310.38}{\Character{27}}%
\TextAt{13}{313.433}{310.38}{x}%
\TextAt{16}{320.101}{310.38}{\Character{106}}%
\TextAt{21}{323.422}{306.893}{w}%
\TextAt{9}{333.532}{310.38}{\Character{61}}%
\TextAt{16}{345.958}{310.38}{\Character{106}}%
\TextAt{13}{349.279}{310.38}{x}%
\TextAt{16}{355.931}{310.38}{\Character{106}}%
\TextAt{21}{359.252}{306.893}{w}%
\TextAt{3}{365.543}{306.893}{\Character{14}}%
\TextAt{21}{369.777}{306.893}{\Character{27}}%
\TextAt{13}{378.722}{310.38}{\Character{60}}%
\TextAt{9}{391.147}{310.38}{1}%
\TextAt{13}{397}{310.38}{\Character{59}}%
\TextAt{8}{79.226}{281.652}{because}%
\TextAt{16}{122.828}{281.652}{\Character{106}}%
\TextAt{13}{126.148}{281.652}{x}%
\TextAt{16}{132.801}{281.652}{\Character{106}}%
\TextAt{21}{136.121}{278.165}{w}%
\TextAt{13}{146.743}{281.652}{\Character{60}}%
\TextAt{9}{159.679}{281.652}{1}%
\TextAt{8}{169.729}{281.652}{means}%
\TextAt{8}{205.858}{281.652}{that}%
\TextAt{13}{231.527}{281.652}{x}%
\TextAt{21}{238.179}{285.99}{n}%
\TextAt{8}{248.013}{281.652}{and}%
\TextAt{8}{271.075}{281.652}{hence}%
\TextAt{8}{303.881}{281.652}{also}%
\TextAt{13}{327.661}{281.652}{\Character{27}}%
\TextAt{13}{334.726}{281.652}{x}%
\TextAt{21}{341.395}{285.99}{n}%
\TextAt{8}{351.229}{281.652}{is}%
\TextAt{8}{363.292}{281.652}{a}%
\TextAt{13}{373.348}{281.652}{w}%
\TextAt{8}{382.037}{281.652}{-nullsequence,}%
\TextAt{8}{457.251}{281.652}{i.e.,}%
\TextAt{16}{479.731}{281.652}{\Character{106}}%
\TextAt{13}{483.052}{281.652}{\Character{27}}%
\TextAt{13}{490.118}{281.652}{x}%
\TextAt{16}{496.786}{281.652}{\Character{106}}%
\TextAt{21}{500.107}{278.165}{w}%
\TextAt{13}{510.728}{281.652}{\Character{60}}%
\TextAt{9}{523.665}{281.652}{1}%
\TextAt{8}{529.518}{281.652}{.}%
\TextAt{8}{79.226}{264.606}{By}%
\TextAt{8}{96.77}{264.606}{§3,}%
\TextAt{8}{114.8}{264.606}{p.}%
\TextAt{8}{127.633}{264.606}{117,}%
\TextAt{8}{151.701}{264.606}{this}%
\TextAt{8}{173.704}{264.606}{implies}%
\TextAt{8}{212.612}{264.606}{that}%
\TextAt{13}{237.161}{264.606}{w}%
\TextAt{8}{248.928}{264.606}{and}%
\TextAt{13}{270.872}{264.606}{w}%
\TextAt{16}{280.542}{264.606}{\Character{14}}%
\TextAt{13}{287.5}{264.606}{\Character{27}}%
\TextAt{8}{297.661}{264.606}{are}%
\TextAt{8}{316.352}{264.606}{equivalent,}%
\TextAt{8}{374.564}{264.606}{and}%
\TextAt{8}{396.507}{264.606}{hence}%
\TextAt{8}{428.193}{264.606}{in}%
\TextAt{8}{441.027}{264.606}{fact}%
\TextAt{8}{463.293}{264.606}{equal}%
\TextAt{8}{493.365}{264.606}{because}%
\TextAt{13}{79.226}{247.56}{w}%
\TextAt{16}{87.916}{247.56}{\Character{106}}%
\TextAt{21}{91.237}{245.766}{K}%
\TextAt{9}{102.76}{247.56}{\Character{61}}%
\TextAt{13}{115.186}{247.56}{w}%
\TextAt{16}{126.532}{247.56}{\Character{14}}%
\TextAt{13}{135.166}{247.56}{\Character{27}}%
\TextAt{16}{142.248}{247.56}{\Character{106}}%
\TextAt{21}{145.569}{245.766}{K}%
\TextAt{8}{153.772}{247.56}{,}%
\TextAt{8}{160.921}{247.56}{so}%
\TextAt{8}{175.291}{247.56}{that}%
\TextAt{13}{200.66}{247.56}{\Character{27}}%
\NativeMathAt{211.063}{247.56}{11.955}{\in}%
\TextAt{13}{222.354}{247.56}{G}%
\TextAt{21}{231.588}{245.766}{w}%
\TextAt{9}{238.377}{247.56}{\Character{40}}%
\TextAt{13}{242.93}{247.56}{L}%
\TextAt{16}{254.215}{247.56}{\Character{106}}%
\TextAt{13}{260.857}{247.56}{K}%
\TextAt{9}{271.674}{247.56}{\Character{41}}%
\TextAt{8}{276.227}{247.56}{.}%
\TextAt{8}{97.159}{227.525}{Since}%
\TextAt{13}{128.316}{227.525}{L}%
\TextAt{8}{140.777}{227.525}{is}%
\TextAt{8}{153.141}{227.525}{dense}%
\TextAt{8}{185.661}{227.525}{in}%
\TextAt{13}{199.913}{227.525}{L}%
\TextAt{21}{207.877}{225.732}{w}%
\TextAt{8}{214.667}{227.525}{,}%
\TextAt{8}{222.565}{227.525}{every}%
\TextAt{13}{254.057}{227.525}{\Character{27}}%
\NativeMathAt{265.482}{227.525}{11.955}{\in}%
\TextAt{13}{277.795}{227.525}{G}%
\TextAt{21}{287.028}{225.732}{w}%
\TextAt{9}{293.818}{227.525}{\Character{40}}%
\TextAt{13}{298.37}{227.525}{L}%
\TextAt{16}{310.677}{227.525}{\Character{106}}%
\TextAt{13}{318.341}{227.525}{K}%
\TextAt{9}{329.158}{227.525}{\Character{41}}%
\TextAt{8}{338.208}{227.525}{extends}%
\TextAt{8}{381.463}{227.525}{uniquely}%
\TextAt{8}{429.537}{227.525}{to}%
\TextAt{8}{444.447}{227.525}{a}%
\TextAt{8}{454.802}{227.525}{continuous}%
\TextAt{13}{514.005}{227.525}{K}%
\TextAt{21}{523.965}{225.732}{v}%
\TextAt{8}{528.872}{227.525}{-}%
\TextAt{8}{79.226}{210.479}{automorphism}%
\TextAt{9}{158.402}{210.479}{\Character{94}}%
\TextAt{13}{157.787}{210.479}{\Character{27}}%
\TextAt{8}{169.859}{210.479}{of}%
\TextAt{13}{184.28}{210.479}{L}%
\TextAt{21}{192.245}{208.686}{w}%
\TextAt{8}{204.023}{210.479}{and}%
\TextAt{8}{227.877}{210.479}{it}%
\TextAt{8}{240.673}{210.479}{is}%
\TextAt{8}{253.529}{210.479}{clear}%
\TextAt{8}{282.583}{210.479}{that}%
\TextAt{9}{309.659}{210.479}{\Character{94}}%
\TextAt{13}{309.044}{210.479}{\Character{27}}%
\NativeMathAt{321.307}{210.479}{11.955}{\in}%
\TextAt{13}{334.459}{210.479}{I}%
\TextAt{9}{340.561}{210.479}{\Character{40}}%
\TextAt{13}{345.114}{210.479}{L}%
\TextAt{21}{353.078}{208.686}{w}%
\TextAt{16}{365.049}{210.479}{\Character{106}}%
\TextAt{13}{373.551}{210.479}{K}%
\TextAt{21}{383.51}{208.686}{v}%
\TextAt{9}{388.417}{210.479}{\Character{41}}%
\TextAt{8}{392.97}{210.479}{,}%
\TextAt{8}{401.483}{210.479}{resp.}%
\TextAt{9}{431.213}{210.479}{\Character{94}}%
\TextAt{13}{430.598}{210.479}{\Character{27}}%
\NativeMathAt{442.861}{210.479}{11.955}{\in}%
\TextAt{13}{456.013}{210.479}{R}%
\TextAt{9}{465.021}{210.479}{\Character{40}}%
\TextAt{13}{469.573}{210.479}{L}%
\TextAt{21}{477.538}{208.686}{w}%
\TextAt{16}{489.508}{210.479}{\Character{106}}%
\TextAt{13}{498.01}{210.479}{K}%
\TextAt{21}{507.969}{208.686}{v}%
\TextAt{9}{512.877}{210.479}{\Character{41}}%
\TextAt{8}{517.429}{210.479}{,}%
\TextAt{8}{525.943}{210.479}{if}%
\TextAt{13}{79.226}{193.433}{\Character{27}}%
\NativeMathAt{91.53}{193.433}{11.955}{\in}%
\TextAt{13}{104.723}{193.433}{I}%
\TextAt{21}{109.885}{191.64}{w}%
\TextAt{9}{116.675}{193.433}{\Character{40}}%
\TextAt{13}{121.227}{193.433}{L}%
\TextAt{16}{134.414}{193.433}{\Character{106}}%
\TextAt{13}{142.956}{193.433}{K}%
\TextAt{9}{153.774}{193.433}{\Character{41}}%
\TextAt{8}{158.326}{193.433}{,}%
\TextAt{8}{166.87}{193.433}{resp.}%
\TextAt{13}{196.009}{193.433}{\Character{27}}%
\NativeMathAt{208.313}{193.433}{11.955}{\in}%
\TextAt{13}{221.505}{193.433}{R}%
\TextAt{21}{230.415}{191.64}{w}%
\TextAt{9}{237.205}{193.433}{\Character{40}}%
\TextAt{13}{241.757}{193.433}{L}%
\TextAt{16}{254.943}{193.433}{\Character{106}}%
\TextAt{13}{263.486}{193.433}{K}%
\TextAt{9}{274.304}{193.433}{\Character{41}}%
\TextAt{8}{278.856}{193.433}{.}%
\TextAt{8}{290.651}{193.433}{This}%
\TextAt{8}{318.486}{193.433}{proves}%
\TextAt{8}{355.766}{193.433}{the}%
\TextAt{8}{377.038}{193.433}{bijectivity}%
\TextAt{8}{433.435}{193.433}{of}%
\TextAt{8}{447.88}{193.433}{the}%
\TextAt{8}{469.152}{193.433}{mappings}%
\TextAt{8}{523.014}{193.433}{in}%
\TextAt{8}{79.226}{176.387}{question}%
\TextAt{8}{125.792}{176.387}{in}%
\TextAt{8}{139.444}{176.387}{all}%
\TextAt{8}{155.703}{176.387}{three}%
\TextAt{8}{185.615}{176.387}{cases.}%
\Vector{q 0 G 0 g 0.399 w 452.334 -597.613 m 452.334 -589.543 l S Q}%
\Vector{q 0 G 0 g 0.399 w 452.534 -589.742 m 459.707 -589.742 l S Q}%
\Vector{q 0 G 0 g 0.399 w 452.534 -597.414 m 459.707 -597.414 l S Q}%
\Vector{q 0 G 0 g 0.399 w 459.906 -597.613 m 459.906 -589.543 l S Q}%
\TextAt{8}{97.159}{147.464}{The}%
\TextAt{8}{120.776}{147.464}{above}%
\TextAt{8}{153.515}{147.464}{proposition}%
\TextAt{8}{214.611}{147.464}{reduces}%
\TextAt{8}{255.851}{147.464}{the}%
\TextAt{8}{275.57}{147.464}{problems}%
\TextAt{8}{325.274}{147.464}{concerning}%
\TextAt{8}{383.369}{147.464}{a}%
\TextAt{8}{392.688}{147.464}{single}%
\TextAt{8}{424.829}{147.464}{valuation}%
\TextAt{8}{475.465}{147.464}{of}%
\TextAt{13}{488.358}{147.464}{K}%
\TextAt{8}{502.637}{147.464}{to}%
\TextAt{8}{516.51}{147.464}{the}%
\TextAt{8}{79.226}{130.418}{local}%
\TextAt{8}{106.53}{130.418}{situation.}%
\TextAt{8}{159.818}{130.418}{We}%
\TextAt{8}{179.626}{130.418}{identify}%
\TextAt{8}{221.886}{130.418}{the}%
\TextAt{8}{241.706}{130.418}{decomposition}%
\TextAt{8}{318.504}{130.418}{group}%
\TextAt{13}{351.343}{130.418}{G}%
\TextAt{21}{360.577}{128.625}{w}%
\TextAt{9}{367.366}{130.418}{\Character{40}}%
\TextAt{13}{371.919}{130.418}{L}%
\TextAt{16}{383.204}{130.418}{\Character{106}}%
\TextAt{13}{389.846}{130.418}{K}%
\TextAt{9}{400.664}{130.418}{\Character{41}}%
\TextAt{8}{408.777}{130.418}{with}%
\TextAt{8}{435.101}{130.418}{the}%
\TextAt{8}{454.921}{130.418}{Galois}%
\TextAt{8}{490.498}{130.418}{group}%
\TextAt{8}{523.337}{130.418}{of}%
\TextAt{13}{79.226}{113.372}{L}%
\TextAt{21}{87.191}{111.579}{w}%
\TextAt{16}{97.301}{113.372}{\Character{106}}%
\TextAt{13}{103.943}{113.372}{K}%
\TextAt{21}{113.902}{111.579}{v}%
\TextAt{8}{122.707}{113.372}{and}%
\TextAt{8}{145.469}{113.372}{write}%
\TextAt{13}{240.602}{96.326}{G}%
\TextAt{21}{249.836}{94.533}{w}%
\TextAt{9}{256.625}{96.326}{\Character{40}}%
\TextAt{13}{261.177}{96.326}{L}%
\TextAt{16}{272.463}{96.326}{\Character{106}}%
\TextAt{13}{279.104}{96.326}{K}%
\TextAt{9}{289.922}{96.326}{\Character{41}}%
\TextAt{9}{297.777}{96.326}{\Character{61}}%
\TextAt{13}{310.221}{96.326}{G}%
\TextAt{9}{319.454}{96.326}{\Character{40}}%
\TextAt{13}{324.007}{96.326}{L}%
\TextAt{21}{331.971}{94.533}{w}%
\TextAt{16}{342.081}{96.326}{\Character{106}}%
\TextAt{13}{348.723}{96.326}{K}%
\TextAt{21}{358.682}{94.533}{v}%
\TextAt{9}{363.59}{96.326}{\Character{41}}%
\TextAt{13}{368.142}{96.326}{\Character{59}}%
\TextAt{8}{79.226}{72.921}{and}%
\TextAt{8}{101.989}{72.921}{similarly}%
\TextAt{13}{149.857}{72.921}{I}%
\TextAt{21}{155.02}{71.128}{w}%
\TextAt{9}{161.81}{72.921}{\Character{40}}%
\TextAt{13}{166.362}{72.921}{L}%
\TextAt{16}{177.647}{72.921}{\Character{106}}%
\TextAt{13}{184.289}{72.921}{K}%
\TextAt{9}{195.107}{72.921}{\Character{41}}%
\TextAt{9}{202.962}{72.921}{\Character{61}}%
\TextAt{13}{215.405}{72.921}{I}%
\TextAt{9}{221.508}{72.921}{\Character{40}}%
\TextAt{13}{226.06}{72.921}{L}%
\TextAt{21}{234.025}{71.128}{w}%
\TextAt{16}{244.135}{72.921}{\Character{106}}%
\TextAt{13}{250.777}{72.921}{K}%
\TextAt{21}{260.736}{71.128}{v}%
\TextAt{9}{265.644}{72.921}{\Character{41}}%
\TextAt{8}{274.093}{72.921}{and}%
\TextAt{13}{296.856}{72.921}{R}%
\TextAt{21}{305.766}{71.128}{w}%
\TextAt{9}{312.556}{72.921}{\Character{40}}%
\TextAt{13}{317.108}{72.921}{L}%
\TextAt{16}{328.393}{72.921}{\Character{106}}%
\TextAt{13}{335.035}{72.921}{K}%
\TextAt{9}{345.853}{72.921}{\Character{41}}%
\TextAt{9}{353.708}{72.921}{\Character{61}}%
\TextAt{13}{366.151}{72.921}{R}%
\TextAt{9}{375.16}{72.921}{\Character{40}}%
\TextAt{13}{379.712}{72.921}{L}%
\TextAt{21}{387.677}{71.128}{w}%
\TextAt{16}{397.787}{72.921}{\Character{106}}%
\TextAt{13}{404.429}{72.921}{K}%
\TextAt{21}{414.388}{71.128}{v}%
\TextAt{9}{419.295}{72.921}{\Character{41}}%
\TextAt{8}{423.848}{72.921}{.}%
\TextAt{8}{297.211}{43.033}{156}%
\Destination{78.23}{809.815}{nav.page.156}%
\Destination{79.36}{776.155}{nav.print.170}%
\Destination{79.36}{499.455}{nav.number.2.9.6}%
\Hotlink{389.845}{634.855}{6.858}{11.047}{nav.section.2.8}%
\Hotlink{418.095}{634.855}{18.574}{11.047}{nav.print.160}%
\Hotlink{101.949}{262.455}{6.858}{11.047}{nav.section.2.3}%
\Hotlink{127.133}{262.455}{18.574}{11.047}{nav.print.117}%
\Hotlink{304.248}{758.955}{25.078}{11.047}{nav.number.2.9.4}%
}

% Page 169 (body page 157)
\Page{%
\TextAt{2}{79.226}{796.943}{II.}%
\TextAt{2}{94.073}{796.943}{The}%
\TextAt{2}{116.492}{796.943}{Theory}%
\TextAt{2}{154.401}{796.943}{of}%
\TextAt{2}{166.826}{796.943}{Valuations}%
\Vector{q 0 G 0 g 0.299 w 7.226 18.729 m 460.769 18.729 l S Q}%
\TextAt{8}{97.159}{761.124}{We}%
\TextAt{8}{117.58}{761.124}{now}%
\TextAt{8}{142.244}{761.124}{explain}%
\TextAt{8}{183.168}{761.124}{the}%
\TextAt{8}{203.601}{761.124}{concrete}%
\TextAt{8}{250.048}{761.124}{meaning}%
\TextAt{8}{297.153}{761.124}{of}%
\TextAt{8}{310.759}{761.124}{the}%
\TextAt{8}{331.192}{761.124}{subgroups}%
\TextAt{13}{386.88}{761.124}{G}%
\TextAt{21}{396.114}{759.33}{w}%
\TextAt{13}{402.903}{761.124}{\Character{59}}%
\TextAt{13}{408.127}{761.124}{I}%
\TextAt{21}{413.31}{759.33}{w}%
\TextAt{13}{420.1}{761.124}{\Character{59}}%
\TextAt{13}{425.324}{761.124}{R}%
\TextAt{21}{434.254}{759.33}{w}%
\TextAt{8}{445.217}{761.124}{of}%
\TextAt{13}{458.823}{761.124}{G}%
\TextAt{9}{471.849}{761.124}{\Character{61}}%
\TextAt{13}{484.745}{761.124}{G}%
\TextAt{9}{493.979}{761.124}{\Character{40}}%
\TextAt{13}{498.531}{761.124}{L}%
\TextAt{16}{510.287}{761.124}{\Character{106}}%
\TextAt{13}{517.399}{761.124}{K}%
\TextAt{9}{528.217}{761.124}{\Character{41}}%
\TextAt{8}{79.226}{744.078}{for}%
\TextAt{8}{97.111}{744.078}{the}%
\TextAt{8}{117.267}{744.078}{field}%
\TextAt{8}{142.624}{744.078}{extension}%
\TextAt{13}{194.39}{744.078}{L}%
\TextAt{16}{205.676}{744.078}{\Character{106}}%
\TextAt{13}{212.317}{744.078}{K}%
\TextAt{8}{223.135}{744.078}{.}%
\TextAt{8}{97.159}{724.043}{The}%
\TextAt{12}{121.422}{724.043}{decomposition}%
\TextAt{12}{210.478}{724.043}{group}%
\TextAt{13}{248.489}{724.043}{G}%
\TextAt{21}{257.723}{722.25}{w}%
\TextAt{8}{268.619}{724.043}{consists—as}%
\TextAt{8}{334.116}{724.043}{was}%
\TextAt{8}{356.825}{724.043}{shown}%
\TextAt{8}{392.541}{724.043}{in}%
\TextAt{8}{406.403}{724.043}{the}%
\TextAt{8}{426.769}{724.043}{proof}%
\TextAt{8}{457.547}{724.043}{of}%
\TextAt{8}{471.087}{724.043}{(9.6)—of}%
\TextAt{8}{520.408}{724.043}{all}%
\TextAt{8}{79.226}{706.997}{automorphisms}%
\TextAt{13}{161.936}{706.997}{\Character{27}}%
\NativeMathAt{173.406}{706.997}{11.955}{\in}%
\TextAt{13}{185.764}{706.997}{G}%
\TextAt{8}{199.521}{706.997}{that}%
\TextAt{8}{225.516}{706.997}{are}%
\TextAt{19}{245.653}{706.997}{continuous}%
\TextAt{8}{305.002}{706.997}{with}%
\TextAt{8}{332.289}{706.997}{respect}%
\TextAt{8}{372.964}{706.997}{to}%
\TextAt{8}{387.901}{706.997}{the}%
\TextAt{8}{408.683}{706.997}{valuation}%
\TextAt{13}{460.381}{706.997}{w}%
\TextAt{8}{469.071}{706.997}{.}%
\TextAt{8}{479.397}{706.997}{It}%
\TextAt{8}{492.696}{706.997}{controls}%
\TextAt{8}{79.226}{689.951}{the}%
\TextAt{8}{100.153}{689.951}{extension}%
\TextAt{8}{152.69}{689.951}{of}%
\TextAt{13}{166.79}{689.951}{v}%
\TextAt{8}{177.547}{689.951}{to}%
\TextAt{13}{192.628}{689.951}{L}%
\TextAt{8}{205.26}{689.951}{in}%
\TextAt{8}{219.683}{689.951}{a}%
\TextAt{8}{230.21}{689.951}{group-theoretic}%
\TextAt{8}{312.933}{689.951}{manner.}%
\TextAt{8}{362.069}{689.951}{Denoting}%
\TextAt{8}{413.41}{689.951}{by}%
\TextAt{13}{430.44}{689.951}{G}%
\TextAt{21}{439.673}{688.158}{w}%
\TextAt{16}{446.463}{689.951}{\Character{110}}%
\TextAt{13}{452.44}{689.951}{G}%
\TextAt{8}{466.342}{689.951}{the}%
\TextAt{8}{487.269}{689.951}{set}%
\TextAt{8}{506.307}{689.951}{of}%
\TextAt{8}{520.408}{689.951}{all}%
\TextAt{8}{79.226}{672.905}{right}%
\TextAt{8}{107.257}{672.905}{cosets}%
\TextAt{13}{140.932}{672.905}{G}%
\TextAt{21}{150.165}{671.112}{w}%
\TextAt{13}{156.955}{672.905}{\Character{27}}%
\TextAt{8}{164.037}{672.905}{,}%
\TextAt{8}{170.973}{672.905}{by}%
\TextAt{13}{186.966}{672.905}{W}%
\TextAt{21}{198.017}{671.112}{v}%
\TextAt{8}{206.555}{672.905}{the}%
\TextAt{8}{226.445}{672.905}{set}%
\TextAt{8}{244.446}{672.905}{of}%
\TextAt{8}{257.51}{672.905}{extensions}%
\TextAt{8}{313.624}{672.905}{of}%
\TextAt{13}{326.687}{672.905}{v}%
\TextAt{8}{336.406}{672.905}{to}%
\TextAt{13}{350.45}{672.905}{L}%
\TextAt{8}{362.045}{672.905}{and}%
\TextAt{8}{384.541}{672.905}{choosing}%
\TextAt{8}{431.821}{672.905}{a}%
\TextAt{8}{441.309}{672.905}{fixed}%
\TextAt{8}{469.329}{672.905}{extension}%
\TextAt{13}{520.828}{672.905}{w}%
\TextAt{8}{529.518}{672.905}{,}%
\TextAt{8}{79.226}{655.859}{we}%
\TextAt{8}{96.454}{655.859}{obtain}%
\TextAt{8}{132.881}{655.859}{a}%
\TextAt{8}{142.636}{655.859}{bijection}%
\TextAt{13}{224.226}{638.813}{G}%
\TextAt{21}{233.459}{637.02}{w}%
\TextAt{16}{240.249}{638.813}{\Character{110}}%
\TextAt{13}{246.226}{638.813}{G}%
\NativeMathAt{261.368}{645.59}{7.97}{\sim}%
\NativeArrowAt{258.781}{638.813}{11.955}{13.175}{\longrightarrow}%
\TextAt{13}{275.275}{638.813}{W}%
\TextAt{21}{286.326}{637.02}{v}%
\TextAt{13}{291.234}{638.813}{\Character{59}}%
\TextAt{13}{320.369}{638.813}{G}%
\TextAt{21}{329.622}{637.02}{w}%
\TextAt{13}{336.411}{638.813}{\Character{27}}%
\NativeArrowAt{346.815}{638.813}{11.955}{19.26}{\longmapsto}%
\TextAt{13}{369.397}{638.813}{w}%
\TextAt{13}{378.065}{638.813}{\Character{27}}%
\TextAt{13}{384.508}{638.813}{\Character{58}}%
\TextAt{8}{79.226}{615.85}{In}%
\TextAt{8}{94.38}{615.85}{particular,}%
\TextAt{8}{152.287}{615.85}{the}%
\TextAt{8}{172.976}{615.85}{number}%
\TextAt{8}{216.105}{615.85}{\char35\relax }%
\TextAt{13}{225.86}{615.85}{W}%
\TextAt{21}{236.912}{614.057}{v}%
\TextAt{8}{246.25}{615.85}{of}%
\TextAt{8}{260.112}{615.85}{extensions}%
\TextAt{8}{317.025}{615.85}{equals}%
\TextAt{8}{353.065}{615.85}{the}%
\TextAt{8}{373.754}{615.85}{index}%
\TextAt{9}{405.824}{615.85}{\Character{40}}%
\TextAt{13}{410.377}{615.85}{G}%
\TextAt{9}{423.839}{615.85}{\Character{58}}%
\TextAt{13}{431.319}{615.85}{G}%
\TextAt{21}{440.552}{614.057}{w}%
\TextAt{9}{447.342}{615.85}{\Character{41}}%
\TextAt{8}{451.894}{615.85}{.}%
\TextAt{8}{461.94}{615.85}{As}%
\TextAt{8}{479.76}{615.85}{mentioned}%
\TextAt{8}{79.226}{598.804}{already}%
\TextAt{8}{121.404}{598.804}{in}%
\TextAt{8}{135.93}{598.804}{chap.}%
\TextAt{8}{167.694}{598.804}{I,}%
\TextAt{8}{179.937}{598.804}{§9—and}%
\TextAt{8}{226.813}{598.804}{left}%
\TextAt{8}{248.165}{598.804}{for}%
\TextAt{8}{266.922}{598.804}{the}%
\TextAt{8}{287.951}{598.804}{reader}%
\TextAt{8}{324.594}{598.804}{to}%
\TextAt{8}{339.777}{598.804}{verify—the}%
\TextAt{8}{401.131}{598.804}{decomposition}%
\TextAt{8}{479.139}{598.804}{group}%
\TextAt{8}{513.187}{598.804}{also}%
\TextAt{8}{79.226}{581.758}{describes}%
\TextAt{8}{129.553}{581.758}{the}%
\TextAt{8}{150.171}{581.758}{way}%
\TextAt{8}{174.375}{581.758}{a}%
\TextAt{8}{184.592}{581.758}{valuation}%
\TextAt{13}{236.126}{581.758}{v}%
\TextAt{8}{246.573}{581.758}{extends}%
\TextAt{8}{289.691}{581.758}{to}%
\TextAt{8}{304.463}{581.758}{an}%
\TextAt{8}{321.184}{581.758}{arbitrary}%
\TextAt{8}{371.414}{581.758}{separable}%
\TextAt{8}{423.319}{581.758}{extension}%
\TextAt{13}{475.547}{581.758}{L}%
\TextAt{16}{487.619}{581.758}{\Character{106}}%
\TextAt{13}{495.047}{581.758}{K}%
\TextAt{8}{505.864}{581.758}{.}%
\TextAt{8}{515.698}{581.758}{For}%
\TextAt{8}{79.226}{564.712}{this,}%
\TextAt{8}{105.92}{564.712}{we}%
\TextAt{8}{123.644}{564.712}{embed}%
\TextAt{13}{161.201}{564.712}{L}%
\TextAt{16}{173.331}{564.712}{\Character{106}}%
\TextAt{13}{180.818}{564.712}{K}%
\TextAt{8}{196.029}{564.712}{into}%
\TextAt{8}{220.268}{564.712}{a}%
\TextAt{8}{230.519}{564.712}{Galois}%
\TextAt{8}{266.929}{564.712}{extension}%
\TextAt{13}{319.19}{564.712}{N}%
\TextAt{16}{333.979}{564.712}{\Character{106}}%
\TextAt{13}{341.466}{564.712}{K}%
\TextAt{8}{352.284}{564.712}{,}%
\TextAt{8}{360.053}{564.712}{choose}%
\TextAt{8}{397.682}{564.712}{an}%
\TextAt{8}{414.437}{564.712}{extension}%
\TextAt{13}{466.698}{564.712}{w}%
\TextAt{8}{479.781}{564.712}{of}%
\TextAt{13}{493.607}{564.712}{v}%
\TextAt{8}{504.089}{564.712}{to}%
\TextAt{13}{518.895}{564.712}{N}%
\TextAt{8}{529.518}{564.712}{,}%
\TextAt{8}{79.226}{547.666}{and}%
\TextAt{8}{101.989}{547.666}{put}%
\TextAt{13}{123.448}{547.666}{G}%
\TextAt{9}{136.003}{547.666}{\Character{61}}%
\TextAt{13}{148.428}{547.666}{G}%
\TextAt{9}{157.662}{547.666}{\Character{40}}%
\TextAt{13}{162.214}{547.666}{N}%
\TextAt{16}{176.158}{547.666}{\Character{106}}%
\TextAt{13}{182.799}{547.666}{K}%
\TextAt{9}{193.617}{547.666}{\Character{41}}%
\TextAt{8}{198.169}{547.666}{,}%
\TextAt{13}{205.319}{547.666}{H}%
\TextAt{9}{219.262}{547.666}{\Character{61}}%
\TextAt{13}{231.687}{547.666}{G}%
\TextAt{9}{240.921}{547.666}{\Character{40}}%
\TextAt{13}{245.473}{547.666}{N}%
\TextAt{16}{259.417}{547.666}{\Character{106}}%
\TextAt{13}{266.059}{547.666}{L}%
\TextAt{9}{274.023}{547.666}{\Character{41}}%
\TextAt{8}{278.575}{547.666}{,}%
\TextAt{13}{285.725}{547.666}{G}%
\TextAt{21}{294.958}{545.873}{w}%
\TextAt{9}{305.069}{547.666}{\Character{61}}%
\TextAt{13}{317.494}{547.666}{G}%
\TextAt{21}{326.728}{545.873}{w}%
\TextAt{9}{333.517}{547.666}{\Character{40}}%
\TextAt{13}{338.069}{547.666}{N}%
\TextAt{16}{352.013}{547.666}{\Character{106}}%
\TextAt{13}{358.655}{547.666}{K}%
\TextAt{9}{369.472}{547.666}{\Character{41}}%
\TextAt{8}{374.025}{547.666}{,}%
\TextAt{8}{381.174}{547.666}{to}%
\TextAt{8}{395.484}{547.666}{get}%
\TextAt{8}{414.995}{547.666}{a}%
\TextAt{8}{424.75}{547.666}{bijection}%
\TextAt{13}{199.917}{519.968}{G}%
\TextAt{21}{209.151}{518.175}{w}%
\TextAt{16}{215.94}{519.968}{\Character{110}}%
\TextAt{13}{221.918}{519.968}{G}%
\TextAt{9}{231.151}{519.968}{\Character{47}}%
\TextAt{13}{237.004}{519.968}{H}%
\NativeMathAt{253.535}{526.745}{7.97}{\sim}%
\NativeArrowAt{250.948}{519.968}{11.955}{13.175}{\longrightarrow}%
\TextAt{13}{267.442}{519.968}{W}%
\TextAt{21}{278.493}{518.175}{v}%
\TextAt{13}{283.401}{519.968}{\Character{59}}%
\TextAt{13}{312.536}{519.968}{G}%
\TextAt{21}{321.789}{518.175}{w}%
\TextAt{13}{328.578}{519.968}{\Character{27}}%
\TextAt{13}{335.644}{519.968}{H}%
\NativeArrowAt{349.604}{519.968}{11.955}{19.26}{\longmapsto}%
\TextAt{13}{372.186}{519.968}{w}%
\TextAt{16}{383.532}{519.968}{\Character{14}}%
\TextAt{13}{392.167}{519.968}{\Character{27}}%
\TextAt{16}{399.249}{519.968}{\Character{106}}%
\TextAt{21}{402.57}{518.175}{L}%
\TextAt{13}{408.827}{519.968}{\Character{58}}%
\TextAt{12}{79.226}{483.676}{Definition}%
\TextAt{12}{142.015}{483.676}{(9.7).}%
\TextAt{8}{177.33}{483.676}{The}%
\TextAt{8}{201.384}{483.676}{fixed}%
\TextAt{8}{229.67}{483.676}{field}%
\TextAt{8}{255.027}{483.676}{of}%
\TextAt{13}{268.357}{483.676}{G}%
\TextAt{21}{277.591}{481.883}{w}%
\TextAt{8}{284.38}{483.676}{,}%
\TextAt{13}{173.38}{455.978}{Z}%
\TextAt{21}{181.393}{454.185}{w}%
\TextAt{9}{191.503}{455.978}{\Character{61}}%
\TextAt{13}{203.928}{455.978}{Z}%
\TextAt{21}{211.941}{454.185}{w}%
\TextAt{9}{218.73}{455.978}{\Character{40}}%
\TextAt{13}{223.283}{455.978}{L}%
\TextAt{16}{234.568}{455.978}{\Character{106}}%
\TextAt{13}{241.21}{455.978}{K}%
\TextAt{9}{252.027}{455.978}{\Character{41}}%
\TextAt{9}{259.882}{455.978}{\Character{61}}%
\TextAt{16}{272.326}{455.978}{\Character{102}}%
\TextAt{13}{278.304}{455.978}{x}%
\NativeMathAt{288.276}{455.978}{11.955}{\in}%
\TextAt{13}{299.567}{455.978}{L}%
\TextAt{16}{310.853}{455.978}{\Character{106}}%
\TextAt{13}{317.494}{455.978}{\Character{27}}%
\TextAt{13}{324.56}{455.978}{x}%
\TextAt{9}{334.55}{455.978}{\Character{61}}%
\TextAt{13}{346.975}{455.978}{x}%
\TextAt{8}{357.525}{455.978}{for}%
\TextAt{8}{375.41}{455.978}{all}%
\TextAt{13}{391.669}{455.978}{\Character{27}}%
\NativeMathAt{402.072}{455.978}{11.955}{\in}%
\TextAt{13}{413.363}{455.978}{G}%
\TextAt{21}{422.596}{454.185}{w}%
\TextAt{16}{429.386}{455.978}{\Character{103}}%
\TextAt{13}{435.364}{455.978}{\Character{59}}%
\TextAt{8}{79.226}{428.28}{is}%
\TextAt{8}{90.99}{428.28}{called}%
\TextAt{8}{124.154}{428.28}{the}%
\TextAt{12}{144.31}{428.28}{decomposition}%
\TextAt{12}{233.125}{428.28}{field}%
\TextAt{8}{261.841}{428.28}{of}%
\TextAt{13}{275.171}{428.28}{w}%
\TextAt{8}{287.758}{428.28}{over}%
\TextAt{13}{312.804}{428.28}{K}%
\TextAt{8}{323.622}{428.28}{.}%
\TextAt{8}{97.159}{402.64}{The}%
\TextAt{8}{120.436}{402.64}{role}%
\TextAt{8}{142.422}{402.64}{of}%
\TextAt{8}{154.976}{402.64}{the}%
\TextAt{8}{174.356}{402.64}{decomposition}%
\TextAt{8}{250.714}{402.64}{field}%
\TextAt{8}{275.294}{402.64}{in}%
\TextAt{8}{288.17}{402.64}{the}%
\TextAt{8}{307.55}{402.64}{extension}%
\TextAt{13}{358.54}{402.64}{L}%
\TextAt{16}{369.825}{402.64}{\Character{106}}%
\TextAt{13}{376.467}{402.64}{K}%
\TextAt{8}{390.405}{402.64}{is}%
\TextAt{8}{401.393}{402.64}{described}%
\TextAt{8}{452.37}{402.64}{by}%
\TextAt{8}{467.852}{402.64}{the}%
\TextAt{8}{487.232}{402.64}{following}%
\TextAt{8}{79.226}{385.594}{proposition.}%
\TextAt{12}{79.226}{359.954}{Proposition}%
\TextAt{12}{151.937}{359.954}{(9.8).}%
\TextAt{8}{89.579}{339.919}{(i)}%
\TextAt{8}{107.918}{339.919}{The}%
\TextAt{8}{131.972}{339.919}{restriction}%
\TextAt{13}{187.97}{339.919}{w}%
\TextAt{21}{196.338}{338.126}{Z}%
\TextAt{8}{207.035}{339.919}{of}%
\TextAt{13}{220.365}{339.919}{w}%
\TextAt{8}{232.952}{339.919}{to}%
\TextAt{13}{247.262}{339.919}{Z}%
\TextAt{21}{255.275}{338.126}{w}%
\TextAt{8}{265.961}{339.919}{extends}%
\TextAt{8}{308.617}{339.919}{uniquely}%
\TextAt{8}{356.091}{339.919}{to}%
\TextAt{13}{370.402}{339.919}{L}%
\TextAt{8}{378.366}{339.919}{.}%
\TextAt{8}{86.327}{322.873}{(ii)}%
\TextAt{8}{107.918}{322.873}{If}%
\TextAt{13}{118.774}{322.873}{v}%
\TextAt{8}{127.923}{322.873}{is}%
\TextAt{8}{138.85}{322.873}{nonarchimedean,}%
\TextAt{13}{228.262}{322.873}{w}%
\TextAt{21}{236.63}{321.08}{Z}%
\TextAt{8}{246.49}{322.873}{has}%
\TextAt{8}{266.527}{322.873}{the}%
\TextAt{8}{285.847}{322.873}{same}%
\TextAt{8}{314.336}{322.873}{residue}%
\TextAt{8}{353.226}{322.873}{class}%
\TextAt{8}{379.827}{322.873}{field}%
\TextAt{8}{404.347}{322.873}{and}%
\TextAt{8}{426.273}{322.873}{the}%
\TextAt{8}{445.592}{322.873}{same}%
\TextAt{8}{474.082}{322.873}{value}%
\TextAt{8}{503.491}{322.873}{group}%
\TextAt{8}{107.918}{305.828}{as}%
\TextAt{13}{122.289}{305.828}{v}%
\TextAt{8}{128.377}{305.828}{.}%
\TextAt{8}{83.076}{288.782}{(iii)}%
\TextAt{13}{107.918}{288.782}{Z}%
\TextAt{21}{115.931}{286.988}{w}%
\TextAt{9}{126.041}{288.782}{\Character{61}}%
\TextAt{13}{138.466}{288.782}{L}%
\TextAt{16}{149.088}{288.782}{\Character{92}}%
\TextAt{13}{159.714}{288.782}{K}%
\TextAt{21}{169.674}{286.988}{v}%
\TextAt{8}{178.479}{288.782}{(the}%
\TextAt{8}{203.19}{288.782}{intersection}%
\TextAt{8}{266.014}{288.782}{is}%
\TextAt{8}{277.778}{288.782}{taken}%
\TextAt{8}{309.651}{288.782}{inside}%
\TextAt{13}{342.874}{288.782}{L}%
\TextAt{21}{350.839}{286.988}{w}%
\TextAt{8}{357.628}{288.782}{).}%
\TextAt{19}{79.226}{254.548}{Proof.}%
\TextAt{8}{116}{254.548}{(i)}%
\TextAt{8}{131.444}{254.548}{An}%
\TextAt{8}{149.804}{254.548}{arbitrary}%
\TextAt{8}{198.758}{254.548}{extension}%
\TextAt{13}{249.708}{254.548}{w}%
\TextAt{3}{258.398}{258.886}{\Character{48}}%
\TextAt{8}{264.274}{254.548}{of}%
\TextAt{13}{276.789}{254.548}{w}%
\TextAt{21}{285.157}{252.754}{Z}%
\TextAt{8}{295.038}{254.548}{to}%
\TextAt{13}{308.533}{254.548}{L}%
\TextAt{8}{319.579}{254.548}{is}%
\TextAt{8}{330.527}{254.548}{conjugate}%
\TextAt{8}{382.721}{254.548}{to}%
\TextAt{13}{396.215}{254.548}{w}%
\TextAt{8}{407.987}{254.548}{over}%
\TextAt{13}{432.217}{254.548}{Z}%
\TextAt{21}{440.229}{252.754}{w}%
\TextAt{8}{447.019}{254.548}{;}%
\TextAt{8}{453.624}{254.548}{thus}%
\TextAt{13}{478.56}{254.548}{w}%
\TextAt{3}{487.249}{258.886}{\Character{48}}%
\TextAt{9}{493.365}{254.548}{\Character{61}}%
\TextAt{13}{505.791}{254.547}{w}%
\TextAt{16}{515.469}{254.548}{\Character{14}}%
\TextAt{13}{522.435}{254.547}{\Character{27}}%
\TextAt{8}{529.518}{254.548}{,}%
\TextAt{8}{79.226}{237.502}{for}%
\TextAt{8}{97.111}{237.502}{some}%
\TextAt{13}{126.437}{237.502}{\Character{27}}%
\NativeMathAt{136.84}{237.502}{11.955}{\in}%
\TextAt{13}{148.131}{237.502}{G}%
\TextAt{9}{157.365}{237.502}{\Character{40}}%
\TextAt{13}{161.917}{237.502}{L}%
\TextAt{16}{173.203}{237.502}{\Character{106}}%
\TextAt{13}{179.844}{237.502}{Z}%
\TextAt{21}{187.857}{235.708}{w}%
\TextAt{9}{194.646}{237.502}{\Character{41}}%
\TextAt{9}{202.501}{237.502}{\Character{61}}%
\TextAt{13}{214.945}{237.502}{G}%
\TextAt{21}{224.178}{235.708}{w}%
\TextAt{8}{230.968}{237.502}{,}%
\TextAt{8}{238.117}{237.502}{i.e.,}%
\TextAt{13}{260.222}{237.502}{w}%
\TextAt{3}{268.912}{241.84}{\Character{48}}%
\TextAt{9}{275.027}{237.502}{\Character{61}}%
\TextAt{13}{287.453}{237.502}{w}%
\TextAt{8}{296.142}{237.502}{.}%
\TextAt{8}{97.159}{217.467}{(iii)}%
\TextAt{8}{119.922}{217.467}{The}%
\TextAt{8}{143.975}{217.467}{identity}%
\TextAt{13}{187.229}{217.467}{Z}%
\TextAt{21}{195.241}{215.674}{w}%
\TextAt{9}{205.352}{217.467}{\Character{61}}%
\TextAt{13}{217.777}{217.467}{L}%
\TextAt{16}{228.398}{217.467}{\Character{92}}%
\TextAt{13}{239.025}{217.467}{K}%
\TextAt{21}{248.984}{215.674}{v}%
\TextAt{8}{257.789}{217.467}{follows}%
\TextAt{8}{296.225}{217.467}{immediately}%
\TextAt{8}{362.887}{217.467}{from}%
\TextAt{13}{390.527}{217.467}{G}%
\TextAt{21}{399.761}{215.674}{w}%
\TextAt{9}{406.55}{217.467}{\Character{40}}%
\TextAt{13}{411.103}{217.467}{L}%
\TextAt{16}{422.388}{217.467}{\Character{106}}%
\TextAt{13}{429.03}{217.467}{K}%
\TextAt{9}{439.848}{217.467}{\Character{41}}%
\NativeMathAt{447.721}{217.467}{11.955}{\cong}%
\TextAt{13}{460.34}{217.467}{G}%
\TextAt{9}{469.574}{217.467}{\Character{40}}%
\TextAt{13}{474.126}{217.467}{L}%
\TextAt{21}{482.09}{215.674}{w}%
\TextAt{16}{492.201}{217.467}{\Character{106}}%
\TextAt{13}{498.842}{217.467}{K}%
\TextAt{21}{508.802}{215.674}{v}%
\TextAt{9}{513.709}{217.467}{\Character{41}}%
\TextAt{8}{518.261}{217.467}{.}%
\TextAt{8}{97.159}{197.432}{(ii)}%
\TextAt{8}{116.093}{197.432}{Since}%
\TextAt{13}{146.074}{197.432}{K}%
\TextAt{21}{156.033}{195.639}{v}%
\TextAt{8}{164.261}{197.432}{has}%
\TextAt{8}{184.558}{197.432}{the}%
\TextAt{8}{204.138}{197.432}{same}%
\TextAt{8}{232.887}{197.432}{residue}%
\TextAt{8}{272.038}{197.432}{class}%
\TextAt{8}{298.898}{197.432}{field}%
\TextAt{8}{323.678}{197.432}{and}%
\TextAt{8}{345.864}{197.432}{the}%
\TextAt{8}{365.444}{197.432}{same}%
\TextAt{8}{394.193}{197.432}{value}%
\TextAt{8}{423.863}{197.432}{group}%
\TextAt{8}{456.462}{197.432}{as}%
\TextAt{13}{470.256}{197.432}{K}%
\TextAt{8}{481.073}{197.432}{,}%
\TextAt{8}{487.761}{197.432}{the}%
\TextAt{8}{507.341}{197.432}{same}%
\TextAt{8}{79.226}{180.386}{holds}%
\TextAt{8}{109.855}{180.386}{true}%
\TextAt{8}{134.567}{180.386}{for}%
\TextAt{13}{152.452}{180.386}{Z}%
\TextAt{21}{160.464}{178.593}{w}%
\TextAt{9}{170.574}{180.386}{\Character{61}}%
\TextAt{13}{183}{180.386}{L}%
\TextAt{16}{193.621}{180.386}{\Character{92}}%
\TextAt{13}{204.248}{180.386}{K}%
\TextAt{21}{214.207}{178.593}{v}%
\TextAt{8}{219.114}{180.386}{.}%
\Vector{q 0 G 0 g 0.399 w 452.334 -593.614 m 452.334 -585.544 l S Q}%
\Vector{q 0 G 0 g 0.399 w 452.534 -585.743 m 459.707 -585.743 l S Q}%
\Vector{q 0 G 0 g 0.399 w 452.534 -593.415 m 459.707 -593.415 l S Q}%
\Vector{q 0 G 0 g 0.399 w 459.906 -593.614 m 459.906 -585.544 l S Q}%
\TextAt{8}{97.159}{151.757}{The}%
\TextAt{12}{122.148}{151.757}{inertia}%
\TextAt{12}{166.071}{151.757}{group}%
\TextAt{13}{204.808}{151.757}{I}%
\TextAt{21}{209.971}{149.964}{w}%
\TextAt{8}{221.593}{151.757}{is}%
\TextAt{8}{234.292}{151.757}{defined}%
\TextAt{8}{275.54}{151.757}{only}%
\TextAt{8}{302.167}{151.757}{if}%
\TextAt{13}{313.826}{151.757}{w}%
\TextAt{8}{327.348}{151.757}{is}%
\TextAt{8}{340.047}{151.757}{a}%
\TextAt{8}{350.738}{151.757}{nonarchimedean}%
\TextAt{8}{438.503}{151.757}{valuation}%
\TextAt{8}{490.511}{151.757}{of}%
\TextAt{13}{504.776}{151.757}{L}%
\TextAt{8}{512.741}{151.757}{.}%
\TextAt{8}{523.994}{151.757}{It}%
\TextAt{8}{79.226}{134.711}{is}%
\TextAt{8}{91.921}{134.711}{the}%
\TextAt{8}{113.009}{134.711}{kernel}%
\TextAt{8}{148.408}{134.711}{of}%
\TextAt{8}{162.669}{134.711}{a}%
\TextAt{8}{173.356}{134.711}{canonical}%
\TextAt{8}{225.671}{134.711}{homomorphism}%
\TextAt{8}{309.272}{134.711}{of}%
\TextAt{13}{323.534}{134.711}{G}%
\TextAt{21}{332.767}{132.918}{w}%
\TextAt{8}{339.557}{134.711}{.}%
\TextAt{8}{350.799}{134.711}{For}%
\TextAt{8}{372.7}{134.711}{if}%
\TextAt{16}{384.355}{134.711}{\Character{79}}%
\TextAt{8}{399.034}{134.711}{is}%
\TextAt{8}{411.729}{134.711}{the}%
\TextAt{8}{432.817}{134.711}{valuation}%
\TextAt{8}{484.821}{134.711}{ring}%
\TextAt{8}{509.818}{134.711}{of}%
\TextAt{13}{524.08}{134.711}{w}%
\TextAt{8}{79.226}{117.665}{and}%
\TextAt{14}{102.833}{117.665}{P}%
\TextAt{8}{117.468}{117.665}{the}%
\TextAt{8}{138.468}{117.665}{maximal}%
\TextAt{8}{187.121}{117.665}{ideal,}%
\TextAt{8}{219.391}{117.665}{then,}%
\TextAt{8}{250.358}{117.665}{since}%
\TextAt{13}{279.87}{117.665}{\Character{27}}%
\TextAt{16}{286.953}{117.665}{\Character{79}}%
\TextAt{9}{301.562}{117.665}{\Character{61}}%
\TextAt{16}{315.426}{117.665}{\Character{79}}%
\TextAt{8}{330.017}{117.665}{and}%
\TextAt{13}{353.624}{117.665}{\Character{27}}%
\TextAt{14}{360.706}{117.665}{P}%
\TextAt{9}{375.359}{117.665}{\Character{61}}%
\TextAt{14}{389.222}{117.665}{P}%
\TextAt{8}{399.116}{117.665}{,}%
\TextAt{8}{407.32}{117.665}{every}%
\TextAt{13}{439.056}{117.665}{\Character{27}}%
\NativeMathAt{450.898}{117.665}{11.955}{\in}%
\TextAt{13}{463.627}{117.665}{G}%
\TextAt{21}{472.861}{115.872}{w}%
\TextAt{8}{484.392}{117.665}{induces}%
\TextAt{8}{526.912}{117.665}{a}%
\TextAt{13}{79.226}{100.619}{\Character{20}}%
\TextAt{8}{85.965}{100.619}{-automorphism}%
\TextAt{9}{181.578}{72.921}{\Character{22}}%
\TextAt{13}{180.963}{72.921}{\Character{27}}%
\TextAt{9}{191.366}{72.921}{\Character{58}}%
\TextAt{16}{197.939}{72.921}{\Character{79}}%
\TextAt{9}{207.788}{72.921}{\Character{47}}%
\TextAt{14}{213.641}{72.921}{P}%
\NativeArrowAt{226.855}{72.921}{11.955}{19.26}{\longrightarrow}%
\TextAt{16}{249.426}{72.921}{\Character{79}}%
\TextAt{9}{259.287}{72.921}{\Character{47}}%
\TextAt{14}{265.14}{72.921}{P}%
\TextAt{13}{275.033}{72.921}{\Character{59}}%
\TextAt{13}{304.168}{72.921}{x}%
\TextAt{8}{314.161}{72.921}{mod}%
\TextAt{14}{339.599}{72.921}{P}%
\NativeArrowAt{352.813}{72.921}{11.955}{19.26}{\longmapsto}%
\TextAt{13}{375.395}{72.921}{\Character{27}}%
\TextAt{13}{382.461}{72.921}{x}%
\TextAt{8}{392.45}{72.921}{mod}%
\TextAt{14}{417.888}{72.921}{P}%
\TextAt{13}{427.781}{72.921}{\Character{59}}%
\TextAt{8}{297.211}{43.033}{157}%
\Destination{79.23}{353.965}{nav.Item.56}%
\Destination{79.23}{336.915}{nav.Item.57}%
\Destination{79.23}{305.705}{nav.Item.58}%
\Destination{78.23}{809.815}{nav.page.157}%
\Destination{79.36}{776.155}{nav.print.171}%
\Destination{79.36}{337.855}{nav.print.172}%
\Destination{79.36}{497.655}{nav.number.2.9.7}%
\Destination{79.36}{373.955}{nav.number.2.9.8}%
\Hotlink{167.194}{596.555}{5.22}{11.047}{nav.chapter.1}%
\Hotlink{185.116}{596.555}{6.858}{11.047}{nav.section.1.9}%
\Hotlink{470.587}{721.855}{25.078}{11.047}{nav.number.2.9.6}%
}

% Page 170 (body page 158)
\Page{%
\TextAt{2}{79.226}{796.943}{II.}%
\TextAt{2}{94.073}{796.943}{The}%
\TextAt{2}{116.492}{796.943}{Theory}%
\TextAt{2}{154.401}{796.943}{of}%
\TextAt{2}{166.826}{796.943}{Valuations}%
\Vector{q 0 G 0 g 0.299 w 7.226 18.729 m 460.769 18.729 l S Q}%
\TextAt{8}{79.226}{761.124}{of}%
\TextAt{8}{92.556}{761.124}{the}%
\TextAt{8}{112.713}{761.124}{residue}%
\TextAt{8}{152.44}{761.124}{class}%
\TextAt{8}{179.877}{761.124}{field}%
\TextAt{13}{205.234}{761.124}{\Character{21}}%
\TextAt{8}{212.062}{761.124}{,}%
\TextAt{8}{219.211}{761.124}{and}%
\TextAt{8}{241.974}{761.124}{we}%
\TextAt{8}{259.201}{761.124}{obtain}%
\TextAt{8}{295.629}{761.124}{a}%
\TextAt{8}{305.384}{761.124}{homomorphism}%
\TextAt{13}{264.462}{737.045}{G}%
\TextAt{21}{273.696}{735.252}{w}%
\NativeArrowAt{283.806}{737.045}{11.955}{19.26}{\longrightarrow}%
\TextAt{8}{306.388}{737.045}{Aut}%
\TextAt{21}{326.222}{735.252}{\Character{20}}%
\TextAt{9}{331.6}{737.045}{\Character{40}}%
\TextAt{13}{336.152}{737.045}{\Character{21}}%
\TextAt{9}{342.981}{737.045}{\Character{41}}%
\TextAt{8}{79.226}{712.967}{with}%
\TextAt{8}{105.886}{712.967}{kernel}%
\TextAt{13}{140.353}{712.967}{I}%
\TextAt{21}{145.516}{711.173}{w}%
\TextAt{8}{152.305}{712.967}{.}%
\TextAt{12}{79.226}{688.361}{Proposition}%
\TextAt{12}{153.012}{688.361}{(9.9).}%
\TextAt{8}{189.262}{688.361}{The}%
\TextAt{8}{214.251}{688.361}{residue}%
\TextAt{8}{254.912}{688.361}{class}%
\TextAt{8}{283.284}{688.361}{field}%
\TextAt{8}{309.575}{688.361}{extension}%
\TextAt{13}{362.276}{688.361}{\Character{21}}%
\TextAt{16}{374.017}{688.361}{\Character{106}}%
\TextAt{13}{382.252}{688.361}{\Character{20}}%
\TextAt{8}{393.822}{688.361}{is}%
\TextAt{8}{406.52}{688.361}{normal,}%
\TextAt{8}{450.62}{688.361}{and}%
\TextAt{8}{474.317}{688.361}{we}%
\TextAt{8}{492.479}{688.361}{have}%
\TextAt{8}{520.408}{688.361}{an}%
\TextAt{8}{79.226}{671.315}{exact}%
\TextAt{8}{110.118}{671.315}{sequence}%
\TextAt{9}{211.792}{654.269}{1}%
\NativeArrowAt{220.965}{654.269}{11.955}{19.26}{\longrightarrow}%
\TextAt{13}{243.548}{654.269}{I}%
\TextAt{21}{248.71}{652.475}{w}%
\NativeArrowAt{258.821}{654.269}{11.955}{19.26}{\longrightarrow}%
\TextAt{13}{281.403}{654.269}{G}%
\TextAt{21}{290.636}{652.475}{w}%
\NativeArrowAt{300.747}{654.269}{11.955}{19.26}{\longrightarrow}%
\TextAt{13}{323.329}{654.269}{G}%
\TextAt{9}{332.562}{654.269}{\Character{40}}%
\TextAt{13}{337.115}{654.269}{\Character{21}}%
\TextAt{16}{347.264}{654.269}{\Character{106}}%
\TextAt{13}{353.906}{654.269}{\Character{20}}%
\TextAt{9}{360.644}{654.269}{\Character{41}}%
\NativeArrowAt{368.517}{654.269}{11.955}{19.26}{\longrightarrow}%
\TextAt{9}{391.099}{654.269}{1}%
\TextAt{13}{396.952}{654.269}{\Character{58}}%
\TextAt{19}{79.226}{617.737}{Proof.}%
\TextAt{8}{116}{617.737}{In}%
\TextAt{8}{131.597}{617.737}{the}%
\TextAt{8}{152.728}{617.737}{case}%
\TextAt{8}{178.475}{617.737}{of}%
\TextAt{8}{192.78}{617.737}{a}%
\TextAt{8}{203.511}{617.737}{finite}%
\TextAt{8}{234.398}{617.737}{Galois}%
\TextAt{8}{271.286}{617.737}{extension,}%
\TextAt{8}{327.523}{617.737}{we}%
\TextAt{8}{345.726}{617.737}{have}%
\TextAt{8}{373.696}{617.737}{proved}%
\TextAt{8}{412.724}{617.737}{this}%
\TextAt{8}{436.522}{617.737}{already}%
\TextAt{8}{478.802}{617.737}{in}%
\TextAt{8}{493.43}{617.737}{chap.}%
\TextAt{8}{525.297}{617.737}{I,}%
\TextAt{8}{79.226}{600.691}{(9.4).}%
\TextAt{8}{111.649}{600.691}{In}%
\TextAt{8}{125.961}{600.691}{the}%
\TextAt{8}{145.808}{600.691}{infinite}%
\TextAt{8}{185.166}{600.691}{case}%
\TextAt{13}{209.627}{600.691}{\Character{21}}%
\TextAt{16}{219.777}{600.691}{\Character{106}}%
\TextAt{13}{226.418}{600.691}{\Character{20}}%
\TextAt{8}{236.745}{600.691}{is}%
\TextAt{8}{248.2}{600.691}{normal}%
\TextAt{8}{287.569}{600.691}{since}%
\TextAt{13}{315.928}{600.691}{L}%
\TextAt{16}{327.214}{600.691}{\Character{106}}%
\TextAt{13}{333.855}{600.691}{K}%
\TextAt{8}{344.673}{600.691}{,}%
\TextAt{8}{351.575}{600.691}{and}%
\TextAt{8}{374.028}{600.691}{hence}%
\TextAt{8}{406.225}{600.691}{also}%
\TextAt{13}{429.395}{600.691}{\Character{21}}%
\TextAt{16}{439.545}{600.691}{\Character{106}}%
\TextAt{13}{446.186}{600.691}{\Character{20}}%
\TextAt{8}{452.925}{600.691}{,}%
\TextAt{8}{459.827}{600.691}{is}%
\TextAt{8}{471.281}{600.691}{the}%
\TextAt{8}{491.128}{600.691}{union}%
\TextAt{8}{523.337}{600.691}{of}%
\TextAt{8}{79.226}{583.645}{the}%
\TextAt{8}{99.63}{583.645}{finite}%
\TextAt{8}{129.79}{583.645}{normal}%
\TextAt{8}{169.717}{583.645}{subextensions.}%
\TextAt{8}{249.337}{583.645}{In}%
\TextAt{8}{264.206}{583.645}{order}%
\TextAt{8}{295.023}{583.645}{to}%
\TextAt{8}{309.581}{583.645}{prove}%
\TextAt{8}{341.379}{583.645}{the}%
\TextAt{8}{361.783}{583.645}{surjectivity}%
\TextAt{8}{423.23}{583.645}{of}%
\TextAt{13}{436.808}{583.645}{f}%
\TextAt{9}{447.597}{583.645}{\Character{58}}%
\TextAt{13}{454.592}{583.645}{G}%
\TextAt{21}{463.826}{581.852}{w}%
\NativeMathAt{474.358}{583.645}{11.955}{\rightarrow}%
\TextAt{13}{490.057}{583.645}{G}%
\TextAt{9}{499.29}{583.645}{\Character{40}}%
\TextAt{13}{503.843}{583.645}{\Character{21}}%
\TextAt{16}{514.414}{583.645}{\Character{106}}%
\TextAt{13}{521.479}{583.645}{\Character{20}}%
\TextAt{9}{528.217}{583.645}{\Character{41}}%
\TextAt{8}{79.226}{566.599}{all}%
\TextAt{8}{94.618}{566.599}{one}%
\TextAt{8}{115.211}{566.599}{has}%
\TextAt{8}{135.218}{566.599}{to}%
\TextAt{8}{148.662}{566.599}{show}%
\TextAt{8}{176.798}{566.599}{is}%
\TextAt{8}{187.695}{566.599}{that}%
\TextAt{13}{212.197}{566.599}{f}%
\TextAt{9}{219.244}{566.599}{\Character{40}}%
\TextAt{13}{223.796}{566.599}{G}%
\TextAt{21}{233.03}{564.806}{w}%
\TextAt{9}{239.819}{566.599}{\Character{41}}%
\TextAt{8}{247.402}{566.599}{is}%
\TextAt{8}{258.299}{566.599}{dense}%
\TextAt{8}{289.353}{566.599}{in}%
\TextAt{13}{302.139}{566.599}{G}%
\TextAt{9}{311.372}{566.599}{\Character{40}}%
\TextAt{13}{315.925}{566.599}{\Character{21}}%
\TextAt{16}{326.074}{566.599}{\Character{106}}%
\TextAt{13}{332.716}{566.599}{\Character{20}}%
\TextAt{9}{339.454}{566.599}{\Character{41}}%
\TextAt{8}{347.037}{566.599}{because}%
\TextAt{13}{389.472}{566.599}{f}%
\TextAt{9}{396.519}{566.599}{\Character{40}}%
\TextAt{13}{401.071}{566.599}{G}%
\TextAt{21}{410.304}{564.806}{w}%
\TextAt{9}{417.094}{566.599}{\Character{41}}%
\TextAt{8}{421.646}{566.599}{,}%
\TextAt{8}{428.102}{566.599}{being}%
\TextAt{8}{458.773}{566.599}{the}%
\TextAt{8}{478.063}{566.599}{continuous}%
\TextAt{8}{79.226}{549.553}{image}%
\TextAt{8}{112.97}{549.553}{of}%
\TextAt{8}{126.222}{549.553}{a}%
\TextAt{8}{135.9}{549.553}{compact}%
\TextAt{8}{182.651}{549.553}{set,}%
\TextAt{8}{204.108}{549.553}{is}%
\TextAt{8}{215.794}{549.553}{compact}%
\TextAt{8}{262.545}{549.553}{and}%
\TextAt{8}{285.23}{549.553}{hence}%
\TextAt{8}{317.658}{549.553}{closed.}%
\TextAt{8}{356.71}{549.553}{Let}%
\TextAt{9}{378.216}{549.553}{\Character{22}}%
\TextAt{13}{377.601}{549.553}{\Character{27}}%
\NativeMathAt{388.005}{549.553}{11.955}{\in}%
\TextAt{13}{399.296}{549.553}{G}%
\TextAt{9}{408.529}{549.553}{\Character{40}}%
\TextAt{13}{413.082}{549.553}{\Character{21}}%
\TextAt{16}{423.231}{549.553}{\Character{106}}%
\TextAt{13}{429.873}{549.553}{\Character{20}}%
\TextAt{9}{436.611}{549.553}{\Character{41}}%
\TextAt{8}{444.983}{549.553}{and}%
\TextAt{9}{468.283}{549.553}{\Character{22}}%
\TextAt{13}{467.668}{549.553}{\Character{27}}%
\TextAt{13}{474.733}{549.553}{G}%
\TextAt{9}{483.984}{549.553}{\Character{40}}%
\TextAt{13}{488.537}{549.553}{\Character{21}}%
\TextAt{16}{498.686}{549.553}{\Character{106}}%
\TextAt{13}{505.328}{549.553}{\Character{22}}%
\TextAt{9}{512.371}{549.553}{\Character{41}}%
\TextAt{8}{520.743}{549.553}{be}%
\TextAt{8}{79.226}{532.507}{a}%
\TextAt{8}{89.599}{532.507}{neighbourhood}%
\TextAt{8}{170.22}{532.507}{of}%
\TextAt{9}{184.782}{532.507}{\Character{22}}%
\TextAt{13}{184.167}{532.507}{\Character{27}}%
\TextAt{8}{191.25}{532.507}{,}%
\TextAt{8}{199.171}{532.507}{where}%
\TextAt{13}{233.597}{532.507}{\Character{22}}%
\TextAt{16}{245.013}{532.507}{\Character{106}}%
\TextAt{13}{252.707}{532.507}{\Character{20}}%
\TextAt{8}{263.96}{532.507}{is}%
\TextAt{8}{276.341}{532.507}{a}%
\TextAt{8}{286.714}{532.507}{finite}%
\TextAt{8}{317.243}{532.507}{Galois}%
\TextAt{8}{353.774}{532.507}{subextension}%
\TextAt{8}{424.102}{532.507}{of}%
\TextAt{13}{438.049}{532.507}{\Character{21}}%
\TextAt{16}{449.25}{532.507}{\Character{106}}%
\TextAt{13}{456.944}{532.507}{\Character{20}}%
\TextAt{8}{463.682}{532.507}{.}%
\TextAt{8}{473.983}{532.507}{We}%
\TextAt{8}{494.744}{532.507}{have}%
\TextAt{8}{522.357}{532.507}{to}%
\TextAt{8}{79.226}{515.461}{show}%
\TextAt{8}{108.392}{515.461}{that}%
\TextAt{8}{133.924}{515.461}{this}%
\TextAt{8}{156.909}{515.461}{neighbourhood}%
\TextAt{8}{237.076}{515.461}{contains}%
\TextAt{8}{283.158}{515.461}{an}%
\TextAt{8}{299.58}{515.461}{element}%
\TextAt{8}{342.984}{515.461}{of}%
\TextAt{8}{356.477}{515.461}{the}%
\TextAt{8}{376.796}{515.461}{image}%
\TextAt{13}{410.78}{515.461}{f}%
\TextAt{9}{417.827}{515.461}{\Character{40}}%
\TextAt{13}{422.379}{515.461}{\Character{28}}%
\TextAt{9}{428.798}{515.461}{\Character{41}}%
\TextAt{8}{433.35}{515.461}{,}%
\TextAt{13}{440.703}{515.461}{\Character{28}}%
\NativeMathAt{450.72}{515.461}{11.955}{\in}%
\TextAt{13}{462.288}{515.461}{G}%
\TextAt{21}{471.522}{513.668}{w}%
\TextAt{8}{478.311}{515.461}{.}%
\TextAt{8}{487.248}{515.461}{Since}%
\TextAt{13}{517.968}{515.461}{Z}%
\TextAt{21}{525.98}{513.668}{w}%
\TextAt{8}{79.226}{498.415}{has}%
\TextAt{8}{100.636}{498.415}{the}%
\TextAt{8}{121.328}{498.415}{residue}%
\TextAt{8}{161.591}{498.415}{class}%
\TextAt{8}{189.565}{498.415}{field}%
\TextAt{13}{215.458}{498.415}{\Character{20}}%
\TextAt{8}{222.196}{498.415}{,}%
\TextAt{8}{230.015}{498.415}{there}%
\TextAt{8}{260.463}{498.415}{exists}%
\TextAt{8}{293.314}{498.415}{a}%
\TextAt{8}{303.606}{498.415}{finite}%
\TextAt{8}{334.053}{498.415}{Galois}%
\TextAt{8}{370.503}{498.415}{subextension}%
\TextAt{13}{440.749}{498.415}{M}%
\TextAt{16}{457.557}{498.415}{\Character{106}}%
\TextAt{13}{465.113}{498.415}{Z}%
\TextAt{21}{473.125}{496.622}{w}%
\TextAt{8}{484.348}{498.415}{of}%
\TextAt{13}{498.214}{498.415}{L}%
\TextAt{16}{510.413}{498.415}{\Character{106}}%
\TextAt{13}{517.968}{498.415}{Z}%
\TextAt{21}{525.98}{496.622}{w}%
\TextAt{8}{79.226}{481.369}{whose}%
\TextAt{8}{113.444}{481.369}{residue}%
\TextAt{8}{152.862}{481.369}{class}%
\TextAt{8}{179.99}{481.369}{field}%
\Vector{q 0 G 0 g 0.478 w 133.038 -282.788 m 145.612 -282.788 l S Q}%
\TextAt{13}{205.038}{481.369}{M}%
\TextAt{8}{221.201}{481.369}{contains}%
\TextAt{8}{266.812}{481.369}{the}%
\TextAt{8}{286.659}{481.369}{field}%
\TextAt{13}{311.707}{481.369}{\Character{22}}%
\TextAt{8}{318.75}{481.369}{.}%
\TextAt{8}{327.096}{481.369}{As}%
\TextAt{13}{344.074}{481.369}{G}%
\TextAt{9}{353.308}{481.369}{\Character{40}}%
\TextAt{13}{357.86}{481.369}{M}%
\TextAt{16}{373.754}{481.369}{\Character{106}}%
\TextAt{13}{380.396}{481.369}{Z}%
\TextAt{21}{388.408}{479.576}{w}%
\TextAt{9}{395.198}{481.369}{\Character{41}}%
\NativeMathAt{403.071}{481.369}{11.955}{\rightarrow}%
\TextAt{13}{418.347}{481.369}{G}%
\TextAt{9}{427.581}{481.369}{\Character{40}}%
\Vector{q 0 G 0 g 0.478 w 360.133 -282.788 m 372.707 -282.788 l S Q}%
\TextAt{13}{432.133}{481.369}{M}%
\TextAt{16}{448.027}{481.369}{\Character{106}}%
\TextAt{13}{454.669}{481.369}{\Character{20}}%
\TextAt{9}{461.408}{481.369}{\Character{41}}%
\TextAt{8}{469.549}{481.369}{is}%
\TextAt{8}{481.004}{481.369}{surjective,}%
\TextAt{8}{79.226}{464.323}{the}%
\TextAt{8}{99.383}{464.323}{composite}%
\TextAt{13}{153.04}{440.245}{G}%
\TextAt{21}{162.273}{438.451}{w}%
\TextAt{9}{172.384}{440.245}{\Character{61}}%
\TextAt{13}{184.809}{440.245}{G}%
\TextAt{9}{194.043}{440.245}{\Character{40}}%
\TextAt{13}{198.595}{440.245}{L}%
\TextAt{16}{209.88}{440.245}{\Character{106}}%
\TextAt{13}{216.522}{440.245}{Z}%
\TextAt{21}{224.534}{438.451}{w}%
\TextAt{9}{231.324}{440.245}{\Character{41}}%
\NativeArrowAt{239.197}{440.245}{11.955}{19.26}{\longrightarrow}%
\TextAt{13}{261.779}{440.245}{G}%
\TextAt{9}{271.012}{440.245}{\Character{40}}%
\TextAt{13}{275.565}{440.245}{M}%
\TextAt{16}{291.459}{440.245}{\Character{106}}%
\TextAt{13}{298.101}{440.245}{Z}%
\TextAt{21}{306.113}{438.451}{w}%
\TextAt{9}{312.903}{440.245}{\Character{41}}%
\NativeArrowAt{320.776}{440.245}{11.955}{19.26}{\longrightarrow}%
\TextAt{13}{343.358}{440.245}{G}%
\TextAt{9}{352.591}{440.245}{\Character{40}}%
\Vector{q 0 G 0 g 0.478 w 285.144 -323.912 m 297.717 -323.912 l S Q}%
\TextAt{13}{357.144}{440.245}{M}%
\TextAt{16}{373.038}{440.245}{\Character{106}}%
\TextAt{13}{379.68}{440.245}{\Character{20}}%
\TextAt{9}{386.418}{440.245}{\Character{41}}%
\NativeArrowAt{394.292}{440.245}{11.955}{19.26}{\longrightarrow}%
\TextAt{13}{416.874}{440.245}{G}%
\TextAt{9}{426.107}{440.245}{\Character{40}}%
\TextAt{13}{430.66}{440.245}{\Character{22}}%
\TextAt{16}{441.023}{440.245}{\Character{106}}%
\TextAt{13}{447.665}{440.245}{\Character{20}}%
\TextAt{9}{454.404}{440.245}{\Character{41}}%
\TextAt{8}{79.226}{416.166}{is}%
\TextAt{8}{90.455}{416.166}{also}%
\TextAt{8}{113.399}{416.166}{surjective,}%
\TextAt{8}{168.634}{416.166}{and}%
\TextAt{8}{190.862}{416.166}{if}%
\TextAt{13}{201.05}{416.166}{\Character{28}}%
\NativeMathAt{210.79}{416.166}{11.955}{\in}%
\TextAt{13}{222.081}{416.166}{G}%
\TextAt{21}{231.315}{414.373}{w}%
\TextAt{8}{241.466}{416.166}{is}%
\TextAt{8}{252.695}{416.166}{mapped}%
\TextAt{8}{296.704}{416.166}{to}%
\TextAt{9}{311.094}{416.166}{\Character{22}}%
\TextAt{13}{310.479}{416.166}{\Character{27}}%
\TextAt{16}{317.562}{416.166}{\Character{106}}%
\TextAt{21}{320.883}{414.373}{\Character{22}}%
\TextAt{8}{326.475}{416.166}{,}%
\TextAt{8}{333.196}{416.166}{then}%
\TextAt{13}{359.321}{416.166}{f}%
\TextAt{9}{366.367}{416.166}{\Character{40}}%
\TextAt{13}{370.92}{416.166}{\Character{28}}%
\TextAt{9}{377.339}{416.166}{\Character{41}}%
\NativeMathAt{385.212}{416.166}{11.955}{\in}%
\TextAt{9}{397.118}{416.166}{\Character{22}}%
\TextAt{13}{396.503}{416.166}{\Character{27}}%
\TextAt{13}{403.569}{416.166}{G}%
\TextAt{9}{412.819}{416.166}{\Character{40}}%
\TextAt{13}{417.371}{416.166}{\Character{21}}%
\TextAt{16}{427.52}{416.166}{\Character{106}}%
\TextAt{13}{434.162}{416.166}{\Character{22}}%
\TextAt{9}{441.205}{416.166}{\Character{41}}%
\TextAt{8}{445.757}{416.166}{,}%
\TextAt{8}{452.478}{416.166}{as}%
\TextAt{8}{466.313}{416.166}{required.}%
\Vector{q 0 G 0 g 0.399 w 452.334 -357.834 m 452.334 -349.764 l S Q}%
\Vector{q 0 G 0 g 0.399 w 452.534 -349.963 m 459.707 -349.963 l S Q}%
\Vector{q 0 G 0 g 0.399 w 452.534 -357.635 m 459.707 -357.635 l S Q}%
\Vector{q 0 G 0 g 0.399 w 459.906 -357.834 m 459.906 -349.764 l S Q}%
\TextAt{12}{79.226}{384}{Definition}%
\TextAt{12}{142.015}{384}{(9.10).}%
\TextAt{8}{184.061}{384}{The}%
\TextAt{8}{208.115}{384}{fixed}%
\TextAt{8}{236.401}{384}{field}%
\TextAt{8}{261.758}{384}{of}%
\TextAt{13}{275.088}{384}{I}%
\TextAt{21}{280.251}{382.207}{w}%
\TextAt{8}{287.04}{384}{,}%
\TextAt{13}{176.567}{359.922}{T}%
\TextAt{21}{183.428}{358.128}{w}%
\TextAt{9}{193.538}{359.922}{\Character{61}}%
\TextAt{13}{205.964}{359.922}{T}%
\TextAt{21}{212.825}{358.128}{w}%
\TextAt{9}{219.614}{359.922}{\Character{40}}%
\TextAt{13}{224.167}{359.922}{L}%
\TextAt{16}{235.452}{359.922}{\Character{106}}%
\TextAt{13}{242.094}{359.922}{K}%
\TextAt{9}{252.911}{359.922}{\Character{41}}%
\TextAt{9}{260.766}{359.922}{\Character{61}}%
\TextAt{16}{273.21}{359.922}{\Character{102}}%
\TextAt{13}{279.188}{359.922}{x}%
\NativeMathAt{289.161}{359.922}{11.955}{\in}%
\TextAt{13}{300.452}{359.922}{L}%
\TextAt{16}{311.737}{359.922}{\Character{106}}%
\TextAt{13}{318.379}{359.922}{\Character{27}}%
\TextAt{13}{325.445}{359.922}{x}%
\TextAt{9}{335.434}{359.922}{\Character{61}}%
\TextAt{13}{347.859}{359.922}{x}%
\TextAt{8}{358.409}{359.922}{for}%
\TextAt{8}{376.294}{359.922}{all}%
\TextAt{13}{392.553}{359.922}{\Character{27}}%
\NativeMathAt{402.956}{359.922}{11.955}{\in}%
\TextAt{13}{414.247}{359.922}{I}%
\TextAt{21}{419.41}{358.128}{w}%
\TextAt{16}{426.199}{359.922}{\Character{103}}%
\TextAt{13}{432.177}{359.922}{\Character{59}}%
\TextAt{8}{79.226}{335.843}{is}%
\TextAt{8}{90.99}{335.843}{called}%
\TextAt{8}{124.154}{335.843}{the}%
\TextAt{12}{144.31}{335.843}{inertia}%
\TextAt{12}{187.157}{335.843}{field}%
\TextAt{8}{215.874}{335.843}{of}%
\TextAt{13}{229.204}{335.843}{w}%
\TextAt{8}{241.791}{335.843}{over}%
\TextAt{13}{266.837}{335.843}{K}%
\TextAt{8}{277.654}{335.843}{.}%
\TextAt{8}{97.159}{311.237}{For}%
\TextAt{8}{118.128}{311.237}{the}%
\TextAt{8}{138.285}{311.237}{inertia}%
\TextAt{8}{175.358}{311.237}{field,}%
\TextAt{8}{203.966}{311.237}{(9.9)}%
\TextAt{8}{231.941}{311.237}{gives}%
\TextAt{8}{260.622}{311.237}{us}%
\TextAt{8}{275.638}{311.237}{the}%
\TextAt{8}{295.794}{311.237}{isomorphism}%
\TextAt{13}{247.091}{287.159}{G}%
\TextAt{9}{256.325}{287.159}{\Character{40}}%
\TextAt{13}{260.877}{287.159}{T}%
\TextAt{21}{267.738}{285.365}{w}%
\TextAt{16}{277.849}{287.159}{\Character{106}}%
\TextAt{13}{284.49}{287.159}{Z}%
\TextAt{21}{292.503}{285.365}{w}%
\TextAt{9}{299.292}{287.159}{\Character{41}}%
\NativeMathAt{307.165}{287.158}{11.955}{\cong}%
\TextAt{13}{319.785}{287.159}{G}%
\TextAt{9}{329.018}{287.159}{\Character{40}}%
\TextAt{13}{333.571}{287.159}{\Character{21}}%
\TextAt{16}{343.72}{287.159}{\Character{106}}%
\TextAt{13}{350.362}{287.159}{\Character{20}}%
\TextAt{9}{357.1}{287.159}{\Character{41}}%
\TextAt{13}{361.652}{287.159}{\Character{58}}%
\TextAt{8}{79.226}{263.08}{It}%
\TextAt{8}{91.899}{263.08}{has}%
\TextAt{8}{112.772}{263.08}{the}%
\TextAt{8}{132.929}{263.08}{following}%
\TextAt{8}{182.363}{263.08}{significance}%
\TextAt{8}{244.208}{263.08}{for}%
\TextAt{8}{262.092}{263.08}{the}%
\TextAt{8}{282.249}{263.08}{extension}%
\TextAt{13}{334.015}{263.08}{L}%
\TextAt{16}{345.3}{263.08}{\Character{106}}%
\TextAt{13}{351.942}{263.08}{K}%
\TextAt{8}{362.759}{263.08}{.}%
\TextAt{12}{79.226}{238.474}{Proposition}%
\TextAt{12}{151.937}{238.474}{(9.11).}%
\TextAt{13}{193.984}{238.474}{T}%
\TextAt{21}{200.845}{236.681}{w}%
\TextAt{16}{210.955}{238.474}{\Character{106}}%
\TextAt{13}{217.597}{238.474}{Z}%
\TextAt{21}{225.609}{236.681}{w}%
\TextAt{8}{236.296}{238.474}{is}%
\TextAt{8}{248.06}{238.474}{the}%
\TextAt{8}{268.216}{238.474}{maximal}%
\TextAt{8}{316.025}{238.474}{unramified}%
\TextAt{8}{374.557}{238.474}{subextension}%
\TextAt{8}{444.268}{238.474}{of}%
\TextAt{13}{457.598}{238.474}{L}%
\TextAt{16}{468.883}{238.474}{\Character{106}}%
\TextAt{13}{475.525}{238.474}{Z}%
\TextAt{21}{483.537}{236.681}{w}%
\TextAt{8}{490.327}{238.474}{.}%
\TextAt{19}{79.226}{206.308}{Proof.}%
\TextAt{8}{116}{206.308}{By}%
\TextAt{8}{134.92}{206.308}{(9.6),}%
\TextAt{8}{166.842}{206.308}{we}%
\TextAt{8}{184.626}{206.308}{may}%
\TextAt{8}{210.552}{206.308}{assume}%
\TextAt{8}{251.553}{206.308}{that}%
\TextAt{13}{277.478}{206.308}{K}%
\TextAt{9}{292.565}{206.308}{\Character{61}}%
\TextAt{13}{305.939}{206.308}{Z}%
\TextAt{21}{313.951}{204.515}{w}%
\TextAt{8}{325.194}{206.308}{is}%
\TextAt{8}{337.515}{206.308}{henselian.}%
\TextAt{8}{394.521}{206.308}{Let}%
\TextAt{13}{416.047}{206.308}{T}%
\TextAt{16}{428.803}{206.308}{\Character{106}}%
\TextAt{13}{436.393}{206.308}{K}%
\TextAt{8}{451.664}{206.308}{be}%
\TextAt{8}{468.145}{206.308}{the}%
\TextAt{8}{488.858}{206.308}{maximal}%
\TextAt{8}{79.226}{189.262}{unramified}%
\TextAt{8}{138.544}{189.262}{subextension}%
\TextAt{8}{209.04}{189.262}{of}%
\TextAt{13}{223.155}{189.262}{L}%
\TextAt{16}{235.779}{189.262}{\Character{106}}%
\TextAt{13}{243.759}{189.262}{K}%
\TextAt{8}{254.577}{189.262}{.}%
\TextAt{8}{265.381}{189.262}{It}%
\TextAt{8}{278.839}{189.262}{is}%
\TextAt{8}{291.388}{189.262}{Galois,}%
\TextAt{8}{331.535}{189.262}{since}%
\TextAt{8}{360.989}{189.262}{the}%
\TextAt{8}{381.93}{189.262}{conjugate}%
\TextAt{8}{435.725}{189.262}{extensions}%
\TextAt{8}{492.891}{189.262}{are}%
\TextAt{8}{513.187}{189.262}{also}%
\TextAt{8}{79.226}{172.216}{unramified.}%
\TextAt{8}{142.31}{172.216}{By}%
\TextAt{8}{160.673}{172.216}{(7.5),}%
\TextAt{13}{191.9}{172.216}{T}%
\TextAt{8}{204.284}{172.216}{has}%
\TextAt{8}{225.158}{172.216}{the}%
\TextAt{8}{245.314}{172.216}{residue}%
\TextAt{8}{285.041}{172.216}{class}%
\TextAt{8}{312.478}{172.216}{field}%
\TextAt{13}{337.835}{172.216}{\Character{21}}%
\TextAt{21}{344.664}{170.423}{s}%
\TextAt{8}{349.078}{172.216}{,}%
\TextAt{8}{356.227}{172.216}{and}%
\TextAt{8}{378.989}{172.216}{we}%
\TextAt{8}{396.217}{172.216}{have}%
\TextAt{8}{423.212}{172.216}{an}%
\TextAt{8}{439.471}{172.216}{isomorphism}%
\TextAt{13}{247.521}{148.138}{G}%
\TextAt{9}{256.755}{148.138}{\Character{40}}%
\TextAt{13}{261.307}{148.138}{T}%
\TextAt{16}{273.115}{148.138}{\Character{106}}%
\TextAt{13}{279.756}{148.138}{K}%
\TextAt{9}{290.574}{148.138}{\Character{41}}%
\NativeMathAt{301.035}{154.915}{7.97}{\sim}%
\NativeArrowAt{298.447}{148.138}{11.955}{13.175}{\longrightarrow}%
\TextAt{13}{314.941}{148.138}{G}%
\TextAt{9}{324.175}{148.138}{\Character{40}}%
\TextAt{13}{328.727}{148.138}{\Character{21}}%
\TextAt{21}{335.555}{146.345}{s}%
\TextAt{16}{343.29}{148.138}{\Character{106}}%
\TextAt{13}{349.932}{148.138}{\Character{20}}%
\TextAt{9}{356.671}{148.138}{\Character{41}}%
\TextAt{13}{361.223}{148.138}{\Character{58}}%
\TextAt{8}{79.226}{124.059}{Surjectivity}%
\TextAt{8}{142.247}{124.059}{follows}%
\TextAt{8}{180.617}{124.059}{from}%
\TextAt{8}{208.191}{124.059}{(9.9)}%
\TextAt{8}{236.1}{124.059}{and}%
\TextAt{8}{258.796}{124.059}{the}%
\TextAt{8}{278.886}{124.059}{injectivity}%
\TextAt{8}{334.101}{124.059}{from}%
\TextAt{8}{361.675}{124.059}{the}%
\TextAt{8}{381.765}{124.059}{fact}%
\TextAt{8}{404.784}{124.059}{that}%
\TextAt{13}{430.087}{124.059}{T}%
\TextAt{16}{441.894}{124.059}{\Character{106}}%
\TextAt{13}{448.536}{124.059}{K}%
\TextAt{8}{463.185}{124.059}{is}%
\TextAt{8}{474.883}{124.059}{unramified:}%
\TextAt{8}{79.226}{107.013}{every}%
\TextAt{8}{110.581}{107.013}{finite}%
\TextAt{8}{140.955}{107.013}{Galois}%
\TextAt{8}{177.331}{107.013}{subextension}%
\TextAt{8}{247.504}{107.013}{has}%
\TextAt{8}{268.841}{107.013}{the}%
\TextAt{8}{289.46}{107.013}{same}%
\TextAt{8}{319.248}{107.013}{degree}%
\TextAt{8}{356.126}{107.013}{as}%
\TextAt{8}{370.959}{107.013}{its}%
\TextAt{8}{387.74}{107.013}{residue}%
\TextAt{8}{427.93}{107.013}{class}%
\TextAt{8}{455.83}{107.013}{field}%
\TextAt{8}{481.649}{107.013}{extension.}%
\TextAt{8}{79.226}{89.967}{An}%
\TextAt{8}{98.374}{89.967}{element}%
\TextAt{13}{141.587}{89.967}{\Character{27}}%
\NativeMathAt{151.991}{89.967}{11.955}{\in}%
\TextAt{13}{163.282}{89.967}{G}%
\TextAt{9}{172.515}{89.967}{\Character{40}}%
\TextAt{13}{177.068}{89.967}{L}%
\TextAt{16}{188.353}{89.967}{\Character{106}}%
\TextAt{13}{194.995}{89.967}{K}%
\TextAt{9}{205.812}{89.967}{\Character{41}}%
\TextAt{8}{214.234}{89.967}{therefore}%
\TextAt{8}{263.305}{89.967}{induces}%
\TextAt{8}{304.953}{89.967}{the}%
\TextAt{8}{325.081}{89.967}{identity}%
\TextAt{8}{368.306}{89.967}{on}%
\TextAt{13}{384.537}{89.967}{\Character{21}}%
\TextAt{21}{391.365}{88.174}{s}%
\TextAt{8}{395.779}{89.967}{,}%
\TextAt{8}{402.906}{89.967}{i.e.,}%
\TextAt{8}{424.988}{89.967}{on}%
\TextAt{13}{441.219}{89.967}{\Character{21}}%
\TextAt{8}{448.048}{89.967}{,}%
\TextAt{8}{455.174}{89.967}{if}%
\TextAt{8}{465.87}{89.967}{and}%
\TextAt{8}{488.604}{89.967}{only}%
\TextAt{8}{514.267}{89.967}{if}%
\TextAt{8}{524.963}{89.967}{it}%
\TextAt{8}{79.226}{72.921}{belongs}%
\TextAt{8}{121.237}{72.921}{to}%
\TextAt{13}{135.547}{72.921}{G}%
\TextAt{9}{144.781}{72.921}{\Character{40}}%
\TextAt{13}{149.333}{72.921}{L}%
\TextAt{16}{160.618}{72.921}{\Character{106}}%
\TextAt{13}{167.26}{72.921}{T}%
\TextAt{9}{175.747}{72.921}{\Character{41}}%
\TextAt{8}{180.299}{72.921}{.}%
\TextAt{8}{188.747}{72.921}{Consequently,}%
\TextAt{13}{263.599}{72.921}{G}%
\TextAt{9}{272.832}{72.921}{\Character{40}}%
\TextAt{13}{277.385}{72.921}{L}%
\TextAt{16}{288.67}{72.921}{\Character{106}}%
\TextAt{13}{295.312}{72.921}{T}%
\TextAt{9}{303.799}{72.921}{\Character{41}}%
\TextAt{9}{311.654}{72.921}{\Character{61}}%
\TextAt{13}{324.097}{72.921}{I}%
\TextAt{21}{329.26}{71.128}{w}%
\TextAt{8}{336.05}{72.921}{,}%
\TextAt{8}{343.199}{72.921}{hence}%
\TextAt{13}{375.705}{72.921}{T}%
\TextAt{9}{387.513}{72.921}{\Character{61}}%
\TextAt{13}{399.938}{72.921}{T}%
\TextAt{21}{406.799}{71.128}{w}%
\TextAt{8}{413.588}{72.921}{.}%
\Vector{q 0 G 0 g 0.399 w 452.334 -701.079 m 452.334 -693.009 l S Q}%
\Vector{q 0 G 0 g 0.399 w 452.534 -693.208 m 459.707 -693.208 l S Q}%
\Vector{q 0 G 0 g 0.399 w 452.534 -700.879 m 459.707 -700.879 l S Q}%
\Vector{q 0 G 0 g 0.399 w 459.906 -701.079 m 459.906 -693.009 l S Q}%
\TextAt{8}{297.211}{43.033}{158}%
\Destination{78.23}{809.815}{nav.page.158}%
\Destination{79.36}{513.455}{nav.print.173}%
\Destination{79.36}{702.355}{nav.number.2.9.9}%
\Destination{79.36}{397.955}{nav.number.2.9.10}%
\Destination{79.36}{252.455}{nav.number.2.9.11}%
\Hotlink{524.797}{615.555}{5.22}{11.047}{nav.chapter.1}%
\Hotlink{78.726}{598.455}{25.078}{11.047}{nav.number.1.9.4}%
\Hotlink{203.466}{309.055}{25.078}{11.047}{nav.number.2.9.9}%
\Hotlink{134.42}{204.155}{25.078}{11.047}{nav.number.2.9.6}%
\Hotlink{160.173}{170.055}{25.078}{11.047}{nav.number.2.7.5}%
\Hotlink{207.691}{121.855}{25.078}{11.047}{nav.number.2.9.9}%
}

% Page 171 (body page 159)
\Page{%
\TextAt{2}{79.226}{796.943}{II.}%
\TextAt{2}{94.073}{796.943}{The}%
\TextAt{2}{116.492}{796.943}{Theory}%
\TextAt{2}{154.401}{796.943}{of}%
\TextAt{2}{166.826}{796.943}{Valuations}%
\Vector{q 0 G 0 g 0.299 w 7.226 18.729 m 460.769 18.729 l S Q}%
\TextAt{8}{97.159}{761.124}{If}%
\TextAt{8}{108.252}{761.124}{in}%
\TextAt{8}{121.307}{761.124}{particular}%
\TextAt{13}{174.697}{761.124}{K}%
\TextAt{8}{188.814}{761.124}{is}%
\TextAt{8}{199.979}{761.124}{a}%
\TextAt{8}{209.136}{761.124}{henselian}%
\TextAt{8}{259.323}{761.124}{field}%
\TextAt{8}{284.081}{761.124}{and}%
\Vector{q 0 G 0 g 0.478 w 234.245 -3.033 m 245.063 -3.033 l S Q}%
\TextAt{13}{306.245}{761.124}{K}%
\TextAt{21}{317.063}{759.33}{s}%
\TextAt{16}{324.798}{761.124}{\Character{106}}%
\TextAt{13}{331.439}{761.124}{K}%
\TextAt{8}{345.556}{761.124}{its}%
\TextAt{8}{361.276}{761.124}{separable}%
\TextAt{8}{412.12}{761.124}{closure,}%
\TextAt{8}{453.975}{761.124}{then}%
\TextAt{8}{480.036}{761.124}{the}%
\TextAt{8}{499.594}{761.124}{inertia}%
\TextAt{8}{79.226}{744.078}{field}%
\TextAt{8}{104.933}{744.078}{of}%
\TextAt{8}{118.612}{744.078}{this}%
\TextAt{8}{141.785}{744.078}{extension}%
\TextAt{8}{193.9}{744.078}{is}%
\TextAt{8}{206.014}{744.078}{the}%
\TextAt{8}{226.52}{744.078}{maximal}%
\TextAt{8}{274.678}{744.078}{unramified}%
\TextAt{8}{333.56}{744.078}{extension}%
\TextAt{13}{385.676}{744.078}{T}%
\TextAt{16}{398.08}{744.078}{\Character{106}}%
\TextAt{13}{405.317}{744.078}{K}%
\TextAt{8}{420.382}{744.078}{and}%
\TextAt{8}{443.494}{744.078}{has}%
\TextAt{8}{464.718}{744.078}{the}%
\TextAt{8}{485.224}{744.078}{separable}%
\TextAt{8}{79.226}{727.032}{closure}%
\TextAt{9}{118.75}{727.032}{\Character{22}}%
\TextAt{13}{118.308}{727.032}{\Character{20}}%
\TextAt{21}{125.046}{725.238}{s}%
\TextAt{16}{132.781}{727.032}{\Character{106}}%
\TextAt{13}{139.423}{727.032}{\Character{20}}%
\TextAt{8}{150.059}{727.032}{as}%
\TextAt{8}{164.429}{727.032}{its}%
\TextAt{8}{180.748}{727.032}{residue}%
\TextAt{8}{220.475}{727.032}{class}%
\TextAt{8}{247.912}{727.032}{field.}%
\TextAt{8}{277.82}{727.032}{The}%
\TextAt{8}{301.873}{727.032}{isomorphism}%
\TextAt{13}{251.129}{698.393}{G}%
\TextAt{9}{260.363}{698.393}{\Character{40}}%
\TextAt{13}{264.915}{698.393}{T}%
\TextAt{16}{276.723}{698.393}{\Character{106}}%
\TextAt{13}{283.364}{698.393}{K}%
\TextAt{9}{294.182}{698.393}{\Character{41}}%
\NativeMathAt{302.055}{698.393}{11.955}{\cong}%
\TextAt{13}{314.675}{698.393}{G}%
\TextAt{9}{323.908}{698.393}{\Character{40}}%
\TextAt{9}{328.893}{698.393}{\Character{22}}%
\TextAt{13}{328.46}{698.393}{\Character{20}}%
\TextAt{21}{335.199}{696.6}{s}%
\TextAt{16}{342.934}{698.393}{\Character{106}}%
\TextAt{13}{349.576}{698.393}{\Character{20}}%
\TextAt{9}{356.314}{698.393}{\Character{41}}%
\TextAt{8}{79.226}{669.754}{shows}%
\TextAt{8}{112.742}{669.754}{that}%
\TextAt{8}{138.009}{669.754}{the}%
\TextAt{8}{158.064}{669.754}{unramified}%
\TextAt{8}{216.494}{669.754}{extensions}%
\TextAt{8}{272.773}{669.754}{of}%
\TextAt{13}{286.001}{669.754}{K}%
\TextAt{8}{300.614}{669.754}{correspond}%
\TextAt{8}{360.085}{669.754}{1–1}%
\TextAt{8}{381.454}{669.754}{to}%
\TextAt{8}{395.663}{669.754}{the}%
\TextAt{8}{415.717}{669.754}{separable}%
\TextAt{8}{467.058}{669.754}{extensions}%
\TextAt{8}{523.337}{669.754}{of}%
\TextAt{13}{79.226}{652.708}{\Character{20}}%
\TextAt{8}{85.965}{652.708}{.}%
\TextAt{8}{97.159}{632.673}{Like}%
\TextAt{8}{123.177}{632.673}{the}%
\TextAt{8}{143.827}{632.673}{inertia}%
\TextAt{8}{181.394}{632.673}{group,}%
\TextAt{8}{218.439}{632.673}{the}%
\TextAt{12}{239.089}{632.673}{ramification}%
\TextAt{12}{314.305}{632.673}{group}%
\TextAt{13}{352.601}{632.673}{R}%
\TextAt{21}{361.511}{630.88}{w}%
\TextAt{8}{372.692}{632.673}{is}%
\TextAt{8}{384.949}{632.673}{the}%
\TextAt{8}{405.6}{632.673}{kernel}%
\TextAt{8}{440.56}{632.673}{of}%
\TextAt{8}{454.384}{632.673}{a}%
\TextAt{8}{464.633}{632.673}{canonical}%
\TextAt{8}{516.51}{632.673}{ho-}%
\TextAt{8}{79.226}{615.627}{momorphism}%
\TextAt{13}{262.853}{598.581}{I}%
\TextAt{21}{268.016}{596.788}{w}%
\NativeArrowAt{278.126}{598.581}{11.955}{19.26}{\longrightarrow}%
\TextAt{13}{300.708}{598.581}{\Character{31}}%
\TextAt{9}{308.042}{598.581}{\Character{40}}%
\TextAt{13}{312.594}{598.581}{L}%
\TextAt{16}{323.879}{598.581}{\Character{106}}%
\TextAt{13}{330.521}{598.581}{K}%
\TextAt{9}{341.339}{598.581}{\Character{41}}%
\TextAt{13}{345.891}{598.581}{\Character{59}}%
\TextAt{8}{79.226}{575.215}{where}%
\TextAt{13}{237.297}{558.169}{\Character{31}}%
\TextAt{9}{244.631}{558.169}{\Character{40}}%
\TextAt{13}{249.183}{558.169}{L}%
\TextAt{16}{260.468}{558.169}{\Character{106}}%
\TextAt{13}{267.11}{558.169}{K}%
\TextAt{9}{277.928}{558.169}{\Character{41}}%
\TextAt{9}{285.783}{558.169}{\Character{61}}%
\TextAt{8}{298.226}{558.169}{Hom}%
\TextAt{9}{322.615}{558.169}{\Character{40}\Character{1}\Character{47}\Character{0}}%
\TextAt{13}{350.089}{558.169}{\Character{59}}%
\TextAt{13}{355.313}{558.169}{\Character{21}}%
\TextAt{3}{362.162}{563.105}{\Character{3}}%
\TextAt{9}{366.894}{558.169}{\Character{41}}%
\TextAt{13}{371.447}{558.169}{\Character{59}}%
\TextAt{8}{79.226}{534.803}{where}%
\TextAt{9}{113.035}{534.803}{\Character{1}}%
\TextAt{9}{126.09}{534.803}{\Character{61}}%
\TextAt{13}{138.537}{534.803}{w}%
\TextAt{9}{147.226}{534.803}{\Character{40}}%
\TextAt{13}{151.779}{534.803}{L}%
\TextAt{3}{159.743}{539.141}{\Character{3}}%
\TextAt{9}{164.475}{534.803}{\Character{41}}%
\TextAt{8}{169.028}{534.803}{,}%
\TextAt{8}{176.177}{534.803}{and}%
\TextAt{9}{198.94}{534.803}{\Character{0}}%
\TextAt{9}{209.556}{534.803}{\Character{61}}%
\TextAt{13}{222}{534.803}{v}%
\TextAt{9}{228.088}{534.803}{\Character{40}}%
\TextAt{13}{232.641}{534.803}{K}%
\TextAt{3}{243.458}{539.141}{\Character{3}}%
\TextAt{9}{248.191}{534.803}{\Character{41}}%
\TextAt{8}{252.743}{534.803}{.}%
\TextAt{8}{261.191}{534.803}{If}%
\TextAt{13}{272.883}{534.803}{\Character{27}}%
\NativeMathAt{283.287}{534.803}{11.955}{\in}%
\TextAt{13}{294.578}{534.803}{I}%
\TextAt{21}{299.74}{533.009}{w}%
\TextAt{8}{306.53}{534.803}{,}%
\TextAt{8}{313.679}{534.803}{then}%
\TextAt{8}{340.339}{534.803}{the}%
\TextAt{8}{360.496}{534.803}{associated}%
\TextAt{8}{416.231}{534.803}{homomorphism}%
\TextAt{13}{264.38}{506.164}{\Character{31}}%
\TextAt{21}{271.713}{504.371}{\Character{27}}%
\TextAt{9}{280.658}{506.164}{\Character{58}}%
\TextAt{9}{287.209}{506.164}{\Character{1}\Character{47}\Character{0}}%
\NativeArrowAt{313.473}{506.164}{11.955}{19.26}{\longrightarrow}%
\TextAt{13}{336.055}{506.164}{\Character{21}}%
\TextAt{3}{342.884}{511.1}{\Character{3}}%
\TextAt{8}{79.226}{477.525}{is}%
\TextAt{8}{90.595}{477.525}{given}%
\TextAt{8}{120.769}{477.525}{as}%
\TextAt{8}{134.743}{477.525}{follows:}%
\TextAt{8}{177.532}{477.525}{for}%
\TextAt{9}{195.562}{480.68}{\Character{22}}%
\TextAt{13}{195.022}{477.525}{\Character{14}}%
\TextAt{9}{203.976}{477.525}{\Character{61}}%
\TextAt{13}{216.401}{477.525}{\Character{14}}%
\TextAt{8}{225.355}{477.525}{mod}%
\TextAt{9}{250.793}{477.525}{\Character{0}}%
\NativeMathAt{261.428}{477.525}{11.955}{\in}%
\TextAt{9}{272.719}{477.525}{\Character{1}\Character{47}\Character{0}}%
\TextAt{8}{295.641}{477.525}{,}%
\TextAt{8}{302.474}{477.525}{choose}%
\TextAt{8}{339.211}{477.525}{an}%
\TextAt{13}{355.075}{477.525}{x}%
\NativeMathAt{365.048}{477.525}{11.955}{\in}%
\TextAt{13}{376.339}{477.525}{L}%
\TextAt{3}{384.303}{481.863}{\Character{3}}%
\TextAt{8}{392.538}{477.525}{such}%
\TextAt{8}{418.539}{477.525}{that}%
\TextAt{13}{443.513}{477.525}{w}%
\TextAt{9}{452.202}{477.525}{\Character{40}}%
\TextAt{13}{456.754}{477.525}{x}%
\TextAt{9}{463.407}{477.525}{\Character{41}}%
\TextAt{9}{471.262}{477.525}{\Character{61}}%
\TextAt{13}{483.705}{477.525}{\Character{14}}%
\TextAt{8}{492.84}{477.525}{and}%
\TextAt{8}{515.207}{477.525}{put}%
\TextAt{13}{255.263}{448.886}{\Character{31}}%
\TextAt{21}{262.596}{447.093}{\Character{27}}%
\TextAt{9}{268.22}{448.886}{\Character{40}}%
\TextAt{9}{273.312}{452.041}{\Character{22}}%
\TextAt{13}{272.772}{448.886}{\Character{14}}%
\TextAt{9}{278.405}{448.886}{\Character{41}}%
\TextAt{9}{286.26}{448.886}{\Character{61}}%
\TextAt{13}{299.899}{456.974}{\Character{27}}%
\TextAt{13}{306.964}{456.974}{x}%
\Vector{q 0 G 0 g 0.478 w 227.899 -322.125 m 241.634 -322.125 l S Q}%
\TextAt{13}{303.44}{440.685}{x}%
\TextAt{8}{318.15}{448.886}{mod}%
\TextAt{14}{343.588}{448.886}{P}%
\TextAt{13}{353.481}{448.886}{\Character{58}}%
\TextAt{8}{79.226}{418.189}{This}%
\TextAt{8}{106.787}{418.189}{definition}%
\TextAt{8}{159.658}{418.189}{is}%
\TextAt{8}{172.263}{418.189}{independent}%
\TextAt{8}{239.432}{418.189}{of}%
\TextAt{8}{253.603}{418.189}{the}%
\TextAt{8}{274.601}{418.189}{choice}%
\TextAt{8}{310.232}{418.189}{of}%
\TextAt{8}{324.404}{418.189}{the}%
\TextAt{8}{345.402}{418.189}{representative}%
\TextAt{13}{421.429}{418.189}{\Character{14}}%
\NativeMathAt{431.817}{418.189}{11.955}{\in}%
\TextAt{9}{445.082}{421.344}{\Character{22}}%
\TextAt{13}{444.542}{418.189}{\Character{14}}%
\TextAt{8}{454.914}{418.189}{and}%
\TextAt{8}{478.518}{418.189}{of}%
\TextAt{13}{492.689}{418.189}{x}%
\NativeMathAt{504.096}{418.189}{11.955}{\in}%
\TextAt{13}{516.821}{418.189}{L}%
\TextAt{3}{524.785}{422.528}{\Character{3}}%
\TextAt{8}{529.518}{418.189}{.}%
\TextAt{8}{79.226}{401.143}{For}%
\TextAt{8}{101.156}{401.143}{if}%
\TextAt{13}{112.839}{401.143}{x}%
\TextAt{3}{119.491}{405.482}{\Character{48}}%
\NativeMathAt{127.243}{401.143}{11.955}{\in}%
\TextAt{13}{140.17}{401.143}{L}%
\TextAt{3}{148.135}{405.482}{\Character{3}}%
\TextAt{8}{157.725}{401.143}{is}%
\TextAt{8}{170.449}{401.143}{an}%
\TextAt{8}{187.668}{401.143}{element}%
\TextAt{8}{231.869}{401.143}{such}%
\TextAt{8}{259.226}{401.143}{that}%
\TextAt{13}{285.555}{401.143}{w}%
\TextAt{9}{294.245}{401.143}{\Character{40}}%
\TextAt{13}{298.797}{401.143}{x}%
\TextAt{3}{305.449}{405.482}{\Character{48}}%
\TextAt{9}{308.244}{401.143}{\Character{41}}%
\NativeMathAt{317.754}{401.143}{11.955}{\equiv}%
\TextAt{13}{332.009}{401.143}{w}%
\TextAt{9}{340.698}{401.143}{\Character{40}}%
\TextAt{13}{345.251}{401.143}{x}%
\TextAt{9}{351.903}{401.143}{\Character{41}}%
\TextAt{8}{359.776}{401.143}{mod}%
\TextAt{9}{385.214}{401.143}{\Character{0}}%
\TextAt{8}{392.528}{401.143}{,}%
\TextAt{8}{400.877}{401.143}{then}%
\TextAt{13}{428.497}{401.143}{w}%
\TextAt{9}{437.187}{401.143}{\Character{40}}%
\TextAt{13}{441.739}{401.143}{x}%
\TextAt{3}{448.391}{405.482}{\Character{48}}%
\TextAt{9}{451.186}{401.143}{\Character{41}}%
\TextAt{9}{460.678}{401.143}{\Character{61}}%
\TextAt{13}{474.757}{401.143}{w}%
\TextAt{9}{483.447}{401.143}{\Character{40}}%
\TextAt{13}{487.999}{401.143}{xa}%
\TextAt{9}{500.796}{401.143}{\Character{41}}%
\TextAt{8}{505.348}{401.143}{,}%
\TextAt{13}{513.698}{401.143}{a}%
\NativeMathAt{524.799}{401.143}{11.955}{\in}%
\TextAt{13}{79.226}{384.097}{K}%
\TextAt{3}{90.044}{388.436}{\Character{3}}%
\TextAt{8}{94.776}{384.097}{.}%
\TextAt{8}{106.323}{384.097}{Then}%
\TextAt{13}{137.913}{384.097}{x}%
\TextAt{3}{144.566}{388.436}{\Character{48}}%
\TextAt{9}{152.442}{384.097}{\Character{61}}%
\TextAt{13}{166.627}{384.097}{xau}%
\TextAt{8}{186.087}{384.097}{,}%
\TextAt{13}{194.527}{384.097}{u}%
\NativeMathAt{206.271}{384.097}{11.955}{\in}%
\TextAt{16}{219.302}{384.097}{\Character{79}}%
\TextAt{3}{229.172}{388.436}{\Character{3}}%
\TextAt{8}{233.904}{384.097}{,}%
\TextAt{8}{242.344}{384.097}{and}%
\TextAt{8}{266.14}{384.097}{since}%
\TextAt{13}{295.841}{384.097}{\Character{27}}%
\TextAt{13}{302.906}{384.097}{u}%
\TextAt{9}{309.586}{384.097}{\Character{47}}%
\TextAt{13}{315.439}{384.097}{u}%
\NativeMathAt{327.183}{384.097}{11.955}{\equiv}%
\TextAt{9}{341.562}{384.097}{1}%
\TextAt{8}{350.736}{384.097}{mod}%
\TextAt{14}{376.174}{384.097}{P}%
\TextAt{8}{390.997}{384.097}{(because}%
\TextAt{13}{439.887}{384.097}{\Character{27}}%
\NativeMathAt{452.05}{384.097}{11.955}{\in}%
\TextAt{13}{465.101}{384.097}{I}%
\TextAt{21}{470.264}{382.304}{w}%
\TextAt{8}{477.054}{384.097}{),}%
\TextAt{8}{490.049}{384.097}{one}%
\TextAt{8}{512.541}{384.097}{gets}%
\TextAt{13}{79.226}{367.051}{\Character{27}}%
\TextAt{13}{86.292}{367.051}{x}%
\TextAt{3}{92.961}{371.39}{\Character{48}}%
\TextAt{9}{95.756}{367.051}{\Character{47}}%
\TextAt{13}{101.609}{367.051}{x}%
\TextAt{3}{108.261}{371.39}{\Character{48}}%
\NativeMathAt{114.492}{367.051}{11.955}{\equiv}%
\TextAt{13}{127.227}{367.051}{\Character{27}}%
\TextAt{13}{134.293}{367.051}{x}%
\TextAt{9}{140.961}{367.051}{\Character{47}}%
\TextAt{13}{146.814}{367.051}{x}%
\TextAt{8}{156.787}{367.051}{mod}%
\TextAt{14}{182.225}{367.051}{P}%
\TextAt{8}{192.118}{367.051}{.}%
\TextAt{8}{200.77}{367.051}{One}%
\TextAt{8}{225.549}{367.051}{sees}%
\TextAt{8}{249.144}{367.051}{immediately}%
\TextAt{8}{315.874}{367.051}{that}%
\TextAt{8}{341.311}{367.051}{mapping}%
\TextAt{13}{389.51}{367.051}{\Character{27}}%
\NativeArrowAt{400.029}{367.051}{11.955}{11.955}{\mapsto}%
\TextAt{13}{415.42}{367.051}{\Character{31}}%
\TextAt{21}{422.754}{365.258}{\Character{27}}%
\TextAt{8}{432.342}{367.051}{is}%
\TextAt{8}{444.174}{367.051}{a}%
\TextAt{8}{453.997}{367.051}{homomorphism}%
\TextAt{13}{79.226}{350.005}{I}%
\TextAt{21}{84.389}{348.212}{w}%
\NativeMathAt{94.499}{350.005}{11.955}{\rightarrow}%
\TextAt{13}{109.775}{350.005}{\Character{31}}%
\TextAt{9}{117.109}{350.005}{\Character{40}}%
\TextAt{13}{121.661}{350.005}{L}%
\TextAt{16}{132.947}{350.005}{\Character{106}}%
\TextAt{13}{139.588}{350.005}{K}%
\TextAt{9}{150.406}{350.005}{\Character{41}}%
\TextAt{8}{158.856}{350.005}{with}%
\TextAt{8}{185.516}{350.005}{kernel}%
\TextAt{13}{219.982}{350.005}{R}%
\TextAt{21}{228.893}{348.212}{w}%
\TextAt{8}{235.682}{350.005}{.}%
\ErrataAt{79.226}{337.05}{455.774}{In \hyperlink{nav.number.2.9.12}{(9.12)}, \hyperlink{nav.number.2.9.14}{(9.14)}, and \hyperlink{nav.number.2.9.15}{(9.15)}, assume $p=\operatorname{char}\lambda>0$.}%
\TextAt{12}{79.226}{324.096}{Proposition}%
\TextAt{12}{151.937}{324.096}{(9.12).}%
\TextAt{13}{193.984}{324.096}{R}%
\TextAt{21}{202.894}{322.303}{w}%
\TextAt{8}{213.581}{324.096}{is}%
\TextAt{8}{225.345}{324.096}{the}%
\TextAt{8}{245.501}{324.096}{unique}%
\TextAt{13}{283.542}{324.096}{p}%
\TextAt{8}{289.418}{324.096}{-Sylow}%
\TextAt{8}{327.136}{324.096}{subgroup}%
\TextAt{8}{377.934}{324.096}{of}%
\TextAt{13}{391.264}{324.096}{I}%
\TextAt{21}{396.427}{322.303}{w}%
\TextAt{8}{403.216}{324.096}{.}%
\TextAt{12}{79.226}{295.251}{Remark.}%
\TextAt{8}{133.729}{295.251}{If}%
\TextAt{13}{145.767}{295.251}{L}%
\TextAt{16}{157.642}{295.25}{\Character{106}}%
\TextAt{13}{164.874}{295.251}{K}%
\TextAt{8}{179.935}{295.251}{is}%
\TextAt{8}{192.045}{295.251}{a}%
\TextAt{8}{202.147}{295.251}{finite}%
\TextAt{8}{232.404}{295.251}{extension,}%
\TextAt{8}{287.855}{295.251}{then}%
\TextAt{8}{314.861}{295.251}{it}%
\TextAt{8}{326.911}{295.251}{is}%
\TextAt{8}{339.021}{295.251}{clear}%
\TextAt{8}{367.33}{295.251}{what}%
\TextAt{8}{396.943}{295.251}{this}%
\TextAt{8}{420.111}{295.251}{means.}%
\TextAt{8}{461.53}{295.251}{In}%
\TextAt{8}{476.497}{295.251}{the}%
\TextAt{8}{497}{295.251}{infinite}%
\TextAt{8}{79.226}{278.205}{case}%
\TextAt{8}{103.444}{278.205}{it}%
\TextAt{8}{114.595}{278.205}{has}%
\TextAt{8}{134.915}{278.205}{to}%
\TextAt{8}{148.672}{278.205}{be}%
\TextAt{8}{164.043}{278.205}{understood}%
\TextAt{8}{224.688}{278.205}{in}%
\TextAt{8}{237.788}{278.205}{the}%
\TextAt{8}{257.391}{278.205}{sense}%
\TextAt{8}{286.869}{278.205}{of}%
\TextAt{8}{299.646}{278.205}{profinite}%
\TextAt{8}{345.921}{278.205}{groups,}%
\TextAt{8}{386.52}{278.205}{i.e.,}%
\TextAt{8}{408.183}{278.205}{all}%
\TextAt{8}{423.889}{278.205}{finite}%
\TextAt{8}{453.247}{278.205}{quotient}%
\TextAt{8}{498.877}{278.205}{groups}%
\TextAt{8}{79.226}{261.159}{of}%
\TextAt{13}{92.105}{261.159}{R}%
\TextAt{21}{101.016}{259.365}{w}%
\TextAt{8}{107.805}{261.159}{,}%
\TextAt{8}{114.594}{261.159}{resp.}%
\TextAt{13}{142.166}{261.159}{I}%
\TextAt{21}{147.329}{259.365}{w}%
\TextAt{9}{154.118}{261.159}{\Character{47}}%
\TextAt{13}{159.971}{261.159}{R}%
\TextAt{21}{168.881}{259.365}{w}%
\TextAt{8}{175.671}{261.159}{,}%
\TextAt{8}{182.459}{261.159}{by}%
\TextAt{8}{198.268}{261.159}{closed}%
\TextAt{8}{232.344}{261.159}{normal}%
\TextAt{8}{271.572}{261.159}{subgroups}%
\TextAt{8}{326.534}{261.159}{have}%
\TextAt{13}{353.078}{261.159}{p}%
\TextAt{8}{358.953}{261.159}{-power}%
\TextAt{8}{396.543}{261.159}{order,}%
\TextAt{8}{430.004}{261.159}{resp.}%
\TextAt{8}{457.576}{261.159}{an}%
\TextAt{8}{473.385}{261.159}{order}%
\TextAt{8}{503.503}{261.159}{prime}%
\TextAt{8}{79.226}{244.113}{to}%
\TextAt{13}{93.226}{244.113}{p}%
\TextAt{8}{99.101}{244.113}{.}%
\TextAt{8}{107.446}{244.113}{In}%
\TextAt{8}{121.756}{244.113}{order}%
\TextAt{8}{152.015}{244.113}{to}%
\TextAt{8}{166.014}{244.113}{understand}%
\TextAt{8}{226.902}{244.113}{this}%
\TextAt{8}{249.414}{244.113}{better,}%
\TextAt{8}{287.207}{244.113}{we}%
\TextAt{8}{304.123}{244.113}{refer}%
\TextAt{8}{330.796}{244.113}{the}%
\TextAt{8}{350.641}{244.113}{reader}%
\TextAt{8}{386.1}{244.113}{to}%
\TextAt{8}{400.1}{244.113}{chap.}%
\TextAt{8}{430.682}{244.113}{IV,}%
\TextAt{8}{450.515}{244.113}{§2,}%
\TextAt{8}{468.953}{244.113}{exercise}%
\TextAt{8}{511.944}{244.113}{3–5.}%
\TextAt{19}{79.226}{212.278}{Proof}%
\TextAt{19}{110.62}{212.278}{of}%
\TextAt{19}{124.369}{212.278}{(9.12).}%
\TextAt{8}{165.016}{212.278}{By}%
\TextAt{8}{184.171}{212.278}{(9.6),}%
\TextAt{8}{216.387}{212.278}{we}%
\TextAt{8}{234.406}{212.278}{may}%
\TextAt{8}{260.566}{212.278}{assume}%
\TextAt{8}{301.802}{212.278}{that}%
\TextAt{13}{327.962}{212.278}{K}%
\TextAt{8}{343.469}{212.278}{is}%
\TextAt{8}{356.024}{212.278}{henselian.}%
\TextAt{8}{413.734}{212.278}{We}%
\TextAt{8}{434.67}{212.278}{restrict}%
\TextAt{8}{475.846}{212.278}{to}%
\TextAt{8}{490.948}{212.278}{the}%
\TextAt{8}{511.896}{212.278}{case}%
\TextAt{8}{79.226}{195.232}{where}%
\TextAt{13}{113.856}{195.232}{L}%
\TextAt{16}{126.54}{195.232}{\Character{106}}%
\TextAt{13}{134.581}{195.232}{K}%
\TextAt{8}{150.117}{195.232}{is}%
\TextAt{8}{162.702}{195.232}{a}%
\TextAt{8}{173.278}{195.232}{finite}%
\TextAt{8}{204.011}{195.232}{extension.}%
\TextAt{8}{262.79}{195.232}{The}%
\TextAt{8}{287.665}{195.232}{infinite}%
\TextAt{8}{328.153}{195.232}{case}%
\TextAt{8}{353.745}{195.232}{of}%
\TextAt{8}{367.896}{195.232}{the}%
\TextAt{8}{388.873}{195.232}{proposition}%
\TextAt{8}{451.227}{195.232}{follows}%
\TextAt{8}{490.484}{195.232}{formally}%
\TextAt{8}{79.226}{178.186}{from}%
\TextAt{8}{106.866}{178.186}{this.}%
\TextAt{8}{97.159}{158.151}{If}%
\TextAt{13}{108.312}{158.151}{R}%
\TextAt{21}{117.223}{156.358}{w}%
\TextAt{8}{127.371}{158.151}{were}%
\TextAt{8}{153.815}{158.151}{not}%
\TextAt{8}{174.09}{158.151}{a}%
\TextAt{13}{183.307}{158.151}{p}%
\TextAt{8}{189.182}{158.151}{-group,}%
\TextAt{8}{229.076}{158.151}{then}%
\TextAt{8}{255.197}{158.151}{we}%
\TextAt{8}{271.886}{158.151}{would}%
\TextAt{8}{305.491}{158.151}{find}%
\TextAt{8}{328.361}{158.151}{an}%
\TextAt{8}{344.081}{158.151}{element}%
\TextAt{13}{386.785}{158.151}{\Character{27}}%
\NativeMathAt{397.188}{158.151}{11.955}{\in}%
\TextAt{13}{408.479}{158.151}{R}%
\TextAt{21}{417.389}{156.358}{w}%
\TextAt{8}{427.537}{158.151}{of}%
\TextAt{8}{440.328}{158.151}{prime}%
\TextAt{8}{472.953}{158.151}{order}%
\TextAt{13}{502.984}{158.151}{\Character{96}}%
\NativeMathAt{511.217}{158.151}{11.955}{\neq}%
\TextAt{13}{523.643}{158.151}{p}%
\TextAt{8}{529.518}{158.151}{.}%
\TextAt{8}{79.226}{141.105}{Let}%
\TextAt{13}{100.578}{141.105}{K}%
\TextAt{3}{111.395}{145.444}{\Character{48}}%
\TextAt{8}{118.47}{141.105}{be}%
\TextAt{8}{134.777}{141.105}{the}%
\TextAt{8}{155.315}{141.105}{fixed}%
\TextAt{8}{183.984}{141.105}{field}%
\TextAt{8}{209.723}{141.105}{of}%
\TextAt{13}{223.435}{141.105}{\Character{27}}%
\TextAt{8}{234.797}{141.105}{and}%
\TextAt{13}{257.942}{141.105}{\Character{20}}%
\TextAt{3}{264.681}{145.444}{\Character{48}}%
\TextAt{8}{271.755}{141.105}{its}%
\TextAt{8}{288.456}{141.105}{residue}%
\TextAt{8}{328.566}{141.105}{class}%
\TextAt{8}{356.385}{141.105}{field.}%
\TextAt{8}{387.44}{141.105}{We}%
\TextAt{8}{407.966}{141.105}{show}%
\TextAt{8}{437.352}{141.105}{that}%
\TextAt{13}{463.103}{141.105}{\Character{20}}%
\TextAt{3}{469.842}{145.444}{\Character{48}}%
\TextAt{9}{476.609}{141.105}{\Character{61}}%
\TextAt{13}{489.686}{141.105}{\Character{21}}%
\TextAt{8}{496.514}{141.105}{.}%
\TextAt{8}{506.109}{141.105}{Since}%
\TextAt{13}{79.226}{124.059}{R}%
\TextAt{21}{88.136}{122.266}{w}%
\NativeMathAt{99.143}{124.059}{11.955}{\subseteq}%
\TextAt{13}{112.659}{124.059}{I}%
\TextAt{21}{117.822}{122.266}{w}%
\TextAt{8}{124.611}{124.059}{,}%
\TextAt{8}{132.418}{124.059}{we}%
\TextAt{8}{150.172}{124.059}{have}%
\TextAt{8}{177.693}{124.059}{that}%
\TextAt{13}{203.588}{124.059}{T}%
\NativeMathAt{216.292}{124.059}{11.955}{\subseteq}%
\TextAt{13}{229.808}{124.059}{K}%
\TextAt{3}{240.625}{128.398}{\Character{48}}%
\TextAt{8}{243.42}{124.059}{.}%
\TextAt{8}{253.447}{124.059}{Thus}%
\TextAt{13}{283.622}{124.059}{\Character{21}}%
\TextAt{21}{290.451}{122.266}{s}%
\NativeMathAt{299.082}{124.059}{11.955}{\subseteq}%
\TextAt{13}{312.598}{124.059}{\Character{20}}%
\TextAt{3}{319.336}{128.398}{\Character{48}}%
\TextAt{8}{322.132}{124.059}{,}%
\TextAt{8}{329.938}{124.059}{so}%
\TextAt{8}{344.835}{124.059}{that}%
\TextAt{13}{370.73}{124.059}{\Character{21}}%
\TextAt{16}{381.775}{124.059}{\Character{106}}%
\TextAt{13}{389.314}{124.059}{\Character{20}}%
\TextAt{3}{396.052}{128.398}{\Character{48}}%
\TextAt{8}{403.271}{124.059}{is}%
\TextAt{8}{415.561}{124.059}{purely}%
\TextAt{8}{452.18}{124.059}{inseparable}%
\TextAt{8}{513.904}{124.059}{and}%
\TextAt{8}{79.226}{107.013}{of}%
\TextAt{13}{93.31}{107.013}{p}%
\TextAt{8}{99.186}{107.013}{-power}%
\TextAt{8}{137.981}{107.013}{degree.}%
\TextAt{8}{181.211}{107.013}{On}%
\TextAt{8}{201.476}{107.013}{the}%
\TextAt{8}{222.386}{107.013}{other}%
\TextAt{8}{253.71}{107.013}{hand,}%
\TextAt{8}{287.171}{107.013}{the}%
\TextAt{8}{308.082}{107.013}{degree}%
\TextAt{8}{345.251}{107.013}{has}%
\TextAt{8}{366.88}{107.013}{to}%
\TextAt{8}{381.944}{107.013}{be}%
\TextAt{8}{398.623}{107.013}{a}%
\TextAt{8}{409.133}{107.013}{power}%
\TextAt{8}{444.031}{107.013}{of}%
\TextAt{13}{458.115}{107.013}{\Character{96}}%
\TextAt{8}{463.027}{107.013}{,}%
\TextAt{8}{471.119}{107.013}{for}%
\TextAt{8}{489.759}{107.013}{if}%
\TextAt{9}{502.396}{107.013}{\Character{22}}%
\TextAt{13}{501.237}{107.013}{\Character{11}}%
\NativeMathAt{513.365}{107.013}{11.955}{\in}%
\TextAt{13}{525.941}{107.013}{\Character{21}}%
\TextAt{8}{79.226}{89.967}{and}%
\TextAt{8}{102.111}{89.967}{if}%
\TextAt{13}{112.958}{89.967}{\Character{11}}%
\NativeMathAt{124.01}{89.967}{11.955}{\in}%
\TextAt{13}{135.51}{89.967}{L}%
\TextAt{8}{147.494}{89.967}{is}%
\TextAt{8}{159.381}{89.967}{a}%
\TextAt{8}{169.259}{89.967}{lifting}%
\TextAt{8}{203.526}{89.967}{of}%
\TextAt{9}{218.138}{89.967}{\Character{22}}%
\TextAt{13}{216.978}{89.967}{\Character{11}}%
\TextAt{8}{224.5}{89.967}{,}%
\TextAt{8}{231.802}{89.967}{and}%
\TextAt{13}{254.688}{89.967}{f}%
\TextAt{9}{261.734}{89.967}{\Character{40}}%
\TextAt{13}{266.287}{89.967}{X}%
\TextAt{9}{276.942}{89.967}{\Character{41}}%
\NativeMathAt{285.024}{89.967}{11.955}{\in}%
\TextAt{13}{296.524}{89.967}{K}%
\TextAt{3}{307.342}{94.306}{\Character{48}}%
\TextAt{9}{310.137}{89.967}{\Character{91}}%
\TextAt{13}{313.388}{89.967}{x}%
\TextAt{9}{320.04}{89.967}{\Character{93}}%
\TextAt{8}{327.312}{89.967}{is}%
\TextAt{8}{339.199}{89.967}{the}%
\TextAt{8}{359.478}{89.967}{minimal}%
\TextAt{8}{405.126}{89.967}{polynomial}%
\TextAt{8}{465.742}{89.967}{of}%
\TextAt{13}{479.194}{89.967}{\Character{11}}%
\TextAt{8}{490.736}{89.967}{over}%
\TextAt{13}{515.905}{89.967}{K}%
\TextAt{3}{526.723}{94.306}{\Character{48}}%
\TextAt{8}{529.518}{89.967}{,}%
\TextAt{8}{79.226}{72.921}{then}%
\TextAt{9}{108.412}{76.076}{\Character{22}}%
\TextAt{13}{105.865}{72.921}{f}%
\TextAt{9}{112.911}{72.921}{\Character{40}}%
\TextAt{13}{117.463}{72.921}{x}%
\TextAt{9}{124.115}{72.921}{\Character{41}}%
\TextAt{9}{131.97}{72.921}{\Character{61}}%
\TextAt{9}{144.81}{72.921}{\Character{22}}%
\TextAt{13}{144.414}{72.921}{g}%
\TextAt{9}{150.448}{72.921}{\Character{40}}%
\TextAt{13}{155.001}{72.921}{x}%
\TextAt{9}{161.653}{72.921}{\Character{41}}%
\TextAt{21}{166.205}{77.26}{m}%
\TextAt{8}{174.194}{72.921}{,}%
\TextAt{8}{181.326}{72.921}{where}%
\TextAt{9}{215.529}{72.921}{\Character{22}}%
\TextAt{13}{215.113}{72.921}{g}%
\TextAt{9}{221.147}{72.921}{\Character{40}}%
\TextAt{13}{225.7}{72.921}{x}%
\TextAt{9}{232.352}{72.921}{\Character{41}}%
\NativeMathAt{240.225}{72.921}{11.955}{\in}%
\TextAt{13}{251.516}{72.921}{\Character{20}}%
\TextAt{3}{258.254}{77.26}{\Character{48}}%
\TextAt{9}{261.05}{72.921}{\Character{91}}%
\TextAt{13}{264.301}{72.921}{x}%
\TextAt{9}{270.953}{72.921}{\Character{93}}%
\TextAt{8}{278.081}{72.921}{is}%
\TextAt{8}{289.823}{72.921}{the}%
\TextAt{8}{309.958}{72.921}{minimal}%
\TextAt{8}{355.461}{72.921}{polynomial}%
\TextAt{8}{415.933}{72.921}{of}%
\TextAt{9}{430.401}{72.921}{\Character{22}}%
\TextAt{13}{429.241}{72.921}{\Character{11}}%
\TextAt{8}{440.638}{72.921}{over}%
\TextAt{13}{465.663}{72.921}{\Character{20}}%
\TextAt{3}{472.401}{77.26}{\Character{48}}%
\TextAt{8}{475.197}{72.921}{,}%
\TextAt{8}{482.328}{72.921}{which}%
\TextAt{8}{515.793}{72.921}{has}%
\TextAt{8}{297.211}{43.033}{159}%
\Destination{78.23}{809.815}{nav.page.159}%
\Destination{79.36}{549.755}{nav.print.174}%
\Destination{79.36}{338.055}{nav.number.2.9.12}%
\Hotlink{511.444}{241.955}{6.858}{11.047}{nav.ex.4.2.3}%
\Hotlink{430.182}{241.955}{13.995}{11.047}{nav.chapter.4}%
\Hotlink{455.694}{241.955}{6.858}{11.047}{nav.section.4.2}%
\Hotlink{123.869}{210.055}{32.083}{11.047}{nav.number.2.9.12}%
\Hotlink{183.671}{210.055}{25.078}{11.047}{nav.number.2.9.6}%
\Hotlink{525.1}{241.95}{8.1}{11.047}{nav.ex.4.2.5}%
\Destination{79.226}{340}{errata.215}%
}

% Page 172 (body page 160)
\Page{%
\TextAt{2}{79.226}{796.943}{II.}%
\TextAt{2}{94.073}{796.943}{The}%
\TextAt{2}{116.492}{796.943}{Theory}%
\TextAt{2}{154.401}{796.943}{of}%
\TextAt{2}{166.826}{796.943}{Valuations}%
\Vector{q 0 G 0 g 0.299 w 7.226 18.729 m 460.769 18.729 l S Q}%
\TextAt{8}{79.226}{761.124}{degree}%
\TextAt{8}{115.717}{761.124}{either}%
\TextAt{8}{148.957}{761.124}{1}%
\TextAt{8}{158.788}{761.124}{or}%
\TextAt{13}{173.174}{761.124}{\Character{96}}%
\TextAt{8}{178.086}{761.124}{,}%
\TextAt{8}{185.329}{761.124}{as}%
\TextAt{8}{199.775}{761.124}{this}%
\TextAt{8}{222.673}{761.124}{is}%
\TextAt{8}{234.513}{761.124}{so}%
\TextAt{8}{248.958}{761.124}{for}%
\TextAt{13}{266.919}{761.124}{f}%
\TextAt{9}{273.966}{761.124}{\Character{40}}%
\TextAt{13}{278.518}{761.124}{x}%
\TextAt{9}{285.17}{761.124}{\Character{41}}%
\TextAt{8}{289.722}{761.124}{.}%
\TextAt{8}{298.397}{761.124}{Thus}%
\TextAt{8}{328.122}{761.124}{we}%
\TextAt{8}{345.425}{761.124}{have}%
\TextAt{8}{372.495}{761.124}{indeed}%
\TextAt{13}{409.632}{761.124}{\Character{20}}%
\TextAt{3}{416.371}{765.462}{\Character{48}}%
\TextAt{9}{422.615}{761.124}{\Character{61}}%
\TextAt{13}{435.17}{761.124}{\Character{21}}%
\TextAt{8}{441.998}{761.124}{,}%
\TextAt{8}{449.242}{761.124}{so}%
\TextAt{8}{463.688}{761.124}{that}%
\TextAt{13}{489.132}{761.124}{L}%
\TextAt{16}{500.547}{761.124}{\Character{106}}%
\TextAt{13}{507.317}{761.124}{K}%
\TextAt{3}{518.135}{765.462}{\Character{48}}%
\TextAt{8}{524.903}{761.124}{is}%
\TextAt{8}{79.226}{744.078}{tamely}%
\TextAt{8}{118.25}{744.078}{ramified}%
\TextAt{8}{164.099}{744.078}{and}%
\TextAt{8}{187.187}{744.078}{by}%
\TextAt{8}{203.77}{744.078}{(7.7)}%
\TextAt{8}{232.07}{744.078}{is}%
\TextAt{8}{244.158}{744.078}{of}%
\TextAt{8}{257.813}{744.078}{the}%
\TextAt{8}{278.294}{744.078}{form}%
\TextAt{13}{306.259}{744.078}{L}%
\TextAt{9}{318.097}{744.078}{\Character{61}}%
\TextAt{13}{331.076}{744.078}{K}%
\TextAt{3}{341.893}{748.416}{\Character{48}}%
\TextAt{9}{344.688}{744.078}{\Character{40}}%
\TextAt{13}{349.241}{744.078}{\Character{11}}%
\TextAt{9}{356.762}{744.078}{\Character{41}}%
\TextAt{8}{365.537}{744.078}{with}%
\TextAt{13}{392.521}{744.078}{\Character{11}}%
\TextAt{9}{403.917}{744.078}{\Character{61}}%
\TextAt{22}{420.216}{748.099}{\Character{96}}%
\TextAt{16}{416.674}{752.689}{\Character{112}}%
\Vector{q 0 G 0 g 0.478 w 354.637 -21.072 m 360.782 -21.072 l S Q}%
\TextAt{13}{426.637}{744.078}{a}%
\TextAt{8}{432.782}{744.078}{,}%
\TextAt{13}{440.336}{744.078}{a}%
\NativeMathAt{450.355}{744.078}{11.955}{\in}%
\TextAt{13}{462.199}{744.078}{K}%
\TextAt{3}{473.017}{748.416}{\Character{48}}%
\TextAt{8}{475.812}{744.078}{.}%
\TextAt{8}{485.234}{744.078}{It}%
\TextAt{8}{498.231}{744.078}{follows}%
\TextAt{8}{79.226}{727.032}{that}%
\TextAt{13}{104.718}{727.032}{\Character{27}}%
\TextAt{13}{111.784}{727.032}{\Character{11}}%
\TextAt{9}{122.852}{727.032}{\Character{61}}%
\TextAt{13}{135.486}{727.032}{\Character{16}}%
\TextAt{13}{141.488}{727.032}{\Character{11}}%
\TextAt{8}{149.024}{727.032}{,}%
\TextAt{8}{156.326}{727.032}{with}%
\TextAt{8}{183.109}{727.032}{a}%
\TextAt{8}{192.987}{727.032}{primitive}%
\TextAt{13}{243.19}{727.032}{\Character{96}}%
\TextAt{8}{248.102}{727.032}{-th}%
\TextAt{8}{267.078}{727.032}{root}%
\TextAt{8}{292.247}{727.032}{of}%
\TextAt{8}{305.7}{727.032}{unity}%
\TextAt{13}{336.392}{727.032}{\Character{16}}%
\NativeMathAt{345.937}{727.032}{11.955}{\in}%
\TextAt{13}{357.437}{727.032}{K}%
\TextAt{3}{368.255}{731.37}{\Character{48}}%
\TextAt{8}{371.05}{727.032}{.}%
\TextAt{8}{379.867}{727.032}{Since}%
\TextAt{13}{410.547}{727.032}{\Character{27}}%
\NativeMathAt{421.159}{727.032}{11.955}{\in}%
\TextAt{13}{432.659}{727.032}{R}%
\TextAt{21}{441.569}{725.238}{w}%
\TextAt{8}{448.359}{727.032}{,}%
\TextAt{8}{455.661}{727.032}{we}%
\TextAt{8}{473.011}{727.032}{have}%
\TextAt{8}{500.129}{727.032}{on}%
\TextAt{8}{516.51}{727.032}{the}%
\TextAt{8}{79.226}{709.986}{other}%
\TextAt{8}{109.796}{709.986}{hand}%
\TextAt{13}{139.062}{709.986}{\Character{27}}%
\TextAt{13}{146.128}{709.986}{\Character{11}}%
\TextAt{9}{153.666}{709.986}{\Character{47}}%
\TextAt{13}{159.519}{709.986}{\Character{11}}%
\TextAt{9}{170.361}{709.986}{\Character{61}}%
\TextAt{13}{182.787}{709.986}{\Character{16}}%
\NativeMathAt{192.123}{709.986}{11.955}{\equiv}%
\TextAt{9}{204.743}{709.986}{1}%
\TextAt{8}{213.916}{709.986}{mod}%
\TextAt{14}{239.354}{709.986}{P}%
\TextAt{8}{249.248}{709.986}{,}%
\TextAt{8}{256.397}{709.986}{a}%
\TextAt{8}{266.152}{709.986}{contradiction.}%
\TextAt{8}{341.932}{709.986}{This}%
\TextAt{8}{368.652}{709.986}{proves}%
\TextAt{8}{404.816}{709.986}{that}%
\TextAt{13}{430.185}{709.986}{R}%
\TextAt{21}{439.095}{708.193}{w}%
\TextAt{8}{449.782}{709.986}{is}%
\TextAt{8}{461.546}{709.986}{a}%
\TextAt{13}{471.301}{709.986}{p}%
\TextAt{8}{477.177}{709.986}{-group.}%
\TextAt{8}{97.159}{689.951}{Since}%
\TextAt{13}{128.813}{689.951}{p}%
\TextAt{9}{139.877}{689.951}{\Character{61}}%
\TextAt{8}{154.171}{689.951}{char}%
\TextAt{9}{176.288}{689.951}{\Character{40}}%
\TextAt{13}{180.84}{689.951}{\Character{21}}%
\TextAt{9}{187.669}{689.951}{\Character{41}}%
\TextAt{8}{192.221}{689.951}{,}%
\TextAt{8}{200.741}{689.951}{the}%
\TextAt{8}{221.993}{689.951}{elements}%
\TextAt{8}{270.946}{689.951}{in}%
\TextAt{13}{285.695}{689.951}{\Character{21}}%
\TextAt{3}{292.524}{694.289}{\Character{3}}%
\TextAt{8}{302.25}{689.951}{have}%
\TextAt{8}{330.341}{689.951}{order}%
\TextAt{8}{362.007}{689.951}{prime}%
\TextAt{8}{396.267}{689.951}{to}%
\TextAt{13}{411.673}{689.951}{p}%
\TextAt{8}{417.549}{689.951}{,}%
\TextAt{8}{426.068}{689.951}{provided}%
\TextAt{8}{475.296}{689.951}{they}%
\TextAt{8}{502.73}{689.951}{are}%
\TextAt{8}{523.337}{689.951}{of}%
\TextAt{8}{79.226}{672.905}{finite}%
\TextAt{8}{109.508}{672.905}{order.}%
\TextAt{8}{145.738}{672.905}{The}%
\TextAt{8}{170.162}{672.905}{group}%
\TextAt{13}{203.708}{672.905}{\Character{31}}%
\TextAt{9}{211.041}{672.905}{\Character{40}}%
\TextAt{13}{215.593}{672.905}{L}%
\TextAt{16}{227.509}{672.905}{\Character{106}}%
\TextAt{13}{234.782}{672.905}{K}%
\TextAt{9}{245.599}{672.905}{\Character{41}}%
\TextAt{9}{254.087}{672.905}{\Character{61}}%
\TextAt{8}{267.159}{672.905}{Hom}%
\TextAt{9}{291.547}{672.905}{\Character{40}\Character{1}\Character{47}\Character{0}}%
\TextAt{13}{319.022}{672.905}{\Character{59}}%
\TextAt{13}{324.246}{672.905}{\Character{21}}%
\TextAt{3}{331.095}{677.243}{\Character{3}}%
\TextAt{9}{335.827}{672.905}{\Character{41}}%
\TextAt{8}{344.647}{672.905}{therefore}%
\TextAt{8}{394.117}{672.905}{has}%
\TextAt{8}{415.36}{672.905}{order}%
\TextAt{8}{446.3}{672.905}{prime}%
\TextAt{8}{479.833}{672.905}{to}%
\TextAt{13}{494.514}{672.905}{p}%
\TextAt{8}{500.389}{672.905}{.}%
\TextAt{8}{509.947}{672.905}{This}%
\TextAt{8}{79.226}{655.859}{also}%
\TextAt{8}{103.479}{655.859}{applies}%
\TextAt{8}{143.333}{655.859}{to}%
\TextAt{8}{158.416}{655.859}{the}%
\TextAt{8}{179.346}{655.859}{group}%
\TextAt{13}{213.294}{655.859}{I}%
\TextAt{21}{218.457}{654.066}{w}%
\TextAt{9}{225.246}{655.859}{\Character{47}}%
\TextAt{13}{231.099}{655.859}{R}%
\TextAt{21}{240.009}{654.066}{w}%
\NativeMathAt{251.437}{655.859}{11.955}{\subseteq}%
\TextAt{13}{265.373}{655.859}{\Character{31}}%
\TextAt{9}{272.706}{655.859}{\Character{40}}%
\TextAt{13}{277.259}{655.859}{L}%
\TextAt{16}{289.861}{655.859}{\Character{106}}%
\TextAt{13}{297.82}{655.859}{K}%
\TextAt{9}{308.637}{655.859}{\Character{41}}%
\TextAt{8}{313.19}{655.859}{,}%
\TextAt{8}{321.305}{655.859}{so}%
\TextAt{8}{336.448}{655.859}{that}%
\TextAt{13}{362.59}{655.859}{R}%
\TextAt{21}{371.5}{654.066}{w}%
\TextAt{8}{382.959}{655.859}{is}%
\TextAt{8}{395.496}{655.859}{indeed}%
\TextAt{8}{433.33}{655.859}{the}%
\TextAt{8}{454.259}{655.859}{unique}%
\TextAt{13}{493.073}{655.859}{p}%
\TextAt{8}{498.948}{655.859}{-Sylow}%
\TextAt{8}{79.226}{638.813}{subgroup.}%
\Vector{q 0 G 0 g 0.399 w 452.334 -135.187 m 452.334 -127.117 l S Q}%
\Vector{q 0 G 0 g 0.399 w 452.534 -127.316 m 459.707 -127.316 l S Q}%
\Vector{q 0 G 0 g 0.399 w 452.534 -134.988 m 459.707 -134.988 l S Q}%
\Vector{q 0 G 0 g 0.399 w 459.906 -135.187 m 459.906 -127.117 l S Q}%
\TextAt{12}{79.226}{605.381}{Definition}%
\TextAt{12}{142.015}{605.381}{(9.13).}%
\TextAt{8}{184.061}{605.381}{The}%
\TextAt{8}{208.115}{605.381}{fixed}%
\TextAt{8}{236.401}{605.381}{field}%
\TextAt{8}{261.758}{605.381}{of}%
\TextAt{13}{275.088}{605.381}{R}%
\TextAt{21}{283.998}{603.588}{w}%
\TextAt{8}{290.787}{605.381}{,}%
\TextAt{13}{174.73}{579.088}{V}%
\TextAt{21}{181.555}{577.294}{w}%
\TextAt{9}{191.665}{579.088}{\Character{61}}%
\TextAt{13}{204.09}{579.088}{V}%
\TextAt{21}{210.915}{577.294}{w}%
\TextAt{9}{217.704}{579.088}{\Character{40}}%
\TextAt{13}{222.256}{579.088}{L}%
\TextAt{16}{233.542}{579.088}{\Character{106}}%
\TextAt{13}{240.183}{579.088}{K}%
\TextAt{9}{251.001}{579.088}{\Character{41}}%
\TextAt{9}{258.856}{579.088}{\Character{61}}%
\TextAt{16}{271.3}{579.088}{\Character{102}}%
\TextAt{13}{277.277}{579.088}{x}%
\NativeMathAt{287.25}{579.088}{11.955}{\in}%
\TextAt{13}{298.541}{579.088}{L}%
\TextAt{16}{309.827}{579.088}{\Character{106}}%
\TextAt{13}{316.468}{579.088}{\Character{27}}%
\TextAt{13}{323.534}{579.088}{x}%
\TextAt{9}{333.524}{579.088}{\Character{61}}%
\TextAt{13}{345.949}{579.088}{x}%
\TextAt{8}{356.499}{579.088}{for}%
\TextAt{8}{374.383}{579.088}{all}%
\TextAt{13}{390.643}{579.088}{\Character{27}}%
\NativeMathAt{401.046}{579.088}{11.955}{\in}%
\TextAt{13}{412.337}{579.088}{R}%
\TextAt{21}{421.247}{577.294}{w}%
\TextAt{16}{428.036}{579.088}{\Character{103}}%
\TextAt{13}{434.014}{579.088}{\Character{59}}%
\TextAt{8}{79.226}{552.794}{is}%
\TextAt{8}{90.99}{552.794}{called}%
\TextAt{8}{124.154}{552.794}{the}%
\TextAt{12}{144.31}{552.794}{ramification}%
\TextAt{12}{218.958}{552.794}{field}%
\TextAt{8}{247.674}{552.794}{of}%
\TextAt{13}{261.005}{552.794}{w}%
\TextAt{8}{273.591}{552.794}{over}%
\TextAt{13}{298.638}{552.794}{K}%
\TextAt{8}{309.455}{552.794}{.}%
\TextAt{12}{79.226}{522.351}{Proposition}%
\TextAt{12}{151.937}{522.351}{(9.14).}%
\TextAt{13}{193.984}{522.351}{V}%
\TextAt{21}{200.808}{520.558}{w}%
\TextAt{16}{210.918}{522.351}{\Character{106}}%
\TextAt{13}{217.56}{522.351}{Z}%
\TextAt{21}{225.572}{520.558}{w}%
\TextAt{8}{236.259}{522.351}{is}%
\TextAt{8}{248.023}{522.351}{the}%
\TextAt{8}{268.18}{522.351}{maximal}%
\TextAt{8}{315.988}{522.351}{tamely}%
\TextAt{8}{354.687}{522.351}{ramified}%
\TextAt{8}{400.212}{522.351}{subextension}%
\TextAt{8}{469.923}{522.351}{of}%
\TextAt{13}{483.253}{522.351}{L}%
\TextAt{16}{494.538}{522.351}{\Character{106}}%
\TextAt{13}{501.18}{522.351}{Z}%
\TextAt{21}{509.192}{520.558}{w}%
\TextAt{8}{515.982}{522.351}{.}%
\TextAt{19}{79.226}{488.919}{Proof.}%
\TextAt{8}{116}{488.919}{By}%
\TextAt{8}{134.701}{488.919}{(9.6)}%
\TextAt{8}{163.014}{488.919}{and}%
\TextAt{8}{186.115}{488.919}{the}%
\TextAt{8}{206.609}{488.919}{fact}%
\TextAt{8}{230.033}{488.919}{that}%
\TextAt{8}{255.739}{488.919}{the}%
\TextAt{8}{276.234}{488.919}{value}%
\TextAt{8}{306.818}{488.919}{group}%
\TextAt{8}{340.332}{488.919}{and}%
\TextAt{8}{363.432}{488.919}{residue}%
\TextAt{8}{403.497}{488.919}{class}%
\TextAt{8}{431.272}{488.919}{field}%
\TextAt{8}{456.967}{488.919}{do}%
\TextAt{8}{473.564}{488.919}{not}%
\TextAt{8}{494.716}{488.919}{change,}%
\TextAt{8}{79.226}{471.873}{we}%
\TextAt{8}{96.825}{471.873}{may}%
\TextAt{8}{122.565}{471.873}{assume}%
\TextAt{8}{163.381}{471.873}{that}%
\TextAt{13}{189.122}{471.873}{K}%
\TextAt{9}{203.893}{471.873}{\Character{61}}%
\TextAt{13}{216.952}{471.873}{Z}%
\TextAt{21}{224.964}{470.08}{w}%
\TextAt{8}{236.022}{471.873}{is}%
\TextAt{8}{248.158}{471.873}{henselian.}%
\TextAt{8}{304.609}{471.873}{Let}%
\TextAt{13}{325.95}{471.873}{V}%
\TextAt{21}{332.774}{470.08}{w}%
\TextAt{8}{343.832}{471.873}{be}%
\TextAt{8}{360.128}{471.873}{the}%
\TextAt{8}{380.656}{471.873}{fixed}%
\TextAt{8}{409.313}{471.873}{field}%
\TextAt{8}{435.042}{471.873}{of}%
\TextAt{13}{448.743}{471.873}{R}%
\TextAt{21}{457.653}{470.08}{w}%
\TextAt{8}{464.443}{471.873}{.}%
\TextAt{8}{474.006}{471.873}{Since}%
\TextAt{13}{504.935}{471.873}{R}%
\TextAt{21}{513.845}{470.08}{w}%
\TextAt{8}{524.903}{471.873}{is}%
\TextAt{8}{79.226}{454.827}{the}%
\TextAt{13}{99.474}{454.827}{p}%
\TextAt{8}{105.349}{454.827}{-Sylow}%
\TextAt{8}{143.16}{454.827}{subgroup}%
\TextAt{8}{194.049}{454.827}{of}%
\TextAt{13}{207.47}{454.827}{I}%
\TextAt{21}{212.633}{453.034}{w}%
\TextAt{8}{219.423}{454.827}{,}%
\TextAt{13}{226.686}{454.827}{V}%
\TextAt{21}{233.511}{453.034}{w}%
\TextAt{8}{244.289}{454.827}{is}%
\TextAt{8}{256.145}{454.827}{the}%
\TextAt{8}{276.393}{454.827}{union}%
\TextAt{8}{309.002}{454.827}{of}%
\TextAt{8}{322.424}{454.827}{all}%
\TextAt{8}{338.775}{454.827}{finite}%
\TextAt{8}{368.778}{454.827}{Galois}%
\TextAt{8}{404.783}{454.827}{subextensions}%
\TextAt{8}{479.2}{454.827}{of}%
\TextAt{13}{492.622}{454.827}{L}%
\TextAt{16}{504.063}{454.827}{\Character{106}}%
\TextAt{13}{510.861}{454.827}{T}%
\TextAt{8}{523.337}{454.827}{of}%
\TextAt{8}{79.226}{437.781}{degree}%
\TextAt{8}{115.627}{437.781}{prime}%
\TextAt{8}{148.776}{437.781}{to}%
\TextAt{13}{163.072}{437.781}{p}%
\TextAt{8}{168.947}{437.781}{.}%
\TextAt{8}{177.391}{437.781}{Therefore}%
\TextAt{13}{230.374}{437.781}{V}%
\TextAt{21}{237.198}{435.988}{w}%
\TextAt{8}{247.87}{437.781}{contains}%
\TextAt{8}{293.776}{437.781}{the}%
\TextAt{8}{313.918}{437.781}{maximal}%
\TextAt{8}{361.712}{437.781}{tamely}%
\TextAt{8}{400.396}{437.781}{ramified}%
\TextAt{8}{445.907}{437.781}{extension}%
\TextAt{13}{497.659}{437.781}{V}%
\TextAt{8}{510.967}{437.781}{of}%
\TextAt{13}{524.283}{437.781}{T}%
\TextAt{8}{79.226}{420.735}{(and}%
\TextAt{8}{107.306}{420.735}{thus}%
\TextAt{8}{133.819}{420.735}{of}%
\TextAt{13}{147.912}{420.735}{Z}%
\TextAt{21}{155.924}{418.942}{w}%
\TextAt{8}{162.713}{420.735}{).}%
\TextAt{8}{178.003}{420.735}{Since}%
\TextAt{8}{209.322}{420.735}{the}%
\TextAt{8}{230.241}{420.735}{degree}%
\TextAt{8}{267.419}{420.735}{of}%
\TextAt{8}{281.511}{420.735}{each}%
\TextAt{8}{308.61}{420.735}{finite}%
\TextAt{8}{339.284}{420.735}{subextension}%
\TextAt{13}{409.757}{420.735}{M}%
\TextAt{16}{426.95}{420.735}{\Character{106}}%
\TextAt{13}{434.89}{420.735}{V}%
\TextAt{8}{448.976}{420.735}{of}%
\TextAt{13}{463.068}{420.735}{V}%
\TextAt{21}{469.892}{418.942}{w}%
\TextAt{16}{481.301}{420.735}{\Character{106}}%
\TextAt{13}{489.242}{420.735}{V}%
\TextAt{8}{503.327}{420.735}{is}%
\TextAt{8}{515.853}{420.735}{not}%
\TextAt{8}{79.226}{403.689}{divisible}%
\TextAt{8}{125.297}{403.689}{by}%
\TextAt{13}{141.72}{403.689}{p}%
\TextAt{8}{147.595}{403.689}{,}%
\TextAt{8}{154.949}{403.689}{the}%
\TextAt{8}{175.268}{403.689}{residue}%
\TextAt{8}{215.159}{403.689}{field}%
\TextAt{8}{240.679}{403.689}{extension}%
\TextAt{8}{292.609}{403.689}{of}%
\TextAt{13}{306.102}{403.689}{M}%
\TextAt{16}{322.275}{403.689}{\Character{106}}%
\TextAt{13}{329.195}{403.689}{V}%
\TextAt{8}{342.682}{403.689}{is}%
\TextAt{8}{354.609}{403.689}{separable}%
\TextAt{8}{406.216}{403.689}{(see}%
\TextAt{8}{429.847}{403.689}{the}%
\TextAt{8}{450.167}{403.689}{argument}%
\TextAt{8}{502.694}{403.689}{in}%
\TextAt{8}{516.51}{403.689}{the}%
\TextAt{8}{79.226}{386.643}{proof}%
\TextAt{8}{109.796}{386.643}{of}%
\TextAt{8}{123.126}{386.643}{(9.12)).}%
\TextAt{8}{166.065}{386.643}{Therefore}%
\TextAt{13}{219.062}{386.643}{V}%
\TextAt{21}{225.886}{384.85}{w}%
\TextAt{16}{235.996}{386.643}{\Character{106}}%
\TextAt{13}{242.638}{386.643}{V}%
\TextAt{8}{255.961}{386.643}{is}%
\TextAt{8}{267.725}{386.643}{tamely}%
\TextAt{8}{306.424}{386.643}{ramified,}%
\TextAt{8}{355.201}{386.643}{and}%
\TextAt{13}{377.964}{386.643}{V}%
\TextAt{21}{384.788}{384.85}{w}%
\TextAt{9}{394.898}{386.643}{\Character{61}}%
\TextAt{13}{407.324}{386.643}{V}%
\TextAt{8}{416.749}{386.643}{.}%
\Vector{q 0 G 0 g 0.399 w 452.334 -387.357 m 452.334 -379.287 l S Q}%
\Vector{q 0 G 0 g 0.399 w 452.534 -379.486 m 459.707 -379.486 l S Q}%
\Vector{q 0 G 0 g 0.399 w 452.534 -387.157 m 459.707 -387.157 l S Q}%
\Vector{q 0 G 0 g 0.399 w 459.906 -387.357 m 459.906 -379.287 l S Q}%
\TextAt{12}{79.226}{353.212}{Corollary}%
\TextAt{12}{139.014}{353.212}{(9.15).}%
\TextAt{8}{181.06}{353.212}{We}%
\TextAt{8}{201.205}{353.212}{have}%
\TextAt{8}{228.2}{353.212}{the}%
\TextAt{8}{248.356}{353.212}{exact}%
\TextAt{8}{279.248}{353.212}{sequence}%
\TextAt{9}{210.296}{326.918}{1}%
\NativeArrowAt{219.47}{326.918}{11.955}{19.26}{\longrightarrow}%
\TextAt{13}{242.052}{326.918}{R}%
\TextAt{21}{250.962}{325.125}{w}%
\NativeArrowAt{261.072}{326.918}{11.955}{19.26}{\longrightarrow}%
\TextAt{13}{283.654}{326.918}{I}%
\TextAt{21}{288.817}{325.125}{w}%
\NativeArrowAt{298.927}{326.918}{11.955}{19.26}{\longrightarrow}%
\TextAt{13}{321.509}{326.918}{\Character{31}}%
\TextAt{9}{328.843}{326.918}{\Character{40}}%
\TextAt{13}{333.395}{326.918}{L}%
\TextAt{16}{344.68}{326.918}{\Character{106}}%
\TextAt{13}{351.322}{326.918}{K}%
\TextAt{9}{362.14}{326.918}{\Character{41}}%
\NativeArrowAt{370.013}{326.918}{11.955}{19.26}{\longrightarrow}%
\TextAt{9}{392.595}{326.918}{1}%
\TextAt{13}{398.448}{326.918}{\Character{58}}%
\TextAt{19}{79.226}{284.238}{Proof.}%
\TextAt{8}{116}{284.238}{By}%
\TextAt{8}{134.38}{284.238}{(9.6)}%
\TextAt{8}{162.371}{284.238}{we}%
\TextAt{8}{179.615}{284.238}{may}%
\TextAt{8}{205}{284.238}{assume,}%
\TextAt{8}{248.717}{284.238}{as}%
\TextAt{8}{263.103}{284.238}{we}%
\TextAt{8}{280.347}{284.238}{have}%
\TextAt{8}{307.358}{284.238}{already}%
\TextAt{8}{348.679}{284.238}{done}%
\TextAt{8}{376.659}{284.238}{several}%
\TextAt{8}{415.111}{284.238}{times}%
\TextAt{8}{446.402}{284.238}{before,}%
\TextAt{8}{484.786}{284.238}{that}%
\TextAt{13}{510.172}{284.238}{K}%
\TextAt{8}{524.903}{284.238}{is}%
\TextAt{8}{79.226}{267.192}{henselian.}%
\TextAt{8}{135.227}{267.192}{We}%
\TextAt{8}{155.592}{267.192}{restrict}%
\TextAt{8}{196.198}{267.192}{to}%
\TextAt{8}{210.73}{267.192}{considering}%
\TextAt{8}{272.15}{267.192}{the}%
\TextAt{8}{292.527}{267.192}{case}%
\TextAt{8}{317.52}{267.192}{of}%
\TextAt{8}{331.071}{267.192}{a}%
\TextAt{8}{341.048}{267.192}{finite}%
\TextAt{8}{371.181}{267.192}{extension}%
\TextAt{13}{423.168}{267.192}{L}%
\TextAt{16}{434.831}{267.192}{\Character{106}}%
\TextAt{13}{441.849}{267.192}{K}%
\TextAt{8}{452.667}{267.192}{.}%
\TextAt{8}{461.779}{267.192}{In}%
\TextAt{8}{476.622}{267.192}{the}%
\TextAt{8}{497}{267.192}{infinite}%
\TextAt{8}{79.226}{250.146}{case}%
\TextAt{8}{103.501}{250.146}{the}%
\TextAt{8}{123.162}{250.146}{proof}%
\TextAt{8}{153.236}{250.146}{follows}%
\TextAt{8}{191.176}{250.146}{as}%
\TextAt{8}{205.05}{250.146}{in}%
\TextAt{8}{218.207}{250.146}{(9.9).}%
\TextAt{8}{250.568}{250.146}{We}%
\TextAt{8}{270.216}{250.146}{have}%
\TextAt{8}{296.715}{250.146}{already}%
\TextAt{8}{337.525}{250.146}{seen}%
\TextAt{8}{362.445}{250.146}{that}%
\TextAt{13}{387.319}{250.146}{R}%
\TextAt{21}{396.229}{248.353}{w}%
\TextAt{8}{406.42}{250.146}{is}%
\TextAt{8}{417.688}{250.146}{the}%
\TextAt{8}{437.348}{250.146}{kernel}%
\TextAt{8}{471.319}{250.146}{of}%
\TextAt{8}{484.153}{250.146}{the}%
\TextAt{8}{503.814}{250.146}{arrow}%
\TextAt{8}{79.226}{233.1}{on}%
\TextAt{8}{95.485}{233.1}{the}%
\TextAt{8}{115.642}{233.1}{right.}%
\TextAt{8}{148.49}{233.1}{It}%
\TextAt{8}{161.163}{233.1}{therefore}%
\TextAt{8}{210.263}{233.1}{suffices}%
\TextAt{8}{250.05}{233.1}{to}%
\TextAt{8}{264.36}{233.1}{show}%
\TextAt{8}{293.363}{233.1}{that}%
\TextAt{9}{216.001}{206.806}{\Character{40}}%
\TextAt{13}{220.553}{206.806}{I}%
\TextAt{21}{225.716}{205.013}{w}%
\TextAt{9}{235.827}{206.806}{\Character{58}}%
\TextAt{13}{242.399}{206.806}{R}%
\TextAt{21}{251.309}{205.013}{w}%
\TextAt{9}{258.099}{206.806}{\Character{41}}%
\TextAt{9}{265.954}{206.806}{\Character{61}}%
\TextAt{9}{278.375}{206.806}{\Character{91}}%
\TextAt{13}{281.649}{206.806}{V}%
\TextAt{21}{288.473}{205.013}{w}%
\TextAt{9}{298.584}{206.806}{\Character{58}}%
\TextAt{13}{305.156}{206.806}{T}%
\TextAt{21}{312.017}{205.013}{w}%
\TextAt{9}{318.807}{206.806}{\Character{93}}%
\TextAt{9}{325.358}{206.806}{\Character{61}}%
\TextAt{8}{337.804}{206.806}{\char35\relax }%
\TextAt{13}{347.56}{206.806}{\Character{31}}%
\TextAt{9}{354.893}{206.806}{\Character{40}}%
\TextAt{13}{359.446}{206.806}{L}%
\TextAt{16}{370.731}{206.806}{\Character{106}}%
\TextAt{13}{377.373}{206.806}{K}%
\TextAt{9}{388.191}{206.806}{\Character{41}}%
\TextAt{13}{392.743}{206.806}{\Character{58}}%
\TextAt{8}{79.226}{180.513}{As}%
\TextAt{13}{96.669}{180.513}{T}%
\TextAt{21}{103.53}{178.719}{w}%
\TextAt{16}{113.905}{180.513}{\Character{106}}%
\TextAt{13}{120.812}{180.513}{K}%
\TextAt{8}{135.682}{180.513}{is}%
\TextAt{8}{147.601}{180.513}{the}%
\TextAt{8}{167.913}{180.513}{maximal}%
\TextAt{8}{215.878}{180.513}{unramified}%
\TextAt{8}{274.565}{180.513}{subextension}%
\TextAt{8}{344.432}{180.513}{of}%
\TextAt{13}{357.917}{180.513}{V}%
\TextAt{21}{364.741}{178.719}{w}%
\TextAt{16}{375.117}{180.513}{\Character{106}}%
\TextAt{13}{382.023}{180.513}{K}%
\TextAt{8}{392.841}{180.513}{,}%
\TextAt{13}{400.184}{180.513}{V}%
\TextAt{21}{407.009}{178.719}{w}%
\TextAt{16}{417.384}{180.513}{\Character{106}}%
\TextAt{13}{424.291}{180.513}{T}%
\TextAt{21}{431.152}{178.719}{w}%
\TextAt{8}{441.994}{180.513}{has}%
\TextAt{8}{463.023}{180.513}{inertia}%
\TextAt{8}{500.251}{180.513}{degree}%
\TextAt{8}{79.226}{163.467}{1.}%
\TextAt{8}{93.532}{163.467}{Thus,}%
\TextAt{8}{126.433}{163.467}{by}%
\TextAt{8}{142.692}{163.467}{(7.7),}%
\TextAt{9}{229.686}{146.421}{\Character{91}}%
\TextAt{13}{232.938}{146.421}{V}%
\TextAt{21}{239.762}{144.627}{w}%
\TextAt{9}{249.873}{146.421}{\Character{58}}%
\TextAt{13}{256.445}{146.421}{T}%
\TextAt{21}{263.306}{144.627}{w}%
\TextAt{9}{270.096}{146.421}{\Character{93}}%
\TextAt{9}{276.647}{146.421}{\Character{61}}%
\TextAt{8}{289.093}{146.421}{\char35\relax }%
\TextAt{20}{298.849}{156.104}{\Character{0}}%
\TextAt{13}{304.328}{146.421}{w}%
\TextAt{9}{313.018}{146.421}{\Character{40}}%
\TextAt{13}{317.57}{146.421}{V}%
\TextAt{3}{326.996}{151.357}{\Character{3}}%
\TextAt{21}{324.395}{143.465}{w}%
\TextAt{9}{331.728}{146.421}{\Character{41}\Character{47}}%
\TextAt{13}{342.134}{146.421}{w}%
\TextAt{9}{350.823}{146.421}{\Character{40}}%
\TextAt{13}{355.375}{146.421}{T}%
\TextAt{3}{363.862}{151.357}{\Character{3}}%
\TextAt{21}{362.236}{143.465}{w}%
\TextAt{9}{369.026}{146.421}{\Character{41}}%
\TextAt{20}{373.578}{156.104}{\Character{1}}%
\TextAt{13}{379.058}{146.421}{\Character{58}}%
\TextAt{8}{79.226}{124.059}{Furthermore,}%
\TextAt{8}{149.824}{124.059}{by}%
\TextAt{8}{165.519}{124.059}{(7.5),}%
\TextAt{8}{196.295}{124.059}{we}%
\TextAt{8}{212.958}{124.059}{have}%
\TextAt{13}{239.389}{124.059}{w}%
\TextAt{9}{248.079}{124.059}{\Character{40}}%
\TextAt{13}{252.631}{124.059}{T}%
\TextAt{3}{261.118}{128.398}{\Character{3}}%
\TextAt{21}{259.492}{121.104}{w}%
\TextAt{9}{266.281}{124.059}{\Character{41}}%
\TextAt{9}{274.136}{124.059}{\Character{61}}%
\TextAt{13}{286.58}{124.059}{v}%
\TextAt{9}{292.668}{124.059}{\Character{40}}%
\TextAt{13}{297.221}{124.059}{K}%
\TextAt{3}{308.038}{128.398}{\Character{3}}%
\TextAt{9}{312.771}{124.059}{\Character{41}}%
\TextAt{9}{320.626}{124.059}{\Character{61}\Character{58}}%
\TextAt{9}{336.287}{124.059}{\Character{0}}%
\TextAt{8}{343.635}{124.059}{,}%
\TextAt{8}{350.333}{124.059}{and}%
\TextAt{8}{372.532}{124.059}{putting}%
\TextAt{9}{413.595}{124.059}{\Character{1}}%
\TextAt{9}{426.65}{124.059}{\Character{61}}%
\TextAt{13}{439.097}{124.059}{w}%
\TextAt{9}{447.786}{124.059}{\Character{40}}%
\TextAt{13}{452.339}{124.059}{L}%
\TextAt{3}{460.303}{128.398}{\Character{3}}%
\TextAt{9}{465.035}{124.059}{\Character{41}}%
\TextAt{8}{469.588}{124.059}{,}%
\TextAt{8}{476.286}{124.059}{we}%
\TextAt{8}{492.949}{124.059}{see}%
\TextAt{8}{511.298}{124.059}{that}%
\TextAt{13}{79.226}{107.013}{w}%
\TextAt{9}{87.916}{107.013}{\Character{40}}%
\TextAt{13}{92.468}{107.013}{V}%
\TextAt{3}{101.894}{111.352}{\Character{3}}%
\TextAt{21}{99.292}{104.058}{w}%
\TextAt{9}{106.626}{107.013}{\Character{41}\Character{47}}%
\TextAt{13}{117.031}{107.013}{v}%
\TextAt{9}{123.12}{107.013}{\Character{40}}%
\TextAt{13}{127.672}{107.013}{K}%
\TextAt{3}{138.49}{111.352}{\Character{3}}%
\TextAt{9}{143.222}{107.013}{\Character{41}}%
\TextAt{8}{151.745}{107.013}{is}%
\TextAt{8}{163.583}{107.013}{the}%
\TextAt{8}{183.813}{107.013}{subgroup}%
\TextAt{9}{234.685}{107.013}{\Character{1}}%
\TextAt{17}{244.439}{111.352}{\Character{40}}%
\TextAt{21}{247.733}{111.352}{p}%
\TextAt{17}{251.996}{111.352}{\Character{41}}%
\TextAt{9}{255.787}{107.013}{\Character{47}\Character{0}}%
\TextAt{8}{272.925}{107.013}{of}%
\TextAt{9}{286.329}{107.013}{\Character{1}\Character{47}\Character{0}}%
\TextAt{8}{313.222}{107.013}{consisting}%
\TextAt{8}{367.405}{107.013}{of}%
\TextAt{8}{380.809}{107.013}{all}%
\TextAt{8}{397.142}{107.013}{elements}%
\TextAt{8}{445.072}{107.013}{of}%
\TextAt{8}{458.476}{107.013}{order}%
\TextAt{8}{489.119}{107.013}{prime}%
\TextAt{8}{522.357}{107.013}{to}%
\TextAt{13}{79.226}{89.967}{p}%
\TextAt{8}{85.101}{89.967}{,}%
\TextAt{8}{92.25}{89.967}{where}%
\TextAt{13}{126.06}{89.967}{p}%
\TextAt{9}{135.256}{89.967}{\Character{61}}%
\TextAt{8}{147.681}{89.967}{char}%
\TextAt{9}{169.798}{89.967}{\Character{40}}%
\TextAt{13}{174.351}{89.967}{\Character{20}}%
\TextAt{9}{181.089}{89.967}{\Character{41}}%
\TextAt{8}{185.641}{89.967}{.}%
\TextAt{8}{194.09}{89.967}{Thus}%
\TextAt{9}{248.104}{72.921}{\Character{91}}%
\TextAt{13}{251.355}{72.921}{V}%
\TextAt{21}{258.18}{71.128}{w}%
\TextAt{9}{268.29}{72.921}{\Character{58}}%
\TextAt{13}{274.862}{72.921}{T}%
\TextAt{21}{281.723}{71.128}{w}%
\TextAt{9}{288.513}{72.921}{\Character{93}}%
\TextAt{9}{295.064}{72.921}{\Character{61}}%
\TextAt{8}{307.511}{72.921}{\char35\relax }%
\TextAt{9}{317.266}{72.921}{\Character{40}\Character{1}}%
\TextAt{17}{331.573}{77.858}{\Character{40}}%
\TextAt{21}{334.867}{77.858}{p}%
\TextAt{17}{339.129}{77.858}{\Character{41}}%
\TextAt{9}{342.921}{72.921}{\Character{47}\Character{0}\Character{41}}%
\TextAt{13}{360.64}{72.921}{\Character{58}}%
\TextAt{8}{297.211}{43.033}{160}%
\Destination{78.23}{809.815}{nav.page.160}%
\Destination{79.36}{620.355}{nav.print.175}%
\Destination{79.36}{619.355}{nav.number.2.9.13}%
\Destination{79.36}{536.355}{nav.number.2.9.14}%
\Destination{79.36}{367.255}{nav.number.2.9.15}%
\Hotlink{203.27}{741.855}{25.078}{11.047}{nav.number.2.7.7}%
\Hotlink{134.201}{486.755}{25.078}{11.047}{nav.number.2.9.6}%
\Hotlink{122.626}{384.455}{30.936}{11.047}{nav.number.2.9.12}%
\Hotlink{133.88}{282.055}{25.078}{11.047}{nav.number.2.9.6}%
\Hotlink{217.707}{247.955}{25.078}{11.047}{nav.number.2.9.9}%
\Hotlink{142.192}{161.255}{25.078}{11.047}{nav.number.2.7.7}%
\Hotlink{165.019}{121.855}{25.078}{11.047}{nav.number.2.7.5}%
}

% Page 173 (body page 161)
\Page{%
\TextAt{2}{79.226}{796.943}{II.}%
\TextAt{2}{94.073}{796.943}{The}%
\TextAt{2}{116.492}{796.943}{Theory}%
\TextAt{2}{154.401}{796.943}{of}%
\TextAt{2}{166.826}{796.943}{Valuations}%
\Vector{q 0 G 0 g 0.299 w 7.226 18.729 m 460.769 18.729 l S Q}%
\TextAt{8}{79.226}{761.124}{Since}%
\TextAt{13}{109.784}{761.124}{\Character{21}}%
\TextAt{3}{116.612}{765.462}{\Character{3}}%
\TextAt{8}{125.242}{761.124}{has}%
\TextAt{8}{146.115}{761.124}{no}%
\TextAt{8}{162.375}{761.124}{elements}%
\TextAt{8}{210.231}{761.124}{of}%
\TextAt{8}{223.561}{761.124}{order}%
\TextAt{8}{254.13}{761.124}{divisible}%
\TextAt{8}{300.038}{761.124}{by}%
\TextAt{13}{316.297}{761.124}{p}%
\TextAt{8}{322.172}{761.124}{,}%
\TextAt{8}{329.322}{761.124}{we}%
\TextAt{8}{346.549}{761.124}{have}%
\TextAt{8}{373.544}{761.124}{on}%
\TextAt{8}{389.803}{761.124}{the}%
\TextAt{8}{409.959}{761.124}{other}%
\TextAt{8}{440.529}{761.124}{hand}%
\TextAt{8}{469.795}{761.124}{that}%
\TextAt{13}{187.14}{732.155}{\Character{31}}%
\TextAt{9}{194.474}{732.155}{\Character{40}}%
\TextAt{13}{199.026}{732.155}{L}%
\TextAt{16}{210.311}{732.155}{\Character{106}}%
\TextAt{13}{216.953}{732.155}{K}%
\TextAt{9}{227.771}{732.155}{\Character{41}}%
\TextAt{9}{235.626}{732.155}{\Character{61}}%
\TextAt{8}{248.069}{732.155}{Hom}%
\TextAt{9}{272.458}{732.155}{\Character{40}\Character{1}\Character{47}\Character{0}}%
\TextAt{13}{299.932}{732.155}{\Character{59}}%
\TextAt{13}{305.156}{732.155}{\Character{21}}%
\TextAt{3}{312.005}{737.092}{\Character{3}}%
\TextAt{9}{316.737}{732.155}{\Character{41}}%
\TextAt{9}{324.592}{732.155}{\Character{61}}%
\TextAt{8}{337.036}{732.155}{Hom}%
\TextAt{9}{361.425}{732.155}{\Character{40}\Character{1}}%
\TextAt{17}{375.732}{737.092}{\Character{40}}%
\TextAt{21}{379.025}{737.092}{p}%
\TextAt{17}{383.288}{737.092}{\Character{41}}%
\TextAt{9}{387.079}{732.155}{\Character{47}\Character{0}}%
\TextAt{13}{400.247}{732.155}{\Character{59}}%
\TextAt{13}{405.471}{732.155}{\Character{21}}%
\TextAt{3}{412.319}{737.092}{\Character{3}}%
\TextAt{9}{417.051}{732.155}{\Character{41}}%
\TextAt{13}{421.604}{732.155}{\Character{58}}%
\TextAt{8}{79.226}{703.187}{But}%
\TextAt{8}{104.12}{703.187}{(7.7)}%
\TextAt{8}{133.747}{703.187}{implies}%
\TextAt{8}{175.127}{703.187}{that}%
\TextAt{13}{202.148}{703.187}{\Character{21}}%
\TextAt{3}{208.977}{707.526}{\Character{3}}%
\TextAt{8}{219.259}{703.187}{contains}%
\TextAt{8}{266.831}{703.187}{the}%
\TextAt{13}{288.641}{703.187}{m}%
\TextAt{8}{298.88}{703.187}{-th}%
\TextAt{8}{319.386}{703.187}{roots}%
\TextAt{8}{350.699}{703.187}{of}%
\TextAt{8}{365.682}{703.187}{unity}%
\TextAt{8}{397.904}{703.187}{whenever}%
\TextAt{9}{450.928}{703.187}{\Character{1}}%
\TextAt{17}{460.682}{707.526}{\Character{40}}%
\TextAt{21}{463.976}{707.526}{p}%
\TextAt{17}{468.239}{707.526}{\Character{41}}%
\TextAt{9}{472.03}{703.187}{\Character{47}\Character{0}}%
\TextAt{8}{490.747}{703.187}{contains}%
\TextAt{8}{79.226}{686.141}{an}%
\TextAt{8}{96.538}{686.141}{element}%
\TextAt{8}{140.832}{686.141}{of}%
\TextAt{8}{155.215}{686.141}{order}%
\TextAt{13}{186.836}{686.141}{m}%
\TextAt{8}{197.076}{686.141}{,}%
\TextAt{8}{205.541}{686.141}{because}%
\TextAt{8}{249.895}{686.141}{then}%
\TextAt{8}{277.607}{686.141}{there}%
\TextAt{8}{308.571}{686.141}{is}%
\TextAt{8}{321.388}{686.141}{a}%
\TextAt{8}{332.196}{686.141}{Galois}%
\TextAt{8}{369.162}{686.141}{extension}%
\TextAt{8}{421.98}{686.141}{generated}%
\TextAt{8}{476.365}{686.141}{by}%
\TextAt{8}{493.676}{686.141}{radicals}%
\TextAt{13}{79.226}{669.095}{T}%
\TextAt{21}{86.087}{667.302}{w}%
\TextAt{9}{92.877}{669.095}{\Character{40}}%
\TextAt{22}{100.75}{673.116}{m}%
\TextAt{16}{100.646}{677.706}{\Character{112}}%
\Vector{q 0 G 0 g 0.478 w 38.609 -96.055 m 44.754 -96.055 l S Q}%
\TextAt{13}{110.609}{669.095}{a}%
\TextAt{9}{116.754}{669.095}{\Character{41}}%
\TextAt{16}{124.628}{669.095}{\Character{106}}%
\TextAt{13}{131.271}{669.095}{T}%
\TextAt{21}{138.132}{667.302}{w}%
\TextAt{8}{148.82}{669.095}{of}%
\TextAt{8}{162.151}{669.095}{degree}%
\TextAt{13}{198.567}{669.095}{m}%
\TextAt{8}{208.806}{669.095}{.}%
\TextAt{8}{217.257}{669.095}{This}%
\TextAt{8}{243.978}{669.095}{shows}%
\TextAt{8}{277.597}{669.095}{that}%
\TextAt{13}{302.966}{669.095}{\Character{31}}%
\TextAt{9}{310.3}{669.095}{\Character{40}}%
\TextAt{13}{314.852}{669.095}{L}%
\TextAt{16}{326.139}{669.095}{\Character{106}}%
\TextAt{13}{332.782}{669.095}{K}%
\TextAt{9}{343.6}{669.095}{\Character{41}}%
\TextAt{8}{352.05}{669.095}{is}%
\TextAt{8}{363.815}{669.095}{the}%
\TextAt{8}{383.972}{669.095}{Pontryagin}%
\TextAt{8}{443.988}{669.095}{dual}%
\TextAt{8}{470.003}{669.095}{of}%
\TextAt{8}{483.334}{669.095}{the}%
\TextAt{8}{503.491}{669.095}{group}%
\TextAt{9}{79.226}{652.049}{\Character{1}}%
\TextAt{17}{88.981}{656.388}{\Character{40}}%
\TextAt{21}{92.274}{656.388}{p}%
\TextAt{17}{96.537}{656.388}{\Character{41}}%
\TextAt{9}{100.329}{652.049}{\Character{47}\Character{0}}%
\TextAt{8}{117.393}{652.049}{so}%
\TextAt{8}{131.763}{652.049}{that}%
\TextAt{8}{157.132}{652.049}{indeed}%
\TextAt{9}{212.761}{623.081}{\Character{91}}%
\TextAt{13}{216.013}{623.081}{V}%
\TextAt{21}{222.837}{621.288}{w}%
\TextAt{9}{232.948}{623.081}{\Character{58}}%
\TextAt{13}{239.52}{623.081}{T}%
\TextAt{21}{246.381}{621.288}{w}%
\TextAt{9}{253.17}{623.081}{\Character{93}}%
\TextAt{9}{259.721}{623.081}{\Character{61}}%
\TextAt{8}{272.168}{623.081}{\char35\relax }%
\TextAt{9}{281.924}{623.081}{\Character{40}\Character{1}}%
\TextAt{17}{296.231}{628.017}{\Character{40}}%
\TextAt{21}{299.524}{628.017}{p}%
\TextAt{17}{303.787}{628.017}{\Character{41}}%
\TextAt{9}{307.579}{623.081}{\Character{47}\Character{0}\Character{41}}%
\TextAt{9}{328.584}{623.081}{\Character{61}}%
\TextAt{8}{341.044}{623.081}{\char35\relax }%
\TextAt{13}{350.8}{623.081}{\Character{31}}%
\TextAt{9}{358.133}{623.081}{\Character{40}}%
\TextAt{13}{362.686}{623.081}{L}%
\TextAt{16}{373.971}{623.081}{\Character{106}}%
\TextAt{13}{380.613}{623.081}{K}%
\TextAt{9}{391.43}{623.081}{\Character{41}}%
\TextAt{13}{395.983}{623.081}{\Character{58}}%
\Vector{q 0 G 0 g 0.399 w 452.334 -179.887 m 452.334 -171.817 l S Q}%
\Vector{q 0 G 0 g 0.399 w 452.534 -172.017 m 459.707 -172.017 l S Q}%
\Vector{q 0 G 0 g 0.399 w 452.534 -179.688 m 459.707 -179.688 l S Q}%
\Vector{q 0 G 0 g 0.399 w 459.906 -179.887 m 459.906 -171.817 l S Q}%
\TextAt{12}{79.226}{559.153}{Exercise}%
\TextAt{12}{132.903}{559.153}{II.9.1.}%
\TextAt{8}{172.29}{559.153}{Let}%
\TextAt{13}{193.673}{559.153}{K}%
\TextAt{8}{208.802}{559.153}{be}%
\TextAt{8}{225.14}{559.153}{a}%
\TextAt{8}{235.31}{559.153}{henselian}%
\TextAt{8}{286.509}{559.153}{field,}%
\TextAt{13}{315.635}{559.153}{L}%
\TextAt{16}{327.625}{559.153}{\Character{106}}%
\TextAt{13}{334.972}{559.153}{K}%
\TextAt{8}{350.101}{559.153}{a}%
\TextAt{8}{360.27}{559.153}{tamely}%
\TextAt{8}{399.383}{559.153}{ramified}%
\TextAt{8}{445.322}{559.153}{Galois}%
\TextAt{8}{481.649}{559.153}{extension,}%
\TextAt{13}{79.226}{542.107}{G}%
\TextAt{9}{91.942}{542.107}{\Character{61}}%
\TextAt{13}{104.528}{542.107}{G}%
\TextAt{9}{113.762}{542.107}{\Character{40}}%
\TextAt{13}{118.314}{542.107}{L}%
\TextAt{16}{129.761}{542.107}{\Character{106}}%
\TextAt{13}{136.564}{542.107}{K}%
\TextAt{9}{147.381}{542.107}{\Character{41}}%
\TextAt{8}{151.934}{542.107}{,}%
\TextAt{13}{159.201}{542.107}{I}%
\TextAt{9}{168.786}{542.107}{\Character{61}}%
\TextAt{13}{181.372}{542.107}{I}%
\TextAt{9}{187.475}{542.107}{\Character{40}}%
\TextAt{13}{192.028}{542.107}{L}%
\TextAt{16}{203.474}{542.107}{\Character{106}}%
\TextAt{13}{210.277}{542.107}{K}%
\TextAt{9}{221.095}{542.107}{\Character{41}}%
\TextAt{8}{229.639}{542.107}{and}%
\TextAt{9}{252.496}{542.107}{\Character{0}}%
\TextAt{9}{263.28}{542.107}{\Character{61}}%
\TextAt{13}{275.879}{542.107}{G}%
\TextAt{9}{285.112}{542.107}{\Character{47}}%
\TextAt{13}{290.965}{542.107}{I}%
\TextAt{9}{300.55}{542.107}{\Character{61}}%
\TextAt{13}{313.137}{542.107}{G}%
\TextAt{9}{322.37}{542.107}{\Character{40}}%
\TextAt{13}{326.923}{542.107}{\Character{21}}%
\TextAt{16}{337.233}{542.107}{\Character{106}}%
\TextAt{13}{344.036}{542.107}{\Character{20}}%
\TextAt{9}{350.775}{542.107}{\Character{41}}%
\TextAt{8}{355.327}{542.107}{.}%
\TextAt{8}{364.059}{542.107}{Then}%
\TextAt{13}{394.711}{542.107}{I}%
\TextAt{8}{404.806}{542.107}{is}%
\TextAt{8}{416.664}{542.107}{abelian}%
\TextAt{8}{457.406}{542.107}{and}%
\TextAt{8}{480.263}{542.107}{becomes}%
\TextAt{8}{526.911}{542.107}{a}%
\TextAt{9}{79.226}{525.061}{\Character{0}}%
\TextAt{8}{86.54}{525.061}{-module}%
\TextAt{8}{131.731}{525.061}{by}%
\TextAt{8}{147.99}{525.061}{letting}%
\TextAt{9}{185.678}{525.061}{\Character{22}}%
\TextAt{13}{185.063}{525.061}{\Character{27}}%
\TextAt{9}{195.466}{525.061}{\Character{61}}%
\TextAt{13}{207.892}{525.061}{\Character{27}}%
\TextAt{13}{214.958}{525.061}{I}%
\NativeMathAt{224.398}{525.061}{11.955}{\in}%
\TextAt{9}{235.689}{525.061}{\Character{0}}%
\TextAt{8}{246.9}{525.061}{operate}%
\TextAt{8}{288.851}{525.061}{on}%
\TextAt{13}{305.11}{525.061}{I}%
\TextAt{8}{315.11}{525.061}{via}%
\TextAt{13}{334.298}{525.061}{\Character{28}}%
\NativeArrowAt{344.038}{525.061}{11.955}{11.955}{\mapsto}%
\TextAt{13}{359.314}{525.061}{\Character{27}}%
\TextAt{13}{366.38}{525.061}{\Character{28}}%
\TextAt{13}{372.8}{525.061}{\Character{27}}%
\TextAt{3}{379.898}{529.399}{\Character{0}}%
\TextAt{17}{386.484}{529.399}{1}%
\TextAt{8}{391.216}{525.061}{.}%
\TextAt{8}{97.159}{505.026}{Show}%
\TextAt{8}{129.911}{505.026}{that}%
\TextAt{8}{157.14}{505.026}{there}%
\TextAt{8}{188.912}{505.026}{is}%
\TextAt{8}{202.535}{505.026}{a}%
\TextAt{8}{214.151}{505.026}{canonical}%
\TextAt{8}{267.394}{505.026}{isomorphism}%
\TextAt{13}{337.673}{505.026}{I}%
\NativeMathAt{350.266}{505.026}{11.955}{\cong}%
\TextAt{13}{366.055}{505.026}{\Character{31}}%
\TextAt{9}{373.389}{505.026}{\Character{40}}%
\TextAt{13}{377.941}{505.026}{L}%
\TextAt{16}{392.396}{505.026}{\Character{106}}%
\TextAt{13}{402.207}{505.026}{K}%
\TextAt{9}{413.025}{505.026}{\Character{41}}%
\TextAt{8}{423.335}{505.026}{of}%
\TextAt{9}{438.525}{505.026}{\Character{0}}%
\TextAt{8}{445.839}{505.026}{-modules.}%
\TextAt{8}{505.775}{505.026}{Show}%
\TextAt{8}{79.226}{487.98}{furthermore}%
\TextAt{8}{143.115}{487.98}{that}%
\TextAt{8}{167.659}{487.98}{every}%
\TextAt{8}{197.726}{487.98}{tamely}%
\TextAt{8}{235.6}{487.98}{ramified}%
\TextAt{8}{280.3}{487.98}{extension}%
\TextAt{8}{331.241}{487.98}{can}%
\TextAt{8}{351.876}{487.98}{be}%
\TextAt{8}{366.975}{487.98}{embedded}%
\TextAt{8}{421.419}{487.98}{into}%
\TextAt{8}{444.337}{487.98}{a}%
\TextAt{8}{453.268}{487.98}{tamely}%
\TextAt{8}{491.142}{487.98}{ramified}%
\TextAt{8}{79.226}{470.934}{extension}%
\TextAt{13}{132.376}{470.934}{L}%
\TextAt{16}{146.018}{470.934}{\Character{106}}%
\TextAt{13}{155.018}{470.934}{K}%
\TextAt{8}{165.836}{470.934}{,}%
\TextAt{8}{174.714}{470.934}{such}%
\TextAt{8}{202.495}{470.934}{that}%
\TextAt{13}{229.247}{470.934}{G}%
\TextAt{8}{243.761}{470.934}{is}%
\TextAt{8}{256.909}{470.934}{the}%
\TextAt{8}{278.449}{470.934}{semi-direct}%
\TextAt{8}{339.716}{470.934}{product}%
\TextAt{8}{384.998}{470.934}{of}%
\TextAt{13}{399.712}{470.934}{\Character{31}}%
\TextAt{9}{407.045}{470.934}{\Character{40}}%
\TextAt{13}{411.598}{470.934}{L}%
\TextAt{16}{425.241}{470.934}{\Character{106}}%
\TextAt{13}{434.24}{470.934}{K}%
\TextAt{9}{445.058}{470.934}{\Character{41}}%
\TextAt{8}{454.891}{470.934}{with}%
\TextAt{13}{482.935}{470.934}{G}%
\TextAt{9}{492.168}{470.934}{\Character{40}}%
\TextAt{13}{496.72}{470.934}{\Character{21}}%
\TextAt{16}{509.227}{470.934}{\Character{106}}%
\TextAt{13}{518.227}{470.934}{\Character{20}}%
\TextAt{9}{524.965}{470.934}{\Character{41}}%
\TextAt{8}{529.518}{470.934}{:}%
\TextAt{13}{79.226}{453.888}{G}%
\NativeMathAt{91.781}{453.888}{11.955}{\cong}%
\TextAt{13}{104.4}{453.888}{\Character{31}}%
\TextAt{9}{111.733}{453.888}{\Character{40}}%
\TextAt{13}{116.286}{453.888}{L}%
\TextAt{16}{127.571}{453.888}{\Character{106}}%
\TextAt{13}{134.213}{453.888}{K}%
\TextAt{9}{145.031}{453.888}{\Character{41}}%
\TextAt{18}{152.24}{453.888}{o}%
\TextAt{13}{164.195}{453.888}{G}%
\TextAt{9}{173.428}{453.888}{\Character{40}}%
\TextAt{13}{177.981}{453.888}{\Character{21}}%
\TextAt{16}{188.13}{453.888}{\Character{106}}%
\TextAt{13}{194.772}{453.888}{\Character{20}}%
\TextAt{9}{201.51}{453.888}{\Character{41}}%
\TextAt{8}{206.063}{453.888}{.}%
\TextAt{12}{97.159}{433.853}{Hint:}%
\TextAt{8}{132.698}{433.853}{Use}%
\TextAt{8}{155.185}{433.853}{(7.7).}%
\TextAt{12}{79.226}{401.882}{Exercise}%
\TextAt{12}{132.902}{401.882}{II.9.2.}%
\TextAt{8}{172.289}{401.882}{The}%
\TextAt{8}{196.756}{401.882}{maximal}%
\TextAt{8}{244.978}{401.882}{tamely}%
\TextAt{8}{284.09}{401.882}{ramified}%
\TextAt{8}{330.029}{401.882}{abelian}%
\TextAt{8}{371.09}{401.882}{extension}%
\TextAt{13}{423.269}{401.882}{V}%
\TextAt{8}{437.005}{401.882}{of}%
\TextAt{18}{450.748}{401.882}{Q}%
\TextAt{21}{460.047}{400.089}{p}%
\TextAt{8}{469.119}{401.882}{is}%
\TextAt{8}{481.296}{401.882}{finite}%
\TextAt{8}{511.621}{401.882}{over}%
\TextAt{8}{79.226}{384.836}{the}%
\TextAt{8}{99.383}{384.836}{maximal}%
\TextAt{8}{147.191}{384.836}{unramified}%
\TextAt{8}{205.724}{384.836}{abelian}%
\TextAt{8}{246.371}{384.836}{extension}%
\TextAt{13}{298.137}{384.836}{T}%
\TextAt{8}{310.521}{384.836}{of}%
\TextAt{18}{323.851}{384.836}{Q}%
\TextAt{21}{333.15}{383.043}{p}%
\TextAt{8}{337.911}{384.836}{.}%
\TextAt{12}{79.226}{352.865}{Exercise}%
\TextAt{12}{133.034}{352.865}{II.9.3.}%
\TextAt{8}{172.537}{352.865}{Show}%
\TextAt{8}{203.957}{352.865}{that}%
\TextAt{8}{229.854}{352.865}{the}%
\TextAt{8}{250.539}{352.865}{maximal}%
\TextAt{8}{298.876}{352.865}{unramified}%
\TextAt{8}{357.937}{352.865}{extension}%
\TextAt{8}{410.232}{352.865}{of}%
\TextAt{8}{424.09}{352.865}{the}%
\TextAt{8}{444.775}{352.865}{power}%
\TextAt{8}{479.447}{352.865}{series}%
\TextAt{8}{511.31}{352.865}{field}%
\ErrataAt{79.22}{335.855}{453.64}{$K=\mathbb F_p((t))$ is $T=\bigcup_{f\ge1}\mathbb F_{p^f}((t))$. Show that the maximal tamely}%
\ErrataAt{79.22}{318.755}{453.64}{ramified extension is $T(\{\sqrt[m]{t}:m\ge1,\ (m,p)=1\})$.}%
\TextAt{12}{79.226}{286.802}{Exercise}%
\TextAt{12}{133.519}{286.802}{II.9.4.}%
\TextAt{8}{173.442}{286.802}{Let}%
\TextAt{13}{195.361}{286.802}{v}%
\TextAt{8}{206.296}{286.802}{be}%
\TextAt{8}{223.17}{286.802}{a}%
\TextAt{8}{233.875}{286.802}{nonarchimedean}%
\TextAt{8}{321.655}{286.802}{valuation}%
\TextAt{8}{373.677}{286.802}{of}%
\TextAt{8}{387.956}{286.802}{the}%
\TextAt{8}{409.062}{286.802}{field}%
\TextAt{13}{435.368}{286.802}{K}%
\TextAt{8}{451.033}{286.802}{and}%
\TextAt{8}{474.745}{286.802}{let}%
\TextAt{9}{493.042}{286.802}{\Character{22}}%
\TextAt{13}{492.599}{286.802}{v}%
\TextAt{8}{503.534}{286.802}{be}%
\TextAt{8}{520.408}{286.802}{an}%
\TextAt{8}{79.226}{269.756}{extension}%
\TextAt{8}{131.007}{269.756}{to}%
\TextAt{8}{145.333}{269.756}{the}%
\TextAt{8}{165.504}{269.756}{separable}%
\TextAt{8}{216.963}{269.756}{closure}%
\Vector{q 0 G 0 g 0.478 w 184.059 -494.401 m 194.877 -494.401 l S Q}%
\TextAt{13}{256.059}{269.756}{K}%
\TextAt{8}{270.789}{269.756}{of}%
\TextAt{13}{284.135}{269.756}{K}%
\TextAt{8}{294.952}{269.756}{.}%
\TextAt{8}{303.446}{269.756}{Then}%
\TextAt{8}{334.019}{269.756}{the}%
\TextAt{8}{354.19}{269.756}{decomposition}%
\TextAt{8}{431.34}{269.756}{field}%
\TextAt{13}{456.712}{269.756}{Z}%
\TextAt{17}{465.047}{267.963}{\Character{22}}%
\TextAt{21}{464.724}{267.963}{v}%
\TextAt{8}{473.545}{269.756}{of}%
\TextAt{9}{487.333}{269.756}{\Character{22}}%
\TextAt{13}{486.89}{269.756}{v}%
\TextAt{8}{496.891}{269.756}{over}%
\TextAt{13}{521.952}{269.756}{K}%
\TextAt{8}{79.226}{252.71}{is}%
\TextAt{8}{90.99}{252.71}{isomorphic}%
\TextAt{8}{150.24}{252.71}{to}%
\TextAt{8}{164.55}{252.71}{the}%
\TextAt{8}{184.707}{252.71}{henselization}%
\TextAt{8}{254.357}{252.71}{of}%
\TextAt{13}{267.687}{252.71}{K}%
\TextAt{8}{282.402}{252.71}{with}%
\TextAt{8}{309.063}{252.71}{respect}%
\TextAt{8}{349.112}{252.71}{to}%
\TextAt{13}{363.423}{252.71}{v}%
\TextAt{8}{369.511}{252.71}{,}%
\TextAt{8}{376.66}{252.71}{in}%
\TextAt{8}{390.313}{252.71}{the}%
\TextAt{8}{410.469}{252.71}{sense}%
\TextAt{8}{440.501}{252.71}{of}%
\TextAt{8}{453.831}{252.71}{§6,}%
\TextAt{8}{472.517}{252.71}{p.}%
\TextAt{8}{486.169}{252.71}{143.}%
\TextAt{15}{79.226}{199.861}{§}%
\TextAt{15}{93.785}{199.861}{10}%
\TextAt{15}{130.385}{199.861}{Higher}%
\TextAt{15}{194.582}{199.861}{Ramification}%
\TextAt{15}{308.6}{199.861}{Groups}%
\TextAt{8}{79.226}{167.975}{The}%
\TextAt{8}{103.277}{167.975}{inertia}%
\TextAt{8}{140.348}{167.975}{group}%
\TextAt{8}{173.521}{167.975}{and}%
\TextAt{8}{196.281}{167.975}{the}%
\TextAt{8}{216.435}{167.975}{ramification}%
\TextAt{8}{281.48}{167.975}{group}%
\TextAt{8}{314.653}{167.975}{inside}%
\TextAt{8}{347.874}{167.975}{the}%
\TextAt{8}{368.028}{167.975}{Galois}%
\TextAt{8}{403.939}{167.975}{group}%
\TextAt{8}{437.112}{167.975}{of}%
\TextAt{8}{450.439}{167.975}{valued}%
\TextAt{8}{487.187}{167.975}{fields}%
\TextAt{8}{517.156}{167.975}{are}%
\TextAt{8}{79.226}{150.929}{only}%
\TextAt{8}{105.457}{150.929}{the}%
\TextAt{8}{126.152}{150.929}{first}%
\TextAt{8}{150.816}{150.929}{terms}%
\TextAt{8}{183.933}{150.929}{in}%
\TextAt{8}{198.125}{150.929}{a}%
\TextAt{8}{208.419}{150.929}{whole}%
\TextAt{8}{242.121}{150.929}{series}%
\TextAt{8}{273.995}{150.929}{of}%
\TextAt{8}{287.863}{150.929}{subgroups}%
\TextAt{8}{343.814}{150.929}{that}%
\TextAt{8}{369.722}{150.929}{we}%
\TextAt{8}{387.488}{150.929}{are}%
\TextAt{8}{407.538}{150.929}{now}%
\TextAt{8}{432.465}{150.929}{going}%
\TextAt{8}{464.231}{150.929}{to}%
\TextAt{8}{479.08}{150.929}{study.}%
\TextAt{8}{516.522}{150.929}{We}%
\TextAt{8}{297.211}{43.033}{161}%
\Destination{78.23}{809.815}{nav.page.161}%
\Destination{79.23}{226.165}{nav.section.2.10}%
\Destination{79.226}{573.155}{nav.ex.2.9.1}%
\Destination{79.226}{415.855}{nav.ex.2.9.2}%
\Destination{79.226}{366.855}{nav.ex.2.9.3}%
\Destination{79.226}{300.755}{nav.ex.2.9.4}%
\Destination{79.36}{718.155}{nav.print.176}%
\Destination{79.23}{226.165}{outline.30}%
\Hotlink{459.01}{250.555}{6.858}{11.047}{nav.section.2.6}%
\Hotlink{485.669}{250.555}{18.574}{11.047}{nav.print.143}%
\Hotlink{103.62}{700.955}{25.078}{11.047}{nav.number.2.7.7}%
\Hotlink{154.685}{431.655}{25.078}{11.047}{nav.number.2.7.7}%
\Destination{79.226}{349.855}{errata.20}%
}

% Page 174 (body page 162)
\Page{%
\TextAt{2}{79.226}{796.943}{II.}%
\TextAt{2}{94.073}{796.943}{The}%
\TextAt{2}{116.492}{796.943}{Theory}%
\TextAt{2}{154.401}{796.943}{of}%
\TextAt{2}{166.826}{796.943}{Valuations}%
\Vector{q 0 G 0 g 0.299 w 7.226 18.729 m 460.769 18.729 l S Q}%
\TextAt{8}{79.226}{761.124}{assume}%
\TextAt{8}{119.709}{761.124}{that}%
\TextAt{13}{145.116}{761.124}{L}%
\TextAt{16}{156.467}{761.124}{\Character{106}}%
\TextAt{13}{163.174}{761.124}{K}%
\TextAt{8}{177.928}{761.124}{is}%
\TextAt{8}{189.73}{761.124}{a}%
\TextAt{8}{199.524}{761.124}{finite}%
\TextAt{8}{229.474}{761.124}{Galois}%
\TextAt{8}{265.426}{761.124}{extension}%
\TextAt{8}{317.23}{761.124}{and}%
\TextAt{8}{340.031}{761.124}{that}%
\TextAt{13}{365.438}{761.124}{v}%
\TextAt{21}{371.098}{759.33}{K}%
\TextAt{8}{383.236}{761.124}{is}%
\TextAt{8}{395.038}{761.124}{a}%
\TextAt{8}{404.832}{761.124}{discrete}%
\TextAt{8}{447.849}{761.124}{normalized}%
\TextAt{8}{507.723}{761.124}{valu-}%
\TextAt{8}{79.226}{744.078}{ation}%
\TextAt{8}{109.499}{744.078}{of}%
\TextAt{13}{123.179}{744.078}{K}%
\TextAt{8}{133.997}{744.078}{,}%
\TextAt{8}{141.583}{744.078}{with}%
\TextAt{8}{168.592}{744.078}{positive}%
\TextAt{8}{212.256}{744.078}{residue}%
\TextAt{8}{252.332}{744.078}{field}%
\TextAt{8}{278.039}{744.078}{characteristic}%
\TextAt{13}{350.322}{744.078}{p}%
\TextAt{8}{356.198}{744.078}{,}%
\TextAt{8}{363.784}{744.078}{which}%
\TextAt{8}{397.62}{744.078}{admits}%
\TextAt{8}{436.405}{744.078}{a}%
\TextAt{8}{446.51}{744.078}{unique}%
\TextAt{8}{484.901}{744.078}{extension}%
\TextAt{13}{79.226}{727.032}{w}%
\TextAt{8}{91.813}{727.032}{to}%
\TextAt{13}{106.123}{727.032}{L}%
\TextAt{8}{114.088}{727.032}{.}%
\TextAt{8}{122.536}{727.032}{We}%
\TextAt{8}{142.681}{727.032}{denote}%
\TextAt{8}{180.399}{727.032}{by}%
\TextAt{13}{196.658}{727.032}{v}%
\TextAt{21}{202.318}{725.238}{L}%
\TextAt{9}{211.896}{727.032}{\Character{61}}%
\TextAt{13}{224.321}{727.032}{ew}%
\TextAt{8}{242.333}{727.032}{the}%
\TextAt{8}{262.49}{727.032}{associated}%
\TextAt{8}{318.225}{727.032}{normalized}%
\TextAt{8}{378.06}{727.032}{valuation}%
\TextAt{8}{429.133}{727.032}{of}%
\TextAt{13}{442.463}{727.032}{L}%
\TextAt{8}{450.427}{727.032}{.}%
\TextAt{12}{79.226}{701.401}{Definition}%
\TextAt{12}{143.769}{701.401}{(10.1).}%
\TextAt{8}{187.341}{701.401}{For}%
\TextAt{8}{209.836}{701.401}{every}%
\TextAt{8}{242.254}{701.401}{real}%
\TextAt{8}{266.542}{701.401}{number}%
\TextAt{13}{310.664}{701.401}{s}%
\NativeMathAt{322.098}{701.401}{11.955}{\geq}%
\TextAt{16}{337.305}{701.401}{\Character{0}}%
\TextAt{9}{346.616}{701.401}{1}%
\TextAt{8}{357.892}{701.401}{we}%
\TextAt{8}{376.644}{701.401}{define}%
\TextAt{8}{411.979}{701.401}{the}%
\TextAt{13}{433.661}{701.401}{s}%
\TextAt{12}{439.175}{701.401}{-th}%
\TextAt{12}{462.605}{701.401}{ramification}%
\TextAt{12}{79.226}{684.355}{group}%
\TextAt{8}{117.028}{684.355}{of}%
\TextAt{13}{130.358}{684.355}{L}%
\TextAt{16}{141.644}{684.355}{\Character{106}}%
\TextAt{13}{148.285}{684.355}{K}%
\TextAt{8}{163.001}{684.355}{by}%
\TextAt{13}{146.941}{656.688}{G}%
\TextAt{21}{156.175}{654.895}{s}%
\TextAt{9}{163.91}{656.688}{\Character{61}}%
\TextAt{13}{176.335}{656.688}{G}%
\TextAt{21}{185.569}{654.895}{s}%
\TextAt{9}{189.983}{656.688}{\Character{40}}%
\TextAt{13}{194.535}{656.688}{L}%
\TextAt{16}{205.821}{656.688}{\Character{106}}%
\TextAt{13}{212.462}{656.688}{K}%
\TextAt{9}{223.28}{656.688}{\Character{41}}%
\TextAt{9}{231.135}{656.688}{\Character{61}}%
\TextAt{16}{243.579}{656.688}{\Character{102}}%
\TextAt{13}{249.556}{656.688}{\Character{27}}%
\NativeMathAt{259.96}{656.688}{11.955}{\in}%
\TextAt{13}{271.251}{656.688}{G}%
\TextAt{16}{283.805}{656.688}{\Character{106}}%
\TextAt{13}{290.447}{656.688}{v}%
\TextAt{21}{296.106}{654.895}{L}%
\TextAt{9}{302.363}{656.688}{\Character{40}}%
\TextAt{13}{306.916}{656.688}{\Character{27}}%
\TextAt{13}{313.982}{656.688}{a}%
\TextAt{16}{322.8}{656.688}{\Character{0}}%
\TextAt{13}{334.755}{656.688}{a}%
\TextAt{9}{340.9}{656.688}{\Character{41}}%
\NativeMathAt{348.773}{656.688}{11.955}{\geq}%
\TextAt{13}{361.392}{656.688}{s}%
\TextAt{9}{369.563}{656.688}{\Character{43}}%
\TextAt{9}{381.315}{656.688}{1}%
\TextAt{8}{391.075}{656.688}{for}%
\TextAt{8}{408.959}{656.688}{all}%
\TextAt{13}{425.218}{656.688}{a}%
\NativeMathAt{434.684}{656.688}{11.955}{\in}%
\TextAt{16}{445.958}{656.688}{\Character{79}}%
\TextAt{16}{455.809}{656.688}{\Character{103}}%
\TextAt{13}{461.803}{656.688}{\Character{58}}%
\ErrataAt{97.06}{620.455}{435.8}{Clearly, $G_{-1}=G$ and $G_0=I(L\mid K)$. If $\lambda\mid\kappa$ is separable, then $G_1=R(L\mid K)$.}%
\ErrataAt{79.22}{603.355}{453.64}{Assume separability also in (10.2) and its consequences. As}%
\TextAt{13}{163}{575.724}{v}%
\TextAt{21}{168.659}{573.931}{L}%
\TextAt{9}{174.916}{575.724}{\Character{40}}%
\TextAt{13}{179.468}{575.724}{\Character{28}}%
\TextAt{3}{185.887}{580.66}{\Character{0}}%
\TextAt{17}{192.474}{580.66}{1}%
\TextAt{13}{197.206}{575.724}{\Character{27}}%
\TextAt{13}{204.272}{575.724}{\Character{28}}%
\TextAt{13}{210.692}{575.724}{a}%
\TextAt{16}{219.509}{575.724}{\Character{0}}%
\TextAt{13}{231.464}{575.724}{a}%
\TextAt{9}{237.609}{575.724}{\Character{41}}%
\TextAt{9}{245.464}{575.724}{\Character{61}}%
\TextAt{13}{257.908}{575.724}{v}%
\TextAt{21}{263.567}{573.931}{L}%
\TextAt{20}{269.824}{585.408}{\Character{0}}%
\TextAt{13}{275.304}{575.724}{\Character{28}}%
\TextAt{3}{281.723}{580.66}{\Character{0}}%
\TextAt{17}{288.309}{580.66}{1}%
\TextAt{9}{293.042}{575.724}{\Character{40}}%
\TextAt{13}{297.594}{575.724}{\Character{27}}%
\TextAt{13}{304.66}{575.724}{\Character{28}}%
\TextAt{13}{311.079}{575.724}{a}%
\TextAt{16}{319.897}{575.724}{\Character{0}}%
\TextAt{13}{331.852}{575.724}{\Character{28}}%
\TextAt{13}{338.26}{575.724}{a}%
\TextAt{9}{344.416}{575.724}{\Character{41}}%
\TextAt{20}{348.968}{585.408}{\Character{1}}%
\TextAt{9}{357.769}{575.724}{\Character{61}}%
\TextAt{13}{370.194}{575.724}{v}%
\TextAt{21}{375.853}{573.931}{L}%
\TextAt{20}{382.11}{585.408}{\Character{0}}%
\TextAt{13}{387.59}{575.724}{\Character{27}}%
\TextAt{9}{394.672}{575.724}{\Character{40}}%
\TextAt{13}{399.225}{575.724}{\Character{28}}%
\TextAt{13}{405.633}{575.724}{a}%
\TextAt{9}{411.789}{575.724}{\Character{41}}%
\TextAt{16}{418.998}{575.724}{\Character{0}}%
\TextAt{13}{430.953}{575.724}{\Character{28}}%
\TextAt{13}{437.361}{575.724}{a}%
\TextAt{20}{443.517}{585.408}{\Character{1}}%
\TextAt{8}{79.226}{548.058}{and}%
\TextAt{13}{101.989}{548.058}{\Character{28}}%
\TextAt{16}{108.408}{548.058}{\Character{79}}%
\TextAt{9}{121.578}{548.058}{\Character{61}}%
\TextAt{16}{134.004}{548.058}{\Character{79}}%
\TextAt{8}{143.854}{548.058}{,}%
\TextAt{8}{151.003}{548.058}{the}%
\TextAt{8}{171.159}{548.058}{ramification}%
\TextAt{8}{236.207}{548.058}{groups}%
\TextAt{8}{273.998}{548.058}{form}%
\TextAt{8}{301.638}{548.058}{a}%
\TextAt{8}{311.393}{548.058}{chain}%
\TextAt{13}{223.429}{520.391}{G}%
\TextAt{9}{235.983}{520.391}{\Character{61}}%
\TextAt{13}{248.409}{520.391}{G}%
\TextAt{3}{257.642}{518.598}{\Character{0}}%
\TextAt{17}{264.229}{518.598}{1}%
\NativeMathAt{272.282}{520.391}{11.955}{\supseteq}%
\TextAt{13}{284.901}{520.391}{G}%
\TextAt{17}{294.135}{518.598}{0}%
\NativeMathAt{302.188}{520.391}{11.955}{\supseteq}%
\TextAt{13}{314.807}{520.391}{G}%
\TextAt{17}{324.041}{518.598}{1}%
\NativeMathAt{332.094}{520.391}{11.955}{\supseteq}%
\TextAt{13}{344.713}{520.391}{G}%
\TextAt{17}{353.947}{518.598}{2}%
\NativeMathAt{362}{520.391}{11.955}{\supseteq}%
\TextAt{16}{374.601}{520.391}{\Character{1}}%
\TextAt{16}{379.909}{520.391}{\Character{1}}%
\TextAt{16}{385.217}{520.391}{\Character{1}}%
\TextAt{8}{79.226}{492.725}{of}%
\TextAt{8}{92.556}{492.725}{normal}%
\TextAt{8}{132.235}{492.725}{subgroups}%
\TextAt{8}{187.648}{492.725}{of}%
\TextAt{13}{200.978}{492.725}{G}%
\TextAt{8}{210.211}{492.725}{.}%
\TextAt{8}{218.66}{492.725}{The}%
\TextAt{8}{242.713}{492.725}{quotients}%
\TextAt{8}{293.511}{492.725}{of}%
\TextAt{8}{306.841}{492.725}{this}%
\TextAt{8}{329.663}{492.725}{chain}%
\TextAt{8}{360.555}{492.725}{satisfy}%
\TextAt{8}{397.102}{492.725}{the}%
\TextAt{12}{79.226}{467.093}{Proposition}%
\TextAt{12}{152.109}{467.093}{(10.2).}%
\TextAt{8}{194.304}{467.093}{Let}%
\TextAt{13}{215.422}{467.093}{\Character{25}}%
\TextAt{21}{222.062}{465.3}{L}%
\NativeMathAt{231.894}{467.093}{11.955}{\in}%
\TextAt{16}{243.431}{467.093}{\Character{79}}%
\TextAt{8}{257.335}{467.093}{be}%
\TextAt{8}{273.408}{467.093}{a}%
\TextAt{8}{283.312}{467.093}{prime}%
\TextAt{8}{316.625}{467.093}{element}%
\TextAt{8}{360.015}{467.093}{of}%
\TextAt{13}{373.494}{467.093}{L}%
\TextAt{8}{381.459}{467.093}{.}%
\TextAt{8}{390.354}{467.093}{For}%
\TextAt{8}{411.472}{467.093}{every}%
\TextAt{8}{442.513}{467.093}{integer}%
\TextAt{13}{481.361}{467.093}{s}%
\NativeMathAt{490.449}{467.093}{11.955}{\geq}%
\TextAt{9}{503.322}{467.093}{0}%
\TextAt{8}{509.175}{467.093}{,}%
\TextAt{8}{516.51}{467.093}{the}%
\TextAt{8}{79.226}{450.048}{mapping}%
\TextAt{13}{198.644}{433.002}{G}%
\TextAt{21}{207.878}{431.208}{s}%
\TextAt{9}{212.292}{433.002}{\Character{47}}%
\TextAt{13}{218.145}{433.002}{G}%
\TextAt{21}{227.379}{431.208}{s}%
\TextAt{17}{231.294}{431.208}{\Character{43}1}%
\NativeArrowAt{245.934}{433.002}{11.955}{19.26}{\longrightarrow}%
\TextAt{13}{268.516}{433.002}{U}%
\TextAt{17}{277.715}{439.112}{\Character{40}}%
\TextAt{21}{281.008}{439.112}{s}%
\TextAt{17}{284.924}{439.112}{\Character{41}}%
\TextAt{21}{276.466}{429.76}{L}%
\TextAt{9}{288.716}{433.002}{\Character{47}}%
\TextAt{13}{294.569}{433.002}{U}%
\TextAt{17}{303.768}{439.112}{\Character{40}}%
\TextAt{21}{307.061}{439.112}{s}%
\TextAt{17}{310.977}{439.112}{\Character{43}1\Character{41}}%
\TextAt{21}{302.519}{429.76}{L}%
\TextAt{13}{325.589}{433.002}{\Character{59}}%
\TextAt{13}{354.724}{433.002}{\Character{27}}%
\NativeArrowAt{365.147}{433.002}{11.955}{19.26}{\longmapsto}%
\TextAt{13}{388.924}{441.089}{\Character{27}}%
\TextAt{13}{395.99}{441.089}{\Character{25}}%
\TextAt{21}{402.647}{439.296}{L}%
\Vector{q 0 G 0 g 0.478 w 316.924 -338.01 m 336.904 -338.01 l S Q}%
\TextAt{13}{392.466}{424.801}{\Character{25}}%
\TextAt{21}{399.106}{423.008}{L}%
\TextAt{13}{410.1}{433.002}{\Character{59}}%
\TextAt{8}{79.226}{404.02}{is}%
\TextAt{8}{91.371}{404.02}{an}%
\TextAt{8}{108.011}{404.02}{injective}%
\TextAt{8}{154.885}{404.02}{homomorphism}%
\TextAt{8}{237.936}{404.02}{which}%
\TextAt{8}{271.804}{404.02}{is}%
\TextAt{8}{283.948}{404.02}{independent}%
\TextAt{8}{350.657}{404.02}{of}%
\TextAt{8}{364.367}{404.02}{the}%
\TextAt{8}{384.905}{404.02}{prime}%
\TextAt{8}{418.449}{404.02}{element}%
\TextAt{13}{462.072}{404.02}{\Character{25}}%
\TextAt{21}{468.712}{402.226}{L}%
\TextAt{8}{474.97}{404.02}{.}%
\TextAt{8}{484.561}{404.02}{Here}%
\TextAt{13}{512.57}{404.02}{U}%
\TextAt{17}{521.769}{410.13}{\Character{40}}%
\TextAt{21}{525.062}{410.13}{s}%
\TextAt{17}{528.978}{410.13}{\Character{41}}%
\TextAt{21}{520.52}{400.779}{L}%
\TextAt{8}{79.226}{386.974}{denotes}%
\TextAt{8}{122.194}{386.974}{the}%
\TextAt{13}{142.985}{386.974}{s}%
\TextAt{8}{148.499}{386.974}{-th}%
\TextAt{8}{167.988}{386.974}{group}%
\TextAt{8}{201.798}{386.974}{of}%
\TextAt{8}{215.763}{386.974}{principal}%
\TextAt{8}{265.175}{386.974}{units}%
\TextAt{8}{295.136}{386.974}{of}%
\TextAt{13}{309.1}{386.974}{L}%
\TextAt{8}{317.065}{386.974}{,}%
\TextAt{8}{325.008}{386.974}{i.e.,}%
\TextAt{13}{347.906}{386.974}{U}%
\TextAt{17}{357.105}{393.084}{\Character{40}0\Character{41}}%
\TextAt{21}{355.857}{383.733}{L}%
\TextAt{9}{372.827}{386.974}{\Character{61}}%
\TextAt{16}{386.334}{386.974}{\Character{79}}%
\TextAt{3}{396.184}{391.312}{\Character{3}}%
\TextAt{8}{405.449}{386.974}{and}%
\TextAt{13}{428.846}{386.974}{U}%
\TextAt{17}{438.045}{393.084}{\Character{40}}%
\TextAt{21}{441.338}{393.084}{s}%
\TextAt{17}{445.254}{393.084}{\Character{41}}%
\TextAt{21}{436.796}{383.733}{L}%
\TextAt{9}{453.448}{386.974}{\Character{61}}%
\TextAt{9}{466.945}{386.974}{1}%
\TextAt{9}{475.876}{386.974}{\Character{43}}%
\TextAt{13}{488.092}{386.974}{\Character{25}}%
\TextAt{21}{495.161}{391.312}{s}%
\TextAt{21}{494.732}{383.953}{L}%
\TextAt{16}{500.989}{386.974}{\Character{79}}%
\TextAt{8}{510.839}{386.974}{,}%
\TextAt{8}{518.782}{386.974}{for}%
\TextAt{13}{79.226}{369.928}{s}%
\NativeMathAt{88.061}{369.928}{11.955}{\geq}%
\TextAt{9}{100.68}{369.928}{1}%
\TextAt{8}{106.533}{369.928}{.}%
\TextAt{8}{97.159}{344.297}{We}%
\TextAt{8}{116.604}{344.297}{leave}%
\TextAt{8}{144.849}{344.297}{the}%
\TextAt{8}{164.306}{344.297}{elementary}%
\TextAt{8}{223.443}{344.297}{proof}%
\TextAt{8}{253.314}{344.297}{to}%
\TextAt{8}{266.925}{344.297}{the}%
\TextAt{8}{286.382}{344.297}{reader.}%
\TextAt{8}{326.47}{344.297}{Observe}%
\TextAt{8}{370.711}{344.297}{that}%
\TextAt{8}{395.381}{344.297}{one}%
\TextAt{8}{416.141}{344.297}{has}%
\TextAt{13}{436.316}{344.297}{U}%
\TextAt{17}{445.515}{350.407}{\Character{40}0\Character{41}}%
\TextAt{21}{444.266}{341.056}{L}%
\TextAt{9}{456.834}{344.297}{\Character{47}}%
\TextAt{13}{462.687}{344.297}{U}%
\TextAt{17}{471.886}{350.407}{\Character{40}1\Character{41}}%
\TextAt{21}{470.637}{341.056}{L}%
\NativeMathAt{486.526}{344.297}{11.955}{\cong}%
\TextAt{13}{499.145}{344.297}{\Character{21}}%
\TextAt{3}{505.973}{348.635}{\Character{3}}%
\TextAt{8}{513.904}{344.297}{and}%
\TextAt{13}{79.226}{327.251}{U}%
\TextAt{17}{88.425}{333.361}{\Character{40}}%
\TextAt{21}{91.718}{333.361}{s}%
\TextAt{17}{95.634}{333.361}{\Character{41}}%
\TextAt{21}{87.176}{324.01}{L}%
\TextAt{9}{99.426}{327.251}{\Character{47}}%
\TextAt{13}{105.279}{327.251}{U}%
\TextAt{17}{114.478}{333.361}{\Character{40}}%
\TextAt{21}{117.771}{333.361}{s}%
\TextAt{17}{121.687}{333.361}{\Character{43}1\Character{41}}%
\TextAt{21}{113.229}{324.01}{L}%
\NativeMathAt{139.62}{327.251}{11.955}{\cong}%
\TextAt{13}{152.239}{327.251}{\Character{21}}%
\TextAt{8}{159.068}{327.251}{,}%
\TextAt{8}{165.901}{327.251}{for}%
\TextAt{13}{183.39}{327.251}{s}%
\NativeMathAt{192.225}{327.251}{11.955}{\geq}%
\TextAt{9}{204.845}{327.251}{1}%
\TextAt{8}{210.697}{327.251}{.}%
\TextAt{8}{219.014}{327.251}{The}%
\TextAt{8}{242.673}{327.251}{factors}%
\TextAt{13}{280.39}{327.251}{G}%
\TextAt{21}{289.624}{325.457}{s}%
\TextAt{9}{294.038}{327.251}{\Character{47}}%
\TextAt{13}{299.891}{327.251}{G}%
\TextAt{21}{309.125}{325.457}{s}%
\TextAt{17}{313.041}{325.457}{\Character{43}1}%
\TextAt{8}{327.862}{327.251}{are}%
\TextAt{8}{346.977}{327.251}{therefore}%
\TextAt{8}{395.682}{327.251}{abelian}%
\TextAt{8}{435.934}{327.251}{groups}%
\TextAt{8}{473.329}{327.251}{of}%
\TextAt{8}{486.264}{327.251}{exponent}%
\TextAt{13}{79.226}{310.205}{p}%
\TextAt{8}{85.101}{310.205}{,}%
\TextAt{8}{93.286}{310.205}{for}%
\TextAt{13}{112}{310.205}{s}%
\NativeMathAt{122.246}{310.205}{11.955}{\geq}%
\TextAt{9}{136.278}{310.205}{1}%
\TextAt{8}{142.131}{310.205}{,}%
\TextAt{8}{150.315}{310.205}{and}%
\TextAt{8}{173.907}{310.205}{of}%
\TextAt{8}{188.065}{310.205}{order}%
\TextAt{8}{219.463}{310.205}{prime}%
\TextAt{8}{253.455}{310.205}{to}%
\TextAt{13}{268.594}{310.205}{p}%
\TextAt{8}{274.469}{310.205}{,}%
\TextAt{8}{282.654}{310.205}{for}%
\TextAt{13}{301.367}{310.205}{s}%
\TextAt{9}{311.614}{310.205}{\Character{61}}%
\TextAt{9}{325.434}{310.205}{0}%
\TextAt{8}{331.304}{310.205}{.}%
\TextAt{8}{342.238}{310.205}{In}%
\TextAt{8}{357.688}{310.205}{particular,}%
\TextAt{8}{415.965}{310.205}{we}%
\TextAt{8}{434.021}{310.205}{find}%
\TextAt{8}{458.258}{310.205}{again}%
\TextAt{8}{490.313}{310.205}{that}%
\TextAt{8}{516.51}{310.205}{the}%
\TextAt{8}{79.226}{293.159}{ramification}%
\TextAt{8}{144.274}{293.159}{group}%
\TextAt{13}{177.45}{293.159}{R}%
\TextAt{9}{189.779}{293.159}{\Character{61}}%
\TextAt{13}{202.205}{293.159}{G}%
\TextAt{17}{211.438}{291.365}{1}%
\TextAt{8}{220.068}{293.159}{is}%
\TextAt{8}{231.832}{293.159}{the}%
\TextAt{8}{251.988}{293.159}{unique}%
\TextAt{13}{290.03}{293.159}{p}%
\TextAt{8}{295.905}{293.159}{-Sylow}%
\TextAt{8}{333.623}{293.159}{subgroup}%
\TextAt{8}{384.421}{293.159}{in}%
\TextAt{8}{398.074}{293.159}{the}%
\TextAt{8}{418.23}{293.159}{inertia}%
\TextAt{8}{455.303}{293.159}{group}%
\TextAt{13}{488.479}{293.159}{I}%
\TextAt{9}{497.902}{293.159}{\Character{61}}%
\TextAt{13}{510.328}{293.159}{G}%
\TextAt{17}{519.561}{291.365}{0}%
\TextAt{8}{524.294}{293.159}{.}%
\TextAt{8}{97.159}{273.124}{We}%
\TextAt{8}{116.16}{273.124}{now}%
\TextAt{8}{139.404}{273.124}{study}%
\TextAt{8}{170.516}{273.124}{the}%
\TextAt{8}{189.528}{273.124}{behaviour}%
\TextAt{8}{242.697}{273.124}{of}%
\TextAt{8}{254.883}{273.124}{the}%
\TextAt{8}{273.895}{273.124}{higher}%
\TextAt{8}{308.522}{273.124}{ramification}%
\TextAt{8}{372.426}{273.124}{groups}%
\TextAt{8}{409.072}{273.124}{under}%
\TextAt{8}{441.092}{273.124}{change}%
\TextAt{8}{478.647}{273.124}{of}%
\TextAt{8}{490.833}{273.124}{fields.}%
\TextAt{8}{524.975}{273.124}{If}%
\TextAt{8}{79.226}{256.078}{only}%
\TextAt{8}{103.707}{256.078}{the}%
\TextAt{8}{122.653}{256.078}{base}%
\TextAt{8}{147.517}{256.078}{field}%
\TextAt{13}{171.663}{256.078}{K}%
\TextAt{8}{185.168}{256.078}{is}%
\TextAt{8}{195.721}{256.078}{changed,}%
\TextAt{8}{243.207}{256.078}{then}%
\TextAt{8}{268.656}{256.078}{we}%
\TextAt{8}{284.673}{256.078}{get}%
\TextAt{8}{302.974}{256.078}{directly}%
\TextAt{8}{344.359}{256.078}{from}%
\TextAt{8}{370.789}{256.078}{the}%
\TextAt{8}{389.735}{256.078}{definition}%
\TextAt{8}{440.553}{256.078}{of}%
\TextAt{8}{452.673}{256.078}{the}%
\TextAt{8}{471.619}{256.078}{ramification}%
\TextAt{8}{79.226}{239.032}{groups}%
\TextAt{8}{117.016}{239.032}{the}%
\TextAt{8}{137.173}{239.032}{following}%
\TextAt{8}{186.607}{239.032}{generalization}%
\TextAt{8}{261.411}{239.032}{of}%
\TextAt{8}{274.741}{239.032}{(9.5).}%
\TextAt{12}{79.226}{213.401}{Proposition}%
\TextAt{12}{151.373}{213.401}{(10.3).}%
\TextAt{8}{192.93}{213.401}{If}%
\TextAt{13}{204.131}{213.401}{K}%
\TextAt{3}{214.949}{217.739}{\Character{48}}%
\TextAt{8}{221.151}{213.401}{is}%
\TextAt{8}{232.425}{213.401}{an}%
\TextAt{8}{248.194}{213.401}{intermediate}%
\TextAt{8}{315.668}{213.401}{field}%
\TextAt{8}{340.535}{213.401}{of}%
\TextAt{13}{353.375}{213.401}{L}%
\TextAt{16}{364.66}{213.401}{\Character{106}}%
\TextAt{13}{371.302}{213.401}{K}%
\TextAt{8}{382.119}{213.401}{,}%
\TextAt{8}{388.876}{213.401}{then}%
\TextAt{8}{415.046}{213.401}{one}%
\TextAt{8}{436.015}{213.401}{has,}%
\TextAt{8}{459.749}{213.401}{for}%
\TextAt{8}{477.143}{213.401}{all}%
\TextAt{13}{492.912}{213.401}{s}%
\NativeMathAt{501.747}{213.401}{11.955}{\geq}%
\TextAt{16}{514.348}{213.401}{\Character{0}}%
\TextAt{9}{523.665}{213.401}{1}%
\TextAt{8}{529.518}{213.401}{,}%
\TextAt{8}{79.226}{196.355}{that}%
\TextAt{13}{212.023}{179.309}{G}%
\TextAt{21}{221.257}{177.516}{s}%
\TextAt{9}{225.671}{179.309}{\Character{40}}%
\TextAt{13}{230.223}{179.309}{L}%
\TextAt{16}{241.509}{179.309}{\Character{106}}%
\TextAt{13}{248.15}{179.309}{K}%
\TextAt{3}{258.968}{184.245}{\Character{48}}%
\TextAt{9}{261.763}{179.309}{\Character{41}}%
\TextAt{9}{269.618}{179.309}{\Character{61}}%
\TextAt{13}{282.062}{179.309}{G}%
\TextAt{21}{291.296}{177.516}{s}%
\TextAt{9}{295.71}{179.309}{\Character{40}}%
\TextAt{13}{300.262}{179.309}{L}%
\TextAt{16}{311.547}{179.309}{\Character{106}}%
\TextAt{13}{318.189}{179.309}{K}%
\TextAt{9}{329.007}{179.309}{\Character{41}}%
\TextAt{16}{336.216}{179.309}{\Character{92}}%
\TextAt{13}{346.842}{179.309}{G}%
\TextAt{9}{356.076}{179.309}{\Character{40}}%
\TextAt{13}{360.628}{179.309}{L}%
\TextAt{16}{371.914}{179.309}{\Character{106}}%
\TextAt{13}{378.555}{179.309}{K}%
\TextAt{3}{389.373}{184.245}{\Character{48}}%
\TextAt{9}{392.168}{179.309}{\Character{41}}%
\TextAt{13}{396.72}{179.309}{\Character{58}}%
\TextAt{8}{297.211}{43.033}{162}%
\Destination{78.23}{809.815}{nav.page.162}%
\Destination{79.36}{635.455}{nav.print.177}%
\Destination{79.36}{715.355}{nav.number.2.10.1}%
\Destination{79.36}{481.055}{nav.number.2.10.2}%
\Destination{79.36}{227.355}{nav.number.2.10.3}%
\Hotlink{274.241}{236.855}{25.078}{11.047}{nav.number.2.9.5}%
%
%
%
%
\Destination{79.226}{634.455}{errata.21}%
}

% Page 175 (body page 163)
\Page{%
\TextAt{2}{79.226}{796.943}{II.}%
\TextAt{2}{94.073}{796.943}{The}%
\TextAt{2}{116.492}{796.943}{Theory}%
\TextAt{2}{154.401}{796.943}{of}%
\TextAt{2}{166.826}{796.943}{Valuations}%
\Vector{q 0 G 0 g 0.299 w 7.226 18.729 m 460.769 18.729 l S Q}%
\TextAt{8}{97.159}{761.124}{Matters}%
\TextAt{8}{140.921}{761.124}{become}%
\TextAt{8}{182.663}{761.124}{much}%
\TextAt{8}{213.68}{761.124}{more}%
\TextAt{8}{242.749}{761.124}{complicated}%
\TextAt{8}{307.588}{761.124}{when}%
\TextAt{8}{337.949}{761.124}{we}%
\TextAt{8}{354.979}{761.124}{pass}%
\TextAt{8}{380.27}{761.124}{from}%
\TextAt{13}{407.713}{761.124}{L}%
\TextAt{16}{418.999}{761.124}{\Character{106}}%
\TextAt{13}{425.64}{761.124}{K}%
\TextAt{8}{440.158}{761.124}{to}%
\TextAt{8}{454.272}{761.124}{a}%
\TextAt{8}{463.83}{761.124}{Galois}%
\TextAt{8}{499.546}{761.124}{subex-}%
\TextAt{8}{79.226}{744.078}{tension}%
\TextAt{13}{121.278}{744.078}{L}%
\TextAt{3}{129.242}{748.416}{\Character{48}}%
\TextAt{16}{138.199}{744.078}{\Character{106}}%
\TextAt{13}{147.681}{744.078}{K}%
\TextAt{8}{158.498}{744.078}{.}%
\TextAt{8}{171.947}{744.078}{It}%
\TextAt{8}{186.287}{744.078}{is}%
\TextAt{8}{199.717}{744.078}{true}%
\TextAt{8}{226.096}{744.078}{that}%
\TextAt{8}{253.131}{744.078}{the}%
\TextAt{8}{274.955}{744.078}{ramification}%
\TextAt{8}{341.67}{744.078}{groups}%
\TextAt{8}{381.127}{744.078}{of}%
\TextAt{13}{396.124}{744.078}{L}%
\TextAt{16}{410.249}{744.078}{\Character{106}}%
\TextAt{13}{419.732}{744.078}{K}%
\TextAt{8}{436.114}{744.078}{are}%
\TextAt{8}{457.291}{744.078}{mapped}%
\TextAt{8}{503.503}{744.078}{under}%
\TextAt{13}{79.226}{727.032}{G}%
\TextAt{9}{88.46}{727.032}{\Character{40}}%
\TextAt{13}{93.012}{727.032}{L}%
\TextAt{16}{104.575}{727.032}{\Character{106}}%
\TextAt{13}{111.495}{727.032}{K}%
\TextAt{9}{122.313}{727.032}{\Character{41}}%
\NativeMathAt{130.464}{727.032}{11.955}{\rightarrow}%
\TextAt{13}{146.018}{727.032}{G}%
\TextAt{9}{155.252}{727.032}{\Character{40}}%
\TextAt{13}{159.804}{727.032}{L}%
\TextAt{3}{167.768}{731.37}{\Character{48}}%
\TextAt{16}{174.162}{727.032}{\Character{106}}%
\TextAt{13}{181.082}{727.032}{K}%
\TextAt{9}{191.9}{727.032}{\Character{41}}%
\TextAt{8}{200.513}{727.032}{into}%
\TextAt{8}{224.419}{727.032}{the}%
\TextAt{8}{244.738}{727.032}{ramification}%
\TextAt{8}{309.949}{727.032}{groups}%
\TextAt{8}{347.903}{727.032}{of}%
\TextAt{13}{361.396}{727.032}{L}%
\TextAt{3}{369.361}{731.37}{\Character{48}}%
\TextAt{16}{375.754}{727.032}{\Character{106}}%
\TextAt{13}{382.674}{727.032}{K}%
\TextAt{8}{393.492}{727.032}{,}%
\TextAt{8}{400.845}{727.032}{but}%
\TextAt{8}{422.468}{727.032}{the}%
\TextAt{8}{442.787}{727.032}{indexing}%
\TextAt{8}{490.101}{727.032}{changes.}%
\TextAt{8}{79.226}{709.986}{For}%
\TextAt{8}{99.713}{709.986}{the}%
\TextAt{8}{119.388}{709.986}{precise}%
\TextAt{8}{157.33}{709.986}{description}%
\TextAt{8}{216.743}{709.986}{of}%
\TextAt{8}{229.591}{709.986}{the}%
\TextAt{8}{249.265}{709.986}{situation}%
\TextAt{8}{297.632}{709.986}{we}%
\TextAt{8}{314.377}{709.986}{need}%
\TextAt{8}{341.201}{709.986}{some}%
\TextAt{8}{370.045}{709.986}{preparation.}%
\TextAt{8}{437.534}{709.986}{We}%
\TextAt{8}{457.197}{709.986}{will}%
\TextAt{8}{478.82}{709.986}{assume}%
\TextAt{8}{518.782}{709.986}{for}%
\TextAt{8}{79.226}{692.94}{the}%
\TextAt{8}{99.383}{692.94}{sequel}%
\TextAt{8}{134.232}{692.94}{that}%
\TextAt{8}{159.601}{692.94}{the}%
\TextAt{8}{179.757}{692.94}{residue}%
\TextAt{8}{219.484}{692.94}{field}%
\TextAt{8}{244.841}{692.94}{extension}%
\TextAt{13}{296.607}{692.94}{\Character{21}}%
\TextAt{16}{306.756}{692.94}{\Character{106}}%
\TextAt{13}{313.398}{692.94}{\Character{20}}%
\TextAt{8}{324.034}{692.94}{of}%
\TextAt{13}{337.364}{692.94}{L}%
\TextAt{16}{348.649}{692.94}{\Character{106}}%
\TextAt{13}{355.291}{692.94}{K}%
\TextAt{8}{370.006}{692.94}{is}%
\TextAt{12}{381.77}{692.94}{separable}%
\TextAt{8}{436.608}{692.94}{.}%
\TextAt{12}{79.226}{666.94}{Lemma}%
\TextAt{12}{127.783}{666.94}{(10.4).}%
\TextAt{8}{170.604}{666.94}{The}%
\TextAt{8}{195.433}{666.94}{ring}%
\TextAt{8}{220.274}{666.94}{extension}%
\TextAt{16}{272.815}{666.94}{\Character{79}}%
\TextAt{8}{287.337}{666.94}{of}%
\TextAt{14}{301.442}{666.94}{o}%
\TextAt{8}{311.959}{666.94}{is}%
\TextAt{8}{324.498}{666.94}{monogenous,}%
\TextAt{8}{395.129}{666.94}{i.e.,}%
\TextAt{8}{418.203}{666.94}{there}%
\TextAt{8}{448.89}{666.94}{exists}%
\TextAt{8}{481.98}{666.94}{an}%
\TextAt{13}{499.014}{666.94}{x}%
\NativeMathAt{510.308}{666.94}{11.955}{\in}%
\TextAt{16}{522.909}{666.94}{\Character{79}}%
\TextAt{8}{79.226}{649.894}{such}%
\TextAt{8}{105.623}{649.894}{that}%
\TextAt{16}{130.992}{649.894}{\Character{79}}%
\TextAt{9}{144.163}{649.894}{\Character{61}}%
\TextAt{14}{156.588}{649.894}{o}%
\TextAt{9}{162.432}{649.894}{\Character{91}}%
\TextAt{13}{165.684}{649.894}{x}%
\TextAt{9}{172.336}{649.894}{\Character{93}}%
\TextAt{8}{175.588}{649.894}{.}%
\TextAt{19}{79.226}{614.94}{Proof.}%
\TextAt{8}{116}{614.94}{As}%
\TextAt{8}{132.064}{614.94}{the}%
\TextAt{8}{150.996}{614.94}{residue}%
\TextAt{8}{189.5}{614.94}{field}%
\TextAt{8}{213.633}{614.94}{extension}%
\TextAt{13}{264.175}{614.94}{\Character{21}}%
\TextAt{16}{274.325}{614.94}{\Character{106}}%
\TextAt{13}{280.966}{614.94}{\Character{20}}%
\TextAt{8}{290.379}{614.94}{is}%
\TextAt{8}{300.919}{614.94}{separable}%
\TextAt{8}{351.138}{614.94}{by}%
\TextAt{8}{366.174}{614.94}{assumption,}%
\TextAt{8}{430.362}{614.94}{it}%
\TextAt{8}{440.843}{614.94}{admits}%
\TextAt{8}{478.055}{614.94}{a}%
\TextAt{8}{486.587}{614.94}{primitive}%
\TextAt{8}{79.226}{597.894}{element}%
\TextAt{9}{122.525}{597.894}{\Character{22}}%
\TextAt{13}{121.801}{597.894}{x}%
\TextAt{8}{128.453}{597.894}{.}%
\TextAt{8}{136.679}{597.894}{Let}%
\TextAt{13}{156.981}{597.894}{f}%
\TextAt{9}{164.027}{597.894}{\Character{40}}%
\TextAt{13}{168.579}{597.894}{X}%
\TextAt{9}{179.234}{597.894}{\Character{41}}%
\NativeMathAt{187.108}{597.894}{11.955}{\in}%
\TextAt{14}{198.399}{597.894}{o}%
\TextAt{9}{204.243}{597.894}{\Character{91}}%
\TextAt{13}{207.494}{597.894}{X}%
\TextAt{9}{218.15}{597.894}{\Character{93}}%
\TextAt{8}{224.631}{597.894}{be}%
\TextAt{8}{239.888}{597.894}{a}%
\TextAt{8}{248.976}{597.894}{lifting}%
\TextAt{8}{282.453}{597.894}{of}%
\TextAt{8}{295.115}{597.894}{the}%
\TextAt{8}{314.605}{597.894}{minimal}%
\TextAt{8}{359.462}{597.894}{polynomial}%
\TextAt{9}{421.836}{601.049}{\Character{22}}%
\TextAt{13}{419.288}{597.894}{f}%
\TextAt{9}{426.335}{597.894}{\Character{40}}%
\TextAt{13}{430.887}{597.894}{X}%
\TextAt{9}{441.542}{597.894}{\Character{41}}%
\TextAt{8}{449.324}{597.894}{of}%
\TextAt{9}{462.712}{597.894}{\Character{22}}%
\TextAt{13}{461.987}{597.894}{x}%
\TextAt{8}{468.639}{597.894}{.}%
\TextAt{8}{476.865}{597.894}{Then}%
\TextAt{8}{506.755}{597.894}{there}%
\TextAt{8}{79.226}{580.848}{is}%
\TextAt{8}{91.314}{580.848}{a}%
\TextAt{8}{101.394}{580.848}{representative}%
\TextAt{13}{176.905}{580.848}{x}%
\NativeMathAt{187.43}{580.848}{11.955}{\in}%
\TextAt{16}{199.266}{580.848}{\Character{79}}%
\TextAt{8}{213.346}{580.848}{of}%
\TextAt{9}{227.725}{580.848}{\Character{22}}%
\TextAt{13}{227}{580.848}{x}%
\TextAt{8}{237.874}{580.848}{such}%
\TextAt{8}{264.595}{580.848}{that}%
\TextAt{13}{290.289}{580.848}{\Character{25}}%
\TextAt{9}{301.231}{580.848}{\Character{61}}%
\TextAt{13}{314.21}{580.848}{f}%
\TextAt{9}{321.256}{580.848}{\Character{40}}%
\TextAt{13}{325.808}{580.848}{x}%
\TextAt{9}{332.461}{580.848}{\Character{41}}%
\TextAt{8}{341.235}{580.848}{is}%
\TextAt{8}{353.323}{580.848}{a}%
\TextAt{8}{363.403}{580.848}{prime}%
\TextAt{8}{396.891}{580.848}{element}%
\TextAt{8}{440.457}{580.848}{of}%
\TextAt{16}{454.111}{580.848}{\Character{79}}%
\TextAt{8}{463.961}{580.848}{.}%
\TextAt{8}{473.383}{580.848}{Indeed,}%
\TextAt{8}{515.069}{580.848}{if}%
\TextAt{13}{526.117}{580.848}{x}%
\TextAt{8}{79.226}{563.802}{is}%
\TextAt{8}{91.215}{563.802}{an}%
\TextAt{8}{107.699}{563.802}{arbitrary}%
\TextAt{8}{157.693}{563.802}{representative,}%
\TextAt{8}{236.412}{563.802}{then}%
\TextAt{8}{263.298}{563.802}{we}%
\TextAt{8}{280.75}{563.802}{certainly}%
\TextAt{8}{329.429}{563.802}{have}%
\TextAt{13}{356.649}{563.802}{v}%
\TextAt{21}{362.308}{562.009}{L}%
\TextAt{9}{368.565}{563.802}{\Character{40}}%
\TextAt{13}{373.118}{563.802}{f}%
\TextAt{9}{380.164}{563.802}{\Character{40}}%
\TextAt{13}{384.716}{563.802}{x}%
\TextAt{9}{391.369}{563.802}{\Character{41}\Character{41}}%
\NativeMathAt{404.177}{563.802}{11.955}{\geq}%
\TextAt{9}{417.18}{563.802}{1}%
\TextAt{8}{427.156}{563.802}{because}%
\TextAt{9}{473.23}{566.957}{\Character{22}}%
\TextAt{13}{470.682}{563.802}{f}%
\TextAt{9}{477.729}{563.802}{\Character{40}}%
\TextAt{9}{482.989}{563.802}{\Character{22}}%
\TextAt{13}{482.281}{563.802}{x}%
\TextAt{9}{488.933}{563.802}{\Character{41}}%
\TextAt{9}{497.17}{563.802}{\Character{61}}%
\TextAt{9}{509.974}{563.802}{0}%
\TextAt{8}{515.851}{563.802}{.}%
\TextAt{8}{524.975}{563.802}{If}%
\TextAt{13}{79.226}{546.756}{x}%
\TextAt{8}{89.766}{546.756}{itself}%
\TextAt{8}{118.101}{546.756}{is}%
\TextAt{8}{129.855}{546.756}{not}%
\TextAt{8}{150.659}{546.756}{as}%
\TextAt{8}{165.019}{546.756}{required,}%
\TextAt{8}{214.111}{546.756}{i.e.,}%
\TextAt{8}{236.208}{546.756}{if}%
\TextAt{13}{246.922}{546.756}{v}%
\TextAt{21}{252.581}{544.963}{L}%
\TextAt{9}{258.838}{546.756}{\Character{40}}%
\TextAt{13}{263.391}{546.756}{f}%
\TextAt{9}{270.437}{546.756}{\Character{40}}%
\TextAt{13}{274.989}{546.756}{x}%
\TextAt{9}{281.641}{546.756}{\Character{41}\Character{41}}%
\NativeMathAt{294.067}{546.756}{11.955}{\geq}%
\TextAt{9}{306.686}{546.756}{2}%
\TextAt{8}{312.539}{546.756}{,}%
\TextAt{8}{319.68}{546.756}{the}%
\TextAt{8}{339.827}{546.756}{representative}%
\TextAt{13}{415.003}{546.756}{x}%
\TextAt{9}{424.291}{546.756}{\Character{43}}%
\TextAt{13}{436.032}{546.756}{\Character{25}}%
\TextAt{21}{442.672}{544.963}{L}%
\TextAt{8}{452.817}{546.756}{will}%
\TextAt{8}{474.912}{546.756}{do.}%
\TextAt{8}{495.718}{546.756}{In}%
\TextAt{8}{510.33}{546.756}{fact,}%
\TextAt{8}{79.226}{529.71}{from}%
\TextAt{8}{106.866}{529.71}{Taylor’s}%
\TextAt{8}{151.483}{529.71}{formula}%
\TextAt{13}{188.138}{500.753}{f}%
\TextAt{9}{195.185}{500.753}{\Character{40}}%
\TextAt{13}{199.737}{500.753}{x}%
\TextAt{9}{209.046}{500.753}{\Character{43}}%
\TextAt{13}{220.807}{500.753}{\Character{25}}%
\TextAt{21}{227.447}{498.96}{L}%
\TextAt{9}{233.705}{500.753}{\Character{41}}%
\TextAt{9}{241.56}{500.753}{\Character{61}}%
\TextAt{13}{254.003}{500.753}{f}%
\TextAt{9}{261.05}{500.753}{\Character{40}}%
\TextAt{13}{265.602}{500.753}{x}%
\TextAt{9}{272.254}{500.753}{\Character{41}}%
\TextAt{9}{279.451}{500.753}{\Character{43}}%
\TextAt{13}{291.224}{500.753}{f}%
\TextAt{3}{298.271}{505.689}{\Character{48}}%
\TextAt{9}{301.066}{500.753}{\Character{40}}%
\TextAt{13}{305.618}{500.753}{x}%
\TextAt{9}{312.27}{500.753}{\Character{41}}%
\TextAt{13}{316.823}{500.753}{\Character{25}}%
\TextAt{21}{323.463}{498.96}{L}%
\TextAt{9}{332.377}{500.753}{\Character{43}}%
\TextAt{13}{344.138}{500.753}{b\Character{25}}%
\TextAt{17}{356.184}{505.689}{2}%
\TextAt{21}{355.755}{497.798}{L}%
\TextAt{13}{362.013}{500.753}{\Character{59}}%
\TextAt{13}{391.148}{500.753}{b}%
\NativeMathAt{399.465}{500.753}{11.955}{\in}%
\TextAt{16}{410.739}{500.753}{\Character{79}}%
\TextAt{13}{420.606}{500.753}{\Character{59}}%
\TextAt{8}{79.226}{471.796}{we}%
\TextAt{8}{96.763}{471.796}{obtain}%
\TextAt{13}{133.5}{471.796}{v}%
\TextAt{21}{139.159}{470.002}{L}%
\TextAt{9}{145.416}{471.796}{\Character{40}}%
\TextAt{13}{149.969}{471.796}{f}%
\TextAt{9}{157.015}{471.796}{\Character{40}}%
\TextAt{13}{161.568}{471.796}{x}%
\TextAt{9}{171.087}{471.796}{\Character{43}}%
\TextAt{13}{183.06}{471.796}{\Character{25}}%
\TextAt{21}{189.7}{470.002}{L}%
\TextAt{9}{195.957}{471.796}{\Character{41}\Character{41}}%
\TextAt{9}{208.881}{471.796}{\Character{61}}%
\TextAt{9}{221.828}{471.796}{1}%
\TextAt{8}{231.923}{471.796}{since}%
\TextAt{13}{260.901}{471.796}{f}%
\TextAt{3}{267.947}{476.134}{\Character{48}}%
\TextAt{9}{270.743}{471.796}{\Character{40}}%
\TextAt{13}{275.295}{471.796}{x}%
\TextAt{9}{281.947}{471.796}{\Character{41}}%
\NativeMathAt{290.348}{471.796}{11.955}{\in}%
\TextAt{16}{302.148}{471.796}{\Character{79}}%
\TextAt{3}{312.016}{476.134}{\Character{3}}%
\TextAt{8}{316.748}{471.796}{,}%
\TextAt{8}{324.284}{471.796}{because}%
\TextAt{9}{370.443}{474.951}{\Character{22}}%
\TextAt{13}{367.895}{471.796}{f}%
\TextAt{3}{374.942}{476.134}{\Character{48}}%
\TextAt{9}{377.737}{471.796}{\Character{40}}%
\TextAt{9}{382.997}{471.796}{\Character{22}}%
\TextAt{13}{382.289}{471.796}{x}%
\TextAt{9}{388.941}{471.796}{\Character{41}}%
\NativeMathAt{397.342}{471.796}{11.955}{\neq}%
\TextAt{9}{410.278}{471.796}{0}%
\TextAt{8}{416.148}{471.796}{.}%
\TextAt{8}{425.525}{471.796}{In}%
\TextAt{8}{440.456}{471.796}{the}%
\TextAt{8}{460.922}{471.796}{proof}%
\TextAt{8}{491.8}{471.796}{of}%
\TextAt{8}{505.44}{471.796}{(6.8),}%
\TextAt{8}{79.226}{454.75}{we}%
\TextAt{8}{96.454}{454.75}{saw}%
\TextAt{8}{118.953}{454.75}{that}%
\TextAt{8}{144.322}{454.75}{the}%
\TextAt{13}{164.695}{425.792}{x}%
\TextAt{21}{171.347}{430.729}{j}%
\TextAt{13}{175.729}{425.792}{\Character{25}}%
\TextAt{21}{182.798}{430.729}{i}%
\TextAt{9}{189.5}{425.792}{\Character{61}}%
\TextAt{13}{201.926}{425.792}{x}%
\TextAt{21}{208.578}{430.729}{j}%
\TextAt{13}{212.96}{425.792}{f}%
\TextAt{9}{220.007}{425.792}{\Character{40}}%
\TextAt{13}{224.559}{425.792}{x}%
\TextAt{9}{231.211}{425.792}{\Character{41}}%
\TextAt{21}{235.763}{430.729}{i}%
\TextAt{13}{239.145}{425.792}{\Character{59}}%
\TextAt{13}{268.28}{425.792}{j}%
\TextAt{9}{277.135}{425.792}{\Character{61}}%
\TextAt{9}{289.545}{425.792}{0}%
\TextAt{13}{295.414}{425.792}{\Character{59}}%
\TextAt{13}{300.638}{425.792}{\Character{58}}%
\TextAt{13}{305.875}{425.792}{\Character{58}}%
\TextAt{13}{311.099}{425.792}{\Character{58}}%
\TextAt{13}{316.336}{425.792}{\Character{59}}%
\TextAt{13}{321.572}{425.792}{f}%
\TextAt{16}{331.338}{425.792}{\Character{0}}%
\TextAt{9}{343.293}{425.792}{1}%
\TextAt{13}{349.146}{425.792}{\Character{59}}%
\TextAt{13}{366.326}{425.792}{i}%
\TextAt{9}{373.66}{425.792}{\Character{61}}%
\TextAt{9}{386.07}{425.792}{0}%
\TextAt{13}{391.938}{425.792}{\Character{59}}%
\TextAt{13}{397.162}{425.792}{\Character{58}}%
\TextAt{13}{402.399}{425.792}{\Character{58}}%
\TextAt{13}{407.623}{425.792}{\Character{58}}%
\TextAt{13}{412.86}{425.792}{\Character{59}}%
\TextAt{13}{418.096}{425.792}{e}%
\TextAt{16}{426.241}{425.792}{\Character{0}}%
\TextAt{9}{438.196}{425.792}{1}%
\TextAt{13}{444.049}{425.792}{\Character{59}}%
\TextAt{8}{79.226}{396.835}{form}%
\TextAt{8}{106.866}{396.835}{an}%
\TextAt{8}{123.126}{396.835}{integral}%
\TextAt{8}{165.734}{396.835}{basis}%
\TextAt{8}{194.474}{396.835}{of}%
\TextAt{16}{207.804}{396.835}{\Character{79}}%
\TextAt{8}{221.551}{396.835}{over}%
\TextAt{14}{246.597}{396.835}{o}%
\TextAt{8}{252.441}{396.835}{.}%
\TextAt{8}{260.89}{396.835}{Hence}%
\TextAt{8}{295.667}{396.835}{indeed}%
\TextAt{16}{332.728}{396.835}{\Character{79}}%
\TextAt{9}{345.899}{396.835}{\Character{61}}%
\TextAt{14}{358.325}{396.835}{o}%
\TextAt{9}{364.169}{396.835}{\Character{91}}%
\TextAt{13}{367.42}{396.835}{x}%
\TextAt{9}{374.072}{396.835}{\Character{93}}%
\TextAt{8}{377.324}{396.835}{.}%
\Vector{q 0 G 0 g 0.399 w 452.334 -377.165 m 452.334 -369.095 l S Q}%
\Vector{q 0 G 0 g 0.399 w 452.534 -369.294 m 459.707 -369.294 l S Q}%
\Vector{q 0 G 0 g 0.399 w 452.534 -376.966 m 459.707 -376.966 l S Q}%
\Vector{q 0 G 0 g 0.399 w 459.906 -377.165 m 459.906 -369.095 l S Q}%
\TextAt{8}{97.159}{367.846}{For}%
\TextAt{8}{118.128}{367.846}{every}%
\TextAt{13}{149.02}{367.846}{\Character{27}}%
\NativeMathAt{159.424}{367.846}{11.955}{\in}%
\TextAt{13}{170.715}{367.846}{G}%
\TextAt{8}{183.846}{367.846}{we}%
\TextAt{8}{201.073}{367.846}{now}%
\TextAt{8}{225.462}{367.846}{put}%
\TextAt{13}{250.242}{338.889}{i}%
\TextAt{21}{254.235}{337.029}{L}%
\TextAt{3}{259.994}{337.029}{\Character{106}}%
\TextAt{21}{262.347}{337.029}{K}%
\TextAt{9}{270.549}{338.889}{\Character{40}}%
\TextAt{13}{275.102}{338.889}{\Character{27}}%
\TextAt{9}{282.184}{338.889}{\Character{41}}%
\TextAt{9}{290.039}{338.889}{\Character{61}}%
\TextAt{13}{302.483}{338.889}{v}%
\TextAt{21}{308.142}{337.096}{L}%
\TextAt{9}{314.399}{338.889}{\Character{40}}%
\TextAt{13}{318.951}{338.889}{\Character{27}}%
\TextAt{13}{326.017}{338.889}{x}%
\TextAt{16}{335.343}{338.889}{\Character{0}}%
\TextAt{13}{347.298}{338.889}{x}%
\TextAt{9}{353.95}{338.889}{\Character{41}}%
\TextAt{13}{358.502}{338.889}{\Character{59}}%
\TextAt{8}{79.226}{309.932}{where}%
\TextAt{16}{112.807}{309.932}{\Character{79}}%
\TextAt{9}{125.977}{309.932}{\Character{61}}%
\TextAt{14}{138.403}{309.932}{o}%
\TextAt{9}{144.247}{309.932}{\Character{91}}%
\TextAt{13}{147.499}{309.932}{x}%
\TextAt{9}{154.151}{309.932}{\Character{93}}%
\TextAt{8}{157.402}{309.932}{.}%
\TextAt{8}{165.774}{309.932}{This}%
\TextAt{8}{192.265}{309.932}{definition}%
\TextAt{8}{244.065}{309.932}{does}%
\TextAt{8}{270.234}{309.932}{not}%
\TextAt{8}{290.819}{309.932}{depend}%
\TextAt{8}{331.226}{309.932}{on}%
\TextAt{8}{347.256}{309.932}{the}%
\TextAt{8}{367.183}{309.932}{choice}%
\TextAt{8}{401.744}{309.932}{of}%
\TextAt{8}{414.845}{309.932}{the}%
\TextAt{8}{434.773}{309.932}{generator}%
\TextAt{13}{486.585}{309.932}{x}%
\TextAt{8}{496.906}{309.932}{and}%
\TextAt{8}{519.439}{309.932}{we}%
\TextAt{8}{79.226}{292.886}{may}%
\TextAt{8}{104.595}{292.886}{write}%
\TextAt{13}{205.218}{275.84}{G}%
\TextAt{21}{214.451}{274.046}{s}%
\TextAt{9}{218.865}{275.84}{\Character{40}}%
\TextAt{13}{223.418}{275.84}{L}%
\TextAt{16}{234.703}{275.84}{\Character{106}}%
\TextAt{13}{241.345}{275.84}{K}%
\TextAt{9}{252.162}{275.84}{\Character{41}}%
\TextAt{9}{260.017}{275.84}{\Character{61}}%
\TextAt{16}{272.461}{275.84}{\Character{102}}%
\TextAt{13}{278.439}{275.84}{\Character{27}}%
\NativeMathAt{288.842}{275.84}{11.955}{\in}%
\TextAt{13}{300.133}{275.84}{G}%
\TextAt{16}{312.687}{275.84}{\Character{106}}%
\TextAt{13}{319.329}{275.84}{i}%
\TextAt{21}{323.323}{273.98}{L}%
\TextAt{3}{329.082}{273.98}{\Character{106}}%
\TextAt{21}{331.434}{273.98}{K}%
\TextAt{9}{339.637}{275.84}{\Character{40}}%
\TextAt{13}{344.189}{275.84}{\Character{27}}%
\TextAt{9}{351.271}{275.84}{\Character{41}}%
\NativeMathAt{359.144}{275.84}{11.955}{\geq}%
\TextAt{13}{371.764}{275.84}{s}%
\TextAt{9}{379.934}{275.84}{\Character{43}}%
\TextAt{9}{391.686}{275.84}{1}%
\TextAt{16}{397.549}{275.84}{\Character{103}}%
\TextAt{13}{403.526}{275.84}{\Character{58}}%
\TextAt{8}{79.226}{252.337}{Passing}%
\TextAt{8}{121.115}{252.337}{to}%
\TextAt{8}{135.077}{252.337}{a}%
\TextAt{8}{144.484}{252.337}{Galois}%
\TextAt{8}{180.049}{252.337}{subextension}%
\TextAt{13}{249.411}{252.337}{L}%
\TextAt{3}{257.375}{256.675}{\Character{48}}%
\TextAt{16}{263.491}{252.337}{\Character{106}}%
\TextAt{13}{270.133}{252.337}{K}%
\TextAt{8}{284.499}{252.337}{of}%
\TextAt{13}{297.481}{252.337}{L}%
\TextAt{16}{308.766}{252.337}{\Character{106}}%
\TextAt{13}{315.408}{252.337}{K}%
\TextAt{8}{326.226}{252.337}{,}%
\TextAt{8}{333.096}{252.337}{the}%
\TextAt{8}{352.904}{252.337}{numbers}%
\TextAt{13}{399.766}{252.337}{i}%
\TextAt{21}{403.76}{250.477}{L}%
\TextAt{3}{409.519}{250.477}{\Character{106}}%
\TextAt{21}{411.871}{250.477}{K}%
\TextAt{9}{420.074}{252.337}{\Character{40}}%
\TextAt{13}{424.626}{252.337}{\Character{27}}%
\TextAt{9}{431.709}{252.337}{\Character{41}}%
\TextAt{8}{439.81}{252.337}{obey}%
\TextAt{8}{467.424}{252.337}{the}%
\TextAt{8}{487.232}{252.337}{following}%
\TextAt{8}{79.226}{235.291}{rule.}%
\TextAt{12}{79.226}{209.291}{Proposition}%
\TextAt{12}{151.937}{209.291}{(10.5).}%
\TextAt{8}{193.984}{209.291}{If}%
\TextAt{13}{205.676}{209.291}{e}%
\TextAt{3}{211.101}{213.629}{\Character{48}}%
\TextAt{9}{217.217}{209.291}{\Character{61}}%
\TextAt{13}{229.643}{209.291}{e}%
\TextAt{21}{235.068}{207.431}{L}%
\TextAt{3}{240.827}{207.431}{\Character{106}}%
\TextAt{21}{243.18}{207.431}{L}%
\TextAt{23}{248.939}{209.701}{\Character{48}}%
\TextAt{8}{256.019}{209.291}{is}%
\TextAt{8}{267.782}{209.291}{the}%
\TextAt{8}{287.939}{209.291}{ramification}%
\TextAt{8}{352.987}{209.291}{index}%
\TextAt{8}{384.525}{209.291}{of}%
\TextAt{13}{397.855}{209.291}{L}%
\TextAt{16}{409.14}{209.291}{\Character{106}}%
\TextAt{13}{415.782}{209.291}{L}%
\TextAt{3}{423.746}{213.629}{\Character{48}}%
\TextAt{8}{426.541}{209.291}{,}%
\TextAt{8}{433.691}{209.291}{then}%
\TextAt{13}{234.909}{176.739}{i}%
\TextAt{21}{238.902}{174.879}{L}%
\TextAt{23}{244.661}{177.148}{\Character{48}}%
\TextAt{3}{247.346}{174.879}{\Character{106}}%
\TextAt{21}{249.698}{174.879}{K}%
\TextAt{9}{257.901}{176.739}{\Character{40}}%
\TextAt{13}{262.453}{176.739}{\Character{27}}%
\TextAt{3}{269.535}{181.675}{\Character{48}}%
\TextAt{9}{272.33}{176.739}{\Character{41}}%
\TextAt{9}{280.185}{176.739}{\Character{61}}%
\TextAt{9}{295.008}{184.827}{1}%
\Vector{q 0 G 0 g 0.478 w 221.825 -594.272 m 230.045 -594.272 l S Q}%
\TextAt{13}{293.825}{168.538}{e}%
\TextAt{3}{299.25}{171.992}{\Character{48}}%
\TextAt{20}{311.656}{188.096}{\Character{88}}%
\TextAt{21}{305.233}{162.193}{\Character{27}}%
\TextAt{3}{310.358}{162.193}{\Character{106}}%
\TextAt{22}{312.711}{159.902}{L}%
\TextAt{23}{317.769}{161.618}{\Character{48}}%
\TextAt{17}{320.952}{162.193}{\Character{61}}%
\TextAt{21}{327.539}{162.193}{\Character{27}}%
\TextAt{23}{332.664}{164.463}{\Character{48}}%
\TextAt{13}{337.341}{176.739}{i}%
\TextAt{21}{341.334}{174.879}{L}%
\TextAt{3}{347.093}{174.879}{\Character{106}}%
\TextAt{21}{349.445}{174.879}{K}%
\TextAt{9}{357.648}{176.739}{\Character{40}}%
\TextAt{13}{362.2}{176.739}{\Character{27}}%
\TextAt{9}{369.283}{176.739}{\Character{41}}%
\TextAt{13}{373.835}{176.739}{\Character{58}}%
\TextAt{19}{79.226}{118.925}{Proof.}%
\TextAt{8}{116}{118.925}{For}%
\TextAt{13}{137.699}{118.925}{\Character{27}}%
\TextAt{3}{144.781}{123.263}{\Character{48}}%
\TextAt{9}{152.139}{118.925}{\Character{61}}%
\TextAt{9}{165.792}{118.925}{1}%
\TextAt{8}{176.287}{118.925}{both}%
\TextAt{8}{204.656}{118.925}{sides}%
\TextAt{8}{233.468}{118.925}{are}%
\TextAt{8}{253.708}{118.925}{infinite.}%
\TextAt{8}{300.113}{118.925}{Let}%
\TextAt{13}{321.811}{118.925}{\Character{27}}%
\TextAt{3}{328.894}{123.263}{\Character{48}}%
\NativeMathAt{336.252}{118.925}{11.955}{\neq}%
\TextAt{9}{349.905}{118.925}{1}%
\TextAt{8}{355.773}{118.925}{,}%
\TextAt{8}{363.834}{118.925}{and}%
\TextAt{8}{387.325}{118.925}{let}%
\TextAt{16}{404.959}{118.925}{\Character{79}}%
\TextAt{9}{419.372}{118.925}{\Character{61}}%
\TextAt{14}{433.04}{118.925}{o}%
\TextAt{9}{438.884}{118.925}{\Character{91}}%
\TextAt{13}{442.136}{118.925}{x}%
\TextAt{9}{448.788}{118.925}{\Character{93}}%
\TextAt{8}{456.666}{118.925}{and}%
\TextAt{16}{480.158}{118.925}{\Character{79}}%
\TextAt{3}{490.007}{123.263}{\Character{48}}%
\TextAt{9}{497.366}{118.925}{\Character{61}}%
\TextAt{14}{511.033}{118.925}{o}%
\TextAt{9}{516.878}{118.925}{\Character{91}}%
\TextAt{13}{520.129}{118.925}{y}%
\TextAt{9}{526.266}{118.925}{\Character{93}}%
\TextAt{8}{529.518}{118.925}{,}%
\TextAt{8}{79.226}{101.879}{with}%
\TextAt{16}{105.886}{101.879}{\Character{79}}%
\TextAt{3}{115.736}{106.217}{\Character{48}}%
\TextAt{8}{122.428}{101.879}{the}%
\TextAt{8}{142.585}{101.879}{valuation}%
\TextAt{8}{193.657}{101.879}{ring}%
\TextAt{8}{217.723}{101.879}{of}%
\TextAt{13}{231.053}{101.879}{L}%
\TextAt{3}{239.018}{106.217}{\Character{48}}%
\TextAt{8}{241.813}{101.879}{.}%
\TextAt{8}{250.261}{101.879}{By}%
\TextAt{8}{268.624}{101.879}{definition,}%
\TextAt{8}{323.905}{101.879}{we}%
\TextAt{8}{341.132}{101.879}{have}%
\TextAt{13}{173.802}{72.921}{e}%
\TextAt{3}{179.228}{77.858}{\Character{48}}%
\TextAt{13}{182.023}{72.921}{i}%
\TextAt{21}{186.016}{71.062}{L}%
\TextAt{23}{191.775}{73.331}{\Character{48}}%
\TextAt{3}{194.46}{71.062}{\Character{106}}%
\TextAt{21}{196.812}{71.062}{K}%
\TextAt{9}{205.015}{72.921}{\Character{40}}%
\TextAt{13}{209.567}{72.921}{\Character{27}}%
\TextAt{3}{216.649}{77.858}{\Character{48}}%
\TextAt{9}{219.445}{72.921}{\Character{41}}%
\TextAt{9}{227.3}{72.921}{\Character{61}}%
\TextAt{13}{239.743}{72.921}{v}%
\TextAt{21}{245.402}{71.128}{L}%
\TextAt{9}{251.66}{72.921}{\Character{40}}%
\TextAt{13}{256.212}{72.921}{\Character{27}}%
\TextAt{3}{263.294}{77.858}{\Character{48}}%
\TextAt{13}{266.089}{72.921}{y}%
\TextAt{16}{274.883}{72.921}{\Character{0}}%
\TextAt{13}{286.838}{72.921}{y}%
\TextAt{9}{292.975}{72.921}{\Character{41}}%
\TextAt{13}{297.527}{72.921}{\Character{59}}%
\TextAt{13}{326.662}{72.921}{i}%
\TextAt{21}{330.675}{71.062}{L}%
\TextAt{3}{336.434}{71.062}{\Character{106}}%
\TextAt{21}{338.786}{71.062}{K}%
\TextAt{9}{346.989}{72.921}{\Character{40}}%
\TextAt{13}{351.541}{72.921}{\Character{27}}%
\TextAt{9}{358.624}{72.921}{\Character{41}}%
\TextAt{9}{366.479}{72.921}{\Character{61}}%
\TextAt{13}{378.922}{72.921}{v}%
\TextAt{21}{384.581}{71.128}{L}%
\TextAt{9}{390.839}{72.921}{\Character{40}}%
\TextAt{13}{395.391}{72.921}{\Character{27}}%
\TextAt{13}{402.457}{72.921}{x}%
\TextAt{16}{411.782}{72.921}{\Character{0}}%
\TextAt{13}{423.737}{72.921}{x}%
\TextAt{9}{430.389}{72.921}{\Character{41}}%
\TextAt{13}{434.942}{72.921}{\Character{58}}%
\TextAt{8}{297.211}{43.033}{163}%
\Destination{78.23}{809.815}{nav.page.163}%
\Destination{79.36}{681.955}{nav.print.178}%
\Destination{79.36}{680.955}{nav.number.2.10.4}%
\Destination{79.36}{223.255}{nav.number.2.10.5}%
\Hotlink{504.94}{469.555}{25.078}{11.047}{nav.number.2.6.8}%
}

% Page 176 (body page 164)
\Page{%
\TextAt{2}{79.226}{796.943}{II.}%
\TextAt{2}{94.073}{796.943}{The}%
\TextAt{2}{116.492}{796.943}{Theory}%
\TextAt{2}{154.401}{796.943}{of}%
\TextAt{2}{166.826}{796.943}{Valuations}%
\Vector{q 0 G 0 g 0.299 w 7.226 18.729 m 460.769 18.729 l S Q}%
\TextAt{8}{79.226}{761.124}{We}%
\TextAt{8}{99.499}{761.124}{choose}%
\TextAt{8}{136.761}{761.124}{a}%
\TextAt{8}{146.645}{761.124}{fixed}%
\TextAt{13}{175.059}{761.124}{\Character{27}}%
\NativeMathAt{185.682}{761.124}{11.955}{\in}%
\TextAt{13}{197.192}{761.124}{G}%
\TextAt{9}{209.966}{761.124}{\Character{61}}%
\TextAt{13}{222.611}{761.124}{G}%
\TextAt{9}{231.844}{761.124}{\Character{40}}%
\TextAt{13}{236.396}{761.124}{L}%
\TextAt{16}{247.901}{761.124}{\Character{106}}%
\TextAt{13}{254.762}{761.124}{K}%
\TextAt{9}{265.58}{761.124}{\Character{41}}%
\TextAt{8}{274.158}{761.124}{such}%
\TextAt{8}{300.684}{761.124}{that}%
\TextAt{13}{326.181}{761.124}{\Character{27}}%
\TextAt{16}{333.264}{761.124}{\Character{106}}%
\TextAt{21}{336.585}{759.33}{L}%
\TextAt{23}{342.344}{761.6}{\Character{48}}%
\TextAt{9}{349.066}{761.124}{\Character{61}}%
\TextAt{13}{361.711}{761.124}{\Character{27}}%
\TextAt{3}{368.793}{765.462}{\Character{48}}%
\TextAt{8}{371.588}{761.124}{.}%
\TextAt{8}{380.423}{761.124}{The}%
\TextAt{8}{404.605}{761.124}{other}%
\TextAt{8}{435.303}{761.124}{elements}%
\TextAt{8}{483.288}{761.124}{of}%
\TextAt{13}{496.747}{761.124}{G}%
\TextAt{8}{510.007}{761.124}{with}%
\TextAt{8}{79.226}{744.078}{image}%
\TextAt{13}{112.738}{744.078}{\Character{27}}%
\TextAt{3}{119.821}{748.416}{\Character{48}}%
\TextAt{8}{126.204}{744.078}{in}%
\TextAt{13}{139.548}{744.078}{G}%
\TextAt{3}{148.782}{748.416}{\Character{48}}%
\TextAt{9}{154.897}{744.078}{\Character{61}}%
\TextAt{13}{167.323}{744.078}{G}%
\TextAt{9}{176.557}{744.078}{\Character{40}}%
\TextAt{13}{181.109}{744.078}{L}%
\TextAt{3}{189.073}{748.416}{\Character{48}}%
\TextAt{16}{195.189}{744.078}{\Character{106}}%
\TextAt{13}{201.831}{744.078}{K}%
\TextAt{9}{212.649}{744.078}{\Character{41}}%
\TextAt{8}{220.789}{744.078}{are}%
\TextAt{8}{239.991}{744.078}{then}%
\TextAt{8}{266.342}{744.078}{given}%
\TextAt{8}{296.603}{744.078}{by}%
\TextAt{13}{312.553}{744.078}{\Character{27}}%
\TextAt{13}{319.619}{744.078}{\Character{28}}%
\TextAt{8}{326.054}{744.078}{,}%
\TextAt{13}{332.956}{744.078}{\Character{28}}%
\NativeMathAt{342.696}{744.078}{11.955}{\in}%
\TextAt{13}{353.987}{744.078}{H}%
\TextAt{9}{367.93}{744.078}{\Character{61}}%
\TextAt{13}{380.356}{744.078}{G}%
\TextAt{9}{389.589}{744.078}{\Character{40}}%
\TextAt{13}{394.141}{744.078}{L}%
\TextAt{16}{405.427}{744.078}{\Character{106}}%
\TextAt{13}{412.069}{744.078}{L}%
\TextAt{3}{420.033}{748.416}{\Character{48}}%
\TextAt{9}{422.828}{744.078}{\Character{41}}%
\TextAt{8}{427.38}{744.078}{.}%
\TextAt{8}{435.726}{744.078}{It}%
\TextAt{8}{448.089}{744.078}{therefore}%
\TextAt{8}{496.88}{744.078}{suffices}%
\TextAt{8}{79.226}{727.032}{to}%
\TextAt{8}{93.536}{727.032}{show}%
\TextAt{8}{122.54}{727.032}{that}%
\TextAt{8}{147.909}{727.032}{the}%
\TextAt{8}{168.065}{727.032}{elements}%
\TextAt{13}{200.156}{698.031}{a}%
\TextAt{9}{209.621}{698.031}{\Character{61}}%
\TextAt{13}{222.047}{698.031}{\Character{27}}%
\TextAt{13}{229.113}{698.031}{y}%
\TextAt{16}{237.923}{698.031}{\Character{0}}%
\TextAt{13}{249.878}{698.031}{y}%
\TextAt{8}{279.925}{698.031}{and}%
\TextAt{13}{322.7}{698.031}{b}%
\TextAt{9}{330.998}{698.031}{\Character{61}}%
\TextAt{20}{344.733}{709.388}{\Character{89}}%
\TextAt{21}{343.424}{684.016}{\Character{28}}%
\NativeMathAt{348.102}{684.016}{7.97}{\in}%
\TextAt{21}{353.748}{684.016}{H}%
\TextAt{9}{361.318}{698.031}{\Character{40}}%
\TextAt{13}{365.87}{698.031}{\Character{27}}%
\TextAt{13}{372.935}{698.031}{\Character{28}}%
\TextAt{13}{379.356}{698.031}{x}%
\TextAt{16}{388.68}{698.031}{\Character{0}}%
\TextAt{13}{400.636}{698.031}{x}%
\TextAt{9}{407.288}{698.031}{\Character{41}}%
\TextAt{8}{79.226}{661.069}{generate}%
\TextAt{8}{126.055}{661.069}{the}%
\TextAt{8}{146.211}{661.069}{same}%
\TextAt{8}{175.537}{661.069}{ideal}%
\TextAt{8}{203.5}{661.069}{in}%
\TextAt{16}{217.153}{661.069}{\Character{79}}%
\TextAt{8}{227.003}{661.069}{.}%
\TextAt{8}{97.159}{641.035}{Let}%
\TextAt{13}{118.074}{641.035}{f}%
\TextAt{9}{125.121}{641.035}{\Character{40}}%
\TextAt{13}{129.673}{641.035}{X}%
\TextAt{9}{140.328}{641.035}{\Character{41}}%
\NativeMathAt{148.201}{641.035}{11.955}{\in}%
\TextAt{16}{159.475}{641.035}{\Character{79}}%
\TextAt{3}{169.342}{645.373}{\Character{48}}%
\TextAt{9}{172.137}{641.035}{\Character{91}}%
\TextAt{13}{175.389}{641.035}{X}%
\TextAt{9}{186.044}{641.035}{\Character{93}}%
\TextAt{8}{193.139}{641.035}{be}%
\TextAt{8}{209.01}{641.035}{the}%
\TextAt{8}{229.112}{641.035}{minimal}%
\TextAt{8}{274.584}{641.035}{polynomial}%
\TextAt{8}{335.023}{641.035}{of}%
\TextAt{13}{348.3}{641.035}{x}%
\TextAt{8}{358.795}{641.035}{over}%
\TextAt{13}{383.788}{641.035}{L}%
\TextAt{3}{391.752}{645.373}{\Character{48}}%
\TextAt{8}{394.547}{641.035}{.}%
\TextAt{8}{402.977}{641.035}{Then}%
\TextAt{13}{433.481}{641.035}{f}%
\TextAt{9}{440.528}{641.035}{\Character{40}}%
\TextAt{13}{445.08}{641.035}{X}%
\TextAt{9}{455.735}{641.035}{\Character{41}}%
\TextAt{9}{463.59}{641.035}{\Character{61}}%
\TextAt{20}{476.034}{650.001}{\Character{81}}%
\TextAt{21}{487.325}{637.548}{\Character{28}}%
\NativeMathAt{492.003}{637.548}{7.97}{\in}%
\TextAt{21}{497.649}{637.548}{H}%
\TextAt{9}{505.717}{641.035}{\Character{40}}%
\TextAt{13}{510.269}{641.035}{X}%
\TextAt{16}{523.471}{641.035}{\Character{0}}%
\TextAt{13}{79.226}{623.989}{\Character{28}}%
\TextAt{13}{85.634}{623.989}{x}%
\TextAt{9}{92.297}{623.989}{\Character{41}}%
\TextAt{8}{96.849}{623.989}{.}%
\TextAt{8}{105.291}{623.989}{Letting}%
\TextAt{13}{146.406}{623.989}{\Character{27}}%
\TextAt{8}{157.365}{623.989}{act}%
\TextAt{8}{176.854}{623.989}{on}%
\TextAt{8}{193.091}{623.989}{the}%
\TextAt{8}{213.225}{623.989}{coefficients}%
\TextAt{8}{272.441}{623.989}{of}%
\TextAt{13}{285.75}{623.989}{f}%
\TextAt{8}{292.796}{623.989}{,}%
\TextAt{8}{299.928}{623.989}{we}%
\TextAt{8}{317.134}{623.989}{get}%
\TextAt{8}{336.623}{623.989}{the}%
\TextAt{8}{356.757}{623.989}{polynomial}%
\TextAt{9}{417.229}{623.989}{\Character{40}}%
\TextAt{13}{421.781}{623.989}{\Character{27}}%
\TextAt{13}{428.846}{623.989}{f}%
\TextAt{9}{435.91}{623.989}{\Character{41}\Character{40}}%
\TextAt{13}{445.014}{623.989}{X}%
\TextAt{9}{455.669}{623.989}{\Character{41}}%
\TextAt{9}{463.524}{623.989}{\Character{61}}%
\TextAt{20}{475.968}{632.955}{\Character{81}}%
\TextAt{21}{487.259}{620.502}{\Character{28}}%
\NativeMathAt{491.937}{620.502}{7.97}{\in}%
\TextAt{21}{497.583}{620.502}{H}%
\TextAt{9}{505.651}{623.989}{\Character{40}}%
\TextAt{13}{510.204}{623.989}{X}%
\TextAt{16}{523.471}{623.989}{\Character{0}}%
\TextAt{13}{79.226}{606.943}{\Character{27}}%
\TextAt{13}{86.292}{606.943}{\Character{28}}%
\TextAt{13}{92.712}{606.943}{x}%
\TextAt{9}{99.38}{606.943}{\Character{41}}%
\TextAt{8}{103.932}{606.943}{.}%
\TextAt{8}{118.007}{606.943}{The}%
\TextAt{8}{143.936}{606.943}{coefficients}%
\TextAt{8}{205.049}{606.943}{of}%
\TextAt{13}{220.254}{606.943}{\Character{27}}%
\TextAt{13}{227.32}{606.943}{f}%
\TextAt{16}{238.318}{606.943}{\Character{0}}%
\TextAt{13}{251.552}{606.943}{f}%
\TextAt{8}{264.371}{606.943}{are}%
\TextAt{8}{285.757}{606.943}{all}%
\TextAt{8}{303.892}{606.943}{divisible}%
\TextAt{8}{351.675}{606.943}{by}%
\TextAt{13}{369.81}{606.943}{a}%
\TextAt{9}{382.471}{606.943}{\Character{61}}%
\TextAt{13}{398.093}{606.943}{\Character{27}}%
\TextAt{13}{405.159}{606.943}{y}%
\TextAt{16}{415.247}{606.943}{\Character{0}}%
\TextAt{13}{428.481}{606.943}{y}%
\TextAt{8}{434.617}{606.943}{.}%
\TextAt{8}{448.692}{606.943}{Hence}%
\TextAt{13}{485.345}{606.943}{a}%
\TextAt{8}{497.263}{606.943}{divides}%
\TextAt{9}{79.226}{589.897}{\Character{40}}%
\TextAt{13}{83.778}{589.897}{\Character{27}}%
\TextAt{13}{90.844}{589.897}{f}%
\TextAt{9}{97.907}{589.897}{\Character{41}\Character{40}}%
\TextAt{13}{107.012}{589.897}{x}%
\TextAt{9}{113.664}{589.897}{\Character{41}}%
\TextAt{16}{120.873}{589.897}{\Character{0}}%
\TextAt{13}{132.828}{589.897}{f}%
\TextAt{9}{139.875}{589.897}{\Character{40}}%
\TextAt{13}{144.427}{589.897}{x}%
\TextAt{9}{151.079}{589.897}{\Character{41}}%
\TextAt{9}{158.934}{589.897}{\Character{61}}%
\NativeMathAt{171.378}{589.897}{11.955}{\pm}%
\TextAt{13}{180.676}{589.897}{b}%
\TextAt{8}{185.653}{589.897}{.}%
\TextAt{8}{97.159}{569.862}{To}%
\TextAt{8}{115.708}{569.862}{show}%
\TextAt{8}{146.033}{569.862}{that}%
\TextAt{8}{172.724}{569.862}{conversely}%
\TextAt{13}{230.044}{569.862}{b}%
\TextAt{8}{240.241}{569.862}{is}%
\TextAt{8}{253.326}{569.862}{a}%
\TextAt{8}{264.404}{569.862}{divisor}%
\TextAt{8}{303.839}{569.862}{of}%
\TextAt{13}{318.491}{569.862}{a}%
\TextAt{8}{324.636}{569.862}{,}%
\TextAt{8}{333.437}{569.862}{we}%
\TextAt{8}{351.987}{569.862}{write}%
\TextAt{13}{383.22}{569.862}{y}%
\TextAt{8}{394.576}{569.862}{as}%
\TextAt{8}{410.268}{569.862}{a}%
\TextAt{8}{421.346}{569.862}{polynomial}%
\TextAt{8}{483.161}{569.862}{in}%
\TextAt{13}{498.135}{569.862}{x}%
\TextAt{8}{510.007}{569.862}{with}%
\TextAt{8}{79.226}{552.816}{coefficients}%
\TextAt{8}{138.464}{552.816}{in}%
\TextAt{14}{152.117}{552.816}{o}%
\TextAt{8}{157.961}{552.816}{,}%
\TextAt{13}{165.11}{552.816}{y}%
\TextAt{9}{174.568}{552.816}{\Character{61}}%
\TextAt{13}{186.993}{552.816}{g}%
\TextAt{9}{193.027}{552.816}{\Character{40}}%
\TextAt{13}{197.58}{552.816}{x}%
\TextAt{9}{204.232}{552.816}{\Character{41}}%
\TextAt{8}{208.784}{552.816}{.}%
\TextAt{8}{217.233}{552.816}{As}%
\TextAt{13}{234.52}{552.816}{x}%
\TextAt{8}{245.069}{552.816}{is}%
\TextAt{8}{256.833}{552.816}{a}%
\TextAt{8}{266.589}{552.816}{zero}%
\TextAt{8}{291.3}{552.816}{of}%
\TextAt{8}{304.63}{552.816}{the}%
\TextAt{8}{324.786}{552.816}{polynomial}%
\TextAt{13}{385.279}{552.816}{g}%
\TextAt{9}{391.314}{552.816}{\Character{40}}%
\TextAt{13}{395.866}{552.816}{X}%
\TextAt{9}{406.521}{552.816}{\Character{41}}%
\TextAt{16}{413.73}{552.816}{\Character{0}}%
\TextAt{13}{425.685}{552.816}{y}%
\NativeMathAt{435.143}{552.816}{11.955}{\in}%
\TextAt{16}{446.417}{552.816}{\Character{79}}%
\TextAt{3}{456.284}{557.154}{\Character{48}}%
\TextAt{9}{459.079}{552.816}{\Character{91}}%
\TextAt{13}{462.33}{552.816}{X}%
\TextAt{9}{472.985}{552.816}{\Character{93}}%
\TextAt{8}{476.237}{552.816}{,}%
\TextAt{8}{483.386}{552.816}{we}%
\TextAt{8}{500.614}{552.816}{have}%
\TextAt{13}{196.541}{523.815}{g}%
\TextAt{9}{202.576}{523.815}{\Character{40}}%
\TextAt{13}{207.128}{523.815}{X}%
\TextAt{9}{217.783}{523.815}{\Character{41}}%
\TextAt{16}{224.992}{523.815}{\Character{0}}%
\TextAt{13}{236.947}{523.815}{y}%
\TextAt{9}{246.405}{523.815}{\Character{61}}%
\TextAt{13}{258.83}{523.815}{f}%
\TextAt{9}{265.877}{523.815}{\Character{40}}%
\TextAt{13}{270.429}{523.815}{X}%
\TextAt{9}{281.084}{523.815}{\Character{41}}%
\TextAt{13}{285.636}{523.815}{h}%
\TextAt{9}{292.375}{523.815}{\Character{40}}%
\TextAt{13}{296.927}{523.815}{X}%
\TextAt{9}{307.582}{523.815}{\Character{41}}%
\TextAt{13}{312.135}{523.815}{\Character{59}}%
\TextAt{13}{341.27}{523.815}{h}%
\TextAt{9}{348.028}{523.815}{\Character{40}}%
\TextAt{13}{352.58}{523.815}{X}%
\TextAt{9}{363.235}{523.815}{\Character{41}}%
\NativeMathAt{371.108}{523.815}{11.955}{\in}%
\TextAt{16}{382.382}{523.815}{\Character{79}}%
\TextAt{3}{392.249}{528.751}{\Character{48}}%
\TextAt{9}{395.044}{523.815}{\Character{91}}%
\TextAt{13}{398.296}{523.815}{X}%
\TextAt{9}{408.951}{523.815}{\Character{93}}%
\TextAt{13}{412.203}{523.815}{\Character{58}}%
\TextAt{8}{79.226}{494.814}{Letting}%
\TextAt{13}{121.231}{494.814}{\Character{27}}%
\TextAt{8}{133.077}{494.814}{operate}%
\TextAt{8}{175.895}{494.814}{on}%
\TextAt{8}{193.021}{494.814}{the}%
\TextAt{8}{214.044}{494.814}{coefficients}%
\TextAt{8}{274.149}{494.814}{of}%
\TextAt{8}{288.346}{494.814}{both}%
\TextAt{8}{316.853}{494.814}{sides}%
\TextAt{8}{345.803}{494.814}{and}%
\TextAt{8}{369.432}{494.814}{then}%
\TextAt{8}{396.959}{494.814}{substituting}%
\TextAt{13}{462.994}{494.814}{X}%
\TextAt{9}{478.447}{494.814}{\Character{61}}%
\TextAt{13}{492.35}{494.814}{x}%
\TextAt{8}{503.766}{494.814}{yields}%
\TextAt{13}{79.226}{477.768}{y}%
\TextAt{16}{88.019}{477.768}{\Character{0}}%
\TextAt{13}{99.975}{477.768}{\Character{27}}%
\TextAt{13}{107.041}{477.768}{y}%
\TextAt{9}{116.515}{477.768}{\Character{61}}%
\TextAt{9}{128.925}{477.768}{\Character{40}}%
\TextAt{13}{133.492}{477.768}{\Character{27}}%
\TextAt{13}{140.558}{477.768}{f}%
\TextAt{9}{147.621}{477.768}{\Character{41}\Character{40}}%
\TextAt{13}{156.726}{477.768}{x}%
\TextAt{9}{163.378}{477.768}{\Character{41}\Character{40}}%
\TextAt{13}{172.483}{477.768}{\Character{27}}%
\TextAt{13}{179.549}{477.768}{h}%
\TextAt{9}{186.304}{477.768}{\Character{41}\Character{40}}%
\TextAt{13}{195.408}{477.768}{x}%
\TextAt{9}{202.06}{477.768}{\Character{41}}%
\TextAt{9}{209.915}{477.768}{\Character{61}}%
\NativeMathAt{222.359}{477.768}{11.955}{\pm}%
\TextAt{13}{231.657}{477.768}{b}%
\TextAt{9}{236.635}{477.768}{\Character{40}}%
\TextAt{13}{241.187}{477.768}{\Character{27}}%
\TextAt{13}{248.253}{477.768}{h}%
\TextAt{9}{255.008}{477.768}{\Character{41}\Character{40}}%
\TextAt{13}{264.112}{477.768}{x}%
\TextAt{9}{270.765}{477.768}{\Character{41}}%
\TextAt{8}{275.317}{477.768}{,}%
\TextAt{8}{282.466}{477.768}{i.e.,}%
\TextAt{13}{304.571}{477.768}{b}%
\TextAt{8}{313.446}{477.768}{divides}%
\TextAt{13}{352.85}{477.768}{a}%
\TextAt{8}{358.995}{477.768}{.}%
\Vector{q 0 G 0 g 0.399 w 452.334 -296.232 m 452.334 -288.162 l S Q}%
\Vector{q 0 G 0 g 0.399 w 452.534 -288.362 m 459.707 -288.362 l S Q}%
\Vector{q 0 G 0 g 0.399 w 452.534 -296.033 m 459.707 -296.033 l S Q}%
\Vector{q 0 G 0 g 0.399 w 459.906 -296.232 m 459.906 -288.162 l S Q}%
\TextAt{8}{97.159}{448.767}{We}%
\TextAt{8}{118.588}{448.767}{now}%
\TextAt{8}{144.262}{448.767}{want}%
\TextAt{8}{174.167}{448.767}{to}%
\TextAt{8}{189.762}{448.767}{show}%
\TextAt{8}{220.051}{448.767}{that}%
\TextAt{8}{246.704}{448.767}{the}%
\TextAt{8}{268.146}{448.767}{ramification}%
\TextAt{8}{334.479}{448.767}{group}%
\TextAt{13}{368.939}{448.767}{G}%
\TextAt{21}{378.173}{446.973}{s}%
\TextAt{9}{382.587}{448.767}{\Character{40}}%
\TextAt{13}{387.139}{448.767}{L}%
\TextAt{16}{400.614}{448.767}{\Character{106}}%
\TextAt{13}{409.445}{448.767}{K}%
\TextAt{9}{420.263}{448.767}{\Character{41}}%
\TextAt{8}{429.998}{448.767}{is}%
\TextAt{8}{443.046}{448.767}{mapped}%
\TextAt{8}{488.876}{448.767}{onto}%
\TextAt{8}{516.51}{448.767}{the}%
\TextAt{8}{79.226}{431.721}{ramification}%
\TextAt{8}{144.274}{431.721}{group}%
\TextAt{13}{177.45}{431.721}{G}%
\TextAt{21}{186.683}{429.927}{t}%
\TextAt{9}{190.24}{431.721}{\Character{40}}%
\TextAt{13}{194.792}{431.721}{L}%
\TextAt{3}{202.756}{436.059}{\Character{48}}%
\TextAt{16}{208.872}{431.721}{\Character{106}}%
\TextAt{13}{215.514}{431.721}{K}%
\TextAt{9}{226.332}{431.721}{\Character{41}}%
\TextAt{8}{234.781}{431.721}{by}%
\TextAt{8}{251.04}{431.721}{the}%
\TextAt{8}{271.197}{431.721}{projection}%
\TextAt{13}{242.94}{402.72}{G}%
\TextAt{9}{252.174}{402.72}{\Character{40}}%
\TextAt{13}{256.726}{402.72}{L}%
\TextAt{16}{268.011}{402.72}{\Character{106}}%
\TextAt{13}{274.653}{402.72}{K}%
\TextAt{9}{285.471}{402.72}{\Character{41}}%
\NativeArrowAt{293.344}{402.72}{11.955}{19.26}{\longrightarrow}%
\TextAt{13}{315.926}{402.72}{G}%
\TextAt{9}{325.159}{402.72}{\Character{40}}%
\TextAt{13}{329.712}{402.72}{L}%
\TextAt{3}{337.676}{407.656}{\Character{48}}%
\TextAt{16}{343.792}{402.72}{\Character{106}}%
\TextAt{13}{350.434}{402.72}{K}%
\TextAt{9}{361.252}{402.72}{\Character{41}}%
\TextAt{13}{365.804}{402.72}{\Character{59}}%
\ErrataAt{79.22}{373.755}{453.64}{where $t=\eta_{L\mid L^{\prime}}(s)$. For a general extension $L\mid K$, define}%
\ErrataAt{226.36}{340.355}{158}{$\displaystyle \eta_{L\mid K}(s)=\int_0^s\frac{dx}{(G_0:G_x)}$.}%
\TextAt{8}{79.226}{306.75}{Here}%
\TextAt{9}{106.886}{306.75}{\Character{40}}%
\TextAt{13}{111.438}{306.75}{G}%
\TextAt{17}{120.672}{304.957}{0}%
\TextAt{9}{128.778}{306.75}{\Character{58}}%
\TextAt{13}{135.404}{306.75}{G}%
\TextAt{21}{144.637}{304.957}{x}%
\TextAt{9}{149.902}{306.75}{\Character{41}}%
\TextAt{8}{158.383}{306.75}{is}%
\TextAt{8}{170.178}{306.75}{meant}%
\TextAt{8}{205.656}{306.75}{to}%
\TextAt{8}{219.998}{306.75}{denote}%
\TextAt{8}{257.748}{306.75}{the}%
\TextAt{8}{277.935}{306.75}{inverse}%
\TextAt{9}{316.725}{306.75}{\Character{40}}%
\TextAt{13}{321.277}{306.75}{G}%
\TextAt{21}{330.511}{304.957}{x}%
\TextAt{9}{339.15}{306.75}{\Character{58}}%
\TextAt{13}{345.776}{306.75}{G}%
\TextAt{17}{355.009}{304.957}{0}%
\TextAt{9}{359.741}{306.75}{\Character{41}}%
\TextAt{3}{364.294}{311.088}{\Character{0}}%
\TextAt{17}{370.88}{311.088}{1}%
\TextAt{8}{379.541}{306.75}{when}%
\TextAt{16}{410.13}{306.75}{\Character{0}}%
\TextAt{9}{419.428}{306.75}{1}%
\NativeMathAt{428.655}{306.75}{11.955}{\leq}%
\TextAt{13}{441.328}{306.75}{x}%
\NativeMathAt{451.354}{306.75}{11.955}{\leq}%
\TextAt{9}{464.026}{306.75}{0}%
\TextAt{8}{469.879}{306.75}{,}%
\TextAt{8}{477.067}{306.75}{i.e.,}%
\TextAt{8}{499.211}{306.75}{simply}%
\TextAt{8}{79.226}{289.704}{1,}%
\TextAt{8}{92.233}{289.704}{if}%
\TextAt{16}{102.957}{289.704}{\Character{0}}%
\TextAt{9}{112.256}{289.704}{1}%
\TextAt{13}{121.429}{289.704}{\Character{60}}%
\TextAt{13}{133.839}{289.704}{x}%
\NativeMathAt{143.828}{289.704}{11.955}{\leq}%
\TextAt{9}{156.447}{289.704}{0}%
\TextAt{8}{162.3}{289.704}{.}%
\TextAt{8}{170.749}{289.704}{For}%
\TextAt{9}{191.718}{289.704}{0}%
\TextAt{13}{200.892}{289.704}{\Character{60}}%
\TextAt{13}{213.302}{289.704}{m}%
\NativeMathAt{226.877}{289.704}{11.955}{\leq}%
\TextAt{13}{239.497}{289.704}{s}%
\NativeMathAt{248.331}{289.704}{11.955}{\leq}%
\TextAt{13}{260.951}{289.704}{m}%
\TextAt{9}{273.847}{289.704}{\Character{43}}%
\TextAt{9}{285.599}{289.704}{1}%
\TextAt{8}{291.461}{289.704}{,}%
\TextAt{13}{298.61}{289.704}{m}%
\NativeMathAt{312.17}{289.704}{11.955}{\in}%
\TextAt{18}{323.461}{289.704}{N}%
\TextAt{8}{332.096}{289.704}{,}%
\TextAt{8}{339.245}{289.704}{we}%
\TextAt{8}{356.472}{289.704}{have}%
\TextAt{8}{383.467}{289.704}{explicitly}%
\TextAt{13}{143.285}{258.522}{\Character{17}}%
\TextAt{21}{149.061}{256.662}{L}%
\TextAt{3}{154.82}{256.662}{\Character{106}}%
\TextAt{21}{157.172}{256.662}{K}%
\TextAt{9}{165.375}{258.522}{\Character{40}}%
\TextAt{13}{169.927}{258.522}{s}%
\TextAt{9}{175.441}{258.522}{\Character{41}}%
\TextAt{9}{183.296}{258.522}{\Character{61}}%
\TextAt{9}{199.178}{266.609}{1}%
\Vector{q 0 G 0 g 0.478 w 124.935 -512.49 m 135.273 -512.49 l S Q}%
\TextAt{13}{196.935}{250.321}{g}%
\TextAt{17}{202.541}{248.528}{0}%
\TextAt{20}{208.469}{268.205}{\Character{0}}%
\TextAt{13}{213.948}{258.522}{g}%
\TextAt{17}{219.553}{256.728}{1}%
\TextAt{9}{226.942}{258.522}{\Character{43}}%
\TextAt{13}{238.704}{258.522}{g}%
\TextAt{17}{244.309}{256.728}{2}%
\TextAt{9}{251.698}{258.522}{\Character{43}}%
\TextAt{16}{263.459}{258.522}{\Character{1}}%
\TextAt{16}{268.755}{258.522}{\Character{1}}%
\TextAt{16}{274.063}{258.522}{\Character{1}}%
\TextAt{9}{280.064}{258.522}{\Character{43}}%
\TextAt{13}{291.825}{258.522}{g}%
\TextAt{21}{297.43}{256.728}{m}%
\TextAt{9}{308.076}{258.522}{\Character{43}}%
\TextAt{9}{319.828}{258.522}{\Character{40}}%
\TextAt{13}{324.389}{258.522}{s}%
\TextAt{16}{332.56}{258.522}{\Character{0}}%
\TextAt{13}{344.515}{258.522}{m}%
\TextAt{9}{354.754}{258.522}{\Character{41}}%
\TextAt{13}{359.307}{258.522}{g}%
\TextAt{21}{364.912}{256.728}{m}%
\TextAt{17}{372.402}{256.728}{\Character{43}1}%
\TextAt{20}{383.721}{268.205}{\Character{1}}%
\TextAt{13}{389.201}{258.522}{\Character{59}}%
\TextAt{13}{418.336}{258.522}{g}%
\TextAt{21}{423.961}{256.728}{i}%
\TextAt{9}{430.663}{258.522}{\Character{61}}%
\TextAt{8}{443.088}{258.522}{\char35\relax }%
\TextAt{13}{452.844}{258.522}{G}%
\TextAt{21}{462.077}{256.728}{i}%
\TextAt{13}{465.459}{258.522}{\Character{58}}%
\TextAt{8}{79.226}{226.078}{The}%
\TextAt{8}{103.28}{226.078}{function}%
\TextAt{13}{149.128}{226.078}{\Character{17}}%
\TextAt{21}{154.904}{224.219}{L}%
\TextAt{3}{160.663}{224.219}{\Character{106}}%
\TextAt{21}{163.015}{224.219}{K}%
\TextAt{8}{175.115}{226.078}{can}%
\TextAt{8}{196.574}{226.078}{be}%
\TextAt{8}{212.499}{226.078}{expressed}%
\TextAt{8}{264.97}{226.078}{in}%
\TextAt{8}{278.623}{226.078}{terms}%
\TextAt{8}{311.201}{226.078}{of}%
\TextAt{8}{324.531}{226.078}{the}%
\TextAt{8}{344.687}{226.078}{numbers}%
\TextAt{13}{391.898}{226.078}{i}%
\TextAt{21}{395.891}{224.219}{L}%
\TextAt{3}{401.651}{224.219}{\Character{106}}%
\TextAt{21}{404.003}{224.219}{K}%
\TextAt{9}{412.206}{226.078}{\Character{40}}%
\TextAt{13}{416.758}{226.078}{\Character{27}}%
\TextAt{9}{423.84}{226.078}{\Character{41}}%
\TextAt{8}{432.29}{226.078}{as}%
\TextAt{8}{446.66}{226.078}{follows:}%
\TextAt{12}{79.226}{200.066}{Proposition}%
\TextAt{12}{151.937}{200.066}{(10.6).}%
\TextAt{13}{201.376}{168.322}{\Character{17}}%
\TextAt{21}{207.151}{166.462}{L}%
\TextAt{3}{212.91}{166.462}{\Character{106}}%
\TextAt{21}{215.263}{166.462}{K}%
\TextAt{9}{223.465}{168.322}{\Character{40}}%
\TextAt{13}{228.018}{168.322}{s}%
\TextAt{9}{233.532}{168.322}{\Character{41}}%
\TextAt{9}{241.387}{168.322}{\Character{61}}%
\TextAt{9}{257.268}{176.41}{1}%
\Vector{q 0 G 0 g 0.478 w 183.026 -602.689 m 193.364 -602.689 l S Q}%
\TextAt{13}{255.026}{160.121}{g}%
\TextAt{17}{260.631}{158.328}{0}%
\TextAt{20}{268.602}{179.679}{\Character{88}}%
\TextAt{21}{268.552}{154.308}{\Character{27}}%
\NativeMathAt{273.677}{154.308}{7.97}{\in}%
\TextAt{21}{279.322}{154.308}{G}%
\TextAt{8}{287.914}{168.322}{min}%
\TextAt{16}{307.424}{168.322}{\Character{102}}%
\TextAt{13}{313.402}{168.322}{i}%
\TextAt{21}{317.395}{166.462}{L}%
\TextAt{3}{323.154}{166.462}{\Character{106}}%
\TextAt{21}{325.507}{166.462}{K}%
\TextAt{9}{333.71}{168.322}{\Character{40}}%
\TextAt{13}{338.262}{168.322}{\Character{27}}%
\TextAt{9}{345.344}{168.322}{\Character{41}}%
\TextAt{13}{349.897}{168.322}{\Character{59}}%
\TextAt{13}{355.121}{168.322}{s}%
\TextAt{9}{363.311}{168.322}{\Character{43}}%
\TextAt{9}{375.063}{168.322}{1}%
\TextAt{16}{380.926}{168.322}{\Character{103}}%
\TextAt{16}{389.558}{168.322}{\Character{0}}%
\TextAt{9}{401.515}{168.322}{1}%
\TextAt{13}{407.368}{168.322}{\Character{58}}%
\TextAt{8}{297.211}{43.033}{164}%
\Destination{78.23}{809.815}{nav.page.164}%
\Destination{79.36}{638.955}{nav.print.179}%
\Destination{79.36}{214.055}{nav.number.2.10.6}%
\Hotlink{451.009}{113.355}{26.464}{9.572}{nav.number.2.10.7}%
\Destination{79.226}{387.755}{errata.43}%
}

% Page 177 (body page 165)
\Page{%
\TextAt{2}{79.226}{796.943}{II.}%
\TextAt{2}{94.073}{796.943}{The}%
\TextAt{2}{116.492}{796.943}{Theory}%
\TextAt{2}{154.401}{796.943}{of}%
\TextAt{2}{166.826}{796.943}{Valuations}%
\Vector{q 0 G 0 g 0.299 w 7.226 18.729 m 460.769 18.729 l S Q}%
\TextAt{19}{79.226}{761.124}{Proof.}%
\TextAt{8}{116}{761.124}{Let}%
\TextAt{13}{137.963}{761.124}{\Character{18}}%
\TextAt{9}{143.743}{761.124}{\Character{40}}%
\TextAt{13}{148.295}{761.124}{s}%
\TextAt{9}{153.809}{761.124}{\Character{41}}%
\TextAt{8}{163.252}{761.124}{be}%
\TextAt{8}{180.17}{761.124}{the}%
\TextAt{8}{201.319}{761.124}{function}%
\TextAt{8}{248.16}{761.124}{on}%
\TextAt{8}{265.412}{761.124}{the}%
\TextAt{8}{286.562}{761.124}{right-hand}%
\TextAt{8}{345.119}{761.124}{side.}%
\TextAt{8}{376.118}{761.124}{It}%
\TextAt{8}{389.783}{761.124}{is}%
\TextAt{8}{402.54}{761.124}{continuous}%
\TextAt{8}{462.137}{761.124}{and}%
\TextAt{8}{485.893}{761.124}{piecewise}%
\TextAt{8}{79.226}{744.078}{linear.}%
\TextAt{8}{116.134}{744.078}{One}%
\TextAt{8}{140.362}{744.078}{has}%
\TextAt{13}{160.753}{744.078}{\Character{18}}%
\TextAt{9}{166.533}{744.078}{\Character{40}0\Character{41}}%
\TextAt{9}{184.777}{744.078}{\Character{61}}%
\TextAt{13}{197.237}{744.078}{\Character{17}}%
\TextAt{21}{203.013}{742.218}{L}%
\TextAt{3}{208.772}{742.218}{\Character{106}}%
\TextAt{21}{211.124}{742.218}{K}%
\TextAt{9}{219.327}{744.078}{\Character{40}0\Character{41}}%
\TextAt{9}{237.571}{744.078}{\Character{61}}%
\TextAt{9}{249.992}{744.078}{0}%
\TextAt{8}{255.884}{744.078}{,}%
\TextAt{8}{262.646}{744.078}{and}%
\TextAt{8}{284.926}{744.078}{if}%
\TextAt{13}{295.167}{744.078}{m}%
\NativeMathAt{308.727}{744.078}{11.955}{\geq}%
\TextAt{16}{321.328}{744.078}{\Character{0}}%
\TextAt{9}{330.645}{744.078}{1}%
\TextAt{8}{339.912}{744.078}{is}%
\TextAt{8}{351.193}{744.078}{an}%
\TextAt{8}{366.969}{744.078}{integer}%
\TextAt{8}{405.184}{744.078}{and}%
\TextAt{13}{427.464}{744.078}{m}%
\TextAt{13}{441.009}{744.078}{\Character{60}}%
\TextAt{13}{453.431}{744.078}{s}%
\TextAt{13}{462.266}{744.078}{\Character{60}}%
\TextAt{13}{474.687}{744.078}{m}%
\TextAt{9}{486.618}{744.078}{\Character{43}}%
\TextAt{9}{497.378}{744.078}{1}%
\TextAt{8}{503.244}{744.078}{,}%
\TextAt{8}{510.007}{744.078}{then}%
\TextAt{13}{139.965}{711.482}{\Character{18}}%
\TextAt{3}{145.745}{716.418}{\Character{48}}%
\TextAt{9}{148.541}{711.482}{\Character{40}}%
\TextAt{13}{153.093}{711.482}{s}%
\TextAt{9}{158.607}{711.482}{\Character{41}}%
\TextAt{9}{166.462}{711.482}{\Character{61}}%
\TextAt{9}{182.343}{719.57}{1}%
\Vector{q 0 G 0 g 0.478 w 108.101 -59.529 m 118.439 -59.529 l S Q}%
\TextAt{13}{180.101}{703.281}{g}%
\TextAt{17}{185.706}{701.488}{0}%
\TextAt{8}{191.634}{711.482}{\char35\relax }%
\TextAt{16}{201.39}{711.482}{\Character{102}}%
\TextAt{13}{207.367}{711.482}{\Character{27}}%
\NativeMathAt{217.77}{711.482}{11.955}{\in}%
\TextAt{13}{229.061}{711.482}{G}%
\TextAt{16}{241.616}{711.482}{\Character{106}}%
\TextAt{13}{248.258}{711.482}{i}%
\TextAt{21}{252.251}{709.622}{L}%
\TextAt{3}{258.01}{709.622}{\Character{106}}%
\TextAt{21}{260.362}{709.622}{K}%
\TextAt{9}{268.565}{711.482}{\Character{40}}%
\TextAt{13}{273.117}{711.482}{\Character{27}}%
\TextAt{9}{280.2}{711.482}{\Character{41}}%
\NativeMathAt{288.073}{711.482}{11.955}{\geq}%
\TextAt{13}{300.692}{711.482}{m}%
\TextAt{9}{313.588}{711.482}{\Character{43}}%
\TextAt{9}{325.34}{711.482}{2}%
\TextAt{16}{331.202}{711.482}{\Character{103}}%
\TextAt{9}{340.501}{711.482}{\Character{61}}%
\TextAt{9}{381.699}{719.57}{1}%
\Vector{q 0 G 0 g 0.478 w 282.122 -59.529 m 343.129 -59.529 l S Q}%
\TextAt{9}{354.122}{703.281}{\Character{40}}%
\TextAt{13}{358.674}{703.281}{G}%
\TextAt{17}{367.908}{701.488}{0}%
\TextAt{9}{375.961}{703.281}{\Character{58}}%
\TextAt{13}{382.534}{703.281}{G}%
\TextAt{21}{391.767}{701.488}{m}%
\TextAt{17}{399.258}{701.488}{\Character{43}1}%
\TextAt{9}{410.577}{703.281}{\Character{41}}%
\TextAt{9}{419.645}{711.482}{\Character{61}}%
\TextAt{13}{432.071}{711.482}{\Character{17}}%
\TextAt{3}{438.275}{716.418}{\Character{48}}%
\TextAt{21}{437.846}{708.526}{L}%
\TextAt{3}{443.605}{708.526}{\Character{106}}%
\TextAt{21}{445.958}{708.526}{K}%
\TextAt{9}{454.16}{711.482}{\Character{40}}%
\TextAt{13}{458.713}{711.482}{s}%
\TextAt{9}{464.227}{711.482}{\Character{41}}%
\TextAt{13}{468.779}{711.482}{\Character{58}}%
\TextAt{8}{79.226}{679.039}{Hence}%
\TextAt{13}{114.004}{679.039}{\Character{18}}%
\TextAt{9}{123.105}{679.039}{\Character{61}}%
\TextAt{13}{135.53}{679.039}{\Character{17}}%
\TextAt{21}{141.306}{677.179}{L}%
\TextAt{3}{147.065}{677.179}{\Character{106}}%
\TextAt{21}{149.417}{677.179}{K}%
\TextAt{8}{157.62}{679.039}{.}%
\Vector{q 0 G 0 g 0.399 w 452.334 -94.961 m 452.334 -86.891 l S Q}%
\Vector{q 0 G 0 g 0.399 w 452.534 -87.091 m 459.707 -87.091 l S Q}%
\Vector{q 0 G 0 g 0.399 w 452.534 -94.762 m 459.707 -94.762 l S Q}%
\Vector{q 0 G 0 g 0.399 w 459.906 -94.961 m 459.906 -86.891 l S Q}%
\TextAt{12}{79.226}{644.06}{Theorem}%
\TextAt{12}{137.963}{644.06}{(10.7)}%
\TextAt{12}{178.576}{644.06}{(}%
\TextAt{26}{183.812}{644.06}{Herbrand}%
\TextAt{12}{243.66}{644.06}{).}%
\TextAt{8}{258.033}{644.06}{Let}%
\TextAt{13}{280.5}{644.06}{L}%
\TextAt{3}{288.465}{648.398}{\Character{48}}%
\TextAt{16}{297.132}{644.06}{\Character{106}}%
\TextAt{13}{306.326}{644.06}{K}%
\TextAt{8}{322.539}{644.06}{be}%
\TextAt{8}{339.961}{644.06}{a}%
\TextAt{8}{351.214}{644.06}{Galois}%
\TextAt{8}{388.625}{644.06}{subextension}%
\TextAt{8}{459.833}{644.06}{of}%
\TextAt{13}{474.66}{644.06}{L}%
\TextAt{16}{488.498}{644.06}{\Character{106}}%
\TextAt{13}{497.692}{644.06}{K}%
\TextAt{8}{513.904}{644.06}{and}%
\TextAt{13}{79.226}{627.014}{H}%
\TextAt{9}{93.17}{627.014}{\Character{61}}%
\TextAt{13}{105.595}{627.014}{G}%
\TextAt{9}{114.829}{627.014}{\Character{40}}%
\TextAt{13}{119.381}{627.014}{L}%
\TextAt{16}{130.666}{627.014}{\Character{106}}%
\TextAt{13}{137.308}{627.014}{L}%
\TextAt{3}{145.273}{631.353}{\Character{48}}%
\TextAt{9}{148.068}{627.014}{\Character{41}}%
\TextAt{8}{152.62}{627.014}{.}%
\TextAt{8}{161.068}{627.014}{Then}%
\TextAt{8}{191.626}{627.014}{one}%
\TextAt{8}{213.085}{627.014}{has}%
\TextAt{13}{172.914}{598.013}{G}%
\TextAt{21}{182.148}{596.22}{s}%
\TextAt{9}{186.562}{598.013}{\Character{40}}%
\TextAt{13}{191.114}{598.013}{L}%
\TextAt{16}{202.399}{598.013}{\Character{106}}%
\TextAt{13}{209.041}{598.013}{K}%
\TextAt{9}{219.859}{598.013}{\Character{41}}%
\TextAt{13}{224.411}{598.013}{H}%
\TextAt{9}{235.034}{598.013}{\Character{47}}%
\TextAt{13}{240.887}{598.013}{H}%
\TextAt{9}{254.83}{598.013}{\Character{61}}%
\TextAt{13}{267.256}{598.013}{G}%
\TextAt{21}{276.489}{596.22}{t}%
\TextAt{9}{280.045}{598.013}{\Character{40}}%
\TextAt{13}{284.598}{598.013}{L}%
\TextAt{3}{292.562}{602.949}{\Character{48}}%
\TextAt{16}{298.678}{598.013}{\Character{106}}%
\TextAt{13}{305.32}{598.013}{K}%
\TextAt{9}{316.137}{598.013}{\Character{41}}%
\TextAt{8}{344.6}{598.013}{where}%
\TextAt{13}{378.409}{598.013}{t}%
\TextAt{9}{385.957}{598.013}{\Character{61}}%
\TextAt{13}{398.383}{598.013}{\Character{17}}%
\TextAt{21}{404.158}{596.153}{L}%
\TextAt{3}{409.917}{596.153}{\Character{106}}%
\TextAt{21}{412.27}{596.153}{L}%
\TextAt{23}{418.029}{598.423}{\Character{48}}%
\TextAt{9}{421.211}{598.013}{\Character{40}}%
\TextAt{13}{425.764}{598.013}{s}%
\TextAt{9}{431.278}{598.013}{\Character{41}}%
\TextAt{13}{435.83}{598.013}{\Character{58}}%
\TextAt{19}{79.226}{551.079}{Proof.}%
\ErrataAt{115.3}{551.079}{419.7}{Let $G=G(L\mid K)$ and $G^{\prime}=G(L^{\prime}\mid K)$. The case $\sigma^{\prime}=1$ is immediate. For $\sigma^{\prime}\ne1$, choose a}%
\ErrataAt{79.19}{534.033}{455.81}{preimage $\sigma\in G$ of maximal value $i_{L\mid K}(\sigma)$ and show that}%
\TextAt{13}{219.906}{505.032}{i}%
\TextAt{21}{223.9}{503.172}{L}%
\TextAt{23}{229.659}{505.442}{\Character{48}}%
\TextAt{3}{232.343}{503.172}{\Character{106}}%
\TextAt{21}{234.695}{503.172}{K}%
\TextAt{9}{242.898}{505.032}{\Character{40}}%
\TextAt{13}{247.45}{505.032}{\Character{27}}%
\TextAt{3}{254.533}{509.968}{\Character{48}}%
\TextAt{9}{257.328}{505.032}{\Character{41}}%
\TextAt{16}{264.537}{505.032}{\Character{0}}%
\TextAt{9}{276.492}{505.032}{1}%
\TextAt{9}{285.65}{505.032}{\Character{61}}%
\TextAt{13}{298.091}{505.032}{\Character{17}}%
\TextAt{21}{303.867}{503.172}{L}%
\TextAt{3}{309.626}{503.172}{\Character{106}}%
\TextAt{21}{311.978}{503.172}{L}%
\TextAt{23}{317.737}{505.442}{\Character{48}}%
\TextAt{20}{320.92}{514.716}{\Character{0}}%
\TextAt{13}{326.399}{505.032}{i}%
\TextAt{21}{330.393}{503.172}{L}%
\TextAt{3}{336.152}{503.172}{\Character{106}}%
\TextAt{21}{338.504}{503.172}{K}%
\TextAt{9}{346.707}{505.032}{\Character{40}}%
\TextAt{13}{351.259}{505.032}{\Character{27}}%
\TextAt{9}{358.341}{505.032}{\Character{41}}%
\TextAt{16}{365.55}{505.032}{\Character{0}}%
\TextAt{9}{377.505}{505.032}{1}%
\TextAt{20}{383.358}{514.716}{\Character{1}}%
\TextAt{13}{388.838}{505.032}{\Character{58}}%
\TextAt{8}{517.802}{505.032}{(*)}%
\TextAt{8}{79.226}{476.031}{Let}%
\TextAt{13}{101.112}{476.031}{m}%
\TextAt{9}{116.234}{476.031}{\Character{61}}%
\TextAt{13}{130.221}{476.031}{i}%
\TextAt{21}{134.215}{474.171}{L}%
\TextAt{3}{139.974}{474.171}{\Character{106}}%
\TextAt{21}{142.326}{474.171}{K}%
\TextAt{9}{150.529}{476.031}{\Character{40}}%
\TextAt{13}{155.081}{476.031}{\Character{27}}%
\TextAt{9}{162.164}{476.031}{\Character{41}}%
\TextAt{8}{166.716}{476.031}{.}%
\TextAt{8}{177.914}{476.031}{If}%
\TextAt{13}{190.522}{476.031}{\Character{28}}%
\NativeMathAt{201.824}{476.031}{11.955}{\in}%
\TextAt{13}{214.677}{476.031}{H}%
\TextAt{8}{230.113}{476.031}{belongs}%
\TextAt{8}{273.04}{476.031}{to}%
\TextAt{13}{288.267}{476.031}{H}%
\TextAt{21}{297.966}{474.238}{m}%
\TextAt{3}{305.457}{474.238}{\Character{0}}%
\TextAt{17}{312.043}{474.238}{1}%
\TextAt{9}{321.658}{476.031}{\Character{61}}%
\TextAt{13}{335.646}{476.031}{G}%
\TextAt{21}{344.879}{474.238}{m}%
\TextAt{3}{352.37}{474.238}{\Character{0}}%
\TextAt{17}{358.956}{474.238}{1}%
\TextAt{9}{363.689}{476.031}{\Character{40}}%
\TextAt{13}{368.241}{476.031}{L}%
\TextAt{16}{381.088}{476.031}{\Character{106}}%
\TextAt{13}{389.292}{476.031}{L}%
\TextAt{3}{397.256}{480.369}{\Character{48}}%
\TextAt{9}{400.051}{476.031}{\Character{41}}%
\TextAt{8}{404.604}{476.031}{,}%
\TextAt{8}{412.899}{476.031}{then}%
\TextAt{13}{440.475}{476.031}{i}%
\TextAt{21}{444.469}{474.171}{L}%
\TextAt{3}{450.228}{474.171}{\Character{106}}%
\TextAt{21}{452.58}{474.171}{K}%
\TextAt{9}{460.783}{476.031}{\Character{40}}%
\TextAt{13}{465.335}{476.031}{\Character{28}}%
\TextAt{9}{471.754}{476.031}{\Character{41}}%
\NativeMathAt{481.189}{476.031}{11.955}{\geq}%
\TextAt{13}{495.37}{476.031}{m}%
\TextAt{8}{505.609}{476.031}{,}%
\TextAt{8}{513.904}{476.031}{and}%
\TextAt{13}{79.226}{458.985}{i}%
\TextAt{21}{83.22}{457.125}{L}%
\TextAt{3}{88.979}{457.125}{\Character{106}}%
\TextAt{21}{91.331}{457.125}{K}%
\TextAt{9}{99.534}{458.985}{\Character{40}}%
\TextAt{13}{104.086}{458.985}{\Character{27}}%
\TextAt{13}{111.152}{458.985}{\Character{28}}%
\TextAt{9}{117.587}{458.985}{\Character{41}}%
\NativeMathAt{126.817}{458.985}{11.955}{\geq}%
\TextAt{13}{140.792}{458.985}{m}%
\TextAt{8}{151.032}{458.985}{,}%
\TextAt{8}{159.176}{458.985}{so}%
\TextAt{8}{174.342}{458.985}{that}%
\TextAt{13}{200.506}{458.985}{i}%
\TextAt{21}{204.5}{457.125}{L}%
\TextAt{3}{210.259}{457.125}{\Character{106}}%
\TextAt{21}{212.611}{457.125}{K}%
\TextAt{9}{220.814}{458.985}{\Character{40}}%
\TextAt{13}{225.366}{458.985}{\Character{27}}%
\TextAt{13}{232.431}{458.985}{\Character{28}}%
\TextAt{9}{238.867}{458.985}{\Character{41}}%
\TextAt{9}{248.085}{458.985}{\Character{61}}%
\TextAt{13}{261.879}{458.985}{m}%
\TextAt{8}{272.118}{458.985}{.}%
\TextAt{8}{282.954}{458.985}{If}%
\TextAt{13}{295.442}{458.985}{\Character{28}}%
\NativeMathAt{306.538}{458.985}{11.955}{\notin}%
\TextAt{13}{319.185}{458.985}{H}%
\TextAt{21}{328.884}{457.192}{m}%
\TextAt{3}{336.375}{457.192}{\Character{0}}%
\TextAt{17}{342.961}{457.192}{1}%
\TextAt{8}{347.693}{458.985}{,}%
\TextAt{8}{355.837}{458.985}{then}%
\TextAt{13}{383.293}{458.985}{i}%
\TextAt{21}{387.287}{457.125}{L}%
\TextAt{3}{393.046}{457.125}{\Character{106}}%
\TextAt{21}{395.398}{457.125}{K}%
\TextAt{9}{403.601}{458.985}{\Character{40}}%
\TextAt{13}{408.153}{458.985}{\Character{28}}%
\TextAt{9}{414.572}{458.985}{\Character{41}}%
\TextAt{13}{423.801}{458.985}{\Character{60}}%
\TextAt{13}{437.573}{458.985}{m}%
\TextAt{8}{452.516}{458.985}{and}%
\TextAt{13}{476.074}{458.985}{i}%
\TextAt{21}{480.068}{457.125}{L}%
\TextAt{3}{485.827}{457.125}{\Character{106}}%
\TextAt{21}{488.179}{457.125}{K}%
\TextAt{9}{496.382}{458.985}{\Character{40}}%
\TextAt{13}{500.934}{458.985}{\Character{27}}%
\TextAt{13}{508}{458.985}{\Character{28}}%
\TextAt{9}{514.435}{458.985}{\Character{41}}%
\TextAt{9}{523.653}{458.985}{\Character{61}}%
\TextAt{13}{79.226}{441.939}{i}%
\TextAt{21}{83.22}{440.079}{L}%
\TextAt{3}{88.979}{440.079}{\Character{106}}%
\TextAt{21}{91.331}{440.079}{K}%
\TextAt{9}{99.534}{441.939}{\Character{40}}%
\TextAt{13}{104.086}{441.939}{\Character{28}}%
\TextAt{9}{110.505}{441.939}{\Character{41}}%
\TextAt{8}{115.057}{441.939}{.}%
\TextAt{8}{128.222}{441.939}{In}%
\TextAt{8}{144.415}{441.939}{both}%
\TextAt{8}{173.627}{441.939}{cases}%
\TextAt{8}{204.585}{441.939}{we}%
\TextAt{8}{223.385}{441.939}{therefore}%
\TextAt{8}{274.057}{441.939}{find}%
\TextAt{8}{299.037}{441.939}{that}%
\TextAt{13}{325.978}{441.939}{i}%
\TextAt{21}{329.972}{440.079}{L}%
\TextAt{3}{335.731}{440.079}{\Character{106}}%
\TextAt{21}{338.083}{440.079}{K}%
\TextAt{9}{346.286}{441.939}{\Character{40}}%
\TextAt{13}{350.838}{441.939}{\Character{27}}%
\TextAt{13}{357.904}{441.939}{\Character{28}}%
\TextAt{9}{364.339}{441.939}{\Character{41}}%
\TextAt{9}{374.872}{441.939}{\Character{61}}%
\TextAt{8}{389.996}{441.939}{min}%
\TextAt{16}{409.507}{441.939}{\Character{102}}%
\TextAt{13}{415.484}{441.939}{i}%
\TextAt{21}{419.478}{440.079}{L}%
\TextAt{3}{425.237}{440.079}{\Character{106}}%
\TextAt{21}{427.589}{440.079}{K}%
\TextAt{9}{435.792}{441.939}{\Character{40}}%
\TextAt{13}{440.344}{441.939}{\Character{28}}%
\TextAt{9}{446.763}{441.939}{\Character{41}}%
\TextAt{13}{451.315}{441.939}{\Character{59}}%
\TextAt{13}{456.539}{441.939}{m}%
\TextAt{16}{466.799}{441.939}{\Character{103}}%
\TextAt{8}{472.777}{441.939}{.}%
\TextAt{8}{485.941}{441.939}{Applying}%
\TextAt{8}{79.226}{424.893}{(10.5),}%
\TextAt{8}{116.311}{424.893}{this}%
\TextAt{8}{139.133}{424.893}{gives}%
\TextAt{13}{217.876}{405.104}{i}%
\TextAt{21}{221.87}{403.244}{L}%
\TextAt{23}{227.629}{405.513}{\Character{48}}%
\TextAt{3}{230.313}{403.244}{\Character{106}}%
\TextAt{21}{232.665}{403.244}{K}%
\TextAt{9}{240.868}{405.104}{\Character{40}}%
\TextAt{13}{245.42}{405.104}{\Character{27}}%
\TextAt{3}{252.503}{410.04}{\Character{48}}%
\TextAt{9}{255.298}{405.104}{\Character{41}}%
\TextAt{9}{263.153}{405.104}{\Character{61}}%
\TextAt{9}{277.976}{413.191}{1}%
\Vector{q 0 G 0 g 0.478 w 204.792 -365.907 m 213.013 -365.907 l S Q}%
\TextAt{13}{276.792}{396.903}{e}%
\TextAt{3}{282.217}{400.357}{\Character{48}}%
\TextAt{20}{288.513}{416.461}{\Character{88}}%
\TextAt{21}{288.201}{391.09}{\Character{28}}%
\NativeMathAt{292.879}{391.09}{7.97}{\in}%
\TextAt{21}{298.524}{391.09}{H}%
\TextAt{8}{308.087}{405.104}{min}%
\TextAt{16}{327.598}{405.104}{\Character{102}}%
\TextAt{13}{333.576}{405.104}{i}%
\TextAt{21}{337.569}{403.244}{L}%
\TextAt{3}{343.328}{403.244}{\Character{106}}%
\TextAt{21}{345.681}{403.244}{K}%
\TextAt{9}{353.883}{405.104}{\Character{40}}%
\TextAt{13}{358.436}{405.104}{\Character{28}}%
\TextAt{9}{364.854}{405.104}{\Character{41}}%
\TextAt{13}{369.407}{405.104}{\Character{59}}%
\TextAt{13}{374.631}{405.104}{m}%
\TextAt{16}{384.89}{405.104}{\Character{103}}%
\TextAt{13}{390.868}{405.104}{\Character{58}}%
\TextAt{8}{79.226}{372.953}{Since}%
\TextAt{13}{109.564}{372.953}{i}%
\TextAt{21}{113.558}{371.093}{L}%
\TextAt{3}{119.317}{371.093}{\Character{106}}%
\TextAt{21}{121.669}{371.093}{K}%
\TextAt{9}{129.872}{372.953}{\Character{40}}%
\TextAt{13}{134.424}{372.953}{\Character{28}}%
\TextAt{9}{140.843}{372.953}{\Character{41}}%
\TextAt{9}{148.698}{372.953}{\Character{61}}%
\TextAt{13}{161.142}{372.953}{i}%
\TextAt{21}{165.135}{371.093}{L}%
\TextAt{3}{170.894}{371.093}{\Character{106}}%
\TextAt{21}{173.247}{371.093}{L}%
\TextAt{23}{179.006}{373.362}{\Character{48}}%
\TextAt{9}{182.188}{372.953}{\Character{40}}%
\TextAt{13}{186.74}{372.953}{\Character{28}}%
\TextAt{9}{193.159}{372.953}{\Character{41}}%
\TextAt{8}{201.39}{372.953}{and}%
\TextAt{13}{223.933}{372.953}{e}%
\TextAt{3}{229.359}{377.291}{\Character{48}}%
\TextAt{9}{235.475}{372.953}{\Character{61}}%
\TextAt{13}{247.9}{372.953}{e}%
\TextAt{21}{253.326}{371.093}{L}%
\TextAt{3}{259.085}{371.093}{\Character{106}}%
\TextAt{21}{261.437}{371.093}{L}%
\TextAt{23}{267.196}{373.362}{\Character{48}}%
\TextAt{9}{273.699}{372.953}{\Character{61}}%
\TextAt{8}{286.125}{372.953}{\char35\relax }%
\TextAt{13}{295.88}{372.953}{H}%
\TextAt{17}{305.579}{371.159}{0}%
\TextAt{8}{310.312}{372.953}{,}%
\TextAt{8}{317.285}{372.953}{(10.6)}%
\TextAt{8}{350.899}{372.953}{gives}%
\TextAt{8}{379.361}{372.953}{the}%
\TextAt{8}{399.298}{372.953}{formula}%
\TextAt{8}{442.01}{372.953}{(*),}%
\TextAt{8}{463.952}{372.953}{which}%
\TextAt{8}{497.219}{372.953}{in}%
\TextAt{8}{510.652}{372.953}{turn}%
\TextAt{8}{79.226}{355.907}{yields}%
\TextAt{13}{175.225}{340.539}{\Character{27}}%
\TextAt{3}{182.308}{345.476}{\Character{48}}%
\NativeMathAt{188.424}{340.539}{11.955}{\in}%
\TextAt{13}{199.715}{340.539}{G}%
\TextAt{21}{208.948}{338.746}{s}%
\TextAt{13}{213.362}{340.539}{H}%
\TextAt{9}{223.985}{340.539}{\Character{47}}%
\TextAt{13}{229.838}{340.539}{H}%
\NativeArrowAt{243.781}{340.539}{11.955}{21.926}{\Longleftrightarrow}%
\TextAt{13}{269.02}{340.539}{i}%
\TextAt{21}{273.013}{338.68}{L}%
\TextAt{3}{278.773}{338.68}{\Character{106}}%
\TextAt{21}{281.125}{338.68}{K}%
\TextAt{9}{289.328}{340.539}{\Character{40}}%
\TextAt{13}{293.88}{340.539}{\Character{27}}%
\TextAt{9}{300.962}{340.539}{\Character{41}}%
\TextAt{16}{308.171}{340.539}{\Character{0}}%
\TextAt{9}{320.126}{340.539}{1}%
\NativeMathAt{329.3}{340.539}{11.955}{\geq}%
\TextAt{13}{341.92}{340.539}{s}%
\NativeArrowAt{243.781}{320.505}{11.955}{21.926}{\Longleftrightarrow}%
\TextAt{13}{269.02}{320.505}{\Character{17}}%
\TextAt{21}{274.796}{318.645}{L}%
\TextAt{3}{280.555}{318.645}{\Character{106}}%
\TextAt{21}{282.907}{318.645}{L}%
\TextAt{23}{288.666}{320.914}{\Character{48}}%
\TextAt{20}{291.848}{330.188}{\Character{0}}%
\TextAt{13}{297.328}{320.505}{i}%
\TextAt{21}{301.321}{318.645}{L}%
\TextAt{3}{307.08}{318.645}{\Character{106}}%
\TextAt{21}{309.433}{318.645}{K}%
\TextAt{9}{317.635}{320.505}{\Character{40}}%
\TextAt{13}{322.188}{320.505}{\Character{27}}%
\TextAt{9}{329.27}{320.505}{\Character{41}}%
\TextAt{16}{336.479}{320.505}{\Character{0}}%
\TextAt{9}{348.434}{320.505}{1}%
\TextAt{20}{354.287}{330.188}{\Character{1}}%
\NativeMathAt{363.088}{320.505}{11.955}{\geq}%
\TextAt{13}{375.707}{320.505}{\Character{17}}%
\TextAt{21}{381.482}{318.645}{L}%
\TextAt{3}{387.241}{318.645}{\Character{106}}%
\TextAt{21}{389.594}{318.645}{L}%
\TextAt{23}{395.353}{320.914}{\Character{48}}%
\TextAt{9}{398.535}{320.505}{\Character{40}}%
\TextAt{13}{403.088}{320.505}{s}%
\TextAt{9}{408.602}{320.505}{\Character{41}}%
\NativeArrowAt{243.781}{300.47}{11.955}{21.926}{\Longleftrightarrow}%
\TextAt{13}{269.02}{300.47}{i}%
\TextAt{21}{273.013}{298.61}{L}%
\TextAt{23}{278.772}{300.879}{\Character{48}}%
\TextAt{3}{281.457}{298.61}{\Character{106}}%
\TextAt{21}{283.809}{298.61}{K}%
\TextAt{9}{292.012}{300.47}{\Character{40}}%
\TextAt{13}{296.564}{300.47}{\Character{27}}%
\TextAt{3}{303.647}{305.406}{\Character{48}}%
\TextAt{9}{306.442}{300.47}{\Character{41}}%
\TextAt{16}{313.651}{300.47}{\Character{0}}%
\TextAt{9}{325.606}{300.47}{1}%
\NativeMathAt{334.78}{300.47}{11.955}{\geq}%
\TextAt{13}{347.399}{300.47}{\Character{17}}%
\TextAt{21}{353.174}{298.61}{L}%
\TextAt{3}{358.933}{298.61}{\Character{106}}%
\TextAt{21}{361.286}{298.61}{L}%
\TextAt{23}{367.045}{300.879}{\Character{48}}%
\TextAt{9}{370.227}{300.47}{\Character{40}}%
\TextAt{13}{374.78}{300.47}{s}%
\TextAt{9}{380.294}{300.47}{\Character{41}}%
\NativeArrowAt{243.781}{280.435}{11.955}{21.926}{\Longleftrightarrow}%
\TextAt{13}{269.02}{280.435}{\Character{27}}%
\TextAt{3}{276.102}{285.371}{\Character{48}}%
\NativeMathAt{282.218}{280.435}{11.955}{\in}%
\TextAt{13}{293.509}{280.435}{G}%
\TextAt{21}{302.743}{278.642}{t}%
\TextAt{9}{306.299}{280.435}{\Character{40}}%
\TextAt{13}{310.851}{280.435}{L}%
\TextAt{3}{318.816}{285.371}{\Character{48}}%
\TextAt{16}{324.932}{280.435}{\Character{106}}%
\TextAt{13}{331.573}{280.435}{K}%
\TextAt{9}{342.391}{280.435}{\Character{41}}%
\TextAt{13}{346.943}{280.435}{\Character{59}}%
\TextAt{13}{376.078}{280.435}{t}%
\TextAt{9}{383.646}{280.435}{\Character{61}}%
\TextAt{13}{396.071}{280.435}{\Character{17}}%
\TextAt{21}{401.847}{278.575}{L}%
\TextAt{3}{407.606}{278.575}{\Character{106}}%
\TextAt{21}{409.958}{278.575}{L}%
\TextAt{23}{415.717}{280.845}{\Character{48}}%
\TextAt{9}{418.9}{280.435}{\Character{40}}%
\TextAt{13}{423.452}{280.435}{s}%
\TextAt{9}{428.966}{280.435}{\Character{41}}%
\TextAt{13}{433.519}{280.435}{\Character{58}}%
\Vector{q 0 G 0 g 0.399 w 452.334 -514.221 m 452.334 -506.151 l S Q}%
\Vector{q 0 G 0 g 0.399 w 452.534 -506.35 m 459.707 -506.35 l S Q}%
\Vector{q 0 G 0 g 0.399 w 452.534 -514.021 m 459.707 -514.021 l S Q}%
\Vector{q 0 G 0 g 0.399 w 459.906 -514.221 m 459.906 -506.151 l S Q}%
\TextAt{8}{97.159}{230.778}{The}%
\TextAt{8}{122.827}{230.778}{function}%
\TextAt{13}{170.29}{230.778}{\Character{17}}%
\TextAt{21}{176.065}{228.919}{L}%
\TextAt{3}{181.824}{228.919}{\Character{106}}%
\TextAt{21}{184.177}{228.919}{K}%
\TextAt{8}{197.891}{230.778}{is}%
\TextAt{8}{211.27}{230.778}{by}%
\TextAt{8}{229.143}{230.778}{definition}%
\TextAt{8}{282.786}{230.778}{strictly}%
\TextAt{8}{324.463}{230.778}{increasing.}%
\TextAt{8}{388.552}{230.778}{Let}%
\TextAt{8}{411.136}{230.778}{the}%
\TextAt{8}{432.907}{230.778}{inverse}%
\TextAt{8}{473.28}{230.778}{function}%
\TextAt{8}{520.743}{230.778}{be}%
\TextAt{13}{79.226}{213.732}{\Character{32}}%
\TextAt{21}{86.836}{211.873}{L}%
\TextAt{3}{92.595}{211.873}{\Character{106}}%
\TextAt{21}{94.948}{211.873}{K}%
\TextAt{9}{109.365}{213.732}{\Character{58}}%
\TextAt{9}{118.81}{213.732}{\Character{91}}%
\TextAt{16}{122.082}{213.732}{\Character{0}}%
\TextAt{9}{131.38}{213.732}{1}%
\TextAt{13}{137.233}{213.732}{\Character{59}}%
\NativeMathAt{142.478}{213.732}{11.955}{\infty}%
\TextAt{9}{154.433}{213.732}{\Character{41}}%
\NativeMathAt{165.199}{213.732}{11.955}{\rightarrow}%
\TextAt{9}{183.369}{213.732}{\Character{91}}%
\TextAt{16}{186.62}{213.732}{\Character{0}}%
\TextAt{9}{195.919}{213.732}{1}%
\TextAt{13}{201.772}{213.732}{\Character{59}}%
\NativeMathAt{207.016}{213.732}{11.955}{\infty}%
\TextAt{9}{218.971}{213.732}{\Character{41}}%
\TextAt{8}{223.523}{213.732}{.}%
\TextAt{8}{237.065}{213.732}{One}%
\TextAt{8}{263.474}{213.732}{defines}%
\TextAt{8}{303.596}{213.732}{the}%
\TextAt{12}{325.45}{213.732}{upper}%
\TextAt{12}{366.305}{213.732}{numbering}%
\TextAt{8}{434.737}{213.732}{of}%
\TextAt{8}{449.765}{213.732}{the}%
\TextAt{8}{471.619}{213.732}{ramification}%
\TextAt{8}{79.226}{196.686}{groups}%
\TextAt{8}{117.016}{196.686}{by}%
\TextAt{13}{185.687}{179.64}{G}%
\TextAt{21}{194.921}{184.577}{t}%
\TextAt{9}{198.477}{179.64}{\Character{40}}%
\TextAt{13}{203.029}{179.64}{L}%
\TextAt{16}{214.314}{179.64}{\Character{106}}%
\TextAt{13}{220.956}{179.64}{K}%
\TextAt{9}{231.774}{179.64}{\Character{41}}%
\TextAt{9}{239.629}{179.64}{\Character{58}\Character{61}}%
\TextAt{13}{255.324}{179.64}{G}%
\TextAt{21}{264.558}{177.847}{s}%
\TextAt{9}{268.972}{179.64}{\Character{40}}%
\TextAt{13}{273.524}{179.64}{L}%
\TextAt{16}{284.809}{179.64}{\Character{106}}%
\TextAt{13}{291.451}{179.64}{K}%
\TextAt{9}{302.269}{179.64}{\Character{41}}%
\TextAt{8}{330.731}{179.64}{where}%
\TextAt{13}{364.541}{179.64}{s}%
\TextAt{9}{373.375}{179.64}{\Character{61}}%
\TextAt{13}{385.801}{179.64}{\Character{32}}%
\TextAt{21}{393.411}{177.781}{L}%
\TextAt{3}{399.17}{177.781}{\Character{106}}%
\TextAt{21}{401.523}{177.781}{K}%
\TextAt{9}{409.725}{179.64}{\Character{40}}%
\TextAt{13}{414.278}{179.64}{t}%
\TextAt{9}{418.505}{179.64}{\Character{41}}%
\TextAt{13}{423.057}{179.64}{\Character{58}}%
\TextAt{8}{79.226}{156.119}{The}%
\TextAt{8}{103.28}{156.119}{functions}%
\TextAt{13}{153.743}{156.119}{\Character{17}}%
\TextAt{21}{159.518}{154.259}{L}%
\TextAt{3}{165.277}{154.259}{\Character{106}}%
\TextAt{21}{167.63}{154.259}{K}%
\TextAt{8}{179.73}{156.119}{and}%
\TextAt{13}{202.492}{156.119}{\Character{32}}%
\TextAt{21}{210.103}{154.259}{L}%
\TextAt{3}{215.862}{154.259}{\Character{106}}%
\TextAt{21}{218.214}{154.259}{K}%
\TextAt{8}{230.314}{156.119}{satisfy}%
\TextAt{8}{266.861}{156.119}{the}%
\TextAt{8}{287.017}{156.119}{following}%
\TextAt{8}{336.452}{156.119}{transitivity}%
\TextAt{8}{397.34}{156.119}{condition:}%
\TextAt{12}{79.226}{130.106}{Proposition}%
\TextAt{12}{151.937}{130.106}{(10.8).}%
\TextAt{8}{193.984}{130.106}{If}%
\TextAt{13}{205.676}{130.106}{L}%
\TextAt{3}{213.64}{134.445}{\Character{48}}%
\TextAt{16}{219.756}{130.106}{\Character{106}}%
\TextAt{13}{226.398}{130.106}{K}%
\TextAt{8}{241.113}{130.106}{is}%
\TextAt{8}{252.877}{130.106}{a}%
\TextAt{8}{262.632}{130.106}{Galois}%
\TextAt{8}{298.546}{130.106}{subextension}%
\TextAt{8}{368.256}{130.106}{of}%
\TextAt{13}{381.586}{130.106}{L}%
\TextAt{16}{392.872}{130.106}{\Character{106}}%
\TextAt{13}{399.513}{130.106}{K}%
\TextAt{8}{410.331}{130.106}{,}%
\TextAt{8}{417.48}{130.106}{then}%
\TextAt{13}{171.548}{101.105}{\Character{17}}%
\TextAt{21}{177.323}{99.246}{L}%
\TextAt{3}{183.082}{99.246}{\Character{106}}%
\TextAt{21}{185.435}{99.246}{K}%
\TextAt{9}{196.958}{101.105}{\Character{61}}%
\TextAt{13}{209.384}{101.105}{\Character{17}}%
\TextAt{21}{215.159}{99.246}{L}%
\TextAt{23}{220.918}{101.515}{\Character{48}}%
\TextAt{3}{223.603}{99.246}{\Character{106}}%
\TextAt{21}{225.955}{99.246}{K}%
\TextAt{16}{236.814}{101.105}{\Character{14}}%
\TextAt{13}{245.448}{101.105}{\Character{17}}%
\TextAt{21}{251.224}{99.246}{L}%
\TextAt{3}{256.983}{99.246}{\Character{106}}%
\TextAt{21}{259.335}{99.246}{L}%
\TextAt{23}{265.094}{101.515}{\Character{48}}%
\TextAt{8}{292.187}{101.105}{and}%
\TextAt{13}{334.963}{101.105}{\Character{32}}%
\TextAt{21}{342.573}{99.246}{L}%
\TextAt{3}{348.332}{99.246}{\Character{106}}%
\TextAt{21}{350.684}{99.246}{K}%
\TextAt{9}{362.208}{101.105}{\Character{61}}%
\TextAt{13}{374.633}{101.105}{\Character{32}}%
\TextAt{21}{382.244}{99.246}{L}%
\TextAt{3}{388.003}{99.246}{\Character{106}}%
\TextAt{21}{390.355}{99.246}{L}%
\TextAt{23}{396.114}{101.515}{\Character{48}}%
\TextAt{16}{401.953}{101.105}{\Character{14}}%
\TextAt{13}{410.588}{101.105}{\Character{32}}%
\TextAt{21}{418.198}{99.246}{L}%
\TextAt{23}{423.957}{101.515}{\Character{48}}%
\TextAt{3}{426.641}{99.246}{\Character{106}}%
\TextAt{21}{428.994}{99.246}{K}%
\TextAt{13}{437.196}{101.105}{\Character{58}}%
\TextAt{8}{297.211}{43.033}{165}%
\Destination{522.36}{522.075}{nav.AMS.21}%
\Destination{78.23}{809.815}{nav.page.165}%
\Destination{79.36}{659.055}{nav.print.180}%
\Destination{79.36}{658.055}{nav.number.2.10.7}%
\Destination{79.36}{144.155}{nav.number.2.10.8}%
\Hotlink{78.726}{422.655}{30.936}{11.047}{nav.number.2.10.5}%
\Hotlink{316.785}{370.755}{30.936}{11.047}{nav.number.2.10.6}%
\Destination{79.226}{517}{nav.local.2.10.star}%
\Destination{79.226}{567}{errata.194}%
\Hotlink{442.01}{370.753}{18}{11.047}{nav.local.2.10.star}%
}

% Page 178 (body page 166)
\Page{%
\TextAt{2}{79.226}{796.943}{II.}%
\TextAt{2}{94.073}{796.943}{The}%
\TextAt{2}{116.492}{796.943}{Theory}%
\TextAt{2}{154.401}{796.943}{of}%
\TextAt{2}{166.826}{796.943}{Valuations}%
\Vector{q 0 G 0 g 0.299 w 7.226 18.729 m 460.769 18.729 l S Q}%
\TextAt{19}{79.226}{761.124}{Proof.}%
\TextAt{8}{116}{761.124}{For}%
\TextAt{8}{138.572}{761.124}{the}%
\TextAt{8}{160.331}{761.124}{ramification}%
\TextAt{8}{226.981}{761.124}{indices}%
\TextAt{8}{267.008}{761.124}{of}%
\TextAt{8}{281.94}{761.124}{the}%
\TextAt{8}{303.699}{761.124}{extensions}%
\TextAt{13}{361.682}{761.124}{L}%
\TextAt{16}{375.698}{761.124}{\Character{106}}%
\TextAt{13}{385.071}{761.124}{K}%
\TextAt{8}{395.888}{761.124}{,}%
\TextAt{13}{405.041}{761.124}{L}%
\TextAt{3}{413.005}{765.462}{\Character{48}}%
\TextAt{16}{421.852}{761.124}{\Character{106}}%
\TextAt{13}{431.224}{761.124}{K}%
\TextAt{8}{442.042}{761.124}{,}%
\TextAt{13}{451.194}{761.124}{L}%
\TextAt{16}{465.21}{761.124}{\Character{106}}%
\TextAt{13}{474.583}{761.124}{L}%
\TextAt{3}{482.547}{765.462}{\Character{48}}%
\TextAt{8}{490.842}{761.124}{we}%
\TextAt{8}{509.672}{761.124}{have}%
\TextAt{13}{79.226}{744.078}{e}%
\TextAt{21}{84.652}{742.218}{L}%
\TextAt{3}{90.411}{742.218}{\Character{106}}%
\TextAt{21}{92.763}{742.218}{K}%
\TextAt{9}{104.286}{744.078}{\Character{61}}%
\TextAt{13}{116.712}{744.078}{e}%
\TextAt{21}{122.137}{742.218}{L}%
\TextAt{23}{127.896}{744.487}{\Character{48}}%
\TextAt{3}{130.581}{742.218}{\Character{106}}%
\TextAt{21}{132.933}{742.218}{K}%
\TextAt{13}{141.136}{744.078}{e}%
\TextAt{21}{146.561}{742.218}{L}%
\TextAt{3}{152.32}{742.218}{\Character{106}}%
\TextAt{21}{154.673}{742.218}{L}%
\TextAt{23}{160.432}{744.487}{\Character{48}}%
\TextAt{8}{163.614}{744.078}{.}%
\TextAt{8}{172.062}{744.078}{From}%
\TextAt{8}{202.787}{744.078}{(10.7),}%
\TextAt{8}{239.872}{744.078}{we}%
\TextAt{8}{257.1}{744.078}{obtain}%
\TextAt{13}{293.527}{744.078}{G}%
\TextAt{21}{302.761}{742.284}{s}%
\TextAt{9}{307.175}{744.078}{\Character{47}}%
\TextAt{13}{313.028}{744.078}{H}%
\TextAt{21}{322.727}{742.284}{s}%
\TextAt{9}{330.462}{744.078}{\Character{61}}%
\TextAt{9}{342.872}{744.078}{\Character{40}}%
\TextAt{13}{347.439}{744.078}{G}%
\TextAt{9}{356.673}{744.078}{\Character{47}}%
\TextAt{13}{362.526}{744.078}{H}%
\TextAt{9}{373.149}{744.078}{\Character{41}}%
\TextAt{21}{377.701}{742.284}{t}%
\TextAt{8}{381.257}{744.078}{,}%
\TextAt{13}{388.406}{744.078}{t}%
\TextAt{9}{395.954}{744.078}{\Character{61}}%
\TextAt{13}{408.38}{744.078}{\Character{17}}%
\TextAt{21}{414.155}{742.218}{L}%
\TextAt{3}{419.914}{742.218}{\Character{106}}%
\TextAt{21}{422.267}{742.218}{L}%
\TextAt{23}{428.026}{744.487}{\Character{48}}%
\TextAt{9}{431.208}{744.078}{\Character{40}}%
\TextAt{13}{435.76}{744.078}{s}%
\TextAt{9}{441.274}{744.078}{\Character{41}}%
\TextAt{8}{445.827}{744.078}{;}%
\TextAt{8}{452.976}{744.078}{thus}%
\TextAt{9}{215.715}{719.57}{1}%
\Vector{q 0 G 0 g 0.478 w 135.771 -59.529 m 157.511 -59.529 l S Q}%
\TextAt{13}{207.771}{703.281}{e}%
\TextAt{21}{213.197}{701.421}{L}%
\TextAt{3}{218.956}{701.421}{\Character{106}}%
\TextAt{21}{221.308}{701.421}{K}%
\TextAt{8}{230.706}{711.482}{\char35\relax }%
\TextAt{13}{240.462}{711.482}{G}%
\TextAt{21}{249.695}{709.689}{s}%
\TextAt{9}{257.43}{711.482}{\Character{61}}%
\TextAt{9}{280.337}{719.57}{1}%
\Vector{q 0 G 0 g 0.478 w 199.051 -59.529 m 223.475 -59.529 l S Q}%
\TextAt{13}{271.051}{703.281}{e}%
\TextAt{21}{276.477}{701.421}{L}%
\TextAt{23}{282.236}{703.691}{\Character{48}}%
\TextAt{3}{284.92}{701.421}{\Character{106}}%
\TextAt{21}{287.272}{701.421}{K}%
\TextAt{8}{296.671}{711.482}{\char35\relax }%
\TextAt{9}{306.426}{711.482}{\Character{40}}%
\TextAt{13}{310.978}{711.482}{G}%
\TextAt{9}{320.212}{711.482}{\Character{47}}%
\TextAt{13}{326.065}{711.482}{H}%
\TextAt{9}{336.688}{711.482}{\Character{41}}%
\TextAt{21}{341.24}{709.689}{t}%
\TextAt{16}{347.453}{711.482}{\Character{1}}%
\TextAt{9}{362.938}{719.57}{1}%
\Vector{q 0 G 0 g 0.478 w 282.626 -59.529 m 305.104 -59.529 l S Q}%
\TextAt{13}{354.626}{703.281}{e}%
\TextAt{21}{360.051}{701.421}{L}%
\TextAt{3}{365.81}{701.421}{\Character{106}}%
\TextAt{21}{368.163}{701.421}{L}%
\TextAt{23}{373.922}{703.691}{\Character{48}}%
\TextAt{8}{378.3}{711.482}{\char35\relax }%
\TextAt{13}{388.055}{711.482}{H}%
\TextAt{21}{397.754}{709.689}{s}%
\TextAt{13}{402.168}{711.482}{\Character{58}}%
\TextAt{8}{79.226}{678.181}{This}%
\TextAt{8}{105.946}{678.181}{equation}%
\TextAt{8}{153.755}{678.181}{is}%
\TextAt{8}{165.519}{678.181}{equivalent}%
\TextAt{8}{221.134}{678.181}{to}%
\TextAt{13}{189.789}{649.179}{\Character{17}}%
\TextAt{3}{195.994}{654.116}{\Character{48}}%
\TextAt{21}{195.565}{646.224}{L}%
\TextAt{3}{201.324}{646.224}{\Character{106}}%
\TextAt{21}{203.676}{646.224}{K}%
\TextAt{9}{211.879}{649.18}{\Character{40}}%
\TextAt{13}{216.431}{649.179}{s}%
\TextAt{9}{221.945}{649.18}{\Character{41}}%
\TextAt{9}{229.8}{649.18}{\Character{61}}%
\TextAt{13}{242.244}{649.179}{\Character{17}}%
\TextAt{3}{248.448}{654.116}{\Character{48}}%
\TextAt{21}{248.019}{646.224}{L}%
\TextAt{23}{253.778}{648.493}{\Character{48}}%
\TextAt{3}{256.463}{646.224}{\Character{106}}%
\TextAt{21}{258.815}{646.224}{K}%
\TextAt{9}{267.018}{649.18}{\Character{40}}%
\TextAt{13}{271.57}{649.179}{t}%
\TextAt{9}{275.797}{649.18}{\Character{41}}%
\TextAt{13}{280.35}{649.179}{\Character{17}}%
\TextAt{3}{286.554}{654.116}{\Character{48}}%
\TextAt{21}{286.125}{646.224}{L}%
\TextAt{3}{291.884}{646.224}{\Character{106}}%
\TextAt{21}{294.236}{646.224}{L}%
\TextAt{23}{299.995}{648.493}{\Character{48}}%
\TextAt{9}{303.178}{649.179}{\Character{40}}%
\TextAt{13}{307.73}{649.18}{s}%
\TextAt{9}{313.244}{649.179}{\Character{41}}%
\TextAt{9}{321.099}{649.179}{\Character{61}}%
\TextAt{9}{333.52}{649.179}{\Character{40}}%
\TextAt{13}{338.095}{649.18}{\Character{17}}%
\TextAt{21}{343.871}{647.32}{L}%
\TextAt{23}{349.63}{649.589}{\Character{48}}%
\TextAt{3}{352.314}{647.32}{\Character{106}}%
\TextAt{21}{354.667}{647.32}{K}%
\TextAt{16}{365.526}{649.179}{\Character{14}}%
\TextAt{13}{374.16}{649.18}{\Character{17}}%
\TextAt{21}{379.936}{647.32}{L}%
\TextAt{3}{385.695}{647.32}{\Character{106}}%
\TextAt{21}{388.047}{647.32}{L}%
\TextAt{23}{393.806}{649.589}{\Character{48}}%
\TextAt{9}{396.989}{649.179}{\Character{41}}%
\TextAt{3}{401.541}{654.116}{\Character{48}}%
\TextAt{9}{404.336}{649.179}{\Character{40}}%
\TextAt{13}{408.888}{649.18}{s}%
\TextAt{9}{414.402}{649.179}{\Character{41}}%
\TextAt{13}{418.955}{649.18}{\Character{58}}%
\TextAt{8}{79.226}{620.178}{As}%
\TextAt{13}{96.789}{620.178}{\Character{17}}%
\TextAt{21}{102.564}{618.319}{L}%
\TextAt{3}{108.323}{618.319}{\Character{106}}%
\TextAt{21}{110.675}{618.319}{K}%
\TextAt{9}{118.878}{620.178}{\Character{40}0\Character{41}}%
\TextAt{9}{137.6}{620.178}{\Character{61}}%
\TextAt{9}{150.488}{620.178}{\Character{40}}%
\TextAt{13}{155.073}{620.178}{\Character{17}}%
\TextAt{21}{160.848}{618.319}{L}%
\TextAt{23}{166.607}{620.588}{\Character{48}}%
\TextAt{3}{169.291}{618.319}{\Character{106}}%
\TextAt{21}{171.644}{618.319}{K}%
\TextAt{16}{182.691}{620.178}{\Character{14}}%
\TextAt{13}{191.513}{620.178}{\Character{17}}%
\TextAt{21}{197.288}{618.319}{L}%
\TextAt{3}{203.047}{618.319}{\Character{106}}%
\TextAt{21}{205.4}{618.319}{L}%
\TextAt{23}{211.159}{620.588}{\Character{48}}%
\TextAt{9}{214.341}{620.178}{\Character{41}\Character{40}0\Character{41}}%
\TextAt{8}{233.851}{620.178}{,}%
\TextAt{8}{241.344}{620.178}{it}%
\TextAt{8}{253.324}{620.178}{follows}%
\TextAt{8}{292.035}{620.178}{that}%
\TextAt{13}{317.679}{620.178}{\Character{17}}%
\TextAt{21}{323.455}{618.319}{L}%
\TextAt{3}{329.214}{618.319}{\Character{106}}%
\TextAt{21}{331.566}{618.319}{K}%
\TextAt{9}{343.558}{620.178}{\Character{61}}%
\TextAt{13}{356.453}{620.178}{\Character{17}}%
\TextAt{21}{362.229}{618.319}{L}%
\TextAt{23}{367.988}{620.588}{\Character{48}}%
\TextAt{3}{370.672}{618.319}{\Character{106}}%
\TextAt{21}{373.024}{618.319}{K}%
\TextAt{16}{384.071}{620.178}{\Character{14}}%
\TextAt{13}{392.893}{620.178}{\Character{17}}%
\TextAt{21}{398.669}{618.319}{L}%
\TextAt{3}{404.428}{618.319}{\Character{106}}%
\TextAt{21}{406.78}{618.319}{L}%
\TextAt{23}{412.539}{620.588}{\Character{48}}%
\TextAt{8}{419.894}{620.178}{and}%
\TextAt{8}{442.932}{620.178}{the}%
\TextAt{8}{463.364}{620.178}{formula}%
\TextAt{8}{506.57}{620.178}{for}%
\TextAt{13}{524.73}{620.178}{\Character{32}}%
\TextAt{8}{79.226}{603.132}{follows.}%
\Vector{q 0 G 0 g 0.399 w 452.334 -170.868 m 452.334 -162.798 l S Q}%
\Vector{q 0 G 0 g 0.399 w 452.534 -162.997 m 459.707 -162.997 l S Q}%
\Vector{q 0 G 0 g 0.399 w 452.534 -170.668 m 459.707 -170.668 l S Q}%
\Vector{q 0 G 0 g 0.399 w 459.906 -170.868 m 459.906 -162.798 l S Q}%
\TextAt{8}{97.159}{574.131}{The}%
\TextAt{8}{120.421}{574.131}{advantage}%
\TextAt{8}{174.934}{574.131}{of}%
\TextAt{8}{187.473}{574.131}{the}%
\TextAt{8}{206.837}{574.131}{upper}%
\TextAt{8}{239.532}{574.131}{numbering}%
\TextAt{8}{296.95}{574.131}{of}%
\TextAt{8}{309.489}{574.131}{the}%
\TextAt{8}{328.854}{574.131}{ramification}%
\TextAt{8}{393.11}{574.131}{groups}%
\TextAt{8}{430.109}{574.131}{is}%
\TextAt{8}{441.081}{574.131}{that}%
\TextAt{8}{465.658}{574.131}{it}%
\TextAt{8}{476.571}{574.131}{is}%
\TextAt{8}{487.543}{574.131}{invariant}%
\TextAt{8}{79.226}{557.085}{when}%
\TextAt{8}{109.784}{557.085}{passing}%
\TextAt{8}{150.885}{557.085}{from}%
\TextAt{13}{178.526}{557.085}{L}%
\TextAt{16}{189.811}{557.085}{\Character{106}}%
\TextAt{13}{196.453}{557.085}{K}%
\TextAt{8}{211.168}{557.085}{to}%
\TextAt{8}{225.478}{557.085}{a}%
\TextAt{8}{235.234}{557.085}{Galois}%
\TextAt{8}{271.147}{557.085}{subextension.}%
\TextAt{12}{79.226}{531.073}{Proposition}%
\TextAt{12}{152.202}{531.073}{(10.9).}%
\TextAt{8}{194.478}{531.073}{Let}%
\TextAt{13}{215.677}{531.073}{L}%
\TextAt{3}{223.642}{535.411}{\Character{48}}%
\TextAt{16}{230.149}{531.073}{\Character{106}}%
\TextAt{13}{237.182}{531.073}{K}%
\TextAt{8}{252.127}{531.073}{be}%
\TextAt{8}{268.281}{531.073}{a}%
\TextAt{8}{278.267}{531.073}{Galois}%
\TextAt{8}{314.41}{531.073}{subextension}%
\TextAt{8}{384.35}{531.073}{of}%
\TextAt{13}{397.91}{531.073}{L}%
\TextAt{16}{409.587}{531.073}{\Character{106}}%
\TextAt{13}{416.62}{531.073}{K}%
\TextAt{8}{431.565}{531.073}{and}%
\TextAt{13}{454.557}{531.073}{H}%
\TextAt{9}{468.892}{531.073}{\Character{61}}%
\TextAt{13}{481.71}{531.073}{G}%
\TextAt{9}{490.943}{531.073}{\Character{40}}%
\TextAt{13}{495.495}{531.073}{L}%
\TextAt{16}{507.172}{531.073}{\Character{106}}%
\TextAt{13}{514.206}{531.073}{L}%
\TextAt{3}{522.17}{535.411}{\Character{48}}%
\TextAt{9}{524.965}{531.073}{\Character{41}}%
\TextAt{8}{529.518}{531.073}{.}%
\TextAt{8}{79.226}{514.027}{Then}%
\TextAt{8}{109.784}{514.027}{one}%
\TextAt{8}{131.243}{514.027}{has}%
\TextAt{13}{230.913}{496.981}{G}%
\TextAt{21}{240.147}{501.917}{t}%
\TextAt{9}{243.703}{496.981}{\Character{40}}%
\TextAt{13}{248.255}{496.981}{L}%
\TextAt{16}{259.54}{496.981}{\Character{106}}%
\TextAt{13}{266.182}{496.981}{K}%
\TextAt{9}{277}{496.981}{\Character{41}}%
\TextAt{13}{281.552}{496.981}{H}%
\TextAt{9}{292.175}{496.981}{\Character{47}}%
\TextAt{13}{298.028}{496.981}{H}%
\TextAt{9}{311.971}{496.981}{\Character{61}}%
\TextAt{13}{324.397}{496.981}{G}%
\TextAt{21}{333.63}{501.917}{t}%
\TextAt{9}{337.186}{496.981}{\Character{40}}%
\TextAt{13}{341.739}{496.981}{L}%
\TextAt{3}{349.703}{501.917}{\Character{48}}%
\TextAt{16}{355.819}{496.981}{\Character{106}}%
\TextAt{13}{362.461}{496.981}{K}%
\TextAt{9}{373.279}{496.981}{\Character{41}}%
\TextAt{13}{377.831}{496.981}{\Character{58}}%
\TextAt{19}{79.226}{455.527}{Proof.}%
\TextAt{8}{116}{455.527}{We}%
\TextAt{8}{136.145}{455.527}{put}%
\TextAt{13}{157.604}{455.527}{s}%
\TextAt{9}{166.439}{455.527}{\Character{61}}%
\TextAt{13}{178.865}{455.527}{\Character{32}}%
\TextAt{21}{186.475}{453.667}{L}%
\TextAt{23}{192.234}{455.936}{\Character{48}}%
\TextAt{3}{194.918}{453.667}{\Character{106}}%
\TextAt{21}{197.271}{453.667}{K}%
\TextAt{9}{205.473}{455.527}{\Character{40}}%
\TextAt{13}{210.026}{455.527}{t}%
\TextAt{9}{214.253}{455.527}{\Character{41}}%
\TextAt{8}{218.805}{455.527}{,}%
\TextAt{13}{225.954}{455.527}{G}%
\TextAt{3}{235.188}{459.865}{\Character{48}}%
\TextAt{9}{241.304}{455.527}{\Character{61}}%
\TextAt{13}{253.729}{455.527}{G}%
\TextAt{9}{262.963}{455.527}{\Character{40}}%
\TextAt{13}{267.515}{455.527}{L}%
\TextAt{3}{275.48}{459.865}{\Character{48}}%
\TextAt{16}{281.596}{455.527}{\Character{106}}%
\TextAt{13}{288.237}{455.527}{K}%
\TextAt{9}{299.055}{455.527}{\Character{41}}%
\TextAt{8}{303.607}{455.527}{,}%
\TextAt{8}{310.757}{455.527}{apply}%
\TextAt{8}{342.952}{455.527}{(10.7)}%
\TextAt{8}{376.785}{455.527}{and}%
\TextAt{8}{399.548}{455.527}{(10.8),}%
\TextAt{8}{436.633}{455.527}{and}%
\TextAt{8}{459.395}{455.527}{get}%
\TextAt{13}{204.274}{426.791}{G}%
\TextAt{21}{213.507}{431.727}{t}%
\TextAt{13}{217.063}{426.791}{H}%
\TextAt{9}{227.686}{426.791}{\Character{47}}%
\TextAt{13}{233.539}{426.791}{H}%
\TextAt{9}{247.482}{426.791}{\Character{61}}%
\TextAt{13}{259.908}{426.791}{G}%
\TextAt{21}{269.142}{424.931}{\Character{32}}%
\TextAt{22}{274.676}{423.193}{L}%
\TextAt{23}{279.735}{423.193}{\Character{106}}%
\TextAt{22}{282.004}{423.193}{K}%
\TextAt{17}{289.047}{424.931}{\Character{40}}%
\TextAt{21}{292.341}{424.931}{t}%
\TextAt{17}{295.399}{424.931}{\Character{41}}%
\TextAt{13}{299.19}{426.791}{H}%
\TextAt{9}{309.813}{426.791}{\Character{47}}%
\TextAt{13}{315.666}{426.791}{H}%
\TextAt{9}{329.609}{426.791}{\Character{61}}%
\TextAt{13}{342.035}{426.791}{G}%
\TextAt{3}{351.268}{431.727}{\Character{48}}%
\TextAt{21}{351.268}{423.835}{\Character{17}}%
\TextAt{22}{355.518}{421.544}{L}%
\TextAt{23}{360.576}{421.544}{\Character{106}}%
\TextAt{22}{362.846}{421.544}{L}%
\TextAt{23}{367.904}{423.26}{\Character{48}}%
\TextAt{17}{371.087}{423.835}{\Character{40}}%
\TextAt{21}{374.38}{423.835}{\Character{32}}%
\TextAt{22}{379.915}{422.097}{L}%
\TextAt{23}{384.974}{422.097}{\Character{106}}%
\TextAt{22}{387.243}{422.097}{K}%
\TextAt{17}{394.286}{423.835}{\Character{40}}%
\TextAt{21}{397.579}{423.835}{t}%
\TextAt{17}{400.637}{423.835}{\Character{41}\Character{41}}%
\TextAt{9}{247.482}{404.132}{\Character{61}}%
\TextAt{13}{259.908}{404.132}{G}%
\TextAt{3}{269.142}{409.068}{\Character{48}}%
\TextAt{21}{269.142}{401.177}{\Character{17}}%
\TextAt{22}{273.391}{398.885}{L}%
\TextAt{23}{278.45}{398.885}{\Character{106}}%
\TextAt{22}{280.719}{398.885}{L}%
\TextAt{23}{285.778}{400.601}{\Character{48}}%
\TextAt{17}{288.96}{401.177}{\Character{40}}%
\TextAt{21}{292.253}{401.177}{\Character{32}}%
\TextAt{22}{297.788}{398.885}{L}%
\TextAt{23}{302.847}{398.885}{\Character{106}}%
\TextAt{22}{305.116}{398.885}{L}%
\TextAt{23}{310.175}{400.601}{\Character{48}}%
\TextAt{17}{313.358}{401.177}{\Character{40}}%
\TextAt{21}{316.651}{401.177}{s}%
\TextAt{17}{320.567}{401.177}{\Character{41}\Character{41}}%
\TextAt{9}{330.972}{404.132}{\Character{61}}%
\TextAt{13}{343.398}{404.132}{G}%
\TextAt{3}{352.631}{409.068}{\Character{48}}%
\TextAt{21}{352.631}{401.177}{s}%
\TextAt{9}{360.366}{404.132}{\Character{61}}%
\TextAt{13}{372.792}{404.132}{G}%
\TextAt{3}{382.025}{409.068}{\Character{48}}%
\TextAt{21}{384.322}{409.068}{t}%
\TextAt{13}{387.878}{404.132}{\Character{58}}%
\Vector{q 0 G 0 g 0.399 w 452.334 -397.63 m 452.334 -389.561 l S Q}%
\Vector{q 0 G 0 g 0.399 w 452.534 -389.76 m 459.707 -389.76 l S Q}%
\Vector{q 0 G 0 g 0.399 w 452.534 -397.431 m 459.707 -397.431 l S Q}%
\Vector{q 0 G 0 g 0.399 w 459.906 -397.63 m 459.906 -389.561 l S Q}%
\TextAt{12}{79.226}{341.391}{Exercise}%
\TextAt{12}{133.667}{341.391}{II.10.1.}%
\TextAt{8}{180.449}{341.391}{Let}%
\TextAt{13}{202.496}{341.391}{K}%
\TextAt{9}{218.472}{341.391}{\Character{61}}%
\TextAt{18}{232.734}{341.391}{Q}%
\TextAt{21}{242.033}{339.598}{p}%
\TextAt{8}{251.769}{341.391}{and}%
\TextAt{13}{275.609}{341.391}{K}%
\TextAt{21}{285.568}{339.598}{n}%
\TextAt{9}{296.362}{341.391}{\Character{61}}%
\TextAt{13}{310.625}{341.391}{K}%
\TextAt{9}{321.442}{341.391}{\Character{40}}%
\TextAt{13}{325.995}{341.391}{\Character{16}}%
\TextAt{9}{332.01}{341.391}{\Character{41}}%
\TextAt{8}{336.563}{341.391}{,}%
\TextAt{8}{345.059}{341.391}{where}%
\TextAt{13}{379.946}{341.391}{\Character{16}}%
\TextAt{8}{390.937}{341.391}{is}%
\TextAt{8}{403.779}{341.391}{a}%
\TextAt{8}{414.612}{341.391}{primitive}%
\TextAt{13}{465.77}{341.391}{p}%
\TextAt{21}{471.645}{345.729}{n}%
\TextAt{8}{477.282}{341.391}{-th}%
\TextAt{8}{497.213}{341.391}{root}%
\TextAt{8}{523.337}{341.391}{of}%
\ErrataAt{79.19}{324.338}{455.81}{unity. Show that for integers $s\ge0$ the ramification groups of $K_n\mid K$ are given as follows:}%
\TextAt{13}{209.286}{296.473}{G}%
\TextAt{21}{218.52}{294.679}{s}%
\TextAt{9}{226.255}{296.473}{\Character{61}}%
\TextAt{13}{238.68}{296.473}{G}%
\TextAt{9}{247.914}{296.473}{\Character{40}}%
\TextAt{13}{252.466}{296.473}{K}%
\TextAt{21}{262.426}{294.679}{n}%
\TextAt{16}{271.383}{296.473}{\Character{106}}%
\TextAt{13}{278.024}{296.473}{K}%
\TextAt{9}{288.842}{296.473}{\Character{41}}%
\TextAt{8}{307.231}{296.473}{for}%
\TextAt{13}{325.116}{296.473}{s}%
\TextAt{9}{333.951}{296.473}{\Character{61}}%
\TextAt{9}{346.361}{296.473}{0}%
\TextAt{13}{352.229}{296.473}{\Character{59}}%
\TextAt{13}{209.287}{276.438}{G}%
\TextAt{21}{218.52}{274.645}{s}%
\TextAt{9}{226.255}{276.438}{\Character{61}}%
\TextAt{13}{238.68}{276.438}{G}%
\TextAt{9}{247.914}{276.438}{\Character{40}}%
\TextAt{13}{252.466}{276.438}{K}%
\TextAt{21}{262.426}{274.645}{n}%
\TextAt{16}{271.383}{276.438}{\Character{106}}%
\TextAt{13}{278.024}{276.438}{K}%
\TextAt{17}{287.984}{274.645}{1}%
\TextAt{9}{292.716}{276.438}{\Character{41}}%
\TextAt{8}{307.231}{276.438}{for}%
\TextAt{9}{325.116}{276.438}{1}%
\NativeMathAt{334.29}{276.438}{11.955}{\leq}%
\TextAt{13}{346.909}{276.438}{s}%
\NativeMathAt{355.744}{276.438}{11.955}{\leq}%
\TextAt{13}{368.363}{276.438}{p}%
\TextAt{16}{376.895}{276.438}{\Character{0}}%
\TextAt{9}{388.85}{276.438}{1}%
\TextAt{13}{394.703}{276.438}{\Character{59}}%
\TextAt{13}{209.286}{256.403}{G}%
\TextAt{21}{218.52}{254.61}{s}%
\TextAt{9}{226.255}{256.403}{\Character{61}}%
\TextAt{13}{238.68}{256.403}{G}%
\TextAt{9}{247.914}{256.403}{\Character{40}}%
\TextAt{13}{252.466}{256.403}{K}%
\TextAt{21}{262.426}{254.61}{n}%
\TextAt{16}{271.383}{256.403}{\Character{106}}%
\TextAt{13}{278.024}{256.403}{K}%
\TextAt{17}{287.984}{254.61}{2}%
\TextAt{9}{292.716}{256.403}{\Character{41}}%
\TextAt{8}{307.231}{256.403}{for}%
\TextAt{13}{325.116}{256.403}{p}%
\NativeMathAt{334.312}{256.403}{11.955}{\leq}%
\TextAt{13}{346.931}{256.403}{s}%
\NativeMathAt{355.766}{256.403}{11.955}{\leq}%
\TextAt{13}{368.385}{256.403}{p}%
\TextAt{17}{374.26}{261.339}{2}%
\TextAt{16}{381.649}{256.403}{\Character{0}}%
\TextAt{9}{393.605}{256.403}{1}%
\TextAt{13}{399.458}{256.403}{\Character{59}}%
\TextAt{16}{228.824}{236.368}{\Character{1}}%
\TextAt{16}{234.12}{236.368}{\Character{1}}%
\TextAt{16}{239.428}{236.368}{\Character{1}}%
\TextAt{13}{209.287}{216.334}{G}%
\TextAt{21}{218.52}{214.54}{s}%
\TextAt{9}{226.255}{216.334}{\Character{61}}%
\TextAt{9}{238.665}{216.334}{1}%
\TextAt{8}{307.231}{216.334}{for}%
\TextAt{13}{325.116}{216.334}{p}%
\TextAt{21}{330.991}{221.27}{n}%
\TextAt{3}{336.129}{221.27}{\Character{0}}%
\TextAt{17}{342.716}{221.27}{1}%
\NativeMathAt{350.769}{216.334}{11.955}{\leq}%
\TextAt{13}{363.388}{216.334}{s\Character{58}}%
\TextAt{12}{79.226}{174.358}{Exercise}%
\TextAt{12}{131.326}{174.358}{II.10.2.}%
\TextAt{8}{176.075}{174.358}{Let}%
\TextAt{13}{196.088}{174.358}{K}%
\TextAt{3}{206.905}{178.696}{\Character{48}}%
\TextAt{8}{212.641}{174.358}{be}%
\TextAt{8}{227.609}{174.358}{an}%
\TextAt{8}{242.912}{174.358}{intermediate}%
\TextAt{8}{309.921}{174.358}{field}%
\TextAt{8}{334.321}{174.358}{of}%
\TextAt{13}{346.695}{174.358}{L}%
\TextAt{16}{357.98}{174.358}{\Character{106}}%
\TextAt{13}{364.622}{174.358}{K}%
\TextAt{8}{375.439}{174.358}{.}%
\TextAt{8}{383.569}{174.358}{Describe}%
\TextAt{8}{430.302}{174.358}{the}%
\TextAt{8}{449.502}{174.358}{relation}%
\TextAt{8}{491.476}{174.358}{between}%
\TextAt{8}{79.226}{157.312}{the}%
\TextAt{8}{99.383}{157.312}{ramification}%
\TextAt{8}{164.431}{157.312}{groups}%
\TextAt{8}{202.221}{157.312}{of}%
\TextAt{13}{215.551}{157.312}{L}%
\TextAt{16}{226.836}{157.312}{\Character{106}}%
\TextAt{13}{233.478}{157.312}{K}%
\TextAt{8}{248.193}{157.312}{and}%
\TextAt{8}{270.956}{157.312}{of}%
\TextAt{13}{284.286}{157.312}{L}%
\TextAt{16}{295.571}{157.312}{\Character{106}}%
\TextAt{13}{302.213}{157.312}{K}%
\TextAt{3}{313.03}{161.65}{\Character{48}}%
\TextAt{8}{319.723}{157.312}{in}%
\TextAt{8}{333.376}{157.312}{the}%
\TextAt{8}{353.532}{157.312}{upper}%
\TextAt{8}{387.019}{157.312}{numbering.}%
\TextAt{8}{297.211}{43.033}{166}%
\Destination{78.23}{809.815}{nav.page.166}%
\Destination{79.226}{355.355}{nav.ex.2.10.1}%
\Destination{79.226}{188.355}{nav.ex.2.10.2}%
\Destination{79.36}{589.155}{nav.print.181}%
\Destination{79.36}{545.055}{nav.number.2.10.9}%
\Hotlink{202.287}{741.855}{30.936}{11.047}{nav.number.2.10.7}%
\Hotlink{342.452}{453.355}{30.936}{11.047}{nav.number.2.10.7}%
\Hotlink{399.048}{453.355}{30.936}{11.047}{nav.number.2.10.8}%
\Destination{79.226}{341.338}{errata.81}%
}


% Chapter 3
% Page 179 (body page 167)
\Page{%
\TextAt{1}{79.226}{673.054}{Chapter}%
\TextAt{1}{188.19}{673.054}{III}%
\TextAt{1}{79.226}{614.873}{Riemann–Roch}%
\TextAt{1}{274.176}{614.873}{Theory}%
\TextAt{15}{79.226}{546.17}{§}%
\TextAt{15}{93.785}{546.17}{1}%
\TextAt{15}{120.693}{546.17}{Primes}%
\TextAt{8}{79.226}{514.284}{Having}%
\TextAt{8}{118.2}{514.284}{set}%
\TextAt{8}{135.439}{514.284}{up}%
\TextAt{8}{151.315}{514.284}{the}%
\TextAt{8}{170.443}{514.284}{general}%
\TextAt{8}{209.74}{514.284}{theory}%
\TextAt{8}{245.462}{514.284}{of}%
\TextAt{8}{257.764}{514.284}{valued}%
\TextAt{8}{293.485}{514.284}{fields,}%
\TextAt{8}{325.886}{514.284}{we}%
\TextAt{8}{342.085}{514.284}{now}%
\TextAt{8}{365.446}{514.284}{return}%
\TextAt{8}{400.187}{514.284}{to}%
\TextAt{8}{413.469}{514.284}{algebraic}%
\TextAt{8}{461.875}{514.284}{number}%
\TextAt{8}{503.443}{514.284}{fields.}%
\TextAt{8}{79.226}{497.238}{We}%
\TextAt{8}{99.921}{497.238}{want}%
\TextAt{8}{129.092}{497.238}{to}%
\TextAt{8}{143.952}{497.238}{develop}%
\TextAt{8}{186.776}{497.238}{their}%
\TextAt{8}{215.289}{497.238}{basic}%
\TextAt{8}{245.166}{497.238}{theory}%
\TextAt{8}{282.466}{497.238}{from}%
\TextAt{8}{310.656}{497.238}{the}%
\TextAt{8}{331.363}{497.238}{valuation-theoretic}%
\TextAt{8}{431.763}{497.238}{point}%
\TextAt{8}{462.882}{497.238}{of}%
\TextAt{8}{476.763}{497.238}{view.}%
\TextAt{8}{509.947}{497.238}{This}%
\TextAt{8}{79.226}{480.192}{approach}%
\TextAt{8}{129.484}{480.192}{has}%
\TextAt{8}{150.202}{480.192}{two}%
\TextAt{8}{172.162}{480.192}{prominent}%
\TextAt{8}{228.266}{480.192}{advantages}%
\TextAt{8}{288.029}{480.192}{compared}%
\TextAt{8}{341.205}{480.192}{to}%
\TextAt{8}{355.358}{480.192}{the}%
\TextAt{8}{375.358}{480.192}{ideal-theoretic}%
\TextAt{8}{451.942}{480.192}{treatment}%
\TextAt{8}{506.097}{480.192}{given}%
\TextAt{8}{79.226}{463.147}{in}%
\TextAt{8}{93.266}{463.147}{the}%
\TextAt{8}{113.81}{463.147}{first}%
\TextAt{8}{138.323}{463.147}{chapter.}%
\TextAt{8}{185.986}{463.147}{The}%
\TextAt{8}{210.427}{463.147}{first}%
\TextAt{8}{234.94}{463.147}{one}%
\TextAt{8}{256.787}{463.147}{results}%
\TextAt{8}{294.367}{463.147}{from}%
\TextAt{8}{322.395}{463.147}{the}%
\TextAt{8}{342.938}{463.147}{possibility}%
\TextAt{8}{399.061}{463.147}{of}%
\TextAt{8}{412.778}{463.147}{passing}%
\TextAt{8}{454.267}{463.147}{to}%
\TextAt{8}{468.965}{463.147}{completions,}%
\TextAt{8}{79.226}{446.101}{thereby}%
\TextAt{8}{121.965}{446.101}{enacting}%
\TextAt{8}{169.259}{446.101}{the}%
\TextAt{8}{189.881}{446.101}{important}%
\TextAt{8}{245.64}{446.101}{“local-to-global}%
\TextAt{8}{328.058}{446.101}{principle”.}%
\TextAt{8}{386.023}{446.101}{This}%
\TextAt{8}{413.208}{446.101}{will}%
\TextAt{8}{435.779}{446.101}{be}%
\TextAt{8}{452.169}{446.101}{done}%
\TextAt{8}{480.598}{446.101}{in}%
\TextAt{8}{494.716}{446.101}{chapter}%
\TextAt{8}{79.226}{429.055}{VI.}%
\TextAt{8}{99.843}{429.055}{The}%
\TextAt{8}{124.369}{429.055}{other}%
\TextAt{8}{155.41}{429.055}{advantage}%
\TextAt{8}{211.187}{429.055}{lies}%
\TextAt{8}{231.876}{429.055}{in}%
\TextAt{8}{246.001}{429.055}{the}%
\TextAt{8}{266.629}{429.055}{fact}%
\TextAt{8}{290.187}{429.055}{that}%
\TextAt{8}{316.028}{429.055}{the}%
\TextAt{8}{336.657}{429.055}{archimedean}%
\TextAt{8}{405.094}{429.055}{valuations}%
\TextAt{8}{461.254}{429.055}{bring}%
\TextAt{8}{492.295}{429.055}{into}%
\TextAt{8}{516.51}{429.055}{the}%
\TextAt{8}{79.226}{412.009}{picture}%
\TextAt{8}{118.997}{412.009}{the}%
\TextAt{8}{139.258}{412.009}{points}%
\TextAt{8}{174.545}{412.009}{at}%
\TextAt{8}{188.96}{412.009}{infinity,}%
\TextAt{8}{231.686}{412.009}{which}%
\TextAt{8}{265.276}{412.009}{were}%
\TextAt{8}{292.363}{412.009}{hitherto}%
\TextAt{8}{337.346}{412.009}{lacking,}%
\TextAt{8}{380.407}{412.009}{as}%
\TextAt{8}{394.881}{412.009}{the}%
\TextAt{8}{415.141}{412.009}{“primes}%
\TextAt{8}{458.547}{412.009}{at}%
\TextAt{8}{472.961}{412.009}{infinity”.}%
\TextAt{8}{522.046}{412.009}{In}%
\TextAt{8}{79.226}{394.963}{this}%
\TextAt{8}{102.115}{394.963}{way}%
\TextAt{8}{125.925}{394.963}{a}%
\TextAt{8}{135.747}{394.963}{perfect}%
\TextAt{8}{174.823}{394.963}{analogy}%
\TextAt{8}{218.156}{394.963}{with}%
\TextAt{8}{244.882}{394.963}{the}%
\TextAt{8}{265.105}{394.963}{function}%
\TextAt{8}{311.02}{394.963}{fields}%
\TextAt{8}{341.058}{394.963}{is}%
\TextAt{8}{352.889}{394.963}{achieved.}%
\TextAt{8}{404.791}{394.963}{This}%
\TextAt{8}{431.577}{394.963}{is}%
\TextAt{8}{443.408}{394.963}{the}%
\TextAt{8}{463.631}{394.963}{task}%
\TextAt{8}{488.803}{394.963}{to}%
\TextAt{8}{503.18}{394.963}{which}%
\TextAt{8}{79.226}{377.917}{we}%
\TextAt{8}{96.454}{377.917}{now}%
\TextAt{8}{120.842}{377.917}{turn.}%
\TextAt{12}{79.226}{352.016}{Definition}%
\TextAt{12}{143.273}{352.016}{(1.1).}%
\TextAt{8}{179.683}{352.016}{A}%
\TextAt{12}{193.448}{352.016}{prime}%
\TextAt{8}{232.499}{352.016}{(or}%
\TextAt{12}{252.459}{352.016}{place}%
\TextAt{8}{282.323}{352.016}{)}%
\TextAt{14}{291.869}{352.016}{p}%
\TextAt{8}{302.843}{352.016}{of}%
\TextAt{8}{317.267}{352.016}{an}%
\TextAt{8}{334.62}{352.016}{algebraic}%
\TextAt{8}{385.149}{352.016}{number}%
\TextAt{8}{428.839}{352.016}{field}%
\TextAt{13}{455.29}{352.016}{K}%
\TextAt{8}{471.099}{352.016}{is}%
\TextAt{8}{483.956}{352.016}{a}%
\TextAt{8}{494.806}{352.016}{class}%
\TextAt{8}{523.337}{352.016}{of}%
\TextAt{8}{79.226}{334.971}{equivalent}%
\TextAt{8}{136.258}{334.971}{valuations}%
\TextAt{8}{193.361}{334.971}{of}%
\TextAt{13}{208.107}{334.971}{K}%
\TextAt{8}{218.925}{334.971}{.}%
\TextAt{8}{231.622}{334.971}{The}%
\TextAt{8}{257.092}{334.971}{nonarchimedean}%
\TextAt{8}{345.338}{334.971}{equivalence}%
\TextAt{8}{408.539}{334.971}{classes}%
\TextAt{8}{447.207}{334.971}{are}%
\TextAt{8}{468.134}{334.971}{called}%
\TextAt{12}{502.714}{334.971}{finite}%
\TextAt{12}{79.226}{317.925}{primes}%
\TextAt{8}{122.492}{317.925}{and}%
\TextAt{8}{145.255}{317.925}{the}%
\TextAt{8}{165.411}{317.925}{archimedean}%
\TextAt{8}{233.376}{317.925}{ones}%
\TextAt{12}{259.45}{317.925}{infinite}%
\TextAt{12}{305.203}{317.925}{primes}%
\TextAt{8}{344.571}{317.925}{.}%
\TextAt{8}{97.159}{292.024}{The}%
\TextAt{8}{121.155}{292.024}{infinite}%
\TextAt{8}{160.765}{292.024}{primes}%
\TextAt{14}{198.487}{292.024}{p}%
\TextAt{8}{208.31}{292.024}{are}%
\TextAt{8}{227.763}{292.024}{obtained,}%
\TextAt{8}{279.101}{292.024}{according}%
\TextAt{8}{331.73}{292.024}{to}%
\TextAt{8}{345.983}{292.024}{chap.}%
\TextAt{8}{378.156}{292.024}{II,}%
\TextAt{8}{394.011}{292.024}{(8.1),}%
\TextAt{8}{425.192}{292.024}{from}%
\TextAt{8}{452.775}{292.024}{the}%
\TextAt{8}{472.874}{292.024}{embeddings}%
\TextAt{13}{79.226}{274.978}{\Character{28}}%
\TextAt{9}{88.966}{274.978}{\Character{58}}%
\TextAt{13}{95.538}{274.978}{K}%
\NativeMathAt{109.677}{274.978}{11.955}{\rightarrow}%
\TextAt{18}{124.953}{274.978}{C}%
\TextAt{8}{133.587}{274.978}{.}%
\TextAt{8}{141.685}{274.978}{There}%
\TextAt{8}{174.443}{274.978}{are}%
\TextAt{8}{192.902}{274.978}{two}%
\TextAt{8}{213.968}{274.978}{sorts}%
\TextAt{8}{241.011}{274.978}{of}%
\TextAt{8}{253.289}{274.978}{these:}%
\TextAt{8}{287.286}{274.978}{the}%
\TextAt{12}{306.391}{274.978}{real}%
\TextAt{12}{331.578}{274.978}{primes}%
\TextAt{8}{370.947}{274.978}{,}%
\TextAt{8}{377.255}{274.978}{which}%
\TextAt{8}{409.689}{274.978}{are}%
\TextAt{8}{428.149}{274.978}{given}%
\TextAt{8}{457.667}{274.978}{by}%
\TextAt{8}{472.874}{274.978}{embeddings}%
\TextAt{13}{79.226}{257.932}{\Character{28}}%
\TextAt{9}{91.838}{257.932}{\Character{58}}%
\TextAt{13}{101.282}{257.932}{K}%
\NativeMathAt{118.292}{257.932}{11.955}{\rightarrow}%
\TextAt{18}{136.44}{257.932}{R}%
\TextAt{8}{145.074}{257.932}{,}%
\TextAt{8}{154.33}{257.932}{and}%
\TextAt{8}{178.777}{257.932}{the}%
\TextAt{12}{200.619}{257.932}{complex}%
\TextAt{12}{255.411}{257.932}{primes}%
\TextAt{8}{294.779}{257.932}{,}%
\TextAt{8}{304.035}{257.932}{which}%
\TextAt{8}{339.207}{257.932}{are}%
\TextAt{8}{360.403}{257.932}{induced}%
\TextAt{8}{405.652}{257.932}{by}%
\TextAt{8}{423.597}{257.932}{the}%
\TextAt{8}{445.438}{257.932}{pairs}%
\TextAt{8}{475.804}{257.932}{of}%
\TextAt{8}{490.819}{257.932}{complex}%
\TextAt{8}{79.226}{240.886}{conjugate}%
\TextAt{8}{133.624}{240.886}{non-real}%
\TextAt{8}{180.537}{240.886}{embeddings}%
\TextAt{13}{245.718}{240.886}{K}%
\NativeMathAt{262.222}{240.886}{11.955}{\rightarrow}%
\TextAt{18}{279.864}{240.886}{C}%
\TextAt{8}{288.498}{240.886}{.}%
\TextAt{14}{301.111}{240.886}{p}%
\TextAt{8}{312.38}{240.886}{is}%
\TextAt{8}{325.532}{240.886}{real}%
\TextAt{8}{349.682}{240.886}{or}%
\TextAt{8}{365.381}{240.886}{complex}%
\TextAt{8}{412.617}{240.886}{depending}%
\TextAt{8}{470.255}{240.886}{whether}%
\TextAt{8}{516.51}{240.886}{the}%
\TextAt{8}{79.226}{223.84}{completion}%
\TextAt{13}{139.178}{223.84}{K}%
\TextAt{27}{149.137}{222.047}{p}%
\TextAt{8}{157.637}{223.84}{is}%
\TextAt{8}{169.517}{223.84}{isomorphic}%
\TextAt{8}{228.883}{223.84}{to}%
\TextAt{18}{243.31}{223.84}{R}%
\TextAt{8}{255.958}{223.84}{or}%
\TextAt{8}{270.384}{223.84}{to}%
\TextAt{18}{284.81}{223.84}{C}%
\TextAt{8}{293.445}{223.84}{.}%
\TextAt{8}{302.241}{223.84}{The}%
\TextAt{8}{326.411}{223.84}{infinite}%
\TextAt{8}{366.195}{223.84}{primes}%
\TextAt{8}{404.089}{223.84}{will}%
\TextAt{8}{426.31}{223.84}{be}%
\TextAt{8}{442.351}{223.84}{referred}%
\TextAt{8}{485.709}{223.84}{to}%
\TextAt{8}{500.135}{223.84}{by}%
\TextAt{8}{516.51}{223.84}{the}%
\TextAt{8}{79.226}{206.795}{formal}%
\TextAt{8}{115.976}{206.795}{notation}%
\TextAt{14}{162.817}{206.795}{p}%
\TextAt{16}{172.121}{206.795}{\Character{106}}%
\NativeMathAt{178.744}{206.795}{11.955}{\infty}%
\TextAt{8}{190.718}{206.795}{,}%
\TextAt{8}{197.867}{206.795}{the}%
\TextAt{8}{218.023}{206.795}{finite}%
\TextAt{8}{247.935}{206.795}{ones}%
\TextAt{8}{274.009}{206.795}{by}%
\TextAt{14}{290.268}{206.795}{p}%
\TextAt{18}{299.572}{206.795}{\Character{45}}%
\NativeMathAt{306.214}{206.795}{11.955}{\infty}%
\TextAt{8}{318.169}{206.795}{.}%
\TextAt{8}{97.159}{186.76}{In}%
\TextAt{8}{112.428}{186.76}{the}%
\TextAt{8}{133.232}{186.76}{case}%
\TextAt{8}{158.651}{186.76}{of}%
\TextAt{8}{172.629}{186.76}{a}%
\TextAt{8}{183.032}{186.76}{finite}%
\TextAt{8}{213.592}{186.76}{prime,}%
\TextAt{8}{250.817}{186.76}{the}%
\TextAt{8}{271.621}{186.76}{letter}%
\TextAt{14}{303.484}{186.76}{p}%
\TextAt{8}{314.012}{186.76}{has}%
\TextAt{8}{335.534}{186.76}{a}%
\TextAt{8}{345.937}{186.76}{multiple}%
\TextAt{8}{392.433}{186.76}{meaning:}%
\TextAt{8}{445.108}{186.76}{it}%
\TextAt{8}{457.459}{186.76}{also}%
\TextAt{8}{481.587}{186.76}{stands}%
\TextAt{8}{518.782}{186.76}{for}%
\TextAt{8}{79.226}{169.714}{the}%
\TextAt{8}{99.579}{169.714}{prime}%
\TextAt{8}{132.938}{169.714}{ideal}%
\TextAt{8}{161.098}{169.714}{of}%
\TextAt{8}{174.624}{169.714}{the}%
\TextAt{8}{194.976}{169.714}{ring}%
\TextAt{14}{219.238}{169.714}{o}%
\TextAt{8}{229.176}{169.714}{of}%
\TextAt{8}{242.702}{169.714}{integers}%
\TextAt{8}{286.211}{169.714}{of}%
\TextAt{13}{299.737}{169.714}{K}%
\TextAt{8}{310.555}{169.714}{,}%
\TextAt{8}{317.949}{169.714}{or}%
\TextAt{8}{332.456}{169.714}{for}%
\TextAt{8}{350.537}{169.714}{the}%
\TextAt{8}{370.889}{169.714}{maximal}%
\TextAt{8}{418.894}{169.714}{ideal}%
\TextAt{8}{447.053}{169.714}{of}%
\TextAt{8}{460.579}{169.714}{the}%
\TextAt{8}{480.932}{169.714}{associated}%
\TextAt{8}{79.226}{152.668}{valuation}%
\TextAt{8}{129.156}{152.668}{ring,}%
\TextAt{8}{155.559}{152.668}{or}%
\TextAt{8}{168.726}{152.668}{even}%
\TextAt{8}{194.243}{152.668}{for}%
\TextAt{8}{210.984}{152.668}{the}%
\TextAt{8}{229.998}{152.668}{maximal}%
\TextAt{8}{276.663}{152.668}{ideal}%
\TextAt{8}{303.483}{152.668}{of}%
\TextAt{8}{315.67}{152.668}{the}%
\TextAt{8}{334.684}{152.668}{completion.}%
\TextAt{8}{398.689}{152.668}{However,}%
\TextAt{8}{448.178}{152.668}{this}%
\TextAt{8}{469.857}{152.668}{will}%
\TextAt{8}{490.819}{152.668}{nowhere}%
\TextAt{8}{79.226}{135.622}{create}%
\TextAt{8}{112.731}{135.622}{any}%
\TextAt{8}{133.886}{135.622}{risk}%
\TextAt{8}{155.424}{135.622}{of}%
\TextAt{8}{167.793}{135.622}{confusion.}%
\TextAt{8}{223.789}{135.622}{We}%
\TextAt{8}{242.971}{135.622}{write}%
\TextAt{14}{271.921}{135.622}{p}%
\TextAt{16}{281.225}{135.622}{\Character{106}}%
\TextAt{13}{287.867}{135.622}{p}%
\TextAt{8}{296.678}{135.622}{if}%
\TextAt{13}{306.44}{135.622}{p}%
\TextAt{8}{315.251}{135.622}{is}%
\TextAt{8}{326.053}{135.622}{the}%
\TextAt{8}{345.247}{135.622}{characteristic}%
\TextAt{8}{416.22}{135.622}{of}%
\TextAt{8}{428.588}{135.622}{the}%
\TextAt{8}{447.783}{135.622}{residue}%
\TextAt{8}{486.548}{135.622}{field}%
\TextAt{13}{510.943}{135.622}{\Character{20}}%
\TextAt{9}{517.682}{135.622}{\Character{40}}%
\TextAt{14}{522.234}{135.622}{p}%
\TextAt{9}{528.217}{135.622}{\Character{41}}%
\TextAt{8}{79.226}{118.576}{of}%
\TextAt{8}{92.669}{118.576}{the}%
\TextAt{8}{112.939}{118.576}{finite}%
\TextAt{8}{142.964}{118.576}{prime}%
\TextAt{14}{176.241}{118.576}{p}%
\TextAt{8}{182.224}{118.576}{.}%
\TextAt{8}{191.011}{118.576}{For}%
\TextAt{8}{212.094}{118.576}{an}%
\TextAt{8}{228.466}{118.576}{infinite}%
\TextAt{8}{268.246}{118.576}{prime}%
\TextAt{8}{301.523}{118.576}{we}%
\TextAt{8}{318.864}{118.576}{adopt}%
\TextAt{8}{352.153}{118.576}{the}%
\TextAt{8}{372.422}{118.576}{convention}%
\TextAt{8}{431.08}{118.576}{that}%
\TextAt{8}{456.562}{118.576}{the}%
\TextAt{8}{476.831}{118.576}{completion}%
\TextAt{13}{79.226}{101.53}{K}%
\TextAt{27}{89.185}{99.737}{p}%
\TextAt{8}{97.569}{101.53}{also}%
\TextAt{8}{121.049}{101.53}{serves}%
\TextAt{8}{154.99}{101.53}{as}%
\TextAt{8}{169.36}{101.53}{its}%
\TextAt{8}{185.679}{101.53}{own}%
\TextAt{8}{210.068}{101.53}{“residue}%
\TextAt{8}{255.318}{101.53}{field,”}%
\TextAt{8}{289.45}{101.53}{i.e.,}%
\TextAt{8}{311.555}{101.53}{we}%
\TextAt{8}{328.782}{101.53}{put}%
\TextAt{13}{232.931}{72.921}{\Character{20}}%
\TextAt{9}{239.669}{72.921}{\Character{40}}%
\TextAt{14}{244.221}{72.921}{p}%
\TextAt{9}{250.204}{72.921}{\Character{41}}%
\TextAt{9}{258.059}{72.921}{\Character{58}\Character{61}}%
\TextAt{13}{273.755}{72.921}{K}%
\TextAt{27}{283.714}{71.128}{p}%
\TextAt{13}{288.201}{72.921}{\Character{59}}%
\TextAt{8}{317.355}{72.921}{when}%
\TextAt{14}{347.913}{72.921}{p}%
\TextAt{16}{357.217}{72.921}{\Character{106}}%
\NativeMathAt{363.84}{72.921}{11.955}{\infty}%
\TextAt{13}{375.813}{72.921}{\Character{58}}%
\TextAt{8}{297.211}{43.033}{167}%
\Destination{79.23}{773.075}{nav.chapter.3}%
\Destination{78.23}{809.815}{nav.page.167}%
\Destination{79.23}{570.065}{nav.section.3.1}%
\Destination{79.36}{629.855}{nav.print.183}%
\Destination{79.36}{366.055}{nav.number.3.1.1}%
\Destination{79.23}{773.075}{outline.31}%
\Destination{79.23}{570.065}{outline.32}%
\Hotlink{377.656}{289.855}{9.751}{11.047}{nav.chapter.2}%
\Hotlink{394.011}{289.824}{28}{11.047}{nav.number.2.8.1}%
}

% Page 180 (body page 168)
\Page{%
\TextAt{2}{79.226}{796.943}{III.}%
\TextAt{2}{98.317}{796.943}{Riemann–Roch}%
\TextAt{2}{174.375}{796.943}{Theory}%
\Vector{q 0 G 0 g 0.299 w 7.226 18.729 m 460.769 18.729 l S Q}%
\TextAt{8}{79.226}{761.124}{To}%
\TextAt{8}{96.454}{761.124}{each}%
\TextAt{8}{122.791}{761.124}{prime}%
\TextAt{14}{155.954}{761.124}{p}%
\TextAt{8}{165.835}{761.124}{of}%
\TextAt{13}{179.165}{761.124}{K}%
\TextAt{8}{193.88}{761.124}{we}%
\TextAt{8}{211.107}{761.124}{now}%
\TextAt{8}{235.496}{761.124}{associate}%
\TextAt{8}{284.727}{761.124}{a}%
\TextAt{8}{294.483}{761.124}{canonical}%
\TextAt{8}{345.866}{761.124}{homomorphism}%
\TextAt{13}{270.935}{732.123}{v}%
\TextAt{27}{276.594}{730.329}{p}%
\TextAt{9}{284.401}{732.123}{\Character{58}}%
\TextAt{13}{290.974}{732.123}{K}%
\TextAt{3}{301.792}{737.059}{\Character{3}}%
\NativeArrowAt{309.845}{732.123}{11.955}{19.26}{\longrightarrow}%
\TextAt{18}{332.427}{732.123}{R}%
\TextAt{8}{79.226}{703.121}{from}%
\TextAt{8}{107.64}{703.121}{the}%
\TextAt{8}{128.57}{703.121}{multiplicative}%
\TextAt{8}{203.167}{703.121}{group}%
\TextAt{13}{237.117}{703.121}{K}%
\TextAt{3}{247.934}{707.46}{\Character{3}}%
\TextAt{8}{257.338}{703.121}{of}%
\TextAt{13}{271.442}{703.121}{K}%
\TextAt{8}{282.259}{703.121}{.}%
\TextAt{8}{293.029}{703.121}{If}%
\TextAt{14}{305.495}{703.121}{p}%
\TextAt{8}{316.149}{703.121}{is}%
\TextAt{8}{328.687}{703.121}{finite,}%
\TextAt{8}{362.818}{703.121}{then}%
\TextAt{13}{390.252}{703.121}{v}%
\TextAt{27}{395.911}{701.328}{p}%
\TextAt{8}{405.069}{703.121}{is}%
\TextAt{8}{417.606}{703.121}{the}%
\TextAt{14}{438.537}{703.121}{p}%
\TextAt{8}{444.52}{703.121}{-adic}%
\TextAt{8}{473.902}{703.121}{exponential}%
\TextAt{8}{79.226}{686.075}{valuation}%
\TextAt{8}{131.204}{686.075}{which}%
\TextAt{8}{165.595}{686.075}{is}%
\TextAt{8}{178.265}{686.075}{normalized}%
\TextAt{8}{239.005}{686.075}{by}%
\TextAt{8}{256.17}{686.075}{the}%
\TextAt{8}{277.231}{686.075}{condition}%
\TextAt{13}{329.52}{686.075}{v}%
\TextAt{27}{335.179}{684.282}{p}%
\TextAt{9}{339.666}{686.075}{\Character{40}}%
\TextAt{13}{344.218}{686.075}{K}%
\TextAt{3}{355.036}{690.414}{\Character{3}}%
\TextAt{9}{359.768}{686.075}{\Character{41}}%
\TextAt{9}{369.165}{686.075}{\Character{61}}%
\TextAt{18}{383.152}{686.075}{Z}%
\TextAt{8}{391.122}{686.075}{.}%
\TextAt{8}{402.286}{686.075}{If}%
\TextAt{14}{414.884}{686.075}{p}%
\TextAt{8}{425.669}{686.075}{is}%
\TextAt{8}{438.338}{686.075}{infinite,}%
\TextAt{8}{482.389}{686.075}{then}%
\TextAt{13}{509.954}{686.075}{v}%
\TextAt{27}{515.614}{684.282}{p}%
\TextAt{8}{524.903}{686.075}{is}%
\TextAt{8}{79.226}{669.029}{given}%
\TextAt{8}{109.796}{669.029}{by}%
\TextAt{13}{260.072}{651.983}{v}%
\TextAt{27}{265.732}{650.19}{p}%
\TextAt{9}{270.219}{651.983}{\Character{40}}%
\TextAt{13}{274.771}{651.983}{a}%
\TextAt{9}{280.916}{651.983}{\Character{41}}%
\TextAt{9}{288.771}{651.983}{\Character{61}}%
\TextAt{16}{301.214}{651.983}{\Character{0}}%
\TextAt{8}{312.505}{651.983}{log}%
\TextAt{16}{329.466}{651.983}{\Character{106}}%
\TextAt{13}{332.787}{651.983}{\Character{28}}%
\TextAt{13}{339.195}{651.983}{a}%
\TextAt{16}{345.351}{651.983}{\Character{106}}%
\TextAt{13}{348.672}{651.983}{\Character{59}}%
\TextAt{8}{79.226}{628.462}{where}%
\TextAt{13}{113.035}{628.462}{\Character{28}}%
\TextAt{9}{122.775}{628.462}{\Character{58}}%
\TextAt{13}{129.348}{628.462}{K}%
\NativeMathAt{143.486}{628.462}{11.955}{\rightarrow}%
\TextAt{18}{158.762}{628.462}{C}%
\TextAt{8}{171.294}{628.462}{is}%
\TextAt{8}{183.058}{628.462}{an}%
\TextAt{8}{199.317}{628.462}{embedding}%
\TextAt{8}{258.495}{628.462}{which}%
\TextAt{8}{291.981}{628.462}{defines}%
\TextAt{14}{330.405}{628.462}{p}%
\TextAt{8}{336.388}{628.462}{.}%
\TextAt{8}{97.159}{608.427}{For}%
\TextAt{8}{118.128}{608.427}{an}%
\TextAt{8}{134.387}{608.427}{arbitrary}%
\TextAt{8}{184.157}{608.427}{prime}%
\TextAt{14}{217.32}{608.427}{p}%
\TextAt{16}{226.624}{608.427}{\Character{106}}%
\TextAt{13}{233.266}{608.427}{p}%
\TextAt{8}{243.038}{608.427}{(}%
\TextAt{13}{247.593}{608.427}{p}%
\TextAt{8}{257.366}{608.427}{prime}%
\TextAt{8}{290.53}{608.427}{number}%
\TextAt{8}{333.126}{608.427}{or}%
\TextAt{13}{347.436}{608.427}{p}%
\TextAt{9}{356.632}{608.427}{\Character{61}}%
\NativeMathAt{369.058}{608.427}{11.955}{\infty}%
\TextAt{8}{381.013}{608.427}{)}%
\TextAt{8}{389.465}{608.427}{we}%
\TextAt{8}{406.693}{608.427}{put}%
\TextAt{8}{428.152}{608.427}{furthermore}%
\TextAt{13}{261.4}{579.426}{f}%
\TextAt{27}{267.17}{577.633}{p}%
\TextAt{9}{274.977}{579.426}{\Character{61}}%
\TextAt{9}{287.387}{579.426}{\Character{91}}%
\TextAt{13}{290.655}{579.426}{\Character{20}}%
\TextAt{9}{297.393}{579.426}{\Character{40}}%
\TextAt{14}{301.945}{579.426}{p}%
\TextAt{9}{307.929}{579.426}{\Character{41}}%
\TextAt{9}{315.784}{579.426}{\Character{58}}%
\TextAt{13}{322.374}{579.426}{\Character{20}}%
\TextAt{9}{329.113}{579.426}{\Character{40}}%
\TextAt{13}{333.665}{579.426}{p}%
\TextAt{9}{339.54}{579.426}{\Character{41}\Character{93}}%
\TextAt{13}{347.344}{579.426}{\Character{59}}%
\TextAt{8}{79.226}{550.425}{so}%
\TextAt{8}{93.596}{550.425}{that}%
\TextAt{13}{118.965}{550.425}{f}%
\TextAt{27}{124.735}{548.631}{p}%
\TextAt{9}{132.543}{550.425}{\Character{61}}%
\TextAt{9}{144.953}{550.425}{\Character{91}}%
\TextAt{13}{148.22}{550.425}{K}%
\TextAt{27}{158.179}{548.631}{p}%
\TextAt{9}{165.987}{550.425}{\Character{58}}%
\TextAt{18}{172.559}{550.425}{R}%
\TextAt{9}{181.194}{550.425}{\Character{93}}%
\TextAt{8}{188.343}{550.425}{if}%
\TextAt{14}{199.066}{550.425}{p}%
\TextAt{16}{208.37}{550.425}{\Character{106}}%
\NativeMathAt{214.993}{550.425}{11.955}{\infty}%
\TextAt{8}{226.967}{550.425}{,}%
\TextAt{8}{234.116}{550.425}{and}%
\TextAt{14}{249.254}{509.665}{N}%
\TextAt{9}{259.202}{509.665}{\Character{40}}%
\TextAt{14}{263.754}{509.665}{p}%
\TextAt{9}{269.737}{509.665}{\Character{41}}%
\TextAt{9}{277.592}{509.665}{\Character{61}}%
\TextAt{20}{290.036}{534.173}{\Character{56}}%
\TextAt{20}{290.036}{523.414}{\Character{60}}%
\TextAt{20}{290.036}{501.894}{\Character{58}}%
\TextAt{13}{300.663}{518.79}{p}%
\TextAt{21}{306.538}{523.129}{f}%
\TextAt{28}{310.606}{522.132}{p}%
\TextAt{13}{315.187}{518.79}{\Character{59}}%
\TextAt{14}{330.394}{518.79}{p}%
\TextAt{18}{339.698}{518.79}{\Character{45}}%
\NativeMathAt{346.339}{518.79}{11.955}{\infty}%
\TextAt{13}{358.295}{518.79}{\Character{59}}%
\TextAt{13}{300.663}{498.335}{e}%
\TextAt{21}{306.088}{502.674}{f}%
\TextAt{28}{310.156}{501.677}{p}%
\TextAt{13}{314.737}{498.335}{\Character{59}}%
\TextAt{14}{330.394}{498.335}{p}%
\TextAt{16}{339.698}{498.335}{\Character{106}}%
\NativeMathAt{346.321}{498.335}{11.955}{\infty}%
\TextAt{13}{358.295}{498.335}{\Character{58}}%
\TextAt{8}{79.226}{469.766}{This}%
\TextAt{8}{106.31}{469.766}{convention}%
\TextAt{8}{165.219}{469.766}{suggests}%
\TextAt{8}{211.3}{469.766}{that}%
\TextAt{8}{237.033}{469.766}{we}%
\TextAt{8}{254.625}{469.766}{consider}%
\TextAt{13}{300.574}{469.766}{e}%
\TextAt{8}{310.261}{469.766}{as}%
\TextAt{8}{324.995}{469.766}{being}%
\TextAt{8}{356.897}{469.766}{an}%
\TextAt{8}{373.521}{469.766}{infinite}%
\TextAt{8}{413.552}{469.766}{prime}%
\TextAt{8}{447.08}{469.766}{number,}%
\TextAt{8}{493.384}{469.766}{and}%
\TextAt{8}{516.51}{469.766}{the}%
\TextAt{8}{79.226}{452.72}{extension}%
\TextAt{18}{131.288}{452.72}{C}%
\TextAt{16}{143.746}{452.72}{\Character{106}}%
\TextAt{18}{150.891}{452.72}{R}%
\TextAt{8}{163.719}{452.72}{as}%
\TextAt{8}{178.384}{452.72}{being}%
\TextAt{19}{210.217}{452.72}{unramified}%
\TextAt{8}{268.663}{452.72}{with}%
\TextAt{8}{295.618}{452.72}{inertia}%
\TextAt{8}{332.987}{452.72}{degree}%
\TextAt{9}{369.698}{452.72}{2}%
\TextAt{8}{375.551}{452.72}{.}%
\TextAt{8}{384.886}{452.72}{We}%
\TextAt{8}{405.326}{452.72}{define}%
\TextAt{8}{439.43}{452.72}{the}%
\TextAt{8}{459.882}{452.72}{absolute}%
\TextAt{8}{506.42}{452.72}{value}%
\TextAt{16}{79.226}{435.674}{\Character{106}}%
\TextAt{16}{86.435}{435.674}{\Character{106}}%
\TextAt{27}{89.765}{433.881}{p}%
\TextAt{9}{97.573}{435.674}{\Character{58}}%
\TextAt{13}{104.145}{435.674}{K}%
\NativeMathAt{118.284}{435.674}{11.955}{\rightarrow}%
\TextAt{18}{133.56}{435.674}{R}%
\TextAt{8}{146.092}{435.674}{by}%
\TextAt{16}{263.782}{418.628}{\Character{106}}%
\TextAt{13}{267.103}{418.628}{a}%
\TextAt{16}{273.248}{418.628}{\Character{106}}%
\TextAt{27}{276.569}{416.835}{p}%
\TextAt{9}{284.377}{418.628}{\Character{61}}%
\TextAt{14}{296.802}{418.628}{N}%
\TextAt{9}{306.75}{418.628}{\Character{40}}%
\TextAt{14}{311.303}{418.628}{p}%
\TextAt{9}{317.286}{418.628}{\Character{41}}%
\TextAt{3}{321.838}{423.565}{\Character{0}}%
\TextAt{21}{328.425}{423.565}{v}%
\TextAt{28}{332.548}{422.568}{p}%
\TextAt{17}{336.631}{423.565}{\Character{40}}%
\TextAt{21}{339.924}{423.565}{a}%
\TextAt{17}{344.422}{423.565}{\Character{41}}%
\TextAt{8}{79.226}{395.107}{for}%
\TextAt{13}{97.142}{395.107}{a}%
\NativeMathAt{106.66}{395.107}{11.955}{\neq}%
\TextAt{9}{119.129}{395.107}{0}%
\TextAt{8}{128.918}{395.107}{and}%
\TextAt{16}{151.712}{395.107}{\Character{106}}%
\TextAt{9}{155.033}{395.107}{0}%
\TextAt{16}{160.886}{395.107}{\Character{106}}%
\TextAt{27}{164.207}{393.313}{p}%
\TextAt{9}{172.066}{395.107}{\Character{61}}%
\TextAt{9}{184.535}{395.107}{0}%
\TextAt{8}{190.397}{395.107}{.}%
\TextAt{8}{198.937}{395.107}{If}%
\TextAt{14}{210.66}{395.107}{p}%
\TextAt{8}{220.571}{395.107}{is}%
\TextAt{8}{232.366}{395.107}{the}%
\TextAt{8}{252.553}{395.107}{infinite}%
\TextAt{8}{292.251}{395.107}{prime}%
\TextAt{8}{325.445}{395.107}{associated}%
\TextAt{8}{381.211}{395.107}{to}%
\TextAt{8}{395.552}{395.107}{the}%
\TextAt{8}{415.739}{395.107}{embedding}%
\TextAt{13}{474.948}{395.107}{\Character{28}}%
\TextAt{9}{484.74}{395.107}{\Character{58}}%
\TextAt{13}{491.364}{395.107}{K}%
\NativeMathAt{505.555}{395.107}{11.955}{\rightarrow}%
\TextAt{18}{520.883}{395.107}{C}%
\TextAt{8}{529.518}{395.107}{,}%
\TextAt{8}{79.226}{378.061}{then}%
\TextAt{8}{105.886}{378.061}{one}%
\TextAt{8}{127.346}{378.061}{finds}%
\TextAt{16}{217.163}{361.015}{\Character{106}}%
\TextAt{13}{220.484}{361.015}{a}%
\TextAt{16}{226.629}{361.015}{\Character{106}}%
\TextAt{27}{229.949}{359.222}{p}%
\TextAt{9}{237.757}{361.015}{\Character{61}}%
\TextAt{16}{250.183}{361.015}{\Character{106}}%
\TextAt{13}{253.503}{361.015}{\Character{28}}%
\TextAt{13}{259.911}{361.015}{a}%
\TextAt{16}{266.067}{361.015}{\Character{106}}%
\TextAt{13}{269.388}{361.015}{\Character{59}}%
\TextAt{8}{298.543}{361.015}{resp.}%
\TextAt{16}{334.623}{361.015}{\Character{106}}%
\TextAt{13}{337.944}{361.015}{a}%
\TextAt{16}{344.089}{361.015}{\Character{106}}%
\TextAt{27}{347.41}{359.222}{p}%
\TextAt{9}{355.218}{361.015}{\Character{61}}%
\TextAt{16}{367.643}{361.015}{\Character{106}}%
\TextAt{13}{370.964}{361.015}{\Character{28}}%
\TextAt{13}{377.372}{361.015}{a}%
\TextAt{16}{383.528}{361.015}{\Character{106}}%
\TextAt{17}{386.849}{365.951}{2}%
\TextAt{13}{391.581}{361.015}{\Character{59}}%
\TextAt{8}{79.226}{337.493}{if}%
\TextAt{14}{89.95}{337.493}{p}%
\TextAt{8}{99.83}{337.493}{is}%
\TextAt{8}{111.594}{337.493}{real,}%
\TextAt{8}{137.609}{337.493}{resp.}%
\TextAt{8}{166.931}{337.493}{complex.}%
\TextAt{8}{97.159}{317.458}{If}%
\TextAt{13}{109.462}{317.458}{L}%
\TextAt{16}{121.787}{317.458}{\Character{106}}%
\TextAt{13}{129.469}{317.458}{K}%
\TextAt{8}{144.795}{317.458}{is}%
\TextAt{8}{157.169}{317.458}{a}%
\TextAt{8}{167.535}{317.458}{finite}%
\TextAt{8}{198.058}{317.458}{extension}%
\TextAt{8}{250.434}{317.458}{of}%
\TextAt{13}{264.375}{317.458}{K}%
\TextAt{8}{275.192}{317.458}{,}%
\TextAt{8}{283.105}{317.458}{then}%
\TextAt{8}{310.375}{317.458}{we}%
\TextAt{8}{328.213}{317.458}{denote}%
\TextAt{8}{366.542}{317.458}{the}%
\TextAt{8}{387.309}{317.458}{primes}%
\TextAt{8}{425.698}{317.458}{of}%
\TextAt{13}{439.639}{317.458}{L}%
\TextAt{8}{452.111}{317.458}{by}%
\TextAt{14}{468.981}{317.458}{P}%
\TextAt{8}{483.382}{317.458}{and}%
\TextAt{8}{506.755}{317.458}{write}%
\TextAt{14}{79.226}{300.412}{P}%
\TextAt{16}{92.496}{300.412}{\Character{106}}%
\TextAt{14}{99.193}{300.412}{p}%
\TextAt{8}{109.106}{300.412}{to}%
\TextAt{8}{123.449}{300.412}{signify}%
\TextAt{8}{160.614}{300.412}{that}%
\TextAt{8}{186.015}{300.412}{the}%
\TextAt{8}{206.204}{300.412}{valuations}%
\TextAt{8}{261.924}{300.412}{in}%
\TextAt{8}{275.609}{300.412}{the}%
\TextAt{8}{295.798}{300.412}{class}%
\TextAt{14}{323.268}{300.412}{P}%
\TextAt{8}{333.161}{300.412}{,}%
\TextAt{8}{340.351}{300.412}{when}%
\TextAt{8}{370.941}{300.412}{restricted}%
\TextAt{8}{423.062}{300.412}{to}%
\TextAt{13}{437.405}{300.412}{K}%
\TextAt{8}{448.222}{300.412}{,}%
\TextAt{8}{455.412}{300.412}{give}%
\TextAt{8}{479.51}{300.412}{those}%
\TextAt{8}{510.172}{300.412}{of}%
\TextAt{14}{523.535}{300.412}{p}%
\TextAt{8}{529.518}{300.412}{.}%
\TextAt{8}{79.226}{283.366}{In}%
\TextAt{8}{93.535}{283.366}{the}%
\TextAt{8}{113.379}{283.366}{case}%
\TextAt{8}{137.837}{283.366}{of}%
\TextAt{8}{150.855}{283.366}{an}%
\TextAt{8}{166.801}{283.366}{infinite}%
\TextAt{8}{206.156}{283.366}{prime}%
\TextAt{14}{239.007}{283.366}{P}%
\TextAt{8}{248.9}{283.366}{,}%
\TextAt{8}{255.799}{283.366}{we}%
\TextAt{8}{272.714}{283.366}{define}%
\TextAt{8}{306.211}{283.366}{the}%
\TextAt{12}{326.055}{283.366}{inertia}%
\TextAt{12}{368.543}{283.366}{degree}%
\TextAt{8}{406.644}{283.366}{,}%
\TextAt{8}{413.543}{283.366}{resp.}%
\TextAt{8}{442.761}{283.366}{the}%
\TextAt{12}{462.605}{283.366}{ramification}%
\TextAt{12}{79.226}{266.32}{index}%
\TextAt{8}{111.146}{266.32}{,}%
\TextAt{8}{118.296}{266.32}{by}%
\TextAt{13}{214.478}{249.274}{f}%
\TextAt{27}{220.248}{247.415}{P}%
\TextAt{3}{226.844}{247.415}{\Character{106}}%
\TextAt{27}{229.196}{247.415}{p}%
\TextAt{9}{237.004}{249.274}{\Character{61}}%
\TextAt{9}{249.414}{249.274}{\Character{91}}%
\TextAt{13}{252.681}{249.274}{L}%
\TextAt{27}{260.645}{247.481}{P}%
\TextAt{9}{271.06}{249.274}{\Character{58}}%
\TextAt{13}{277.632}{249.274}{K}%
\TextAt{27}{287.591}{247.481}{p}%
\TextAt{9}{292.078}{249.274}{\Character{93}}%
\TextAt{13}{295.33}{249.274}{\Character{59}}%
\TextAt{8}{324.484}{249.274}{resp.}%
\TextAt{13}{353.806}{249.274}{e}%
\TextAt{27}{359.232}{247.415}{P}%
\TextAt{3}{365.827}{247.415}{\Character{106}}%
\TextAt{27}{368.18}{247.415}{p}%
\TextAt{9}{375.987}{249.274}{\Character{61}}%
\TextAt{9}{388.397}{249.274}{1}%
\TextAt{13}{394.266}{249.274}{\Character{58}}%
\TextAt{8}{79.226}{225.753}{For}%
\TextAt{8}{100.196}{225.753}{arbitrary}%
\TextAt{8}{149.965}{225.753}{primes}%
\TextAt{14}{187.743}{225.753}{P}%
\TextAt{16}{200.957}{225.753}{\Character{106}}%
\TextAt{14}{207.599}{225.753}{p}%
\TextAt{8}{217.479}{225.753}{we}%
\TextAt{8}{234.707}{225.753}{then}%
\TextAt{8}{261.367}{225.753}{have}%
\TextAt{8}{288.362}{225.753}{the}%
\TextAt{12}{79.226}{199.74}{Proposition}%
\TextAt{12}{151.937}{199.74}{(1.2).}%
\TextAt{8}{89.579}{179.706}{(i)}%
\TextAt{20}{107.918}{191.063}{\Character{88}}%
\TextAt{27}{110.085}{165.16}{P}%
\TextAt{3}{116.68}{165.16}{\Character{106}}%
\TextAt{27}{119.032}{165.16}{p}%
\TextAt{13}{127.18}{179.706}{e}%
\TextAt{27}{132.605}{177.846}{P}%
\TextAt{3}{139.201}{177.846}{\Character{106}}%
\TextAt{27}{141.553}{177.846}{p}%
\TextAt{13}{146.04}{179.706}{f}%
\TextAt{27}{151.81}{177.846}{P}%
\TextAt{3}{158.405}{177.846}{\Character{106}}%
\TextAt{27}{160.757}{177.846}{p}%
\TextAt{9}{168.565}{179.706}{\Character{61}}%
\TextAt{20}{180.991}{191.063}{\Character{88}}%
\TextAt{27}{183.157}{165.16}{P}%
\TextAt{3}{189.752}{165.16}{\Character{106}}%
\TextAt{27}{192.104}{165.16}{p}%
\TextAt{9}{198.259}{179.706}{\Character{91}}%
\TextAt{13}{201.511}{179.706}{L}%
\TextAt{27}{209.475}{177.912}{P}%
\TextAt{9}{219.89}{179.706}{\Character{58}}%
\TextAt{13}{226.462}{179.706}{K}%
\TextAt{27}{236.421}{177.912}{p}%
\TextAt{9}{240.908}{179.706}{\Character{93}}%
\TextAt{9}{247.459}{179.706}{\Character{61}}%
\TextAt{9}{259.881}{179.706}{\Character{91}}%
\TextAt{13}{263.158}{179.706}{L}%
\TextAt{9}{274.443}{179.706}{\Character{58}}%
\TextAt{13}{281.016}{179.706}{K}%
\TextAt{9}{291.833}{179.706}{\Character{93}}%
\TextAt{8}{295.085}{179.706}{,}%
\TextAt{8}{86.327}{150.922}{(ii)}%
\TextAt{14}{107.918}{150.922}{N}%
\TextAt{9}{117.867}{150.922}{\Character{40}}%
\TextAt{14}{122.419}{150.922}{P}%
\TextAt{9}{132.312}{150.922}{\Character{41}}%
\TextAt{9}{140.167}{150.922}{\Character{61}}%
\TextAt{14}{152.611}{150.922}{N}%
\TextAt{9}{162.559}{150.922}{\Character{40}}%
\TextAt{14}{167.111}{150.922}{p}%
\TextAt{9}{173.094}{150.922}{\Character{41}}%
\TextAt{21}{177.647}{155.441}{f}%
\TextAt{28}{181.715}{153.703}{P}%
\TextAt{23}{187.414}{153.703}{\Character{106}}%
\TextAt{28}{189.683}{153.703}{p}%
\TextAt{8}{194.264}{150.922}{,}%
\TextAt{8}{83.076}{133.876}{(iii)}%
\TextAt{13}{107.918}{133.876}{v}%
\TextAt{27}{113.578}{132.083}{P}%
\TextAt{9}{120.671}{133.876}{\Character{40}}%
\TextAt{13}{125.224}{133.876}{a}%
\TextAt{9}{131.369}{133.876}{\Character{41}}%
\TextAt{9}{139.224}{133.876}{\Character{61}}%
\TextAt{13}{151.667}{133.876}{e}%
\TextAt{27}{157.093}{132.016}{P}%
\TextAt{3}{163.688}{132.016}{\Character{106}}%
\TextAt{27}{166.041}{132.016}{p}%
\TextAt{13}{170.527}{133.876}{v}%
\TextAt{27}{176.187}{132.083}{p}%
\TextAt{9}{180.673}{133.876}{\Character{40}}%
\TextAt{13}{185.226}{133.876}{a}%
\TextAt{9}{191.371}{133.876}{\Character{41}}%
\TextAt{8}{199.82}{133.876}{for}%
\TextAt{13}{217.705}{133.876}{a}%
\NativeMathAt{227.171}{133.876}{11.955}{\in}%
\TextAt{13}{238.462}{133.876}{K}%
\TextAt{3}{249.28}{138.214}{\Character{3}}%
\TextAt{8}{254.012}{133.876}{,}%
\TextAt{8}{83.398}{116.83}{(iv)}%
\TextAt{13}{107.918}{116.83}{v}%
\TextAt{27}{113.578}{115.037}{p}%
\TextAt{9}{118.065}{116.83}{\Character{40}}%
\TextAt{13}{122.617}{116.83}{N}%
\TextAt{21}{131.991}{114.97}{L}%
\TextAt{28}{137.75}{113.579}{P}%
\TextAt{3}{143.947}{114.97}{\Character{106}}%
\TextAt{21}{146.299}{114.97}{K}%
\TextAt{28}{153.441}{113.974}{p}%
\TextAt{9}{158.021}{116.83}{\Character{40}}%
\TextAt{13}{162.574}{116.83}{a}%
\TextAt{9}{168.719}{116.83}{\Character{41}\Character{41}}%
\TextAt{9}{181.117}{116.83}{\Character{61}}%
\TextAt{13}{193.57}{116.83}{f}%
\TextAt{27}{199.34}{114.97}{P}%
\TextAt{3}{205.935}{114.97}{\Character{106}}%
\TextAt{27}{208.288}{114.97}{p}%
\TextAt{13}{212.774}{116.83}{v}%
\TextAt{27}{218.434}{115.037}{P}%
\TextAt{9}{225.527}{116.83}{\Character{40}}%
\TextAt{13}{230.08}{116.83}{a}%
\TextAt{9}{236.225}{116.83}{\Character{41}}%
\TextAt{8}{244.674}{116.83}{for}%
\TextAt{13}{262.559}{116.83}{a}%
\NativeMathAt{272.025}{116.83}{11.955}{\in}%
\TextAt{13}{283.316}{116.83}{L}%
\TextAt{3}{291.28}{121.168}{\Character{3}}%
\TextAt{8}{296.013}{116.83}{,}%
\TextAt{8}{86.65}{99.784}{(v)}%
\TextAt{16}{107.918}{99.784}{\Character{106}}%
\TextAt{13}{111.239}{99.784}{a}%
\TextAt{16}{117.384}{99.784}{\Character{106}}%
\TextAt{27}{120.705}{97.991}{P}%
\TextAt{9}{131.12}{99.784}{\Character{61}}%
\TextAt{16}{143.545}{99.784}{\Character{106}}%
\TextAt{13}{146.866}{99.784}{N}%
\TextAt{21}{156.24}{97.924}{L}%
\TextAt{28}{161.999}{96.533}{P}%
\TextAt{3}{168.196}{97.924}{\Character{106}}%
\TextAt{21}{170.548}{97.924}{K}%
\TextAt{28}{177.69}{96.928}{p}%
\TextAt{9}{182.271}{99.784}{\Character{40}}%
\TextAt{13}{186.823}{99.784}{a}%
\TextAt{9}{192.968}{99.784}{\Character{41}}%
\TextAt{16}{197.52}{99.784}{\Character{106}}%
\TextAt{27}{200.841}{97.991}{p}%
\TextAt{8}{209.225}{99.784}{for}%
\TextAt{13}{227.11}{99.784}{a}%
\NativeMathAt{236.576}{99.784}{11.955}{\in}%
\TextAt{13}{247.867}{99.784}{L}%
\TextAt{8}{255.831}{99.784}{.}%
\TextAt{8}{297.211}{43.033}{168}%
\Destination{79.23}{193.755}{nav.Item.59}%
\Destination{79.23}{161.975}{nav.Item.60}%
\Destination{79.23}{147.925}{nav.Item.61}%
\Destination{79.23}{130.025}{nav.Item.62}%
\Destination{79.23}{112.445}{nav.Item.63}%
\Destination{78.23}{809.815}{nav.page.168}%
\Destination{79.36}{776.155}{nav.print.184}%
\Destination{79.36}{213.755}{nav.number.3.1.2}%
}

% Page 181 (body page 169)
\Page{%
\TextAt{2}{79.226}{796.943}{III.}%
\TextAt{2}{98.317}{796.943}{Riemann–Roch}%
\TextAt{2}{174.375}{796.943}{Theory}%
\Vector{q 0 G 0 g 0.299 w 7.226 18.729 m 460.769 18.729 l S Q}%
\TextAt{8}{97.159}{761.124}{The}%
\TextAt{8}{120.936}{761.124}{normalized}%
\TextAt{8}{180.494}{761.124}{valuations}%
\TextAt{16}{235.904}{761.124}{\Character{106}}%
\TextAt{16}{242.826}{761.124}{\Character{106}}%
\TextAt{27}{246.166}{759.33}{p}%
\TextAt{8}{254.273}{761.124}{satisfy}%
\TextAt{8}{290.543}{761.124}{the}%
\TextAt{8}{310.422}{761.124}{following}%
\TextAt{12}{359.58}{761.124}{product}%
\TextAt{12}{409.975}{761.124}{formula}%
\TextAt{8}{454.914}{761.124}{,}%
\TextAt{8}{461.842}{761.124}{which}%
\TextAt{8}{495.051}{761.124}{demon-}%
\TextAt{8}{79.226}{744.078}{strates}%
\TextAt{8}{117.076}{744.078}{how}%
\TextAt{8}{141.465}{744.078}{important}%
\TextAt{8}{196.757}{744.078}{it}%
\TextAt{8}{208.462}{744.078}{is}%
\TextAt{8}{220.225}{744.078}{to}%
\TextAt{8}{234.536}{744.078}{include}%
\TextAt{8}{274.849}{744.078}{the}%
\TextAt{8}{295.005}{744.078}{infinite}%
\TextAt{8}{334.672}{744.078}{primes.}%
\TextAt{12}{79.226}{718.065}{Proposition}%
\TextAt{12}{151.937}{718.065}{(1.3).}%
\TextAt{8}{187.253}{718.065}{Given}%
\TextAt{8}{221.146}{718.065}{any}%
\TextAt{13}{243.263}{718.065}{a}%
\NativeMathAt{252.729}{718.065}{11.955}{\in}%
\TextAt{13}{264.02}{718.065}{K}%
\TextAt{3}{274.837}{722.404}{\Character{3}}%
\TextAt{8}{279.57}{718.065}{,}%
\TextAt{8}{286.719}{718.065}{one}%
\TextAt{8}{308.178}{718.065}{has}%
\TextAt{16}{329.052}{718.065}{\Character{106}}%
\TextAt{13}{332.373}{718.065}{a}%
\TextAt{16}{338.518}{718.065}{\Character{106}}%
\TextAt{27}{341.839}{716.272}{p}%
\TextAt{9}{349.647}{718.065}{\Character{61}}%
\TextAt{9}{362.057}{718.065}{1}%
\TextAt{8}{371.822}{718.065}{for}%
\TextAt{8}{389.707}{718.065}{almost}%
\TextAt{8}{427.498}{718.065}{all}%
\TextAt{14}{443.757}{718.065}{p}%
\TextAt{8}{449.74}{718.065}{,}%
\TextAt{8}{456.889}{718.065}{and}%
\TextAt{20}{276.301}{700.422}{\Character{89}}%
\TextAt{27}{281.945}{675.316}{p}%
\TextAt{16}{293.57}{689.064}{\Character{106}}%
\TextAt{13}{296.891}{689.064}{a}%
\TextAt{16}{303.036}{689.064}{\Character{106}}%
\TextAt{27}{306.357}{687.271}{p}%
\TextAt{9}{314.164}{689.064}{\Character{61}}%
\TextAt{9}{326.574}{689.064}{1}%
\TextAt{13}{332.443}{689.064}{\Character{58}}%
\TextAt{19}{79.226}{632.763}{Proof.}%
\TextAt{8}{116}{632.763}{We}%
\TextAt{8}{136.399}{632.763}{have}%
\TextAt{13}{163.648}{632.763}{v}%
\TextAt{27}{169.307}{630.97}{p}%
\TextAt{9}{173.794}{632.763}{\Character{40}}%
\TextAt{13}{178.346}{632.763}{a}%
\TextAt{9}{184.491}{632.763}{\Character{41}}%
\TextAt{9}{192.776}{632.763}{\Character{61}}%
\TextAt{9}{205.628}{632.763}{0}%
\TextAt{8}{215.66}{632.763}{and}%
\TextAt{8}{238.676}{632.763}{therefore}%
\TextAt{16}{288.03}{632.763}{\Character{106}}%
\TextAt{13}{291.351}{632.763}{a}%
\TextAt{16}{297.496}{632.763}{\Character{106}}%
\TextAt{27}{300.817}{630.97}{p}%
\TextAt{9}{309.058}{632.763}{\Character{61}}%
\TextAt{9}{321.898}{632.763}{1}%
\TextAt{8}{331.92}{632.763}{for}%
\TextAt{8}{350.059}{632.763}{all}%
\TextAt{14}{366.572}{632.763}{p}%
\TextAt{18}{376.309}{632.763}{\Character{45}}%
\NativeMathAt{383.384}{632.763}{11.955}{\infty}%
\TextAt{8}{399.49}{632.763}{which}%
\TextAt{8}{433.231}{632.763}{do}%
\TextAt{8}{449.744}{632.763}{not}%
\TextAt{8}{470.812}{632.763}{occur}%
\TextAt{8}{502.604}{632.763}{in}%
\TextAt{8}{516.51}{632.763}{the}%
\TextAt{8}{79.226}{615.717}{prime}%
\TextAt{8}{113.312}{615.717}{decomposition}%
\TextAt{8}{191.368}{615.717}{of}%
\TextAt{8}{205.62}{615.717}{the}%
\TextAt{8}{226.698}{615.717}{principal}%
\TextAt{8}{276.397}{615.717}{ideal}%
\TextAt{9}{305.282}{615.717}{\Character{40}}%
\TextAt{13}{309.835}{615.717}{a}%
\TextAt{9}{315.979}{615.717}{\Character{41}}%
\TextAt{8}{325.351}{615.717}{(see}%
\TextAt{8}{349.741}{615.717}{chap.}%
\TextAt{8}{384.698}{615.717}{I,}%
\TextAt{8}{396.989}{615.717}{§11,}%
\TextAt{8}{422.685}{615.717}{p.69).}%
\TextAt{8}{459.925}{615.717}{This}%
\TextAt{8}{487.567}{615.717}{therefore}%
\TextAt{8}{79.226}{598.671}{holds}%
\TextAt{8}{109.855}{598.671}{for}%
\TextAt{8}{127.74}{598.671}{almost}%
\TextAt{8}{165.531}{598.671}{all}%
\TextAt{14}{181.79}{598.671}{p}%
\TextAt{8}{187.773}{598.671}{.}%
\TextAt{8}{196.221}{598.671}{From}%
\TextAt{8}{226.946}{598.671}{(1.2)}%
\TextAt{8}{254.921}{598.671}{and}%
\TextAt{8}{277.683}{598.671}{formula}%
\TextAt{8}{320.615}{598.671}{(8.4)}%
\TextAt{8}{348.59}{598.671}{of}%
\TextAt{8}{361.92}{598.671}{chap.}%
\TextAt{8}{394.111}{598.671}{II,}%
\TextAt{13}{244.184}{569.67}{N}%
\TextAt{21}{253.558}{567.81}{K}%
\TextAt{3}{261.262}{567.81}{\Character{106}}%
\TextAt{25}{263.615}{567.81}{Q}%
\TextAt{9}{270.312}{569.67}{\Character{40}}%
\TextAt{13}{274.864}{569.67}{a}%
\TextAt{9}{281.009}{569.67}{\Character{41}}%
\TextAt{9}{288.864}{569.67}{\Character{61}}%
\TextAt{20}{301.308}{581.027}{\Character{89}}%
\TextAt{27}{303.644}{555.124}{p}%
\TextAt{3}{307.632}{555.124}{\Character{106}}%
\TextAt{21}{309.985}{555.124}{p}%
\TextAt{13}{318.576}{569.67}{N}%
\TextAt{21}{327.95}{567.81}{K}%
\TextAt{28}{335.091}{566.814}{p}%
\TextAt{3}{339.174}{567.81}{\Character{106}}%
\TextAt{25}{341.526}{567.81}{Q}%
\TextAt{22}{347.725}{566.814}{p}%
\TextAt{9}{352.562}{569.67}{\Character{40}}%
\TextAt{13}{357.114}{569.67}{a}%
\TextAt{9}{363.259}{569.67}{\Character{41}}%
\TextAt{8}{79.226}{530.018}{(which}%
\TextAt{8}{117.094}{530.018}{includes}%
\TextAt{8}{161.847}{530.018}{the}%
\TextAt{8}{181.83}{530.018}{case}%
\TextAt{13}{206.427}{530.018}{p}%
\TextAt{9}{215.623}{530.018}{\Character{61}}%
\NativeMathAt{228.049}{530.018}{11.955}{\infty}%
\TextAt{8}{240.004}{530.018}{,}%
\TextAt{18}{247.014}{530.018}{Q}%
\TextAt{21}{256.313}{528.225}{p}%
\TextAt{9}{264.394}{530.018}{\Character{61}}%
\TextAt{18}{276.82}{530.018}{R}%
\TextAt{8}{285.454}{530.018}{),}%
\TextAt{8}{297.019}{530.018}{we}%
\TextAt{8}{314.073}{530.018}{obtain}%
\TextAt{8}{350.326}{530.018}{the}%
\TextAt{8}{370.309}{530.018}{product}%
\TextAt{8}{414.034}{530.018}{formula}%
\TextAt{8}{456.792}{530.018}{for}%
\TextAt{13}{474.503}{530.018}{K}%
\TextAt{8}{489.044}{530.018}{from}%
\TextAt{8}{516.51}{530.018}{the}%
\TextAt{8}{79.226}{512.972}{product}%
\TextAt{8}{123.126}{512.972}{formula}%
\TextAt{8}{166.057}{512.972}{for}%
\TextAt{18}{183.941}{512.972}{Q}%
\TextAt{8}{193.24}{512.972}{,}%
\TextAt{8}{200.389}{512.972}{which}%
\TextAt{8}{233.876}{512.972}{was}%
\TextAt{8}{256.375}{512.972}{proved}%
\TextAt{8}{294.429}{512.972}{already}%
\TextAt{8}{335.734}{512.972}{in}%
\TextAt{8}{349.386}{512.972}{chap.}%
\TextAt{8}{381.578}{512.972}{II,}%
\TextAt{8}{397.49}{512.972}{(2.1):}%
\TextAt{20}{144.164}{495.329}{\Character{89}}%
\TextAt{27}{149.808}{470.223}{p}%
\TextAt{16}{161.433}{483.971}{\Character{106}}%
\TextAt{13}{164.754}{483.971}{a}%
\TextAt{16}{170.899}{483.971}{\Character{106}}%
\TextAt{27}{174.22}{482.178}{p}%
\TextAt{9}{182.027}{483.971}{\Character{61}}%
\TextAt{20}{194.453}{495.329}{\Character{89}}%
\TextAt{21}{199.96}{470.223}{p}%
\TextAt{20}{211.721}{495.329}{\Character{89}}%
\TextAt{27}{214.058}{469.426}{p}%
\TextAt{3}{218.046}{469.426}{\Character{106}}%
\TextAt{21}{220.399}{469.426}{p}%
\TextAt{16}{228.99}{483.971}{\Character{106}}%
\TextAt{13}{232.311}{483.971}{a}%
\TextAt{16}{238.456}{483.971}{\Character{106}}%
\TextAt{27}{241.777}{482.178}{p}%
\TextAt{9}{249.584}{483.971}{\Character{61}}%
\TextAt{20}{262.01}{495.329}{\Character{89}}%
\TextAt{21}{267.517}{470.223}{p}%
\TextAt{20}{279.279}{495.329}{\Character{89}}%
\TextAt{27}{281.615}{469.426}{p}%
\TextAt{3}{285.603}{469.426}{\Character{106}}%
\TextAt{21}{287.956}{469.426}{p}%
\TextAt{16}{296.547}{483.971}{\Character{106}}%
\TextAt{13}{299.868}{483.971}{N}%
\TextAt{21}{309.242}{482.112}{K}%
\TextAt{28}{316.383}{481.115}{p}%
\TextAt{3}{320.466}{482.112}{\Character{106}}%
\TextAt{25}{322.818}{482.112}{Q}%
\TextAt{22}{329.017}{481.115}{p}%
\TextAt{9}{333.854}{483.971}{\Character{40}}%
\TextAt{13}{338.406}{483.971}{a}%
\TextAt{9}{344.551}{483.971}{\Character{41}}%
\TextAt{16}{349.104}{483.971}{\Character{106}}%
\TextAt{21}{352.424}{482.178}{p}%
\TextAt{9}{360.506}{483.971}{\Character{61}}%
\TextAt{20}{372.932}{495.329}{\Character{89}}%
\TextAt{21}{378.438}{470.223}{p}%
\TextAt{16}{390.2}{483.971}{\Character{106}}%
\TextAt{13}{393.521}{483.971}{N}%
\TextAt{21}{402.895}{482.112}{K}%
\TextAt{3}{410.6}{482.112}{\Character{106}}%
\TextAt{25}{412.952}{482.112}{Q}%
\TextAt{9}{419.649}{483.971}{\Character{40}}%
\TextAt{13}{424.201}{483.971}{a}%
\TextAt{9}{430.346}{483.971}{\Character{41}}%
\TextAt{16}{434.899}{483.971}{\Character{106}}%
\TextAt{21}{438.219}{482.178}{p}%
\TextAt{9}{446.301}{483.971}{\Character{61}}%
\TextAt{9}{458.711}{483.971}{1}%
\TextAt{13}{464.58}{483.971}{\Character{58}}%
\Vector{q 0 G 0 g 0.399 w 452.334 -328.784 m 452.334 -320.714 l S Q}%
\Vector{q 0 G 0 g 0.399 w 452.534 -320.913 m 459.707 -320.913 l S Q}%
\Vector{q 0 G 0 g 0.399 w 452.534 -328.584 m 459.707 -328.584 l S Q}%
\Vector{q 0 G 0 g 0.399 w 459.906 -328.784 m 459.906 -320.714 l S Q}%
\TextAt{8}{97.159}{416.215}{We}%
\TextAt{8}{119.107}{416.215}{denote}%
\TextAt{8}{158.629}{416.215}{by}%
\TextAt{13}{176.691}{416.215}{J}%
\TextAt{9}{184.268}{416.215}{\Character{40}}%
\TextAt{14}{188.821}{416.215}{o}%
\TextAt{9}{194.665}{416.215}{\Character{41}}%
\TextAt{8}{204.918}{416.215}{the}%
\TextAt{8}{226.878}{416.215}{group}%
\TextAt{8}{261.857}{416.215}{of}%
\TextAt{8}{276.99}{416.215}{fractional}%
\TextAt{8}{331.157}{416.215}{ideals}%
\TextAt{8}{365.538}{416.215}{of}%
\TextAt{13}{380.672}{416.215}{K}%
\TextAt{8}{391.489}{416.215}{,}%
\TextAt{8}{400.893}{416.215}{by}%
\TextAt{13}{418.955}{416.215}{P}%
\TextAt{9}{428.126}{416.215}{\Character{40}}%
\TextAt{14}{432.679}{416.215}{o}%
\TextAt{9}{438.523}{416.215}{\Character{41}}%
\TextAt{8}{448.776}{416.215}{the}%
\TextAt{8}{470.736}{416.215}{subgroup}%
\TextAt{8}{523.337}{416.215}{of}%
\TextAt{8}{79.226}{399.169}{fractional}%
\TextAt{8}{131.59}{399.169}{principal}%
\TextAt{8}{180.367}{399.169}{ideals,}%
\TextAt{8}{216.197}{399.169}{and}%
\TextAt{8}{238.959}{399.169}{by}%
\TextAt{8}{256.193}{370.168}{Pic}%
\TextAt{9}{272.608}{370.168}{\Character{40}}%
\TextAt{14}{277.16}{370.168}{o}%
\TextAt{9}{283.005}{370.168}{\Character{41}}%
\TextAt{9}{290.86}{370.168}{\Character{61}}%
\TextAt{13}{303.303}{370.168}{J}%
\TextAt{9}{310.88}{370.168}{\Character{40}}%
\TextAt{14}{315.433}{370.168}{o}%
\TextAt{9}{321.277}{370.168}{\Character{41}\Character{47}}%
\TextAt{13}{331.682}{370.168}{P}%
\TextAt{9}{340.853}{370.168}{\Character{40}}%
\TextAt{14}{345.406}{370.168}{o}%
\TextAt{9}{351.25}{370.168}{\Character{41}}%
\TextAt{8}{79.226}{341.167}{the}%
\TextAt{8}{99.383}{341.167}{ideal}%
\TextAt{8}{127.346}{341.167}{class}%
\TextAt{8}{154.783}{341.167}{group}%
\TextAt{8}{187.958}{341.167}{Cl}%
\TextAt{21}{199.663}{339.374}{K}%
\TextAt{8}{211.763}{341.167}{of}%
\TextAt{13}{225.093}{341.167}{K}%
\TextAt{8}{235.91}{341.167}{.}%
\TextAt{8}{97.159}{321.132}{Let}%
\TextAt{8}{118.116}{321.132}{us}%
\TextAt{8}{133.12}{321.132}{now}%
\TextAt{8}{157.497}{321.132}{extend}%
\TextAt{8}{195.526}{321.132}{the}%
\TextAt{8}{215.671}{321.132}{notion}%
\TextAt{8}{252.087}{321.132}{of}%
\TextAt{8}{265.405}{321.132}{fractional}%
\TextAt{8}{317.757}{321.132}{ideal}%
\TextAt{8}{345.708}{321.132}{of}%
\TextAt{13}{359.026}{321.132}{K}%
\TextAt{8}{373.729}{321.132}{by}%
\TextAt{8}{389.976}{321.132}{taking}%
\TextAt{8}{426.069}{321.132}{into}%
\TextAt{8}{449.8}{321.132}{account}%
\TextAt{8}{493.042}{321.132}{also}%
\TextAt{8}{516.51}{321.132}{the}%
\TextAt{8}{79.226}{304.086}{infinite}%
\TextAt{8}{118.893}{304.086}{primes.}%
\TextAt{12}{79.226}{278.074}{Definition}%
\TextAt{12}{142.015}{278.074}{(1.4).}%
\TextAt{8}{177.33}{278.074}{A}%
\TextAt{12}{190.003}{278.074}{replete}%
\TextAt{12}{234.835}{278.074}{ideal}%
\TextAt{8}{266.36}{278.074}{of}%
\TextAt{13}{279.69}{278.074}{K}%
\TextAt{8}{294.405}{278.074}{is}%
\TextAt{8}{306.169}{278.074}{an}%
\TextAt{8}{322.428}{278.074}{element}%
\TextAt{8}{365.67}{278.074}{of}%
\TextAt{8}{379}{278.074}{the}%
\TextAt{8}{399.157}{278.074}{group}%
\TextAt{13}{248.547}{249.073}{J}%
\TextAt{9}{256.124}{249.073}{\Character{40}}%
\Vector{q 0 G 0 g 0.478 w 188.677 -517.571 m 194.521 -517.571 l S Q}%
\TextAt{14}{260.677}{249.073}{o}%
\TextAt{9}{266.521}{249.073}{\Character{41}}%
\TextAt{9}{274.376}{249.073}{\Character{58}\Character{61}}%
\TextAt{13}{290.071}{249.073}{J}%
\TextAt{9}{297.648}{249.073}{\Character{40}}%
\TextAt{14}{302.201}{249.073}{o}%
\TextAt{9}{308.045}{249.073}{\Character{41}}%
\TextAt{16}{315.254}{249.073}{\Character{2}}%
\TextAt{20}{327.209}{260.43}{\Character{89}}%
\TextAt{27}{327.442}{234.527}{p}%
\TextAt{3}{331.431}{234.527}{\Character{106}}%
\NativeMathAt{333.782}{234.527}{7.97}{\infty}%
\TextAt{18}{344.478}{249.073}{R}%
\TextAt{3}{353.112}{254.009}{\Character{3}}%
\TextAt{17}{353.112}{246.117}{\Character{43}}%
\TextAt{13}{360.197}{249.073}{\Character{59}}%
\TextAt{8}{79.226}{209.971}{where}%
\TextAt{18}{113.035}{209.971}{R}%
\TextAt{3}{121.67}{214.31}{\Character{3}}%
\TextAt{17}{121.67}{207.016}{\Character{43}}%
\TextAt{8}{132.652}{209.971}{denotes}%
\TextAt{8}{174.985}{209.971}{the}%
\TextAt{8}{195.141}{209.971}{multiplicative}%
\TextAt{8}{268.965}{209.971}{group}%
\TextAt{8}{302.14}{209.971}{of}%
\TextAt{8}{315.47}{209.971}{positive}%
\TextAt{8}{358.784}{209.971}{real}%
\TextAt{8}{381.546}{209.971}{numbers.}%
\TextAt{8}{97.159}{183.959}{In}%
\TextAt{8}{111.022}{183.959}{order}%
\TextAt{8}{140.834}{183.959}{to}%
\TextAt{8}{154.387}{183.959}{unify}%
\TextAt{8}{183.541}{183.959}{notation,}%
\TextAt{8}{233.027}{183.959}{we}%
\TextAt{8}{249.496}{183.959}{put,}%
\TextAt{8}{273.601}{183.959}{for}%
\TextAt{8}{290.729}{183.959}{any}%
\TextAt{8}{312.088}{183.959}{infinite}%
\TextAt{8}{350.997}{183.959}{prime}%
\TextAt{14}{383.403}{183.959}{p}%
\TextAt{8}{392.526}{183.959}{and}%
\TextAt{8}{414.531}{183.959}{any}%
\TextAt{8}{435.89}{183.959}{real}%
\TextAt{8}{457.895}{183.959}{number}%
\TextAt{13}{499.734}{183.959}{\Character{23}}%
\NativeMathAt{509.592}{183.959}{11.955}{\in}%
\TextAt{18}{520.883}{183.959}{R}%
\TextAt{8}{529.518}{183.959}{,}%
\TextAt{14}{268.821}{154.958}{p}%
\TextAt{21}{274.804}{159.894}{\Character{23}}%
\TextAt{9}{283.307}{154.958}{\Character{58}\Character{61}}%
\TextAt{13}{298.984}{154.958}{e}%
\TextAt{21}{304.41}{159.894}{\Character{23}}%
\NativeMathAt{312.913}{154.958}{11.955}{\in}%
\TextAt{18}{324.204}{154.958}{R}%
\TextAt{3}{332.838}{159.894}{\Character{3}}%
\TextAt{17}{332.838}{152.002}{\Character{43}}%
\TextAt{13}{339.923}{154.958}{\Character{58}}%
\TextAt{8}{79.226}{125.957}{Given}%
\TextAt{8}{113.119}{125.957}{a}%
\TextAt{8}{122.874}{125.957}{system}%
\TextAt{8}{161.693}{125.957}{of}%
\TextAt{8}{175.023}{125.957}{real}%
\TextAt{8}{197.786}{125.957}{numbers}%
\TextAt{13}{244.997}{125.957}{\Character{23}}%
\TextAt{27}{250.78}{124.164}{p}%
\TextAt{8}{255.267}{125.957}{,}%
\TextAt{14}{262.416}{125.957}{p}%
\TextAt{16}{271.72}{125.957}{\Character{106}}%
\NativeMathAt{278.343}{125.957}{11.955}{\infty}%
\TextAt{8}{290.317}{125.957}{,}%
\TextAt{8}{297.466}{125.957}{let}%
\TextAt{20}{314.371}{134.923}{\Character{81}}%
\TextAt{27}{325.662}{122.47}{p}%
\TextAt{3}{329.651}{122.47}{\Character{106}}%
\NativeMathAt{332.002}{122.47}{7.97}{\infty}%
\TextAt{14}{342.962}{125.957}{p}%
\TextAt{21}{348.945}{130.295}{\Character{23}}%
\TextAt{28}{353.108}{129.299}{p}%
\TextAt{8}{361.586}{125.957}{always}%
\TextAt{8}{399.054}{125.957}{denote}%
\TextAt{8}{436.772}{125.957}{the}%
\TextAt{8}{456.929}{125.957}{vector}%
\TextAt{20}{225.084}{106.33}{\Character{89}}%
\TextAt{27}{225.318}{80.427}{p}%
\TextAt{3}{229.306}{80.427}{\Character{106}}%
\NativeMathAt{231.657}{80.427}{7.97}{\infty}%
\TextAt{14}{242.353}{94.973}{p}%
\TextAt{21}{248.336}{99.909}{\Character{23}}%
\TextAt{28}{252.499}{98.913}{p}%
\TextAt{9}{260.4}{94.973}{\Character{61}}%
\TextAt{9}{272.81}{94.973}{\Character{40}}%
\TextAt{13}{277.378}{94.973}{\Character{58}}%
\TextAt{13}{282.602}{94.973}{\Character{58}}%
\TextAt{13}{287.839}{94.973}{\Character{58}}%
\TextAt{13}{293.075}{94.973}{\Character{59}}%
\TextAt{13}{298.3}{94.973}{e}%
\TextAt{21}{303.78}{99.909}{\Character{23}}%
\TextAt{28}{307.943}{98.913}{p}%
\TextAt{13}{312.524}{94.973}{\Character{59}}%
\TextAt{13}{317.748}{94.973}{\Character{58}}%
\TextAt{13}{322.985}{94.973}{\Character{58}}%
\TextAt{13}{328.209}{94.973}{\Character{58}}%
\TextAt{9}{331.508}{94.973}{\Character{41}}%
\NativeMathAt{339.381}{94.973}{11.955}{\in}%
\TextAt{20}{350.672}{106.33}{\Character{89}}%
\TextAt{27}{350.905}{80.427}{p}%
\TextAt{3}{354.894}{80.427}{\Character{106}}%
\NativeMathAt{357.245}{80.427}{7.97}{\infty}%
\TextAt{18}{367.941}{94.973}{R}%
\TextAt{3}{376.575}{99.909}{\Character{3}}%
\TextAt{17}{376.575}{92.017}{\Character{43}}%
\TextAt{13}{383.66}{94.973}{\Character{59}}%
\TextAt{8}{297.211}{43.033}{169}%
\Destination{78.23}{809.815}{nav.page.169}%
\Destination{79.36}{776.155}{nav.print.185}%
\Destination{79.36}{732.055}{nav.number.3.1.3}%
\Destination{79.36}{292.055}{nav.number.3.1.4}%
\Hotlink{384.198}{613.555}{5.22}{11.047}{nav.chapter.1}%
\Hotlink{402.168}{613.555}{12.716}{11.047}{nav.section.1.11}%
\Hotlink{431.94}{613.555}{12.716}{11.047}{nav.print.69}%
\Hotlink{393.611}{596.455}{9.751}{11.047}{nav.chapter.2}%
\Hotlink{381.078}{510.755}{9.751}{11.047}{nav.chapter.2}%
\Hotlink{226.446}{596.455}{25.078}{11.047}{nav.number.3.1.2}%
\Hotlink{320.115}{596.455}{25.078}{11.047}{nav.number.2.8.4}%
\Hotlink{396.99}{510.755}{25.078}{11.047}{nav.number.2.2.1}%
}

% Page 182 (body page 170)
\Page{%
\TextAt{2}{79.226}{796.943}{III.}%
\TextAt{2}{98.317}{796.943}{Riemann–Roch}%
\TextAt{2}{174.375}{796.943}{Theory}%
\Vector{q 0 G 0 g 0.299 w 7.226 18.729 m 460.769 18.729 l S Q}%
\TextAt{8}{79.226}{761.124}{and}%
\TextAt{19}{102}{761.124}{not}%
\TextAt{8}{122.837}{761.124}{the}%
\TextAt{8}{143.005}{761.124}{product}%
\TextAt{8}{186.915}{761.124}{of}%
\TextAt{8}{200.256}{761.124}{the}%
\TextAt{8}{220.424}{761.124}{quantities}%
\TextAt{13}{274.484}{761.124}{e}%
\TextAt{21}{279.91}{765.462}{\Character{23}}%
\TextAt{28}{284.073}{764.466}{p}%
\TextAt{8}{292.562}{761.124}{in}%
\TextAt{18}{306.226}{761.124}{R}%
\TextAt{8}{314.86}{761.124}{.}%
\TextAt{8}{323.342}{761.124}{Then}%
\TextAt{8}{353.911}{761.124}{every}%
\TextAt{8}{384.814}{761.124}{replete}%
\TextAt{8}{423.189}{761.124}{ideal}%
\TextAt{14}{451.164}{761.124}{a}%
\NativeMathAt{460.486}{761.124}{11.955}{\in}%
\TextAt{13}{471.796}{761.124}{J}%
\TextAt{9}{479.374}{761.124}{\Character{40}}%
\Vector{q 0 G 0 g 0.478 w 411.926 -5.52 m 417.77 -5.52 l S Q}%
\TextAt{14}{483.926}{761.124}{o}%
\TextAt{9}{489.77}{761.124}{\Character{41}}%
\TextAt{8}{498.231}{761.124}{admits}%
\TextAt{8}{79.226}{744.078}{the}%
\TextAt{8}{99.383}{744.078}{unique}%
\TextAt{8}{137.424}{744.078}{product}%
\TextAt{8}{181.323}{744.078}{representation}%
\TextAt{14}{228.71}{716.237}{a}%
\TextAt{9}{238.014}{716.237}{\Character{61}}%
\TextAt{20}{250.439}{727.595}{\Character{89}}%
\TextAt{27}{250.742}{701.763}{p}%
\TextAt{25}{254.73}{701.763}{\Character{45}}%
\NativeMathAt{256.944}{701.763}{7.97}{\infty}%
\TextAt{14}{267.708}{716.237}{p}%
\TextAt{21}{273.691}{721.174}{\Character{23}}%
\TextAt{28}{277.853}{720.177}{p}%
\TextAt{16}{285.091}{716.237}{\Character{2}}%
\TextAt{20}{297.046}{727.595}{\Character{89}}%
\TextAt{27}{297.279}{701.692}{p}%
\TextAt{3}{301.268}{701.692}{\Character{106}}%
\NativeMathAt{303.619}{701.692}{7.97}{\infty}%
\TextAt{14}{314.315}{716.237}{p}%
\TextAt{21}{320.298}{721.174}{\Character{23}}%
\TextAt{28}{324.46}{720.177}{p}%
\TextAt{9}{332.362}{716.237}{\Character{61}\Character{58}}%
\TextAt{20}{348.039}{727.595}{\Character{89}}%
\TextAt{27}{353.683}{702.489}{p}%
\TextAt{14}{365.308}{716.237}{p}%
\TextAt{21}{371.291}{721.173}{\Character{23}}%
\TextAt{28}{375.454}{720.177}{p}%
\TextAt{13}{380.034}{716.237}{\Character{59}}%
\TextAt{8}{79.226}{677.747}{where}%
\TextAt{13}{113.035}{677.747}{\Character{23}}%
\TextAt{27}{118.819}{675.953}{p}%
\NativeMathAt{126.627}{677.747}{11.955}{\in}%
\TextAt{18}{137.918}{677.747}{Z}%
\TextAt{8}{149.785}{677.747}{for}%
\TextAt{14}{167.67}{677.747}{p}%
\TextAt{18}{176.974}{677.747}{\Character{45}}%
\NativeMathAt{183.616}{677.747}{11.955}{\infty}%
\TextAt{8}{195.571}{677.747}{,}%
\TextAt{8}{202.72}{677.747}{and}%
\TextAt{13}{225.483}{677.747}{\Character{23}}%
\TextAt{27}{231.267}{675.953}{p}%
\NativeMathAt{239.074}{677.747}{11.955}{\in}%
\TextAt{18}{250.365}{677.747}{R}%
\TextAt{8}{262.897}{677.747}{for}%
\TextAt{14}{280.782}{677.747}{p}%
\TextAt{16}{290.086}{677.747}{\Character{106}}%
\NativeMathAt{296.709}{677.747}{11.955}{\infty}%
\TextAt{8}{308.683}{677.747}{.}%
\TextAt{8}{317.131}{677.747}{Put}%
\TextAt{14}{210.099}{649.906}{a}%
\TextAt{21}{216.082}{648.113}{f}%
\TextAt{9}{224.848}{649.906}{\Character{61}}%
\TextAt{20}{237.273}{661.264}{\Character{89}}%
\TextAt{27}{237.576}{635.432}{p}%
\TextAt{25}{241.564}{635.432}{\Character{45}}%
\NativeMathAt{243.778}{635.432}{7.97}{\infty}%
\TextAt{14}{254.542}{649.906}{p}%
\TextAt{21}{260.525}{654.843}{\Character{23}}%
\TextAt{28}{264.687}{653.846}{p}%
\TextAt{8}{293.179}{649.906}{and}%
\TextAt{14}{335.954}{649.906}{a}%
\NativeMathAt{341.937}{648.113}{7.97}{\infty}%
\TextAt{9}{354.225}{649.906}{\Character{61}}%
\TextAt{20}{366.65}{661.264}{\Character{89}}%
\TextAt{27}{366.883}{635.361}{p}%
\TextAt{3}{370.872}{635.361}{\Character{106}}%
\NativeMathAt{373.223}{635.361}{7.97}{\infty}%
\TextAt{14}{383.919}{649.906}{p}%
\TextAt{21}{389.902}{654.843}{\Character{23}}%
\TextAt{28}{394.064}{653.846}{p}%
\TextAt{13}{398.645}{649.906}{\Character{59}}%
\TextAt{8}{79.226}{611.416}{so}%
\TextAt{8}{93.681}{611.416}{that}%
\TextAt{14}{119.134}{611.416}{a}%
\TextAt{9}{128.582}{611.416}{\Character{61}}%
\TextAt{14}{141.152}{611.416}{a}%
\TextAt{21}{147.135}{609.622}{f}%
\TextAt{16}{155.294}{611.416}{\Character{2}}%
\TextAt{14}{167.307}{611.416}{a}%
\NativeMathAt{173.29}{609.622}{7.97}{\infty}%
\TextAt{8}{182.256}{611.416}{.}%
\TextAt{14}{190.958}{611.416}{a}%
\TextAt{21}{196.941}{609.622}{f}%
\TextAt{8}{206.368}{611.416}{is}%
\TextAt{8}{218.216}{611.416}{a}%
\TextAt{8}{228.056}{611.416}{fractional}%
\TextAt{8}{280.504}{611.416}{ideal}%
\TextAt{8}{308.552}{611.416}{of}%
\TextAt{13}{321.967}{611.416}{K}%
\TextAt{8}{332.784}{611.416}{,}%
\TextAt{8}{340.039}{611.416}{and}%
\TextAt{14}{362.886}{611.416}{a}%
\NativeMathAt{368.869}{609.622}{7.97}{\infty}%
\TextAt{8}{381.818}{611.416}{is}%
\TextAt{8}{393.666}{611.416}{an}%
\TextAt{8}{410.01}{611.416}{element}%
\TextAt{8}{453.336}{611.416}{of}%
\TextAt{20}{466.75}{620.382}{\Character{81}}%
\TextAt{27}{478.042}{607.929}{p}%
\TextAt{3}{482.03}{607.929}{\Character{106}}%
\NativeMathAt{484.381}{607.929}{7.97}{\infty}%
\TextAt{18}{495.342}{611.416}{R}%
\TextAt{3}{503.976}{615.754}{\Character{3}}%
\TextAt{17}{503.976}{608.46}{\Character{43}}%
\TextAt{8}{511.06}{611.416}{.}%
\TextAt{8}{519.762}{611.416}{At}%
\TextAt{8}{79.226}{594.37}{the}%
\TextAt{8}{99.383}{594.37}{same}%
\TextAt{8}{128.709}{594.37}{time,}%
\TextAt{8}{158.62}{594.37}{we}%
\TextAt{8}{175.848}{594.37}{view}%
\TextAt{14}{202.831}{594.37}{a}%
\TextAt{21}{208.814}{592.576}{f}%
\TextAt{8}{214.259}{594.37}{,}%
\TextAt{8}{221.408}{594.37}{resp.}%
\TextAt{14}{250.73}{594.37}{a}%
\NativeMathAt{256.713}{592.576}{7.97}{\infty}%
\TextAt{8}{265.679}{594.37}{,}%
\TextAt{8}{272.829}{594.37}{as}%
\TextAt{8}{287.199}{594.37}{replete}%
\TextAt{8}{325.563}{594.37}{ideals}%
\TextAt{14}{228.294}{566.529}{a}%
\TextAt{21}{234.277}{564.736}{f}%
\TextAt{16}{242.378}{566.529}{\Character{2}}%
\TextAt{20}{254.333}{577.887}{\Character{89}}%
\TextAt{27}{254.567}{551.984}{p}%
\TextAt{3}{258.555}{551.984}{\Character{106}}%
\NativeMathAt{260.906}{551.984}{7.97}{\infty}%
\TextAt{9}{271.602}{566.529}{1}%
\TextAt{13}{277.455}{566.529}{\Character{59}}%
\TextAt{8}{306.609}{566.529}{resp.}%
\TextAt{9}{335.931}{566.529}{\Character{40}1\Character{41}}%
\TextAt{16}{353.546}{566.529}{\Character{2}}%
\TextAt{14}{365.501}{566.529}{a}%
\NativeMathAt{371.484}{564.736}{7.97}{\infty}%
\TextAt{13}{380.451}{566.529}{\Character{58}}%
\TextAt{8}{79.226}{528.039}{Thus}%
\TextAt{8}{108.875}{528.039}{for}%
\TextAt{8}{126.76}{528.039}{all}%
\TextAt{8}{143.019}{528.039}{elements}%
\TextAt{8}{190.875}{528.039}{of}%
\TextAt{13}{204.206}{528.039}{J}%
\TextAt{9}{211.783}{528.039}{\Character{40}}%
\Vector{q 0 G 0 g 0.478 w 144.335 -238.605 m 150.179 -238.605 l S Q}%
\TextAt{14}{216.335}{528.039}{o}%
\TextAt{9}{222.179}{528.039}{\Character{41}}%
\TextAt{8}{230.629}{528.039}{the}%
\TextAt{8}{250.785}{528.039}{decomposition}%
\TextAt{14}{277.627}{500.198}{a}%
\TextAt{9}{286.931}{500.198}{\Character{61}}%
\TextAt{14}{299.357}{500.198}{a}%
\TextAt{21}{305.34}{498.405}{f}%
\TextAt{16}{313.441}{500.198}{\Character{1}}%
\TextAt{14}{319.419}{500.198}{a}%
\NativeMathAt{325.402}{498.405}{7.97}{\infty}%
\TextAt{8}{79.226}{472.358}{applies.}%
\TextAt{8}{122.859}{472.358}{To}%
\TextAt{13}{140.086}{472.358}{a}%
\NativeMathAt{149.552}{472.358}{11.955}{\in}%
\TextAt{13}{160.843}{472.358}{K}%
\TextAt{3}{171.66}{476.696}{\Character{3}}%
\TextAt{8}{180.29}{472.358}{we}%
\TextAt{8}{197.517}{472.358}{associate}%
\TextAt{8}{246.749}{472.358}{the}%
\TextAt{12}{266.905}{472.358}{replete}%
\TextAt{12}{311.737}{472.358}{principal}%
\TextAt{12}{367.879}{472.358}{ideal}%
\TextAt{9}{224.33}{444.518}{\Character{91}}%
\TextAt{13}{227.582}{444.518}{a}%
\TextAt{9}{233.727}{444.518}{\Character{93}}%
\TextAt{9}{240.278}{444.518}{\Character{61}}%
\TextAt{20}{252.725}{455.875}{\Character{89}}%
\TextAt{27}{258.369}{430.769}{p}%
\TextAt{14}{269.993}{444.518}{p}%
\TextAt{21}{275.977}{449.454}{v}%
\TextAt{28}{280.1}{448.457}{p}%
\TextAt{17}{284.183}{449.454}{\Character{40}}%
\TextAt{21}{287.476}{449.454}{a}%
\TextAt{17}{291.974}{449.454}{\Character{41}}%
\TextAt{9}{299.086}{444.518}{\Character{61}}%
\TextAt{9}{311.495}{444.518}{\Character{40}}%
\TextAt{13}{316.064}{444.518}{a}%
\TextAt{9}{322.209}{444.518}{\Character{41}}%
\TextAt{16}{329.418}{444.518}{\Character{2}}%
\TextAt{20}{341.373}{455.875}{\Character{89}}%
\TextAt{27}{341.607}{429.972}{p}%
\TextAt{3}{345.595}{429.972}{\Character{106}}%
\NativeMathAt{347.946}{429.972}{7.97}{\infty}%
\TextAt{14}{358.642}{444.518}{p}%
\TextAt{21}{364.625}{449.454}{v}%
\TextAt{28}{368.748}{448.458}{p}%
\TextAt{17}{372.831}{449.454}{\Character{40}}%
\TextAt{21}{376.124}{449.454}{a}%
\TextAt{17}{380.622}{449.454}{\Character{41}}%
\TextAt{13}{384.414}{444.518}{\Character{58}}%
\TextAt{8}{79.226}{406.027}{These}%
\TextAt{8}{113.095}{406.027}{replete}%
\TextAt{8}{151.459}{406.027}{ideals}%
\TextAt{8}{184.037}{406.027}{form}%
\TextAt{8}{211.677}{406.027}{a}%
\TextAt{8}{221.433}{406.027}{subgroup}%
\TextAt{13}{272.23}{406.027}{P}%
\TextAt{9}{281.402}{406.027}{\Character{40}}%
\Vector{q 0 G 0 g 0.478 w 213.954 -360.617 m 219.798 -360.617 l S Q}%
\TextAt{14}{285.954}{406.027}{o}%
\TextAt{9}{291.798}{406.027}{\Character{41}}%
\TextAt{8}{300.248}{406.027}{of}%
\TextAt{13}{313.578}{406.027}{J}%
\TextAt{9}{321.155}{406.027}{\Character{40}}%
\Vector{q 0 G 0 g 0.478 w 253.707 -360.617 m 259.551 -360.617 l S Q}%
\TextAt{14}{325.707}{406.027}{o}%
\TextAt{9}{331.551}{406.027}{\Character{41}}%
\TextAt{8}{336.104}{406.027}{.}%
\TextAt{8}{344.552}{406.027}{The}%
\TextAt{8}{368.606}{406.027}{factor}%
\TextAt{8}{402.104}{406.027}{group}%
\TextAt{8}{256.193}{378.187}{Pic}%
\TextAt{9}{272.608}{378.187}{\Character{40}}%
\Vector{q 0 G 0 g 0.478 w 205.16 -388.457 m 211.005 -388.457 l S Q}%
\TextAt{14}{277.16}{378.187}{o}%
\TextAt{9}{283.005}{378.187}{\Character{41}}%
\TextAt{9}{290.86}{378.187}{\Character{61}}%
\TextAt{13}{303.303}{378.187}{J}%
\TextAt{9}{310.88}{378.187}{\Character{40}}%
\Vector{q 0 G 0 g 0.478 w 243.433 -388.457 m 249.277 -388.457 l S Q}%
\TextAt{14}{315.433}{378.187}{o}%
\TextAt{9}{321.277}{378.187}{\Character{41}\Character{47}}%
\TextAt{13}{331.682}{378.187}{P}%
\TextAt{9}{340.853}{378.187}{\Character{40}}%
\Vector{q 0 G 0 g 0.478 w 273.406 -388.457 m 279.25 -388.457 l S Q}%
\TextAt{14}{345.406}{378.187}{o}%
\TextAt{9}{351.25}{378.187}{\Character{41}}%
\TextAt{8}{79.226}{350.346}{is}%
\TextAt{8}{90.99}{350.346}{called}%
\TextAt{8}{124.154}{350.346}{the}%
\TextAt{12}{144.31}{350.346}{replete}%
\TextAt{12}{189.142}{350.346}{ideal}%
\TextAt{12}{221.254}{350.346}{class}%
\TextAt{12}{252.612}{350.346}{group}%
\TextAt{8}{286.517}{350.346}{,}%
\TextAt{8}{293.666}{350.346}{or}%
\TextAt{12}{307.976}{350.346}{replete}%
\TextAt{12}{352.808}{350.346}{Picard}%
\TextAt{12}{395.715}{350.346}{group}%
\TextAt{8}{429.62}{350.346}{.}%
\TextAt{12}{79.226}{324.665}{Definition}%
\TextAt{12}{141.801}{324.665}{(1.5).}%
\TextAt{8}{176.932}{324.665}{The}%
\TextAt{12}{200.8}{324.665}{absolute}%
\TextAt{12}{253.703}{324.665}{norm}%
\TextAt{8}{288.331}{324.665}{of}%
\TextAt{8}{301.476}{324.665}{a}%
\TextAt{8}{311.046}{324.665}{replete}%
\TextAt{8}{349.225}{324.665}{ideal}%
\TextAt{14}{377.002}{324.665}{a}%
\TextAt{9}{386.306}{324.666}{\Character{61}}%
\TextAt{20}{398.732}{333.632}{\Character{81}}%
\TextAt{27}{410.023}{321.179}{p}%
\TextAt{14}{416.502}{324.665}{p}%
\TextAt{21}{422.485}{329.004}{\Character{23}}%
\TextAt{28}{426.648}{328.008}{p}%
\TextAt{8}{434.941}{324.665}{is}%
\TextAt{8}{446.519}{324.665}{defined}%
\TextAt{8}{486.647}{324.665}{to}%
\TextAt{8}{500.772}{324.665}{be}%
\TextAt{8}{516.51}{324.665}{the}%
\TextAt{8}{79.226}{307.62}{positive}%
\TextAt{8}{122.54}{307.62}{real}%
\TextAt{8}{145.302}{307.62}{number}%
\TextAt{14}{258.457}{290.574}{N}%
\TextAt{9}{268.405}{290.574}{\Character{40}}%
\TextAt{14}{272.957}{290.574}{a}%
\TextAt{9}{278.94}{290.574}{\Character{41}}%
\TextAt{9}{286.795}{290.574}{\Character{61}}%
\TextAt{20}{299.239}{301.931}{\Character{89}}%
\TextAt{27}{304.883}{276.825}{p}%
\TextAt{14}{316.508}{290.574}{N}%
\TextAt{9}{326.456}{290.574}{\Character{40}}%
\TextAt{14}{331.008}{290.574}{p}%
\TextAt{9}{336.991}{290.574}{\Character{41}}%
\TextAt{21}{341.544}{295.51}{\Character{23}}%
\TextAt{28}{345.706}{294.513}{p}%
\TextAt{13}{350.287}{290.574}{\Character{58}}%
\TextAt{8}{97.159}{250.217}{The}%
\TextAt{8}{121.213}{250.217}{absolute}%
\TextAt{8}{167.455}{250.217}{norm}%
\TextAt{8}{198.025}{250.217}{is}%
\TextAt{8}{209.789}{250.217}{multiplicative}%
\TextAt{8}{283.612}{250.217}{and}%
\TextAt{8}{306.374}{250.217}{induces}%
\TextAt{8}{348.05}{250.217}{a}%
\TextAt{8}{357.805}{250.217}{surjective}%
\TextAt{8}{410.217}{250.217}{homomorphism}%
\TextAt{14}{262.377}{222.376}{N}%
\TextAt{9}{275.646}{222.376}{\Character{58}}%
\TextAt{13}{282.219}{222.376}{J}%
\TextAt{9}{289.796}{222.376}{\Character{40}}%
\Vector{q 0 G 0 g 0.478 w 222.348 -544.267 m 228.192 -544.267 l S Q}%
\TextAt{14}{294.348}{222.376}{o}%
\TextAt{9}{300.192}{222.376}{\Character{41}}%
\NativeArrowAt{308.066}{222.376}{11.955}{19.26}{\longrightarrow}%
\TextAt{18}{330.648}{222.376}{R}%
\TextAt{3}{339.282}{227.312}{\Character{3}}%
\TextAt{17}{339.282}{219.421}{\Character{43}}%
\TextAt{13}{346.367}{222.376}{\Character{58}}%
\TextAt{8}{79.226}{194.536}{The}%
\TextAt{8}{104.286}{194.536}{absolute}%
\TextAt{8}{151.534}{194.536}{norm}%
\TextAt{8}{183.11}{194.536}{of}%
\TextAt{8}{197.446}{194.536}{a}%
\TextAt{8}{208.207}{194.536}{replete}%
\TextAt{8}{247.577}{194.536}{principal}%
\TextAt{8}{297.36}{194.536}{ideal}%
\TextAt{9}{326.329}{194.536}{\Character{91}}%
\TextAt{13}{329.581}{194.536}{a}%
\TextAt{9}{335.726}{194.536}{\Character{93}}%
\TextAt{8}{343.881}{194.536}{is}%
\TextAt{8}{356.651}{194.536}{equal}%
\TextAt{8}{388.549}{194.536}{to}%
\TextAt{9}{403.865}{194.536}{1}%
\TextAt{8}{414.622}{194.536}{in}%
\TextAt{8}{429.28}{194.536}{view}%
\TextAt{8}{457.269}{194.536}{of}%
\TextAt{8}{471.605}{194.536}{the}%
\TextAt{8}{492.767}{194.536}{product}%
\TextAt{8}{79.226}{177.49}{formula}%
\TextAt{8}{122.157}{177.49}{(1.3),}%
\TextAt{14}{210.242}{160.444}{N}%
\TextAt{9}{220.19}{160.444}{\Character{40}\Character{91}}%
\TextAt{13}{227.994}{160.444}{a}%
\TextAt{9}{234.139}{160.444}{\Character{93}\Character{41}}%
\TextAt{9}{245.233}{160.444}{\Character{61}}%
\TextAt{20}{257.689}{171.801}{\Character{89}}%
\TextAt{27}{263.333}{146.695}{p}%
\TextAt{14}{274.958}{160.444}{N}%
\TextAt{9}{284.906}{160.444}{\Character{40}}%
\TextAt{14}{289.458}{160.444}{p}%
\TextAt{9}{295.441}{160.444}{\Character{41}}%
\TextAt{21}{299.994}{165.38}{v}%
\TextAt{28}{304.117}{164.384}{p}%
\TextAt{17}{308.2}{165.38}{\Character{40}}%
\TextAt{21}{311.493}{165.38}{a}%
\TextAt{17}{315.991}{165.38}{\Character{41}}%
\TextAt{9}{323.103}{160.444}{\Character{61}}%
\TextAt{20}{335.529}{171.801}{\Character{89}}%
\TextAt{27}{341.172}{146.695}{p}%
\TextAt{16}{352.797}{160.444}{\Character{106}}%
\TextAt{13}{356.118}{160.444}{a}%
\TextAt{16}{362.263}{160.444}{\Character{106}}%
\TextAt{3}{365.584}{165.38}{\Character{0}}%
\TextAt{17}{372.171}{165.38}{1}%
\TextAt{27}{365.584}{157.488}{p}%
\TextAt{9}{380.224}{160.444}{\Character{61}}%
\TextAt{9}{392.634}{160.444}{1}%
\TextAt{13}{398.502}{160.444}{\Character{58}}%
\TextAt{8}{79.226}{128.602}{We}%
\TextAt{8}{99.371}{128.602}{therefore}%
\TextAt{8}{148.47}{128.602}{obtain}%
\TextAt{8}{184.898}{128.602}{a}%
\TextAt{8}{194.653}{128.602}{surjective}%
\TextAt{8}{247.065}{128.602}{homomorphism}%
\TextAt{14}{259.585}{100.762}{N}%
\TextAt{9}{272.853}{100.762}{\Character{58}}%
\TextAt{8}{279.426}{100.762}{Pic}%
\TextAt{9}{295.84}{100.762}{\Character{40}}%
\Vector{q 0 G 0 g 0.478 w 228.393 -665.882 m 234.237 -665.882 l S Q}%
\TextAt{14}{300.393}{100.762}{o}%
\TextAt{9}{306.237}{100.762}{\Character{41}}%
\NativeArrowAt{314.11}{100.762}{11.955}{19.26}{\longrightarrow}%
\TextAt{18}{336.692}{100.762}{R}%
\TextAt{3}{345.326}{105.698}{\Character{3}}%
\TextAt{17}{345.326}{97.806}{\Character{43}}%
\TextAt{8}{79.226}{72.921}{on}%
\TextAt{8}{95.485}{72.921}{the}%
\TextAt{8}{115.642}{72.921}{replete}%
\TextAt{8}{154.006}{72.921}{Picard}%
\TextAt{8}{191.234}{72.921}{group.}%
\TextAt{8}{297.211}{43.033}{170}%
\Destination{78.23}{809.815}{nav.page.170}%
\Destination{79.36}{776.155}{nav.print.186}%
\Destination{79.36}{338.655}{nav.number.3.1.5}%
\Hotlink{121.657}{175.255}{25.078}{11.047}{nav.number.3.1.3}%
}

% Page 183 (body page 171)
\Page{%
\TextAt{2}{79.226}{796.943}{III.}%
\TextAt{2}{98.317}{796.943}{Riemann–Roch}%
\TextAt{2}{174.375}{796.943}{Theory}%
\Vector{q 0 G 0 g 0.299 w 7.226 18.729 m 460.769 18.729 l S Q}%
\TextAt{8}{97.159}{761.124}{The}%
\TextAt{8}{120.621}{761.124}{relations}%
\TextAt{8}{167.574}{761.124}{between}%
\TextAt{8}{212.173}{761.124}{the}%
\TextAt{8}{231.737}{761.124}{replete}%
\TextAt{8}{269.509}{761.124}{ideals}%
\TextAt{8}{301.495}{761.124}{of}%
\TextAt{8}{314.233}{761.124}{a}%
\TextAt{8}{323.396}{761.124}{number}%
\TextAt{8}{365.401}{761.124}{field}%
\TextAt{13}{390.165}{761.124}{K}%
\TextAt{8}{404.288}{761.124}{and}%
\TextAt{8}{426.459}{761.124}{those}%
\TextAt{8}{456.496}{761.124}{of}%
\TextAt{8}{469.234}{761.124}{an}%
\TextAt{8}{484.901}{761.124}{extension}%
\TextAt{8}{79.226}{744.078}{field}%
\TextAt{13}{104.583}{744.078}{L}%
\TextAt{8}{116.445}{744.078}{are}%
\TextAt{8}{135.956}{744.078}{afforded}%
\TextAt{8}{181.158}{744.078}{by}%
\TextAt{8}{197.417}{744.078}{the}%
\TextAt{8}{217.574}{744.078}{two}%
\TextAt{8}{239.691}{744.078}{homomorphisms}%
\TextAt{13}{255.116}{708.028}{J}%
\TextAt{9}{262.693}{708.028}{\Character{40}}%
\NativeDiagram{%
\path[draw=black,line width=0.478bp] (195.245,-58.616) -- (201.09,-58.616);%
}%
\TextAt{14}{267.245}{708.028}{o}%
\TextAt{21}{273.09}{706.234}{K}%
\TextAt{9}{281.292}{708.028}{\Character{41}}%
\TextAt{13}{324.843}{708.028}{J}%
\TextAt{9}{332.42}{708.028}{\Character{40}}%
\NativeDiagram{%
\path[draw=black,line width=0.478bp] (264.972,-58.616) -- (270.816,-58.616);%
}%
\TextAt{14}{336.972}{708.028}{o}%
\TextAt{21}{342.816}{706.234}{L}%
\TextAt{9}{349.073}{708.028}{\Character{41}}%
\NativeDiagram{%
\path[draw=black,line width=0.478bp] (219.198,-58.583) .. controls (230.049,-56.36) and (236.675,-56.408) .. (247.025,-58.687);%
\path[draw=black,line width=0.478bp,line cap=round,line join=round] (245.447,-55.402) .. controls (245.52,-57.181) and (246.334,-58.192) .. (247.258,-58.739) .. controls (246.19,-58.846) and (245.027,-59.422) .. (244.213,-61.006);%
}%
\TextAt{21}{296.73}{723.088}{i}%
\TextAt{22}{299.613}{721.35}{L}%
\TextAt{23}{304.672}{721.35}{\Character{106}}%
\TextAt{22}{306.941}{721.35}{K}%
\NativeDiagram{%
\path[draw=black,line width=0.478bp] (247.492,-67.176) .. controls (236.675,-69.558) and (230.049,-69.606) .. (219.666,-67.479);%
\path[draw=black,line width=0.478bp,line cap=round,line join=round] (221.292,-70.74) .. controls (221.192,-68.963) and (220.364,-67.963) .. (219.432,-67.431) .. controls (220.498,-67.307) and (221.653,-66.714) .. (222.444,-65.119);%
}%
\TextAt{21}{294.811}{696.732}{N}%
\TextAt{22}{301.532}{694.994}{L}%
\TextAt{23}{306.591}{694.994}{\Character{106}}%
\TextAt{22}{308.86}{694.994}{K}%
\TextAt{13}{358.779}{708.028}{\Character{59}}%
\TextAt{8}{79.226}{671.801}{which}%
\TextAt{8}{112.713}{671.801}{are}%
\TextAt{8}{132.223}{671.801}{defined}%
\TextAt{8}{172.536}{671.801}{by}%
\TextAt{13}{126.796}{637.734}{i}%
\TextAt{21}{130.79}{635.875}{L}%
\TextAt{3}{136.549}{635.875}{\Character{106}}%
\TextAt{21}{138.901}{635.875}{K}%
\TextAt{20}{149.096}{658.178}{\Character{32}}%
\TextAt{20}{158.561}{649.092}{\Character{89}}%
\TextAt{27}{164.205}{623.986}{p}%
\TextAt{14}{175.829}{637.734}{p}%
\TextAt{21}{181.813}{642.671}{\Character{23}}%
\TextAt{28}{185.975}{641.674}{p}%
\TextAt{20}{190.556}{658.178}{\Character{33}}%
\TextAt{9}{203.341}{637.734}{\Character{61}}%
\TextAt{20}{215.767}{649.092}{\Character{89}}%
\TextAt{27}{221.411}{623.986}{p}%
\TextAt{20}{233.035}{649.092}{\Character{89}}%
\TextAt{27}{234.205}{623.189}{P}%
\TextAt{3}{240.801}{623.189}{\Character{106}}%
\TextAt{27}{243.153}{623.189}{p}%
\TextAt{14}{250.304}{637.734}{P}%
\TextAt{21}{260.197}{642.671}{e}%
\TextAt{28}{264.126}{640.933}{P}%
\TextAt{23}{269.825}{640.933}{\Character{106}}%
\TextAt{28}{272.094}{640.933}{p}%
\TextAt{21}{276.177}{642.671}{\Character{23}}%
\TextAt{28}{280.339}{641.674}{p}%
\TextAt{13}{284.92}{637.734}{\Character{59}}%
\TextAt{13}{314.055}{637.734}{N}%
\TextAt{21}{323.449}{635.875}{L}%
\TextAt{3}{329.208}{635.875}{\Character{106}}%
\TextAt{21}{331.56}{635.875}{K}%
\TextAt{20}{341.755}{658.178}{\Character{32}}%
\TextAt{20}{351.22}{649.092}{\Character{89}}%
\TextAt{27}{355.56}{623.656}{P}%
\TextAt{14}{368.488}{637.734}{P}%
\TextAt{21}{378.381}{642.671}{\Character{23}}%
\TextAt{28}{382.544}{641.28}{P}%
\TextAt{20}{389.239}{658.178}{\Character{33}}%
\TextAt{9}{402.025}{637.734}{\Character{61}}%
\TextAt{20}{414.45}{649.092}{\Character{89}}%
\TextAt{27}{420.094}{623.986}{p}%
\TextAt{20}{431.719}{649.092}{\Character{89}}%
\TextAt{27}{432.889}{623.189}{P}%
\TextAt{3}{439.484}{623.189}{\Character{106}}%
\TextAt{27}{441.836}{623.189}{p}%
\TextAt{14}{448.987}{637.734}{p}%
\TextAt{21}{454.971}{642.671}{f}%
\TextAt{28}{459.039}{640.933}{P}%
\TextAt{23}{464.738}{640.933}{\Character{106}}%
\TextAt{28}{467.007}{640.933}{p}%
\TextAt{21}{471.09}{642.671}{\Character{23}}%
\TextAt{28}{475.252}{641.279}{P}%
\TextAt{13}{481.948}{637.734}{\Character{58}}%
\TextAt{8}{79.226}{600.998}{Here}%
\TextAt{8}{108.783}{600.998}{the}%
\TextAt{8}{130.867}{600.998}{various}%
\TextAt{8}{172.869}{600.998}{product}%
\TextAt{8}{218.697}{600.998}{signs}%
\TextAt{8}{249.366}{600.998}{have}%
\TextAt{8}{278.289}{600.998}{to}%
\TextAt{8}{294.527}{600.998}{be}%
\TextAt{8}{312.38}{600.998}{read}%
\TextAt{8}{340.322}{600.998}{according}%
\TextAt{8}{394.937}{600.998}{to}%
\TextAt{8}{411.176}{600.998}{our}%
\TextAt{8}{433.918}{600.998}{convention.}%
\TextAt{8}{502.798}{600.998}{These}%
\TextAt{8}{79.226}{583.952}{homomorphisms}%
\TextAt{8}{166.511}{583.952}{satisfy}%
\TextAt{8}{203.058}{583.952}{the}%
\TextAt{12}{79.226}{558.581}{Proposition}%
\TextAt{12}{151.937}{558.581}{(1.6).}%
\TextAt{8}{89.579}{538.547}{(i)}%
\TextAt{8}{107.918}{538.547}{For}%
\TextAt{8}{130.417}{538.547}{a}%
\TextAt{8}{141.701}{538.547}{chain}%
\TextAt{8}{174.123}{538.547}{of}%
\TextAt{8}{188.982}{538.547}{fields}%
\TextAt{13}{220.482}{538.547}{K}%
\NativeMathAt{237.227}{538.547}{11.955}{\subseteq}%
\TextAt{13}{252.452}{538.547}{L}%
\NativeMathAt{266.343}{538.547}{11.955}{\subseteq}%
\TextAt{13}{281.568}{538.547}{M}%
\TextAt{8}{294.141}{538.547}{,}%
\TextAt{8}{303.202}{538.547}{one}%
\TextAt{8}{326.191}{538.547}{has}%
\TextAt{13}{348.593}{538.547}{N}%
\TextAt{21}{357.967}{536.687}{M}%
\TextAt{3}{366.949}{536.687}{\Character{106}}%
\TextAt{21}{369.301}{536.687}{K}%
\TextAt{9}{383.43}{538.547}{\Character{61}}%
\TextAt{13}{398.462}{538.547}{N}%
\TextAt{21}{407.835}{536.687}{L}%
\TextAt{3}{413.594}{536.687}{\Character{106}}%
\TextAt{21}{415.947}{536.687}{K}%
\TextAt{16}{427.848}{538.547}{\Character{14}}%
\TextAt{13}{437.525}{538.547}{N}%
\TextAt{21}{446.899}{536.687}{M}%
\TextAt{3}{455.881}{536.687}{\Character{106}}%
\TextAt{21}{458.233}{536.687}{L}%
\TextAt{8}{469.916}{538.547}{and}%
\TextAt{13}{494.208}{538.547}{i}%
\TextAt{21}{498.202}{536.687}{M}%
\TextAt{3}{507.183}{536.687}{\Character{106}}%
\TextAt{21}{509.536}{536.687}{K}%
\TextAt{9}{523.665}{538.547}{\Character{61}}%
\TextAt{13}{107.918}{521.501}{i}%
\TextAt{21}{111.912}{519.641}{M}%
\TextAt{3}{120.894}{519.641}{\Character{106}}%
\TextAt{21}{123.246}{519.641}{L}%
\TextAt{16}{132.16}{521.501}{\Character{14}}%
\TextAt{13}{140.794}{521.501}{i}%
\TextAt{21}{144.787}{519.641}{L}%
\TextAt{3}{150.547}{519.641}{\Character{106}}%
\TextAt{21}{152.899}{519.641}{K}%
\TextAt{8}{161.102}{521.501}{.}%
\TextAt{8}{86.327}{504.455}{(ii)}%
\TextAt{13}{107.918}{504.455}{N}%
\TextAt{21}{117.292}{502.595}{L}%
\TextAt{3}{123.051}{502.595}{\Character{106}}%
\TextAt{21}{125.404}{502.595}{K}%
\TextAt{9}{133.606}{504.455}{\Character{40}}%
\TextAt{13}{138.159}{504.455}{i}%
\TextAt{21}{142.152}{502.595}{L}%
\TextAt{3}{147.911}{502.595}{\Character{106}}%
\TextAt{21}{150.263}{502.595}{K}%
\TextAt{14}{158.466}{504.455}{a}%
\TextAt{9}{164.449}{504.455}{\Character{41}}%
\TextAt{9}{172.304}{504.455}{\Character{61}}%
\TextAt{14}{184.748}{504.455}{a}%
\TextAt{17}{190.731}{508.793}{\Character{91}}%
\TextAt{21}{193.083}{508.793}{L}%
\TextAt{17}{198.842}{508.793}{\Character{58}}%
\TextAt{21}{201.195}{508.793}{K}%
\TextAt{17}{208.899}{508.793}{\Character{93}}%
\TextAt{8}{215.647}{504.455}{for}%
\TextAt{14}{233.532}{504.455}{a}%
\NativeMathAt{242.836}{504.455}{11.955}{\in}%
\TextAt{13}{254.127}{504.455}{J}%
\TextAt{9}{261.704}{504.455}{\Character{40}}%
\Vector{q 0 G 0 g 0.478 w 194.256 -262.189 m 200.1 -262.189 l S Q}%
\TextAt{14}{266.256}{504.455}{o}%
\TextAt{21}{272.1}{502.661}{K}%
\TextAt{9}{280.303}{504.455}{\Character{41}}%
\TextAt{8}{284.855}{504.455}{.}%
\TextAt{8}{83.076}{487.409}{(iii)}%
\TextAt{14}{107.918}{487.409}{N}%
\TextAt{9}{117.867}{487.409}{\Character{40}}%
\TextAt{13}{122.419}{487.409}{N}%
\TextAt{21}{131.793}{485.549}{L}%
\TextAt{3}{137.552}{485.549}{\Character{106}}%
\TextAt{21}{139.904}{485.549}{K}%
\TextAt{9}{148.107}{487.409}{\Character{40}}%
\TextAt{14}{152.659}{487.409}{A}%
\TextAt{9}{161.238}{487.409}{\Character{41}\Character{41}}%
\TextAt{9}{173.636}{487.409}{\Character{61}}%
\TextAt{14}{186.089}{487.409}{N}%
\TextAt{9}{196.037}{487.409}{\Character{40}}%
\TextAt{14}{200.59}{487.409}{A}%
\TextAt{9}{209.169}{487.409}{\Character{41}}%
\TextAt{8}{217.618}{487.409}{for}%
\TextAt{14}{235.503}{487.409}{A}%
\NativeMathAt{247.403}{487.409}{11.955}{\in}%
\TextAt{13}{258.694}{487.409}{J}%
\TextAt{9}{266.271}{487.409}{\Character{40}}%
\Vector{q 0 G 0 g 0.478 w 198.824 -279.235 m 204.668 -279.235 l S Q}%
\TextAt{14}{270.824}{487.409}{o}%
\TextAt{21}{276.668}{485.615}{L}%
\TextAt{9}{282.925}{487.409}{\Character{41}}%
\TextAt{8}{287.477}{487.409}{.}%
\TextAt{8}{83.398}{470.363}{(iv)}%
\TextAt{8}{107.918}{470.363}{If}%
\TextAt{13}{119.571}{470.363}{L}%
\TextAt{16}{130.856}{470.363}{\Character{106}}%
\TextAt{13}{137.498}{470.363}{K}%
\TextAt{8}{152.173}{470.363}{is}%
\TextAt{8}{163.897}{470.363}{Galois}%
\TextAt{8}{199.77}{470.363}{with}%
\TextAt{8}{226.39}{470.363}{Galois}%
\TextAt{8}{262.264}{470.363}{group}%
\TextAt{13}{295.399}{470.363}{G}%
\TextAt{8}{304.633}{470.363}{,}%
\TextAt{8}{311.75}{470.363}{then}%
\TextAt{8}{338.37}{470.363}{for}%
\TextAt{8}{356.215}{470.363}{every}%
\TextAt{8}{387.067}{470.363}{prime}%
\TextAt{8}{420.191}{470.363}{ideal}%
\TextAt{14}{448.114}{470.363}{P}%
\TextAt{8}{461.865}{470.363}{of}%
\TextAt{14}{475.155}{470.363}{o}%
\TextAt{21}{480.999}{468.57}{L}%
\TextAt{8}{487.256}{470.363}{,}%
\TextAt{8}{494.374}{470.363}{one}%
\TextAt{8}{515.793}{470.363}{has}%
\TextAt{13}{107.918}{453.317}{N}%
\TextAt{21}{117.292}{451.457}{L}%
\TextAt{3}{123.051}{451.457}{\Character{106}}%
\TextAt{21}{125.404}{451.457}{K}%
\TextAt{9}{133.606}{453.317}{\Character{40}}%
\TextAt{14}{138.159}{453.317}{P}%
\TextAt{9}{148.052}{453.317}{\Character{41}}%
\TextAt{14}{152.604}{453.317}{o}%
\TextAt{21}{158.448}{451.524}{L}%
\TextAt{9}{168.026}{453.317}{\Character{61}}%
\TextAt{20}{180.452}{462.283}{\Character{81}}%
\TextAt{21}{191.743}{449.83}{\Character{27}}%
\NativeMathAt{196.868}{449.83}{7.97}{\in}%
\TextAt{21}{202.514}{449.83}{G}%
\TextAt{13}{211.603}{453.317}{\Character{27}}%
\TextAt{14}{218.686}{453.317}{P}%
\TextAt{8}{228.579}{453.317}{.}%
\TextAt{8}{86.65}{436.271}{(v)}%
\TextAt{8}{107.918}{436.271}{For}%
\TextAt{8}{128.888}{436.271}{any}%
\TextAt{8}{151.005}{436.271}{replete}%
\TextAt{8}{189.369}{436.271}{principal}%
\TextAt{8}{238.146}{436.271}{ideal}%
\TextAt{9}{266.109}{436.271}{\Character{91}}%
\TextAt{13}{269.361}{436.271}{a}%
\TextAt{9}{275.506}{436.271}{\Character{93}}%
\TextAt{8}{282.655}{436.271}{of}%
\TextAt{13}{295.985}{436.271}{K}%
\TextAt{8}{306.803}{436.271}{,}%
\TextAt{8}{313.952}{436.271}{resp.}%
\TextAt{13}{343.274}{436.271}{L}%
\TextAt{8}{351.238}{436.271}{,}%
\TextAt{8}{358.388}{436.271}{one}%
\TextAt{8}{379.847}{436.271}{has}%
\TextAt{13}{198.938}{409.515}{i}%
\TextAt{21}{202.932}{407.656}{L}%
\TextAt{3}{208.691}{407.656}{\Character{106}}%
\TextAt{21}{211.043}{407.656}{K}%
\TextAt{9}{219.246}{409.515}{\Character{40}\Character{91}}%
\TextAt{13}{227.05}{409.515}{a}%
\TextAt{9}{233.195}{409.515}{\Character{93}\Character{41}}%
\TextAt{9}{244.289}{409.515}{\Character{61}}%
\TextAt{9}{256.711}{409.515}{\Character{91}}%
\TextAt{13}{259.997}{409.515}{a}%
\TextAt{9}{266.142}{409.515}{\Character{93}}%
\TextAt{13}{269.393}{409.515}{\Character{59}}%
\TextAt{8}{298.548}{409.515}{resp.}%
\TextAt{13}{327.87}{409.515}{N}%
\TextAt{21}{337.244}{407.656}{L}%
\TextAt{3}{343.003}{407.656}{\Character{106}}%
\TextAt{21}{345.355}{407.656}{K}%
\TextAt{9}{353.558}{409.515}{\Character{40}\Character{91}}%
\TextAt{13}{361.362}{409.515}{a}%
\TextAt{9}{367.507}{409.515}{\Character{93}\Character{41}}%
\TextAt{9}{378.601}{409.515}{\Character{61}}%
\TextAt{9}{391.023}{409.515}{\Character{91}}%
\TextAt{13}{394.309}{409.515}{N}%
\TextAt{21}{403.683}{407.656}{L}%
\TextAt{3}{409.442}{407.656}{\Character{106}}%
\TextAt{21}{411.794}{407.656}{K}%
\TextAt{9}{419.997}{409.515}{\Character{40}}%
\TextAt{13}{424.549}{409.515}{a}%
\TextAt{9}{430.694}{409.515}{\Character{41}\Character{93}}%
\TextAt{13}{438.498}{409.515}{\Character{58}}%
\TextAt{8}{83.398}{382.76}{(vi)}%
\TextAt{13}{107.918}{382.76}{N}%
\TextAt{21}{117.292}{380.9}{L}%
\TextAt{3}{123.051}{380.9}{\Character{106}}%
\TextAt{21}{125.404}{380.9}{K}%
\TextAt{9}{133.606}{382.76}{\Character{40}}%
\TextAt{14}{138.159}{382.76}{A}%
\TextAt{21}{146.738}{380.967}{f}%
\TextAt{9}{152.183}{382.76}{\Character{41}}%
\TextAt{9}{162.716}{382.76}{\Character{61}}%
\TextAt{13}{177.832}{382.76}{N}%
\TextAt{21}{187.205}{380.9}{L}%
\TextAt{3}{192.964}{380.9}{\Character{106}}%
\TextAt{21}{195.317}{380.9}{K}%
\TextAt{9}{203.519}{382.76}{\Character{40}}%
\TextAt{14}{208.072}{382.76}{A}%
\TextAt{9}{216.651}{382.76}{\Character{41}}%
\TextAt{21}{221.203}{380.967}{f}%
\TextAt{8}{232.115}{382.76}{is}%
\TextAt{8}{245.449}{382.76}{the}%
\TextAt{8}{267.175}{382.76}{ideal}%
\TextAt{8}{296.708}{382.76}{of}%
\TextAt{13}{311.608}{382.76}{K}%
\TextAt{8}{327.893}{382.76}{generated}%
\TextAt{8}{382.795}{382.76}{by}%
\TextAt{8}{400.623}{382.76}{the}%
\TextAt{8}{422.35}{382.76}{norms}%
\TextAt{13}{459.103}{382.76}{N}%
\TextAt{21}{468.477}{380.9}{L}%
\TextAt{3}{474.236}{380.9}{\Character{106}}%
\TextAt{21}{476.589}{380.9}{K}%
\TextAt{9}{484.791}{382.76}{\Character{40}}%
\TextAt{13}{489.344}{382.76}{a}%
\TextAt{9}{495.489}{382.76}{\Character{41}}%
\TextAt{8}{505.508}{382.76}{of}%
\TextAt{8}{520.408}{382.76}{all}%
\TextAt{13}{107.918}{365.714}{a}%
\NativeMathAt{117.384}{365.714}{11.955}{\in}%
\TextAt{14}{128.675}{365.714}{A}%
\TextAt{21}{137.254}{363.921}{f}%
\TextAt{8}{142.699}{365.714}{.}%
\TextAt{19}{79.226}{332.018}{Proof.}%
\TextAt{8}{116}{332.018}{(i)}%
\TextAt{8}{131.781}{332.018}{is}%
\TextAt{8}{143.066}{332.018}{based}%
\TextAt{8}{175.165}{332.018}{on}%
\TextAt{8}{190.946}{332.018}{the}%
\TextAt{8}{210.624}{332.018}{transitivity}%
\TextAt{8}{271.033}{332.018}{of}%
\TextAt{8}{283.884}{332.018}{inertia}%
\TextAt{8}{320.479}{332.018}{degree}%
\TextAt{8}{356.416}{332.018}{and}%
\TextAt{8}{378.7}{332.018}{ramification}%
\TextAt{8}{443.27}{332.018}{index.}%
\TextAt{8}{479.199}{332.018}{(ii)}%
\TextAt{8}{498.231}{332.018}{follows}%
\TextAt{8}{79.226}{314.972}{from}%
\TextAt{8}{106.684}{314.972}{(1.2)}%
\TextAt{8}{134.476}{314.972}{in}%
\TextAt{8}{147.946}{314.972}{view}%
\TextAt{8}{174.746}{314.972}{of}%
\TextAt{8}{187.893}{314.972}{the}%
\TextAt{8}{207.867}{314.972}{fundamental}%
\TextAt{8}{275.327}{314.972}{identity}%
\TextAt{20}{318.398}{323.939}{\Character{80}}%
\TextAt{27}{331.017}{311.485}{P}%
\TextAt{3}{337.612}{311.485}{\Character{106}}%
\TextAt{27}{339.965}{311.485}{p}%
\TextAt{13}{346.444}{314.972}{f}%
\TextAt{27}{352.214}{313.113}{P}%
\TextAt{3}{358.81}{313.113}{\Character{106}}%
\TextAt{27}{361.162}{313.113}{p}%
\TextAt{13}{365.649}{314.972}{e}%
\TextAt{27}{371.074}{313.113}{P}%
\TextAt{3}{377.67}{313.113}{\Character{106}}%
\TextAt{27}{380.022}{313.113}{p}%
\TextAt{9}{387.83}{314.972}{\Character{61}}%
\TextAt{9}{400.24}{314.972}{\Character{91}}%
\TextAt{13}{403.507}{314.972}{L}%
\TextAt{9}{414.792}{314.972}{\Character{58}}%
\TextAt{13}{421.365}{314.972}{K}%
\TextAt{9}{432.182}{314.972}{\Character{93}}%
\TextAt{8}{435.434}{314.972}{.}%
\TextAt{8}{443.821}{314.972}{(iii)}%
\TextAt{8}{466.401}{314.972}{holds}%
\TextAt{8}{496.848}{314.972}{for}%
\TextAt{8}{514.55}{314.972}{any}%
\TextAt{8}{79.226}{297.926}{replete}%
\TextAt{8}{117.59}{297.926}{“prime}%
\TextAt{8}{156.277}{297.926}{ideal”}%
\TextAt{14}{189.764}{297.926}{P}%
\TextAt{8}{203.554}{297.926}{of}%
\TextAt{13}{216.884}{297.926}{L}%
\TextAt{8}{224.849}{297.926}{,}%
\TextAt{8}{231.998}{297.926}{by}%
\TextAt{8}{248.257}{297.926}{(1.2):}%
\TextAt{14}{192.757}{271.171}{N}%
\TextAt{9}{202.705}{271.171}{\Character{40}}%
\TextAt{13}{207.258}{271.171}{N}%
\TextAt{21}{216.631}{269.311}{L}%
\TextAt{3}{222.39}{269.311}{\Character{106}}%
\TextAt{21}{224.743}{269.311}{K}%
\TextAt{9}{232.946}{271.171}{\Character{40}}%
\TextAt{14}{237.498}{271.171}{P}%
\TextAt{9}{247.391}{271.171}{\Character{41}\Character{41}}%
\TextAt{9}{259.788}{271.171}{\Character{61}}%
\TextAt{14}{272.242}{271.171}{N}%
\TextAt{9}{282.19}{271.171}{\Character{40}}%
\TextAt{14}{286.742}{271.171}{p}%
\TextAt{21}{292.726}{276.107}{f}%
\TextAt{28}{296.794}{274.369}{P}%
\TextAt{23}{302.493}{274.369}{\Character{106}}%
\TextAt{28}{304.762}{274.369}{p}%
\TextAt{9}{309.343}{271.171}{\Character{41}}%
\TextAt{9}{317.198}{271.171}{\Character{61}}%
\TextAt{14}{329.641}{271.171}{N}%
\TextAt{9}{339.59}{271.171}{\Character{40}}%
\TextAt{14}{344.142}{271.171}{p}%
\TextAt{9}{350.125}{271.171}{\Character{41}}%
\TextAt{21}{354.677}{276.107}{f}%
\TextAt{28}{358.745}{274.369}{P}%
\TextAt{23}{364.444}{274.369}{\Character{106}}%
\TextAt{28}{366.714}{274.369}{p}%
\TextAt{9}{374.615}{271.171}{\Character{61}}%
\TextAt{14}{387.041}{271.171}{N}%
\TextAt{9}{396.989}{271.171}{\Character{40}}%
\TextAt{14}{401.541}{271.171}{P}%
\TextAt{9}{411.434}{271.171}{\Character{41}}%
\TextAt{13}{415.987}{271.171}{\Character{59}}%
\TextAt{8}{79.226}{244.415}{and}%
\TextAt{8}{101.989}{244.415}{therefore}%
\TextAt{8}{151.089}{244.415}{for}%
\TextAt{8}{168.974}{244.415}{all}%
\TextAt{8}{185.233}{244.415}{replete}%
\TextAt{8}{223.597}{244.415}{ideals}%
\TextAt{8}{256.175}{244.415}{of}%
\TextAt{13}{269.505}{244.415}{L}%
\TextAt{8}{277.469}{244.415}{.}%
\TextAt{8}{97.159}{224.381}{(iv)}%
\TextAt{8}{120.466}{224.381}{The}%
\TextAt{8}{145.388}{224.381}{prime}%
\TextAt{8}{179.419}{224.381}{ideal}%
\TextAt{14}{208.25}{224.381}{p}%
\TextAt{8}{218.999}{224.381}{lying}%
\TextAt{8}{248.81}{224.381}{below}%
\TextAt{14}{282.841}{224.381}{P}%
\TextAt{8}{297.5}{224.381}{decomposes}%
\TextAt{8}{361.897}{224.381}{in}%
\TextAt{8}{376.418}{224.381}{the}%
\TextAt{8}{397.442}{224.381}{ring}%
\TextAt{16}{422.375}{224.381}{\Character{79}}%
\TextAt{8}{436.99}{224.381}{of}%
\TextAt{8}{451.188}{224.381}{integers}%
\TextAt{8}{495.369}{224.381}{of}%
\TextAt{13}{509.567}{224.381}{L}%
\TextAt{8}{522.297}{224.381}{as}%
\TextAt{14}{79.226}{207.335}{p}%
\TextAt{9}{88.53}{207.335}{\Character{61}}%
\TextAt{9}{100.94}{207.335}{\Character{40}}%
\TextAt{14}{105.508}{207.335}{P}%
\TextAt{17}{115.401}{205.541}{1}%
\TextAt{16}{122.126}{207.335}{\Character{1}}%
\TextAt{16}{127.422}{207.335}{\Character{1}}%
\TextAt{16}{132.73}{207.335}{\Character{1}}%
\TextAt{14}{138.066}{207.335}{P}%
\TextAt{21}{147.959}{205.541}{r}%
\TextAt{9}{152.514}{207.335}{\Character{41}}%
\TextAt{21}{157.066}{211.673}{e}%
\TextAt{8}{161.493}{207.335}{,}%
\TextAt{8}{168.32}{207.335}{with}%
\TextAt{8}{194.577}{207.335}{prime}%
\TextAt{8}{227.338}{207.335}{ideals}%
\TextAt{14}{259.513}{207.335}{P}%
\TextAt{21}{269.406}{205.541}{i}%
\TextAt{9}{276.108}{207.335}{\Character{61}}%
\TextAt{13}{288.534}{207.335}{\Character{27}}%
\TextAt{21}{295.187}{205.541}{i}%
\TextAt{14}{298.569}{207.335}{P}%
\TextAt{8}{308.462}{207.335}{,}%
\TextAt{13}{315.289}{207.335}{\Character{27}}%
\TextAt{21}{321.942}{205.541}{i}%
\NativeMathAt{328.644}{207.335}{11.955}{\in}%
\TextAt{13}{339.935}{207.335}{G}%
\TextAt{9}{349.169}{207.335}{\Character{47}}%
\TextAt{13}{355.022}{207.335}{G}%
\TextAt{27}{364.256}{205.541}{P}%
\TextAt{8}{371.349}{207.335}{,}%
\TextAt{8}{378.176}{207.335}{which}%
\TextAt{8}{411.26}{207.335}{are}%
\TextAt{8}{430.368}{207.335}{conjugates}%
\TextAt{8}{487.589}{207.335}{of}%
\TextAt{14}{500.516}{207.335}{P}%
\TextAt{8}{513.904}{207.335}{and}%
\TextAt{8}{79.226}{190.289}{thus}%
\TextAt{8}{104.978}{190.289}{have}%
\TextAt{8}{131.972}{190.289}{the}%
\TextAt{8}{152.129}{190.289}{same}%
\TextAt{8}{181.455}{190.289}{inertia}%
\TextAt{8}{218.528}{190.289}{degree}%
\TextAt{13}{254.943}{190.289}{f}%
\TextAt{8}{261.99}{190.289}{.}%
\TextAt{8}{270.438}{190.289}{Therefore}%
\TextAt{13}{164.173}{157.561}{N}%
\TextAt{21}{173.546}{155.701}{L}%
\TextAt{3}{179.305}{155.701}{\Character{106}}%
\TextAt{21}{181.658}{155.701}{K}%
\TextAt{9}{189.86}{157.561}{\Character{40}}%
\TextAt{14}{194.413}{157.561}{P}%
\TextAt{9}{204.306}{157.561}{\Character{41}}%
\TextAt{16}{208.858}{157.561}{\Character{79}}%
\TextAt{9}{222.029}{157.561}{\Character{61}}%
\TextAt{14}{234.454}{157.561}{p}%
\TextAt{21}{240.437}{162.497}{f}%
\TextAt{16}{245.882}{157.561}{\Character{79}}%
\TextAt{9}{259.053}{157.561}{\Character{61}}%
\TextAt{21}{277.088}{172.505}{r}%
\TextAt{20}{271.478}{168.918}{\Character{89}}%
\TextAt{21}{272.265}{143.724}{i}%
\TextAt{17}{275.148}{143.724}{\Character{61}1}%
\TextAt{14}{288.747}{157.561}{P}%
\TextAt{21}{298.64}{163.229}{ef}%
\TextAt{21}{298.64}{154.497}{i}%
\TextAt{9}{311.334}{157.561}{\Character{61}}%
\TextAt{21}{329.37}{172.505}{r}%
\TextAt{20}{323.76}{168.918}{\Character{89}}%
\TextAt{21}{324.546}{143.724}{i}%
\TextAt{17}{327.429}{143.724}{\Character{61}1}%
\TextAt{20}{344.95}{168.918}{\Character{89}}%
\TextAt{21}{341.028}{143.547}{\Character{28}}%
\NativeMathAt{345.707}{143.547}{7.97}{\in}%
\TextAt{21}{351.352}{143.547}{G}%
\TextAt{28}{357.951}{142.156}{P}%
\TextAt{13}{366.141}{157.561}{\Character{27}}%
\TextAt{21}{372.794}{155.768}{i}%
\TextAt{13}{376.175}{157.561}{\Character{28}}%
\TextAt{14}{382.594}{157.561}{P}%
\TextAt{9}{395.808}{157.561}{\Character{61}}%
\TextAt{20}{409.281}{168.918}{\Character{89}}%
\TextAt{21}{408.234}{143.547}{\Character{27}}%
\NativeMathAt{413.359}{143.547}{7.97}{\in}%
\TextAt{21}{419.005}{143.547}{G}%
\TextAt{13}{427.596}{157.561}{\Character{27}}%
\TextAt{14}{434.678}{157.561}{P}%
\TextAt{13}{444.572}{157.561}{\Character{58}}%
\TextAt{8}{79.226}{120.157}{(v)}%
\TextAt{8}{98.414}{120.157}{For}%
\TextAt{8}{119.384}{120.157}{any}%
\TextAt{8}{141.501}{120.157}{element}%
\TextAt{13}{184.742}{120.157}{a}%
\NativeMathAt{194.208}{120.157}{11.955}{\in}%
\TextAt{13}{205.499}{120.157}{K}%
\TextAt{3}{216.317}{124.496}{\Character{3}}%
\TextAt{8}{221.049}{120.157}{,}%
\TextAt{8}{228.198}{120.157}{(1.2)}%
\TextAt{8}{256.174}{120.157}{implies}%
\TextAt{8}{295.901}{120.157}{that}%
\TextAt{13}{321.269}{120.157}{v}%
\TextAt{27}{326.929}{118.364}{P}%
\TextAt{9}{334.022}{120.157}{\Character{40}}%
\TextAt{13}{338.575}{120.157}{a}%
\TextAt{9}{344.72}{120.157}{\Character{41}}%
\TextAt{9}{352.575}{120.157}{\Character{61}}%
\TextAt{13}{365.018}{120.157}{e}%
\TextAt{27}{370.444}{118.297}{P}%
\TextAt{3}{377.039}{118.297}{\Character{106}}%
\TextAt{27}{379.391}{118.297}{p}%
\TextAt{13}{383.878}{120.157}{v}%
\TextAt{27}{389.538}{118.364}{p}%
\TextAt{9}{394.024}{120.157}{\Character{40}}%
\TextAt{13}{398.577}{120.157}{a}%
\TextAt{9}{404.722}{120.157}{\Character{41}}%
\TextAt{8}{409.274}{120.157}{.}%
\TextAt{8}{417.722}{120.157}{Hence}%
\TextAt{13}{138.758}{84.677}{i}%
\TextAt{21}{142.752}{82.818}{L}%
\TextAt{3}{148.511}{82.818}{\Character{106}}%
\TextAt{21}{150.863}{82.818}{K}%
\TextAt{9}{159.066}{84.677}{\Character{40}\Character{91}}%
\TextAt{13}{166.87}{84.677}{a}%
\TextAt{9}{173.015}{84.677}{\Character{93}\Character{41}}%
\TextAt{9}{184.109}{84.677}{\Character{61}}%
\TextAt{13}{196.565}{84.677}{i}%
\TextAt{21}{200.559}{82.818}{L}%
\TextAt{3}{206.318}{82.818}{\Character{106}}%
\TextAt{21}{208.67}{82.818}{K}%
\TextAt{20}{218.865}{105.121}{\Character{32}}%
\TextAt{20}{228.33}{96.035}{\Character{89}}%
\TextAt{27}{233.973}{70.929}{p}%
\TextAt{14}{245.598}{84.677}{p}%
\TextAt{21}{251.581}{89.614}{v}%
\TextAt{28}{255.705}{88.617}{p}%
\TextAt{17}{259.787}{89.614}{\Character{40}}%
\TextAt{21}{263.081}{89.614}{a}%
\TextAt{17}{267.579}{89.614}{\Character{41}}%
\TextAt{20}{271.37}{105.121}{\Character{33}}%
\TextAt{9}{284.155}{84.677}{\Character{61}}%
\TextAt{20}{296.581}{96.035}{\Character{89}}%
\TextAt{27}{302.225}{70.929}{p}%
\TextAt{20}{313.85}{96.035}{\Character{89}}%
\TextAt{27}{315.019}{70.132}{P}%
\TextAt{3}{321.615}{70.132}{\Character{106}}%
\TextAt{27}{323.967}{70.132}{p}%
\TextAt{14}{331.118}{84.677}{P}%
\TextAt{21}{341.011}{89.614}{e}%
\TextAt{28}{344.94}{87.876}{P}%
\TextAt{23}{350.639}{87.876}{\Character{106}}%
\TextAt{28}{352.908}{87.876}{p}%
\TextAt{21}{356.991}{89.614}{v}%
\TextAt{28}{361.114}{88.617}{p}%
\TextAt{17}{365.197}{89.614}{\Character{40}}%
\TextAt{21}{368.49}{89.614}{a}%
\TextAt{17}{372.988}{89.614}{\Character{41}}%
\TextAt{9}{380.1}{84.677}{\Character{61}}%
\TextAt{20}{392.526}{96.035}{\Character{89}}%
\TextAt{27}{396.866}{70.599}{P}%
\TextAt{14}{409.795}{84.677}{P}%
\TextAt{21}{419.688}{89.614}{v}%
\TextAt{28}{423.811}{88.223}{P}%
\TextAt{17}{430.008}{89.614}{\Character{40}}%
\TextAt{21}{433.302}{89.614}{a}%
\TextAt{17}{437.8}{89.614}{\Character{41}}%
\TextAt{9}{444.912}{84.677}{\Character{61}}%
\TextAt{9}{457.322}{84.677}{\Character{91}}%
\TextAt{13}{460.589}{84.677}{a}%
\TextAt{9}{466.734}{84.677}{\Character{93}}%
\TextAt{13}{469.986}{84.677}{\Character{58}}%
\TextAt{8}{297.211}{43.033}{171}%
\Destination{79.23}{552.595}{nav.Item.64}%
\Destination{79.23}{517.645}{nav.Item.65}%
\Destination{79.23}{500.605}{nav.Item.66}%
\Destination{79.23}{483.555}{nav.Item.67}%
\Destination{79.23}{449.325}{nav.Item.68}%
\Destination{79.23}{395.955}{nav.Item.69}%
\Destination{78.23}{809.815}{nav.page.171}%
\Destination{79.36}{776.155}{nav.print.187}%
\Destination{79.36}{572.555}{nav.number.3.1.6}%
\Hotlink{106.184}{312.755}{25.078}{11.047}{nav.number.3.1.2}%
\Hotlink{247.757}{295.755}{25.078}{11.047}{nav.number.3.1.2}%
\Hotlink{227.698}{117.955}{25.078}{11.047}{nav.number.3.1.2}%
}

% Page 184 (body page 172)
\Page{%
\TextAt{2}{79.226}{796.943}{III.}%
\TextAt{2}{98.317}{796.943}{Riemann–Roch}%
\TextAt{2}{174.375}{796.943}{Theory}%
\Vector{q 0 G 0 g 0.299 w 7.226 18.729 m 460.769 18.729 l S Q}%
\TextAt{8}{79.226}{761.124}{If,}%
\TextAt{8}{94.982}{761.124}{on}%
\TextAt{8}{111.89}{761.124}{the}%
\TextAt{8}{132.696}{761.124}{other}%
\TextAt{8}{163.914}{761.124}{hand,}%
\TextAt{13}{197.244}{761.124}{a}%
\NativeMathAt{207.816}{761.124}{11.955}{\in}%
\TextAt{13}{220.213}{761.124}{L}%
\TextAt{3}{228.178}{765.462}{\Character{3}}%
\TextAt{8}{232.91}{761.124}{,}%
\TextAt{8}{240.871}{761.124}{then}%
\TextAt{8}{268.18}{761.124}{(1.2)}%
\TextAt{8}{296.804}{761.124}{and}%
\TextAt{8}{320.216}{761.124}{chap.}%
\TextAt{8}{354.355}{761.124}{II,}%
\TextAt{8}{370.917}{761.124}{(8.4)}%
\TextAt{8}{399.541}{761.124}{imply}%
\TextAt{8}{433.031}{761.124}{that}%
\TextAt{13}{459.049}{761.124}{v}%
\TextAt{27}{464.709}{759.33}{p}%
\TextAt{9}{469.196}{761.124}{\Character{40}}%
\TextAt{13}{473.748}{761.124}{N}%
\TextAt{21}{483.122}{759.264}{L}%
\TextAt{3}{488.881}{759.264}{\Character{106}}%
\TextAt{21}{491.233}{759.264}{K}%
\TextAt{9}{499.436}{761.124}{\Character{40}}%
\TextAt{13}{503.988}{761.124}{a}%
\TextAt{9}{510.133}{761.124}{\Character{41}\Character{41}}%
\TextAt{9}{523.642}{761.124}{\Character{61}}%
\TextAt{20}{79.226}{753.044}{\Character{80}}%
\TextAt{27}{91.846}{740.591}{P}%
\TextAt{3}{98.441}{740.591}{\Character{106}}%
\TextAt{27}{100.793}{740.591}{p}%
\TextAt{13}{107.273}{744.078}{f}%
\TextAt{27}{113.043}{742.218}{P}%
\TextAt{3}{119.638}{742.218}{\Character{106}}%
\TextAt{27}{121.99}{742.218}{p}%
\TextAt{13}{126.477}{744.078}{v}%
\TextAt{27}{132.137}{742.284}{P}%
\TextAt{9}{139.23}{744.078}{\Character{40}}%
\TextAt{13}{143.782}{744.078}{a}%
\TextAt{9}{149.927}{744.078}{\Character{41}}%
\TextAt{8}{154.48}{744.078}{.}%
\TextAt{8}{162.928}{744.078}{Hence}%
\TextAt{13}{150.115}{704.798}{N}%
\TextAt{21}{159.489}{702.938}{L}%
\TextAt{3}{165.248}{702.938}{\Character{106}}%
\TextAt{21}{167.6}{702.938}{K}%
\TextAt{9}{175.803}{704.798}{\Character{40}\Character{91}}%
\TextAt{13}{183.607}{704.798}{a}%
\TextAt{9}{189.752}{704.798}{\Character{93}\Character{41}}%
\TextAt{9}{200.846}{704.798}{\Character{61}}%
\TextAt{13}{213.302}{704.798}{N}%
\TextAt{21}{222.676}{702.938}{L}%
\TextAt{3}{228.435}{702.938}{\Character{106}}%
\TextAt{21}{230.787}{702.938}{K}%
\TextAt{20}{240.982}{725.241}{\Character{32}}%
\TextAt{20}{250.447}{716.155}{\Character{89}}%
\TextAt{27}{254.787}{690.719}{P}%
\TextAt{14}{267.715}{704.798}{P}%
\TextAt{21}{277.609}{709.734}{v}%
\TextAt{28}{281.732}{708.343}{P}%
\TextAt{17}{287.929}{709.734}{\Character{40}}%
\TextAt{21}{291.222}{709.734}{a}%
\TextAt{17}{295.72}{709.734}{\Character{41}}%
\TextAt{20}{299.512}{725.241}{\Character{33}}%
\TextAt{9}{312.297}{704.798}{\Character{61}}%
\TextAt{20}{324.723}{716.155}{\Character{89}}%
\TextAt{27}{330.366}{691.049}{p}%
\TextAt{14}{341.991}{704.798}{p}%
\TextAt{21}{347.974}{709.734}{v}%
\TextAt{28}{352.098}{708.738}{p}%
\TextAt{17}{356.181}{709.734}{\Character{40}}%
\TextAt{21}{359.474}{709.734}{N}%
\TextAt{22}{366.195}{707.996}{L}%
\TextAt{23}{371.254}{707.996}{\Character{106}}%
\TextAt{22}{373.523}{707.996}{K}%
\TextAt{17}{380.566}{709.734}{\Character{40}}%
\TextAt{21}{383.86}{709.734}{a}%
\TextAt{17}{388.358}{709.734}{\Character{41}\Character{41}}%
\TextAt{9}{398.763}{704.798}{\Character{61}}%
\TextAt{9}{411.173}{704.798}{\Character{91}}%
\TextAt{13}{414.44}{704.798}{N}%
\TextAt{21}{423.814}{702.938}{L}%
\TextAt{3}{429.573}{702.938}{\Character{106}}%
\TextAt{21}{431.925}{702.938}{K}%
\TextAt{9}{440.128}{704.798}{\Character{40}}%
\TextAt{13}{444.68}{704.798}{a}%
\TextAt{9}{450.825}{704.798}{\Character{41}\Character{93}}%
\TextAt{13}{458.629}{704.798}{\Character{58}}%
\TextAt{8}{79.226}{666.172}{(vi)}%
\TextAt{8}{102.349}{666.172}{Let}%
\TextAt{14}{124.002}{666.172}{a}%
\TextAt{21}{129.985}{664.378}{f}%
\TextAt{8}{140.01}{666.172}{be}%
\TextAt{8}{156.618}{666.172}{the}%
\TextAt{8}{177.458}{666.172}{ideal}%
\TextAt{8}{206.104}{666.172}{of}%
\TextAt{13}{220.117}{666.172}{K}%
\TextAt{8}{235.515}{666.172}{which}%
\TextAt{8}{269.685}{666.172}{is}%
\TextAt{8}{282.132}{666.172}{generated}%
\TextAt{8}{336.147}{666.172}{by}%
\TextAt{8}{353.09}{666.172}{all}%
\TextAt{13}{370.032}{666.172}{N}%
\TextAt{21}{379.406}{664.312}{L}%
\TextAt{3}{385.165}{664.312}{\Character{106}}%
\TextAt{21}{387.517}{664.312}{K}%
\TextAt{9}{395.72}{666.172}{\Character{40}}%
\TextAt{13}{400.272}{666.172}{a}%
\TextAt{9}{406.417}{666.172}{\Character{41}}%
\TextAt{8}{410.969}{666.172}{,}%
\TextAt{8}{418.973}{666.172}{with}%
\TextAt{13}{446.316}{666.172}{a}%
\NativeMathAt{456.946}{666.172}{11.955}{\in}%
\TextAt{14}{469.401}{666.172}{A}%
\TextAt{21}{477.98}{664.378}{f}%
\TextAt{8}{483.425}{666.172}{.}%
\TextAt{8}{493.923}{666.172}{If}%
\TextAt{14}{506.298}{666.172}{A}%
\TextAt{21}{514.878}{664.378}{f}%
\TextAt{8}{524.903}{666.172}{is}%
\TextAt{8}{79.226}{649.126}{a}%
\TextAt{8}{89.164}{649.126}{principal}%
\TextAt{8}{138.124}{649.126}{ideal}%
\TextAt{9}{166.27}{649.126}{\Character{40}}%
\TextAt{13}{170.822}{649.126}{a}%
\TextAt{9}{176.967}{649.126}{\Character{41}}%
\TextAt{8}{181.52}{649.126}{,}%
\TextAt{8}{188.897}{649.126}{then}%
\TextAt{14}{215.74}{649.126}{a}%
\TextAt{21}{221.723}{647.332}{f}%
\TextAt{9}{230.8}{649.126}{\Character{61}}%
\TextAt{9}{243.52}{649.126}{\Character{40}}%
\TextAt{13}{248.089}{649.126}{N}%
\TextAt{21}{257.463}{647.266}{L}%
\TextAt{3}{263.222}{647.266}{\Character{106}}%
\TextAt{21}{265.575}{647.266}{K}%
\TextAt{9}{273.777}{649.126}{\Character{40}}%
\TextAt{13}{278.33}{649.126}{a}%
\TextAt{9}{284.475}{649.126}{\Character{41}\Character{41}}%
\TextAt{9}{297.183}{649.126}{\Character{61}}%
\TextAt{13}{309.948}{649.126}{N}%
\TextAt{21}{319.322}{647.266}{L}%
\TextAt{3}{325.081}{647.266}{\Character{106}}%
\TextAt{21}{327.434}{647.266}{K}%
\TextAt{9}{335.636}{649.126}{\Character{40}}%
\TextAt{14}{340.189}{649.126}{A}%
\TextAt{21}{348.768}{647.332}{f}%
\TextAt{9}{354.212}{649.126}{\Character{41}}%
\TextAt{8}{358.765}{649.126}{,}%
\TextAt{8}{366.142}{649.126}{by}%
\TextAt{8}{382.584}{649.126}{(v).}%
\TextAt{8}{406.871}{649.126}{But}%
\TextAt{8}{430.295}{649.126}{the}%
\TextAt{8}{450.634}{649.126}{argument}%
\TextAt{8}{503.18}{649.126}{which}%
\TextAt{8}{79.226}{632.08}{yielded}%
\TextAt{8}{120.609}{632.08}{(v)}%
\TextAt{8}{141.189}{632.08}{applies}%
\TextAt{8}{181.663}{632.08}{equally}%
\TextAt{8}{223.38}{632.08}{well}%
\TextAt{8}{248.504}{632.08}{to}%
\TextAt{8}{264.206}{632.08}{the}%
\TextAt{8}{285.755}{632.08}{localizations}%
\TextAt{16}{353.881}{632.08}{\Character{79}}%
\TextAt{27}{363.399}{630.286}{p}%
\TextAt{16}{373.58}{632.08}{\Character{106}}%
\TextAt{14}{382.594}{632.08}{o}%
\TextAt{27}{388.438}{630.286}{p}%
\TextAt{8}{398.215}{632.08}{of}%
\TextAt{8}{412.938}{632.08}{the}%
\TextAt{8}{434.486}{632.08}{extension}%
\TextAt{16}{487.645}{632.08}{\Character{79}}%
\TextAt{16}{503.187}{632.08}{\Character{106}}%
\TextAt{14}{512.203}{632.08}{o}%
\TextAt{8}{523.337}{632.08}{of}%
\TextAt{8}{79.226}{615.034}{maximal}%
\TextAt{8}{127.104}{615.034}{orders}%
\TextAt{8}{162.356}{615.034}{of}%
\TextAt{13}{175.755}{615.034}{L}%
\TextAt{16}{187.158}{615.034}{\Character{106}}%
\TextAt{13}{193.917}{615.034}{K}%
\TextAt{8}{204.734}{615.034}{.}%
\TextAt{16}{213.389}{615.034}{\Character{79}}%
\TextAt{27}{222.906}{613.24}{p}%
\TextAt{8}{231.359}{615.034}{has}%
\TextAt{8}{252.302}{615.034}{only}%
\TextAt{8}{278.062}{615.034}{a}%
\TextAt{8}{287.886}{615.034}{finite}%
\TextAt{8}{317.867}{615.034}{number}%
\TextAt{8}{360.532}{615.034}{of}%
\TextAt{8}{373.931}{615.034}{prime}%
\TextAt{8}{407.163}{615.034}{ideals,}%
\TextAt{8}{443.079}{615.034}{and}%
\TextAt{8}{465.91}{615.034}{is}%
\TextAt{8}{477.743}{615.034}{therefore}%
\TextAt{8}{526.911}{615.034}{a}%
\TextAt{8}{79.226}{597.988}{principal}%
\TextAt{8}{128.003}{597.988}{ideal}%
\TextAt{8}{155.966}{597.988}{domain}%
\TextAt{8}{197.594}{597.988}{(see}%
\TextAt{8}{221.062}{597.988}{chap.}%
\TextAt{8}{253.254}{597.988}{I,}%
\TextAt{8}{264.623}{597.988}{§3,}%
\TextAt{8}{283.309}{597.988}{exercise}%
\TextAt{8}{326.61}{597.988}{4).}%
\TextAt{8}{345.472}{597.988}{We}%
\TextAt{8}{365.616}{597.988}{thus}%
\TextAt{8}{391.368}{597.988}{get}%
\TextAt{9}{219.886}{569.059}{\Character{40}}%
\TextAt{14}{224.438}{569.059}{a}%
\TextAt{21}{230.422}{567.266}{f}%
\TextAt{9}{235.866}{569.059}{\Character{41}}%
\TextAt{27}{240.419}{567.266}{p}%
\TextAt{9}{248.226}{569.059}{\Character{61}}%
\TextAt{13}{260.652}{569.059}{N}%
\TextAt{21}{270.026}{567.199}{L}%
\TextAt{3}{275.785}{567.199}{\Character{106}}%
\TextAt{21}{278.137}{567.199}{K}%
\TextAt{9}{286.34}{569.059}{\Character{40}\Character{40}}%
\TextAt{14}{295.444}{569.059}{A}%
\TextAt{21}{304.024}{567.266}{f}%
\TextAt{9}{309.468}{569.059}{\Character{41}}%
\TextAt{27}{314.021}{567.266}{p}%
\TextAt{9}{318.508}{569.059}{\Character{41}}%
\TextAt{9}{326.363}{569.059}{\Character{61}}%
\TextAt{13}{338.806}{569.059}{N}%
\TextAt{21}{348.18}{567.199}{L}%
\TextAt{3}{353.939}{567.199}{\Character{106}}%
\TextAt{21}{356.291}{567.199}{K}%
\TextAt{9}{364.494}{569.059}{\Character{40}}%
\TextAt{14}{369.046}{569.059}{A}%
\TextAt{21}{377.625}{567.266}{f}%
\TextAt{9}{383.07}{569.059}{\Character{41}}%
\TextAt{27}{387.623}{567.266}{p}%
\TextAt{8}{79.226}{540.13}{for}%
\TextAt{8}{97.111}{540.13}{all}%
\TextAt{8}{113.37}{540.13}{prime}%
\TextAt{8}{146.534}{540.13}{ideals}%
\TextAt{14}{179.112}{540.13}{p}%
\TextAt{8}{188.992}{540.13}{of}%
\TextAt{14}{202.322}{540.13}{o}%
\TextAt{8}{208.166}{540.13}{,}%
\TextAt{8}{215.316}{540.13}{and}%
\TextAt{8}{238.078}{540.13}{consequently}%
\TextAt{14}{307.406}{540.13}{a}%
\TextAt{21}{313.389}{538.337}{f}%
\TextAt{9}{322.155}{540.131}{\Character{61}}%
\TextAt{13}{334.58}{540.131}{N}%
\TextAt{21}{343.954}{538.271}{L}%
\TextAt{3}{349.713}{538.271}{\Character{106}}%
\TextAt{21}{352.066}{538.271}{K}%
\TextAt{9}{360.268}{540.131}{\Character{40}}%
\TextAt{14}{364.821}{540.131}{A}%
\TextAt{21}{373.4}{538.337}{f}%
\TextAt{9}{378.845}{540.131}{\Character{41}}%
\TextAt{8}{383.397}{540.13}{.}%
\Vector{q 0 G 0 g 0.399 w 452.334 -233.87 m 452.334 -225.8 l S Q}%
\Vector{q 0 G 0 g 0.399 w 452.534 -225.999 m 459.707 -225.999 l S Q}%
\Vector{q 0 G 0 g 0.399 w 452.534 -233.67 m 459.707 -233.67 l S Q}%
\Vector{q 0 G 0 g 0.399 w 459.906 -233.87 m 459.906 -225.8 l S Q}%
\TextAt{8}{97.159}{511.15}{Since}%
\TextAt{8}{127.545}{511.15}{the}%
\TextAt{8}{147.529}{511.15}{homomorphisms}%
\TextAt{13}{234.642}{511.15}{i}%
\TextAt{21}{238.636}{509.29}{L}%
\TextAt{3}{244.395}{509.29}{\Character{106}}%
\TextAt{21}{246.747}{509.29}{K}%
\TextAt{8}{258.676}{511.15}{and}%
\TextAt{13}{281.267}{511.15}{N}%
\TextAt{21}{290.64}{509.29}{L}%
\TextAt{3}{296.399}{509.29}{\Character{106}}%
\TextAt{21}{298.752}{509.29}{K}%
\TextAt{8}{310.68}{511.15}{map}%
\TextAt{8}{336.523}{511.15}{replete}%
\TextAt{8}{374.715}{511.15}{principal}%
\TextAt{8}{423.321}{511.15}{ideals}%
\TextAt{8}{455.727}{511.15}{to}%
\TextAt{8}{469.866}{511.15}{replete}%
\TextAt{8}{508.058}{511.15}{prin-}%
\TextAt{8}{79.226}{494.104}{cipal}%
\TextAt{8}{107.7}{494.104}{ideals,}%
\TextAt{8}{144.169}{494.104}{they}%
\TextAt{8}{171.017}{494.104}{induce}%
\TextAt{8}{208.59}{494.104}{homomorphisms}%
\TextAt{8}{296.385}{494.104}{of}%
\TextAt{8}{310.227}{494.104}{the}%
\TextAt{8}{330.894}{494.104}{replete}%
\TextAt{8}{369.769}{494.104}{Picard}%
\TextAt{8}{407.509}{494.104}{groups}%
\TextAt{8}{445.81}{494.104}{of}%
\TextAt{13}{459.652}{494.104}{K}%
\TextAt{8}{474.878}{494.104}{and}%
\TextAt{13}{498.152}{494.104}{L}%
\TextAt{8}{506.116}{494.104}{,}%
\TextAt{8}{513.904}{494.104}{and}%
\TextAt{8}{79.226}{477.058}{we}%
\TextAt{8}{96.454}{477.058}{obtain}%
\TextAt{8}{132.881}{477.058}{the}%
\TextAt{12}{79.226}{451.067}{Proposition}%
\TextAt{12}{153.383}{451.067}{(1.7).}%
\TextAt{8}{189.955}{451.067}{For}%
\TextAt{8}{212.181}{451.067}{every}%
\TextAt{8}{244.33}{451.067}{finite}%
\TextAt{8}{275.499}{451.067}{extension}%
\TextAt{13}{328.521}{451.067}{L}%
\TextAt{16}{341.948}{451.067}{\Character{106}}%
\TextAt{13}{350.731}{451.067}{K}%
\TextAt{8}{361.549}{451.067}{,}%
\TextAt{8}{370.269}{451.067}{the}%
\TextAt{8}{391.682}{451.067}{following}%
\TextAt{8}{442.374}{451.067}{two}%
\TextAt{8}{465.747}{451.067}{diagrams}%
\TextAt{8}{517.156}{451.067}{are}%
\TextAt{8}{79.226}{434.021}{commutative:}%
\TextAt{8}{251.793}{419.013}{Pic}%
\TextAt{9}{268.207}{419.013}{\Character{40}}%
\NativeDiagram{%
\path[draw=black,line width=0.478bp] (200.76,-347.63) -- (206.604,-347.63);%
}%
\TextAt{14}{272.76}{419.013}{o}%
\TextAt{21}{278.604}{417.22}{L}%
\TextAt{9}{284.861}{419.013}{\Character{41}}%
\TextAt{18}{343.731}{419.013}{R}%
\TextAt{3}{352.365}{423.352}{\Character{3}}%
\TextAt{17}{352.365}{416.058}{\Character{43}}%
\TextAt{8}{250.821}{365.223}{Pic}%
\TextAt{9}{267.235}{365.223}{\Character{40}}%
\NativeDiagram{%
\path[draw=black,line width=0.478bp] (199.788,-401.421) -- (205.632,-401.421);%
}%
\TextAt{14}{271.788}{365.223}{o}%
\TextAt{21}{277.632}{363.429}{K}%
\TextAt{9}{285.835}{365.223}{\Character{41}}%
\TextAt{18}{343.731}{365.223}{R}%
\TextAt{3}{352.365}{369.561}{\Character{3}}%
\TextAt{17}{352.365}{362.267}{\Character{43}}%
\NativeDiagram{%
\path[draw=black,line width=0.478bp] (222.767,-351.997) -- (265.902,-351.997);%
\path[draw=black,line width=0.478bp,line cap=round,line join=round] (263.654,-349.128) .. controls (264.109,-350.849) and (265.121,-351.662) .. (266.141,-351.997) .. controls (265.121,-352.332) and (264.109,-353.144) .. (263.654,-354.866);%
}%
\TextAt{27}{313.257}{424.819}{N}%
\NativeDiagram{%
\path[draw=black,line width=0.478bp] (206.333,-362.554) -- (206.333,-394.753);%
\path[draw=black,line width=0.478bp,line cap=round,line join=round] (209.202,-392.506) .. controls (207.481,-392.96) and (206.668,-393.972) .. (206.333,-394.992) .. controls (205.998,-393.972) and (205.186,-392.96) .. (203.464,-392.506);%
}%
\TextAt{21}{281.149}{394.001}{N}%
\TextAt{22}{287.871}{392.263}{L}%
\TextAt{23}{292.93}{392.263}{\Character{106}}%
\TextAt{22}{295.199}{392.263}{K}%
\NativeDiagram{%
\path[draw=black,line width=0.478bp] (287.321,-363.351) -- (287.321,-395.484);%
\path[draw=black,line width=0.478bp,line cap=round,line join=round] (290.19,-393.236) .. controls (288.468,-393.691) and (287.655,-394.703) .. (287.321,-395.723) .. controls (286.986,-394.703) and (286.173,-393.691) .. (284.452,-393.236);%
}%
\TextAt{7}{362.136}{391.618}{id}%
\NativeDiagram{%
\path[draw=black,line width=0.478bp] (223.739,-405.788) -- (265.902,-405.788);%
\path[draw=black,line width=0.478bp,line cap=round,line join=round] (263.654,-402.919) .. controls (264.109,-404.641) and (265.121,-405.453) .. (266.141,-405.788) .. controls (265.121,-406.123) and (264.109,-406.936) .. (263.654,-408.657);%
}%
\TextAt{27}{313.744}{371.027}{N}%
\NativeDiagram{%
\path[draw=black,line width=0.478bp] (190.875,-395.231) -- (190.875,-363.032);%
\path[draw=black,line width=0.478bp,line cap=round,line join=round] (188.006,-365.279) .. controls (189.727,-364.825) and (190.54,-363.813) .. (190.875,-362.793) .. controls (191.21,-363.813) and (192.023,-364.825) .. (193.744,-365.279);%
}%
\TextAt{21}{242.805}{394.089}{i}%
\TextAt{22}{245.688}{392.351}{L}%
\TextAt{23}{250.747}{392.351}{\Character{106}}%
\TextAt{22}{253.017}{392.351}{K}%
\NativeDiagram{%
\path[draw=black,line width=0.478bp] (271.863,-395.962) -- (271.863,-363.829);%
\path[draw=black,line width=0.478bp,line cap=round,line join=round] (268.993,-366.076) .. controls (270.715,-365.622) and (271.528,-364.61) .. (271.863,-363.59) .. controls (272.197,-364.61) and (273.01,-365.622) .. (274.732,-366.076);%
}%
\TextAt{17}{320.526}{392.351}{\Character{91}}%
\TextAt{21}{322.879}{392.351}{L}%
\TextAt{17}{328.638}{392.351}{\Character{58}}%
\TextAt{21}{330.99}{392.351}{K}%
\TextAt{17}{338.695}{392.351}{\Character{93}}%
\TextAt{8}{97.159}{332.254}{Let}%
\TextAt{8}{118.495}{332.254}{us}%
\TextAt{8}{133.877}{332.254}{now}%
\TextAt{8}{158.632}{332.254}{translate}%
\TextAt{8}{207.848}{332.254}{the}%
\TextAt{8}{228.371}{332.254}{notions}%
\TextAt{8}{269.78}{332.254}{we}%
\TextAt{8}{287.374}{332.254}{have}%
\TextAt{8}{314.735}{332.254}{introduced}%
\TextAt{8}{373.634}{332.254}{into}%
\TextAt{8}{397.744}{332.254}{the}%
\TextAt{8}{418.267}{332.254}{function-theoretic}%
\TextAt{8}{513.259}{332.254}{lan-}%
\TextAt{8}{79.226}{315.208}{guage}%
\TextAt{8}{113.011}{315.208}{of}%
\TextAt{8}{126.949}{315.208}{divisors.}%
\TextAt{8}{176.055}{315.208}{In}%
\TextAt{8}{191.285}{315.208}{chap.}%
\TextAt{8}{225.302}{315.208}{I,}%
\TextAt{8}{237.281}{315.208}{§12,}%
\TextAt{8}{262.586}{315.208}{we}%
\TextAt{8}{280.422}{315.208}{defined}%
\TextAt{8}{321.344}{315.208}{the}%
\TextAt{8}{342.109}{315.208}{divisor}%
\TextAt{8}{380.831}{315.208}{group}%
\TextAt{8}{414.615}{315.208}{Div}%
\TextAt{9}{432.99}{315.208}{\Character{40}}%
\TextAt{14}{437.543}{315.208}{o}%
\TextAt{9}{443.387}{315.208}{\Character{41}}%
\TextAt{8}{452.445}{315.208}{to}%
\TextAt{8}{467.364}{315.208}{consist}%
\TextAt{8}{506.469}{315.208}{of}%
\TextAt{8}{520.408}{315.208}{all}%
\TextAt{8}{79.226}{298.163}{formal}%
\TextAt{8}{115.976}{298.163}{sums}%
\TextAt{13}{273.719}{281.117}{D}%
\TextAt{9}{287.084}{281.117}{\Character{61}}%
\TextAt{20}{299.51}{292.474}{\Character{88}}%
\TextAt{27}{300.809}{266.642}{p}%
\TextAt{25}{304.797}{266.642}{\Character{45}}%
\NativeMathAt{307.011}{266.642}{7.97}{\infty}%
\TextAt{13}{318.771}{281.117}{\Character{23}}%
\TextAt{27}{324.555}{279.323}{p}%
\TextAt{14}{329.042}{281.117}{p}%
\TextAt{13}{335.025}{281.117}{\Character{59}}%
\TextAt{8}{79.226}{247.035}{where}%
\TextAt{13}{113.849}{247.035}{\Character{23}}%
\TextAt{27}{119.633}{245.242}{p}%
\NativeMathAt{128.827}{247.035}{11.955}{\in}%
\TextAt{18}{141.505}{247.035}{Z}%
\TextAt{8}{149.475}{247.035}{,}%
\TextAt{8}{157.641}{247.035}{and}%
\TextAt{13}{181.217}{247.035}{\Character{23}}%
\TextAt{27}{187.001}{245.242}{p}%
\TextAt{9}{196.195}{247.035}{\Character{61}}%
\TextAt{9}{209.991}{247.035}{0}%
\TextAt{8}{220.571}{247.035}{for}%
\TextAt{8}{239.27}{247.035}{almost}%
\TextAt{8}{277.874}{247.035}{all}%
\TextAt{14}{294.947}{247.035}{p}%
\TextAt{8}{300.93}{247.035}{.}%
\TextAt{8}{311.819}{247.035}{Contained}%
\TextAt{8}{368.894}{247.035}{in}%
\TextAt{8}{383.36}{247.035}{this}%
\TextAt{8}{406.996}{247.035}{group}%
\TextAt{8}{440.985}{247.035}{is}%
\TextAt{8}{453.563}{247.035}{the}%
\TextAt{8}{474.533}{247.035}{group}%
\TextAt{16}{508.522}{247.035}{\Character{80}}%
\TextAt{9}{517.821}{247.035}{\Character{40}}%
\TextAt{14}{522.373}{247.035}{o}%
\TextAt{9}{528.217}{247.035}{\Character{41}}%
\TextAt{8}{79.226}{229.989}{of}%
\TextAt{8}{93.348}{229.989}{principal}%
\TextAt{8}{142.917}{229.989}{divisors}%
\TextAt{8}{186.436}{229.989}{div}%
\TextAt{9}{202.372}{229.989}{\Character{40}}%
\TextAt{13}{206.925}{229.989}{f}%
\TextAt{9}{213.971}{229.989}{\Character{41}}%
\TextAt{9}{223.177}{229.989}{\Character{61}}%
\TextAt{20}{236.968}{238.956}{\Character{80}}%
\TextAt{27}{249.588}{226.502}{p}%
\TextAt{25}{253.576}{226.502}{\Character{45}}%
\NativeMathAt{255.79}{226.502}{7.97}{\infty}%
\TextAt{13}{266.749}{229.989}{v}%
\TextAt{27}{272.409}{228.196}{p}%
\TextAt{9}{276.895}{229.989}{\Character{40}}%
\TextAt{13}{281.448}{229.989}{f}%
\TextAt{9}{288.494}{229.989}{\Character{41}}%
\TextAt{14}{293.047}{229.989}{p}%
\TextAt{8}{299.03}{229.989}{,}%
\TextAt{8}{307.169}{229.989}{which}%
\TextAt{8}{341.447}{229.989}{allowed}%
\TextAt{8}{383.866}{229.989}{us}%
\TextAt{8}{399.674}{229.989}{to}%
\TextAt{8}{414.776}{229.989}{define}%
\TextAt{8}{449.377}{229.989}{the}%
\TextAt{8}{470.325}{229.989}{divisor}%
\TextAt{8}{509.23}{229.989}{class}%
\TextAt{8}{79.226}{212.943}{group}%
\TextAt{13}{245.018}{195.897}{C}%
\TextAt{13}{254.236}{195.897}{H}%
\TextAt{17}{264.874}{200.834}{1}%
\TextAt{9}{269.606}{195.897}{\Character{40}}%
\TextAt{14}{274.159}{195.897}{o}%
\TextAt{9}{280.003}{195.897}{\Character{41}}%
\TextAt{9}{287.858}{195.897}{\Character{61}}%
\TextAt{8}{300.302}{195.897}{Div}%
\TextAt{9}{318.677}{195.897}{\Character{40}}%
\TextAt{14}{323.229}{195.897}{o}%
\TextAt{9}{329.073}{195.897}{\Character{41}\Character{47}}%
\TextAt{16}{339.479}{195.897}{\Character{80}}%
\TextAt{9}{348.777}{195.897}{\Character{40}}%
\TextAt{14}{353.329}{195.897}{o}%
\TextAt{9}{359.174}{195.897}{\Character{41}}%
\TextAt{13}{363.726}{195.897}{\Character{58}}%
\TextAt{8}{79.226}{172.407}{It}%
\TextAt{8}{91.485}{172.407}{follows}%
\TextAt{8}{129.507}{172.407}{from}%
\TextAt{8}{156.734}{172.407}{the}%
\TextAt{8}{176.477}{172.407}{main}%
\TextAt{8}{205.329}{172.407}{theorem}%
\TextAt{8}{250.441}{172.407}{of}%
\TextAt{8}{263.357}{172.407}{ideal}%
\TextAt{8}{290.907}{172.407}{theory,}%
\TextAt{8}{329.598}{172.407}{chap.}%
\TextAt{8}{361.651}{172.407}{I,}%
\TextAt{8}{372.607}{172.407}{(3.9),}%
\TextAt{8}{403.503}{172.407}{that}%
\TextAt{8}{428.458}{172.407}{this}%
\TextAt{8}{450.867}{172.407}{group}%
\TextAt{8}{483.629}{172.407}{is}%
\TextAt{8}{494.979}{172.407}{isomor-}%
\TextAt{8}{79.226}{155.361}{phic}%
\TextAt{8}{104.749}{155.361}{to}%
\TextAt{8}{119.226}{155.361}{the}%
\TextAt{8}{139.549}{155.361}{ideal}%
\TextAt{8}{167.678}{155.361}{class}%
\TextAt{8}{195.281}{155.361}{group}%
\TextAt{8}{228.623}{155.361}{Cl}%
\TextAt{21}{240.327}{153.568}{K}%
\TextAt{8}{248.53}{155.361}{,}%
\TextAt{8}{255.887}{155.361}{which}%
\TextAt{8}{289.54}{155.361}{is}%
\TextAt{8}{301.47}{155.361}{a}%
\TextAt{8}{311.392}{155.361}{finite}%
\TextAt{8}{341.47}{155.361}{group}%
\TextAt{8}{374.812}{155.361}{(see}%
\TextAt{8}{398.446}{155.361}{chap.}%
\TextAt{8}{431.136}{155.361}{I,}%
\TextAt{8}{442.672}{155.361}{(12.14)).}%
\TextAt{8}{491.968}{155.361}{We}%
\TextAt{8}{512.278}{155.361}{now}%
\TextAt{8}{79.226}{138.315}{extend}%
\TextAt{8}{117.267}{138.315}{these}%
\TextAt{8}{147.239}{138.315}{concepts}%
\TextAt{8}{194.773}{138.315}{as}%
\TextAt{8}{209.143}{138.315}{follows.}%
\TextAt{12}{79.226}{112.323}{Definition}%
\TextAt{12}{142.015}{112.323}{(1.8).}%
\TextAt{8}{177.33}{112.323}{A}%
\TextAt{12}{190.003}{112.323}{replete}%
\TextAt{12}{234.835}{112.323}{divisor}%
\TextAt{8}{278.328}{112.323}{(or}%
\TextAt{12}{297.193}{112.323}{Arakelov}%
\TextAt{12}{353.944}{112.323}{divisor}%
\TextAt{8}{393.539}{112.323}{)}%
\TextAt{8}{401.992}{112.323}{of}%
\TextAt{13}{415.322}{112.323}{K}%
\TextAt{8}{430.037}{112.323}{is}%
\TextAt{8}{441.801}{112.323}{a}%
\TextAt{8}{451.556}{112.323}{formal}%
\TextAt{8}{488.306}{112.323}{sum}%
\TextAt{13}{273.719}{83.395}{D}%
\TextAt{9}{287.084}{83.395}{\Character{61}}%
\TextAt{20}{299.51}{94.752}{\Character{88}}%
\TextAt{27}{306.15}{69.646}{p}%
\TextAt{13}{318.771}{83.395}{\Character{23}}%
\TextAt{27}{324.555}{81.601}{p}%
\TextAt{14}{329.042}{83.395}{p}%
\TextAt{13}{335.025}{83.395}{\Character{59}}%
\TextAt{8}{297.211}{43.033}{172}%
\Destination{78.23}{809.815}{nav.page.172}%
\Destination{79.36}{776.155}{nav.print.188}%
\Destination{79.36}{262.055}{nav.print.189}%
\Destination{79.36}{465.055}{nav.number.3.1.7}%
\Destination{79.36}{126.355}{nav.number.3.1.8}%
\Hotlink{353.855}{758.955}{9.763}{11.047}{nav.chapter.2}%
\Hotlink{326.11}{595.755}{6.858}{11.047}{nav.ex.1.3.4}%
\Hotlink{252.754}{595.755}{5.22}{11.047}{nav.chapter.1}%
\Hotlink{269.802}{595.755}{6.858}{11.047}{nav.section.1.3}%
\Hotlink{224.802}{313.055}{5.22}{11.047}{nav.chapter.1}%
\Hotlink{242.46}{313.055}{12.716}{11.047}{nav.section.1.12}%
\Hotlink{361.151}{170.255}{5.22}{11.047}{nav.chapter.1}%
\Hotlink{430.636}{153.155}{5.22}{11.047}{nav.chapter.1}%
\Hotlink{267.68}{758.955}{25.078}{11.047}{nav.number.3.1.2}%
\Hotlink{370.417}{758.955}{25.078}{11.047}{nav.number.2.8.4}%
\Hotlink{372.107}{170.255}{25.078}{11.047}{nav.number.1.3.9}%
\Hotlink{442.172}{153.155}{36.794}{11.047}{nav.number.1.12.14}%
}

% Page 185 (body page 173)
\Page{%
\TextAt{2}{79.226}{796.943}{III.}%
\TextAt{2}{98.317}{796.943}{Riemann–Roch}%
\TextAt{2}{174.375}{796.943}{Theory}%
\Vector{q 0 G 0 g 0.299 w 7.226 18.729 m 460.769 18.729 l S Q}%
\TextAt{8}{79.226}{761.124}{where}%
\TextAt{13}{113.035}{761.124}{\Character{23}}%
\TextAt{27}{118.819}{759.33}{p}%
\NativeMathAt{126.627}{761.124}{11.955}{\in}%
\TextAt{18}{137.918}{761.124}{Z}%
\TextAt{8}{149.785}{761.124}{for}%
\TextAt{14}{167.67}{761.124}{p}%
\TextAt{18}{176.974}{761.124}{\Character{45}}%
\NativeMathAt{183.616}{761.124}{11.955}{\infty}%
\TextAt{8}{195.571}{761.124}{,}%
\TextAt{13}{202.72}{761.124}{\Character{23}}%
\TextAt{27}{208.504}{759.33}{p}%
\NativeMathAt{216.312}{761.124}{11.955}{\in}%
\TextAt{18}{227.603}{761.124}{R}%
\TextAt{8}{240.135}{761.124}{for}%
\TextAt{14}{258.019}{761.124}{p}%
\TextAt{16}{267.323}{761.124}{\Character{106}}%
\NativeMathAt{273.946}{761.124}{11.955}{\infty}%
\TextAt{8}{285.92}{761.124}{,}%
\TextAt{8}{293.069}{761.124}{and}%
\TextAt{13}{315.832}{761.124}{\Character{23}}%
\TextAt{27}{321.616}{759.33}{p}%
\TextAt{9}{329.424}{761.124}{\Character{61}}%
\TextAt{9}{341.834}{761.124}{0}%
\TextAt{8}{351.599}{761.124}{for}%
\TextAt{8}{369.484}{761.124}{almost}%
\TextAt{8}{407.275}{761.124}{all}%
\TextAt{14}{423.534}{761.124}{p}%
\TextAt{8}{429.517}{761.124}{.}%
\TextAt{8}{97.159}{735.111}{The}%
\TextAt{8}{121.15}{735.111}{replete}%
\TextAt{8}{159.451}{735.111}{divisors}%
\TextAt{8}{202.116}{735.111}{form}%
\TextAt{8}{229.693}{735.111}{a}%
\TextAt{8}{239.386}{735.111}{group,}%
\TextAt{8}{275.763}{735.111}{which}%
\TextAt{8}{309.187}{735.111}{is}%
\TextAt{8}{320.888}{735.111}{denoted}%
\TextAt{8}{365.047}{735.111}{by}%
\TextAt{8}{381.243}{735.111}{Div}%
\TextAt{9}{399.618}{735.111}{\Character{40}}%
\Vector{q 0 G 0 g 0.478 w 332.171 -31.532 m 338.015 -31.532 l S Q}%
\TextAt{14}{404.171}{735.111}{o}%
\TextAt{9}{410.015}{735.111}{\Character{41}}%
\TextAt{8}{414.567}{735.111}{.}%
\TextAt{8}{422.994}{735.111}{It}%
\TextAt{8}{435.604}{735.111}{admits}%
\TextAt{8}{473.977}{735.111}{a}%
\TextAt{8}{483.67}{735.111}{decompo-}%
\TextAt{8}{79.226}{718.065}{sition}%
\TextAt{8}{238.434}{701.019}{Div}%
\TextAt{9}{256.809}{701.019}{\Character{40}}%
\Vector{q 0 G 0 g 0.478 w 189.362 -65.624 m 195.206 -65.624 l S Q}%
\TextAt{14}{261.362}{701.019}{o}%
\TextAt{9}{267.206}{701.019}{\Character{41}}%
\NativeMathAt{275.079}{701.019}{11.955}{\cong}%
\TextAt{8}{287.698}{701.019}{Div}%
\TextAt{9}{306.073}{701.019}{\Character{40}}%
\TextAt{14}{310.626}{701.019}{o}%
\TextAt{9}{316.47}{701.019}{\Character{41}}%
\TextAt{16}{323.679}{701.019}{\Character{2}}%
\TextAt{20}{335.634}{712.377}{\Character{77}}%
\TextAt{27}{337.262}{686.474}{p}%
\TextAt{3}{341.251}{686.474}{\Character{106}}%
\NativeMathAt{343.602}{686.474}{7.97}{\infty}%
\TextAt{18}{355.692}{701.019}{R}%
\TextAt{14}{364.327}{701.019}{p}%
\TextAt{13}{370.31}{701.019}{\Character{58}}%
\TextAt{8}{79.226}{667.397}{On}%
\TextAt{8}{99.778}{667.397}{the}%
\TextAt{8}{120.974}{667.397}{right-hand}%
\TextAt{8}{179.579}{667.397}{side,}%
\TextAt{8}{207.6}{667.397}{the}%
\TextAt{8}{228.797}{667.397}{second}%
\TextAt{8}{267.616}{667.397}{factor}%
\TextAt{8}{302.154}{667.397}{is}%
\TextAt{8}{314.959}{667.397}{endowed}%
\TextAt{8}{363.473}{667.397}{with}%
\TextAt{8}{391.174}{667.397}{the}%
\TextAt{8}{412.371}{667.397}{canonical}%
\TextAt{8}{464.795}{667.397}{topology,}%
\TextAt{8}{516.51}{667.397}{the}%
\TextAt{8}{79.226}{650.351}{first}%
\TextAt{8}{102.941}{650.351}{one}%
\TextAt{8}{123.989}{650.351}{with}%
\TextAt{8}{150.238}{650.351}{the}%
\TextAt{8}{169.983}{650.351}{discrete}%
\TextAt{8}{212.551}{650.351}{topology.}%
\TextAt{8}{264.128}{650.351}{On}%
\TextAt{8}{283.227}{650.351}{the}%
\TextAt{8}{302.973}{650.351}{product}%
\TextAt{8}{346.461}{650.351}{we}%
\TextAt{8}{363.277}{650.351}{have}%
\TextAt{8}{389.861}{650.351}{the}%
\TextAt{8}{409.606}{650.351}{product}%
\TextAt{8}{453.094}{650.351}{topology,}%
\TextAt{8}{503.18}{650.351}{which}%
\TextAt{8}{79.226}{633.305}{makes}%
\TextAt{8}{114.41}{633.305}{Div}%
\TextAt{9}{132.785}{633.305}{\Character{40}}%
\Vector{q 0 G 0 g 0.478 w 65.338 -133.338 m 71.182 -133.338 l S Q}%
\TextAt{14}{137.338}{633.305}{o}%
\TextAt{9}{143.182}{633.305}{\Character{41}}%
\TextAt{8}{151.632}{633.305}{into}%
\TextAt{8}{175.375}{633.305}{a}%
\TextAt{8}{185.13}{633.305}{locally}%
\TextAt{8}{222.203}{633.305}{compact}%
\TextAt{8}{269.031}{633.305}{topological}%
\TextAt{8}{328.556}{633.305}{group.}%
\TextAt{8}{97.159}{613.271}{We}%
\TextAt{8}{117.303}{613.271}{now}%
\TextAt{8}{141.692}{613.271}{study}%
\TextAt{8}{173.947}{613.271}{the}%
\TextAt{8}{194.103}{613.271}{canonical}%
\TextAt{8}{245.487}{613.271}{homomorphism}%
\TextAt{8}{189.804}{584.27}{div}%
\TextAt{9}{209.061}{584.27}{\Character{58}}%
\TextAt{13}{215.634}{584.27}{K}%
\TextAt{3}{226.451}{589.206}{\Character{3}}%
\NativeArrowAt{234.504}{584.27}{11.955}{19.26}{\longrightarrow}%
\TextAt{8}{257.086}{584.27}{Div}%
\TextAt{9}{275.462}{584.27}{\Character{40}}%
\Vector{q 0 G 0 g 0.478 w 208.014 -182.374 m 213.858 -182.374 l S Q}%
\TextAt{14}{280.014}{584.27}{o}%
\TextAt{9}{285.858}{584.27}{\Character{41}}%
\TextAt{13}{290.41}{584.27}{\Character{59}}%
\TextAt{8}{319.565}{584.27}{div}%
\TextAt{9}{335.501}{584.27}{\Character{40}}%
\TextAt{13}{340.053}{584.27}{f}%
\TextAt{9}{347.1}{584.27}{\Character{41}}%
\TextAt{9}{354.955}{584.27}{\Character{61}}%
\TextAt{20}{367.399}{595.627}{\Character{88}}%
\TextAt{27}{374.038}{570.521}{p}%
\TextAt{13}{386.66}{584.27}{v}%
\TextAt{27}{392.319}{582.476}{p}%
\TextAt{9}{396.806}{584.27}{\Character{40}}%
\TextAt{13}{401.358}{584.27}{f}%
\TextAt{9}{408.405}{584.27}{\Character{41}}%
\TextAt{14}{412.957}{584.27}{p}%
\TextAt{13}{418.94}{584.27}{\Character{58}}%
\TextAt{8}{79.226}{545.901}{The}%
\TextAt{8}{103.28}{545.901}{elements}%
\TextAt{8}{151.136}{545.901}{of}%
\TextAt{8}{164.467}{545.901}{the}%
\TextAt{8}{184.623}{545.901}{form}%
\TextAt{8}{212.263}{545.901}{div}%
\TextAt{9}{228.2}{545.901}{\Character{40}}%
\TextAt{13}{232.752}{545.901}{f}%
\TextAt{9}{239.798}{545.901}{\Character{41}}%
\TextAt{8}{248.248}{545.901}{are}%
\TextAt{8}{267.759}{545.901}{called}%
\TextAt{12}{300.922}{545.901}{replete}%
\TextAt{12}{345.754}{545.901}{principal}%
\TextAt{12}{401.896}{545.901}{divisors}%
\TextAt{8}{446.799}{545.901}{.}%
\TextAt{12}{79.226}{522.877}{Remark.}%
\TextAt{8}{133.383}{522.877}{The}%
\TextAt{8}{157.437}{522.877}{composite}%
\TextAt{8}{212.455}{522.877}{of}%
\TextAt{8}{225.785}{522.877}{the}%
\TextAt{8}{245.941}{522.877}{mapping}%
\TextAt{8}{294.072}{522.877}{div}%
\TextAt{9}{313.33}{522.877}{\Character{58}}%
\TextAt{13}{319.902}{522.877}{K}%
\TextAt{3}{330.72}{527.216}{\Character{3}}%
\NativeMathAt{338.773}{522.877}{11.955}{\rightarrow}%
\TextAt{8}{354.049}{522.877}{Div}%
\TextAt{9}{372.424}{522.877}{\Character{40}}%
\Vector{q 0 G 0 g 0.478 w 304.976 -243.766 m 310.821 -243.766 l S Q}%
\TextAt{14}{376.976}{522.877}{o}%
\TextAt{9}{382.821}{522.877}{\Character{41}}%
\TextAt{8}{391.27}{522.877}{with}%
\TextAt{8}{417.93}{522.877}{the}%
\TextAt{8}{438.087}{522.877}{mapping}%
\TextAt{8}{194.051}{493.876}{Div}%
\TextAt{9}{212.426}{493.876}{\Character{40}}%
\Vector{q 0 G 0 g 0.478 w 144.979 -272.767 m 150.823 -272.767 l S Q}%
\TextAt{14}{216.979}{493.876}{o}%
\TextAt{9}{222.823}{493.876}{\Character{41}}%
\NativeArrowAt{230.696}{493.876}{11.955}{19.26}{\longrightarrow}%
\TextAt{20}{253.278}{505.234}{\Character{89}}%
\TextAt{27}{253.511}{479.331}{p}%
\TextAt{3}{257.5}{479.331}{\Character{106}}%
\NativeMathAt{259.851}{479.331}{7.97}{\infty}%
\TextAt{18}{270.547}{493.876}{R}%
\TextAt{13}{279.181}{493.876}{\Character{59}}%
\TextAt{20}{308.335}{505.234}{\Character{88}}%
\TextAt{27}{314.975}{480.128}{p}%
\TextAt{13}{327.597}{493.876}{\Character{23}}%
\TextAt{27}{333.38}{492.083}{p}%
\TextAt{14}{337.867}{493.876}{p}%
\NativeArrowAt{347.171}{493.876}{11.955}{19.26}{\longmapsto}%
\TextAt{9}{369.753}{493.876}{\Character{40}}%
\TextAt{13}{374.305}{493.876}{\Character{23}}%
\TextAt{27}{380.089}{492.083}{p}%
\TextAt{13}{384.576}{493.876}{f}%
\TextAt{27}{390.346}{492.083}{p}%
\TextAt{9}{394.833}{493.876}{\Character{41}}%
\TextAt{27}{399.385}{492.017}{p}%
\TextAt{3}{403.374}{492.016}{\Character{106}}%
\NativeMathAt{405.725}{492.016}{7.97}{\infty}%
\TextAt{13}{414.693}{493.876}{\Character{59}}%
\TextAt{8}{79.226}{454.894}{is}%
\TextAt{8}{90.99}{454.894}{equal,}%
\TextAt{8}{125.134}{454.894}{up}%
\TextAt{8}{142.039}{454.894}{to}%
\TextAt{8}{156.349}{454.894}{sign,}%
\TextAt{8}{183.726}{454.894}{to}%
\TextAt{8}{198.037}{454.894}{the}%
\TextAt{8}{218.193}{454.894}{logarithm}%
\TextAt{8}{271.537}{454.894}{map}%
\TextAt{13}{186.293}{425.893}{\Character{21}}%
\TextAt{9}{196.442}{425.893}{\Character{58}}%
\TextAt{13}{203.014}{425.893}{K}%
\TextAt{3}{213.832}{430.829}{\Character{3}}%
\NativeArrowAt{221.885}{425.893}{11.955}{19.26}{\longrightarrow}%
\TextAt{20}{244.467}{437.251}{\Character{89}}%
\TextAt{27}{244.701}{411.348}{p}%
\TextAt{3}{248.689}{411.348}{\Character{106}}%
\NativeMathAt{251.04}{411.348}{7.97}{\infty}%
\TextAt{18}{261.736}{425.893}{R}%
\TextAt{13}{270.37}{425.893}{\Character{59}}%
\TextAt{13}{299.505}{425.893}{\Character{21}}%
\TextAt{9}{306.353}{425.893}{\Character{40}}%
\TextAt{13}{310.905}{425.893}{f}%
\TextAt{9}{317.952}{425.893}{\Character{41}}%
\TextAt{9}{325.807}{425.893}{\Character{61}}%
\TextAt{9}{338.228}{425.893}{\Character{40}}%
\TextAt{13}{342.803}{425.893}{\Character{58}}%
\TextAt{13}{348.027}{425.893}{\Character{58}}%
\TextAt{13}{353.264}{425.893}{\Character{58}}%
\TextAt{13}{358.488}{425.893}{\Character{59}}%
\TextAt{8}{363.779}{425.893}{log}%
\TextAt{16}{380.74}{425.893}{\Character{106}}%
\TextAt{13}{384.061}{425.893}{f}%
\TextAt{16}{391.107}{425.893}{\Character{106}}%
\TextAt{27}{394.428}{424.1}{p}%
\TextAt{13}{398.915}{425.893}{\Character{59}}%
\TextAt{13}{404.139}{425.893}{\Character{58}}%
\TextAt{13}{409.376}{425.893}{\Character{58}}%
\TextAt{13}{414.6}{425.893}{\Character{58}}%
\TextAt{9}{417.899}{425.893}{\Character{41}}%
\TextAt{13}{422.451}{425.893}{\Character{59}}%
\TextAt{8}{79.226}{386.242}{of}%
\TextAt{8}{94.173}{386.242}{Minkowski}%
\TextAt{8}{154.06}{386.242}{theory}%
\TextAt{8}{192.427}{386.242}{(see}%
\TextAt{8}{217.512}{386.242}{chap.}%
\TextAt{8}{254.554}{386.242}{I,}%
\TextAt{8}{267.541}{386.242}{§7,}%
\TextAt{8}{288.248}{386.242}{p.39,}%
\TextAt{8}{318.89}{386.242}{and}%
\TextAt{8}{343.269}{386.242}{chap.}%
\TextAt{8}{380.312}{386.242}{III,}%
\TextAt{8}{402.384}{386.242}{§3,}%
\TextAt{8}{423.091}{386.242}{p.211).}%
\TextAt{8}{468.275}{386.242}{By}%
\TextAt{8}{488.255}{386.242}{chap.}%
\TextAt{8}{525.297}{386.242}{I,}%
\TextAt{8}{79.226}{369.196}{(7.3),}%
\TextAt{8}{111.706}{369.196}{it}%
\TextAt{8}{124.412}{369.196}{maps}%
\TextAt{8}{156.043}{369.196}{the}%
\TextAt{8}{177.201}{369.196}{unit}%
\TextAt{8}{202.915}{369.196}{group}%
\TextAt{14}{237.092}{369.196}{o}%
\TextAt{3}{242.936}{373.534}{\Character{3}}%
\TextAt{8}{252.568}{369.196}{onto}%
\TextAt{8}{279.919}{369.196}{a}%
\TextAt{8}{290.677}{369.196}{complete}%
\TextAt{8}{341.101}{369.196}{lattice}%
\TextAt{9}{377.873}{369.196}{\Character{0}}%
\TextAt{9}{390.199}{369.196}{\Character{61}}%
\TextAt{13}{404.349}{369.196}{\Character{21}}%
\TextAt{9}{411.177}{369.196}{\Character{40}}%
\TextAt{14}{415.729}{369.196}{o}%
\TextAt{3}{421.574}{373.534}{\Character{3}}%
\TextAt{9}{426.306}{369.196}{\Character{41}}%
\TextAt{8}{435.758}{369.196}{in}%
\TextAt{8}{450.413}{369.196}{trace-zero}%
\TextAt{8}{505.392}{369.196}{space}%
\TextAt{13}{79.226}{352.15}{H}%
\TextAt{9}{93.17}{352.15}{\Character{61}}%
\TextAt{16}{105.595}{352.15}{\Character{102}}%
\TextAt{9}{111.573}{352.15}{\Character{40}}%
\TextAt{13}{116.125}{352.15}{x}%
\TextAt{27}{122.777}{350.357}{p}%
\TextAt{9}{127.264}{352.15}{\Character{41}}%
\NativeMathAt{135.137}{352.15}{11.955}{\in}%
\TextAt{20}{146.428}{361.116}{\Character{81}}%
\TextAt{27}{157.719}{348.663}{p}%
\TextAt{3}{161.708}{348.663}{\Character{106}}%
\NativeMathAt{164.059}{348.663}{7.97}{\infty}%
\TextAt{18}{175.019}{352.15}{R}%
\TextAt{16}{186.974}{352.15}{\Character{106}}%
\TextAt{20}{193.616}{361.116}{\Character{80}}%
\TextAt{27}{206.235}{348.663}{p}%
\TextAt{3}{210.224}{348.663}{\Character{106}}%
\NativeMathAt{212.575}{348.663}{7.97}{\infty}%
\TextAt{13}{223.535}{352.15}{x}%
\TextAt{27}{230.188}{350.357}{p}%
\TextAt{9}{237.995}{352.15}{\Character{61}}%
\TextAt{9}{250.405}{352.15}{0}%
\TextAt{16}{256.274}{352.15}{\Character{103}}%
\TextAt{8}{262.251}{352.15}{.}%
\TextAt{12}{79.226}{323.149}{Proposition}%
\TextAt{12}{151.72}{323.149}{(1.9).}%
\TextAt{8}{186.847}{323.149}{The}%
\TextAt{8}{210.712}{323.149}{kernel}%
\TextAt{8}{244.99}{323.149}{of}%
\TextAt{8}{258.131}{323.149}{div}%
\TextAt{9}{277.388}{323.149}{\Character{58}}%
\TextAt{13}{283.961}{323.149}{K}%
\TextAt{3}{294.778}{327.487}{\Character{3}}%
\NativeMathAt{302.831}{323.149}{11.955}{\rightarrow}%
\TextAt{8}{318.107}{323.149}{Div}%
\TextAt{9}{336.483}{323.149}{\Character{40}}%
\Vector{q 0 G 0 g 0.478 w 269.035 -443.495 m 274.879 -443.495 l S Q}%
\TextAt{14}{341.035}{323.149}{o}%
\TextAt{9}{346.879}{323.149}{\Character{41}}%
\TextAt{8}{355.14}{323.149}{is}%
\TextAt{8}{366.715}{323.149}{the}%
\TextAt{8}{386.683}{323.149}{group}%
\TextAt{13}{419.669}{323.149}{\Character{22}}%
\TextAt{9}{426.712}{323.149}{\Character{40}}%
\TextAt{13}{431.265}{323.149}{K}%
\TextAt{9}{442.082}{323.149}{\Character{41}}%
\TextAt{8}{450.343}{323.149}{of}%
\TextAt{8}{463.484}{323.149}{roots}%
\TextAt{8}{492.956}{323.149}{of}%
\TextAt{8}{506.097}{323.149}{unity}%
\TextAt{8}{79.226}{306.103}{in}%
\TextAt{13}{92.879}{306.103}{K}%
\TextAt{8}{103.697}{306.103}{,}%
\TextAt{8}{110.846}{306.103}{and}%
\TextAt{8}{133.608}{306.103}{its}%
\TextAt{8}{149.927}{306.103}{image}%
\TextAt{16}{183.748}{306.103}{\Character{80}}%
\TextAt{9}{193.047}{306.103}{\Character{40}}%
\Vector{q 0 G 0 g 0.478 w 125.599 -460.541 m 131.444 -460.541 l S Q}%
\TextAt{14}{197.599}{306.103}{o}%
\TextAt{9}{203.444}{306.103}{\Character{41}}%
\TextAt{8}{211.893}{306.103}{is}%
\TextAt{8}{223.657}{306.103}{a}%
\TextAt{8}{233.413}{306.103}{discrete}%
\TextAt{8}{276.391}{306.103}{subgroup}%
\TextAt{8}{327.189}{306.103}{of}%
\TextAt{8}{340.519}{306.103}{Div}%
\TextAt{9}{358.894}{306.103}{\Character{40}}%
\Vector{q 0 G 0 g 0.478 w 291.446 -460.541 m 297.291 -460.541 l S Q}%
\TextAt{14}{363.446}{306.103}{o}%
\TextAt{9}{369.291}{306.103}{\Character{41}}%
\TextAt{8}{373.843}{306.103}{.}%
\TextAt{19}{79.226}{271.124}{Proof.}%
\TextAt{8}{116}{271.124}{By}%
\TextAt{8}{136.291}{271.124}{the}%
\TextAt{8}{158.375}{271.124}{above}%
\TextAt{8}{193.478}{271.124}{remark,}%
\TextAt{8}{239.141}{271.124}{the}%
\TextAt{8}{261.225}{271.124}{composite}%
\TextAt{8}{318.17}{271.124}{of}%
\TextAt{8}{333.428}{271.124}{div}%
\TextAt{8}{355.189}{271.124}{with}%
\TextAt{8}{383.776}{271.124}{the}%
\TextAt{8}{405.86}{271.124}{map}%
\TextAt{8}{433.802}{271.124}{Div}%
\TextAt{9}{452.177}{271.124}{\Character{40}}%
\Vector{q 0 G 0 g 0.478 w 384.729 -495.52 m 390.574 -495.52 l S Q}%
\TextAt{14}{456.729}{271.124}{o}%
\TextAt{9}{462.574}{271.124}{\Character{41}}%
\NativeMathAt{473.732}{271.124}{11.955}{\rightarrow}%
\TextAt{20}{492.292}{280.09}{\Character{81}}%
\TextAt{27}{503.583}{267.637}{p}%
\TextAt{3}{507.572}{267.637}{\Character{106}}%
\NativeMathAt{509.923}{267.637}{7.97}{\infty}%
\TextAt{18}{520.883}{271.124}{R}%
\TextAt{8}{529.518}{271.124}{,}%
\TextAt{20}{79.226}{263.044}{\Character{80}}%
\TextAt{27}{91.846}{250.591}{p}%
\TextAt{13}{98.325}{254.078}{\Character{23}}%
\TextAt{27}{104.109}{252.285}{p}%
\TextAt{14}{108.596}{254.078}{p}%
\NativeArrowAt{118.522}{254.078}{11.955}{11.955}{\mapsto}%
\TextAt{9}{134.421}{254.078}{\Character{40}}%
\TextAt{13}{138.974}{254.078}{\Character{23}}%
\TextAt{27}{144.758}{252.285}{p}%
\TextAt{13}{149.245}{254.078}{f}%
\TextAt{27}{155.014}{252.285}{p}%
\TextAt{9}{159.501}{254.078}{\Character{41}}%
\TextAt{27}{164.054}{252.218}{p}%
\TextAt{3}{168.042}{252.218}{\Character{106}}%
\NativeMathAt{170.393}{252.218}{7.97}{\infty}%
\TextAt{8}{179.361}{254.078}{,}%
\TextAt{8}{186.967}{254.078}{yields,}%
\TextAt{8}{223.577}{254.078}{up}%
\TextAt{8}{240.847}{254.078}{to}%
\TextAt{8}{255.523}{254.078}{sign,}%
\TextAt{8}{283.357}{254.078}{the}%
\TextAt{8}{303.879}{254.078}{homomorphism}%
\TextAt{13}{386.915}{254.078}{\Character{21}}%
\TextAt{9}{397.687}{254.078}{\Character{58}}%
\TextAt{13}{404.883}{254.078}{K}%
\TextAt{3}{415.7}{258.416}{\Character{3}}%
\NativeMathAt{424.377}{254.078}{11.955}{\rightarrow}%
\TextAt{20}{440.276}{263.044}{\Character{81}}%
\TextAt{27}{451.567}{250.591}{p}%
\TextAt{3}{455.555}{250.591}{\Character{106}}%
\NativeMathAt{457.906}{250.591}{7.97}{\infty}%
\TextAt{18}{468.867}{254.078}{R}%
\TextAt{8}{477.501}{254.078}{.}%
\TextAt{8}{487.046}{254.078}{By}%
\TextAt{8}{505.775}{254.078}{chap.}%
\TextAt{8}{79.226}{237.032}{I,}%
\TextAt{8}{90.596}{237.032}{(7.1),}%
\TextAt{8}{121.822}{237.032}{the}%
\TextAt{8}{141.979}{237.032}{latter}%
\TextAt{8}{173.851}{237.032}{fits}%
\TextAt{8}{193.422}{237.032}{into}%
\TextAt{8}{217.165}{237.032}{the}%
\TextAt{8}{237.321}{237.032}{exact}%
\TextAt{8}{268.213}{237.032}{sequence}%
\TextAt{9}{227.938}{208.031}{1}%
\NativeArrowAt{237.112}{208.031}{11.955}{19.26}{\longrightarrow}%
\TextAt{13}{259.694}{208.031}{\Character{22}}%
\TextAt{9}{266.737}{208.031}{\Character{40}}%
\TextAt{13}{271.289}{208.031}{K}%
\TextAt{9}{282.107}{208.031}{\Character{41}}%
\NativeArrowAt{289.98}{208.031}{11.955}{19.26}{\longrightarrow}%
\TextAt{14}{312.562}{208.031}{o}%
\TextAt{3}{318.406}{212.967}{\Character{3}}%
\TextAt{21}{329.262}{214.808}{\Character{21}}%
\NativeArrowAt{326.46}{208.031}{11.955}{11.955}{\longrightarrow}%
\TextAt{9}{341.736}{208.031}{\Character{0}}%
\NativeArrowAt{352.371}{208.031}{11.955}{19.26}{\longrightarrow}%
\TextAt{9}{374.953}{208.031}{0}%
\TextAt{13}{380.806}{208.031}{\Character{59}}%
\TextAt{8}{79.226}{179.03}{where}%
\TextAt{9}{111.859}{179.03}{\Character{0}}%
\TextAt{8}{121.895}{179.03}{is}%
\TextAt{8}{132.483}{179.03}{a}%
\TextAt{8}{141.062}{179.03}{complete}%
\TextAt{8}{189.309}{179.03}{lattice}%
\TextAt{8}{223.903}{179.03}{in}%
\TextAt{8}{236.379}{179.03}{trace-zero}%
\TextAt{8}{289.181}{179.03}{space}%
\TextAt{13}{319.28}{179.03}{H}%
\NativeMathAt{333.223}{179.03}{11.955}{\subseteq}%
\TextAt{20}{345.842}{187.996}{\Character{81}}%
\TextAt{27}{357.133}{175.543}{p}%
\TextAt{3}{361.122}{175.543}{\Character{106}}%
\NativeMathAt{363.473}{175.543}{7.97}{\infty}%
\TextAt{18}{374.433}{179.03}{R}%
\TextAt{8}{383.068}{179.03}{.}%
\TextAt{8}{391.124}{179.03}{Therefore}%
\TextAt{13}{442.945}{179.03}{\Character{22}}%
\TextAt{9}{449.988}{179.03}{\Character{40}}%
\TextAt{13}{454.541}{179.03}{K}%
\TextAt{9}{465.358}{179.03}{\Character{41}}%
\TextAt{8}{472.632}{179.03}{is}%
\TextAt{8}{483.22}{179.03}{the}%
\TextAt{8}{502.2}{179.03}{kernel}%
\TextAt{8}{79.226}{161.984}{of}%
\TextAt{8}{91.32}{161.984}{div.}%
\TextAt{8}{115.292}{161.984}{Since}%
\TextAt{9}{144.613}{161.984}{\Character{0}}%
\TextAt{8}{154.588}{161.984}{is}%
\TextAt{8}{165.115}{161.984}{a}%
\TextAt{8}{173.634}{161.984}{lattice,}%
\TextAt{8}{211.666}{161.984}{there}%
\TextAt{8}{240.341}{161.984}{exists}%
\TextAt{8}{271.419}{161.984}{a}%
\TextAt{8}{279.938}{161.984}{neighbourhood}%
\TextAt{13}{358.706}{161.984}{U}%
\TextAt{8}{370.565}{161.984}{of}%
\TextAt{9}{382.659}{161.984}{0}%
\TextAt{8}{391.173}{161.984}{in}%
\TextAt{20}{403.589}{170.95}{\Character{81}}%
\TextAt{27}{414.88}{158.497}{p}%
\TextAt{3}{418.868}{158.497}{\Character{106}}%
\NativeMathAt{421.219}{158.497}{7.97}{\infty}%
\TextAt{18}{432.18}{161.984}{R}%
\TextAt{8}{443.475}{161.984}{which}%
\TextAt{8}{475.725}{161.984}{contains}%
\TextAt{8}{520.408}{161.984}{no}%
\TextAt{8}{79.226}{144.938}{element}%
\TextAt{8}{122.779}{144.938}{of}%
\TextAt{9}{136.42}{144.938}{\Character{0}}%
\TextAt{8}{147.942}{144.938}{except}%
\TextAt{9}{184.992}{144.938}{0}%
\TextAt{8}{190.845}{144.938}{.}%
\TextAt{8}{200.226}{144.938}{Considering}%
\TextAt{8}{264.987}{144.938}{the}%
\TextAt{8}{285.454}{144.938}{isomorphism}%
\TextAt{13}{354.185}{144.938}{\Character{11}}%
\TextAt{9}{365.557}{144.938}{\Character{58}}%
\TextAt{20}{372.66}{153.904}{\Character{81}}%
\TextAt{27}{383.951}{141.451}{p}%
\TextAt{3}{387.939}{141.451}{\Character{106}}%
\NativeMathAt{390.29}{141.451}{7.97}{\infty}%
\TextAt{18}{401.251}{144.938}{R}%
\NativeMathAt{413.736}{144.938}{11.955}{\rightarrow}%
\TextAt{20}{429.542}{153.904}{\Character{76}}%
\TextAt{27}{442.825}{141.451}{p}%
\TextAt{3}{446.814}{141.451}{\Character{106}}%
\NativeMathAt{449.165}{141.451}{7.97}{\infty}%
\TextAt{18}{460.125}{144.938}{R}%
\TextAt{14}{468.76}{144.938}{p}%
\TextAt{8}{474.743}{144.938}{,}%
\TextAt{9}{482.281}{144.938}{\Character{40}}%
\TextAt{13}{486.833}{144.938}{\Character{23}}%
\TextAt{27}{492.617}{143.145}{p}%
\TextAt{9}{497.104}{144.938}{\Character{41}}%
\TextAt{27}{501.656}{143.078}{p}%
\TextAt{3}{505.645}{143.078}{\Character{106}}%
\NativeMathAt{507.996}{143.078}{7.97}{\infty}%
\NativeArrowAt{520.814}{144.938}{11.955}{11.955}{\mapsto}%
\TextAt{20}{79.226}{136.858}{\Character{80}}%
\TextAt{27}{91.846}{124.405}{p}%
\TextAt{3}{95.834}{124.405}{\Character{106}}%
\NativeMathAt{98.185}{124.405}{7.97}{\infty}%
\TextAt{9}{107.153}{127.892}{\Character{40}}%
\TextAt{13}{111.705}{127.892}{\Character{23}}%
\TextAt{27}{117.489}{126.099}{p}%
\TextAt{9}{121.976}{127.892}{\Character{47}}%
\TextAt{13}{127.829}{127.892}{f}%
\TextAt{27}{133.599}{126.099}{p}%
\TextAt{9}{138.086}{127.892}{\Character{41}}%
\TextAt{14}{142.638}{127.892}{p}%
\TextAt{8}{148.621}{127.892}{,}%
\TextAt{8}{158.126}{127.892}{the}%
\TextAt{8}{180.167}{127.892}{set}%
\TextAt{16}{200.319}{127.892}{\Character{102}}%
\TextAt{9}{206.297}{127.892}{0}%
\TextAt{16}{212.15}{127.892}{\Character{103}}%
\TextAt{16}{222.061}{127.892}{\Character{2}}%
\TextAt{13}{235.309}{127.892}{\Character{11}U}%
\NativeMathAt{258.562}{127.892}{11.955}{\subset}%
\TextAt{8}{274.393}{127.892}{Div}%
\TextAt{9}{292.768}{127.892}{\Character{40}}%
\TextAt{14}{297.32}{127.892}{o}%
\TextAt{9}{303.165}{127.892}{\Character{41}}%
\TextAt{16}{311.658}{127.892}{\Character{2}}%
\TextAt{20}{324.898}{136.858}{\Character{76}}%
\TextAt{27}{338.182}{124.405}{p}%
\TextAt{3}{342.17}{124.405}{\Character{106}}%
\NativeMathAt{344.521}{124.405}{7.97}{\infty}%
\TextAt{18}{355.482}{127.892}{R}%
\TextAt{14}{364.116}{127.892}{p}%
\TextAt{9}{376.631}{127.892}{\Character{61}}%
\TextAt{8}{392.269}{127.892}{Div}%
\TextAt{9}{410.644}{127.892}{\Character{40}}%
\Vector{q 0 G 0 g 0.478 w 343.196 -638.752 m 349.04 -638.752 l S Q}%
\TextAt{14}{415.196}{127.892}{o}%
\TextAt{9}{421.04}{127.892}{\Character{41}}%
\TextAt{8}{431.374}{127.892}{is}%
\TextAt{8}{445.023}{127.892}{a}%
\TextAt{8}{456.663}{127.892}{neighbourhood}%
\TextAt{8}{79.226}{110.846}{of}%
\TextAt{9}{94.17}{110.846}{0}%
\TextAt{8}{105.534}{110.846}{in}%
\TextAt{8}{120.801}{110.846}{Div}%
\TextAt{9}{139.176}{110.846}{\Character{40}}%
\Vector{q 0 G 0 g 0.478 w 71.728 -655.798 m 77.572 -655.798 l S Q}%
\TextAt{14}{143.728}{110.846}{o}%
\TextAt{9}{149.572}{110.846}{\Character{41}}%
\TextAt{8}{159.636}{110.846}{which}%
\TextAt{8}{194.736}{110.846}{contains}%
\TextAt{8}{242.27}{110.846}{no}%
\TextAt{8}{260.142}{110.846}{replete}%
\TextAt{8}{300.12}{110.846}{principal}%
\TextAt{8}{350.511}{110.846}{divisor}%
\TextAt{8}{390.238}{110.846}{except}%
\TextAt{9}{428.59}{110.846}{0}%
\TextAt{8}{434.443}{110.846}{.}%
\TextAt{8}{447.733}{110.846}{This}%
\TextAt{8}{476.066}{110.846}{shows}%
\TextAt{8}{511.298}{110.846}{that}%
\TextAt{16}{79.226}{93.8}{\Character{80}}%
\TextAt{9}{88.525}{93.8}{\Character{40}}%
\Vector{q 0 G 0 g 0.478 w 21.077 -672.844 m 26.921 -672.844 l S Q}%
\TextAt{14}{93.077}{93.8}{o}%
\TextAt{9}{98.921}{93.8}{\Character{41}}%
\TextAt{9}{106.776}{93.8}{\Character{61}}%
\TextAt{8}{119.22}{93.8}{div}%
\TextAt{9}{135.156}{93.8}{\Character{40}}%
\TextAt{13}{139.708}{93.8}{K}%
\TextAt{3}{150.526}{98.138}{\Character{3}}%
\TextAt{9}{155.258}{93.8}{\Character{41}}%
\TextAt{8}{163.708}{93.8}{lies}%
\TextAt{8}{183.924}{93.8}{discretely}%
\TextAt{8}{236.336}{93.8}{in}%
\TextAt{8}{249.989}{93.8}{Div}%
\TextAt{9}{268.364}{93.8}{\Character{40}}%
\Vector{q 0 G 0 g 0.478 w 200.916 -672.844 m 206.76 -672.844 l S Q}%
\TextAt{14}{272.916}{93.8}{o}%
\TextAt{9}{278.76}{93.8}{\Character{41}}%
\TextAt{8}{283.313}{93.8}{.}%
\Vector{q 0 G 0 g 0.399 w 452.334 -680.2 m 452.334 -672.13 l S Q}%
\Vector{q 0 G 0 g 0.399 w 452.534 -672.33 m 459.707 -672.33 l S Q}%
\Vector{q 0 G 0 g 0.399 w 452.534 -680.001 m 459.707 -680.001 l S Q}%
\Vector{q 0 G 0 g 0.399 w 459.906 -680.2 m 459.906 -672.13 l S Q}%
\TextAt{8}{297.211}{43.033}{173}%
\Destination{78.23}{809.815}{nav.page.173}%
\Destination{79.36}{338.155}{nav.print.190}%
\Destination{79.36}{337.155}{nav.number.3.1.9}%
\Hotlink{254.054}{384.055}{5.22}{11.047}{nav.chapter.1}%
\Hotlink{379.812}{384.055}{14.294}{11.047}{nav.chapter.3}%
\Hotlink{524.797}{384.055}{5.22}{11.047}{nav.chapter.1}%
\Hotlink{272.72}{384.055}{6.858}{11.047}{nav.section.1.7}%
\Hotlink{407.563}{384.055}{6.858}{11.047}{nav.section.3.3}%
\Hotlink{297.503}{384.055}{12.716}{11.047}{nav.print.39}%
\Hotlink{432.346}{384.055}{18.574}{11.047}{nav.print.211}%
\Hotlink{78.726}{366.955}{25.078}{11.047}{nav.number.1.7.3}%
\Hotlink{90.096}{234.855}{25.078}{11.047}{nav.number.1.7.1}%
}

% Page 186 (body page 174)
\Page{%
\TextAt{2}{79.226}{796.943}{III.}%
\TextAt{2}{98.317}{796.943}{Riemann–Roch}%
\TextAt{2}{174.375}{796.943}{Theory}%
\Vector{q 0 G 0 g 0.299 w 7.226 18.729 m 460.769 18.729 l S Q}%
\TextAt{12}{79.226}{761.124}{Definition}%
\TextAt{12}{142.015}{761.124}{(1.10).}%
\TextAt{8}{184.061}{761.124}{The}%
\TextAt{8}{208.115}{761.124}{factor}%
\TextAt{8}{241.613}{761.124}{group}%
\TextAt{13}{246.644}{732.994}{C}%
\TextAt{13}{255.862}{732.994}{H}%
\TextAt{17}{266.5}{737.931}{1}%
\TextAt{9}{271.232}{732.994}{\Character{40}}%
\Vector{q 0 G 0 g 0.478 w 203.785 -33.649 m 209.629 -33.649 l S Q}%
\TextAt{14}{275.785}{732.994}{o}%
\TextAt{9}{281.629}{732.994}{\Character{41}}%
\TextAt{9}{289.484}{732.994}{\Character{61}}%
\TextAt{8}{301.927}{732.994}{Div}%
\TextAt{9}{320.303}{732.994}{\Character{40}}%
\Vector{q 0 G 0 g 0.478 w 252.855 -33.649 m 258.699 -33.649 l S Q}%
\TextAt{14}{324.855}{732.994}{o}%
\TextAt{9}{330.699}{732.994}{\Character{41}\Character{47}}%
\TextAt{16}{341.104}{732.994}{\Character{80}}%
\TextAt{9}{350.403}{732.994}{\Character{40}}%
\Vector{q 0 G 0 g 0.478 w 282.955 -33.649 m 288.8 -33.649 l S Q}%
\TextAt{14}{354.955}{732.994}{o}%
\TextAt{9}{360.8}{732.994}{\Character{41}}%
\TextAt{8}{79.226}{704.865}{is}%
\TextAt{8}{90.99}{704.865}{called}%
\TextAt{8}{124.154}{704.865}{the}%
\TextAt{12}{144.31}{704.865}{replete}%
\TextAt{12}{189.142}{704.865}{divisor}%
\TextAt{12}{233.221}{704.865}{class}%
\TextAt{12}{264.579}{704.865}{group}%
\TextAt{8}{302.381}{704.865}{(or}%
\TextAt{12}{321.247}{704.865}{Arakelov}%
\TextAt{12}{377.998}{704.865}{class}%
\TextAt{12}{409.356}{704.865}{group}%
\TextAt{8}{443.261}{704.865}{)}%
\TextAt{8}{451.713}{704.865}{of}%
\TextAt{13}{465.043}{704.865}{K}%
\TextAt{8}{475.861}{704.865}{.}%
\TextAt{8}{97.159}{679.102}{As}%
\TextAt{16}{114.903}{679.102}{\Character{80}}%
\TextAt{9}{124.202}{679.102}{\Character{40}}%
\Vector{q 0 G 0 g 0.478 w 56.754 -87.542 m 62.598 -87.542 l S Q}%
\TextAt{14}{128.754}{679.102}{o}%
\TextAt{9}{134.598}{679.102}{\Character{41}}%
\TextAt{8}{143.505}{679.102}{is}%
\TextAt{8}{155.726}{679.102}{discrete}%
\TextAt{8}{199.162}{679.102}{in}%
\TextAt{8}{213.272}{679.102}{Div}%
\TextAt{9}{231.647}{679.102}{\Character{40}}%
\Vector{q 0 G 0 g 0.478 w 164.2 -87.542 m 170.044 -87.542 l S Q}%
\TextAt{14}{236.2}{679.102}{o}%
\TextAt{9}{242.044}{679.102}{\Character{41}}%
\TextAt{8}{246.596}{679.102}{,}%
\TextAt{8}{254.317}{679.102}{and}%
\TextAt{8}{277.537}{679.102}{is}%
\TextAt{8}{289.758}{679.102}{therefore}%
\TextAt{8}{339.315}{679.102}{in}%
\TextAt{8}{353.425}{679.102}{particular}%
\TextAt{8}{407.872}{679.102}{closed,}%
\TextAt{13}{446.221}{679.102}{C}%
\TextAt{13}{455.439}{679.102}{H}%
\TextAt{17}{466.078}{683.44}{1}%
\TextAt{9}{470.81}{679.102}{\Character{40}}%
\Vector{q 0 G 0 g 0.478 w 403.362 -87.542 m 409.207 -87.542 l S Q}%
\TextAt{14}{475.362}{679.102}{o}%
\TextAt{9}{481.207}{679.102}{\Character{41}}%
\TextAt{8}{490.113}{679.102}{becomes}%
\TextAt{8}{79.226}{662.056}{a}%
\TextAt{8}{89.406}{662.056}{locally}%
\TextAt{8}{126.904}{662.056}{compact}%
\TextAt{8}{174.157}{662.056}{Hausdorff}%
\TextAt{8}{227.974}{662.056}{topological}%
\TextAt{8}{287.923}{662.056}{group}%
\TextAt{8}{321.523}{662.056}{with}%
\TextAt{8}{348.608}{662.056}{respect}%
\TextAt{8}{389.083}{662.056}{to}%
\TextAt{8}{403.818}{662.056}{the}%
\TextAt{8}{424.399}{662.056}{quotient}%
\TextAt{8}{471.006}{662.056}{topology.}%
\TextAt{8}{523.994}{662.056}{It}%
\TextAt{8}{79.226}{645.01}{is}%
\TextAt{8}{91.828}{645.01}{the}%
\TextAt{8}{112.822}{645.01}{correct}%
\TextAt{8}{152.681}{645.01}{analogue}%
\TextAt{8}{202.308}{645.01}{of}%
\TextAt{8}{216.476}{645.01}{the}%
\TextAt{8}{237.47}{645.01}{divisor}%
\TextAt{8}{276.421}{645.01}{class}%
\TextAt{8}{304.696}{645.01}{group}%
\TextAt{8}{338.709}{645.01}{of}%
\TextAt{8}{352.877}{645.01}{a}%
\TextAt{8}{363.47}{645.01}{function}%
\TextAt{8}{410.156}{645.01}{field}%
\TextAt{8}{436.351}{645.01}{(see}%
\TextAt{8}{460.656}{645.01}{chap.}%
\TextAt{8}{495.361}{645.01}{I,}%
\TextAt{8}{507.568}{645.01}{§14).}%
\TextAt{8}{79.226}{627.964}{For}%
\TextAt{8}{100.593}{627.964}{the}%
\TextAt{8}{121.148}{627.964}{latter}%
\TextAt{8}{153.418}{627.964}{we}%
\TextAt{8}{171.044}{627.964}{introduced}%
\TextAt{8}{229.974}{627.964}{the}%
\TextAt{8}{250.528}{627.964}{degree}%
\TextAt{8}{287.342}{627.964}{map}%
\TextAt{8}{313.754}{627.964}{onto}%
\TextAt{8}{340.501}{627.964}{the}%
\TextAt{8}{361.055}{627.964}{group}%
\TextAt{18}{394.629}{627.964}{Z}%
\TextAt{8}{402.599}{627.964}{;}%
\TextAt{8}{410.345}{627.964}{for}%
\TextAt{13}{428.628}{627.964}{C}%
\TextAt{13}{437.846}{627.964}{H}%
\TextAt{17}{448.484}{632.302}{1}%
\TextAt{9}{453.217}{627.964}{\Character{40}}%
\Vector{q 0 G 0 g 0.478 w 385.769 -138.68 m 391.613 -138.68 l S Q}%
\TextAt{14}{457.769}{627.964}{o}%
\TextAt{9}{463.613}{627.964}{\Character{41}}%
\TextAt{8}{472.461}{627.964}{we}%
\TextAt{8}{490.086}{627.964}{obtain}%
\TextAt{8}{526.911}{627.964}{a}%
\TextAt{12}{79.226}{610.918}{degree}%
\TextAt{12}{121.81}{610.918}{map}%
\TextAt{8}{150.933}{610.918}{onto}%
\TextAt{8}{177.282}{610.918}{the}%
\TextAt{8}{197.439}{610.918}{group}%
\TextAt{18}{230.614}{610.918}{R}%
\TextAt{8}{239.249}{610.918}{.}%
\TextAt{8}{247.697}{610.918}{It}%
\TextAt{8}{260.37}{610.918}{is}%
\TextAt{8}{272.133}{610.918}{induced}%
\TextAt{8}{315.698}{610.918}{by}%
\TextAt{8}{331.957}{610.918}{the}%
\TextAt{8}{352.114}{610.918}{continuous}%
\TextAt{8}{410.718}{610.918}{homomorphism}%
\TextAt{8}{258.34}{582.789}{deg}%
\TextAt{9}{279.222}{582.789}{\Character{58}}%
\TextAt{8}{285.795}{582.789}{Div}%
\TextAt{9}{304.17}{582.789}{\Character{40}}%
\Vector{q 0 G 0 g 0.478 w 236.722 -183.855 m 242.567 -183.855 l S Q}%
\TextAt{14}{308.722}{582.789}{o}%
\TextAt{9}{314.567}{582.789}{\Character{41}}%
\NativeArrowAt{322.44}{582.789}{11.955}{19.26}{\longrightarrow}%
\TextAt{18}{345.022}{582.789}{R}%
\TextAt{8}{79.226}{554.659}{which}%
\TextAt{8}{112.713}{554.659}{sends}%
\TextAt{8}{144.047}{554.659}{a}%
\TextAt{8}{153.803}{554.659}{replete}%
\TextAt{8}{192.167}{554.659}{divisor}%
\TextAt{13}{230.28}{554.659}{D}%
\TextAt{9}{243.645}{554.659}{\Character{61}}%
\TextAt{20}{256.071}{563.626}{\Character{80}}%
\TextAt{27}{268.69}{551.172}{p}%
\TextAt{13}{275.169}{554.659}{\Character{23}}%
\TextAt{27}{280.953}{552.866}{p}%
\TextAt{14}{285.44}{554.659}{p}%
\TextAt{8}{295.32}{554.659}{to}%
\TextAt{8}{309.631}{554.659}{the}%
\TextAt{8}{329.787}{554.659}{real}%
\TextAt{8}{352.55}{554.659}{number}%
\TextAt{8}{189.045}{516.664}{deg}%
\TextAt{9}{206.607}{516.664}{\Character{40}}%
\TextAt{13}{211.159}{516.664}{D}%
\TextAt{9}{221.204}{516.664}{\Character{41}}%
\TextAt{9}{229.059}{516.664}{\Character{61}}%
\TextAt{20}{241.502}{528.022}{\Character{88}}%
\TextAt{27}{248.142}{502.916}{p}%
\TextAt{13}{260.764}{516.664}{\Character{23}}%
\TextAt{27}{266.547}{514.871}{p}%
\TextAt{8}{273.027}{516.664}{log}%
\TextAt{14}{289.987}{516.664}{N}%
\TextAt{9}{299.935}{516.664}{\Character{40}}%
\TextAt{14}{304.487}{516.664}{p}%
\TextAt{9}{310.471}{516.664}{\Character{41}}%
\TextAt{9}{318.326}{516.664}{\Character{61}}%
\TextAt{8}{330.769}{516.664}{log}%
\TextAt{20}{347.73}{537.108}{\Character{32}}%
\TextAt{20}{357.194}{528.022}{\Character{89}}%
\TextAt{27}{362.838}{502.916}{p}%
\TextAt{14}{374.463}{516.664}{N}%
\TextAt{9}{384.411}{516.664}{\Character{40}}%
\TextAt{14}{388.963}{516.664}{p}%
\TextAt{9}{394.946}{516.664}{\Character{41}}%
\TextAt{21}{399.499}{521.6}{\Character{23}}%
\TextAt{28}{403.661}{520.604}{p}%
\TextAt{20}{408.242}{537.108}{\Character{33}}%
\TextAt{13}{419.699}{516.664}{\Character{58}}%
\TextAt{8}{79.226}{479.167}{From}%
\TextAt{8}{110.447}{479.167}{the}%
\TextAt{8}{131.099}{479.167}{product}%
\TextAt{8}{175.495}{479.167}{formula}%
\TextAt{8}{218.922}{479.167}{(1.3),}%
\TextAt{8}{250.769}{479.167}{we}%
\TextAt{8}{268.492}{479.167}{find}%
\TextAt{8}{292.396}{479.167}{for}%
\TextAt{8}{310.777}{479.167}{any}%
\TextAt{8}{333.39}{479.167}{replete}%
\TextAt{8}{372.25}{479.167}{principal}%
\TextAt{8}{421.523}{479.167}{divisor}%
\TextAt{8}{460.132}{479.167}{div}%
\TextAt{9}{476.069}{479.167}{\Character{40}}%
\TextAt{13}{480.621}{479.167}{f}%
\TextAt{9}{487.667}{479.167}{\Character{41}}%
\NativeMathAt{496.386}{479.167}{11.955}{\in}%
\TextAt{16}{508.509}{479.167}{\Character{80}}%
\TextAt{9}{517.821}{479.167}{\Character{40}}%
\Vector{q 0 G 0 g 0.478 w 450.373 -287.476 m 456.217 -287.476 l S Q}%
\TextAt{14}{522.373}{479.167}{o}%
\TextAt{9}{528.217}{479.167}{\Character{41}}%
\TextAt{8}{79.226}{462.121}{that}%
\TextAt{8}{168.619}{440.084}{deg}%
\TextAt{9}{186.181}{440.084}{\Character{40}}%
\TextAt{8}{190.733}{440.084}{div}%
\TextAt{9}{206.67}{440.084}{\Character{40}}%
\TextAt{13}{211.222}{440.084}{f}%
\TextAt{9}{218.268}{440.084}{\Character{41}\Character{41}}%
\TextAt{9}{230.666}{440.084}{\Character{61}}%
\TextAt{20}{243.119}{451.441}{\Character{88}}%
\TextAt{27}{249.759}{426.335}{p}%
\TextAt{8}{262.38}{440.084}{log}%
\TextAt{14}{279.341}{440.084}{N}%
\TextAt{9}{289.289}{440.084}{\Character{40}}%
\TextAt{14}{293.841}{440.084}{p}%
\TextAt{9}{299.824}{440.084}{\Character{41}}%
\TextAt{21}{304.377}{445.02}{v}%
\TextAt{28}{308.5}{444.024}{p}%
\TextAt{17}{312.583}{445.02}{\Character{40}}%
\TextAt{21}{315.876}{445.02}{f}%
\TextAt{17}{320.823}{445.02}{\Character{41}}%
\TextAt{9}{327.935}{440.084}{\Character{61}}%
\TextAt{8}{340.361}{440.084}{log}%
\TextAt{20}{357.321}{460.527}{\Character{32}}%
\TextAt{20}{366.786}{451.441}{\Character{89}}%
\TextAt{27}{372.429}{426.335}{p}%
\TextAt{16}{384.054}{440.084}{\Character{106}}%
\TextAt{13}{387.375}{440.084}{f}%
\TextAt{16}{394.421}{440.084}{\Character{106}}%
\TextAt{3}{397.742}{445.02}{\Character{0}}%
\TextAt{17}{404.329}{445.02}{1}%
\TextAt{27}{397.742}{437.128}{p}%
\TextAt{20}{409.061}{460.527}{\Character{33}}%
\TextAt{9}{421.847}{440.084}{\Character{61}}%
\TextAt{9}{434.257}{440.084}{0}%
\TextAt{13}{440.125}{440.084}{\Character{58}}%
\TextAt{8}{79.226}{408.118}{Thus}%
\TextAt{8}{108.875}{408.118}{we}%
\TextAt{8}{126.102}{408.118}{obtain}%
\TextAt{8}{162.53}{408.118}{a}%
\TextAt{8}{172.285}{408.118}{well-defined}%
\TextAt{8}{236.329}{408.118}{continuous}%
\TextAt{8}{294.933}{408.118}{homomorphism}%
\TextAt{8}{253.607}{379.989}{deg}%
\TextAt{9}{274.49}{379.989}{\Character{58}}%
\TextAt{13}{281.062}{379.989}{C}%
\TextAt{13}{290.28}{379.989}{H}%
\TextAt{17}{300.919}{384.925}{1}%
\TextAt{9}{305.651}{379.989}{\Character{40}}%
\Vector{q 0 G 0 g 0.478 w 238.203 -386.655 m 244.048 -386.655 l S Q}%
\TextAt{14}{310.203}{379.989}{o}%
\TextAt{9}{316.048}{379.989}{\Character{41}}%
\NativeArrowAt{323.921}{379.989}{11.955}{19.26}{\longrightarrow}%
\TextAt{18}{346.503}{379.989}{R}%
\TextAt{13}{355.137}{379.989}{\Character{58}}%
\TextAt{8}{79.226}{351.86}{The}%
\TextAt{8}{103.881}{351.86}{kernel}%
\TextAt{13}{138.95}{351.86}{C}%
\TextAt{13}{148.168}{351.86}{H}%
\TextAt{17}{158.806}{356.198}{1}%
\TextAt{9}{163.538}{351.86}{\Character{40}}%
\Vector{q 0 G 0 g 0.478 w 96.091 -414.784 m 101.935 -414.784 l S Q}%
\TextAt{14}{168.091}{351.86}{o}%
\TextAt{9}{173.935}{351.86}{\Character{41}}%
\TextAt{17}{178.487}{356.198}{0}%
\TextAt{8}{187.718}{351.86}{of}%
\TextAt{8}{201.65}{351.86}{this}%
\TextAt{8}{225.074}{351.86}{map}%
\TextAt{8}{251.69}{351.86}{is}%
\TextAt{8}{264.055}{351.86}{made}%
\TextAt{8}{295.872}{351.86}{up}%
\TextAt{8}{313.378}{351.86}{from}%
\TextAt{8}{341.62}{351.86}{the}%
\TextAt{8}{362.378}{351.86}{unit}%
\TextAt{8}{387.69}{351.86}{group}%
\TextAt{14}{421.468}{351.86}{o}%
\TextAt{3}{427.312}{356.198}{\Character{3}}%
\TextAt{8}{436.543}{351.86}{and}%
\TextAt{8}{459.907}{351.86}{the}%
\TextAt{8}{480.665}{351.86}{ideal}%
\TextAt{8}{509.23}{351.86}{class}%
\TextAt{8}{79.226}{334.814}{group}%
\TextAt{8}{112.402}{334.814}{Cl}%
\TextAt{21}{124.106}{333.02}{K}%
\TextAt{8}{136.206}{334.814}{as}%
\TextAt{8}{150.576}{334.814}{follows.}%
\TextAt{12}{79.226}{309.05}{Proposition}%
\TextAt{12}{153.114}{309.05}{(1.11).}%
\TextAt{8}{196.183}{309.05}{Let}%
\TextAt{9}{218.176}{309.05}{\Character{0}}%
\TextAt{9}{230.538}{309.05}{\Character{61}}%
\TextAt{13}{244.723}{309.05}{\Character{21}}%
\TextAt{9}{251.551}{309.05}{\Character{40}}%
\TextAt{14}{256.103}{309.05}{o}%
\TextAt{3}{261.948}{313.389}{\Character{3}}%
\TextAt{9}{266.68}{309.05}{\Character{41}}%
\TextAt{8}{276.153}{309.05}{denote}%
\TextAt{8}{314.894}{309.05}{the}%
\TextAt{8}{336.073}{309.05}{complete}%
\TextAt{8}{386.519}{309.05}{lattice}%
\TextAt{8}{423.312}{309.05}{of}%
\TextAt{8}{437.665}{309.05}{units}%
\TextAt{8}{468.014}{309.05}{in}%
\TextAt{8}{482.689}{309.05}{trace-zero}%
\TextAt{8}{79.226}{292.004}{space}%
\TextAt{13}{110.501}{292.004}{H}%
\TextAt{9}{124.444}{292.004}{\Character{61}}%
\TextAt{16}{136.87}{292.004}{\Character{102}}%
\TextAt{9}{142.847}{292.004}{\Character{40}}%
\TextAt{13}{147.4}{292.004}{x}%
\TextAt{27}{154.052}{290.211}{p}%
\TextAt{9}{158.539}{292.004}{\Character{41}}%
\NativeMathAt{166.412}{292.004}{11.955}{\in}%
\TextAt{20}{177.703}{300.971}{\Character{81}}%
\TextAt{27}{188.994}{288.517}{p}%
\TextAt{3}{192.983}{288.517}{\Character{106}}%
\NativeMathAt{195.334}{288.517}{7.97}{\infty}%
\TextAt{18}{206.294}{292.004}{R}%
\TextAt{16}{218.249}{292.004}{\Character{106}}%
\TextAt{20}{224.891}{300.971}{\Character{80}}%
\TextAt{27}{237.51}{288.517}{p}%
\TextAt{3}{241.499}{288.517}{\Character{106}}%
\NativeMathAt{243.85}{288.517}{7.97}{\infty}%
\TextAt{13}{254.81}{292.004}{x}%
\TextAt{27}{261.462}{290.211}{p}%
\TextAt{9}{269.27}{292.004}{\Character{61}}%
\TextAt{9}{281.68}{292.004}{0}%
\TextAt{16}{287.548}{292.004}{\Character{103}}%
\TextAt{8}{293.526}{292.004}{.}%
\TextAt{8}{301.974}{292.004}{There}%
\TextAt{8}{335.784}{292.004}{is}%
\TextAt{8}{347.547}{292.004}{an}%
\TextAt{8}{363.806}{292.004}{exact}%
\TextAt{8}{394.699}{292.004}{sequence}%
\TextAt{9}{192.915}{263.875}{0}%
\NativeArrowAt{202.089}{263.875}{11.955}{19.26}{\longrightarrow}%
\TextAt{13}{224.671}{263.875}{H}%
\TextAt{9}{235.293}{263.875}{\Character{47}\Character{0}}%
\NativeArrowAt{251.781}{263.875}{11.955}{19.26}{\longrightarrow}%
\TextAt{13}{274.363}{263.875}{C}%
\TextAt{13}{283.581}{263.875}{H}%
\TextAt{17}{294.219}{268.811}{1}%
\TextAt{9}{298.952}{263.875}{\Character{40}}%
\Vector{q 0 G 0 g 0.478 w 231.504 -502.768 m 237.348 -502.768 l S Q}%
\TextAt{14}{303.504}{263.875}{o}%
\TextAt{9}{309.348}{263.875}{\Character{41}}%
\TextAt{17}{313.901}{268.811}{0}%
\NativeArrowAt{321.954}{263.875}{11.955}{19.26}{\longrightarrow}%
\TextAt{13}{344.536}{263.875}{C}%
\TextAt{13}{353.754}{263.875}{H}%
\TextAt{17}{364.392}{268.811}{1}%
\TextAt{9}{369.124}{263.875}{\Character{40}}%
\TextAt{14}{373.677}{263.875}{o}%
\TextAt{9}{379.521}{263.875}{\Character{41}}%
\NativeArrowAt{387.394}{263.875}{11.955}{19.26}{\longrightarrow}%
\TextAt{9}{409.976}{263.875}{0}%
\TextAt{13}{415.829}{263.875}{\Character{58}}%
\TextAt{19}{79.226}{218.311}{Proof.}%
\TextAt{8}{116}{218.311}{Let}%
\TextAt{8}{136.97}{218.311}{Div}%
\TextAt{9}{155.345}{218.311}{\Character{40}}%
\Vector{q 0 G 0 g 0.478 w 87.897 -548.332 m 93.741 -548.332 l S Q}%
\TextAt{14}{159.897}{218.311}{o}%
\TextAt{9}{165.741}{218.311}{\Character{41}}%
\TextAt{17}{170.294}{222.65}{0}%
\TextAt{8}{178.923}{218.311}{be}%
\TextAt{8}{194.848}{218.311}{the}%
\TextAt{8}{215.004}{218.311}{kernel}%
\TextAt{8}{249.471}{218.311}{of}%
\TextAt{8}{262.801}{218.311}{deg}%
\TextAt{9}{283.684}{218.311}{\Character{58}}%
\TextAt{8}{290.256}{218.311}{Div}%
\TextAt{9}{308.631}{218.311}{\Character{40}}%
\Vector{q 0 G 0 g 0.478 w 241.184 -548.332 m 247.028 -548.332 l S Q}%
\TextAt{14}{313.184}{218.311}{o}%
\TextAt{9}{319.028}{218.311}{\Character{41}}%
\NativeMathAt{326.901}{218.311}{11.955}{\rightarrow}%
\TextAt{18}{342.177}{218.311}{R}%
\TextAt{8}{350.811}{218.311}{.}%
\TextAt{8}{359.26}{218.311}{Consider}%
\TextAt{8}{408.097}{218.311}{the}%
\TextAt{8}{428.253}{218.311}{exact}%
\TextAt{8}{459.145}{218.311}{sequence}%
\TextAt{9}{204.056}{190.182}{0}%
\NativeArrowAt{213.23}{190.182}{11.955}{19.26}{\longrightarrow}%
\TextAt{20}{235.812}{201.539}{\Character{89}}%
\TextAt{27}{236.045}{175.636}{p}%
\TextAt{3}{240.034}{175.636}{\Character{106}}%
\NativeMathAt{242.385}{175.636}{7.97}{\infty}%
\TextAt{18}{253.08}{190.182}{R}%
\TextAt{21}{267.623}{196.959}{\Character{11}}%
\NativeArrowAt{265.035}{190.182}{11.955}{12.027}{\longrightarrow}%
\TextAt{8}{280.382}{190.182}{Div}%
\TextAt{9}{298.757}{190.182}{\Character{40}}%
\Vector{q 0 G 0 g 0.478 w 231.309 -576.462 m 237.153 -576.462 l S Q}%
\TextAt{14}{303.309}{190.182}{o}%
\TextAt{9}{309.153}{190.182}{\Character{41}}%
\NativeArrowAt{317.026}{190.182}{11.955}{19.26}{\longrightarrow}%
\TextAt{8}{339.608}{190.182}{Div}%
\TextAt{9}{357.983}{190.182}{\Character{40}}%
\TextAt{14}{362.536}{190.182}{o}%
\TextAt{9}{368.38}{190.182}{\Character{41}}%
\NativeArrowAt{376.253}{190.182}{11.955}{19.26}{\longrightarrow}%
\TextAt{9}{398.835}{190.182}{0}%
\TextAt{13}{404.688}{190.182}{\Character{59}}%
\TextAt{8}{79.226}{150.894}{where}%
\TextAt{13}{114.813}{150.894}{\Character{11}}%
\TextAt{9}{122.335}{150.894}{\Character{40}\Character{40}}%
\TextAt{13}{131.44}{150.894}{\Character{23}}%
\TextAt{27}{137.223}{149.101}{p}%
\TextAt{9}{141.71}{150.894}{\Character{41}\Character{41}}%
\TextAt{9}{157.144}{150.894}{\Character{61}}%
\TextAt{20}{172.621}{159.86}{\Character{80}}%
\TextAt{27}{185.24}{147.407}{p}%
\TextAt{3}{189.229}{147.407}{\Character{106}}%
\NativeMathAt{191.58}{147.407}{7.97}{\infty}%
\TextAt{9}{200.548}{150.894}{\Character{40}}%
\TextAt{13}{205.1}{150.894}{\Character{23}}%
\TextAt{27}{210.884}{149.101}{p}%
\TextAt{9}{215.371}{150.894}{\Character{47}}%
\TextAt{13}{221.224}{150.894}{f}%
\TextAt{27}{226.994}{149.101}{p}%
\TextAt{9}{231.48}{150.894}{\Character{41}}%
\TextAt{14}{236.033}{150.894}{p}%
\TextAt{8}{242.016}{150.894}{.}%
\TextAt{8}{255.798}{150.894}{Restricting}%
\TextAt{8}{317.639}{150.894}{to}%
\TextAt{8}{333.727}{150.894}{Div}%
\TextAt{9}{352.102}{150.894}{\Character{40}}%
\Vector{q 0 G 0 g 0.478 w 284.654 -615.75 m 290.498 -615.75 l S Q}%
\TextAt{14}{356.654}{150.894}{o}%
\TextAt{9}{362.498}{150.894}{\Character{41}}%
\TextAt{17}{367.051}{155.232}{0}%
\TextAt{8}{377.458}{150.894}{yields}%
\TextAt{8}{412.137}{150.894}{the}%
\TextAt{8}{434.071}{150.894}{exact}%
\TextAt{8}{466.741}{150.894}{commutative}%
\TextAt{8}{79.226}{133.848}{diagram}%
\TextAt{9}{201.895}{117.055}{0}%
\TextAt{13}{232.4}{117.055}{\Character{21}}%
\TextAt{9}{239.228}{117.055}{\Character{40}}%
\TextAt{14}{243.781}{117.055}{o}%
\TextAt{3}{249.625}{121.393}{\Character{3}}%
\TextAt{9}{254.357}{117.055}{\Character{41}}%
\TextAt{16}{290.465}{117.055}{\Character{80}}%
\TextAt{9}{299.764}{117.055}{\Character{40}}%
\NativeDiagram{%
\path[draw=black,line width=0.478bp] (232.316,-649.589) -- (238.16,-649.589);%
}%
\TextAt{14}{304.316}{117.055}{o}%
\TextAt{9}{310.16}{117.055}{\Character{41}}%
\TextAt{16}{350.807}{117.055}{\Character{80}}%
\TextAt{9}{360.105}{117.055}{\Character{40}}%
\TextAt{14}{364.658}{117.055}{o}%
\TextAt{9}{370.502}{117.055}{\Character{41}}%
\TextAt{9}{404.244}{117.055}{0}%
\TextAt{9}{201.895}{74.312}{0}%
\TextAt{13}{240.344}{74.312}{H}%
\TextAt{8}{283.561}{74.312}{Div}%
\TextAt{9}{301.936}{74.312}{\Character{40}}%
\NativeDiagram{%
\path[draw=black,line width=0.478bp] (234.489,-692.331) -- (240.333,-692.331);%
}%
\TextAt{14}{306.489}{74.312}{o}%
\TextAt{9}{312.333}{74.312}{\Character{41}}%
\TextAt{17}{316.885}{78.651}{0}%
\TextAt{8}{346.269}{74.312}{Div}%
\TextAt{9}{364.644}{74.312}{\Character{40}}%
\TextAt{14}{369.196}{74.312}{o}%
\TextAt{9}{375.041}{74.312}{\Character{41}}%
\TextAt{9}{404.244}{74.312}{0}%
\NativeDiagram{%
\path[draw=black,line width=0.478bp] (141.101,-653.957) -- (154.57,-653.957);%
\path[draw=black,line width=0.478bp,line cap=round,line join=round] (152.323,-651.087) .. controls (152.777,-652.809) and (153.79,-653.622) .. (154.81,-653.957) .. controls (153.79,-654.291) and (152.777,-655.104) .. (152.323,-656.826);%
\path[draw=black,line width=0.478bp] (192.263,-653.957) -- (212.637,-653.957);%
\path[draw=black,line width=0.478bp,line cap=round,line join=round] (210.389,-651.087) .. controls (210.844,-652.809) and (211.856,-653.622) .. (212.876,-653.957) .. controls (211.856,-654.291) and (210.844,-655.104) .. (210.389,-656.826);%
}%
\TextAt{21}{271.969}{122.859}{\Character{11}}%
\NativeDiagram{%
\path[draw=black,line width=0.478bp] (173.656,-664.513) -- (173.656,-686.462);%
\path[draw=black,line width=0.478bp,line cap=round,line join=round] (176.525,-684.214) .. controls (174.803,-684.668) and (173.99,-685.681) .. (173.656,-686.701) .. controls (173.321,-685.681) and (172.508,-684.668) .. (170.786,-684.214);%
\path[draw=black,line width=0.478bp] (248.067,-653.957) -- (272.979,-653.957);%
\path[draw=black,line width=0.478bp,line cap=round,line join=round] (270.732,-651.087) .. controls (271.186,-652.809) and (272.198,-653.622) .. (273.218,-653.957) .. controls (272.198,-654.291) and (271.186,-655.104) .. (270.732,-656.826);%
\path[draw=black,line width=0.478bp] (230.591,-664.513) -- (230.591,-685.156);%
\path[draw=black,line width=0.478bp,line cap=round,line join=round] (233.46,-682.909) .. controls (231.738,-683.363) and (230.926,-684.375) .. (230.591,-685.395) .. controls (230.256,-684.375) and (229.443,-683.363) .. (227.722,-682.909);%
\path[draw=black,line width=0.478bp] (308.409,-653.957) -- (326.417,-653.957);%
\path[draw=black,line width=0.478bp,line cap=round,line join=round] (324.17,-651.087) .. controls (324.624,-652.809) and (325.636,-653.622) .. (326.656,-653.957) .. controls (325.636,-654.291) and (324.624,-655.104) .. (324.17,-656.826);%
\path[draw=black,line width=0.478bp] (290.933,-664.513) -- (290.933,-685.665);%
\path[draw=black,line width=0.478bp,line cap=round,line join=round] (293.802,-683.417) .. controls (292.081,-683.871) and (291.268,-684.884) .. (290.933,-685.904) .. controls (290.598,-684.884) and (289.786,-683.871) .. (288.064,-683.417);%
\path[draw=black,line width=0.478bp] (141.101,-696.699) -- (162.514,-696.699);%
\path[draw=black,line width=0.478bp,line cap=round,line join=round] (160.267,-693.83) .. controls (160.721,-695.552) and (161.733,-696.365) .. (162.753,-696.699) .. controls (161.733,-697.034) and (160.721,-697.847) .. (160.267,-699.569);%
\path[draw=black,line width=0.478bp] (184.319,-696.699) -- (205.732,-696.699);%
\path[draw=black,line width=0.478bp,line cap=round,line join=round] (203.485,-693.83) .. controls (203.939,-695.552) and (204.951,-696.365) .. (205.971,-696.699) .. controls (204.951,-697.034) and (203.939,-697.847) .. (203.485,-699.569);%
}%
\TextAt{21}{264.545}{80.116}{\Character{11}}%
\NativeDiagram{%
\path[draw=black,line width=0.478bp] (254.971,-696.699) -- (268.441,-696.699);%
\path[draw=black,line width=0.478bp,line cap=round,line join=round] (266.193,-693.83) .. controls (266.648,-695.552) and (267.66,-696.365) .. (268.68,-696.699) .. controls (267.66,-697.034) and (266.648,-697.847) .. (266.193,-699.569);%
\path[draw=black,line width=0.478bp] (312.947,-696.699) -- (326.417,-696.699);%
\path[draw=black,line width=0.478bp,line cap=round,line join=round] (324.17,-693.83) .. controls (324.624,-695.552) and (325.636,-696.365) .. (326.656,-696.699) .. controls (325.636,-697.034) and (324.624,-697.847) .. (324.17,-699.569);%
}%
\TextAt{8}{297.211}{43.033}{174}%
\Destination{78.23}{809.815}{nav.page.174}%
\Destination{79.36}{324.055}{nav.print.191}%
\Destination{79.36}{775.155}{nav.number.3.1.10}%
\Destination{79.36}{323.055}{nav.number.3.1.11}%
\Hotlink{494.861}{642.855}{5.22}{11.047}{nav.chapter.1}%
\Hotlink{512.747}{642.855}{12.716}{11.047}{nav.section.1.14}%
\Hotlink{218.422}{476.955}{25.078}{11.047}{nav.number.3.1.3}%
}

% Page 187 (body page 175)
\Page{%
\TextAt{2}{79.226}{796.943}{III.}%
\TextAt{2}{98.317}{796.943}{Riemann–Roch}%
\TextAt{2}{174.375}{796.943}{Theory}%
\Vector{q 0 G 0 g 0.299 w 7.226 18.729 m 460.769 18.729 l S Q}%
\TextAt{8}{79.226}{761.124}{Via}%
\TextAt{8}{101.008}{761.124}{the}%
\TextAt{8}{121.165}{761.124}{snake}%
\TextAt{8}{153.097}{761.124}{lemma}%
\TextAt{8}{190.816}{761.124}{(see}%
\TextAt{8}{214.284}{761.124}{[23],}%
\TextAt{8}{239.653}{761.124}{chap.}%
\TextAt{8}{271.844}{761.124}{III,}%
\TextAt{8}{292.299}{761.124}{§3,}%
\TextAt{8}{310.985}{761.124}{(3.3)),}%
\TextAt{8}{346.767}{761.124}{this}%
\TextAt{8}{369.589}{761.124}{gives}%
\TextAt{8}{398.27}{761.124}{rise}%
\TextAt{8}{419.789}{761.124}{to}%
\TextAt{8}{434.099}{761.124}{the}%
\TextAt{8}{454.256}{761.124}{exact}%
\TextAt{8}{485.148}{761.124}{sequence}%
\TextAt{9}{183.317}{734.207}{0}%
\NativeArrowAt{192.491}{734.207}{11.955}{19.26}{\longrightarrow}%
\TextAt{13}{215.073}{734.207}{H}%
\TextAt{9}{225.696}{734.207}{\Character{47}}%
\TextAt{13}{231.548}{734.207}{\Character{21}}%
\TextAt{9}{238.377}{734.207}{\Character{40}}%
\TextAt{14}{242.929}{734.207}{o}%
\TextAt{3}{248.774}{739.143}{\Character{3}}%
\TextAt{9}{253.506}{734.207}{\Character{41}}%
\NativeArrowAt{261.379}{734.207}{11.955}{19.26}{\longrightarrow}%
\TextAt{13}{283.961}{734.207}{C}%
\TextAt{13}{293.178}{734.207}{H}%
\TextAt{17}{303.817}{739.143}{1}%
\TextAt{9}{308.55}{734.207}{\Character{40}}%
\Vector{q 0 G 0 g 0.478 w 241.102 -32.436 m 246.946 -32.436 l S Q}%
\TextAt{14}{313.102}{734.207}{o}%
\TextAt{9}{318.946}{734.207}{\Character{41}}%
\TextAt{17}{323.498}{739.143}{0}%
\NativeArrowAt{331.552}{734.207}{11.955}{19.26}{\longrightarrow}%
\TextAt{13}{354.134}{734.207}{C}%
\TextAt{13}{363.351}{734.207}{H}%
\TextAt{17}{373.99}{739.143}{1}%
\TextAt{9}{378.722}{734.207}{\Character{40}}%
\TextAt{14}{383.274}{734.207}{o}%
\TextAt{9}{389.119}{734.207}{\Character{41}}%
\NativeArrowAt{396.992}{734.207}{11.955}{19.26}{\longrightarrow}%
\TextAt{9}{419.574}{734.207}{0}%
\TextAt{13}{425.427}{734.207}{\Character{58}}%
\Vector{q 0 G 0 g 0.399 w 452.334 -66.709 m 452.334 -58.64 l S Q}%
\Vector{q 0 G 0 g 0.399 w 452.534 -58.839 m 459.707 -58.839 l S Q}%
\Vector{q 0 G 0 g 0.399 w 452.534 -66.51 m 459.707 -66.51 l S Q}%
\Vector{q 0 G 0 g 0.399 w 459.906 -66.709 m 459.906 -58.64 l S Q}%
\TextAt{8}{97.159}{678.885}{The}%
\TextAt{8}{122.983}{678.885}{two}%
\TextAt{8}{146.87}{678.885}{fundamental}%
\TextAt{8}{216.283}{678.885}{facts}%
\TextAt{8}{245.754}{678.885}{of}%
\TextAt{8}{260.854}{678.885}{algebraic}%
\TextAt{8}{312.059}{678.885}{number}%
\TextAt{8}{356.426}{678.885}{theory,}%
\TextAt{8}{397.66}{678.885}{the}%
\TextAt{8}{419.587}{678.885}{finiteness}%
\TextAt{8}{472.203}{678.885}{of}%
\TextAt{8}{487.303}{678.885}{the}%
\TextAt{8}{509.23}{678.885}{class}%
\TextAt{8}{79.226}{661.839}{number}%
\TextAt{8}{123.467}{661.839}{and}%
\TextAt{8}{147.875}{661.839}{Dirichlet’s}%
\TextAt{8}{205.673}{661.839}{unit}%
\TextAt{8}{232.029}{661.839}{theorem,}%
\TextAt{8}{282.862}{661.839}{now}%
\TextAt{8}{308.896}{661.839}{merge}%
\TextAt{8}{345.007}{661.839}{into}%
\TextAt{8}{370.395}{661.839}{(and}%
\TextAt{8}{399.358}{661.839}{are}%
\TextAt{8}{420.513}{661.839}{in}%
\TextAt{8}{435.811}{661.839}{fact}%
\TextAt{8}{460.541}{661.839}{equivalent}%
\TextAt{8}{517.802}{661.839}{to)}%
\TextAt{8}{79.226}{644.793}{the}%
\TextAt{8}{100.046}{644.793}{simple}%
\TextAt{8}{137.184}{644.793}{statement}%
\TextAt{8}{192.219}{644.793}{that}%
\TextAt{8}{218.251}{644.793}{the}%
\TextAt{8}{239.07}{644.793}{kernel}%
\TextAt{13}{274.2}{644.793}{C}%
\TextAt{13}{283.418}{644.793}{H}%
\TextAt{17}{294.056}{649.132}{1}%
\TextAt{9}{298.789}{644.793}{\Character{40}}%
\Vector{q 0 G 0 g 0.478 w 231.341 -121.85 m 237.185 -121.85 l S Q}%
\TextAt{14}{303.341}{644.793}{o}%
\TextAt{9}{309.185}{644.793}{\Character{41}}%
\TextAt{17}{313.738}{649.132}{0}%
\TextAt{8}{323.03}{644.793}{of}%
\TextAt{8}{337.023}{644.793}{the}%
\TextAt{8}{357.843}{644.793}{degree}%
\TextAt{8}{394.921}{644.793}{map}%
\TextAt{8}{421.599}{644.793}{deg}%
\TextAt{9}{443.612}{644.793}{\Character{58}}%
\TextAt{13}{451.314}{644.793}{C}%
\TextAt{13}{460.531}{644.793}{H}%
\TextAt{17}{471.17}{649.132}{1}%
\TextAt{9}{475.903}{644.793}{\Character{40}}%
\Vector{q 0 G 0 g 0.478 w 408.455 -121.85 m 414.299 -121.85 l S Q}%
\TextAt{14}{480.455}{644.793}{o}%
\TextAt{9}{486.299}{644.793}{\Character{41}}%
\NativeMathAt{495.302}{644.793}{11.955}{\rightarrow}%
\TextAt{18}{511.708}{644.793}{R}%
\TextAt{8}{524.903}{644.793}{is}%
\TextAt{8}{79.226}{627.747}{compact.}%
\TextAt{12}{79.226}{602.331}{Theorem}%
\TextAt{12}{136.24}{602.331}{(1.12).}%
\TextAt{8}{178.287}{602.331}{The}%
\TextAt{8}{202.34}{602.331}{group}%
\TextAt{13}{235.516}{602.331}{C}%
\TextAt{13}{244.734}{602.331}{H}%
\TextAt{17}{255.372}{606.669}{1}%
\TextAt{9}{260.105}{602.331}{\Character{40}}%
\Vector{q 0 G 0 g 0.478 w 192.657 -164.313 m 198.501 -164.313 l S Q}%
\TextAt{14}{264.657}{602.331}{o}%
\TextAt{9}{270.501}{602.331}{\Character{41}}%
\TextAt{17}{275.053}{606.669}{0}%
\TextAt{8}{283.683}{602.331}{is}%
\TextAt{8}{295.447}{602.331}{compact.}%
\TextAt{19}{79.226}{568.543}{Proof.}%
\TextAt{8}{116}{568.543}{This}%
\TextAt{8}{142.72}{568.543}{follows}%
\TextAt{8}{181.156}{568.543}{immediately}%
\TextAt{8}{247.818}{568.543}{from}%
\TextAt{8}{275.458}{568.543}{the}%
\TextAt{8}{295.615}{568.543}{exact}%
\TextAt{8}{326.507}{568.543}{sequence}%
\TextAt{9}{192.915}{541.627}{0}%
\NativeArrowAt{202.089}{541.627}{11.955}{19.26}{\longrightarrow}%
\TextAt{13}{224.671}{541.627}{H}%
\TextAt{9}{235.293}{541.627}{\Character{47}\Character{0}}%
\NativeArrowAt{251.781}{541.627}{11.955}{19.26}{\longrightarrow}%
\TextAt{13}{274.363}{541.627}{C}%
\TextAt{13}{283.581}{541.627}{H}%
\TextAt{17}{294.219}{546.563}{1}%
\TextAt{9}{298.952}{541.627}{\Character{40}}%
\Vector{q 0 G 0 g 0.478 w 231.504 -225.017 m 237.348 -225.017 l S Q}%
\TextAt{14}{303.504}{541.627}{o}%
\TextAt{9}{309.348}{541.627}{\Character{41}}%
\TextAt{17}{313.901}{546.563}{0}%
\NativeArrowAt{321.954}{541.627}{11.955}{19.26}{\longrightarrow}%
\TextAt{13}{344.536}{541.627}{C}%
\TextAt{13}{353.754}{541.627}{H}%
\TextAt{17}{364.392}{546.563}{1}%
\TextAt{9}{369.124}{541.627}{\Character{40}}%
\TextAt{14}{373.677}{541.627}{o}%
\TextAt{9}{379.521}{541.627}{\Character{41}}%
\NativeArrowAt{387.394}{541.627}{11.955}{19.26}{\longrightarrow}%
\TextAt{9}{409.976}{541.627}{0}%
\TextAt{13}{415.829}{541.627}{\Character{58}}%
\TextAt{8}{79.226}{514.71}{As}%
\TextAt{9}{96.282}{514.71}{\Character{0}}%
\TextAt{8}{107.262}{514.71}{is}%
\TextAt{8}{118.794}{514.71}{a}%
\TextAt{8}{128.318}{514.71}{complete}%
\TextAt{8}{177.509}{514.71}{lattice}%
\TextAt{8}{213.047}{514.71}{in}%
\TextAt{8}{226.468}{514.71}{the}%
\TextAt{18}{246.393}{514.71}{R}%
\TextAt{8}{255.027}{514.71}{-vector}%
\TextAt{8}{293.817}{514.71}{space}%
\TextAt{13}{324.861}{514.71}{H}%
\TextAt{8}{335.483}{514.71}{,}%
\TextAt{8}{342.447}{514.71}{the}%
\TextAt{8}{362.372}{514.71}{quotient}%
\TextAt{13}{408.323}{514.71}{H}%
\TextAt{9}{418.946}{514.71}{\Character{47}\Character{0}}%
\TextAt{8}{435.779}{514.71}{is}%
\TextAt{8}{447.311}{514.71}{a}%
\TextAt{8}{456.835}{514.71}{compact}%
\TextAt{8}{503.431}{514.71}{torus.}%
\TextAt{8}{79.226}{497.664}{In}%
\TextAt{8}{93.716}{497.664}{view}%
\TextAt{8}{120.567}{497.664}{of}%
\TextAt{8}{133.765}{497.664}{the}%
\TextAt{8}{153.79}{497.664}{finiteness}%
\TextAt{8}{204.503}{497.664}{of}%
\TextAt{13}{217.701}{497.664}{C}%
\TextAt{13}{226.919}{497.664}{H}%
\TextAt{17}{237.558}{502.003}{1}%
\TextAt{9}{242.29}{497.664}{\Character{40}}%
\TextAt{14}{246.842}{497.664}{o}%
\TextAt{9}{252.687}{497.664}{\Character{41}}%
\TextAt{8}{257.239}{497.664}{,}%
\TextAt{8}{264.283}{497.664}{we}%
\TextAt{8}{281.378}{497.664}{obtain}%
\TextAt{13}{317.674}{497.664}{C}%
\TextAt{13}{326.892}{497.664}{H}%
\TextAt{17}{337.53}{502.003}{1}%
\TextAt{9}{342.263}{497.664}{\Character{40}}%
\Vector{q 0 G 0 g 0.478 w 274.815 -268.979 m 280.659 -268.979 l S Q}%
\TextAt{14}{346.815}{497.664}{o}%
\TextAt{9}{352.659}{497.664}{\Character{41}}%
\TextAt{17}{357.211}{502.003}{0}%
\TextAt{8}{365.709}{497.664}{as}%
\TextAt{8}{379.948}{497.664}{the}%
\TextAt{8}{399.972}{497.664}{union}%
\TextAt{8}{432.359}{497.664}{of}%
\TextAt{8}{445.557}{497.664}{the}%
\TextAt{8}{465.582}{497.664}{finitely}%
\TextAt{8}{504.794}{497.664}{many}%
\TextAt{8}{79.226}{480.618}{compact}%
\TextAt{8}{126.055}{480.618}{cosets}%
\TextAt{8}{159.995}{480.618}{of}%
\TextAt{13}{173.325}{480.618}{H}%
\TextAt{9}{183.948}{480.618}{\Character{47}\Character{0}}%
\TextAt{8}{201.012}{480.618}{in}%
\TextAt{13}{214.665}{480.618}{C}%
\TextAt{13}{223.883}{480.618}{H}%
\TextAt{17}{234.521}{484.957}{1}%
\TextAt{9}{239.254}{480.618}{\Character{40}}%
\Vector{q 0 G 0 g 0.478 w 171.806 -286.025 m 177.65 -286.025 l S Q}%
\TextAt{14}{243.806}{480.618}{o}%
\TextAt{9}{249.65}{480.618}{\Character{41}}%
\TextAt{17}{254.203}{484.957}{0}%
\TextAt{8}{258.935}{480.618}{.}%
\TextAt{8}{267.383}{480.618}{Thus}%
\TextAt{13}{297.032}{480.618}{C}%
\TextAt{13}{306.25}{480.618}{H}%
\TextAt{17}{316.888}{484.957}{1}%
\TextAt{9}{321.621}{480.618}{\Character{40}}%
\Vector{q 0 G 0 g 0.478 w 254.173 -286.025 m 260.017 -286.025 l S Q}%
\TextAt{14}{326.173}{480.618}{o}%
\TextAt{9}{332.017}{480.618}{\Character{41}}%
\TextAt{17}{336.57}{484.957}{0}%
\TextAt{8}{345.199}{480.618}{itself}%
\TextAt{8}{373.545}{480.618}{is}%
\TextAt{8}{385.309}{480.618}{compact.}%
\Vector{q 0 G 0 g 0.399 w 452.334 -293.382 m 452.334 -285.312 l S Q}%
\Vector{q 0 G 0 g 0.399 w 452.534 -285.511 m 459.707 -285.511 l S Q}%
\Vector{q 0 G 0 g 0.399 w 452.534 -293.183 m 459.707 -293.183 l S Q}%
\Vector{q 0 G 0 g 0.399 w 459.906 -293.382 m 459.906 -285.312 l S Q}%
\TextAt{8}{97.159}{452.213}{The}%
\TextAt{8}{122.29}{452.213}{correspondence}%
\TextAt{8}{205.044}{452.213}{between}%
\TextAt{8}{251.312}{452.213}{replete}%
\TextAt{8}{290.753}{452.213}{ideals}%
\TextAt{8}{324.408}{452.213}{and}%
\TextAt{8}{348.247}{452.213}{replete}%
\TextAt{8}{387.688}{452.213}{divisors}%
\TextAt{8}{431.493}{452.213}{is}%
\TextAt{8}{444.334}{452.213}{given}%
\TextAt{8}{475.98}{452.213}{by}%
\TextAt{8}{493.316}{452.213}{the}%
\TextAt{8}{514.55}{452.213}{two}%
\TextAt{8}{79.226}{435.167}{mutually}%
\TextAt{8}{128.661}{435.167}{inverse}%
\TextAt{8}{167.419}{435.167}{mappings}%
\TextAt{8}{171.297}{400.927}{div}%
\TextAt{9}{190.554}{400.927}{\Character{58}}%
\TextAt{13}{197.127}{400.927}{J}%
\TextAt{9}{204.704}{400.927}{\Character{40}}%
\Vector{q 0 G 0 g 0.478 w 137.256 -365.716 m 143.1 -365.716 l S Q}%
\TextAt{14}{209.256}{400.927}{o}%
\TextAt{9}{215.1}{400.927}{\Character{41}}%
\NativeArrowAt{222.974}{400.927}{11.955}{19.26}{\longrightarrow}%
\TextAt{8}{245.556}{400.927}{Div}%
\TextAt{9}{263.931}{400.927}{\Character{40}}%
\Vector{q 0 G 0 g 0.478 w 196.483 -365.716 m 202.327 -365.716 l S Q}%
\TextAt{14}{268.483}{400.927}{o}%
\TextAt{9}{274.327}{400.927}{\Character{41}}%
\TextAt{13}{278.88}{400.927}{\Character{59}}%
\TextAt{8}{308.034}{400.927}{div}%
\TextAt{20}{325.963}{421.371}{\Character{32}}%
\TextAt{20}{335.427}{412.285}{\Character{89}}%
\TextAt{27}{341.071}{387.179}{p}%
\TextAt{14}{352.696}{400.927}{p}%
\TextAt{21}{358.679}{405.864}{\Character{23}}%
\TextAt{28}{362.842}{404.867}{p}%
\TextAt{20}{367.423}{421.371}{\Character{33}}%
\TextAt{9}{380.208}{400.927}{\Character{61}}%
\TextAt{20}{392.633}{412.285}{\Character{88}}%
\TextAt{27}{399.273}{387.179}{p}%
\TextAt{16}{411.895}{400.927}{\Character{0}}%
\TextAt{13}{421.193}{400.927}{\Character{23}}%
\TextAt{27}{426.977}{399.134}{p}%
\TextAt{14}{431.464}{400.927}{p}%
\TextAt{13}{437.447}{400.927}{\Character{59}}%
\TextAt{8}{198.919}{361.057}{Div}%
\TextAt{9}{217.294}{361.057}{\Character{40}}%
\Vector{q 0 G 0 g 0.478 w 149.846 -405.587 m 155.69 -405.587 l S Q}%
\TextAt{14}{221.846}{361.057}{o}%
\TextAt{9}{227.69}{361.057}{\Character{41}}%
\NativeArrowAt{235.563}{361.057}{11.955}{19.26}{\longrightarrow}%
\TextAt{13}{258.145}{361.057}{J}%
\TextAt{9}{265.723}{361.057}{\Character{40}}%
\Vector{q 0 G 0 g 0.478 w 198.275 -405.587 m 204.119 -405.587 l S Q}%
\TextAt{14}{270.275}{361.057}{o}%
\TextAt{9}{276.119}{361.057}{\Character{41}}%
\TextAt{13}{280.671}{361.057}{\Character{59}}%
\TextAt{20}{309.826}{372.414}{\Character{88}}%
\TextAt{27}{316.466}{347.308}{p}%
\TextAt{13}{329.087}{361.057}{\Character{23}}%
\TextAt{27}{334.871}{359.264}{p}%
\TextAt{14}{339.358}{361.057}{p}%
\NativeArrowAt{348.662}{361.057}{11.955}{19.26}{\longmapsto}%
\TextAt{20}{371.244}{372.414}{\Character{89}}%
\TextAt{27}{376.887}{347.308}{p}%
\TextAt{14}{388.512}{361.057}{p}%
\TextAt{3}{394.495}{365.993}{\Character{0}}%
\TextAt{21}{401.082}{365.993}{\Character{23}}%
\TextAt{28}{405.245}{364.997}{p}%
\TextAt{13}{409.825}{361.057}{\Character{58}}%
\TextAt{8}{79.226}{329.73}{These}%
\TextAt{8}{113.095}{329.73}{are}%
\TextAt{19}{132.606}{329.73}{topological}%
\TextAt{19}{189.094}{329.73}{isomorphisms}%
\TextAt{8}{262.618}{329.73}{once}%
\TextAt{8}{289.278}{329.73}{we}%
\TextAt{8}{306.506}{329.73}{equip}%
\TextAt{13}{251.799}{302.814}{J}%
\TextAt{9}{259.376}{302.814}{\Character{40}}%
\Vector{q 0 G 0 g 0.478 w 191.928 -463.83 m 197.773 -463.83 l S Q}%
\TextAt{14}{263.928}{302.814}{o}%
\TextAt{9}{269.773}{302.814}{\Character{41}}%
\TextAt{9}{277.628}{302.814}{\Character{61}}%
\TextAt{13}{290.071}{302.814}{J}%
\TextAt{9}{297.648}{302.814}{\Character{40}}%
\TextAt{14}{302.201}{302.814}{o}%
\TextAt{9}{308.045}{302.814}{\Character{41}}%
\TextAt{16}{315.254}{302.814}{\Character{2}}%
\TextAt{20}{327.209}{314.171}{\Character{89}}%
\TextAt{27}{327.442}{288.268}{p}%
\TextAt{3}{331.431}{288.268}{\Character{106}}%
\NativeMathAt{333.782}{288.268}{7.97}{\infty}%
\TextAt{18}{344.478}{302.814}{R}%
\TextAt{3}{353.112}{307.75}{\Character{3}}%
\TextAt{17}{353.112}{299.858}{\Character{43}}%
\TextAt{8}{79.226}{265.247}{with}%
\TextAt{8}{106.567}{265.247}{the}%
\TextAt{8}{127.403}{265.247}{product}%
\TextAt{8}{171.983}{265.247}{topology}%
\TextAt{8}{220.807}{265.247}{of}%
\TextAt{8}{234.817}{265.247}{the}%
\TextAt{8}{255.654}{265.247}{discrete}%
\TextAt{8}{299.313}{265.247}{topology}%
\TextAt{8}{348.137}{265.247}{on}%
\TextAt{13}{365.076}{265.247}{J}%
\TextAt{9}{372.653}{265.247}{\Character{40}}%
\TextAt{14}{377.206}{265.247}{o}%
\TextAt{9}{383.05}{265.247}{\Character{41}}%
\TextAt{8}{392.18}{265.247}{and}%
\TextAt{8}{415.623}{265.247}{the}%
\TextAt{8}{436.46}{265.247}{canonical}%
\TextAt{8}{488.523}{265.247}{topology}%
\TextAt{8}{79.226}{248.201}{on}%
\TextAt{20}{95.485}{257.168}{\Character{81}}%
\TextAt{27}{106.776}{244.714}{p}%
\TextAt{3}{110.765}{244.714}{\Character{106}}%
\NativeMathAt{113.116}{244.714}{7.97}{\infty}%
\TextAt{18}{124.076}{248.201}{R}%
\TextAt{3}{132.711}{252.54}{\Character{3}}%
\TextAt{17}{132.711}{245.246}{\Character{43}}%
\TextAt{8}{139.795}{248.201}{.}%
\TextAt{8}{148.244}{248.201}{The}%
\TextAt{8}{172.297}{248.201}{image}%
\TextAt{8}{206.118}{248.201}{of}%
\TextAt{8}{219.448}{248.201}{a}%
\TextAt{8}{229.204}{248.201}{divisor}%
\TextAt{13}{267.317}{248.201}{D}%
\TextAt{9}{280.682}{248.201}{\Character{61}}%
\TextAt{20}{293.108}{257.168}{\Character{80}}%
\TextAt{27}{305.727}{244.714}{p}%
\TextAt{13}{312.207}{248.201}{\Character{23}}%
\TextAt{27}{317.99}{246.408}{p}%
\TextAt{14}{322.477}{248.201}{p}%
\TextAt{8}{332.358}{248.201}{is}%
\TextAt{8}{344.122}{248.201}{also}%
\TextAt{8}{367.602}{248.201}{written}%
\TextAt{8}{408.572}{248.201}{as}%
\TextAt{14}{264.711}{219.302}{o}%
\TextAt{9}{270.556}{219.302}{\Character{40}}%
\TextAt{13}{275.108}{219.302}{D}%
\TextAt{9}{285.152}{219.302}{\Character{41}}%
\TextAt{9}{293.007}{219.302}{\Character{61}}%
\TextAt{20}{305.451}{230.659}{\Character{89}}%
\TextAt{27}{311.095}{205.553}{p}%
\TextAt{14}{322.72}{219.302}{p}%
\TextAt{3}{328.703}{224.238}{\Character{0}}%
\TextAt{21}{335.289}{224.238}{\Character{23}}%
\TextAt{28}{339.452}{223.242}{p}%
\TextAt{13}{344.033}{219.302}{\Character{58}}%
\TextAt{8}{79.226}{183.568}{The}%
\TextAt{8}{103.178}{183.568}{minus}%
\TextAt{8}{137.281}{183.568}{sign}%
\TextAt{8}{161.305}{183.568}{here}%
\TextAt{8}{186.56}{183.568}{is}%
\TextAt{8}{198.223}{183.568}{motivated}%
\TextAt{8}{253.091}{183.568}{by}%
\TextAt{8}{269.249}{183.568}{classical}%
\TextAt{8}{314.147}{183.568}{usage}%
\TextAt{8}{345.977}{183.568}{in}%
\TextAt{8}{359.529}{183.568}{function}%
\TextAt{8}{405.275}{183.568}{theory.}%
\TextAt{8}{445.562}{183.568}{Replete}%
\TextAt{8}{487.89}{183.568}{principal}%
\TextAt{8}{79.226}{166.522}{ideals}%
\TextAt{8}{112.926}{166.522}{correspond}%
\TextAt{8}{173.62}{166.522}{to}%
\TextAt{8}{189.052}{166.522}{replete}%
\TextAt{8}{228.538}{166.522}{principal}%
\TextAt{8}{278.436}{166.522}{divisors}%
\TextAt{8}{322.286}{166.522}{in}%
\TextAt{8}{337.06}{166.522}{such}%
\TextAt{8}{364.579}{166.522}{a}%
\TextAt{8}{375.456}{166.522}{way}%
\TextAt{8}{400.32}{166.522}{that}%
\TextAt{13}{426.811}{166.522}{P}%
\TextAt{9}{435.982}{166.522}{\Character{40}}%
\Vector{q 0 G 0 g 0.478 w 368.534 -600.122 m 374.379 -600.122 l S Q}%
\TextAt{14}{440.534}{166.522}{o}%
\TextAt{9}{446.379}{166.522}{\Character{41}}%
\TextAt{8}{455.95}{166.522}{becomes}%
\TextAt{8}{503.625}{166.522}{a}%
\TextAt{8}{514.502}{166.522}{dis-}%
\TextAt{8}{79.226}{149.476}{crete}%
\TextAt{8}{108.2}{149.476}{subgroup}%
\TextAt{8}{159.362}{149.476}{of}%
\TextAt{13}{173.057}{149.476}{J}%
\TextAt{9}{180.634}{149.476}{\Character{40}}%
\Vector{q 0 G 0 g 0.478 w 113.187 -617.168 m 119.031 -617.168 l S Q}%
\TextAt{14}{185.187}{149.476}{o}%
\TextAt{9}{191.031}{149.476}{\Character{41}}%
\TextAt{8}{199.845}{149.476}{by}%
\TextAt{8}{216.469}{149.476}{(1.9),}%
\TextAt{8}{248.152}{149.476}{and}%
\TextAt{8}{271.28}{149.476}{Pic}%
\TextAt{9}{287.694}{149.476}{\Character{40}}%
\Vector{q 0 G 0 g 0.478 w 220.247 -617.168 m 226.091 -617.168 l S Q}%
\TextAt{14}{292.247}{149.476}{o}%
\TextAt{9}{298.091}{149.476}{\Character{41}}%
\TextAt{9}{306.567}{149.476}{\Character{61}}%
\TextAt{13}{319.633}{149.476}{J}%
\TextAt{9}{327.211}{149.476}{\Character{40}}%
\Vector{q 0 G 0 g 0.478 w 259.763 -617.168 m 265.607 -617.168 l S Q}%
\TextAt{14}{331.763}{149.476}{o}%
\TextAt{9}{337.607}{149.476}{\Character{41}\Character{47}}%
\TextAt{13}{348.012}{149.476}{P}%
\TextAt{9}{357.184}{149.476}{\Character{40}}%
\Vector{q 0 G 0 g 0.478 w 289.736 -617.168 m 295.58 -617.168 l S Q}%
\TextAt{14}{361.736}{149.476}{o}%
\TextAt{9}{367.58}{149.476}{\Character{41}}%
\TextAt{8}{376.395}{149.476}{is}%
\TextAt{8}{388.524}{149.476}{a}%
\TextAt{8}{398.644}{149.476}{locally}%
\TextAt{8}{436.082}{149.476}{compact}%
\TextAt{8}{483.275}{149.476}{Hausdorff}%
\TextAt{8}{79.226}{132.43}{topological}%
\TextAt{8}{138.751}{132.43}{group.}%
\TextAt{8}{176.477}{132.43}{We}%
\TextAt{8}{196.622}{132.43}{obtain}%
\TextAt{8}{233.049}{132.43}{the}%
\TextAt{8}{253.206}{132.43}{following}%
\TextAt{8}{302.64}{132.43}{extension}%
\TextAt{8}{354.406}{132.43}{of}%
\TextAt{8}{367.736}{132.43}{chap.}%
\TextAt{8}{399.928}{132.43}{I,}%
\TextAt{8}{411.297}{132.43}{(12.14).}%
\TextAt{12}{79.226}{107.013}{Proposition}%
\TextAt{12}{152.895}{107.013}{(1.13).}%
\TextAt{8}{195.774}{107.013}{The}%
\TextAt{8}{220.661}{107.013}{mapping}%
\TextAt{8}{269.625}{107.013}{div}%
\TextAt{9}{290.301}{107.013}{\Character{58}}%
\TextAt{13}{298.293}{107.013}{J}%
\TextAt{9}{305.87}{107.013}{\Character{40}}%
\Vector{q 0 G 0 g 0.478 w 238.422 -659.63 m 244.266 -659.63 l S Q}%
\TextAt{14}{310.422}{107.013}{o}%
\TextAt{9}{316.266}{107.013}{\Character{41}}%
\NativeMathAt{328.146}{113.79}{7.97}{\sim}%
\NativeArrowAt{325.558}{107.013}{11.955}{13.175}{\longrightarrow}%
\TextAt{8}{343.471}{107.013}{Div}%
\TextAt{9}{361.846}{107.013}{\Character{40}}%
\Vector{q 0 G 0 g 0.478 w 294.399 -659.63 m 300.243 -659.63 l S Q}%
\TextAt{14}{366.399}{107.013}{o}%
\TextAt{9}{372.243}{107.013}{\Character{41}}%
\TextAt{8}{381.525}{107.013}{induces}%
\TextAt{8}{424.034}{107.013}{a}%
\TextAt{8}{434.622}{107.013}{topological}%
\TextAt{8}{494.979}{107.013}{isomor-}%
\TextAt{8}{79.226}{89.967}{phism}%
\TextAt{8}{246.099}{72.921}{div}%
\TextAt{9}{265.357}{72.921}{\Character{58}}%
\TextAt{8}{271.929}{72.921}{Pic}%
\TextAt{9}{288.344}{72.921}{\Character{40}}%
\Vector{q 0 G 0 g 0.478 w 220.896 -693.722 m 226.74 -693.722 l S Q}%
\TextAt{14}{292.896}{72.921}{o}%
\TextAt{9}{298.74}{72.921}{\Character{41}}%
\NativeMathAt{309.201}{79.698}{7.97}{\sim}%
\NativeArrowAt{306.613}{72.921}{11.955}{13.175}{\longrightarrow}%
\TextAt{13}{323.107}{72.921}{C}%
\TextAt{13}{332.325}{72.921}{H}%
\TextAt{17}{342.963}{77.858}{1}%
\TextAt{9}{347.696}{72.921}{\Character{40}}%
\Vector{q 0 G 0 g 0.478 w 280.248 -693.722 m 286.092 -693.722 l S Q}%
\TextAt{14}{352.248}{72.921}{o}%
\TextAt{9}{358.092}{72.921}{\Character{41}}%
\TextAt{13}{362.644}{72.921}{\Character{58}}%
\TextAt{8}{297.211}{43.033}{175}%
\Destination{78.23}{809.815}{nav.page.175}%
\Destination{79.36}{467.255}{nav.print.192}%
\Destination{79.36}{616.355}{nav.number.3.1.12}%
\Destination{79.36}{121.055}{nav.number.3.1.13}%
%
%
\Hotlink{213.784}{758.955}{19.22}{11.047}{nav.cite.N023}%
\Hotlink{399.428}{130.255}{5.22}{11.047}{nav.chapter.1}%
%
\Hotlink{215.969}{147.255}{25.078}{11.047}{nav.number.3.1.9}%
\Hotlink{410.797}{130.255}{36.794}{11.047}{nav.number.1.12.14}%
}

% Page 188 (body page 176)
\Page{%
\TextAt{2}{79.226}{796.943}{III.}%
\TextAt{2}{98.317}{796.943}{Riemann–Roch}%
\TextAt{2}{174.375}{796.943}{Theory}%
\Vector{q 0 G 0 g 0.299 w 7.226 18.729 m 460.769 18.729 l S Q}%
\TextAt{8}{97.159}{761.124}{On}%
\TextAt{8}{117.562}{761.124}{the}%
\TextAt{8}{138.61}{761.124}{group}%
\TextAt{13}{172.678}{761.124}{J}%
\TextAt{9}{180.255}{761.124}{\Character{40}}%
\Vector{q 0 G 0 g 0.478 w 112.807 -5.52 m 118.651 -5.52 l S Q}%
\TextAt{14}{184.807}{761.124}{o}%
\TextAt{9}{190.651}{761.124}{\Character{41}}%
\TextAt{8}{199.993}{761.124}{we}%
\TextAt{8}{218.112}{761.124}{have}%
\TextAt{8}{245.999}{761.124}{the}%
\TextAt{8}{267.048}{761.124}{homomorphism}%
\TextAt{14}{350.61}{761.124}{N}%
\TextAt{9}{365.399}{761.124}{\Character{58}}%
\TextAt{13}{373.491}{761.124}{J}%
\TextAt{9}{381.068}{761.124}{\Character{40}}%
\Vector{q 0 G 0 g 0.478 w 313.621 -5.52 m 319.465 -5.52 l S Q}%
\TextAt{14}{385.621}{761.124}{o}%
\TextAt{9}{391.465}{761.124}{\Character{41}}%
\NativeMathAt{400.858}{761.124}{11.955}{\rightarrow}%
\TextAt{18}{417.654}{761.124}{R}%
\TextAt{3}{426.288}{765.462}{\Character{3}}%
\TextAt{17}{426.288}{758.168}{\Character{43}}%
\TextAt{8}{433.373}{761.124}{,}%
\TextAt{8}{441.637}{761.124}{and}%
\TextAt{8}{465.292}{761.124}{on}%
\TextAt{8}{482.443}{761.124}{the}%
\TextAt{8}{503.491}{761.124}{group}%
\TextAt{8}{79.226}{744.078}{Div}%
\TextAt{9}{97.601}{744.078}{\Character{40}}%
\Vector{q 0 G 0 g 0.478 w 30.154 -22.566 m 35.998 -22.566 l S Q}%
\TextAt{14}{102.154}{744.078}{o}%
\TextAt{9}{107.998}{744.078}{\Character{41}}%
\TextAt{8}{117.53}{744.078}{there}%
\TextAt{8}{148.525}{744.078}{is}%
\TextAt{8}{161.371}{744.078}{the}%
\TextAt{8}{182.61}{744.078}{degree}%
\TextAt{8}{220.109}{744.078}{map}%
\TextAt{8}{247.206}{744.078}{deg}%
\TextAt{9}{269.934}{744.078}{\Character{58}}%
\TextAt{8}{278.351}{744.078}{Div}%
\TextAt{9}{296.726}{744.078}{\Character{40}}%
\Vector{q 0 G 0 g 0.478 w 229.279 -22.566 m 235.123 -22.566 l S Q}%
\TextAt{14}{301.279}{744.078}{o}%
\TextAt{9}{307.123}{744.078}{\Character{41}}%
\NativeMathAt{316.841}{744.078}{11.955}{\rightarrow}%
\TextAt{18}{333.962}{744.078}{R}%
\TextAt{8}{342.597}{744.078}{.}%
\TextAt{8}{354.293}{744.078}{They}%
\TextAt{8}{385.61}{744.078}{are}%
\TextAt{8}{406.204}{744.078}{obviously}%
\TextAt{8}{459.064}{744.078}{related}%
\TextAt{8}{499.169}{744.078}{by}%
\TextAt{8}{516.51}{744.078}{the}%
\TextAt{8}{79.226}{727.032}{formula}%
\TextAt{8}{241.01}{709.986}{deg}%
\TextAt{9}{258.572}{709.986}{\Character{40}}%
\TextAt{8}{263.124}{709.986}{div}%
\TextAt{9}{279.061}{709.986}{\Character{40}}%
\TextAt{14}{283.613}{709.986}{a}%
\TextAt{9}{289.596}{709.986}{\Character{41}\Character{41}}%
\TextAt{9}{301.994}{709.986}{\Character{61}}%
\TextAt{16}{314.447}{709.986}{\Character{0}}%
\TextAt{8}{325.738}{709.986}{log}%
\TextAt{14}{342.698}{709.986}{N}%
\TextAt{9}{352.646}{709.986}{\Character{40}}%
\TextAt{14}{357.199}{709.986}{a}%
\TextAt{9}{363.182}{709.986}{\Character{41}}%
\TextAt{13}{367.734}{709.986}{\Character{59}}%
\TextAt{8}{79.226}{687.254}{and}%
\TextAt{8}{101.989}{687.254}{we}%
\TextAt{8}{119.216}{687.254}{get}%
\TextAt{8}{138.727}{687.254}{a}%
\TextAt{8}{148.482}{687.254}{commutative}%
\TextAt{8}{218.408}{687.254}{diagram}%
\TextAt{9}{201.077}{659.803}{0}%
\TextAt{8}{235.669}{659.803}{Pic}%
\TextAt{9}{252.083}{659.803}{\Character{40}}%
\NativeDiagram{%
\path[draw=black,line width=0.478bp] (184.636,-106.841) -- (190.48,-106.841);%
}%
\TextAt{14}{256.636}{659.803}{o}%
\TextAt{9}{262.48}{659.803}{\Character{41}}%
\TextAt{17}{267.032}{664.141}{0}%
\TextAt{8}{304.59}{659.803}{Pic}%
\TextAt{9}{321.004}{659.803}{\Character{40}}%
\NativeDiagram{%
\path[draw=black,line width=0.478bp] (253.557,-106.841) -- (259.401,-106.841);%
}%
\TextAt{14}{325.557}{659.803}{o}%
\TextAt{9}{331.401}{659.803}{\Character{41}}%
\TextAt{18}{364.692}{659.803}{R}%
\TextAt{3}{373.326}{664.141}{\Character{3}}%
\TextAt{17}{373.326}{656.847}{\Character{43}}%
\TextAt{9}{405.062}{659.803}{0}%
\TextAt{9}{201.077}{616.263}{0}%
\TextAt{13}{231.582}{616.263}{C}%
\TextAt{13}{240.8}{616.263}{H}%
\TextAt{17}{251.438}{620.602}{1}%
\TextAt{9}{256.17}{616.263}{\Character{40}}%
\NativeDiagram{%
\path[draw=black,line width=0.478bp] (188.723,-150.38) -- (194.567,-150.38);%
}%
\TextAt{14}{260.723}{616.263}{o}%
\TextAt{9}{266.567}{616.263}{\Character{41}}%
\TextAt{17}{271.119}{620.602}{0}%
\TextAt{13}{300.503}{616.263}{C}%
\TextAt{13}{309.721}{616.263}{H}%
\TextAt{17}{320.359}{620.602}{1}%
\TextAt{9}{325.091}{616.263}{\Character{40}}%
\NativeDiagram{%
\path[draw=black,line width=0.478bp] (257.644,-150.38) -- (263.488,-150.38);%
}%
\TextAt{14}{329.644}{616.263}{o}%
\TextAt{9}{335.488}{616.263}{\Character{41}}%
\TextAt{18}{368.235}{616.263}{R}%
\TextAt{9}{405.062}{616.263}{0}%
\NativeDiagram{%
\path[draw=black,line width=0.478bp] (140.283,-111.208) -- (157.839,-111.208);%
\path[draw=black,line width=0.478bp,line cap=round,line join=round] (155.592,-108.339) .. controls (156.046,-110.061) and (157.059,-110.873) .. (158.079,-111.208) .. controls (157.059,-111.543) and (156.046,-112.356) .. (155.592,-114.077);%
\path[draw=black,line width=0.478bp] (205.118,-111.208) -- (226.762,-111.208);%
\path[draw=black,line width=0.478bp,line cap=round,line join=round] (224.514,-108.339) .. controls (224.968,-110.061) and (225.981,-110.873) .. (227.001,-111.208) .. controls (225.981,-111.543) and (224.968,-112.356) .. (224.514,-114.077);%
\path[draw=black,line width=0.478bp] (181.718,-121.765) -- (181.718,-143.205);%
\path[draw=black,line width=0.478bp,line cap=round,line join=round] (184.587,-140.957) .. controls (182.865,-141.412) and (182.052,-142.424) .. (181.718,-143.444) .. controls (181.383,-142.424) and (180.57,-141.412) .. (178.849,-140.957);%
}%
\NativeMathAt{256.533}{641.028}{7.97}{\sim}%
\TextAt{17}{256.533}{638.513}{\Character{61}}%
\NativeDiagram{%
\path[draw=black,line width=0.478bp] (269.307,-111.208) -- (286.864,-111.208);%
\path[draw=black,line width=0.478bp,line cap=round,line join=round] (284.617,-108.339) .. controls (285.071,-110.061) and (286.083,-110.873) .. (287.103,-111.208) .. controls (286.083,-111.543) and (285.071,-112.356) .. (284.617,-114.077);%
}%
\TextAt{27}{347.009}{665.607}{N}%
\NativeDiagram{%
\path[draw=black,line width=0.478bp] (248.274,-121.765) -- (248.274,-143.205);%
\path[draw=black,line width=0.478bp,line cap=round,line join=round] (251.143,-140.957) .. controls (249.421,-141.412) and (248.608,-142.424) .. (248.274,-143.444) .. controls (247.939,-142.424) and (247.126,-141.412) .. (245.404,-140.957);%
}%
\TextAt{7}{323.089}{638.55}{div}%
\NativeMathAt{310.872}{641.028}{7.97}{\sim}%
\TextAt{17}{310.872}{638.513}{\Character{61}}%
\NativeDiagram{%
\path[draw=black,line width=0.478bp] (313.765,-111.208) -- (327.235,-111.208);%
\path[draw=black,line width=0.478bp,line cap=round,line join=round] (324.988,-108.339) .. controls (325.442,-110.061) and (326.454,-110.873) .. (327.474,-111.208) .. controls (326.454,-111.543) and (325.442,-112.356) .. (324.988,-114.077);%
\path[draw=black,line width=0.478bp] (300.554,-122.562) -- (300.554,-144.444);%
\path[draw=black,line width=0.478bp,line cap=round,line join=round] (303.423,-142.196) .. controls (301.702,-142.651) and (300.889,-143.663) .. (300.554,-144.683) .. controls (300.219,-143.663) and (299.406,-142.651) .. (297.685,-142.196);%
}%
\TextAt{3}{375.369}{638.309}{\Character{0}}%
\TextAt{7}{383.367}{638.309}{log}%
\NativeMathAt{363.152}{640.01}{7.97}{\sim}%
\TextAt{17}{363.152}{637.495}{\Character{61}}%
\NativeDiagram{%
\path[draw=black,line width=0.478bp] (140.283,-154.748) -- (153.752,-154.748);%
\path[draw=black,line width=0.478bp,line cap=round,line join=round] (151.505,-151.879) .. controls (151.959,-153.601) and (152.971,-154.413) .. (153.991,-154.748) .. controls (152.971,-155.083) and (151.959,-155.896) .. (151.505,-157.617);%
\path[draw=black,line width=0.478bp] (209.205,-154.748) -- (222.674,-154.748);%
\path[draw=black,line width=0.478bp,line cap=round,line join=round] (220.427,-151.879) .. controls (220.881,-153.601) and (221.894,-154.413) .. (222.913,-154.748) .. controls (221.894,-155.083) and (220.881,-155.896) .. (220.427,-157.617);%
\path[draw=black,line width=0.478bp] (273.395,-154.748) -- (290.407,-154.748);%
\path[draw=black,line width=0.478bp,line cap=round,line join=round] (288.159,-151.879) .. controls (288.613,-153.601) and (289.626,-154.413) .. (290.646,-154.748) .. controls (289.626,-155.083) and (288.613,-155.896) .. (288.159,-157.617);%
}%
\TextAt{7}{347.791}{623.701}{deg}%
\NativeDiagram{%
\path[draw=black,line width=0.478bp] (310.223,-154.748) -- (327.235,-154.748);%
\path[draw=black,line width=0.478bp,line cap=round,line join=round] (324.988,-151.879) .. controls (325.442,-153.601) and (326.454,-154.413) .. (327.474,-154.748) .. controls (326.454,-155.083) and (325.442,-155.896) .. (324.988,-157.617);%
}%
\TextAt{8}{79.226}{588.82}{with}%
\TextAt{8}{105.886}{588.82}{exact}%
\TextAt{8}{136.778}{588.82}{rows.}%
\TextAt{8}{168.384}{588.82}{(1.12)}%
\TextAt{8}{202.217}{588.82}{now}%
\TextAt{8}{226.606}{588.82}{yields}%
\TextAt{8}{259.506}{588.82}{the}%
\TextAt{12}{79.226}{563.334}{Corollary}%
\TextAt{12}{139.014}{563.334}{(1.14).}%
\TextAt{8}{181.06}{563.334}{The}%
\TextAt{8}{205.114}{563.334}{group}%
\TextAt{8}{216.57}{536.176}{Pic}%
\TextAt{9}{232.984}{536.176}{\Character{40}}%
\Vector{q 0 G 0 g 0.478 w 165.536 -230.468 m 171.381 -230.468 l S Q}%
\TextAt{14}{237.536}{536.176}{o}%
\TextAt{9}{243.381}{536.176}{\Character{41}}%
\TextAt{17}{247.933}{541.112}{0}%
\TextAt{9}{255.986}{536.176}{\Character{61}}%
\TextAt{16}{268.412}{536.176}{\Character{102}}%
\TextAt{9}{274.389}{536.176}{\Character{91}}%
\TextAt{14}{277.641}{536.176}{a}%
\TextAt{9}{283.624}{536.176}{\Character{93}}%
\NativeMathAt{290.196}{536.176}{11.955}{\in}%
\TextAt{8}{301.487}{536.176}{Pic}%
\TextAt{9}{317.902}{536.176}{\Character{40}}%
\Vector{q 0 G 0 g 0.478 w 250.454 -230.468 m 256.298 -230.468 l S Q}%
\TextAt{14}{322.454}{536.176}{o}%
\TextAt{9}{328.298}{536.176}{\Character{41}}%
\TextAt{16}{336.172}{536.176}{\Character{106}}%
\TextAt{14}{342.813}{536.176}{N}%
\TextAt{9}{352.761}{536.176}{\Character{40}}%
\TextAt{14}{357.314}{536.176}{a}%
\TextAt{9}{363.297}{536.176}{\Character{41}}%
\TextAt{9}{371.152}{536.176}{\Character{61}}%
\TextAt{9}{383.573}{536.176}{1}%
\TextAt{16}{389.448}{536.176}{\Character{103}}%
\TextAt{8}{79.226}{509.017}{is}%
\TextAt{8}{90.99}{509.017}{compact.}%
\TextAt{8}{97.159}{483.531}{The}%
\TextAt{8}{120.064}{483.531}{preceding}%
\TextAt{8}{171.59}{483.531}{isomorphism}%
\TextAt{8}{238.861}{483.531}{result}%
\TextAt{8}{270.29}{483.531}{(1.13)}%
\TextAt{8}{302.974}{483.531}{invites}%
\TextAt{8}{338.958}{483.531}{a}%
\TextAt{8}{347.565}{483.531}{philosophical}%
\TextAt{8}{416.725}{483.531}{reflection.}%
\TextAt{8}{471.619}{483.531}{Throughout}%
\TextAt{8}{79.226}{466.485}{the}%
\TextAt{8}{100.142}{466.485}{historical}%
\TextAt{8}{151.699}{466.485}{development}%
\TextAt{8}{220.423}{466.485}{of}%
\TextAt{8}{234.512}{466.485}{algebraic}%
\TextAt{8}{284.706}{466.485}{number}%
\TextAt{8}{328.062}{466.485}{theory,}%
\TextAt{8}{368.033}{466.485}{a}%
\TextAt{8}{378.547}{466.485}{controversy}%
\TextAt{8}{441.498}{466.485}{persisted}%
\TextAt{8}{491.476}{466.485}{between}%
\TextAt{8}{79.226}{449.439}{the}%
\TextAt{8}{99.372}{449.439}{followers}%
\TextAt{8}{147.229}{449.439}{of}%
\TextAt{8}{160.548}{449.439}{Dedekind’s}%
\TextAt{8}{220.588}{449.439}{ideal-theoretic}%
\TextAt{8}{297.317}{449.439}{approach,}%
\TextAt{8}{350.975}{449.439}{and}%
\TextAt{8}{373.727}{449.439}{the}%
\TextAt{8}{393.872}{449.439}{divisor-theoretic}%
\TextAt{8}{480.752}{449.439}{method}%
\TextAt{8}{523.337}{449.439}{of}%
\TextAt{8}{79.226}{432.393}{building}%
\TextAt{8}{125.189}{432.393}{up}%
\TextAt{8}{142.531}{432.393}{the}%
\TextAt{8}{163.125}{432.393}{theory}%
\TextAt{8}{200.313}{432.393}{from}%
\TextAt{8}{228.391}{432.393}{the}%
\TextAt{8}{248.984}{432.393}{valuation-theoretic}%
\TextAt{8}{349.272}{432.393}{notion}%
\TextAt{8}{386.136}{432.393}{of}%
\TextAt{8}{399.904}{432.393}{primes.}%
\TextAt{8}{443.546}{432.393}{Both}%
\TextAt{8}{473.082}{432.393}{theories}%
\TextAt{8}{517.156}{432.393}{are}%
\TextAt{8}{79.226}{415.347}{equivalent}%
\TextAt{8}{134.78}{415.347}{in}%
\TextAt{8}{148.371}{415.347}{the}%
\TextAt{8}{168.466}{415.347}{sense}%
\TextAt{8}{198.436}{415.347}{of}%
\TextAt{8}{211.704}{415.347}{the}%
\TextAt{8}{231.799}{415.347}{above}%
\TextAt{8}{264.913}{415.347}{isomorphism}%
\TextAt{8}{333.271}{415.347}{result,}%
\TextAt{8}{369.051}{415.347}{but}%
\TextAt{8}{390.449}{415.347}{they}%
\TextAt{8}{416.724}{415.347}{are}%
\TextAt{8}{436.174}{415.347}{also}%
\TextAt{8}{459.592}{415.347}{fundamentally}%
\TextAt{8}{79.226}{398.301}{different}%
\TextAt{8}{125.555}{398.301}{in}%
\TextAt{8}{139.366}{398.301}{nature.}%
\TextAt{8}{181.464}{398.301}{The}%
\TextAt{8}{205.676}{398.301}{controversy}%
\TextAt{8}{268.025}{398.301}{has}%
\TextAt{8}{289.057}{398.301}{finally}%
\TextAt{8}{324.662}{398.301}{given}%
\TextAt{8}{355.389}{398.301}{way}%
\TextAt{8}{379.29}{398.301}{to}%
\TextAt{8}{393.759}{398.301}{the}%
\TextAt{8}{414.073}{398.301}{realization}%
\TextAt{8}{471.473}{398.301}{that}%
\TextAt{8}{497}{398.301}{neither}%
\TextAt{8}{79.226}{381.255}{approach}%
\TextAt{8}{130.434}{381.255}{is}%
\TextAt{8}{142.99}{381.255}{dominant,}%
\TextAt{8}{199.596}{381.255}{that}%
\TextAt{8}{225.758}{381.255}{each}%
\TextAt{8}{252.888}{381.255}{one}%
\TextAt{8}{275.14}{381.255}{instead}%
\TextAt{8}{316.317}{381.255}{emanates}%
\TextAt{8}{368.553}{381.255}{from}%
\TextAt{8}{396.986}{381.255}{its}%
\TextAt{8}{414.097}{381.255}{own}%
\TextAt{8}{439.278}{381.255}{proper}%
\TextAt{8}{477.466}{381.255}{world,}%
\TextAt{8}{513.904}{381.255}{and}%
\TextAt{8}{79.226}{364.209}{that}%
\TextAt{8}{104.833}{364.209}{the}%
\TextAt{8}{125.228}{364.209}{relation}%
\TextAt{8}{168.397}{364.209}{between}%
\TextAt{8}{213.825}{364.209}{these}%
\TextAt{8}{244.035}{364.209}{worlds}%
\TextAt{8}{281.083}{364.209}{is}%
\TextAt{8}{293.085}{364.209}{expressed}%
\TextAt{8}{345.795}{364.209}{by}%
\TextAt{8}{362.292}{364.209}{an}%
\TextAt{8}{378.789}{364.209}{important}%
\TextAt{8}{434.32}{364.209}{mathematical}%
\TextAt{8}{508.058}{364.209}{prin-}%
\TextAt{8}{79.226}{347.163}{ciple.}%
\TextAt{8}{111.76}{347.163}{However,}%
\TextAt{8}{162.446}{347.163}{all}%
\TextAt{8}{178.931}{347.163}{this}%
\TextAt{8}{201.979}{347.163}{becomes}%
\TextAt{8}{248.758}{347.163}{evident}%
\TextAt{8}{289.954}{347.163}{only}%
\TextAt{8}{315.872}{347.163}{in}%
\TextAt{8}{329.751}{347.163}{higher}%
\TextAt{8}{365.746}{347.163}{dimensional}%
\TextAt{8}{430.423}{347.163}{arithmetic}%
\TextAt{8}{487.232}{347.163}{algebraic}%
\TextAt{8}{79.226}{330.117}{geometry.}%
\TextAt{8}{134.682}{330.117}{There,}%
\TextAt{8}{172.087}{330.117}{on}%
\TextAt{8}{188.621}{330.117}{an}%
\TextAt{8}{205.155}{330.117}{algebraic}%
\TextAt{18}{254.864}{330.117}{Z}%
\TextAt{8}{262.834}{330.117}{-scheme}%
\TextAt{13}{307.056}{330.117}{X}%
\TextAt{8}{317.711}{330.117}{,}%
\TextAt{8}{325.204}{330.117}{one}%
\TextAt{8}{346.939}{330.117}{studies}%
\TextAt{8}{386.355}{330.117}{on}%
\TextAt{8}{402.889}{330.117}{the}%
\TextAt{8}{423.32}{330.117}{one}%
\TextAt{8}{445.055}{330.117}{hand}%
\TextAt{8}{474.596}{330.117}{the}%
\TextAt{8}{495.027}{330.117}{totality}%
\TextAt{8}{79.226}{313.071}{of}%
\TextAt{8}{92.452}{313.071}{all}%
\TextAt{19}{108.608}{313.071}{vector}%
\TextAt{19}{143.023}{313.071}{bundles}%
\TextAt{8}{180.98}{313.071}{,}%
\TextAt{8}{188.047}{313.071}{and}%
\TextAt{8}{210.705}{313.071}{on}%
\TextAt{8}{226.861}{313.071}{the}%
\TextAt{8}{246.913}{313.071}{other,}%
\TextAt{8}{280.651}{313.071}{that}%
\TextAt{8}{305.917}{313.071}{of}%
\TextAt{8}{319.143}{313.071}{all}%
\TextAt{19}{335.298}{313.071}{irreducible}%
\TextAt{19}{391.986}{313.071}{subschemes}%
\TextAt{8}{453.594}{313.071}{of}%
\TextAt{13}{466.821}{313.071}{X}%
\TextAt{8}{477.476}{313.071}{.}%
\TextAt{8}{485.889}{313.071}{From}%
\TextAt{8}{516.51}{313.071}{the}%
\TextAt{8}{79.226}{296.025}{first,}%
\TextAt{8}{107.723}{296.025}{one}%
\TextAt{8}{130.078}{296.025}{constructs}%
\TextAt{8}{187.031}{296.025}{a}%
\TextAt{8}{197.681}{296.025}{series}%
\TextAt{8}{229.911}{296.025}{of}%
\TextAt{8}{244.137}{296.025}{groups}%
\TextAt{13}{282.822}{296.025}{K}%
\TextAt{21}{292.781}{294.232}{i}%
\TextAt{9}{296.163}{296.025}{\Character{40}}%
\TextAt{13}{300.715}{296.025}{X}%
\TextAt{9}{311.37}{296.025}{\Character{41}}%
\TextAt{8}{320.715}{296.025}{which}%
\TextAt{8}{355.097}{296.025}{constitute}%
\TextAt{8}{410.687}{296.025}{the}%
\TextAt{8}{431.739}{296.025}{subject}%
\TextAt{8}{473.007}{296.025}{of}%
\TextAt{8}{487.232}{296.025}{algebraic}%
\TextAt{13}{79.226}{278.979}{K}%
\TextAt{12}{90.044}{278.979}{-theory}%
\TextAt{8}{132.7}{278.979}{.}%
\TextAt{8}{144.243}{278.979}{From}%
\TextAt{8}{175.999}{278.979}{the}%
\TextAt{8}{197.188}{278.979}{second}%
\TextAt{8}{235.997}{278.979}{is}%
\TextAt{8}{248.793}{278.979}{constructed}%
\TextAt{8}{312.972}{278.979}{a}%
\TextAt{8}{323.759}{278.979}{series}%
\TextAt{8}{356.125}{278.979}{of}%
\TextAt{8}{370.487}{278.979}{groups}%
\TextAt{13}{409.309}{278.979}{C}%
\TextAt{13}{418.526}{278.979}{H}%
\TextAt{21}{429.165}{283.318}{i}%
\TextAt{9}{432.546}{278.979}{\Character{40}}%
\TextAt{13}{437.098}{278.979}{X}%
\TextAt{9}{447.754}{278.979}{\Character{41}}%
\TextAt{8}{452.306}{278.979}{,}%
\TextAt{8}{460.745}{278.979}{the}%
\TextAt{8}{481.933}{278.979}{subject}%
\TextAt{8}{523.337}{278.979}{of}%
\TextAt{12}{79.226}{261.934}{Chow}%
\TextAt{12}{117.044}{261.934}{theory}%
\TextAt{8}{155.217}{261.934}{.}%
\TextAt{8}{163.83}{261.934}{Vector}%
\TextAt{8}{200.946}{261.934}{bundles}%
\TextAt{8}{243.98}{261.934}{are}%
\TextAt{8}{263.546}{261.934}{by}%
\TextAt{8}{279.86}{261.934}{definition}%
\TextAt{8}{331.943}{261.934}{locally}%
\TextAt{8}{369.071}{261.934}{free}%
\TextAt{16}{391.554}{261.934}{\Character{79}}%
\TextAt{21}{401.072}{260.14}{X}%
\TextAt{8}{409.157}{261.934}{-modules.}%
\TextAt{8}{463.678}{261.934}{In}%
\TextAt{8}{478.354}{261.934}{the}%
\TextAt{8}{498.566}{261.934}{special}%
\TextAt{8}{79.226}{244.888}{case}%
\TextAt{13}{105.272}{244.888}{X}%
\TextAt{9}{121.419}{244.888}{\Character{61}}%
\TextAt{8}{136.017}{244.888}{Spec}%
\TextAt{9}{159.425}{244.888}{\Character{40}}%
\TextAt{14}{163.977}{244.888}{o}%
\TextAt{9}{169.822}{244.888}{\Character{41}}%
\TextAt{8}{179.546}{244.888}{this}%
\TextAt{8}{203.643}{244.888}{includes}%
\TextAt{8}{249.845}{244.888}{the}%
\TextAt{8}{271.275}{244.888}{fractional}%
\TextAt{8}{324.914}{244.888}{ideals.}%
\TextAt{8}{365.866}{244.888}{The}%
\TextAt{8}{391.194}{244.888}{irreducible}%
\TextAt{8}{450.343}{244.888}{subschemes}%
\TextAt{8}{513.904}{244.888}{and}%
\TextAt{8}{79.226}{227.842}{their}%
\TextAt{8}{107.766}{227.842}{formal}%
\TextAt{8}{145.093}{227.842}{linear}%
\TextAt{8}{178.188}{227.842}{combinations,}%
\TextAt{8}{253.45}{227.842}{i.e.,}%
\TextAt{8}{276.276}{227.842}{the}%
\TextAt{12}{297.009}{227.842}{cycles}%
\TextAt{8}{335.723}{227.842}{of}%
\TextAt{13}{349.63}{227.842}{X}%
\TextAt{8}{360.285}{227.842}{,}%
\TextAt{8}{368.156}{227.842}{are}%
\TextAt{8}{388.243}{227.842}{to}%
\TextAt{8}{403.13}{227.842}{be}%
\TextAt{8}{419.632}{227.842}{considered}%
\TextAt{8}{477.498}{227.842}{as}%
\TextAt{8}{492.445}{227.842}{general-}%
\TextAt{8}{79.226}{210.796}{izations}%
\TextAt{8}{123.216}{210.796}{of}%
\TextAt{8}{137.546}{210.796}{the}%
\TextAt{8}{158.702}{210.796}{primes}%
\TextAt{8}{197.479}{210.796}{and}%
\TextAt{8}{221.241}{210.796}{divisors.}%
\TextAt{8}{271.518}{210.796}{The}%
\TextAt{8}{296.571}{210.796}{isomorphism}%
\TextAt{8}{365.99}{210.796}{between}%
\TextAt{8}{412.18}{210.796}{divisor}%
\TextAt{8}{451.293}{210.796}{class}%
\TextAt{8}{479.729}{210.796}{group}%
\TextAt{8}{513.904}{210.796}{and}%
\TextAt{8}{79.226}{193.75}{ideal}%
\TextAt{8}{107.734}{193.75}{class}%
\TextAt{8}{135.715}{193.75}{group}%
\TextAt{8}{169.435}{193.75}{extends}%
\TextAt{8}{212.636}{193.75}{to}%
\TextAt{8}{227.491}{193.75}{the}%
\TextAt{8}{248.191}{193.75}{general}%
\TextAt{8}{289.061}{193.75}{setting}%
\TextAt{8}{328.041}{193.75}{as}%
\TextAt{8}{342.955}{193.75}{a}%
\TextAt{8}{353.255}{193.75}{homomorphic}%
\TextAt{8}{427.3}{193.75}{relation}%
\TextAt{8}{470.775}{193.75}{between}%
\TextAt{8}{516.51}{193.75}{the}%
\TextAt{8}{79.226}{176.704}{groups}%
\TextAt{13}{116.478}{176.704}{C}%
\TextAt{13}{125.696}{176.704}{H}%
\TextAt{21}{136.334}{181.042}{i}%
\TextAt{9}{139.715}{176.704}{\Character{40}}%
\TextAt{13}{144.267}{176.704}{X}%
\TextAt{9}{154.923}{176.704}{\Character{41}}%
\TextAt{8}{162.833}{176.704}{and}%
\TextAt{13}{185.057}{176.704}{K}%
\TextAt{21}{195.016}{174.91}{i}%
\TextAt{9}{198.398}{176.704}{\Character{40}}%
\TextAt{13}{202.95}{176.704}{X}%
\TextAt{9}{213.605}{176.704}{\Character{41}}%
\TextAt{8}{218.157}{176.704}{.}%
\TextAt{8}{226.426}{176.704}{Thus}%
\TextAt{8}{255.536}{176.704}{the}%
\TextAt{8}{275.154}{176.704}{initial}%
\TextAt{8}{308.436}{176.704}{controversy}%
\TextAt{8}{370.088}{176.704}{has}%
\TextAt{8}{390.423}{176.704}{been}%
\TextAt{8}{417.512}{176.704}{resolved}%
\TextAt{8}{461.913}{176.704}{into}%
\TextAt{8}{485.117}{176.704}{a}%
\TextAt{8}{494.334}{176.704}{seminal}%
\TextAt{8}{79.226}{159.658}{mathematical}%
\TextAt{8}{152.727}{159.658}{theory}%
\TextAt{8}{189.477}{159.658}{(for}%
\TextAt{8}{211.917}{159.658}{further}%
\TextAt{8}{251.261}{159.658}{reading,}%
\TextAt{8}{296.141}{159.658}{see}%
\TextAt{8}{315.054}{159.658}{[13]).}%
\TextAt{12}{79.226}{134.172}{Exercise}%
\TextAt{12}{131.779}{134.172}{III.1.1.}%
\TextAt{12}{175.569}{134.172}{(Strong}%
\TextAt{12}{223.782}{134.172}{Approximation}%
\TextAt{12}{316.361}{134.172}{Theorem).}%
\TextAt{8}{382.879}{134.172}{Let}%
\TextAt{13}{403.286}{134.172}{S}%
\TextAt{8}{414.515}{134.172}{be}%
\TextAt{8}{429.876}{134.172}{a}%
\TextAt{8}{439.069}{134.172}{finite}%
\TextAt{8}{468.417}{134.172}{set}%
\TextAt{8}{486.122}{134.172}{of}%
\TextAt{8}{498.889}{134.172}{primes}%
\TextAt{8}{79.226}{117.126}{and}%
\TextAt{8}{103.19}{117.126}{let}%
\TextAt{14}{121.297}{117.126}{p}%
\TextAt{17}{127.28}{115.333}{0}%
\TextAt{8}{137.111}{117.126}{be}%
\TextAt{8}{154.237}{117.126}{another}%
\TextAt{8}{198.37}{117.126}{prime}%
\TextAt{8}{232.735}{117.126}{of}%
\TextAt{13}{247.267}{117.126}{K}%
\TextAt{8}{263.184}{117.126}{which}%
\TextAt{8}{297.872}{117.126}{does}%
\TextAt{8}{325.471}{117.126}{not}%
\TextAt{8}{347.486}{117.126}{belong}%
\TextAt{8}{386.084}{117.126}{to}%
\TextAt{13}{401.596}{117.126}{S}%
\TextAt{8}{409.491}{117.126}{.}%
\TextAt{8}{421.544}{117.126}{Let}%
\TextAt{13}{443.715}{117.126}{a}%
\TextAt{27}{449.86}{115.333}{p}%
\NativeMathAt{459.716}{117.126}{11.955}{\in}%
\TextAt{13}{473.055}{117.126}{K}%
\TextAt{8}{488.972}{117.126}{be}%
\TextAt{8}{506.098}{117.126}{given}%
\TextAt{8}{79.226}{100.08}{numbers,}%
\TextAt{8}{129.689}{100.08}{for}%
\TextAt{14}{147.574}{100.08}{p}%
\NativeMathAt{156.878}{100.08}{11.955}{\in}%
\TextAt{13}{168.169}{100.08}{S}%
\TextAt{8}{176.064}{100.08}{.}%
\TextAt{8}{184.512}{100.08}{Then}%
\TextAt{8}{215.07}{100.08}{for}%
\TextAt{8}{232.955}{100.08}{every}%
\TextAt{13}{263.847}{100.08}{\Character{34}}%
\TextAt{13}{272.634}{100.08}{\Character{62}}%
\TextAt{9}{285.071}{100.08}{0}%
\TextAt{8}{290.924}{100.08}{,}%
\TextAt{8}{298.073}{100.08}{there}%
\TextAt{8}{327.985}{100.08}{exists}%
\TextAt{8}{360.3}{100.08}{an}%
\TextAt{13}{376.559}{100.08}{x}%
\NativeMathAt{386.532}{100.08}{11.955}{\in}%
\TextAt{13}{397.823}{100.08}{K}%
\TextAt{8}{412.538}{100.08}{such}%
\TextAt{8}{438.935}{100.08}{that}%
\TextAt{16}{169.959}{72.921}{\Character{106}}%
\TextAt{13}{173.28}{72.921}{x}%
\TextAt{16}{182.589}{72.921}{\Character{0}}%
\TextAt{13}{194.544}{72.921}{a}%
\TextAt{27}{200.689}{71.128}{p}%
\TextAt{16}{205.176}{72.921}{\Character{106}}%
\TextAt{27}{208.497}{71.128}{p}%
\TextAt{13}{216.304}{72.921}{\Character{60}}%
\TextAt{13}{228.714}{72.921}{\Character{34}}%
\TextAt{8}{238.105}{72.921}{for}%
\TextAt{14}{255.99}{72.921}{p}%
\NativeMathAt{265.294}{72.921}{11.955}{\in}%
\TextAt{13}{276.585}{72.921}{S\Character{59}}%
\TextAt{16}{312.984}{72.921}{\Character{106}}%
\TextAt{13}{316.305}{72.921}{x}%
\TextAt{16}{322.957}{72.921}{\Character{106}}%
\TextAt{27}{326.278}{71.128}{p}%
\NativeMathAt{334.086}{72.921}{11.955}{\leq}%
\TextAt{9}{346.705}{72.921}{1}%
\TextAt{8}{356.456}{72.921}{for}%
\TextAt{14}{374.341}{72.921}{p}%
\NativeMathAt{383.645}{72.921}{11.955}{\notin}%
\TextAt{13}{394.936}{72.921}{S}%
\TextAt{16}{405.488}{72.921}{\Character{91}}%
\TextAt{16}{416.104}{72.921}{\Character{102}}%
\TextAt{14}{422.092}{72.921}{p}%
\TextAt{17}{428.075}{71.128}{0}%
\TextAt{16}{432.807}{72.921}{\Character{103}}%
\TextAt{13}{438.785}{72.921}{\Character{58}}%
\TextAt{8}{297.211}{43.033}{176}%
\Destination{78.23}{809.815}{nav.page.176}%
\Destination{79.226}{148.155}{nav.ex.3.1.1}%
\Destination{79.36}{551.155}{nav.print.193}%
\Destination{79.36}{577.355}{nav.number.3.1.14}%
\Hotlink{314.554}{157.455}{19.22}{11.047}{nav.cite.N013}%
\Hotlink{167.884}{586.655}{30.936}{11.047}{nav.number.3.1.12}%
\Hotlink{269.79}{481.355}{30.936}{11.047}{nav.number.3.1.13}%
\Destination{79.226}{702}{errata.203}%
}

% Page 189 (body page 177)
\Page{%
\TextAt{2}{79.226}{796.943}{III.}%
\TextAt{2}{98.317}{796.943}{Riemann–Roch}%
\TextAt{2}{174.375}{796.943}{Theory}%
\Vector{q 0 G 0 g 0.299 w 7.226 18.729 m 460.769 18.729 l S Q}%
\TextAt{12}{79.226}{761.124}{Exercise}%
\TextAt{12}{133.019}{761.124}{III.1.2.}%
\TextAt{8}{177.887}{761.124}{Let}%
\TextAt{13}{199.371}{761.124}{K}%
\TextAt{8}{214.6}{761.124}{be}%
\TextAt{8}{231.039}{761.124}{totally}%
\TextAt{8}{268.962}{761.124}{real,}%
\TextAt{8}{295.619}{761.124}{i.e.,}%
\TextAt{13}{318.367}{761.124}{K}%
\TextAt{27}{328.327}{759.33}{p}%
\TextAt{9}{337.011}{761.124}{\Character{61}}%
\TextAt{18}{350.314}{761.124}{R}%
\TextAt{8}{363.36}{761.124}{for}%
\TextAt{8}{381.759}{761.124}{all}%
\TextAt{14}{398.533}{761.124}{p}%
\TextAt{16}{408.714}{761.124}{\Character{106}}%
\NativeMathAt{416.222}{761.124}{11.955}{\infty}%
\TextAt{8}{428.187}{761.124}{.}%
\TextAt{8}{438.18}{761.124}{Let}%
\TextAt{13}{459.663}{761.124}{T}%
\TextAt{8}{472.562}{761.124}{be}%
\TextAt{8}{489.001}{761.124}{a}%
\TextAt{8}{499.271}{761.124}{proper}%
\TextAt{8}{79.226}{744.078}{nonempty}%
\TextAt{8}{134.659}{744.078}{subset}%
\TextAt{8}{171.347}{744.078}{of}%
\TextAt{8}{185.475}{744.078}{Hom}%
\TextAt{9}{209.863}{744.078}{\Character{40}}%
\TextAt{13}{214.416}{744.078}{K}%
\TextAt{13}{224.578}{744.078}{\Character{59}}%
\TextAt{18}{229.827}{744.078}{R}%
\TextAt{9}{238.462}{744.078}{\Character{41}}%
\TextAt{8}{243.014}{744.078}{.}%
\TextAt{8}{253.856}{744.078}{Then}%
\TextAt{8}{285.212}{744.078}{there}%
\TextAt{8}{315.922}{744.078}{exists}%
\TextAt{8}{349.034}{744.078}{a}%
\TextAt{8}{359.588}{744.078}{unit}%
\TextAt{13}{385.097}{744.078}{\Character{34}}%
\TextAt{8}{395.271}{744.078}{of}%
\TextAt{13}{409.399}{744.078}{K}%
\TextAt{8}{424.912}{744.078}{satisfying}%
\TextAt{13}{477.87}{744.078}{\Character{28}}%
\TextAt{13}{484.278}{744.078}{\Character{34}}%
\TextAt{13}{494.44}{744.078}{\Character{62}}%
\TextAt{9}{508.234}{744.078}{1}%
\TextAt{8}{518.782}{744.078}{for}%
\TextAt{13}{79.226}{727.032}{\Character{28}}%
\NativeMathAt{88.966}{727.032}{11.955}{\in}%
\TextAt{13}{100.257}{727.032}{T}%
\TextAt{8}{112.641}{727.032}{and}%
\TextAt{9}{135.404}{727.032}{0}%
\TextAt{13}{144.578}{727.032}{\Character{60}}%
\TextAt{13}{156.988}{727.032}{\Character{28}}%
\TextAt{13}{163.407}{727.032}{\Character{34}}%
\TextAt{13}{172.207}{727.032}{\Character{60}}%
\TextAt{9}{184.646}{727.032}{1}%
\TextAt{8}{194.397}{727.032}{for}%
\TextAt{13}{212.282}{727.032}{\Character{28}}%
\NativeMathAt{222.021}{727.032}{11.955}{\notin}%
\TextAt{13}{233.312}{727.032}{T}%
\TextAt{8}{241.799}{727.032}{.}%
\TextAt{12}{79.226}{696.887}{Exercise}%
\TextAt{12}{132.427}{696.887}{III.1.3.}%
\TextAt{8}{176.78}{696.887}{Show}%
\TextAt{8}{207.672}{696.887}{that}%
\TextAt{8}{233.041}{696.887}{the}%
\TextAt{8}{253.198}{696.887}{absolute}%
\TextAt{8}{299.44}{696.887}{norm}%
\TextAt{14}{330.01}{696.887}{N}%
\TextAt{9}{343.279}{696.887}{\Character{58}}%
\TextAt{8}{349.851}{696.887}{Pic}%
\TextAt{9}{366.266}{696.887}{\Character{40}}%
\Vector{q 0 G 0 g 0.478 w 298.818 -67.204 m 306.788 -67.204 l S Q}%
\TextAt{18}{370.818}{696.887}{Z}%
\TextAt{9}{378.788}{696.887}{\Character{41}}%
\NativeMathAt{386.661}{696.887}{11.955}{\rightarrow}%
\TextAt{18}{401.937}{696.887}{R}%
\TextAt{3}{410.572}{701.225}{\Character{3}}%
\TextAt{17}{410.572}{693.931}{\Character{43}}%
\TextAt{8}{421.554}{696.887}{is}%
\TextAt{8}{433.317}{696.887}{an}%
\TextAt{8}{449.577}{696.887}{isomorphism.}%
\TextAt{12}{79.226}{666.741}{Exercise}%
\TextAt{12}{131.518}{666.741}{III.1.4.}%
\TextAt{8}{175.081}{666.741}{Let}%
\TextAt{13}{195.26}{666.741}{K}%
\TextAt{8}{209.185}{666.741}{and}%
\TextAt{13}{231.158}{666.741}{L}%
\TextAt{8}{242.229}{666.741}{be}%
\TextAt{8}{257.363}{666.741}{number}%
\TextAt{8}{299.169}{666.741}{fields,}%
\TextAt{8}{331.761}{666.741}{and}%
\TextAt{8}{353.733}{666.741}{let}%
\TextAt{13}{369.847}{666.741}{\Character{28}}%
\TextAt{9}{379.587}{666.741}{\Character{58}}%
\TextAt{13}{386.16}{666.741}{K}%
\NativeMathAt{400.298}{666.741}{11.955}{\rightarrow}%
\TextAt{13}{415.574}{666.741}{L}%
\TextAt{8}{426.646}{666.741}{be}%
\TextAt{8}{441.78}{666.741}{a}%
\TextAt{8}{450.745}{666.741}{homomorphism.}%
\TextAt{8}{79.226}{649.696}{Given}%
\TextAt{8}{113.119}{649.696}{any}%
\TextAt{8}{135.236}{649.696}{replete}%
\TextAt{8}{173.6}{649.696}{divisor}%
\TextAt{13}{211.713}{649.696}{D}%
\TextAt{9}{225.079}{649.696}{\Character{61}}%
\TextAt{20}{237.504}{658.662}{\Character{80}}%
\TextAt{27}{250.124}{646.209}{P}%
\TextAt{13}{259.21}{649.696}{\Character{23}}%
\TextAt{27}{264.993}{647.902}{P}%
\TextAt{14}{272.087}{649.696}{P}%
\TextAt{8}{285.878}{649.696}{of}%
\TextAt{13}{299.208}{649.696}{L}%
\TextAt{8}{307.172}{649.696}{,}%
\TextAt{8}{314.321}{649.696}{define}%
\TextAt{8}{348.131}{649.696}{a}%
\TextAt{8}{357.886}{649.696}{replete}%
\TextAt{8}{396.25}{649.696}{divisor}%
\TextAt{8}{434.363}{649.696}{of}%
\TextAt{13}{447.693}{649.696}{K}%
\TextAt{8}{462.408}{649.696}{by}%
\TextAt{8}{478.667}{649.696}{the}%
\TextAt{8}{498.824}{649.696}{rule}%
\TextAt{13}{232.266}{610.47}{\Character{28}}%
\TextAt{3}{237.349}{608.677}{\Character{3}}%
\TextAt{9}{242.082}{610.47}{\Character{40}}%
\TextAt{13}{246.634}{610.47}{D}%
\TextAt{9}{256.678}{610.47}{\Character{41}}%
\TextAt{9}{264.533}{610.47}{\Character{61}}%
\TextAt{20}{276.977}{621.828}{\Character{88}}%
\TextAt{27}{283.617}{596.722}{p}%
\TextAt{20}{296.238}{634.501}{\Character{48}}%
\TextAt{20}{296.238}{612.981}{\Character{64}}%
\TextAt{20}{306.699}{621.828}{\Character{88}}%
\TextAt{27}{308.865}{595.925}{P}%
\TextAt{3}{315.46}{595.925}{\Character{106}}%
\TextAt{27}{317.813}{595.925}{p}%
\TextAt{13}{325.96}{610.47}{\Character{23}}%
\TextAt{27}{331.744}{608.677}{P}%
\TextAt{13}{338.837}{610.47}{f}%
\TextAt{27}{344.607}{608.611}{P}%
\TextAt{3}{351.203}{608.611}{\Character{106}}%
\TextAt{27}{353.555}{608.611}{p}%
\TextAt{20}{358.042}{634.501}{\Character{49}}%
\TextAt{20}{358.042}{612.981}{\Character{65}}%
\TextAt{14}{370.495}{610.47}{p}%
\TextAt{13}{376.478}{610.47}{\Character{59}}%
\TextAt{8}{79.226}{573.25}{where}%
\TextAt{13}{113.316}{573.25}{f}%
\TextAt{27}{119.086}{571.391}{P}%
\TextAt{3}{125.681}{571.391}{\Character{106}}%
\TextAt{27}{128.033}{571.391}{p}%
\TextAt{8}{136.698}{573.25}{is}%
\TextAt{8}{148.742}{573.25}{the}%
\TextAt{8}{169.178}{573.25}{inertia}%
\TextAt{8}{206.531}{573.25}{degree}%
\TextAt{8}{243.227}{573.25}{of}%
\TextAt{14}{256.837}{573.25}{P}%
\TextAt{8}{270.908}{573.25}{over}%
\TextAt{13}{296.234}{573.25}{\Character{28}}%
\TextAt{13}{302.642}{573.25}{K}%
\TextAt{8}{317.648}{573.25}{and}%
\TextAt{14}{340.691}{573.25}{P}%
\TextAt{16}{354.383}{573.25}{\Character{106}}%
\TextAt{14}{361.502}{573.25}{p}%
\TextAt{8}{371.662}{573.25}{signifies}%
\TextAt{13}{415.639}{573.25}{\Character{28}}%
\TextAt{14}{422.058}{573.25}{p}%
\TextAt{9}{431.839}{573.25}{\Character{61}}%
\TextAt{14}{444.742}{573.25}{P}%
\TextAt{16}{454.635}{573.25}{\Character{106}}%
\TextAt{21}{457.956}{571.457}{\Character{28}}%
\TextAt{21}{462.627}{571.457}{K}%
\TextAt{8}{470.837}{573.25}{.}%
\TextAt{8}{480.126}{573.25}{Show}%
\TextAt{8}{511.298}{573.25}{that}%
\TextAt{13}{79.226}{556.204}{\Character{28}}%
\TextAt{3}{84.31}{554.411}{\Character{3}}%
\TextAt{8}{92.94}{556.204}{induces}%
\TextAt{8}{134.615}{556.204}{a}%
\TextAt{8}{144.371}{556.204}{homomorphism}%
\TextAt{13}{234.799}{530.432}{\Character{28}}%
\TextAt{3}{239.882}{528.638}{\Character{3}}%
\TextAt{9}{247.935}{530.432}{\Character{58}}%
\TextAt{13}{254.508}{530.432}{C}%
\TextAt{13}{263.726}{530.432}{H}%
\TextAt{17}{274.364}{535.368}{1}%
\TextAt{9}{279.096}{530.432}{\Character{40}}%
\Vector{q 0 G 0 g 0.478 w 211.649 -236.212 m 217.493 -236.212 l S Q}%
\TextAt{14}{283.649}{530.432}{o}%
\TextAt{21}{289.493}{528.638}{L}%
\TextAt{9}{295.75}{530.432}{\Character{41}}%
\NativeArrowAt{303.623}{530.432}{11.955}{19.26}{\longrightarrow}%
\TextAt{13}{326.205}{530.432}{C}%
\TextAt{13}{335.423}{530.432}{H}%
\TextAt{17}{346.062}{535.368}{1}%
\TextAt{9}{350.794}{530.432}{\Character{40}}%
\Vector{q 0 G 0 g 0.478 w 283.346 -236.212 m 289.19 -236.212 l S Q}%
\TextAt{14}{355.346}{530.432}{o}%
\TextAt{21}{361.19}{528.638}{K}%
\TextAt{9}{369.393}{530.432}{\Character{41}}%
\TextAt{13}{373.945}{530.432}{\Character{58}}%
\TextAt{12}{79.226}{491.56}{Exercise}%
\TextAt{12}{132.221}{491.56}{III.1.5.}%
\TextAt{8}{176.397}{491.56}{Given}%
\TextAt{8}{210.111}{491.56}{any}%
\TextAt{8}{232.05}{491.56}{replete}%
\TextAt{8}{270.236}{491.56}{divisor}%
\TextAt{13}{308.17}{491.56}{D}%
\TextAt{9}{321.536}{491.56}{\Character{61}}%
\TextAt{20}{333.961}{500.526}{\Character{80}}%
\TextAt{27}{346.581}{488.073}{p}%
\TextAt{13}{353.06}{491.56}{\Character{23}}%
\TextAt{27}{358.844}{489.766}{p}%
\TextAt{14}{363.331}{491.56}{p}%
\TextAt{8}{373.033}{491.56}{of}%
\TextAt{13}{386.184}{491.56}{K}%
\TextAt{8}{397.002}{491.56}{,}%
\TextAt{8}{404.009}{491.56}{define}%
\TextAt{8}{437.639}{491.56}{a}%
\TextAt{8}{447.216}{491.56}{replete}%
\TextAt{8}{485.402}{491.56}{divisor}%
\TextAt{8}{523.337}{491.56}{of}%
\TextAt{13}{79.226}{474.514}{L}%
\TextAt{8}{91.088}{474.514}{by}%
\TextAt{8}{107.347}{474.514}{the}%
\TextAt{8}{127.503}{474.514}{rule}%
\TextAt{13}{242.576}{457.468}{\Character{28}}%
\TextAt{3}{248.994}{462.404}{\Character{3}}%
\TextAt{9}{253.727}{457.468}{\Character{40}}%
\TextAt{13}{258.279}{457.468}{D}%
\TextAt{9}{268.324}{457.468}{\Character{41}}%
\TextAt{9}{276.179}{457.468}{\Character{61}}%
\TextAt{20}{288.622}{468.825}{\Character{88}}%
\TextAt{27}{295.262}{443.719}{p}%
\TextAt{20}{307.883}{468.825}{\Character{88}}%
\TextAt{27}{310.049}{442.922}{P}%
\TextAt{3}{316.645}{442.922}{\Character{106}}%
\TextAt{27}{318.997}{442.922}{p}%
\TextAt{13}{327.144}{457.468}{\Character{23}}%
\TextAt{27}{332.928}{455.675}{p}%
\TextAt{13}{337.415}{457.468}{e}%
\TextAt{27}{342.841}{455.608}{P}%
\TextAt{3}{349.436}{455.608}{\Character{106}}%
\TextAt{27}{351.788}{455.608}{p}%
\TextAt{14}{356.275}{457.468}{P}%
\TextAt{13}{366.168}{457.468}{\Character{59}}%
\TextAt{8}{79.226}{425.229}{where}%
\TextAt{13}{113.733}{425.229}{e}%
\TextAt{27}{119.158}{423.37}{P}%
\TextAt{3}{125.754}{423.37}{\Character{106}}%
\TextAt{27}{128.106}{423.37}{p}%
\TextAt{8}{137.188}{425.229}{denotes}%
\TextAt{8}{180.218}{425.229}{the}%
\TextAt{8}{201.072}{425.229}{ramification}%
\TextAt{8}{266.817}{425.229}{index}%
\TextAt{8}{299.053}{425.229}{of}%
\TextAt{14}{313.08}{425.229}{P}%
\TextAt{8}{327.568}{425.229}{over}%
\TextAt{13}{353.311}{425.229}{K}%
\TextAt{8}{364.129}{425.229}{.}%
\TextAt{8}{374.669}{425.229}{Show}%
\TextAt{8}{406.259}{425.229}{that}%
\TextAt{13}{432.325}{425.229}{\Character{28}}%
\TextAt{3}{438.744}{429.568}{\Character{3}}%
\TextAt{8}{448.071}{425.229}{induces}%
\TextAt{8}{490.444}{425.229}{a}%
\TextAt{8}{500.897}{425.229}{homo-}%
\TextAt{8}{79.226}{408.183}{morphism}%
\TextAt{13}{234.131}{391.137}{\Character{28}}%
\TextAt{3}{240.55}{396.074}{\Character{3}}%
\TextAt{9}{248.603}{391.137}{\Character{58}}%
\TextAt{13}{255.175}{391.137}{C}%
\TextAt{13}{264.393}{391.137}{H}%
\TextAt{17}{275.032}{396.074}{1}%
\TextAt{9}{279.764}{391.137}{\Character{40}}%
\Vector{q 0 G 0 g 0.478 w 212.316 -375.506 m 218.161 -375.506 l S Q}%
\TextAt{14}{284.316}{391.137}{o}%
\TextAt{21}{290.161}{389.344}{K}%
\TextAt{9}{298.363}{391.137}{\Character{41}}%
\NativeArrowAt{306.236}{391.137}{11.955}{19.26}{\longrightarrow}%
\TextAt{13}{328.818}{391.137}{C}%
\TextAt{13}{338.036}{391.137}{H}%
\TextAt{17}{348.675}{396.074}{1}%
\TextAt{9}{353.407}{391.137}{\Character{40}}%
\Vector{q 0 G 0 g 0.478 w 285.959 -375.506 m 291.804 -375.506 l S Q}%
\TextAt{14}{357.959}{391.137}{o}%
\TextAt{21}{363.804}{389.344}{L}%
\TextAt{9}{370.061}{391.137}{\Character{41}}%
\TextAt{13}{374.613}{391.137}{\Character{58}}%
\TextAt{12}{79.226}{355.9}{Exercise}%
\TextAt{12}{132.427}{355.9}{III.1.6.}%
\TextAt{8}{176.78}{355.9}{Show}%
\TextAt{8}{207.672}{355.9}{that}%
\TextAt{13}{233.041}{355.9}{\Character{28}}%
\TextAt{3}{238.125}{354.107}{\Character{3}}%
\TextAt{16}{245.514}{355.9}{\Character{14}}%
\TextAt{13}{254.148}{355.9}{\Character{28}}%
\TextAt{3}{260.567}{360.239}{\Character{3}}%
\TextAt{9}{268.62}{355.9}{\Character{61}}%
\TextAt{9}{281.03}{355.9}{\Character{91}}%
\TextAt{13}{284.298}{355.9}{L}%
\TextAt{9}{295.583}{355.9}{\Character{58}}%
\TextAt{13}{302.155}{355.9}{K}%
\TextAt{9}{312.973}{355.9}{\Character{93}}%
\TextAt{8}{320.122}{355.9}{and}%
\TextAt{8}{342.885}{355.9}{that}%
\TextAt{8}{171.557}{330.127}{deg}%
\TextAt{9}{189.119}{330.127}{\Character{40}}%
\TextAt{13}{193.671}{330.127}{\Character{28}}%
\TextAt{3}{198.755}{328.334}{\Character{3}}%
\TextAt{13}{203.487}{330.127}{D}%
\TextAt{9}{213.532}{330.127}{\Character{41}}%
\TextAt{9}{221.387}{330.127}{\Character{61}}%
\TextAt{8}{233.83}{330.127}{deg}%
\TextAt{9}{251.392}{330.127}{\Character{40}}%
\TextAt{13}{255.945}{330.127}{D}%
\TextAt{9}{265.989}{330.127}{\Character{41}}%
\TextAt{13}{270.541}{330.127}{\Character{59}}%
\TextAt{8}{299.696}{330.127}{deg}%
\TextAt{9}{317.258}{330.127}{\Character{40}}%
\TextAt{13}{321.81}{330.127}{\Character{28}}%
\TextAt{3}{328.229}{335.064}{\Character{3}}%
\TextAt{13}{332.962}{330.127}{D}%
\TextAt{9}{343.006}{330.127}{\Character{41}}%
\TextAt{9}{350.861}{330.127}{\Character{61}}%
\TextAt{9}{363.282}{330.127}{\Character{91}}%
\TextAt{13}{366.556}{330.127}{L}%
\TextAt{9}{377.842}{330.127}{\Character{58}}%
\TextAt{13}{384.414}{330.127}{K}%
\TextAt{9}{395.232}{330.127}{\Character{93}}%
\TextAt{8}{400.476}{330.127}{deg}%
\TextAt{9}{418.038}{330.127}{\Character{40}}%
\TextAt{13}{422.591}{330.127}{D}%
\TextAt{9}{432.635}{330.127}{\Character{41}}%
\TextAt{13}{437.187}{330.127}{\Character{58}}%
\TextAt{15}{79.226}{269.937}{§}%
\TextAt{15}{93.785}{269.937}{2}%
\TextAt{15}{120.693}{269.937}{Different}%
\TextAt{15}{202.019}{269.937}{and}%
\TextAt{15}{239.411}{269.937}{Discriminant}%
\TextAt{8}{79.226}{238.051}{It}%
\TextAt{8}{91.581}{238.051}{is}%
\TextAt{8}{103.028}{238.051}{our}%
\TextAt{8}{123.525}{238.051}{intention}%
\TextAt{8}{172.642}{238.051}{to}%
\TextAt{8}{186.635}{238.051}{develop}%
\TextAt{8}{228.592}{238.051}{a}%
\TextAt{8}{238.03}{238.051}{framework}%
\TextAt{8}{295.277}{238.051}{for}%
\TextAt{8}{312.845}{238.051}{the}%
\TextAt{8}{332.684}{238.051}{theory}%
\TextAt{8}{369.117}{238.051}{of}%
\TextAt{8}{382.13}{238.051}{algebraic}%
\TextAt{8}{431.247}{238.051}{number}%
\TextAt{8}{473.526}{238.051}{fields}%
\TextAt{8}{503.18}{238.051}{which}%
\TextAt{8}{79.226}{221.005}{shows}%
\TextAt{8}{111.907}{221.005}{the}%
\TextAt{8}{131.127}{221.005}{complete}%
\TextAt{8}{179.613}{221.005}{analogy}%
\TextAt{8}{221.942}{221.005}{with}%
\TextAt{8}{247.666}{221.005}{the}%
\TextAt{8}{266.885}{221.005}{theory}%
\TextAt{8}{302.699}{221.005}{of}%
\TextAt{8}{315.092}{221.005}{function}%
\TextAt{8}{360.003}{221.005}{fields.}%
\TextAt{8}{394.214}{221.005}{This}%
\TextAt{8}{419.997}{221.005}{goal}%
\TextAt{8}{443.784}{221.005}{leads}%
\TextAt{8}{472.173}{221.005}{us}%
\TextAt{8}{486.252}{221.005}{naturally}%
\TextAt{8}{79.226}{203.959}{to}%
\TextAt{8}{92.743}{203.959}{the}%
\TextAt{8}{112.106}{203.959}{notions}%
\TextAt{8}{152.354}{203.959}{of}%
\TextAt{8}{164.89}{203.959}{different}%
\TextAt{8}{210.268}{203.959}{and}%
\TextAt{8}{232.236}{203.959}{discriminant,}%
\TextAt{8}{302.233}{203.959}{as}%
\TextAt{8}{315.809}{203.959}{we}%
\TextAt{8}{332.243}{203.959}{shall}%
\TextAt{8}{358.827}{203.959}{explain}%
\TextAt{8}{398.68}{203.959}{in}%
\TextAt{8}{411.54}{203.959}{§3}%
\TextAt{8}{426.18}{203.959}{and}%
\TextAt{8}{448.149}{203.959}{§7.}%
\TextAt{8}{467.869}{203.959}{They}%
\TextAt{8}{497.31}{203.959}{control}%
\TextAt{8}{79.226}{186.913}{the}%
\TextAt{8}{99.383}{186.913}{ramification}%
\TextAt{8}{164.431}{186.913}{behaviour}%
\TextAt{8}{218.743}{186.913}{of}%
\TextAt{8}{232.073}{186.913}{an}%
\TextAt{8}{248.332}{186.913}{extension}%
\TextAt{8}{300.098}{186.913}{of}%
\TextAt{8}{313.428}{186.913}{valued}%
\TextAt{8}{350.178}{186.913}{fields.}%
\TextAt{8}{97.159}{166.878}{Let}%
\TextAt{13}{118.447}{166.878}{L}%
\TextAt{16}{130.276}{166.878}{\Character{106}}%
\TextAt{13}{137.462}{166.878}{K}%
\TextAt{8}{152.496}{166.878}{be}%
\TextAt{8}{168.74}{166.878}{a}%
\TextAt{8}{178.814}{166.878}{finite}%
\TextAt{8}{209.045}{166.878}{separable}%
\TextAt{8}{260.807}{166.878}{field}%
\TextAt{8}{286.483}{166.878}{extension,}%
\TextAt{14}{341.9}{166.878}{o}%
\NativeMathAt{351.609}{166.878}{11.955}{\subseteq}%
\TextAt{13}{364.772}{166.878}{K}%
\TextAt{8}{379.806}{166.878}{a}%
\TextAt{8}{389.88}{166.878}{Dedekind}%
\TextAt{8}{442.384}{166.878}{domain}%
\TextAt{8}{484.331}{166.878}{with}%
\TextAt{8}{511.31}{166.878}{field}%
\TextAt{8}{79.226}{149.832}{of}%
\TextAt{8}{93.175}{149.832}{fractions}%
\TextAt{13}{141.662}{149.832}{K}%
\TextAt{8}{152.479}{149.832}{,}%
\TextAt{8}{160.402}{149.832}{and}%
\TextAt{8}{183.783}{149.832}{let}%
\TextAt{16}{201.306}{149.832}{\Character{79}}%
\NativeMathAt{215.521}{149.832}{11.955}{\subseteq}%
\TextAt{13}{229.204}{149.832}{L}%
\TextAt{8}{241.685}{149.832}{be}%
\TextAt{8}{258.228}{149.832}{its}%
\TextAt{8}{275.165}{149.832}{integral}%
\TextAt{8}{318.392}{149.832}{closure}%
\TextAt{8}{358.092}{149.832}{in}%
\TextAt{13}{372.363}{149.832}{L}%
\TextAt{8}{380.328}{149.832}{.}%
\TextAt{8}{390.632}{149.832}{Throughout}%
\TextAt{8}{456.298}{149.832}{this}%
\TextAt{8}{479.739}{149.832}{section}%
\TextAt{8}{519.439}{149.832}{we}%
\TextAt{8}{79.226}{132.786}{assume}%
\TextAt{8}{119.858}{132.786}{systematically}%
\TextAt{8}{196.272}{132.786}{that}%
\TextAt{8}{221.829}{132.786}{the}%
\TextAt{8}{242.173}{132.786}{residue}%
\TextAt{8}{282.088}{132.786}{field}%
\TextAt{8}{307.633}{132.786}{extensions}%
\TextAt{13}{364.202}{132.786}{\Character{21}}%
\TextAt{16}{374.671}{132.786}{\Character{106}}%
\TextAt{13}{381.633}{132.786}{\Character{20}}%
\TextAt{8}{392.457}{132.786}{of}%
\TextAt{16}{405.975}{132.786}{\Character{79}}%
\TextAt{16}{419.461}{132.786}{\Character{106}}%
\TextAt{14}{426.427}{132.786}{o}%
\TextAt{8}{436.357}{132.786}{are}%
\TextAt{8}{456.055}{132.786}{separable.}%
\TextAt{8}{512.613}{132.786}{The}%
\TextAt{8}{79.226}{115.74}{theory}%
\TextAt{8}{114.943}{115.74}{of}%
\TextAt{8}{127.24}{115.74}{the}%
\TextAt{8}{146.363}{115.74}{different}%
\TextAt{8}{191.501}{115.74}{originates}%
\TextAt{8}{243.871}{115.74}{from}%
\TextAt{8}{270.478}{115.74}{the}%
\TextAt{8}{289.601}{115.74}{fact}%
\TextAt{8}{311.653}{115.74}{that}%
\TextAt{8}{335.989}{115.74}{we}%
\TextAt{8}{352.183}{115.74}{are}%
\TextAt{8}{370.661}{115.74}{given}%
\TextAt{8}{400.197}{115.74}{a}%
\TextAt{8}{408.919}{115.74}{canonical}%
\TextAt{8}{459.269}{115.74}{nondegenerate}%
\TextAt{8}{79.226}{98.694}{symmetric}%
\TextAt{8}{136.193}{98.694}{bilinear}%
\TextAt{8}{178.466}{98.694}{form}%
\TextAt{8}{206.106}{98.694}{on}%
\TextAt{8}{222.365}{98.694}{the}%
\TextAt{13}{242.522}{98.694}{K}%
\TextAt{8}{253.34}{98.694}{-vector}%
\TextAt{8}{292.361}{98.694}{space}%
\TextAt{13}{323.636}{98.694}{L}%
\TextAt{8}{331.6}{98.694}{,}%
\TextAt{8}{338.75}{98.694}{viz.,}%
\TextAt{8}{363.784}{98.694}{the}%
\TextAt{12}{383.94}{98.694}{trace}%
\TextAt{12}{417.809}{98.694}{form}%
\TextAt{13}{262.862}{72.921}{T}%
\TextAt{9}{271.349}{72.921}{\Character{40}}%
\TextAt{13}{275.901}{72.921}{x\Character{59}}%
\TextAt{13}{287.773}{72.921}{y}%
\TextAt{9}{293.934}{72.921}{\Character{41}}%
\TextAt{9}{301.789}{72.921}{\Character{61}}%
\TextAt{8}{314.233}{72.921}{Tr}%
\TextAt{9}{327.24}{72.921}{\Character{40}}%
\TextAt{13}{331.792}{72.921}{xy}%
\TextAt{9}{344.581}{72.921}{\Character{41}}%
\TextAt{8}{297.211}{43.033}{177}%
\Destination{78.23}{809.815}{nav.page.177}%
\Destination{79.23}{295.795}{nav.section.3.2}%
\Destination{79.226}{775.155}{nav.ex.3.1.2}%
\Destination{79.226}{710.855}{nav.ex.3.1.3}%
\Destination{79.226}{680.755}{nav.ex.3.1.4}%
\Destination{79.226}{505.555}{nav.ex.3.1.5}%
\Destination{79.226}{369.855}{nav.ex.3.1.6}%
\Destination{79.36}{664.655}{nav.print.194}%
\Destination{79.23}{295.795}{outline.33}%
\Hotlink{416.719}{201.755}{6.858}{11.047}{nav.section.3.3}%
\Hotlink{453.328}{201.755}{6.858}{11.047}{nav.section.3.7}%
}

% Page 190 (body page 178)
\Page{%
\TextAt{2}{79.226}{796.943}{III.}%
\TextAt{2}{98.317}{796.943}{Riemann–Roch}%
\TextAt{2}{174.375}{796.943}{Theory}%
\Vector{q 0 G 0 g 0.299 w 7.226 18.729 m 460.769 18.729 l S Q}%
\TextAt{8}{79.226}{761.124}{(see}%
\TextAt{8}{104.142}{761.124}{chap.}%
\TextAt{8}{140.678}{761.124}{I,}%
\TextAt{8}{153.496}{761.124}{§2).}%
\TextAt{8}{182.38}{761.124}{It}%
\TextAt{8}{196.501}{761.124}{allows}%
\TextAt{8}{232.811}{761.124}{us}%
\TextAt{8}{249.275}{761.124}{to}%
\TextAt{8}{265.033}{761.124}{associate}%
\TextAt{8}{315.713}{761.124}{to}%
\TextAt{8}{331.471}{761.124}{every}%
\TextAt{8}{363.812}{761.124}{fractional}%
\TextAt{8}{417.623}{761.124}{ideal}%
\TextAt{14}{447.035}{761.124}{A}%
\TextAt{8}{460.959}{761.124}{of}%
\TextAt{13}{475.738}{761.124}{L}%
\TextAt{8}{489.048}{761.124}{the}%
\TextAt{8}{510.652}{761.124}{dual}%
\TextAt{16}{79.226}{744.078}{\Character{79}}%
\TextAt{8}{89.076}{744.078}{-module}%
\TextAt{3}{234.706}{731.968}{\Character{3}}%
\TextAt{14}{239.439}{727.032}{A}%
\TextAt{9}{251.339}{727.032}{\Character{61}}%
\TextAt{16}{263.764}{727.032}{\Character{102}}%
\TextAt{13}{269.742}{727.032}{x}%
\NativeMathAt{279.715}{727.032}{11.955}{\in}%
\TextAt{13}{291.006}{727.032}{L}%
\TextAt{16}{302.291}{727.032}{\Character{106}}%
\TextAt{8}{308.933}{727.032}{Tr}%
\TextAt{9}{321.94}{727.032}{\Character{40}}%
\TextAt{13}{326.492}{727.032}{x}%
\TextAt{14}{333.144}{727.032}{A}%
\TextAt{9}{341.723}{727.032}{\Character{41}}%
\NativeMathAt{349.597}{727.032}{11.955}{\subseteq}%
\TextAt{14}{362.216}{727.032}{o}%
\TextAt{16}{368.06}{727.032}{\Character{103}}%
\TextAt{13}{374.038}{727.032}{\Character{58}}%
\TextAt{8}{79.226}{704.484}{It}%
\TextAt{8}{94.023}{704.484}{is}%
\TextAt{8}{107.911}{704.484}{again}%
\TextAt{8}{141.262}{704.484}{a}%
\TextAt{8}{153.142}{704.484}{fractional}%
\TextAt{8}{207.63}{704.484}{ideal.}%
\TextAt{8}{246.516}{704.484}{For}%
\TextAt{8}{269.61}{704.484}{if}%
\TextAt{13}{282.458}{704.484}{\Character{11}}%
\TextAt{17}{289.904}{702.69}{1}%
\TextAt{13}{294.636}{704.484}{\Character{59}}%
\TextAt{13}{299.86}{704.484}{\Character{58}}%
\TextAt{13}{305.097}{704.484}{\Character{58}}%
\TextAt{13}{310.321}{704.484}{\Character{58}}%
\TextAt{13}{315.558}{704.484}{\Character{59}}%
\TextAt{13}{320.794}{704.484}{\Character{11}}%
\TextAt{21}{328.302}{702.69}{n}%
\NativeMathAt{340.879}{704.484}{11.955}{\in}%
\TextAt{16}{355.775}{704.484}{\Character{79}}%
\TextAt{8}{371.662}{704.484}{is}%
\TextAt{8}{385.55}{704.484}{a}%
\TextAt{8}{397.43}{704.484}{basis}%
\TextAt{8}{428.294}{704.484}{of}%
\TextAt{13}{443.748}{704.484}{L}%
\TextAt{16}{458.653}{704.484}{\Character{106}}%
\TextAt{13}{468.915}{704.484}{K}%
\TextAt{8}{485.754}{704.484}{and}%
\TextAt{13}{510.641}{704.484}{d}%
\TextAt{9}{523.665}{704.484}{\Character{61}}%
\TextAt{8}{79.226}{687.438}{det}%
\TextAt{9}{95.485}{687.438}{\Character{40}}%
\TextAt{8}{100.038}{687.438}{Tr}%
\TextAt{9}{113.045}{687.438}{\Character{40}}%
\TextAt{13}{117.597}{687.438}{\Character{11}}%
\TextAt{21}{125.043}{685.644}{i}%
\TextAt{13}{128.424}{687.438}{\Character{11}}%
\TextAt{21}{135.87}{685.644}{j}%
\TextAt{9}{140.252}{687.438}{\Character{41}\Character{41}}%
\TextAt{8}{154.388}{687.438}{its}%
\TextAt{8}{171.841}{687.438}{discriminant,}%
\TextAt{8}{243.89}{687.438}{then}%
\TextAt{13}{271.684}{687.438}{ad}%
\TextAt{3}{285.904}{691.776}{\Character{3}}%
\TextAt{14}{290.637}{687.438}{A}%
\NativeMathAt{304.469}{687.438}{11.955}{\subseteq}%
\TextAt{16}{319.007}{687.438}{\Character{79}}%
\TextAt{8}{333.903}{687.438}{for}%
\TextAt{8}{352.922}{687.438}{every}%
\TextAt{8}{384.948}{687.438}{nonzero}%
\TextAt{13}{429.659}{687.438}{a}%
\NativeMathAt{441.058}{687.438}{11.955}{\in}%
\TextAt{14}{454.282}{687.438}{A}%
\TextAt{16}{466.291}{687.438}{\Character{92}}%
\TextAt{14}{477.691}{687.438}{o}%
\TextAt{8}{483.535}{687.438}{.}%
\TextAt{8}{495.386}{687.438}{Indeed,}%
\TextAt{8}{79.226}{670.392}{if}%
\TextAt{13}{90.768}{670.392}{x}%
\TextAt{9}{102.134}{670.392}{\Character{61}}%
\TextAt{13}{115.953}{670.392}{x}%
\TextAt{17}{122.605}{668.598}{1}%
\TextAt{13}{127.338}{670.392}{\Character{11}}%
\TextAt{17}{134.783}{668.598}{1}%
\TextAt{9}{142.73}{670.392}{\Character{43}}%
\TextAt{16}{155.048}{670.392}{\Character{1}}%
\TextAt{16}{160.344}{670.392}{\Character{1}}%
\TextAt{16}{165.652}{670.392}{\Character{1}}%
\TextAt{9}{172.21}{670.392}{\Character{43}}%
\TextAt{13}{184.529}{670.392}{x}%
\TextAt{21}{191.181}{668.598}{n}%
\TextAt{13}{196.817}{670.392}{\Character{11}}%
\TextAt{21}{204.263}{668.598}{n}%
\NativeMathAt{214.614}{670.392}{11.955}{\in}%
\TextAt{3}{227.298}{674.73}{\Character{3}}%
\TextAt{14}{232.031}{670.392}{A}%
\TextAt{8}{240.61}{670.392}{,}%
\TextAt{8}{248.781}{670.392}{with}%
\TextAt{13}{276.259}{670.392}{x}%
\TextAt{21}{282.911}{668.598}{i}%
\NativeMathAt{291.007}{670.392}{11.955}{\in}%
\TextAt{13}{303.691}{670.392}{K}%
\TextAt{8}{314.509}{670.392}{,}%
\TextAt{8}{322.68}{670.392}{then}%
\TextAt{8}{350.158}{670.392}{the}%
\TextAt{13}{371.132}{670.392}{ax}%
\TextAt{21}{383.929}{668.598}{i}%
\TextAt{8}{392.026}{670.392}{satisfy}%
\TextAt{8}{429.391}{670.392}{the}%
\TextAt{8}{450.365}{670.392}{system}%
\TextAt{8}{490.001}{670.392}{of}%
\TextAt{8}{504.149}{670.392}{linear}%
\TextAt{8}{79.226}{653.346}{equations}%
\TextAt{20}{131.65}{662.312}{\Character{80}}%
\TextAt{21}{144.269}{659.157}{n}%
\TextAt{21}{144.269}{649.859}{i}%
\TextAt{17}{147.152}{649.859}{\Character{61}1}%
\TextAt{13}{160.463}{653.346}{ax}%
\TextAt{21}{173.26}{651.552}{i}%
\TextAt{8}{178.634}{653.346}{Tr}%
\TextAt{9}{191.641}{653.346}{\Character{40}}%
\TextAt{13}{196.194}{653.346}{\Character{11}}%
\TextAt{21}{203.639}{651.552}{i}%
\TextAt{13}{207.021}{653.346}{\Character{11}}%
\TextAt{21}{214.466}{651.552}{j}%
\TextAt{9}{218.848}{653.346}{\Character{41}}%
\TextAt{9}{226.703}{653.346}{\Character{61}}%
\TextAt{8}{239.147}{653.346}{Tr}%
\TextAt{9}{252.154}{653.346}{\Character{40}}%
\TextAt{13}{256.707}{653.346}{xa\Character{11}}%
\TextAt{21}{276.949}{651.552}{j}%
\TextAt{9}{281.331}{653.346}{\Character{41}}%
\NativeMathAt{289.205}{653.346}{11.955}{\in}%
\TextAt{14}{300.496}{653.346}{o}%
\TextAt{8}{306.34}{653.346}{.}%
\TextAt{8}{314.788}{653.346}{This}%
\TextAt{8}{341.508}{653.346}{implies}%
\TextAt{13}{381.235}{653.346}{dax}%
\TextAt{21}{400.115}{651.552}{i}%
\NativeMathAt{406.817}{653.346}{11.955}{\in}%
\TextAt{14}{418.108}{653.346}{o}%
\TextAt{8}{427.849}{653.346}{and}%
\TextAt{8}{450.612}{653.346}{thus}%
\TextAt{13}{476.363}{653.346}{dax}%
\NativeMathAt{498.564}{653.346}{11.955}{\in}%
\TextAt{16}{509.838}{653.346}{\Character{79}}%
\TextAt{8}{519.705}{653.346}{.}%
\TextAt{8}{97.159}{633.311}{The}%
\TextAt{8}{121.213}{633.311}{notion}%
\TextAt{8}{157.64}{633.311}{of}%
\TextAt{8}{170.97}{633.311}{duality}%
\TextAt{8}{210.649}{633.311}{is}%
\TextAt{8}{222.413}{633.311}{justified}%
\TextAt{8}{267.018}{633.311}{by}%
\TextAt{8}{283.277}{633.311}{the}%
\TextAt{8}{303.433}{633.311}{isomorphism}%
\TextAt{3}{189.227}{611.518}{\Character{3}}%
\TextAt{14}{193.96}{606.581}{A}%
\NativeMathAt{208.447}{613.359}{7.97}{\sim}%
\NativeArrowAt{205.859}{606.581}{11.955}{13.175}{\longrightarrow}%
\TextAt{8}{222.353}{606.581}{Hom}%
\TextAt{27}{246.742}{604.788}{o}%
\TextAt{9}{251.136}{606.581}{\Character{40}}%
\TextAt{14}{255.688}{606.581}{A}%
\TextAt{13}{264.268}{606.581}{\Character{59}}%
\TextAt{14}{269.512}{606.581}{o}%
\TextAt{9}{275.356}{606.581}{\Character{41}}%
\TextAt{13}{279.908}{606.581}{\Character{59}}%
\TextAt{13}{309.043}{606.581}{x}%
\NativeArrowAt{319.036}{606.581}{11.955}{19.26}{\longmapsto}%
\TextAt{20}{341.618}{616.265}{\Character{0}}%
\TextAt{13}{347.097}{606.581}{y}%
\NativeArrowAt{356.555}{606.581}{11.955}{19.26}{\longmapsto}%
\TextAt{8}{379.137}{606.581}{Tr}%
\TextAt{9}{392.144}{606.581}{\Character{40}}%
\TextAt{13}{396.696}{606.581}{xy}%
\TextAt{9}{409.485}{606.581}{\Character{41}}%
\TextAt{20}{414.037}{616.265}{\Character{1}}%
\TextAt{13}{419.517}{606.581}{\Character{58}}%
\TextAt{8}{79.226}{579.852}{Indeed,}%
\TextAt{8}{120.394}{579.852}{since}%
\TextAt{8}{148.92}{579.852}{every}%
\TextAt{14}{179.67}{579.852}{o}%
\TextAt{8}{185.514}{579.852}{-homomorphism}%
\TextAt{13}{271.939}{579.852}{f}%
\TextAt{9}{282.307}{579.852}{\Character{58}}%
\TextAt{14}{288.879}{579.852}{A}%
\NativeMathAt{300.779}{579.852}{11.955}{\rightarrow}%
\TextAt{14}{316.055}{579.852}{o}%
\TextAt{8}{325.655}{579.852}{extends}%
\TextAt{8}{368.169}{579.852}{uniquely}%
\TextAt{8}{415.5}{579.852}{to}%
\TextAt{8}{429.669}{579.852}{a}%
\TextAt{13}{439.282}{579.852}{K}%
\TextAt{8}{450.099}{579.852}{-homomorphism}%
\TextAt{13}{79.226}{562.806}{f}%
\TextAt{9}{93.309}{562.806}{\Character{58}}%
\TextAt{13}{103.596}{562.806}{L}%
\NativeMathAt{118.597}{562.806}{11.955}{\rightarrow}%
\TextAt{13}{137.588}{562.806}{K}%
\TextAt{8}{154.483}{562.806}{in}%
\TextAt{8}{170.316}{562.806}{view}%
\TextAt{8}{199.479}{562.806}{of}%
\TextAt{14}{214.989}{562.806}{A}%
\TextAt{13}{223.568}{562.806}{K}%
\TextAt{9}{241.422}{562.806}{\Character{61}}%
\TextAt{13}{257.563}{562.806}{L}%
\TextAt{8}{265.527}{562.806}{,}%
\TextAt{8}{275.402}{562.806}{we}%
\TextAt{8}{294.809}{562.806}{may}%
\TextAt{8}{322.358}{562.806}{consider}%
\TextAt{8}{370.123}{562.806}{Hom}%
\TextAt{27}{394.512}{561.013}{o}%
\TextAt{9}{398.906}{562.806}{\Character{40}}%
\TextAt{14}{403.459}{562.806}{A}%
\TextAt{13}{412.038}{562.806}{\Character{59}}%
\TextAt{14}{417.282}{562.806}{o}%
\TextAt{9}{423.126}{562.806}{\Character{41}}%
\TextAt{8}{433.756}{562.806}{as}%
\TextAt{8}{450.306}{562.806}{a}%
\TextAt{8}{462.242}{562.806}{submodule}%
\TextAt{8}{523.337}{562.806}{of}%
\TextAt{8}{79.226}{545.76}{Hom}%
\TextAt{21}{103.615}{543.967}{K}%
\TextAt{9}{111.817}{545.76}{\Character{40}}%
\TextAt{13}{116.37}{545.76}{L\Character{59}}%
\TextAt{13}{129.557}{545.76}{K}%
\TextAt{9}{140.396}{545.76}{\Character{41}}%
\TextAt{8}{144.948}{545.76}{,}%
\TextAt{8}{153.497}{545.76}{namely,}%
\TextAt{8}{197.815}{545.76}{the}%
\TextAt{8}{219.091}{545.76}{image}%
\TextAt{8}{254.031}{545.76}{of}%
\TextAt{3}{268.48}{550.099}{\Character{3}}%
\TextAt{14}{273.213}{545.76}{A}%
\TextAt{8}{286.808}{545.76}{with}%
\TextAt{8}{314.588}{545.76}{respect}%
\TextAt{8}{355.757}{545.76}{to}%
\TextAt{13}{371.186}{545.76}{L}%
\NativeMathAt{384.379}{545.76}{11.955}{\rightarrow}%
\TextAt{8}{401.563}{545.76}{Hom}%
\TextAt{21}{425.951}{543.967}{K}%
\TextAt{9}{434.154}{545.76}{\Character{40}}%
\TextAt{13}{438.706}{545.76}{L\Character{59}}%
\TextAt{13}{451.893}{545.76}{K}%
\TextAt{9}{462.733}{545.76}{\Character{41}}%
\TextAt{8}{467.285}{545.76}{,}%
\TextAt{13}{475.833}{545.76}{x}%
\NativeArrowAt{487.714}{545.76}{11.955}{11.955}{\mapsto}%
\TextAt{9}{504.897}{545.76}{\Character{40}}%
\TextAt{13}{509.449}{545.76}{y}%
\NativeArrowAt{520.814}{545.76}{11.955}{11.955}{\mapsto}%
\TextAt{8}{79.226}{528.714}{Tr}%
\TextAt{9}{92.233}{528.714}{\Character{40}}%
\TextAt{13}{96.786}{528.714}{xy}%
\TextAt{9}{109.574}{528.714}{\Character{41}\Character{41}}%
\TextAt{8}{118.679}{528.714}{.}%
\TextAt{8}{127.127}{528.714}{The}%
\TextAt{8}{151.181}{528.714}{module}%
\TextAt{8}{192.474}{528.714}{dual}%
\TextAt{8}{218.489}{528.714}{to}%
\TextAt{16}{232.799}{528.714}{\Character{79}}%
\TextAt{8}{242.649}{528.714}{,}%
\TextAt{3}{259.698}{506.921}{\Character{3}}%
\TextAt{16}{264.43}{501.985}{\Character{79}}%
\NativeMathAt{277.601}{501.984}{11.955}{\cong}%
\TextAt{8}{290.22}{501.985}{Hom}%
\TextAt{27}{314.609}{500.191}{o}%
\TextAt{9}{319.003}{501.985}{\Character{40}}%
\TextAt{16}{323.555}{501.985}{\Character{79}}%
\TextAt{13}{333.405}{501.985}{\Character{59}}%
\TextAt{14}{338.649}{501.985}{o}%
\TextAt{9}{344.494}{501.985}{\Character{41}}%
\TextAt{13}{349.046}{501.985}{\Character{59}}%
\TextAt{8}{79.226}{475.255}{will}%
\TextAt{8}{101.331}{475.255}{obviously}%
\TextAt{8}{153.109}{475.255}{occupy}%
\TextAt{8}{192.454}{475.255}{a}%
\TextAt{8}{202.209}{475.255}{distinguished}%
\TextAt{8}{273.223}{475.255}{place}%
\TextAt{8}{303.135}{475.255}{in}%
\TextAt{8}{316.787}{475.255}{this}%
\TextAt{8}{339.61}{475.255}{theory.}%
\TextAt{12}{79.226}{449.892}{Definition}%
\TextAt{12}{142.015}{449.892}{(2.1).}%
\TextAt{8}{177.33}{449.892}{The}%
\TextAt{8}{201.384}{449.892}{fractional}%
\TextAt{8}{253.748}{449.892}{ideal}%
\TextAt{14}{216.682}{423.162}{C}%
\TextAt{3}{224.005}{421.303}{\Character{79}}%
\TextAt{3}{230.939}{421.303}{\Character{106}}%
\TextAt{27}{233.3}{421.303}{o}%
\TextAt{9}{241.015}{423.162}{\Character{61}}%
\TextAt{3}{253.441}{428.099}{\Character{3}}%
\TextAt{16}{258.173}{423.162}{\Character{79}}%
\TextAt{9}{271.344}{423.162}{\Character{61}}%
\TextAt{16}{283.769}{423.162}{\Character{102}}%
\TextAt{13}{289.747}{423.162}{x}%
\NativeMathAt{299.72}{423.162}{11.955}{\in}%
\TextAt{13}{311.011}{423.162}{L}%
\TextAt{16}{322.296}{423.162}{\Character{106}}%
\TextAt{8}{328.938}{423.162}{Tr}%
\TextAt{9}{341.945}{423.162}{\Character{40}}%
\TextAt{13}{346.498}{423.162}{x}%
\TextAt{16}{353.15}{423.162}{\Character{79}}%
\TextAt{9}{362.999}{423.162}{\Character{41}}%
\NativeMathAt{370.873}{423.162}{11.955}{\subseteq}%
\TextAt{14}{383.492}{423.162}{o}%
\TextAt{16}{389.336}{423.162}{\Character{103}}%
\TextAt{8}{79.226}{396.433}{is}%
\TextAt{8}{90.99}{396.433}{called}%
\TextAt{12}{124.154}{396.433}{Dedekind’s}%
\TextAt{12}{193.518}{396.433}{complementary}%
\TextAt{12}{288.095}{396.433}{module}%
\TextAt{8}{331.229}{396.433}{,}%
\TextAt{8}{338.378}{396.433}{or}%
\TextAt{8}{352.689}{396.433}{the}%
\TextAt{12}{372.845}{396.433}{inverse}%
\TextAt{12}{417.976}{396.433}{different}%
\TextAt{8}{467.136}{396.433}{.}%
\TextAt{8}{475.584}{396.433}{Its}%
\TextAt{8}{492.871}{396.433}{inverse,}%
\TextAt{14}{274.177}{369.703}{D}%
\TextAt{3}{284.119}{367.844}{\Character{79}}%
\TextAt{3}{291.053}{367.844}{\Character{106}}%
\TextAt{27}{293.414}{367.844}{o}%
\TextAt{9}{301.129}{369.703}{\Character{61}}%
\TextAt{14}{313.554}{369.703}{C}%
\TextAt{3}{320.878}{374.652}{\Character{0}}%
\TextAt{17}{327.464}{374.652}{1}%
\TextAt{3}{320.878}{365.931}{\Character{79}}%
\TextAt{3}{327.812}{365.931}{\Character{106}}%
\TextAt{27}{330.173}{365.931}{o}%
\TextAt{13}{334.567}{369.703}{\Character{59}}%
\TextAt{8}{79.226}{342.974}{is}%
\TextAt{8}{90.99}{342.974}{called}%
\TextAt{8}{124.154}{342.974}{the}%
\TextAt{12}{144.31}{342.974}{different}%
\TextAt{8}{197.367}{342.974}{of}%
\TextAt{16}{210.697}{342.974}{\Character{79}}%
\TextAt{16}{223.86}{342.974}{\Character{106}}%
\TextAt{14}{230.509}{342.974}{o}%
\TextAt{8}{236.354}{342.974}{.}%
\TextAt{8}{97.159}{317.611}{As}%
\TextAt{14}{114.172}{317.611}{C}%
\TextAt{3}{121.496}{315.751}{\Character{79}}%
\TextAt{3}{128.43}{315.751}{\Character{106}}%
\TextAt{27}{130.791}{315.751}{o}%
\NativeMathAt{138.506}{317.611}{11.955}{\supseteq}%
\TextAt{16}{151.107}{317.611}{\Character{79}}%
\TextAt{8}{160.975}{317.611}{,}%
\TextAt{8}{167.905}{317.611}{the}%
\TextAt{8}{187.788}{317.611}{ideal}%
\TextAt{14}{215.478}{317.611}{D}%
\TextAt{3}{225.419}{315.751}{\Character{79}}%
\TextAt{3}{232.353}{315.751}{\Character{106}}%
\TextAt{27}{234.715}{315.751}{o}%
\NativeMathAt{242.43}{317.611}{11.955}{\subseteq}%
\TextAt{16}{255.031}{317.611}{\Character{79}}%
\TextAt{8}{268.523}{317.611}{is}%
\TextAt{8}{280.013}{317.611}{actually}%
\TextAt{8}{324.296}{317.611}{an}%
\TextAt{8}{340.282}{317.611}{integral}%
\TextAt{8}{382.616}{317.611}{ideal}%
\TextAt{8}{410.306}{317.611}{of}%
\TextAt{13}{423.362}{317.611}{L}%
\TextAt{8}{431.327}{317.611}{.}%
\TextAt{8}{439.684}{317.611}{We}%
\TextAt{8}{459.555}{317.611}{will}%
\TextAt{8}{481.386}{317.611}{frequently}%
\TextAt{8}{79.226}{300.565}{denote}%
\TextAt{8}{118.301}{300.565}{it}%
\TextAt{8}{131.362}{300.565}{by}%
\TextAt{14}{148.977}{300.565}{D}%
\TextAt{21}{158.919}{298.705}{L}%
\TextAt{3}{164.678}{298.705}{\Character{106}}%
\TextAt{21}{167.03}{298.705}{K}%
\TextAt{8}{175.233}{300.565}{,}%
\TextAt{8}{184.078}{300.565}{provided}%
\TextAt{8}{233.566}{300.565}{the}%
\TextAt{8}{255.079}{300.565}{intended}%
\TextAt{8}{304.232}{300.565}{subrings}%
\TextAt{14}{351.891}{300.565}{o}%
\TextAt{13}{357.735}{300.565}{\Character{59}}%
\TextAt{16}{362.979}{300.565}{\Character{79}}%
\TextAt{8}{378.083}{300.565}{are}%
\TextAt{8}{398.95}{300.565}{evident}%
\TextAt{8}{441.277}{300.565}{from}%
\TextAt{8}{470.274}{300.565}{the}%
\TextAt{8}{491.787}{300.565}{context.}%
\TextAt{8}{79.226}{283.519}{In}%
\TextAt{8}{94.564}{283.519}{the}%
\TextAt{8}{115.438}{283.519}{same}%
\TextAt{8}{145.481}{283.519}{way,}%
\TextAt{8}{172.392}{283.519}{we}%
\TextAt{8}{190.336}{283.519}{write}%
\TextAt{14}{220.965}{283.519}{C}%
\TextAt{21}{228.288}{281.659}{L}%
\TextAt{3}{234.047}{281.659}{\Character{106}}%
\TextAt{21}{236.399}{281.659}{K}%
\TextAt{8}{249.216}{283.519}{instead}%
\TextAt{8}{290.318}{283.519}{of}%
\TextAt{14}{304.365}{283.519}{C}%
\TextAt{3}{311.688}{281.659}{\Character{79}}%
\TextAt{3}{318.622}{281.659}{\Character{106}}%
\TextAt{27}{320.983}{281.659}{o}%
\TextAt{8}{325.378}{283.519}{.}%
\TextAt{8}{335.977}{283.519}{The}%
\TextAt{8}{360.748}{283.519}{different}%
\TextAt{8}{407.636}{283.519}{behaves}%
\TextAt{8}{451.989}{283.519}{in}%
\TextAt{8}{466.359}{283.519}{the}%
\TextAt{8}{487.232}{283.519}{following}%
\TextAt{8}{79.226}{266.473}{manner}%
\TextAt{8}{121.5}{266.473}{under}%
\TextAt{8}{154.663}{266.473}{change}%
\TextAt{8}{193.362}{266.473}{of}%
\TextAt{8}{206.692}{266.473}{rings}%
\TextAt{14}{235.373}{266.473}{o}%
\TextAt{8}{245.114}{266.473}{and}%
\TextAt{16}{267.877}{266.473}{\Character{79}}%
\TextAt{8}{277.727}{266.473}{.}%
\TextAt{12}{79.226}{241.109}{Proposition}%
\TextAt{12}{151.937}{241.109}{(2.2).}%
\TextAt{8}{89.579}{221.075}{(i)}%
\TextAt{8}{107.918}{221.075}{For}%
\TextAt{8}{128.888}{221.075}{a}%
\TextAt{8}{138.643}{221.075}{tower}%
\TextAt{8}{170.516}{221.075}{of}%
\TextAt{8}{183.846}{221.075}{fields}%
\TextAt{13}{213.817}{221.075}{K}%
\NativeMathAt{227.956}{221.075}{11.955}{\subseteq}%
\TextAt{13}{240.575}{221.075}{L}%
\NativeMathAt{251.861}{221.075}{11.955}{\subseteq}%
\TextAt{13}{264.48}{221.075}{M}%
\TextAt{8}{277.053}{221.075}{,}%
\TextAt{8}{284.203}{221.075}{one}%
\TextAt{8}{305.662}{221.075}{has}%
\TextAt{14}{326.536}{221.075}{D}%
\TextAt{21}{336.478}{219.215}{M}%
\TextAt{3}{345.459}{219.215}{\Character{106}}%
\TextAt{21}{347.812}{219.215}{K}%
\TextAt{9}{359.335}{221.075}{\Character{61}}%
\TextAt{14}{371.761}{221.075}{D}%
\TextAt{21}{381.702}{219.215}{M}%
\TextAt{3}{390.684}{219.215}{\Character{106}}%
\TextAt{21}{393.036}{219.215}{L}%
\TextAt{14}{399.293}{221.075}{D}%
\TextAt{21}{409.235}{219.215}{L}%
\TextAt{3}{414.994}{219.215}{\Character{106}}%
\TextAt{21}{417.347}{219.215}{K}%
\TextAt{8}{425.549}{221.075}{.}%
\TextAt{8}{86.327}{204.029}{(ii)}%
\TextAt{8}{107.918}{204.029}{For}%
\TextAt{8}{128.888}{204.029}{any}%
\TextAt{8}{151.005}{204.029}{multiplicative}%
\TextAt{8}{224.828}{204.029}{subset}%
\TextAt{13}{260.717}{204.029}{S}%
\TextAt{8}{272.51}{204.029}{of}%
\TextAt{14}{285.84}{204.029}{o}%
\TextAt{8}{291.684}{204.029}{,}%
\TextAt{8}{298.834}{204.029}{one}%
\TextAt{8}{320.293}{204.029}{has}%
\TextAt{14}{341.167}{204.029}{D}%
\TextAt{21}{351.109}{202.025}{S}%
\TextAt{23}{356.704}{204.294}{\Character{0}}%
\TextAt{24}{362.46}{204.294}{1}%
\TextAt{3}{366.612}{202.025}{\Character{79}}%
\TextAt{3}{373.546}{202.025}{\Character{106}}%
\TextAt{21}{375.907}{202.025}{S}%
\TextAt{23}{381.503}{204.294}{\Character{0}}%
\TextAt{24}{387.259}{204.294}{1}%
\TextAt{27}{391.41}{202.025}{o}%
\TextAt{9}{399.125}{204.029}{\Character{61}}%
\TextAt{13}{411.55}{204.029}{S}%
\TextAt{3}{419.446}{208.367}{\Character{0}}%
\TextAt{17}{426.032}{208.367}{1}%
\TextAt{14}{430.764}{204.029}{D}%
\TextAt{3}{440.706}{202.169}{\Character{79}}%
\TextAt{3}{447.64}{202.169}{\Character{106}}%
\TextAt{27}{450.001}{202.169}{o}%
\TextAt{8}{454.396}{204.029}{.}%
\TextAt{8}{83.076}{186.983}{(iii)}%
\TextAt{8}{107.918}{186.983}{If}%
\TextAt{14}{119.729}{186.983}{P}%
\TextAt{16}{133.144}{186.983}{\Character{106}}%
\TextAt{14}{139.987}{186.983}{p}%
\TextAt{8}{149.986}{186.983}{are}%
\TextAt{8}{169.615}{186.983}{prime}%
\TextAt{8}{202.896}{186.983}{ideals}%
\TextAt{8}{235.592}{186.983}{of}%
\TextAt{16}{249.04}{186.983}{\Character{79}}%
\TextAt{8}{258.89}{186.983}{,}%
\TextAt{8}{266.187}{186.983}{resp.}%
\TextAt{14}{295.863}{186.983}{o}%
\TextAt{8}{301.708}{186.983}{,}%
\TextAt{8}{309.004}{186.983}{and}%
\TextAt{16}{331.885}{186.983}{\Character{79}}%
\TextAt{27}{341.403}{185.189}{P}%
\TextAt{16}{352.019}{186.983}{\Character{106}}%
\TextAt{14}{358.862}{186.983}{o}%
\TextAt{27}{364.706}{185.189}{p}%
\TextAt{8}{373.208}{186.983}{are}%
\TextAt{8}{392.837}{186.983}{the}%
\TextAt{8}{413.112}{186.983}{associated}%
\TextAt{8}{468.965}{186.983}{completions,}%
\TextAt{8}{107.918}{169.937}{then}%
\TextAt{14}{273.886}{152.891}{D}%
\TextAt{3}{283.827}{151.031}{\Character{79}}%
\TextAt{3}{290.761}{151.031}{\Character{106}}%
\TextAt{27}{293.123}{151.031}{o}%
\TextAt{16}{297.517}{152.891}{\Character{79}}%
\TextAt{27}{307.035}{151.098}{P}%
\TextAt{9}{317.449}{152.891}{\Character{61}}%
\TextAt{14}{329.875}{152.891}{D}%
\TextAt{3}{339.816}{151.031}{\Character{79}}%
\TextAt{28}{346.524}{149.64}{P}%
\TextAt{3}{352.721}{151.031}{\Character{106}}%
\TextAt{27}{355.074}{151.031}{o}%
\TextAt{28}{358.97}{150.035}{p}%
\TextAt{13}{363.551}{152.891}{\Character{58}}%
\TextAt{19}{79.226}{116.697}{Proof.}%
\TextAt{8}{116}{116.697}{(i)}%
\TextAt{8}{132.236}{116.697}{Let}%
\TextAt{13}{153.183}{116.697}{A}%
\TextAt{9}{165.279}{116.697}{\Character{61}}%
\TextAt{14}{177.705}{116.697}{o}%
\NativeMathAt{186.87}{116.697}{11.955}{\subseteq}%
\TextAt{13}{199.489}{116.697}{K}%
\TextAt{8}{210.307}{116.697}{,}%
\TextAt{8}{217.438}{116.697}{and}%
\TextAt{8}{240.178}{116.697}{let}%
\TextAt{13}{257.059}{116.697}{B}%
\NativeMathAt{269.877}{116.697}{11.955}{\subseteq}%
\TextAt{13}{282.496}{116.697}{L}%
\TextAt{8}{290.46}{116.697}{,}%
\TextAt{13}{297.591}{116.697}{C}%
\NativeMathAt{310.146}{116.697}{11.955}{\subseteq}%
\TextAt{13}{322.765}{116.697}{M}%
\TextAt{8}{339.213}{116.697}{be}%
\TextAt{8}{355.115}{116.697}{the}%
\TextAt{8}{375.248}{116.697}{integral}%
\TextAt{8}{417.834}{116.697}{closure}%
\TextAt{8}{456.892}{116.697}{of}%
\TextAt{14}{470.2}{116.697}{o}%
\TextAt{8}{479.919}{116.697}{in}%
\TextAt{13}{493.549}{116.697}{L}%
\TextAt{8}{501.513}{116.697}{,}%
\TextAt{8}{508.644}{116.697}{resp.}%
\TextAt{13}{79.226}{99.651}{M}%
\TextAt{8}{91.8}{99.651}{.}%
\TextAt{8}{100.248}{99.651}{It}%
\TextAt{8}{112.921}{99.651}{then}%
\TextAt{8}{139.581}{99.651}{suffices}%
\TextAt{8}{179.367}{99.651}{to}%
\TextAt{8}{193.678}{99.651}{show}%
\TextAt{8}{222.681}{99.651}{that}%
\TextAt{14}{261.517}{72.921}{C}%
\TextAt{21}{268.84}{71.062}{C}%
\TextAt{3}{275.438}{71.062}{\Character{106}}%
\TextAt{21}{277.791}{71.062}{A}%
\TextAt{9}{287.953}{72.921}{\Character{61}}%
\TextAt{14}{300.378}{72.921}{C}%
\TextAt{21}{307.701}{71.062}{C}%
\TextAt{3}{314.3}{71.062}{\Character{106}}%
\TextAt{21}{316.652}{71.062}{B}%
\TextAt{14}{323.931}{72.921}{C}%
\TextAt{21}{331.254}{71.062}{B}%
\TextAt{3}{338.034}{71.062}{\Character{106}}%
\TextAt{21}{340.386}{71.062}{A}%
\TextAt{13}{347.227}{72.921}{\Character{58}}%
\TextAt{8}{297.211}{43.033}{178}%
\Destination{79.23}{235.115}{nav.Item.70}%
\Destination{79.23}{217.225}{nav.Item.71}%
\Destination{79.23}{200.035}{nav.Item.72}%
\Destination{78.23}{809.815}{nav.page.178}%
\Destination{79.36}{621.555}{nav.print.195}%
\Destination{79.36}{463.855}{nav.number.3.2.1}%
\Destination{79.36}{255.155}{nav.number.3.2.2}%
\Hotlink{140.178}{758.955}{5.22}{11.047}{nav.chapter.1}%
\Hotlink{158.675}{758.955}{6.858}{11.047}{nav.section.1.2}%
}

% Page 191 (body page 179)
\Page{%
\TextAt{2}{79.226}{796.943}{III.}%
\TextAt{2}{98.317}{796.943}{Riemann–Roch}%
\TextAt{2}{174.375}{796.943}{Theory}%
\Vector{q 0 G 0 g 0.299 w 7.226 18.729 m 460.769 18.729 l S Q}%
\TextAt{8}{79.226}{761.124}{The}%
\TextAt{8}{103.28}{761.124}{inclusion}%
\NativeMathAt{152.117}{761.124}{11.955}{\supseteq}%
\TextAt{8}{165.313}{761.124}{follows}%
\TextAt{8}{203.749}{761.124}{from}%
\TextAt{8}{94.823}{735.282}{Tr}%
\TextAt{21}{107.831}{733.422}{M}%
\TextAt{3}{116.812}{733.422}{\Character{106}}%
\TextAt{21}{119.165}{733.422}{K}%
\TextAt{9}{127.367}{735.282}{\Character{40}}%
\TextAt{14}{131.92}{735.282}{C}%
\TextAt{21}{139.243}{733.422}{C}%
\TextAt{3}{145.841}{733.422}{\Character{106}}%
\TextAt{21}{148.194}{733.422}{B}%
\TextAt{14}{155.472}{735.282}{C}%
\TextAt{21}{162.795}{733.422}{B}%
\TextAt{3}{169.575}{733.422}{\Character{106}}%
\TextAt{21}{171.928}{733.422}{A}%
\TextAt{13}{178.769}{735.282}{C}%
\TextAt{9}{188.002}{735.282}{\Character{41}}%
\TextAt{9}{195.857}{735.282}{\Character{61}}%
\TextAt{8}{208.301}{735.282}{Tr}%
\TextAt{21}{221.308}{733.422}{L}%
\TextAt{3}{227.067}{733.422}{\Character{106}}%
\TextAt{21}{229.42}{733.422}{K}%
\TextAt{8}{239.615}{735.282}{Tr}%
\TextAt{21}{252.622}{733.422}{M}%
\TextAt{3}{261.604}{733.422}{\Character{106}}%
\TextAt{21}{263.956}{733.422}{L}%
\TextAt{9}{270.213}{735.282}{\Character{40}}%
\TextAt{14}{274.765}{735.282}{C}%
\TextAt{21}{282.089}{733.422}{C}%
\TextAt{3}{288.687}{733.422}{\Character{106}}%
\TextAt{21}{291.039}{733.422}{B}%
\TextAt{14}{298.318}{735.282}{C}%
\TextAt{21}{305.641}{733.422}{B}%
\TextAt{3}{312.421}{733.422}{\Character{106}}%
\TextAt{21}{314.773}{733.422}{A}%
\TextAt{13}{321.615}{735.282}{C}%
\TextAt{9}{330.848}{735.282}{\Character{41}}%
\TextAt{9}{338.703}{735.282}{\Character{61}}%
\TextAt{8}{351.147}{735.282}{Tr}%
\TextAt{21}{364.154}{733.422}{L}%
\TextAt{3}{369.913}{733.422}{\Character{106}}%
\TextAt{21}{372.265}{733.422}{K}%
\TextAt{20}{380.468}{744.965}{\Character{0}}%
\TextAt{14}{385.948}{735.282}{C}%
\TextAt{21}{393.271}{733.422}{B}%
\TextAt{3}{400.051}{733.422}{\Character{106}}%
\TextAt{21}{402.403}{733.422}{A}%
\TextAt{8}{411.237}{735.282}{Tr}%
\TextAt{21}{424.244}{733.422}{M}%
\TextAt{3}{433.226}{733.422}{\Character{106}}%
\TextAt{21}{435.578}{733.422}{L}%
\TextAt{9}{441.835}{735.282}{\Character{40}}%
\TextAt{14}{446.387}{735.282}{C}%
\TextAt{21}{453.711}{733.422}{C}%
\TextAt{3}{460.309}{733.422}{\Character{106}}%
\TextAt{21}{462.662}{733.422}{B}%
\TextAt{13}{469.94}{735.282}{C}%
\TextAt{9}{479.173}{735.282}{\Character{41}}%
\TextAt{20}{483.726}{744.965}{\Character{1}}%
\NativeMathAt{492.526}{735.282}{11.955}{\subseteq}%
\TextAt{13}{505.145}{735.282}{A\Character{58}}%
\TextAt{8}{79.226}{709.44}{In}%
\TextAt{8}{93.847}{709.44}{view}%
\TextAt{8}{120.83}{709.44}{of}%
\TextAt{13}{134.16}{709.44}{B}%
\TextAt{13}{143.641}{709.44}{C}%
\TextAt{9}{156.211}{709.44}{\Character{61}}%
\TextAt{13}{168.636}{709.44}{C}%
\TextAt{8}{177.87}{709.44}{,}%
\TextAt{8}{185.019}{709.44}{the}%
\TextAt{8}{205.176}{709.44}{inclusion}%
\NativeMathAt{254.012}{709.44}{11.955}{\subseteq}%
\TextAt{8}{267.208}{709.44}{is}%
\TextAt{8}{278.972}{709.44}{derived}%
\TextAt{8}{319.943}{709.44}{as}%
\TextAt{8}{334.313}{709.44}{follows:}%
\TextAt{8}{185.232}{683.598}{Tr}%
\TextAt{21}{198.239}{681.738}{M}%
\TextAt{3}{207.221}{681.738}{\Character{106}}%
\TextAt{21}{209.573}{681.738}{K}%
\TextAt{9}{217.776}{683.598}{\Character{40}}%
\TextAt{14}{222.328}{683.598}{C}%
\TextAt{21}{229.651}{681.738}{C}%
\TextAt{3}{236.25}{681.738}{\Character{106}}%
\TextAt{21}{238.602}{681.738}{A}%
\TextAt{13}{245.443}{683.598}{C}%
\TextAt{9}{254.677}{683.598}{\Character{41}}%
\TextAt{9}{262.532}{683.598}{\Character{61}}%
\TextAt{8}{274.976}{683.598}{Tr}%
\TextAt{21}{287.983}{681.738}{L}%
\TextAt{3}{293.742}{681.738}{\Character{106}}%
\TextAt{21}{296.094}{681.738}{K}%
\TextAt{20}{304.297}{693.282}{\Character{0}}%
\TextAt{13}{309.776}{683.598}{B}%
\TextAt{8}{321.265}{683.598}{Tr}%
\TextAt{21}{334.272}{681.738}{M}%
\TextAt{3}{343.254}{681.738}{\Character{106}}%
\TextAt{21}{345.606}{681.738}{L}%
\TextAt{9}{351.864}{683.598}{\Character{40}}%
\TextAt{14}{356.416}{683.598}{C}%
\TextAt{21}{363.739}{681.738}{C}%
\TextAt{3}{370.338}{681.738}{\Character{106}}%
\TextAt{21}{372.69}{681.738}{A}%
\TextAt{13}{379.531}{683.598}{C}%
\TextAt{9}{388.765}{683.598}{\Character{41}}%
\TextAt{20}{393.317}{693.282}{\Character{1}}%
\NativeMathAt{402.117}{683.598}{11.955}{\subseteq}%
\TextAt{13}{414.737}{683.598}{A\Character{59}}%
\TextAt{8}{79.226}{657.756}{so}%
\TextAt{8}{93.596}{657.756}{that}%
\TextAt{8}{118.965}{657.756}{Tr}%
\TextAt{21}{131.972}{655.896}{M}%
\TextAt{3}{140.954}{655.896}{\Character{106}}%
\TextAt{21}{143.306}{655.896}{L}%
\TextAt{9}{149.564}{657.756}{\Character{40}}%
\TextAt{14}{154.116}{657.756}{C}%
\TextAt{21}{161.439}{655.896}{C}%
\TextAt{3}{168.038}{655.896}{\Character{106}}%
\TextAt{21}{170.39}{655.896}{A}%
\TextAt{13}{177.231}{657.756}{C}%
\TextAt{9}{186.465}{657.756}{\Character{41}}%
\NativeMathAt{194.338}{657.756}{11.955}{\subseteq}%
\TextAt{14}{206.957}{657.756}{C}%
\TextAt{21}{214.28}{655.896}{B}%
\TextAt{3}{221.06}{655.896}{\Character{106}}%
\TextAt{21}{223.413}{655.896}{A}%
\TextAt{8}{230.254}{657.756}{,}%
\TextAt{8}{237.403}{657.756}{and}%
\TextAt{8}{260.166}{657.756}{thus}%
\TextAt{8}{187.436}{631.914}{Tr}%
\TextAt{21}{200.443}{630.054}{M}%
\TextAt{3}{209.425}{630.054}{\Character{106}}%
\TextAt{21}{211.777}{630.054}{L}%
\TextAt{9}{218.034}{631.914}{\Character{40}}%
\TextAt{14}{222.587}{631.914}{C}%
\TextAt{3}{229.91}{636.862}{\Character{0}}%
\TextAt{17}{236.496}{636.862}{1}%
\TextAt{21}{229.91}{628.141}{B}%
\TextAt{3}{236.69}{628.141}{\Character{106}}%
\TextAt{21}{239.042}{628.141}{A}%
\TextAt{14}{245.883}{631.914}{C}%
\TextAt{21}{253.206}{630.054}{C}%
\TextAt{3}{259.805}{630.054}{\Character{106}}%
\TextAt{21}{262.157}{630.054}{A}%
\TextAt{13}{268.998}{631.914}{C}%
\TextAt{9}{278.232}{631.914}{\Character{41}}%
\TextAt{9}{286.087}{631.914}{\Character{61}}%
\TextAt{14}{298.531}{631.914}{C}%
\TextAt{3}{305.854}{636.862}{\Character{0}}%
\TextAt{17}{312.44}{636.862}{1}%
\TextAt{21}{305.854}{628.141}{B}%
\TextAt{3}{312.634}{628.141}{\Character{106}}%
\TextAt{21}{314.986}{628.141}{A}%
\TextAt{8}{323.82}{631.914}{Tr}%
\TextAt{21}{336.827}{630.054}{M}%
\TextAt{3}{345.809}{630.054}{\Character{106}}%
\TextAt{21}{348.161}{630.054}{L}%
\TextAt{9}{354.418}{631.914}{\Character{40}}%
\TextAt{14}{358.971}{631.914}{C}%
\TextAt{21}{366.294}{630.054}{C}%
\TextAt{3}{372.892}{630.054}{\Character{106}}%
\TextAt{21}{375.245}{630.054}{A}%
\TextAt{13}{382.086}{631.914}{C}%
\TextAt{9}{391.319}{631.914}{\Character{41}}%
\NativeMathAt{399.193}{631.914}{11.955}{\subseteq}%
\TextAt{13}{411.812}{631.914}{B}%
\TextAt{13}{421.293}{631.914}{\Character{58}}%
\TextAt{8}{79.226}{606.072}{This}%
\TextAt{8}{105.946}{606.072}{does}%
\TextAt{8}{132.343}{606.072}{indeed}%
\TextAt{8}{169.404}{606.072}{imply}%
\TextAt{14}{202.245}{606.072}{C}%
\TextAt{3}{209.568}{611.02}{\Character{0}}%
\TextAt{17}{216.154}{611.02}{1}%
\TextAt{21}{209.568}{602.299}{B}%
\TextAt{3}{216.348}{602.299}{\Character{106}}%
\TextAt{21}{218.7}{602.299}{A}%
\TextAt{14}{225.542}{606.072}{C}%
\TextAt{21}{232.865}{604.212}{C}%
\TextAt{3}{239.463}{604.212}{\Character{106}}%
\TextAt{21}{241.816}{604.212}{A}%
\NativeMathAt{251.978}{606.072}{11.955}{\subseteq}%
\TextAt{14}{264.597}{606.072}{C}%
\TextAt{21}{271.92}{604.212}{C}%
\TextAt{3}{278.519}{604.212}{\Character{106}}%
\TextAt{21}{280.871}{604.212}{B}%
\TextAt{8}{288.149}{606.072}{,}%
\TextAt{8}{295.298}{606.072}{and}%
\TextAt{8}{318.061}{606.072}{so}%
\TextAt{14}{332.431}{606.072}{C}%
\TextAt{21}{339.754}{604.212}{C}%
\TextAt{3}{346.353}{604.212}{\Character{106}}%
\TextAt{21}{348.705}{604.212}{A}%
\NativeMathAt{358.867}{606.072}{11.955}{\subseteq}%
\TextAt{14}{371.487}{606.072}{C}%
\TextAt{21}{378.81}{604.212}{C}%
\TextAt{3}{385.408}{604.212}{\Character{106}}%
\TextAt{21}{387.761}{604.212}{B}%
\TextAt{14}{395.039}{606.072}{C}%
\TextAt{21}{402.362}{604.212}{B}%
\TextAt{3}{409.142}{604.212}{\Character{106}}%
\TextAt{21}{411.495}{604.212}{A}%
\TextAt{8}{418.336}{606.072}{.}%
\TextAt{8}{97.159}{586.037}{(ii)}%
\TextAt{8}{116.67}{586.037}{is}%
\TextAt{8}{128.434}{586.037}{trivial.}%
\TextAt{8}{97.159}{566.002}{(iii)}%
\TextAt{8}{120.203}{566.002}{By}%
\TextAt{8}{138.847}{566.002}{(ii)}%
\TextAt{8}{158.639}{566.002}{we}%
\TextAt{8}{176.148}{566.002}{may}%
\TextAt{8}{201.798}{566.002}{assume}%
\TextAt{8}{242.523}{566.002}{that}%
\TextAt{14}{268.173}{566.002}{o}%
\TextAt{8}{278.196}{566.002}{is}%
\TextAt{8}{290.241}{566.002}{a}%
\TextAt{8}{300.278}{566.002}{discrete}%
\TextAt{8}{343.538}{566.002}{valuation}%
\TextAt{8}{394.891}{566.002}{ring.}%
\TextAt{8}{424.352}{566.002}{We}%
\TextAt{8}{444.777}{566.002}{show}%
\TextAt{8}{474.062}{566.002}{that}%
\TextAt{14}{499.712}{566.002}{C}%
\TextAt{3}{507.035}{564.143}{\Character{79}}%
\TextAt{3}{513.969}{564.143}{\Character{106}}%
\TextAt{27}{516.33}{564.143}{o}%
\TextAt{8}{524.903}{566.002}{is}%
\TextAt{8}{79.226}{548.956}{dense}%
\TextAt{8}{111.146}{548.956}{in}%
\TextAt{14}{124.799}{548.956}{C}%
\TextAt{3}{132.122}{547.097}{\Character{79}}%
\TextAt{28}{138.83}{545.706}{P}%
\TextAt{3}{145.027}{547.097}{\Character{106}}%
\TextAt{27}{147.38}{547.097}{o}%
\TextAt{28}{151.276}{546.1}{p}%
\TextAt{8}{155.857}{548.956}{.}%
\TextAt{8}{164.305}{548.956}{In}%
\TextAt{8}{178.926}{548.956}{order}%
\TextAt{8}{209.495}{548.956}{to}%
\TextAt{8}{223.806}{548.956}{do}%
\TextAt{8}{240.065}{548.956}{this,}%
\TextAt{8}{266.139}{548.956}{we}%
\TextAt{8}{283.366}{548.956}{use}%
\TextAt{8}{303.583}{548.956}{the}%
\TextAt{8}{323.739}{548.956}{formula}%
\TextAt{8}{254.314}{523.114}{Tr}%
\TextAt{21}{267.322}{521.255}{L}%
\TextAt{3}{273.081}{521.255}{\Character{106}}%
\TextAt{21}{275.433}{521.255}{K}%
\TextAt{9}{286.957}{523.114}{\Character{61}}%
\TextAt{20}{299.382}{534.472}{\Character{88}}%
\TextAt{27}{301.548}{508.569}{P}%
\TextAt{3}{308.144}{508.569}{\Character{106}}%
\TextAt{27}{310.496}{508.569}{p}%
\TextAt{8}{318.643}{523.114}{Tr}%
\TextAt{21}{331.65}{521.255}{L}%
\TextAt{28}{337.409}{519.864}{P}%
\TextAt{3}{343.607}{521.255}{\Character{106}}%
\TextAt{21}{345.959}{521.255}{K}%
\TextAt{28}{353.1}{520.258}{p}%
\TextAt{8}{79.226}{486.622}{(see}%
\TextAt{8}{102.742}{486.622}{chap.}%
\TextAt{8}{135.078}{486.622}{II,}%
\TextAt{8}{151.038}{486.622}{(8.4)).}%
\TextAt{8}{188.263}{486.622}{Let}%
\TextAt{13}{209.281}{486.622}{x}%
\NativeMathAt{219.335}{486.622}{11.955}{\in}%
\TextAt{14}{230.708}{486.622}{C}%
\TextAt{3}{238.031}{484.763}{\Character{79}}%
\TextAt{3}{244.965}{484.763}{\Character{106}}%
\TextAt{27}{247.327}{484.763}{o}%
\TextAt{8}{255.666}{486.622}{and}%
\TextAt{13}{278.477}{486.622}{y}%
\NativeMathAt{288.016}{486.622}{11.955}{\in}%
\TextAt{16}{299.373}{486.622}{\Character{79}}%
\TextAt{27}{308.907}{484.829}{P}%
\TextAt{8}{316.001}{486.622}{.}%
\TextAt{8}{324.593}{486.622}{The}%
\TextAt{8}{348.695}{486.622}{approximation}%
\TextAt{8}{426.81}{486.622}{theorem}%
\TextAt{8}{472.384}{486.622}{allows}%
\TextAt{8}{507.293}{486.622}{us}%
\TextAt{8}{522.357}{486.622}{to}%
\TextAt{8}{79.226}{469.576}{find}%
\TextAt{8}{102.877}{469.576}{an}%
\TextAt{13}{119.379}{469.576}{\Character{17}}%
\TextAt{8}{129.724}{469.576}{in}%
\TextAt{13}{143.62}{469.576}{L}%
\TextAt{8}{155.725}{469.576}{which}%
\TextAt{8}{189.454}{469.576}{is}%
\TextAt{8}{201.461}{469.576}{close}%
\TextAt{8}{229.727}{469.576}{to}%
\TextAt{13}{244.281}{469.576}{y}%
\TextAt{8}{254.558}{469.576}{with}%
\TextAt{8}{281.461}{469.576}{respect}%
\TextAt{8}{321.754}{469.576}{to}%
\TextAt{13}{336.307}{469.576}{v}%
\TextAt{27}{341.966}{467.783}{P}%
\TextAt{8}{349.06}{469.576}{,}%
\TextAt{8}{356.513}{469.576}{and}%
\TextAt{8}{379.519}{469.576}{close}%
\TextAt{8}{407.784}{469.576}{to}%
\TextAt{9}{422.338}{469.576}{0}%
\TextAt{8}{432.331}{469.576}{with}%
\TextAt{8}{459.234}{469.576}{respect}%
\TextAt{8}{499.527}{469.576}{to}%
\TextAt{13}{514.08}{469.576}{v}%
\TextAt{27}{519.74}{467.783}{P}%
\TextAt{23}{526.335}{470.052}{\Character{48}}%
\TextAt{8}{529.518}{469.576}{,}%
\TextAt{8}{79.226}{452.53}{for}%
\TextAt{14}{97.111}{452.53}{P}%
\TextAt{3}{107.004}{456.869}{\Character{48}}%
\TextAt{16}{113.12}{452.53}{\Character{106}}%
\TextAt{14}{119.762}{452.53}{p}%
\TextAt{8}{125.745}{452.53}{,}%
\TextAt{14}{132.894}{452.53}{P}%
\TextAt{3}{142.787}{456.869}{\Character{48}}%
\NativeMathAt{148.903}{452.53}{11.955}{\neq}%
\TextAt{14}{161.329}{452.53}{P}%
\TextAt{8}{171.222}{452.53}{.}%
\TextAt{8}{179.67}{452.53}{The}%
\TextAt{8}{203.724}{452.53}{left-hand}%
\TextAt{8}{253.47}{452.53}{side}%
\TextAt{8}{276.938}{452.53}{of}%
\TextAt{8}{290.268}{452.53}{the}%
\TextAt{8}{310.424}{452.53}{equation}%
\TextAt{8}{190.706}{426.688}{Tr}%
\TextAt{21}{203.713}{424.829}{L}%
\TextAt{3}{209.472}{424.829}{\Character{106}}%
\TextAt{21}{211.824}{424.829}{K}%
\TextAt{9}{220.027}{426.688}{\Character{40}}%
\TextAt{13}{224.579}{426.688}{x\Character{17}}%
\TextAt{9}{237.436}{426.688}{\Character{41}}%
\TextAt{9}{245.291}{426.688}{\Character{61}}%
\TextAt{8}{257.735}{426.688}{Tr}%
\TextAt{21}{270.742}{424.829}{L}%
\TextAt{28}{276.501}{423.438}{P}%
\TextAt{3}{282.698}{424.829}{\Character{106}}%
\TextAt{21}{285.05}{424.829}{K}%
\TextAt{28}{292.192}{423.832}{p}%
\TextAt{9}{296.773}{426.688}{\Character{40}}%
\TextAt{13}{301.325}{426.688}{x\Character{17}}%
\TextAt{9}{314.182}{426.688}{\Character{41}}%
\TextAt{9}{321.379}{426.688}{\Character{43}}%
\TextAt{20}{335.748}{438.046}{\Character{88}}%
\TextAt{27}{333.152}{412.53}{P}%
\TextAt{23}{339.747}{414.799}{\Character{48}}%
\NativeMathAt{342.432}{412.53}{7.97}{\neq}%
\TextAt{27}{349.018}{412.53}{P}%
\TextAt{8}{357.606}{426.688}{Tr}%
\TextAt{21}{370.613}{424.829}{L}%
\TextAt{28}{376.372}{422.537}{P}%
\TextAt{23}{382.071}{424.253}{\Character{48}}%
\TextAt{3}{385.254}{424.829}{\Character{106}}%
\TextAt{21}{387.606}{424.829}{K}%
\TextAt{28}{394.748}{423.832}{p}%
\TextAt{9}{399.329}{426.688}{\Character{40}}%
\TextAt{13}{403.881}{426.688}{x\Character{17}}%
\TextAt{9}{416.737}{426.688}{\Character{41}}%
\TextAt{8}{79.226}{391.026}{then}%
\TextAt{8}{107.579}{391.026}{belongs}%
\TextAt{8}{151.282}{391.026}{to}%
\TextAt{14}{167.285}{391.026}{o}%
\TextAt{27}{173.13}{389.233}{p}%
\TextAt{8}{177.616}{391.026}{,}%
\TextAt{8}{186.882}{391.026}{since}%
\TextAt{8}{217.243}{391.026}{Tr}%
\TextAt{21}{230.25}{389.167}{L}%
\TextAt{3}{236.009}{389.167}{\Character{106}}%
\TextAt{21}{238.362}{389.167}{K}%
\TextAt{9}{246.564}{391.026}{\Character{40}}%
\TextAt{13}{251.117}{391.026}{x\Character{17}}%
\TextAt{9}{263.973}{391.026}{\Character{41}}%
\NativeMathAt{274.731}{391.026}{11.955}{\in}%
\TextAt{14}{288.907}{391.026}{o}%
\NativeMathAt{300.957}{391.026}{11.955}{\subseteq}%
\TextAt{14}{316.461}{391.026}{o}%
\TextAt{27}{322.305}{389.233}{p}%
\TextAt{8}{326.792}{391.026}{,}%
\TextAt{8}{336.057}{391.026}{but}%
\TextAt{8}{359.209}{391.026}{the}%
\TextAt{8}{381.059}{391.026}{same}%
\TextAt{8}{412.077}{391.026}{is}%
\TextAt{8}{425.534}{391.026}{true}%
\TextAt{8}{451.938}{391.026}{of}%
\TextAt{8}{466.961}{391.026}{the}%
\TextAt{8}{488.81}{391.026}{elements}%
\TextAt{8}{79.226}{373.98}{Tr}%
\TextAt{21}{92.233}{372.121}{L}%
\TextAt{28}{97.992}{369.829}{P}%
\TextAt{23}{103.691}{371.545}{\Character{48}}%
\TextAt{3}{106.874}{372.121}{\Character{106}}%
\TextAt{21}{109.226}{372.121}{K}%
\TextAt{28}{116.368}{371.124}{p}%
\TextAt{9}{120.949}{373.98}{\Character{40}}%
\TextAt{13}{125.501}{373.98}{x\Character{17}}%
\TextAt{9}{138.357}{373.98}{\Character{41}}%
\TextAt{8}{145.751}{373.98}{because}%
\TextAt{8}{187.996}{373.98}{they}%
\TextAt{8}{213.277}{373.98}{are}%
\TextAt{8}{231.732}{373.98}{close}%
\TextAt{8}{258.699}{373.98}{to}%
\TextAt{8}{271.953}{373.98}{zero}%
\TextAt{8}{295.608}{373.98}{with}%
\TextAt{8}{321.212}{373.98}{respect}%
\TextAt{8}{360.206}{373.98}{to}%
\TextAt{13}{373.46}{373.98}{v}%
\TextAt{27}{379.119}{372.187}{p}%
\TextAt{8}{383.606}{373.98}{.}%
\TextAt{8}{391.702}{373.98}{Therefore}%
\TextAt{8}{443.643}{373.98}{Tr}%
\TextAt{21}{456.651}{372.121}{L}%
\TextAt{28}{462.41}{370.73}{P}%
\TextAt{3}{468.607}{372.121}{\Character{106}}%
\TextAt{21}{470.959}{372.121}{K}%
\TextAt{28}{478.101}{371.124}{p}%
\TextAt{9}{482.681}{373.98}{\Character{40}}%
\TextAt{13}{487.234}{373.98}{xy}%
\TextAt{9}{500.022}{373.98}{\Character{41}}%
\NativeMathAt{507.896}{373.98}{11.955}{\in}%
\TextAt{14}{519.187}{373.98}{o}%
\TextAt{27}{525.031}{372.187}{p}%
\TextAt{8}{529.518}{373.98}{.}%
\TextAt{8}{79.226}{356.934}{This}%
\TextAt{8}{105.946}{356.934}{shows}%
\TextAt{8}{139.564}{356.934}{that}%
\TextAt{14}{164.933}{356.934}{C}%
\TextAt{3}{172.256}{355.075}{\Character{79}}%
\TextAt{3}{179.19}{355.075}{\Character{106}}%
\TextAt{27}{181.551}{355.075}{o}%
\NativeMathAt{189.266}{356.934}{11.955}{\subseteq}%
\TextAt{14}{201.886}{356.934}{C}%
\TextAt{3}{209.209}{355.075}{\Character{79}}%
\TextAt{28}{215.916}{353.684}{P}%
\TextAt{3}{222.113}{355.075}{\Character{106}}%
\TextAt{27}{224.466}{355.075}{o}%
\TextAt{28}{228.362}{354.078}{p}%
\TextAt{8}{232.943}{356.934}{.}%
\TextAt{8}{97.159}{336.9}{If}%
\TextAt{8}{109.599}{336.9}{on}%
\TextAt{8}{126.605}{336.9}{the}%
\TextAt{8}{147.509}{336.9}{other}%
\TextAt{8}{178.826}{336.9}{hand}%
\TextAt{13}{208.84}{336.9}{x}%
\NativeMathAt{220.086}{336.9}{11.955}{\in}%
\TextAt{14}{232.651}{336.9}{C}%
\TextAt{3}{239.974}{335.04}{\Character{79}}%
\TextAt{28}{246.682}{333.649}{P}%
\TextAt{3}{252.879}{335.04}{\Character{106}}%
\TextAt{27}{255.232}{335.04}{o}%
\TextAt{28}{259.128}{334.044}{p}%
\TextAt{8}{263.709}{336.9}{,}%
\TextAt{8}{271.792}{336.9}{and}%
\TextAt{8}{295.302}{336.9}{if}%
\TextAt{13}{306.774}{336.9}{\Character{24}}%
\NativeMathAt{317.059}{336.9}{11.955}{\in}%
\TextAt{13}{329.624}{336.9}{L}%
\TextAt{8}{342.233}{336.9}{is}%
\TextAt{8}{354.744}{336.9}{sufficiently}%
\TextAt{8}{414.084}{336.9}{close}%
\TextAt{8}{442.855}{336.9}{to}%
\TextAt{13}{457.913}{336.9}{x}%
\TextAt{8}{469.21}{336.9}{with}%
\TextAt{8}{496.617}{336.9}{respect}%
\TextAt{8}{79.226}{319.854}{to}%
\TextAt{13}{93.749}{319.854}{v}%
\TextAt{27}{99.409}{318.06}{P}%
\TextAt{8}{106.502}{319.854}{,}%
\TextAt{8}{113.917}{319.854}{and}%
\TextAt{8}{136.893}{319.854}{sufficiently}%
\TextAt{8}{195.698}{319.854}{close}%
\TextAt{8}{223.933}{319.854}{to}%
\TextAt{9}{238.457}{319.854}{0}%
\TextAt{8}{248.42}{319.854}{with}%
\TextAt{8}{275.292}{319.854}{respect}%
\TextAt{8}{315.555}{319.854}{to}%
\TextAt{13}{330.078}{319.854}{v}%
\TextAt{27}{335.737}{318.06}{P}%
\TextAt{23}{342.333}{320.33}{\Character{48}}%
\TextAt{8}{345.515}{319.854}{,}%
\TextAt{8}{352.931}{319.854}{for}%
\TextAt{14}{371.028}{319.854}{P}%
\TextAt{3}{380.921}{324.192}{\Character{48}}%
\NativeMathAt{387.4}{319.854}{11.955}{\neq}%
\TextAt{14}{400.188}{319.854}{P}%
\TextAt{8}{410.081}{319.854}{,}%
\TextAt{8}{417.496}{319.854}{then}%
\TextAt{13}{444.369}{319.854}{\Character{24}}%
\NativeMathAt{453.743}{319.854}{11.955}{\in}%
\TextAt{14}{465.397}{319.854}{C}%
\TextAt{3}{472.72}{317.994}{\Character{79}}%
\TextAt{3}{479.654}{317.994}{\Character{106}}%
\TextAt{27}{482.015}{317.994}{o}%
\TextAt{8}{486.409}{319.854}{.}%
\TextAt{8}{495.496}{319.854}{In}%
\TextAt{8}{510.33}{319.854}{fact,}%
\TextAt{8}{79.226}{302.808}{if}%
\TextAt{13}{90.399}{302.808}{y}%
\NativeMathAt{100.621}{302.808}{11.955}{\in}%
\TextAt{16}{112.66}{302.808}{\Character{79}}%
\TextAt{8}{122.527}{302.808}{,}%
\TextAt{8}{130.237}{302.808}{then}%
\TextAt{8}{157.346}{302.808}{Tr}%
\TextAt{21}{170.354}{300.948}{L}%
\TextAt{28}{176.113}{299.557}{P}%
\TextAt{3}{182.31}{300.948}{\Character{106}}%
\TextAt{21}{184.662}{300.948}{K}%
\TextAt{28}{191.804}{299.952}{p}%
\TextAt{9}{196.384}{302.808}{\Character{40}}%
\TextAt{13}{200.937}{302.808}{\Character{24}}%
\TextAt{13}{206.616}{302.808}{y}%
\TextAt{9}{212.764}{302.808}{\Character{41}}%
\NativeMathAt{221.402}{302.808}{11.955}{\in}%
\TextAt{14}{233.458}{302.808}{o}%
\TextAt{27}{239.302}{301.014}{p}%
\TextAt{8}{243.789}{302.808}{,}%
\TextAt{8}{251.499}{302.808}{since}%
\TextAt{8}{280.617}{302.808}{Tr}%
\TextAt{21}{293.624}{300.948}{L}%
\TextAt{28}{299.383}{299.557}{P}%
\TextAt{3}{305.58}{300.948}{\Character{106}}%
\TextAt{21}{307.933}{300.948}{K}%
\TextAt{28}{315.074}{299.952}{p}%
\TextAt{9}{319.655}{302.808}{\Character{40}}%
\TextAt{13}{324.207}{302.808}{xy}%
\TextAt{9}{336.996}{302.808}{\Character{41}}%
\NativeMathAt{345.634}{302.808}{11.955}{\in}%
\TextAt{14}{357.69}{302.808}{o}%
\TextAt{27}{363.534}{301.014}{p}%
\TextAt{8}{368.021}{302.808}{.}%
\TextAt{8}{377.816}{302.808}{Likewise}%
\TextAt{8}{425.309}{302.808}{Tr}%
\TextAt{21}{438.316}{300.948}{L}%
\TextAt{28}{444.075}{298.657}{P}%
\TextAt{23}{449.774}{300.372}{\Character{48}}%
\TextAt{3}{452.957}{300.948}{\Character{106}}%
\TextAt{21}{455.309}{300.948}{K}%
\TextAt{28}{462.45}{299.952}{p}%
\TextAt{9}{467.031}{302.808}{\Character{40}}%
\TextAt{13}{471.583}{302.808}{\Character{24}}%
\TextAt{13}{477.262}{302.808}{y}%
\TextAt{9}{483.411}{302.808}{\Character{41}}%
\NativeMathAt{492.049}{302.808}{11.955}{\in}%
\TextAt{14}{504.105}{302.808}{o}%
\TextAt{27}{509.949}{301.014}{p}%
\TextAt{8}{518.782}{302.808}{for}%
\TextAt{14}{79.226}{285.762}{P}%
\TextAt{3}{89.119}{290.1}{\Character{48}}%
\TextAt{16}{96.738}{285.762}{\Character{106}}%
\TextAt{14}{104.883}{285.762}{p}%
\TextAt{8}{110.866}{285.762}{,}%
\TextAt{8}{119.117}{285.762}{because}%
\TextAt{8}{163.301}{285.762}{these}%
\TextAt{8}{194.154}{285.762}{elements}%
\TextAt{8}{242.893}{285.762}{are}%
\TextAt{8}{263.286}{285.762}{close}%
\TextAt{8}{292.19}{285.762}{to}%
\TextAt{9}{307.383}{285.762}{0}%
\TextAt{8}{313.236}{285.762}{.}%
\TextAt{8}{324.33}{285.762}{Therefore}%
\TextAt{8}{378.209}{285.762}{Tr}%
\TextAt{21}{391.216}{283.902}{L}%
\TextAt{3}{396.975}{283.902}{\Character{106}}%
\TextAt{21}{399.327}{283.902}{K}%
\TextAt{9}{407.53}{285.762}{\Character{40}}%
\TextAt{13}{412.082}{285.762}{\Character{24}}%
\TextAt{13}{417.761}{285.762}{y}%
\TextAt{9}{423.91}{285.762}{\Character{41}}%
\NativeMathAt{433.286}{285.762}{11.955}{\in}%
\TextAt{14}{446.079}{285.762}{o}%
\TextAt{27}{451.924}{283.968}{p}%
\TextAt{16}{459.668}{285.762}{\Character{92}}%
\TextAt{13}{470.896}{285.762}{K}%
\TextAt{9}{486.538}{285.762}{\Character{61}}%
\TextAt{14}{500.466}{285.762}{o}%
\TextAt{8}{506.31}{285.762}{,}%
\TextAt{8}{514.562}{285.762}{i.e.,}%
\TextAt{13}{79.226}{268.716}{\Character{24}}%
\NativeMathAt{88.702}{268.716}{11.955}{\in}%
\TextAt{14}{100.457}{268.716}{C}%
\TextAt{3}{107.78}{266.856}{\Character{79}}%
\TextAt{3}{114.714}{266.856}{\Character{106}}%
\TextAt{27}{117.075}{266.856}{o}%
\TextAt{8}{121.469}{268.716}{.}%
\TextAt{8}{130.735}{268.716}{This}%
\TextAt{8}{157.727}{268.716}{shows}%
\TextAt{8}{191.617}{268.716}{that}%
\TextAt{14}{217.258}{268.716}{C}%
\TextAt{3}{224.581}{266.856}{\Character{79}}%
\TextAt{3}{231.515}{266.856}{\Character{106}}%
\TextAt{27}{233.877}{266.856}{o}%
\TextAt{8}{242.441}{268.716}{is}%
\TextAt{8}{254.477}{268.716}{dense}%
\TextAt{8}{286.67}{268.716}{in}%
\TextAt{14}{300.595}{268.716}{C}%
\TextAt{3}{307.918}{266.856}{\Character{79}}%
\TextAt{28}{314.625}{265.465}{P}%
\TextAt{3}{320.823}{266.856}{\Character{106}}%
\TextAt{27}{323.175}{266.856}{o}%
\TextAt{28}{327.071}{265.86}{p}%
\TextAt{8}{331.652}{268.716}{,}%
\TextAt{8}{339.142}{268.716}{in}%
\TextAt{8}{353.067}{268.716}{other}%
\TextAt{8}{383.909}{268.716}{words,}%
\TextAt{14}{421.059}{268.716}{C}%
\TextAt{3}{428.382}{266.856}{\Character{79}}%
\TextAt{3}{435.316}{266.856}{\Character{106}}%
\TextAt{27}{437.677}{266.856}{o}%
\TextAt{16}{442.072}{268.716}{\Character{79}}%
\TextAt{27}{451.589}{266.923}{P}%
\TextAt{9}{462.468}{268.716}{\Character{61}}%
\TextAt{14}{475.357}{268.716}{C}%
\TextAt{3}{482.68}{266.856}{\Character{79}}%
\TextAt{28}{489.388}{265.465}{P}%
\TextAt{3}{495.585}{266.856}{\Character{106}}%
\TextAt{27}{497.938}{266.856}{o}%
\TextAt{28}{501.834}{265.86}{p}%
\TextAt{8}{506.415}{268.716}{,}%
\TextAt{8}{513.904}{268.716}{and}%
\TextAt{8}{79.226}{251.67}{so}%
\TextAt{14}{93.596}{251.67}{D}%
\TextAt{3}{103.538}{249.81}{\Character{79}}%
\TextAt{3}{110.472}{249.81}{\Character{106}}%
\TextAt{27}{112.833}{249.81}{o}%
\TextAt{16}{117.227}{251.67}{\Character{79}}%
\TextAt{27}{126.745}{249.877}{P}%
\TextAt{9}{137.16}{251.67}{\Character{61}}%
\TextAt{14}{149.585}{251.67}{D}%
\TextAt{3}{159.527}{249.81}{\Character{79}}%
\TextAt{28}{166.234}{248.419}{P}%
\TextAt{3}{172.432}{249.81}{\Character{106}}%
\TextAt{27}{174.784}{249.81}{o}%
\TextAt{28}{178.68}{248.814}{p}%
\TextAt{8}{183.261}{251.67}{.}%
\Vector{q 0 G 0 g 0.399 w 452.334 -522.33 m 452.334 -514.26 l S Q}%
\Vector{q 0 G 0 g 0.399 w 452.534 -514.46 m 459.707 -514.46 l S Q}%
\Vector{q 0 G 0 g 0.399 w 452.534 -522.131 m 459.707 -522.131 l S Q}%
\Vector{q 0 G 0 g 0.399 w 459.906 -522.33 m 459.906 -514.26 l S Q}%
\TextAt{8}{97.159}{223.571}{If}%
\TextAt{8}{108.601}{223.571}{we}%
\TextAt{8}{125.578}{223.571}{put}%
\TextAt{14}{146.787}{223.571}{D}%
\TextAt{9}{160.05}{223.571}{\Character{61}}%
\TextAt{14}{172.475}{223.571}{D}%
\TextAt{21}{182.417}{221.712}{L}%
\TextAt{3}{188.176}{221.712}{\Character{106}}%
\TextAt{21}{190.528}{221.712}{K}%
\TextAt{8}{202.378}{223.571}{and}%
\TextAt{14}{224.891}{223.571}{D}%
\TextAt{27}{234.832}{221.778}{P}%
\TextAt{9}{245.247}{223.571}{\Character{61}}%
\TextAt{14}{257.672}{223.571}{D}%
\TextAt{21}{267.614}{221.712}{L}%
\TextAt{28}{273.373}{220.321}{P}%
\TextAt{3}{279.57}{221.712}{\Character{106}}%
\TextAt{21}{281.922}{221.712}{K}%
\TextAt{28}{289.064}{220.715}{p}%
\TextAt{8}{293.645}{223.571}{,}%
\TextAt{8}{300.594}{223.571}{and}%
\TextAt{8}{323.106}{223.571}{consider}%
\TextAt{14}{368.441}{223.571}{D}%
\TextAt{27}{378.382}{221.778}{P}%
\TextAt{8}{389.123}{223.571}{at}%
\TextAt{8}{403.183}{223.571}{the}%
\TextAt{8}{423.09}{223.571}{same}%
\TextAt{8}{452.165}{223.571}{time}%
\TextAt{8}{478.575}{223.571}{as}%
\TextAt{8}{492.695}{223.571}{an}%
\TextAt{8}{508.704}{223.571}{ideal}%
\TextAt{8}{79.226}{206.525}{of}%
\TextAt{16}{92.556}{206.525}{\Character{79}}%
\TextAt{8}{106.303}{206.525}{(i.e.,}%
\TextAt{8}{132.963}{206.525}{as}%
\TextAt{8}{147.333}{206.525}{the}%
\TextAt{8}{167.49}{206.525}{ideal}%
\TextAt{16}{195.453}{206.525}{\Character{79}}%
\TextAt{16}{207.958}{206.525}{\Character{92}}%
\TextAt{14}{218.586}{206.525}{D}%
\TextAt{27}{228.528}{204.732}{P}%
\TextAt{8}{235.622}{206.525}{),}%
\TextAt{8}{247.326}{206.525}{then}%
\TextAt{8}{273.986}{206.525}{(2.2),}%
\TextAt{8}{305.213}{206.525}{(iii)}%
\TextAt{8}{327.975}{206.525}{gives}%
\TextAt{8}{356.656}{206.525}{us}%
\TextAt{8}{371.671}{206.525}{the}%
\TextAt{12}{79.226}{181.416}{Corollary}%
\TextAt{12}{139.014}{181.416}{(2.3).}%
\TextAt{14}{174.33}{181.416}{D}%
\TextAt{9}{187.592}{181.416}{\Character{61}}%
\TextAt{20}{200.017}{190.382}{\Character{81}}%
\TextAt{27}{211.309}{177.929}{P}%
\TextAt{14}{220.395}{181.416}{D}%
\TextAt{27}{230.336}{179.622}{P}%
\TextAt{8}{237.43}{181.416}{.}%
\TextAt{8}{97.159}{156.306}{The}%
\TextAt{8}{122.691}{156.306}{name}%
\TextAt{8}{155.384}{156.306}{“different”}%
\TextAt{8}{214.08}{156.306}{is}%
\TextAt{8}{227.322}{156.306}{explained}%
\TextAt{8}{281.152}{156.306}{by}%
\TextAt{8}{298.889}{156.306}{the}%
\TextAt{8}{320.523}{156.306}{following}%
\TextAt{8}{371.436}{156.306}{explicit}%
\TextAt{8}{414.208}{156.306}{description,}%
\TextAt{8}{479.203}{156.306}{which}%
\TextAt{8}{514.167}{156.306}{was}%
\TextAt{8}{79.226}{139.26}{Dedekind’s}%
\TextAt{8}{140.219}{139.26}{original}%
\TextAt{8}{183.446}{139.26}{way}%
\TextAt{8}{208.131}{139.26}{to}%
\TextAt{8}{223.383}{139.26}{define}%
\TextAt{8}{258.134}{139.26}{it.}%
\TextAt{8}{277.214}{139.26}{Let}%
\TextAt{13}{299.125}{139.26}{\Character{11}}%
\NativeMathAt{311.572}{139.26}{11.955}{\in}%
\TextAt{16}{324.46}{139.26}{\Character{79}}%
\TextAt{8}{339.157}{139.26}{and}%
\TextAt{8}{362.862}{139.26}{let}%
\TextAt{13}{380.708}{139.26}{f}%
\TextAt{9}{387.754}{139.26}{\Character{40}}%
\TextAt{13}{392.307}{139.26}{X}%
\TextAt{9}{402.962}{139.26}{\Character{41}}%
\NativeMathAt{412.44}{139.26}{11.955}{\in}%
\TextAt{14}{425.336}{139.26}{o}%
\TextAt{9}{431.18}{139.26}{\Character{91}}%
\TextAt{13}{434.432}{139.26}{X}%
\TextAt{9}{445.087}{139.26}{\Character{93}}%
\TextAt{8}{453.177}{139.26}{be}%
\TextAt{8}{470.043}{139.26}{the}%
\TextAt{8}{491.142}{139.26}{minimal}%
\TextAt{8}{79.226}{122.214}{polynomial}%
\TextAt{8}{139.719}{122.214}{of}%
\TextAt{13}{153.049}{122.214}{\Character{11}}%
\TextAt{8}{160.571}{122.214}{.}%
\TextAt{8}{169.019}{122.214}{We}%
\TextAt{8}{189.164}{122.214}{define}%
\TextAt{8}{222.973}{122.214}{the}%
\TextAt{12}{243.129}{122.214}{different}%
\TextAt{12}{296.772}{122.214}{of}%
\TextAt{12}{312.099}{122.214}{the}%
\TextAt{12}{335.423}{122.214}{element}%
\TextAt{13}{385.599}{122.214}{\Character{11}}%
\TextAt{8}{397.018}{122.214}{by}%
\TextAt{13}{221.374}{85.475}{\Character{14}}%
\TextAt{21}{226.559}{83.615}{L}%
\TextAt{3}{232.318}{83.615}{\Character{106}}%
\TextAt{21}{234.67}{83.615}{K}%
\TextAt{9}{242.873}{85.475}{\Character{40}}%
\TextAt{13}{247.425}{85.475}{\Character{11}}%
\TextAt{9}{254.947}{85.475}{\Character{41}}%
\TextAt{9}{262.802}{85.475}{\Character{61}}%
\TextAt{20}{275.246}{109.983}{\Character{56}}%
\TextAt{20}{275.246}{99.223}{\Character{60}}%
\TextAt{20}{275.246}{77.704}{\Character{58}}%
\TextAt{13}{285.872}{94.6}{f}%
\TextAt{3}{292.919}{98.938}{\Character{48}}%
\TextAt{9}{295.714}{94.6}{\Character{40}}%
\TextAt{13}{300.266}{94.6}{\Character{11}}%
\TextAt{9}{307.788}{94.6}{\Character{41}}%
\TextAt{8}{324.296}{94.6}{if}%
\TextAt{13}{335.019}{94.6}{L}%
\TextAt{9}{346.305}{94.6}{\Character{61}}%
\TextAt{13}{358.73}{94.6}{K}%
\TextAt{9}{369.548}{94.6}{\Character{40}}%
\TextAt{13}{374.1}{94.6}{\Character{11}}%
\TextAt{9}{381.622}{94.6}{\Character{41}}%
\TextAt{13}{386.174}{94.6}{\Character{59}}%
\TextAt{9}{285.872}{74.145}{0}%
\TextAt{8}{324.296}{74.145}{if}%
\TextAt{13}{335.019}{74.145}{L}%
\NativeMathAt{346.305}{74.145}{11.955}{\neq}%
\TextAt{13}{358.73}{74.145}{K}%
\TextAt{9}{369.548}{74.145}{\Character{40}}%
\TextAt{13}{374.1}{74.145}{\Character{11}}%
\TextAt{9}{381.622}{74.145}{\Character{41}}%
\TextAt{13}{386.174}{74.145}{\Character{58}}%
\TextAt{8}{297.211}{43.033}{179}%
\Destination{78.23}{809.815}{nav.page.179}%
\Destination{79.36}{724.455}{nav.print.196}%
\Destination{79.36}{195.455}{nav.number.3.2.3}%
\Hotlink{134.578}{484.455}{9.751}{11.047}{nav.chapter.2}%
\Hotlink{150.538}{484.455}{25.078}{11.047}{nav.number.2.8.4}%
\Hotlink{273.486}{204.355}{25.078}{11.047}{nav.number.3.2.2}%
}

% Page 192 (body page 180)
\Page{%
\TextAt{2}{79.226}{796.943}{III.}%
\TextAt{2}{98.317}{796.943}{Riemann–Roch}%
\TextAt{2}{174.375}{796.943}{Theory}%
\Vector{q 0 G 0 g 0.299 w 7.226 18.729 m 460.769 18.729 l S Q}%
\TextAt{8}{79.226}{761.124}{In}%
\TextAt{8}{93.847}{761.124}{the}%
\TextAt{8}{114.004}{761.124}{special}%
\TextAt{8}{152.105}{761.124}{case}%
\TextAt{8}{176.876}{761.124}{where}%
\TextAt{16}{210.685}{761.124}{\Character{79}}%
\TextAt{9}{223.856}{761.124}{\Character{61}}%
\TextAt{14}{236.281}{761.124}{o}%
\TextAt{9}{242.126}{761.124}{\Character{91}}%
\TextAt{13}{245.377}{761.124}{\Character{11}}%
\TextAt{9}{252.899}{761.124}{\Character{93}}%
\TextAt{8}{260.048}{761.124}{we}%
\TextAt{8}{277.275}{761.124}{then}%
\TextAt{8}{303.935}{761.124}{obtain:}%
\TextAt{12}{79.226}{736.247}{Proposition}%
\TextAt{12}{151.937}{736.247}{(2.4).}%
\TextAt{8}{187.253}{736.247}{If}%
\TextAt{16}{198.945}{736.247}{\Character{79}}%
\TextAt{9}{212.116}{736.247}{\Character{61}}%
\TextAt{14}{224.541}{736.247}{o}%
\TextAt{9}{230.386}{736.247}{\Character{91}}%
\TextAt{13}{233.637}{736.247}{\Character{11}}%
\TextAt{9}{241.159}{736.247}{\Character{93}}%
\TextAt{8}{244.411}{736.247}{,}%
\TextAt{8}{251.56}{736.247}{then}%
\TextAt{8}{278.22}{736.247}{the}%
\TextAt{8}{298.376}{736.247}{different}%
\TextAt{8}{344.547}{736.247}{is}%
\TextAt{8}{356.311}{736.247}{the}%
\TextAt{8}{376.467}{736.247}{principal}%
\TextAt{8}{425.245}{736.247}{ideal}%
\TextAt{14}{259.756}{711.219}{D}%
\TextAt{21}{269.698}{709.359}{L}%
\TextAt{3}{275.457}{709.359}{\Character{106}}%
\TextAt{21}{277.809}{709.359}{K}%
\TextAt{9}{289.333}{711.219}{\Character{61}}%
\TextAt{9}{301.743}{711.219}{\Character{40}}%
\TextAt{13}{306.31}{711.219}{\Character{14}}%
\TextAt{21}{311.495}{709.359}{L}%
\TextAt{3}{317.254}{709.359}{\Character{106}}%
\TextAt{21}{319.606}{709.359}{K}%
\TextAt{9}{327.809}{711.219}{\Character{40}}%
\TextAt{13}{332.361}{711.219}{\Character{11}}%
\TextAt{9}{339.883}{711.219}{\Character{41}\Character{41}}%
\TextAt{13}{348.988}{711.219}{\Character{58}}%
\TextAt{19}{79.226}{670.529}{Proof.}%
\TextAt{8}{116}{670.529}{Let}%
\TextAt{13}{136.97}{670.529}{f}%
\TextAt{9}{144.016}{670.529}{\Character{40}}%
\TextAt{13}{148.568}{670.529}{X}%
\TextAt{9}{159.223}{670.529}{\Character{41}}%
\TextAt{9}{167.078}{670.529}{\Character{61}}%
\TextAt{13}{179.522}{670.529}{a}%
\TextAt{17}{185.667}{668.736}{0}%
\TextAt{9}{193.056}{670.529}{\Character{43}}%
\TextAt{13}{204.817}{670.529}{a}%
\TextAt{17}{210.962}{668.736}{1}%
\TextAt{13}{215.695}{670.529}{X}%
\TextAt{9}{229.006}{670.529}{\Character{43}}%
\TextAt{16}{240.768}{670.529}{\Character{1}}%
\TextAt{16}{246.064}{670.529}{\Character{1}}%
\TextAt{16}{251.372}{670.529}{\Character{1}}%
\TextAt{9}{257.372}{670.529}{\Character{43}}%
\TextAt{13}{269.133}{670.529}{a}%
\TextAt{21}{275.278}{668.736}{n}%
\TextAt{13}{280.915}{670.529}{X}%
\TextAt{21}{291.57}{674.868}{n}%
\TextAt{8}{301.103}{670.529}{be}%
\TextAt{8}{317.028}{670.529}{the}%
\TextAt{8}{337.184}{670.529}{minimal}%
\TextAt{8}{382.709}{670.529}{polynomial}%
\TextAt{8}{443.203}{670.529}{of}%
\TextAt{13}{456.533}{670.529}{\Character{11}}%
\TextAt{8}{467.952}{670.529}{and}%
\TextAt{13}{215.185}{649.596}{f}%
\TextAt{9}{222.231}{649.596}{\Character{40}}%
\TextAt{13}{226.784}{649.596}{X}%
\TextAt{9}{237.439}{649.596}{\Character{41}}%
\Vector{q 0 G 0 g 0.478 w 140.194 -129.503 m 172.982 -129.503 l S Q}%
\TextAt{13}{212.194}{633.308}{X}%
\TextAt{16}{225.505}{633.308}{\Character{0}}%
\TextAt{13}{237.461}{633.308}{\Character{11}}%
\TextAt{9}{249.499}{641.508}{\Character{61}}%
\TextAt{13}{261.924}{641.508}{b}%
\TextAt{17}{266.901}{639.715}{0}%
\TextAt{9}{274.29}{641.508}{\Character{43}}%
\TextAt{13}{286.051}{641.508}{b}%
\TextAt{17}{291.029}{639.715}{1}%
\TextAt{13}{295.761}{641.508}{X}%
\TextAt{9}{309.073}{641.508}{\Character{43}}%
\TextAt{16}{320.834}{641.508}{\Character{1}}%
\TextAt{16}{326.13}{641.508}{\Character{1}}%
\TextAt{16}{331.438}{641.508}{\Character{1}}%
\TextAt{9}{337.438}{641.508}{\Character{43}}%
\TextAt{13}{349.2}{641.508}{b}%
\TextAt{21}{354.177}{639.715}{n}%
\TextAt{3}{359.315}{639.715}{\Character{0}}%
\TextAt{17}{365.901}{639.715}{1}%
\TextAt{13}{370.634}{641.508}{X}%
\TextAt{21}{381.289}{646.445}{n}%
\TextAt{3}{386.427}{646.445}{\Character{0}}%
\TextAt{17}{393.014}{646.445}{1}%
\TextAt{13}{397.746}{641.508}{\Character{58}}%
\TextAt{8}{79.226}{613.859}{The}%
\TextAt{8}{103.28}{613.859}{dual}%
\TextAt{8}{129.294}{613.859}{basis}%
\TextAt{8}{158.035}{613.859}{of}%
\TextAt{9}{171.365}{613.859}{1}%
\TextAt{13}{177.218}{613.859}{\Character{59}}%
\TextAt{13}{182.442}{613.859}{\Character{11}\Character{59}}%
\TextAt{13}{195.187}{613.859}{\Character{58}}%
\TextAt{13}{200.423}{613.859}{\Character{58}}%
\TextAt{13}{205.648}{613.859}{\Character{58}}%
\TextAt{13}{210.884}{613.859}{\Character{59}}%
\TextAt{13}{216.12}{613.859}{\Character{11}}%
\TextAt{21}{223.726}{618.197}{n}%
\TextAt{3}{228.864}{618.197}{\Character{0}}%
\TextAt{17}{235.451}{618.197}{1}%
\TextAt{8}{244.08}{613.859}{with}%
\TextAt{8}{270.74}{613.859}{respect}%
\TextAt{8}{310.79}{613.859}{to}%
\TextAt{8}{325.101}{613.859}{Tr}%
\TextAt{9}{338.108}{613.859}{\Character{40}}%
\TextAt{13}{342.66}{613.859}{xy}%
\TextAt{9}{355.449}{613.859}{\Character{41}}%
\TextAt{8}{363.899}{613.859}{is}%
\TextAt{8}{375.662}{613.859}{then}%
\TextAt{8}{402.323}{613.859}{given}%
\TextAt{8}{432.892}{613.859}{by}%
\TextAt{13}{271.977}{593.59}{b}%
\TextAt{17}{276.955}{591.797}{0}%
\Vector{q 0 G 0 g 0.478 w 191.598 -185.509 m 218.066 -185.509 l S Q}%
\TextAt{13}{263.598}{577.302}{f}%
\TextAt{3}{270.645}{580.755}{\Character{48}}%
\TextAt{9}{273.44}{577.302}{\Character{40}}%
\TextAt{13}{277.992}{577.302}{\Character{11}}%
\TextAt{9}{285.514}{577.302}{\Character{41}}%
\TextAt{13}{291.262}{585.502}{\Character{59}}%
\TextAt{13}{296.486}{585.502}{\Character{58}}%
\TextAt{13}{301.723}{585.502}{\Character{58}}%
\TextAt{13}{306.947}{585.502}{\Character{58}}%
\TextAt{13}{312.184}{585.502}{\Character{59}}%
\TextAt{13}{321.195}{593.59}{b}%
\TextAt{21}{326.172}{591.797}{n}%
\TextAt{3}{331.31}{591.797}{\Character{0}}%
\TextAt{17}{337.897}{591.797}{1}%
\Vector{q 0 G 0 g 0.478 w 246.678 -185.509 m 273.146 -185.509 l S Q}%
\TextAt{13}{318.678}{577.302}{f}%
\TextAt{3}{325.724}{580.755}{\Character{48}}%
\TextAt{9}{328.519}{577.302}{\Character{40}}%
\TextAt{13}{333.072}{577.302}{\Character{11}}%
\TextAt{9}{340.593}{577.302}{\Character{41}}%
\TextAt{13}{346.341}{585.502}{\Character{58}}%
\TextAt{8}{79.226}{556.918}{For}%
\TextAt{8}{100.196}{556.918}{if}%
\TextAt{13}{110.919}{556.918}{\Character{11}}%
\TextAt{17}{118.365}{555.125}{1}%
\TextAt{13}{123.097}{556.918}{\Character{59}}%
\TextAt{13}{128.321}{556.918}{\Character{58}}%
\TextAt{13}{133.558}{556.918}{\Character{58}}%
\TextAt{13}{138.794}{556.918}{\Character{58}}%
\TextAt{13}{144.019}{556.918}{\Character{59}}%
\TextAt{13}{149.255}{556.918}{\Character{11}}%
\TextAt{21}{156.764}{555.125}{n}%
\TextAt{8}{166.297}{556.918}{are}%
\TextAt{8}{185.808}{556.918}{the}%
\TextAt{8}{205.965}{556.918}{roots}%
\TextAt{8}{235.625}{556.918}{of}%
\TextAt{13}{248.955}{556.918}{f}%
\TextAt{8}{256.002}{556.918}{,}%
\TextAt{8}{263.151}{556.918}{then}%
\TextAt{8}{289.811}{556.918}{one}%
\TextAt{8}{311.271}{556.918}{has}%
\TextAt{21}{200.034}{540.989}{n}%
\TextAt{20}{193.969}{537.402}{\Character{88}}%
\TextAt{21}{195.751}{512.208}{i}%
\TextAt{17}{198.634}{512.208}{\Character{61}1}%
\TextAt{13}{219.069}{534.133}{f}%
\TextAt{9}{226.115}{534.133}{\Character{40}}%
\TextAt{13}{230.668}{534.133}{X}%
\TextAt{9}{241.323}{534.133}{\Character{41}}%
\Vector{q 0 G 0 g 0.478 w 142.425 -244.966 m 178.519 -244.966 l S Q}%
\TextAt{13}{214.425}{517.844}{X}%
\TextAt{16}{227.737}{517.844}{\Character{0}}%
\TextAt{13}{239.692}{517.844}{\Character{11}}%
\TextAt{21}{247.138}{516.051}{i}%
\TextAt{13}{261.758}{534.133}{\Character{11}}%
\TextAt{21}{269.28}{538.471}{r}%
\TextAt{21}{269.204}{531.177}{i}%
\Vector{q 0 G 0 g 0.478 w 180.91 -244.966 m 210.683 -244.966 l S Q}%
\TextAt{13}{252.91}{517.844}{f}%
\TextAt{3}{259.956}{521.298}{\Character{48}}%
\TextAt{9}{262.752}{517.844}{\Character{40}}%
\TextAt{13}{267.304}{517.844}{\Character{11}}%
\TextAt{21}{274.749}{516.051}{i}%
\TextAt{9}{278.131}{517.844}{\Character{41}}%
\TextAt{9}{287.199}{526.045}{\Character{61}}%
\TextAt{13}{299.625}{526.045}{X}%
\TextAt{21}{310.28}{530.981}{r}%
\TextAt{13}{314.835}{526.045}{\Character{59}}%
\TextAt{9}{343.989}{526.045}{0}%
\NativeMathAt{353.163}{526.045}{11.955}{\leq}%
\TextAt{13}{365.782}{526.045}{r}%
\NativeMathAt{374.704}{526.045}{11.955}{\leq}%
\TextAt{13}{387.323}{526.045}{n}%
\TextAt{16}{396.967}{526.045}{\Character{0}}%
\TextAt{9}{408.922}{526.045}{1}%
\TextAt{13}{414.775}{526.045}{\Character{59}}%
\TextAt{8}{79.226}{493.618}{as}%
\TextAt{8}{93.398}{493.618}{the}%
\TextAt{8}{113.356}{493.618}{difference}%
\TextAt{8}{165.497}{493.618}{of}%
\TextAt{8}{178.628}{493.618}{the}%
\TextAt{8}{198.586}{493.618}{two}%
\TextAt{8}{220.505}{493.618}{sides}%
\TextAt{8}{248.389}{493.618}{is}%
\TextAt{8}{259.954}{493.618}{a}%
\TextAt{8}{269.511}{493.618}{polynomial}%
\TextAt{8}{329.806}{493.618}{of}%
\TextAt{8}{342.938}{493.618}{degree}%
\NativeMathAt{379.154}{493.618}{11.955}{\leq}%
\TextAt{13}{391.774}{493.618}{n}%
\TextAt{16}{401.012}{493.618}{\Character{0}}%
\TextAt{9}{412.561}{493.618}{1}%
\TextAt{8}{422.113}{493.618}{with}%
\TextAt{8}{448.575}{493.618}{roots}%
\TextAt{13}{478.037}{493.618}{\Character{11}}%
\TextAt{17}{485.483}{491.825}{1}%
\TextAt{13}{490.215}{493.618}{\Character{59}}%
\TextAt{13}{495.439}{493.618}{\Character{58}}%
\TextAt{13}{500.676}{493.618}{\Character{58}}%
\TextAt{13}{505.9}{493.618}{\Character{58}}%
\TextAt{13}{511.137}{493.618}{\Character{59}}%
\TextAt{13}{516.373}{493.618}{\Character{11}}%
\TextAt{21}{523.881}{491.825}{n}%
\TextAt{8}{529.518}{493.618}{,}%
\TextAt{8}{79.226}{476.572}{so}%
\TextAt{8}{93.596}{476.572}{is}%
\TextAt{8}{105.36}{476.572}{identically}%
\TextAt{8}{161.944}{476.572}{zero.}%
\TextAt{8}{191.206}{476.572}{We}%
\TextAt{8}{211.351}{476.572}{may}%
\TextAt{8}{236.72}{476.572}{write}%
\TextAt{8}{266.631}{476.572}{this}%
\TextAt{8}{289.454}{476.572}{equation}%
\TextAt{8}{337.263}{476.572}{in}%
\TextAt{8}{350.915}{476.572}{the}%
\TextAt{8}{371.072}{476.572}{form}%
\TextAt{8}{243.065}{447.82}{Tr}%
\TextAt{20}{258.065}{464.677}{\Character{20}}%
\TextAt{13}{268.561}{455.908}{f}%
\TextAt{9}{275.608}{455.908}{\Character{40}}%
\TextAt{13}{280.16}{455.908}{X}%
\TextAt{9}{290.815}{455.908}{\Character{41}}%
\Vector{q 0 G 0 g 0.478 w 193.57 -323.191 m 226.359 -323.191 l S Q}%
\TextAt{13}{265.57}{439.62}{X}%
\TextAt{16}{278.882}{439.62}{\Character{0}}%
\TextAt{13}{290.837}{439.62}{\Character{11}}%
\TextAt{13}{307.945}{455.908}{\Character{11}}%
\TextAt{21}{315.467}{460.246}{r}%
\Vector{q 0 G 0 g 0.478 w 228.75 -323.191 m 255.218 -323.191 l S Q}%
\TextAt{13}{300.75}{439.62}{f}%
\TextAt{3}{307.796}{443.073}{\Character{48}}%
\TextAt{9}{310.591}{439.62}{\Character{40}}%
\TextAt{13}{315.144}{439.62}{\Character{11}}%
\TextAt{9}{322.665}{439.62}{\Character{41}}%
\TextAt{20}{328.413}{464.677}{\Character{21}}%
\TextAt{9}{338.044}{447.82}{\Character{61}}%
\TextAt{13}{350.469}{447.82}{X}%
\TextAt{21}{361.124}{452.756}{r}%
\TextAt{13}{365.679}{447.82}{\Character{58}}%
\TextAt{8}{79.226}{419.068}{Considering}%
\TextAt{8}{143.676}{419.068}{now}%
\TextAt{8}{168.065}{419.068}{the}%
\TextAt{8}{188.221}{419.068}{coefficient}%
\TextAt{8}{242.845}{419.068}{of}%
\TextAt{8}{256.175}{419.068}{each}%
\TextAt{8}{282.512}{419.068}{of}%
\TextAt{8}{295.842}{419.068}{the}%
\TextAt{8}{315.998}{419.068}{powers}%
\TextAt{8}{354.757}{419.068}{of}%
\TextAt{13}{368.087}{419.068}{X}%
\TextAt{8}{378.742}{419.068}{,}%
\TextAt{8}{385.891}{419.068}{we}%
\TextAt{8}{403.119}{419.068}{obtain}%
\TextAt{8}{255.719}{390.304}{Tr}%
\TextAt{20}{270.718}{407.161}{\Character{18}}%
\TextAt{13}{279.519}{390.305}{\Character{11}}%
\TextAt{21}{287.04}{395.241}{i}%
\TextAt{13}{300.171}{398.392}{b}%
\TextAt{21}{305.149}{396.599}{j}%
\Vector{q 0 G 0 g 0.478 w 219.617 -380.707 m 246.085 -380.707 l S Q}%
\TextAt{13}{291.617}{382.104}{f}%
\TextAt{3}{298.664}{385.558}{\Character{48}}%
\TextAt{9}{301.459}{382.104}{\Character{40}}%
\TextAt{13}{306.011}{382.104}{\Character{11}}%
\TextAt{9}{313.533}{382.104}{\Character{41}}%
\TextAt{20}{319.281}{407.161}{\Character{19}}%
\TextAt{9}{331.402}{390.304}{\Character{61}}%
\TextAt{13}{343.827}{390.304}{\Character{14}}%
\TextAt{21}{349.012}{388.511}{ij}%
\TextAt{8}{79.226}{361.553}{and}%
\TextAt{8}{101.989}{361.553}{the}%
\TextAt{8}{122.145}{361.553}{claim}%
\TextAt{8}{153.36}{361.553}{follows.}%
\TextAt{8}{97.159}{341.518}{As}%
\TextAt{16}{114.446}{341.518}{\Character{79}}%
\TextAt{9}{127.617}{341.518}{\Character{61}}%
\TextAt{14}{140.042}{341.518}{o}%
\TextAt{9}{148.543}{341.518}{\Character{43}}%
\TextAt{14}{160.304}{341.518}{o}%
\TextAt{13}{166.149}{341.518}{\Character{11}}%
\TextAt{9}{176.327}{341.518}{\Character{43}}%
\TextAt{16}{188.088}{341.518}{\Character{1}}%
\TextAt{16}{193.384}{341.518}{\Character{1}}%
\TextAt{16}{198.692}{341.518}{\Character{1}}%
\TextAt{9}{204.693}{341.518}{\Character{43}}%
\TextAt{14}{216.454}{341.518}{o}%
\TextAt{13}{222.298}{341.518}{\Character{11}}%
\TextAt{21}{229.82}{345.856}{n}%
\TextAt{3}{234.958}{345.856}{\Character{0}}%
\TextAt{17}{241.545}{345.856}{1}%
\TextAt{8}{246.277}{341.518}{,}%
\TextAt{8}{253.426}{341.518}{we}%
\TextAt{8}{270.654}{341.518}{get}%
\TextAt{14}{219.739}{316.49}{C}%
\TextAt{3}{227.062}{314.631}{\Character{79}}%
\TextAt{3}{233.996}{314.631}{\Character{106}}%
\TextAt{27}{236.357}{314.631}{o}%
\TextAt{9}{244.072}{316.49}{\Character{61}}%
\TextAt{13}{256.498}{316.49}{f}%
\TextAt{3}{263.544}{321.426}{\Character{48}}%
\TextAt{9}{266.339}{316.49}{\Character{40}}%
\TextAt{13}{270.892}{316.49}{\Character{11}}%
\TextAt{9}{278.414}{316.49}{\Character{41}}%
\TextAt{3}{282.966}{321.426}{\Character{0}}%
\TextAt{17}{289.552}{321.426}{1}%
\TextAt{9}{294.285}{316.49}{\Character{40}}%
\TextAt{14}{298.837}{316.49}{o}%
\TextAt{13}{304.681}{316.49}{b}%
\TextAt{17}{309.658}{314.697}{0}%
\TextAt{9}{317.047}{316.49}{\Character{43}}%
\TextAt{16}{328.809}{316.49}{\Character{1}}%
\TextAt{16}{334.105}{316.49}{\Character{1}}%
\TextAt{16}{339.413}{316.49}{\Character{1}}%
\TextAt{9}{345.413}{316.49}{\Character{43}}%
\TextAt{14}{357.174}{316.49}{o}%
\TextAt{13}{363.019}{316.49}{b}%
\TextAt{21}{367.996}{314.697}{n}%
\TextAt{3}{373.134}{314.697}{\Character{0}}%
\TextAt{17}{379.72}{314.697}{1}%
\TextAt{9}{384.453}{316.49}{\Character{41}}%
\TextAt{13}{389.005}{316.49}{\Character{58}}%
\TextAt{8}{79.226}{291.463}{From}%
\TextAt{8}{109.951}{291.463}{the}%
\TextAt{8}{130.107}{291.463}{recursive}%
\TextAt{8}{178.944}{291.463}{formulas}%
\TextAt{13}{181.831}{266.435}{b}%
\TextAt{21}{186.808}{264.642}{n}%
\TextAt{3}{191.946}{264.642}{\Character{0}}%
\TextAt{17}{198.533}{264.642}{1}%
\TextAt{9}{206.586}{266.435}{\Character{61}}%
\TextAt{9}{218.996}{266.435}{1}%
\TextAt{13}{224.864}{266.435}{\Character{59}}%
\TextAt{13}{253.999}{266.435}{b}%
\TextAt{21}{258.996}{264.642}{n}%
\TextAt{3}{264.134}{264.642}{\Character{0}}%
\TextAt{17}{270.72}{264.642}{2}%
\TextAt{16}{278.109}{266.435}{\Character{0}}%
\TextAt{13}{290.065}{266.435}{\Character{11}b}%
\TextAt{21}{302.563}{264.642}{n}%
\TextAt{3}{307.702}{264.642}{\Character{0}}%
\TextAt{17}{314.288}{264.642}{1}%
\TextAt{9}{322.341}{266.435}{\Character{61}}%
\TextAt{13}{334.767}{266.435}{a}%
\TextAt{21}{340.912}{264.642}{n}%
\TextAt{3}{346.05}{264.642}{\Character{0}}%
\TextAt{17}{352.636}{264.642}{1}%
\TextAt{13}{357.369}{266.435}{\Character{59}}%
\TextAt{16}{386.523}{266.435}{\Character{1}}%
\TextAt{16}{391.819}{266.435}{\Character{1}}%
\TextAt{16}{397.127}{266.435}{\Character{1}}%
\TextAt{9}{403.792}{266.435}{\Character{61}}%
\TextAt{16}{416.217}{266.435}{\Character{1}}%
\TextAt{16}{421.513}{266.435}{\Character{1}}%
\TextAt{16}{426.821}{266.435}{\Character{1}}%
\TextAt{8}{79.226}{241.407}{it}%
\TextAt{8}{90.93}{241.407}{follows}%
\TextAt{8}{129.366}{241.407}{that}%
\TextAt{13}{208.796}{224.361}{b}%
\TextAt{21}{213.773}{222.568}{n}%
\TextAt{3}{218.911}{222.568}{\Character{0}}%
\TextAt{21}{225.498}{222.568}{i}%
\TextAt{9}{232.2}{224.361}{\Character{61}}%
\TextAt{13}{244.625}{224.361}{\Character{11}}%
\TextAt{21}{252.147}{229.298}{i}%
\TextAt{3}{255.03}{229.298}{\Character{0}}%
\TextAt{17}{261.617}{229.298}{1}%
\TextAt{9}{269.006}{224.361}{\Character{43}}%
\TextAt{13}{280.767}{224.361}{a}%
\TextAt{21}{286.912}{222.568}{n}%
\TextAt{3}{292.05}{222.568}{\Character{0}}%
\TextAt{17}{298.637}{222.568}{1}%
\TextAt{13}{303.369}{224.361}{\Character{11}}%
\TextAt{21}{310.891}{229.298}{i}%
\TextAt{3}{313.774}{229.298}{\Character{0}}%
\TextAt{17}{320.36}{229.298}{2}%
\TextAt{9}{327.749}{224.361}{\Character{43}}%
\TextAt{16}{339.511}{224.361}{\Character{1}}%
\TextAt{16}{344.807}{224.361}{\Character{1}}%
\TextAt{16}{350.115}{224.361}{\Character{1}}%
\TextAt{9}{356.115}{224.361}{\Character{43}}%
\TextAt{13}{367.876}{224.361}{a}%
\TextAt{21}{374.021}{222.568}{n}%
\TextAt{3}{379.16}{222.568}{\Character{0}}%
\TextAt{21}{385.746}{222.568}{i}%
\TextAt{17}{388.629}{222.568}{\Character{43}1}%
\TextAt{13}{399.948}{224.361}{\Character{59}}%
\TextAt{8}{79.226}{202.543}{so}%
\TextAt{8}{92.984}{202.543}{that}%
\TextAt{14}{117.74}{202.543}{o}%
\TextAt{13}{123.584}{202.543}{b}%
\TextAt{17}{128.561}{200.749}{0}%
\TextAt{9}{134.698}{202.543}{\Character{43}}%
\TextAt{16}{145.207}{202.543}{\Character{1}}%
\TextAt{16}{150.503}{202.543}{\Character{1}}%
\TextAt{16}{155.811}{202.543}{\Character{1}}%
\TextAt{9}{160.558}{202.543}{\Character{43}}%
\TextAt{14}{171.067}{202.543}{o}%
\TextAt{13}{176.912}{202.543}{b}%
\TextAt{21}{181.889}{200.749}{n}%
\TextAt{3}{187.027}{200.749}{\Character{0}}%
\TextAt{17}{193.613}{200.749}{1}%
\TextAt{9}{201.666}{202.543}{\Character{61}}%
\TextAt{14}{214.092}{202.543}{o}%
\TextAt{9}{219.936}{202.543}{\Character{91}}%
\TextAt{13}{223.188}{202.543}{\Character{11}}%
\TextAt{9}{230.71}{202.543}{\Character{93}}%
\TextAt{9}{237.261}{202.543}{\Character{61}}%
\TextAt{16}{249.708}{202.543}{\Character{79}}%
\TextAt{8}{259.557}{202.543}{;}%
\TextAt{8}{266.298}{202.543}{then}%
\TextAt{14}{292.346}{202.543}{C}%
\TextAt{3}{299.669}{200.683}{\Character{79}}%
\TextAt{3}{306.603}{200.683}{\Character{106}}%
\TextAt{27}{308.964}{200.683}{o}%
\TextAt{9}{316.679}{202.543}{\Character{61}}%
\TextAt{13}{329.105}{202.543}{f}%
\TextAt{3}{336.151}{206.881}{\Character{48}}%
\TextAt{9}{338.946}{202.543}{\Character{40}}%
\TextAt{13}{343.499}{202.543}{\Character{11}}%
\TextAt{9}{351.02}{202.543}{\Character{41}}%
\TextAt{3}{355.573}{206.881}{\Character{0}}%
\TextAt{17}{362.159}{206.881}{1}%
\TextAt{16}{366.891}{202.543}{\Character{79}}%
\TextAt{8}{376.741}{202.543}{,}%
\TextAt{8}{383.4}{202.543}{and}%
\TextAt{8}{405.551}{202.543}{thus}%
\TextAt{14}{430.689}{202.543}{D}%
\TextAt{21}{440.631}{200.683}{L}%
\TextAt{3}{446.39}{200.683}{\Character{106}}%
\TextAt{21}{448.743}{200.683}{K}%
\TextAt{9}{460.266}{202.543}{\Character{61}}%
\TextAt{9}{472.675}{202.543}{\Character{40}}%
\TextAt{13}{477.244}{202.543}{f}%
\TextAt{3}{484.29}{206.881}{\Character{48}}%
\TextAt{9}{487.085}{202.543}{\Character{40}}%
\TextAt{13}{491.638}{202.543}{\Character{11}}%
\TextAt{9}{499.159}{202.543}{\Character{41}\Character{41}}%
\TextAt{8}{508.264}{202.543}{.}%
\Vector{q 0 G 0 g 0.399 w 452.334 -571.457 m 452.334 -563.388 l S Q}%
\Vector{q 0 G 0 g 0.399 w 452.534 -563.587 m 459.707 -563.587 l S Q}%
\Vector{q 0 G 0 g 0.399 w 452.534 -571.258 m 459.707 -571.258 l S Q}%
\Vector{q 0 G 0 g 0.399 w 459.906 -571.457 m 459.906 -563.388 l S Q}%
\TextAt{8}{97.159}{174.677}{The}%
\TextAt{8}{121.81}{174.677}{proof}%
\TextAt{8}{152.977}{174.677}{shows}%
\TextAt{8}{187.192}{174.677}{that}%
\TextAt{8}{213.158}{174.677}{the}%
\TextAt{8}{233.912}{174.677}{module}%
\TextAt{3}{275.802}{179.015}{\Character{3}}%
\TextAt{14}{280.534}{174.677}{o}%
\TextAt{9}{286.378}{174.677}{\Character{91}}%
\TextAt{13}{289.63}{174.677}{\Character{11}}%
\TextAt{9}{297.152}{174.677}{\Character{93}}%
\TextAt{9}{304.72}{174.677}{\Character{61}}%
\TextAt{16}{318.185}{174.677}{\Character{102}}%
\TextAt{13}{324.163}{174.677}{x}%
\NativeMathAt{335.154}{174.677}{11.955}{\in}%
\TextAt{13}{347.462}{174.677}{L}%
\TextAt{16}{359.766}{174.677}{\Character{106}}%
\TextAt{8}{367.425}{174.677}{Tr}%
\TextAt{21}{380.432}{172.817}{L}%
\TextAt{3}{386.191}{172.817}{\Character{106}}%
\TextAt{21}{388.544}{172.817}{K}%
\TextAt{9}{396.746}{174.677}{\Character{40}}%
\TextAt{13}{401.299}{174.677}{x}%
\TextAt{14}{407.951}{174.677}{o}%
\TextAt{9}{413.795}{174.677}{\Character{91}}%
\TextAt{13}{417.047}{174.677}{\Character{11}}%
\TextAt{9}{424.568}{174.677}{\Character{93}\Character{41}}%
\NativeMathAt{436.711}{174.677}{11.955}{\subseteq}%
\TextAt{14}{450.348}{174.677}{o}%
\TextAt{16}{456.192}{174.677}{\Character{103}}%
\TextAt{8}{462.17}{174.677}{,}%
\TextAt{8}{470.066}{174.677}{which}%
\TextAt{8}{504.149}{174.677}{is}%
\TextAt{8}{516.51}{174.677}{the}%
\TextAt{8}{79.226}{157.631}{dual}%
\TextAt{8}{104.469}{157.631}{of}%
\TextAt{8}{117.028}{157.631}{the}%
\TextAt{14}{136.413}{157.631}{o}%
\TextAt{8}{142.258}{157.631}{-module}%
\TextAt{14}{186.677}{157.631}{o}%
\TextAt{9}{192.521}{157.631}{\Character{91}}%
\TextAt{13}{195.773}{157.631}{\Character{11}}%
\TextAt{9}{203.295}{157.631}{\Character{93}}%
\TextAt{8}{206.546}{157.631}{,}%
\TextAt{8}{213.079}{157.631}{always}%
\TextAt{8}{249.775}{157.631}{admits}%
\TextAt{8}{287.44}{157.631}{the}%
\TextAt{14}{306.825}{157.631}{o}%
\TextAt{8}{312.669}{157.631}{-basis}%
\TextAt{13}{344.535}{157.631}{\Character{11}}%
\TextAt{21}{352.057}{161.969}{i}%
\TextAt{9}{355.438}{157.631}{\Character{47}}%
\TextAt{13}{361.291}{157.631}{f}%
\TextAt{3}{368.338}{161.969}{\Character{48}}%
\TextAt{9}{371.133}{157.631}{\Character{40}}%
\TextAt{13}{375.685}{157.631}{\Character{11}}%
\TextAt{9}{383.207}{157.631}{\Character{41}}%
\TextAt{8}{387.759}{157.631}{,}%
\TextAt{13}{394.292}{157.631}{i}%
\TextAt{9}{401.606}{157.631}{\Character{61}}%
\TextAt{9}{414.016}{157.631}{0}%
\TextAt{13}{419.884}{157.631}{\Character{59}}%
\TextAt{13}{425.108}{157.631}{\Character{58}}%
\TextAt{13}{430.345}{157.631}{\Character{58}}%
\TextAt{13}{435.581}{157.631}{\Character{58}}%
\TextAt{13}{440.806}{157.631}{\Character{59}}%
\TextAt{13}{446.042}{157.631}{n}%
\TextAt{16}{454.172}{157.631}{\Character{0}}%
\TextAt{9}{464.55}{157.631}{1}%
\TextAt{8}{470.403}{157.631}{.}%
\TextAt{8}{478.595}{157.631}{We}%
\TextAt{8}{497.968}{157.631}{exploit}%
\TextAt{8}{79.226}{140.585}{this}%
\TextAt{8}{102.633}{140.585}{for}%
\TextAt{8}{121.102}{140.585}{the}%
\TextAt{8}{141.842}{140.585}{following}%
\TextAt{8}{191.861}{140.585}{characterization}%
\TextAt{8}{277.985}{140.585}{of}%
\TextAt{8}{291.899}{140.585}{the}%
\TextAt{8}{312.64}{140.585}{different}%
\TextAt{8}{359.395}{140.585}{in}%
\TextAt{8}{373.632}{140.585}{the}%
\TextAt{8}{394.372}{140.585}{general}%
\TextAt{8}{435.281}{140.585}{case}%
\TextAt{8}{460.636}{140.585}{where}%
\TextAt{16}{495.03}{140.585}{\Character{79}}%
\TextAt{8}{509.361}{140.585}{need}%
\TextAt{8}{79.226}{123.539}{not}%
\TextAt{8}{100.04}{123.539}{be}%
\TextAt{8}{115.964}{123.539}{monogenous.}%
\TextAt{12}{79.226}{98.662}{Theorem}%
\TextAt{12}{136.855}{98.662}{(2.5).}%
\TextAt{8}{172.705}{98.662}{The}%
\TextAt{8}{197.293}{98.662}{different}%
\TextAt{14}{243.999}{98.662}{D}%
\TextAt{21}{253.94}{96.802}{L}%
\TextAt{3}{259.699}{96.802}{\Character{106}}%
\TextAt{21}{262.052}{96.802}{K}%
\TextAt{8}{274.686}{98.662}{is}%
\TextAt{8}{286.984}{98.662}{the}%
\TextAt{8}{307.675}{98.662}{ideal}%
\TextAt{8}{336.173}{98.662}{generated}%
\TextAt{8}{390.039}{98.662}{by}%
\TextAt{8}{406.833}{98.662}{all}%
\TextAt{8}{423.626}{98.662}{differents}%
\TextAt{8}{474.946}{98.662}{of}%
\TextAt{8}{488.81}{98.662}{elements}%
\TextAt{13}{79.226}{81.616}{\Character{14}}%
\TextAt{21}{84.411}{79.756}{L}%
\TextAt{3}{90.17}{79.756}{\Character{106}}%
\TextAt{21}{92.522}{79.756}{K}%
\TextAt{9}{100.725}{81.616}{\Character{40}}%
\TextAt{13}{105.277}{81.616}{\Character{11}}%
\TextAt{9}{112.799}{81.616}{\Character{41}}%
\TextAt{8}{121.249}{81.616}{for}%
\TextAt{13}{139.134}{81.616}{\Character{11}}%
\NativeMathAt{149.976}{81.616}{11.955}{\in}%
\TextAt{16}{161.25}{81.616}{\Character{79}}%
\TextAt{8}{171.117}{81.616}{.}%
\TextAt{8}{297.211}{43.033}{180}%
\Destination{78.23}{809.815}{nav.page.180}%
\Destination{79.36}{776.155}{nav.print.197}%
\Destination{79.36}{189.655}{nav.print.198}%
\Destination{79.36}{750.255}{nav.number.3.2.4}%
\Destination{79.36}{112.655}{nav.number.3.2.5}%
}

% Page 193 (body page 181)
\Page{%
\TextAt{2}{79.226}{796.943}{III.}%
\TextAt{2}{98.317}{796.943}{Riemann–Roch}%
\TextAt{2}{174.375}{796.943}{Theory}%
\Vector{q 0 G 0 g 0.299 w 7.226 18.729 m 460.769 18.729 l S Q}%
\TextAt{19}{79.226}{761.124}{Proof.}%
\TextAt{8}{116}{761.124}{Let}%
\TextAt{13}{136.768}{761.124}{\Character{11}}%
\NativeMathAt{147.611}{761.124}{11.955}{\in}%
\TextAt{16}{158.885}{761.124}{\Character{79}}%
\TextAt{8}{172.448}{761.124}{such}%
\TextAt{8}{198.643}{761.124}{that}%
\TextAt{13}{223.811}{761.124}{L}%
\TextAt{9}{235.096}{761.124}{\Character{61}}%
\TextAt{13}{247.522}{761.124}{K}%
\TextAt{9}{258.339}{761.124}{\Character{40}}%
\TextAt{13}{262.892}{761.124}{\Character{11}}%
\TextAt{9}{270.413}{761.124}{\Character{41}}%
\TextAt{8}{274.966}{761.124}{,}%
\TextAt{8}{281.954}{761.124}{and}%
\TextAt{8}{304.515}{761.124}{let}%
\TextAt{13}{321.218}{761.124}{f}%
\TextAt{9}{328.265}{761.124}{\Character{40}}%
\TextAt{13}{332.817}{761.124}{X}%
\TextAt{9}{343.472}{761.124}{\Character{41}}%
\TextAt{8}{351.72}{761.124}{be}%
\TextAt{8}{367.443}{761.124}{the}%
\TextAt{8}{387.398}{761.124}{minimal}%
\TextAt{8}{432.722}{761.124}{polynomial}%
\TextAt{8}{493.014}{761.124}{of}%
\TextAt{13}{506.143}{761.124}{\Character{11}}%
\TextAt{8}{513.665}{761.124}{.}%
\TextAt{8}{522.046}{761.124}{In}%
\TextAt{8}{79.226}{744.078}{order}%
\TextAt{8}{109.904}{744.078}{to}%
\TextAt{8}{124.323}{744.078}{show}%
\TextAt{8}{153.434}{744.078}{that}%
\TextAt{13}{178.912}{744.078}{f}%
\TextAt{3}{185.958}{748.416}{\Character{48}}%
\TextAt{9}{188.753}{744.078}{\Character{40}}%
\TextAt{13}{193.305}{744.078}{\Character{11}}%
\TextAt{9}{200.827}{744.078}{\Character{41}}%
\NativeMathAt{208.885}{744.078}{11.955}{\in}%
\TextAt{14}{220.361}{744.078}{D}%
\TextAt{21}{230.303}{742.218}{L}%
\TextAt{3}{236.062}{742.218}{\Character{106}}%
\TextAt{21}{238.414}{742.218}{K}%
\TextAt{8}{246.617}{744.078}{,}%
\TextAt{8}{253.901}{744.078}{we}%
\TextAt{8}{271.237}{744.078}{consider}%
\TextAt{8}{316.931}{744.078}{the}%
\TextAt{8}{337.195}{744.078}{“conductor”}%
\TextAt{14}{402.986}{744.078}{f}%
\TextAt{9}{410.392}{744.078}{\Character{61}}%
\TextAt{14}{423.002}{744.078}{f}%
\TextAt{27}{426.902}{742.218}{o}%
\TextAt{17}{430.798}{742.218}{\Character{91}}%
\TextAt{21}{433.151}{742.218}{\Character{11}}%
\TextAt{17}{438.59}{742.218}{\Character{93}}%
\TextAt{9}{444.946}{744.078}{\Character{61}}%
\TextAt{16}{457.556}{744.078}{\Character{102}}%
\TextAt{13}{463.534}{744.078}{x}%
\NativeMathAt{473.691}{744.078}{11.955}{\in}%
\TextAt{13}{485.167}{744.078}{L}%
\TextAt{16}{496.637}{744.078}{\Character{106}}%
\TextAt{13}{503.464}{744.078}{x}%
\TextAt{16}{510.116}{744.078}{\Character{79}}%
\NativeMathAt{523.458}{744.078}{11.955}{\subseteq}%
\TextAt{14}{79.226}{727.032}{o}%
\TextAt{9}{85.07}{727.032}{\Character{91}}%
\TextAt{13}{88.322}{727.032}{\Character{11}}%
\TextAt{9}{95.844}{727.032}{\Character{93}}%
\TextAt{16}{99.095}{727.032}{\Character{103}}%
\TextAt{8}{108.97}{727.032}{of}%
\TextAt{14}{122.3}{727.032}{o}%
\TextAt{9}{128.145}{727.032}{\Character{91}}%
\TextAt{13}{131.396}{727.032}{\Character{11}}%
\TextAt{9}{138.918}{727.032}{\Character{93}}%
\TextAt{8}{146.067}{727.032}{(see}%
\TextAt{8}{169.535}{727.032}{chap.}%
\TextAt{8}{201.726}{727.032}{I,}%
\TextAt{8}{213.096}{727.032}{§12,}%
\TextAt{8}{237.64}{727.032}{p.79).}%
\TextAt{8}{272.114}{727.032}{On}%
\TextAt{8}{291.625}{727.032}{putting}%
\TextAt{13}{333.253}{727.032}{b}%
\TextAt{9}{341.551}{727.032}{\Character{61}}%
\TextAt{13}{353.977}{727.032}{f}%
\TextAt{3}{361.023}{731.37}{\Character{48}}%
\TextAt{9}{363.818}{727.032}{\Character{40}}%
\TextAt{13}{368.37}{727.032}{\Character{11}}%
\TextAt{9}{375.892}{727.032}{\Character{41}}%
\TextAt{8}{380.444}{727.032}{,}%
\TextAt{8}{387.594}{727.032}{we}%
\TextAt{8}{404.821}{727.032}{have}%
\TextAt{8}{431.816}{727.032}{for}%
\TextAt{13}{449.701}{727.032}{x}%
\NativeMathAt{459.674}{727.032}{11.955}{\in}%
\TextAt{13}{470.965}{727.032}{L}%
\TextAt{8}{478.929}{727.032}{:}%
\TextAt{13}{142.485}{703.578}{x}%
\NativeMathAt{152.458}{703.578}{11.955}{\in}%
\TextAt{14}{163.749}{703.578}{f}%
\NativeArrowAt{174.291}{703.578}{11.955}{21.926}{\Longleftrightarrow}%
\TextAt{13}{202.85}{703.578}{x}%
\TextAt{16}{209.503}{703.578}{\Character{79}}%
\NativeMathAt{222.666}{703.578}{11.955}{\subseteq}%
\TextAt{14}{235.293}{703.578}{o}%
\TextAt{9}{241.137}{703.578}{\Character{91}}%
\TextAt{13}{244.388}{703.578}{\Character{11}}%
\TextAt{9}{251.91}{703.578}{\Character{93}}%
\NativeArrowAt{261.803}{703.578}{11.955}{21.926}{\Longleftrightarrow}%
\TextAt{13}{290.363}{703.578}{b}%
\TextAt{3}{295.34}{708.514}{\Character{0}}%
\TextAt{17}{301.927}{708.514}{1}%
\TextAt{13}{306.659}{703.578}{x}%
\TextAt{16}{313.311}{703.578}{\Character{79}}%
\NativeMathAt{326.474}{703.578}{11.955}{\subseteq}%
\TextAt{13}{339.101}{703.578}{b}%
\TextAt{3}{344.078}{708.514}{\Character{0}}%
\TextAt{17}{350.665}{708.514}{1}%
\TextAt{14}{355.397}{703.578}{o}%
\TextAt{9}{361.241}{703.578}{\Character{91}}%
\TextAt{13}{364.493}{703.578}{\Character{11}}%
\TextAt{9}{372.015}{703.578}{\Character{93}}%
\TextAt{9}{378.566}{703.578}{\Character{61}}%
\TextAt{3}{391.012}{708.514}{\Character{3}}%
\TextAt{14}{395.745}{703.578}{o}%
\TextAt{9}{401.589}{703.578}{\Character{91}}%
\TextAt{13}{404.841}{703.578}{\Character{11}}%
\TextAt{9}{412.362}{703.578}{\Character{93}}%
\NativeArrowAt{174.291}{683.543}{11.955}{21.926}{\Longleftrightarrow}%
\TextAt{8}{202.85}{683.543}{Tr}%
\TextAt{9}{215.858}{683.543}{\Character{40}}%
\TextAt{13}{220.41}{683.543}{b}%
\TextAt{3}{225.387}{688.479}{\Character{0}}%
\TextAt{17}{231.974}{688.479}{1}%
\TextAt{13}{236.706}{683.543}{x}%
\TextAt{16}{243.358}{683.543}{\Character{79}}%
\TextAt{9}{253.208}{683.543}{\Character{41}}%
\NativeMathAt{261.081}{683.543}{11.955}{\subseteq}%
\TextAt{14}{273.7}{683.543}{o}%
\NativeArrowAt{286.186}{683.543}{11.955}{21.926}{\Longleftrightarrow}%
\TextAt{13}{314.746}{683.543}{b}%
\TextAt{3}{319.723}{688.479}{\Character{0}}%
\TextAt{17}{326.309}{688.479}{1}%
\TextAt{13}{331.042}{683.543}{x}%
\NativeMathAt{341.015}{683.543}{11.955}{\in}%
\TextAt{14}{352.306}{683.543}{D}%
\TextAt{3}{362.247}{688.491}{\Character{0}}%
\TextAt{17}{368.834}{688.491}{1}%
\TextAt{21}{362.247}{679.77}{L}%
\TextAt{3}{368.006}{679.77}{\Character{106}}%
\TextAt{21}{370.359}{679.77}{K}%
\NativeArrowAt{385.203}{683.543}{11.955}{21.926}{\Longleftrightarrow}%
\TextAt{13}{413.763}{683.543}{x}%
\NativeMathAt{423.735}{683.543}{11.955}{\in}%
\TextAt{13}{435.026}{683.543}{b}%
\TextAt{14}{440.004}{683.543}{D}%
\TextAt{3}{449.945}{688.491}{\Character{0}}%
\TextAt{17}{456.532}{688.491}{1}%
\TextAt{21}{449.945}{679.77}{L}%
\TextAt{3}{455.704}{679.77}{\Character{106}}%
\TextAt{21}{458.057}{679.77}{K}%
\TextAt{13}{466.259}{683.543}{\Character{58}}%
\TextAt{8}{79.226}{660.089}{Therefore}%
\TextAt{9}{132.664}{660.089}{\Character{40}}%
\TextAt{13}{137.216}{660.089}{f}%
\TextAt{3}{144.263}{664.427}{\Character{48}}%
\TextAt{9}{147.058}{660.089}{\Character{40}}%
\TextAt{13}{151.61}{660.089}{\Character{11}}%
\TextAt{9}{159.132}{660.089}{\Character{41}\Character{41}}%
\TextAt{9}{172.283}{660.089}{\Character{61}}%
\TextAt{14}{185.484}{660.089}{f}%
\TextAt{27}{189.385}{658.229}{o}%
\TextAt{17}{193.281}{658.229}{\Character{91}}%
\TextAt{21}{195.633}{658.229}{\Character{11}}%
\TextAt{17}{201.072}{658.229}{\Character{93}}%
\TextAt{14}{203.922}{660.089}{D}%
\TextAt{21}{213.864}{658.229}{L}%
\TextAt{3}{219.623}{658.229}{\Character{106}}%
\TextAt{21}{221.976}{658.229}{K}%
\TextAt{8}{230.178}{660.089}{,}%
\TextAt{8}{237.878}{660.089}{so}%
\TextAt{8}{252.689}{660.089}{in}%
\TextAt{8}{266.782}{660.089}{particular,}%
\TextAt{13}{324.574}{660.089}{f}%
\TextAt{3}{331.62}{664.427}{\Character{48}}%
\TextAt{9}{334.416}{660.089}{\Character{40}}%
\TextAt{13}{338.968}{660.089}{\Character{11}}%
\TextAt{9}{346.49}{660.089}{\Character{41}}%
\NativeMathAt{355.113}{660.089}{11.955}{\in}%
\TextAt{14}{367.155}{660.089}{D}%
\TextAt{21}{377.097}{658.229}{L}%
\TextAt{3}{382.856}{658.229}{\Character{106}}%
\TextAt{21}{385.208}{658.229}{K}%
\TextAt{8}{393.411}{660.089}{.}%
\TextAt{14}{403.181}{660.089}{D}%
\TextAt{21}{413.122}{658.229}{L}%
\TextAt{3}{418.881}{658.229}{\Character{106}}%
\TextAt{21}{421.234}{658.229}{K}%
\TextAt{8}{433.774}{660.089}{thus}%
\TextAt{8}{459.966}{660.089}{divides}%
\TextAt{8}{499.811}{660.089}{all}%
\TextAt{8}{516.511}{660.089}{the}%
\TextAt{8}{79.226}{643.043}{differents}%
\TextAt{8}{129.491}{643.043}{of}%
\TextAt{8}{142.301}{643.043}{elements}%
\TextAt{13}{189.637}{643.043}{\Character{14}}%
\TextAt{21}{194.822}{641.183}{L}%
\TextAt{3}{200.581}{641.183}{\Character{106}}%
\TextAt{21}{202.933}{641.183}{K}%
\TextAt{9}{211.136}{643.043}{\Character{40}}%
\TextAt{13}{215.688}{643.043}{\Character{11}}%
\TextAt{9}{223.21}{643.043}{\Character{41}}%
\TextAt{8}{227.762}{643.043}{.}%
\TextAt{8}{236.037}{643.043}{In}%
\TextAt{8}{250.138}{643.043}{order}%
\TextAt{8}{280.187}{643.043}{to}%
\TextAt{8}{293.977}{643.043}{prove}%
\TextAt{8}{325.006}{643.043}{that}%
\TextAt{14}{349.855}{643.043}{D}%
\TextAt{21}{359.797}{641.183}{L}%
\TextAt{3}{365.556}{641.183}{\Character{106}}%
\TextAt{21}{367.908}{641.183}{K}%
\TextAt{8}{379.488}{643.043}{is}%
\TextAt{8}{390.731}{643.043}{in}%
\TextAt{8}{403.864}{643.043}{fact}%
\TextAt{8}{426.429}{643.043}{the}%
\TextAt{8}{446.065}{643.043}{greatest}%
\TextAt{8}{489.838}{643.043}{common}%
\TextAt{8}{79.226}{625.997}{divisor}%
\TextAt{8}{118.271}{625.997}{of}%
\TextAt{8}{132.533}{625.997}{all}%
\TextAt{13}{149.724}{625.997}{\Character{14}}%
\TextAt{21}{154.908}{624.137}{L}%
\TextAt{3}{160.667}{624.137}{\Character{106}}%
\TextAt{21}{163.02}{624.137}{K}%
\TextAt{9}{171.222}{625.997}{\Character{40}}%
\TextAt{13}{175.775}{625.997}{\Character{11}}%
\TextAt{9}{183.296}{625.997}{\Character{41}}%
\TextAt{8}{187.849}{625.997}{,}%
\TextAt{8}{196.163}{625.997}{it}%
\TextAt{8}{208.798}{625.997}{suffices}%
\TextAt{8}{249.517}{625.997}{to}%
\TextAt{8}{264.759}{625.997}{show}%
\TextAt{8}{294.694}{625.997}{that,}%
\TextAt{8}{324.48}{625.997}{for}%
\TextAt{8}{343.296}{625.997}{every}%
\TextAt{8}{375.12}{625.997}{prime}%
\TextAt{8}{409.216}{625.997}{ideal}%
\TextAt{14}{438.111}{625.997}{P}%
\TextAt{8}{448.004}{625.997}{,}%
\TextAt{8}{456.318}{625.997}{there}%
\TextAt{8}{487.161}{625.997}{exists}%
\TextAt{8}{520.408}{625.997}{an}%
\TextAt{13}{79.226}{608.951}{\Character{11}}%
\NativeMathAt{90.069}{608.951}{11.955}{\in}%
\TextAt{16}{101.343}{608.951}{\Character{79}}%
\TextAt{8}{115.107}{608.951}{such}%
\TextAt{8}{141.504}{608.951}{that}%
\TextAt{13}{166.873}{608.951}{L}%
\TextAt{9}{178.158}{608.951}{\Character{61}}%
\TextAt{13}{190.584}{608.951}{K}%
\TextAt{9}{201.401}{608.951}{\Character{40}}%
\TextAt{13}{205.954}{608.951}{\Character{11}}%
\TextAt{9}{213.475}{608.951}{\Character{41}}%
\TextAt{8}{221.925}{608.951}{and}%
\TextAt{13}{244.688}{608.951}{v}%
\TextAt{27}{250.347}{607.158}{P}%
\TextAt{9}{257.441}{608.951}{\Character{40}}%
\TextAt{14}{261.993}{608.951}{D}%
\TextAt{21}{271.934}{607.091}{L}%
\TextAt{3}{277.694}{607.091}{\Character{106}}%
\TextAt{21}{280.046}{607.091}{K}%
\TextAt{9}{288.249}{608.951}{\Character{41}}%
\TextAt{9}{296.104}{608.951}{\Character{61}}%
\TextAt{13}{308.547}{608.951}{v}%
\TextAt{27}{314.206}{607.158}{P}%
\TextAt{9}{321.3}{608.951}{\Character{40}}%
\TextAt{13}{325.852}{608.951}{f}%
\TextAt{3}{332.899}{613.289}{\Character{48}}%
\TextAt{9}{335.694}{608.951}{\Character{40}}%
\TextAt{13}{340.246}{608.951}{\Character{11}}%
\TextAt{9}{347.768}{608.951}{\Character{41}\Character{41}}%
\TextAt{8}{356.873}{608.951}{.}%
\TextAt{8}{97.159}{588.916}{We}%
\TextAt{8}{117.858}{588.916}{think}%
\TextAt{8}{149.304}{588.916}{of}%
\TextAt{13}{163.189}{588.916}{L}%
\TextAt{8}{175.605}{588.916}{as}%
\TextAt{8}{190.529}{588.916}{embedded}%
\TextAt{8}{246.352}{588.916}{into}%
\TextAt{8}{270.65}{588.916}{the}%
\TextAt{8}{291.361}{588.916}{separable}%
\TextAt{8}{343.358}{588.916}{closure}%
\Vector{q 0 G 0 g 0.478 w 310.994 -175.241 m 321.812 -175.241 l S Q}%
\TextAt{13}{382.994}{588.916}{K}%
\TextAt{27}{393.812}{587.123}{p}%
\TextAt{8}{402.75}{588.916}{of}%
\TextAt{13}{416.635}{588.916}{K}%
\TextAt{27}{426.594}{587.123}{p}%
\TextAt{8}{435.532}{588.916}{in}%
\TextAt{8}{449.739}{588.916}{such}%
\TextAt{8}{476.691}{588.916}{a}%
\TextAt{8}{487.001}{588.916}{way}%
\TextAt{8}{511.298}{588.916}{that}%
\TextAt{8}{79.226}{571.87}{the}%
\TextAt{8}{99.383}{571.87}{absolute}%
\TextAt{8}{145.625}{571.87}{value}%
\TextAt{16}{175.872}{571.87}{\Character{106}}%
\TextAt{16}{183.069}{571.87}{\Character{106}}%
\TextAt{8}{190.308}{571.87}{of}%
\Vector{q 0 G 0 g 0.478 w 131.638 -192.287 m 142.456 -192.287 l S Q}%
\TextAt{13}{203.638}{571.87}{K}%
\TextAt{27}{214.456}{570.077}{p}%
\TextAt{8}{222.84}{571.87}{defines}%
\TextAt{8}{261.264}{571.87}{the}%
\TextAt{8}{281.421}{571.87}{prime}%
\TextAt{14}{314.584}{571.87}{P}%
\TextAt{8}{324.477}{571.87}{.}%
\TextAt{8}{97.159}{551.835}{By}%
\TextAt{8}{115.625}{551.835}{chap.}%
\TextAt{8}{148.124}{551.835}{II,}%
\TextAt{8}{164.139}{551.835}{(10.4),}%
\TextAt{8}{201.352}{551.835}{we}%
\TextAt{8}{218.682}{551.835}{find}%
\TextAt{8}{242.193}{551.835}{an}%
\TextAt{8}{258.555}{551.835}{element}%
\TextAt{13}{301.899}{551.835}{\Character{12}}%
\TextAt{8}{313.17}{551.835}{in}%
\TextAt{8}{326.926}{551.835}{the}%
\TextAt{8}{347.185}{551.835}{valuation}%
\TextAt{8}{398.36}{551.835}{ring}%
\TextAt{16}{422.528}{551.835}{\Character{79}}%
\TextAt{27}{432.046}{550.042}{P}%
\TextAt{8}{443.14}{551.835}{of}%
\TextAt{8}{456.572}{551.835}{the}%
\TextAt{8}{476.831}{551.835}{completion}%
\TextAt{13}{79.226}{534.789}{L}%
\TextAt{27}{87.191}{532.996}{P}%
\TextAt{8}{98.499}{534.789}{satisfying}%
\TextAt{16}{150.976}{534.789}{\Character{79}}%
\TextAt{27}{160.493}{532.996}{P}%
\TextAt{9}{171.448}{534.789}{\Character{61}}%
\TextAt{14}{184.413}{534.789}{o}%
\TextAt{27}{190.258}{532.996}{p}%
\TextAt{9}{194.744}{534.789}{\Character{91}}%
\TextAt{13}{197.996}{534.789}{\Character{12}}%
\TextAt{9}{205.267}{534.789}{\Character{93}}%
\TextAt{8}{208.519}{534.789}{,}%
\TextAt{8}{216.064}{534.789}{and}%
\TextAt{8}{239.144}{534.789}{the}%
\TextAt{8}{259.617}{534.789}{proof}%
\TextAt{19}{290.503}{534.789}{loc.}%
\TextAt{19}{314.438}{534.789}{cit.}%
\TextAt{8}{337.024}{534.789}{shows}%
\TextAt{8}{370.958}{534.789}{that,}%
\TextAt{8}{399.975}{534.789}{for}%
\TextAt{8}{418.177}{534.789}{every}%
\TextAt{8}{449.386}{534.789}{element}%
\TextAt{13}{492.945}{534.789}{\Character{11}}%
\NativeMathAt{504.327}{534.789}{11.955}{\in}%
\TextAt{16}{516.139}{534.789}{\Character{79}}%
\TextAt{27}{525.676}{532.996}{P}%
\TextAt{8}{79.226}{517.743}{which}%
\TextAt{8}{113.443}{517.743}{is}%
\TextAt{8}{125.938}{517.743}{sufficiently}%
\TextAt{8}{185.261}{517.743}{close}%
\TextAt{8}{214.015}{517.743}{to}%
\TextAt{13}{229.056}{517.743}{\Character{12}}%
\TextAt{8}{236.327}{517.743}{,}%
\TextAt{8}{244.39}{517.743}{one}%
\TextAt{8}{266.58}{517.743}{also}%
\TextAt{8}{290.791}{517.743}{has}%
\TextAt{16}{312.395}{517.743}{\Character{79}}%
\TextAt{27}{321.913}{515.95}{P}%
\TextAt{9}{333.572}{517.743}{\Character{61}}%
\TextAt{14}{347.243}{517.743}{o}%
\TextAt{27}{353.087}{515.95}{p}%
\TextAt{9}{357.574}{517.743}{\Character{91}}%
\TextAt{13}{360.826}{517.743}{\Character{11}}%
\TextAt{9}{368.348}{517.743}{\Character{93}}%
\TextAt{8}{371.599}{517.743}{.}%
\TextAt{8}{382.24}{517.743}{From}%
\TextAt{8}{413.695}{517.743}{(2.2),}%
\TextAt{8}{445.836}{517.743}{(iii)}%
\TextAt{8}{469.329}{517.743}{and}%
\TextAt{8}{492.822}{517.743}{(2.4),}%
\TextAt{8}{524.963}{517.743}{it}%
\TextAt{8}{79.226}{500.697}{follows}%
\TextAt{8}{117.662}{500.697}{that}%
\TextAt{13}{200.804}{483.651}{v}%
\TextAt{27}{206.464}{481.858}{P}%
\TextAt{9}{213.557}{483.651}{\Character{40}}%
\TextAt{14}{218.11}{483.651}{D}%
\TextAt{21}{228.051}{481.792}{L}%
\TextAt{3}{233.81}{481.792}{\Character{106}}%
\TextAt{21}{236.163}{481.792}{K}%
\TextAt{9}{244.365}{483.651}{\Character{41}}%
\TextAt{9}{252.22}{483.651}{\Character{61}}%
\TextAt{13}{264.664}{483.651}{v}%
\TextAt{27}{270.323}{481.858}{P}%
\TextAt{9}{277.417}{483.651}{\Character{40}}%
\TextAt{14}{281.969}{483.651}{D}%
\TextAt{21}{291.911}{481.792}{L}%
\TextAt{28}{297.67}{480.401}{P}%
\TextAt{3}{303.867}{481.792}{\Character{106}}%
\TextAt{21}{306.219}{481.792}{K}%
\TextAt{28}{313.361}{480.796}{p}%
\TextAt{9}{317.942}{483.651}{\Character{41}}%
\TextAt{9}{325.797}{483.651}{\Character{61}}%
\TextAt{13}{338.24}{483.651}{v}%
\TextAt{27}{343.9}{481.858}{P}%
\TextAt{9}{350.993}{483.651}{\Character{40}}%
\TextAt{13}{355.545}{483.651}{\Character{14}}%
\TextAt{21}{360.73}{481.792}{L}%
\TextAt{28}{366.489}{480.401}{P}%
\TextAt{3}{372.686}{481.792}{\Character{106}}%
\TextAt{21}{375.039}{481.792}{K}%
\TextAt{28}{382.18}{480.796}{p}%
\TextAt{9}{386.761}{483.651}{\Character{40}}%
\TextAt{13}{391.313}{483.651}{\Character{11}}%
\TextAt{9}{398.835}{483.651}{\Character{41}\Character{41}}%
\TextAt{13}{407.94}{483.651}{\Character{58}}%
\TextAt{8}{79.226}{462.507}{It}%
\TextAt{8}{91.734}{462.507}{therefore}%
\TextAt{8}{140.669}{462.507}{suffices}%
\TextAt{8}{180.292}{462.507}{to}%
\TextAt{8}{194.438}{462.507}{show}%
\TextAt{8}{223.276}{462.507}{that}%
\TextAt{8}{248.481}{462.507}{we}%
\TextAt{8}{265.544}{462.507}{can}%
\TextAt{8}{286.839}{462.507}{find}%
\TextAt{8}{310.082}{462.507}{an}%
\TextAt{8}{326.177}{462.507}{element}%
\TextAt{13}{369.254}{462.507}{\Character{11}}%
\TextAt{8}{380.509}{462.507}{in}%
\TextAt{16}{393.997}{462.507}{\Character{79}}%
\TextAt{8}{407.58}{462.507}{such}%
\TextAt{8}{433.812}{462.507}{that}%
\TextAt{13}{459.016}{462.507}{L}%
\TextAt{9}{470.302}{462.507}{\Character{61}}%
\TextAt{13}{482.727}{462.507}{K}%
\TextAt{9}{493.545}{462.507}{\Character{40}}%
\TextAt{13}{498.097}{462.507}{\Character{11}}%
\TextAt{9}{505.619}{462.507}{\Character{41}}%
\TextAt{8}{513.904}{462.507}{and}%
\TextAt{13}{231.658}{439.053}{v}%
\TextAt{27}{237.317}{437.26}{P}%
\TextAt{9}{244.411}{439.053}{\Character{40}}%
\TextAt{13}{248.963}{439.053}{\Character{14}}%
\TextAt{21}{254.148}{437.193}{L}%
\TextAt{28}{259.907}{435.802}{P}%
\TextAt{3}{266.104}{437.193}{\Character{106}}%
\TextAt{21}{268.456}{437.193}{K}%
\TextAt{28}{275.598}{436.197}{p}%
\TextAt{9}{280.179}{439.053}{\Character{40}}%
\TextAt{13}{284.731}{439.053}{\Character{11}}%
\TextAt{9}{292.253}{439.053}{\Character{41}\Character{41}}%
\TextAt{9}{304.651}{439.053}{\Character{61}}%
\TextAt{13}{317.104}{439.053}{v}%
\TextAt{27}{322.763}{437.26}{P}%
\TextAt{9}{329.856}{439.053}{\Character{40}}%
\TextAt{13}{334.409}{439.053}{\Character{14}}%
\TextAt{21}{339.593}{437.193}{L}%
\TextAt{3}{345.352}{437.193}{\Character{106}}%
\TextAt{21}{347.705}{437.193}{K}%
\TextAt{9}{355.907}{439.053}{\Character{40}}%
\TextAt{13}{360.46}{439.053}{\Character{11}}%
\TextAt{9}{367.981}{439.053}{\Character{41}\Character{41}}%
\TextAt{13}{377.086}{439.053}{\Character{58}}%
\TextAt{8}{79.226}{415.599}{For}%
\TextAt{8}{99.726}{415.599}{this,}%
\TextAt{8}{125.424}{415.599}{let}%
\TextAt{13}{141.858}{415.599}{\Character{27}}%
\TextAt{17}{148.512}{413.806}{2}%
\TextAt{13}{153.244}{415.599}{\Character{59}}%
\TextAt{13}{158.468}{415.599}{\Character{58}}%
\TextAt{13}{163.705}{415.599}{\Character{58}}%
\TextAt{13}{168.941}{415.599}{\Character{58}}%
\TextAt{13}{174.166}{415.599}{\Character{59}}%
\TextAt{13}{179.402}{415.599}{\Character{27}}%
\TextAt{21}{186.118}{413.806}{r}%
\TextAt{9}{193.994}{415.599}{\Character{58}}%
\TextAt{13}{200.566}{415.599}{L}%
\NativeMathAt{211.852}{415.599}{11.955}{\rightarrow}%
\Vector{q 0 G 0 g 0.478 w 155.128 -348.558 m 165.945 -348.558 l S Q}%
\TextAt{13}{227.128}{415.599}{K}%
\TextAt{27}{237.945}{413.806}{p}%
\TextAt{8}{245.86}{415.599}{be}%
\TextAt{13}{261.314}{415.599}{K}%
\TextAt{8}{272.132}{415.599}{-embeddings}%
\TextAt{8}{339.352}{415.599}{giving}%
\TextAt{8}{373.683}{415.599}{the}%
\TextAt{8}{393.37}{415.599}{primes}%
\TextAt{14}{430.678}{415.599}{P}%
\TextAt{21}{440.571}{413.806}{i}%
\TextAt{16}{447.273}{415.599}{\Character{106}}%
\TextAt{14}{453.915}{415.599}{p}%
\TextAt{8}{463.326}{415.599}{different}%
\TextAt{8}{509.027}{415.599}{from}%
\TextAt{14}{79.226}{398.553}{P}%
\TextAt{8}{89.119}{398.553}{.}%
\TextAt{8}{97.568}{398.553}{Let}%
\TextAt{13}{118.537}{398.553}{a}%
\NativeMathAt{128.003}{398.553}{11.955}{\in}%
\TextAt{14}{139.294}{398.553}{o}%
\TextAt{27}{145.138}{396.76}{p}%
\TextAt{8}{153.522}{398.553}{be}%
\TextAt{8}{169.447}{398.553}{an}%
\TextAt{8}{185.706}{398.553}{element}%
\TextAt{8}{228.947}{398.553}{such}%
\TextAt{8}{255.344}{398.553}{that}%
\TextAt{16}{189.726}{375.099}{\Character{106}}%
\TextAt{13}{193.047}{375.099}{\Character{28}}%
\TextAt{13}{199.455}{375.099}{\Character{12}}%
\TextAt{16}{209.394}{375.099}{\Character{0}}%
\TextAt{13}{221.349}{375.099}{a}%
\TextAt{16}{227.494}{375.099}{\Character{106}}%
\TextAt{9}{234.136}{375.099}{\Character{61}}%
\TextAt{9}{246.546}{375.099}{1}%
\TextAt{8}{276.325}{375.099}{for}%
\TextAt{8}{294.21}{375.099}{all}%
\TextAt{13}{310.469}{375.099}{\Character{28}}%
\NativeMathAt{320.209}{375.099}{11.955}{\in}%
\TextAt{13}{331.5}{375.099}{G}%
\TextAt{27}{340.733}{373.306}{p}%
\TextAt{9}{348.541}{375.099}{\Character{61}}%
\TextAt{13}{360.966}{375.099}{G}%
\TextAt{9}{370.2}{375.099}{\Character{40}}%
\Vector{q 0 G 0 g 0.478 w 302.752 -389.058 m 313.57 -389.058 l S Q}%
\TextAt{13}{374.752}{375.099}{K}%
\TextAt{27}{385.57}{373.306}{p}%
\TextAt{16}{393.378}{375.099}{\Character{106}}%
\TextAt{13}{400.019}{375.099}{K}%
\TextAt{27}{409.978}{373.306}{p}%
\TextAt{9}{414.465}{375.099}{\Character{41}}%
\TextAt{13}{419.018}{375.099}{\Character{58}}%
\TextAt{8}{517.802}{375.099}{(*)}%
\TextAt{8}{79.226}{351.645}{(Choose}%
\TextAt{13}{123.266}{351.645}{a}%
\TextAt{9}{132.731}{351.645}{\Character{61}}%
\TextAt{9}{145.141}{351.645}{1}%
\TextAt{8}{151.01}{351.645}{,}%
\TextAt{8}{157.181}{351.645}{resp.}%
\TextAt{13}{186.095}{351.645}{a}%
\TextAt{9}{195.561}{351.645}{\Character{61}}%
\TextAt{9}{207.971}{351.645}{0}%
\TextAt{8}{213.84}{351.645}{,}%
\TextAt{8}{220.011}{351.645}{according}%
\TextAt{8}{271.474}{351.645}{as}%
\TextAt{8}{284.622}{351.645}{the}%
\TextAt{8}{303.555}{351.645}{residue}%
\TextAt{8}{342.06}{351.645}{classes}%
\TextAt{13}{378.089}{351.645}{\Character{28}}%
\TextAt{13}{384.497}{351.645}{\Character{12}}%
\TextAt{8}{395.1}{351.645}{mod}%
\TextAt{14}{420.538}{351.645}{P}%
\TextAt{8}{433.106}{351.645}{which}%
\TextAt{8}{465.37}{351.645}{are}%
\TextAt{8}{483.658}{351.645}{conjugate}%
\TextAt{8}{79.226}{334.599}{over}%
\TextAt{14}{104.178}{334.599}{o}%
\TextAt{27}{110.022}{332.805}{p}%
\TextAt{9}{114.509}{334.599}{\Character{47}}%
\TextAt{14}{120.362}{334.599}{p}%
\TextAt{8}{130.149}{334.599}{are}%
\TextAt{8}{149.566}{334.599}{zero}%
\TextAt{8}{174.183}{334.599}{or}%
\TextAt{8}{188.399}{334.599}{not.)}%
\TextAt{8}{218.288}{334.599}{Using}%
\TextAt{8}{251.094}{334.599}{the}%
\TextAt{8}{271.157}{334.599}{Chinese}%
\TextAt{8}{314.687}{334.599}{remainder}%
\TextAt{8}{369.874}{334.599}{theorem,}%
\TextAt{8}{418.576}{334.599}{we}%
\TextAt{8}{435.709}{334.599}{now}%
\TextAt{8}{460.004}{334.599}{pick}%
\TextAt{8}{484.621}{334.599}{an}%
\TextAt{13}{500.786}{334.599}{\Character{11}}%
\NativeMathAt{511.629}{334.599}{11.955}{\in}%
\TextAt{16}{522.903}{334.599}{\Character{79}}%
\TextAt{8}{79.226}{317.553}{such}%
\TextAt{8}{105.152}{317.553}{that}%
\TextAt{16}{130.049}{317.553}{\Character{106}}%
\TextAt{13}{133.37}{317.553}{\Character{11}}%
\TextAt{16}{142.584}{317.553}{\Character{0}}%
\TextAt{13}{153.576}{317.553}{\Character{12}}%
\TextAt{16}{160.847}{317.553}{\Character{106}}%
\TextAt{8}{167.594}{317.553}{and}%
\TextAt{16}{189.885}{317.553}{\Character{106}}%
\TextAt{13}{193.206}{317.553}{\Character{27}}%
\TextAt{21}{199.859}{315.759}{i}%
\TextAt{13}{203.24}{317.553}{\Character{11}}%
\TextAt{16}{212.455}{317.553}{\Character{0}}%
\TextAt{13}{223.446}{317.553}{a}%
\TextAt{16}{229.591}{317.553}{\Character{106}}%
\TextAt{8}{232.912}{317.553}{,}%
\TextAt{8}{239.684}{317.553}{for}%
\TextAt{13}{257.097}{317.553}{i}%
\TextAt{9}{264.412}{317.553}{\Character{61}}%
\TextAt{9}{276.822}{317.553}{2}%
\TextAt{13}{282.69}{317.553}{\Character{59}}%
\TextAt{13}{287.914}{317.553}{\Character{58}}%
\TextAt{13}{293.151}{317.553}{\Character{58}}%
\TextAt{13}{298.375}{317.553}{\Character{58}}%
\TextAt{13}{303.612}{317.553}{\Character{59}}%
\TextAt{13}{308.848}{317.553}{r}%
\TextAt{8}{314.511}{317.553}{,}%
\TextAt{8}{321.283}{317.553}{are}%
\TextAt{8}{340.323}{317.553}{very}%
\TextAt{8}{365.543}{317.553}{small.}%
\TextAt{8}{400.566}{317.553}{We}%
\TextAt{8}{420.239}{317.553}{may}%
\TextAt{8}{445.136}{317.553}{even}%
\TextAt{8}{471.325}{317.553}{assume}%
\TextAt{8}{511.298}{317.553}{that}%
\TextAt{13}{79.226}{300.507}{L}%
\TextAt{9}{91.435}{300.507}{\Character{61}}%
\TextAt{13}{104.784}{300.507}{K}%
\TextAt{9}{115.602}{300.507}{\Character{40}}%
\TextAt{13}{120.154}{300.507}{\Character{11}}%
\TextAt{9}{127.676}{300.507}{\Character{41}}%
\TextAt{8}{136.668}{300.507}{(if}%
\TextAt{8}{152.488}{300.507}{not,}%
\TextAt{8}{177.232}{300.507}{modify}%
\TextAt{8}{217.118}{300.507}{by}%
\TextAt{13}{233.919}{300.507}{\Character{11}}%
\TextAt{9}{244.467}{300.507}{\Character{43}}%
\TextAt{13}{256.598}{300.507}{\Character{25}}%
\TextAt{21}{263.667}{304.845}{\Character{23}}%
\TextAt{13}{268.849}{300.507}{\Character{13}}%
\TextAt{8}{275.571}{300.507}{,}%
\TextAt{13}{283.398}{300.507}{\Character{25}}%
\NativeMathAt{294.712}{300.507}{11.955}{\in}%
\TextAt{14}{306.927}{300.507}{p}%
\TextAt{8}{312.91}{300.507}{,}%
\TextAt{8}{320.736}{300.507}{for}%
\TextAt{13}{339.163}{300.507}{\Character{23}}%
\TextAt{8}{350.141}{300.507}{big,}%
\TextAt{13}{373.581}{300.507}{\Character{13}}%
\NativeMathAt{384.548}{300.507}{11.955}{\in}%
\TextAt{16}{396.742}{300.507}{\Character{79}}%
\TextAt{8}{406.612}{300.507}{,}%
\TextAt{13}{414.439}{300.507}{L}%
\TextAt{9}{426.648}{300.507}{\Character{61}}%
\TextAt{13}{439.997}{300.507}{K}%
\TextAt{9}{450.814}{300.507}{\Character{40}}%
\TextAt{13}{455.367}{300.507}{\Character{13}}%
\TextAt{9}{462.089}{300.507}{\Character{41}}%
\TextAt{8}{466.641}{300.507}{;}%
\TextAt{8}{474.604}{300.507}{for}%
\TextAt{8}{493.031}{300.507}{suitable}%
\TextAt{13}{79.226}{283.461}{\Character{23}}%
\NativeMathAt{89.085}{283.461}{11.955}{\neq}%
\TextAt{13}{101.51}{283.461}{\Character{22}}%
\TextAt{8}{108.553}{283.461}{,}%
\TextAt{8}{115.572}{283.461}{one}%
\TextAt{8}{136.869}{283.461}{then}%
\TextAt{8}{163.366}{283.461}{finds}%
\TextAt{13}{191.225}{283.461}{K}%
\TextAt{9}{202.043}{283.461}{\Character{40}}%
\TextAt{13}{206.595}{283.461}{\Character{11}}%
\TextAt{9}{216.44}{283.461}{\Character{43}}%
\TextAt{13}{227.868}{283.461}{\Character{25}}%
\TextAt{21}{234.938}{287.799}{\Character{23}}%
\TextAt{13}{240.12}{283.461}{\Character{13}}%
\TextAt{9}{246.842}{283.461}{\Character{41}}%
\TextAt{9}{254.697}{283.461}{\Character{61}}%
\TextAt{13}{267.14}{283.461}{K}%
\TextAt{9}{277.958}{283.461}{\Character{40}}%
\TextAt{13}{282.51}{283.461}{\Character{11}}%
\TextAt{9}{292.355}{283.461}{\Character{43}}%
\TextAt{13}{303.783}{283.461}{\Character{25}}%
\TextAt{21}{310.853}{287.799}{\Character{22}}%
\TextAt{13}{316.445}{283.461}{\Character{13}}%
\TextAt{9}{323.167}{283.461}{\Character{41}}%
\TextAt{9}{331.022}{283.461}{\Character{61}}%
\TextAt{13}{343.466}{283.461}{K}%
\TextAt{9}{354.284}{283.461}{\Character{40}}%
\TextAt{13}{358.836}{283.461}{\Character{13}}%
\TextAt{9}{365.558}{283.461}{\Character{41}}%
\TextAt{8}{370.11}{283.461}{).}%
\TextAt{8}{383.059}{283.461}{Since}%
\TextAt{13}{413.454}{283.461}{\Character{11}}%
\TextAt{8}{424.71}{283.461}{is}%
\TextAt{8}{436.311}{283.461}{close}%
\TextAt{8}{464.17}{283.461}{to}%
\TextAt{13}{478.318}{283.461}{\Character{12}}%
\TextAt{8}{485.589}{283.461}{,}%
\TextAt{8}{492.608}{283.461}{we}%
\TextAt{8}{509.672}{283.461}{have}%
\TextAt{16}{79.226}{266.415}{\Character{79}}%
\TextAt{27}{88.744}{264.622}{P}%
\TextAt{9}{99.158}{266.415}{\Character{61}}%
\TextAt{14}{111.584}{266.415}{o}%
\TextAt{27}{117.428}{264.622}{p}%
\TextAt{9}{121.915}{266.415}{\Character{91}}%
\TextAt{13}{125.166}{266.415}{\Character{11}}%
\TextAt{9}{132.688}{266.415}{\Character{93}}%
\TextAt{8}{135.94}{266.415}{.}%
\TextAt{8}{144.388}{266.415}{Now}%
\TextAt{13}{242.239}{249.369}{\Character{14}}%
\TextAt{21}{247.424}{247.509}{L}%
\TextAt{28}{253.183}{246.118}{P}%
\TextAt{3}{259.38}{247.509}{\Character{106}}%
\TextAt{21}{261.732}{247.509}{K}%
\TextAt{28}{268.874}{246.513}{p}%
\TextAt{9}{273.454}{249.369}{\Character{40}}%
\TextAt{13}{278.007}{249.369}{\Character{11}}%
\TextAt{9}{285.528}{249.369}{\Character{41}}%
\TextAt{9}{293.383}{249.369}{\Character{61}}%
\TextAt{20}{305.939}{260.726}{\Character{89}}%
\TextAt{21}{305.827}{235.266}{\Character{28}}%
\NativeMathAt{310.505}{235.266}{7.97}{\neq}%
\TextAt{17}{317.088}{235.266}{1}%
\TextAt{9}{321.326}{249.369}{\Character{40}}%
\TextAt{13}{325.878}{249.369}{\Character{11}}%
\TextAt{16}{336.057}{249.369}{\Character{0}}%
\TextAt{13}{348.012}{249.369}{\Character{28}}%
\TextAt{13}{354.42}{249.369}{\Character{11}}%
\TextAt{9}{361.953}{249.369}{\Character{41}}%
\TextAt{13}{366.505}{249.369}{\Character{59}}%
\TextAt{8}{79.226}{216.866}{where}%
\TextAt{13}{113.035}{216.866}{\Character{28}}%
\TextAt{8}{123.352}{216.866}{runs}%
\TextAt{8}{149.426}{216.866}{through}%
\TextAt{8}{193.66}{216.866}{the}%
\TextAt{13}{213.816}{216.866}{K}%
\TextAt{27}{223.776}{215.073}{p}%
\TextAt{8}{228.262}{216.866}{-embeddings}%
\TextAt{13}{295.953}{216.866}{L}%
\TextAt{27}{303.917}{215.073}{P}%
\NativeMathAt{314.332}{216.866}{11.955}{\rightarrow}%
\Vector{q 0 G 0 g 0.478 w 257.608 -547.291 m 268.425 -547.291 l S Q}%
\TextAt{13}{329.608}{216.866}{K}%
\TextAt{27}{340.425}{215.073}{p}%
\TextAt{8}{348.809}{216.866}{different}%
\TextAt{8}{394.98}{216.866}{from}%
\TextAt{9}{422.621}{216.866}{1}%
\TextAt{8}{428.474}{216.866}{.}%
\TextAt{8}{436.922}{216.866}{Furthermore,}%
\TextAt{13}{160.49}{186.59}{\Character{14}}%
\TextAt{21}{165.674}{184.731}{L}%
\TextAt{3}{171.433}{184.731}{\Character{106}}%
\TextAt{21}{173.786}{184.731}{K}%
\TextAt{9}{181.988}{186.59}{\Character{40}}%
\TextAt{13}{186.541}{186.59}{\Character{11}}%
\TextAt{9}{194.062}{186.59}{\Character{41}}%
\TextAt{9}{201.917}{186.59}{\Character{61}}%
\TextAt{20}{214.696}{197.948}{\Character{89}}%
\TextAt{21}{214.361}{172.488}{\Character{27}}%
\NativeMathAt{219.486}{172.488}{7.97}{\neq}%
\TextAt{17}{226.069}{172.488}{1}%
\TextAt{9}{230.307}{186.59}{\Character{40}}%
\TextAt{13}{234.859}{186.59}{\Character{11}}%
\TextAt{16}{245.038}{186.59}{\Character{0}}%
\TextAt{13}{256.993}{186.59}{\Character{27}}%
\TextAt{13}{264.058}{186.59}{\Character{11}}%
\TextAt{9}{271.597}{186.59}{\Character{41}}%
\TextAt{9}{279.452}{186.59}{\Character{61}}%
\TextAt{20}{292.007}{197.948}{\Character{89}}%
\TextAt{21}{291.896}{172.488}{\Character{28}}%
\NativeMathAt{296.574}{172.488}{7.97}{\neq}%
\TextAt{17}{303.157}{172.488}{1}%
\TextAt{9}{307.395}{186.59}{\Character{40}}%
\TextAt{13}{311.947}{186.59}{\Character{11}}%
\TextAt{16}{322.125}{186.59}{\Character{0}}%
\TextAt{13}{334.081}{186.59}{\Character{28}}%
\TextAt{13}{340.489}{186.59}{\Character{11}}%
\TextAt{9}{348.021}{186.59}{\Character{41}}%
\TextAt{21}{360.176}{201.534}{r}%
\TextAt{20}{354.566}{197.948}{\Character{89}}%
\TextAt{21}{355.352}{172.753}{i}%
\TextAt{17}{358.235}{172.753}{\Character{61}2}%
\TextAt{20}{371.835}{197.948}{\Character{89}}%
\TextAt{21}{377.531}{172.753}{j}%
\TextAt{9}{387.111}{186.59}{\Character{40}}%
\TextAt{13}{391.663}{186.59}{\Character{11}}%
\TextAt{16}{401.841}{186.59}{\Character{0}}%
\TextAt{13}{413.797}{186.59}{\Character{28}}%
\TextAt{21}{418.88}{184.797}{ij}%
\TextAt{13}{426.146}{186.59}{\Character{27}}%
\TextAt{21}{432.799}{184.797}{i}%
\TextAt{13}{436.18}{186.59}{\Character{11}}%
\TextAt{9}{443.702}{186.59}{\Character{41}}%
\TextAt{13}{448.254}{186.59}{\Character{59}}%
\TextAt{8}{79.226}{153.922}{where}%
\TextAt{13}{112.822}{153.922}{\Character{27}}%
\TextAt{8}{123.589}{153.922}{runs}%
\TextAt{8}{149.45}{153.922}{over}%
\TextAt{8}{174.283}{153.922}{the}%
\TextAt{13}{194.226}{153.922}{K}%
\TextAt{8}{205.044}{153.922}{-embeddings}%
\TextAt{8}{272.521}{153.922}{different}%
\TextAt{8}{318.479}{153.922}{from}%
\TextAt{9}{345.906}{153.922}{1}%
\TextAt{8}{351.759}{153.922}{,}%
\TextAt{8}{358.738}{153.922}{and}%
\TextAt{8}{381.288}{153.922}{the}%
\TextAt{13}{401.231}{153.922}{\Character{28}}%
\TextAt{21}{406.315}{152.128}{ij}%
\TextAt{8}{417.264}{153.922}{are}%
\TextAt{8}{436.562}{153.922}{certain}%
\TextAt{8}{475.371}{153.922}{elements}%
\TextAt{8}{523.014}{153.922}{in}%
\TextAt{13}{79.226}{136.876}{G}%
\TextAt{27}{88.46}{135.082}{p}%
\TextAt{8}{92.947}{136.876}{.}%
\TextAt{8}{101.395}{136.876}{But}%
\TextAt{8}{124.636}{136.876}{now}%
\TextAt{16}{163.756}{113.421}{\Character{106}}%
\TextAt{13}{167.077}{113.421}{\Character{11}}%
\TextAt{16}{177.256}{113.421}{\Character{0}}%
\TextAt{13}{189.211}{113.421}{\Character{28}}%
\TextAt{21}{194.294}{111.628}{ij}%
\TextAt{13}{201.56}{113.421}{\Character{27}}%
\TextAt{21}{208.213}{111.628}{i}%
\TextAt{13}{211.594}{113.421}{\Character{11}}%
\TextAt{16}{219.116}{113.421}{\Character{106}}%
\TextAt{9}{225.758}{113.421}{\Character{61}}%
\TextAt{16}{238.183}{113.421}{\Character{106}}%
\TextAt{13}{241.504}{113.421}{\Character{28}}%
\TextAt{3}{247.923}{118.37}{\Character{0}}%
\TextAt{17}{254.51}{118.37}{1}%
\TextAt{21}{246.588}{110.358}{ij}%
\TextAt{13}{259.242}{113.421}{\Character{11}}%
\TextAt{16}{269.42}{113.421}{\Character{0}}%
\TextAt{13}{281.376}{113.421}{\Character{27}}%
\TextAt{21}{288.029}{111.628}{i}%
\TextAt{13}{291.41}{113.421}{\Character{11}}%
\TextAt{16}{298.932}{113.421}{\Character{106}}%
\TextAt{9}{305.574}{113.421}{\Character{61}}%
\TextAt{16}{317.999}{113.421}{\Character{106}}%
\TextAt{13}{321.32}{113.421}{\Character{28}}%
\TextAt{3}{327.739}{118.37}{\Character{0}}%
\TextAt{17}{334.326}{118.37}{1}%
\TextAt{21}{326.404}{110.358}{ij}%
\TextAt{13}{339.058}{113.421}{\Character{11}}%
\TextAt{16}{349.236}{113.421}{\Character{0}}%
\TextAt{13}{361.191}{113.421}{a}%
\TextAt{9}{369.993}{113.421}{\Character{43}}%
\TextAt{13}{381.754}{113.421}{a}%
\TextAt{16}{390.556}{113.421}{\Character{0}}%
\TextAt{13}{402.511}{113.421}{\Character{27}}%
\TextAt{21}{409.165}{111.628}{i}%
\TextAt{13}{412.546}{113.421}{\Character{11}}%
\TextAt{16}{420.068}{113.421}{\Character{106}}%
\TextAt{9}{426.709}{113.421}{\Character{61}}%
\TextAt{9}{439.119}{113.421}{1}%
\TextAt{13}{444.988}{113.421}{\Character{59}}%
\TextAt{8}{79.226}{89.967}{since}%
\TextAt{16}{110.366}{89.967}{\Character{106}}%
\TextAt{13}{113.687}{89.967}{a}%
\TextAt{16}{124.173}{89.967}{\Character{0}}%
\TextAt{13}{137.813}{89.967}{\Character{27}}%
\TextAt{21}{144.467}{88.174}{i}%
\TextAt{13}{147.848}{89.967}{\Character{11}}%
\TextAt{16}{155.37}{89.967}{\Character{106}}%
\TextAt{8}{165.06}{89.967}{is}%
\TextAt{8}{179.295}{89.967}{very}%
\TextAt{8}{207.458}{89.967}{small,}%
\TextAt{8}{244.429}{89.967}{and}%
\TextAt{13}{269.663}{89.967}{\Character{28}}%
\TextAt{3}{276.082}{94.915}{\Character{0}}%
\TextAt{17}{282.668}{94.915}{1}%
\TextAt{21}{274.747}{86.904}{ij}%
\TextAt{13}{287.401}{89.967}{\Character{11}}%
\TextAt{8}{301.291}{89.967}{is}%
\TextAt{8}{315.527}{89.967}{very}%
\TextAt{8}{343.69}{89.967}{close}%
\TextAt{8}{374.184}{89.967}{to}%
\TextAt{13}{390.966}{89.967}{\Character{28}}%
\TextAt{3}{397.385}{94.915}{\Character{0}}%
\TextAt{17}{403.972}{94.915}{1}%
\TextAt{21}{396.05}{86.904}{ij}%
\TextAt{13}{408.704}{89.967}{\Character{12}}%
\TextAt{8}{422.344}{89.967}{(see}%
\TextAt{8}{448.284}{89.967}{(*)).}%
\TextAt{8}{483.67}{89.967}{Therefore}%
\TextAt{13}{79.226}{72.921}{v}%
\TextAt{27}{84.885}{71.128}{P}%
\TextAt{9}{91.979}{72.921}{\Character{40}}%
\TextAt{13}{96.531}{72.921}{\Character{14}}%
\TextAt{21}{101.716}{71.062}{L}%
\TextAt{3}{107.475}{71.062}{\Character{106}}%
\TextAt{21}{109.827}{71.062}{K}%
\TextAt{9}{118.03}{72.921}{\Character{40}}%
\TextAt{13}{122.582}{72.921}{\Character{11}}%
\TextAt{9}{130.104}{72.921}{\Character{41}\Character{41}}%
\TextAt{9}{142.502}{72.921}{\Character{61}}%
\TextAt{13}{154.955}{72.921}{v}%
\TextAt{27}{160.614}{71.128}{P}%
\TextAt{9}{167.708}{72.921}{\Character{40}}%
\TextAt{20}{172.26}{81.888}{\Character{81}}%
\TextAt{21}{183.551}{69.434}{\Character{28}}%
\NativeMathAt{188.23}{69.434}{7.97}{\neq}%
\TextAt{17}{194.813}{69.434}{1}%
\TextAt{9}{199.548}{72.921}{\Character{40}}%
\TextAt{13}{204.101}{72.921}{\Character{11}}%
\TextAt{16}{214.279}{72.921}{\Character{0}}%
\TextAt{13}{226.234}{72.921}{\Character{28}}%
\TextAt{13}{232.642}{72.921}{\Character{11}}%
\TextAt{9}{240.175}{72.921}{\Character{41}\Character{41}}%
\TextAt{9}{252.573}{72.921}{\Character{61}}%
\TextAt{13}{265.026}{72.921}{v}%
\TextAt{27}{270.685}{71.128}{P}%
\TextAt{9}{277.779}{72.921}{\Character{40}}%
\TextAt{13}{282.331}{72.921}{\Character{14}}%
\TextAt{21}{287.516}{71.062}{L}%
\TextAt{28}{293.275}{69.671}{P}%
\TextAt{3}{299.472}{71.062}{\Character{106}}%
\TextAt{21}{301.824}{71.062}{K}%
\TextAt{28}{308.966}{70.065}{p}%
\TextAt{9}{313.547}{72.921}{\Character{40}}%
\TextAt{13}{318.099}{72.921}{\Character{11}}%
\TextAt{9}{325.621}{72.921}{\Character{41}\Character{41}}%
\TextAt{8}{334.725}{72.921}{,}%
\TextAt{8}{341.874}{72.921}{as}%
\TextAt{8}{356.245}{72.921}{required.}%
\Vector{q 0 G 0 g 0.399 w 452.334 -701.079 m 452.334 -693.009 l S Q}%
\Vector{q 0 G 0 g 0.399 w 452.534 -693.208 m 459.707 -693.208 l S Q}%
\Vector{q 0 G 0 g 0.399 w 452.534 -700.879 m 459.707 -700.879 l S Q}%
\Vector{q 0 G 0 g 0.399 w 459.906 -701.079 m 459.906 -693.009 l S Q}%
\TextAt{8}{297.211}{43.033}{181}%
\Destination{522.36}{392.145}{nav.AMS.22}%
\Destination{78.23}{809.815}{nav.page.181}%
\Destination{79.36}{366.655}{nav.print.199}%
\Hotlink{201.226}{724.855}{5.22}{11.047}{nav.chapter.1}%
\Hotlink{218.275}{724.855}{12.716}{11.047}{nav.section.1.12}%
\Hotlink{246.895}{724.855}{12.716}{11.047}{nav.print.79}%
\Hotlink{147.624}{549.655}{9.751}{11.047}{nav.chapter.2}%
\Hotlink{163.639}{549.655}{30.936}{11.047}{nav.number.2.10.4}%
\Hotlink{413.195}{515.555}{25.078}{11.047}{nav.number.3.2.2}%
\Hotlink{492.322}{515.555}{25.078}{11.047}{nav.number.3.2.4}%
}

% Page 194 (body page 182)
\Page{%
\TextAt{2}{79.226}{796.943}{III.}%
\TextAt{2}{98.317}{796.943}{Riemann–Roch}%
\TextAt{2}{174.375}{796.943}{Theory}%
\Vector{q 0 G 0 g 0.299 w 7.226 18.729 m 460.769 18.729 l S Q}%
\TextAt{8}{97.159}{761.124}{The}%
\TextAt{8}{119.984}{761.124}{different}%
\TextAt{8}{164.927}{761.124}{characterizes}%
\TextAt{8}{233.026}{761.124}{the}%
\TextAt{8}{251.954}{761.124}{ramification}%
\TextAt{8}{315.773}{761.124}{behaviour}%
\TextAt{8}{368.857}{761.124}{of}%
\TextAt{8}{380.959}{761.124}{the}%
\TextAt{8}{399.887}{761.124}{extension}%
\TextAt{13}{450.424}{761.124}{L}%
\TextAt{16}{461.709}{761.124}{\Character{106}}%
\TextAt{13}{468.351}{761.124}{K}%
\TextAt{8}{481.838}{761.124}{as}%
\TextAt{8}{494.979}{761.124}{follows.}%
\TextAt{12}{79.226}{735.111}{Theorem}%
\TextAt{12}{136.24}{735.111}{(2.6).}%
\TextAt{8}{171.556}{735.111}{A}%
\TextAt{8}{184.228}{735.111}{prime}%
\TextAt{8}{217.392}{735.111}{ideal}%
\TextAt{14}{245.355}{735.111}{P}%
\TextAt{8}{259.146}{735.111}{of}%
\TextAt{13}{272.476}{735.111}{L}%
\TextAt{8}{284.338}{735.111}{is}%
\TextAt{8}{296.102}{735.111}{ramified}%
\TextAt{8}{341.627}{735.111}{over}%
\TextAt{13}{366.673}{735.111}{K}%
\TextAt{8}{381.388}{735.111}{if}%
\TextAt{8}{392.112}{735.111}{and}%
\TextAt{8}{414.874}{735.111}{only}%
\TextAt{8}{440.566}{735.111}{if}%
\TextAt{14}{451.29}{735.111}{P}%
\TextAt{16}{464.504}{735.111}{\Character{106}}%
\TextAt{14}{471.146}{735.111}{D}%
\TextAt{21}{481.087}{733.252}{L}%
\TextAt{3}{486.846}{733.252}{\Character{106}}%
\TextAt{21}{489.199}{733.252}{K}%
\TextAt{8}{497.401}{735.111}{.}%
\TextAt{8}{97.159}{715.077}{Let}%
\TextAt{14}{117.806}{715.077}{P}%
\TextAt{21}{127.7}{719.415}{s}%
\TextAt{8}{135.689}{715.077}{be}%
\TextAt{8}{151.292}{715.077}{the}%
\TextAt{8}{171.127}{715.077}{maximal}%
\TextAt{8}{218.613}{715.077}{power}%
\TextAt{8}{252.436}{715.077}{of}%
\TextAt{14}{265.444}{715.077}{P}%
\TextAt{8}{278.913}{715.077}{dividing}%
\TextAt{14}{323.793}{715.077}{D}%
\TextAt{21}{333.735}{713.217}{L}%
\TextAt{3}{339.494}{713.217}{\Character{106}}%
\TextAt{21}{341.847}{713.217}{K}%
\TextAt{8}{350.049}{715.077}{,}%
\TextAt{8}{356.941}{715.077}{and}%
\TextAt{8}{379.382}{715.077}{let}%
\TextAt{13}{395.965}{715.077}{e}%
\TextAt{8}{404.966}{715.077}{be}%
\TextAt{8}{420.568}{715.077}{the}%
\TextAt{8}{440.403}{715.077}{ramification}%
\TextAt{8}{505.129}{715.077}{index}%
\TextAt{8}{79.226}{698.031}{of}%
\TextAt{14}{92.556}{698.031}{P}%
\TextAt{8}{106.347}{698.031}{over}%
\TextAt{13}{131.393}{698.031}{K}%
\TextAt{8}{142.211}{698.031}{.}%
\TextAt{8}{150.659}{698.031}{Then}%
\TextAt{8}{181.216}{698.031}{one}%
\TextAt{8}{202.676}{698.031}{has}%
\TextAt{13}{213.895}{669.029}{s}%
\TextAt{9}{222.729}{669.029}{\Character{61}}%
\TextAt{13}{235.155}{669.029}{e}%
\TextAt{16}{243.237}{669.029}{\Character{0}}%
\TextAt{9}{255.192}{669.029}{1}%
\TextAt{13}{261.045}{669.029}{\Character{59}}%
\TextAt{8}{278.244}{669.029}{if}%
\TextAt{14}{288.968}{669.029}{P}%
\TextAt{8}{302.759}{669.029}{is}%
\TextAt{8}{314.523}{669.029}{tamely}%
\TextAt{8}{353.222}{669.029}{ramified}%
\TextAt{13}{394.849}{669.029}{\Character{59}}%
\TextAt{13}{184.219}{640.028}{e}%
\NativeMathAt{192.965}{640.028}{11.955}{\leq}%
\TextAt{13}{205.585}{640.028}{s}%
\NativeMathAt{214.419}{640.028}{11.955}{\leq}%
\TextAt{13}{227.039}{640.028}{e}%
\TextAt{16}{235.121}{640.028}{\Character{0}}%
\TextAt{9}{247.076}{640.028}{1}%
\TextAt{9}{255.576}{640.028}{\Character{43}}%
\TextAt{13}{267.347}{640.028}{v}%
\TextAt{27}{273.006}{638.235}{P}%
\TextAt{9}{280.1}{640.028}{\Character{40}}%
\TextAt{13}{284.652}{640.028}{e}%
\TextAt{9}{290.078}{640.028}{\Character{41}}%
\TextAt{13}{294.63}{640.028}{\Character{59}}%
\TextAt{8}{311.829}{640.028}{if}%
\TextAt{14}{322.553}{640.028}{P}%
\TextAt{8}{336.344}{640.028}{is}%
\TextAt{8}{348.108}{640.028}{wildly}%
\TextAt{8}{382.897}{640.028}{ramified}%
\TextAt{13}{424.525}{640.028}{\Character{58}}%
\TextAt{19}{79.226}{598.574}{Proof.}%
\TextAt{8}{116}{598.574}{By}%
\TextAt{8}{135.135}{598.574}{(2.2),}%
\TextAt{8}{167.327}{598.574}{(iii),}%
\TextAt{8}{194.306}{598.574}{we}%
\TextAt{8}{212.305}{598.574}{may}%
\TextAt{8}{238.446}{598.574}{assume}%
\TextAt{8}{279.662}{598.574}{that}%
\TextAt{14}{305.803}{598.574}{o}%
\TextAt{8}{316.317}{598.574}{is}%
\TextAt{8}{328.852}{598.574}{a}%
\TextAt{8}{339.38}{598.574}{complete}%
\TextAt{8}{389.574}{598.574}{discrete}%
\TextAt{8}{433.325}{598.574}{valuation}%
\TextAt{8}{485.169}{598.574}{ring}%
\TextAt{8}{510.007}{598.574}{with}%
\TextAt{8}{79.226}{581.528}{maximal}%
\TextAt{8}{127.237}{581.528}{ideal}%
\TextAt{14}{155.402}{581.528}{p}%
\TextAt{8}{161.385}{581.528}{.}%
\TextAt{8}{170.439}{581.528}{Then,}%
\TextAt{8}{204.501}{581.528}{by}%
\TextAt{8}{220.962}{581.528}{chap.}%
\TextAt{8}{253.759}{581.528}{II,}%
\TextAt{8}{269.873}{581.528}{(10.4),}%
\TextAt{8}{307.21}{581.528}{we}%
\TextAt{8}{324.64}{581.528}{have}%
\TextAt{16}{351.836}{581.528}{\Character{79}}%
\TextAt{9}{365.351}{581.528}{\Character{61}}%
\TextAt{14}{378.121}{581.528}{o}%
\TextAt{9}{383.965}{581.528}{\Character{91}}%
\TextAt{13}{387.217}{581.528}{\Character{11}}%
\TextAt{9}{394.738}{581.528}{\Character{93}}%
\TextAt{8}{402.089}{581.528}{for}%
\TextAt{8}{420.176}{581.528}{a}%
\TextAt{8}{430.134}{581.528}{suitable}%
\TextAt{13}{473.972}{581.528}{\Character{11}}%
\NativeMathAt{485.159}{581.528}{11.955}{\in}%
\TextAt{16}{496.78}{581.528}{\Character{79}}%
\TextAt{8}{506.643}{581.528}{.}%
\TextAt{8}{515.698}{581.528}{Let}%
\TextAt{13}{79.226}{564.482}{f}%
\TextAt{9}{86.273}{564.482}{\Character{40}}%
\TextAt{13}{90.825}{564.482}{X}%
\TextAt{9}{101.48}{564.482}{\Character{41}}%
\TextAt{8}{110.943}{564.482}{be}%
\TextAt{8}{127.88}{564.482}{the}%
\TextAt{8}{149.049}{564.482}{minimal}%
\TextAt{8}{195.588}{564.482}{polynomial}%
\TextAt{8}{257.094}{564.482}{of}%
\TextAt{13}{271.437}{564.482}{\Character{11}}%
\TextAt{8}{278.959}{564.482}{.}%
\TextAt{8}{290.446}{564.482}{(2.4)}%
\TextAt{8}{319.434}{564.482}{says}%
\TextAt{8}{345.29}{564.482}{that}%
\TextAt{13}{371.672}{564.482}{s}%
\TextAt{9}{382.233}{564.482}{\Character{61}}%
\TextAt{13}{396.385}{564.482}{v}%
\TextAt{27}{402.044}{562.689}{P}%
\TextAt{9}{409.138}{564.482}{\Character{40}}%
\TextAt{13}{413.69}{564.482}{f}%
\TextAt{3}{420.737}{568.82}{\Character{48}}%
\TextAt{9}{423.532}{564.482}{\Character{40}}%
\TextAt{13}{428.084}{564.482}{\Character{11}}%
\TextAt{9}{435.606}{564.482}{\Character{41}\Character{41}}%
\TextAt{8}{444.71}{564.482}{.}%
\TextAt{8}{456.198}{564.482}{Assume}%
\TextAt{13}{500.572}{564.482}{L}%
\TextAt{16}{513.584}{564.482}{\Character{106}}%
\TextAt{13}{521.952}{564.482}{K}%
\TextAt{8}{79.226}{547.436}{is}%
\TextAt{8}{92.568}{547.436}{unramified.}%
\TextAt{8}{160.387}{547.436}{Then}%
\Vector{q 0 G 0 g 0.478 w 120.523 -219.743 m 128.044 -219.743 l S Q}%
\TextAt{13}{192.523}{547.436}{\Character{11}}%
\TextAt{9}{206.055}{547.436}{\Character{61}}%
\TextAt{13}{221.17}{547.436}{\Character{11}}%
\TextAt{8}{232.013}{547.436}{mod}%
\TextAt{14}{257.451}{547.436}{P}%
\TextAt{8}{272.82}{547.436}{is}%
\TextAt{8}{286.162}{547.436}{a}%
\TextAt{8}{297.496}{547.436}{simple}%
\TextAt{8}{335.549}{547.436}{zero}%
\TextAt{8}{361.839}{547.436}{of}%
\Vector{q 0 G 0 g 0.478 w 304.747 -216.588 m 311.794 -216.588 l S Q}%
\TextAt{13}{376.747}{547.436}{f}%
\TextAt{9}{383.794}{547.436}{\Character{40}}%
\TextAt{13}{388.346}{547.436}{X}%
\TextAt{9}{399.001}{547.436}{\Character{41}}%
\TextAt{9}{409.546}{547.436}{\Character{61}}%
\TextAt{13}{424.679}{547.436}{f}%
\TextAt{9}{431.726}{547.436}{\Character{40}}%
\TextAt{13}{436.278}{547.436}{X}%
\TextAt{9}{446.933}{547.436}{\Character{41}}%
\TextAt{8}{454.806}{547.436}{mod}%
\TextAt{14}{480.244}{547.436}{p}%
\TextAt{8}{486.227}{547.436}{,}%
\TextAt{8}{495.35}{547.436}{so}%
\TextAt{8}{511.298}{547.436}{that}%
\TextAt{13}{79.226}{530.39}{f}%
\TextAt{3}{86.273}{534.728}{\Character{48}}%
\TextAt{9}{89.068}{530.39}{\Character{40}}%
\TextAt{13}{93.62}{530.39}{\Character{11}}%
\TextAt{9}{101.142}{530.39}{\Character{41}}%
\NativeMathAt{109.015}{530.39}{11.955}{\in}%
\TextAt{16}{120.289}{530.39}{\Character{79}}%
\TextAt{3}{130.156}{534.728}{\Character{3}}%
\TextAt{8}{138.785}{530.39}{and}%
\TextAt{8}{161.548}{530.39}{thus}%
\TextAt{13}{187.299}{530.39}{s}%
\TextAt{9}{196.134}{530.39}{\Character{61}}%
\TextAt{9}{208.544}{530.39}{0}%
\TextAt{9}{217.713}{530.39}{\Character{61}}%
\TextAt{13}{230.159}{530.39}{e}%
\TextAt{16}{238.241}{530.39}{\Character{0}}%
\TextAt{9}{250.196}{530.39}{1}%
\TextAt{8}{256.049}{530.39}{.}%
\TextAt{8}{97.159}{510.355}{By}%
\TextAt{8}{114.554}{510.355}{(2.2),}%
\TextAt{8}{145.007}{510.355}{(i)}%
\TextAt{8}{160.299}{510.355}{and}%
\TextAt{8}{182.094}{510.355}{chap.}%
\TextAt{8}{213.963}{510.355}{II,}%
\TextAt{8}{228.907}{510.355}{(7.5),}%
\TextAt{8}{259.36}{510.355}{we}%
\TextAt{8}{275.62}{510.355}{may}%
\TextAt{8}{300.021}{510.355}{now}%
\TextAt{8}{323.442}{510.355}{pass}%
\TextAt{8}{347.963}{510.355}{to}%
\TextAt{8}{361.306}{510.355}{the}%
\TextAt{8}{380.495}{510.355}{maximal}%
\TextAt{8}{427.336}{510.355}{unramified}%
\TextAt{8}{484.901}{510.355}{extension}%
\TextAt{8}{79.226}{493.309}{and}%
\TextAt{8}{101.979}{493.309}{assume}%
\TextAt{8}{142.414}{493.309}{that}%
\TextAt{13}{167.773}{493.309}{L}%
\TextAt{16}{179.059}{493.309}{\Character{106}}%
\TextAt{13}{185.701}{493.309}{K}%
\TextAt{8}{200.406}{493.309}{is}%
\TextAt{8}{212.16}{493.309}{totally}%
\TextAt{8}{249.559}{493.309}{ramified.}%
\TextAt{8}{299.632}{493.309}{Then}%
\TextAt{13}{330.18}{493.309}{\Character{11}}%
\TextAt{8}{341.589}{493.309}{may}%
\TextAt{8}{366.949}{493.309}{be}%
\TextAt{8}{382.863}{493.309}{chosen}%
\TextAt{8}{420.309}{493.309}{to}%
\TextAt{8}{434.61}{493.309}{be}%
\TextAt{8}{450.525}{493.309}{a}%
\TextAt{8}{460.271}{493.309}{prime}%
\TextAt{8}{493.425}{493.309}{element}%
\TextAt{8}{79.226}{476.263}{of}%
\TextAt{16}{92.556}{476.263}{\Character{79}}%
\TextAt{8}{102.406}{476.263}{.}%
\TextAt{8}{110.854}{476.263}{In}%
\TextAt{8}{125.475}{476.263}{this}%
\TextAt{8}{148.298}{476.263}{case}%
\TextAt{8}{173.069}{476.263}{the}%
\TextAt{8}{193.225}{476.263}{minimal}%
\TextAt{8}{238.751}{476.263}{polynomial}%
\TextAt{13}{187.024}{447.262}{f}%
\TextAt{9}{194.071}{447.262}{\Character{40}}%
\TextAt{13}{198.623}{447.262}{X}%
\TextAt{9}{209.278}{447.262}{\Character{41}}%
\TextAt{9}{217.133}{447.262}{\Character{61}}%
\TextAt{13}{229.577}{447.262}{a}%
\TextAt{17}{235.722}{445.469}{0}%
\TextAt{13}{240.454}{447.262}{X}%
\TextAt{21}{251.109}{452.198}{e}%
\TextAt{9}{258.192}{447.262}{\Character{43}}%
\TextAt{13}{269.954}{447.262}{a}%
\TextAt{17}{276.099}{445.469}{1}%
\TextAt{13}{280.831}{447.262}{X}%
\TextAt{21}{291.486}{452.198}{e}%
\TextAt{3}{295.415}{452.198}{\Character{0}}%
\TextAt{17}{302.001}{452.198}{1}%
\TextAt{9}{309.39}{447.262}{\Character{43}}%
\TextAt{16}{321.151}{447.262}{\Character{1}}%
\TextAt{16}{326.447}{447.262}{\Character{1}}%
\TextAt{16}{331.755}{447.262}{\Character{1}}%
\TextAt{9}{337.756}{447.262}{\Character{43}}%
\TextAt{13}{349.517}{447.262}{a}%
\TextAt{21}{355.662}{445.469}{e}%
\TextAt{13}{360.088}{447.262}{\Character{59}}%
\TextAt{13}{389.223}{447.262}{a}%
\TextAt{17}{395.388}{445.469}{0}%
\TextAt{9}{403.441}{447.262}{\Character{61}}%
\TextAt{9}{415.851}{447.262}{1}%
\TextAt{13}{421.72}{447.262}{\Character{59}}%
\TextAt{8}{79.226}{418.261}{is}%
\TextAt{8}{90.99}{418.261}{an}%
\TextAt{8}{107.249}{418.261}{Eisenstein}%
\TextAt{8}{162.805}{418.261}{polynomial.}%
\TextAt{8}{227.849}{418.261}{Let}%
\TextAt{8}{248.818}{418.261}{us}%
\TextAt{8}{263.834}{418.261}{look}%
\TextAt{8}{289.203}{418.261}{at}%
\TextAt{8}{303.513}{418.261}{the}%
\TextAt{8}{323.67}{418.261}{derivative}%
\TextAt{13}{190.112}{389.26}{f}%
\TextAt{3}{197.158}{394.196}{\Character{48}}%
\TextAt{9}{199.953}{389.26}{\Character{40}}%
\TextAt{13}{204.505}{389.26}{\Character{11}}%
\TextAt{9}{212.027}{389.26}{\Character{41}}%
\TextAt{9}{219.882}{389.26}{\Character{61}}%
\TextAt{13}{232.326}{389.26}{ea}%
\TextAt{17}{243.896}{387.467}{0}%
\TextAt{13}{248.629}{389.26}{\Character{11}}%
\TextAt{21}{256.15}{394.196}{e}%
\TextAt{3}{260.079}{394.196}{\Character{0}}%
\TextAt{17}{266.665}{394.196}{1}%
\TextAt{9}{274.054}{389.26}{\Character{43}}%
\TextAt{9}{285.806}{389.26}{\Character{40}}%
\TextAt{13}{290.368}{389.26}{e}%
\TextAt{16}{298.45}{389.26}{\Character{0}}%
\TextAt{9}{310.405}{389.26}{1\Character{41}}%
\TextAt{13}{320.81}{389.26}{a}%
\TextAt{17}{326.955}{387.467}{1}%
\TextAt{13}{331.688}{389.26}{\Character{11}}%
\TextAt{21}{339.209}{394.196}{e}%
\TextAt{3}{343.138}{394.196}{\Character{0}}%
\TextAt{17}{349.724}{394.196}{2}%
\TextAt{9}{357.113}{389.26}{\Character{43}}%
\TextAt{16}{368.874}{389.26}{\Character{1}}%
\TextAt{16}{374.17}{389.26}{\Character{1}}%
\TextAt{16}{379.478}{389.26}{\Character{1}}%
\TextAt{9}{385.479}{389.26}{\Character{43}}%
\TextAt{13}{397.24}{389.26}{a}%
\TextAt{21}{403.385}{387.467}{e}%
\TextAt{3}{407.313}{387.467}{\Character{0}}%
\TextAt{17}{413.9}{387.467}{1}%
\TextAt{13}{418.632}{389.26}{\Character{58}}%
\ErrataAt{79.19}{360.259}{455.81}{For $0\le i<e$ with $(e-i)a_i\ne0$ we find}%
\TextAt{13}{118.858}{331.258}{v}%
\TextAt{27}{124.517}{329.464}{P}%
\TextAt{9}{131.611}{331.258}{\Character{40}\Character{40}}%
\TextAt{13}{140.715}{331.258}{e}%
\TextAt{16}{148.797}{331.258}{\Character{0}}%
\TextAt{13}{160.752}{331.258}{i}%
\TextAt{9}{164.746}{331.258}{\Character{41}}%
\TextAt{13}{169.298}{331.258}{a}%
\TextAt{21}{175.443}{329.464}{i}%
\TextAt{13}{178.824}{331.258}{\Character{11}}%
\TextAt{21}{186.346}{336.194}{e}%
\TextAt{3}{190.275}{336.194}{\Character{0}}%
\TextAt{21}{196.861}{336.194}{i}%
\TextAt{3}{199.744}{336.194}{\Character{0}}%
\TextAt{17}{206.331}{336.194}{1}%
\TextAt{9}{211.063}{331.258}{\Character{41}}%
\TextAt{9}{218.918}{331.258}{\Character{61}}%
\TextAt{13}{231.362}{331.258}{ev}%
\TextAt{27}{242.446}{329.464}{p}%
\TextAt{9}{246.933}{331.258}{\Character{40}}%
\TextAt{13}{251.485}{331.258}{e}%
\TextAt{16}{259.568}{331.258}{\Character{0}}%
\TextAt{13}{271.523}{331.258}{i}%
\TextAt{9}{275.516}{331.258}{\Character{41}}%
\TextAt{9}{282.713}{331.258}{\Character{43}}%
\TextAt{13}{294.486}{331.258}{ev}%
\TextAt{27}{305.571}{329.464}{p}%
\TextAt{9}{310.058}{331.258}{\Character{40}}%
\TextAt{13}{314.61}{331.258}{a}%
\TextAt{21}{320.755}{329.464}{i}%
\TextAt{9}{324.137}{331.258}{\Character{41}}%
\TextAt{9}{331.334}{331.258}{\Character{43}}%
\TextAt{13}{343.107}{331.258}{e}%
\TextAt{16}{351.189}{331.258}{\Character{0}}%
\TextAt{13}{363.144}{331.258}{i}%
\TextAt{16}{369.794}{331.258}{\Character{0}}%
\TextAt{9}{381.749}{331.258}{1}%
\NativeMathAt{390.923}{331.258}{11.955}{\equiv}%
\TextAt{16}{403.524}{331.258}{\Character{0}}%
\TextAt{13}{412.841}{331.258}{i}%
\TextAt{16}{419.491}{331.258}{\Character{0}}%
\TextAt{9}{431.446}{331.258}{1}%
\TextAt{9}{449.247}{331.258}{\Character{40}}%
\TextAt{8}{453.807}{331.258}{mod}%
\TextAt{13}{479.909}{331.258}{e}%
\TextAt{9}{485.334}{331.258}{\Character{41}}%
\TextAt{13}{489.886}{331.258}{\Character{59}}%
\ErrataAt{79.19}{302.256}{455.81}{so that the nonzero terms of $f^{\prime}(\alpha)$ have distinct valuations. Therefore}%
\TextAt{13}{195.24}{273.255}{s}%
\TextAt{9}{204.075}{273.255}{\Character{61}}%
\TextAt{13}{216.5}{273.255}{v}%
\TextAt{27}{222.16}{271.462}{P}%
\TextAt{9}{229.253}{273.255}{\Character{40}}%
\TextAt{13}{233.806}{273.255}{f}%
\TextAt{3}{240.852}{278.192}{\Character{48}}%
\TextAt{9}{243.647}{273.255}{\Character{40}}%
\TextAt{13}{248.2}{273.255}{\Character{11}}%
\TextAt{9}{255.721}{273.255}{\Character{41}\Character{41}}%
\TextAt{9}{268.119}{273.255}{\Character{61}}%
\TextAt{8}{282.926}{273.255}{min}%
\TextAt{17}{280.572}{265.934}{0}%
\NativeMathAt{284.806}{265.934}{7.97}{\leq}%
\TextAt{21}{291.393}{265.934}{i\Character{60}e}%
\TextAt{16}{304.791}{273.255}{\Character{102}}%
\TextAt{13}{310.769}{273.255}{v}%
\TextAt{27}{316.428}{271.462}{P}%
\TextAt{9}{323.521}{273.255}{\Character{40}\Character{40}}%
\TextAt{13}{332.626}{273.255}{e}%
\TextAt{16}{340.708}{273.255}{\Character{0}}%
\TextAt{13}{352.663}{273.255}{i}%
\TextAt{9}{356.657}{273.255}{\Character{41}}%
\TextAt{13}{361.209}{273.255}{a}%
\TextAt{21}{367.354}{271.462}{i}%
\TextAt{13}{370.735}{273.255}{\Character{11}}%
\TextAt{21}{378.257}{278.192}{e}%
\TextAt{3}{382.185}{278.191}{\Character{0}}%
\TextAt{21}{388.772}{278.192}{i}%
\TextAt{3}{391.655}{278.191}{\Character{0}}%
\TextAt{17}{398.242}{278.192}{1}%
\TextAt{9}{402.974}{273.255}{\Character{41}}%
\TextAt{16}{407.526}{273.255}{\Character{103}}%
\TextAt{13}{413.504}{273.255}{\Character{58}}%
\TextAt{8}{79.226}{241.477}{If}%
\TextAt{8}{91.364}{241.477}{now}%
\TextAt{13}{116.199}{241.477}{L}%
\TextAt{16}{128.245}{241.477}{\Character{106}}%
\TextAt{13}{135.647}{241.477}{K}%
\TextAt{8}{150.808}{241.477}{is}%
\TextAt{8}{163.018}{241.477}{tamely}%
\TextAt{8}{202.163}{241.477}{ramified,}%
\TextAt{8}{251.498}{241.477}{i.e.,}%
\TextAt{8}{274.161}{241.477}{if}%
\TextAt{13}{285.331}{241.477}{v}%
\TextAt{27}{290.99}{239.683}{p}%
\TextAt{9}{295.477}{241.477}{\Character{40}}%
\TextAt{13}{300.029}{241.477}{e}%
\TextAt{9}{305.455}{241.477}{\Character{41}}%
\TextAt{9}{314.075}{241.477}{\Character{61}}%
\TextAt{9}{327.249}{241.477}{0}%
\TextAt{8}{333.127}{241.477}{,}%
\TextAt{8}{340.834}{241.477}{then}%
\TextAt{8}{367.94}{241.477}{the}%
\TextAt{8}{388.543}{241.477}{minimum}%
\TextAt{8}{441.341}{241.477}{is}%
\TextAt{8}{453.551}{241.477}{obviously}%
\TextAt{8}{505.775}{241.477}{equal}%
\TextAt{8}{79.226}{224.431}{to}%
\TextAt{13}{93.536}{224.431}{e}%
\TextAt{16}{101.619}{224.431}{\Character{0}}%
\TextAt{9}{113.574}{224.431}{1}%
\TextAt{8}{119.427}{224.431}{,}%
\TextAt{8}{126.576}{224.431}{and}%
\TextAt{8}{149.339}{224.431}{for}%
\TextAt{13}{167.224}{224.431}{v}%
\TextAt{27}{172.883}{222.637}{p}%
\TextAt{9}{177.37}{224.431}{\Character{40}}%
\TextAt{13}{181.922}{224.431}{e}%
\TextAt{9}{187.347}{224.431}{\Character{41}}%
\NativeMathAt{195.221}{224.431}{11.955}{\geq}%
\TextAt{9}{207.84}{224.431}{1}%
\TextAt{8}{213.693}{224.431}{,}%
\TextAt{8}{220.842}{224.431}{we}%
\TextAt{8}{238.069}{224.431}{deduce}%
\TextAt{8}{277.079}{224.431}{that}%
\TextAt{13}{302.448}{224.431}{e}%
\NativeMathAt{311.194}{224.431}{11.955}{\leq}%
\TextAt{13}{323.814}{224.431}{s}%
\NativeMathAt{332.649}{224.431}{11.955}{\leq}%
\TextAt{13}{345.268}{224.431}{v}%
\TextAt{27}{350.927}{222.637}{P}%
\TextAt{9}{358.021}{224.431}{\Character{40}}%
\TextAt{13}{362.573}{224.431}{e}%
\TextAt{9}{367.998}{224.431}{\Character{41}}%
\TextAt{9}{375.195}{224.431}{\Character{43}}%
\TextAt{13}{386.969}{224.431}{e}%
\TextAt{16}{395.051}{224.431}{\Character{0}}%
\TextAt{9}{407.006}{224.431}{1}%
\TextAt{8}{412.859}{224.431}{.}%
\Vector{q 0 G 0 g 0.399 w 452.334 -549.569 m 452.334 -541.499 l S Q}%
\Vector{q 0 G 0 g 0.399 w 452.534 -541.699 m 459.707 -541.699 l S Q}%
\Vector{q 0 G 0 g 0.399 w 452.534 -549.37 m 459.707 -549.37 l S Q}%
\Vector{q 0 G 0 g 0.399 w 459.906 -549.569 m 459.906 -541.499 l S Q}%
\TextAt{8}{97.159}{195.43}{The}%
\TextAt{8}{121.963}{195.43}{geometric}%
\TextAt{8}{176.045}{195.43}{significance}%
\TextAt{8}{238.639}{195.43}{of}%
\TextAt{8}{252.719}{195.43}{the}%
\TextAt{8}{273.625}{195.43}{different,}%
\TextAt{8}{323.986}{195.43}{and}%
\TextAt{8}{347.498}{195.43}{thus}%
\TextAt{8}{374}{195.43}{also}%
\TextAt{8}{398.23}{195.43}{the}%
\TextAt{8}{419.136}{195.43}{way}%
\TextAt{8}{443.629}{195.43}{it}%
\TextAt{8}{456.083}{195.43}{fits}%
\TextAt{8}{476.404}{195.43}{into}%
\TextAt{8}{500.897}{195.43}{higher}%
\TextAt{8}{79.226}{178.384}{dimensional}%
\TextAt{8}{144.21}{178.384}{algebraic}%
\TextAt{8}{194.179}{178.384}{geometry,}%
\TextAt{8}{248.179}{178.384}{is}%
\TextAt{8}{260.476}{178.384}{brought}%
\TextAt{8}{304.922}{178.384}{out}%
\TextAt{8}{326.27}{178.384}{by}%
\TextAt{8}{343.063}{178.384}{the}%
\TextAt{8}{363.753}{178.384}{following}%
\TextAt{8}{413.722}{178.384}{characterization,}%
\TextAt{8}{503.18}{178.384}{which}%
\TextAt{8}{79.226}{161.338}{is}%
\TextAt{8}{90.945}{161.338}{due}%
\TextAt{8}{113.005}{161.338}{to}%
\TextAt{8}{127.271}{161.338}{E.}%
\TextAt{8}{142.337}{161.338}{Kähler.}%
\TextAt{8}{185.237}{161.338}{For}%
\TextAt{8}{206.162}{161.338}{an}%
\TextAt{8}{222.376}{161.338}{arbitrary}%
\TextAt{8}{272.1}{161.338}{extension}%
\TextAt{13}{323.821}{161.338}{B}%
\TextAt{16}{336.638}{161.338}{\Character{106}}%
\TextAt{13}{343.28}{161.338}{A}%
\TextAt{8}{355.908}{161.338}{of}%
\TextAt{8}{369.193}{161.338}{commutative}%
\TextAt{8}{439.074}{161.338}{rings,}%
\TextAt{8}{470.97}{161.338}{consider}%
\TextAt{8}{516.51}{161.338}{the}%
\TextAt{8}{79.226}{144.292}{homomorphism}%
\TextAt{13}{210.358}{127.246}{\Character{22}}%
\TextAt{9}{220.722}{127.246}{\Character{58}}%
\TextAt{13}{227.295}{127.246}{B}%
\TextAt{16}{239.447}{127.246}{\Character{10}}%
\TextAt{21}{248.746}{125.452}{A}%
\TextAt{13}{258.244}{127.246}{B}%
\NativeArrowAt{271.061}{127.246}{11.955}{19.26}{\longrightarrow}%
\TextAt{13}{293.643}{127.246}{B}%
\TextAt{13}{303.124}{127.246}{\Character{59}}%
\TextAt{13}{332.27}{127.246}{x}%
\TextAt{16}{341.602}{127.246}{\Character{10}}%
\TextAt{13}{353.558}{127.246}{y}%
\NativeArrowAt{363.015}{127.246}{11.955}{19.26}{\longmapsto}%
\TextAt{13}{385.597}{127.246}{xy}%
\TextAt{13}{398.365}{127.246}{\Character{59}}%
\TextAt{8}{79.226}{103.724}{whose}%
\TextAt{8}{113.753}{103.724}{kernel}%
\TextAt{8}{148.219}{103.724}{we}%
\TextAt{8}{165.447}{103.724}{denote}%
\TextAt{8}{203.165}{103.724}{by}%
\TextAt{13}{219.424}{103.724}{I}%
\TextAt{8}{225.527}{103.724}{.}%
\TextAt{8}{233.976}{103.724}{Then}%
\TextAt{9}{239.588}{74.723}{\Character{10}}%
\TextAt{17}{248.043}{79.659}{1}%
\TextAt{21}{248.043}{71.767}{B}%
\TextAt{3}{254.823}{71.767}{\Character{106}}%
\TextAt{21}{257.175}{71.767}{A}%
\TextAt{9}{267.337}{74.723}{\Character{58}\Character{61}}%
\TextAt{13}{283.014}{74.723}{I}%
\TextAt{9}{289.117}{74.723}{\Character{47}}%
\TextAt{13}{294.97}{74.723}{I}%
\TextAt{17}{301.073}{79.659}{2}%
\TextAt{9}{309.126}{74.723}{\Character{61}}%
\TextAt{13}{321.551}{74.723}{I}%
\TextAt{16}{330.311}{74.723}{\Character{10}}%
\TextAt{21}{339.609}{72.93}{B}%
\TextAt{3}{346.39}{72.93}{\Character{10}}%
\TextAt{21}{352.976}{72.93}{B}%
\TextAt{13}{362.911}{74.723}{B}%
\TextAt{8}{297.211}{43.033}{182}%
\Destination{78.23}{809.815}{nav.page.182}%
\Destination{79.36}{525.355}{nav.print.200}%
\Destination{79.36}{749.155}{nav.number.3.2.6}%
\Hotlink{253.259}{579.355}{9.751}{11.047}{nav.chapter.2}%
\Hotlink{213.463}{508.155}{9.751}{11.047}{nav.chapter.2}%
\Hotlink{134.635}{596.355}{25.078}{11.047}{nav.number.3.2.2}%
\Hotlink{269.373}{579.355}{30.936}{11.047}{nav.number.2.10.4}%
\Hotlink{289.946}{562.255}{25.078}{11.047}{nav.number.3.2.4}%
\Hotlink{114.054}{508.155}{25.078}{11.047}{nav.number.3.2.2}%
\Hotlink{228.407}{508.155}{25.078}{11.047}{nav.number.2.7.5}%
\Destination{79.226}{377.259}{errata.97}%
}

% Page 195 (body page 183)
\Page{%
\TextAt{2}{79.226}{796.943}{III.}%
\TextAt{2}{98.317}{796.943}{Riemann–Roch}%
\TextAt{2}{174.375}{796.943}{Theory}%
\Vector{q 0 G 0 g 0.299 w 7.226 18.729 m 460.769 18.729 l S Q}%
\TextAt{8}{79.226}{761.124}{is}%
\TextAt{8}{91.75}{761.124}{a}%
\TextAt{13}{102.266}{761.124}{B}%
\TextAt{16}{114.937}{761.124}{\Character{10}}%
\TextAt{13}{127.411}{761.124}{B}%
\TextAt{8}{136.907}{761.124}{-module,}%
\TextAt{8}{186.3}{761.124}{and}%
\TextAt{8}{209.823}{761.124}{hence}%
\TextAt{8}{243.089}{761.124}{in}%
\TextAt{8}{257.503}{761.124}{particular}%
\TextAt{8}{312.253}{761.124}{also}%
\TextAt{8}{336.493}{761.124}{a}%
\TextAt{13}{347.009}{761.124}{B}%
\TextAt{8}{356.505}{761.124}{-module,}%
\TextAt{8}{405.898}{761.124}{via}%
\TextAt{8}{425.846}{761.124}{the}%
\TextAt{8}{446.763}{761.124}{embedding}%
\TextAt{13}{506.701}{761.124}{B}%
\NativeMathAt{520.814}{761.124}{11.955}{\rightarrow}%
\TextAt{13}{79.226}{744.078}{B}%
\TextAt{16}{91.661}{744.078}{\Character{10}}%
\TextAt{13}{103.898}{744.078}{B}%
\TextAt{8}{113.394}{744.078}{,}%
\TextAt{13}{121.061}{744.078}{b}%
\NativeArrowAt{130.063}{744.078}{11.955}{11.955}{\mapsto}%
\TextAt{13}{146.044}{744.078}{b}%
\TextAt{16}{153.96}{744.078}{\Character{10}}%
\TextAt{9}{166.197}{744.078}{1}%
\TextAt{8}{172.05}{744.078}{.}%
\TextAt{8}{181.739}{744.078}{It}%
\TextAt{8}{194.825}{744.078}{is}%
\TextAt{8}{207.003}{744.078}{called}%
\TextAt{8}{240.58}{744.078}{the}%
\TextAt{12}{261.15}{744.078}{module}%
\TextAt{12}{309.243}{744.078}{of}%
\TextAt{12}{325.046}{744.078}{differentials}%
\TextAt{8}{397.848}{744.078}{of}%
\TextAt{13}{411.592}{744.078}{B}%
\TextAt{16}{425.114}{744.078}{\Character{106}}%
\TextAt{13}{432.46}{744.078}{A}%
\TextAt{8}{441.236}{744.078}{,}%
\TextAt{8}{448.902}{744.078}{and}%
\TextAt{8}{472.078}{744.078}{its}%
\TextAt{8}{488.81}{744.078}{elements}%
\TextAt{8}{79.226}{727.032}{are}%
\TextAt{8}{98.737}{727.032}{called}%
\TextAt{12}{131.901}{727.032}{Kähler}%
\TextAt{12}{176.302}{727.032}{differentials}%
\TextAt{8}{244.793}{727.032}{.}%
\TextAt{8}{253.242}{727.032}{If}%
\TextAt{8}{264.934}{727.032}{we}%
\TextAt{8}{282.161}{727.032}{put}%
\TextAt{8}{235.701}{698.031}{d}%
\TextAt{13}{242.205}{698.031}{x}%
\TextAt{9}{252.178}{698.031}{\Character{61}}%
\TextAt{13}{264.603}{698.031}{x}%
\TextAt{16}{273.912}{698.031}{\Character{10}}%
\TextAt{9}{285.867}{698.031}{1}%
\TextAt{16}{294.377}{698.031}{\Character{0}}%
\TextAt{9}{306.332}{698.031}{1}%
\TextAt{16}{314.842}{698.031}{\Character{10}}%
\TextAt{13}{326.797}{698.031}{x}%
\TextAt{8}{336.77}{698.031}{mod}%
\TextAt{13}{362.208}{698.031}{I}%
\TextAt{17}{368.31}{702.967}{2}%
\TextAt{13}{373.043}{698.031}{\Character{59}}%
\TextAt{8}{79.226}{669.029}{then}%
\TextAt{8}{105.886}{669.029}{we}%
\TextAt{8}{123.114}{669.029}{obtain}%
\TextAt{8}{159.541}{669.029}{a}%
\TextAt{8}{169.296}{669.029}{mapping}%
\TextAt{8}{267.886}{651.983}{d}%
\TextAt{9}{277.71}{651.983}{\Character{58}}%
\TextAt{13}{284.283}{651.983}{B}%
\NativeArrowAt{297.1}{651.983}{11.955}{19.26}{\longrightarrow}%
\TextAt{9}{319.682}{651.983}{\Character{10}}%
\TextAt{17}{328.136}{656.92}{1}%
\TextAt{21}{328.136}{649.028}{B}%
\TextAt{3}{334.916}{649.028}{\Character{106}}%
\TextAt{21}{337.269}{649.028}{A}%
\TextAt{8}{79.226}{628.462}{satisfying}%
\TextAt{8}{196.448}{611.416}{d}%
\TextAt{9}{202.952}{611.416}{\Character{40}}%
\TextAt{13}{207.504}{611.416}{xy}%
\TextAt{9}{220.293}{611.416}{\Character{41}}%
\TextAt{9}{228.148}{611.416}{\Character{61}}%
\TextAt{13}{240.591}{611.416}{x}%
\TextAt{8}{249.236}{611.416}{d}%
\TextAt{13}{255.74}{611.416}{y}%
\TextAt{9}{264.533}{611.416}{\Character{43}}%
\TextAt{13}{276.294}{611.416}{y}%
\TextAt{8}{284.423}{611.416}{d}%
\TextAt{13}{290.927}{611.416}{x\Character{59}}%
\TextAt{8}{326.734}{611.416}{d}%
\TextAt{13}{333.237}{611.416}{a}%
\TextAt{9}{342.703}{611.416}{\Character{61}}%
\TextAt{9}{355.113}{611.416}{0}%
\TextAt{8}{364.879}{611.416}{for}%
\TextAt{13}{382.764}{611.416}{a}%
\NativeMathAt{392.23}{611.416}{11.955}{\in}%
\TextAt{13}{403.521}{611.416}{A\Character{58}}%
\TextAt{8}{79.226}{587.894}{Such}%
\TextAt{8}{107.493}{587.894}{a}%
\TextAt{8}{117.229}{587.894}{map}%
\TextAt{8}{143.224}{587.894}{is}%
\TextAt{8}{154.969}{587.894}{called}%
\TextAt{8}{188.113}{587.894}{a}%
\TextAt{12}{197.849}{587.894}{derivation}%
\TextAt{8}{260.642}{587.894}{of}%
\TextAt{13}{273.953}{587.894}{B}%
\TextAt{16}{286.77}{587.894}{\Character{106}}%
\TextAt{13}{293.412}{587.894}{A}%
\TextAt{8}{302.187}{587.894}{.}%
\TextAt{8}{310.629}{587.894}{One}%
\TextAt{8}{335.321}{587.894}{can}%
\TextAt{8}{356.762}{587.894}{show}%
\TextAt{8}{385.746}{587.894}{that}%
\TextAt{8}{411.095}{587.894}{d}%
\TextAt{8}{421.477}{587.894}{is}%
\TextAt{8}{433.221}{587.894}{universal}%
\TextAt{8}{482.697}{587.894}{among}%
\TextAt{8}{520.408}{587.894}{all}%
\TextAt{8}{79.226}{570.848}{derivations}%
\TextAt{8}{139.077}{570.848}{of}%
\TextAt{13}{152.674}{570.848}{B}%
\TextAt{16}{165.945}{570.848}{\Character{106}}%
\TextAt{13}{173.041}{570.848}{A}%
\TextAt{8}{185.98}{570.848}{with}%
\TextAt{8}{212.907}{570.848}{values}%
\TextAt{8}{248.034}{570.848}{in}%
\TextAt{13}{261.954}{570.848}{B}%
\TextAt{8}{271.45}{570.848}{-modules.}%
\TextAt{9}{326.606}{570.848}{\Character{10}}%
\TextAt{17}{335.06}{575.187}{1}%
\TextAt{21}{335.06}{567.296}{B}%
\TextAt{3}{341.84}{567.296}{\Character{106}}%
\TextAt{21}{344.192}{567.296}{A}%
\TextAt{8}{355.197}{570.848}{consists}%
\TextAt{8}{398.574}{570.848}{of}%
\TextAt{8}{412.171}{570.848}{the}%
\TextAt{8}{432.594}{570.848}{linear}%
\TextAt{8}{465.378}{570.848}{combinations}%
\TextAt{20}{79.226}{562.769}{\Character{80}}%
\TextAt{13}{93.838}{553.802}{y}%
\TextAt{21}{99.546}{552.009}{i}%
\TextAt{8}{104.92}{553.802}{d}%
\TextAt{13}{111.423}{553.802}{x}%
\TextAt{21}{118.075}{552.009}{i}%
\TextAt{8}{121.456}{553.802}{.}%
\TextAt{8}{129.905}{553.802}{The}%
\TextAt{8}{153.959}{553.802}{link}%
\TextAt{8}{177.044}{553.802}{with}%
\TextAt{8}{203.704}{553.802}{the}%
\TextAt{8}{223.86}{553.802}{different}%
\TextAt{8}{270.031}{553.802}{is}%
\TextAt{8}{281.795}{553.802}{now}%
\TextAt{8}{306.184}{553.802}{this.}%
\TextAt{12}{79.226}{527.79}{Proposition}%
\TextAt{12}{151.937}{527.79}{(2.7).}%
\TextAt{8}{187.253}{527.79}{The}%
\TextAt{8}{211.307}{527.79}{different}%
\TextAt{14}{257.478}{527.79}{D}%
\TextAt{3}{267.419}{525.93}{\Character{79}}%
\TextAt{3}{274.353}{525.93}{\Character{106}}%
\TextAt{27}{276.715}{525.93}{o}%
\TextAt{8}{285.006}{527.79}{is}%
\TextAt{8}{296.77}{527.79}{the}%
\TextAt{12}{316.927}{527.79}{annihilator}%
\TextAt{8}{385.011}{527.79}{of}%
\TextAt{8}{398.341}{527.79}{the}%
\TextAt{16}{418.498}{527.79}{\Character{79}}%
\TextAt{8}{428.348}{527.79}{-module}%
\TextAt{9}{473.538}{527.79}{\Character{10}}%
\TextAt{17}{481.992}{532.128}{1}%
\TextAt{3}{481.992}{524.238}{\Character{79}}%
\TextAt{3}{488.926}{524.238}{\Character{106}}%
\TextAt{27}{491.288}{524.238}{o}%
\TextAt{8}{495.682}{527.79}{,}%
\TextAt{8}{502.831}{527.79}{i.e.,}%
\TextAt{14}{202.406}{498.789}{D}%
\TextAt{3}{212.347}{496.929}{\Character{79}}%
\TextAt{3}{219.281}{496.929}{\Character{106}}%
\TextAt{27}{221.643}{496.929}{o}%
\TextAt{9}{229.358}{498.789}{\Character{61}}%
\TextAt{16}{241.783}{498.789}{\Character{102}}%
\TextAt{13}{247.761}{498.789}{x}%
\NativeMathAt{257.734}{498.789}{11.955}{\in}%
\TextAt{16}{269.008}{498.789}{\Character{79}}%
\TextAt{16}{282.182}{498.789}{\Character{106}}%
\TextAt{13}{288.837}{498.789}{x}%
\TextAt{8}{297.482}{498.789}{d}%
\TextAt{13}{303.985}{498.789}{y}%
\TextAt{9}{313.443}{498.789}{\Character{61}}%
\TextAt{9}{325.853}{498.789}{0}%
\TextAt{8}{335.619}{498.789}{for}%
\TextAt{8}{353.504}{498.789}{all}%
\TextAt{13}{369.763}{498.789}{y}%
\NativeMathAt{379.22}{498.789}{11.955}{\in}%
\TextAt{16}{390.494}{498.789}{\Character{79}}%
\TextAt{16}{400.345}{498.789}{\Character{103}}%
\TextAt{13}{406.338}{498.789}{\Character{58}}%
\TextAt{19}{79.226}{451.855}{Proof.}%
\TextAt{8}{116}{451.855}{For}%
\TextAt{8}{137.223}{451.855}{greater}%
\TextAt{8}{177.155}{451.855}{notational}%
\TextAt{8}{233.358}{451.855}{clarity,}%
\TextAt{8}{272.373}{451.855}{let}%
\TextAt{8}{289.531}{451.855}{us}%
\TextAt{8}{304.8}{451.855}{put}%
\TextAt{16}{326.512}{451.855}{\Character{79}}%
\TextAt{9}{340.114}{451.855}{\Character{61}}%
\TextAt{13}{352.971}{451.855}{B}%
\TextAt{8}{366.617}{451.855}{and}%
\TextAt{14}{389.633}{451.855}{o}%
\TextAt{9}{399.229}{451.855}{\Character{61}}%
\TextAt{13}{412.086}{451.855}{A}%
\TextAt{8}{420.861}{451.855}{.}%
\TextAt{8}{430.069}{451.855}{If}%
\TextAt{13}{442.014}{451.855}{A}%
\TextAt{3}{450.789}{456.193}{\Character{48}}%
\TextAt{8}{457.735}{451.855}{is}%
\TextAt{8}{469.752}{451.855}{any}%
\TextAt{8}{492.122}{451.855}{commu-}%
\TextAt{8}{79.226}{434.809}{tative}%
\TextAt{13}{112.121}{434.809}{A}%
\TextAt{8}{120.896}{434.809}{-algebra}%
\TextAt{8}{165.495}{434.809}{and}%
\TextAt{13}{187.977}{434.809}{B}%
\TextAt{3}{197.473}{439.147}{\Character{48}}%
\TextAt{9}{203.589}{434.809}{\Character{61}}%
\TextAt{13}{216.015}{434.809}{B}%
\TextAt{16}{227.594}{434.809}{\Character{10}}%
\TextAt{21}{236.892}{433.016}{A}%
\TextAt{13}{245.816}{434.809}{A}%
\TextAt{3}{254.591}{439.147}{\Character{48}}%
\TextAt{8}{257.386}{434.809}{,}%
\TextAt{8}{264.311}{434.809}{then}%
\TextAt{8}{290.69}{434.809}{it}%
\TextAt{8}{302.113}{434.809}{is}%
\TextAt{8}{313.597}{434.809}{easy}%
\TextAt{8}{339.067}{434.809}{to}%
\TextAt{8}{353.097}{434.809}{see}%
\TextAt{8}{371.729}{434.809}{that}%
\TextAt{9}{396.817}{434.809}{\Character{10}}%
\TextAt{17}{405.272}{439.147}{1}%
\TextAt{21}{405.272}{431.257}{B}%
\TextAt{23}{412.052}{433.526}{\Character{48}}%
\TextAt{3}{414.736}{431.257}{\Character{106}}%
\TextAt{21}{417.089}{431.257}{A}%
\TextAt{23}{423.432}{433.526}{\Character{48}}%
\TextAt{9}{429.935}{434.809}{\Character{61}}%
\TextAt{9}{442.345}{434.809}{\Character{10}}%
\TextAt{17}{450.815}{439.147}{1}%
\TextAt{21}{450.815}{431.257}{B}%
\TextAt{3}{457.595}{431.257}{\Character{106}}%
\TextAt{21}{459.947}{431.257}{A}%
\TextAt{16}{468.871}{434.809}{\Character{10}}%
\TextAt{21}{478.169}{433.016}{A}%
\TextAt{13}{487.093}{434.809}{A}%
\TextAt{3}{495.868}{439.147}{\Character{48}}%
\TextAt{8}{498.663}{434.809}{.}%
\TextAt{8}{507.018}{434.809}{Thus}%
\TextAt{8}{79.226}{417.763}{the}%
\TextAt{8}{100.131}{417.763}{module}%
\TextAt{8}{142.173}{417.763}{of}%
\TextAt{8}{156.251}{417.763}{differentials}%
\TextAt{8}{220.147}{417.763}{is}%
\TextAt{8}{232.66}{417.763}{preserved}%
\TextAt{8}{285.497}{417.763}{under}%
\TextAt{8}{319.409}{417.763}{localization}%
\TextAt{8}{382.277}{417.763}{and}%
\TextAt{8}{405.788}{417.763}{completion,}%
\TextAt{8}{469.811}{417.763}{and}%
\TextAt{8}{493.322}{417.763}{we}%
\TextAt{8}{511.298}{417.763}{may}%
\TextAt{8}{79.226}{400.717}{therefore}%
\TextAt{8}{128.362}{400.717}{assume}%
\TextAt{8}{168.843}{400.717}{that}%
\TextAt{13}{194.248}{400.717}{A}%
\TextAt{8}{206.957}{400.717}{is}%
\TextAt{8}{218.757}{400.717}{a}%
\TextAt{8}{228.549}{400.717}{complete}%
\TextAt{8}{278.008}{400.717}{discrete}%
\TextAt{8}{321.023}{400.717}{valuation}%
\TextAt{8}{372.132}{400.717}{ring.}%
\TextAt{8}{400.857}{400.717}{Then}%
\TextAt{8}{431.451}{400.717}{we}%
\TextAt{8}{448.715}{400.717}{find}%
\TextAt{8}{472.159}{400.717}{by}%
\TextAt{8}{488.454}{400.717}{chap.}%
\TextAt{8}{520.755}{400.717}{II,}%
\TextAt{8}{79.226}{383.671}{(10.4),}%
\TextAt{8}{116.549}{383.671}{that}%
\TextAt{13}{142.109}{383.671}{B}%
\TextAt{9}{155.251}{383.671}{\Character{61}}%
\TextAt{13}{168.001}{383.671}{A}%
\TextAt{9}{176.776}{383.671}{\Character{91}}%
\TextAt{13}{180.028}{383.671}{x}%
\TextAt{9}{186.68}{383.671}{\Character{93}}%
\TextAt{8}{189.931}{383.671}{,}%
\TextAt{8}{197.319}{383.671}{and}%
\TextAt{8}{220.272}{383.671}{if}%
\TextAt{13}{231.186}{383.671}{f}%
\TextAt{9}{238.233}{383.671}{\Character{40}}%
\TextAt{13}{242.785}{383.671}{X}%
\TextAt{9}{253.44}{383.671}{\Character{41}}%
\NativeMathAt{261.638}{383.671}{11.955}{\in}%
\TextAt{13}{273.254}{383.671}{A}%
\TextAt{9}{282.029}{383.671}{\Character{91}}%
\TextAt{13}{285.281}{383.671}{X}%
\TextAt{9}{295.936}{383.671}{\Character{93}}%
\TextAt{8}{303.276}{383.671}{is}%
\TextAt{8}{315.23}{383.671}{the}%
\TextAt{8}{335.577}{383.671}{minimal}%
\TextAt{8}{381.293}{383.671}{polynomial}%
\TextAt{8}{441.977}{383.671}{of}%
\TextAt{13}{455.497}{383.671}{x}%
\TextAt{8}{462.149}{383.671}{,}%
\TextAt{8}{469.537}{383.671}{then}%
\TextAt{9}{496.387}{383.671}{\Character{10}}%
\TextAt{17}{504.841}{388.009}{1}%
\TextAt{21}{504.841}{380.119}{B}%
\TextAt{3}{511.622}{380.119}{\Character{106}}%
\TextAt{21}{513.974}{380.119}{A}%
\TextAt{8}{524.903}{383.671}{is}%
\TextAt{8}{79.226}{366.625}{generated}%
\TextAt{8}{132.35}{366.625}{by}%
\TextAt{8}{148.402}{366.625}{d}%
\TextAt{13}{154.905}{366.625}{x}%
\TextAt{8}{165.247}{366.625}{(exercise}%
\TextAt{8}{212.896}{366.625}{3).}%
\TextAt{8}{231.688}{366.625}{The}%
\TextAt{8}{255.534}{366.625}{annihilator}%
\TextAt{8}{315.174}{366.625}{of}%
\TextAt{8}{328.296}{366.625}{d}%
\TextAt{13}{334.8}{366.625}{x}%
\TextAt{8}{345.142}{366.625}{is}%
\TextAt{13}{356.698}{366.625}{f}%
\TextAt{3}{363.744}{370.963}{\Character{48}}%
\TextAt{9}{366.539}{366.625}{\Character{40}}%
\TextAt{13}{371.092}{366.625}{x}%
\TextAt{9}{377.744}{366.625}{\Character{41}}%
\TextAt{8}{382.296}{366.625}{.}%
\TextAt{8}{390.675}{366.625}{On}%
\TextAt{8}{409.978}{366.625}{the}%
\TextAt{8}{429.927}{366.625}{other}%
\TextAt{8}{460.289}{366.625}{hand,}%
\TextAt{8}{492.64}{366.625}{by}%
\TextAt{8}{508.692}{366.625}{(2.4)}%
\TextAt{8}{79.226}{349.579}{we}%
\TextAt{8}{96.454}{349.579}{have}%
\TextAt{14}{123.448}{349.579}{D}%
\TextAt{21}{133.39}{347.719}{B}%
\TextAt{3}{140.17}{347.719}{\Character{106}}%
\TextAt{21}{142.523}{347.719}{A}%
\TextAt{9}{152.684}{349.579}{\Character{61}}%
\TextAt{9}{165.094}{349.579}{\Character{40}}%
\TextAt{13}{169.662}{349.579}{f}%
\TextAt{3}{176.709}{353.917}{\Character{48}}%
\TextAt{9}{179.504}{349.579}{\Character{40}}%
\TextAt{13}{184.056}{349.579}{x}%
\TextAt{9}{190.708}{349.579}{\Character{41}\Character{41}}%
\TextAt{8}{199.813}{349.579}{.}%
\TextAt{8}{208.261}{349.579}{This}%
\TextAt{8}{234.981}{349.579}{proves}%
\TextAt{8}{271.145}{349.579}{the}%
\TextAt{8}{291.302}{349.579}{claim.}%
\Vector{q 0 G 0 g 0.399 w 452.334 -424.421 m 452.334 -416.351 l S Q}%
\Vector{q 0 G 0 g 0.399 w 452.534 -416.55 m 459.707 -416.55 l S Q}%
\Vector{q 0 G 0 g 0.399 w 452.534 -424.222 m 459.707 -424.222 l S Q}%
\Vector{q 0 G 0 g 0.399 w 459.906 -424.421 m 459.906 -416.351 l S Q}%
\TextAt{8}{97.159}{320.578}{A}%
\TextAt{8}{109.376}{320.578}{most}%
\TextAt{8}{137.601}{320.578}{intimate}%
\TextAt{8}{183.651}{320.578}{connection}%
\TextAt{8}{241.728}{320.578}{holds}%
\TextAt{8}{271.902}{320.578}{between}%
\TextAt{8}{316.637}{320.578}{the}%
\TextAt{8}{336.338}{320.578}{different}%
\TextAt{8}{382.054}{320.578}{and}%
\TextAt{8}{404.361}{320.578}{the}%
\TextAt{8}{424.062}{320.578}{discriminant}%
\TextAt{8}{490.986}{320.578}{of}%
\TextAt{16}{503.861}{320.578}{\Character{79}}%
\TextAt{16}{517.024}{320.578}{\Character{106}}%
\TextAt{14}{523.673}{320.578}{o}%
\TextAt{8}{529.518}{320.578}{.}%
\TextAt{8}{79.226}{303.532}{The}%
\TextAt{8}{103.28}{303.532}{latter}%
\TextAt{8}{135.152}{303.532}{is}%
\TextAt{8}{146.916}{303.532}{defined}%
\TextAt{8}{187.229}{303.532}{as}%
\TextAt{8}{201.599}{303.532}{follows.}%
\TextAt{12}{79.226}{277.52}{Definition}%
\TextAt{12}{143.384}{277.52}{(2.8).}%
\TextAt{8}{179.889}{277.52}{The}%
\TextAt{12}{205.133}{277.52}{discriminant}%
\TextAt{14}{283.266}{277.52}{d}%
\TextAt{3}{289.227}{275.66}{\Character{79}}%
\TextAt{3}{296.161}{275.66}{\Character{106}}%
\TextAt{27}{298.522}{275.66}{o}%
\TextAt{8}{308.004}{277.52}{is}%
\TextAt{8}{320.958}{277.52}{the}%
\TextAt{8}{342.304}{277.52}{ideal}%
\TextAt{8}{371.457}{277.52}{of}%
\TextAt{14}{385.977}{277.52}{o}%
\TextAt{8}{396.909}{277.52}{which}%
\TextAt{8}{431.585}{277.52}{is}%
\TextAt{8}{444.539}{277.52}{generated}%
\TextAt{8}{499.061}{277.52}{by}%
\TextAt{8}{516.51}{277.52}{the}%
\TextAt{8}{79.226}{260.474}{discriminants}%
\TextAt{13}{150.946}{260.474}{d}%
\TextAt{9}{157.029}{260.474}{\Character{40}}%
\TextAt{13}{161.581}{260.474}{\Character{11}}%
\TextAt{17}{169.027}{258.68}{1}%
\TextAt{13}{173.759}{260.474}{\Character{59}}%
\TextAt{13}{178.983}{260.474}{\Character{58}}%
\TextAt{13}{184.22}{260.474}{\Character{58}}%
\TextAt{13}{189.456}{260.474}{\Character{58}}%
\TextAt{13}{194.681}{260.474}{\Character{59}}%
\TextAt{13}{199.917}{260.474}{\Character{11}}%
\TextAt{21}{207.426}{258.68}{n}%
\TextAt{9}{213.062}{260.474}{\Character{41}}%
\TextAt{8}{221.238}{260.474}{of}%
\TextAt{8}{234.294}{260.474}{all}%
\TextAt{8}{250.279}{260.474}{the}%
\TextAt{8}{270.161}{260.474}{bases}%
\TextAt{13}{300.576}{260.474}{\Character{11}}%
\TextAt{17}{308.022}{258.68}{1}%
\TextAt{13}{312.754}{260.474}{\Character{59}}%
\TextAt{13}{317.978}{260.474}{\Character{58}}%
\TextAt{13}{323.215}{260.474}{\Character{58}}%
\TextAt{13}{328.451}{260.474}{\Character{58}}%
\TextAt{13}{333.676}{260.474}{\Character{59}}%
\TextAt{13}{338.912}{260.474}{\Character{11}}%
\TextAt{21}{346.42}{258.68}{n}%
\TextAt{8}{355.68}{260.474}{of}%
\TextAt{13}{368.736}{260.474}{L}%
\TextAt{16}{380.022}{260.474}{\Character{106}}%
\TextAt{13}{386.663}{260.474}{K}%
\TextAt{8}{401.104}{260.474}{which}%
\TextAt{8}{434.317}{260.474}{are}%
\TextAt{8}{453.554}{260.474}{contained}%
\TextAt{8}{506.289}{260.474}{in}%
\TextAt{16}{519.668}{260.474}{\Character{79}}%
\TextAt{8}{529.518}{260.474}{.}%
\TextAt{8}{97.159}{234.461}{We}%
\TextAt{8}{118.823}{234.461}{will}%
\TextAt{8}{142.447}{234.461}{frequently}%
\TextAt{8}{199.247}{234.461}{write}%
\TextAt{14}{230.678}{234.461}{d}%
\TextAt{21}{236.639}{232.602}{L}%
\TextAt{3}{242.398}{232.601}{\Character{106}}%
\TextAt{21}{244.75}{232.602}{K}%
\TextAt{8}{258.37}{234.461}{instead}%
\TextAt{8}{300.273}{234.461}{of}%
\TextAt{14}{315.123}{234.461}{d}%
\TextAt{3}{321.083}{232.602}{\Character{79}}%
\TextAt{3}{328.017}{232.602}{\Character{106}}%
\TextAt{27}{330.379}{232.601}{o}%
\TextAt{8}{334.773}{234.461}{.}%
\TextAt{8}{347.779}{234.461}{If}%
\TextAt{13}{360.991}{234.461}{\Character{11}}%
\TextAt{17}{368.436}{232.668}{1}%
\TextAt{13}{373.169}{234.461}{\Character{59}}%
\TextAt{13}{378.393}{234.461}{\Character{58}}%
\TextAt{13}{383.63}{234.461}{\Character{58}}%
\TextAt{13}{388.854}{234.461}{\Character{58}}%
\TextAt{13}{394.091}{234.461}{\Character{59}}%
\TextAt{13}{399.327}{234.461}{\Character{11}}%
\TextAt{21}{406.835}{232.668}{n}%
\TextAt{8}{417.888}{234.461}{is}%
\TextAt{8}{431.171}{234.461}{an}%
\TextAt{8}{448.95}{234.461}{integral}%
\TextAt{8}{493.077}{234.461}{basis}%
\TextAt{8}{523.337}{234.461}{of}%
\TextAt{16}{79.226}{217.415}{\Character{79}}%
\TextAt{16}{94.445}{217.415}{\Character{106}}%
\TextAt{14}{103.138}{217.415}{o}%
\TextAt{8}{108.982}{217.415}{,}%
\TextAt{8}{117.635}{217.415}{then}%
\TextAt{14}{145.498}{217.415}{d}%
\TextAt{21}{151.458}{215.556}{L}%
\TextAt{3}{157.217}{215.556}{\Character{106}}%
\TextAt{21}{159.569}{215.556}{K}%
\TextAt{8}{172.872}{217.415}{is}%
\TextAt{8}{185.839}{217.415}{the}%
\TextAt{8}{207.198}{217.415}{principal}%
\TextAt{8}{257.178}{217.415}{ideal}%
\TextAt{8}{286.344}{217.415}{generated}%
\TextAt{8}{340.879}{217.415}{by}%
\TextAt{13}{358.34}{217.415}{d}%
\TextAt{9}{364.423}{217.415}{\Character{40}}%
\TextAt{13}{368.975}{217.415}{\Character{11}}%
\TextAt{17}{376.421}{215.622}{1}%
\TextAt{13}{381.153}{217.415}{\Character{59}}%
\TextAt{13}{386.377}{217.415}{\Character{58}}%
\TextAt{13}{391.614}{217.415}{\Character{58}}%
\TextAt{13}{396.85}{217.415}{\Character{58}}%
\TextAt{13}{402.075}{217.415}{\Character{59}}%
\TextAt{13}{407.311}{217.415}{\Character{11}}%
\TextAt{21}{414.82}{215.622}{n}%
\TextAt{9}{420.456}{217.415}{\Character{41}}%
\TextAt{9}{430.367}{217.415}{\Character{61}}%
\TextAt{13}{444.854}{217.415}{d}%
\TextAt{21}{450.937}{215.556}{L}%
\TextAt{3}{456.696}{215.556}{\Character{106}}%
\TextAt{21}{459.048}{215.556}{K}%
\TextAt{8}{467.251}{217.415}{,}%
\TextAt{8}{475.903}{217.415}{because}%
\TextAt{8}{520.408}{217.415}{all}%
\TextAt{8}{79.226}{200.369}{other}%
\TextAt{8}{109.702}{200.369}{bases}%
\TextAt{8}{140.297}{200.369}{contained}%
\TextAt{8}{193.212}{200.369}{in}%
\TextAt{16}{206.771}{200.369}{\Character{79}}%
\TextAt{8}{220.424}{200.369}{are}%
\TextAt{8}{239.841}{200.369}{transforms}%
\TextAt{8}{298.088}{200.369}{of}%
\TextAt{8}{311.324}{200.369}{the}%
\TextAt{8}{331.387}{200.369}{given}%
\TextAt{8}{361.862}{200.369}{one}%
\TextAt{8}{383.228}{200.369}{by}%
\TextAt{8}{399.393}{200.369}{matrices}%
\TextAt{8}{446.187}{200.369}{with}%
\TextAt{8}{472.753}{200.369}{entries}%
\TextAt{8}{510.115}{200.369}{in}%
\TextAt{14}{523.673}{200.369}{o}%
\TextAt{8}{529.518}{200.369}{.}%
\TextAt{8}{79.226}{183.323}{The}%
\TextAt{8}{103.28}{183.323}{discriminant}%
\TextAt{8}{170.659}{183.323}{is}%
\TextAt{8}{182.423}{183.323}{obtained}%
\TextAt{8}{230.555}{183.323}{from}%
\TextAt{8}{258.195}{183.323}{the}%
\TextAt{8}{278.351}{183.323}{different}%
\TextAt{8}{324.522}{183.323}{by}%
\TextAt{8}{340.781}{183.323}{taking}%
\TextAt{8}{376.886}{183.323}{the}%
\TextAt{8}{397.042}{183.323}{norm}%
\TextAt{13}{427.612}{183.323}{N}%
\TextAt{21}{436.986}{181.464}{L}%
\TextAt{3}{442.745}{181.464}{\Character{106}}%
\TextAt{21}{445.097}{181.464}{K}%
\TextAt{8}{457.197}{183.323}{(see}%
\TextAt{8}{480.665}{183.323}{§1).}%
\TextAt{12}{79.226}{157.311}{Theorem}%
\TextAt{12}{135.344}{157.311}{(2.9).}%
\TextAt{8}{169.88}{157.311}{The}%
\TextAt{8}{193.154}{157.311}{following}%
\TextAt{8}{241.809}{157.311}{relation}%
\TextAt{8}{283.961}{157.311}{exists}%
\TextAt{8}{315.496}{157.311}{between}%
\TextAt{8}{359.907}{157.311}{the}%
\TextAt{8}{379.284}{157.311}{discriminant}%
\TextAt{8}{445.884}{157.311}{and}%
\TextAt{8}{467.867}{157.311}{the}%
\TextAt{8}{487.244}{157.311}{different:}%
\TextAt{14}{254.837}{128.31}{d}%
\TextAt{21}{260.798}{126.45}{L}%
\TextAt{3}{266.557}{126.45}{\Character{106}}%
\TextAt{21}{268.909}{126.45}{K}%
\TextAt{9}{280.433}{128.31}{\Character{61}}%
\TextAt{13}{292.858}{128.31}{N}%
\TextAt{21}{302.232}{126.45}{L}%
\TextAt{3}{307.991}{126.45}{\Character{106}}%
\TextAt{21}{310.343}{126.45}{K}%
\TextAt{9}{318.546}{128.31}{\Character{40}}%
\TextAt{14}{323.098}{128.31}{D}%
\TextAt{21}{333.04}{126.45}{L}%
\TextAt{3}{338.799}{126.45}{\Character{106}}%
\TextAt{21}{341.152}{126.45}{K}%
\TextAt{9}{349.354}{128.31}{\Character{41}}%
\TextAt{13}{353.907}{128.31}{\Character{58}}%
\TextAt{8}{297.211}{43.033}{183}%
\Destination{78.23}{809.815}{nav.page.183}%
\Destination{79.36}{542.755}{nav.print.201}%
\Destination{79.36}{541.755}{nav.number.3.2.7}%
\Destination{79.36}{291.555}{nav.number.3.2.8}%
\Destination{79.36}{171.355}{nav.number.3.2.9}%
\Hotlink{520.255}{398.555}{9.751}{11.047}{nav.chapter.2}%
\Hotlink{212.396}{364.455}{6.858}{11.047}{nav.ex.3.2.3}%
\Hotlink{485.844}{181.155}{6.858}{11.047}{nav.section.3.1}%
\Hotlink{508.192}{364.455}{25.078}{11.047}{nav.number.3.2.4}%
\Hotlink{78.76}{381.5}{29.21}{11.047}{nav.number.2.10.4}%
}

% Page 196 (body page 184)
\Page{%
\TextAt{2}{79.226}{796.943}{III.}%
\TextAt{2}{98.317}{796.943}{Riemann–Roch}%
\TextAt{2}{174.375}{796.943}{Theory}%
\Vector{q 0 G 0 g 0.299 w 7.226 18.729 m 460.769 18.729 l S Q}%
\TextAt{19}{79.226}{761.124}{Proof.}%
\TextAt{8}{116}{761.124}{If}%
\TextAt{13}{130.67}{761.124}{S}%
\TextAt{8}{145.44}{761.124}{is}%
\TextAt{8}{160.181}{761.124}{a}%
\TextAt{8}{172.914}{761.124}{multiplicative}%
\TextAt{8}{249.715}{761.124}{subset}%
\TextAt{8}{288.582}{761.124}{of}%
\TextAt{14}{304.889}{761.124}{o}%
\TextAt{8}{310.733}{761.124}{,}%
\TextAt{8}{321.604}{761.124}{then}%
\TextAt{8}{351.242}{761.124}{clearly}%
\TextAt{14}{391.615}{761.124}{d}%
\TextAt{21}{397.576}{759.12}{S}%
\TextAt{23}{403.171}{761.389}{\Character{0}}%
\TextAt{24}{408.927}{761.389}{1}%
\TextAt{3}{413.078}{759.12}{\Character{79}}%
\TextAt{3}{420.012}{759.12}{\Character{106}}%
\TextAt{21}{422.374}{759.12}{S}%
\TextAt{23}{427.969}{761.389}{\Character{0}}%
\TextAt{24}{433.726}{761.389}{1}%
\TextAt{27}{437.877}{759.12}{o}%
\TextAt{9}{450.666}{761.124}{\Character{61}}%
\TextAt{13}{468.165}{761.124}{S}%
\TextAt{3}{476.061}{765.462}{\Character{0}}%
\TextAt{17}{482.647}{765.462}{1}%
\TextAt{14}{487.379}{761.124}{d}%
\TextAt{3}{493.34}{759.264}{\Character{79}}%
\TextAt{3}{500.274}{759.264}{\Character{106}}%
\TextAt{27}{502.635}{759.264}{o}%
\TextAt{8}{513.904}{761.124}{and}%
\TextAt{14}{79.226}{744.078}{D}%
\TextAt{21}{89.168}{742.074}{S}%
\TextAt{23}{94.764}{744.343}{\Character{0}}%
\TextAt{24}{100.52}{744.343}{1}%
\TextAt{3}{104.671}{742.074}{\Character{79}}%
\TextAt{3}{111.605}{742.074}{\Character{106}}%
\TextAt{21}{113.966}{742.074}{S}%
\TextAt{23}{119.562}{744.343}{\Character{0}}%
\TextAt{24}{125.318}{744.343}{1}%
\TextAt{27}{129.469}{742.074}{o}%
\TextAt{9}{139.542}{744.078}{\Character{61}}%
\TextAt{13}{154.326}{744.078}{S}%
\TextAt{3}{162.222}{748.416}{\Character{0}}%
\TextAt{17}{168.808}{748.416}{1}%
\TextAt{14}{173.54}{744.078}{D}%
\TextAt{3}{183.482}{742.218}{\Character{79}}%
\TextAt{3}{190.416}{742.218}{\Character{106}}%
\TextAt{27}{192.777}{742.218}{o}%
\TextAt{8}{197.172}{744.078}{.}%
\TextAt{8}{209.772}{744.078}{We}%
\TextAt{8}{231.3}{744.078}{may}%
\TextAt{8}{258.053}{744.078}{therefore}%
\TextAt{8}{308.537}{744.078}{assume}%
\TextAt{8}{350.365}{744.078}{that}%
\TextAt{14}{377.118}{744.078}{o}%
\TextAt{8}{388.243}{744.078}{is}%
\TextAt{8}{401.391}{744.078}{a}%
\TextAt{8}{412.53}{744.078}{discrete}%
\TextAt{8}{456.893}{744.078}{valuation}%
\TextAt{8}{509.349}{744.078}{ring.}%
\TextAt{8}{79.226}{727.032}{Then,}%
\TextAt{8}{113.092}{727.032}{since}%
\TextAt{14}{141.806}{727.032}{o}%
\TextAt{8}{151.593}{727.032}{is}%
\TextAt{8}{163.402}{727.032}{a}%
\TextAt{8}{173.203}{727.032}{principal}%
\TextAt{8}{222.026}{727.032}{ideal}%
\TextAt{8}{250.034}{727.032}{domain,}%
\TextAt{8}{294.971}{727.032}{so}%
\TextAt{8}{309.386}{727.032}{is}%
\TextAt{16}{321.196}{727.032}{\Character{79}}%
\TextAt{8}{334.988}{727.032}{(see}%
\TextAt{8}{358.502}{727.032}{chap.}%
\TextAt{8}{390.829}{727.032}{I,}%
\TextAt{8}{402.244}{727.032}{§3,}%
\TextAt{8}{420.987}{727.032}{exercise}%
\TextAt{8}{464.334}{727.032}{4),}%
\TextAt{8}{481.953}{727.032}{and}%
\TextAt{8}{504.761}{727.032}{it}%
\TextAt{8}{516.51}{727.032}{ad-}%
\TextAt{8}{79.226}{709.986}{mits}%
\TextAt{8}{105.018}{709.986}{an}%
\TextAt{8}{120.995}{709.986}{integral}%
\TextAt{8}{163.321}{709.986}{basis}%
\TextAt{13}{191.779}{709.986}{\Character{11}}%
\TextAt{17}{199.224}{708.193}{1}%
\TextAt{13}{203.957}{709.986}{\Character{59}}%
\TextAt{13}{209.181}{709.986}{\Character{58}}%
\TextAt{13}{214.418}{709.986}{\Character{58}}%
\TextAt{13}{219.642}{709.986}{\Character{58}}%
\TextAt{13}{224.879}{709.986}{\Character{59}}%
\TextAt{13}{230.115}{709.986}{\Character{11}}%
\TextAt{21}{237.623}{708.193}{n}%
\TextAt{8}{246.875}{709.986}{(see}%
\TextAt{8}{270.06}{709.986}{chap.}%
\TextAt{8}{302.158}{709.986}{I,}%
\TextAt{8}{313.245}{709.986}{(2.10)).}%
\TextAt{8}{356.09}{709.986}{So}%
\TextAt{8}{372.067}{709.986}{we}%
\TextAt{8}{389.012}{709.986}{have}%
\TextAt{14}{415.724}{709.986}{d}%
\TextAt{21}{421.685}{708.126}{L}%
\TextAt{3}{427.444}{708.126}{\Character{106}}%
\TextAt{21}{429.796}{708.126}{K}%
\TextAt{9}{441.32}{709.986}{\Character{61}}%
\TextAt{9}{453.73}{709.986}{\Character{40}}%
\TextAt{13}{458.297}{709.986}{d}%
\TextAt{9}{464.38}{709.986}{\Character{40}}%
\TextAt{13}{468.932}{709.986}{\Character{11}}%
\TextAt{17}{476.378}{708.193}{1}%
\TextAt{13}{481.11}{709.986}{\Character{59}}%
\TextAt{13}{486.334}{709.986}{\Character{58}}%
\TextAt{13}{491.571}{709.986}{\Character{58}}%
\TextAt{13}{496.807}{709.986}{\Character{58}}%
\TextAt{13}{502.032}{709.986}{\Character{59}}%
\TextAt{13}{507.268}{709.986}{\Character{11}}%
\TextAt{21}{514.777}{708.193}{n}%
\TextAt{9}{520.413}{709.986}{\Character{41}\Character{41}}%
\TextAt{8}{529.518}{709.986}{.}%
\TextAt{8}{79.226}{692.94}{Dedekind’s}%
\TextAt{8}{139.65}{692.94}{complementary}%
\TextAt{8}{221.976}{692.94}{module}%
\TextAt{14}{263.642}{692.94}{C}%
\TextAt{3}{270.965}{691.08}{\Character{79}}%
\TextAt{3}{277.899}{691.08}{\Character{106}}%
\TextAt{27}{280.26}{691.08}{o}%
\TextAt{8}{288.925}{692.94}{is}%
\TextAt{8}{301.062}{692.94}{generated}%
\TextAt{8}{354.767}{692.94}{by}%
\TextAt{8}{371.399}{692.94}{the}%
\TextAt{8}{391.929}{692.94}{dual}%
\TextAt{8}{418.316}{692.94}{basis}%
\TextAt{13}{447.429}{692.94}{\Character{11}}%
\TextAt{3}{454.951}{697.278}{\Character{48}}%
\TextAt{17}{454.875}{689.984}{1}%
\TextAt{13}{459.607}{692.94}{\Character{59}}%
\TextAt{13}{464.831}{692.94}{\Character{58}}%
\TextAt{13}{470.068}{692.94}{\Character{58}}%
\TextAt{13}{475.304}{692.94}{\Character{58}}%
\TextAt{13}{480.529}{692.94}{\Character{59}}%
\TextAt{13}{485.765}{692.94}{\Character{11}}%
\TextAt{3}{493.35}{697.278}{\Character{48}}%
\TextAt{21}{493.274}{689.984}{n}%
\TextAt{8}{503.18}{692.94}{which}%
\TextAt{8}{79.226}{675.894}{is}%
\TextAt{8}{91.755}{675.894}{characterized}%
\TextAt{8}{163.737}{675.894}{by}%
\TextAt{8}{180.762}{675.894}{Tr}%
\TextAt{21}{193.769}{674.034}{L}%
\TextAt{3}{199.528}{674.034}{\Character{106}}%
\TextAt{21}{201.88}{674.034}{K}%
\TextAt{9}{210.083}{675.894}{\Character{40}}%
\TextAt{13}{214.635}{675.894}{\Character{11}}%
\TextAt{21}{222.081}{674.101}{i}%
\TextAt{13}{225.462}{675.894}{\Character{11}}%
\TextAt{3}{232.984}{680.232}{\Character{48}}%
\TextAt{21}{232.908}{672.938}{j}%
\TextAt{9}{237.29}{675.894}{\Character{41}}%
\TextAt{9}{246.448}{675.894}{\Character{61}}%
\TextAt{13}{260.197}{675.894}{\Character{14}}%
\TextAt{21}{265.381}{674.101}{ij}%
\TextAt{8}{272.647}{675.894}{.}%
\TextAt{8}{283.391}{675.894}{On}%
\TextAt{8}{303.667}{675.894}{the}%
\TextAt{8}{324.589}{675.894}{other}%
\TextAt{8}{355.923}{675.894}{hand,}%
\TextAt{14}{389.398}{675.894}{C}%
\TextAt{3}{396.721}{674.034}{\Character{79}}%
\TextAt{3}{403.655}{674.034}{\Character{106}}%
\TextAt{27}{406.016}{674.034}{o}%
\TextAt{8}{415.073}{675.894}{is}%
\TextAt{8}{427.602}{675.894}{a}%
\TextAt{8}{438.123}{675.894}{principal}%
\TextAt{8}{487.665}{675.894}{ideal}%
\TextAt{9}{516.394}{675.894}{\Character{40}}%
\TextAt{13}{520.946}{675.894}{\Character{12}}%
\TextAt{9}{528.217}{675.894}{\Character{41}}%
\TextAt{8}{79.226}{658.848}{and}%
\TextAt{8}{101.989}{658.848}{admits}%
\TextAt{8}{140.425}{658.848}{the}%
\TextAt{14}{160.581}{658.848}{o}%
\TextAt{8}{166.425}{658.848}{-basis}%
\TextAt{13}{199.063}{658.848}{\Character{12}}%
\TextAt{13}{206.332}{658.848}{\Character{11}}%
\TextAt{17}{213.78}{657.055}{1}%
\TextAt{13}{218.512}{658.848}{\Character{59}}%
\TextAt{13}{223.736}{658.848}{\Character{58}}%
\TextAt{13}{228.973}{658.848}{\Character{58}}%
\TextAt{13}{234.209}{658.848}{\Character{58}}%
\TextAt{13}{239.434}{658.848}{\Character{59}}%
\TextAt{13}{244.67}{658.848}{\Character{12}}%
\TextAt{13}{251.939}{658.848}{\Character{11}}%
\TextAt{21}{259.45}{657.055}{n}%
\TextAt{8}{268.984}{658.848}{of}%
\TextAt{8}{282.314}{658.848}{discriminant}%
\TextAt{13}{199.162}{630.522}{d}%
\TextAt{9}{205.244}{630.522}{\Character{40}}%
\TextAt{13}{209.797}{630.522}{\Character{12}}%
\TextAt{13}{217.066}{630.522}{\Character{11}}%
\TextAt{17}{224.513}{628.729}{1}%
\TextAt{13}{229.246}{630.522}{\Character{59}}%
\TextAt{13}{234.47}{630.522}{\Character{58}}%
\TextAt{13}{239.707}{630.522}{\Character{58}}%
\TextAt{13}{244.931}{630.522}{\Character{58}}%
\TextAt{13}{250.168}{630.522}{\Character{59}}%
\TextAt{13}{255.404}{630.522}{\Character{12}}%
\TextAt{13}{262.673}{630.522}{\Character{11}}%
\TextAt{21}{270.183}{628.729}{n}%
\TextAt{9}{275.82}{630.522}{\Character{41}}%
\TextAt{9}{283.675}{630.522}{\Character{61}}%
\TextAt{13}{296.118}{630.522}{N}%
\TextAt{21}{305.492}{628.662}{L}%
\TextAt{3}{311.251}{628.662}{\Character{106}}%
\TextAt{21}{313.604}{628.662}{K}%
\TextAt{9}{321.806}{630.522}{\Character{40}}%
\TextAt{13}{326.359}{630.522}{\Character{12}}%
\TextAt{9}{333.63}{630.522}{\Character{41}}%
\TextAt{17}{338.182}{635.458}{2}%
\TextAt{13}{342.914}{630.522}{d}%
\TextAt{9}{348.997}{630.522}{\Character{40}}%
\TextAt{13}{353.549}{630.522}{\Character{11}}%
\TextAt{17}{360.995}{628.729}{1}%
\TextAt{13}{365.727}{630.522}{\Character{59}}%
\TextAt{13}{370.951}{630.522}{\Character{58}}%
\TextAt{13}{376.188}{630.522}{\Character{58}}%
\TextAt{13}{381.424}{630.522}{\Character{58}}%
\TextAt{13}{386.649}{630.522}{\Character{59}}%
\TextAt{13}{391.885}{630.522}{\Character{11}}%
\TextAt{21}{399.394}{628.729}{n}%
\TextAt{9}{405.03}{630.522}{\Character{41}}%
\TextAt{13}{409.582}{630.522}{\Character{58}}%
\TextAt{8}{79.226}{602.196}{But}%
\TextAt{9}{104.085}{602.196}{\Character{40}}%
\TextAt{13}{108.638}{602.196}{N}%
\TextAt{21}{118.011}{600.337}{L}%
\TextAt{3}{123.771}{600.337}{\Character{106}}%
\TextAt{21}{126.123}{600.337}{K}%
\TextAt{9}{134.326}{602.196}{\Character{40}}%
\TextAt{13}{138.878}{602.196}{\Character{12}}%
\TextAt{9}{146.149}{602.196}{\Character{41}\Character{41}}%
\TextAt{9}{161.308}{602.196}{\Character{61}}%
\TextAt{13}{176.516}{602.196}{N}%
\TextAt{21}{185.89}{600.337}{L}%
\TextAt{3}{191.649}{600.337}{\Character{106}}%
\TextAt{21}{194.001}{600.337}{K}%
\TextAt{9}{202.204}{602.196}{\Character{40}}%
\TextAt{14}{206.756}{602.196}{C}%
\TextAt{3}{214.079}{600.337}{\Character{79}}%
\TextAt{3}{221.013}{600.337}{\Character{106}}%
\TextAt{27}{223.374}{600.337}{o}%
\TextAt{9}{227.769}{602.196}{\Character{41}}%
\TextAt{9}{238.385}{602.196}{\Character{61}}%
\TextAt{13}{253.583}{602.196}{N}%
\TextAt{21}{262.957}{600.337}{L}%
\TextAt{3}{268.716}{600.337}{\Character{106}}%
\TextAt{21}{271.068}{600.337}{K}%
\TextAt{9}{279.271}{602.196}{\Character{40}}%
\TextAt{14}{283.823}{602.196}{D}%
\TextAt{3}{293.765}{607.145}{\Character{0}}%
\TextAt{17}{300.351}{607.145}{1}%
\TextAt{21}{293.765}{598.424}{L}%
\TextAt{3}{299.524}{598.424}{\Character{106}}%
\TextAt{21}{301.876}{598.424}{K}%
\TextAt{9}{310.079}{602.196}{\Character{41}}%
\TextAt{9}{320.695}{602.196}{\Character{61}}%
\TextAt{13}{335.893}{602.196}{N}%
\TextAt{21}{345.267}{600.337}{L}%
\TextAt{3}{351.026}{600.337}{\Character{106}}%
\TextAt{21}{353.379}{600.337}{K}%
\TextAt{9}{361.581}{602.196}{\Character{40}}%
\TextAt{14}{366.134}{602.196}{D}%
\TextAt{21}{376.075}{600.337}{L}%
\TextAt{3}{381.834}{600.337}{\Character{106}}%
\TextAt{21}{384.187}{600.337}{K}%
\TextAt{9}{392.389}{602.196}{\Character{41}}%
\TextAt{3}{396.942}{606.535}{\Character{0}}%
\TextAt{17}{403.528}{606.535}{1}%
\TextAt{8}{408.26}{602.196}{,}%
\TextAt{8}{417.432}{602.196}{and}%
\TextAt{9}{441.813}{602.196}{\Character{40}}%
\TextAt{13}{446.366}{602.196}{d}%
\TextAt{9}{452.449}{602.196}{\Character{40}}%
\TextAt{13}{457.001}{602.196}{\Character{11}}%
\TextAt{17}{464.446}{600.403}{1}%
\TextAt{13}{469.179}{602.196}{\Character{59}}%
\TextAt{13}{474.403}{602.196}{\Character{58}}%
\TextAt{13}{479.64}{602.196}{\Character{58}}%
\TextAt{13}{484.864}{602.196}{\Character{58}}%
\TextAt{13}{490.101}{602.196}{\Character{59}}%
\TextAt{13}{495.337}{602.196}{\Character{11}}%
\TextAt{21}{502.845}{600.403}{n}%
\TextAt{9}{508.482}{602.196}{\Character{41}\Character{41}}%
\TextAt{9}{523.641}{602.196}{\Character{61}}%
\TextAt{14}{79.226}{585.151}{d}%
\TextAt{21}{85.187}{583.291}{L}%
\TextAt{3}{90.946}{583.291}{\Character{106}}%
\TextAt{21}{93.298}{583.291}{K}%
\TextAt{8}{101.501}{585.151}{.}%
\TextAt{8}{116.43}{585.151}{One}%
\TextAt{8}{143.302}{585.151}{has}%
\TextAt{13}{166.336}{585.151}{d}%
\TextAt{9}{172.419}{585.151}{\Character{40}}%
\TextAt{13}{176.971}{585.151}{\Character{11}}%
\TextAt{17}{184.417}{583.357}{1}%
\TextAt{13}{189.149}{585.151}{\Character{59}}%
\TextAt{13}{194.373}{585.151}{\Character{58}}%
\TextAt{13}{199.61}{585.151}{\Character{58}}%
\TextAt{13}{204.834}{585.151}{\Character{58}}%
\TextAt{13}{210.071}{585.151}{\Character{59}}%
\TextAt{13}{215.307}{585.151}{\Character{11}}%
\TextAt{21}{222.815}{583.357}{n}%
\TextAt{9}{228.452}{585.151}{\Character{41}}%
\TextAt{9}{239.989}{585.151}{\Character{61}}%
\TextAt{8}{256.114}{585.151}{det}%
\TextAt{9}{272.373}{585.151}{\Character{40}\Character{40}}%
\TextAt{13}{281.477}{585.151}{\Character{27}}%
\TextAt{21}{288.131}{583.357}{i}%
\TextAt{13}{291.512}{585.151}{\Character{11}}%
\TextAt{21}{298.958}{583.357}{j}%
\TextAt{9}{303.34}{585.151}{\Character{41}\Character{41}}%
\TextAt{17}{312.444}{589.489}{2}%
\TextAt{8}{317.177}{585.151}{,}%
\TextAt{13}{327.026}{585.151}{d}%
\TextAt{9}{333.109}{585.151}{\Character{40}}%
\TextAt{13}{337.661}{585.151}{\Character{11}}%
\TextAt{3}{345.183}{589.489}{\Character{48}}%
\TextAt{17}{345.107}{582.195}{1}%
\TextAt{13}{349.839}{585.151}{\Character{59}}%
\TextAt{13}{355.063}{585.151}{\Character{58}}%
\TextAt{13}{360.3}{585.151}{\Character{58}}%
\TextAt{13}{365.536}{585.151}{\Character{58}}%
\TextAt{13}{370.761}{585.151}{\Character{59}}%
\TextAt{13}{375.997}{585.151}{\Character{11}}%
\TextAt{3}{383.582}{589.489}{\Character{48}}%
\TextAt{21}{383.506}{582.195}{n}%
\TextAt{9}{389.142}{585.151}{\Character{41}}%
\TextAt{9}{400.679}{585.151}{\Character{61}}%
\TextAt{8}{416.804}{585.151}{det}%
\TextAt{9}{433.063}{585.151}{\Character{40}\Character{40}}%
\TextAt{13}{442.168}{585.151}{\Character{27}}%
\TextAt{21}{448.821}{583.357}{i}%
\TextAt{13}{452.202}{585.151}{\Character{11}}%
\TextAt{3}{459.724}{589.489}{\Character{48}}%
\TextAt{21}{459.648}{582.195}{j}%
\TextAt{9}{464.03}{585.151}{\Character{41}\Character{41}}%
\TextAt{17}{473.135}{589.489}{2}%
\TextAt{8}{477.867}{585.151}{,}%
\TextAt{8}{487.717}{585.151}{for}%
\TextAt{13}{507.762}{585.151}{\Character{27}}%
\TextAt{21}{514.416}{583.357}{i}%
\NativeMathAt{524.799}{585.151}{11.955}{\in}%
\TextAt{8}{79.226}{568.105}{Hom}%
\TextAt{21}{103.615}{566.311}{K}%
\TextAt{9}{111.817}{568.105}{\Character{40}}%
\TextAt{13}{116.37}{568.105}{L\Character{59}}%
\Vector{q 0 G 0 g 0.478 w 57.578 -196.052 m 68.396 -196.052 l S Q}%
\TextAt{13}{129.578}{568.105}{K}%
\TextAt{9}{140.396}{568.105}{\Character{41}}%
\TextAt{8}{144.948}{568.105}{,}%
\TextAt{8}{153.41}{568.105}{and}%
\TextAt{8}{177.223}{568.105}{Tr}%
\TextAt{9}{190.231}{568.105}{\Character{40}}%
\TextAt{13}{194.783}{568.105}{\Character{11}}%
\TextAt{21}{202.229}{566.311}{i}%
\TextAt{13}{205.61}{568.105}{\Character{11}}%
\TextAt{3}{213.132}{572.443}{\Character{48}}%
\TextAt{21}{213.056}{565.149}{j}%
\TextAt{9}{217.438}{568.105}{\Character{41}}%
\TextAt{9}{227.086}{568.105}{\Character{61}}%
\TextAt{13}{241.316}{568.105}{\Character{14}}%
\TextAt{21}{246.501}{566.311}{ij}%
\TextAt{8}{253.766}{568.105}{.}%
\TextAt{8}{265.366}{568.105}{Then}%
\TextAt{13}{296.973}{568.105}{d}%
\TextAt{9}{303.056}{568.105}{\Character{40}}%
\TextAt{13}{307.608}{568.105}{\Character{11}}%
\TextAt{17}{315.054}{566.311}{1}%
\TextAt{13}{319.786}{568.105}{\Character{59}}%
\TextAt{13}{325.01}{568.105}{\Character{58}}%
\TextAt{13}{330.247}{568.105}{\Character{58}}%
\TextAt{13}{335.483}{568.105}{\Character{58}}%
\TextAt{13}{340.708}{568.105}{\Character{59}}%
\TextAt{13}{345.944}{568.105}{\Character{11}}%
\TextAt{21}{353.453}{566.311}{n}%
\TextAt{9}{359.089}{568.105}{\Character{41}}%
\TextAt{16}{367.014}{568.105}{\Character{1}}%
\TextAt{13}{373.707}{568.105}{d}%
\TextAt{9}{379.79}{568.105}{\Character{40}}%
\TextAt{13}{384.343}{568.105}{\Character{11}}%
\TextAt{3}{391.864}{572.443}{\Character{48}}%
\TextAt{17}{391.788}{565.149}{1}%
\TextAt{13}{396.52}{568.105}{\Character{59}}%
\TextAt{13}{401.744}{568.105}{\Character{58}}%
\TextAt{13}{406.981}{568.105}{\Character{58}}%
\TextAt{13}{412.217}{568.105}{\Character{58}}%
\TextAt{13}{417.442}{568.105}{\Character{59}}%
\TextAt{13}{422.678}{568.105}{\Character{11}}%
\TextAt{3}{430.263}{572.443}{\Character{48}}%
\TextAt{21}{430.187}{565.149}{n}%
\TextAt{9}{435.823}{568.105}{\Character{41}}%
\TextAt{9}{445.471}{568.105}{\Character{61}}%
\TextAt{9}{459.686}{568.105}{1}%
\TextAt{8}{465.555}{568.105}{.}%
\TextAt{8}{477.154}{568.105}{Combining}%
\TextAt{8}{79.226}{551.059}{these}%
\TextAt{8}{109.198}{551.059}{yields}%
\TextAt{14}{144.651}{522.733}{d}%
\TextAt{3}{150.612}{527.681}{\Character{0}}%
\TextAt{17}{157.198}{527.681}{1}%
\TextAt{21}{150.612}{518.96}{L}%
\TextAt{3}{156.371}{518.96}{\Character{106}}%
\TextAt{21}{158.723}{518.96}{K}%
\TextAt{9}{170.247}{522.733}{\Character{61}}%
\TextAt{9}{182.657}{522.733}{\Character{40}}%
\TextAt{13}{187.225}{522.733}{d}%
\TextAt{9}{193.307}{522.733}{\Character{40}}%
\TextAt{13}{197.86}{522.733}{\Character{11}}%
\TextAt{17}{205.305}{520.94}{1}%
\TextAt{13}{210.038}{522.733}{\Character{59}}%
\TextAt{13}{215.262}{522.733}{\Character{58}}%
\TextAt{13}{220.499}{522.733}{\Character{58}}%
\TextAt{13}{225.723}{522.733}{\Character{58}}%
\TextAt{13}{230.96}{522.733}{\Character{59}}%
\TextAt{13}{236.196}{522.733}{\Character{11}}%
\TextAt{21}{243.704}{520.94}{n}%
\TextAt{9}{249.34}{522.733}{\Character{41}}%
\TextAt{3}{253.893}{527.669}{\Character{0}}%
\TextAt{17}{260.479}{527.669}{1}%
\TextAt{9}{265.211}{522.733}{\Character{41}}%
\TextAt{9}{273.066}{522.733}{\Character{61}}%
\TextAt{9}{285.487}{522.733}{\Character{40}}%
\TextAt{13}{290.062}{522.733}{d}%
\TextAt{9}{296.145}{522.733}{\Character{40}}%
\TextAt{13}{300.697}{522.733}{\Character{11}}%
\TextAt{3}{308.219}{527.669}{\Character{48}}%
\TextAt{17}{308.143}{519.777}{1}%
\TextAt{13}{312.875}{522.733}{\Character{59}}%
\TextAt{13}{318.099}{522.733}{\Character{58}}%
\TextAt{13}{323.336}{522.733}{\Character{58}}%
\TextAt{13}{328.572}{522.733}{\Character{58}}%
\TextAt{13}{333.797}{522.733}{\Character{59}}%
\TextAt{13}{339.033}{522.733}{\Character{11}}%
\TextAt{3}{346.618}{527.669}{\Character{48}}%
\TextAt{21}{346.542}{519.777}{n}%
\TextAt{9}{352.178}{522.733}{\Character{41}\Character{41}}%
\TextAt{9}{364.576}{522.733}{\Character{61}}%
\TextAt{9}{376.997}{522.733}{\Character{40}}%
\TextAt{13}{381.581}{522.733}{d}%
\TextAt{9}{387.664}{522.733}{\Character{40}}%
\TextAt{13}{392.216}{522.733}{\Character{12}}%
\TextAt{13}{399.485}{522.733}{\Character{11}}%
\TextAt{17}{406.933}{520.94}{1}%
\TextAt{13}{411.666}{522.733}{\Character{59}}%
\TextAt{13}{416.89}{522.733}{\Character{58}}%
\TextAt{13}{422.127}{522.733}{\Character{58}}%
\TextAt{13}{427.351}{522.733}{\Character{58}}%
\TextAt{13}{432.588}{522.733}{\Character{59}}%
\TextAt{13}{437.824}{522.733}{\Character{12}}%
\TextAt{13}{445.093}{522.733}{\Character{11}}%
\TextAt{21}{452.603}{520.94}{n}%
\TextAt{9}{458.24}{522.733}{\Character{41}\Character{41}}%
\TextAt{9}{170.247}{501.051}{\Character{61}}%
\TextAt{13}{182.672}{501.051}{N}%
\TextAt{21}{192.046}{499.191}{L}%
\TextAt{3}{197.805}{499.191}{\Character{106}}%
\TextAt{21}{200.157}{499.191}{K}%
\TextAt{9}{208.36}{501.051}{\Character{40}}%
\TextAt{14}{212.912}{501.051}{D}%
\TextAt{21}{222.854}{499.191}{L}%
\TextAt{3}{228.613}{499.191}{\Character{106}}%
\TextAt{21}{230.965}{499.191}{K}%
\TextAt{9}{239.168}{501.051}{\Character{41}}%
\TextAt{3}{243.72}{505.987}{\Character{0}}%
\TextAt{17}{250.307}{505.987}{2}%
\TextAt{14}{255.039}{501.051}{d}%
\TextAt{21}{261}{499.191}{L}%
\TextAt{3}{266.759}{499.191}{\Character{106}}%
\TextAt{21}{269.111}{499.191}{K}%
\TextAt{8}{79.226}{472.725}{and}%
\TextAt{8}{101.989}{472.725}{hence}%
\TextAt{13}{134.495}{472.725}{N}%
\TextAt{21}{143.869}{470.865}{L}%
\TextAt{3}{149.628}{470.865}{\Character{106}}%
\TextAt{21}{151.98}{470.865}{K}%
\TextAt{9}{160.183}{472.725}{\Character{40}}%
\TextAt{14}{164.735}{472.725}{D}%
\TextAt{21}{174.677}{470.865}{L}%
\TextAt{3}{180.436}{470.865}{\Character{106}}%
\TextAt{21}{182.788}{470.865}{K}%
\TextAt{9}{190.991}{472.725}{\Character{41}}%
\TextAt{9}{198.846}{472.725}{\Character{61}}%
\TextAt{14}{211.289}{472.725}{d}%
\TextAt{21}{217.25}{470.865}{L}%
\TextAt{3}{223.009}{470.865}{\Character{106}}%
\TextAt{21}{225.361}{470.865}{K}%
\TextAt{8}{233.564}{472.725}{.}%
\Vector{q 0 G 0 g 0.399 w 452.334 -301.275 m 452.334 -293.205 l S Q}%
\Vector{q 0 G 0 g 0.399 w 452.534 -293.404 m 459.707 -293.404 l S Q}%
\Vector{q 0 G 0 g 0.399 w 452.534 -301.076 m 459.707 -301.076 l S Q}%
\Vector{q 0 G 0 g 0.399 w 459.906 -301.275 m 459.906 -293.205 l S Q}%
\TextAt{12}{79.226}{438.132}{Corollary}%
\TextAt{12}{139.014}{438.132}{(2.10).}%
\TextAt{8}{181.06}{438.132}{For}%
\TextAt{8}{202.03}{438.132}{a}%
\TextAt{8}{211.785}{438.132}{tower}%
\TextAt{8}{243.658}{438.132}{of}%
\TextAt{8}{256.988}{438.132}{fields}%
\TextAt{13}{286.959}{438.132}{K}%
\NativeMathAt{301.098}{438.132}{11.955}{\subseteq}%
\TextAt{13}{313.717}{438.132}{L}%
\NativeMathAt{325.002}{438.132}{11.955}{\subseteq}%
\TextAt{13}{337.622}{438.132}{M}%
\TextAt{8}{350.195}{438.132}{,}%
\TextAt{8}{357.345}{438.132}{one}%
\TextAt{8}{378.804}{438.132}{has}%
\TextAt{14}{240.45}{409.807}{d}%
\TextAt{21}{246.41}{407.947}{M}%
\TextAt{3}{255.392}{407.947}{\Character{106}}%
\TextAt{21}{257.744}{407.947}{K}%
\TextAt{9}{269.268}{409.807}{\Character{61}}%
\TextAt{14}{281.693}{409.807}{d}%
\TextAt{17}{287.654}{415.917}{\Character{91}}%
\TextAt{21}{290.006}{415.917}{M}%
\TextAt{17}{298.988}{415.917}{\Character{58}}%
\TextAt{21}{301.34}{415.917}{L}%
\TextAt{17}{307.099}{415.917}{\Character{93}}%
\TextAt{21}{287.654}{406.034}{L}%
\TextAt{3}{293.413}{406.034}{\Character{106}}%
\TextAt{21}{295.765}{406.034}{K}%
\TextAt{13}{309.95}{409.807}{N}%
\TextAt{21}{319.324}{407.947}{L}%
\TextAt{3}{325.083}{407.947}{\Character{106}}%
\TextAt{21}{327.435}{407.947}{K}%
\TextAt{9}{335.638}{409.807}{\Character{40}}%
\TextAt{14}{340.19}{409.807}{d}%
\TextAt{21}{346.15}{407.947}{M}%
\TextAt{3}{355.132}{407.947}{\Character{106}}%
\TextAt{21}{357.484}{407.947}{L}%
\TextAt{9}{363.742}{409.807}{\Character{41}}%
\TextAt{13}{368.294}{409.807}{\Character{58}}%
\TextAt{19}{79.226}{363.934}{Proof.}%
\TextAt{8}{116}{363.934}{Applying}%
\TextAt{8}{166.726}{363.934}{to}%
\TextAt{14}{181.036}{363.934}{D}%
\TextAt{21}{190.978}{362.074}{M}%
\TextAt{3}{199.96}{362.074}{\Character{106}}%
\TextAt{21}{202.312}{362.074}{K}%
\TextAt{9}{213.836}{363.934}{\Character{61}}%
\TextAt{14}{226.261}{363.934}{D}%
\TextAt{21}{236.203}{362.074}{M}%
\TextAt{3}{245.184}{362.074}{\Character{106}}%
\TextAt{21}{247.537}{362.074}{L}%
\TextAt{14}{253.794}{363.934}{D}%
\TextAt{21}{263.736}{362.074}{L}%
\TextAt{3}{269.495}{362.074}{\Character{106}}%
\TextAt{21}{271.847}{362.074}{K}%
\TextAt{8}{283.947}{363.934}{the}%
\TextAt{8}{304.103}{363.934}{norm}%
\TextAt{13}{334.673}{363.934}{N}%
\TextAt{21}{344.047}{362.074}{M}%
\TextAt{3}{353.028}{362.074}{\Character{106}}%
\TextAt{21}{355.381}{362.074}{K}%
\TextAt{9}{366.904}{363.934}{\Character{61}}%
\TextAt{13}{379.33}{363.934}{N}%
\TextAt{21}{388.703}{362.074}{L}%
\TextAt{3}{394.462}{362.074}{\Character{106}}%
\TextAt{21}{396.815}{362.074}{K}%
\TextAt{16}{407.674}{363.934}{\Character{14}}%
\TextAt{13}{416.308}{363.934}{N}%
\TextAt{21}{425.682}{362.074}{M}%
\TextAt{3}{434.664}{362.074}{\Character{106}}%
\TextAt{21}{437.016}{362.074}{L}%
\TextAt{8}{443.273}{363.934}{,}%
\TextAt{8}{450.423}{363.934}{(1.6)}%
\TextAt{8}{478.398}{363.934}{gives}%
\TextAt{14}{169.89}{335.608}{d}%
\TextAt{21}{175.85}{333.749}{M}%
\TextAt{3}{184.832}{333.749}{\Character{106}}%
\TextAt{21}{187.184}{333.749}{K}%
\TextAt{9}{198.708}{335.608}{\Character{61}}%
\TextAt{13}{211.133}{335.608}{N}%
\TextAt{21}{220.507}{333.749}{L}%
\TextAt{3}{226.266}{333.749}{\Character{106}}%
\TextAt{21}{228.618}{333.749}{K}%
\TextAt{9}{236.821}{335.608}{\Character{40}}%
\TextAt{14}{241.373}{335.608}{d}%
\TextAt{21}{247.334}{333.749}{M}%
\TextAt{3}{256.316}{333.749}{\Character{106}}%
\TextAt{21}{258.668}{333.749}{L}%
\TextAt{9}{264.925}{335.608}{\Character{41}}%
\TextAt{13}{269.477}{335.608}{N}%
\TextAt{21}{278.851}{333.749}{L}%
\TextAt{3}{284.61}{333.749}{\Character{106}}%
\TextAt{21}{286.963}{333.749}{K}%
\TextAt{9}{295.165}{335.608}{\Character{40}}%
\TextAt{14}{299.718}{335.608}{D}%
\TextAt{17}{309.659}{341.719}{\Character{91}}%
\TextAt{21}{312.012}{341.719}{M}%
\TextAt{17}{320.993}{341.719}{\Character{58}}%
\TextAt{21}{323.346}{341.719}{L}%
\TextAt{17}{329.105}{341.719}{\Character{93}}%
\TextAt{21}{309.659}{331.836}{L}%
\TextAt{3}{315.418}{331.836}{\Character{106}}%
\TextAt{21}{317.771}{331.836}{K}%
\TextAt{9}{331.955}{335.608}{\Character{41}}%
\TextAt{9}{339.81}{335.608}{\Character{61}}%
\TextAt{13}{352.254}{335.608}{N}%
\TextAt{21}{361.628}{333.749}{L}%
\TextAt{3}{367.387}{333.749}{\Character{106}}%
\TextAt{21}{369.739}{333.749}{K}%
\TextAt{9}{377.942}{335.608}{\Character{40}}%
\TextAt{14}{382.494}{335.608}{d}%
\TextAt{21}{388.454}{333.749}{M}%
\TextAt{3}{397.436}{333.749}{\Character{106}}%
\TextAt{21}{399.788}{333.749}{L}%
\TextAt{9}{406.046}{335.608}{\Character{41}}%
\TextAt{14}{410.598}{335.608}{d}%
\TextAt{17}{416.558}{341.719}{\Character{91}}%
\TextAt{21}{418.911}{341.719}{M}%
\TextAt{17}{427.892}{341.719}{\Character{58}}%
\TextAt{21}{430.245}{341.719}{L}%
\TextAt{17}{436.004}{341.719}{\Character{93}}%
\TextAt{21}{416.558}{331.836}{L}%
\TextAt{3}{422.317}{331.836}{\Character{106}}%
\TextAt{21}{424.67}{331.836}{K}%
\TextAt{13}{438.854}{335.608}{\Character{58}}%
\Vector{q 0 G 0 g 0.399 w 452.334 -466.717 m 452.334 -458.647 l S Q}%
\Vector{q 0 G 0 g 0.399 w 452.534 -458.847 m 459.707 -458.847 l S Q}%
\Vector{q 0 G 0 g 0.399 w 452.534 -466.518 m 459.707 -466.518 l S Q}%
\Vector{q 0 G 0 g 0.399 w 459.906 -466.717 m 459.906 -458.647 l S Q}%
\TextAt{8}{97.159}{278.475}{Putting}%
\TextAt{14}{140.335}{278.475}{d}%
\TextAt{9}{149.768}{278.475}{\Character{61}}%
\TextAt{14}{162.346}{278.475}{d}%
\TextAt{21}{168.307}{276.615}{L}%
\TextAt{3}{174.066}{276.615}{\Character{106}}%
\TextAt{21}{176.418}{276.615}{K}%
\TextAt{8}{188.607}{278.475}{and}%
\TextAt{14}{211.459}{278.475}{d}%
\TextAt{27}{217.42}{276.681}{P}%
\TextAt{9}{227.987}{278.475}{\Character{61}}%
\TextAt{14}{240.564}{278.475}{d}%
\TextAt{21}{246.525}{276.615}{L}%
\TextAt{28}{252.284}{275.224}{P}%
\TextAt{3}{258.481}{276.615}{\Character{106}}%
\TextAt{21}{260.833}{276.615}{K}%
\TextAt{28}{267.975}{275.619}{p}%
\TextAt{8}{272.556}{278.475}{,}%
\TextAt{8}{279.816}{278.475}{and}%
\TextAt{8}{302.668}{278.475}{viewing}%
\TextAt{14}{345.354}{278.475}{d}%
\TextAt{27}{351.315}{276.681}{P}%
\TextAt{8}{362.395}{278.475}{also}%
\TextAt{8}{385.964}{278.475}{as}%
\TextAt{8}{400.424}{278.475}{the}%
\TextAt{8}{420.67}{278.475}{ideal}%
\TextAt{14}{448.722}{278.475}{d}%
\TextAt{27}{454.683}{276.681}{P}%
\TextAt{16}{464.494}{278.475}{\Character{92}}%
\TextAt{14}{475.181}{278.475}{o}%
\TextAt{8}{485.012}{278.475}{of}%
\TextAt{13}{498.432}{278.475}{K}%
\TextAt{8}{509.25}{278.475}{,}%
\TextAt{8}{516.51}{278.475}{the}%
\TextAt{8}{79.226}{261.429}{product}%
\TextAt{8}{123.126}{261.429}{formula}%
\TextAt{8}{166.057}{261.429}{(2.3)}%
\TextAt{8}{194.032}{261.429}{for}%
\TextAt{8}{211.917}{261.429}{the}%
\TextAt{8}{232.073}{261.429}{different,}%
\TextAt{8}{281.496}{261.429}{together}%
\TextAt{8}{327.678}{261.429}{with}%
\TextAt{8}{354.339}{261.429}{theorem}%
\TextAt{8}{399.864}{261.429}{(2.9),}%
\TextAt{8}{431.091}{261.429}{yields:}%
\TextAt{12}{79.226}{235.609}{Corollary}%
\TextAt{12}{139.014}{235.609}{(2.11).}%
\TextAt{14}{181.06}{235.609}{d}%
\TextAt{9}{190.342}{235.609}{\Character{61}}%
\TextAt{20}{202.767}{244.576}{\Character{81}}%
\TextAt{27}{214.058}{232.122}{P}%
\TextAt{14}{223.144}{235.609}{d}%
\TextAt{27}{229.105}{233.816}{P}%
\TextAt{8}{236.198}{235.609}{.}%
\TextAt{8}{97.159}{209.79}{The}%
\TextAt{8}{121.63}{209.79}{extension}%
\TextAt{13}{173.813}{209.79}{L}%
\TextAt{16}{185.809}{209.79}{\Character{106}}%
\TextAt{13}{193.161}{209.79}{K}%
\TextAt{8}{208.293}{209.79}{is}%
\TextAt{8}{220.474}{209.79}{called}%
\TextAt{12}{254.055}{209.79}{unramified}%
\TextAt{8}{321.385}{209.79}{if}%
\TextAt{8}{332.526}{209.79}{all}%
\TextAt{8}{349.202}{209.79}{prime}%
\TextAt{8}{382.783}{209.79}{ideals}%
\TextAt{14}{415.778}{209.79}{p}%
\TextAt{8}{426.075}{209.79}{of}%
\TextAt{13}{439.823}{209.79}{K}%
\TextAt{8}{454.955}{209.79}{are}%
\TextAt{8}{474.883}{209.79}{unramified.}%
\TextAt{8}{79.226}{192.744}{This}%
\TextAt{8}{105.55}{192.744}{amounts}%
\TextAt{8}{152.378}{192.744}{to}%
\TextAt{8}{166.293}{192.744}{requiring}%
\TextAt{8}{215.655}{192.744}{that}%
\TextAt{8}{240.628}{192.744}{all}%
\TextAt{8}{256.492}{192.744}{primes}%
\TextAt{8}{293.875}{192.744}{of}%
\TextAt{13}{306.809}{192.744}{K}%
\TextAt{8}{321.129}{192.744}{be}%
\TextAt{8}{336.658}{192.744}{unramified.}%
\TextAt{8}{399.609}{192.744}{In}%
\TextAt{8}{413.835}{192.744}{fact,}%
\TextAt{8}{439.856}{192.744}{the}%
\TextAt{8}{459.617}{192.744}{infinite}%
\TextAt{8}{498.889}{192.744}{primes}%
\TextAt{8}{79.226}{175.698}{are}%
\TextAt{8}{98.737}{175.698}{always}%
\TextAt{8}{136.204}{175.698}{to}%
\TextAt{8}{150.515}{175.698}{be}%
\TextAt{8}{166.439}{175.698}{regarded}%
\TextAt{8}{214.571}{175.698}{as}%
\TextAt{8}{228.941}{175.698}{unramified}%
\TextAt{8}{287.473}{175.698}{because}%
\TextAt{13}{330.775}{175.698}{e}%
\TextAt{27}{336.2}{173.838}{P}%
\TextAt{3}{342.796}{173.838}{\Character{106}}%
\TextAt{27}{345.148}{173.838}{p}%
\TextAt{9}{352.956}{175.698}{\Character{61}}%
\TextAt{9}{365.365}{175.698}{1}%
\TextAt{8}{371.234}{175.698}{.}%
\TextAt{12}{79.226}{149.879}{Corollary}%
\TextAt{12}{138.288}{149.879}{(2.12).}%
\TextAt{8}{179.702}{149.879}{A}%
\TextAt{8}{191.743}{149.879}{prime}%
\TextAt{8}{224.275}{149.879}{ideal}%
\TextAt{14}{251.607}{149.879}{p}%
\TextAt{8}{260.856}{149.879}{of}%
\TextAt{13}{273.554}{149.879}{K}%
\TextAt{8}{287.638}{149.879}{is}%
\TextAt{8}{298.77}{149.879}{ramified}%
\TextAt{8}{343.664}{149.879}{in}%
\TextAt{13}{356.685}{149.879}{L}%
\TextAt{8}{367.916}{149.879}{if}%
\TextAt{8}{378.008}{149.879}{and}%
\TextAt{8}{400.139}{149.879}{only}%
\TextAt{8}{425.2}{149.879}{if}%
\TextAt{14}{435.292}{149.879}{p}%
\TextAt{16}{444.596}{149.879}{\Character{106}}%
\TextAt{14}{451.238}{149.879}{d}%
\TextAt{8}{457.198}{149.879}{.}%
\TextAt{8}{465.436}{149.879}{In}%
\TextAt{8}{479.426}{149.879}{particular,}%
\TextAt{8}{79.226}{132.833}{the}%
\TextAt{8}{99.383}{132.833}{extension}%
\TextAt{13}{151.148}{132.833}{L}%
\TextAt{16}{162.434}{132.833}{\Character{106}}%
\TextAt{13}{169.076}{132.833}{K}%
\TextAt{8}{183.791}{132.833}{is}%
\TextAt{8}{195.554}{132.833}{unramified}%
\TextAt{8}{254.087}{132.833}{if}%
\TextAt{8}{264.811}{132.833}{the}%
\TextAt{8}{284.967}{132.833}{discriminant}%
\TextAt{14}{352.347}{132.833}{d}%
\TextAt{9}{361.628}{132.833}{\Character{61}}%
\TextAt{9}{374.038}{132.833}{\Character{40}1\Character{41}}%
\TextAt{8}{389.011}{132.833}{.}%
\TextAt{8}{97.159}{107.013}{Combining}%
\TextAt{8}{158.078}{107.013}{this}%
\TextAt{8}{182.307}{107.013}{result}%
\TextAt{8}{216.291}{107.013}{with}%
\TextAt{8}{244.357}{107.013}{Minkowski}%
\TextAt{8}{304.033}{107.013}{theory}%
\TextAt{8}{342.19}{107.013}{leads}%
\TextAt{8}{372.922}{107.013}{to}%
\TextAt{8}{388.639}{107.013}{two}%
\TextAt{8}{412.162}{107.013}{important}%
\TextAt{8}{468.861}{107.013}{theorems}%
\TextAt{8}{520.408}{107.013}{on}%
\TextAt{8}{79.226}{89.967}{unramified}%
\TextAt{8}{138.919}{89.967}{extensions}%
\TextAt{8}{196.461}{89.967}{of}%
\TextAt{8}{210.951}{89.967}{number}%
\TextAt{8}{254.708}{89.967}{fields}%
\TextAt{8}{285.84}{89.967}{which}%
\TextAt{8}{320.487}{89.967}{belong}%
\TextAt{8}{359.044}{89.967}{to}%
\TextAt{8}{374.515}{89.967}{the}%
\TextAt{8}{395.832}{89.967}{classical}%
\TextAt{8}{441.992}{89.967}{body}%
\TextAt{8}{472.742}{89.967}{of}%
\TextAt{8}{487.232}{89.967}{algebraic}%
\TextAt{8}{79.226}{72.921}{number}%
\TextAt{8}{121.822}{72.921}{theory.}%
\TextAt{8}{162.143}{72.921}{The}%
\TextAt{8}{186.197}{72.921}{first}%
\TextAt{8}{210.323}{72.921}{of}%
\TextAt{8}{223.653}{72.921}{these}%
\TextAt{8}{253.624}{72.921}{results}%
\TextAt{8}{290.817}{72.921}{is}%
\TextAt{8}{302.581}{72.921}{the}%
\TextAt{8}{322.737}{72.921}{following.}%
\TextAt{8}{297.211}{43.033}{184}%
\Destination{78.23}{809.815}{nav.page.184}%
\Destination{79.36}{487.755}{nav.print.202}%
\Destination{79.36}{452.155}{nav.number.3.2.10}%
\Destination{79.36}{249.655}{nav.number.3.2.11}%
\Destination{79.36}{163.855}{nav.number.3.2.12}%
\Hotlink{463.834}{724.855}{6.858}{11.047}{nav.ex.1.3.4}%
\Hotlink{390.329}{724.855}{5.22}{11.047}{nav.chapter.1}%
\Hotlink{407.423}{724.855}{6.858}{11.047}{nav.section.1.3}%
\Hotlink{301.658}{707.755}{5.22}{11.047}{nav.chapter.1}%
\Hotlink{312.745}{707.755}{30.936}{11.047}{nav.number.1.2.10}%
\Hotlink{449.923}{361.755}{25.078}{11.047}{nav.number.3.1.6}%
\Hotlink{165.557}{259.255}{25.078}{11.047}{nav.number.3.2.3}%
\Hotlink{399.364}{259.255}{25.078}{11.047}{nav.number.3.2.9}%
}

% Page 197 (body page 185)
\Page{%
\TextAt{2}{79.226}{796.943}{III.}%
\TextAt{2}{98.317}{796.943}{Riemann–Roch}%
\TextAt{2}{174.375}{796.943}{Theory}%
\Vector{q 0 G 0 g 0.299 w 7.226 18.729 m 460.769 18.729 l S Q}%
\TextAt{12}{79.226}{761.124}{Theorem}%
\TextAt{12}{136.019}{761.124}{(2.13).}%
\TextAt{8}{177.873}{761.124}{Let}%
\TextAt{13}{198.649}{761.124}{K}%
\TextAt{8}{213.172}{761.124}{be}%
\TextAt{8}{228.904}{761.124}{an}%
\TextAt{8}{244.97}{761.124}{algebraic}%
\TextAt{8}{294.212}{761.124}{number}%
\TextAt{8}{336.616}{761.124}{field}%
\TextAt{8}{361.78}{761.124}{and}%
\TextAt{8}{384.35}{761.124}{let}%
\TextAt{13}{401.062}{761.124}{S}%
\TextAt{8}{412.662}{761.124}{be}%
\TextAt{8}{428.394}{761.124}{a}%
\TextAt{8}{437.957}{761.124}{finite}%
\TextAt{8}{467.676}{761.124}{set}%
\TextAt{8}{485.751}{761.124}{of}%
\TextAt{8}{498.889}{761.124}{primes}%
\TextAt{8}{79.226}{744.078}{of}%
\TextAt{13}{92.877}{744.078}{K}%
\TextAt{8}{103.695}{744.078}{.}%
\TextAt{8}{113.105}{744.078}{Then}%
\TextAt{8}{143.984}{744.078}{there}%
\TextAt{8}{174.216}{744.078}{exist}%
\TextAt{8}{202.237}{744.078}{only}%
\TextAt{8}{228.25}{744.078}{finitely}%
\TextAt{8}{267.915}{744.078}{many}%
\TextAt{8}{300.108}{744.078}{extensions}%
\TextAt{13}{356.81}{744.078}{L}%
\TextAt{16}{368.642}{744.078}{\Character{106}}%
\TextAt{13}{375.83}{744.078}{K}%
\TextAt{8}{390.866}{744.078}{of}%
\TextAt{8}{404.517}{744.078}{given}%
\TextAt{8}{435.407}{744.078}{degree}%
\TextAt{13}{472.143}{744.078}{n}%
\TextAt{8}{483.349}{744.078}{which}%
\TextAt{8}{517.156}{744.078}{are}%
\TextAt{8}{79.226}{727.032}{unramified}%
\TextAt{8}{137.759}{727.032}{outside}%
\TextAt{8}{178.143}{727.032}{of}%
\TextAt{13}{191.473}{727.032}{S}%
\TextAt{8}{199.369}{727.032}{.}%
\TextAt{19}{79.226}{692.053}{Proof.}%
\TextAt{8}{116}{692.053}{If}%
\TextAt{13}{128.602}{692.053}{L}%
\TextAt{16}{141.438}{692.053}{\Character{106}}%
\TextAt{13}{149.63}{692.053}{K}%
\TextAt{8}{165.254}{692.053}{is}%
\TextAt{8}{177.928}{692.053}{an}%
\TextAt{8}{195.097}{692.053}{extension}%
\TextAt{8}{247.772}{692.053}{of}%
\TextAt{8}{262.012}{692.053}{degree}%
\TextAt{13}{299.337}{692.053}{n}%
\TextAt{8}{311.132}{692.053}{which}%
\TextAt{8}{345.528}{692.053}{is}%
\TextAt{8}{358.201}{692.053}{unramified}%
\TextAt{8}{417.643}{692.053}{outside}%
\TextAt{8}{458.938}{692.053}{of}%
\TextAt{13}{473.177}{692.053}{S}%
\TextAt{8}{481.073}{692.053}{,}%
\TextAt{8}{489.359}{692.053}{then,}%
\TextAt{8}{520.408}{692.053}{by}%
\TextAt{8}{79.226}{675.007}{(2.12)}%
\TextAt{8}{113.245}{675.007}{and}%
\TextAt{8}{136.194}{675.007}{(2.6),}%
\TextAt{8}{167.653}{675.007}{its}%
\TextAt{8}{184.158}{675.007}{discriminant}%
\TextAt{14}{251.724}{675.007}{d}%
\TextAt{21}{257.684}{673.147}{L}%
\TextAt{3}{263.443}{673.147}{\Character{106}}%
\TextAt{21}{265.795}{673.147}{K}%
\TextAt{8}{278.081}{675.007}{is}%
\TextAt{8}{290.031}{675.007}{one}%
\TextAt{8}{311.677}{675.007}{of}%
\TextAt{8}{325.193}{675.007}{the}%
\TextAt{8}{345.535}{675.007}{finite}%
\TextAt{8}{375.633}{675.007}{number}%
\TextAt{8}{418.415}{675.007}{of}%
\TextAt{8}{431.931}{675.007}{divisors}%
\TextAt{8}{474.845}{675.007}{of}%
\TextAt{8}{488.361}{675.007}{the}%
\TextAt{8}{508.704}{675.007}{ideal}%
\ErrataAt{79.22}{657.955}{453.64}{$\mathfrak a=\prod_{\mathfrak p\in S,\,\mathfrak p\nmid\infty}\mathfrak p^{\,n(1+v_{\mathfrak p}(n!))}$. It therefore suffices to show that there are only finitely many}%
\TextAt{8}{79.226}{640.915}{extensions}%
\TextAt{13}{136.342}{640.915}{L}%
\TextAt{16}{148.879}{640.915}{\Character{106}}%
\TextAt{13}{156.773}{640.915}{K}%
\TextAt{8}{172.223}{640.915}{of}%
\TextAt{8}{186.288}{640.915}{degree}%
\TextAt{13}{223.438}{640.915}{n}%
\TextAt{8}{235.058}{640.915}{with}%
\TextAt{8}{262.453}{640.915}{given}%
\TextAt{8}{293.757}{640.915}{discriminant.}%
\TextAt{8}{367.892}{640.915}{We}%
\TextAt{8}{388.771}{640.915}{may}%
\TextAt{8}{414.875}{640.915}{assume}%
\TextAt{8}{456.054}{640.915}{without}%
\TextAt{8}{500.365}{640.915}{loss}%
\TextAt{8}{523.337}{640.915}{of}%
\TextAt{8}{79.226}{623.869}{generality}%
\TextAt{8}{134.174}{623.869}{that}%
\TextAt{13}{160.501}{623.869}{K}%
\TextAt{9}{176.273}{623.869}{\Character{61}}%
\TextAt{18}{190.332}{623.869}{Q}%
\TextAt{8}{199.63}{623.869}{.}%
\TextAt{8}{210.954}{623.869}{For}%
\TextAt{8}{232.882}{623.869}{if}%
\TextAt{13}{244.564}{623.869}{L}%
\TextAt{16}{257.483}{623.869}{\Character{106}}%
\TextAt{13}{265.758}{623.869}{K}%
\TextAt{8}{281.431}{623.869}{is}%
\TextAt{8}{294.154}{623.869}{an}%
\TextAt{8}{311.371}{623.869}{extension}%
\TextAt{8}{364.095}{623.869}{of}%
\TextAt{8}{378.384}{623.869}{degree}%
\TextAt{13}{415.758}{623.869}{n}%
\TextAt{8}{427.601}{623.869}{with}%
\TextAt{8}{455.219}{623.869}{discriminant}%
\TextAt{14}{523.557}{623.869}{d}%
\TextAt{8}{529.518}{623.869}{,}%
\TextAt{8}{79.226}{606.823}{then}%
\TextAt{13}{105.916}{606.823}{L}%
\TextAt{16}{117.253}{606.823}{\Character{106}}%
\TextAt{18}{123.946}{606.823}{Q}%
\TextAt{8}{137.172}{606.823}{is}%
\TextAt{8}{148.966}{606.823}{an}%
\TextAt{8}{165.255}{606.823}{extension}%
\TextAt{8}{217.051}{606.823}{of}%
\TextAt{8}{230.411}{606.823}{degree}%
\TextAt{13}{266.856}{606.823}{m}%
\TextAt{9}{280.468}{606.823}{\Character{61}}%
\TextAt{13}{292.944}{606.823}{n}%
\TextAt{9}{299.932}{606.823}{\Character{91}}%
\TextAt{13}{303.184}{606.823}{K}%
\TextAt{9}{317.373}{606.823}{\Character{58}}%
\TextAt{18}{323.997}{606.823}{Q}%
\TextAt{9}{333.296}{606.823}{\Character{93}}%
\TextAt{8}{340.475}{606.823}{with}%
\TextAt{8}{367.165}{606.823}{discriminant}%
\TextAt{9}{434.574}{606.823}{\Character{40}}%
\TextAt{13}{439.127}{606.823}{d}%
\TextAt{9}{445.209}{606.823}{\Character{41}}%
\TextAt{9}{453.123}{606.823}{\Character{61}}%
\TextAt{14}{465.61}{606.823}{d}%
\TextAt{21}{471.571}{611.162}{n}%
\TextAt{21}{471.571}{603.271}{K}%
\TextAt{3}{479.275}{603.271}{\Character{106}}%
\TextAt{25}{481.628}{603.271}{Q}%
\TextAt{13}{488.325}{606.823}{N}%
\TextAt{21}{497.699}{604.963}{K}%
\TextAt{3}{505.403}{604.963}{\Character{106}}%
\TextAt{25}{507.755}{604.963}{Q}%
\TextAt{9}{514.453}{606.823}{\Character{40}}%
\TextAt{14}{519.005}{606.823}{d}%
\TextAt{9}{524.965}{606.823}{\Character{41}}%
\TextAt{8}{529.518}{606.823}{.}%
\ErrataAt{79.19}{589.777}{455.81}{Adjoining $\sqrt{-1}$ at most doubles the degree and adds no finite ramification away from $2$.}%
\ErrataAt{79.19}{572.731}{455.81}{The preceding bound gives finitely many resulting discriminants, and each resulting field}%
\ErrataAt{79.16}{555.701}{455.84}{has finitely many subfields. It remains to count fields $K\mid\mathbb Q$ containing $\sqrt{-1}$ at fixed degree}%
\ErrataAt{79.16}{538.655}{455.84}{and discriminant. Its embeddings $\tau:K\to\mathbb C$ are all complex. Choose one, $\tau_0$. In the Minkowski space}%
\TextAt{13}{266.55}{500.655}{K}%
\TextAt{25}{276.509}{498.862}{R}%
\TextAt{9}{286.085}{500.655}{\Character{61}}%
\TextAt{20}{298.51}{521.098}{\Character{34}}%
\TextAt{20}{305.484}{512.012}{\Character{89}}%
\TextAt{21}{310.783}{486.906}{\Character{28}}%
\TextAt{18}{322.753}{500.655}{C}%
\TextAt{20}{331.387}{521.098}{\Character{35}}%
\TextAt{17}{338.361}{518.422}{\Character{43}}%
\TextAt{8}{79.226}{463.793}{(see}%
\TextAt{8}{102.694}{463.793}{chap.}%
\TextAt{8}{134.885}{463.793}{I,}%
\TextAt{8}{146.255}{463.793}{§5)}%
\TextAt{8}{166.244}{463.793}{consider}%
\TextAt{8}{211.829}{463.793}{the}%
\TextAt{8}{231.985}{463.793}{convex,}%
\TextAt{8}{273.613}{463.793}{centrally}%
\TextAt{8}{321.745}{463.793}{symmetric}%
\TextAt{8}{378.711}{463.793}{subset}%
\TextAt{13}{110.18}{434.792}{X}%
\TextAt{9}{124.156}{434.792}{\Character{61}}%
\TextAt{16}{136.581}{434.792}{\Character{102}}%
\TextAt{9}{142.559}{434.792}{\Character{40}}%
\TextAt{13}{147.111}{434.792}{z}%
\TextAt{21}{152.549}{432.999}{\Character{28}}%
\TextAt{9}{157.726}{434.792}{\Character{41}}%
\NativeMathAt{165.599}{434.792}{11.955}{\in}%
\TextAt{13}{176.89}{434.792}{K}%
\TextAt{25}{186.849}{432.999}{R}%
\TextAt{16}{196.424}{434.792}{\Character{106}}%
\TextAt{16}{203.047}{434.792}{\Character{106}}%
\TextAt{8}{208.379}{434.792}{Im}%
\TextAt{9}{222.355}{434.792}{\Character{40}}%
\TextAt{13}{226.907}{434.792}{z}%
\TextAt{21}{232.345}{432.999}{\Character{28}}%
\TextAt{24}{236.103}{431.892}{0}%
\TextAt{9}{240.752}{434.792}{\Character{41}}%
\TextAt{16}{245.305}{434.792}{\Character{106}}%
\TextAt{13}{251.946}{434.792}{\Character{60}}%
\TextAt{13}{264.356}{434.792}{C}%
\TextAt{20}{273.606}{445.597}{\Character{112}}%
\Vector{q 0 G 0 g 0.478 w 213.561 -328.164 m 226.285 -328.164 l S Q}%
\TextAt{16}{285.561}{434.792}{\Character{106}}%
\TextAt{13}{288.882}{434.792}{d}%
\TextAt{16}{294.964}{434.792}{\Character{106}}%
\TextAt{13}{298.285}{434.792}{\Character{59}}%
\TextAt{16}{307.427}{434.792}{\Character{106}}%
\TextAt{8}{312.74}{434.792}{Re}%
\TextAt{9}{326.56}{434.792}{\Character{40}}%
\TextAt{13}{331.113}{434.792}{z}%
\TextAt{21}{336.551}{432.999}{\Character{28}}%
\TextAt{24}{340.309}{431.892}{0}%
\TextAt{9}{344.958}{434.792}{\Character{41}}%
\TextAt{16}{349.51}{434.792}{\Character{106}}%
\TextAt{13}{356.152}{434.792}{\Character{60}}%
\TextAt{9}{368.577}{434.792}{1}%
\TextAt{13}{374.43}{434.792}{\Character{59}}%
\TextAt{16}{383.572}{434.792}{\Character{106}}%
\TextAt{13}{386.893}{434.792}{z}%
\TextAt{21}{392.331}{432.999}{\Character{28}}%
\TextAt{16}{397.507}{434.792}{\Character{106}}%
\TextAt{13}{404.149}{434.792}{\Character{60}}%
\TextAt{9}{416.574}{434.792}{1}%
\TextAt{8}{426.325}{434.792}{for}%
\TextAt{13}{444.21}{434.792}{\Character{28}}%
\NativeMathAt{453.95}{434.792}{11.955}{\neq}%
\TextAt{13}{466.375}{434.792}{\Character{28}}%
\TextAt{17}{471.459}{432.999}{0}%
\TextAt{13}{476.191}{434.792}{\Character{59}}%
\Vector{q 0 G 0 g 0.478 w 409.435 -332.387 m 415.854 -332.387 l S Q}%
\TextAt{13}{481.435}{434.792}{\Character{28}}%
\TextAt{17}{487.854}{432.999}{0}%
\TextAt{16}{492.586}{434.792}{\Character{103}}%
\TextAt{13}{498.564}{434.792}{\Character{59}}%
\TextAt{8}{79.226}{405.791}{where}%
\TextAt{13}{113.032}{405.791}{C}%
\TextAt{8}{126.159}{405.791}{is}%
\TextAt{8}{137.919}{405.791}{an}%
\TextAt{8}{154.174}{405.791}{arbitrarily}%
\TextAt{8}{210.443}{405.791}{big}%
\TextAt{8}{229.951}{405.791}{constant}%
\TextAt{8}{277.17}{405.791}{which}%
\TextAt{8}{310.652}{405.791}{depends}%
\TextAt{8}{355.899}{405.791}{only}%
\TextAt{8}{381.587}{405.791}{on}%
\TextAt{13}{397.842}{405.791}{n}%
\TextAt{8}{404.83}{405.791}{.}%
\TextAt{8}{413.277}{405.791}{For}%
\TextAt{8}{434.242}{405.791}{a}%
\TextAt{8}{443.994}{405.791}{convenient}%
\TextAt{8}{501.877}{405.791}{choice}%
\TextAt{8}{79.226}{388.745}{of}%
\TextAt{13}{92.556}{388.745}{C}%
\TextAt{8}{101.79}{388.745}{,}%
\TextAt{8}{108.939}{388.745}{the}%
\TextAt{8}{129.095}{388.745}{volume}%
\TextAt{8}{169.42}{388.745}{will}%
\TextAt{8}{191.525}{388.745}{satisfy}%
\TextAt{8}{227.054}{359.744}{vol}%
\TextAt{9}{242.345}{359.744}{\Character{40}}%
\TextAt{13}{246.897}{359.744}{X}%
\TextAt{9}{257.552}{359.744}{\Character{41}}%
\TextAt{13}{265.425}{359.744}{\Character{62}}%
\TextAt{9}{277.851}{359.744}{2}%
\TextAt{21}{283.704}{364.68}{n}%
\TextAt{20}{289.34}{370.549}{\Character{112}}%
\Vector{q 0 G 0 g 0.478 w 229.295 -403.212 m 242.02 -403.212 l S Q}%
\TextAt{16}{301.295}{359.744}{\Character{106}}%
\TextAt{13}{304.616}{359.744}{d}%
\TextAt{16}{310.699}{359.744}{\Character{106}}%
\TextAt{9}{317.341}{359.744}{\Character{61}}%
\TextAt{9}{329.75}{359.744}{2}%
\TextAt{21}{335.619}{364.68}{n}%
\TextAt{8}{343.248}{359.744}{vol}%
\TextAt{9}{358.538}{359.744}{\Character{40}}%
\TextAt{14}{363.091}{359.744}{o}%
\TextAt{21}{368.935}{357.95}{K}%
\TextAt{9}{377.138}{359.744}{\Character{41}}%
\TextAt{13}{381.69}{359.744}{\Character{59}}%
\TextAt{8}{79.226}{330.742}{where}%
\TextAt{8}{112.487}{330.742}{vol}%
\TextAt{9}{127.777}{330.742}{\Character{40}}%
\TextAt{14}{132.33}{330.742}{o}%
\TextAt{21}{138.174}{328.949}{K}%
\TextAt{9}{146.377}{330.742}{\Character{41}}%
\TextAt{8}{154.278}{330.742}{is}%
\TextAt{8}{165.493}{330.742}{the}%
\TextAt{8}{185.101}{330.742}{volume}%
\TextAt{8}{224.877}{330.742}{of}%
\TextAt{8}{237.658}{330.742}{a}%
\TextAt{8}{246.865}{330.742}{fundamental}%
\TextAt{8}{313.959}{330.742}{mesh}%
\TextAt{8}{343.381}{330.742}{of}%
\TextAt{8}{356.163}{330.742}{the}%
\TextAt{8}{375.771}{330.742}{lattice}%
\TextAt{13}{410.992}{330.742}{j}%
\TextAt{14}{416.507}{330.742}{o}%
\TextAt{21}{422.352}{328.949}{K}%
\TextAt{8}{433.903}{330.742}{in}%
\TextAt{13}{447.007}{330.742}{K}%
\TextAt{25}{456.966}{328.949}{R}%
\TextAt{8}{463.221}{330.742}{—see}%
\TextAt{8}{493.289}{330.742}{chap.}%
\TextAt{8}{525.297}{330.742}{I,}%
\TextAt{8}{79.226}{313.696}{(5.2).}%
\TextAt{8}{111.713}{313.696}{By}%
\TextAt{8}{129.958}{313.696}{Minkowski’s}%
\TextAt{8}{195.975}{313.696}{lattice}%
\TextAt{8}{231.627}{313.696}{point}%
\TextAt{8}{262.078}{313.696}{theorem}%
\TextAt{8}{307.485}{313.696}{(chap.}%
\TextAt{8}{344.191}{313.696}{I,}%
\TextAt{8}{355.443}{313.696}{(4.4)),}%
\TextAt{8}{391.13}{313.696}{we}%
\TextAt{8}{408.239}{313.696}{thus}%
\TextAt{8}{433.872}{313.696}{find}%
\TextAt{13}{457.162}{313.696}{\Character{11}}%
\NativeMathAt{468.004}{313.696}{11.955}{\in}%
\TextAt{14}{479.295}{313.696}{o}%
\TextAt{21}{485.139}{311.903}{K}%
\TextAt{8}{493.342}{313.696}{,}%
\TextAt{13}{500.397}{313.696}{\Character{11}}%
\NativeMathAt{511.239}{313.696}{11.955}{\neq}%
\TextAt{9}{523.649}{313.696}{0}%
\TextAt{8}{529.518}{313.696}{,}%
\TextAt{8}{79.226}{296.65}{such}%
\TextAt{8}{105.623}{296.65}{that}%
\TextAt{13}{130.992}{296.65}{j}%
\TextAt{13}{136.503}{296.65}{\Character{11}}%
\TextAt{9}{147.35}{296.65}{\Character{61}}%
\TextAt{9}{159.76}{296.65}{\Character{40}}%
\TextAt{13}{164.328}{296.65}{\Character{28}}%
\TextAt{13}{170.736}{296.65}{\Character{11}}%
\TextAt{9}{178.269}{296.65}{\Character{41}}%
\NativeMathAt{186.142}{296.65}{11.955}{\in}%
\TextAt{13}{197.433}{296.65}{X}%
\TextAt{8}{208.088}{296.65}{,}%
\TextAt{8}{215.237}{296.65}{that}%
\TextAt{8}{240.606}{296.65}{is,}%
\TextAt{16}{146.398}{267.649}{\Character{106}}%
\TextAt{8}{151.712}{267.649}{Im}%
\TextAt{9}{165.687}{267.649}{\Character{40}}%
\TextAt{13}{170.24}{267.649}{\Character{28}}%
\TextAt{17}{175.323}{265.856}{0}%
\TextAt{13}{180.056}{267.649}{\Character{11}}%
\TextAt{9}{187.577}{267.649}{\Character{41}}%
\TextAt{16}{192.13}{267.649}{\Character{106}}%
\TextAt{13}{198.771}{267.649}{\Character{60}}%
\TextAt{13}{211.181}{267.649}{C}%
\TextAt{20}{220.43}{278.455}{\Character{112}}%
\Vector{q 0 G 0 g 0.478 w 160.386 -495.306 m 173.11 -495.306 l S Q}%
\TextAt{16}{232.386}{267.649}{\Character{106}}%
\TextAt{13}{235.707}{267.649}{d}%
\TextAt{16}{241.789}{267.649}{\Character{106}}%
\TextAt{13}{245.11}{267.649}{\Character{59}}%
\TextAt{16}{262.309}{267.649}{\Character{106}}%
\TextAt{8}{267.623}{267.649}{Re}%
\TextAt{9}{281.443}{267.649}{\Character{40}}%
\TextAt{13}{285.995}{267.649}{\Character{28}}%
\TextAt{17}{291.079}{265.856}{0}%
\TextAt{13}{295.811}{267.649}{\Character{11}}%
\TextAt{9}{303.333}{267.649}{\Character{41}}%
\TextAt{16}{307.885}{267.649}{\Character{106}}%
\TextAt{13}{314.527}{267.649}{\Character{60}}%
\TextAt{9}{326.953}{267.649}{1}%
\TextAt{13}{332.806}{267.649}{\Character{59}}%
\TextAt{16}{350.005}{267.649}{\Character{106}}%
\TextAt{13}{353.326}{267.649}{\Character{28}}%
\TextAt{13}{359.734}{267.649}{\Character{11}}%
\TextAt{16}{367.267}{267.649}{\Character{106}}%
\TextAt{13}{373.908}{267.649}{\Character{60}}%
\TextAt{9}{386.334}{267.649}{1}%
\TextAt{8}{396.084}{267.649}{for}%
\TextAt{13}{413.969}{267.649}{\Character{28}}%
\NativeMathAt{423.709}{267.649}{11.955}{\neq}%
\TextAt{13}{436.134}{267.649}{\Character{28}}%
\TextAt{17}{441.218}{265.856}{0}%
\TextAt{13}{445.95}{267.649}{\Character{59}}%
\Vector{q 0 G 0 g 0.478 w 379.194 -499.53 m 385.613 -499.53 l S Q}%
\TextAt{13}{451.194}{267.649}{\Character{28}}%
\TextAt{17}{457.613}{265.856}{0}%
\TextAt{13}{462.346}{267.649}{\Character{58}}%
\TextAt{8}{517.802}{267.649}{(*)}%
\TextAt{8}{79.226}{238.648}{This}%
\TextAt{13}{104.736}{238.648}{\Character{11}}%
\TextAt{8}{114.946}{238.648}{is}%
\TextAt{8}{125.5}{238.648}{a}%
\TextAt{8}{134.045}{238.648}{primitive}%
\TextAt{8}{182.916}{238.648}{element}%
\TextAt{8}{224.948}{238.648}{of}%
\TextAt{13}{237.068}{238.648}{K}%
\TextAt{8}{247.886}{238.648}{,}%
\TextAt{8}{254.067}{238.648}{i.e.,}%
\TextAt{8}{275.204}{238.648}{one}%
\TextAt{8}{295.454}{238.648}{has}%
\TextAt{13}{315.118}{238.648}{K}%
\TextAt{9}{329.257}{238.648}{\Character{61}}%
\TextAt{18}{341.682}{238.648}{Q}%
\TextAt{9}{350.981}{238.648}{\Character{40}}%
\TextAt{13}{355.533}{238.648}{\Character{11}}%
\TextAt{9}{363.055}{238.648}{\Character{41}}%
\TextAt{8}{367.607}{238.648}{.}%
\TextAt{8}{375.652}{238.648}{Indeed,}%
\TextAt{16}{415.965}{238.648}{\Character{106}}%
\TextAt{13}{419.286}{238.648}{N}%
\TextAt{21}{428.66}{236.788}{K}%
\TextAt{3}{436.365}{236.788}{\Character{106}}%
\TextAt{25}{438.717}{236.788}{Q}%
\TextAt{9}{445.414}{238.648}{\Character{40}}%
\TextAt{13}{449.966}{238.648}{\Character{11}}%
\TextAt{9}{457.488}{238.648}{\Character{41}}%
\TextAt{16}{462.041}{238.648}{\Character{106}}%
\TextAt{9}{468.682}{238.648}{\Character{61}}%
\TextAt{20}{481.108}{247.615}{\Character{81}}%
\TextAt{21}{492.399}{235.161}{\Character{28}}%
\TextAt{16}{499.568}{238.648}{\Character{106}}%
\TextAt{13}{502.889}{238.648}{\Character{28}}%
\TextAt{13}{509.297}{238.648}{\Character{11}}%
\TextAt{16}{516.829}{238.648}{\Character{106}}%
\NativeMathAt{523.452}{238.648}{11.955}{\geq}%
\ErrataAt{79.19}{221.602}{455.81}{$1$ implies $|\tau_0\alpha|\ge1$. Thus $\operatorname{Im}(\tau_0\alpha)\ne0$, so the conjugates $\tau_0\alpha$ and $\overline{\tau_0}\alpha$ have to}%
\TextAt{8}{79.226}{204.556}{be}%
\TextAt{8}{95.484}{204.556}{distinct.}%
\TextAt{8}{143.37}{204.556}{Since}%
\TextAt{16}{174.261}{204.556}{\Character{106}}%
\TextAt{13}{177.582}{204.556}{\Character{28}}%
\TextAt{13}{183.99}{204.556}{\Character{11}}%
\TextAt{16}{191.523}{204.556}{\Character{106}}%
\TextAt{13}{198.733}{204.556}{\Character{60}}%
\TextAt{9}{211.728}{204.556}{1}%
\TextAt{8}{221.812}{204.556}{for}%
\TextAt{13}{240.031}{204.556}{\Character{28}}%
\NativeMathAt{250.339}{204.556}{11.955}{\neq}%
\TextAt{13}{263.334}{204.556}{\Character{28}}%
\TextAt{17}{268.417}{202.763}{0}%
\TextAt{13}{273.149}{204.556}{\Character{59}}%
\Vector{q 0 G 0 g 0.478 w 206.394 -562.623 m 212.813 -562.623 l S Q}%
\TextAt{13}{278.394}{204.556}{\Character{28}}%
\TextAt{17}{284.813}{202.763}{0}%
\TextAt{8}{289.545}{204.556}{,}%
\TextAt{8}{297.111}{204.556}{one}%
\TextAt{8}{318.905}{204.556}{has}%
\TextAt{13}{340.112}{204.556}{\Character{28}}%
\TextAt{17}{345.196}{202.763}{0}%
\TextAt{13}{349.928}{204.556}{\Character{11}}%
\NativeMathAt{361.34}{204.556}{11.955}{\neq}%
\TextAt{13}{374.334}{204.556}{\Character{28}}%
\TextAt{13}{380.742}{204.556}{\Character{11}}%
\TextAt{8}{392.506}{204.556}{for}%
\TextAt{8}{410.725}{204.556}{all}%
\TextAt{13}{427.317}{204.556}{\Character{28}}%
\NativeMathAt{437.626}{204.556}{11.955}{\neq}%
\TextAt{13}{450.62}{204.556}{\Character{28}}%
\TextAt{17}{455.704}{202.763}{0}%
\TextAt{8}{460.436}{204.556}{.}%
\TextAt{8}{469.886}{204.556}{This}%
\TextAt{8}{496.94}{204.556}{implies}%
\TextAt{8}{297.211}{43.033}{185}%
\Destination{522.36}{284.695}{nav.AMS.23}%
\Destination{78.23}{809.815}{nav.page.185}%
\Destination{79.36}{776.155}{nav.print.203}%
\Destination{79.36}{775.155}{nav.number.3.2.13}%
\Hotlink{134.385}{461.555}{5.22}{11.047}{nav.chapter.1}%
\Hotlink{151.434}{461.555}{6.858}{11.047}{nav.section.1.5}%
\Hotlink{524.797}{328.555}{5.22}{11.047}{nav.chapter.1}%
\Hotlink{343.691}{311.455}{5.22}{11.047}{nav.chapter.1}%
\Hotlink{78.726}{672.855}{30.936}{11.047}{nav.number.3.2.12}%
\Hotlink{135.694}{672.855}{25.078}{11.047}{nav.number.3.2.6}%
\Hotlink{78.726}{311.455}{25.078}{11.047}{nav.number.1.5.2}%
\Hotlink{221.837}{166.155}{21.483}{9.572}{nav.number.3.2.6}%
\Destination{79.226}{671.955}{errata.22}%
\Destination{79.226}{606.777}{errata.91}%
\Destination{79.226}{238.602}{errata.92}%
\Hotlink{354.97}{311.45}{23.35}{11.047}{nav.number.1.4.4}%
}

% Page 198 (body page 186)
\Page{%
\TextAt{2}{79.226}{796.943}{III.}%
\TextAt{2}{98.317}{796.943}{Riemann–Roch}%
\TextAt{2}{174.375}{796.943}{Theory}%
\Vector{q 0 G 0 g 0.299 w 7.226 18.729 m 460.769 18.729 l S Q}%
\TextAt{13}{79.226}{761.124}{K}%
\TextAt{9}{93.572}{761.124}{\Character{61}}%
\TextAt{18}{106.205}{761.124}{Q}%
\TextAt{9}{115.503}{761.124}{\Character{40}}%
\TextAt{13}{120.056}{761.124}{\Character{11}}%
\TextAt{9}{127.577}{761.124}{\Character{41}}%
\TextAt{8}{132.13}{761.124}{,}%
\TextAt{8}{139.431}{761.124}{because}%
\TextAt{8}{182.854}{761.124}{if}%
\TextAt{18}{193.7}{761.124}{Q}%
\TextAt{9}{202.998}{761.124}{\Character{40}}%
\TextAt{13}{207.55}{761.124}{\Character{11}}%
\TextAt{9}{215.072}{761.124}{\Character{41}}%
\TextAt{18}{223.153}{761.124}{\Character{40}}%
\TextAt{13}{235.979}{761.124}{K}%
\TextAt{8}{250.816}{761.124}{then}%
\TextAt{8}{277.598}{761.124}{the}%
\TextAt{8}{297.876}{761.124}{restriction}%
\TextAt{13}{353.995}{761.124}{\Character{28}}%
\TextAt{17}{359.079}{759.33}{0}%
\TextAt{16}{363.811}{761.124}{\Character{106}}%
\TextAt{25}{367.132}{759.264}{Q}%
\TextAt{17}{373.331}{759.264}{\Character{40}}%
\TextAt{21}{376.625}{759.264}{\Character{11}}%
\TextAt{17}{382.064}{759.264}{\Character{41}}%
\TextAt{8}{389.874}{761.124}{would}%
\TextAt{8}{424.14}{761.124}{admit}%
\TextAt{8}{458.082}{761.124}{an}%
\TextAt{8}{474.463}{761.124}{extension}%
\TextAt{13}{526.351}{761.124}{\Character{28}}%
\TextAt{8}{79.226}{744.078}{different}%
\TextAt{8}{125.397}{744.078}{from}%
\TextAt{13}{153.037}{744.078}{\Character{28}}%
\TextAt{17}{158.121}{742.284}{0}%
\TextAt{8}{162.853}{744.078}{,}%
\TextAt{8}{170.003}{744.078}{contradicting}%
\TextAt{13}{241.231}{744.078}{\Character{28}}%
\TextAt{17}{246.315}{742.284}{0}%
\TextAt{13}{251.047}{744.078}{\Character{11}}%
\NativeMathAt{261.89}{744.078}{11.955}{\neq}%
\TextAt{13}{274.315}{744.078}{\Character{28}}%
\TextAt{13}{280.723}{744.078}{\Character{11}}%
\TextAt{8}{288.256}{744.078}{.}%
\TextAt{8}{97.159}{724.043}{Since}%
\TextAt{8}{127.719}{724.043}{the}%
\TextAt{8}{147.879}{724.043}{conjugates}%
\TextAt{13}{205.505}{724.043}{\Character{28}}%
\TextAt{13}{211.913}{724.043}{\Character{11}}%
\TextAt{8}{223.347}{724.043}{of}%
\TextAt{13}{236.679}{724.043}{\Character{11}}%
\TextAt{8}{248.102}{724.043}{are}%
\TextAt{8}{267.615}{724.043}{subject}%
\TextAt{8}{307.991}{724.043}{to}%
\TextAt{8}{322.304}{724.043}{the}%
\TextAt{8}{342.464}{724.043}{conditions}%
\TextAt{8}{398.465}{724.043}{(*),}%
\TextAt{8}{420.585}{724.043}{which}%
\TextAt{8}{454.075}{724.043}{only}%
\TextAt{8}{479.769}{724.043}{depend}%
\TextAt{8}{520.408}{724.043}{on}%
\TextAt{13}{79.226}{706.997}{d}%
\TextAt{8}{89.775}{706.997}{and}%
\TextAt{13}{113.106}{706.997}{n}%
\TextAt{8}{120.093}{706.997}{,}%
\TextAt{8}{127.953}{706.997}{the}%
\TextAt{8}{148.678}{706.997}{coefficients}%
\TextAt{8}{208.485}{706.997}{of}%
\TextAt{8}{222.383}{706.997}{the}%
\TextAt{8}{243.108}{706.997}{minimal}%
\TextAt{8}{289.202}{706.997}{polynomial}%
\TextAt{8}{350.264}{706.997}{of}%
\TextAt{13}{364.162}{706.997}{\Character{11}}%
\TextAt{8}{376.15}{706.997}{are}%
\TextAt{8}{396.229}{706.997}{bounded}%
\TextAt{8}{444.594}{706.997}{once}%
\TextAt{13}{471.823}{706.997}{d}%
\TextAt{8}{482.371}{706.997}{and}%
\TextAt{13}{505.703}{706.997}{n}%
\TextAt{8}{517.156}{706.997}{are}%
\ErrataAt{79.19}{689.951}{455.81}{fixed. As the coefficients are integers, only finitely many such minimal polynomials occur.}%
\ErrataAt{79.19}{672.905}{455.81}{Every field under consideration is generated by a root of one of these polynomials, so only}%
\ErrataAt{79.19}{655.859}{455.81}{finitely many fields occur.}%
\Vector{q 0 G 0 g 0.399 w 452.334 -118.141 m 452.334 -110.071 l S Q}%
\Vector{q 0 G 0 g 0.399 w 452.534 -110.27 m 459.707 -110.27 l S Q}%
\Vector{q 0 G 0 g 0.399 w 452.534 -117.942 m 459.707 -117.942 l S Q}%
\Vector{q 0 G 0 g 0.399 w 459.906 -118.141 m 459.906 -110.071 l S Q}%
\TextAt{8}{97.159}{628.479}{The}%
\TextAt{8}{121.115}{628.479}{second}%
\TextAt{8}{158.796}{628.479}{theorem}%
\TextAt{8}{204.224}{628.479}{alluded}%
\TextAt{8}{245.097}{628.479}{to}%
\TextAt{8}{259.31}{628.479}{above}%
\TextAt{8}{292.389}{628.479}{is}%
\TextAt{8}{304.055}{628.479}{in}%
\TextAt{8}{317.611}{628.479}{fact}%
\TextAt{8}{340.599}{628.479}{a}%
\TextAt{8}{350.257}{628.479}{strengthening}%
\TextAt{8}{423.72}{628.479}{of}%
\TextAt{8}{436.953}{628.479}{the}%
\TextAt{8}{457.012}{628.479}{first.}%
\TextAt{8}{485.656}{628.479}{It}%
\TextAt{8}{498.231}{628.479}{follows}%
\TextAt{8}{79.226}{611.433}{from}%
\TextAt{8}{106.866}{611.433}{the}%
\TextAt{8}{127.023}{611.433}{following}%
\TextAt{19}{176.458}{611.433}{bound}%
\TextAt{19}{210.231}{611.433}{on}%
\TextAt{19}{226.968}{611.433}{the}%
\TextAt{19}{246.395}{611.433}{discriminant}%
\TextAt{8}{310.81}{611.433}{.}%
\TextAt{12}{79.226}{587.042}{Proposition}%
\TextAt{12}{150.697}{587.042}{(2.14).}%
\TextAt{8}{191.665}{587.042}{The}%
\TextAt{8}{214.64}{587.042}{discriminant}%
\TextAt{8}{280.941}{587.042}{of}%
\TextAt{8}{293.192}{587.042}{an}%
\TextAt{8}{308.373}{587.042}{algebraic}%
\TextAt{8}{356.729}{587.042}{number}%
\TextAt{8}{398.247}{587.042}{field}%
\TextAt{13}{422.525}{587.042}{K}%
\TextAt{8}{436.162}{587.042}{of}%
\TextAt{8}{448.413}{587.042}{degree}%
\TextAt{13}{483.75}{587.042}{n}%
\TextAt{8}{493.557}{587.042}{satisfies}%
\TextAt{16}{250.916}{560.192}{\Character{106}}%
\TextAt{13}{254.236}{560.192}{d}%
\TextAt{21}{260.319}{558.398}{K}%
\TextAt{16}{268.522}{560.192}{\Character{106}}%
\TextAt{17}{271.843}{565.128}{1\Character{47}2}%
\NativeMathAt{288.364}{560.192}{11.955}{\geq}%
\TextAt{13}{302.179}{568.279}{n}%
\TextAt{21}{309.167}{572.618}{n}%
\Vector{q 0 G 0 g 0.478 w 230.179 -210.82 m 242.803 -210.82 l S Q}%
\TextAt{13}{303.371}{551.991}{n}%
\TextAt{9}{310.359}{551.991}{\Character{33}}%
\TextAt{20}{317.991}{573.462}{\Character{16}}%
\TextAt{13}{326.326}{568.279}{\Character{25}}%
\Vector{q 0 G 0 g 0.478 w 254.326 -210.82 m 261.396 -210.82 l S Q}%
\TextAt{9}{326.935}{551.991}{4}%
\TextAt{20}{334.591}{573.462}{\Character{17}}%
\TextAt{21}{341.731}{570.785}{n}%
\TextAt{17}{346.869}{570.785}{\Character{47}2}%
\TextAt{13}{357.828}{560.192}{\Character{58}}%
\TextAt{19}{79.226}{520.868}{Proof.}%
\TextAt{8}{116}{520.868}{In}%
\TextAt{8}{130.012}{520.868}{Minkowski}%
\TextAt{8}{187.672}{520.868}{space}%
\TextAt{13}{218.337}{520.868}{K}%
\TextAt{25}{228.296}{519.075}{R}%
\TextAt{9}{237.872}{520.868}{\Character{61}}%
\TextAt{9}{250.282}{520.868}{\Character{91}}%
\TextAt{20}{253.549}{529.835}{\Character{81}}%
\TextAt{21}{264.84}{517.381}{\Character{28}}%
\TextAt{18}{272.009}{520.868}{C}%
\TextAt{9}{280.643}{520.868}{\Character{93}}%
\TextAt{17}{283.895}{525.207}{\Character{43}}%
\TextAt{8}{290.979}{520.868}{,}%
\TextAt{13}{297.641}{520.868}{\Character{28}}%
\NativeMathAt{307.381}{520.868}{11.955}{\in}%
\TextAt{8}{318.672}{520.868}{Hom}%
\TextAt{9}{343.06}{520.868}{\Character{40}}%
\TextAt{13}{347.613}{520.868}{K}%
\TextAt{13}{357.775}{520.868}{\Character{59}}%
\TextAt{18}{363.024}{520.868}{C}%
\TextAt{9}{371.658}{520.868}{\Character{41}}%
\TextAt{8}{376.211}{520.868}{,}%
\TextAt{8}{382.872}{520.868}{consider}%
\TextAt{8}{427.848}{520.868}{the}%
\TextAt{8}{447.395}{520.868}{convex,}%
\TextAt{8}{488.535}{520.868}{centrally}%
\TextAt{8}{79.226}{503.822}{symmetric}%
\TextAt{8}{136.193}{503.822}{subset}%
\TextAt{13}{225.29}{486.776}{X}%
\TextAt{21}{235.005}{484.983}{t}%
\TextAt{9}{241.882}{486.777}{\Character{61}}%
\TextAt{16}{254.308}{486.777}{\Character{102}}%
\TextAt{9}{260.285}{486.777}{\Character{40}}%
\TextAt{13}{264.837}{486.777}{z}%
\TextAt{21}{270.275}{484.983}{\Character{28}}%
\TextAt{9}{275.452}{486.777}{\Character{41}}%
\NativeMathAt{283.325}{486.777}{11.955}{\in}%
\TextAt{13}{294.616}{486.777}{K}%
\TextAt{25}{304.575}{484.983}{R}%
\TextAt{16}{314.15}{486.777}{\Character{106}}%
\TextAt{20}{320.792}{498.134}{\Character{88}}%
\TextAt{21}{327.087}{473.028}{\Character{28}}%
\TextAt{16}{340.053}{486.777}{\Character{106}}%
\TextAt{13}{343.374}{486.777}{z}%
\TextAt{21}{348.812}{484.983}{\Character{28}}%
\TextAt{16}{353.988}{486.777}{\Character{106}}%
\NativeMathAt{360.611}{486.777}{11.955}{\leq}%
\TextAt{13}{373.249}{486.777}{t}%
\TextAt{16}{377.477}{486.777}{\Character{103}}%
\TextAt{13}{383.454}{486.777}{\Character{58}}%
\TextAt{8}{79.226}{458.495}{Its}%
\TextAt{8}{96.513}{458.495}{volume}%
\TextAt{8}{136.838}{458.495}{is}%
\TextAt{8}{260.405}{441.449}{vol}%
\TextAt{9}{275.695}{441.449}{\Character{40}}%
\TextAt{13}{280.248}{441.449}{X}%
\TextAt{21}{289.963}{439.656}{t}%
\TextAt{9}{293.519}{441.449}{\Character{41}}%
\TextAt{9}{301.374}{441.449}{\Character{61}}%
\TextAt{9}{313.795}{441.449}{2}%
\TextAt{21}{319.671}{446.385}{r}%
\TextAt{13}{324.226}{441.449}{\Character{25}}%
\TextAt{21}{331.295}{446.385}{s}%
\TextAt{13}{337.092}{449.537}{t}%
\TextAt{21}{341.319}{453.875}{n}%
\Vector{q 0 G 0 g 0.478 w 264.904 -329.562 m 275.144 -329.562 l S Q}%
\TextAt{13}{336.904}{433.248}{n}%
\TextAt{9}{343.892}{433.248}{\Character{33}}%
\TextAt{13}{348.339}{441.449}{\Character{58}}%
\TextAt{8}{79.226}{419.242}{Leaving}%
\TextAt{8}{122.587}{419.242}{aside}%
\TextAt{8}{151.529}{419.242}{the}%
\TextAt{8}{171.302}{419.242}{proof}%
\TextAt{8}{201.488}{419.242}{of}%
\TextAt{8}{214.435}{419.242}{this}%
\TextAt{8}{236.874}{419.242}{formula}%
\TextAt{8}{279.421}{419.242}{for}%
\TextAt{8}{296.923}{419.242}{the}%
\TextAt{8}{316.696}{419.242}{moment}%
\TextAt{8}{361.515}{419.242}{(which}%
\TextAt{8}{399.173}{419.242}{incidentally}%
\TextAt{8}{461.877}{419.242}{was}%
\TextAt{8}{483.993}{419.242}{exercise}%
\TextAt{8}{526.911}{419.242}{2}%
\TextAt{8}{79.226}{402.196}{of}%
\TextAt{8}{92.416}{402.196}{chap.}%
\TextAt{8}{124.56}{402.196}{I,}%
\TextAt{8}{135.789}{402.196}{§5),}%
\TextAt{8}{158.917}{402.196}{we}%
\TextAt{8}{176.004}{402.196}{deduce}%
\TextAt{8}{214.873}{402.196}{the}%
\TextAt{8}{234.889}{402.196}{proposition}%
\TextAt{8}{296.282}{402.196}{from}%
\TextAt{8}{323.782}{402.196}{Minkowski’s}%
\TextAt{8}{389.777}{402.196}{lattice}%
\TextAt{8}{425.406}{402.196}{point}%
\TextAt{8}{455.835}{402.196}{theorem}%
\TextAt{8}{501.22}{402.196}{(chap.}%
\TextAt{8}{79.226}{385.15}{I,}%
\TextAt{8}{90.269}{385.15}{(4.4))}%
\TextAt{8}{122.473}{385.15}{as}%
\TextAt{8}{136.516}{385.15}{follows.}%
\TextAt{8}{179.394}{385.15}{Consider}%
\TextAt{8}{227.905}{385.15}{in}%
\TextAt{13}{241.231}{385.15}{K}%
\TextAt{25}{251.19}{383.356}{R}%
\TextAt{8}{261.015}{385.15}{the}%
\TextAt{8}{280.845}{385.15}{lattice}%
\TextAt{9}{316.289}{385.15}{\Character{0}}%
\TextAt{9}{326.905}{385.15}{\Character{61}}%
\TextAt{13}{339.349}{385.15}{j}%
\TextAt{14}{344.865}{385.15}{o}%
\TextAt{8}{354.28}{385.15}{defined}%
\TextAt{8}{394.266}{385.15}{by}%
\TextAt{14}{410.199}{385.15}{o}%
\TextAt{8}{416.043}{385.15}{.}%
\TextAt{8}{424.383}{385.15}{By}%
\TextAt{8}{442.419}{385.15}{chap.}%
\TextAt{8}{474.502}{385.15}{I,}%
\TextAt{8}{485.545}{385.15}{(5.2),}%
\TextAt{8}{516.51}{385.15}{the}%
\TextAt{8}{79.226}{368.104}{volume}%
\TextAt{8}{119.551}{368.104}{of}%
\TextAt{8}{132.881}{368.104}{a}%
\TextAt{8}{142.636}{368.104}{fundamental}%
\TextAt{8}{210.279}{368.104}{mesh}%
\TextAt{8}{240.25}{368.104}{is}%
\TextAt{8}{252.014}{368.104}{vol}%
\TextAt{9}{267.305}{368.104}{\Character{40}\Character{0}\Character{41}}%
\TextAt{9}{287.007}{368.104}{\Character{61}}%
\TextAt{20}{299.47}{378.325}{\Character{112}}%
\Vector{q 0 G 0 g 0.478 w 239.425 -395.436 m 260.352 -395.436 l S Q}%
\TextAt{16}{311.425}{368.104}{\Character{106}}%
\TextAt{13}{314.746}{368.104}{d}%
\TextAt{21}{320.829}{366.31}{K}%
\TextAt{16}{329.031}{368.104}{\Character{106}}%
\TextAt{8}{332.352}{368.104}{.}%
\TextAt{8}{340.801}{368.104}{The}%
\TextAt{8}{364.854}{368.104}{inequality}%
\TextAt{8}{256.696}{344.776}{vol}%
\TextAt{9}{271.986}{344.776}{\Character{40}}%
\TextAt{13}{276.539}{344.776}{X}%
\TextAt{21}{286.254}{342.983}{t}%
\TextAt{9}{289.81}{344.776}{\Character{41}}%
\TextAt{13}{297.683}{344.776}{\Character{62}}%
\TextAt{9}{310.109}{344.776}{2}%
\TextAt{21}{315.962}{349.712}{n}%
\TextAt{8}{323.59}{344.776}{vol}%
\TextAt{9}{338.881}{344.776}{\Character{40}\Character{0}\Character{41}}%
\TextAt{8}{79.226}{321.448}{therefore}%
\TextAt{8}{128.326}{321.448}{holds}%
\TextAt{8}{158.955}{321.448}{if}%
\TextAt{8}{169.679}{321.448}{and}%
\TextAt{8}{192.442}{321.448}{only}%
\TextAt{8}{218.133}{321.448}{if}%
\TextAt{9}{228.857}{321.448}{2}%
\TextAt{21}{234.71}{325.787}{r}%
\TextAt{13}{239.265}{321.448}{\Character{25}}%
\TextAt{21}{246.334}{325.787}{s}%
\TextAt{13}{250.748}{321.448}{t}%
\TextAt{21}{254.975}{325.787}{n}%
\TextAt{9}{260.612}{321.448}{\Character{47}}%
\TextAt{13}{266.465}{321.448}{n}%
\TextAt{9}{273.452}{321.448}{\Character{33}}%
\TextAt{13}{280.025}{321.448}{\Character{62}}%
\TextAt{9}{292.45}{321.448}{2}%
\TextAt{21}{298.303}{325.787}{n}%
\TextAt{20}{303.939}{331.67}{\Character{112}}%
\Vector{q 0 G 0 g 0.478 w 243.895 -442.091 m 264.822 -442.091 l S Q}%
\TextAt{16}{315.895}{321.448}{\Character{106}}%
\TextAt{13}{319.216}{321.448}{d}%
\TextAt{21}{325.298}{319.655}{K}%
\TextAt{16}{333.501}{321.448}{\Character{106}}%
\TextAt{8}{336.822}{321.448}{,}%
\TextAt{8}{343.971}{321.448}{or}%
\TextAt{8}{358.281}{321.448}{equivalently}%
\TextAt{8}{423.329}{321.448}{if}%
\TextAt{13}{242.328}{292.912}{t}%
\TextAt{21}{246.555}{297.848}{n}%
\TextAt{9}{255.513}{292.912}{\Character{61}}%
\TextAt{13}{267.938}{292.912}{n}%
\TextAt{9}{274.926}{292.912}{\Character{33}}%
\TextAt{20}{280.17}{309.769}{\Character{18}}%
\TextAt{9}{290.774}{301}{4}%
\Vector{q 0 G 0 g 0.478 w 218.166 -478.099 m 225.235 -478.099 l S Q}%
\TextAt{13}{290.166}{284.711}{\Character{25}}%
\TextAt{20}{298.43}{309.769}{\Character{19}}%
\TextAt{21}{307.231}{307.092}{s}%
\TextAt{20}{313.637}{303.717}{\Character{112}}%
\Vector{q 0 G 0 g 0.478 w 253.593 -470.043 m 274.52 -470.043 l S Q}%
\TextAt{16}{325.593}{292.912}{\Character{106}}%
\TextAt{13}{328.913}{292.912}{d}%
\TextAt{21}{334.996}{291.119}{K}%
\TextAt{16}{343.199}{292.912}{\Character{106}}%
\TextAt{9}{349.176}{292.912}{\Character{43}}%
\TextAt{13}{360.938}{292.912}{\Character{34}\Character{59}}%
\TextAt{8}{79.226}{265.741}{for}%
\TextAt{8}{96.389}{265.741}{some}%
\TextAt{13}{124.993}{265.741}{\Character{34}}%
\TextAt{13}{133.78}{265.741}{\Character{62}}%
\TextAt{9}{146.217}{265.741}{0}%
\TextAt{8}{152.07}{265.741}{.}%
\TextAt{8}{160.278}{265.741}{If}%
\TextAt{8}{171.248}{265.741}{this}%
\TextAt{8}{193.348}{265.741}{is}%
\TextAt{8}{204.39}{265.741}{the}%
\TextAt{8}{223.824}{265.741}{case,}%
\TextAt{8}{251.269}{265.741}{there}%
\TextAt{8}{280.459}{265.741}{exists}%
\TextAt{8}{312.051}{265.741}{an}%
\TextAt{13}{327.588}{265.741}{\Character{11}}%
\NativeMathAt{338.431}{265.741}{11.955}{\in}%
\TextAt{14}{349.722}{265.741}{o}%
\TextAt{8}{355.566}{265.741}{,}%
\TextAt{13}{362.137}{265.741}{\Character{11}}%
\NativeMathAt{372.98}{265.741}{11.955}{\neq}%
\TextAt{9}{385.389}{265.741}{0}%
\TextAt{8}{391.258}{265.741}{,}%
\TextAt{8}{397.83}{265.741}{such}%
\TextAt{8}{423.505}{265.741}{that}%
\TextAt{13}{448.151}{265.741}{j}%
\TextAt{13}{453.662}{265.741}{\Character{11}}%
\NativeMathAt{464.51}{265.741}{11.955}{\in}%
\TextAt{13}{475.8}{265.741}{X}%
\TextAt{21}{485.516}{263.947}{t}%
\TextAt{8}{489.072}{265.741}{.}%
\TextAt{8}{497.279}{265.741}{As}%
\TextAt{8}{513.844}{265.741}{this}%
\TextAt{8}{79.226}{248.695}{holds}%
\TextAt{8}{109.537}{248.695}{for}%
\TextAt{8}{127.103}{248.695}{all}%
\TextAt{13}{143.044}{248.695}{\Character{34}}%
\TextAt{13}{151.831}{248.695}{\Character{62}}%
\TextAt{9}{164.268}{248.695}{0}%
\TextAt{8}{170.121}{248.695}{,}%
\TextAt{8}{177.015}{248.695}{and}%
\TextAt{8}{199.459}{248.695}{since}%
\TextAt{13}{227.809}{248.695}{X}%
\TextAt{21}{237.525}{246.901}{t}%
\TextAt{8}{244.66}{248.695}{contains}%
\TextAt{8}{290.261}{248.695}{only}%
\TextAt{8}{315.634}{248.695}{finitely}%
\TextAt{8}{354.66}{248.695}{many}%
\TextAt{8}{386.214}{248.695}{lattice}%
\TextAt{8}{421.665}{248.695}{points,}%
\TextAt{8}{459.846}{248.695}{it}%
\TextAt{8}{471.232}{248.695}{continues}%
\TextAt{8}{522.357}{248.695}{to}%
\TextAt{8}{79.226}{231.649}{hold}%
\TextAt{8}{105.241}{231.649}{for}%
\TextAt{13}{123.126}{231.649}{\Character{34}}%
\TextAt{9}{131.924}{231.649}{\Character{61}}%
\TextAt{9}{144.334}{231.649}{0}%
\TextAt{8}{150.203}{231.649}{.}%
\TextAt{8}{158.651}{231.649}{Applying}%
\TextAt{8}{209.377}{231.649}{now}%
\TextAt{8}{233.766}{231.649}{the}%
\TextAt{8}{253.922}{231.649}{inequality}%
\TextAt{8}{308.234}{231.649}{between}%
\TextAt{8}{353.425}{231.649}{arithmetic}%
\TextAt{8}{410.009}{231.649}{and}%
\TextAt{8}{432.771}{231.649}{geometric}%
\TextAt{8}{486.103}{231.649}{means,}%
\TextAt{9}{239.445}{206.263}{1}%
\Vector{q 0 G 0 g 0.478 w 166.878 -572.836 m 173.865 -572.836 l S Q}%
\TextAt{13}{238.878}{189.975}{n}%
\TextAt{20}{249.053}{209.533}{\Character{88}}%
\TextAt{21}{255.348}{184.427}{\Character{28}}%
\TextAt{16}{268.314}{198.175}{\Character{106}}%
\TextAt{13}{271.635}{198.175}{z}%
\TextAt{21}{277.073}{196.382}{\Character{28}}%
\TextAt{16}{282.25}{198.175}{\Character{106}}%
\NativeMathAt{288.873}{198.175}{11.955}{\geq}%
\TextAt{20}{301.511}{218.619}{\Character{32}}%
\TextAt{20}{310.975}{209.533}{\Character{89}}%
\TextAt{21}{316.274}{184.427}{\Character{28}}%
\TextAt{16}{328.244}{198.175}{\Character{106}}%
\TextAt{13}{331.565}{198.175}{z}%
\TextAt{21}{337.003}{196.382}{\Character{28}}%
\TextAt{16}{342.179}{198.175}{\Character{106}}%
\TextAt{20}{345.5}{218.619}{\Character{33}}%
\TextAt{17}{354.965}{215.942}{1\Character{47}}%
\TextAt{21}{363.433}{215.942}{n}%
\TextAt{13}{371.062}{198.175}{\Character{59}}%
\TextAt{8}{79.226}{167.656}{we}%
\TextAt{8}{96.454}{167.656}{obtain}%
\TextAt{8}{132.881}{167.656}{the}%
\TextAt{8}{153.037}{167.656}{desired}%
\TextAt{8}{192.764}{167.656}{result:}%
\TextAt{9}{129.897}{139.06}{1}%
\NativeMathAt{139.071}{139.06}{11.955}{\leq}%
\TextAt{16}{151.672}{139.06}{\Character{106}}%
\TextAt{13}{155.011}{139.06}{N}%
\TextAt{21}{164.385}{137.201}{K}%
\TextAt{3}{172.089}{137.201}{\Character{106}}%
\TextAt{25}{174.442}{137.201}{Q}%
\TextAt{9}{181.139}{139.06}{\Character{40}}%
\TextAt{13}{185.691}{139.06}{\Character{11}}%
\TextAt{9}{193.213}{139.06}{\Character{41}}%
\TextAt{16}{197.765}{139.06}{\Character{106}}%
\TextAt{9}{204.407}{139.06}{\Character{61}}%
\TextAt{20}{216.832}{150.418}{\Character{89}}%
\TextAt{21}{222.131}{125.312}{\Character{28}}%
\TextAt{16}{234.101}{139.06}{\Character{106}}%
\TextAt{13}{237.422}{139.06}{\Character{28}}%
\TextAt{13}{243.83}{139.06}{\Character{11}}%
\TextAt{16}{251.362}{139.06}{\Character{106}}%
\NativeMathAt{257.985}{139.06}{11.955}{\leq}%
\TextAt{9}{275.204}{147.148}{1}%
\Vector{q 0 G 0 g 0.478 w 199.819 -631.951 m 212.443 -631.951 l S Q}%
\TextAt{13}{271.819}{130.86}{n}%
\TextAt{21}{278.807}{134.313}{n}%
\TextAt{20}{287.631}{159.504}{\Character{32}}%
\TextAt{20}{297.095}{150.418}{\Character{88}}%
\TextAt{21}{303.391}{125.312}{\Character{28}}%
\TextAt{16}{316.357}{139.06}{\Character{106}}%
\TextAt{13}{319.677}{139.06}{\Character{28}}%
\TextAt{13}{326.085}{139.06}{\Character{11}}%
\TextAt{16}{333.618}{139.06}{\Character{106}}%
\TextAt{20}{336.939}{159.504}{\Character{33}}%
\TextAt{21}{346.404}{156.827}{n}%
\NativeMathAt{355.361}{139.06}{11.955}{\leq}%
\TextAt{13}{370.556}{147.148}{t}%
\TextAt{21}{374.783}{151.487}{n}%
\Vector{q 0 G 0 g 0.478 w 297.176 -631.951 m 309.8 -631.951 l S Q}%
\TextAt{13}{369.176}{130.86}{n}%
\TextAt{21}{376.163}{134.313}{n}%
\TextAt{9}{386.316}{139.06}{\Character{61}}%
\TextAt{13}{401.129}{147.148}{n}%
\TextAt{9}{408.117}{147.148}{\Character{33}}%
\Vector{q 0 G 0 g 0.478 w 327.937 -631.951 m 340.561 -631.951 l S Q}%
\TextAt{13}{399.937}{130.86}{n}%
\TextAt{21}{406.924}{134.313}{n}%
\TextAt{20}{415.749}{155.917}{\Character{18}}%
\TextAt{9}{426.353}{147.148}{4}%
\Vector{q 0 G 0 g 0.478 w 353.745 -631.951 m 360.814 -631.951 l S Q}%
\TextAt{13}{425.745}{130.86}{\Character{25}}%
\TextAt{20}{434.009}{155.917}{\Character{19}}%
\TextAt{21}{442.81}{153.241}{s}%
\TextAt{20}{449.216}{149.866}{\Character{112}}%
\Vector{q 0 G 0 g 0.478 w 389.172 -623.895 m 410.099 -623.895 l S Q}%
\TextAt{16}{461.172}{139.06}{\Character{106}}%
\TextAt{13}{464.492}{139.06}{d}%
\TextAt{21}{470.575}{137.267}{K}%
\TextAt{16}{478.778}{139.06}{\Character{106}}%
\NativeMathAt{204.407}{99.973}{11.955}{\leq}%
\TextAt{13}{219.414}{108.061}{n}%
\TextAt{9}{226.402}{108.061}{\Character{33}}%
\Vector{q 0 G 0 g 0.478 w 146.222 -671.038 m 158.846 -671.038 l S Q}%
\TextAt{13}{218.222}{91.773}{n}%
\TextAt{21}{225.209}{95.226}{n}%
\TextAt{20}{234.034}{116.83}{\Character{18}}%
\TextAt{9}{244.638}{108.061}{4}%
\Vector{q 0 G 0 g 0.478 w 172.03 -671.038 m 179.099 -671.038 l S Q}%
\TextAt{13}{244.03}{91.773}{\Character{25}}%
\TextAt{20}{252.294}{116.83}{\Character{19}}%
\TextAt{21}{261.095}{114.154}{n}%
\TextAt{17}{266.233}{114.154}{\Character{47}2}%
\TextAt{20}{277.192}{110.779}{\Character{112}}%
\Vector{q 0 G 0 g 0.478 w 217.147 -662.982 m 238.074 -662.982 l S Q}%
\TextAt{16}{289.147}{99.973}{\Character{106}}%
\TextAt{13}{292.468}{99.973}{d}%
\TextAt{21}{298.551}{98.18}{K}%
\TextAt{16}{306.753}{99.973}{\Character{106}}%
\TextAt{13}{310.074}{99.973}{\Character{58}}%
\TextAt{8}{79.226}{72.921}{Given}%
\TextAt{8}{113.119}{72.921}{this,}%
\TextAt{8}{139.193}{72.921}{it}%
\TextAt{8}{150.897}{72.921}{remains}%
\TextAt{8}{194.534}{72.921}{to}%
\TextAt{8}{208.844}{72.921}{prove}%
\TextAt{8}{240.394}{72.921}{the}%
\TextAt{8}{260.55}{72.921}{following}%
\TextAt{8}{309.985}{72.921}{lemma.}%
\Vector{q 0 G 0 g 0.399 w 452.334 -701.079 m 452.334 -693.009 l S Q}%
\Vector{q 0 G 0 g 0.399 w 452.534 -693.208 m 459.707 -693.208 l S Q}%
\Vector{q 0 G 0 g 0.399 w 452.534 -700.879 m 459.707 -700.879 l S Q}%
\Vector{q 0 G 0 g 0.399 w 459.906 -701.079 m 459.906 -693.009 l S Q}%
\TextAt{8}{297.211}{43.033}{186}%
\Destination{78.23}{809.815}{nav.page.186}%
\Destination{79.36}{721.955}{nav.print.204}%
\Destination{79.36}{601.055}{nav.number.3.2.14}%
\Hotlink{526.411}{417.055}{6.858}{11.047}{nav.ex.1.5.2}%
\Hotlink{124.06}{399.955}{5.22}{11.047}{nav.chapter.1}%
\Hotlink{140.968}{399.955}{6.858}{11.047}{nav.section.1.5}%
\Hotlink{474.002}{382.955}{5.22}{11.047}{nav.chapter.1}%
\Hotlink{89.769}{382.955}{25.078}{11.047}{nav.number.1.4.4}%
\Hotlink{485.045}{382.955}{25.078}{11.047}{nav.number.1.5.2}%
\Destination{79.226}{706.951}{errata.93}%
}

% Page 199 (body page 187)
\Page{%
\TextAt{2}{79.226}{796.943}{III.}%
\TextAt{2}{98.317}{796.943}{Riemann–Roch}%
\TextAt{2}{174.375}{796.943}{Theory}%
\Vector{q 0 G 0 g 0.299 w 7.226 18.729 m 460.769 18.729 l S Q}%
\TextAt{12}{79.226}{761.124}{Lemma}%
\TextAt{12}{126.891}{761.124}{(2.15).}%
\TextAt{8}{168.938}{761.124}{In}%
\TextAt{8}{183.559}{761.124}{Minkowski}%
\TextAt{8}{241.828}{761.124}{space}%
\TextAt{13}{273.103}{761.124}{K}%
\TextAt{25}{283.062}{759.33}{R}%
\TextAt{9}{292.637}{761.124}{\Character{61}}%
\TextAt{9}{305.047}{761.124}{\Character{91}}%
\TextAt{20}{308.315}{770.09}{\Character{81}}%
\TextAt{21}{319.606}{757.637}{\Character{28}}%
\TextAt{18}{326.775}{761.124}{C}%
\TextAt{9}{335.409}{761.124}{\Character{93}}%
\TextAt{17}{338.661}{765.462}{\Character{43}}%
\TextAt{8}{345.745}{761.124}{,}%
\TextAt{8}{352.894}{761.124}{the}%
\TextAt{8}{373.051}{761.124}{domain}%
\TextAt{13}{227.013}{733.192}{X}%
\TextAt{21}{236.728}{731.399}{t}%
\TextAt{9}{243.605}{733.192}{\Character{61}}%
\TextAt{16}{256.03}{733.192}{\Character{102}}%
\TextAt{9}{262.008}{733.192}{\Character{40}}%
\TextAt{13}{266.56}{733.192}{z}%
\TextAt{21}{271.998}{731.399}{\Character{28}}%
\TextAt{9}{277.175}{733.192}{\Character{41}}%
\NativeMathAt{285.048}{733.192}{11.955}{\in}%
\TextAt{13}{296.339}{733.192}{K}%
\TextAt{25}{306.298}{731.399}{R}%
\TextAt{16}{315.873}{733.192}{\Character{106}}%
\TextAt{20}{322.515}{744.549}{\Character{88}}%
\TextAt{21}{328.81}{719.443}{\Character{28}}%
\TextAt{16}{341.776}{733.192}{\Character{106}}%
\TextAt{13}{345.097}{733.192}{z}%
\TextAt{21}{350.535}{731.399}{\Character{28}}%
\TextAt{16}{355.711}{733.192}{\Character{106}}%
\TextAt{13}{362.353}{733.192}{\Character{60}}%
\TextAt{13}{374.762}{733.192}{t}%
\TextAt{16}{379.006}{733.192}{\Character{103}}%
\TextAt{8}{79.226}{698.069}{has}%
\TextAt{8}{100.1}{698.069}{volume}%
\TextAt{8}{260.405}{681.023}{vol}%
\TextAt{9}{275.695}{681.023}{\Character{40}}%
\TextAt{13}{280.248}{681.023}{X}%
\TextAt{21}{289.963}{679.23}{t}%
\TextAt{9}{293.519}{681.023}{\Character{41}}%
\TextAt{9}{301.374}{681.023}{\Character{61}}%
\TextAt{9}{313.795}{681.023}{2}%
\TextAt{21}{319.671}{685.959}{r}%
\TextAt{13}{324.226}{681.023}{\Character{25}}%
\TextAt{21}{331.295}{685.959}{s}%
\TextAt{13}{337.092}{689.111}{t}%
\TextAt{21}{341.319}{693.449}{n}%
\Vector{q 0 G 0 g 0.478 w 264.904 -89.988 m 275.144 -89.988 l S Q}%
\TextAt{13}{336.904}{672.822}{n}%
\TextAt{9}{343.892}{672.822}{\Character{33}}%
\TextAt{13}{348.339}{681.023}{\Character{58}}%
\TextAt{19}{79.226}{639.52}{Proof.}%
\TextAt{8}{116}{639.52}{vol}%
\TextAt{9}{131.291}{639.52}{\Character{40}}%
\TextAt{13}{135.843}{639.52}{X}%
\TextAt{21}{145.558}{637.727}{t}%
\TextAt{9}{149.115}{639.52}{\Character{41}}%
\TextAt{8}{157.298}{639.52}{is}%
\TextAt{9}{168.796}{639.52}{2}%
\TextAt{21}{174.649}{643.859}{s}%
\TextAt{8}{182.695}{639.52}{times}%
\TextAt{8}{213.704}{639.52}{the}%
\TextAt{8}{233.594}{639.52}{Lebesgue}%
\TextAt{8}{283.947}{639.52}{volume}%
\TextAt{8}{324.005}{639.52}{Vol}%
\TextAt{9}{341.89}{639.52}{\Character{40}}%
\TextAt{13}{346.443}{639.52}{f}%
\TextAt{9}{353.489}{639.52}{\Character{40}}%
\TextAt{13}{358.041}{639.52}{X}%
\TextAt{21}{367.757}{637.727}{t}%
\TextAt{9}{371.313}{639.52}{\Character{41}\Character{41}}%
\TextAt{8}{384.049}{639.52}{of}%
\TextAt{8}{397.113}{639.52}{the}%
\TextAt{8}{417.004}{639.52}{image}%
\TextAt{13}{450.559}{639.52}{f}%
\TextAt{9}{457.605}{639.52}{\Character{40}}%
\TextAt{13}{462.158}{639.52}{X}%
\TextAt{21}{471.873}{637.727}{t}%
\TextAt{9}{475.429}{639.52}{\Character{41}}%
\TextAt{8}{483.613}{639.52}{under}%
\TextAt{8}{516.51}{639.52}{the}%
\TextAt{8}{79.226}{622.474}{mapping}%
\TextAt{8}{127.358}{622.474}{chap.}%
\TextAt{8}{159.549}{622.474}{I,}%
\TextAt{8}{170.918}{622.474}{(5.1),}%
\TextAt{13}{214.038}{594.542}{f}%
\TextAt{9}{224.405}{594.542}{\Character{58}}%
\TextAt{13}{230.977}{594.542}{K}%
\TextAt{25}{240.937}{592.749}{R}%
\NativeArrowAt{250.512}{594.542}{11.955}{19.26}{\longrightarrow}%
\TextAt{20}{273.094}{605.9}{\Character{89}}%
\TextAt{21}{278.393}{580.794}{\Character{28}}%
\TextAt{18}{290.362}{594.542}{R}%
\TextAt{13}{298.997}{594.542}{\Character{59}}%
\TextAt{9}{328.151}{594.542}{\Character{40}}%
\TextAt{13}{332.703}{594.542}{z}%
\TextAt{21}{338.141}{592.749}{\Character{28}}%
\TextAt{9}{343.318}{594.542}{\Character{41}}%
\NativeArrowAt{351.191}{594.542}{11.955}{19.26}{\longmapsto}%
\TextAt{9}{373.773}{594.542}{\Character{40}}%
\TextAt{13}{378.325}{594.542}{x}%
\TextAt{21}{384.977}{592.749}{\Character{28}}%
\TextAt{9}{390.154}{594.542}{\Character{41}}%
\TextAt{13}{394.706}{594.542}{\Character{59}}%
\TextAt{8}{79.226}{558.75}{where}%
\TextAt{13}{113.74}{558.75}{x}%
\TextAt{21}{120.392}{556.957}{\Character{26}}%
\TextAt{9}{129.748}{558.75}{\Character{61}}%
\TextAt{13}{143.373}{558.75}{z}%
\TextAt{21}{148.811}{556.957}{\Character{26}}%
\TextAt{8}{153.646}{558.75}{,}%
\TextAt{13}{161.676}{558.75}{x}%
\TextAt{21}{168.328}{556.957}{\Character{27}}%
\TextAt{9}{178.472}{558.75}{\Character{61}}%
\TextAt{8}{192.098}{558.75}{Re}%
\TextAt{9}{205.918}{558.75}{\Character{40}}%
\TextAt{13}{210.471}{558.75}{z}%
\TextAt{21}{215.909}{556.957}{\Character{27}}%
\TextAt{9}{221.532}{558.75}{\Character{41}}%
\TextAt{8}{226.084}{558.75}{,}%
\TextAt{13}{234.114}{558.75}{x}%
\Vector{q 0 G 0 g 0.359 w 168.766 -212.356 m 174.389 -212.356 l S Q}%
\TextAt{21}{240.766}{556.957}{\Character{27}}%
\TextAt{9}{250.911}{558.75}{\Character{61}}%
\TextAt{8}{264.536}{558.75}{Im}%
\TextAt{9}{278.512}{558.75}{\Character{40}}%
\TextAt{13}{283.064}{558.75}{z}%
\TextAt{21}{288.502}{556.957}{\Character{27}}%
\TextAt{9}{294.126}{558.75}{\Character{41}}%
\TextAt{8}{298.678}{558.75}{.}%
\TextAt{8}{309.239}{558.75}{Substituting}%
\TextAt{13}{377}{558.75}{x}%
\TextAt{21}{383.652}{556.957}{i}%
\TextAt{8}{387.034}{558.75}{,}%
\TextAt{13}{395.063}{558.75}{i}%
\TextAt{9}{403.578}{558.75}{\Character{61}}%
\TextAt{9}{417.195}{558.75}{1}%
\TextAt{13}{423.057}{558.75}{\Character{59}}%
\TextAt{13}{428.281}{558.75}{\Character{58}}%
\TextAt{13}{433.518}{558.75}{\Character{58}}%
\TextAt{13}{438.742}{558.75}{\Character{58}}%
\TextAt{13}{443.979}{558.75}{\Character{59}}%
\TextAt{13}{449.215}{558.75}{r}%
\TextAt{8}{454.878}{558.75}{,}%
\TextAt{8}{462.908}{558.75}{instead}%
\TextAt{8}{503.996}{558.75}{of}%
\TextAt{13}{518.031}{558.75}{x}%
\TextAt{21}{524.683}{556.957}{\Character{26}}%
\TextAt{8}{529.518}{558.75}{,}%
\TextAt{8}{79.226}{541.704}{and}%
\TextAt{13}{101.738}{541.704}{y}%
\TextAt{21}{107.446}{539.911}{j}%
\TextAt{13}{111.828}{541.704}{\Character{59}}%
\TextAt{13}{117.052}{541.704}{z}%
\TextAt{21}{122.51}{539.911}{j}%
\TextAt{8}{126.892}{541.704}{,}%
\TextAt{13}{133.841}{541.704}{j}%
\TextAt{9}{142.677}{541.704}{\Character{61}}%
\TextAt{9}{155.087}{541.704}{1}%
\TextAt{13}{160.956}{541.704}{\Character{59}}%
\TextAt{13}{166.18}{541.704}{\Character{58}}%
\TextAt{13}{171.417}{541.704}{\Character{58}}%
\TextAt{13}{176.641}{541.704}{\Character{58}}%
\TextAt{13}{181.878}{541.704}{\Character{59}}%
\TextAt{13}{187.114}{541.704}{s}%
\TextAt{8}{192.691}{541.704}{,}%
\TextAt{8}{199.639}{541.704}{instead}%
\TextAt{8}{239.773}{541.704}{of}%
\TextAt{13}{252.853}{541.704}{x}%
\TextAt{21}{259.505}{539.911}{\Character{27}}%
\TextAt{13}{265.128}{541.704}{\Character{59}}%
\TextAt{13}{270.352}{541.704}{x}%
\Vector{q 0 G 0 g 0.359 w 205.025 -229.402 m 210.648 -229.402 l S Q}%
\TextAt{21}{277.025}{539.911}{\Character{27}}%
\TextAt{8}{282.648}{541.704}{,}%
\TextAt{8}{289.597}{541.704}{we}%
\TextAt{8}{306.574}{541.704}{see}%
\TextAt{8}{325.236}{541.704}{that}%
\TextAt{13}{350.354}{541.704}{f}%
\TextAt{9}{357.401}{541.704}{\Character{40}}%
\TextAt{13}{361.953}{541.704}{X}%
\TextAt{21}{371.668}{539.911}{t}%
\TextAt{9}{375.224}{541.704}{\Character{41}}%
\TextAt{8}{383.424}{541.704}{is}%
\TextAt{8}{394.937}{541.704}{described}%
\TextAt{8}{446.44}{541.704}{by}%
\TextAt{8}{462.449}{541.704}{the}%
\TextAt{8}{482.355}{541.704}{inequality}%
\TextAt{16}{172.657}{510.958}{\Character{106}}%
\TextAt{13}{175.978}{510.958}{x}%
\TextAt{17}{182.63}{509.164}{1}%
\TextAt{16}{187.362}{510.958}{\Character{106}}%
\TextAt{9}{193.34}{510.958}{\Character{43}}%
\TextAt{16}{205.101}{510.958}{\Character{1}}%
\TextAt{16}{210.397}{510.958}{\Character{1}}%
\TextAt{16}{215.705}{510.958}{\Character{1}}%
\TextAt{9}{221.706}{510.958}{\Character{43}}%
\TextAt{16}{233.467}{510.958}{\Character{106}}%
\TextAt{13}{236.788}{510.958}{x}%
\TextAt{21}{243.44}{509.164}{r}%
\TextAt{16}{247.995}{510.958}{\Character{106}}%
\TextAt{9}{253.972}{510.958}{\Character{43}}%
\TextAt{9}{265.724}{510.958}{2}%
\TextAt{20}{271.587}{525.522}{\Character{113}}%
\Vector{q 0 G 0 g 0.478 w 211.542 -248.238 m 247.532 -248.238 l S Q}%
\TextAt{13}{283.542}{510.958}{y}%
\TextAt{17}{289.678}{515.076}{2}%
\TextAt{17}{289.249}{508.027}{1}%
\TextAt{9}{297.067}{510.958}{\Character{43}}%
\TextAt{13}{308.829}{510.958}{z}%
\TextAt{17}{314.799}{515.076}{2}%
\TextAt{17}{314.267}{508.027}{1}%
\TextAt{9}{322.188}{510.958}{\Character{43}}%
\TextAt{16}{333.95}{510.958}{\Character{1}}%
\TextAt{16}{339.246}{510.958}{\Character{1}}%
\TextAt{16}{344.554}{510.958}{\Character{1}}%
\TextAt{9}{350.554}{510.958}{\Character{43}}%
\TextAt{9}{362.306}{510.958}{2}%
\TextAt{20}{368.168}{521.592}{\Character{112}}%
\Vector{q 0 G 0 g 0.478 w 308.123 -252.169 m 344.113 -252.169 l S Q}%
\TextAt{13}{380.123}{510.958}{y}%
\TextAt{17}{386.26}{514.411}{2}%
\TextAt{21}{385.831}{508.002}{s}%
\TextAt{9}{393.649}{510.958}{\Character{43}}%
\TextAt{13}{405.41}{510.958}{z}%
\TextAt{17}{411.381}{514.411}{2}%
\TextAt{21}{410.848}{508.002}{s}%
\TextAt{13}{419.434}{510.958}{\Character{60}}%
\TextAt{13}{431.844}{510.958}{t\Character{58}}%
\TextAt{8}{79.226}{483.026}{The}%
\TextAt{8}{102.702}{483.026}{factor}%
\TextAt{9}{135.622}{483.026}{2}%
\TextAt{8}{144.794}{483.026}{occurs}%
\TextAt{8}{180.368}{483.026}{because}%
\TextAt{16}{223.091}{483.026}{\Character{106}}%
\TextAt{13}{226.412}{483.026}{z}%
\Vector{q 0 G 0 g 0.359 w 159.85 -288.081 m 165.473 -288.081 l S Q}%
\TextAt{21}{231.85}{481.233}{\Character{27}}%
\TextAt{16}{237.473}{483.026}{\Character{106}}%
\TextAt{9}{244.115}{483.026}{\Character{61}}%
\TextAt{16}{256.541}{483.026}{\Character{106}}%
\Vector{q 0 G 0 g 0.478 w 187.861 -284.153 m 193.832 -284.153 l S Q}%
\TextAt{13}{259.861}{483.026}{z}%
\TextAt{21}{265.832}{481.233}{\Character{27}}%
\TextAt{16}{271.455}{483.026}{\Character{106}}%
\TextAt{9}{278.097}{483.026}{\Character{61}}%
\TextAt{16}{290.523}{483.026}{\Character{106}}%
\TextAt{13}{293.844}{483.026}{z}%
\TextAt{21}{299.281}{481.233}{\Character{27}}%
\TextAt{16}{304.905}{483.026}{\Character{106}}%
\TextAt{8}{308.226}{483.026}{.}%
\TextAt{8}{316.481}{483.026}{Passing}%
\TextAt{8}{358.141}{483.026}{to}%
\TextAt{8}{371.873}{483.026}{polar}%
\TextAt{8}{401.541}{483.026}{coordinates}%
\TextAt{13}{463.141}{483.026}{y}%
\TextAt{21}{468.849}{481.233}{j}%
\TextAt{9}{476.552}{483.026}{\Character{61}}%
\TextAt{13}{488.978}{483.026}{u}%
\TextAt{21}{495.64}{481.233}{j}%
\TextAt{8}{502.015}{483.026}{cos}%
\TextAt{13}{519.68}{483.026}{\Character{18}}%
\TextAt{21}{525.136}{481.233}{j}%
\TextAt{8}{529.518}{483.026}{,}%
\TextAt{13}{79.226}{465.98}{z}%
\TextAt{21}{84.664}{464.187}{j}%
\TextAt{9}{93.121}{465.98}{\Character{61}}%
\TextAt{13}{106.3}{465.98}{u}%
\TextAt{21}{112.962}{464.187}{j}%
\TextAt{8}{119.337}{465.98}{sin}%
\TextAt{13}{135.7}{465.98}{\Character{18}}%
\TextAt{21}{141.155}{464.187}{j}%
\TextAt{8}{145.537}{465.98}{,}%
\TextAt{8}{153.239}{465.98}{where}%
\TextAt{9}{187.491}{465.98}{0}%
\NativeMathAt{197.418}{465.98}{11.955}{\leq}%
\TextAt{13}{210.791}{465.98}{\Character{18}}%
\TextAt{21}{216.247}{464.187}{j}%
\NativeMathAt{224.703}{465.98}{11.955}{\leq}%
\TextAt{9}{238.076}{465.98}{2}%
\TextAt{13}{243.929}{465.98}{\Character{25}}%
\TextAt{8}{250.999}{465.98}{,}%
\TextAt{9}{258.701}{465.98}{0}%
\NativeMathAt{268.628}{465.98}{11.955}{\leq}%
\TextAt{13}{282.001}{465.98}{u}%
\TextAt{21}{288.664}{464.187}{j}%
\TextAt{8}{293.046}{465.98}{,}%
\TextAt{8}{300.748}{465.98}{one}%
\TextAt{8}{322.65}{465.98}{sees}%
\TextAt{8}{346.62}{465.98}{that}%
\TextAt{8}{372.431}{465.98}{Vol}%
\TextAt{9}{390.316}{465.98}{\Character{40}}%
\TextAt{13}{394.869}{465.98}{f}%
\TextAt{9}{401.915}{465.98}{\Character{40}}%
\TextAt{13}{406.467}{465.98}{X}%
\TextAt{21}{416.183}{464.187}{t}%
\TextAt{9}{419.739}{465.98}{\Character{41}\Character{41}}%
\TextAt{8}{433.183}{465.98}{is}%
\TextAt{8}{445.389}{465.98}{computed}%
\TextAt{8}{499.809}{465.98}{by}%
\TextAt{8}{516.51}{465.98}{the}%
\TextAt{8}{79.226}{448.934}{integral}%
\TextAt{13}{181.266}{429.214}{I}%
\TextAt{9}{187.369}{429.214}{\Character{40}}%
\TextAt{13}{191.921}{429.214}{t}%
\TextAt{9}{196.148}{429.214}{\Character{41}}%
\TextAt{9}{204.003}{429.214}{\Character{61}}%
\TextAt{20}{216.447}{445.487}{\Character{90}}%
\TextAt{13}{230.395}{429.214}{u}%
\TextAt{17}{237.057}{427.421}{1}%
\TextAt{16}{243.782}{429.214}{\Character{1}}%
\TextAt{16}{249.078}{429.214}{\Character{1}}%
\TextAt{16}{254.386}{429.214}{\Character{1}}%
\TextAt{13}{259.722}{429.214}{u}%
\TextAt{21}{266.384}{427.421}{s}%
\TextAt{8}{272.791}{429.214}{d}%
\TextAt{13}{279.295}{429.214}{x}%
\TextAt{17}{285.947}{427.421}{1}%
\TextAt{16}{292.671}{429.214}{\Character{1}}%
\TextAt{16}{297.967}{429.214}{\Character{1}}%
\TextAt{16}{303.275}{429.214}{\Character{1}}%
\TextAt{8}{308.612}{429.214}{d}%
\TextAt{13}{315.115}{429.214}{x}%
\TextAt{21}{321.767}{427.421}{r}%
\TextAt{8}{326.322}{429.214}{d}%
\TextAt{13}{332.826}{429.214}{u}%
\TextAt{17}{339.488}{427.421}{1}%
\TextAt{16}{346.213}{429.214}{\Character{1}}%
\TextAt{16}{351.509}{429.214}{\Character{1}}%
\TextAt{16}{356.817}{429.214}{\Character{1}}%
\TextAt{8}{362.153}{429.214}{d}%
\TextAt{13}{368.657}{429.214}{u}%
\TextAt{21}{375.319}{427.421}{s}%
\TextAt{8}{379.733}{429.214}{d}%
\TextAt{13}{386.237}{429.214}{\Character{18}}%
\TextAt{17}{391.692}{427.421}{1}%
\TextAt{16}{398.417}{429.214}{\Character{1}}%
\TextAt{16}{403.713}{429.214}{\Character{1}}%
\TextAt{16}{409.021}{429.214}{\Character{1}}%
\TextAt{8}{414.357}{429.214}{d}%
\TextAt{13}{420.861}{429.214}{\Character{18}}%
\TextAt{21}{426.316}{427.421}{s}%
\TextAt{8}{79.226}{403.609}{extended}%
\TextAt{8}{128.972}{403.609}{over}%
\TextAt{8}{154.018}{403.609}{the}%
\TextAt{8}{174.174}{403.609}{domain}%
\TextAt{16}{209.27}{375.677}{\Character{106}}%
\TextAt{13}{212.591}{375.677}{x}%
\TextAt{17}{219.243}{373.884}{1}%
\TextAt{16}{223.975}{375.677}{\Character{106}}%
\TextAt{9}{229.953}{375.677}{\Character{43}}%
\TextAt{16}{241.714}{375.677}{\Character{1}}%
\TextAt{16}{247.01}{375.677}{\Character{1}}%
\TextAt{16}{252.318}{375.677}{\Character{1}}%
\TextAt{9}{258.318}{375.677}{\Character{43}}%
\TextAt{16}{270.08}{375.677}{\Character{106}}%
\TextAt{13}{273.4}{375.677}{x}%
\TextAt{21}{280.053}{373.884}{r}%
\TextAt{16}{284.607}{375.677}{\Character{106}}%
\TextAt{9}{290.585}{375.677}{\Character{43}}%
\TextAt{9}{302.337}{375.677}{2}%
\TextAt{13}{308.199}{375.677}{u}%
\TextAt{17}{314.862}{373.884}{1}%
\TextAt{9}{322.251}{375.677}{\Character{43}}%
\TextAt{16}{334.012}{375.677}{\Character{1}}%
\TextAt{16}{339.308}{375.677}{\Character{1}}%
\TextAt{16}{344.616}{375.677}{\Character{1}}%
\TextAt{9}{350.616}{375.677}{\Character{43}}%
\TextAt{9}{362.368}{375.677}{2}%
\TextAt{13}{368.231}{375.677}{u}%
\TextAt{21}{374.893}{373.884}{s}%
\NativeMathAt{382.628}{375.677}{11.955}{\leq}%
\TextAt{13}{395.247}{375.677}{t\Character{58}}%
\TextAt{8}{79.226}{347.745}{Restricting}%
\TextAt{8}{139.491}{347.745}{this}%
\TextAt{8}{162.515}{347.745}{domain}%
\TextAt{8}{204.344}{347.745}{of}%
\TextAt{8}{217.876}{347.745}{integration}%
\TextAt{8}{277.602}{347.745}{to}%
\TextAt{13}{292.114}{347.745}{x}%
\TextAt{21}{298.766}{345.952}{i}%
\NativeMathAt{305.812}{347.745}{11.955}{\geq}%
\TextAt{9}{318.775}{347.745}{0}%
\TextAt{8}{324.628}{347.745}{,}%
\TextAt{8}{332.029}{347.745}{the}%
\TextAt{8}{352.387}{347.745}{integral}%
\TextAt{8}{395.197}{347.745}{gets}%
\TextAt{8}{419.524}{347.745}{divided}%
\TextAt{8}{461.018}{347.745}{by}%
\TextAt{9}{477.479}{347.745}{2}%
\TextAt{21}{483.332}{352.084}{r}%
\TextAt{8}{487.887}{347.745}{.}%
\TextAt{8}{496.94}{347.745}{Substi-}%
\TextAt{8}{79.226}{330.699}{tuting}%
\TextAt{9}{114.35}{330.699}{2}%
\TextAt{13}{120.203}{330.699}{u}%
\TextAt{21}{126.866}{328.906}{j}%
\TextAt{9}{134.569}{330.699}{\Character{61}}%
\TextAt{13}{146.994}{330.699}{w}%
\TextAt{21}{155.362}{328.906}{j}%
\TextAt{8}{163.642}{330.699}{gives}%
\TextAt{13}{245.507}{313.653}{I}%
\TextAt{9}{251.61}{313.653}{\Character{40}}%
\TextAt{13}{256.162}{313.653}{t}%
\TextAt{9}{260.389}{313.653}{\Character{41}}%
\TextAt{9}{268.244}{313.653}{\Character{61}}%
\TextAt{9}{280.665}{313.653}{2}%
\TextAt{21}{286.541}{318.589}{r}%
\TextAt{9}{291.095}{313.653}{4}%
\TextAt{3}{296.948}{318.589}{\Character{0}}%
\TextAt{21}{303.535}{318.589}{s}%
\TextAt{9}{307.949}{313.653}{\Character{40}2}%
\TextAt{13}{318.354}{313.653}{\Character{25}}%
\TextAt{9}{325.424}{313.653}{\Character{41}}%
\TextAt{21}{329.976}{318.589}{s}%
\TextAt{13}{334.39}{313.653}{I}%
\TextAt{21}{339.553}{311.86}{r}%
\TextAt{21}{343.139}{311.86}{\Character{59}s}%
\TextAt{9}{349.905}{313.653}{\Character{40}}%
\TextAt{13}{354.458}{313.653}{t}%
\TextAt{9}{358.685}{313.653}{\Character{41}}%
\TextAt{13}{363.237}{313.653}{\Character{59}}%
\TextAt{8}{79.226}{290.59}{where}%
\TextAt{8}{113.035}{290.59}{the}%
\TextAt{8}{133.192}{290.59}{integral}%
\TextAt{13}{198.647}{270.87}{I}%
\TextAt{21}{203.81}{269.077}{r}%
\TextAt{21}{207.397}{269.077}{\Character{59}s}%
\TextAt{9}{214.163}{270.87}{\Character{40}}%
\TextAt{13}{218.715}{270.87}{t}%
\TextAt{9}{222.942}{270.87}{\Character{41}}%
\TextAt{9}{230.797}{270.87}{\Character{61}}%
\TextAt{20}{243.241}{287.143}{\Character{90}}%
\TextAt{13}{257.188}{270.87}{w}%
\TextAt{17}{265.556}{269.077}{1}%
\TextAt{16}{272.281}{270.87}{\Character{1}}%
\TextAt{16}{277.577}{270.87}{\Character{1}}%
\TextAt{16}{282.885}{270.87}{\Character{1}}%
\TextAt{13}{288.221}{270.87}{w}%
\TextAt{21}{296.589}{269.077}{s}%
\TextAt{8}{302.996}{270.87}{d}%
\TextAt{13}{309.499}{270.87}{x}%
\TextAt{17}{316.151}{269.077}{1}%
\TextAt{16}{322.876}{270.87}{\Character{1}}%
\TextAt{16}{328.172}{270.87}{\Character{1}}%
\TextAt{16}{333.48}{270.87}{\Character{1}}%
\TextAt{8}{338.816}{270.87}{d}%
\TextAt{13}{345.32}{270.87}{x}%
\TextAt{21}{351.972}{269.077}{r}%
\TextAt{8}{356.527}{270.87}{d}%
\TextAt{13}{363.03}{270.87}{w}%
\TextAt{17}{371.398}{269.077}{1}%
\TextAt{16}{378.123}{270.87}{\Character{1}}%
\TextAt{16}{383.419}{270.87}{\Character{1}}%
\TextAt{16}{388.727}{270.87}{\Character{1}}%
\TextAt{8}{394.063}{270.87}{d}%
\TextAt{13}{400.567}{270.87}{w}%
\TextAt{21}{408.934}{269.077}{s}%
\TextAt{8}{79.226}{245.265}{has}%
\TextAt{8}{100.1}{245.265}{to}%
\TextAt{8}{114.41}{245.265}{be}%
\TextAt{8}{130.335}{245.265}{taken}%
\TextAt{8}{162.207}{245.265}{over}%
\TextAt{8}{187.253}{245.265}{the}%
\TextAt{8}{207.409}{245.265}{domain}%
\TextAt{13}{249.037}{245.265}{x}%
\TextAt{21}{255.689}{243.472}{i}%
\NativeMathAt{262.392}{245.265}{11.955}{\geq}%
\TextAt{9}{275.011}{245.265}{0}%
\TextAt{8}{280.864}{245.265}{,}%
\TextAt{13}{288.013}{245.265}{w}%
\TextAt{21}{296.381}{243.472}{j}%
\NativeMathAt{304.084}{245.265}{11.955}{\geq}%
\TextAt{9}{316.703}{245.265}{0}%
\TextAt{8}{326.454}{245.265}{and}%
\TextAt{13}{220.059}{217.333}{x}%
\TextAt{17}{226.711}{215.54}{1}%
\TextAt{9}{234.1}{217.333}{\Character{43}}%
\TextAt{16}{245.861}{217.333}{\Character{1}}%
\TextAt{16}{251.157}{217.333}{\Character{1}}%
\TextAt{16}{256.465}{217.333}{\Character{1}}%
\TextAt{9}{262.466}{217.333}{\Character{43}}%
\TextAt{13}{274.227}{217.333}{x}%
\TextAt{21}{280.879}{215.54}{r}%
\TextAt{9}{288.091}{217.333}{\Character{43}}%
\TextAt{13}{299.852}{217.333}{w}%
\TextAt{17}{308.22}{215.54}{1}%
\TextAt{9}{315.609}{217.333}{\Character{43}}%
\TextAt{16}{327.37}{217.333}{\Character{1}}%
\TextAt{16}{332.666}{217.333}{\Character{1}}%
\TextAt{16}{337.974}{217.333}{\Character{1}}%
\TextAt{9}{343.974}{217.333}{\Character{43}}%
\TextAt{13}{355.736}{217.333}{w}%
\TextAt{21}{364.104}{215.54}{s}%
\NativeMathAt{371.838}{217.333}{11.955}{\leq}%
\TextAt{13}{384.458}{217.333}{t\Character{58}}%
\TextAt{8}{517.802}{217.333}{(*)}%
\TextAt{8}{79.226}{189.401}{Clearly}%
\TextAt{13}{120.397}{189.401}{I}%
\TextAt{21}{125.56}{187.608}{r}%
\TextAt{21}{129.147}{187.608}{\Character{59}s}%
\TextAt{9}{135.912}{189.401}{\Character{40}}%
\TextAt{13}{140.464}{189.401}{t}%
\TextAt{9}{144.692}{189.401}{\Character{41}}%
\TextAt{9}{153.443}{189.401}{\Character{61}}%
\TextAt{13}{166.772}{189.401}{t}%
\TextAt{21}{171}{193.74}{r}%
\TextAt{17}{175.056}{193.74}{\Character{43}2}%
\TextAt{21}{185.877}{193.74}{s}%
\TextAt{13}{190.291}{189.401}{I}%
\TextAt{21}{195.454}{187.608}{r}%
\TextAt{21}{199.041}{187.608}{\Character{59}s}%
\TextAt{9}{205.806}{189.401}{\Character{40}1\Character{41}}%
\TextAt{9}{224.946}{189.401}{\Character{61}}%
\TextAt{13}{238.292}{189.401}{t}%
\TextAt{21}{242.52}{193.74}{n}%
\TextAt{13}{248.156}{189.401}{I}%
\TextAt{21}{253.319}{187.608}{r}%
\TextAt{21}{256.906}{187.608}{\Character{59}s}%
\TextAt{9}{263.671}{189.401}{\Character{40}1\Character{41}}%
\TextAt{8}{278.629}{189.401}{.}%
\TextAt{8}{288.646}{189.401}{Writing}%
\TextAt{13}{332.088}{189.401}{x}%
\TextAt{17}{338.74}{187.608}{2}%
\TextAt{9}{346.486}{189.401}{\Character{43}}%
\TextAt{16}{358.603}{189.401}{\Character{1}}%
\TextAt{16}{363.899}{189.401}{\Character{1}}%
\TextAt{16}{369.207}{189.401}{\Character{1}}%
\TextAt{9}{375.564}{189.401}{\Character{43}}%
\TextAt{13}{387.682}{189.401}{x}%
\TextAt{21}{394.334}{187.608}{r}%
\TextAt{9}{401.902}{189.401}{\Character{43}}%
\TextAt{13}{414.019}{189.401}{w}%
\TextAt{17}{422.387}{187.608}{1}%
\TextAt{9}{430.133}{189.401}{\Character{43}}%
\TextAt{16}{442.25}{189.401}{\Character{1}}%
\TextAt{16}{447.546}{189.401}{\Character{1}}%
\TextAt{16}{452.854}{189.401}{\Character{1}}%
\TextAt{9}{459.211}{189.401}{\Character{43}}%
\TextAt{13}{471.329}{189.401}{w}%
\TextAt{21}{479.697}{187.608}{s}%
\NativeMathAt{488.323}{189.401}{11.955}{\leq}%
\TextAt{13}{501.833}{189.401}{t}%
\TextAt{16}{509.074}{189.401}{\Character{0}}%
\TextAt{13}{521.385}{189.401}{x}%
\TextAt{17}{528.037}{187.608}{1}%
\TextAt{8}{79.226}{172.355}{instead}%
\TextAt{8}{119.611}{172.355}{of}%
\TextAt{8}{132.941}{172.355}{(*),}%
\TextAt{8}{155.058}{172.355}{Fubini’s}%
\TextAt{8}{199.495}{172.355}{theorem}%
\TextAt{8}{245.02}{172.355}{yields}%
\TextAt{13}{114.595}{139.218}{I}%
\TextAt{21}{119.758}{137.425}{r}%
\TextAt{21}{123.345}{137.425}{\Character{59}s}%
\TextAt{9}{130.111}{139.218}{\Character{40}1\Character{41}}%
\TextAt{9}{148.355}{139.218}{\Character{61}}%
\TextAt{20}{160.815}{155.491}{\Character{90}}%
\TextAt{17}{172.77}{152.336}{1}%
\TextAt{17}{167.457}{128.425}{0}%
\TextAt{13}{179.495}{139.218}{I}%
\TextAt{21}{184.658}{137.425}{r}%
\TextAt{3}{188.714}{137.425}{\Character{0}}%
\TextAt{17}{195.301}{137.425}{1}%
\TextAt{21}{199.535}{137.425}{\Character{59}s}%
\TextAt{9}{206.301}{139.218}{\Character{40}1}%
\TextAt{16}{219.363}{139.218}{\Character{0}}%
\TextAt{13}{231.319}{139.218}{x}%
\TextAt{17}{237.971}{137.425}{1}%
\TextAt{9}{242.703}{139.218}{\Character{41}}%
\TextAt{8}{249.248}{139.218}{d}%
\TextAt{13}{255.751}{139.218}{x}%
\TextAt{17}{262.403}{137.425}{1}%
\TextAt{9}{270.457}{139.218}{\Character{61}}%
\TextAt{20}{282.882}{155.491}{\Character{90}}%
\TextAt{17}{294.837}{152.336}{1}%
\TextAt{17}{289.524}{128.425}{0}%
\TextAt{9}{299.57}{139.218}{\Character{40}1}%
\TextAt{16}{312.632}{139.218}{\Character{0}}%
\TextAt{13}{324.587}{139.218}{x}%
\TextAt{17}{331.239}{137.425}{1}%
\TextAt{9}{335.971}{139.218}{\Character{41}}%
\TextAt{21}{340.523}{144.155}{n}%
\TextAt{3}{345.662}{144.155}{\Character{0}}%
\TextAt{17}{352.248}{144.155}{1}%
\TextAt{8}{358.973}{139.218}{d}%
\TextAt{13}{365.477}{139.218}{x}%
\TextAt{17}{372.129}{137.425}{1}%
\TextAt{16}{379.518}{139.218}{\Character{1}}%
\TextAt{13}{385.495}{139.218}{I}%
\TextAt{21}{390.658}{137.425}{r}%
\TextAt{3}{394.715}{137.425}{\Character{0}}%
\TextAt{17}{401.301}{137.425}{1}%
\TextAt{21}{405.535}{137.425}{\Character{59}s}%
\TextAt{9}{412.302}{139.218}{\Character{40}1\Character{41}}%
\TextAt{9}{430.546}{139.218}{\Character{61}}%
\TextAt{9}{444.769}{147.306}{1}%
\Vector{q 0 G 0 g 0.478 w 372.201 -631.793 m 379.189 -631.793 l S Q}%
\TextAt{13}{444.201}{131.018}{n}%
\TextAt{13}{452.384}{139.218}{I}%
\TextAt{21}{457.547}{137.425}{r}%
\TextAt{3}{461.604}{137.425}{\Character{0}}%
\TextAt{17}{468.19}{137.425}{1}%
\TextAt{21}{472.425}{137.425}{\Character{59}s}%
\TextAt{9}{479.191}{139.218}{\Character{40}1\Character{41}}%
\TextAt{13}{494.149}{139.218}{\Character{58}}%
\TextAt{8}{79.226}{108.246}{By}%
\TextAt{8}{97.589}{108.246}{induction,}%
\TextAt{8}{152.87}{108.246}{this}%
\TextAt{8}{175.692}{108.246}{implies}%
\TextAt{8}{215.419}{108.246}{that}%
\TextAt{13}{205.48}{78.133}{I}%
\TextAt{21}{210.643}{76.34}{r}%
\TextAt{21}{214.23}{76.34}{\Character{59}s}%
\TextAt{9}{220.995}{78.133}{\Character{40}1\Character{41}}%
\TextAt{9}{239.239}{78.133}{\Character{61}}%
\TextAt{9}{308.995}{86.221}{1}%
\Vector{q 0 G 0 g 0.478 w 180.895 -692.878 m 298.948 -692.878 l S Q}%
\TextAt{13}{252.895}{69.933}{n}%
\TextAt{9}{259.882}{69.933}{\Character{40}}%
\TextAt{13}{264.435}{69.933}{n}%
\TextAt{16}{274.079}{69.933}{\Character{0}}%
\TextAt{9}{286.034}{69.933}{1\Character{41}}%
\TextAt{16}{298.432}{69.933}{\Character{1}}%
\TextAt{16}{303.728}{69.933}{\Character{1}}%
\TextAt{16}{309.036}{69.933}{\Character{1}}%
\TextAt{9}{314.372}{69.933}{\Character{40}}%
\TextAt{13}{318.924}{69.933}{n}%
\TextAt{16}{328.569}{69.933}{\Character{0}}%
\TextAt{13}{340.524}{69.933}{r}%
\TextAt{9}{348.781}{69.933}{\Character{43}}%
\TextAt{9}{360.533}{69.933}{1\Character{41}}%
\TextAt{13}{372.143}{78.133}{I}%
\TextAt{17}{377.306}{76.34}{0}%
\TextAt{21}{381.54}{76.34}{\Character{59}s}%
\TextAt{9}{388.307}{78.133}{\Character{40}1\Character{41}}%
\TextAt{13}{403.264}{78.133}{\Character{58}}%
\TextAt{8}{297.211}{43.033}{187}%
\Destination{522.36}{234.375}{nav.AMS.24}%
\Destination{78.23}{809.815}{nav.page.187}%
\Destination{79.36}{654.555}{nav.print.205}%
\Destination{79.36}{775.155}{nav.number.3.2.15}%
\Hotlink{159.049}{620.255}{5.22}{11.047}{nav.chapter.1}%
\Hotlink{170.418}{620.255}{25.078}{11.047}{nav.number.1.5.1}%
}

% Page 200 (body page 188)
\Page{%
\TextAt{2}{79.226}{796.943}{III.}%
\TextAt{2}{98.317}{796.943}{Riemann–Roch}%
\TextAt{2}{174.375}{796.943}{Theory}%
\Vector{q 0 G 0 g 0.299 w 7.226 18.729 m 460.769 18.729 l S Q}%
\TextAt{8}{79.226}{761.124}{In}%
\TextAt{8}{93.847}{761.124}{the}%
\TextAt{8}{114.004}{761.124}{same}%
\TextAt{8}{143.33}{761.124}{way,}%
\TextAt{8}{169.344}{761.124}{one}%
\TextAt{8}{190.804}{761.124}{gets}%
\TextAt{13}{202.21}{728.259}{I}%
\TextAt{17}{207.373}{726.466}{0}%
\TextAt{21}{211.607}{726.466}{\Character{59}s}%
\TextAt{9}{218.373}{728.259}{\Character{40}1\Character{41}}%
\TextAt{9}{236.617}{728.259}{\Character{61}}%
\TextAt{20}{249.077}{744.532}{\Character{90}}%
\TextAt{17}{261.033}{741.377}{1}%
\TextAt{17}{255.719}{717.466}{0}%
\TextAt{13}{267.757}{728.259}{w}%
\TextAt{17}{276.125}{726.466}{1}%
\TextAt{9}{280.857}{728.259}{\Character{40}1}%
\TextAt{16}{293.919}{728.259}{\Character{0}}%
\TextAt{13}{305.875}{728.259}{w}%
\TextAt{17}{314.242}{726.466}{1}%
\TextAt{9}{318.975}{728.259}{\Character{41}}%
\TextAt{17}{323.527}{733.196}{2}%
\TextAt{21}{327.761}{733.196}{s}%
\TextAt{3}{331.677}{733.196}{\Character{0}}%
\TextAt{17}{338.264}{733.196}{2}%
\TextAt{8}{344.989}{728.259}{d}%
\TextAt{13}{351.492}{728.259}{w}%
\TextAt{17}{359.86}{726.466}{1}%
\TextAt{13}{364.592}{728.259}{I}%
\TextAt{17}{369.755}{726.466}{0}%
\TextAt{21}{373.989}{726.466}{\Character{59}s}%
\TextAt{3}{380.258}{726.466}{\Character{0}}%
\TextAt{17}{386.844}{726.466}{1}%
\TextAt{9}{391.576}{728.259}{\Character{40}1\Character{41}}%
\TextAt{13}{406.534}{728.259}{\Character{59}}%
\TextAt{8}{79.226}{697.01}{and,}%
\TextAt{8}{105.241}{697.01}{doing}%
\TextAt{8}{137.113}{697.01}{the}%
\TextAt{8}{157.269}{697.01}{integration,}%
\TextAt{8}{220.046}{697.01}{induction}%
\TextAt{8}{272.075}{697.01}{shows}%
\TextAt{8}{305.693}{697.01}{that}%
\TextAt{13}{231.231}{666.607}{I}%
\TextAt{17}{236.394}{664.814}{0}%
\TextAt{21}{240.628}{664.814}{\Character{59}s}%
\TextAt{9}{247.395}{666.607}{\Character{40}1\Character{41}}%
\TextAt{9}{265.639}{666.607}{\Character{61}}%
\TextAt{9}{288.229}{674.695}{1}%
\Vector{q 0 G 0 g 0.478 w 207.294 -104.404 m 231.017 -104.404 l S Q}%
\TextAt{9}{279.294}{658.407}{\Character{40}2}%
\TextAt{13}{289.699}{658.407}{s}%
\TextAt{9}{295.213}{658.407}{\Character{41}\Character{33}}%
\TextAt{13}{304.213}{666.607}{I}%
\TextAt{17}{309.376}{664.814}{0}%
\TextAt{21}{313.61}{664.814}{\Character{59}}%
\TextAt{17}{315.962}{664.814}{0}%
\TextAt{9}{320.695}{666.607}{\Character{40}1\Character{41}}%
\TextAt{9}{338.939}{666.607}{\Character{61}}%
\TextAt{9}{361.529}{674.695}{1}%
\Vector{q 0 G 0 g 0.478 w 280.594 -104.404 m 304.317 -104.404 l S Q}%
\TextAt{9}{352.594}{658.407}{\Character{40}2}%
\TextAt{13}{362.999}{658.407}{s}%
\TextAt{9}{368.513}{658.407}{\Character{41}\Character{33}}%
\TextAt{13}{377.513}{666.607}{\Character{58}}%
\TextAt{8}{79.226}{634.292}{Together,}%
\TextAt{8}{131.578}{634.292}{this}%
\TextAt{8}{154.4}{634.292}{gives}%
\TextAt{13}{183.081}{634.292}{I}%
\TextAt{21}{188.244}{632.499}{r}%
\TextAt{21}{191.831}{632.499}{\Character{59}s}%
\TextAt{9}{198.596}{634.292}{\Character{40}1\Character{41}}%
\TextAt{9}{216.84}{634.292}{\Character{61}}%
\TextAt{9}{229.261}{634.292}{1\Character{47}}%
\TextAt{13}{241.006}{634.292}{n}%
\TextAt{9}{247.994}{634.292}{\Character{33}}%
\TextAt{8}{255.143}{634.292}{and}%
\TextAt{8}{277.905}{634.292}{therefore}%
\TextAt{8}{327.005}{634.292}{indeed}%
\TextAt{8}{158.358}{602.932}{vol}%
\TextAt{9}{173.648}{602.932}{\Character{40}}%
\TextAt{13}{178.201}{602.932}{X}%
\TextAt{21}{187.916}{601.138}{t}%
\TextAt{9}{191.472}{602.932}{\Character{41}}%
\TextAt{9}{199.327}{602.932}{\Character{61}}%
\TextAt{9}{211.748}{602.932}{2}%
\TextAt{21}{217.624}{607.868}{s}%
\TextAt{8}{224.03}{602.932}{Vol}%
\TextAt{9}{241.915}{602.932}{\Character{40}}%
\TextAt{13}{246.468}{602.932}{f}%
\TextAt{9}{253.514}{602.932}{\Character{40}}%
\TextAt{13}{258.066}{602.932}{X}%
\TextAt{21}{267.782}{601.138}{t}%
\TextAt{9}{271.338}{602.932}{\Character{41}\Character{41}}%
\TextAt{9}{283.736}{602.932}{\Character{61}}%
\TextAt{9}{296.157}{602.932}{2}%
\TextAt{21}{302.042}{607.868}{s}%
\TextAt{9}{306.456}{602.932}{2}%
\TextAt{21}{312.309}{607.868}{r}%
\TextAt{9}{316.863}{602.932}{4}%
\TextAt{3}{322.716}{607.868}{\Character{0}}%
\TextAt{21}{329.303}{607.868}{s}%
\TextAt{9}{333.717}{602.932}{\Character{40}2}%
\TextAt{13}{344.122}{602.932}{\Character{25}}%
\TextAt{9}{351.191}{602.932}{\Character{41}}%
\TextAt{21}{355.744}{607.868}{s}%
\TextAt{13}{360.158}{602.932}{t}%
\TextAt{21}{364.385}{607.868}{n}%
\TextAt{13}{370.021}{602.932}{I}%
\TextAt{21}{375.184}{601.138}{r}%
\TextAt{21}{378.771}{601.138}{\Character{59}s}%
\TextAt{9}{385.537}{602.932}{\Character{40}1\Character{41}}%
\TextAt{9}{403.781}{602.932}{\Character{61}}%
\TextAt{9}{417.436}{611.019}{2}%
\TextAt{21}{423.289}{615.358}{r}%
\TextAt{13}{427.844}{611.019}{\Character{25}}%
\TextAt{21}{434.913}{615.358}{s}%
\Vector{q 0 G 0 g 0.478 w 345.436 -168.08 m 367.327 -168.08 l S Q}%
\TextAt{13}{423.262}{594.731}{n}%
\TextAt{9}{430.25}{594.731}{\Character{33}}%
\TextAt{13}{440.523}{602.932}{t}%
\TextAt{21}{444.75}{607.868}{n}%
\TextAt{13}{450.386}{602.932}{\Character{58}}%
\Vector{q 0 G 0 g 0.399 w 452.334 -199.278 m 452.334 -191.208 l S Q}%
\Vector{q 0 G 0 g 0.399 w 452.534 -191.407 m 459.707 -191.407 l S Q}%
\Vector{q 0 G 0 g 0.399 w 452.534 -199.079 m 459.707 -199.079 l S Q}%
\Vector{q 0 G 0 g 0.399 w 459.906 -199.278 m 459.906 -191.208 l S Q}%
\TextAt{8}{97.159}{545.947}{If}%
\TextAt{8}{108.851}{545.947}{we}%
\TextAt{8}{126.078}{545.947}{combine}%
\TextAt{8}{171.927}{545.947}{Stirling’s}%
\TextAt{8}{221.421}{545.947}{formula,}%
\TextAt{13}{203.468}{517.738}{n}%
\TextAt{9}{210.456}{517.738}{\Character{33}}%
\TextAt{9}{217.007}{517.738}{\Character{61}}%
\TextAt{16}{229.454}{528.211}{\Character{112}}%
\Vector{q 0 G 0 g 0.478 w 167.416 -245.55 m 187.326 -245.55 l S Q}%
\TextAt{9}{239.416}{517.738}{2}%
\TextAt{13}{245.269}{517.738}{\Character{25}}%
\TextAt{13}{252.323}{517.738}{n}%
\TextAt{20}{261.319}{531.008}{\Character{16}}%
\TextAt{13}{269.654}{525.826}{n}%
\Vector{q 0 G 0 g 0.478 w 197.654 -253.273 m 204.642 -253.273 l S Q}%
\TextAt{13}{270.435}{509.537}{e}%
\TextAt{20}{277.837}{531.008}{\Character{17}}%
\TextAt{21}{284.977}{528.332}{n}%
\TextAt{13}{292.606}{517.738}{e}%
\TextAt{21}{298.031}{522.674}{\Character{18}}%
\TextAt{17}{302.217}{522.674}{\Character{47}\Character{40}12}%
\TextAt{21}{318.213}{522.674}{n}%
\TextAt{17}{323.351}{522.674}{\Character{41}}%
\TextAt{13}{327.142}{517.738}{\Character{59}}%
\TextAt{9}{356.297}{517.738}{0}%
\TextAt{13}{365.471}{517.738}{\Character{60}}%
\TextAt{13}{377.881}{517.738}{\Character{18}}%
\TextAt{13}{386.978}{517.738}{\Character{60}}%
\TextAt{9}{399.423}{517.738}{1}%
\TextAt{13}{405.276}{517.738}{\Character{59}}%
\TextAt{8}{79.226}{488.411}{with}%
\TextAt{8}{105.886}{488.411}{the}%
\TextAt{8}{126.043}{488.411}{inequality}%
\TextAt{8}{180.355}{488.411}{(2.14),}%
\TextAt{8}{217.44}{488.411}{we}%
\TextAt{8}{234.667}{488.411}{obtain}%
\TextAt{8}{271.095}{488.411}{the}%
\TextAt{8}{291.251}{488.411}{inequality}%
\TextAt{16}{234.54}{457.459}{\Character{106}}%
\TextAt{13}{237.86}{457.459}{d}%
\TextAt{21}{243.943}{455.666}{K}%
\TextAt{16}{252.146}{457.459}{\Character{106}}%
\TextAt{13}{258.787}{457.459}{\Character{62}}%
\TextAt{20}{271.213}{470.729}{\Character{16}}%
\TextAt{13}{279.548}{465.547}{\Character{25}}%
\Vector{q 0 G 0 g 0.478 w 207.548 -313.552 m 214.618 -313.552 l S Q}%
\TextAt{9}{280.157}{449.258}{4}%
\TextAt{20}{287.813}{470.729}{\Character{17}}%
\TextAt{17}{294.953}{468.053}{2}%
\TextAt{21}{299.187}{468.053}{s}%
\TextAt{9}{313.818}{465.547}{1}%
\Vector{q 0 G 0 g 0.478 w 234.789 -313.552 m 254.699 -313.552 l S Q}%
\TextAt{9}{306.789}{449.258}{2}%
\TextAt{13}{312.642}{449.258}{\Character{25}}%
\TextAt{13}{319.696}{449.258}{n}%
\TextAt{13}{327.895}{457.459}{e}%
\TextAt{17}{333.32}{462.395}{2}%
\TextAt{21}{337.554}{462.395}{n}%
\TextAt{3}{342.693}{462.395}{\Character{0}}%
\TextAt{17}{349.279}{462.395}{1\Character{47}\Character{40}6}%
\TextAt{21}{365.275}{462.395}{n}%
\TextAt{17}{370.413}{462.395}{\Character{41}}%
\TextAt{13}{374.204}{457.459}{\Character{58}}%
\TextAt{8}{79.226}{429.25}{This}%
\TextAt{8}{105.658}{429.25}{shows}%
\TextAt{8}{138.988}{429.25}{that}%
\TextAt{8}{164.069}{429.25}{the}%
\TextAt{8}{183.938}{429.25}{absolute}%
\TextAt{8}{229.893}{429.25}{value}%
\TextAt{8}{259.852}{429.25}{of}%
\TextAt{8}{272.894}{429.25}{the}%
\TextAt{8}{292.763}{429.25}{discriminant}%
\TextAt{8}{359.854}{429.25}{of}%
\TextAt{8}{372.896}{429.25}{an}%
\TextAt{8}{388.868}{429.25}{algebraic}%
\TextAt{8}{438.014}{429.25}{number}%
\TextAt{8}{480.323}{429.25}{field}%
\TextAt{8}{505.392}{429.25}{tends}%
\TextAt{8}{79.226}{412.204}{to}%
\TextAt{8}{92.991}{412.204}{infinity}%
\TextAt{8}{132.77}{412.204}{with}%
\TextAt{8}{158.884}{412.204}{the}%
\TextAt{8}{178.495}{412.204}{degree.}%
\TextAt{8}{219.279}{412.204}{In}%
\TextAt{8}{233.354}{412.204}{the}%
\TextAt{8}{252.965}{412.204}{proof}%
\TextAt{8}{282.989}{412.204}{of}%
\TextAt{8}{295.773}{412.204}{(2.13)}%
\TextAt{8}{329.06}{412.204}{we}%
\TextAt{8}{345.742}{412.204}{saw}%
\TextAt{8}{367.696}{412.204}{that}%
\TextAt{8}{392.519}{412.204}{there}%
\TextAt{8}{421.885}{412.204}{are}%
\TextAt{8}{440.85}{412.204}{only}%
\TextAt{8}{465.996}{412.204}{finitely}%
\TextAt{8}{504.794}{412.204}{many}%
\TextAt{8}{79.226}{395.158}{number}%
\TextAt{8}{121.048}{395.158}{fields}%
\TextAt{8}{150.246}{395.158}{with}%
\TextAt{8}{176.132}{395.158}{bounded}%
\TextAt{8}{223.154}{395.158}{degree}%
\TextAt{8}{258.796}{395.158}{and}%
\TextAt{8}{280.784}{395.158}{discriminant.}%
\TextAt{8}{352.456}{395.158}{So}%
\TextAt{8}{367.941}{395.158}{now,}%
\TextAt{8}{394.962}{395.158}{since}%
\TextAt{8}{422.857}{395.158}{the}%
\TextAt{8}{442.239}{395.158}{degree}%
\TextAt{8}{477.88}{395.158}{is}%
\TextAt{8}{488.87}{395.158}{bounded}%
\TextAt{8}{79.226}{378.112}{if}%
\TextAt{8}{89.95}{378.112}{the}%
\TextAt{8}{110.106}{378.112}{discriminant}%
\TextAt{8}{177.486}{378.112}{is,}%
\TextAt{8}{192.501}{378.112}{we}%
\TextAt{8}{209.729}{378.112}{may}%
\TextAt{8}{235.098}{378.112}{strengthen}%
\TextAt{8}{293.044}{378.112}{(2.13),}%
\TextAt{8}{330.129}{378.112}{obtaining}%
\TextAt{12}{79.226}{352.326}{Hermite’s}%
\TextAt{12}{140.314}{352.326}{Theorem}%
\TextAt{12}{196.405}{352.326}{(2.16).}%
\TextAt{8}{237.648}{352.326}{There}%
\TextAt{8}{270.654}{352.326}{exist}%
\TextAt{8}{297.551}{352.326}{only}%
\TextAt{8}{322.44}{352.326}{finitely}%
\TextAt{8}{360.982}{352.326}{many}%
\TextAt{8}{392.051}{352.326}{number}%
\TextAt{8}{433.844}{352.326}{fields}%
\TextAt{8}{463.013}{352.326}{with}%
\TextAt{8}{488.87}{352.326}{bounded}%
\TextAt{8}{79.226}{335.28}{discriminant.}%
\TextAt{8}{97.159}{309.493}{The}%
\TextAt{8}{121.213}{309.493}{expression}%
\TextAt{13}{177.593}{309.493}{a}%
\TextAt{21}{183.738}{307.7}{n}%
\TextAt{9}{192.695}{309.493}{\Character{61}}%
\TextAt{9}{205.105}{309.493}{\Character{40}}%
\TextAt{13}{209.673}{309.493}{n}%
\TextAt{21}{216.661}{313.832}{n}%
\TextAt{9}{222.297}{309.493}{\Character{47}}%
\TextAt{13}{228.15}{309.493}{n}%
\TextAt{9}{235.138}{309.493}{\Character{33}\Character{41}\Character{40}}%
\TextAt{13}{247.494}{309.493}{\Character{25}}%
\TextAt{9}{254.563}{309.493}{\Character{47}4\Character{41}}%
\TextAt{21}{270.822}{313.832}{n}%
\TextAt{17}{275.96}{313.832}{\Character{47}2}%
\TextAt{8}{288.824}{309.493}{satisfies}%
\TextAt{13}{228.488}{284.821}{a}%
\TextAt{21}{234.633}{283.028}{n}%
\TextAt{17}{239.771}{283.028}{\Character{43}1}%
\Vector{q 0 G 0 g 0.478 w 156.488 -494.278 m 179.09 -494.278 l S Q}%
\TextAt{13}{233.899}{268.533}{a}%
\TextAt{21}{240.043}{266.739}{n}%
\TextAt{9}{255.607}{276.733}{\Character{61}}%
\TextAt{20}{268.032}{290.004}{\Character{16}}%
\TextAt{13}{276.367}{284.821}{\Character{25}}%
\Vector{q 0 G 0 g 0.478 w 204.367 -494.278 m 211.437 -494.278 l S Q}%
\TextAt{9}{276.976}{268.533}{4}%
\TextAt{20}{284.632}{290.004}{\Character{17}}%
\TextAt{17}{291.772}{287.327}{1\Character{47}2}%
\TextAt{20}{306.965}{293.59}{\Character{18}}%
\TextAt{9}{315.766}{276.733}{1}%
\TextAt{9}{324.266}{276.733}{\Character{43}}%
\TextAt{9}{337.799}{284.821}{1}%
\Vector{q 0 G 0 g 0.478 w 265.232 -494.278 m 272.22 -494.278 l S Q}%
\TextAt{13}{337.232}{268.533}{n}%
\TextAt{20}{345.415}{293.59}{\Character{19}}%
\TextAt{21}{354.216}{290.913}{n}%
\TextAt{13}{363.173}{276.733}{\Character{62}}%
\TextAt{9}{375.598}{276.733}{1}%
\TextAt{13}{381.451}{276.733}{\Character{59}}%
\TextAt{8}{79.226}{244.25}{i.e.,}%
\TextAt{13}{101.331}{244.25}{a}%
\TextAt{21}{107.476}{242.456}{n}%
\TextAt{17}{112.614}{242.456}{\Character{43}1}%
\TextAt{13}{127.254}{244.25}{\Character{62}}%
\TextAt{13}{139.664}{244.25}{a}%
\TextAt{21}{145.824}{242.456}{n}%
\TextAt{8}{151.461}{244.25}{.}%
\TextAt{8}{159.909}{244.25}{Since}%
\TextAt{13}{190.467}{244.25}{a}%
\TextAt{17}{196.611}{242.456}{2}%
\TextAt{9}{204.665}{244.25}{\Character{61}}%
\TextAt{13}{217.09}{244.25}{\Character{25}}%
\TextAt{9}{224.159}{244.25}{\Character{47}2}%
\TextAt{13}{239.186}{244.25}{\Character{62}}%
\TextAt{9}{251.612}{244.25}{1}%
\TextAt{8}{257.465}{244.25}{,}%
\TextAt{8}{264.614}{244.25}{(2.14)}%
\TextAt{8}{298.447}{244.25}{yields}%
\TextAt{12}{79.226}{218.464}{Minkowski’s}%
\TextAt{12}{155.478}{218.464}{Theorem}%
\TextAt{12}{212.279}{218.464}{(2.17).}%
\TextAt{8}{254.14}{218.464}{The}%
\TextAt{8}{278.009}{218.464}{discriminant}%
\TextAt{8}{345.203}{218.464}{of}%
\TextAt{8}{358.348}{218.464}{a}%
\TextAt{8}{367.918}{218.464}{number}%
\TextAt{8}{410.329}{218.464}{field}%
\TextAt{13}{435.501}{218.464}{K}%
\TextAt{8}{450.03}{218.464}{different}%
\TextAt{8}{496.016}{218.464}{from}%
\TextAt{18}{523.471}{218.464}{Q}%
\TextAt{8}{79.226}{201.418}{is}%
\NativeMathAt{90.99}{201.418}{11.955}{\neq}%
\NativeMathAt{103.416}{201.418}{11.955}{\pm}%
\TextAt{9}{112.714}{201.418}{1}%
\TextAt{8}{118.567}{201.418}{.}%
\TextAt{8}{97.159}{175.631}{Combining}%
\TextAt{8}{156.672}{175.631}{this}%
\TextAt{8}{179.494}{175.631}{result}%
\TextAt{8}{212.072}{175.631}{with}%
\TextAt{8}{238.732}{175.631}{corollary}%
\TextAt{8}{287.198}{175.631}{(2.12),}%
\TextAt{8}{324.283}{175.631}{we}%
\TextAt{8}{341.511}{175.631}{obtain}%
\TextAt{8}{377.938}{175.631}{the}%
\TextAt{12}{79.226}{149.845}{Theorem}%
\TextAt{12}{136.24}{149.845}{(2.18).}%
\TextAt{8}{178.287}{149.845}{The}%
\TextAt{8}{202.34}{149.845}{field}%
\TextAt{18}{227.697}{149.845}{Q}%
\TextAt{8}{240.893}{149.845}{does}%
\TextAt{8}{267.29}{149.845}{not}%
\TextAt{8}{288.104}{149.845}{admit}%
\TextAt{8}{321.925}{149.845}{any}%
\TextAt{8}{344.042}{149.845}{unramified}%
\TextAt{8}{402.575}{149.845}{extensions.}%
\TextAt{8}{97.159}{124.059}{These}%
\TextAt{8}{131.1}{124.059}{last}%
\TextAt{8}{153.349}{124.059}{theorems}%
\TextAt{8}{203.561}{124.059}{are}%
\TextAt{8}{223.144}{124.059}{of}%
\TextAt{8}{236.546}{124.059}{fundamental}%
\TextAt{8}{304.26}{124.059}{importance}%
\TextAt{8}{365.794}{124.059}{for}%
\TextAt{8}{383.751}{124.059}{number}%
\TextAt{8}{426.419}{124.059}{theory.}%
\TextAt{8}{466.956}{124.059}{Their}%
\TextAt{8}{498.889}{124.059}{signifi-}%
\TextAt{8}{79.226}{107.013}{cance}%
\TextAt{8}{110.775}{107.013}{is}%
\TextAt{8}{122.227}{107.013}{seen}%
\TextAt{8}{147.332}{107.013}{especially}%
\TextAt{8}{199.754}{107.013}{clearly}%
\TextAt{8}{236.838}{107.013}{in}%
\TextAt{8}{250.179}{107.013}{the}%
\TextAt{8}{270.024}{107.013}{light}%
\TextAt{8}{296.707}{107.013}{of}%
\TextAt{8}{309.725}{107.013}{higher}%
\TextAt{8}{345.183}{107.013}{dimensional}%
\TextAt{8}{409.322}{107.013}{analogues.}%
\TextAt{8}{467.172}{107.013}{For}%
\TextAt{8}{487.83}{107.013}{instance,}%
\TextAt{8}{79.226}{89.967}{let}%
\TextAt{8}{95.652}{89.967}{us}%
\TextAt{8}{110.19}{89.967}{replace}%
\TextAt{8}{149.379}{89.967}{the}%
\TextAt{8}{169.057}{89.967}{finite}%
\TextAt{8}{198.49}{89.967}{field}%
\TextAt{8}{223.369}{89.967}{extensions}%
\TextAt{13}{279.271}{89.967}{L}%
\TextAt{16}{290.556}{89.967}{\Character{106}}%
\TextAt{13}{297.198}{89.967}{K}%
\TextAt{8}{311.435}{89.967}{of}%
\TextAt{8}{324.286}{89.967}{a}%
\TextAt{8}{333.563}{89.967}{number}%
\TextAt{8}{375.681}{89.967}{field}%
\TextAt{13}{400.56}{89.967}{K}%
\TextAt{8}{414.796}{89.967}{by}%
\TextAt{8}{430.577}{89.967}{all}%
\TextAt{8}{446.358}{89.967}{smooth}%
\TextAt{8}{487.244}{89.967}{complete}%
\TextAt{8}{79.226}{72.921}{(i.e.,}%
\TextAt{8}{105.439}{72.921}{proper)}%
\TextAt{8}{146.831}{72.921}{algebraic}%
\TextAt{8}{195.707}{72.921}{curves}%
\TextAt{8}{230.978}{72.921}{defined}%
\TextAt{8}{270.732}{72.921}{over}%
\TextAt{13}{295.219}{72.921}{K}%
\TextAt{8}{309.375}{72.921}{of}%
\TextAt{8}{322.147}{72.921}{a}%
\TextAt{8}{331.343}{72.921}{fixed}%
\TextAt{8}{359.07}{72.921}{genus}%
\TextAt{13}{390.767}{72.921}{g}%
\TextAt{8}{396.801}{72.921}{.}%
\TextAt{8}{405.063}{72.921}{If}%
\TextAt{14}{416.196}{72.921}{p}%
\TextAt{8}{425.518}{72.921}{is}%
\TextAt{8}{436.723}{72.921}{a}%
\TextAt{8}{445.92}{72.921}{prime}%
\TextAt{8}{478.524}{72.921}{ideal}%
\TextAt{8}{505.929}{72.921}{of}%
\TextAt{13}{518.7}{72.921}{K}%
\TextAt{8}{529.518}{72.921}{,}%
\TextAt{8}{297.211}{43.033}{188}%
\Destination{78.23}{809.815}{nav.page.188}%
\Destination{79.36}{776.155}{nav.print.206}%
\Destination{79.36}{233.455}{nav.print.207}%
\Destination{79.36}{163.855}{nav.number.3.2.18}%
\Hotlink{179.855}{486.255}{30.936}{11.047}{nav.number.3.2.14}%
\Hotlink{295.273}{409.955}{30.936}{11.047}{nav.number.3.2.13}%
\Hotlink{292.544}{375.955}{30.936}{11.047}{nav.number.3.2.13}%
\Hotlink{264.114}{242.055}{30.936}{11.047}{nav.number.3.2.14}%
\Hotlink{286.698}{173.455}{30.936}{11.047}{nav.number.3.2.12}%
\Destination{79.23}{225.5}{nav.number.3.2.17}%
}

% Page 201 (body page 189)
\Page{%
\TextAt{2}{79.226}{796.943}{III.}%
\TextAt{2}{98.317}{796.943}{Riemann–Roch}%
\TextAt{2}{174.375}{796.943}{Theory}%
\Vector{q 0 G 0 g 0.299 w 7.226 18.729 m 460.769 18.729 l S Q}%
\TextAt{8}{79.226}{761.124}{then}%
\TextAt{8}{105.426}{761.124}{for}%
\TextAt{8}{122.852}{761.124}{any}%
\TextAt{8}{144.509}{761.124}{such}%
\TextAt{8}{170.446}{761.124}{curve}%
\TextAt{13}{201.202}{761.124}{X}%
\TextAt{8}{211.857}{761.124}{,}%
\TextAt{8}{218.638}{761.124}{one}%
\TextAt{8}{239.638}{761.124}{may}%
\TextAt{8}{264.547}{761.124}{define}%
\TextAt{8}{297.897}{761.124}{the}%
\TextAt{8}{317.593}{761.124}{“reduction}%
\TextAt{8}{374.686}{761.124}{mod}%
\TextAt{14}{400.563}{761.124}{p}%
\TextAt{8}{406.546}{761.124}{”.}%
\TextAt{8}{418.739}{761.124}{This}%
\TextAt{8}{444.999}{761.124}{is}%
\TextAt{8}{456.303}{761.124}{a}%
\TextAt{8}{465.599}{761.124}{curve}%
\TextAt{8}{496.354}{761.124}{defined}%
\TextAt{8}{79.226}{744.078}{over}%
\TextAt{8}{104.261}{744.078}{the}%
\TextAt{8}{124.407}{744.078}{residue}%
\TextAt{8}{164.123}{744.078}{class}%
\TextAt{8}{191.549}{744.078}{field}%
\TextAt{8}{216.895}{744.078}{of}%
\TextAt{14}{230.214}{744.078}{p}%
\TextAt{8}{236.197}{744.078}{.}%
\TextAt{8}{244.642}{744.078}{One}%
\TextAt{8}{269.342}{744.078}{says}%
\TextAt{8}{294.174}{744.078}{that}%
\TextAt{13}{319.532}{744.078}{X}%
\TextAt{8}{334.074}{744.078}{has}%
\TextAt{19}{354.936}{744.078}{good}%
\TextAt{19}{381.226}{744.078}{reduction}%
\TextAt{8}{432.168}{744.078}{at}%
\TextAt{8}{446.468}{744.078}{the}%
\TextAt{8}{466.613}{744.078}{prime}%
\TextAt{14}{499.766}{744.078}{p}%
\TextAt{8}{509.635}{744.078}{if}%
\TextAt{8}{520.348}{744.078}{its}%
\TextAt{8}{79.226}{727.032}{reduction}%
\TextAt{8}{131.353}{727.032}{mod}%
\TextAt{14}{157.788}{727.032}{p}%
\TextAt{8}{167.766}{727.032}{is}%
\TextAt{8}{179.627}{727.032}{again}%
\TextAt{8}{210.952}{727.032}{a}%
\TextAt{8}{220.805}{727.032}{smooth}%
\TextAt{8}{262.267}{727.032}{curve.}%
\TextAt{8}{298.326}{727.032}{This}%
\TextAt{8}{325.144}{727.032}{corresponds}%
\TextAt{8}{389.428}{727.032}{to}%
\TextAt{8}{403.836}{727.032}{an}%
\TextAt{8}{420.193}{727.032}{extension}%
\TextAt{13}{472.057}{727.032}{L}%
\TextAt{16}{483.508}{727.032}{\Character{106}}%
\TextAt{13}{490.316}{727.032}{K}%
\TextAt{8}{505.129}{727.032}{being}%
\TextAt{8}{79.226}{709.986}{unramified.}%
\TextAt{8}{142.654}{709.986}{In}%
\TextAt{8}{157.39}{709.986}{analogy}%
\TextAt{8}{200.771}{709.986}{to}%
\TextAt{8}{215.196}{709.986}{Hermite’s}%
\TextAt{8}{268.368}{709.986}{theorem,}%
\TextAt{8}{317.289}{709.986}{the}%
\TextAt{8}{337.56}{709.986}{Russian}%
\TextAt{8}{381.216}{709.986}{mathematician}%
\TextAt{8}{461.335}{709.986}{I.}%
\TextAt{8}{472.819}{709.986}{S.}%
\TextAt{8}{486.587}{709.986}{Šafarevič}%
\ErrataAt{79.19}{692.94}{455.81}{formulated the conjecture that up to $K$-isomorphism there are only finitely many smooth proper}%
\ErrataAt{79.19}{675.894}{455.81}{geometrically connected curves of genus $g\ge2$ over $K$ with good reduction outside $S$. This conjecture}%
\TextAt{8}{79.226}{658.848}{was}%
\TextAt{8}{102.462}{658.848}{proved}%
\TextAt{8}{141.252}{658.848}{in}%
\TextAt{8}{155.641}{658.848}{1983}%
\TextAt{8}{183.706}{658.848}{by}%
\TextAt{8}{200.702}{658.848}{the}%
\TextAt{8}{221.594}{658.848}{mathematician}%
\TextAt{8}{302.335}{658.848}{Gerd}%
\TextAt{8}{332.409}{658.848}{Faltings}%
\TextAt{8}{377.594}{658.848}{(see}%
\TextAt{8}{401.799}{658.848}{[35]).}%
\TextAt{8}{435.23}{658.848}{The}%
\TextAt{8}{460.02}{658.848}{impact}%
\TextAt{8}{499.778}{658.848}{of}%
\TextAt{8}{513.844}{658.848}{this}%
\TextAt{8}{79.226}{641.802}{result}%
\TextAt{8}{111.97}{641.802}{can}%
\TextAt{8}{133.595}{641.802}{be}%
\TextAt{8}{149.685}{641.802}{gauged}%
\TextAt{8}{189.531}{641.802}{by}%
\TextAt{8}{205.955}{641.802}{the}%
\TextAt{8}{226.278}{641.802}{non-expert}%
\TextAt{8}{285.622}{641.802}{from}%
\TextAt{8}{313.428}{641.802}{the}%
\TextAt{8}{333.75}{641.802}{fact}%
\TextAt{8}{357.002}{641.802}{that}%
\TextAt{8}{382.536}{641.802}{it}%
\TextAt{8}{394.406}{641.802}{was}%
\TextAt{8}{417.072}{641.802}{the}%
\TextAt{8}{437.394}{641.802}{basis}%
\TextAt{8}{466.3}{641.802}{for}%
\TextAt{8}{484.351}{641.802}{Faltings’s}%
\TextAt{8}{79.226}{624.756}{proof}%
\TextAt{8}{109.796}{624.756}{of}%
\TextAt{8}{123.126}{624.756}{the}%
\TextAt{8}{143.282}{624.756}{famous}%
\TextAt{12}{183.344}{624.756}{Mordell}%
\TextAt{12}{233.902}{624.756}{Conjecture}%
\TextAt{8}{298.388}{624.756}{:}%
\TextAt{8}{97.159}{604.721}{Every}%
\TextAt{8}{130.813}{604.721}{algebraic}%
\TextAt{8}{180.247}{604.721}{equation}%
\TextAt{13}{278.106}{587.675}{f}%
\TextAt{9}{285.153}{587.675}{\Character{40}}%
\TextAt{13}{289.705}{587.675}{x\Character{59}}%
\TextAt{13}{301.577}{587.675}{y}%
\TextAt{9}{307.738}{587.675}{\Character{41}}%
\TextAt{9}{315.593}{587.675}{\Character{61}}%
\TextAt{9}{328.014}{587.675}{0}%
\TextAt{8}{79.226}{564.153}{of}%
\TextAt{8}{92.556}{564.153}{genus}%
\TextAt{13}{124.811}{564.153}{g}%
\TextAt{13}{134.148}{564.153}{\Character{62}}%
\TextAt{9}{146.592}{564.153}{1}%
\TextAt{8}{156.342}{564.153}{with}%
\TextAt{8}{183.002}{564.153}{coefficients}%
\TextAt{8}{242.24}{564.153}{in}%
\TextAt{13}{255.893}{564.153}{K}%
\TextAt{8}{270.608}{564.153}{admits}%
\TextAt{8}{309.044}{564.153}{only}%
\TextAt{12}{334.735}{564.153}{finitely}%
\TextAt{12}{380.117}{564.153}{many}%
\TextAt{12}{416.557}{564.153}{solutions}%
\TextAt{8}{472.196}{564.153}{in}%
\TextAt{13}{485.849}{564.153}{K}%
\TextAt{8}{496.666}{564.153}{.}%
\TextAt{8}{97.159}{544.119}{A}%
\TextAt{8}{110.42}{544.119}{1-dimensional}%
\TextAt{8}{185.214}{544.119}{analogue}%
\TextAt{8}{234.592}{544.119}{of}%
\TextAt{8}{248.511}{544.119}{Minkowski’s}%
\TextAt{8}{315.236}{544.119}{theorem}%
\TextAt{8}{361.35}{544.119}{(2.18)}%
\TextAt{8}{395.772}{544.119}{was}%
\TextAt{8}{418.86}{544.119}{proved}%
\TextAt{8}{457.503}{544.119}{in}%
\TextAt{8}{471.744}{544.119}{1985}%
\TextAt{8}{499.663}{544.119}{by}%
\TextAt{8}{516.51}{544.119}{the}%
\TextAt{8}{79.226}{527.073}{French}%
\TextAt{8}{116.481}{527.073}{mathematician}%
\TextAt{8}{195.542}{527.073}{J.-M.}%
\TextAt{8}{225.636}{527.073}{Fontaine:}%
\TextAt{8}{277.68}{527.073}{over}%
\TextAt{8}{301.783}{527.073}{the}%
\TextAt{8}{320.998}{527.073}{field}%
\TextAt{18}{345.412}{527.073}{Q}%
\TextAt{8}{354.711}{527.073}{,}%
\TextAt{8}{361.106}{527.073}{there}%
\TextAt{8}{390.076}{527.073}{are}%
\TextAt{8}{408.644}{527.073}{no}%
\TextAt{8}{423.961}{527.073}{smooth}%
\TextAt{8}{464.384}{527.073}{proper}%
\TextAt{8}{500.837}{527.073}{curves}%
\EditionChange{\BodyAt{79.226}{510.055}{with good reduction mod $p$ for all prime numbers $p$ (see \hyperlink{nav.cite.N039}{[39]}).}}{\BodyAt{79.226}{510.055}{of positive genus with good reduction mod $p$ for all prime numbers $p$ (see \hyperlink{nav.cite.N039}{[39]}).}}%
\TextAt{12}{79.226}{484.014}{Exercise}%
\TextAt{12}{133.642}{484.014}{III.2.1.}%
\TextAt{8}{179.052}{484.014}{Let}%
\TextAt{13}{201.078}{484.014}{d}%
\TextAt{9}{207.16}{484.014}{\Character{40}}%
\TextAt{13}{211.713}{484.014}{\Character{11}}%
\TextAt{9}{219.234}{484.014}{\Character{41}}%
\TextAt{9}{228.894}{484.014}{\Character{61}}%
\TextAt{13}{243.133}{484.014}{d}%
\TextAt{9}{249.216}{484.014}{\Character{40}1}%
\TextAt{13}{259.621}{484.014}{\Character{59}}%
\TextAt{13}{264.845}{484.014}{\Character{11}\Character{59}}%
\TextAt{13}{277.59}{484.014}{\Character{58}}%
\TextAt{13}{282.826}{484.014}{\Character{58}}%
\TextAt{13}{288.062}{484.014}{\Character{58}}%
\TextAt{13}{293.287}{484.014}{\Character{59}}%
\TextAt{13}{298.523}{484.014}{\Character{11}}%
\TextAt{21}{306.13}{488.353}{n}%
\TextAt{3}{311.268}{488.353}{\Character{0}}%
\TextAt{17}{317.854}{488.353}{1}%
\TextAt{9}{322.587}{484.014}{\Character{41}}%
\TextAt{8}{327.139}{484.014}{,}%
\TextAt{8}{335.609}{484.014}{for}%
\TextAt{8}{354.55}{484.014}{an}%
\TextAt{8}{371.865}{484.014}{element}%
\TextAt{13}{416.163}{484.014}{\Character{11}}%
\NativeMathAt{428.806}{484.014}{11.955}{\in}%
\TextAt{16}{441.885}{484.014}{\Character{79}}%
\TextAt{8}{456.701}{484.014}{such}%
\TextAt{8}{484.154}{484.014}{that}%
\TextAt{13}{510.579}{484.014}{L}%
\TextAt{9}{523.665}{484.014}{\Character{61}}%
\TextAt{13}{79.226}{466.968}{K}%
\TextAt{9}{90.044}{466.968}{\Character{40}}%
\TextAt{13}{94.596}{466.968}{\Character{11}}%
\TextAt{9}{102.118}{466.968}{\Character{41}}%
\TextAt{8}{106.67}{466.968}{.}%
\TextAt{8}{114.908}{466.968}{Show}%
\TextAt{8}{145.168}{466.968}{that}%
\TextAt{14}{169.905}{466.968}{d}%
\TextAt{21}{175.865}{465.109}{L}%
\TextAt{3}{181.624}{465.109}{\Character{106}}%
\TextAt{21}{183.976}{465.109}{K}%
\TextAt{8}{195.444}{466.968}{is}%
\TextAt{8}{206.576}{466.968}{generated}%
\TextAt{8}{259.276}{466.968}{by}%
\TextAt{8}{274.903}{466.968}{all}%
\TextAt{8}{290.53}{466.968}{discriminants}%
\TextAt{8}{361.892}{466.968}{of}%
\TextAt{8}{374.589}{466.968}{elements}%
\TextAt{13}{421.814}{466.968}{d}%
\TextAt{9}{427.896}{466.968}{\Character{40}}%
\TextAt{13}{432.449}{466.968}{\Character{11}}%
\TextAt{9}{439.971}{466.968}{\Character{41}}%
\TextAt{8}{447.788}{466.968}{if}%
\TextAt{14}{457.88}{466.968}{o}%
\TextAt{8}{466.989}{466.968}{is}%
\TextAt{8}{478.121}{466.968}{a}%
\TextAt{8}{487.244}{466.968}{complete}%
\TextAt{8}{79.226}{449.922}{discrete}%
\TextAt{8}{122.334}{449.922}{valuation}%
\TextAt{8}{173.535}{449.922}{ring}%
\TextAt{8}{197.729}{449.922}{and}%
\TextAt{8}{220.621}{449.922}{the}%
\TextAt{8}{240.906}{449.922}{residue}%
\TextAt{8}{280.762}{449.922}{field}%
\TextAt{8}{306.247}{449.922}{extension}%
\TextAt{13}{358.142}{449.922}{\Character{21}}%
\TextAt{16}{368.51}{449.923}{\Character{106}}%
\TextAt{13}{375.372}{449.922}{\Character{20}}%
\TextAt{8}{386.136}{449.922}{is}%
\TextAt{8}{398.029}{449.922}{separable.}%
\TextAt{8}{454.409}{449.922}{In}%
\TextAt{8}{469.159}{449.922}{other}%
\TextAt{8}{499.857}{449.922}{words,}%
\TextAt{14}{79.226}{432.877}{d}%
\TextAt{21}{85.187}{431.017}{L}%
\TextAt{3}{90.946}{431.017}{\Character{106}}%
\TextAt{21}{93.298}{431.017}{K}%
\TextAt{8}{107.068}{432.877}{equals}%
\TextAt{8}{144.245}{432.877}{the}%
\TextAt{8}{166.072}{432.877}{gcd}%
\TextAt{8}{189.201}{432.877}{of}%
\TextAt{8}{204.201}{432.877}{all}%
\TextAt{8}{222.131}{432.877}{discriminants}%
\TextAt{8}{295.795}{432.877}{of}%
\TextAt{8}{310.795}{432.877}{individual}%
\TextAt{8}{367.423}{432.877}{elements.}%
\TextAt{8}{424.841}{432.877}{This}%
\TextAt{8}{453.231}{432.877}{fails}%
\TextAt{8}{479.349}{432.877}{to}%
\TextAt{8}{495.329}{432.877}{hold}%
\TextAt{8}{523.014}{432.877}{in}%
\TextAt{8}{79.226}{415.831}{general.}%
\TextAt{8}{124.109}{415.831}{Counterexample:}%
\TextAt{13}{215.818}{415.831}{K}%
\TextAt{9}{229.961}{415.831}{\Character{61}}%
\TextAt{18}{242.39}{415.831}{Q}%
\TextAt{8}{251.689}{415.831}{,}%
\TextAt{13}{258.841}{415.831}{L}%
\TextAt{9}{270.131}{415.831}{\Character{61}}%
\TextAt{18}{282.56}{415.831}{Q}%
\TextAt{9}{291.859}{415.831}{\Character{40}}%
\TextAt{13}{296.411}{415.831}{\Character{11}}%
\TextAt{9}{303.933}{415.831}{\Character{41}}%
\TextAt{8}{308.485}{415.831}{,}%
\TextAt{13}{315.637}{415.831}{\Character{11}}%
\TextAt{17}{323.159}{420.169}{3}%
\TextAt{16}{330.55}{415.831}{\Character{0}}%
\TextAt{13}{342.507}{415.831}{\Character{11}}%
\TextAt{17}{350.028}{420.169}{2}%
\TextAt{16}{357.419}{415.831}{\Character{0}}%
\TextAt{9}{369.376}{415.831}{2}%
\TextAt{13}{375.229}{415.831}{\Character{11}}%
\TextAt{16}{385.409}{415.831}{\Character{0}}%
\TextAt{9}{397.366}{415.831}{8}%
\TextAt{9}{406.536}{415.831}{\Character{61}}%
\TextAt{9}{418.957}{415.831}{0}%
\TextAt{8}{424.826}{415.831}{.}%
\TextAt{8}{433.282}{415.831}{(See}%
\TextAt{8}{458.641}{415.831}{[60],}%
\TextAt{8}{484.013}{415.831}{chap.}%
\TextAt{8}{516.212}{415.831}{III,}%
\TextAt{8}{79.226}{398.785}{§25,}%
\TextAt{8}{104.443}{398.785}{p.443.}%
\TextAt{8}{141.836}{398.785}{The}%
\TextAt{8}{166.428}{398.785}{untranslatable}%
\TextAt{8}{244.113}{398.785}{German}%
\TextAt{8}{289.603}{398.785}{catch}%
\TextAt{8}{321.033}{398.785}{phrase}%
\TextAt{8}{358.704}{398.785}{for}%
\TextAt{8}{377.127}{398.785}{this}%
\TextAt{8}{400.488}{398.785}{phenomenon}%
\TextAt{8}{469.314}{398.785}{is:}%
\TextAt{8}{486.706}{398.785}{there}%
\TextAt{8}{517.156}{398.785}{are}%
\TextAt{8}{79.226}{381.739}{“außerwesentliche}%
\TextAt{8}{173.588}{381.739}{Diskriminantenteiler”.)}%
\TextAt{12}{79.226}{349.749}{Exercise}%
\TextAt{12}{133.959}{349.749}{III.2.2.}%
\TextAt{8}{179.645}{349.749}{Let}%
\TextAt{13}{201.946}{349.749}{L}%
\TextAt{16}{215.502}{349.749}{\Character{106}}%
\TextAt{13}{224.414}{349.749}{K}%
\TextAt{8}{240.461}{349.749}{be}%
\TextAt{8}{257.717}{349.749}{a}%
\TextAt{8}{268.805}{349.749}{Galois}%
\TextAt{8}{306.05}{349.749}{extension}%
\TextAt{8}{359.148}{349.749}{of}%
\TextAt{8}{373.81}{349.749}{henselian}%
\TextAt{8}{425.928}{349.749}{fields}%
\TextAt{8}{457.232}{349.749}{with}%
\TextAt{8}{485.224}{349.749}{separable}%
\TextAt{8}{79.226}{332.703}{residue}%
\TextAt{8}{119.428}{332.703}{field}%
\TextAt{8}{145.26}{332.703}{extension}%
\TextAt{13}{197.501}{332.703}{\Character{21}}%
\TextAt{16}{208.459}{332.703}{\Character{106}}%
\TextAt{13}{215.91}{332.703}{\Character{20}}%
\TextAt{8}{222.649}{332.703}{,}%
\TextAt{8}{230.391}{332.703}{and}%
\TextAt{8}{253.629}{332.703}{let}%
\TextAt{13}{271.009}{332.703}{G}%
\TextAt{21}{280.242}{330.909}{i}%
\TextAt{8}{283.623}{332.703}{,}%
\TextAt{13}{291.366}{332.703}{i}%
\NativeMathAt{299.49}{332.703}{11.955}{\geq}%
\TextAt{9}{312.918}{332.703}{0}%
\TextAt{8}{318.771}{332.703}{,}%
\TextAt{8}{326.514}{332.703}{be}%
\TextAt{8}{342.913}{332.703}{the}%
\TextAt{13}{363.545}{332.703}{i}%
\TextAt{8}{367.538}{332.703}{-th}%
\TextAt{8}{386.866}{332.703}{ramification}%
\TextAt{8}{452.389}{332.703}{group.}%
\TextAt{8}{491.54}{332.703}{Then,}%
\TextAt{8}{525.943}{332.703}{if}%
\TextAt{14}{79.226}{315.657}{D}%
\TextAt{21}{89.168}{313.797}{L}%
\TextAt{3}{94.927}{313.797}{\Character{106}}%
\TextAt{21}{97.279}{313.797}{K}%
\TextAt{9}{108.803}{315.657}{\Character{61}}%
\TextAt{14}{121.228}{315.657}{P}%
\TextAt{21}{131.121}{319.995}{s}%
\TextAt{8}{135.535}{315.657}{,}%
\TextAt{8}{142.685}{315.657}{one}%
\TextAt{8}{164.144}{315.657}{has}%
\TextAt{13}{259.138}{291.237}{s}%
\TextAt{9}{267.972}{291.237}{\Character{61}}%
\NativeMathAt{284.798}{306.181}{7.97}{\infty}%
\TextAt{20}{280.398}{302.595}{\Character{88}}%
\TextAt{21}{282.18}{277.4}{i}%
\TextAt{17}{285.063}{277.4}{\Character{61}0}%
\TextAt{9}{297.667}{291.237}{\Character{40}}%
\TextAt{8}{302.219}{291.237}{\char35\relax }%
\TextAt{13}{311.974}{291.237}{G}%
\TextAt{21}{321.208}{289.444}{i}%
\TextAt{16}{327.246}{291.237}{\Character{0}}%
\TextAt{9}{339.201}{291.237}{1\Character{41}}%
\TextAt{13}{349.606}{291.237}{\Character{58}}%
\TextAt{12}{97.159}{256.777}{Hint.}%
\TextAt{8}{132.485}{256.777}{If}%
\TextAt{16}{143.54}{256.777}{\Character{79}}%
\TextAt{9}{156.711}{256.778}{\Character{61}}%
\TextAt{14}{169.136}{256.777}{o}%
\TextAt{9}{174.98}{256.778}{\Character{91}}%
\TextAt{13}{178.232}{256.777}{x}%
\TextAt{9}{184.884}{256.778}{\Character{93}}%
\TextAt{8}{191.396}{256.777}{(see}%
\TextAt{8}{214.226}{256.777}{chap.}%
\TextAt{8}{246.205}{256.777}{II,}%
\TextAt{8}{261.48}{256.777}{(10.4)),}%
\TextAt{8}{302.61}{256.777}{then}%
\TextAt{13}{328.633}{256.777}{s}%
\TextAt{9}{337.468}{256.778}{\Character{61}}%
\TextAt{13}{349.893}{256.777}{v}%
\TextAt{21}{355.552}{254.984}{L}%
\TextAt{9}{361.809}{256.778}{\Character{40}}%
\TextAt{13}{366.362}{256.778}{\Character{14}}%
\TextAt{21}{371.546}{254.918}{L}%
\TextAt{3}{377.305}{254.918}{\Character{106}}%
\TextAt{21}{379.658}{254.918}{K}%
\TextAt{9}{387.86}{256.778}{\Character{40}}%
\TextAt{13}{392.413}{256.778}{x}%
\TextAt{9}{399.065}{256.778}{\Character{41}\Character{41}}%
\TextAt{9}{411.463}{256.778}{\Character{61}}%
\TextAt{20}{423.916}{265.744}{\Character{80}}%
\TextAt{21}{436.535}{253.291}{\Character{27}}%
\NativeMathAt{441.661}{253.291}{7.97}{\in}%
\TextAt{21}{447.306}{253.291}{G\Character{59}}%
\TextAt{21}{457.659}{253.291}{\Character{27}}%
\NativeMathAt{462.794}{253.291}{7.97}{\neq}%
\TextAt{17}{469.377}{253.291}{1}%
\TextAt{13}{476.105}{256.778}{v}%
\TextAt{21}{481.764}{254.984}{L}%
\TextAt{9}{488.022}{256.778}{\Character{40}}%
\TextAt{13}{492.574}{256.778}{x}%
\TextAt{16}{500.579}{256.778}{\Character{0}}%
\TextAt{13}{511.231}{256.778}{\Character{27}}%
\TextAt{13}{518.297}{256.778}{x}%
\TextAt{9}{524.965}{256.778}{\Character{41}}%
\TextAt{8}{529.518}{256.777}{.}%
\TextAt{12}{79.226}{224.788}{Exercise}%
\TextAt{12}{132.845}{224.788}{III.2.3.}%
\TextAt{8}{177.563}{224.788}{The}%
\TextAt{8}{201.981}{224.788}{module}%
\TextAt{8}{243.638}{224.788}{of}%
\TextAt{8}{257.332}{224.788}{differentials}%
\TextAt{9}{320.843}{224.788}{\Character{10}}%
\TextAt{17}{329.297}{229.126}{1}%
\TextAt{3}{329.297}{221.236}{\Character{79}}%
\TextAt{3}{336.231}{221.236}{\Character{106}}%
\TextAt{27}{338.593}{221.236}{o}%
\TextAt{8}{347.248}{224.788}{is}%
\TextAt{8}{359.376}{224.788}{generated}%
\TextAt{8}{413.072}{224.788}{by}%
\TextAt{8}{429.695}{224.788}{a}%
\TextAt{8}{439.814}{224.788}{single}%
\TextAt{8}{472.756}{224.788}{element}%
\TextAt{8}{516.362}{224.788}{d}%
\TextAt{13}{522.866}{224.788}{x}%
\TextAt{8}{529.518}{224.788}{,}%
\TextAt{13}{79.226}{207.742}{x}%
\NativeMathAt{89.199}{207.742}{11.955}{\in}%
\TextAt{16}{100.473}{207.742}{\Character{79}}%
\TextAt{8}{110.34}{207.742}{,}%
\TextAt{8}{117.489}{207.742}{and}%
\TextAt{8}{140.252}{207.742}{there}%
\TextAt{8}{170.163}{207.742}{is}%
\TextAt{8}{181.927}{207.742}{an}%
\TextAt{8}{198.186}{207.742}{exact}%
\TextAt{8}{229.079}{207.742}{sequence}%
\TextAt{8}{277.581}{207.742}{of}%
\TextAt{16}{290.911}{207.742}{\Character{79}}%
\TextAt{8}{300.761}{207.742}{-modules}%
\TextAt{9}{218.901}{178.74}{0}%
\NativeArrowAt{228.075}{178.74}{11.955}{19.26}{\longrightarrow}%
\TextAt{14}{250.657}{178.74}{D}%
\TextAt{3}{260.598}{176.881}{\Character{79}}%
\TextAt{3}{267.532}{176.881}{\Character{106}}%
\TextAt{27}{269.894}{176.881}{o}%
\NativeArrowAt{277.609}{178.74}{11.955}{19.26}{\longrightarrow}%
\TextAt{16}{300.18}{178.74}{\Character{79}}%
\NativeArrowAt{313.343}{178.74}{11.955}{19.26}{\longrightarrow}%
\TextAt{9}{335.943}{178.74}{\Character{10}}%
\TextAt{17}{344.398}{183.677}{1}%
\TextAt{3}{344.398}{175.785}{\Character{79}}%
\TextAt{3}{351.332}{175.785}{\Character{106}}%
\TextAt{27}{353.693}{175.785}{o}%
\NativeArrowAt{361.408}{178.74}{11.955}{19.26}{\longrightarrow}%
\TextAt{9}{383.99}{178.74}{0}%
\TextAt{13}{389.843}{178.74}{\Character{58}}%
\TextAt{8}{297.211}{43.033}{189}%
\Destination{79.36}{529.055}{errata.4}%
\Destination{78.23}{809.815}{nav.page.189}%
\Destination{79.226}{498.055}{nav.ex.3.2.1}%
\Destination{79.226}{363.755}{nav.ex.3.2.2}%
\Destination{79.226}{238.755}{nav.ex.3.2.3}%
\Destination{79.36}{364.755}{nav.print.208}%
\Hotlink{401.299}{656.655}{19.22}{11.047}{nav.cite.N035}%
%
%
\Hotlink{458.141}{413.655}{19.22}{11.047}{nav.cite.N060}%
%
\Hotlink{245.705}{254.555}{9.763}{11.047}{nav.chapter.2}%
\Hotlink{348.189}{138.355}{16.502}{9.572}{nav.cite.N039}%
\Hotlink{360.85}{541.955}{30.936}{11.047}{nav.number.3.2.18}%
\Hotlink{260.98}{254.555}{30.936}{11.047}{nav.number.2.10.4}%
\Destination{79.226}{709.94}{errata.176}%
}

% Page 202 (body page 190)
\Page{%
\TextAt{2}{79.226}{796.943}{III.}%
\TextAt{2}{98.317}{796.943}{Riemann–Roch}%
\TextAt{2}{174.375}{796.943}{Theory}%
\Vector{q 0 G 0 g 0.299 w 7.226 18.729 m 460.769 18.729 l S Q}%
\TextAt{12}{79.226}{761.124}{Exercise}%
\TextAt{12}{132.772}{761.124}{III.2.4.}%
\TextAt{8}{177.426}{761.124}{For}%
\TextAt{8}{198.696}{761.124}{a}%
\TextAt{8}{208.752}{761.124}{tower}%
\TextAt{13}{240.924}{761.124}{M}%
\NativeMathAt{257.331}{761.124}{11.955}{\supseteq}%
\TextAt{13}{270.462}{761.124}{L}%
\NativeMathAt{282.259}{761.124}{11.955}{\supseteq}%
\TextAt{13}{295.39}{761.124}{K}%
\TextAt{8}{310.405}{761.124}{of}%
\TextAt{8}{324.036}{761.124}{algebraic}%
\TextAt{8}{373.771}{761.124}{number}%
\TextAt{8}{416.667}{761.124}{fields}%
\TextAt{8}{446.939}{761.124}{there}%
\TextAt{8}{477.151}{761.124}{is}%
\TextAt{8}{489.215}{761.124}{an}%
\TextAt{8}{505.775}{761.124}{exact}%
\TextAt{8}{79.226}{744.078}{sequence}%
\TextAt{8}{127.728}{744.078}{of}%
\TextAt{14}{141.058}{744.078}{o}%
\TextAt{21}{146.903}{742.284}{M}%
\TextAt{8}{156.382}{744.078}{-modules}%
\TextAt{9}{192.343}{715.164}{0}%
\NativeArrowAt{201.517}{715.164}{11.955}{19.26}{\longrightarrow}%
\TextAt{9}{224.099}{715.164}{\Character{10}}%
\TextAt{17}{232.553}{720.1}{1}%
\TextAt{21}{232.553}{712.208}{L}%
\TextAt{3}{238.312}{712.209}{\Character{106}}%
\TextAt{21}{240.664}{712.208}{K}%
\TextAt{16}{251.524}{715.164}{\Character{10}}%
\TextAt{14}{263.479}{715.164}{o}%
\TextAt{21}{269.323}{713.371}{M}%
\NativeArrowAt{282.124}{715.164}{11.955}{19.26}{\longrightarrow}%
\TextAt{9}{304.706}{715.164}{\Character{10}}%
\TextAt{17}{313.16}{720.1}{1}%
\TextAt{21}{313.16}{712.208}{M}%
\TextAt{3}{322.142}{712.209}{\Character{106}}%
\TextAt{21}{324.494}{712.208}{K}%
\NativeArrowAt{336.018}{715.164}{11.955}{19.26}{\longrightarrow}%
\TextAt{9}{358.6}{715.164}{\Character{10}}%
\TextAt{17}{367.054}{720.1}{1}%
\TextAt{21}{367.054}{712.208}{M}%
\TextAt{3}{376.036}{712.209}{\Character{106}}%
\TextAt{21}{378.388}{712.208}{L}%
\NativeArrowAt{387.966}{715.164}{11.955}{19.26}{\longrightarrow}%
\TextAt{9}{410.548}{715.164}{0}%
\TextAt{13}{416.401}{715.164}{\Character{58}}%
\TextAt{12}{79.226}{671.356}{Exercise}%
\TextAt{12}{132.427}{671.356}{III.2.5.}%
\TextAt{8}{176.78}{671.356}{If}%
\TextAt{13}{188.472}{671.356}{\Character{16}}%
\TextAt{8}{198.385}{671.356}{is}%
\TextAt{8}{210.149}{671.356}{a}%
\TextAt{8}{219.905}{671.356}{primitive}%
\TextAt{13}{269.985}{671.356}{p}%
\TextAt{21}{275.86}{675.695}{n}%
\TextAt{8}{281.496}{671.356}{-th}%
\TextAt{8}{300.35}{671.356}{root}%
\TextAt{8}{325.396}{671.356}{of}%
\TextAt{8}{338.726}{671.356}{unity,}%
\TextAt{8}{371.567}{671.356}{then}%
\TextAt{14}{240.29}{642.443}{D}%
\TextAt{25}{250.232}{640.583}{Q}%
\TextAt{22}{256.431}{639.587}{p}%
\TextAt{17}{260.769}{640.583}{\Character{40}}%
\TextAt{21}{264.063}{640.583}{\Character{16}}%
\TextAt{17}{268.331}{640.583}{\Character{41}}%
\TextAt{3}{271.625}{640.583}{\Character{106}}%
\TextAt{25}{273.977}{640.583}{Q}%
\TextAt{22}{280.176}{639.587}{p}%
\TextAt{9}{288.333}{642.443}{\Character{61}}%
\TextAt{14}{300.759}{642.443}{P}%
\TextAt{21}{310.652}{647.379}{p}%
\TextAt{22}{314.915}{650.192}{n}%
\TextAt{23}{319.516}{650.192}{\Character{0}}%
\TextAt{24}{325.272}{650.192}{1}%
\TextAt{17}{329.423}{647.379}{\Character{40}}%
\TextAt{21}{332.716}{647.379}{pn}%
\TextAt{3}{342.117}{647.379}{\Character{0}}%
\TextAt{21}{348.704}{647.379}{n}%
\TextAt{3}{353.842}{647.379}{\Character{0}}%
\TextAt{17}{360.428}{647.379}{1\Character{41}}%
\TextAt{13}{368.454}{642.443}{\Character{58}}%
\TextAt{15}{79.226}{577.749}{§}%
\TextAt{15}{93.785}{577.749}{3}%
\TextAt{15}{120.693}{577.749}{Riemann–Roch}%
\TextAt{8}{79.226}{545.863}{The}%
\TextAt{8}{104.383}{545.863}{notion}%
\TextAt{8}{141.913}{545.863}{of}%
\TextAt{8}{156.345}{545.863}{replete}%
\TextAt{8}{195.812}{545.863}{divisor}%
\TextAt{8}{235.028}{545.863}{introduced}%
\TextAt{8}{294.663}{545.863}{into}%
\TextAt{8}{319.509}{545.863}{our}%
\TextAt{8}{341.425}{545.863}{development}%
\TextAt{8}{410.493}{545.863}{of}%
\TextAt{8}{424.926}{545.863}{number}%
\TextAt{8}{468.624}{545.863}{theory}%
\TextAt{8}{506.477}{545.863}{in}%
\TextAt{8}{521.233}{545.863}{§1}%
\TextAt{8}{79.226}{528.817}{is}%
\TextAt{8}{91.992}{528.817}{a}%
\TextAt{8}{102.75}{528.817}{terminology}%
\TextAt{8}{168.477}{528.817}{reminiscent}%
\TextAt{8}{231.646}{528.817}{of}%
\TextAt{8}{245.978}{528.817}{the}%
\TextAt{8}{267.137}{528.817}{function-theoretic}%
\TextAt{8}{362.764}{528.817}{model.}%
\TextAt{8}{405.111}{528.817}{We}%
\TextAt{8}{426.257}{528.817}{now}%
\TextAt{8}{451.648}{528.817}{have}%
\TextAt{8}{479.645}{528.817}{to}%
\TextAt{8}{494.957}{528.817}{ask}%
\TextAt{8}{516.51}{528.817}{the}%
\TextAt{8}{79.226}{511.771}{question}%
\TextAt{8}{127.291}{511.771}{to}%
\TextAt{8}{143.101}{511.771}{what}%
\TextAt{8}{173.867}{511.771}{extent}%
\TextAt{8}{211.137}{511.771}{this}%
\TextAt{8}{235.459}{511.771}{point}%
\TextAt{8}{267.528}{511.771}{of}%
\TextAt{8}{282.357}{511.771}{view}%
\TextAt{8}{310.84}{511.771}{does}%
\TextAt{8}{338.736}{511.771}{justice}%
\TextAt{8}{377.034}{511.771}{to}%
\TextAt{8}{392.844}{511.771}{our}%
\TextAt{8}{415.157}{511.771}{goal}%
\TextAt{8}{441.38}{511.771}{to}%
\TextAt{8}{457.19}{511.771}{also}%
\TextAt{8}{482.17}{511.771}{couch}%
\TextAt{8}{516.51}{511.771}{the}%
\TextAt{8}{79.226}{494.725}{number-theoretic}%
\TextAt{8}{170.299}{494.725}{results}%
\TextAt{8}{207.192}{494.725}{in}%
\TextAt{8}{220.545}{494.725}{a}%
\TextAt{8}{230}{494.725}{geometric}%
\TextAt{8}{283.032}{494.725}{function-theoretic}%
\TextAt{8}{377.357}{494.725}{fashion,}%
\TextAt{8}{420.43}{494.725}{and}%
\TextAt{8}{442.893}{494.725}{conversely}%
\TextAt{8}{498.591}{494.725}{to}%
\TextAt{8}{512.601}{494.725}{give}%
\TextAt{8}{79.226}{477.679}{arithmetic}%
\TextAt{8}{136.885}{477.679}{significance}%
\TextAt{8}{199.805}{477.679}{to}%
\TextAt{8}{215.19}{477.679}{the}%
\TextAt{8}{236.422}{477.679}{classical}%
\TextAt{8}{282.497}{477.679}{theorems}%
\TextAt{8}{333.712}{477.679}{of}%
\TextAt{8}{348.118}{477.679}{function}%
\TextAt{8}{395.041}{477.679}{theory.}%
\TextAt{8}{438.588}{477.679}{Among}%
\TextAt{8}{480.311}{477.679}{the}%
\TextAt{8}{501.543}{477.679}{latter,}%
\TextAt{8}{79.226}{460.633}{the}%
\TextAt{8}{99.401}{460.633}{Riemann–Roch}%
\TextAt{8}{181.05}{460.633}{theorem}%
\TextAt{8}{226.594}{460.633}{stands}%
\TextAt{8}{263.16}{460.633}{out}%
\TextAt{8}{283.992}{460.633}{as}%
\TextAt{8}{298.381}{460.633}{the}%
\TextAt{8}{318.556}{460.633}{most}%
\TextAt{8}{347.255}{460.633}{important}%
\TextAt{8}{402.567}{460.633}{representative.}%
\TextAt{8}{482.36}{460.633}{If}%
\TextAt{8}{494.071}{460.633}{number}%
\TextAt{8}{79.226}{443.587}{theory}%
\TextAt{8}{116.755}{443.587}{is}%
\TextAt{8}{129.297}{443.587}{to}%
\TextAt{8}{144.386}{443.587}{proceed}%
\TextAt{8}{188.407}{443.587}{in}%
\TextAt{8}{202.838}{443.587}{a}%
\TextAt{8}{213.372}{443.587}{geometric}%
\TextAt{8}{267.483}{443.587}{manner,}%
\TextAt{8}{313.981}{443.587}{it}%
\TextAt{8}{326.464}{443.587}{must}%
\TextAt{8}{356.246}{443.587}{work}%
\TextAt{8}{385.645}{443.587}{towards}%
\TextAt{8}{430.072}{443.587}{finding}%
\TextAt{8}{469.872}{443.587}{an}%
\TextAt{8}{486.909}{443.587}{adequate}%
\TextAt{8}{79.226}{426.542}{way}%
\TextAt{8}{102.969}{426.542}{to}%
\TextAt{8}{117.279}{426.542}{incorporate}%
\TextAt{8}{179.399}{426.542}{this}%
\TextAt{8}{202.221}{426.542}{result}%
\TextAt{8}{234.799}{426.542}{as}%
\TextAt{8}{249.169}{426.542}{well.}%
\TextAt{8}{277.451}{426.542}{This}%
\TextAt{8}{304.171}{426.542}{is}%
\TextAt{8}{315.934}{426.542}{the}%
\TextAt{8}{336.091}{426.542}{task}%
\TextAt{8}{361.197}{426.542}{we}%
\TextAt{8}{378.424}{426.542}{are}%
\TextAt{8}{397.935}{426.542}{now}%
\TextAt{8}{422.324}{426.542}{going}%
\TextAt{8}{453.55}{426.542}{to}%
\TextAt{8}{467.861}{426.542}{tackle.}%
\TextAt{8}{97.159}{406.507}{First}%
\TextAt{8}{124.745}{406.507}{recall}%
\TextAt{8}{155.033}{406.507}{the}%
\TextAt{8}{174.263}{406.507}{classical}%
\TextAt{8}{218.335}{406.507}{situation}%
\TextAt{8}{266.257}{406.507}{in}%
\TextAt{8}{278.983}{406.507}{function}%
\TextAt{8}{323.905}{406.507}{theory.}%
\TextAt{8}{363.916}{406.507}{There}%
\TextAt{8}{396.799}{406.507}{the}%
\TextAt{8}{416.028}{406.507}{basic}%
\TextAt{8}{444.428}{406.507}{data}%
\TextAt{8}{470.173}{406.507}{is}%
\TextAt{8}{481.01}{406.507}{a}%
\TextAt{8}{489.838}{406.507}{compact}%
\TextAt{8}{79.226}{389.461}{Riemann}%
\TextAt{8}{129.787}{389.461}{surface}%
\TextAt{13}{170.162}{389.461}{X}%
\TextAt{8}{185.685}{389.461}{with}%
\TextAt{8}{213.316}{389.461}{the}%
\TextAt{8}{234.443}{389.461}{sheaf}%
\TextAt{14}{265.063}{389.461}{o}%
\TextAt{21}{270.907}{387.668}{X}%
\TextAt{8}{283.86}{389.461}{of}%
\TextAt{8}{298.161}{389.461}{holomorphic}%
\TextAt{8}{366.129}{389.461}{functions.}%
\TextAt{8}{424.055}{389.461}{To}%
\TextAt{8}{442.253}{389.461}{each}%
\TextAt{8}{469.561}{389.461}{divisor}%
\TextAt{13}{508.645}{389.461}{D}%
\TextAt{9}{523.665}{389.461}{\Character{61}}%
\TextAt{20}{79.226}{381.381}{\Character{80}}%
\TextAt{21}{91.846}{368.928}{P}%
\NativeMathAt{98.391}{368.928}{7.97}{\in}%
\TextAt{21}{104.036}{368.928}{X}%
\TextAt{13}{114.114}{372.415}{\Character{23}}%
\TextAt{21}{119.897}{370.622}{P}%
\TextAt{13}{126.941}{372.415}{P}%
\TextAt{8}{140.794}{372.415}{on}%
\TextAt{13}{157.839}{372.415}{X}%
\TextAt{8}{173.177}{372.415}{corresponds}%
\TextAt{8}{238.15}{372.415}{a}%
\TextAt{12}{248.69}{372.415}{line}%
\TextAt{12}{275.166}{372.415}{bundle}%
\TextAt{14}{319.612}{372.415}{o}%
\TextAt{9}{325.456}{372.415}{\Character{40}}%
\TextAt{13}{330.008}{372.415}{D}%
\TextAt{9}{340.053}{372.415}{\Character{41}}%
\TextAt{8}{344.605}{372.415}{,}%
\TextAt{8}{352.736}{372.415}{i.e.,}%
\TextAt{8}{375.823}{372.415}{an}%
\TextAt{14}{392.867}{372.415}{o}%
\TextAt{21}{398.712}{370.622}{X}%
\TextAt{8}{406.797}{372.415}{-module}%
\TextAt{8}{452.773}{372.415}{which}%
\TextAt{8}{487.045}{372.415}{is}%
\TextAt{8}{499.594}{372.415}{locally}%
\TextAt{8}{79.226}{355.369}{free}%
\TextAt{8}{101.057}{355.369}{of}%
\TextAt{8}{113.791}{355.369}{rank}%
\TextAt{9}{140.189}{355.369}{1}%
\TextAt{8}{146.042}{355.369}{.}%
\TextAt{8}{154.292}{355.369}{If}%
\TextAt{13}{165.387}{355.369}{U}%
\TextAt{8}{177.887}{355.369}{is}%
\TextAt{8}{189.054}{355.369}{an}%
\TextAt{8}{204.717}{355.369}{open}%
\TextAt{8}{232.406}{355.369}{subset}%
\TextAt{8}{267.699}{355.369}{of}%
\TextAt{13}{280.433}{355.369}{X}%
\TextAt{8}{294.388}{355.369}{and}%
\TextAt{13}{316.555}{355.369}{K}%
\TextAt{9}{327.372}{355.369}{\Character{40}}%
\TextAt{13}{331.925}{355.369}{U}%
\TextAt{9}{341.124}{355.369}{\Character{41}}%
\TextAt{8}{348.977}{355.369}{is}%
\TextAt{8}{360.144}{355.369}{the}%
\TextAt{8}{379.704}{355.369}{ring}%
\TextAt{8}{403.173}{355.369}{of}%
\TextAt{8}{415.907}{355.369}{meromorphic}%
\TextAt{8}{486.204}{355.369}{functions}%
\TextAt{8}{79.226}{338.323}{on}%
\TextAt{13}{95.485}{338.323}{U}%
\TextAt{8}{104.684}{338.323}{,}%
\TextAt{8}{111.833}{338.323}{then}%
\TextAt{8}{138.493}{338.323}{the}%
\TextAt{8}{158.65}{338.323}{vector}%
\TextAt{8}{193.774}{338.323}{space}%
\TextAt{14}{225.049}{338.323}{o}%
\TextAt{9}{230.893}{338.323}{\Character{40}}%
\TextAt{13}{235.445}{338.323}{D}%
\TextAt{9}{245.49}{338.323}{\Character{41}\Character{40}}%
\TextAt{13}{254.595}{338.323}{U}%
\TextAt{9}{263.794}{338.323}{\Character{41}}%
\TextAt{8}{272.243}{338.323}{of}%
\TextAt{8}{285.573}{338.323}{sections}%
\TextAt{8}{329.27}{338.323}{of}%
\TextAt{8}{342.6}{338.323}{the}%
\TextAt{8}{362.756}{338.323}{sheaf}%
\TextAt{14}{392.405}{338.323}{o}%
\TextAt{9}{398.249}{338.323}{\Character{40}}%
\TextAt{13}{402.801}{338.323}{D}%
\TextAt{9}{412.846}{338.323}{\Character{41}}%
\TextAt{8}{421.296}{338.323}{over}%
\TextAt{13}{446.342}{338.323}{U}%
\TextAt{8}{459.438}{338.323}{is}%
\TextAt{8}{471.202}{338.323}{given}%
\TextAt{8}{501.771}{338.323}{as}%
\TextAt{14}{163.902}{309.409}{o}%
\TextAt{9}{169.746}{309.409}{\Character{40}}%
\TextAt{13}{174.298}{309.409}{D}%
\TextAt{9}{184.343}{309.409}{\Character{41}\Character{40}}%
\TextAt{13}{193.447}{309.409}{U}%
\TextAt{9}{202.647}{309.409}{\Character{41}}%
\TextAt{9}{210.502}{309.409}{\Character{61}}%
\TextAt{16}{222.945}{309.409}{\Character{102}}%
\TextAt{13}{228.923}{309.409}{f}%
\NativeMathAt{239.29}{309.409}{11.955}{\in}%
\TextAt{13}{250.581}{309.409}{K}%
\TextAt{9}{261.399}{309.409}{\Character{40}}%
\TextAt{13}{265.951}{309.409}{U}%
\TextAt{9}{275.15}{309.409}{\Character{41}}%
\TextAt{16}{283.023}{309.409}{\Character{106}}%
\TextAt{8}{289.665}{309.409}{ord}%
\TextAt{21}{306.582}{307.616}{P}%
\TextAt{9}{313.625}{309.409}{\Character{40}}%
\TextAt{13}{318.177}{309.409}{f}%
\TextAt{9}{325.223}{309.409}{\Character{41}}%
\NativeMathAt{333.097}{309.409}{11.955}{\geq}%
\TextAt{16}{345.698}{309.409}{\Character{0}}%
\TextAt{13}{355.014}{309.409}{\Character{23}}%
\TextAt{21}{360.798}{307.616}{P}%
\TextAt{8}{371.739}{309.409}{for}%
\TextAt{8}{389.624}{309.409}{all}%
\TextAt{13}{405.883}{309.409}{P}%
\NativeMathAt{418.375}{309.409}{11.955}{\in}%
\TextAt{13}{429.666}{309.409}{U}%
\TextAt{16}{438.865}{309.409}{\Character{103}}%
\TextAt{13}{444.842}{309.409}{\Character{58}}%
\TextAt{8}{79.226}{280.495}{The}%
\TextAt{8}{103.28}{280.495}{Riemann–Roch}%
\TextAt{8}{184.91}{280.495}{problem}%
\TextAt{8}{230.435}{280.495}{is}%
\TextAt{8}{242.199}{280.495}{to}%
\TextAt{8}{256.509}{280.495}{calculate}%
\TextAt{8}{305.286}{280.495}{the}%
\TextAt{8}{325.443}{280.495}{dimension}%
\TextAt{13}{242.991}{251.582}{\Character{96}}%
\TextAt{9}{247.904}{251.582}{\Character{40}}%
\TextAt{13}{252.456}{251.582}{D}%
\TextAt{9}{262.5}{251.582}{\Character{41}}%
\TextAt{9}{270.355}{251.582}{\Character{61}}%
\TextAt{8}{282.799}{251.582}{dim}%
\TextAt{13}{304.302}{251.582}{H}%
\TextAt{17}{314.925}{256.518}{0}%
\TextAt{9}{319.657}{251.582}{\Character{40}}%
\TextAt{13}{324.21}{251.582}{X}%
\TextAt{13}{334.204}{251.582}{\Character{59}}%
\TextAt{14}{339.459}{251.582}{o}%
\TextAt{9}{345.303}{251.582}{\Character{40}}%
\TextAt{13}{349.855}{251.582}{D}%
\TextAt{9}{359.9}{251.582}{\Character{41}\Character{41}}%
\TextAt{8}{79.226}{222.668}{of}%
\TextAt{8}{92.556}{222.668}{the}%
\TextAt{8}{112.713}{222.668}{vector}%
\TextAt{8}{147.837}{222.668}{space}%
\TextAt{8}{179.112}{222.668}{of}%
\TextAt{8}{192.442}{222.668}{global}%
\TextAt{8}{226.92}{222.668}{sections}%
\TextAt{13}{241.771}{193.754}{H}%
\TextAt{17}{252.394}{198.691}{0}%
\TextAt{9}{257.126}{193.754}{\Character{40}}%
\TextAt{13}{261.679}{193.754}{X}%
\TextAt{13}{271.673}{193.754}{\Character{59}}%
\TextAt{14}{276.928}{193.754}{o}%
\TextAt{9}{282.772}{193.754}{\Character{40}}%
\TextAt{13}{287.324}{193.754}{D}%
\TextAt{9}{297.369}{193.754}{\Character{41}\Character{41}}%
\TextAt{9}{309.767}{193.754}{\Character{61}}%
\TextAt{14}{322.22}{193.754}{o}%
\TextAt{9}{328.064}{193.754}{\Character{40}}%
\TextAt{13}{332.616}{193.754}{D}%
\TextAt{9}{342.661}{193.754}{\Character{41}\Character{40}}%
\TextAt{13}{351.765}{193.754}{X}%
\TextAt{9}{362.42}{193.754}{\Character{41}}%
\TextAt{13}{366.973}{193.754}{\Character{58}}%
\TextAt{8}{79.226}{164.841}{In}%
\TextAt{8}{92.872}{164.841}{its}%
\TextAt{8}{108.215}{164.841}{first}%
\TextAt{8}{131.364}{164.841}{version}%
\TextAt{8}{170.128}{164.841}{the}%
\TextAt{8}{189.308}{164.841}{Riemann–Roch}%
\TextAt{8}{269.962}{164.841}{theorem}%
\TextAt{8}{314.512}{164.841}{does}%
\TextAt{8}{339.933}{164.841}{not}%
\TextAt{8}{359.771}{164.841}{provide}%
\TextAt{8}{400.424}{164.841}{a}%
\TextAt{8}{409.203}{164.841}{formula}%
\TextAt{8}{451.158}{164.841}{for}%
\TextAt{13}{468.068}{164.841}{H}%
\TextAt{17}{478.69}{169.179}{0}%
\TextAt{9}{483.423}{164.841}{\Character{40}}%
\TextAt{13}{487.975}{164.841}{X}%
\TextAt{13}{497.97}{164.841}{\Character{59}}%
\TextAt{14}{503.224}{164.841}{o}%
\TextAt{9}{509.068}{164.841}{\Character{40}}%
\TextAt{13}{513.62}{164.841}{D}%
\TextAt{9}{523.665}{164.841}{\Character{41}\Character{41}}%
\TextAt{8}{79.226}{147.795}{itself,}%
\TextAt{8}{110.824}{147.795}{but}%
\TextAt{8}{132.283}{147.795}{for}%
\TextAt{8}{150.168}{147.795}{the}%
\TextAt{12}{170.325}{147.795}{Euler–Poincaré}%
\TextAt{12}{264.125}{147.795}{characteristic}%
\TextAt{13}{182.272}{118.881}{\Character{31}}%
\TextAt{9}{189.605}{118.881}{\Character{40}}%
\TextAt{14}{194.158}{118.881}{o}%
\TextAt{9}{200.002}{118.881}{\Character{40}}%
\TextAt{13}{204.554}{118.881}{D}%
\TextAt{9}{214.599}{118.881}{\Character{41}\Character{41}}%
\TextAt{9}{226.997}{118.881}{\Character{61}}%
\TextAt{8}{239.45}{118.881}{dim}%
\TextAt{13}{260.953}{118.881}{H}%
\TextAt{17}{271.576}{123.817}{0}%
\TextAt{9}{276.308}{118.881}{\Character{40}}%
\TextAt{13}{280.86}{118.881}{X}%
\TextAt{13}{290.855}{118.881}{\Character{59}}%
\TextAt{14}{296.109}{118.881}{o}%
\TextAt{9}{301.954}{118.881}{\Character{40}}%
\TextAt{13}{306.506}{118.881}{D}%
\TextAt{9}{316.55}{118.881}{\Character{41}\Character{41}}%
\TextAt{16}{328.312}{118.881}{\Character{0}}%
\TextAt{8}{340.267}{118.881}{dim}%
\TextAt{13}{361.77}{118.881}{H}%
\TextAt{17}{372.393}{123.817}{1}%
\TextAt{9}{377.125}{118.881}{\Character{40}}%
\TextAt{13}{381.677}{118.881}{X}%
\TextAt{13}{391.672}{118.881}{\Character{59}}%
\TextAt{14}{396.926}{118.881}{o}%
\TextAt{9}{402.771}{118.881}{\Character{40}}%
\TextAt{13}{407.323}{118.881}{D}%
\TextAt{9}{417.367}{118.881}{\Character{41}\Character{41}}%
\TextAt{13}{426.472}{118.881}{\Character{58}}%
\TextAt{8}{79.226}{89.967}{The}%
\TextAt{8}{103.28}{89.967}{formula}%
\TextAt{8}{146.211}{89.967}{reads}%
\TextAt{13}{236.969}{72.921}{\Character{31}}%
\TextAt{9}{244.302}{72.921}{\Character{40}}%
\TextAt{14}{248.855}{72.921}{o}%
\TextAt{9}{254.699}{72.921}{\Character{40}}%
\TextAt{13}{259.251}{72.921}{D}%
\TextAt{9}{269.296}{72.921}{\Character{41}\Character{41}}%
\TextAt{9}{281.694}{72.921}{\Character{61}}%
\TextAt{8}{294.147}{72.921}{deg}%
\TextAt{9}{311.709}{72.921}{\Character{40}}%
\TextAt{13}{316.261}{72.921}{D}%
\TextAt{9}{326.306}{72.921}{\Character{41}}%
\TextAt{9}{333.503}{72.921}{\Character{43}}%
\TextAt{9}{345.255}{72.921}{1}%
\TextAt{16}{353.786}{72.921}{\Character{0}}%
\TextAt{13}{365.741}{72.921}{g}%
\TextAt{13}{371.755}{72.921}{\Character{59}}%
\TextAt{8}{297.211}{43.033}{190}%
\Destination{78.23}{809.815}{nav.page.190}%
\Destination{79.23}{602.585}{nav.section.3.3}%
\Destination{79.226}{775.155}{nav.ex.3.2.4}%
\Destination{79.226}{685.355}{nav.ex.3.2.5}%
\Destination{79.23}{602.585}{outline.34}%
\Hotlink{526.412}{543.655}{6.858}{11.047}{nav.section.3.1}%
}

% Page 203 (body page 191)
\Page{%
\TextAt{2}{79.226}{796.943}{III.}%
\TextAt{2}{98.317}{796.943}{Riemann–Roch}%
\TextAt{2}{174.375}{796.943}{Theory}%
\Vector{q 0 G 0 g 0.299 w 7.226 18.729 m 460.769 18.729 l S Q}%
\TextAt{8}{79.226}{761.124}{where}%
\TextAt{13}{113.237}{761.124}{g}%
\TextAt{8}{123.371}{761.124}{is}%
\TextAt{8}{135.337}{761.124}{the}%
\TextAt{12}{155.695}{761.124}{genus}%
\TextAt{8}{192.54}{761.124}{of}%
\TextAt{13}{206.072}{761.124}{X}%
\TextAt{8}{216.727}{761.124}{.}%
\TextAt{8}{225.781}{761.124}{For}%
\TextAt{8}{246.953}{761.124}{the}%
\TextAt{8}{267.311}{761.124}{divisor}%
\TextAt{13}{305.626}{761.124}{D}%
\TextAt{9}{319.336}{761.124}{\Character{61}}%
\TextAt{9}{332.092}{761.124}{0}%
\TextAt{8}{337.958}{761.124}{,}%
\TextAt{8}{345.36}{761.124}{one}%
\TextAt{8}{367.022}{761.124}{has}%
\TextAt{14}{388.097}{761.124}{o}%
\TextAt{9}{393.942}{761.124}{\Character{40}}%
\TextAt{13}{398.494}{761.124}{D}%
\TextAt{9}{408.538}{761.124}{\Character{41}}%
\TextAt{9}{416.739}{761.124}{\Character{61}}%
\TextAt{14}{429.525}{761.124}{o}%
\TextAt{21}{435.37}{759.33}{X}%
\TextAt{8}{447.554}{761.124}{and}%
\TextAt{8}{470.519}{761.124}{deg}%
\TextAt{9}{488.081}{761.124}{\Character{40}}%
\TextAt{13}{492.633}{761.124}{D}%
\TextAt{9}{502.678}{761.124}{\Character{41}}%
\TextAt{9}{510.879}{761.124}{\Character{61}}%
\TextAt{9}{523.647}{761.124}{0}%
\TextAt{8}{529.518}{761.124}{,}%
\TextAt{8}{79.226}{744.078}{so}%
\TextAt{8}{93.596}{744.078}{that}%
\TextAt{13}{118.965}{744.078}{\Character{31}}%
\TextAt{9}{126.299}{744.078}{\Character{40}}%
\TextAt{14}{130.851}{744.078}{o}%
\TextAt{21}{136.695}{742.284}{X}%
\TextAt{9}{144.78}{744.078}{\Character{41}}%
\TextAt{9}{152.635}{744.078}{\Character{61}}%
\TextAt{9}{165.056}{744.078}{1}%
\TextAt{16}{173.589}{744.078}{\Character{0}}%
\TextAt{13}{185.544}{744.078}{g}%
\TextAt{8}{191.578}{744.078}{;}%
\TextAt{8}{198.727}{744.078}{then}%
\TextAt{8}{225.387}{744.078}{this}%
\TextAt{8}{248.21}{744.078}{equation}%
\TextAt{8}{296.018}{744.078}{may}%
\TextAt{8}{321.387}{744.078}{also}%
\TextAt{8}{344.867}{744.078}{be}%
\TextAt{8}{360.792}{744.078}{replaced}%
\TextAt{8}{406.962}{744.078}{by}%
\TextAt{13}{235.035}{715.1}{\Character{31}}%
\TextAt{9}{242.368}{715.1}{\Character{40}}%
\TextAt{14}{246.921}{715.1}{o}%
\TextAt{9}{252.765}{715.1}{\Character{40}}%
\TextAt{13}{257.317}{715.1}{D}%
\TextAt{9}{267.362}{715.1}{\Character{41}\Character{41}}%
\TextAt{9}{279.76}{715.1}{\Character{61}}%
\TextAt{8}{292.213}{715.1}{deg}%
\TextAt{9}{309.775}{715.1}{\Character{40}}%
\TextAt{13}{314.327}{715.1}{D}%
\TextAt{9}{324.371}{715.1}{\Character{41}}%
\TextAt{9}{331.568}{715.1}{\Character{43}}%
\TextAt{13}{343.342}{715.1}{\Character{31}}%
\TextAt{9}{350.675}{715.1}{\Character{40}}%
\TextAt{14}{355.228}{715.1}{o}%
\TextAt{21}{361.072}{713.307}{X}%
\TextAt{9}{369.157}{715.1}{\Character{41}}%
\TextAt{13}{373.709}{715.1}{\Character{58}}%
\TextAt{8}{79.226}{686.122}{The}%
\TextAt{8}{103.28}{686.122}{classical}%
\TextAt{8}{148.279}{686.122}{Riemann–Roch}%
\TextAt{8}{229.909}{686.122}{formula}%
\TextAt{13}{215.996}{657.145}{\Character{96}}%
\TextAt{9}{220.909}{657.145}{\Character{40}}%
\TextAt{13}{225.461}{657.145}{D}%
\TextAt{9}{235.505}{657.145}{\Character{41}}%
\TextAt{16}{242.714}{657.145}{\Character{0}}%
\TextAt{13}{254.67}{657.145}{\Character{96}}%
\TextAt{9}{259.582}{657.145}{\Character{40}}%
\TextAt{16}{264.134}{657.145}{\Character{75}}%
\TextAt{16}{276.053}{657.145}{\Character{0}}%
\TextAt{13}{288.028}{657.145}{D}%
\TextAt{9}{298.072}{657.145}{\Character{41}}%
\TextAt{9}{305.927}{657.145}{\Character{61}}%
\TextAt{8}{318.371}{657.145}{deg}%
\TextAt{9}{335.933}{657.145}{\Character{40}}%
\TextAt{13}{340.485}{657.145}{D}%
\TextAt{9}{350.53}{657.145}{\Character{41}}%
\TextAt{9}{357.727}{657.145}{\Character{43}}%
\TextAt{9}{369.479}{657.145}{1}%
\TextAt{16}{378.01}{657.145}{\Character{0}}%
\TextAt{13}{389.965}{657.145}{g}%
\TextAt{8}{79.226}{628.167}{is}%
\TextAt{8}{93.069}{628.167}{then}%
\TextAt{8}{121.807}{628.167}{obtained}%
\TextAt{8}{172.018}{628.167}{by}%
\TextAt{8}{190.355}{628.167}{using}%
\TextAt{12}{223.063}{628.167}{Serre}%
\TextAt{12}{260.674}{628.167}{duality}%
\TextAt{8}{301.609}{628.167}{,}%
\TextAt{8}{311.356}{628.167}{which}%
\TextAt{8}{346.921}{628.167}{states}%
\TextAt{8}{382.295}{628.167}{that}%
\TextAt{13}{409.743}{628.167}{H}%
\TextAt{17}{420.365}{632.506}{1}%
\TextAt{9}{425.098}{628.167}{\Character{40}}%
\TextAt{13}{429.65}{628.167}{X}%
\TextAt{13}{439.644}{628.167}{\Character{59}}%
\TextAt{14}{444.899}{628.167}{o}%
\TextAt{9}{450.743}{628.167}{\Character{40}}%
\TextAt{13}{455.296}{628.167}{D}%
\TextAt{9}{465.34}{628.167}{\Character{41}\Character{41}}%
\TextAt{8}{480.421}{628.167}{is}%
\TextAt{8}{494.263}{628.167}{dual}%
\TextAt{8}{522.357}{628.167}{to}%
\TextAt{13}{79.226}{611.121}{H}%
\TextAt{17}{89.849}{615.46}{0}%
\TextAt{9}{94.581}{611.121}{\Character{40}}%
\TextAt{13}{99.133}{611.121}{X}%
\TextAt{13}{109.128}{611.121}{\Character{59}}%
\TextAt{13}{114.352}{611.121}{\Character{33}}%
\TextAt{16}{126.401}{611.121}{\Character{10}}%
\TextAt{14}{140}{611.121}{o}%
\TextAt{9}{145.844}{611.121}{\Character{40}}%
\TextAt{16}{150.397}{611.121}{\Character{0}}%
\TextAt{13}{159.695}{611.121}{D}%
\TextAt{9}{169.74}{611.121}{\Character{41}\Character{41}}%
\TextAt{8}{178.844}{611.121}{,}%
\TextAt{8}{189.008}{611.121}{where}%
\TextAt{13}{225.229}{611.121}{\Character{33}}%
\TextAt{9}{240.378}{611.121}{\Character{61}}%
\TextAt{9}{256.9}{611.121}{\Character{10}}%
\TextAt{17}{265.367}{615.46}{1}%
\TextAt{21}{265.367}{608.101}{X}%
\TextAt{8}{279.761}{611.121}{is}%
\TextAt{8}{293.937}{611.121}{the}%
\TextAt{8}{316.505}{611.121}{so-called}%
\TextAt{12}{366.45}{611.121}{canonical}%
\TextAt{12}{427.901}{611.121}{module}%
\TextAt{8}{477.344}{611.121}{of}%
\TextAt{13}{493.085}{611.121}{X}%
\TextAt{8}{503.741}{611.121}{,}%
\TextAt{8}{513.904}{611.121}{and}%
\TextAt{16}{79.226}{594.075}{\Character{75}}%
\TextAt{9}{91.829}{594.075}{\Character{61}}%
\TextAt{8}{104.254}{594.075}{div}%
\TextAt{9}{120.191}{594.075}{\Character{40}}%
\TextAt{13}{124.743}{594.075}{\Character{33}}%
\TextAt{9}{132.461}{594.075}{\Character{41}}%
\TextAt{8}{140.598}{594.075}{is}%
\TextAt{8}{152.05}{594.075}{the}%
\TextAt{8}{171.893}{594.075}{associated}%
\TextAt{8}{227.316}{594.075}{divisor}%
\TextAt{8}{265.116}{594.075}{(see}%
\TextAt{8}{288.272}{594.075}{for}%
\TextAt{8}{305.844}{594.075}{instance}%
\TextAt{8}{351.116}{594.075}{[51],}%
\TextAt{8}{376.235}{594.075}{chap.}%
\TextAt{8}{408.322}{594.075}{III,}%
\TextAt{8}{428.465}{594.075}{7.12.1}%
\TextAt{8}{461.985}{594.075}{and}%
\TextAt{8}{484.435}{594.075}{chap.}%
\TextAt{8}{516.522}{594.075}{IV,}%
\TextAt{8}{79.226}{577.029}{1.1.3).}%
\TextAt{8}{97.159}{556.995}{In}%
\TextAt{8}{110.848}{556.995}{order}%
\TextAt{8}{140.485}{556.995}{to}%
\TextAt{8}{153.863}{556.995}{mimic}%
\TextAt{8}{188.043}{556.995}{this}%
\TextAt{8}{209.934}{556.995}{state}%
\TextAt{8}{237.682}{556.995}{of}%
\TextAt{8}{250.08}{556.995}{affairs}%
\TextAt{8}{284.009}{556.995}{in}%
\TextAt{8}{296.729}{556.995}{number}%
\TextAt{8}{338.393}{556.995}{theory,}%
\TextAt{8}{376.669}{556.995}{let}%
\TextAt{8}{392.642}{556.995}{us}%
\TextAt{8}{406.725}{556.995}{recall}%
\TextAt{8}{437.008}{556.995}{the}%
\TextAt{8}{456.232}{556.995}{explanations}%
\TextAt{8}{523.337}{556.995}{of}%
\TextAt{8}{79.226}{539.949}{chap.}%
\TextAt{8}{111.168}{539.949}{I,}%
\TextAt{8}{121.79}{539.949}{§14}%
\TextAt{8}{142.335}{539.949}{and}%
\TextAt{8}{164.35}{539.949}{chap.}%
\TextAt{8}{196.292}{539.949}{III,}%
\TextAt{8}{216}{539.949}{§1.}%
\TextAt{8}{235.736}{539.949}{We}%
\TextAt{8}{255.133}{539.949}{endow}%
\TextAt{8}{290.478}{539.949}{the}%
\TextAt{8}{309.887}{539.949}{places}%
\TextAt{14}{343.666}{539.949}{p}%
\TextAt{8}{352.799}{539.949}{of}%
\TextAt{8}{365.382}{539.949}{an}%
\TextAt{8}{380.894}{539.949}{algebraic}%
\TextAt{8}{429.581}{539.949}{number}%
\TextAt{8}{471.43}{539.949}{field}%
\TextAt{13}{496.039}{539.949}{K}%
\TextAt{8}{510.007}{539.949}{with}%
\TextAt{8}{79.226}{522.903}{the}%
\TextAt{8}{99.01}{522.903}{role}%
\TextAt{8}{121.401}{522.903}{of}%
\TextAt{8}{134.359}{522.903}{points}%
\TextAt{8}{169.171}{522.903}{of}%
\TextAt{8}{182.129}{522.903}{a}%
\TextAt{8}{191.512}{522.903}{space}%
\TextAt{13}{222.415}{522.903}{X}%
\TextAt{8}{236.595}{522.903}{which}%
\TextAt{8}{269.709}{522.903}{should}%
\TextAt{8}{306.47}{522.903}{be}%
\TextAt{8}{322.022}{522.903}{conceived}%
\TextAt{8}{374.325}{522.903}{of}%
\TextAt{8}{387.283}{522.903}{as}%
\TextAt{8}{401.281}{522.903}{the}%
\TextAt{8}{421.065}{522.903}{analogue}%
\TextAt{8}{469.482}{522.903}{of}%
\TextAt{8}{482.44}{522.903}{a}%
\TextAt{19}{491.823}{522.903}{compact}%
\TextAt{8}{79.226}{505.857}{Riemann}%
\TextAt{8}{128.254}{505.857}{surface.}%
\TextAt{8}{172.022}{505.857}{The}%
\TextAt{8}{195.513}{505.857}{elements}%
\TextAt{13}{242.808}{505.857}{f}%
\NativeMathAt{253.175}{505.857}{11.955}{\in}%
\TextAt{13}{264.466}{505.857}{K}%
\TextAt{3}{275.284}{510.195}{\Character{3}}%
\TextAt{8}{283.351}{505.857}{will}%
\TextAt{8}{304.894}{505.857}{be}%
\TextAt{8}{320.256}{505.857}{given}%
\TextAt{8}{350.263}{505.857}{the}%
\TextAt{8}{369.857}{505.857}{role}%
\TextAt{8}{392.058}{505.857}{of}%
\TextAt{8}{404.826}{505.857}{“meromorphic}%
\TextAt{8}{480.681}{505.857}{functions”}%
\ErrataAt{79.19}{488.811}{455.81}{on this space $X$. At a finite prime $\mathfrak p$ the order of vanishing of $f$ is the integer $v_{\mathfrak p}(f)$.}%
\ErrataAt{79.19}{471.765}{455.81}{It is negative at a pole. For $\mathfrak p\mid\infty$ put $v_{\mathfrak p}(f)=-\log|\tau f|$. In}%
\TextAt{8}{79.226}{454.719}{this}%
\TextAt{8}{102.049}{454.719}{way}%
\TextAt{8}{125.792}{454.719}{we}%
\TextAt{8}{143.019}{454.719}{associate}%
\TextAt{8}{192.25}{454.719}{to}%
\TextAt{8}{206.561}{454.719}{each}%
\TextAt{13}{232.898}{454.719}{f}%
\NativeMathAt{243.265}{454.719}{11.955}{\in}%
\TextAt{13}{254.556}{454.719}{K}%
\TextAt{3}{265.374}{459.057}{\Character{3}}%
\TextAt{8}{274.004}{454.719}{the}%
\TextAt{8}{294.16}{454.719}{replete}%
\TextAt{8}{332.524}{454.719}{divisor}%
\TextAt{8}{230.717}{425.741}{div}%
\TextAt{9}{246.653}{425.741}{\Character{40}}%
\TextAt{13}{251.205}{425.741}{f}%
\TextAt{9}{258.252}{425.741}{\Character{41}}%
\TextAt{9}{266.107}{425.741}{\Character{61}}%
\TextAt{20}{278.55}{437.099}{\Character{88}}%
\TextAt{27}{285.19}{411.993}{p}%
\TextAt{13}{297.811}{425.741}{v}%
\TextAt{27}{303.471}{423.948}{p}%
\TextAt{9}{307.957}{425.741}{\Character{40}}%
\TextAt{13}{312.51}{425.741}{f}%
\TextAt{9}{319.556}{425.741}{\Character{41}}%
\TextAt{14}{324.109}{425.741}{p}%
\NativeMathAt{333.412}{425.741}{11.955}{\in}%
\TextAt{8}{344.703}{425.741}{Div}%
\TextAt{9}{363.079}{425.741}{\Character{40}}%
\Vector{q 0 G 0 g 0.478 w 295.631 -340.902 m 301.475 -340.902 l S Q}%
\TextAt{14}{367.631}{425.741}{o}%
\TextAt{9}{373.475}{425.741}{\Character{41}}%
\TextAt{13}{378.027}{425.741}{\Character{58}}%
\TextAt{8}{79.226}{387.396}{More}%
\TextAt{8}{109.461}{387.396}{precisely,}%
\TextAt{8}{159.589}{387.396}{for}%
\TextAt{8}{177.474}{387.396}{a}%
\TextAt{8}{187.229}{387.396}{given}%
\TextAt{8}{217.799}{387.396}{divisor}%
\TextAt{13}{255.912}{387.396}{D}%
\TextAt{9}{269.277}{387.396}{\Character{61}}%
\TextAt{20}{281.702}{396.362}{\Character{80}}%
\TextAt{27}{294.322}{383.909}{p}%
\TextAt{13}{300.801}{387.396}{\Character{23}}%
\TextAt{27}{306.585}{385.603}{p}%
\TextAt{14}{311.072}{387.396}{p}%
\TextAt{8}{317.055}{387.396}{,}%
\TextAt{8}{324.204}{387.396}{we}%
\TextAt{8}{341.431}{387.396}{are}%
\TextAt{8}{360.942}{387.396}{interested}%
\TextAt{8}{414.657}{387.396}{in}%
\TextAt{8}{428.31}{387.396}{the}%
\TextAt{8}{448.466}{387.396}{replete}%
\TextAt{8}{486.83}{387.396}{ideal}%
\TextAt{14}{266.337}{356.921}{o}%
\TextAt{9}{272.181}{356.921}{\Character{40}}%
\TextAt{13}{276.734}{356.921}{D}%
\TextAt{9}{286.778}{356.921}{\Character{41}}%
\TextAt{9}{294.633}{356.921}{\Character{61}}%
\TextAt{20}{307.077}{368.279}{\Character{89}}%
\TextAt{27}{312.72}{343.173}{p}%
\TextAt{14}{324.345}{356.921}{p}%
\TextAt{3}{330.328}{361.857}{\Character{0}}%
\TextAt{21}{336.915}{361.857}{\Character{23}}%
\TextAt{28}{341.078}{360.861}{p}%
\TextAt{8}{79.226}{319.246}{and}%
\TextAt{8}{101.989}{319.246}{the}%
\TextAt{8}{122.145}{319.246}{set}%
\TextAt{13}{165.177}{290.268}{H}%
\TextAt{17}{175.799}{295.204}{0}%
\TextAt{9}{180.532}{290.268}{\Character{40}}%
\TextAt{14}{185.084}{290.268}{o}%
\TextAt{9}{190.928}{290.268}{\Character{40}}%
\TextAt{13}{195.481}{290.268}{D}%
\TextAt{9}{205.525}{290.268}{\Character{41}\Character{41}}%
\TextAt{9}{217.923}{290.268}{\Character{61}}%
\TextAt{16}{230.376}{290.268}{\Character{102}}%
\TextAt{13}{236.354}{290.268}{f}%
\NativeMathAt{246.721}{290.268}{11.955}{\in}%
\TextAt{13}{258.012}{290.268}{K}%
\TextAt{3}{268.829}{295.204}{\Character{3}}%
\TextAt{16}{276.883}{290.268}{\Character{106}}%
\TextAt{8}{283.524}{290.268}{div}%
\TextAt{9}{299.461}{290.268}{\Character{40}}%
\TextAt{13}{304.013}{290.268}{f}%
\TextAt{9}{311.059}{290.268}{\Character{41}}%
\NativeMathAt{318.932}{290.268}{11.955}{\geq}%
\TextAt{16}{331.533}{290.268}{\Character{0}}%
\TextAt{13}{340.85}{290.268}{D}%
\TextAt{16}{350.895}{290.268}{\Character{103}}%
\TextAt{9}{217.95}{270.233}{\Character{61}}%
\TextAt{16}{230.376}{270.233}{\Character{102}}%
\TextAt{13}{236.354}{270.233}{f}%
\NativeMathAt{246.721}{270.233}{11.955}{\in}%
\TextAt{14}{258.012}{270.233}{o}%
\TextAt{9}{263.856}{270.233}{\Character{40}}%
\TextAt{13}{268.408}{270.233}{D}%
\TextAt{9}{278.453}{270.233}{\Character{41}}%
\TextAt{21}{283.005}{268.44}{f}%
\TextAt{16}{291.771}{270.233}{\Character{106}}%
\TextAt{9}{298.413}{270.233}{0}%
\NativeMathAt{307.586}{270.233}{11.955}{\neq}%
\TextAt{16}{320.012}{270.233}{\Character{106}}%
\TextAt{13}{323.333}{270.233}{f}%
\TextAt{16}{330.379}{270.233}{\Character{106}}%
\TextAt{27}{333.7}{268.44}{p}%
\NativeMathAt{341.508}{270.233}{11.955}{\leq}%
\TextAt{14}{354.127}{270.233}{N}%
\TextAt{9}{364.075}{270.233}{\Character{40}}%
\TextAt{14}{368.628}{270.233}{p}%
\TextAt{9}{374.611}{270.233}{\Character{41}}%
\TextAt{21}{379.163}{275.169}{\Character{23}}%
\TextAt{28}{383.326}{274.173}{p}%
\TextAt{8}{391.804}{270.233}{for}%
\TextAt{14}{409.689}{270.233}{p}%
\TextAt{16}{418.993}{270.233}{\Character{106}}%
\NativeMathAt{425.616}{270.233}{11.955}{\infty}%
\TextAt{16}{437.571}{270.233}{\Character{103}}%
\TextAt{13}{443.567}{270.233}{\Character{59}}%
\TextAt{8}{79.226}{241.256}{where}%
\TextAt{8}{114.274}{241.256}{the}%
\TextAt{8}{135.669}{241.256}{relation}%
\TextAt{13}{179.838}{241.256}{D}%
\TextAt{3}{189.882}{245.594}{\Character{48}}%
\NativeMathAt{198.109}{241.256}{11.955}{\geq}%
\TextAt{13}{212.838}{241.256}{D}%
\TextAt{8}{228.019}{241.256}{between}%
\TextAt{8}{274.447}{241.256}{divisors}%
\TextAt{13}{318.414}{241.256}{D}%
\TextAt{3}{328.458}{245.594}{\Character{48}}%
\TextAt{9}{336.684}{241.256}{\Character{61}}%
\TextAt{20}{351.22}{250.222}{\Character{80}}%
\TextAt{27}{363.84}{237.769}{p}%
\TextAt{13}{370.319}{241.256}{\Character{23}}%
\TextAt{3}{376.857}{245.594}{\Character{48}}%
\TextAt{27}{376.103}{238.3}{p}%
\TextAt{14}{380.59}{241.256}{p}%
\TextAt{8}{391.708}{241.256}{and}%
\TextAt{13}{415.709}{241.256}{D}%
\TextAt{9}{431.185}{241.256}{\Character{61}}%
\TextAt{20}{445.721}{250.222}{\Character{80}}%
\TextAt{27}{458.34}{237.769}{p}%
\TextAt{13}{464.82}{241.256}{\Character{23}}%
\TextAt{27}{470.603}{239.462}{p}%
\TextAt{14}{475.09}{241.256}{p}%
\TextAt{8}{486.209}{241.256}{is}%
\TextAt{8}{499.211}{241.256}{simply}%
\TextAt{8}{79.226}{224.21}{defined}%
\TextAt{8}{120.078}{224.21}{to}%
\TextAt{8}{134.927}{224.21}{mean}%
\TextAt{13}{166.681}{224.21}{\Character{23}}%
\TextAt{3}{173.219}{228.548}{\Character{48}}%
\TextAt{27}{172.465}{221.254}{p}%
\NativeMathAt{181.191}{224.21}{11.955}{\geq}%
\TextAt{13}{194.728}{224.21}{\Character{23}}%
\TextAt{27}{200.512}{222.416}{p}%
\TextAt{8}{209.435}{224.21}{for}%
\TextAt{8}{227.859}{224.21}{all}%
\TextAt{14}{244.657}{224.21}{p}%
\TextAt{8}{250.64}{224.21}{.}%
\TextAt{8}{260.705}{224.21}{Note}%
\TextAt{8}{289.53}{224.21}{that}%
\TextAt{13}{315.438}{224.21}{H}%
\TextAt{17}{326.06}{228.548}{0}%
\TextAt{9}{330.793}{224.21}{\Character{40}}%
\TextAt{14}{335.345}{224.21}{o}%
\TextAt{9}{341.189}{224.21}{\Character{40}}%
\TextAt{13}{345.742}{224.21}{D}%
\TextAt{9}{355.786}{224.21}{\Character{41}\Character{41}}%
\TextAt{8}{369.327}{224.21}{is}%
\TextAt{8}{381.63}{224.21}{no}%
\TextAt{8}{398.428}{224.21}{longer}%
\TextAt{8}{434.091}{224.21}{a}%
\TextAt{8}{444.385}{224.21}{vector}%
\TextAt{8}{480.048}{224.21}{space.}%
\TextAt{8}{517.491}{224.21}{An}%
\TextAt{8}{79.226}{207.164}{analogue}%
\TextAt{8}{127.575}{207.164}{of}%
\TextAt{13}{140.464}{207.164}{H}%
\TextAt{17}{151.087}{211.502}{1}%
\TextAt{9}{155.819}{207.164}{\Character{40}}%
\TextAt{13}{160.372}{207.164}{X}%
\TextAt{13}{170.367}{207.164}{\Character{59}}%
\TextAt{14}{175.62}{207.164}{o}%
\TextAt{9}{181.465}{207.164}{\Character{40}}%
\TextAt{13}{186.017}{207.164}{D}%
\TextAt{9}{196.062}{207.164}{\Character{41}\Character{41}}%
\TextAt{8}{208.623}{207.164}{is}%
\TextAt{8}{219.947}{207.164}{completely}%
\TextAt{8}{278.361}{207.164}{missing.}%
\TextAt{8}{324.513}{207.164}{Instead}%
\TextAt{8}{365.425}{207.164}{of}%
\TextAt{8}{378.315}{207.164}{attacking}%
\TextAt{8}{429.27}{207.164}{directly}%
\TextAt{8}{471.426}{207.164}{the}%
\TextAt{8}{491.142}{207.164}{problem}%
\TextAt{8}{79.226}{190.118}{of}%
\TextAt{8}{94.182}{190.118}{measuring}%
\TextAt{8}{151.805}{190.118}{the}%
\TextAt{8}{173.587}{190.118}{size}%
\TextAt{8}{197.377}{190.118}{of}%
\TextAt{13}{212.333}{190.118}{H}%
\TextAt{17}{222.955}{194.456}{0}%
\TextAt{9}{227.688}{190.118}{\Character{40}}%
\TextAt{14}{232.24}{190.118}{o}%
\TextAt{9}{238.084}{190.118}{\Character{40}}%
\TextAt{13}{242.636}{190.118}{D}%
\TextAt{9}{252.681}{190.118}{\Character{41}\Character{41}}%
\TextAt{8}{261.786}{190.118}{,}%
\TextAt{8}{270.967}{190.118}{we}%
\TextAt{8}{289.819}{190.118}{proceed}%
\TextAt{8}{334.687}{190.118}{as}%
\TextAt{8}{350.682}{190.118}{in}%
\TextAt{8}{365.96}{190.118}{the}%
\TextAt{8}{387.742}{190.118}{function-theoretic}%
\TextAt{8}{483.993}{190.118}{model}%
\TextAt{8}{520.408}{190.118}{by}%
\TextAt{8}{79.226}{173.072}{looking}%
\TextAt{8}{120.491}{173.072}{at}%
\TextAt{8}{135.083}{173.072}{the}%
\TextAt{8}{155.522}{173.072}{“Euler–Poincaré}%
\TextAt{8}{242.622}{173.072}{characteristic”}%
\TextAt{8}{320.362}{173.072}{of}%
\TextAt{8}{333.974}{173.072}{the}%
\TextAt{8}{354.412}{173.072}{replete}%
\TextAt{8}{393.059}{173.072}{ideal}%
\TextAt{14}{421.304}{173.072}{o}%
\TextAt{9}{427.148}{173.072}{\Character{40}}%
\TextAt{13}{431.7}{173.072}{D}%
\TextAt{9}{441.745}{173.072}{\Character{41}}%
\TextAt{8}{446.297}{173.072}{.}%
\TextAt{8}{455.592}{173.072}{Before}%
\TextAt{8}{492.445}{173.072}{defining}%
\TextAt{8}{79.226}{156.026}{this,}%
\TextAt{8}{107.448}{156.026}{we}%
\TextAt{8}{126.394}{156.026}{want}%
\TextAt{8}{156.733}{156.026}{to}%
\TextAt{8}{172.762}{156.026}{establish}%
\TextAt{8}{222.731}{156.026}{the}%
\TextAt{8}{244.606}{156.026}{relation}%
\TextAt{8}{289.255}{156.026}{between}%
\TextAt{8}{336.164}{156.026}{the}%
\TextAt{8}{358.038}{156.026}{Minkowski}%
\TextAt{8}{418.026}{156.026}{space}%
\TextAt{13}{451.019}{156.026}{K}%
\TextAt{25}{460.978}{154.232}{R}%
\TextAt{9}{473.482}{156.026}{\Character{61}}%
\TextAt{9}{488.821}{156.026}{\Character{91}}%
\TextAt{20}{492.087}{164.992}{\Character{81}}%
\TextAt{21}{503.378}{152.539}{\Character{28}}%
\TextAt{18}{510.547}{156.026}{C}%
\TextAt{9}{519.181}{156.026}{\Character{93}}%
\TextAt{17}{522.433}{160.364}{\Character{43}}%
\TextAt{8}{529.518}{156.026}{,}%
\TextAt{13}{79.226}{138.98}{\Character{28}}%
\NativeMathAt{91.711}{138.98}{11.955}{\in}%
\TextAt{8}{105.747}{138.98}{Hom}%
\TextAt{9}{130.136}{138.98}{\Character{40}}%
\TextAt{13}{134.688}{138.98}{K}%
\TextAt{13}{144.85}{138.98}{\Character{59}}%
\TextAt{18}{150.1}{138.98}{C}%
\TextAt{9}{158.734}{138.98}{\Character{41}}%
\TextAt{8}{163.286}{138.98}{,}%
\TextAt{8}{172.449}{138.98}{and}%
\TextAt{8}{196.823}{138.98}{the}%
\TextAt{8}{218.59}{138.98}{product}%
\TextAt{20}{264.101}{147.946}{\Character{81}}%
\TextAt{27}{275.392}{135.493}{p}%
\TextAt{3}{279.38}{135.493}{\Character{106}}%
\NativeMathAt{281.731}{135.493}{7.97}{\infty}%
\TextAt{13}{292.692}{138.98}{K}%
\TextAt{27}{302.651}{137.186}{p}%
\TextAt{8}{307.138}{138.98}{.}%
\TextAt{8}{320.419}{138.98}{The}%
\TextAt{8}{346.084}{138.98}{reader}%
\TextAt{8}{383.464}{138.98}{will}%
\TextAt{8}{407.18}{138.98}{allow}%
\TextAt{8}{439.038}{138.98}{us}%
\TextAt{8}{455.665}{138.98}{to}%
\TextAt{8}{471.586}{138.98}{explain}%
\TextAt{8}{513.844}{138.98}{this}%
\TextAt{8}{79.226}{121.934}{simple}%
\TextAt{8}{115.701}{121.934}{situation}%
\TextAt{8}{164.55}{121.934}{in}%
\TextAt{8}{178.203}{121.934}{the}%
\TextAt{8}{198.359}{121.934}{following}%
\TextAt{8}{247.794}{121.934}{sketch.}%
\TextAt{8}{97.159}{101.899}{We}%
\TextAt{8}{117.303}{101.899}{have}%
\TextAt{8}{144.298}{101.899}{the}%
\TextAt{8}{164.455}{101.899}{correspondences}%
\TextAt{13}{170.373}{72.921}{\Character{26}}%
\TextAt{9}{179.731}{72.921}{\Character{58}}%
\TextAt{13}{186.304}{72.921}{K}%
\NativeMathAt{200.442}{72.921}{11.955}{\rightarrow}%
\TextAt{18}{215.718}{72.921}{R}%
\NativeArrowAt{239.629}{72.921}{11.955}{19.26}{\longmapsto}%
\TextAt{14}{274.166}{72.921}{p}%
\TextAt{8}{284.046}{72.921}{real}%
\TextAt{8}{306.809}{72.921}{prime}%
\TextAt{13}{336.075}{72.921}{\Character{59}}%
\TextAt{13}{353.255}{72.921}{\Character{26}}%
\TextAt{9}{362.633}{72.921}{\Character{61}}%
\TextAt{13}{375.058}{72.921}{\Character{26}}%
\TextAt{27}{381.096}{71.128}{p}%
\TextAt{9}{388.904}{72.921}{\Character{58}}%
\TextAt{13}{395.476}{72.921}{K}%
\TextAt{27}{405.435}{71.128}{p}%
\NativeMathAt{415.831}{79.698}{7.97}{\sim}%
\NativeArrowAt{413.243}{72.921}{11.955}{13.175}{\longrightarrow}%
\TextAt{18}{429.737}{72.921}{R}%
\TextAt{13}{438.371}{72.921}{\Character{59}}%
\TextAt{8}{297.211}{43.033}{191}%
\Destination{78.23}{809.815}{nav.page.191}%
\Destination{79.36}{776.155}{nav.print.209}%
\Destination{79.36}{256.255}{nav.print.210}%
%
%
\Hotlink{350.616}{591.855}{19.22}{11.047}{nav.cite.N051}%
\Hotlink{110.668}{537.755}{5.22}{11.047}{nav.chapter.1}%
\Hotlink{195.792}{537.755}{14.306}{11.047}{nav.chapter.3}%
\Hotlink{126.969}{537.755}{12.716}{11.047}{nav.section.1.14}%
\Hotlink{221.179}{537.755}{6.858}{11.047}{nav.section.3.1}%
\Destination{79.226}{505.811}{errata.98}%
}

% Page 204 (body page 192)
\Page{%
\TextAt{2}{79.226}{796.943}{III.}%
\TextAt{2}{98.317}{796.943}{Riemann–Roch}%
\TextAt{2}{174.375}{796.943}{Theory}%
\Vector{q 0 G 0 g 0.299 w 7.226 18.729 m 460.769 18.729 l S Q}%
\TextAt{13}{151.639}{761.124}{\Character{27}}%
\TextAt{13}{158.071}{761.124}{\Character{59}}%
\Vector{q 0 G 0 g 0.478 w 91.316 -6.055 m 98.398 -6.055 l S Q}%
\TextAt{13}{163.316}{761.124}{\Character{27}}%
\TextAt{9}{173.719}{761.124}{\Character{58}}%
\TextAt{13}{180.291}{761.124}{K}%
\NativeMathAt{194.43}{761.124}{11.955}{\rightarrow}%
\TextAt{18}{209.706}{761.124}{C}%
\NativeArrowAt{233.616}{761.124}{11.955}{19.26}{\longmapsto}%
\TextAt{14}{268.153}{761.124}{p}%
\TextAt{8}{278.034}{761.124}{complex}%
\TextAt{8}{323.882}{761.124}{prime}%
\TextAt{13}{353.148}{761.124}{\Character{59}}%
\TextAt{13}{370.328}{761.124}{\Character{27}}%
\TextAt{9}{380.751}{761.124}{\Character{61}}%
\TextAt{13}{393.176}{761.124}{\Character{27}}%
\TextAt{27}{399.83}{759.33}{p}%
\TextAt{9}{407.637}{761.124}{\Character{58}}%
\TextAt{13}{414.21}{761.124}{K}%
\TextAt{27}{424.169}{759.33}{p}%
\NativeMathAt{434.564}{767.901}{7.97}{\sim}%
\NativeArrowAt{431.977}{761.124}{11.955}{13.175}{\longrightarrow}%
\TextAt{18}{448.47}{761.124}{C}%
\TextAt{13}{457.105}{761.124}{\Character{58}}%
\TextAt{8}{79.226}{737.602}{There}%
\TextAt{8}{113.035}{737.602}{are}%
\TextAt{8}{132.546}{737.602}{the}%
\TextAt{8}{152.703}{737.602}{following}%
\TextAt{8}{202.137}{737.602}{isomorphisms}%
\TextAt{13}{203.443}{708.601}{K}%
\TextAt{16}{216.918}{708.601}{\Character{10}}%
\TextAt{25}{226.216}{706.808}{Q}%
\TextAt{18}{235.57}{708.601}{R}%
\NativeMathAt{250.113}{715.378}{7.97}{\sim}%
\NativeArrowAt{247.525}{708.601}{11.955}{13.175}{\longrightarrow}%
\TextAt{13}{264.019}{708.601}{K}%
\TextAt{25}{273.978}{706.808}{R}%
\TextAt{13}{280.233}{708.601}{\Character{59}}%
\TextAt{13}{309.368}{708.601}{a}%
\TextAt{16}{318.189}{708.601}{\Character{10}}%
\TextAt{13}{330.144}{708.601}{x}%
\NativeArrowAt{340.117}{708.601}{11.955}{19.26}{\longmapsto}%
\TextAt{9}{362.699}{708.601}{\Character{40}\Character{40}}%
\TextAt{13}{371.803}{708.601}{\Character{28}}%
\TextAt{13}{378.211}{708.601}{a}%
\TextAt{9}{384.367}{708.601}{\Character{41}}%
\TextAt{13}{388.92}{708.601}{x}%
\TextAt{9}{395.572}{708.601}{\Character{41}}%
\TextAt{21}{400.124}{706.808}{\Character{28}}%
\TextAt{13}{405.301}{708.601}{\Character{59}}%
\TextAt{13}{189.052}{679.6}{K}%
\TextAt{16}{202.526}{679.6}{\Character{10}}%
\TextAt{25}{211.824}{677.806}{Q}%
\TextAt{18}{221.178}{679.6}{R}%
\NativeMathAt{235.721}{686.377}{7.97}{\sim}%
\NativeArrowAt{233.133}{679.6}{11.955}{13.175}{\longrightarrow}%
\TextAt{20}{249.627}{690.957}{\Character{89}}%
\TextAt{27}{249.861}{665.054}{p}%
\TextAt{3}{253.849}{665.054}{\Character{106}}%
\NativeMathAt{256.2}{665.054}{7.97}{\infty}%
\TextAt{13}{266.896}{679.6}{K}%
\TextAt{27}{276.855}{677.806}{p}%
\TextAt{13}{281.342}{679.6}{\Character{59}}%
\TextAt{13}{310.477}{679.6}{a}%
\TextAt{16}{319.298}{679.6}{\Character{10}}%
\TextAt{13}{331.253}{679.6}{x}%
\NativeArrowAt{341.226}{679.6}{11.955}{19.26}{\longmapsto}%
\TextAt{9}{363.808}{679.6}{\Character{40}\Character{40}}%
\TextAt{13}{372.913}{679.6}{\Character{28}}%
\TextAt{27}{377.996}{677.806}{p}%
\TextAt{13}{382.483}{679.6}{a}%
\TextAt{9}{388.628}{679.6}{\Character{41}}%
\TextAt{13}{393.18}{679.6}{x}%
\TextAt{9}{399.832}{679.6}{\Character{41}}%
\TextAt{27}{404.385}{677.74}{p}%
\TextAt{3}{408.374}{677.74}{\Character{106}}%
\NativeMathAt{410.725}{677.74}{7.97}{\infty}%
\TextAt{13}{419.692}{679.6}{\Character{59}}%
\TextAt{13}{79.226}{645.428}{\Character{28}}%
\TextAt{27}{84.31}{643.635}{p}%
\TextAt{8}{94.469}{645.428}{being}%
\TextAt{8}{127.781}{645.428}{the}%
\TextAt{8}{149.712}{645.428}{canonical}%
\TextAt{8}{202.87}{645.428}{embedding}%
\TextAt{13}{263.823}{645.428}{K}%
\NativeMathAt{280.985}{645.428}{11.955}{\rightarrow}%
\TextAt{13}{299.286}{645.428}{K}%
\TextAt{27}{309.245}{643.635}{p}%
\TextAt{8}{319.404}{645.428}{(see}%
\TextAt{8}{344.646}{645.428}{chap.}%
\TextAt{8}{382.161}{645.428}{II,}%
\TextAt{8}{399.848}{645.428}{(8.3)).}%
\TextAt{8}{442.253}{645.428}{They}%
\TextAt{8}{474.262}{645.428}{fit}%
\TextAt{8}{490.993}{645.428}{into}%
\TextAt{8}{516.51}{645.428}{the}%
\TextAt{8}{79.226}{628.382}{commutative}%
\TextAt{8}{149.152}{628.382}{diagram}%
\TextAt{13}{176.55}{596.046}{K}%
\TextAt{16}{190.025}{596.046}{\Character{10}}%
\TextAt{18}{201.98}{596.046}{R}%
\TextAt{13}{250.189}{596.046}{K}%
\TextAt{25}{260.148}{594.253}{R}%
\TextAt{20}{311.473}{607.403}{\Character{89}}%
\TextAt{21}{316.943}{582.297}{\Character{26}}%
\TextAt{18}{328.741}{596.046}{R}%
\TextAt{20}{374.696}{607.403}{\Character{89}}%
\TextAt{21}{379.771}{582.297}{\Character{27}}%
\TextAt{9}{389.972}{596.046}{\Character{91}}%
\TextAt{18}{393.223}{596.046}{C}%
\TextAt{16}{404.514}{596.046}{\Character{2}}%
\TextAt{18}{416.469}{596.046}{C}%
\TextAt{9}{425.104}{596.046}{\Character{93}}%
\TextAt{17}{428.355}{600.982}{\Character{43}}%
\TextAt{13}{176.55}{510.419}{K}%
\TextAt{16}{190.025}{510.419}{\Character{10}}%
\TextAt{18}{201.98}{510.419}{R}%
\TextAt{20}{242.439}{521.776}{\Character{89}}%
\TextAt{27}{242.673}{495.873}{p}%
\TextAt{3}{246.661}{495.873}{\Character{106}}%
\NativeMathAt{249.012}{495.873}{7.97}{\infty}%
\TextAt{13}{259.708}{510.419}{K}%
\TextAt{27}{269.667}{508.626}{p}%
\TextAt{20}{308.567}{521.776}{\Character{89}}%
\TextAt{27}{305.978}{496.32}{p}%
\TextAt{7}{312.788}{496.32}{real}%
\TextAt{13}{328.425}{510.419}{K}%
\TextAt{27}{338.384}{508.626}{p}%
\TextAt{20}{389.211}{521.776}{\Character{89}}%
\TextAt{27}{378.277}{496.32}{p}%
\TextAt{7}{385.087}{496.32}{complex}%
\TextAt{13}{417.413}{510.419}{K}%
\TextAt{27}{427.373}{508.626}{p}%
\NativeMathAt{225.754}{595.433}{11.955}{\cong}%
\NativeDiagram{%
\path[draw=black,line width=2.797bp] (121.583,-183.528) -- (121.583,-250.766);%
\path[draw=black!0,line width=1.841bp] (121.583,-183.528) -- (121.583,-250.766);%
}%
\TextAt{9}{284.388}{596.879}{\Character{61}}%
\TextAt{16}{351.39}{596.047}{\Character{2}}%
\NativeMathAt{221.879}{509.805}{11.955}{\cong}%
\TextAt{9}{285.516}{511.251}{\Character{61}}%
\NativeDiagram{%
\path[draw=black,line width=0.478bp] (186.298,-246.448) -- (186.298,-184.803);%
\path[draw=black,line width=0.478bp,line cap=round,line join=round] (183.429,-187.051) .. controls (185.15,-186.596) and (185.963,-185.584) .. (186.298,-184.564) .. controls (186.633,-185.584) and (187.446,-186.596) .. (189.167,-187.051);%
}%
\TextAt{3}{261.113}{556.621}{\Character{111}}%
\TextAt{16}{355.928}{510.419}{\Character{2}}%
\NativeDiagram{%
\path[draw=black,line width=0.478bp] (252.428,-246.448) -- (252.428,-199.504);%
\path[draw=black,line width=0.478bp,line cap=round,line join=round] (249.558,-201.751) .. controls (251.28,-201.297) and (252.093,-200.285) .. (252.428,-199.265) .. controls (252.762,-200.285) and (253.575,-201.297) .. (255.297,-201.751);%
}%
\TextAt{30}{303.783}{555.316}{\Character{81}}%
\TextAt{21}{313.192}{549.338}{\Character{26}}%
\TextAt{28}{317.529}{548.342}{p}%
\TextAt{3}{327.243}{549.27}{\Character{111}}%
\NativeDiagram{%
\path[draw=black,line width=0.478bp] (333.072,-246.448) -- (333.072,-197.954);%
\path[draw=black,line width=0.478bp,line cap=round,line join=round] (330.203,-200.202) .. controls (331.925,-199.747) and (332.737,-198.735) .. (333.072,-197.715) .. controls (333.407,-198.735) and (334.22,-199.747) .. (335.941,-200.202);%
}%
\TextAt{30}{364.837}{556.091}{\Character{81}}%
\TextAt{17}{372.835}{550.113}{\Character{91}}%
\TextAt{21}{375.187}{550.113}{\Character{27}}%
\TextAt{28}{380.026}{549.117}{p}%
\TextAt{3}{384.109}{550.113}{\Character{2}}%
\NativeDiagram{%
\path[draw=black,line width=0.359bp] (318.696,-219.2) -- (323.821,-219.2);%
}%
\TextAt{21}{390.696}{550.113}{\Character{27}}%
\TextAt{28}{395.821}{549.117}{p}%
\TextAt{17}{399.904}{550.113}{\Character{93}}%
\TextAt{3}{407.888}{550.045}{\Character{111}}%
\TextAt{8}{79.226}{466.389}{where}%
\TextAt{8}{112.824}{466.389}{the}%
\TextAt{8}{132.769}{466.389}{arrow}%
\TextAt{8}{165.41}{466.389}{on}%
\TextAt{8}{181.458}{466.389}{the}%
\TextAt{8}{201.403}{466.389}{right}%
\TextAt{8}{229.489}{466.389}{is}%
\TextAt{8}{241.042}{466.389}{given}%
\TextAt{8}{271.4}{466.389}{by}%
\TextAt{13}{287.447}{466.389}{a}%
\NativeArrowAt{296.913}{466.389}{11.955}{11.955}{\mapsto}%
\TextAt{9}{312.189}{466.389}{\Character{40}}%
\TextAt{13}{316.742}{466.389}{\Character{27}}%
\TextAt{13}{323.808}{466.389}{a\Character{59}}%
\Vector{q 0 G 0 g 0.478 w 263.213 -300.79 m 270.295 -300.79 l S Q}%
\TextAt{13}{335.213}{466.389}{\Character{27}}%
\TextAt{13}{342.279}{466.389}{a}%
\TextAt{9}{348.44}{466.389}{\Character{41}}%
\TextAt{8}{352.993}{466.389}{.}%
\TextAt{8}{361.371}{466.389}{Via}%
\TextAt{8}{382.941}{466.389}{this}%
\TextAt{8}{405.553}{466.389}{isomorphism,}%
\TextAt{8}{477.055}{466.389}{we}%
\TextAt{8}{494.071}{466.389}{identify}%
\TextAt{13}{79.226}{449.343}{K}%
\TextAt{25}{89.185}{447.549}{R}%
\TextAt{8}{99.337}{449.343}{with}%
\TextAt{20}{125.997}{458.309}{\Character{81}}%
\TextAt{27}{137.288}{445.856}{p}%
\TextAt{3}{141.277}{445.856}{\Character{106}}%
\NativeMathAt{143.628}{445.856}{7.97}{\infty}%
\TextAt{13}{154.588}{449.343}{K}%
\TextAt{27}{164.547}{447.549}{p}%
\TextAt{8}{169.034}{449.343}{:}%
\TextAt{13}{272.535}{430.314}{K}%
\TextAt{25}{282.494}{428.521}{R}%
\TextAt{9}{292.069}{430.314}{\Character{61}}%
\TextAt{20}{304.495}{441.671}{\Character{89}}%
\TextAt{27}{304.728}{415.768}{p}%
\TextAt{3}{308.717}{415.768}{\Character{106}}%
\NativeMathAt{311.068}{415.768}{7.97}{\infty}%
\TextAt{13}{321.763}{430.314}{K}%
\TextAt{27}{331.722}{428.521}{p}%
\TextAt{13}{336.209}{430.314}{\Character{58}}%
\TextAt{8}{79.226}{396.142}{The}%
\TextAt{8}{103.28}{396.142}{scalar}%
\TextAt{8}{136.515}{396.142}{product}%
\TextAt{16}{180.415}{396.142}{\Character{104}}%
\TextAt{13}{185.064}{396.142}{x\Character{59}}%
\TextAt{13}{196.936}{396.142}{y}%
\TextAt{16}{203.097}{396.142}{\Character{105}}%
\TextAt{9}{211.067}{396.142}{\Character{61}}%
\TextAt{20}{223.492}{405.108}{\Character{80}}%
\TextAt{21}{236.112}{392.655}{\Character{28}}%
\TextAt{13}{243.281}{396.142}{x}%
\TextAt{21}{249.933}{394.349}{\Character{28}}%
\Vector{q 0 G 0 g 0.478 w 183.109 -371.037 m 189.246 -371.037 l S Q}%
\TextAt{13}{255.109}{396.142}{y}%
\TextAt{21}{261.246}{393.319}{\Character{28}}%
\TextAt{8}{270.32}{396.142}{on}%
\TextAt{13}{286.579}{396.142}{K}%
\TextAt{25}{296.538}{394.349}{R}%
\TextAt{8}{306.69}{396.142}{is}%
\TextAt{8}{318.454}{396.142}{then}%
\TextAt{8}{345.114}{396.142}{transformed}%
\TextAt{8}{410.544}{396.142}{into}%
\TextAt{16}{201.852}{367.141}{\Character{104}}%
\TextAt{13}{206.501}{367.141}{x\Character{59}}%
\TextAt{13}{218.373}{367.141}{y}%
\TextAt{16}{224.534}{367.141}{\Character{105}}%
\TextAt{9}{232.504}{367.141}{\Character{61}}%
\TextAt{20}{246.523}{378.498}{\Character{88}}%
\TextAt{27}{244.93}{353.041}{p}%
\TextAt{7}{251.74}{353.041}{real}%
\TextAt{13}{267.377}{367.141}{x}%
\TextAt{27}{274.029}{365.347}{p}%
\TextAt{13}{278.516}{367.141}{y}%
\TextAt{27}{284.224}{365.347}{p}%
\TextAt{9}{291.368}{367.141}{\Character{43}}%
\TextAt{20}{313.067}{378.498}{\Character{88}}%
\TextAt{27}{303.129}{353.041}{p}%
\TextAt{7}{309.939}{353.041}{complex}%
\TextAt{9}{340.273}{367.141}{\Character{40}}%
\TextAt{13}{344.826}{367.141}{x}%
\TextAt{27}{351.478}{365.347}{p}%
\Vector{q 0 G 0 g 0.478 w 283.965 -400.038 m 290.101 -400.038 l S Q}%
\TextAt{13}{355.965}{367.141}{y}%
\TextAt{27}{362.101}{364.318}{p}%
\TextAt{9}{369.245}{367.141}{\Character{43}}%
\Vector{q 0 G 0 g 0.478 w 309.006 -400.038 m 315.658 -400.038 l S Q}%
\TextAt{13}{381.006}{367.141}{x}%
\TextAt{27}{387.658}{365.347}{p}%
\TextAt{13}{392.145}{367.141}{y}%
\TextAt{27}{397.853}{365.347}{p}%
\TextAt{9}{402.34}{367.141}{\Character{41}}%
\TextAt{13}{406.892}{367.141}{\Character{58}}%
\TextAt{8}{79.226}{328.382}{The}%
\TextAt{8}{103.28}{328.382}{Haar}%
\TextAt{8}{132.223}{328.382}{measure}%
\TextAt{13}{177.808}{328.382}{\Character{22}}%
\TextAt{8}{188.749}{328.382}{on}%
\TextAt{13}{205.008}{328.382}{K}%
\TextAt{25}{214.967}{326.589}{R}%
\TextAt{8}{225.119}{328.382}{which}%
\TextAt{8}{258.605}{328.382}{is}%
\TextAt{8}{270.369}{328.382}{determined}%
\TextAt{8}{331.496}{328.382}{by}%
\TextAt{16}{347.755}{328.382}{\Character{104}}%
\TextAt{13}{352.404}{328.382}{x\Character{59}}%
\TextAt{13}{364.276}{328.382}{y}%
\TextAt{16}{370.437}{328.382}{\Character{105}}%
\TextAt{8}{378.984}{328.382}{becomes}%
\TextAt{8}{425.537}{328.382}{the}%
\TextAt{8}{445.693}{328.382}{product}%
\TextAt{8}{489.593}{328.382}{measure}%
\TextAt{13}{278.578}{299.381}{\Character{22}}%
\TextAt{9}{288.942}{299.381}{\Character{61}}%
\TextAt{20}{301.367}{310.738}{\Character{89}}%
\TextAt{27}{301.601}{284.835}{p}%
\TextAt{3}{305.589}{284.835}{\Character{106}}%
\NativeMathAt{307.94}{284.835}{7.97}{\infty}%
\TextAt{13}{318.636}{299.381}{\Character{22}}%
\TextAt{27}{325.679}{297.588}{p}%
\TextAt{13}{330.166}{299.381}{\Character{59}}%
\TextAt{8}{79.226}{260.399}{where}%
\TextAt{13}{186.755}{243.353}{\Character{22}}%
\TextAt{27}{193.798}{241.56}{p}%
\TextAt{9}{201.606}{243.353}{\Character{61}}%
\TextAt{8}{214.031}{243.353}{Lebesgue}%
\TextAt{8}{264.649}{243.353}{measure}%
\TextAt{8}{310.234}{243.353}{on}%
\TextAt{13}{326.493}{243.353}{K}%
\TextAt{27}{336.453}{241.56}{p}%
\TextAt{9}{344.26}{243.353}{\Character{61}}%
\TextAt{18}{356.686}{243.353}{R}%
\TextAt{13}{365.32}{243.353}{\Character{59}}%
\TextAt{8}{382.519}{243.353}{if}%
\TextAt{14}{393.243}{243.353}{p}%
\TextAt{8}{403.124}{243.353}{real}%
\TextAt{13}{421.989}{243.353}{\Character{59}}%
\TextAt{13}{170.337}{219.831}{\Character{22}}%
\TextAt{27}{177.38}{218.038}{p}%
\TextAt{9}{185.188}{219.831}{\Character{61}}%
\TextAt{9}{197.598}{219.831}{2}%
\TextAt{8}{207.364}{219.831}{Lebesgue}%
\TextAt{8}{257.982}{219.831}{measure}%
\TextAt{8}{303.567}{219.831}{on}%
\TextAt{13}{319.826}{219.831}{K}%
\TextAt{27}{329.785}{218.038}{p}%
\TextAt{9}{337.593}{219.831}{\Character{61}}%
\TextAt{18}{350.018}{219.831}{C}%
\TextAt{13}{358.653}{219.831}{\Character{59}}%
\TextAt{8}{375.852}{219.831}{if}%
\TextAt{14}{386.576}{219.831}{p}%
\TextAt{8}{396.456}{219.831}{complex}%
\TextAt{13}{438.407}{219.831}{\Character{58}}%
\TextAt{8}{79.226}{196.31}{Indeed,}%
\TextAt{8}{120.056}{196.31}{the}%
\TextAt{8}{139.648}{196.31}{system}%
\TextAt{9}{177.902}{196.31}{1\Character{47}}%
\TextAt{16}{189.608}{206.199}{\Character{112}}%
\Vector{q 0 G 0 g 0.478 w 127.571 -567.562 m 133.424 -567.562 l S Q}%
\TextAt{9}{199.571}{196.31}{2}%
\TextAt{13}{205.424}{196.31}{\Character{59}}%
\TextAt{13}{210.648}{196.31}{i}%
\TextAt{9}{214.661}{196.31}{\Character{47}}%
\TextAt{16}{220.514}{206.199}{\Character{112}}%
\Vector{q 0 G 0 g 0.478 w 158.477 -567.562 m 164.33 -567.562 l S Q}%
\TextAt{9}{230.477}{196.31}{2}%
\TextAt{8}{239.663}{196.31}{is}%
\TextAt{8}{250.862}{196.31}{an}%
\TextAt{8}{266.557}{196.31}{orthonormal}%
\TextAt{8}{333.001}{196.31}{basis}%
\TextAt{8}{361.177}{196.31}{with}%
\TextAt{8}{387.273}{196.31}{respect}%
\TextAt{8}{426.758}{196.31}{to}%
\TextAt{8}{440.504}{196.31}{the}%
\TextAt{8}{460.096}{196.31}{scalar}%
\TextAt{8}{492.767}{196.31}{product}%
\TextAt{13}{79.226}{179.264}{x}%
\Vector{q 0 G 0 g 0.478 w 13.878 -587.915 m 20.015 -587.915 l S Q}%
\TextAt{13}{85.878}{179.264}{y}%
\TextAt{9}{94.878}{179.264}{\Character{43}}%
\Vector{q 0 G 0 g 0.478 w 34.845 -587.915 m 41.497 -587.915 l S Q}%
\TextAt{13}{106.845}{179.264}{xy}%
\TextAt{8}{123.834}{179.264}{on}%
\TextAt{13}{140.395}{179.264}{K}%
\TextAt{27}{150.354}{177.47}{p}%
\TextAt{9}{158.677}{179.264}{\Character{61}}%
\TextAt{18}{171.618}{179.264}{C}%
\TextAt{8}{180.253}{179.264}{.}%
\TextAt{8}{189.608}{179.264}{Hence}%
\TextAt{8}{224.688}{179.264}{the}%
\TextAt{8}{245.147}{179.264}{square}%
\TextAt{13}{282.26}{179.264}{Q}%
\TextAt{9}{295.367}{179.264}{\Character{61}}%
\TextAt{16}{308.308}{179.264}{\Character{102}}%
\TextAt{13}{314.286}{179.264}{z}%
\TextAt{9}{324.092}{179.264}{\Character{61}}%
\TextAt{13}{337.033}{179.264}{x}%
\TextAt{9}{346.548}{179.264}{\Character{43}}%
\TextAt{13}{358.516}{179.264}{iy}%
\TextAt{16}{372.482}{179.264}{\Character{106}}%
\TextAt{9}{379.639}{179.264}{0}%
\NativeMathAt{389.328}{179.264}{11.955}{\leq}%
\TextAt{13}{402.463}{179.264}{x\Character{59}}%
\TextAt{13}{414.334}{179.264}{y}%
\NativeMathAt{424.332}{179.264}{11.955}{\leq}%
\TextAt{9}{437.467}{179.264}{1\Character{47}}%
\TextAt{16}{449.173}{189.153}{\Character{112}}%
\Vector{q 0 G 0 g 0.478 w 387.136 -584.608 m 392.989 -584.608 l S Q}%
\TextAt{9}{459.136}{179.264}{2}%
\TextAt{16}{464.989}{179.264}{\Character{103}}%
\TextAt{8}{475.166}{179.264}{has}%
\TextAt{8}{496.342}{179.264}{volume}%
\TextAt{13}{79.226}{162.218}{\Character{22}}%
\TextAt{27}{86.269}{160.424}{p}%
\TextAt{9}{90.756}{162.218}{\Character{40}}%
\TextAt{13}{95.308}{162.218}{Q}%
\TextAt{9}{104.58}{162.218}{\Character{41}}%
\TextAt{9}{112.435}{162.218}{\Character{61}}%
\TextAt{9}{124.856}{162.218}{1}%
\TextAt{8}{130.731}{162.218}{,}%
\TextAt{8}{137.88}{162.218}{but}%
\TextAt{8}{159.34}{162.218}{Lebesgue}%
\TextAt{8}{209.958}{162.218}{volume}%
\TextAt{9}{250.283}{162.218}{1\Character{47}2}%
\TextAt{8}{267.842}{162.218}{.}%
\TextAt{8}{97.159}{142.183}{Finally,}%
\TextAt{8}{139.265}{142.183}{the}%
\TextAt{8}{159.421}{142.183}{logarithm}%
\TextAt{8}{212.765}{142.183}{map}%
\TextAt{13}{179.635}{104.198}{\Character{96}}%
\TextAt{9}{187.868}{104.198}{\Character{58}}%
\TextAt{20}{194.441}{124.642}{\Character{34}}%
\TextAt{20}{201.415}{115.556}{\Character{89}}%
\TextAt{21}{206.714}{90.45}{\Character{28}}%
\TextAt{18}{218.683}{104.198}{C}%
\TextAt{3}{227.318}{109.135}{\Character{3}}%
\TextAt{20}{232.05}{124.642}{\Character{35}}%
\TextAt{17}{239.024}{121.965}{\Character{43}}%
\NativeArrowAt{249.429}{104.198}{11.955}{19.26}{\longrightarrow}%
\TextAt{20}{272.011}{124.642}{\Character{34}}%
\TextAt{20}{278.985}{115.556}{\Character{89}}%
\TextAt{21}{284.284}{90.45}{\Character{28}}%
\TextAt{18}{296.254}{104.198}{R}%
\TextAt{20}{304.888}{124.642}{\Character{35}}%
\TextAt{17}{311.862}{121.965}{\Character{43}}%
\TextAt{13}{320.939}{104.198}{\Character{59}}%
\TextAt{13}{350.074}{104.198}{x}%
\NativeArrowAt{360.066}{104.198}{11.955}{19.26}{\longmapsto}%
\TextAt{9}{382.648}{104.198}{\Character{40}}%
\TextAt{8}{387.201}{104.198}{log}%
\TextAt{16}{404.161}{104.198}{\Character{106}}%
\TextAt{13}{407.482}{104.198}{x}%
\TextAt{21}{414.134}{102.405}{\Character{28}}%
\TextAt{16}{419.311}{104.198}{\Character{106}}%
\TextAt{9}{422.631}{104.198}{\Character{41}}%
\TextAt{21}{427.184}{102.405}{\Character{28}}%
\TextAt{8}{297.211}{43.033}{192}%
\Destination{78.23}{809.815}{nav.page.192}%
\Destination{79.36}{481.355}{nav.print.211}%
\Hotlink{381.661}{643.255}{9.763}{11.047}{nav.chapter.2}%
\Hotlink{399.348}{643.255}{25.078}{11.047}{nav.number.2.8.3}%
}

% Page 205 (body page 193)
\Page{%
\TextAt{2}{79.226}{796.943}{III.}%
\TextAt{2}{98.317}{796.943}{Riemann–Roch}%
\TextAt{2}{174.375}{796.943}{Theory}%
\Vector{q 0 G 0 g 0.299 w 7.226 18.729 m 460.769 18.729 l S Q}%
\TextAt{8}{79.226}{761.124}{studied}%
\TextAt{8}{120.256}{761.124}{in}%
\TextAt{8}{133.909}{761.124}{Minkowski}%
\TextAt{8}{192.179}{761.124}{theory}%
\TextAt{8}{228.929}{761.124}{is}%
\TextAt{8}{240.693}{761.124}{transformed}%
\TextAt{8}{306.123}{761.124}{into}%
\TextAt{8}{329.866}{761.124}{the}%
\TextAt{8}{350.023}{761.124}{mapping}%
\TextAt{13}{199.636}{733.785}{\Character{96}}%
\TextAt{9}{207.869}{733.785}{\Character{58}}%
\TextAt{20}{214.442}{745.142}{\Character{89}}%
\TextAt{27}{214.675}{719.239}{p}%
\TextAt{3}{218.664}{719.239}{\Character{106}}%
\NativeMathAt{221.015}{719.239}{7.97}{\infty}%
\TextAt{13}{231.71}{733.785}{K}%
\TextAt{3}{242.528}{738.721}{\Character{3}}%
\TextAt{27}{241.669}{730.829}{p}%
\NativeArrowAt{250.581}{733.785}{11.955}{19.26}{\longrightarrow}%
\TextAt{20}{273.163}{745.142}{\Character{89}}%
\TextAt{27}{273.396}{719.239}{p}%
\TextAt{3}{277.385}{719.239}{\Character{106}}%
\NativeMathAt{279.736}{719.239}{7.97}{\infty}%
\TextAt{18}{290.432}{733.785}{R}%
\TextAt{13}{299.066}{733.785}{\Character{59}}%
\TextAt{13}{328.201}{733.785}{x}%
\NativeArrowAt{338.193}{733.785}{11.955}{19.26}{\longmapsto}%
\TextAt{9}{360.775}{733.785}{\Character{40}}%
\TextAt{8}{365.328}{733.785}{log}%
\TextAt{16}{382.288}{733.785}{\Character{106}}%
\TextAt{13}{385.609}{733.785}{x}%
\TextAt{27}{392.261}{731.991}{p}%
\TextAt{16}{396.748}{733.785}{\Character{106}}%
\TextAt{27}{400.069}{731.991}{p}%
\TextAt{9}{404.556}{733.785}{\Character{41}}%
\TextAt{13}{409.108}{733.785}{\Character{59}}%
\TextAt{8}{79.226}{696.345}{for}%
\TextAt{8}{97.111}{696.345}{one}%
\TextAt{8}{118.571}{696.345}{has}%
\TextAt{8}{139.444}{696.345}{the}%
\TextAt{8}{159.601}{696.345}{commutative}%
\TextAt{8}{229.527}{696.345}{diagram}%
\TextAt{13}{249.251}{668.69}{K}%
\TextAt{3}{260.069}{673.028}{\Character{3}}%
\TextAt{25}{259.21}{665.625}{R}%
\TextAt{9}{332.773}{668.69}{\Character{91}}%
\TextAt{20}{336.025}{677.656}{\Character{81}}%
\TextAt{21}{347.316}{665.203}{\Character{28}}%
\TextAt{18}{354.485}{668.69}{R}%
\TextAt{9}{363.119}{668.69}{\Character{93}}%
\TextAt{17}{366.37}{673.028}{\Character{43}}%
\TextAt{20}{235.287}{624.164}{\Character{81}}%
\TextAt{27}{246.578}{611.711}{p}%
\TextAt{3}{250.567}{611.711}{\Character{106}}%
\NativeMathAt{252.918}{611.711}{7.97}{\infty}%
\TextAt{13}{263.878}{615.198}{K}%
\TextAt{3}{274.696}{619.536}{\Character{3}}%
\TextAt{27}{273.837}{612.242}{p}%
\TextAt{20}{334.501}{624.164}{\Character{81}}%
\TextAt{27}{345.792}{611.711}{p}%
\TextAt{3}{349.781}{611.711}{\Character{106}}%
\NativeMathAt{352.132}{611.711}{7.97}{\infty}%
\TextAt{18}{363.092}{615.198}{R}%
\NativeDiagram{%
\path[draw=black,line width=0.478bp] (198.817,-102.321) -- (254.944,-102.321);%
\path[draw=black,line width=0.478bp,line cap=round,line join=round] (252.696,-99.452) .. controls (253.15,-101.173) and (254.163,-101.986) .. (255.183,-102.321) .. controls (254.163,-102.656) and (253.15,-103.469) .. (252.696,-105.19);%
}%
\TextAt{21}{297.39}{674.494}{\Character{96}}%
\NativeDiagram{%
\path[draw=black,line width=0.478bp] (185.358,-112.954) -- (185.358,-144.779);%
\path[draw=black,line width=0.478bp,line cap=round,line join=round] (188.227,-142.531) .. controls (186.506,-142.985) and (185.693,-143.998) .. (185.358,-145.018) .. controls (185.024,-143.998) and (184.211,-142.985) .. (182.489,-142.531);%
\path[draw=black,line width=0.478bp] (281.115,-113.376) -- (281.115,-144.779);%
\path[draw=black,line width=0.478bp,line cap=round,line join=round] (283.984,-142.531) .. controls (282.263,-142.985) and (281.45,-143.998) .. (281.115,-145.018) .. controls (280.78,-143.998) and (279.968,-142.985) .. (278.246,-142.531);%
\path[draw=black,line width=0.478bp] (212.781,-155.814) -- (256.672,-155.814);%
\path[draw=black,line width=0.478bp,line cap=round,line join=round] (254.425,-152.944) .. controls (254.879,-154.666) and (255.891,-155.479) .. (256.911,-155.814) .. controls (255.891,-156.148) and (254.879,-156.961) .. (254.425,-158.683);%
}%
\TextAt{21}{305.236}{621.002}{\Character{96}}%
\TextAt{13}{378.608}{643.781}{\Character{59}}%
\TextAt{8}{79.226}{585.752}{where}%
\TextAt{8}{113.035}{585.752}{the}%
\TextAt{8}{133.192}{585.752}{arrow}%
\TextAt{8}{166.045}{585.752}{on}%
\TextAt{8}{182.304}{585.752}{the}%
\TextAt{8}{202.46}{585.752}{right,}%
\TextAt{20}{196.499}{569.873}{\Character{34}}%
\TextAt{20}{203.473}{560.787}{\Character{89}}%
\TextAt{21}{208.771}{535.681}{\Character{28}}%
\TextAt{18}{220.741}{549.43}{R}%
\TextAt{20}{229.375}{569.873}{\Character{35}}%
\TextAt{17}{236.349}{567.196}{\Character{43}}%
\TextAt{9}{246.755}{549.43}{\Character{61}}%
\TextAt{20}{259.18}{560.787}{\Character{89}}%
\TextAt{21}{264.65}{535.681}{\Character{26}}%
\TextAt{18}{276.449}{549.43}{R}%
\TextAt{16}{287.74}{549.43}{\Character{2}}%
\TextAt{20}{299.695}{560.787}{\Character{89}}%
\TextAt{21}{304.77}{535.681}{\Character{27}}%
\TextAt{9}{314.971}{549.43}{\Character{91}}%
\TextAt{18}{318.223}{549.43}{R}%
\TextAt{16}{329.514}{549.43}{\Character{2}}%
\TextAt{18}{341.469}{549.43}{R}%
\TextAt{9}{350.103}{549.43}{\Character{93}}%
\TextAt{17}{353.355}{554.366}{\Character{43}}%
\NativeArrowAt{363.76}{549.43}{11.955}{19.26}{\longrightarrow}%
\TextAt{20}{386.342}{560.787}{\Character{89}}%
\TextAt{27}{386.576}{534.884}{p}%
\TextAt{3}{390.564}{534.884}{\Character{106}}%
\NativeMathAt{392.915}{534.884}{7.97}{\infty}%
\TextAt{18}{403.611}{549.43}{R}%
\TextAt{13}{412.245}{549.43}{\Character{59}}%
\TextAt{8}{79.226}{511.44}{is}%
\TextAt{8}{91.435}{511.44}{defined}%
\TextAt{8}{132.194}{511.44}{by}%
\TextAt{13}{148.898}{511.44}{x}%
\NativeArrowAt{159.63}{511.44}{11.955}{11.955}{\mapsto}%
\TextAt{13}{175.665}{511.44}{x}%
\TextAt{8}{186.659}{511.44}{for}%
\TextAt{13}{204.99}{511.44}{\Character{26}}%
\NativeMathAt{215.107}{511.44}{11.955}{\leftrightarrow}%
\TextAt{14}{231.142}{511.44}{p}%
\TextAt{8}{237.125}{511.44}{,}%
\TextAt{8}{244.831}{511.44}{and}%
\TextAt{8}{268.039}{511.44}{by}%
\TextAt{9}{284.743}{511.44}{\Character{40}}%
\TextAt{13}{289.296}{511.44}{x\Character{59}}%
\TextAt{13}{301.168}{511.44}{x}%
\TextAt{9}{307.844}{511.44}{\Character{41}}%
\NativeArrowAt{316.476}{511.44}{11.955}{11.955}{\mapsto}%
\TextAt{9}{332.511}{511.44}{2}%
\TextAt{13}{338.364}{511.44}{x}%
\TextAt{8}{349.359}{511.44}{for}%
\TextAt{13}{367.689}{511.44}{\Character{27}}%
\NativeMathAt{378.851}{511.44}{11.955}{\leftrightarrow}%
\TextAt{14}{394.886}{511.44}{p}%
\TextAt{8}{400.869}{511.44}{.}%
\TextAt{8}{410.653}{511.44}{This}%
\TextAt{8}{437.819}{511.44}{isomorphism}%
\TextAt{8}{506.683}{511.44}{takes}%
\TextAt{8}{79.226}{494.394}{the}%
\TextAt{8}{99.762}{494.394}{trace}%
\TextAt{8}{129.408}{494.394}{map}%
\TextAt{13}{155.803}{494.394}{x}%
\NativeArrowAt{166.423}{494.394}{11.955}{11.955}{\mapsto}%
\TextAt{20}{182.346}{503.361}{\Character{80}}%
\TextAt{21}{194.965}{490.907}{\Character{28}}%
\TextAt{13}{202.134}{494.394}{x}%
\TextAt{21}{208.786}{492.601}{\Character{28}}%
\TextAt{8}{218.24}{494.394}{on}%
\TextAt{9}{234.879}{494.394}{\Character{91}}%
\TextAt{20}{238.13}{503.361}{\Character{81}}%
\TextAt{21}{249.421}{490.907}{\Character{28}}%
\TextAt{18}{256.59}{494.394}{R}%
\TextAt{9}{265.225}{494.394}{\Character{93}}%
\TextAt{17}{268.476}{498.733}{\Character{43}}%
\TextAt{8}{279.838}{494.394}{into}%
\TextAt{8}{303.961}{494.394}{the}%
\TextAt{8}{324.497}{494.394}{trace}%
\TextAt{8}{354.143}{494.394}{map}%
\TextAt{13}{380.537}{494.394}{x}%
\NativeArrowAt{391.157}{494.394}{11.955}{11.955}{\mapsto}%
\TextAt{20}{407.081}{503.361}{\Character{80}}%
\TextAt{27}{419.7}{490.907}{p}%
\TextAt{3}{423.689}{490.907}{\Character{106}}%
\NativeMathAt{426.04}{490.907}{7.97}{\infty}%
\TextAt{13}{437}{494.394}{x}%
\TextAt{27}{443.652}{492.601}{p}%
\TextAt{8}{452.416}{494.394}{on}%
\TextAt{20}{469.055}{503.361}{\Character{81}}%
\TextAt{27}{480.346}{490.907}{p}%
\TextAt{3}{484.335}{490.907}{\Character{106}}%
\NativeMathAt{486.686}{490.907}{7.97}{\infty}%
\TextAt{18}{497.646}{494.394}{R}%
\TextAt{8}{506.28}{494.394}{,}%
\TextAt{8}{513.904}{494.394}{and}%
\TextAt{8}{79.226}{477.348}{hence}%
\TextAt{8}{111.732}{477.348}{the}%
\TextAt{8}{131.889}{477.348}{trace-zero}%
\TextAt{8}{185.866}{477.348}{space}%
\TextAt{13}{217.141}{477.348}{H}%
\TextAt{9}{231.084}{477.348}{\Character{61}}%
\TextAt{16}{243.51}{477.348}{\Character{102}}%
\TextAt{13}{249.488}{477.348}{x}%
\NativeMathAt{259.46}{477.348}{11.955}{\in}%
\TextAt{9}{270.751}{477.348}{\Character{91}}%
\TextAt{20}{274.003}{486.315}{\Character{81}}%
\TextAt{21}{285.294}{473.861}{\Character{28}}%
\TextAt{18}{292.463}{477.348}{R}%
\TextAt{9}{301.097}{477.348}{\Character{93}}%
\TextAt{17}{304.349}{481.687}{\Character{43}}%
\TextAt{16}{314.755}{477.348}{\Character{106}}%
\TextAt{20}{321.396}{486.315}{\Character{80}}%
\TextAt{21}{334.016}{473.861}{\Character{28}}%
\TextAt{13}{341.185}{477.348}{x}%
\TextAt{21}{347.837}{475.555}{\Character{28}}%
\TextAt{9}{356.334}{477.348}{\Character{61}}%
\TextAt{9}{368.744}{477.348}{0}%
\TextAt{16}{374.612}{477.348}{\Character{103}}%
\TextAt{8}{384.487}{477.348}{into}%
\TextAt{8}{408.23}{477.348}{the}%
\TextAt{8}{428.387}{477.348}{trace-zero}%
\TextAt{8}{482.364}{477.348}{space}%
\TextAt{13}{230.646}{450.009}{H}%
\TextAt{9}{244.589}{450.009}{\Character{61}}%
\TextAt{16}{257.014}{450.009}{\Character{102}}%
\TextAt{13}{262.992}{450.009}{x}%
\NativeMathAt{272.965}{450.009}{11.955}{\in}%
\TextAt{20}{284.256}{461.367}{\Character{89}}%
\TextAt{27}{284.489}{435.464}{p}%
\TextAt{3}{288.478}{435.464}{\Character{106}}%
\NativeMathAt{290.829}{435.464}{7.97}{\infty}%
\TextAt{18}{301.525}{450.009}{R}%
\TextAt{16}{313.48}{450.009}{\Character{106}}%
\TextAt{20}{320.121}{461.367}{\Character{88}}%
\TextAt{27}{321.351}{435.464}{p}%
\TextAt{3}{325.34}{435.464}{\Character{106}}%
\NativeMathAt{327.691}{435.464}{7.97}{\infty}%
\TextAt{13}{339.383}{450.009}{x}%
\TextAt{27}{346.035}{448.216}{p}%
\TextAt{9}{353.842}{450.009}{\Character{61}}%
\TextAt{9}{366.252}{450.009}{0}%
\TextAt{16}{372.121}{450.009}{\Character{103}}%
\TextAt{13}{378.098}{450.009}{\Character{58}}%
\TextAt{8}{79.226}{412.57}{In}%
\TextAt{8}{93.664}{412.57}{this}%
\TextAt{8}{116.303}{412.57}{way}%
\TextAt{8}{139.862}{412.57}{we}%
\TextAt{8}{156.906}{412.57}{have}%
\TextAt{8}{183.718}{412.57}{translated}%
\TextAt{8}{238.887}{412.57}{all}%
\TextAt{8}{254.962}{412.57}{necessary}%
\TextAt{8}{306.604}{412.57}{invariants}%
\TextAt{8}{360.159}{412.57}{of}%
\TextAt{8}{373.306}{412.57}{the}%
\TextAt{8}{393.279}{412.57}{Minkowski}%
\TextAt{8}{451.365}{412.57}{space}%
\TextAt{13}{482.456}{412.57}{K}%
\TextAt{25}{492.415}{410.777}{R}%
\TextAt{8}{502.384}{412.57}{to}%
\TextAt{8}{516.51}{412.57}{the}%
\TextAt{8}{79.226}{395.524}{product}%
\TextAt{20}{123.126}{404.49}{\Character{81}}%
\TextAt{27}{134.417}{392.037}{p}%
\TextAt{3}{138.405}{392.037}{\Character{106}}%
\NativeMathAt{140.756}{392.037}{7.97}{\infty}%
\TextAt{13}{151.717}{395.524}{K}%
\TextAt{27}{161.676}{393.731}{p}%
\TextAt{8}{166.163}{395.524}{.}%
\TextAt{8}{97.159}{375.489}{To}%
\TextAt{8}{114.386}{375.489}{a}%
\TextAt{8}{124.142}{375.489}{given}%
\TextAt{8}{154.711}{375.489}{replete}%
\TextAt{8}{193.075}{375.489}{ideal}%
\TextAt{14}{230.453}{348.15}{a}%
\TextAt{9}{239.757}{348.15}{\Character{61}}%
\TextAt{14}{252.182}{348.15}{a}%
\TextAt{21}{258.165}{346.357}{f}%
\TextAt{16}{266.267}{348.15}{\Character{1}}%
\TextAt{14}{272.245}{348.15}{a}%
\NativeMathAt{278.228}{346.357}{7.97}{\infty}%
\TextAt{9}{290.515}{348.15}{\Character{61}}%
\TextAt{20}{302.94}{359.507}{\Character{89}}%
\TextAt{27}{303.243}{333.675}{p}%
\TextAt{25}{307.232}{333.675}{\Character{45}}%
\NativeMathAt{309.446}{333.675}{7.97}{\infty}%
\TextAt{14}{320.209}{348.15}{p}%
\TextAt{21}{326.192}{353.086}{\Character{23}}%
\TextAt{28}{330.355}{352.09}{p}%
\TextAt{16}{337.592}{348.15}{\Character{2}}%
\TextAt{20}{349.548}{359.507}{\Character{89}}%
\TextAt{27}{349.781}{333.604}{p}%
\TextAt{3}{353.77}{333.604}{\Character{106}}%
\NativeMathAt{356.121}{333.604}{7.97}{\infty}%
\TextAt{14}{366.816}{348.15}{p}%
\TextAt{21}{372.799}{353.086}{\Character{23}}%
\TextAt{28}{376.962}{352.09}{p}%
\TextAt{8}{79.226}{310.711}{we}%
\TextAt{8}{96.444}{310.711}{now}%
\TextAt{8}{120.823}{310.711}{associate}%
\TextAt{8}{170.045}{310.711}{the}%
\TextAt{8}{190.191}{310.711}{following}%
\TextAt{8}{239.616}{310.711}{complete}%
\TextAt{8}{289.029}{310.711}{lattice}%
\TextAt{13}{324.79}{310.711}{j}%
\TextAt{14}{330.305}{310.711}{a}%
\TextAt{8}{340.176}{310.711}{in}%
\TextAt{13}{353.819}{310.711}{K}%
\TextAt{25}{363.778}{308.917}{R}%
\TextAt{8}{370.033}{310.711}{.}%
\TextAt{8}{378.478}{310.711}{The}%
\TextAt{8}{402.522}{310.711}{fractional}%
\TextAt{8}{454.876}{310.711}{ideal}%
\TextAt{14}{482.829}{310.711}{a}%
\TextAt{21}{488.813}{308.917}{f}%
\NativeMathAt{497.578}{310.711}{11.955}{\subseteq}%
\TextAt{13}{510.198}{310.711}{K}%
\TextAt{8}{524.903}{310.711}{is}%
\TextAt{8}{79.226}{293.665}{mapped}%
\TextAt{8}{123.338}{293.665}{by}%
\TextAt{8}{139.164}{293.665}{the}%
\TextAt{8}{158.887}{293.665}{embedding}%
\TextAt{13}{217.632}{293.665}{j}%
\TextAt{9}{226.469}{293.665}{\Character{58}}%
\TextAt{13}{233.041}{293.665}{K}%
\NativeMathAt{247.18}{293.665}{11.955}{\rightarrow}%
\TextAt{13}{262.456}{293.665}{K}%
\TextAt{25}{272.415}{291.871}{R}%
\TextAt{8}{282.133}{293.665}{onto}%
\TextAt{8}{308.05}{293.665}{a}%
\TextAt{8}{317.372}{293.665}{complete}%
\TextAt{8}{366.361}{293.665}{lattice}%
\TextAt{13}{401.698}{293.665}{j}%
\TextAt{14}{407.214}{293.665}{a}%
\TextAt{21}{413.197}{291.871}{f}%
\TextAt{8}{422.106}{293.665}{of}%
\TextAt{13}{435.003}{293.665}{K}%
\TextAt{25}{444.962}{291.871}{R}%
\TextAt{9}{454.537}{293.665}{\Character{61}}%
\TextAt{20}{466.963}{302.631}{\Character{81}}%
\TextAt{27}{478.254}{290.178}{p}%
\TextAt{3}{482.242}{290.178}{\Character{106}}%
\NativeMathAt{484.593}{290.178}{7.97}{\infty}%
\TextAt{13}{495.554}{293.665}{K}%
\TextAt{27}{505.513}{291.871}{p}%
\TextAt{8}{510}{293.665}{.}%
\TextAt{8}{518.304}{293.665}{By}%
\TextAt{8}{79.226}{276.619}{componentwise}%
\TextAt{8}{160.258}{276.619}{multiplication,}%
\TextAt{14}{238.636}{276.619}{a}%
\NativeMathAt{244.619}{274.825}{7.97}{\infty}%
\TextAt{9}{256.907}{276.619}{\Character{61}}%
\TextAt{20}{269.332}{285.585}{\Character{81}}%
\TextAt{27}{280.623}{273.132}{p}%
\TextAt{3}{284.612}{273.132}{\Character{106}}%
\NativeMathAt{286.963}{273.132}{7.97}{\infty}%
\TextAt{14}{297.923}{276.619}{p}%
\TextAt{21}{303.906}{280.957}{\Character{23}}%
\TextAt{28}{308.069}{279.961}{p}%
\TextAt{9}{315.971}{276.619}{\Character{61}}%
\TextAt{9}{328.381}{276.619}{\Character{40}}%
\TextAt{13}{332.949}{276.619}{\Character{58}}%
\TextAt{13}{338.173}{276.619}{\Character{58}}%
\TextAt{13}{343.41}{276.619}{\Character{58}}%
\TextAt{13}{348.634}{276.619}{\Character{59}}%
\TextAt{13}{353.871}{276.619}{e}%
\TextAt{21}{359.351}{280.957}{\Character{23}}%
\TextAt{28}{363.513}{279.961}{p}%
\TextAt{13}{368.094}{276.619}{\Character{59}}%
\TextAt{13}{373.318}{276.619}{\Character{58}}%
\TextAt{13}{378.555}{276.619}{\Character{58}}%
\TextAt{13}{383.791}{276.619}{\Character{58}}%
\TextAt{9}{387.078}{276.619}{\Character{41}}%
\TextAt{27}{391.631}{274.759}{p}%
\TextAt{3}{395.619}{274.759}{\Character{106}}%
\NativeMathAt{397.97}{274.759}{7.97}{\infty}%
\TextAt{8}{410.836}{276.619}{yields}%
\TextAt{8}{443.736}{276.619}{an}%
\TextAt{8}{459.995}{276.619}{isomorphism}%
\TextAt{14}{192.621}{249.279}{a}%
\NativeMathAt{198.604}{247.486}{7.97}{\infty}%
\TextAt{9}{210.892}{249.279}{\Character{58}}%
\TextAt{13}{217.464}{249.279}{K}%
\TextAt{25}{227.423}{247.486}{R}%
\NativeArrowAt{236.999}{249.279}{11.955}{19.26}{\longrightarrow}%
\TextAt{13}{259.581}{249.279}{K}%
\TextAt{25}{269.54}{247.486}{R}%
\TextAt{13}{275.794}{249.279}{\Character{59}}%
\TextAt{9}{304.948}{249.279}{\Character{40}}%
\TextAt{13}{309.501}{249.279}{x}%
\TextAt{27}{316.153}{247.486}{p}%
\TextAt{9}{320.64}{249.279}{\Character{41}}%
\TextAt{27}{325.192}{247.42}{p}%
\TextAt{3}{329.181}{247.42}{\Character{106}}%
\NativeMathAt{331.532}{247.42}{7.97}{\infty}%
\NativeArrowAt{343.82}{249.279}{11.955}{19.26}{\longmapsto}%
\TextAt{9}{366.402}{249.279}{\Character{40}}%
\TextAt{13}{370.955}{249.279}{e}%
\TextAt{21}{376.38}{254.216}{\Character{23}}%
\TextAt{28}{380.543}{253.219}{p}%
\TextAt{13}{385.124}{249.279}{x}%
\TextAt{27}{391.776}{247.486}{p}%
\TextAt{9}{396.263}{249.279}{\Character{41}}%
\TextAt{27}{400.815}{247.42}{p}%
\TextAt{3}{404.804}{247.42}{\Character{106}}%
\NativeMathAt{407.155}{247.42}{7.97}{\infty}%
\TextAt{13}{416.123}{249.279}{\Character{59}}%
\TextAt{8}{79.226}{221.94}{with}%
\TextAt{8}{105.886}{221.94}{determinant}%
\TextAt{8}{198.277}{204.894}{det}%
\TextAt{9}{214.536}{204.894}{\Character{40}}%
\TextAt{14}{219.088}{204.894}{a}%
\NativeMathAt{225.071}{203.101}{7.97}{\infty}%
\TextAt{9}{234.037}{204.894}{\Character{41}}%
\TextAt{9}{241.892}{204.894}{\Character{61}}%
\TextAt{20}{254.336}{216.252}{\Character{89}}%
\TextAt{27}{254.569}{190.349}{p}%
\TextAt{3}{258.558}{190.349}{\Character{106}}%
\NativeMathAt{260.909}{190.349}{7.97}{\infty}%
\TextAt{13}{271.605}{204.894}{e}%
\TextAt{21}{277.03}{209.831}{\Character{23}}%
\TextAt{28}{281.193}{208.834}{p}%
\TextAt{21}{285.276}{209.831}{f}%
\TextAt{28}{289.344}{208.834}{p}%
\TextAt{9}{297.245}{204.894}{\Character{61}}%
\TextAt{20}{309.671}{216.252}{\Character{89}}%
\TextAt{27}{309.904}{190.349}{p}%
\TextAt{3}{313.893}{190.349}{\Character{106}}%
\NativeMathAt{316.244}{190.349}{7.97}{\infty}%
\TextAt{14}{326.939}{204.894}{N}%
\TextAt{9}{336.888}{204.894}{\Character{40}}%
\TextAt{14}{341.44}{204.894}{p}%
\TextAt{9}{347.423}{204.894}{\Character{41}}%
\TextAt{21}{351.975}{209.831}{\Character{23}}%
\TextAt{28}{356.138}{208.834}{p}%
\TextAt{9}{364.04}{204.894}{\Character{61}}%
\TextAt{14}{376.465}{204.894}{N}%
\TextAt{9}{386.413}{204.894}{\Character{40}}%
\TextAt{14}{390.966}{204.894}{a}%
\NativeMathAt{396.949}{203.101}{7.97}{\infty}%
\TextAt{9}{405.915}{204.894}{\Character{41}}%
\TextAt{13}{410.468}{204.894}{\Character{58}}%
\TextAt{8}{517.802}{204.894}{(*)}%
\TextAt{8}{79.226}{171.985}{The}%
\TextAt{8}{103.28}{171.985}{image}%
\TextAt{8}{137.101}{171.985}{of}%
\TextAt{8}{150.431}{171.985}{the}%
\TextAt{8}{170.588}{171.985}{lattice}%
\TextAt{13}{206.357}{171.985}{j}%
\TextAt{14}{211.873}{171.985}{a}%
\TextAt{21}{217.856}{170.191}{f}%
\TextAt{8}{227.198}{171.985}{under}%
\TextAt{8}{260.362}{171.985}{this}%
\TextAt{8}{283.184}{171.985}{map}%
\TextAt{8}{309.199}{171.985}{is}%
\TextAt{8}{320.963}{171.985}{a}%
\TextAt{8}{330.718}{171.985}{complete}%
\TextAt{8}{380.141}{171.985}{lattice}%
\TextAt{13}{273.177}{144.646}{j}%
\TextAt{14}{278.693}{144.646}{a}%
\TextAt{9}{287.997}{144.646}{\Character{58}\Character{61}}%
\TextAt{14}{303.674}{144.646}{a}%
\NativeMathAt{309.657}{142.852}{7.97}{\infty}%
\TextAt{13}{318.623}{144.646}{j}%
\TextAt{14}{324.139}{144.646}{a}%
\TextAt{21}{330.122}{142.852}{f}%
\TextAt{13}{335.567}{144.646}{\Character{58}}%
\TextAt{8}{79.226}{117.306}{Let}%
\TextAt{8}{99.409}{117.306}{vol}%
\TextAt{9}{114.7}{117.306}{\Character{40}}%
\TextAt{14}{119.252}{117.306}{a}%
\TextAt{9}{125.235}{117.306}{\Character{41}}%
\TextAt{8}{132.898}{117.306}{denote}%
\TextAt{8}{169.83}{117.306}{the}%
\TextAt{8}{189.2}{117.306}{volume}%
\TextAt{8}{228.738}{117.306}{of}%
\TextAt{8}{241.281}{117.306}{a}%
\TextAt{12}{250.25}{117.306}{fundamental}%
\TextAt{12}{326.863}{117.306}{mesh}%
\TextAt{8}{360.101}{117.306}{of}%
\TextAt{13}{372.644}{117.306}{j}%
\TextAt{14}{378.159}{117.306}{a}%
\TextAt{8}{387.253}{117.306}{with}%
\TextAt{8}{413.127}{117.306}{respect}%
\TextAt{8}{452.39}{117.306}{to}%
\TextAt{8}{465.914}{117.306}{the}%
\TextAt{8}{485.284}{117.306}{canonical}%
\TextAt{8}{79.226}{100.26}{measure.}%
\TextAt{8}{129.362}{100.26}{By}%
\TextAt{8}{147.725}{100.26}{(*),}%
\TextAt{8}{169.842}{100.26}{we}%
\TextAt{8}{187.07}{100.26}{then}%
\TextAt{8}{213.73}{100.26}{have}%
\TextAt{8}{245.401}{72.921}{vol}%
\TextAt{9}{260.691}{72.921}{\Character{40}}%
\TextAt{14}{265.244}{72.921}{a}%
\TextAt{9}{271.227}{72.921}{\Character{41}}%
\TextAt{9}{279.082}{72.921}{\Character{61}}%
\TextAt{14}{291.525}{72.921}{N}%
\TextAt{9}{301.473}{72.921}{\Character{40}}%
\TextAt{14}{306.026}{72.921}{a}%
\NativeMathAt{312.009}{71.128}{7.97}{\infty}%
\TextAt{9}{320.975}{72.921}{\Character{41}}%
\TextAt{8}{327.52}{72.921}{vol}%
\TextAt{9}{342.811}{72.921}{\Character{40}}%
\TextAt{14}{347.363}{72.921}{a}%
\TextAt{21}{353.346}{71.128}{f}%
\TextAt{9}{358.791}{72.921}{\Character{41}}%
\TextAt{13}{363.343}{72.921}{\Character{58}}%
\TextAt{8}{297.211}{43.033}{193}%
\Destination{522.36}{221.945}{nav.AMS.25}%
\Destination{78.23}{809.815}{nav.page.193}%
\Destination{79.36}{186.955}{nav.print.212}%
}

% Page 206 (body page 194)
\Page{%
\TextAt{2}{79.226}{796.943}{III.}%
\TextAt{2}{98.317}{796.943}{Riemann–Roch}%
\TextAt{2}{174.375}{796.943}{Theory}%
\Vector{q 0 G 0 g 0.299 w 7.226 18.729 m 460.769 18.729 l S Q}%
\TextAt{12}{79.226}{761.124}{Definition}%
\TextAt{12}{142.015}{761.124}{(3.1).}%
\TextAt{8}{177.33}{761.124}{If}%
\TextAt{14}{189.022}{761.124}{a}%
\TextAt{8}{198.903}{761.124}{is}%
\TextAt{8}{210.667}{761.124}{a}%
\TextAt{8}{220.422}{761.124}{replete}%
\TextAt{8}{258.786}{761.124}{ideal}%
\TextAt{8}{286.749}{761.124}{of}%
\TextAt{13}{300.08}{761.124}{K}%
\TextAt{8}{310.897}{761.124}{,}%
\TextAt{8}{318.046}{761.124}{then}%
\TextAt{8}{344.706}{761.124}{the}%
\TextAt{8}{364.863}{761.124}{real}%
\TextAt{8}{387.625}{761.124}{number}%
\TextAt{13}{257.599}{732.999}{\Character{31}}%
\TextAt{9}{264.933}{732.999}{\Character{40}}%
\TextAt{14}{269.485}{732.999}{a}%
\TextAt{9}{275.468}{732.999}{\Character{41}}%
\TextAt{9}{283.323}{732.999}{\Character{61}}%
\TextAt{16}{295.767}{732.999}{\Character{0}}%
\TextAt{8}{307.058}{732.999}{log}%
\TextAt{8}{324.011}{732.999}{vol}%
\TextAt{9}{339.309}{732.999}{\Character{40}}%
\TextAt{14}{343.861}{732.999}{a}%
\TextAt{9}{349.844}{732.999}{\Character{41}}%
\TextAt{8}{79.226}{704.874}{will}%
\TextAt{8}{101.331}{704.874}{be}%
\TextAt{8}{117.256}{704.874}{called}%
\TextAt{8}{150.419}{704.874}{the}%
\TextAt{12}{170.576}{704.874}{Euler–Minkowski}%
\TextAt{12}{276.391}{704.874}{characteristic}%
\TextAt{8}{358.798}{704.874}{of}%
\TextAt{14}{372.128}{704.874}{a}%
\TextAt{8}{378.111}{704.874}{.}%
\TextAt{8}{97.159}{679.112}{The}%
\TextAt{8}{121.213}{679.112}{reason}%
\TextAt{8}{157.7}{679.112}{for}%
\TextAt{8}{175.585}{679.112}{this}%
\TextAt{8}{198.407}{679.112}{terminology}%
\TextAt{8}{263.133}{679.112}{will}%
\TextAt{8}{285.238}{679.112}{become}%
\TextAt{8}{327.176}{679.112}{clear}%
\TextAt{8}{355.139}{679.112}{in}%
\TextAt{8}{368.792}{679.112}{§8.}%
\TextAt{12}{79.226}{653.351}{Proposition}%
\TextAt{12}{152.018}{653.351}{(3.2).}%
\TextAt{8}{187.404}{653.351}{The}%
\TextAt{8}{211.528}{653.351}{Euler–Minkowski}%
\TextAt{8}{303.198}{653.351}{characteristic}%
\TextAt{13}{375.203}{653.351}{\Character{31}}%
\TextAt{9}{382.536}{653.351}{\Character{40}}%
\TextAt{14}{387.089}{653.351}{a}%
\TextAt{9}{393.072}{653.351}{\Character{41}}%
\TextAt{8}{401.592}{653.351}{only}%
\TextAt{8}{427.354}{653.351}{depends}%
\TextAt{8}{472.674}{653.351}{on}%
\TextAt{8}{489.003}{653.351}{the}%
\TextAt{8}{509.23}{653.351}{class}%
\TextAt{8}{79.226}{636.305}{of}%
\TextAt{14}{92.556}{636.305}{a}%
\TextAt{8}{102.437}{636.305}{in}%
\TextAt{8}{116.089}{636.305}{Pic}%
\TextAt{9}{132.504}{636.305}{\Character{40}}%
\Vector{q 0 G 0 g 0.478 w 65.056 -130.339 m 70.9 -130.339 l S Q}%
\TextAt{14}{137.056}{636.305}{o}%
\TextAt{9}{142.9}{636.305}{\Character{41}}%
\TextAt{9}{150.755}{636.305}{\Character{61}}%
\TextAt{13}{163.199}{636.305}{J}%
\TextAt{9}{170.776}{636.305}{\Character{40}}%
\Vector{q 0 G 0 g 0.478 w 103.329 -130.339 m 109.173 -130.339 l S Q}%
\TextAt{14}{175.329}{636.305}{o}%
\TextAt{9}{181.173}{636.305}{\Character{41}\Character{47}}%
\TextAt{13}{191.578}{636.305}{P}%
\TextAt{9}{200.749}{636.305}{\Character{40}}%
\Vector{q 0 G 0 g 0.478 w 133.302 -130.339 m 139.146 -130.339 l S Q}%
\TextAt{14}{205.302}{636.305}{o}%
\TextAt{9}{211.146}{636.305}{\Character{41}}%
\TextAt{8}{215.698}{636.305}{.}%
\TextAt{19}{79.226}{601.827}{Proof.}%
\TextAt{8}{116}{601.827}{Let}%
\TextAt{9}{136.97}{601.827}{\Character{91}}%
\TextAt{13}{140.221}{601.827}{a}%
\TextAt{9}{146.366}{601.827}{\Character{93}}%
\TextAt{9}{152.917}{601.827}{\Character{61}}%
\TextAt{9}{165.339}{601.827}{\Character{91}}%
\TextAt{13}{168.616}{601.827}{a}%
\TextAt{9}{174.761}{601.827}{\Character{93}}%
\TextAt{21}{178.012}{600.033}{f}%
\TextAt{16}{186.114}{601.827}{\Character{1}}%
\TextAt{9}{192.092}{601.827}{\Character{91}}%
\TextAt{13}{195.343}{601.827}{a}%
\TextAt{9}{201.488}{601.827}{\Character{93}}%
\NativeMathAt{204.74}{600.033}{7.97}{\infty}%
\TextAt{9}{217.027}{601.827}{\Character{61}}%
\TextAt{9}{229.437}{601.827}{\Character{40}}%
\TextAt{13}{234.005}{601.827}{a}%
\TextAt{9}{240.15}{601.827}{\Character{41}}%
\TextAt{16}{247.359}{601.827}{\Character{2}}%
\TextAt{9}{259.314}{601.827}{\Character{91}}%
\TextAt{13}{262.566}{601.827}{a}%
\TextAt{9}{268.711}{601.827}{\Character{93}}%
\NativeMathAt{271.962}{600.033}{7.97}{\infty}%
\TextAt{8}{284.826}{601.827}{be}%
\TextAt{8}{300.751}{601.827}{a}%
\TextAt{8}{310.506}{601.827}{replete}%
\TextAt{8}{348.87}{601.827}{principal}%
\TextAt{8}{397.647}{601.827}{ideal.}%
\TextAt{8}{430.161}{601.827}{Then}%
\TextAt{8}{460.719}{601.827}{one}%
\TextAt{8}{482.178}{601.827}{has}%
\TextAt{9}{252.809}{573.702}{\Character{91}}%
\TextAt{13}{256.06}{573.702}{a}%
\TextAt{9}{262.205}{573.702}{\Character{93}}%
\TextAt{14}{265.457}{573.702}{a}%
\TextAt{9}{274.761}{573.702}{\Character{61}}%
\TextAt{13}{287.186}{573.702}{a}%
\TextAt{14}{293.331}{573.702}{a}%
\TextAt{21}{299.314}{571.909}{f}%
\TextAt{16}{307.416}{573.702}{\Character{2}}%
\TextAt{9}{319.371}{573.702}{\Character{91}}%
\TextAt{13}{322.623}{573.702}{a}%
\TextAt{9}{328.768}{573.702}{\Character{93}}%
\NativeMathAt{332.019}{571.909}{7.97}{\infty}%
\TextAt{14}{340.986}{573.702}{a}%
\NativeMathAt{346.969}{571.909}{7.97}{\infty}%
\TextAt{13}{355.935}{573.702}{\Character{58}}%
\TextAt{8}{79.226}{545.577}{The}%
\TextAt{8}{103.958}{545.577}{lattice}%
\TextAt{13}{140.405}{545.577}{j}%
\TextAt{9}{145.921}{545.577}{\Character{40}}%
\TextAt{13}{150.473}{545.577}{a}%
\TextAt{14}{156.618}{545.577}{a}%
\TextAt{21}{162.601}{543.784}{f}%
\TextAt{9}{168.046}{545.577}{\Character{41}}%
\TextAt{8}{177.174}{545.577}{is}%
\TextAt{8}{189.615}{545.577}{the}%
\TextAt{8}{210.449}{545.577}{image}%
\TextAt{8}{244.948}{545.577}{of}%
\TextAt{8}{258.956}{545.577}{the}%
\TextAt{8}{279.79}{545.577}{lattice}%
\TextAt{13}{316.238}{545.577}{j}%
\TextAt{14}{321.753}{545.577}{a}%
\TextAt{21}{327.737}{543.784}{f}%
\TextAt{8}{337.757}{545.577}{under}%
\TextAt{8}{371.598}{545.577}{the}%
\TextAt{8}{392.432}{545.577}{linear}%
\TextAt{8}{425.628}{545.577}{map}%
\TextAt{13}{452.32}{545.577}{j}%
\TextAt{13}{457.831}{545.577}{a}%
\TextAt{9}{468.456}{545.577}{\Character{58}}%
\TextAt{13}{476.184}{545.577}{K}%
\TextAt{25}{486.143}{543.784}{R}%
\NativeMathAt{496.873}{545.577}{11.955}{\rightarrow}%
\TextAt{13}{513.304}{545.577}{K}%
\TextAt{25}{523.263}{543.784}{R}%
\TextAt{8}{529.518}{545.577}{,}%
\TextAt{9}{79.226}{528.531}{\Character{40}}%
\TextAt{13}{83.778}{528.531}{x}%
\TextAt{27}{90.431}{526.738}{p}%
\TextAt{9}{94.917}{528.531}{\Character{41}}%
\TextAt{27}{99.47}{526.672}{p}%
\TextAt{3}{103.458}{526.672}{\Character{106}}%
\NativeMathAt{105.809}{526.672}{7.97}{\infty}%
\NativeArrowAt{118.646}{528.531}{11.955}{11.955}{\mapsto}%
\TextAt{9}{134.47}{528.531}{\Character{40}}%
\TextAt{13}{139.023}{528.531}{ax}%
\TextAt{27}{151.82}{526.738}{p}%
\TextAt{9}{156.306}{528.531}{\Character{41}}%
\TextAt{27}{160.859}{526.672}{p}%
\TextAt{3}{164.847}{526.672}{\Character{106}}%
\NativeMathAt{167.198}{526.672}{7.97}{\infty}%
\TextAt{8}{176.166}{528.531}{.}%
\TextAt{8}{185.579}{528.531}{The}%
\TextAt{8}{209.955}{528.531}{absolute}%
\TextAt{8}{256.519}{528.531}{value}%
\TextAt{8}{287.087}{528.531}{of}%
\TextAt{8}{300.739}{528.531}{the}%
\TextAt{8}{321.217}{528.531}{determinant}%
\TextAt{8}{387.555}{528.531}{of}%
\TextAt{8}{401.206}{528.531}{this}%
\TextAt{8}{424.35}{528.531}{mapping}%
\TextAt{8}{472.804}{528.531}{is}%
\TextAt{8}{484.889}{528.531}{obviously}%
\TextAt{8}{79.226}{511.485}{given}%
\TextAt{8}{109.796}{511.485}{by}%
\TextAt{16}{180.319}{494.439}{\Character{106}}%
\TextAt{8}{185.633}{494.439}{det}%
\TextAt{9}{201.892}{494.439}{\Character{40}}%
\TextAt{13}{206.444}{494.439}{j}%
\TextAt{13}{211.955}{494.439}{a}%
\TextAt{9}{218.104}{494.439}{\Character{41}}%
\TextAt{16}{222.657}{494.439}{\Character{106}}%
\TextAt{9}{229.298}{494.439}{\Character{61}}%
\TextAt{20}{241.724}{505.797}{\Character{89}}%
\TextAt{27}{241.957}{479.894}{p}%
\TextAt{3}{245.946}{479.894}{\Character{106}}%
\NativeMathAt{248.297}{479.894}{7.97}{\infty}%
\TextAt{16}{258.992}{494.439}{\Character{106}}%
\TextAt{13}{262.313}{494.439}{a}%
\TextAt{16}{268.458}{494.439}{\Character{106}}%
\TextAt{27}{271.779}{492.646}{p}%
\TextAt{9}{279.587}{494.439}{\Character{61}}%
\TextAt{20}{292.012}{505.797}{\Character{89}}%
\TextAt{27}{292.246}{479.894}{p}%
\TextAt{3}{296.234}{479.894}{\Character{106}}%
\NativeMathAt{298.585}{479.894}{7.97}{\infty}%
\TextAt{14}{309.281}{494.439}{N}%
\TextAt{9}{319.229}{494.439}{\Character{40}}%
\TextAt{14}{323.781}{494.439}{p}%
\TextAt{9}{329.764}{494.439}{\Character{41}}%
\TextAt{3}{334.317}{499.376}{\Character{0}}%
\TextAt{21}{340.903}{499.376}{v}%
\TextAt{28}{345.027}{498.379}{p}%
\TextAt{17}{349.109}{499.376}{\Character{40}}%
\TextAt{21}{352.403}{499.376}{a}%
\TextAt{17}{356.901}{499.376}{\Character{41}}%
\TextAt{9}{364.013}{494.439}{\Character{61}}%
\TextAt{14}{376.438}{494.439}{N}%
\TextAt{9}{386.387}{494.439}{\Character{40}\Character{91}}%
\TextAt{13}{394.191}{494.439}{a}%
\TextAt{9}{400.335}{494.439}{\Character{93}}%
\NativeMathAt{403.587}{492.646}{7.97}{\infty}%
\TextAt{9}{412.554}{494.439}{\Character{41}}%
\TextAt{3}{417.106}{499.376}{\Character{0}}%
\TextAt{17}{423.692}{499.376}{1}%
\TextAt{13}{428.425}{494.439}{\Character{58}}%
\TextAt{8}{79.226}{461.193}{For}%
\TextAt{8}{100.196}{461.193}{the}%
\TextAt{8}{120.352}{461.193}{canonical}%
\TextAt{8}{171.735}{461.193}{measure,}%
\TextAt{8}{220.572}{461.193}{we}%
\TextAt{8}{237.8}{461.193}{therefore}%
\TextAt{8}{286.899}{461.193}{have}%
\TextAt{8}{230.614}{433.068}{vol}%
\TextAt{9}{245.904}{433.068}{\Character{40}}%
\TextAt{13}{250.457}{433.068}{a}%
\TextAt{14}{256.602}{433.068}{a}%
\TextAt{21}{262.585}{431.275}{f}%
\TextAt{9}{268.03}{433.068}{\Character{41}}%
\TextAt{9}{275.885}{433.068}{\Character{61}}%
\TextAt{14}{288.328}{433.068}{N}%
\TextAt{9}{298.276}{433.068}{\Character{40}\Character{91}}%
\TextAt{13}{306.08}{433.068}{a}%
\TextAt{9}{312.225}{433.068}{\Character{93}}%
\NativeMathAt{315.477}{431.275}{7.97}{\infty}%
\TextAt{9}{324.443}{433.068}{\Character{41}}%
\TextAt{3}{328.996}{438.005}{\Character{0}}%
\TextAt{17}{335.582}{438.005}{1}%
\TextAt{8}{342.307}{433.068}{vol}%
\TextAt{9}{357.598}{433.068}{\Character{40}}%
\TextAt{14}{362.15}{433.068}{a}%
\TextAt{21}{368.133}{431.275}{f}%
\TextAt{9}{373.578}{433.068}{\Character{41}}%
\TextAt{13}{378.13}{433.068}{\Character{58}}%
\TextAt{8}{79.226}{404.944}{Taken}%
\TextAt{8}{114.016}{404.944}{together}%
\TextAt{8}{160.198}{404.944}{with}%
\TextAt{8}{186.859}{404.944}{(*),}%
\TextAt{8}{208.976}{404.944}{this}%
\TextAt{8}{231.798}{404.944}{yields}%
\TextAt{8}{158.352}{376.819}{vol}%
\TextAt{9}{173.643}{376.819}{\Character{40}\Character{91}}%
\TextAt{13}{181.447}{376.819}{a}%
\TextAt{9}{187.592}{376.819}{\Character{93}}%
\TextAt{14}{190.843}{376.819}{a}%
\TextAt{9}{196.826}{376.819}{\Character{41}}%
\TextAt{9}{204.681}{376.819}{\Character{61}}%
\TextAt{14}{217.125}{376.819}{N}%
\TextAt{9}{227.073}{376.819}{\Character{40}\Character{91}}%
\TextAt{13}{234.877}{376.819}{a}%
\TextAt{9}{241.022}{376.819}{\Character{93}}%
\NativeMathAt{244.274}{375.026}{7.97}{\infty}%
\TextAt{14}{253.24}{376.819}{a}%
\NativeMathAt{259.223}{375.026}{7.97}{\infty}%
\TextAt{9}{268.19}{376.819}{\Character{41}}%
\TextAt{8}{274.735}{376.819}{vol}%
\TextAt{9}{290.025}{376.819}{\Character{40}}%
\TextAt{13}{294.578}{376.819}{a}%
\TextAt{14}{300.723}{376.819}{a}%
\TextAt{21}{306.706}{375.026}{f}%
\TextAt{9}{312.151}{376.819}{\Character{41}}%
\TextAt{9}{320.006}{376.819}{\Character{61}}%
\TextAt{14}{332.449}{376.819}{N}%
\TextAt{9}{342.397}{376.819}{\Character{40}}%
\TextAt{14}{346.95}{376.819}{a}%
\NativeMathAt{352.933}{375.026}{7.97}{\infty}%
\TextAt{9}{361.899}{376.819}{\Character{41}}%
\TextAt{8}{368.444}{376.819}{vol}%
\TextAt{9}{383.735}{376.819}{\Character{40}}%
\TextAt{14}{388.287}{376.819}{a}%
\TextAt{21}{394.27}{375.026}{f}%
\TextAt{9}{399.715}{376.819}{\Character{41}}%
\TextAt{9}{407.57}{376.819}{\Character{61}}%
\TextAt{8}{420.014}{376.819}{vol}%
\TextAt{9}{435.304}{376.819}{\Character{40}}%
\TextAt{14}{439.857}{376.819}{a}%
\TextAt{9}{445.84}{376.819}{\Character{41}}%
\TextAt{13}{450.392}{376.819}{\Character{59}}%
\TextAt{8}{79.226}{348.694}{so}%
\TextAt{8}{93.596}{348.694}{that}%
\TextAt{13}{118.965}{348.694}{\Character{31}}%
\TextAt{9}{126.299}{348.694}{\Character{40}\Character{91}}%
\TextAt{13}{134.103}{348.694}{a}%
\TextAt{9}{140.248}{348.694}{\Character{93}}%
\TextAt{14}{143.499}{348.694}{a}%
\TextAt{9}{149.482}{348.694}{\Character{41}}%
\TextAt{9}{157.337}{348.694}{\Character{61}}%
\TextAt{13}{169.781}{348.694}{\Character{31}}%
\TextAt{9}{177.115}{348.694}{\Character{40}}%
\TextAt{14}{181.667}{348.694}{a}%
\TextAt{9}{187.65}{348.694}{\Character{41}}%
\TextAt{8}{192.202}{348.694}{.}%
\Vector{q 0 G 0 g 0.399 w 452.334 -425.306 m 452.334 -417.236 l S Q}%
\Vector{q 0 G 0 g 0.399 w 452.534 -417.435 m 459.707 -417.435 l S Q}%
\Vector{q 0 G 0 g 0.399 w 452.534 -425.106 m 459.707 -425.106 l S Q}%
\Vector{q 0 G 0 g 0.399 w 459.906 -425.306 m 459.906 -417.236 l S Q}%
\TextAt{8}{97.159}{319.944}{The}%
\TextAt{8}{121.246}{319.944}{explicit}%
\TextAt{8}{162.573}{319.944}{evaluation}%
\TextAt{8}{218.88}{319.944}{of}%
\TextAt{8}{232.244}{319.944}{the}%
\TextAt{8}{252.434}{319.944}{Euler–Minkowski}%
\TextAt{8}{344.068}{319.944}{characteristic}%
\TextAt{8}{416.036}{319.944}{results}%
\TextAt{8}{453.262}{319.944}{from}%
\TextAt{8}{480.936}{319.944}{a}%
\TextAt{8}{490.725}{319.944}{result}%
\TextAt{8}{523.337}{319.944}{of}%
\TextAt{8}{79.226}{302.898}{Minkowski}%
\TextAt{8}{137.496}{302.898}{theory,}%
\TextAt{8}{176.517}{302.898}{viz.,}%
\TextAt{8}{201.551}{302.898}{proposition}%
\TextAt{8}{263.085}{302.898}{(5.2)}%
\TextAt{8}{291.06}{302.898}{of}%
\TextAt{8}{304.39}{302.898}{chap.}%
\TextAt{8}{336.581}{302.898}{I.}%
\TextAt{12}{79.226}{277.136}{Proposition}%
\TextAt{12}{151.937}{277.136}{(3.3).}%
\TextAt{8}{187.253}{277.136}{For}%
\TextAt{8}{208.222}{277.136}{every}%
\TextAt{8}{239.115}{277.136}{replete}%
\TextAt{8}{277.479}{277.136}{ideal}%
\TextAt{14}{305.442}{277.136}{a}%
\TextAt{8}{315.322}{277.136}{of}%
\TextAt{13}{328.652}{277.136}{K}%
\TextAt{8}{343.367}{277.136}{one}%
\TextAt{8}{364.827}{277.136}{has}%
\TextAt{8}{252.351}{249.011}{vol}%
\TextAt{9}{267.641}{249.011}{\Character{40}}%
\TextAt{14}{272.194}{249.011}{a}%
\TextAt{9}{278.177}{249.011}{\Character{41}}%
\TextAt{9}{286.032}{249.011}{\Character{61}}%
\TextAt{20}{298.475}{259.817}{\Character{112}}%
\Vector{q 0 G 0 g 0.478 w 238.43 -513.944 m 259.358 -513.944 l S Q}%
\TextAt{16}{310.43}{249.011}{\Character{106}}%
\TextAt{13}{313.751}{249.011}{d}%
\TextAt{21}{319.834}{247.218}{K}%
\TextAt{16}{328.037}{249.011}{\Character{106}}%
\TextAt{14}{331.358}{249.011}{N}%
\TextAt{9}{341.306}{249.011}{\Character{40}}%
\TextAt{14}{345.858}{249.011}{a}%
\TextAt{9}{351.841}{249.011}{\Character{41}}%
\TextAt{13}{356.393}{249.011}{\Character{58}}%
\TextAt{19}{79.226}{203.455}{Proof.}%
\TextAt{8}{116}{203.455}{Multiplying}%
\TextAt{8}{180.112}{203.455}{by}%
\TextAt{8}{196.751}{203.455}{a}%
\TextAt{8}{206.885}{203.455}{replete}%
\TextAt{8}{245.629}{203.455}{principal}%
\TextAt{8}{294.785}{203.455}{ideal}%
\TextAt{9}{323.127}{203.455}{\Character{91}}%
\TextAt{13}{326.379}{203.455}{a}%
\TextAt{9}{332.524}{203.455}{\Character{93}}%
\TextAt{8}{340.052}{203.455}{we}%
\TextAt{8}{357.658}{203.455}{may}%
\TextAt{8}{383.406}{203.455}{assume,}%
\TextAt{8}{427.577}{203.455}{as}%
\TextAt{8}{442.326}{203.455}{vol}%
\TextAt{9}{457.616}{203.455}{\Character{40}\Character{91}}%
\TextAt{13}{465.42}{203.455}{a}%
\TextAt{9}{471.565}{203.455}{\Character{93}}%
\TextAt{14}{474.817}{203.455}{a}%
\TextAt{9}{480.8}{203.455}{\Character{41}}%
\TextAt{9}{489.3}{203.455}{\Character{61}}%
\TextAt{8}{502.391}{203.455}{vol}%
\TextAt{9}{517.682}{203.455}{\Character{40}}%
\TextAt{14}{522.234}{203.455}{a}%
\TextAt{9}{528.217}{203.455}{\Character{41}}%
\TextAt{8}{79.226}{186.409}{and}%
\TextAt{14}{102.139}{186.409}{N}%
\TextAt{9}{112.087}{186.409}{\Character{40}\Character{91}}%
\TextAt{13}{119.891}{186.409}{a}%
\TextAt{9}{126.036}{186.409}{\Character{93}}%
\TextAt{14}{129.288}{186.409}{a}%
\TextAt{9}{135.271}{186.409}{\Character{41}}%
\TextAt{9}{143.389}{186.409}{\Character{61}}%
\TextAt{14}{156.083}{186.409}{N}%
\TextAt{9}{166.031}{186.409}{\Character{40}}%
\TextAt{14}{170.583}{186.409}{a}%
\TextAt{9}{176.566}{186.409}{\Character{41}}%
\TextAt{8}{181.119}{186.409}{,}%
\TextAt{8}{188.456}{186.409}{that}%
\TextAt{14}{213.976}{186.409}{a}%
\TextAt{21}{219.959}{184.615}{f}%
\TextAt{8}{229.452}{186.409}{is}%
\TextAt{8}{241.366}{186.409}{an}%
\TextAt{8}{257.776}{186.409}{integral}%
\TextAt{8}{300.534}{186.409}{ideal}%
\TextAt{8}{328.648}{186.409}{of}%
\TextAt{13}{342.129}{186.409}{K}%
\TextAt{8}{352.946}{186.409}{.}%
\TextAt{8}{361.846}{186.409}{By}%
\TextAt{8}{380.36}{186.409}{chap.}%
\TextAt{8}{413.003}{186.409}{I,}%
\TextAt{8}{424.523}{186.409}{(5.2)}%
\TextAt{8}{452.649}{186.409}{the}%
\TextAt{8}{472.956}{186.409}{volume}%
\TextAt{8}{513.431}{186.409}{of}%
\TextAt{8}{526.911}{186.409}{a}%
\TextAt{8}{79.226}{169.363}{fundamental}%
\TextAt{8}{146.868}{169.363}{mesh}%
\TextAt{8}{176.84}{169.363}{of}%
\TextAt{14}{190.17}{169.363}{a}%
\TextAt{21}{196.153}{167.569}{f}%
\TextAt{8}{205.495}{169.363}{is}%
\TextAt{8}{217.259}{169.363}{given}%
\TextAt{8}{247.829}{169.363}{by}%
\TextAt{8}{244.011}{141.238}{vol}%
\TextAt{9}{259.302}{141.238}{\Character{40}}%
\TextAt{14}{263.854}{141.238}{a}%
\TextAt{21}{269.837}{139.445}{f}%
\TextAt{9}{275.282}{141.238}{\Character{41}}%
\TextAt{9}{283.137}{141.238}{\Character{61}}%
\TextAt{20}{295.581}{152.043}{\Character{112}}%
\Vector{q 0 G 0 g 0.478 w 235.536 -621.717 m 256.463 -621.717 l S Q}%
\TextAt{16}{307.536}{141.238}{\Character{106}}%
\TextAt{13}{310.857}{141.238}{d}%
\TextAt{21}{316.939}{139.445}{K}%
\TextAt{16}{325.142}{141.238}{\Character{106}}%
\TextAt{9}{328.463}{141.238}{\Character{40}}%
\TextAt{14}{333.015}{141.238}{o}%
\TextAt{9}{342.18}{141.238}{\Character{58}}%
\TextAt{14}{348.753}{141.238}{a}%
\TextAt{21}{354.736}{139.445}{f}%
\TextAt{9}{360.181}{141.238}{\Character{41}}%
\TextAt{13}{364.733}{141.238}{\Character{58}}%
\TextAt{8}{79.226}{113.113}{Hence}%
\TextAt{8}{152.012}{96.067}{vol}%
\TextAt{9}{167.303}{96.067}{\Character{40}}%
\TextAt{14}{171.855}{96.067}{a}%
\TextAt{9}{177.839}{96.067}{\Character{41}}%
\TextAt{9}{185.694}{96.067}{\Character{61}}%
\TextAt{14}{198.137}{96.067}{N}%
\TextAt{9}{208.085}{96.067}{\Character{40}}%
\TextAt{14}{212.638}{96.067}{a}%
\NativeMathAt{218.621}{94.274}{7.97}{\infty}%
\TextAt{9}{227.587}{96.067}{\Character{41}}%
\TextAt{8}{234.132}{96.067}{vol}%
\TextAt{9}{249.423}{96.067}{\Character{40}}%
\TextAt{14}{253.975}{96.067}{a}%
\TextAt{21}{259.958}{94.274}{f}%
\TextAt{9}{265.403}{96.067}{\Character{41}}%
\TextAt{9}{273.258}{96.067}{\Character{61}}%
\TextAt{14}{285.702}{96.067}{N}%
\TextAt{9}{295.65}{96.067}{\Character{40}}%
\TextAt{14}{300.202}{96.067}{a}%
\NativeMathAt{306.185}{94.274}{7.97}{\infty}%
\TextAt{9}{315.152}{96.067}{\Character{41}}%
\TextAt{20}{319.704}{106.873}{\Character{112}}%
\Vector{q 0 G 0 g 0.478 w 259.659 -666.888 m 280.586 -666.888 l S Q}%
\TextAt{16}{331.659}{96.067}{\Character{106}}%
\TextAt{13}{334.98}{96.067}{d}%
\TextAt{21}{341.063}{94.274}{K}%
\TextAt{16}{349.265}{96.067}{\Character{106}}%
\TextAt{14}{352.586}{96.067}{N}%
\TextAt{9}{362.534}{96.067}{\Character{40}}%
\TextAt{14}{367.087}{96.067}{a}%
\TextAt{21}{373.07}{94.274}{f}%
\TextAt{9}{378.515}{96.067}{\Character{41}}%
\TextAt{9}{386.37}{96.067}{\Character{61}}%
\TextAt{20}{398.813}{106.873}{\Character{112}}%
\Vector{q 0 G 0 g 0.478 w 338.769 -666.888 m 359.696 -666.888 l S Q}%
\TextAt{16}{410.769}{96.067}{\Character{106}}%
\TextAt{13}{414.089}{96.067}{d}%
\TextAt{21}{420.172}{94.274}{K}%
\TextAt{16}{428.375}{96.067}{\Character{106}}%
\TextAt{14}{431.696}{96.067}{N}%
\TextAt{9}{441.644}{96.067}{\Character{40}}%
\TextAt{14}{446.196}{96.067}{a}%
\TextAt{9}{452.179}{96.067}{\Character{41}}%
\TextAt{13}{456.732}{96.067}{\Character{58}}%
\Vector{q 0 G 0 g 0.399 w 452.334 -701.079 m 452.334 -693.009 l S Q}%
\Vector{q 0 G 0 g 0.399 w 452.534 -693.208 m 459.707 -693.208 l S Q}%
\Vector{q 0 G 0 g 0.399 w 452.534 -700.879 m 459.707 -700.879 l S Q}%
\Vector{q 0 G 0 g 0.399 w 459.906 -701.079 m 459.906 -693.009 l S Q}%
\TextAt{8}{297.211}{43.033}{194}%
\Destination{78.23}{809.815}{nav.page.194}%
\Destination{79.36}{775.155}{nav.number.3.3.1}%
\Destination{79.36}{667.355}{nav.number.3.3.2}%
\Destination{79.36}{291.155}{nav.number.3.3.3}%
\Hotlink{373.971}{676.955}{6.858}{11.047}{nav.section.3.8}%
\Hotlink{336.081}{300.655}{5.22}{11.047}{nav.chapter.1}%
\Hotlink{412.503}{184.255}{5.22}{11.047}{nav.chapter.1}%
\Hotlink{262.585}{300.655}{25.078}{11.047}{nav.number.1.5.2}%
\Hotlink{424.023}{184.255}{25.078}{11.047}{nav.number.1.5.2}%
}

% Page 207 (body page 195)
\Page{%
\TextAt{2}{79.226}{796.943}{III.}%
\TextAt{2}{98.317}{796.943}{Riemann–Roch}%
\TextAt{2}{174.375}{796.943}{Theory}%
\Vector{q 0 G 0 g 0.299 w 7.226 18.729 m 460.769 18.729 l S Q}%
\TextAt{8}{97.159}{761.124}{In}%
\TextAt{8}{112.072}{761.124}{view}%
\TextAt{8}{139.347}{761.124}{of}%
\TextAt{8}{152.969}{761.124}{the}%
\TextAt{8}{173.417}{761.124}{commutative}%
\TextAt{8}{243.634}{761.124}{diagram}%
\TextAt{8}{289.464}{761.124}{in}%
\TextAt{8}{303.408}{761.124}{§1,}%
\TextAt{8}{322.459}{761.124}{p.192,}%
\TextAt{8}{357.303}{761.124}{we}%
\TextAt{8}{374.822}{761.124}{will}%
\TextAt{8}{397.219}{761.124}{now}%
\TextAt{8}{421.899}{761.124}{introduce}%
\TextAt{8}{474.22}{761.124}{the}%
\TextAt{12}{494.668}{761.124}{degree}%
\TextAt{8}{79.226}{744.078}{of}%
\TextAt{8}{92.556}{744.078}{the}%
\TextAt{8}{112.713}{744.078}{replete}%
\TextAt{8}{151.077}{744.078}{ideal}%
\TextAt{14}{179.04}{744.078}{a}%
\TextAt{8}{188.92}{744.078}{to}%
\TextAt{8}{203.231}{744.078}{be}%
\TextAt{8}{219.155}{744.078}{the}%
\TextAt{8}{239.311}{744.078}{real}%
\TextAt{8}{262.074}{744.078}{number}%
\TextAt{8}{216.812}{715.077}{deg}%
\TextAt{9}{234.374}{715.077}{\Character{40}}%
\TextAt{14}{238.926}{715.077}{a}%
\TextAt{9}{244.909}{715.077}{\Character{41}}%
\TextAt{9}{252.764}{715.077}{\Character{61}}%
\TextAt{16}{265.208}{715.077}{\Character{0}}%
\TextAt{8}{276.499}{715.077}{log}%
\TextAt{14}{293.459}{715.077}{N}%
\TextAt{9}{303.407}{715.077}{\Character{40}}%
\TextAt{14}{307.96}{715.077}{a}%
\TextAt{9}{313.943}{715.077}{\Character{41}}%
\TextAt{9}{321.798}{715.077}{\Character{61}}%
\TextAt{8}{334.241}{715.077}{deg}%
\TextAt{9}{351.804}{715.077}{\Character{40}}%
\TextAt{8}{356.356}{715.077}{div}%
\TextAt{9}{372.292}{715.077}{\Character{40}}%
\TextAt{14}{376.845}{715.077}{a}%
\TextAt{9}{382.828}{715.077}{\Character{41}\Character{41}}%
\TextAt{13}{391.932}{715.077}{\Character{58}}%
\TextAt{8}{79.226}{686.075}{Observing}%
\TextAt{8}{134.901}{686.075}{that}%
\TextAt{13}{254.791}{669.029}{\Character{31}}%
\TextAt{9}{262.124}{669.029}{\Character{40}}%
\TextAt{14}{266.677}{669.029}{o}%
\TextAt{9}{272.521}{669.029}{\Character{41}}%
\TextAt{9}{280.376}{669.029}{\Character{61}}%
\TextAt{16}{292.819}{669.029}{\Character{0}}%
\TextAt{8}{304.111}{669.029}{log}%
\TextAt{20}{321.071}{679.835}{\Character{112}}%
\Vector{q 0 G 0 g 0.478 w 261.026 -93.926 m 281.953 -93.926 l S Q}%
\TextAt{16}{333.026}{669.029}{\Character{106}}%
\TextAt{13}{336.347}{669.029}{d}%
\TextAt{21}{342.43}{667.236}{K}%
\TextAt{16}{350.632}{669.029}{\Character{106}}%
\TextAt{13}{353.953}{669.029}{\Character{59}}%
\TextAt{8}{79.226}{645.508}{we}%
\TextAt{8}{96.454}{645.508}{deduce}%
\TextAt{8}{135.463}{645.508}{from}%
\TextAt{8}{163.104}{645.508}{proposition}%
\TextAt{8}{224.637}{645.508}{(3.3)}%
\TextAt{8}{252.612}{645.508}{the}%
\TextAt{8}{272.768}{645.508}{first}%
\TextAt{8}{296.894}{645.508}{version}%
\TextAt{8}{336.633}{645.508}{of}%
\TextAt{8}{349.963}{645.508}{the}%
\TextAt{8}{370.119}{645.508}{Riemann–Roch}%
\TextAt{8}{451.749}{645.508}{theorem:}%
\TextAt{12}{79.226}{619.495}{Proposition}%
\TextAt{12}{151.937}{619.495}{(3.4).}%
\TextAt{8}{187.253}{619.495}{For}%
\TextAt{8}{208.222}{619.495}{every}%
\TextAt{8}{239.115}{619.495}{replete}%
\TextAt{8}{277.479}{619.495}{ideal}%
\TextAt{8}{305.442}{619.495}{of}%
\TextAt{13}{318.772}{619.495}{K}%
\TextAt{8}{333.487}{619.495}{we}%
\TextAt{8}{350.714}{619.495}{have}%
\TextAt{8}{377.709}{619.495}{the}%
\TextAt{8}{397.866}{619.495}{formula}%
\TextAt{13}{250.613}{590.494}{\Character{31}}%
\TextAt{9}{257.947}{590.494}{\Character{40}}%
\TextAt{14}{262.499}{590.494}{a}%
\TextAt{9}{268.482}{590.494}{\Character{41}}%
\TextAt{9}{276.337}{590.494}{\Character{61}}%
\TextAt{8}{288.781}{590.494}{deg}%
\TextAt{9}{306.343}{590.494}{\Character{40}}%
\TextAt{14}{310.895}{590.494}{a}%
\TextAt{9}{316.878}{590.494}{\Character{41}}%
\TextAt{9}{324.075}{590.494}{\Character{43}}%
\TextAt{13}{335.849}{590.494}{\Character{31}}%
\TextAt{9}{343.182}{590.494}{\Character{40}}%
\TextAt{14}{347.734}{590.494}{o}%
\TextAt{9}{353.579}{590.494}{\Character{41}}%
\TextAt{13}{358.131}{590.494}{\Character{58}}%
\TextAt{8}{97.159}{552.527}{In}%
\TextAt{8}{110.961}{552.527}{function}%
\TextAt{8}{155.989}{552.527}{theory}%
\TextAt{8}{191.92}{552.527}{there}%
\TextAt{8}{221.012}{552.527}{is}%
\TextAt{8}{231.957}{552.527}{the}%
\TextAt{8}{251.294}{552.527}{following}%
\TextAt{8}{299.909}{552.527}{relationship}%
\TextAt{8}{362.894}{552.527}{between}%
\TextAt{8}{407.265}{552.527}{the}%
\TextAt{8}{426.602}{552.527}{Euler–Poincaré}%
\TextAt{8}{507.078}{552.527}{char-}%
\TextAt{8}{79.226}{535.481}{acteristic}%
\TextAt{8}{129.366}{535.481}{and}%
\TextAt{8}{152.129}{535.481}{the}%
\TextAt{8}{172.285}{535.481}{genus}%
\TextAt{13}{204.54}{535.481}{g}%
\TextAt{8}{214.472}{535.481}{of}%
\TextAt{8}{227.802}{535.481}{the}%
\TextAt{8}{247.958}{535.481}{Riemann}%
\TextAt{8}{297.548}{535.481}{surface}%
\TextAt{13}{336.953}{535.481}{X}%
\TextAt{8}{351.505}{535.481}{in}%
\TextAt{8}{365.158}{535.481}{question:}%
\TextAt{13}{181.788}{506.48}{\Character{31}}%
\TextAt{9}{189.121}{506.48}{\Character{40}}%
\TextAt{14}{193.674}{506.48}{o}%
\TextAt{9}{199.518}{506.48}{\Character{41}}%
\TextAt{9}{207.373}{506.48}{\Character{61}}%
\TextAt{8}{219.817}{506.48}{dim}%
\TextAt{13}{241.32}{506.48}{H}%
\TextAt{17}{251.943}{511.416}{0}%
\TextAt{9}{256.675}{506.48}{\Character{40}}%
\TextAt{13}{261.227}{506.48}{X}%
\TextAt{13}{271.221}{506.48}{\Character{59}}%
\TextAt{14}{276.476}{506.48}{o}%
\TextAt{21}{282.32}{504.686}{X}%
\TextAt{9}{290.405}{506.48}{\Character{41}}%
\TextAt{16}{297.614}{506.48}{\Character{0}}%
\TextAt{8}{309.57}{506.48}{dim}%
\TextAt{13}{331.073}{506.48}{H}%
\TextAt{17}{341.696}{511.416}{1}%
\TextAt{9}{346.428}{506.48}{\Character{40}}%
\TextAt{13}{350.98}{506.48}{X}%
\TextAt{13}{360.975}{506.48}{\Character{59}}%
\TextAt{14}{366.229}{506.48}{o}%
\TextAt{21}{372.073}{504.686}{X}%
\TextAt{9}{380.158}{506.48}{\Character{41}}%
\TextAt{9}{388.013}{506.48}{\Character{61}}%
\TextAt{9}{400.434}{506.48}{1}%
\TextAt{16}{408.967}{506.48}{\Character{0}}%
\TextAt{13}{420.922}{506.48}{g}%
\TextAt{13}{426.935}{506.48}{\Character{58}}%
\TextAt{8}{79.226}{477.479}{There}%
\TextAt{8}{111.948}{477.479}{is}%
\TextAt{8}{122.625}{477.479}{no}%
\TextAt{8}{137.797}{477.479}{immediate}%
\TextAt{8}{193.94}{477.479}{analogue}%
\TextAt{8}{241.642}{477.479}{of}%
\TextAt{13}{253.884}{477.479}{H}%
\TextAt{17}{264.507}{481.817}{1}%
\TextAt{9}{269.239}{477.479}{\Character{40}}%
\TextAt{13}{273.792}{477.479}{X}%
\TextAt{13}{283.786}{477.479}{\Character{59}}%
\TextAt{14}{289.041}{477.479}{o}%
\TextAt{21}{294.885}{475.685}{X}%
\TextAt{9}{302.97}{477.479}{\Character{41}}%
\TextAt{8}{310.333}{477.479}{in}%
\TextAt{8}{322.898}{477.479}{arithmetic.}%
\TextAt{8}{383.671}{477.479}{However,}%
\TextAt{8}{433.204}{477.479}{there}%
\TextAt{8}{462.029}{477.479}{is}%
\TextAt{8}{472.706}{477.479}{an}%
\TextAt{8}{487.878}{477.479}{analogue}%
\TextAt{8}{79.226}{460.433}{of}%
\TextAt{13}{92.556}{460.433}{H}%
\TextAt{17}{103.179}{464.771}{0}%
\TextAt{9}{107.911}{460.433}{\Character{40}}%
\TextAt{13}{112.463}{460.433}{X}%
\TextAt{13}{122.458}{460.433}{\Character{59}}%
\TextAt{14}{127.712}{460.433}{o}%
\TextAt{21}{133.557}{458.639}{X}%
\TextAt{9}{141.642}{460.433}{\Character{41}}%
\TextAt{8}{146.194}{460.433}{.}%
\TextAt{8}{154.642}{460.433}{For}%
\TextAt{8}{175.612}{460.433}{each}%
\TextAt{8}{201.949}{460.433}{replete}%
\TextAt{8}{240.313}{460.433}{ideal}%
\TextAt{14}{268.276}{460.433}{a}%
\TextAt{9}{277.58}{460.433}{\Character{61}}%
\TextAt{20}{290.006}{469.399}{\Character{81}}%
\TextAt{27}{301.297}{456.946}{p}%
\TextAt{14}{307.776}{460.433}{p}%
\TextAt{21}{313.759}{464.771}{\Character{23}}%
\TextAt{28}{317.922}{463.775}{p}%
\TextAt{8}{326.4}{460.433}{of}%
\TextAt{8}{339.73}{460.433}{the}%
\TextAt{8}{359.886}{460.433}{number}%
\TextAt{8}{402.483}{460.433}{field}%
\TextAt{13}{427.84}{460.433}{K}%
\TextAt{8}{438.657}{460.433}{,}%
\TextAt{8}{445.807}{460.433}{we}%
\TextAt{8}{463.034}{460.433}{define}%
\TextAt{13}{203.448}{431.431}{H}%
\TextAt{17}{214.071}{436.368}{0}%
\TextAt{9}{218.803}{431.431}{\Character{40}}%
\TextAt{14}{223.356}{431.431}{a}%
\TextAt{9}{229.339}{431.431}{\Character{41}}%
\TextAt{9}{237.194}{431.431}{\Character{61}}%
\TextAt{16}{249.637}{431.431}{\Character{102}}%
\TextAt{13}{255.615}{431.431}{f}%
\NativeMathAt{265.982}{431.431}{11.955}{\in}%
\TextAt{13}{277.273}{431.431}{K}%
\TextAt{3}{288.091}{436.368}{\Character{3}}%
\TextAt{16}{296.144}{431.431}{\Character{106}}%
\TextAt{13}{302.786}{431.431}{v}%
\TextAt{27}{308.445}{429.638}{p}%
\TextAt{9}{312.932}{431.431}{\Character{40}}%
\TextAt{13}{317.484}{431.431}{f}%
\TextAt{9}{324.531}{431.431}{\Character{41}}%
\NativeMathAt{332.404}{431.431}{11.955}{\geq}%
\TextAt{13}{345.023}{431.431}{\Character{23}}%
\TextAt{27}{350.807}{429.638}{p}%
\TextAt{8}{359.191}{431.431}{for}%
\TextAt{8}{377.076}{431.431}{all}%
\TextAt{14}{393.335}{431.431}{p}%
\TextAt{16}{399.318}{431.431}{\Character{103}}%
\TextAt{13}{405.296}{431.431}{\Character{58}}%
\TextAt{8}{79.226}{402.43}{This}%
\TextAt{8}{104.956}{402.43}{is}%
\TextAt{8}{115.73}{402.43}{a}%
\TextAt{8}{124.496}{402.43}{finite}%
\TextAt{8}{153.418}{402.43}{set}%
\TextAt{8}{170.695}{402.43}{because}%
\TextAt{13}{213.007}{402.43}{j}%
\TextAt{13}{218.518}{402.43}{H}%
\TextAt{17}{229.145}{406.769}{0}%
\TextAt{9}{233.878}{402.43}{\Character{40}}%
\TextAt{14}{238.43}{402.43}{a}%
\TextAt{9}{244.413}{402.43}{\Character{41}}%
\TextAt{8}{251.873}{402.43}{lies}%
\TextAt{8}{271.099}{402.43}{in}%
\TextAt{8}{283.762}{402.43}{the}%
\TextAt{8}{302.929}{402.43}{part}%
\TextAt{8}{327.308}{402.43}{of}%
\TextAt{8}{339.648}{402.43}{the}%
\TextAt{8}{358.815}{402.43}{lattice}%
\TextAt{13}{393.595}{402.43}{j}%
\TextAt{14}{399.11}{402.43}{a}%
\TextAt{21}{405.093}{400.637}{f}%
\NativeMathAt{413.859}{402.43}{11.955}{\subseteq}%
\TextAt{13}{426.478}{402.43}{K}%
\TextAt{25}{436.438}{400.637}{R}%
\TextAt{8}{445.599}{402.43}{which}%
\TextAt{8}{478.096}{402.43}{is}%
\TextAt{8}{488.87}{402.43}{bounded}%
\TextAt{8}{79.226}{385.384}{by}%
\TextAt{8}{95.151}{385.384}{the}%
\TextAt{8}{114.973}{385.384}{conditions}%
\TextAt{16}{170.637}{385.384}{\Character{106}}%
\TextAt{13}{173.958}{385.384}{f}%
\TextAt{16}{181.004}{385.384}{\Character{106}}%
\TextAt{27}{184.325}{383.591}{p}%
\NativeMathAt{192.133}{385.384}{11.955}{\leq}%
\TextAt{13}{204.752}{385.384}{e}%
\TextAt{3}{210.178}{389.723}{\Character{0}}%
\TextAt{21}{216.764}{389.723}{\Character{23}}%
\TextAt{28}{220.927}{388.726}{p}%
\TextAt{21}{225.01}{389.723}{f}%
\TextAt{28}{229.078}{388.726}{p}%
\TextAt{8}{233.659}{385.384}{,}%
\TextAt{14}{240.541}{385.384}{p}%
\TextAt{16}{249.845}{385.384}{\Character{106}}%
\NativeMathAt{256.468}{385.384}{11.955}{\infty}%
\TextAt{8}{268.441}{385.384}{.}%
\TextAt{8}{276.778}{385.384}{As}%
\TextAt{8}{293.731}{385.384}{the}%
\TextAt{8}{313.554}{385.384}{analogue}%
\TextAt{8}{362.009}{385.384}{of}%
\TextAt{8}{375.004}{385.384}{the}%
\TextAt{8}{394.827}{385.384}{dimension,}%
\TextAt{8}{453.152}{385.384}{we}%
\TextAt{8}{470.045}{385.384}{put}%
\TextAt{13}{491.17}{385.384}{\Character{96}}%
\TextAt{9}{496.082}{385.384}{\Character{40}}%
\TextAt{14}{500.635}{385.384}{a}%
\TextAt{9}{506.618}{385.384}{\Character{41}}%
\TextAt{9}{514.473}{385.384}{\Character{61}}%
\TextAt{9}{526.894}{385.384}{0}%
\TextAt{8}{79.226}{368.338}{if}%
\TextAt{13}{89.95}{368.338}{H}%
\TextAt{17}{100.573}{372.677}{0}%
\TextAt{9}{105.305}{368.338}{\Character{40}}%
\TextAt{14}{109.857}{368.338}{a}%
\TextAt{9}{115.84}{368.338}{\Character{41}}%
\TextAt{9}{123.695}{368.338}{\Character{61}}%
\TextAt{18}{136.139}{368.338}{\Character{63}}%
\TextAt{8}{145.437}{368.338}{,}%
\TextAt{8}{152.587}{368.338}{and}%
\TextAt{8}{175.349}{368.338}{in}%
\TextAt{8}{189.002}{368.338}{all}%
\TextAt{8}{205.261}{368.338}{other}%
\TextAt{8}{235.83}{368.338}{cases}%
\TextAt{13}{255.098}{334.836}{\Character{96}}%
\TextAt{9}{260.01}{334.836}{\Character{40}}%
\TextAt{14}{264.563}{334.836}{a}%
\TextAt{9}{270.546}{334.836}{\Character{41}}%
\TextAt{9}{278.401}{334.836}{\Character{58}\Character{61}}%
\TextAt{8}{294.096}{334.836}{log}%
\TextAt{8}{312.252}{342.923}{\char35\relax }%
\TextAt{13}{322.007}{342.923}{H}%
\TextAt{17}{332.63}{347.262}{0}%
\TextAt{9}{337.362}{342.923}{\Character{40}}%
\TextAt{14}{341.915}{342.923}{a}%
\TextAt{9}{347.898}{342.923}{\Character{41}}%
\Vector{q 0 G 0 g 0.478 w 240.252 -436.176 m 280.45 -436.176 l S Q}%
\TextAt{8}{313.815}{326.635}{vol}%
\TextAt{9}{329.105}{326.635}{\Character{40}}%
\TextAt{13}{333.658}{326.635}{W}%
\TextAt{9}{346.335}{326.635}{\Character{41}}%
\TextAt{13}{353.646}{334.836}{\Character{59}}%
\TextAt{8}{79.226}{301.728}{where}%
\TextAt{8}{113.035}{301.728}{the}%
\TextAt{8}{133.192}{301.728}{normalizing}%
\TextAt{8}{196.937}{301.728}{factor}%
\TextAt{8}{230.435}{301.728}{vol}%
\TextAt{9}{245.726}{301.728}{\Character{40}}%
\TextAt{13}{250.278}{301.728}{W}%
\TextAt{9}{262.955}{301.728}{\Character{41}}%
\TextAt{8}{271.405}{301.728}{is}%
\TextAt{8}{283.169}{301.728}{the}%
\TextAt{8}{303.325}{301.728}{volume}%
\TextAt{8}{343.65}{301.728}{of}%
\TextAt{8}{356.98}{301.728}{the}%
\TextAt{8}{377.137}{301.728}{set}%
\TextAt{13}{206.786}{272.727}{W}%
\TextAt{9}{222.784}{272.727}{\Character{61}}%
\TextAt{16}{235.21}{272.727}{\Character{102}}%
\TextAt{9}{241.187}{272.727}{\Character{40}}%
\TextAt{13}{245.74}{272.727}{z}%
\TextAt{21}{251.177}{270.934}{\Character{28}}%
\TextAt{9}{256.354}{272.727}{\Character{41}}%
\NativeMathAt{264.227}{272.727}{11.955}{\in}%
\TextAt{13}{275.518}{272.727}{K}%
\TextAt{25}{285.477}{270.934}{R}%
\TextAt{9}{295.052}{272.727}{\Character{61}}%
\TextAt{9}{307.462}{272.727}{\Character{91}}%
\TextAt{20}{310.729}{284.085}{\Character{89}}%
\TextAt{21}{316.028}{258.979}{\Character{28}}%
\TextAt{18}{327.998}{272.727}{C}%
\TextAt{9}{336.632}{272.727}{\Character{93}}%
\TextAt{17}{339.884}{277.663}{\Character{43}}%
\TextAt{16}{350.289}{272.727}{\Character{106}}%
\TextAt{16}{356.912}{272.727}{\Character{106}}%
\TextAt{13}{360.252}{272.727}{z}%
\TextAt{21}{365.69}{270.934}{\Character{28}}%
\TextAt{16}{370.866}{272.727}{\Character{106}}%
\NativeMathAt{377.489}{272.727}{11.955}{\leq}%
\TextAt{9}{390.127}{272.727}{1}%
\TextAt{16}{395.98}{272.727}{\Character{103}}%
\TextAt{13}{401.958}{272.727}{\Character{58}}%
\TextAt{8}{79.226}{236.535}{This}%
\TextAt{8}{105.946}{236.535}{volume}%
\TextAt{8}{146.271}{236.535}{is}%
\TextAt{8}{158.035}{236.535}{given}%
\TextAt{8}{188.604}{236.535}{explicitly}%
\TextAt{8}{239.33}{236.535}{by}%
\TextAt{8}{259.538}{207.534}{vol}%
\TextAt{9}{274.829}{207.534}{\Character{40}}%
\TextAt{13}{279.381}{207.534}{W}%
\TextAt{9}{292.058}{207.534}{\Character{41}}%
\TextAt{9}{299.913}{207.534}{\Character{61}}%
\TextAt{9}{312.334}{207.534}{2}%
\TextAt{21}{318.21}{212.47}{r}%
\TextAt{9}{322.765}{207.534}{\Character{40}2}%
\TextAt{13}{333.17}{207.534}{\Character{25}}%
\TextAt{9}{340.239}{207.534}{\Character{41}}%
\TextAt{21}{344.792}{212.47}{s}%
\TextAt{13}{349.206}{207.534}{\Character{59}}%
\TextAt{8}{79.226}{178.533}{where}%
\TextAt{13}{112.555}{178.533}{r}%
\TextAt{8}{118.155}{178.533}{,}%
\TextAt{8}{124.92}{178.533}{resp.}%
\TextAt{13}{154.082}{178.533}{s}%
\TextAt{8}{159.596}{178.533}{,}%
\TextAt{8}{166.36}{178.533}{is}%
\TextAt{8}{177.644}{178.533}{the}%
\TextAt{8}{197.32}{178.533}{number}%
\TextAt{8}{239.435}{178.533}{of}%
\TextAt{8}{252.285}{178.533}{real,}%
\TextAt{8}{277.915}{178.533}{resp.}%
\TextAt{8}{307.076}{178.533}{complex,}%
\TextAt{8}{355.792}{178.533}{primes}%
\TextAt{8}{393.09}{178.533}{of}%
\TextAt{13}{405.939}{178.533}{K}%
\TextAt{8}{420.173}{178.533}{(see}%
\TextAt{8}{443.161}{178.533}{the}%
\TextAt{8}{462.837}{178.533}{proof}%
\TextAt{8}{492.925}{178.533}{of}%
\TextAt{8}{505.775}{178.533}{chap.}%
\TextAt{8}{79.226}{161.487}{I,}%
\TextAt{8}{90.596}{161.487}{(5.3)).}%
\TextAt{8}{127.676}{161.487}{In}%
\TextAt{8}{142.298}{161.487}{particular,}%
\TextAt{8}{199.539}{161.487}{one}%
\TextAt{8}{220.999}{161.487}{has}%
\TextAt{13}{183.485}{128.48}{H}%
\TextAt{17}{194.108}{133.416}{0}%
\TextAt{9}{198.84}{128.48}{\Character{40}}%
\TextAt{14}{203.392}{128.48}{o}%
\TextAt{9}{209.236}{128.48}{\Character{41}}%
\TextAt{9}{217.091}{128.48}{\Character{61}}%
\TextAt{13}{229.535}{128.48}{\Character{22}}%
\TextAt{9}{236.578}{128.48}{\Character{40}}%
\TextAt{13}{241.13}{128.48}{K}%
\TextAt{9}{251.948}{128.48}{\Character{41}}%
\TextAt{13}{256.5}{128.48}{\Character{59}}%
\TextAt{8}{285.655}{128.48}{so}%
\TextAt{8}{300.025}{128.48}{that}%
\TextAt{13}{333.452}{128.48}{\Character{96}}%
\TextAt{9}{338.364}{128.48}{\Character{40}}%
\TextAt{14}{342.916}{128.48}{o}%
\TextAt{9}{348.76}{128.48}{\Character{41}}%
\TextAt{9}{356.615}{128.48}{\Character{61}}%
\TextAt{8}{369.059}{128.48}{log}%
\TextAt{8}{387.279}{136.568}{\char35\relax }%
\TextAt{13}{397.034}{136.568}{\Character{22}}%
\TextAt{9}{404.077}{136.568}{\Character{40}}%
\TextAt{13}{408.63}{136.568}{K}%
\TextAt{9}{419.447}{136.568}{\Character{41}}%
\Vector{q 0 G 0 g 0.478 w 315.215 -642.531 m 352.063 -642.531 l S Q}%
\TextAt{9}{387.215}{120.28}{2}%
\TextAt{21}{393.068}{123.733}{r}%
\TextAt{9}{397.623}{120.28}{\Character{40}2}%
\TextAt{13}{408.028}{120.28}{\Character{25}}%
\TextAt{9}{415.097}{120.28}{\Character{41}}%
\TextAt{21}{419.649}{123.733}{s}%
\TextAt{13}{425.259}{128.48}{\Character{59}}%
\TextAt{8}{79.226}{94.865}{because}%
\TextAt{16}{123.001}{94.865}{\Character{106}}%
\TextAt{13}{126.322}{94.865}{f}%
\TextAt{16}{133.369}{94.865}{\Character{106}}%
\TextAt{27}{136.69}{93.071}{p}%
\NativeMathAt{145.304}{94.865}{11.955}{\leq}%
\TextAt{9}{158.731}{94.865}{1}%
\TextAt{8}{168.955}{94.865}{for}%
\TextAt{8}{187.314}{94.865}{all}%
\TextAt{14}{204.046}{94.865}{p}%
\TextAt{8}{210.029}{94.865}{,}%
\TextAt{8}{217.771}{94.865}{and}%
\TextAt{20}{241.007}{103.831}{\Character{81}}%
\TextAt{27}{252.298}{91.378}{p}%
\TextAt{16}{258.777}{94.865}{\Character{106}}%
\TextAt{13}{262.098}{94.865}{f}%
\TextAt{16}{269.145}{94.865}{\Character{106}}%
\TextAt{27}{272.465}{93.071}{p}%
\TextAt{9}{281.08}{94.865}{\Character{61}}%
\TextAt{9}{294.303}{94.865}{1}%
\TextAt{8}{304.537}{94.865}{implies}%
\TextAt{16}{344.738}{94.865}{\Character{106}}%
\TextAt{13}{348.058}{94.865}{f}%
\TextAt{16}{355.105}{94.865}{\Character{106}}%
\TextAt{27}{358.426}{93.071}{p}%
\TextAt{9}{367.041}{94.865}{\Character{61}}%
\TextAt{9}{380.264}{94.865}{1}%
\TextAt{8}{390.497}{94.865}{for}%
\TextAt{8}{408.856}{94.865}{all}%
\TextAt{14}{425.589}{94.865}{p}%
\TextAt{8}{431.572}{94.865}{,}%
\TextAt{8}{439.313}{94.865}{so}%
\TextAt{8}{454.157}{94.865}{that}%
\TextAt{13}{479.999}{94.865}{H}%
\TextAt{17}{490.622}{99.203}{0}%
\TextAt{9}{495.354}{94.865}{\Character{40}}%
\TextAt{14}{499.906}{94.865}{o}%
\TextAt{9}{505.751}{94.865}{\Character{41}}%
\TextAt{8}{514.674}{94.865}{is}%
\TextAt{8}{526.911}{94.865}{a}%
\TextAt{8}{79.226}{77.819}{finite}%
\TextAt{8}{108.511}{77.819}{subgroup}%
\TextAt{8}{158.682}{77.819}{of}%
\TextAt{13}{171.385}{77.819}{K}%
\TextAt{3}{182.203}{82.157}{\Character{3}}%
\TextAt{8}{190.206}{77.819}{and}%
\TextAt{8}{212.342}{77.819}{thus}%
\TextAt{8}{237.467}{77.819}{must}%
\TextAt{8}{265.843}{77.819}{consist}%
\TextAt{8}{303.712}{77.819}{of}%
\TextAt{8}{316.416}{77.819}{all}%
\TextAt{8}{332.048}{77.819}{roots}%
\TextAt{8}{361.082}{77.819}{of}%
\TextAt{8}{373.785}{77.819}{unity.}%
\TextAt{8}{407.716}{77.819}{This}%
\TextAt{8}{433.81}{77.819}{normalization}%
\TextAt{8}{507.341}{77.819}{leads}%
\TextAt{8}{297.211}{43.033}{195}%
\Destination{78.23}{809.815}{nav.page.195}%
\Destination{79.36}{776.155}{nav.print.213}%
\Destination{79.36}{316.755}{nav.print.214}%
\Destination{79.36}{633.455}{nav.number.3.3.4}%
\Hotlink{308.587}{758.955}{6.858}{11.047}{nav.section.3.1}%
\Hotlink{331.714}{758.955}{18.574}{11.047}{nav.print.192}%
\Hotlink{224.137}{643.355}{25.078}{11.047}{nav.number.3.3.3}%
\Hotlink{90.096}{159.255}{25.078}{11.047}{nav.number.1.5.3}%
}

% Page 208 (body page 196)
\Page{%
\TextAt{2}{79.226}{796.943}{III.}%
\TextAt{2}{98.317}{796.943}{Riemann–Roch}%
\TextAt{2}{174.375}{796.943}{Theory}%
\Vector{q 0 G 0 g 0.299 w 7.226 18.729 m 460.769 18.729 l S Q}%
\TextAt{8}{79.226}{761.124}{us}%
\TextAt{8}{93.605}{761.124}{necessarily}%
\TextAt{8}{151.298}{761.124}{to}%
\TextAt{8}{164.972}{761.124}{the}%
\TextAt{8}{184.492}{761.124}{following}%
\TextAt{8}{233.29}{761.124}{definition}%
\TextAt{8}{284.682}{761.124}{of}%
\TextAt{8}{297.376}{761.124}{the}%
\TextAt{8}{316.896}{761.124}{genus}%
\TextAt{8}{348.514}{761.124}{of}%
\TextAt{8}{361.208}{761.124}{a}%
\TextAt{8}{370.326}{761.124}{number}%
\TextAt{8}{412.286}{761.124}{field,}%
\TextAt{8}{440.386}{761.124}{which}%
\TextAt{8}{473.236}{761.124}{had}%
\TextAt{8}{495.362}{761.124}{already}%
\TextAt{8}{79.226}{744.078}{been}%
\TextAt{8}{106.855}{744.078}{proposed}%
\TextAt{19}{156.672}{744.078}{ad}%
\TextAt{19}{172.811}{744.078}{hoc}%
\TextAt{8}{193.577}{744.078}{by}%
\TextAt{8}{209.836}{744.078}{the}%
\TextAt{8}{229.993}{744.078}{French}%
\TextAt{8}{268.19}{744.078}{mathematician}%
\TextAt{8}{348.194}{744.078}{André}%
\TextAt{8}{383.629}{744.078}{Weil}%
\TextAt{8}{410.277}{744.078}{in}%
\TextAt{8}{423.93}{744.078}{1939}%
\TextAt{8}{451.259}{744.078}{(see}%
\TextAt{8}{474.727}{744.078}{[138]).}%
\TextAt{12}{79.226}{718.246}{Definition}%
\TextAt{12}{142.015}{718.246}{(3.5).}%
\TextAt{8}{177.33}{718.246}{The}%
\TextAt{12}{201.384}{718.246}{genus}%
\TextAt{8}{238.027}{718.246}{of}%
\TextAt{8}{251.357}{718.246}{a}%
\TextAt{8}{261.112}{718.246}{number}%
\TextAt{8}{303.708}{718.246}{field}%
\TextAt{13}{329.065}{718.246}{K}%
\TextAt{8}{343.78}{718.246}{is}%
\TextAt{8}{355.544}{718.246}{defined}%
\TextAt{8}{395.857}{718.246}{to}%
\TextAt{8}{410.167}{718.246}{be}%
\TextAt{8}{426.092}{718.246}{the}%
\TextAt{8}{446.248}{718.246}{real}%
\TextAt{8}{469.011}{718.246}{number}%
\TextAt{13}{212.754}{683.441}{g}%
\TextAt{9}{222.109}{683.441}{\Character{61}}%
\TextAt{13}{234.534}{683.441}{\Character{96}}%
\TextAt{9}{239.446}{683.441}{\Character{40}}%
\TextAt{14}{243.999}{683.441}{o}%
\TextAt{9}{249.843}{683.441}{\Character{41}}%
\TextAt{16}{257.052}{683.441}{\Character{0}}%
\TextAt{13}{269.007}{683.441}{\Character{31}}%
\TextAt{9}{276.341}{683.441}{\Character{40}}%
\TextAt{14}{280.893}{683.441}{o}%
\TextAt{9}{286.737}{683.441}{\Character{41}}%
\TextAt{9}{294.592}{683.441}{\Character{61}}%
\TextAt{8}{307.036}{683.441}{log}%
\TextAt{8}{325.192}{691.75}{\char35\relax }%
\TextAt{13}{334.947}{691.75}{\Character{22}}%
\TextAt{9}{341.99}{691.75}{\Character{40}}%
\TextAt{13}{346.542}{691.75}{K}%
\TextAt{9}{357.36}{691.75}{\Character{41}}%
\TextAt{20}{361.912}{701.972}{\Character{112}}%
\Vector{q 0 G 0 g 0.478 w 301.868 -71.789 m 322.795 -71.789 l S Q}%
\TextAt{16}{373.868}{691.75}{\Character{106}}%
\TextAt{13}{377.189}{691.75}{d}%
\TextAt{21}{383.271}{689.957}{K}%
\TextAt{16}{391.474}{691.75}{\Character{106}}%
\Vector{q 0 G 0 g 0.478 w 253.192 -87.57 m 322.795 -87.57 l S Q}%
\TextAt{9}{341.569}{675.24}{2}%
\TextAt{21}{347.422}{678.694}{r}%
\TextAt{9}{351.977}{675.24}{\Character{40}2}%
\TextAt{13}{362.382}{675.24}{\Character{25}}%
\TextAt{9}{369.451}{675.24}{\Character{41}}%
\TextAt{21}{374.004}{678.694}{s}%
\TextAt{13}{395.99}{683.441}{\Character{58}}%
\TextAt{8}{97.159}{642.732}{Observe}%
\TextAt{8}{142.307}{642.732}{that}%
\TextAt{8}{167.885}{642.732}{the}%
\TextAt{8}{188.25}{642.732}{genus}%
\TextAt{8}{220.714}{642.732}{of}%
\TextAt{8}{234.253}{642.732}{the}%
\TextAt{8}{254.618}{642.732}{field}%
\TextAt{8}{280.184}{642.732}{of}%
\TextAt{8}{293.723}{642.732}{rational}%
\TextAt{8}{337.52}{642.732}{numbers}%
\TextAt{18}{384.94}{642.732}{Q}%
\TextAt{8}{398.345}{642.732}{is}%
\TextAt{9}{410.318}{642.732}{0}%
\TextAt{8}{416.171}{642.732}{.}%
\TextAt{8}{425.245}{642.732}{Using}%
\TextAt{8}{458.355}{642.732}{this}%
\TextAt{8}{481.386}{642.732}{definition,}%
\TextAt{8}{79.226}{625.686}{the}%
\TextAt{8}{99.383}{625.686}{Riemann–Roch}%
\TextAt{8}{181.012}{625.686}{formula}%
\TextAt{8}{223.943}{625.686}{(3.4)}%
\TextAt{8}{251.919}{625.686}{takes}%
\TextAt{8}{281.902}{625.686}{the}%
\TextAt{8}{302.059}{625.686}{following}%
\TextAt{8}{351.493}{625.686}{shape:}%
\TextAt{12}{79.226}{599.855}{Proposition}%
\TextAt{12}{151.937}{599.855}{(3.6).}%
\TextAt{8}{187.253}{599.855}{For}%
\TextAt{8}{208.222}{599.855}{every}%
\TextAt{8}{239.115}{599.855}{replete}%
\TextAt{8}{277.479}{599.855}{ideal}%
\TextAt{14}{305.442}{599.855}{a}%
\TextAt{8}{315.322}{599.855}{of}%
\TextAt{13}{328.652}{599.855}{K}%
\TextAt{8}{343.367}{599.855}{one}%
\TextAt{8}{364.827}{599.855}{has}%
\TextAt{13}{241.501}{571.488}{\Character{31}}%
\TextAt{9}{248.834}{571.488}{\Character{40}}%
\TextAt{14}{253.387}{571.488}{a}%
\TextAt{9}{259.37}{571.488}{\Character{41}}%
\TextAt{9}{267.225}{571.488}{\Character{61}}%
\TextAt{8}{279.668}{571.488}{deg}%
\TextAt{9}{297.23}{571.488}{\Character{40}}%
\TextAt{14}{301.783}{571.488}{a}%
\TextAt{9}{307.766}{571.488}{\Character{41}}%
\TextAt{9}{314.963}{571.488}{\Character{43}}%
\TextAt{13}{326.736}{571.488}{\Character{96}}%
\TextAt{9}{331.648}{571.488}{\Character{40}}%
\TextAt{14}{336.201}{571.488}{o}%
\TextAt{9}{342.045}{571.488}{\Character{41}}%
\TextAt{16}{349.254}{571.488}{\Character{0}}%
\TextAt{13}{361.209}{571.488}{g}%
\TextAt{13}{367.223}{571.488}{\Character{58}}%
\TextAt{8}{97.159}{534.335}{The}%
\TextAt{8}{121.213}{534.335}{analogue}%
\TextAt{8}{170.002}{534.335}{of}%
\TextAt{8}{183.332}{534.335}{the}%
\TextAt{8}{203.488}{534.335}{strong}%
\TextAt{8}{239.33}{534.335}{Riemann–Roch}%
\TextAt{8}{320.96}{534.335}{formula}%
\TextAt{13}{214.468}{505.968}{\Character{96}}%
\TextAt{9}{219.38}{505.968}{\Character{40}}%
\TextAt{13}{223.932}{505.968}{D}%
\TextAt{9}{233.976}{505.968}{\Character{41}}%
\TextAt{9}{241.831}{505.968}{\Character{61}}%
\TextAt{8}{254.275}{505.968}{deg}%
\TextAt{9}{271.837}{505.968}{\Character{40}}%
\TextAt{13}{276.39}{505.968}{D}%
\TextAt{9}{286.434}{505.968}{\Character{41}}%
\TextAt{9}{293.631}{505.968}{\Character{43}}%
\TextAt{9}{305.383}{505.968}{1}%
\TextAt{16}{313.914}{505.968}{\Character{0}}%
\TextAt{13}{325.869}{505.968}{g}%
\TextAt{9}{334.56}{505.968}{\Character{43}}%
\TextAt{13}{346.321}{505.968}{\Character{96}}%
\TextAt{9}{351.234}{505.968}{\Character{40}}%
\TextAt{16}{355.786}{505.968}{\Character{75}}%
\TextAt{16}{367.705}{505.968}{\Character{0}}%
\TextAt{13}{379.68}{505.968}{D}%
\TextAt{9}{389.724}{505.968}{\Character{41}}%
\TextAt{13}{394.276}{505.968}{\Character{59}}%
\TextAt{8}{79.226}{477.6}{hinges}%
\TextAt{8}{115.757}{477.6}{on}%
\TextAt{8}{132.718}{477.6}{the}%
\TextAt{8}{153.575}{477.6}{following}%
\TextAt{8}{203.711}{477.6}{deep}%
\TextAt{8}{231.718}{477.6}{theorem}%
\TextAt{8}{277.945}{477.6}{of}%
\TextAt{8}{291.976}{477.6}{Minkowski}%
\TextAt{8}{350.947}{477.6}{theory,}%
\TextAt{8}{390.846}{477.6}{which}%
\TextAt{8}{425.033}{477.6}{is}%
\TextAt{8}{437.499}{477.6}{due}%
\TextAt{8}{460.305}{477.6}{to}%
\TextAt{8}{475.317}{477.6}{Serge}%
\TextAt{8}{507.233}{477.6}{Lang}%
\TextAt{8}{79.226}{460.554}{and}%
\TextAt{8}{101.439}{460.554}{which}%
\TextAt{8}{134.375}{460.554}{reflects}%
\TextAt{8}{173.553}{460.554}{an}%
\TextAt{8}{189.262}{460.554}{arithmetic}%
\TextAt{8}{245.296}{460.554}{analogue}%
\TextAt{8}{293.535}{460.554}{of}%
\TextAt{8}{306.315}{460.554}{Serre}%
\TextAt{8}{335.677}{460.554}{duality.}%
\TextAt{8}{378.743}{460.554}{As}%
\TextAt{8}{395.481}{460.554}{usual,}%
\TextAt{8}{428.922}{460.554}{let}%
\TextAt{13}{445.276}{460.554}{r}%
\TextAt{8}{450.877}{460.554}{,}%
\TextAt{8}{457.586}{460.554}{resp.}%
\TextAt{13}{486.725}{460.554}{s}%
\TextAt{8}{492.239}{460.554}{,}%
\TextAt{8}{498.948}{460.554}{denote}%
\TextAt{8}{79.226}{443.508}{the}%
\TextAt{8}{99.383}{443.508}{number}%
\TextAt{8}{141.979}{443.508}{of}%
\TextAt{8}{155.309}{443.508}{real,}%
\TextAt{8}{181.323}{443.508}{resp.}%
\TextAt{8}{210.645}{443.508}{complex,}%
\TextAt{8}{259.745}{443.508}{primes,}%
\TextAt{8}{300.775}{443.508}{and}%
\TextAt{13}{323.538}{443.508}{n}%
\TextAt{9}{333.846}{443.508}{\Character{61}}%
\TextAt{9}{346.256}{443.508}{\Character{91}}%
\TextAt{13}{349.524}{443.508}{K}%
\TextAt{9}{363.662}{443.508}{\Character{58}}%
\TextAt{18}{370.235}{443.508}{Q}%
\TextAt{9}{379.533}{443.508}{\Character{93}}%
\TextAt{8}{382.785}{443.508}{.}%
\TextAt{12}{79.226}{417.677}{Theorem}%
\TextAt{12}{136.24}{417.677}{(3.7)}%
\TextAt{12}{168.4}{417.677}{(S.}%
\TextAt{12}{189.333}{417.677}{Lang).}%
\TextAt{8}{231.033}{417.677}{For}%
\TextAt{8}{252.002}{417.677}{replete}%
\TextAt{8}{290.366}{417.677}{ideals}%
\TextAt{14}{322.944}{417.677}{a}%
\TextAt{9}{332.248}{417.677}{\Character{61}}%
\TextAt{20}{344.674}{426.644}{\Character{81}}%
\TextAt{27}{355.965}{414.19}{p}%
\TextAt{14}{362.444}{417.677}{p}%
\TextAt{21}{368.427}{422.016}{\Character{23}}%
\TextAt{28}{372.59}{421.019}{p}%
\NativeMathAt{380.491}{417.677}{11.955}{\in}%
\TextAt{13}{391.782}{417.677}{J}%
\TextAt{9}{399.36}{417.677}{\Character{40}}%
\Vector{q 0 G 0 g 0.478 w 331.912 -348.966 m 337.756 -348.966 l S Q}%
\TextAt{14}{403.912}{417.677}{o}%
\TextAt{9}{409.756}{417.677}{\Character{41}}%
\TextAt{8}{414.308}{417.677}{,}%
\TextAt{8}{421.458}{417.677}{one}%
\TextAt{8}{442.917}{417.677}{has}%
\TextAt{8}{198.851}{383.312}{\char35\relax }%
\TextAt{13}{208.606}{383.312}{H}%
\TextAt{17}{219.229}{388.248}{0}%
\TextAt{9}{223.961}{383.312}{\Character{40}}%
\TextAt{14}{228.514}{383.312}{a}%
\TextAt{3}{234.497}{388.248}{\Character{0}}%
\TextAt{17}{241.083}{388.248}{1}%
\TextAt{9}{245.815}{383.312}{\Character{41}}%
\TextAt{9}{253.67}{383.312}{\Character{61}}%
\TextAt{9}{267.31}{391.399}{2}%
\TextAt{21}{273.163}{395.738}{r}%
\TextAt{9}{277.717}{391.399}{\Character{40}2}%
\TextAt{13}{288.123}{391.399}{\Character{25}}%
\TextAt{9}{295.192}{391.399}{\Character{41}}%
\TextAt{21}{299.744}{395.738}{s}%
\Vector{q 0 G 0 g 0.478 w 195.31 -387.7 m 232.158 -387.7 l S Q}%
\TextAt{20}{269.293}{383.67}{\Character{112}}%
\Vector{q 0 G 0 g 0.478 w 209.248 -390.09 m 230.175 -390.09 l S Q}%
\TextAt{16}{281.248}{373.449}{\Character{106}}%
\TextAt{13}{284.569}{373.449}{d}%
\TextAt{21}{290.652}{371.655}{K}%
\TextAt{16}{298.854}{373.449}{\Character{106}}%
\TextAt{14}{305.354}{383.312}{N}%
\TextAt{9}{315.302}{383.312}{\Character{40}}%
\TextAt{14}{319.854}{383.312}{a}%
\TextAt{9}{325.837}{383.312}{\Character{41}}%
\TextAt{9}{333.034}{383.312}{\Character{43}}%
\TextAt{13}{344.808}{383.312}{O}%
\TextAt{9}{354.079}{383.312}{\Character{40}}%
\TextAt{14}{358.631}{383.312}{N}%
\TextAt{9}{368.579}{383.312}{\Character{40}}%
\TextAt{14}{373.132}{383.312}{a}%
\TextAt{9}{379.115}{383.312}{\Character{41}}%
\TextAt{17}{383.667}{388.248}{1}%
\TextAt{3}{387.901}{388.248}{\Character{0}}%
\TextAt{17}{394.488}{388.248}{1\Character{47}}%
\TextAt{21}{402.956}{388.248}{n}%
\TextAt{9}{408.592}{383.312}{\Character{41}}%
\TextAt{8}{79.226}{348.518}{if}%
\TextAt{14}{89.894}{348.518}{N}%
\TextAt{9}{99.842}{348.518}{\Character{40}}%
\TextAt{14}{104.394}{348.518}{a}%
\TextAt{9}{110.378}{348.518}{\Character{41}}%
\NativeMathAt{118.251}{348.518}{11.955}{\rightarrow}%
\NativeMathAt{133.518}{348.518}{11.955}{\infty}%
\TextAt{8}{145.482}{348.518}{.}%
\TextAt{8}{153.912}{348.518}{Here,}%
\TextAt{8}{184.747}{348.518}{as}%
\TextAt{8}{199.061}{348.518}{usual,}%
\TextAt{13}{232.897}{348.518}{O}%
\TextAt{9}{242.169}{348.518}{\Character{40}}%
\TextAt{13}{246.721}{348.518}{t}%
\TextAt{9}{250.948}{348.518}{\Character{41}}%
\TextAt{8}{259.342}{348.518}{denotes}%
\TextAt{8}{301.619}{348.518}{a}%
\TextAt{8}{311.319}{348.518}{function}%
\TextAt{8}{357.111}{348.518}{such}%
\TextAt{8}{383.452}{348.518}{that}%
\TextAt{13}{408.765}{348.518}{O}%
\TextAt{9}{418.036}{348.518}{\Character{40}}%
\TextAt{13}{422.589}{348.518}{t}%
\TextAt{9}{426.816}{348.518}{\Character{41}\Character{47}}%
\TextAt{13}{437.221}{348.518}{t}%
\TextAt{8}{445.29}{348.518}{remains}%
\TextAt{8}{488.87}{348.518}{bounded}%
\TextAt{8}{79.226}{331.472}{as}%
\TextAt{13}{93.596}{331.472}{t}%
\NativeMathAt{101.144}{331.472}{11.955}{\rightarrow}%
\NativeMathAt{116.411}{331.472}{11.955}{\infty}%
\TextAt{8}{128.375}{331.472}{.}%
\TextAt{8}{97.159}{305.641}{For}%
\TextAt{8}{118.128}{305.641}{the}%
\TextAt{8}{138.285}{305.641}{proof}%
\TextAt{8}{168.854}{305.641}{of}%
\TextAt{8}{182.184}{305.641}{the}%
\TextAt{8}{202.341}{305.641}{theorem}%
\TextAt{8}{247.866}{305.641}{we}%
\TextAt{8}{265.093}{305.641}{need}%
\TextAt{8}{292.399}{305.641}{the}%
\TextAt{8}{312.555}{305.641}{following}%
\TextAt{12}{79.226}{279.81}{Lemma}%
\TextAt{12}{128.174}{279.81}{(3.8).}%
\TextAt{8}{164.604}{279.81}{Let}%
\TextAt{14}{186.689}{279.81}{a}%
\TextAt{17}{192.672}{278.017}{1}%
\TextAt{13}{197.404}{279.81}{\Character{59}}%
\TextAt{13}{202.628}{279.81}{\Character{58}}%
\TextAt{13}{207.865}{279.81}{\Character{58}}%
\TextAt{13}{213.101}{279.81}{\Character{58}}%
\TextAt{13}{218.326}{279.81}{\Character{59}}%
\TextAt{14}{223.625}{279.81}{a}%
\TextAt{21}{229.608}{278.017}{h}%
\TextAt{8}{239.998}{279.81}{be}%
\TextAt{8}{257.038}{279.81}{fractional}%
\TextAt{8}{310.516}{279.81}{ideals}%
\TextAt{8}{344.209}{279.81}{representing}%
\TextAt{8}{411.4}{279.81}{the}%
\TextAt{8}{432.671}{279.81}{classes}%
\TextAt{8}{471.039}{279.81}{of}%
\TextAt{8}{485.484}{279.81}{the}%
\TextAt{8}{506.755}{279.81}{finite}%
\TextAt{8}{79.226}{262.764}{ideal}%
\TextAt{8}{107.189}{262.764}{class}%
\TextAt{8}{134.626}{262.764}{group}%
\TextAt{8}{167.802}{262.764}{Pic}%
\TextAt{9}{184.216}{262.764}{\Character{40}}%
\TextAt{14}{188.769}{262.764}{o}%
\TextAt{9}{194.613}{262.764}{\Character{41}}%
\TextAt{8}{199.165}{262.764}{.}%
\TextAt{8}{207.614}{262.764}{Let}%
\TextAt{13}{228.583}{262.764}{c}%
\TextAt{8}{237.518}{262.764}{be}%
\TextAt{8}{253.443}{262.764}{a}%
\TextAt{8}{263.198}{262.764}{positive}%
\TextAt{8}{306.512}{262.764}{constant}%
\TextAt{8}{353.735}{262.764}{and}%
\TextAt{16}{155.625}{234.396}{\Character{65}}%
\TextAt{21}{165.171}{232.603}{i}%
\TextAt{9}{171.873}{234.396}{\Character{61}}%
\TextAt{16}{184.299}{234.396}{\Character{102}}%
\TextAt{14}{190.276}{234.396}{a}%
\TextAt{9}{199.58}{234.396}{\Character{61}}%
\TextAt{20}{212.006}{245.754}{\Character{89}}%
\TextAt{27}{217.649}{220.648}{p}%
\TextAt{14}{229.274}{234.396}{p}%
\TextAt{21}{235.257}{239.333}{\Character{23}}%
\TextAt{28}{239.42}{238.336}{p}%
\TextAt{16}{247.322}{234.396}{\Character{106}}%
\TextAt{14}{253.963}{234.396}{a}%
\TextAt{21}{259.946}{232.603}{f}%
\TextAt{9}{268.712}{234.396}{\Character{61}}%
\TextAt{14}{281.138}{234.396}{a}%
\TextAt{21}{287.121}{232.603}{i}%
\TextAt{13}{290.502}{234.396}{\Character{59}}%
\TextAt{14}{299.644}{234.396}{N}%
\TextAt{9}{309.592}{234.396}{\Character{40}}%
\TextAt{14}{314.144}{234.396}{p}%
\TextAt{9}{320.127}{234.396}{\Character{41}}%
\TextAt{21}{324.679}{239.333}{\Character{23}}%
\TextAt{28}{328.842}{238.336}{p}%
\NativeMathAt{336.744}{234.396}{11.955}{\leq}%
\TextAt{13}{349.363}{234.396}{c}%
\TextAt{14}{354.401}{234.396}{N}%
\TextAt{9}{364.349}{234.396}{\Character{40}}%
\TextAt{14}{368.902}{234.396}{a}%
\TextAt{9}{374.885}{234.396}{\Character{41}}%
\TextAt{21}{379.437}{239.333}{f}%
\TextAt{28}{383.505}{238.336}{p}%
\TextAt{17}{387.588}{239.333}{\Character{47}}%
\TextAt{21}{391.822}{239.333}{n}%
\TextAt{8}{401.356}{234.396}{for}%
\TextAt{14}{419.241}{234.396}{p}%
\TextAt{16}{428.545}{234.396}{\Character{106}}%
\NativeMathAt{435.168}{234.396}{11.955}{\infty}%
\TextAt{16}{447.123}{234.396}{\Character{103}}%
\TextAt{13}{453.119}{234.396}{\Character{58}}%
\TextAt{8}{79.226}{197.331}{Then}%
\TextAt{8}{109.784}{197.331}{the}%
\TextAt{8}{129.94}{197.331}{constant}%
\TextAt{13}{177.163}{197.331}{c}%
\TextAt{8}{186.098}{197.331}{can}%
\TextAt{8}{207.558}{197.331}{be}%
\TextAt{8}{223.482}{197.331}{chosen}%
\TextAt{8}{260.938}{197.331}{in}%
\TextAt{8}{274.59}{197.331}{such}%
\TextAt{8}{300.987}{197.331}{a}%
\TextAt{8}{310.743}{197.331}{way}%
\TextAt{8}{334.486}{197.331}{that}%
\TextAt{13}{258.864}{160.9}{J}%
\TextAt{9}{266.441}{160.9}{\Character{40}}%
\Vector{q 0 G 0 g 0.478 w 198.994 -605.743 m 204.838 -605.743 l S Q}%
\TextAt{14}{270.994}{160.9}{o}%
\TextAt{9}{276.838}{160.9}{\Character{41}}%
\TextAt{9}{284.693}{160.9}{\Character{61}}%
\TextAt{21}{301.548}{175.844}{h}%
\TextAt{20}{297.347}{172.258}{\Character{91}}%
\TextAt{21}{297.136}{147.063}{i}%
\TextAt{17}{300.02}{147.063}{\Character{61}1}%
\TextAt{16}{312.833}{160.9}{\Character{65}}%
\TextAt{21}{322.379}{159.107}{i}%
\TextAt{13}{325.76}{160.9}{P}%
\TextAt{9}{334.931}{160.9}{\Character{40}}%
\Vector{q 0 G 0 g 0.478 w 267.483 -605.743 m 273.328 -605.743 l S Q}%
\TextAt{14}{339.483}{160.9}{o}%
\TextAt{9}{345.328}{160.9}{\Character{41}}%
\TextAt{13}{349.88}{160.9}{\Character{58}}%
\TextAt{19}{79.226}{107.013}{Proof.}%
\TextAt{8}{116}{107.013}{Let}%
\TextAt{16}{137.571}{107.013}{\Character{66}}%
\TextAt{21}{145.423}{105.22}{i}%
\TextAt{9}{153.149}{107.013}{\Character{61}}%
\TextAt{16}{166.599}{107.013}{\Character{102}}%
\TextAt{14}{172.576}{107.013}{a}%
\NativeMathAt{182.905}{107.013}{11.955}{\in}%
\TextAt{13}{195.22}{107.013}{J}%
\TextAt{9}{202.797}{107.013}{\Character{40}}%
\Vector{q 0 G 0 g 0.478 w 135.349 -659.63 m 141.193 -659.63 l S Q}%
\TextAt{14}{207.349}{107.013}{o}%
\TextAt{9}{213.193}{107.013}{\Character{41}}%
\TextAt{16}{222.091}{107.013}{\Character{106}}%
\TextAt{14}{229.757}{107.013}{a}%
\TextAt{21}{235.74}{105.22}{f}%
\TextAt{9}{245.53}{107.013}{\Character{61}}%
\TextAt{14}{258.979}{107.013}{a}%
\TextAt{21}{264.962}{105.22}{i}%
\TextAt{16}{268.344}{107.013}{\Character{103}}%
\TextAt{8}{274.321}{107.013}{.}%
\TextAt{8}{284.573}{107.013}{Multiplying}%
\TextAt{8}{348.907}{107.013}{by}%
\TextAt{8}{365.767}{107.013}{a}%
\TextAt{8}{376.123}{107.013}{suitable}%
\TextAt{8}{420.361}{107.013}{replete}%
\TextAt{8}{459.326}{107.013}{principal}%
\TextAt{8}{508.704}{107.013}{ideal}%
\TextAt{9}{79.226}{89.967}{\Character{91}}%
\TextAt{13}{82.478}{89.967}{a}%
\TextAt{9}{88.623}{89.967}{\Character{93}}%
\TextAt{8}{91.874}{89.967}{,}%
\TextAt{8}{99.765}{89.967}{every}%
\TextAt{14}{131.251}{89.967}{a}%
\NativeMathAt{141.566}{89.967}{11.955}{\in}%
\TextAt{13}{153.868}{89.967}{J}%
\TextAt{9}{161.445}{89.967}{\Character{40}}%
\Vector{q 0 G 0 g 0.478 w 93.997 -676.676 m 99.842 -676.676 l S Q}%
\TextAt{14}{165.997}{89.967}{o}%
\TextAt{9}{171.842}{89.967}{\Character{41}}%
\TextAt{8}{180.885}{89.967}{may}%
\TextAt{8}{206.847}{89.967}{be}%
\TextAt{8}{223.364}{89.967}{transformed}%
\TextAt{8}{289.388}{89.967}{into}%
\TextAt{8}{313.725}{89.967}{a}%
\TextAt{8}{324.073}{89.967}{replete}%
\TextAt{8}{363.031}{89.967}{ideal}%
\TextAt{14}{391.587}{89.967}{a}%
\TextAt{3}{397.57}{94.306}{\Character{48}}%
\TextAt{9}{404.697}{89.967}{\Character{61}}%
\TextAt{14}{418.134}{89.967}{a}%
\TextAt{9}{424.117}{89.967}{\Character{91}}%
\TextAt{13}{427.369}{89.967}{a}%
\TextAt{9}{433.514}{89.967}{\Character{93}}%
\TextAt{8}{441.256}{89.967}{such}%
\TextAt{8}{468.246}{89.967}{that}%
\TextAt{14}{494.209}{89.967}{a}%
\TextAt{3}{500.192}{94.306}{\Character{48}}%
\TextAt{21}{500.192}{86.858}{f}%
\TextAt{9}{509.969}{89.967}{\Character{61}}%
\TextAt{14}{523.405}{89.967}{a}%
\TextAt{21}{529.388}{88.174}{i}%
\TextAt{8}{79.226}{72.921}{for}%
\TextAt{8}{97.219}{72.921}{some}%
\TextAt{13}{126.654}{72.921}{i}%
\TextAt{8}{130.647}{72.921}{.}%
\TextAt{8}{139.421}{72.921}{Consequently,}%
\TextAt{8}{214.408}{72.921}{one}%
\TextAt{8}{235.975}{72.921}{has}%
\TextAt{13}{256.958}{72.921}{J}%
\TextAt{9}{264.535}{72.921}{\Character{40}}%
\Vector{q 0 G 0 g 0.478 w 197.087 -693.722 m 202.931 -693.722 l S Q}%
\TextAt{14}{269.087}{72.921}{o}%
\TextAt{9}{274.931}{72.921}{\Character{41}}%
\TextAt{9}{282.977}{72.921}{\Character{61}}%
\TextAt{20}{295.599}{81.888}{\Character{83}}%
\TextAt{21}{305.562}{78.733}{h}%
\TextAt{21}{305.562}{69.434}{i}%
\TextAt{17}{308.445}{69.434}{\Character{61}1}%
\TextAt{16}{321.756}{72.921}{\Character{66}}%
\TextAt{21}{329.609}{71.128}{i}%
\TextAt{13}{332.99}{72.921}{P}%
\TextAt{9}{342.161}{72.921}{\Character{40}}%
\Vector{q 0 G 0 g 0.478 w 274.713 -693.722 m 280.558 -693.722 l S Q}%
\TextAt{14}{346.713}{72.921}{o}%
\TextAt{9}{352.558}{72.921}{\Character{41}}%
\TextAt{8}{357.11}{72.921}{.}%
\TextAt{8}{365.883}{72.921}{It}%
\TextAt{8}{378.664}{72.921}{therefore}%
\TextAt{8}{427.872}{72.921}{suffices}%
\TextAt{8}{467.768}{72.921}{to}%
\TextAt{8}{482.186}{72.921}{show}%
\TextAt{8}{511.298}{72.921}{that}%
\TextAt{8}{297.211}{43.033}{196}%
\Destination{78.23}{809.815}{nav.page.196}%
\Destination{79.36}{320.655}{nav.print.215}%
\Destination{79.36}{732.255}{nav.number.3.3.5}%
\Destination{79.36}{613.855}{nav.number.3.3.6}%
\Destination{79.36}{431.655}{nav.number.3.3.7}%
\Destination{79.36}{293.855}{nav.number.3.3.8}%
\Hotlink{474.227}{741.855}{25.078}{11.047}{nav.cite.N138}%
\Hotlink{223.443}{623.455}{25.078}{11.047}{nav.number.3.3.4}%
}

% Page 209 (body page 197)
\Page{%
\TextAt{2}{79.226}{796.943}{III.}%
\TextAt{2}{98.317}{796.943}{Riemann–Roch}%
\TextAt{2}{174.375}{796.943}{Theory}%
\Vector{q 0 G 0 g 0.299 w 7.226 18.729 m 460.769 18.729 l S Q}%
\TextAt{16}{79.226}{761.124}{\Character{66}}%
\TextAt{21}{87.078}{759.33}{i}%
\NativeMathAt{95.47}{761.124}{11.955}{\subseteq}%
\TextAt{16}{109.768}{761.124}{\Character{65}}%
\TextAt{21}{119.324}{759.33}{i}%
\TextAt{13}{122.706}{761.124}{P}%
\TextAt{9}{131.877}{761.124}{\Character{40}}%
\Vector{q 0 G 0 g 0.478 w 64.429 -5.52 m 70.273 -5.52 l S Q}%
\TextAt{14}{136.429}{761.124}{o}%
\TextAt{9}{142.273}{761.124}{\Character{41}}%
\TextAt{8}{151.714}{761.124}{for}%
\TextAt{13}{170.591}{761.124}{i}%
\TextAt{9}{179.594}{761.124}{\Character{61}}%
\TextAt{9}{193.701}{761.124}{1}%
\TextAt{13}{199.562}{761.124}{\Character{59}}%
\TextAt{13}{204.786}{761.124}{\Character{58}}%
\TextAt{13}{210.023}{761.124}{\Character{58}}%
\TextAt{13}{215.259}{761.124}{\Character{58}}%
\TextAt{13}{220.484}{761.124}{\Character{59}}%
\TextAt{13}{225.72}{761.124}{h}%
\TextAt{8}{232.521}{761.124}{,}%
\TextAt{8}{240.91}{761.124}{if}%
\TextAt{8}{252.625}{761.124}{the}%
\TextAt{8}{273.773}{761.124}{constant}%
\TextAt{13}{321.987}{761.124}{c}%
\TextAt{8}{331.914}{761.124}{is}%
\TextAt{8}{344.669}{761.124}{chosen}%
\TextAt{8}{383.116}{761.124}{conveniently.}%
\TextAt{8}{456.98}{761.124}{To}%
\TextAt{8}{475.199}{761.124}{do}%
\TextAt{8}{492.449}{761.124}{this,}%
\TextAt{8}{519.762}{761.124}{let}%
\TextAt{14}{79.226}{744.078}{a}%
\TextAt{9}{88.53}{744.078}{\Character{61}}%
\TextAt{14}{100.956}{744.078}{a}%
\TextAt{21}{106.939}{742.284}{i}%
\TextAt{14}{110.32}{744.078}{a}%
\NativeMathAt{116.303}{742.284}{7.97}{\infty}%
\NativeMathAt{128.59}{744.078}{11.955}{\in}%
\TextAt{16}{139.864}{744.078}{\Character{66}}%
\TextAt{21}{147.734}{742.284}{i}%
\TextAt{8}{151.115}{744.078}{,}%
\TextAt{14}{158.264}{744.078}{a}%
\NativeMathAt{164.247}{742.284}{7.97}{\infty}%
\TextAt{9}{176.534}{744.078}{\Character{61}}%
\TextAt{20}{188.96}{753.044}{\Character{81}}%
\TextAt{27}{200.251}{740.591}{p}%
\TextAt{3}{204.24}{740.591}{\Character{106}}%
\NativeMathAt{206.591}{740.591}{7.97}{\infty}%
\TextAt{14}{217.551}{744.078}{p}%
\TextAt{21}{223.534}{748.416}{\Character{23}}%
\TextAt{28}{227.697}{747.42}{p}%
\NativeMathAt{235.598}{744.078}{11.955}{\in}%
\TextAt{20}{246.889}{753.044}{\Character{81}}%
\TextAt{27}{258.18}{740.591}{p}%
\TextAt{3}{262.169}{740.591}{\Character{106}}%
\NativeMathAt{264.52}{740.591}{7.97}{\infty}%
\TextAt{18}{275.481}{744.078}{R}%
\TextAt{3}{284.115}{748.416}{\Character{3}}%
\TextAt{17}{284.115}{741.122}{\Character{43}}%
\TextAt{8}{291.199}{744.078}{.}%
\TextAt{8}{299.648}{744.078}{Then}%
\TextAt{8}{330.205}{744.078}{we}%
\TextAt{8}{347.433}{744.078}{find}%
\TextAt{8}{370.841}{744.078}{for}%
\TextAt{8}{388.726}{744.078}{the}%
\TextAt{8}{408.882}{744.078}{replete}%
\TextAt{8}{447.246}{744.078}{ideal}%
\TextAt{14}{230.332}{713.693}{a}%
\TextAt{3}{236.315}{718.629}{\Character{48}}%
\NativeMathAt{236.315}{710.737}{7.97}{\infty}%
\TextAt{9}{248.602}{713.693}{\Character{61}}%
\TextAt{14}{261.028}{713.693}{a}%
\NativeMathAt{267.011}{711.9}{7.97}{\infty}%
\TextAt{14}{275.977}{713.693}{N}%
\TextAt{9}{285.925}{713.693}{\Character{40}}%
\TextAt{14}{290.478}{713.693}{a}%
\NativeMathAt{296.461}{711.9}{7.97}{\infty}%
\TextAt{9}{305.427}{713.693}{\Character{41}}%
\TextAt{3}{309.98}{718.629}{\Character{0}}%
\TextAt{17}{316.566}{718.629}{1\Character{47}}%
\TextAt{21}{325.034}{718.629}{n}%
\TextAt{9}{333.992}{713.693}{\Character{61}}%
\TextAt{20}{346.417}{725.05}{\Character{89}}%
\TextAt{27}{346.65}{699.147}{p}%
\TextAt{3}{350.639}{699.147}{\Character{106}}%
\NativeMathAt{352.99}{699.147}{7.97}{\infty}%
\TextAt{14}{363.686}{713.693}{p}%
\TextAt{21}{369.669}{718.629}{\Character{23}}%
\TextAt{23}{374.353}{721.442}{\Character{48}}%
\TextAt{28}{373.831}{716.637}{p}%
\TextAt{13}{378.412}{713.693}{\Character{59}}%
\TextAt{8}{79.226}{673.764}{where}%
\TextAt{13}{113.035}{673.764}{\Character{23}}%
\TextAt{3}{119.573}{678.102}{\Character{48}}%
\TextAt{27}{118.819}{670.808}{p}%
\TextAt{9}{126.627}{673.764}{\Character{61}}%
\TextAt{13}{139.052}{673.764}{\Character{23}}%
\TextAt{27}{144.836}{671.97}{p}%
\TextAt{16}{151.98}{673.764}{\Character{0}}%
\TextAt{17}{165.582}{678.471}{1}%
\Vector{q 0 G 0 g 0.478 w 93.13 -97.247 m 98.269 -97.247 l S Q}%
\TextAt{21}{165.13}{669.641}{n}%
\TextAt{20}{173.457}{682.73}{\Character{80}}%
\TextAt{27}{186.076}{670.277}{q}%
\TextAt{3}{189.974}{670.277}{\Character{106}}%
\NativeMathAt{192.325}{670.277}{7.97}{\infty}%
\TextAt{13}{203.286}{673.764}{f}%
\TextAt{27}{209.056}{671.97}{q}%
\TextAt{13}{213.452}{673.764}{\Character{23}}%
\TextAt{27}{219.236}{671.97}{q}%
\TextAt{8}{223.632}{673.764}{,}%
\TextAt{8}{230.781}{673.764}{that}%
\TextAt{14}{256.15}{673.764}{N}%
\TextAt{9}{266.098}{673.764}{\Character{40}}%
\TextAt{14}{270.651}{673.764}{a}%
\TextAt{3}{276.634}{678.102}{\Character{48}}%
\NativeMathAt{276.634}{670.808}{7.97}{\infty}%
\TextAt{9}{285.6}{673.764}{\Character{41}}%
\TextAt{9}{293.455}{673.764}{\Character{61}}%
\TextAt{9}{305.876}{673.764}{1}%
\TextAt{8}{311.752}{673.764}{,}%
\TextAt{8}{318.901}{673.764}{and}%
\TextAt{8}{341.664}{673.764}{thus}%
\TextAt{20}{367.415}{682.73}{\Character{80}}%
\TextAt{27}{380.035}{670.277}{p}%
\TextAt{3}{384.023}{670.277}{\Character{106}}%
\NativeMathAt{386.374}{670.277}{7.97}{\infty}%
\TextAt{13}{397.335}{673.764}{f}%
\TextAt{27}{403.105}{671.97}{p}%
\TextAt{13}{407.592}{673.764}{\Character{23}}%
\TextAt{3}{414.129}{678.102}{\Character{48}}%
\TextAt{27}{413.375}{670.808}{p}%
\TextAt{9}{421.183}{673.764}{\Character{61}}%
\TextAt{9}{433.593}{673.764}{0}%
\TextAt{8}{439.461}{673.764}{.}%
\TextAt{8}{447.91}{673.764}{The}%
\TextAt{8}{471.964}{673.764}{vector}%
\TextAt{9}{250.944}{643.379}{\Character{40}}%
\TextAt{13}{255.496}{643.379}{\Character{58}}%
\TextAt{13}{260.72}{643.379}{\Character{58}}%
\TextAt{13}{265.957}{643.379}{\Character{58}}%
\TextAt{13}{271.193}{643.379}{\Character{59}}%
\TextAt{13}{276.418}{643.379}{f}%
\TextAt{27}{282.243}{641.586}{p}%
\TextAt{13}{286.73}{643.379}{\Character{23}}%
\TextAt{3}{293.268}{648.315}{\Character{48}}%
\TextAt{27}{292.514}{640.423}{p}%
\TextAt{13}{297}{643.379}{\Character{59}}%
\TextAt{13}{302.224}{643.379}{\Character{58}}%
\TextAt{13}{307.461}{643.379}{\Character{58}}%
\TextAt{13}{312.697}{643.379}{\Character{58}}%
\TextAt{9}{315.985}{643.379}{\Character{41}}%
\NativeMathAt{323.858}{643.379}{11.955}{\in}%
\TextAt{20}{335.149}{654.736}{\Character{89}}%
\TextAt{27}{335.382}{628.833}{p}%
\TextAt{3}{339.371}{628.833}{\Character{106}}%
\NativeMathAt{341.722}{628.833}{7.97}{\infty}%
\TextAt{18}{352.417}{643.379}{R}%
\TextAt{8}{79.226}{604.327}{therefore}%
\TextAt{8}{128.913}{604.327}{lies}%
\TextAt{8}{149.717}{604.327}{in}%
\TextAt{8}{163.957}{604.327}{the}%
\TextAt{8}{184.701}{604.327}{trace-zero}%
\TextAt{8}{239.265}{604.327}{space}%
\TextAt{13}{271.127}{604.327}{H}%
\TextAt{9}{286.072}{604.327}{\Character{61}}%
\TextAt{16}{299.498}{604.327}{\Character{102}}%
\TextAt{9}{305.476}{604.327}{\Character{40}}%
\TextAt{13}{310.028}{604.327}{x}%
\TextAt{27}{316.68}{602.533}{p}%
\TextAt{9}{321.167}{604.327}{\Character{41}}%
\NativeMathAt{330.041}{604.327}{11.955}{\in}%
\TextAt{20}{342.333}{613.293}{\Character{81}}%
\TextAt{27}{353.624}{600.84}{p}%
\TextAt{3}{357.612}{600.84}{\Character{106}}%
\NativeMathAt{359.963}{600.84}{7.97}{\infty}%
\TextAt{18}{370.924}{604.327}{R}%
\TextAt{16}{383.88}{604.327}{\Character{106}}%
\TextAt{20}{391.522}{613.293}{\Character{80}}%
\TextAt{27}{404.141}{600.84}{p}%
\TextAt{3}{408.13}{600.84}{\Character{106}}%
\NativeMathAt{410.481}{600.84}{7.97}{\infty}%
\TextAt{13}{421.441}{604.327}{x}%
\TextAt{27}{428.094}{602.533}{p}%
\TextAt{9}{436.902}{604.327}{\Character{61}}%
\TextAt{9}{450.316}{604.327}{0}%
\TextAt{16}{456.181}{604.327}{\Character{103}}%
\TextAt{8}{462.159}{604.327}{.}%
\TextAt{8}{472.369}{604.327}{Inside}%
\TextAt{8}{507.148}{604.327}{it}%
\TextAt{8}{519.439}{604.327}{we}%
\TextAt{8}{79.226}{587.281}{have—see}%
\TextAt{8}{132.203}{587.281}{chap.}%
\TextAt{8}{164.149}{587.281}{I,}%
\TextAt{8}{174.78}{587.281}{(7.3)—the}%
\TextAt{8}{229.981}{587.281}{complete}%
\TextAt{8}{278.666}{587.281}{unit}%
\TextAt{8}{302.64}{587.281}{lattice}%
\TextAt{13}{337.672}{587.281}{\Character{21}}%
\TextAt{9}{344.501}{587.281}{\Character{40}}%
\TextAt{14}{349.053}{587.281}{o}%
\TextAt{3}{354.897}{591.619}{\Character{3}}%
\TextAt{9}{359.63}{587.281}{\Character{41}}%
\TextAt{8}{364.182}{587.281}{.}%
\TextAt{8}{372.384}{587.281}{Thus}%
\TextAt{8}{401.296}{587.281}{there}%
\TextAt{8}{430.47}{587.281}{exists}%
\TextAt{8}{462.047}{587.281}{a}%
\TextAt{8}{471.065}{587.281}{lattice}%
\TextAt{8}{506.097}{587.281}{point}%
\TextAt{13}{79.226}{570.235}{\Character{21}}%
\TextAt{9}{86.055}{570.235}{\Character{40}}%
\TextAt{13}{90.607}{570.235}{u}%
\TextAt{9}{97.269}{570.235}{\Character{41}}%
\TextAt{9}{105.124}{570.235}{\Character{61}}%
\TextAt{9}{117.545}{570.235}{\Character{40}}%
\TextAt{13}{122.12}{570.235}{\Character{58}}%
\TextAt{13}{127.344}{570.235}{\Character{58}}%
\TextAt{13}{132.581}{570.235}{\Character{58}}%
\TextAt{13}{137.817}{570.235}{\Character{59}}%
\TextAt{16}{143.097}{570.235}{\Character{0}}%
\TextAt{13}{152.395}{570.235}{f}%
\TextAt{27}{158.165}{568.442}{p}%
\TextAt{13}{162.652}{570.235}{v}%
\TextAt{27}{168.312}{568.442}{p}%
\TextAt{9}{172.798}{570.235}{\Character{40}}%
\TextAt{13}{177.351}{570.235}{u}%
\TextAt{9}{184.013}{570.235}{\Character{41}}%
\TextAt{13}{188.566}{570.235}{\Character{59}}%
\TextAt{13}{193.79}{570.235}{\Character{58}}%
\TextAt{13}{199.027}{570.235}{\Character{58}}%
\TextAt{13}{204.251}{570.235}{\Character{58}}%
\TextAt{9}{207.55}{570.235}{\Character{41}}%
\TextAt{27}{212.102}{568.375}{p}%
\TextAt{3}{216.091}{568.375}{\Character{106}}%
\NativeMathAt{218.442}{568.375}{7.97}{\infty}%
\TextAt{8}{227.41}{570.235}{,}%
\TextAt{13}{234.559}{570.235}{u}%
\NativeMathAt{244.542}{570.235}{11.955}{\in}%
\TextAt{14}{255.833}{570.235}{o}%
\TextAt{3}{261.677}{574.573}{\Character{3}}%
\TextAt{8}{266.41}{570.235}{,}%
\TextAt{8}{273.559}{570.235}{such}%
\TextAt{8}{299.956}{570.235}{that}%
\TextAt{16}{249.039}{541.833}{\Character{106}}%
\TextAt{13}{252.359}{541.833}{f}%
\TextAt{27}{258.129}{540.04}{p}%
\TextAt{13}{262.616}{541.833}{\Character{23}}%
\TextAt{3}{269.154}{546.769}{\Character{48}}%
\TextAt{27}{268.4}{538.877}{p}%
\TextAt{16}{275.544}{541.833}{\Character{0}}%
\TextAt{13}{287.499}{541.833}{f}%
\TextAt{27}{293.269}{540.04}{p}%
\TextAt{13}{297.756}{541.833}{v}%
\TextAt{27}{303.415}{540.04}{p}%
\TextAt{9}{307.902}{541.833}{\Character{40}}%
\TextAt{13}{312.454}{541.833}{u}%
\TextAt{9}{319.117}{541.833}{\Character{41}}%
\TextAt{16}{323.669}{541.833}{\Character{106}}%
\NativeMathAt{330.292}{541.833}{11.955}{\leq}%
\TextAt{13}{342.93}{541.833}{f}%
\TextAt{27}{348.7}{540.04}{p}%
\TextAt{13}{353.187}{541.833}{c}%
\TextAt{17}{358.225}{540.04}{0}%
\ErrataAt{79.19}{513.431}{455.81}{Choose the lattice bound $c_0$ also to satisfy $c_0\ge n^{-1}\max_i\log\mathrm N(\mathfrak a_i)$. Then}%
\TextAt{13}{116.748}{482.298}{\Character{23}}%
\TextAt{27}{122.532}{480.504}{p}%
\TextAt{16}{129.676}{482.298}{\Character{0}}%
\TextAt{13}{141.631}{482.298}{v}%
\TextAt{27}{147.29}{480.504}{p}%
\TextAt{9}{151.777}{482.298}{\Character{40}}%
\TextAt{13}{156.329}{482.298}{u}%
\TextAt{9}{162.992}{482.298}{\Character{41}}%
\TextAt{9}{170.847}{482.298}{\Character{61}}%
\TextAt{13}{183.29}{482.298}{\Character{23}}%
\TextAt{3}{189.828}{487.234}{\Character{48}}%
\TextAt{27}{189.074}{479.342}{p}%
\TextAt{9}{196.218}{482.298}{\Character{43}}%
\TextAt{9}{209.742}{490.385}{1}%
\Vector{q 0 G 0 g 0.478 w 137.175 -288.713 m 144.162 -288.713 l S Q}%
\TextAt{13}{209.175}{474.097}{n}%
\TextAt{20}{219.35}{493.655}{\Character{88}}%
\TextAt{27}{220.625}{467.752}{q}%
\TextAt{3}{224.523}{467.752}{\Character{106}}%
\NativeMathAt{226.874}{467.752}{7.97}{\infty}%
\TextAt{13}{238.611}{482.298}{f}%
\TextAt{27}{244.381}{480.504}{q}%
\TextAt{13}{248.778}{482.298}{\Character{23}}%
\TextAt{27}{254.562}{480.504}{q}%
\TextAt{16}{261.615}{482.298}{\Character{0}}%
\TextAt{13}{273.57}{482.298}{v}%
\TextAt{27}{279.229}{480.504}{p}%
\TextAt{9}{283.716}{482.298}{\Character{40}}%
\TextAt{13}{288.268}{482.298}{u}%
\TextAt{9}{294.931}{482.298}{\Character{41}}%
\NativeMathAt{302.804}{482.298}{11.955}{\leq}%
\TextAt{9}{317.186}{490.385}{1}%
\Vector{q 0 G 0 g 0.478 w 244.619 -288.713 m 251.606 -288.713 l S Q}%
\TextAt{13}{316.619}{474.097}{n}%
\TextAt{8}{326.794}{482.298}{log}%
\TextAt{14}{343.755}{482.298}{N}%
\TextAt{9}{353.703}{482.298}{\Character{40}}%
\TextAt{14}{358.255}{482.298}{a}%
\NativeMathAt{364.238}{480.504}{7.97}{\infty}%
\TextAt{9}{373.205}{482.298}{\Character{41}}%
\TextAt{9}{380.402}{482.298}{\Character{43}}%
\TextAt{13}{392.175}{482.298}{c}%
\TextAt{17}{397.213}{480.504}{0}%
\TextAt{9}{405.266}{482.298}{\Character{61}}%
\TextAt{9}{419.455}{490.385}{1}%
\Vector{q 0 G 0 g 0.478 w 346.887 -288.713 m 353.875 -288.713 l S Q}%
\TextAt{13}{418.887}{474.097}{n}%
\TextAt{8}{429.063}{482.298}{log}%
\TextAt{14}{446.023}{482.298}{N}%
\TextAt{9}{455.971}{482.298}{\Character{40}}%
\TextAt{14}{460.524}{482.298}{a}%
\TextAt{9}{466.507}{482.298}{\Character{41}}%
\TextAt{9}{473.704}{482.298}{\Character{43}}%
\TextAt{13}{485.477}{482.298}{c}%
\TextAt{17}{490.515}{480.504}{1}%
\TextAt{8}{79.226}{442.368}{with}%
\TextAt{13}{106.66}{442.368}{c}%
\TextAt{17}{111.698}{440.575}{1}%
\TextAt{9}{121.071}{442.368}{\Character{61}}%
\TextAt{13}{134.816}{442.368}{c}%
\TextAt{17}{139.854}{440.575}{0}%
\TextAt{16}{147.771}{442.368}{\Character{0}}%
\TextAt{17}{161.901}{447.076}{1}%
\Vector{q 0 G 0 g 0.478 w 89.449 -328.643 m 94.588 -328.643 l S Q}%
\TextAt{21}{161.449}{438.246}{n}%
\TextAt{8}{169.776}{442.368}{log}%
\TextAt{14}{186.736}{442.368}{N}%
\TextAt{9}{196.684}{442.368}{\Character{40}}%
\TextAt{14}{201.236}{442.368}{a}%
\TextAt{21}{207.219}{440.575}{i}%
\TextAt{9}{210.601}{442.368}{\Character{41}}%
\TextAt{8}{215.153}{442.368}{.}%
\TextAt{8}{225.924}{442.368}{Putting}%
\TextAt{8}{269.785}{442.368}{now}%
\TextAt{14}{294.948}{442.368}{b}%
\TextAt{9}{305.723}{442.368}{\Character{61}}%
\TextAt{14}{319.468}{442.368}{a}%
\TextAt{9}{325.451}{442.368}{\Character{91}}%
\TextAt{13}{328.703}{442.368}{u}%
\TextAt{3}{335.365}{446.707}{\Character{0}}%
\TextAt{17}{341.952}{446.707}{1}%
\TextAt{9}{346.684}{442.368}{\Character{93}}%
\TextAt{9}{354.562}{442.368}{\Character{61}}%
\TextAt{20}{368.321}{451.335}{\Character{81}}%
\TextAt{27}{379.612}{438.881}{p}%
\TextAt{14}{386.091}{442.368}{p}%
\TextAt{21}{392.074}{446.707}{\Character{17}}%
\TextAt{28}{396.324}{445.711}{p}%
\TextAt{8}{400.905}{442.368}{,}%
\TextAt{8}{409.022}{442.368}{we}%
\TextAt{8}{427.023}{442.368}{get}%
\TextAt{14}{447.308}{442.368}{b}%
\TextAt{21}{453.443}{440.575}{f}%
\TextAt{9}{463.528}{442.368}{\Character{61}}%
\TextAt{14}{477.273}{442.368}{a}%
\TextAt{21}{483.257}{440.575}{i}%
\TextAt{8}{486.638}{442.368}{.}%
\TextAt{8}{497.409}{442.368}{This}%
\TextAt{8}{524.903}{442.368}{is}%
\TextAt{8}{79.226}{425.323}{because}%
\TextAt{9}{122.528}{425.323}{\Character{91}}%
\TextAt{13}{125.779}{425.322}{u}%
\TextAt{9}{132.442}{425.323}{\Character{93}}%
\TextAt{21}{135.694}{423.529}{f}%
\TextAt{9}{144.459}{425.322}{\Character{61}}%
\TextAt{9}{156.869}{425.322}{\Character{40}}%
\TextAt{13}{161.437}{425.322}{u}%
\TextAt{9}{168.1}{425.322}{\Character{41}}%
\TextAt{9}{175.955}{425.322}{\Character{61}}%
\TextAt{9}{188.376}{425.322}{\Character{40}1\Character{41}}%
\TextAt{8}{207.253}{425.323}{and}%
\TextAt{13}{199.284}{393.237}{f}%
\TextAt{27}{205.054}{391.444}{p}%
\TextAt{13}{209.541}{393.237}{\Character{17}}%
\TextAt{27}{215.317}{391.444}{p}%
\TextAt{9}{223.124}{393.237}{\Character{61}}%
\TextAt{13}{235.55}{393.237}{f}%
\TextAt{27}{241.32}{391.444}{p}%
\TextAt{9}{245.807}{393.237}{\Character{40}}%
\TextAt{13}{250.359}{393.237}{\Character{23}}%
\TextAt{27}{256.143}{391.444}{p}%
\TextAt{16}{263.286}{393.237}{\Character{0}}%
\TextAt{13}{275.242}{393.237}{v}%
\TextAt{27}{280.901}{391.444}{p}%
\TextAt{9}{285.388}{393.237}{\Character{40}}%
\TextAt{13}{289.94}{393.237}{u}%
\TextAt{9}{296.602}{393.237}{\Character{41}\Character{41}}%
\NativeMathAt{309.028}{393.237}{11.955}{\leq}%
\TextAt{13}{322.843}{401.325}{f}%
\TextAt{27}{328.613}{399.532}{p}%
\Vector{q 0 G 0 g 0.478 w 250.843 -377.774 m 261.1 -377.774 l S Q}%
\TextAt{13}{324.477}{385.037}{n}%
\TextAt{8}{336.288}{393.237}{log}%
\TextAt{14}{353.248}{393.237}{N}%
\TextAt{9}{363.196}{393.237}{\Character{40}}%
\TextAt{14}{367.748}{393.237}{a}%
\TextAt{9}{373.731}{393.237}{\Character{41}}%
\TextAt{9}{380.928}{393.237}{\Character{43}}%
\TextAt{13}{392.702}{393.237}{nc}%
\TextAt{17}{404.727}{391.444}{1}%
\TextAt{13}{409.46}{393.237}{\Character{59}}%
\TextAt{8}{79.226}{362.368}{so}%
\TextAt{8}{93.891}{362.368}{that}%
\TextAt{14}{119.554}{362.368}{N}%
\TextAt{9}{129.502}{362.368}{\Character{40}}%
\TextAt{14}{134.054}{362.368}{p}%
\TextAt{9}{140.037}{362.368}{\Character{41}}%
\TextAt{21}{144.59}{366.707}{\Character{17}}%
\TextAt{28}{148.839}{365.711}{p}%
\NativeMathAt{157.243}{362.368}{11.955}{\leq}%
\TextAt{13}{170.364}{362.368}{e}%
\TextAt{21}{175.789}{366.707}{nc}%
\TextAt{24}{184.595}{365.6}{1}%
\TextAt{14}{189.244}{362.368}{N}%
\TextAt{9}{199.192}{362.368}{\Character{40}}%
\TextAt{14}{203.745}{362.368}{a}%
\TextAt{9}{209.728}{362.368}{\Character{41}}%
\TextAt{21}{214.28}{366.707}{f}%
\TextAt{28}{218.348}{365.711}{p}%
\TextAt{17}{222.431}{366.707}{\Character{47}}%
\TextAt{21}{226.665}{366.707}{n}%
\TextAt{8}{236.493}{362.368}{for}%
\TextAt{14}{254.673}{362.368}{p}%
\TextAt{16}{264.478}{362.368}{\Character{106}}%
\NativeMathAt{271.603}{362.368}{11.955}{\infty}%
\TextAt{8}{283.577}{362.368}{;}%
\TextAt{8}{291.168}{362.368}{then}%
\TextAt{14}{318.122}{362.368}{b}%
\NativeMathAt{328.08}{362.368}{11.955}{\in}%
\TextAt{16}{339.856}{362.368}{\Character{65}}%
\TextAt{21}{349.418}{360.575}{i}%
\TextAt{8}{352.8}{362.368}{,}%
\TextAt{8}{360.317}{362.368}{so}%
\TextAt{8}{374.981}{362.368}{that}%
\TextAt{14}{400.645}{362.368}{a}%
\TextAt{9}{410.45}{362.368}{\Character{61}}%
\TextAt{14}{423.378}{362.368}{b}%
\TextAt{9}{429.512}{362.368}{\Character{91}}%
\TextAt{13}{432.764}{362.368}{u}%
\TextAt{9}{439.426}{362.368}{\Character{93}}%
\NativeMathAt{446.501}{362.368}{11.955}{\in}%
\TextAt{16}{458.277}{362.368}{\Character{65}}%
\TextAt{21}{467.839}{360.575}{i}%
\TextAt{13}{471.221}{362.368}{P}%
\TextAt{9}{480.392}{362.368}{\Character{40}}%
\Vector{q 0 G 0 g 0.478 w 412.944 -404.275 m 418.788 -404.275 l S Q}%
\TextAt{14}{484.944}{362.368}{o}%
\TextAt{9}{490.788}{362.368}{\Character{41}}%
\TextAt{8}{495.341}{362.368}{,}%
\TextAt{8}{502.858}{362.368}{where}%
\TextAt{13}{79.226}{345.322}{c}%
\TextAt{9}{87.585}{345.322}{\Character{61}}%
\TextAt{13}{100.01}{345.322}{e}%
\TextAt{21}{105.436}{349.661}{nc}%
\TextAt{24}{114.242}{348.554}{1}%
\TextAt{8}{118.891}{345.322}{.}%
\Vector{q 0 G 0 g 0.399 w 452.334 -428.678 m 452.334 -420.608 l S Q}%
\Vector{q 0 G 0 g 0.399 w 452.534 -420.807 m 459.707 -420.807 l S Q}%
\Vector{q 0 G 0 g 0.399 w 452.534 -428.478 m 459.707 -428.478 l S Q}%
\Vector{q 0 G 0 g 0.399 w 459.906 -428.678 m 459.906 -420.608 l S Q}%
\TextAt{19}{79.226}{316.493}{Proof}%
\TextAt{19}{110.62}{316.493}{of}%
\TextAt{19}{124.369}{316.493}{(3.7).}%
\TextAt{8}{159.039}{316.493}{As}%
\TextAt{13}{176.326}{316.493}{O}%
\TextAt{9}{185.597}{316.493}{\Character{40}}%
\TextAt{13}{190.15}{316.493}{t}%
\TextAt{9}{194.377}{316.493}{\Character{41}}%
\TextAt{9}{202.232}{316.493}{\Character{61}}%
\TextAt{13}{214.675}{316.493}{O}%
\TextAt{9}{223.947}{316.493}{\Character{40}}%
\TextAt{13}{228.499}{316.493}{t}%
\TextAt{9}{232.726}{316.493}{\Character{41}}%
\TextAt{16}{239.935}{316.493}{\Character{0}}%
\TextAt{9}{251.89}{316.493}{1}%
\TextAt{8}{257.743}{316.493}{,}%
\TextAt{8}{264.892}{316.493}{we}%
\TextAt{8}{282.12}{316.493}{may}%
\TextAt{8}{307.489}{316.493}{replace}%
\TextAt{13}{347.156}{316.493}{H}%
\TextAt{17}{357.779}{320.831}{0}%
\TextAt{9}{362.511}{316.493}{\Character{40}}%
\TextAt{14}{367.063}{316.493}{a}%
\TextAt{3}{373.046}{320.831}{\Character{0}}%
\TextAt{17}{379.633}{320.831}{1}%
\TextAt{9}{384.365}{316.493}{\Character{41}}%
\TextAt{8}{392.815}{316.493}{by}%
\Vector{q 0 G 0 g 0.478 w 70.091 -476.066 m 80.713 -476.066 l S Q}%
\TextAt{13}{142.091}{288.091}{H}%
\TextAt{17}{152.713}{295.496}{0}%
\TextAt{9}{157.445}{288.091}{\Character{40}}%
\TextAt{14}{161.998}{288.091}{a}%
\TextAt{3}{167.981}{293.027}{\Character{0}}%
\TextAt{17}{174.567}{293.027}{1}%
\TextAt{9}{179.3}{288.091}{\Character{41}}%
\TextAt{9}{187.155}{288.091}{\Character{61}}%
\TextAt{13}{199.598}{288.091}{H}%
\TextAt{17}{210.221}{293.027}{0}%
\TextAt{9}{214.953}{288.091}{\Character{40}}%
\TextAt{14}{219.506}{288.091}{a}%
\TextAt{3}{225.489}{293.027}{\Character{0}}%
\TextAt{17}{232.075}{293.027}{1}%
\TextAt{9}{236.808}{288.091}{\Character{41}}%
\TextAt{16}{244.017}{288.091}{\Character{91}}%
\TextAt{16}{254.633}{288.091}{\Character{102}}%
\TextAt{9}{260.621}{288.091}{0}%
\TextAt{16}{266.474}{288.091}{\Character{103}}%
\TextAt{9}{275.772}{288.091}{\Character{61}}%
\TextAt{16}{288.198}{288.091}{\Character{102}}%
\TextAt{13}{294.175}{288.091}{f}%
\NativeMathAt{304.543}{288.091}{11.955}{\in}%
\TextAt{14}{315.834}{288.091}{a}%
\TextAt{3}{321.817}{293.039}{\Character{0}}%
\TextAt{17}{328.403}{293.039}{1}%
\TextAt{21}{321.817}{284.761}{f}%
\TextAt{16}{336.456}{288.091}{\Character{106}}%
\TextAt{16}{343.079}{288.091}{\Character{106}}%
\TextAt{13}{346.419}{288.091}{f}%
\TextAt{16}{353.465}{288.091}{\Character{106}}%
\TextAt{27}{356.786}{286.297}{p}%
\NativeMathAt{364.594}{288.091}{11.955}{\leq}%
\TextAt{14}{377.213}{288.091}{N}%
\TextAt{9}{387.162}{288.091}{\Character{40}}%
\TextAt{14}{391.714}{288.091}{p}%
\TextAt{9}{397.697}{288.091}{\Character{41}}%
\TextAt{21}{402.249}{293.027}{\Character{23}}%
\TextAt{28}{406.412}{292.031}{p}%
\TextAt{8}{414.89}{288.091}{for}%
\TextAt{14}{432.775}{288.091}{p}%
\TextAt{16}{442.079}{288.091}{\Character{106}}%
\NativeMathAt{448.702}{288.091}{11.955}{\infty}%
\TextAt{16}{460.657}{288.091}{\Character{103}}%
\TextAt{13}{466.654}{288.091}{\Character{58}}%
\TextAt{8}{79.226}{259.689}{We}%
\TextAt{8}{99.371}{259.689}{have}%
\TextAt{8}{126.365}{259.689}{to}%
\TextAt{8}{140.676}{259.689}{show}%
\TextAt{8}{169.679}{259.689}{that}%
\TextAt{8}{195.048}{259.689}{there}%
\TextAt{8}{224.96}{259.689}{are}%
\TextAt{8}{244.47}{259.689}{constants}%
\TextAt{13}{296.308}{259.689}{C}%
\TextAt{13}{304.88}{259.689}{\Character{59}}%
\TextAt{13}{310.116}{259.689}{C}%
\TextAt{3}{319.369}{264.027}{\Character{48}}%
\TextAt{8}{326.062}{259.689}{such}%
\TextAt{8}{352.459}{259.689}{that}%
\TextAt{20}{199.243}{245.012}{\Character{12}}%
\TextAt{20}{199.243}{237.839}{\Character{12}}%
\TextAt{20}{199.243}{230.665}{\Character{12}}%
\TextAt{20}{199.243}{223.492}{\Character{12}}%
\TextAt{20}{199.243}{216.319}{\Character{12}}%
\TextAt{8}{203.228}{224.09}{\char35\relax }%
\Vector{q 0 G 0 g 0.478 w 140.984 -540.067 m 151.606 -540.067 l S Q}%
\TextAt{13}{212.984}{224.09}{H}%
\TextAt{17}{223.606}{231.496}{0}%
\TextAt{9}{228.339}{224.09}{\Character{40}}%
\TextAt{14}{232.891}{224.09}{a}%
\TextAt{3}{238.874}{229.026}{\Character{0}}%
\TextAt{17}{245.46}{229.026}{1}%
\TextAt{9}{250.193}{224.09}{\Character{41}}%
\TextAt{16}{257.402}{224.09}{\Character{0}}%
\TextAt{9}{270.552}{232.178}{2}%
\TextAt{21}{276.405}{236.516}{r}%
\TextAt{9}{280.96}{232.178}{\Character{40}2}%
\TextAt{13}{291.365}{232.178}{\Character{25}}%
\TextAt{9}{298.435}{232.178}{\Character{41}}%
\TextAt{21}{302.987}{236.516}{s}%
\Vector{q 0 G 0 g 0.478 w 198.552 -546.921 m 235.401 -546.921 l S Q}%
\TextAt{20}{272.536}{224.449}{\Character{112}}%
\Vector{q 0 G 0 g 0.478 w 212.491 -549.312 m 233.418 -549.312 l S Q}%
\TextAt{16}{284.491}{214.227}{\Character{106}}%
\TextAt{13}{287.812}{214.227}{d}%
\TextAt{21}{293.894}{212.434}{K}%
\TextAt{16}{302.097}{214.227}{\Character{106}}%
\TextAt{14}{308.597}{224.09}{N}%
\TextAt{9}{318.545}{224.09}{\Character{40}}%
\TextAt{14}{323.097}{224.09}{a}%
\TextAt{9}{329.08}{224.09}{\Character{41}}%
\TextAt{20}{333.632}{245.012}{\Character{12}}%
\TextAt{20}{333.632}{237.839}{\Character{12}}%
\TextAt{20}{333.632}{230.665}{\Character{12}}%
\TextAt{20}{333.632}{223.492}{\Character{12}}%
\TextAt{20}{333.632}{216.319}{\Character{12}}%
\NativeMathAt{340.938}{224.09}{11.955}{\leq}%
\TextAt{13}{353.558}{224.09}{C}%
\TextAt{14}{362.791}{224.09}{N}%
\TextAt{9}{372.739}{224.09}{\Character{40}}%
\TextAt{14}{377.292}{224.09}{a}%
\TextAt{9}{383.275}{224.09}{\Character{41}}%
\TextAt{17}{387.827}{229.026}{1}%
\TextAt{3}{392.061}{229.026}{\Character{0}}%
\TextAt{17}{398.648}{229.026}{1\Character{47}}%
\TextAt{21}{407.116}{229.026}{n}%
\TextAt{8}{517.802}{224.09}{(*)}%
\TextAt{8}{79.226}{184.252}{for}%
\TextAt{8}{97.19}{184.252}{all}%
\TextAt{14}{113.527}{184.252}{a}%
\NativeMathAt{122.965}{184.252}{11.955}{\in}%
\TextAt{13}{134.39}{184.252}{J}%
\TextAt{9}{141.967}{184.252}{\Character{40}}%
\Vector{q 0 G 0 g 0.478 w 74.519 -582.392 m 80.363 -582.392 l S Q}%
\TextAt{14}{146.519}{184.252}{o}%
\TextAt{9}{152.363}{184.252}{\Character{41}}%
\TextAt{8}{160.892}{184.252}{satisfying}%
\TextAt{14}{213.131}{184.252}{N}%
\TextAt{9}{223.079}{184.252}{\Character{40}}%
\TextAt{14}{227.631}{184.252}{a}%
\TextAt{9}{233.614}{184.252}{\Character{41}}%
\NativeMathAt{241.621}{184.252}{11.955}{\geq}%
\TextAt{13}{254.374}{184.252}{C}%
\TextAt{3}{263.608}{188.59}{\Character{48}}%
\TextAt{8}{266.403}{184.252}{.}%
\TextAt{8}{275.087}{184.252}{For}%
\TextAt{13}{296.135}{184.252}{a}%
\NativeMathAt{305.734}{184.252}{11.955}{\in}%
\TextAt{13}{317.159}{184.252}{K}%
\TextAt{3}{327.977}{188.59}{\Character{3}}%
\TextAt{8}{332.709}{184.252}{,}%
\TextAt{8}{339.956}{184.252}{the}%
\TextAt{8}{360.191}{184.252}{set}%
\Vector{q 0 G 0 g 0.478 w 306.537 -579.905 m 317.16 -579.905 l S Q}%
\TextAt{13}{378.537}{184.252}{H}%
\TextAt{17}{389.16}{191.657}{0}%
\TextAt{9}{393.892}{184.252}{\Character{40}}%
\TextAt{14}{398.445}{184.252}{a}%
\TextAt{3}{404.428}{188.59}{\Character{0}}%
\TextAt{17}{411.014}{188.59}{1}%
\TextAt{9}{415.747}{184.252}{\Character{41}}%
\TextAt{8}{424.275}{184.252}{is}%
\TextAt{8}{436.117}{184.252}{mapped}%
\TextAt{8}{480.741}{184.252}{bijectively}%
\TextAt{8}{79.226}{167.206}{via}%
\TextAt{13}{98.27}{167.206}{x}%
\NativeArrowAt{108.243}{167.206}{11.955}{11.955}{\mapsto}%
\TextAt{13}{123.519}{167.206}{ax}%
\TextAt{8}{140.07}{167.206}{onto}%
\TextAt{8}{166.275}{167.206}{the}%
\TextAt{8}{186.287}{167.206}{set}%
\Vector{q 0 G 0 g 0.478 w 132.411 -596.951 m 143.033 -596.951 l S Q}%
\TextAt{13}{204.411}{167.206}{H}%
\TextAt{17}{215.033}{174.611}{0}%
\TextAt{9}{219.766}{167.206}{\Character{40}\Character{91}}%
\TextAt{13}{227.57}{167.206}{a}%
\TextAt{9}{233.715}{167.206}{\Character{93}}%
\TextAt{14}{236.966}{167.206}{a}%
\TextAt{3}{242.949}{171.544}{\Character{0}}%
\TextAt{17}{249.536}{171.544}{1}%
\TextAt{9}{254.268}{167.206}{\Character{41}}%
\TextAt{8}{258.821}{167.206}{.}%
\TextAt{8}{267.221}{167.206}{The}%
\TextAt{8}{291.131}{167.206}{numbers}%
\TextAt{8}{338.198}{167.206}{\char35\relax }%
\Vector{q 0 G 0 g 0.478 w 275.953 -596.951 m 286.576 -596.951 l S Q}%
\TextAt{13}{347.953}{167.206}{H}%
\TextAt{17}{358.576}{174.611}{0}%
\TextAt{9}{363.308}{167.206}{\Character{40}}%
\TextAt{14}{367.86}{167.206}{a}%
\TextAt{3}{373.844}{171.544}{\Character{0}}%
\TextAt{17}{380.43}{171.544}{1}%
\TextAt{9}{385.162}{167.206}{\Character{41}}%
\TextAt{8}{393.468}{167.206}{and}%
\TextAt{14}{416.087}{167.206}{N}%
\TextAt{9}{426.035}{167.206}{\Character{40}}%
\TextAt{14}{430.587}{167.206}{a}%
\TextAt{9}{436.57}{167.206}{\Character{41}}%
\TextAt{8}{444.876}{167.206}{thus}%
\TextAt{8}{470.484}{167.206}{depend}%
\TextAt{8}{510.975}{167.206}{only}%
\TextAt{8}{79.226}{150.16}{on}%
\TextAt{8}{95.924}{150.16}{the}%
\TextAt{8}{116.52}{150.16}{class}%
\TextAt{14}{144.396}{150.16}{a}%
\TextAt{8}{153.7}{150.16}{mod}%
\TextAt{13}{179.138}{150.16}{P}%
\TextAt{9}{188.309}{150.16}{\Character{40}}%
\Vector{q 0 G 0 g 0.478 w 120.862 -616.484 m 126.706 -616.484 l S Q}%
\TextAt{14}{192.862}{150.16}{o}%
\TextAt{9}{198.706}{150.16}{\Character{41}}%
\TextAt{8}{203.258}{150.16}{.}%
\TextAt{8}{213.024}{150.16}{As}%
\TextAt{8}{230.75}{150.16}{by}%
\TextAt{8}{247.449}{150.16}{the}%
\TextAt{8}{268.044}{150.16}{preceding}%
\TextAt{8}{321.158}{150.16}{lemma}%
\TextAt{13}{359.316}{150.16}{J}%
\TextAt{9}{366.893}{150.16}{\Character{40}}%
\Vector{q 0 G 0 g 0.478 w 299.445 -616.484 m 305.289 -616.484 l S Q}%
\TextAt{14}{371.445}{150.16}{o}%
\TextAt{9}{377.289}{150.16}{\Character{41}}%
\TextAt{9}{385.897}{150.16}{\Character{61}}%
\TextAt{20}{399.085}{159.126}{\Character{83}}%
\TextAt{21}{409.047}{155.972}{h}%
\TextAt{21}{409.047}{146.673}{i}%
\TextAt{17}{411.931}{146.673}{\Character{61}1}%
\TextAt{16}{425.242}{150.16}{\Character{65}}%
\TextAt{21}{434.788}{148.367}{i}%
\TextAt{13}{438.169}{150.16}{P}%
\TextAt{9}{447.34}{150.16}{\Character{40}}%
\Vector{q 0 G 0 g 0.478 w 379.893 -616.484 m 385.737 -616.484 l S Q}%
\TextAt{14}{451.893}{150.16}{o}%
\TextAt{9}{457.737}{150.16}{\Character{41}}%
\TextAt{8}{462.289}{150.16}{,}%
\TextAt{8}{469.987}{150.16}{it}%
\TextAt{8}{482.131}{150.16}{suffices}%
\TextAt{8}{522.357}{150.16}{to}%
\TextAt{8}{79.226}{133.114}{show}%
\TextAt{8}{108.229}{133.114}{(*)}%
\TextAt{8}{127.095}{133.114}{for}%
\TextAt{14}{144.98}{133.114}{a}%
\TextAt{8}{154.86}{133.114}{ranging}%
\TextAt{8}{197.145}{133.114}{over}%
\TextAt{8}{222.192}{133.114}{the}%
\TextAt{8}{242.348}{133.114}{set}%
\TextAt{16}{260.615}{133.114}{\Character{65}}%
\TextAt{21}{270.161}{131.321}{i}%
\TextAt{8}{273.543}{133.114}{.}%
\TextAt{8}{97.159}{113.079}{For}%
\TextAt{8}{118.128}{113.079}{this,}%
\TextAt{8}{144.202}{113.079}{we}%
\TextAt{8}{161.43}{113.079}{shall}%
\TextAt{8}{188.807}{113.079}{use}%
\TextAt{8}{209.023}{113.079}{the}%
\TextAt{8}{229.18}{113.079}{identification}%
\TextAt{8}{299.751}{113.079}{of}%
\TextAt{8}{313.081}{113.079}{Minkowski}%
\TextAt{8}{371.351}{113.079}{space}%
\TextAt{13}{274.161}{84.677}{K}%
\TextAt{25}{284.12}{82.884}{R}%
\TextAt{9}{293.695}{84.677}{\Character{61}}%
\TextAt{20}{306.12}{96.035}{\Character{89}}%
\TextAt{27}{306.354}{70.132}{p}%
\TextAt{3}{310.342}{70.132}{\Character{106}}%
\NativeMathAt{312.693}{70.132}{7.97}{\infty}%
\TextAt{13}{323.389}{84.677}{K}%
\TextAt{27}{333.348}{82.884}{p}%
\TextAt{8}{297.211}{43.033}{197}%
\Destination{522.36}{241.135}{nav.AMS.26}%
\Destination{199.24}{241.135}{nav.equation.3.3.1}%
\Destination{78.23}{809.815}{nav.page.197}%
\Destination{79.36}{199.255}{nav.print.216}%
\Hotlink{163.649}{585.055}{5.22}{11.047}{nav.chapter.1}%
\Hotlink{174.28}{585.055}{25.078}{11.047}{nav.number.1.7.3}%
\Hotlink{123.869}{314.255}{26.106}{11.047}{nav.number.3.3.7}%
\Destination{79.226}{530.431}{errata.99}%
}

% Page 210 (body page 198)
\Page{%
\TextAt{2}{79.226}{796.943}{III.}%
\TextAt{2}{98.317}{796.943}{Riemann–Roch}%
\TextAt{2}{174.375}{796.943}{Theory}%
\Vector{q 0 G 0 g 0.299 w 7.226 18.729 m 460.769 18.729 l S Q}%
\TextAt{8}{79.226}{761.124}{with}%
\TextAt{8}{105.886}{761.124}{its}%
\TextAt{8}{122.205}{761.124}{canonical}%
\TextAt{8}{173.588}{761.124}{measure.}%
\TextAt{8}{223.724}{761.124}{Since}%
\TextAt{14}{254.282}{761.124}{a}%
\TextAt{21}{260.265}{759.33}{f}%
\TextAt{9}{269.031}{761.124}{\Character{61}}%
\TextAt{14}{281.456}{761.124}{a}%
\TextAt{21}{287.439}{759.33}{i}%
\TextAt{8}{294.718}{761.124}{for}%
\TextAt{14}{312.603}{761.124}{a}%
\TextAt{9}{321.907}{761.124}{\Character{61}}%
\TextAt{20}{334.332}{770.09}{\Character{81}}%
\TextAt{27}{345.623}{757.637}{p}%
\TextAt{14}{352.102}{761.124}{p}%
\TextAt{21}{358.086}{765.462}{\Character{23}}%
\TextAt{28}{362.248}{764.466}{p}%
\NativeMathAt{370.15}{761.124}{11.955}{\in}%
\TextAt{16}{381.424}{761.124}{\Character{65}}%
\TextAt{21}{390.987}{759.33}{i}%
\TextAt{8}{394.368}{761.124}{,}%
\TextAt{8}{401.517}{761.124}{we}%
\TextAt{8}{418.745}{761.124}{have}%
\Vector{q 0 G 0 g 0.478 w 114.39 -32.86 m 125.013 -32.86 l S Q}%
\TextAt{13}{186.39}{731.297}{H}%
\TextAt{17}{197.013}{738.703}{0}%
\TextAt{9}{201.745}{731.297}{\Character{40}}%
\TextAt{14}{206.298}{731.297}{a}%
\TextAt{3}{212.281}{736.233}{\Character{0}}%
\TextAt{17}{218.867}{736.233}{1}%
\TextAt{9}{223.599}{731.297}{\Character{41}}%
\TextAt{9}{231.454}{731.297}{\Character{61}}%
\TextAt{16}{243.898}{731.297}{\Character{102}}%
\TextAt{13}{249.876}{731.297}{f}%
\NativeMathAt{260.243}{731.297}{11.955}{\in}%
\TextAt{14}{271.534}{731.297}{a}%
\TextAt{3}{277.517}{736.245}{\Character{0}}%
\TextAt{17}{284.104}{736.245}{1}%
\TextAt{21}{277.517}{728.233}{i}%
\TextAt{16}{292.157}{731.297}{\Character{106}}%
\TextAt{16}{298.78}{731.297}{\Character{106}}%
\TextAt{13}{302.119}{731.297}{f}%
\TextAt{16}{309.166}{731.297}{\Character{106}}%
\TextAt{27}{312.487}{729.504}{p}%
\NativeMathAt{320.294}{731.297}{11.955}{\leq}%
\TextAt{14}{332.914}{731.297}{N}%
\TextAt{9}{342.862}{731.297}{\Character{40}}%
\TextAt{14}{347.414}{731.297}{p}%
\TextAt{9}{353.397}{731.297}{\Character{41}}%
\TextAt{21}{357.949}{736.233}{\Character{23}}%
\TextAt{28}{362.112}{735.237}{p}%
\TextAt{8}{370.59}{731.297}{for}%
\TextAt{14}{388.475}{731.297}{p}%
\TextAt{16}{397.779}{731.297}{\Character{106}}%
\NativeMathAt{404.402}{731.297}{11.955}{\infty}%
\TextAt{16}{416.357}{731.297}{\Character{103}}%
\TextAt{13}{422.354}{731.297}{\Character{58}}%
\TextAt{8}{79.226}{702.957}{We}%
\TextAt{8}{98.596}{702.957}{therefore}%
\TextAt{8}{146.921}{702.957}{have}%
\TextAt{8}{173.142}{702.957}{to}%
\TextAt{8}{186.677}{702.957}{count}%
\TextAt{8}{218.098}{702.957}{the}%
\TextAt{8}{237.48}{702.957}{lattice}%
\TextAt{8}{272.475}{702.957}{points}%
\TextAt{8}{306.885}{702.957}{in}%
\TextAt{9}{319.763}{702.957}{\Character{0}}%
\TextAt{9}{330.379}{702.957}{\Character{61}}%
\TextAt{13}{342.824}{702.957}{j}%
\TextAt{14}{348.339}{702.957}{a}%
\TextAt{3}{354.322}{707.905}{\Character{0}}%
\TextAt{17}{360.909}{707.905}{1}%
\TextAt{21}{354.322}{699.893}{i}%
\NativeMathAt{368.962}{702.957}{11.955}{\subseteq}%
\TextAt{13}{381.581}{702.957}{K}%
\TextAt{25}{391.54}{701.163}{R}%
\TextAt{8}{400.918}{702.957}{which}%
\TextAt{8}{433.63}{702.957}{fall}%
\TextAt{8}{452.689}{702.957}{into}%
\TextAt{8}{475.657}{702.957}{the}%
\TextAt{8}{495.039}{702.957}{domain}%
\TextAt{13}{276.371}{674.616}{P}%
\TextAt{27}{283.916}{672.823}{a}%
\TextAt{9}{291.724}{674.616}{\Character{61}}%
\TextAt{20}{304.15}{685.973}{\Character{89}}%
\TextAt{27}{304.383}{660.07}{p}%
\TextAt{3}{308.372}{660.07}{\Character{106}}%
\NativeMathAt{310.723}{660.07}{7.97}{\infty}%
\TextAt{13}{321.418}{674.616}{D}%
\TextAt{27}{331.138}{672.823}{p}%
\TextAt{8}{79.226}{635.625}{where}%
\TextAt{13}{113.312}{635.625}{D}%
\TextAt{27}{123.031}{633.832}{p}%
\TextAt{9}{131.31}{635.625}{\Character{61}}%
\TextAt{16}{144.207}{635.625}{\Character{102}}%
\TextAt{13}{150.184}{635.625}{x}%
\NativeMathAt{160.628}{635.625}{11.955}{\in}%
\TextAt{13}{172.391}{635.625}{K}%
\TextAt{27}{182.35}{633.832}{p}%
\TextAt{16}{190.628}{635.625}{\Character{106}}%
\TextAt{16}{197.729}{635.625}{\Character{106}}%
\TextAt{13}{201.062}{635.625}{x}%
\TextAt{16}{207.714}{635.625}{\Character{106}}%
\TextAt{27}{211.035}{633.832}{p}%
\NativeMathAt{219.314}{635.625}{11.955}{\leq}%
\TextAt{14}{232.405}{635.625}{N}%
\TextAt{9}{242.353}{635.625}{\Character{40}}%
\TextAt{14}{246.905}{635.625}{p}%
\TextAt{9}{252.888}{635.625}{\Character{41}}%
\TextAt{21}{257.44}{639.963}{\Character{23}}%
\TextAt{28}{261.603}{638.967}{p}%
\TextAt{16}{266.184}{635.625}{\Character{103}}%
\TextAt{8}{272.162}{635.625}{.}%
\TextAt{8}{281.439}{635.625}{Let}%
\TextAt{13}{302.685}{635.625}{F}%
\TextAt{8}{316.063}{635.625}{be}%
\TextAt{8}{332.264}{635.625}{a}%
\TextAt{8}{342.295}{635.625}{fundamental}%
\TextAt{8}{410.214}{635.625}{mesh}%
\TextAt{8}{440.462}{635.625}{of}%
\TextAt{9}{454.069}{635.625}{\Character{0}}%
\TextAt{8}{461.383}{635.625}{.}%
\TextAt{8}{470.661}{635.625}{We}%
\TextAt{8}{491.082}{635.625}{consider}%
\TextAt{8}{79.226}{618.579}{the}%
\TextAt{8}{99.383}{618.579}{sets}%
\TextAt{13}{227.911}{590.238}{X}%
\TextAt{9}{241.887}{590.238}{\Character{61}}%
\TextAt{16}{254.312}{590.238}{\Character{102}}%
\TextAt{13}{260.29}{590.238}{\Character{13}}%
\NativeMathAt{270.333}{590.238}{11.955}{\in}%
\TextAt{9}{281.624}{590.238}{\Character{0}}%
\TextAt{16}{292.259}{590.238}{\Character{106}}%
\TextAt{9}{298.901}{590.238}{\Character{40}}%
\TextAt{13}{303.453}{590.238}{F}%
\TextAt{9}{315.313}{590.238}{\Character{43}}%
\TextAt{13}{327.075}{590.238}{\Character{13}}%
\TextAt{9}{333.797}{590.238}{\Character{41}}%
\TextAt{16}{341.006}{590.238}{\Character{92}}%
\TextAt{13}{351.633}{590.238}{P}%
\TextAt{27}{359.178}{588.445}{a}%
\NativeMathAt{366.986}{590.238}{11.955}{\neq}%
\TextAt{18}{379.411}{590.238}{\Character{63}}%
\TextAt{16}{388.71}{590.238}{\Character{103}}%
\TextAt{13}{394.687}{590.238}{\Character{59}}%
\TextAt{13}{229.173}{570.204}{Y}%
\TextAt{9}{241.887}{570.204}{\Character{61}}%
\TextAt{16}{254.312}{570.204}{\Character{102}}%
\TextAt{13}{260.29}{570.204}{\Character{13}}%
\NativeMathAt{270.333}{570.204}{11.955}{\in}%
\TextAt{9}{281.624}{570.204}{\Character{0}}%
\TextAt{16}{292.259}{570.204}{\Character{106}}%
\TextAt{13}{298.901}{570.204}{F}%
\TextAt{9}{310.761}{570.204}{\Character{43}}%
\TextAt{13}{322.522}{570.204}{\Character{13}}%
\NativeMathAt{332.565}{570.204}{11.955}{\subseteq}%
\TextAt{9}{346.358}{573.226}{\Character{23}}%
\TextAt{13}{345.185}{570.204}{P}%
\TextAt{27}{352.73}{568.41}{a}%
\TextAt{16}{357.217}{570.204}{\Character{103}}%
\TextAt{13}{363.195}{570.204}{\Character{59}}%
\TextAt{13}{207.227}{550.169}{X}%
\TextAt{16}{220.539}{550.169}{\Character{110}}%
\TextAt{13}{229.173}{550.169}{Y}%
\TextAt{9}{241.887}{550.169}{\Character{61}}%
\TextAt{16}{254.312}{550.169}{\Character{102}}%
\TextAt{13}{260.29}{550.169}{\Character{13}}%
\NativeMathAt{270.333}{550.169}{11.955}{\in}%
\TextAt{9}{281.624}{550.169}{\Character{0}}%
\TextAt{16}{292.259}{550.169}{\Character{106}}%
\TextAt{9}{298.901}{550.169}{\Character{40}}%
\TextAt{13}{303.453}{550.169}{F}%
\TextAt{9}{315.313}{550.169}{\Character{43}}%
\TextAt{13}{327.075}{550.169}{\Character{13}}%
\TextAt{9}{333.797}{550.169}{\Character{41}}%
\TextAt{16}{341.006}{550.169}{\Character{92}}%
\TextAt{13}{351.633}{550.169}{\Character{64}}%
\TextAt{13}{358.447}{550.169}{P}%
\TextAt{27}{366.008}{548.376}{a}%
\NativeMathAt{373.816}{550.169}{11.955}{\neq}%
\TextAt{18}{386.241}{550.169}{\Character{63}}%
\TextAt{16}{395.54}{550.169}{\Character{103}}%
\TextAt{13}{401.517}{550.169}{\Character{58}}%
\TextAt{8}{79.226}{520.222}{As}%
\TextAt{13}{96.513}{520.222}{Y}%
\NativeMathAt{109.227}{520.222}{11.955}{\subseteq}%
\TextAt{9}{121.847}{520.222}{\Character{0}}%
\TextAt{16}{131.817}{520.222}{\Character{92}}%
\TextAt{13}{142.444}{520.222}{P}%
\TextAt{27}{149.99}{518.429}{a}%
\TextAt{9}{157.797}{520.222}{\Character{61}}%
\Vector{q 0 G 0 g 0.478 w 98.223 -243.935 m 108.845 -243.935 l S Q}%
\TextAt{13}{170.223}{520.222}{H}%
\TextAt{17}{180.845}{527.628}{0}%
\TextAt{9}{185.578}{520.222}{\Character{40}}%
\TextAt{14}{190.13}{520.222}{a}%
\TextAt{3}{196.113}{524.561}{\Character{0}}%
\TextAt{17}{202.7}{524.561}{1}%
\TextAt{9}{207.432}{520.222}{\Character{41}}%
\NativeMathAt{215.305}{520.222}{11.955}{\subseteq}%
\TextAt{13}{227.924}{520.222}{X}%
\TextAt{8}{242.477}{520.222}{and}%
\TextAt{8}{265.24}{520.222}{as}%
\TextAt{20}{279.61}{529.189}{\Character{83}}%
\TextAt{21}{289.572}{516.735}{\Character{13}}%
\NativeMathAt{294.404}{516.735}{7.97}{\in}%
\TextAt{21}{300.05}{516.735}{Y}%
\TextAt{9}{307.338}{520.222}{\Character{40}}%
\TextAt{13}{311.89}{520.222}{F}%
\TextAt{9}{323.751}{520.222}{\Character{43}}%
\TextAt{13}{335.512}{520.222}{\Character{13}}%
\TextAt{9}{342.234}{520.222}{\Character{41}}%
\NativeMathAt{350.107}{520.222}{11.955}{\subseteq}%
\TextAt{13}{362.727}{520.222}{P}%
\TextAt{27}{370.272}{518.429}{a}%
\NativeMathAt{378.08}{520.222}{11.955}{\subseteq}%
\TextAt{20}{390.699}{529.189}{\Character{83}}%
\TextAt{21}{400.662}{516.735}{\Character{13}}%
\NativeMathAt{405.494}{516.735}{7.97}{\in}%
\TextAt{21}{411.139}{516.735}{X}%
\TextAt{9}{419.224}{520.222}{\Character{40}}%
\TextAt{13}{423.777}{520.222}{F}%
\TextAt{9}{435.637}{520.222}{\Character{43}}%
\TextAt{13}{447.398}{520.222}{\Character{13}}%
\TextAt{9}{454.121}{520.222}{\Character{41}}%
\TextAt{8}{458.673}{520.222}{,}%
\TextAt{8}{465.822}{520.222}{one}%
\TextAt{8}{487.282}{520.222}{has}%
\TextAt{8}{244.52}{490.353}{\char35\relax }%
\TextAt{13}{254.275}{490.353}{Y}%
\NativeMathAt{266.989}{490.353}{11.955}{\leq}%
\TextAt{8}{279.608}{490.353}{\char35\relax }%
\Vector{q 0 G 0 g 0.478 w 217.364 -273.804 m 227.986 -273.804 l S Q}%
\TextAt{13}{289.364}{490.353}{H}%
\TextAt{17}{299.986}{497.758}{0}%
\TextAt{9}{304.719}{490.353}{\Character{40}}%
\TextAt{14}{309.271}{490.353}{a}%
\TextAt{3}{315.254}{495.289}{\Character{0}}%
\TextAt{17}{321.841}{495.289}{1}%
\TextAt{9}{326.573}{490.353}{\Character{41}}%
\NativeMathAt{334.446}{490.353}{11.955}{\leq}%
\TextAt{8}{347.066}{490.353}{\char35\relax }%
\TextAt{13}{356.821}{490.353}{X}%
\TextAt{8}{79.226}{462.012}{as}%
\TextAt{8}{93.596}{462.012}{well}%
\TextAt{8}{117.327}{462.012}{as}%
\TextAt{8}{214.847}{444.966}{\char35\relax }%
\TextAt{13}{224.603}{444.966}{Y}%
\TextAt{8}{235.988}{444.966}{vol}%
\TextAt{9}{251.279}{444.966}{\Character{40}}%
\TextAt{13}{255.831}{444.966}{F}%
\TextAt{9}{265.035}{444.966}{\Character{41}}%
\NativeMathAt{272.908}{444.966}{11.955}{\leq}%
\TextAt{8}{285.527}{444.966}{vol}%
\TextAt{9}{300.818}{444.966}{\Character{40}}%
\TextAt{13}{305.37}{444.966}{P}%
\TextAt{27}{312.916}{443.173}{a}%
\TextAt{9}{317.402}{444.966}{\Character{41}}%
\NativeMathAt{325.276}{444.966}{11.955}{\leq}%
\TextAt{8}{337.895}{444.966}{\char35\relax }%
\TextAt{13}{347.65}{444.966}{X}%
\TextAt{8}{360.298}{444.966}{vol}%
\TextAt{9}{375.589}{444.966}{\Character{40}}%
\TextAt{13}{380.141}{444.966}{F}%
\TextAt{9}{389.345}{444.966}{\Character{41}}%
\TextAt{13}{393.897}{444.966}{\Character{58}}%
\TextAt{8}{79.226}{421.728}{This}%
\TextAt{8}{105.946}{421.728}{implies}%
\TextAt{20}{179.885}{418.412}{\Character{12}}%
\TextAt{20}{179.885}{411.239}{\Character{12}}%
\TextAt{20}{179.885}{404.066}{\Character{12}}%
\TextAt{20}{179.885}{396.893}{\Character{12}}%
\TextAt{8}{183.87}{401.077}{\char35\relax }%
\Vector{q 0 G 0 g 0.478 w 121.626 -363.08 m 132.248 -363.08 l S Q}%
\TextAt{13}{193.626}{401.077}{H}%
\TextAt{17}{204.248}{408.482}{0}%
\TextAt{9}{208.98}{401.077}{\Character{40}}%
\TextAt{14}{213.533}{401.077}{a}%
\TextAt{3}{219.516}{406.013}{\Character{0}}%
\TextAt{17}{226.102}{406.013}{1}%
\TextAt{9}{230.835}{401.077}{\Character{41}}%
\TextAt{16}{238.044}{401.077}{\Character{0}}%
\TextAt{8}{251.194}{409.165}{vol}%
\TextAt{9}{266.485}{409.165}{\Character{40}}%
\TextAt{13}{271.037}{409.165}{P}%
\TextAt{27}{278.583}{407.371}{a}%
\TextAt{9}{283.069}{409.165}{\Character{41}}%
\Vector{q 0 G 0 g 0.478 w 179.194 -369.934 m 215.622 -369.934 l S Q}%
\TextAt{8}{252.609}{392.876}{vol}%
\TextAt{9}{267.899}{392.876}{\Character{40}}%
\TextAt{13}{272.452}{392.876}{F}%
\TextAt{9}{281.655}{392.876}{\Character{41}}%
\TextAt{20}{288.817}{418.412}{\Character{12}}%
\TextAt{20}{288.817}{411.239}{\Character{12}}%
\TextAt{20}{288.817}{404.066}{\Character{12}}%
\TextAt{20}{288.817}{396.893}{\Character{12}}%
\NativeMathAt{296.123}{401.077}{11.955}{\leq}%
\TextAt{8}{308.742}{401.077}{\char35\relax }%
\TextAt{13}{318.498}{401.077}{X}%
\TextAt{16}{331.81}{401.077}{\Character{0}}%
\TextAt{8}{343.765}{401.077}{\char35\relax }%
\TextAt{13}{353.52}{401.077}{Y}%
\TextAt{9}{366.234}{401.077}{\Character{61}}%
\TextAt{8}{378.66}{401.077}{\char35\relax }%
\TextAt{9}{388.415}{401.077}{\Character{40}}%
\TextAt{13}{392.967}{401.077}{X}%
\TextAt{16}{406.279}{401.077}{\Character{110}}%
\TextAt{13}{414.914}{401.077}{Y}%
\TextAt{9}{424.307}{401.077}{\Character{41}}%
\TextAt{13}{428.859}{401.077}{\Character{58}}%
\TextAt{8}{79.226}{373.564}{For}%
\TextAt{8}{100.196}{373.564}{the}%
\TextAt{8}{120.352}{373.564}{set}%
\TextAt{13}{138.619}{373.564}{P}%
\TextAt{27}{146.165}{371.771}{a}%
\TextAt{9}{153.972}{373.564}{\Character{61}}%
\TextAt{20}{166.398}{382.53}{\Character{81}}%
\TextAt{27}{177.689}{370.077}{p}%
\TextAt{3}{181.678}{370.077}{\Character{106}}%
\NativeMathAt{184.029}{370.077}{7.97}{\infty}%
\TextAt{13}{194.989}{373.564}{D}%
\TextAt{27}{204.708}{371.771}{p}%
\TextAt{8}{209.195}{373.564}{,}%
\TextAt{8}{216.344}{373.564}{we}%
\TextAt{8}{233.572}{373.564}{now}%
\TextAt{8}{257.96}{373.564}{have}%
\TextAt{8}{161.657}{343.241}{vol}%
\TextAt{9}{176.948}{343.241}{\Character{40}}%
\TextAt{13}{181.5}{343.241}{P}%
\TextAt{27}{189.045}{341.447}{a}%
\TextAt{9}{193.532}{343.241}{\Character{41}}%
\TextAt{9}{201.387}{343.241}{\Character{61}}%
\TextAt{20}{216.42}{354.598}{\Character{89}}%
\TextAt{27}{213.831}{329.141}{p}%
\TextAt{7}{220.641}{329.141}{real}%
\TextAt{9}{236.278}{343.241}{2}%
\TextAt{14}{242.131}{343.241}{N}%
\TextAt{9}{252.079}{343.241}{\Character{40}}%
\TextAt{14}{256.631}{343.241}{p}%
\TextAt{9}{262.615}{343.241}{\Character{41}}%
\TextAt{21}{267.167}{348.177}{\Character{23}}%
\TextAt{28}{271.33}{347.18}{p}%
\TextAt{20}{288.837}{354.598}{\Character{89}}%
\TextAt{27}{277.903}{329.141}{p}%
\TextAt{7}{284.713}{329.141}{complex}%
\TextAt{9}{317.04}{343.241}{2}%
\TextAt{13}{322.893}{343.241}{\Character{25}}%
\TextAt{14}{329.962}{343.241}{N}%
\TextAt{9}{339.91}{343.241}{\Character{40}}%
\TextAt{14}{344.463}{343.241}{p}%
\TextAt{9}{350.446}{343.241}{\Character{41}}%
\TextAt{21}{354.998}{348.177}{\Character{23}}%
\TextAt{28}{359.161}{347.18}{p}%
\TextAt{9}{367.062}{343.241}{\Character{61}}%
\TextAt{9}{379.471}{343.241}{2}%
\TextAt{21}{385.341}{348.177}{r}%
\TextAt{9}{389.896}{343.241}{\Character{40}2}%
\TextAt{13}{400.301}{343.241}{\Character{25}}%
\TextAt{9}{407.37}{343.241}{\Character{41}}%
\TextAt{21}{411.922}{348.177}{s}%
\TextAt{14}{416.336}{343.241}{N}%
\TextAt{9}{426.285}{343.241}{\Character{40}}%
\TextAt{14}{430.837}{343.241}{a}%
\NativeMathAt{436.82}{341.447}{7.97}{\infty}%
\TextAt{9}{445.786}{343.241}{\Character{41}}%
\TextAt{8}{79.226}{304.634}{(observe}%
\TextAt{8}{126.222}{304.634}{here}%
\TextAt{8}{152.332}{304.634}{that,}%
\TextAt{8}{181.894}{304.634}{under}%
\TextAt{8}{215.811}{304.634}{the}%
\TextAt{8}{236.721}{304.634}{identification}%
\TextAt{13}{308.046}{304.634}{K}%
\TextAt{27}{318.005}{302.841}{p}%
\TextAt{9}{327.096}{304.634}{\Character{61}}%
\TextAt{18}{340.805}{304.634}{C}%
\TextAt{8}{349.44}{304.634}{,}%
\TextAt{8}{357.531}{304.634}{one}%
\TextAt{8}{379.743}{304.634}{has}%
\TextAt{8}{401.37}{304.634}{the}%
\TextAt{8}{422.28}{304.634}{equation}%
\TextAt{16}{470.842}{304.634}{\Character{106}}%
\TextAt{13}{474.163}{304.634}{x}%
\TextAt{16}{480.815}{304.634}{\Character{106}}%
\TextAt{27}{484.136}{302.841}{p}%
\TextAt{9}{493.227}{304.634}{\Character{61}}%
\TextAt{16}{506.937}{304.634}{\Character{106}}%
\TextAt{13}{510.257}{304.634}{x}%
\TextAt{16}{516.91}{304.634}{\Character{106}}%
\TextAt{17}{520.23}{308.972}{2}%
\TextAt{8}{524.963}{304.634}{).}%
\TextAt{8}{79.226}{287.588}{For}%
\TextAt{8}{100.196}{287.588}{the}%
\TextAt{8}{120.352}{287.588}{fundamental}%
\TextAt{8}{187.994}{287.588}{mesh}%
\TextAt{13}{217.966}{287.588}{F}%
\TextAt{8}{227.17}{287.588}{,}%
\TextAt{8}{234.319}{287.588}{(3.3)}%
\TextAt{8}{262.294}{287.588}{yields}%
\TextAt{8}{245.081}{259.247}{vol}%
\TextAt{9}{260.372}{259.247}{\Character{40}}%
\TextAt{13}{264.924}{259.247}{F}%
\TextAt{9}{274.127}{259.247}{\Character{41}}%
\TextAt{9}{281.982}{259.247}{\Character{61}}%
\TextAt{20}{294.426}{270.053}{\Character{112}}%
\Vector{q 0 G 0 g 0.478 w 234.381 -503.708 m 255.308 -503.708 l S Q}%
\TextAt{16}{306.381}{259.247}{\Character{106}}%
\TextAt{13}{309.702}{259.247}{d}%
\TextAt{21}{315.785}{257.454}{K}%
\TextAt{16}{323.988}{259.247}{\Character{106}}%
\TextAt{14}{327.308}{259.247}{N}%
\TextAt{9}{337.257}{259.247}{\Character{40}}%
\TextAt{14}{341.809}{259.247}{a}%
\TextAt{3}{347.792}{264.195}{\Character{0}}%
\TextAt{17}{354.378}{264.195}{1}%
\TextAt{21}{347.792}{255.918}{f}%
\TextAt{9}{359.111}{259.247}{\Character{41}}%
\TextAt{13}{363.663}{259.247}{\Character{58}}%
\TextAt{8}{79.226}{230.907}{From}%
\TextAt{8}{109.951}{230.907}{this}%
\TextAt{8}{132.773}{230.907}{we}%
\TextAt{8}{150.001}{230.907}{get}%
\TextAt{20}{202.115}{227.46}{\Character{12}}%
\TextAt{20}{202.115}{220.286}{\Character{12}}%
\TextAt{20}{202.115}{213.113}{\Character{12}}%
\TextAt{20}{202.115}{205.94}{\Character{12}}%
\TextAt{20}{202.115}{198.767}{\Character{12}}%
\TextAt{8}{206.1}{206.538}{\char35\relax }%
\Vector{q 0 G 0 g 0.478 w 143.856 -557.619 m 154.478 -557.619 l S Q}%
\TextAt{13}{215.856}{206.538}{H}%
\TextAt{17}{226.478}{213.943}{0}%
\TextAt{9}{231.211}{206.538}{\Character{40}}%
\TextAt{14}{235.763}{206.538}{a}%
\TextAt{3}{241.746}{211.474}{\Character{0}}%
\TextAt{17}{248.332}{211.474}{1}%
\TextAt{9}{253.065}{206.538}{\Character{41}}%
\TextAt{16}{260.274}{206.538}{\Character{0}}%
\TextAt{9}{273.424}{214.626}{2}%
\TextAt{21}{279.277}{218.964}{r}%
\TextAt{9}{283.832}{214.626}{\Character{40}2}%
\TextAt{13}{294.237}{214.626}{\Character{25}}%
\TextAt{9}{301.307}{214.626}{\Character{41}}%
\TextAt{21}{305.859}{218.964}{s}%
\Vector{q 0 G 0 g 0.478 w 201.424 -564.473 m 238.273 -564.473 l S Q}%
\TextAt{20}{275.408}{206.897}{\Character{112}}%
\Vector{q 0 G 0 g 0.478 w 215.363 -566.864 m 236.29 -566.864 l S Q}%
\TextAt{16}{287.363}{196.675}{\Character{106}}%
\TextAt{13}{290.684}{196.675}{d}%
\TextAt{21}{296.766}{194.882}{K}%
\TextAt{16}{304.969}{196.675}{\Character{106}}%
\TextAt{14}{311.469}{206.538}{N}%
\TextAt{9}{321.417}{206.538}{\Character{40}}%
\TextAt{14}{325.969}{206.538}{a}%
\TextAt{9}{331.952}{206.538}{\Character{41}}%
\TextAt{20}{336.504}{227.46}{\Character{12}}%
\TextAt{20}{336.504}{220.286}{\Character{12}}%
\TextAt{20}{336.504}{213.113}{\Character{12}}%
\TextAt{20}{336.504}{205.94}{\Character{12}}%
\TextAt{20}{336.504}{198.767}{\Character{12}}%
\NativeMathAt{343.81}{206.538}{11.955}{\leq}%
\TextAt{8}{356.43}{206.538}{\char35\relax }%
\TextAt{9}{366.185}{206.538}{\Character{40}}%
\TextAt{13}{370.737}{206.538}{X}%
\TextAt{16}{384.049}{206.538}{\Character{110}}%
\TextAt{13}{392.683}{206.538}{Y}%
\TextAt{9}{402.077}{206.538}{\Character{41}}%
\TextAt{13}{406.629}{206.538}{\Character{58}}%
\TextAt{8}{79.226}{175.639}{Having}%
\TextAt{8}{119.54}{175.639}{obtained}%
\TextAt{8}{167.983}{175.639}{this}%
\TextAt{8}{191.117}{175.639}{inequality,}%
\TextAt{8}{248.09}{175.639}{it}%
\TextAt{8}{260.105}{175.639}{suffices}%
\TextAt{8}{300.204}{175.639}{to}%
\TextAt{8}{314.826}{175.639}{show}%
\TextAt{8}{344.14}{175.639}{that}%
\TextAt{8}{369.821}{175.639}{there}%
\TextAt{8}{400.044}{175.639}{exist}%
\TextAt{8}{428.056}{175.639}{constants}%
\TextAt{13}{480.205}{175.639}{C}%
\TextAt{13}{488.777}{175.639}{\Character{59}}%
\TextAt{13}{494.013}{175.639}{C}%
\TextAt{3}{503.266}{179.977}{\Character{48}}%
\TextAt{8}{510.27}{175.639}{such}%
\TextAt{8}{79.226}{158.593}{that}%
\TextAt{8}{155.352}{141.547}{\char35\relax }%
\TextAt{9}{165.107}{141.547}{\Character{40}}%
\TextAt{13}{169.659}{141.547}{X}%
\TextAt{16}{182.971}{141.547}{\Character{110}}%
\TextAt{13}{191.605}{141.547}{Y}%
\TextAt{9}{200.999}{141.547}{\Character{41}}%
\TextAt{9}{208.854}{141.547}{\Character{61}}%
\TextAt{8}{221.297}{141.547}{\char35\relax }%
\TextAt{16}{231.053}{141.547}{\Character{102}}%
\TextAt{13}{237.03}{141.547}{\Character{13}}%
\NativeMathAt{247.073}{141.547}{11.955}{\in}%
\TextAt{9}{258.364}{141.547}{\Character{0}}%
\TextAt{16}{268.999}{141.547}{\Character{106}}%
\TextAt{9}{275.641}{141.547}{\Character{40}}%
\TextAt{13}{280.193}{141.547}{F}%
\TextAt{9}{292.054}{141.547}{\Character{43}}%
\TextAt{13}{303.815}{141.547}{\Character{13}}%
\TextAt{9}{310.537}{141.547}{\Character{41}}%
\TextAt{16}{317.746}{141.547}{\Character{92}}%
\TextAt{13}{328.373}{141.547}{\Character{64}}%
\TextAt{13}{335.188}{141.547}{P}%
\TextAt{27}{342.748}{139.753}{a}%
\NativeMathAt{350.556}{141.547}{11.955}{\neq}%
\TextAt{18}{362.981}{141.547}{\Character{63}}%
\TextAt{16}{372.28}{141.547}{\Character{103}}%
\NativeMathAt{381.569}{141.547}{11.955}{\leq}%
\TextAt{13}{394.198}{141.547}{C}%
\TextAt{14}{403.431}{141.547}{N}%
\TextAt{9}{413.379}{141.547}{\Character{40}}%
\TextAt{14}{417.932}{141.547}{a}%
\TextAt{9}{423.915}{141.547}{\Character{41}}%
\TextAt{17}{428.467}{146.483}{1}%
\TextAt{3}{432.701}{146.483}{\Character{0}}%
\TextAt{17}{439.288}{146.483}{1\Character{47}}%
\TextAt{21}{447.756}{146.483}{n}%
\TextAt{13}{453.392}{141.547}{\Character{59}}%
\TextAt{8}{79.226}{118.308}{for}%
\TextAt{8}{96.217}{118.308}{all}%
\TextAt{14}{111.583}{118.308}{a}%
\NativeMathAt{120.887}{118.308}{11.955}{\in}%
\TextAt{16}{132.161}{118.308}{\Character{65}}%
\TextAt{21}{141.724}{116.515}{i}%
\TextAt{8}{148.108}{118.308}{with}%
\TextAt{14}{173.875}{118.308}{N}%
\TextAt{9}{183.823}{118.308}{\Character{40}}%
\TextAt{14}{188.375}{118.308}{a}%
\TextAt{9}{194.358}{118.308}{\Character{41}}%
\NativeMathAt{202.232}{118.308}{11.955}{\geq}%
\TextAt{13}{214.851}{118.308}{C}%
\TextAt{3}{224.084}{122.646}{\Character{48}}%
\TextAt{8}{226.88}{118.308}{.}%
\TextAt{8}{235.03}{118.308}{We}%
\TextAt{8}{254.281}{118.308}{choose}%
\TextAt{13}{290.52}{118.308}{C}%
\TextAt{3}{299.753}{122.646}{\Character{48}}%
\TextAt{9}{305.869}{118.308}{\Character{61}}%
\TextAt{9}{318.279}{118.308}{1}%
\TextAt{8}{327.152}{118.308}{and}%
\TextAt{8}{349.02}{118.308}{find}%
\TextAt{8}{371.535}{118.308}{the}%
\TextAt{8}{390.798}{118.308}{constant}%
\TextAt{13}{437.127}{118.308}{C}%
\TextAt{8}{449.364}{118.308}{in}%
\TextAt{8}{462.123}{118.308}{the}%
\TextAt{8}{481.386}{118.308}{remainder}%
\TextAt{8}{79.226}{101.262}{of}%
\TextAt{8}{92.556}{101.262}{the}%
\TextAt{8}{112.713}{101.262}{proof.}%
\TextAt{8}{147.833}{101.262}{We}%
\TextAt{8}{167.977}{101.262}{parametrize}%
\TextAt{8}{232.368}{101.262}{the}%
\TextAt{8}{252.524}{101.262}{set}%
\TextAt{13}{270.792}{101.262}{P}%
\TextAt{27}{278.337}{99.469}{a}%
\TextAt{9}{286.145}{101.262}{\Character{61}}%
\TextAt{20}{298.57}{110.228}{\Character{81}}%
\TextAt{27}{309.861}{97.775}{p}%
\TextAt{3}{313.85}{97.775}{\Character{106}}%
\NativeMathAt{316.201}{97.775}{7.97}{\infty}%
\TextAt{13}{327.161}{101.262}{D}%
\TextAt{27}{336.881}{99.469}{p}%
\TextAt{8}{345.265}{101.262}{via}%
\TextAt{8}{364.453}{101.262}{the}%
\TextAt{8}{384.609}{101.262}{mapping}%
\TextAt{13}{270.754}{72.921}{\Character{39}}%
\TextAt{9}{281.743}{72.921}{\Character{58}}%
\TextAt{13}{288.316}{72.921}{I}%
\TextAt{21}{294.419}{77.858}{n}%
\NativeArrowAt{303.376}{72.921}{11.955}{19.26}{\longrightarrow}%
\TextAt{13}{325.958}{72.921}{P}%
\TextAt{27}{333.503}{71.128}{a}%
\TextAt{13}{337.99}{72.921}{\Character{59}}%
\TextAt{8}{297.211}{43.033}{198}%
\Destination{78.23}{809.815}{nav.page.198}%
\Destination{79.36}{319.655}{nav.print.217}%
\Hotlink{233.819}{285.355}{25.078}{11.047}{nav.number.3.3.3}%
}

% Page 211 (body page 199)
\Page{%
\TextAt{2}{79.226}{796.943}{III.}%
\TextAt{2}{98.317}{796.943}{Riemann–Roch}%
\TextAt{2}{174.375}{796.943}{Theory}%
\Vector{q 0 G 0 g 0.299 w 7.226 18.729 m 460.769 18.729 l S Q}%
\TextAt{8}{79.226}{761.124}{where}%
\TextAt{13}{113.035}{761.124}{I}%
\TextAt{9}{122.459}{761.124}{\Character{61}}%
\TextAt{9}{134.869}{761.124}{\Character{91}0}%
\TextAt{13}{143.989}{761.124}{\Character{59}}%
\TextAt{9}{149.233}{761.124}{1\Character{93}}%
\TextAt{8}{158.338}{761.124}{,}%
\TextAt{8}{165.487}{761.124}{which}%
\TextAt{8}{198.974}{761.124}{is}%
\TextAt{8}{210.737}{761.124}{given}%
\TextAt{8}{241.307}{761.124}{by}%
\TextAt{13}{122.579}{734.919}{I}%
\NativeArrowAt{132.002}{734.919}{11.955}{19.26}{\longrightarrow}%
\TextAt{13}{154.584}{734.919}{D}%
\TextAt{27}{164.304}{733.126}{p}%
\TextAt{13}{168.791}{734.919}{\Character{59}}%
\TextAt{13}{232.602}{734.919}{t}%
\NativeArrowAt{240.15}{734.919}{11.955}{19.26}{\longmapsto}%
\TextAt{9}{262.732}{734.919}{2}%
\TextAt{13}{268.585}{734.919}{\Character{11}}%
\TextAt{27}{276.031}{733.126}{p}%
\TextAt{20}{280.518}{744.603}{\Character{0}}%
\TextAt{13}{285.997}{734.919}{t}%
\TextAt{16}{292.881}{734.919}{\Character{0}}%
\TextAt{17}{306.032}{739.626}{1}%
\Vector{q 0 G 0 g 0.478 w 234.032 -36.092 m 238.266 -36.092 l S Q}%
\TextAt{17}{306.032}{730.797}{2}%
\TextAt{20}{311.462}{744.603}{\Character{1}}%
\TextAt{13}{316.941}{734.919}{\Character{59}}%
\TextAt{8}{428.342}{734.919}{if}%
\TextAt{14}{439.066}{734.919}{p}%
\TextAt{8}{448.947}{734.919}{real}%
\TextAt{13}{467.812}{734.919}{\Character{59}}%
\TextAt{13}{117.847}{714.884}{I}%
\TextAt{17}{123.949}{719.821}{2}%
\NativeArrowAt{132.002}{714.884}{11.955}{19.26}{\longrightarrow}%
\TextAt{13}{154.584}{714.884}{D}%
\TextAt{27}{164.304}{713.091}{p}%
\TextAt{13}{168.791}{714.884}{\Character{59}}%
\TextAt{9}{210.663}{714.884}{\Character{40}}%
\TextAt{13}{215.215}{714.884}{\Character{26}\Character{59}}%
\TextAt{13}{226.477}{714.884}{\Character{18}}%
\TextAt{9}{232.277}{714.884}{\Character{41}}%
\NativeArrowAt{240.15}{714.884}{11.955}{19.26}{\longmapsto}%
\TextAt{16}{262.732}{722.429}{\Character{112}}%
\Vector{q 0 G 0 g 0.478 w 200.695 -51.332 m 212.628 -51.332 l S Q}%
\TextAt{13}{272.695}{714.884}{\Character{11}}%
\TextAt{27}{280.141}{713.091}{p}%
\TextAt{9}{284.628}{714.884}{\Character{40}}%
\TextAt{13}{289.18}{714.884}{\Character{26}}%
\TextAt{8}{297.21}{714.884}{cos}%
\TextAt{9}{314.876}{714.884}{2}%
\TextAt{13}{320.729}{714.884}{\Character{25}}%
\TextAt{13}{327.783}{714.884}{\Character{18}}%
\TextAt{13}{333.569}{714.884}{\Character{59}}%
\TextAt{13}{338.793}{714.884}{\Character{26}}%
\TextAt{8}{346.853}{714.884}{sin}%
\TextAt{9}{363.215}{714.884}{2}%
\TextAt{13}{369.068}{714.884}{\Character{25}}%
\TextAt{13}{376.122}{714.884}{\Character{18}}%
\TextAt{9}{381.918}{714.884}{\Character{41}}%
\TextAt{13}{386.47}{714.884}{\Character{59}}%
\TextAt{8}{428.342}{714.884}{if}%
\TextAt{14}{439.066}{714.884}{p}%
\TextAt{8}{448.947}{714.884}{complex}%
\TextAt{13}{490.897}{714.884}{\Character{59}}%
\TextAt{8}{79.226}{688.68}{where}%
\TextAt{13}{114.107}{688.68}{\Character{11}}%
\TextAt{27}{121.552}{686.887}{p}%
\TextAt{9}{131.186}{688.68}{\Character{61}}%
\TextAt{14}{145.437}{688.68}{N}%
\TextAt{9}{155.385}{688.68}{\Character{40}}%
\TextAt{14}{159.938}{688.68}{p}%
\TextAt{9}{165.921}{688.68}{\Character{41}}%
\TextAt{21}{170.473}{693.018}{\Character{23}}%
\TextAt{28}{174.636}{692.022}{p}%
\TextAt{8}{179.217}{688.68}{.}%
\TextAt{8}{190.879}{688.68}{We}%
\TextAt{8}{212.095}{688.68}{bound}%
\TextAt{8}{249.259}{688.68}{the}%
\TextAt{8}{270.487}{688.68}{norm}%
\TextAt{16}{302.128}{688.68}{\Character{107}}%
\TextAt{13}{308.105}{688.68}{d\Character{39}}%
\TextAt{9}{321.856}{688.68}{\Character{40}}%
\TextAt{13}{326.409}{688.68}{x}%
\TextAt{9}{333.061}{688.68}{\Character{41}}%
\TextAt{16}{337.613}{688.68}{\Character{107}}%
\TextAt{8}{348.56}{688.68}{of}%
\TextAt{8}{362.961}{688.68}{the}%
\TextAt{8}{384.189}{688.68}{derivative}%
\TextAt{13}{438.927}{688.68}{d\Character{39}}%
\TextAt{9}{452.678}{688.68}{\Character{40}}%
\TextAt{13}{457.23}{688.68}{x}%
\TextAt{9}{463.882}{688.68}{\Character{41}}%
\TextAt{9}{473.566}{688.68}{\Character{58}}%
\TextAt{18}{481.98}{688.68}{R}%
\TextAt{21}{490.614}{693.018}{n}%
\NativeMathAt{501.397}{688.68}{11.955}{\rightarrow}%
\TextAt{18}{518.499}{688.68}{R}%
\TextAt{21}{527.133}{693.018}{n}%
\TextAt{8}{79.226}{671.634}{(}%
\TextAt{13}{83.781}{671.634}{x}%
\NativeMathAt{94.572}{671.634}{11.955}{\in}%
\TextAt{13}{106.682}{671.634}{I}%
\TextAt{21}{112.784}{675.972}{n}%
\TextAt{8}{118.421}{671.634}{).}%
\TextAt{8}{132.865}{671.634}{If}%
\TextAt{13}{145.037}{671.634}{d\Character{39}}%
\TextAt{9}{158.788}{671.634}{\Character{40}}%
\TextAt{13}{163.34}{671.634}{x}%
\TextAt{9}{169.992}{671.634}{\Character{41}}%
\TextAt{9}{178.672}{671.634}{\Character{61}}%
\TextAt{9}{191.906}{671.634}{\Character{40}}%
\TextAt{13}{196.48}{671.634}{a}%
\TextAt{21}{202.625}{669.841}{ik}%
\TextAt{9}{210.628}{671.634}{\Character{41}}%
\TextAt{8}{215.18}{671.634}{,}%
\TextAt{8}{222.93}{671.634}{then}%
\TextAt{16}{250.07}{671.634}{\Character{107}}%
\TextAt{13}{256.048}{671.634}{d\Character{39}}%
\TextAt{9}{269.799}{671.634}{\Character{40}}%
\TextAt{13}{274.351}{671.634}{x}%
\TextAt{9}{281.003}{671.634}{\Character{41}}%
\TextAt{16}{285.555}{671.634}{\Character{107}}%
\NativeMathAt{295.669}{671.634}{11.955}{\leq}%
\TextAt{13}{309.11}{671.634}{n}%
\TextAt{8}{318.09}{671.634}{max}%
\TextAt{16}{341.877}{671.634}{\Character{106}}%
\TextAt{13}{345.198}{671.634}{a}%
\TextAt{21}{351.342}{669.841}{ik}%
\TextAt{16}{359.345}{671.634}{\Character{106}}%
\TextAt{8}{362.666}{671.634}{.}%
\TextAt{8}{372.555}{671.634}{Every}%
\TextAt{8}{406.689}{671.634}{partial}%
\TextAt{8}{444.9}{671.634}{derivative}%
\TextAt{8}{499.047}{671.634}{of}%
\TextAt{13}{512.857}{671.634}{\Character{39}}%
\TextAt{8}{524.903}{671.634}{is}%
\TextAt{8}{79.226}{654.588}{now}%
\TextAt{8}{103.4}{654.588}{bounded}%
\TextAt{8}{150.982}{654.588}{by}%
\TextAt{9}{167.027}{654.588}{2}%
\TextAt{13}{172.88}{654.588}{\Character{11}}%
\TextAt{27}{180.325}{652.795}{p}%
\TextAt{8}{184.812}{654.588}{,}%
\TextAt{8}{191.79}{654.588}{resp.}%
\TextAt{9}{221.04}{654.588}{2}%
\TextAt{13}{226.893}{654.588}{\Character{25}}%
\TextAt{16}{233.963}{661.549}{\Character{112}}%
\Vector{q 0 G 0 g 0.478 w 171.925 -112.212 m 183.858 -112.212 l S Q}%
\TextAt{13}{243.925}{654.588}{\Character{11}}%
\TextAt{27}{251.371}{652.795}{p}%
\TextAt{8}{255.858}{654.588}{.}%
\TextAt{8}{264.234}{654.588}{Since}%
\TextAt{14}{294.577}{654.588}{a}%
\NativeMathAt{303.881}{654.588}{11.955}{\in}%
\TextAt{16}{315.155}{654.588}{\Character{65}}%
\TextAt{21}{324.718}{652.795}{i}%
\TextAt{8}{328.099}{654.588}{,}%
\TextAt{8}{335.077}{654.588}{we}%
\TextAt{8}{352.09}{654.588}{have}%
\TextAt{8}{378.87}{654.588}{that}%
\TextAt{13}{404.024}{654.588}{\Character{11}}%
\TextAt{27}{411.47}{652.795}{p}%
\TextAt{9}{419.278}{654.588}{\Character{61}}%
\TextAt{14}{431.703}{654.588}{N}%
\TextAt{9}{441.651}{654.588}{\Character{40}}%
\TextAt{14}{446.203}{654.588}{p}%
\TextAt{9}{452.187}{654.588}{\Character{41}}%
\TextAt{21}{456.739}{658.926}{\Character{23}}%
\TextAt{28}{460.902}{657.93}{p}%
\NativeMathAt{468.803}{654.588}{11.955}{\leq}%
\TextAt{13}{481.423}{654.588}{c}%
\TextAt{14}{486.461}{654.588}{N}%
\TextAt{9}{496.409}{654.588}{\Character{40}}%
\TextAt{14}{500.961}{654.588}{a}%
\TextAt{9}{506.944}{654.588}{\Character{41}}%
\TextAt{21}{511.496}{658.926}{f}%
\TextAt{28}{515.564}{657.93}{p}%
\TextAt{17}{519.647}{658.926}{\Character{47}}%
\TextAt{21}{523.881}{658.926}{n}%
\TextAt{8}{529.518}{654.588}{,}%
\TextAt{8}{79.226}{637.542}{for}%
\TextAt{8}{97.111}{637.542}{all}%
\TextAt{14}{113.37}{637.542}{p}%
\TextAt{16}{122.674}{637.542}{\Character{106}}%
\NativeMathAt{129.297}{637.542}{11.955}{\infty}%
\TextAt{8}{141.271}{637.542}{.}%
\TextAt{8}{149.719}{637.542}{It}%
\TextAt{8}{162.392}{637.542}{follows}%
\TextAt{8}{200.828}{637.542}{that}%
\TextAt{16}{208.081}{611.338}{\Character{107}}%
\TextAt{13}{214.058}{611.337}{d\Character{39}}%
\TextAt{9}{227.81}{611.338}{\Character{40}}%
\TextAt{13}{232.362}{611.337}{x}%
\TextAt{9}{239.014}{611.338}{\Character{41}}%
\TextAt{16}{243.566}{611.337}{\Character{107}}%
\NativeMathAt{252.855}{611.337}{11.955}{\leq}%
\TextAt{9}{265.484}{611.338}{2}%
\TextAt{13}{271.337}{611.337}{\Character{25}}%
\TextAt{13}{278.391}{611.337}{n}%
\TextAt{8}{287.386}{611.338}{max}%
\TextAt{13}{311.173}{611.337}{\Character{11}}%
\TextAt{17}{318.695}{617.583}{1\Character{47}}%
\TextAt{21}{327.163}{617.583}{f}%
\TextAt{28}{331.231}{616.586}{p}%
\TextAt{27}{318.619}{609.367}{p}%
\NativeMathAt{339.133}{611.338}{11.955}{\leq}%
\TextAt{13}{351.752}{611.337}{c}%
\TextAt{17}{356.79}{609.544}{1}%
\TextAt{14}{361.523}{611.338}{N}%
\TextAt{9}{371.471}{611.338}{\Character{40}}%
\TextAt{14}{376.023}{611.338}{a}%
\TextAt{9}{382.006}{611.338}{\Character{41}}%
\TextAt{17}{386.559}{616.274}{1\Character{47}}%
\TextAt{21}{395.027}{616.274}{n}%
\TextAt{13}{400.663}{611.338}{\Character{58}}%
\TextAt{8}{79.226}{585.133}{The}%
\TextAt{8}{103.28}{585.133}{mean}%
\TextAt{8}{134.495}{585.133}{value}%
\TextAt{8}{164.741}{585.133}{theorem}%
\TextAt{8}{210.267}{585.133}{implies}%
\TextAt{8}{249.994}{585.133}{that}%
\TextAt{16}{215.818}{558.929}{\Character{107}}%
\TextAt{13}{221.795}{558.928}{\Character{39}}%
\TextAt{9}{229.464}{558.929}{\Character{40}}%
\TextAt{13}{234.016}{558.928}{x}%
\TextAt{9}{240.668}{558.929}{\Character{41}}%
\TextAt{16}{247.877}{558.928}{\Character{0}}%
\TextAt{13}{259.832}{558.929}{\Character{39}}%
\TextAt{9}{267.501}{558.928}{\Character{40}}%
\TextAt{13}{272.053}{558.929}{y}%
\TextAt{9}{278.19}{558.928}{\Character{41}}%
\TextAt{16}{282.742}{558.929}{\Character{107}}%
\NativeMathAt{292.031}{558.929}{11.955}{\leq}%
\TextAt{13}{304.66}{558.928}{c}%
\TextAt{17}{309.698}{557.135}{1}%
\TextAt{14}{314.43}{558.929}{N}%
\TextAt{9}{324.378}{558.929}{\Character{40}}%
\TextAt{14}{328.93}{558.929}{a}%
\TextAt{9}{334.914}{558.929}{\Character{41}}%
\TextAt{17}{339.466}{563.865}{1\Character{47}}%
\TextAt{21}{347.934}{563.865}{n}%
\TextAt{16}{353.571}{558.929}{\Character{107}}%
\TextAt{13}{359.548}{558.929}{x}%
\TextAt{16}{368.857}{558.929}{\Character{0}}%
\TextAt{13}{380.812}{558.929}{y}%
\TextAt{16}{386.949}{558.929}{\Character{107}}%
\TextAt{13}{392.926}{558.929}{\Character{59}}%
\TextAt{8}{511.944}{558.929}{(**)}%
\TextAt{8}{79.226}{532.724}{where}%
\TextAt{16}{113.035}{532.724}{\Character{107}}%
\TextAt{16}{122.91}{532.724}{\Character{107}}%
\TextAt{8}{132.785}{532.724}{is}%
\TextAt{8}{144.549}{532.724}{the}%
\TextAt{8}{164.706}{532.724}{euclidean}%
\TextAt{8}{216.077}{532.724}{norm.}%
\TextAt{8}{251.197}{532.724}{The}%
\TextAt{8}{275.251}{532.724}{boundary}%
\TextAt{8}{327.938}{532.724}{of}%
\TextAt{13}{341.268}{532.724}{P}%
\TextAt{27}{348.813}{530.931}{a}%
\TextAt{8}{353.3}{532.724}{,}%
\TextAt{13}{237.944}{498.659}{\Character{64}}%
\TextAt{13}{244.759}{498.659}{P}%
\TextAt{27}{252.319}{496.865}{a}%
\TextAt{9}{260.126}{498.659}{\Character{61}}%
\TextAt{20}{272.552}{510.016}{\Character{91}}%
\TextAt{27}{277.199}{484.91}{p}%
\TextAt{20}{287.828}{519.102}{\Character{34}}%
\TextAt{13}{294.802}{498.659}{\Character{64}}%
\TextAt{13}{301.617}{498.659}{D}%
\TextAt{27}{311.351}{496.865}{p}%
\TextAt{16}{318.495}{498.659}{\Character{2}}%
\TextAt{20}{330.45}{510.016}{\Character{89}}%
\TextAt{27}{330.851}{484.556}{q}%
\NativeMathAt{334.749}{484.556}{7.97}{\neq}%
\TextAt{27}{341.336}{484.556}{p}%
\TextAt{13}{347.718}{498.659}{D}%
\TextAt{27}{357.438}{496.865}{q}%
\TextAt{20}{361.834}{519.102}{\Character{35}}%
\TextAt{13}{370.8}{498.659}{\Character{59}}%
\TextAt{8}{79.226}{462.181}{is}%
\TextAt{8}{91.881}{462.181}{parametrized}%
\TextAt{8}{163.666}{462.181}{by}%
\TextAt{8}{180.816}{462.181}{a}%
\TextAt{8}{191.462}{462.181}{finite}%
\TextAt{8}{222.265}{462.181}{number}%
\TextAt{8}{265.752}{462.181}{of}%
\TextAt{8}{279.973}{462.181}{boundary}%
\TextAt{8}{333.551}{462.181}{cubes}%
\TextAt{13}{366.685}{462.181}{I}%
\TextAt{21}{372.788}{466.52}{n}%
\TextAt{3}{377.926}{466.52}{\Character{0}}%
\TextAt{17}{384.512}{466.52}{1}%
\TextAt{8}{394.033}{462.181}{of}%
\TextAt{13}{408.254}{462.181}{I}%
\TextAt{21}{414.357}{466.52}{n}%
\TextAt{8}{419.993}{462.181}{.}%
\TextAt{8}{431.114}{462.181}{We}%
\TextAt{8}{452.15}{462.181}{subdivide}%
\TextAt{8}{505.775}{462.181}{every}%
\TextAt{8}{79.226}{445.135}{edge}%
\TextAt{8}{105.809}{445.135}{of}%
\TextAt{13}{119.062}{445.135}{I}%
\TextAt{21}{125.165}{449.474}{n}%
\TextAt{3}{130.303}{449.474}{\Character{0}}%
\TextAt{17}{136.89}{449.474}{1}%
\TextAt{8}{145.442}{445.135}{into}%
\TextAt{13}{169.108}{445.135}{m}%
\TextAt{9}{182.668}{445.135}{\Character{61}}%
\TextAt{9}{195.078}{445.135}{\Character{91}}%
\TextAt{14}{198.346}{445.135}{N}%
\TextAt{9}{208.294}{445.135}{\Character{40}}%
\TextAt{14}{212.846}{445.135}{a}%
\TextAt{9}{218.829}{445.135}{\Character{41}}%
\TextAt{17}{223.381}{449.474}{1\Character{47}}%
\TextAt{21}{231.85}{449.474}{n}%
\TextAt{9}{237.486}{445.135}{\Character{93}}%
\NativeMathAt{244.059}{445.135}{11.955}{\geq}%
\TextAt{13}{256.678}{445.135}{C}%
\TextAt{3}{265.912}{449.474}{\Character{48}}%
\TextAt{9}{272.027}{445.135}{\Character{61}}%
\TextAt{9}{284.437}{445.135}{1}%
\TextAt{8}{294.126}{445.135}{segments}%
\TextAt{8}{343.926}{445.135}{of}%
\TextAt{8}{357.179}{445.135}{equal}%
\TextAt{8}{387.994}{445.135}{length}%
\TextAt{8}{423.687}{445.135}{and}%
\TextAt{8}{446.373}{445.135}{obtain}%
\TextAt{8}{482.723}{445.135}{for}%
\TextAt{13}{500.531}{445.135}{I}%
\TextAt{21}{506.634}{449.474}{n}%
\TextAt{3}{511.772}{449.474}{\Character{0}}%
\TextAt{17}{518.359}{449.474}{1}%
\TextAt{8}{526.911}{445.135}{a}%
\TextAt{8}{79.226}{428.089}{decomposition}%
\TextAt{8}{156.446}{428.089}{into}%
\TextAt{13}{180.275}{428.089}{m}%
\TextAt{21}{190.514}{432.428}{n}%
\TextAt{3}{195.652}{432.428}{\Character{0}}%
\TextAt{17}{202.239}{432.428}{1}%
\TextAt{8}{210.954}{428.089}{small}%
\TextAt{8}{241.669}{428.089}{cubes}%
\TextAt{8}{273.997}{428.089}{of}%
\TextAt{8}{287.413}{428.089}{diameter}%
\NativeMathAt{336.275}{428.089}{11.955}{\leq}%
\TextAt{9}{349.04}{428.089}{\Character{40}}%
\TextAt{13}{353.593}{428.089}{n}%
\TextAt{16}{363.295}{428.089}{\Character{0}}%
\TextAt{9}{375.309}{428.089}{1\Character{41}}%
\TextAt{17}{385.714}{432.428}{1\Character{47}2}%
\TextAt{9}{398.915}{428.089}{\Character{47}}%
\TextAt{13}{404.768}{428.089}{m}%
\TextAt{8}{415.007}{428.089}{.}%
\TextAt{8}{423.712}{428.089}{From}%
\TextAt{8}{454.522}{428.089}{(**),}%
\TextAt{8}{482.604}{428.089}{the}%
\TextAt{8}{502.846}{428.089}{image}%
\TextAt{8}{79.226}{411.043}{of}%
\TextAt{8}{92.556}{411.043}{such}%
\TextAt{8}{118.953}{411.043}{a}%
\TextAt{8}{128.709}{411.043}{small}%
\TextAt{8}{159.338}{411.043}{cube}%
\TextAt{8}{186.966}{411.043}{under}%
\TextAt{13}{220.13}{411.043}{\Character{39}}%
\TextAt{8}{231.696}{411.043}{has}%
\TextAt{8}{252.569}{411.043}{a}%
\TextAt{8}{262.325}{411.043}{diameter}%
\NativeMathAt{137.307}{380.16}{11.955}{\leq}%
\TextAt{9}{151.122}{388.248}{\Character{40}}%
\TextAt{13}{155.674}{388.248}{n}%
\TextAt{16}{165.318}{388.248}{\Character{0}}%
\TextAt{9}{177.273}{388.248}{1\Character{41}}%
\TextAt{17}{187.679}{392.586}{1\Character{47}2}%
\Vector{q 0 G 0 g 0.478 w 79.122 -390.851 m 128.879 -390.851 l S Q}%
\TextAt{13}{170.881}{371.959}{m}%
\TextAt{13}{202.075}{380.16}{c}%
\TextAt{17}{207.113}{378.367}{1}%
\TextAt{14}{211.845}{380.16}{N}%
\TextAt{9}{221.793}{380.16}{\Character{40}}%
\TextAt{14}{226.346}{380.16}{a}%
\TextAt{9}{232.329}{380.16}{\Character{41}}%
\TextAt{17}{236.881}{385.096}{1\Character{47}}%
\TextAt{21}{245.349}{385.096}{n}%
\NativeMathAt{254.307}{380.16}{11.955}{\leq}%
\TextAt{9}{266.926}{380.16}{\Character{40}}%
\TextAt{13}{271.478}{380.16}{n}%
\TextAt{16}{281.123}{380.16}{\Character{0}}%
\TextAt{9}{293.078}{380.16}{1\Character{41}}%
\TextAt{17}{303.483}{385.096}{1\Character{47}2}%
\TextAt{13}{316.684}{380.16}{c}%
\TextAt{17}{321.722}{378.367}{1}%
\TextAt{13}{327.65}{388.248}{m}%
\TextAt{9}{340.545}{388.248}{\Character{43}}%
\TextAt{9}{352.297}{388.248}{1}%
\Vector{q 0 G 0 g 0.478 w 255.65 -390.851 m 286.16 -390.851 l S Q}%
\TextAt{13}{337.785}{371.959}{m}%
\NativeMathAt{362.676}{380.16}{11.955}{\leq}%
\TextAt{9}{375.295}{380.16}{\Character{40}}%
\TextAt{13}{379.848}{380.16}{n}%
\TextAt{16}{389.492}{380.16}{\Character{0}}%
\TextAt{9}{401.447}{380.16}{1\Character{41}}%
\TextAt{17}{411.853}{385.096}{1\Character{47}2}%
\TextAt{13}{425.053}{380.16}{c}%
\TextAt{17}{430.091}{378.367}{1}%
\TextAt{9}{436.816}{380.16}{2}%
\TextAt{9}{445.974}{380.16}{\Character{61}\Character{58}}%
\TextAt{13}{461.667}{380.16}{c}%
\TextAt{17}{466.705}{378.367}{2}%
\TextAt{13}{471.437}{380.16}{\Character{58}}%
\TextAt{8}{79.226}{353.956}{The}%
\TextAt{8}{103.638}{353.956}{number}%
\TextAt{8}{146.592}{353.956}{of}%
\TextAt{8}{160.28}{353.956}{translates}%
\TextAt{13}{214.101}{353.956}{F}%
\TextAt{9}{226.205}{353.956}{\Character{43}}%
\TextAt{13}{238.21}{353.956}{\Character{13}}%
\TextAt{8}{244.932}{353.956}{,}%
\TextAt{13}{252.529}{353.956}{\Character{13}}%
\NativeMathAt{263.181}{353.956}{11.955}{\in}%
\TextAt{9}{275.082}{353.956}{\Character{0}}%
\TextAt{8}{282.396}{353.956}{,}%
\TextAt{8}{289.993}{353.956}{meeting}%
\TextAt{8}{334.573}{353.956}{a}%
\TextAt{8}{344.686}{353.956}{domain}%
\TextAt{8}{386.671}{353.956}{of}%
\TextAt{8}{400.359}{353.956}{diameter}%
\NativeMathAt{449.494}{353.956}{11.955}{\leq}%
\TextAt{13}{462.723}{353.956}{c}%
\TextAt{17}{467.761}{352.162}{2}%
\TextAt{8}{476.748}{353.956}{is}%
\TextAt{8}{488.87}{353.956}{bounded}%
\TextAt{8}{79.226}{336.91}{by}%
\TextAt{8}{95.862}{336.91}{a}%
\TextAt{8}{105.995}{336.91}{constant}%
\TextAt{13}{153.595}{336.91}{c}%
\TextAt{17}{158.633}{335.116}{3}%
\TextAt{8}{167.64}{336.91}{which}%
\TextAt{8}{201.504}{336.91}{depends}%
\TextAt{8}{247.131}{336.91}{only}%
\TextAt{8}{273.2}{336.91}{on}%
\TextAt{13}{289.837}{336.91}{c}%
\TextAt{17}{294.875}{335.116}{2}%
\TextAt{8}{303.882}{336.91}{and}%
\TextAt{8}{327.022}{336.91}{the}%
\TextAt{8}{347.555}{336.91}{fundamental}%
\TextAt{8}{415.575}{336.91}{mesh}%
\TextAt{13}{445.924}{336.91}{F}%
\TextAt{8}{455.127}{336.91}{.}%
\TextAt{8}{464.707}{336.91}{The}%
\TextAt{8}{489.138}{336.91}{image}%
\TextAt{8}{523.337}{336.91}{of}%
\TextAt{8}{79.226}{319.864}{a}%
\TextAt{8}{89.47}{319.864}{small}%
\TextAt{8}{120.588}{319.864}{cube}%
\TextAt{8}{148.705}{319.864}{under}%
\TextAt{13}{182.357}{319.864}{\Character{39}}%
\TextAt{8}{194.412}{319.864}{thus}%
\TextAt{8}{220.652}{319.864}{meets}%
\TextAt{8}{254.364}{319.864}{at}%
\TextAt{8}{269.163}{319.864}{most}%
\TextAt{13}{298.332}{319.864}{c}%
\TextAt{17}{303.37}{318.07}{3}%
\TextAt{8}{312.488}{319.864}{translates}%
\TextAt{13}{366.44}{319.864}{F}%
\TextAt{9}{378.634}{319.864}{\Character{43}}%
\TextAt{13}{390.728}{319.864}{\Character{13}}%
\TextAt{8}{397.45}{319.864}{.}%
\TextAt{8}{407.364}{319.864}{Since}%
\TextAt{8}{438.41}{319.864}{there}%
\TextAt{8}{468.811}{319.864}{are}%
\TextAt{8}{488.81}{319.864}{precisely}%
\TextAt{13}{79.226}{302.818}{m}%
\TextAt{21}{89.465}{307.156}{n}%
\TextAt{3}{94.604}{307.156}{\Character{0}}%
\TextAt{17}{101.19}{307.156}{1}%
\TextAt{9}{109.243}{302.818}{\Character{61}}%
\TextAt{9}{121.653}{302.818}{\Character{91}}%
\TextAt{14}{124.92}{302.818}{N}%
\TextAt{9}{134.869}{302.818}{\Character{40}}%
\TextAt{14}{139.421}{302.818}{a}%
\TextAt{9}{145.404}{302.818}{\Character{41}}%
\TextAt{17}{149.956}{307.156}{1\Character{47}}%
\TextAt{21}{158.425}{307.156}{n}%
\TextAt{9}{164.061}{302.818}{\Character{93}}%
\TextAt{21}{167.313}{307.156}{n}%
\TextAt{3}{172.451}{307.156}{\Character{0}}%
\TextAt{17}{179.037}{307.156}{1}%
\TextAt{8}{186.489}{302.818}{cubes}%
\TextAt{8}{217.555}{302.818}{in}%
\TextAt{13}{230.03}{302.818}{\Character{39}}%
\TextAt{9}{237.698}{302.818}{\Character{40}}%
\TextAt{13}{242.25}{302.818}{I}%
\TextAt{21}{248.353}{307.156}{n}%
\TextAt{3}{253.491}{307.156}{\Character{0}}%
\TextAt{17}{260.078}{307.156}{1}%
\TextAt{9}{264.81}{302.818}{\Character{41}}%
\TextAt{8}{269.363}{302.818}{,}%
\TextAt{8}{275.57}{302.818}{we}%
\TextAt{8}{291.619}{302.818}{see}%
\TextAt{8}{309.354}{302.818}{that}%
\TextAt{13}{333.546}{302.818}{\Character{39}}%
\TextAt{9}{341.214}{302.818}{\Character{40}}%
\TextAt{13}{345.766}{302.818}{I}%
\TextAt{21}{351.869}{307.156}{n}%
\TextAt{3}{357.007}{307.156}{\Character{0}}%
\TextAt{17}{363.594}{307.156}{1}%
\TextAt{9}{368.326}{302.818}{\Character{41}}%
\TextAt{8}{375.598}{302.818}{meets}%
\TextAt{8}{407.644}{302.818}{at}%
\TextAt{8}{420.776}{302.818}{most}%
\TextAt{13}{448.279}{302.818}{c}%
\TextAt{17}{453.317}{301.024}{3}%
\TextAt{9}{458.049}{302.818}{\Character{91}}%
\TextAt{14}{461.301}{302.818}{N}%
\TextAt{9}{471.249}{302.818}{\Character{40}}%
\TextAt{14}{475.801}{302.818}{a}%
\TextAt{9}{481.785}{302.818}{\Character{41}}%
\TextAt{17}{486.337}{307.156}{1\Character{47}}%
\TextAt{21}{494.805}{307.156}{n}%
\TextAt{9}{500.442}{302.818}{\Character{93}}%
\TextAt{21}{503.693}{307.156}{n}%
\TextAt{3}{508.831}{307.156}{\Character{0}}%
\TextAt{17}{515.418}{307.156}{1}%
\NativeMathAt{523.471}{302.818}{11.955}{\leq}%
\TextAt{13}{79.226}{285.772}{c}%
\TextAt{17}{84.264}{283.978}{3}%
\TextAt{14}{88.996}{285.772}{N}%
\TextAt{9}{98.945}{285.772}{\Character{40}}%
\TextAt{14}{103.497}{285.772}{a}%
\TextAt{9}{109.48}{285.772}{\Character{41}}%
\TextAt{17}{114.032}{290.11}{1}%
\TextAt{3}{118.266}{290.11}{\Character{0}}%
\TextAt{17}{124.853}{290.11}{1\Character{47}}%
\TextAt{21}{133.321}{290.11}{n}%
\TextAt{8}{143.134}{285.772}{translates,}%
\TextAt{8}{200.199}{285.772}{and}%
\TextAt{8}{223.241}{285.772}{since}%
\TextAt{8}{252.189}{285.772}{the}%
\TextAt{8}{272.624}{285.772}{boundary}%
\TextAt{13}{325.59}{285.772}{\Character{64}}%
\TextAt{13}{332.405}{285.772}{P}%
\TextAt{27}{339.965}{283.978}{a}%
\TextAt{8}{348.629}{285.772}{is}%
\TextAt{8}{360.672}{285.772}{covered}%
\TextAt{8}{402.902}{285.772}{by}%
\TextAt{8}{419.44}{285.772}{at}%
\TextAt{8}{434.03}{285.772}{most}%
\TextAt{9}{462.99}{285.772}{2}%
\TextAt{13}{468.843}{285.772}{n}%
\TextAt{8}{480.007}{285.772}{such}%
\TextAt{8}{506.683}{285.772}{parts}%
\TextAt{13}{79.226}{268.726}{\Character{39}}%
\TextAt{9}{86.895}{268.726}{\Character{40}}%
\TextAt{13}{91.447}{268.726}{I}%
\TextAt{21}{97.55}{273.064}{n}%
\TextAt{3}{102.688}{273.064}{\Character{0}}%
\TextAt{17}{109.274}{273.064}{1}%
\TextAt{9}{114.007}{268.726}{\Character{41}}%
\TextAt{8}{118.559}{268.726}{,}%
\TextAt{8}{125.708}{268.726}{we}%
\TextAt{8}{142.936}{268.726}{do}%
\TextAt{8}{159.195}{268.726}{indeed}%
\TextAt{8}{196.256}{268.726}{find}%
\TextAt{8}{219.664}{268.726}{that}%
\TextAt{8}{189.95}{242.521}{\char35\relax }%
\TextAt{16}{199.706}{242.521}{\Character{102}}%
\TextAt{13}{205.683}{242.521}{\Character{13}}%
\NativeMathAt{215.726}{242.521}{11.955}{\in}%
\TextAt{9}{227.017}{242.521}{\Character{0}}%
\TextAt{16}{237.652}{242.521}{\Character{106}}%
\TextAt{9}{244.294}{242.521}{\Character{40}}%
\TextAt{13}{248.846}{242.521}{F}%
\TextAt{9}{260.707}{242.521}{\Character{43}}%
\TextAt{13}{272.468}{242.521}{\Character{13}}%
\TextAt{9}{279.19}{242.521}{\Character{41}}%
\TextAt{16}{286.399}{242.521}{\Character{92}}%
\TextAt{13}{297.026}{242.521}{\Character{64}}%
\TextAt{13}{303.841}{242.521}{P}%
\TextAt{27}{311.401}{240.728}{a}%
\NativeMathAt{319.209}{242.521}{11.955}{\neq}%
\TextAt{18}{331.634}{242.521}{\Character{63}}%
\TextAt{16}{340.933}{242.521}{\Character{103}}%
\NativeMathAt{350.222}{242.521}{11.955}{\leq}%
\TextAt{13}{362.851}{242.521}{C}%
\TextAt{14}{372.084}{242.521}{N}%
\TextAt{9}{382.032}{242.521}{\Character{40}}%
\TextAt{14}{386.585}{242.521}{a}%
\TextAt{9}{392.568}{242.521}{\Character{41}}%
\TextAt{17}{397.12}{247.457}{1}%
\TextAt{3}{401.354}{247.457}{\Character{0}}%
\TextAt{17}{407.941}{247.457}{1\Character{47}}%
\TextAt{21}{416.409}{247.457}{n}%
\TextAt{8}{79.226}{216.317}{for}%
\TextAt{8}{98.231}{216.317}{all}%
\TextAt{14}{115.61}{216.317}{a}%
\NativeMathAt{126.823}{216.317}{11.955}{\in}%
\TextAt{16}{140.01}{216.317}{\Character{65}}%
\TextAt{21}{149.568}{214.523}{i}%
\TextAt{8}{157.967}{216.317}{with}%
\TextAt{14}{185.747}{216.317}{N}%
\TextAt{9}{195.695}{216.317}{\Character{40}}%
\TextAt{14}{200.247}{216.317}{a}%
\TextAt{9}{206.23}{216.317}{\Character{41}}%
\NativeMathAt{216.012}{216.317}{11.955}{\geq}%
\TextAt{9}{230.54}{216.317}{1}%
\TextAt{8}{236.393}{216.317}{,}%
\TextAt{8}{244.942}{216.317}{where}%
\TextAt{13}{279.872}{216.317}{C}%
\TextAt{9}{294.335}{216.317}{\Character{61}}%
\TextAt{9}{308.657}{216.317}{2}%
\TextAt{13}{314.522}{216.317}{nc}%
\TextAt{17}{326.547}{214.523}{3}%
\TextAt{8}{336.297}{216.317}{is}%
\TextAt{8}{349.181}{216.317}{a}%
\TextAt{8}{360.056}{216.317}{constant}%
\TextAt{8}{408.399}{216.317}{which}%
\TextAt{8}{443.006}{216.317}{is}%
\TextAt{8}{455.89}{216.317}{independent}%
\TextAt{8}{523.337}{216.317}{of}%
\TextAt{14}{79.226}{199.271}{a}%
\NativeMathAt{88.53}{199.271}{11.955}{\in}%
\TextAt{16}{99.804}{199.271}{\Character{65}}%
\TextAt{21}{109.367}{197.477}{i}%
\TextAt{8}{112.748}{199.271}{,}%
\TextAt{8}{119.897}{199.271}{as}%
\TextAt{8}{134.267}{199.271}{required.}%
\Vector{q 0 G 0 g 0.399 w 452.334 -574.729 m 452.334 -566.659 l S Q}%
\Vector{q 0 G 0 g 0.399 w 452.534 -566.859 m 459.707 -566.859 l S Q}%
\Vector{q 0 G 0 g 0.399 w 452.534 -574.53 m 459.707 -574.53 l S Q}%
\Vector{q 0 G 0 g 0.399 w 459.906 -574.729 m 459.906 -566.659 l S Q}%
\TextAt{8}{97.159}{171.069}{From}%
\TextAt{8}{129.79}{171.069}{the}%
\TextAt{8}{151.852}{171.069}{theorem}%
\TextAt{8}{199.284}{171.069}{we}%
\TextAt{8}{218.417}{171.069}{have}%
\TextAt{8}{247.318}{171.069}{just}%
\TextAt{8}{272.369}{171.069}{proved,}%
\TextAt{8}{316.057}{171.069}{we}%
\TextAt{8}{335.191}{171.069}{now}%
\TextAt{8}{361.485}{171.069}{obtain}%
\TextAt{8}{399.819}{171.069}{the}%
\TextAt{8}{421.881}{171.069}{strong}%
\TextAt{8}{459.629}{171.069}{version}%
\TextAt{8}{501.274}{171.069}{of}%
\TextAt{8}{516.51}{171.069}{the}%
\TextAt{8}{79.226}{154.023}{Riemann–Roch}%
\TextAt{8}{161.269}{154.023}{theorem.}%
\TextAt{8}{212.584}{154.023}{We}%
\TextAt{8}{233.142}{154.023}{want}%
\TextAt{8}{262.175}{154.023}{to}%
\TextAt{8}{276.899}{154.023}{formulate}%
\TextAt{8}{329.998}{154.023}{it}%
\TextAt{8}{342.115}{154.023}{in}%
\TextAt{8}{356.181}{154.023}{the}%
\TextAt{8}{376.751}{154.023}{language}%
\TextAt{8}{425.953}{154.023}{of}%
\TextAt{8}{439.696}{154.023}{divisors.}%
\TextAt{8}{488.213}{154.023}{Let}%
\TextAt{13}{509.596}{154.023}{D}%
\TextAt{9}{523.665}{154.023}{\Character{61}}%
\TextAt{20}{79.226}{145.943}{\Character{80}}%
\TextAt{27}{91.846}{133.49}{p}%
\TextAt{13}{98.325}{136.977}{\Character{23}}%
\TextAt{27}{104.109}{135.183}{p}%
\TextAt{14}{108.596}{136.977}{p}%
\TextAt{8}{118.476}{136.977}{be}%
\TextAt{8}{134.4}{136.977}{a}%
\TextAt{8}{144.156}{136.977}{replete}%
\TextAt{8}{182.52}{136.977}{divisor}%
\TextAt{8}{220.633}{136.977}{of}%
\TextAt{13}{233.963}{136.977}{K}%
\TextAt{8}{244.781}{136.977}{,}%
\TextAt{13}{186.181}{110.772}{H}%
\TextAt{17}{196.804}{115.708}{0}%
\TextAt{9}{201.536}{110.772}{\Character{40}}%
\TextAt{13}{206.088}{110.772}{D}%
\TextAt{9}{216.133}{110.772}{\Character{41}}%
\TextAt{9}{223.988}{110.772}{\Character{61}}%
\TextAt{13}{236.431}{110.772}{H}%
\TextAt{17}{247.054}{115.708}{0}%
\TextAt{9}{251.786}{110.772}{\Character{40}}%
\TextAt{14}{256.339}{110.772}{o}%
\TextAt{9}{262.183}{110.772}{\Character{40}}%
\TextAt{13}{266.735}{110.772}{D}%
\TextAt{9}{276.78}{110.772}{\Character{41}\Character{41}}%
\TextAt{9}{289.178}{110.772}{\Character{61}}%
\TextAt{16}{301.631}{110.772}{\Character{102}}%
\TextAt{13}{307.608}{110.772}{f}%
\NativeMathAt{317.975}{110.772}{11.955}{\in}%
\TextAt{13}{329.266}{110.772}{K}%
\TextAt{3}{340.084}{115.708}{\Character{3}}%
\TextAt{16}{348.137}{110.772}{\Character{106}}%
\TextAt{13}{354.779}{110.772}{v}%
\TextAt{27}{360.438}{108.979}{p}%
\TextAt{9}{364.925}{110.772}{\Character{40}}%
\TextAt{13}{369.477}{110.772}{f}%
\TextAt{9}{376.524}{110.772}{\Character{41}}%
\NativeMathAt{384.397}{110.772}{11.955}{\geq}%
\TextAt{16}{396.998}{110.772}{\Character{0}}%
\TextAt{13}{406.315}{110.772}{\Character{23}}%
\TextAt{27}{412.099}{108.979}{p}%
\TextAt{16}{416.585}{110.772}{\Character{103}}%
\TextAt{13}{422.563}{110.772}{\Character{59}}%
\ErrataAt{145}{84.111}{342}{$\ell(D)=\ell(\mathfrak o(D)),\qquad \chi(D)=\chi(\mathfrak o(D)).$}%
\TextAt{8}{297.211}{43.033}{199}%
\Destination{516.5}{575.975}{nav.AMS.27}%
\Destination{78.23}{809.815}{nav.page.199}%
\Destination{79.36}{443.055}{nav.print.218}%
\Destination{79.226}{101.111}{errata.94}%
}

% Page 212 (body page 200)
\Page{%
\TextAt{2}{79.226}{796.943}{III.}%
\TextAt{2}{98.317}{796.943}{Riemann–Roch}%
\TextAt{2}{174.375}{796.943}{Theory}%
\Vector{q 0 G 0 g 0.299 w 7.226 18.729 m 460.769 18.729 l S Q}%
\TextAt{8}{79.226}{761.124}{We}%
\TextAt{8}{99.371}{761.124}{call}%
\TextAt{8}{120.83}{761.124}{the}%
\TextAt{8}{140.987}{761.124}{number}%
\TextAt{13}{253.976}{744.078}{i}%
\TextAt{9}{257.969}{744.078}{\Character{40}}%
\TextAt{13}{262.521}{744.078}{D}%
\TextAt{9}{272.566}{744.078}{\Character{41}}%
\TextAt{9}{280.421}{744.078}{\Character{61}}%
\TextAt{13}{292.864}{744.078}{\Character{96}}%
\TextAt{9}{297.776}{744.078}{\Character{40}}%
\TextAt{13}{302.329}{744.078}{D}%
\TextAt{9}{312.373}{744.078}{\Character{41}}%
\TextAt{16}{319.582}{744.078}{\Character{0}}%
\TextAt{13}{331.537}{744.078}{\Character{31}}%
\TextAt{9}{338.871}{744.078}{\Character{40}}%
\TextAt{13}{343.423}{744.078}{D}%
\TextAt{9}{353.468}{744.078}{\Character{41}}%
\TextAt{8}{79.226}{720.556}{the}%
\TextAt{12}{99.383}{720.556}{index}%
\TextAt{12}{135.786}{720.556}{of}%
\TextAt{12}{151.113}{720.556}{specialty}%
\TextAt{8}{206.262}{720.556}{of}%
\TextAt{13}{219.592}{720.556}{D}%
\TextAt{8}{233.534}{720.556}{and}%
\TextAt{8}{256.296}{720.556}{get}%
\TextAt{8}{275.807}{720.556}{the}%
\TextAt{12}{79.226}{694.544}{Theorem}%
\TextAt{12}{136.745}{694.544}{(3.9)}%
\TextAt{12}{169.409}{694.544}{(Riemann–Roch).}%
\TextAt{8}{277.504}{694.544}{For}%
\TextAt{8}{298.913}{694.544}{every}%
\TextAt{8}{330.244}{694.544}{replete}%
\TextAt{8}{369.047}{694.544}{divisor}%
\TextAt{13}{407.599}{694.544}{D}%
\NativeMathAt{421.712}{694.544}{11.955}{\in}%
\TextAt{8}{433.75}{694.544}{Div}%
\TextAt{9}{452.126}{694.544}{\Character{40}}%
\Vector{q 0 G 0 g 0.478 w 384.678 -72.1 m 390.522 -72.1 l S Q}%
\TextAt{14}{456.678}{694.544}{o}%
\TextAt{9}{462.522}{694.544}{\Character{41}}%
\TextAt{8}{471.411}{694.544}{we}%
\TextAt{8}{489.077}{694.544}{have}%
\TextAt{8}{516.51}{694.544}{the}%
\TextAt{8}{79.226}{677.498}{formula}%
\TextAt{13}{219.87}{660.452}{\Character{96}}%
\TextAt{9}{224.782}{660.452}{\Character{40}}%
\TextAt{13}{229.334}{660.452}{D}%
\TextAt{9}{239.379}{660.452}{\Character{41}}%
\TextAt{9}{247.234}{660.452}{\Character{61}}%
\TextAt{8}{259.677}{660.452}{deg}%
\TextAt{9}{277.239}{660.452}{\Character{40}}%
\TextAt{13}{281.792}{660.452}{D}%
\TextAt{9}{291.836}{660.452}{\Character{41}}%
\TextAt{9}{299.033}{660.452}{\Character{43}}%
\TextAt{13}{310.807}{660.452}{\Character{96}}%
\TextAt{9}{315.719}{660.452}{\Character{40}}%
\TextAt{14}{320.271}{660.452}{o}%
\TextAt{9}{326.115}{660.452}{\Character{41}}%
\TextAt{16}{333.324}{660.452}{\Character{0}}%
\TextAt{13}{345.279}{660.452}{g}%
\TextAt{9}{353.97}{660.452}{\Character{43}}%
\TextAt{13}{365.732}{660.452}{i}%
\TextAt{9}{369.725}{660.452}{\Character{40}}%
\TextAt{13}{374.277}{660.452}{D}%
\TextAt{9}{384.322}{660.452}{\Character{41}}%
\TextAt{13}{388.874}{660.452}{\Character{58}}%
\TextAt{8}{79.226}{636.93}{The}%
\TextAt{8}{103.28}{636.93}{index}%
\TextAt{8}{134.818}{636.93}{of}%
\TextAt{8}{148.148}{636.93}{specialty}%
\TextAt{13}{196.662}{636.93}{i}%
\TextAt{9}{200.655}{636.93}{\Character{40}}%
\TextAt{13}{205.208}{636.93}{D}%
\TextAt{9}{215.252}{636.93}{\Character{41}}%
\TextAt{8}{223.702}{636.93}{satisfies}%
\TextAt{13}{252.05}{607.929}{i}%
\TextAt{9}{256.044}{607.929}{\Character{40}}%
\TextAt{13}{260.596}{607.929}{D}%
\TextAt{9}{270.64}{607.929}{\Character{41}}%
\TextAt{9}{278.495}{607.929}{\Character{61}}%
\TextAt{13}{290.939}{607.929}{O}%
\TextAt{9}{300.21}{607.929}{\Character{40}}%
\TextAt{13}{304.763}{607.929}{e}%
\TextAt{3}{310.188}{612.865}{\Character{0}}%
\TextAt{24}{318.444}{616.124}{1}%
\Vector{q 0 G 0 g 0.359 w 245.97 -159.142 m 250.571 -159.142 l S Q}%
\TextAt{22}{317.97}{609.74}{n}%
\TextAt{7}{325.178}{612.865}{deg}%
\TextAt{17}{337.874}{612.865}{\Character{40}}%
\TextAt{21}{341.167}{612.865}{D}%
\TextAt{17}{348.35}{612.865}{\Character{41}}%
\TextAt{9}{352.141}{607.929}{\Character{41}}%
\TextAt{13}{356.694}{607.929}{\Character{59}}%
\TextAt{8}{79.226}{578.928}{in}%
\TextAt{8}{92.879}{578.928}{particular,}%
\TextAt{13}{150.12}{578.928}{i}%
\TextAt{9}{154.114}{578.928}{\Character{40}}%
\TextAt{13}{158.666}{578.928}{D}%
\TextAt{9}{168.711}{578.928}{\Character{41}}%
\NativeMathAt{176.584}{578.928}{11.955}{\rightarrow}%
\TextAt{9}{191.86}{578.928}{0}%
\TextAt{8}{201.61}{578.928}{for}%
\TextAt{8}{219.495}{578.928}{deg}%
\TextAt{9}{237.057}{578.928}{\Character{40}}%
\TextAt{13}{241.61}{578.928}{D}%
\TextAt{9}{251.654}{578.928}{\Character{41}}%
\NativeMathAt{259.527}{578.928}{11.955}{\rightarrow}%
\NativeMathAt{274.794}{578.928}{11.955}{\infty}%
\TextAt{8}{286.758}{578.928}{.}%
\TextAt{19}{79.226}{543.949}{Proof.}%
\TextAt{8}{116}{543.949}{The}%
\TextAt{8}{140.004}{543.949}{formula}%
\TextAt{8}{182.886}{543.949}{for}%
\TextAt{13}{200.721}{543.949}{\Character{96}}%
\TextAt{9}{205.633}{543.949}{\Character{40}}%
\TextAt{13}{210.185}{543.949}{D}%
\TextAt{9}{220.23}{543.949}{\Character{41}}%
\TextAt{8}{228.63}{543.949}{follows}%
\TextAt{8}{267.016}{543.949}{from}%
\TextAt{13}{294.607}{543.949}{\Character{31}}%
\TextAt{9}{301.94}{543.949}{\Character{40}}%
\TextAt{13}{306.492}{543.949}{D}%
\TextAt{9}{316.537}{543.949}{\Character{41}}%
\TextAt{9}{324.392}{543.949}{\Character{61}}%
\TextAt{8}{336.835}{543.949}{deg}%
\TextAt{9}{354.398}{543.949}{\Character{40}}%
\TextAt{13}{358.95}{543.949}{D}%
\TextAt{9}{368.994}{543.949}{\Character{41}}%
\TextAt{9}{376.083}{543.949}{\Character{43}}%
\TextAt{13}{387.761}{543.949}{\Character{96}}%
\TextAt{9}{392.674}{543.949}{\Character{40}}%
\TextAt{14}{397.226}{543.949}{o}%
\TextAt{9}{403.07}{543.949}{\Character{41}}%
\TextAt{16}{410.177}{543.949}{\Character{0}}%
\TextAt{13}{422.031}{543.949}{g}%
\TextAt{8}{431.913}{543.949}{and}%
\TextAt{13}{454.626}{543.949}{\Character{31}}%
\TextAt{9}{461.959}{543.949}{\Character{40}}%
\TextAt{13}{466.512}{543.949}{D}%
\TextAt{9}{476.556}{543.949}{\Character{41}}%
\TextAt{9}{484.411}{543.949}{\Character{61}}%
\TextAt{13}{496.855}{543.949}{\Character{96}}%
\TextAt{9}{501.767}{543.949}{\Character{40}}%
\TextAt{13}{506.319}{543.949}{D}%
\TextAt{9}{516.364}{543.949}{\Character{41}}%
\TextAt{16}{523.471}{543.949}{\Character{0}}%
\TextAt{13}{79.226}{526.903}{i}%
\TextAt{9}{83.22}{526.903}{\Character{40}}%
\TextAt{13}{87.772}{526.903}{D}%
\TextAt{9}{97.816}{526.903}{\Character{41}}%
\TextAt{8}{102.369}{526.903}{.}%
\TextAt{8}{110.817}{526.903}{Putting}%
\TextAt{14}{153.903}{526.903}{a}%
\TextAt{3}{159.887}{531.242}{\Character{0}}%
\TextAt{17}{166.473}{531.242}{1}%
\TextAt{9}{174.526}{526.903}{\Character{61}}%
\TextAt{14}{186.952}{526.903}{o}%
\TextAt{9}{192.796}{526.903}{\Character{40}}%
\TextAt{13}{197.348}{526.903}{D}%
\TextAt{9}{207.393}{526.903}{\Character{41}}%
\TextAt{8}{211.945}{526.903}{,}%
\TextAt{8}{219.094}{526.903}{we}%
\TextAt{8}{236.322}{526.903}{find}%
\TextAt{8}{259.73}{526.903}{by}%
\TextAt{8}{275.989}{526.903}{(3.7)}%
\TextAt{8}{303.964}{526.903}{that}%
\TextAt{8}{203.247}{501.476}{\char35\relax }%
\TextAt{13}{213.003}{501.476}{H}%
\TextAt{17}{223.625}{505.815}{0}%
\TextAt{9}{228.358}{501.476}{\Character{40}}%
\TextAt{14}{232.91}{501.476}{a}%
\TextAt{3}{238.893}{505.815}{\Character{0}}%
\TextAt{17}{245.479}{505.815}{1}%
\TextAt{9}{250.212}{501.476}{\Character{41}}%
\Vector{q 0 G 0 g 0.478 w 131.247 -277.623 m 182.764 -277.623 l S Q}%
\TextAt{9}{210.581}{485.188}{2}%
\TextAt{21}{216.434}{488.641}{r}%
\TextAt{9}{220.989}{485.188}{\Character{40}2}%
\TextAt{13}{231.394}{485.188}{\Character{25}}%
\TextAt{9}{238.464}{485.188}{\Character{41}}%
\TextAt{21}{243.016}{488.641}{s}%
\TextAt{9}{259.28}{493.388}{\Character{61}}%
\TextAt{14}{276.825}{501.476}{N}%
\TextAt{9}{286.773}{501.476}{\Character{40}}%
\TextAt{14}{291.325}{501.476}{a}%
\TextAt{9}{297.308}{501.476}{\Character{41}}%
\Vector{q 0 G 0 g 0.478 w 200.901 -277.623 m 233.784 -277.623 l S Q}%
\TextAt{20}{272.901}{493.747}{\Character{112}}%
\Vector{q 0 G 0 g 0.478 w 212.857 -280.014 m 233.784 -280.014 l S Q}%
\TextAt{16}{284.857}{483.525}{\Character{106}}%
\TextAt{13}{288.177}{483.525}{d}%
\TextAt{21}{294.26}{481.732}{K}%
\TextAt{16}{302.463}{483.525}{\Character{106}}%
\TextAt{20}{306.979}{503.072}{\Character{0}}%
\TextAt{9}{312.459}{493.388}{1}%
\TextAt{9}{320.959}{493.388}{\Character{43}}%
\TextAt{13}{332.73}{493.388}{\Character{39}}%
\TextAt{9}{340.398}{493.388}{\Character{40}}%
\TextAt{14}{344.95}{493.388}{a}%
\TextAt{9}{350.933}{493.388}{\Character{41}}%
\TextAt{14}{355.486}{493.388}{N}%
\TextAt{9}{365.434}{493.388}{\Character{40}}%
\TextAt{14}{369.986}{493.388}{a}%
\TextAt{9}{375.969}{493.388}{\Character{41}}%
\TextAt{3}{380.522}{498.325}{\Character{0}}%
\TextAt{17}{387.108}{498.325}{1\Character{47}}%
\TextAt{21}{395.577}{498.325}{n}%
\TextAt{20}{401.213}{503.072}{\Character{1}}%
\TextAt{13}{406.692}{493.388}{\Character{59}}%
\TextAt{8}{79.226}{457.961}{for}%
\TextAt{8}{99.735}{457.961}{some}%
\TextAt{8}{131.684}{457.961}{function}%
\TextAt{13}{180.156}{457.961}{\Character{39}}%
\TextAt{9}{187.824}{457.961}{\Character{40}}%
\TextAt{14}{192.377}{457.961}{a}%
\TextAt{9}{198.36}{457.961}{\Character{41}}%
\TextAt{8}{209.433}{457.961}{which}%
\TextAt{8}{245.543}{457.961}{remains}%
\TextAt{8}{291.803}{457.961}{bounded}%
\TextAt{8}{342.223}{457.961}{as}%
\TextAt{14}{359.217}{457.961}{N}%
\TextAt{9}{369.165}{457.961}{\Character{40}}%
\TextAt{14}{373.718}{457.961}{a}%
\TextAt{9}{379.701}{457.961}{\Character{41}}%
\NativeMathAt{392.045}{457.961}{11.955}{\rightarrow}%
\NativeMathAt{411.783}{457.961}{11.955}{\infty}%
\TextAt{8}{423.747}{457.961}{,}%
\TextAt{8}{434.176}{457.961}{so}%
\TextAt{8}{451.169}{457.961}{that}%
\TextAt{8}{479.162}{457.961}{deg}%
\TextAt{9}{496.724}{457.961}{\Character{40}}%
\TextAt{13}{501.276}{457.961}{D}%
\TextAt{9}{511.321}{457.961}{\Character{41}}%
\TextAt{9}{523.647}{457.961}{\Character{61}}%
\TextAt{16}{79.226}{440.915}{\Character{0}}%
\TextAt{8}{90.517}{440.915}{log}%
\TextAt{14}{107.478}{440.915}{N}%
\TextAt{9}{117.426}{440.915}{\Character{40}}%
\TextAt{14}{121.978}{440.915}{a}%
\TextAt{3}{127.961}{445.254}{\Character{0}}%
\TextAt{17}{134.548}{445.254}{1}%
\TextAt{9}{139.28}{440.915}{\Character{41}}%
\TextAt{9}{147.135}{440.915}{\Character{61}}%
\TextAt{8}{159.579}{440.915}{log}%
\TextAt{14}{176.539}{440.915}{N}%
\TextAt{9}{186.487}{440.915}{\Character{40}}%
\TextAt{14}{191.039}{440.915}{a}%
\TextAt{9}{197.022}{440.915}{\Character{41}}%
\NativeMathAt{204.896}{440.915}{11.955}{\rightarrow}%
\NativeMathAt{220.163}{440.915}{11.955}{\infty}%
\TextAt{8}{232.127}{440.915}{.}%
\TextAt{8}{240.199}{440.915}{Taking}%
\TextAt{8}{278.092}{440.915}{logarithms}%
\TextAt{8}{334.922}{440.915}{and}%
\TextAt{8}{356.555}{440.915}{observing}%
\TextAt{8}{407.85}{440.915}{that}%
\TextAt{8}{432.09}{440.915}{log}%
\TextAt{9}{447.058}{440.915}{\Character{40}1}%
\TextAt{9}{457.794}{440.915}{\Character{43}}%
\TextAt{13}{467.265}{440.915}{O}%
\TextAt{9}{476.536}{440.915}{\Character{40}}%
\TextAt{13}{481.088}{440.915}{t}%
\TextAt{9}{485.315}{440.915}{\Character{41}\Character{41}}%
\TextAt{9}{497.713}{440.915}{\Character{61}}%
\TextAt{13}{510.166}{440.915}{O}%
\TextAt{9}{519.438}{440.915}{\Character{40}}%
\TextAt{13}{523.99}{440.915}{t}%
\TextAt{9}{528.217}{440.915}{\Character{41}}%
\TextAt{8}{79.226}{423.869}{and}%
\TextAt{14}{101.989}{423.869}{N}%
\TextAt{9}{111.937}{423.869}{\Character{40}}%
\TextAt{14}{116.489}{423.869}{a}%
\TextAt{9}{122.472}{423.869}{\Character{41}}%
\TextAt{3}{127.025}{428.208}{\Character{0}}%
\TextAt{17}{133.611}{428.208}{1\Character{47}}%
\TextAt{21}{142.08}{428.208}{n}%
\TextAt{9}{151.037}{423.869}{\Character{61}}%
\TextAt{8}{163.462}{423.869}{exp}%
\TextAt{9}{181.347}{423.869}{\Character{40}}%
\TextAt{16}{185.899}{423.869}{\Character{0}}%
\TextAt{17}{196.845}{428.576}{1}%
\Vector{q 0 G 0 g 0.478 w 124.393 -347.142 m 129.532 -347.142 l S Q}%
\TextAt{21}{196.393}{419.747}{n}%
\TextAt{8}{204.72}{423.869}{deg}%
\TextAt{13}{224.274}{423.869}{D}%
\TextAt{9}{234.319}{423.869}{\Character{41}}%
\TextAt{8}{238.871}{423.869}{,}%
\TextAt{8}{246.02}{423.869}{we}%
\TextAt{8}{263.248}{423.869}{obtain}%
\TextAt{13}{170.074}{394.868}{\Character{96}}%
\TextAt{9}{174.986}{394.868}{\Character{40}}%
\TextAt{13}{179.538}{394.868}{D}%
\TextAt{9}{189.583}{394.868}{\Character{41}}%
\TextAt{9}{197.438}{394.868}{\Character{61}}%
\TextAt{13}{209.881}{394.868}{\Character{96}}%
\TextAt{9}{214.793}{394.868}{\Character{40}}%
\TextAt{14}{219.346}{394.868}{a}%
\TextAt{3}{225.329}{399.804}{\Character{0}}%
\TextAt{17}{231.915}{399.804}{1}%
\TextAt{9}{236.648}{394.868}{\Character{41}}%
\TextAt{9}{244.503}{394.868}{\Character{61}}%
\TextAt{16}{256.946}{394.868}{\Character{0}}%
\TextAt{8}{268.237}{394.868}{log}%
\TextAt{20}{283.205}{404.552}{\Character{0}}%
\TextAt{20}{288.685}{405.674}{\Character{112}}%
\Vector{q 0 G 0 g 0.478 w 228.64 -368.087 m 249.567 -368.087 l S Q}%
\TextAt{16}{300.64}{394.868}{\Character{106}}%
\TextAt{13}{303.961}{394.868}{d}%
\TextAt{21}{310.043}{393.075}{K}%
\TextAt{16}{318.246}{394.868}{\Character{106}}%
\TextAt{14}{321.567}{394.868}{N}%
\TextAt{9}{331.515}{394.868}{\Character{40}}%
\TextAt{14}{336.067}{394.868}{a}%
\TextAt{3}{342.05}{399.804}{\Character{0}}%
\TextAt{17}{348.637}{399.804}{1}%
\TextAt{9}{353.369}{394.868}{\Character{41}}%
\TextAt{20}{357.922}{404.552}{\Character{1}}%
\TextAt{9}{366.058}{394.868}{\Character{43}}%
\TextAt{13}{377.819}{394.868}{O}%
\TextAt{9}{387.09}{394.868}{\Character{40}}%
\TextAt{14}{391.643}{394.868}{N}%
\TextAt{9}{401.591}{394.868}{\Character{40}}%
\TextAt{14}{406.143}{394.868}{a}%
\TextAt{9}{412.126}{394.868}{\Character{41}}%
\TextAt{3}{416.678}{399.804}{\Character{0}}%
\TextAt{17}{423.265}{399.804}{1\Character{47}}%
\TextAt{21}{431.733}{399.804}{n}%
\TextAt{9}{437.37}{394.868}{\Character{41}}%
\TextAt{9}{244.521}{373.722}{\Character{61}}%
\TextAt{13}{256.946}{373.722}{\Character{31}}%
\TextAt{9}{264.28}{373.722}{\Character{40}}%
\TextAt{13}{268.832}{373.722}{D}%
\TextAt{9}{278.877}{373.722}{\Character{41}}%
\TextAt{9}{286.074}{373.722}{\Character{43}}%
\TextAt{13}{297.847}{373.722}{O}%
\TextAt{9}{307.118}{373.722}{\Character{40}}%
\TextAt{13}{311.67}{373.722}{e}%
\TextAt{3}{317.096}{378.659}{\Character{0}}%
\TextAt{24}{325.352}{381.917}{1}%
\Vector{q 0 G 0 g 0.359 w 252.878 -393.349 m 257.479 -393.349 l S Q}%
\TextAt{22}{324.878}{375.533}{n}%
\TextAt{7}{332.086}{378.659}{deg}%
\TextAt{21}{346.193}{378.659}{D}%
\TextAt{9}{353.874}{373.722}{\Character{41}}%
\TextAt{13}{358.426}{373.722}{\Character{58}}%
\TextAt{8}{79.226}{344.721}{Hence}%
\TextAt{13}{114.004}{344.721}{i}%
\TextAt{9}{117.997}{344.721}{\Character{40}}%
\TextAt{13}{122.549}{344.721}{D}%
\TextAt{9}{132.594}{344.721}{\Character{41}}%
\TextAt{9}{140.449}{344.721}{\Character{61}}%
\TextAt{13}{152.893}{344.721}{\Character{96}}%
\TextAt{9}{157.805}{344.721}{\Character{40}}%
\TextAt{13}{162.357}{344.721}{D}%
\TextAt{9}{172.401}{344.721}{\Character{41}}%
\TextAt{16}{179.61}{344.721}{\Character{0}}%
\TextAt{13}{191.566}{344.721}{\Character{31}}%
\TextAt{9}{198.899}{344.721}{\Character{40}}%
\TextAt{13}{203.451}{344.721}{D}%
\TextAt{9}{213.496}{344.721}{\Character{41}}%
\TextAt{9}{221.351}{344.721}{\Character{61}}%
\TextAt{13}{233.795}{344.721}{O}%
\TextAt{9}{243.066}{344.721}{\Character{40}}%
\TextAt{13}{247.618}{344.721}{e}%
\TextAt{3}{253.044}{349.134}{\Character{0}}%
\TextAt{24}{261.3}{352.392}{1}%
\Vector{q 0 G 0 g 0.359 w 188.826 -422.874 m 193.426 -422.874 l S Q}%
\TextAt{22}{260.826}{346.008}{n}%
\TextAt{7}{268.033}{349.134}{deg}%
\TextAt{21}{282.141}{349.134}{D}%
\TextAt{9}{289.822}{344.721}{\Character{41}}%
\TextAt{8}{294.374}{344.721}{.}%
\Vector{q 0 G 0 g 0.399 w 452.334 -429.279 m 452.334 -421.209 l S Q}%
\Vector{q 0 G 0 g 0.399 w 452.534 -421.408 m 459.707 -421.408 l S Q}%
\Vector{q 0 G 0 g 0.399 w 452.534 -429.079 m 459.707 -429.079 l S Q}%
\Vector{q 0 G 0 g 0.399 w 459.906 -429.279 m 459.906 -421.209 l S Q}%
\TextAt{8}{97.159}{315.72}{To}%
\TextAt{8}{115.188}{315.72}{conclude}%
\TextAt{8}{164.108}{315.72}{this}%
\TextAt{8}{187.732}{315.72}{section,}%
\TextAt{8}{231.067}{315.72}{let}%
\TextAt{8}{248.773}{315.72}{us}%
\TextAt{8}{264.59}{315.72}{study}%
\TextAt{8}{297.646}{315.72}{the}%
\TextAt{8}{318.604}{315.72}{variation}%
\TextAt{8}{368.529}{315.72}{of}%
\TextAt{8}{382.66}{315.72}{the}%
\TextAt{8}{403.618}{315.72}{Euler–Minkowski}%
\TextAt{8}{496.019}{315.72}{charac-}%
\TextAt{8}{79.226}{298.674}{teristic}%
\TextAt{8}{119.131}{298.674}{and}%
\TextAt{8}{142.717}{298.674}{of}%
\TextAt{8}{156.87}{298.674}{the}%
\TextAt{8}{177.849}{298.674}{genus}%
\TextAt{8}{210.928}{298.674}{when}%
\TextAt{8}{242.308}{298.674}{we}%
\TextAt{8}{260.359}{298.674}{change}%
\TextAt{8}{299.881}{298.674}{the}%
\TextAt{8}{320.861}{298.674}{field}%
\TextAt{13}{347.041}{298.674}{K}%
\TextAt{8}{357.858}{298.674}{.}%
\TextAt{8}{368.776}{298.674}{Let}%
\TextAt{13}{390.569}{298.674}{L}%
\TextAt{16}{403.257}{298.674}{\Character{106}}%
\TextAt{13}{411.302}{298.674}{K}%
\TextAt{8}{426.84}{298.674}{be}%
\TextAt{8}{443.587}{298.674}{a}%
\TextAt{8}{454.166}{298.674}{finite}%
\TextAt{8}{484.901}{298.674}{extension}%
\TextAt{8}{79.226}{281.628}{and}%
\TextAt{14}{103.049}{281.628}{o}%
\TextAt{8}{108.893}{281.628}{,}%
\TextAt{8}{117.367}{281.628}{resp.}%
\TextAt{16}{149.87}{281.628}{\Character{79}}%
\TextAt{8}{159.719}{281.628}{,}%
\TextAt{8}{168.194}{281.628}{the}%
\TextAt{8}{189.41}{281.628}{ring}%
\TextAt{8}{214.536}{281.628}{of}%
\TextAt{8}{228.926}{281.628}{integers}%
\TextAt{8}{273.3}{281.628}{of}%
\TextAt{13}{287.69}{281.628}{K}%
\TextAt{8}{298.507}{281.628}{,}%
\TextAt{8}{306.982}{281.628}{resp.}%
\TextAt{13}{339.484}{281.628}{L}%
\TextAt{8}{347.448}{281.628}{.}%
\TextAt{8}{359.077}{281.628}{In}%
\TextAt{8}{374.758}{281.628}{§2}%
\TextAt{8}{391.252}{281.628}{we}%
\TextAt{8}{409.54}{281.628}{considered}%
\TextAt{12}{467.889}{281.628}{Dedekind’s}%
\TextAt{12}{79.226}{264.582}{complementary}%
\TextAt{12}{173.803}{264.582}{module}%
\TextAt{14}{191.525}{235.581}{C}%
\TextAt{21}{198.848}{233.721}{L}%
\TextAt{3}{204.607}{233.721}{\Character{106}}%
\TextAt{21}{206.96}{233.721}{K}%
\TextAt{9}{218.483}{235.581}{\Character{61}}%
\TextAt{16}{230.909}{235.581}{\Character{102}}%
\TextAt{13}{236.886}{235.581}{x}%
\NativeMathAt{246.859}{235.581}{11.955}{\in}%
\TextAt{13}{258.15}{235.581}{L}%
\TextAt{16}{269.436}{235.581}{\Character{106}}%
\TextAt{8}{276.077}{235.581}{Tr}%
\TextAt{9}{289.084}{235.581}{\Character{40}}%
\TextAt{13}{293.637}{235.581}{x}%
\TextAt{16}{300.289}{235.581}{\Character{79}}%
\TextAt{9}{310.139}{235.581}{\Character{41}}%
\NativeMathAt{318.012}{235.581}{11.955}{\subseteq}%
\TextAt{14}{330.631}{235.581}{o}%
\TextAt{16}{336.475}{235.581}{\Character{103}}%
\NativeMathAt{345.774}{235.581}{11.955}{\cong}%
\TextAt{8}{358.393}{235.581}{Hom}%
\TextAt{27}{382.782}{233.788}{o}%
\TextAt{9}{387.176}{235.581}{\Character{40}}%
\TextAt{16}{391.728}{235.581}{\Character{79}}%
\TextAt{13}{401.578}{235.581}{\Character{59}}%
\TextAt{14}{406.822}{235.581}{o}%
\TextAt{9}{412.666}{235.581}{\Character{41}}%
\TextAt{13}{417.219}{235.581}{\Character{58}}%
\TextAt{8}{79.226}{206.58}{It}%
\TextAt{8}{91.899}{206.58}{is}%
\TextAt{8}{103.662}{206.58}{a}%
\TextAt{8}{113.418}{206.58}{fractional}%
\TextAt{8}{165.781}{206.58}{ideal}%
\TextAt{8}{193.745}{206.58}{in}%
\TextAt{13}{207.397}{206.58}{L}%
\TextAt{8}{219.259}{206.58}{whose}%
\TextAt{8}{253.786}{206.58}{inverse}%
\TextAt{8}{292.544}{206.58}{is}%
\TextAt{8}{304.308}{206.58}{the}%
\TextAt{8}{324.465}{206.58}{different}%
\TextAt{14}{370.635}{206.58}{D}%
\TextAt{21}{380.577}{204.72}{L}%
\TextAt{3}{386.336}{204.72}{\Character{106}}%
\TextAt{21}{388.688}{204.72}{K}%
\TextAt{8}{396.891}{206.58}{.}%
\TextAt{8}{405.339}{206.58}{From}%
\TextAt{8}{436.064}{206.58}{(2.6),}%
\TextAt{8}{467.291}{206.58}{it}%
\TextAt{8}{478.995}{206.58}{is}%
\TextAt{8}{490.759}{206.58}{divisible}%
\TextAt{8}{79.226}{189.534}{only}%
\TextAt{8}{104.918}{189.534}{by}%
\TextAt{8}{121.177}{189.534}{the}%
\TextAt{8}{141.333}{189.534}{prime}%
\TextAt{8}{174.497}{189.534}{ideals}%
\TextAt{8}{207.075}{189.534}{of}%
\TextAt{13}{220.405}{189.534}{L}%
\TextAt{8}{232.267}{189.534}{which}%
\TextAt{8}{265.753}{189.534}{are}%
\TextAt{8}{285.264}{189.534}{ramified}%
\TextAt{8}{330.789}{189.534}{over}%
\TextAt{13}{355.835}{189.534}{K}%
\TextAt{8}{366.653}{189.534}{.}%
\TextAt{12}{79.226}{163.522}{Definition}%
\TextAt{12}{142.015}{163.522}{(3.10).}%
\TextAt{8}{184.061}{163.522}{The}%
\TextAt{8}{208.115}{163.522}{fractional}%
\TextAt{8}{260.478}{163.522}{ideal}%
\TextAt{13}{241.188}{134.52}{\Character{33}}%
\TextAt{21}{248.478}{132.727}{K}%
\TextAt{9}{260.001}{134.52}{\Character{61}}%
\TextAt{14}{272.427}{134.52}{C}%
\TextAt{21}{279.75}{132.661}{K}%
\TextAt{3}{287.454}{132.661}{\Character{106}}%
\TextAt{25}{289.807}{132.661}{Q}%
\NativeMathAt{299.825}{134.52}{11.955}{\cong}%
\TextAt{8}{312.444}{134.52}{Hom}%
\TextAt{25}{336.832}{132.727}{Z}%
\TextAt{9}{342.644}{134.52}{\Character{40}}%
\TextAt{14}{347.196}{134.52}{o}%
\TextAt{13}{353.041}{134.52}{\Character{59}}%
\TextAt{18}{358.285}{134.52}{Z}%
\TextAt{9}{366.255}{134.52}{\Character{41}}%
\TextAt{8}{79.226}{105.519}{is}%
\TextAt{8}{90.99}{105.519}{called}%
\TextAt{8}{124.154}{105.519}{the}%
\TextAt{12}{144.31}{105.519}{canonical}%
\TextAt{12}{202.986}{105.519}{module}%
\TextAt{8}{250.018}{105.519}{of}%
\TextAt{8}{263.348}{105.519}{the}%
\TextAt{8}{283.504}{105.519}{number}%
\TextAt{8}{326.1}{105.519}{field}%
\TextAt{13}{351.457}{105.519}{K}%
\TextAt{8}{362.275}{105.519}{.}%
\TextAt{8}{97.159}{79.507}{By}%
\TextAt{8}{115.522}{79.507}{(2.2)}%
\TextAt{8}{143.497}{79.507}{we}%
\TextAt{8}{160.725}{79.507}{have}%
\TextAt{8}{187.719}{79.507}{the}%
\TextAt{8}{297.211}{43.033}{200}%
\Destination{78.23}{809.815}{nav.page.200}%
\Destination{79.36}{472.955}{nav.print.219}%
\Destination{79.36}{708.555}{nav.number.3.3.9}%
\Destination{79.36}{177.555}{nav.number.3.3.10}%
\Hotlink{379.937}{279.455}{6.858}{11.047}{nav.section.3.2}%
\Hotlink{275.489}{524.655}{25.078}{11.047}{nav.number.3.3.7}%
\Hotlink{435.564}{204.355}{25.078}{11.047}{nav.number.3.2.6}%
\Hotlink{115.022}{77.355}{25.078}{11.047}{nav.number.3.2.2}%
}

% Page 213 (body page 201)
\Page{%
\TextAt{2}{79.226}{796.943}{III.}%
\TextAt{2}{98.317}{796.943}{Riemann–Roch}%
\TextAt{2}{174.375}{796.943}{Theory}%
\Vector{q 0 G 0 g 0.299 w 7.226 18.729 m 460.769 18.729 l S Q}%
\TextAt{12}{79.226}{761.124}{Proposition}%
\TextAt{12}{151.937}{761.124}{(3.11).}%
\TextAt{8}{193.984}{761.124}{The}%
\TextAt{8}{218.038}{761.124}{canonical}%
\TextAt{8}{269.421}{761.124}{modules}%
\TextAt{8}{315.329}{761.124}{of}%
\TextAt{13}{328.659}{761.124}{L}%
\TextAt{8}{340.521}{761.124}{and}%
\TextAt{13}{363.283}{761.124}{K}%
\TextAt{8}{377.998}{761.124}{satisfy}%
\TextAt{8}{414.545}{761.124}{the}%
\TextAt{8}{434.702}{761.124}{relation}%
\TextAt{13}{270.161}{732.123}{\Character{33}}%
\TextAt{21}{277.45}{730.329}{L}%
\TextAt{9}{287.028}{732.123}{\Character{61}}%
\TextAt{14}{299.454}{732.123}{C}%
\TextAt{21}{306.777}{730.263}{L}%
\TextAt{3}{312.536}{730.263}{\Character{106}}%
\TextAt{21}{314.888}{730.263}{K}%
\TextAt{13}{323.091}{732.123}{\Character{33}}%
\TextAt{21}{330.38}{730.329}{K}%
\TextAt{13}{338.583}{732.123}{\Character{58}}%
\TextAt{8}{97.159}{694.155}{The}%
\TextAt{8}{121.727}{694.155}{canonical}%
\TextAt{8}{173.625}{694.155}{module}%
\TextAt{13}{215.432}{694.155}{\Character{33}}%
\TextAt{21}{222.722}{692.362}{K}%
\TextAt{8}{235.336}{694.155}{is}%
\TextAt{8}{247.615}{694.155}{related}%
\TextAt{8}{287.151}{694.155}{to}%
\TextAt{8}{301.975}{694.155}{the}%
\TextAt{8}{322.646}{694.155}{Euler–Minkowski}%
\TextAt{8}{414.761}{694.155}{characteristic}%
\TextAt{13}{487.21}{694.155}{\Character{31}}%
\TextAt{9}{494.544}{694.155}{\Character{40}}%
\TextAt{14}{499.096}{694.155}{o}%
\TextAt{9}{504.94}{694.155}{\Character{41}}%
\TextAt{8}{513.904}{694.155}{and}%
\TextAt{8}{79.226}{677.109}{the}%
\TextAt{8}{99.383}{677.109}{genus}%
\TextAt{13}{131.638}{677.109}{g}%
\TextAt{8}{141.569}{677.109}{of}%
\TextAt{13}{154.899}{677.109}{K}%
\TextAt{8}{169.614}{677.109}{in}%
\TextAt{8}{183.267}{677.109}{the}%
\TextAt{8}{203.424}{677.109}{following}%
\TextAt{8}{252.858}{677.109}{way,}%
\TextAt{8}{278.873}{677.109}{by}%
\TextAt{8}{295.132}{677.109}{formula}%
\TextAt{8}{338.063}{677.109}{(3.3):}%
\TextAt{8}{264.938}{648.108}{vol}%
\TextAt{9}{280.229}{648.108}{\Character{40}}%
\TextAt{14}{284.781}{648.108}{o}%
\TextAt{9}{290.625}{648.108}{\Character{41}}%
\TextAt{9}{298.48}{648.108}{\Character{61}}%
\TextAt{20}{310.924}{658.913}{\Character{112}}%
\Vector{q 0 G 0 g 0.478 w 250.879 -114.848 m 271.806 -114.848 l S Q}%
\TextAt{16}{322.879}{648.108}{\Character{106}}%
\TextAt{13}{326.2}{648.108}{d}%
\TextAt{21}{332.283}{646.315}{K}%
\TextAt{16}{340.485}{648.108}{\Character{106}}%
\TextAt{13}{343.806}{648.108}{\Character{58}}%
\TextAt{12}{79.226}{610.14}{Proposition}%
\TextAt{12}{151.937}{610.14}{(3.12).}%
\TextAt{8}{193.984}{610.14}{deg}%
\TextAt{13}{213.538}{610.14}{\Character{33}}%
\TextAt{21}{220.828}{608.347}{K}%
\TextAt{9}{232.351}{610.14}{\Character{61}}%
\TextAt{16}{244.777}{610.14}{\Character{0}}%
\TextAt{9}{254.075}{610.14}{2}%
\TextAt{13}{259.928}{610.14}{\Character{31}}%
\TextAt{9}{267.262}{610.14}{\Character{40}}%
\TextAt{14}{271.814}{610.14}{o}%
\TextAt{9}{277.658}{610.14}{\Character{41}}%
\TextAt{9}{285.513}{610.14}{\Character{61}}%
\TextAt{9}{297.934}{610.14}{2}%
\TextAt{13}{303.81}{610.14}{g}%
\TextAt{16}{312.501}{610.14}{\Character{0}}%
\TextAt{9}{324.456}{610.14}{2}%
\TextAt{13}{330.309}{610.14}{\Character{96}}%
\TextAt{9}{335.221}{610.14}{\Character{40}}%
\TextAt{14}{339.773}{610.14}{o}%
\TextAt{9}{345.618}{610.14}{\Character{41}}%
\TextAt{8}{350.17}{610.14}{.}%
\TextAt{19}{79.226}{575.162}{Proof.}%
\TextAt{8}{116}{575.162}{By}%
\TextAt{8}{135.476}{575.162}{(2.9)}%
\TextAt{8}{164.564}{575.162}{we}%
\TextAt{8}{182.904}{575.162}{know}%
\TextAt{8}{214.586}{575.162}{that}%
\TextAt{13}{241.068}{575.162}{N}%
\TextAt{21}{250.442}{573.302}{K}%
\TextAt{3}{258.147}{573.302}{\Character{106}}%
\TextAt{25}{260.499}{573.302}{Q}%
\TextAt{9}{267.196}{575.162}{\Character{40}}%
\TextAt{14}{271.748}{575.162}{D}%
\TextAt{21}{281.69}{573.302}{K}%
\TextAt{3}{289.395}{573.302}{\Character{106}}%
\TextAt{25}{291.747}{573.302}{Q}%
\TextAt{9}{298.444}{575.162}{\Character{41}}%
\TextAt{8}{308.006}{575.162}{is}%
\TextAt{8}{320.883}{575.162}{the}%
\TextAt{8}{342.152}{575.162}{discriminant}%
\TextAt{8}{410.645}{575.162}{ideal}%
\TextAt{14}{439.72}{575.162}{d}%
\TextAt{21}{445.681}{573.302}{K}%
\TextAt{3}{453.385}{573.302}{\Character{106}}%
\TextAt{25}{455.738}{573.302}{Q}%
\TextAt{9}{467.652}{575.162}{\Character{61}}%
\TextAt{9}{481.962}{575.162}{\Character{40}}%
\TextAt{13}{486.526}{575.162}{d}%
\TextAt{21}{492.609}{573.368}{K}%
\TextAt{9}{500.812}{575.162}{\Character{41}}%
\TextAt{8}{505.364}{575.162}{,}%
\TextAt{8}{513.904}{575.162}{and}%
\TextAt{8}{79.226}{558.116}{therefore}%
\TextAt{8}{128.326}{558.116}{by}%
\TextAt{8}{144.585}{558.116}{(1.6),}%
\TextAt{8}{175.812}{558.116}{(iii):}%
\TextAt{14}{192.281}{529.115}{N}%
\TextAt{9}{202.229}{529.115}{\Character{40}}%
\TextAt{13}{206.781}{529.115}{\Character{33}}%
\TextAt{21}{214.07}{527.321}{K}%
\TextAt{9}{222.273}{529.115}{\Character{41}}%
\TextAt{9}{230.128}{529.115}{\Character{61}}%
\TextAt{14}{242.572}{529.115}{N}%
\TextAt{9}{252.52}{529.115}{\Character{40}}%
\TextAt{14}{257.072}{529.115}{D}%
\TextAt{21}{267.014}{527.255}{K}%
\TextAt{3}{274.718}{527.255}{\Character{106}}%
\TextAt{25}{277.071}{527.255}{Q}%
\TextAt{9}{283.768}{529.115}{\Character{41}}%
\TextAt{3}{288.32}{534.051}{\Character{0}}%
\TextAt{17}{294.906}{534.051}{1}%
\TextAt{9}{302.96}{529.115}{\Character{61}}%
\TextAt{14}{315.385}{529.115}{N}%
\TextAt{9}{325.333}{529.115}{\Character{40}}%
\TextAt{14}{329.886}{529.115}{d}%
\TextAt{21}{335.846}{527.255}{K}%
\TextAt{3}{343.551}{527.255}{\Character{106}}%
\TextAt{25}{345.903}{527.255}{Q}%
\TextAt{9}{352.6}{529.115}{\Character{41}}%
\TextAt{3}{357.152}{534.051}{\Character{0}}%
\TextAt{17}{363.739}{534.051}{1}%
\TextAt{9}{371.792}{529.115}{\Character{61}}%
\TextAt{16}{384.217}{529.115}{\Character{106}}%
\TextAt{13}{387.538}{529.115}{d}%
\TextAt{21}{393.621}{527.321}{K}%
\TextAt{16}{401.824}{529.115}{\Character{106}}%
\TextAt{3}{405.145}{534.051}{\Character{0}}%
\TextAt{17}{411.731}{534.051}{1}%
\TextAt{13}{416.463}{529.115}{\Character{59}}%
\TextAt{8}{79.226}{500.113}{so}%
\TextAt{8}{93.596}{500.113}{that,}%
\TextAt{8}{122.217}{500.113}{as}%
\TextAt{8}{136.587}{500.113}{vol}%
\TextAt{9}{151.878}{500.113}{\Character{40}}%
\TextAt{14}{156.43}{500.113}{o}%
\TextAt{9}{162.274}{500.113}{\Character{41}}%
\TextAt{9}{170.129}{500.113}{\Character{61}}%
\TextAt{20}{182.573}{510.335}{\Character{112}}%
\Vector{q 0 G 0 g 0.478 w 122.528 -263.426 m 143.455 -263.426 l S Q}%
\TextAt{16}{194.528}{500.113}{\Character{106}}%
\TextAt{13}{197.849}{500.113}{d}%
\TextAt{21}{203.932}{498.32}{K}%
\TextAt{16}{212.134}{500.113}{\Character{106}}%
\TextAt{8}{215.455}{500.113}{,}%
\TextAt{8}{222.604}{500.113}{we}%
\TextAt{8}{239.832}{500.113}{have}%
\TextAt{8}{266.827}{500.113}{indeed}%
\TextAt{8}{186.991}{471.112}{deg}%
\TextAt{13}{206.545}{471.112}{\Character{33}}%
\TextAt{21}{213.834}{469.319}{K}%
\TextAt{9}{225.358}{471.112}{\Character{61}}%
\TextAt{16}{237.783}{471.112}{\Character{0}}%
\TextAt{8}{249.074}{471.112}{log}%
\TextAt{14}{266.035}{471.112}{N}%
\TextAt{9}{275.983}{471.112}{\Character{40}}%
\TextAt{13}{280.535}{471.112}{\Character{33}}%
\TextAt{21}{287.824}{469.319}{K}%
\TextAt{9}{296.027}{471.112}{\Character{41}}%
\TextAt{9}{303.882}{471.112}{\Character{61}}%
\TextAt{8}{316.326}{471.112}{log}%
\TextAt{16}{333.286}{471.112}{\Character{106}}%
\TextAt{13}{336.607}{471.112}{d}%
\TextAt{21}{342.69}{469.319}{K}%
\TextAt{16}{350.892}{471.112}{\Character{106}}%
\TextAt{9}{357.534}{471.112}{\Character{61}}%
\TextAt{9}{369.944}{471.112}{2}%
\TextAt{8}{377.805}{471.112}{log}%
\TextAt{8}{394.757}{471.112}{vol}%
\TextAt{9}{410.056}{471.112}{\Character{40}}%
\TextAt{14}{414.609}{471.112}{o}%
\TextAt{9}{420.453}{471.112}{\Character{41}}%
\TextAt{9}{225.358}{451.078}{\Character{61}}%
\TextAt{16}{237.783}{451.078}{\Character{0}}%
\TextAt{9}{247.082}{451.078}{2}%
\TextAt{13}{252.935}{451.078}{\Character{31}}%
\TextAt{9}{260.268}{451.078}{\Character{40}}%
\TextAt{14}{264.821}{451.078}{o}%
\TextAt{9}{270.665}{451.078}{\Character{41}}%
\TextAt{9}{278.52}{451.078}{\Character{61}}%
\TextAt{9}{290.941}{451.078}{2}%
\TextAt{13}{296.817}{451.078}{g}%
\TextAt{16}{305.508}{451.078}{\Character{0}}%
\TextAt{9}{317.463}{451.078}{2}%
\TextAt{13}{323.316}{451.078}{\Character{96}}%
\TextAt{9}{328.228}{451.078}{\Character{40}}%
\TextAt{14}{332.78}{451.078}{o}%
\TextAt{9}{338.624}{451.078}{\Character{41}}%
\TextAt{13}{343.177}{451.078}{\Character{58}}%
\Vector{q 0 G 0 g 0.399 w 452.334 -351.924 m 452.334 -343.854 l S Q}%
\Vector{q 0 G 0 g 0.399 w 452.534 -344.053 m 459.707 -344.053 l S Q}%
\Vector{q 0 G 0 g 0.399 w 452.534 -351.724 m 459.707 -351.724 l S Q}%
\Vector{q 0 G 0 g 0.399 w 459.906 -351.924 m 459.906 -343.854 l S Q}%
\TextAt{8}{97.159}{393.075}{As}%
\TextAt{8}{115.047}{393.075}{for}%
\TextAt{8}{133.533}{393.075}{the}%
\TextAt{8}{154.291}{393.075}{genus,}%
\TextAt{8}{190.549}{393.075}{we}%
\TextAt{8}{208.378}{393.075}{now}%
\TextAt{8}{233.368}{393.075}{obtain}%
\TextAt{8}{270.396}{393.075}{the}%
\TextAt{8}{291.154}{393.075}{following}%
\TextAt{8}{341.19}{393.075}{analogue}%
\TextAt{8}{390.58}{393.075}{of}%
\TextAt{8}{404.511}{393.075}{the}%
\TextAt{12}{425.269}{393.075}{Riemann–Hurwitz}%
\TextAt{12}{79.226}{376.029}{formula}%
\TextAt{8}{128.063}{376.029}{of}%
\TextAt{8}{141.393}{376.029}{function}%
\TextAt{8}{187.241}{376.029}{theory.}%
\TextAt{12}{79.226}{350.017}{Proposition}%
\TextAt{12}{152.418}{350.017}{(3.13).}%
\TextAt{8}{194.883}{350.017}{Let}%
\TextAt{13}{216.27}{350.017}{L}%
\TextAt{16}{228.268}{350.017}{\Character{106}}%
\TextAt{13}{235.623}{350.017}{K}%
\TextAt{8}{250.756}{350.017}{be}%
\TextAt{8}{267.098}{350.017}{a}%
\TextAt{8}{277.272}{350.017}{finite}%
\TextAt{8}{307.602}{350.017}{extension}%
\TextAt{8}{359.786}{350.017}{and}%
\TextAt{13}{382.967}{350.017}{g}%
\TextAt{21}{388.572}{348.224}{L}%
\TextAt{8}{394.829}{350.017}{,}%
\TextAt{8}{402.501}{350.017}{resp.}%
\TextAt{13}{433.077}{350.017}{g}%
\TextAt{21}{438.683}{348.224}{K}%
\TextAt{8}{446.885}{350.017}{,}%
\TextAt{8}{454.557}{350.017}{the}%
\TextAt{8}{475.132}{350.017}{genus}%
\TextAt{8}{507.805}{350.017}{of}%
\TextAt{13}{521.553}{350.017}{L}%
\TextAt{8}{529.518}{350.017}{,}%
\TextAt{8}{79.226}{332.971}{resp.}%
\TextAt{13}{108.548}{332.971}{K}%
\TextAt{8}{119.366}{332.971}{.}%
\TextAt{8}{127.814}{332.971}{Then}%
\TextAt{8}{158.372}{332.971}{one}%
\TextAt{8}{179.831}{332.971}{has}%
\TextAt{13}{186.705}{303.97}{g}%
\TextAt{21}{192.311}{302.177}{L}%
\TextAt{16}{201.225}{303.97}{\Character{0}}%
\TextAt{13}{213.18}{303.97}{\Character{96}}%
\TextAt{9}{218.092}{303.97}{\Character{40}}%
\TextAt{14}{222.644}{303.97}{o}%
\TextAt{21}{228.488}{302.177}{L}%
\TextAt{9}{234.746}{303.97}{\Character{41}}%
\TextAt{9}{242.601}{303.97}{\Character{61}}%
\TextAt{9}{255.022}{303.97}{\Character{91}}%
\TextAt{13}{258.296}{303.97}{L}%
\TextAt{9}{269.581}{303.97}{\Character{58}}%
\TextAt{13}{276.154}{303.97}{K}%
\TextAt{9}{286.971}{303.97}{\Character{93}\Character{40}}%
\TextAt{13}{294.775}{303.97}{g}%
\TextAt{21}{300.381}{302.177}{K}%
\TextAt{16}{311.24}{303.97}{\Character{0}}%
\TextAt{13}{323.195}{303.97}{\Character{96}}%
\TextAt{9}{328.107}{303.97}{\Character{40}}%
\TextAt{14}{332.66}{303.97}{o}%
\TextAt{21}{338.504}{302.177}{K}%
\TextAt{9}{346.706}{303.97}{\Character{41}\Character{41}}%
\TextAt{9}{358.446}{303.97}{\Character{43}}%
\TextAt{17}{371.425}{308.677}{1}%
\Vector{q 0 G 0 g 0.478 w 299.425 -467.041 m 303.659 -467.041 l S Q}%
\TextAt{17}{371.425}{299.847}{2}%
\TextAt{8}{378.847}{303.97}{deg}%
\TextAt{14}{398.401}{303.97}{C}%
\TextAt{21}{405.725}{302.11}{L}%
\TextAt{3}{411.484}{302.11}{\Character{106}}%
\TextAt{21}{413.836}{302.11}{K}%
\TextAt{13}{422.039}{303.97}{\Character{58}}%
\TextAt{8}{79.226}{274.969}{In}%
\TextAt{8}{93.847}{274.969}{particular,}%
\TextAt{8}{151.089}{274.969}{in}%
\TextAt{8}{164.741}{274.969}{the}%
\TextAt{8}{184.898}{274.969}{case}%
\TextAt{8}{209.669}{274.969}{of}%
\TextAt{8}{222.999}{274.969}{an}%
\TextAt{8}{239.258}{274.969}{unramified}%
\TextAt{8}{297.791}{274.969}{extension}%
\TextAt{13}{349.556}{274.969}{L}%
\TextAt{16}{360.842}{274.969}{\Character{106}}%
\TextAt{13}{367.484}{274.969}{K}%
\TextAt{8}{378.301}{274.969}{:}%
\TextAt{13}{249.397}{245.968}{\Character{31}}%
\TextAt{9}{256.731}{245.968}{\Character{40}}%
\TextAt{14}{261.283}{245.968}{o}%
\TextAt{21}{267.127}{244.174}{L}%
\TextAt{9}{273.384}{245.968}{\Character{41}}%
\TextAt{9}{281.239}{245.968}{\Character{61}}%
\TextAt{9}{293.66}{245.968}{\Character{91}}%
\TextAt{13}{296.935}{245.968}{L}%
\TextAt{9}{308.22}{245.968}{\Character{58}}%
\TextAt{13}{314.792}{245.968}{K}%
\TextAt{9}{325.61}{245.968}{\Character{93}}%
\TextAt{13}{328.862}{245.968}{\Character{31}}%
\TextAt{9}{336.195}{245.968}{\Character{40}}%
\TextAt{14}{340.748}{245.968}{o}%
\TextAt{21}{346.592}{244.174}{K}%
\TextAt{9}{354.795}{245.968}{\Character{41}}%
\TextAt{13}{359.347}{245.968}{\Character{58}}%
\TextAt{19}{79.226}{199.034}{Proof.}%
\TextAt{8}{116}{199.034}{Since}%
\TextAt{13}{146.558}{199.034}{\Character{33}}%
\TextAt{21}{153.847}{197.24}{L}%
\TextAt{9}{163.425}{199.034}{\Character{61}}%
\TextAt{14}{175.85}{199.034}{C}%
\TextAt{21}{183.173}{197.174}{L}%
\TextAt{3}{188.932}{197.174}{\Character{106}}%
\TextAt{21}{191.285}{197.174}{K}%
\TextAt{13}{199.488}{199.034}{\Character{33}}%
\TextAt{21}{206.777}{197.24}{K}%
\TextAt{8}{214.979}{199.034}{,}%
\TextAt{8}{222.129}{199.034}{one}%
\TextAt{8}{243.588}{199.034}{has}%
\TextAt{14}{174.428}{170.033}{N}%
\TextAt{9}{184.377}{170.033}{\Character{40}}%
\TextAt{13}{188.929}{170.033}{\Character{33}}%
\TextAt{21}{196.218}{168.239}{L}%
\TextAt{9}{202.475}{170.033}{\Character{41}}%
\TextAt{9}{210.33}{170.033}{\Character{61}}%
\TextAt{14}{222.774}{170.033}{N}%
\TextAt{9}{232.722}{170.033}{\Character{40}}%
\TextAt{13}{237.274}{170.033}{i}%
\TextAt{21}{241.268}{168.173}{L}%
\TextAt{3}{247.027}{168.173}{\Character{106}}%
\TextAt{21}{249.379}{168.173}{K}%
\TextAt{13}{257.582}{170.033}{\Character{33}}%
\TextAt{21}{264.871}{168.239}{K}%
\TextAt{9}{273.074}{170.033}{\Character{41}}%
\TextAt{14}{277.626}{170.033}{N}%
\TextAt{9}{287.574}{170.033}{\Character{40}}%
\TextAt{14}{292.127}{170.033}{C}%
\TextAt{21}{299.45}{168.173}{L}%
\TextAt{3}{305.209}{168.173}{\Character{106}}%
\TextAt{21}{307.561}{168.173}{K}%
\TextAt{9}{315.764}{170.033}{\Character{41}}%
\TextAt{9}{323.619}{170.033}{\Character{61}}%
\TextAt{14}{336.062}{170.033}{N}%
\TextAt{9}{346.01}{170.033}{\Character{40}}%
\TextAt{13}{350.563}{170.033}{\Character{33}}%
\TextAt{21}{357.852}{168.239}{K}%
\TextAt{9}{366.055}{170.033}{\Character{41}}%
\TextAt{17}{370.607}{174.969}{\Character{91}}%
\TextAt{21}{372.959}{174.969}{L}%
\TextAt{17}{378.718}{174.969}{\Character{58}}%
\TextAt{21}{381.071}{174.969}{K}%
\TextAt{17}{388.775}{174.969}{\Character{93}}%
\TextAt{14}{391.626}{170.033}{N}%
\TextAt{9}{401.574}{170.033}{\Character{40}}%
\TextAt{14}{406.126}{170.033}{C}%
\TextAt{21}{413.449}{168.173}{L}%
\TextAt{3}{419.208}{168.173}{\Character{106}}%
\TextAt{21}{421.561}{168.173}{K}%
\TextAt{9}{429.763}{170.033}{\Character{41}}%
\TextAt{13}{434.316}{170.033}{\Character{59}}%
\TextAt{8}{79.226}{141.031}{so}%
\TextAt{8}{93.596}{141.031}{that}%
\TextAt{8}{215.035}{123.985}{deg}%
\TextAt{13}{234.589}{123.985}{\Character{33}}%
\TextAt{21}{241.878}{122.192}{L}%
\TextAt{9}{251.456}{123.985}{\Character{61}}%
\TextAt{9}{263.865}{123.985}{\Character{91}}%
\TextAt{13}{267.134}{123.985}{L}%
\TextAt{9}{278.419}{123.985}{\Character{58}}%
\TextAt{13}{284.991}{123.985}{K}%
\TextAt{9}{295.809}{123.985}{\Character{93}}%
\TextAt{8}{301.053}{123.985}{deg}%
\TextAt{13}{320.608}{123.985}{\Character{33}}%
\TextAt{21}{327.897}{122.192}{K}%
\TextAt{9}{338.756}{123.985}{\Character{43}}%
\TextAt{8}{350.518}{123.985}{deg}%
\TextAt{14}{370.072}{123.985}{C}%
\TextAt{21}{377.395}{122.126}{L}%
\TextAt{3}{383.154}{122.126}{\Character{106}}%
\TextAt{21}{385.507}{122.126}{K}%
\TextAt{13}{393.71}{123.985}{\Character{58}}%
\TextAt{8}{79.226}{100.464}{Thus}%
\TextAt{8}{108.875}{100.464}{the}%
\TextAt{8}{129.031}{100.464}{proposition}%
\TextAt{8}{190.565}{100.464}{follows}%
\TextAt{8}{229.001}{100.464}{from}%
\TextAt{8}{256.641}{100.464}{(3.12).}%
\Vector{q 0 G 0 g 0.399 w 452.334 -673.536 m 452.334 -665.467 l S Q}%
\Vector{q 0 G 0 g 0.399 w 452.534 -665.666 m 459.707 -665.666 l S Q}%
\Vector{q 0 G 0 g 0.399 w 452.534 -673.337 m 459.707 -673.337 l S Q}%
\Vector{q 0 G 0 g 0.399 w 459.906 -673.536 m 459.906 -665.467 l S Q}%
\TextAt{8}{297.211}{43.033}{201}%
\Destination{78.23}{809.815}{nav.page.201}%
\Destination{79.36}{408.055}{nav.print.220}%
\Destination{79.36}{775.155}{nav.number.3.3.11}%
\Destination{79.36}{624.155}{nav.number.3.3.12}%
\Destination{79.36}{364.055}{nav.number.3.3.13}%
\Hotlink{337.563}{674.955}{25.078}{11.047}{nav.number.3.3.3}%
\Hotlink{134.976}{572.955}{25.078}{11.047}{nav.number.3.2.9}%
\Hotlink{144.085}{555.955}{25.078}{11.047}{nav.number.3.1.6}%
\Hotlink{256.141}{98.255}{30.936}{11.047}{nav.number.3.3.12}%
}

% Page 214 (body page 202)
\Page{%
\TextAt{2}{79.226}{796.943}{III.}%
\TextAt{2}{98.317}{796.943}{Riemann–Roch}%
\TextAt{2}{174.375}{796.943}{Theory}%
\Vector{q 0 G 0 g 0.299 w 7.226 18.729 m 460.769 18.729 l S Q}%
\TextAt{8}{97.159}{761.124}{The}%
\TextAt{8}{121.521}{761.124}{Riemann–Hurwitz}%
\TextAt{8}{218.571}{761.124}{formula}%
\TextAt{8}{261.811}{761.124}{tells}%
\TextAt{8}{286.89}{761.124}{us}%
\TextAt{8}{302.215}{761.124}{in}%
\TextAt{8}{316.176}{761.124}{particular}%
\TextAt{8}{370.474}{761.124}{that,}%
\TextAt{8}{399.481}{761.124}{in}%
\TextAt{8}{413.442}{761.124}{the}%
\TextAt{8}{433.907}{761.124}{decision}%
\TextAt{8}{478.497}{761.124}{we}%
\TextAt{8}{496.033}{761.124}{took}%
\TextAt{8}{523.014}{761.124}{in}%
\TextAt{8}{79.226}{744.078}{§1,}%
\TextAt{8}{97.914}{744.078}{we}%
\TextAt{8}{115.143}{744.078}{really}%
\TextAt{8}{147.339}{744.078}{had}%
\TextAt{8}{170.103}{744.078}{no}%
\TextAt{8}{186.364}{744.078}{choice}%
\TextAt{8}{221.155}{744.078}{but}%
\TextAt{8}{242.616}{744.078}{to}%
\TextAt{8}{256.927}{744.078}{consider}%
\TextAt{8}{302.514}{744.078}{the}%
\TextAt{8}{322.672}{744.078}{extension}%
\TextAt{18}{374.439}{744.078}{C}%
\TextAt{16}{386.396}{744.078}{\Character{106}}%
\TextAt{18}{393.041}{744.078}{R}%
\TextAt{8}{405.574}{744.078}{as}%
\TextAt{19}{419.945}{744.078}{unramified}%
\TextAt{8}{473.6}{744.078}{.}%
\TextAt{8}{482.052}{744.078}{In}%
\TextAt{8}{496.675}{744.078}{fact,}%
\TextAt{8}{523.014}{744.078}{in}%
\TextAt{8}{79.226}{727.032}{function}%
\TextAt{8}{125.24}{727.032}{theory}%
\TextAt{8}{162.155}{727.032}{the}%
\TextAt{8}{182.477}{727.032}{module}%
\TextAt{8}{223.936}{727.032}{corresponding}%
\TextAt{8}{299.287}{727.032}{by}%
\TextAt{8}{315.712}{727.032}{analogy}%
\TextAt{8}{359.143}{727.032}{to}%
\TextAt{8}{373.618}{727.032}{the}%
\TextAt{8}{393.94}{727.032}{ideal}%
\TextAt{14}{422.069}{727.032}{C}%
\TextAt{21}{429.392}{725.172}{L}%
\TextAt{3}{435.151}{725.172}{\Character{106}}%
\TextAt{21}{437.503}{725.172}{K}%
\TextAt{8}{449.769}{727.032}{takes}%
\TextAt{8}{479.918}{727.032}{account}%
\TextAt{8}{523.337}{727.032}{of}%
\TextAt{8}{79.226}{709.986}{precisely}%
\TextAt{8}{127.395}{709.986}{the}%
\TextAt{8}{147.863}{709.986}{branch}%
\TextAt{8}{186.874}{709.986}{points}%
\TextAt{8}{222.369}{709.986}{of}%
\TextAt{8}{236.011}{709.986}{the}%
\TextAt{8}{256.48}{709.986}{covering}%
\TextAt{8}{302.651}{709.986}{of}%
\TextAt{8}{316.293}{709.986}{Riemann}%
\TextAt{8}{366.195}{709.986}{surfaces}%
\TextAt{8}{410.526}{709.986}{in}%
\TextAt{8}{424.491}{709.986}{question.}%
\TextAt{8}{476.542}{709.986}{In}%
\TextAt{8}{491.475}{709.986}{order}%
\TextAt{8}{522.357}{709.986}{to}%
\TextAt{8}{79.226}{692.94}{obtain}%
\TextAt{8}{116.272}{692.94}{the}%
\TextAt{8}{137.046}{692.94}{same}%
\TextAt{8}{166.991}{692.94}{phenomenon}%
\TextAt{8}{235.897}{692.94}{in}%
\TextAt{8}{250.168}{692.94}{number}%
\TextAt{8}{293.383}{692.94}{theory}%
\TextAt{8}{330.751}{692.94}{we}%
\TextAt{8}{348.597}{692.94}{are}%
\TextAt{8}{368.726}{692.94}{forced}%
\TextAt{8}{404.134}{692.94}{to}%
\TextAt{8}{419.062}{692.94}{declare}%
\TextAt{8}{459.348}{692.94}{all}%
\TextAt{8}{476.225}{692.94}{the}%
\TextAt{8}{497}{692.94}{infinite}%
\TextAt{8}{79.226}{675.894}{primes}%
\TextAt{14}{117.004}{675.894}{P}%
\TextAt{8}{130.795}{675.894}{of}%
\TextAt{13}{144.125}{675.894}{L}%
\TextAt{8}{155.987}{675.894}{unramified,}%
\TextAt{8}{217.771}{675.894}{since}%
\TextAt{8}{246.44}{675.894}{they}%
\TextAt{8}{272.777}{675.894}{do}%
\TextAt{8}{289.036}{675.894}{not}%
\TextAt{8}{309.85}{675.894}{occur}%
\TextAt{8}{341.388}{675.894}{in}%
\TextAt{8}{355.041}{675.894}{the}%
\TextAt{8}{375.197}{675.894}{ideal}%
\TextAt{14}{403.16}{675.894}{C}%
\TextAt{21}{410.483}{674.034}{L}%
\TextAt{3}{416.242}{674.034}{\Character{106}}%
\TextAt{21}{418.595}{674.034}{K}%
\TextAt{8}{426.797}{675.894}{.}%
\TextAt{8}{97.159}{655.859}{Thus}%
\TextAt{8}{126.992}{655.859}{the}%
\TextAt{8}{147.333}{655.859}{fact}%
\TextAt{8}{170.603}{655.859}{that}%
\TextAt{18}{196.157}{655.859}{C}%
\TextAt{16}{208.426}{655.859}{\Character{106}}%
\TextAt{18}{215.382}{655.859}{R}%
\TextAt{8}{228.099}{655.859}{is}%
\TextAt{8}{240.047}{655.859}{unramified}%
\TextAt{8}{298.764}{655.859}{appears}%
\TextAt{8}{342.262}{655.859}{to}%
\TextAt{8}{356.757}{655.859}{be}%
\TextAt{8}{372.866}{655.859}{forced}%
\TextAt{8}{407.84}{655.859}{by}%
\TextAt{8}{424.283}{655.859}{nature}%
\TextAt{8}{461.541}{655.859}{itself.}%
\TextAt{8}{494.991}{655.859}{Investi-}%
\TextAt{8}{79.226}{638.813}{gating}%
\TextAt{8}{115.472}{638.813}{the}%
\TextAt{8}{136.093}{638.813}{matter}%
\TextAt{8}{174.933}{638.813}{a}%
\TextAt{8}{185.153}{638.813}{little}%
\TextAt{8}{213.581}{638.813}{more}%
\TextAt{8}{243.311}{638.813}{closely,}%
\TextAt{8}{283.619}{638.813}{however,}%
\TextAt{8}{332.331}{638.813}{this}%
\TextAt{8}{355.617}{638.813}{turns}%
\TextAt{8}{386.711}{638.813}{out}%
\TextAt{8}{407.989}{638.813}{not}%
\TextAt{8}{429.267}{638.813}{to}%
\TextAt{8}{444.042}{638.813}{be}%
\TextAt{8}{460.43}{638.813}{the}%
\TextAt{8}{481.051}{638.813}{case.}%
\TextAt{8}{511.766}{638.813}{It}%
\TextAt{8}{524.903}{638.813}{is}%
\TextAt{8}{79.226}{621.767}{rather}%
\TextAt{8}{114.177}{621.767}{a}%
\TextAt{8}{123.758}{621.767}{consequence}%
\TextAt{8}{189.649}{621.767}{of}%
\TextAt{8}{202.805}{621.767}{a}%
\TextAt{8}{212.386}{621.767}{well-hidden}%
\TextAt{8}{274.308}{621.767}{initial}%
\TextAt{8}{307.955}{621.767}{choice}%
\TextAt{8}{342.571}{621.767}{that}%
\TextAt{8}{367.766}{621.767}{we}%
\TextAt{8}{384.82}{621.767}{made.}%
\TextAt{8}{420.527}{621.767}{In}%
\TextAt{8}{434.975}{621.767}{fact,}%
\TextAt{8}{461.173}{621.767}{in}%
\TextAt{8}{474.652}{621.767}{chap.}%
\TextAt{8}{506.785}{621.767}{I,}%
\TextAt{8}{517.981}{621.767}{§5,}%
\TextAt{8}{79.226}{604.721}{we}%
\TextAt{8}{96.454}{604.721}{equipped}%
\TextAt{8}{146.522}{604.721}{the}%
\TextAt{8}{166.678}{604.721}{Minkowski}%
\TextAt{8}{224.948}{604.721}{space}%
\TextAt{13}{266.55}{567.097}{K}%
\TextAt{25}{276.509}{565.304}{R}%
\TextAt{9}{286.085}{567.097}{\Character{61}}%
\TextAt{20}{298.51}{587.541}{\Character{34}}%
\TextAt{20}{305.484}{578.455}{\Character{89}}%
\TextAt{21}{310.783}{553.349}{\Character{28}}%
\TextAt{18}{322.753}{567.097}{C}%
\TextAt{20}{331.387}{587.541}{\Character{35}}%
\TextAt{17}{338.361}{584.864}{\Character{43}}%
\TextAt{8}{79.226}{531.134}{with}%
\TextAt{8}{105.886}{531.134}{the}%
\TextAt{8}{126.043}{531.134}{“canonical}%
\TextAt{8}{182.949}{531.134}{metric”}%
\TextAt{16}{261.632}{514.088}{\Character{104}}%
\TextAt{13}{266.281}{514.088}{x\Character{59}}%
\TextAt{13}{278.153}{514.088}{y}%
\TextAt{16}{284.314}{514.088}{\Character{105}}%
\TextAt{9}{292.284}{514.088}{\Character{61}}%
\TextAt{20}{304.709}{525.445}{\Character{88}}%
\TextAt{21}{311.005}{500.339}{\Character{28}}%
\TextAt{13}{323.971}{514.088}{x}%
\TextAt{21}{330.623}{512.295}{\Character{28}}%
\Vector{q 0 G 0 g 0.478 w 263.799 -253.091 m 269.936 -253.091 l S Q}%
\TextAt{13}{335.799}{514.088}{y}%
\TextAt{21}{341.936}{511.265}{\Character{28}}%
\TextAt{13}{347.112}{514.088}{\Character{58}}%
\TextAt{8}{79.226}{483.354}{Replacing}%
\TextAt{8}{133.371}{483.354}{it,}%
\TextAt{8}{148.327}{483.354}{for}%
\TextAt{8}{166.212}{483.354}{instance,}%
\TextAt{8}{215.049}{483.354}{by}%
\TextAt{8}{231.308}{483.354}{the}%
\TextAt{8}{251.464}{483.354}{“Minkowski}%
\TextAt{8}{315.257}{483.354}{metric”}%
\TextAt{9}{255.418}{454.581}{\Character{40}}%
\TextAt{13}{259.97}{454.581}{x\Character{59}}%
\TextAt{13}{271.841}{454.581}{y}%
\TextAt{9}{278.003}{454.581}{\Character{41}}%
\TextAt{9}{285.858}{454.581}{\Character{61}}%
\TextAt{20}{298.301}{465.939}{\Character{88}}%
\TextAt{21}{304.597}{440.833}{\Character{28}}%
\TextAt{13}{317.563}{454.581}{\Character{11}}%
\TextAt{21}{325.008}{452.788}{\Character{28}}%
\TextAt{13}{330.185}{454.581}{x}%
\TextAt{21}{336.837}{452.788}{\Character{28}}%
\Vector{q 0 G 0 g 0.478 w 270.013 -312.598 m 276.15 -312.598 l S Q}%
\TextAt{13}{342.013}{454.581}{y}%
\TextAt{21}{348.15}{451.759}{\Character{28}}%
\TextAt{13}{353.326}{454.581}{\Character{59}}%
\TextAt{13}{79.226}{417.072}{\Character{11}}%
\TextAt{21}{86.672}{415.278}{\Character{28}}%
\TextAt{9}{95.57}{417.072}{\Character{61}}%
\TextAt{9}{108.386}{417.072}{1}%
\TextAt{8}{118.383}{417.072}{if}%
\TextAt{13}{129.343}{417.072}{\Character{28}}%
\TextAt{9}{139.484}{417.072}{\Character{61}}%
\Vector{q 0 G 0 g 0.478 w 80.311 -350.107 m 86.73 -350.107 l S Q}%
\TextAt{13}{152.311}{417.072}{\Character{28}}%
\TextAt{8}{158.73}{417.072}{,}%
\TextAt{13}{166.173}{417.072}{\Character{11}}%
\TextAt{21}{173.619}{415.278}{\Character{28}}%
\TextAt{9}{182.518}{417.072}{\Character{61}}%
\TextAt{17}{196.54}{421.779}{1}%
\Vector{q 0 G 0 g 0.478 w 124.54 -353.94 m 128.774 -353.94 l S Q}%
\TextAt{17}{196.54}{412.949}{2}%
\TextAt{8}{206.103}{417.072}{if}%
\TextAt{13}{217.062}{417.072}{\Character{28}}%
\NativeMathAt{227.203}{417.072}{11.955}{\neq}%
\Vector{q 0 G 0 g 0.478 w 168.03 -350.107 m 174.449 -350.107 l S Q}%
\TextAt{13}{240.03}{417.072}{\Character{28}}%
\TextAt{8}{246.449}{417.072}{,}%
\TextAt{8}{253.893}{417.072}{changes}%
\TextAt{8}{297.442}{417.072}{the}%
\TextAt{8}{317.834}{417.072}{whole}%
\TextAt{8}{351.233}{417.072}{picture.}%
\TextAt{8}{396.158}{417.072}{The}%
\TextAt{8}{420.447}{417.072}{Haar}%
\TextAt{8}{449.626}{417.072}{measures}%
\TextAt{8}{500.061}{417.072}{on}%
\TextAt{13}{516.556}{417.072}{K}%
\TextAt{25}{526.515}{415.278}{R}%
\TextAt{8}{79.226}{400.026}{belonging}%
\TextAt{8}{132.235}{400.026}{to}%
\TextAt{16}{146.546}{400.026}{\Character{104}}%
\TextAt{13}{155.092}{400.026}{\Character{59}}%
\TextAt{16}{164.234}{400.026}{\Character{105}}%
\TextAt{8}{172.781}{400.026}{and}%
\TextAt{9}{195.543}{400.026}{\Character{40}}%
\TextAt{13}{203.993}{400.026}{\Character{59}}%
\TextAt{9}{213.134}{400.026}{\Character{41}}%
\TextAt{8}{221.584}{400.026}{are}%
\TextAt{8}{241.095}{400.026}{related}%
\TextAt{8}{280.117}{400.026}{as}%
\TextAt{8}{294.487}{400.026}{follows:}%
\TextAt{8}{218.003}{371.253}{vol}%
\TextAt{7}{233.293}{369.46}{canonical}%
\TextAt{9}{268.119}{371.253}{\Character{40}}%
\TextAt{13}{272.671}{371.253}{X}%
\TextAt{9}{283.326}{371.253}{\Character{41}}%
\TextAt{9}{291.181}{371.253}{\Character{61}}%
\TextAt{9}{303.602}{371.253}{2}%
\TextAt{21}{309.478}{376.19}{s}%
\TextAt{8}{315.884}{371.253}{vol}%
\TextAt{7}{331.175}{369.46}{Minkowski}%
\TextAt{9}{370.982}{371.253}{\Character{40}}%
\TextAt{13}{375.534}{371.253}{X}%
\TextAt{9}{386.189}{371.253}{\Character{41}}%
\TextAt{13}{390.741}{371.253}{\Character{58}}%
\TextAt{8}{79.226}{342.481}{Distinguishing}%
\TextAt{8}{158.091}{342.481}{the}%
\TextAt{8}{179.75}{342.481}{invariants}%
\TextAt{8}{234.991}{342.481}{of}%
\TextAt{8}{249.824}{342.481}{Riemann–Roch}%
\TextAt{8}{332.957}{342.481}{theory}%
\TextAt{8}{371.21}{342.481}{with}%
\TextAt{8}{399.372}{342.481}{respect}%
\TextAt{8}{440.925}{342.481}{to}%
\TextAt{8}{456.738}{342.481}{the}%
\TextAt{8}{478.397}{342.481}{Minkowski}%
\TextAt{8}{79.226}{325.435}{measure}%
\TextAt{8}{124.811}{325.435}{by}%
\TextAt{8}{141.07}{325.435}{a}%
\TextAt{8}{150.826}{325.435}{tilde,}%
\TextAt{8}{180.738}{325.435}{we}%
\TextAt{8}{197.965}{325.435}{get}%
\TextAt{8}{217.476}{325.435}{the}%
\TextAt{8}{237.632}{325.435}{relations}%
\TextAt{20}{192.604}{296.663}{\Character{101}}%
\TextAt{13}{191.608}{296.663}{\Character{31}}%
\TextAt{9}{198.941}{296.663}{\Character{40}}%
\TextAt{14}{203.494}{296.663}{a}%
\TextAt{9}{209.477}{296.663}{\Character{41}}%
\TextAt{9}{217.332}{296.663}{\Character{61}}%
\TextAt{13}{229.775}{296.663}{\Character{31}}%
\TextAt{9}{237.109}{296.663}{\Character{40}}%
\TextAt{14}{241.661}{296.663}{a}%
\TextAt{9}{247.644}{296.663}{\Character{41}}%
\TextAt{9}{254.841}{296.663}{\Character{43}}%
\TextAt{8}{266.615}{296.663}{log}%
\TextAt{9}{283.575}{296.663}{2}%
\TextAt{21}{289.428}{301.599}{s}%
\TextAt{13}{293.842}{296.663}{\Character{59}}%
\TextAt{20}{323.432}{299.818}{\Character{101}}%
\TextAt{13}{322.997}{296.663}{\Character{96}}%
\TextAt{9}{327.909}{296.663}{\Character{40}}%
\TextAt{14}{332.461}{296.663}{a}%
\TextAt{9}{338.444}{296.663}{\Character{41}}%
\TextAt{9}{346.299}{296.663}{\Character{61}}%
\TextAt{13}{358.743}{296.663}{\Character{96}}%
\TextAt{9}{363.655}{296.663}{\Character{40}}%
\TextAt{14}{368.207}{296.663}{a}%
\TextAt{9}{374.19}{296.663}{\Character{41}}%
\TextAt{9}{381.387}{296.663}{\Character{43}}%
\TextAt{8}{393.161}{296.663}{log}%
\TextAt{9}{410.121}{296.663}{2}%
\TextAt{21}{415.974}{301.599}{s}%
\TextAt{8}{79.226}{267.891}{(the}%
\TextAt{8}{103.91}{267.891}{latter}%
\TextAt{8}{135.756}{267.891}{in}%
\TextAt{8}{149.382}{267.891}{case}%
\TextAt{8}{174.126}{267.891}{that}%
\TextAt{13}{199.467}{267.891}{H}%
\TextAt{17}{210.09}{272.23}{0}%
\TextAt{9}{214.822}{267.891}{\Character{40}}%
\TextAt{14}{219.375}{267.891}{a}%
\TextAt{9}{225.358}{267.891}{\Character{41}}%
\NativeMathAt{233.231}{267.891}{11.955}{\neq}%
\TextAt{18}{245.656}{267.891}{\Character{63}}%
\TextAt{8}{254.955}{267.891}{),}%
\TextAt{8}{266.637}{267.891}{whereas}%
\TextAt{8}{310.892}{267.891}{the}%
\TextAt{8}{331.021}{267.891}{genus}%
\TextAt{8}{363.249}{267.891}{remains}%
\TextAt{8}{406.859}{267.891}{unchanged.}%
\TextAt{8}{469.61}{267.891}{Substituting}%
\TextAt{8}{79.226}{250.845}{this}%
\TextAt{8}{103.325}{250.845}{into}%
\TextAt{8}{128.344}{250.845}{the}%
\TextAt{8}{149.777}{250.845}{Riemann–Hurwitz}%
\TextAt{8}{247.794}{250.845}{formula}%
\TextAt{8}{292.001}{250.845}{(3.13)}%
\TextAt{8}{327.111}{250.845}{preserves}%
\TextAt{8}{378.587}{250.845}{its}%
\TextAt{8}{396.182}{250.845}{shape}%
\TextAt{8}{430.358}{250.845}{only}%
\TextAt{8}{457.326}{250.845}{if}%
\TextAt{8}{469.326}{250.845}{one}%
\TextAt{8}{492.062}{250.845}{enriches}%
\TextAt{14}{79.226}{233.799}{C}%
\TextAt{21}{86.549}{231.939}{L}%
\TextAt{3}{92.308}{231.939}{\Character{106}}%
\TextAt{21}{94.661}{231.939}{K}%
\TextAt{8}{107.335}{233.799}{into}%
\TextAt{8}{131.653}{233.799}{a}%
\TextAt{8}{141.982}{233.799}{replete}%
\TextAt{8}{180.921}{233.799}{ideal}%
\TextAt{8}{209.459}{233.799}{in}%
\TextAt{8}{223.686}{233.799}{which}%
\TextAt{8}{257.747}{233.799}{all}%
\TextAt{8}{274.58}{233.799}{infinite}%
\TextAt{8}{314.822}{233.799}{primes}%
\TextAt{14}{353.175}{233.799}{P}%
\TextAt{8}{367.54}{233.799}{such}%
\TextAt{8}{394.511}{233.799}{that}%
\TextAt{13}{420.455}{233.799}{L}%
\TextAt{27}{428.419}{232.006}{P}%
\NativeMathAt{439.813}{233.799}{11.955}{\neq}%
\TextAt{13}{453.217}{233.799}{K}%
\TextAt{27}{463.176}{232.006}{p}%
\TextAt{8}{472.135}{233.799}{occur.}%
\TextAt{8}{509.947}{233.799}{This}%
\TextAt{8}{79.226}{216.753}{forces}%
\TextAt{8}{112.619}{216.753}{us}%
\TextAt{8}{128.128}{216.753}{to}%
\TextAt{8}{142.931}{216.753}{consider}%
\TextAt{8}{189.009}{216.753}{the}%
\TextAt{8}{209.658}{216.753}{extension}%
\TextAt{18}{261.917}{216.753}{C}%
\TextAt{16}{274.711}{216.753}{\Character{106}}%
\TextAt{18}{282.193}{216.753}{R}%
\TextAt{8}{295.217}{216.753}{as}%
\TextAt{19}{310.08}{216.753}{ramified}%
\TextAt{8}{350.883}{216.753}{,}%
\TextAt{8}{358.648}{216.753}{to}%
\TextAt{8}{373.451}{216.753}{put}%
\TextAt{20}{395.446}{216.753}{\Character{101}}%
\TextAt{13}{395.403}{216.753}{e}%
\TextAt{27}{400.829}{214.894}{P}%
\TextAt{3}{407.424}{214.893}{\Character{106}}%
\TextAt{27}{409.777}{214.894}{p}%
\TextAt{9}{418.424}{216.753}{\Character{61}}%
\TextAt{9}{431.682}{216.753}{\Character{91}}%
\TextAt{13}{434.941}{216.753}{L}%
\TextAt{27}{442.905}{214.96}{P}%
\TextAt{9}{454.159}{216.753}{\Character{58}}%
\TextAt{13}{461.572}{216.753}{K}%
\TextAt{27}{471.531}{214.96}{p}%
\TextAt{9}{476.018}{216.753}{\Character{93}}%
\TextAt{8}{479.269}{216.753}{,}%
\TextAt{20}{489.188}{219.908}{\Character{101}}%
\TextAt{13}{487.034}{216.753}{f}%
\TextAt{27}{492.804}{214.893}{P}%
\TextAt{3}{499.4}{214.893}{\Character{106}}%
\TextAt{27}{501.752}{214.893}{p}%
\TextAt{9}{510.4}{216.753}{\Character{61}}%
\TextAt{9}{523.646}{216.753}{1}%
\TextAt{8}{529.518}{216.753}{,}%
\TextAt{8}{79.226}{199.707}{and}%
\TextAt{8}{101.989}{199.707}{in}%
\TextAt{8}{115.642}{199.707}{particular}%
\TextAt{20}{241.341}{182.661}{\Character{101}}%
\TextAt{13}{241.299}{182.661}{e}%
\TextAt{27}{246.724}{180.868}{p}%
\TextAt{9}{254.532}{182.661}{\Character{61}}%
\TextAt{9}{266.942}{182.661}{\Character{91}}%
\TextAt{13}{270.209}{182.661}{K}%
\TextAt{27}{280.168}{180.868}{p}%
\TextAt{9}{287.976}{182.661}{\Character{58}}%
\TextAt{18}{294.548}{182.661}{R}%
\TextAt{9}{303.183}{182.661}{\Character{93}}%
\TextAt{13}{306.434}{182.661}{\Character{59}}%
\TextAt{20}{337.742}{185.816}{\Character{101}}%
\TextAt{13}{335.589}{182.661}{f}%
\TextAt{27}{341.359}{180.868}{p}%
\TextAt{9}{349.167}{182.661}{\Character{61}}%
\TextAt{9}{361.577}{182.661}{1}%
\TextAt{13}{367.445}{182.661}{\Character{58}}%
\TextAt{8}{79.226}{159.238}{The}%
\TextAt{8}{103.28}{159.238}{following}%
\TextAt{8}{152.715}{159.238}{modifications}%
\TextAt{8}{224.649}{159.238}{ensue}%
\TextAt{8}{256.569}{159.238}{from}%
\TextAt{8}{284.209}{159.238}{this.}%
\TextAt{8}{311.583}{159.238}{For}%
\TextAt{8}{332.552}{159.238}{an}%
\TextAt{8}{348.811}{159.238}{infinite}%
\TextAt{8}{388.478}{159.238}{prime}%
\TextAt{14}{421.642}{159.238}{p}%
\TextAt{8}{431.523}{159.238}{one}%
\TextAt{8}{452.982}{159.238}{has}%
\TextAt{8}{473.856}{159.238}{to}%
\TextAt{8}{488.166}{159.238}{define}%
\TextAt{20}{178.024}{130.466}{\Character{101}}%
\TextAt{13}{177.976}{130.466}{v}%
\TextAt{27}{183.635}{128.672}{p}%
\TextAt{9}{188.122}{130.466}{\Character{40}}%
\TextAt{13}{192.674}{130.466}{a}%
\TextAt{9}{198.819}{130.466}{\Character{41}}%
\TextAt{9}{206.674}{130.466}{\Character{61}}%
\TextAt{16}{219.118}{130.466}{\Character{0}}%
\TextAt{20}{228.459}{130.466}{\Character{101}}%
\TextAt{13}{228.416}{130.466}{e}%
\TextAt{27}{233.842}{128.672}{p}%
\TextAt{8}{240.321}{130.466}{log}%
\TextAt{16}{257.282}{130.466}{\Character{106}}%
\TextAt{13}{260.603}{130.466}{\Character{28}}%
\TextAt{13}{267.011}{130.466}{a}%
\TextAt{16}{273.166}{130.466}{\Character{106}}%
\TextAt{13}{276.487}{130.466}{\Character{59}}%
\TextAt{14}{305.642}{130.466}{p}%
\TextAt{21}{311.625}{135.402}{\Character{23}}%
\TextAt{9}{320.128}{130.466}{\Character{61}}%
\TextAt{13}{332.553}{130.466}{e}%
\TextAt{21}{337.979}{135.402}{\Character{23}}%
\TextAt{17}{342.663}{135.402}{\Character{47}}%
\TextAt{30}{346.979}{135.402}{\Character{101}}%
\TextAt{21}{346.897}{135.402}{e}%
\TextAt{28}{350.825}{134.406}{p}%
\TextAt{13}{355.406}{130.466}{\Character{59}}%
\TextAt{20}{386.214}{133.584}{\Character{101}}%
\TextAt{14}{384.56}{130.466}{N}%
\TextAt{9}{394.509}{130.466}{\Character{40}}%
\TextAt{14}{399.061}{130.466}{p}%
\TextAt{9}{405.044}{130.466}{\Character{41}}%
\TextAt{9}{412.899}{130.466}{\Character{61}}%
\TextAt{13}{425.343}{130.466}{e\Character{58}}%
\TextAt{8}{79.226}{101.693}{The}%
\TextAt{8}{103.28}{101.693}{absolute}%
\TextAt{8}{149.523}{101.693}{norm}%
\TextAt{8}{180.092}{101.693}{as}%
\TextAt{8}{194.462}{101.693}{well}%
\TextAt{8}{218.193}{101.693}{as}%
\TextAt{8}{232.563}{101.693}{the}%
\TextAt{8}{252.72}{101.693}{degree}%
\TextAt{8}{289.135}{101.693}{of}%
\TextAt{8}{302.465}{101.693}{a}%
\TextAt{8}{312.22}{101.693}{replete}%
\TextAt{8}{350.585}{101.693}{ideal}%
\TextAt{14}{378.548}{101.693}{a}%
\TextAt{8}{388.428}{101.693}{remain}%
\TextAt{8}{427.45}{101.693}{unaltered:}%
\TextAt{20}{183.499}{76.04}{\Character{101}}%
\TextAt{14}{181.846}{72.921}{N}%
\TextAt{9}{191.794}{72.921}{\Character{40}}%
\TextAt{14}{196.346}{72.921}{a}%
\TextAt{9}{202.329}{72.921}{\Character{41}}%
\TextAt{9}{210.184}{72.921}{\Character{61}}%
\TextAt{14}{222.628}{72.921}{N}%
\TextAt{9}{232.576}{72.921}{\Character{40}}%
\TextAt{14}{237.129}{72.921}{a}%
\TextAt{9}{243.112}{72.921}{\Character{41}}%
\TextAt{13}{247.664}{72.921}{\Character{59}}%
\TextAt{20}{276.965}{76.071}{\Character{103}}%
\TextAt{8}{276.818}{72.921}{deg}%
\TextAt{9}{294.381}{72.921}{\Character{40}}%
\TextAt{14}{298.933}{72.921}{a}%
\TextAt{9}{304.916}{72.921}{\Character{41}}%
\TextAt{9}{312.771}{72.921}{\Character{61}}%
\TextAt{16}{325.215}{72.921}{\Character{0}}%
\TextAt{8}{336.506}{72.921}{log}%
\TextAt{20}{355.119}{76.04}{\Character{101}}%
\TextAt{14}{353.466}{72.921}{N}%
\TextAt{9}{363.414}{72.921}{\Character{40}}%
\TextAt{14}{367.966}{72.921}{a}%
\TextAt{9}{373.95}{72.921}{\Character{41}}%
\TextAt{9}{381.805}{72.921}{\Character{61}}%
\TextAt{8}{394.248}{72.921}{deg}%
\TextAt{9}{411.81}{72.921}{\Character{40}}%
\TextAt{14}{416.363}{72.921}{a}%
\TextAt{9}{422.346}{72.921}{\Character{41}}%
\TextAt{13}{426.898}{72.921}{\Character{58}}%
\TextAt{8}{297.211}{43.033}{202}%
\Destination{78.23}{809.815}{nav.page.202}%
\Destination{79.36}{776.155}{nav.print.221}%
\Hotlink{84.405}{741.855}{6.858}{11.047}{nav.section.3.1}%
\Hotlink{506.285}{619.555}{5.22}{11.047}{nav.chapter.1}%
\Hotlink{523.16}{619.555}{6.858}{11.047}{nav.section.1.5}%
\Hotlink{291.501}{248.655}{30.936}{11.047}{nav.number.3.3.13}%
}

% Page 215 (body page 203)
\Page{%
\TextAt{2}{79.226}{796.943}{III.}%
\TextAt{2}{98.317}{796.943}{Riemann–Roch}%
\TextAt{2}{174.375}{796.943}{Theory}%
\Vector{q 0 G 0 g 0.299 w 7.226 18.729 m 460.769 18.729 l S Q}%
\TextAt{8}{79.226}{761.124}{The}%
\TextAt{8}{103.28}{761.124}{canonical}%
\TextAt{8}{154.663}{761.124}{module}%
\TextAt{13}{195.956}{761.124}{\Character{33}}%
\TextAt{21}{203.246}{759.33}{K}%
\TextAt{8}{215.346}{761.124}{however}%
\TextAt{8}{260.225}{761.124}{has}%
\TextAt{8}{281.099}{761.124}{to}%
\TextAt{8}{295.409}{761.124}{be}%
\TextAt{8}{311.334}{761.124}{changed:}%
\TextAt{20}{246.687}{732.123}{\Character{101}}%
\TextAt{13}{246.148}{732.123}{\Character{33}}%
\TextAt{21}{253.438}{730.329}{K}%
\TextAt{9}{264.961}{732.123}{\Character{61}}%
\TextAt{13}{277.387}{732.123}{\Character{33}}%
\TextAt{21}{284.676}{730.329}{K}%
\TextAt{20}{305.805}{743.48}{\Character{89}}%
\TextAt{27}{294.871}{718.023}{p}%
\TextAt{7}{301.681}{718.023}{complex}%
\TextAt{14}{334.008}{732.123}{p}%
\TextAt{17}{339.991}{737.059}{2}%
\TextAt{7}{345.637}{737.059}{log}%
\TextAt{17}{357.863}{737.059}{2}%
\TextAt{13}{362.596}{732.123}{\Character{59}}%
\TextAt{8}{79.226}{693.914}{in}%
\TextAt{8}{92.879}{693.914}{order}%
\TextAt{8}{123.448}{693.914}{for}%
\TextAt{8}{141.333}{693.914}{the}%
\TextAt{8}{161.49}{693.914}{equation}%
\TextAt{8}{227.905}{664.913}{deg}%
\TextAt{20}{247.998}{664.913}{\Character{101}}%
\TextAt{13}{247.459}{664.913}{\Character{33}}%
\TextAt{21}{254.749}{663.119}{K}%
\TextAt{9}{266.272}{664.913}{\Character{61}}%
\TextAt{16}{278.698}{664.913}{\Character{0}}%
\TextAt{9}{287.996}{664.913}{2}%
\TextAt{20}{294.845}{664.913}{\Character{101}}%
\TextAt{13}{293.849}{664.913}{\Character{31}}%
\TextAt{9}{301.183}{664.913}{\Character{40}}%
\TextAt{14}{305.735}{664.913}{o}%
\TextAt{9}{311.579}{664.913}{\Character{41}}%
\TextAt{9}{319.434}{664.913}{\Character{61}}%
\TextAt{9}{331.855}{664.913}{2}%
\TextAt{13}{337.731}{664.913}{g}%
\TextAt{16}{346.422}{664.913}{\Character{0}}%
\TextAt{9}{358.377}{664.913}{2}%
\TextAt{20}{364.666}{668.067}{\Character{101}}%
\TextAt{13}{364.23}{664.913}{\Character{96}}%
\TextAt{9}{369.142}{664.913}{\Character{40}}%
\TextAt{14}{373.694}{664.913}{o}%
\TextAt{9}{379.539}{664.913}{\Character{41}}%
\TextAt{8}{79.226}{635.911}{to}%
\TextAt{8}{93.536}{635.911}{hold.}%
\TextAt{8}{124.102}{635.911}{By}%
\TextAt{8}{142.465}{635.911}{the}%
\TextAt{8}{162.621}{635.911}{same}%
\TextAt{8}{191.947}{635.911}{token,}%
\TextAt{8}{227.072}{635.911}{the}%
\TextAt{8}{247.228}{635.911}{ideal}%
\TextAt{14}{275.191}{635.911}{C}%
\TextAt{21}{282.514}{634.052}{L}%
\TextAt{3}{288.273}{634.052}{\Character{106}}%
\TextAt{21}{290.626}{634.052}{K}%
\TextAt{8}{302.726}{635.911}{has}%
\TextAt{8}{323.6}{635.911}{to}%
\TextAt{8}{337.91}{635.911}{be}%
\TextAt{8}{353.834}{635.911}{replaced}%
\TextAt{8}{400.005}{635.911}{by}%
\TextAt{8}{416.264}{635.911}{the}%
\TextAt{8}{436.42}{635.911}{replete}%
\TextAt{8}{474.785}{635.911}{ideal}%
\TextAt{20}{243.187}{610.029}{\Character{101}}%
\TextAt{14}{242.846}{606.91}{C}%
\TextAt{21}{250.169}{605.051}{L}%
\TextAt{3}{255.928}{605.051}{\Character{106}}%
\TextAt{21}{258.28}{605.051}{K}%
\TextAt{9}{269.804}{606.91}{\Character{61}}%
\TextAt{14}{282.229}{606.91}{C}%
\TextAt{21}{289.553}{605.051}{L}%
\TextAt{3}{295.312}{605.051}{\Character{106}}%
\TextAt{21}{297.664}{605.051}{K}%
\TextAt{20}{313.621}{618.268}{\Character{89}}%
\TextAt{27}{312.551}{592.365}{P}%
\TextAt{3}{319.146}{592.365}{\Character{106}}%
\NativeMathAt{321.497}{592.365}{7.97}{\infty}%
\TextAt{30}{307.941}{583.54}{\Character{101}}%
\TextAt{21}{307.859}{583.54}{e}%
\TextAt{28}{311.787}{581.802}{P}%
\TextAt{23}{317.486}{581.802}{\Character{106}}%
\TextAt{28}{319.756}{581.802}{p}%
\NativeMathAt{323.838}{583.54}{7.97}{\neq}%
\TextAt{17}{330.421}{583.54}{1}%
\TextAt{14}{336.652}{606.91}{P}%
\TextAt{17}{346.545}{611.847}{2}%
\TextAt{7}{352.19}{611.846}{log}%
\TextAt{17}{364.417}{611.847}{2}%
\TextAt{8}{79.226}{557.864}{so}%
\TextAt{8}{93.596}{557.864}{that}%
\TextAt{20}{254.997}{540.818}{\Character{101}}%
\TextAt{13}{254.459}{540.818}{\Character{33}}%
\TextAt{21}{261.748}{539.025}{L}%
\TextAt{9}{271.326}{540.818}{\Character{61}}%
\TextAt{20}{284.092}{543.937}{\Character{101}}%
\TextAt{14}{283.752}{540.818}{C}%
\TextAt{21}{291.075}{538.958}{L}%
\TextAt{3}{296.834}{538.958}{\Character{106}}%
\TextAt{21}{299.186}{538.958}{K}%
\TextAt{13}{309.381}{540.818}{i}%
\TextAt{21}{313.375}{538.958}{L}%
\TextAt{3}{319.134}{538.958}{\Character{106}}%
\TextAt{21}{321.486}{538.958}{K}%
\TextAt{9}{329.689}{540.818}{\Character{40}}%
\TextAt{20}{334.779}{540.818}{\Character{101}}%
\TextAt{13}{334.241}{540.818}{\Character{33}}%
\TextAt{21}{341.53}{539.025}{K}%
\TextAt{9}{349.733}{540.818}{\Character{41}}%
\TextAt{13}{354.285}{540.818}{\Character{58}}%
\TextAt{8}{79.226}{517.296}{In}%
\TextAt{8}{93.847}{517.296}{the}%
\TextAt{8}{114.004}{517.296}{same}%
\TextAt{8}{143.33}{517.296}{way}%
\TextAt{8}{167.073}{517.296}{as}%
\TextAt{8}{181.443}{517.296}{in}%
\TextAt{8}{195.096}{517.296}{(3.13),}%
\TextAt{8}{232.181}{517.296}{this}%
\TextAt{8}{255.003}{517.296}{yields}%
\TextAt{8}{287.904}{517.296}{the}%
\TextAt{8}{308.06}{517.296}{Riemann–Hurwitz}%
\TextAt{8}{404.801}{517.296}{formula}%
\TextAt{13}{186.705}{488.295}{g}%
\TextAt{21}{192.311}{486.502}{L}%
\TextAt{16}{201.225}{488.295}{\Character{0}}%
\TextAt{20}{213.616}{491.45}{\Character{101}}%
\TextAt{13}{213.18}{488.295}{\Character{96}}%
\TextAt{9}{218.092}{488.295}{\Character{40}}%
\TextAt{14}{222.644}{488.295}{o}%
\TextAt{21}{228.488}{486.502}{L}%
\TextAt{9}{234.746}{488.295}{\Character{41}}%
\TextAt{9}{242.601}{488.295}{\Character{61}}%
\TextAt{9}{255.022}{488.295}{\Character{91}}%
\TextAt{13}{258.296}{488.295}{L}%
\TextAt{9}{269.581}{488.295}{\Character{58}}%
\TextAt{13}{276.154}{488.295}{K}%
\TextAt{9}{286.971}{488.295}{\Character{93}\Character{40}}%
\TextAt{13}{294.775}{488.295}{g}%
\TextAt{21}{300.381}{486.502}{K}%
\TextAt{16}{311.24}{488.295}{\Character{0}}%
\TextAt{20}{323.631}{491.45}{\Character{101}}%
\TextAt{13}{323.195}{488.295}{\Character{96}}%
\TextAt{9}{328.107}{488.295}{\Character{40}}%
\TextAt{14}{332.66}{488.295}{o}%
\TextAt{21}{338.504}{486.502}{K}%
\TextAt{9}{346.706}{488.295}{\Character{41}\Character{41}}%
\TextAt{9}{358.446}{488.295}{\Character{43}}%
\TextAt{17}{371.425}{493.002}{1}%
\Vector{q 0 G 0 g 0.478 w 299.425 -282.716 m 303.659 -282.716 l S Q}%
\TextAt{17}{371.425}{484.173}{2}%
\TextAt{8}{378.847}{488.295}{deg}%
\TextAt{20}{398.742}{491.414}{\Character{101}}%
\TextAt{14}{398.401}{488.295}{C}%
\TextAt{21}{405.725}{486.435}{L}%
\TextAt{3}{411.484}{486.435}{\Character{106}}%
\TextAt{21}{413.836}{486.435}{K}%
\TextAt{13}{422.039}{488.295}{\Character{58}}%
\TextAt{8}{97.159}{456.305}{In}%
\TextAt{8}{112.755}{456.305}{view}%
\TextAt{8}{140.712}{456.305}{of}%
\TextAt{8}{155.017}{456.305}{this}%
\TextAt{8}{178.814}{456.305}{sensitivity}%
\TextAt{8}{235.524}{456.305}{to}%
\TextAt{8}{250.809}{456.305}{the}%
\TextAt{8}{271.94}{456.305}{chosen}%
\TextAt{8}{310.37}{456.305}{metric}%
\TextAt{8}{347.761}{456.305}{on}%
\TextAt{8}{364.994}{456.305}{Minkowski}%
\TextAt{8}{424.238}{456.305}{space}%
\TextAt{13}{456.488}{456.305}{K}%
\TextAt{25}{466.447}{454.512}{R}%
\TextAt{8}{472.701}{456.305}{,}%
\TextAt{8}{481.069}{456.305}{the}%
\TextAt{8}{502.2}{456.305}{math-}%
\TextAt{8}{79.226}{439.259}{ematician}%
\TextAt{26}{133.339}{439.259}{Uwe}%
\TextAt{26}{165.16}{439.259}{Jannsen}%
\TextAt{8}{218.077}{439.259}{proposes}%
\TextAt{8}{266.786}{439.259}{as}%
\TextAt{8}{281.937}{439.259}{analogues}%
\TextAt{8}{336.121}{439.259}{of}%
\TextAt{8}{350.232}{439.259}{the}%
\TextAt{8}{371.169}{439.259}{function}%
\TextAt{8}{417.798}{439.259}{fields}%
\TextAt{8}{448.55}{439.259}{not}%
\TextAt{8}{470.145}{439.259}{just}%
\TextAt{8}{494.071}{439.259}{number}%
\TextAt{8}{79.226}{422.213}{fields}%
\TextAt{13}{109.198}{422.213}{K}%
\TextAt{8}{123.913}{422.213}{by}%
\TextAt{8}{140.172}{422.213}{themselves}%
\TextAt{8}{198.824}{422.213}{but}%
\TextAt{8}{220.283}{422.213}{number}%
\TextAt{8}{262.88}{422.213}{fields}%
\TextAt{8}{292.851}{422.213}{equipped}%
\TextAt{8}{342.92}{422.213}{with}%
\TextAt{8}{369.58}{422.213}{a}%
\TextAt{8}{379.335}{422.213}{metric}%
\TextAt{8}{415.75}{422.213}{of}%
\TextAt{8}{429.081}{422.213}{the}%
\TextAt{8}{449.237}{422.213}{type}%
\TextAt{16}{251.219}{393.212}{\Character{104}}%
\TextAt{13}{255.869}{393.212}{x\Character{59}}%
\TextAt{13}{267.74}{393.212}{y}%
\TextAt{16}{273.902}{393.212}{\Character{105}}%
\TextAt{21}{278.551}{391.419}{K}%
\TextAt{9}{290.074}{393.212}{\Character{61}}%
\TextAt{20}{302.5}{404.57}{\Character{88}}%
\TextAt{21}{308.795}{379.464}{\Character{28}}%
\TextAt{13}{321.761}{393.212}{\Character{11}}%
\TextAt{21}{329.206}{391.419}{\Character{28}}%
\TextAt{13}{334.383}{393.212}{x}%
\TextAt{21}{341.035}{391.419}{\Character{28}}%
\Vector{q 0 G 0 g 0.478 w 274.212 -373.967 m 280.348 -373.967 l S Q}%
\TextAt{13}{346.212}{393.212}{y}%
\TextAt{21}{352.348}{390.389}{\Character{28}}%
\TextAt{13}{357.525}{393.212}{\Character{59}}%
\TextAt{13}{79.226}{356.948}{\Character{11}}%
\TextAt{21}{86.672}{355.155}{\Character{28}}%
\TextAt{13}{95.169}{356.948}{\Character{62}}%
\TextAt{9}{107.595}{356.948}{0}%
\TextAt{8}{113.448}{356.948}{,}%
\TextAt{13}{120.088}{356.948}{\Character{11}}%
\TextAt{21}{127.533}{355.155}{\Character{28}}%
\TextAt{9}{136.031}{356.948}{\Character{61}}%
\TextAt{13}{148.456}{356.948}{\Character{11}}%
\Vector{q 0 G 0 g 0.359 w 83.902 -414.158 m 89.078 -414.158 l S Q}%
\TextAt{21}{155.902}{355.155}{\Character{28}}%
\TextAt{8}{161.078}{356.948}{,}%
\TextAt{8}{167.718}{356.948}{on}%
\TextAt{13}{183.341}{356.948}{K}%
\TextAt{25}{193.3}{355.155}{R}%
\TextAt{8}{199.554}{356.948}{.}%
\TextAt{8}{207.791}{356.948}{Let}%
\TextAt{8}{228.124}{356.948}{these}%
\TextAt{8}{257.459}{356.948}{new}%
\TextAt{8}{280.876}{356.948}{objects}%
\TextAt{8}{319.967}{356.948}{be}%
\TextAt{8}{335.255}{356.948}{called}%
\TextAt{12}{367.782}{356.948}{metrized}%
\TextAt{12}{422.941}{356.948}{number}%
\TextAt{12}{471.584}{356.948}{fields}%
\TextAt{8}{501.711}{356.948}{.}%
\TextAt{8}{509.947}{356.948}{This}%
\TextAt{8}{79.226}{339.902}{idea}%
\TextAt{8}{104.255}{339.902}{does}%
\TextAt{8}{130.97}{339.902}{indeed}%
\TextAt{8}{168.349}{339.902}{do}%
\TextAt{8}{184.926}{339.902}{justice}%
\TextAt{8}{222.042}{339.902}{to}%
\TextAt{8}{236.67}{339.902}{the}%
\TextAt{8}{257.144}{339.902}{situation}%
\TextAt{8}{306.311}{339.902}{in}%
\TextAt{8}{320.281}{339.902}{question}%
\TextAt{8}{367.165}{339.902}{in}%
\TextAt{8}{381.135}{339.902}{a}%
\TextAt{8}{391.208}{339.902}{very}%
\TextAt{8}{417.218}{339.902}{precise}%
\TextAt{8}{455.96}{339.902}{manner,}%
\TextAt{8}{501.882}{339.902}{and}%
\TextAt{8}{524.963}{339.902}{it}%
\TextAt{8}{79.226}{322.856}{is}%
\TextAt{8}{91.331}{322.856}{of}%
\TextAt{8}{105.001}{322.856}{fundamental}%
\TextAt{8}{172.984}{322.856}{importance}%
\TextAt{8}{234.786}{322.856}{for}%
\TextAt{8}{253.012}{322.856}{algebraic}%
\TextAt{8}{302.787}{322.856}{number}%
\TextAt{8}{345.724}{322.856}{theory.}%
\TextAt{8}{387.066}{322.856}{We}%
\TextAt{8}{407.551}{322.856}{denote}%
\TextAt{8}{445.61}{322.856}{metrized}%
\TextAt{8}{494.071}{322.856}{number}%
\TextAt{8}{79.226}{305.81}{fields}%
\TextAt{9}{109.198}{305.81}{\Character{40}}%
\TextAt{13}{113.75}{305.81}{K}%
\TextAt{13}{123.912}{305.81}{\Character{59}}%
\TextAt{16}{129.162}{305.81}{\Character{104}}%
\TextAt{13}{137.708}{305.81}{\Character{59}}%
\TextAt{16}{146.85}{305.81}{\Character{105}}%
\TextAt{21}{151.499}{304.017}{K}%
\TextAt{9}{159.702}{305.81}{\Character{41}}%
\TextAt{8}{168.151}{305.81}{as}%
\TextAt{20}{185.26}{308.832}{\Character{101}}%
\TextAt{13}{182.522}{305.81}{K}%
\TextAt{8}{197.237}{305.81}{and}%
\TextAt{8}{219.999}{305.81}{attach}%
\TextAt{8}{256.104}{305.81}{to}%
\TextAt{8}{270.414}{305.81}{them}%
\TextAt{8}{300.326}{305.81}{the}%
\TextAt{8}{320.482}{305.81}{following}%
\TextAt{8}{369.917}{305.81}{invariants.}%
\TextAt{8}{428.206}{305.81}{Let}%
\TextAt{16}{251.219}{276.809}{\Character{104}}%
\TextAt{13}{255.869}{276.809}{x\Character{59}}%
\TextAt{13}{267.74}{276.809}{y}%
\TextAt{16}{273.902}{276.809}{\Character{105}}%
\TextAt{21}{278.551}{275.016}{K}%
\TextAt{9}{290.074}{276.809}{\Character{61}}%
\TextAt{20}{302.5}{288.166}{\Character{88}}%
\TextAt{21}{308.795}{263.06}{\Character{28}}%
\TextAt{13}{321.761}{276.809}{\Character{11}}%
\TextAt{21}{329.206}{275.016}{\Character{28}}%
\TextAt{13}{334.383}{276.809}{x}%
\TextAt{21}{341.035}{275.016}{\Character{28}}%
\Vector{q 0 G 0 g 0.478 w 274.212 -490.37 m 280.348 -490.37 l S Q}%
\TextAt{13}{346.212}{276.809}{y}%
\TextAt{21}{352.348}{273.986}{\Character{28}}%
\TextAt{13}{357.525}{276.809}{\Character{58}}%
\TextAt{8}{79.226}{240.378}{Let}%
\TextAt{14}{100.344}{240.378}{p}%
\TextAt{9}{109.901}{240.378}{\Character{61}}%
\TextAt{14}{122.58}{240.378}{p}%
\TextAt{21}{128.563}{238.584}{\Character{28}}%
\TextAt{8}{137.785}{240.378}{be}%
\TextAt{8}{153.858}{240.378}{the}%
\TextAt{8}{174.163}{240.378}{infinite}%
\TextAt{8}{213.979}{240.378}{prime}%
\TextAt{8}{247.291}{240.378}{corresponding}%
\TextAt{8}{322.625}{240.378}{to}%
\TextAt{13}{337.084}{240.378}{\Character{28}}%
\TextAt{9}{347.077}{240.378}{\Character{58}}%
\TextAt{13}{353.903}{240.378}{K}%
\NativeMathAt{368.294}{240.378}{11.955}{\rightarrow}%
\TextAt{18}{383.824}{240.378}{C}%
\TextAt{8}{392.458}{240.378}{.}%
\TextAt{8}{401.352}{240.378}{We}%
\TextAt{8}{421.645}{240.378}{then}%
\TextAt{8}{448.453}{240.378}{put}%
\TextAt{13}{470.061}{240.378}{\Character{11}}%
\TextAt{27}{477.507}{238.584}{p}%
\TextAt{9}{485.568}{240.378}{\Character{61}}%
\TextAt{13}{498.246}{240.378}{\Character{11}}%
\TextAt{21}{505.692}{238.584}{\Character{28}}%
\TextAt{8}{510.868}{240.378}{.}%
\TextAt{8}{519.762}{240.378}{At}%
\TextAt{8}{79.226}{223.332}{the}%
\TextAt{8}{99.383}{223.332}{same}%
\TextAt{8}{128.709}{223.332}{time,}%
\TextAt{8}{158.62}{223.332}{we}%
\TextAt{8}{175.848}{223.332}{also}%
\TextAt{8}{199.328}{223.332}{use}%
\TextAt{8}{219.544}{223.332}{the}%
\TextAt{8}{239.7}{223.332}{letter}%
\TextAt{14}{270.915}{223.332}{p}%
\TextAt{8}{280.796}{223.332}{for}%
\TextAt{8}{298.681}{223.332}{the}%
\TextAt{8}{318.837}{223.332}{positive}%
\TextAt{8}{362.151}{223.332}{real}%
\TextAt{8}{384.913}{223.332}{number}%
\TextAt{14}{270.619}{194.33}{p}%
\TextAt{9}{279.923}{194.33}{\Character{61}}%
\TextAt{13}{292.349}{194.33}{e}%
\TextAt{21}{297.774}{199.267}{\Character{11}}%
\TextAt{28}{303.213}{198.27}{p}%
\NativeMathAt{311.115}{194.33}{11.955}{\in}%
\TextAt{18}{322.406}{194.33}{R}%
\TextAt{3}{331.04}{199.267}{\Character{3}}%
\TextAt{17}{331.04}{191.375}{\Character{43}}%
\TextAt{13}{338.125}{194.33}{\Character{59}}%
\TextAt{8}{79.226}{165.329}{which}%
\TextAt{8}{112.217}{165.329}{we}%
\TextAt{8}{128.948}{165.329}{interpret}%
\TextAt{8}{176.906}{165.329}{as}%
\TextAt{8}{190.78}{165.329}{the}%
\TextAt{8}{210.441}{165.329}{replete}%
\TextAt{8}{248.309}{165.329}{ideal}%
\TextAt{9}{275.776}{165.329}{\Character{40}1\Character{41}}%
\TextAt{16}{292.376}{165.329}{\Character{2}}%
\TextAt{9}{303.317}{165.329}{\Character{40}1}%
\TextAt{13}{313.722}{165.329}{\Character{59}}%
\TextAt{13}{318.946}{165.329}{\Character{58}}%
\TextAt{13}{324.183}{165.329}{\Character{58}}%
\TextAt{13}{329.407}{165.329}{\Character{58}}%
\TextAt{13}{334.644}{165.329}{\Character{59}}%
\TextAt{9}{339.943}{165.329}{1}%
\TextAt{13}{345.796}{165.329}{\Character{59}}%
\TextAt{13}{351.02}{165.329}{e}%
\TextAt{21}{356.465}{169.668}{\Character{11}}%
\TextAt{28}{361.904}{168.671}{p}%
\TextAt{13}{366.485}{165.329}{\Character{59}}%
\TextAt{9}{371.729}{165.329}{1}%
\TextAt{13}{377.582}{165.329}{\Character{59}}%
\TextAt{13}{382.806}{165.329}{\Character{58}}%
\TextAt{13}{388.043}{165.329}{\Character{58}}%
\TextAt{13}{393.279}{165.329}{\Character{58}}%
\TextAt{13}{398.504}{165.329}{\Character{59}}%
\TextAt{9}{403.803}{165.329}{1\Character{41}}%
\NativeMathAt{417.529}{165.329}{11.955}{\in}%
\TextAt{13}{428.82}{165.329}{J}%
\TextAt{9}{436.397}{165.329}{\Character{40}}%
\TextAt{14}{440.95}{165.329}{o}%
\TextAt{9}{446.794}{165.329}{\Character{41}}%
\TextAt{16}{452.989}{165.329}{\Character{2}}%
\TextAt{20}{463.929}{174.296}{\Character{81}}%
\TextAt{27}{475.22}{161.842}{p}%
\TextAt{3}{479.209}{161.842}{\Character{106}}%
\NativeMathAt{481.56}{161.842}{7.97}{\infty}%
\TextAt{18}{492.521}{165.329}{R}%
\TextAt{3}{501.155}{169.668}{\Character{3}}%
\TextAt{17}{501.155}{162.374}{\Character{43}}%
\TextAt{8}{508.239}{165.329}{.}%
\TextAt{8}{516.522}{165.329}{We}%
\TextAt{8}{79.226}{148.283}{put}%
\TextAt{13}{207.674}{131.237}{e}%
\TextAt{27}{213.1}{129.444}{p}%
\TextAt{9}{220.907}{131.237}{\Character{61}}%
\TextAt{9}{233.317}{131.237}{1\Character{47}}%
\TextAt{13}{245.039}{131.237}{\Character{11}}%
\TextAt{27}{252.484}{129.444}{p}%
\TextAt{8}{280.882}{131.237}{and}%
\TextAt{13}{323.657}{131.237}{f}%
\TextAt{27}{329.427}{129.444}{p}%
\TextAt{9}{337.235}{131.237}{\Character{61}}%
\TextAt{13}{349.66}{131.237}{\Character{11}}%
\TextAt{27}{357.106}{129.444}{p}%
\TextAt{9}{361.593}{131.237}{\Character{91}}%
\TextAt{13}{364.844}{131.237}{K}%
\TextAt{27}{374.804}{129.444}{p}%
\TextAt{9}{382.611}{131.237}{\Character{58}}%
\TextAt{18}{389.184}{131.237}{R}%
\TextAt{9}{397.818}{131.237}{\Character{93}}%
\TextAt{13}{401.07}{131.237}{\Character{59}}%
\TextAt{8}{79.226}{107.716}{and}%
\TextAt{8}{101.989}{107.716}{we}%
\TextAt{8}{119.216}{107.716}{define}%
\TextAt{8}{153.025}{107.716}{the}%
\TextAt{8}{173.182}{107.716}{valuation}%
\TextAt{13}{224.254}{107.716}{v}%
\TextAt{27}{229.914}{105.922}{p}%
\TextAt{8}{238.298}{107.716}{of}%
\TextAt{13}{251.628}{107.716}{K}%
\TextAt{3}{262.446}{112.054}{\Character{3}}%
\TextAt{8}{271.075}{107.716}{associated}%
\TextAt{8}{326.81}{107.716}{to}%
\TextAt{14}{341.121}{107.716}{p}%
\TextAt{8}{351.001}{107.716}{by}%
\TextAt{13}{255.116}{78.715}{v}%
\TextAt{27}{260.776}{76.921}{p}%
\TextAt{9}{265.262}{78.715}{\Character{40}}%
\TextAt{13}{269.815}{78.715}{a}%
\TextAt{9}{275.96}{78.715}{\Character{41}}%
\TextAt{9}{283.815}{78.715}{\Character{61}}%
\TextAt{16}{296.258}{78.715}{\Character{0}}%
\TextAt{13}{305.557}{78.715}{e}%
\TextAt{27}{310.982}{76.921}{p}%
\TextAt{8}{317.462}{78.715}{log}%
\TextAt{16}{334.422}{78.715}{\Character{106}}%
\TextAt{13}{337.743}{78.715}{\Character{28}}%
\TextAt{13}{344.151}{78.715}{a}%
\TextAt{16}{350.307}{78.715}{\Character{106}}%
\TextAt{13}{353.628}{78.715}{\Character{58}}%
\TextAt{8}{297.211}{43.033}{203}%
\Destination{78.23}{809.815}{nav.page.203}%
\Destination{79.36}{437.255}{nav.print.222}%
\Hotlink{194.596}{515.055}{30.936}{11.047}{nav.number.3.3.13}%
}

% Page 216 (body page 204)
\Page{%
\TextAt{2}{79.226}{796.943}{III.}%
\TextAt{2}{98.317}{796.943}{Riemann–Roch}%
\TextAt{2}{174.375}{796.943}{Theory}%
\Vector{q 0 G 0 g 0.299 w 7.226 18.729 m 460.769 18.729 l S Q}%
\TextAt{8}{79.226}{761.124}{Further,}%
\TextAt{8}{124.907}{761.124}{we}%
\TextAt{8}{142.134}{761.124}{put}%
\TextAt{14}{201.385}{744.078}{N}%
\TextAt{9}{211.333}{744.078}{\Character{40}}%
\TextAt{14}{215.886}{744.078}{p}%
\TextAt{9}{221.869}{744.078}{\Character{41}}%
\TextAt{9}{229.724}{744.078}{\Character{61}}%
\TextAt{13}{242.167}{744.078}{e}%
\TextAt{21}{247.593}{749.014}{f}%
\TextAt{28}{251.661}{748.018}{p}%
\TextAt{8}{280.152}{744.078}{and}%
\TextAt{16}{322.928}{744.078}{\Character{106}}%
\TextAt{13}{326.249}{744.078}{a}%
\TextAt{16}{332.394}{744.078}{\Character{106}}%
\TextAt{27}{335.714}{742.284}{p}%
\TextAt{9}{343.522}{744.078}{\Character{61}}%
\TextAt{14}{355.948}{744.078}{N}%
\TextAt{9}{365.896}{744.078}{\Character{40}}%
\TextAt{14}{370.448}{744.078}{p}%
\TextAt{9}{376.431}{744.078}{\Character{41}}%
\TextAt{3}{380.983}{749.014}{\Character{0}}%
\TextAt{21}{387.57}{749.014}{v}%
\TextAt{28}{391.693}{748.018}{p}%
\TextAt{17}{395.776}{749.014}{\Character{40}}%
\TextAt{21}{399.069}{749.014}{a}%
\TextAt{17}{403.567}{749.014}{\Character{41}}%
\TextAt{13}{407.359}{744.078}{\Character{59}}%
\TextAt{8}{79.226}{720.556}{so}%
\TextAt{8}{94.209}{720.556}{that}%
\TextAt{8}{120.191}{720.556}{again}%
\TextAt{16}{152.03}{720.556}{\Character{106}}%
\TextAt{13}{155.351}{720.556}{a}%
\TextAt{16}{161.496}{720.556}{\Character{106}}%
\TextAt{27}{164.817}{718.763}{p}%
\TextAt{9}{173.669}{720.556}{\Character{61}}%
\TextAt{16}{187.138}{720.556}{\Character{106}}%
\TextAt{13}{190.459}{720.556}{\Character{28}}%
\TextAt{13}{196.867}{720.556}{a}%
\TextAt{16}{203.023}{720.556}{\Character{106}}%
\TextAt{8}{210.854}{720.556}{if}%
\TextAt{14}{222.191}{720.556}{p}%
\TextAt{8}{232.684}{720.556}{is}%
\TextAt{8}{245.06}{720.556}{real,}%
\TextAt{8}{271.841}{720.556}{and}%
\TextAt{16}{295.216}{720.556}{\Character{106}}%
\TextAt{13}{298.537}{720.556}{a}%
\TextAt{16}{304.682}{720.556}{\Character{106}}%
\TextAt{27}{308.003}{718.763}{p}%
\TextAt{9}{316.855}{720.556}{\Character{61}}%
\TextAt{16}{330.324}{720.556}{\Character{106}}%
\TextAt{13}{333.645}{720.556}{\Character{28}}%
\TextAt{13}{340.053}{720.556}{a}%
\TextAt{16}{346.209}{720.556}{\Character{106}}%
\TextAt{17}{349.53}{724.894}{2}%
\TextAt{8}{358.772}{720.556}{if}%
\TextAt{14}{370.109}{720.556}{p}%
\TextAt{8}{380.602}{720.556}{is}%
\TextAt{8}{392.979}{720.556}{complex.}%
\TextAt{8}{445.216}{720.556}{For}%
\TextAt{8}{466.798}{720.556}{every}%
\TextAt{8}{498.303}{720.556}{replete}%
\TextAt{8}{79.226}{703.51}{ideal}%
\TextAt{14}{107.291}{703.51}{a}%
\TextAt{8}{117.274}{703.51}{of}%
\TextAt{20}{133.444}{706.532}{\Character{101}}%
\TextAt{13}{130.706}{703.51}{K}%
\TextAt{8}{141.524}{703.51}{,}%
\TextAt{8}{148.8}{703.51}{there}%
\TextAt{8}{178.814}{703.51}{is}%
\TextAt{8}{190.68}{703.51}{a}%
\TextAt{8}{200.538}{703.51}{unique}%
\TextAt{8}{238.681}{703.51}{representation}%
\TextAt{14}{315.272}{703.51}{a}%
\TextAt{9}{324.75}{703.51}{\Character{61}}%
\TextAt{20}{337.35}{712.477}{\Character{81}}%
\TextAt{27}{348.641}{700.023}{p}%
\TextAt{14}{355.12}{703.51}{p}%
\TextAt{21}{361.103}{707.848}{\Character{23}}%
\TextAt{28}{365.266}{706.852}{p}%
\TextAt{8}{369.847}{703.51}{,}%
\TextAt{8}{377.123}{703.51}{which}%
\TextAt{8}{410.712}{703.51}{gives}%
\TextAt{8}{439.494}{703.51}{the}%
\TextAt{8}{459.753}{703.51}{absolute}%
\TextAt{8}{506.098}{703.51}{norm}%
\TextAt{14}{79.226}{686.464}{N}%
\TextAt{9}{89.174}{686.464}{\Character{40}}%
\TextAt{14}{93.727}{686.464}{a}%
\TextAt{9}{99.71}{686.464}{\Character{41}}%
\TextAt{9}{107.565}{686.464}{\Character{61}}%
\TextAt{20}{120.008}{695.431}{\Character{81}}%
\TextAt{27}{131.299}{682.977}{p}%
\TextAt{14}{137.779}{686.464}{N}%
\TextAt{9}{147.727}{686.464}{\Character{40}}%
\TextAt{14}{152.279}{686.464}{p}%
\TextAt{9}{158.262}{686.464}{\Character{41}}%
\TextAt{21}{162.815}{690.803}{\Character{23}}%
\TextAt{28}{166.977}{689.806}{p}%
\TextAt{8}{171.558}{686.464}{,}%
\TextAt{8}{178.707}{686.464}{and}%
\TextAt{8}{201.47}{686.464}{the}%
\TextAt{8}{221.626}{686.464}{degree}%
\TextAt{8}{249.429}{657.463}{deg}%
\TextAt{30}{268.962}{655.825}{\Character{101}}%
\TextAt{21}{266.991}{653.81}{K}%
\TextAt{9}{275.194}{657.463}{\Character{40}}%
\TextAt{14}{279.746}{657.463}{a}%
\TextAt{9}{285.729}{657.463}{\Character{41}}%
\TextAt{9}{293.584}{657.463}{\Character{61}}%
\TextAt{16}{306.028}{657.463}{\Character{0}}%
\TextAt{8}{317.319}{657.463}{log}%
\TextAt{14}{334.279}{657.463}{N}%
\TextAt{9}{344.227}{657.463}{\Character{40}}%
\TextAt{14}{348.78}{657.463}{a}%
\TextAt{9}{354.763}{657.463}{\Character{41}}%
\TextAt{13}{359.315}{657.463}{\Character{58}}%
\TextAt{8}{79.226}{628.462}{The}%
\TextAt{19}{103.28}{628.462}{canonical}%
\TextAt{19}{154.687}{628.462}{module}%
\TextAt{8}{194.235}{628.462}{of}%
\TextAt{20}{210.303}{631.484}{\Character{101}}%
\TextAt{13}{207.565}{628.462}{K}%
\TextAt{8}{222.28}{628.462}{is}%
\TextAt{8}{234.044}{628.462}{defined}%
\TextAt{8}{274.357}{628.462}{to}%
\TextAt{8}{288.667}{628.462}{be}%
\TextAt{8}{304.591}{628.462}{the}%
\TextAt{8}{324.748}{628.462}{replete}%
\TextAt{8}{363.112}{628.462}{ideal}%
\TextAt{13}{207.057}{599.461}{\Character{33}}%
\TextAt{30}{216.317}{597.822}{\Character{101}}%
\TextAt{21}{214.346}{595.808}{K}%
\TextAt{9}{225.87}{599.461}{\Character{61}}%
\TextAt{13}{238.295}{599.461}{\Character{33}}%
\TextAt{21}{245.585}{597.667}{K}%
\TextAt{16}{256.444}{599.461}{\Character{1}}%
\TextAt{13}{262.422}{599.461}{\Character{33}}%
\NativeMathAt{269.711}{597.667}{7.97}{\infty}%
\NativeMathAt{281.998}{599.461}{11.955}{\in}%
\TextAt{13}{293.289}{599.461}{J}%
\TextAt{9}{300.866}{599.461}{\Character{40}}%
\Vector{q 0 G 0 g 0.478 w 233.419 -167.183 m 239.263 -167.183 l S Q}%
\TextAt{14}{305.419}{599.461}{o}%
\TextAt{9}{311.263}{599.461}{\Character{41}}%
\TextAt{9}{319.118}{599.461}{\Character{61}}%
\TextAt{13}{331.561}{599.461}{J}%
\TextAt{9}{339.139}{599.461}{\Character{40}}%
\TextAt{14}{343.691}{599.461}{o}%
\TextAt{9}{349.535}{599.461}{\Character{41}}%
\TextAt{16}{356.744}{599.461}{\Character{2}}%
\TextAt{20}{368.699}{610.818}{\Character{89}}%
\TextAt{27}{368.933}{584.915}{p}%
\TextAt{3}{372.921}{584.915}{\Character{106}}%
\NativeMathAt{375.272}{584.915}{7.97}{\infty}%
\TextAt{18}{385.968}{599.461}{R}%
\TextAt{3}{394.602}{604.397}{\Character{3}}%
\TextAt{17}{394.602}{596.505}{\Character{43}}%
\TextAt{13}{401.687}{599.461}{\Character{59}}%
\TextAt{8}{79.226}{559.809}{where}%
\TextAt{13}{113.035}{559.809}{\Character{33}}%
\TextAt{21}{120.325}{558.016}{K}%
\TextAt{8}{132.425}{559.809}{is}%
\TextAt{8}{144.189}{559.809}{the}%
\TextAt{8}{164.345}{559.809}{inverse}%
\TextAt{8}{203.104}{559.809}{of}%
\TextAt{8}{216.434}{559.809}{the}%
\TextAt{8}{236.59}{559.809}{different}%
\TextAt{14}{282.761}{559.809}{D}%
\TextAt{21}{292.703}{557.95}{K}%
\TextAt{3}{300.407}{557.95}{\Character{106}}%
\TextAt{25}{302.759}{557.95}{Q}%
\TextAt{8}{313.354}{559.809}{of}%
\TextAt{13}{326.684}{559.809}{K}%
\TextAt{16}{340.823}{559.809}{\Character{106}}%
\TextAt{18}{347.464}{559.809}{Q}%
\TextAt{8}{356.763}{559.809}{,}%
\TextAt{8}{363.912}{559.809}{and}%
\TextAt{13}{242.945}{530.808}{\Character{33}}%
\NativeMathAt{250.234}{529.015}{7.97}{\infty}%
\TextAt{9}{262.522}{530.808}{\Character{61}}%
\TextAt{9}{274.932}{530.808}{\Character{40}}%
\TextAt{13}{279.499}{530.808}{\Character{11}}%
\TextAt{3}{287.021}{535.744}{\Character{0}}%
\TextAt{17}{293.608}{535.744}{1}%
\TextAt{27}{286.945}{527.853}{p}%
\TextAt{9}{298.34}{530.808}{\Character{41}}%
\TextAt{27}{302.892}{528.948}{p}%
\TextAt{3}{306.881}{528.948}{\Character{106}}%
\NativeMathAt{309.232}{528.948}{7.97}{\infty}%
\NativeMathAt{321.521}{530.808}{11.955}{\in}%
\TextAt{20}{332.812}{542.166}{\Character{89}}%
\TextAt{27}{333.045}{516.263}{p}%
\TextAt{3}{337.034}{516.263}{\Character{106}}%
\NativeMathAt{339.385}{516.263}{7.97}{\infty}%
\TextAt{18}{350.08}{530.808}{R}%
\TextAt{3}{358.714}{535.744}{\Character{3}}%
\TextAt{17}{358.714}{527.853}{\Character{43}}%
\TextAt{13}{365.799}{530.808}{\Character{58}}%
\TextAt{8}{79.226}{491.707}{The}%
\TextAt{8}{102.562}{491.707}{Riemann–Roch}%
\TextAt{8}{183.473}{491.707}{theory}%
\TextAt{8}{219.505}{491.707}{may}%
\TextAt{8}{244.155}{491.707}{be}%
\TextAt{8}{259.361}{491.707}{transferred}%
\TextAt{8}{318.215}{491.707}{without}%
\TextAt{8}{361.074}{491.707}{any}%
\TextAt{8}{382.472}{491.707}{problem,}%
\TextAt{8}{430.675}{491.707}{using}%
\TextAt{8}{460.585}{491.707}{the}%
\TextAt{8}{480.023}{491.707}{definitions}%
\TextAt{8}{79.226}{474.661}{given}%
\TextAt{8}{110.175}{474.661}{above,}%
\TextAt{8}{147.076}{474.661}{to}%
\TextAt{8}{161.765}{474.661}{metrized}%
\TextAt{8}{210.264}{474.661}{number}%
\TextAt{8}{253.239}{474.661}{fields}%
\TextAt{20}{286.328}{477.683}{\Character{101}}%
\TextAt{13}{283.59}{474.661}{K}%
\TextAt{9}{298.374}{474.661}{\Character{61}}%
\TextAt{9}{311.429}{474.661}{\Character{40}}%
\TextAt{13}{315.998}{474.661}{K}%
\TextAt{13}{326.16}{474.661}{\Character{59}}%
\TextAt{16}{331.409}{474.661}{\Character{104}}%
\TextAt{13}{340.335}{474.661}{\Character{59}}%
\TextAt{16}{349.856}{474.661}{\Character{105}}%
\TextAt{21}{354.505}{472.867}{K}%
\TextAt{9}{362.708}{474.661}{\Character{41}}%
\TextAt{8}{367.26}{474.661}{.}%
\TextAt{8}{376.845}{474.661}{Distinguishing}%
\TextAt{8}{454.586}{474.661}{their}%
\TextAt{8}{482.928}{474.661}{invariants}%
\TextAt{8}{79.226}{457.615}{by}%
\TextAt{8}{95.485}{457.615}{the}%
\TextAt{8}{115.642}{457.615}{suffix}%
\TextAt{20}{149.332}{460.637}{\Character{101}}%
\TextAt{13}{146.594}{457.615}{K}%
\TextAt{8}{161.309}{457.615}{yields}%
\TextAt{8}{194.209}{457.615}{the}%
\TextAt{8}{214.366}{457.615}{relations}%
\TextAt{8}{236.424}{428.614}{vol}%
\TextAt{30}{253.685}{426.975}{\Character{101}}%
\TextAt{21}{251.715}{424.961}{K}%
\TextAt{9}{259.917}{428.614}{\Character{40}}%
\TextAt{13}{264.47}{428.614}{X}%
\TextAt{9}{275.125}{428.614}{\Character{41}}%
\TextAt{9}{282.98}{428.614}{\Character{61}}%
\TextAt{20}{295.424}{439.971}{\Character{89}}%
\TextAt{21}{300.722}{414.865}{\Character{28}}%
\TextAt{16}{312.692}{436.912}{\Character{112}}%
\Vector{q 0 G 0 g 0.478 w 250.655 -336.849 m 263.277 -336.849 l S Q}%
\TextAt{13}{322.655}{428.614}{\Character{11}}%
\TextAt{21}{330.1}{426.82}{\Character{28}}%
\TextAt{8}{337.269}{428.614}{vol}%
\TextAt{9}{352.56}{428.614}{\Character{40}}%
\TextAt{13}{357.112}{428.614}{X}%
\TextAt{9}{367.768}{428.614}{\Character{41}}%
\TextAt{13}{372.32}{428.614}{\Character{59}}%
\TextAt{8}{79.226}{391.752}{because}%
\TextAt{13}{121.673}{391.752}{T}%
\TextAt{21}{128.534}{389.959}{\Character{11}}%
\TextAt{9}{137.792}{391.752}{\Character{58}}%
\TextAt{9}{144.343}{391.752}{\Character{40}}%
\TextAt{13}{148.917}{391.752}{K}%
\TextAt{25}{158.876}{389.959}{R}%
\TextAt{13}{165.13}{391.752}{\Character{59}}%
\TextAt{16}{170.375}{391.752}{\Character{104}}%
\TextAt{13}{178.067}{391.752}{\Character{59}}%
\TextAt{16}{186.354}{391.752}{\Character{105}}%
\TextAt{21}{191.003}{389.959}{K}%
\TextAt{9}{199.206}{391.752}{\Character{41}}%
\NativeMathAt{207.079}{391.752}{11.955}{\rightarrow}%
\TextAt{9}{222.355}{391.752}{\Character{40}}%
\TextAt{13}{226.907}{391.752}{K}%
\TextAt{25}{236.866}{389.959}{R}%
\TextAt{13}{243.121}{391.752}{\Character{59}}%
\TextAt{16}{248.365}{391.752}{\Character{104}}%
\TextAt{13}{256.057}{391.752}{\Character{59}}%
\TextAt{16}{264.344}{391.752}{\Character{105}}%
\TextAt{9}{268.993}{391.752}{\Character{41}}%
\TextAt{8}{273.545}{391.752}{,}%
\TextAt{9}{280.011}{391.752}{\Character{40}}%
\TextAt{13}{284.563}{391.752}{x}%
\TextAt{21}{291.215}{389.959}{\Character{28}}%
\TextAt{9}{296.392}{391.752}{\Character{41}}%
\NativeArrowAt{304.265}{391.752}{11.955}{11.955}{\mapsto}%
\TextAt{9}{319.541}{391.752}{\Character{40}}%
\TextAt{16}{324.093}{399.466}{\Character{112}}%
\Vector{q 0 G 0 g 0.478 w 262.056 -374.295 m 274.678 -374.295 l S Q}%
\TextAt{13}{334.056}{391.752}{\Character{11}}%
\TextAt{21}{341.502}{389.959}{\Character{28}}%
\TextAt{13}{348.671}{391.752}{x}%
\TextAt{21}{355.323}{389.959}{\Character{28}}%
\TextAt{9}{360.499}{391.752}{\Character{41}}%
\TextAt{8}{365.052}{391.752}{,}%
\TextAt{8}{371.517}{391.752}{is}%
\TextAt{8}{382.426}{391.752}{an}%
\TextAt{8}{397.831}{391.752}{isometry}%
\TextAt{8}{444.845}{391.752}{with}%
\TextAt{8}{470.65}{391.752}{determinant}%
\TextAt{20}{79.226}{383.672}{\Character{81}}%
\TextAt{21}{90.517}{371.219}{\Character{28}}%
\TextAt{16}{97.686}{382.42}{\Character{112}}%
\Vector{q 0 G 0 g 0.478 w 35.649 -391.341 m 48.271 -391.341 l S Q}%
\TextAt{13}{107.649}{374.706}{\Character{11}}%
\TextAt{21}{115.094}{372.913}{\Character{28}}%
\TextAt{8}{120.271}{374.706}{,}%
\TextAt{8}{127.42}{374.706}{and}%
\TextAt{8}{150.183}{374.706}{therefore}%
\TextAt{13}{180.878}{345.705}{\Character{31}}%
\TextAt{30}{190.182}{344.066}{\Character{101}}%
\TextAt{21}{188.212}{342.052}{K}%
\TextAt{9}{196.415}{345.705}{\Character{40}}%
\TextAt{14}{200.967}{345.705}{o}%
\TextAt{21}{206.811}{343.911}{K}%
\TextAt{9}{215.014}{345.705}{\Character{41}}%
\TextAt{9}{222.869}{345.705}{\Character{61}}%
\TextAt{16}{235.312}{345.705}{\Character{0}}%
\TextAt{8}{246.603}{345.705}{log}%
\TextAt{8}{263.556}{345.705}{vol}%
\TextAt{30}{280.825}{344.066}{\Character{101}}%
\TextAt{21}{278.854}{342.052}{K}%
\TextAt{9}{287.057}{345.705}{\Character{40}}%
\TextAt{14}{291.609}{345.705}{o}%
\TextAt{21}{297.454}{343.911}{K}%
\TextAt{9}{305.656}{345.705}{\Character{41}}%
\TextAt{9}{313.511}{345.705}{\Character{61}}%
\TextAt{13}{325.955}{345.705}{\Character{31}}%
\TextAt{9}{333.289}{345.705}{\Character{40}}%
\TextAt{14}{337.841}{345.705}{o}%
\TextAt{21}{343.685}{343.911}{K}%
\TextAt{9}{351.888}{345.705}{\Character{41}}%
\TextAt{16}{359.097}{345.705}{\Character{0}}%
\TextAt{8}{371.052}{345.705}{log}%
\TextAt{20}{388.012}{357.062}{\Character{89}}%
\TextAt{21}{393.311}{331.956}{\Character{28}}%
\TextAt{16}{405.281}{354.003}{\Character{112}}%
\Vector{q 0 G 0 g 0.478 w 343.244 -419.758 m 355.866 -419.758 l S Q}%
\TextAt{13}{415.244}{345.705}{\Character{11}}%
\TextAt{21}{422.689}{343.911}{\Character{28}}%
\TextAt{13}{427.866}{345.705}{\Character{59}}%
\TextAt{13}{183.3}{309.213}{\Character{96}}%
\TextAt{30}{190.182}{307.575}{\Character{101}}%
\TextAt{21}{188.212}{305.56}{K}%
\TextAt{9}{196.415}{309.213}{\Character{40}}%
\TextAt{14}{200.967}{309.213}{o}%
\TextAt{21}{206.811}{307.42}{K}%
\TextAt{9}{215.014}{309.213}{\Character{41}}%
\TextAt{9}{222.869}{309.213}{\Character{61}}%
\TextAt{8}{235.312}{309.213}{log}%
\TextAt{8}{253.468}{317.301}{\char35\relax }%
\TextAt{13}{263.224}{317.301}{H}%
\TextAt{17}{273.846}{321.639}{0}%
\TextAt{9}{278.579}{317.301}{\Character{40}}%
\TextAt{14}{283.131}{317.301}{o}%
\TextAt{21}{288.975}{315.508}{K}%
\TextAt{9}{297.178}{317.301}{\Character{41}}%
\Vector{q 0 G 0 g 0.478 w 181.468 -461.798 m 229.73 -461.798 l S Q}%
\TextAt{8}{254.962}{301.012}{vol}%
\TextAt{30}{272.223}{299.374}{\Character{101}}%
\TextAt{21}{270.252}{297.359}{K}%
\TextAt{9}{278.455}{301.012}{\Character{40}}%
\TextAt{13}{283.007}{301.012}{W}%
\TextAt{9}{295.685}{301.012}{\Character{41}}%
\TextAt{9}{306.247}{309.213}{\Character{61}}%
\TextAt{13}{318.672}{309.213}{\Character{96}}%
\TextAt{9}{323.584}{309.213}{\Character{40}}%
\TextAt{14}{328.136}{309.213}{o}%
\TextAt{21}{333.981}{307.42}{K}%
\TextAt{9}{342.183}{309.213}{\Character{41}}%
\TextAt{16}{349.392}{309.213}{\Character{0}}%
\TextAt{8}{361.348}{309.213}{log}%
\TextAt{20}{378.308}{320.571}{\Character{89}}%
\TextAt{21}{383.607}{295.465}{\Character{28}}%
\TextAt{16}{395.576}{317.511}{\Character{112}}%
\Vector{q 0 G 0 g 0.478 w 333.539 -456.25 m 346.161 -456.25 l S Q}%
\TextAt{13}{405.539}{309.213}{\Character{11}}%
\TextAt{21}{412.985}{307.42}{\Character{28}}%
\TextAt{13}{418.161}{309.213}{\Character{58}}%
\TextAt{8}{79.226}{273.021}{The}%
\TextAt{8}{103.28}{273.021}{genus}%
\TextAt{13}{149.634}{238.132}{g}%
\TextAt{30}{157.21}{236.493}{\Character{101}}%
\TextAt{21}{155.239}{234.479}{K}%
\TextAt{9}{166.763}{238.132}{\Character{61}}%
\TextAt{13}{179.188}{238.131}{\Character{96}}%
\TextAt{30}{186.071}{236.493}{\Character{101}}%
\TextAt{21}{184.1}{234.479}{K}%
\TextAt{9}{192.303}{238.132}{\Character{40}}%
\TextAt{14}{196.855}{238.131}{o}%
\TextAt{21}{202.699}{236.338}{K}%
\TextAt{9}{210.902}{238.132}{\Character{41}}%
\TextAt{16}{218.111}{238.132}{\Character{0}}%
\TextAt{13}{230.066}{238.132}{\Character{31}}%
\TextAt{30}{239.37}{236.493}{\Character{101}}%
\TextAt{21}{237.4}{234.479}{K}%
\TextAt{9}{245.603}{238.132}{\Character{40}}%
\TextAt{14}{250.155}{238.131}{o}%
\TextAt{21}{255.999}{236.338}{K}%
\TextAt{9}{264.202}{238.132}{\Character{41}}%
\TextAt{9}{272.057}{238.132}{\Character{61}}%
\TextAt{13}{284.5}{238.132}{\Character{96}}%
\TextAt{9}{289.412}{238.132}{\Character{40}}%
\TextAt{14}{293.965}{238.132}{o}%
\TextAt{21}{299.809}{236.338}{K}%
\TextAt{9}{308.012}{238.132}{\Character{41}}%
\TextAt{16}{315.221}{238.132}{\Character{0}}%
\TextAt{13}{327.176}{238.132}{\Character{31}}%
\TextAt{9}{334.509}{238.132}{\Character{40}}%
\TextAt{14}{339.062}{238.132}{o}%
\TextAt{21}{344.906}{236.338}{K}%
\TextAt{9}{353.109}{238.132}{\Character{41}}%
\TextAt{9}{360.964}{238.132}{\Character{61}}%
\TextAt{8}{373.407}{238.132}{log}%
\TextAt{8}{391.563}{246.44}{\char35\relax }%
\TextAt{13}{401.319}{246.44}{\Character{22}}%
\TextAt{9}{408.362}{246.44}{\Character{40}}%
\TextAt{13}{412.914}{246.44}{K}%
\TextAt{9}{423.732}{246.44}{\Character{41}}%
\TextAt{20}{428.284}{256.662}{\Character{112}}%
\Vector{q 0 G 0 g 0.478 w 368.239 -517.099 m 389.166 -517.099 l S Q}%
\TextAt{16}{440.239}{246.44}{\Character{106}}%
\TextAt{13}{443.56}{246.44}{d}%
\TextAt{21}{449.643}{244.647}{K}%
\TextAt{16}{457.845}{246.44}{\Character{106}}%
\Vector{q 0 G 0 g 0.478 w 319.563 -532.88 m 389.166 -532.88 l S Q}%
\TextAt{9}{407.94}{229.931}{2}%
\TextAt{21}{413.793}{233.385}{r}%
\TextAt{9}{418.348}{229.931}{\Character{40}2}%
\TextAt{13}{428.753}{229.931}{\Character{25}}%
\TextAt{9}{435.823}{229.931}{\Character{41}}%
\TextAt{21}{440.375}{233.385}{s}%
\TextAt{8}{79.226}{205.574}{does}%
\TextAt{8}{105.623}{205.574}{not}%
\TextAt{8}{126.437}{205.574}{depend}%
\TextAt{8}{167.073}{205.574}{on}%
\TextAt{8}{183.332}{205.574}{the}%
\TextAt{8}{203.488}{205.574}{choice}%
\TextAt{8}{238.278}{205.574}{of}%
\TextAt{8}{251.608}{205.574}{metric.}%
\TextAt{8}{97.159}{185.539}{Just}%
\TextAt{8}{122.336}{185.539}{as}%
\TextAt{8}{136.3}{185.539}{in}%
\TextAt{8}{149.546}{185.539}{function}%
\TextAt{8}{194.987}{185.539}{theory,}%
\TextAt{8}{233.683}{185.539}{there}%
\TextAt{8}{263.188}{185.539}{is}%
\TextAt{8}{274.545}{185.539}{then}%
\TextAt{8}{300.799}{185.539}{no}%
\TextAt{8}{316.651}{185.539}{longer}%
\TextAt{8}{351.368}{185.539}{one}%
\TextAt{8}{372.421}{185.539}{smallest}%
\TextAt{8}{417.014}{185.539}{field,}%
\TextAt{8}{445.297}{185.539}{but}%
\TextAt{18}{466.35}{185.539}{Q}%
\TextAt{8}{479.139}{185.539}{is}%
\TextAt{8}{490.496}{185.539}{replaced}%
\TextAt{8}{79.226}{168.493}{by}%
\TextAt{8}{96.479}{168.493}{the}%
\TextAt{8}{117.628}{168.493}{continuous}%
\TextAt{8}{177.226}{168.493}{family}%
\TextAt{8}{213.989}{168.493}{of}%
\TextAt{8}{228.313}{168.493}{metrized}%
\TextAt{8}{277.425}{168.493}{fields}%
\TextAt{9}{308.39}{168.493}{\Character{40}}%
\TextAt{18}{312.943}{168.493}{Q}%
\TextAt{13}{322.241}{168.493}{\Character{59}}%
\TextAt{13}{327.465}{168.493}{\Character{11}xy}%
\TextAt{9}{347.796}{168.493}{\Character{41}}%
\TextAt{8}{352.348}{168.493}{,}%
\TextAt{13}{360.739}{168.493}{\Character{11}}%
\NativeMathAt{373.274}{168.493}{11.955}{\in}%
\TextAt{18}{386.258}{168.493}{R}%
\TextAt{3}{394.893}{172.832}{\Character{3}}%
\TextAt{17}{394.893}{165.538}{\Character{43}}%
\TextAt{8}{401.977}{168.493}{,}%
\TextAt{8}{410.368}{168.493}{all}%
\TextAt{8}{427.621}{168.493}{of}%
\TextAt{8}{441.944}{168.493}{which}%
\TextAt{8}{476.424}{168.493}{have}%
\TextAt{8}{504.412}{168.493}{genus}%
\TextAt{13}{79.226}{151.447}{g}%
\TextAt{9}{88.581}{151.447}{\Character{61}}%
\TextAt{9}{100.991}{151.447}{0}%
\TextAt{8}{106.86}{151.447}{.}%
\TextAt{8}{115.308}{151.447}{One}%
\TextAt{8}{140.019}{151.447}{even}%
\TextAt{8}{166.679}{151.447}{has}%
\TextAt{8}{187.553}{151.447}{the}%
\TextAt{12}{79.226}{125.435}{Proposition}%
\TextAt{12}{152.622}{125.435}{(3.14).}%
\TextAt{8}{195.264}{125.435}{The}%
\TextAt{8}{219.913}{125.435}{metrized}%
\TextAt{8}{268.627}{125.435}{fields}%
\TextAt{9}{299.194}{125.435}{\Character{40}}%
\TextAt{18}{303.747}{125.435}{Q}%
\TextAt{13}{313.045}{125.435}{\Character{59}}%
\TextAt{13}{318.269}{125.435}{\Character{11}xy}%
\TextAt{9}{338.6}{125.435}{\Character{41}}%
\TextAt{8}{347.645}{125.435}{are}%
\TextAt{8}{367.751}{125.435}{the}%
\TextAt{8}{388.502}{125.435}{only}%
\TextAt{8}{414.789}{125.435}{metrized}%
\TextAt{8}{463.504}{125.435}{number}%
\TextAt{8}{506.695}{125.435}{fields}%
\TextAt{8}{79.226}{108.389}{of}%
\TextAt{8}{92.556}{108.389}{genus}%
\TextAt{9}{124.811}{108.389}{0}%
\TextAt{8}{130.664}{108.389}{.}%
\TextAt{8}{297.211}{43.033}{204}%
\Destination{78.23}{809.815}{nav.page.204}%
\Destination{79.36}{506.755}{nav.print.223}%
\Destination{79.36}{139.455}{nav.number.3.3.14}%
}

% Page 217 (body page 205)
\Page{%
\TextAt{2}{79.226}{796.943}{III.}%
\TextAt{2}{98.317}{796.943}{Riemann–Roch}%
\TextAt{2}{174.375}{796.943}{Theory}%
\Vector{q 0 G 0 g 0.299 w 7.226 18.729 m 460.769 18.729 l S Q}%
\TextAt{19}{79.226}{761.124}{Proof.}%
\TextAt{8}{116}{761.124}{We}%
\TextAt{8}{136.145}{761.124}{have}%
\TextAt{13}{146.984}{729.154}{g}%
\TextAt{30}{154.56}{727.516}{\Character{101}}%
\TextAt{21}{152.589}{725.501}{K}%
\TextAt{9}{164.113}{729.154}{\Character{61}}%
\TextAt{8}{176.538}{729.154}{log}%
\TextAt{8}{194.694}{737.463}{\char35\relax }%
\TextAt{13}{204.45}{737.463}{\Character{22}}%
\TextAt{9}{211.493}{737.463}{\Character{40}}%
\TextAt{13}{216.045}{737.463}{K}%
\TextAt{9}{226.863}{737.463}{\Character{41}}%
\TextAt{20}{231.415}{747.685}{\Character{112}}%
\Vector{q 0 G 0 g 0.478 w 171.37 -26.076 m 192.297 -26.076 l S Q}%
\TextAt{16}{243.37}{737.463}{\Character{106}}%
\TextAt{13}{246.691}{737.463}{d}%
\TextAt{21}{252.774}{735.67}{K}%
\TextAt{16}{260.976}{737.463}{\Character{106}}%
\Vector{q 0 G 0 g 0.478 w 122.694 -41.857 m 192.297 -41.857 l S Q}%
\TextAt{9}{211.071}{720.953}{2}%
\TextAt{21}{216.924}{724.407}{r}%
\TextAt{9}{221.479}{720.953}{\Character{40}2}%
\TextAt{13}{231.884}{720.953}{\Character{25}}%
\TextAt{9}{238.954}{720.953}{\Character{41}}%
\TextAt{21}{243.506}{724.407}{s}%
\TextAt{9}{268.814}{729.154}{\Character{61}}%
\TextAt{9}{281.224}{729.154}{0}%
\NativeArrowAt{302.368}{729.154}{11.955}{21.926}{\Longleftrightarrow}%
\TextAt{8}{339.562}{729.154}{\char35\relax }%
\TextAt{13}{349.317}{729.154}{\Character{22}}%
\TextAt{9}{356.36}{729.154}{\Character{40}}%
\TextAt{13}{360.913}{729.154}{K}%
\TextAt{9}{371.73}{729.154}{\Character{41}}%
\TextAt{20}{376.283}{739.959}{\Character{112}}%
\Vector{q 0 G 0 g 0.478 w 316.238 -33.802 m 337.165 -33.802 l S Q}%
\TextAt{16}{388.238}{729.154}{\Character{106}}%
\TextAt{13}{391.559}{729.154}{d}%
\TextAt{21}{397.641}{727.361}{K}%
\TextAt{16}{405.844}{729.154}{\Character{106}}%
\TextAt{9}{412.486}{729.154}{\Character{61}}%
\TextAt{9}{424.896}{729.154}{2}%
\TextAt{21}{430.764}{734.09}{r}%
\TextAt{9}{435.319}{729.154}{\Character{40}2}%
\TextAt{13}{445.724}{729.154}{\Character{25}}%
\TextAt{9}{452.794}{729.154}{\Character{41}}%
\TextAt{21}{457.346}{734.09}{s}%
\TextAt{13}{461.76}{729.154}{\Character{58}}%
\TextAt{8}{79.226}{698.954}{Since}%
\TextAt{13}{109.822}{698.954}{\Character{25}}%
\TextAt{8}{120.828}{698.954}{is}%
\TextAt{8}{132.631}{698.954}{transcendental,}%
\TextAt{8}{214.369}{698.954}{one}%
\TextAt{8}{235.868}{698.954}{has}%
\TextAt{13}{256.78}{698.954}{s}%
\TextAt{9}{265.682}{698.954}{\Character{61}}%
\TextAt{9}{278.163}{698.954}{0}%
\TextAt{8}{284.026}{698.954}{,}%
\TextAt{8}{291.224}{698.954}{i.e.,}%
\TextAt{13}{313.378}{698.954}{K}%
\TextAt{8}{328.132}{698.954}{is}%
\TextAt{8}{339.935}{698.954}{totally}%
\TextAt{8}{377.382}{698.954}{real.}%
\TextAt{8}{404.812}{698.954}{Thus}%
\TextAt{8}{434.5}{698.954}{\char35\relax }%
\TextAt{13}{444.255}{698.954}{\Character{22}}%
\TextAt{9}{451.298}{698.954}{\Character{40}}%
\TextAt{13}{455.851}{698.954}{K}%
\TextAt{9}{466.668}{698.954}{\Character{41}}%
\TextAt{9}{474.594}{698.954}{\Character{61}}%
\TextAt{9}{487.076}{698.954}{2}%
\TextAt{8}{496.889}{698.954}{so}%
\TextAt{8}{511.298}{698.954}{that}%
\TextAt{16}{79.226}{681.908}{\Character{106}}%
\TextAt{13}{82.547}{681.908}{d}%
\TextAt{21}{88.63}{680.115}{K}%
\TextAt{16}{96.832}{681.908}{\Character{106}}%
\TextAt{9}{103.474}{681.908}{\Character{61}}%
\TextAt{9}{115.884}{681.908}{4}%
\TextAt{21}{121.753}{686.247}{n}%
\TextAt{3}{126.891}{686.247}{\Character{0}}%
\TextAt{17}{133.477}{686.247}{1}%
\TextAt{8}{138.21}{681.908}{,}%
\TextAt{8}{145.359}{681.908}{where}%
\TextAt{13}{179.168}{681.908}{n}%
\TextAt{9}{189.476}{681.908}{\Character{61}}%
\TextAt{13}{201.902}{681.908}{r}%
\TextAt{9}{210.823}{681.908}{\Character{61}}%
\TextAt{9}{223.233}{681.908}{\Character{91}}%
\TextAt{13}{226.5}{681.908}{K}%
\TextAt{9}{240.639}{681.908}{\Character{58}}%
\TextAt{18}{247.211}{681.908}{Q}%
\TextAt{9}{256.51}{681.908}{\Character{93}}%
\TextAt{8}{259.762}{681.908}{.}%
\TextAt{8}{268.21}{681.908}{In}%
\TextAt{8}{282.831}{681.908}{view}%
\TextAt{8}{309.814}{681.908}{of}%
\TextAt{8}{323.144}{681.908}{the}%
\TextAt{8}{343.3}{681.908}{bound}%
\TextAt{8}{379.393}{681.908}{(2.14)}%
\TextAt{8}{413.226}{681.908}{on}%
\TextAt{8}{429.485}{681.908}{the}%
\TextAt{8}{449.642}{681.908}{discriminant}%
\TextAt{16}{250.916}{652.292}{\Character{106}}%
\TextAt{13}{254.236}{652.292}{d}%
\TextAt{21}{260.319}{650.499}{K}%
\TextAt{16}{268.522}{652.292}{\Character{106}}%
\TextAt{17}{271.843}{657.229}{1\Character{47}2}%
\NativeMathAt{288.364}{652.292}{11.955}{\geq}%
\TextAt{13}{302.179}{660.38}{n}%
\TextAt{21}{309.167}{664.719}{n}%
\Vector{q 0 G 0 g 0.478 w 230.179 -118.719 m 242.803 -118.719 l S Q}%
\TextAt{13}{303.371}{644.092}{n}%
\TextAt{9}{310.359}{644.092}{\Character{33}}%
\TextAt{20}{317.991}{665.563}{\Character{16}}%
\TextAt{13}{326.326}{660.38}{\Character{25}}%
\Vector{q 0 G 0 g 0.478 w 254.326 -118.719 m 261.396 -118.719 l S Q}%
\TextAt{9}{326.935}{644.092}{4}%
\TextAt{20}{334.591}{665.563}{\Character{17}}%
\TextAt{21}{341.731}{662.886}{n}%
\TextAt{17}{346.869}{662.886}{\Character{47}2}%
\TextAt{13}{357.828}{652.292}{\Character{59}}%
\TextAt{8}{79.226}{626.199}{this}%
\TextAt{8}{101.37}{626.199}{can}%
\TextAt{8}{122.151}{626.199}{only}%
\TextAt{8}{147.164}{626.199}{happen}%
\TextAt{8}{187.778}{626.199}{if}%
\TextAt{13}{197.823}{626.199}{n}%
\NativeMathAt{208.132}{626.199}{11.955}{\leq}%
\TextAt{9}{220.751}{626.199}{6}%
\TextAt{8}{226.604}{626.199}{.}%
\TextAt{8}{234.826}{626.199}{But}%
\TextAt{8}{257.389}{626.199}{for}%
\TextAt{8}{274.595}{626.199}{this}%
\TextAt{8}{296.739}{626.199}{case}%
\TextAt{8}{320.831}{626.199}{one}%
\TextAt{8}{341.612}{626.199}{has}%
\TextAt{8}{361.807}{626.199}{sharper}%
\TextAt{8}{403.139}{626.199}{bounds,}%
\TextAt{8}{446.555}{626.199}{due}%
\TextAt{8}{467.982}{626.199}{to}%
\TextAt{26}{481.613}{626.199}{Odlyzko}%
\TextAt{8}{79.226}{609.153}{(see}%
\TextAt{8}{102.694}{609.153}{[111],}%
\TextAt{8}{133.921}{609.153}{table}%
\TextAt{8}{163.187}{609.153}{2):}%
\Vector{q 0 G 0 g 0.956 w 134.346 -178.289 m 333.65 -178.289 l S Q}%
\TextAt{13}{240.368}{579.952}{n}%
\TextAt{9}{259.407}{579.952}{\Character{61}}%
\TextAt{9}{295.512}{579.952}{3}%
\TextAt{9}{328.269}{579.952}{4}%
\TextAt{9}{361.026}{579.952}{5}%
\TextAt{9}{393.783}{579.952}{6}%
\Vector{q 0 G 0 g 0.598 w 134.346 -201.522 m 333.65 -201.522 l S Q}%
\TextAt{16}{212.323}{556.898}{\Character{106}}%
\TextAt{13}{215.644}{556.898}{d}%
\TextAt{21}{221.727}{555.104}{K}%
\TextAt{16}{229.93}{556.898}{\Character{106}}%
\TextAt{17}{233.251}{561.236}{1\Character{47}}%
\TextAt{21}{241.719}{561.236}{n}%
\NativeMathAt{259.31}{556.898}{11.955}{\geq}%
\TextAt{9}{280.564}{556.898}{3}%
\TextAt{13}{286.417}{556.898}{\Character{59}}%
\TextAt{9}{289.669}{556.898}{09}%
\TextAt{9}{313.304}{556.898}{4}%
\TextAt{13}{319.183}{556.898}{\Character{59}}%
\TextAt{9}{322.435}{556.898}{21}%
\TextAt{9}{346.07}{556.898}{5}%
\TextAt{13}{351.949}{556.898}{\Character{59}}%
\TextAt{9}{355.2}{556.898}{30}%
\TextAt{9}{378.847}{556.898}{6}%
\TextAt{13}{384.715}{556.898}{\Character{59}}%
\TextAt{9}{387.966}{556.898}{35}%
\Vector{q 0 G 0 g 0.956 w 134.346 -224.755 m 333.65 -224.755 l S Q}%
\TextAt{8}{79.226}{527.178}{This}%
\TextAt{8}{106.405}{527.178}{is}%
\TextAt{8}{118.628}{527.178}{not}%
\TextAt{8}{139.901}{527.178}{compatible}%
\TextAt{8}{200.196}{527.178}{with}%
\TextAt{16}{227.315}{527.178}{\Character{106}}%
\TextAt{13}{230.636}{527.178}{d}%
\TextAt{21}{236.719}{525.385}{K}%
\TextAt{16}{244.921}{527.178}{\Character{106}}%
\TextAt{17}{248.242}{531.517}{1\Character{47}}%
\TextAt{21}{256.711}{531.517}{n}%
\TextAt{9}{266.45}{527.178}{\Character{61}}%
\TextAt{9}{279.649}{527.178}{4}%
\TextAt{22}{286.707}{534.952}{n}%
\TextAt{23}{291.307}{534.952}{\Character{0}}%
\TextAt{24}{297.064}{534.952}{1}%
\Vector{q 0 G 0 g 0.359 w 214.707 -240.417 m 228.716 -240.417 l S Q}%
\TextAt{22}{291.411}{528.465}{n}%
\TextAt{8}{302.41}{527.178}{,}%
\TextAt{8}{310.133}{527.178}{so}%
\TextAt{8}{324.962}{527.178}{we}%
\TextAt{8}{342.649}{527.178}{may}%
\TextAt{8}{368.477}{527.178}{conclude}%
\TextAt{8}{417.056}{527.178}{that}%
\TextAt{13}{442.884}{527.178}{n}%
\NativeMathAt{453.975}{527.178}{11.955}{\leq}%
\TextAt{9}{467.376}{527.178}{2}%
\TextAt{8}{473.229}{527.178}{.}%
\TextAt{8}{483.055}{527.178}{But}%
\TextAt{8}{506.755}{527.178}{there}%
\TextAt{8}{79.226}{510.132}{is}%
\TextAt{8}{91.259}{510.132}{no}%
\TextAt{8}{107.786}{510.132}{real}%
\TextAt{8}{130.817}{510.132}{quadratic}%
\TextAt{8}{183.45}{510.132}{field}%
\TextAt{8}{209.075}{510.132}{with}%
\TextAt{8}{236.004}{510.132}{discriminant}%
\TextAt{16}{303.651}{510.132}{\Character{106}}%
\TextAt{13}{306.972}{510.132}{d}%
\TextAt{21}{313.055}{508.339}{K}%
\TextAt{16}{321.258}{510.132}{\Character{106}}%
\TextAt{9}{328.357}{510.132}{\Character{61}}%
\TextAt{9}{341.233}{510.132}{4}%
\TextAt{8}{351.259}{510.132}{(see}%
\TextAt{8}{374.996}{510.132}{chap.}%
\TextAt{8}{407.993}{510.132}{I,}%
\TextAt{8}{419.631}{510.132}{§2,}%
\TextAt{8}{438.652}{510.132}{exercise}%
\TextAt{8}{482.222}{510.132}{4).}%
\TextAt{8}{501.889}{510.132}{Hence}%
\TextAt{13}{79.226}{493.086}{n}%
\TextAt{9}{89.535}{493.086}{\Character{61}}%
\TextAt{9}{101.945}{493.086}{1}%
\TextAt{8}{107.813}{493.086}{,}%
\TextAt{8}{114.962}{493.086}{so}%
\TextAt{8}{129.332}{493.086}{that}%
\TextAt{13}{154.701}{493.086}{K}%
\TextAt{9}{168.84}{493.086}{\Character{61}}%
\TextAt{18}{181.265}{493.086}{Q}%
\TextAt{8}{190.564}{493.086}{.}%
\Vector{q 0 G 0 g 0.399 w 452.334 -280.914 m 452.334 -272.844 l S Q}%
\Vector{q 0 G 0 g 0.399 w 452.534 -273.043 m 459.707 -273.043 l S Q}%
\Vector{q 0 G 0 g 0.399 w 452.534 -280.714 m 459.707 -280.714 l S Q}%
\Vector{q 0 G 0 g 0.399 w 459.906 -280.914 m 459.906 -272.844 l S Q}%
\TextAt{8}{97.159}{464.916}{An}%
\TextAt{19}{115.922}{464.916}{extension}%
\TextAt{19}{167.184}{464.916}{of}%
\TextAt{19}{180.489}{464.916}{metrized}%
\TextAt{19}{227.125}{464.916}{number}%
\TextAt{19}{268.381}{464.916}{fields}%
\TextAt{8}{297.701}{464.916}{is}%
\TextAt{8}{309.051}{464.916}{a}%
\TextAt{8}{318.394}{464.916}{pair}%
\TextAt{20}{344.784}{467.938}{\Character{101}}%
\TextAt{13}{342.046}{464.916}{K}%
\TextAt{9}{356.185}{464.916}{\Character{61}}%
\TextAt{9}{368.595}{464.916}{\Character{40}}%
\TextAt{13}{373.162}{464.916}{K}%
\TextAt{13}{383.324}{464.916}{\Character{59}}%
\TextAt{16}{388.574}{464.916}{\Character{104}}%
\TextAt{13}{396.707}{464.916}{\Character{59}}%
\TextAt{16}{405.436}{464.916}{\Character{105}}%
\TextAt{21}{410.085}{463.123}{K}%
\TextAt{9}{418.288}{464.916}{\Character{41}}%
\TextAt{8}{422.84}{464.916}{,}%
\TextAt{20}{430.645}{467.938}{\Character{101}}%
\TextAt{13}{429.659}{464.916}{L}%
\TextAt{9}{440.944}{464.916}{\Character{61}}%
\TextAt{9}{453.353}{464.916}{\Character{40}}%
\TextAt{13}{457.922}{464.916}{L\Character{59}}%
\TextAt{16}{471.131}{464.916}{\Character{104}}%
\TextAt{13}{479.264}{464.916}{\Character{59}}%
\TextAt{16}{487.992}{464.916}{\Character{105}}%
\TextAt{21}{492.642}{463.123}{L}%
\TextAt{9}{498.899}{464.916}{\Character{41}}%
\TextAt{8}{503.451}{464.916}{,}%
\TextAt{8}{510.27}{464.916}{such}%
\TextAt{8}{79.226}{447.87}{that}%
\TextAt{13}{104.595}{447.87}{K}%
\NativeMathAt{118.734}{447.87}{11.955}{\subseteq}%
\TextAt{13}{131.353}{447.87}{L}%
\TextAt{8}{143.215}{447.87}{and}%
\TextAt{8}{165.977}{447.87}{the}%
\TextAt{8}{186.134}{447.87}{metrics}%
\TextAt{16}{185.836}{421.776}{\Character{104}}%
\TextAt{13}{190.485}{421.776}{x\Character{59}}%
\TextAt{13}{202.357}{421.776}{y}%
\TextAt{16}{208.518}{421.776}{\Character{105}}%
\TextAt{21}{213.168}{419.983}{K}%
\TextAt{9}{224.691}{421.776}{\Character{61}}%
\TextAt{20}{237.117}{433.134}{\Character{88}}%
\TextAt{21}{243.412}{408.028}{\Character{28}}%
\TextAt{13}{256.378}{421.776}{\Character{11}}%
\TextAt{21}{263.823}{419.983}{\Character{28}}%
\TextAt{13}{269}{421.776}{x}%
\TextAt{21}{275.652}{419.983}{\Character{28}}%
\Vector{q 0 G 0 g 0.478 w 208.828 -345.403 m 214.965 -345.403 l S Q}%
\TextAt{13}{280.828}{421.776}{y}%
\TextAt{21}{286.965}{418.954}{\Character{28}}%
\TextAt{13}{292.141}{421.776}{\Character{59}}%
\TextAt{16}{321.296}{421.776}{\Character{104}}%
\TextAt{13}{325.945}{421.776}{x\Character{59}}%
\TextAt{13}{337.817}{421.776}{y}%
\TextAt{16}{343.978}{421.776}{\Character{105}}%
\TextAt{21}{348.627}{419.983}{L}%
\TextAt{9}{358.205}{421.776}{\Character{61}}%
\TextAt{20}{370.631}{433.134}{\Character{88}}%
\TextAt{21}{376.703}{408.028}{\Character{27}}%
\TextAt{13}{389.892}{421.776}{\Character{12}}%
\TextAt{21}{396.5}{419.983}{\Character{27}}%
\TextAt{13}{402.124}{421.776}{x}%
\TextAt{21}{408.776}{419.983}{\Character{27}}%
\Vector{q 0 G 0 g 0.478 w 342.399 -345.403 m 348.536 -345.403 l S Q}%
\TextAt{13}{414.399}{421.776}{y}%
\TextAt{21}{420.536}{418.954}{\Character{27}}%
\TextAt{8}{79.226}{387.822}{satisfy}%
\TextAt{8}{116.158}{387.822}{the}%
\TextAt{8}{136.7}{387.822}{relation}%
\TextAt{13}{180.016}{387.822}{\Character{11}}%
\TextAt{21}{187.462}{386.029}{\Character{28}}%
\NativeMathAt{196.615}{387.822}{11.955}{\geq}%
\TextAt{13}{209.891}{387.822}{\Character{12}}%
\TextAt{21}{216.499}{386.029}{\Character{27}}%
\TextAt{8}{226.405}{387.822}{whenever}%
\TextAt{13}{278.162}{387.822}{\Character{28}}%
\TextAt{9}{288.558}{387.822}{\Character{61}}%
\TextAt{13}{301.64}{387.822}{\Character{27}}%
\TextAt{16}{308.722}{387.822}{\Character{106}}%
\TextAt{21}{312.043}{386.029}{K}%
\TextAt{8}{320.246}{387.822}{.}%
\TextAt{8}{329.85}{387.822}{If}%
\TextAt{14}{341.927}{387.822}{P}%
\TextAt{16}{355.797}{387.822}{\Character{106}}%
\TextAt{14}{363.095}{387.822}{p}%
\TextAt{8}{373.361}{387.822}{are}%
\TextAt{8}{393.257}{387.822}{infinite}%
\TextAt{8}{433.31}{387.822}{primes}%
\TextAt{8}{471.473}{387.822}{of}%
\TextAt{13}{485.188}{387.822}{L}%
\TextAt{16}{497.13}{387.822}{\Character{106}}%
\TextAt{13}{504.428}{387.822}{K}%
\TextAt{8}{515.246}{387.822}{,}%
\TextAt{14}{522.876}{387.822}{P}%
\TextAt{8}{79.226}{370.776}{belonging}%
\TextAt{8}{132.235}{370.776}{to}%
\TextAt{13}{146.546}{370.776}{\Character{27}}%
\TextAt{8}{157.525}{370.776}{and}%
\TextAt{14}{180.288}{370.776}{p}%
\TextAt{8}{190.169}{370.776}{to}%
\TextAt{13}{204.479}{370.776}{\Character{28}}%
\TextAt{9}{214.219}{370.776}{\Character{61}}%
\TextAt{13}{226.644}{370.776}{\Character{27}}%
\TextAt{16}{233.727}{370.776}{\Character{106}}%
\TextAt{21}{237.048}{368.983}{K}%
\TextAt{8}{245.25}{370.776}{,}%
\TextAt{8}{252.399}{370.776}{we}%
\TextAt{8}{269.627}{370.776}{define}%
\TextAt{8}{303.436}{370.776}{the}%
\TextAt{19}{323.592}{370.776}{ramification}%
\TextAt{19}{388.605}{370.776}{index}%
\TextAt{19}{419.688}{370.776}{and}%
\TextAt{19}{442.403}{370.776}{inertia}%
\TextAt{19}{480.516}{370.776}{degree}%
\TextAt{8}{515.138}{370.776}{by}%
\TextAt{13}{182.593}{344.683}{e}%
\TextAt{27}{188.018}{342.823}{P}%
\TextAt{3}{194.614}{342.823}{\Character{106}}%
\TextAt{27}{196.966}{342.823}{p}%
\TextAt{9}{204.774}{344.683}{\Character{61}}%
\TextAt{13}{217.199}{344.683}{\Character{11}}%
\TextAt{21}{224.645}{342.889}{\Character{28}}%
\TextAt{9}{229.821}{344.683}{\Character{47}}%
\TextAt{13}{235.674}{344.683}{\Character{12}}%
\TextAt{21}{242.283}{342.889}{\Character{27}}%
\TextAt{8}{271.817}{344.683}{and}%
\TextAt{13}{314.592}{344.683}{f}%
\TextAt{27}{320.362}{342.823}{P}%
\TextAt{3}{326.958}{342.823}{\Character{106}}%
\TextAt{27}{329.31}{342.823}{p}%
\TextAt{9}{337.118}{344.683}{\Character{61}}%
\TextAt{13}{349.543}{344.683}{\Character{12}}%
\TextAt{21}{356.152}{342.889}{\Character{27}}%
\TextAt{9}{361.775}{344.683}{\Character{47}}%
\TextAt{13}{367.628}{344.683}{\Character{11}}%
\TextAt{21}{375.074}{342.889}{\Character{28}}%
\TextAt{9}{380.25}{344.683}{\Character{91}}%
\TextAt{13}{383.502}{344.683}{L}%
\TextAt{27}{391.466}{342.889}{P}%
\TextAt{9}{401.881}{344.683}{\Character{58}}%
\TextAt{13}{408.453}{344.683}{K}%
\TextAt{27}{418.412}{342.889}{p}%
\TextAt{9}{422.899}{344.683}{\Character{93}}%
\TextAt{13}{426.151}{344.683}{\Character{58}}%
\TextAt{8}{79.226}{318.589}{Thus}%
\TextAt{8}{108.875}{318.589}{the}%
\TextAt{8}{129.031}{318.589}{fundamental}%
\TextAt{8}{196.674}{318.589}{identity}%
\TextAt{20}{251.872}{303.853}{\Character{88}}%
\TextAt{27}{254.038}{277.95}{P}%
\TextAt{3}{260.634}{277.95}{\Character{106}}%
\TextAt{27}{262.986}{277.95}{p}%
\TextAt{13}{271.133}{292.495}{e}%
\TextAt{27}{276.559}{290.636}{P}%
\TextAt{3}{283.154}{290.636}{\Character{106}}%
\TextAt{27}{285.507}{290.636}{p}%
\TextAt{13}{289.994}{292.495}{f}%
\TextAt{27}{295.764}{290.636}{P}%
\TextAt{3}{302.359}{290.636}{\Character{106}}%
\TextAt{27}{304.711}{290.636}{p}%
\TextAt{9}{312.519}{292.495}{\Character{61}}%
\TextAt{9}{324.928}{292.495}{\Character{91}}%
\TextAt{13}{328.196}{292.495}{L}%
\TextAt{9}{339.481}{292.495}{\Character{58}}%
\TextAt{13}{346.054}{292.495}{K}%
\TextAt{9}{356.872}{292.495}{\Character{93}}%
\TextAt{8}{79.226}{256.182}{is}%
\TextAt{8}{92.087}{256.182}{preserved.}%
\TextAt{8}{152.017}{256.182}{Also}%
\TextAt{14}{179.511}{256.182}{P}%
\TextAt{8}{194.398}{256.182}{is}%
\TextAt{8}{207.259}{256.182}{unramified}%
\TextAt{8}{266.888}{256.182}{if}%
\TextAt{8}{278.709}{256.182}{and}%
\TextAt{8}{302.568}{256.182}{only}%
\TextAt{8}{329.357}{256.182}{if}%
\TextAt{13}{341.177}{256.182}{\Character{11}}%
\TextAt{21}{348.623}{254.389}{\Character{28}}%
\TextAt{9}{358.99}{256.182}{\Character{61}}%
\TextAt{13}{373.284}{256.182}{\Character{12}}%
\TextAt{21}{379.893}{254.389}{\Character{27}}%
\TextAt{8}{385.516}{256.182}{.}%
\TextAt{8}{397.255}{256.182}{For}%
\TextAt{8}{419.321}{256.182}{“replete}%
\TextAt{8}{464.305}{256.182}{prime}%
\TextAt{8}{498.566}{256.182}{ideals”}%
\TextAt{14}{79.226}{239.136}{p}%
\TextAt{9}{88.53}{239.136}{\Character{61}}%
\TextAt{13}{100.956}{239.136}{e}%
\TextAt{21}{106.381}{243.474}{\Character{11}}%
\TextAt{22}{111.82}{242.478}{\Character{28}}%
\TextAt{8}{116.987}{239.136}{,}%
\TextAt{14}{124.137}{239.136}{P}%
\TextAt{9}{137.351}{239.136}{\Character{61}}%
\TextAt{13}{149.776}{239.136}{e}%
\TextAt{21}{155.202}{243.474}{\Character{12}}%
\TextAt{22}{159.953}{242.478}{\Character{27}}%
\TextAt{8}{165.447}{239.136}{,}%
\TextAt{8}{172.596}{239.136}{we}%
\TextAt{8}{189.824}{239.136}{put}%
\TextAt{13}{200.888}{213.042}{i}%
\TextAt{21}{204.882}{211.183}{L}%
\TextAt{3}{210.641}{211.183}{\Character{106}}%
\TextAt{21}{212.993}{211.183}{K}%
\TextAt{9}{221.196}{213.042}{\Character{40}}%
\TextAt{14}{225.748}{213.042}{p}%
\TextAt{9}{231.731}{213.042}{\Character{41}}%
\TextAt{9}{239.586}{213.042}{\Character{61}}%
\TextAt{20}{252.03}{224.4}{\Character{89}}%
\TextAt{27}{253.199}{198.497}{P}%
\TextAt{3}{259.795}{198.497}{\Character{106}}%
\TextAt{27}{262.147}{198.497}{p}%
\TextAt{14}{269.298}{213.042}{P}%
\TextAt{21}{279.191}{217.978}{e}%
\TextAt{28}{283.12}{216.24}{P}%
\TextAt{23}{288.819}{216.241}{\Character{106}}%
\TextAt{28}{291.088}{216.24}{p}%
\TextAt{13}{295.669}{213.042}{\Character{59}}%
\TextAt{13}{324.804}{213.042}{N}%
\TextAt{21}{334.197}{211.183}{L}%
\TextAt{3}{339.956}{211.183}{\Character{106}}%
\TextAt{21}{342.309}{211.183}{K}%
\TextAt{9}{350.511}{213.042}{\Character{40}}%
\TextAt{14}{355.064}{213.042}{P}%
\TextAt{9}{364.957}{213.042}{\Character{41}}%
\TextAt{9}{372.812}{213.042}{\Character{61}}%
\TextAt{14}{385.256}{213.042}{p}%
\TextAt{21}{391.239}{217.978}{f}%
\TextAt{28}{395.307}{216.24}{P}%
\TextAt{23}{401.006}{216.241}{\Character{106}}%
\TextAt{28}{403.275}{216.24}{p}%
\TextAt{13}{407.856}{213.042}{\Character{58}}%
\TextAt{8}{79.226}{173.608}{Finally}%
\TextAt{8}{119.061}{173.608}{we}%
\TextAt{8}{136.288}{173.608}{define}%
\TextAt{8}{170.097}{173.608}{the}%
\TextAt{19}{190.254}{173.608}{different}%
\TextAt{8}{236.939}{173.608}{of}%
\TextAt{20}{251.255}{176.63}{\Character{101}}%
\TextAt{13}{250.269}{173.608}{L}%
\TextAt{16}{261.554}{173.608}{\Character{106}}%
\TextAt{20}{270.934}{176.63}{\Character{101}}%
\TextAt{13}{268.196}{173.608}{K}%
\TextAt{8}{282.911}{173.608}{to}%
\TextAt{8}{297.221}{173.608}{be}%
\TextAt{8}{313.146}{173.608}{the}%
\TextAt{8}{333.302}{173.608}{replete}%
\TextAt{8}{371.666}{173.608}{ideal}%
\TextAt{14}{187.64}{147.515}{D}%
\TextAt{30}{198.344}{145.877}{\Character{101}}%
\TextAt{21}{197.582}{143.862}{L}%
\TextAt{3}{203.341}{143.862}{\Character{106}}%
\TextAt{30}{207.663}{145.877}{\Character{101}}%
\TextAt{21}{205.693}{143.862}{K}%
\TextAt{9}{217.217}{147.515}{\Character{61}}%
\TextAt{14}{229.642}{147.515}{D}%
\TextAt{21}{239.584}{145.655}{L}%
\TextAt{3}{245.343}{145.655}{\Character{106}}%
\TextAt{21}{247.695}{145.655}{K}%
\TextAt{16}{258.554}{147.515}{\Character{1}}%
\TextAt{14}{264.532}{147.515}{D}%
\NativeMathAt{274.474}{145.722}{7.97}{\infty}%
\NativeMathAt{286.761}{147.515}{11.955}{\in}%
\TextAt{13}{298.052}{147.515}{J}%
\TextAt{9}{305.629}{147.515}{\Character{40}}%
\Vector{q 0 G 0 g 0.478 w 238.181 -619.129 m 244.026 -619.129 l S Q}%
\TextAt{14}{310.181}{147.515}{o}%
\TextAt{21}{316.026}{145.722}{L}%
\TextAt{9}{322.283}{147.515}{\Character{41}}%
\TextAt{9}{330.138}{147.515}{\Character{61}}%
\TextAt{13}{342.581}{147.515}{J}%
\TextAt{9}{350.159}{147.515}{\Character{40}}%
\TextAt{14}{354.711}{147.515}{o}%
\TextAt{21}{360.555}{145.722}{L}%
\TextAt{9}{366.812}{147.515}{\Character{41}}%
\TextAt{16}{374.021}{147.515}{\Character{2}}%
\TextAt{20}{387.046}{158.872}{\Character{89}}%
\TextAt{27}{385.976}{132.969}{P}%
\TextAt{3}{392.572}{132.969}{\Character{106}}%
\NativeMathAt{394.923}{132.969}{7.97}{\infty}%
\TextAt{18}{405.385}{147.515}{R}%
\TextAt{3}{414.019}{152.451}{\Character{3}}%
\TextAt{17}{414.019}{144.559}{\Character{43}}%
\TextAt{13}{421.104}{147.515}{\Character{59}}%
\TextAt{8}{79.226}{110.771}{where}%
\TextAt{14}{113.035}{110.771}{D}%
\TextAt{21}{122.977}{108.911}{L}%
\TextAt{3}{128.736}{108.911}{\Character{106}}%
\TextAt{21}{131.088}{108.911}{K}%
\TextAt{8}{143.188}{110.771}{is}%
\TextAt{8}{154.952}{110.771}{the}%
\TextAt{8}{175.109}{110.771}{different}%
\TextAt{8}{221.28}{110.771}{of}%
\TextAt{13}{234.61}{110.771}{L}%
\TextAt{16}{245.895}{110.771}{\Character{106}}%
\TextAt{13}{252.537}{110.771}{K}%
\TextAt{8}{267.252}{110.771}{and}%
\TextAt{14}{232.922}{84.677}{D}%
\NativeMathAt{242.863}{82.884}{7.97}{\infty}%
\TextAt{9}{255.151}{84.677}{\Character{61}}%
\TextAt{9}{267.56}{84.677}{\Character{40}}%
\TextAt{13}{272.129}{84.677}{\Character{12}}%
\TextAt{27}{278.737}{82.884}{P}%
\TextAt{9}{285.831}{84.677}{\Character{47}}%
\TextAt{13}{291.684}{84.677}{\Character{11}}%
\TextAt{27}{299.129}{82.884}{p}%
\TextAt{9}{303.616}{84.677}{\Character{41}}%
\TextAt{27}{308.168}{82.818}{P}%
\TextAt{3}{314.764}{82.818}{\Character{106}}%
\NativeMathAt{317.115}{82.818}{7.97}{\infty}%
\NativeMathAt{329.404}{84.677}{11.955}{\in}%
\TextAt{20}{341.765}{96.035}{\Character{89}}%
\TextAt{27}{340.695}{70.132}{P}%
\TextAt{3}{347.29}{70.132}{\Character{106}}%
\NativeMathAt{349.641}{70.132}{7.97}{\infty}%
\TextAt{18}{360.103}{84.677}{R}%
\TextAt{3}{368.738}{89.614}{\Character{3}}%
\TextAt{17}{368.738}{81.722}{\Character{43}}%
\TextAt{13}{375.822}{84.677}{\Character{59}}%
\TextAt{8}{297.211}{43.033}{205}%
\Destination{78.23}{809.815}{nav.page.205}%
\Destination{79.36}{479.955}{nav.print.224}%
\Hotlink{102.194}{606.955}{25.078}{11.047}{nav.cite.N111}%
\Hotlink{407.493}{507.955}{5.22}{11.047}{nav.chapter.1}%
\Hotlink{424.81}{507.955}{6.858}{11.047}{nav.section.1.2}%
\Hotlink{378.893}{679.755}{30.936}{11.047}{nav.number.3.2.14}%
}

% Page 218 (body page 206)
\Page{%
\TextAt{2}{79.226}{796.943}{III.}%
\TextAt{2}{98.317}{796.943}{Riemann–Roch}%
\TextAt{2}{174.375}{796.943}{Theory}%
\Vector{q 0 G 0 g 0.299 w 7.226 18.729 m 460.769 18.729 l S Q}%
\TextAt{8}{79.226}{761.124}{where}%
\TextAt{13}{113.412}{761.124}{\Character{12}}%
\TextAt{27}{120.02}{759.33}{P}%
\TextAt{9}{131.076}{761.124}{\Character{61}}%
\TextAt{13}{144.142}{761.124}{\Character{12}}%
\TextAt{21}{150.751}{759.33}{\Character{27}}%
\TextAt{8}{160.648}{761.124}{and}%
\TextAt{13}{183.786}{761.124}{\Character{11}}%
\TextAt{27}{191.232}{759.33}{p}%
\TextAt{9}{199.681}{761.124}{\Character{61}}%
\TextAt{13}{212.747}{761.124}{\Character{11}}%
\TextAt{21}{220.193}{759.33}{\Character{28}}%
\TextAt{8}{229.643}{761.124}{(}%
\TextAt{14}{234.198}{761.124}{P}%
\TextAt{8}{248.365}{761.124}{belongs}%
\TextAt{8}{290.751}{761.124}{to}%
\TextAt{13}{305.438}{761.124}{\Character{27}}%
\TextAt{8}{316.794}{761.124}{and}%
\TextAt{14}{339.933}{761.124}{p}%
\TextAt{8}{350.189}{761.124}{to}%
\TextAt{13}{364.876}{761.124}{\Character{28}}%
\TextAt{9}{375.257}{761.124}{\Character{61}}%
\TextAt{13}{388.323}{761.124}{\Character{27}}%
\TextAt{16}{395.406}{761.124}{\Character{106}}%
\TextAt{21}{398.727}{759.33}{K}%
\TextAt{8}{406.929}{761.124}{).}%
\TextAt{8}{421.061}{761.124}{With}%
\TextAt{8}{451.672}{761.124}{this}%
\TextAt{8}{474.871}{761.124}{convention,}%
\TextAt{8}{79.226}{744.078}{we}%
\TextAt{8}{96.454}{744.078}{obtain}%
\TextAt{8}{132.881}{744.078}{the}%
\TextAt{8}{153.037}{744.078}{general}%
\TextAt{19}{193.362}{744.078}{Riemann–Hurwitz}%
\TextAt{19}{289.338}{744.078}{formula}%
\TextAt{13}{177.142}{717.893}{g}%
\TextAt{30}{183.51}{716.255}{\Character{101}}%
\TextAt{21}{182.747}{714.24}{L}%
\TextAt{16}{191.661}{717.893}{\Character{0}}%
\TextAt{13}{203.616}{717.893}{\Character{96}}%
\TextAt{30}{209.291}{716.255}{\Character{101}}%
\TextAt{21}{208.529}{714.24}{L}%
\TextAt{9}{214.786}{717.893}{\Character{40}}%
\TextAt{14}{219.338}{717.893}{o}%
\TextAt{21}{225.182}{716.1}{L}%
\TextAt{9}{231.439}{717.893}{\Character{41}}%
\TextAt{9}{239.294}{717.893}{\Character{61}}%
\TextAt{9}{251.715}{717.893}{\Character{91}}%
\TextAt{13}{254.99}{717.893}{L}%
\TextAt{9}{266.275}{717.893}{\Character{58}}%
\TextAt{13}{272.848}{717.893}{K}%
\TextAt{9}{283.665}{717.893}{\Character{93}}%
\TextAt{20}{286.917}{727.577}{\Character{0}}%
\TextAt{13}{292.396}{717.893}{g}%
\TextAt{30}{299.972}{716.255}{\Character{101}}%
\TextAt{21}{298.002}{714.24}{K}%
\TextAt{16}{308.861}{717.893}{\Character{0}}%
\TextAt{13}{320.816}{717.893}{\Character{96}}%
\TextAt{30}{327.699}{716.255}{\Character{101}}%
\TextAt{21}{325.728}{714.24}{K}%
\TextAt{9}{333.931}{717.893}{\Character{40}}%
\TextAt{14}{338.483}{717.893}{o}%
\TextAt{21}{344.328}{716.1}{K}%
\TextAt{9}{352.53}{717.893}{\Character{41}}%
\TextAt{20}{357.083}{727.577}{\Character{1}}%
\TextAt{16}{365.219}{717.893}{\Character{0}}%
\TextAt{17}{378.369}{722.6}{1}%
\Vector{q 0 G 0 g 0.478 w 306.369 -53.118 m 310.604 -53.118 l S Q}%
\TextAt{17}{378.369}{713.77}{2}%
\TextAt{8}{385.792}{717.893}{deg}%
\TextAt{14}{405.346}{717.893}{D}%
\TextAt{30}{416.05}{716.255}{\Character{101}}%
\TextAt{21}{415.288}{714.24}{L}%
\TextAt{3}{421.047}{714.24}{\Character{106}}%
\TextAt{30}{425.37}{716.255}{\Character{101}}%
\TextAt{21}{423.399}{714.24}{K}%
\TextAt{13}{431.602}{717.893}{\Character{58}}%
\TextAt{8}{79.226}{691.709}{If}%
\TextAt{8}{91.981}{691.709}{we}%
\TextAt{8}{110.271}{691.709}{consider}%
\TextAt{8}{156.919}{691.709}{only}%
\TextAt{8}{183.674}{691.709}{number}%
\TextAt{8}{227.333}{691.709}{fields}%
\TextAt{8}{258.367}{691.709}{endowed}%
\TextAt{8}{306.904}{691.709}{with}%
\TextAt{8}{334.627}{691.709}{the}%
\TextAt{8}{355.846}{691.709}{Minkowski}%
\TextAt{8}{415.178}{691.709}{metric,}%
\TextAt{8}{456.174}{691.709}{then}%
\TextAt{13}{483.897}{691.709}{L}%
\TextAt{27}{491.861}{689.915}{P}%
\NativeMathAt{504.087}{691.709}{11.955}{\neq}%
\TextAt{13}{518.323}{691.709}{K}%
\TextAt{27}{528.283}{689.915}{p}%
\TextAt{8}{79.226}{674.663}{is}%
\TextAt{8}{92.408}{674.663}{always}%
\TextAt{8}{131.293}{674.663}{ramified.}%
\TextAt{8}{185.622}{674.663}{In}%
\TextAt{8}{201.661}{674.663}{this}%
\TextAt{8}{225.901}{674.663}{way}%
\TextAt{8}{251.062}{674.663}{the}%
\TextAt{8}{272.636}{674.663}{convention}%
\TextAt{8}{332.598}{674.663}{found}%
\TextAt{8}{366.857}{674.663}{in}%
\TextAt{8}{381.927}{674.663}{many}%
\TextAt{8}{415.217}{674.663}{textbooks}%
\TextAt{8}{470.684}{674.663}{is}%
\TextAt{8}{483.866}{674.663}{no}%
\TextAt{8}{501.543}{674.663}{longer}%
\TextAt{8}{79.226}{657.617}{incompatible}%
\TextAt{8}{148.817}{657.617}{with}%
\TextAt{8}{175.477}{657.617}{the}%
\TextAt{8}{195.634}{657.617}{customs}%
\TextAt{8}{240.633}{657.617}{introduced}%
\TextAt{8}{299.165}{657.617}{in}%
\TextAt{8}{312.818}{657.617}{the}%
\TextAt{8}{332.975}{657.617}{present}%
\TextAt{8}{373.682}{657.617}{book.}%
\TextAt{15}{79.226}{611.147}{§}%
\TextAt{15}{93.785}{611.147}{4}%
\TextAt{15}{120.693}{611.147}{Metrized}%
\TextAt{32}{203.396}{611.147}{o}%
\TextAt{15}{211.812}{611.147}{-Modules}%
\TextAt{8}{79.226}{579.261}{The}%
\TextAt{8}{105.152}{579.261}{Riemann–Roch}%
\TextAt{8}{188.654}{579.261}{theory}%
\TextAt{8}{227.276}{579.261}{which}%
\TextAt{8}{262.634}{579.261}{was}%
\TextAt{8}{287.006}{579.261}{presented}%
\TextAt{8}{341.289}{579.261}{in}%
\TextAt{8}{356.814}{579.261}{the}%
\TextAt{8}{378.843}{579.261}{preceding}%
\TextAt{8}{433.389}{579.261}{section}%
\TextAt{8}{474.343}{579.261}{in}%
\TextAt{8}{489.867}{579.261}{the}%
\TextAt{8}{511.896}{579.261}{case}%
\TextAt{8}{79.226}{562.215}{of}%
\TextAt{8}{93.713}{562.215}{replete}%
\TextAt{8}{133.234}{562.215}{ideals}%
\TextAt{8}{166.969}{562.215}{is}%
\TextAt{8}{179.89}{562.215}{embedded}%
\TextAt{8}{236.315}{562.215}{in}%
\TextAt{8}{251.125}{562.215}{a}%
\TextAt{8}{262.037}{562.215}{much}%
\TextAt{8}{294.409}{562.215}{more}%
\TextAt{8}{324.832}{562.215}{far-reaching}%
\TextAt{8}{390.38}{562.215}{theory}%
\TextAt{8}{428.287}{562.215}{which}%
\TextAt{8}{462.93}{562.215}{treats}%
\TextAt{8}{497.322}{562.215}{finitely}%
\TextAt{8}{79.226}{545.169}{generated}%
\TextAt{14}{133.238}{545.169}{o}%
\TextAt{8}{139.083}{545.169}{-modules.}%
\TextAt{8}{195.48}{545.169}{It}%
\TextAt{8}{208.833}{545.169}{is}%
\TextAt{8}{221.277}{545.169}{only}%
\TextAt{8}{247.649}{545.169}{in}%
\TextAt{8}{261.982}{545.169}{this}%
\TextAt{8}{285.485}{545.169}{setting}%
\TextAt{8}{324.601}{545.169}{that}%
\TextAt{8}{350.65}{545.169}{the}%
\TextAt{8}{371.487}{545.169}{theory}%
\TextAt{8}{408.917}{545.169}{displays}%
\TextAt{8}{453.951}{545.169}{its}%
\TextAt{8}{470.95}{545.169}{true}%
\TextAt{8}{496.342}{545.169}{nature,}%
\TextAt{8}{79.226}{528.123}{and}%
\TextAt{8}{101.274}{528.123}{becomes}%
\TextAt{8}{147.113}{528.123}{susceptible}%
\TextAt{8}{205.696}{528.123}{to}%
\TextAt{8}{219.292}{528.123}{the}%
\TextAt{8}{238.734}{528.123}{most}%
\TextAt{8}{266.699}{528.123}{comprehensive}%
\TextAt{8}{344.088}{528.123}{generalization.}%
\TextAt{8}{423.204}{528.123}{This}%
\TextAt{8}{449.209}{528.123}{theory}%
\TextAt{8}{485.245}{528.123}{is}%
\TextAt{8}{496.294}{528.123}{subject}%
\TextAt{8}{79.226}{511.077}{to}%
\TextAt{8}{93.91}{511.077}{a}%
\TextAt{8}{104.04}{511.077}{formalism}%
\TextAt{8}{158.786}{511.077}{which}%
\TextAt{8}{192.646}{511.077}{has}%
\TextAt{8}{213.894}{511.077}{been}%
\TextAt{8}{241.896}{511.077}{constructed}%
\TextAt{8}{305.417}{511.077}{by}%
\TextAt{26}{322.05}{511.077}{Alexander}%
\TextAt{26}{392.852}{511.077}{Grothendieck}%
\TextAt{8}{482.638}{511.077}{for}%
\TextAt{8}{500.897}{511.077}{higher}%
\TextAt{8}{79.226}{494.031}{dimensional}%
\TextAt{8}{142.905}{494.031}{algebraic}%
\TextAt{8}{191.569}{494.031}{varieties,}%
\TextAt{8}{240.123}{494.031}{and}%
\TextAt{8}{262.115}{494.031}{which}%
\TextAt{8}{294.83}{494.031}{we}%
\TextAt{8}{311.286}{494.031}{will}%
\TextAt{8}{332.62}{494.031}{now}%
\TextAt{8}{356.238}{494.031}{develop}%
\TextAt{8}{397.74}{494.031}{for}%
\TextAt{8}{414.854}{494.031}{number}%
\TextAt{8}{456.679}{494.031}{fields.}%
\TextAt{8}{490.944}{494.031}{In}%
\TextAt{8}{504.794}{494.031}{doing}%
\TextAt{8}{79.226}{476.985}{so,}%
\TextAt{8}{96.456}{476.985}{our}%
\TextAt{8}{116.781}{476.985}{principal}%
\TextAt{8}{165.069}{476.985}{attention}%
\TextAt{8}{214.994}{476.985}{will}%
\TextAt{8}{236.61}{476.985}{be}%
\TextAt{8}{252.045}{476.985}{focused}%
\TextAt{8}{293.231}{476.985}{as}%
\TextAt{8}{307.111}{476.985}{before}%
\TextAt{8}{341.734}{476.985}{on}%
\TextAt{8}{357.504}{476.985}{the}%
\TextAt{8}{377.171}{476.985}{kind}%
\TextAt{8}{403.019}{476.985}{of}%
\TextAt{8}{415.859}{476.985}{compactification}%
\TextAt{8}{503.18}{476.985}{which}%
\TextAt{8}{79.226}{459.939}{is}%
\TextAt{8}{91.844}{459.939}{accomplished}%
\TextAt{8}{164.298}{459.939}{by}%
\TextAt{8}{181.411}{459.939}{taking}%
\TextAt{8}{218.37}{459.939}{into}%
\TextAt{8}{242.968}{459.939}{account}%
\TextAt{8}{287.076}{459.939}{the}%
\TextAt{8}{308.086}{459.939}{infinite}%
\TextAt{8}{348.608}{459.939}{places.}%
\TextAt{8}{390.248}{459.939}{The}%
\TextAt{8}{415.156}{459.939}{effect}%
\TextAt{8}{446.891}{459.939}{is}%
\TextAt{8}{459.509}{459.939}{that}%
\TextAt{8}{485.732}{459.939}{a}%
\TextAt{8}{496.342}{459.939}{leading}%
\TextAt{8}{79.226}{442.893}{rôle}%
\TextAt{8}{102.394}{442.893}{is}%
\TextAt{8}{114.563}{442.893}{claimed}%
\TextAt{8}{157.888}{442.893}{by}%
\TextAt{8}{174.552}{442.893}{linear}%
\TextAt{8}{207.475}{442.893}{algebra—for}%
\TextAt{8}{274.555}{442.893}{which}%
\TextAt{8}{308.446}{442.893}{we}%
\TextAt{8}{326.079}{442.893}{refer}%
\TextAt{8}{353.467}{442.893}{to}%
\TextAt{8}{368.183}{442.893}{[15].}%
\TextAt{8}{396.067}{442.893}{Our}%
\TextAt{8}{420.538}{442.893}{treatment}%
\TextAt{8}{475.255}{442.893}{is}%
\TextAt{8}{487.425}{442.893}{based}%
\TextAt{8}{520.408}{442.893}{on}%
\TextAt{8}{79.226}{425.847}{a}%
\TextAt{8}{90.408}{425.847}{course}%
\TextAt{8}{127.665}{425.847}{on}%
\TextAt{8}{145.35}{425.847}{“Arakelov}%
\TextAt{8}{201.412}{425.847}{Theory}%
\TextAt{8}{243.486}{425.847}{and}%
\TextAt{8}{267.676}{425.847}{Grothendieck–Riemann–Roch”}%
\TextAt{8}{430.485}{425.847}{taught}%
\TextAt{8}{469.319}{425.847}{by}%
\TextAt{26}{487.005}{425.847}{Günter}%
\TextAt{26}{79.226}{408.801}{Tamme}%
\TextAt{8}{119.431}{408.801}{.}%
\TextAt{8}{128.777}{408.801}{There,}%
\TextAt{8}{166.212}{408.801}{however,}%
\TextAt{8}{214.717}{408.801}{proofs}%
\TextAt{8}{250.201}{408.801}{were}%
\TextAt{8}{277.483}{408.801}{not}%
\TextAt{8}{298.596}{408.801}{given}%
\TextAt{8}{329.464}{408.801}{directly,}%
\TextAt{8}{374.706}{408.801}{as}%
\TextAt{8}{389.375}{408.801}{we}%
\TextAt{8}{406.901}{408.801}{will}%
\TextAt{8}{429.306}{408.801}{do}%
\TextAt{8}{445.864}{408.801}{here,}%
\TextAt{8}{474.846}{408.801}{but}%
\TextAt{8}{496.605}{408.801}{usually}%
\TextAt{8}{79.226}{391.755}{deduced}%
\TextAt{8}{124.739}{391.755}{as}%
\TextAt{8}{139.11}{391.755}{special}%
\TextAt{8}{177.211}{391.755}{cases}%
\TextAt{8}{206.597}{391.755}{from}%
\TextAt{8}{234.237}{391.755}{the}%
\TextAt{8}{254.393}{391.755}{general}%
\TextAt{8}{294.718}{391.755}{abstract}%
\TextAt{8}{340.315}{391.755}{theory.}%
\TextAt{8}{97.159}{371.72}{Let}%
\TextAt{13}{118.379}{371.72}{K}%
\TextAt{8}{133.344}{371.72}{be}%
\TextAt{8}{149.519}{371.72}{an}%
\TextAt{8}{166.028}{371.72}{algebraic}%
\TextAt{8}{215.713}{371.72}{number}%
\TextAt{8}{258.559}{371.72}{field}%
\TextAt{8}{284.167}{371.72}{and}%
\TextAt{14}{307.179}{371.72}{o}%
\TextAt{8}{317.171}{371.72}{the}%
\TextAt{8}{337.578}{371.72}{ring}%
\TextAt{8}{361.894}{371.72}{of}%
\TextAt{8}{375.474}{371.72}{integers}%
\TextAt{8}{419.038}{371.72}{of}%
\TextAt{13}{432.619}{371.72}{K}%
\TextAt{8}{443.436}{371.72}{.}%
\TextAt{8}{452.635}{371.72}{For}%
\TextAt{8}{473.855}{371.72}{the}%
\TextAt{8}{494.262}{371.72}{passage}%
\TextAt{8}{79.226}{354.674}{from}%
\TextAt{13}{106.866}{354.674}{K}%
\TextAt{8}{121.582}{354.674}{to}%
\TextAt{18}{135.892}{354.674}{R}%
\TextAt{8}{148.424}{354.674}{and}%
\TextAt{18}{171.186}{354.674}{C}%
\TextAt{8}{179.821}{354.674}{,}%
\TextAt{8}{186.97}{354.674}{we}%
\TextAt{8}{204.197}{354.674}{start}%
\TextAt{8}{232.232}{354.674}{by}%
\TextAt{8}{248.491}{354.674}{considering}%
\TextAt{8}{309.69}{354.674}{the}%
\TextAt{8}{329.846}{354.674}{ring}%
\TextAt{13}{268.012}{328.49}{K}%
\TextAt{25}{277.971}{326.697}{C}%
\TextAt{9}{287.546}{328.49}{\Character{61}}%
\TextAt{13}{299.971}{328.49}{K}%
\TextAt{16}{313.446}{328.49}{\Character{10}}%
\TextAt{25}{322.744}{326.697}{Q}%
\TextAt{18}{332.098}{328.49}{C}%
\TextAt{13}{340.732}{328.49}{\Character{58}}%
\TextAt{8}{517.802}{328.49}{(1)}%
\TextAt{8}{79.226}{302.305}{It}%
\TextAt{8}{91.909}{302.305}{admits}%
\TextAt{8}{130.354}{302.305}{the}%
\TextAt{8}{150.521}{302.305}{following}%
\TextAt{8}{199.965}{302.305}{two}%
\TextAt{8}{222.092}{302.305}{further}%
\TextAt{8}{261.447}{302.305}{interpretations,}%
\TextAt{8}{343.806}{302.305}{between}%
\TextAt{8}{389.007}{302.305}{which}%
\TextAt{8}{422.503}{302.305}{we}%
\TextAt{8}{439.74}{302.305}{will}%
\TextAt{8}{461.855}{302.305}{freely}%
\TextAt{8}{493.726}{302.305}{go}%
\TextAt{8}{509.349}{302.305}{back}%
\TextAt{8}{79.226}{285.259}{and}%
\TextAt{8}{101.989}{285.259}{forth}%
\TextAt{8}{130.932}{285.259}{in}%
\TextAt{8}{144.585}{285.259}{the}%
\TextAt{8}{164.741}{285.259}{sequel}%
\TextAt{8}{199.591}{285.259}{without}%
\TextAt{8}{243.167}{285.259}{further}%
\TextAt{8}{282.512}{285.259}{explanation.}%
\TextAt{8}{350.485}{285.259}{The}%
\TextAt{8}{374.539}{285.259}{set}%
\TextAt{13}{255.158}{259.075}{X}%
\TextAt{9}{265.813}{259.075}{\Character{40}}%
\TextAt{18}{270.366}{259.075}{C}%
\TextAt{9}{279}{259.075}{\Character{41}}%
\TextAt{9}{286.855}{259.075}{\Character{61}}%
\TextAt{8}{299.298}{259.075}{Hom}%
\TextAt{9}{323.687}{259.075}{\Character{40}}%
\TextAt{13}{328.239}{259.075}{K}%
\TextAt{13}{338.401}{259.075}{\Character{59}}%
\TextAt{18}{343.651}{259.075}{C}%
\TextAt{9}{352.285}{259.075}{\Character{41}}%
\TextAt{8}{79.226}{232.89}{induces}%
\TextAt{8}{120.902}{232.89}{a}%
\TextAt{8}{130.657}{232.89}{canonical}%
\TextAt{8}{182.041}{232.89}{decomposition}%
\TextAt{8}{259.175}{232.89}{of}%
\TextAt{8}{272.505}{232.89}{rings}%
\TextAt{13}{200.794}{206.706}{K}%
\TextAt{25}{210.753}{204.912}{C}%
\NativeMathAt{220.328}{206.705}{11.955}{\cong}%
\TextAt{20}{239.265}{218.063}{\Character{77}}%
\TextAt{21}{232.947}{192.16}{\Character{27}}%
\NativeMathAt{238.073}{192.16}{7.97}{\in}%
\TextAt{21}{243.718}{192.16}{X}%
\TextAt{17}{251.305}{192.16}{\Character{40}}%
\TextAt{25}{254.598}{192.16}{C}%
\TextAt{17}{260.355}{192.16}{\Character{41}}%
\TextAt{18}{265.64}{206.706}{C}%
\TextAt{13}{274.275}{206.706}{\Character{59}}%
\TextAt{13}{303.41}{206.706}{a}%
\TextAt{16}{312.231}{206.706}{\Character{10}}%
\TextAt{13}{324.186}{206.706}{z}%
\NativeArrowAt{333.477}{206.706}{11.955}{19.26}{\longmapsto}%
\TextAt{20}{362.377}{218.063}{\Character{77}}%
\TextAt{21}{356.059}{192.16}{\Character{27}}%
\NativeMathAt{361.185}{192.16}{7.97}{\in}%
\TextAt{21}{366.83}{192.16}{X}%
\TextAt{17}{374.417}{192.16}{\Character{40}}%
\TextAt{25}{377.71}{192.16}{C}%
\TextAt{17}{383.467}{192.16}{\Character{41}}%
\TextAt{13}{388.752}{206.706}{z}%
\TextAt{13}{394.706}{206.706}{\Character{27}}%
\TextAt{13}{401.783}{206.706}{a\Character{58}}%
\TextAt{8}{517.802}{206.706}{(2)}%
\TextAt{8}{79.226}{168.521}{Alternatively,}%
\TextAt{8}{152.16}{168.521}{the}%
\TextAt{8}{172.026}{168.521}{right-hand}%
\TextAt{8}{229.3}{168.521}{side}%
\TextAt{8}{252.477}{168.521}{may}%
\TextAt{8}{277.556}{168.521}{be}%
\TextAt{8}{293.19}{168.521}{viewed}%
\TextAt{8}{331.264}{168.521}{as}%
\TextAt{8}{345.344}{168.521}{the}%
\TextAt{8}{365.21}{168.521}{set}%
\TextAt{18}{383.187}{168.521}{C}%
\TextAt{21}{391.821}{172.86}{X}%
\TextAt{17}{399.408}{172.86}{\Character{40}}%
\TextAt{25}{402.701}{172.86}{C}%
\TextAt{17}{408.458}{172.86}{\Character{41}}%
\TextAt{9}{415.57}{168.521}{\Character{61}}%
\TextAt{8}{427.995}{168.521}{Hom}%
\TextAt{9}{452.384}{168.521}{\Character{40}}%
\TextAt{13}{456.936}{168.521}{X}%
\TextAt{9}{467.591}{168.521}{\Character{40}}%
\TextAt{18}{472.144}{168.521}{C}%
\TextAt{9}{480.778}{168.521}{\Character{41}}%
\TextAt{13}{485.33}{168.521}{\Character{59}}%
\TextAt{18}{490.574}{168.521}{C}%
\TextAt{9}{499.209}{168.521}{\Character{41}}%
\TextAt{8}{507.368}{168.521}{of}%
\TextAt{8}{520.408}{168.521}{all}%
\TextAt{8}{79.226}{151.475}{functions}%
\TextAt{13}{129.689}{151.475}{x}%
\TextAt{9}{139.662}{151.475}{\Character{58}}%
\TextAt{13}{146.234}{151.475}{X}%
\TextAt{9}{156.889}{151.475}{\Character{40}}%
\TextAt{18}{161.442}{151.475}{C}%
\TextAt{9}{170.076}{151.475}{\Character{41}}%
\NativeMathAt{177.949}{151.475}{11.955}{\rightarrow}%
\TextAt{18}{193.225}{151.475}{C}%
\TextAt{8}{201.86}{151.475}{,}%
\TextAt{8}{209.009}{151.475}{i.e.,}%
\TextAt{13}{250.412}{125.291}{K}%
\TextAt{25}{260.371}{123.497}{C}%
\NativeMathAt{269.947}{125.29}{11.955}{\cong}%
\TextAt{8}{282.566}{125.291}{Hom}%
\TextAt{9}{306.954}{125.291}{\Character{40}}%
\TextAt{13}{311.507}{125.291}{X}%
\TextAt{9}{322.162}{125.291}{\Character{40}}%
\TextAt{18}{326.714}{125.291}{C}%
\TextAt{9}{335.349}{125.291}{\Character{41}}%
\TextAt{13}{339.901}{125.291}{\Character{59}}%
\TextAt{18}{345.145}{125.291}{C}%
\TextAt{9}{353.779}{125.291}{\Character{41}}%
\TextAt{13}{358.332}{125.291}{\Character{58}}%
\TextAt{8}{517.802}{125.291}{(3)}%
\TextAt{8}{79.226}{99.106}{The}%
\TextAt{8}{103.28}{99.106}{field}%
\TextAt{13}{128.637}{99.106}{K}%
\TextAt{8}{143.352}{99.106}{is}%
\TextAt{8}{155.116}{99.106}{embedded}%
\TextAt{8}{210.385}{99.106}{in}%
\TextAt{13}{224.037}{99.106}{K}%
\TextAt{25}{233.996}{97.313}{C}%
\TextAt{8}{244.148}{99.106}{via}%
\TextAt{13}{221.725}{72.921}{K}%
\NativeArrowAt{235.864}{72.921}{11.955}{19.26}{\longrightarrow}%
\TextAt{13}{258.446}{72.921}{K}%
\TextAt{16}{271.92}{72.921}{\Character{10}}%
\TextAt{25}{281.219}{71.128}{Q}%
\TextAt{18}{290.572}{72.921}{C}%
\TextAt{13}{299.207}{72.921}{\Character{59}}%
\TextAt{13}{328.342}{72.921}{a}%
\NativeArrowAt{337.827}{72.921}{11.955}{19.26}{\longmapsto}%
\TextAt{13}{360.409}{72.921}{a}%
\TextAt{16}{369.211}{72.921}{\Character{10}}%
\TextAt{9}{381.166}{72.921}{1}%
\TextAt{13}{387.019}{72.921}{\Character{59}}%
\TextAt{8}{297.211}{43.033}{206}%
\Destination{522.36}{345.535}{nav.AMS.28}%
\Destination{522.36}{223.755}{nav.AMS.29}%
\Destination{522.36}{142.335}{nav.AMS.30}%
\Destination{268.01}{345.535}{nav.equation.3.4.1}%
\Destination{78.23}{809.815}{nav.page.206}%
\Destination{79.23}{637.675}{nav.section.3.4}%
\Destination{79.36}{560.155}{nav.print.225}%
\Destination{79.23}{637.675}{outline.35}%
\Hotlink{367.683}{440.655}{19.22}{11.047}{nav.cite.N015}%
}

% Page 219 (body page 207)
\Page{%
\TextAt{2}{79.226}{796.943}{III.}%
\TextAt{2}{98.317}{796.943}{Riemann–Roch}%
\TextAt{2}{174.375}{796.943}{Theory}%
\Vector{q 0 G 0 g 0.299 w 7.226 18.729 m 460.769 18.729 l S Q}%
\TextAt{8}{79.226}{761.124}{and}%
\TextAt{8}{102.198}{761.124}{we}%
\TextAt{8}{119.635}{761.124}{identify}%
\TextAt{8}{162.441}{761.124}{it}%
\TextAt{8}{174.355}{761.124}{with}%
\TextAt{8}{201.225}{761.124}{its}%
\TextAt{8}{217.753}{761.124}{image.}%
\TextAt{8}{256.754}{761.124}{In}%
\TextAt{8}{271.585}{761.124}{the}%
\TextAt{8}{291.951}{761.124}{interpretation}%
\TextAt{8}{366.642}{761.124}{(2),}%
\TextAt{8}{389.021}{761.124}{the}%
\TextAt{8}{409.387}{761.124}{image}%
\TextAt{8}{443.418}{761.124}{of}%
\TextAt{13}{456.957}{761.124}{a}%
\NativeMathAt{466.78}{761.124}{11.955}{\in}%
\TextAt{13}{478.429}{761.124}{K}%
\TextAt{8}{493.353}{761.124}{appears}%
\TextAt{8}{79.226}{744.078}{as}%
\TextAt{8}{95.374}{744.078}{the}%
\TextAt{8}{117.308}{744.078}{tuple}%
\TextAt{20}{148.997}{753.044}{\Character{76}}%
\TextAt{21}{162.281}{740.591}{\Character{27}}%
\TextAt{13}{169.897}{744.078}{\Character{27}}%
\TextAt{13}{176.963}{744.078}{a}%
\TextAt{8}{188.799}{744.078}{of}%
\TextAt{8}{203.906}{744.078}{conjugates}%
\TextAt{8}{263.308}{744.078}{of}%
\TextAt{13}{278.416}{744.078}{a}%
\TextAt{8}{284.56}{744.078}{,}%
\TextAt{8}{293.932}{744.078}{and}%
\TextAt{8}{318.472}{744.078}{in}%
\TextAt{8}{333.902}{744.078}{the}%
\TextAt{8}{355.836}{744.078}{interpretation}%
\TextAt{8}{432.094}{744.078}{(3)}%
\TextAt{8}{452.737}{744.078}{as}%
\TextAt{8}{468.885}{744.078}{the}%
\TextAt{8}{490.819}{744.078}{function}%
\TextAt{13}{79.226}{727.032}{x}%
\TextAt{9}{85.878}{727.032}{\Character{40}}%
\TextAt{13}{90.431}{727.032}{\Character{27}}%
\TextAt{9}{97.513}{727.032}{\Character{41}}%
\TextAt{9}{105.368}{727.032}{\Character{61}}%
\TextAt{13}{117.812}{727.032}{\Character{27}}%
\TextAt{13}{124.878}{727.032}{a}%
\TextAt{8}{131.039}{727.032}{.}%
\TextAt{8}{97.159}{706.997}{We}%
\TextAt{8}{117.548}{706.997}{denote}%
\TextAt{8}{155.511}{706.997}{the}%
\TextAt{8}{175.913}{706.997}{generator}%
\TextAt{8}{228.198}{706.997}{of}%
\TextAt{8}{241.773}{706.997}{the}%
\TextAt{8}{262.174}{706.997}{Galois}%
\TextAt{8}{298.332}{706.997}{group}%
\TextAt{13}{331.753}{706.997}{G}%
\TextAt{9}{340.986}{706.997}{\Character{40}}%
\TextAt{18}{345.539}{706.997}{C}%
\TextAt{16}{357.911}{706.997}{\Character{106}}%
\TextAt{18}{364.97}{706.997}{R}%
\TextAt{9}{373.604}{706.997}{\Character{41}}%
\TextAt{8}{382.298}{706.997}{by}%
\TextAt{13}{398.802}{706.997}{F}%
\NativeMathAt{406.38}{705.204}{7.97}{\infty}%
\TextAt{8}{415.346}{706.997}{,}%
\TextAt{8}{422.802}{706.997}{or}%
\TextAt{8}{437.357}{706.997}{simply}%
\TextAt{8}{475.057}{706.997}{by}%
\TextAt{13}{491.561}{706.997}{F}%
\TextAt{8}{500.764}{706.997}{.}%
\TextAt{8}{509.947}{706.997}{This}%
\TextAt{8}{79.226}{689.951}{underlines}%
\TextAt{8}{135.097}{689.951}{the}%
\TextAt{8}{155.138}{689.951}{fact}%
\TextAt{8}{178.108}{689.951}{that}%
\TextAt{8}{203.361}{689.951}{it}%
\TextAt{8}{214.95}{689.951}{has}%
\TextAt{8}{235.708}{689.951}{a}%
\TextAt{8}{245.348}{689.951}{position}%
\TextAt{8}{289.849}{689.951}{analogous}%
\TextAt{8}{343.795}{689.951}{to}%
\TextAt{8}{357.99}{689.951}{the}%
\TextAt{8}{378.031}{689.951}{Frobenius}%
\TextAt{8}{431.786}{689.951}{automorphism}%
\TextAt{13}{509.14}{689.951}{F}%
\TextAt{21}{516.718}{688.158}{p}%
\NativeMathAt{524.799}{689.951}{11.955}{\in}%
\TextAt{13}{79.226}{672.905}{G}%
\TextAt{9}{88.46}{672.905}{\Character{40}}%
\Vector{q 0 G 0 g 0.478 w 21.012 -91.186 m 28.318 -91.186 l S Q}%
\TextAt{18}{93.012}{672.905}{F}%
\TextAt{21}{100.318}{671.112}{p}%
\TextAt{16}{108.4}{672.905}{\Character{106}}%
\TextAt{18}{115.041}{672.905}{F}%
\TextAt{21}{122.347}{671.112}{p}%
\TextAt{9}{127.108}{672.905}{\Character{41}}%
\TextAt{8}{131.661}{672.905}{,}%
\TextAt{8}{137.845}{672.905}{in}%
\TextAt{8}{150.291}{672.905}{accordance}%
\TextAt{8}{208.92}{672.905}{with}%
\TextAt{8}{234.374}{672.905}{our}%
\TextAt{8}{253.981}{672.905}{decision}%
\TextAt{8}{297.056}{672.905}{of}%
\TextAt{8}{309.18}{672.905}{§1}%
\TextAt{8}{323.408}{672.905}{to}%
\TextAt{8}{336.511}{672.905}{view}%
\TextAt{8}{362.288}{672.905}{the}%
\TextAt{8}{381.238}{672.905}{extension}%
\TextAt{18}{431.797}{672.905}{C}%
\TextAt{16}{443.752}{672.905}{\Character{106}}%
\TextAt{18}{450.394}{672.905}{R}%
\TextAt{8}{461.719}{672.905}{as}%
\TextAt{8}{474.883}{672.905}{unramified.}%
\TextAt{13}{79.226}{655.859}{F}%
\TextAt{8}{92.673}{655.859}{induces}%
\TextAt{8}{134.695}{655.859}{an}%
\TextAt{8}{151.3}{655.859}{involution}%
\TextAt{13}{206.616}{655.859}{F}%
\TextAt{8}{220.063}{655.859}{on}%
\TextAt{13}{236.669}{655.859}{K}%
\TextAt{25}{246.628}{654.066}{C}%
\TextAt{8}{257.126}{655.859}{which,}%
\TextAt{8}{294.296}{655.859}{in}%
\TextAt{8}{308.295}{655.859}{the}%
\TextAt{8}{328.798}{655.859}{representation}%
\TextAt{13}{405.633}{655.859}{K}%
\TextAt{25}{415.592}{654.066}{C}%
\TextAt{9}{425.757}{655.859}{\Character{61}}%
\TextAt{8}{438.773}{655.859}{Hom}%
\TextAt{9}{463.161}{655.859}{\Character{40}}%
\TextAt{13}{467.714}{655.859}{X}%
\TextAt{9}{478.369}{655.859}{\Character{40}}%
\TextAt{18}{482.921}{655.859}{C}%
\TextAt{9}{491.555}{655.859}{\Character{41}}%
\TextAt{13}{496.108}{655.859}{\Character{59}}%
\TextAt{18}{501.352}{655.859}{C}%
\TextAt{9}{509.986}{655.859}{\Character{41}}%
\TextAt{8}{518.782}{655.859}{for}%
\TextAt{13}{79.226}{638.813}{x}%
\TextAt{9}{89.199}{638.813}{\Character{58}}%
\TextAt{13}{95.772}{638.813}{X}%
\TextAt{9}{106.427}{638.813}{\Character{40}}%
\TextAt{18}{110.979}{638.813}{C}%
\TextAt{9}{119.613}{638.813}{\Character{41}}%
\NativeMathAt{127.486}{638.813}{11.955}{\rightarrow}%
\TextAt{18}{142.762}{638.813}{C}%
\TextAt{8}{151.397}{638.813}{,}%
\TextAt{8}{158.546}{638.813}{is}%
\TextAt{8}{170.31}{638.813}{given}%
\TextAt{8}{200.879}{638.813}{by}%
\TextAt{9}{264.506}{621.767}{\Character{40}}%
\TextAt{13}{269.058}{621.767}{F}%
\TextAt{13}{278.24}{621.767}{x}%
\TextAt{9}{284.914}{621.767}{\Character{41}\Character{40}}%
\TextAt{13}{294.018}{621.767}{\Character{27}}%
\TextAt{9}{301.101}{621.767}{\Character{41}}%
\TextAt{9}{308.956}{621.767}{\Character{61}}%
\Vector{q 0 G 0 g 0.478 w 249.399 -141.593 m 272.238 -141.593 l S Q}%
\TextAt{13}{321.399}{621.767}{x}%
\TextAt{9}{328.051}{621.767}{\Character{40}}%
\Vector{q 0 G 0 g 0.478 w 260.604 -145.412 m 267.686 -145.412 l S Q}%
\TextAt{13}{332.604}{621.767}{\Character{27}}%
\TextAt{9}{339.686}{621.767}{\Character{41}}%
\TextAt{13}{344.238}{621.767}{\Character{58}}%
\TextAt{13}{79.226}{598.245}{F}%
\TextAt{8}{91.517}{598.245}{is}%
\TextAt{8}{102.471}{598.245}{an}%
\TextAt{8}{117.92}{598.245}{automorphism}%
\TextAt{8}{194.579}{598.245}{of}%
\TextAt{8}{207.099}{598.245}{the}%
\TextAt{18}{226.445}{598.245}{R}%
\TextAt{8}{235.08}{598.245}{-algebra}%
\TextAt{13}{279.149}{598.245}{K}%
\TextAt{25}{289.108}{596.452}{C}%
\TextAt{8}{295.363}{598.245}{.}%
\TextAt{8}{303.541}{598.245}{It}%
\TextAt{8}{315.403}{598.245}{is}%
\TextAt{8}{326.357}{598.245}{called}%
\TextAt{8}{358.711}{598.245}{the}%
\TextAt{12}{378.057}{598.245}{Frobenius}%
\TextAt{12}{439.137}{598.245}{correspondence}%
\TextAt{8}{529.518}{598.245}{.}%
\TextAt{8}{79.226}{581.199}{Sometimes}%
\TextAt{8}{137.818}{581.199}{we}%
\TextAt{8}{155.046}{581.199}{also}%
\TextAt{8}{178.526}{581.199}{consider,}%
\TextAt{8}{227.363}{581.199}{besides}%
\TextAt{13}{267.472}{581.199}{F}%
\TextAt{8}{276.676}{581.199}{,}%
\TextAt{8}{283.825}{581.199}{the}%
\TextAt{8}{303.982}{581.199}{involution}%
\TextAt{13}{358.951}{581.199}{z}%
\NativeArrowAt{368.243}{581.199}{11.955}{11.955}{\mapsto}%
\Vector{q 0 G 0 g 0.478 w 311.519 -185.98 m 317.489 -185.98 l S Q}%
\TextAt{13}{383.519}{581.199}{z}%
\TextAt{8}{393.387}{581.199}{on}%
\TextAt{13}{409.646}{581.199}{K}%
\TextAt{25}{419.605}{579.406}{C}%
\TextAt{8}{429.757}{581.199}{which}%
\TextAt{8}{463.243}{581.199}{is}%
\TextAt{8}{475.007}{581.199}{given}%
\TextAt{8}{505.576}{581.199}{by}%
\Vector{q 0 G 0 g 0.478 w 202.341 -214.981 m 208.312 -214.981 l S Q}%
\TextAt{13}{274.341}{552.198}{z}%
\TextAt{9}{280.312}{552.198}{\Character{40}}%
\TextAt{13}{284.864}{552.198}{\Character{27}}%
\TextAt{9}{291.946}{552.198}{\Character{41}}%
\TextAt{9}{299.801}{552.198}{\Character{61}}%
\Vector{q 0 G 0 g 0.478 w 240.245 -211.162 m 262.403 -211.162 l S Q}%
\TextAt{13}{312.245}{552.198}{z}%
\TextAt{9}{318.216}{552.198}{\Character{40}}%
\TextAt{13}{322.768}{552.198}{\Character{27}}%
\TextAt{9}{329.85}{552.198}{\Character{41}}%
\TextAt{13}{334.403}{552.198}{\Character{58}}%
\TextAt{8}{79.226}{523.197}{We}%
\TextAt{8}{100.68}{523.197}{call}%
\TextAt{8}{123.45}{523.197}{it}%
\TextAt{8}{136.464}{523.197}{the}%
\TextAt{12}{157.93}{523.197}{conjugation}%
\TextAt{8}{226.146}{523.197}{.}%
\TextAt{8}{238.524}{523.197}{Finally,}%
\TextAt{8}{282.267}{523.197}{we}%
\TextAt{8}{300.804}{523.197}{call}%
\TextAt{8}{323.574}{523.197}{an}%
\TextAt{8}{341.142}{523.197}{element}%
\TextAt{13}{385.694}{523.197}{x}%
\NativeMathAt{397.899}{523.197}{11.955}{\in}%
\TextAt{13}{411.422}{523.197}{K}%
\TextAt{25}{421.381}{521.404}{C}%
\TextAt{8}{427.636}{523.197}{,}%
\TextAt{8}{436.422}{523.197}{that}%
\TextAt{8}{463.101}{523.197}{is,}%
\TextAt{8}{479.754}{523.197}{a}%
\TextAt{8}{490.819}{523.197}{function}%
\TextAt{13}{79.226}{506.151}{x}%
\TextAt{9}{90.284}{506.151}{\Character{58}}%
\TextAt{13}{97.942}{506.151}{X}%
\TextAt{9}{108.597}{506.151}{\Character{40}}%
\TextAt{18}{113.149}{506.151}{C}%
\TextAt{9}{121.783}{506.151}{\Character{41}}%
\NativeMathAt{130.741}{506.151}{11.955}{\rightarrow}%
\TextAt{18}{147.103}{506.151}{C}%
\TextAt{8}{155.737}{506.151}{,}%
\TextAt{12}{163.682}{506.151}{positive}%
\TextAt{8}{213.682}{506.151}{(written}%
\TextAt{13}{259.844}{506.151}{x}%
\TextAt{13}{270.891}{506.151}{\Character{62}}%
\TextAt{9}{284.412}{506.151}{0}%
\TextAt{8}{290.265}{506.151}{)}%
\TextAt{8}{299.354}{506.151}{if}%
\TextAt{8}{310.714}{506.151}{it}%
\TextAt{8}{323.055}{506.151}{takes}%
\TextAt{8}{353.676}{506.151}{real}%
\TextAt{8}{377.075}{506.151}{values,}%
\TextAt{8}{415.984}{506.151}{and}%
\TextAt{8}{439.383}{506.151}{if}%
\TextAt{13}{450.744}{506.151}{x}%
\TextAt{9}{457.396}{506.151}{\Character{40}}%
\TextAt{13}{461.948}{506.151}{\Character{27}}%
\TextAt{9}{469.03}{506.151}{\Character{41}}%
\TextAt{13}{477.989}{506.151}{\Character{62}}%
\TextAt{9}{491.499}{506.151}{0}%
\TextAt{8}{501.886}{506.151}{for}%
\TextAt{8}{520.408}{506.151}{all}%
\TextAt{13}{79.226}{489.105}{\Character{27}}%
\NativeMathAt{89.629}{489.105}{11.955}{\in}%
\TextAt{13}{100.92}{489.105}{X}%
\TextAt{9}{111.575}{489.105}{\Character{40}}%
\TextAt{18}{116.128}{489.105}{C}%
\TextAt{9}{124.762}{489.105}{\Character{41}}%
\TextAt{8}{129.314}{489.105}{.}%
\TextAt{8}{97.159}{469.07}{By}%
\TextAt{8}{115.32}{469.07}{convention}%
\TextAt{8}{173.662}{469.07}{every}%
\TextAt{14}{204.352}{469.07}{o}%
\TextAt{8}{210.196}{469.07}{-module}%
\TextAt{8}{255.185}{469.07}{considered}%
\TextAt{8}{312.272}{469.07}{in}%
\TextAt{8}{325.722}{469.07}{the}%
\TextAt{8}{345.677}{469.07}{sequel}%
\TextAt{8}{380.324}{469.07}{will}%
\TextAt{8}{402.227}{469.07}{be}%
\TextAt{8}{417.949}{469.07}{supposed}%
\TextAt{8}{468.269}{469.07}{to}%
\TextAt{8}{482.377}{469.07}{be}%
\TextAt{19}{498.099}{469.07}{finitely}%
\TextAt{19}{79.226}{452.025}{generated}%
\TextAt{8}{126.903}{452.025}{.}%
\TextAt{8}{135.352}{452.025}{For}%
\TextAt{8}{156.321}{452.025}{every}%
\TextAt{8}{187.213}{452.025}{such}%
\TextAt{14}{213.61}{452.025}{o}%
\TextAt{8}{219.454}{452.025}{-module}%
\TextAt{13}{264.645}{452.025}{M}%
\TextAt{8}{277.219}{452.025}{,}%
\TextAt{8}{284.368}{452.025}{we}%
\TextAt{8}{301.595}{452.025}{put}%
\TextAt{13}{266.894}{423.023}{M}%
\TextAt{25}{278.218}{421.23}{C}%
\TextAt{9}{287.794}{423.023}{\Character{61}}%
\TextAt{13}{300.219}{423.023}{M}%
\TextAt{16}{315.449}{423.023}{\Character{10}}%
\TextAt{25}{324.748}{421.23}{Z}%
\TextAt{18}{333.216}{423.023}{C}%
\TextAt{13}{341.85}{423.023}{\Character{58}}%
\TextAt{8}{79.226}{394.022}{This}%
\TextAt{8}{106.463}{394.022}{is}%
\TextAt{8}{118.744}{394.022}{a}%
\TextAt{8}{129.017}{394.022}{module}%
\TextAt{8}{170.827}{394.022}{over}%
\TextAt{8}{196.39}{394.022}{the}%
\TextAt{8}{217.064}{394.022}{ring}%
\TextAt{13}{241.646}{394.022}{K}%
\TextAt{25}{251.606}{392.229}{C}%
\TextAt{9}{262.062}{394.022}{\Character{61}}%
\TextAt{14}{275.369}{394.022}{o}%
\TextAt{16}{284.222}{394.022}{\Character{10}}%
\TextAt{25}{293.521}{392.229}{Z}%
\TextAt{18}{302.341}{394.022}{C}%
\TextAt{8}{310.976}{394.022}{,}%
\TextAt{8}{318.771}{394.022}{and}%
\TextAt{8}{342.051}{394.022}{viewing}%
\TextAt{14}{385.164}{394.022}{o}%
\TextAt{8}{395.423}{394.022}{as}%
\TextAt{8}{410.31}{394.022}{a}%
\TextAt{8}{420.583}{394.022}{subring}%
\TextAt{8}{462.788}{394.022}{of}%
\TextAt{13}{476.635}{394.022}{K}%
\TextAt{25}{486.594}{392.229}{C}%
\TextAt{8}{492.848}{394.022}{—as}%
\TextAt{8}{519.439}{394.022}{we}%
\TextAt{8}{79.226}{376.976}{agreed}%
\TextAt{8}{116.299}{376.976}{above—we}%
\TextAt{8}{174.509}{376.976}{may}%
\TextAt{8}{199.878}{376.976}{also}%
\TextAt{8}{223.358}{376.976}{write}%
\TextAt{13}{265.439}{347.975}{M}%
\TextAt{25}{276.763}{346.182}{C}%
\TextAt{9}{286.338}{347.975}{\Character{61}}%
\TextAt{13}{298.764}{347.975}{M}%
\TextAt{16}{313.994}{347.975}{\Character{10}}%
\TextAt{27}{323.293}{346.182}{o}%
\TextAt{13}{330.344}{347.975}{K}%
\TextAt{25}{340.303}{346.182}{C}%
\TextAt{8}{79.226}{318.974}{as}%
\TextAt{13}{93.596}{318.974}{M}%
\TextAt{16}{108.827}{318.974}{\Character{10}}%
\TextAt{25}{118.125}{317.181}{Z}%
\TextAt{18}{126.593}{318.974}{C}%
\TextAt{9}{138.548}{318.974}{\Character{61}}%
\TextAt{13}{150.974}{318.974}{M}%
\TextAt{16}{166.204}{318.974}{\Character{10}}%
\TextAt{27}{175.503}{317.181}{o}%
\TextAt{9}{182.554}{318.974}{\Character{40}}%
\TextAt{14}{187.106}{318.974}{o}%
\TextAt{16}{195.607}{318.974}{\Character{10}}%
\TextAt{25}{204.905}{317.181}{Z}%
\TextAt{18}{213.374}{318.974}{C}%
\TextAt{9}{222.008}{318.974}{\Character{41}}%
\TextAt{8}{226.56}{318.974}{.}%
\TextAt{8}{235.008}{318.974}{The}%
\TextAt{8}{259.062}{318.974}{involution}%
\TextAt{13}{314.032}{318.974}{x}%
\NativeArrowAt{324.005}{318.974}{11.955}{11.955}{\mapsto}%
\TextAt{13}{339.281}{318.974}{F}%
\TextAt{13}{348.475}{318.974}{x}%
\TextAt{8}{359.034}{318.974}{on}%
\TextAt{13}{375.293}{318.974}{K}%
\TextAt{25}{385.252}{317.181}{C}%
\TextAt{8}{395.404}{318.974}{induces}%
\TextAt{8}{437.08}{318.974}{the}%
\TextAt{8}{457.236}{318.974}{involution}%
\TextAt{13}{256.96}{289.973}{F}%
\TextAt{9}{266.163}{289.973}{\Character{40}}%
\TextAt{13}{270.716}{289.973}{a}%
\TextAt{16}{279.517}{289.973}{\Character{10}}%
\TextAt{13}{291.473}{289.973}{x}%
\TextAt{9}{298.125}{289.973}{\Character{41}}%
\TextAt{9}{305.98}{289.973}{\Character{61}}%
\TextAt{13}{318.423}{289.973}{a}%
\TextAt{16}{327.225}{289.973}{\Character{10}}%
\TextAt{13}{339.18}{289.973}{F}%
\TextAt{13}{348.374}{289.973}{x}%
\TextAt{8}{79.226}{260.972}{on}%
\TextAt{13}{95.485}{260.972}{M}%
\TextAt{25}{106.81}{259.178}{C}%
\TextAt{8}{113.064}{260.972}{.}%
\TextAt{8}{121.513}{260.972}{In}%
\TextAt{8}{136.134}{260.972}{the}%
\TextAt{8}{156.29}{260.972}{representation}%
\TextAt{13}{232.779}{260.972}{M}%
\TextAt{25}{244.104}{259.178}{C}%
\TextAt{9}{253.679}{260.972}{\Character{61}}%
\TextAt{13}{266.105}{260.972}{M}%
\TextAt{16}{281.335}{260.972}{\Character{10}}%
\TextAt{25}{290.634}{259.178}{Z}%
\TextAt{18}{299.102}{260.972}{C}%
\TextAt{8}{307.736}{260.972}{,}%
\TextAt{8}{314.885}{260.972}{one}%
\TextAt{8}{336.345}{260.972}{clearly}%
\TextAt{8}{373.741}{260.972}{has}%
\TextAt{13}{260.617}{231.971}{F}%
\TextAt{9}{269.821}{231.971}{\Character{40}}%
\TextAt{13}{274.373}{231.971}{a}%
\TextAt{16}{283.175}{231.971}{\Character{10}}%
\TextAt{13}{295.13}{231.971}{z}%
\TextAt{9}{301.101}{231.971}{\Character{41}}%
\TextAt{9}{308.956}{231.971}{\Character{61}}%
\TextAt{13}{321.399}{231.971}{a}%
\TextAt{16}{330.201}{231.971}{\Character{10}}%
\Vector{q 0 G 0 g 0.478 w 270.156 -535.208 m 276.127 -535.208 l S Q}%
\TextAt{13}{342.156}{231.971}{z}%
\TextAt{13}{348.11}{231.971}{\Character{58}}%
\TextAt{12}{79.226}{194.003}{Definition}%
\TextAt{12}{141.229}{194.003}{(4.1).}%
\TextAt{8}{175.862}{194.003}{A}%
\TextAt{12}{187.851}{194.003}{hermitian}%
\TextAt{12}{248.599}{194.003}{metric}%
\TextAt{8}{289.616}{194.003}{on}%
\TextAt{8}{305.192}{194.003}{the}%
\TextAt{13}{324.665}{194.003}{K}%
\TextAt{25}{334.625}{192.21}{C}%
\TextAt{8}{340.879}{194.003}{-module}%
\TextAt{13}{385.387}{194.003}{M}%
\TextAt{25}{396.711}{192.21}{C}%
\TextAt{8}{406.18}{194.003}{is}%
\TextAt{8}{417.261}{194.003}{a}%
\TextAt{8}{426.334}{194.003}{sesquilinear}%
\TextAt{8}{488.535}{194.003}{mapping}%
\TextAt{16}{237.574}{165.002}{\Character{104}}%
\TextAt{13}{246.12}{165.002}{\Character{59}}%
\TextAt{16}{255.262}{165.002}{\Character{105}}%
\TextAt{21}{259.911}{163.209}{M}%
\TextAt{9}{272.712}{165.002}{\Character{58}}%
\TextAt{13}{279.284}{165.002}{M}%
\TextAt{25}{290.609}{163.209}{C}%
\TextAt{16}{299.52}{165.002}{\Character{2}}%
\TextAt{13}{311.475}{165.002}{M}%
\TextAt{25}{322.8}{163.209}{C}%
\NativeArrowAt{332.375}{165.002}{11.955}{19.26}{\longrightarrow}%
\TextAt{13}{354.957}{165.002}{K}%
\TextAt{25}{364.916}{163.209}{C}%
\TextAt{13}{371.171}{165.002}{\Character{59}}%
\TextAt{8}{79.226}{136.001}{i.e.,}%
\TextAt{8}{101.331}{136.001}{a}%
\TextAt{13}{111.087}{136.001}{K}%
\TextAt{25}{121.046}{134.208}{C}%
\TextAt{8}{127.3}{136.001}{-linear}%
\TextAt{8}{163.716}{136.001}{form}%
\TextAt{16}{191.356}{136.001}{\Character{104}}%
\TextAt{13}{196.005}{136.001}{x\Character{59}}%
\TextAt{13}{207.877}{136.001}{y}%
\TextAt{16}{214.038}{136.001}{\Character{105}}%
\TextAt{21}{218.687}{134.208}{M}%
\TextAt{8}{232.065}{136.001}{in}%
\TextAt{8}{245.717}{136.001}{the}%
\TextAt{8}{265.874}{136.001}{first}%
\TextAt{8}{289.999}{136.001}{variable}%
\TextAt{8}{333.911}{136.001}{satisfying}%
\Vector{q 0 G 0 g 0.478 w 187.688 -656.36 m 224.499 -656.36 l S Q}%
\TextAt{16}{259.688}{107}{\Character{104}}%
\TextAt{13}{264.337}{107}{x\Character{59}}%
\TextAt{13}{276.209}{107}{y}%
\TextAt{16}{282.37}{107}{\Character{105}}%
\TextAt{21}{287.019}{105.206}{M}%
\TextAt{9}{299.82}{107}{\Character{61}}%
\TextAt{16}{312.245}{107}{\Character{104}}%
\TextAt{13}{316.894}{107}{y}%
\TextAt{13}{323.015}{107}{\Character{59}}%
\TextAt{13}{328.251}{107}{x}%
\TextAt{16}{334.927}{107}{\Character{105}}%
\TextAt{21}{339.577}{105.206}{M}%
\TextAt{13}{349.056}{107}{\Character{59}}%
\TextAt{8}{297.211}{43.033}{207}%
\Destination{78.23}{809.815}{nav.page.207}%
\Destination{79.36}{670.855}{nav.print.226}%
\Destination{79.36}{207.955}{nav.number.3.4.1}%
\Hotlink{314.359}{670.755}{6.858}{11.047}{nav.section.3.1}%
}

% Page 220 (body page 208)
\Page{%
\TextAt{2}{79.226}{796.943}{III.}%
\TextAt{2}{98.317}{796.943}{Riemann–Roch}%
\TextAt{2}{174.375}{796.943}{Theory}%
\Vector{q 0 G 0 g 0.299 w 7.226 18.729 m 460.769 18.729 l S Q}%
\ErrataAt{79.19}{761.124}{455.81}{such that $\langle x,x\rangle_M$ has nonnegative coordinates for all $x$ and is zero only when $x=0$.}%
\TextAt{8}{97.159}{741.089}{The}%
\TextAt{8}{121.213}{741.089}{metric}%
\TextAt{16}{157.628}{741.089}{\Character{104}}%
\TextAt{13}{166.175}{741.089}{\Character{59}}%
\TextAt{16}{175.316}{741.089}{\Character{105}}%
\TextAt{21}{179.966}{739.296}{M}%
\TextAt{8}{193.343}{741.089}{is}%
\TextAt{8}{205.107}{741.089}{called}%
\TextAt{13}{238.27}{741.089}{F}%
\TextAt{12}{247.474}{741.089}{-invariant}%
\TextAt{8}{307.704}{741.089}{if}%
\TextAt{8}{318.428}{741.089}{we}%
\TextAt{8}{335.655}{741.089}{have}%
\TextAt{8}{362.65}{741.089}{furthermore}%
\TextAt{13}{245.882}{713.981}{F}%
\TextAt{16}{255.086}{713.981}{\Character{104}}%
\TextAt{13}{259.735}{713.981}{x\Character{59}}%
\TextAt{13}{271.607}{713.981}{y}%
\TextAt{16}{277.768}{713.981}{\Character{105}}%
\TextAt{21}{282.417}{712.187}{M}%
\TextAt{9}{295.218}{713.981}{\Character{61}}%
\TextAt{16}{307.643}{713.981}{\Character{104}}%
\TextAt{13}{312.293}{713.981}{F}%
\TextAt{13}{321.475}{713.981}{x\Character{59}}%
\TextAt{13}{333.358}{713.981}{F}%
\TextAt{13}{342.552}{713.981}{y}%
\TextAt{16}{348.733}{713.981}{\Character{105}}%
\TextAt{21}{353.382}{712.187}{M}%
\TextAt{13}{362.862}{713.981}{\Character{58}}%
\TextAt{8}{97.159}{678.447}{This}%
\TextAt{8}{123.862}{678.447}{notion}%
\TextAt{8}{160.272}{678.447}{may}%
\TextAt{8}{185.624}{678.447}{be}%
\TextAt{8}{201.531}{678.447}{immediately}%
\TextAt{8}{268.176}{678.447}{reduced}%
\TextAt{8}{311.723}{678.447}{to}%
\TextAt{8}{326.017}{678.447}{the}%
\TextAt{8}{346.156}{678.447}{usual}%
\TextAt{8}{376.768}{678.447}{notion}%
\TextAt{8}{413.178}{678.447}{of}%
\TextAt{8}{426.491}{678.447}{a}%
\TextAt{8}{436.23}{678.447}{hermitian}%
\TextAt{8}{489.545}{678.447}{metric}%
\TextAt{8}{525.943}{678.447}{if}%
\TextAt{8}{79.226}{661.401}{we}%
\TextAt{8}{95.978}{661.401}{view}%
\TextAt{8}{122.485}{661.401}{the}%
\TextAt{13}{142.165}{661.401}{K}%
\TextAt{25}{152.124}{659.607}{C}%
\TextAt{8}{158.379}{661.401}{-module}%
\TextAt{13}{203.093}{661.401}{M}%
\TextAt{25}{214.418}{659.607}{C}%
\TextAt{8}{220.673}{661.401}{,}%
\TextAt{8}{227.441}{661.401}{by}%
\TextAt{8}{243.224}{661.401}{means}%
\TextAt{8}{278.578}{661.401}{of}%
\TextAt{8}{291.432}{661.401}{the}%
\TextAt{8}{311.113}{661.401}{decomposition}%
\TextAt{13}{387.772}{661.401}{K}%
\TextAt{25}{397.731}{659.607}{C}%
\TextAt{9}{407.306}{661.401}{\Character{61}}%
\TextAt{20}{419.732}{670.367}{\Character{76}}%
\TextAt{21}{433.015}{657.914}{\Character{27}}%
\TextAt{18}{440.631}{661.401}{C}%
\TextAt{8}{449.265}{661.401}{,}%
\TextAt{8}{456.034}{661.401}{as}%
\TextAt{8}{469.928}{661.401}{a}%
\TextAt{8}{479.208}{661.401}{direct}%
\TextAt{8}{511.896}{661.401}{sum}%
\TextAt{13}{232.745}{634.292}{M}%
\TextAt{25}{244.07}{632.499}{C}%
\TextAt{9}{253.645}{634.292}{\Character{61}}%
\TextAt{13}{266.07}{634.292}{M}%
\TextAt{16}{281.3}{634.292}{\Character{10}}%
\TextAt{27}{290.599}{632.499}{o}%
\TextAt{13}{297.65}{634.292}{K}%
\TextAt{25}{307.609}{632.499}{C}%
\TextAt{9}{317.184}{634.292}{\Character{61}}%
\TextAt{20}{335.927}{645.65}{\Character{77}}%
\TextAt{21}{329.61}{619.747}{\Character{27}}%
\NativeMathAt{334.735}{619.747}{7.97}{\in}%
\TextAt{21}{340.381}{619.747}{X}%
\TextAt{17}{347.967}{619.747}{\Character{40}}%
\TextAt{25}{351.261}{619.747}{C}%
\TextAt{17}{357.017}{619.747}{\Character{41}}%
\TextAt{13}{362.303}{634.292}{M}%
\TextAt{21}{373.627}{632.499}{\Character{27}}%
\TextAt{8}{79.226}{597.083}{of}%
\TextAt{18}{92.556}{597.083}{C}%
\TextAt{8}{101.19}{597.083}{-vector}%
\TextAt{8}{140.212}{597.083}{spaces}%
\TextAt{13}{264.179}{580.037}{M}%
\TextAt{21}{275.504}{578.244}{\Character{27}}%
\TextAt{9}{284.448}{580.037}{\Character{61}}%
\TextAt{13}{296.873}{580.037}{M}%
\TextAt{16}{312.104}{580.037}{\Character{10}}%
\TextAt{27}{321.402}{578.244}{o}%
\TextAt{21}{325.298}{578.244}{\Character{59}\Character{27}}%
\TextAt{18}{335.931}{580.037}{C}%
\TextAt{13}{344.565}{580.037}{\Character{58}}%
\TextAt{8}{79.226}{557.327}{The}%
\TextAt{8}{103.28}{557.327}{hermitian}%
\TextAt{8}{156.612}{557.327}{metric}%
\TextAt{16}{193.027}{557.327}{\Character{104}}%
\TextAt{13}{201.574}{557.327}{\Character{59}}%
\TextAt{16}{210.716}{557.327}{\Character{105}}%
\TextAt{21}{215.365}{555.534}{M}%
\TextAt{8}{228.742}{557.327}{then}%
\TextAt{8}{255.402}{557.327}{splits}%
\TextAt{8}{286.091}{557.327}{into}%
\TextAt{8}{309.834}{557.327}{the}%
\TextAt{8}{329.99}{557.327}{direct}%
\TextAt{8}{363.154}{557.327}{sum}%
\TextAt{16}{238.481}{530.219}{\Character{104}}%
\TextAt{13}{243.13}{530.219}{x\Character{59}}%
\TextAt{13}{255.002}{530.219}{y}%
\TextAt{16}{261.163}{530.219}{\Character{105}}%
\TextAt{21}{265.812}{528.425}{M}%
\TextAt{9}{278.613}{530.219}{\Character{61}}%
\TextAt{20}{297.356}{541.576}{\Character{77}}%
\TextAt{21}{291.038}{515.673}{\Character{27}}%
\NativeMathAt{296.164}{515.673}{7.97}{\in}%
\TextAt{21}{301.809}{515.673}{X}%
\TextAt{17}{309.396}{515.673}{\Character{40}}%
\TextAt{25}{312.689}{515.673}{C}%
\TextAt{17}{318.446}{515.673}{\Character{41}}%
\TextAt{16}{321.739}{530.219}{\Character{104}}%
\TextAt{13}{326.388}{530.219}{x}%
\TextAt{21}{333.04}{528.425}{\Character{27}}%
\TextAt{13}{338.664}{530.219}{\Character{59}}%
\TextAt{13}{343.888}{530.219}{y}%
\TextAt{21}{349.616}{528.425}{\Character{27}}%
\TextAt{16}{355.239}{530.219}{\Character{105}}%
\TextAt{21}{359.888}{528.425}{M}%
\TextAt{22}{368.021}{527.429}{\Character{27}}%
\TextAt{8}{79.226}{492.46}{of}%
\TextAt{8}{93.096}{492.46}{hermitian}%
\TextAt{8}{146.967}{492.46}{scalar}%
\TextAt{8}{180.742}{492.46}{products}%
\TextAt{16}{229.796}{492.46}{\Character{104}}%
\TextAt{13}{238.882}{492.46}{\Character{59}}%
\TextAt{16}{248.564}{492.46}{\Character{105}}%
\TextAt{21}{253.213}{490.667}{M}%
\TextAt{22}{261.346}{489.67}{\Character{27}}%
\TextAt{8}{271.277}{492.46}{on}%
\TextAt{8}{288.075}{492.46}{the}%
\TextAt{18}{308.771}{492.46}{C}%
\TextAt{8}{317.406}{492.46}{-vector}%
\TextAt{8}{356.967}{492.46}{spaces}%
\TextAt{13}{393.396}{492.46}{M}%
\TextAt{21}{404.721}{490.667}{\Character{27}}%
\TextAt{8}{410.344}{492.46}{.}%
\TextAt{8}{420.412}{492.46}{In}%
\TextAt{8}{435.572}{492.46}{this}%
\TextAt{8}{458.934}{492.46}{interpretation,}%
\TextAt{8}{79.226}{475.414}{the}%
\TextAt{13}{99.383}{475.414}{F}%
\TextAt{8}{108.586}{475.414}{-invariance}%
\TextAt{8}{167.776}{475.414}{of}%
\TextAt{16}{181.106}{475.414}{\Character{104}}%
\TextAt{13}{185.755}{475.414}{x\Character{59}}%
\TextAt{13}{197.627}{475.414}{y}%
\TextAt{16}{203.788}{475.414}{\Character{105}}%
\TextAt{21}{208.438}{473.621}{M}%
\TextAt{8}{221.815}{475.414}{amounts}%
\TextAt{8}{269.038}{475.414}{to}%
\TextAt{8}{283.348}{475.414}{the}%
\TextAt{8}{303.505}{475.414}{commutativity}%
\TextAt{8}{382.217}{475.414}{of}%
\TextAt{8}{395.547}{475.414}{the}%
\TextAt{8}{415.704}{475.414}{diagrams}%
\TextAt{13}{248.712}{443.406}{M}%
\TextAt{21}{260.037}{441.613}{\Character{27}}%
\TextAt{16}{268.317}{443.406}{\Character{2}}%
\TextAt{13}{280.272}{443.406}{M}%
\TextAt{21}{291.597}{441.613}{\Character{27}}%
\TextAt{18}{352.19}{443.406}{C}%
\TextAt{13}{248.712}{392.338}{M}%
\NativeDiagram{%
\path[draw=black,line width=0.359bp] (188.037,-378.768) -- (193.66,-378.768);%
}%
\TextAt{21}{260.037}{390.545}{\Character{27}}%
\TextAt{16}{268.317}{392.338}{\Character{2}}%
\TextAt{13}{280.272}{392.338}{M}%
\NativeDiagram{%
\path[draw=black,line width=0.359bp] (219.597,-378.768) -- (225.22,-378.768);%
}%
\TextAt{21}{291.597}{390.545}{\Character{27}}%
\TextAt{18}{350.564}{392.338}{C}%
\TextAt{13}{359.198}{392.338}{\Character{58}}%
\NativeDiagram{%
\path[draw=black,line width=0.478bp] (230.573,-327.605) -- (274.361,-327.605);%
\path[draw=black,line width=0.478bp,line cap=round,line join=round] (272.114,-324.736) .. controls (272.568,-326.458) and (273.581,-327.271) .. (274.601,-327.605) .. controls (273.581,-327.94) and (272.568,-328.753) .. (272.114,-330.475);%
}%
\TextAt{3}{310.121}{451.546}{\Character{104}}%
\TextAt{21}{317.312}{451.546}{\Character{59}}%
\TextAt{3}{323.561}{451.546}{\Character{105}}%
\TextAt{22}{326.855}{450.207}{M}%
\TextAt{22}{333.798}{449.211}{\Character{27}}%
\NativeDiagram{%
\path[draw=black,line width=0.478bp] (200.967,-336.967) -- (200.967,-368.436);%
\path[draw=black,line width=0.478bp,line cap=round,line join=round] (203.836,-366.188) .. controls (202.114,-366.642) and (201.301,-367.655) .. (200.967,-368.675) .. controls (200.632,-367.655) and (199.819,-366.642) .. (198.097,-366.188);%
}%
\TextAt{21}{250.494}{418.752}{F}%
\TextAt{3}{257.029}{418.752}{\Character{2}}%
\TextAt{21}{263.616}{418.752}{F}%
\NativeDiagram{%
\path[draw=black,line width=0.478bp] (284.509,-335.173) -- (284.509,-368.369);%
\path[draw=black,line width=0.478bp,line cap=round,line join=round] (287.378,-366.122) .. controls (285.656,-366.576) and (284.844,-367.588) .. (284.509,-368.608) .. controls (284.174,-367.588) and (283.361,-366.576) .. (281.64,-366.122);%
}%
\TextAt{21}{359.324}{419.267}{F}%
\NativeDiagram{%
\path[draw=black,line width=0.478bp] (230.573,-378.674) -- (272.736,-378.674);%
\path[draw=black,line width=0.478bp,line cap=round,line join=round] (270.488,-375.804) .. controls (270.942,-377.526) and (271.955,-378.339) .. (272.975,-378.674) .. controls (271.955,-379.008) and (270.942,-379.821) .. (270.488,-381.543);%
}%
\TextAt{3}{309.308}{386.533}{\Character{104}}%
\TextAt{21}{316.499}{386.533}{\Character{59}}%
\TextAt{3}{322.748}{386.533}{\Character{105}}%
\TextAt{22}{326.042}{385.194}{M}%
\NativeDiagram{%
\path[draw=black,line width=0.29bp] (260.985,-387.183) -- (265.981,-387.183);%
}%
\TextAt{22}{332.985}{383.228}{\Character{27}}%
\TextAt{12}{79.226}{353.923}{Definition}%
\TextAt{12}{142.837}{353.923}{(4.2).}%
\TextAt{8}{178.868}{353.923}{A}%
\TextAt{12}{192.254}{353.923}{metrized}%
\TextAt{14}{248.967}{353.923}{o}%
\TextAt{12}{254.812}{353.923}{-module}%
\TextAt{8}{307.041}{353.923}{is}%
\TextAt{8}{319.52}{353.923}{a}%
\TextAt{8}{329.991}{353.923}{finitely}%
\TextAt{8}{370.05}{353.923}{generated}%
\TextAt{14}{424.097}{353.923}{o}%
\TextAt{8}{429.941}{353.923}{-module}%
\TextAt{13}{475.847}{353.923}{M}%
\TextAt{8}{493.033}{353.923}{with}%
\TextAt{8}{520.408}{353.923}{an}%
\TextAt{13}{79.226}{336.877}{F}%
\TextAt{8}{88.43}{336.877}{-invariant}%
\TextAt{8}{141.451}{336.877}{hermitian}%
\TextAt{8}{194.783}{336.877}{metric}%
\TextAt{8}{231.198}{336.877}{on}%
\TextAt{13}{247.457}{336.877}{M}%
\TextAt{25}{258.782}{335.084}{C}%
\TextAt{8}{265.037}{336.877}{.}%
\TextAt{12}{79.226}{294.939}{Example}%
\TextAt{12}{134.884}{294.939}{1.}%
\TextAt{8}{157.312}{294.939}{Every}%
\TextAt{8}{191.095}{294.939}{fractional}%
\TextAt{8}{243.589}{294.939}{ideal}%
\TextAt{14}{271.683}{294.939}{a}%
\NativeMathAt{281.208}{294.939}{11.955}{\subseteq}%
\TextAt{13}{294.049}{294.939}{K}%
\TextAt{8}{308.894}{294.939}{of}%
\TextAt{14}{322.355}{294.939}{o}%
\TextAt{8}{328.199}{294.939}{,}%
\TextAt{8}{335.511}{294.939}{in}%
\TextAt{8}{349.294}{294.939}{particular}%
\TextAt{14}{403.413}{294.939}{o}%
\TextAt{8}{413.285}{294.939}{itself,}%
\TextAt{8}{445.045}{294.939}{may}%
\TextAt{8}{470.544}{294.939}{be}%
\TextAt{8}{486.599}{294.939}{equipped}%
\TextAt{8}{79.226}{277.893}{with}%
\TextAt{8}{105.886}{277.893}{the}%
\TextAt{12}{126.043}{277.893}{trivial}%
\TextAt{12}{166.128}{277.893}{metric}%
\TextAt{16}{241.827}{260.847}{\Character{104}}%
\TextAt{13}{246.476}{260.847}{x\Character{59}}%
\TextAt{13}{258.348}{260.847}{y}%
\TextAt{16}{264.509}{260.847}{\Character{105}}%
\TextAt{9}{272.479}{260.847}{\Character{61}}%
\TextAt{13}{284.905}{260.847}{x}%
\Vector{q 0 G 0 g 0.478 w 219.557 -506.332 m 225.694 -506.332 l S Q}%
\TextAt{13}{291.557}{260.847}{y}%
\TextAt{9}{301.014}{260.847}{\Character{61}}%
\TextAt{20}{319.757}{272.205}{\Character{77}}%
\TextAt{21}{313.44}{246.302}{\Character{27}}%
\NativeMathAt{318.565}{246.302}{7.97}{\in}%
\TextAt{21}{324.211}{246.302}{X}%
\TextAt{17}{331.798}{246.302}{\Character{40}}%
\TextAt{25}{335.091}{246.302}{C}%
\TextAt{17}{340.847}{246.302}{\Character{41}}%
\TextAt{13}{346.133}{260.847}{x}%
\TextAt{21}{352.785}{259.054}{\Character{27}}%
\Vector{q 0 G 0 g 0.478 w 286.408 -506.332 m 292.545 -506.332 l S Q}%
\TextAt{13}{358.408}{260.847}{y}%
\TextAt{21}{364.545}{258.025}{\Character{27}}%
\TextAt{8}{79.226}{227.917}{on}%
\TextAt{14}{95.485}{227.917}{a}%
\TextAt{16}{104.125}{227.917}{\Character{10}}%
\TextAt{25}{113.423}{226.124}{Z}%
\TextAt{18}{121.892}{227.917}{C}%
\TextAt{9}{133.847}{227.917}{\Character{61}}%
\TextAt{13}{146.272}{227.917}{K}%
\TextAt{16}{159.747}{227.917}{\Character{10}}%
\TextAt{25}{169.045}{226.124}{Q}%
\TextAt{18}{178.399}{227.917}{C}%
\TextAt{9}{190.354}{227.917}{\Character{61}}%
\TextAt{13}{202.78}{227.917}{K}%
\TextAt{25}{212.739}{226.124}{C}%
\TextAt{8}{218.993}{227.917}{.}%
\TextAt{8}{227.441}{227.917}{All}%
\TextAt{8}{246.617}{227.917}{the}%
\TextAt{13}{266.774}{227.917}{F}%
\TextAt{8}{275.977}{227.917}{-invariant}%
\TextAt{8}{328.999}{227.917}{hermitian}%
\TextAt{8}{382.331}{227.917}{metrics}%
\TextAt{8}{423.361}{227.917}{on}%
\TextAt{14}{439.62}{227.917}{a}%
\TextAt{8}{449.5}{227.917}{are}%
\TextAt{8}{469.011}{227.917}{obtained}%
\TextAt{8}{517.143}{227.917}{as}%
\TextAt{13}{227.143}{200.809}{\Character{11}}%
\TextAt{16}{234.664}{200.809}{\Character{104}}%
\TextAt{13}{239.314}{200.809}{x\Character{59}}%
\TextAt{13}{251.186}{200.809}{y}%
\TextAt{16}{257.347}{200.809}{\Character{105}}%
\TextAt{9}{265.317}{200.809}{\Character{61}}%
\TextAt{13}{277.742}{200.809}{\Character{11}x}%
\Vector{q 0 G 0 g 0.478 w 219.916 -566.37 m 226.053 -566.37 l S Q}%
\TextAt{13}{291.916}{200.809}{y}%
\TextAt{9}{301.373}{200.809}{\Character{61}}%
\TextAt{20}{313.799}{212.166}{\Character{77}}%
\TextAt{21}{320.269}{187.06}{\Character{27}}%
\TextAt{13}{333.857}{200.809}{\Character{11}}%
\TextAt{9}{341.379}{200.809}{\Character{40}}%
\TextAt{13}{345.931}{200.809}{\Character{27}}%
\TextAt{9}{353.013}{200.809}{\Character{41}}%
\TextAt{13}{357.566}{200.809}{x}%
\TextAt{21}{364.218}{199.015}{\Character{27}}%
\Vector{q 0 G 0 g 0.478 w 297.841 -566.37 m 303.978 -566.37 l S Q}%
\TextAt{13}{369.841}{200.809}{y}%
\TextAt{21}{375.978}{197.986}{\Character{27}}%
\TextAt{13}{381.601}{200.809}{\Character{59}}%
\TextAt{8}{79.226}{165.84}{where}%
\TextAt{13}{113.035}{165.84}{\Character{11}}%
\NativeMathAt{123.878}{165.84}{11.955}{\in}%
\TextAt{13}{135.169}{165.84}{K}%
\TextAt{25}{145.128}{164.046}{C}%
\TextAt{8}{155.28}{165.84}{varies}%
\TextAt{8}{188.192}{165.84}{over}%
\TextAt{8}{213.238}{165.84}{the}%
\TextAt{8}{233.395}{165.84}{functions}%
\TextAt{13}{283.858}{165.84}{\Character{11}}%
\TextAt{9}{294.7}{165.84}{\Character{58}}%
\TextAt{13}{301.273}{165.84}{X}%
\TextAt{9}{311.928}{165.84}{\Character{40}}%
\TextAt{18}{316.48}{165.84}{C}%
\TextAt{9}{325.114}{165.84}{\Character{41}}%
\NativeMathAt{332.987}{165.84}{11.955}{\rightarrow}%
\TextAt{18}{348.264}{165.84}{R}%
\TextAt{3}{356.898}{170.178}{\Character{3}}%
\TextAt{17}{356.898}{162.884}{\Character{43}}%
\TextAt{8}{367.88}{165.84}{such}%
\TextAt{8}{394.277}{165.84}{that}%
\TextAt{13}{419.646}{165.84}{\Character{11}}%
\TextAt{9}{427.167}{165.84}{\Character{40}}%
\Vector{q 0 G 0 g 0.478 w 359.72 -601.339 m 366.802 -601.339 l S Q}%
\TextAt{13}{431.72}{165.84}{\Character{27}}%
\TextAt{9}{438.802}{165.84}{\Character{41}}%
\TextAt{9}{446.657}{165.84}{\Character{61}}%
\TextAt{13}{459.101}{165.84}{\Character{11}}%
\TextAt{9}{466.623}{165.84}{\Character{40}}%
\TextAt{13}{471.175}{165.84}{\Character{27}}%
\TextAt{9}{478.257}{165.84}{\Character{41}}%
\TextAt{8}{482.81}{165.84}{.}%
\TextAt{8}{297.211}{43.033}{208}%
\Destination{78.23}{809.815}{nav.page.208}%
\Destination{79.23}{311.985}{nav.sectionz2az.31}%
\Destination{79.36}{693.455}{nav.print.227}%
\Destination{79.36}{367.955}{nav.number.3.4.2}%
\Destination{79.226}{778.124}{errata.85}%
}

% Page 221 (body page 209)
\Page{%
\TextAt{2}{79.226}{796.943}{III.}%
\TextAt{2}{98.317}{796.943}{Riemann–Roch}%
\TextAt{2}{174.375}{796.943}{Theory}%
\Vector{q 0 G 0 g 0.299 w 7.226 18.729 m 460.769 18.729 l S Q}%
\TextAt{12}{79.226}{761.124}{Example}%
\TextAt{12}{133.802}{761.124}{2.}%
\TextAt{8}{156.23}{761.124}{Let}%
\TextAt{13}{176.389}{761.124}{L}%
\TextAt{16}{187.674}{761.124}{\Character{106}}%
\TextAt{13}{194.316}{761.124}{K}%
\TextAt{8}{208.22}{761.124}{be}%
\TextAt{8}{223.334}{761.124}{a}%
\TextAt{8}{232.279}{761.124}{finite}%
\TextAt{8}{261.381}{761.124}{extension}%
\TextAt{8}{312.336}{761.124}{and}%
\TextAt{14}{334.289}{761.124}{A}%
\TextAt{8}{345.955}{761.124}{a}%
\TextAt{8}{354.9}{761.124}{fractional}%
\TextAt{8}{406.453}{761.124}{ideal}%
\TextAt{8}{433.606}{761.124}{of}%
\TextAt{13}{446.126}{761.124}{L}%
\TextAt{8}{454.09}{761.124}{,}%
\TextAt{8}{460.591}{761.124}{which}%
\TextAt{8}{493.267}{761.124}{we}%
\TextAt{8}{509.684}{761.124}{view}%
\TextAt{8}{79.226}{744.078}{as}%
\TextAt{8}{94.096}{744.078}{an}%
\TextAt{14}{110.856}{744.078}{o}%
\TextAt{8}{116.7}{744.078}{-module}%
\TextAt{13}{162.39}{744.078}{M}%
\TextAt{8}{174.964}{744.078}{.}%
\TextAt{8}{184.913}{744.078}{If}%
\TextAt{13}{197.105}{744.078}{Y}%
\TextAt{9}{206.498}{744.078}{\Character{40}}%
\TextAt{18}{211.05}{744.078}{C}%
\TextAt{9}{219.685}{744.078}{\Character{41}}%
\TextAt{9}{228.4}{744.078}{\Character{61}}%
\TextAt{8}{241.688}{744.078}{Hom}%
\TextAt{9}{266.076}{744.078}{\Character{40}}%
\TextAt{13}{270.629}{744.078}{L\Character{59}}%
\TextAt{18}{283.837}{744.078}{C}%
\TextAt{9}{292.472}{744.078}{\Character{41}}%
\TextAt{8}{297.024}{744.078}{,}%
\TextAt{8}{304.798}{744.078}{we}%
\TextAt{8}{322.526}{744.078}{have}%
\TextAt{8}{350.021}{744.078}{the}%
\TextAt{8}{370.677}{744.078}{restriction}%
\TextAt{8}{427.176}{744.078}{map}%
\TextAt{13}{453.69}{744.078}{Y}%
\TextAt{9}{463.083}{744.078}{\Character{40}}%
\TextAt{18}{467.636}{744.078}{C}%
\TextAt{9}{476.27}{744.078}{\Character{41}}%
\NativeMathAt{484.995}{744.078}{11.955}{\rightarrow}%
\TextAt{13}{501.124}{744.078}{X}%
\TextAt{9}{511.779}{744.078}{\Character{40}}%
\TextAt{18}{516.331}{744.078}{C}%
\TextAt{9}{524.965}{744.078}{\Character{41}}%
\TextAt{8}{529.518}{744.078}{,}%
\TextAt{13}{79.226}{727.032}{\Character{28}}%
\NativeArrowAt{88.966}{727.032}{11.955}{11.955}{\mapsto}%
\TextAt{13}{104.242}{727.032}{\Character{28}}%
\TextAt{16}{110.661}{727.032}{\Character{106}}%
\TextAt{21}{113.982}{725.238}{K}%
\TextAt{8}{122.184}{727.032}{,}%
\TextAt{8}{129.155}{727.032}{and}%
\TextAt{8}{151.694}{727.032}{we}%
\TextAt{8}{168.698}{727.032}{write}%
\TextAt{13}{198.386}{727.032}{\Character{28}}%
\TextAt{16}{208.126}{727.032}{\Character{106}}%
\TextAt{13}{214.768}{727.032}{\Character{27}}%
\TextAt{8}{225.524}{727.032}{if}%
\TextAt{13}{236.025}{727.032}{\Character{27}}%
\TextAt{9}{246.428}{727.032}{\Character{61}}%
\TextAt{13}{258.853}{727.032}{\Character{28}}%
\TextAt{16}{265.272}{727.032}{\Character{106}}%
\TextAt{21}{268.593}{725.238}{K}%
\TextAt{8}{276.796}{727.032}{.}%
\TextAt{8}{285.17}{727.032}{For}%
\TextAt{8}{305.916}{727.032}{the}%
\TextAt{8}{325.849}{727.032}{complexification}%
\TextAt{13}{412.456}{727.032}{M}%
\TextAt{25}{423.78}{725.238}{C}%
\TextAt{9}{433.356}{727.032}{\Character{61}}%
\TextAt{14}{445.781}{727.032}{A}%
\TextAt{16}{456.56}{727.032}{\Character{10}}%
\TextAt{25}{465.858}{725.238}{Z}%
\TextAt{18}{473.87}{727.032}{C}%
\TextAt{9}{485.825}{727.032}{\Character{61}}%
\TextAt{13}{498.25}{727.032}{L}%
\TextAt{25}{506.215}{725.238}{C}%
\TextAt{8}{512.469}{727.032}{,}%
\TextAt{8}{519.439}{727.032}{we}%
\TextAt{8}{79.226}{709.986}{obtain}%
\TextAt{8}{115.654}{709.986}{the}%
\TextAt{8}{135.81}{709.986}{decomposition}%
\TextAt{13}{234.974}{692.94}{M}%
\TextAt{25}{246.299}{691.147}{C}%
\TextAt{9}{255.874}{692.94}{\Character{61}}%
\TextAt{20}{273.995}{704.297}{\Character{77}}%
\TextAt{21}{268.299}{678.394}{\Character{28}}%
\NativeMathAt{272.978}{678.394}{7.97}{\in}%
\TextAt{21}{278.623}{678.394}{Y}%
\TextAt{17}{285.413}{678.394}{\Character{40}}%
\TextAt{25}{288.706}{678.394}{C}%
\TextAt{17}{294.463}{678.394}{\Character{41}}%
\TextAt{18}{299.748}{692.94}{C}%
\TextAt{9}{311.703}{692.94}{\Character{61}}%
\TextAt{20}{330.446}{704.297}{\Character{77}}%
\TextAt{21}{324.129}{678.394}{\Character{27}}%
\NativeMathAt{329.254}{678.394}{7.97}{\in}%
\TextAt{21}{334.9}{678.394}{X}%
\TextAt{17}{342.487}{678.394}{\Character{40}}%
\TextAt{25}{345.78}{678.394}{C}%
\TextAt{17}{351.536}{678.394}{\Character{41}}%
\TextAt{13}{356.822}{692.94}{M}%
\TextAt{21}{368.147}{691.147}{\Character{27}}%
\TextAt{13}{373.77}{692.94}{\Character{59}}%
\TextAt{8}{79.226}{658.768}{where}%
\TextAt{13}{114.372}{658.768}{M}%
\TextAt{21}{125.697}{656.975}{\Character{27}}%
\TextAt{9}{136.919}{658.768}{\Character{61}}%
\TextAt{20}{151.623}{667.734}{\Character{76}}%
\TextAt{21}{164.906}{655.281}{\Character{28}}%
\TextAt{3}{169.585}{655.281}{\Character{106}}%
\TextAt{21}{171.937}{655.281}{\Character{27}}%
\TextAt{18}{179.553}{658.768}{C}%
\TextAt{8}{188.187}{658.768}{.}%
\TextAt{13}{200.646}{658.768}{M}%
\TextAt{8}{218.454}{658.768}{is}%
\TextAt{8}{231.554}{658.768}{turned}%
\TextAt{8}{270.61}{658.768}{into}%
\TextAt{8}{295.69}{658.768}{a}%
\TextAt{8}{306.782}{658.768}{metrized}%
\TextAt{14}{356.238}{658.768}{o}%
\TextAt{8}{362.082}{658.768}{-module}%
\TextAt{8}{408.61}{658.768}{by}%
\TextAt{8}{426.205}{658.768}{fixing}%
\TextAt{8}{459.738}{658.768}{the}%
\TextAt{12}{481.231}{658.768}{standard}%
\TextAt{12}{79.226}{641.722}{metrics}%
\TextAt{16}{256.444}{624.676}{\Character{104}}%
\TextAt{13}{261.093}{624.676}{x\Character{59}}%
\TextAt{13}{272.964}{624.676}{y}%
\TextAt{16}{279.126}{624.676}{\Character{105}}%
\TextAt{21}{283.776}{622.883}{M}%
\TextAt{22}{291.908}{621.886}{\Character{27}}%
\TextAt{9}{300.723}{624.676}{\Character{61}}%
\TextAt{20}{313.149}{636.033}{\Character{88}}%
\TextAt{21}{315.705}{610.13}{\Character{28}}%
\TextAt{3}{320.383}{610.13}{\Character{106}}%
\TextAt{21}{322.736}{610.13}{\Character{27}}%
\TextAt{13}{332.41}{624.676}{x}%
\TextAt{21}{339.062}{622.883}{\Character{28}}%
\Vector{q 0 G 0 g 0.478 w 272.238 -142.503 m 278.375 -142.503 l S Q}%
\TextAt{13}{344.238}{624.676}{y}%
\TextAt{21}{350.375}{621.853}{\Character{28}}%
\TextAt{8}{79.226}{590.504}{on}%
\TextAt{8}{95.485}{590.504}{the}%
\TextAt{9}{115.642}{590.504}{\Character{91}}%
\TextAt{13}{118.893}{590.504}{L}%
\TextAt{9}{130.179}{590.504}{\Character{58}}%
\TextAt{13}{136.751}{590.504}{K}%
\TextAt{9}{147.569}{590.504}{\Character{93}}%
\TextAt{8}{150.82}{590.504}{-dimensional}%
\TextAt{18}{219.168}{590.504}{C}%
\TextAt{8}{227.802}{590.504}{-vector}%
\TextAt{8}{266.824}{590.504}{spaces}%
\TextAt{13}{302.714}{590.504}{M}%
\TextAt{21}{314.038}{588.711}{\Character{27}}%
\TextAt{8}{319.662}{590.504}{.}%
\TextAt{8}{97.159}{570.469}{If}%
\TextAt{13}{108.8}{570.469}{M}%
\TextAt{8}{125.22}{570.469}{and}%
\TextAt{13}{147.932}{570.469}{M}%
\TextAt{3}{160.506}{574.808}{\Character{48}}%
\TextAt{8}{167.147}{570.469}{are}%
\TextAt{8}{186.607}{570.469}{metrized}%
\TextAt{14}{234.676}{570.469}{o}%
\TextAt{8}{240.52}{570.469}{-modules,}%
\TextAt{8}{293.536}{570.469}{then}%
\TextAt{8}{320.145}{570.469}{so}%
\TextAt{8}{334.465}{570.469}{is}%
\TextAt{8}{346.178}{570.469}{their}%
\TextAt{8}{374.09}{570.469}{direct}%
\TextAt{8}{407.203}{570.469}{sum}%
\TextAt{13}{431.923}{570.469}{M}%
\TextAt{16}{447.049}{570.469}{\Character{8}}%
\TextAt{13}{458.9}{570.469}{M}%
\TextAt{3}{471.474}{574.808}{\Character{48}}%
\TextAt{8}{474.269}{570.469}{,}%
\TextAt{8}{481.377}{570.469}{the}%
\TextAt{8}{501.483}{570.469}{tensor}%
\TextAt{8}{79.226}{553.423}{product}%
\TextAt{13}{123.048}{553.423}{M}%
\TextAt{16}{138.12}{553.423}{\Character{10}}%
\TextAt{13}{149.917}{553.423}{M}%
\TextAt{3}{162.49}{557.762}{\Character{48}}%
\TextAt{8}{165.285}{553.423}{,}%
\TextAt{8}{172.373}{553.423}{the}%
\TextAt{8}{192.452}{553.423}{dual}%
\TextAt{9}{222.725}{556.445}{\Character{20}}%
\TextAt{13}{218.389}{553.423}{M}%
\TextAt{9}{234.283}{553.423}{\Character{61}}%
\TextAt{8}{246.709}{553.423}{Hom}%
\TextAt{27}{271.097}{551.63}{o}%
\TextAt{9}{275.491}{553.423}{\Character{40}}%
\TextAt{13}{280.044}{553.423}{M}%
\TextAt{13}{291.963}{553.423}{\Character{59}}%
\TextAt{14}{297.211}{553.423}{o}%
\TextAt{9}{303.055}{553.423}{\Character{41}}%
\TextAt{8}{311.428}{553.423}{and}%
\TextAt{8}{334.113}{553.423}{the}%
\TextAt{13}{354.192}{553.423}{n}%
\TextAt{8}{361.18}{553.423}{-th}%
\TextAt{8}{379.955}{553.423}{exterior}%
\TextAt{8}{423.132}{553.423}{power}%
\TextAt{20}{457.198}{562.39}{\Character{86}}%
\TextAt{21}{467.161}{559.235}{n}%
\TextAt{13}{474.79}{553.423}{M}%
\TextAt{8}{487.363}{553.423}{.}%
\TextAt{8}{495.786}{553.423}{In}%
\TextAt{8}{510.33}{553.423}{fact,}%
\TextAt{8}{79.226}{536.377}{we}%
\TextAt{8}{96.454}{536.377}{have}%
\TextAt{8}{123.448}{536.377}{that}%
\TextAt{9}{128.599}{507.376}{\Character{40}}%
\TextAt{13}{133.151}{507.376}{M}%
\TextAt{16}{148.381}{507.376}{\Character{8}}%
\TextAt{13}{160.336}{507.376}{M}%
\TextAt{3}{172.91}{512.312}{\Character{48}}%
\TextAt{9}{175.705}{507.376}{\Character{41}}%
\TextAt{25}{180.257}{505.583}{C}%
\TextAt{9}{189.833}{507.376}{\Character{61}}%
\TextAt{13}{202.258}{507.376}{M}%
\TextAt{25}{213.583}{505.583}{C}%
\TextAt{16}{222.494}{507.376}{\Character{8}}%
\TextAt{13}{234.449}{507.376}{M}%
\TextAt{3}{247.023}{512.312}{\Character{48}}%
\TextAt{25}{245.774}{504.421}{C}%
\TextAt{13}{252.028}{507.376}{\Character{59}}%
\TextAt{9}{339.867}{507.376}{\Character{40}}%
\TextAt{13}{344.419}{507.376}{M}%
\TextAt{16}{359.649}{507.376}{\Character{10}}%
\TextAt{27}{368.948}{505.583}{o}%
\TextAt{13}{375.999}{507.376}{M}%
\TextAt{3}{388.572}{512.312}{\Character{48}}%
\TextAt{9}{391.368}{507.376}{\Character{41}}%
\TextAt{25}{395.92}{505.583}{C}%
\TextAt{9}{405.495}{507.376}{\Character{61}}%
\TextAt{13}{417.92}{507.376}{M}%
\TextAt{25}{429.245}{505.583}{C}%
\TextAt{16}{438.156}{507.376}{\Character{10}}%
\TextAt{21}{447.455}{505.583}{K}%
\TextAt{29}{454.596}{504.21}{C}%
\TextAt{13}{462.566}{507.376}{M}%
\TextAt{3}{475.14}{512.312}{\Character{48}}%
\TextAt{25}{473.891}{504.421}{C}%
\TextAt{13}{480.145}{507.376}{\Character{59}}%
\TextAt{9}{173.268}{490.363}{\Character{20}}%
\TextAt{13}{168.933}{487.341}{M}%
\TextAt{25}{180.257}{485.548}{C}%
\TextAt{9}{189.832}{487.341}{\Character{61}}%
\TextAt{8}{202.258}{487.341}{Hom}%
\TextAt{21}{226.647}{485.548}{K}%
\TextAt{29}{233.788}{484.175}{C}%
\TextAt{9}{239.101}{487.341}{\Character{40}}%
\TextAt{13}{243.654}{487.341}{M}%
\TextAt{25}{254.978}{485.548}{C}%
\TextAt{13}{261.233}{487.341}{\Character{59}}%
\TextAt{13}{266.457}{487.341}{K}%
\TextAt{25}{276.436}{485.548}{C}%
\TextAt{9}{282.69}{487.341}{\Character{41}}%
\TextAt{13}{287.243}{487.341}{\Character{59}}%
\TextAt{9}{356.65}{487.341}{\Character{40}}%
\TextAt{20}{361.202}{496.308}{\Character{86}}%
\TextAt{21}{371.165}{493.153}{n}%
\TextAt{13}{378.794}{487.341}{M}%
\TextAt{9}{391.368}{487.341}{\Character{41}}%
\TextAt{25}{395.92}{485.548}{C}%
\TextAt{9}{405.495}{487.341}{\Character{61}}%
\TextAt{20}{417.92}{496.308}{\Character{86}}%
\TextAt{21}{427.883}{493.153}{n}%
\TextAt{21}{427.883}{483.854}{K}%
\TextAt{29}{435.025}{482.482}{C}%
\TextAt{13}{442.33}{487.341}{M}%
\TextAt{25}{453.655}{485.548}{C}%
\TextAt{13}{459.91}{487.341}{\Character{59}}%
\TextAt{8}{79.226}{458.34}{and}%
\TextAt{8}{101.989}{458.34}{the}%
\TextAt{8}{122.145}{458.34}{metrics}%
\TextAt{8}{163.175}{458.34}{on}%
\TextAt{8}{179.434}{458.34}{these}%
\TextAt{13}{209.406}{458.34}{K}%
\TextAt{25}{219.365}{456.547}{C}%
\TextAt{8}{225.619}{458.34}{-modules}%
\TextAt{8}{275.425}{458.34}{are}%
\TextAt{8}{294.936}{458.34}{given}%
\TextAt{8}{325.505}{458.34}{by}%
\TextAt{16}{190.645}{429.339}{\Character{104}}%
\TextAt{13}{195.295}{429.339}{x}%
\TextAt{16}{204.603}{429.339}{\Character{8}}%
\TextAt{13}{216.559}{429.339}{x}%
\TextAt{3}{223.211}{434.275}{\Character{48}}%
\TextAt{13}{226.006}{429.339}{\Character{59}}%
\TextAt{13}{231.23}{429.339}{y}%
\TextAt{16}{240.043}{429.339}{\Character{8}}%
\TextAt{13}{251.998}{429.339}{y}%
\TextAt{3}{258.135}{434.275}{\Character{48}}%
\TextAt{16}{260.93}{429.339}{\Character{105}}%
\TextAt{21}{265.579}{427.546}{M}%
\TextAt{3}{274.561}{427.546}{\Character{8}}%
\TextAt{21}{281.147}{427.546}{M}%
\TextAt{23}{290.129}{429.815}{\Character{48}}%
\TextAt{9}{296.632}{429.339}{\Character{61}}%
\TextAt{16}{309.058}{429.339}{\Character{104}}%
\TextAt{13}{313.707}{429.339}{x\Character{59}}%
\TextAt{13}{325.579}{429.339}{y}%
\TextAt{16}{331.74}{429.339}{\Character{105}}%
\TextAt{21}{336.389}{427.546}{M}%
\TextAt{9}{348.526}{429.339}{\Character{43}}%
\TextAt{16}{360.287}{429.339}{\Character{104}}%
\TextAt{13}{364.936}{429.339}{x}%
\TextAt{3}{371.588}{434.275}{\Character{48}}%
\TextAt{13}{374.384}{429.339}{\Character{59}}%
\TextAt{13}{379.608}{429.339}{y}%
\TextAt{3}{385.764}{434.275}{\Character{48}}%
\TextAt{16}{388.559}{429.339}{\Character{105}}%
\TextAt{21}{393.209}{427.546}{M}%
\TextAt{23}{402.19}{429.815}{\Character{48}}%
\TextAt{13}{405.373}{429.339}{\Character{59}}%
\TextAt{8}{458.634}{429.339}{resp.}%
\TextAt{16}{190.645}{409.304}{\Character{104}}%
\TextAt{13}{195.295}{409.304}{x}%
\TextAt{16}{204.603}{409.304}{\Character{10}}%
\TextAt{13}{216.559}{409.304}{x}%
\TextAt{3}{223.211}{414.24}{\Character{48}}%
\TextAt{13}{226.006}{409.304}{\Character{59}}%
\TextAt{13}{231.23}{409.304}{y}%
\TextAt{16}{240.043}{409.304}{\Character{10}}%
\TextAt{13}{251.998}{409.304}{y}%
\TextAt{3}{258.135}{414.24}{\Character{48}}%
\TextAt{16}{260.93}{409.304}{\Character{105}}%
\TextAt{21}{265.579}{407.511}{M}%
\TextAt{3}{274.561}{407.511}{\Character{10}}%
\TextAt{21}{281.147}{407.511}{M}%
\TextAt{23}{290.129}{409.78}{\Character{48}}%
\TextAt{9}{296.632}{409.304}{\Character{61}}%
\TextAt{16}{309.058}{409.304}{\Character{104}}%
\TextAt{13}{313.707}{409.304}{x\Character{59}}%
\TextAt{13}{325.579}{409.304}{y}%
\TextAt{16}{331.74}{409.304}{\Character{105}}%
\TextAt{21}{336.389}{407.511}{M}%
\TextAt{16}{348.526}{409.304}{\Character{1}}%
\TextAt{16}{354.492}{409.304}{\Character{104}}%
\TextAt{13}{359.153}{409.304}{x}%
\TextAt{3}{365.805}{414.24}{\Character{48}}%
\TextAt{13}{368.6}{409.304}{\Character{59}}%
\TextAt{13}{373.824}{409.304}{y}%
\TextAt{3}{379.981}{414.24}{\Character{48}}%
\TextAt{16}{382.776}{409.304}{\Character{105}}%
\TextAt{21}{387.425}{407.511}{M}%
\TextAt{23}{396.407}{409.78}{\Character{48}}%
\TextAt{13}{399.589}{409.304}{\Character{59}}%
\TextAt{8}{458.634}{409.304}{resp.}%
\TextAt{16}{256.5}{389.269}{\Character{104}}%
\TextAt{9}{261.874}{389.27}{\Character{20}}%
\TextAt{13}{261.15}{389.269}{x\Character{59}}%
\TextAt{9}{273.838}{389.27}{\Character{20}}%
\TextAt{13}{273.046}{389.269}{y}%
\TextAt{16}{279.183}{389.27}{\Character{105}}%
\TextAt{17}{286.911}{388.378}{\Character{20}}%
\TextAt{21}{283.832}{386.364}{M}%
\TextAt{9}{296.632}{389.269}{\Character{61}}%
\Vector{q 0 G 0 g 0.478 w 237.058 -374.09 m 273.869 -374.09 l S Q}%
\TextAt{16}{309.058}{389.269}{\Character{104}}%
\TextAt{13}{313.707}{389.27}{x\Character{59}}%
\TextAt{13}{325.579}{389.27}{y}%
\TextAt{16}{331.74}{389.269}{\Character{105}}%
\TextAt{21}{336.389}{387.476}{M}%
\TextAt{13}{345.869}{389.27}{\Character{59}}%
\TextAt{8}{458.634}{389.269}{resp.}%
\TextAt{16}{129.236}{369.235}{\Character{104}}%
\TextAt{13}{133.885}{369.235}{x}%
\TextAt{17}{140.537}{367.441}{1}%
\TextAt{16}{147.926}{369.235}{\Character{94}}%
\TextAt{16}{158.542}{369.235}{\Character{1}}%
\TextAt{16}{163.838}{369.235}{\Character{1}}%
\TextAt{16}{169.147}{369.235}{\Character{1}}%
\TextAt{16}{175.112}{369.235}{\Character{94}}%
\TextAt{13}{185.784}{369.235}{x}%
\TextAt{21}{192.436}{367.441}{n}%
\TextAt{13}{198.073}{369.235}{\Character{59}}%
\TextAt{13}{203.297}{369.235}{y}%
\TextAt{17}{209.024}{367.441}{1}%
\TextAt{16}{216.413}{369.235}{\Character{94}}%
\TextAt{16}{227.029}{369.235}{\Character{1}}%
\TextAt{16}{232.325}{369.235}{\Character{1}}%
\TextAt{16}{237.634}{369.235}{\Character{1}}%
\TextAt{16}{243.599}{369.235}{\Character{94}}%
\TextAt{13}{254.271}{369.235}{y}%
\TextAt{21}{259.979}{367.441}{n}%
\TextAt{16}{265.615}{369.235}{\Character{105}}%
\TextAt{30}{270.265}{373.242}{\Character{86}}%
\TextAt{22}{277.322}{370.779}{n}%
\TextAt{21}{283.832}{367.264}{M}%
\TextAt{9}{296.632}{369.235}{\Character{61}}%
\TextAt{8}{309.058}{369.235}{det}%
\TextAt{9}{325.317}{369.235}{\Character{40}}%
\TextAt{16}{329.869}{369.235}{\Character{104}}%
\TextAt{13}{334.519}{369.235}{x}%
\TextAt{21}{341.171}{367.441}{i}%
\TextAt{13}{344.552}{369.235}{\Character{59}}%
\TextAt{13}{349.776}{369.235}{y}%
\TextAt{21}{355.504}{367.441}{j}%
\TextAt{16}{359.886}{369.235}{\Character{105}}%
\TextAt{21}{364.535}{367.441}{M}%
\TextAt{9}{374.015}{369.235}{\Character{41}}%
\TextAt{13}{378.567}{369.235}{\Character{58}}%
\TextAt{8}{79.226}{340.234}{Here}%
\TextAt{9}{106.497}{340.234}{\Character{20}}%
\TextAt{13}{105.772}{340.234}{x}%
\TextAt{8}{112.424}{340.234}{,}%
\TextAt{8}{118.707}{340.234}{in}%
\TextAt{8}{131.278}{340.234}{the}%
\TextAt{8}{150.352}{340.234}{case}%
\TextAt{8}{174.04}{340.234}{of}%
\TextAt{8}{186.288}{340.234}{the}%
\TextAt{8}{205.362}{340.234}{module}%
\TextAt{9}{249.908}{343.256}{\Character{20}}%
\TextAt{13}{245.572}{340.234}{M}%
\TextAt{25}{256.897}{338.44}{C}%
\TextAt{8}{263.152}{340.234}{,}%
\TextAt{8}{269.435}{340.234}{denotes}%
\TextAt{8}{310.686}{340.234}{the}%
\TextAt{8}{329.76}{340.234}{homomorphism}%
\TextAt{9}{412.072}{340.234}{\Character{20}}%
\TextAt{13}{411.347}{340.234}{x}%
\TextAt{9}{421.32}{340.234}{\Character{61}}%
\TextAt{16}{433.745}{340.234}{\Character{104}}%
\TextAt{13}{441.21}{340.234}{\Character{59}}%
\TextAt{13}{446.434}{340.234}{x}%
\TextAt{16}{453.106}{340.234}{\Character{105}}%
\TextAt{21}{457.755}{338.44}{M}%
\TextAt{9}{470.556}{340.234}{\Character{58}}%
\TextAt{13}{477.128}{340.234}{M}%
\TextAt{25}{488.453}{338.44}{C}%
\NativeMathAt{498.028}{340.234}{11.955}{\rightarrow}%
\TextAt{13}{513.304}{340.234}{K}%
\TextAt{25}{523.263}{338.44}{C}%
\TextAt{8}{529.518}{340.234}{.}%
\TextAt{8}{97.159}{320.199}{Among}%
\TextAt{8}{137.458}{320.199}{all}%
\TextAt{14}{153.369}{320.199}{o}%
\TextAt{8}{159.214}{320.199}{-modules}%
\TextAt{13}{208.671}{320.199}{M}%
\TextAt{8}{224.794}{320.199}{the}%
\TextAt{12}{244.602}{320.199}{projective}%
\TextAt{8}{306.289}{320.199}{ones}%
\TextAt{8}{332.015}{320.199}{play}%
\TextAt{8}{357.036}{320.199}{a}%
\TextAt{8}{366.443}{320.199}{special}%
\TextAt{8}{404.196}{320.199}{rôle.}%
\TextAt{8}{431.394}{320.199}{They}%
\TextAt{8}{461.28}{320.199}{are}%
\TextAt{8}{480.443}{320.199}{defined}%
\TextAt{8}{520.408}{320.199}{by}%
\ErrataAt{79.19}{303.153}{455.81}{the condition that for every short exact sequence $0\to F^{\prime}\to F\to F^{\prime\prime}\to0$ the sequence}%
\ErrataAt{150}{274.152}{366}{$0\longrightarrow\operatorname{Hom}_{\mathfrak o}(M,F^{\prime})\longrightarrow\operatorname{Hom}_{\mathfrak o}(M,F)\longrightarrow\operatorname{Hom}_{\mathfrak o}(M,F^{\prime\prime})\longrightarrow0$}%
\TextAt{8}{79.226}{245.151}{is}%
\TextAt{8}{91.091}{245.151}{also}%
\TextAt{8}{114.673}{245.151}{exact.}%
\TextAt{8}{150.42}{245.151}{This}%
\TextAt{8}{177.241}{245.151}{is}%
\TextAt{8}{189.107}{245.151}{equivalent}%
\TextAt{8}{244.824}{245.151}{to}%
\TextAt{8}{259.235}{245.151}{any}%
\TextAt{8}{281.454}{245.151}{of}%
\TextAt{8}{294.885}{245.151}{the}%
\TextAt{8}{315.143}{245.151}{following}%
\TextAt{8}{364.679}{245.151}{conditions}%
\TextAt{8}{420.779}{245.151}{(the}%
\TextAt{8}{445.591}{245.151}{last}%
\TextAt{8}{467.87}{245.151}{two,}%
\TextAt{8}{493.365}{245.151}{because}%
\TextAt{14}{79.226}{228.105}{o}%
\TextAt{8}{89.99}{228.105}{is}%
\TextAt{8}{102.777}{228.105}{a}%
\TextAt{8}{113.555}{228.105}{Dedekind}%
\TextAt{8}{166.762}{228.105}{domain).}%
\TextAt{8}{220.563}{228.105}{For}%
\TextAt{8}{242.555}{228.105}{the}%
\TextAt{8}{263.734}{228.105}{proof,}%
\TextAt{8}{298.834}{228.105}{we}%
\TextAt{8}{317.084}{228.105}{refer}%
\TextAt{8}{345.089}{228.105}{the}%
\TextAt{8}{366.268}{228.105}{reader}%
\TextAt{8}{403.061}{228.105}{to}%
\TextAt{8}{418.394}{228.105}{standard}%
\TextAt{8}{468.265}{228.105}{textbooks}%
\TextAt{8}{523.337}{228.105}{of}%
\TextAt{8}{79.226}{211.059}{commutative}%
\TextAt{8}{149.152}{211.059}{algebra}%
\TextAt{8}{190.134}{211.059}{(see}%
\TextAt{8}{213.602}{211.059}{for}%
\TextAt{8}{231.487}{211.059}{instance}%
\TextAt{8}{277.072}{211.059}{[90],}%
\TextAt{8}{302.441}{211.059}{chap.}%
\TextAt{8}{334.632}{211.059}{IV,}%
\TextAt{8}{354.777}{211.059}{§3,}%
\TextAt{8}{373.463}{211.059}{or}%
\TextAt{8}{387.773}{211.059}{[16],}%
\TextAt{8}{413.142}{211.059}{chap.}%
\TextAt{8}{445.333}{211.059}{7,}%
\TextAt{8}{458.341}{211.059}{§4).}%
\TextAt{12}{79.226}{185.046}{Proposition}%
\TextAt{12}{152.045}{185.046}{(4.3).}%
\TextAt{8}{187.455}{185.046}{For}%
\TextAt{8}{208.518}{185.046}{any}%
\TextAt{8}{230.729}{185.046}{finitely}%
\TextAt{8}{270.167}{185.046}{generated}%
\TextAt{14}{323.592}{185.046}{o}%
\TextAt{8}{329.437}{185.046}{-module}%
\TextAt{13}{374.721}{185.046}{M}%
\TextAt{8}{391.286}{185.046}{the}%
\TextAt{8}{411.536}{185.046}{following}%
\TextAt{8}{461.064}{185.046}{conditions}%
\TextAt{8}{517.156}{185.046}{are}%
\TextAt{8}{79.226}{168}{equivalent:}%
\TextAt{8}{89.579}{147.966}{(i)}%
\TextAt{13}{107.918}{147.966}{M}%
\TextAt{8}{124.389}{147.966}{is}%
\TextAt{8}{136.153}{147.966}{projective,}%
\TextAt{8}{86.327}{130.92}{(ii)}%
\TextAt{13}{107.918}{130.92}{M}%
\TextAt{8}{124.389}{130.92}{is}%
\TextAt{8}{136.153}{130.92}{a}%
\TextAt{8}{145.909}{130.92}{direct}%
\TextAt{8}{179.072}{130.92}{summand}%
\TextAt{8}{232.464}{130.92}{of}%
\TextAt{8}{245.794}{130.92}{a}%
\TextAt{8}{255.55}{130.92}{finitely}%
\TextAt{8}{294.894}{130.92}{generated}%
\TextAt{8}{348.226}{130.92}{free}%
\TextAt{14}{370.654}{130.92}{o}%
\TextAt{8}{376.498}{130.92}{-module,}%
\TextAt{8}{83.076}{113.874}{(iii)}%
\TextAt{13}{107.918}{113.874}{M}%
\TextAt{8}{124.389}{113.874}{is}%
\TextAt{8}{136.153}{113.874}{locally}%
\TextAt{8}{173.226}{113.874}{free,}%
\TextAt{8}{198.906}{113.874}{i.e.,}%
\TextAt{13}{221.011}{113.874}{M}%
\TextAt{16}{236.241}{113.874}{\Character{10}}%
\TextAt{27}{245.54}{112.08}{o}%
\TextAt{14}{252.591}{113.874}{o}%
\TextAt{27}{258.435}{112.08}{p}%
\TextAt{8}{266.819}{113.874}{is}%
\TextAt{8}{278.583}{113.874}{a}%
\TextAt{8}{288.339}{113.874}{free}%
\TextAt{14}{310.767}{113.874}{o}%
\TextAt{27}{316.611}{112.08}{p}%
\TextAt{8}{321.098}{113.874}{-module}%
\TextAt{8}{366.288}{113.874}{for}%
\TextAt{8}{384.173}{113.874}{any}%
\TextAt{8}{406.29}{113.874}{prime}%
\TextAt{8}{439.454}{113.874}{ideal}%
\TextAt{14}{467.417}{113.874}{p}%
\TextAt{8}{473.4}{113.874}{,}%
\TextAt{8}{83.398}{96.828}{(iv)}%
\TextAt{13}{107.918}{96.828}{M}%
\TextAt{8}{124.389}{96.828}{is}%
\TextAt{8}{136.153}{96.828}{torsion}%
\TextAt{8}{175.247}{96.828}{free,}%
\TextAt{8}{200.926}{96.828}{i.e.,}%
\TextAt{8}{223.032}{96.828}{the}%
\TextAt{8}{243.188}{96.828}{map}%
\TextAt{13}{269.202}{96.828}{M}%
\NativeMathAt{285.097}{96.828}{11.955}{\rightarrow}%
\TextAt{13}{300.373}{96.828}{M}%
\TextAt{8}{312.947}{96.828}{,}%
\TextAt{13}{320.096}{96.828}{x}%
\NativeArrowAt{330.069}{96.828}{11.955}{11.955}{\mapsto}%
\TextAt{13}{345.345}{96.828}{ax}%
\TextAt{8}{358.142}{96.828}{,}%
\TextAt{8}{365.291}{96.828}{is}%
\TextAt{8}{377.055}{96.828}{injective}%
\TextAt{8}{423.548}{96.828}{for}%
\TextAt{8}{441.433}{96.828}{all}%
\TextAt{8}{457.692}{96.828}{nonzero}%
\TextAt{13}{501.269}{96.828}{a}%
\NativeMathAt{510.735}{96.828}{11.955}{\in}%
\TextAt{14}{522.026}{96.828}{o}%
\TextAt{8}{527.87}{96.828}{,}%
\TextAt{8}{86.65}{79.782}{(v)}%
\TextAt{13}{107.918}{79.782}{M}%
\NativeMathAt{123.813}{79.782}{11.955}{\cong}%
\TextAt{14}{136.432}{79.782}{a}%
\TextAt{16}{145.072}{79.782}{\Character{8}}%
\TextAt{14}{157.027}{79.782}{o}%
\TextAt{21}{162.871}{84.12}{n}%
\TextAt{8}{172.405}{79.782}{for}%
\TextAt{8}{190.29}{79.782}{some}%
\TextAt{8}{219.616}{79.782}{ideal}%
\TextAt{14}{247.579}{79.782}{a}%
\TextAt{8}{257.46}{79.782}{of}%
\TextAt{14}{270.79}{79.782}{o}%
\TextAt{8}{280.531}{79.782}{and}%
\TextAt{8}{303.294}{79.782}{some}%
\TextAt{8}{332.62}{79.782}{integer}%
\TextAt{13}{371.319}{79.782}{n}%
\NativeMathAt{381.627}{79.782}{11.955}{\geq}%
\TextAt{9}{394.247}{79.782}{0}%
\TextAt{8}{400.1}{79.782}{.}%
\TextAt{8}{297.211}{43.033}{209}%
\Destination{79.23}{162.695}{nav.Item.73}%
\Destination{79.23}{144.965}{nav.Item.74}%
\Destination{79.23}{127.915}{nav.Item.75}%
\Destination{79.23}{110.575}{nav.Item.76}%
\Destination{79.23}{93.825}{nav.Item.77}%
\Destination{78.23}{809.815}{nav.page.209}%
\Destination{79.23}{778.165}{nav.sectionz2az.32}%
\Destination{79.36}{673.755}{nav.print.228}%
\Destination{79.36}{199.055}{nav.number.3.4.3}%
\Hotlink{334.132}{208.855}{13.995}{11.047}{nav.chapter.4}%
\Hotlink{359.956}{208.855}{6.858}{11.047}{nav.section.4.3}%
\Hotlink{463.52}{208.855}{6.858}{11.047}{nav.section.4.4}%
\Hotlink{276.572}{208.855}{19.22}{11.047}{nav.cite.N090}%
\Hotlink{387.273}{208.855}{19.22}{11.047}{nav.cite.N016}%
\Destination{79.226}{320.153}{errata.88}%
}

% Page 222 (body page 210)
\Page{%
\TextAt{2}{79.226}{796.943}{III.}%
\TextAt{2}{98.317}{796.943}{Riemann–Roch}%
\TextAt{2}{174.375}{796.943}{Theory}%
\Vector{q 0 G 0 g 0.299 w 7.226 18.729 m 460.769 18.729 l S Q}%
\TextAt{8}{97.159}{761.124}{In}%
\TextAt{8}{112.9}{761.124}{order}%
\TextAt{8}{144.589}{761.124}{to}%
\TextAt{8}{160.018}{761.124}{distinguish}%
\TextAt{8}{220.448}{761.124}{them}%
\TextAt{8}{251.479}{761.124}{from}%
\TextAt{8}{280.239}{761.124}{projective}%
\TextAt{14}{335.336}{761.124}{o}%
\TextAt{8}{341.18}{761.124}{-modules,}%
\TextAt{8}{395.637}{761.124}{we}%
\TextAt{8}{413.984}{761.124}{will}%
\TextAt{8}{437.208}{761.124}{henceforth}%
\TextAt{8}{495.88}{761.124}{call}%
\TextAt{8}{518.459}{761.124}{ar-}%
\TextAt{8}{79.226}{744.078}{bitrary}%
\TextAt{8}{119.65}{744.078}{finitely}%
\TextAt{8}{160.062}{744.078}{generated}%
\TextAt{14}{214.461}{744.078}{o}%
\TextAt{8}{220.305}{744.078}{-modules}%
\TextAt{19}{271.178}{744.078}{coherent}%
\TextAt{8}{313.475}{744.078}{.}%
\TextAt{8}{325.126}{744.078}{The}%
\TextAt{12}{350.247}{744.078}{rank}%
\TextAt{8}{381.824}{744.078}{of}%
\TextAt{8}{396.221}{744.078}{a}%
\TextAt{8}{407.044}{744.078}{coherent}%
\TextAt{14}{455.262}{744.078}{o}%
\TextAt{8}{461.107}{744.078}{-module}%
\TextAt{13}{507.365}{744.078}{M}%
\TextAt{8}{524.903}{744.078}{is}%
\TextAt{8}{79.226}{727.032}{defined}%
\TextAt{8}{119.539}{727.032}{to}%
\TextAt{8}{133.849}{727.032}{be}%
\TextAt{8}{149.774}{727.032}{the}%
\TextAt{8}{169.93}{727.032}{dimension}%
\TextAt{8}{240.684}{698.031}{rk}%
\TextAt{9}{251.42}{698.031}{\Character{40}}%
\TextAt{13}{255.972}{698.031}{M}%
\TextAt{9}{268.546}{698.031}{\Character{41}}%
\TextAt{9}{276.401}{698.031}{\Character{61}}%
\TextAt{8}{288.844}{698.031}{dim}%
\TextAt{21}{308.355}{696.237}{K}%
\TextAt{9}{316.558}{698.031}{\Character{40}}%
\TextAt{13}{321.11}{698.031}{M}%
\TextAt{16}{336.34}{698.031}{\Character{10}}%
\TextAt{27}{345.639}{696.237}{o}%
\TextAt{13}{352.69}{698.031}{K}%
\TextAt{9}{363.508}{698.031}{\Character{41}}%
\TextAt{13}{368.06}{698.031}{\Character{58}}%
\TextAt{8}{79.226}{669.029}{The}%
\TextAt{8}{103.258}{669.029}{projective}%
\TextAt{14}{157.214}{669.029}{o}%
\TextAt{8}{163.058}{669.029}{-modules}%
\TextAt{13}{212.841}{669.029}{L}%
\TextAt{8}{224.681}{669.029}{of}%
\TextAt{8}{237.989}{669.029}{rank}%
\TextAt{9}{264.962}{669.029}{1}%
\TextAt{8}{274.691}{669.029}{are}%
\TextAt{8}{294.18}{669.029}{called}%
\TextAt{12}{327.321}{669.029}{invertible}%
\TextAt{14}{386.729}{669.029}{o}%
\TextAt{8}{392.573}{669.029}{-modules,}%
\TextAt{8}{445.612}{669.029}{because}%
\TextAt{8}{488.892}{669.029}{for}%
\TextAt{8}{506.755}{669.029}{them}%
\TextAt{13}{79.226}{651.983}{L}%
\TextAt{16}{90.877}{651.983}{\Character{10}}%
\TextAt{27}{100.176}{650.19}{o}%
\TextAt{9}{109.638}{655.005}{\Character{20}}%
\TextAt{13}{108.257}{651.983}{L}%
\NativeMathAt{122.117}{651.983}{11.955}{\rightarrow}%
\TextAt{14}{139.968}{651.983}{o}%
\TextAt{8}{145.812}{651.983}{,}%
\TextAt{13}{154.85}{651.983}{a}%
\TextAt{16}{164.682}{651.983}{\Character{10}}%
\TextAt{9}{177.813}{651.983}{\Character{20}}%
\TextAt{13}{177.667}{651.983}{a}%
\NativeArrowAt{189.707}{651.983}{11.955}{11.955}{\mapsto}%
\TextAt{9}{207.704}{651.983}{\Character{20}}%
\TextAt{13}{207.558}{651.983}{a}%
\TextAt{9}{213.703}{651.983}{\Character{40}}%
\TextAt{13}{218.256}{651.983}{a}%
\TextAt{9}{224.401}{651.983}{\Character{41}}%
\TextAt{8}{228.953}{651.983}{,}%
\TextAt{8}{237.991}{651.983}{is}%
\TextAt{8}{251.266}{651.983}{an}%
\TextAt{8}{269.036}{651.983}{isomorphism.}%
\TextAt{8}{346.539}{651.983}{The}%
\TextAt{8}{372.104}{651.983}{invertible}%
\TextAt{14}{425.321}{651.983}{o}%
\TextAt{8}{431.165}{651.983}{-modules}%
\TextAt{8}{482.481}{651.983}{are}%
\TextAt{8}{503.503}{651.983}{either}%
\TextAt{8}{79.226}{634.938}{fractional}%
\TextAt{8}{132.335}{634.938}{ideals}%
\TextAt{8}{165.659}{634.938}{of}%
\TextAt{13}{179.734}{634.938}{K}%
\TextAt{8}{190.552}{634.938}{,}%
\TextAt{8}{198.633}{634.938}{or}%
\TextAt{8}{213.689}{634.938}{isomorphic}%
\TextAt{8}{273.685}{634.938}{to}%
\TextAt{8}{288.741}{634.938}{a}%
\TextAt{8}{299.242}{634.938}{fractional}%
\TextAt{8}{352.351}{634.938}{ideal}%
\TextAt{8}{381.06}{634.938}{as}%
\TextAt{14}{396.176}{634.938}{o}%
\TextAt{8}{402.02}{634.938}{-modules.}%
\TextAt{8}{458.613}{634.938}{Indeed,}%
\TextAt{8}{500.826}{634.938}{if}%
\TextAt{13}{512.295}{634.938}{L}%
\TextAt{8}{524.903}{634.938}{is}%
\TextAt{8}{79.226}{617.892}{projective}%
\TextAt{8}{133.204}{617.892}{of}%
\TextAt{8}{146.534}{617.892}{rank}%
\TextAt{9}{173.529}{617.892}{1}%
\TextAt{8}{183.279}{617.892}{and}%
\TextAt{13}{206.042}{617.892}{\Character{11}}%
\NativeMathAt{216.884}{617.892}{11.955}{\in}%
\TextAt{13}{228.175}{617.892}{L}%
\TextAt{8}{236.14}{617.892}{,}%
\TextAt{13}{243.289}{617.892}{\Character{11}}%
\NativeMathAt{254.131}{617.892}{11.955}{\neq}%
\TextAt{9}{266.541}{617.892}{0}%
\TextAt{8}{272.41}{617.892}{,}%
\TextAt{8}{279.559}{617.892}{then,}%
\TextAt{8}{309.471}{617.892}{by}%
\TextAt{8}{325.73}{617.892}{(4.3),}%
\TextAt{8}{356.957}{617.892}{(iv),}%
\TextAt{8}{382.648}{617.892}{mapping}%
\TextAt{13}{175.915}{588.89}{L}%
\NativeArrowAt{187.2}{588.89}{11.955}{19.26}{\longrightarrow}%
\TextAt{13}{209.782}{588.89}{L}%
\TextAt{16}{220.403}{588.89}{\Character{10}}%
\TextAt{27}{229.702}{587.097}{o}%
\TextAt{13}{236.753}{588.89}{K}%
\TextAt{9}{250.891}{588.89}{\Character{61}}%
\TextAt{13}{263.317}{588.89}{K}%
\TextAt{9}{274.134}{588.89}{\Character{40}}%
\TextAt{13}{278.687}{588.89}{\Character{11}}%
\TextAt{16}{288.865}{588.89}{\Character{10}}%
\TextAt{9}{300.82}{588.89}{1\Character{41}}%
\TextAt{13}{311.226}{588.89}{\Character{59}}%
\TextAt{13}{340.361}{588.89}{x}%
\NativeArrowAt{350.353}{588.89}{11.955}{19.26}{\longmapsto}%
\TextAt{13}{372.935}{588.89}{f}%
\TextAt{9}{379.981}{588.89}{\Character{40}}%
\TextAt{13}{384.534}{588.89}{x}%
\TextAt{9}{391.186}{588.89}{\Character{41}\Character{40}}%
\TextAt{13}{400.29}{588.89}{\Character{11}}%
\TextAt{16}{410.469}{588.89}{\Character{10}}%
\TextAt{9}{422.424}{588.89}{1\Character{41}}%
\TextAt{13}{432.829}{588.89}{\Character{59}}%
\TextAt{8}{79.226}{559.889}{gives}%
\TextAt{8}{108.839}{559.889}{an}%
\TextAt{8}{126.03}{559.889}{injective}%
\TextAt{14}{173.456}{559.889}{o}%
\TextAt{8}{179.3}{559.889}{-module}%
\TextAt{8}{225.423}{559.889}{homomorphism}%
\TextAt{13}{309.025}{559.889}{L}%
\NativeMathAt{321.899}{559.889}{11.955}{\rightarrow}%
\TextAt{13}{338.763}{559.889}{K}%
\TextAt{8}{349.581}{559.889}{,}%
\TextAt{13}{357.895}{559.889}{x}%
\NativeArrowAt{369.457}{559.889}{11.955}{11.955}{\mapsto}%
\TextAt{13}{386.321}{559.889}{f}%
\TextAt{9}{393.368}{559.889}{\Character{40}}%
\TextAt{13}{397.92}{559.889}{x}%
\TextAt{9}{404.572}{559.889}{\Character{41}}%
\TextAt{8}{409.125}{559.889}{,}%
\TextAt{8}{417.439}{559.889}{onto}%
\TextAt{8}{444.72}{559.889}{a}%
\TextAt{8}{455.408}{559.889}{fractional}%
\TextAt{8}{508.704}{559.889}{ideal}%
\TextAt{14}{79.226}{542.843}{a}%
\NativeMathAt{88.53}{542.843}{11.955}{\subseteq}%
\TextAt{13}{101.149}{542.843}{K}%
\TextAt{8}{111.967}{542.843}{.}%
\TextAt{8}{97.159}{522.809}{To}%
\TextAt{8}{115.199}{522.809}{see}%
\TextAt{8}{134.925}{522.809}{the}%
\TextAt{8}{155.895}{522.809}{connection}%
\TextAt{8}{215.24}{522.809}{with}%
\TextAt{8}{242.713}{522.809}{the}%
\TextAt{8}{263.682}{522.809}{Riemann–Roch}%
\TextAt{8}{346.125}{522.809}{theory}%
\TextAt{8}{383.688}{522.809}{of}%
\TextAt{8}{397.831}{522.809}{the}%
\TextAt{8}{418.801}{522.809}{last}%
\TextAt{8}{441.791}{522.809}{section,}%
\TextAt{8}{485.14}{522.809}{which}%
\TextAt{8}{519.439}{522.809}{we}%
\TextAt{8}{79.226}{505.763}{are}%
\TextAt{8}{98.737}{505.763}{about}%
\TextAt{8}{132.235}{505.763}{to}%
\TextAt{8}{146.546}{505.763}{generalize,}%
\TextAt{8}{203.775}{505.763}{we}%
\TextAt{8}{221.002}{505.763}{observe}%
\TextAt{8}{262.69}{505.763}{that}%
\TextAt{8}{288.059}{505.763}{every}%
\TextAt{8}{318.951}{505.763}{replete}%
\TextAt{8}{357.315}{505.763}{ideal}%
\TextAt{14}{241.08}{476.761}{a}%
\TextAt{9}{250.384}{476.761}{\Character{61}}%
\TextAt{20}{262.809}{488.119}{\Character{89}}%
\TextAt{27}{263.112}{462.287}{p}%
\TextAt{25}{267.1}{462.287}{\Character{45}}%
\NativeMathAt{269.314}{462.287}{7.97}{\infty}%
\TextAt{14}{280.078}{476.761}{p}%
\TextAt{21}{286.061}{481.698}{\Character{23}}%
\TextAt{28}{290.224}{480.701}{p}%
\TextAt{20}{296.797}{488.119}{\Character{89}}%
\TextAt{27}{297.03}{462.216}{p}%
\TextAt{3}{301.019}{462.216}{\Character{106}}%
\NativeMathAt{303.37}{462.216}{7.97}{\infty}%
\TextAt{14}{314.065}{476.761}{p}%
\TextAt{21}{320.049}{481.698}{\Character{23}}%
\TextAt{28}{324.211}{480.701}{p}%
\TextAt{9}{332.113}{476.761}{\Character{61}}%
\TextAt{14}{344.538}{476.761}{a}%
\TextAt{21}{350.521}{474.968}{f}%
\TextAt{14}{355.966}{476.761}{a}%
\NativeMathAt{361.949}{474.968}{7.97}{\infty}%
\TextAt{8}{79.226}{437.11}{of}%
\TextAt{13}{92.605}{437.11}{K}%
\TextAt{8}{107.368}{437.11}{defines}%
\TextAt{8}{145.841}{437.11}{an}%
\TextAt{8}{162.148}{437.11}{invertible,}%
\TextAt{8}{217.167}{437.11}{metrized}%
\TextAt{14}{265.335}{437.11}{o}%
\TextAt{8}{271.179}{437.11}{-module.}%
\TextAt{8}{321.066}{437.11}{In}%
\TextAt{8}{335.736}{437.11}{fact,}%
\TextAt{8}{362.134}{437.11}{the}%
\TextAt{8}{382.339}{437.11}{identity}%
\TextAt{14}{425.641}{437.11}{a}%
\NativeMathAt{431.624}{435.317}{7.97}{\infty}%
\TextAt{9}{443.994}{437.11}{\Character{61}}%
\TextAt{20}{456.503}{446.076}{\Character{81}}%
\TextAt{27}{467.794}{433.623}{p}%
\TextAt{3}{471.782}{433.623}{\Character{106}}%
\NativeMathAt{474.133}{433.623}{7.97}{\infty}%
\TextAt{14}{485.094}{437.11}{p}%
\TextAt{21}{491.077}{441.448}{\Character{23}}%
\TextAt{28}{495.239}{440.452}{p}%
\TextAt{8}{503.766}{437.11}{yields}%
\TextAt{8}{79.226}{420.064}{the}%
\TextAt{8}{99.383}{420.064}{function}%
\TextAt{13}{214.652}{403.018}{\Character{11}}%
\TextAt{9}{225.495}{403.018}{\Character{58}}%
\TextAt{13}{232.067}{403.018}{X}%
\TextAt{9}{242.722}{403.018}{\Character{40}}%
\TextAt{18}{247.275}{403.018}{C}%
\TextAt{9}{255.909}{403.018}{\Character{41}}%
\NativeArrowAt{263.782}{403.018}{11.955}{19.26}{\longrightarrow}%
\TextAt{18}{286.364}{403.018}{R}%
\TextAt{17}{294.998}{401.225}{\Character{43}}%
\TextAt{13}{302.083}{403.018}{\Character{59}}%
\TextAt{13}{331.218}{403.018}{\Character{11}}%
\TextAt{9}{338.759}{403.018}{\Character{40}}%
\TextAt{13}{343.312}{403.018}{\Character{27}}%
\TextAt{9}{350.394}{403.018}{\Character{41}}%
\TextAt{9}{358.249}{403.018}{\Character{61}}%
\TextAt{13}{370.693}{403.018}{e}%
\TextAt{17}{376.118}{407.954}{2}%
\TextAt{21}{380.352}{407.954}{\Character{23}}%
\TextAt{28}{384.515}{406.958}{p}%
\TextAt{22}{388.1}{405.962}{\Character{27}}%
\TextAt{13}{394.092}{403.018}{\Character{59}}%
\TextAt{8}{79.226}{379.496}{where}%
\TextAt{14}{113.365}{379.496}{p}%
\TextAt{21}{119.348}{377.703}{\Character{27}}%
\TextAt{8}{129.199}{379.496}{denotes}%
\TextAt{8}{171.862}{379.496}{as}%
\TextAt{8}{186.563}{379.496}{before}%
\TextAt{8}{222.005}{379.496}{the}%
\TextAt{8}{242.491}{379.496}{infinite}%
\TextAt{8}{282.489}{379.496}{place}%
\TextAt{8}{312.73}{379.496}{defined}%
\TextAt{8}{353.373}{379.496}{by}%
\TextAt{13}{369.962}{379.496}{\Character{27}}%
\TextAt{9}{380.928}{379.496}{\Character{58}}%
\TextAt{13}{388.063}{379.496}{K}%
\NativeMathAt{402.763}{379.496}{11.955}{\rightarrow}%
\TextAt{18}{418.602}{379.496}{C}%
\TextAt{8}{427.236}{379.496}{.}%
\TextAt{8}{436.674}{379.496}{Since}%
\TextAt{14}{467.562}{379.496}{p}%
\Vector{q 0 G 0 g 0.359 w 401.545 -391.61 m 407.168 -391.61 l S Q}%
\TextAt{21}{473.545}{377.703}{\Character{27}}%
\TextAt{9}{483.051}{379.496}{\Character{61}}%
\TextAt{14}{496.039}{379.496}{p}%
\TextAt{21}{502.022}{377.703}{\Character{27}}%
\TextAt{8}{507.646}{379.496}{,}%
\TextAt{8}{515.207}{379.496}{one}%
\TextAt{8}{79.226}{362.45}{has}%
\TextAt{13}{100.1}{362.45}{\Character{11}}%
\TextAt{9}{107.622}{362.45}{\Character{40}}%
\Vector{q 0 G 0 g 0.478 w 40.174 -404.729 m 47.256 -404.729 l S Q}%
\TextAt{13}{112.174}{362.45}{\Character{27}}%
\TextAt{9}{119.256}{362.45}{\Character{41}}%
\TextAt{9}{127.111}{362.45}{\Character{61}}%
\TextAt{13}{139.555}{362.45}{\Character{11}}%
\TextAt{9}{147.077}{362.45}{\Character{40}}%
\TextAt{13}{151.629}{362.45}{\Character{27}}%
\TextAt{9}{158.711}{362.45}{\Character{41}}%
\TextAt{8}{163.264}{362.45}{,}%
\TextAt{8}{170.413}{362.45}{and}%
\TextAt{8}{193.176}{362.45}{we}%
\TextAt{8}{210.403}{362.45}{obtain}%
\TextAt{8}{246.83}{362.45}{on}%
\TextAt{8}{263.089}{362.45}{the}%
\TextAt{8}{283.246}{362.45}{complexification}%
\TextAt{14}{253.31}{333.449}{a}%
\TextAt{21}{259.293}{331.656}{f}%
\TextAt{25}{264.24}{331.656}{C}%
\TextAt{9}{273.815}{333.449}{\Character{61}}%
\TextAt{14}{286.24}{333.449}{a}%
\TextAt{21}{292.223}{331.656}{f}%
\TextAt{16}{300.325}{333.449}{\Character{10}}%
\TextAt{25}{309.624}{331.656}{Z}%
\TextAt{18}{318.092}{333.449}{C}%
\TextAt{9}{330.047}{333.449}{\Character{61}}%
\TextAt{13}{342.472}{333.449}{K}%
\TextAt{25}{352.431}{331.656}{C}%
\TextAt{8}{79.226}{304.448}{the}%
\TextAt{13}{99.383}{304.448}{F}%
\TextAt{8}{108.586}{304.448}{-invariant}%
\TextAt{8}{161.607}{304.448}{hermitian}%
\TextAt{8}{214.939}{304.448}{metric}%
\TextAt{16}{224.123}{275.447}{\Character{104}}%
\TextAt{13}{228.773}{275.447}{x\Character{59}}%
\TextAt{13}{240.645}{275.447}{y}%
\TextAt{16}{246.806}{275.447}{\Character{105}}%
\TextAt{27}{251.455}{273.654}{a}%
\TextAt{9}{259.262}{275.447}{\Character{61}}%
\TextAt{13}{271.688}{275.447}{\Character{11}x}%
\Vector{q 0 G 0 g 0.478 w 213.862 -491.732 m 219.998 -491.732 l S Q}%
\TextAt{13}{285.862}{275.447}{y}%
\TextAt{9}{295.319}{275.447}{\Character{61}}%
\TextAt{20}{314.062}{286.804}{\Character{77}}%
\TextAt{21}{307.745}{260.901}{\Character{27}}%
\NativeMathAt{312.87}{260.901}{7.97}{\in}%
\TextAt{21}{318.516}{260.901}{X}%
\TextAt{17}{326.102}{260.901}{\Character{40}}%
\TextAt{25}{329.396}{260.901}{C}%
\TextAt{17}{335.152}{260.901}{\Character{41}}%
\TextAt{13}{340.438}{275.447}{e}%
\TextAt{17}{345.863}{280.383}{2}%
\TextAt{21}{350.097}{280.383}{\Character{23}}%
\TextAt{28}{354.26}{279.387}{p}%
\TextAt{22}{357.845}{278.391}{\Character{27}}%
\TextAt{13}{363.837}{275.447}{x}%
\TextAt{21}{370.489}{273.654}{\Character{27}}%
\Vector{q 0 G 0 g 0.478 w 304.112 -491.732 m 310.249 -491.732 l S Q}%
\TextAt{13}{376.112}{275.447}{y}%
\TextAt{21}{382.249}{272.624}{\Character{27}}%
\TextAt{8}{79.226}{235.796}{(see}%
\TextAt{8}{102.694}{235.796}{example}%
\TextAt{8}{148.542}{235.796}{1,}%
\TextAt{8}{161.549}{235.796}{p.227).}%
\TextAt{8}{201.882}{235.796}{We}%
\TextAt{8}{222.027}{235.796}{denote}%
\TextAt{8}{259.745}{235.796}{the}%
\TextAt{8}{279.902}{235.796}{metrized}%
\TextAt{14}{328.021}{235.796}{o}%
\TextAt{8}{333.865}{235.796}{-module}%
\TextAt{8}{379.056}{235.796}{thus}%
\TextAt{8}{404.807}{235.796}{obtained}%
\TextAt{8}{452.939}{235.796}{by}%
\TextAt{13}{469.198}{235.796}{L}%
\TextAt{9}{477.162}{235.796}{\Character{40}}%
\TextAt{14}{481.715}{235.796}{a}%
\TextAt{9}{487.698}{235.796}{\Character{41}}%
\TextAt{8}{492.25}{235.796}{.}%
\TextAt{8}{97.159}{215.761}{The}%
\TextAt{8}{122.678}{215.761}{ordinary}%
\TextAt{8}{171.306}{215.761}{fractional}%
\TextAt{8}{225.134}{215.761}{ideals,}%
\TextAt{8}{262.795}{215.761}{i.e.,}%
\TextAt{8}{286.731}{215.761}{the}%
\TextAt{8}{308.352}{215.761}{replete}%
\TextAt{8}{348.181}{215.761}{ideals}%
\TextAt{14}{382.224}{215.761}{a}%
\TextAt{8}{393.569}{215.761}{such}%
\TextAt{8}{421.431}{215.761}{that}%
\TextAt{14}{448.265}{215.761}{a}%
\NativeMathAt{454.248}{213.968}{7.97}{\infty}%
\TextAt{9}{469.032}{215.761}{\Character{61}}%
\TextAt{9}{483.94}{215.761}{1}%
\TextAt{8}{489.806}{215.761}{,}%
\TextAt{8}{498.787}{215.761}{and}%
\TextAt{8}{523.014}{215.761}{in}%
\TextAt{8}{79.226}{198.715}{particular}%
\TextAt{14}{133.216}{198.715}{o}%
\TextAt{8}{142.957}{198.715}{itself,}%
\TextAt{8}{174.555}{198.715}{are}%
\TextAt{8}{194.066}{198.715}{thus}%
\TextAt{8}{219.817}{198.715}{equipped}%
\TextAt{8}{269.885}{198.715}{with}%
\TextAt{8}{296.545}{198.715}{the}%
\TextAt{19}{316.702}{198.715}{trivial}%
\TextAt{8}{351.551}{198.715}{metric}%
\TextAt{16}{387.967}{198.715}{\Character{104}}%
\TextAt{13}{392.616}{198.715}{x\Character{59}}%
\TextAt{13}{404.488}{198.715}{y}%
\TextAt{16}{410.649}{198.715}{\Character{105}}%
\TextAt{27}{415.298}{196.922}{a}%
\TextAt{9}{423.106}{198.715}{\Character{61}}%
\TextAt{16}{435.531}{198.715}{\Character{104}}%
\TextAt{13}{440.18}{198.715}{x\Character{59}}%
\TextAt{13}{452.052}{198.715}{y}%
\TextAt{16}{458.213}{198.715}{\Character{105}}%
\TextAt{9}{466.183}{198.715}{\Character{61}}%
\TextAt{13}{478.609}{198.715}{x}%
\Vector{q 0 G 0 g 0.478 w 413.261 -568.464 m 419.398 -568.464 l S Q}%
\TextAt{13}{485.261}{198.715}{y}%
\TextAt{8}{491.398}{198.715}{.}%
\TextAt{12}{79.226}{172.702}{Definition}%
\TextAt{12}{141.018}{172.702}{(4.4).}%
\TextAt{8}{175.467}{172.702}{Two}%
\TextAt{8}{199.958}{172.702}{metrized}%
\TextAt{14}{247.211}{172.702}{o}%
\TextAt{8}{253.055}{172.702}{-modules}%
\TextAt{13}{301.994}{172.702}{M}%
\TextAt{8}{317.598}{172.702}{and}%
\TextAt{13}{339.494}{172.702}{M}%
\TextAt{3}{352.068}{177.041}{\Character{48}}%
\TextAt{8}{357.894}{172.702}{are}%
\TextAt{8}{376.538}{172.702}{called}%
\TextAt{12}{408.835}{172.702}{isometric}%
\TextAt{8}{465.449}{172.702}{if}%
\TextAt{8}{475.307}{172.702}{there}%
\TextAt{8}{504.352}{172.702}{exists}%
\TextAt{8}{79.226}{155.657}{an}%
\TextAt{8}{95.485}{155.657}{isomorphism}%
\TextAt{13}{270.605}{138.611}{f}%
\TextAt{9}{280.973}{138.611}{\Character{58}}%
\TextAt{13}{287.545}{138.611}{M}%
\NativeArrowAt{303.44}{138.611}{11.955}{19.26}{\longrightarrow}%
\TextAt{13}{326.022}{138.611}{M}%
\TextAt{3}{338.595}{143.547}{\Character{48}}%
\TextAt{8}{79.226}{115.089}{of}%
\TextAt{14}{92.556}{115.089}{o}%
\TextAt{8}{98.4}{115.089}{-modules}%
\TextAt{8}{148.206}{115.089}{which}%
\TextAt{8}{181.692}{115.089}{induces}%
\TextAt{8}{223.368}{115.089}{an}%
\TextAt{8}{239.627}{115.089}{isometry}%
\TextAt{13}{287.495}{115.089}{f}%
\TextAt{25}{293.265}{113.296}{C}%
\TextAt{9}{302.84}{115.089}{\Character{58}}%
\TextAt{13}{309.413}{115.089}{M}%
\TextAt{25}{320.738}{113.296}{C}%
\NativeMathAt{330.313}{115.089}{11.955}{\rightarrow}%
\TextAt{13}{345.589}{115.089}{M}%
\TextAt{3}{358.163}{119.427}{\Character{48}}%
\TextAt{25}{356.914}{112.024}{C}%
\TextAt{8}{363.168}{115.089}{.}%
\TextAt{12}{79.226}{83.099}{Proposition}%
\TextAt{12}{151.937}{83.099}{(4.5).}%
\TextAt{8}{297.211}{43.033}{210}%
\Destination{78.23}{809.815}{nav.page.210}%
\Destination{79.36}{776.155}{nav.print.229}%
\Destination{79.36}{186.655}{nav.number.3.4.4}%
\Destination{79.36}{97.055}{nav.number.3.4.5}%
\Hotlink{170.804}{233.555}{18.574}{11.047}{nav.print.227}%
\Hotlink{325.23}{615.655}{25.078}{11.047}{nav.number.3.4.3}%
}

% Page 223 (body page 211)
\Page{%
\TextAt{2}{79.226}{796.943}{III.}%
\TextAt{2}{98.317}{796.943}{Riemann–Roch}%
\TextAt{2}{174.375}{796.943}{Theory}%
\Vector{q 0 G 0 g 0.299 w 7.226 18.729 m 460.769 18.729 l S Q}%
\TextAt{8}{89.579}{761.124}{(i)}%
\TextAt{8}{107.918}{761.124}{Two}%
\TextAt{8}{132.687}{761.124}{replete}%
\TextAt{8}{170.462}{761.124}{ideals}%
\TextAt{14}{202.451}{761.124}{a}%
\TextAt{8}{211.742}{761.124}{and}%
\TextAt{14}{233.916}{761.124}{b}%
\TextAt{8}{243.36}{761.124}{define}%
\TextAt{8}{276.58}{761.124}{isometric}%
\TextAt{8}{326.131}{761.124}{metrized}%
\TextAt{14}{373.662}{761.124}{o}%
\TextAt{8}{379.506}{761.124}{-modules}%
\TextAt{13}{428.722}{761.124}{L}%
\TextAt{9}{436.687}{761.124}{\Character{40}}%
\TextAt{14}{441.239}{761.124}{a}%
\TextAt{9}{447.222}{761.124}{\Character{41}}%
\TextAt{8}{455.083}{761.124}{and}%
\TextAt{13}{477.257}{761.124}{L}%
\TextAt{9}{485.221}{761.124}{\Character{40}}%
\TextAt{14}{489.774}{761.124}{b}%
\TextAt{9}{495.908}{761.124}{\Character{41}}%
\TextAt{8}{503.769}{761.124}{if}%
\TextAt{8}{513.904}{761.124}{and}%
\TextAt{8}{107.918}{744.078}{only}%
\TextAt{8}{133.61}{744.078}{if}%
\TextAt{8}{144.334}{744.078}{they}%
\TextAt{8}{170.671}{744.078}{differ}%
\TextAt{8}{200.906}{744.078}{by}%
\TextAt{8}{217.165}{744.078}{a}%
\TextAt{8}{226.92}{744.078}{replete}%
\TextAt{8}{265.284}{744.078}{principal}%
\TextAt{8}{314.062}{744.078}{ideal}%
\TextAt{9}{342.025}{744.078}{\Character{91}}%
\TextAt{13}{345.276}{744.078}{a}%
\TextAt{9}{351.421}{744.078}{\Character{93}}%
\TextAt{8}{354.673}{744.078}{:}%
\TextAt{14}{363.121}{744.078}{a}%
\TextAt{9}{372.425}{744.078}{\Character{61}}%
\TextAt{14}{384.851}{744.078}{b}%
\TextAt{9}{390.985}{744.078}{\Character{91}}%
\TextAt{13}{394.237}{744.078}{a}%
\TextAt{9}{400.382}{744.078}{\Character{93}}%
\TextAt{8}{403.634}{744.078}{.}%
\TextAt{8}{86.327}{727.032}{(ii)}%
\TextAt{8}{107.918}{727.032}{Every}%
\TextAt{8}{141.572}{727.032}{invertible}%
\TextAt{8}{193.278}{727.032}{metrized}%
\TextAt{14}{241.398}{727.032}{o}%
\TextAt{8}{247.242}{727.032}{-module}%
\TextAt{8}{292.433}{727.032}{is}%
\TextAt{8}{304.197}{727.032}{isometric}%
\TextAt{8}{354.337}{727.032}{to}%
\TextAt{8}{368.647}{727.032}{an}%
\TextAt{14}{384.906}{727.032}{o}%
\TextAt{8}{390.75}{727.032}{-module}%
\TextAt{13}{435.941}{727.032}{L}%
\TextAt{9}{443.905}{727.032}{\Character{40}}%
\TextAt{14}{448.458}{727.032}{a}%
\TextAt{9}{454.441}{727.032}{\Character{41}}%
\TextAt{8}{458.993}{727.032}{.}%
\TextAt{8}{83.076}{709.986}{(iii)}%
\TextAt{13}{107.918}{709.986}{L}%
\TextAt{9}{115.883}{709.986}{\Character{40}}%
\TextAt{14}{120.435}{709.986}{ab}%
\TextAt{9}{132.553}{709.986}{\Character{41}}%
\NativeMathAt{140.426}{709.986}{11.955}{\cong}%
\TextAt{13}{153.046}{709.986}{L}%
\TextAt{9}{161.01}{709.986}{\Character{40}}%
\TextAt{14}{165.563}{709.986}{a}%
\TextAt{9}{171.546}{709.986}{\Character{41}}%
\TextAt{16}{178.755}{709.986}{\Character{10}}%
\TextAt{27}{188.053}{708.193}{o}%
\TextAt{13}{195.104}{709.986}{L}%
\TextAt{9}{203.069}{709.986}{\Character{40}}%
\TextAt{14}{207.621}{709.986}{b}%
\TextAt{9}{213.756}{709.986}{\Character{41}}%
\TextAt{8}{218.308}{709.986}{,}%
\TextAt{13}{225.457}{709.986}{L}%
\TextAt{9}{233.422}{709.986}{\Character{40}}%
\TextAt{14}{237.974}{709.986}{a}%
\TextAt{3}{243.957}{714.324}{\Character{0}}%
\TextAt{17}{250.544}{714.324}{1}%
\TextAt{9}{255.276}{709.986}{\Character{41}}%
\TextAt{9}{263.131}{709.986}{\Character{61}}%
\TextAt{9}{276.956}{713.008}{\Character{20}}%
\TextAt{13}{275.575}{709.986}{L}%
\TextAt{9}{283.539}{709.986}{\Character{40}}%
\TextAt{14}{288.091}{709.986}{a}%
\TextAt{9}{294.075}{709.986}{\Character{41}}%
\TextAt{8}{298.627}{709.986}{.}%
\TextAt{19}{79.226}{675.117}{Proof.}%
\TextAt{8}{116}{675.117}{(i)}%
\TextAt{8}{132.259}{675.117}{Let}%
\TextAt{14}{153.229}{675.117}{a}%
\TextAt{9}{162.533}{675.117}{\Character{61}}%
\TextAt{20}{174.958}{684.084}{\Character{81}}%
\TextAt{27}{186.249}{671.63}{p}%
\TextAt{14}{192.728}{675.117}{p}%
\TextAt{21}{198.712}{679.456}{\Character{23}}%
\TextAt{28}{202.874}{678.459}{p}%
\TextAt{8}{207.455}{675.117}{,}%
\TextAt{14}{214.604}{675.117}{b}%
\TextAt{9}{224.06}{675.117}{\Character{61}}%
\TextAt{20}{236.485}{684.084}{\Character{81}}%
\TextAt{27}{247.776}{671.63}{p}%
\TextAt{14}{254.256}{675.117}{p}%
\TextAt{21}{260.239}{679.456}{\Character{22}}%
\TextAt{28}{265.333}{678.459}{p}%
\TextAt{8}{269.914}{675.117}{,}%
\TextAt{9}{277.063}{675.117}{\Character{91}}%
\TextAt{13}{280.315}{675.117}{a}%
\TextAt{9}{286.46}{675.117}{\Character{93}}%
\TextAt{9}{293.011}{675.117}{\Character{61}}%
\TextAt{20}{305.458}{684.084}{\Character{81}}%
\TextAt{27}{316.749}{671.63}{p}%
\TextAt{14}{323.228}{675.117}{p}%
\TextAt{21}{329.211}{679.456}{v}%
\TextAt{28}{333.335}{678.459}{p}%
\TextAt{17}{337.418}{679.456}{\Character{40}}%
\TextAt{21}{340.711}{679.456}{a}%
\TextAt{17}{345.209}{679.456}{\Character{41}}%
\TextAt{8}{349}{675.117}{,}%
\TextAt{8}{356.149}{675.117}{and}%
\TextAt{8}{378.912}{675.117}{let}%
\TextAt{13}{175.473}{646.309}{\Character{11}}%
\TextAt{9}{182.995}{646.309}{\Character{40}}%
\TextAt{13}{187.547}{646.309}{\Character{27}}%
\TextAt{9}{194.629}{646.309}{\Character{41}}%
\TextAt{9}{202.484}{646.309}{\Character{61}}%
\TextAt{13}{214.928}{646.309}{e}%
\TextAt{17}{220.353}{651.245}{2}%
\TextAt{21}{224.588}{651.245}{\Character{23}}%
\TextAt{28}{228.75}{650.249}{p}%
\TextAt{22}{232.335}{649.252}{\Character{27}}%
\TextAt{13}{238.327}{646.309}{\Character{59}}%
\TextAt{13}{267.462}{646.309}{\Character{12}}%
\TextAt{9}{274.753}{646.309}{\Character{40}}%
\TextAt{13}{279.305}{646.309}{\Character{27}}%
\TextAt{9}{286.387}{646.309}{\Character{41}}%
\TextAt{9}{294.242}{646.309}{\Character{61}}%
\TextAt{13}{306.686}{646.309}{e}%
\TextAt{17}{312.111}{651.245}{2}%
\TextAt{21}{316.346}{651.245}{\Character{22}}%
\TextAt{28}{321.44}{650.249}{p}%
\TextAt{22}{325.025}{649.252}{\Character{27}}%
\TextAt{13}{331.017}{646.309}{\Character{59}}%
\TextAt{13}{360.152}{646.309}{\Character{13}}%
\TextAt{9}{366.893}{646.309}{\Character{40}}%
\TextAt{13}{371.446}{646.309}{\Character{27}}%
\TextAt{9}{378.528}{646.309}{\Character{41}}%
\TextAt{9}{386.383}{646.309}{\Character{61}}%
\TextAt{13}{398.827}{646.309}{e}%
\TextAt{17}{404.252}{651.245}{2}%
\TextAt{21}{408.486}{651.245}{v}%
\TextAt{28}{412.61}{650.249}{p}%
\TextAt{22}{416.195}{649.252}{\Character{27}}%
\TextAt{17}{421.688}{651.245}{\Character{40}}%
\TextAt{21}{424.982}{651.245}{a}%
\TextAt{17}{429.48}{651.245}{\Character{41}}%
\TextAt{13}{433.271}{646.309}{\Character{58}}%
\TextAt{8}{79.226}{617.5}{If}%
\TextAt{14}{89.975}{617.5}{a}%
\TextAt{9}{99.279}{617.5}{\Character{61}}%
\TextAt{14}{111.705}{617.5}{b}%
\TextAt{9}{117.839}{617.5}{\Character{91}}%
\TextAt{13}{121.091}{617.5}{a}%
\TextAt{9}{127.236}{617.5}{\Character{93}}%
\TextAt{8}{130.488}{617.5}{,}%
\TextAt{8}{136.882}{617.5}{then}%
\TextAt{13}{162.599}{617.5}{\Character{23}}%
\TextAt{27}{168.383}{615.707}{p}%
\TextAt{9}{176.191}{617.5}{\Character{61}}%
\TextAt{13}{188.616}{617.5}{\Character{22}}%
\TextAt{27}{195.659}{615.707}{p}%
\TextAt{9}{200.874}{617.5}{\Character{43}}%
\TextAt{13}{210.707}{617.5}{v}%
\TextAt{27}{216.366}{615.707}{p}%
\TextAt{9}{220.853}{617.5}{\Character{40}}%
\TextAt{13}{225.406}{617.5}{a}%
\TextAt{9}{231.55}{617.5}{\Character{41}}%
\TextAt{8}{236.103}{617.5}{;}%
\TextAt{8}{242.623}{617.5}{thus}%
\TextAt{13}{267.432}{617.5}{\Character{11}}%
\TextAt{9}{278.274}{617.5}{\Character{61}}%
\TextAt{13}{290.7}{617.5}{\Character{12}}%
\TextAt{13}{297.969}{617.5}{\Character{13}}%
\TextAt{8}{304.693}{617.5}{,}%
\TextAt{8}{311.088}{617.5}{and}%
\TextAt{14}{332.907}{617.5}{a}%
\TextAt{21}{338.891}{615.707}{f}%
\TextAt{9}{347.656}{617.5}{\Character{61}}%
\TextAt{14}{360.082}{617.5}{b}%
\TextAt{21}{366.217}{615.707}{f}%
\TextAt{9}{371.662}{617.5}{\Character{40}}%
\TextAt{13}{376.214}{617.5}{a}%
\TextAt{9}{382.359}{617.5}{\Character{41}}%
\TextAt{8}{386.911}{617.5}{.}%
\TextAt{8}{395.045}{617.5}{The}%
\TextAt{14}{418.156}{617.5}{o}%
\TextAt{8}{424}{617.5}{-module}%
\TextAt{8}{468.247}{617.5}{isomorphism}%
\TextAt{14}{79.226}{600.454}{b}%
\TextAt{21}{85.361}{598.661}{f}%
\NativeMathAt{94.127}{600.454}{11.955}{\rightarrow}%
\TextAt{14}{109.403}{600.454}{a}%
\TextAt{21}{115.386}{598.661}{f}%
\TextAt{8}{120.831}{600.454}{,}%
\TextAt{13}{127.691}{600.454}{x}%
\NativeArrowAt{137.664}{600.454}{11.955}{11.955}{\mapsto}%
\TextAt{13}{152.94}{600.454}{ax}%
\TextAt{8}{165.737}{600.454}{,}%
\TextAt{8}{172.597}{600.454}{takes}%
\TextAt{8}{202.219}{600.454}{the}%
\TextAt{8}{222.014}{600.454}{form}%
\TextAt{16}{249.292}{600.454}{\Character{104}}%
\TextAt{13}{257.477}{600.454}{\Character{59}}%
\TextAt{16}{266.257}{600.454}{\Character{105}}%
\TextAt{27}{270.907}{598.661}{b}%
\TextAt{8}{279.031}{600.454}{to}%
\TextAt{8}{292.979}{600.454}{the}%
\TextAt{8}{312.774}{600.454}{form}%
\TextAt{16}{340.053}{600.454}{\Character{104}}%
\TextAt{13}{348.238}{600.454}{\Character{59}}%
\TextAt{16}{357.018}{600.454}{\Character{105}}%
\TextAt{27}{361.667}{598.661}{a}%
\TextAt{8}{366.154}{600.454}{.}%
\TextAt{8}{374.482}{600.454}{Indeed,}%
\TextAt{8}{415.474}{600.454}{viewing}%
\TextAt{13}{457.709}{600.454}{a}%
\TextAt{8}{467.39}{600.454}{as}%
\TextAt{8}{481.398}{600.454}{embedded}%
\TextAt{8}{79.226}{583.408}{in}%
\TextAt{13}{92.879}{583.408}{K}%
\TextAt{25}{102.838}{581.615}{C}%
\TextAt{8}{109.092}{583.408}{,}%
\TextAt{8}{116.242}{583.408}{we}%
\TextAt{8}{133.469}{583.408}{find}%
\TextAt{13}{156.877}{583.408}{a}%
\TextAt{9}{166.343}{583.408}{\Character{61}}%
\TextAt{20}{178.768}{592.375}{\Character{76}}%
\TextAt{21}{192.052}{579.921}{\Character{27}}%
\TextAt{13}{199.668}{583.408}{\Character{27}}%
\TextAt{13}{206.734}{583.408}{a}%
\TextAt{8}{216.793}{583.408}{and}%
\TextAt{13}{242.916}{554.6}{a}%
\Vector{q 0 G 0 g 0.478 w 177.061 -212.579 m 183.206 -212.579 l S Q}%
\TextAt{13}{249.061}{554.6}{a}%
\TextAt{9}{258.526}{554.6}{\Character{61}}%
\TextAt{20}{270.952}{565.957}{\Character{77}}%
\TextAt{21}{277.422}{540.851}{\Character{27}}%
\TextAt{13}{291.01}{554.6}{e}%
\TextAt{3}{296.435}{559.536}{\Character{0}}%
\TextAt{17}{303.022}{559.536}{2}%
\TextAt{21}{307.256}{559.536}{v}%
\TextAt{28}{311.38}{558.54}{p}%
\TextAt{22}{314.964}{557.544}{\Character{27}}%
\TextAt{17}{320.458}{559.536}{\Character{40}}%
\TextAt{21}{323.751}{559.536}{a}%
\TextAt{17}{328.249}{559.536}{\Character{41}}%
\TextAt{9}{335.362}{554.6}{\Character{61}}%
\TextAt{13}{347.787}{554.6}{\Character{13}}%
\TextAt{3}{354.509}{559.536}{\Character{0}}%
\TextAt{17}{361.096}{559.536}{1}%
\TextAt{13}{365.828}{554.6}{\Character{59}}%
\TextAt{8}{79.226}{517.931}{because}%
\TextAt{13}{122.528}{517.931}{v}%
\TextAt{27}{128.187}{516.137}{p}%
\TextAt{22}{132.176}{515.141}{\Character{27}}%
\TextAt{9}{137.67}{517.931}{\Character{40}}%
\TextAt{13}{142.222}{517.931}{a}%
\TextAt{9}{148.367}{517.931}{\Character{41}}%
\TextAt{9}{156.222}{517.931}{\Character{61}}%
\TextAt{16}{168.666}{517.931}{\Character{0}}%
\TextAt{8}{179.957}{517.931}{log}%
\TextAt{16}{196.917}{517.931}{\Character{106}}%
\TextAt{13}{200.238}{517.931}{\Character{27}}%
\TextAt{13}{207.304}{517.931}{a}%
\TextAt{16}{213.465}{517.931}{\Character{106}}%
\TextAt{8}{216.786}{517.931}{,}%
\TextAt{8}{223.935}{517.931}{so}%
\TextAt{8}{238.305}{517.931}{that}%
\TextAt{16}{167.546}{489.122}{\Character{104}}%
\TextAt{13}{172.195}{489.122}{ax\Character{59}}%
\TextAt{13}{190.212}{489.122}{ay}%
\TextAt{16}{202.518}{489.122}{\Character{105}}%
\TextAt{27}{207.167}{487.329}{a}%
\TextAt{9}{214.975}{489.122}{\Character{61}}%
\TextAt{13}{227.4}{489.122}{\Character{11}}%
\TextAt{16}{234.922}{489.122}{\Character{104}}%
\TextAt{13}{239.571}{489.122}{ax\Character{59}}%
\TextAt{13}{257.587}{489.122}{ay}%
\TextAt{16}{269.894}{489.122}{\Character{105}}%
\TextAt{9}{277.864}{489.122}{\Character{61}}%
\TextAt{13}{290.29}{489.122}{\Character{11}\Character{13}}%
\TextAt{3}{304.534}{494.058}{\Character{0}}%
\TextAt{17}{311.12}{494.058}{1}%
\TextAt{16}{315.852}{489.122}{\Character{104}}%
\TextAt{13}{320.502}{489.122}{x\Character{59}}%
\TextAt{13}{332.374}{489.122}{y}%
\TextAt{16}{338.534}{489.122}{\Character{105}}%
\TextAt{9}{346.505}{489.122}{\Character{61}}%
\TextAt{13}{358.93}{489.122}{\Character{12}}%
\TextAt{16}{366.201}{489.122}{\Character{104}}%
\TextAt{13}{370.851}{489.122}{x\Character{59}}%
\TextAt{13}{382.723}{489.122}{y}%
\TextAt{16}{388.883}{489.122}{\Character{105}}%
\TextAt{9}{396.854}{489.122}{\Character{61}}%
\TextAt{16}{409.279}{489.122}{\Character{104}}%
\TextAt{13}{413.928}{489.122}{x\Character{59}}%
\TextAt{13}{425.8}{489.122}{y}%
\TextAt{16}{431.961}{489.122}{\Character{105}}%
\TextAt{27}{436.61}{487.329}{b}%
\TextAt{13}{441.199}{489.122}{\Character{58}}%
\TextAt{8}{79.226}{460.314}{Therefore}%
\TextAt{14}{132.223}{460.314}{b}%
\TextAt{21}{138.358}{458.521}{f}%
\NativeMathAt{147.124}{460.314}{11.955}{\rightarrow}%
\TextAt{14}{162.4}{460.314}{a}%
\TextAt{21}{168.383}{458.521}{f}%
\TextAt{8}{173.828}{460.314}{,}%
\TextAt{13}{180.977}{460.314}{x}%
\NativeArrowAt{190.95}{460.314}{11.955}{11.955}{\mapsto}%
\TextAt{13}{206.226}{460.314}{ax}%
\TextAt{8}{219.023}{460.314}{,}%
\TextAt{8}{226.172}{460.314}{gives}%
\TextAt{8}{254.853}{460.314}{an}%
\TextAt{8}{271.112}{460.314}{isometry}%
\TextAt{13}{318.98}{460.314}{L}%
\TextAt{9}{326.945}{460.314}{\Character{40}}%
\TextAt{14}{331.497}{460.314}{a}%
\TextAt{9}{337.48}{460.314}{\Character{41}}%
\NativeMathAt{345.353}{460.314}{11.955}{\cong}%
\TextAt{13}{357.973}{460.314}{L}%
\TextAt{9}{365.937}{460.314}{\Character{40}}%
\TextAt{14}{370.49}{460.314}{b}%
\TextAt{9}{376.624}{460.314}{\Character{41}}%
\TextAt{8}{381.177}{460.314}{.}%
\TextAt{8}{97.159}{440.279}{Conversely,}%
\TextAt{8}{158.68}{440.279}{let}%
\TextAt{13}{175.585}{440.279}{g}%
\TextAt{9}{184.94}{440.279}{\Character{58}}%
\TextAt{13}{191.512}{440.279}{L}%
\TextAt{9}{199.477}{440.279}{\Character{40}}%
\TextAt{14}{204.029}{440.279}{b}%
\TextAt{9}{210.164}{440.279}{\Character{41}}%
\NativeMathAt{218.037}{440.279}{11.955}{\rightarrow}%
\TextAt{13}{233.313}{440.279}{L}%
\TextAt{9}{241.278}{440.279}{\Character{40}}%
\TextAt{14}{245.83}{440.279}{a}%
\TextAt{9}{251.813}{440.279}{\Character{41}}%
\TextAt{8}{260.263}{440.279}{be}%
\TextAt{8}{276.187}{440.279}{an}%
\TextAt{8}{292.446}{440.279}{isometry.}%
\TextAt{8}{343.885}{440.279}{Then}%
\TextAt{8}{374.443}{440.279}{the}%
\TextAt{14}{394.599}{440.279}{o}%
\TextAt{8}{400.443}{440.279}{-module}%
\TextAt{8}{445.634}{440.279}{homomorphism}%
\TextAt{13}{273.579}{411.471}{g}%
\TextAt{9}{282.934}{411.471}{\Character{58}}%
\TextAt{14}{289.506}{411.471}{b}%
\TextAt{21}{295.641}{409.677}{f}%
\NativeArrowAt{304.407}{411.471}{11.955}{19.26}{\longrightarrow}%
\TextAt{14}{326.989}{411.471}{a}%
\TextAt{21}{332.972}{409.677}{f}%
\TextAt{8}{79.226}{382.662}{is}%
\TextAt{8}{90.99}{382.662}{given}%
\TextAt{8}{121.559}{382.662}{as}%
\TextAt{8}{135.93}{382.662}{multiplication}%
\TextAt{8}{211.056}{382.662}{by}%
\TextAt{8}{227.315}{382.662}{some}%
\TextAt{8}{256.641}{382.662}{element}%
\TextAt{13}{299.883}{382.662}{a}%
\NativeMathAt{309.348}{382.662}{11.955}{\in}%
\TextAt{14}{320.639}{382.662}{b}%
\TextAt{3}{326.774}{387.61}{\Character{0}}%
\TextAt{17}{333.361}{387.61}{1}%
\TextAt{21}{326.774}{379.333}{f}%
\TextAt{14}{338.093}{382.662}{a}%
\TextAt{21}{344.076}{380.869}{f}%
\NativeMathAt{352.842}{382.662}{11.955}{\cong}%
\TextAt{8}{365.461}{382.662}{Hom}%
\TextAt{27}{389.85}{380.869}{o}%
\TextAt{9}{394.244}{382.662}{\Character{40}}%
\TextAt{14}{398.796}{382.662}{b}%
\TextAt{21}{404.931}{380.869}{f}%
\TextAt{13}{410.376}{382.662}{\Character{59}}%
\TextAt{14}{415.62}{382.662}{a}%
\TextAt{21}{421.603}{380.869}{f}%
\TextAt{9}{427.048}{382.662}{\Character{41}}%
\TextAt{8}{431.601}{382.662}{.}%
\TextAt{8}{440.049}{382.662}{The}%
\TextAt{8}{464.103}{382.662}{identity}%
\TextAt{13}{160.178}{353.854}{\Character{12}}%
\TextAt{16}{167.449}{353.854}{\Character{104}}%
\TextAt{13}{172.098}{353.854}{x\Character{59}}%
\TextAt{13}{183.97}{353.854}{y}%
\TextAt{16}{190.131}{353.854}{\Character{105}}%
\TextAt{9}{198.101}{353.854}{\Character{61}}%
\TextAt{16}{210.527}{353.854}{\Character{104}}%
\TextAt{13}{215.176}{353.854}{x\Character{59}}%
\TextAt{13}{227.048}{353.854}{y}%
\TextAt{16}{233.209}{353.854}{\Character{105}}%
\TextAt{27}{237.858}{352.06}{b}%
\TextAt{9}{245.767}{353.854}{\Character{61}}%
\TextAt{16}{258.192}{353.854}{\Character{104}}%
\TextAt{13}{262.842}{353.854}{g}%
\TextAt{9}{268.876}{353.854}{\Character{40}}%
\TextAt{13}{273.428}{353.854}{x}%
\TextAt{9}{280.08}{353.854}{\Character{41}}%
\TextAt{13}{284.633}{353.854}{\Character{59}}%
\TextAt{13}{289.857}{353.854}{g}%
\TextAt{9}{295.911}{353.854}{\Character{40}}%
\TextAt{13}{300.463}{353.854}{y}%
\TextAt{9}{306.6}{353.854}{\Character{41}}%
\TextAt{16}{311.152}{353.854}{\Character{105}}%
\TextAt{27}{315.802}{352.06}{a}%
\TextAt{9}{323.609}{353.854}{\Character{61}}%
\TextAt{13}{336.035}{353.854}{\Character{11}}%
\TextAt{16}{343.556}{353.854}{\Character{104}}%
\TextAt{13}{348.206}{353.854}{ax\Character{59}}%
\TextAt{13}{366.223}{353.854}{ay}%
\TextAt{16}{378.529}{353.854}{\Character{105}}%
\TextAt{9}{386.499}{353.854}{\Character{61}}%
\TextAt{13}{398.924}{353.854}{\Character{11}\Character{13}}%
\TextAt{3}{413.168}{358.79}{\Character{0}}%
\TextAt{17}{419.755}{358.79}{1}%
\TextAt{16}{424.487}{353.854}{\Character{104}}%
\TextAt{13}{429.136}{353.854}{x\Character{59}}%
\TextAt{13}{441.007}{353.854}{y}%
\TextAt{16}{447.169}{353.854}{\Character{105}}%
\TextAt{8}{79.226}{325.045}{then}%
\TextAt{8}{106.339}{325.045}{implies}%
\TextAt{8}{146.519}{325.045}{that}%
\TextAt{13}{172.341}{325.045}{\Character{11}}%
\TextAt{9}{183.955}{325.045}{\Character{61}}%
\TextAt{13}{197.152}{325.045}{\Character{12}}%
\TextAt{13}{204.421}{325.045}{\Character{13}}%
\TextAt{8}{211.146}{325.045}{,}%
\TextAt{8}{218.861}{325.045}{so}%
\TextAt{8}{233.684}{325.045}{that}%
\TextAt{13}{259.506}{325.045}{\Character{23}}%
\TextAt{27}{265.289}{323.252}{p}%
\TextAt{9}{273.869}{325.045}{\Character{61}}%
\TextAt{13}{287.066}{325.045}{\Character{22}}%
\TextAt{27}{294.109}{323.252}{p}%
\TextAt{9}{301.561}{325.045}{\Character{43}}%
\TextAt{13}{313.631}{325.045}{v}%
\TextAt{27}{319.29}{323.252}{p}%
\TextAt{9}{323.777}{325.045}{\Character{40}}%
\TextAt{13}{328.33}{325.045}{a}%
\TextAt{9}{334.475}{325.045}{\Character{41}}%
\TextAt{8}{343.377}{325.045}{for}%
\TextAt{8}{361.715}{325.045}{all}%
\TextAt{14}{378.427}{325.045}{p}%
\TextAt{16}{388.503}{325.045}{\Character{106}}%
\NativeMathAt{395.903}{325.045}{11.955}{\infty}%
\TextAt{8}{407.871}{325.045}{.}%
\TextAt{8}{417.678}{325.045}{In}%
\TextAt{8}{432.752}{325.045}{view}%
\TextAt{8}{460.188}{325.045}{of}%
\TextAt{14}{473.971}{325.045}{a}%
\TextAt{21}{479.954}{323.252}{f}%
\TextAt{9}{489.491}{325.045}{\Character{61}}%
\TextAt{14}{502.688}{325.045}{b}%
\TextAt{21}{508.823}{323.252}{f}%
\TextAt{9}{514.268}{325.045}{\Character{40}}%
\TextAt{13}{518.82}{325.045}{a}%
\TextAt{9}{524.965}{325.045}{\Character{41}}%
\TextAt{8}{529.518}{325.045}{,}%
\TextAt{8}{79.226}{307.999}{this}%
\TextAt{8}{102.049}{307.999}{yields}%
\TextAt{14}{134.949}{307.999}{a}%
\TextAt{9}{144.253}{307.999}{\Character{61}}%
\TextAt{14}{156.679}{307.999}{b}%
\TextAt{9}{162.813}{307.999}{\Character{91}}%
\TextAt{13}{166.065}{307.999}{a}%
\TextAt{9}{172.21}{307.999}{\Character{93}}%
\TextAt{8}{175.462}{307.999}{.}%
\TextAt{8}{97.159}{287.965}{(ii)}%
\TextAt{8}{118.41}{287.965}{Let}%
\TextAt{13}{141.119}{287.965}{L}%
\TextAt{8}{154.72}{287.965}{be}%
\TextAt{8}{172.384}{287.965}{an}%
\TextAt{8}{190.383}{287.965}{invertible}%
\TextAt{8}{243.829}{287.965}{metrized}%
\TextAt{14}{293.689}{287.965}{o}%
\TextAt{8}{299.533}{287.965}{-module.}%
\TextAt{8}{354.494}{287.965}{As}%
\TextAt{8}{373.521}{287.965}{mentioned}%
\TextAt{8}{432.167}{287.965}{before,}%
\TextAt{8}{472.706}{287.965}{we}%
\TextAt{8}{491.673}{287.965}{have}%
\TextAt{8}{520.408}{287.965}{an}%
\TextAt{8}{79.226}{270.919}{isomorphism}%
\TextAt{13}{275.386}{253.873}{g}%
\TextAt{9}{284.741}{253.873}{\Character{58}}%
\TextAt{13}{291.314}{253.873}{L}%
\NativeArrowAt{302.599}{253.873}{11.955}{19.26}{\longrightarrow}%
\TextAt{14}{325.181}{253.873}{a}%
\TextAt{21}{331.164}{252.079}{f}%
\TextAt{8}{79.226}{230.434}{for}%
\TextAt{8}{96.463}{230.434}{the}%
\TextAt{8}{115.97}{230.434}{underlying}%
\TextAt{14}{173.532}{230.434}{o}%
\TextAt{8}{179.376}{230.434}{-module}%
\TextAt{8}{223.918}{230.434}{onto}%
\TextAt{8}{249.618}{230.434}{a}%
\TextAt{8}{258.725}{230.434}{fractional}%
\TextAt{8}{310.44}{230.434}{ideal}%
\TextAt{14}{337.755}{230.434}{a}%
\TextAt{21}{343.738}{228.64}{f}%
\TextAt{8}{349.183}{230.434}{.}%
\TextAt{8}{357.415}{230.434}{The}%
\TextAt{8}{380.82}{230.434}{isomorphism}%
\TextAt{13}{448.591}{230.434}{g}%
\TextAt{25}{454.196}{228.64}{C}%
\TextAt{9}{463.772}{230.434}{\Character{58}}%
\TextAt{13}{470.344}{230.434}{L}%
\TextAt{25}{478.309}{228.64}{C}%
\NativeMathAt{487.884}{230.434}{11.955}{\rightarrow}%
\TextAt{14}{503.16}{230.434}{a}%
\TextAt{21}{509.143}{228.64}{f}%
\TextAt{25}{514.09}{228.64}{C}%
\TextAt{9}{523.665}{230.434}{\Character{61}}%
\TextAt{13}{79.226}{213.388}{K}%
\TextAt{25}{89.185}{211.594}{C}%
\TextAt{8}{99.497}{213.388}{gives}%
\TextAt{8}{128.338}{213.388}{us}%
\TextAt{8}{143.515}{213.388}{the}%
\TextAt{13}{163.832}{213.388}{F}%
\TextAt{8}{173.035}{213.388}{-invariant}%
\TextAt{8}{226.217}{213.388}{hermitian}%
\TextAt{8}{279.709}{213.388}{metric}%
\TextAt{13}{316.285}{213.388}{h}%
\TextAt{9}{323.024}{213.388}{\Character{40}}%
\TextAt{13}{327.576}{213.388}{x\Character{59}}%
\TextAt{13}{339.447}{213.388}{y}%
\TextAt{9}{345.609}{213.388}{\Character{41}}%
\TextAt{9}{353.738}{213.388}{\Character{61}}%
\TextAt{16}{366.455}{213.388}{\Character{104}}%
\TextAt{13}{371.104}{213.388}{g}%
\TextAt{3}{377.138}{218.336}{\Character{0}}%
\TextAt{17}{383.725}{218.336}{1}%
\TextAt{25}{376.709}{210.102}{C}%
\TextAt{9}{388.457}{213.388}{\Character{40}}%
\TextAt{13}{393.009}{213.388}{x}%
\TextAt{9}{399.662}{213.388}{\Character{41}}%
\TextAt{13}{404.214}{213.388}{\Character{59}}%
\TextAt{13}{409.438}{213.388}{g}%
\TextAt{3}{415.492}{218.336}{\Character{0}}%
\TextAt{17}{422.079}{218.336}{1}%
\TextAt{25}{415.063}{210.102}{C}%
\TextAt{9}{426.811}{213.388}{\Character{40}}%
\TextAt{13}{431.363}{213.388}{y}%
\TextAt{9}{437.5}{213.388}{\Character{41}}%
\TextAt{16}{442.052}{213.388}{\Character{105}}%
\TextAt{21}{446.702}{211.594}{L}%
\TextAt{8}{457.017}{213.388}{on}%
\TextAt{13}{473.436}{213.388}{K}%
\TextAt{25}{483.395}{211.594}{C}%
\TextAt{8}{489.65}{213.388}{.}%
\TextAt{8}{498.58}{213.388}{It}%
\TextAt{8}{511.413}{213.388}{is}%
\TextAt{8}{523.337}{213.388}{of}%
\TextAt{8}{79.226}{196.342}{the}%
\TextAt{8}{99.383}{196.342}{form}%
\TextAt{13}{271.031}{179.296}{h}%
\TextAt{9}{277.77}{179.296}{\Character{40}}%
\TextAt{13}{282.322}{179.296}{x\Character{59}}%
\TextAt{13}{294.194}{179.296}{y}%
\TextAt{9}{300.355}{179.296}{\Character{41}}%
\TextAt{9}{308.21}{179.296}{\Character{61}}%
\TextAt{13}{320.654}{179.296}{\Character{11}x}%
\Vector{q 0 G 0 g 0.478 w 262.828 -587.883 m 268.964 -587.883 l S Q}%
\TextAt{13}{334.828}{179.296}{y}%
\TextAt{8}{79.226}{155.857}{for}%
\TextAt{8}{97.903}{155.857}{some}%
\TextAt{8}{128.021}{155.857}{function}%
\TextAt{13}{174.66}{155.857}{\Character{11}}%
\TextAt{9}{186.852}{155.857}{\Character{58}}%
\TextAt{13}{194.774}{155.857}{X}%
\TextAt{9}{205.429}{155.857}{\Character{40}}%
\TextAt{18}{209.981}{155.857}{C}%
\TextAt{9}{218.616}{155.857}{\Character{41}}%
\NativeMathAt{227.838}{155.857}{11.955}{\rightarrow}%
\TextAt{18}{244.463}{155.857}{R}%
\TextAt{3}{253.098}{160.195}{\Character{3}}%
\TextAt{17}{253.098}{152.901}{\Character{43}}%
\TextAt{8}{264.871}{155.857}{such}%
\TextAt{8}{292.06}{155.857}{that}%
\TextAt{13}{318.221}{155.857}{\Character{11}}%
\TextAt{9}{325.742}{155.857}{\Character{40}}%
\Vector{q 0 G 0 g 0.478 w 258.295 -611.322 m 265.377 -611.322 l S Q}%
\TextAt{13}{330.295}{155.857}{\Character{27}}%
\TextAt{9}{337.377}{155.857}{\Character{41}}%
\TextAt{9}{346.583}{155.857}{\Character{61}}%
\TextAt{13}{360.374}{155.857}{\Character{11}}%
\TextAt{9}{367.896}{155.857}{\Character{40}}%
\TextAt{13}{372.448}{155.857}{\Character{27}}%
\TextAt{9}{379.531}{155.857}{\Character{41}}%
\TextAt{8}{384.083}{155.857}{.}%
\TextAt{8}{394.907}{155.857}{Putting}%
\TextAt{8}{438.785}{155.857}{now}%
\TextAt{13}{463.965}{155.857}{\Character{11}}%
\TextAt{9}{471.487}{155.857}{\Character{40}}%
\TextAt{13}{476.039}{155.857}{\Character{27}}%
\TextAt{9}{483.122}{155.857}{\Character{41}}%
\TextAt{9}{492.327}{155.857}{\Character{61}}%
\TextAt{13}{506.119}{155.857}{e}%
\TextAt{17}{511.544}{160.195}{2}%
\TextAt{21}{515.778}{160.195}{\Character{23}}%
\TextAt{28}{519.941}{159.199}{p}%
\TextAt{22}{523.526}{158.202}{\Character{27}}%
\TextAt{8}{529.518}{155.857}{,}%
\TextAt{8}{79.226}{138.811}{with}%
\TextAt{13}{105.925}{138.811}{\Character{23}}%
\TextAt{27}{111.709}{137.017}{p}%
\TextAt{22}{115.698}{136.021}{\Character{27}}%
\NativeMathAt{124.58}{138.811}{11.955}{\in}%
\TextAt{18}{135.937}{138.811}{R}%
\TextAt{8}{144.572}{138.811}{,}%
\TextAt{8}{151.77}{138.811}{makes}%
\TextAt{14}{186.993}{138.811}{a}%
\TextAt{21}{192.976}{137.017}{f}%
\TextAt{8}{202.358}{138.811}{with}%
\TextAt{8}{229.057}{138.811}{the}%
\TextAt{8}{249.253}{138.811}{metric}%
\TextAt{13}{285.708}{138.811}{h}%
\TextAt{8}{296.383}{138.811}{into}%
\TextAt{8}{320.165}{138.811}{the}%
\TextAt{8}{340.361}{138.811}{metrized}%
\TextAt{14}{388.519}{138.811}{o}%
\TextAt{8}{394.364}{138.811}{-module}%
\TextAt{13}{439.593}{138.811}{L}%
\TextAt{9}{447.558}{138.811}{\Character{40}}%
\TextAt{14}{452.11}{138.811}{a}%
\TextAt{9}{458.093}{138.811}{\Character{41}}%
\TextAt{8}{466.582}{138.811}{associated}%
\TextAt{8}{522.357}{138.811}{to}%
\TextAt{8}{79.226}{121.765}{the}%
\TextAt{8}{99.383}{121.765}{replete}%
\TextAt{8}{137.747}{121.765}{ideal}%
\TextAt{14}{165.71}{121.765}{a}%
\TextAt{9}{175.014}{121.765}{\Character{61}}%
\TextAt{14}{187.439}{121.765}{a}%
\TextAt{21}{193.422}{119.971}{f}%
\TextAt{20}{200.86}{130.731}{\Character{81}}%
\TextAt{27}{212.151}{118.278}{p}%
\TextAt{3}{216.139}{118.278}{\Character{106}}%
\NativeMathAt{218.49}{118.278}{7.97}{\infty}%
\TextAt{14}{229.451}{121.765}{p}%
\TextAt{21}{235.434}{126.103}{\Character{23}}%
\TextAt{28}{239.597}{125.107}{p}%
\TextAt{8}{244.177}{121.765}{,}%
\TextAt{8}{251.327}{121.765}{and}%
\TextAt{13}{274.089}{121.765}{L}%
\TextAt{8}{285.951}{121.765}{is}%
\TextAt{8}{297.715}{121.765}{isometric}%
\TextAt{8}{347.855}{121.765}{to}%
\TextAt{13}{362.165}{121.765}{L}%
\TextAt{9}{370.13}{121.765}{\Character{40}}%
\TextAt{14}{374.682}{121.765}{a}%
\TextAt{9}{380.665}{121.765}{\Character{41}}%
\TextAt{8}{385.218}{121.765}{.}%
\TextAt{8}{97.159}{101.73}{(iii)}%
\TextAt{8}{119.922}{101.73}{Let}%
\TextAt{14}{140.891}{101.73}{a}%
\TextAt{9}{150.195}{101.73}{\Character{61}}%
\TextAt{20}{162.62}{110.696}{\Character{81}}%
\TextAt{27}{173.911}{98.243}{p}%
\TextAt{14}{180.391}{101.73}{p}%
\TextAt{21}{186.374}{106.068}{\Character{23}}%
\TextAt{28}{190.536}{105.072}{p}%
\TextAt{8}{195.117}{101.73}{,}%
\TextAt{14}{202.267}{101.73}{b}%
\TextAt{9}{211.722}{101.73}{\Character{61}}%
\TextAt{20}{224.148}{110.696}{\Character{81}}%
\TextAt{27}{235.439}{98.243}{p}%
\TextAt{14}{241.918}{101.73}{p}%
\TextAt{21}{247.901}{106.068}{\Character{22}}%
\TextAt{28}{252.995}{105.072}{p}%
\TextAt{8}{257.576}{101.73}{,}%
\TextAt{13}{264.726}{101.73}{\Character{11}}%
\TextAt{9}{272.247}{101.73}{\Character{40}}%
\TextAt{13}{276.8}{101.73}{\Character{27}}%
\TextAt{9}{283.882}{101.73}{\Character{41}}%
\TextAt{9}{291.737}{101.73}{\Character{61}}%
\TextAt{13}{304.181}{101.73}{e}%
\TextAt{17}{309.606}{106.068}{2}%
\TextAt{21}{313.84}{106.068}{\Character{23}}%
\TextAt{28}{318.003}{105.072}{p}%
\TextAt{22}{321.588}{104.076}{\Character{27}}%
\TextAt{8}{327.58}{101.73}{,}%
\TextAt{13}{334.729}{101.73}{\Character{12}}%
\TextAt{9}{342}{101.73}{\Character{40}}%
\TextAt{13}{346.552}{101.73}{\Character{27}}%
\TextAt{9}{353.635}{101.73}{\Character{41}}%
\TextAt{9}{361.49}{101.73}{\Character{61}}%
\TextAt{13}{373.933}{101.73}{e}%
\TextAt{17}{379.359}{106.068}{2}%
\TextAt{21}{383.593}{106.068}{\Character{22}}%
\TextAt{28}{388.687}{105.072}{p}%
\TextAt{22}{392.272}{104.076}{\Character{27}}%
\TextAt{8}{398.264}{101.73}{.}%
\TextAt{8}{406.713}{101.73}{The}%
\TextAt{8}{430.766}{101.73}{isomorphism}%
\TextAt{14}{212.953}{72.921}{a}%
\TextAt{21}{218.936}{71.128}{f}%
\TextAt{16}{227.038}{72.921}{\Character{10}}%
\TextAt{27}{236.336}{71.128}{o}%
\TextAt{14}{243.387}{72.921}{b}%
\TextAt{21}{249.522}{71.128}{f}%
\NativeArrowAt{258.288}{72.921}{11.955}{19.26}{\longrightarrow}%
\TextAt{14}{280.87}{72.921}{a}%
\TextAt{21}{286.853}{71.128}{f}%
\TextAt{14}{292.298}{72.921}{b}%
\TextAt{21}{298.433}{71.128}{f}%
\TextAt{13}{303.878}{72.921}{\Character{59}}%
\TextAt{13}{333.013}{72.921}{a}%
\TextAt{16}{341.834}{72.921}{\Character{10}}%
\TextAt{13}{353.789}{72.921}{b}%
\NativeArrowAt{362.087}{72.921}{11.955}{19.26}{\longmapsto}%
\TextAt{13}{384.669}{72.921}{ab\Character{59}}%
\TextAt{8}{297.211}{43.033}{211}%
\Destination{79.23}{773.075}{nav.Item.78}%
\Destination{79.23}{741.085}{nav.Item.79}%
\Destination{79.23}{724.035}{nav.Item.80}%
\Destination{78.23}{809.815}{nav.page.211}%
\Destination{79.36}{776.155}{nav.print.230}%
}

% Page 224 (body page 212)
\Page{%
\TextAt{2}{79.226}{796.943}{III.}%
\TextAt{2}{98.317}{796.943}{Riemann–Roch}%
\TextAt{2}{174.375}{796.943}{Theory}%
\Vector{q 0 G 0 g 0.299 w 7.226 18.729 m 460.769 18.729 l S Q}%
\TextAt{8}{79.226}{761.124}{between}%
\TextAt{8}{124.417}{761.124}{the}%
\TextAt{14}{144.573}{761.124}{o}%
\TextAt{8}{150.417}{761.124}{-modules}%
\TextAt{8}{200.223}{761.124}{underlying}%
\TextAt{13}{258.432}{761.124}{L}%
\TextAt{9}{266.397}{761.124}{\Character{40}}%
\TextAt{14}{270.949}{761.124}{a}%
\TextAt{9}{276.932}{761.124}{\Character{41}}%
\TextAt{16}{284.141}{761.124}{\Character{10}}%
\TextAt{27}{293.44}{759.33}{o}%
\TextAt{13}{300.491}{761.124}{L}%
\TextAt{9}{308.455}{761.124}{\Character{40}}%
\TextAt{14}{313.007}{761.124}{b}%
\TextAt{9}{319.142}{761.124}{\Character{41}}%
\TextAt{8}{327.592}{761.124}{and}%
\TextAt{13}{350.355}{761.124}{L}%
\TextAt{9}{358.319}{761.124}{\Character{40}}%
\TextAt{14}{362.872}{761.124}{ab}%
\TextAt{9}{374.989}{761.124}{\Character{41}}%
\TextAt{8}{383.439}{761.124}{then}%
\TextAt{8}{410.099}{761.124}{yields,}%
\TextAt{8}{446.252}{761.124}{as}%
\TextAt{16}{165.109}{733.965}{\Character{104}}%
\TextAt{13}{169.758}{733.965}{ab\Character{59}}%
\TextAt{13}{186.101}{733.965}{a}%
\TextAt{3}{192.269}{738.901}{\Character{48}}%
\TextAt{13}{195.064}{733.965}{b}%
\TextAt{3}{200.041}{738.901}{\Character{48}}%
\TextAt{16}{202.836}{733.965}{\Character{105}}%
\TextAt{27}{207.486}{732.171}{ab}%
\TextAt{9}{219.383}{733.965}{\Character{61}}%
\TextAt{13}{231.809}{733.965}{\Character{11}\Character{12}}%
\TextAt{13}{246.586}{733.965}{ab}%
\Vector{q 0 G 0 g 0.478 w 185.724 -30.06 m 202.436 -30.06 l S Q}%
\TextAt{13}{257.724}{733.965}{a}%
\TextAt{3}{263.869}{737.418}{\Character{48}}%
\TextAt{13}{266.664}{733.965}{b}%
\TextAt{3}{271.641}{737.418}{\Character{48}}%
\TextAt{9}{277.757}{733.965}{\Character{61}}%
\TextAt{13}{290.182}{733.965}{\Character{11}}%
\TextAt{16}{297.704}{733.965}{\Character{104}}%
\TextAt{13}{302.353}{733.965}{a\Character{59}}%
\TextAt{13}{313.722}{733.965}{a}%
\TextAt{3}{319.887}{738.901}{\Character{48}}%
\TextAt{16}{322.682}{733.965}{\Character{105}}%
\TextAt{13}{327.332}{733.965}{\Character{12}}%
\TextAt{16}{334.603}{733.965}{\Character{104}}%
\TextAt{13}{339.252}{733.965}{b\Character{59}}%
\TextAt{13}{349.45}{733.965}{b}%
\TextAt{3}{354.45}{738.901}{\Character{48}}%
\TextAt{16}{357.246}{733.965}{\Character{105}}%
\TextAt{9}{365.216}{733.965}{\Character{61}}%
\TextAt{16}{377.641}{733.965}{\Character{104}}%
\TextAt{13}{382.29}{733.965}{a\Character{59}}%
\TextAt{13}{393.659}{733.965}{a}%
\TextAt{3}{399.824}{738.901}{\Character{48}}%
\TextAt{16}{402.619}{733.965}{\Character{105}}%
\TextAt{27}{407.269}{732.171}{a}%
\TextAt{16}{411.756}{733.965}{\Character{104}}%
\TextAt{13}{416.405}{733.965}{b\Character{59}}%
\TextAt{13}{426.603}{733.965}{b}%
\TextAt{3}{431.603}{738.901}{\Character{48}}%
\TextAt{16}{434.398}{733.965}{\Character{105}}%
\TextAt{27}{439.048}{732.171}{b}%
\TextAt{13}{443.636}{733.965}{\Character{59}}%
\TextAt{8}{79.226}{706.805}{an}%
\TextAt{8}{95.485}{706.805}{isometry}%
\TextAt{13}{143.354}{706.805}{L}%
\TextAt{9}{151.318}{706.805}{\Character{40}}%
\TextAt{14}{155.871}{706.805}{a}%
\TextAt{9}{161.854}{706.805}{\Character{41}}%
\TextAt{16}{169.063}{706.805}{\Character{10}}%
\TextAt{27}{178.361}{705.012}{o}%
\TextAt{13}{185.412}{706.805}{L}%
\TextAt{9}{193.377}{706.805}{\Character{40}}%
\TextAt{14}{197.929}{706.805}{b}%
\TextAt{9}{204.064}{706.805}{\Character{41}}%
\NativeMathAt{211.937}{706.805}{11.955}{\cong}%
\TextAt{13}{224.556}{706.805}{L}%
\TextAt{9}{232.521}{706.805}{\Character{40}}%
\TextAt{14}{237.073}{706.805}{ab}%
\TextAt{9}{249.191}{706.805}{\Character{41}}%
\TextAt{8}{253.743}{706.805}{.}%
\TextAt{8}{97.159}{686.771}{The}%
\TextAt{14}{121.213}{686.771}{o}%
\TextAt{8}{127.057}{686.771}{-module}%
\TextAt{8}{172.247}{686.771}{Hom}%
\TextAt{27}{196.636}{684.977}{o}%
\TextAt{9}{201.03}{686.771}{\Character{40}}%
\TextAt{14}{205.583}{686.771}{a}%
\TextAt{21}{211.566}{684.977}{f}%
\TextAt{13}{217.011}{686.771}{\Character{59}}%
\TextAt{14}{222.255}{686.771}{o}%
\TextAt{9}{228.099}{686.771}{\Character{41}}%
\TextAt{8}{236.549}{686.771}{underlying}%
\TextAt{9}{296.139}{689.793}{\Character{20}}%
\TextAt{13}{294.758}{686.771}{L}%
\TextAt{9}{302.723}{686.771}{\Character{40}}%
\TextAt{14}{307.275}{686.771}{a}%
\TextAt{9}{313.258}{686.771}{\Character{41}}%
\TextAt{8}{321.708}{686.771}{is}%
\TextAt{8}{333.472}{686.771}{isomorphic,}%
\TextAt{8}{395.974}{686.771}{via}%
\TextAt{8}{415.162}{686.771}{the}%
\TextAt{8}{435.318}{686.771}{isomorphism}%
\TextAt{13}{173.912}{659.611}{g}%
\TextAt{9}{183.267}{659.611}{\Character{58}}%
\TextAt{14}{189.839}{659.611}{a}%
\TextAt{3}{195.822}{664.56}{\Character{0}}%
\TextAt{17}{202.409}{664.56}{1}%
\TextAt{21}{195.822}{656.282}{f}%
\NativeArrowAt{210.462}{659.611}{11.955}{19.26}{\longrightarrow}%
\TextAt{8}{233.044}{659.611}{Hom}%
\TextAt{27}{257.433}{657.818}{o}%
\TextAt{9}{261.827}{659.611}{\Character{40}}%
\TextAt{14}{266.379}{659.611}{a}%
\TextAt{21}{272.362}{657.818}{f}%
\TextAt{13}{277.807}{659.611}{\Character{59}}%
\TextAt{14}{283.051}{659.611}{o}%
\TextAt{9}{288.896}{659.611}{\Character{41}}%
\TextAt{13}{293.448}{659.611}{\Character{59}}%
\TextAt{13}{322.583}{659.611}{a}%
\NativeArrowAt{332.068}{659.611}{11.955}{19.26}{\longmapsto}%
\TextAt{20}{354.65}{669.295}{\Character{0}}%
\TextAt{13}{360.13}{659.611}{g}%
\TextAt{9}{366.164}{659.611}{\Character{40}}%
\TextAt{13}{370.716}{659.611}{a}%
\TextAt{9}{376.861}{659.611}{\Character{41}}%
\TextAt{9}{384.716}{659.611}{\Character{58}}%
\TextAt{13}{391.307}{659.611}{x}%
\NativeArrowAt{401.28}{659.611}{11.955}{11.955}{\mapsto}%
\TextAt{13}{416.556}{659.611}{ax}%
\TextAt{20}{429.353}{669.295}{\Character{1}}%
\TextAt{13}{434.832}{659.611}{\Character{59}}%
\TextAt{8}{79.226}{632.452}{to}%
\TextAt{8}{93.536}{632.452}{the}%
\TextAt{8}{113.693}{632.452}{fractional}%
\TextAt{8}{166.057}{632.452}{ideal}%
\TextAt{14}{194.02}{632.452}{a}%
\TextAt{3}{200.003}{637.4}{\Character{0}}%
\TextAt{17}{206.589}{637.4}{1}%
\TextAt{21}{200.003}{629.123}{f}%
\TextAt{8}{211.322}{632.452}{.}%
\TextAt{8}{219.77}{632.452}{For}%
\TextAt{8}{240.739}{632.452}{the}%
\TextAt{8}{260.896}{632.452}{induced}%
\TextAt{13}{304.46}{632.452}{K}%
\TextAt{25}{314.419}{630.659}{C}%
\TextAt{8}{320.674}{632.452}{-isomorphism}%
\TextAt{13}{232.254}{605.293}{g}%
\TextAt{25}{237.859}{603.5}{C}%
\TextAt{9}{247.434}{605.293}{\Character{58}}%
\TextAt{13}{254.007}{605.293}{K}%
\TextAt{25}{263.966}{603.5}{C}%
\NativeArrowAt{273.541}{605.293}{11.955}{19.26}{\longrightarrow}%
\TextAt{8}{296.123}{605.293}{Hom}%
\TextAt{21}{320.511}{603.5}{K}%
\TextAt{29}{327.653}{602.127}{C}%
\TextAt{9}{332.966}{605.293}{\Character{40}}%
\TextAt{13}{337.519}{605.293}{K}%
\TextAt{25}{347.478}{603.5}{C}%
\TextAt{13}{353.732}{605.293}{\Character{59}}%
\TextAt{13}{358.956}{605.293}{K}%
\TextAt{25}{368.935}{603.5}{C}%
\TextAt{9}{375.19}{605.293}{\Character{41}}%
\TextAt{8}{79.226}{578.134}{we}%
\TextAt{8}{96.454}{578.134}{have}%
\TextAt{13}{201.852}{561.088}{g}%
\TextAt{25}{207.457}{559.295}{C}%
\TextAt{9}{213.711}{561.088}{\Character{40}}%
\TextAt{13}{218.264}{561.088}{x}%
\TextAt{9}{224.916}{561.088}{\Character{41}\Character{40}}%
\TextAt{13}{234.02}{561.088}{y}%
\TextAt{9}{240.157}{561.088}{\Character{41}}%
\TextAt{9}{248.012}{561.088}{\Character{61}}%
\TextAt{13}{260.456}{561.088}{xy}%
\TextAt{9}{276.565}{561.088}{\Character{61}}%
\TextAt{13}{288.991}{561.088}{\Character{11}}%
\TextAt{3}{296.512}{566.024}{\Character{0}}%
\TextAt{17}{303.099}{566.024}{1}%
\TextAt{13}{307.831}{561.088}{\Character{11}xy}%
\TextAt{9}{331.463}{561.088}{\Character{61}}%
\TextAt{13}{343.888}{561.088}{\Character{11}}%
\TextAt{3}{351.41}{566.024}{\Character{0}}%
\TextAt{17}{357.996}{566.024}{1}%
\TextAt{16}{362.729}{561.088}{\Character{104}}%
\TextAt{13}{367.378}{561.088}{y}%
\TextAt{13}{373.499}{561.088}{\Character{59}}%
\Vector{q 0 G 0 g 0.478 w 306.759 -206.091 m 313.411 -206.091 l S Q}%
\TextAt{13}{378.759}{561.088}{x}%
\TextAt{16}{385.411}{561.088}{\Character{105}}%
\TextAt{21}{390.06}{559.228}{L}%
\TextAt{17}{395.819}{559.228}{\Character{40}}%
\TextAt{27}{399.112}{559.228}{a}%
\TextAt{17}{403.101}{559.228}{\Character{41}}%
\TextAt{13}{406.892}{561.088}{\Character{59}}%
\TextAt{8}{79.226}{538.356}{so}%
\TextAt{8}{93.596}{538.356}{that}%
\TextAt{13}{118.965}{538.356}{g}%
\TextAt{25}{124.57}{536.562}{C}%
\TextAt{9}{130.825}{538.356}{\Character{40}}%
\TextAt{13}{135.377}{538.356}{x}%
\TextAt{9}{142.029}{538.356}{\Character{41}}%
\TextAt{9}{149.884}{538.356}{\Character{61}}%
\TextAt{13}{162.328}{538.356}{\Character{11}}%
\TextAt{3}{169.85}{542.694}{\Character{0}}%
\TextAt{17}{176.436}{542.694}{1}%
\TextAt{9}{181.568}{540.747}{\Character{20}}%
\Vector{q 0 G 0 g 0.478 w 109.168 -228.823 m 115.82 -228.823 l S Q}%
\TextAt{13}{181.168}{538.356}{x}%
\TextAt{8}{187.82}{538.356}{,}%
\TextAt{8}{194.97}{538.356}{and}%
\TextAt{8}{217.732}{538.356}{thus}%
\TextAt{16}{184.201}{511.197}{\Character{104}}%
\TextAt{13}{188.85}{511.197}{g}%
\TextAt{25}{194.456}{509.403}{C}%
\TextAt{9}{200.71}{511.197}{\Character{40}}%
\TextAt{13}{205.262}{511.197}{x}%
\TextAt{9}{211.914}{511.197}{\Character{41}}%
\TextAt{13}{216.467}{511.197}{\Character{59}}%
\TextAt{13}{221.691}{511.197}{g}%
\TextAt{25}{227.316}{509.403}{C}%
\TextAt{9}{233.57}{511.197}{\Character{40}}%
\TextAt{13}{238.123}{511.197}{y}%
\TextAt{9}{244.259}{511.197}{\Character{41}}%
\TextAt{16}{248.812}{511.197}{\Character{105}}%
\TextAt{17}{254.459}{510.305}{\Character{20}}%
\TextAt{21}{253.461}{508.291}{L}%
\TextAt{17}{259.22}{508.291}{\Character{40}}%
\TextAt{27}{262.513}{508.291}{a}%
\TextAt{17}{266.502}{508.291}{\Character{41}}%
\TextAt{9}{273.614}{511.197}{\Character{61}}%
\TextAt{13}{286.04}{511.197}{\Character{11}}%
\TextAt{3}{293.561}{516.133}{\Character{0}}%
\TextAt{17}{300.148}{516.133}{2}%
\TextAt{16}{304.88}{511.197}{\Character{104}}%
\TextAt{9}{309.929}{513.587}{\Character{20}}%
\Vector{q 0 G 0 g 0.478 w 237.529 -255.982 m 244.182 -255.982 l S Q}%
\TextAt{13}{309.529}{511.197}{x\Character{59}}%
\TextAt{9}{321.568}{513.587}{\Character{20}}%
\Vector{q 0 G 0 g 0.478 w 249.426 -255.982 m 255.562 -255.982 l S Q}%
\TextAt{13}{321.426}{511.197}{y}%
\TextAt{16}{327.562}{511.197}{\Character{105}}%
\TextAt{17}{333.209}{510.305}{\Character{20}}%
\TextAt{21}{332.212}{508.291}{L}%
\TextAt{17}{337.971}{508.291}{\Character{40}}%
\TextAt{27}{341.264}{508.291}{a}%
\TextAt{17}{345.253}{508.291}{\Character{41}}%
\TextAt{9}{352.365}{511.197}{\Character{61}}%
\TextAt{13}{364.79}{511.197}{\Character{11}}%
\TextAt{3}{372.312}{516.133}{\Character{0}}%
\TextAt{17}{378.899}{516.133}{2}%
\Vector{q 0 G 0 g 0.478 w 311.631 -252.163 m 355.795 -252.163 l S Q}%
\TextAt{16}{383.631}{511.197}{\Character{104}}%
\Vector{q 0 G 0 g 0.478 w 316.28 -255.982 m 322.932 -255.982 l S Q}%
\TextAt{13}{388.28}{511.197}{x\Character{59}}%
\Vector{q 0 G 0 g 0.478 w 328.176 -255.982 m 334.313 -255.982 l S Q}%
\TextAt{13}{400.176}{511.197}{y}%
\TextAt{16}{406.313}{511.197}{\Character{105}}%
\TextAt{21}{410.962}{509.337}{L}%
\TextAt{17}{416.721}{509.337}{\Character{40}}%
\TextAt{27}{420.015}{509.337}{a}%
\TextAt{17}{424.003}{509.337}{\Character{41}}%
\TextAt{9}{273.614}{491.162}{\Character{61}}%
\TextAt{13}{286.04}{491.162}{\Character{11}}%
\TextAt{3}{293.561}{496.098}{\Character{0}}%
\TextAt{17}{300.148}{496.098}{1}%
\TextAt{13}{304.88}{491.162}{x}%
\Vector{q 0 G 0 g 0.478 w 239.532 -276.017 m 245.669 -276.017 l S Q}%
\TextAt{13}{311.532}{491.162}{y}%
\TextAt{9}{320.99}{491.162}{\Character{61}}%
\TextAt{16}{333.415}{491.162}{\Character{104}}%
\TextAt{13}{338.065}{491.162}{x\Character{59}}%
\TextAt{13}{349.937}{491.162}{y}%
\TextAt{16}{356.097}{491.162}{\Character{105}}%
\TextAt{21}{360.747}{489.158}{L}%
\TextAt{17}{366.506}{489.158}{\Character{40}}%
\TextAt{27}{369.799}{489.158}{a}%
\TextAt{23}{373.788}{491.427}{\Character{0}}%
\TextAt{24}{379.544}{491.427}{1}%
\TextAt{17}{383.695}{489.158}{\Character{41}}%
\TextAt{13}{387.486}{491.162}{\Character{58}}%
\TextAt{8}{79.226}{464.003}{Thus}%
\TextAt{13}{108.875}{464.003}{g}%
\TextAt{8}{118.807}{464.003}{gives}%
\TextAt{8}{147.487}{464.003}{an}%
\TextAt{8}{163.746}{464.003}{isometry}%
\TextAt{9}{212.996}{467.025}{\Character{20}}%
\TextAt{13}{211.615}{464.003}{L}%
\TextAt{9}{219.579}{464.003}{\Character{40}}%
\TextAt{14}{224.131}{464.003}{a}%
\TextAt{9}{230.115}{464.003}{\Character{41}}%
\NativeMathAt{237.988}{464.002}{11.955}{\cong}%
\TextAt{13}{250.607}{464.003}{L}%
\TextAt{9}{258.571}{464.003}{\Character{40}}%
\TextAt{14}{263.124}{464.003}{a}%
\TextAt{3}{269.107}{468.341}{\Character{0}}%
\TextAt{17}{275.693}{468.341}{1}%
\TextAt{9}{280.426}{464.003}{\Character{41}}%
\TextAt{8}{284.978}{464.003}{.}%
\Vector{q 0 G 0 g 0.399 w 452.334 -309.997 m 452.334 -301.928 l S Q}%
\Vector{q 0 G 0 g 0.399 w 452.534 -302.127 m 459.707 -302.127 l S Q}%
\Vector{q 0 G 0 g 0.399 w 452.534 -309.798 m 459.707 -309.798 l S Q}%
\Vector{q 0 G 0 g 0.399 w 459.906 -309.997 m 459.906 -301.928 l S Q}%
\TextAt{12}{79.226}{430.076}{Definition}%
\TextAt{12}{142.015}{430.076}{(4.6).}%
\TextAt{8}{177.33}{430.076}{A}%
\TextAt{12}{190.003}{430.076}{short}%
\TextAt{12}{224.732}{430.076}{exact}%
\TextAt{12}{260.203}{430.076}{sequence}%
\TextAt{9}{232.806}{402.917}{0}%
\NativeArrowAt{241.98}{402.917}{11.955}{19.26}{\longrightarrow}%
\TextAt{13}{264.562}{402.917}{M}%
\TextAt{3}{277.135}{407.853}{\Character{48}}%
\TextAt{21}{285.839}{409.694}{\Character{11}}%
\NativeArrowAt{283.251}{402.917}{11.955}{12.027}{\longrightarrow}%
\TextAt{13}{298.598}{402.917}{M}%
\TextAt{21}{317.183}{410.181}{\Character{12}}%
\NativeArrowAt{314.492}{402.917}{11.955}{11.955}{\longrightarrow}%
\TextAt{13}{329.768}{402.917}{M}%
\TextAt{3}{342.342}{407.853}{\Character{48}\Character{48}}%
\NativeArrowAt{350.755}{402.917}{11.955}{19.26}{\longrightarrow}%
\TextAt{9}{373.337}{402.917}{0}%
\TextAt{8}{79.226}{375.758}{of}%
\TextAt{8}{92.08}{375.758}{metrized}%
\TextAt{14}{139.724}{375.758}{o}%
\TextAt{8}{145.568}{375.758}{-modules}%
\TextAt{8}{194.897}{375.758}{is}%
\TextAt{8}{206.185}{375.758}{by}%
\TextAt{8}{221.968}{375.758}{definition}%
\TextAt{8}{273.521}{375.758}{a}%
\TextAt{8}{282.8}{375.758}{short}%
\TextAt{8}{312.307}{375.758}{exact}%
\TextAt{8}{342.723}{375.758}{sequence}%
\TextAt{8}{390.749}{375.758}{of}%
\TextAt{8}{403.603}{375.758}{the}%
\TextAt{8}{423.284}{375.758}{underlying}%
\TextAt{14}{481.017}{375.758}{o}%
\TextAt{8}{486.862}{375.758}{-modules}%
\TextAt{8}{79.226}{358.712}{which}%
\TextAt{19}{112.713}{358.712}{splits}%
\TextAt{19}{142.899}{358.712}{isometrically}%
\TextAt{8}{207.912}{358.712}{,}%
\TextAt{8}{215.061}{358.712}{i.e.,}%
\TextAt{8}{237.166}{358.712}{in}%
\TextAt{8}{250.819}{358.712}{the}%
\TextAt{8}{270.975}{358.712}{sequence}%
\TextAt{9}{223.065}{331.553}{0}%
\NativeArrowAt{232.239}{331.553}{11.955}{19.26}{\longrightarrow}%
\TextAt{13}{254.821}{331.553}{M}%
\TextAt{3}{267.395}{336.489}{\Character{48}}%
\TextAt{25}{266.146}{328.597}{C}%
\TextAt{21}{278.309}{338.64}{\Character{11}}%
\TextAt{29}{283.748}{337.267}{C}%
\NativeArrowAt{275.721}{331.553}{11.955}{16.833}{\longrightarrow}%
\TextAt{13}{295.883}{331.553}{M}%
\TextAt{25}{307.208}{329.76}{C}%
\TextAt{21}{319.37}{338.817}{\Character{12}}%
\TextAt{29}{324.122}{337.445}{C}%
\NativeArrowAt{316.783}{331.553}{11.955}{16.151}{\longrightarrow}%
\TextAt{13}{336.257}{331.553}{M}%
\TextAt{3}{348.831}{336.489}{\Character{48}\Character{48}}%
\TextAt{25}{347.582}{328.597}{C}%
\NativeArrowAt{357.243}{331.553}{11.955}{19.26}{\longrightarrow}%
\TextAt{9}{379.825}{331.553}{0}%
\TextAt{13}{385.678}{331.553}{\Character{59}}%
\TextAt{13}{79.226}{304.394}{M}%
\TextAt{3}{91.8}{308.732}{\Character{48}}%
\TextAt{25}{90.551}{301.329}{C}%
\TextAt{8}{102.145}{304.394}{is}%
\TextAt{8}{115.352}{304.394}{mapped}%
\TextAt{8}{161.34}{304.394}{isometrically}%
\TextAt{8}{231.465}{304.394}{onto}%
\TextAt{13}{259.257}{304.394}{\Character{11}}%
\TextAt{25}{266.703}{302.601}{C}%
\TextAt{13}{272.957}{304.394}{M}%
\TextAt{3}{285.531}{308.732}{\Character{48}}%
\TextAt{25}{284.282}{301.329}{C}%
\TextAt{8}{290.536}{304.394}{,}%
\TextAt{8}{299.489}{304.394}{and}%
\TextAt{8}{323.694}{304.394}{the}%
\TextAt{8}{345.294}{304.394}{orthogonal}%
\TextAt{8}{405.293}{304.394}{complement}%
\TextAt{9}{472.094}{304.394}{\Character{40}}%
\TextAt{13}{476.647}{304.394}{\Character{11}}%
\TextAt{25}{484.092}{302.601}{C}%
\TextAt{13}{490.347}{304.394}{M}%
\TextAt{3}{502.92}{308.732}{\Character{48}}%
\TextAt{25}{501.672}{301.329}{C}%
\TextAt{9}{507.926}{304.394}{\Character{41}}%
\TextAt{3}{512.478}{308.732}{\Character{63}}%
\TextAt{8}{524.903}{304.394}{is}%
\TextAt{8}{79.226}{287.348}{mapped}%
\TextAt{8}{123.771}{287.348}{isometrically}%
\TextAt{8}{192.454}{287.348}{onto}%
\TextAt{13}{218.803}{287.348}{M}%
\TextAt{3}{231.376}{291.686}{\Character{48}\Character{48}}%
\TextAt{25}{230.128}{284.283}{C}%
\TextAt{8}{236.468}{287.348}{.}%
\TextAt{8}{97.159}{261.862}{The}%
\TextAt{8}{121.352}{261.862}{homomorphisms}%
\TextAt{13}{208.777}{261.862}{\Character{11}\Character{59}}%
\TextAt{13}{221.509}{261.862}{\Character{12}}%
\TextAt{8}{232.851}{261.862}{in}%
\TextAt{8}{246.643}{261.862}{a}%
\TextAt{8}{256.538}{261.862}{short}%
\TextAt{8}{286.661}{261.862}{exact}%
\TextAt{8}{317.693}{261.862}{sequence}%
\TextAt{8}{366.335}{261.862}{of}%
\TextAt{8}{379.805}{261.862}{metrized}%
\TextAt{14}{428.064}{261.862}{o}%
\TextAt{8}{433.908}{261.862}{-modules}%
\TextAt{8}{483.853}{261.862}{are}%
\TextAt{8}{503.503}{261.862}{called}%
\TextAt{8}{79.226}{244.816}{an}%
\TextAt{12}{95.485}{244.816}{admissible}%
\TextAt{12}{160.641}{244.816}{monomorphism}%
\TextAt{8}{251.44}{244.816}{,}%
\TextAt{8}{258.59}{244.816}{resp.}%
\TextAt{12}{287.912}{244.816}{epimorphism}%
\TextAt{8}{363.911}{244.816}{.}%
\TextAt{8}{97.159}{224.781}{To}%
\TextAt{8}{115.507}{224.781}{each}%
\TextAt{8}{142.966}{224.781}{projective}%
\TextAt{8}{198.065}{224.781}{metrized}%
\TextAt{14}{247.305}{224.781}{o}%
\TextAt{8}{253.15}{224.781}{-module}%
\TextAt{13}{299.461}{224.781}{M}%
\TextAt{8}{317.053}{224.781}{is}%
\TextAt{8}{329.939}{224.781}{associated}%
\TextAt{8}{386.795}{224.781}{its}%
\TextAt{12}{404.235}{224.781}{determinant}%
\TextAt{8}{481.032}{224.781}{det}%
\TextAt{13}{499.284}{224.781}{M}%
\TextAt{8}{511.857}{224.781}{,}%
\TextAt{8}{520.408}{224.781}{an}%
\TextAt{8}{79.226}{207.735}{invertible}%
\TextAt{8}{130.932}{207.735}{metrized}%
\TextAt{14}{179.052}{207.735}{o}%
\TextAt{8}{184.896}{207.735}{-module.}%
\TextAt{8}{234.638}{207.735}{The}%
\TextAt{8}{258.691}{207.735}{determinant}%
\TextAt{8}{324.708}{207.735}{is}%
\TextAt{8}{336.472}{207.735}{the}%
\TextAt{8}{356.628}{207.735}{highest}%
\TextAt{8}{397.013}{207.735}{exterior}%
\TextAt{8}{440.266}{207.735}{power}%
\TextAt{8}{474.41}{207.735}{of}%
\TextAt{13}{487.74}{207.735}{M}%
\TextAt{8}{500.314}{207.735}{,}%
\TextAt{8}{507.463}{207.735}{i.e.,}%
\TextAt{8}{225.335}{174.604}{det}%
\TextAt{13}{243.587}{174.604}{M}%
\TextAt{9}{259.481}{174.604}{\Character{61}}%
\TextAt{21}{275.98}{189.548}{n}%
\TextAt{20}{271.907}{185.961}{\Character{94}}%
\TextAt{13}{287.183}{174.604}{M}%
\TextAt{13}{299.102}{174.604}{\Character{59}}%
\TextAt{13}{328.237}{174.604}{n}%
\TextAt{9}{338.569}{174.604}{\Character{61}}%
\TextAt{8}{350.995}{174.604}{rk}%
\TextAt{9}{361.73}{174.604}{\Character{40}}%
\TextAt{13}{366.283}{174.604}{M}%
\TextAt{9}{378.856}{174.604}{\Character{41}}%
\TextAt{13}{383.409}{174.604}{\Character{58}}%
\TextAt{12}{79.226}{134.484}{Proposition}%
\TextAt{12}{152.124}{134.484}{(4.7).}%
\TextAt{8}{187.601}{134.484}{If}%
\TextAt{9}{199.455}{134.484}{0}%
\NativeMathAt{208.904}{134.484}{11.955}{\rightarrow}%
\TextAt{13}{224.456}{134.484}{M}%
\TextAt{3}{237.03}{138.823}{\Character{48}}%
\TextAt{21}{246.009}{141.261}{\Character{11}}%
\NativeArrowAt{243.421}{134.484}{11.955}{12.027}{\longrightarrow}%
\TextAt{13}{259.043}{134.484}{M}%
\TextAt{21}{277.905}{141.748}{\Character{12}}%
\NativeArrowAt{275.213}{134.484}{11.955}{11.955}{\longrightarrow}%
\TextAt{13}{290.765}{134.484}{M}%
\TextAt{3}{303.339}{138.823}{\Character{48}\Character{48}}%
\NativeMathAt{312.027}{134.484}{11.955}{\rightarrow}%
\TextAt{9}{327.579}{134.484}{0}%
\TextAt{8}{337.491}{134.484}{is}%
\TextAt{8}{349.417}{134.484}{a}%
\TextAt{8}{359.334}{134.484}{short}%
\TextAt{8}{389.48}{134.484}{exact}%
\TextAt{8}{420.534}{134.484}{sequence}%
\TextAt{8}{469.197}{134.484}{of}%
\TextAt{8}{482.689}{134.484}{projective}%
\TextAt{8}{79.226}{117.438}{metrized}%
\TextAt{14}{127.346}{117.438}{o}%
\TextAt{8}{133.19}{117.438}{-modules,}%
\TextAt{8}{186.247}{117.438}{we}%
\TextAt{8}{203.474}{117.438}{have}%
\TextAt{8}{230.469}{117.438}{a}%
\TextAt{8}{240.225}{117.438}{canonical}%
\TextAt{8}{291.608}{117.438}{isometry}%
\TextAt{8}{237.043}{90.279}{det}%
\TextAt{13}{255.294}{90.279}{M}%
\TextAt{3}{267.868}{95.215}{\Character{48}}%
\TextAt{16}{273.32}{90.279}{\Character{10}}%
\TextAt{27}{282.618}{88.486}{o}%
\TextAt{8}{289.669}{90.279}{det}%
\TextAt{13}{307.921}{90.279}{M}%
\TextAt{3}{320.494}{95.215}{\Character{48}\Character{48}}%
\NativeMathAt{328.907}{90.279}{11.955}{\cong}%
\TextAt{8}{341.526}{90.279}{det}%
\TextAt{13}{359.778}{90.279}{M}%
\TextAt{13}{371.697}{90.279}{\Character{58}}%
\TextAt{8}{297.211}{43.033}{212}%
\Destination{78.23}{809.815}{nav.page.212}%
\Destination{79.36}{776.155}{nav.print.231}%
\Destination{79.36}{444.055}{nav.number.3.4.6}%
\Destination{79.36}{148.455}{nav.number.3.4.7}%
}

% Page 225 (body page 213)
\Page{%
\TextAt{2}{79.226}{796.943}{III.}%
\TextAt{2}{98.317}{796.943}{Riemann–Roch}%
\TextAt{2}{174.375}{796.943}{Theory}%
\Vector{q 0 G 0 g 0.299 w 7.226 18.729 m 460.769 18.729 l S Q}%
\TextAt{19}{79.226}{761.124}{Proof.}%
\TextAt{8}{116}{761.124}{Let}%
\TextAt{13}{136.97}{761.124}{n}%
\TextAt{3}{143.957}{765.462}{\Character{48}}%
\TextAt{9}{150.073}{761.124}{\Character{61}}%
\TextAt{8}{162.499}{761.124}{rk}%
\TextAt{9}{173.234}{761.124}{\Character{40}}%
\TextAt{13}{177.787}{761.124}{M}%
\TextAt{3}{190.36}{765.462}{\Character{48}}%
\TextAt{9}{193.155}{761.124}{\Character{41}}%
\TextAt{8}{201.605}{761.124}{and}%
\TextAt{13}{224.368}{761.124}{n}%
\TextAt{3}{231.355}{765.462}{\Character{48}\Character{48}}%
\TextAt{9}{239.768}{761.124}{\Character{61}}%
\TextAt{8}{252.194}{761.124}{rk}%
\TextAt{9}{262.929}{761.124}{\Character{40}}%
\TextAt{13}{267.482}{761.124}{M}%
\TextAt{3}{280.055}{765.462}{\Character{48}\Character{48}}%
\TextAt{9}{285.147}{761.124}{\Character{41}}%
\TextAt{8}{289.7}{761.124}{.}%
\TextAt{8}{298.148}{761.124}{We}%
\TextAt{8}{318.292}{761.124}{obtain}%
\TextAt{8}{354.72}{761.124}{an}%
\TextAt{8}{370.979}{761.124}{isomorphism}%
\TextAt{13}{228.09}{735.171}{\Character{20}}%
\TextAt{9}{238.15}{735.171}{\Character{58}}%
\TextAt{8}{244.722}{735.171}{det}%
\TextAt{13}{262.974}{735.171}{M}%
\TextAt{3}{275.547}{740.107}{\Character{48}}%
\TextAt{16}{280.999}{735.171}{\Character{10}}%
\TextAt{27}{290.297}{733.378}{o}%
\TextAt{8}{297.348}{735.171}{det}%
\TextAt{13}{315.6}{735.171}{M}%
\TextAt{3}{328.174}{740.107}{\Character{48}\Character{48}}%
\NativeMathAt{339.174}{741.948}{7.97}{\sim}%
\NativeArrowAt{336.586}{735.171}{11.955}{13.175}{\longrightarrow}%
\TextAt{8}{353.08}{735.171}{det}%
\TextAt{13}{371.332}{735.171}{M}%
\TextAt{8}{79.226}{709.218}{of}%
\TextAt{8}{92.556}{709.218}{projective}%
\TextAt{14}{146.534}{709.218}{o}%
\TextAt{8}{152.378}{709.218}{-modules}%
\TextAt{8}{202.183}{709.218}{of}%
\TextAt{8}{215.513}{709.218}{rank}%
\TextAt{9}{242.508}{709.218}{1}%
\TextAt{8}{252.258}{709.218}{by}%
\TextAt{8}{268.517}{709.218}{mapping}%
\TextAt{9}{89.189}{680.277}{\Character{40}}%
\TextAt{13}{93.741}{680.277}{m}%
\TextAt{3}{103.98}{685.213}{\Character{48}}%
\TextAt{17}{103.98}{677.321}{1}%
\TextAt{16}{111.369}{680.277}{\Character{94}}%
\TextAt{16}{121.985}{680.277}{\Character{1}}%
\TextAt{16}{127.281}{680.277}{\Character{1}}%
\TextAt{16}{132.59}{680.277}{\Character{1}}%
\TextAt{16}{138.555}{680.277}{\Character{94}}%
\TextAt{13}{149.227}{680.277}{m}%
\TextAt{3}{159.467}{685.213}{\Character{48}}%
\TextAt{21}{159.467}{677.321}{n}%
\TextAt{23}{164.605}{679.59}{\Character{48}}%
\TextAt{9}{167.787}{680.277}{\Character{41}}%
\TextAt{16}{174.996}{680.277}{\Character{10}}%
\TextAt{9}{186.951}{680.277}{\Character{40}}%
\TextAt{13}{191.504}{680.277}{m}%
\TextAt{3}{201.743}{685.213}{\Character{48}\Character{48}}%
\TextAt{17}{201.743}{677.321}{1}%
\TextAt{16}{209.492}{680.277}{\Character{94}}%
\TextAt{16}{220.108}{680.277}{\Character{1}}%
\TextAt{16}{225.404}{680.277}{\Character{1}}%
\TextAt{16}{230.713}{680.277}{\Character{1}}%
\TextAt{16}{236.678}{680.277}{\Character{94}}%
\TextAt{13}{247.35}{680.277}{m}%
\TextAt{3}{257.589}{685.213}{\Character{48}\Character{48}}%
\TextAt{21}{257.589}{677.321}{n}%
\TextAt{23}{262.727}{679.59}{\Character{48}\Character{48}}%
\TextAt{9}{268.096}{680.277}{\Character{41}}%
\NativeArrowAt{318.008}{660.242}{11.955}{19.26}{\longmapsto}%
\TextAt{13}{340.59}{660.242}{\Character{11}m}%
\TextAt{3}{358.351}{665.178}{\Character{48}}%
\TextAt{17}{358.351}{657.286}{1}%
\TextAt{16}{365.74}{660.242}{\Character{94}}%
\TextAt{16}{376.356}{660.242}{\Character{1}}%
\TextAt{16}{381.652}{660.242}{\Character{1}}%
\TextAt{16}{386.961}{660.242}{\Character{1}}%
\TextAt{16}{392.926}{660.242}{\Character{94}}%
\TextAt{13}{403.598}{660.242}{\Character{11}m}%
\TextAt{3}{421.359}{665.178}{\Character{48}}%
\TextAt{21}{421.359}{657.286}{n}%
\TextAt{23}{426.497}{659.556}{\Character{48}}%
\TextAt{16}{432.336}{660.242}{\Character{94}}%
\TextAt{20}{444.762}{660.242}{\Character{101}}%
\TextAt{13}{442.963}{660.242}{m}%
\TextAt{3}{453.202}{665.178}{\Character{48}\Character{48}}%
\TextAt{17}{453.202}{657.286}{1}%
\TextAt{16}{460.951}{660.242}{\Character{94}}%
\TextAt{16}{471.567}{660.242}{\Character{1}}%
\TextAt{16}{476.863}{660.242}{\Character{1}}%
\TextAt{16}{482.172}{660.242}{\Character{1}}%
\TextAt{16}{488.137}{660.242}{\Character{94}}%
\TextAt{20}{500.608}{660.242}{\Character{101}}%
\TextAt{13}{498.809}{660.242}{m}%
\TextAt{3}{509.048}{665.178}{\Character{48}\Character{48}}%
\TextAt{21}{509.048}{657.286}{n}%
\TextAt{23}{514.186}{659.556}{\Character{48}\Character{48}}%
\TextAt{13}{519.555}{660.242}{\Character{59}}%
\TextAt{8}{79.226}{634.289}{where}%
\TextAt{20}{114.282}{634.289}{\Character{101}}%
\TextAt{13}{112.483}{634.289}{m}%
\TextAt{3}{122.722}{638.627}{\Character{48}\Character{48}}%
\TextAt{17}{122.722}{631.334}{1}%
\TextAt{13}{127.814}{634.289}{\Character{59}}%
\TextAt{13}{133.038}{634.289}{\Character{58}}%
\TextAt{13}{138.275}{634.289}{\Character{58}}%
\TextAt{13}{143.511}{634.289}{\Character{58}}%
\TextAt{13}{148.736}{634.289}{\Character{59}}%
\TextAt{20}{155.834}{634.289}{\Character{101}}%
\TextAt{13}{154.035}{634.289}{m}%
\TextAt{3}{164.274}{638.627}{\Character{48}\Character{48}}%
\TextAt{21}{164.274}{631.125}{n}%
\TextAt{23}{169.412}{633.394}{\Character{48}\Character{48}}%
\TextAt{8}{178.126}{634.289}{are}%
\TextAt{8}{197.084}{634.289}{preimages}%
\TextAt{8}{251.227}{634.289}{of}%
\TextAt{13}{264.004}{634.289}{m}%
\TextAt{3}{274.244}{638.627}{\Character{48}\Character{48}}%
\TextAt{17}{274.244}{631.334}{1}%
\TextAt{13}{279.336}{634.289}{\Character{59}}%
\TextAt{13}{284.56}{634.289}{\Character{58}}%
\TextAt{13}{289.797}{634.289}{\Character{58}}%
\TextAt{13}{295.021}{634.289}{\Character{58}}%
\TextAt{13}{300.258}{634.289}{\Character{59}}%
\TextAt{13}{305.494}{634.289}{m}%
\TextAt{3}{315.796}{638.627}{\Character{48}\Character{48}}%
\TextAt{21}{315.796}{631.125}{n}%
\TextAt{23}{320.934}{633.394}{\Character{48}\Character{48}}%
\TextAt{8}{329.648}{634.289}{under}%
\TextAt{13}{362.259}{634.289}{\Character{12}}%
\TextAt{9}{372.851}{634.289}{\Character{58}}%
\TextAt{13}{379.423}{634.289}{M}%
\NativeMathAt{395.318}{634.289}{11.955}{\rightarrow}%
\TextAt{13}{410.594}{634.289}{M}%
\TextAt{3}{423.167}{638.627}{\Character{48}\Character{48}}%
\TextAt{8}{428.259}{634.289}{.}%
\TextAt{8}{436.524}{634.289}{This}%
\TextAt{8}{462.691}{634.289}{mapping}%
\TextAt{8}{510.27}{634.289}{does}%
\TextAt{8}{79.226}{617.243}{not}%
\TextAt{8}{100.3}{617.243}{depend}%
\TextAt{8}{141.195}{617.243}{on}%
\TextAt{8}{157.713}{617.243}{the}%
\TextAt{8}{178.129}{617.243}{choice}%
\TextAt{8}{213.179}{617.243}{of}%
\TextAt{8}{226.768}{617.243}{the}%
\TextAt{8}{247.184}{617.243}{preimages,}%
\TextAt{8}{305.455}{617.243}{for}%
\TextAt{8}{323.6}{617.243}{if,}%
\TextAt{8}{337.9}{617.243}{say,}%
\TextAt{20}{362.523}{617.243}{\Character{101}}%
\TextAt{13}{360.724}{617.243}{m}%
\TextAt{3}{370.963}{621.582}{\Character{48}\Character{48}}%
\TextAt{17}{370.963}{614.288}{1}%
\TextAt{9}{378.889}{617.243}{\Character{43}}%
\TextAt{13}{390.827}{617.243}{\Character{11}m}%
\TextAt{3}{408.588}{621.582}{\Character{48}}%
\TextAt{21}{408.588}{614.079}{n}%
\TextAt{23}{413.726}{616.348}{\Character{48}}%
\TextAt{17}{416.411}{614.079}{\Character{43}1}%
\TextAt{8}{427.729}{617.243}{,}%
\TextAt{8}{435.203}{617.243}{where}%
\TextAt{13}{469.272}{617.243}{m}%
\TextAt{3}{479.511}{621.582}{\Character{48}}%
\TextAt{21}{479.511}{614.079}{n}%
\TextAt{23}{484.649}{616.348}{\Character{48}}%
\TextAt{17}{487.334}{614.079}{\Character{43}1}%
\NativeMathAt{502.416}{617.243}{11.955}{\in}%
\TextAt{13}{514.149}{617.243}{M}%
\TextAt{3}{526.723}{621.582}{\Character{48}}%
\TextAt{8}{529.518}{617.243}{,}%
\TextAt{8}{79.226}{600.197}{is}%
\TextAt{8}{90.99}{600.197}{another}%
\TextAt{8}{133.921}{600.197}{preimage}%
\TextAt{8}{184.001}{600.197}{of}%
\TextAt{13}{197.331}{600.197}{m}%
\TextAt{3}{207.571}{604.536}{\Character{48}\Character{48}}%
\TextAt{17}{207.571}{597.242}{1}%
\TextAt{8}{212.663}{600.197}{,}%
\TextAt{8}{219.812}{600.197}{then}%
\TextAt{13}{89.189}{571.256}{\Character{11}m}%
\TextAt{3}{106.95}{576.192}{\Character{48}}%
\TextAt{17}{106.95}{568.3}{1}%
\TextAt{16}{114.339}{571.256}{\Character{94}}%
\TextAt{16}{124.955}{571.256}{\Character{1}}%
\TextAt{16}{130.251}{571.256}{\Character{1}}%
\TextAt{16}{135.56}{571.256}{\Character{1}}%
\TextAt{16}{141.525}{571.256}{\Character{94}}%
\TextAt{13}{152.197}{571.256}{\Character{11}m}%
\TextAt{3}{169.958}{576.192}{\Character{48}}%
\TextAt{21}{169.958}{568.3}{n}%
\TextAt{23}{175.096}{570.569}{\Character{48}}%
\TextAt{16}{180.935}{571.256}{\Character{94}}%
\TextAt{9}{191.562}{571.256}{\Character{40}}%
\TextAt{20}{197.913}{571.256}{\Character{101}}%
\TextAt{13}{196.114}{571.256}{m}%
\TextAt{3}{206.353}{576.192}{\Character{48}\Character{48}}%
\TextAt{17}{206.353}{568.3}{1}%
\TextAt{9}{214.102}{571.256}{\Character{43}}%
\TextAt{13}{225.863}{571.256}{\Character{11}m}%
\TextAt{3}{243.624}{576.192}{\Character{48}}%
\TextAt{21}{243.624}{568.3}{n}%
\TextAt{23}{248.763}{570.569}{\Character{48}}%
\TextAt{17}{251.447}{568.3}{\Character{43}1}%
\TextAt{9}{262.766}{571.256}{\Character{41}}%
\TextAt{16}{269.975}{571.256}{\Character{94}}%
\TextAt{20}{282.4}{571.256}{\Character{101}}%
\TextAt{13}{280.602}{571.256}{m}%
\TextAt{3}{290.841}{576.192}{\Character{48}\Character{48}}%
\TextAt{17}{290.841}{568.3}{2}%
\TextAt{16}{298.59}{571.256}{\Character{94}}%
\TextAt{16}{309.206}{571.256}{\Character{1}}%
\TextAt{16}{314.502}{571.256}{\Character{1}}%
\TextAt{16}{319.81}{571.256}{\Character{1}}%
\TextAt{16}{325.776}{571.256}{\Character{94}}%
\TextAt{20}{338.246}{571.256}{\Character{101}}%
\TextAt{13}{336.448}{571.256}{m}%
\TextAt{3}{346.687}{576.192}{\Character{48}\Character{48}}%
\TextAt{21}{346.687}{568.3}{n}%
\TextAt{23}{351.825}{570.569}{\Character{48}\Character{48}}%
\TextAt{9}{331.416}{551.221}{\Character{61}}%
\TextAt{13}{343.842}{551.221}{\Character{11}m}%
\TextAt{3}{361.603}{556.157}{\Character{48}}%
\TextAt{17}{361.603}{548.265}{1}%
\TextAt{16}{368.992}{551.221}{\Character{94}}%
\TextAt{16}{379.608}{551.221}{\Character{1}}%
\TextAt{16}{384.904}{551.221}{\Character{1}}%
\TextAt{16}{390.213}{551.221}{\Character{1}}%
\TextAt{16}{396.178}{551.221}{\Character{94}}%
\TextAt{13}{406.85}{551.221}{\Character{11}m}%
\TextAt{3}{424.611}{556.157}{\Character{48}}%
\TextAt{21}{424.611}{548.265}{n}%
\TextAt{23}{429.749}{550.535}{\Character{48}}%
\TextAt{16}{435.588}{551.221}{\Character{94}}%
\TextAt{20}{448.014}{551.221}{\Character{101}}%
\TextAt{13}{446.215}{551.221}{m}%
\TextAt{3}{456.454}{556.157}{\Character{48}\Character{48}}%
\TextAt{17}{456.454}{548.265}{1}%
\TextAt{16}{464.203}{551.221}{\Character{94}}%
\TextAt{16}{474.819}{551.221}{\Character{1}}%
\TextAt{16}{480.115}{551.221}{\Character{1}}%
\TextAt{16}{485.423}{551.221}{\Character{1}}%
\TextAt{16}{491.389}{551.221}{\Character{94}}%
\TextAt{20}{503.859}{551.221}{\Character{101}}%
\TextAt{13}{502.061}{551.221}{m}%
\TextAt{3}{512.3}{556.157}{\Character{48}\Character{48}}%
\TextAt{21}{512.3}{548.265}{n}%
\TextAt{23}{517.438}{550.535}{\Character{48}\Character{48}}%
\TextAt{8}{79.226}{525.268}{since}%
\TextAt{13}{108.236}{525.268}{\Character{11}m}%
\TextAt{3}{125.997}{529.606}{\Character{48}}%
\TextAt{17}{125.997}{522.312}{1}%
\TextAt{16}{133.619}{525.268}{\Character{94}}%
\TextAt{16}{144.462}{525.268}{\Character{1}}%
\TextAt{16}{149.771}{525.268}{\Character{1}}%
\TextAt{16}{155.079}{525.268}{\Character{1}}%
\TextAt{16}{161.271}{525.268}{\Character{94}}%
\TextAt{13}{172.176}{525.268}{\Character{11}m}%
\TextAt{3}{189.937}{529.606}{\Character{48}}%
\TextAt{21}{189.937}{522.104}{n}%
\TextAt{23}{195.075}{524.373}{\Character{48}}%
\TextAt{16}{201.147}{525.268}{\Character{94}}%
\TextAt{13}{212.007}{525.268}{\Character{11}m}%
\TextAt{3}{229.768}{529.606}{\Character{48}}%
\TextAt{21}{229.768}{522.104}{n}%
\TextAt{23}{234.906}{524.373}{\Character{48}}%
\TextAt{17}{237.591}{522.104}{\Character{43}1}%
\TextAt{9}{252.813}{525.268}{\Character{61}}%
\TextAt{9}{265.808}{525.268}{0}%
\TextAt{8}{271.673}{525.268}{.}%
\TextAt{8}{281.147}{525.268}{We}%
\TextAt{8}{301.633}{525.268}{show}%
\TextAt{8}{330.978}{525.268}{that}%
\TextAt{8}{356.689}{525.268}{the}%
\TextAt{14}{377.187}{525.268}{o}%
\TextAt{8}{383.031}{525.268}{-module}%
\TextAt{8}{428.563}{525.268}{isomorphism}%
\TextAt{13}{497.325}{525.268}{\Character{20}}%
\TextAt{8}{508.302}{525.268}{is}%
\TextAt{8}{520.408}{525.268}{an}%
\TextAt{8}{79.226}{508.222}{isometry.}%
\TextAt{8}{130.665}{508.222}{According}%
\TextAt{8}{185.946}{508.222}{to}%
\TextAt{8}{200.256}{508.222}{the}%
\TextAt{8}{220.413}{508.222}{rules}%
\TextAt{8}{248.436}{508.222}{of}%
\TextAt{8}{261.766}{508.222}{multilinear}%
\TextAt{8}{321.278}{508.222}{algebra}%
\TextAt{8}{362.261}{508.222}{it}%
\TextAt{8}{373.965}{508.222}{induces}%
\TextAt{8}{415.641}{508.222}{an}%
\TextAt{8}{431.9}{508.222}{isomorphism}%
\TextAt{13}{220.452}{482.269}{\Character{20}}%
\TextAt{9}{230.511}{482.269}{\Character{58}}%
\TextAt{8}{237.084}{482.269}{det}%
\TextAt{13}{255.335}{482.269}{M}%
\TextAt{3}{267.909}{487.205}{\Character{48}}%
\TextAt{25}{266.66}{479.314}{C}%
\TextAt{16}{275.571}{482.269}{\Character{10}}%
\TextAt{21}{284.87}{480.476}{K}%
\TextAt{29}{292.011}{479.103}{C}%
\TextAt{8}{299.981}{482.269}{det}%
\TextAt{13}{318.233}{482.269}{M}%
\TextAt{3}{330.806}{487.205}{\Character{48}\Character{48}}%
\TextAt{25}{329.557}{479.314}{C}%
\NativeMathAt{341.807}{489.046}{7.97}{\sim}%
\NativeArrowAt{339.219}{482.269}{11.955}{13.175}{\longrightarrow}%
\TextAt{8}{355.713}{482.269}{det}%
\TextAt{13}{373.964}{482.269}{M}%
\TextAt{25}{385.289}{480.476}{C}%
\TextAt{8}{79.226}{456.317}{of}%
\TextAt{13}{93.068}{456.317}{K}%
\TextAt{25}{103.027}{454.523}{C}%
\TextAt{8}{109.282}{456.317}{-modules.}%
\TextAt{8}{165.174}{456.317}{Let}%
\TextAt{13}{186.655}{456.317}{x}%
\TextAt{3}{193.307}{460.655}{\Character{48}}%
\TextAt{21}{193.307}{453.361}{i}%
\TextAt{13}{196.688}{456.317}{\Character{59}}%
\TextAt{13}{201.912}{456.317}{y}%
\TextAt{3}{208.069}{460.655}{\Character{48}}%
\TextAt{21}{207.64}{453.361}{i}%
\NativeMathAt{215.215}{456.317}{11.955}{\in}%
\TextAt{13}{227.378}{456.316}{M}%
\TextAt{3}{239.952}{460.655}{\Character{48}}%
\TextAt{25}{238.703}{453.252}{C}%
\TextAt{8}{244.957}{456.317}{,}%
\TextAt{13}{252.746}{456.317}{i}%
\TextAt{9}{260.933}{456.316}{\Character{61}}%
\TextAt{9}{274.215}{456.316}{1}%
\TextAt{13}{280.084}{456.317}{\Character{59}}%
\TextAt{13}{285.308}{456.317}{\Character{58}}%
\TextAt{13}{290.545}{456.317}{\Character{58}}%
\TextAt{13}{295.769}{456.317}{\Character{58}}%
\TextAt{13}{301.006}{456.317}{\Character{59}}%
\TextAt{13}{306.242}{456.317}{n}%
\TextAt{3}{313.292}{460.655}{\Character{48}}%
\TextAt{8}{316.087}{456.317}{,}%
\TextAt{8}{323.876}{456.317}{and}%
\TextAt{13}{347.151}{456.317}{x}%
\TextAt{21}{353.803}{454.523}{j}%
\TextAt{13}{358.185}{456.316}{\Character{59}}%
\TextAt{13}{363.409}{456.316}{y}%
\TextAt{21}{369.137}{454.523}{j}%
\NativeMathAt{377.712}{456.316}{11.955}{\in}%
\TextAt{9}{389.876}{456.316}{\Character{40}}%
\TextAt{13}{394.428}{456.316}{\Character{11}M}%
\TextAt{3}{414.523}{460.655}{\Character{48}}%
\TextAt{25}{413.275}{453.252}{C}%
\TextAt{9}{419.529}{456.316}{\Character{41}}%
\TextAt{3}{424.081}{460.655}{\Character{63}}%
\TextAt{8}{431.166}{456.317}{,}%
\TextAt{13}{438.955}{456.317}{j}%
\TextAt{9}{448.664}{456.316}{\Character{61}}%
\TextAt{9}{461.946}{456.316}{1}%
\TextAt{13}{467.815}{456.317}{\Character{59}}%
\TextAt{13}{473.039}{456.317}{\Character{58}}%
\TextAt{13}{478.276}{456.317}{\Character{58}}%
\TextAt{13}{483.5}{456.317}{\Character{58}}%
\TextAt{13}{488.737}{456.317}{\Character{59}}%
\TextAt{13}{493.973}{456.317}{n}%
\TextAt{3}{501.023}{460.655}{\Character{48}\Character{48}}%
\TextAt{8}{506.115}{456.317}{,}%
\TextAt{8}{513.904}{456.317}{and}%
\TextAt{8}{79.226}{439.271}{furthermore}%
\TextAt{13}{180.821}{422.225}{x}%
\TextAt{3}{187.473}{427.161}{\Character{48}}%
\TextAt{9}{193.589}{422.225}{\Character{61}}%
\TextAt{20}{206.015}{433.582}{\Character{94}}%
\TextAt{21}{211.215}{408.388}{i}%
\TextAt{13}{221.291}{422.225}{x}%
\TextAt{3}{227.943}{427.161}{\Character{48}}%
\TextAt{21}{227.943}{419.269}{i}%
\TextAt{13}{231.324}{422.225}{\Character{59}}%
\TextAt{13}{248.504}{422.225}{y}%
\TextAt{3}{254.66}{427.161}{\Character{48}}%
\TextAt{9}{260.776}{422.225}{\Character{61}}%
\TextAt{20}{273.202}{433.582}{\Character{94}}%
\TextAt{21}{278.402}{408.388}{i}%
\TextAt{13}{288.478}{422.225}{y}%
\TextAt{3}{294.614}{427.161}{\Character{48}}%
\TextAt{21}{294.185}{419.269}{i}%
\TextAt{13}{297.567}{422.225}{\Character{59}}%
\TextAt{13}{314.747}{422.225}{x}%
\TextAt{9}{324.739}{422.225}{\Character{61}}%
\TextAt{20}{337.164}{433.582}{\Character{94}}%
\TextAt{21}{341.864}{408.388}{j}%
\TextAt{13}{352.44}{422.225}{x}%
\TextAt{21}{359.092}{420.431}{j}%
\TextAt{13}{363.475}{422.225}{\Character{59}}%
\TextAt{13}{380.655}{422.225}{y}%
\TextAt{9}{390.131}{422.225}{\Character{61}}%
\TextAt{20}{402.557}{433.582}{\Character{94}}%
\TextAt{21}{407.257}{408.388}{j}%
\TextAt{13}{417.833}{422.225}{y}%
\TextAt{21}{423.541}{420.431}{j}%
\TextAt{13}{427.923}{422.225}{\Character{58}}%
\TextAt{8}{79.226}{391.18}{Then}%
\TextAt{8}{109.784}{391.18}{we}%
\TextAt{8}{127.011}{391.18}{have}%
\TextAt{16}{171.236}{365.227}{\Character{104}}%
\TextAt{13}{175.885}{365.227}{\Character{20}}%
\TextAt{9}{182.624}{365.227}{\Character{40}}%
\TextAt{13}{187.176}{365.227}{x}%
\TextAt{3}{193.828}{370.163}{\Character{48}}%
\TextAt{16}{199.28}{365.227}{\Character{10}}%
\TextAt{13}{211.235}{365.227}{\Character{12}}%
\TextAt{13}{218.504}{365.227}{x}%
\TextAt{9}{225.158}{365.227}{\Character{41}}%
\TextAt{13}{229.711}{365.227}{\Character{59}}%
\TextAt{13}{234.935}{365.227}{\Character{20}}%
\TextAt{9}{241.693}{365.227}{\Character{40}}%
\TextAt{13}{246.246}{365.227}{y}%
\TextAt{3}{252.382}{370.163}{\Character{48}}%
\TextAt{16}{257.834}{365.227}{\Character{10}}%
\TextAt{13}{269.789}{365.227}{\Character{12}}%
\TextAt{13}{277.058}{365.227}{y}%
\TextAt{9}{283.197}{365.227}{\Character{41}}%
\TextAt{16}{287.749}{365.227}{\Character{105}}%
\TextAt{7}{292.399}{363.434}{det}%
\TextAt{21}{305.566}{363.434}{M}%
\TextAt{9}{318.367}{365.227}{\Character{61}}%
\TextAt{16}{330.792}{365.227}{\Character{104}}%
\TextAt{13}{335.441}{365.227}{\Character{11}x}%
\TextAt{3}{349.615}{370.163}{\Character{48}}%
\TextAt{16}{355.067}{365.227}{\Character{94}}%
\TextAt{13}{365.694}{365.227}{x\Character{59}}%
\TextAt{13}{377.565}{365.227}{\Character{11}y}%
\TextAt{3}{391.248}{370.163}{\Character{48}}%
\TextAt{16}{396.7}{365.227}{\Character{94}}%
\TextAt{13}{407.327}{365.227}{y}%
\TextAt{16}{413.464}{365.227}{\Character{105}}%
\TextAt{7}{418.113}{363.434}{det}%
\TextAt{21}{431.28}{363.434}{M}%
\TextAt{9}{198.467}{335.207}{\Character{61}}%
\TextAt{8}{210.892}{335.207}{det}%
\TextAt{20}{229.144}{355.65}{\Character{32}}%
\TextAt{9}{238.609}{343.309}{\Character{40}}%
\TextAt{16}{243.161}{343.309}{\Character{104}}%
\TextAt{13}{247.81}{343.309}{x}%
\TextAt{3}{254.462}{347.648}{\Character{48}}%
\TextAt{21}{254.462}{340.354}{i}%
\TextAt{13}{257.844}{343.309}{\Character{59}}%
\TextAt{13}{263.068}{343.309}{y}%
\TextAt{3}{269.224}{347.648}{\Character{48}}%
\TextAt{21}{268.795}{340.2}{k}%
\TextAt{16}{273.915}{343.309}{\Character{105}}%
\TextAt{21}{278.564}{341.516}{M}%
\TextAt{23}{287.546}{343.785}{\Character{48}}%
\TextAt{9}{290.729}{343.309}{\Character{41}}%
\TextAt{9}{338.937}{343.309}{0}%
\TextAt{9}{264.018}{326.263}{0}%
\TextAt{9}{305.228}{326.263}{\Character{40}}%
\TextAt{16}{309.796}{326.263}{\Character{104}}%
\TextAt{13}{314.445}{326.263}{\Character{12}}%
\TextAt{13}{321.714}{326.263}{x}%
\TextAt{21}{328.368}{324.47}{j}%
\TextAt{13}{332.751}{326.263}{\Character{59}}%
\TextAt{13}{337.975}{326.263}{\Character{12}}%
\TextAt{13}{345.256}{326.263}{y}%
\TextAt{21}{350.974}{324.47}{\Character{96}}%
\TextAt{16}{354.931}{326.263}{\Character{105}}%
\TextAt{21}{359.58}{324.47}{M}%
\TextAt{23}{368.562}{326.739}{\Character{48}\Character{48}}%
\TextAt{9}{373.931}{326.263}{\Character{41}}%
\TextAt{20}{378.483}{355.65}{\Character{33}}%
\TextAt{9}{198.467}{304.345}{\Character{61}}%
\TextAt{8}{210.893}{304.345}{det}%
\TextAt{9}{227.152}{304.345}{\Character{40}}%
\TextAt{16}{231.704}{304.345}{\Character{104}}%
\TextAt{13}{236.353}{304.345}{x}%
\TextAt{3}{243.005}{309.281}{\Character{48}}%
\TextAt{21}{243.005}{301.39}{i}%
\TextAt{13}{246.387}{304.345}{\Character{59}}%
\TextAt{13}{251.611}{304.345}{y}%
\TextAt{3}{257.767}{309.281}{\Character{48}}%
\TextAt{21}{257.338}{301.39}{k}%
\TextAt{16}{262.458}{304.345}{\Character{105}}%
\TextAt{21}{267.107}{302.552}{M}%
\TextAt{23}{276.089}{304.821}{\Character{48}}%
\TextAt{9}{279.272}{304.345}{\Character{41}}%
\TextAt{8}{285.816}{304.345}{det}%
\TextAt{9}{302.075}{304.345}{\Character{40}}%
\TextAt{16}{306.628}{304.345}{\Character{104}}%
\TextAt{13}{311.277}{304.345}{\Character{12}}%
\TextAt{13}{318.546}{304.345}{x}%
\TextAt{21}{325.2}{302.552}{j}%
\TextAt{13}{329.583}{304.345}{\Character{59}}%
\TextAt{13}{334.807}{304.345}{\Character{12}}%
\TextAt{13}{342.088}{304.345}{y}%
\TextAt{21}{347.806}{302.552}{\Character{96}}%
\TextAt{16}{351.763}{304.345}{\Character{105}}%
\TextAt{21}{356.412}{302.552}{M}%
\TextAt{23}{365.394}{304.821}{\Character{48}\Character{48}}%
\TextAt{9}{370.763}{304.345}{\Character{41}}%
\TextAt{9}{198.467}{284.311}{\Character{61}}%
\TextAt{16}{210.893}{284.311}{\Character{104}}%
\TextAt{13}{215.542}{284.311}{x}%
\TextAt{3}{222.194}{289.247}{\Character{48}}%
\TextAt{13}{224.989}{284.311}{\Character{59}}%
\TextAt{13}{230.213}{284.311}{y}%
\TextAt{3}{236.37}{289.247}{\Character{48}}%
\TextAt{16}{239.165}{284.311}{\Character{105}}%
\TextAt{7}{243.814}{282.517}{det}%
\TextAt{21}{256.981}{282.517}{M}%
\TextAt{23}{265.963}{284.787}{\Character{48}}%
\TextAt{16}{269.146}{284.311}{\Character{104}}%
\TextAt{13}{273.795}{284.311}{\Character{12}}%
\TextAt{13}{281.064}{284.311}{x\Character{59}}%
\TextAt{13}{292.947}{284.311}{\Character{12}}%
\TextAt{13}{300.216}{284.311}{y}%
\TextAt{16}{306.37}{284.311}{\Character{105}}%
\TextAt{7}{311.02}{282.517}{det}%
\TextAt{21}{324.187}{282.517}{M}%
\TextAt{23}{333.169}{284.787}{\Character{48}\Character{48}}%
\TextAt{9}{198.467}{264.276}{\Character{61}}%
\TextAt{16}{210.893}{264.276}{\Character{104}}%
\TextAt{13}{215.542}{264.276}{x}%
\TextAt{3}{222.194}{269.212}{\Character{48}}%
\TextAt{16}{227.646}{264.276}{\Character{10}}%
\TextAt{13}{239.601}{264.276}{\Character{12}}%
\TextAt{13}{246.87}{264.276}{x\Character{59}}%
\TextAt{13}{258.753}{264.276}{y}%
\TextAt{3}{264.905}{269.212}{\Character{48}}%
\TextAt{16}{270.357}{264.276}{\Character{10}}%
\TextAt{13}{282.312}{264.276}{\Character{12}}%
\TextAt{13}{289.581}{264.276}{y}%
\TextAt{16}{295.72}{264.276}{\Character{105}}%
\TextAt{7}{300.369}{262.482}{det}%
\TextAt{21}{313.536}{262.482}{M}%
\TextAt{23}{322.518}{264.752}{\Character{48}}%
\TextAt{3}{325.203}{262.482}{\Character{10}}%
\TextAt{7}{331.789}{262.482}{det}%
\TextAt{21}{344.956}{262.482}{M}%
\TextAt{23}{353.938}{264.752}{\Character{48}\Character{48}}%
\TextAt{13}{359.307}{264.276}{\Character{58}}%
\TextAt{8}{79.226}{238.323}{Thus}%
\TextAt{13}{108.875}{238.323}{\Character{20}}%
\TextAt{8}{119.511}{238.323}{is}%
\TextAt{8}{131.275}{238.323}{an}%
\TextAt{8}{147.534}{238.323}{isometry.}%
\Vector{q 0 G 0 g 0.399 w 452.334 -535.677 m 452.334 -527.607 l S Q}%
\Vector{q 0 G 0 g 0.399 w 452.534 -527.806 m 459.707 -527.806 l S Q}%
\Vector{q 0 G 0 g 0.399 w 452.534 -535.478 m 459.707 -535.478 l S Q}%
\Vector{q 0 G 0 g 0.399 w 459.906 -535.677 m 459.906 -527.607 l S Q}%
\TextAt{12}{79.226}{205.086}{Exercise}%
\TextAt{12}{132.427}{205.086}{III.4.1.}%
\TextAt{8}{176.78}{205.086}{If}%
\TextAt{13}{188.472}{205.086}{M}%
\TextAt{13}{200.391}{205.086}{\Character{59}}%
\TextAt{13}{205.616}{205.086}{N}%
\TextAt{13}{215.586}{205.086}{\Character{59}}%
\TextAt{13}{220.823}{205.086}{L}%
\TextAt{8}{232.718}{205.086}{are}%
\TextAt{8}{252.229}{205.086}{metrized}%
\TextAt{14}{300.349}{205.086}{o}%
\TextAt{8}{306.193}{205.086}{-modules,}%
\TextAt{8}{359.25}{205.086}{then}%
\TextAt{8}{385.91}{205.086}{one}%
\TextAt{8}{407.369}{205.086}{has}%
\TextAt{8}{428.243}{205.086}{canonical}%
\TextAt{8}{479.627}{205.086}{isometries}%
\TextAt{13}{153.7}{179.133}{M}%
\TextAt{16}{168.93}{179.133}{\Character{10}}%
\TextAt{27}{178.229}{177.34}{o}%
\TextAt{13}{185.28}{179.133}{N}%
\NativeMathAt{199.223}{179.133}{11.955}{\cong}%
\TextAt{13}{211.842}{179.133}{N}%
\TextAt{16}{225.122}{179.133}{\Character{10}}%
\TextAt{27}{234.42}{177.34}{o}%
\TextAt{13}{241.471}{179.133}{M}%
\TextAt{13}{253.39}{179.133}{\Character{59}}%
\TextAt{9}{282.549}{179.133}{\Character{40}}%
\TextAt{13}{287.101}{179.133}{M}%
\TextAt{16}{302.331}{179.133}{\Character{10}}%
\TextAt{27}{311.63}{177.34}{o}%
\TextAt{13}{318.681}{179.133}{N}%
\TextAt{9}{329.303}{179.133}{\Character{41}}%
\TextAt{16}{336.512}{179.133}{\Character{10}}%
\TextAt{27}{345.811}{177.34}{o}%
\TextAt{13}{352.862}{179.133}{L}%
\NativeMathAt{364.147}{179.133}{11.955}{\cong}%
\TextAt{13}{376.767}{179.133}{M}%
\TextAt{16}{391.997}{179.133}{\Character{10}}%
\TextAt{27}{401.295}{177.34}{o}%
\TextAt{9}{408.346}{179.133}{\Character{40}}%
\TextAt{13}{412.899}{179.133}{N}%
\TextAt{16}{426.178}{179.133}{\Character{10}}%
\TextAt{27}{435.476}{177.34}{o}%
\TextAt{13}{442.527}{179.133}{L}%
\TextAt{9}{450.492}{179.133}{\Character{41}}%
\TextAt{13}{455.044}{179.133}{\Character{59}}%
\TextAt{13}{202.176}{153.181}{M}%
\TextAt{16}{217.407}{153.181}{\Character{10}}%
\TextAt{27}{226.705}{151.387}{o}%
\TextAt{9}{233.756}{153.181}{\Character{40}}%
\TextAt{13}{238.308}{153.181}{N}%
\TextAt{16}{251.588}{153.181}{\Character{8}}%
\TextAt{13}{263.543}{153.181}{L}%
\TextAt{9}{271.507}{153.181}{\Character{41}}%
\NativeMathAt{279.381}{153.181}{11.955}{\cong}%
\TextAt{9}{292}{153.181}{\Character{40}}%
\TextAt{13}{296.552}{153.181}{M}%
\TextAt{16}{311.782}{153.181}{\Character{10}}%
\TextAt{27}{321.081}{151.387}{o}%
\TextAt{13}{328.132}{153.181}{N}%
\TextAt{9}{338.755}{153.181}{\Character{41}}%
\TextAt{16}{345.964}{153.181}{\Character{8}}%
\TextAt{9}{357.919}{153.181}{\Character{40}}%
\TextAt{13}{362.471}{153.181}{M}%
\TextAt{16}{377.701}{153.181}{\Character{10}}%
\TextAt{27}{387}{151.387}{o}%
\TextAt{13}{394.051}{153.181}{L}%
\TextAt{9}{402.015}{153.181}{\Character{41}}%
\TextAt{13}{406.568}{153.181}{\Character{58}}%
\TextAt{12}{79.226}{117.763}{Exercise}%
\TextAt{12}{132.427}{117.763}{III.4.2.}%
\TextAt{8}{176.78}{117.763}{For}%
\TextAt{8}{197.75}{117.763}{any}%
\TextAt{8}{219.867}{117.763}{two}%
\TextAt{8}{241.984}{117.763}{projective}%
\TextAt{8}{295.961}{117.763}{metrized}%
\TextAt{14}{344.081}{117.763}{o}%
\TextAt{8}{349.925}{117.763}{-modules}%
\TextAt{13}{399.73}{117.763}{M}%
\TextAt{13}{411.649}{117.763}{\Character{59}}%
\TextAt{13}{416.874}{117.763}{N}%
\TextAt{8}{427.52}{117.763}{,}%
\TextAt{8}{434.67}{117.763}{one}%
\TextAt{8}{456.129}{117.763}{has}%
\TextAt{21}{222.986}{98.47}{n}%
\TextAt{20}{218.913}{94.883}{\Character{94}}%
\TextAt{9}{232.197}{83.526}{\Character{40}}%
\TextAt{13}{236.749}{83.526}{M}%
\TextAt{16}{251.98}{83.526}{\Character{8}}%
\TextAt{13}{263.935}{83.526}{N}%
\TextAt{9}{274.557}{83.526}{\Character{41}}%
\NativeMathAt{282.431}{83.526}{11.955}{\cong}%
\TextAt{20}{298.556}{94.883}{\Character{77}}%
\TextAt{21}{295.05}{69.689}{i}%
\TextAt{17}{297.933}{69.689}{\Character{43}}%
\TextAt{21}{304.519}{69.689}{j}%
\TextAt{17}{308.403}{69.689}{\Character{61}}%
\TextAt{21}{314.99}{69.689}{n}%
\TextAt{21}{327.321}{98.47}{i}%
\TextAt{20}{322.121}{94.883}{\Character{94}}%
\TextAt{13}{337.397}{83.526}{M}%
\TextAt{16}{352.627}{83.526}{\Character{10}}%
\TextAt{21}{369.282}{98.957}{j}%
\TextAt{20}{364.582}{94.883}{\Character{94}}%
\TextAt{13}{379.858}{83.526}{N}%
\TextAt{13}{389.829}{83.526}{\Character{58}}%
\TextAt{8}{297.211}{43.033}{213}%
\Destination{78.23}{809.815}{nav.page.213}%
\Destination{79.226}{219.055}{nav.ex.3.4.1}%
\Destination{79.226}{131.755}{nav.ex.3.4.2}%
\Destination{79.36}{649.255}{nav.print.232}%
}

% Page 226 (body page 214)
\Page{%
\TextAt{2}{79.226}{796.943}{III.}%
\TextAt{2}{98.317}{796.943}{Riemann–Roch}%
\TextAt{2}{174.375}{796.943}{Theory}%
\Vector{q 0 G 0 g 0.299 w 7.226 18.729 m 460.769 18.729 l S Q}%
\TextAt{12}{79.226}{761.124}{Exercise}%
\TextAt{12}{132.427}{761.124}{III.4.3.}%
\TextAt{8}{176.78}{761.124}{For}%
\TextAt{8}{197.75}{761.124}{any}%
\TextAt{8}{219.867}{761.124}{two}%
\TextAt{8}{241.984}{761.124}{projective}%
\TextAt{8}{295.961}{761.124}{metrized}%
\TextAt{14}{344.081}{761.124}{o}%
\TextAt{8}{349.925}{761.124}{-modules}%
\TextAt{13}{399.73}{761.124}{M}%
\TextAt{13}{411.649}{761.124}{\Character{59}}%
\TextAt{13}{416.874}{761.124}{N}%
\TextAt{8}{427.52}{761.124}{,}%
\TextAt{8}{434.67}{761.124}{one}%
\TextAt{8}{456.129}{761.124}{has}%
\TextAt{8}{183.032}{732.123}{det}%
\TextAt{9}{199.291}{732.123}{\Character{40}}%
\TextAt{13}{203.844}{732.123}{M}%
\TextAt{16}{219.074}{732.123}{\Character{10}}%
\TextAt{27}{228.372}{730.329}{o}%
\TextAt{13}{235.423}{732.123}{N}%
\TextAt{9}{246.046}{732.123}{\Character{41}}%
\NativeMathAt{253.919}{732.122}{11.955}{\cong}%
\TextAt{9}{266.538}{732.123}{\Character{40}}%
\TextAt{8}{271.091}{732.123}{det}%
\TextAt{13}{289.342}{732.123}{M}%
\TextAt{9}{301.916}{732.123}{\Character{41}}%
\TextAt{3}{306.468}{737.059}{\Character{10}}%
\TextAt{7}{314.466}{737.059}{rk}%
\TextAt{17}{322.237}{737.059}{\Character{40}}%
\TextAt{21}{325.53}{737.059}{N}%
\TextAt{17}{333.1}{737.059}{\Character{41}}%
\TextAt{16}{339.548}{732.123}{\Character{10}}%
\TextAt{27}{348.847}{730.329}{o}%
\TextAt{9}{355.898}{732.123}{\Character{40}}%
\TextAt{8}{360.45}{732.123}{det}%
\TextAt{13}{378.702}{732.123}{N}%
\TextAt{9}{389.324}{732.123}{\Character{41}}%
\TextAt{3}{393.877}{737.059}{\Character{10}}%
\TextAt{7}{401.875}{737.059}{rk}%
\TextAt{17}{409.645}{737.059}{\Character{40}}%
\TextAt{21}{412.939}{737.059}{M}%
\TextAt{17}{421.92}{737.059}{\Character{41}}%
\TextAt{13}{425.712}{732.123}{\Character{58}}%
\TextAt{12}{79.226}{688.177}{Exercise}%
\TextAt{12}{133.597}{688.177}{III.4.4.}%
\TextAt{8}{178.968}{688.177}{If}%
\TextAt{13}{191.677}{688.177}{M}%
\TextAt{8}{209.165}{688.177}{is}%
\TextAt{8}{221.947}{688.177}{a}%
\TextAt{8}{232.719}{688.177}{projective}%
\TextAt{8}{287.714}{688.177}{metrized}%
\TextAt{14}{336.851}{688.177}{o}%
\TextAt{8}{342.695}{688.177}{-module}%
\TextAt{8}{388.903}{688.177}{of}%
\TextAt{8}{403.25}{688.177}{rank}%
\TextAt{13}{431.262}{688.177}{n}%
\TextAt{8}{438.25}{688.177}{,}%
\TextAt{8}{446.671}{688.177}{and}%
\TextAt{13}{470.451}{688.177}{p}%
\NativeMathAt{481.38}{688.177}{11.955}{\geq}%
\TextAt{9}{495.733}{688.177}{0}%
\TextAt{8}{501.586}{688.177}{,}%
\TextAt{8}{510.007}{688.177}{then}%
\TextAt{8}{79.226}{671.131}{there}%
\TextAt{8}{109.138}{671.131}{is}%
\TextAt{8}{120.902}{671.131}{a}%
\TextAt{8}{130.657}{671.131}{canonical}%
\TextAt{8}{182.041}{671.131}{isometry}%
\TextAt{8}{228.822}{634.819}{det}%
\TextAt{20}{247.074}{655.263}{\Character{32}}%
\TextAt{21}{261.049}{650.251}{p}%
\TextAt{20}{256.538}{646.177}{\Character{94}}%
\TextAt{13}{271.814}{634.82}{M}%
\TextAt{20}{284.388}{655.263}{\Character{33}}%
\NativeMathAt{297.173}{634.819}{11.955}{\cong}%
\TextAt{9}{309.793}{634.819}{\Character{40}}%
\TextAt{8}{314.345}{634.82}{det}%
\TextAt{13}{332.597}{634.819}{M}%
\TextAt{9}{345.17}{634.82}{\Character{41}}%
\TextAt{3}{349.723}{640.394}{\Character{10}}%
\TextAt{9}{356.309}{640.394}{\Character{40}}%
\TextAt{22}{360.861}{644.04}{n}%
\TextAt{23}{365.462}{644.04}{\Character{0}}%
\TextAt{24}{371.218}{644.04}{1}%
\TextAt{22}{361.241}{637.269}{p}%
\TextAt{23}{365.082}{637.269}{\Character{0}}%
\TextAt{24}{370.838}{637.269}{1}%
\TextAt{9}{374.871}{640.394}{\Character{41}}%
\TextAt{13}{379.922}{634.819}{\Character{58}}%
\TextAt{15}{79.226}{567.649}{§}%
\TextAt{15}{93.785}{567.649}{5}%
\TextAt{15}{120.693}{567.649}{Grothendieck}%
\TextAt{15}{241.149}{567.649}{Groups}%
\TextAt{8}{79.226}{535.763}{We}%
\TextAt{8}{98.307}{535.763}{will}%
\TextAt{8}{119.348}{535.763}{now}%
\TextAt{8}{142.673}{535.763}{manufacture}%
\TextAt{8}{209.251}{535.763}{two}%
\TextAt{8}{230.304}{535.763}{abelian}%
\TextAt{8}{269.888}{535.763}{groups}%
\TextAt{8}{306.614}{535.763}{from}%
\TextAt{8}{333.191}{535.763}{the}%
\TextAt{8}{352.284}{535.763}{collection}%
\TextAt{8}{403.249}{535.763}{of}%
\TextAt{8}{415.515}{535.763}{all}%
\TextAt{8}{430.71}{535.763}{metrized}%
\TextAt{14}{477.766}{535.763}{o}%
\TextAt{8}{483.61}{535.763}{-modules,}%
\TextAt{8}{79.226}{518.717}{resp.}%
\TextAt{8}{108.362}{518.717}{the}%
\TextAt{8}{127.959}{518.717}{collection}%
\TextAt{8}{179.428}{518.717}{of}%
\TextAt{8}{192.199}{518.717}{all}%
\TextAt{8}{207.898}{518.717}{projective}%
\TextAt{8}{261.317}{518.717}{metrized}%
\TextAt{14}{308.877}{518.717}{o}%
\TextAt{8}{314.721}{518.717}{-modules.}%
\TextAt{8}{368.891}{518.717}{We}%
\TextAt{8}{388.476}{518.717}{denote}%
\TextAt{8}{425.635}{518.717}{by}%
\TextAt{16}{441.335}{518.717}{\Character{102}}%
\TextAt{13}{447.312}{518.717}{M}%
\TextAt{16}{459.886}{518.717}{\Character{103}}%
\TextAt{8}{469.201}{518.717}{the}%
\TextAt{8}{488.798}{518.717}{isometry}%
\TextAt{8}{79.226}{501.671}{class}%
\TextAt{8}{106.663}{501.671}{of}%
\TextAt{8}{119.993}{501.671}{a}%
\TextAt{8}{129.749}{501.671}{metrized}%
\TextAt{14}{177.868}{501.671}{o}%
\TextAt{8}{183.712}{501.671}{-module}%
\TextAt{13}{228.903}{501.671}{M}%
\TextAt{8}{245.374}{501.671}{and}%
\TextAt{8}{268.137}{501.671}{form}%
\TextAt{8}{295.777}{501.671}{the}%
\TextAt{8}{315.933}{501.671}{free}%
\TextAt{8}{338.361}{501.671}{abelian}%
\TextAt{8}{379.009}{501.671}{group}%
\TextAt{13}{169.402}{472.67}{F}%
\TextAt{17}{176.979}{470.876}{0}%
\TextAt{9}{181.712}{472.67}{\Character{40}}%
\Vector{q 0 G 0 g 0.478 w 114.264 -293.974 m 120.108 -293.974 l S Q}%
\TextAt{14}{186.264}{472.67}{o}%
\TextAt{9}{192.108}{472.67}{\Character{41}}%
\TextAt{9}{199.963}{472.67}{\Character{61}}%
\TextAt{20}{212.407}{484.027}{\Character{77}}%
\TextAt{3}{212.715}{458.124}{\Character{102}}%
\TextAt{21}{216.949}{458.124}{M}%
\TextAt{3}{225.93}{458.124}{\Character{103}}%
\TextAt{18}{232.465}{472.67}{Z}%
\TextAt{16}{240.435}{472.67}{\Character{102}}%
\TextAt{13}{246.413}{472.67}{M}%
\TextAt{16}{258.986}{472.67}{\Character{103}}%
\TextAt{13}{264.964}{472.67}{\Character{59}}%
\TextAt{8}{294.118}{472.67}{resp.}%
\TextAt{13}{342.154}{472.67}{F}%
\TextAt{17}{351.358}{477.606}{0}%
\TextAt{9}{356.09}{472.67}{\Character{40}}%
\Vector{q 0 G 0 g 0.478 w 288.643 -293.974 m 294.487 -293.974 l S Q}%
\TextAt{14}{360.643}{472.67}{o}%
\TextAt{9}{366.487}{472.67}{\Character{41}}%
\TextAt{9}{374.342}{472.67}{\Character{61}}%
\TextAt{20}{386.785}{484.027}{\Character{77}}%
\TextAt{3}{387.093}{458.124}{\Character{102}}%
\TextAt{21}{391.327}{458.124}{M}%
\TextAt{3}{400.309}{458.124}{\Character{103}}%
\TextAt{18}{406.844}{472.67}{Z}%
\TextAt{16}{414.814}{472.67}{\Character{102}}%
\TextAt{13}{420.791}{472.67}{M}%
\TextAt{16}{433.365}{472.67}{\Character{103}}%
\TextAt{13}{439.343}{472.67}{\Character{59}}%
\TextAt{8}{79.226}{433.568}{on}%
\TextAt{8}{95.693}{433.568}{the}%
\TextAt{8}{116.057}{433.568}{isometry}%
\TextAt{8}{164.134}{433.568}{classes}%
\TextAt{8}{201.594}{433.568}{of}%
\TextAt{8}{215.132}{433.568}{projective,}%
\TextAt{8}{272.621}{433.568}{resp.}%
\TextAt{8}{302.566}{433.568}{coherent,}%
\TextAt{8}{353.229}{433.568}{metrized}%
\TextAt{14}{401.556}{433.568}{o}%
\TextAt{8}{407.401}{433.568}{-modules.}%
\TextAt{8}{462.38}{433.568}{In}%
\TextAt{8}{477.209}{433.568}{this}%
\TextAt{8}{500.239}{433.568}{group,}%
\TextAt{8}{79.226}{416.522}{we}%
\TextAt{8}{96.454}{416.522}{consider}%
\TextAt{8}{142.039}{416.522}{the}%
\TextAt{8}{162.195}{416.522}{subgroup}%
\TextAt{13}{193.124}{387.521}{R}%
\TextAt{17}{202.034}{385.728}{0}%
\TextAt{9}{206.767}{387.521}{\Character{40}}%
\Vector{q 0 G 0 g 0.478 w 139.319 -379.123 m 145.163 -379.123 l S Q}%
\TextAt{14}{211.319}{387.521}{o}%
\TextAt{9}{217.163}{387.521}{\Character{41}}%
\NativeMathAt{225.036}{387.521}{11.955}{\subseteq}%
\TextAt{13}{237.656}{387.521}{F}%
\TextAt{17}{245.233}{385.728}{0}%
\TextAt{9}{249.966}{387.521}{\Character{40}}%
\Vector{q 0 G 0 g 0.478 w 182.518 -379.123 m 188.362 -379.123 l S Q}%
\TextAt{14}{254.518}{387.521}{o}%
\TextAt{9}{260.362}{387.521}{\Character{41}}%
\TextAt{13}{264.915}{387.521}{\Character{59}}%
\TextAt{8}{294.069}{387.521}{resp.}%
\TextAt{13}{342.105}{387.521}{R}%
\TextAt{17}{351.114}{392.457}{0}%
\TextAt{9}{355.846}{387.521}{\Character{40}}%
\Vector{q 0 G 0 g 0.478 w 288.398 -379.123 m 294.242 -379.123 l S Q}%
\TextAt{14}{360.398}{387.521}{o}%
\TextAt{9}{366.242}{387.521}{\Character{41}}%
\NativeMathAt{374.116}{387.521}{11.955}{\subseteq}%
\TextAt{13}{386.735}{387.521}{F}%
\TextAt{17}{395.939}{392.457}{0}%
\TextAt{9}{400.671}{387.521}{\Character{40}}%
\Vector{q 0 G 0 g 0.478 w 333.223 -379.123 m 339.067 -379.123 l S Q}%
\TextAt{14}{405.223}{387.521}{o}%
\TextAt{9}{411.067}{387.521}{\Character{41}}%
\TextAt{13}{415.62}{387.521}{\Character{59}}%
\TextAt{8}{79.226}{358.52}{generated}%
\TextAt{8}{132.558}{358.52}{by}%
\TextAt{8}{148.817}{358.52}{all}%
\TextAt{8}{165.076}{358.52}{elements}%
\TextAt{16}{212.933}{358.52}{\Character{102}}%
\TextAt{13}{218.91}{358.52}{M}%
\TextAt{3}{231.484}{362.858}{\Character{48}}%
\TextAt{16}{234.279}{358.52}{\Character{103}}%
\TextAt{16}{242.911}{358.52}{\Character{0}}%
\TextAt{16}{254.854}{358.52}{\Character{102}}%
\TextAt{13}{260.846}{358.52}{M}%
\TextAt{16}{273.42}{358.52}{\Character{103}}%
\TextAt{9}{282.054}{358.52}{\Character{43}}%
\TextAt{16}{293.815}{358.52}{\Character{102}}%
\TextAt{13}{299.793}{358.52}{M}%
\TextAt{3}{312.367}{362.858}{\Character{48}\Character{48}}%
\TextAt{16}{317.459}{358.52}{\Character{103}}%
\TextAt{8}{327.334}{358.52}{which}%
\TextAt{8}{360.82}{358.52}{arise}%
\TextAt{8}{388.197}{358.52}{from}%
\TextAt{8}{415.838}{358.52}{a}%
\TextAt{8}{425.593}{358.52}{short}%
\TextAt{8}{455.577}{358.52}{exact}%
\TextAt{8}{486.469}{358.52}{sequence}%
\TextAt{9}{225.535}{329.519}{0}%
\NativeArrowAt{234.709}{329.519}{11.955}{19.26}{\longrightarrow}%
\TextAt{13}{257.291}{329.519}{M}%
\TextAt{3}{269.865}{334.455}{\Character{48}}%
\NativeArrowAt{275.981}{329.519}{11.955}{19.26}{\longrightarrow}%
\TextAt{13}{298.563}{329.519}{M}%
\NativeArrowAt{314.457}{329.519}{11.955}{19.26}{\longrightarrow}%
\TextAt{13}{337.039}{329.519}{M}%
\TextAt{3}{349.613}{334.455}{\Character{48}\Character{48}}%
\NativeArrowAt{358.025}{329.519}{11.955}{19.26}{\longrightarrow}%
\TextAt{9}{380.607}{329.519}{0}%
\TextAt{8}{79.226}{300.518}{of}%
\TextAt{8}{92.556}{300.518}{projective,}%
\TextAt{8}{149.786}{300.518}{resp.}%
\TextAt{8}{179.108}{300.518}{coherent,}%
\TextAt{8}{229.511}{300.518}{metrized}%
\TextAt{14}{277.63}{300.518}{o}%
\TextAt{8}{283.474}{300.518}{-modules.}%
\TextAt{12}{79.226}{274.505}{Definition}%
\TextAt{12}{142.015}{274.505}{(5.1).}%
\TextAt{8}{177.33}{274.505}{The}%
\TextAt{8}{201.384}{274.505}{quotient}%
\TextAt{8}{247.567}{274.505}{groups}%
\TextAt{13}{159.021}{245.504}{K}%
\TextAt{17}{168.98}{243.711}{0}%
\TextAt{9}{173.713}{245.504}{\Character{40}}%
\Vector{q 0 G 0 g 0.478 w 106.265 -521.139 m 112.109 -521.139 l S Q}%
\TextAt{14}{178.265}{245.504}{o}%
\TextAt{9}{184.109}{245.504}{\Character{41}}%
\TextAt{9}{191.964}{245.504}{\Character{61}}%
\TextAt{13}{204.408}{245.504}{F}%
\TextAt{17}{211.986}{243.711}{0}%
\TextAt{9}{216.718}{245.504}{\Character{40}}%
\Vector{q 0 G 0 g 0.478 w 149.27 -521.139 m 155.115 -521.139 l S Q}%
\TextAt{14}{221.27}{245.504}{o}%
\TextAt{9}{227.115}{245.504}{\Character{41}\Character{47}}%
\TextAt{13}{237.52}{245.504}{R}%
\TextAt{17}{246.43}{243.711}{0}%
\TextAt{9}{251.162}{245.504}{\Character{40}}%
\Vector{q 0 G 0 g 0.478 w 183.715 -521.139 m 189.559 -521.139 l S Q}%
\TextAt{14}{255.715}{245.504}{o}%
\TextAt{9}{261.559}{245.504}{\Character{41}}%
\TextAt{13}{266.111}{245.504}{\Character{59}}%
\TextAt{8}{295.266}{245.504}{resp.}%
\TextAt{13}{343.302}{245.504}{K}%
\TextAt{17}{354.119}{250.44}{0}%
\TextAt{9}{358.852}{245.504}{\Character{40}}%
\Vector{q 0 G 0 g 0.478 w 291.404 -521.139 m 297.248 -521.139 l S Q}%
\TextAt{14}{363.404}{245.504}{o}%
\TextAt{9}{369.248}{245.504}{\Character{41}}%
\TextAt{9}{377.103}{245.504}{\Character{61}}%
\TextAt{13}{389.547}{245.504}{F}%
\TextAt{17}{398.75}{250.44}{0}%
\TextAt{9}{403.483}{245.504}{\Character{40}}%
\Vector{q 0 G 0 g 0.478 w 336.035 -521.139 m 341.879 -521.139 l S Q}%
\TextAt{14}{408.035}{245.504}{o}%
\TextAt{9}{413.879}{245.504}{\Character{41}\Character{47}}%
\TextAt{13}{424.285}{245.504}{R}%
\TextAt{17}{433.293}{250.44}{0}%
\TextAt{9}{438.025}{245.504}{\Character{40}}%
\Vector{q 0 G 0 g 0.478 w 370.578 -521.139 m 376.422 -521.139 l S Q}%
\TextAt{14}{442.578}{245.504}{o}%
\TextAt{9}{448.422}{245.504}{\Character{41}}%
\TextAt{8}{79.226}{216.503}{are}%
\TextAt{8}{100.665}{216.503}{called}%
\TextAt{8}{135.756}{216.503}{the}%
\TextAt{12}{157.841}{216.503}{replete}%
\TextAt{12}{204.89}{216.503}{(or}%
\TextAt{12}{229.057}{216.503}{compactified)}%
\TextAt{12}{314.961}{216.503}{Grothendieck}%
\TextAt{12}{400.829}{216.503}{groups}%
\TextAt{8}{445.867}{216.503}{of}%
\TextAt{14}{461.125}{216.503}{o}%
\TextAt{8}{466.969}{216.503}{.}%
\TextAt{8}{481.201}{216.503}{If}%
\TextAt{13}{494.821}{216.503}{M}%
\TextAt{8}{513.22}{216.503}{is}%
\TextAt{8}{526.911}{216.503}{a}%
\TextAt{8}{79.226}{199.457}{metrized}%
\TextAt{14}{127.346}{199.457}{o}%
\TextAt{8}{133.19}{199.457}{-module,}%
\TextAt{8}{181.632}{199.457}{then}%
\TextAt{9}{208.292}{199.457}{\Character{91}}%
\TextAt{13}{211.544}{199.457}{M}%
\TextAt{9}{224.118}{199.457}{\Character{93}}%
\TextAt{8}{231.267}{199.457}{denotes}%
\TextAt{8}{273.6}{199.457}{the}%
\TextAt{8}{293.756}{199.457}{class}%
\TextAt{8}{321.193}{199.457}{it}%
\TextAt{8}{332.898}{199.457}{defines}%
\TextAt{8}{371.321}{199.457}{in}%
\TextAt{13}{384.974}{199.457}{K}%
\TextAt{17}{394.933}{197.664}{0}%
\TextAt{9}{399.666}{199.457}{\Character{40}}%
\Vector{q 0 G 0 g 0.478 w 332.218 -567.187 m 338.062 -567.187 l S Q}%
\TextAt{14}{404.218}{199.457}{o}%
\TextAt{9}{410.062}{199.457}{\Character{41}}%
\TextAt{8}{414.615}{199.457}{,}%
\TextAt{8}{421.764}{199.457}{resp.}%
\TextAt{13}{451.086}{199.457}{K}%
\TextAt{17}{461.904}{203.795}{0}%
\TextAt{9}{466.636}{199.457}{\Character{40}}%
\Vector{q 0 G 0 g 0.478 w 399.188 -567.187 m 405.032 -567.187 l S Q}%
\TextAt{14}{471.188}{199.457}{o}%
\TextAt{9}{477.032}{199.457}{\Character{41}}%
\TextAt{8}{481.585}{199.457}{.}%
\TextAt{8}{97.159}{173.445}{The}%
\TextAt{8}{121.213}{173.445}{construction}%
\TextAt{8}{188.269}{173.445}{of}%
\TextAt{8}{201.599}{173.445}{the}%
\TextAt{8}{221.756}{173.445}{Grothendieck}%
\TextAt{8}{294.025}{173.445}{groups}%
\TextAt{8}{331.815}{173.445}{is}%
\TextAt{8}{343.579}{173.445}{such}%
\TextAt{8}{369.976}{173.445}{that}%
\TextAt{8}{395.345}{173.445}{a}%
\TextAt{8}{405.1}{173.445}{short}%
\TextAt{8}{435.084}{173.445}{exact}%
\TextAt{8}{465.976}{173.445}{sequence}%
\TextAt{9}{225.535}{144.444}{0}%
\NativeArrowAt{234.709}{144.444}{11.955}{19.26}{\longrightarrow}%
\TextAt{13}{257.291}{144.444}{M}%
\TextAt{3}{269.865}{149.38}{\Character{48}}%
\NativeArrowAt{275.981}{144.444}{11.955}{19.26}{\longrightarrow}%
\TextAt{13}{298.563}{144.444}{M}%
\NativeArrowAt{314.457}{144.444}{11.955}{19.26}{\longrightarrow}%
\TextAt{13}{337.039}{144.444}{M}%
\TextAt{3}{349.613}{149.38}{\Character{48}\Character{48}}%
\NativeArrowAt{358.025}{144.444}{11.955}{19.26}{\longrightarrow}%
\TextAt{9}{380.607}{144.444}{0}%
\TextAt{8}{79.226}{115.442}{of}%
\TextAt{8}{92.556}{115.442}{metrized}%
\TextAt{14}{140.676}{115.442}{o}%
\TextAt{8}{146.52}{115.442}{-modules}%
\TextAt{8}{196.325}{115.442}{becomes}%
\TextAt{8}{242.879}{115.442}{an}%
\TextAt{8}{259.138}{115.442}{additive}%
\TextAt{8}{304.017}{115.442}{decomposition}%
\TextAt{8}{381.152}{115.442}{in}%
\TextAt{8}{394.805}{115.442}{the}%
\TextAt{8}{414.961}{115.442}{group:}%
\TextAt{9}{256.731}{86.441}{\Character{91}}%
\TextAt{13}{259.983}{86.441}{M}%
\TextAt{9}{272.556}{86.441}{\Character{93}}%
\TextAt{9}{279.107}{86.441}{\Character{61}}%
\TextAt{9}{291.529}{86.441}{\Character{91}}%
\TextAt{13}{294.806}{86.441}{M}%
\TextAt{3}{307.379}{91.377}{\Character{48}}%
\TextAt{9}{310.174}{86.441}{\Character{93}}%
\TextAt{9}{316.068}{86.441}{\Character{43}}%
\TextAt{9}{327.82}{86.441}{\Character{91}}%
\TextAt{13}{331.096}{86.441}{M}%
\TextAt{3}{343.669}{91.377}{\Character{48}\Character{48}}%
\TextAt{9}{348.761}{86.441}{\Character{93}}%
\TextAt{13}{352.013}{86.441}{\Character{58}}%
\TextAt{8}{297.211}{43.033}{214}%
\Destination{78.23}{809.815}{nav.page.214}%
\Destination{79.23}{583.905}{nav.section.3.5}%
\Destination{79.226}{775.155}{nav.ex.3.4.3}%
\Destination{79.226}{702.155}{nav.ex.3.4.4}%
\Destination{79.36}{776.155}{nav.print.233}%
\Destination{79.36}{288.455}{nav.number.3.5.1}%
\Destination{79.23}{583.905}{outline.36}%
}

% Page 227 (body page 215)
\Page{%
\TextAt{2}{79.226}{796.943}{III.}%
\TextAt{2}{98.317}{796.943}{Riemann–Roch}%
\TextAt{2}{174.375}{796.943}{Theory}%
\Vector{q 0 G 0 g 0.299 w 7.226 18.729 m 460.769 18.729 l S Q}%
\TextAt{8}{79.226}{761.124}{In}%
\TextAt{8}{93.847}{761.124}{particular,}%
\TextAt{8}{151.089}{761.124}{one}%
\TextAt{8}{172.548}{761.124}{has}%
\TextAt{9}{239.195}{744.078}{\Character{91}}%
\TextAt{13}{242.446}{744.078}{M}%
\TextAt{3}{255.02}{749.014}{\Character{48}}%
\TextAt{16}{260.472}{744.078}{\Character{8}}%
\TextAt{13}{272.427}{744.078}{M}%
\TextAt{3}{285}{749.014}{\Character{48}\Character{48}}%
\TextAt{9}{290.092}{744.078}{\Character{93}}%
\TextAt{9}{296.643}{744.078}{\Character{61}}%
\TextAt{9}{309.065}{744.078}{\Character{91}}%
\TextAt{13}{312.342}{744.078}{M}%
\TextAt{3}{324.916}{749.014}{\Character{48}}%
\TextAt{9}{327.711}{744.078}{\Character{93}}%
\TextAt{9}{333.605}{744.078}{\Character{43}}%
\TextAt{9}{345.357}{744.078}{\Character{91}}%
\TextAt{13}{348.632}{744.078}{M}%
\TextAt{3}{361.206}{749.014}{\Character{48}\Character{48}}%
\TextAt{9}{366.298}{744.078}{\Character{93}}%
\TextAt{13}{369.549}{744.078}{\Character{58}}%
\TextAt{8}{79.226}{720.556}{The}%
\TextAt{8}{103.9}{720.556}{tensor}%
\TextAt{8}{139.703}{720.556}{product}%
\TextAt{8}{184.222}{720.556}{even}%
\TextAt{8}{211.502}{720.556}{induces}%
\TextAt{8}{253.797}{720.556}{a}%
\TextAt{8}{264.173}{720.556}{ring}%
\TextAt{8}{288.858}{720.556}{structure}%
\TextAt{8}{339.618}{720.556}{on}%
\TextAt{13}{356.496}{720.556}{K}%
\TextAt{17}{366.455}{718.763}{0}%
\TextAt{9}{371.188}{720.556}{\Character{40}}%
\Vector{q 0 G 0 g 0.478 w 303.74 -46.088 m 309.584 -46.088 l S Q}%
\TextAt{14}{375.74}{720.556}{o}%
\TextAt{9}{381.584}{720.556}{\Character{41}}%
\TextAt{8}{386.137}{720.556}{,}%
\TextAt{8}{394.06}{720.556}{and}%
\TextAt{13}{417.443}{720.556}{K}%
\TextAt{17}{428.26}{724.894}{0}%
\TextAt{9}{432.993}{720.556}{\Character{40}}%
\Vector{q 0 G 0 g 0.478 w 365.545 -46.088 m 371.389 -46.088 l S Q}%
\TextAt{14}{437.545}{720.556}{o}%
\TextAt{9}{443.389}{720.556}{\Character{41}}%
\TextAt{8}{452.459}{720.556}{then}%
\TextAt{8}{479.738}{720.556}{becomes}%
\TextAt{8}{526.911}{720.556}{a}%
\TextAt{13}{79.226}{703.51}{K}%
\TextAt{17}{89.185}{701.717}{0}%
\TextAt{9}{93.918}{703.51}{\Character{40}}%
\Vector{q 0 G 0 g 0.478 w 26.47 -63.133 m 32.314 -63.133 l S Q}%
\TextAt{14}{98.47}{703.51}{o}%
\TextAt{9}{104.314}{703.51}{\Character{41}}%
\TextAt{8}{108.866}{703.51}{-module:}%
\TextAt{8}{158.608}{703.51}{extending}%
\TextAt{8}{212.263}{703.51}{the}%
\TextAt{8}{232.419}{703.51}{product}%
\TextAt{16}{241.121}{674.509}{\Character{102}}%
\TextAt{13}{247.098}{674.509}{M}%
\TextAt{16}{259.672}{674.509}{\Character{103}\Character{102}}%
\TextAt{13}{271.627}{674.509}{M}%
\TextAt{3}{284.201}{679.445}{\Character{48}}%
\TextAt{16}{286.996}{674.509}{\Character{103}}%
\TextAt{9}{296.294}{674.509}{\Character{58}\Character{61}}%
\TextAt{16}{311.971}{674.509}{\Character{102}}%
\TextAt{13}{317.949}{674.509}{M}%
\TextAt{16}{333.179}{674.509}{\Character{10}}%
\TextAt{27}{342.478}{672.716}{o}%
\TextAt{13}{349.529}{674.509}{M}%
\TextAt{3}{362.102}{679.445}{\Character{48}}%
\TextAt{16}{364.897}{674.509}{\Character{103}}%
\TextAt{8}{79.226}{645.508}{by}%
\TextAt{8}{96.02}{645.508}{linearity,}%
\TextAt{8}{145.142}{645.508}{and}%
\TextAt{8}{168.439}{645.508}{observing}%
\TextAt{8}{221.397}{645.508}{that}%
\TextAt{13}{247.3}{645.508}{N}%
\TextAt{16}{260.944}{645.508}{\Character{10}}%
\TextAt{13}{273.263}{645.508}{M}%
\NativeMathAt{290.068}{645.508}{11.955}{\cong}%
\TextAt{13}{303.598}{645.508}{M}%
\TextAt{16}{319.193}{645.508}{\Character{10}}%
\TextAt{13}{331.512}{645.508}{N}%
\TextAt{8}{346.567}{645.508}{and}%
\TextAt{9}{369.864}{645.508}{\Character{40}}%
\TextAt{13}{374.416}{645.508}{M}%
\TextAt{16}{390.011}{645.508}{\Character{10}}%
\TextAt{13}{402.33}{645.508}{N}%
\TextAt{9}{412.953}{645.508}{\Character{41}}%
\TextAt{16}{420.526}{645.508}{\Character{10}}%
\TextAt{13}{432.845}{645.508}{L}%
\NativeMathAt{445.041}{645.508}{11.955}{\cong}%
\TextAt{13}{458.572}{645.508}{M}%
\TextAt{16}{474.166}{645.508}{\Character{10}}%
\TextAt{9}{486.486}{645.508}{\Character{40}}%
\TextAt{13}{491.038}{645.508}{N}%
\TextAt{16}{504.681}{645.508}{\Character{10}}%
\TextAt{13}{517.001}{645.508}{L}%
\TextAt{9}{524.965}{645.508}{\Character{41}}%
\TextAt{8}{529.518}{645.508}{,}%
\TextAt{8}{79.226}{628.462}{we}%
\TextAt{8}{96.181}{628.462}{find}%
\TextAt{8}{119.317}{628.462}{right}%
\TextAt{8}{147.342}{628.462}{away}%
\TextAt{8}{176.347}{628.462}{that}%
\TextAt{13}{201.444}{628.462}{F}%
\TextAt{17}{210.647}{632.8}{0}%
\TextAt{9}{215.38}{628.462}{\Character{40}}%
\Vector{q 0 G 0 g 0.478 w 147.932 -138.182 m 153.776 -138.182 l S Q}%
\TextAt{14}{219.932}{628.462}{o}%
\TextAt{9}{225.776}{628.462}{\Character{41}}%
\TextAt{8}{233.953}{628.462}{is}%
\TextAt{8}{245.444}{628.462}{a}%
\TextAt{8}{254.927}{628.462}{commutative}%
\TextAt{8}{324.58}{628.462}{ring}%
\TextAt{8}{348.374}{628.462}{and}%
\TextAt{13}{370.864}{628.462}{F}%
\TextAt{17}{378.441}{626.669}{0}%
\TextAt{9}{383.174}{628.462}{\Character{40}}%
\Vector{q 0 G 0 g 0.478 w 315.726 -138.182 m 321.57 -138.182 l S Q}%
\TextAt{14}{387.726}{628.462}{o}%
\TextAt{9}{393.57}{628.462}{\Character{41}}%
\TextAt{8}{401.747}{628.462}{is}%
\TextAt{8}{413.239}{628.462}{a}%
\TextAt{8}{422.721}{628.462}{subring.}%
\TextAt{8}{468.869}{628.462}{Furthermore}%
\TextAt{8}{79.226}{611.416}{the}%
\TextAt{8}{99.191}{611.416}{subgroups}%
\TextAt{13}{154.411}{611.416}{R}%
\TextAt{17}{163.321}{609.623}{0}%
\TextAt{9}{168.053}{611.416}{\Character{40}}%
\Vector{q 0 G 0 g 0.478 w 100.606 -155.228 m 106.45 -155.228 l S Q}%
\TextAt{14}{172.606}{611.416}{o}%
\TextAt{9}{178.45}{611.416}{\Character{41}}%
\NativeMathAt{186.323}{611.416}{11.955}{\subseteq}%
\TextAt{13}{198.942}{611.416}{F}%
\TextAt{17}{206.52}{609.623}{0}%
\TextAt{9}{211.252}{611.416}{\Character{40}}%
\Vector{q 0 G 0 g 0.478 w 143.805 -155.228 m 149.649 -155.228 l S Q}%
\TextAt{14}{215.805}{611.416}{o}%
\TextAt{9}{221.649}{611.416}{\Character{41}}%
\TextAt{8}{229.907}{611.416}{and}%
\TextAt{13}{252.477}{611.416}{R}%
\TextAt{17}{261.486}{615.754}{0}%
\TextAt{9}{266.218}{611.416}{\Character{40}}%
\Vector{q 0 G 0 g 0.478 w 198.771 -155.228 m 204.615 -155.228 l S Q}%
\TextAt{14}{270.771}{611.416}{o}%
\TextAt{9}{276.615}{611.416}{\Character{41}}%
\NativeMathAt{284.488}{611.416}{11.955}{\subseteq}%
\TextAt{13}{297.107}{611.416}{F}%
\TextAt{17}{306.311}{615.754}{0}%
\TextAt{9}{311.043}{611.416}{\Character{40}}%
\Vector{q 0 G 0 g 0.478 w 243.596 -155.228 m 249.44 -155.228 l S Q}%
\TextAt{14}{315.596}{611.416}{o}%
\TextAt{9}{321.44}{611.416}{\Character{41}}%
\TextAt{8}{329.697}{611.416}{turn}%
\TextAt{8}{355.52}{611.416}{out}%
\TextAt{8}{376.142}{611.416}{to}%
\TextAt{8}{390.26}{611.416}{be}%
\TextAt{13}{405.993}{611.416}{F}%
\TextAt{17}{413.57}{609.623}{0}%
\TextAt{9}{418.303}{611.416}{\Character{40}}%
\Vector{q 0 G 0 g 0.478 w 350.855 -155.228 m 356.699 -155.228 l S Q}%
\TextAt{14}{422.855}{611.416}{o}%
\TextAt{9}{428.699}{611.416}{\Character{41}}%
\TextAt{8}{433.252}{611.416}{-submodules.}%
\TextAt{8}{505.166}{611.416}{For}%
\TextAt{8}{525.943}{611.416}{if}%
\TextAt{9}{225.535}{582.415}{0}%
\NativeArrowAt{234.709}{582.415}{11.955}{19.26}{\longrightarrow}%
\TextAt{13}{257.291}{582.415}{M}%
\TextAt{3}{269.865}{587.351}{\Character{48}}%
\NativeArrowAt{275.981}{582.415}{11.955}{19.26}{\longrightarrow}%
\TextAt{13}{298.563}{582.415}{M}%
\NativeArrowAt{314.457}{582.415}{11.955}{19.26}{\longrightarrow}%
\TextAt{13}{337.039}{582.415}{M}%
\TextAt{3}{349.613}{587.351}{\Character{48}\Character{48}}%
\NativeArrowAt{358.025}{582.415}{11.955}{19.26}{\longrightarrow}%
\TextAt{9}{380.607}{582.415}{0}%
\TextAt{8}{79.226}{553.414}{is}%
\TextAt{8}{90.855}{553.414}{a}%
\TextAt{8}{100.476}{553.414}{short}%
\TextAt{8}{130.325}{553.414}{exact}%
\TextAt{8}{161.083}{553.414}{sequence}%
\TextAt{8}{209.45}{553.414}{of}%
\TextAt{8}{222.646}{553.414}{coherent}%
\TextAt{8}{269.662}{553.414}{metrized}%
\TextAt{14}{317.647}{553.414}{o}%
\TextAt{8}{323.491}{553.414}{-modules,}%
\TextAt{8}{376.441}{553.414}{and}%
\TextAt{13}{399.069}{553.414}{N}%
\TextAt{8}{413.454}{553.414}{is}%
\TextAt{8}{425.084}{553.414}{a}%
\TextAt{8}{434.704}{553.414}{projective}%
\TextAt{8}{488.547}{553.414}{metrized}%
\TextAt{14}{79.226}{536.368}{o}%
\TextAt{8}{85.07}{536.368}{-module,}%
\TextAt{8}{133.513}{536.368}{then}%
\TextAt{8}{160.173}{536.368}{it}%
\TextAt{8}{171.877}{536.368}{is}%
\TextAt{8}{183.641}{536.368}{clear}%
\TextAt{8}{211.604}{536.368}{that}%
\TextAt{9}{187.684}{507.366}{0}%
\NativeArrowAt{196.857}{507.366}{11.955}{19.26}{\longrightarrow}%
\TextAt{13}{219.439}{507.366}{N}%
\TextAt{16}{232.719}{507.366}{\Character{10}}%
\TextAt{13}{244.674}{507.366}{M}%
\TextAt{3}{257.247}{512.303}{\Character{48}}%
\NativeArrowAt{263.363}{507.366}{11.955}{19.26}{\longrightarrow}%
\TextAt{13}{285.945}{507.366}{N}%
\TextAt{16}{299.225}{507.366}{\Character{10}}%
\TextAt{13}{311.18}{507.366}{M}%
\NativeArrowAt{327.074}{507.366}{11.955}{19.26}{\longrightarrow}%
\TextAt{13}{349.656}{507.366}{N}%
\TextAt{16}{362.935}{507.366}{\Character{10}}%
\TextAt{13}{374.891}{507.366}{M}%
\TextAt{3}{387.464}{512.303}{\Character{48}\Character{48}}%
\NativeArrowAt{395.877}{507.366}{11.955}{19.26}{\longrightarrow}%
\TextAt{9}{418.459}{507.366}{0}%
\TextAt{8}{79.226}{478.365}{is}%
\TextAt{8}{91.34}{478.365}{a}%
\TextAt{8}{101.444}{478.365}{short}%
\TextAt{8}{131.778}{478.365}{exact}%
\TextAt{8}{163.019}{478.365}{sequence}%
\TextAt{8}{211.871}{478.365}{of}%
\TextAt{8}{225.55}{478.365}{metrized}%
\TextAt{14}{274.02}{478.365}{o}%
\TextAt{8}{279.864}{478.365}{-modules}%
\TextAt{8}{330.019}{478.365}{as}%
\TextAt{8}{344.738}{478.365}{well,}%
\TextAt{8}{372.158}{478.365}{so}%
\TextAt{8}{386.877}{478.365}{that,}%
\TextAt{8}{415.935}{478.365}{along}%
\TextAt{8}{447.512}{478.365}{with}%
\TextAt{8}{474.521}{478.365}{a}%
\TextAt{8}{484.626}{478.365}{generator}%
\TextAt{16}{79.226}{461.319}{\Character{102}}%
\TextAt{13}{85.204}{461.319}{M}%
\TextAt{3}{97.777}{465.658}{\Character{48}}%
\TextAt{16}{100.572}{461.319}{\Character{103}}%
\TextAt{16}{109.204}{461.319}{\Character{0}}%
\TextAt{16}{121.147}{461.319}{\Character{102}}%
\TextAt{13}{127.14}{461.319}{M}%
\TextAt{16}{139.713}{461.319}{\Character{103}}%
\TextAt{9}{148.347}{461.319}{\Character{43}}%
\TextAt{16}{160.109}{461.319}{\Character{102}}%
\TextAt{13}{166.086}{461.319}{M}%
\TextAt{3}{178.66}{465.658}{\Character{48}\Character{48}}%
\TextAt{16}{183.752}{461.319}{\Character{103}}%
\TextAt{8}{189.73}{461.319}{,}%
\TextAt{8}{196.879}{461.319}{the}%
\TextAt{8}{217.035}{461.319}{element}%
\TextAt{16}{133.928}{432.318}{\Character{102}}%
\TextAt{13}{139.906}{432.318}{N}%
\TextAt{16}{150.529}{432.318}{\Character{103}}%
\TextAt{9}{156.506}{432.318}{\Character{40}}%
\TextAt{16}{161.059}{432.318}{\Character{102}}%
\TextAt{13}{167.036}{432.318}{M}%
\TextAt{3}{179.61}{437.254}{\Character{48}}%
\TextAt{16}{182.405}{432.318}{\Character{103}}%
\TextAt{16}{191.037}{432.318}{\Character{0}}%
\TextAt{16}{202.98}{432.318}{\Character{102}}%
\TextAt{13}{208.972}{432.318}{M}%
\TextAt{16}{221.546}{432.318}{\Character{103}}%
\TextAt{9}{230.18}{432.318}{\Character{43}}%
\TextAt{16}{241.941}{432.318}{\Character{102}}%
\TextAt{13}{247.919}{432.318}{M}%
\TextAt{3}{260.492}{437.254}{\Character{48}\Character{48}}%
\TextAt{16}{265.584}{432.318}{\Character{103}}%
\TextAt{9}{271.562}{432.318}{\Character{41}}%
\TextAt{9}{279.417}{432.318}{\Character{61}}%
\TextAt{16}{291.861}{432.318}{\Character{102}}%
\TextAt{13}{297.838}{432.318}{N}%
\TextAt{16}{311.118}{432.318}{\Character{10}}%
\TextAt{13}{323.073}{432.318}{M}%
\TextAt{3}{335.646}{437.254}{\Character{48}}%
\TextAt{16}{338.441}{432.318}{\Character{103}}%
\TextAt{16}{347.073}{432.318}{\Character{0}}%
\TextAt{16}{359.016}{432.318}{\Character{102}}%
\TextAt{13}{365.008}{432.318}{N}%
\TextAt{16}{378.288}{432.318}{\Character{10}}%
\TextAt{13}{390.243}{432.318}{M}%
\TextAt{16}{402.817}{432.318}{\Character{103}}%
\TextAt{9}{411.451}{432.318}{\Character{43}}%
\TextAt{16}{423.212}{432.318}{\Character{102}}%
\TextAt{13}{429.19}{432.318}{N}%
\TextAt{16}{442.469}{432.318}{\Character{10}}%
\TextAt{13}{454.424}{432.318}{M}%
\TextAt{3}{466.998}{437.254}{\Character{48}\Character{48}}%
\TextAt{16}{472.09}{432.318}{\Character{103}}%
\TextAt{8}{79.226}{403.317}{will}%
\TextAt{8}{102.043}{403.317}{also}%
\TextAt{8}{126.235}{403.317}{belong}%
\TextAt{8}{164.343}{403.317}{to}%
\TextAt{13}{179.365}{403.317}{R}%
\TextAt{17}{188.275}{401.524}{0}%
\TextAt{9}{193.007}{403.317}{\Character{40}}%
\Vector{q 0 G 0 g 0.478 w 125.56 -363.326 m 131.404 -363.326 l S Q}%
\TextAt{14}{197.56}{403.317}{o}%
\TextAt{9}{203.404}{403.317}{\Character{41}}%
\TextAt{8}{207.956}{403.317}{,}%
\TextAt{8}{215.995}{403.317}{resp.}%
\TextAt{13}{247.453}{403.317}{R}%
\TextAt{17}{256.462}{407.655}{0}%
\TextAt{9}{261.194}{403.317}{\Character{40}}%
\Vector{q 0 G 0 g 0.478 w 193.746 -363.326 m 199.591 -363.326 l S Q}%
\TextAt{14}{265.746}{403.317}{o}%
\TextAt{9}{271.591}{403.317}{\Character{41}}%
\TextAt{8}{276.143}{403.317}{.}%
\TextAt{8}{286.727}{403.317}{This}%
\TextAt{8}{314.159}{403.317}{is}%
\TextAt{8}{326.635}{403.317}{why}%
\TextAt{13}{352.058}{403.317}{K}%
\TextAt{17}{362.017}{401.524}{0}%
\TextAt{9}{366.749}{403.317}{\Character{40}}%
\Vector{q 0 G 0 g 0.478 w 299.302 -363.326 m 305.146 -363.326 l S Q}%
\TextAt{14}{371.302}{403.317}{o}%
\TextAt{9}{377.146}{403.317}{\Character{41}}%
\TextAt{9}{386.22}{403.317}{\Character{61}}%
\TextAt{13}{399.871}{403.317}{F}%
\TextAt{17}{407.449}{401.524}{0}%
\TextAt{9}{412.181}{403.317}{\Character{40}}%
\Vector{q 0 G 0 g 0.478 w 344.733 -363.326 m 350.577 -363.326 l S Q}%
\TextAt{14}{416.733}{403.317}{o}%
\TextAt{9}{422.577}{403.317}{\Character{41}\Character{47}}%
\TextAt{13}{432.983}{403.317}{R}%
\TextAt{17}{441.893}{401.524}{0}%
\TextAt{9}{446.625}{403.317}{\Character{40}}%
\Vector{q 0 G 0 g 0.478 w 379.178 -363.326 m 385.022 -363.326 l S Q}%
\TextAt{14}{451.178}{403.317}{o}%
\TextAt{9}{457.022}{403.317}{\Character{41}}%
\TextAt{8}{466.183}{403.317}{is}%
\TextAt{8}{478.659}{403.317}{a}%
\TextAt{8}{489.127}{403.317}{ring}%
\TextAt{8}{513.904}{403.317}{and}%
\TextAt{13}{79.226}{386.271}{K}%
\TextAt{17}{90.044}{390.61}{0}%
\TextAt{9}{94.776}{386.271}{\Character{40}}%
\Vector{q 0 G 0 g 0.478 w 27.328 -380.372 m 33.173 -380.372 l S Q}%
\TextAt{14}{99.328}{386.271}{o}%
\TextAt{9}{105.173}{386.271}{\Character{41}}%
\TextAt{9}{113.028}{386.271}{\Character{61}}%
\TextAt{13}{125.471}{386.271}{F}%
\TextAt{17}{134.675}{390.61}{0}%
\TextAt{9}{139.407}{386.271}{\Character{40}}%
\Vector{q 0 G 0 g 0.478 w 71.96 -380.372 m 77.804 -380.372 l S Q}%
\TextAt{14}{143.96}{386.271}{o}%
\TextAt{9}{149.804}{386.271}{\Character{41}\Character{47}}%
\TextAt{13}{160.209}{386.271}{R}%
\TextAt{17}{169.218}{390.61}{0}%
\TextAt{9}{173.95}{386.271}{\Character{40}}%
\Vector{q 0 G 0 g 0.478 w 106.502 -380.372 m 112.347 -380.372 l S Q}%
\TextAt{14}{178.502}{386.271}{o}%
\TextAt{9}{184.347}{386.271}{\Character{41}}%
\TextAt{8}{192.796}{386.271}{is}%
\TextAt{8}{204.56}{386.271}{a}%
\TextAt{13}{214.316}{386.271}{K}%
\TextAt{17}{224.275}{384.478}{0}%
\TextAt{9}{229.007}{386.271}{\Character{40}}%
\Vector{q 0 G 0 g 0.478 w 161.559 -380.372 m 167.404 -380.372 l S Q}%
\TextAt{14}{233.559}{386.271}{o}%
\TextAt{9}{239.404}{386.271}{\Character{41}}%
\TextAt{8}{243.956}{386.271}{-module.}%
\TextAt{8}{97.159}{366.236}{Associating}%
\TextAt{8}{160.287}{366.236}{to}%
\TextAt{8}{175.165}{366.236}{the}%
\TextAt{8}{195.888}{366.236}{class}%
\TextAt{9}{223.892}{366.236}{\Character{91}}%
\TextAt{13}{227.144}{366.236}{M}%
\TextAt{9}{239.717}{366.236}{\Character{93}}%
\TextAt{8}{247.433}{366.236}{of}%
\TextAt{8}{261.33}{366.236}{a}%
\TextAt{8}{271.652}{366.236}{projective}%
\TextAt{14}{326.197}{366.236}{o}%
\TextAt{8}{332.041}{366.236}{-module}%
\TextAt{13}{377.799}{366.236}{M}%
\TextAt{8}{394.837}{366.236}{in}%
\TextAt{13}{409.056}{366.236}{K}%
\TextAt{17}{419.016}{364.443}{0}%
\TextAt{9}{423.748}{366.236}{\Character{40}}%
\Vector{q 0 G 0 g 0.478 w 356.3 -400.407 m 362.144 -400.407 l S Q}%
\TextAt{14}{428.3}{366.236}{o}%
\TextAt{9}{434.144}{366.236}{\Character{41}}%
\TextAt{8}{443.161}{366.236}{its}%
\TextAt{8}{460.047}{366.236}{class}%
\TextAt{8}{488.051}{366.236}{in}%
\TextAt{13}{502.271}{366.236}{K}%
\TextAt{17}{513.088}{370.575}{0}%
\TextAt{9}{517.821}{366.236}{\Character{40}}%
\Vector{q 0 G 0 g 0.478 w 450.373 -400.407 m 456.217 -400.407 l S Q}%
\TextAt{14}{522.373}{366.236}{o}%
\TextAt{9}{528.217}{366.236}{\Character{41}}%
\TextAt{8}{79.226}{349.19}{(which}%
\TextAt{8}{117.267}{349.19}{again}%
\TextAt{8}{148.494}{349.19}{is}%
\TextAt{8}{160.258}{349.19}{denoted}%
\TextAt{8}{204.48}{349.19}{by}%
\TextAt{9}{220.739}{349.19}{\Character{91}}%
\TextAt{13}{223.991}{349.19}{M}%
\TextAt{9}{236.565}{349.19}{\Character{93}}%
\TextAt{8}{239.816}{349.19}{),}%
\TextAt{8}{251.521}{349.19}{defines}%
\TextAt{8}{289.944}{349.19}{a}%
\TextAt{8}{299.7}{349.19}{homomorphism}%
\TextAt{13}{261.351}{320.189}{K}%
\TextAt{17}{271.31}{318.396}{0}%
\TextAt{9}{276.042}{320.189}{\Character{40}}%
\Vector{q 0 G 0 g 0.478 w 208.595 -446.454 m 214.439 -446.454 l S Q}%
\TextAt{14}{280.595}{320.189}{o}%
\TextAt{9}{286.439}{320.189}{\Character{41}}%
\NativeArrowAt{294.312}{320.189}{11.955}{19.26}{\longrightarrow}%
\TextAt{13}{316.894}{320.189}{K}%
\TextAt{17}{327.712}{325.125}{0}%
\TextAt{9}{332.444}{320.189}{\Character{40}}%
\Vector{q 0 G 0 g 0.478 w 264.996 -446.454 m 270.841 -446.454 l S Q}%
\TextAt{14}{336.996}{320.189}{o}%
\TextAt{9}{342.841}{320.189}{\Character{41}}%
\TextAt{13}{347.393}{320.189}{\Character{58}}%
\TextAt{8}{79.226}{291.188}{It}%
\TextAt{8}{91.753}{291.188}{is}%
\TextAt{8}{103.372}{291.188}{called}%
\TextAt{8}{136.39}{291.188}{the}%
\TextAt{12}{156.402}{291.188}{Poincaré}%
\TextAt{12}{211.635}{291.188}{homomorphism}%
\TextAt{8}{302.434}{291.188}{.}%
\TextAt{8}{310.834}{291.188}{We}%
\TextAt{8}{330.833}{291.188}{will}%
\TextAt{8}{352.793}{291.188}{show}%
\TextAt{8}{381.651}{291.188}{next}%
\TextAt{8}{407.843}{291.188}{that}%
\TextAt{8}{433.067}{291.188}{the}%
\TextAt{8}{453.078}{291.188}{Poincaré}%
\TextAt{8}{500.897}{291.188}{homo-}%
\TextAt{8}{79.226}{274.142}{morphism}%
\TextAt{8}{133.921}{274.142}{is}%
\TextAt{8}{145.685}{274.142}{an}%
\TextAt{8}{161.944}{274.142}{isomorphism.}%
\TextAt{8}{234.914}{274.142}{The}%
\TextAt{8}{258.968}{274.142}{proof}%
\TextAt{8}{289.537}{274.142}{is}%
\TextAt{8}{301.301}{274.142}{based}%
\TextAt{8}{333.879}{274.142}{on}%
\TextAt{8}{350.138}{274.142}{the}%
\TextAt{8}{370.295}{274.142}{following}%
\TextAt{8}{419.729}{274.142}{two}%
\TextAt{8}{441.846}{274.142}{lemmas.}%
\TextAt{12}{79.226}{248.13}{Lemma}%
\TextAt{12}{126.78}{248.13}{(5.2).}%
\TextAt{8}{161.998}{248.13}{All}%
\TextAt{8}{181.077}{248.13}{coherent}%
\TextAt{8}{228.131}{248.13}{metrized}%
\TextAt{14}{276.154}{248.13}{o}%
\TextAt{8}{281.998}{248.13}{-modules}%
\TextAt{13}{331.706}{248.13}{M}%
\TextAt{8}{348.08}{248.13}{admit}%
\TextAt{8}{381.804}{248.13}{a}%
\TextAt{8}{391.462}{248.13}{“metrized}%
\TextAt{8}{445.008}{248.13}{projective}%
\TextAt{8}{498.889}{248.13}{resolu-}%
\TextAt{8}{79.226}{231.084}{tion”,}%
\TextAt{8}{112.067}{231.084}{i.e.,}%
\TextAt{8}{134.172}{231.084}{a}%
\TextAt{8}{143.928}{231.084}{short}%
\TextAt{8}{173.911}{231.084}{exact}%
\TextAt{8}{204.803}{231.084}{sequence}%
\TextAt{9}{232.767}{202.083}{0}%
\NativeArrowAt{241.941}{202.083}{11.955}{19.26}{\longrightarrow}%
\TextAt{13}{264.523}{202.083}{E}%
\NativeArrowAt{277.21}{202.083}{11.955}{19.26}{\longrightarrow}%
\TextAt{13}{299.792}{202.083}{F}%
\NativeArrowAt{312.317}{202.083}{11.955}{19.26}{\longrightarrow}%
\TextAt{13}{334.899}{202.083}{M}%
\NativeArrowAt{350.793}{202.083}{11.955}{19.26}{\longrightarrow}%
\TextAt{9}{373.375}{202.083}{0}%
\TextAt{8}{79.226}{173.082}{of}%
\TextAt{8}{92.556}{173.082}{metrized}%
\TextAt{14}{140.676}{173.082}{o}%
\TextAt{8}{146.52}{173.082}{-modules}%
\TextAt{8}{196.325}{173.082}{in}%
\TextAt{8}{209.978}{173.082}{which}%
\TextAt{13}{243.464}{173.082}{E}%
\TextAt{8}{256.728}{173.082}{and}%
\TextAt{13}{279.491}{173.082}{F}%
\TextAt{8}{292.592}{173.082}{are}%
\TextAt{8}{312.102}{173.082}{projective.}%
\TextAt{19}{79.226}{138.103}{Proof.}%
\TextAt{8}{116}{138.103}{If}%
\TextAt{13}{127.981}{138.103}{\Character{11}}%
\TextAt{17}{135.427}{136.31}{1}%
\TextAt{13}{140.159}{138.103}{\Character{59}}%
\TextAt{13}{145.383}{138.103}{\Character{58}}%
\TextAt{13}{150.62}{138.103}{\Character{58}}%
\TextAt{13}{155.856}{138.103}{\Character{58}}%
\TextAt{13}{161.081}{138.103}{\Character{59}}%
\TextAt{13}{166.317}{138.103}{\Character{11}}%
\TextAt{21}{173.826}{136.31}{n}%
\TextAt{8}{183.649}{138.103}{is}%
\TextAt{8}{195.702}{138.103}{a}%
\TextAt{8}{205.746}{138.103}{system}%
\TextAt{8}{244.854}{138.103}{of}%
\TextAt{8}{258.473}{138.103}{generators}%
\TextAt{8}{315.417}{138.103}{of}%
\TextAt{13}{329.037}{138.103}{M}%
\TextAt{8}{341.61}{138.103}{,}%
\TextAt{8}{349.121}{138.103}{and}%
\TextAt{13}{372.172}{138.103}{F}%
\TextAt{8}{385.563}{138.103}{is}%
\TextAt{8}{397.615}{138.103}{the}%
\TextAt{8}{418.061}{138.103}{free}%
\TextAt{14}{440.778}{138.103}{o}%
\TextAt{8}{446.622}{138.103}{-module}%
\TextAt{13}{492.102}{138.103}{F}%
\TextAt{9}{505.119}{138.103}{\Character{61}}%
\TextAt{14}{518.037}{138.103}{o}%
\TextAt{21}{523.881}{142.441}{n}%
\TextAt{8}{529.518}{138.103}{,}%
\TextAt{8}{79.226}{121.057}{then}%
\TextAt{13}{203.769}{100.37}{F}%
\NativeArrowAt{216.293}{100.37}{11.955}{19.26}{\longrightarrow}%
\TextAt{13}{238.875}{100.37}{M}%
\TextAt{13}{250.794}{100.37}{\Character{59}}%
\TextAt{9}{279.953}{100.37}{\Character{40}}%
\TextAt{13}{284.505}{100.37}{x}%
\TextAt{17}{291.157}{98.577}{1}%
\TextAt{13}{295.89}{100.37}{\Character{59}}%
\TextAt{13}{301.114}{100.37}{\Character{58}}%
\TextAt{13}{306.351}{100.37}{\Character{58}}%
\TextAt{13}{311.575}{100.37}{\Character{58}}%
\TextAt{13}{316.812}{100.37}{\Character{59}}%
\TextAt{13}{322.048}{100.37}{x}%
\TextAt{21}{328.762}{98.577}{n}%
\TextAt{9}{334.399}{100.37}{\Character{41}}%
\NativeArrowAt{342.272}{100.37}{11.955}{19.26}{\longmapsto}%
\TextAt{21}{370.919}{115.314}{n}%
\TextAt{20}{364.854}{111.727}{\Character{88}}%
\TextAt{21}{366.636}{86.533}{i}%
\TextAt{17}{369.52}{86.533}{\Character{61}1}%
\TextAt{13}{384.115}{100.37}{x}%
\TextAt{21}{390.767}{98.577}{i}%
\TextAt{13}{394.148}{100.37}{\Character{11}}%
\TextAt{21}{401.594}{98.577}{i}%
\TextAt{13}{404.975}{100.37}{\Character{59}}%
\TextAt{8}{297.211}{43.033}{215}%
\Destination{78.23}{809.815}{nav.page.215}%
\Destination{79.36}{735.555}{nav.print.234}%
\Destination{79.36}{262.155}{nav.number.3.5.2}%
}

% Page 228 (body page 216)
\Page{%
\TextAt{2}{79.226}{796.943}{III.}%
\TextAt{2}{98.317}{796.943}{Riemann–Roch}%
\TextAt{2}{174.375}{796.943}{Theory}%
\Vector{q 0 G 0 g 0.299 w 7.226 18.729 m 460.769 18.729 l S Q}%
\TextAt{8}{79.226}{761.124}{is}%
\TextAt{8}{89.825}{761.124}{a}%
\TextAt{8}{98.415}{761.124}{surjective}%
\TextAt{14}{149.662}{761.124}{o}%
\TextAt{8}{155.506}{761.124}{-module}%
\TextAt{8}{199.531}{761.124}{homomorphism.}%
\TextAt{8}{286.364}{761.124}{Its}%
\TextAt{8}{302.486}{761.124}{kernel}%
\TextAt{13}{335.787}{761.124}{E}%
\TextAt{8}{347.886}{761.124}{is}%
\TextAt{8}{358.485}{761.124}{torsion}%
\TextAt{8}{396.413}{761.124}{free,}%
\TextAt{8}{421.16}{761.124}{and}%
\TextAt{8}{442.758}{761.124}{hence}%
\TextAt{8}{474.099}{761.124}{a}%
\TextAt{8}{482.689}{761.124}{projective}%
\TextAt{14}{79.226}{744.078}{o}%
\TextAt{8}{85.07}{744.078}{-module}%
\TextAt{8}{130.261}{744.078}{by}%
\TextAt{8}{146.52}{744.078}{(4.3).}%
\TextAt{8}{179.046}{744.078}{In}%
\TextAt{8}{193.667}{744.078}{the}%
\TextAt{8}{213.824}{744.078}{exact}%
\TextAt{8}{244.716}{744.078}{sequence}%
\TextAt{9}{227.201}{715.24}{0}%
\NativeArrowAt{236.375}{715.24}{11.955}{19.26}{\longrightarrow}%
\TextAt{13}{258.957}{715.24}{E}%
\TextAt{25}{267.622}{713.447}{C}%
\NativeArrowAt{277.197}{715.24}{11.955}{19.26}{\longrightarrow}%
\TextAt{13}{299.779}{715.24}{F}%
\TextAt{25}{307.357}{713.447}{C}%
\TextAt{21}{319.731}{722.504}{f}%
\NativeArrowAt{316.932}{715.24}{11.955}{11.955}{\longrightarrow}%
\TextAt{13}{332.208}{715.24}{M}%
\TextAt{25}{343.533}{713.447}{C}%
\NativeArrowAt{353.108}{715.24}{11.955}{19.26}{\longrightarrow}%
\TextAt{9}{375.69}{715.24}{0}%
\TextAt{13}{381.543}{715.24}{\Character{59}}%
\ErrataAt{79.19}{686.402}{455.81}{we choose an $F$-equivariant section $s:M_{\mathbb C}\to F_{\mathbb C}$ of $f$, so that $F_{\mathbb C}=E_{\mathbb C}\oplus sM_{\mathbb C}$. We give}%
\ErrataAt{79.19}{669.356}{455.81}{$F_{\mathbb C}$ the transported metric on $sM_{\mathbb C}$ and an $F$-invariant Hermitian metric on $E_{\mathbb C}$. This}%
\TextAt{8}{79.226}{652.31}{makes}%
\TextAt{9}{113.865}{652.31}{0}%
\NativeMathAt{123.039}{652.31}{11.955}{\rightarrow}%
\TextAt{13}{138.315}{652.31}{E}%
\NativeMathAt{151.002}{652.31}{11.955}{\rightarrow}%
\TextAt{13}{166.278}{652.31}{F}%
\NativeMathAt{178.803}{652.31}{11.955}{\rightarrow}%
\TextAt{13}{194.079}{652.31}{M}%
\NativeMathAt{209.973}{652.31}{11.955}{\rightarrow}%
\TextAt{9}{225.249}{652.31}{0}%
\TextAt{8}{234.455}{652.31}{into}%
\TextAt{8}{257.653}{652.31}{a}%
\TextAt{8}{266.864}{652.31}{short}%
\TextAt{8}{296.303}{652.31}{exact}%
\TextAt{8}{326.65}{652.31}{sequence}%
\TextAt{8}{374.608}{652.31}{of}%
\TextAt{8}{387.393}{652.31}{metrized}%
\TextAt{14}{434.968}{652.31}{o}%
\TextAt{8}{440.812}{652.31}{-modules}%
\TextAt{8}{490.072}{652.31}{in}%
\TextAt{8}{503.18}{652.31}{which}%
\TextAt{13}{79.226}{635.264}{E}%
\TextAt{8}{92.49}{635.264}{and}%
\TextAt{13}{115.252}{635.264}{F}%
\TextAt{8}{128.353}{635.264}{are}%
\TextAt{8}{147.864}{635.264}{projective.}%
\Vector{q 0 G 0 g 0.399 w 452.334 -138.736 m 452.334 -130.666 l S Q}%
\Vector{q 0 G 0 g 0.399 w 452.534 -130.865 m 459.707 -130.865 l S Q}%
\Vector{q 0 G 0 g 0.399 w 452.534 -138.536 m 459.707 -138.536 l S Q}%
\Vector{q 0 G 0 g 0.399 w 459.906 -138.736 m 459.906 -130.666 l S Q}%
\TextAt{8}{97.159}{606.31}{In}%
\TextAt{8}{111.78}{606.31}{a}%
\TextAt{8}{121.536}{606.31}{diagram}%
\TextAt{8}{167.073}{606.31}{of}%
\TextAt{8}{180.403}{606.31}{metrized}%
\TextAt{8}{228.522}{606.31}{projective}%
\TextAt{8}{282.5}{606.31}{resolutions}%
\TextAt{8}{341.164}{606.31}{of}%
\TextAt{13}{354.494}{606.31}{M}%
\TextAt{9}{175.088}{578.522}{0}%
\TextAt{13}{235.683}{578.522}{E}%
\TextAt{13}{301.188}{578.522}{F}%
\TextAt{13}{365.133}{578.522}{M}%
\TextAt{9}{431.051}{578.522}{0}%
\TextAt{9}{175.088}{528.717}{0}%
\TextAt{13}{234.285}{528.717}{E}%
\TextAt{3}{243.651}{533.056}{\Character{48}}%
\TextAt{13}{299.791}{528.717}{F}%
\TextAt{3}{308.995}{533.056}{\Character{48}}%
\TextAt{13}{365.133}{528.717}{M}%
\TextAt{9}{431.051}{528.717}{0}%
\NativeDiagram{%
\path[draw=black,line width=0.478bp] (114.294,-192.489) -- (157.854,-192.489);%
\path[draw=black,line width=0.478bp,line cap=round,line join=round] (155.606,-189.62) .. controls (156.061,-191.342) and (157.073,-192.154) .. (158.093,-192.489) .. controls (157.073,-192.824) and (156.061,-193.637) .. (155.606,-195.358);%
\path[draw=black,line width=0.478bp] (178.402,-192.489) -- (223.36,-192.489);%
\path[draw=black,line width=0.478bp,line cap=round,line join=round] (221.112,-189.62) .. controls (221.567,-191.342) and (222.579,-192.154) .. (223.599,-192.489) .. controls (222.579,-192.824) and (221.567,-193.637) .. (221.112,-195.358);%
\path[draw=black,line width=0.478bp] (168.367,-200.057) -- (168.367,-231.46);%
\path[draw=black,line width=0.478bp,line cap=round,line join=round] (171.236,-229.212) .. controls (169.515,-229.667) and (168.702,-230.679) .. (168.367,-231.699) .. controls (168.032,-230.679) and (167.219,-229.667) .. (165.498,-229.212);%
\path[draw=black,line width=0.478bp] (243.746,-192.489) -- (287.306,-192.489);%
\path[draw=black,line width=0.478bp,line cap=round,line join=round] (285.058,-189.62) .. controls (285.512,-191.342) and (286.525,-192.154) .. (287.545,-192.489) .. controls (286.525,-192.824) and (285.512,-193.637) .. (285.058,-195.358);%
\path[draw=black,line width=0.478bp] (233.792,-200.057) -- (233.792,-231.46);%
\path[draw=black,line width=0.478bp,line cap=round,line join=round] (236.661,-229.212) .. controls (234.939,-229.667) and (234.127,-230.679) .. (233.792,-231.699) .. controls (233.457,-230.679) and (232.644,-229.667) .. (230.923,-229.212);%
\path[draw=black,line width=0.478bp] (311.062,-192.489) -- (353.224,-192.489);%
\path[draw=black,line width=0.478bp,line cap=round,line join=round] (350.977,-189.62) .. controls (351.431,-191.342) and (352.443,-192.154) .. (353.463,-192.489) .. controls (352.443,-192.824) and (351.431,-193.637) .. (350.977,-195.358);%
\path[draw=black,line width=2.797bp] (299.423,-200.057) -- (299.423,-232.535);%
\path[draw=black!0,line width=1.841bp] (299.423,-200.057) -- (299.423,-232.535);%
\path[draw=black,line width=0.478bp] (114.294,-242.295) -- (156.456,-242.295);%
\path[draw=black,line width=0.478bp,line cap=round,line join=round] (154.209,-239.425) .. controls (154.663,-241.147) and (155.675,-241.96) .. (156.695,-242.295) .. controls (155.675,-242.629) and (154.663,-243.442) .. (154.209,-245.164);%
\path[draw=black,line width=0.478bp] (179.8,-242.295) -- (221.962,-242.295);%
\path[draw=black,line width=0.478bp,line cap=round,line join=round] (219.715,-239.425) .. controls (220.169,-241.147) and (221.181,-241.96) .. (222.201,-242.295) .. controls (221.181,-242.629) and (220.169,-243.442) .. (219.715,-245.164);%
\path[draw=black,line width=0.478bp] (245.143,-242.295) -- (287.306,-242.295);%
\path[draw=black,line width=0.478bp,line cap=round,line join=round] (285.058,-239.425) .. controls (285.512,-241.147) and (286.525,-241.96) .. (287.545,-242.295) .. controls (286.525,-242.629) and (285.512,-243.442) .. (285.058,-245.164);%
\path[draw=black,line width=0.478bp] (311.062,-242.295) -- (353.224,-242.295);%
\path[draw=black,line width=0.478bp,line cap=round,line join=round] (350.977,-239.425) .. controls (351.431,-241.147) and (352.443,-241.96) .. (353.463,-242.295) .. controls (352.443,-242.629) and (351.431,-243.442) .. (350.977,-245.164);%
}%
\TextAt{8}{79.226}{503.133}{the}%
\TextAt{8}{99.157}{503.133}{resolution}%
\TextAt{8}{152.981}{503.133}{in}%
\TextAt{8}{166.408}{503.133}{the}%
\TextAt{8}{186.339}{503.133}{top}%
\TextAt{8}{206.927}{503.133}{line}%
\TextAt{8}{228.807}{503.133}{will}%
\TextAt{8}{250.686}{503.133}{be}%
\TextAt{8}{266.385}{503.133}{called}%
\TextAt{19}{299.323}{503.133}{dominant}%
\TextAt{8}{351.605}{503.133}{if}%
\TextAt{8}{362.103}{503.133}{the}%
\TextAt{8}{382.034}{503.133}{vertical}%
\TextAt{8}{423.436}{503.133}{arrows}%
\TextAt{8}{460.678}{503.133}{are}%
\TextAt{8}{479.963}{503.133}{admissible}%
\TextAt{8}{79.226}{486.087}{epimorphisms.}%
\TextAt{12}{79.226}{460.121}{Lemma}%
\TextAt{12}{126.891}{460.121}{(5.3).}%
\TextAt{8}{162.207}{460.121}{Let}%
\TextAt{9}{141.023}{429.739}{0}%
\NativeArrowAt{150.197}{429.739}{11.955}{19.26}{\longrightarrow}%
\TextAt{13}{172.779}{429.739}{E}%
\TextAt{3}{182.145}{434.676}{\Character{48}}%
\NativeArrowAt{188.261}{429.739}{11.955}{19.26}{\longrightarrow}%
\TextAt{13}{210.843}{429.739}{F}%
\TextAt{3}{220.047}{434.676}{\Character{48}}%
\TextAt{21}{228.75}{437.003}{f}%
\TextAt{23}{233.697}{439.816}{\Character{48}}%
\NativeArrowAt{226.163}{429.739}{11.955}{14.215}{\longrightarrow}%
\TextAt{13}{243.701}{429.739}{M}%
\NativeArrowAt{259.595}{429.739}{11.955}{19.26}{\longrightarrow}%
\TextAt{9}{282.178}{429.739}{0}%
\TextAt{13}{288.031}{429.739}{\Character{59}}%
\TextAt{9}{317.185}{429.739}{0}%
\NativeArrowAt{326.359}{429.739}{11.955}{19.26}{\longrightarrow}%
\TextAt{13}{348.941}{429.739}{E}%
\TextAt{3}{358.307}{434.676}{\Character{48}\Character{48}}%
\NativeArrowAt{366.72}{429.739}{11.955}{19.26}{\longrightarrow}%
\TextAt{13}{389.302}{429.739}{F}%
\TextAt{3}{398.506}{434.676}{\Character{48}\Character{48}}%
\TextAt{21}{409.506}{437.003}{f}%
\TextAt{23}{414.453}{439.816}{\Character{48}\Character{48}}%
\NativeArrowAt{406.918}{429.739}{11.955}{16.402}{\longrightarrow}%
\TextAt{13}{426.643}{429.739}{M}%
\NativeArrowAt{442.538}{429.739}{11.955}{19.26}{\longrightarrow}%
\TextAt{9}{465.12}{429.739}{0}%
\TextAt{8}{79.226}{400.902}{be}%
\TextAt{8}{96.196}{400.902}{two}%
\TextAt{8}{119.359}{400.902}{metrized}%
\TextAt{8}{168.525}{400.902}{projective}%
\TextAt{8}{223.549}{400.902}{resolutions}%
\TextAt{8}{283.259}{400.902}{of}%
\TextAt{8}{297.635}{400.902}{the}%
\TextAt{8}{318.837}{400.902}{metrized}%
\TextAt{14}{368.003}{400.902}{o}%
\TextAt{8}{373.847}{400.902}{-module}%
\TextAt{13}{420.083}{400.902}{M}%
\TextAt{8}{432.657}{400.902}{.}%
\TextAt{8}{444.243}{400.902}{Then,}%
\TextAt{8}{479.36}{400.902}{taking}%
\TextAt{8}{516.51}{400.902}{the}%
\TextAt{14}{79.226}{383.856}{o}%
\TextAt{8}{85.07}{383.856}{-module}%
\TextAt{13}{165.405}{366.81}{F}%
\TextAt{9}{177.93}{366.81}{\Character{61}}%
\TextAt{13}{190.355}{366.81}{F}%
\TextAt{3}{199.559}{371.746}{\Character{48}}%
\TextAt{16}{205.011}{366.81}{\Character{2}}%
\TextAt{21}{214.309}{365.016}{M}%
\TextAt{13}{226.446}{366.81}{F}%
\TextAt{3}{235.649}{371.746}{\Character{48}\Character{48}}%
\TextAt{9}{244.062}{366.81}{\Character{61}}%
\TextAt{16}{256.488}{366.81}{\Character{102}}%
\TextAt{9}{262.465}{366.81}{\Character{40}}%
\TextAt{13}{267.018}{366.81}{x}%
\TextAt{3}{273.67}{371.746}{\Character{48}}%
\TextAt{13}{276.465}{366.81}{\Character{59}}%
\TextAt{13}{281.689}{366.81}{x}%
\TextAt{3}{288.361}{371.746}{\Character{48}\Character{48}}%
\TextAt{9}{293.453}{366.81}{\Character{41}}%
\NativeMathAt{301.326}{366.81}{11.955}{\in}%
\TextAt{13}{312.617}{366.81}{F}%
\TextAt{3}{321.821}{371.746}{\Character{48}}%
\TextAt{16}{327.272}{366.81}{\Character{2}}%
\TextAt{13}{339.228}{366.81}{F}%
\TextAt{3}{348.431}{371.746}{\Character{48}\Character{48}}%
\TextAt{16}{356.844}{366.81}{\Character{106}}%
\TextAt{13}{363.486}{366.81}{f}%
\TextAt{3}{370.532}{371.746}{\Character{48}}%
\TextAt{9}{373.327}{366.81}{\Character{40}}%
\TextAt{13}{377.88}{366.81}{x}%
\TextAt{3}{384.532}{371.746}{\Character{48}}%
\TextAt{9}{387.327}{366.81}{\Character{41}}%
\TextAt{9}{395.182}{366.81}{\Character{61}}%
\TextAt{13}{407.625}{366.81}{f}%
\TextAt{3}{414.672}{371.746}{\Character{48}\Character{48}}%
\TextAt{9}{419.764}{366.81}{\Character{40}}%
\TextAt{13}{424.316}{366.81}{x}%
\TextAt{3}{430.968}{371.746}{\Character{48}\Character{48}}%
\TextAt{9}{436.06}{366.81}{\Character{41}}%
\TextAt{16}{440.613}{366.81}{\Character{103}}%
\TextAt{8}{79.226}{343.358}{and}%
\TextAt{8}{102.588}{343.358}{the}%
\TextAt{8}{123.344}{343.358}{mapping}%
\TextAt{13}{172.076}{343.358}{f}%
\TextAt{9}{183.465}{343.358}{\Character{58}}%
\TextAt{13}{191.059}{343.358}{F}%
\NativeMathAt{204.605}{343.358}{11.955}{\rightarrow}%
\TextAt{13}{220.903}{343.358}{M}%
\TextAt{8}{233.477}{343.358}{,}%
\TextAt{9}{241.376}{343.358}{\Character{40}}%
\TextAt{13}{245.928}{343.358}{x}%
\TextAt{3}{252.58}{347.696}{\Character{48}}%
\TextAt{13}{255.375}{343.358}{\Character{59}}%
\TextAt{13}{260.599}{343.358}{x}%
\TextAt{3}{267.271}{347.696}{\Character{48}\Character{48}}%
\TextAt{9}{272.363}{343.358}{\Character{41}}%
\NativeArrowAt{281.258}{343.358}{11.955}{11.955}{\mapsto}%
\TextAt{13}{297.556}{343.358}{f}%
\TextAt{3}{304.603}{347.696}{\Character{48}}%
\TextAt{9}{307.398}{343.358}{\Character{40}}%
\TextAt{13}{311.95}{343.358}{x}%
\TextAt{3}{318.602}{347.696}{\Character{48}}%
\TextAt{9}{321.397}{343.358}{\Character{41}}%
\TextAt{9}{330.28}{343.358}{\Character{61}}%
\TextAt{13}{343.74}{343.358}{f}%
\TextAt{3}{350.786}{347.696}{\Character{48}\Character{48}}%
\TextAt{9}{355.878}{343.358}{\Character{40}}%
\TextAt{13}{360.43}{343.358}{x}%
\TextAt{3}{367.082}{347.696}{\Character{48}\Character{48}}%
\TextAt{9}{372.174}{343.358}{\Character{41}}%
\TextAt{8}{376.727}{343.358}{,}%
\TextAt{8}{384.625}{343.358}{one}%
\TextAt{8}{406.685}{343.358}{obtains}%
\TextAt{8}{448.326}{343.358}{a}%
\TextAt{8}{458.681}{343.358}{third}%
\TextAt{8}{488.547}{343.358}{metrized}%
\TextAt{8}{79.226}{326.312}{projective}%
\TextAt{8}{133.204}{326.312}{resolution}%
\TextAt{9}{232.767}{309.266}{0}%
\NativeArrowAt{241.941}{309.266}{11.955}{19.26}{\longrightarrow}%
\TextAt{13}{264.523}{309.266}{E}%
\NativeArrowAt{277.21}{309.266}{11.955}{19.26}{\longrightarrow}%
\TextAt{13}{299.792}{309.266}{F}%
\NativeArrowAt{312.317}{309.266}{11.955}{19.26}{\longrightarrow}%
\TextAt{13}{334.899}{309.266}{M}%
\NativeArrowAt{350.793}{309.266}{11.955}{19.26}{\longrightarrow}%
\TextAt{9}{373.375}{309.266}{0}%
\TextAt{8}{79.226}{285.814}{with}%
\TextAt{8}{105.886}{285.814}{kernel}%
\TextAt{13}{140.353}{285.814}{E}%
\TextAt{9}{153.04}{285.814}{\Character{61}}%
\TextAt{13}{165.465}{285.814}{E}%
\TextAt{3}{174.832}{290.153}{\Character{48}}%
\TextAt{16}{180.283}{285.814}{\Character{2}}%
\TextAt{13}{192.239}{285.814}{E}%
\TextAt{3}{201.605}{290.153}{\Character{48}\Character{48}}%
\TextAt{8}{210.594}{285.814}{which}%
\TextAt{8}{244.081}{285.814}{dominates}%
\TextAt{8}{300.079}{285.814}{both}%
\TextAt{8}{327.719}{285.814}{given}%
\TextAt{8}{358.288}{285.814}{ones.}%
\TextAt{19}{79.226}{250.929}{Proof.}%
\TextAt{8}{116}{250.929}{Since}%
\TextAt{13}{146.793}{250.929}{F}%
\TextAt{3}{155.997}{255.268}{\Character{48}}%
\TextAt{16}{161.609}{250.929}{\Character{8}}%
\TextAt{13}{173.724}{250.929}{F}%
\TextAt{3}{182.928}{255.268}{\Character{48}\Character{48}}%
\TextAt{8}{192.152}{250.929}{is}%
\TextAt{8}{204.152}{250.929}{projective,}%
\TextAt{8}{261.675}{250.929}{so}%
\TextAt{8}{276.28}{250.929}{is}%
\TextAt{13}{288.28}{250.929}{F}%
\TextAt{8}{297.483}{250.929}{,}%
\TextAt{8}{304.926}{250.929}{being}%
\TextAt{8}{336.699}{250.929}{the}%
\TextAt{8}{357.091}{250.929}{kernel}%
\TextAt{8}{391.793}{250.929}{of}%
\TextAt{8}{405.358}{250.929}{the}%
\TextAt{8}{425.75}{250.929}{homomorphism}%
\TextAt{13}{508.655}{250.929}{F}%
\TextAt{3}{517.859}{255.268}{\Character{48}}%
\TextAt{16}{523.471}{250.929}{\Character{8}}%
\TextAt{13}{79.226}{233.883}{F}%
\TextAt{3}{88.43}{238.222}{\Character{48}\Character{48}}%
\TextAt{21}{100.512}{241.147}{f}%
\TextAt{23}{105.459}{243.96}{\Character{48}}%
\TextAt{3}{108.143}{241.147}{\Character{0}}%
\TextAt{21}{114.729}{241.147}{f}%
\TextAt{23}{119.676}{243.96}{\Character{48}\Character{48}}%
\NativeArrowAt{97.924}{233.883}{11.955}{30.605}{\longrightarrow}%
\TextAt{13}{132.948}{233.883}{M}%
\TextAt{8}{145.522}{233.883}{.}%
\TextAt{8}{155.874}{233.883}{Thus}%
\TextAt{13}{186.158}{233.883}{E}%
\TextAt{8}{200.056}{233.883}{is}%
\TextAt{8}{212.455}{233.883}{also}%
\TextAt{8}{236.569}{233.883}{projective,}%
\TextAt{8}{294.592}{233.883}{being}%
\TextAt{8}{326.764}{233.883}{the}%
\TextAt{8}{347.556}{233.883}{kernel}%
\TextAt{8}{382.657}{233.883}{of}%
\TextAt{13}{396.622}{233.883}{F}%
\NativeMathAt{410.228}{233.883}{11.955}{\rightarrow}%
\TextAt{13}{426.585}{233.883}{M}%
\TextAt{8}{439.159}{233.883}{.}%
\TextAt{8}{449.512}{233.883}{We}%
\TextAt{8}{470.291}{233.883}{consider}%
\TextAt{8}{516.51}{233.883}{the}%
\TextAt{8}{79.226}{216.837}{commutative}%
\TextAt{8}{149.152}{216.837}{diagram}%
\TextAt{9}{168.431}{182.442}{0}%
\TextAt{13}{227.628}{182.442}{E}%
\TextAt{3}{236.995}{186.78}{\Character{48}}%
\TextAt{25}{236.294}{179.377}{C}%
\TextAt{13}{296.124}{182.442}{F}%
\TextAt{3}{305.328}{186.78}{\Character{48}}%
\TextAt{25}{303.702}{179.377}{C}%
\TextAt{13}{363.531}{182.442}{M}%
\TextAt{25}{374.856}{180.649}{C}%
\TextAt{9}{436.08}{182.442}{0}%
\TextAt{9}{168.431}{130.169}{0}%
\TextAt{13}{227.628}{130.169}{E}%
\TextAt{25}{236.294}{128.376}{C}%
\TextAt{13}{296.124}{130.169}{F}%
\TextAt{25}{303.702}{128.376}{C}%
\TextAt{13}{363.531}{130.169}{M}%
\TextAt{25}{374.856}{128.376}{C}%
\TextAt{9}{436.08}{130.169}{0}%
\TextAt{9}{168.431}{78.571}{0}%
\TextAt{13}{227.628}{78.571}{E}%
\TextAt{3}{236.995}{82.91}{\Character{48}\Character{48}}%
\TextAt{25}{236.294}{75.506}{C}%
\TextAt{13}{295.892}{78.571}{F}%
\TextAt{3}{305.096}{82.91}{\Character{48}\Character{48}}%
\TextAt{25}{303.47}{75.506}{C}%
\TextAt{13}{363.531}{78.571}{M}%
\TextAt{25}{374.856}{76.778}{C}%
\TextAt{9}{434.455}{78.571}{0}%
\TextAt{13}{440.308}{78.571}{\Character{59}}%
\NativeDiagram{%
\path[draw=black,line width=0.478bp] (107.637,-588.568) -- (149.8,-588.568);%
\path[draw=black,line width=0.478bp,line cap=round,line join=round] (147.552,-585.699) .. controls (148.007,-587.42) and (149.019,-588.233) .. (150.039,-588.568) .. controls (149.019,-588.903) and (148.007,-589.716) .. (147.552,-591.437);%
\path[draw=black,line width=0.478bp] (175.902,-588.568) -- (218.296,-588.568);%
\path[draw=black,line width=0.478bp,line cap=round,line join=round] (216.049,-585.699) .. controls (216.503,-587.42) and (217.515,-588.233) .. (218.535,-588.568) .. controls (217.515,-588.903) and (216.503,-589.716) .. (216.049,-591.437);%
\path[draw=black,line width=0.478bp] (243.311,-585.476) -- (285.705,-585.476);%
\path[draw=black,line width=0.478bp,line cap=round,line join=round] (283.457,-582.607) .. controls (283.912,-584.329) and (284.924,-585.141) .. (285.944,-585.476) .. controls (284.924,-585.811) and (283.912,-586.624) .. (283.457,-588.345);%
}%
\TextAt{21}{332.931}{192.889}{f}%
\TextAt{23}{337.878}{195.702}{\Character{48}}%
\NativeDiagram{%
\path[draw=black,line width=0.478bp,dash pattern=on 3.347bp off 1.913bp] (286.183,-591.66) -- (243.789,-591.66);%
\path[draw=black,line width=0.478bp,line cap=round,line join=round] (246.036,-594.529) .. controls (245.582,-592.807) and (244.57,-591.994) .. (243.55,-591.66) .. controls (244.57,-591.325) and (245.582,-590.512) .. (246.036,-588.79);%
}%
\TextAt{21}{333.447}{173.391}{s}%
\TextAt{23}{337.363}{176.204}{\Character{48}}%
\NativeDiagram{%
\path[draw=black,line width=0.478bp] (314.466,-588.568) -- (358.255,-588.568);%
\path[draw=black,line width=0.478bp,line cap=round,line join=round] (356.007,-585.699) .. controls (356.462,-587.42) and (357.474,-588.233) .. (358.494,-588.568) .. controls (357.474,-588.903) and (356.462,-589.716) .. (356.007,-591.437);%
\path[draw=black,line width=0.478bp] (107.637,-640.841) -- (149.8,-640.841);%
\path[draw=black,line width=0.478bp,line cap=round,line join=round] (147.552,-637.972) .. controls (148.007,-639.694) and (149.019,-640.507) .. (150.039,-640.841) .. controls (149.019,-641.176) and (148.007,-641.989) .. (147.552,-643.71);%
\path[draw=black,line width=0.478bp] (175.902,-640.841) -- (218.296,-640.841);%
\path[draw=black,line width=0.478bp,line cap=round,line join=round] (216.049,-637.972) .. controls (216.503,-639.694) and (217.515,-640.507) .. (218.535,-640.841) .. controls (217.515,-641.176) and (216.503,-641.989) .. (216.049,-643.71);%
\path[draw=black,line width=0.478bp] (163.09,-631.082) -- (163.09,-599.679);%
\path[draw=black,line width=0.478bp,line cap=round,line join=round] (160.221,-601.926) .. controls (161.942,-601.472) and (162.755,-600.46) .. (163.09,-599.44) .. controls (163.425,-600.46) and (164.238,-601.472) .. (165.959,-601.926);%
\path[draw=black,line width=0.478bp] (163.09,-650.202) -- (163.09,-681.605);%
\path[draw=black,line width=0.478bp,line cap=round,line join=round] (165.959,-679.358) .. controls (164.238,-679.812) and (163.425,-680.824) .. (163.09,-681.844) .. controls (162.755,-680.824) and (161.942,-679.812) .. (160.221,-679.358);%
\path[draw=black,line width=0.478bp] (243.311,-637.75) -- (285.705,-637.75);%
\path[draw=black,line width=0.478bp,line cap=round,line join=round] (283.457,-634.88) .. controls (283.912,-636.602) and (284.924,-637.415) .. (285.944,-637.75) .. controls (284.924,-638.084) and (283.912,-638.897) .. (283.457,-640.619);%
}%
\TextAt{21}{334.273}{140.616}{f}%
\NativeDiagram{%
\path[draw=black,line width=0.478bp] (231.042,-631.082) -- (231.042,-599.679);%
\path[draw=black,line width=0.478bp,line cap=round,line join=round] (228.173,-601.926) .. controls (229.895,-601.472) and (230.708,-600.46) .. (231.042,-599.44) .. controls (231.377,-600.46) and (232.19,-601.472) .. (233.912,-601.926);%
\path[draw=black,line width=0.478bp] (231.042,-650.202) -- (231.042,-681.605);%
\path[draw=black,line width=0.478bp,line cap=round,line join=round] (233.912,-679.358) .. controls (232.19,-679.812) and (231.377,-680.824) .. (231.042,-681.844) .. controls (230.708,-680.824) and (229.895,-679.812) .. (228.173,-679.358);%
\path[draw=black,line width=0.478bp,dash pattern=on 3.347bp off 1.913bp] (286.183,-643.933) -- (243.789,-643.933);%
\path[draw=black,line width=0.478bp,line cap=round,line join=round] (246.036,-646.802) .. controls (245.582,-645.081) and (244.57,-644.268) .. (243.55,-643.933) .. controls (244.57,-643.598) and (245.582,-642.785) .. (246.036,-641.064);%
}%
\TextAt{21}{334.789}{123.82}{s}%
\NativeDiagram{%
\path[draw=black,line width=0.478bp] (314.466,-640.841) -- (358.255,-640.841);%
\path[draw=black,line width=0.478bp,line cap=round,line join=round] (356.007,-637.972) .. controls (356.462,-639.694) and (357.474,-640.507) .. (358.494,-640.841) .. controls (357.474,-641.176) and (356.462,-641.989) .. (356.007,-643.71);%
\path[draw=black,line width=2.797bp] (300.325,-631.082) -- (300.325,-597.929);%
\path[draw=black!0,line width=1.841bp] (300.325,-631.082) -- (300.325,-597.929);%
\path[draw=black,line width=2.797bp] (300.325,-650.202) -- (300.325,-682.68);%
\path[draw=black!0,line width=1.841bp] (300.325,-650.202) -- (300.325,-682.68);%
\path[draw=black,line width=0.478bp] (107.637,-692.44) -- (149.8,-692.44);%
\path[draw=black,line width=0.478bp,line cap=round,line join=round] (147.552,-689.571) .. controls (148.007,-691.292) and (149.019,-692.105) .. (150.039,-692.44) .. controls (149.019,-692.775) and (148.007,-693.588) .. (147.552,-695.309);%
\path[draw=black,line width=0.478bp] (175.902,-692.44) -- (218.064,-692.44);%
\path[draw=black,line width=0.478bp,line cap=round,line join=round] (215.817,-689.571) .. controls (216.271,-691.292) and (217.283,-692.105) .. (218.303,-692.44) .. controls (217.283,-692.775) and (216.271,-693.588) .. (215.817,-695.309);%
\path[draw=black,line width=0.478bp] (243.542,-689.348) -- (285.705,-689.348);%
\path[draw=black,line width=0.478bp,line cap=round,line join=round] (283.457,-686.479) .. controls (283.912,-688.201) and (284.924,-689.014) .. (285.944,-689.348) .. controls (284.924,-689.683) and (283.912,-690.496) .. (283.457,-692.217);%
}%
\TextAt{21}{331.954}{89.017}{f}%
\TextAt{23}{336.901}{91.83}{\Character{48}\Character{48}}%
\NativeDiagram{%
\path[draw=black,line width=0.478bp,dash pattern=on 3.347bp off 1.913bp] (286.183,-695.532) -- (244.02,-695.532);%
\path[draw=black,line width=0.478bp,line cap=round,line join=round] (246.268,-698.401) .. controls (245.814,-696.679) and (244.801,-695.866) .. (243.781,-695.532) .. controls (244.801,-695.197) and (245.814,-694.384) .. (246.268,-692.663);%
}%
\TextAt{21}{332.469}{69.519}{s}%
\TextAt{23}{336.385}{72.332}{\Character{48}\Character{48}}%
\NativeDiagram{%
\path[draw=black,line width=0.478bp] (314.466,-692.44) -- (356.629,-692.44);%
\path[draw=black,line width=0.478bp,line cap=round,line join=round] (354.381,-689.571) .. controls (354.836,-691.292) and (355.848,-692.105) .. (356.868,-692.44) .. controls (355.848,-692.775) and (354.836,-693.588) .. (354.381,-695.309);%
}%
\TextAt{8}{297.211}{43.033}{216}%
\Destination{78.23}{809.815}{nav.page.216}%
\Destination{79.36}{776.155}{nav.print.235}%
\Destination{79.36}{474.155}{nav.number.3.5.3}%
\Hotlink{146.02}{741.855}{25.078}{11.047}{nav.number.3.4.3}%
\Destination{79.226}{703.402}{errata.95}%
}

% Page 229 (body page 217)
\Page{%
\TextAt{2}{79.226}{796.943}{III.}%
\TextAt{2}{98.317}{796.943}{Riemann–Roch}%
\TextAt{2}{174.375}{796.943}{Theory}%
\Vector{q 0 G 0 g 0.299 w 7.226 18.729 m 460.769 18.729 l S Q}%
\TextAt{8}{79.226}{761.124}{where}%
\TextAt{8}{113.035}{761.124}{the}%
\TextAt{8}{133.192}{761.124}{vertical}%
\TextAt{8}{174.82}{761.124}{arrows}%
\TextAt{8}{212.287}{761.124}{are}%
\TextAt{8}{231.798}{761.124}{induced}%
\TextAt{8}{275.363}{761.124}{by}%
\TextAt{8}{291.622}{761.124}{the}%
\TextAt{8}{311.778}{761.124}{surjective}%
\TextAt{8}{364.19}{761.124}{projections}%
\TextAt{13}{245.289}{735.043}{F}%
\TextAt{21}{260.401}{741.82}{\Character{25}}%
\TextAt{23}{265.55}{744.632}{\Character{48}}%
\NativeArrowAt{257.814}{735.043}{11.955}{14.418}{\longrightarrow}%
\TextAt{13}{275.554}{735.043}{F}%
\TextAt{3}{284.758}{739.979}{\Character{48}}%
\TextAt{13}{287.553}{735.043}{\Character{59}}%
\TextAt{13}{316.688}{735.043}{F}%
\TextAt{21}{331.82}{741.82}{\Character{25}}%
\TextAt{23}{336.968}{744.632}{\Character{48}\Character{48}}%
\NativeArrowAt{329.232}{735.043}{11.955}{16.606}{\longrightarrow}%
\TextAt{13}{349.159}{735.043}{F}%
\TextAt{3}{358.362}{739.979}{\Character{48}\Character{48}}%
\TextAt{13}{363.454}{735.043}{\Character{58}}%
\TextAt{8}{79.226}{708.961}{The}%
\TextAt{8}{103.28}{708.961}{canonical}%
\TextAt{8}{154.663}{708.961}{isometries}%
\TextAt{13}{201.551}{682.88}{s}%
\TextAt{3}{207.065}{687.817}{\Character{48}}%
\TextAt{9}{213.181}{682.88}{\Character{58}}%
\TextAt{13}{219.753}{682.88}{M}%
\TextAt{25}{231.078}{681.087}{C}%
\NativeArrowAt{240.653}{682.88}{11.955}{19.26}{\longrightarrow}%
\TextAt{13}{263.235}{682.88}{s}%
\TextAt{3}{268.749}{687.817}{\Character{48}}%
\TextAt{13}{271.544}{682.88}{M}%
\TextAt{25}{282.869}{681.087}{C}%
\TextAt{13}{289.124}{682.88}{\Character{59}}%
\TextAt{13}{318.259}{682.88}{s}%
\TextAt{3}{323.792}{687.817}{\Character{48}\Character{48}}%
\TextAt{9}{332.205}{682.88}{\Character{58}}%
\TextAt{13}{338.777}{682.88}{M}%
\TextAt{25}{350.102}{681.087}{C}%
\NativeArrowAt{359.677}{682.88}{11.955}{19.26}{\longrightarrow}%
\TextAt{13}{382.259}{682.88}{s}%
\TextAt{3}{387.773}{687.817}{\Character{48}\Character{48}}%
\TextAt{13}{392.865}{682.88}{M}%
\TextAt{25}{404.19}{681.087}{C}%
\TextAt{8}{79.226}{656.799}{give}%
\TextAt{8}{103.292}{656.799}{a}%
\TextAt{8}{113.047}{656.799}{section}%
\TextAt{13}{216.194}{639.753}{s}%
\TextAt{9}{225.029}{639.753}{\Character{58}}%
\TextAt{13}{231.601}{639.753}{M}%
\TextAt{25}{242.926}{637.96}{C}%
\NativeArrowAt{252.501}{639.753}{11.955}{19.26}{\longrightarrow}%
\TextAt{13}{275.083}{639.753}{F}%
\TextAt{25}{282.661}{637.96}{C}%
\TextAt{13}{288.915}{639.753}{\Character{59}}%
\TextAt{13}{318.05}{639.753}{sx}%
\TextAt{9}{333.557}{639.753}{\Character{61}}%
\TextAt{9}{345.966}{639.753}{\Character{40}}%
\TextAt{13}{350.534}{639.753}{s}%
\TextAt{3}{356.048}{644.689}{\Character{48}}%
\TextAt{13}{358.844}{639.753}{x\Character{59}}%
\TextAt{13}{370.716}{639.753}{s}%
\TextAt{3}{376.254}{644.689}{\Character{48}\Character{48}}%
\TextAt{13}{381.346}{639.753}{x}%
\TextAt{9}{387.998}{639.753}{\Character{41}}%
\TextAt{13}{392.55}{639.753}{\Character{59}}%
\TextAt{8}{79.226}{617.483}{of}%
\NativeMathAt{93.048}{617.483}{11.955}{f}%
\TextAt{8}{106.64}{617.483}{which}%
\TextAt{8}{140.618}{617.483}{transfers}%
\TextAt{8}{189.038}{617.483}{the}%
\TextAt{8}{209.686}{617.483}{metric}%
\TextAt{8}{246.593}{617.483}{on}%
\TextAt{13}{263.344}{617.483}{M}%
\TextAt{25}{274.669}{615.69}{C}%
\TextAt{8}{285.312}{617.483}{to}%
\TextAt{8}{300.114}{617.483}{a}%
\TextAt{8}{310.361}{617.483}{metric}%
\TextAt{8}{347.268}{617.483}{on}%
\TextAt{13}{364.018}{617.483}{sM}%
\TextAt{25}{380.857}{615.69}{C}%
\TextAt{8}{387.111}{617.483}{.}%
\TextAt{13}{397.034}{617.483}{E}%
\TextAt{25}{405.7}{615.69}{C}%
\TextAt{9}{416.113}{617.483}{\Character{61}}%
\TextAt{13}{429.376}{617.483}{E}%
\TextAt{3}{438.742}{621.821}{\Character{48}}%
\TextAt{25}{438.041}{614.418}{C}%
\TextAt{16}{447.287}{617.483}{\Character{2}}%
\TextAt{13}{459.578}{617.483}{E}%
\TextAt{3}{468.944}{621.821}{\Character{48}\Character{48}}%
\TextAt{25}{468.243}{614.418}{C}%
\TextAt{8}{478.886}{617.483}{carries}%
\TextAt{8}{516.51}{617.483}{the}%
\TextAt{8}{79.226}{600.437}{sum}%
\TextAt{8}{103.997}{600.437}{of}%
\TextAt{8}{117.327}{600.437}{the}%
\TextAt{8}{137.484}{600.437}{metrics}%
\TextAt{8}{178.514}{600.437}{of}%
\TextAt{13}{191.844}{600.437}{E}%
\TextAt{3}{201.21}{604.775}{\Character{48}}%
\TextAt{25}{200.509}{597.372}{C}%
\TextAt{13}{206.764}{600.437}{\Character{59}}%
\TextAt{13}{211.988}{600.437}{E}%
\TextAt{3}{221.374}{604.775}{\Character{48}\Character{48}}%
\TextAt{25}{220.673}{597.372}{C}%
\TextAt{8}{226.927}{600.437}{,}%
\TextAt{8}{234.077}{600.437}{so}%
\TextAt{8}{248.447}{600.437}{that}%
\TextAt{13}{273.816}{600.437}{F}%
\TextAt{25}{281.393}{598.644}{C}%
\TextAt{9}{290.969}{600.437}{\Character{61}}%
\TextAt{13}{303.394}{600.437}{E}%
\TextAt{25}{312.059}{598.644}{C}%
\TextAt{16}{320.97}{600.437}{\Character{8}}%
\TextAt{13}{332.926}{600.437}{sM}%
\TextAt{25}{349.764}{598.644}{C}%
\TextAt{8}{359.916}{600.437}{also}%
\TextAt{8}{383.396}{600.437}{receives}%
\TextAt{8}{426.375}{600.437}{a}%
\TextAt{8}{436.13}{600.437}{metric,}%
\TextAt{8}{475.798}{600.437}{and}%
\TextAt{9}{232.767}{574.356}{0}%
\NativeArrowAt{241.941}{574.356}{11.955}{19.26}{\longrightarrow}%
\TextAt{13}{264.523}{574.356}{E}%
\NativeArrowAt{277.21}{574.356}{11.955}{19.26}{\longrightarrow}%
\TextAt{13}{299.792}{574.356}{F}%
\NativeArrowAt{312.317}{574.356}{11.955}{19.26}{\longrightarrow}%
\TextAt{13}{334.899}{574.356}{M}%
\NativeArrowAt{350.793}{574.356}{11.955}{19.26}{\longrightarrow}%
\TextAt{9}{373.375}{574.356}{0}%
\TextAt{8}{79.226}{548.275}{becomes}%
\TextAt{8}{125.674}{548.275}{a}%
\TextAt{8}{135.324}{548.275}{metrized}%
\TextAt{8}{183.339}{548.275}{projective}%
\TextAt{8}{237.211}{548.275}{resolution}%
\TextAt{8}{291.155}{548.275}{of}%
\TextAt{13}{304.38}{548.275}{M}%
\TextAt{8}{316.953}{548.275}{.}%
\TextAt{8}{325.367}{548.275}{It}%
\TextAt{8}{337.934}{548.275}{is}%
\TextAt{8}{349.592}{548.275}{trivial}%
\TextAt{8}{384.289}{548.275}{that}%
\TextAt{8}{409.552}{548.275}{the}%
\TextAt{8}{429.603}{548.275}{projections}%
\TextAt{13}{489.393}{548.275}{E}%
\NativeMathAt{502.08}{548.275}{11.955}{\rightarrow}%
\TextAt{13}{517.356}{548.275}{E}%
\TextAt{3}{526.723}{552.613}{\Character{48}}%
\TextAt{8}{529.518}{548.275}{,}%
\TextAt{8}{79.226}{531.229}{and}%
\TextAt{13}{101.399}{531.229}{E}%
\NativeMathAt{114.086}{531.229}{11.955}{\rightarrow}%
\TextAt{13}{129.362}{531.229}{E}%
\TextAt{3}{138.728}{535.567}{\Character{48}\Character{48}}%
\TextAt{8}{147.128}{531.229}{are}%
\TextAt{8}{166.049}{531.229}{admissible}%
\TextAt{8}{222.163}{531.229}{epimorphisms,}%
\TextAt{8}{299.208}{531.229}{and}%
\TextAt{8}{321.381}{531.229}{it}%
\TextAt{8}{332.496}{531.229}{remains}%
\TextAt{8}{375.543}{531.229}{to}%
\TextAt{8}{389.263}{531.229}{show}%
\TextAt{8}{417.677}{531.229}{this}%
\TextAt{8}{439.91}{531.229}{for}%
\TextAt{8}{457.205}{531.229}{the}%
\TextAt{8}{476.771}{531.229}{projections}%
\TextAt{13}{79.226}{514.183}{\Character{25}}%
\TextAt{3}{86.295}{518.521}{\Character{48}}%
\TextAt{9}{92.411}{514.183}{\Character{58}}%
\TextAt{13}{98.984}{514.183}{F}%
\NativeMathAt{111.508}{514.183}{11.955}{\rightarrow}%
\TextAt{13}{126.784}{514.183}{F}%
\TextAt{3}{135.988}{518.521}{\Character{48}}%
\TextAt{8}{138.783}{514.183}{,}%
\TextAt{13}{145.932}{514.183}{\Character{25}}%
\TextAt{3}{153.001}{518.521}{\Character{48}\Character{48}}%
\TextAt{9}{161.414}{514.183}{\Character{58}}%
\TextAt{13}{167.987}{514.183}{F}%
\NativeMathAt{180.511}{514.183}{11.955}{\rightarrow}%
\TextAt{13}{195.787}{514.183}{F}%
\TextAt{3}{204.991}{518.521}{\Character{48}\Character{48}}%
\TextAt{8}{210.083}{514.183}{.}%
\TextAt{8}{218.531}{514.183}{But}%
\TextAt{8}{241.772}{514.183}{we}%
\TextAt{8}{258.999}{514.183}{clearly}%
\TextAt{8}{296.395}{514.183}{have}%
\TextAt{8}{323.39}{514.183}{the}%
\TextAt{8}{343.546}{514.183}{exact}%
\TextAt{8}{374.439}{514.183}{sequence}%
\TextAt{8}{422.941}{514.183}{of}%
\TextAt{14}{436.271}{514.183}{o}%
\TextAt{8}{442.115}{514.183}{-modules}%
\TextAt{9}{201.891}{488.102}{0}%
\NativeArrowAt{211.064}{488.102}{11.955}{19.26}{\longrightarrow}%
\TextAt{13}{233.646}{488.102}{E}%
\TextAt{3}{243.013}{493.038}{\Character{48}\Character{48}}%
\TextAt{21}{255.256}{494.879}{i}%
\NativeArrowAt{251.425}{488.102}{11.955}{11.955}{\longrightarrow}%
\TextAt{13}{266.702}{488.102}{F}%
\TextAt{9}{279.226}{488.102}{\Character{61}}%
\TextAt{13}{291.652}{488.102}{F}%
\TextAt{3}{300.855}{493.038}{\Character{48}}%
\TextAt{16}{306.307}{488.102}{\Character{2}}%
\TextAt{21}{315.605}{486.309}{M}%
\TextAt{13}{327.742}{488.102}{F}%
\TextAt{3}{336.946}{493.038}{\Character{48}\Character{48}}%
\TextAt{21}{347.946}{494.879}{\Character{25}}%
\TextAt{23}{353.095}{497.692}{\Character{48}}%
\NativeArrowAt{345.358}{488.102}{11.955}{14.418}{\longrightarrow}%
\TextAt{13}{363.099}{488.102}{F}%
\TextAt{3}{372.303}{493.038}{\Character{48}}%
\NativeArrowAt{378.418}{488.102}{11.955}{19.26}{\longrightarrow}%
\TextAt{9}{401}{488.102}{0}%
\TextAt{13}{406.853}{488.102}{\Character{59}}%
\TextAt{8}{79.226}{462.021}{where}%
\TextAt{13}{113.893}{462.021}{ix}%
\TextAt{3}{124.538}{466.359}{\Character{48}\Character{48}}%
\TextAt{9}{134.413}{462.021}{\Character{61}}%
\TextAt{9}{148.281}{462.021}{\Character{40}0}%
\TextAt{13}{158.705}{462.021}{\Character{59}}%
\TextAt{13}{163.929}{462.021}{x}%
\TextAt{3}{170.601}{466.359}{\Character{48}\Character{48}}%
\TextAt{9}{175.693}{462.021}{\Character{41}}%
\TextAt{8}{180.246}{462.021}{.}%
\TextAt{8}{191.267}{462.021}{As}%
\TextAt{8}{209.411}{462.021}{the}%
\TextAt{8}{230.425}{462.021}{restriction}%
\TextAt{8}{287.281}{462.021}{of}%
\TextAt{8}{301.469}{462.021}{the}%
\TextAt{8}{322.482}{462.021}{metric}%
\TextAt{8}{359.756}{462.021}{of}%
\TextAt{13}{373.943}{462.021}{F}%
\TextAt{8}{387.902}{462.021}{to}%
\TextAt{13}{403.07}{462.021}{E}%
\TextAt{9}{417.218}{462.021}{\Character{61}}%
\TextAt{13}{431.105}{462.021}{E}%
\TextAt{3}{440.471}{466.359}{\Character{48}}%
\TextAt{16}{446.507}{462.021}{\Character{2}}%
\TextAt{13}{459.047}{462.021}{E}%
\TextAt{3}{468.413}{466.359}{\Character{48}\Character{48}}%
\TextAt{8}{478.26}{462.021}{is}%
\TextAt{8}{490.882}{462.021}{the}%
\TextAt{8}{511.896}{462.021}{sum}%
\TextAt{8}{79.226}{444.975}{of}%
\TextAt{8}{93.073}{444.975}{the}%
\TextAt{8}{113.746}{444.975}{metrics}%
\TextAt{8}{155.294}{444.975}{on}%
\TextAt{13}{172.07}{444.975}{E}%
\TextAt{3}{181.436}{449.313}{\Character{48}}%
\TextAt{8}{188.645}{444.975}{and}%
\TextAt{13}{211.925}{444.975}{E}%
\TextAt{3}{221.291}{449.313}{\Character{48}\Character{48}}%
\TextAt{8}{226.383}{444.975}{,}%
\TextAt{8}{234.178}{444.975}{we}%
\TextAt{8}{251.923}{444.975}{see}%
\TextAt{8}{271.353}{444.975}{that}%
\TextAt{13}{297.239}{444.975}{i}%
\TextAt{9}{305.434}{444.975}{\Character{58}}%
\TextAt{13}{312.887}{444.975}{E}%
\TextAt{3}{322.253}{449.313}{\Character{48}\Character{48}}%
\TextAt{25}{321.553}{441.91}{C}%
\NativeMathAt{332.009}{444.975}{11.955}{\rightarrow}%
\TextAt{13}{348.166}{444.975}{iE}%
\TextAt{3}{361.525}{449.313}{\Character{48}\Character{48}}%
\TextAt{25}{360.824}{441.91}{C}%
\TextAt{8}{371.493}{444.975}{is}%
\TextAt{8}{383.774}{444.975}{an}%
\TextAt{8}{400.55}{444.975}{isometry.}%
\TextAt{8}{453.54}{444.975}{The}%
\TextAt{8}{478.111}{444.975}{orthogonal}%
\TextAt{8}{79.226}{427.929}{complement}%
\TextAt{8}{144.585}{427.929}{of}%
\TextAt{13}{157.915}{427.929}{iE}%
\TextAt{3}{171.275}{432.267}{\Character{48}\Character{48}}%
\TextAt{25}{170.574}{424.864}{C}%
\TextAt{8}{180.726}{427.929}{in}%
\TextAt{13}{194.378}{427.929}{F}%
\TextAt{25}{201.956}{426.135}{C}%
\TextAt{8}{212.108}{427.929}{is}%
\TextAt{8}{223.872}{427.929}{the}%
\TextAt{8}{244.028}{427.929}{space}%
\TextAt{13}{169.466}{401.848}{F}%
\TextAt{3}{178.67}{406.784}{\Character{48}}%
\TextAt{25}{177.044}{398.892}{C}%
\TextAt{16}{185.955}{401.848}{\Character{2}}%
\TextAt{21}{195.253}{400.054}{M}%
\TextAt{29}{203.386}{398.682}{C}%
\TextAt{13}{211.356}{401.848}{s}%
\TextAt{3}{216.87}{406.784}{\Character{48}\Character{48}}%
\TextAt{13}{221.962}{401.848}{M}%
\TextAt{25}{233.287}{400.054}{C}%
\TextAt{9}{242.862}{401.848}{\Character{61}}%
\TextAt{16}{255.288}{401.848}{\Character{102}}%
\TextAt{9}{261.265}{401.848}{\Character{40}}%
\TextAt{13}{265.818}{401.848}{x}%
\TextAt{3}{272.47}{406.784}{\Character{48}}%
\TextAt{13}{275.265}{401.848}{\Character{59}}%
\TextAt{13}{280.489}{401.848}{s}%
\TextAt{3}{286.023}{406.784}{\Character{48}\Character{48}}%
\TextAt{13}{291.115}{401.848}{a}%
\TextAt{9}{297.26}{401.848}{\Character{41}}%
\NativeMathAt{305.133}{401.848}{11.955}{\in}%
\TextAt{13}{316.424}{401.848}{F}%
\TextAt{3}{325.628}{406.784}{\Character{48}}%
\TextAt{25}{324.002}{398.892}{C}%
\TextAt{16}{332.913}{401.848}{\Character{2}}%
\TextAt{13}{344.868}{401.848}{s}%
\TextAt{3}{350.382}{406.784}{\Character{48}\Character{48}}%
\TextAt{13}{355.474}{401.848}{M}%
\TextAt{25}{366.799}{400.054}{C}%
\TextAt{16}{376.374}{401.848}{\Character{106}}%
\TextAt{13}{383.016}{401.848}{f}%
\TextAt{3}{390.062}{406.784}{\Character{48}}%
\TextAt{9}{392.857}{401.848}{\Character{40}}%
\TextAt{13}{397.41}{401.848}{x}%
\TextAt{3}{404.062}{406.784}{\Character{48}}%
\TextAt{9}{406.857}{401.848}{\Character{41}}%
\TextAt{9}{414.712}{401.848}{\Character{61}}%
\TextAt{13}{427.155}{401.848}{a}%
\TextAt{16}{433.3}{401.848}{\Character{103}}%
\TextAt{13}{439.278}{401.848}{\Character{58}}%
\TextAt{8}{79.226}{375.767}{Indeed,}%
\TextAt{8}{120.609}{375.767}{on}%
\TextAt{8}{136.949}{375.767}{the}%
\TextAt{8}{157.187}{375.767}{one}%
\TextAt{8}{178.727}{375.767}{hand}%
\TextAt{8}{208.075}{375.767}{it}%
\TextAt{8}{219.86}{375.767}{is}%
\TextAt{8}{231.705}{375.767}{clearly}%
\TextAt{8}{269.182}{375.767}{mapped}%
\TextAt{8}{313.808}{375.767}{bijectively}%
\TextAt{8}{369.815}{375.767}{onto}%
\TextAt{13}{396.246}{375.767}{F}%
\TextAt{3}{405.449}{380.105}{\Character{48}}%
\TextAt{25}{403.824}{372.702}{C}%
\TextAt{8}{410.078}{375.767}{,}%
\TextAt{8}{417.329}{375.767}{and}%
\TextAt{8}{440.172}{375.767}{on}%
\TextAt{8}{456.512}{375.767}{the}%
\TextAt{8}{476.75}{375.767}{other}%
\TextAt{8}{507.401}{375.767}{hand}%
\TextAt{8}{79.226}{358.721}{it}%
\TextAt{8}{90.93}{358.721}{is}%
\TextAt{8}{102.694}{358.721}{orthogonal}%
\TextAt{8}{161.251}{358.721}{to}%
\TextAt{13}{175.561}{358.721}{iE}%
\TextAt{3}{188.921}{363.059}{\Character{48}\Character{48}}%
\TextAt{25}{188.22}{355.656}{C}%
\TextAt{8}{194.474}{358.721}{.}%
\TextAt{8}{202.922}{358.721}{For}%
\TextAt{8}{223.892}{358.721}{if}%
\TextAt{8}{234.615}{358.721}{we}%
\TextAt{8}{251.843}{358.721}{write}%
\TextAt{13}{281.755}{358.721}{x}%
\TextAt{3}{288.407}{363.059}{\Character{48}}%
\TextAt{9}{294.523}{358.721}{\Character{61}}%
\TextAt{13}{306.948}{358.721}{s}%
\TextAt{3}{312.462}{363.059}{\Character{48}}%
\TextAt{13}{315.257}{358.721}{a}%
\TextAt{9}{324.059}{358.721}{\Character{43}}%
\TextAt{13}{335.82}{358.721}{e}%
\TextAt{3}{341.246}{363.059}{\Character{48}}%
\TextAt{8}{344.041}{358.721}{,}%
\TextAt{8}{351.19}{358.721}{with}%
\TextAt{13}{377.85}{358.721}{e}%
\TextAt{3}{383.275}{363.059}{\Character{48}}%
\NativeMathAt{389.391}{358.721}{11.955}{\in}%
\TextAt{13}{400.682}{358.721}{E}%
\TextAt{3}{410.048}{363.059}{\Character{48}}%
\TextAt{25}{409.348}{355.656}{C}%
\TextAt{8}{415.602}{358.721}{,}%
\TextAt{8}{422.751}{358.721}{then}%
\TextAt{9}{248.976}{332.639}{\Character{40}}%
\TextAt{13}{253.528}{332.639}{x}%
\TextAt{3}{260.18}{337.576}{\Character{48}}%
\TextAt{13}{262.975}{332.639}{\Character{59}}%
\TextAt{13}{268.199}{332.639}{s}%
\TextAt{3}{273.733}{337.576}{\Character{48}\Character{48}}%
\TextAt{13}{278.825}{332.639}{a}%
\TextAt{9}{284.97}{332.639}{\Character{41}}%
\TextAt{9}{292.825}{332.639}{\Character{61}}%
\TextAt{13}{305.269}{332.639}{sa}%
\TextAt{9}{319.585}{332.639}{\Character{43}}%
\TextAt{9}{331.337}{332.639}{\Character{40}}%
\TextAt{13}{335.898}{332.639}{e}%
\TextAt{3}{341.324}{337.576}{\Character{48}}%
\TextAt{13}{344.119}{332.639}{\Character{59}}%
\TextAt{9}{349.363}{332.639}{0\Character{41}}%
\TextAt{13}{359.768}{332.639}{\Character{59}}%
\TextAt{8}{79.226}{306.558}{where}%
\TextAt{9}{113.035}{306.558}{\Character{40}}%
\TextAt{13}{117.588}{306.558}{e}%
\TextAt{3}{123.013}{310.897}{\Character{48}}%
\TextAt{13}{125.808}{306.558}{\Character{59}}%
\TextAt{9}{131.052}{306.558}{0\Character{41}}%
\NativeMathAt{144.779}{306.558}{11.955}{\in}%
\TextAt{13}{156.069}{306.558}{E}%
\TextAt{25}{164.735}{304.765}{C}%
\TextAt{8}{174.887}{306.558}{and}%
\TextAt{8}{197.649}{306.558}{we}%
\TextAt{8}{214.877}{306.558}{find}%
\TextAt{8}{238.285}{306.558}{that,}%
\TextAt{8}{266.905}{306.558}{for}%
\TextAt{8}{284.79}{306.558}{all}%
\TextAt{13}{301.049}{306.558}{x}%
\TextAt{3}{307.702}{310.897}{\Character{48}\Character{48}}%
\NativeMathAt{316.114}{306.558}{11.955}{\in}%
\TextAt{13}{327.405}{306.558}{E}%
\TextAt{3}{336.772}{310.897}{\Character{48}\Character{48}}%
\TextAt{25}{336.071}{303.493}{C}%
\TextAt{8}{342.325}{306.558}{,}%
\TextAt{16}{165.939}{280.477}{\Character{104}}%
\TextAt{13}{170.588}{280.477}{ix}%
\TextAt{3}{181.233}{285.413}{\Character{48}\Character{48}}%
\TextAt{13}{186.325}{280.477}{\Character{59}}%
\TextAt{9}{191.569}{280.477}{\Character{40}}%
\TextAt{13}{196.122}{280.477}{x}%
\TextAt{3}{202.774}{285.413}{\Character{48}}%
\TextAt{13}{205.569}{280.477}{\Character{59}}%
\TextAt{13}{210.793}{280.477}{s}%
\TextAt{3}{216.327}{285.413}{\Character{48}\Character{48}}%
\TextAt{13}{221.419}{280.477}{a}%
\TextAt{9}{227.564}{280.477}{\Character{41}}%
\TextAt{16}{232.116}{280.477}{\Character{105}}%
\TextAt{21}{236.766}{278.684}{F}%
\TextAt{9}{247.12}{280.477}{\Character{61}}%
\TextAt{16}{259.545}{280.477}{\Character{104}}%
\TextAt{9}{264.195}{280.477}{\Character{40}0}%
\TextAt{13}{274.6}{280.477}{\Character{59}}%
\TextAt{13}{279.824}{280.477}{x}%
\TextAt{3}{286.496}{285.413}{\Character{48}\Character{48}}%
\TextAt{9}{291.588}{280.477}{\Character{41}}%
\TextAt{13}{296.141}{280.477}{\Character{59}}%
\TextAt{13}{301.365}{280.477}{sa}%
\TextAt{16}{313.044}{280.477}{\Character{105}}%
\TextAt{21}{317.693}{278.684}{F}%
\TextAt{9}{327.383}{280.477}{\Character{43}}%
\TextAt{16}{339.144}{280.477}{\Character{104}}%
\TextAt{9}{343.794}{280.477}{\Character{40}0}%
\TextAt{13}{354.199}{280.477}{\Character{59}}%
\TextAt{13}{359.423}{280.477}{x}%
\TextAt{3}{366.095}{285.413}{\Character{48}\Character{48}}%
\TextAt{9}{371.187}{280.477}{\Character{41}}%
\TextAt{13}{375.739}{280.477}{\Character{59}}%
\TextAt{9}{380.984}{280.477}{\Character{40}}%
\TextAt{13}{385.536}{280.477}{e}%
\TextAt{3}{390.961}{285.413}{\Character{48}}%
\TextAt{13}{393.756}{280.477}{\Character{59}}%
\TextAt{9}{399.001}{280.477}{0\Character{41}}%
\TextAt{16}{409.406}{280.477}{\Character{105}}%
\TextAt{21}{414.055}{278.684}{E}%
\TextAt{9}{424.527}{280.477}{\Character{61}}%
\TextAt{9}{436.937}{280.477}{0}%
\TextAt{13}{442.806}{280.477}{\Character{58}}%
\TextAt{8}{79.226}{254.396}{Finally,}%
\TextAt{8}{123.436}{254.396}{the}%
\TextAt{8}{145.275}{254.396}{projection}%
\TextAt{13}{202.238}{254.396}{F}%
\TextAt{3}{211.442}{258.735}{\Character{48}}%
\TextAt{25}{209.816}{251.331}{C}%
\TextAt{16}{219.874}{254.396}{\Character{2}}%
\TextAt{21}{229.172}{252.603}{M}%
\TextAt{29}{237.305}{251.23}{C}%
\TextAt{13}{246.422}{254.396}{s}%
\TextAt{3}{251.936}{258.735}{\Character{48}\Character{48}}%
\TextAt{13}{257.028}{254.396}{M}%
\TextAt{25}{268.353}{252.603}{C}%
\NativeMathAt{280.795}{254.396}{11.955}{\rightarrow}%
\TextAt{13}{298.939}{254.396}{F}%
\TextAt{3}{308.142}{258.735}{\Character{48}}%
\TextAt{25}{306.517}{251.331}{C}%
\TextAt{8}{318.351}{254.396}{is}%
\TextAt{8}{331.797}{254.396}{an}%
\TextAt{8}{349.739}{254.396}{isometry,}%
\TextAt{8}{401.982}{254.396}{for}%
\TextAt{8}{421.55}{254.396}{if}%
\TextAt{9}{433.956}{254.396}{\Character{40}}%
\TextAt{13}{438.509}{254.396}{x}%
\TextAt{3}{445.161}{258.735}{\Character{48}}%
\TextAt{13}{447.956}{254.396}{\Character{59}}%
\TextAt{13}{453.18}{254.396}{s}%
\TextAt{3}{458.714}{258.735}{\Character{48}\Character{48}}%
\TextAt{13}{463.806}{254.396}{a}%
\TextAt{9}{469.951}{254.396}{\Character{41}}%
\TextAt{13}{474.503}{254.396}{\Character{59}}%
\TextAt{9}{479.747}{254.396}{\Character{40}}%
\TextAt{13}{484.3}{254.396}{y}%
\TextAt{3}{490.436}{258.735}{\Character{48}}%
\TextAt{13}{493.232}{254.396}{\Character{59}}%
\TextAt{13}{498.456}{254.396}{s}%
\TextAt{3}{503.99}{258.735}{\Character{48}\Character{48}}%
\TextAt{13}{509.082}{254.396}{b}%
\TextAt{9}{514.059}{254.396}{\Character{41}}%
\NativeMathAt{524.799}{254.396}{11.955}{\in}%
\TextAt{13}{79.226}{237.35}{F}%
\TextAt{3}{88.43}{241.689}{\Character{48}}%
\TextAt{25}{86.804}{234.285}{C}%
\TextAt{16}{95.715}{237.35}{\Character{2}}%
\TextAt{21}{105.013}{235.557}{M}%
\TextAt{29}{113.146}{234.184}{C}%
\TextAt{13}{121.116}{237.35}{s}%
\TextAt{3}{126.63}{241.689}{\Character{48}\Character{48}}%
\TextAt{13}{131.722}{237.35}{M}%
\TextAt{25}{143.047}{235.557}{C}%
\TextAt{8}{153.199}{237.35}{and}%
\TextAt{13}{175.961}{237.35}{x}%
\TextAt{3}{182.614}{241.689}{\Character{48}}%
\TextAt{9}{188.729}{237.35}{\Character{61}}%
\TextAt{13}{201.155}{237.35}{s}%
\TextAt{3}{206.669}{241.689}{\Character{48}}%
\TextAt{13}{209.464}{237.35}{a}%
\TextAt{9}{218.266}{237.35}{\Character{43}}%
\TextAt{13}{230.027}{237.35}{e}%
\TextAt{3}{235.452}{241.689}{\Character{48}}%
\TextAt{8}{238.247}{237.35}{,}%
\TextAt{13}{245.397}{237.35}{y}%
\TextAt{3}{251.533}{241.689}{\Character{48}}%
\TextAt{9}{257.649}{237.35}{\Character{61}}%
\TextAt{13}{270.075}{237.35}{s}%
\TextAt{3}{275.589}{241.689}{\Character{48}}%
\TextAt{13}{278.384}{237.35}{b}%
\TextAt{9}{286.018}{237.35}{\Character{43}}%
\TextAt{13}{297.779}{237.35}{d}%
\TextAt{3}{303.862}{241.689}{\Character{48}}%
\TextAt{8}{306.657}{237.35}{,}%
\TextAt{8}{313.806}{237.35}{with}%
\TextAt{13}{340.466}{237.35}{e}%
\TextAt{3}{345.891}{241.689}{\Character{48}}%
\TextAt{13}{348.686}{237.35}{\Character{59}}%
\TextAt{13}{353.91}{237.35}{d}%
\TextAt{3}{360.013}{241.689}{\Character{48}}%
\NativeMathAt{366.129}{237.35}{11.955}{\in}%
\TextAt{13}{377.42}{237.35}{E}%
\TextAt{3}{386.786}{241.689}{\Character{48}}%
\TextAt{25}{386.085}{234.285}{C}%
\TextAt{8}{392.34}{237.35}{,}%
\TextAt{8}{399.489}{237.35}{then}%
\TextAt{8}{426.149}{237.35}{we}%
\TextAt{8}{443.376}{237.35}{get}%
\TextAt{9}{181.725}{211.269}{\Character{40}}%
\TextAt{13}{186.277}{211.269}{x}%
\TextAt{3}{192.929}{216.205}{\Character{48}}%
\TextAt{13}{195.725}{211.269}{\Character{59}}%
\TextAt{13}{200.949}{211.269}{s}%
\TextAt{3}{206.483}{216.205}{\Character{48}\Character{48}}%
\TextAt{13}{211.575}{211.269}{a}%
\TextAt{9}{217.72}{211.269}{\Character{41}}%
\TextAt{9}{225.575}{211.269}{\Character{61}}%
\TextAt{13}{238.018}{211.269}{sa}%
\TextAt{9}{252.334}{211.269}{\Character{43}}%
\TextAt{9}{264.086}{211.269}{\Character{40}}%
\TextAt{13}{268.648}{211.269}{e}%
\TextAt{3}{274.073}{216.205}{\Character{48}}%
\TextAt{13}{276.868}{211.269}{\Character{59}}%
\TextAt{9}{282.112}{211.269}{0\Character{41}}%
\TextAt{13}{292.517}{211.269}{\Character{59}}%
\TextAt{9}{321.672}{211.269}{\Character{40}}%
\TextAt{13}{326.224}{211.269}{y}%
\TextAt{3}{332.361}{216.205}{\Character{48}}%
\TextAt{13}{335.156}{211.269}{\Character{59}}%
\TextAt{13}{340.38}{211.269}{s}%
\TextAt{3}{345.914}{216.205}{\Character{48}\Character{48}}%
\TextAt{13}{351.006}{211.269}{b}%
\TextAt{9}{355.983}{211.269}{\Character{41}}%
\TextAt{9}{363.838}{211.269}{\Character{61}}%
\TextAt{13}{376.282}{211.269}{sb}%
\TextAt{9}{389.43}{211.269}{\Character{43}}%
\TextAt{9}{401.182}{211.269}{\Character{40}}%
\TextAt{13}{405.743}{211.269}{d}%
\TextAt{3}{411.826}{216.205}{\Character{48}}%
\TextAt{13}{414.621}{211.269}{\Character{59}}%
\TextAt{9}{419.865}{211.269}{0\Character{41}}%
\TextAt{8}{79.226}{185.188}{and}%
\TextAt{16}{140.291}{159.107}{\Character{104}}%
\TextAt{9}{144.941}{159.107}{\Character{40}}%
\TextAt{13}{149.493}{159.107}{x}%
\TextAt{3}{156.145}{164.043}{\Character{48}}%
\TextAt{13}{158.94}{159.107}{\Character{59}}%
\TextAt{13}{164.164}{159.107}{s}%
\TextAt{3}{169.698}{164.043}{\Character{48}\Character{48}}%
\TextAt{13}{174.79}{159.107}{a}%
\TextAt{9}{180.935}{159.107}{\Character{41}}%
\TextAt{13}{185.487}{159.107}{\Character{59}}%
\TextAt{9}{190.732}{159.107}{\Character{40}}%
\TextAt{13}{195.284}{159.107}{y}%
\TextAt{3}{201.421}{164.043}{\Character{48}}%
\TextAt{13}{204.216}{159.107}{\Character{59}}%
\TextAt{13}{209.44}{159.107}{s}%
\TextAt{3}{214.974}{164.043}{\Character{48}\Character{48}}%
\TextAt{13}{220.066}{159.107}{b}%
\TextAt{9}{225.043}{159.107}{\Character{41}}%
\TextAt{16}{229.595}{159.107}{\Character{105}}%
\TextAt{21}{234.245}{157.313}{F}%
\TextAt{9}{244.599}{159.107}{\Character{61}}%
\TextAt{16}{257.024}{159.107}{\Character{104}}%
\TextAt{13}{261.673}{159.107}{sa\Character{59}}%
\TextAt{13}{278.554}{159.107}{sb}%
\TextAt{16}{289.068}{159.107}{\Character{105}}%
\TextAt{21}{293.717}{157.313}{F}%
\TextAt{9}{303.407}{159.107}{\Character{43}}%
\TextAt{16}{315.168}{159.107}{\Character{104}}%
\TextAt{13}{319.818}{159.107}{sa\Character{59}}%
\TextAt{9}{336.721}{159.107}{\Character{40}}%
\TextAt{13}{341.273}{159.107}{d}%
\TextAt{3}{347.356}{164.043}{\Character{48}}%
\TextAt{13}{350.151}{159.107}{\Character{59}}%
\TextAt{9}{355.395}{159.107}{0\Character{41}}%
\TextAt{16}{365.8}{159.107}{\Character{105}}%
\TextAt{21}{370.45}{157.313}{F}%
\TextAt{9}{380.14}{159.107}{\Character{43}}%
\TextAt{16}{391.901}{159.107}{\Character{104}}%
\TextAt{9}{396.55}{159.107}{\Character{40}}%
\TextAt{13}{401.103}{159.107}{e}%
\TextAt{3}{406.528}{164.043}{\Character{48}}%
\TextAt{13}{409.323}{159.107}{\Character{59}}%
\TextAt{9}{414.567}{159.107}{0\Character{41}}%
\TextAt{13}{424.973}{159.107}{\Character{59}}%
\TextAt{13}{430.197}{159.107}{sb}%
\TextAt{16}{440.708}{159.107}{\Character{105}}%
\TextAt{21}{445.357}{157.313}{F}%
\TextAt{9}{255.89}{139.072}{\Character{43}}%
\TextAt{16}{267.651}{139.072}{\Character{104}}%
\TextAt{9}{272.3}{139.072}{\Character{40}}%
\TextAt{13}{276.853}{139.072}{e}%
\TextAt{3}{282.278}{144.008}{\Character{48}}%
\TextAt{13}{285.073}{139.072}{\Character{59}}%
\TextAt{9}{290.317}{139.072}{0\Character{41}}%
\TextAt{13}{300.723}{139.072}{\Character{59}}%
\TextAt{9}{305.967}{139.072}{\Character{40}}%
\TextAt{13}{310.519}{139.072}{d}%
\TextAt{3}{316.602}{144.008}{\Character{48}}%
\TextAt{13}{319.397}{139.072}{\Character{59}}%
\TextAt{9}{324.641}{139.072}{0\Character{41}}%
\TextAt{16}{335.046}{139.072}{\Character{105}}%
\TextAt{21}{339.696}{137.279}{E}%
\TextAt{9}{244.599}{119.037}{\Character{61}}%
\TextAt{16}{257.024}{119.037}{\Character{104}}%
\TextAt{13}{261.674}{119.037}{a\Character{59}}%
\TextAt{13}{273.043}{119.037}{b}%
\TextAt{16}{278.04}{119.037}{\Character{105}}%
\TextAt{21}{282.689}{117.244}{M}%
\TextAt{9}{294.826}{119.037}{\Character{43}}%
\TextAt{16}{306.587}{119.037}{\Character{104}}%
\TextAt{13}{311.236}{119.037}{e}%
\TextAt{3}{316.662}{123.973}{\Character{48}}%
\TextAt{13}{319.457}{119.037}{\Character{59}}%
\TextAt{13}{324.681}{119.037}{d}%
\TextAt{3}{330.783}{123.973}{\Character{48}}%
\TextAt{16}{333.579}{119.037}{\Character{105}}%
\TextAt{21}{338.228}{117.244}{E}%
\TextAt{23}{344.881}{119.513}{\Character{48}}%
\TextAt{9}{351.384}{119.037}{\Character{61}}%
\TextAt{16}{363.81}{119.037}{\Character{104}}%
\TextAt{13}{368.459}{119.037}{s}%
\TextAt{3}{373.973}{123.973}{\Character{48}}%
\TextAt{13}{376.768}{119.037}{a\Character{59}}%
\TextAt{13}{388.137}{119.037}{s}%
\TextAt{3}{393.671}{123.973}{\Character{48}}%
\TextAt{13}{396.466}{119.037}{b}%
\TextAt{16}{401.443}{119.037}{\Character{105}}%
\TextAt{21}{406.092}{117.244}{F}%
\TextAt{23}{412.628}{119.513}{\Character{48}}%
\TextAt{9}{418.467}{119.037}{\Character{43}}%
\TextAt{16}{430.228}{119.037}{\Character{104}}%
\TextAt{13}{434.877}{119.037}{e}%
\TextAt{3}{440.303}{123.973}{\Character{48}}%
\TextAt{13}{443.098}{119.037}{\Character{59}}%
\TextAt{13}{448.322}{119.037}{d}%
\TextAt{3}{454.425}{123.973}{\Character{48}}%
\TextAt{16}{457.22}{119.037}{\Character{105}}%
\TextAt{21}{461.869}{117.244}{E}%
\TextAt{23}{468.522}{119.513}{\Character{48}}%
\TextAt{9}{244.599}{99.002}{\Character{61}}%
\TextAt{16}{257.024}{99.002}{\Character{104}}%
\TextAt{13}{261.674}{99.002}{s}%
\TextAt{3}{267.188}{103.939}{\Character{48}}%
\TextAt{13}{269.983}{99.002}{a}%
\TextAt{9}{278.784}{99.002}{\Character{43}}%
\TextAt{13}{290.546}{99.002}{e}%
\TextAt{3}{295.971}{103.939}{\Character{48}}%
\TextAt{13}{298.766}{99.002}{\Character{59}}%
\TextAt{13}{303.99}{99.002}{s}%
\TextAt{3}{309.524}{103.939}{\Character{48}}%
\TextAt{13}{312.319}{99.002}{b}%
\TextAt{9}{319.953}{99.002}{\Character{43}}%
\TextAt{13}{331.714}{99.002}{d}%
\TextAt{3}{337.797}{103.939}{\Character{48}}%
\TextAt{16}{340.592}{99.002}{\Character{105}}%
\TextAt{21}{345.241}{97.209}{F}%
\TextAt{23}{351.777}{99.478}{\Character{48}}%
\TextAt{9}{358.28}{99.002}{\Character{61}}%
\TextAt{16}{370.706}{99.002}{\Character{104}}%
\TextAt{13}{375.355}{99.002}{x}%
\TextAt{3}{382.007}{103.939}{\Character{48}}%
\TextAt{13}{384.802}{99.002}{\Character{59}}%
\TextAt{13}{390.026}{99.002}{y}%
\TextAt{3}{396.183}{103.939}{\Character{48}}%
\TextAt{16}{398.978}{99.002}{\Character{105}}%
\TextAt{21}{403.627}{97.209}{F}%
\TextAt{23}{410.162}{99.478}{\Character{48}}%
\TextAt{13}{413.345}{99.002}{\Character{58}}%
\Vector{q 0 G 0 g 0.399 w 452.334 -701.079 m 452.334 -693.009 l S Q}%
\Vector{q 0 G 0 g 0.399 w 452.534 -693.208 m 459.707 -693.208 l S Q}%
\Vector{q 0 G 0 g 0.399 w 452.534 -700.879 m 459.707 -700.879 l S Q}%
\Vector{q 0 G 0 g 0.399 w 459.906 -701.079 m 459.906 -693.009 l S Q}%
\TextAt{8}{297.211}{43.033}{217}%
\Destination{78.23}{809.815}{nav.page.217}%
\Destination{79.36}{776.155}{nav.print.236}%
\Destination{79.226}{634.483}{errata.96}%
}

% Page 230 (body page 218)
\Page{%
\TextAt{2}{79.226}{796.943}{III.}%
\TextAt{2}{98.317}{796.943}{Riemann–Roch}%
\TextAt{2}{174.375}{796.943}{Theory}%
\Vector{q 0 G 0 g 0.299 w 7.226 18.729 m 460.769 18.729 l S Q}%
\TextAt{12}{79.226}{761.124}{Theorem}%
\TextAt{12}{136.24}{761.124}{(5.4).}%
\TextAt{8}{171.556}{761.124}{The}%
\TextAt{8}{195.61}{761.124}{Poincaré}%
\TextAt{8}{243.574}{761.124}{homomorphism}%
\TextAt{13}{262.977}{734.528}{K}%
\TextAt{17}{272.936}{732.734}{0}%
\TextAt{9}{277.668}{734.528}{\Character{40}}%
\Vector{q 0 G 0 g 0.478 w 210.221 -32.116 m 216.065 -32.116 l S Q}%
\TextAt{14}{282.221}{734.528}{o}%
\TextAt{9}{288.065}{734.528}{\Character{41}}%
\NativeArrowAt{295.938}{734.528}{11.955}{19.26}{\longrightarrow}%
\TextAt{13}{318.52}{734.528}{K}%
\TextAt{17}{329.338}{739.464}{0}%
\TextAt{9}{334.07}{734.528}{\Character{40}}%
\Vector{q 0 G 0 g 0.478 w 266.622 -32.116 m 272.467 -32.116 l S Q}%
\TextAt{14}{338.622}{734.528}{o}%
\TextAt{9}{344.467}{734.528}{\Character{41}}%
\TextAt{8}{79.226}{707.932}{is}%
\TextAt{8}{90.99}{707.932}{an}%
\TextAt{8}{107.249}{707.932}{isomorphism.}%
\TextAt{19}{79.226}{674.328}{Proof.}%
\TextAt{8}{116}{674.328}{We}%
\TextAt{8}{136.145}{674.328}{define}%
\TextAt{8}{169.954}{674.328}{a}%
\TextAt{8}{179.709}{674.328}{mapping}%
\TextAt{13}{255.303}{657.282}{\Character{25}}%
\TextAt{9}{265.693}{657.282}{\Character{58}}%
\TextAt{13}{272.265}{657.282}{F}%
\TextAt{17}{281.469}{662.218}{0}%
\TextAt{9}{286.201}{657.282}{\Character{40}}%
\Vector{q 0 G 0 g 0.478 w 218.753 -109.362 m 224.598 -109.362 l S Q}%
\TextAt{14}{290.753}{657.282}{o}%
\TextAt{9}{296.598}{657.282}{\Character{41}}%
\NativeArrowAt{304.471}{657.282}{11.955}{19.26}{\longrightarrow}%
\TextAt{13}{327.053}{657.282}{K}%
\TextAt{17}{337.012}{655.488}{0}%
\TextAt{9}{341.744}{657.282}{\Character{40}}%
\Vector{q 0 G 0 g 0.478 w 274.297 -109.362 m 280.141 -109.362 l S Q}%
\TextAt{14}{346.297}{657.282}{o}%
\TextAt{9}{352.141}{657.282}{\Character{41}}%
\TextAt{8}{79.226}{634.791}{by}%
\TextAt{8}{95.485}{634.791}{choosing,}%
\TextAt{8}{146.283}{634.791}{for}%
\TextAt{8}{164.168}{634.791}{every}%
\TextAt{8}{195.06}{634.791}{coherent}%
\TextAt{8}{242.211}{634.791}{metrized}%
\TextAt{14}{290.331}{634.791}{o}%
\TextAt{8}{296.175}{634.791}{-module}%
\TextAt{13}{341.365}{634.791}{M}%
\TextAt{8}{353.939}{634.791}{,}%
\TextAt{8}{361.088}{634.791}{a}%
\TextAt{8}{370.844}{634.791}{metrized}%
\TextAt{8}{418.963}{634.791}{projective}%
\TextAt{8}{472.941}{634.791}{resolution}%
\TextAt{9}{232.767}{608.195}{0}%
\NativeArrowAt{241.941}{608.195}{11.955}{19.26}{\longrightarrow}%
\TextAt{13}{264.523}{608.195}{E}%
\NativeArrowAt{277.21}{608.195}{11.955}{19.26}{\longrightarrow}%
\TextAt{13}{299.792}{608.195}{F}%
\NativeArrowAt{312.317}{608.195}{11.955}{19.26}{\longrightarrow}%
\TextAt{13}{334.899}{608.195}{M}%
\NativeArrowAt{350.793}{608.195}{11.955}{19.26}{\longrightarrow}%
\TextAt{9}{373.375}{608.195}{0}%
\TextAt{8}{79.226}{581.599}{and}%
\TextAt{8}{102.147}{581.599}{associating}%
\TextAt{8}{161.949}{581.599}{to}%
\TextAt{8}{176.417}{581.599}{the}%
\TextAt{8}{196.731}{581.599}{class}%
\TextAt{16}{224.326}{581.599}{\Character{102}}%
\TextAt{13}{230.303}{581.599}{M}%
\TextAt{16}{242.877}{581.599}{\Character{103}}%
\TextAt{8}{252.91}{581.599}{in}%
\TextAt{13}{266.72}{581.599}{F}%
\TextAt{17}{275.924}{585.937}{0}%
\TextAt{9}{280.656}{581.599}{\Character{40}}%
\Vector{q 0 G 0 g 0.478 w 213.208 -185.045 m 219.053 -185.045 l S Q}%
\TextAt{14}{285.208}{581.599}{o}%
\TextAt{9}{291.053}{581.599}{\Character{41}}%
\TextAt{8}{299.66}{581.599}{the}%
\TextAt{8}{319.974}{581.599}{difference}%
\TextAt{9}{372.472}{581.599}{\Character{91}}%
\TextAt{13}{375.723}{581.599}{F}%
\TextAt{9}{384.927}{581.599}{\Character{93}}%
\TextAt{16}{390.943}{581.599}{\Character{0}}%
\TextAt{9}{403.006}{581.599}{\Character{91}}%
\TextAt{13}{406.257}{581.599}{E}%
\TextAt{9}{415.623}{581.599}{\Character{93}}%
\TextAt{8}{422.93}{581.599}{of}%
\TextAt{8}{436.418}{581.599}{the}%
\TextAt{8}{456.732}{581.599}{classes}%
\TextAt{9}{494.142}{581.599}{\Character{91}}%
\TextAt{13}{497.394}{581.599}{F}%
\TextAt{9}{506.597}{581.599}{\Character{93}}%
\TextAt{8}{513.904}{581.599}{and}%
\TextAt{9}{79.226}{564.553}{\Character{91}}%
\TextAt{13}{82.478}{564.553}{E}%
\TextAt{9}{91.844}{564.553}{\Character{93}}%
\TextAt{8}{98.738}{564.553}{in}%
\TextAt{13}{112.136}{564.553}{K}%
\TextAt{17}{122.096}{562.76}{0}%
\TextAt{9}{126.828}{564.553}{\Character{40}}%
\Vector{q 0 G 0 g 0.478 w 59.38 -202.091 m 65.224 -202.091 l S Q}%
\TextAt{14}{131.38}{564.553}{o}%
\TextAt{9}{137.224}{564.553}{\Character{41}}%
\TextAt{8}{141.777}{564.553}{.}%
\TextAt{8}{150.14}{564.553}{To}%
\TextAt{8}{167.113}{564.553}{see}%
\TextAt{8}{185.771}{564.553}{that}%
\TextAt{8}{210.885}{564.553}{this}%
\TextAt{8}{233.453}{564.553}{mapping}%
\TextAt{8}{281.33}{564.553}{is}%
\TextAt{8}{292.839}{564.553}{well-defined}%
\TextAt{8}{356.628}{564.553}{let}%
\TextAt{8}{373.278}{564.553}{us}%
\TextAt{8}{388.039}{564.553}{first}%
\TextAt{8}{411.91}{564.553}{consider}%
\TextAt{8}{457.24}{564.553}{a}%
\TextAt{8}{466.741}{564.553}{commutative}%
\TextAt{8}{79.226}{547.507}{diagram}%
\TextAt{9}{175.088}{531.511}{0}%
\TextAt{13}{235.683}{531.511}{E}%
\TextAt{13}{301.188}{531.511}{F}%
\TextAt{13}{365.133}{531.511}{M}%
\TextAt{9}{431.051}{531.511}{0}%
\TextAt{9}{175.088}{481.706}{0}%
\TextAt{13}{234.285}{481.706}{E}%
\TextAt{3}{243.651}{486.044}{\Character{48}}%
\TextAt{13}{299.791}{481.706}{F}%
\TextAt{3}{308.995}{486.044}{\Character{48}}%
\TextAt{13}{365.133}{481.706}{M}%
\TextAt{9}{431.051}{481.706}{0}%
\NativeDiagram{%
\path[draw=black,line width=0.478bp] (114.294,-239.5) -- (157.854,-239.5);%
\path[draw=black,line width=0.478bp,line cap=round,line join=round] (155.606,-236.631) .. controls (156.061,-238.353) and (157.073,-239.165) .. (158.093,-239.5) .. controls (157.073,-239.835) and (156.061,-240.648) .. (155.606,-242.369);%
\path[draw=black,line width=0.478bp] (178.402,-239.5) -- (223.36,-239.5);%
\path[draw=black,line width=0.478bp,line cap=round,line join=round] (221.112,-236.631) .. controls (221.567,-238.353) and (222.579,-239.165) .. (223.599,-239.5) .. controls (222.579,-239.835) and (221.567,-240.648) .. (221.112,-242.369);%
\path[draw=black,line width=0.478bp] (168.367,-247.068) -- (168.367,-278.471);%
\path[draw=black,line width=0.478bp,line cap=round,line join=round] (171.236,-276.223) .. controls (169.515,-276.677) and (168.702,-277.69) .. (168.367,-278.71) .. controls (168.032,-277.69) and (167.219,-276.677) .. (165.498,-276.223);%
}%
\TextAt{21}{243.182}{509.276}{\Character{11}}%
\NativeDiagram{%
\path[draw=black,line width=0.478bp] (243.746,-239.5) -- (287.306,-239.5);%
\path[draw=black,line width=0.478bp,line cap=round,line join=round] (285.058,-236.631) .. controls (285.512,-238.353) and (286.525,-239.165) .. (287.545,-239.5) .. controls (286.525,-239.835) and (285.512,-240.648) .. (285.058,-242.369);%
\path[draw=black,line width=0.478bp] (233.792,-247.068) -- (233.792,-278.471);%
\path[draw=black,line width=0.478bp,line cap=round,line join=round] (236.661,-276.223) .. controls (234.939,-276.677) and (234.127,-277.69) .. (233.792,-278.71) .. controls (233.457,-277.69) and (232.644,-276.677) .. (230.923,-276.223);%
}%
\TextAt{21}{308.607}{508.999}{\Character{12}}%
\NativeDiagram{%
\path[draw=black,line width=0.478bp] (311.062,-239.5) -- (353.224,-239.5);%
\path[draw=black,line width=0.478bp,line cap=round,line join=round] (350.977,-236.631) .. controls (351.431,-238.353) and (352.443,-239.165) .. (353.463,-239.5) .. controls (352.443,-239.835) and (351.431,-240.648) .. (350.977,-242.369);%
\path[draw=black,line width=2.797bp] (299.423,-247.068) -- (299.423,-279.546);%
\path[draw=black!0,line width=1.841bp] (299.423,-247.068) -- (299.423,-279.546);%
\path[draw=black,line width=0.478bp] (114.294,-289.306) -- (156.456,-289.306);%
\path[draw=black,line width=0.478bp,line cap=round,line join=round] (154.209,-286.436) .. controls (154.663,-288.158) and (155.675,-288.971) .. (156.695,-289.306) .. controls (155.675,-289.64) and (154.663,-290.453) .. (154.209,-292.175);%
\path[draw=black,line width=0.478bp] (179.8,-289.306) -- (221.962,-289.306);%
\path[draw=black,line width=0.478bp,line cap=round,line join=round] (219.715,-286.436) .. controls (220.169,-288.158) and (221.181,-288.971) .. (222.201,-289.306) .. controls (221.181,-289.64) and (220.169,-290.453) .. (219.715,-292.175);%
\path[draw=black,line width=0.478bp] (245.143,-289.306) -- (287.306,-289.306);%
\path[draw=black,line width=0.478bp,line cap=round,line join=round] (285.058,-286.436) .. controls (285.512,-288.158) and (286.525,-288.971) .. (287.545,-289.306) .. controls (286.525,-289.64) and (285.512,-290.453) .. (285.058,-292.175);%
\path[draw=black,line width=0.478bp] (311.062,-289.306) -- (353.224,-289.306);%
\path[draw=black,line width=0.478bp,line cap=round,line join=round] (350.977,-286.436) .. controls (351.431,-288.158) and (352.443,-288.971) .. (353.463,-289.306) .. controls (352.443,-289.64) and (351.431,-290.453) .. (350.977,-292.175);%
}%
\TextAt{8}{79.226}{462.468}{of}%
\TextAt{8}{93.751}{462.468}{two}%
\TextAt{8}{117.063}{462.468}{metrized}%
\TextAt{8}{166.377}{462.468}{projective}%
\TextAt{8}{221.55}{462.468}{resolutions}%
\TextAt{8}{281.409}{462.468}{of}%
\TextAt{13}{295.934}{462.468}{M}%
\TextAt{8}{308.507}{462.468}{,}%
\TextAt{8}{317.15}{462.468}{with}%
\TextAt{8}{345.005}{462.468}{the}%
\TextAt{8}{366.356}{462.468}{top}%
\TextAt{8}{388.365}{462.468}{one}%
\TextAt{8}{411.019}{462.468}{dominating}%
\TextAt{8}{474.01}{462.468}{the}%
\TextAt{8}{495.362}{462.468}{bottom}%
\TextAt{8}{79.226}{445.422}{one.}%
\TextAt{8}{106.463}{445.422}{Then}%
\TextAt{13}{137.43}{445.422}{E}%
\NativeMathAt{150.814}{445.422}{11.955}{\rightarrow}%
\TextAt{13}{166.787}{445.422}{F}%
\TextAt{8}{180.296}{445.422}{induces}%
\TextAt{8}{222.381}{445.422}{an}%
\TextAt{8}{239.049}{445.422}{isometry}%
\TextAt{8}{287.326}{445.422}{ker}%
\TextAt{9}{303.263}{445.422}{\Character{40}}%
\TextAt{13}{307.815}{445.422}{\Character{11}}%
\TextAt{9}{315.337}{445.422}{\Character{41}}%
\NativeMathAt{326.494}{452.2}{7.97}{\sim}%
\NativeArrowAt{323.907}{445.422}{11.955}{13.175}{\longrightarrow}%
\TextAt{8}{341.097}{445.422}{ker}%
\TextAt{9}{357.034}{445.422}{\Character{40}}%
\TextAt{13}{361.586}{445.422}{\Character{12}}%
\TextAt{9}{368.857}{445.422}{\Character{41}}%
\TextAt{8}{373.41}{445.422}{,}%
\TextAt{8}{381.07}{445.422}{so}%
\TextAt{8}{395.849}{445.422}{that}%
\TextAt{8}{421.627}{445.422}{we}%
\TextAt{8}{439.263}{445.422}{have}%
\TextAt{8}{466.667}{445.422}{the}%
\TextAt{8}{487.232}{445.422}{following}%
\TextAt{8}{79.226}{428.377}{identity}%
\TextAt{8}{122.48}{428.377}{in}%
\TextAt{13}{136.133}{428.377}{K}%
\TextAt{17}{146.092}{426.583}{0}%
\TextAt{9}{150.824}{428.377}{\Character{40}}%
\Vector{q 0 G 0 g 0.478 w 83.377 -338.267 m 89.221 -338.267 l S Q}%
\TextAt{14}{155.377}{428.377}{o}%
\TextAt{9}{161.221}{428.377}{\Character{41}}%
\TextAt{8}{165.773}{428.377}{:}%
\TextAt{9}{160.298}{401.781}{\Character{91}}%
\TextAt{13}{163.549}{401.781}{F}%
\TextAt{9}{172.753}{401.781}{\Character{93}}%
\TextAt{16}{178.661}{401.781}{\Character{0}}%
\TextAt{9}{190.616}{401.781}{\Character{91}}%
\TextAt{13}{193.868}{401.781}{E}%
\TextAt{9}{203.234}{401.781}{\Character{93}}%
\TextAt{9}{209.785}{401.781}{\Character{61}}%
\TextAt{9}{222.207}{401.781}{\Character{91}}%
\TextAt{13}{225.484}{401.781}{F}%
\TextAt{3}{234.687}{406.717}{\Character{48}}%
\TextAt{9}{237.482}{401.781}{\Character{93}}%
\TextAt{9}{243.376}{401.781}{\Character{43}}%
\TextAt{9}{255.128}{401.781}{\Character{91}}%
\TextAt{8}{258.404}{401.781}{ker}%
\TextAt{9}{274.34}{401.781}{\Character{40}}%
\TextAt{13}{278.892}{401.781}{\Character{12}}%
\TextAt{9}{286.164}{401.781}{\Character{41}\Character{93}}%
\TextAt{16}{296.624}{401.781}{\Character{0}}%
\TextAt{9}{308.579}{401.781}{\Character{91}}%
\TextAt{13}{311.831}{401.781}{E}%
\TextAt{3}{321.197}{406.717}{\Character{48}}%
\TextAt{9}{323.992}{401.781}{\Character{93}}%
\TextAt{16}{329.901}{401.781}{\Character{0}}%
\TextAt{9}{341.856}{401.781}{\Character{91}}%
\TextAt{8}{345.108}{401.781}{ker}%
\TextAt{9}{361.044}{401.781}{\Character{40}}%
\TextAt{13}{365.596}{401.781}{\Character{11}}%
\TextAt{9}{373.118}{401.781}{\Character{41}\Character{93}}%
\TextAt{9}{384.212}{401.781}{\Character{61}}%
\TextAt{9}{396.634}{401.781}{\Character{91}}%
\TextAt{13}{399.92}{401.781}{F}%
\TextAt{3}{409.123}{406.717}{\Character{48}}%
\TextAt{9}{411.918}{401.781}{\Character{93}}%
\TextAt{16}{417.827}{401.781}{\Character{0}}%
\TextAt{9}{429.782}{401.781}{\Character{91}}%
\TextAt{13}{433.034}{401.781}{E}%
\TextAt{3}{442.4}{406.717}{\Character{48}}%
\TextAt{9}{445.195}{401.781}{\Character{93}}%
\TextAt{13}{448.447}{401.781}{\Character{58}}%
\TextAt{8}{79.226}{375.185}{If}%
\TextAt{8}{90.634}{375.185}{now}%
\TextAt{9}{114.739}{375.185}{0}%
\NativeMathAt{123.913}{375.185}{11.955}{\rightarrow}%
\TextAt{13}{139.189}{375.185}{E}%
\TextAt{3}{148.555}{379.523}{\Character{48}}%
\NativeMathAt{154.671}{375.185}{11.955}{\rightarrow}%
\TextAt{13}{169.947}{375.185}{F}%
\TextAt{3}{179.151}{379.523}{\Character{48}}%
\NativeMathAt{185.267}{375.185}{11.955}{\rightarrow}%
\TextAt{13}{200.543}{375.185}{M}%
\NativeMathAt{216.437}{375.185}{11.955}{\rightarrow}%
\TextAt{9}{231.713}{375.185}{0}%
\TextAt{8}{237.566}{375.185}{,}%
\TextAt{8}{244.488}{375.185}{and}%
\TextAt{9}{266.967}{375.185}{0}%
\NativeMathAt{276.141}{375.185}{11.955}{\rightarrow}%
\TextAt{13}{291.417}{375.185}{E}%
\TextAt{3}{300.783}{379.523}{\Character{48}\Character{48}}%
\NativeMathAt{309.196}{375.185}{11.955}{\rightarrow}%
\TextAt{13}{324.472}{375.185}{F}%
\TextAt{3}{333.676}{379.523}{\Character{48}\Character{48}}%
\NativeMathAt{342.089}{375.185}{11.955}{\rightarrow}%
\TextAt{13}{357.365}{375.185}{M}%
\NativeMathAt{373.259}{375.185}{11.955}{\rightarrow}%
\TextAt{9}{388.535}{375.185}{0}%
\TextAt{8}{398.002}{375.185}{are}%
\TextAt{8}{417.229}{375.185}{two}%
\TextAt{8}{439.062}{375.185}{arbitrary}%
\TextAt{8}{488.547}{375.185}{metrized}%
\TextAt{8}{79.226}{358.139}{projective}%
\TextAt{8}{133.667}{358.139}{resolutions}%
\TextAt{8}{192.794}{358.139}{of}%
\TextAt{13}{206.587}{358.139}{M}%
\TextAt{8}{219.16}{358.139}{,}%
\TextAt{8}{226.888}{358.139}{then}%
\TextAt{8}{254.011}{358.139}{by}%
\TextAt{8}{270.733}{358.139}{(5.3)}%
\TextAt{8}{299.171}{358.139}{we}%
\TextAt{8}{316.861}{358.139}{find}%
\TextAt{8}{340.732}{358.139}{a}%
\TextAt{8}{350.951}{358.139}{third}%
\TextAt{8}{380.68}{358.139}{one,}%
\TextAt{9}{405.97}{358.139}{0}%
\NativeMathAt{415.933}{358.139}{11.955}{\rightarrow}%
\TextAt{13}{431.997}{358.139}{E}%
\NativeMathAt{445.473}{358.139}{11.955}{\rightarrow}%
\TextAt{13}{461.538}{358.139}{F}%
\NativeMathAt{474.852}{358.139}{11.955}{\rightarrow}%
\TextAt{13}{490.916}{358.139}{M}%
\NativeMathAt{507.6}{358.139}{11.955}{\rightarrow}%
\TextAt{9}{523.665}{358.139}{0}%
\TextAt{8}{529.518}{358.139}{,}%
\TextAt{8}{79.226}{341.093}{dominating}%
\TextAt{8}{141.022}{341.093}{both,}%
\TextAt{8}{171.915}{341.093}{such}%
\TextAt{8}{198.312}{341.093}{that}%
\TextAt{9}{211.456}{314.497}{\Character{91}}%
\TextAt{13}{214.708}{314.497}{F}%
\TextAt{3}{223.911}{319.433}{\Character{48}}%
\TextAt{9}{226.707}{314.497}{\Character{93}}%
\TextAt{16}{232.615}{314.497}{\Character{0}}%
\TextAt{9}{244.57}{314.497}{\Character{91}}%
\TextAt{13}{247.822}{314.497}{E}%
\TextAt{3}{257.188}{319.433}{\Character{48}}%
\TextAt{9}{259.983}{314.497}{\Character{93}}%
\TextAt{9}{266.534}{314.497}{\Character{61}}%
\TextAt{9}{278.956}{314.497}{\Character{91}}%
\TextAt{13}{282.233}{314.497}{F}%
\TextAt{9}{291.436}{314.497}{\Character{93}}%
\TextAt{16}{297.345}{314.497}{\Character{0}}%
\TextAt{9}{309.3}{314.497}{\Character{91}}%
\TextAt{13}{312.551}{314.497}{E}%
\TextAt{9}{321.918}{314.497}{\Character{93}}%
\TextAt{9}{328.469}{314.497}{\Character{61}}%
\TextAt{9}{340.891}{314.497}{\Character{91}}%
\TextAt{13}{344.167}{314.497}{F}%
\TextAt{3}{353.371}{319.433}{\Character{48}\Character{48}}%
\TextAt{9}{358.463}{314.497}{\Character{93}}%
\TextAt{16}{364.371}{314.497}{\Character{0}}%
\TextAt{9}{376.326}{314.497}{\Character{91}}%
\TextAt{13}{379.578}{314.497}{E}%
\TextAt{3}{388.944}{319.433}{\Character{48}\Character{48}}%
\TextAt{9}{394.036}{314.497}{\Character{93}}%
\TextAt{13}{397.288}{314.497}{\Character{58}}%
\TextAt{8}{79.226}{287.901}{This}%
\TextAt{8}{107.705}{287.901}{shows}%
\TextAt{8}{143.082}{287.901}{that}%
\TextAt{8}{170.209}{287.901}{the}%
\TextAt{8}{192.125}{287.901}{map}%
\TextAt{13}{219.898}{287.901}{\Character{25}}%
\TextAt{9}{233.285}{287.901}{\Character{58}}%
\TextAt{13}{242.855}{287.901}{F}%
\TextAt{17}{252.059}{292.239}{0}%
\TextAt{9}{256.791}{287.901}{\Character{40}}%
\Vector{q 0 G 0 g 0.478 w 189.343 -478.743 m 195.188 -478.743 l S Q}%
\TextAt{14}{261.343}{287.901}{o}%
\TextAt{9}{267.188}{287.901}{\Character{41}}%
\NativeMathAt{278.058}{287.901}{11.955}{\rightarrow}%
\TextAt{13}{296.332}{287.901}{K}%
\TextAt{17}{306.291}{286.108}{0}%
\TextAt{9}{311.023}{287.901}{\Character{40}}%
\Vector{q 0 G 0 g 0.478 w 243.575 -478.743 m 249.42 -478.743 l S Q}%
\TextAt{14}{315.575}{287.901}{o}%
\TextAt{9}{321.42}{287.901}{\Character{41}}%
\TextAt{8}{331.628}{287.901}{is}%
\TextAt{8}{345.151}{287.901}{well-defined.}%
\TextAt{8}{419.022}{287.901}{We}%
\TextAt{8}{440.926}{287.901}{now}%
\TextAt{8}{467.073}{287.901}{show}%
\TextAt{8}{497.835}{287.901}{that}%
\TextAt{8}{524.963}{287.901}{it}%
\TextAt{8}{79.226}{270.855}{factorizes}%
\TextAt{8}{131.923}{270.855}{via}%
\TextAt{13}{152.042}{270.855}{K}%
\TextAt{17}{162.859}{275.193}{0}%
\TextAt{9}{167.592}{270.855}{\Character{40}}%
\Vector{q 0 G 0 g 0.478 w 100.144 -495.789 m 105.988 -495.789 l S Q}%
\TextAt{14}{172.144}{270.855}{o}%
\TextAt{9}{177.988}{270.855}{\Character{41}}%
\TextAt{9}{187.433}{270.855}{\Character{61}}%
\TextAt{13}{201.46}{270.855}{F}%
\TextAt{17}{210.663}{275.193}{0}%
\TextAt{9}{215.396}{270.855}{\Character{40}}%
\Vector{q 0 G 0 g 0.478 w 147.948 -495.789 m 153.792 -495.789 l S Q}%
\TextAt{14}{219.948}{270.855}{o}%
\TextAt{9}{225.792}{270.855}{\Character{41}\Character{47}}%
\TextAt{13}{236.197}{270.855}{R}%
\TextAt{17}{245.206}{275.193}{0}%
\TextAt{9}{249.938}{270.855}{\Character{40}}%
\Vector{q 0 G 0 g 0.478 w 182.491 -495.789 m 188.335 -495.789 l S Q}%
\TextAt{14}{254.491}{270.855}{o}%
\TextAt{9}{260.335}{270.855}{\Character{41}}%
\TextAt{8}{264.887}{270.855}{.}%
\TextAt{8}{276.128}{270.855}{Let}%
\TextAt{9}{298.028}{270.855}{0}%
\NativeMathAt{308.789}{270.855}{11.955}{\rightarrow}%
\TextAt{13}{325.651}{270.855}{M}%
\TextAt{3}{338.224}{275.193}{\Character{48}}%
\NativeMathAt{345.927}{270.855}{11.955}{\rightarrow}%
\TextAt{13}{362.789}{270.855}{M}%
\TextAt{21}{382.857}{277.632}{\Character{11}}%
\NativeArrowAt{380.27}{270.855}{11.955}{12.027}{\longrightarrow}%
\TextAt{13}{397.202}{270.855}{M}%
\TextAt{3}{409.776}{275.193}{\Character{48}\Character{48}}%
\NativeMathAt{419.775}{270.855}{11.955}{\rightarrow}%
\TextAt{9}{436.638}{270.855}{0}%
\TextAt{8}{447.319}{270.855}{be}%
\TextAt{8}{464.174}{270.855}{a}%
\TextAt{8}{474.86}{270.855}{short}%
\TextAt{8}{505.775}{270.855}{exact}%
\TextAt{8}{79.226}{253.809}{sequence}%
\TextAt{8}{128.802}{253.809}{of}%
\TextAt{8}{143.205}{253.809}{metrized}%
\TextAt{8}{192.398}{253.809}{coherent}%
\TextAt{14}{240.623}{253.809}{o}%
\TextAt{8}{246.467}{253.809}{-modules.}%
\TextAt{8}{304.044}{253.809}{By}%
\TextAt{8}{323.48}{253.809}{(5.2),}%
\TextAt{8}{356.049}{253.809}{we}%
\TextAt{8}{374.35}{253.809}{can}%
\TextAt{8}{396.883}{253.809}{pick}%
\TextAt{8}{422.667}{253.809}{a}%
\TextAt{8}{433.496}{253.809}{metrized}%
\TextAt{8}{482.689}{253.809}{projective}%
\TextAt{8}{79.226}{235.219}{resolution}%
\TextAt{9}{133.895}{235.219}{0}%
\NativeMathAt{144.126}{235.219}{11.955}{\rightarrow}%
\TextAt{13}{160.458}{235.219}{E}%
\NativeMathAt{174.202}{235.219}{11.955}{\rightarrow}%
\TextAt{13}{190.534}{235.219}{F}%
\TextAt{21}{206.914}{242.483}{f}%
\NativeArrowAt{204.115}{235.219}{11.955}{11.955}{\longrightarrow}%
\TextAt{13}{220.448}{235.219}{M}%
\NativeMathAt{237.399}{235.219}{11.955}{\rightarrow}%
\TextAt{9}{253.731}{235.219}{0}%
\TextAt{8}{259.584}{235.219}{.}%
\TextAt{8}{269.892}{235.219}{Then}%
\TextAt{8}{301.07}{235.219}{clearly}%
\TextAt{9}{339.085}{235.219}{0}%
\NativeMathAt{349.316}{235.219}{11.955}{\rightarrow}%
\TextAt{13}{365.648}{235.219}{E}%
\TextAt{3}{375.014}{239.557}{\Character{48}\Character{48}}%
\NativeMathAt{384.484}{235.219}{11.955}{\rightarrow}%
\TextAt{13}{400.816}{235.219}{F}%
\TextAt{21}{416.985}{242.483}{f}%
\TextAt{23}{421.931}{245.296}{\Character{48}\Character{48}}%
\NativeArrowAt{414.397}{235.219}{11.955}{16.403}{\longrightarrow}%
\TextAt{13}{435.178}{235.219}{M}%
\TextAt{3}{447.752}{239.557}{\Character{48}\Character{48}}%
\NativeMathAt{457.221}{235.219}{11.955}{\rightarrow}%
\TextAt{9}{473.554}{235.219}{0}%
\TextAt{8}{483.924}{235.219}{is}%
\TextAt{8}{496.308}{235.219}{a}%
\TextAt{8}{506.683}{235.219}{short}%
\TextAt{8}{79.226}{218.173}{exact}%
\TextAt{8}{109.093}{218.173}{sequence}%
\TextAt{8}{156.57}{218.173}{of}%
\TextAt{8}{168.875}{218.173}{metrized}%
\TextAt{14}{215.969}{218.173}{o}%
\TextAt{8}{221.813}{218.173}{-modules}%
\TextAt{8}{270.593}{218.173}{as}%
\TextAt{8}{283.938}{218.173}{well,}%
\TextAt{8}{310.101}{218.173}{where}%
\TextAt{8}{342.885}{218.173}{we}%
\TextAt{8}{359.087}{218.173}{write}%
\TextAt{13}{387.974}{218.173}{f}%
\TextAt{3}{395.02}{222.511}{\Character{48}\Character{48}}%
\TextAt{9}{403.433}{218.173}{\Character{61}}%
\TextAt{13}{415.858}{218.173}{\Character{11}}%
\TextAt{16}{423.94}{218.173}{\Character{14}}%
\TextAt{13}{430.478}{218.173}{f}%
\TextAt{8}{440.396}{218.173}{and}%
\TextAt{13}{462.134}{218.173}{E}%
\TextAt{3}{471.5}{222.511}{\Character{48}\Character{48}}%
\TextAt{9}{479.913}{218.173}{\Character{61}}%
\TextAt{8}{492.338}{218.173}{ker}%
\TextAt{9}{508.275}{218.173}{\Character{40}}%
\TextAt{13}{512.827}{218.173}{f}%
\TextAt{3}{519.873}{222.511}{\Character{48}\Character{48}}%
\TextAt{9}{524.965}{218.173}{\Character{41}}%
\TextAt{8}{529.518}{218.173}{.}%
\TextAt{8}{79.226}{201.127}{We}%
\TextAt{8}{99.371}{201.127}{thus}%
\TextAt{8}{125.122}{201.127}{get}%
\TextAt{8}{144.633}{201.127}{the}%
\TextAt{8}{164.789}{201.127}{commutative}%
\TextAt{8}{234.715}{201.127}{diagram}%
\TextAt{9}{172.791}{172.665}{0}%
\TextAt{13}{234.535}{172.665}{E}%
\TextAt{13}{299.791}{172.665}{F}%
\TextAt{13}{364.884}{172.665}{M}%
\TextAt{9}{433.348}{172.665}{0}%
\TextAt{9}{172.791}{122.86}{0}%
\TextAt{13}{231.989}{122.86}{E}%
\TextAt{3}{241.355}{127.198}{\Character{48}\Character{48}}%
\TextAt{13}{299.791}{122.86}{F}%
\TextAt{13}{362.338}{122.86}{M}%
\TextAt{3}{374.912}{127.198}{\Character{48}\Character{48}}%
\TextAt{9}{433.348}{122.86}{0}%
\NativeDiagram{%
\path[draw=black,line width=0.478bp] (111.997,-598.346) -- (156.705,-598.346);%
\path[draw=black,line width=0.478bp,line cap=round,line join=round] (154.458,-595.477) .. controls (154.912,-597.199) and (155.924,-598.011) .. (156.944,-598.346) .. controls (155.924,-598.681) and (154.912,-599.494) .. (154.458,-601.215);%
\path[draw=black,line width=0.478bp] (177.254,-598.346) -- (221.962,-598.346);%
\path[draw=black,line width=0.478bp,line cap=round,line join=round] (219.715,-595.477) .. controls (220.169,-597.199) and (221.181,-598.011) .. (222.201,-598.346) .. controls (221.181,-598.681) and (220.169,-599.494) .. (219.715,-601.215);%
\path[draw=black,line width=0.478bp] (167.219,-605.914) -- (167.219,-637.317);%
\path[draw=black,line width=0.478bp,line cap=round,line join=round] (170.088,-635.069) .. controls (168.366,-635.524) and (167.553,-636.536) .. (167.219,-637.556) .. controls (166.884,-636.536) and (166.071,-635.524) .. (164.349,-635.069);%
\path[draw=black,line width=0.478bp] (242.348,-598.346) -- (287.057,-598.346);%
\path[draw=black,line width=0.478bp,line cap=round,line join=round] (284.809,-595.477) .. controls (285.263,-597.199) and (286.276,-598.011) .. (287.296,-598.346) .. controls (286.276,-598.681) and (285.263,-599.494) .. (284.809,-601.215);%
}%
\TextAt{21}{334.468}{180.019}{f}%
\NativeDiagram{%
\path[draw=black,line width=0.478bp] (232.394,-605.914) -- (232.394,-637.914);%
\path[draw=black,line width=0.478bp,line cap=round,line join=round] (235.263,-635.666) .. controls (233.542,-636.12) and (232.729,-637.133) .. (232.394,-638.153) .. controls (232.06,-637.133) and (231.247,-636.12) .. (229.525,-635.666);%
}%
\TextAt{7}{307.21}{149.121}{id}%
\NativeDiagram{%
\path[draw=black,line width=0.478bp] (310.813,-598.346) -- (355.521,-598.346);%
\path[draw=black,line width=0.478bp,line cap=round,line join=round] (353.274,-595.477) .. controls (353.728,-597.199) and (354.74,-598.011) .. (355.76,-598.346) .. controls (354.74,-598.681) and (353.728,-599.494) .. (353.274,-601.215);%
\path[draw=black,line width=0.478bp] (299.174,-605.914) -- (299.174,-637.317);%
\path[draw=black,line width=0.478bp,line cap=round,line join=round] (302.043,-635.069) .. controls (300.321,-635.524) and (299.509,-636.536) .. (299.174,-637.556) .. controls (298.839,-636.536) and (298.026,-635.524) .. (296.305,-635.069);%
}%
\TextAt{21}{373.989}{150.43}{\Character{11}}%
\NativeDiagram{%
\path[draw=black,line width=0.478bp] (111.997,-648.152) -- (154.159,-648.152);%
\path[draw=black,line width=0.478bp,line cap=round,line join=round] (151.912,-645.282) .. controls (152.366,-647.004) and (153.378,-647.817) .. (154.398,-648.152) .. controls (153.378,-648.486) and (152.366,-649.299) .. (151.912,-651.021);%
\path[draw=black,line width=0.478bp] (179.8,-648.152) -- (221.962,-648.152);%
\path[draw=black,line width=0.478bp,line cap=round,line join=round] (219.715,-645.282) .. controls (220.169,-647.004) and (221.181,-647.817) .. (222.201,-648.152) .. controls (221.181,-648.486) and (220.169,-649.299) .. (219.715,-651.021);%
\path[draw=black,line width=0.478bp] (242.348,-648.152) -- (284.511,-648.152);%
\path[draw=black,line width=0.478bp,line cap=round,line join=round] (282.263,-645.282) .. controls (282.717,-647.004) and (283.73,-647.817) .. (284.75,-648.152) .. controls (283.73,-648.486) and (282.717,-649.299) .. (282.263,-651.021);%
}%
\TextAt{21}{330.76}{130.214}{f}%
\TextAt{23}{335.707}{133.027}{\Character{48}\Character{48}}%
\NativeDiagram{%
\path[draw=black,line width=0.478bp] (313.359,-648.152) -- (355.521,-648.152);%
\path[draw=black,line width=0.478bp,line cap=round,line join=round] (353.274,-645.282) .. controls (353.728,-647.004) and (354.74,-647.817) .. (355.76,-648.152) .. controls (354.74,-648.486) and (353.728,-649.299) .. (353.274,-651.021);%
}%
\TextAt{8}{79.226}{99.517}{and}%
\TextAt{8}{101.989}{99.517}{the}%
\TextAt{8}{122.145}{99.517}{snake}%
\TextAt{8}{154.077}{99.517}{lemma}%
\TextAt{8}{191.796}{99.517}{gives}%
\TextAt{8}{220.476}{99.517}{the}%
\TextAt{8}{240.633}{99.517}{exact}%
\TextAt{8}{271.525}{99.517}{sequence}%
\TextAt{8}{320.027}{99.517}{of}%
\TextAt{14}{333.357}{99.517}{o}%
\TextAt{8}{339.201}{99.517}{-modules}%
\TextAt{9}{230.77}{72.921}{0}%
\NativeArrowAt{239.943}{72.921}{11.955}{19.26}{\longrightarrow}%
\TextAt{13}{262.526}{72.921}{E}%
\NativeArrowAt{275.213}{72.921}{11.955}{19.26}{\longrightarrow}%
\TextAt{13}{297.795}{72.921}{E}%
\TextAt{3}{307.161}{77.858}{\Character{48}\Character{48}}%
\TextAt{21}{318.372}{80.185}{f}%
\NativeArrowAt{315.574}{72.921}{11.955}{11.955}{\longrightarrow}%
\TextAt{13}{330.85}{72.921}{M}%
\TextAt{3}{343.423}{77.858}{\Character{48}}%
\NativeArrowAt{349.539}{72.921}{11.955}{19.26}{\longrightarrow}%
\TextAt{9}{372.121}{72.921}{0}%
\TextAt{13}{377.974}{72.921}{\Character{58}}%
\TextAt{8}{297.211}{43.033}{218}%
\Destination{78.23}{809.815}{nav.page.218}%
\Destination{79.36}{722.955}{nav.print.237}%
\Destination{79.36}{390.155}{nav.print.238}%
\Destination{79.36}{775.155}{nav.number.3.5.4}%
\Hotlink{270.233}{355.955}{25.078}{11.047}{nav.number.3.5.3}%
\Hotlink{322.98}{251.655}{25.078}{11.047}{nav.number.3.5.2}%
}

% Page 231 (body page 219)
\Page{%
\TextAt{2}{79.226}{796.943}{III.}%
\TextAt{2}{98.317}{796.943}{Riemann–Roch}%
\TextAt{2}{174.375}{796.943}{Theory}%
\Vector{q 0 G 0 g 0.299 w 7.226 18.729 m 460.769 18.729 l S Q}%
\ErrataAt{79.19}{761.124}{455.81}{It is a short exact sequence of metrized $\mathfrak o$-modules. In fact $E^{\prime\prime}_{\mathbb C}=f^{-1}(M^{\prime}_{\mathbb C})$, so the}%
\ErrataAt{79.19}{744.078}{455.81}{isometry $f:E_{\mathbb C}^{\perp}\to M_{\mathbb C}$ carries $E^{\prime\prime}_{\mathbb C}\cap E_{\mathbb C}^{\perp}$ isometrically onto $M^{\prime}_{\mathbb C}$. We therefore}%
\ErrataAt{79.19}{727.032}{455.81}{obtain in $K_0(\overline{\mathfrak o})$ the identity}%
\Vector{q 0 G 0 g 0.478 w 76.55 -39.612 m 82.394 -39.612 l S Q}%
\TextAt{13}{138.746}{698.031}{\Character{25}}%
\TextAt{16}{145.815}{698.031}{\Character{102}}%
\TextAt{13}{151.793}{698.031}{M}%
\TextAt{3}{164.367}{702.967}{\Character{48}}%
\TextAt{16}{167.162}{698.031}{\Character{103}}%
\TextAt{16}{175.794}{698.031}{\Character{0}}%
\TextAt{13}{187.751}{698.031}{\Character{25}}%
\TextAt{16}{194.82}{698.031}{\Character{102}}%
\TextAt{13}{200.798}{698.031}{M}%
\TextAt{16}{213.372}{698.031}{\Character{103}}%
\TextAt{9}{222.006}{698.031}{\Character{43}}%
\TextAt{13}{233.767}{698.031}{\Character{25}}%
\TextAt{16}{240.837}{698.031}{\Character{102}}%
\TextAt{13}{246.814}{698.031}{M}%
\TextAt{3}{259.388}{702.967}{\Character{48}\Character{48}}%
\TextAt{16}{264.48}{698.031}{\Character{103}}%
\TextAt{9}{273.778}{698.031}{\Character{61}}%
\TextAt{9}{286.188}{698.031}{\Character{91}}%
\TextAt{13}{289.455}{698.031}{E}%
\TextAt{3}{298.821}{702.967}{\Character{48}\Character{48}}%
\TextAt{9}{303.913}{698.031}{\Character{93}}%
\TextAt{16}{309.822}{698.031}{\Character{0}}%
\TextAt{9}{321.777}{698.031}{\Character{91}}%
\TextAt{13}{325.029}{698.031}{E}%
\TextAt{9}{334.395}{698.031}{\Character{93}}%
\TextAt{16}{340.303}{698.031}{\Character{0}}%
\TextAt{9}{352.258}{698.031}{\Character{40}\Character{91}}%
\TextAt{13}{360.062}{698.031}{F}%
\TextAt{9}{369.266}{698.031}{\Character{93}}%
\TextAt{16}{375.174}{698.031}{\Character{0}}%
\TextAt{9}{387.129}{698.031}{\Character{91}}%
\TextAt{13}{390.381}{698.031}{E}%
\TextAt{9}{399.747}{698.031}{\Character{93}\Character{41}}%
\TextAt{9}{410.184}{698.031}{\Character{43}}%
\TextAt{9}{421.936}{698.031}{\Character{91}}%
\TextAt{13}{425.221}{698.031}{F}%
\TextAt{9}{434.424}{698.031}{\Character{93}}%
\TextAt{16}{440.333}{698.031}{\Character{0}}%
\TextAt{9}{452.288}{698.031}{\Character{91}}%
\TextAt{13}{455.54}{698.031}{E}%
\TextAt{3}{464.906}{702.967}{\Character{48}\Character{48}}%
\TextAt{9}{469.998}{698.031}{\Character{93}}%
\TextAt{9}{273.778}{677.996}{\Character{61}}%
\TextAt{9}{286.188}{677.996}{0}%
\TextAt{13}{292.057}{677.996}{\Character{58}}%
\TextAt{8}{79.226}{648.995}{It}%
\TextAt{8}{91.899}{648.995}{shows}%
\TextAt{8}{125.517}{648.995}{that}%
\TextAt{13}{150.885}{648.995}{\Character{25}}%
\TextAt{9}{161.276}{648.995}{\Character{58}}%
\TextAt{13}{167.848}{648.995}{F}%
\TextAt{17}{177.052}{653.333}{0}%
\TextAt{9}{181.784}{648.995}{\Character{40}}%
\Vector{q 0 G 0 g 0.478 w 114.336 -117.649 m 120.181 -117.649 l S Q}%
\TextAt{14}{186.336}{648.995}{o}%
\TextAt{9}{192.181}{648.995}{\Character{41}}%
\NativeMathAt{200.054}{648.995}{11.955}{\rightarrow}%
\TextAt{13}{215.33}{648.995}{K}%
\TextAt{17}{225.289}{647.201}{0}%
\TextAt{9}{230.021}{648.995}{\Character{40}}%
\Vector{q 0 G 0 g 0.478 w 162.573 -117.649 m 168.418 -117.649 l S Q}%
\TextAt{14}{234.573}{648.995}{o}%
\TextAt{9}{240.418}{648.995}{\Character{41}}%
\TextAt{8}{248.867}{648.995}{does}%
\TextAt{8}{275.264}{648.995}{indeed}%
\TextAt{8}{312.325}{648.995}{factorize}%
\TextAt{8}{359.477}{648.995}{via}%
\TextAt{8}{378.665}{648.995}{a}%
\TextAt{8}{388.42}{648.995}{homomorphism}%
\TextAt{13}{261.351}{619.994}{K}%
\TextAt{17}{272.169}{624.93}{0}%
\TextAt{9}{276.901}{619.994}{\Character{40}}%
\Vector{q 0 G 0 g 0.478 w 209.453 -146.65 m 215.298 -146.65 l S Q}%
\TextAt{14}{281.453}{619.994}{o}%
\TextAt{9}{287.298}{619.994}{\Character{41}}%
\NativeArrowAt{295.171}{619.994}{11.955}{19.26}{\longrightarrow}%
\TextAt{13}{317.753}{619.994}{K}%
\TextAt{17}{327.712}{618.2}{0}%
\TextAt{9}{332.444}{619.994}{\Character{40}}%
\Vector{q 0 G 0 g 0.478 w 264.996 -146.65 m 270.841 -146.65 l S Q}%
\TextAt{14}{336.996}{619.994}{o}%
\TextAt{9}{342.841}{619.994}{\Character{41}}%
\TextAt{13}{347.393}{619.994}{\Character{58}}%
\TextAt{8}{79.226}{590.992}{It}%
\TextAt{8}{91.899}{590.992}{is}%
\TextAt{8}{103.663}{590.992}{the}%
\TextAt{8}{123.819}{590.992}{inverse}%
\TextAt{8}{162.578}{590.992}{of}%
\TextAt{8}{175.908}{590.992}{the}%
\TextAt{8}{196.064}{590.992}{Poincaré}%
\TextAt{8}{244.028}{590.992}{homomorphism}%
\TextAt{8}{326.698}{590.992}{because}%
\TextAt{8}{370}{590.992}{the}%
\TextAt{8}{390.156}{590.992}{composed}%
\TextAt{8}{443.871}{590.992}{maps}%
\TextAt{13}{130.64}{561.991}{K}%
\TextAt{17}{140.599}{560.198}{0}%
\TextAt{9}{145.332}{561.991}{\Character{40}}%
\Vector{q 0 G 0 g 0.478 w 77.884 -204.652 m 83.728 -204.652 l S Q}%
\TextAt{14}{149.884}{561.991}{o}%
\TextAt{9}{155.728}{561.991}{\Character{41}}%
\NativeArrowAt{163.601}{561.991}{11.955}{19.26}{\longrightarrow}%
\TextAt{13}{186.184}{561.991}{K}%
\TextAt{17}{197.001}{566.927}{0}%
\TextAt{9}{201.734}{561.991}{\Character{40}}%
\Vector{q 0 G 0 g 0.478 w 134.286 -204.652 m 140.13 -204.652 l S Q}%
\TextAt{14}{206.286}{561.991}{o}%
\TextAt{9}{212.13}{561.991}{\Character{41}}%
\NativeArrowAt{220.003}{561.991}{11.955}{19.26}{\longrightarrow}%
\TextAt{13}{242.585}{561.991}{K}%
\TextAt{17}{252.544}{560.198}{0}%
\TextAt{9}{257.277}{561.991}{\Character{40}}%
\Vector{q 0 G 0 g 0.478 w 189.829 -204.652 m 195.673 -204.652 l S Q}%
\TextAt{14}{261.829}{561.991}{o}%
\TextAt{9}{267.673}{561.991}{\Character{41}}%
\TextAt{8}{296.136}{561.991}{and}%
\TextAt{13}{338.911}{561.991}{K}%
\TextAt{17}{349.729}{566.927}{0}%
\TextAt{9}{354.461}{561.991}{\Character{40}}%
\Vector{q 0 G 0 g 0.478 w 287.014 -204.652 m 292.858 -204.652 l S Q}%
\TextAt{14}{359.014}{561.991}{o}%
\TextAt{9}{364.858}{561.991}{\Character{41}}%
\NativeArrowAt{372.731}{561.991}{11.955}{19.26}{\longrightarrow}%
\TextAt{13}{395.313}{561.991}{K}%
\TextAt{17}{405.272}{560.198}{0}%
\TextAt{9}{410.005}{561.991}{\Character{40}}%
\Vector{q 0 G 0 g 0.478 w 342.557 -204.652 m 348.401 -204.652 l S Q}%
\TextAt{14}{414.557}{561.991}{o}%
\TextAt{9}{420.401}{561.991}{\Character{41}}%
\NativeArrowAt{428.274}{561.991}{11.955}{19.26}{\longrightarrow}%
\TextAt{13}{450.856}{561.991}{K}%
\TextAt{17}{461.674}{566.927}{0}%
\TextAt{9}{466.406}{561.991}{\Character{40}}%
\Vector{q 0 G 0 g 0.478 w 398.959 -204.652 m 404.803 -204.652 l S Q}%
\TextAt{14}{470.959}{561.991}{o}%
\TextAt{9}{476.803}{561.991}{\Character{41}}%
\TextAt{8}{79.226}{532.99}{are}%
\TextAt{8}{100.105}{532.99}{the}%
\TextAt{8}{121.63}{532.99}{identity}%
\TextAt{8}{166.252}{532.99}{homomorphisms.}%
\TextAt{8}{262.193}{532.99}{Indeed,}%
\TextAt{8}{305.184}{532.99}{if}%
\TextAt{9}{317.276}{532.99}{0}%
\NativeMathAt{328.782}{532.99}{11.955}{\rightarrow}%
\TextAt{13}{346.39}{532.99}{E}%
\NativeMathAt{361.409}{532.99}{11.955}{\rightarrow}%
\TextAt{13}{379.017}{532.99}{F}%
\NativeMathAt{393.873}{532.99}{11.955}{\rightarrow}%
\TextAt{13}{411.481}{532.99}{M}%
\NativeMathAt{429.707}{532.99}{11.955}{\rightarrow}%
\TextAt{9}{447.315}{532.99}{0}%
\TextAt{8}{458.433}{532.99}{is}%
\TextAt{8}{471.566}{532.99}{a}%
\TextAt{8}{482.689}{532.99}{projective}%
\TextAt{8}{79.226}{515.944}{resolution}%
\TextAt{8}{132.958}{515.944}{of}%
\TextAt{13}{145.971}{515.944}{M}%
\TextAt{8}{158.544}{515.944}{,}%
\TextAt{8}{165.44}{515.944}{and}%
\TextAt{13}{187.885}{515.944}{M}%
\TextAt{8}{204.039}{515.944}{is}%
\TextAt{8}{215.485}{515.944}{projective,}%
\TextAt{8}{272.461}{515.944}{resp.}%
\TextAt{8}{301.677}{515.944}{coherent,}%
\TextAt{8}{351.826}{515.944}{then}%
\TextAt{8}{378.169}{515.944}{in}%
\TextAt{13}{391.505}{515.944}{K}%
\TextAt{17}{401.464}{514.151}{0}%
\TextAt{9}{406.196}{515.944}{\Character{40}}%
\Vector{q 0 G 0 g 0.478 w 338.748 -250.699 m 344.593 -250.699 l S Q}%
\TextAt{14}{410.748}{515.944}{o}%
\TextAt{9}{416.593}{515.944}{\Character{41}}%
\TextAt{8}{421.145}{515.944}{,}%
\TextAt{8}{428.04}{515.944}{resp.}%
\TextAt{13}{457.257}{515.944}{K}%
\TextAt{17}{468.074}{520.283}{0}%
\TextAt{9}{472.807}{515.944}{\Character{40}}%
\Vector{q 0 G 0 g 0.478 w 405.359 -250.699 m 411.203 -250.699 l S Q}%
\TextAt{14}{477.359}{515.944}{o}%
\TextAt{9}{483.203}{515.944}{\Character{41}}%
\TextAt{8}{487.756}{515.944}{,}%
\TextAt{8}{494.651}{515.944}{one}%
\TextAt{8}{515.793}{515.944}{has}%
\TextAt{8}{79.226}{498.898}{the}%
\TextAt{8}{99.383}{498.898}{identity}%
\TextAt{9}{142.636}{498.898}{\Character{91}}%
\TextAt{13}{145.888}{498.898}{M}%
\TextAt{9}{158.462}{498.898}{\Character{93}}%
\TextAt{9}{165.013}{498.898}{\Character{61}}%
\TextAt{9}{177.435}{498.898}{\Character{91}}%
\TextAt{13}{180.711}{498.898}{F}%
\TextAt{9}{189.915}{498.898}{\Character{93}}%
\TextAt{16}{195.823}{498.898}{\Character{0}}%
\TextAt{9}{207.778}{498.898}{\Character{91}}%
\TextAt{13}{211.03}{498.898}{E}%
\TextAt{9}{220.396}{498.898}{\Character{93}}%
\TextAt{8}{223.648}{498.898}{.}%
\Vector{q 0 G 0 g 0.399 w 452.334 -275.102 m 452.334 -267.032 l S Q}%
\Vector{q 0 G 0 g 0.399 w 452.534 -267.231 m 459.707 -267.231 l S Q}%
\Vector{q 0 G 0 g 0.399 w 452.534 -274.903 m 459.707 -274.903 l S Q}%
\Vector{q 0 G 0 g 0.399 w 459.906 -275.102 m 459.906 -267.032 l S Q}%
\TextAt{8}{97.159}{469.897}{The}%
\TextAt{8}{120.305}{469.897}{preceding}%
\TextAt{8}{172.071}{469.897}{theorem}%
\TextAt{8}{216.689}{469.897}{shows}%
\TextAt{8}{249.399}{469.897}{that}%
\TextAt{8}{273.86}{469.897}{the}%
\TextAt{8}{293.108}{469.897}{Grothendieck}%
\TextAt{8}{364.47}{469.897}{group}%
\TextAt{13}{396.737}{469.897}{K}%
\TextAt{17}{406.696}{468.104}{0}%
\TextAt{9}{411.429}{469.897}{\Character{40}}%
\Vector{q 0 G 0 g 0.478 w 343.981 -296.746 m 349.825 -296.746 l S Q}%
\TextAt{14}{415.981}{469.897}{o}%
\TextAt{9}{421.825}{469.897}{\Character{41}}%
\TextAt{8}{429.367}{469.897}{does}%
\TextAt{8}{454.856}{469.897}{not}%
\TextAt{8}{474.762}{469.897}{just}%
\TextAt{8}{497}{469.897}{accom-}%
\TextAt{8}{79.226}{452.851}{modate}%
\TextAt{8}{122.031}{452.851}{all}%
\TextAt{8}{139.143}{452.851}{projective}%
\TextAt{8}{193.974}{452.851}{metrized}%
\TextAt{14}{242.948}{452.851}{o}%
\TextAt{8}{248.792}{452.851}{-modules,}%
\TextAt{8}{302.916}{452.851}{but}%
\TextAt{8}{325.229}{452.851}{in}%
\TextAt{8}{339.736}{452.851}{fact}%
\TextAt{8}{363.675}{452.851}{all}%
\TextAt{8}{380.788}{452.851}{coherent}%
\TextAt{8}{428.792}{452.851}{metrized}%
\TextAt{14}{477.766}{452.851}{o}%
\TextAt{8}{483.61}{452.851}{-modules.}%
\TextAt{8}{79.226}{435.805}{This}%
\TextAt{8}{106.672}{435.805}{fact}%
\TextAt{8}{130.483}{435.805}{has}%
\TextAt{8}{152.082}{435.805}{fundamental}%
\TextAt{8}{220.45}{435.805}{significance.}%
\TextAt{8}{289.023}{435.805}{For}%
\TextAt{8}{310.718}{435.805}{when}%
\TextAt{8}{342.001}{435.805}{dealing}%
\TextAt{8}{383.051}{435.805}{with}%
\TextAt{8}{410.437}{435.805}{projective}%
\TextAt{8}{465.14}{435.805}{modules,}%
\TextAt{8}{515.207}{435.805}{one}%
\TextAt{8}{79.226}{418.759}{is}%
\TextAt{8}{90.897}{418.759}{led}%
\TextAt{8}{109.657}{418.759}{very}%
\TextAt{8}{135.256}{418.759}{quickly}%
\TextAt{8}{175.487}{418.759}{to}%
\TextAt{8}{189.705}{418.759}{non-projective}%
\TextAt{8}{266.352}{418.759}{modules,}%
\TextAt{8}{315.437}{418.759}{for}%
\TextAt{8}{333.229}{418.759}{instance,}%
\TextAt{8}{381.991}{418.759}{to}%
\TextAt{8}{396.209}{418.759}{the}%
\TextAt{8}{416.272}{418.759}{residue}%
\TextAt{8}{455.906}{418.759}{class}%
\TextAt{8}{483.25}{418.759}{rings}%
\TextAt{14}{511.837}{418.759}{o}%
\TextAt{9}{517.682}{418.759}{\Character{47}}%
\TextAt{14}{523.535}{418.759}{a}%
\TextAt{8}{529.518}{418.759}{.}%
\TextAt{8}{79.226}{401.713}{The}%
\TextAt{8}{102.501}{401.713}{corresponding}%
\TextAt{8}{176.908}{401.713}{classes}%
\TextAt{8}{213.381}{401.713}{in}%
\TextAt{13}{226.255}{401.713}{K}%
\TextAt{17}{237.072}{406.052}{0}%
\TextAt{9}{241.805}{401.713}{\Character{40}}%
\Vector{q 0 G 0 g 0.478 w 174.357 -364.93 m 180.201 -364.93 l S Q}%
\TextAt{14}{246.357}{401.713}{o}%
\TextAt{9}{252.201}{401.713}{\Character{41}}%
\TextAt{8}{256.754}{401.713}{,}%
\TextAt{8}{263.28}{401.713}{however,}%
\TextAt{8}{310.788}{401.713}{can}%
\TextAt{8}{331.468}{401.713}{act}%
\TextAt{8}{350.2}{401.713}{out}%
\TextAt{8}{370.235}{401.713}{their}%
\TextAt{8}{397.419}{401.713}{important}%
\TextAt{8}{451.933}{401.713}{rôles}%
\TextAt{8}{478.531}{401.713}{only}%
\TextAt{8}{503.443}{401.713}{inside}%
\TextAt{8}{79.226}{384.667}{the}%
\TextAt{8}{99.385}{384.667}{ring}%
\TextAt{13}{123.453}{384.667}{K}%
\TextAt{17}{133.412}{382.874}{0}%
\TextAt{9}{138.144}{384.667}{\Character{40}}%
\Vector{q 0 G 0 g 0.478 w 70.697 -381.976 m 76.541 -381.976 l S Q}%
\TextAt{14}{142.697}{384.667}{o}%
\TextAt{9}{148.541}{384.667}{\Character{41}}%
\TextAt{8}{153.093}{384.667}{,}%
\TextAt{8}{160.245}{384.667}{because}%
\TextAt{8}{203.549}{384.667}{only}%
\TextAt{8}{229.243}{384.667}{this}%
\TextAt{8}{252.068}{384.667}{ring}%
\TextAt{8}{276.136}{384.667}{can}%
\TextAt{8}{297.598}{384.667}{be}%
\TextAt{8}{313.524}{384.667}{immediately}%
\TextAt{8}{380.188}{384.667}{subjected}%
\TextAt{8}{432.267}{384.667}{to}%
\TextAt{8}{446.58}{384.667}{a}%
\TextAt{8}{456.338}{384.667}{more}%
\TextAt{8}{485.606}{384.667}{advanced}%
\TextAt{8}{79.226}{367.621}{theory.}%
\TextAt{8}{97.159}{347.586}{The}%
\TextAt{8}{120.766}{347.586}{following}%
\TextAt{8}{169.753}{347.586}{relationship}%
\TextAt{8}{233.111}{347.586}{holds}%
\TextAt{8}{263.293}{347.586}{between}%
\TextAt{8}{308.037}{347.586}{the}%
\TextAt{8}{327.746}{347.586}{Grothendieck}%
\TextAt{8}{399.568}{347.586}{ring}%
\TextAt{13}{423.187}{347.586}{K}%
\TextAt{17}{433.146}{345.793}{0}%
\TextAt{9}{437.878}{347.586}{\Character{40}}%
\Vector{q 0 G 0 g 0.478 w 370.431 -419.057 m 376.275 -419.057 l S Q}%
\TextAt{14}{442.431}{347.586}{o}%
\TextAt{9}{448.275}{347.586}{\Character{41}}%
\TextAt{8}{456.278}{347.586}{and}%
\TextAt{8}{478.593}{347.586}{the}%
\TextAt{8}{498.303}{347.586}{replete}%
\TextAt{8}{79.226}{330.54}{Picard}%
\TextAt{8}{116.455}{330.54}{group}%
\TextAt{8}{149.63}{330.54}{Pic}%
\TextAt{9}{166.045}{330.541}{\Character{40}}%
\Vector{q 0 G 0 g 0.478 w 98.597 -436.103 m 104.441 -436.103 l S Q}%
\TextAt{14}{170.597}{330.54}{o}%
\TextAt{9}{176.441}{330.541}{\Character{41}}%
\TextAt{8}{180.993}{330.54}{,}%
\TextAt{8}{188.143}{330.54}{which}%
\TextAt{8}{221.629}{330.54}{was}%
\TextAt{8}{244.129}{330.54}{introduced}%
\TextAt{8}{302.661}{330.54}{in}%
\TextAt{8}{316.314}{330.54}{§1.}%
\TextAt{12}{79.226}{304.528}{Proposition}%
\TextAt{12}{153.011}{304.528}{(5.5).}%
\TextAt{8}{189.26}{304.528}{Associating}%
\TextAt{8}{252.755}{304.528}{to}%
\TextAt{8}{267.998}{304.528}{a}%
\TextAt{8}{278.687}{304.528}{replete}%
\TextAt{8}{317.985}{304.528}{ideal}%
\TextAt{14}{346.881}{304.528}{a}%
\TextAt{8}{357.695}{304.528}{of}%
\TextAt{13}{371.958}{304.528}{K}%
\TextAt{8}{387.607}{304.528}{the}%
\TextAt{8}{408.696}{304.528}{metrized}%
\TextAt{14}{457.749}{304.528}{o}%
\TextAt{8}{463.593}{304.528}{-module}%
\TextAt{13}{509.717}{304.528}{L}%
\TextAt{9}{517.682}{304.528}{\Character{40}}%
\TextAt{14}{522.234}{304.528}{a}%
\TextAt{9}{528.217}{304.528}{\Character{41}}%
\TextAt{8}{79.226}{287.482}{yields}%
\TextAt{8}{112.127}{287.482}{a}%
\TextAt{8}{121.882}{287.482}{homomorphism}%
\TextAt{8}{210.003}{258.481}{Pic}%
\TextAt{9}{226.417}{258.481}{\Character{40}}%
\Vector{q 0 G 0 g 0.478 w 158.97 -508.163 m 164.814 -508.163 l S Q}%
\TextAt{14}{230.97}{258.481}{o}%
\TextAt{9}{236.814}{258.481}{\Character{41}}%
\NativeArrowAt{244.687}{258.481}{11.955}{19.26}{\longrightarrow}%
\TextAt{13}{267.269}{258.481}{K}%
\TextAt{17}{277.228}{256.688}{0}%
\TextAt{9}{281.961}{258.481}{\Character{40}}%
\Vector{q 0 G 0 g 0.478 w 214.513 -508.163 m 220.357 -508.163 l S Q}%
\TextAt{14}{286.513}{258.481}{o}%
\TextAt{9}{292.357}{258.481}{\Character{41}}%
\TextAt{3}{296.909}{263.417}{\Character{3}}%
\TextAt{13}{301.642}{258.481}{\Character{59}}%
\TextAt{9}{330.796}{258.481}{\Character{91}}%
\TextAt{14}{334.048}{258.481}{a}%
\TextAt{9}{340.031}{258.481}{\Character{93}}%
\NativeArrowAt{346.603}{258.481}{11.955}{19.26}{\longmapsto}%
\TextAt{9}{369.186}{258.481}{\Character{91}}%
\TextAt{13}{372.437}{258.481}{L}%
\TextAt{9}{380.402}{258.481}{\Character{40}}%
\TextAt{14}{384.954}{258.481}{a}%
\TextAt{9}{390.937}{258.481}{\Character{41}\Character{93}}%
\TextAt{13}{398.741}{258.481}{\Character{59}}%
\TextAt{8}{79.226}{229.48}{into}%
\TextAt{8}{102.969}{229.48}{the}%
\TextAt{8}{123.126}{229.48}{unit}%
\TextAt{8}{147.837}{229.48}{group}%
\TextAt{8}{181.012}{229.48}{of}%
\TextAt{8}{194.342}{229.48}{the}%
\TextAt{8}{214.499}{229.48}{ring}%
\TextAt{13}{238.565}{229.48}{K}%
\TextAt{17}{248.524}{227.687}{0}%
\TextAt{9}{253.256}{229.48}{\Character{40}}%
\Vector{q 0 G 0 g 0.478 w 185.808 -537.164 m 191.653 -537.164 l S Q}%
\TextAt{14}{257.808}{229.48}{o}%
\TextAt{9}{263.653}{229.48}{\Character{41}}%
\TextAt{8}{268.205}{229.48}{.}%
\TextAt{19}{79.226}{194.501}{Proof.}%
\TextAt{8}{116}{194.501}{The}%
\TextAt{8}{141.069}{194.501}{correspondence}%
\TextAt{9}{223.761}{194.501}{\Character{91}}%
\TextAt{14}{227.013}{194.501}{a}%
\TextAt{9}{232.996}{194.501}{\Character{93}}%
\NativeArrowAt{241.297}{194.501}{11.955}{11.955}{\mapsto}%
\TextAt{9}{258.302}{194.501}{\Character{91}}%
\TextAt{13}{261.554}{194.501}{L}%
\TextAt{9}{269.519}{194.501}{\Character{40}}%
\TextAt{14}{274.071}{194.501}{a}%
\TextAt{9}{280.054}{194.501}{\Character{41}\Character{93}}%
\TextAt{8}{292.77}{194.501}{is}%
\TextAt{8}{305.549}{194.501}{independent}%
\TextAt{8}{372.89}{194.501}{of}%
\TextAt{8}{387.235}{194.501}{the}%
\TextAt{8}{408.406}{194.501}{choice}%
\TextAt{8}{444.21}{194.501}{of}%
\TextAt{8}{458.555}{194.501}{a}%
\TextAt{8}{469.325}{194.501}{replete}%
\TextAt{8}{508.704}{194.501}{ideal}%
\TextAt{14}{79.226}{177.455}{a}%
\TextAt{8}{90.893}{177.455}{inside}%
\TextAt{8}{125.903}{177.455}{the}%
\TextAt{8}{147.846}{177.455}{class}%
\TextAt{9}{177.07}{177.455}{\Character{91}}%
\TextAt{14}{180.322}{177.455}{a}%
\TextAt{9}{186.305}{177.455}{\Character{93}}%
\NativeMathAt{195.922}{177.455}{11.955}{\in}%
\TextAt{8}{210.257}{177.455}{Pic}%
\TextAt{9}{226.672}{177.455}{\Character{40}}%
\Vector{q 0 G 0 g 0.478 w 159.224 -589.188 m 165.068 -589.188 l S Q}%
\TextAt{14}{231.224}{177.455}{o}%
\TextAt{9}{237.068}{177.455}{\Character{41}}%
\TextAt{8}{241.621}{177.455}{.}%
\TextAt{8}{255.429}{177.455}{Indeed,}%
\TextAt{8}{298.943}{177.455}{if}%
\TextAt{14}{311.454}{177.455}{b}%
\TextAt{8}{323.272}{177.455}{is}%
\TextAt{8}{336.823}{177.455}{another}%
\TextAt{8}{381.54}{177.455}{representative,}%
\TextAt{8}{462.212}{177.455}{then}%
\TextAt{8}{490.658}{177.455}{we}%
\TextAt{8}{509.672}{177.455}{have}%
\TextAt{14}{79.226}{160.409}{a}%
\TextAt{9}{88.95}{160.409}{\Character{61}}%
\TextAt{14}{101.795}{160.409}{b}%
\TextAt{9}{107.93}{160.409}{\Character{91}}%
\TextAt{13}{111.181}{160.409}{a}%
\TextAt{9}{117.326}{160.409}{\Character{93}}%
\TextAt{8}{120.578}{160.409}{,}%
\TextAt{8}{128.035}{160.409}{for}%
\TextAt{8}{146.166}{160.409}{some}%
\TextAt{8}{175.739}{160.409}{replete}%
\TextAt{8}{214.349}{160.409}{principal}%
\TextAt{8}{263.373}{160.409}{ideal}%
\TextAt{9}{291.582}{160.409}{\Character{91}}%
\TextAt{13}{294.834}{160.409}{a}%
\TextAt{9}{300.979}{160.409}{\Character{93}}%
\TextAt{8}{304.23}{160.409}{,}%
\TextAt{8}{311.687}{160.409}{and}%
\TextAt{8}{334.696}{160.409}{the}%
\TextAt{8}{355.099}{160.409}{metrized}%
\TextAt{14}{403.465}{160.409}{o}%
\TextAt{8}{409.309}{160.409}{-modules}%
\TextAt{13}{459.361}{160.409}{L}%
\TextAt{9}{467.325}{160.409}{\Character{40}}%
\TextAt{14}{471.877}{160.409}{a}%
\TextAt{9}{477.861}{160.409}{\Character{41}}%
\TextAt{8}{486.557}{160.409}{and}%
\TextAt{13}{509.566}{160.409}{L}%
\TextAt{9}{517.53}{160.409}{\Character{40}}%
\TextAt{14}{522.082}{160.409}{b}%
\TextAt{9}{528.217}{160.409}{\Character{41}}%
\TextAt{8}{79.226}{143.363}{are}%
\TextAt{8}{98.419}{143.363}{isometric}%
\TextAt{8}{148.241}{143.363}{by}%
\TextAt{8}{164.182}{143.363}{(4.5),}%
\TextAt{8}{195.154}{143.363}{(i),}%
\TextAt{8}{214.411}{143.363}{so}%
\TextAt{8}{228.463}{143.363}{that}%
\TextAt{9}{253.514}{143.363}{\Character{91}}%
\TextAt{13}{256.765}{143.363}{L}%
\TextAt{9}{264.73}{143.363}{\Character{40}}%
\TextAt{14}{269.282}{143.363}{a}%
\TextAt{9}{275.265}{143.363}{\Character{41}\Character{93}}%
\TextAt{9}{286.359}{143.363}{\Character{61}}%
\TextAt{9}{298.781}{143.363}{\Character{91}}%
\TextAt{13}{302.067}{143.363}{L}%
\TextAt{9}{310.032}{143.363}{\Character{40}}%
\TextAt{14}{314.584}{143.363}{b}%
\TextAt{9}{320.719}{143.363}{\Character{41}\Character{93}}%
\TextAt{8}{328.523}{143.363}{.}%
\TextAt{8}{336.865}{143.363}{The}%
\TextAt{8}{360.601}{143.363}{correspondence}%
\TextAt{8}{441.96}{143.363}{is}%
\TextAt{8}{453.406}{143.363}{a}%
\TextAt{8}{462.844}{143.363}{multiplicative}%
\TextAt{8}{79.226}{126.317}{homomorphism}%
\TextAt{8}{161.896}{126.317}{as}%
\TextAt{9}{205.266}{109.271}{\Character{91}}%
\TextAt{13}{208.518}{109.271}{L}%
\TextAt{9}{216.482}{109.271}{\Character{40}}%
\TextAt{14}{221.035}{109.271}{ab}%
\TextAt{9}{233.153}{109.271}{\Character{41}\Character{93}}%
\TextAt{9}{244.247}{109.271}{\Character{61}}%
\TextAt{9}{256.669}{109.271}{\Character{91}}%
\TextAt{13}{259.955}{109.271}{L}%
\TextAt{9}{267.919}{109.271}{\Character{40}}%
\TextAt{14}{272.471}{109.271}{a}%
\TextAt{9}{278.455}{109.271}{\Character{41}}%
\TextAt{16}{285.664}{109.271}{\Character{10}}%
\TextAt{27}{294.962}{107.478}{o}%
\TextAt{13}{302.013}{109.271}{L}%
\TextAt{9}{309.977}{109.271}{\Character{40}}%
\TextAt{14}{314.53}{109.271}{b}%
\TextAt{9}{320.665}{109.271}{\Character{41}\Character{93}}%
\TextAt{9}{331.759}{109.271}{\Character{61}}%
\TextAt{9}{344.181}{109.271}{\Character{91}}%
\TextAt{13}{347.467}{109.271}{L}%
\TextAt{9}{355.431}{109.271}{\Character{40}}%
\TextAt{14}{359.983}{109.271}{a}%
\TextAt{9}{365.966}{109.271}{\Character{41}\Character{93}\Character{91}}%
\TextAt{13}{377.022}{109.271}{L}%
\TextAt{9}{384.987}{109.271}{\Character{40}}%
\TextAt{14}{389.539}{109.271}{b}%
\TextAt{9}{395.674}{109.271}{\Character{41}\Character{93}}%
\TextAt{13}{403.478}{109.271}{\Character{58}}%
\Vector{q 0 G 0 g 0.399 w 452.334 -688.25 m 452.334 -680.181 l S Q}%
\Vector{q 0 G 0 g 0.399 w 452.534 -680.38 m 459.707 -680.38 l S Q}%
\Vector{q 0 G 0 g 0.399 w 452.534 -688.051 m 459.707 -688.051 l S Q}%
\Vector{q 0 G 0 g 0.399 w 459.906 -688.25 m 459.906 -680.181 l S Q}%
\TextAt{8}{297.211}{43.033}{219}%
\Destination{78.23}{809.815}{nav.page.219}%
\Destination{79.36}{450.855}{nav.print.239}%
\Destination{79.36}{318.555}{nav.number.3.5.5}%
\Hotlink{321.493}{328.355}{6.858}{11.047}{nav.section.3.1}%
\Hotlink{163.682}{141.155}{25.078}{11.047}{nav.number.3.4.5}%
\Destination{79.226}{778.124}{errata.90}%
}

% Page 232 (body page 220)
\Page{%
\TextAt{2}{79.226}{796.943}{III.}%
\TextAt{2}{98.317}{796.943}{Riemann–Roch}%
\TextAt{2}{174.375}{796.943}{Theory}%
\Vector{q 0 G 0 g 0.299 w 7.226 18.729 m 460.769 18.729 l S Q}%
\TextAt{8}{97.159}{761.124}{In}%
\TextAt{8}{112.628}{761.124}{the}%
\TextAt{8}{133.633}{761.124}{sequel,}%
\TextAt{8}{172.794}{761.124}{we}%
\TextAt{8}{190.87}{761.124}{simply}%
\TextAt{8}{229.173}{761.124}{denote}%
\TextAt{8}{267.74}{761.124}{the}%
\TextAt{8}{288.745}{761.124}{class}%
\TextAt{8}{317.03}{761.124}{of}%
\TextAt{8}{331.208}{761.124}{a}%
\TextAt{8}{341.811}{761.124}{metrized}%
\TextAt{8}{390.779}{761.124}{invertible}%
\TextAt{14}{443.333}{761.124}{o}%
\TextAt{8}{449.178}{761.124}{-module}%
\TextAt{13}{495.216}{761.124}{L}%
\TextAt{9}{503.181}{761.124}{\Character{40}}%
\TextAt{14}{507.733}{761.124}{a}%
\TextAt{9}{513.716}{761.124}{\Character{41}}%
\TextAt{8}{523.014}{761.124}{in}%
\TextAt{13}{79.226}{744.078}{K}%
\TextAt{17}{89.185}{742.284}{0}%
\TextAt{9}{93.918}{744.078}{\Character{40}}%
\Vector{q 0 G 0 g 0.478 w 26.47 -22.566 m 32.314 -22.566 l S Q}%
\TextAt{14}{98.47}{744.078}{o}%
\TextAt{9}{104.314}{744.078}{\Character{41}}%
\TextAt{8}{112.545}{744.078}{by}%
\TextAt{9}{128.585}{744.078}{\Character{91}}%
\TextAt{14}{131.837}{744.078}{a}%
\TextAt{9}{137.82}{744.078}{\Character{93}}%
\TextAt{8}{141.071}{744.078}{.}%
\TextAt{8}{149.447}{744.078}{In}%
\TextAt{8}{163.849}{744.078}{particular,}%
\TextAt{8}{220.915}{744.078}{to}%
\TextAt{8}{235.006}{744.078}{the}%
\TextAt{8}{254.944}{744.078}{replete}%
\TextAt{8}{293.089}{744.078}{ideal}%
\TextAt{14}{320.833}{744.078}{o}%
\TextAt{9}{329.998}{744.078}{\Character{61}}%
\TextAt{20}{342.424}{753.044}{\Character{81}}%
\TextAt{27}{353.715}{740.591}{p}%
\TextAt{14}{360.194}{744.078}{p}%
\TextAt{17}{366.177}{748.416}{0}%
\TextAt{8}{374.588}{744.078}{corresponds}%
\TextAt{8}{438.556}{744.078}{the}%
\TextAt{8}{458.493}{744.078}{class}%
\TextAt{9}{485.712}{744.078}{1}%
\TextAt{9}{494.87}{744.078}{\Character{61}}%
\TextAt{9}{507.291}{744.078}{\Character{91}}%
\TextAt{14}{510.563}{744.078}{o}%
\TextAt{9}{516.407}{744.078}{\Character{93}}%
\TextAt{8}{523.337}{744.078}{of}%
\TextAt{8}{79.226}{727.032}{the}%
\TextAt{14}{99.383}{727.032}{o}%
\TextAt{8}{105.227}{727.032}{-module}%
\TextAt{14}{150.417}{727.032}{o}%
\TextAt{8}{160.159}{727.032}{equipped}%
\TextAt{8}{210.227}{727.032}{with}%
\TextAt{8}{236.887}{727.032}{the}%
\TextAt{8}{257.044}{727.032}{trivial}%
\TextAt{8}{291.845}{727.032}{metric.}%
\TextAt{12}{79.226}{701.019}{Proposition}%
\TextAt{12}{151.937}{701.019}{(5.6).}%
\TextAt{13}{187.253}{701.019}{K}%
\TextAt{17}{197.212}{699.226}{0}%
\TextAt{9}{201.944}{701.019}{\Character{40}}%
\Vector{q 0 G 0 g 0.478 w 134.497 -65.624 m 140.341 -65.624 l S Q}%
\TextAt{14}{206.497}{701.019}{o}%
\TextAt{9}{212.341}{701.019}{\Character{41}}%
\TextAt{8}{220.791}{701.019}{is}%
\TextAt{8}{232.555}{701.019}{generated}%
\TextAt{8}{285.887}{701.019}{as}%
\TextAt{8}{300.257}{701.019}{an}%
\TextAt{8}{316.516}{701.019}{additive}%
\TextAt{8}{361.395}{701.019}{group}%
\TextAt{8}{394.571}{701.019}{by}%
\TextAt{8}{410.83}{701.019}{the}%
\TextAt{8}{430.987}{701.019}{elements}%
\TextAt{9}{478.843}{701.019}{\Character{91}}%
\TextAt{14}{482.095}{701.019}{a}%
\TextAt{9}{488.078}{701.019}{\Character{93}}%
\TextAt{8}{491.329}{701.019}{.}%
\TextAt{19}{79.226}{666.041}{Proof.}%
\TextAt{8}{116}{666.041}{Let}%
\TextAt{13}{138.545}{666.041}{M}%
\TextAt{8}{156.591}{666.041}{be}%
\TextAt{8}{174.091}{666.041}{a}%
\TextAt{8}{185.422}{666.041}{projective}%
\TextAt{8}{240.974}{666.041}{metrized}%
\TextAt{14}{290.669}{666.041}{o}%
\TextAt{8}{296.513}{666.041}{-module.}%
\TextAt{8}{350.981}{666.041}{By}%
\TextAt{8}{370.919}{666.041}{(4.3),}%
\TextAt{8}{404.115}{666.041}{the}%
\TextAt{8}{425.847}{666.041}{underlying}%
\TextAt{14}{485.632}{666.041}{o}%
\TextAt{8}{491.476}{666.041}{-module}%
\TextAt{8}{79.226}{648.995}{admits}%
\TextAt{8}{117.662}{648.995}{as}%
\TextAt{8}{132.032}{648.995}{quotient}%
\TextAt{8}{178.215}{648.995}{a}%
\TextAt{8}{187.97}{648.995}{fractional}%
\TextAt{8}{240.334}{648.995}{ideal}%
\TextAt{14}{268.297}{648.995}{a}%
\TextAt{21}{274.28}{647.201}{f}%
\TextAt{8}{279.725}{648.995}{,}%
\TextAt{8}{286.874}{648.995}{i.e.,}%
\TextAt{8}{308.979}{648.995}{we}%
\TextAt{8}{326.207}{648.995}{have}%
\TextAt{8}{353.202}{648.995}{an}%
\TextAt{8}{369.461}{648.995}{exact}%
\TextAt{8}{400.353}{648.995}{sequence}%
\TextAt{9}{231.027}{619.994}{0}%
\NativeArrowAt{240.201}{619.994}{11.955}{19.26}{\longrightarrow}%
\TextAt{13}{262.783}{619.994}{N}%
\NativeArrowAt{276.726}{619.994}{11.955}{19.26}{\longrightarrow}%
\TextAt{13}{299.308}{619.994}{M}%
\NativeArrowAt{315.203}{619.994}{11.955}{19.26}{\longrightarrow}%
\TextAt{14}{337.785}{619.994}{a}%
\TextAt{21}{343.768}{618.2}{f}%
\NativeArrowAt{352.534}{619.994}{11.955}{19.26}{\longrightarrow}%
\TextAt{9}{375.116}{619.994}{0}%
\TextAt{8}{79.226}{590.992}{of}%
\TextAt{14}{92.026}{590.992}{o}%
\TextAt{8}{97.87}{590.992}{-modules.}%
\TextAt{8}{152.049}{590.992}{This}%
\TextAt{8}{178.238}{590.992}{becomes}%
\TextAt{8}{224.261}{590.992}{an}%
\TextAt{8}{239.99}{590.992}{exact}%
\TextAt{8}{270.351}{590.992}{sequence}%
\TextAt{8}{318.323}{590.992}{of}%
\TextAt{8}{331.122}{590.992}{metrized}%
\TextAt{14}{378.711}{590.992}{o}%
\TextAt{8}{384.555}{590.992}{-modules}%
\TextAt{8}{433.83}{590.992}{once}%
\TextAt{8}{459.96}{590.992}{we}%
\TextAt{8}{476.656}{590.992}{restrict}%
\TextAt{8}{516.51}{590.992}{the}%
\TextAt{8}{79.226}{573.946}{metric}%
\TextAt{8}{114.857}{573.946}{from}%
\TextAt{13}{141.713}{573.946}{M}%
\TextAt{8}{157.399}{573.946}{to}%
\TextAt{13}{170.925}{573.946}{N}%
\TextAt{8}{184.66}{573.946}{and}%
\TextAt{8}{206.638}{573.946}{choose}%
\TextAt{8}{242.986}{573.946}{on}%
\TextAt{14}{258.46}{573.946}{a}%
\TextAt{21}{264.443}{572.153}{f}%
\TextAt{8}{273.001}{573.946}{the}%
\TextAt{8}{292.373}{573.946}{metric}%
\TextAt{8}{328.003}{573.946}{which}%
\TextAt{8}{360.705}{573.946}{is}%
\TextAt{8}{371.684}{573.946}{transferred}%
\TextAt{8}{430.472}{573.946}{via}%
\TextAt{8}{448.876}{573.946}{the}%
\TextAt{8}{468.247}{573.946}{isomorphism}%
\TextAt{13}{79.226}{556.9}{N}%
\TextAt{3}{89.849}{561.239}{\Character{63}}%
\TextAt{25}{88.6}{553.836}{C}%
\NativeMathAt{101.537}{556.9}{11.955}{\cong}%
\TextAt{14}{115.44}{556.9}{a}%
\TextAt{21}{121.423}{555.107}{f}%
\TextAt{25}{126.37}{555.107}{C}%
\TextAt{8}{132.624}{556.9}{.}%
\TextAt{8}{143.332}{556.9}{Thus}%
\TextAt{14}{173.734}{556.9}{a}%
\TextAt{21}{179.717}{555.107}{f}%
\TextAt{8}{189.812}{556.9}{becomes}%
\TextAt{8}{237.118}{556.9}{the}%
\TextAt{8}{258.028}{556.9}{metrized}%
\TextAt{14}{306.9}{556.9}{o}%
\TextAt{8}{312.745}{556.9}{-module}%
\TextAt{13}{358.688}{556.9}{L}%
\TextAt{9}{366.653}{556.9}{\Character{40}}%
\TextAt{14}{371.205}{556.9}{a}%
\TextAt{9}{377.188}{556.9}{\Character{41}}%
\TextAt{8}{386.391}{556.9}{corresponding}%
\TextAt{8}{462.33}{556.9}{to}%
\TextAt{8}{477.393}{556.9}{the}%
\TextAt{8}{498.303}{556.9}{replete}%
\TextAt{8}{79.226}{539.855}{ideal}%
\TextAt{14}{107.06}{539.855}{a}%
\TextAt{8}{116.811}{539.855}{of}%
\TextAt{13}{130.012}{539.855}{K}%
\TextAt{8}{140.829}{539.855}{,}%
\TextAt{8}{147.875}{539.855}{so}%
\TextAt{8}{162.116}{539.855}{that}%
\TextAt{8}{187.355}{539.855}{we}%
\TextAt{8}{204.453}{539.855}{get}%
\TextAt{8}{223.835}{539.855}{the}%
\TextAt{8}{243.862}{539.855}{identity}%
\TextAt{9}{286.986}{539.855}{\Character{91}}%
\TextAt{13}{290.238}{539.855}{M}%
\TextAt{9}{302.811}{539.855}{\Character{93}}%
\TextAt{9}{309.362}{539.855}{\Character{61}}%
\TextAt{9}{321.784}{539.855}{\Character{91}}%
\TextAt{13}{325.061}{539.855}{N}%
\TextAt{9}{335.684}{539.855}{\Character{93}}%
\TextAt{9}{341.315}{539.855}{\Character{43}}%
\TextAt{9}{352.804}{539.855}{\Character{91}}%
\TextAt{14}{356.076}{539.855}{a}%
\TextAt{9}{362.059}{539.855}{\Character{93}}%
\TextAt{8}{369.079}{539.855}{in}%
\TextAt{13}{382.602}{539.855}{K}%
\TextAt{17}{392.561}{538.061}{0}%
\TextAt{9}{397.293}{539.855}{\Character{40}}%
\Vector{q 0 G 0 g 0.478 w 329.846 -226.789 m 335.69 -226.789 l S Q}%
\TextAt{14}{401.846}{539.855}{o}%
\TextAt{9}{407.69}{539.855}{\Character{41}}%
\TextAt{8}{412.242}{539.855}{.}%
\TextAt{8}{420.648}{539.855}{Induction}%
\TextAt{8}{473.515}{539.855}{on}%
\TextAt{8}{489.645}{539.855}{the}%
\TextAt{8}{509.672}{539.855}{rank}%
\TextAt{8}{79.226}{522.809}{shows}%
\TextAt{8}{112.844}{522.809}{that}%
\TextAt{8}{138.213}{522.809}{for}%
\TextAt{8}{156.098}{522.809}{every}%
\TextAt{8}{186.99}{522.809}{projective}%
\TextAt{8}{240.968}{522.809}{metrized}%
\TextAt{14}{289.087}{522.809}{o}%
\TextAt{8}{294.931}{522.809}{-module}%
\TextAt{13}{340.122}{522.809}{M}%
\TextAt{8}{352.696}{522.809}{,}%
\TextAt{8}{359.845}{522.809}{there}%
\TextAt{8}{389.757}{522.809}{is}%
\TextAt{8}{401.52}{522.809}{a}%
\TextAt{8}{411.276}{522.809}{decomposition}%
\TextAt{9}{248.439}{493.807}{\Character{91}}%
\TextAt{13}{251.69}{493.807}{M}%
\TextAt{9}{264.264}{493.807}{\Character{93}}%
\TextAt{9}{270.815}{493.807}{\Character{61}}%
\TextAt{9}{283.237}{493.807}{\Character{91}}%
\TextAt{14}{286.514}{493.807}{a}%
\TextAt{17}{292.497}{492.014}{1}%
\TextAt{9}{297.229}{493.807}{\Character{93}}%
\TextAt{9}{303.123}{493.807}{\Character{43}}%
\TextAt{16}{314.899}{493.807}{\Character{1}}%
\TextAt{16}{320.195}{493.807}{\Character{1}}%
\TextAt{16}{325.503}{493.807}{\Character{1}}%
\TextAt{9}{331.503}{493.807}{\Character{43}}%
\TextAt{9}{343.255}{493.807}{\Character{91}}%
\TextAt{14}{346.516}{493.807}{a}%
\TextAt{21}{352.499}{492.014}{r}%
\TextAt{9}{357.054}{493.807}{\Character{93}}%
\TextAt{13}{360.305}{493.807}{\Character{58}}%
\Vector{q 0 G 0 g 0.399 w 452.334 -309.194 m 452.334 -301.124 l S Q}%
\Vector{q 0 G 0 g 0.399 w 452.534 -301.323 m 459.707 -301.323 l S Q}%
\Vector{q 0 G 0 g 0.399 w 452.534 -308.994 m 459.707 -308.994 l S Q}%
\Vector{q 0 G 0 g 0.399 w 459.906 -309.194 m 459.906 -301.124 l S Q}%
\TextAt{8}{97.159}{435.805}{The}%
\TextAt{8}{121.213}{435.805}{elements}%
\TextAt{9}{169.069}{435.805}{\Character{91}}%
\TextAt{14}{172.321}{435.805}{a}%
\TextAt{9}{178.304}{435.805}{\Character{93}}%
\TextAt{8}{185.453}{435.805}{in}%
\TextAt{13}{199.106}{435.805}{K}%
\TextAt{17}{209.065}{434.012}{0}%
\TextAt{9}{213.797}{435.805}{\Character{40}}%
\Vector{q 0 G 0 g 0.478 w 146.35 -330.838 m 152.194 -330.838 l S Q}%
\TextAt{14}{218.35}{435.805}{o}%
\TextAt{9}{224.194}{435.805}{\Character{41}}%
\TextAt{8}{232.644}{435.805}{satisfy}%
\TextAt{8}{269.191}{435.805}{the}%
\TextAt{8}{289.347}{435.805}{following}%
\TextAt{8}{338.782}{435.805}{remarkable}%
\TextAt{8}{399.598}{435.805}{relation.}%
\TextAt{12}{79.226}{409.793}{Proposition}%
\TextAt{12}{150.809}{409.793}{(5.7).}%
\TextAt{8}{185.143}{409.793}{For}%
\TextAt{8}{205.131}{409.793}{any}%
\TextAt{8}{226.267}{409.793}{two}%
\TextAt{8}{247.402}{409.793}{replete}%
\TextAt{8}{284.785}{409.793}{ideals}%
\TextAt{14}{316.382}{409.793}{a}%
\TextAt{8}{325.281}{409.793}{and}%
\TextAt{14}{347.062}{409.793}{b}%
\TextAt{8}{356.113}{409.793}{of}%
\TextAt{13}{368.462}{409.793}{K}%
\TextAt{8}{382.196}{409.793}{we}%
\TextAt{8}{398.442}{409.793}{have}%
\TextAt{8}{424.455}{409.793}{in}%
\TextAt{13}{437.127}{409.793}{K}%
\TextAt{17}{447.086}{408}{0}%
\TextAt{9}{451.818}{409.793}{\Character{40}}%
\Vector{q 0 G 0 g 0.478 w 384.37 -356.851 m 390.215 -356.851 l S Q}%
\TextAt{14}{456.37}{409.793}{o}%
\TextAt{9}{462.215}{409.793}{\Character{41}}%
\TextAt{8}{469.683}{409.793}{the}%
\TextAt{8}{488.858}{409.793}{equation}%
\TextAt{9}{251.441}{380.792}{\Character{40}\Character{91}}%
\TextAt{14}{259.245}{380.792}{a}%
\TextAt{9}{265.228}{380.792}{\Character{93}}%
\TextAt{16}{271.136}{380.792}{\Character{0}}%
\TextAt{9}{283.091}{380.792}{1\Character{41}\Character{40}\Character{91}}%
\TextAt{14}{301.3}{380.792}{b}%
\TextAt{9}{307.435}{380.792}{\Character{93}}%
\TextAt{16}{313.344}{380.792}{\Character{0}}%
\TextAt{9}{325.299}{380.792}{1\Character{41}}%
\TextAt{9}{339}{380.792}{\Character{61}}%
\TextAt{9}{351.421}{380.792}{0}%
\TextAt{13}{357.303}{380.792}{\Character{58}}%
\TextAt{19}{79.226}{333.858}{Proof}%
\TextAt{19}{110.62}{333.858}{(Tamme).}%
\TextAt{8}{167.706}{333.858}{For}%
\TextAt{8}{189.876}{333.858}{every}%
\TextAt{8}{221.969}{333.858}{function}%
\TextAt{13}{269.018}{333.858}{\Character{11}}%
\TextAt{9}{281.906}{333.858}{\Character{58}}%
\TextAt{13}{290.524}{333.858}{X}%
\TextAt{9}{301.18}{333.858}{\Character{40}}%
\TextAt{18}{305.732}{333.858}{C}%
\TextAt{9}{314.366}{333.858}{\Character{41}}%
\NativeMathAt{324.285}{333.858}{11.955}{\rightarrow}%
\TextAt{18}{341.607}{333.858}{C}%
\TextAt{8}{355.339}{333.858}{let}%
\TextAt{8}{373.444}{333.858}{us}%
\TextAt{8}{389.661}{333.858}{consider}%
\TextAt{8}{436.446}{333.858}{on}%
\TextAt{8}{453.906}{333.858}{the}%
\TextAt{13}{475.263}{333.858}{K}%
\TextAt{25}{485.222}{332.064}{C}%
\TextAt{8}{491.476}{333.858}{-module}%
\TextAt{13}{79.226}{316.812}{K}%
\TextAt{25}{89.185}{315.019}{C}%
\TextAt{9}{98.76}{316.812}{\Character{61}}%
\TextAt{20}{111.186}{325.778}{\Character{76}}%
\TextAt{21}{124.469}{313.325}{\Character{27}}%
\NativeMathAt{129.595}{313.325}{7.97}{\in}%
\TextAt{21}{135.24}{313.325}{X}%
\TextAt{17}{142.827}{313.325}{\Character{40}}%
\TextAt{25}{146.121}{313.325}{C}%
\TextAt{17}{151.877}{313.325}{\Character{41}}%
\TextAt{18}{157.661}{316.812}{C}%
\TextAt{8}{170.192}{316.812}{the}%
\TextAt{8}{190.349}{316.812}{form}%
\TextAt{13}{252.442}{297.783}{\Character{11}x}%
\Vector{q 0 G 0 g 0.478 w 194.616 -469.396 m 200.753 -469.396 l S Q}%
\TextAt{13}{266.616}{297.783}{y}%
\TextAt{9}{276.074}{297.783}{\Character{61}}%
\TextAt{20}{288.499}{309.14}{\Character{77}}%
\TextAt{21}{294.969}{284.034}{\Character{27}}%
\TextAt{13}{308.557}{297.783}{\Character{11}}%
\TextAt{9}{316.079}{297.783}{\Character{40}}%
\TextAt{13}{320.631}{297.783}{\Character{27}}%
\TextAt{9}{327.714}{297.783}{\Character{41}}%
\TextAt{13}{332.266}{297.783}{x}%
\TextAt{21}{338.918}{295.99}{\Character{27}}%
\Vector{q 0 G 0 g 0.478 w 272.542 -469.396 m 278.678 -469.396 l S Q}%
\TextAt{13}{344.542}{297.783}{y}%
\TextAt{21}{350.678}{294.96}{\Character{27}}%
\TextAt{13}{356.302}{297.783}{\Character{58}}%
\TextAt{8}{79.226}{254.445}{For}%
\TextAt{8}{99.592}{254.445}{every}%
\TextAt{8}{129.881}{254.445}{matrix}%
\TextAt{13}{167.331}{254.445}{A}%
\TextAt{9}{179.427}{254.445}{\Character{61}}%
\TextAt{20}{191.853}{274.889}{\Character{32}}%
\TextAt{13}{201.317}{262.548}{\Character{11}}%
\TextAt{13}{219.059}{262.548}{\Character{13}}%
\TextAt{13}{202.261}{245.502}{\Character{14}}%
\TextAt{13}{218.783}{245.502}{\Character{12}}%
\TextAt{20}{226.073}{274.889}{\Character{33}}%
\TextAt{8}{238.831}{254.445}{of}%
\TextAt{8}{251.558}{254.445}{such}%
\TextAt{8}{277.352}{254.445}{functions,}%
\TextAt{8}{330.583}{254.445}{we}%
\TextAt{8}{347.207}{254.445}{consider}%
\TextAt{8}{392.189}{254.445}{on}%
\TextAt{8}{407.845}{254.445}{the}%
\TextAt{13}{427.398}{254.445}{K}%
\TextAt{25}{437.357}{252.652}{C}%
\TextAt{8}{443.611}{254.445}{-module}%
\TextAt{13}{488.199}{254.445}{K}%
\TextAt{25}{498.158}{252.652}{C}%
\TextAt{16}{505.835}{254.445}{\Character{8}}%
\TextAt{13}{516.556}{254.445}{K}%
\TextAt{25}{526.515}{252.652}{C}%
\TextAt{8}{79.226}{230.088}{the}%
\TextAt{8}{99.383}{230.088}{form}%
\TextAt{16}{189.243}{213.042}{\Character{104}}%
\TextAt{13}{193.892}{213.042}{x}%
\TextAt{16}{203.201}{213.042}{\Character{8}}%
\TextAt{13}{215.156}{213.042}{y}%
\TextAt{13}{221.277}{213.042}{\Character{59}}%
\TextAt{13}{226.513}{213.042}{x}%
\TextAt{3}{233.189}{217.979}{\Character{48}}%
\TextAt{16}{238.641}{213.042}{\Character{8}}%
\TextAt{13}{250.596}{213.042}{y}%
\TextAt{3}{256.732}{217.979}{\Character{48}}%
\TextAt{16}{259.527}{213.042}{\Character{105}}%
\TextAt{21}{264.177}{211.249}{A}%
\TextAt{9}{274.339}{213.042}{\Character{61}}%
\TextAt{13}{286.764}{213.042}{\Character{11}x}%
\Vector{q 0 G 0 g 0.478 w 228.938 -551.402 m 238.385 -551.402 l S Q}%
\TextAt{13}{300.938}{213.042}{x}%
\TextAt{3}{307.59}{216.496}{\Character{48}}%
\TextAt{9}{313.042}{213.042}{\Character{43}}%
\TextAt{13}{324.803}{213.042}{\Character{13}}%
\TextAt{13}{331.51}{213.042}{x}%
\Vector{q 0 G 0 g 0.478 w 266.177 -551.402 m 275.109 -551.402 l S Q}%
\TextAt{13}{338.177}{213.042}{y}%
\TextAt{3}{344.314}{216.496}{\Character{48}}%
\TextAt{9}{349.766}{213.042}{\Character{43}}%
\TextAt{13}{361.527}{213.042}{\Character{14}}%
\TextAt{13}{367.146}{213.042}{y}%
\Vector{q 0 G 0 g 0.478 w 301.297 -551.402 m 310.744 -551.402 l S Q}%
\TextAt{13}{373.297}{213.042}{x}%
\TextAt{3}{379.949}{216.496}{\Character{48}}%
\TextAt{9}{385.4}{213.042}{\Character{43}}%
\TextAt{13}{397.162}{213.042}{\Character{12}}%
\TextAt{13}{404.431}{213.042}{y}%
\Vector{q 0 G 0 g 0.478 w 338.57 -551.402 m 347.501 -551.402 l S Q}%
\TextAt{13}{410.57}{213.042}{y}%
\TextAt{3}{416.706}{216.496}{\Character{48}}%
\TextAt{13}{419.501}{213.042}{\Character{58}}%
\TextAt{13}{79.226}{189.521}{\Character{11}x}%
\Vector{q 0 G 0 g 0.478 w 21.4 -577.658 m 27.537 -577.658 l S Q}%
\TextAt{13}{93.4}{189.521}{y}%
\TextAt{8}{99.537}{189.521}{,}%
\TextAt{8}{107.556}{189.521}{resp.}%
\TextAt{16}{138.966}{189.521}{\Character{104}}%
\TextAt{13}{148.208}{189.521}{\Character{59}}%
\TextAt{16}{158.046}{189.521}{\Character{105}}%
\TextAt{21}{162.695}{187.727}{A}%
\TextAt{8}{169.536}{189.521}{,}%
\TextAt{8}{177.555}{189.521}{is}%
\TextAt{8}{190.015}{189.521}{an}%
\TextAt{13}{206.97}{189.521}{F}%
\TextAt{8}{216.174}{189.521}{-invariant}%
\TextAt{8}{269.891}{189.521}{metric}%
\TextAt{8}{307.003}{189.521}{on}%
\TextAt{13}{323.958}{189.521}{K}%
\TextAt{25}{333.917}{187.727}{C}%
\TextAt{8}{340.171}{189.521}{,}%
\TextAt{8}{348.19}{189.521}{resp.}%
\TextAt{8}{379.6}{189.521}{on}%
\TextAt{13}{396.555}{189.521}{K}%
\TextAt{25}{406.514}{187.727}{C}%
\TextAt{16}{415.9}{189.521}{\Character{8}}%
\TextAt{13}{428.329}{189.521}{K}%
\TextAt{25}{438.288}{187.727}{C}%
\TextAt{8}{444.543}{189.521}{,}%
\TextAt{8}{452.562}{189.521}{if}%
\TextAt{8}{463.982}{189.521}{and}%
\TextAt{8}{487.44}{189.521}{only}%
\TextAt{8}{513.828}{189.521}{if}%
\TextAt{13}{525.248}{189.521}{\Character{11}}%
\TextAt{8}{79.226}{172.475}{is}%
\TextAt{13}{91.366}{172.475}{F}%
\TextAt{8}{100.569}{172.475}{-invariant}%
\TextAt{8}{153.966}{172.475}{(i.e.,}%
\TextAt{13}{181.096}{172.475}{\Character{11}}%
\TextAt{9}{188.618}{172.475}{\Character{40}}%
\TextAt{13}{193.17}{172.475}{\Character{27}}%
\TextAt{9}{200.253}{172.475}{\Character{41}}%
\TextAt{9}{208.753}{172.475}{\Character{61}}%
\Vector{q 0 G 0 g 0.478 w 149.832 -590.885 m 173.541 -590.885 l S Q}%
\TextAt{13}{221.832}{172.475}{\Character{11}}%
\TextAt{9}{229.354}{172.475}{\Character{40}}%
\Vector{q 0 G 0 g 0.478 w 161.906 -594.704 m 168.989 -594.704 l S Q}%
\TextAt{13}{233.906}{172.475}{\Character{27}}%
\TextAt{9}{240.989}{172.475}{\Character{41}}%
\TextAt{8}{245.541}{172.475}{)}%
\TextAt{8}{254.369}{172.475}{and}%
\TextAt{13}{277.508}{172.475}{\Character{11}}%
\TextAt{9}{285.029}{172.475}{\Character{40}}%
\TextAt{13}{289.582}{172.475}{\Character{27}}%
\TextAt{9}{296.664}{172.475}{\Character{41}}%
\NativeMathAt{305.178}{172.475}{11.955}{\in}%
\TextAt{18}{317.109}{172.475}{R}%
\TextAt{3}{325.744}{176.813}{\Character{3}}%
\TextAt{17}{325.744}{169.519}{\Character{43}}%
\TextAt{8}{332.828}{172.475}{,}%
\TextAt{8}{340.447}{172.475}{resp.}%
\TextAt{8}{370.897}{172.475}{if}%
\TextAt{8}{381.996}{172.475}{all}%
\TextAt{8}{398.631}{172.475}{the}%
\TextAt{8}{419.164}{172.475}{functions}%
\TextAt{13}{470.002}{172.475}{\Character{11}\Character{59}}%
\TextAt{13}{482.746}{172.475}{\Character{12}}%
\TextAt{13}{490.015}{172.475}{\Character{59}}%
\TextAt{13}{495.251}{172.475}{\Character{13}}%
\TextAt{13}{501.958}{172.475}{\Character{59}}%
\TextAt{13}{507.195}{172.475}{\Character{14}}%
\TextAt{8}{517.156}{172.475}{are}%
\TextAt{13}{79.226}{155.429}{F}%
\TextAt{8}{88.43}{155.429}{-invariant,}%
\TextAt{13}{144.854}{155.429}{\Character{11}}%
\TextAt{9}{152.376}{155.429}{\Character{40}}%
\TextAt{13}{156.928}{155.429}{\Character{27}}%
\TextAt{9}{164.01}{155.429}{\Character{41}}%
\TextAt{13}{168.563}{155.429}{\Character{59}}%
\TextAt{13}{173.787}{155.429}{\Character{12}}%
\TextAt{9}{181.078}{155.429}{\Character{40}}%
\TextAt{13}{185.631}{155.429}{\Character{27}}%
\TextAt{9}{192.713}{155.429}{\Character{41}}%
\NativeMathAt{200.792}{155.429}{11.955}{\in}%
\TextAt{18}{212.29}{155.429}{R}%
\TextAt{3}{220.924}{159.767}{\Character{3}}%
\TextAt{17}{220.924}{152.473}{\Character{43}}%
\TextAt{8}{232.027}{155.429}{and}%
\TextAt{13}{254.911}{155.429}{\Character{14}}%
\TextAt{9}{264.071}{155.429}{\Character{61}}%
\Vector{q 0 G 0 g 0.478 w 204.702 -611.75 m 211.425 -611.75 l S Q}%
\TextAt{13}{276.702}{155.429}{\Character{13}}%
\TextAt{8}{283.425}{155.429}{,}%
\TextAt{8}{290.725}{155.429}{and}%
\TextAt{8}{313.609}{155.429}{if}%
\TextAt{8}{324.454}{155.429}{moreover}%
\TextAt{8}{374.99}{155.429}{det}%
\TextAt{13}{393.241}{155.429}{A}%
\TextAt{9}{405.544}{155.429}{\Character{61}}%
\TextAt{13}{418.175}{155.429}{\Character{11}\Character{12}}%
\TextAt{16}{435.708}{155.429}{\Character{0}}%
\TextAt{13}{447.745}{155.429}{\Character{13}}%
\Vector{q 0 G 0 g 0.478 w 382.467 -611.75 m 389.19 -611.75 l S Q}%
\TextAt{13}{454.467}{155.429}{\Character{13}}%
\TextAt{13}{464.701}{155.429}{\Character{62}}%
\TextAt{9}{477.348}{155.429}{0}%
\TextAt{8}{483.201}{155.429}{.}%
\TextAt{8}{492.013}{155.429}{We}%
\TextAt{8}{512.278}{155.429}{now}%
\TextAt{8}{79.226}{138.383}{assume}%
\TextAt{8}{119.67}{138.383}{this}%
\TextAt{8}{142.493}{138.383}{in}%
\TextAt{8}{156.146}{138.383}{what}%
\TextAt{8}{185.412}{138.383}{follows.}%
\TextAt{8}{97.159}{118.348}{Let}%
\TextAt{14}{118.128}{118.348}{a}%
\TextAt{8}{128.009}{118.348}{and}%
\TextAt{14}{150.771}{118.348}{b}%
\TextAt{8}{160.804}{118.348}{be}%
\TextAt{8}{176.728}{118.348}{fractional}%
\TextAt{8}{229.092}{118.348}{ideals}%
\TextAt{8}{261.669}{118.348}{of}%
\TextAt{13}{274.999}{118.348}{K}%
\TextAt{8}{285.817}{118.348}{.}%
\TextAt{8}{294.265}{118.348}{We}%
\TextAt{8}{314.41}{118.348}{have}%
\TextAt{8}{341.405}{118.348}{to}%
\TextAt{8}{355.715}{118.348}{prove}%
\TextAt{8}{387.265}{118.348}{the}%
\TextAt{8}{407.421}{118.348}{formula}%
\TextAt{9}{257.281}{89.347}{\Character{91}}%
\TextAt{14}{260.533}{89.347}{a}%
\TextAt{9}{266.516}{89.347}{\Character{93}}%
\TextAt{9}{272.41}{89.347}{\Character{43}}%
\TextAt{9}{284.162}{89.347}{\Character{91}}%
\TextAt{14}{287.438}{89.347}{b}%
\TextAt{9}{293.572}{89.347}{\Character{93}}%
\TextAt{9}{300.123}{89.347}{\Character{61}}%
\TextAt{9}{312.545}{89.347}{\Character{91}}%
\TextAt{14}{315.822}{89.347}{ab}%
\TextAt{9}{327.94}{89.347}{\Character{93}}%
\TextAt{9}{333.834}{89.347}{\Character{43}}%
\TextAt{9}{345.586}{89.347}{1}%
\TextAt{13}{351.463}{89.347}{\Character{58}}%
\TextAt{8}{297.211}{43.033}{220}%
\Destination{78.23}{809.815}{nav.page.220}%
\Destination{79.36}{537.855}{nav.print.240}%
\Destination{79.36}{715.055}{nav.number.3.5.6}%
\Destination{79.36}{423.755}{nav.number.3.5.7}%
\Hotlink{370.419}{663.855}{25.078}{11.047}{nav.number.3.4.3}%
}

% Page 233 (body page 221)
\Page{%
\TextAt{2}{79.226}{796.943}{III.}%
\TextAt{2}{98.317}{796.943}{Riemann–Roch}%
\TextAt{2}{174.375}{796.943}{Theory}%
\Vector{q 0 G 0 g 0.299 w 7.226 18.729 m 460.769 18.729 l S Q}%
\TextAt{8}{79.226}{761.124}{We}%
\TextAt{8}{98.745}{761.124}{may}%
\TextAt{8}{123.489}{761.124}{assume}%
\TextAt{8}{163.308}{761.124}{that}%
\TextAt{14}{188.052}{761.124}{a}%
\TextAt{21}{194.035}{759.33}{f}%
\TextAt{8}{202.752}{761.124}{and}%
\TextAt{14}{224.89}{761.124}{b}%
\TextAt{21}{231.024}{759.33}{f}%
\TextAt{8}{239.742}{761.124}{are}%
\TextAt{19}{258.627}{761.124}{integral}%
\TextAt{19}{299.656}{761.124}{ideals}%
\TextAt{19}{331.263}{761.124}{relatively}%
\TextAt{19}{380.362}{761.124}{prime}%
\TextAt{8}{413.163}{761.124}{to}%
\TextAt{8}{426.848}{761.124}{one}%
\TextAt{8}{447.683}{761.124}{another,}%
\TextAt{8}{493.365}{761.124}{because}%
\TextAt{8}{79.226}{744.078}{if}%
\TextAt{8}{90.508}{744.078}{necessary}%
\TextAt{8}{142.891}{744.078}{we}%
\TextAt{8}{160.677}{744.078}{may}%
\TextAt{8}{186.603}{744.078}{pass}%
\TextAt{8}{212.65}{744.078}{to}%
\TextAt{8}{227.518}{744.078}{replete}%
\TextAt{8}{266.44}{744.078}{ideals}%
\TextAt{14}{299.575}{744.078}{a}%
\TextAt{3}{305.559}{748.416}{\Character{48}}%
\TextAt{9}{312.625}{744.078}{\Character{61}}%
\TextAt{14}{326.001}{744.078}{a}%
\TextAt{9}{331.984}{744.078}{\Character{91}}%
\TextAt{13}{335.236}{744.078}{a}%
\TextAt{9}{341.381}{744.078}{\Character{93}}%
\TextAt{8}{344.633}{744.078}{,}%
\TextAt{14}{352.479}{744.078}{b}%
\TextAt{3}{358.614}{748.416}{\Character{48}}%
\TextAt{9}{365.681}{744.078}{\Character{61}}%
\TextAt{14}{379.057}{744.078}{b}%
\TextAt{9}{385.192}{744.078}{\Character{91}}%
\TextAt{13}{388.443}{744.078}{b}%
\TextAt{9}{393.42}{744.078}{\Character{93}}%
\TextAt{8}{401.127}{744.078}{with}%
\TextAt{8}{428.345}{744.078}{corresponding}%
\TextAt{8}{504.089}{744.078}{ideals}%
\TextAt{14}{79.226}{727.032}{a}%
\TextAt{3}{85.209}{731.37}{\Character{48}}%
\TextAt{21}{85.209}{723.923}{f}%
\TextAt{9}{95.654}{727.032}{\Character{61}}%
\TextAt{14}{109.759}{727.032}{a}%
\TextAt{21}{115.742}{725.238}{f}%
\TextAt{13}{121.187}{727.032}{a}%
\TextAt{8}{127.332}{727.032}{,}%
\TextAt{14}{135.713}{727.032}{b}%
\TextAt{3}{141.848}{731.37}{\Character{48}}%
\TextAt{21}{141.848}{723.923}{f}%
\TextAt{9}{152.293}{727.032}{\Character{61}}%
\TextAt{14}{166.398}{727.032}{b}%
\TextAt{21}{172.533}{725.238}{f}%
\TextAt{13}{177.978}{727.032}{b}%
\TextAt{8}{187.837}{727.032}{without}%
\TextAt{8}{232.399}{727.032}{changing}%
\TextAt{8}{282.497}{727.032}{the}%
\TextAt{8}{303.639}{727.032}{classes}%
\TextAt{9}{341.876}{727.032}{\Character{91}}%
\TextAt{14}{345.128}{727.032}{a}%
\TextAt{9}{351.111}{727.032}{\Character{93}}%
\TextAt{13}{354.363}{727.032}{\Character{59}}%
\TextAt{9}{359.607}{727.032}{\Character{91}}%
\TextAt{14}{362.859}{727.032}{b}%
\TextAt{9}{368.993}{727.032}{\Character{93}}%
\TextAt{13}{372.245}{727.032}{\Character{59}}%
\TextAt{9}{377.489}{727.032}{\Character{91}}%
\TextAt{14}{380.741}{727.032}{ab}%
\TextAt{9}{392.859}{727.032}{\Character{93}}%
\TextAt{8}{400.993}{727.032}{in}%
\TextAt{13}{415.632}{727.032}{K}%
\TextAt{17}{425.591}{725.238}{0}%
\TextAt{9}{430.323}{727.032}{\Character{40}}%
\Vector{q 0 G 0 g 0.478 w 362.875 -39.612 m 368.72 -39.612 l S Q}%
\TextAt{14}{434.875}{727.032}{o}%
\TextAt{9}{440.72}{727.032}{\Character{41}}%
\TextAt{8}{445.272}{727.032}{.}%
\TextAt{8}{456.677}{727.032}{We}%
\TextAt{8}{477.807}{727.032}{denote}%
\TextAt{8}{516.511}{727.032}{the}%
\TextAt{14}{79.226}{709.986}{o}%
\TextAt{8}{85.07}{709.986}{-module}%
\TextAt{14}{130.627}{709.986}{a}%
\TextAt{21}{136.61}{708.193}{f}%
\TextAt{8}{142.055}{709.986}{,}%
\TextAt{8}{149.662}{709.986}{when}%
\TextAt{8}{180.585}{709.986}{metrized}%
\TextAt{8}{229.071}{709.986}{by}%
\TextAt{13}{245.696}{709.986}{\Character{11}x}%
\Vector{q 0 G 0 g 0.478 w 187.87 -57.193 m 194.007 -57.193 l S Q}%
\TextAt{13}{259.87}{709.986}{y}%
\TextAt{8}{266.007}{709.986}{,}%
\TextAt{8}{273.614}{709.986}{by}%
\TextAt{9}{290.239}{709.986}{\Character{40}}%
\TextAt{14}{294.791}{709.986}{a}%
\TextAt{21}{300.774}{708.193}{f}%
\TextAt{13}{306.219}{709.986}{\Character{59}}%
\TextAt{13}{311.443}{709.986}{\Character{11}}%
\TextAt{9}{318.985}{709.986}{\Character{41}}%
\TextAt{8}{323.537}{709.986}{,}%
\TextAt{8}{331.144}{709.986}{and}%
\TextAt{8}{354.273}{709.986}{the}%
\TextAt{14}{374.796}{709.986}{o}%
\TextAt{8}{380.64}{709.986}{-module}%
\TextAt{14}{426.197}{709.986}{a}%
\TextAt{21}{432.18}{708.193}{f}%
\TextAt{16}{440.531}{709.986}{\Character{8}}%
\TextAt{14}{452.736}{709.986}{b}%
\TextAt{21}{458.87}{708.193}{f}%
\TextAt{8}{464.315}{709.986}{,}%
\TextAt{8}{471.922}{709.986}{metrized}%
\TextAt{8}{520.408}{709.986}{by}%
\TextAt{16}{79.226}{684.725}{\Character{104}}%
\TextAt{13}{88.923}{684.725}{\Character{59}}%
\TextAt{16}{99.214}{684.725}{\Character{105}}%
\TextAt{21}{103.863}{682.932}{A}%
\TextAt{8}{115.752}{684.725}{for}%
\TextAt{13}{134.787}{684.725}{A}%
\TextAt{9}{148.842}{684.725}{\Character{61}}%
\TextAt{20}{163.227}{705.168}{\Character{32}}%
\TextAt{13}{172.692}{692.827}{\Character{11}}%
\TextAt{13}{190.434}{692.827}{\Character{13}}%
\Vector{q 0 G 0 g 0.478 w 101.092 -91.398 m 107.814 -91.398 l S Q}%
\TextAt{13}{173.092}{675.781}{\Character{13}}%
\TextAt{13}{190.164}{675.781}{\Character{12}}%
\TextAt{20}{197.447}{705.168}{\Character{33}}%
\TextAt{8}{206.912}{684.725}{,}%
\TextAt{8}{215.499}{684.725}{by}%
\TextAt{9}{232.907}{684.725}{\Character{40}}%
\TextAt{14}{237.46}{684.725}{a}%
\TextAt{21}{243.443}{682.932}{f}%
\TextAt{16}{252.328}{684.725}{\Character{8}}%
\TextAt{14}{265.067}{684.725}{b}%
\TextAt{21}{271.202}{682.932}{f}%
\TextAt{13}{276.647}{684.725}{\Character{59}}%
\TextAt{13}{281.871}{684.725}{A}%
\TextAt{9}{290.667}{684.725}{\Character{41}}%
\TextAt{8}{295.219}{684.725}{.}%
\TextAt{8}{307.117}{684.725}{Given}%
\TextAt{8}{342.16}{684.725}{any}%
\TextAt{8}{365.427}{684.725}{two}%
\TextAt{8}{388.693}{684.725}{matrices}%
\TextAt{13}{436.732}{684.725}{A}%
\TextAt{9}{450.787}{684.725}{\Character{61}}%
\TextAt{20}{465.172}{705.168}{\Character{32}}%
\TextAt{13}{474.637}{692.827}{\Character{11}}%
\TextAt{13}{492.379}{692.827}{\Character{13}}%
\Vector{q 0 G 0 g 0.478 w 403.037 -91.398 m 409.759 -91.398 l S Q}%
\TextAt{13}{475.037}{675.781}{\Character{13}}%
\TextAt{13}{492.109}{675.781}{\Character{12}}%
\TextAt{20}{499.393}{705.168}{\Character{33}}%
\TextAt{8}{513.904}{684.725}{and}%
\TextAt{13}{79.226}{647.863}{A}%
\TextAt{3}{88.002}{652.201}{\Character{48}}%
\TextAt{9}{94.117}{647.863}{\Character{61}}%
\TextAt{20}{106.543}{668.306}{\Character{32}}%
\TextAt{13}{116.007}{655.965}{\Character{11}}%
\TextAt{3}{123.529}{660.304}{\Character{48}}%
\TextAt{13}{136.561}{655.965}{\Character{13}}%
\TextAt{3}{143.284}{660.304}{\Character{48}}%
\Vector{q 0 G 0 g 0.478 w 44.407 -125.525 m 53.924 -125.525 l S Q}%
\TextAt{13}{116.407}{638.919}{\Character{13}}%
\TextAt{3}{123.129}{642.373}{\Character{48}}%
\TextAt{13}{136.287}{638.919}{\Character{12}}%
\TextAt{3}{143.558}{643.258}{\Character{48}}%
\TextAt{20}{146.353}{668.306}{\Character{33}}%
\TextAt{8}{159.715}{647.863}{we}%
\TextAt{8}{176.943}{647.863}{write}%
\TextAt{13}{286.229}{622.558}{A}%
\NativeMathAt{298.325}{622.558}{11.955}{\sim}%
\TextAt{13}{310.945}{622.558}{A}%
\TextAt{3}{319.72}{627.494}{\Character{48}}%
\TextAt{13}{322.515}{622.558}{\Character{59}}%
\TextAt{8}{79.226}{599.857}{if}%
\TextAt{9}{89.778}{599.857}{\Character{91}\Character{40}}%
\TextAt{14}{97.582}{599.857}{a}%
\TextAt{21}{103.565}{598.064}{f}%
\TextAt{16}{111.315}{599.857}{\Character{8}}%
\TextAt{14}{122.919}{599.857}{b}%
\TextAt{21}{129.054}{598.064}{f}%
\TextAt{9}{134.499}{599.857}{\Character{41}}%
\TextAt{13}{139.051}{599.857}{\Character{59}}%
\TextAt{13}{144.275}{599.857}{A}%
\TextAt{9}{153.071}{599.857}{\Character{93}}%
\TextAt{9}{159.622}{599.857}{\Character{61}}%
\TextAt{9}{172.044}{599.857}{\Character{91}\Character{40}}%
\TextAt{14}{179.873}{599.857}{a}%
\TextAt{21}{185.856}{598.064}{f}%
\TextAt{16}{193.606}{599.857}{\Character{8}}%
\TextAt{14}{205.21}{599.857}{b}%
\TextAt{21}{211.344}{598.064}{f}%
\TextAt{9}{216.789}{599.857}{\Character{41}}%
\TextAt{13}{221.342}{599.857}{\Character{59}}%
\TextAt{13}{226.566}{599.857}{A}%
\TextAt{3}{235.361}{604.195}{\Character{48}}%
\TextAt{9}{238.156}{599.857}{\Character{93}}%
\TextAt{8}{245.134}{599.857}{in}%
\TextAt{13}{258.614}{599.857}{K}%
\TextAt{17}{268.574}{598.064}{0}%
\TextAt{9}{273.306}{599.857}{\Character{40}}%
\Vector{q 0 G 0 g 0.478 w 205.858 -166.787 m 211.702 -166.787 l S Q}%
\TextAt{14}{277.858}{599.857}{o}%
\TextAt{9}{283.702}{599.857}{\Character{41}}%
\TextAt{8}{288.255}{599.857}{.}%
\TextAt{8}{296.646}{599.857}{We}%
\TextAt{8}{316.618}{599.857}{now}%
\TextAt{8}{340.835}{599.857}{consider}%
\TextAt{8}{386.248}{599.857}{the}%
\TextAt{8}{406.233}{599.857}{canonical}%
\TextAt{8}{457.444}{599.857}{exact}%
\TextAt{8}{488.165}{599.857}{sequence}%
\TextAt{9}{216.4}{572.77}{0}%
\NativeArrowAt{225.574}{572.77}{11.955}{19.26}{\longrightarrow}%
\TextAt{14}{248.156}{572.77}{a}%
\TextAt{21}{254.139}{570.977}{f}%
\NativeArrowAt{262.904}{572.77}{11.955}{19.26}{\longrightarrow}%
\TextAt{14}{285.486}{572.77}{a}%
\TextAt{21}{291.469}{570.977}{f}%
\TextAt{16}{299.571}{572.77}{\Character{8}}%
\TextAt{14}{311.526}{572.77}{b}%
\TextAt{21}{317.661}{570.977}{f}%
\NativeArrowAt{326.427}{572.77}{11.955}{19.26}{\longrightarrow}%
\TextAt{14}{349.009}{572.77}{b}%
\TextAt{21}{355.144}{570.977}{f}%
\NativeArrowAt{363.909}{572.77}{11.955}{19.26}{\longrightarrow}%
\TextAt{9}{386.491}{572.77}{0}%
\TextAt{13}{392.344}{572.77}{\Character{58}}%
\TextAt{8}{79.226}{537.468}{Once}%
\TextAt{8}{109.13}{537.468}{we}%
\TextAt{8}{126.35}{537.468}{equip}%
\TextAt{14}{157.88}{537.468}{a}%
\TextAt{21}{163.864}{535.674}{f}%
\TextAt{16}{171.95}{537.468}{\Character{8}}%
\TextAt{14}{183.889}{537.468}{b}%
\TextAt{21}{190.024}{535.674}{f}%
\TextAt{8}{199.359}{537.468}{with}%
\TextAt{8}{226.011}{537.468}{the}%
\TextAt{8}{246.16}{537.468}{metric}%
\TextAt{16}{282.568}{537.468}{\Character{104}}%
\TextAt{13}{291.107}{537.468}{\Character{59}}%
\TextAt{16}{300.241}{537.468}{\Character{105}}%
\TextAt{21}{304.891}{535.674}{A}%
\TextAt{8}{315.622}{537.468}{which}%
\TextAt{8}{349.1}{537.468}{is}%
\TextAt{8}{360.857}{537.468}{given}%
\TextAt{8}{391.419}{537.468}{by}%
\TextAt{13}{407.67}{537.468}{A}%
\TextAt{9}{419.766}{537.468}{\Character{61}}%
\TextAt{20}{432.192}{557.911}{\Character{32}}%
\TextAt{13}{441.656}{545.57}{\Character{11}}%
\TextAt{13}{459.398}{545.57}{\Character{13}}%
\Vector{q 0 G 0 g 0.478 w 370.056 -238.655 m 376.778 -238.655 l S Q}%
\TextAt{13}{442.056}{528.524}{\Character{13}}%
\TextAt{13}{459.128}{528.524}{\Character{12}}%
\TextAt{20}{466.412}{557.911}{\Character{33}}%
\TextAt{8}{475.876}{537.468}{,}%
\TextAt{8}{483.02}{537.468}{we}%
\TextAt{8}{500.24}{537.468}{obtain}%
\TextAt{8}{79.226}{513.111}{the}%
\TextAt{8}{99.383}{513.111}{following}%
\TextAt{8}{148.817}{513.111}{exact}%
\TextAt{8}{179.709}{513.111}{sequence}%
\TextAt{8}{228.211}{513.111}{of}%
\TextAt{8}{241.542}{513.111}{metrized}%
\TextAt{14}{289.661}{513.111}{o}%
\TextAt{8}{295.505}{513.111}{-modules:}%
\TextAt{9}{163.621}{482.288}{0}%
\NativeArrowAt{172.794}{482.288}{11.955}{19.26}{\longrightarrow}%
\TextAt{9}{195.376}{482.288}{\Character{40}}%
\TextAt{14}{199.929}{482.288}{a}%
\TextAt{21}{205.912}{480.494}{f}%
\TextAt{13}{211.357}{482.288}{\Character{59}}%
\TextAt{13}{216.581}{482.288}{\Character{11}}%
\TextAt{9}{224.123}{482.288}{\Character{41}}%
\NativeArrowAt{231.996}{482.288}{11.955}{19.26}{\longrightarrow}%
\TextAt{9}{254.578}{482.288}{\Character{40}}%
\TextAt{14}{259.13}{482.288}{a}%
\TextAt{21}{265.113}{480.494}{f}%
\TextAt{16}{273.215}{482.288}{\Character{8}}%
\TextAt{14}{285.17}{482.288}{b}%
\TextAt{21}{291.305}{480.494}{f}%
\TextAt{13}{296.75}{482.288}{\Character{59}}%
\TextAt{13}{301.974}{482.288}{A}%
\TextAt{9}{310.769}{482.288}{\Character{41}}%
\NativeArrowAt{318.642}{482.288}{11.955}{19.26}{\longrightarrow}%
\TextAt{20}{341.224}{499.144}{\Character{18}}%
\TextAt{14}{350.025}{482.288}{b}%
\TextAt{21}{356.16}{480.494}{f}%
\TextAt{13}{361.604}{482.288}{\Character{59}}%
\TextAt{13}{366.828}{482.288}{\Character{12}}%
\TextAt{16}{376.777}{482.288}{\Character{0}}%
\TextAt{13}{389.927}{490.375}{\Character{13}}%
\Vector{q 0 G 0 g 0.478 w 324.649 -276.804 m 331.372 -276.804 l S Q}%
\TextAt{13}{396.649}{490.375}{\Character{13}}%
\Vector{q 0 G 0 g 0.478 w 317.927 -288.724 m 331.372 -288.724 l S Q}%
\TextAt{13}{392.889}{474.087}{\Character{11}}%
\TextAt{20}{404.567}{499.144}{\Character{19}}%
\NativeArrowAt{416.688}{482.288}{11.955}{19.26}{\longrightarrow}%
\TextAt{9}{439.27}{482.288}{0}%
\TextAt{13}{445.123}{482.288}{\Character{58}}%
\TextAt{8}{517.802}{482.288}{(*)}%
\TextAt{8}{79.226}{451.596}{Indeed,}%
\TextAt{8}{120.507}{451.596}{in}%
\TextAt{8}{134.16}{451.596}{the}%
\TextAt{8}{154.317}{451.596}{exact}%
\TextAt{8}{185.209}{451.596}{sequence}%
\TextAt{9}{206.98}{424.509}{0}%
\NativeArrowAt{216.154}{424.509}{11.955}{19.26}{\longrightarrow}%
\TextAt{13}{238.736}{424.509}{K}%
\TextAt{25}{248.695}{422.715}{C}%
\NativeArrowAt{258.271}{424.509}{11.955}{19.26}{\longrightarrow}%
\TextAt{13}{280.853}{424.509}{K}%
\TextAt{25}{290.812}{422.715}{C}%
\TextAt{16}{299.723}{424.509}{\Character{8}}%
\TextAt{13}{311.678}{424.509}{K}%
\TextAt{25}{321.637}{422.715}{C}%
\NativeArrowAt{331.212}{424.509}{11.955}{19.26}{\longrightarrow}%
\TextAt{13}{353.794}{424.509}{K}%
\TextAt{25}{363.753}{422.715}{C}%
\NativeArrowAt{373.328}{424.509}{11.955}{19.26}{\longrightarrow}%
\TextAt{9}{395.91}{424.509}{0}%
\TextAt{13}{401.763}{424.509}{\Character{59}}%
\TextAt{8}{79.226}{397.422}{the}%
\TextAt{8}{100.135}{397.422}{restriction}%
\TextAt{8}{156.886}{397.422}{of}%
\TextAt{16}{170.969}{397.422}{\Character{104}}%
\TextAt{13}{180.269}{397.422}{\Character{59}}%
\TextAt{16}{190.163}{397.422}{\Character{105}}%
\TextAt{21}{194.813}{395.628}{A}%
\TextAt{8}{206.304}{397.422}{to}%
\TextAt{13}{221.367}{397.422}{K}%
\TextAt{25}{231.327}{395.628}{C}%
\TextAt{16}{240.751}{397.422}{\Character{8}}%
\TextAt{16}{253.208}{397.422}{\Character{102}}%
\TextAt{9}{259.197}{397.422}{0}%
\TextAt{16}{265.05}{397.422}{\Character{103}}%
\TextAt{8}{275.678}{397.422}{yields}%
\TextAt{8}{309.331}{397.422}{the}%
\TextAt{8}{330.241}{397.422}{metric}%
\TextAt{13}{367.409}{397.422}{\Character{11}x}%
\Vector{q 0 G 0 g 0.478 w 309.583 -369.757 m 315.719 -369.757 l S Q}%
\TextAt{13}{381.583}{397.422}{y}%
\TextAt{8}{392.37}{397.422}{on}%
\TextAt{13}{409.382}{397.422}{K}%
\TextAt{25}{419.341}{395.628}{C}%
\TextAt{8}{425.595}{397.422}{,}%
\TextAt{8}{433.686}{397.422}{and}%
\TextAt{8}{457.201}{397.422}{the}%
\TextAt{8}{478.111}{397.422}{orthogonal}%
\TextAt{8}{79.226}{380.376}{complement}%
\TextAt{13}{144.585}{380.376}{V}%
\TextAt{8}{157.908}{380.376}{of}%
\TextAt{13}{171.238}{380.376}{K}%
\TextAt{25}{181.197}{378.582}{C}%
\TextAt{16}{190.108}{380.376}{\Character{8}}%
\TextAt{16}{202.051}{380.376}{\Character{102}}%
\TextAt{9}{208.041}{380.376}{0}%
\TextAt{16}{213.894}{380.376}{\Character{103}}%
\TextAt{8}{223.769}{380.376}{consists}%
\TextAt{8}{266.879}{380.376}{of}%
\TextAt{8}{280.209}{380.376}{all}%
\TextAt{8}{296.468}{380.376}{elements}%
\TextAt{13}{344.325}{380.376}{a}%
\TextAt{9}{353.127}{380.376}{\Character{43}}%
\TextAt{13}{364.888}{380.376}{b}%
\NativeMathAt{373.186}{380.376}{11.955}{\in}%
\TextAt{13}{384.477}{380.376}{K}%
\TextAt{25}{394.436}{378.582}{C}%
\TextAt{16}{403.347}{380.376}{\Character{8}}%
\TextAt{13}{415.302}{380.376}{K}%
\TextAt{25}{425.261}{378.582}{C}%
\TextAt{8}{435.413}{380.376}{such}%
\TextAt{8}{461.81}{380.376}{that}%
\TextAt{16}{225.458}{353.288}{\Character{104}}%
\TextAt{13}{230.108}{353.288}{x}%
\TextAt{16}{239.416}{353.288}{\Character{8}}%
\TextAt{9}{251.372}{353.288}{0}%
\TextAt{13}{257.225}{353.288}{\Character{59}}%
\TextAt{13}{262.449}{353.288}{a}%
\TextAt{16}{271.27}{353.288}{\Character{8}}%
\TextAt{13}{283.225}{353.288}{b}%
\TextAt{16}{288.203}{353.288}{\Character{105}}%
\TextAt{9}{296.173}{353.288}{\Character{61}}%
\TextAt{13}{308.598}{353.288}{\Character{11}x}%
\Vector{q 0 G 0 g 0.478 w 250.772 -413.89 m 256.917 -413.89 l S Q}%
\TextAt{13}{322.772}{353.288}{a}%
\TextAt{9}{331.574}{353.288}{\Character{43}}%
\TextAt{13}{343.335}{353.288}{\Character{13}}%
\TextAt{13}{350.042}{353.288}{x}%
\Vector{q 0 G 0 g 0.478 w 284.709 -410.736 m 289.686 -410.736 l S Q}%
\TextAt{13}{356.709}{353.288}{b}%
\TextAt{9}{365.007}{353.288}{\Character{61}}%
\TextAt{9}{377.417}{353.288}{0}%
\TextAt{13}{383.286}{353.288}{\Character{59}}%
\TextAt{8}{79.226}{326.201}{for}%
\TextAt{8}{97.111}{326.201}{all}%
\TextAt{13}{113.37}{326.201}{x}%
\NativeMathAt{123.343}{326.201}{11.955}{\in}%
\TextAt{13}{134.634}{326.201}{K}%
\TextAt{25}{144.593}{324.408}{C}%
\TextAt{8}{150.847}{326.201}{,}%
\TextAt{8}{157.997}{326.201}{so}%
\TextAt{8}{172.367}{326.201}{that}%
\TextAt{13}{231.393}{309.155}{V}%
\TextAt{9}{244.139}{309.155}{\Character{61}}%
\TextAt{16}{256.565}{309.155}{\Character{102}}%
\TextAt{9}{262.543}{309.155}{\Character{40}}%
\TextAt{16}{267.095}{309.155}{\Character{0}}%
\Vector{q 0 G 0 g 0.478 w 204.393 -458.024 m 211.116 -458.024 l S Q}%
\TextAt{13}{276.393}{309.155}{\Character{13}}%
\TextAt{9}{283.116}{309.155}{\Character{47}}%
\TextAt{13}{288.969}{309.155}{\Character{11}}%
\TextAt{9}{296.49}{309.155}{\Character{41}}%
\TextAt{13}{301.043}{309.155}{b}%
\TextAt{16}{308.676}{309.155}{\Character{8}}%
\TextAt{13}{320.632}{309.155}{b}%
\TextAt{16}{328.929}{309.155}{\Character{106}}%
\TextAt{13}{335.571}{309.155}{b}%
\NativeMathAt{343.869}{309.155}{11.955}{\in}%
\TextAt{13}{355.16}{309.155}{K}%
\TextAt{25}{365.119}{307.362}{C}%
\TextAt{16}{371.374}{309.155}{\Character{103}}%
\TextAt{13}{377.351}{309.155}{\Character{58}}%
\TextAt{8}{79.226}{286.454}{The}%
\TextAt{8}{103.683}{286.454}{isomorphism}%
\TextAt{13}{172.505}{286.454}{V}%
\TextAt{21}{188.636}{293.231}{\Character{25}}%
\NativeArrowAt{185.938}{286.454}{11.955}{11.955}{\longrightarrow}%
\TextAt{13}{201.901}{286.454}{K}%
\TextAt{25}{211.86}{284.661}{C}%
\TextAt{8}{218.115}{286.454}{,}%
\TextAt{9}{225.768}{286.454}{\Character{40}}%
\TextAt{16}{230.32}{286.454}{\Character{0}}%
\Vector{q 0 G 0 g 0.478 w 167.618 -480.725 m 174.341 -480.725 l S Q}%
\TextAt{13}{239.618}{286.454}{\Character{13}}%
\TextAt{9}{246.341}{286.454}{\Character{47}}%
\TextAt{13}{252.194}{286.454}{\Character{11}}%
\TextAt{9}{259.715}{286.454}{\Character{41}}%
\TextAt{13}{264.268}{286.454}{b}%
\TextAt{16}{272.176}{286.454}{\Character{8}}%
\TextAt{13}{284.406}{286.454}{b}%
\NativeArrowAt{293.39}{286.454}{11.955}{11.955}{\mapsto}%
\TextAt{13}{309.353}{286.454}{b}%
\TextAt{8}{314.33}{286.454}{,}%
\TextAt{8}{321.983}{286.454}{transfers}%
\TextAt{8}{370.314}{286.454}{the}%
\TextAt{8}{390.874}{286.454}{metric}%
\TextAt{16}{427.692}{286.454}{\Character{104}}%
\TextAt{13}{436.642}{286.454}{\Character{59}}%
\TextAt{16}{446.186}{286.454}{\Character{105}}%
\TextAt{21}{450.835}{284.661}{A}%
\TextAt{8}{461.977}{286.454}{on}%
\TextAt{13}{478.639}{286.454}{V}%
\TextAt{8}{492.365}{286.454}{into}%
\TextAt{8}{516.511}{286.454}{the}%
\TextAt{8}{79.226}{269.408}{metric}%
\TextAt{13}{115.642}{269.408}{\Character{14}}%
\TextAt{13}{121.261}{269.408}{x}%
\Vector{q 0 G 0 g 0.478 w 55.927 -497.771 m 62.063 -497.771 l S Q}%
\TextAt{13}{127.927}{269.408}{y}%
\TextAt{8}{134.063}{269.408}{,}%
\TextAt{8}{141.212}{269.408}{where}%
\TextAt{13}{175.022}{269.408}{\Character{14}}%
\TextAt{8}{184.552}{269.408}{is}%
\TextAt{8}{196.316}{269.408}{determined}%
\TextAt{8}{257.443}{269.408}{by}%
\TextAt{8}{273.702}{269.408}{the}%
\TextAt{8}{293.858}{269.408}{rule}%
\TextAt{13}{167.888}{242.321}{\Character{14}}%
\TextAt{9}{176.841}{242.321}{\Character{61}}%
\TextAt{16}{189.267}{242.321}{\Character{104}}%
\TextAt{13}{193.916}{242.321}{\Character{25}}%
\TextAt{3}{200.985}{247.257}{\Character{0}}%
\TextAt{17}{207.572}{247.257}{1}%
\TextAt{9}{212.304}{242.321}{\Character{40}1\Character{41}}%
\TextAt{13}{227.262}{242.321}{\Character{59}}%
\TextAt{13}{232.486}{242.321}{\Character{25}}%
\TextAt{3}{239.575}{247.257}{\Character{0}}%
\TextAt{17}{246.162}{247.257}{1}%
\TextAt{9}{250.894}{242.321}{\Character{40}1\Character{41}}%
\TextAt{16}{265.852}{242.321}{\Character{105}}%
\TextAt{21}{270.501}{240.528}{A}%
\TextAt{9}{280.663}{242.321}{\Character{61}}%
\TextAt{16}{293.088}{242.321}{\Character{104}}%
\TextAt{9}{297.738}{242.321}{\Character{40}}%
\TextAt{16}{302.29}{242.321}{\Character{0}}%
\Vector{q 0 G 0 g 0.478 w 239.589 -524.858 m 246.311 -524.858 l S Q}%
\TextAt{13}{311.589}{242.321}{\Character{13}}%
\TextAt{9}{318.311}{242.321}{\Character{47}}%
\TextAt{13}{324.164}{242.321}{\Character{11}}%
\TextAt{9}{331.685}{242.321}{\Character{41}1}%
\TextAt{16}{344.747}{242.321}{\Character{8}}%
\TextAt{9}{356.703}{242.321}{1}%
\TextAt{13}{362.556}{242.321}{\Character{59}}%
\TextAt{9}{367.8}{242.321}{\Character{40}}%
\TextAt{16}{372.352}{242.321}{\Character{0}}%
\Vector{q 0 G 0 g 0.478 w 309.651 -524.858 m 316.373 -524.858 l S Q}%
\TextAt{13}{381.651}{242.321}{\Character{13}}%
\TextAt{9}{388.373}{242.321}{\Character{47}}%
\TextAt{13}{394.226}{242.321}{\Character{11}}%
\TextAt{9}{401.748}{242.321}{\Character{41}1}%
\TextAt{16}{414.81}{242.321}{\Character{8}}%
\TextAt{9}{426.765}{242.321}{1}%
\TextAt{16}{432.618}{242.321}{\Character{105}}%
\TextAt{21}{437.267}{240.528}{A}%
\TextAt{9}{176.841}{217.596}{\Character{61}}%
\TextAt{13}{189.267}{217.596}{\Character{11}}%
\TextAt{13}{197.984}{225.684}{\Character{13}}%
\Vector{q 0 G 0 g 0.478 w 132.706 -541.495 m 139.429 -541.495 l S Q}%
\TextAt{13}{204.706}{225.684}{\Character{13}}%
\Vector{q 0 G 0 g 0.478 w 125.984 -553.415 m 139.429 -553.415 l S Q}%
\TextAt{13}{198.579}{209.395}{\Character{11}}%
\TextAt{17}{206.101}{212.849}{2}%
\TextAt{16}{215.281}{217.596}{\Character{0}}%
\TextAt{13}{227.236}{217.596}{\Character{13}}%
\Vector{q 0 G 0 g 0.478 w 163.553 -541.495 m 170.276 -541.495 l S Q}%
\TextAt{13}{235.553}{225.684}{\Character{13}}%
\Vector{q 0 G 0 g 0.478 w 163.154 -553.415 m 170.675 -553.415 l S Q}%
\TextAt{13}{235.154}{209.395}{\Character{11}}%
\TextAt{16}{246.528}{217.596}{\Character{0}}%
\Vector{q 0 G 0 g 0.478 w 186.483 -549.583 m 193.205 -549.583 l S Q}%
\TextAt{13}{258.483}{217.596}{\Character{13}}%
\TextAt{13}{266.8}{225.684}{\Character{13}}%
\Vector{q 0 G 0 g 0.478 w 194.401 -553.415 m 201.922 -553.415 l S Q}%
\TextAt{13}{266.401}{209.395}{\Character{11}}%
\TextAt{9}{277.774}{217.596}{\Character{43}}%
\TextAt{13}{289.536}{217.596}{\Character{12}}%
\TextAt{9}{300.128}{217.596}{\Character{61}}%
\TextAt{13}{312.553}{217.596}{\Character{12}}%
\TextAt{16}{322.481}{217.596}{\Character{0}}%
\TextAt{13}{335.632}{225.684}{\Character{13}}%
\Vector{q 0 G 0 g 0.478 w 270.354 -541.495 m 277.076 -541.495 l S Q}%
\TextAt{13}{342.354}{225.684}{\Character{13}}%
\Vector{q 0 G 0 g 0.478 w 263.632 -553.415 m 277.076 -553.415 l S Q}%
\TextAt{13}{338.593}{209.395}{\Character{11}}%
\TextAt{13}{350.272}{217.596}{\Character{58}}%
\TextAt{8}{79.226}{189.392}{This}%
\TextAt{8}{105.946}{189.392}{shows}%
\TextAt{8}{139.564}{189.392}{that}%
\TextAt{8}{164.933}{189.392}{(*)}%
\TextAt{8}{183.798}{189.392}{is}%
\TextAt{8}{195.562}{189.392}{a}%
\TextAt{8}{205.317}{189.392}{short}%
\TextAt{8}{235.301}{189.392}{exact}%
\TextAt{8}{266.193}{189.392}{sequence}%
\TextAt{8}{314.695}{189.392}{of}%
\TextAt{8}{328.025}{189.392}{metrized}%
\TextAt{14}{376.145}{189.392}{o}%
\TextAt{8}{381.989}{189.392}{-modules,}%
\TextAt{8}{435.046}{189.392}{i.e.,}%
\TextAt{20}{238.387}{174.875}{\Character{32}}%
\TextAt{13}{247.852}{162.534}{\Character{11}}%
\TextAt{13}{265.594}{162.534}{\Character{13}}%
\Vector{q 0 G 0 g 0.478 w 176.252 -621.69 m 182.974 -621.69 l S Q}%
\TextAt{13}{248.252}{145.489}{\Character{13}}%
\TextAt{13}{265.324}{145.489}{\Character{12}}%
\TextAt{20}{272.608}{174.875}{\Character{33}}%
\NativeMathAt{285.393}{154.432}{11.955}{\sim}%
\TextAt{20}{298.012}{174.875}{\Character{32}}%
\TextAt{13}{307.477}{162.534}{\Character{11}}%
\TextAt{9}{339.004}{162.534}{0}%
\TextAt{9}{308.311}{145.488}{0}%
\TextAt{13}{324.961}{145.489}{\Character{12}}%
\TextAt{16}{334.889}{145.488}{\Character{0}}%
\TextAt{21}{348.04}{150.744}{\Character{13}}%
\Vector{q 0 G 0 g 0.359 w 280.872 -618.569 m 285.704 -618.569 l S Q}%
\TextAt{21}{352.872}{150.744}{\Character{13}}%
\Vector{q 0 G 0 g 0.478 w 276.04 -625.523 m 285.704 -625.523 l S Q}%
\TextAt{21}{350.152}{141.366}{\Character{11}}%
\TextAt{20}{358.9}{174.875}{\Character{33}}%
\TextAt{13}{370.357}{154.432}{\Character{58}}%
\TextAt{8}{79.226}{117.97}{Replacing}%
\TextAt{13}{133.371}{117.97}{\Character{12}}%
\TextAt{8}{144.54}{117.97}{by}%
\TextAt{13}{160.799}{117.97}{\Character{12}}%
\TextAt{9}{170.727}{117.97}{\Character{43}}%
\TextAt{21}{183.684}{123.225}{\Character{13}}%
\Vector{q 0 G 0 g 0.359 w 116.516 -646.088 m 121.348 -646.088 l S Q}%
\TextAt{21}{188.516}{123.225}{\Character{13}}%
\Vector{q 0 G 0 g 0.478 w 111.684 -653.042 m 121.348 -653.042 l S Q}%
\TextAt{21}{185.796}{113.847}{\Character{11}}%
\TextAt{8}{194.543}{117.97}{,}%
\TextAt{8}{201.693}{117.97}{we}%
\TextAt{8}{218.92}{117.97}{get}%
\TextAt{20}{238.484}{102.331}{\Character{32}}%
\TextAt{13}{247.949}{89.991}{\Character{11}}%
\TextAt{13}{278.925}{89.991}{\Character{13}}%
\Vector{q 0 G 0 g 0.478 w 176.349 -694.234 m 183.071 -694.234 l S Q}%
\TextAt{13}{248.349}{72.945}{\Character{13}}%
\TextAt{13}{265.421}{72.945}{\Character{12}}%
\TextAt{9}{275.361}{72.945}{\Character{43}}%
\TextAt{21}{288.318}{78.2}{\Character{13}}%
\Vector{q 0 G 0 g 0.359 w 221.15 -691.113 m 225.982 -691.113 l S Q}%
\TextAt{21}{293.15}{78.2}{\Character{13}}%
\Vector{q 0 G 0 g 0.478 w 216.318 -698.067 m 225.982 -698.067 l S Q}%
\TextAt{21}{290.431}{68.822}{\Character{11}}%
\TextAt{20}{299.178}{102.331}{\Character{33}}%
\NativeMathAt{311.963}{81.888}{11.955}{\sim}%
\TextAt{20}{324.582}{102.331}{\Character{32}}%
\TextAt{13}{334.047}{89.991}{\Character{11}}%
\TextAt{9}{352.241}{89.991}{0}%
\TextAt{9}{334.881}{72.945}{0}%
\TextAt{13}{351.531}{72.945}{\Character{12}}%
\TextAt{20}{358.803}{102.331}{\Character{33}}%
\TextAt{13}{370.26}{81.888}{\Character{58}}%
\TextAt{8}{297.211}{43.033}{221}%
\Destination{522.36}{499.335}{nav.AMS.33}%
\Destination{163.62}{499.335}{nav.equation.3.5.1}%
\Destination{78.23}{809.815}{nav.page.221}%
\Destination{79.36}{552.455}{nav.print.241}%
}

% Page 234 (body page 222)
\Page{%
\TextAt{2}{79.226}{796.943}{III.}%
\TextAt{2}{98.317}{796.943}{Riemann–Roch}%
\TextAt{2}{174.375}{796.943}{Theory}%
\Vector{q 0 G 0 g 0.299 w 7.226 18.729 m 460.769 18.729 l S Q}%
\TextAt{8}{79.226}{761.124}{Applying}%
\TextAt{8}{129.635}{761.124}{the}%
\TextAt{8}{149.474}{761.124}{same}%
\TextAt{8}{178.483}{761.124}{procedure}%
\TextAt{8}{232.467}{761.124}{to}%
\TextAt{8}{246.46}{761.124}{the}%
\TextAt{8}{266.3}{761.124}{exact}%
\TextAt{8}{296.875}{761.124}{sequence}%
\TextAt{9}{345.06}{761.124}{0}%
\NativeMathAt{354.234}{761.124}{11.955}{\rightarrow}%
\TextAt{14}{369.51}{761.124}{b}%
\TextAt{21}{375.645}{759.33}{f}%
\NativeMathAt{384.41}{761.124}{11.955}{\rightarrow}%
\TextAt{14}{399.686}{761.124}{a}%
\TextAt{21}{405.67}{759.33}{f}%
\TextAt{16}{413.123}{761.124}{\Character{8}}%
\TextAt{14}{424.43}{761.124}{b}%
\TextAt{21}{430.565}{759.33}{f}%
\NativeMathAt{439.331}{761.124}{11.955}{\rightarrow}%
\TextAt{14}{454.607}{761.124}{a}%
\TextAt{21}{460.59}{759.33}{f}%
\NativeMathAt{469.355}{761.124}{11.955}{\rightarrow}%
\TextAt{9}{484.631}{761.124}{0}%
\TextAt{8}{494.065}{761.124}{and}%
\TextAt{8}{516.51}{761.124}{the}%
\TextAt{8}{79.226}{735.863}{metric}%
\TextAt{20}{115.642}{756.306}{\Character{32}}%
\TextAt{13}{125.106}{743.965}{\Character{11}}%
\TextAt{3}{132.628}{748.304}{\Character{48}}%
\TextAt{13}{147.058}{743.965}{\Character{13}}%
\Vector{q 0 G 0 g 0.478 w 54.903 -40.26 m 61.626 -40.26 l S Q}%
\TextAt{13}{126.903}{726.919}{\Character{13}}%
\TextAt{13}{145.374}{726.919}{\Character{12}}%
\TextAt{3}{152.657}{731.258}{\Character{48}}%
\TextAt{20}{155.452}{756.306}{\Character{33}}%
\TextAt{8}{168.814}{735.863}{on}%
\TextAt{14}{185.073}{735.863}{a}%
\TextAt{21}{191.056}{734.069}{f}%
\TextAt{16}{199.158}{735.863}{\Character{8}}%
\TextAt{14}{211.113}{735.863}{b}%
\TextAt{21}{217.248}{734.069}{f}%
\TextAt{8}{222.692}{735.863}{,}%
\TextAt{8}{229.842}{735.863}{we}%
\TextAt{8}{247.069}{735.863}{obtain}%
\TextAt{20}{232.894}{709.726}{\Character{32}}%
\TextAt{13}{242.359}{697.664}{\Character{11}}%
\TextAt{3}{249.88}{702.003}{\Character{48}}%
\TextAt{9}{255.332}{697.664}{\Character{43}}%
\TextAt{21}{268.289}{702.92}{\Character{13}}%
\Vector{q 0 G 0 g 0.359 w 201.121 -66.393 m 205.953 -66.393 l S Q}%
\TextAt{21}{273.121}{702.92}{\Character{13}}%
\Vector{q 0 G 0 g 0.478 w 196.289 -73.347 m 205.953 -73.347 l S Q}%
\TextAt{21}{269.198}{693.542}{\Character{12}}%
\TextAt{23}{274.36}{695.811}{\Character{48}}%
\TextAt{13}{290.783}{697.664}{\Character{13}}%
\Vector{q 0 G 0 g 0.478 w 185.393 -87.119 m 192.115 -87.119 l S Q}%
\TextAt{13}{257.393}{680.06}{\Character{13}}%
\TextAt{13}{289.098}{680.06}{\Character{12}}%
\TextAt{3}{296.383}{684.398}{\Character{48}}%
\TextAt{20}{299.178}{709.726}{\Character{33}}%
\NativeMathAt{311.963}{689.282}{11.955}{\sim}%
\TextAt{20}{324.582}{709.726}{\Character{32}}%
\TextAt{13}{334.047}{697.385}{\Character{11}}%
\TextAt{3}{341.569}{701.723}{\Character{48}}%
\TextAt{9}{356.433}{697.385}{0}%
\TextAt{9}{336.279}{680.339}{0}%
\TextAt{13}{354.326}{680.339}{\Character{12}}%
\TextAt{3}{361.598}{684.677}{\Character{48}}%
\TextAt{20}{364.393}{709.726}{\Character{33}}%
\TextAt{13}{375.85}{689.282}{\Character{58}}%
\TextAt{8}{79.226}{655.207}{Choosing}%
\TextAt{13}{227.027}{635.37}{\Character{12}}%
\TextAt{3}{234.298}{640.306}{\Character{48}}%
\TextAt{9}{240.414}{635.37}{\Character{61}}%
\TextAt{13}{252.839}{635.37}{\Character{12}}%
\TextAt{9}{262.767}{635.37}{\Character{43}}%
\TextAt{13}{275.724}{643.458}{\Character{13}}%
\Vector{q 0 G 0 g 0.478 w 210.446 -123.721 m 217.169 -123.721 l S Q}%
\TextAt{13}{282.446}{643.458}{\Character{13}}%
\Vector{q 0 G 0 g 0.478 w 203.724 -135.641 m 217.169 -135.641 l S Q}%
\TextAt{13}{278.685}{627.169}{\Character{11}}%
\TextAt{13}{290.364}{635.37}{\Character{59}}%
\TextAt{13}{319.499}{635.37}{\Character{11}}%
\TextAt{3}{327.04}{640.306}{\Character{48}}%
\TextAt{9}{333.156}{635.37}{\Character{61}}%
\TextAt{13}{356.253}{643.458}{\Character{11}\Character{12}}%
\Vector{q 0 G 0 g 0.478 w 274.777 -135.641 m 308.522 -135.641 l S Q}%
\TextAt{13}{346.777}{626.205}{\Character{12}}%
\TextAt{9}{356.705}{626.205}{\Character{43}}%
\TextAt{21}{369.662}{631.46}{\Character{13}}%
\Vector{q 0 G 0 g 0.359 w 302.494 -137.853 m 307.326 -137.853 l S Q}%
\TextAt{21}{374.494}{631.46}{\Character{13}}%
\Vector{q 0 G 0 g 0.478 w 297.662 -144.807 m 307.326 -144.807 l S Q}%
\TextAt{21}{371.775}{622.082}{\Character{11}}%
\TextAt{13}{381.717}{635.37}{\Character{59}}%
\TextAt{8}{79.226}{607.152}{makes}%
\TextAt{8}{114.41}{607.152}{the}%
\TextAt{8}{134.567}{607.152}{matrices}%
\TextAt{8}{181.455}{607.152}{on}%
\TextAt{8}{197.714}{607.152}{the}%
\TextAt{8}{217.87}{607.152}{left}%
\TextAt{8}{238.349}{607.152}{equal,}%
\TextAt{8}{272.493}{607.152}{and}%
\TextAt{8}{295.256}{607.152}{yields}%
\TextAt{20}{223.181}{588.092}{\Character{32}}%
\TextAt{13}{232.646}{575.751}{\Character{11}}%
\TextAt{9}{250.839}{575.751}{0}%
\TextAt{9}{233.48}{558.705}{0}%
\TextAt{13}{250.13}{558.705}{\Character{12}}%
\TextAt{20}{257.401}{588.092}{\Character{33}}%
\NativeMathAt{270.187}{567.648}{11.955}{\sim}%
\TextAt{20}{282.806}{591.678}{\Character{48}}%
\TextAt{20}{282.806}{570.159}{\Character{64}}%
\TextAt{13}{303.938}{585.697}{\Character{11}\Character{12}}%
\Vector{q 0 G 0 g 0.478 w 222.462 -193.402 m 256.207 -193.402 l S Q}%
\TextAt{13}{294.462}{568.443}{\Character{12}}%
\TextAt{9}{304.39}{568.443}{\Character{43}}%
\TextAt{21}{317.347}{573.699}{\Character{13}}%
\Vector{q 0 G 0 g 0.359 w 250.179 -195.614 m 255.011 -195.614 l S Q}%
\TextAt{21}{322.179}{573.699}{\Character{13}}%
\Vector{q 0 G 0 g 0.478 w 245.347 -202.568 m 255.011 -202.568 l S Q}%
\TextAt{21}{319.46}{564.321}{\Character{11}}%
\TextAt{9}{353.311}{577.609}{0}%
\TextAt{9}{308.408}{552.389}{0}%
\TextAt{13}{339.365}{552.389}{\Character{12}}%
\TextAt{9}{349.293}{552.389}{\Character{43}}%
\TextAt{21}{362.25}{557.644}{\Character{13}}%
\Vector{q 0 G 0 g 0.359 w 295.082 -211.669 m 299.914 -211.669 l S Q}%
\TextAt{21}{367.082}{557.644}{\Character{13}}%
\Vector{q 0 G 0 g 0.478 w 290.25 -218.623 m 299.914 -218.623 l S Q}%
\TextAt{21}{364.362}{548.266}{\Character{11}}%
\TextAt{20}{373.11}{591.678}{\Character{49}}%
\TextAt{20}{373.11}{570.159}{\Character{65}}%
\TextAt{13}{385.563}{567.648}{\Character{59}}%
\TextAt{8}{79.226}{526.08}{or,}%
\TextAt{8}{96.788}{526.08}{if}%
\TextAt{8}{107.512}{526.08}{we}%
\TextAt{8}{124.739}{526.08}{put}%
\TextAt{13}{146.199}{526.08}{\Character{14}}%
\TextAt{9}{155.153}{526.08}{\Character{61}}%
\TextAt{13}{167.578}{526.08}{\Character{12}}%
\TextAt{9}{177.506}{526.08}{\Character{43}}%
\TextAt{21}{190.463}{531.335}{\Character{13}}%
\Vector{q 0 G 0 g 0.359 w 123.295 -237.978 m 128.127 -237.978 l S Q}%
\TextAt{21}{195.295}{531.335}{\Character{13}}%
\Vector{q 0 G 0 g 0.478 w 118.463 -244.932 m 128.127 -244.932 l S Q}%
\TextAt{21}{192.576}{521.957}{\Character{11}}%
\TextAt{8}{201.323}{526.08}{,}%
\TextAt{20}{249.695}{520.482}{\Character{32}}%
\TextAt{13}{259.16}{508.142}{\Character{11}}%
\TextAt{9}{277.353}{508.142}{0}%
\TextAt{9}{259.994}{491.096}{0}%
\TextAt{13}{276.644}{491.096}{\Character{12}}%
\TextAt{20}{283.916}{520.482}{\Character{33}}%
\NativeMathAt{296.701}{500.039}{11.955}{\sim}%
\TextAt{20}{309.32}{520.482}{\Character{32}}%
\TextAt{21}{319.98}{513.397}{\Character{11}\Character{12}}%
\Vector{q 0 G 0 g 0.478 w 247.98 -262.87 m 258.58 -262.87 l S Q}%
\TextAt{21}{323.257}{504.019}{\Character{14}}%
\TextAt{9}{341.739}{508.142}{0}%
\TextAt{9}{322.354}{491.096}{0}%
\TextAt{13}{341.849}{491.096}{\Character{14}}%
\TextAt{20}{347.592}{520.482}{\Character{33}}%
\TextAt{13}{359.049}{500.039}{\Character{59}}%
\TextAt{8}{511.944}{500.039}{(**)}%
\TextAt{8}{79.226}{469.615}{which}%
\TextAt{8}{112.93}{469.615}{is}%
\TextAt{8}{124.911}{469.615}{valid}%
\TextAt{8}{153.426}{469.615}{for}%
\TextAt{8}{171.528}{469.615}{any}%
\TextAt{13}{193.863}{469.615}{F}%
\TextAt{8}{203.066}{469.615}{-invariant}%
\TextAt{8}{256.305}{469.615}{function}%
\TextAt{13}{302.37}{469.615}{\Character{14}}%
\TextAt{9}{311.694}{469.615}{\Character{58}}%
\TextAt{13}{318.637}{469.615}{X}%
\TextAt{9}{329.292}{469.615}{\Character{40}}%
\TextAt{18}{333.844}{469.615}{C}%
\TextAt{9}{342.479}{469.615}{\Character{41}}%
\NativeMathAt{350.722}{469.615}{11.955}{\rightarrow}%
\TextAt{18}{366.368}{469.615}{R}%
\TextAt{8}{379.117}{469.615}{such}%
\TextAt{8}{405.732}{469.615}{that}%
\TextAt{13}{431.318}{469.615}{\Character{14}}%
\NativeMathAt{440.642}{469.615}{11.955}{\geq}%
\TextAt{13}{453.631}{469.615}{\Character{12}}%
\TextAt{8}{460.903}{469.615}{.}%
\TextAt{8}{470.003}{469.615}{This}%
\TextAt{8}{496.94}{469.615}{implies}%
\TextAt{8}{79.226}{452.569}{furthermore}%
\TextAt{20}{250.292}{450.975}{\Character{32}}%
\TextAt{21}{260.952}{443.89}{\Character{11}\Character{12}}%
\Vector{q 0 G 0 g 0.478 w 188.952 -332.377 m 199.552 -332.377 l S Q}%
\TextAt{21}{264.229}{434.512}{\Character{14}}%
\TextAt{9}{282.71}{438.634}{0}%
\TextAt{9}{263.325}{421.588}{0}%
\TextAt{13}{282.82}{421.588}{\Character{14}}%
\TextAt{20}{288.563}{450.975}{\Character{33}}%
\NativeMathAt{301.349}{430.532}{11.955}{\sim}%
\TextAt{20}{313.968}{450.975}{\Character{32}}%
\TextAt{21}{324.628}{443.89}{\Character{11}\Character{12}}%
\Vector{q 0 G 0 g 0.478 w 252.628 -332.377 m 263.228 -332.377 l S Q}%
\TextAt{21}{327.973}{434.512}{\Character{34}}%
\TextAt{9}{346.386}{438.634}{0}%
\TextAt{9}{327.002}{421.588}{0}%
\TextAt{13}{346.574}{421.588}{\Character{34}}%
\TextAt{20}{352.239}{450.975}{\Character{33}}%
\TextAt{8}{506.086}{430.532}{(***)}%
\TextAt{8}{79.226}{400.108}{for}%
\TextAt{8}{95.929}{400.108}{any}%
\TextAt{8}{116.863}{400.108}{two}%
\TextAt{13}{137.798}{400.108}{F}%
\TextAt{8}{147.002}{400.108}{-invariant}%
\TextAt{8}{198.841}{400.108}{functions}%
\TextAt{13}{248.121}{400.108}{\Character{14}}%
\TextAt{13}{253.106}{400.108}{\Character{59}}%
\TextAt{13}{258.319}{400.108}{\Character{34}}%
\TextAt{9}{267.147}{400.108}{\Character{58}}%
\TextAt{13}{273.719}{400.108}{X}%
\TextAt{9}{284.374}{400.108}{\Character{40}}%
\TextAt{18}{288.927}{400.108}{C}%
\TextAt{9}{297.561}{400.108}{\Character{41}}%
\NativeMathAt{305.434}{400.108}{11.955}{\rightarrow}%
\TextAt{18}{320.71}{400.108}{R}%
\TextAt{3}{329.344}{404.446}{\Character{3}}%
\TextAt{17}{329.344}{397.152}{\Character{43}}%
\TextAt{8}{336.429}{400.108}{.}%
\TextAt{8}{344.483}{400.108}{For}%
\TextAt{8}{364.27}{400.108}{if}%
\TextAt{13}{373.812}{400.108}{\Character{20}}%
\TextAt{9}{383.871}{400.108}{\Character{58}}%
\TextAt{13}{390.444}{400.108}{X}%
\TextAt{9}{401.099}{400.108}{\Character{40}}%
\TextAt{18}{405.651}{400.108}{C}%
\TextAt{9}{414.285}{400.108}{\Character{41}}%
\NativeMathAt{422.158}{400.108}{11.955}{\rightarrow}%
\TextAt{18}{437.434}{400.108}{R}%
\TextAt{8}{448.784}{400.108}{is}%
\TextAt{8}{459.365}{400.108}{an}%
\TextAt{13}{474.442}{400.108}{F}%
\TextAt{8}{483.646}{400.108}{-invariant}%
\TextAt{8}{79.226}{383.062}{function}%
\TextAt{8}{125.074}{383.062}{such}%
\TextAt{8}{151.471}{383.062}{that}%
\TextAt{13}{176.84}{383.062}{\Character{20}}%
\NativeMathAt{186.899}{383.062}{11.955}{\geq}%
\TextAt{13}{199.519}{383.062}{\Character{14}}%
\TextAt{8}{205.152}{383.062}{,}%
\TextAt{13}{212.301}{383.062}{\Character{20}}%
\NativeMathAt{222.36}{383.062}{11.955}{\geq}%
\TextAt{13}{234.98}{383.062}{\Character{34}}%
\TextAt{8}{240.458}{383.062}{,}%
\TextAt{8}{247.607}{383.062}{then}%
\TextAt{8}{274.267}{383.062}{(**)}%
\TextAt{8}{298.99}{383.062}{gives}%
\TextAt{20}{215.389}{368.868}{\Character{32}}%
\TextAt{21}{226.049}{361.783}{\Character{11}\Character{12}}%
\Vector{q 0 G 0 g 0.478 w 154.049 -414.484 m 164.649 -414.484 l S Q}%
\TextAt{21}{229.326}{352.405}{\Character{14}}%
\TextAt{9}{247.807}{356.527}{0}%
\TextAt{9}{228.423}{339.481}{0}%
\TextAt{13}{247.917}{339.482}{\Character{14}}%
\TextAt{20}{253.66}{368.868}{\Character{33}}%
\NativeMathAt{266.446}{348.425}{11.955}{\sim}%
\TextAt{20}{279.065}{368.868}{\Character{32}}%
\TextAt{21}{289.725}{361.783}{\Character{11}\Character{12}}%
\Vector{q 0 G 0 g 0.478 w 217.725 -414.484 m 228.325 -414.484 l S Q}%
\TextAt{21}{292.585}{352.405}{\Character{20}}%
\TextAt{9}{311.926}{356.527}{0}%
\TextAt{9}{292.099}{339.481}{0}%
\TextAt{13}{311.483}{339.482}{\Character{20}}%
\TextAt{20}{318.222}{368.868}{\Character{33}}%
\NativeMathAt{331.007}{348.425}{11.955}{\sim}%
\TextAt{20}{343.627}{368.868}{\Character{32}}%
\TextAt{21}{354.287}{361.783}{\Character{11}\Character{12}}%
\Vector{q 0 G 0 g 0.478 w 282.287 -414.484 m 292.887 -414.484 l S Q}%
\TextAt{21}{357.632}{352.405}{\Character{34}}%
\TextAt{9}{376.045}{356.527}{0}%
\TextAt{9}{356.66}{339.481}{0}%
\TextAt{13}{376.233}{339.482}{\Character{34}}%
\TextAt{20}{381.898}{368.868}{\Character{33}}%
\TextAt{13}{393.355}{348.425}{\Character{58}}%
\TextAt{8}{79.226}{313.8}{Now}%
\TextAt{8}{105.886}{313.8}{putting}%
\TextAt{13}{147.514}{313.8}{\Character{14}}%
\TextAt{9}{156.468}{313.8}{\Character{61}}%
\TextAt{13}{168.893}{313.8}{\Character{12}}%
\TextAt{8}{180.062}{313.8}{and}%
\TextAt{13}{202.825}{313.8}{\Character{34}}%
\TextAt{9}{211.624}{313.8}{\Character{61}}%
\TextAt{9}{224.034}{313.8}{1}%
\TextAt{8}{233.799}{313.8}{in}%
\TextAt{8}{247.452}{313.8}{(***),}%
\TextAt{8}{281.285}{313.8}{we}%
\TextAt{8}{298.513}{313.8}{find}%
\TextAt{9}{209.75}{287.035}{\Character{91}\Character{40}}%
\TextAt{14}{217.554}{287.035}{a}%
\TextAt{21}{223.537}{285.242}{f}%
\TextAt{13}{228.982}{287.035}{\Character{59}}%
\TextAt{13}{234.206}{287.035}{\Character{11}}%
\TextAt{9}{241.748}{287.035}{\Character{41}\Character{93}}%
\TextAt{9}{252.185}{287.035}{\Character{43}}%
\TextAt{9}{263.937}{287.035}{\Character{91}\Character{40}}%
\TextAt{14}{271.774}{287.035}{b}%
\TextAt{21}{277.909}{285.242}{f}%
\TextAt{13}{283.354}{287.035}{\Character{59}}%
\TextAt{13}{288.578}{287.035}{\Character{12}}%
\TextAt{9}{295.869}{287.035}{\Character{41}\Character{93}}%
\TextAt{9}{306.963}{287.035}{\Character{61}}%
\TextAt{9}{319.385}{287.035}{\Character{91}\Character{40}}%
\TextAt{14}{327.224}{287.035}{a}%
\TextAt{21}{333.207}{285.242}{f}%
\TextAt{13}{338.652}{287.035}{\Character{59}}%
\TextAt{13}{343.876}{287.035}{\Character{11}\Character{12}}%
\TextAt{9}{358.689}{287.035}{\Character{41}\Character{93}}%
\TextAt{9}{369.126}{287.035}{\Character{43}}%
\TextAt{9}{380.878}{287.035}{\Character{91}}%
\TextAt{14}{384.162}{287.035}{b}%
\TextAt{21}{390.297}{285.242}{f}%
\TextAt{9}{395.742}{287.035}{\Character{93}}%
\TextAt{13}{398.994}{287.035}{\Character{58}}%
\TextAt{8}{79.226}{260.271}{For}%
\TextAt{8}{100.196}{260.271}{the}%
\TextAt{8}{120.352}{260.271}{replete}%
\TextAt{8}{158.716}{260.271}{ideals}%
\TextAt{14}{191.294}{260.271}{a}%
\TextAt{9}{200.598}{260.271}{\Character{61}}%
\TextAt{20}{213.023}{269.238}{\Character{81}}%
\TextAt{27}{224.314}{256.784}{p}%
\TextAt{14}{230.794}{260.271}{p}%
\TextAt{21}{236.777}{264.61}{\Character{23}}%
\TextAt{28}{240.939}{263.613}{p}%
\TextAt{8}{245.52}{260.271}{,}%
\TextAt{14}{252.67}{260.271}{b}%
\TextAt{9}{262.125}{260.271}{\Character{61}}%
\TextAt{20}{274.551}{269.238}{\Character{81}}%
\TextAt{27}{285.842}{256.784}{p}%
\TextAt{14}{292.321}{260.271}{p}%
\TextAt{21}{298.304}{264.61}{\Character{22}}%
\TextAt{28}{303.398}{263.613}{p}%
\TextAt{8}{307.979}{260.271}{,}%
\TextAt{8}{315.129}{260.271}{this}%
\TextAt{8}{337.951}{260.271}{means}%
\TextAt{9}{246.683}{233.507}{\Character{91}}%
\TextAt{14}{249.935}{233.507}{a}%
\TextAt{9}{255.918}{233.507}{\Character{93}}%
\TextAt{9}{261.812}{233.507}{\Character{43}}%
\TextAt{9}{273.564}{233.507}{\Character{91}}%
\TextAt{14}{276.839}{233.507}{b}%
\TextAt{9}{282.974}{233.507}{\Character{93}}%
\TextAt{9}{289.525}{233.507}{\Character{61}}%
\TextAt{9}{301.947}{233.507}{\Character{91}}%
\TextAt{14}{305.224}{233.507}{ab}%
\NativeMathAt{317.342}{231.714}{7.97}{\infty}%
\TextAt{9}{326.308}{233.507}{\Character{93}}%
\TextAt{9}{332.202}{233.507}{\Character{43}}%
\TextAt{9}{343.954}{233.507}{\Character{91}}%
\TextAt{14}{347.229}{233.507}{b}%
\TextAt{21}{353.364}{231.714}{f}%
\TextAt{9}{358.809}{233.507}{\Character{93}}%
\TextAt{13}{362.061}{233.507}{\Character{59}}%
\TextAt{8}{517.802}{233.507}{(1)}%
\TextAt{8}{79.226}{206.743}{for}%
\TextAt{8}{97.111}{206.743}{if}%
\TextAt{8}{107.835}{206.743}{we}%
\TextAt{8}{125.062}{206.743}{put}%
\TextAt{13}{146.522}{206.743}{\Character{11}}%
\TextAt{9}{154.044}{206.743}{\Character{40}}%
\TextAt{13}{158.596}{206.743}{\Character{27}}%
\TextAt{9}{165.678}{206.743}{\Character{41}}%
\TextAt{9}{173.533}{206.743}{\Character{61}}%
\TextAt{13}{185.977}{206.743}{e}%
\TextAt{17}{191.402}{211.081}{2}%
\TextAt{21}{195.637}{211.081}{\Character{23}}%
\TextAt{28}{199.799}{210.085}{p}%
\TextAt{22}{203.384}{209.089}{\Character{27}}%
\TextAt{8}{209.376}{206.743}{,}%
\TextAt{13}{216.525}{206.743}{\Character{12}}%
\TextAt{9}{223.796}{206.743}{\Character{40}}%
\TextAt{13}{228.349}{206.743}{\Character{27}}%
\TextAt{9}{235.431}{206.743}{\Character{41}}%
\TextAt{9}{243.286}{206.743}{\Character{61}}%
\TextAt{13}{255.73}{206.743}{e}%
\TextAt{17}{261.155}{211.081}{2}%
\TextAt{21}{265.389}{211.081}{\Character{22}}%
\TextAt{28}{270.484}{210.085}{p}%
\TextAt{22}{274.068}{209.089}{\Character{27}}%
\TextAt{8}{280.06}{206.743}{,}%
\TextAt{8}{287.21}{206.743}{then}%
\TextAt{8}{313.87}{206.743}{we}%
\TextAt{8}{331.097}{206.743}{have}%
\TextAt{9}{167.814}{179.978}{\Character{40}}%
\TextAt{14}{172.367}{179.978}{a}%
\TextAt{21}{178.35}{178.185}{f}%
\TextAt{13}{183.795}{179.978}{\Character{59}}%
\TextAt{13}{189.019}{179.978}{\Character{11}}%
\TextAt{9}{196.56}{179.978}{\Character{41}}%
\TextAt{9}{204.415}{179.978}{\Character{61}}%
\TextAt{13}{216.859}{179.978}{L}%
\TextAt{9}{224.824}{179.978}{\Character{40}}%
\TextAt{14}{229.376}{179.978}{a}%
\TextAt{9}{235.359}{179.978}{\Character{41}}%
\TextAt{13}{239.911}{179.978}{\Character{59}}%
\TextAt{9}{257.111}{179.978}{\Character{40}}%
\TextAt{14}{261.663}{179.978}{b}%
\TextAt{21}{267.798}{178.185}{f}%
\TextAt{13}{273.243}{179.978}{\Character{59}}%
\TextAt{13}{278.467}{179.978}{\Character{12}}%
\TextAt{9}{285.758}{179.978}{\Character{41}}%
\TextAt{9}{293.613}{179.978}{\Character{61}}%
\TextAt{13}{306.057}{179.978}{L}%
\TextAt{9}{314.021}{179.978}{\Character{40}}%
\TextAt{14}{318.574}{179.978}{b}%
\TextAt{9}{324.708}{179.978}{\Character{41}}%
\TextAt{13}{329.261}{179.978}{\Character{59}}%
\TextAt{9}{346.46}{179.978}{\Character{40}}%
\TextAt{14}{351.012}{179.978}{a}%
\TextAt{21}{356.995}{178.185}{f}%
\TextAt{13}{362.44}{179.978}{\Character{59}}%
\TextAt{13}{367.664}{179.978}{\Character{11}\Character{12}}%
\TextAt{9}{382.478}{179.978}{\Character{41}}%
\TextAt{9}{390.333}{179.978}{\Character{61}}%
\TextAt{13}{402.776}{179.978}{L}%
\TextAt{9}{410.741}{179.978}{\Character{40}}%
\TextAt{14}{415.293}{179.978}{ab}%
\NativeMathAt{427.411}{178.185}{7.97}{\infty}%
\TextAt{9}{436.377}{179.978}{\Character{41}}%
\TextAt{13}{440.93}{179.978}{\Character{58}}%
\TextAt{8}{79.226}{153.214}{On}%
\TextAt{8}{98.737}{153.214}{the}%
\TextAt{8}{118.893}{153.214}{other}%
\TextAt{8}{149.463}{153.214}{hand,}%
\TextAt{8}{181.981}{153.214}{we}%
\TextAt{8}{199.208}{153.214}{obtain}%
\TextAt{8}{235.636}{153.214}{the}%
\TextAt{8}{255.792}{153.214}{formula}%
\TextAt{9}{253.462}{126.45}{\Character{91}}%
\TextAt{14}{256.714}{126.45}{a}%
\TextAt{9}{262.697}{126.45}{\Character{93}}%
\TextAt{9}{268.591}{126.45}{\Character{43}}%
\TextAt{9}{280.343}{126.45}{\Character{91}}%
\TextAt{14}{283.618}{126.45}{b}%
\TextAt{21}{289.753}{124.657}{f}%
\TextAt{9}{295.198}{126.45}{\Character{93}}%
\TextAt{9}{301.749}{126.45}{\Character{61}}%
\TextAt{9}{314.171}{126.45}{\Character{91}}%
\TextAt{14}{317.448}{126.45}{ab}%
\TextAt{21}{329.566}{124.657}{f}%
\TextAt{9}{335.011}{126.45}{\Character{93}}%
\TextAt{9}{340.905}{126.45}{\Character{43}}%
\TextAt{9}{352.657}{126.45}{1}%
\TextAt{8}{517.802}{126.45}{(2)}%
\TextAt{8}{79.226}{99.686}{in}%
\TextAt{8}{92.879}{99.686}{the}%
\TextAt{8}{113.035}{99.686}{following}%
\TextAt{8}{162.47}{99.686}{manner.}%
\TextAt{8}{209.294}{99.686}{We}%
\TextAt{8}{229.439}{99.686}{have}%
\TextAt{8}{256.434}{99.686}{two}%
\TextAt{8}{278.551}{99.686}{exact}%
\TextAt{8}{309.443}{99.686}{sequences}%
\TextAt{8}{362.56}{99.686}{of}%
\TextAt{19}{375.89}{99.686}{coherent}%
\TextAt{8}{422.575}{99.686}{metrized}%
\TextAt{14}{470.694}{99.686}{o}%
\TextAt{8}{476.538}{99.686}{-modules:}%
\TextAt{9}{187.481}{72.921}{0}%
\NativeArrowAt{196.654}{72.921}{11.955}{19.26}{\longrightarrow}%
\TextAt{9}{219.236}{72.921}{\Character{40}}%
\TextAt{14}{223.789}{72.921}{a}%
\TextAt{21}{229.772}{71.128}{f}%
\TextAt{14}{235.217}{72.921}{b}%
\TextAt{21}{241.352}{71.128}{f}%
\TextAt{13}{246.796}{72.921}{\Character{59}}%
\TextAt{13}{252.02}{72.921}{\Character{11}}%
\TextAt{9}{259.562}{72.921}{\Character{41}}%
\NativeArrowAt{267.436}{72.921}{11.955}{19.26}{\longrightarrow}%
\TextAt{9}{290.018}{72.921}{\Character{40}}%
\TextAt{14}{294.57}{72.921}{a}%
\TextAt{21}{300.553}{71.128}{f}%
\TextAt{13}{305.998}{72.921}{\Character{59}}%
\TextAt{13}{311.222}{72.921}{\Character{11}}%
\TextAt{9}{318.764}{72.921}{\Character{41}}%
\NativeArrowAt{326.637}{72.921}{11.955}{19.26}{\longrightarrow}%
\TextAt{14}{349.219}{72.921}{a}%
\TextAt{21}{355.202}{71.128}{f}%
\TextAt{9}{360.647}{72.921}{\Character{47}}%
\TextAt{14}{366.5}{72.921}{a}%
\TextAt{21}{372.483}{71.128}{f}%
\TextAt{14}{377.928}{72.921}{b}%
\TextAt{21}{384.063}{71.128}{f}%
\NativeArrowAt{392.828}{72.921}{11.955}{19.26}{\longrightarrow}%
\TextAt{9}{415.41}{72.921}{0}%
\TextAt{13}{421.263}{72.921}{\Character{59}}%
\TextAt{8}{297.211}{43.033}{222}%
\Destination{516.5}{517.085}{nav.AMS.34}%
\Destination{510.64}{447.575}{nav.AMS.35}%
\Destination{522.36}{250.555}{nav.AMS.36}%
\Destination{522.36}{143.495}{nav.AMS.37}%
\Destination{78.23}{809.815}{nav.page.222}%
\Destination{79.36}{622.155}{nav.print.242}%
}

% Page 235 (body page 223)
\Page{%
\TextAt{2}{79.226}{796.943}{III.}%
\TextAt{2}{98.317}{796.943}{Riemann–Roch}%
\TextAt{2}{174.375}{796.943}{Theory}%
\Vector{q 0 G 0 g 0.299 w 7.226 18.729 m 460.769 18.729 l S Q}%
\TextAt{9}{206.161}{761.124}{0}%
\NativeArrowAt{215.335}{761.124}{11.955}{19.26}{\longrightarrow}%
\TextAt{9}{237.917}{761.124}{\Character{40}}%
\TextAt{14}{242.469}{761.124}{b}%
\TextAt{21}{248.604}{759.33}{f}%
\TextAt{13}{254.049}{761.124}{\Character{59}}%
\TextAt{9}{259.293}{761.124}{1\Character{41}}%
\NativeArrowAt{273.019}{761.124}{11.955}{19.26}{\longrightarrow}%
\TextAt{9}{295.601}{761.124}{\Character{40}}%
\TextAt{14}{300.154}{761.124}{o}%
\TextAt{13}{305.998}{761.124}{\Character{59}}%
\TextAt{9}{311.242}{761.124}{1\Character{41}}%
\NativeArrowAt{324.968}{761.124}{11.955}{19.26}{\longrightarrow}%
\TextAt{14}{347.55}{761.124}{o}%
\TextAt{9}{353.394}{761.124}{\Character{47}}%
\TextAt{14}{359.247}{761.124}{b}%
\TextAt{21}{365.382}{759.33}{f}%
\NativeArrowAt{374.148}{761.124}{11.955}{19.26}{\longrightarrow}%
\TextAt{9}{396.73}{761.124}{0}%
\TextAt{13}{402.583}{761.124}{\Character{58}}%
\TextAt{8}{79.226}{737.636}{As}%
\TextAt{14}{96.513}{737.636}{a}%
\TextAt{21}{102.496}{735.842}{f}%
\TextAt{8}{111.839}{737.636}{and}%
\TextAt{14}{134.601}{737.636}{b}%
\TextAt{21}{140.736}{735.842}{f}%
\TextAt{8}{150.078}{737.636}{are}%
\TextAt{8}{169.589}{737.636}{relatively}%
\TextAt{8}{220.65}{737.636}{prime,}%
\TextAt{8}{257.065}{737.636}{i.e.,}%
\TextAt{14}{279.17}{737.636}{a}%
\TextAt{21}{285.153}{735.842}{f}%
\TextAt{9}{293.255}{737.636}{\Character{43}}%
\TextAt{14}{305.016}{737.636}{b}%
\TextAt{21}{311.151}{735.842}{f}%
\TextAt{9}{319.917}{737.636}{\Character{61}}%
\TextAt{14}{332.342}{737.636}{o}%
\TextAt{8}{338.187}{737.636}{,}%
\TextAt{8}{345.336}{737.636}{it}%
\TextAt{8}{357.04}{737.636}{follows}%
\TextAt{8}{395.476}{737.636}{that}%
\TextAt{14}{261.264}{708.713}{a}%
\TextAt{21}{267.247}{706.92}{f}%
\TextAt{9}{272.692}{708.713}{\Character{47}}%
\TextAt{14}{278.545}{708.713}{a}%
\TextAt{21}{284.528}{706.92}{f}%
\TextAt{14}{289.973}{708.713}{b}%
\TextAt{21}{296.107}{706.92}{f}%
\NativeArrowAt{304.873}{708.713}{11.955}{19.26}{\longrightarrow}%
\TextAt{14}{327.455}{708.713}{o}%
\TextAt{9}{333.299}{708.713}{\Character{47}}%
\TextAt{14}{339.152}{708.713}{b}%
\TextAt{21}{345.287}{706.92}{f}%
\TextAt{8}{79.226}{679.79}{is}%
\TextAt{8}{91.675}{679.79}{an}%
\TextAt{8}{108.62}{679.79}{isomorphism,}%
\TextAt{8}{181.147}{679.79}{so}%
\TextAt{8}{196.203}{679.79}{that}%
\TextAt{8}{222.257}{679.79}{in}%
\TextAt{8}{236.595}{679.79}{the}%
\TextAt{8}{257.437}{679.79}{group}%
\TextAt{13}{291.297}{679.79}{K}%
\TextAt{17}{302.115}{684.129}{0}%
\TextAt{9}{306.847}{679.79}{\Character{40}}%
\Vector{q 0 G 0 g 0.478 w 239.4 -86.853 m 245.244 -86.853 l S Q}%
\TextAt{14}{311.4}{679.79}{o}%
\TextAt{9}{317.244}{679.79}{\Character{41}}%
\TextAt{8}{326.379}{679.79}{one}%
\TextAt{8}{348.524}{679.79}{has}%
\TextAt{8}{370.083}{679.79}{the}%
\TextAt{8}{390.924}{679.79}{identity}%
\TextAt{9}{434.864}{679.79}{\Character{91}}%
\TextAt{14}{438.115}{679.79}{a}%
\TextAt{21}{444.098}{677.997}{f}%
\TextAt{9}{449.543}{679.79}{\Character{47}}%
\TextAt{14}{455.396}{679.79}{a}%
\TextAt{21}{461.379}{677.997}{f}%
\TextAt{14}{466.824}{679.79}{b}%
\TextAt{21}{472.959}{677.997}{f}%
\TextAt{9}{478.404}{679.79}{\Character{93}}%
\TextAt{9}{486.127}{679.79}{\Character{61}}%
\TextAt{9}{499.708}{679.79}{\Character{91}}%
\TextAt{14}{502.989}{679.79}{o}%
\TextAt{9}{508.833}{679.79}{\Character{47}}%
\TextAt{14}{514.686}{679.79}{b}%
\TextAt{21}{520.821}{677.997}{f}%
\TextAt{9}{526.266}{679.79}{\Character{93}}%
\TextAt{8}{529.518}{679.79}{,}%
\TextAt{8}{79.226}{662.744}{and}%
\TextAt{8}{101.989}{662.744}{therefore}%
\TextAt{9}{200.878}{645.698}{\Character{91}\Character{40}}%
\TextAt{14}{208.682}{645.698}{a}%
\TextAt{21}{214.665}{643.905}{f}%
\TextAt{13}{220.11}{645.698}{\Character{59}}%
\TextAt{13}{225.334}{645.698}{\Character{11}}%
\TextAt{9}{232.876}{645.698}{\Character{41}\Character{93}}%
\TextAt{16}{243.337}{645.698}{\Character{0}}%
\TextAt{9}{255.292}{645.698}{\Character{91}\Character{40}}%
\TextAt{14}{263.096}{645.698}{a}%
\TextAt{21}{269.079}{643.905}{f}%
\TextAt{14}{274.524}{645.698}{b}%
\TextAt{21}{280.659}{643.905}{f}%
\TextAt{13}{286.104}{645.698}{\Character{59}}%
\TextAt{13}{291.328}{645.698}{\Character{11}}%
\TextAt{9}{298.869}{645.698}{\Character{41}\Character{93}}%
\TextAt{9}{309.963}{645.698}{\Character{61}}%
\TextAt{9}{322.385}{645.698}{\Character{91}\Character{40}}%
\TextAt{14}{330.224}{645.698}{o}%
\TextAt{13}{336.068}{645.698}{\Character{59}}%
\TextAt{9}{341.312}{645.698}{1\Character{41}\Character{93}}%
\TextAt{16}{357.626}{645.698}{\Character{0}}%
\TextAt{9}{369.581}{645.698}{\Character{91}\Character{40}}%
\TextAt{14}{377.385}{645.698}{b}%
\TextAt{21}{383.52}{643.905}{f}%
\TextAt{13}{388.965}{645.698}{\Character{59}}%
\TextAt{9}{394.209}{645.698}{1\Character{41}\Character{93}}%
\TextAt{13}{407.866}{645.698}{\Character{59}}%
\TextAt{8}{79.226}{622.21}{and}%
\TextAt{8}{101.989}{622.21}{so}%
\TextAt{9}{251.643}{605.164}{\Character{91}}%
\TextAt{14}{254.894}{605.164}{a}%
\TextAt{9}{260.877}{605.164}{\Character{93}}%
\TextAt{16}{266.786}{605.164}{\Character{0}}%
\TextAt{9}{278.741}{605.164}{\Character{91}}%
\TextAt{14}{281.993}{605.164}{ab}%
\TextAt{21}{294.111}{603.371}{f}%
\TextAt{9}{299.555}{605.164}{\Character{93}}%
\TextAt{9}{306.106}{605.164}{\Character{61}}%
\TextAt{9}{318.528}{605.164}{1}%
\TextAt{16}{327.063}{605.164}{\Character{0}}%
\TextAt{9}{339.018}{605.164}{\Character{91}}%
\TextAt{14}{342.27}{605.164}{b}%
\TextAt{21}{348.405}{603.371}{f}%
\TextAt{9}{353.85}{605.164}{\Character{93}}%
\TextAt{13}{357.101}{605.164}{\Character{58}}%
\TextAt{8}{79.226}{581.676}{From}%
\TextAt{8}{109.951}{581.676}{(1)}%
\TextAt{8}{128.816}{581.676}{and}%
\TextAt{8}{151.579}{581.676}{(2)}%
\TextAt{8}{170.444}{581.676}{it}%
\TextAt{8}{182.148}{581.676}{now}%
\TextAt{8}{206.537}{581.676}{follows}%
\TextAt{8}{244.973}{581.676}{that}%
\TextAt{9}{181.772}{552.754}{\Character{91}}%
\TextAt{14}{185.023}{552.754}{a}%
\TextAt{9}{191.006}{552.754}{\Character{93}}%
\TextAt{9}{196.9}{552.754}{\Character{43}}%
\TextAt{9}{208.652}{552.754}{\Character{91}}%
\TextAt{14}{211.928}{552.754}{b}%
\TextAt{9}{218.062}{552.754}{\Character{93}}%
\TextAt{9}{224.613}{552.754}{\Character{61}}%
\TextAt{9}{237.035}{552.754}{\Character{91}}%
\TextAt{14}{240.312}{552.754}{ab}%
\NativeMathAt{252.43}{550.961}{7.97}{\infty}%
\TextAt{9}{261.396}{552.754}{\Character{93}}%
\TextAt{9}{267.29}{552.754}{\Character{43}}%
\TextAt{9}{279.042}{552.754}{\Character{91}}%
\TextAt{14}{282.318}{552.754}{b}%
\TextAt{21}{288.453}{550.961}{f}%
\TextAt{9}{293.898}{552.754}{\Character{93}}%
\TextAt{9}{300.449}{552.754}{\Character{61}}%
\TextAt{9}{312.871}{552.754}{\Character{91}}%
\TextAt{14}{316.147}{552.754}{ab}%
\NativeMathAt{328.265}{550.961}{7.97}{\infty}%
\TextAt{14}{337.232}{552.754}{b}%
\TextAt{21}{343.366}{550.961}{f}%
\TextAt{9}{348.811}{552.754}{\Character{93}}%
\TextAt{9}{354.705}{552.754}{\Character{43}}%
\TextAt{9}{366.457}{552.754}{1}%
\TextAt{9}{375.627}{552.754}{\Character{61}}%
\TextAt{9}{388.048}{552.754}{\Character{91}}%
\TextAt{14}{391.332}{552.754}{ab}%
\TextAt{9}{403.45}{552.754}{\Character{93}}%
\TextAt{9}{409.344}{552.754}{\Character{43}}%
\TextAt{9}{421.096}{552.754}{1}%
\TextAt{13}{426.972}{552.754}{\Character{58}}%
\TextAt{8}{79.226}{523.831}{In}%
\TextAt{8}{93.847}{523.831}{view}%
\TextAt{8}{120.83}{523.831}{of}%
\TextAt{8}{134.16}{523.831}{the}%
\TextAt{8}{154.317}{523.831}{isomorphism}%
\TextAt{13}{222.736}{523.831}{K}%
\TextAt{17}{232.695}{522.038}{0}%
\TextAt{9}{237.427}{523.831}{\Character{40}}%
\Vector{q 0 G 0 g 0.478 w 169.98 -242.812 m 175.824 -242.812 l S Q}%
\TextAt{14}{241.98}{523.831}{o}%
\TextAt{9}{247.824}{523.831}{\Character{41}}%
\NativeMathAt{255.697}{523.831}{11.955}{\cong}%
\TextAt{13}{268.316}{523.831}{K}%
\TextAt{17}{279.134}{528.17}{0}%
\TextAt{9}{283.866}{523.831}{\Character{40}}%
\Vector{q 0 G 0 g 0.478 w 216.419 -242.812 m 222.263 -242.812 l S Q}%
\TextAt{14}{288.419}{523.831}{o}%
\TextAt{9}{294.263}{523.831}{\Character{41}}%
\TextAt{8}{298.815}{523.831}{,}%
\TextAt{8}{305.965}{523.831}{this}%
\TextAt{8}{328.787}{523.831}{is}%
\TextAt{8}{340.551}{523.831}{indeed}%
\TextAt{8}{377.612}{523.831}{an}%
\TextAt{8}{393.871}{523.831}{identity}%
\TextAt{8}{437.125}{523.831}{in}%
\TextAt{13}{450.778}{523.831}{K}%
\TextAt{17}{460.737}{522.038}{0}%
\TextAt{9}{465.469}{523.831}{\Character{40}}%
\Vector{q 0 G 0 g 0.478 w 398.021 -242.812 m 403.866 -242.812 l S Q}%
\TextAt{14}{470.021}{523.831}{o}%
\TextAt{9}{475.866}{523.831}{\Character{41}}%
\TextAt{8}{480.418}{523.831}{.}%
\Vector{q 0 G 0 g 0.399 w 452.334 -250.169 m 452.334 -242.099 l S Q}%
\Vector{q 0 G 0 g 0.399 w 452.534 -242.298 m 459.707 -242.298 l S Q}%
\Vector{q 0 G 0 g 0.399 w 452.534 -249.97 m 459.707 -249.97 l S Q}%
\Vector{q 0 G 0 g 0.399 w 459.906 -250.169 m 459.906 -242.099 l S Q}%
\TextAt{15}{79.226}{476.957}{§}%
\TextAt{15}{93.785}{476.957}{6}%
\TextAt{15}{120.693}{476.957}{The}%
\TextAt{15}{160.202}{476.957}{Chern}%
\TextAt{15}{218.924}{476.957}{Character}%
\TextAt{8}{79.226}{445.071}{The}%
\TextAt{8}{103.28}{445.071}{Grothendieck}%
\TextAt{8}{175.549}{445.071}{ring}%
\TextAt{13}{199.615}{445.071}{K}%
\TextAt{17}{209.574}{443.278}{0}%
\TextAt{9}{214.306}{445.071}{\Character{40}}%
\Vector{q 0 G 0 g 0.478 w 146.858 -321.573 m 152.703 -321.573 l S Q}%
\TextAt{14}{218.858}{445.071}{o}%
\TextAt{9}{224.703}{445.071}{\Character{41}}%
\TextAt{8}{233.152}{445.071}{is}%
\TextAt{8}{244.916}{445.071}{equipped}%
\TextAt{8}{294.985}{445.071}{with}%
\TextAt{8}{321.645}{445.071}{a}%
\TextAt{8}{331.4}{445.071}{canonical}%
\TextAt{8}{382.783}{445.071}{surjective}%
\TextAt{8}{435.195}{445.071}{homomorphism}%
\TextAt{8}{262.301}{416.148}{rk}%
\TextAt{9}{276.357}{416.148}{\Character{58}}%
\TextAt{13}{282.93}{416.148}{K}%
\TextAt{17}{292.889}{414.355}{0}%
\TextAt{9}{297.621}{416.148}{\Character{40}}%
\Vector{q 0 G 0 g 0.478 w 230.174 -350.495 m 236.018 -350.495 l S Q}%
\TextAt{14}{302.174}{416.148}{o}%
\TextAt{9}{308.018}{416.148}{\Character{41}}%
\NativeArrowAt{315.891}{416.148}{11.955}{19.26}{\longrightarrow}%
\TextAt{18}{338.473}{416.148}{Z}%
\TextAt{13}{346.443}{416.148}{\Character{58}}%
\TextAt{8}{79.226}{387.226}{Indeed,}%
\TextAt{8}{122.64}{387.226}{the}%
\TextAt{8}{144.503}{387.226}{rule}%
\TextAt{8}{169.617}{387.226}{which}%
\TextAt{8}{204.81}{387.226}{associates}%
\TextAt{8}{260.362}{387.226}{to}%
\TextAt{8}{276.379}{387.226}{every}%
\TextAt{8}{308.977}{387.226}{isometry}%
\TextAt{8}{358.552}{387.226}{class}%
\TextAt{16}{387.695}{387.226}{\Character{102}}%
\TextAt{13}{393.673}{387.226}{M}%
\TextAt{16}{406.246}{387.226}{\Character{103}}%
\TextAt{8}{417.827}{387.226}{of}%
\TextAt{8}{432.864}{387.226}{projective}%
\TextAt{8}{488.547}{387.226}{metrized}%
\TextAt{14}{79.226}{370.18}{o}%
\TextAt{8}{85.07}{370.18}{-modules}%
\TextAt{8}{134.876}{370.18}{the}%
\TextAt{8}{155.032}{370.18}{rank}%
\TextAt{8}{240.885}{353.134}{rk}%
\TextAt{16}{251.62}{353.134}{\Character{102}}%
\TextAt{13}{257.598}{353.134}{M}%
\TextAt{16}{270.172}{353.134}{\Character{103}}%
\TextAt{9}{279.47}{353.134}{\Character{61}}%
\TextAt{8}{291.895}{353.134}{dim}%
\TextAt{21}{311.406}{351.34}{K}%
\TextAt{9}{319.609}{353.134}{\Character{40}}%
\TextAt{13}{324.161}{353.134}{M}%
\TextAt{16}{339.392}{353.134}{\Character{10}}%
\TextAt{27}{348.69}{351.34}{o}%
\TextAt{13}{355.741}{353.134}{K}%
\TextAt{9}{366.559}{353.134}{\Character{41}}%
\TextAt{8}{79.226}{329.646}{extends}%
\TextAt{8}{122.983}{329.646}{by}%
\TextAt{8}{140.343}{329.646}{linearity}%
\TextAt{8}{187.627}{329.646}{to}%
\TextAt{8}{203.038}{329.646}{a}%
\TextAt{8}{213.894}{329.646}{ring}%
\TextAt{8}{239.061}{329.646}{homomorphism}%
\TextAt{13}{322.832}{329.646}{F}%
\TextAt{17}{330.409}{327.852}{0}%
\TextAt{9}{335.142}{329.646}{\Character{40}}%
\Vector{q 0 G 0 g 0.478 w 267.694 -436.998 m 273.538 -436.998 l S Q}%
\TextAt{14}{339.694}{329.646}{o}%
\TextAt{9}{345.538}{329.646}{\Character{41}}%
\NativeMathAt{355.288}{329.646}{11.955}{\rightarrow}%
\TextAt{18}{372.44}{329.646}{Z}%
\TextAt{8}{380.41}{329.646}{.}%
\TextAt{8}{392.161}{329.646}{For}%
\TextAt{8}{414.231}{329.646}{a}%
\TextAt{8}{425.087}{329.646}{short}%
\TextAt{8}{456.172}{329.646}{exact}%
\TextAt{8}{488.165}{329.646}{sequence}%
\TextAt{9}{79.226}{312.6}{0}%
\NativeMathAt{89.794}{312.6}{11.955}{\rightarrow}%
\TextAt{13}{106.465}{312.6}{M}%
\TextAt{3}{119.038}{316.938}{\Character{48}}%
\NativeMathAt{126.548}{312.6}{11.955}{\rightarrow}%
\TextAt{13}{143.219}{312.6}{M}%
\NativeMathAt{160.507}{312.6}{11.955}{\rightarrow}%
\TextAt{13}{177.178}{312.6}{M}%
\TextAt{3}{189.751}{316.938}{\Character{48}\Character{48}}%
\NativeMathAt{199.558}{312.6}{11.955}{\rightarrow}%
\TextAt{9}{216.229}{312.6}{0}%
\TextAt{8}{226.797}{312.6}{of}%
\TextAt{8}{240.945}{312.6}{metrized}%
\TextAt{14}{289.883}{312.6}{o}%
\TextAt{8}{295.727}{312.6}{-modules}%
\TextAt{8}{346.351}{312.6}{one}%
\TextAt{8}{368.628}{312.6}{has}%
\TextAt{8}{390.32}{312.6}{rk}%
\TextAt{9}{401.056}{312.6}{\Character{40}}%
\TextAt{13}{405.608}{312.6}{M}%
\TextAt{9}{418.182}{312.6}{\Character{41}}%
\TextAt{9}{427.435}{312.6}{\Character{61}}%
\TextAt{8}{441.269}{312.6}{rk}%
\TextAt{9}{452.005}{312.6}{\Character{40}}%
\TextAt{13}{456.557}{312.6}{M}%
\TextAt{3}{469.131}{316.938}{\Character{48}}%
\TextAt{9}{471.926}{312.6}{\Character{41}}%
\TextAt{9}{479.673}{312.6}{\Character{43}}%
\TextAt{8}{492.012}{312.6}{rk}%
\TextAt{9}{502.747}{312.6}{\Character{40}}%
\TextAt{13}{507.3}{312.6}{M}%
\TextAt{3}{519.873}{316.938}{\Character{48}\Character{48}}%
\TextAt{9}{524.965}{312.6}{\Character{41}}%
\TextAt{8}{529.518}{312.6}{,}%
\TextAt{8}{79.226}{295.554}{and}%
\TextAt{8}{102.875}{295.554}{so}%
\TextAt{8}{118.132}{295.554}{rk}%
\TextAt{9}{128.868}{295.554}{\Character{40}}%
\TextAt{16}{133.42}{295.554}{\Character{102}}%
\TextAt{13}{139.398}{295.554}{M}%
\TextAt{3}{151.972}{299.892}{\Character{48}}%
\TextAt{16}{154.767}{295.554}{\Character{103}}%
\TextAt{16}{163.996}{295.554}{\Character{0}}%
\TextAt{16}{176.549}{295.554}{\Character{102}}%
\TextAt{13}{182.543}{295.554}{M}%
\TextAt{16}{195.116}{295.554}{\Character{103}}%
\TextAt{9}{204.355}{295.554}{\Character{43}}%
\TextAt{16}{216.721}{295.554}{\Character{102}}%
\TextAt{13}{222.698}{295.554}{M}%
\TextAt{3}{235.272}{299.892}{\Character{48}\Character{48}}%
\TextAt{16}{240.364}{295.554}{\Character{103}}%
\TextAt{9}{246.341}{295.554}{\Character{41}}%
\TextAt{9}{255.714}{295.554}{\Character{61}}%
\TextAt{9}{269.642}{295.554}{0}%
\TextAt{8}{275.515}{295.554}{.}%
\TextAt{8}{286.624}{295.554}{Thus}%
\TextAt{8}{317.159}{295.554}{rk}%
\TextAt{8}{332.679}{295.554}{is}%
\TextAt{8}{345.33}{295.554}{zero}%
\TextAt{8}{370.928}{295.554}{on}%
\TextAt{8}{388.074}{295.554}{the}%
\TextAt{8}{409.117}{295.554}{ideal}%
\TextAt{13}{437.966}{295.554}{R}%
\TextAt{17}{446.877}{293.76}{0}%
\TextAt{9}{451.609}{295.554}{\Character{40}}%
\Vector{q 0 G 0 g 0.478 w 384.161 -471.09 m 390.005 -471.09 l S Q}%
\TextAt{14}{456.161}{295.554}{o}%
\TextAt{9}{462.005}{295.554}{\Character{41}}%
\TextAt{8}{471.342}{295.554}{and}%
\TextAt{8}{494.991}{295.554}{induces}%
\TextAt{8}{79.226}{278.508}{therefore}%
\TextAt{8}{128.528}{278.508}{a}%
\TextAt{8}{138.485}{278.508}{homomorphism}%
\TextAt{13}{221.356}{278.508}{K}%
\TextAt{17}{231.315}{276.714}{0}%
\TextAt{9}{236.048}{278.508}{\Character{40}}%
\Vector{q 0 G 0 g 0.478 w 168.6 -488.136 m 174.444 -488.136 l S Q}%
\TextAt{14}{240.6}{278.508}{o}%
\TextAt{9}{246.444}{278.508}{\Character{41}}%
\NativeMathAt{254.661}{278.508}{11.955}{\rightarrow}%
\TextAt{18}{270.28}{278.508}{Z}%
\TextAt{8}{278.251}{278.508}{.}%
\TextAt{8}{287.304}{278.508}{It}%
\TextAt{8}{300.178}{278.508}{is}%
\TextAt{8}{312.143}{278.508}{called}%
\TextAt{8}{345.508}{278.508}{the}%
\TextAt{12}{365.866}{278.508}{augmentation}%
\TextAt{8}{450.113}{278.508}{of}%
\TextAt{13}{463.645}{278.508}{K}%
\TextAt{17}{473.604}{276.714}{0}%
\TextAt{9}{478.336}{278.508}{\Character{40}}%
\Vector{q 0 G 0 g 0.478 w 410.888 -488.136 m 416.733 -488.136 l S Q}%
\TextAt{14}{482.888}{278.508}{o}%
\TextAt{9}{488.733}{278.508}{\Character{41}}%
\TextAt{8}{497.384}{278.508}{and}%
\TextAt{8}{520.348}{278.508}{its}%
\TextAt{8}{79.226}{261.462}{kernel}%
\TextAt{13}{113.693}{261.462}{I}%
\TextAt{9}{123.117}{261.462}{\Character{61}}%
\TextAt{8}{135.542}{261.462}{ker}%
\TextAt{9}{151.478}{261.462}{\Character{40}}%
\TextAt{8}{156.031}{261.462}{rk}%
\TextAt{9}{166.766}{261.462}{\Character{41}}%
\TextAt{8}{175.216}{261.462}{is}%
\TextAt{8}{186.98}{261.462}{called}%
\TextAt{8}{220.144}{261.462}{the}%
\TextAt{12}{240.3}{261.462}{augmentation}%
\TextAt{12}{324.931}{261.462}{ideal}%
\TextAt{8}{352.559}{261.462}{.}%
\TextAt{12}{79.226}{235.472}{Proposition}%
\TextAt{12}{153.778}{235.472}{(6.1).}%
\TextAt{8}{190.694}{235.472}{The}%
\TextAt{8}{216.348}{235.472}{ideal}%
\TextAt{13}{245.911}{235.472}{I}%
\TextAt{8}{252.014}{235.472}{,}%
\TextAt{8}{261.163}{235.472}{resp.}%
\TextAt{13}{295.286}{235.472}{I}%
\TextAt{17}{301.388}{239.81}{2}%
\TextAt{8}{306.121}{235.472}{,}%
\TextAt{8}{315.27}{235.472}{is}%
\TextAt{8}{328.634}{235.472}{generated}%
\TextAt{8}{383.566}{235.472}{as}%
\TextAt{8}{399.537}{235.472}{an}%
\TextAt{8}{417.396}{235.472}{additive}%
\TextAt{8}{463.876}{235.472}{group}%
\TextAt{8}{498.651}{235.472}{by}%
\TextAt{8}{516.51}{235.472}{the}%
\TextAt{8}{79.226}{218.426}{elements}%
\TextAt{9}{127.083}{218.426}{\Character{91}}%
\TextAt{14}{130.334}{218.426}{a}%
\TextAt{9}{136.317}{218.426}{\Character{93}}%
\TextAt{16}{142.226}{218.426}{\Character{0}}%
\TextAt{9}{154.181}{218.426}{1}%
\TextAt{8}{160.034}{218.426}{,}%
\TextAt{8}{167.183}{218.426}{resp.}%
\TextAt{9}{196.505}{218.426}{\Character{40}\Character{91}}%
\TextAt{14}{204.309}{218.426}{a}%
\TextAt{9}{210.292}{218.426}{\Character{93}}%
\TextAt{16}{216.201}{218.426}{\Character{0}}%
\TextAt{9}{228.156}{218.426}{1\Character{41}\Character{40}\Character{91}}%
\TextAt{14}{246.365}{218.426}{b}%
\TextAt{9}{252.5}{218.426}{\Character{93}}%
\TextAt{16}{258.408}{218.426}{\Character{0}}%
\TextAt{9}{270.363}{218.426}{1\Character{41}}%
\TextAt{8}{280.769}{218.426}{,}%
\TextAt{8}{287.918}{218.426}{where}%
\TextAt{14}{321.727}{218.426}{a}%
\TextAt{13}{327.71}{218.426}{\Character{59}}%
\TextAt{14}{332.954}{218.426}{b}%
\TextAt{8}{342.987}{218.426}{vary}%
\TextAt{8}{369.013}{218.426}{over}%
\TextAt{8}{394.059}{218.426}{the}%
\TextAt{8}{414.215}{218.426}{replete}%
\TextAt{8}{452.58}{218.426}{ideals}%
\TextAt{8}{485.157}{218.426}{of}%
\TextAt{13}{498.487}{218.426}{K}%
\TextAt{8}{509.305}{218.426}{.}%
\TextAt{19}{79.226}{183.492}{Proof.}%
\TextAt{8}{116}{183.492}{By}%
\TextAt{8}{134.363}{183.492}{(5.6),}%
\TextAt{8}{165.59}{183.492}{every}%
\TextAt{8}{196.482}{183.492}{element}%
\TextAt{13}{239.724}{183.492}{\Character{24}}%
\NativeMathAt{248.736}{183.492}{11.955}{\in}%
\TextAt{13}{260.027}{183.492}{K}%
\TextAt{17}{269.986}{181.699}{0}%
\TextAt{9}{274.718}{183.492}{\Character{40}}%
\Vector{q 0 G 0 g 0.478 w 207.27 -583.152 m 213.115 -583.152 l S Q}%
\TextAt{14}{279.27}{183.492}{o}%
\TextAt{9}{285.115}{183.492}{\Character{41}}%
\TextAt{8}{293.564}{183.492}{is}%
\TextAt{8}{305.328}{183.492}{of}%
\TextAt{8}{318.658}{183.492}{the}%
\TextAt{8}{338.815}{183.492}{form}%
\TextAt{13}{270.905}{148.047}{\Character{24}}%
\TextAt{9}{279.916}{148.047}{\Character{61}}%
\TextAt{21}{298.948}{162.991}{r}%
\TextAt{20}{292.342}{159.405}{\Character{88}}%
\TextAt{21}{294.124}{134.21}{i}%
\TextAt{17}{297.007}{134.21}{\Character{61}1}%
\TextAt{13}{311.603}{148.047}{n}%
\TextAt{21}{318.59}{146.254}{i}%
\TextAt{9}{321.972}{148.047}{\Character{91}}%
\TextAt{14}{325.223}{148.047}{a}%
\TextAt{21}{331.206}{146.254}{i}%
\TextAt{9}{334.588}{148.047}{\Character{93}}%
\TextAt{13}{337.839}{148.047}{\Character{58}}%
\TextAt{8}{79.226}{110.899}{If}%
\TextAt{13}{90.918}{110.899}{\Character{24}}%
\NativeMathAt{99.93}{110.899}{11.955}{\in}%
\TextAt{13}{111.221}{110.899}{I}%
\TextAt{8}{117.323}{110.899}{,}%
\TextAt{8}{124.473}{110.899}{then}%
\TextAt{8}{151.133}{110.899}{rk}%
\TextAt{9}{161.868}{110.899}{\Character{40}}%
\TextAt{13}{166.421}{110.899}{\Character{24}}%
\TextAt{9}{172.111}{110.899}{\Character{41}}%
\TextAt{9}{179.966}{110.899}{\Character{61}}%
\TextAt{20}{192.41}{119.865}{\Character{80}}%
\TextAt{21}{205.029}{116.71}{r}%
\TextAt{21}{205.029}{107.412}{i}%
\TextAt{17}{207.912}{107.412}{\Character{61}1}%
\TextAt{13}{221.224}{110.899}{n}%
\TextAt{21}{228.211}{109.106}{i}%
\TextAt{9}{234.913}{110.899}{\Character{61}}%
\TextAt{9}{247.323}{110.899}{0}%
\TextAt{8}{253.192}{110.899}{,}%
\TextAt{8}{260.341}{110.899}{and}%
\TextAt{8}{283.104}{110.899}{thus}%
\TextAt{13}{203.377}{81.976}{\Character{24}}%
\TextAt{9}{212.388}{81.976}{\Character{61}}%
\TextAt{20}{224.814}{93.334}{\Character{88}}%
\TextAt{21}{232.007}{68.139}{i}%
\TextAt{13}{244.075}{81.976}{n}%
\TextAt{21}{251.063}{80.183}{i}%
\TextAt{9}{254.444}{81.976}{\Character{91}}%
\TextAt{14}{257.696}{81.976}{a}%
\TextAt{21}{263.679}{80.183}{i}%
\TextAt{9}{267.06}{81.976}{\Character{93}}%
\TextAt{16}{272.968}{81.976}{\Character{0}}%
\TextAt{20}{284.923}{93.334}{\Character{88}}%
\TextAt{21}{292.116}{68.139}{i}%
\TextAt{13}{304.185}{81.976}{n}%
\TextAt{21}{311.172}{80.183}{i}%
\TextAt{9}{317.874}{81.976}{\Character{61}}%
\TextAt{20}{330.3}{93.334}{\Character{88}}%
\TextAt{21}{337.492}{68.139}{i}%
\TextAt{13}{349.561}{81.976}{n}%
\TextAt{21}{356.548}{80.183}{i}%
\TextAt{9}{359.93}{81.976}{\Character{40}\Character{91}}%
\TextAt{14}{367.734}{81.976}{a}%
\TextAt{21}{373.717}{80.183}{i}%
\TextAt{9}{377.098}{81.976}{\Character{93}}%
\TextAt{16}{383.006}{81.976}{\Character{0}}%
\TextAt{9}{394.962}{81.976}{1\Character{41}}%
\TextAt{13}{405.367}{81.976}{\Character{58}}%
\TextAt{8}{297.211}{43.033}{223}%
\Destination{78.23}{809.815}{nav.page.223}%
\Destination{79.23}{502.815}{nav.section.3.6}%
\Destination{79.36}{694.755}{nav.print.243}%
\Destination{79.36}{249.455}{nav.number.3.6.1}%
\Destination{79.23}{502.815}{outline.37}%
\Hotlink{133.863}{181.255}{25.078}{11.047}{nav.number.3.5.6}%
}

% Page 236 (body page 224)
\Page{%
\TextAt{2}{79.226}{796.943}{III.}%
\TextAt{2}{98.317}{796.943}{Riemann–Roch}%
\TextAt{2}{174.375}{796.943}{Theory}%
\Vector{q 0 G 0 g 0.299 w 7.226 18.729 m 460.769 18.729 l S Q}%
\TextAt{8}{79.226}{761.124}{The}%
\TextAt{8}{103.28}{761.124}{ideal}%
\TextAt{13}{131.243}{761.124}{I}%
\TextAt{17}{137.346}{765.462}{2}%
\TextAt{8}{145.976}{761.124}{is}%
\TextAt{8}{157.739}{761.124}{therefore}%
\TextAt{8}{206.839}{761.124}{generated}%
\TextAt{8}{260.171}{761.124}{by}%
\TextAt{8}{276.43}{761.124}{the}%
\TextAt{8}{296.587}{761.124}{elements}%
\TextAt{9}{344.443}{761.124}{\Character{40}\Character{91}}%
\TextAt{14}{352.247}{761.124}{a}%
\TextAt{9}{358.23}{761.124}{\Character{93}}%
\TextAt{16}{364.139}{761.124}{\Character{0}}%
\TextAt{9}{376.094}{761.124}{1\Character{41}\Character{40}\Character{91}}%
\TextAt{14}{394.303}{761.124}{b}%
\TextAt{9}{400.438}{761.124}{\Character{93}}%
\TextAt{16}{406.346}{761.124}{\Character{0}}%
\TextAt{9}{418.302}{761.124}{1\Character{41}}%
\TextAt{8}{428.707}{761.124}{.}%
\TextAt{8}{437.155}{761.124}{As}%
\TextAt{9}{171.178}{732.123}{\Character{91}}%
\TextAt{14}{174.43}{732.123}{c}%
\TextAt{9}{179.086}{732.123}{\Character{93}\Character{40}\Character{91}}%
\TextAt{14}{190.141}{732.123}{a}%
\TextAt{9}{196.124}{732.123}{\Character{93}}%
\TextAt{16}{202.033}{732.123}{\Character{0}}%
\TextAt{9}{213.988}{732.123}{1\Character{41}\Character{40}\Character{91}}%
\TextAt{14}{232.197}{732.123}{b}%
\TextAt{9}{238.332}{732.123}{\Character{93}}%
\TextAt{16}{244.24}{732.123}{\Character{0}}%
\TextAt{9}{256.195}{732.123}{1\Character{41}}%
\TextAt{9}{269.896}{732.123}{\Character{61}}%
\TextAt{20}{282.347}{741.806}{\Character{0}}%
\TextAt{9}{287.827}{732.123}{\Character{40}\Character{91}}%
\TextAt{14}{295.631}{732.123}{ca}%
\TextAt{9}{306.27}{732.123}{\Character{93}}%
\TextAt{16}{312.178}{732.123}{\Character{0}}%
\TextAt{9}{324.133}{732.123}{1\Character{41}}%
\TextAt{16}{337.195}{732.123}{\Character{0}}%
\TextAt{9}{349.15}{732.123}{\Character{40}\Character{91}}%
\TextAt{14}{356.954}{732.123}{c}%
\TextAt{9}{361.61}{732.123}{\Character{93}}%
\TextAt{16}{367.519}{732.123}{\Character{0}}%
\TextAt{9}{379.474}{732.123}{1\Character{41}}%
\TextAt{20}{389.879}{741.806}{\Character{1}}%
\TextAt{9}{395.359}{732.123}{\Character{40}\Character{91}}%
\TextAt{14}{403.163}{732.123}{b}%
\TextAt{9}{409.297}{732.123}{\Character{93}}%
\TextAt{16}{415.206}{732.123}{\Character{0}}%
\TextAt{9}{427.161}{732.123}{1\Character{41}}%
\TextAt{13}{437.566}{732.123}{\Character{59}}%
\TextAt{8}{79.226}{703.121}{these}%
\TextAt{8}{109.198}{703.121}{elements}%
\TextAt{8}{157.054}{703.121}{already}%
\TextAt{8}{198.359}{703.121}{form}%
\TextAt{8}{226}{703.121}{a}%
\TextAt{8}{235.755}{703.121}{system}%
\TextAt{8}{274.574}{703.121}{of}%
\TextAt{8}{287.904}{703.121}{generators}%
\TextAt{8}{344.559}{703.121}{of}%
\TextAt{8}{357.889}{703.121}{the}%
\TextAt{8}{378.046}{703.121}{abelian}%
\TextAt{8}{418.693}{703.121}{group}%
\TextAt{13}{451.869}{703.121}{I}%
\TextAt{17}{457.972}{707.46}{2}%
\TextAt{8}{462.704}{703.121}{.}%
\Vector{q 0 G 0 g 0.399 w 452.334 -70.879 m 452.334 -62.809 l S Q}%
\Vector{q 0 G 0 g 0.399 w 452.534 -63.008 m 459.707 -63.008 l S Q}%
\Vector{q 0 G 0 g 0.399 w 452.534 -70.679 m 459.707 -70.679 l S Q}%
\Vector{q 0 G 0 g 0.399 w 459.906 -70.879 m 459.906 -62.809 l S Q}%
\TextAt{8}{97.159}{674.12}{By}%
\TextAt{8}{115.522}{674.12}{(5.7),}%
\TextAt{8}{146.749}{674.12}{this}%
\TextAt{8}{169.571}{674.12}{gives}%
\TextAt{8}{198.252}{674.12}{us}%
\TextAt{8}{213.268}{674.12}{the}%
\TextAt{12}{79.226}{648.108}{Corollary}%
\TextAt{12}{139.014}{648.108}{(6.2).}%
\TextAt{13}{174.33}{648.108}{I}%
\TextAt{17}{180.432}{652.446}{2}%
\TextAt{9}{188.485}{648.108}{\Character{61}}%
\TextAt{9}{200.895}{648.108}{0}%
\TextAt{8}{206.764}{648.108}{.}%
\TextAt{8}{97.159}{622.096}{We}%
\TextAt{8}{117.303}{622.096}{now}%
\TextAt{8}{141.692}{622.096}{define}%
\TextAt{8}{262.759}{605.05}{gr}%
\TextAt{13}{275.165}{605.05}{K}%
\TextAt{17}{285.124}{603.256}{0}%
\TextAt{9}{289.856}{605.05}{\Character{40}}%
\Vector{q 0 G 0 g 0.478 w 222.409 -161.594 m 228.253 -161.594 l S Q}%
\TextAt{14}{294.409}{605.05}{o}%
\TextAt{9}{300.253}{605.05}{\Character{41}}%
\TextAt{9}{308.108}{605.05}{\Character{61}}%
\TextAt{18}{320.551}{605.05}{Z}%
\TextAt{16}{331.178}{605.05}{\Character{8}}%
\TextAt{13}{343.133}{605.05}{I}%
\TextAt{8}{79.226}{581.528}{and}%
\TextAt{8}{101.989}{581.528}{turn}%
\TextAt{8}{128.003}{581.528}{this}%
\TextAt{8}{150.826}{581.528}{additive}%
\TextAt{8}{195.705}{581.528}{group}%
\TextAt{8}{228.881}{581.528}{into}%
\TextAt{8}{252.624}{581.528}{a}%
\TextAt{8}{262.379}{581.528}{ring}%
\TextAt{8}{286.445}{581.528}{by}%
\TextAt{8}{302.704}{581.528}{putting}%
\TextAt{13}{344.332}{581.528}{xy}%
\TextAt{9}{360.442}{581.528}{\Character{61}}%
\TextAt{9}{372.851}{581.528}{0}%
\TextAt{8}{382.617}{581.528}{for}%
\TextAt{13}{400.502}{581.528}{x\Character{59}}%
\TextAt{13}{412.374}{581.528}{y}%
\NativeMathAt{421.856}{581.528}{11.955}{\in}%
\TextAt{13}{433.147}{581.528}{I}%
\TextAt{8}{439.25}{581.528}{.}%
\TextAt{12}{79.226}{555.516}{Definition}%
\TextAt{12}{142.015}{555.516}{(6.3).}%
\TextAt{8}{177.33}{555.516}{The}%
\TextAt{8}{201.384}{555.516}{additive}%
\TextAt{8}{246.264}{555.516}{homomorphism}%
\TextAt{13}{207.693}{526.514}{c}%
\TextAt{17}{212.731}{524.721}{1}%
\TextAt{9}{220.785}{526.514}{\Character{58}}%
\TextAt{13}{227.357}{526.514}{K}%
\TextAt{17}{237.316}{524.721}{0}%
\TextAt{9}{242.048}{526.514}{\Character{40}}%
\Vector{q 0 G 0 g 0.478 w 174.601 -240.129 m 180.445 -240.129 l S Q}%
\TextAt{14}{246.601}{526.514}{o}%
\TextAt{9}{252.445}{526.514}{\Character{41}}%
\NativeArrowAt{260.318}{526.514}{11.955}{19.26}{\longrightarrow}%
\TextAt{13}{282.9}{526.514}{I}%
\TextAt{13}{288.985}{526.514}{\Character{59}}%
\TextAt{13}{318.132}{526.514}{c}%
\TextAt{17}{323.195}{524.721}{1}%
\TextAt{9}{327.928}{526.514}{\Character{40}}%
\TextAt{13}{332.48}{526.514}{\Character{24}}%
\TextAt{9}{338.171}{526.514}{\Character{41}}%
\TextAt{9}{346.026}{526.514}{\Character{61}}%
\TextAt{13}{358.469}{526.514}{\Character{24}}%
\TextAt{16}{366.816}{526.514}{\Character{0}}%
\TextAt{8}{378.771}{526.514}{rk}%
\TextAt{9}{389.507}{526.514}{\Character{40}}%
\TextAt{13}{394.059}{526.514}{\Character{24}}%
\TextAt{9}{399.75}{526.514}{\Character{41}}%
\TextAt{8}{79.226}{497.513}{is}%
\TextAt{8}{90.99}{497.513}{called}%
\TextAt{8}{124.154}{497.513}{the}%
\TextAt{12}{144.31}{497.513}{first}%
\TextAt{12}{172.309}{497.513}{Chern}%
\TextAt{12}{213.088}{497.513}{class}%
\TextAt{8}{239.963}{497.513}{.}%
\TextAt{8}{248.412}{497.513}{The}%
\TextAt{8}{272.466}{497.513}{mapping}%
\TextAt{8}{176.822}{468.512}{ch}%
\TextAt{9}{191.847}{468.512}{\Character{58}}%
\TextAt{13}{198.419}{468.512}{K}%
\TextAt{17}{208.378}{466.719}{0}%
\TextAt{9}{213.111}{468.512}{\Character{40}}%
\Vector{q 0 G 0 g 0.478 w 145.663 -298.131 m 151.507 -298.131 l S Q}%
\TextAt{14}{217.663}{468.512}{o}%
\TextAt{9}{223.507}{468.512}{\Character{41}}%
\NativeArrowAt{231.38}{468.512}{11.955}{19.26}{\longrightarrow}%
\TextAt{8}{253.962}{468.512}{gr}%
\TextAt{13}{266.368}{468.512}{K}%
\TextAt{17}{276.327}{466.719}{0}%
\TextAt{9}{281.059}{468.512}{\Character{40}}%
\Vector{q 0 G 0 g 0.478 w 213.612 -298.131 m 219.456 -298.131 l S Q}%
\TextAt{14}{285.612}{468.512}{o}%
\TextAt{9}{291.456}{468.512}{\Character{41}}%
\TextAt{13}{296.008}{468.512}{\Character{59}}%
\TextAt{8}{325.163}{468.512}{ch}%
\TextAt{9}{336.867}{468.512}{\Character{40}}%
\TextAt{13}{341.419}{468.512}{\Character{24}}%
\TextAt{9}{347.109}{468.512}{\Character{41}}%
\TextAt{9}{354.964}{468.512}{\Character{61}}%
\TextAt{8}{367.408}{468.512}{rk}%
\TextAt{9}{378.144}{468.512}{\Character{40}}%
\TextAt{13}{382.696}{468.512}{\Character{24}}%
\TextAt{9}{388.387}{468.512}{\Character{41}}%
\TextAt{9}{395.584}{468.512}{\Character{43}}%
\TextAt{13}{407.357}{468.512}{c}%
\TextAt{17}{412.395}{466.719}{1}%
\TextAt{9}{417.127}{468.512}{\Character{40}}%
\TextAt{13}{421.679}{468.512}{\Character{24}}%
\TextAt{9}{427.37}{468.512}{\Character{41}}%
\TextAt{13}{431.922}{468.512}{\Character{59}}%
\TextAt{8}{79.226}{439.511}{is}%
\TextAt{8}{90.99}{439.511}{called}%
\TextAt{8}{124.154}{439.511}{the}%
\TextAt{12}{144.31}{439.511}{Chern}%
\TextAt{12}{185.089}{439.511}{character}%
\TextAt{8}{243.49}{439.511}{of}%
\TextAt{13}{256.82}{439.511}{K}%
\TextAt{17}{266.779}{437.718}{0}%
\TextAt{9}{271.512}{439.511}{\Character{40}}%
\Vector{q 0 G 0 g 0.478 w 204.064 -327.133 m 209.908 -327.133 l S Q}%
\TextAt{14}{276.064}{439.511}{o}%
\TextAt{9}{281.908}{439.511}{\Character{41}}%
\TextAt{8}{286.461}{439.511}{.}%
\TextAt{12}{79.226}{407.521}{Proposition}%
\TextAt{12}{151.937}{407.521}{(6.4).}%
\TextAt{8}{187.253}{407.521}{The}%
\TextAt{8}{211.307}{407.521}{Chern}%
\TextAt{8}{246.419}{407.521}{character}%
\TextAt{8}{246.405}{378.52}{ch}%
\TextAt{9}{261.43}{378.52}{\Character{58}}%
\TextAt{13}{268.002}{378.52}{K}%
\TextAt{17}{277.961}{376.727}{0}%
\TextAt{9}{282.693}{378.52}{\Character{40}}%
\Vector{q 0 G 0 g 0.478 w 215.246 -388.124 m 221.09 -388.124 l S Q}%
\TextAt{14}{287.246}{378.52}{o}%
\TextAt{9}{293.09}{378.52}{\Character{41}}%
\NativeArrowAt{300.963}{378.52}{11.955}{19.26}{\longrightarrow}%
\TextAt{8}{323.545}{378.52}{gr}%
\TextAt{13}{335.951}{378.52}{K}%
\TextAt{17}{345.91}{376.727}{0}%
\TextAt{9}{350.642}{378.52}{\Character{40}}%
\Vector{q 0 G 0 g 0.478 w 283.194 -388.124 m 289.039 -388.124 l S Q}%
\TextAt{14}{355.194}{378.52}{o}%
\TextAt{9}{361.039}{378.52}{\Character{41}}%
\TextAt{8}{79.226}{349.519}{is}%
\TextAt{8}{90.99}{349.519}{an}%
\TextAt{8}{107.249}{349.519}{isomorphism}%
\TextAt{8}{175.669}{349.519}{of}%
\TextAt{8}{188.999}{349.519}{rings.}%
\ErrataAt{79.22}{314.555}{453.64}{\textbf{Proof.} Both $\mathrm{rk}$ and $c_1$ are additive, and $\mathrm{rk}$ is multiplicative. We show that}%
\ErrataAt{79.22}{297.455}{453.64}{$c_1(xy)=\mathrm{rk}(x)c_1(y)+\mathrm{rk}(y)c_1(x)$. By additivity it suffices to check the generators}%
\ErrataAt{79.22}{280.455}{453.64}{$x=[\mathfrak a]$ and $y=[\mathfrak b]$, which both have rank one. This works because}%
\TextAt{13}{129.47}{251.447}{c}%
\TextAt{17}{134.508}{249.654}{1}%
\TextAt{9}{139.24}{251.447}{\Character{40}}%
\TextAt{13}{143.793}{251.447}{xy}%
\TextAt{9}{156.581}{251.447}{\Character{41}}%
\TextAt{9}{164.436}{251.447}{\Character{61}}%
\TextAt{13}{176.88}{251.447}{xy}%
\TextAt{16}{192.325}{251.447}{\Character{0}}%
\TextAt{9}{204.281}{251.447}{1}%
\TextAt{9}{213.439}{251.447}{\Character{61}}%
\TextAt{9}{225.86}{251.447}{\Character{40}}%
\TextAt{13}{230.432}{251.447}{x}%
\TextAt{16}{239.741}{251.447}{\Character{0}}%
\TextAt{9}{251.696}{251.447}{1\Character{41}}%
\TextAt{9}{264.739}{251.447}{\Character{43}}%
\TextAt{9}{276.491}{251.447}{\Character{40}}%
\TextAt{13}{281.072}{251.447}{y}%
\TextAt{16}{289.865}{251.447}{\Character{0}}%
\TextAt{9}{301.82}{251.447}{1\Character{41}}%
\TextAt{9}{314.863}{251.447}{\Character{43}}%
\TextAt{9}{326.615}{251.447}{\Character{40}}%
\TextAt{13}{331.196}{251.447}{x}%
\TextAt{16}{340.505}{251.447}{\Character{0}}%
\TextAt{9}{352.46}{251.447}{1\Character{41}\Character{40}}%
\TextAt{13}{367.417}{251.447}{y}%
\TextAt{16}{376.211}{251.447}{\Character{0}}%
\TextAt{9}{388.166}{251.447}{1\Character{41}}%
\TextAt{9}{401.867}{251.447}{\Character{61}}%
\TextAt{13}{414.318}{251.447}{c}%
\TextAt{17}{419.356}{249.654}{1}%
\TextAt{9}{424.088}{251.447}{\Character{40}}%
\TextAt{13}{428.64}{251.447}{x}%
\TextAt{9}{435.292}{251.447}{\Character{41}}%
\TextAt{9}{442.489}{251.447}{\Character{43}}%
\TextAt{13}{454.263}{251.447}{c}%
\TextAt{17}{459.301}{249.654}{1}%
\TextAt{9}{464.033}{251.447}{\Character{40}}%
\TextAt{13}{468.585}{251.447}{y}%
\TextAt{9}{474.722}{251.447}{\Character{41}}%
\TextAt{13}{479.274}{251.447}{\Character{59}}%
\TextAt{8}{79.226}{222.446}{because}%
\TextAt{9}{122.211}{222.446}{\Character{40}}%
\TextAt{13}{126.763}{222.446}{x}%
\TextAt{16}{135.423}{222.446}{\Character{0}}%
\TextAt{9}{146.73}{222.446}{1\Character{41}\Character{40}}%
\TextAt{13}{161.687}{222.446}{y}%
\TextAt{16}{169.832}{222.446}{\Character{0}}%
\TextAt{9}{181.139}{222.446}{1\Character{41}}%
\TextAt{9}{194.84}{222.446}{\Character{61}}%
\TextAt{9}{207.261}{222.446}{0}%
\TextAt{8}{216.723}{222.446}{by}%
\TextAt{8}{232.665}{222.446}{(5.7).}%
\TextAt{8}{265.086}{222.446}{Therefore}%
\TextAt{8}{317.766}{222.446}{ch}%
\TextAt{8}{333.05}{222.446}{is}%
\TextAt{8}{344.497}{222.446}{a}%
\TextAt{8}{353.935}{222.446}{ring}%
\TextAt{8}{377.684}{222.446}{homomorphism.}%
\TextAt{8}{464.799}{222.446}{The}%
\TextAt{8}{488.535}{222.446}{mapping}%
\TextAt{18}{207.536}{193.445}{Z}%
\TextAt{16}{218.163}{193.445}{\Character{8}}%
\TextAt{13}{230.118}{193.445}{I}%
\NativeArrowAt{239.542}{193.445}{11.955}{19.26}{\longrightarrow}%
\TextAt{13}{262.124}{193.445}{K}%
\TextAt{17}{272.083}{191.651}{0}%
\TextAt{9}{276.815}{193.445}{\Character{40}}%
\Vector{q 0 G 0 g 0.478 w 209.368 -573.199 m 215.212 -573.199 l S Q}%
\TextAt{14}{281.368}{193.445}{o}%
\TextAt{9}{287.212}{193.445}{\Character{41}}%
\TextAt{13}{291.764}{193.445}{\Character{59}}%
\TextAt{13}{320.899}{193.445}{n}%
\TextAt{16}{330.563}{193.445}{\Character{8}}%
\TextAt{13}{342.518}{193.445}{\Character{24}}%
\NativeArrowAt{351.53}{193.445}{11.955}{19.26}{\longmapsto}%
\TextAt{13}{374.112}{193.445}{\Character{24}}%
\TextAt{9}{382.459}{193.445}{\Character{43}}%
\TextAt{13}{394.22}{193.445}{n\Character{59}}%
\TextAt{8}{79.226}{164.444}{is}%
\TextAt{8}{90.99}{164.444}{obviously}%
\TextAt{8}{142.768}{164.444}{an}%
\TextAt{8}{159.027}{164.444}{inverse}%
\TextAt{8}{197.786}{164.444}{mapping,}%
\TextAt{8}{249.169}{164.444}{so}%
\TextAt{8}{263.539}{164.444}{that}%
\TextAt{8}{288.908}{164.444}{ch}%
\TextAt{8}{304.509}{164.444}{is}%
\TextAt{8}{316.273}{164.444}{even}%
\TextAt{8}{342.933}{164.444}{an}%
\TextAt{8}{359.192}{164.444}{isomorphism.}%
\Vector{q 0 G 0 g 0.399 w 452.334 -609.556 m 452.334 -601.487 l S Q}%
\Vector{q 0 G 0 g 0.399 w 452.534 -601.686 m 459.707 -601.686 l S Q}%
\Vector{q 0 G 0 g 0.399 w 459.906 -609.556 m 459.906 -601.487 l S Q}%
\TextAt{8}{297.211}{43.033}{224}%
\Destination{78.23}{809.815}{nav.page.224}%
\Destination{79.36}{776.155}{nav.print.244}%
\Destination{79.36}{662.155}{nav.number.3.6.2}%
\Destination{79.36}{569.555}{nav.number.3.6.3}%
\Destination{79.36}{421.555}{nav.number.3.6.4}%
\Hotlink{115.022}{671.955}{25.078}{11.047}{nav.number.3.5.7}%
\Hotlink{232.165}{220.255}{25.078}{11.047}{nav.number.3.5.7}%
\Hotlink{147.774}{112.455}{21.483}{9.572}{nav.number.3.6.2}%
\Destination{79.226}{328.555}{errata.23}%
}

% Page 237 (body page 225)
\Page{%
\TextAt{2}{79.226}{796.943}{III.}%
\TextAt{2}{98.317}{796.943}{Riemann–Roch}%
\TextAt{2}{174.375}{796.943}{Theory}%
\Vector{q 0 G 0 g 0.299 w 7.226 18.729 m 460.769 18.729 l S Q}%
\TextAt{8}{97.159}{761.124}{We}%
\TextAt{8}{117.614}{761.124}{obtain}%
\TextAt{8}{154.352}{761.124}{a}%
\TextAt{8}{164.419}{761.124}{complete}%
\TextAt{8}{214.152}{761.124}{and}%
\TextAt{8}{237.226}{761.124}{explicit}%
\TextAt{8}{278.83}{761.124}{description}%
\TextAt{8}{339.036}{761.124}{of}%
\TextAt{8}{352.677}{761.124}{the}%
\TextAt{8}{373.144}{761.124}{Chern}%
\TextAt{8}{408.567}{761.124}{character}%
\TextAt{8}{459.939}{761.124}{by}%
\TextAt{8}{476.508}{761.124}{taking}%
\TextAt{8}{512.924}{761.124}{into}%
\TextAt{8}{79.226}{744.078}{account}%
\TextAt{8}{122.48}{744.078}{another}%
\TextAt{8}{165.411}{744.078}{homomorphism,}%
\TextAt{8}{251.333}{744.078}{as}%
\TextAt{8}{265.703}{744.078}{well}%
\TextAt{8}{289.434}{744.078}{as}%
\TextAt{8}{303.804}{744.078}{the}%
\TextAt{8}{323.96}{744.078}{homomorphism}%
\TextAt{8}{406.63}{744.078}{rk}%
\TextAt{9}{420.687}{744.078}{\Character{58}}%
\TextAt{13}{427.259}{744.078}{K}%
\TextAt{17}{437.219}{742.284}{0}%
\TextAt{9}{441.951}{744.078}{\Character{40}}%
\Vector{q 0 G 0 g 0.478 w 374.503 -22.566 m 380.347 -22.566 l S Q}%
\TextAt{14}{446.503}{744.078}{o}%
\TextAt{9}{452.347}{744.078}{\Character{41}}%
\NativeMathAt{460.221}{744.078}{11.955}{\rightarrow}%
\TextAt{18}{475.497}{744.078}{Z}%
\TextAt{8}{483.467}{744.078}{,}%
\TextAt{8}{490.616}{744.078}{namely}%
\TextAt{8}{249.468}{718.062}{det}%
\TextAt{9}{269.048}{718.062}{\Character{58}}%
\TextAt{13}{275.621}{718.062}{K}%
\TextAt{17}{285.58}{716.269}{0}%
\TextAt{9}{290.312}{718.062}{\Character{40}}%
\Vector{q 0 G 0 g 0.478 w 222.864 -48.582 m 228.709 -48.582 l S Q}%
\TextAt{14}{294.864}{718.062}{o}%
\TextAt{9}{300.709}{718.062}{\Character{41}}%
\NativeArrowAt{308.582}{718.062}{11.955}{19.26}{\longrightarrow}%
\TextAt{8}{331.164}{718.062}{Pic}%
\TextAt{9}{347.578}{718.062}{\Character{40}}%
\Vector{q 0 G 0 g 0.478 w 280.131 -48.582 m 285.975 -48.582 l S Q}%
\TextAt{14}{352.131}{718.062}{o}%
\TextAt{9}{357.975}{718.062}{\Character{41}}%
\TextAt{8}{79.226}{692.046}{which}%
\TextAt{8}{112.344}{692.046}{is}%
\TextAt{8}{123.739}{692.046}{induced}%
\TextAt{8}{166.935}{692.046}{by}%
\TextAt{8}{182.825}{692.046}{taking}%
\TextAt{8}{218.561}{692.046}{determinants}%
\TextAt{8}{288.823}{692.046}{det}%
\TextAt{13}{307.075}{692.046}{M}%
\TextAt{8}{323.177}{692.046}{of}%
\TextAt{8}{336.139}{692.046}{projective}%
\TextAt{14}{389.747}{692.046}{o}%
\TextAt{8}{395.592}{692.046}{-modules}%
\TextAt{13}{445.028}{692.046}{M}%
\TextAt{8}{461.13}{692.046}{as}%
\TextAt{8}{475.132}{692.046}{follows}%
\TextAt{8}{513.199}{692.046}{(see}%
\TextAt{8}{79.226}{675}{§4).}%
\TextAt{8}{105.408}{675}{det}%
\TextAt{13}{123.659}{675}{M}%
\TextAt{8}{140.678}{675}{is}%
\TextAt{8}{152.989}{675}{an}%
\TextAt{8}{169.795}{675}{invertible}%
\TextAt{8}{222.049}{675}{metrized}%
\TextAt{14}{270.715}{675}{o}%
\TextAt{8}{276.56}{675}{-module,}%
\TextAt{8}{325.686}{675}{and}%
\TextAt{8}{348.996}{675}{therefore}%
\TextAt{8}{398.643}{675}{of}%
\TextAt{8}{412.52}{675}{the}%
\TextAt{8}{433.224}{675}{form}%
\TextAt{13}{461.412}{675}{L}%
\TextAt{9}{469.376}{675}{\Character{40}}%
\TextAt{14}{473.929}{675}{a}%
\TextAt{9}{479.912}{675}{\Character{41}}%
\TextAt{8}{488.909}{675}{for}%
\TextAt{8}{507.341}{675}{some}%
\TextAt{8}{79.226}{657.954}{replete}%
\TextAt{8}{116.68}{657.954}{ideal}%
\TextAt{14}{143.732}{657.954}{a}%
\TextAt{8}{149.715}{657.954}{,}%
\TextAt{8}{156.136}{657.954}{which}%
\TextAt{8}{188.711}{657.954}{is}%
\TextAt{8}{199.564}{657.954}{well}%
\TextAt{8}{222.385}{657.954}{determined}%
\TextAt{8}{282.601}{657.954}{up}%
\TextAt{8}{298.595}{657.954}{to}%
\TextAt{8}{311.994}{657.954}{isomorphism.}%
\TextAt{8}{384.661}{657.954}{Denoting}%
\TextAt{8}{434.321}{657.954}{by}%
\TextAt{9}{449.669}{657.954}{\Character{91}}%
\TextAt{8}{452.921}{657.954}{det}%
\TextAt{13}{471.172}{657.954}{M}%
\TextAt{9}{483.746}{657.954}{\Character{93}}%
\TextAt{8}{489.984}{657.954}{the}%
\TextAt{8}{509.23}{657.954}{class}%
\TextAt{8}{79.226}{640.908}{of}%
\TextAt{14}{92.556}{640.908}{a}%
\TextAt{8}{102.437}{640.908}{in}%
\TextAt{8}{116.089}{640.908}{Pic}%
\TextAt{9}{132.504}{640.908}{\Character{40}}%
\Vector{q 0 G 0 g 0.478 w 65.056 -125.735 m 70.9 -125.735 l S Q}%
\TextAt{14}{137.056}{640.908}{o}%
\TextAt{9}{142.9}{640.908}{\Character{41}}%
\TextAt{8}{147.453}{640.908}{,}%
\TextAt{8}{154.602}{640.908}{the}%
\TextAt{8}{174.758}{640.908}{linear}%
\TextAt{8}{207.276}{640.908}{extension}%
\TextAt{8}{259.042}{640.908}{of}%
\TextAt{8}{272.372}{640.908}{the}%
\TextAt{8}{292.529}{640.908}{map}%
\TextAt{16}{318.543}{640.908}{\Character{102}}%
\TextAt{13}{324.521}{640.908}{M}%
\TextAt{16}{337.094}{640.908}{\Character{103}}%
\NativeArrowAt{346.383}{640.908}{11.955}{11.955}{\mapsto}%
\TextAt{9}{361.669}{640.908}{\Character{91}}%
\TextAt{8}{364.921}{640.908}{det}%
\TextAt{13}{383.172}{640.908}{M}%
\TextAt{9}{395.746}{640.908}{\Character{93}}%
\TextAt{8}{402.895}{640.908}{gives}%
\TextAt{8}{431.575}{640.908}{a}%
\TextAt{8}{441.331}{640.908}{homomorphism}%
\TextAt{8}{249.033}{614.893}{det}%
\TextAt{9}{268.613}{614.893}{\Character{58}}%
\TextAt{13}{275.186}{614.893}{F}%
\TextAt{17}{282.763}{613.099}{0}%
\TextAt{9}{287.496}{614.893}{\Character{40}}%
\Vector{q 0 G 0 g 0.478 w 220.048 -151.751 m 225.892 -151.751 l S Q}%
\TextAt{14}{292.048}{614.893}{o}%
\TextAt{9}{297.892}{614.893}{\Character{41}}%
\NativeArrowAt{305.765}{614.893}{11.955}{19.26}{\longrightarrow}%
\TextAt{8}{328.347}{614.893}{Pic}%
\TextAt{9}{344.762}{614.893}{\Character{40}}%
\Vector{q 0 G 0 g 0.478 w 277.314 -151.751 m 283.158 -151.751 l S Q}%
\TextAt{14}{349.314}{614.893}{o}%
\TextAt{9}{355.158}{614.893}{\Character{41}}%
\TextAt{13}{359.711}{614.893}{\Character{58}}%
\TextAt{8}{79.226}{588.877}{It}%
\TextAt{8}{91.742}{588.877}{maps}%
\TextAt{8}{122.215}{588.877}{the}%
\TextAt{8}{142.216}{588.877}{subgroup}%
\TextAt{13}{192.857}{588.877}{R}%
\TextAt{17}{201.767}{587.084}{0}%
\TextAt{9}{206.5}{588.877}{\Character{40}}%
\Vector{q 0 G 0 g 0.478 w 139.052 -177.767 m 144.896 -177.767 l S Q}%
\TextAt{14}{211.052}{588.877}{o}%
\TextAt{9}{216.896}{588.877}{\Character{41}}%
\TextAt{8}{225.19}{588.877}{to}%
\TextAt{9}{239.344}{588.877}{1}%
\TextAt{8}{245.197}{588.877}{,}%
\TextAt{8}{252.221}{588.877}{because}%
\TextAt{8}{295.367}{588.877}{it}%
\TextAt{8}{306.915}{588.877}{is}%
\TextAt{8}{318.522}{588.877}{generated}%
\TextAt{8}{371.698}{588.877}{by}%
\TextAt{8}{387.801}{588.877}{the}%
\TextAt{8}{407.801}{588.877}{elements}%
\TextAt{16}{455.502}{588.877}{\Character{102}}%
\TextAt{13}{461.479}{588.877}{M}%
\TextAt{3}{474.053}{593.215}{\Character{48}}%
\TextAt{16}{476.848}{588.877}{\Character{103}}%
\TextAt{16}{485.157}{588.877}{\Character{0}}%
\TextAt{16}{496.789}{588.877}{\Character{102}}%
\TextAt{13}{502.776}{588.877}{M}%
\TextAt{16}{515.35}{588.877}{\Character{103}}%
\TextAt{9}{523.665}{588.877}{\Character{43}}%
\TextAt{16}{79.226}{571.831}{\Character{102}}%
\TextAt{13}{85.204}{571.831}{M}%
\TextAt{3}{97.777}{576.169}{\Character{48}\Character{48}}%
\TextAt{16}{102.869}{571.831}{\Character{103}}%
\TextAt{8}{112.744}{571.831}{which}%
\TextAt{8}{146.231}{571.831}{arise}%
\TextAt{8}{173.608}{571.831}{from}%
\TextAt{8}{201.249}{571.831}{short}%
\TextAt{8}{231.232}{571.831}{exact}%
\TextAt{8}{262.124}{571.831}{sequences}%
\TextAt{9}{225.535}{545.815}{0}%
\NativeArrowAt{234.709}{545.815}{11.955}{19.26}{\longrightarrow}%
\TextAt{13}{257.291}{545.815}{M}%
\TextAt{3}{269.865}{550.751}{\Character{48}}%
\NativeArrowAt{275.981}{545.815}{11.955}{19.26}{\longrightarrow}%
\TextAt{13}{298.563}{545.815}{M}%
\NativeArrowAt{314.457}{545.815}{11.955}{19.26}{\longrightarrow}%
\TextAt{13}{337.039}{545.815}{M}%
\TextAt{3}{349.613}{550.751}{\Character{48}\Character{48}}%
\NativeArrowAt{358.025}{545.815}{11.955}{19.26}{\longrightarrow}%
\TextAt{9}{380.607}{545.815}{0}%
\TextAt{8}{79.226}{519.799}{of}%
\TextAt{8}{92.556}{519.799}{projective}%
\TextAt{8}{146.534}{519.799}{metrized}%
\TextAt{14}{194.653}{519.799}{o}%
\TextAt{8}{200.498}{519.799}{-modules}%
\TextAt{8}{250.303}{519.799}{and}%
\TextAt{8}{273.065}{519.799}{which,}%
\TextAt{8}{309.804}{519.799}{by}%
\TextAt{8}{326.063}{519.799}{(4.7),}%
\TextAt{8}{357.29}{519.799}{satisfy}%
\TextAt{8}{181.232}{493.784}{det}%
\TextAt{16}{197.491}{493.784}{\Character{102}}%
\TextAt{13}{203.469}{493.784}{M}%
\TextAt{16}{216.042}{493.784}{\Character{103}}%
\TextAt{9}{225.341}{493.784}{\Character{61}}%
\TextAt{9}{237.751}{493.784}{\Character{91}}%
\TextAt{8}{241.018}{493.784}{det}%
\TextAt{13}{259.269}{493.784}{M}%
\TextAt{9}{271.843}{493.784}{\Character{93}}%
\TextAt{9}{278.394}{493.784}{\Character{61}}%
\TextAt{9}{290.816}{493.784}{\Character{91}}%
\TextAt{8}{294.093}{493.784}{det}%
\TextAt{13}{312.344}{493.784}{M}%
\TextAt{3}{324.918}{498.72}{\Character{48}}%
\TextAt{16}{330.37}{493.784}{\Character{10}}%
\TextAt{8}{342.325}{493.784}{det}%
\TextAt{13}{360.576}{493.784}{M}%
\TextAt{3}{373.15}{498.72}{\Character{48}\Character{48}}%
\TextAt{9}{378.242}{493.784}{\Character{93}}%
\TextAt{9}{225.341}{473.749}{\Character{61}}%
\TextAt{9}{237.751}{473.749}{\Character{91}}%
\TextAt{8}{241.018}{473.749}{det}%
\TextAt{13}{259.269}{473.749}{M}%
\TextAt{3}{271.843}{478.685}{\Character{48}}%
\TextAt{9}{274.638}{473.749}{\Character{93}\Character{91}}%
\TextAt{8}{281.141}{473.749}{det}%
\TextAt{13}{299.393}{473.749}{M}%
\TextAt{3}{311.967}{478.685}{\Character{48}\Character{48}}%
\TextAt{9}{317.059}{473.749}{\Character{93}}%
\TextAt{9}{323.61}{473.749}{\Character{61}}%
\TextAt{8}{336.057}{473.749}{det}%
\TextAt{16}{352.316}{473.749}{\Character{102}}%
\TextAt{13}{358.293}{473.749}{M}%
\TextAt{3}{370.867}{478.685}{\Character{48}}%
\TextAt{16}{373.662}{473.749}{\Character{103}}%
\TextAt{8}{381.632}{473.749}{det}%
\TextAt{16}{397.891}{473.749}{\Character{102}}%
\TextAt{13}{403.869}{473.749}{M}%
\TextAt{3}{416.442}{478.685}{\Character{48}\Character{48}}%
\TextAt{16}{421.534}{473.749}{\Character{103}}%
\TextAt{13}{427.512}{473.749}{\Character{58}}%
\TextAt{8}{79.226}{447.733}{Thus}%
\TextAt{8}{109.352}{447.733}{we}%
\TextAt{8}{127.057}{447.733}{get}%
\TextAt{8}{147.046}{447.733}{an}%
\TextAt{8}{163.782}{447.733}{induced}%
\TextAt{8}{207.824}{447.733}{homomorphism}%
\TextAt{8}{290.972}{447.733}{det}%
\TextAt{9}{311.365}{447.733}{\Character{58}}%
\TextAt{13}{318.752}{447.733}{K}%
\TextAt{17}{328.711}{445.94}{0}%
\TextAt{9}{333.443}{447.733}{\Character{40}}%
\Vector{q 0 G 0 g 0.478 w 265.995 -318.911 m 271.84 -318.911 l S Q}%
\TextAt{14}{337.995}{447.733}{o}%
\TextAt{9}{343.84}{447.733}{\Character{41}}%
\NativeMathAt{352.527}{447.733}{11.955}{\rightarrow}%
\TextAt{8}{368.616}{447.733}{Pic}%
\TextAt{9}{385.031}{447.733}{\Character{40}}%
\Vector{q 0 G 0 g 0.478 w 317.583 -318.911 m 323.427 -318.911 l S Q}%
\TextAt{14}{389.583}{447.733}{o}%
\TextAt{9}{395.427}{447.733}{\Character{41}}%
\TextAt{8}{399.98}{447.733}{.}%
\TextAt{8}{409.86}{447.733}{It}%
\TextAt{8}{423.01}{447.733}{satisfies}%
\TextAt{8}{466.598}{447.733}{the}%
\TextAt{8}{487.232}{447.733}{following}%
\TextAt{8}{79.226}{430.687}{proposition.}%
\TextAt{12}{79.226}{405.528}{Proposition}%
\TextAt{12}{151.937}{405.528}{(6.5).}%
\TextAt{8}{187.253}{405.528}{(i)}%
\TextAt{8}{203.512}{405.528}{The}%
\TextAt{8}{227.566}{405.528}{canonical}%
\TextAt{8}{278.949}{405.528}{homomorphism}%
\TextAt{8}{260.178}{379.512}{Pic}%
\TextAt{9}{276.593}{379.512}{\Character{40}}%
\Vector{q 0 G 0 g 0.478 w 209.145 -387.132 m 214.989 -387.132 l S Q}%
\TextAt{14}{281.145}{379.512}{o}%
\TextAt{9}{286.989}{379.512}{\Character{41}}%
\NativeArrowAt{294.863}{379.512}{11.955}{19.26}{\longrightarrow}%
\TextAt{13}{317.445}{379.512}{K}%
\TextAt{17}{327.404}{377.719}{0}%
\TextAt{9}{332.136}{379.512}{\Character{40}}%
\Vector{q 0 G 0 g 0.478 w 264.688 -387.132 m 270.533 -387.132 l S Q}%
\TextAt{14}{336.688}{379.512}{o}%
\TextAt{9}{342.533}{379.512}{\Character{41}}%
\TextAt{3}{347.085}{384.448}{\Character{3}}%
\TextAt{8}{79.226}{353.496}{is}%
\TextAt{8}{90.99}{353.496}{injective.}%
\TextAt{8}{97.159}{333.461}{(ii)}%
\TextAt{8}{116.67}{333.461}{The}%
\TextAt{8}{140.724}{333.461}{restriction}%
\TextAt{8}{196.722}{333.461}{of}%
\TextAt{8}{210.052}{333.461}{det}%
\TextAt{8}{230.208}{333.461}{to}%
\TextAt{13}{244.518}{333.461}{I}%
\TextAt{8}{250.621}{333.461}{,}%
\TextAt{8}{259.611}{307.446}{det}%
\TextAt{9}{279.191}{307.446}{\Character{58}}%
\TextAt{13}{285.764}{307.446}{I}%
\NativeArrowAt{295.187}{307.446}{11.955}{19.26}{\longrightarrow}%
\TextAt{8}{317.769}{307.446}{Pic}%
\TextAt{9}{334.184}{307.446}{\Character{40}}%
\Vector{q 0 G 0 g 0.478 w 266.736 -459.198 m 272.58 -459.198 l S Q}%
\TextAt{14}{338.736}{307.446}{o}%
\TextAt{9}{344.58}{307.446}{\Character{41}}%
\TextAt{13}{349.133}{307.446}{\Character{59}}%
\TextAt{8}{79.226}{281.43}{is}%
\TextAt{8}{90.99}{281.43}{an}%
\TextAt{8}{107.249}{281.43}{isomorphism.}%
\TextAt{19}{79.226}{248.157}{Proof.}%
\TextAt{8}{116}{248.157}{(i)}%
\TextAt{8}{132.259}{248.157}{The}%
\TextAt{8}{156.313}{248.157}{composite}%
\TextAt{8}{211.331}{248.157}{of}%
\TextAt{8}{224.661}{248.157}{both}%
\TextAt{8}{252.301}{248.157}{mappings}%
\TextAt{8}{232.005}{222.141}{Pic}%
\TextAt{9}{248.419}{222.141}{\Character{40}}%
\Vector{q 0 G 0 g 0.478 w 180.971 -544.502 m 186.816 -544.502 l S Q}%
\TextAt{14}{252.971}{222.141}{o}%
\TextAt{9}{258.816}{222.141}{\Character{41}}%
\NativeArrowAt{266.689}{222.141}{11.955}{19.26}{\longrightarrow}%
\TextAt{13}{289.271}{222.141}{K}%
\TextAt{17}{299.23}{220.348}{0}%
\TextAt{9}{303.962}{222.141}{\Character{40}}%
\Vector{q 0 G 0 g 0.478 w 236.515 -544.502 m 242.359 -544.502 l S Q}%
\TextAt{14}{308.515}{222.141}{o}%
\TextAt{9}{314.359}{222.141}{\Character{41}}%
\TextAt{3}{318.911}{227.078}{\Character{3}}%
\TextAt{7}{329.552}{228.918}{det}%
\NativeArrowAt{326.964}{222.141}{11.955}{18.339}{\longrightarrow}%
\TextAt{8}{348.628}{222.141}{Pic}%
\TextAt{9}{365.042}{222.141}{\Character{40}}%
\Vector{q 0 G 0 g 0.478 w 297.594 -544.502 m 303.439 -544.502 l S Q}%
\TextAt{14}{369.594}{222.141}{o}%
\TextAt{9}{375.439}{222.141}{\Character{41}}%
\TextAt{8}{79.226}{196.126}{is}%
\TextAt{8}{91.024}{196.126}{the}%
\TextAt{8}{111.215}{196.126}{identity,}%
\TextAt{8}{156.782}{196.126}{since}%
\TextAt{8}{185.485}{196.126}{for}%
\TextAt{8}{203.404}{196.126}{an}%
\TextAt{8}{219.697}{196.126}{invertible}%
\TextAt{8}{271.437}{196.126}{metrized}%
\TextAt{14}{319.591}{196.126}{o}%
\TextAt{8}{325.435}{196.126}{-module}%
\TextAt{13}{370.66}{196.126}{M}%
\TextAt{8}{383.233}{196.126}{,}%
\TextAt{8}{390.425}{196.126}{one}%
\TextAt{8}{411.919}{196.126}{clearly}%
\TextAt{8}{449.349}{196.126}{has}%
\TextAt{8}{470.256}{196.126}{det}%
\TextAt{13}{488.508}{196.126}{M}%
\TextAt{9}{504.46}{196.126}{\Character{61}}%
\TextAt{13}{516.944}{196.126}{M}%
\TextAt{8}{529.518}{196.126}{.}%
\TextAt{8}{79.226}{179.08}{This}%
\TextAt{8}{105.946}{179.08}{gives}%
\TextAt{8}{134.626}{179.08}{(i).}%
\TextAt{8}{97.159}{159.045}{(ii)}%
\TextAt{8}{116.67}{159.045}{Next,}%
\TextAt{8}{148.53}{159.045}{viewing}%
\TextAt{8}{191.127}{159.045}{the}%
\TextAt{8}{211.283}{159.045}{elements}%
\TextAt{8}{259.14}{159.045}{of}%
\TextAt{8}{272.47}{159.045}{Pic}%
\TextAt{9}{288.884}{159.045}{\Character{40}}%
\Vector{q 0 G 0 g 0.478 w 221.436 -607.599 m 227.281 -607.599 l S Q}%
\TextAt{14}{293.436}{159.045}{o}%
\TextAt{9}{299.281}{159.045}{\Character{41}}%
\TextAt{8}{307.73}{159.045}{as}%
\TextAt{8}{322.1}{159.045}{elements}%
\TextAt{8}{369.957}{159.045}{of}%
\TextAt{13}{383.287}{159.045}{K}%
\TextAt{17}{393.246}{157.252}{0}%
\TextAt{9}{397.978}{159.045}{\Character{40}}%
\Vector{q 0 G 0 g 0.478 w 330.531 -607.599 m 336.375 -607.599 l S Q}%
\TextAt{14}{402.531}{159.045}{o}%
\TextAt{9}{408.375}{159.045}{\Character{41}}%
\TextAt{8}{412.927}{159.045}{,}%
\TextAt{13}{218.221}{133.029}{\Character{14}}%
\TextAt{9}{227.174}{133.029}{\Character{58}}%
\TextAt{8}{233.747}{133.029}{Pic}%
\TextAt{9}{250.161}{133.029}{\Character{40}}%
\Vector{q 0 G 0 g 0.478 w 182.714 -633.614 m 188.558 -633.614 l S Q}%
\TextAt{14}{254.714}{133.029}{o}%
\TextAt{9}{260.558}{133.029}{\Character{41}}%
\NativeArrowAt{268.431}{133.029}{11.955}{19.26}{\longrightarrow}%
\TextAt{13}{291.013}{133.029}{I}%
\TextAt{13}{297.098}{133.029}{\Character{59}}%
\TextAt{13}{326.245}{133.029}{\Character{14}}%
\TextAt{9}{331.903}{133.029}{\Character{40}}%
\TextAt{13}{336.456}{133.029}{x}%
\TextAt{9}{343.108}{133.029}{\Character{41}}%
\TextAt{9}{350.963}{133.029}{\Character{61}}%
\TextAt{13}{363.406}{133.029}{x}%
\TextAt{16}{372.715}{133.029}{\Character{0}}%
\TextAt{9}{384.67}{133.029}{1}%
\TextAt{13}{390.523}{133.029}{\Character{59}}%
\TextAt{8}{79.226}{107.013}{gives}%
\TextAt{8}{109.312}{107.013}{us}%
\TextAt{8}{125.733}{107.013}{an}%
\TextAt{8}{143.398}{107.013}{inverse}%
\TextAt{8}{183.562}{107.013}{mapping}%
\TextAt{8}{233.099}{107.013}{to}%
\TextAt{8}{248.815}{107.013}{det}%
\TextAt{9}{270.79}{107.013}{\Character{58}}%
\TextAt{13}{279.758}{107.013}{I}%
\NativeMathAt{291.577}{107.013}{11.955}{\rightarrow}%
\TextAt{8}{309.248}{107.013}{Pic}%
\TextAt{9}{325.662}{107.013}{\Character{40}}%
\Vector{q 0 G 0 g 0.478 w 258.215 -659.63 m 264.059 -659.63 l S Q}%
\TextAt{14}{330.215}{107.013}{o}%
\TextAt{9}{336.059}{107.013}{\Character{41}}%
\TextAt{8}{340.611}{107.013}{.}%
\TextAt{8}{353.276}{107.013}{In}%
\TextAt{8}{369.303}{107.013}{fact,}%
\TextAt{8}{397.397}{107.013}{one}%
\TextAt{8}{420.262}{107.013}{has}%
\TextAt{8}{442.541}{107.013}{det}%
\TextAt{16}{460.793}{107.013}{\Character{14}}%
\TextAt{13}{466.77}{107.013}{\Character{14}}%
\TextAt{9}{478.119}{107.013}{\Character{61}}%
\TextAt{8}{492.94}{107.013}{id}%
\TextAt{8}{507.998}{107.013}{since}%
\TextAt{8}{79.226}{89.967}{det}%
\TextAt{9}{95.485}{89.967}{\Character{40}\Character{91}}%
\TextAt{14}{103.289}{89.967}{a}%
\TextAt{9}{109.272}{89.967}{\Character{93}}%
\TextAt{16}{112.887}{89.967}{\Character{0}}%
\TextAt{9}{122.549}{89.967}{1\Character{41}}%
\TextAt{9}{136.25}{89.967}{\Character{61}}%
\TextAt{8}{148.701}{89.967}{det}%
\TextAt{9}{164.96}{89.967}{\Character{91}}%
\TextAt{14}{168.211}{89.967}{a}%
\TextAt{9}{174.195}{89.967}{\Character{93}}%
\TextAt{9}{180.746}{89.967}{\Character{61}}%
\TextAt{9}{193.168}{89.967}{\Character{91}}%
\TextAt{14}{196.444}{89.967}{a}%
\TextAt{9}{202.427}{89.967}{\Character{93}}%
\TextAt{8}{205.679}{89.967}{,}%
\TextAt{8}{211.931}{89.967}{and}%
\TextAt{13}{233.572}{89.967}{\Character{14}}%
\TextAt{16}{239.568}{89.967}{\Character{14}}%
\TextAt{8}{245.909}{89.967}{det}%
\TextAt{9}{265.489}{89.967}{\Character{61}}%
\TextAt{8}{277.915}{89.967}{id}%
\TextAt{8}{290.446}{89.967}{since}%
\TextAt{13}{317.993}{89.967}{\Character{14}}%
\TextAt{9}{323.626}{89.967}{\Character{40}}%
\TextAt{8}{328.178}{89.967}{det}%
\TextAt{9}{344.437}{89.967}{\Character{40}\Character{91}}%
\TextAt{14}{352.241}{89.967}{a}%
\TextAt{9}{358.225}{89.967}{\Character{93}}%
\TextAt{16}{361.84}{89.967}{\Character{0}}%
\TextAt{9}{371.501}{89.967}{1\Character{41}\Character{41}}%
\TextAt{9}{389.745}{89.967}{\Character{61}}%
\TextAt{13}{402.205}{89.967}{\Character{14}}%
\TextAt{9}{407.838}{89.967}{\Character{40}}%
\TextAt{8}{412.391}{89.967}{det}%
\TextAt{9}{428.65}{89.967}{\Character{91}}%
\TextAt{14}{431.901}{89.967}{a}%
\TextAt{9}{437.884}{89.967}{\Character{93}\Character{41}}%
\TextAt{9}{448.978}{89.967}{\Character{61}}%
\TextAt{13}{461.435}{89.967}{\Character{14}}%
\TextAt{9}{467.068}{89.967}{\Character{40}\Character{91}}%
\TextAt{14}{474.872}{89.967}{a}%
\TextAt{9}{480.855}{89.967}{\Character{93}\Character{41}}%
\TextAt{9}{491.949}{89.967}{\Character{61}}%
\TextAt{9}{504.371}{89.967}{\Character{91}}%
\TextAt{14}{507.657}{89.967}{a}%
\TextAt{9}{513.64}{89.967}{\Character{93}}%
\TextAt{16}{517.255}{89.967}{\Character{0}}%
\TextAt{9}{526.917}{89.967}{1}%
\TextAt{8}{79.226}{72.921}{and}%
\TextAt{8}{101.578}{72.921}{because}%
\TextAt{8}{144.469}{72.921}{of}%
\TextAt{8}{157.388}{72.921}{the}%
\TextAt{8}{177.133}{72.921}{fact}%
\TextAt{8}{199.808}{72.921}{that}%
\TextAt{13}{224.766}{72.921}{I}%
\TextAt{8}{234.355}{72.921}{is}%
\TextAt{8}{245.708}{72.921}{generated}%
\TextAt{8}{298.629}{72.921}{by}%
\TextAt{8}{314.477}{72.921}{elements}%
\TextAt{8}{361.923}{72.921}{of}%
\TextAt{8}{374.842}{72.921}{the}%
\TextAt{8}{394.588}{72.921}{form}%
\TextAt{9}{421.817}{72.921}{\Character{91}}%
\TextAt{14}{425.069}{72.921}{a}%
\TextAt{9}{431.052}{72.921}{\Character{93}}%
\TextAt{16}{436.12}{72.921}{\Character{0}}%
\TextAt{9}{447.235}{72.921}{1}%
\TextAt{8}{456.574}{72.921}{(see}%
\TextAt{8}{479.631}{72.921}{(6.1)).}%
\Vector{q 0 G 0 g 0.399 w 452.334 -701.079 m 452.334 -693.009 l S Q}%
\Vector{q 0 G 0 g 0.399 w 452.534 -693.208 m 459.707 -693.208 l S Q}%
\Vector{q 0 G 0 g 0.399 w 452.534 -700.879 m 459.707 -700.879 l S Q}%
\Vector{q 0 G 0 g 0.399 w 459.906 -701.079 m 459.906 -693.009 l S Q}%
\TextAt{8}{297.211}{43.033}{225}%
\Destination{78.23}{809.815}{nav.page.225}%
\Destination{79.36}{776.155}{nav.print.245}%
\Destination{79.36}{419.555}{nav.number.3.6.5}%
\Hotlink{84.405}{672.755}{6.858}{11.047}{nav.section.3.4}%
\Hotlink{325.563}{517.555}{25.078}{11.047}{nav.number.3.4.7}%
\Hotlink{479.131}{70.755}{25.078}{11.047}{nav.number.3.6.1}%
}

% Page 238 (body page 226)
\Page{%
\TextAt{2}{79.226}{796.943}{III.}%
\TextAt{2}{98.317}{796.943}{Riemann–Roch}%
\TextAt{2}{174.375}{796.943}{Theory}%
\Vector{q 0 G 0 g 0.299 w 7.226 18.729 m 460.769 18.729 l S Q}%
\TextAt{8}{97.159}{761.124}{From}%
\TextAt{8}{127.884}{761.124}{the}%
\TextAt{8}{148.04}{761.124}{isomorphism}%
\TextAt{8}{216.46}{761.124}{det}%
\TextAt{9}{236.039}{761.124}{\Character{58}}%
\TextAt{13}{242.612}{761.124}{I}%
\NativeMathAt{254.623}{767.901}{7.97}{\sim}%
\NativeArrowAt{252.035}{761.124}{11.955}{13.175}{\longrightarrow}%
\TextAt{8}{268.529}{761.124}{Pic}%
\TextAt{9}{284.944}{761.124}{\Character{40}}%
\Vector{q 0 G 0 g 0.478 w 217.496 -5.52 m 223.34 -5.52 l S Q}%
\TextAt{14}{289.496}{761.124}{o}%
\TextAt{9}{295.34}{761.124}{\Character{41}}%
\TextAt{8}{299.893}{761.124}{,}%
\TextAt{8}{307.042}{761.124}{we}%
\TextAt{8}{324.269}{761.124}{now}%
\TextAt{8}{348.658}{761.124}{obtain}%
\TextAt{8}{385.085}{761.124}{an}%
\TextAt{8}{401.344}{761.124}{isomorphism}%
\TextAt{8}{248.095}{733.294}{gr}%
\TextAt{13}{260.5}{733.294}{K}%
\TextAt{17}{270.46}{731.501}{0}%
\TextAt{9}{275.192}{733.294}{\Character{40}}%
\Vector{q 0 G 0 g 0.478 w 207.744 -33.349 m 213.588 -33.349 l S Q}%
\TextAt{14}{279.744}{733.294}{o}%
\TextAt{9}{285.588}{733.294}{\Character{41}}%
\NativeMathAt{296.049}{740.071}{7.97}{\sim}%
\NativeArrowAt{293.462}{733.294}{11.955}{13.175}{\longrightarrow}%
\TextAt{18}{309.955}{733.294}{Z}%
\TextAt{16}{320.582}{733.294}{\Character{8}}%
\TextAt{8}{332.537}{733.294}{Pic}%
\TextAt{9}{348.952}{733.294}{\Character{40}}%
\Vector{q 0 G 0 g 0.478 w 281.504 -33.349 m 287.348 -33.349 l S Q}%
\TextAt{14}{353.504}{733.294}{o}%
\TextAt{9}{359.348}{733.294}{\Character{41}}%
\TextAt{8}{79.226}{705.465}{and}%
\TextAt{8}{101.989}{705.465}{the}%
\TextAt{8}{122.145}{705.465}{composite}%
\TextAt{13}{211.612}{688.419}{K}%
\TextAt{17}{221.572}{686.626}{0}%
\TextAt{9}{226.304}{688.419}{\Character{40}}%
\Vector{q 0 G 0 g 0.478 w 158.856 -78.224 m 164.7 -78.224 l S Q}%
\TextAt{14}{230.856}{688.419}{o}%
\TextAt{9}{236.7}{688.419}{\Character{41}}%
\TextAt{7}{247.161}{695.196}{ch}%
\NativeArrowAt{244.574}{688.419}{11.955}{15.052}{\longrightarrow}%
\TextAt{8}{262.945}{688.419}{gr}%
\TextAt{13}{275.351}{688.419}{K}%
\TextAt{17}{285.31}{686.626}{0}%
\TextAt{9}{290.042}{688.419}{\Character{40}}%
\Vector{q 0 G 0 g 0.478 w 222.594 -78.224 m 228.439 -78.224 l S Q}%
\TextAt{14}{294.594}{688.419}{o}%
\TextAt{9}{300.439}{688.419}{\Character{41}}%
\TextAt{7}{310.899}{695.196}{id}%
\TextAt{3}{319.364}{695.196}{\Character{8}}%
\TextAt{7}{327.362}{695.196}{det}%
\NativeArrowAt{308.312}{688.419}{11.955}{34.778}{\longrightarrow}%
\TextAt{18}{346.438}{688.419}{Z}%
\TextAt{16}{357.065}{688.419}{\Character{8}}%
\TextAt{8}{369.02}{688.419}{Pic}%
\TextAt{9}{385.434}{688.419}{\Character{40}}%
\Vector{q 0 G 0 g 0.478 w 317.987 -78.224 m 323.831 -78.224 l S Q}%
\TextAt{14}{389.987}{688.419}{o}%
\TextAt{9}{395.831}{688.419}{\Character{41}}%
\TextAt{8}{79.226}{665.4}{will}%
\TextAt{8}{101.811}{665.4}{again}%
\TextAt{8}{133.518}{665.4}{be}%
\TextAt{8}{149.922}{665.4}{called}%
\TextAt{8}{183.565}{665.4}{the}%
\TextAt{8}{204.201}{665.4}{Chern}%
\TextAt{8}{239.793}{665.4}{character}%
\TextAt{8}{291.334}{665.4}{of}%
\TextAt{13}{305.143}{665.4}{K}%
\TextAt{17}{315.103}{663.606}{0}%
\TextAt{9}{319.835}{665.4}{\Character{40}}%
\Vector{q 0 G 0 g 0.478 w 252.387 -101.244 m 258.231 -101.244 l S Q}%
\TextAt{14}{324.387}{665.4}{o}%
\TextAt{9}{330.231}{665.4}{\Character{41}}%
\TextAt{8}{334.784}{665.4}{.}%
\TextAt{8}{344.671}{665.4}{Observing}%
\TextAt{8}{400.826}{665.4}{that}%
\TextAt{8}{426.675}{665.4}{det}%
\TextAt{9}{442.934}{665.4}{\Character{40}}%
\TextAt{13}{447.486}{665.4}{c}%
\TextAt{17}{452.524}{663.606}{1}%
\TextAt{9}{457.257}{665.4}{\Character{40}}%
\TextAt{13}{461.809}{665.4}{\Character{24}}%
\TextAt{9}{467.499}{665.4}{\Character{41}\Character{41}}%
\TextAt{9}{480.721}{665.4}{\Character{61}}%
\TextAt{8}{493.985}{665.4}{det}%
\TextAt{9}{510.245}{665.4}{\Character{40}}%
\TextAt{13}{514.797}{665.4}{\Character{24}}%
\TextAt{16}{523.471}{665.4}{\Character{0}}%
\TextAt{8}{79.226}{648.354}{rk}%
\TextAt{9}{89.962}{648.354}{\Character{40}}%
\TextAt{13}{94.514}{648.354}{\Character{24}}%
\TextAt{9}{100.205}{648.354}{\Character{41}}%
\TextAt{16}{107.414}{648.354}{\Character{1}}%
\TextAt{9}{113.391}{648.354}{1\Character{41}}%
\TextAt{9}{127.092}{648.354}{\Character{61}}%
\TextAt{8}{139.543}{648.354}{det}%
\TextAt{9}{155.802}{648.354}{\Character{40}}%
\TextAt{13}{160.354}{648.354}{\Character{24}}%
\TextAt{9}{166.045}{648.354}{\Character{41}}%
\TextAt{8}{170.597}{648.354}{,}%
\TextAt{8}{177.746}{648.354}{this}%
\TextAt{8}{200.569}{648.354}{yields}%
\TextAt{8}{233.469}{648.354}{the}%
\TextAt{8}{253.626}{648.354}{explicit}%
\TextAt{8}{294.919}{648.354}{description}%
\TextAt{8}{354.814}{648.354}{of}%
\TextAt{8}{368.144}{648.354}{the}%
\TextAt{8}{388.301}{648.354}{Grothendieck}%
\TextAt{8}{460.57}{648.354}{group}%
\TextAt{13}{493.745}{648.354}{K}%
\TextAt{17}{503.704}{646.56}{0}%
\TextAt{9}{508.437}{648.354}{\Character{40}}%
\Vector{q 0 G 0 g 0.478 w 440.989 -118.29 m 446.833 -118.29 l S Q}%
\TextAt{14}{512.989}{648.354}{o}%
\TextAt{9}{518.833}{648.354}{\Character{41}}%
\TextAt{8}{523.386}{648.354}{:}%
\TextAt{12}{79.226}{622.676}{Theorem}%
\TextAt{12}{136.24}{622.676}{(6.6).}%
\TextAt{8}{171.556}{622.676}{The}%
\TextAt{8}{195.61}{622.676}{Chern}%
\TextAt{8}{230.722}{622.676}{character}%
\TextAt{8}{281.783}{622.676}{gives}%
\TextAt{8}{310.463}{622.676}{an}%
\TextAt{8}{326.722}{622.676}{isomorphism}%
\TextAt{8}{170.575}{594.847}{ch}%
\TextAt{9}{185.6}{594.847}{\Character{58}}%
\TextAt{13}{192.172}{594.847}{K}%
\TextAt{17}{202.131}{593.054}{0}%
\TextAt{9}{206.864}{594.847}{\Character{40}}%
\Vector{q 0 G 0 g 0.478 w 139.416 -171.797 m 145.26 -171.797 l S Q}%
\TextAt{14}{211.416}{594.847}{o}%
\TextAt{9}{217.26}{594.847}{\Character{41}}%
\NativeMathAt{227.721}{601.624}{7.97}{\sim}%
\NativeArrowAt{225.133}{594.847}{11.955}{13.175}{\longrightarrow}%
\TextAt{18}{241.627}{594.847}{Z}%
\TextAt{16}{252.254}{594.847}{\Character{8}}%
\TextAt{8}{264.209}{594.847}{Pic}%
\TextAt{9}{280.624}{594.847}{\Character{40}}%
\Vector{q 0 G 0 g 0.478 w 213.176 -171.797 m 219.02 -171.797 l S Q}%
\TextAt{14}{285.176}{594.847}{o}%
\TextAt{9}{291.02}{594.847}{\Character{41}}%
\TextAt{13}{295.572}{594.847}{\Character{59}}%
\TextAt{8}{324.727}{594.847}{ch}%
\TextAt{9}{336.431}{594.847}{\Character{40}}%
\TextAt{13}{340.983}{594.847}{\Character{24}}%
\TextAt{9}{346.674}{594.847}{\Character{41}}%
\TextAt{9}{354.529}{594.847}{\Character{61}}%
\TextAt{8}{366.972}{594.847}{rk}%
\TextAt{9}{377.708}{594.847}{\Character{40}}%
\TextAt{13}{382.261}{594.847}{\Character{24}}%
\TextAt{9}{387.951}{594.847}{\Character{41}}%
\TextAt{16}{395.16}{594.847}{\Character{8}}%
\TextAt{8}{407.115}{594.847}{det}%
\TextAt{9}{423.374}{594.847}{\Character{40}}%
\TextAt{13}{427.926}{594.847}{\Character{24}}%
\TextAt{9}{433.617}{594.847}{\Character{41}}%
\TextAt{13}{438.169}{594.847}{\Character{58}}%
\TextAt{8}{97.159}{558.386}{The}%
\TextAt{8}{121.308}{558.386}{expert}%
\TextAt{8}{157.82}{558.386}{should}%
\TextAt{8}{195.048}{558.386}{note}%
\TextAt{8}{221.159}{558.386}{that}%
\TextAt{8}{246.623}{558.386}{this}%
\TextAt{8}{269.541}{558.386}{homomorphism}%
\TextAt{8}{352.307}{558.386}{is}%
\TextAt{8}{364.167}{558.386}{a}%
\TextAt{8}{374.018}{558.386}{realization}%
\TextAt{8}{431.355}{558.386}{map}%
\TextAt{8}{457.465}{558.386}{from}%
\TextAt{13}{485.202}{558.386}{K}%
\TextAt{8}{496.019}{558.386}{-theory}%
\TextAt{8}{79.226}{541.34}{into}%
\TextAt{8}{104.045}{541.34}{Chow-theory.}%
\TextAt{8}{180.436}{541.34}{Identifying}%
\TextAt{8}{240.69}{541.34}{Pic}%
\TextAt{9}{257.104}{541.34}{\Character{40}}%
\Vector{q 0 G 0 g 0.478 w 189.657 -225.303 m 195.501 -225.303 l S Q}%
\TextAt{14}{261.657}{541.34}{o}%
\TextAt{9}{267.501}{541.34}{\Character{41}}%
\TextAt{8}{277.027}{541.34}{with}%
\TextAt{8}{304.763}{541.34}{the}%
\TextAt{8}{325.996}{541.34}{divisor}%
\TextAt{8}{365.185}{541.34}{class}%
\TextAt{8}{393.698}{541.34}{group}%
\TextAt{13}{427.95}{541.34}{C}%
\TextAt{13}{437.168}{541.34}{H}%
\TextAt{17}{447.806}{545.679}{1}%
\TextAt{9}{452.539}{541.34}{\Character{40}}%
\Vector{q 0 G 0 g 0.478 w 385.091 -225.303 m 390.935 -225.303 l S Q}%
\TextAt{14}{457.091}{541.34}{o}%
\TextAt{9}{462.935}{541.34}{\Character{41}}%
\TextAt{8}{467.487}{541.34}{,}%
\TextAt{8}{475.982}{541.34}{we}%
\TextAt{8}{494.286}{541.34}{have}%
\TextAt{8}{522.357}{541.34}{to}%
\TextAt{8}{79.226}{524.294}{view}%
\TextAt{18}{106.209}{524.294}{Z}%
\TextAt{16}{116.836}{524.294}{\Character{8}}%
\TextAt{8}{128.791}{524.294}{Pic}%
\TextAt{9}{145.205}{524.294}{\Character{40}}%
\Vector{q 0 G 0 g 0.478 w 77.758 -242.349 m 83.602 -242.349 l S Q}%
\TextAt{14}{149.758}{524.294}{o}%
\TextAt{9}{155.602}{524.294}{\Character{41}}%
\TextAt{8}{164.052}{524.294}{as}%
\TextAt{8}{178.422}{524.294}{the}%
\TextAt{8}{198.578}{524.294}{“replete”}%
\TextAt{19}{247.989}{524.294}{Chow}%
\TextAt{19}{280.268}{524.294}{ring}%
\TextAt{13}{304.979}{524.294}{C}%
\TextAt{13}{314.197}{524.294}{H}%
\TextAt{9}{324.835}{524.294}{\Character{40}}%
\Vector{q 0 G 0 g 0.478 w 257.388 -242.349 m 263.232 -242.349 l S Q}%
\TextAt{14}{329.388}{524.294}{o}%
\TextAt{9}{335.232}{524.294}{\Character{41}}%
\TextAt{8}{339.784}{524.294}{.}%
\TextAt{15}{79.226}{477.581}{§}%
\TextAt{15}{93.785}{477.581}{7}%
\TextAt{15}{120.693}{477.581}{Grothendieck–Riemann–Roch}%
\TextAt{8}{79.226}{445.695}{We}%
\TextAt{8}{98.313}{445.695}{now}%
\TextAt{8}{121.644}{445.695}{consider}%
\TextAt{8}{166.172}{445.695}{a}%
\TextAt{8}{174.87}{445.695}{finite}%
\TextAt{8}{203.724}{445.695}{extension}%
\TextAt{13}{254.432}{445.695}{L}%
\TextAt{16}{265.718}{445.695}{\Character{106}}%
\TextAt{13}{272.359}{445.695}{K}%
\TextAt{8}{286.017}{445.695}{of}%
\TextAt{8}{298.29}{445.695}{algebraic}%
\TextAt{8}{346.667}{445.695}{number}%
\TextAt{8}{388.205}{445.695}{fields}%
\TextAt{8}{417.12}{445.695}{and}%
\TextAt{8}{438.825}{445.695}{study}%
\TextAt{8}{470.022}{445.695}{the}%
\TextAt{8}{489.121}{445.695}{relations}%
\TextAt{8}{79.226}{428.649}{between}%
\TextAt{8}{124.622}{428.649}{the}%
\TextAt{8}{144.983}{428.649}{Grothendieck}%
\TextAt{8}{217.457}{428.649}{groups}%
\TextAt{8}{255.453}{428.649}{of}%
\TextAt{13}{268.988}{428.649}{L}%
\TextAt{8}{281.055}{428.649}{and}%
\TextAt{13}{304.023}{428.649}{K}%
\TextAt{8}{314.84}{428.649}{.}%
\TextAt{8}{323.904}{428.649}{Let}%
\TextAt{14}{345.079}{428.649}{o}%
\TextAt{8}{350.923}{428.649}{,}%
\TextAt{8}{358.328}{428.649}{resp.}%
\TextAt{16}{388.266}{428.649}{\Character{79}}%
\TextAt{8}{398.116}{428.649}{,}%
\TextAt{8}{405.521}{428.649}{be}%
\TextAt{8}{421.651}{428.649}{the}%
\TextAt{8}{442.012}{428.649}{ring}%
\TextAt{8}{466.283}{428.649}{of}%
\TextAt{8}{479.818}{428.649}{integers}%
\TextAt{8}{523.337}{428.649}{of}%
\TextAt{13}{79.226}{411.603}{K}%
\TextAt{8}{90.044}{411.603}{,}%
\TextAt{8}{97.28}{411.603}{resp.}%
\TextAt{13}{126.811}{411.603}{L}%
\TextAt{8}{134.775}{411.603}{,}%
\TextAt{8}{142.012}{411.603}{and}%
\TextAt{8}{164.844}{411.603}{write}%
\TextAt{13}{194.825}{411.603}{X}%
\TextAt{9}{205.481}{411.603}{\Character{40}}%
\TextAt{18}{210.033}{411.603}{C}%
\TextAt{9}{218.667}{411.603}{\Character{41}}%
\TextAt{9}{226.641}{411.603}{\Character{61}}%
\TextAt{8}{239.203}{411.603}{Hom}%
\TextAt{9}{263.592}{411.603}{\Character{40}}%
\TextAt{13}{268.144}{411.603}{K}%
\TextAt{13}{278.306}{411.603}{\Character{59}}%
\TextAt{18}{283.555}{411.603}{C}%
\TextAt{9}{292.19}{411.603}{\Character{41}}%
\TextAt{8}{296.742}{411.603}{,}%
\TextAt{13}{303.978}{411.603}{Y}%
\TextAt{9}{313.371}{411.603}{\Character{40}}%
\TextAt{18}{317.924}{411.603}{C}%
\TextAt{9}{326.558}{411.603}{\Character{41}}%
\TextAt{9}{334.532}{411.603}{\Character{61}}%
\TextAt{8}{347.094}{411.603}{Hom}%
\TextAt{9}{371.483}{411.603}{\Character{40}}%
\TextAt{13}{376.035}{411.603}{L\Character{59}}%
\TextAt{18}{389.244}{411.603}{C}%
\TextAt{9}{397.878}{411.603}{\Character{41}}%
\TextAt{8}{402.43}{411.603}{.}%
\TextAt{8}{411.087}{411.603}{The}%
\TextAt{8}{435.211}{411.603}{inclusion}%
\TextAt{13}{484.117}{411.603}{i}%
\TextAt{9}{491.55}{411.603}{\Character{58}}%
\TextAt{14}{498.241}{411.603}{o}%
\NativeMathAt{507.525}{411.603}{11.955}{\rightarrow}%
\TextAt{16}{522.911}{411.603}{\Character{79}}%
\TextAt{8}{79.226}{394.557}{and}%
\TextAt{8}{101.989}{394.557}{the}%
\TextAt{8}{122.145}{394.557}{surjection}%
\TextAt{13}{175.86}{394.557}{Y}%
\TextAt{9}{185.253}{394.558}{\Character{40}}%
\TextAt{18}{189.805}{394.557}{C}%
\TextAt{9}{198.44}{394.558}{\Character{41}}%
\NativeMathAt{206.313}{394.557}{11.955}{\rightarrow}%
\TextAt{13}{221.589}{394.558}{X}%
\TextAt{9}{232.244}{394.557}{\Character{40}}%
\TextAt{18}{236.796}{394.558}{C}%
\TextAt{9}{245.431}{394.557}{\Character{41}}%
\TextAt{8}{249.983}{394.557}{,}%
\TextAt{13}{257.132}{394.557}{\Character{27}}%
\NativeArrowAt{267.535}{394.558}{11.955}{11.955}{\mapsto}%
\TextAt{13}{282.811}{394.557}{\Character{27}}%
\TextAt{16}{289.894}{394.558}{\Character{106}}%
\TextAt{21}{293.215}{392.764}{K}%
\TextAt{8}{301.417}{394.557}{,}%
\TextAt{8}{308.566}{394.557}{give}%
\TextAt{8}{332.632}{394.557}{two}%
\TextAt{8}{354.749}{394.557}{canonical}%
\TextAt{8}{406.133}{394.557}{homomorphisms}%
\TextAt{13}{163.221}{366.728}{i}%
\TextAt{3}{167.214}{371.664}{\Character{3}}%
\TextAt{9}{175.268}{366.728}{\Character{58}}%
\TextAt{13}{181.84}{366.728}{K}%
\TextAt{17}{191.799}{364.935}{0}%
\TextAt{9}{196.531}{366.728}{\Character{40}}%
\Vector{q 0 G 0 g 0.478 w 129.084 -399.915 m 134.928 -399.915 l S Q}%
\TextAt{14}{201.084}{366.728}{o}%
\TextAt{9}{206.928}{366.728}{\Character{41}}%
\NativeArrowAt{214.801}{366.728}{11.955}{19.26}{\longrightarrow}%
\TextAt{13}{237.383}{366.728}{K}%
\TextAt{17}{247.342}{364.935}{0}%
\TextAt{9}{252.075}{366.728}{\Character{40}}%
\Vector{q 0 G 0 g 0.478 w 184.627 -397.429 m 194.477 -397.429 l S Q}%
\TextAt{16}{256.627}{366.728}{\Character{79}}%
\TextAt{9}{266.477}{366.728}{\Character{41}}%
\TextAt{8}{294.939}{366.728}{and}%
\TextAt{13}{337.715}{366.728}{i}%
\TextAt{3}{341.708}{364.935}{\Character{3}}%
\TextAt{9}{349.762}{366.728}{\Character{58}}%
\TextAt{13}{356.334}{366.728}{K}%
\TextAt{17}{366.293}{364.935}{0}%
\TextAt{9}{371.025}{366.728}{\Character{40}}%
\Vector{q 0 G 0 g 0.478 w 303.578 -397.429 m 313.428 -397.429 l S Q}%
\TextAt{16}{375.578}{366.728}{\Character{79}}%
\TextAt{9}{385.428}{366.728}{\Character{41}}%
\NativeArrowAt{393.301}{366.728}{11.955}{19.26}{\longrightarrow}%
\TextAt{13}{415.883}{366.728}{K}%
\TextAt{17}{425.842}{364.935}{0}%
\TextAt{9}{430.574}{366.728}{\Character{40}}%
\Vector{q 0 G 0 g 0.478 w 363.126 -399.915 m 368.971 -399.915 l S Q}%
\TextAt{14}{435.126}{366.728}{o}%
\TextAt{9}{440.971}{366.728}{\Character{41}}%
\TextAt{13}{445.523}{366.728}{\Character{59}}%
\TextAt{8}{79.226}{338.899}{defined}%
\TextAt{8}{119.539}{338.899}{as}%
\TextAt{8}{133.909}{338.899}{follows.}%
\TextAt{8}{97.159}{318.864}{If}%
\TextAt{13}{108.851}{318.864}{M}%
\TextAt{8}{125.322}{318.864}{is}%
\TextAt{8}{137.086}{318.864}{a}%
\TextAt{8}{146.841}{318.864}{projective}%
\TextAt{8}{200.819}{318.864}{metrized}%
\TextAt{14}{248.939}{318.864}{o}%
\TextAt{8}{254.783}{318.864}{-module,}%
\TextAt{8}{303.225}{318.864}{then}%
\TextAt{13}{329.885}{318.864}{M}%
\TextAt{16}{345.115}{318.864}{\Character{10}}%
\TextAt{27}{354.414}{317.071}{o}%
\TextAt{16}{361.465}{318.864}{\Character{79}}%
\TextAt{8}{375.212}{318.864}{is}%
\TextAt{8}{386.976}{318.864}{a}%
\TextAt{8}{396.731}{318.864}{projective}%
\TextAt{16}{450.709}{318.864}{\Character{79}}%
\TextAt{8}{460.559}{318.864}{-module.}%
\TextAt{8}{510.3}{318.864}{As}%
\TextAt{9}{195.556}{291.035}{\Character{40}}%
\TextAt{13}{200.108}{291.035}{M}%
\TextAt{16}{215.338}{291.035}{\Character{10}}%
\TextAt{27}{224.637}{289.242}{o}%
\TextAt{16}{231.688}{291.035}{\Character{79}}%
\TextAt{9}{241.537}{291.035}{\Character{41}}%
\TextAt{25}{246.09}{289.242}{C}%
\TextAt{9}{255.665}{291.035}{\Character{61}}%
\TextAt{13}{268.09}{291.035}{M}%
\TextAt{16}{283.321}{291.035}{\Character{10}}%
\TextAt{27}{292.619}{289.242}{o}%
\TextAt{16}{299.67}{291.035}{\Character{79}}%
\TextAt{16}{312.175}{291.035}{\Character{10}}%
\TextAt{25}{321.475}{289.242}{Z}%
\TextAt{18}{329.943}{291.035}{C}%
\TextAt{9}{341.898}{291.035}{\Character{61}}%
\TextAt{13}{354.324}{291.035}{M}%
\TextAt{25}{365.649}{289.242}{C}%
\TextAt{16}{374.56}{291.035}{\Character{10}}%
\TextAt{21}{383.858}{289.242}{K}%
\TextAt{29}{391}{287.869}{C}%
\TextAt{13}{398.97}{291.035}{L}%
\TextAt{25}{406.934}{289.242}{C}%
\TextAt{13}{413.188}{291.035}{\Character{59}}%
\TextAt{8}{79.226}{263.206}{the}%
\TextAt{8}{99.013}{263.206}{hermitian}%
\TextAt{8}{151.974}{263.206}{metric}%
\TextAt{8}{188.02}{263.206}{on}%
\TextAt{8}{203.909}{263.206}{the}%
\TextAt{13}{223.695}{263.206}{K}%
\TextAt{25}{233.654}{261.412}{C}%
\TextAt{8}{239.909}{263.206}{-module}%
\TextAt{13}{284.729}{263.206}{M}%
\TextAt{25}{296.054}{261.412}{C}%
\TextAt{8}{305.836}{263.206}{extends}%
\TextAt{8}{348.122}{263.206}{canonically}%
\TextAt{8}{408.567}{263.206}{to}%
\TextAt{8}{422.508}{263.206}{an}%
\TextAt{13}{438.397}{263.206}{F}%
\TextAt{8}{447.6}{263.206}{-invariant}%
\TextAt{8}{500.251}{263.206}{metric}%
\TextAt{8}{79.226}{246.16}{of}%
\TextAt{8}{92.511}{246.16}{the}%
\TextAt{13}{112.623}{246.16}{L}%
\TextAt{25}{120.588}{244.366}{C}%
\TextAt{8}{126.842}{246.16}{-module}%
\TextAt{9}{171.988}{246.16}{\Character{40}}%
\TextAt{13}{176.54}{246.16}{M}%
\TextAt{16}{191.679}{246.16}{\Character{10}}%
\TextAt{27}{200.978}{244.366}{o}%
\TextAt{16}{207.937}{246.16}{\Character{79}}%
\TextAt{9}{217.787}{246.16}{\Character{41}}%
\TextAt{25}{222.339}{244.366}{C}%
\TextAt{8}{228.594}{246.16}{.}%
\TextAt{8}{237.027}{246.16}{Therefore}%
\TextAt{13}{289.98}{246.16}{M}%
\TextAt{16}{305.119}{246.16}{\Character{10}}%
\TextAt{27}{314.417}{244.366}{o}%
\TextAt{16}{321.377}{246.16}{\Character{79}}%
\TextAt{8}{335.079}{246.16}{is}%
\TextAt{8}{346.799}{246.16}{automatically}%
\TextAt{8}{420.589}{246.16}{a}%
\TextAt{8}{430.3}{246.16}{metrized}%
\TextAt{16}{478.375}{246.16}{\Character{79}}%
\TextAt{8}{488.225}{246.16}{-module,}%
\TextAt{8}{79.226}{229.114}{which}%
\TextAt{8}{112.713}{229.114}{we}%
\TextAt{8}{129.94}{229.114}{denote}%
\TextAt{8}{167.659}{229.114}{by}%
\TextAt{13}{183.918}{229.114}{i}%
\TextAt{3}{187.911}{233.452}{\Character{3}}%
\TextAt{13}{192.643}{229.114}{M}%
\TextAt{8}{205.217}{229.114}{.}%
\TextAt{8}{213.665}{229.114}{If}%
\TextAt{9}{225.535}{201.284}{0}%
\NativeArrowAt{234.709}{201.284}{11.955}{19.26}{\longrightarrow}%
\TextAt{13}{257.291}{201.284}{M}%
\TextAt{3}{269.865}{206.221}{\Character{48}}%
\NativeArrowAt{275.981}{201.284}{11.955}{19.26}{\longrightarrow}%
\TextAt{13}{298.563}{201.284}{M}%
\NativeArrowAt{314.457}{201.284}{11.955}{19.26}{\longrightarrow}%
\TextAt{13}{337.039}{201.284}{M}%
\TextAt{3}{349.613}{206.221}{\Character{48}\Character{48}}%
\NativeArrowAt{358.025}{201.284}{11.955}{19.26}{\longrightarrow}%
\TextAt{9}{380.607}{201.284}{0}%
\TextAt{8}{79.226}{173.455}{is}%
\TextAt{8}{90.99}{173.455}{a}%
\TextAt{8}{100.745}{173.455}{short}%
\TextAt{8}{130.729}{173.455}{exact}%
\TextAt{8}{161.621}{173.455}{sequence}%
\TextAt{8}{210.123}{173.455}{of}%
\TextAt{8}{223.453}{173.455}{projective}%
\TextAt{8}{277.431}{173.455}{metrized}%
\TextAt{14}{325.55}{173.455}{o}%
\TextAt{8}{331.395}{173.455}{-modules,}%
\TextAt{8}{384.452}{173.455}{then}%
\TextAt{9}{182.251}{145.626}{0}%
\NativeArrowAt{191.425}{145.626}{11.955}{19.26}{\longrightarrow}%
\TextAt{13}{214.007}{145.626}{M}%
\TextAt{3}{226.581}{150.562}{\Character{48}}%
\TextAt{16}{232.033}{145.626}{\Character{10}}%
\TextAt{27}{241.331}{143.833}{o}%
\TextAt{16}{248.382}{145.626}{\Character{79}}%
\NativeArrowAt{261.545}{145.626}{11.955}{19.26}{\longrightarrow}%
\TextAt{13}{284.135}{145.626}{M}%
\TextAt{16}{299.365}{145.626}{\Character{10}}%
\TextAt{27}{308.663}{143.833}{o}%
\TextAt{16}{315.714}{145.626}{\Character{79}}%
\NativeArrowAt{328.877}{145.626}{11.955}{19.26}{\longrightarrow}%
\TextAt{13}{351.467}{145.626}{M}%
\TextAt{3}{364.041}{150.562}{\Character{48}\Character{48}}%
\TextAt{16}{371.789}{145.626}{\Character{10}}%
\TextAt{27}{381.088}{143.833}{o}%
\TextAt{16}{388.139}{145.626}{\Character{79}}%
\NativeArrowAt{401.302}{145.626}{11.955}{19.26}{\longrightarrow}%
\TextAt{9}{423.891}{145.626}{0}%
\TextAt{8}{79.226}{117.797}{is}%
\TextAt{8}{90.689}{117.797}{a}%
\TextAt{8}{100.143}{117.797}{short}%
\TextAt{8}{129.826}{117.797}{exact}%
\TextAt{8}{160.417}{117.797}{sequence}%
\TextAt{8}{208.618}{117.797}{of}%
\TextAt{8}{221.646}{117.797}{metrized}%
\TextAt{16}{269.465}{117.797}{\Character{79}}%
\TextAt{8}{279.315}{117.797}{-modules,}%
\TextAt{8}{332.131}{117.797}{because}%
\TextAt{16}{375.131}{117.797}{\Character{79}}%
\TextAt{8}{388.577}{117.797}{is}%
\TextAt{8}{400.04}{117.797}{a}%
\TextAt{8}{409.494}{117.797}{projective}%
\TextAt{14}{463.171}{117.797}{o}%
\TextAt{8}{469.015}{117.797}{-module}%
\TextAt{8}{513.904}{117.797}{and}%
\TextAt{8}{79.226}{100.751}{the}%
\TextAt{8}{99.383}{100.751}{metrics}%
\TextAt{8}{140.413}{100.751}{in}%
\TextAt{8}{154.066}{100.751}{the}%
\TextAt{8}{174.222}{100.751}{sequence}%
\TextAt{9}{221.927}{72.921}{0}%
\NativeArrowAt{231.101}{72.921}{11.955}{19.26}{\longrightarrow}%
\TextAt{13}{253.683}{72.921}{M}%
\TextAt{3}{266.257}{77.858}{\Character{48}}%
\TextAt{25}{265.008}{69.966}{C}%
\NativeArrowAt{274.583}{72.921}{11.955}{19.26}{\longrightarrow}%
\TextAt{13}{297.165}{72.921}{M}%
\TextAt{25}{308.49}{71.128}{C}%
\NativeArrowAt{318.065}{72.921}{11.955}{19.26}{\longrightarrow}%
\TextAt{13}{340.647}{72.921}{M}%
\TextAt{3}{353.221}{77.858}{\Character{48}\Character{48}}%
\TextAt{25}{351.972}{69.966}{C}%
\NativeArrowAt{361.633}{72.921}{11.955}{19.26}{\longrightarrow}%
\TextAt{9}{384.215}{72.921}{0}%
\TextAt{8}{297.211}{43.033}{226}%
\Destination{78.23}{809.815}{nav.page.226}%
\Destination{79.23}{503.445}{nav.section.3.7}%
\Destination{79.36}{776.155}{nav.print.246}%
\Destination{79.36}{636.655}{nav.number.3.6.6}%
\Destination{79.23}{503.445}{outline.38}%
}

% Page 239 (body page 227)
\Page{%
\TextAt{2}{79.226}{796.943}{III.}%
\TextAt{2}{98.317}{796.943}{Riemann–Roch}%
\TextAt{2}{174.375}{796.943}{Theory}%
\Vector{q 0 G 0 g 0.299 w 7.226 18.729 m 460.769 18.729 l S Q}%
\TextAt{8}{79.226}{761.124}{simply}%
\TextAt{8}{116.682}{761.124}{extend}%
\TextAt{13}{154.723}{761.124}{L}%
\TextAt{25}{162.688}{759.33}{C}%
\TextAt{8}{168.942}{761.124}{-sesquilinearly}%
\TextAt{8}{245.156}{761.124}{to}%
\TextAt{8}{259.466}{761.124}{metrics}%
\TextAt{8}{300.497}{761.124}{in}%
\TextAt{8}{314.149}{761.124}{the}%
\TextAt{8}{334.306}{761.124}{sequence}%
\TextAt{8}{382.808}{761.124}{of}%
\TextAt{13}{396.138}{761.124}{L}%
\TextAt{25}{404.102}{759.33}{C}%
\TextAt{8}{410.357}{761.124}{-modules}%
\TextAt{9}{158.373}{738.51}{0}%
\NativeArrowAt{167.547}{738.51}{11.955}{19.26}{\longrightarrow}%
\TextAt{13}{190.129}{738.51}{M}%
\TextAt{3}{202.703}{743.446}{\Character{48}}%
\TextAt{25}{201.454}{735.554}{C}%
\TextAt{16}{210.365}{738.51}{\Character{10}}%
\TextAt{21}{219.663}{736.717}{K}%
\TextAt{29}{226.805}{735.344}{C}%
\TextAt{13}{234.775}{738.51}{L}%
\TextAt{25}{242.739}{736.717}{C}%
\NativeArrowAt{252.314}{738.51}{11.955}{19.26}{\longrightarrow}%
\TextAt{13}{274.896}{738.51}{M}%
\TextAt{25}{286.221}{736.717}{C}%
\TextAt{16}{295.132}{738.51}{\Character{10}}%
\TextAt{21}{304.431}{736.717}{K}%
\TextAt{29}{311.572}{735.344}{C}%
\TextAt{13}{319.542}{738.51}{L}%
\TextAt{25}{327.507}{736.717}{C}%
\NativeArrowAt{337.082}{738.51}{11.955}{19.26}{\longrightarrow}%
\TextAt{13}{359.664}{738.51}{M}%
\TextAt{3}{372.237}{743.446}{\Character{48}\Character{48}}%
\TextAt{25}{370.989}{735.554}{C}%
\TextAt{16}{379.986}{738.51}{\Character{10}}%
\TextAt{21}{389.285}{736.717}{K}%
\TextAt{29}{396.426}{735.344}{C}%
\TextAt{13}{404.396}{738.51}{L}%
\TextAt{25}{412.361}{736.717}{C}%
\NativeArrowAt{421.936}{738.51}{11.955}{19.26}{\longrightarrow}%
\TextAt{9}{444.518}{738.51}{0}%
\TextAt{13}{450.371}{738.51}{\Character{58}}%
\TextAt{8}{79.226}{715.896}{This}%
\TextAt{8}{105.376}{715.896}{is}%
\TextAt{8}{116.571}{715.896}{why}%
\TextAt{8}{140.713}{715.896}{mapping,}%
\TextAt{8}{191.64}{715.896}{in}%
\TextAt{8}{204.724}{715.896}{the}%
\TextAt{8}{224.31}{715.896}{usual}%
\TextAt{8}{254.37}{715.896}{way}%
\TextAt{8}{277.543}{715.896}{(i.e.,}%
\TextAt{8}{303.748}{715.896}{via}%
\TextAt{8}{322.366}{715.896}{the}%
\TextAt{8}{341.953}{715.896}{representation}%
\TextAt{13}{417.873}{715.896}{K}%
\TextAt{17}{427.832}{714.103}{0}%
\TextAt{9}{432.564}{715.896}{\Character{40}}%
\Vector{q 0 G 0 g 0.478 w 365.117 -50.748 m 370.961 -50.748 l S Q}%
\TextAt{14}{437.117}{715.896}{o}%
\TextAt{9}{442.961}{715.896}{\Character{41}}%
\TextAt{9}{450.816}{715.896}{\Character{61}}%
\TextAt{13}{463.259}{715.896}{F}%
\TextAt{17}{470.837}{714.103}{0}%
\TextAt{9}{475.57}{715.896}{\Character{40}}%
\Vector{q 0 G 0 g 0.478 w 408.122 -50.748 m 413.966 -50.748 l S Q}%
\TextAt{14}{480.122}{715.896}{o}%
\TextAt{9}{485.966}{715.896}{\Character{41}\Character{47}}%
\TextAt{13}{496.371}{715.896}{R}%
\TextAt{17}{505.282}{714.103}{0}%
\TextAt{9}{510.014}{715.896}{\Character{40}}%
\Vector{q 0 G 0 g 0.478 w 442.566 -50.748 m 448.41 -50.748 l S Q}%
\TextAt{14}{514.566}{715.896}{o}%
\TextAt{9}{520.41}{715.896}{\Character{41}}%
\TextAt{8}{524.963}{715.896}{),}%
\TextAt{13}{241.019}{693.282}{M}%
\NativeArrowAt{256.913}{693.282}{11.955}{19.26}{\longmapsto}%
\TextAt{9}{279.495}{693.282}{\Character{91}}%
\TextAt{13}{282.747}{693.282}{i}%
\TextAt{3}{286.74}{698.218}{\Character{3}}%
\TextAt{13}{291.473}{693.282}{M}%
\TextAt{9}{304.046}{693.282}{\Character{93}}%
\TextAt{9}{310.597}{693.282}{\Character{61}}%
\TextAt{9}{323.019}{693.282}{\Character{91}}%
\TextAt{13}{326.296}{693.282}{M}%
\TextAt{16}{341.526}{693.282}{\Character{10}}%
\TextAt{27}{350.825}{691.489}{o}%
\TextAt{16}{357.876}{693.282}{\Character{79}}%
\TextAt{9}{367.725}{693.282}{\Character{93}}%
\TextAt{8}{79.226}{670.668}{gives}%
\TextAt{8}{107.907}{670.668}{a}%
\TextAt{8}{117.662}{670.668}{well-defined}%
\TextAt{8}{181.706}{670.668}{homomorphism}%
\TextAt{13}{250.468}{648.054}{i}%
\TextAt{3}{254.461}{652.99}{\Character{3}}%
\TextAt{9}{262.515}{648.054}{\Character{58}}%
\TextAt{13}{269.087}{648.054}{K}%
\TextAt{17}{279.046}{646.261}{0}%
\TextAt{9}{283.778}{648.054}{\Character{40}}%
\Vector{q 0 G 0 g 0.478 w 216.331 -118.589 m 222.175 -118.589 l S Q}%
\TextAt{14}{288.331}{648.054}{o}%
\TextAt{9}{294.175}{648.054}{\Character{41}}%
\NativeArrowAt{302.048}{648.054}{11.955}{19.26}{\longrightarrow}%
\TextAt{13}{324.63}{648.054}{K}%
\TextAt{17}{334.589}{646.261}{0}%
\TextAt{9}{339.322}{648.054}{\Character{40}}%
\Vector{q 0 G 0 g 0.478 w 271.874 -116.103 m 281.724 -116.103 l S Q}%
\TextAt{16}{343.874}{648.054}{\Character{79}}%
\TextAt{9}{353.724}{648.054}{\Character{41}}%
\TextAt{13}{358.276}{648.054}{\Character{58}}%
\TextAt{8}{79.226}{625.44}{The}%
\TextAt{8}{103.28}{625.44}{reader}%
\TextAt{8}{139.05}{625.44}{may}%
\TextAt{8}{164.419}{625.44}{verify}%
\TextAt{8}{196.937}{625.44}{for}%
\TextAt{8}{214.822}{625.44}{himself}%
\TextAt{8}{254.871}{625.44}{that}%
\TextAt{8}{280.24}{625.44}{this}%
\TextAt{8}{303.063}{625.44}{is}%
\TextAt{8}{314.827}{625.44}{in}%
\TextAt{8}{328.479}{625.44}{fact}%
\TextAt{8}{351.565}{625.44}{a}%
\TextAt{8}{361.32}{625.44}{ring}%
\TextAt{8}{385.386}{625.44}{homomorphism.}%
\TextAt{8}{97.159}{605.406}{On}%
\TextAt{8}{117.964}{605.406}{the}%
\TextAt{8}{139.415}{605.406}{other}%
\TextAt{8}{171.278}{605.406}{hand,}%
\TextAt{8}{205.414}{605.406}{if}%
\TextAt{13}{217.432}{605.406}{M}%
\TextAt{8}{235.197}{605.406}{is}%
\TextAt{8}{248.255}{605.406}{a}%
\TextAt{8}{259.305}{605.406}{projective}%
\TextAt{8}{314.576}{605.406}{metrized}%
\TextAt{16}{363.99}{605.406}{\Character{79}}%
\TextAt{8}{373.84}{605.406}{-module,}%
\TextAt{8}{423.9}{605.406}{then}%
\TextAt{13}{451.854}{605.406}{M}%
\TextAt{8}{469.619}{605.406}{is}%
\TextAt{8}{482.677}{605.406}{automati-}%
\TextAt{8}{79.226}{588.36}{cally}%
\TextAt{8}{107.756}{588.36}{also}%
\TextAt{8}{132.126}{588.36}{a}%
\TextAt{8}{142.771}{588.36}{projective}%
\TextAt{14}{197.639}{588.36}{o}%
\TextAt{8}{203.483}{588.36}{-module.}%
\TextAt{8}{255.894}{588.36}{For}%
\TextAt{8}{277.753}{588.36}{the}%
\TextAt{8}{298.799}{588.36}{complexification}%
\TextAt{13}{386.519}{588.36}{M}%
\TextAt{25}{397.844}{586.566}{C}%
\TextAt{9}{408.935}{588.36}{\Character{61}}%
\TextAt{13}{422.877}{588.36}{M}%
\TextAt{16}{438.714}{588.36}{\Character{10}}%
\TextAt{25}{448.013}{586.566}{Z}%
\TextAt{18}{457.087}{588.36}{C}%
\TextAt{8}{470.509}{588.36}{we}%
\TextAt{8}{488.626}{588.36}{have}%
\TextAt{8}{516.51}{588.36}{the}%
\TextAt{8}{79.226}{571.314}{decomposition}%
\TextAt{13}{188.541}{554.268}{M}%
\TextAt{25}{199.866}{552.474}{C}%
\TextAt{9}{209.441}{554.268}{\Character{61}}%
\TextAt{20}{227.562}{565.625}{\Character{77}}%
\TextAt{21}{221.866}{539.722}{\Character{28}}%
\NativeMathAt{226.545}{539.722}{7.97}{\in}%
\TextAt{21}{232.19}{539.722}{Y}%
\TextAt{17}{238.98}{539.722}{\Character{40}}%
\TextAt{25}{242.273}{539.722}{C}%
\TextAt{17}{248.03}{539.722}{\Character{41}}%
\TextAt{13}{253.315}{554.268}{M}%
\TextAt{21}{264.64}{552.474}{\Character{28}}%
\TextAt{9}{273.138}{554.268}{\Character{61}}%
\TextAt{20}{291.88}{565.625}{\Character{77}}%
\TextAt{21}{285.563}{539.722}{\Character{27}}%
\NativeMathAt{290.688}{539.722}{7.97}{\in}%
\TextAt{21}{296.334}{539.722}{X}%
\TextAt{17}{303.921}{539.722}{\Character{40}}%
\TextAt{25}{307.214}{539.722}{C}%
\TextAt{17}{312.97}{539.722}{\Character{41}}%
\TextAt{20}{318.256}{565.625}{\Character{77}}%
\TextAt{21}{321.211}{539.722}{\Character{28}}%
\TextAt{3}{325.889}{539.722}{\Character{106}}%
\TextAt{21}{328.242}{539.722}{\Character{27}}%
\TextAt{13}{338.314}{554.268}{M}%
\TextAt{21}{349.639}{552.474}{\Character{28}}%
\TextAt{9}{358.136}{554.268}{\Character{61}}%
\TextAt{20}{376.879}{565.625}{\Character{77}}%
\TextAt{21}{370.562}{539.722}{\Character{27}}%
\NativeMathAt{375.687}{539.722}{7.97}{\in}%
\TextAt{21}{381.333}{539.722}{X}%
\TextAt{17}{388.92}{539.722}{\Character{40}}%
\TextAt{25}{392.213}{539.722}{C}%
\TextAt{17}{397.969}{539.722}{\Character{41}}%
\TextAt{13}{403.255}{554.268}{M}%
\TextAt{21}{414.58}{552.474}{\Character{27}}%
\TextAt{13}{420.203}{554.268}{\Character{59}}%
\TextAt{8}{79.226}{523.503}{where}%
\TextAt{13}{113.035}{523.503}{M}%
\TextAt{21}{124.36}{521.709}{\Character{28}}%
\TextAt{9}{132.857}{523.503}{\Character{61}}%
\TextAt{13}{145.283}{523.503}{M}%
\TextAt{16}{160.513}{523.503}{\Character{10}}%
\TextAt{3}{169.812}{521.709}{\Character{79}}%
\TextAt{21}{176.755}{521.709}{\Character{59}\Character{28}}%
\TextAt{18}{186.94}{523.503}{C}%
\TextAt{8}{199.472}{523.503}{and}%
\TextAt{13}{238.026}{500.889}{M}%
\TextAt{21}{249.351}{499.095}{\Character{27}}%
\TextAt{9}{258.295}{500.889}{\Character{61}}%
\TextAt{13}{270.721}{500.889}{M}%
\TextAt{16}{285.951}{500.889}{\Character{10}}%
\TextAt{27}{295.249}{499.095}{o}%
\TextAt{21}{299.146}{499.095}{\Character{59}\Character{27}}%
\TextAt{18}{309.778}{500.889}{C}%
\TextAt{9}{321.733}{500.889}{\Character{61}}%
\TextAt{20}{334.159}{512.246}{\Character{77}}%
\TextAt{21}{337.113}{486.343}{\Character{28}}%
\TextAt{3}{341.792}{486.343}{\Character{106}}%
\TextAt{21}{344.144}{486.343}{\Character{27}}%
\TextAt{13}{354.217}{500.889}{M}%
\TextAt{21}{365.541}{499.095}{\Character{28}}%
\TextAt{13}{370.718}{500.889}{\Character{58}}%
\TextAt{8}{79.226}{467.625}{The}%
\TextAt{18}{102.793}{467.625}{C}%
\TextAt{8}{111.427}{467.625}{-vector}%
\TextAt{8}{149.961}{467.625}{spaces}%
\TextAt{13}{185.364}{467.625}{M}%
\TextAt{21}{196.688}{465.831}{\Character{28}}%
\TextAt{8}{205.275}{467.625}{carry}%
\TextAt{8}{235.034}{467.625}{hermitian}%
\TextAt{8}{287.879}{467.625}{metrics}%
\TextAt{16}{328.422}{467.625}{\Character{104}}%
\TextAt{13}{336.482}{467.625}{\Character{59}}%
\TextAt{16}{345.136}{467.625}{\Character{105}}%
\TextAt{21}{349.785}{465.831}{M}%
\TextAt{22}{357.918}{464.835}{\Character{28}}%
\TextAt{8}{363.085}{467.625}{,}%
\TextAt{8}{369.845}{467.625}{and}%
\TextAt{8}{392.12}{467.625}{we}%
\TextAt{8}{408.86}{467.625}{define}%
\TextAt{8}{442.182}{467.625}{the}%
\TextAt{8}{461.852}{467.625}{metric}%
\TextAt{16}{497.78}{467.625}{\Character{104}}%
\TextAt{13}{505.839}{467.625}{\Character{59}}%
\TextAt{16}{514.494}{467.625}{\Character{105}}%
\TextAt{21}{519.143}{465.831}{M}%
\TextAt{22}{527.276}{464.835}{\Character{27}}%
\TextAt{8}{79.226}{450.579}{on}%
\TextAt{8}{95.485}{450.579}{the}%
\TextAt{18}{115.642}{450.579}{C}%
\TextAt{8}{124.276}{450.579}{-vector}%
\TextAt{8}{163.298}{450.579}{space}%
\TextAt{13}{194.572}{450.579}{M}%
\TextAt{21}{205.897}{448.785}{\Character{27}}%
\TextAt{8}{215.418}{450.579}{to}%
\TextAt{8}{229.728}{450.579}{be}%
\TextAt{8}{245.653}{450.579}{the}%
\TextAt{8}{265.809}{450.579}{orthogonal}%
\TextAt{8}{324.365}{450.579}{sum}%
\TextAt{16}{242.108}{427.965}{\Character{104}}%
\TextAt{13}{246.757}{427.965}{x\Character{59}}%
\TextAt{13}{258.629}{427.965}{y}%
\TextAt{16}{264.79}{427.965}{\Character{105}}%
\TextAt{21}{269.439}{426.171}{M}%
\TextAt{22}{277.572}{425.175}{\Character{27}}%
\TextAt{9}{286.387}{427.965}{\Character{61}}%
\TextAt{20}{298.812}{439.322}{\Character{88}}%
\TextAt{21}{301.369}{413.419}{\Character{28}}%
\TextAt{3}{306.047}{413.419}{\Character{106}}%
\TextAt{21}{308.399}{413.419}{\Character{27}}%
\TextAt{16}{316.081}{427.965}{\Character{104}}%
\TextAt{13}{320.73}{427.965}{x}%
\TextAt{21}{327.382}{426.171}{\Character{28}}%
\TextAt{13}{332.559}{427.965}{\Character{59}}%
\TextAt{13}{337.783}{427.965}{y}%
\TextAt{21}{343.511}{426.171}{\Character{28}}%
\TextAt{16}{348.687}{427.965}{\Character{105}}%
\TextAt{21}{353.336}{426.171}{M}%
\TextAt{22}{361.469}{425.175}{\Character{28}}%
\TextAt{13}{366.636}{427.965}{\Character{58}}%
\TextAt{8}{79.226}{395.37}{This}%
\TextAt{8}{105.738}{395.37}{gives}%
\TextAt{8}{134.211}{395.37}{a}%
\TextAt{8}{143.759}{395.37}{hermitian}%
\TextAt{8}{196.883}{395.37}{metric}%
\TextAt{8}{233.09}{395.37}{on}%
\TextAt{8}{249.142}{395.37}{the}%
\TextAt{13}{269.09}{395.37}{K}%
\TextAt{25}{279.049}{393.577}{C}%
\TextAt{8}{285.304}{395.37}{-module}%
\TextAt{13}{330.287}{395.37}{M}%
\TextAt{25}{341.611}{393.577}{C}%
\TextAt{8}{347.866}{395.37}{,}%
\TextAt{8}{354.849}{395.37}{whose}%
\TextAt{13}{389.167}{395.37}{F}%
\TextAt{8}{398.371}{395.37}{-invariance}%
\TextAt{8}{457.353}{395.37}{is}%
\TextAt{8}{468.909}{395.37}{clearly}%
\TextAt{8}{506.097}{395.37}{guar-}%
\TextAt{8}{79.226}{378.324}{anteed}%
\TextAt{8}{115.618}{378.324}{by}%
\TextAt{8}{130.873}{378.324}{the}%
\TextAt{13}{150.025}{378.324}{F}%
\TextAt{8}{159.229}{378.324}{-invariance}%
\TextAt{8}{217.415}{378.324}{of}%
\TextAt{8}{229.741}{378.324}{the}%
\TextAt{8}{248.893}{378.324}{original}%
\TextAt{8}{290.175}{378.324}{metric}%
\TextAt{16}{325.586}{378.324}{\Character{104}}%
\TextAt{13}{333.129}{378.324}{\Character{59}}%
\TextAt{16}{341.266}{378.324}{\Character{105}}%
\TextAt{21}{345.916}{376.531}{M}%
\TextAt{8}{355.396}{378.324}{.}%
\TextAt{8}{363.509}{378.324}{We}%
\TextAt{8}{382.65}{378.324}{denote}%
\TextAt{8}{419.364}{378.324}{the}%
\TextAt{8}{438.517}{378.324}{metrized}%
\TextAt{14}{485.632}{378.324}{o}%
\TextAt{8}{491.476}{378.324}{-module}%
\TextAt{13}{79.226}{361.278}{M}%
\TextAt{8}{95.697}{361.278}{thus}%
\TextAt{8}{121.449}{361.278}{constructed}%
\TextAt{8}{184.596}{361.278}{by}%
\TextAt{13}{200.855}{361.278}{i}%
\TextAt{3}{204.848}{359.485}{\Character{3}}%
\TextAt{13}{209.581}{361.278}{M}%
\TextAt{8}{222.154}{361.278}{.}%
\TextAt{8}{97.159}{341.243}{If}%
\TextAt{9}{110.046}{341.243}{0}%
\NativeMathAt{121.257}{341.243}{11.955}{\rightarrow}%
\TextAt{13}{138.57}{341.243}{M}%
\TextAt{3}{151.144}{345.582}{\Character{48}}%
\NativeMathAt{159.296}{341.243}{11.955}{\rightarrow}%
\TextAt{13}{176.609}{341.243}{M}%
\NativeMathAt{194.541}{341.243}{11.955}{\rightarrow}%
\TextAt{13}{211.854}{341.243}{M}%
\TextAt{3}{224.427}{345.582}{\Character{48}\Character{48}}%
\NativeMathAt{234.877}{341.243}{11.955}{\rightarrow}%
\TextAt{9}{252.19}{341.243}{0}%
\TextAt{8}{263.136}{341.243}{is}%
\TextAt{8}{276.095}{341.243}{a}%
\TextAt{8}{287.046}{341.243}{short}%
\TextAt{8}{318.224}{341.243}{exact}%
\TextAt{8}{350.312}{341.243}{sequence}%
\TextAt{8}{400.009}{341.243}{of}%
\TextAt{8}{414.535}{341.243}{projective}%
\TextAt{8}{469.707}{341.243}{metrized}%
\TextAt{16}{519.022}{341.243}{\Character{79}}%
\TextAt{8}{528.872}{341.243}{-}%
\TextAt{8}{79.226}{324.197}{modules,}%
\TextAt{8}{128.386}{324.197}{then}%
\TextAt{9}{212.447}{307.151}{0}%
\NativeArrowAt{221.62}{307.151}{11.955}{19.26}{\longrightarrow}%
\TextAt{13}{244.202}{307.151}{i}%
\TextAt{3}{248.196}{305.358}{\Character{3}}%
\TextAt{13}{252.928}{307.151}{M}%
\TextAt{3}{265.502}{312.088}{\Character{48}}%
\NativeArrowAt{271.618}{307.151}{11.955}{19.26}{\longrightarrow}%
\TextAt{13}{294.2}{307.151}{i}%
\TextAt{3}{298.193}{305.358}{\Character{3}}%
\TextAt{13}{302.925}{307.151}{M}%
\NativeArrowAt{318.82}{307.151}{11.955}{19.26}{\longrightarrow}%
\TextAt{13}{341.402}{307.151}{i}%
\TextAt{3}{345.395}{305.358}{\Character{3}}%
\TextAt{13}{350.128}{307.151}{M}%
\TextAt{3}{362.701}{312.088}{\Character{48}\Character{48}}%
\NativeArrowAt{371.114}{307.151}{11.955}{19.26}{\longrightarrow}%
\TextAt{9}{393.696}{307.151}{0}%
\TextAt{8}{79.226}{286.367}{is}%
\TextAt{8}{91.107}{286.367}{clearly}%
\TextAt{8}{128.619}{286.367}{an}%
\TextAt{8}{144.995}{286.367}{exact}%
\TextAt{8}{176.003}{286.367}{sequence}%
\TextAt{8}{224.622}{286.367}{of}%
\TextAt{8}{238.069}{286.367}{projective}%
\TextAt{8}{292.163}{286.367}{metrized}%
\TextAt{14}{340.399}{286.367}{o}%
\TextAt{8}{346.243}{286.367}{-modules.}%
\TextAt{8}{400.949}{286.367}{As}%
\TextAt{8}{418.353}{286.367}{before,}%
\TextAt{8}{456.863}{286.367}{this}%
\TextAt{8}{479.802}{286.367}{is}%
\TextAt{8}{491.683}{286.367}{why}%
\TextAt{8}{516.51}{286.367}{the}%
\TextAt{8}{79.226}{269.321}{correspondence}%
\TextAt{13}{272.858}{252.275}{M}%
\NativeArrowAt{288.753}{252.275}{11.955}{19.26}{\longmapsto}%
\TextAt{9}{311.335}{252.275}{\Character{91}}%
\TextAt{13}{314.586}{252.275}{i}%
\TextAt{3}{318.58}{250.482}{\Character{3}}%
\TextAt{13}{323.312}{252.275}{M}%
\TextAt{9}{335.886}{252.275}{\Character{93}}%
\TextAt{8}{79.226}{231.491}{gives}%
\TextAt{8}{107.907}{231.491}{us}%
\TextAt{8}{122.922}{231.491}{a}%
\TextAt{8}{132.678}{231.491}{well-defined}%
\TextAt{8}{196.722}{231.491}{(additive)}%
\TextAt{8}{250.711}{231.491}{homomorphism}%
\TextAt{13}{250.468}{208.877}{i}%
\TextAt{3}{254.461}{207.084}{\Character{3}}%
\TextAt{9}{262.515}{208.877}{\Character{58}}%
\TextAt{13}{269.087}{208.877}{K}%
\TextAt{17}{279.046}{207.084}{0}%
\TextAt{9}{283.778}{208.877}{\Character{40}}%
\Vector{q 0 G 0 g 0.478 w 216.331 -555.28 m 226.181 -555.28 l S Q}%
\TextAt{16}{288.331}{208.877}{\Character{79}}%
\TextAt{9}{298.181}{208.877}{\Character{41}}%
\NativeArrowAt{306.054}{208.877}{11.955}{19.26}{\longrightarrow}%
\TextAt{13}{328.636}{208.877}{K}%
\TextAt{17}{338.595}{207.084}{0}%
\TextAt{9}{343.327}{208.877}{\Character{40}}%
\Vector{q 0 G 0 g 0.478 w 275.879 -557.767 m 281.724 -557.767 l S Q}%
\TextAt{14}{347.879}{208.877}{o}%
\TextAt{9}{353.724}{208.877}{\Character{41}}%
\TextAt{13}{358.276}{208.877}{\Character{58}}%
\TextAt{12}{79.226}{179.122}{Proposition}%
\TextAt{12}{151.937}{179.122}{(7.1)}%
\TextAt{12}{184.097}{179.122}{(Projection}%
\TextAt{12}{255.015}{179.122}{Formula).}%
\TextAt{8}{316.046}{179.122}{The}%
\TextAt{8}{340.1}{179.122}{diagram}%
\ErrataAt{347.8}{161.8}{14.2}{$\cdot$}%
\TextAt{13}{202.183}{154.616}{K}%
\TextAt{17}{212.143}{152.823}{0}%
\TextAt{9}{216.875}{154.616}{\Character{40}}%
\NativeDiagram{%
\path[draw=black,line width=0.478bp] (149.427,-609.541) -- (159.277,-609.541);%
}%
\TextAt{16}{221.427}{154.616}{\Character{79}}%
\TextAt{9}{231.277}{154.616}{\Character{41}}%
\TextAt{13}{289.173}{154.616}{K}%
\TextAt{17}{299.132}{152.823}{0}%
\TextAt{9}{303.865}{154.616}{\Character{40}}%
\NativeDiagram{%
\path[draw=black,line width=0.478bp] (236.417,-609.541) -- (246.267,-609.541);%
}%
\TextAt{16}{308.417}{154.616}{\Character{79}}%
\TextAt{9}{318.267}{154.616}{\Character{41}}%
\TextAt{13}{376.163}{154.616}{K}%
\TextAt{17}{386.122}{152.823}{0}%
\TextAt{9}{390.855}{154.616}{\Character{40}}%
\NativeDiagram{%
\path[draw=black,line width=0.478bp] (323.407,-609.541) -- (333.257,-609.541);%
}%
\TextAt{16}{395.407}{154.616}{\Character{79}}%
\TextAt{9}{405.257}{154.616}{\Character{41}}%
\ErrataAt{347.8}{108.8}{14.2}{$\cdot$}%
\TextAt{13}{204.186}{101.623}{K}%
\TextAt{17}{214.146}{99.829}{0}%
\TextAt{9}{218.878}{101.623}{\Character{40}}%
\NativeDiagram{%
\path[draw=black,line width=0.478bp] (151.43,-665.021) -- (157.274,-665.021);%
}%
\TextAt{14}{223.43}{101.623}{o}%
\TextAt{9}{229.274}{101.623}{\Character{41}}%
\TextAt{13}{291.176}{101.623}{K}%
\TextAt{17}{301.135}{99.829}{0}%
\TextAt{9}{305.868}{101.623}{\Character{40}}%
\NativeDiagram{%
\path[draw=black,line width=0.478bp] (238.42,-665.021) -- (244.264,-665.021);%
}%
\TextAt{14}{310.42}{101.623}{o}%
\TextAt{9}{316.264}{101.623}{\Character{41}}%
\TextAt{13}{378.166}{101.623}{K}%
\TextAt{17}{388.125}{99.829}{0}%
\TextAt{9}{392.858}{101.623}{\Character{40}}%
\NativeDiagram{%
\path[draw=black,line width=0.478bp] (325.41,-665.021) -- (331.254,-665.021);%
}%
\TextAt{14}{397.41}{101.623}{o}%
\TextAt{9}{403.254}{101.623}{\Character{41}}%
\TextAt{16}{257.853}{154.617}{\Character{2}}%
\NativeDiagram{%
\path[draw=black,line width=0.478bp] (147.007,-626.951) -- (147.007,-658.353);%
\path[draw=black,line width=0.478bp,line cap=round,line join=round] (149.876,-656.106) .. controls (148.155,-656.56) and (147.342,-657.572) .. (147.007,-658.592) .. controls (146.672,-657.572) and (145.859,-656.56) .. (144.138,-656.106);%
}%
\TextAt{21}{208.991}{128.973}{i}%
\TextAt{23}{211.874}{127.976}{\Character{3}}%
\NativeDiagram{%
\path[draw=black,line width=0.478bp] (256.173,-616.394) -- (298.336,-616.394);%
\path[draw=black,line width=0.478bp,line cap=round,line join=round] (296.088,-613.525) .. controls (296.543,-615.246) and (297.555,-616.059) .. (298.575,-616.394) .. controls (297.555,-616.729) and (296.543,-617.541) .. (296.088,-619.263);%
\path[draw=black,line width=0.478bp] (320.989,-626.951) -- (320.989,-658.353);%
\path[draw=black,line width=0.478bp,line cap=round,line join=round] (323.858,-656.106) .. controls (322.137,-656.56) and (321.324,-657.572) .. (320.989,-658.592) .. controls (320.654,-657.572) and (319.841,-656.56) .. (318.12,-656.106);%
}%
\TextAt{21}{395.804}{128.973}{i}%
\TextAt{23}{398.687}{127.976}{\Character{3}}%
\TextAt{16}{257.853}{101.623}{\Character{2}}%
\NativeDiagram{%
\path[draw=black,line width=0.478bp] (233.998,-658.831) -- (233.998,-627.429);%
\path[draw=black,line width=0.478bp,line cap=round,line join=round] (231.129,-629.676) .. controls (232.85,-629.222) and (233.663,-628.21) .. (233.998,-627.19) .. controls (234.333,-628.21) and (235.146,-629.222) .. (236.867,-629.676);%
}%
\TextAt{21}{295.982}{128.312}{i}%
\TextAt{23}{298.865}{131.124}{\Character{3}}%
\NativeDiagram{%
\path[draw=black,line width=0.478bp] (254.17,-669.388) -- (300.338,-669.388);%
\path[draw=black,line width=0.478bp,line cap=round,line join=round] (298.091,-666.519) .. controls (298.545,-668.24) and (299.558,-669.053) .. (300.577,-669.388) .. controls (299.558,-669.723) and (298.545,-670.536) .. (298.091,-672.257);%
}%
\TextAt{8}{79.226}{79.393}{is}%
\TextAt{8}{90.99}{79.393}{commutative,}%
\TextAt{8}{164.168}{79.393}{where}%
\TextAt{8}{197.977}{79.393}{the}%
\TextAt{8}{218.133}{79.393}{horizontal}%
\TextAt{8}{273.103}{79.393}{arrows}%
\TextAt{8}{310.571}{79.393}{are}%
\TextAt{8}{330.081}{79.393}{multiplication.}%
\TextAt{8}{297.211}{43.033}{227}%
\Destination{78.23}{809.815}{nav.page.227}%
\Destination{79.36}{776.155}{nav.print.247}%
\Destination{79.36}{193.155}{nav.number.3.7.1}%
\Destination{79.226}{190}{errata.204}%
}

% Page 240 (body page 228)
\Page{%
\TextAt{2}{79.226}{796.943}{III.}%
\TextAt{2}{98.317}{796.943}{Riemann–Roch}%
\TextAt{2}{174.375}{796.943}{Theory}%
\Vector{q 0 G 0 g 0.299 w 7.226 18.729 m 460.769 18.729 l S Q}%
\TextAt{19}{79.226}{761.124}{Proof.}%
\TextAt{8}{116}{761.124}{If}%
\TextAt{13}{128.732}{761.124}{M}%
\TextAt{8}{141.305}{761.124}{,}%
\TextAt{8}{149.753}{761.124}{resp.}%
\TextAt{13}{182.193}{761.124}{N}%
\TextAt{8}{192.816}{761.124}{,}%
\TextAt{8}{201.264}{761.124}{is}%
\TextAt{8}{214.067}{761.124}{a}%
\TextAt{8}{224.862}{761.124}{projective}%
\TextAt{8}{279.878}{761.124}{metrized}%
\TextAt{16}{329.037}{761.124}{\Character{79}}%
\TextAt{8}{338.887}{761.124}{-module,}%
\TextAt{8}{388.628}{761.124}{resp.}%
\TextAt{14}{421.068}{761.124}{o}%
\TextAt{8}{426.912}{761.124}{-module,}%
\TextAt{8}{476.654}{761.124}{there}%
\TextAt{8}{507.605}{761.124}{is}%
\TextAt{8}{520.408}{761.124}{an}%
\TextAt{8}{79.226}{744.078}{isometry}%
\TextAt{13}{236.661}{727.032}{i}%
\TextAt{3}{240.655}{725.238}{\Character{3}}%
\TextAt{9}{245.387}{727.032}{\Character{40}}%
\TextAt{13}{249.939}{727.032}{M}%
\TextAt{16}{265.169}{727.032}{\Character{10}}%
\TextAt{3}{274.468}{725.238}{\Character{79}}%
\TextAt{13}{284.566}{727.032}{i}%
\TextAt{3}{288.559}{731.968}{\Character{3}}%
\TextAt{13}{293.291}{727.032}{N}%
\TextAt{9}{303.914}{727.032}{\Character{41}}%
\NativeMathAt{311.787}{727.032}{11.955}{\cong}%
\TextAt{13}{324.406}{727.032}{i}%
\TextAt{3}{328.4}{725.238}{\Character{3}}%
\TextAt{13}{333.132}{727.032}{M}%
\TextAt{16}{348.363}{727.032}{\Character{10}}%
\TextAt{27}{357.661}{725.238}{o}%
\TextAt{13}{364.712}{727.032}{N}%
\TextAt{8}{79.226}{704.388}{of}%
\TextAt{8}{94.061}{704.388}{projective}%
\TextAt{8}{149.543}{704.388}{metrized}%
\TextAt{14}{199.167}{704.388}{o}%
\TextAt{8}{205.011}{704.388}{-modules.}%
\TextAt{8}{263.881}{704.388}{Indeed,}%
\TextAt{8}{307.043}{704.388}{we}%
\TextAt{8}{325.775}{704.388}{have}%
\TextAt{8}{354.274}{704.388}{an}%
\TextAt{8}{372.038}{704.388}{isomorphism}%
\TextAt{8}{441.962}{704.388}{of}%
\TextAt{8}{456.796}{704.388}{the}%
\TextAt{8}{478.457}{704.388}{underlying}%
\TextAt{14}{79.226}{687.342}{o}%
\TextAt{8}{85.07}{687.342}{-modules}%
\TextAt{13}{163.862}{670.296}{M}%
\TextAt{16}{179.093}{670.296}{\Character{10}}%
\TextAt{3}{188.391}{668.503}{\Character{79}}%
\TextAt{9}{198.489}{670.296}{\Character{40}}%
\TextAt{13}{203.041}{670.296}{N}%
\TextAt{16}{216.32}{670.296}{\Character{10}}%
\TextAt{27}{225.619}{668.503}{o}%
\TextAt{16}{232.67}{670.296}{\Character{79}}%
\TextAt{9}{242.52}{670.296}{\Character{41}}%
\NativeMathAt{250.393}{670.296}{11.955}{\cong}%
\TextAt{13}{263.012}{670.296}{M}%
\TextAt{16}{278.242}{670.296}{\Character{10}}%
\TextAt{27}{287.541}{668.503}{o}%
\TextAt{13}{294.592}{670.296}{N}%
\TextAt{13}{304.563}{670.296}{\Character{59}}%
\TextAt{13}{333.697}{670.296}{a}%
\TextAt{16}{342.52}{670.296}{\Character{10}}%
\TextAt{9}{354.475}{670.296}{\Character{40}}%
\TextAt{13}{359.028}{670.296}{b}%
\TextAt{16}{366.662}{670.296}{\Character{10}}%
\TextAt{13}{378.617}{670.296}{c}%
\TextAt{9}{383.655}{670.296}{\Character{41}}%
\NativeArrowAt{391.528}{670.296}{11.955}{19.26}{\longmapsto}%
\TextAt{13}{414.11}{670.296}{ca}%
\TextAt{16}{427.949}{670.296}{\Character{10}}%
\TextAt{13}{439.905}{670.296}{b\Character{58}}%
\TextAt{8}{79.226}{647.652}{Tensoring}%
\TextAt{8}{132.941}{647.652}{with}%
\TextAt{18}{159.601}{647.652}{C}%
\TextAt{8}{168.235}{647.652}{,}%
\TextAt{8}{175.384}{647.652}{it}%
\TextAt{8}{187.088}{647.652}{induces}%
\TextAt{8}{228.764}{647.652}{an}%
\TextAt{8}{245.023}{647.652}{isomorphism}%
\TextAt{13}{211.624}{620.699}{M}%
\TextAt{25}{222.949}{618.905}{C}%
\TextAt{16}{231.86}{620.699}{\Character{10}}%
\TextAt{21}{241.159}{618.905}{L}%
\TextAt{29}{246.918}{617.533}{C}%
\TextAt{9}{254.888}{620.699}{\Character{40}}%
\TextAt{13}{259.44}{620.699}{N}%
\TextAt{25}{268.814}{618.905}{C}%
\TextAt{16}{277.725}{620.699}{\Character{10}}%
\TextAt{21}{287.023}{618.905}{K}%
\TextAt{29}{294.165}{617.533}{C}%
\TextAt{13}{302.135}{620.699}{L}%
\TextAt{25}{310.099}{618.905}{C}%
\TextAt{9}{316.353}{620.699}{\Character{41}}%
\NativeMathAt{324.227}{620.699}{11.955}{\cong}%
\TextAt{13}{336.846}{620.699}{M}%
\TextAt{25}{348.171}{618.905}{C}%
\TextAt{16}{357.082}{620.699}{\Character{10}}%
\TextAt{21}{366.38}{618.905}{K}%
\TextAt{29}{373.522}{617.533}{C}%
\TextAt{13}{381.492}{620.699}{N}%
\TextAt{25}{390.865}{618.905}{C}%
\TextAt{13}{397.12}{620.699}{\Character{58}}%
\TextAt{8}{79.226}{593.746}{That}%
\TextAt{8}{108.492}{593.746}{this}%
\TextAt{8}{131.315}{593.746}{is}%
\TextAt{8}{143.079}{593.746}{an}%
\TextAt{8}{159.338}{593.746}{isometry}%
\TextAt{8}{207.206}{593.746}{of}%
\TextAt{8}{220.536}{593.746}{metrized}%
\TextAt{13}{268.656}{593.746}{K}%
\TextAt{25}{278.615}{591.952}{C}%
\TextAt{8}{284.869}{593.746}{-modules}%
\TextAt{8}{334.674}{593.746}{results}%
\TextAt{8}{371.867}{593.746}{from}%
\TextAt{8}{399.507}{593.746}{the}%
\TextAt{8}{419.664}{593.746}{distributivity}%
\TextAt{20}{199.209}{567.252}{\Character{88}}%
\TextAt{21}{201.765}{541.349}{\Character{28}}%
\TextAt{3}{206.443}{541.349}{\Character{106}}%
\TextAt{21}{208.796}{541.349}{\Character{27}}%
\TextAt{16}{216.477}{555.895}{\Character{104}}%
\TextAt{13}{225.024}{555.895}{\Character{59}}%
\TextAt{16}{234.165}{555.895}{\Character{105}}%
\TextAt{21}{238.815}{554.102}{M}%
\TextAt{22}{246.947}{553.105}{\Character{28}}%
\TextAt{16}{252.115}{555.895}{\Character{104}}%
\TextAt{13}{260.662}{555.895}{\Character{59}}%
\TextAt{16}{269.803}{555.895}{\Character{105}}%
\TextAt{21}{274.452}{554.102}{N}%
\TextAt{22}{281.174}{553.105}{\Character{27}}%
\TextAt{9}{289.989}{555.895}{\Character{61}}%
\TextAt{20}{302.414}{579.925}{\Character{48}}%
\TextAt{20}{302.414}{558.406}{\Character{64}}%
\TextAt{20}{312.875}{567.252}{\Character{88}}%
\TextAt{21}{315.431}{541.349}{\Character{28}}%
\TextAt{3}{320.109}{541.349}{\Character{106}}%
\TextAt{21}{322.462}{541.349}{\Character{27}}%
\TextAt{16}{330.143}{555.895}{\Character{104}}%
\TextAt{13}{338.69}{555.895}{\Character{59}}%
\TextAt{16}{347.832}{555.895}{\Character{105}}%
\TextAt{21}{352.481}{554.102}{M}%
\TextAt{22}{360.614}{553.105}{\Character{28}}%
\TextAt{20}{365.781}{579.925}{\Character{49}}%
\TextAt{20}{365.781}{558.406}{\Character{65}}%
\TextAt{16}{378.234}{555.895}{\Character{104}}%
\TextAt{13}{386.781}{555.895}{\Character{59}}%
\TextAt{16}{395.923}{555.895}{\Character{105}}%
\TextAt{21}{400.572}{554.102}{N}%
\TextAt{22}{407.293}{553.105}{\Character{27}}%
\TextAt{8}{79.226}{518.164}{by}%
\TextAt{8}{95.485}{518.164}{applying}%
\TextAt{8}{143.294}{518.164}{mathematical}%
\TextAt{8}{216.794}{518.164}{grammar.}%
\Vector{q 0 G 0 g 0.399 w 452.334 -255.836 m 452.334 -247.766 l S Q}%
\Vector{q 0 G 0 g 0.399 w 452.534 -247.966 m 459.707 -247.966 l S Q}%
\Vector{q 0 G 0 g 0.399 w 452.534 -255.637 m 459.707 -255.637 l S Q}%
\Vector{q 0 G 0 g 0.399 w 459.906 -255.836 m 459.906 -247.766 l S Q}%
\TextAt{8}{97.159}{489.748}{The}%
\TextAt{8}{121.721}{489.748}{Riemann–Roch}%
\TextAt{8}{203.859}{489.748}{problem}%
\TextAt{8}{249.892}{489.748}{in}%
\TextAt{8}{264.053}{489.748}{Grothendieck’s}%
\TextAt{8}{344.697}{489.748}{perspective}%
\TextAt{8}{406.391}{489.748}{is}%
\TextAt{8}{418.663}{489.748}{the}%
\TextAt{8}{439.328}{489.748}{task}%
\TextAt{8}{464.942}{489.748}{of}%
\TextAt{8}{478.78}{489.748}{computing}%
\TextAt{8}{79.226}{472.702}{the}%
\TextAt{8}{98.19}{472.702}{Chern}%
\TextAt{8}{132.11}{472.702}{character}%
\TextAt{8}{181.979}{472.702}{ch}%
\TextAt{9}{193.683}{472.702}{\Character{40}}%
\TextAt{13}{198.235}{472.702}{i}%
\TextAt{3}{202.228}{470.909}{\Character{3}}%
\TextAt{13}{206.961}{472.702}{M}%
\TextAt{9}{219.534}{472.702}{\Character{41}}%
\TextAt{8}{226.792}{472.702}{for}%
\TextAt{8}{243.484}{472.702}{a}%
\TextAt{8}{252.047}{472.702}{projective}%
\TextAt{8}{304.833}{472.702}{metrized}%
\TextAt{16}{351.76}{472.702}{\Character{79}}%
\TextAt{8}{361.61}{472.702}{-module}%
\TextAt{13}{405.608}{472.702}{M}%
\TextAt{8}{420.887}{472.702}{in}%
\TextAt{8}{433.347}{472.702}{terms}%
\TextAt{8}{464.733}{472.702}{of}%
\TextAt{8}{476.87}{472.702}{ch}%
\TextAt{9}{488.575}{472.702}{\Character{40}}%
\TextAt{13}{493.127}{472.702}{M}%
\TextAt{9}{505.7}{472.702}{\Character{41}}%
\TextAt{8}{510.253}{472.702}{.}%
\TextAt{8}{518.304}{472.702}{By}%
\TextAt{8}{79.226}{455.656}{(6.6),}%
\TextAt{8}{110.263}{455.656}{this}%
\TextAt{8}{132.848}{455.656}{amounts}%
\TextAt{8}{179.834}{455.656}{to}%
\TextAt{8}{193.907}{455.656}{computing}%
\TextAt{8}{251.557}{455.656}{det}%
\TextAt{9}{267.816}{455.656}{\Character{40}}%
\TextAt{13}{272.368}{455.656}{i}%
\TextAt{3}{276.362}{453.863}{\Character{3}}%
\TextAt{13}{281.094}{455.656}{M}%
\TextAt{9}{293.668}{455.656}{\Character{41}}%
\TextAt{8}{301.88}{455.656}{in}%
\TextAt{8}{315.296}{455.656}{terms}%
\TextAt{8}{347.636}{455.656}{of}%
\TextAt{8}{360.729}{455.656}{det}%
\TextAt{13}{378.981}{455.656}{M}%
\TextAt{8}{391.554}{455.656}{.}%
\TextAt{8}{399.923}{455.656}{But}%
\TextAt{8}{422.927}{455.656}{det}%
\TextAt{13}{441.179}{455.656}{M}%
\TextAt{8}{457.412}{455.656}{is}%
\TextAt{8}{468.939}{455.656}{an}%
\TextAt{8}{484.961}{455.656}{invertible}%
\TextAt{8}{79.226}{438.61}{metrized}%
\TextAt{16}{127.021}{438.61}{\Character{79}}%
\TextAt{8}{136.871}{438.61}{-module}%
\TextAt{8}{181.736}{438.61}{and}%
\TextAt{8}{204.174}{438.61}{is}%
\TextAt{8}{215.613}{438.61}{therefore}%
\TextAt{8}{264.389}{438.61}{isometric}%
\TextAt{8}{314.204}{438.61}{by}%
\TextAt{8}{330.138}{438.61}{(4.5)}%
\TextAt{8}{357.788}{438.61}{to}%
\TextAt{8}{371.774}{438.61}{the}%
\TextAt{8}{391.606}{438.61}{metrized}%
\TextAt{16}{439.4}{438.61}{\Character{79}}%
\TextAt{8}{449.25}{438.61}{-module}%
\TextAt{13}{494.116}{438.61}{L}%
\TextAt{9}{502.081}{438.61}{\Character{40}}%
\TextAt{14}{506.633}{438.61}{A}%
\TextAt{9}{515.212}{438.61}{\Character{41}}%
\TextAt{8}{523.337}{438.61}{of}%
\TextAt{8}{79.226}{421.564}{a}%
\TextAt{8}{88.982}{421.564}{replete}%
\TextAt{8}{127.346}{421.564}{ideal}%
\TextAt{14}{155.309}{421.564}{A}%
\TextAt{8}{167.785}{421.564}{of}%
\TextAt{13}{181.115}{421.564}{L}%
\TextAt{8}{189.08}{421.564}{.}%
\TextAt{13}{197.528}{421.564}{N}%
\TextAt{21}{206.902}{419.704}{L}%
\TextAt{3}{212.661}{419.704}{\Character{106}}%
\TextAt{21}{215.013}{419.704}{K}%
\TextAt{9}{223.216}{421.564}{\Character{40}}%
\TextAt{14}{227.768}{421.564}{A}%
\TextAt{9}{236.347}{421.564}{\Character{41}}%
\TextAt{8}{244.797}{421.564}{is}%
\TextAt{8}{256.561}{421.564}{then}%
\TextAt{8}{283.221}{421.564}{a}%
\TextAt{8}{292.976}{421.564}{replete}%
\TextAt{8}{331.341}{421.564}{ideal}%
\TextAt{8}{359.304}{421.564}{of}%
\TextAt{13}{372.634}{421.564}{K}%
\TextAt{8}{383.451}{421.564}{,}%
\TextAt{8}{390.601}{421.564}{and}%
\TextAt{8}{413.363}{421.564}{we}%
\TextAt{8}{430.591}{421.564}{put}%
\TextAt{13}{231.844}{394.611}{N}%
\TextAt{21}{241.218}{392.751}{L}%
\TextAt{3}{246.977}{392.751}{\Character{106}}%
\TextAt{21}{249.329}{392.751}{K}%
\TextAt{9}{257.532}{394.611}{\Character{40}}%
\TextAt{8}{262.084}{394.611}{det}%
\TextAt{13}{280.336}{394.611}{M}%
\TextAt{9}{292.909}{394.611}{\Character{41}}%
\TextAt{9}{300.764}{394.611}{\Character{58}\Character{61}}%
\TextAt{13}{316.459}{394.611}{L}%
\TextAt{9}{324.424}{394.611}{\Character{40}}%
\TextAt{13}{328.976}{394.611}{N}%
\TextAt{21}{338.35}{392.751}{L}%
\TextAt{3}{344.109}{392.751}{\Character{106}}%
\TextAt{21}{346.461}{392.751}{K}%
\TextAt{9}{354.664}{394.611}{\Character{40}}%
\TextAt{14}{359.216}{394.611}{A}%
\TextAt{9}{367.795}{394.611}{\Character{41}\Character{41}}%
\TextAt{13}{376.9}{394.611}{\Character{58}}%
\TextAt{8}{79.226}{367.658}{This}%
\TextAt{8}{106.492}{367.658}{is}%
\TextAt{8}{118.802}{367.658}{an}%
\TextAt{8}{135.607}{367.658}{invertible}%
\TextAt{8}{187.859}{367.658}{metrized}%
\TextAt{14}{236.525}{367.658}{o}%
\TextAt{8}{242.369}{367.658}{-module}%
\TextAt{8}{288.106}{367.658}{which}%
\TextAt{8}{322.138}{367.658}{is}%
\TextAt{8}{334.448}{367.658}{well}%
\TextAt{8}{358.725}{367.658}{determined}%
\TextAt{8}{420.398}{367.658}{by}%
\TextAt{13}{437.203}{367.658}{M}%
\TextAt{8}{454.22}{367.658}{up}%
\TextAt{8}{471.671}{367.658}{to}%
\TextAt{8}{486.527}{367.658}{isometry.}%
\TextAt{8}{79.226}{350.612}{With}%
\TextAt{8}{109.461}{350.612}{this}%
\TextAt{8}{132.283}{350.612}{notation}%
\TextAt{8}{179.124}{350.612}{we}%
\TextAt{8}{196.351}{350.612}{first}%
\TextAt{8}{220.476}{350.612}{establish}%
\TextAt{8}{268.728}{350.612}{the}%
\TextAt{8}{288.884}{350.612}{following}%
\TextAt{8}{338.319}{350.612}{theorem.}%
\TextAt{12}{79.226}{325.185}{Theorem}%
\TextAt{12}{136.24}{325.185}{(7.2).}%
\TextAt{8}{171.556}{325.185}{For}%
\TextAt{8}{192.525}{325.185}{any}%
\TextAt{8}{214.642}{325.185}{projective}%
\TextAt{8}{268.62}{325.185}{metrized}%
\TextAt{16}{316.739}{325.185}{\Character{79}}%
\TextAt{8}{326.589}{325.185}{-module}%
\TextAt{13}{371.78}{325.185}{M}%
\TextAt{8}{388.251}{325.185}{one}%
\TextAt{8}{409.71}{325.185}{has:}%
\TextAt{8}{243.881}{298.231}{rk}%
\TextAt{9}{254.616}{298.231}{\Character{40}}%
\TextAt{13}{259.169}{298.231}{i}%
\TextAt{3}{263.162}{296.438}{\Character{3}}%
\TextAt{13}{267.894}{298.231}{M}%
\TextAt{9}{280.468}{298.231}{\Character{41}}%
\TextAt{9}{288.323}{298.231}{\Character{61}}%
\TextAt{8}{300.767}{298.231}{rk}%
\TextAt{9}{311.502}{298.231}{\Character{40}}%
\TextAt{13}{316.055}{298.231}{M}%
\TextAt{9}{328.628}{298.231}{\Character{41}}%
\TextAt{8}{335.173}{298.231}{rk}%
\TextAt{9}{345.909}{298.231}{\Character{40}}%
\TextAt{16}{350.461}{298.231}{\Character{79}}%
\TextAt{9}{360.311}{298.231}{\Character{41}}%
\TextAt{13}{364.863}{298.231}{\Character{59}}%
\TextAt{8}{195.874}{271.278}{det}%
\TextAt{9}{212.133}{271.278}{\Character{40}}%
\TextAt{13}{216.685}{271.278}{i}%
\TextAt{3}{220.679}{269.485}{\Character{3}}%
\TextAt{13}{225.411}{271.278}{M}%
\TextAt{9}{237.985}{271.278}{\Character{41}}%
\NativeMathAt{245.858}{271.278}{11.955}{\cong}%
\TextAt{13}{258.477}{271.278}{N}%
\TextAt{21}{267.851}{269.419}{L}%
\TextAt{3}{273.61}{269.419}{\Character{106}}%
\TextAt{21}{275.962}{269.419}{K}%
\TextAt{9}{284.165}{271.278}{\Character{40}}%
\TextAt{8}{288.717}{271.278}{det}%
\TextAt{13}{306.969}{271.278}{M}%
\TextAt{9}{319.543}{271.278}{\Character{41}}%
\TextAt{16}{326.752}{271.278}{\Character{10}}%
\TextAt{27}{336.05}{269.485}{o}%
\TextAt{9}{343.101}{271.278}{\Character{40}}%
\TextAt{8}{347.653}{271.278}{det}%
\TextAt{13}{365.905}{271.278}{i}%
\TextAt{3}{369.898}{269.485}{\Character{3}}%
\TextAt{16}{374.631}{271.278}{\Character{79}}%
\TextAt{9}{384.48}{271.278}{\Character{41}}%
\TextAt{7}{389.033}{276.215}{rk}%
\TextAt{17}{396.804}{276.214}{\Character{40}}%
\TextAt{21}{400.097}{276.215}{M}%
\TextAt{17}{409.079}{276.214}{\Character{41}}%
\TextAt{13}{412.87}{271.278}{\Character{58}}%
\TextAt{8}{79.226}{248.634}{Here}%
\TextAt{8}{106.855}{248.634}{we}%
\TextAt{8}{124.082}{248.634}{have}%
\TextAt{8}{151.077}{248.634}{rk}%
\TextAt{9}{161.812}{248.634}{\Character{40}}%
\TextAt{16}{166.365}{248.634}{\Character{79}}%
\TextAt{9}{176.215}{248.634}{\Character{41}}%
\TextAt{9}{184.07}{248.634}{\Character{61}}%
\TextAt{9}{196.491}{248.634}{\Character{91}}%
\TextAt{13}{199.765}{248.634}{L}%
\TextAt{9}{211.05}{248.634}{\Character{58}}%
\TextAt{13}{217.623}{248.634}{K}%
\TextAt{9}{228.44}{248.634}{\Character{93}}%
\TextAt{8}{231.692}{248.634}{.}%
\TextAt{19}{79.226}{214.826}{Proof.}%
\TextAt{8}{116}{214.826}{One}%
\TextAt{8}{140.712}{214.826}{has}%
\TextAt{13}{161.585}{214.826}{M}%
\TextAt{21}{172.91}{213.033}{K}%
\TextAt{9}{184.434}{214.826}{\Character{58}\Character{61}}%
\TextAt{13}{200.111}{214.826}{M}%
\TextAt{16}{215.341}{214.826}{\Character{10}}%
\TextAt{27}{224.64}{213.033}{o}%
\TextAt{13}{231.69}{214.826}{K}%
\TextAt{9}{245.829}{214.826}{\Character{61}}%
\TextAt{13}{258.254}{214.826}{M}%
\TextAt{16}{273.485}{214.826}{\Character{10}}%
\TextAt{3}{282.783}{213.033}{\Character{79}}%
\TextAt{16}{292.881}{214.826}{\Character{79}}%
\TextAt{16}{305.374}{214.826}{\Character{10}}%
\TextAt{27}{314.686}{213.033}{o}%
\TextAt{13}{321.737}{214.826}{K}%
\TextAt{9}{335.875}{214.826}{\Character{61}}%
\TextAt{13}{348.301}{214.826}{M}%
\TextAt{16}{363.531}{214.826}{\Character{10}}%
\TextAt{3}{372.83}{213.033}{\Character{79}}%
\TextAt{13}{382.927}{214.826}{L}%
\TextAt{9}{394.213}{214.826}{\Character{61}\Character{58}}%
\TextAt{13}{409.89}{214.826}{M}%
\TextAt{21}{421.215}{213.033}{L}%
\TextAt{8}{431.369}{214.826}{and}%
\TextAt{8}{454.132}{214.826}{therefore}%
\TextAt{8}{121.072}{187.873}{rk}%
\TextAt{9}{131.807}{187.873}{\Character{40}}%
\TextAt{13}{136.36}{187.873}{i}%
\TextAt{3}{140.353}{186.079}{\Character{3}}%
\TextAt{13}{145.085}{187.873}{M}%
\TextAt{9}{157.659}{187.873}{\Character{41}}%
\TextAt{9}{165.514}{187.873}{\Character{61}}%
\TextAt{8}{177.958}{187.873}{dim}%
\TextAt{21}{197.468}{186.079}{K}%
\TextAt{9}{205.671}{187.873}{\Character{40}}%
\TextAt{13}{210.223}{187.873}{M}%
\TextAt{21}{221.548}{186.079}{K}%
\TextAt{9}{229.751}{187.873}{\Character{41}}%
\TextAt{9}{237.606}{187.873}{\Character{61}}%
\TextAt{8}{250.05}{187.873}{dim}%
\TextAt{21}{269.56}{186.079}{K}%
\TextAt{9}{277.763}{187.873}{\Character{40}}%
\TextAt{13}{282.315}{187.873}{M}%
\TextAt{21}{293.64}{186.079}{L}%
\TextAt{9}{299.897}{187.873}{\Character{41}}%
\TextAt{9}{307.752}{187.873}{\Character{61}}%
\TextAt{8}{320.196}{187.873}{dim}%
\TextAt{21}{339.707}{186.079}{L}%
\TextAt{9}{345.964}{187.873}{\Character{40}}%
\TextAt{13}{350.516}{187.873}{M}%
\TextAt{21}{361.841}{186.079}{L}%
\TextAt{9}{368.098}{187.873}{\Character{41}\Character{91}}%
\TextAt{13}{375.902}{187.873}{L}%
\TextAt{9}{387.188}{187.873}{\Character{58}}%
\TextAt{13}{393.76}{187.873}{K}%
\TextAt{9}{404.578}{187.873}{\Character{93}}%
\TextAt{9}{411.129}{187.873}{\Character{61}}%
\TextAt{8}{423.576}{187.873}{rk}%
\TextAt{9}{434.312}{187.873}{\Character{40}}%
\TextAt{13}{438.864}{187.873}{M}%
\TextAt{9}{451.437}{187.873}{\Character{41}}%
\TextAt{8}{457.982}{187.873}{rk}%
\TextAt{9}{468.718}{187.873}{\Character{40}}%
\TextAt{16}{473.27}{187.873}{\Character{79}}%
\TextAt{9}{483.12}{187.873}{\Character{41}}%
\TextAt{13}{487.672}{187.873}{\Character{58}}%
\TextAt{8}{79.226}{160.92}{In}%
\TextAt{8}{93.847}{160.92}{order}%
\TextAt{8}{124.417}{160.92}{to}%
\TextAt{8}{138.727}{160.92}{prove}%
\TextAt{8}{170.277}{160.92}{the}%
\TextAt{8}{190.433}{160.92}{second}%
\TextAt{8}{228.211}{160.92}{equation,}%
\TextAt{8}{279.272}{160.92}{we}%
\TextAt{8}{296.499}{160.92}{first}%
\TextAt{8}{320.625}{160.92}{reduce}%
\TextAt{8}{357.686}{160.92}{to}%
\TextAt{8}{371.996}{160.92}{a}%
\TextAt{8}{381.752}{160.92}{special}%
\TextAt{8}{419.853}{160.92}{case.}%
\TextAt{8}{449.175}{160.92}{Let}%
\TextAt{13}{126.643}{133.966}{\Character{21}}%
\TextAt{9}{133.472}{133.966}{\Character{40}}%
\TextAt{13}{138.024}{133.966}{M}%
\TextAt{9}{150.598}{133.966}{\Character{41}}%
\TextAt{9}{158.453}{133.966}{\Character{61}}%
\TextAt{8}{170.897}{133.966}{det}%
\TextAt{9}{187.156}{133.966}{\Character{40}}%
\TextAt{13}{191.708}{133.966}{i}%
\TextAt{3}{195.701}{132.173}{\Character{3}}%
\TextAt{13}{200.434}{133.966}{M}%
\TextAt{9}{213.007}{133.966}{\Character{41}}%
\TextAt{8}{241.47}{133.966}{and}%
\TextAt{13}{284.246}{133.966}{\Character{26}}%
\TextAt{9}{290.283}{133.966}{\Character{40}}%
\TextAt{13}{294.836}{133.966}{M}%
\TextAt{9}{307.409}{133.966}{\Character{41}}%
\TextAt{9}{315.264}{133.966}{\Character{61}}%
\TextAt{13}{327.708}{133.966}{N}%
\TextAt{21}{337.082}{132.107}{L}%
\TextAt{3}{342.841}{132.107}{\Character{106}}%
\TextAt{21}{345.193}{132.107}{K}%
\TextAt{9}{353.396}{133.966}{\Character{40}}%
\TextAt{8}{357.948}{133.966}{det}%
\TextAt{13}{376.2}{133.966}{M}%
\TextAt{9}{388.773}{133.966}{\Character{41}}%
\TextAt{16}{395.982}{133.966}{\Character{10}}%
\TextAt{27}{405.281}{132.173}{o}%
\TextAt{9}{412.332}{133.966}{\Character{40}}%
\TextAt{8}{416.884}{133.966}{det}%
\TextAt{13}{435.135}{133.966}{i}%
\TextAt{3}{439.129}{132.173}{\Character{3}}%
\TextAt{16}{443.861}{133.966}{\Character{79}}%
\TextAt{9}{453.711}{133.966}{\Character{41}}%
\TextAt{7}{458.263}{138.903}{rk}%
\TextAt{17}{466.034}{138.903}{\Character{40}}%
\TextAt{21}{469.327}{138.903}{M}%
\TextAt{17}{478.309}{138.903}{\Character{41}}%
\TextAt{13}{482.1}{133.966}{\Character{58}}%
\TextAt{8}{79.226}{107.013}{If}%
\TextAt{9}{90.957}{107.013}{0}%
\NativeMathAt{100.197}{107.013}{11.955}{\rightarrow}%
\TextAt{13}{115.54}{107.013}{M}%
\TextAt{3}{128.113}{111.352}{\Character{48}}%
\NativeMathAt{134.295}{107.013}{11.955}{\rightarrow}%
\TextAt{13}{149.638}{107.013}{M}%
\NativeMathAt{165.598}{107.013}{11.955}{\rightarrow}%
\TextAt{13}{180.94}{107.013}{M}%
\TextAt{3}{193.514}{111.352}{\Character{48}\Character{48}}%
\NativeMathAt{201.993}{107.013}{11.955}{\rightarrow}%
\TextAt{9}{217.335}{107.013}{0}%
\TextAt{8}{227.125}{107.013}{is}%
\TextAt{8}{238.927}{107.013}{a}%
\TextAt{8}{248.722}{107.013}{short}%
\TextAt{8}{278.744}{107.013}{exact}%
\TextAt{8}{309.675}{107.013}{sequence}%
\TextAt{8}{358.216}{107.013}{of}%
\TextAt{8}{371.585}{107.013}{projective}%
\TextAt{8}{425.602}{107.013}{metrized}%
\TextAt{16}{473.76}{107.013}{\Character{79}}%
\TextAt{8}{483.61}{107.013}{-modules,}%
\TextAt{8}{79.226}{89.967}{one}%
\TextAt{8}{100.686}{89.967}{has}%
\TextAt{13}{143.862}{72.921}{\Character{21}}%
\TextAt{9}{150.69}{72.921}{\Character{40}}%
\TextAt{13}{155.242}{72.921}{M}%
\TextAt{9}{167.816}{72.921}{\Character{41}}%
\NativeMathAt{175.689}{72.921}{11.955}{\cong}%
\TextAt{13}{188.308}{72.921}{\Character{21}}%
\TextAt{9}{195.137}{72.921}{\Character{40}}%
\TextAt{13}{199.689}{72.921}{M}%
\TextAt{3}{212.263}{77.858}{\Character{48}}%
\TextAt{9}{215.058}{72.921}{\Character{41}}%
\TextAt{16}{222.267}{72.921}{\Character{10}}%
\TextAt{27}{231.565}{71.128}{o}%
\TextAt{13}{238.616}{72.921}{\Character{21}}%
\TextAt{9}{245.445}{72.921}{\Character{40}}%
\TextAt{13}{249.997}{72.921}{M}%
\TextAt{3}{262.571}{77.858}{\Character{48}\Character{48}}%
\TextAt{9}{267.663}{72.921}{\Character{41}}%
\TextAt{8}{296.126}{72.921}{and}%
\TextAt{13}{338.901}{72.921}{\Character{26}}%
\TextAt{9}{344.939}{72.921}{\Character{40}}%
\TextAt{13}{349.491}{72.921}{M}%
\TextAt{9}{362.065}{72.921}{\Character{41}}%
\NativeMathAt{369.938}{72.921}{11.955}{\cong}%
\TextAt{13}{382.557}{72.921}{\Character{26}}%
\TextAt{9}{388.595}{72.921}{\Character{40}}%
\TextAt{13}{393.147}{72.921}{M}%
\TextAt{3}{405.721}{77.858}{\Character{48}}%
\TextAt{9}{408.516}{72.921}{\Character{41}}%
\TextAt{16}{415.725}{72.921}{\Character{10}}%
\TextAt{27}{425.023}{71.128}{o}%
\TextAt{13}{432.074}{72.921}{\Character{26}}%
\TextAt{9}{438.112}{72.921}{\Character{40}}%
\TextAt{13}{442.664}{72.921}{M}%
\TextAt{3}{455.238}{77.858}{\Character{48}\Character{48}}%
\TextAt{9}{460.33}{72.921}{\Character{41}}%
\TextAt{13}{464.882}{72.921}{\Character{58}}%
\TextAt{8}{517.802}{72.921}{(*)}%
\TextAt{8}{297.211}{43.033}{228}%
\Destination{522.36}{89.965}{nav.AMS.38}%
\Destination{143.86}{89.965}{nav.equation.3.7.1}%
\Destination{78.23}{809.815}{nav.page.228}%
\Destination{79.36}{776.155}{nav.print.248}%
\Destination{79.36}{339.155}{nav.number.3.7.2}%
\Hotlink{78.726}{453.455}{25.078}{11.047}{nav.number.3.6.6}%
\Hotlink{329.638}{436.455}{25.078}{11.047}{nav.number.3.4.5}%
}

% Page 241 (body page 229)
\Page{%
\TextAt{2}{79.226}{796.943}{III.}%
\TextAt{2}{98.317}{796.943}{Riemann–Roch}%
\TextAt{2}{174.375}{796.943}{Theory}%
\Vector{q 0 G 0 g 0.299 w 7.226 18.729 m 460.769 18.729 l S Q}%
\TextAt{8}{79.226}{761.124}{The}%
\TextAt{8}{102.264}{761.124}{isomorphism}%
\TextAt{8}{169.668}{761.124}{on}%
\TextAt{8}{184.912}{761.124}{the}%
\TextAt{8}{204.052}{761.124}{left}%
\TextAt{8}{223.516}{761.124}{follows}%
\TextAt{8}{260.936}{761.124}{from}%
\TextAt{8}{287.561}{761.124}{the}%
\TextAt{8}{306.702}{761.124}{exact}%
\TextAt{8}{336.578}{761.124}{sequence}%
\TextAt{9}{384.065}{761.124}{0}%
\NativeMathAt{393.239}{761.124}{11.955}{\rightarrow}%
\TextAt{13}{408.515}{761.124}{i}%
\TextAt{3}{412.508}{759.33}{\Character{3}}%
\TextAt{13}{417.24}{761.124}{M}%
\TextAt{3}{429.814}{765.462}{\Character{48}}%
\NativeMathAt{435.93}{761.124}{11.955}{\rightarrow}%
\TextAt{13}{451.206}{761.124}{i}%
\TextAt{3}{455.199}{759.33}{\Character{3}}%
\TextAt{13}{459.932}{761.124}{M}%
\NativeMathAt{475.826}{761.124}{11.955}{\rightarrow}%
\TextAt{13}{491.102}{761.124}{i}%
\TextAt{3}{495.096}{759.33}{\Character{3}}%
\TextAt{13}{499.828}{761.124}{M}%
\TextAt{3}{512.401}{765.462}{\Character{48}\Character{48}}%
\NativeMathAt{520.814}{761.124}{11.955}{\rightarrow}%
\TextAt{9}{79.226}{744.078}{0}%
\TextAt{8}{88.633}{744.078}{by}%
\TextAt{8}{104.549}{744.078}{(4.7),}%
\TextAt{8}{135.501}{744.078}{and}%
\TextAt{8}{157.921}{744.078}{the}%
\TextAt{8}{177.734}{744.078}{one}%
\TextAt{8}{198.85}{744.078}{on}%
\TextAt{8}{214.766}{744.078}{the}%
\TextAt{8}{234.579}{744.078}{right}%
\TextAt{8}{262.534}{744.078}{from}%
\TextAt{8}{289.831}{744.078}{(4.7)}%
\TextAt{8}{317.463}{744.078}{also,}%
\TextAt{8}{343.92}{744.078}{from}%
\TextAt{8}{371.217}{744.078}{the}%
\TextAt{8}{391.03}{744.078}{multiplicativity}%
\TextAt{8}{473.297}{744.078}{of}%
\TextAt{8}{486.284}{744.078}{the}%
\TextAt{8}{506.097}{744.078}{norm}%
\TextAt{13}{79.226}{727.032}{N}%
\TextAt{21}{88.6}{725.172}{L}%
\TextAt{3}{94.359}{725.172}{\Character{106}}%
\TextAt{21}{96.711}{725.172}{K}%
\TextAt{8}{108.783}{727.032}{and}%
\TextAt{8}{131.518}{727.032}{the}%
\TextAt{8}{151.647}{727.032}{additivity}%
\TextAt{8}{205.286}{727.032}{of}%
\TextAt{8}{218.588}{727.032}{the}%
\TextAt{8}{238.716}{727.032}{rank}%
\TextAt{8}{265.683}{727.032}{rk.}%
\TextAt{8}{284.858}{727.032}{As}%
\TextAt{8}{302.117}{727.032}{in}%
\TextAt{8}{315.742}{727.032}{the}%
\TextAt{8}{335.87}{727.032}{proof}%
\TextAt{8}{366.412}{727.032}{of}%
\TextAt{8}{379.714}{727.032}{(5.6),}%
\TextAt{8}{410.919}{727.032}{we}%
\TextAt{8}{428.118}{727.032}{now}%
\TextAt{8}{452.479}{727.032}{make}%
\TextAt{8}{483.02}{727.032}{use}%
\TextAt{8}{503.208}{727.032}{of}%
\TextAt{8}{516.51}{727.032}{the}%
\TextAt{8}{79.226}{709.986}{fact}%
\TextAt{8}{102.166}{709.986}{that}%
\TextAt{8}{127.39}{709.986}{every}%
\TextAt{8}{158.137}{709.986}{projective}%
\TextAt{8}{211.969}{709.986}{metrized}%
\TextAt{16}{259.944}{709.986}{\Character{79}}%
\TextAt{8}{269.793}{709.986}{-module}%
\TextAt{13}{314.839}{709.986}{M}%
\TextAt{8}{331.165}{709.986}{projects}%
\TextAt{8}{375.301}{709.986}{via}%
\TextAt{8}{394.344}{709.986}{an}%
\TextAt{8}{410.458}{709.986}{admissible}%
\TextAt{8}{467.016}{709.986}{epimorphism}%
\TextAt{8}{79.226}{692.94}{onto}%
\TextAt{8}{105.744}{692.94}{a}%
\TextAt{8}{115.667}{692.94}{suitable}%
\TextAt{16}{159.472}{692.94}{\Character{79}}%
\TextAt{8}{169.322}{692.94}{-module}%
\TextAt{8}{214.681}{692.94}{of}%
\TextAt{8}{228.179}{692.94}{the}%
\TextAt{8}{248.504}{692.94}{form}%
\TextAt{13}{276.313}{692.94}{L}%
\TextAt{9}{284.277}{692.94}{\Character{40}}%
\TextAt{14}{288.83}{692.94}{A}%
\TextAt{9}{297.409}{692.94}{\Character{41}}%
\TextAt{8}{306.027}{692.94}{for}%
\TextAt{8}{324.08}{692.94}{some}%
\TextAt{8}{353.574}{692.94}{replete}%
\TextAt{8}{392.107}{692.94}{ideal}%
\TextAt{14}{420.238}{692.94}{A}%
\TextAt{8}{428.817}{692.94}{.}%
\TextAt{8}{437.771}{692.94}{Thus}%
\TextAt{8}{467.588}{692.94}{(*)}%
\TextAt{8}{486.622}{692.94}{allows}%
\TextAt{8}{521.651}{692.94}{us}%
\TextAt{8}{79.226}{675.894}{to}%
\TextAt{8}{93.614}{675.894}{reduce}%
\TextAt{8}{130.752}{675.894}{by}%
\TextAt{8}{147.089}{675.894}{induction}%
\TextAt{8}{199.195}{675.894}{on}%
\TextAt{8}{215.532}{675.894}{rk}%
\TextAt{9}{226.267}{675.894}{\Character{40}}%
\TextAt{13}{230.82}{675.894}{M}%
\TextAt{9}{243.393}{675.894}{\Character{41}}%
\TextAt{8}{251.92}{675.894}{to}%
\TextAt{8}{266.308}{675.894}{the}%
\TextAt{8}{286.542}{675.894}{case}%
\TextAt{13}{311.39}{675.894}{M}%
\TextAt{9}{327.417}{675.894}{\Character{61}}%
\TextAt{13}{339.974}{675.894}{L}%
\TextAt{9}{347.939}{675.894}{\Character{40}}%
\TextAt{14}{352.491}{675.894}{A}%
\TextAt{9}{361.07}{675.894}{\Character{41}}%
\TextAt{8}{365.622}{675.894}{.}%
\TextAt{8}{374.303}{675.894}{Here}%
\TextAt{8}{402.009}{675.894}{rk}%
\TextAt{9}{412.745}{675.894}{\Character{40}}%
\TextAt{13}{417.297}{675.894}{M}%
\TextAt{9}{429.871}{675.894}{\Character{41}}%
\TextAt{9}{437.857}{675.894}{\Character{61}}%
\TextAt{9}{450.41}{675.894}{1}%
\TextAt{8}{456.286}{675.894}{,}%
\TextAt{8}{463.532}{675.894}{so}%
\TextAt{8}{477.98}{675.894}{we}%
\TextAt{8}{495.284}{675.894}{have}%
\TextAt{8}{522.357}{675.894}{to}%
\TextAt{8}{79.226}{658.848}{establish}%
\TextAt{8}{127.477}{658.848}{the}%
\TextAt{8}{147.634}{658.848}{isomorphism}%
\TextAt{8}{212.856}{631.288}{det}%
\TextAt{9}{229.115}{631.288}{\Character{40}}%
\TextAt{13}{233.667}{631.288}{i}%
\TextAt{3}{237.661}{629.495}{\Character{3}}%
\TextAt{13}{242.393}{631.288}{L}%
\TextAt{9}{250.358}{631.288}{\Character{40}}%
\TextAt{14}{254.91}{631.288}{A}%
\TextAt{9}{263.489}{631.288}{\Character{41}\Character{41}}%
\TextAt{9}{275.887}{631.288}{\Character{61}}%
\TextAt{13}{288.34}{631.288}{L}%
\TextAt{9}{296.304}{631.288}{\Character{40}}%
\TextAt{13}{300.857}{631.288}{N}%
\TextAt{21}{310.231}{629.429}{L}%
\TextAt{3}{315.99}{629.429}{\Character{106}}%
\TextAt{21}{318.342}{629.429}{K}%
\TextAt{9}{326.545}{631.288}{\Character{40}}%
\TextAt{14}{331.097}{631.288}{A}%
\TextAt{9}{339.676}{631.288}{\Character{41}\Character{41}}%
\TextAt{16}{351.437}{631.288}{\Character{10}}%
\TextAt{27}{360.736}{629.495}{o}%
\TextAt{8}{367.787}{631.288}{det}%
\TextAt{27}{373.968}{623.996}{o}%
\TextAt{16}{386.038}{631.288}{\Character{79}}%
\TextAt{13}{395.888}{631.288}{\Character{58}}%
\TextAt{8}{79.226}{603.729}{For}%
\TextAt{8}{100.196}{603.729}{the}%
\TextAt{8}{120.352}{603.729}{underlying}%
\TextAt{14}{178.562}{603.729}{o}%
\TextAt{8}{184.406}{603.729}{-modules}%
\TextAt{8}{234.211}{603.729}{this}%
\TextAt{8}{257.034}{603.729}{amounts}%
\TextAt{8}{304.256}{603.729}{to}%
\TextAt{8}{318.567}{603.729}{the}%
\TextAt{8}{338.723}{603.729}{identity}%
\TextAt{8}{240.906}{576.169}{det}%
\TextAt{27}{247.087}{568.876}{o}%
\TextAt{14}{259.157}{576.169}{A}%
\TextAt{21}{267.737}{574.376}{f}%
\TextAt{9}{276.502}{576.169}{\Character{61}}%
\TextAt{13}{288.928}{576.169}{N}%
\TextAt{21}{298.302}{574.309}{L}%
\TextAt{3}{304.061}{574.309}{\Character{106}}%
\TextAt{21}{306.413}{574.309}{K}%
\TextAt{9}{314.616}{576.169}{\Character{40}}%
\TextAt{14}{319.168}{576.169}{A}%
\TextAt{21}{327.747}{574.376}{f}%
\TextAt{9}{333.192}{576.169}{\Character{41}}%
\TextAt{8}{339.737}{576.169}{det}%
\TextAt{27}{345.918}{568.876}{o}%
\TextAt{16}{357.988}{576.169}{\Character{79}}%
\TextAt{13}{367.838}{576.169}{\Character{59}}%
\TextAt{8}{511.944}{576.169}{(**)}%
\TextAt{8}{79.226}{548.609}{which}%
\TextAt{8}{114.144}{548.609}{has}%
\TextAt{8}{136.45}{548.609}{to}%
\TextAt{8}{152.192}{548.609}{be}%
\TextAt{8}{169.548}{548.609}{viewed}%
\TextAt{8}{209.344}{548.609}{as}%
\TextAt{8}{225.146}{548.609}{inside}%
\TextAt{8}{259.802}{548.609}{det}%
\TextAt{21}{276.061}{546.816}{K}%
\TextAt{13}{286.256}{548.609}{L}%
\TextAt{8}{299.55}{548.609}{and}%
\TextAt{8}{323.744}{548.609}{which}%
\TextAt{8}{358.662}{548.609}{is}%
\TextAt{8}{371.858}{548.609}{proved}%
\TextAt{8}{411.343}{548.609}{as}%
\TextAt{8}{427.145}{548.609}{follows.}%
\TextAt{8}{474.428}{548.609}{If}%
\TextAt{16}{487.552}{548.609}{\Character{79}}%
\TextAt{8}{502.731}{548.609}{and}%
\TextAt{14}{526.925}{548.609}{o}%
\TextAt{8}{79.226}{531.563}{were}%
\TextAt{8}{106.508}{531.563}{principal}%
\TextAt{8}{155.584}{531.563}{ideal}%
\TextAt{8}{183.845}{531.563}{domains,}%
\TextAt{8}{233.713}{531.563}{it}%
\TextAt{8}{245.716}{531.563}{would}%
\TextAt{8}{280.159}{531.563}{be}%
\TextAt{8}{296.382}{531.563}{obvious.}%
\TextAt{8}{344.174}{531.563}{In}%
\TextAt{8}{359.094}{531.563}{fact,}%
\TextAt{8}{385.805}{531.563}{in}%
\TextAt{8}{399.757}{531.563}{that}%
\TextAt{8}{425.424}{531.563}{case}%
\TextAt{8}{450.494}{531.563}{we}%
\TextAt{8}{468.02}{531.563}{could}%
\TextAt{8}{499.534}{531.563}{choose}%
\TextAt{8}{79.226}{514.518}{a}%
\TextAt{8}{89.841}{514.518}{generator}%
\TextAt{13}{142.741}{514.518}{\Character{11}}%
\TextAt{8}{155.02}{514.518}{of}%
\TextAt{14}{169.21}{514.518}{A}%
\TextAt{21}{177.789}{512.724}{f}%
\TextAt{8}{187.99}{514.518}{and}%
\TextAt{8}{211.613}{514.518}{an}%
\TextAt{8}{228.731}{514.518}{integral}%
\TextAt{8}{272.199}{514.518}{basis}%
\TextAt{13}{301.799}{514.518}{\Character{33}}%
\TextAt{17}{309.088}{512.724}{1}%
\TextAt{13}{313.82}{514.517}{\Character{59}}%
\TextAt{13}{319.044}{514.517}{\Character{58}}%
\TextAt{13}{324.281}{514.517}{\Character{58}}%
\TextAt{13}{329.517}{514.517}{\Character{58}}%
\TextAt{13}{334.742}{514.517}{\Character{59}}%
\TextAt{13}{339.978}{514.517}{\Character{33}}%
\TextAt{21}{347.33}{512.724}{n}%
\TextAt{8}{357.724}{514.518}{of}%
\TextAt{16}{371.913}{514.518}{\Character{79}}%
\TextAt{8}{386.52}{514.518}{over}%
\TextAt{14}{412.425}{514.518}{o}%
\TextAt{8}{418.27}{514.518}{.}%
\TextAt{8}{429.296}{514.518}{Since}%
\TextAt{13}{460.713}{514.518}{N}%
\TextAt{21}{470.087}{512.658}{L}%
\TextAt{3}{475.846}{512.658}{\Character{106}}%
\TextAt{21}{478.199}{512.658}{K}%
\TextAt{9}{486.401}{514.517}{\Character{40}}%
\TextAt{13}{490.954}{514.518}{\Character{11}}%
\TextAt{9}{498.475}{514.517}{\Character{41}}%
\TextAt{8}{507.784}{514.518}{is}%
\TextAt{8}{520.408}{514.518}{by}%
\TextAt{8}{79.226}{497.472}{definition}%
\TextAt{8}{131.334}{497.472}{the}%
\TextAt{8}{151.569}{497.472}{determinant}%
\TextAt{8}{217.664}{497.472}{det}%
\TextAt{9}{233.923}{497.472}{\Character{40}}%
\TextAt{13}{238.475}{497.472}{T}%
\TextAt{21}{245.336}{495.678}{\Character{11}}%
\TextAt{9}{251.273}{497.472}{\Character{41}}%
\TextAt{8}{259.802}{497.472}{of}%
\TextAt{8}{273.21}{497.472}{the}%
\TextAt{8}{293.446}{497.472}{transformation}%
\TextAt{13}{373.277}{497.472}{T}%
\TextAt{21}{380.138}{495.678}{\Character{11}}%
\TextAt{9}{389.53}{497.472}{\Character{58}}%
\TextAt{13}{396.237}{497.472}{L}%
\NativeMathAt{407.656}{497.472}{11.955}{\rightarrow}%
\TextAt{13}{423.066}{497.472}{L}%
\TextAt{8}{431.031}{497.472}{,}%
\TextAt{13}{438.278}{497.472}{x}%
\NativeArrowAt{448.385}{497.472}{11.955}{11.955}{\mapsto}%
\TextAt{13}{463.795}{497.472}{\Character{11}x}%
\TextAt{8}{477.969}{497.472}{,}%
\TextAt{8}{485.217}{497.472}{we}%
\TextAt{8}{502.523}{497.472}{would}%
\TextAt{8}{79.226}{480.426}{get}%
\TextAt{8}{98.737}{480.426}{the}%
\TextAt{8}{118.893}{480.426}{equation}%
\TextAt{13}{197.806}{463.38}{\Character{11}\Character{33}}%
\TextAt{17}{212.617}{461.586}{1}%
\TextAt{16}{220.006}{463.38}{\Character{94}}%
\TextAt{16}{230.622}{463.38}{\Character{1}}%
\TextAt{16}{235.918}{463.38}{\Character{1}}%
\TextAt{16}{241.227}{463.38}{\Character{1}}%
\TextAt{16}{247.192}{463.38}{\Character{94}}%
\TextAt{13}{257.864}{463.38}{\Character{11}\Character{33}}%
\TextAt{21}{272.675}{461.586}{n}%
\TextAt{9}{281.632}{463.38}{\Character{61}}%
\TextAt{13}{294.057}{463.38}{N}%
\TextAt{21}{303.431}{461.52}{L}%
\TextAt{3}{309.19}{461.52}{\Character{106}}%
\TextAt{21}{311.543}{461.52}{K}%
\TextAt{9}{319.745}{463.38}{\Character{40}}%
\TextAt{13}{324.298}{463.38}{\Character{11}}%
\TextAt{9}{331.819}{463.38}{\Character{41}\Character{40}}%
\TextAt{13}{340.924}{463.38}{\Character{33}}%
\TextAt{17}{348.213}{461.586}{1}%
\TextAt{16}{355.602}{463.38}{\Character{94}}%
\TextAt{16}{366.218}{463.38}{\Character{1}}%
\TextAt{16}{371.514}{463.38}{\Character{1}}%
\TextAt{16}{376.822}{463.38}{\Character{1}}%
\TextAt{16}{382.788}{463.38}{\Character{94}}%
\TextAt{13}{393.46}{463.38}{\Character{33}}%
\TextAt{21}{400.749}{461.586}{n}%
\TextAt{9}{406.386}{463.38}{\Character{41}}%
\TextAt{13}{410.938}{463.38}{\Character{59}}%
\TextAt{8}{79.226}{440.476}{the}%
\TextAt{8}{99.277}{440.476}{left-hand}%
\TextAt{8}{148.917}{440.476}{side,}%
\TextAt{8}{175.553}{440.476}{resp.}%
\TextAt{8}{204.84}{440.476}{right-hand}%
\TextAt{8}{262.298}{440.476}{side,}%
\TextAt{8}{288.934}{440.476}{of}%
\TextAt{8}{302.159}{440.476}{which}%
\TextAt{8}{335.54}{440.476}{would,}%
\TextAt{8}{372.851}{440.476}{by}%
\TextAt{8}{389.005}{440.476}{(1.6),}%
\TextAt{8}{420.147}{440.476}{generate}%
\TextAt{8}{466.87}{440.476}{the}%
\TextAt{8}{486.921}{440.476}{left-hand}%
\TextAt{8}{79.226}{423.43}{side,}%
\TextAt{8}{105.851}{423.43}{resp.}%
\TextAt{8}{135.133}{423.43}{right-hand}%
\TextAt{8}{192.578}{423.43}{side,}%
\TextAt{8}{219.202}{423.43}{of}%
\TextAt{8}{232.413}{423.43}{(**).}%
\TextAt{8}{261.647}{423.43}{But}%
\TextAt{8}{284.769}{423.43}{we}%
\TextAt{8}{301.877}{423.43}{may}%
\TextAt{8}{327.127}{423.43}{always}%
\TextAt{8}{364.475}{423.43}{produce}%
\TextAt{8}{408.901}{423.43}{a}%
\TextAt{8}{418.537}{423.43}{principal}%
\TextAt{8}{467.195}{423.43}{ideal}%
\TextAt{8}{495.039}{423.43}{domain}%
\TextAt{8}{79.226}{406.384}{as}%
\TextAt{8}{94.085}{406.384}{desired}%
\TextAt{8}{134.301}{406.384}{by}%
\TextAt{8}{151.049}{406.384}{passing}%
\TextAt{8}{192.639}{406.384}{from}%
\TextAt{16}{220.768}{406.384}{\Character{79}}%
\TextAt{16}{234.768}{406.384}{\Character{106}}%
\TextAt{14}{242.246}{406.384}{o}%
\TextAt{8}{252.477}{406.384}{to}%
\TextAt{8}{267.276}{406.384}{the}%
\TextAt{8}{287.921}{406.384}{localization}%
\TextAt{16}{350.529}{406.384}{\Character{79}}%
\TextAt{27}{360.046}{404.59}{p}%
\TextAt{16}{368.687}{406.384}{\Character{106}}%
\TextAt{14}{376.161}{406.384}{o}%
\TextAt{27}{382.006}{404.59}{p}%
\TextAt{8}{390.879}{406.384}{for}%
\TextAt{8}{409.252}{406.384}{every}%
\TextAt{8}{440.633}{406.384}{prime}%
\TextAt{8}{474.285}{406.384}{ideal}%
\TextAt{14}{502.737}{406.384}{p}%
\TextAt{8}{513.107}{406.384}{of}%
\TextAt{14}{526.925}{406.384}{o}%
\TextAt{8}{79.226}{389.338}{(see}%
\TextAt{8}{102.694}{389.338}{chap.}%
\TextAt{8}{134.885}{389.338}{I,}%
\TextAt{8}{146.255}{389.338}{§11}%
\TextAt{8}{167.547}{389.338}{and}%
\TextAt{8}{190.31}{389.338}{§3,}%
\TextAt{8}{208.996}{389.338}{exercise}%
\TextAt{8}{252.297}{389.338}{4).}%
\TextAt{8}{271.158}{389.338}{The}%
\TextAt{8}{295.212}{389.338}{preceding}%
\TextAt{8}{347.887}{389.338}{argument}%
\TextAt{8}{400.25}{389.338}{then}%
\TextAt{8}{426.91}{389.338}{shows}%
\TextAt{8}{460.528}{389.338}{that}%
\TextAt{9}{205.659}{361.778}{\Character{40}}%
\TextAt{8}{210.211}{361.778}{det}%
\TextAt{27}{216.393}{354.486}{o}%
\TextAt{14}{228.463}{361.778}{A}%
\TextAt{21}{237.042}{359.985}{f}%
\TextAt{9}{242.487}{361.778}{\Character{41}}%
\TextAt{27}{247.039}{359.985}{p}%
\TextAt{9}{254.847}{361.778}{\Character{61}}%
\TextAt{8}{267.272}{361.778}{det}%
\TextAt{27}{271.412}{354.486}{o}%
\TextAt{28}{275.309}{353.489}{p}%
\TextAt{14}{285.524}{361.778}{A}%
\TextAt{21}{294.103}{359.985}{f}%
\TextAt{27}{299.05}{359.985}{p}%
\TextAt{9}{306.857}{361.778}{\Character{61}}%
\TextAt{13}{319.283}{361.778}{N}%
\TextAt{21}{328.657}{359.918}{L}%
\TextAt{3}{334.416}{359.918}{\Character{106}}%
\TextAt{21}{336.768}{359.918}{K}%
\TextAt{9}{344.971}{361.778}{\Character{40}}%
\TextAt{14}{349.523}{361.778}{A}%
\TextAt{21}{358.102}{359.985}{f}%
\TextAt{27}{363.049}{359.985}{p}%
\TextAt{9}{367.536}{361.778}{\Character{41}}%
\TextAt{8}{374.081}{361.778}{det}%
\TextAt{27}{378.221}{354.486}{o}%
\TextAt{28}{382.117}{353.489}{p}%
\TextAt{16}{392.332}{361.778}{\Character{79}}%
\TextAt{27}{401.85}{359.985}{p}%
\TextAt{9}{254.847}{335.246}{\Character{61}}%
\TextAt{9}{267.257}{335.246}{\Character{40}}%
\TextAt{13}{271.825}{335.246}{N}%
\TextAt{21}{281.198}{333.386}{L}%
\TextAt{3}{286.957}{333.386}{\Character{106}}%
\TextAt{21}{289.31}{333.386}{K}%
\TextAt{9}{297.512}{335.246}{\Character{40}}%
\TextAt{14}{302.065}{335.246}{A}%
\TextAt{21}{310.644}{333.452}{f}%
\TextAt{9}{316.089}{335.246}{\Character{41}}%
\TextAt{8}{322.634}{335.246}{det}%
\TextAt{27}{328.815}{327.953}{o}%
\TextAt{16}{340.885}{335.246}{\Character{79}}%
\TextAt{9}{350.735}{335.246}{\Character{41}}%
\TextAt{27}{355.287}{333.452}{p}%
\TextAt{13}{359.774}{335.246}{\Character{59}}%
\TextAt{8}{79.226}{306.281}{and}%
\TextAt{8}{101.989}{306.281}{since}%
\TextAt{8}{130.657}{306.281}{this}%
\TextAt{8}{153.48}{306.281}{identity}%
\TextAt{8}{196.734}{306.281}{is}%
\TextAt{8}{208.497}{306.281}{valid}%
\TextAt{8}{236.795}{306.281}{for}%
\TextAt{8}{254.68}{306.281}{all}%
\TextAt{8}{270.939}{306.281}{prime}%
\TextAt{8}{304.103}{306.281}{ideals}%
\TextAt{14}{336.681}{306.281}{p}%
\TextAt{8}{346.561}{306.281}{of}%
\TextAt{14}{359.891}{306.281}{o}%
\TextAt{8}{365.735}{306.281}{,}%
\TextAt{8}{372.885}{306.281}{we}%
\TextAt{8}{390.112}{306.281}{deduce}%
\TextAt{8}{429.122}{306.281}{the}%
\TextAt{8}{449.278}{306.281}{equality}%
\TextAt{8}{493.835}{306.281}{(**).}%
\TextAt{8}{97.159}{286.246}{In}%
\TextAt{8}{111.222}{286.246}{order}%
\TextAt{8}{141.234}{286.246}{to}%
\TextAt{8}{154.986}{286.246}{prove}%
\TextAt{8}{185.978}{286.246}{that}%
\TextAt{8}{210.788}{286.246}{the}%
\TextAt{8}{230.387}{286.246}{metrics}%
\TextAt{8}{270.859}{286.246}{agree}%
\TextAt{8}{300.871}{286.246}{on}%
\TextAt{8}{316.572}{286.246}{both}%
\TextAt{8}{343.654}{286.246}{sides}%
\TextAt{8}{371.179}{286.246}{of}%
\TextAt{8}{383.951}{286.246}{(**),}%
\TextAt{8}{411.479}{286.246}{we}%
\TextAt{8}{428.149}{286.246}{put}%
\TextAt{13}{449.05}{286.246}{M}%
\TextAt{9}{464.945}{286.246}{\Character{61}}%
\TextAt{13}{477.37}{286.246}{L}%
\TextAt{9}{485.335}{286.246}{\Character{40}}%
\TextAt{14}{489.887}{286.246}{A}%
\TextAt{9}{498.466}{286.246}{\Character{41}}%
\TextAt{8}{503.019}{286.246}{,}%
\TextAt{13}{509.721}{286.246}{N}%
\TextAt{9}{523.665}{286.246}{\Character{61}}%
\TextAt{13}{79.226}{269.2}{L}%
\TextAt{9}{87.191}{269.2}{\Character{40}}%
\TextAt{16}{91.743}{269.2}{\Character{79}}%
\TextAt{9}{101.593}{269.2}{\Character{41}}%
\TextAt{8}{106.145}{269.2}{,}%
\TextAt{14}{112.824}{269.2}{a}%
\TextAt{9}{122.128}{269.2}{\Character{61}}%
\TextAt{13}{134.554}{269.2}{N}%
\TextAt{21}{143.927}{267.341}{L}%
\TextAt{3}{149.686}{267.341}{\Character{106}}%
\TextAt{21}{152.039}{267.341}{K}%
\TextAt{9}{160.241}{269.2}{\Character{40}}%
\TextAt{14}{164.794}{269.2}{A}%
\TextAt{9}{173.373}{269.2}{\Character{41}}%
\TextAt{8}{181.235}{269.2}{and}%
\TextAt{8}{203.41}{269.2}{we}%
\TextAt{8}{220.05}{269.2}{view}%
\TextAt{13}{246.445}{269.2}{M}%
\TextAt{13}{258.364}{269.2}{\Character{59}}%
\TextAt{13}{263.589}{269.2}{N}%
\TextAt{13}{273.559}{269.2}{\Character{59}}%
\TextAt{14}{278.829}{269.2}{a}%
\TextAt{8}{288.122}{269.2}{as}%
\TextAt{8}{301.904}{269.2}{metrized}%
\TextAt{14}{349.436}{269.2}{o}%
\TextAt{8}{355.281}{269.2}{-modules.}%
\TextAt{8}{409.441}{269.2}{One}%
\TextAt{8}{433.565}{269.2}{has}%
\TextAt{13}{453.851}{269.2}{M}%
\TextAt{25}{465.176}{267.407}{C}%
\TextAt{9}{474.751}{269.2}{\Character{61}}%
\TextAt{13}{487.176}{269.2}{N}%
\TextAt{25}{496.55}{267.407}{C}%
\TextAt{9}{506.125}{269.2}{\Character{61}}%
\TextAt{13}{518.551}{269.2}{L}%
\TextAt{25}{526.515}{267.407}{C}%
\TextAt{8}{79.226}{252.154}{and}%
\TextAt{14}{101.989}{252.154}{a}%
\TextAt{25}{107.972}{250.361}{C}%
\TextAt{9}{117.547}{252.154}{\Character{61}}%
\TextAt{13}{129.973}{252.154}{K}%
\TextAt{25}{139.932}{250.361}{C}%
\TextAt{8}{146.186}{252.154}{,}%
\TextAt{8}{153.335}{252.154}{and}%
\TextAt{8}{176.098}{252.154}{we}%
\TextAt{8}{193.325}{252.154}{consider}%
\TextAt{8}{238.91}{252.154}{the}%
\TextAt{8}{259.067}{252.154}{metrics}%
\TextAt{8}{300.097}{252.154}{on}%
\TextAt{8}{316.356}{252.154}{the}%
\TextAt{8}{336.512}{252.154}{components}%
\TextAt{13}{196.812}{224.595}{M}%
\TextAt{21}{208.137}{222.802}{\Character{27}}%
\TextAt{9}{217.081}{224.595}{\Character{61}}%
\TextAt{20}{229.507}{235.952}{\Character{77}}%
\TextAt{21}{232.462}{210.049}{\Character{28}}%
\TextAt{3}{237.14}{210.049}{\Character{106}}%
\TextAt{21}{239.492}{210.049}{\Character{27}}%
\TextAt{18}{249.565}{224.595}{C}%
\TextAt{13}{258.199}{224.595}{\Character{59}}%
\TextAt{13}{287.334}{224.595}{N}%
\TextAt{21}{296.728}{222.802}{\Character{27}}%
\TextAt{9}{305.672}{224.595}{\Character{61}}%
\TextAt{20}{318.097}{235.952}{\Character{77}}%
\TextAt{21}{321.052}{210.049}{\Character{28}}%
\TextAt{3}{325.731}{210.049}{\Character{106}}%
\TextAt{21}{328.083}{210.049}{\Character{27}}%
\TextAt{18}{338.156}{224.595}{C}%
\TextAt{13}{346.79}{224.595}{\Character{59}}%
\TextAt{14}{375.944}{224.595}{a}%
\TextAt{21}{381.927}{222.802}{\Character{27}}%
\TextAt{9}{390.872}{224.595}{\Character{61}}%
\TextAt{18}{403.297}{224.595}{C}%
\TextAt{13}{411.932}{224.595}{\Character{59}}%
\TextAt{8}{79.226}{186.385}{where}%
\TextAt{13}{113.235}{186.385}{\Character{27}}%
\NativeMathAt{123.977}{186.385}{11.955}{\in}%
\TextAt{8}{135.608}{186.385}{Hom}%
\TextAt{9}{159.996}{186.385}{\Character{40}}%
\TextAt{13}{164.548}{186.385}{K}%
\TextAt{13}{174.71}{186.385}{\Character{59}}%
\TextAt{18}{179.96}{186.385}{C}%
\TextAt{9}{188.594}{186.385}{\Character{41}}%
\TextAt{8}{197.243}{186.385}{and}%
\TextAt{13}{220.205}{186.385}{\Character{28}}%
\NativeMathAt{230.284}{186.385}{11.955}{\in}%
\TextAt{8}{241.914}{186.385}{Hom}%
\TextAt{9}{266.303}{186.385}{\Character{40}}%
\TextAt{13}{270.855}{186.385}{L\Character{59}}%
\TextAt{18}{284.064}{186.385}{C}%
\TextAt{9}{292.698}{186.385}{\Character{41}}%
\TextAt{8}{301.347}{186.385}{is}%
\TextAt{8}{313.31}{186.385}{such}%
\TextAt{8}{339.906}{186.385}{that}%
\TextAt{13}{365.474}{186.385}{\Character{28}}%
\TextAt{16}{375.554}{186.385}{\Character{106}}%
\TextAt{13}{382.535}{186.385}{\Character{27}}%
\TextAt{8}{389.617}{186.385}{.}%
\TextAt{8}{398.663}{186.385}{We}%
\TextAt{8}{419.007}{186.385}{have}%
\TextAt{8}{446.2}{186.385}{to}%
\TextAt{8}{460.71}{186.385}{show}%
\TextAt{8}{489.912}{186.385}{that,}%
\TextAt{8}{518.782}{186.385}{for}%
\TextAt{13}{79.226}{169.339}{\Character{24}}%
\TextAt{13}{84.905}{169.339}{\Character{59}}%
\TextAt{13}{90.141}{169.339}{\Character{17}}%
\NativeMathAt{99.686}{169.339}{11.955}{\in}%
\TextAt{8}{110.977}{169.339}{det}%
\TextAt{25}{127.236}{167.546}{C}%
\TextAt{13}{135.483}{169.339}{M}%
\TextAt{21}{146.808}{167.546}{\Character{27}}%
\TextAt{8}{156.328}{169.339}{and}%
\TextAt{13}{179.091}{169.339}{a\Character{59}}%
\TextAt{13}{190.46}{169.339}{b}%
\NativeMathAt{198.778}{169.339}{11.955}{\in}%
\TextAt{18}{210.069}{169.339}{C}%
\TextAt{8}{218.703}{169.339}{,}%
\TextAt{8}{225.853}{169.339}{one}%
\TextAt{8}{247.312}{169.339}{has}%
\TextAt{8}{268.186}{169.339}{the}%
\TextAt{8}{288.342}{169.339}{identity}%
\TextAt{16}{220.838}{141.779}{\Character{104}}%
\TextAt{13}{225.487}{141.779}{a\Character{24}}%
\TextAt{13}{237.311}{141.779}{\Character{59}}%
\TextAt{13}{242.547}{141.779}{b\Character{17}}%
\TextAt{16}{253.749}{141.779}{\Character{105}}%
\TextAt{7}{258.398}{139.986}{det}%
\TextAt{21}{271.565}{139.986}{M}%
\TextAt{22}{279.698}{138.99}{\Character{27}}%
\TextAt{9}{288.513}{141.779}{\Character{61}}%
\TextAt{16}{300.938}{141.779}{\Character{104}}%
\TextAt{13}{305.587}{141.779}{a\Character{59}}%
\TextAt{13}{316.956}{141.779}{b}%
\TextAt{16}{321.954}{141.779}{\Character{105}}%
\TextAt{27}{326.603}{139.986}{a}%
\TextAt{22}{330.592}{138.99}{\Character{27}}%
\TextAt{16}{336.086}{141.779}{\Character{104}}%
\TextAt{13}{340.735}{141.779}{\Character{24}}%
\TextAt{13}{346.414}{141.779}{\Character{59}}%
\TextAt{13}{351.638}{141.779}{\Character{17}}%
\TextAt{16}{357.874}{141.779}{\Character{105}}%
\TextAt{7}{362.523}{139.986}{det}%
\TextAt{21}{375.691}{139.986}{N}%
\TextAt{22}{382.412}{138.99}{\Character{27}}%
\TextAt{13}{387.906}{141.779}{\Character{58}}%
\TextAt{8}{79.226}{114.22}{For}%
\TextAt{8}{100.196}{114.22}{this,}%
\TextAt{8}{126.27}{114.22}{let}%
\TextAt{14}{143.174}{114.22}{A}%
\NativeMathAt{151.753}{112.427}{7.97}{\infty}%
\TextAt{9}{164.041}{114.22}{\Character{61}}%
\TextAt{20}{176.466}{123.186}{\Character{81}}%
\TextAt{27}{187.757}{110.733}{P}%
\TextAt{3}{194.353}{110.733}{\Character{106}}%
\NativeMathAt{196.704}{110.733}{7.97}{\infty}%
\TextAt{14}{207.664}{114.22}{P}%
\TextAt{21}{217.557}{118.558}{\Character{23}}%
\TextAt{28}{221.72}{117.167}{P}%
\TextAt{8}{228.415}{114.22}{,}%
\TextAt{8}{235.564}{114.22}{so}%
\TextAt{8}{249.935}{114.22}{that}%
\TextAt{8}{275.303}{114.22}{one}%
\TextAt{8}{296.763}{114.22}{gets}%
\TextAt{14}{240.61}{84.677}{a}%
\NativeMathAt{246.593}{82.884}{7.97}{\infty}%
\TextAt{9}{258.88}{84.677}{\Character{61}}%
\TextAt{13}{271.306}{84.677}{N}%
\TextAt{21}{280.68}{82.818}{L}%
\TextAt{3}{286.439}{82.818}{\Character{106}}%
\TextAt{21}{288.791}{82.818}{K}%
\TextAt{9}{296.994}{84.677}{\Character{40}}%
\TextAt{14}{301.546}{84.677}{A}%
\NativeMathAt{310.125}{82.884}{7.97}{\infty}%
\TextAt{9}{319.092}{84.677}{\Character{41}}%
\TextAt{9}{326.947}{84.677}{\Character{61}}%
\TextAt{20}{339.39}{96.035}{\Character{89}}%
\TextAt{27}{339.624}{70.132}{p}%
\TextAt{3}{343.612}{70.132}{\Character{106}}%
\NativeMathAt{345.963}{70.132}{7.97}{\infty}%
\TextAt{14}{356.659}{84.677}{p}%
\TextAt{21}{362.642}{89.614}{\Character{23}}%
\TextAt{28}{366.805}{88.617}{p}%
\TextAt{8}{297.211}{43.033}{229}%
\Destination{516.5}{593.215}{nav.AMS.39}%
\Destination{78.23}{809.815}{nav.page.229}%
\Destination{79.36}{776.155}{nav.print.249}%
\Destination{79.36}{301.255}{nav.print.250}%
\Hotlink{134.385}{387.155}{5.22}{11.047}{nav.chapter.1}%
\Hotlink{151.434}{387.155}{12.716}{11.047}{nav.section.1.11}%
\Hotlink{195.489}{387.155}{6.858}{11.047}{nav.section.1.3}%
\Hotlink{104.049}{741.855}{25.078}{11.047}{nav.number.3.4.7}%
\Hotlink{289.331}{741.855}{25.078}{11.047}{nav.number.3.4.7}%
\Hotlink{379.214}{724.855}{25.078}{11.047}{nav.number.3.5.6}%
\Hotlink{388.505}{438.255}{25.078}{11.047}{nav.number.3.1.6}%
}

% Page 242 (body page 230)
\Page{%
\TextAt{2}{79.226}{796.943}{III.}%
\TextAt{2}{98.317}{796.943}{Riemann–Roch}%
\TextAt{2}{174.375}{796.943}{Theory}%
\Vector{q 0 G 0 g 0.299 w 7.226 18.729 m 460.769 18.729 l S Q}%
\TextAt{8}{79.226}{761.124}{with}%
\TextAt{13}{105.886}{761.124}{\Character{23}}%
\TextAt{27}{111.67}{759.33}{p}%
\TextAt{9}{119.478}{761.124}{\Character{61}}%
\TextAt{20}{131.903}{770.09}{\Character{80}}%
\TextAt{27}{144.523}{757.637}{P}%
\TextAt{3}{151.118}{757.637}{\Character{106}}%
\TextAt{27}{153.47}{757.637}{p}%
\TextAt{13}{159.95}{761.124}{f}%
\TextAt{27}{165.72}{759.264}{P}%
\TextAt{3}{172.315}{759.264}{\Character{106}}%
\TextAt{27}{174.667}{759.264}{p}%
\TextAt{13}{179.154}{761.124}{\Character{23}}%
\TextAt{27}{184.938}{759.33}{P}%
\TextAt{8}{192.032}{761.124}{.}%
\TextAt{8}{200.48}{761.124}{Then}%
\TextAt{16}{140.254}{733.305}{\Character{104}}%
\TextAt{13}{144.903}{733.305}{x\Character{59}}%
\TextAt{13}{156.775}{733.305}{y}%
\TextAt{16}{162.936}{733.305}{\Character{105}}%
\TextAt{21}{167.585}{731.512}{N}%
\TextAt{22}{174.307}{730.516}{\Character{27}}%
\TextAt{9}{183.121}{733.305}{\Character{61}}%
\TextAt{20}{195.547}{744.663}{\Character{88}}%
\TextAt{21}{198.103}{718.76}{\Character{28}}%
\TextAt{3}{202.782}{718.76}{\Character{106}}%
\TextAt{21}{205.134}{718.76}{\Character{27}}%
\TextAt{13}{214.808}{733.306}{x}%
\TextAt{21}{221.46}{731.512}{\Character{28}}%
\Vector{q 0 G 0 g 0.478 w 154.636 -33.873 m 160.773 -33.873 l S Q}%
\TextAt{13}{226.636}{733.305}{y}%
\TextAt{21}{232.773}{730.483}{\Character{28}}%
\TextAt{13}{237.95}{733.305}{\Character{59}}%
\TextAt{16}{302.229}{733.305}{\Character{104}}%
\TextAt{13}{306.878}{733.305}{x\Character{59}}%
\TextAt{13}{318.75}{733.305}{y}%
\TextAt{16}{324.911}{733.305}{\Character{105}}%
\TextAt{21}{329.56}{731.512}{M}%
\TextAt{22}{337.693}{730.516}{\Character{27}}%
\TextAt{9}{346.508}{733.305}{\Character{61}}%
\TextAt{20}{358.933}{744.663}{\Character{88}}%
\TextAt{21}{361.49}{718.76}{\Character{28}}%
\TextAt{3}{366.168}{718.76}{\Character{106}}%
\TextAt{21}{368.52}{718.76}{\Character{27}}%
\TextAt{13}{378.194}{733.306}{e}%
\TextAt{17}{383.62}{738.242}{2}%
\TextAt{21}{387.854}{738.242}{\Character{23}}%
\TextAt{28}{392.017}{736.851}{P}%
\TextAt{22}{397.716}{735.854}{\Character{28}}%
\TextAt{13}{403.381}{733.305}{x}%
\TextAt{21}{410.033}{731.512}{\Character{28}}%
\Vector{q 0 G 0 g 0.478 w 343.21 -33.873 m 349.347 -33.873 l S Q}%
\TextAt{13}{415.21}{733.305}{y}%
\TextAt{21}{421.347}{730.483}{\Character{28}}%
\TextAt{13}{426.523}{733.305}{\Character{59}}%
\TextAt{16}{144.653}{699.034}{\Character{104}}%
\TextAt{13}{149.302}{699.034}{a\Character{59}}%
\TextAt{13}{160.671}{699.034}{b}%
\TextAt{16}{165.669}{699.034}{\Character{105}}%
\TextAt{27}{170.318}{697.241}{a}%
\TextAt{22}{174.307}{696.244}{\Character{27}}%
\TextAt{9}{183.121}{699.034}{\Character{61}}%
\TextAt{13}{195.547}{699.034}{e}%
\TextAt{17}{200.972}{703.97}{2}%
\TextAt{21}{205.206}{703.97}{\Character{23}}%
\TextAt{28}{209.369}{702.974}{p}%
\TextAt{22}{212.954}{701.978}{\Character{27}}%
\TextAt{13}{218.946}{699.034}{a}%
\Vector{q 0 G 0 g 0.478 w 153.091 -64.99 m 158.068 -64.99 l S Q}%
\TextAt{13}{225.091}{699.034}{b\Character{59}}%
\TextAt{13}{327.921}{699.034}{\Character{23}}%
\TextAt{27}{333.704}{697.241}{p}%
\TextAt{22}{337.693}{696.244}{\Character{27}}%
\TextAt{9}{346.508}{699.034}{\Character{61}}%
\TextAt{20}{359.265}{710.391}{\Character{88}}%
\TextAt{27}{358.933}{684.488}{P}%
\TextAt{3}{365.529}{684.488}{\Character{106}}%
\TextAt{27}{367.881}{684.488}{p}%
\TextAt{22}{371.87}{683.492}{\Character{27}}%
\TextAt{13}{378.858}{699.034}{f}%
\TextAt{27}{384.628}{697.174}{P}%
\TextAt{3}{391.224}{697.174}{\Character{106}}%
\TextAt{27}{393.576}{697.174}{p}%
\TextAt{22}{397.565}{696.178}{\Character{27}}%
\TextAt{13}{403.059}{699.034}{\Character{23}}%
\TextAt{27}{408.843}{697.241}{P}%
\TextAt{9}{419.257}{699.034}{\Character{61}}%
\TextAt{20}{431.682}{710.391}{\Character{88}}%
\TextAt{21}{434.239}{684.488}{\Character{28}}%
\TextAt{3}{438.917}{684.488}{\Character{106}}%
\TextAt{21}{441.269}{684.488}{\Character{27}}%
\TextAt{13}{450.944}{699.034}{\Character{23}}%
\TextAt{27}{456.727}{697.241}{P}%
\TextAt{22}{463.323}{696.244}{\Character{28}}%
\TextAt{13}{468.49}{699.034}{\Character{58}}%
\TextAt{8}{79.226}{662.548}{Let}%
\TextAt{13}{100.126}{662.548}{\Character{24}}%
\TextAt{9}{109.138}{662.548}{\Character{61}}%
\TextAt{13}{121.563}{662.548}{x}%
\TextAt{17}{128.215}{660.755}{1}%
\TextAt{16}{135.463}{662.548}{\Character{94}}%
\TextAt{16}{145.936}{662.548}{\Character{1}}%
\TextAt{16}{151.232}{662.548}{\Character{1}}%
\TextAt{16}{156.54}{662.548}{\Character{1}}%
\TextAt{16}{162.374}{662.548}{\Character{94}}%
\TextAt{13}{172.896}{662.548}{x}%
\TextAt{21}{179.548}{660.755}{n}%
\TextAt{8}{185.184}{662.548}{,}%
\TextAt{13}{192.278}{662.548}{\Character{17}}%
\TextAt{9}{201.803}{662.548}{\Character{61}}%
\TextAt{13}{214.229}{662.548}{y}%
\TextAt{17}{219.936}{660.755}{1}%
\TextAt{16}{227.184}{662.548}{\Character{94}}%
\TextAt{16}{237.657}{662.548}{\Character{1}}%
\TextAt{16}{242.965}{662.548}{\Character{1}}%
\TextAt{16}{248.261}{662.548}{\Character{1}}%
\TextAt{16}{254.095}{662.548}{\Character{94}}%
\TextAt{13}{264.617}{662.548}{y}%
\TextAt{21}{270.325}{660.755}{n}%
\TextAt{8}{275.961}{662.548}{.}%
\TextAt{8}{284.386}{662.548}{We}%
\TextAt{8}{304.461}{662.548}{number}%
\TextAt{8}{346.988}{662.548}{the}%
\TextAt{8}{367.076}{662.548}{embeddings}%
\TextAt{13}{430.799}{662.548}{\Character{28}}%
\TextAt{16}{440.539}{662.548}{\Character{106}}%
\TextAt{13}{447.181}{662.548}{\Character{27}}%
\TextAt{8}{454.263}{662.548}{,}%
\TextAt{13}{461.357}{662.548}{\Character{28}}%
\TextAt{17}{466.441}{660.755}{1}%
\TextAt{13}{471.173}{662.548}{\Character{59}}%
\TextAt{13}{476.397}{662.548}{\Character{58}}%
\TextAt{13}{481.634}{662.548}{\Character{58}}%
\TextAt{13}{486.858}{662.548}{\Character{58}}%
\TextAt{13}{492.095}{662.548}{\Character{59}}%
\TextAt{13}{497.331}{662.548}{\Character{28}}%
\TextAt{21}{502.477}{660.755}{n}%
\TextAt{8}{508.114}{662.548}{,}%
\TextAt{8}{515.207}{662.548}{put}%
\TextAt{13}{79.226}{645.502}{\Character{23}}%
\TextAt{21}{85.01}{643.709}{k}%
\TextAt{9}{93.451}{645.502}{\Character{61}}%
\TextAt{13}{105.876}{645.502}{\Character{23}}%
\TextAt{27}{111.66}{643.709}{P}%
\TextAt{22}{118.255}{642.713}{\Character{28}}%
\TextAt{22}{121.735}{640.621}{k}%
\TextAt{8}{131.201}{645.502}{and}%
\TextAt{8}{153.964}{645.502}{form}%
\TextAt{8}{181.604}{645.502}{the}%
\TextAt{8}{201.761}{645.502}{matrices}%
\TextAt{13}{167.227}{600.69}{A}%
\TextAt{9}{179.323}{600.69}{\Character{61}}%
\TextAt{9}{191.733}{600.69}{\Character{40}}%
\TextAt{13}{196.301}{600.69}{x}%
\TextAt{21}{202.953}{598.896}{i\Character{28}}%
\TextAt{22}{209.595}{597.49}{k}%
\TextAt{9}{214.666}{600.69}{\Character{41}}%
\TextAt{13}{219.218}{600.69}{\Character{59}}%
\TextAt{13}{248.353}{600.69}{B}%
\TextAt{9}{261.19}{600.69}{\Character{61}}%
\TextAt{9}{273.6}{600.69}{\Character{40}}%
\Vector{q 0 G 0 g 0.478 w 206.168 -166.489 m 212.304 -166.489 l S Q}%
\TextAt{13}{278.168}{600.69}{y}%
\TextAt{21}{284.304}{597.867}{i\Character{28}}%
\TextAt{22}{290.946}{596.461}{k}%
\TextAt{9}{296.017}{600.69}{\Character{41}}%
\TextAt{13}{300.569}{600.69}{\Character{59}}%
\TextAt{13}{329.704}{600.69}{D}%
\TextAt{9}{343.089}{600.69}{\Character{61}}%
\TextAt{20}{355.514}{628.306}{\Character{48}}%
\TextAt{20}{355.514}{607.265}{\Character{66}}%
\TextAt{20}{355.514}{599.614}{\Character{64}}%
\TextAt{13}{365.975}{618.903}{e}%
\TextAt{21}{371.401}{623.241}{\Character{23}}%
\TextAt{24}{375.563}{622.134}{1}%
\TextAt{9}{418.544}{618.903}{0}%
\TextAt{8}{390.839}{605.655}{.}%
\TextAt{8}{395.419}{602.667}{.}%
\TextAt{8}{400}{599.678}{.}%
\TextAt{9}{370.167}{581.636}{0}%
\TextAt{13}{413.878}{581.636}{e}%
\TextAt{21}{419.304}{585.974}{\Character{23}}%
\TextAt{22}{423.466}{584.978}{n}%
\TextAt{20}{429.063}{628.306}{\Character{49}}%
\TextAt{20}{429.063}{607.265}{\Character{67}}%
\TextAt{20}{429.063}{599.614}{\Character{65}}%
\TextAt{13}{441.517}{600.69}{\Character{58}}%
\TextAt{8}{79.226}{558.439}{Then,}%
\TextAt{8}{113.035}{558.439}{observing}%
\TextAt{8}{165.459}{558.439}{that}%
\TextAt{8}{205.694}{532.604}{det}%
\TextAt{9}{221.953}{532.604}{\Character{40}}%
\TextAt{13}{226.505}{532.604}{D}%
\TextAt{9}{236.55}{532.604}{\Character{41}}%
\TextAt{9}{244.405}{532.604}{\Character{61}}%
\TextAt{20}{256.848}{543.961}{\Character{89}}%
\TextAt{21}{262.147}{518.855}{\Character{28}}%
\TextAt{13}{274.117}{532.604}{e}%
\TextAt{21}{279.542}{537.54}{\Character{23}}%
\TextAt{28}{283.705}{536.149}{P}%
\TextAt{22}{289.404}{535.153}{\Character{28}}%
\TextAt{9}{298.39}{532.604}{\Character{61}}%
\TextAt{20}{312.144}{543.961}{\Character{89}}%
\TextAt{27}{310.816}{518.058}{P}%
\TextAt{3}{317.411}{518.058}{\Character{106}}%
\TextAt{27}{319.764}{518.058}{p}%
\TextAt{22}{323.752}{517.062}{\Character{27}}%
\TextAt{13}{330.741}{532.604}{e}%
\TextAt{21}{336.166}{537.54}{f}%
\TextAt{28}{340.234}{535.802}{P}%
\TextAt{23}{345.933}{535.802}{\Character{106}}%
\TextAt{28}{348.203}{535.802}{p}%
\TextAt{22}{351.787}{534.806}{\Character{27}}%
\TextAt{21}{357.281}{537.54}{\Character{23}}%
\TextAt{28}{361.444}{536.149}{P}%
\TextAt{9}{371.46}{532.604}{\Character{61}}%
\TextAt{13}{383.885}{532.604}{e}%
\TextAt{21}{389.311}{537.54}{\Character{23}}%
\TextAt{28}{393.474}{536.544}{p}%
\TextAt{22}{397.058}{535.548}{\Character{27}}%
\TextAt{13}{403.05}{532.604}{\Character{59}}%
\TextAt{8}{79.226}{496.788}{we}%
\TextAt{8}{96.454}{496.788}{do}%
\TextAt{8}{112.713}{496.788}{indeed}%
\TextAt{8}{149.774}{496.788}{get}%
\TextAt{16}{156.626}{470.952}{\Character{104}}%
\TextAt{13}{161.275}{470.952}{a\Character{24}}%
\TextAt{13}{173.099}{470.952}{\Character{59}}%
\TextAt{13}{178.335}{470.952}{b\Character{17}}%
\TextAt{16}{189.536}{470.952}{\Character{105}}%
\TextAt{7}{194.186}{469.159}{det}%
\TextAt{21}{207.353}{469.159}{M}%
\TextAt{22}{215.486}{468.163}{\Character{27}}%
\TextAt{9}{224.3}{470.952}{\Character{61}}%
\TextAt{13}{236.726}{470.952}{a}%
\Vector{q 0 G 0 g 0.478 w 170.871 -293.072 m 175.848 -293.072 l S Q}%
\TextAt{13}{242.871}{470.952}{b}%
\TextAt{16}{247.848}{470.952}{\Character{104}}%
\TextAt{13}{252.497}{470.952}{\Character{24}}%
\TextAt{13}{258.176}{470.952}{\Character{59}}%
\TextAt{13}{263.412}{470.952}{\Character{17}}%
\TextAt{16}{269.636}{470.952}{\Character{105}}%
\TextAt{7}{274.285}{469.159}{det}%
\TextAt{21}{287.453}{469.159}{M}%
\TextAt{22}{295.586}{468.163}{\Character{27}}%
\TextAt{9}{224.3}{450.918}{\Character{61}}%
\TextAt{13}{236.726}{450.918}{a}%
\Vector{q 0 G 0 g 0.478 w 170.871 -313.107 m 175.848 -313.107 l S Q}%
\TextAt{13}{242.871}{450.918}{b}%
\TextAt{8}{249.84}{450.918}{det}%
\TextAt{20}{266.099}{460.601}{\Character{0}}%
\TextAt{9}{271.579}{450.918}{\Character{40}}%
\TextAt{13}{276.131}{450.918}{AD}%
\TextAt{9}{294.951}{450.918}{\Character{41}\Character{40}}%
\TextAt{13}{304.056}{450.918}{B}%
\TextAt{13}{313.536}{450.918}{D}%
\TextAt{9}{323.596}{450.918}{\Character{41}}%
\TextAt{21}{328.149}{455.854}{t}%
\TextAt{20}{331.705}{460.601}{\Character{1}}%
\TextAt{9}{340.505}{450.918}{\Character{61}}%
\TextAt{13}{352.931}{450.918}{a}%
\Vector{q 0 G 0 g 0.478 w 287.076 -313.107 m 292.053 -313.107 l S Q}%
\TextAt{13}{359.076}{450.918}{b}%
\TextAt{9}{364.053}{450.918}{\Character{40}}%
\TextAt{8}{368.605}{450.918}{det}%
\TextAt{13}{386.857}{450.918}{D}%
\TextAt{9}{396.901}{450.918}{\Character{41}}%
\TextAt{17}{401.453}{455.854}{2}%
\TextAt{8}{408.178}{450.918}{det}%
\TextAt{9}{424.437}{450.918}{\Character{40}}%
\TextAt{13}{428.99}{450.918}{AB}%
\TextAt{21}{447.261}{455.854}{t}%
\TextAt{9}{450.817}{450.918}{\Character{41}}%
\TextAt{9}{224.3}{430.883}{\Character{61}}%
\TextAt{13}{236.726}{430.883}{e}%
\TextAt{17}{242.151}{435.819}{2}%
\TextAt{21}{246.386}{435.819}{\Character{23}}%
\TextAt{28}{250.548}{434.823}{p}%
\TextAt{22}{254.133}{433.826}{\Character{27}}%
\TextAt{13}{260.125}{430.883}{a}%
\Vector{q 0 G 0 g 0.478 w 194.27 -333.141 m 199.247 -333.141 l S Q}%
\TextAt{13}{266.27}{430.883}{b}%
\TextAt{16}{271.247}{430.883}{\Character{104}}%
\TextAt{13}{275.896}{430.883}{\Character{24}}%
\TextAt{13}{281.575}{430.883}{\Character{59}}%
\TextAt{13}{286.811}{430.883}{\Character{17}}%
\TextAt{16}{293.035}{430.883}{\Character{105}}%
\TextAt{7}{297.684}{429.09}{det}%
\TextAt{21}{310.852}{429.09}{N}%
\TextAt{22}{317.573}{428.093}{\Character{27}}%
\TextAt{9}{326.388}{430.883}{\Character{61}}%
\TextAt{16}{338.814}{430.883}{\Character{104}}%
\TextAt{13}{343.463}{430.883}{a\Character{59}}%
\TextAt{13}{354.832}{430.883}{b}%
\TextAt{16}{359.829}{430.883}{\Character{105}}%
\TextAt{27}{364.478}{429.09}{a}%
\TextAt{22}{368.467}{428.093}{\Character{27}}%
\TextAt{16}{373.961}{430.883}{\Character{104}}%
\TextAt{13}{378.61}{430.883}{\Character{24}}%
\TextAt{13}{384.289}{430.883}{\Character{59}}%
\TextAt{13}{389.525}{430.883}{\Character{17}}%
\TextAt{16}{395.749}{430.883}{\Character{105}}%
\TextAt{7}{400.398}{429.09}{det}%
\TextAt{21}{413.566}{429.09}{N}%
\TextAt{22}{420.287}{428.093}{\Character{27}}%
\TextAt{13}{425.781}{430.883}{\Character{58}}%
\TextAt{8}{79.226}{405.047}{This}%
\TextAt{8}{105.946}{405.047}{proves}%
\TextAt{8}{142.11}{405.047}{our}%
\TextAt{8}{162.924}{405.047}{theorem.}%
\Vector{q 0 G 0 g 0.399 w 452.334 -368.953 m 452.334 -360.883 l S Q}%
\Vector{q 0 G 0 g 0.399 w 452.534 -361.082 m 459.707 -361.082 l S Q}%
\Vector{q 0 G 0 g 0.399 w 452.534 -368.753 m 459.707 -368.753 l S Q}%
\Vector{q 0 G 0 g 0.399 w 459.906 -368.953 m 459.906 -360.883 l S Q}%
\TextAt{8}{97.159}{376.951}{Extending}%
\TextAt{8}{153.575}{376.951}{the}%
\TextAt{8}{173.732}{376.951}{formulas}%
\TextAt{8}{221.277}{376.951}{of}%
\TextAt{8}{234.607}{376.951}{(7.2)}%
\TextAt{8}{262.583}{376.951}{to}%
\TextAt{8}{276.893}{376.951}{the}%
\TextAt{8}{297.049}{376.951}{free}%
\TextAt{8}{319.477}{376.951}{abelian}%
\TextAt{8}{360.125}{376.951}{group}%
\TextAt{13}{256.214}{351.115}{F}%
\TextAt{17}{263.792}{349.322}{0}%
\TextAt{9}{268.524}{351.115}{\Character{40}}%
\Vector{q 0 G 0 g 0.478 w 201.076 -413.042 m 210.926 -413.042 l S Q}%
\TextAt{16}{273.076}{351.115}{\Character{79}}%
\TextAt{9}{282.926}{351.115}{\Character{41}}%
\TextAt{9}{290.781}{351.115}{\Character{61}}%
\TextAt{20}{303.225}{362.473}{\Character{77}}%
\TextAt{3}{303.532}{336.57}{\Character{102}}%
\TextAt{21}{307.767}{336.57}{M}%
\TextAt{3}{316.748}{336.57}{\Character{103}}%
\TextAt{18}{323.283}{351.115}{Z}%
\TextAt{16}{331.253}{351.115}{\Character{102}}%
\TextAt{13}{337.231}{351.115}{M}%
\TextAt{16}{349.804}{351.115}{\Character{103}}%
\TextAt{8}{79.226}{313.036}{by}%
\TextAt{8}{94.517}{313.036}{linearity,}%
\TextAt{8}{142.197}{313.036}{and}%
\TextAt{8}{163.991}{313.036}{passing}%
\TextAt{8}{204.124}{313.036}{to}%
\TextAt{8}{217.466}{313.036}{the}%
\TextAt{8}{236.654}{313.036}{quotient}%
\TextAt{8}{281.869}{313.036}{group}%
\TextAt{13}{314.076}{313.036}{K}%
\TextAt{17}{324.035}{311.243}{0}%
\TextAt{9}{328.768}{313.036}{\Character{40}}%
\Vector{q 0 G 0 g 0.478 w 261.32 -451.121 m 271.17 -451.121 l S Q}%
\TextAt{16}{333.32}{313.036}{\Character{79}}%
\TextAt{9}{343.17}{313.036}{\Character{41}}%
\TextAt{9}{351.025}{313.036}{\Character{61}}%
\TextAt{13}{363.468}{313.036}{F}%
\TextAt{17}{371.046}{311.243}{0}%
\TextAt{9}{375.779}{313.036}{\Character{40}}%
\Vector{q 0 G 0 g 0.478 w 308.331 -451.121 m 318.181 -451.121 l S Q}%
\TextAt{16}{380.331}{313.036}{\Character{79}}%
\TextAt{9}{390.181}{313.036}{\Character{41}\Character{47}}%
\TextAt{13}{400.586}{313.036}{R}%
\TextAt{17}{409.496}{311.243}{0}%
\TextAt{9}{414.228}{313.036}{\Character{40}}%
\Vector{q 0 G 0 g 0.478 w 346.781 -451.121 m 356.63 -451.121 l S Q}%
\TextAt{16}{418.781}{313.036}{\Character{79}}%
\TextAt{9}{428.63}{313.036}{\Character{41}}%
\TextAt{8}{436.112}{313.036}{yields}%
\TextAt{8}{468.044}{313.036}{the}%
\TextAt{8}{487.232}{313.036}{following}%
\TextAt{8}{79.226}{295.99}{corollary.}%
\TextAt{12}{79.226}{270.882}{Corollary}%
\TextAt{12}{139.014}{270.882}{(7.3).}%
\TextAt{8}{174.33}{270.882}{For}%
\TextAt{8}{195.299}{270.882}{every}%
\TextAt{8}{226.191}{270.882}{class}%
\TextAt{13}{253.628}{270.882}{\Character{24}}%
\NativeMathAt{262.639}{270.882}{11.955}{\in}%
\TextAt{13}{273.93}{270.882}{K}%
\TextAt{17}{283.89}{269.089}{0}%
\TextAt{9}{288.622}{270.882}{\Character{40}}%
\Vector{q 0 G 0 g 0.478 w 221.174 -493.275 m 231.024 -493.275 l S Q}%
\TextAt{16}{293.174}{270.882}{\Character{79}}%
\TextAt{9}{303.024}{270.882}{\Character{41}}%
\TextAt{8}{307.576}{270.882}{,}%
\TextAt{8}{314.725}{270.882}{one}%
\TextAt{8}{336.185}{270.882}{has}%
\TextAt{8}{357.059}{270.882}{the}%
\TextAt{8}{377.215}{270.882}{formulas}%
\TextAt{8}{248.019}{245.047}{rk}%
\TextAt{9}{258.755}{245.047}{\Character{40}}%
\TextAt{13}{263.308}{245.047}{i}%
\TextAt{3}{267.301}{243.253}{\Character{3}}%
\TextAt{13}{272.033}{245.047}{\Character{24}}%
\TextAt{9}{277.724}{245.047}{\Character{41}}%
\TextAt{9}{285.579}{245.047}{\Character{61}}%
\TextAt{9}{298}{245.047}{\Character{91}}%
\TextAt{13}{301.274}{245.047}{L}%
\TextAt{9}{312.559}{245.047}{\Character{58}}%
\TextAt{13}{319.132}{245.047}{K}%
\TextAt{9}{329.95}{245.047}{\Character{93}}%
\TextAt{8}{335.194}{245.047}{rk}%
\TextAt{9}{345.929}{245.047}{\Character{40}}%
\TextAt{13}{350.482}{245.047}{\Character{24}}%
\TextAt{9}{356.172}{245.047}{\Character{41}}%
\TextAt{13}{360.725}{245.047}{\Character{59}}%
\TextAt{8}{216.132}{219.211}{det}%
\TextAt{9}{232.391}{219.211}{\Character{40}}%
\TextAt{13}{236.943}{219.211}{i}%
\TextAt{3}{240.937}{217.418}{\Character{3}}%
\TextAt{13}{245.669}{219.211}{\Character{24}}%
\TextAt{9}{251.36}{219.211}{\Character{41}}%
\TextAt{9}{259.215}{219.211}{\Character{61}}%
\TextAt{9}{271.636}{219.211}{\Character{91}}%
\TextAt{8}{274.91}{219.211}{det}%
\TextAt{13}{293.161}{219.211}{i}%
\TextAt{3}{297.155}{217.418}{\Character{3}}%
\TextAt{16}{301.887}{219.211}{\Character{79}}%
\TextAt{9}{311.737}{219.211}{\Character{93}}%
\TextAt{7}{314.989}{224.148}{rk}%
\TextAt{17}{322.759}{224.148}{\Character{40}}%
\TextAt{21}{326.053}{224.148}{\Character{24}}%
\TextAt{17}{330.086}{224.148}{\Character{41}}%
\TextAt{13}{333.878}{219.211}{N}%
\TextAt{21}{343.251}{217.352}{L}%
\TextAt{3}{349.01}{217.352}{\Character{106}}%
\TextAt{21}{351.363}{217.352}{K}%
\TextAt{9}{359.565}{219.211}{\Character{40}}%
\TextAt{8}{364.118}{219.211}{det}%
\TextAt{13}{382.369}{219.211}{\Character{24}}%
\TextAt{9}{388.06}{219.211}{\Character{41}}%
\TextAt{13}{392.612}{219.211}{\Character{58}}%
\TextAt{8}{97.159}{188.985}{The}%
\TextAt{8}{121.141}{188.985}{square}%
\TextAt{8}{157.879}{188.985}{of}%
\TextAt{8}{171.137}{188.985}{the}%
\TextAt{8}{191.221}{188.985}{metrized}%
\TextAt{14}{239.269}{188.985}{o}%
\TextAt{8}{245.113}{188.985}{-module}%
\TextAt{8}{290.232}{188.985}{det}%
\TextAt{13}{308.483}{188.985}{i}%
\TextAt{3}{312.476}{187.191}{\Character{3}}%
\TextAt{16}{317.209}{188.985}{\Character{79}}%
\TextAt{8}{330.884}{188.985}{appearing}%
\TextAt{8}{385.124}{188.985}{in}%
\TextAt{8}{398.705}{188.985}{the}%
\TextAt{8}{418.79}{188.985}{second}%
\TextAt{8}{456.496}{188.985}{formula}%
\TextAt{8}{499.355}{188.985}{can}%
\TextAt{8}{520.743}{188.985}{be}%
\TextAt{8}{79.226}{171.939}{computed}%
\TextAt{8}{132.326}{171.939}{to}%
\TextAt{8}{145.759}{171.939}{be}%
\TextAt{8}{160.805}{171.939}{the}%
\TextAt{19}{180.084}{171.939}{discriminant}%
\TextAt{14}{248.008}{171.939}{d}%
\TextAt{21}{253.969}{170.079}{L}%
\TextAt{3}{259.728}{170.079}{\Character{106}}%
\TextAt{21}{262.08}{170.079}{K}%
\TextAt{8}{273.303}{171.939}{of}%
\TextAt{8}{285.755}{171.939}{the}%
\TextAt{8}{305.034}{171.939}{extension}%
\TextAt{13}{355.922}{171.939}{L}%
\TextAt{16}{367.207}{171.939}{\Character{106}}%
\TextAt{13}{373.849}{171.939}{K}%
\TextAt{8}{384.667}{171.939}{,}%
\TextAt{8}{391.114}{171.939}{which}%
\TextAt{8}{423.722}{171.939}{we}%
\TextAt{8}{440.072}{171.939}{view}%
\TextAt{8}{466.177}{171.939}{as}%
\TextAt{8}{479.67}{171.939}{a}%
\TextAt{8}{488.547}{171.939}{metrized}%
\TextAt{14}{79.226}{154.893}{o}%
\TextAt{8}{85.07}{154.893}{-module}%
\TextAt{8}{130.261}{154.893}{with}%
\TextAt{8}{156.921}{154.893}{the}%
\TextAt{19}{177.077}{154.893}{trivial}%
\TextAt{8}{211.927}{154.893}{metric.}%
\TextAt{12}{79.226}{129.785}{Proposition}%
\TextAt{12}{151.937}{129.785}{(7.4).}%
\TextAt{8}{187.253}{129.785}{There}%
\TextAt{8}{221.062}{129.785}{is}%
\TextAt{8}{232.826}{129.785}{a}%
\TextAt{8}{242.582}{129.785}{canonical}%
\TextAt{8}{293.965}{129.785}{isomorphism}%
\TextAt{9}{258.265}{103.949}{\Character{40}}%
\TextAt{8}{262.818}{103.949}{det}%
\TextAt{13}{281.069}{103.949}{i}%
\TextAt{3}{285.063}{102.156}{\Character{3}}%
\TextAt{16}{289.795}{103.949}{\Character{79}}%
\TextAt{9}{299.645}{103.949}{\Character{41}}%
\TextAt{3}{304.197}{108.886}{\Character{10}}%
\TextAt{17}{310.783}{108.886}{2}%
\NativeMathAt{318.837}{103.949}{11.955}{\cong}%
\TextAt{14}{331.456}{103.949}{d}%
\TextAt{21}{337.416}{102.09}{L}%
\TextAt{3}{343.175}{102.09}{\Character{106}}%
\TextAt{21}{345.528}{102.09}{K}%
\TextAt{8}{79.226}{78.114}{of}%
\TextAt{8}{92.556}{78.114}{metrized}%
\TextAt{14}{140.676}{78.114}{o}%
\TextAt{8}{146.52}{78.114}{-modules.}%
\TextAt{8}{297.211}{43.033}{230}%
\Destination{78.23}{809.815}{nav.page.230}%
\Destination{79.36}{203.955}{nav.print.251}%
\Destination{79.36}{284.855}{nav.number.3.7.3}%
\Destination{79.36}{143.755}{nav.number.3.7.4}%
\Hotlink{234.107}{374.755}{25.078}{11.047}{nav.number.3.7.2}%
}

% Page 243 (body page 231)
\Page{%
\TextAt{2}{79.226}{796.943}{III.}%
\TextAt{2}{98.317}{796.943}{Riemann–Roch}%
\TextAt{2}{174.375}{796.943}{Theory}%
\Vector{q 0 G 0 g 0.299 w 7.226 18.729 m 460.769 18.729 l S Q}%
\TextAt{19}{79.226}{761.124}{Proof.}%
\TextAt{8}{116}{761.124}{Consider}%
\TextAt{8}{164.837}{761.124}{on}%
\TextAt{16}{181.096}{761.124}{\Character{79}}%
\TextAt{8}{194.843}{761.124}{the}%
\TextAt{8}{215}{761.124}{bilinear}%
\TextAt{8}{257.273}{761.124}{trace}%
\TextAt{8}{286.539}{761.124}{map}%
\TextAt{13}{195.448}{733.669}{T}%
\TextAt{9}{207.256}{733.669}{\Character{58}}%
\TextAt{16}{213.828}{733.669}{\Character{79}}%
\TextAt{16}{226.333}{733.669}{\Character{2}}%
\TextAt{16}{238.276}{733.669}{\Character{79}}%
\NativeArrowAt{251.439}{733.669}{11.955}{19.272}{\longrightarrow}%
\TextAt{14}{274.042}{733.669}{o}%
\TextAt{13}{279.887}{733.669}{\Character{59}}%
\TextAt{9}{309.041}{733.669}{\Character{40}}%
\TextAt{13}{313.593}{733.669}{x\Character{59}}%
\TextAt{13}{325.464}{733.669}{y}%
\TextAt{9}{331.626}{733.669}{\Character{41}}%
\NativeArrowAt{339.499}{733.669}{11.955}{19.26}{\longmapsto}%
\TextAt{8}{362.081}{733.669}{Tr}%
\TextAt{21}{375.089}{731.809}{L}%
\TextAt{3}{380.848}{731.809}{\Character{106}}%
\TextAt{21}{383.2}{731.809}{K}%
\TextAt{9}{391.403}{733.669}{\Character{40}}%
\TextAt{13}{395.955}{733.669}{xy}%
\TextAt{9}{408.744}{733.669}{\Character{41}}%
\TextAt{13}{413.296}{733.669}{\Character{58}}%
\TextAt{8}{79.226}{706.214}{It}%
\TextAt{8}{91.899}{706.214}{induces}%
\TextAt{8}{133.574}{706.214}{an}%
\TextAt{14}{149.833}{706.214}{o}%
\TextAt{8}{155.678}{706.214}{-module}%
\TextAt{8}{200.868}{706.214}{homomorphism}%
\TextAt{13}{243.901}{678.76}{T}%
\TextAt{9}{255.709}{678.76}{\Character{58}}%
\TextAt{8}{262.281}{678.76}{det}%
\TextAt{16}{280.533}{678.76}{\Character{79}}%
\TextAt{16}{293.026}{678.76}{\Character{10}}%
\TextAt{8}{304.994}{678.76}{det}%
\TextAt{16}{323.246}{678.76}{\Character{79}}%
\NativeArrowAt{336.409}{678.76}{11.955}{19.26}{\longrightarrow}%
\TextAt{14}{358.999}{678.76}{o}%
\TextAt{13}{364.843}{678.76}{\Character{59}}%
\TextAt{8}{79.226}{651.305}{given}%
\TextAt{8}{109.796}{651.305}{by}%
\TextAt{13}{162.624}{634.259}{T}%
\TextAt{20}{171.111}{643.943}{\Character{0}}%
\TextAt{9}{176.591}{634.259}{\Character{40}}%
\TextAt{13}{181.143}{634.259}{\Character{11}}%
\TextAt{17}{188.589}{632.466}{1}%
\TextAt{16}{195.978}{634.259}{\Character{94}}%
\TextAt{16}{206.594}{634.259}{\Character{1}}%
\TextAt{16}{211.89}{634.259}{\Character{1}}%
\TextAt{16}{217.199}{634.259}{\Character{1}}%
\TextAt{16}{223.164}{634.259}{\Character{94}}%
\TextAt{13}{233.836}{634.259}{\Character{11}}%
\TextAt{21}{241.281}{632.466}{n}%
\TextAt{9}{246.918}{634.259}{\Character{41}}%
\TextAt{16}{254.127}{634.259}{\Character{10}}%
\TextAt{9}{266.082}{634.259}{\Character{40}}%
\TextAt{13}{270.634}{634.259}{\Character{12}}%
\TextAt{17}{277.243}{632.466}{1}%
\TextAt{16}{284.632}{634.259}{\Character{94}}%
\TextAt{16}{295.248}{634.259}{\Character{1}}%
\TextAt{16}{300.544}{634.259}{\Character{1}}%
\TextAt{16}{305.853}{634.259}{\Character{1}}%
\TextAt{16}{311.818}{634.259}{\Character{94}}%
\TextAt{13}{322.489}{634.259}{\Character{12}}%
\TextAt{21}{329.098}{632.466}{n}%
\TextAt{9}{334.734}{634.259}{\Character{41}}%
\TextAt{20}{339.287}{643.943}{\Character{1}}%
\TextAt{9}{348.087}{634.259}{\Character{61}}%
\TextAt{8}{360.512}{634.259}{det}%
\TextAt{9}{376.771}{634.259}{\Character{40}}%
\TextAt{8}{381.324}{634.259}{Tr}%
\TextAt{21}{394.331}{632.399}{L}%
\TextAt{3}{400.09}{632.399}{\Character{106}}%
\TextAt{21}{402.442}{632.399}{K}%
\TextAt{9}{410.645}{634.259}{\Character{40}}%
\TextAt{13}{415.197}{634.259}{\Character{11}}%
\TextAt{21}{422.643}{632.466}{i}%
\TextAt{13}{426.024}{634.259}{\Character{12}}%
\TextAt{21}{432.633}{632.466}{j}%
\TextAt{9}{437.015}{634.259}{\Character{41}\Character{41}}%
\TextAt{13}{446.12}{634.259}{\Character{58}}%
\TextAt{8}{79.226}{611.4}{The}%
\TextAt{8}{103.748}{611.4}{image}%
\TextAt{8}{138.038}{611.4}{of}%
\TextAt{13}{151.836}{611.4}{T}%
\TextAt{8}{164.689}{611.4}{is}%
\TextAt{8}{176.921}{611.4}{the}%
\TextAt{8}{197.545}{611.4}{discriminant}%
\TextAt{8}{265.393}{611.4}{ideal}%
\TextAt{14}{293.825}{611.4}{d}%
\TextAt{21}{299.785}{609.54}{L}%
\TextAt{3}{305.544}{609.54}{\Character{106}}%
\TextAt{21}{307.896}{609.54}{K}%
\TextAt{8}{316.099}{611.4}{,}%
\TextAt{8}{323.834}{611.4}{which,}%
\TextAt{8}{361.157}{611.4}{by}%
\TextAt{8}{377.885}{611.4}{definition,}%
\TextAt{8}{433.751}{611.4}{is}%
\TextAt{8}{445.983}{611.4}{generated}%
\TextAt{8}{499.783}{611.4}{by}%
\TextAt{8}{516.51}{611.4}{the}%
\TextAt{8}{79.226}{594.354}{discriminants}%
\TextAt{13}{221.881}{577.308}{d}%
\TextAt{9}{227.964}{577.308}{\Character{40}}%
\TextAt{13}{232.516}{577.308}{\Character{33}}%
\TextAt{17}{239.806}{575.515}{1}%
\TextAt{13}{244.538}{577.308}{\Character{59}}%
\TextAt{13}{249.762}{577.308}{\Character{58}}%
\TextAt{13}{254.999}{577.308}{\Character{58}}%
\TextAt{13}{260.223}{577.308}{\Character{58}}%
\TextAt{13}{265.46}{577.308}{\Character{59}}%
\TextAt{13}{270.696}{577.308}{\Character{33}}%
\TextAt{21}{278.048}{575.515}{n}%
\TextAt{9}{283.684}{577.308}{\Character{41}}%
\TextAt{9}{291.539}{577.308}{\Character{61}}%
\TextAt{8}{303.983}{577.308}{det}%
\TextAt{9}{320.242}{577.308}{\Character{40}}%
\TextAt{8}{324.794}{577.308}{Tr}%
\TextAt{21}{337.801}{575.448}{L}%
\TextAt{3}{343.56}{575.448}{\Character{106}}%
\TextAt{21}{345.913}{575.448}{K}%
\TextAt{9}{354.116}{577.308}{\Character{40}}%
\TextAt{13}{358.668}{577.308}{\Character{33}}%
\TextAt{21}{365.957}{575.515}{i}%
\TextAt{13}{369.338}{577.308}{\Character{33}}%
\TextAt{21}{376.628}{575.515}{j}%
\TextAt{9}{381.01}{577.308}{\Character{41}\Character{41}}%
\TextAt{8}{79.226}{554.449}{of}%
\TextAt{8}{92.341}{554.449}{all}%
\TextAt{8}{108.385}{554.449}{bases}%
\TextAt{8}{138.859}{554.449}{of}%
\TextAt{13}{151.974}{554.449}{L}%
\TextAt{16}{163.259}{554.449}{\Character{106}}%
\TextAt{13}{169.901}{554.449}{K}%
\TextAt{8}{184.401}{554.449}{which}%
\TextAt{8}{217.672}{554.449}{are}%
\TextAt{8}{236.968}{554.449}{contained}%
\TextAt{8}{289.762}{554.449}{in}%
\TextAt{16}{303.2}{554.449}{\Character{79}}%
\TextAt{8}{313.05}{554.449}{.}%
\TextAt{8}{321.426}{554.449}{This}%
\TextAt{8}{347.931}{554.449}{is}%
\TextAt{8}{359.48}{554.449}{clear}%
\TextAt{8}{387.228}{554.449}{if}%
\TextAt{16}{397.737}{554.449}{\Character{79}}%
\TextAt{8}{411.269}{554.449}{admits}%
\TextAt{8}{449.49}{554.449}{an}%
\TextAt{8}{465.534}{554.449}{integral}%
\TextAt{8}{507.927}{554.449}{basis}%
\TextAt{8}{79.226}{537.403}{over}%
\TextAt{14}{104.472}{537.403}{o}%
\TextAt{8}{110.316}{537.403}{,}%
\TextAt{8}{117.716}{537.403}{since}%
\TextAt{8}{146.584}{537.403}{the}%
\TextAt{13}{166.941}{537.403}{\Character{11}}%
\TextAt{21}{174.386}{535.61}{i}%
\TextAt{8}{181.865}{537.403}{and}%
\TextAt{13}{204.828}{537.403}{\Character{12}}%
\TextAt{21}{211.436}{535.61}{i}%
\TextAt{8}{218.915}{537.403}{can}%
\TextAt{8}{240.574}{537.403}{be}%
\TextAt{8}{256.699}{537.403}{written}%
\TextAt{8}{297.869}{537.403}{in}%
\TextAt{8}{311.722}{537.403}{terms}%
\TextAt{8}{344.5}{537.403}{of}%
\TextAt{8}{358.03}{537.403}{such}%
\TextAt{8}{384.627}{537.403}{a}%
\TextAt{8}{394.582}{537.403}{basis}%
\TextAt{8}{423.523}{537.403}{with}%
\TextAt{8}{450.383}{537.403}{coefficients}%
\TextAt{8}{509.821}{537.403}{in}%
\TextAt{14}{523.673}{537.403}{o}%
\TextAt{8}{529.518}{537.403}{.}%
\TextAt{8}{79.226}{520.357}{If}%
\TextAt{8}{90.991}{520.357}{there}%
\TextAt{8}{120.976}{520.357}{is}%
\TextAt{8}{132.813}{520.357}{no}%
\TextAt{8}{149.145}{520.357}{such}%
\TextAt{8}{175.615}{520.357}{integral}%
\TextAt{8}{218.297}{520.357}{basis,}%
\TextAt{8}{250.38}{520.357}{it}%
\TextAt{8}{262.157}{520.357}{will}%
\TextAt{8}{284.335}{520.357}{exist}%
\TextAt{8}{312.108}{520.357}{after}%
\TextAt{8}{339.822}{520.357}{localizing}%
\TextAt{16}{391.601}{520.357}{\Character{79}}%
\TextAt{27}{401.119}{518.564}{p}%
\TextAt{16}{409.051}{520.357}{\Character{106}}%
\TextAt{14}{415.817}{520.357}{o}%
\TextAt{27}{421.661}{518.564}{p}%
\TextAt{8}{430.118}{520.357}{at}%
\TextAt{8}{444.502}{520.357}{every}%
\TextAt{8}{475.467}{520.357}{prime}%
\TextAt{8}{508.704}{520.357}{ideal}%
\TextAt{14}{79.226}{503.311}{p}%
\TextAt{8}{89.107}{503.311}{(see}%
\TextAt{8}{112.575}{503.311}{chap.}%
\TextAt{8}{144.766}{503.311}{I,}%
\TextAt{8}{156.135}{503.311}{(2.10)).}%
\TextAt{8}{199.074}{503.311}{The}%
\TextAt{8}{223.128}{503.311}{image}%
\TextAt{8}{256.949}{503.311}{of}%
\TextAt{13}{228.594}{475.856}{T}%
\TextAt{27}{235.455}{474.063}{p}%
\TextAt{9}{243.262}{475.856}{\Character{58}}%
\TextAt{9}{249.813}{475.856}{\Character{40}}%
\TextAt{8}{254.387}{475.856}{det}%
\TextAt{16}{272.639}{475.856}{\Character{79}}%
\TextAt{27}{282.156}{474.063}{p}%
\TextAt{9}{286.643}{475.856}{\Character{41}}%
\TextAt{16}{293.852}{475.856}{\Character{10}}%
\TextAt{9}{305.807}{475.856}{\Character{40}}%
\TextAt{8}{310.36}{475.856}{det}%
\TextAt{16}{328.611}{475.856}{\Character{79}}%
\TextAt{27}{338.129}{474.063}{p}%
\TextAt{9}{342.616}{475.856}{\Character{41}}%
\NativeArrowAt{350.489}{475.856}{11.955}{19.26}{\longrightarrow}%
\TextAt{14}{373.071}{475.856}{o}%
\TextAt{27}{378.915}{474.063}{p}%
\TextAt{8}{79.226}{448.402}{is}%
\TextAt{8}{89.726}{448.402}{therefore}%
\TextAt{8}{137.562}{448.402}{the}%
\TextAt{8}{156.454}{448.402}{discriminant}%
\TextAt{8}{222.569}{448.402}{ideal}%
\TextAt{8}{249.268}{448.402}{of}%
\TextAt{16}{261.334}{448.402}{\Character{79}}%
\TextAt{27}{270.851}{446.608}{p}%
\TextAt{16}{278.659}{448.402}{\Character{106}}%
\TextAt{14}{285.301}{448.402}{o}%
\TextAt{27}{291.145}{446.608}{p}%
\TextAt{8}{298.265}{448.402}{and}%
\TextAt{8}{319.764}{448.402}{at}%
\TextAt{8}{332.81}{448.402}{the}%
\TextAt{8}{351.702}{448.402}{same}%
\TextAt{8}{379.764}{448.402}{time}%
\TextAt{8}{405.16}{448.402}{the}%
\TextAt{8}{424.052}{448.402}{localization}%
\TextAt{8}{484.907}{448.402}{of}%
\TextAt{8}{496.973}{448.402}{the}%
\TextAt{8}{515.865}{448.402}{im-}%
\TextAt{8}{79.226}{431.356}{age}%
\TextAt{8}{99.231}{431.356}{of}%
\TextAt{13}{111.751}{431.356}{T}%
\TextAt{8}{120.238}{431.356}{.}%
\TextAt{8}{128.416}{431.356}{Since}%
\TextAt{8}{158.164}{431.356}{two}%
\TextAt{8}{179.472}{431.356}{ideals}%
\TextAt{8}{211.24}{431.356}{are}%
\TextAt{8}{229.941}{431.356}{equal}%
\TextAt{8}{260.024}{431.356}{when}%
\TextAt{8}{289.772}{431.356}{their}%
\TextAt{8}{316.926}{431.356}{localizations}%
\TextAt{8}{382.85}{431.356}{are,}%
\TextAt{8}{404.965}{431.356}{we}%
\TextAt{8}{421.383}{431.356}{find}%
\TextAt{8}{443.982}{431.356}{image}%
\TextAt{9}{473.905}{431.356}{\Character{40}}%
\TextAt{13}{478.458}{431.356}{T}%
\TextAt{9}{486.945}{431.356}{\Character{41}}%
\TextAt{9}{494.8}{431.356}{\Character{61}}%
\TextAt{14}{507.243}{431.356}{d}%
\TextAt{21}{513.204}{429.496}{L}%
\TextAt{3}{518.963}{429.496}{\Character{106}}%
\TextAt{21}{521.315}{429.496}{K}%
\TextAt{8}{529.518}{431.356}{.}%
\TextAt{8}{79.226}{414.31}{Furthermore,}%
\TextAt{13}{150.89}{414.31}{T}%
\TextAt{8}{163.766}{414.31}{has}%
\TextAt{8}{185.131}{414.31}{to}%
\TextAt{8}{199.932}{414.31}{be}%
\TextAt{8}{216.348}{414.31}{injective}%
\TextAt{8}{263.333}{414.31}{since}%
\TextAt{9}{292.493}{414.31}{\Character{40}}%
\TextAt{8}{297.045}{414.31}{det}%
\TextAt{16}{315.297}{414.31}{\Character{79}}%
\TextAt{9}{325.147}{414.31}{\Character{41}}%
\TextAt{3}{329.699}{418.648}{\Character{10}}%
\TextAt{17}{336.285}{418.648}{2}%
\TextAt{8}{345.407}{414.31}{is}%
\TextAt{8}{357.662}{414.31}{an}%
\TextAt{8}{374.412}{414.31}{invertible}%
\TextAt{14}{426.61}{414.31}{o}%
\TextAt{8}{432.454}{414.31}{-module.}%
\TextAt{8}{483.67}{414.31}{Therefore}%
\TextAt{13}{79.226}{397.264}{T}%
\TextAt{8}{91.61}{397.264}{is}%
\TextAt{8}{103.374}{397.264}{an}%
\TextAt{14}{119.633}{397.264}{o}%
\TextAt{8}{125.478}{397.264}{-module}%
\TextAt{8}{170.668}{397.264}{isomorphism.}%
\TextAt{8}{97.159}{377.229}{We}%
\TextAt{8}{117.303}{377.229}{now}%
\TextAt{8}{141.692}{377.229}{check}%
\TextAt{8}{173.23}{377.229}{that}%
\TextAt{13}{231.78}{360.183}{T}%
\TextAt{25}{238.641}{358.39}{C}%
\TextAt{9}{248.216}{360.183}{\Character{58}}%
\TextAt{9}{254.767}{360.183}{\Character{40}}%
\TextAt{8}{259.341}{360.183}{det}%
\TextAt{9}{275.6}{360.183}{\Character{40}}%
\TextAt{16}{280.152}{360.183}{\Character{79}}%
\TextAt{9}{290.002}{360.183}{\Character{41}}%
\TextAt{3}{294.554}{365.119}{\Character{10}}%
\TextAt{17}{301.141}{365.119}{2}%
\TextAt{9}{305.873}{360.183}{\Character{41}}%
\TextAt{25}{310.425}{358.39}{C}%
\NativeArrowAt{320}{360.183}{11.955}{19.26}{\longrightarrow}%
\TextAt{9}{342.582}{360.183}{\Character{40}}%
\TextAt{14}{347.135}{360.183}{d}%
\TextAt{21}{353.095}{358.323}{L}%
\TextAt{3}{358.854}{358.323}{\Character{106}}%
\TextAt{21}{361.207}{358.323}{K}%
\TextAt{9}{369.409}{360.183}{\Character{41}}%
\TextAt{25}{373.962}{358.39}{C}%
\TextAt{8}{79.226}{337.324}{is}%
\TextAt{8}{90.99}{337.324}{indeed}%
\TextAt{8}{128.051}{337.324}{an}%
\TextAt{8}{144.31}{337.324}{isometry.}%
\TextAt{8}{195.749}{337.324}{For}%
\TextAt{16}{216.719}{337.324}{\Character{79}}%
\TextAt{25}{226.236}{335.531}{C}%
\TextAt{9}{235.811}{337.324}{\Character{61}}%
\TextAt{16}{248.237}{337.324}{\Character{79}}%
\TextAt{16}{260.73}{337.324}{\Character{10}}%
\TextAt{25}{270.042}{335.531}{Z}%
\TextAt{18}{278.51}{337.324}{C}%
\TextAt{8}{287.144}{337.324}{,}%
\TextAt{8}{294.294}{337.324}{we}%
\TextAt{8}{311.521}{337.324}{obtain}%
\TextAt{8}{347.948}{337.324}{the}%
\TextAt{13}{368.105}{337.324}{K}%
\TextAt{25}{378.064}{335.531}{C}%
\TextAt{8}{384.318}{337.324}{-module}%
\TextAt{8}{429.509}{337.324}{decomposition}%
\TextAt{16}{271.013}{309.869}{\Character{79}}%
\TextAt{25}{280.531}{308.076}{C}%
\TextAt{9}{290.106}{309.869}{\Character{61}}%
\TextAt{20}{302.532}{321.227}{\Character{77}}%
\TextAt{21}{309.002}{296.121}{\Character{27}}%
\TextAt{16}{322.59}{309.869}{\Character{79}}%
\TextAt{21}{332.107}{308.076}{\Character{27}}%
\TextAt{13}{337.731}{309.869}{\Character{59}}%
\TextAt{8}{79.226}{274.554}{where}%
\TextAt{13}{113.035}{274.554}{\Character{27}}%
\TextAt{8}{124.015}{274.554}{varies}%
\TextAt{8}{156.928}{274.554}{over}%
\TextAt{8}{181.974}{274.554}{the}%
\TextAt{8}{202.13}{274.554}{set}%
\TextAt{8}{220.398}{274.554}{Hom}%
\TextAt{9}{244.786}{274.554}{\Character{40}}%
\TextAt{13}{249.339}{274.554}{K}%
\TextAt{13}{259.501}{274.554}{\Character{59}}%
\TextAt{18}{264.75}{274.554}{C}%
\TextAt{9}{273.384}{274.554}{\Character{41}}%
\TextAt{8}{277.937}{274.554}{,}%
\TextAt{8}{285.086}{274.554}{and}%
\TextAt{8}{307.849}{274.554}{the}%
\TextAt{8}{328.005}{274.554}{direct}%
\TextAt{8}{361.169}{274.554}{sum}%
\TextAt{16}{230.966}{247.099}{\Character{79}}%
\TextAt{21}{240.483}{245.306}{\Character{27}}%
\TextAt{9}{249.428}{247.099}{\Character{61}}%
\TextAt{20}{261.853}{258.457}{\Character{77}}%
\TextAt{21}{264.808}{232.554}{\Character{28}}%
\TextAt{3}{269.486}{232.554}{\Character{106}}%
\TextAt{21}{271.839}{232.554}{\Character{27}}%
\TextAt{9}{279.919}{247.099}{\Character{40}}%
\TextAt{16}{284.471}{247.099}{\Character{79}}%
\TextAt{16}{296.976}{247.099}{\Character{10}}%
\TextAt{3}{306.276}{245.306}{\Character{79}}%
\TextAt{21}{313.219}{245.306}{\Character{59}\Character{28}}%
\TextAt{18}{323.405}{247.099}{C}%
\TextAt{9}{332.039}{247.099}{\Character{41}}%
\TextAt{9}{339.894}{247.099}{\Character{61}}%
\TextAt{20}{352.337}{258.457}{\Character{77}}%
\TextAt{21}{355.292}{232.554}{\Character{28}}%
\TextAt{3}{359.971}{232.554}{\Character{106}}%
\TextAt{21}{362.323}{232.554}{\Character{27}}%
\TextAt{18}{372.396}{247.099}{C}%
\TextAt{8}{79.226}{208.994}{is}%
\TextAt{8}{91.187}{208.994}{taken}%
\TextAt{8}{123.256}{208.994}{over}%
\TextAt{8}{148.498}{208.994}{all}%
\TextAt{13}{164.954}{208.994}{\Character{28}}%
\NativeMathAt{175.028}{208.994}{11.955}{\in}%
\TextAt{8}{186.654}{208.994}{Hom}%
\TextAt{9}{211.043}{208.994}{\Character{40}}%
\TextAt{13}{215.595}{208.994}{L\Character{59}}%
\TextAt{18}{228.804}{208.994}{C}%
\TextAt{9}{237.438}{208.994}{\Character{41}}%
\TextAt{8}{246.084}{208.994}{such}%
\TextAt{8}{272.678}{208.994}{that}%
\TextAt{13}{298.243}{208.994}{\Character{28}}%
\TextAt{16}{304.662}{208.994}{\Character{106}}%
\TextAt{21}{307.983}{207.201}{K}%
\TextAt{9}{319.842}{208.994}{\Character{61}}%
\TextAt{13}{332.602}{208.994}{\Character{27}}%
\TextAt{8}{339.684}{208.994}{.}%
\TextAt{8}{348.722}{208.994}{The}%
\TextAt{8}{372.973}{208.994}{mapping}%
\TextAt{16}{421.301}{208.994}{\Character{79}}%
\TextAt{25}{430.818}{207.201}{C}%
\NativeMathAt{440.728}{208.994}{11.955}{\rightarrow}%
\TextAt{13}{456.339}{208.994}{K}%
\TextAt{25}{466.298}{207.201}{C}%
\TextAt{8}{476.647}{208.994}{induced}%
\TextAt{8}{520.408}{208.994}{by}%
\TextAt{8}{79.226}{191.948}{Tr}%
\TextAt{21}{92.233}{190.089}{L}%
\TextAt{3}{97.992}{190.089}{\Character{106}}%
\TextAt{21}{100.345}{190.089}{K}%
\TextAt{9}{111.868}{191.948}{\Character{58}}%
\TextAt{16}{118.441}{191.948}{\Character{79}}%
\NativeMathAt{131.604}{191.948}{11.955}{\rightarrow}%
\TextAt{14}{146.887}{191.948}{o}%
\TextAt{8}{156.629}{191.948}{is}%
\TextAt{8}{168.393}{191.948}{given,}%
\TextAt{8}{202.214}{191.948}{for}%
\TextAt{13}{220.099}{191.948}{x}%
\TextAt{9}{230.072}{191.948}{\Character{61}}%
\TextAt{20}{242.497}{200.915}{\Character{76}}%
\TextAt{21}{255.781}{188.461}{\Character{27}}%
\TextAt{13}{263.397}{191.948}{x}%
\TextAt{21}{270.049}{190.155}{\Character{27}}%
\TextAt{8}{275.672}{191.948}{,}%
\TextAt{13}{282.822}{191.948}{x}%
\TextAt{21}{289.474}{190.155}{\Character{27}}%
\NativeMathAt{298.418}{191.948}{11.955}{\in}%
\TextAt{16}{309.692}{191.948}{\Character{79}}%
\TextAt{21}{319.227}{190.155}{\Character{27}}%
\TextAt{8}{324.85}{191.948}{,}%
\TextAt{8}{331.999}{191.948}{by}%
\TextAt{8}{244.324}{164.494}{Tr}%
\TextAt{21}{257.331}{162.634}{L}%
\TextAt{3}{263.09}{162.634}{\Character{106}}%
\TextAt{21}{265.442}{162.634}{K}%
\TextAt{9}{273.645}{164.494}{\Character{40}}%
\TextAt{13}{278.197}{164.494}{x}%
\TextAt{9}{284.85}{164.494}{\Character{41}}%
\TextAt{9}{292.705}{164.494}{\Character{61}}%
\TextAt{20}{305.148}{175.851}{\Character{88}}%
\TextAt{21}{311.22}{150.745}{\Character{27}}%
\TextAt{8}{324.409}{164.494}{Tr}%
\TextAt{21}{337.416}{162.7}{\Character{27}}%
\TextAt{9}{343.04}{164.494}{\Character{40}}%
\TextAt{13}{347.592}{164.494}{x}%
\TextAt{21}{354.244}{162.7}{\Character{27}}%
\TextAt{9}{359.868}{164.494}{\Character{41}}%
\TextAt{13}{364.42}{164.494}{\Character{59}}%
\TextAt{8}{79.226}{129.178}{where}%
\TextAt{8}{114.11}{129.178}{Tr}%
\TextAt{21}{127.117}{127.385}{\Character{27}}%
\TextAt{9}{132.741}{129.178}{\Character{40}}%
\TextAt{13}{137.293}{129.178}{x}%
\TextAt{21}{143.945}{127.385}{\Character{27}}%
\TextAt{9}{149.568}{129.178}{\Character{41}}%
\TextAt{9}{159.252}{129.178}{\Character{61}}%
\TextAt{20}{173.529}{138.145}{\Character{80}}%
\TextAt{21}{186.149}{125.691}{\Character{28}}%
\TextAt{3}{190.827}{125.691}{\Character{106}}%
\TextAt{21}{193.179}{125.691}{\Character{27}}%
\TextAt{13}{200.795}{129.178}{x}%
\TextAt{21}{207.447}{127.385}{\Character{27}}%
\TextAt{21}{212.102}{127.385}{\Character{59}\Character{28}}%
\TextAt{8}{219.631}{129.178}{,}%
\TextAt{8}{228.123}{129.178}{the}%
\TextAt{13}{249.354}{129.178}{x}%
\TextAt{21}{256.006}{127.385}{\Character{27}}%
\TextAt{21}{260.661}{127.385}{\Character{59}\Character{28}}%
\NativeMathAt{273.342}{129.178}{11.955}{\in}%
\TextAt{18}{286.464}{129.178}{C}%
\TextAt{8}{300.07}{129.178}{being}%
\TextAt{8}{332.682}{129.178}{the}%
\TextAt{8}{353.913}{129.178}{components}%
\TextAt{8}{419.438}{129.178}{of}%
\TextAt{13}{433.842}{129.178}{x}%
\TextAt{21}{440.494}{127.385}{\Character{27}}%
\TextAt{8}{446.118}{129.178}{.}%
\TextAt{8}{457.79}{129.178}{The}%
\TextAt{8}{482.918}{129.178}{metric}%
\TextAt{8}{520.408}{129.178}{on}%
\TextAt{9}{79.226}{112.132}{\Character{40}}%
\TextAt{13}{83.778}{112.132}{i}%
\TextAt{3}{87.772}{110.339}{\Character{3}}%
\TextAt{16}{92.504}{112.132}{\Character{79}}%
\TextAt{9}{102.354}{112.132}{\Character{41}}%
\TextAt{25}{106.906}{110.339}{C}%
\TextAt{9}{116.481}{112.132}{\Character{61}}%
\TextAt{16}{128.907}{112.132}{\Character{79}}%
\TextAt{25}{138.425}{110.339}{C}%
\TextAt{8}{148.576}{112.132}{is}%
\TextAt{8}{160.34}{112.132}{the}%
\TextAt{8}{180.497}{112.132}{orthogonal}%
\TextAt{8}{239.053}{112.132}{sum}%
\TextAt{8}{263.824}{112.132}{of}%
\TextAt{8}{277.154}{112.132}{the}%
\TextAt{8}{297.311}{112.132}{standard}%
\TextAt{8}{346.159}{112.132}{metrics}%
\TextAt{16}{232.311}{84.677}{\Character{104}}%
\TextAt{13}{236.96}{84.677}{x\Character{59}}%
\TextAt{13}{248.832}{84.677}{y}%
\TextAt{16}{254.993}{84.677}{\Character{105}}%
\TextAt{21}{259.642}{82.884}{\Character{27}}%
\TextAt{9}{268.586}{84.677}{\Character{61}}%
\TextAt{20}{281.012}{96.035}{\Character{88}}%
\TextAt{21}{283.568}{70.132}{\Character{28}}%
\TextAt{3}{288.247}{70.132}{\Character{106}}%
\TextAt{21}{290.599}{70.132}{\Character{27}}%
\TextAt{13}{300.273}{84.677}{x}%
\TextAt{21}{306.925}{82.884}{\Character{28}}%
\Vector{q 0 G 0 g 0.478 w 240.102 -682.502 m 246.238 -682.502 l S Q}%
\TextAt{13}{312.102}{84.677}{y}%
\TextAt{21}{318.238}{81.855}{\Character{28}}%
\TextAt{9}{326.735}{84.677}{\Character{61}}%
\TextAt{8}{339.161}{84.677}{Tr}%
\TextAt{21}{352.168}{82.884}{\Character{27}}%
\TextAt{9}{357.792}{84.677}{\Character{40}}%
\TextAt{13}{362.344}{84.677}{x}%
\Vector{q 0 G 0 g 0.478 w 296.996 -682.502 m 303.133 -682.502 l S Q}%
\TextAt{13}{368.996}{84.677}{y}%
\TextAt{9}{375.133}{84.677}{\Character{41}}%
\TextAt{8}{297.211}{43.033}{231}%
\Destination{78.23}{809.815}{nav.page.231}%
\Destination{79.36}{352.355}{nav.print.252}%
\Hotlink{144.266}{501.155}{5.22}{11.047}{nav.chapter.1}%
\Hotlink{155.635}{501.155}{30.936}{11.047}{nav.number.1.2.10}%
}

% Page 244 (body page 232)
\Page{%
\TextAt{2}{79.226}{796.943}{III.}%
\TextAt{2}{98.317}{796.943}{Riemann–Roch}%
\TextAt{2}{174.375}{796.943}{Theory}%
\Vector{q 0 G 0 g 0.299 w 7.226 18.729 m 460.769 18.729 l S Q}%
\TextAt{8}{79.226}{761.124}{on}%
\TextAt{8}{96.24}{761.124}{the}%
\TextAt{18}{117.152}{761.124}{C}%
\TextAt{8}{125.787}{761.124}{-vector}%
\TextAt{8}{165.564}{761.124}{spaces}%
\TextAt{9}{202.208}{761.124}{\Character{40}}%
\TextAt{13}{206.761}{761.124}{i}%
\TextAt{3}{210.754}{759.33}{\Character{3}}%
\TextAt{16}{215.486}{761.124}{\Character{79}}%
\TextAt{9}{225.336}{761.124}{\Character{41}}%
\TextAt{21}{229.888}{759.33}{\Character{27}}%
\TextAt{9}{240.12}{761.124}{\Character{61}}%
\TextAt{16}{253.833}{761.124}{\Character{79}}%
\TextAt{21}{263.35}{759.33}{\Character{27}}%
\TextAt{9}{273.582}{761.124}{\Character{61}}%
\TextAt{20}{287.294}{770.09}{\Character{76}}%
\TextAt{21}{300.578}{757.637}{\Character{28}}%
\TextAt{3}{305.256}{757.637}{\Character{106}}%
\TextAt{21}{307.609}{757.637}{\Character{27}}%
\TextAt{18}{315.224}{761.124}{C}%
\TextAt{8}{323.859}{761.124}{.}%
\TextAt{8}{334.573}{761.124}{Now}%
\TextAt{8}{361.988}{761.124}{let}%
\TextAt{13}{379.648}{761.124}{x}%
\TextAt{21}{386.3}{759.33}{i}%
\TextAt{13}{389.682}{761.124}{\Character{59}}%
\TextAt{13}{394.906}{761.124}{y}%
\TextAt{21}{400.634}{759.33}{i}%
\NativeMathAt{408.623}{761.124}{11.955}{\in}%
\TextAt{16}{421.188}{761.124}{\Character{79}}%
\TextAt{21}{430.719}{759.33}{\Character{27}}%
\TextAt{8}{436.342}{761.124}{,}%
\TextAt{13}{444.435}{761.124}{i}%
\TextAt{9}{453.037}{761.124}{\Character{61}}%
\TextAt{9}{466.738}{761.124}{1}%
\TextAt{13}{472.603}{761.124}{\Character{59}}%
\TextAt{13}{477.827}{761.124}{\Character{58}}%
\TextAt{13}{483.064}{761.124}{\Character{58}}%
\TextAt{13}{488.288}{761.124}{\Character{58}}%
\TextAt{13}{493.525}{761.124}{\Character{59}}%
\TextAt{13}{498.761}{761.124}{n}%
\TextAt{8}{505.811}{761.124}{,}%
\TextAt{8}{513.904}{761.124}{and}%
\TextAt{8}{79.226}{744.078}{write}%
\TextAt{13}{109.554}{744.078}{x}%
\TextAt{9}{120.237}{744.078}{\Character{61}}%
\TextAt{13}{133.372}{744.078}{x}%
\TextAt{17}{140.024}{742.284}{1}%
\TextAt{16}{147.697}{744.078}{\Character{94}}%
\TextAt{16}{158.588}{744.078}{\Character{1}}%
\TextAt{16}{163.896}{744.078}{\Character{1}}%
\TextAt{16}{169.204}{744.078}{\Character{1}}%
\TextAt{16}{175.457}{744.078}{\Character{94}}%
\TextAt{13}{186.407}{744.078}{x}%
\TextAt{21}{193.059}{742.284}{n}%
\TextAt{8}{198.695}{744.078}{,}%
\TextAt{13}{206.365}{744.078}{y}%
\TextAt{9}{216.532}{744.078}{\Character{61}}%
\TextAt{13}{229.667}{744.078}{y}%
\TextAt{17}{235.375}{742.284}{1}%
\TextAt{16}{243.048}{744.078}{\Character{94}}%
\TextAt{16}{253.939}{744.078}{\Character{1}}%
\TextAt{16}{259.247}{744.078}{\Character{1}}%
\TextAt{16}{264.555}{744.078}{\Character{1}}%
\TextAt{16}{270.808}{744.078}{\Character{94}}%
\TextAt{13}{281.758}{744.078}{y}%
\TextAt{21}{287.465}{742.284}{n}%
\NativeMathAt{297.132}{744.078}{11.955}{\in}%
\TextAt{8}{309.133}{744.078}{det}%
\TextAt{9}{325.392}{744.078}{\Character{40}}%
\TextAt{16}{329.944}{744.078}{\Character{79}}%
\TextAt{21}{339.462}{742.284}{\Character{27}}%
\TextAt{9}{345.085}{744.078}{\Character{41}}%
\TextAt{8}{349.638}{744.078}{.}%
\TextAt{8}{359.335}{744.078}{The}%
\TextAt{8}{383.806}{744.078}{map}%
\TextAt{13}{410.236}{744.078}{T}%
\TextAt{25}{417.097}{742.284}{C}%
\TextAt{8}{427.666}{744.078}{splits}%
\TextAt{8}{458.771}{744.078}{into}%
\TextAt{8}{482.93}{744.078}{the}%
\TextAt{8}{503.503}{744.078}{direct}%
\TextAt{8}{79.226}{727.032}{sum}%
\TextAt{13}{103.997}{727.032}{T}%
\TextAt{25}{110.858}{725.238}{C}%
\TextAt{9}{120.433}{727.032}{\Character{61}}%
\TextAt{20}{132.859}{735.998}{\Character{76}}%
\TextAt{21}{146.142}{723.545}{\Character{27}}%
\TextAt{13}{153.758}{727.032}{T}%
\TextAt{21}{160.619}{725.238}{\Character{27}}%
\TextAt{8}{170.14}{727.032}{of}%
\TextAt{8}{183.47}{727.032}{the}%
\TextAt{8}{203.627}{727.032}{maps}%
\TextAt{13}{200.228}{698.764}{T}%
\TextAt{21}{207.089}{696.971}{\Character{27}}%
\TextAt{9}{216.033}{698.764}{\Character{58}}%
\TextAt{8}{222.606}{698.764}{det}%
\TextAt{9}{238.865}{698.764}{\Character{40}}%
\TextAt{16}{243.417}{698.764}{\Character{79}}%
\TextAt{21}{252.935}{696.971}{\Character{27}}%
\TextAt{9}{258.558}{698.764}{\Character{41}}%
\TextAt{16}{265.767}{698.764}{\Character{10}}%
\TextAt{25}{275.066}{696.971}{C}%
\TextAt{8}{283.977}{698.764}{det}%
\TextAt{9}{300.236}{698.764}{\Character{40}}%
\TextAt{16}{304.788}{698.764}{\Character{79}}%
\TextAt{21}{314.306}{696.971}{\Character{27}}%
\TextAt{9}{319.929}{698.764}{\Character{41}}%
\NativeArrowAt{327.802}{698.764}{11.955}{19.26}{\longrightarrow}%
\TextAt{9}{350.384}{698.764}{\Character{40}}%
\TextAt{14}{354.937}{698.764}{d}%
\TextAt{21}{360.897}{696.905}{L}%
\TextAt{3}{366.656}{696.905}{\Character{106}}%
\TextAt{21}{369.009}{696.905}{K}%
\TextAt{9}{377.211}{698.764}{\Character{41}}%
\TextAt{21}{381.764}{696.971}{\Character{27}}%
\TextAt{9}{390.708}{698.764}{\Character{61}}%
\TextAt{18}{403.133}{698.764}{C}%
\TextAt{8}{79.226}{670.497}{which}%
\TextAt{8}{112.713}{670.497}{are}%
\TextAt{8}{132.223}{670.497}{given}%
\TextAt{8}{162.793}{670.497}{by}%
\TextAt{13}{235.393}{653.451}{T}%
\TextAt{21}{242.254}{651.658}{\Character{27}}%
\TextAt{9}{247.877}{653.451}{\Character{40}}%
\TextAt{13}{252.43}{653.451}{x}%
\TextAt{16}{261.738}{653.451}{\Character{10}}%
\TextAt{13}{273.694}{653.451}{y}%
\TextAt{9}{279.83}{653.451}{\Character{41}}%
\TextAt{9}{287.685}{653.451}{\Character{61}}%
\TextAt{8}{300.129}{653.451}{det}%
\TextAt{9}{316.388}{653.451}{\Character{40}}%
\TextAt{8}{320.94}{653.451}{Tr}%
\TextAt{21}{333.947}{651.658}{\Character{27}}%
\TextAt{9}{339.571}{653.451}{\Character{40}}%
\TextAt{13}{344.123}{653.451}{x}%
\TextAt{21}{350.775}{651.658}{i}%
\TextAt{13}{354.157}{653.451}{y}%
\TextAt{21}{359.864}{651.658}{j}%
\TextAt{9}{364.246}{653.451}{\Character{41}\Character{41}}%
\TextAt{13}{373.351}{653.451}{\Character{58}}%
\TextAt{8}{79.226}{630.244}{For}%
\TextAt{8}{100.196}{630.244}{any}%
\TextAt{8}{122.313}{630.244}{two}%
\TextAt{13}{144.43}{630.244}{n}%
\TextAt{8}{151.417}{630.244}{-tuples}%
\TextAt{13}{189.841}{630.244}{x}%
\TextAt{3}{196.493}{634.582}{\Character{48}}%
\TextAt{21}{196.493}{627.288}{i}%
\TextAt{13}{199.875}{630.244}{\Character{59}}%
\TextAt{13}{205.099}{630.244}{y}%
\TextAt{3}{211.255}{634.582}{\Character{48}}%
\TextAt{21}{210.826}{627.288}{i}%
\NativeMathAt{217.528}{630.244}{11.955}{\in}%
\TextAt{16}{228.802}{630.244}{\Character{79}}%
\TextAt{21}{238.337}{628.45}{\Character{27}}%
\TextAt{8}{247.858}{630.244}{we}%
\TextAt{8}{265.085}{630.244}{form}%
\TextAt{8}{292.726}{630.244}{the}%
\TextAt{8}{312.882}{630.244}{matrices}%
\TextAt{13}{211.894}{601.976}{A}%
\TextAt{9}{223.99}{601.976}{\Character{61}}%
\TextAt{9}{236.4}{601.976}{\Character{40}}%
\TextAt{8}{240.968}{601.976}{Tr}%
\TextAt{21}{253.975}{600.183}{\Character{27}}%
\TextAt{9}{259.598}{601.976}{\Character{40}}%
\TextAt{13}{264.151}{601.976}{x}%
\TextAt{21}{270.803}{600.183}{i}%
\TextAt{13}{274.184}{601.976}{y}%
\TextAt{21}{279.892}{600.183}{j}%
\TextAt{9}{284.274}{601.976}{\Character{41}\Character{41}}%
\TextAt{13}{293.379}{601.976}{\Character{59}}%
\TextAt{13}{310.559}{601.976}{A}%
\TextAt{3}{319.353}{606.912}{\Character{48}}%
\TextAt{9}{325.469}{601.976}{\Character{61}}%
\TextAt{9}{337.879}{601.976}{\Character{40}}%
\TextAt{8}{342.447}{601.976}{Tr}%
\TextAt{21}{355.454}{600.183}{\Character{27}}%
\TextAt{9}{361.078}{601.976}{\Character{40}}%
\Vector{q 0 G 0 g 0.478 w 293.63 -161.805 m 303.663 -161.805 l S Q}%
\TextAt{13}{365.63}{601.976}{x}%
\TextAt{3}{372.282}{606.094}{\Character{48}}%
\TextAt{21}{372.282}{598.912}{i}%
\Vector{q 0 G 0 g 0.478 w 305.656 -161.805 m 315.746 -161.805 l S Q}%
\TextAt{13}{377.656}{601.976}{y}%
\TextAt{3}{383.792}{606.094}{\Character{48}}%
\TextAt{21}{383.363}{598.912}{j}%
\TextAt{9}{387.746}{601.976}{\Character{41}\Character{41}}%
\TextAt{13}{396.85}{601.976}{\Character{59}}%
\TextAt{13}{212.169}{581.941}{B}%
\TextAt{9}{224.986}{581.941}{\Character{61}}%
\TextAt{9}{237.396}{581.941}{\Character{40}}%
\TextAt{8}{241.964}{581.941}{Tr}%
\TextAt{21}{254.971}{580.148}{\Character{27}}%
\TextAt{9}{260.595}{581.941}{\Character{40}}%
\TextAt{13}{265.147}{581.941}{x}%
\TextAt{21}{271.799}{580.148}{i}%
\Vector{q 0 G 0 g 0.478 w 203.18 -181.839 m 214.215 -181.839 l S Q}%
\TextAt{13}{275.18}{581.941}{x}%
\TextAt{3}{281.832}{586.059}{\Character{48}}%
\TextAt{21}{281.832}{578.877}{j}%
\TextAt{9}{286.215}{581.941}{\Character{41}\Character{41}}%
\TextAt{13}{295.319}{581.941}{\Character{59}}%
\TextAt{13}{312.499}{581.941}{B}%
\TextAt{3}{322.015}{586.877}{\Character{48}}%
\TextAt{9}{328.131}{581.941}{\Character{61}}%
\TextAt{9}{340.541}{581.941}{\Character{40}}%
\TextAt{8}{345.109}{581.941}{Tr}%
\TextAt{21}{358.116}{580.148}{\Character{27}}%
\TextAt{9}{363.739}{581.941}{\Character{40}}%
\TextAt{13}{368.291}{581.941}{y}%
\TextAt{21}{373.999}{580.148}{i}%
\Vector{q 0 G 0 g 0.478 w 305.38 -181.839 m 315.47 -181.839 l S Q}%
\TextAt{13}{377.38}{581.941}{y}%
\TextAt{3}{383.517}{586.059}{\Character{48}}%
\TextAt{21}{383.088}{578.877}{j}%
\TextAt{9}{387.47}{581.941}{\Character{41}\Character{41}}%
\TextAt{13}{396.575}{581.941}{\Character{58}}%
\ErrataAt{79.22}{553.655}{453.64}{Then one has $\det A\,\det A^{\prime}=\det B\,\det B^{\prime}$, and we obtain}%
\TextAt{16}{143.406}{525.406}{\Character{104}}%
\TextAt{13}{148.056}{525.406}{T}%
\TextAt{21}{154.917}{523.613}{\Character{27}}%
\TextAt{9}{160.54}{525.406}{\Character{40}}%
\TextAt{13}{165.092}{525.406}{x}%
\TextAt{16}{174.401}{525.406}{\Character{10}}%
\TextAt{13}{186.356}{525.406}{y}%
\TextAt{9}{192.493}{525.406}{\Character{41}}%
\TextAt{13}{197.045}{525.406}{\Character{59}}%
\TextAt{13}{202.269}{525.406}{T}%
\TextAt{21}{209.15}{523.613}{\Character{27}}%
\TextAt{9}{214.774}{525.406}{\Character{40}}%
\TextAt{13}{219.326}{525.406}{x}%
\TextAt{3}{225.978}{530.342}{\Character{48}}%
\TextAt{16}{231.43}{525.406}{\Character{10}}%
\TextAt{13}{243.385}{525.406}{y}%
\TextAt{3}{249.522}{530.342}{\Character{48}}%
\TextAt{9}{252.317}{525.406}{\Character{41}}%
\TextAt{16}{256.869}{525.406}{\Character{105}}%
\TextAt{17}{261.518}{523.547}{\Character{40}}%
\TextAt{27}{264.812}{523.547}{d}%
\TextAt{22}{268.785}{521.809}{L}%
\TextAt{23}{273.844}{521.809}{\Character{106}}%
\TextAt{22}{276.113}{521.809}{K}%
\TextAt{17}{283.156}{523.547}{\Character{41}}%
\TextAt{22}{286.45}{522.55}{\Character{27}}%
\TextAt{9}{295.264}{525.406}{\Character{61}}%
\TextAt{13}{307.69}{525.406}{T}%
\TextAt{21}{314.551}{523.613}{\Character{27}}%
\TextAt{9}{320.174}{525.406}{\Character{40}}%
\TextAt{13}{324.727}{525.406}{x}%
\TextAt{16}{334.035}{525.406}{\Character{10}}%
\TextAt{13}{345.99}{525.406}{y}%
\TextAt{9}{352.127}{525.406}{\Character{41}}%
\Vector{q 0 G 0 g 0.478 w 284.679 -237.954 m 339.259 -237.954 l S Q}%
\TextAt{13}{356.679}{525.406}{T}%
\TextAt{21}{363.541}{523.613}{\Character{27}}%
\TextAt{9}{369.164}{525.406}{\Character{40}}%
\TextAt{13}{373.716}{525.406}{x}%
\TextAt{3}{380.368}{528.86}{\Character{48}}%
\TextAt{16}{385.82}{525.406}{\Character{10}}%
\TextAt{13}{397.775}{525.406}{y}%
\TextAt{3}{403.912}{528.86}{\Character{48}}%
\TextAt{9}{406.707}{525.406}{\Character{41}}%
\TextAt{9}{158.682}{505.372}{\Character{61}}%
\TextAt{8}{171.108}{505.372}{det}%
\TextAt{9}{187.367}{505.372}{\Character{40}}%
\TextAt{8}{191.919}{505.372}{Tr}%
\TextAt{21}{204.926}{503.578}{\Character{27}}%
\TextAt{9}{210.55}{505.372}{\Character{40}}%
\TextAt{13}{215.102}{505.372}{x}%
\TextAt{21}{221.754}{503.578}{i}%
\TextAt{13}{225.135}{505.372}{y}%
\TextAt{21}{230.843}{503.578}{j}%
\TextAt{9}{235.225}{505.372}{\Character{41}\Character{41}}%
\TextAt{8}{246.322}{505.372}{det}%
\TextAt{9}{262.581}{505.372}{\Character{40}}%
\TextAt{8}{267.134}{505.372}{Tr}%
\TextAt{21}{280.141}{503.578}{\Character{27}}%
\TextAt{9}{285.764}{505.372}{\Character{40}}%
\Vector{q 0 G 0 g 0.478 w 218.317 -258.409 m 228.35 -258.409 l S Q}%
\TextAt{13}{290.317}{505.372}{x}%
\TextAt{3}{296.969}{509.489}{\Character{48}}%
\TextAt{21}{296.969}{502.308}{i}%
\Vector{q 0 G 0 g 0.478 w 230.343 -258.409 m 240.432 -258.409 l S Q}%
\TextAt{13}{302.343}{505.372}{y}%
\TextAt{3}{308.479}{509.489}{\Character{48}}%
\TextAt{21}{308.05}{502.308}{j}%
\TextAt{9}{312.432}{505.372}{\Character{41}\Character{41}}%
\TextAt{9}{324.829}{505.372}{\Character{61}}%
\TextAt{8}{337.283}{505.372}{det}%
\TextAt{9}{353.542}{505.372}{\Character{40}}%
\TextAt{13}{358.095}{505.372}{AA}%
\TextAt{3}{375.646}{510.308}{\Character{48}}%
\TextAt{9}{378.441}{505.372}{\Character{41}}%
\TextAt{9}{386.296}{505.372}{\Character{61}}%
\TextAt{8}{398.739}{505.372}{det}%
\TextAt{9}{414.998}{505.372}{\Character{40}}%
\TextAt{13}{419.551}{505.372}{B}%
\TextAt{13}{429.031}{505.372}{B}%
\TextAt{3}{438.543}{510.308}{\Character{48}}%
\TextAt{9}{441.338}{505.372}{\Character{41}}%
\TextAt{9}{158.682}{485.337}{\Character{61}}%
\TextAt{8}{171.108}{485.337}{det}%
\TextAt{9}{187.367}{485.337}{\Character{40}}%
\TextAt{8}{191.919}{485.337}{Tr}%
\TextAt{21}{204.926}{483.543}{\Character{27}}%
\TextAt{9}{210.55}{485.337}{\Character{40}}%
\TextAt{13}{215.102}{485.337}{x}%
\TextAt{21}{221.754}{483.543}{i}%
\Vector{q 0 G 0 g 0.478 w 153.135 -278.444 m 164.17 -278.444 l S Q}%
\TextAt{13}{225.135}{485.337}{x}%
\TextAt{3}{231.788}{489.455}{\Character{48}}%
\TextAt{21}{231.788}{482.273}{j}%
\TextAt{9}{236.17}{485.337}{\Character{41}\Character{41}}%
\TextAt{8}{247.267}{485.337}{det}%
\TextAt{9}{263.526}{485.337}{\Character{40}}%
\TextAt{8}{268.078}{485.337}{Tr}%
\TextAt{21}{281.085}{483.544}{\Character{27}}%
\TextAt{9}{286.709}{485.337}{\Character{40}}%
\TextAt{13}{291.261}{485.337}{y}%
\TextAt{21}{296.969}{483.544}{i}%
\Vector{q 0 G 0 g 0.478 w 228.35 -278.444 m 238.44 -278.444 l S Q}%
\TextAt{13}{300.35}{485.337}{y}%
\TextAt{3}{306.487}{489.455}{\Character{48}}%
\TextAt{21}{306.058}{482.273}{j}%
\TextAt{9}{310.44}{485.337}{\Character{41}\Character{41}}%
\TextAt{9}{322.838}{485.337}{\Character{61}}%
\TextAt{8}{335.291}{485.337}{det}%
\TextAt{9}{351.55}{485.337}{\Character{40}}%
\TextAt{16}{356.102}{485.337}{\Character{104}}%
\TextAt{13}{360.752}{485.337}{x}%
\TextAt{21}{367.404}{483.544}{i}%
\TextAt{13}{370.785}{485.337}{\Character{59}}%
\TextAt{13}{376.009}{485.337}{x}%
\TextAt{3}{382.681}{490.273}{\Character{48}}%
\TextAt{21}{382.681}{482.381}{j}%
\TextAt{16}{387.063}{485.337}{\Character{105}}%
\TextAt{21}{391.713}{483.544}{\Character{27}}%
\TextAt{9}{397.336}{485.337}{\Character{41}}%
\TextAt{8}{403.881}{485.337}{det}%
\TextAt{9}{420.14}{485.337}{\Character{40}}%
\TextAt{16}{424.692}{485.337}{\Character{104}}%
\TextAt{13}{429.341}{485.337}{y}%
\TextAt{21}{435.049}{483.544}{i}%
\TextAt{13}{438.43}{485.337}{\Character{59}}%
\TextAt{13}{443.654}{485.337}{y}%
\TextAt{3}{449.811}{490.273}{\Character{48}}%
\TextAt{21}{449.382}{482.381}{j}%
\TextAt{16}{453.764}{485.337}{\Character{105}}%
\TextAt{21}{458.414}{483.544}{\Character{27}}%
\TextAt{9}{464.037}{485.337}{\Character{41}}%
\TextAt{9}{158.682}{465.302}{\Character{61}}%
\TextAt{16}{171.108}{465.302}{\Character{104}}%
\TextAt{13}{175.757}{465.302}{x\Character{59}}%
\TextAt{13}{187.629}{465.302}{x}%
\TextAt{3}{194.305}{470.238}{\Character{48}}%
\TextAt{16}{197.1}{465.302}{\Character{105}}%
\TextAt{7}{201.75}{463.509}{det}%
\TextAt{3}{214.917}{463.509}{\Character{79}}%
\TextAt{22}{221.625}{462.512}{\Character{27}}%
\TextAt{16}{227.119}{465.302}{\Character{104}}%
\TextAt{13}{231.768}{465.302}{y}%
\TextAt{13}{237.889}{465.302}{\Character{59}}%
\TextAt{13}{243.125}{465.302}{y}%
\TextAt{3}{249.285}{470.238}{\Character{48}}%
\TextAt{16}{252.08}{465.302}{\Character{105}}%
\TextAt{7}{256.73}{463.509}{det}%
\TextAt{3}{269.897}{463.509}{\Character{79}}%
\TextAt{22}{276.605}{462.512}{\Character{27}}%
\TextAt{9}{285.42}{465.302}{\Character{61}}%
\TextAt{16}{297.845}{465.302}{\Character{104}}%
\TextAt{13}{302.494}{465.302}{x}%
\TextAt{16}{311.803}{465.302}{\Character{10}}%
\TextAt{13}{323.758}{465.302}{y}%
\TextAt{13}{329.879}{465.302}{\Character{59}}%
\TextAt{13}{335.115}{465.302}{x}%
\TextAt{3}{341.791}{470.238}{\Character{48}}%
\TextAt{16}{347.243}{465.302}{\Character{10}}%
\TextAt{13}{359.198}{465.302}{y}%
\TextAt{3}{365.335}{470.238}{\Character{48}}%
\TextAt{16}{368.13}{465.302}{\Character{105}}%
\TextAt{17}{372.779}{463.298}{\Character{40}}%
\TextAt{7}{376.072}{463.298}{det}%
\TextAt{3}{389.24}{463.298}{\Character{79}}%
\TextAt{22}{395.947}{462.302}{\Character{27}}%
\TextAt{17}{400.943}{463.298}{\Character{41}}%
\TextAt{23}{404.236}{465.568}{\Character{10}}%
\TextAt{24}{409.993}{465.568}{2}%
\TextAt{13}{414.642}{465.302}{\Character{58}}%
\TextAt{8}{79.226}{437.034}{This}%
\TextAt{8}{105.946}{437.034}{shows}%
\TextAt{8}{139.564}{437.034}{that}%
\TextAt{13}{164.933}{437.034}{T}%
\TextAt{25}{171.794}{435.241}{C}%
\TextAt{8}{181.945}{437.034}{is}%
\TextAt{8}{193.709}{437.034}{an}%
\TextAt{8}{209.968}{437.034}{isometry.}%
\Vector{q 0 G 0 g 0.399 w 452.334 -336.966 m 452.334 -328.896 l S Q}%
\Vector{q 0 G 0 g 0.399 w 452.534 -329.095 m 459.707 -329.095 l S Q}%
\Vector{q 0 G 0 g 0.399 w 452.534 -336.766 m 459.707 -336.766 l S Q}%
\Vector{q 0 G 0 g 0.399 w 459.906 -336.966 m 459.906 -328.896 l S Q}%
\TextAt{8}{97.159}{408.243}{We}%
\TextAt{8}{117.812}{408.243}{now}%
\TextAt{8}{142.71}{408.243}{set}%
\TextAt{8}{161.486}{408.243}{out}%
\TextAt{8}{182.809}{408.243}{to}%
\TextAt{8}{197.628}{408.243}{rewrite}%
\TextAt{8}{237.804}{408.243}{the}%
\TextAt{8}{258.469}{408.243}{results}%
\TextAt{8}{296.171}{408.243}{obtained}%
\TextAt{8}{344.811}{408.243}{in}%
\TextAt{8}{358.973}{408.243}{(7.2)}%
\TextAt{8}{387.457}{408.243}{and}%
\TextAt{8}{410.728}{408.243}{(7.4)}%
\TextAt{8}{439.212}{408.243}{in}%
\TextAt{8}{453.374}{408.243}{the}%
\TextAt{8}{474.039}{408.243}{language}%
\TextAt{8}{523.337}{408.243}{of}%
\TextAt{26}{79.226}{391.197}{Grothendieck}%
\TextAt{8}{164.741}{391.197}{’s}%
\TextAt{8}{175.727}{391.197}{general}%
\TextAt{8}{215.274}{391.197}{formalism.}%
\TextAt{8}{273.937}{391.197}{For}%
\TextAt{8}{294.129}{391.197}{the}%
\TextAt{8}{313.507}{391.197}{homomorphism}%
\TextAt{13}{395.399}{391.197}{i}%
\TextAt{3}{399.392}{389.404}{\Character{3}}%
\TextAt{8}{407.244}{391.197}{there}%
\TextAt{8}{436.377}{391.197}{is}%
\TextAt{8}{447.363}{391.197}{the}%
\TextAt{8}{466.741}{391.197}{commutative}%
\TextAt{8}{79.226}{374.151}{diagram}%
\TextAt{13}{249.317}{355.763}{K}%
\TextAt{17}{259.276}{353.97}{0}%
\TextAt{9}{264.008}{355.763}{\Character{40}}%
\NativeDiagram{%
\path[draw=black,line width=0.478bp] (196.561,-408.394) -- (206.411,-408.394);%
}%
\TextAt{16}{268.561}{355.763}{\Character{79}}%
\TextAt{9}{278.411}{355.763}{\Character{41}}%
\TextAt{18}{337.933}{355.763}{Z}%
\TextAt{13}{251.32}{302.77}{K}%
\TextAt{17}{261.279}{300.977}{0}%
\TextAt{9}{266.011}{302.77}{\Character{40}}%
\NativeDiagram{%
\path[draw=black,line width=0.478bp] (198.564,-463.874) -- (204.408,-463.874);%
}%
\TextAt{14}{270.564}{302.77}{o}%
\TextAt{9}{276.408}{302.77}{\Character{41}}%
\TextAt{18}{336.307}{302.77}{Z}%
\TextAt{13}{344.277}{302.77}{\Character{59}}%
\NativeDiagram{%
\path[draw=black,line width=0.478bp] (216.315,-415.247) -- (260.103,-415.247);%
\path[draw=black,line width=0.478bp,line cap=round,line join=round] (257.856,-412.378) .. controls (258.31,-414.099) and (259.322,-414.912) .. (260.342,-415.247) .. controls (259.322,-415.582) and (258.31,-416.394) .. (257.856,-418.116);%
}%
\TextAt{7}{306.563}{361.569}{rk}%
\NativeDiagram{%
\path[draw=black,line width=0.478bp] (194.14,-425.804) -- (194.14,-457.206);%
\path[draw=black,line width=0.478bp,line cap=round,line join=round] (197.009,-454.959) .. controls (195.287,-455.413) and (194.474,-456.425) .. (194.14,-457.445) .. controls (193.805,-456.425) and (192.992,-455.413) .. (191.271,-454.959);%
}%
\TextAt{21}{256.124}{330.12}{i}%
\TextAt{23}{259.008}{329.124}{\Character{3}}%
\NativeDiagram{%
\path[draw=black,line width=0.478bp] (269.918,-422.815) -- (269.918,-457.937);%
\path[draw=black,line width=0.478bp,line cap=round,line join=round] (272.788,-455.689) .. controls (271.066,-456.144) and (270.253,-457.156) .. (269.918,-458.176) .. controls (269.584,-457.156) and (268.771,-456.144) .. (267.049,-455.689);%
}%
\TextAt{17}{344.734}{331.393}{\Character{91}}%
\TextAt{21}{347.087}{331.393}{L}%
\TextAt{17}{352.846}{331.393}{\Character{58}}%
\TextAt{21}{355.198}{331.393}{K}%
\TextAt{17}{362.903}{331.393}{\Character{93}}%
\NativeDiagram{%
\path[draw=black,line width=0.478bp] (214.312,-468.241) -- (258.477,-468.241);%
\path[draw=black,line width=0.478bp,line cap=round,line join=round] (256.23,-465.372) .. controls (256.684,-467.093) and (257.696,-467.906) .. (258.716,-468.241) .. controls (257.696,-468.576) and (256.684,-469.389) .. (256.23,-471.11);%
}%
\TextAt{7}{304.748}{308.574}{rk}%
\TextAt{8}{79.226}{279.278}{because}%
\TextAt{9}{122.767}{279.278}{\Character{91}}%
\TextAt{13}{126.019}{279.278}{L}%
\TextAt{9}{137.712}{279.278}{\Character{58}}%
\TextAt{13}{144.692}{279.278}{K}%
\TextAt{9}{155.51}{279.278}{\Character{93}}%
\TextAt{8}{162.898}{279.278}{times}%
\TextAt{8}{194.412}{279.278}{the}%
\TextAt{8}{214.808}{279.278}{rank}%
\TextAt{8}{242.042}{279.278}{of}%
\TextAt{8}{255.611}{279.278}{an}%
\TextAt{16}{272.109}{279.278}{\Character{79}}%
\TextAt{8}{281.959}{279.278}{-module}%
\TextAt{13}{327.389}{279.278}{M}%
\TextAt{8}{344.099}{279.278}{is}%
\TextAt{8}{356.102}{279.278}{its}%
\TextAt{8}{372.66}{279.278}{rank}%
\TextAt{8}{399.894}{279.278}{as}%
\TextAt{14}{414.504}{279.278}{o}%
\TextAt{8}{420.348}{279.278}{-module.}%
\TextAt{8}{470.807}{279.278}{Therefore}%
\TextAt{13}{524.044}{279.278}{i}%
\TextAt{3}{528.037}{277.484}{\Character{3}}%
\TextAt{8}{79.226}{262.232}{induces}%
\TextAt{8}{120.902}{262.232}{a}%
\TextAt{8}{130.657}{262.232}{homomorphism}%
\TextAt{13}{260.682}{245.186}{i}%
\TextAt{3}{264.676}{243.392}{\Character{3}}%
\TextAt{9}{272.729}{245.186}{\Character{58}}%
\TextAt{13}{279.301}{245.186}{I}%
\TextAt{9}{285.404}{245.186}{\Character{40}}%
\Vector{q 0 G 0 g 0.478 w 217.957 -518.971 m 227.806 -518.971 l S Q}%
\TextAt{16}{289.957}{245.186}{\Character{79}}%
\TextAt{9}{299.806}{245.186}{\Character{41}}%
\NativeArrowAt{307.68}{245.186}{11.955}{19.26}{\longrightarrow}%
\TextAt{13}{330.262}{245.186}{I}%
\TextAt{9}{336.364}{245.186}{\Character{40}}%
\Vector{q 0 G 0 g 0.478 w 268.917 -521.458 m 274.761 -521.458 l S Q}%
\TextAt{14}{340.917}{245.186}{o}%
\TextAt{9}{346.761}{245.186}{\Character{41}}%
\TextAt{8}{79.226}{221.978}{between}%
\TextAt{8}{124.417}{221.978}{the}%
\TextAt{8}{144.573}{221.978}{kernels}%
\TextAt{8}{183.655}{221.978}{of}%
\TextAt{8}{196.985}{221.978}{both}%
\TextAt{8}{224.625}{221.978}{rank}%
\TextAt{8}{251.62}{221.978}{homomorphisms,}%
\TextAt{8}{342.156}{221.978}{so}%
\TextAt{8}{356.526}{221.978}{that}%
\TextAt{8}{381.895}{221.978}{there}%
\TextAt{8}{411.807}{221.978}{is}%
\TextAt{8}{423.571}{221.978}{a}%
\TextAt{8}{433.326}{221.978}{homomorphism}%
\TextAt{13}{238.062}{193.711}{i}%
\TextAt{3}{242.056}{191.918}{\Character{3}}%
\TextAt{9}{250.109}{193.711}{\Character{58}}%
\TextAt{8}{256.682}{193.711}{gr}%
\TextAt{13}{269.087}{193.711}{K}%
\TextAt{17}{279.046}{191.918}{0}%
\TextAt{9}{283.778}{193.711}{\Character{40}}%
\Vector{q 0 G 0 g 0.478 w 216.331 -570.446 m 226.181 -570.446 l S Q}%
\TextAt{16}{288.331}{193.711}{\Character{79}}%
\TextAt{9}{298.181}{193.711}{\Character{41}}%
\NativeArrowAt{306.054}{193.711}{11.955}{19.26}{\longrightarrow}%
\TextAt{8}{328.636}{193.711}{gr}%
\TextAt{13}{341.041}{193.711}{K}%
\TextAt{17}{351}{191.918}{0}%
\TextAt{9}{355.733}{193.711}{\Character{40}}%
\Vector{q 0 G 0 g 0.478 w 288.285 -572.933 m 294.129 -572.933 l S Q}%
\TextAt{14}{360.285}{193.711}{o}%
\TextAt{9}{366.129}{193.711}{\Character{41}}%
\TextAt{13}{370.681}{193.711}{\Character{58}}%
\TextAt{8}{79.226}{165.443}{It}%
\TextAt{8}{91.891}{165.443}{is}%
\TextAt{8}{103.648}{165.443}{called}%
\TextAt{8}{136.804}{165.443}{the}%
\TextAt{12}{156.953}{165.443}{Gysin}%
\TextAt{12}{195.62}{165.443}{map}%
\TextAt{8}{220.845}{165.443}{.}%
\TextAt{8}{229.291}{165.443}{(7.3)}%
\TextAt{8}{257.259}{165.443}{immediately}%
\TextAt{8}{323.913}{165.443}{gives}%
\TextAt{8}{352.587}{165.443}{the}%
\TextAt{8}{372.736}{165.443}{following}%
\TextAt{8}{422.163}{165.443}{explicit}%
\TextAt{8}{463.449}{165.443}{description}%
\TextAt{8}{523.337}{165.443}{of}%
\TextAt{8}{79.226}{148.397}{it.}%
\TextAt{8}{297.211}{43.033}{232}%
\Destination{78.23}{809.815}{nav.page.232}%
\Destination{79.36}{406.155}{nav.print.253}%
\Hotlink{358.473}{406.055}{25.078}{11.047}{nav.number.3.7.2}%
\Hotlink{410.228}{406.055}{25.078}{11.047}{nav.number.3.7.4}%
\Hotlink{228.791}{163.255}{25.078}{11.047}{nav.number.3.7.3}%
\Destination{79.226}{567.655}{errata.24}%
}

% Page 245 (body page 233)
\Page{%
\TextAt{2}{79.226}{796.943}{III.}%
\TextAt{2}{98.317}{796.943}{Riemann–Roch}%
\TextAt{2}{174.375}{796.943}{Theory}%
\Vector{q 0 G 0 g 0.299 w 7.226 18.729 m 460.769 18.729 l S Q}%
\TextAt{12}{79.226}{761.124}{Corollary}%
\TextAt{12}{139.014}{761.124}{(7.5).}%
\TextAt{8}{174.33}{761.124}{The}%
\TextAt{8}{198.383}{761.124}{diagram}%
\TextAt{8}{217.592}{730.231}{gr}%
\TextAt{13}{229.997}{730.231}{K}%
\TextAt{17}{239.956}{728.438}{0}%
\TextAt{9}{244.688}{730.231}{\Character{40}}%
\NativeDiagram{%
\path[draw=black,line width=0.478bp] (177.241,-33.926) -- (187.091,-33.926);%
}%
\TextAt{16}{249.241}{730.231}{\Character{79}}%
\TextAt{9}{259.091}{730.231}{\Character{41}}%
\TextAt{18}{316.987}{730.231}{Z}%
\TextAt{16}{327.613}{730.231}{\Character{8}}%
\TextAt{8}{339.569}{730.231}{Pic}%
\TextAt{9}{355.983}{730.231}{\Character{40}}%
\NativeDiagram{%
\path[draw=black,line width=0.478bp] (288.535,-33.926) -- (298.385,-33.926);%
}%
\TextAt{16}{360.535}{730.231}{\Character{79}}%
\TextAt{9}{370.385}{730.231}{\Character{41}}%
\TextAt{8}{219.595}{677.237}{gr}%
\TextAt{13}{232}{677.237}{K}%
\TextAt{17}{241.959}{675.444}{0}%
\TextAt{9}{246.691}{677.237}{\Character{40}}%
\NativeDiagram{%
\path[draw=black,line width=0.478bp] (179.244,-89.406) -- (185.088,-89.406);%
}%
\TextAt{14}{251.244}{677.237}{o}%
\TextAt{9}{257.088}{677.237}{\Character{41}}%
\TextAt{18}{318.99}{677.237}{Z}%
\TextAt{16}{329.616}{677.237}{\Character{8}}%
\TextAt{8}{341.572}{677.237}{Pic}%
\TextAt{9}{357.986}{677.237}{\Character{40}}%
\NativeDiagram{%
\path[draw=black,line width=0.478bp] (290.538,-89.406) -- (296.383,-89.406);%
}%
\TextAt{14}{362.538}{677.237}{o}%
\TextAt{9}{368.383}{677.237}{\Character{41}}%
\NativeDiagram{%
\path[draw=black,line width=0.478bp] (196.996,-40.78) -- (239.158,-40.78);%
\path[draw=black,line width=0.478bp,line cap=round,line join=round] (236.911,-37.911) .. controls (237.365,-39.632) and (238.377,-40.445) .. (239.397,-40.78) .. controls (238.377,-41.115) and (237.365,-41.927) .. (236.911,-43.649);%
}%
\TextAt{7}{276.207}{736.866}{id}%
\TextAt{3}{284.672}{736.866}{\Character{8}}%
\TextAt{7}{292.67}{736.866}{det}%
\NativeMathAt{287.023}{727.392}{7.97}{\sim}%
\NativeDiagram{%
\path[draw=black,line width=0.478bp] (168.618,-51.337) -- (168.618,-82.739);%
\path[draw=black,line width=0.478bp,line cap=round,line join=round] (171.487,-80.492) .. controls (169.766,-80.946) and (168.953,-81.958) .. (168.618,-82.978) .. controls (168.283,-81.958) and (167.47,-80.946) .. (165.749,-80.492);%
}%
\TextAt{21}{230.602}{704.587}{i}%
\TextAt{23}{233.485}{703.591}{\Character{3}}%
\NativeDiagram{%
\path[draw=black,line width=0.478bp] (273.964,-51.337) -- (273.964,-82.739);%
\path[draw=black,line width=0.478bp,line cap=round,line join=round] (276.833,-80.492) .. controls (275.112,-80.946) and (274.299,-81.958) .. (273.964,-82.978) .. controls (273.63,-81.958) and (272.817,-80.946) .. (271.095,-80.492);%
}%
\TextAt{17}{348.78}{705.35}{\Character{91}}%
\TextAt{21}{351.133}{705.35}{L}%
\TextAt{17}{356.892}{705.35}{\Character{58}}%
\TextAt{21}{359.244}{705.35}{K}%
\TextAt{17}{366.948}{705.35}{\Character{93}}%
\TextAt{3}{369.301}{705.35}{\Character{8}}%
\TextAt{21}{375.887}{705.35}{N}%
\TextAt{22}{382.609}{703.612}{L}%
\TextAt{23}{387.667}{703.612}{\Character{106}}%
\TextAt{22}{389.937}{703.612}{K}%
\NativeDiagram{%
\path[draw=black,line width=0.478bp] (194.993,-93.774) -- (241.161,-93.774);%
\path[draw=black,line width=0.478bp,line cap=round,line join=round] (238.914,-90.905) .. controls (239.368,-92.626) and (240.38,-93.439) .. (241.4,-93.774) .. controls (240.38,-94.109) and (239.368,-94.922) .. (238.914,-96.643);%
}%
\TextAt{7}{276.207}{683.872}{id}%
\TextAt{3}{284.672}{683.872}{\Character{8}}%
\TextAt{7}{292.67}{683.872}{det}%
\NativeMathAt{287.023}{674.398}{7.97}{\sim}%
\TextAt{8}{79.226}{649.111}{is}%
\TextAt{8}{90.99}{649.111}{commutative.}%
\TextAt{8}{97.159}{623.099}{We}%
\TextAt{8}{117.303}{623.099}{now}%
\TextAt{8}{141.692}{623.099}{consider}%
\TextAt{8}{187.277}{623.099}{the}%
\TextAt{8}{207.433}{623.099}{following}%
\TextAt{8}{256.868}{623.099}{diagram}%
\TextAt{13}{239.476}{592.756}{K}%
\TextAt{17}{249.436}{590.963}{0}%
\TextAt{9}{254.168}{592.756}{\Character{40}}%
\NativeDiagram{%
\path[draw=black,line width=0.478bp] (186.72,-171.401) -- (196.57,-171.401);%
}%
\TextAt{16}{258.72}{592.756}{\Character{79}}%
\TextAt{9}{268.57}{592.756}{\Character{41}}%
\TextAt{8}{326.466}{592.756}{gr}%
\TextAt{13}{338.871}{592.756}{K}%
\TextAt{17}{348.831}{590.963}{0}%
\TextAt{9}{353.563}{592.756}{\Character{40}}%
\NativeDiagram{%
\path[draw=black,line width=0.478bp] (286.115,-171.401) -- (295.965,-171.401);%
}%
\TextAt{16}{358.115}{592.756}{\Character{79}}%
\TextAt{9}{367.965}{592.756}{\Character{41}}%
\TextAt{13}{241.479}{539.762}{K}%
\TextAt{17}{251.439}{537.969}{0}%
\TextAt{9}{256.171}{539.762}{\Character{40}}%
\NativeDiagram{%
\path[draw=black,line width=0.478bp] (188.723,-226.881) -- (194.567,-226.881);%
}%
\TextAt{14}{260.723}{539.762}{o}%
\TextAt{9}{266.567}{539.762}{\Character{41}}%
\TextAt{8}{328.469}{539.762}{gr}%
\TextAt{13}{340.874}{539.762}{K}%
\TextAt{17}{350.834}{537.969}{0}%
\TextAt{9}{355.566}{539.762}{\Character{40}}%
\NativeDiagram{%
\path[draw=black,line width=0.478bp] (288.118,-226.881) -- (293.962,-226.881);%
}%
\TextAt{14}{360.118}{539.762}{o}%
\TextAt{9}{365.962}{539.762}{\Character{41}}%
\NativeDiagram{%
\path[draw=black,line width=0.478bp] (206.475,-178.255) -- (248.637,-178.255);%
\path[draw=black,line width=0.478bp,line cap=round,line join=round] (246.39,-175.386) .. controls (246.844,-177.107) and (247.857,-177.92) .. (248.877,-178.255) .. controls (247.857,-178.59) and (246.844,-179.402) .. (246.39,-181.124);%
}%
\TextAt{7}{295.563}{598.64}{ch}%
\NativeDiagram{%
\path[draw=black,line width=0.478bp] (184.3,-188.812) -- (184.3,-220.214);%
\path[draw=black,line width=0.478bp,line cap=round,line join=round] (187.169,-217.967) .. controls (185.447,-218.421) and (184.634,-219.433) .. (184.3,-220.453) .. controls (183.965,-219.433) and (183.152,-218.421) .. (181.431,-217.967);%
}%
\TextAt{21}{246.284}{567.112}{i}%
\TextAt{23}{249.167}{566.116}{\Character{3}}%
\NativeDiagram{%
\path[draw=black,line width=0.478bp] (277.493,-188.812) -- (277.493,-220.214);%
\path[draw=black,line width=0.478bp,line cap=round,line join=round] (280.363,-217.967) .. controls (278.641,-218.421) and (277.828,-219.433) .. (277.494,-220.453) .. controls (277.159,-219.433) and (276.346,-218.421) .. (274.624,-217.967);%
}%
\TextAt{21}{352.309}{567.112}{i}%
\TextAt{23}{355.192}{566.116}{\Character{3}}%
\NativeDiagram{%
\path[draw=black,line width=0.478bp] (204.472,-231.249) -- (250.64,-231.249);%
\path[draw=black,line width=0.478bp,line cap=round,line join=round] (248.393,-228.38) .. controls (248.847,-230.101) and (249.859,-230.914) .. (250.879,-231.249) .. controls (249.859,-231.584) and (248.847,-232.397) .. (248.393,-234.118);%
}%
\TextAt{7}{295.563}{545.646}{ch}%
\TextAt{8}{79.226}{510.476}{where}%
\TextAt{8}{113.188}{510.476}{the}%
\TextAt{8}{133.498}{510.476}{Gysin}%
\TextAt{8}{167.281}{510.476}{map}%
\TextAt{13}{193.448}{510.476}{i}%
\TextAt{3}{197.441}{508.683}{\Character{3}}%
\TextAt{8}{206.224}{510.476}{on}%
\TextAt{8}{222.636}{510.476}{the}%
\TextAt{8}{242.946}{510.476}{right}%
\TextAt{8}{271.396}{510.476}{is}%
\TextAt{8}{283.313}{510.476}{explicitly}%
\TextAt{8}{334.192}{510.476}{given}%
\TextAt{8}{364.914}{510.476}{by}%
\TextAt{8}{381.326}{510.476}{(7.5),}%
\TextAt{8}{412.744}{510.476}{whereas}%
\TextAt{8}{457.179}{510.476}{the}%
\TextAt{8}{477.489}{510.476}{determina-}%
\TextAt{8}{79.226}{493.43}{tion}%
\TextAt{8}{103.498}{493.43}{of}%
\TextAt{8}{117.035}{493.43}{the}%
\TextAt{8}{137.398}{493.43}{composite}%
\TextAt{8}{192.622}{493.43}{ch}%
\TextAt{16}{206.319}{493.43}{\Character{14}}%
\TextAt{13}{212.296}{493.43}{i}%
\TextAt{3}{216.29}{491.637}{\Character{3}}%
\TextAt{8}{225.126}{493.43}{is}%
\TextAt{8}{237.096}{493.43}{precisely}%
\TextAt{8}{285.159}{493.43}{the}%
\TextAt{8}{305.522}{493.43}{Riemann–Roch}%
\TextAt{8}{387.358}{493.43}{problem.}%
\TextAt{8}{438.054}{493.43}{The}%
\TextAt{8}{462.314}{493.43}{difficulty}%
\TextAt{8}{511.298}{493.43}{that}%
\TextAt{8}{79.226}{476.384}{confronts}%
\TextAt{8}{130.186}{476.384}{us}%
\TextAt{8}{145.364}{476.384}{here}%
\TextAt{8}{170.883}{476.384}{lies}%
\TextAt{8}{191.262}{476.384}{in}%
\TextAt{8}{205.077}{476.384}{the}%
\TextAt{8}{225.395}{476.384}{fact}%
\TextAt{8}{248.643}{476.384}{that}%
\TextAt{8}{274.174}{476.384}{the}%
\TextAt{8}{294.493}{476.384}{diagram}%
\TextAt{8}{340.192}{476.384}{is}%
\TextAt{19}{352.119}{476.384}{not}%
\TextAt{19}{372.916}{476.384}{commutative}%
\TextAt{8}{437.175}{476.384}{.}%
\TextAt{8}{446.11}{476.384}{In}%
\TextAt{8}{460.893}{476.384}{order}%
\TextAt{8}{491.625}{476.384}{to}%
\TextAt{8}{506.097}{476.384}{make}%
\TextAt{8}{79.226}{459.338}{it}%
\TextAt{8}{91.147}{459.338}{commute,}%
\TextAt{8}{145.073}{459.338}{we}%
\TextAt{8}{162.517}{459.338}{need}%
\TextAt{8}{190.04}{459.338}{a}%
\TextAt{8}{200.012}{459.338}{correction,}%
\TextAt{8}{258.17}{459.338}{which}%
\TextAt{8}{291.873}{459.338}{will}%
\TextAt{8}{314.195}{459.338}{be}%
\TextAt{8}{330.336}{459.338}{provided}%
\TextAt{8}{378.685}{459.338}{via}%
\TextAt{8}{398.09}{459.338}{the}%
\TextAt{8}{418.463}{459.338}{module}%
\TextAt{8}{459.973}{459.338}{of}%
\TextAt{8}{473.52}{459.338}{differentials}%
\TextAt{8}{79.226}{442.292}{(with}%
\TextAt{8}{110.441}{442.292}{trivial}%
\TextAt{8}{145.243}{442.292}{metric),}%
\TextAt{8}{189.465}{442.292}{by}%
\TextAt{8}{205.724}{442.292}{the}%
\TextAt{19}{225.88}{442.292}{Todd}%
\TextAt{19}{254.871}{442.292}{class}%
\TextAt{8}{278.782}{442.292}{,}%
\TextAt{8}{285.931}{442.292}{which}%
\TextAt{8}{319.417}{442.292}{is}%
\TextAt{8}{331.181}{442.292}{defined}%
\TextAt{8}{371.494}{442.292}{as}%
\TextAt{8}{385.864}{442.292}{follows.}%
\TextAt{8}{97.159}{422.258}{The}%
\TextAt{8}{121.849}{422.258}{module}%
\TextAt{9}{163.778}{422.258}{\Character{10}}%
\TextAt{17}{172.232}{426.596}{1}%
\TextAt{3}{172.232}{418.706}{\Character{79}}%
\TextAt{3}{179.166}{418.706}{\Character{106}}%
\TextAt{27}{181.527}{418.706}{o}%
\TextAt{8}{190.455}{422.258}{of}%
\TextAt{8}{204.421}{422.258}{differentials}%
\TextAt{8}{268.204}{422.258}{is}%
\TextAt{8}{280.604}{422.258}{only}%
\TextAt{8}{306.932}{422.258}{a}%
\TextAt{8}{317.323}{422.258}{coherent,}%
\TextAt{8}{368.521}{422.258}{and}%
\TextAt{8}{391.92}{422.258}{not}%
\TextAt{8}{413.37}{422.258}{a}%
\TextAt{8}{423.761}{422.258}{projective}%
\TextAt{16}{478.375}{422.258}{\Character{79}}%
\TextAt{8}{488.225}{422.258}{-module.}%
\TextAt{8}{79.226}{405.212}{But}%
\TextAt{8}{102.467}{405.212}{its}%
\TextAt{8}{118.786}{405.212}{class}%
\TextAt{9}{146.223}{405.212}{\Character{91}\Character{10}}%
\TextAt{17}{157.929}{409.55}{1}%
\TextAt{3}{157.929}{401.66}{\Character{79}}%
\TextAt{3}{164.863}{401.66}{\Character{106}}%
\TextAt{27}{167.224}{401.66}{o}%
\TextAt{9}{171.618}{405.212}{\Character{93}}%
\TextAt{8}{178.768}{405.212}{is}%
\TextAt{8}{190.531}{405.212}{viewed}%
\TextAt{8}{228.896}{405.212}{as}%
\TextAt{8}{243.266}{405.212}{an}%
\TextAt{8}{259.525}{405.212}{element}%
\TextAt{8}{302.767}{405.212}{of}%
\TextAt{13}{316.097}{405.212}{K}%
\TextAt{17}{326.056}{403.418}{0}%
\TextAt{9}{330.788}{405.212}{\Character{40}}%
\Vector{q 0 G 0 g 0.478 w 263.34 -358.945 m 273.19 -358.945 l S Q}%
\TextAt{16}{335.34}{405.212}{\Character{79}}%
\TextAt{9}{345.19}{405.212}{\Character{41}}%
\TextAt{8}{353.64}{405.212}{via}%
\TextAt{8}{372.828}{405.212}{the}%
\TextAt{8}{392.984}{405.212}{Poincaré}%
\TextAt{8}{440.948}{405.212}{isomorphism}%
\TextAt{13}{260.39}{376.21}{K}%
\TextAt{17}{270.349}{374.417}{0}%
\TextAt{9}{275.081}{376.211}{\Character{40}}%
\Vector{q 0 G 0 g 0.478 w 207.633 -387.946 m 217.483 -387.946 l S Q}%
\TextAt{16}{279.633}{376.21}{\Character{79}}%
\TextAt{9}{289.483}{376.211}{\Character{41}}%
\NativeMathAt{299.944}{382.988}{7.97}{\sim}%
\NativeArrowAt{297.356}{376.21}{11.955}{13.175}{\longrightarrow}%
\TextAt{13}{313.85}{376.211}{K}%
\TextAt{17}{324.668}{381.147}{0}%
\TextAt{9}{329.4}{376.21}{\Character{40}}%
\Vector{q 0 G 0 g 0.478 w 261.952 -387.946 m 271.802 -387.946 l S Q}%
\TextAt{16}{333.952}{376.21}{\Character{79}}%
\TextAt{9}{343.802}{376.211}{\Character{41}}%
\TextAt{13}{348.354}{376.21}{\Character{59}}%
\TextAt{8}{79.226}{347.209}{and}%
\TextAt{8}{101.989}{347.209}{since}%
\TextAt{8}{130.657}{347.209}{rk}%
\TextAt{3}{141.393}{345.416}{\Character{79}}%
\TextAt{9}{148.834}{347.209}{\Character{40}\Character{10}}%
\TextAt{17}{161.841}{351.548}{1}%
\TextAt{3}{161.841}{343.657}{\Character{79}}%
\TextAt{3}{168.775}{343.657}{\Character{106}}%
\TextAt{27}{171.136}{343.657}{o}%
\TextAt{9}{175.53}{347.209}{\Character{41}}%
\TextAt{9}{183.385}{347.209}{\Character{61}}%
\TextAt{9}{195.806}{347.209}{0}%
\TextAt{8}{201.682}{347.209}{,}%
\TextAt{8}{208.831}{347.209}{it}%
\TextAt{8}{220.535}{347.209}{lies}%
\TextAt{8}{240.751}{347.209}{in}%
\TextAt{13}{254.404}{347.209}{I}%
\TextAt{9}{260.507}{347.209}{\Character{40}}%
\Vector{q 0 G 0 g 0.478 w 193.059 -416.948 m 202.909 -416.948 l S Q}%
\TextAt{16}{265.059}{347.209}{\Character{79}}%
\TextAt{9}{274.909}{347.209}{\Character{41}}%
\TextAt{8}{279.461}{347.209}{.}%
\TextAt{12}{79.226}{321.197}{Definition}%
\TextAt{12}{142.015}{321.197}{(7.6).}%
\TextAt{8}{177.33}{321.197}{The}%
\TextAt{12}{201.384}{321.197}{Todd}%
\TextAt{12}{236.138}{321.197}{class}%
\TextAt{8}{266.91}{321.197}{of}%
\TextAt{16}{280.24}{321.197}{\Character{79}}%
\TextAt{16}{293.403}{321.197}{\Character{106}}%
\TextAt{14}{300.053}{321.197}{o}%
\TextAt{8}{309.794}{321.197}{is}%
\TextAt{8}{321.558}{321.197}{defined}%
\TextAt{8}{361.871}{321.197}{to}%
\TextAt{8}{376.181}{321.197}{be}%
\TextAt{8}{392.106}{321.197}{the}%
\TextAt{8}{412.262}{321.197}{element}%
\TextAt{8}{150.406}{292.196}{Td}%
\TextAt{9}{165.362}{292.196}{\Character{40}}%
\TextAt{16}{169.914}{292.196}{\Character{79}}%
\TextAt{16}{183.077}{292.196}{\Character{106}}%
\TextAt{14}{189.726}{292.196}{o}%
\TextAt{9}{195.57}{292.196}{\Character{41}}%
\TextAt{9}{203.425}{292.196}{\Character{61}}%
\TextAt{9}{215.846}{292.196}{1}%
\TextAt{16}{224.379}{292.196}{\Character{0}}%
\TextAt{17}{237.529}{296.903}{1}%
\Vector{q 0 G 0 g 0.478 w 165.529 -478.815 m 169.764 -478.815 l S Q}%
\TextAt{17}{237.529}{288.073}{2}%
\TextAt{13}{242.959}{292.196}{c}%
\TextAt{17}{247.997}{290.403}{1}%
\TextAt{9}{252.729}{292.196}{\Character{40}\Character{91}\Character{10}}%
\TextAt{17}{268.988}{297.132}{1}%
\TextAt{3}{268.988}{289.24}{\Character{79}}%
\TextAt{3}{275.922}{289.24}{\Character{106}}%
\TextAt{27}{278.283}{289.24}{o}%
\TextAt{9}{282.677}{292.196}{\Character{93}\Character{41}}%
\TextAt{9}{293.771}{292.196}{\Character{61}}%
\TextAt{9}{306.193}{292.196}{1}%
\TextAt{16}{314.737}{292.196}{\Character{0}}%
\TextAt{17}{327.888}{296.903}{1}%
\Vector{q 0 G 0 g 0.478 w 255.888 -478.815 m 260.122 -478.815 l S Q}%
\TextAt{17}{327.888}{288.073}{2}%
\TextAt{9}{333.318}{292.196}{\Character{91}\Character{10}}%
\TextAt{17}{345.024}{297.132}{1}%
\TextAt{3}{345.024}{289.24}{\Character{79}}%
\TextAt{3}{351.958}{289.24}{\Character{106}}%
\TextAt{27}{354.319}{289.24}{o}%
\TextAt{9}{358.713}{292.196}{\Character{93}}%
\NativeMathAt{365.286}{292.196}{11.955}{\in}%
\TextAt{8}{376.577}{292.196}{gr}%
\TextAt{13}{388.982}{292.196}{K}%
\TextAt{17}{398.941}{290.403}{0}%
\TextAt{9}{403.673}{292.196}{\Character{40}}%
\Vector{q 0 G 0 g 0.478 w 336.226 -471.961 m 346.076 -471.961 l S Q}%
\TextAt{16}{408.226}{292.196}{\Character{79}}%
\TextAt{9}{418.076}{292.196}{\Character{41}}%
\TextAt{16}{425.285}{292.196}{\Character{10}}%
\TextAt{18}{437.24}{292.196}{Z}%
\TextAt{9}{445.21}{292.196}{\Character{91}}%
\TextAt{17}{449.657}{296.903}{1}%
\Vector{q 0 G 0 g 0.478 w 377.657 -478.815 m 381.891 -478.815 l S Q}%
\TextAt{17}{449.657}{288.073}{2}%
\TextAt{9}{455.087}{292.196}{\Character{93}}%
\TextAt{13}{458.338}{292.196}{\Character{58}}%
\ErrataAt{97.1}{254.228}{437.9}{The formula defines the Todd class canonically in $\operatorname{gr}K_0(\overline O)\otimes_{\mathbb Z}\mathbb Z[1/2]$.}%
\ErrataAt{79.19}{237.182}{455.81}{The module of differentials $\Omega_{O\mid\mathfrak o}^{1}$ is connected}%
\TextAt{8}{79.226}{220.136}{with}%
\TextAt{8}{105.886}{220.136}{the}%
\TextAt{12}{126.043}{220.136}{different}%
\TextAt{14}{179.1}{220.136}{D}%
\TextAt{21}{189.041}{218.277}{L}%
\TextAt{3}{194.8}{218.277}{\Character{106}}%
\TextAt{21}{197.153}{218.277}{K}%
\TextAt{8}{209.253}{220.136}{of}%
\TextAt{8}{222.583}{220.136}{the}%
\TextAt{8}{242.739}{220.136}{extension}%
\TextAt{13}{294.505}{220.136}{L}%
\TextAt{16}{305.79}{220.136}{\Character{106}}%
\TextAt{13}{312.432}{220.136}{K}%
\TextAt{8}{327.147}{220.136}{by}%
\TextAt{8}{343.406}{220.136}{the}%
\TextAt{8}{363.563}{220.136}{exact}%
\TextAt{8}{394.455}{220.136}{sequence}%
\TextAt{9}{219.214}{191.135}{0}%
\NativeArrowAt{228.388}{191.135}{11.955}{19.26}{\longrightarrow}%
\TextAt{14}{250.97}{191.135}{D}%
\TextAt{21}{260.912}{189.276}{L}%
\TextAt{3}{266.671}{189.276}{\Character{106}}%
\TextAt{21}{269.023}{189.276}{K}%
\NativeArrowAt{280.547}{191.135}{11.955}{19.26}{\longrightarrow}%
\TextAt{16}{303.118}{191.135}{\Character{79}}%
\NativeArrowAt{316.281}{191.135}{11.955}{19.26}{\longrightarrow}%
\TextAt{9}{338.881}{191.135}{\Character{10}}%
\TextAt{17}{347.336}{196.071}{1}%
\TextAt{3}{347.336}{188.18}{\Character{79}}%
\TextAt{3}{354.27}{188.18}{\Character{106}}%
\TextAt{27}{356.631}{188.18}{o}%
\NativeArrowAt{364.346}{191.135}{11.955}{19.26}{\longrightarrow}%
\TextAt{9}{386.928}{191.135}{0}%
\TextAt{8}{79.226}{162.134}{of}%
\TextAt{16}{93.135}{162.134}{\Character{79}}%
\TextAt{8}{102.985}{162.134}{-modules}%
\TextAt{8}{153.37}{162.134}{(with}%
\TextAt{8}{185.164}{162.134}{trivial}%
\TextAt{8}{220.544}{162.134}{metrics)}%
\TextAt{8}{266.709}{162.134}{(see}%
\TextAt{8}{290.756}{162.134}{§2,}%
\TextAt{8}{310.166}{162.134}{exercise}%
\TextAt{8}{354.047}{162.134}{3).}%
\TextAt{8}{374.646}{162.134}{This}%
\TextAt{8}{401.945}{162.134}{implies}%
\TextAt{8}{442.251}{162.134}{that}%
\TextAt{9}{468.199}{162.134}{\Character{91}\Character{10}}%
\TextAt{17}{479.905}{166.473}{1}%
\TextAt{3}{479.905}{158.582}{\Character{79}}%
\TextAt{3}{486.839}{158.582}{\Character{106}}%
\TextAt{27}{489.2}{158.582}{o}%
\TextAt{9}{493.594}{162.134}{\Character{93}}%
\TextAt{9}{501.138}{162.134}{\Character{61}}%
\TextAt{9}{514.54}{162.134}{1}%
\TextAt{16}{523.471}{162.134}{\Character{0}}%
\TextAt{9}{79.226}{145.088}{\Character{91}}%
\TextAt{14}{82.478}{145.088}{D}%
\TextAt{21}{92.419}{143.228}{L}%
\TextAt{3}{98.179}{143.228}{\Character{106}}%
\TextAt{21}{100.531}{143.228}{K}%
\TextAt{9}{108.734}{145.088}{\Character{93}}%
\TextAt{8}{111.985}{145.088}{.}%
\TextAt{8}{120.433}{145.088}{We}%
\TextAt{8}{140.578}{145.088}{may}%
\TextAt{8}{165.947}{145.088}{therefore}%
\TextAt{8}{215.047}{145.088}{describe}%
\TextAt{8}{260.297}{145.088}{the}%
\TextAt{8}{280.453}{145.088}{Todd}%
\TextAt{8}{311.011}{145.088}{class}%
\TextAt{8}{338.448}{145.088}{also}%
\TextAt{8}{361.928}{145.088}{by}%
\TextAt{8}{378.187}{145.088}{the}%
\TextAt{8}{398.343}{145.088}{different:}%
\TextAt{8}{184.775}{116.087}{Td}%
\TextAt{9}{199.731}{116.087}{\Character{40}}%
\TextAt{16}{204.283}{116.087}{\Character{79}}%
\TextAt{16}{217.446}{116.087}{\Character{106}}%
\TextAt{14}{224.095}{116.087}{o}%
\TextAt{9}{229.94}{116.087}{\Character{41}}%
\TextAt{9}{237.795}{116.087}{\Character{61}}%
\TextAt{9}{250.216}{116.087}{1}%
\TextAt{9}{258.716}{116.087}{\Character{43}}%
\TextAt{17}{271.705}{120.794}{1}%
\Vector{q 0 G 0 g 0.478 w 199.705 -654.924 m 203.939 -654.924 l S Q}%
\TextAt{17}{271.705}{111.964}{2}%
\TextAt{9}{277.134}{116.087}{\Character{40}\Character{91}}%
\TextAt{14}{284.938}{116.087}{D}%
\TextAt{21}{294.88}{114.227}{L}%
\TextAt{3}{300.639}{114.227}{\Character{106}}%
\TextAt{21}{302.992}{114.227}{K}%
\TextAt{9}{311.194}{116.087}{\Character{93}}%
\TextAt{16}{317.103}{116.087}{\Character{0}}%
\TextAt{9}{329.058}{116.087}{1\Character{41}}%
\TextAt{9}{342.759}{116.087}{\Character{61}}%
\TextAt{17}{356.405}{120.794}{1}%
\Vector{q 0 G 0 g 0.478 w 284.405 -654.924 m 288.639 -654.924 l S Q}%
\TextAt{17}{356.405}{111.964}{2}%
\TextAt{9}{361.834}{116.087}{\Character{40}1}%
\TextAt{9}{374.877}{116.087}{\Character{43}}%
\TextAt{9}{386.629}{116.087}{\Character{91}}%
\TextAt{14}{389.909}{116.087}{D}%
\TextAt{21}{399.851}{114.227}{L}%
\TextAt{3}{405.61}{114.227}{\Character{106}}%
\TextAt{21}{407.962}{114.227}{K}%
\TextAt{9}{416.165}{116.087}{\Character{93}\Character{41}}%
\TextAt{13}{423.969}{116.087}{\Character{58}}%
\TextAt{8}{79.226}{87.086}{The}%
\TextAt{8}{103.28}{87.086}{main}%
\TextAt{8}{132.546}{87.086}{result}%
\TextAt{8}{165.124}{87.086}{now}%
\TextAt{8}{189.513}{87.086}{follows}%
\TextAt{8}{227.948}{87.086}{from}%
\TextAt{8}{255.589}{87.086}{(7.3)}%
\TextAt{8}{283.564}{87.086}{using}%
\TextAt{8}{314.193}{87.086}{the}%
\TextAt{8}{334.349}{87.086}{Todd}%
\TextAt{8}{364.907}{87.086}{class.}%
\TextAt{8}{297.211}{43.033}{233}%
\Destination{78.23}{809.815}{nav.page.233}%
\Destination{79.36}{437.255}{nav.print.254}%
\Destination{79.36}{775.155}{nav.number.3.7.5}%
\Destination{79.36}{335.155}{nav.number.3.7.6}%
\Hotlink{295.935}{159.955}{6.858}{11.047}{nav.section.3.2}%
\Hotlink{380.826}{508.255}{25.078}{11.047}{nav.number.3.7.5}%
\Hotlink{255.089}{84.855}{25.078}{11.047}{nav.number.3.7.3}%
\Destination{79.226}{271.228}{errata.89}%
}

% Page 246 (body page 234)
\Page{%
\TextAt{2}{79.226}{796.943}{III.}%
\TextAt{2}{98.317}{796.943}{Riemann–Roch}%
\TextAt{2}{174.375}{796.943}{Theory}%
\Vector{q 0 G 0 g 0.299 w 7.226 18.729 m 460.769 18.729 l S Q}%
\TextAt{12}{79.226}{761.124}{Theorem}%
\TextAt{12}{136.24}{761.124}{(7.7)}%
\TextAt{12}{168.4}{761.124}{(Grothendieck–Riemann–Roch).}%
\TextAt{8}{361.954}{761.124}{The}%
\TextAt{8}{386.008}{761.124}{diagram}%
\TextAt{13}{239.476}{731.361}{K}%
\TextAt{17}{249.436}{729.568}{0}%
\TextAt{9}{254.168}{731.361}{\Character{40}}%
\NativeDiagram{%
\path[draw=black,line width=0.478bp] (186.72,-32.796) -- (196.57,-32.796);%
}%
\TextAt{16}{258.72}{731.361}{\Character{79}}%
\TextAt{9}{268.57}{731.361}{\Character{41}}%
\TextAt{8}{326.466}{731.361}{gr}%
\TextAt{13}{338.871}{731.361}{K}%
\TextAt{17}{348.831}{729.568}{0}%
\TextAt{9}{353.563}{731.361}{\Character{40}}%
\NativeDiagram{%
\path[draw=black,line width=0.478bp] (286.115,-32.796) -- (295.965,-32.796);%
}%
\TextAt{16}{358.115}{731.361}{\Character{79}}%
\TextAt{9}{367.965}{731.361}{\Character{41}}%
\ErrataAt{372.4}{731.361}{124.6}{$\otimes_{\mathbb Z}\mathbb Z[1/2]$}%
\TextAt{13}{241.479}{678.367}{K}%
\TextAt{17}{251.439}{676.574}{0}%
\TextAt{9}{256.171}{678.367}{\Character{40}}%
\NativeDiagram{%
\path[draw=black,line width=0.478bp] (188.723,-88.276) -- (194.567,-88.276);%
}%
\TextAt{14}{260.723}{678.367}{o}%
\TextAt{9}{266.567}{678.367}{\Character{41}}%
\TextAt{8}{328.469}{678.367}{gr}%
\TextAt{13}{340.874}{678.367}{K}%
\TextAt{17}{350.834}{676.574}{0}%
\TextAt{9}{355.566}{678.367}{\Character{40}}%
\NativeDiagram{%
\path[draw=black,line width=0.478bp] (288.118,-88.276) -- (293.962,-88.276);%
}%
\TextAt{14}{360.118}{678.367}{o}%
\TextAt{9}{365.962}{678.367}{\Character{41}}%
\NativeDiagram{%
\path[draw=black,line width=0.478bp] (206.475,-39.65) -- (248.637,-39.65);%
\path[draw=black,line width=0.478bp,line cap=round,line join=round] (246.39,-36.781) .. controls (246.844,-38.502) and (247.857,-39.315) .. (248.877,-39.65) .. controls (247.857,-39.985) and (246.844,-40.797) .. (246.39,-42.519);%
}%
\TextAt{7}{279.56}{739.158}{Td}%
\TextAt{17}{290.375}{739.158}{\Character{40}}%
\TextAt{3}{293.668}{739.158}{\Character{79}}%
\TextAt{3}{300.602}{739.158}{\Character{106}}%
\TextAt{27}{302.964}{739.158}{o}%
\TextAt{17}{306.86}{739.158}{\Character{41}}%
\TextAt{7}{311.565}{739.158}{ch}%
\NativeDiagram{%
\path[draw=black,line width=0.478bp] (184.3,-50.207) -- (184.3,-81.609);%
\path[draw=black,line width=0.478bp,line cap=round,line join=round] (187.169,-79.362) .. controls (185.447,-79.816) and (184.634,-80.828) .. (184.3,-81.848) .. controls (183.965,-80.828) and (183.152,-79.816) .. (181.431,-79.362);%
}%
\TextAt{21}{246.284}{705.717}{i}%
\TextAt{23}{249.167}{704.721}{\Character{3}}%
\NativeDiagram{%
\path[draw=black,line width=0.478bp] (277.493,-50.207) -- (277.493,-81.609);%
\path[draw=black,line width=0.478bp,line cap=round,line join=round] (280.363,-79.362) .. controls (278.641,-79.816) and (277.828,-80.828) .. (277.494,-81.848) .. controls (277.159,-80.828) and (276.346,-79.816) .. (274.624,-79.362);%
}%
\ErrataAt{352.29}{705.717}{47.71}{$\scriptstyle i_*\otimes1$}%
\NativeDiagram{%
\path[draw=black,line width=0.478bp] (204.472,-92.644) -- (250.64,-92.644);%
\path[draw=black,line width=0.478bp,line cap=round,line join=round] (248.393,-89.775) .. controls (248.847,-91.496) and (249.859,-92.309) .. (250.879,-92.644) .. controls (249.859,-92.979) and (248.847,-93.792) .. (248.393,-95.513);%
}%
\ErrataAt{292.2}{684.251}{33.8}{$\scriptstyle\mathrm{ch}_{1/2}$}%
\ErrataAt{370.4}{678.367}{126.6}{$\otimes_{\mathbb Z}\mathbb Z[1/2]$}%
\TextAt{8}{79.226}{650.806}{is}%
\TextAt{8}{90.99}{650.806}{commutative.}%
\TextAt{19}{79.226}{617.279}{Proof.}%
\TextAt{8}{116}{617.279}{For}%
\TextAt{13}{136.97}{617.279}{\Character{24}}%
\NativeMathAt{145.981}{617.279}{11.955}{\in}%
\TextAt{13}{157.272}{617.279}{K}%
\TextAt{17}{167.231}{615.486}{0}%
\TextAt{9}{171.963}{617.279}{\Character{40}}%
\Vector{q 0 G 0 g 0.478 w 104.516 -146.878 m 114.365 -146.878 l S Q}%
\TextAt{16}{176.516}{617.279}{\Character{79}}%
\TextAt{9}{186.365}{617.279}{\Character{41}}%
\TextAt{8}{190.918}{617.279}{,}%
\TextAt{8}{198.067}{617.279}{we}%
\TextAt{8}{215.294}{617.279}{have}%
\TextAt{8}{242.289}{617.279}{to}%
\TextAt{8}{256.599}{617.279}{show}%
\TextAt{8}{285.603}{617.279}{the}%
\TextAt{8}{305.759}{617.279}{identity}%
\TextAt{8}{230.867}{590.819}{ch}%
\TextAt{9}{242.571}{590.819}{\Character{40}}%
\TextAt{13}{247.123}{590.819}{i}%
\TextAt{3}{251.117}{589.025}{\Character{3}}%
\TextAt{13}{255.849}{590.819}{\Character{24}}%
\TextAt{9}{261.539}{590.819}{\Character{41}}%
\TextAt{9}{269.394}{590.819}{\Character{61}}%
\TextAt{13}{281.838}{590.819}{i}%
\TextAt{3}{285.831}{589.025}{\Character{3}}%
\TextAt{9}{290.564}{590.819}{\Character{40}}%
\TextAt{8}{295.116}{590.819}{Td}%
\TextAt{9}{310.072}{590.819}{\Character{40}}%
\TextAt{16}{314.624}{590.819}{\Character{79}}%
\TextAt{16}{327.787}{590.819}{\Character{106}}%
\TextAt{14}{334.437}{590.819}{o}%
\TextAt{9}{340.281}{590.819}{\Character{41}}%
\TextAt{8}{346.826}{590.819}{ch}%
\TextAt{9}{358.53}{590.819}{\Character{40}}%
\TextAt{13}{363.082}{590.819}{\Character{24}}%
\TextAt{9}{368.773}{590.819}{\Character{41}\Character{41}}%
\TextAt{13}{377.877}{590.819}{\Character{58}}%
\TextAt{8}{79.226}{564.358}{Decomposing}%
\TextAt{8}{150.993}{564.358}{ch}%
\TextAt{9}{162.697}{564.358}{\Character{40}}%
\TextAt{13}{167.249}{564.358}{i}%
\TextAt{3}{171.243}{562.565}{\Character{3}}%
\TextAt{13}{175.975}{564.358}{\Character{24}}%
\TextAt{9}{181.666}{564.358}{\Character{41}}%
\TextAt{9}{189.521}{564.358}{\Character{61}}%
\TextAt{8}{201.964}{564.358}{rk}%
\TextAt{9}{212.7}{564.358}{\Character{40}}%
\TextAt{13}{217.252}{564.358}{i}%
\TextAt{3}{221.246}{562.565}{\Character{3}}%
\TextAt{13}{225.978}{564.358}{\Character{24}}%
\TextAt{9}{231.669}{564.358}{\Character{41}}%
\TextAt{16}{238.878}{564.358}{\Character{8}}%
\TextAt{13}{250.833}{564.358}{c}%
\TextAt{17}{255.871}{562.565}{1}%
\TextAt{9}{260.603}{564.358}{\Character{40}}%
\TextAt{13}{265.155}{564.358}{i}%
\TextAt{3}{269.149}{562.565}{\Character{3}}%
\TextAt{13}{273.881}{564.358}{\Character{24}}%
\TextAt{9}{279.572}{564.358}{\Character{41}}%
\TextAt{8}{288.021}{564.358}{and}%
\TextAt{8}{310.784}{564.358}{ch}%
\TextAt{9}{322.488}{564.358}{\Character{40}}%
\TextAt{13}{327.04}{564.358}{\Character{24}}%
\TextAt{9}{332.731}{564.358}{\Character{41}}%
\TextAt{9}{340.586}{564.358}{\Character{61}}%
\TextAt{8}{353.029}{564.358}{rk}%
\TextAt{9}{363.765}{564.358}{\Character{40}}%
\TextAt{13}{368.317}{564.358}{\Character{24}}%
\TextAt{9}{374.008}{564.358}{\Character{41}}%
\TextAt{16}{381.217}{564.358}{\Character{8}}%
\TextAt{13}{393.172}{564.358}{c}%
\TextAt{17}{398.21}{562.565}{1}%
\TextAt{9}{402.942}{564.358}{\Character{40}}%
\TextAt{13}{407.495}{564.358}{\Character{24}}%
\TextAt{9}{413.185}{564.358}{\Character{41}}%
\TextAt{8}{421.635}{564.358}{and}%
\TextAt{8}{444.397}{564.358}{observing}%
\TextAt{8}{496.821}{564.358}{that}%
\TextAt{8}{166.438}{537.898}{Td}%
\TextAt{9}{181.394}{537.898}{\Character{40}}%
\TextAt{16}{185.946}{537.898}{\Character{79}}%
\TextAt{16}{199.109}{537.898}{\Character{106}}%
\TextAt{14}{205.759}{537.898}{o}%
\TextAt{9}{211.603}{537.898}{\Character{41}}%
\TextAt{8}{218.148}{537.898}{ch}%
\TextAt{9}{229.852}{537.898}{\Character{40}}%
\TextAt{13}{234.404}{537.898}{\Character{24}}%
\TextAt{9}{240.095}{537.898}{\Character{41}}%
\TextAt{9}{247.95}{537.898}{\Character{61}}%
\TextAt{20}{260.393}{547.581}{\Character{0}}%
\TextAt{9}{265.873}{537.898}{1}%
\TextAt{9}{274.373}{537.898}{\Character{43}}%
\TextAt{17}{287.339}{542.605}{1}%
\Vector{q 0 G 0 g 0.478 w 215.339 -233.114 m 219.573 -233.114 l S Q}%
\TextAt{17}{287.339}{533.775}{2}%
\TextAt{9}{292.769}{537.898}{\Character{40}\Character{91}}%
\TextAt{14}{300.573}{537.898}{D}%
\TextAt{21}{310.515}{536.038}{L}%
\TextAt{3}{316.274}{536.038}{\Character{106}}%
\TextAt{21}{318.626}{536.038}{K}%
\TextAt{9}{326.829}{537.898}{\Character{93}}%
\TextAt{16}{332.737}{537.898}{\Character{0}}%
\TextAt{9}{344.692}{537.898}{1\Character{41}}%
\TextAt{20}{355.097}{547.581}{\Character{1}}%
\TextAt{9}{360.577}{537.898}{\Character{40}}%
\TextAt{8}{365.129}{537.898}{rk}%
\TextAt{9}{375.865}{537.898}{\Character{40}}%
\TextAt{13}{380.417}{537.898}{\Character{24}}%
\TextAt{9}{386.108}{537.898}{\Character{41}}%
\TextAt{9}{393.305}{537.898}{\Character{43}}%
\TextAt{13}{405.078}{537.898}{c}%
\TextAt{17}{410.116}{536.104}{1}%
\TextAt{9}{414.848}{537.898}{\Character{40}}%
\TextAt{13}{419.401}{537.898}{\Character{24}}%
\TextAt{9}{425.091}{537.898}{\Character{41}\Character{41}}%
\TextAt{9}{247.968}{517.863}{\Character{61}}%
\TextAt{8}{260.393}{517.863}{rk}%
\TextAt{9}{271.129}{517.863}{\Character{40}}%
\TextAt{13}{275.681}{517.863}{\Character{24}}%
\TextAt{9}{281.372}{517.863}{\Character{41}}%
\TextAt{9}{288.569}{517.863}{\Character{43}}%
\TextAt{9}{300.321}{517.863}{\Character{91}}%
\TextAt{13}{303.594}{517.863}{c}%
\TextAt{17}{308.632}{516.07}{1}%
\TextAt{9}{313.364}{517.863}{\Character{40}}%
\TextAt{13}{317.916}{517.863}{\Character{24}}%
\TextAt{9}{323.607}{517.863}{\Character{41}}%
\TextAt{9}{330.804}{517.863}{\Character{43}}%
\TextAt{17}{343.773}{522.57}{1}%
\Vector{q 0 G 0 g 0.478 w 271.773 -253.148 m 276.007 -253.148 l S Q}%
\TextAt{17}{343.773}{513.74}{2}%
\TextAt{8}{351.195}{517.863}{rk}%
\TextAt{9}{361.931}{517.863}{\Character{40}}%
\TextAt{13}{366.483}{517.863}{\Character{24}}%
\TextAt{9}{372.173}{517.863}{\Character{41}\Character{40}\Character{91}}%
\TextAt{14}{384.53}{517.863}{D}%
\TextAt{21}{394.471}{516.003}{L}%
\TextAt{3}{400.23}{516.003}{\Character{106}}%
\TextAt{21}{402.583}{516.003}{K}%
\TextAt{9}{410.785}{517.863}{\Character{93}}%
\TextAt{16}{416.694}{517.863}{\Character{0}}%
\TextAt{9}{428.649}{517.863}{1\Character{41}\Character{93}}%
\TextAt{13}{442.306}{517.863}{\Character{59}}%
\TextAt{8}{79.226}{491.402}{it}%
\TextAt{8}{90.93}{491.402}{suffices}%
\TextAt{8}{130.717}{491.402}{to}%
\TextAt{8}{145.027}{491.402}{check}%
\TextAt{8}{176.565}{491.402}{the}%
\TextAt{8}{196.722}{491.402}{equations}%
\TextAt{8}{217.847}{464.942}{rk}%
\TextAt{9}{228.583}{464.942}{\Character{40}}%
\TextAt{13}{233.135}{464.942}{i}%
\TextAt{3}{237.129}{463.148}{\Character{3}}%
\TextAt{13}{241.861}{464.942}{\Character{24}}%
\TextAt{9}{247.551}{464.942}{\Character{41}}%
\TextAt{9}{255.406}{464.942}{\Character{61}}%
\TextAt{8}{267.85}{464.942}{rk}%
\TextAt{9}{278.586}{464.942}{\Character{40}}%
\TextAt{13}{283.138}{464.942}{\Character{24}}%
\TextAt{9}{288.829}{464.942}{\Character{41}}%
\TextAt{8}{295.373}{464.942}{rk}%
\TextAt{9}{306.109}{464.942}{\Character{40}}%
\TextAt{13}{310.661}{464.942}{i}%
\TextAt{3}{314.655}{463.148}{\Character{3}}%
\TextAt{9}{319.387}{464.942}{\Character{91}}%
\TextAt{16}{322.639}{464.942}{\Character{79}}%
\TextAt{9}{332.489}{464.942}{\Character{93}\Character{41}}%
\TextAt{13}{340.293}{464.942}{\Character{59}}%
\TextAt{8}{517.802}{464.942}{(a)}%
\TextAt{13}{218.813}{444.907}{c}%
\TextAt{17}{223.851}{443.114}{1}%
\TextAt{9}{228.583}{444.907}{\Character{40}}%
\TextAt{13}{233.135}{444.907}{i}%
\TextAt{3}{237.129}{443.114}{\Character{3}}%
\TextAt{13}{241.861}{444.907}{\Character{24}}%
\TextAt{9}{247.551}{444.907}{\Character{41}}%
\TextAt{9}{255.406}{444.907}{\Character{61}}%
\TextAt{13}{267.85}{444.907}{i}%
\TextAt{3}{271.843}{443.114}{\Character{3}}%
\TextAt{9}{276.576}{444.907}{\Character{40}}%
\TextAt{13}{281.128}{444.907}{c}%
\TextAt{17}{286.166}{443.114}{1}%
\TextAt{9}{290.898}{444.907}{\Character{40}}%
\TextAt{13}{295.451}{444.907}{\Character{24}}%
\TextAt{9}{301.141}{444.907}{\Character{41}\Character{41}}%
\TextAt{9}{312.881}{444.907}{\Character{43}}%
\TextAt{8}{324.664}{444.907}{rk}%
\TextAt{9}{335.4}{444.907}{\Character{40}}%
\TextAt{13}{339.952}{444.907}{\Character{24}}%
\TextAt{9}{345.642}{444.907}{\Character{41}}%
\TextAt{13}{350.195}{444.907}{c}%
\TextAt{17}{355.233}{443.114}{1}%
\TextAt{9}{359.965}{444.907}{\Character{40}}%
\TextAt{13}{364.517}{444.907}{i}%
\TextAt{3}{368.511}{443.114}{\Character{3}}%
\TextAt{9}{373.243}{444.907}{\Character{91}}%
\TextAt{16}{376.495}{444.907}{\Character{79}}%
\TextAt{9}{386.344}{444.907}{\Character{93}\Character{41}}%
\TextAt{8}{517.156}{444.907}{(b)}%
\TextAt{8}{79.226}{418.446}{and}%
\TextAt{8}{242.582}{391.986}{rk}%
\TextAt{9}{253.317}{391.986}{\Character{40}}%
\TextAt{13}{257.87}{391.986}{i}%
\TextAt{3}{261.863}{390.192}{\Character{3}}%
\TextAt{9}{266.596}{391.986}{\Character{91}}%
\TextAt{16}{269.847}{391.986}{\Character{79}}%
\TextAt{9}{279.697}{391.986}{\Character{93}\Character{41}}%
\TextAt{9}{290.791}{391.986}{\Character{61}}%
\TextAt{13}{303.247}{391.986}{i}%
\TextAt{3}{307.241}{390.192}{\Character{3}}%
\TextAt{9}{311.973}{391.986}{\Character{40}1\Character{41}}%
\TextAt{13}{326.931}{391.986}{\Character{59}}%
\TextAt{8}{518.459}{391.986}{(c)}%
\TextAt{9}{237.694}{371.951}{2}%
\TextAt{13}{243.547}{371.951}{c}%
\TextAt{17}{248.585}{370.158}{1}%
\TextAt{9}{253.317}{371.951}{\Character{40}}%
\TextAt{13}{257.87}{371.951}{i}%
\TextAt{3}{261.863}{370.158}{\Character{3}}%
\TextAt{9}{266.596}{371.951}{\Character{91}}%
\TextAt{16}{269.847}{371.951}{\Character{79}}%
\TextAt{9}{279.697}{371.951}{\Character{93}\Character{41}}%
\TextAt{9}{290.791}{371.951}{\Character{61}}%
\TextAt{13}{303.247}{371.951}{i}%
\TextAt{3}{307.241}{370.158}{\Character{3}}%
\TextAt{9}{311.973}{371.951}{\Character{40}\Character{91}}%
\TextAt{14}{319.777}{371.951}{D}%
\TextAt{21}{329.719}{370.091}{L}%
\TextAt{3}{335.478}{370.091}{\Character{106}}%
\TextAt{21}{337.83}{370.091}{K}%
\TextAt{9}{346.033}{371.951}{\Character{93}}%
\TextAt{16}{351.941}{371.951}{\Character{0}}%
\TextAt{9}{363.896}{371.951}{1\Character{41}}%
\TextAt{8}{517.156}{371.951}{(d)}%
\TextAt{8}{79.226}{345.49}{in}%
\TextAt{8}{92.465}{345.49}{gr}%
\TextAt{13}{104.87}{345.49}{K}%
\TextAt{17}{114.829}{343.697}{0}%
\TextAt{9}{119.562}{345.49}{\Character{40}}%
\Vector{q 0 G 0 g 0.478 w 52.114 -421.153 m 57.958 -421.153 l S Q}%
\TextAt{14}{124.114}{345.49}{o}%
\TextAt{9}{129.958}{345.49}{\Character{41}}%
\TextAt{8}{134.51}{345.49}{.}%
\TextAt{8}{142.821}{345.49}{The}%
\TextAt{8}{166.46}{345.49}{equations}%
\TextAt{8}{218.469}{345.49}{(a)}%
\TextAt{8}{236.921}{345.49}{and}%
\TextAt{8}{259.269}{345.49}{(c)}%
\TextAt{8}{277.062}{345.49}{are}%
\TextAt{8}{296.159}{345.49}{clear}%
\TextAt{8}{323.708}{345.49}{because}%
\TextAt{8}{366.595}{345.49}{of}%
\TextAt{8}{379.511}{345.49}{rk}%
\TextAt{9}{390.247}{345.49}{\Character{40}}%
\TextAt{13}{394.799}{345.49}{i}%
\TextAt{3}{398.793}{343.697}{\Character{3}}%
\TextAt{9}{403.525}{345.49}{\Character{91}}%
\TextAt{16}{406.777}{345.49}{\Character{79}}%
\TextAt{9}{416.626}{345.49}{\Character{93}\Character{41}}%
\TextAt{9}{427.72}{345.49}{\Character{61}}%
\TextAt{8}{440.177}{345.49}{rk}%
\TextAt{9}{450.912}{345.49}{\Character{40}}%
\TextAt{13}{455.465}{345.49}{i}%
\TextAt{3}{459.458}{343.697}{\Character{3}}%
\TextAt{16}{464.19}{345.49}{\Character{79}}%
\TextAt{9}{474.04}{345.49}{\Character{41}}%
\TextAt{9}{481.895}{345.49}{\Character{61}}%
\TextAt{9}{494.316}{345.49}{\Character{91}}%
\TextAt{13}{497.591}{345.49}{L}%
\TextAt{9}{508.876}{345.49}{\Character{58}}%
\TextAt{13}{515.448}{345.49}{K}%
\TextAt{9}{526.266}{345.49}{\Character{93}}%
\TextAt{8}{529.518}{345.49}{.}%
\TextAt{8}{79.226}{328.444}{To}%
\TextAt{8}{97.339}{328.444}{show}%
\TextAt{8}{127.227}{328.444}{(b)}%
\TextAt{8}{147.623}{328.444}{and}%
\TextAt{8}{171.27}{328.444}{(d),}%
\TextAt{8}{195.139}{328.444}{we}%
\TextAt{8}{213.252}{328.444}{apply}%
\TextAt{8}{246.332}{328.444}{det}%
\TextAt{8}{267.374}{328.444}{to}%
\TextAt{8}{282.569}{328.444}{both}%
\TextAt{8}{311.094}{328.444}{sides}%
\TextAt{8}{340.062}{328.444}{and}%
\TextAt{8}{363.71}{328.444}{are}%
\TextAt{8}{384.106}{328.444}{reduced}%
\TextAt{8}{428.556}{328.444}{by}%
\TextAt{8}{445.7}{328.444}{the}%
\TextAt{8}{466.741}{328.444}{commutative}%
\TextAt{8}{79.226}{311.398}{diagram}%
\TextAt{8}{124.763}{311.398}{(7.5)}%
\TextAt{8}{152.738}{311.398}{to}%
\TextAt{8}{167.049}{311.398}{the}%
\TextAt{8}{187.205}{311.398}{equations}%
\TextAt{8}{224.867}{284.938}{det}%
\TextAt{9}{241.126}{284.938}{\Character{40}}%
\TextAt{13}{245.679}{284.938}{i}%
\TextAt{3}{249.672}{283.145}{\Character{3}}%
\TextAt{13}{254.404}{284.938}{\Character{24}}%
\TextAt{9}{260.095}{284.938}{\Character{41}}%
\TextAt{9}{267.95}{284.938}{\Character{61}}%
\TextAt{13}{280.393}{284.938}{N}%
\TextAt{21}{289.767}{283.078}{L}%
\TextAt{3}{295.526}{283.078}{\Character{106}}%
\TextAt{21}{297.879}{283.078}{K}%
\TextAt{9}{306.081}{284.938}{\Character{40}}%
\TextAt{8}{310.634}{284.938}{det}%
\TextAt{13}{328.885}{284.938}{\Character{24}}%
\TextAt{9}{334.576}{284.938}{\Character{41}\Character{91}}%
\TextAt{8}{342.38}{284.938}{det}%
\TextAt{13}{360.631}{284.938}{i}%
\TextAt{3}{364.625}{283.145}{\Character{3}}%
\TextAt{16}{369.357}{284.938}{\Character{79}}%
\TextAt{9}{379.207}{284.938}{\Character{93}}%
\TextAt{7}{382.458}{289.874}{rk}%
\TextAt{17}{390.229}{289.874}{\Character{40}}%
\TextAt{21}{393.522}{289.874}{\Character{24}}%
\TextAt{17}{397.556}{289.874}{\Character{41}}%
\TextAt{13}{401.347}{284.938}{\Character{59}}%
\TextAt{8}{518.459}{284.938}{(e)}%
\TextAt{9}{207.397}{264.903}{\Character{40}}%
\TextAt{8}{211.949}{264.903}{det}%
\TextAt{13}{230.2}{264.903}{i}%
\TextAt{3}{234.194}{263.11}{\Character{3}}%
\TextAt{16}{238.926}{264.903}{\Character{79}}%
\TextAt{9}{248.776}{264.903}{\Character{41}}%
\TextAt{3}{253.328}{269.839}{\Character{10}}%
\TextAt{17}{259.915}{269.839}{2}%
\TextAt{9}{267.968}{264.903}{\Character{61}}%
\TextAt{13}{280.393}{264.903}{N}%
\TextAt{21}{289.767}{263.043}{L}%
\TextAt{3}{295.526}{263.043}{\Character{106}}%
\TextAt{21}{297.879}{263.043}{K}%
\TextAt{9}{306.081}{264.903}{\Character{40}}%
\TextAt{8}{310.634}{264.903}{det}%
\TextAt{14}{328.885}{264.903}{D}%
\TextAt{21}{338.827}{263.043}{L}%
\TextAt{3}{344.586}{263.043}{\Character{106}}%
\TextAt{21}{346.938}{263.043}{K}%
\TextAt{9}{355.141}{264.903}{\Character{41}}%
\TextAt{13}{359.693}{264.903}{\Character{58}}%
\TextAt{8}{519.571}{264.903}{(f)}%
\TextAt{8}{79.226}{238.443}{But}%
\TextAt{8}{102.467}{238.443}{(e)}%
\TextAt{8}{120.675}{238.443}{is}%
\TextAt{8}{132.439}{238.443}{the}%
\TextAt{8}{152.595}{238.443}{second}%
\TextAt{8}{190.373}{238.443}{identity}%
\TextAt{8}{233.627}{238.443}{of}%
\TextAt{8}{246.957}{238.443}{(7.3),}%
\TextAt{8}{278.184}{238.443}{and}%
\TextAt{8}{300.947}{238.443}{(f)}%
\TextAt{8}{317.851}{238.443}{follows}%
\TextAt{8}{356.287}{238.443}{from}%
\TextAt{8}{383.928}{238.443}{(7.4)}%
\TextAt{8}{411.903}{238.443}{and}%
\TextAt{8}{434.665}{238.443}{(2.9).}%
\Vector{q 0 G 0 g 0.399 w 452.334 -535.557 m 452.334 -527.488 l S Q}%
\Vector{q 0 G 0 g 0.399 w 452.534 -527.687 m 459.707 -527.687 l S Q}%
\Vector{q 0 G 0 g 0.399 w 452.534 -535.358 m 459.707 -535.358 l S Q}%
\Vector{q 0 G 0 g 0.399 w 459.906 -535.557 m 459.906 -527.488 l S Q}%
\TextAt{8}{97.159}{210.167}{With}%
\TextAt{8}{127.124}{210.167}{this}%
\TextAt{8}{149.677}{210.167}{final}%
\TextAt{8}{175.422}{210.167}{theorem,}%
\TextAt{8}{223.983}{210.167}{the}%
\TextAt{8}{243.87}{210.167}{theory}%
\TextAt{8}{280.351}{210.167}{of}%
\TextAt{8}{293.412}{210.167}{algebraic}%
\TextAt{8}{342.577}{210.167}{integers}%
\TextAt{8}{385.621}{210.167}{can}%
\TextAt{8}{406.811}{210.167}{be}%
\TextAt{8}{422.466}{210.167}{integrated}%
\TextAt{8}{477.812}{210.167}{completely}%
\TextAt{8}{79.226}{193.121}{into}%
\TextAt{8}{103.023}{193.121}{a}%
\TextAt{8}{112.832}{193.121}{general}%
\TextAt{8}{153.21}{193.121}{programme}%
\TextAt{8}{215.06}{193.121}{of}%
\TextAt{8}{228.444}{193.121}{algebraic}%
\TextAt{8}{277.932}{193.121}{geometry}%
\TextAt{8}{329.047}{193.121}{as}%
\TextAt{8}{343.47}{193.121}{a}%
\TextAt{8}{353.279}{193.121}{special}%
\TextAt{8}{391.434}{193.121}{case.}%
\TextAt{8}{420.917}{193.121}{What}%
\TextAt{8}{453.812}{193.121}{is}%
\TextAt{8}{465.629}{193.121}{needed}%
\TextAt{8}{504.693}{193.121}{is}%
\TextAt{8}{516.51}{193.121}{the}%
\TextAt{8}{79.226}{176.075}{use}%
\TextAt{8}{98.891}{176.075}{of}%
\TextAt{8}{111.671}{176.075}{the}%
\TextAt{8}{131.276}{176.075}{geometric}%
\TextAt{8}{184.057}{176.075}{language}%
\TextAt{8}{232.295}{176.075}{for}%
\TextAt{8}{249.629}{176.075}{the}%
\TextAt{8}{269.235}{176.075}{objects}%
\TextAt{8}{308.411}{176.075}{considered.}%
\TextAt{8}{370.067}{176.075}{Thus}%
\TextAt{8}{399.165}{176.075}{the}%
\TextAt{8}{418.771}{176.075}{ring}%
\TextAt{14}{442.286}{176.075}{o}%
\TextAt{8}{451.476}{176.075}{is}%
\TextAt{8}{462.689}{176.075}{interpreted}%
\TextAt{8}{522.297}{176.075}{as}%
\TextAt{8}{79.226}{159.029}{the}%
\TextAt{8}{99.769}{159.029}{scheme}%
\TextAt{13}{140.204}{159.029}{X}%
\TextAt{9}{154.838}{159.029}{\Character{61}}%
\TextAt{8}{167.921}{159.029}{Spec}%
\TextAt{9}{191.33}{159.029}{\Character{40}}%
\TextAt{14}{195.882}{159.029}{o}%
\TextAt{9}{201.726}{159.029}{\Character{41}}%
\TextAt{8}{206.279}{159.029}{,}%
\TextAt{8}{213.91}{159.029}{the}%
\TextAt{8}{234.453}{159.029}{projective}%
\TextAt{8}{288.816}{159.029}{metrized}%
\TextAt{14}{337.322}{159.029}{o}%
\TextAt{8}{343.166}{159.029}{-modules}%
\TextAt{8}{393.357}{159.029}{as}%
\TextAt{8}{408.113}{159.029}{metrized}%
\TextAt{19}{456.619}{159.029}{vector}%
\TextAt{19}{491.56}{159.029}{bundles}%
\TextAt{8}{529.518}{159.029}{,}%
\TextAt{8}{79.226}{141.983}{the}%
\TextAt{8}{100.025}{141.983}{invertible}%
\TextAt{14}{152.374}{141.983}{o}%
\TextAt{8}{158.219}{141.983}{-modules}%
\TextAt{8}{208.667}{141.983}{as}%
\TextAt{19}{223.68}{141.983}{line}%
\TextAt{19}{247.085}{141.983}{bundles}%
\TextAt{8}{285.042}{141.983}{,}%
\TextAt{8}{292.995}{141.983}{the}%
\TextAt{8}{313.794}{141.983}{inclusion}%
\TextAt{13}{363.274}{141.983}{i}%
\TextAt{9}{371.684}{141.983}{\Character{58}}%
\TextAt{14}{379.352}{141.983}{o}%
\NativeMathAt{389.612}{141.983}{11.955}{\rightarrow}%
\TextAt{16}{405.979}{141.983}{\Character{79}}%
\TextAt{8}{420.374}{141.983}{as}%
\TextAt{8}{435.387}{141.983}{morphism}%
\TextAt{13}{490.725}{141.983}{f}%
\TextAt{9}{502.187}{141.983}{\Character{58}}%
\TextAt{13}{509.855}{141.983}{Y}%
\TextAt{9}{523.665}{141.983}{\Character{61}}%
\TextAt{8}{79.226}{124.937}{Spec}%
\TextAt{9}{102.634}{124.937}{\Character{40}}%
\TextAt{16}{107.187}{124.937}{\Character{79}}%
\TextAt{9}{117.036}{124.937}{\Character{41}}%
\NativeMathAt{125.555}{124.937}{11.955}{\rightarrow}%
\TextAt{13}{141.476}{124.937}{X}%
\TextAt{8}{156.407}{124.937}{of}%
\TextAt{8}{170.116}{124.937}{schemes,}%
\TextAt{8}{218.505}{124.937}{the}%
\TextAt{8}{239.041}{124.937}{class}%
\TextAt{9}{266.856}{124.937}{\Character{10}}%
\TextAt{17}{275.311}{129.276}{1}%
\TextAt{3}{275.311}{121.386}{\Character{79}}%
\TextAt{3}{282.245}{121.386}{\Character{106}}%
\TextAt{27}{284.606}{121.386}{o}%
\TextAt{8}{293.276}{124.937}{as}%
\TextAt{8}{308.025}{124.937}{the}%
\TextAt{19}{328.56}{124.937}{cotangent}%
\TextAt{19}{381.569}{124.937}{element}%
\TextAt{8}{420.722}{124.937}{,}%
\TextAt{8}{428.345}{124.937}{etc.}%
\TextAt{8}{452.885}{124.937}{In}%
\TextAt{8}{467.885}{124.937}{this}%
\TextAt{8}{491.086}{124.937}{way}%
\TextAt{8}{515.207}{124.937}{one}%
\TextAt{8}{79.226}{107.891}{realizes}%
\TextAt{8}{119.871}{107.891}{in}%
\TextAt{8}{133.138}{107.891}{the}%
\TextAt{8}{152.908}{107.891}{present}%
\TextAt{8}{193.23}{107.891}{context}%
\TextAt{8}{234.472}{107.891}{the}%
\TextAt{8}{254.243}{107.891}{old}%
\TextAt{8}{273.368}{107.891}{idea}%
\TextAt{8}{297.693}{107.891}{of}%
\TextAt{8}{310.638}{107.891}{viewing}%
\TextAt{8}{352.848}{107.891}{number}%
\TextAt{8}{395.059}{107.891}{theory}%
\TextAt{8}{431.423}{107.891}{as}%
\TextAt{8}{445.407}{107.891}{part}%
\TextAt{8}{470.391}{107.891}{of}%
\TextAt{8}{483.335}{107.891}{geometry.}%
\TextAt{8}{297.211}{43.033}{234}%
\Destination{522.36}{484.975}{nav.AMS.42}%
\Destination{521.71}{464.945}{nav.AMS.43}%
\Destination{523.01}{412.025}{nav.AMS.46}%
\Destination{521.71}{391.985}{nav.AMS.47}%
\Destination{523.01}{304.975}{nav.AMS.50}%
\Destination{524.13}{284.935}{nav.AMS.51}%
\Destination{78.23}{809.815}{nav.page.234}%
\Destination{79.36}{579.355}{nav.print.255}%
\Destination{79.36}{775.155}{nav.number.3.7.7}%
\Hotlink{124.263}{309.155}{25.078}{11.047}{nav.number.3.7.5}%
\Hotlink{246.457}{236.255}{25.078}{11.047}{nav.number.3.7.3}%
\Hotlink{383.428}{236.255}{25.078}{11.047}{nav.number.3.7.4}%
\Hotlink{434.165}{236.255}{25.078}{11.047}{nav.number.3.2.9}%
\Hotlink{214.864}{70.755}{21.483}{9.572}{nav.number.3.7.6}%
\Destination{79.226}{748.361}{errata.86}%
}

% Page 247 (body page 235)
\Page{%
\TextAt{2}{79.226}{796.943}{III.}%
\TextAt{2}{98.317}{796.943}{Riemann–Roch}%
\TextAt{2}{174.375}{796.943}{Theory}%
\Vector{q 0 G 0 g 0.299 w 7.226 18.729 m 460.769 18.729 l S Q}%
\TextAt{15}{79.226}{760.804}{§}%
\TextAt{15}{93.785}{760.804}{8}%
\TextAt{15}{120.693}{760.804}{The}%
\TextAt{15}{160.202}{760.804}{Euler–Minkowski}%
\TextAt{15}{312.576}{760.804}{Characteristic}%
\TextAt{8}{79.226}{728.918}{Considering}%
\TextAt{8}{143.35}{728.918}{the}%
\TextAt{8}{163.181}{728.918}{theorem}%
\TextAt{8}{208.38}{728.918}{of}%
\TextAt{8}{221.384}{728.918}{Grothendieck–Riemann–Roch}%
\TextAt{8}{376.917}{728.918}{in}%
\TextAt{8}{390.244}{728.918}{the}%
\TextAt{8}{410.075}{728.918}{special}%
\TextAt{8}{447.85}{728.918}{case}%
\TextAt{8}{472.295}{728.918}{of}%
\TextAt{8}{485.299}{728.918}{an}%
\TextAt{8}{501.232}{728.918}{exten-}%
\TextAt{8}{79.226}{711.872}{sion}%
\TextAt{13}{103.332}{711.872}{K}%
\TextAt{16}{117.47}{711.872}{\Character{106}}%
\TextAt{18}{124.112}{711.872}{Q}%
\TextAt{8}{133.411}{711.872}{,}%
\TextAt{8}{140.544}{711.872}{amounts}%
\TextAt{8}{187.747}{711.872}{to}%
\TextAt{8}{202.037}{711.872}{revisiting}%
\TextAt{8}{253.138}{711.872}{the}%
\TextAt{8}{273.274}{711.872}{Riemann–Roch}%
\TextAt{8}{354.884}{711.872}{theory}%
\TextAt{8}{391.615}{711.872}{of}%
\TextAt{8}{404.925}{711.872}{§3}%
\TextAt{8}{420.339}{711.872}{from}%
\TextAt{8}{447.959}{711.872}{our}%
\TextAt{8}{468.753}{711.872}{new}%
\TextAt{8}{492.787}{711.872}{point}%
\TextAt{8}{523.337}{711.872}{of}%
\TextAt{8}{79.226}{694.826}{view.}%
\TextAt{8}{110.76}{694.826}{At}%
\TextAt{8}{127.665}{694.826}{the}%
\TextAt{8}{147.821}{694.826}{center}%
\TextAt{8}{182.61}{694.826}{of}%
\TextAt{8}{195.94}{694.826}{that}%
\TextAt{8}{221.309}{694.826}{theory}%
\TextAt{8}{258.06}{694.826}{was}%
\TextAt{8}{280.559}{694.826}{the}%
\TextAt{8}{300.716}{694.826}{Euler–Minkowski}%
\TextAt{8}{392.316}{694.826}{characteristic}%
\TextAt{13}{257.599}{665.825}{\Character{31}}%
\TextAt{9}{264.933}{665.825}{\Character{40}}%
\TextAt{14}{269.485}{665.825}{a}%
\TextAt{9}{275.468}{665.825}{\Character{41}}%
\TextAt{9}{283.323}{665.825}{\Character{61}}%
\TextAt{16}{295.767}{665.825}{\Character{0}}%
\TextAt{8}{307.058}{665.825}{log}%
\TextAt{8}{324.011}{665.825}{vol}%
\TextAt{9}{339.309}{665.825}{\Character{40}}%
\TextAt{14}{343.861}{665.825}{a}%
\TextAt{9}{349.844}{665.825}{\Character{41}}%
\TextAt{8}{79.226}{636.824}{of}%
\TextAt{8}{92.885}{636.824}{replete}%
\TextAt{8}{131.577}{636.824}{ideals}%
\TextAt{14}{164.484}{636.824}{a}%
\TextAt{8}{174.693}{636.824}{of}%
\TextAt{13}{188.351}{636.824}{K}%
\TextAt{8}{199.169}{636.824}{.}%
\TextAt{8}{208.603}{636.824}{Here,}%
\TextAt{8}{239.893}{636.824}{vol}%
\TextAt{9}{255.184}{636.824}{\Character{40}}%
\TextAt{14}{259.736}{636.824}{a}%
\TextAt{9}{265.719}{636.824}{\Character{41}}%
\TextAt{8}{274.498}{636.824}{was}%
\TextAt{8}{297.326}{636.824}{the}%
\TextAt{19}{317.811}{636.824}{canonical}%
\TextAt{19}{369.571}{636.824}{measure}%
\TextAt{8}{414.982}{636.824}{of}%
\TextAt{8}{428.641}{636.824}{a}%
\TextAt{8}{438.724}{636.824}{fundamental}%
\TextAt{8}{506.695}{636.824}{mesh}%
\TextAt{8}{79.226}{619.778}{of}%
\TextAt{8}{92.847}{619.778}{the}%
\TextAt{8}{113.294}{619.778}{lattice}%
\TextAt{8}{149.354}{619.778}{in}%
\TextAt{8}{163.298}{619.778}{Minkowski}%
\TextAt{8}{221.858}{619.778}{space}%
\TextAt{13}{253.423}{619.778}{K}%
\TextAt{25}{263.382}{617.985}{R}%
\TextAt{9}{273.452}{619.778}{\Character{61}}%
\TextAt{14}{286.373}{619.778}{a}%
\TextAt{16}{295.211}{619.778}{\Character{10}}%
\TextAt{25}{304.509}{617.985}{Z}%
\TextAt{18}{313.176}{619.778}{R}%
\TextAt{8}{325.998}{619.778}{defined}%
\TextAt{8}{366.601}{619.778}{by}%
\TextAt{14}{383.151}{619.778}{a}%
\TextAt{8}{389.134}{619.778}{.}%
\TextAt{8}{398.454}{619.778}{This}%
\TextAt{8}{425.465}{619.778}{definition}%
\TextAt{8}{477.784}{619.778}{is}%
\TextAt{8}{489.839}{619.778}{properly}%
\TextAt{8}{79.226}{602.732}{explained}%
\TextAt{8}{132.394}{602.732}{in}%
\TextAt{8}{146.863}{602.732}{the}%
\TextAt{8}{167.836}{602.732}{theory}%
\TextAt{8}{205.403}{602.732}{of}%
\TextAt{8}{219.549}{602.732}{metrized}%
\TextAt{8}{268.485}{602.732}{modules}%
\TextAt{8}{315.209}{602.732}{of}%
\TextAt{8}{329.355}{602.732}{higher}%
\TextAt{8}{365.942}{602.732}{rank.}%
\TextAt{8}{399.936}{602.732}{More}%
\TextAt{8}{430.987}{602.732}{precisely,}%
\TextAt{8}{482.136}{602.732}{instead}%
\TextAt{8}{523.337}{602.732}{of}%
\TextAt{8}{79.226}{585.686}{considering}%
\TextAt{14}{141.57}{585.686}{a}%
\TextAt{8}{152.597}{585.686}{as}%
\TextAt{8}{168.113}{585.686}{a}%
\TextAt{8}{179.014}{585.686}{metrized}%
\TextAt{14}{228.279}{585.686}{o}%
\TextAt{8}{234.124}{585.686}{-module}%
\TextAt{8}{280.46}{585.686}{of}%
\TextAt{8}{294.936}{585.686}{rank}%
\TextAt{9}{323.076}{585.686}{1}%
\TextAt{8}{328.929}{585.686}{,}%
\TextAt{8}{337.511}{585.686}{it}%
\TextAt{8}{350.361}{585.686}{should}%
\TextAt{8}{388.639}{585.686}{be}%
\TextAt{8}{405.709}{585.686}{viewed}%
\TextAt{8}{445.219}{585.686}{as}%
\TextAt{8}{460.735}{585.686}{a}%
\TextAt{8}{471.637}{585.686}{metrized}%
\TextAt{18}{520.902}{585.686}{Z}%
\TextAt{8}{528.872}{585.686}{-}%
\TextAt{8}{79.226}{568.64}{module}%
\TextAt{8}{120.119}{568.64}{of}%
\TextAt{8}{133.048}{568.64}{rank}%
\TextAt{9}{159.642}{568.64}{\Character{91}}%
\TextAt{13}{162.894}{568.64}{K}%
\TextAt{9}{177.032}{568.64}{\Character{58}}%
\TextAt{18}{183.605}{568.64}{Q}%
\TextAt{9}{192.903}{568.64}{\Character{93}}%
\TextAt{8}{196.155}{568.64}{.}%
\TextAt{8}{204.47}{568.64}{This}%
\TextAt{8}{230.789}{568.64}{point}%
\TextAt{8}{260.958}{568.64}{of}%
\TextAt{8}{273.887}{568.64}{view}%
\TextAt{8}{300.469}{568.64}{leads}%
\TextAt{8}{329.395}{568.64}{us}%
\TextAt{8}{344.01}{568.64}{necessarily}%
\TextAt{8}{401.938}{568.64}{to}%
\TextAt{8}{415.848}{568.64}{the}%
\TextAt{8}{435.604}{568.64}{following}%
\TextAt{8}{484.638}{568.64}{definition}%
\TextAt{8}{79.226}{551.594}{of}%
\TextAt{8}{92.556}{551.594}{the}%
\TextAt{8}{112.713}{551.594}{Euler–Minkowski}%
\TextAt{8}{204.313}{551.594}{characteristic.}%
\TextAt{12}{79.226}{525.582}{Proposition}%
\TextAt{12}{151.937}{525.582}{(8.1).}%
\TextAt{8}{187.253}{525.582}{The}%
\TextAt{8}{211.307}{525.582}{degree}%
\TextAt{8}{247.722}{525.582}{map}%
\TextAt{8}{180.821}{496.581}{deg}%
\TextAt{21}{198.383}{493.632}{K}%
\TextAt{9}{209.906}{496.581}{\Character{58}}%
\TextAt{8}{216.479}{496.581}{Pic}%
\TextAt{9}{232.893}{496.581}{\Character{40}}%
\Vector{q 0 G 0 g 0.478 w 165.446 -270.063 m 171.29 -270.063 l S Q}%
\TextAt{14}{237.446}{496.581}{o}%
\TextAt{9}{243.29}{496.581}{\Character{41}}%
\NativeArrowAt{251.163}{496.581}{11.955}{19.26}{\longrightarrow}%
\TextAt{18}{273.745}{496.581}{R}%
\TextAt{13}{282.379}{496.581}{\Character{59}}%
\TextAt{8}{311.534}{496.581}{deg}%
\TextAt{21}{329.096}{493.632}{K}%
\TextAt{9}{337.299}{496.581}{\Character{40}\Character{91}}%
\TextAt{14}{345.103}{496.581}{a}%
\TextAt{9}{351.086}{496.581}{\Character{93}\Character{41}}%
\TextAt{9}{362.18}{496.581}{\Character{61}}%
\TextAt{16}{374.636}{496.581}{\Character{0}}%
\TextAt{8}{385.927}{496.581}{log}%
\TextAt{14}{402.887}{496.581}{N}%
\TextAt{9}{412.836}{496.581}{\Character{40}}%
\TextAt{14}{417.388}{496.581}{a}%
\TextAt{9}{423.371}{496.581}{\Character{41}}%
\TextAt{13}{427.923}{496.581}{\Character{59}}%
\TextAt{8}{79.226}{467.58}{extends}%
\TextAt{8}{121.882}{467.58}{uniquely}%
\TextAt{8}{169.356}{467.58}{to}%
\TextAt{8}{183.667}{467.58}{a}%
\TextAt{8}{193.422}{467.58}{homomorphism}%
\TextAt{13}{261.194}{438.578}{\Character{31}}%
\TextAt{21}{268.528}{436.785}{K}%
\TextAt{9}{280.051}{438.578}{\Character{58}}%
\TextAt{13}{286.624}{438.578}{K}%
\TextAt{17}{296.583}{436.785}{0}%
\TextAt{9}{301.315}{438.578}{\Character{40}}%
\Vector{q 0 G 0 g 0.478 w 233.868 -328.065 m 239.712 -328.065 l S Q}%
\TextAt{14}{305.868}{438.578}{o}%
\TextAt{9}{311.712}{438.578}{\Character{41}}%
\NativeArrowAt{319.585}{438.578}{11.955}{19.26}{\longrightarrow}%
\TextAt{18}{342.167}{438.578}{R}%
\TextAt{8}{79.226}{409.577}{on}%
\TextAt{13}{95.485}{409.577}{K}%
\TextAt{17}{105.444}{407.784}{0}%
\TextAt{9}{110.177}{409.577}{\Character{40}}%
\Vector{q 0 G 0 g 0.478 w 42.729 -357.066 m 48.573 -357.066 l S Q}%
\TextAt{14}{114.729}{409.577}{o}%
\TextAt{9}{120.573}{409.577}{\Character{41}}%
\TextAt{8}{125.125}{409.577}{,}%
\TextAt{8}{132.275}{409.577}{and}%
\TextAt{8}{155.037}{409.577}{thereby}%
\TextAt{8}{197.311}{409.577}{on}%
\TextAt{13}{213.57}{409.577}{K}%
\TextAt{17}{224.388}{413.916}{0}%
\TextAt{9}{229.12}{409.577}{\Character{40}}%
\Vector{q 0 G 0 g 0.478 w 161.672 -357.066 m 167.516 -357.066 l S Q}%
\TextAt{14}{233.672}{409.577}{o}%
\TextAt{9}{239.516}{409.577}{\Character{41}}%
\TextAt{8}{244.069}{409.577}{.}%
\TextAt{8}{252.517}{409.577}{It}%
\TextAt{8}{265.19}{409.577}{is}%
\TextAt{8}{276.953}{409.577}{given}%
\TextAt{8}{307.523}{409.577}{by}%
\TextAt{13}{268.465}{380.576}{\Character{31}}%
\TextAt{21}{275.798}{378.783}{K}%
\TextAt{9}{287.322}{380.576}{\Character{61}}%
\TextAt{8}{299.747}{380.576}{deg}%
\TextAt{16}{319.302}{380.576}{\Character{14}}%
\TextAt{8}{327.272}{380.576}{det}%
\TextAt{8}{79.226}{351.575}{and}%
\TextAt{8}{101.989}{351.575}{called}%
\TextAt{8}{135.152}{351.575}{the}%
\TextAt{12}{155.309}{351.575}{Euler–Minkowski}%
\TextAt{12}{261.124}{351.575}{characteristic}%
\TextAt{12}{344.117}{351.575}{over}%
\TextAt{13}{373.323}{351.575}{K}%
\TextAt{8}{384.141}{351.575}{.}%
\TextAt{19}{79.226}{316.596}{Proof.}%
\TextAt{8}{116}{316.596}{Since,}%
\TextAt{8}{151.463}{316.596}{by}%
\TextAt{8}{169.044}{316.596}{(5.6),}%
\TextAt{13}{201.924}{316.596}{K}%
\TextAt{17}{211.883}{314.803}{0}%
\TextAt{9}{216.615}{316.596}{\Character{40}}%
\Vector{q 0 G 0 g 0.478 w 149.168 -450.047 m 155.012 -450.047 l S Q}%
\TextAt{14}{221.168}{316.596}{o}%
\TextAt{9}{227.012}{316.596}{\Character{41}}%
\TextAt{8}{236.784}{316.596}{is}%
\TextAt{8}{249.87}{316.596}{generated}%
\TextAt{8}{304.525}{316.596}{as}%
\TextAt{8}{320.217}{316.596}{an}%
\TextAt{8}{337.799}{316.596}{additive}%
\TextAt{8}{384.001}{316.596}{group}%
\TextAt{8}{418.499}{316.596}{by}%
\TextAt{8}{436.081}{316.596}{the}%
\TextAt{8}{457.559}{316.596}{elements}%
\TextAt{9}{506.738}{316.596}{\Character{91}}%
\TextAt{14}{509.99}{316.596}{a}%
\TextAt{9}{515.973}{316.596}{\Character{93}}%
\NativeMathAt{524.799}{316.596}{11.955}{\in}%
\TextAt{8}{79.226}{299.55}{Pic}%
\TextAt{9}{95.641}{299.55}{\Character{40}}%
\Vector{q 0 G 0 g 0.478 w 28.193 -467.093 m 34.037 -467.093 l S Q}%
\TextAt{14}{100.193}{299.55}{o}%
\TextAt{9}{106.037}{299.55}{\Character{41}}%
\TextAt{8}{110.589}{299.55}{,}%
\TextAt{8}{118.388}{299.55}{the}%
\TextAt{8}{139.064}{299.55}{map}%
\TextAt{8}{165.598}{299.55}{deg}%
\TextAt{21}{183.16}{296.601}{K}%
\TextAt{8}{195.78}{299.55}{on}%
\TextAt{8}{212.558}{299.55}{Pic}%
\TextAt{9}{228.973}{299.55}{\Character{40}}%
\Vector{q 0 G 0 g 0.478 w 161.525 -467.093 m 167.369 -467.093 l S Q}%
\TextAt{14}{233.525}{299.55}{o}%
\TextAt{9}{239.369}{299.55}{\Character{41}}%
\TextAt{8}{248.339}{299.55}{determines}%
\TextAt{8}{308.096}{299.55}{a}%
\TextAt{8}{318.371}{299.55}{unique}%
\TextAt{8}{356.932}{299.55}{homomorphism}%
\TextAt{13}{440.121}{299.55}{K}%
\TextAt{17}{450.08}{297.757}{0}%
\TextAt{9}{454.813}{299.55}{\Character{40}}%
\Vector{q 0 G 0 g 0.478 w 387.365 -467.093 m 393.209 -467.093 l S Q}%
\TextAt{14}{459.365}{299.55}{o}%
\TextAt{9}{465.209}{299.55}{\Character{41}}%
\NativeMathAt{473.968}{299.55}{11.955}{\rightarrow}%
\TextAt{18}{490.129}{299.55}{R}%
\TextAt{8}{503.18}{299.55}{which}%
\TextAt{8}{79.226}{282.504}{extends}%
\TextAt{8}{122.903}{282.504}{deg}%
\TextAt{21}{140.465}{279.555}{K}%
\TextAt{8}{148.668}{282.504}{.}%
\TextAt{8}{160.18}{282.504}{But}%
\TextAt{8}{184.442}{282.504}{such}%
\TextAt{8}{211.86}{282.504}{a}%
\TextAt{8}{222.637}{282.504}{homomorphism}%
\TextAt{8}{306.328}{282.504}{is}%
\TextAt{8}{319.113}{282.504}{given}%
\TextAt{8}{350.703}{282.504}{by}%
\TextAt{8}{367.983}{282.504}{the}%
\TextAt{8}{389.161}{282.504}{composite}%
\TextAt{8}{445.2}{282.504}{of}%
\TextAt{8}{459.551}{282.504}{the}%
\TextAt{8}{480.729}{282.504}{homomor-}%
\TextAt{8}{79.226}{265.458}{phisms}%
\TextAt{13}{244.099}{248.412}{K}%
\TextAt{17}{254.058}{246.619}{0}%
\TextAt{9}{258.79}{248.412}{\Character{40}}%
\Vector{q 0 G 0 g 0.478 w 191.342 -518.231 m 197.187 -518.231 l S Q}%
\TextAt{14}{263.342}{248.412}{o}%
\TextAt{9}{269.187}{248.412}{\Character{41}}%
\TextAt{7}{279.647}{255.19}{det}%
\NativeArrowAt{277.06}{248.412}{11.955}{18.339}{\longrightarrow}%
\TextAt{8}{298.723}{248.412}{Pic}%
\TextAt{9}{315.138}{248.412}{\Character{40}}%
\Vector{q 0 G 0 g 0.478 w 247.69 -518.231 m 253.534 -518.231 l S Q}%
\TextAt{14}{319.69}{248.412}{o}%
\TextAt{9}{325.534}{248.412}{\Character{41}}%
\TextAt{7}{335.995}{255.761}{deg}%
\NativeArrowAt{333.407}{248.412}{11.955}{19.272}{\longrightarrow}%
\TextAt{18}{356.011}{248.412}{R}%
\TextAt{13}{364.645}{248.412}{\Character{59}}%
\TextAt{8}{79.226}{224.891}{as}%
\TextAt{8}{93.596}{224.891}{the}%
\TextAt{8}{113.753}{224.891}{composite}%
\TextAt{8}{168.77}{224.891}{Pic}%
\TextAt{9}{185.185}{224.891}{\Character{40}}%
\Vector{q 0 G 0 g 0.478 w 117.737 -541.753 m 123.581 -541.753 l S Q}%
\TextAt{14}{189.737}{224.891}{o}%
\TextAt{9}{195.581}{224.891}{\Character{41}}%
\NativeArrowAt{203.455}{224.891}{11.955}{13.214}{\hookrightarrow}%
\TextAt{13}{219.99}{224.891}{K}%
\TextAt{17}{229.949}{223.098}{0}%
\TextAt{9}{234.681}{224.891}{\Character{40}}%
\Vector{q 0 G 0 g 0.478 w 167.233 -541.753 m 173.078 -541.753 l S Q}%
\TextAt{14}{239.233}{224.891}{o}%
\TextAt{9}{245.078}{224.891}{\Character{41}}%
\TextAt{7}{255.538}{231.668}{det}%
\NativeArrowAt{252.951}{224.891}{11.955}{18.339}{\longrightarrow}%
\TextAt{8}{274.614}{224.891}{Pic}%
\TextAt{9}{291.029}{224.891}{\Character{40}}%
\Vector{q 0 G 0 g 0.478 w 223.581 -541.753 m 229.425 -541.753 l S Q}%
\TextAt{14}{295.581}{224.891}{o}%
\TextAt{9}{301.425}{224.891}{\Character{41}}%
\TextAt{8}{309.875}{224.891}{is}%
\TextAt{8}{321.639}{224.891}{the}%
\TextAt{8}{341.795}{224.891}{identity.}%
\Vector{q 0 G 0 g 0.399 w 452.334 -549.109 m 452.334 -541.039 l S Q}%
\Vector{q 0 G 0 g 0.399 w 452.534 -541.239 m 459.707 -541.239 l S Q}%
\Vector{q 0 G 0 g 0.399 w 452.534 -548.91 m 459.707 -548.91 l S Q}%
\Vector{q 0 G 0 g 0.399 w 459.906 -549.109 m 459.906 -541.039 l S Q}%
\TextAt{8}{97.159}{195.89}{Via}%
\TextAt{8}{119.33}{195.89}{the}%
\TextAt{8}{139.876}{195.89}{Poincaré}%
\TextAt{8}{188.229}{195.89}{isomorphism}%
\TextAt{13}{257.038}{195.89}{K}%
\TextAt{17}{266.997}{194.096}{0}%
\TextAt{9}{271.729}{195.89}{\Character{40}}%
\Vector{q 0 G 0 g 0.478 w 204.281 -570.754 m 210.126 -570.754 l S Q}%
\TextAt{14}{276.281}{195.89}{o}%
\TextAt{9}{282.126}{195.89}{\Character{41}}%
\NativeMathAt{293.25}{202.667}{7.97}{\sim}%
\NativeArrowAt{290.662}{195.89}{11.955}{13.175}{\longrightarrow}%
\TextAt{13}{307.819}{195.89}{K}%
\TextAt{17}{318.637}{200.228}{0}%
\TextAt{9}{323.369}{195.89}{\Character{40}}%
\Vector{q 0 G 0 g 0.478 w 255.921 -570.754 m 261.765 -570.754 l S Q}%
\TextAt{14}{327.921}{195.89}{o}%
\TextAt{9}{333.765}{195.89}{\Character{41}}%
\TextAt{8}{338.318}{195.89}{,}%
\TextAt{8}{345.953}{195.89}{we}%
\TextAt{8}{363.57}{195.89}{transfer}%
\TextAt{8}{407.273}{195.89}{the}%
\TextAt{8}{427.818}{195.89}{maps}%
\TextAt{8}{458.836}{195.89}{det}%
\TextAt{8}{479.382}{195.89}{and}%
\TextAt{13}{502.534}{195.89}{\Character{31}}%
\TextAt{21}{509.867}{194.096}{K}%
\TextAt{8}{522.357}{195.89}{to}%
\TextAt{8}{79.226}{178.844}{the}%
\TextAt{8}{99.074}{178.844}{Grothendieck}%
\TextAt{8}{171.035}{178.844}{group}%
\TextAt{13}{203.902}{178.844}{K}%
\TextAt{17}{214.719}{183.182}{0}%
\TextAt{9}{219.452}{178.844}{\Character{40}}%
\Vector{q 0 G 0 g 0.478 w 152.004 -587.8 m 157.848 -587.8 l S Q}%
\TextAt{14}{224.004}{178.844}{o}%
\TextAt{9}{229.848}{178.844}{\Character{41}}%
\TextAt{8}{237.99}{178.844}{of}%
\TextAt{8}{251.011}{178.844}{coherent}%
\TextAt{8}{297.854}{178.844}{metrized}%
\TextAt{14}{345.665}{178.844}{o}%
\TextAt{8}{351.509}{178.844}{-modules.}%
\TextAt{8}{405.763}{178.844}{Then}%
\TextAt{8}{436.012}{178.844}{proposition}%
\TextAt{8}{497.236}{178.844}{(8.1)}%
\TextAt{8}{524.903}{178.844}{is}%
\TextAt{8}{79.226}{161.798}{equally}%
\TextAt{8}{119.688}{161.798}{valid}%
\TextAt{8}{148.124}{161.798}{for}%
\TextAt{13}{166.146}{161.798}{K}%
\TextAt{17}{176.964}{166.136}{0}%
\TextAt{9}{181.696}{161.798}{\Character{40}}%
\Vector{q 0 G 0 g 0.478 w 114.249 -604.846 m 120.093 -604.846 l S Q}%
\TextAt{14}{186.249}{161.798}{o}%
\TextAt{9}{192.093}{161.798}{\Character{41}}%
\TextAt{8}{200.68}{161.798}{as}%
\TextAt{8}{215.188}{161.798}{for}%
\TextAt{13}{233.21}{161.798}{K}%
\TextAt{17}{243.169}{160.004}{0}%
\TextAt{9}{247.901}{161.798}{\Character{40}}%
\Vector{q 0 G 0 g 0.478 w 180.454 -604.846 m 186.298 -604.846 l S Q}%
\TextAt{14}{252.454}{161.798}{o}%
\TextAt{9}{258.298}{161.798}{\Character{41}}%
\TextAt{8}{262.85}{161.798}{.}%
\TextAt{8}{271.711}{161.798}{We}%
\TextAt{8}{291.993}{161.798}{define}%
\TextAt{8}{325.94}{161.798}{in}%
\TextAt{8}{339.73}{161.798}{what}%
\TextAt{8}{369.134}{161.798}{follows}%
\TextAt{13}{407.707}{161.798}{\Character{31}}%
\TextAt{21}{415.041}{160.004}{K}%
\TextAt{9}{423.243}{161.798}{\Character{40}}%
\TextAt{13}{427.796}{161.798}{M}%
\TextAt{9}{440.369}{161.798}{\Character{41}}%
\TextAt{9}{448.463}{161.798}{\Character{61}}%
\TextAt{13}{461.136}{161.798}{\Character{31}}%
\TextAt{21}{468.47}{160.004}{K}%
\TextAt{9}{476.673}{161.798}{\Character{40}\Character{91}}%
\TextAt{13}{484.477}{161.798}{M}%
\TextAt{9}{497.05}{161.798}{\Character{93}\Character{41}}%
\TextAt{8}{508.889}{161.798}{for}%
\TextAt{8}{526.911}{161.798}{a}%
\TextAt{8}{79.226}{144.752}{metrized}%
\TextAt{14}{127.463}{144.752}{o}%
\TextAt{8}{133.308}{144.752}{-module}%
\TextAt{13}{178.616}{144.752}{M}%
\TextAt{8}{191.19}{144.752}{.}%
\TextAt{8}{199.991}{144.752}{If}%
\TextAt{13}{211.801}{144.752}{L}%
\TextAt{16}{223.287}{144.752}{\Character{106}}%
\TextAt{13}{230.13}{144.752}{K}%
\TextAt{8}{244.963}{144.752}{is}%
\TextAt{8}{256.844}{144.752}{an}%
\TextAt{8}{273.221}{144.752}{extension}%
\TextAt{8}{325.105}{144.752}{of}%
\TextAt{8}{338.553}{144.752}{algebraic}%
\TextAt{8}{388.105}{144.752}{number}%
\TextAt{8}{430.819}{144.752}{fields}%
\TextAt{8}{460.909}{144.752}{and}%
\TextAt{13}{483.789}{144.752}{i}%
\TextAt{9}{491.304}{144.752}{\Character{58}}%
\TextAt{14}{498.077}{144.752}{o}%
\NativeMathAt{507.443}{144.752}{11.955}{\rightarrow}%
\TextAt{16}{522.913}{144.752}{\Character{79}}%
\TextAt{8}{79.226}{127.706}{the}%
\TextAt{8}{100.223}{127.706}{inclusion}%
\TextAt{8}{149.9}{127.706}{of}%
\TextAt{8}{164.071}{127.706}{the}%
\TextAt{8}{185.067}{127.706}{maximal}%
\TextAt{8}{233.717}{127.706}{orders}%
\TextAt{8}{269.741}{127.706}{of}%
\TextAt{13}{283.911}{127.706}{K}%
\TextAt{8}{294.729}{127.706}{,}%
\TextAt{8}{302.929}{127.706}{resp.}%
\TextAt{13}{334.772}{127.706}{L}%
\TextAt{8}{342.736}{127.706}{,}%
\TextAt{8}{350.936}{127.706}{then}%
\TextAt{8}{378.437}{127.706}{applying}%
\TextAt{8}{427.086}{127.706}{deg}%
\TextAt{21}{444.648}{124.757}{K}%
\TextAt{8}{457.588}{127.706}{to}%
\TextAt{8}{472.739}{127.706}{the}%
\TextAt{8}{493.736}{127.706}{formula}%
\TextAt{8}{79.226}{110.66}{(7.2)}%
\TextAt{8}{107.201}{110.66}{and}%
\TextAt{8}{129.964}{110.66}{using}%
\TextAt{8}{149.227}{81.659}{deg}%
\TextAt{21}{166.789}{78.71}{L}%
\TextAt{9}{173.046}{81.659}{\Character{40}}%
\TextAt{14}{177.599}{81.659}{A}%
\TextAt{9}{186.178}{81.659}{\Character{41}}%
\TextAt{9}{194.033}{81.659}{\Character{61}}%
\TextAt{16}{206.476}{81.659}{\Character{0}}%
\TextAt{8}{217.767}{81.659}{log}%
\TextAt{14}{234.728}{81.659}{N}%
\TextAt{9}{244.676}{81.659}{\Character{40}}%
\TextAt{14}{249.228}{81.659}{A}%
\TextAt{9}{257.807}{81.659}{\Character{41}}%
\TextAt{9}{265.662}{81.659}{\Character{61}}%
\TextAt{16}{278.106}{81.659}{\Character{0}}%
\TextAt{8}{289.397}{81.659}{log}%
\TextAt{14}{306.357}{81.659}{N}%
\TextAt{9}{316.305}{81.659}{\Character{40}}%
\TextAt{13}{320.858}{81.659}{N}%
\TextAt{21}{330.231}{79.799}{L}%
\TextAt{3}{335.99}{79.799}{\Character{106}}%
\TextAt{21}{338.343}{79.799}{K}%
\TextAt{9}{346.545}{81.659}{\Character{40}}%
\TextAt{14}{351.098}{81.659}{A}%
\TextAt{9}{359.677}{81.659}{\Character{41}\Character{41}}%
\TextAt{9}{372.075}{81.659}{\Character{61}}%
\TextAt{8}{384.528}{81.659}{deg}%
\TextAt{21}{402.09}{78.71}{K}%
\TextAt{9}{410.293}{81.659}{\Character{40}}%
\TextAt{13}{414.845}{81.659}{N}%
\TextAt{21}{424.219}{79.799}{L}%
\TextAt{3}{429.978}{79.799}{\Character{106}}%
\TextAt{21}{432.33}{79.799}{K}%
\TextAt{9}{440.533}{81.659}{\Character{40}}%
\TextAt{14}{445.085}{81.659}{A}%
\TextAt{9}{453.664}{81.659}{\Character{41}\Character{41}}%
\TextAt{8}{297.211}{43.033}{235}%
\Destination{78.23}{809.815}{nav.page.235}%
\Destination{79.23}{773.075}{nav.section.3.8}%
\Destination{79.36}{709.855}{nav.print.256}%
\Destination{79.36}{539.555}{nav.number.3.8.1}%
\Destination{79.23}{773.075}{outline.39}%
\Hotlink{410.104}{709.655}{6.858}{11.047}{nav.section.3.3}%
\Hotlink{168.544}{314.355}{25.078}{11.047}{nav.number.3.5.6}%
\Hotlink{496.736}{176.655}{25.078}{11.047}{nav.number.3.8.1}%
\Hotlink{78.726}{108.455}{25.078}{11.047}{nav.number.3.7.2}%
}

% Page 248 (body page 236)
\Page{%
\TextAt{2}{79.226}{796.943}{III.}%
\TextAt{2}{98.317}{796.943}{Riemann–Roch}%
\TextAt{2}{174.375}{796.943}{Theory}%
\Vector{q 0 G 0 g 0.299 w 7.226 18.729 m 460.769 18.729 l S Q}%
\TextAt{8}{79.226}{761.124}{(see}%
\TextAt{8}{102.694}{761.124}{(1.6),}%
\TextAt{8}{133.921}{761.124}{(iii))}%
\TextAt{8}{161.239}{761.124}{gives}%
\TextAt{8}{189.919}{761.124}{the}%
\TextAt{12}{79.226}{735.973}{Theorem}%
\TextAt{12}{136.24}{735.973}{(8.2).}%
\TextAt{8}{171.556}{735.973}{For}%
\TextAt{8}{192.525}{735.973}{every}%
\TextAt{8}{223.417}{735.973}{coherent}%
\TextAt{16}{270.569}{735.973}{\Character{79}}%
\TextAt{8}{280.418}{735.973}{-module}%
\TextAt{13}{325.609}{735.973}{M}%
\TextAt{8}{338.183}{735.973}{,}%
\TextAt{8}{345.332}{735.973}{the}%
\TextAt{8}{365.488}{735.973}{Riemann–Roch}%
\TextAt{8}{447.118}{735.973}{formula}%
\TextAt{13}{198.256}{709.99}{\Character{31}}%
\TextAt{21}{205.589}{708.196}{K}%
\TextAt{9}{213.792}{709.99}{\Character{40}}%
\TextAt{13}{218.344}{709.99}{i}%
\TextAt{3}{222.338}{708.196}{\Character{3}}%
\TextAt{13}{227.07}{709.99}{M}%
\TextAt{9}{239.644}{709.99}{\Character{41}}%
\TextAt{9}{247.499}{709.99}{\Character{61}}%
\TextAt{8}{259.942}{709.99}{deg}%
\TextAt{21}{277.504}{707.041}{L}%
\TextAt{9}{283.762}{709.99}{\Character{40}}%
\TextAt{8}{288.314}{709.99}{det}%
\TextAt{13}{306.566}{709.99}{M}%
\TextAt{9}{319.139}{709.99}{\Character{41}}%
\TextAt{9}{326.336}{709.99}{\Character{43}}%
\TextAt{8}{338.109}{709.99}{rk}%
\TextAt{9}{348.845}{709.99}{\Character{40}}%
\TextAt{13}{353.397}{709.99}{M}%
\TextAt{9}{365.971}{709.99}{\Character{41}}%
\TextAt{13}{370.523}{709.99}{\Character{31}}%
\TextAt{21}{377.857}{708.196}{K}%
\TextAt{9}{386.06}{709.99}{\Character{40}}%
\TextAt{13}{390.612}{709.99}{i}%
\TextAt{3}{394.605}{708.196}{\Character{3}}%
\TextAt{16}{399.338}{709.99}{\Character{79}}%
\TextAt{9}{409.187}{709.99}{\Character{41}}%
\TextAt{8}{79.226}{684.006}{is}%
\TextAt{8}{90.99}{684.006}{valid,}%
\TextAt{8}{122.54}{684.006}{and}%
\TextAt{8}{145.302}{684.006}{in}%
\TextAt{8}{158.955}{684.006}{particular,}%
\TextAt{8}{216.197}{684.006}{for}%
\TextAt{8}{234.081}{684.006}{an}%
\TextAt{8}{250.34}{684.006}{invertible}%
\TextAt{8}{302.047}{684.006}{metrized}%
\TextAt{16}{350.166}{684.006}{\Character{79}}%
\TextAt{8}{360.016}{684.006}{-module}%
\TextAt{13}{405.206}{684.006}{M}%
\TextAt{8}{417.78}{684.006}{,}%
\TextAt{8}{424.929}{684.006}{we}%
\TextAt{8}{442.157}{684.006}{have}%
\TextAt{13}{221.963}{658.022}{\Character{31}}%
\TextAt{21}{229.296}{656.229}{K}%
\TextAt{9}{237.499}{658.022}{\Character{40}}%
\TextAt{13}{242.051}{658.022}{i}%
\TextAt{3}{246.045}{656.229}{\Character{3}}%
\TextAt{13}{250.777}{658.022}{M}%
\TextAt{9}{263.351}{658.022}{\Character{41}}%
\TextAt{9}{271.206}{658.022}{\Character{61}}%
\TextAt{8}{283.649}{658.022}{deg}%
\TextAt{21}{301.211}{655.073}{L}%
\TextAt{9}{307.469}{658.022}{\Character{40}}%
\TextAt{13}{312.021}{658.022}{M}%
\TextAt{9}{324.595}{658.022}{\Character{41}}%
\TextAt{9}{331.792}{658.022}{\Character{43}}%
\TextAt{13}{343.565}{658.022}{\Character{31}}%
\TextAt{21}{350.898}{656.229}{K}%
\TextAt{9}{359.101}{658.022}{\Character{40}}%
\TextAt{13}{363.653}{658.022}{i}%
\TextAt{3}{367.647}{656.229}{\Character{3}}%
\TextAt{16}{372.379}{658.022}{\Character{79}}%
\TextAt{9}{382.229}{658.022}{\Character{41}}%
\TextAt{13}{386.781}{658.022}{\Character{58}}%
\TextAt{8}{97.159}{623.934}{We}%
\TextAt{8}{117.554}{623.934}{now}%
\TextAt{8}{142.193}{623.934}{specialize}%
\TextAt{8}{194.197}{623.934}{to}%
\TextAt{8}{208.758}{623.934}{the}%
\TextAt{8}{229.165}{623.934}{case}%
\TextAt{8}{254.187}{623.934}{of}%
\TextAt{8}{267.767}{623.934}{the}%
\TextAt{8}{288.174}{623.934}{base}%
\TextAt{8}{314.499}{623.934}{field}%
\TextAt{13}{340.107}{623.934}{K}%
\TextAt{9}{354.672}{623.934}{\Character{61}}%
\TextAt{18}{367.525}{623.934}{Q}%
\TextAt{8}{376.823}{623.934}{,}%
\TextAt{8}{384.285}{623.934}{that}%
\TextAt{8}{409.905}{623.934}{is,}%
\TextAt{8}{425.234}{623.934}{we}%
\TextAt{8}{442.712}{623.934}{consider}%
\TextAt{8}{488.547}{623.934}{metrized}%
\TextAt{18}{79.226}{606.888}{Z}%
\TextAt{8}{87.196}{606.888}{-modules.}%
\TextAt{8}{141.751}{606.888}{Such}%
\TextAt{8}{170.103}{606.888}{a}%
\TextAt{8}{179.924}{606.888}{module}%
\TextAt{8}{221.284}{606.888}{is}%
\TextAt{8}{233.114}{606.888}{simply}%
\TextAt{8}{270.635}{606.888}{a}%
\TextAt{8}{280.457}{606.888}{finitely}%
\TextAt{8}{319.867}{606.888}{generated}%
\TextAt{8}{373.265}{606.888}{abelian}%
\TextAt{8}{413.979}{606.888}{group}%
\TextAt{13}{447.221}{606.888}{M}%
\TextAt{8}{463.758}{606.888}{together}%
\TextAt{8}{510.007}{606.888}{with}%
\TextAt{8}{79.226}{589.842}{a}%
\TextAt{8}{88.982}{589.842}{euclidean}%
\TextAt{8}{140.353}{589.842}{metric}%
\TextAt{8}{176.768}{589.842}{on}%
\TextAt{8}{193.027}{589.842}{the}%
\TextAt{8}{213.184}{589.842}{real}%
\TextAt{8}{235.946}{589.842}{vector}%
\TextAt{8}{271.071}{589.842}{space}%
\TextAt{13}{266.894}{563.858}{M}%
\TextAt{25}{278.218}{562.065}{R}%
\TextAt{9}{287.794}{563.858}{\Character{61}}%
\TextAt{13}{300.219}{563.858}{M}%
\TextAt{16}{315.449}{563.858}{\Character{10}}%
\TextAt{25}{324.748}{562.065}{Z}%
\TextAt{18}{333.216}{563.858}{R}%
\TextAt{13}{341.85}{563.858}{\Character{58}}%
\TextAt{8}{79.226}{537.874}{Indeed,}%
\TextAt{8}{119.983}{537.874}{since}%
\TextAt{18}{147.996}{537.874}{Q}%
\TextAt{8}{160.537}{537.874}{has}%
\TextAt{8}{180.755}{537.874}{only}%
\TextAt{8}{205.792}{537.874}{a}%
\TextAt{8}{214.892}{537.874}{single}%
\TextAt{8}{246.814}{537.874}{embedding}%
\TextAt{8}{305.337}{537.874}{into}%
\TextAt{18}{328.425}{537.874}{C}%
\TextAt{8}{337.059}{537.874}{,}%
\TextAt{8}{343.684}{537.874}{i.e.,}%
\TextAt{18}{365.265}{537.874}{Q}%
\TextAt{25}{374.563}{536.081}{C}%
\TextAt{9}{384.138}{537.874}{\Character{61}}%
\TextAt{18}{396.564}{537.874}{C}%
\TextAt{8}{405.198}{537.874}{,}%
\TextAt{8}{411.823}{537.874}{a}%
\TextAt{8}{420.923}{537.874}{metric}%
\TextAt{8}{456.683}{537.874}{on}%
\TextAt{13}{472.287}{537.874}{M}%
\TextAt{8}{488.103}{537.874}{is}%
\TextAt{8}{499.211}{537.874}{simply}%
\TextAt{8}{79.226}{520.828}{given}%
\TextAt{8}{110.492}{520.828}{by}%
\TextAt{8}{127.448}{520.828}{a}%
\TextAt{8}{137.9}{520.828}{hermitian}%
\TextAt{8}{191.929}{520.828}{scalar}%
\TextAt{8}{225.861}{520.828}{product}%
\TextAt{8}{270.457}{520.828}{on}%
\TextAt{8}{287.413}{520.828}{the}%
\TextAt{18}{308.266}{520.828}{C}%
\TextAt{8}{316.9}{520.828}{-vector}%
\TextAt{8}{356.619}{520.828}{space}%
\TextAt{13}{388.59}{520.828}{M}%
\TextAt{25}{399.915}{519.035}{C}%
\TextAt{9}{410.678}{520.828}{\Character{61}}%
\TextAt{13}{424.29}{520.828}{M}%
\TextAt{25}{435.615}{519.035}{R}%
\TextAt{16}{445.001}{520.828}{\Character{10}}%
\TextAt{18}{457.431}{520.828}{C}%
\TextAt{8}{466.066}{520.828}{.}%
\TextAt{8}{476.604}{520.828}{Restricting}%
\TextAt{8}{79.226}{503.782}{this}%
\TextAt{8}{102.965}{503.782}{to}%
\TextAt{13}{118.192}{503.782}{M}%
\TextAt{25}{129.517}{501.989}{R}%
\TextAt{8}{140.585}{503.782}{gives}%
\TextAt{8}{170.182}{503.782}{a}%
\TextAt{8}{180.854}{503.782}{euclidean}%
\TextAt{8}{233.142}{503.782}{metric}%
\TextAt{8}{270.474}{503.782}{the}%
\TextAt{8}{291.546}{503.782}{sesquilinear}%
\TextAt{8}{355.347}{503.782}{extension}%
\TextAt{8}{408.029}{503.782}{of}%
\TextAt{8}{422.276}{503.782}{which}%
\TextAt{8}{456.679}{503.782}{reproduces}%
\TextAt{8}{516.51}{503.782}{the}%
\TextAt{8}{79.226}{486.736}{original}%
\TextAt{8}{121.512}{486.736}{metric.}%
\TextAt{8}{97.159}{466.701}{If}%
\TextAt{13}{109.907}{466.701}{M}%
\TextAt{8}{127.433}{466.701}{is}%
\TextAt{8}{140.253}{466.701}{a}%
\TextAt{19}{151.064}{466.701}{projective}%
\TextAt{8}{204.375}{466.701}{metrized}%
\TextAt{18}{253.551}{466.701}{Z}%
\TextAt{8}{261.521}{466.701}{-module,}%
\TextAt{8}{311.283}{466.701}{then}%
\TextAt{8}{338.998}{466.701}{the}%
\TextAt{8}{360.21}{466.701}{underlying}%
\TextAt{18}{419.476}{466.701}{Z}%
\TextAt{8}{427.446}{466.701}{-module}%
\TextAt{8}{473.692}{466.701}{is}%
\TextAt{8}{486.511}{466.701}{a}%
\TextAt{8}{497.322}{466.701}{finitely}%
\TextAt{8}{79.226}{449.656}{generated}%
\TextAt{8}{132.721}{449.656}{free}%
\TextAt{8}{155.311}{449.656}{abelian}%
\TextAt{8}{196.121}{449.656}{group.}%
\TextAt{8}{234.336}{449.656}{The}%
\TextAt{8}{258.552}{449.656}{canonical}%
\TextAt{8}{310.098}{449.656}{map}%
\TextAt{13}{336.275}{449.656}{M}%
\NativeMathAt{352.447}{449.656}{11.955}{\rightarrow}%
\TextAt{13}{368}{449.656}{M}%
\TextAt{16}{383.341}{449.656}{\Character{10}}%
\TextAt{18}{395.407}{449.656}{R}%
\TextAt{8}{404.041}{449.656}{,}%
\TextAt{13}{411.394}{449.656}{a}%
\NativeArrowAt{421.137}{449.656}{11.955}{11.955}{\mapsto}%
\TextAt{13}{436.69}{449.656}{a}%
\TextAt{16}{445.602}{449.656}{\Character{10}}%
\TextAt{9}{457.668}{449.656}{1}%
\TextAt{8}{463.521}{449.656}{,}%
\TextAt{8}{470.874}{449.656}{identifies}%
\TextAt{13}{520.196}{449.656}{M}%
\TextAt{8}{79.226}{432.61}{with}%
\TextAt{8}{105.886}{432.61}{a}%
\TextAt{8}{115.642}{432.61}{complete}%
\TextAt{8}{165.064}{432.61}{lattice}%
\TextAt{8}{200.834}{432.61}{in}%
\TextAt{13}{214.487}{432.61}{M}%
\TextAt{25}{225.812}{430.816}{R}%
\TextAt{8}{232.066}{432.61}{.}%
\TextAt{8}{240.514}{432.61}{If}%
\TextAt{13}{252.207}{432.61}{\Character{11}}%
\TextAt{17}{259.652}{430.816}{1}%
\TextAt{13}{264.385}{432.61}{\Character{59}}%
\TextAt{13}{269.609}{432.61}{\Character{58}}%
\TextAt{13}{274.846}{432.61}{\Character{58}}%
\TextAt{13}{280.07}{432.61}{\Character{58}}%
\TextAt{13}{285.307}{432.61}{\Character{59}}%
\TextAt{13}{290.543}{432.61}{\Character{11}}%
\TextAt{21}{298.051}{430.816}{n}%
\TextAt{8}{307.585}{432.61}{is}%
\TextAt{8}{319.349}{432.61}{a}%
\TextAt{18}{329.104}{432.61}{Z}%
\TextAt{8}{337.074}{432.61}{-basis}%
\TextAt{8}{369.712}{432.61}{of}%
\TextAt{13}{383.042}{432.61}{M}%
\TextAt{8}{395.615}{432.61}{,}%
\TextAt{8}{402.765}{432.61}{then}%
\TextAt{8}{429.425}{432.61}{the}%
\TextAt{8}{449.581}{432.61}{set}%
\TextAt{9}{189.157}{406.626}{\Character{8}}%
\TextAt{9}{200.921}{406.626}{\Character{61}}%
\TextAt{16}{213.358}{406.626}{\Character{102}}%
\TextAt{13}{219.335}{406.626}{x}%
\TextAt{17}{225.987}{404.832}{1}%
\TextAt{13}{230.72}{406.626}{\Character{11}}%
\TextAt{17}{238.165}{404.832}{1}%
\TextAt{9}{245.554}{406.626}{\Character{43}}%
\TextAt{16}{257.316}{406.626}{\Character{1}}%
\TextAt{16}{262.612}{406.626}{\Character{1}}%
\TextAt{16}{267.92}{406.626}{\Character{1}}%
\TextAt{9}{273.92}{406.626}{\Character{43}}%
\TextAt{13}{285.681}{406.626}{x}%
\TextAt{21}{292.333}{404.832}{n}%
\TextAt{13}{297.97}{406.626}{\Character{11}}%
\TextAt{21}{305.415}{404.832}{n}%
\TextAt{16}{314.372}{406.626}{\Character{106}}%
\TextAt{13}{321.014}{406.626}{x}%
\TextAt{21}{327.666}{404.832}{i}%
\NativeMathAt{334.368}{406.626}{11.955}{\in}%
\TextAt{18}{345.659}{406.626}{R}%
\TextAt{13}{354.294}{406.626}{\Character{59}}%
\TextAt{9}{363.435}{406.626}{0}%
\NativeMathAt{372.609}{406.626}{11.955}{\leq}%
\TextAt{13}{385.228}{406.626}{x}%
\TextAt{21}{391.88}{404.832}{i}%
\TextAt{13}{398.583}{406.626}{\Character{60}}%
\TextAt{9}{411.008}{406.626}{1}%
\TextAt{16}{416.861}{406.626}{\Character{103}}%
\TextAt{8}{79.226}{380.642}{is}%
\TextAt{8}{92.326}{380.642}{a}%
\TextAt{19}{103.417}{380.642}{fundamental}%
\TextAt{19}{171.8}{380.642}{mesh}%
\TextAt{8}{203.011}{380.642}{of}%
\TextAt{8}{217.677}{380.642}{the}%
\TextAt{8}{239.169}{380.642}{lattice}%
\TextAt{13}{276.275}{380.642}{M}%
\TextAt{8}{288.848}{380.642}{.}%
\TextAt{8}{301.304}{380.642}{The}%
\TextAt{8}{326.693}{380.642}{euclidean}%
\TextAt{8}{379.4}{380.642}{metric}%
\TextAt{16}{417.151}{380.642}{\Character{104}}%
\TextAt{13}{427.033}{380.642}{\Character{59}}%
\TextAt{16}{437.511}{380.642}{\Character{105}}%
\TextAt{21}{442.16}{378.849}{M}%
\TextAt{8}{456.873}{380.642}{defines}%
\TextAt{8}{496.632}{380.642}{a}%
\TextAt{8}{507.723}{380.642}{Haar}%
\TextAt{8}{79.226}{363.596}{measure}%
\TextAt{8}{123.787}{363.596}{on}%
\TextAt{13}{139.021}{363.596}{M}%
\TextAt{25}{150.346}{361.803}{R}%
\TextAt{8}{156.6}{363.596}{.}%
\TextAt{8}{164.707}{363.596}{Once}%
\TextAt{8}{193.594}{363.596}{we}%
\TextAt{8}{209.797}{363.596}{choose}%
\TextAt{8}{245.905}{363.596}{an}%
\TextAt{8}{261.139}{363.596}{orthonormal}%
\TextAt{8}{327.123}{363.596}{basis}%
\TextAt{13}{354.838}{363.596}{e}%
\TextAt{17}{360.264}{361.803}{1}%
\TextAt{13}{364.996}{363.596}{\Character{59}}%
\TextAt{13}{370.22}{363.596}{\Character{58}}%
\TextAt{13}{375.457}{363.596}{\Character{58}}%
\TextAt{13}{380.693}{363.596}{\Character{58}}%
\TextAt{13}{385.918}{363.596}{\Character{59}}%
\TextAt{13}{391.154}{363.596}{e}%
\TextAt{21}{396.642}{361.803}{n}%
\TextAt{8}{405.151}{363.596}{of}%
\TextAt{13}{417.457}{363.596}{M}%
\TextAt{25}{428.782}{361.803}{R}%
\TextAt{8}{435.036}{363.596}{,}%
\TextAt{8}{441.365}{363.596}{this}%
\TextAt{8}{463.163}{363.596}{Haar}%
\TextAt{8}{491.082}{363.596}{measure}%
\TextAt{8}{79.226}{346.55}{can}%
\TextAt{8}{100.017}{346.55}{be}%
\TextAt{8}{115.273}{346.55}{expressed,}%
\TextAt{8}{170.462}{346.55}{via}%
\TextAt{8}{188.982}{346.55}{the}%
\TextAt{8}{208.47}{346.55}{isomorphism}%
\TextAt{13}{276.221}{346.55}{M}%
\TextAt{25}{287.546}{344.757}{R}%
\NativeMathAt{299.709}{353.327}{7.97}{\sim}%
\NativeArrowAt{297.121}{346.55}{11.955}{13.175}{\longrightarrow}%
\TextAt{18}{313.615}{346.55}{R}%
\TextAt{21}{322.249}{350.888}{n}%
\TextAt{8}{327.886}{346.55}{,}%
\TextAt{13}{334.5}{346.55}{x}%
\TextAt{17}{341.152}{344.757}{1}%
\TextAt{13}{345.885}{346.55}{e}%
\TextAt{17}{351.31}{344.757}{1}%
\TextAt{9}{357.332}{346.55}{\Character{43}}%
\TextAt{16}{367.727}{346.55}{\Character{1}}%
\TextAt{16}{373.023}{346.55}{\Character{1}}%
\TextAt{16}{378.331}{346.55}{\Character{1}}%
\TextAt{9}{382.965}{346.55}{\Character{43}}%
\TextAt{13}{393.36}{346.55}{x}%
\TextAt{21}{400.012}{344.757}{n}%
\TextAt{13}{405.648}{346.55}{e}%
\TextAt{21}{411.074}{344.757}{n}%
\NativeArrowAt{420.031}{346.55}{11.955}{11.955}{\mapsto}%
\TextAt{9}{435.307}{346.55}{\Character{40}}%
\TextAt{13}{439.859}{346.55}{x}%
\TextAt{17}{446.511}{344.757}{1}%
\TextAt{13}{451.244}{346.55}{\Character{59}}%
\TextAt{13}{456.468}{346.55}{\Character{58}}%
\TextAt{13}{461.705}{346.55}{\Character{58}}%
\TextAt{13}{466.929}{346.55}{\Character{58}}%
\TextAt{13}{472.166}{346.55}{\Character{59}}%
\TextAt{13}{477.402}{346.55}{x}%
\TextAt{21}{484.116}{344.757}{n}%
\TextAt{9}{489.753}{346.55}{\Character{41}}%
\TextAt{8}{494.305}{346.55}{,}%
\TextAt{8}{500.92}{346.55}{by}%
\TextAt{8}{516.51}{346.55}{the}%
\TextAt{8}{79.226}{329.504}{Lebesgue}%
\TextAt{8}{130.133}{329.504}{measure}%
\TextAt{8}{176.006}{329.504}{on}%
\TextAt{18}{192.553}{329.504}{R}%
\TextAt{21}{201.187}{333.842}{n}%
\TextAt{8}{206.824}{329.504}{.}%
\TextAt{8}{216.137}{329.504}{With}%
\TextAt{8}{246.66}{329.504}{respect}%
\TextAt{8}{286.998}{329.504}{to}%
\TextAt{8}{301.596}{329.504}{this}%
\TextAt{8}{324.707}{329.504}{measure,}%
\TextAt{8}{373.904}{329.504}{the}%
\TextAt{8}{394.349}{329.504}{volume}%
\TextAt{8}{434.962}{329.504}{of}%
\TextAt{8}{448.58}{329.504}{the}%
\TextAt{8}{469.025}{329.504}{fundamental}%
\TextAt{8}{79.226}{312.458}{mesh}%
\TextAt{9}{109.198}{312.458}{\Character{8}}%
\TextAt{8}{121.549}{312.458}{is}%
\TextAt{8}{133.313}{312.458}{given}%
\TextAt{8}{163.883}{312.458}{by}%
\TextAt{8}{237.876}{295.412}{vol}%
\TextAt{9}{253.167}{295.412}{\Character{40}\Character{8}\Character{41}}%
\TextAt{9}{274.017}{295.412}{\Character{61}}%
\TextAt{16}{286.472}{295.412}{\Character{106}}%
\TextAt{8}{291.785}{295.412}{det}%
\TextAt{9}{308.044}{295.412}{\Character{40}}%
\TextAt{16}{312.597}{295.412}{\Character{104}}%
\TextAt{13}{317.246}{295.412}{\Character{11}}%
\TextAt{21}{324.692}{293.619}{i}%
\TextAt{13}{328.073}{295.412}{\Character{59}}%
\TextAt{13}{333.297}{295.412}{\Character{11}}%
\TextAt{21}{340.763}{293.619}{j}%
\TextAt{16}{345.145}{295.412}{\Character{105}}%
\TextAt{9}{349.794}{295.412}{\Character{41}}%
\TextAt{16}{354.346}{295.412}{\Character{106}}%
\TextAt{17}{357.667}{300.348}{1\Character{47}2}%
\TextAt{13}{370.868}{295.412}{\Character{58}}%
\TextAt{8}{79.226}{273.183}{It}%
\TextAt{8}{93.416}{273.183}{will}%
\TextAt{8}{117.038}{273.183}{be}%
\TextAt{8}{134.48}{273.183}{denoted}%
\TextAt{8}{180.219}{273.183}{by}%
\TextAt{8}{197.995}{273.183}{vol}%
\TextAt{9}{213.286}{273.183}{\Character{40}}%
\TextAt{13}{217.838}{273.183}{M}%
\TextAt{9}{230.412}{273.183}{\Character{41}}%
\TextAt{8}{240.379}{273.183}{for}%
\TextAt{8}{259.781}{273.183}{short.}%
\TextAt{8}{298.867}{273.183}{It}%
\TextAt{8}{313.057}{273.183}{does}%
\TextAt{8}{340.971}{273.183}{not}%
\TextAt{8}{363.302}{273.183}{depend}%
\TextAt{8}{405.455}{273.183}{on}%
\TextAt{8}{423.231}{273.183}{the}%
\TextAt{8}{444.905}{273.183}{choice}%
\TextAt{8}{481.212}{273.183}{of}%
\TextAt{18}{496.059}{273.183}{Z}%
\TextAt{8}{504.029}{273.183}{-basis}%
\TextAt{13}{79.226}{256.138}{\Character{11}}%
\TextAt{17}{86.672}{254.344}{1}%
\TextAt{13}{91.404}{256.137}{\Character{59}}%
\TextAt{13}{96.628}{256.137}{\Character{58}}%
\TextAt{13}{101.865}{256.137}{\Character{58}}%
\TextAt{13}{107.089}{256.137}{\Character{58}}%
\TextAt{13}{112.326}{256.137}{\Character{59}}%
\TextAt{13}{117.562}{256.137}{\Character{11}}%
\TextAt{21}{125.07}{254.344}{n}%
\TextAt{8}{134.216}{256.138}{because}%
\TextAt{8}{177.129}{256.138}{a}%
\TextAt{8}{186.496}{256.138}{different}%
\TextAt{8}{232.279}{256.138}{choice}%
\TextAt{8}{266.68}{256.138}{is}%
\TextAt{8}{278.056}{256.138}{linked}%
\TextAt{8}{312.134}{256.138}{to}%
\TextAt{8}{326.057}{256.138}{the}%
\TextAt{8}{345.825}{256.138}{original}%
\TextAt{8}{387.722}{256.138}{one}%
\TextAt{8}{408.793}{256.138}{by}%
\TextAt{8}{424.664}{256.138}{a}%
\TextAt{8}{434.031}{256.138}{matrix}%
\TextAt{8}{471.696}{256.138}{with}%
\TextAt{8}{497.968}{256.138}{integer}%
\TextAt{8}{79.226}{239.092}{coefficients}%
\TextAt{8}{138.853}{239.092}{which}%
\TextAt{8}{172.729}{239.092}{also}%
\TextAt{8}{196.599}{239.092}{has}%
\TextAt{8}{217.862}{239.092}{an}%
\TextAt{8}{234.51}{239.092}{inverse}%
\TextAt{8}{273.658}{239.092}{with}%
\TextAt{8}{300.708}{239.092}{integer}%
\TextAt{8}{339.796}{239.092}{coefficients,}%
\TextAt{8}{402.772}{239.092}{hence}%
\TextAt{8}{435.668}{239.092}{has}%
\TextAt{8}{456.931}{239.092}{determinant}%
\TextAt{8}{523.337}{239.092}{of}%
\TextAt{8}{79.226}{222.046}{absolute}%
\TextAt{8}{125.469}{222.046}{value}%
\TextAt{9}{155.715}{222.046}{1}%
\TextAt{8}{161.568}{222.046}{.}%
\TextAt{8}{97.159}{202.011}{A}%
\TextAt{8}{110.106}{202.011}{more}%
\TextAt{8}{139.647}{202.011}{elegant}%
\TextAt{8}{179.924}{202.011}{definition}%
\TextAt{8}{232.227}{202.011}{of}%
\TextAt{8}{245.832}{202.011}{vol}%
\TextAt{9}{261.123}{202.011}{\Character{40}}%
\TextAt{13}{265.675}{202.011}{M}%
\TextAt{9}{278.249}{202.011}{\Character{41}}%
\TextAt{8}{286.973}{202.011}{can}%
\TextAt{8}{308.707}{202.011}{be}%
\TextAt{8}{324.906}{202.011}{given}%
\TextAt{8}{355.75}{202.011}{in}%
\TextAt{8}{369.678}{202.011}{terms}%
\TextAt{8}{402.531}{202.011}{of}%
\TextAt{8}{416.135}{202.011}{the}%
\TextAt{8}{436.566}{202.011}{invertible}%
\TextAt{8}{488.547}{202.011}{metrized}%
\TextAt{18}{79.226}{184.965}{Z}%
\TextAt{8}{87.196}{184.965}{-module}%
\TextAt{8}{134.08}{184.965}{det}%
\TextAt{13}{152.332}{184.965}{M}%
\TextAt{8}{164.905}{184.965}{.}%
\TextAt{8}{178.434}{184.965}{det}%
\TextAt{13}{196.686}{184.965}{M}%
\TextAt{25}{208.011}{183.172}{R}%
\TextAt{8}{219.856}{184.965}{is}%
\TextAt{8}{233.313}{184.965}{a}%
\TextAt{8}{244.762}{184.965}{one-dimensional}%
\TextAt{18}{332.366}{184.965}{R}%
\TextAt{8}{341}{184.965}{-vector}%
\TextAt{8}{381.715}{184.965}{space}%
\TextAt{8}{414.684}{184.965}{with}%
\TextAt{8}{443.037}{184.965}{metric}%
\TextAt{16}{481.146}{184.965}{\Character{104}}%
\TextAt{13}{491.386}{184.965}{\Character{59}}%
\TextAt{16}{502.221}{184.965}{\Character{105}}%
\TextAt{7}{506.871}{183.172}{det}%
\TextAt{21}{520.038}{183.172}{M}%
\TextAt{8}{529.518}{184.965}{,}%
\TextAt{8}{79.226}{167.919}{and}%
\TextAt{8}{102.755}{167.919}{with}%
\TextAt{8}{130.182}{167.919}{the}%
\TextAt{8}{151.105}{167.919}{lattice}%
\TextAt{8}{187.641}{167.919}{det}%
\TextAt{13}{205.892}{167.919}{M}%
\TextAt{8}{223.13}{167.919}{isomorphic}%
\TextAt{8}{283.146}{167.919}{to}%
\TextAt{18}{298.223}{167.919}{Z}%
\TextAt{8}{306.193}{167.919}{.}%
\TextAt{8}{316.941}{167.919}{If}%
\TextAt{13}{329.399}{167.919}{x}%
\NativeMathAt{340.678}{167.919}{11.955}{\in}%
\TextAt{8}{353.275}{167.919}{det}%
\TextAt{13}{371.527}{167.919}{M}%
\TextAt{8}{388.764}{167.919}{is}%
\TextAt{8}{401.295}{167.919}{a}%
\TextAt{8}{411.816}{167.919}{generator}%
\TextAt{8}{464.624}{167.919}{(for}%
\TextAt{8}{487.83}{167.919}{instance,}%
\TextAt{13}{79.226}{150.873}{x}%
\TextAt{9}{89.199}{150.873}{\Character{61}}%
\TextAt{13}{101.625}{150.873}{\Character{11}}%
\TextAt{17}{109.07}{149.08}{1}%
\TextAt{16}{116.459}{150.873}{\Character{94}}%
\TextAt{16}{127.075}{150.873}{\Character{1}}%
\TextAt{16}{132.371}{150.873}{\Character{1}}%
\TextAt{16}{137.68}{150.873}{\Character{1}}%
\TextAt{16}{143.645}{150.873}{\Character{94}}%
\TextAt{13}{154.317}{150.873}{\Character{11}}%
\TextAt{21}{161.763}{149.08}{n}%
\TextAt{8}{167.399}{150.873}{),}%
\TextAt{8}{179.103}{150.873}{then}%
\TextAt{8}{218.289}{124.889}{vol}%
\TextAt{9}{233.58}{124.889}{\Character{40}}%
\TextAt{13}{238.132}{124.889}{M}%
\TextAt{9}{250.706}{124.889}{\Character{41}}%
\TextAt{9}{258.561}{124.889}{\Character{61}}%
\TextAt{16}{271.005}{124.889}{\Character{107}}%
\TextAt{13}{276.982}{124.889}{x}%
\TextAt{16}{283.634}{124.889}{\Character{107}}%
\TextAt{7}{289.612}{123.096}{det}%
\TextAt{21}{302.779}{123.096}{M}%
\TextAt{9}{315.58}{124.889}{\Character{61}}%
\TextAt{20}{328.005}{135.694}{\Character{112}}%
\Vector{q 0 G 0 g 0.478 w 267.961 -638.066 m 318.455 -638.066 l S Q}%
\TextAt{16}{339.961}{124.889}{\Character{104}}%
\TextAt{13}{344.61}{124.889}{x\Character{59}}%
\TextAt{13}{356.482}{124.889}{x}%
\TextAt{16}{363.158}{124.889}{\Character{105}}%
\TextAt{7}{367.807}{123.096}{det}%
\TextAt{21}{380.975}{123.096}{M}%
\TextAt{13}{390.455}{124.889}{\Character{58}}%
\TextAt{8}{79.226}{98.905}{In}%
\TextAt{8}{93.847}{98.905}{the}%
\TextAt{8}{114.004}{98.905}{present}%
\TextAt{8}{154.711}{98.905}{case,}%
\TextAt{8}{182.734}{98.905}{where}%
\TextAt{8}{216.543}{98.905}{the}%
\TextAt{8}{236.7}{98.905}{base}%
\TextAt{8}{262.774}{98.905}{field}%
\TextAt{8}{288.131}{98.905}{is}%
\TextAt{18}{299.895}{98.905}{Q}%
\TextAt{8}{309.193}{98.905}{,}%
\TextAt{8}{316.342}{98.905}{the}%
\TextAt{8}{336.499}{98.905}{degree}%
\TextAt{8}{372.914}{98.905}{map}%
\TextAt{8}{258.257}{72.921}{deg}%
\TextAt{9}{279.14}{72.921}{\Character{58}}%
\TextAt{8}{285.712}{72.921}{Pic}%
\TextAt{9}{302.127}{72.921}{\Character{40}}%
\Vector{q 0 G 0 g 0.478 w 234.679 -691.169 m 242.649 -691.169 l S Q}%
\TextAt{18}{306.679}{72.921}{Z}%
\TextAt{9}{314.649}{72.921}{\Character{41}}%
\NativeArrowAt{322.522}{72.921}{11.955}{19.26}{\longrightarrow}%
\TextAt{18}{345.104}{72.921}{R}%
\TextAt{8}{297.211}{43.033}{236}%
\Destination{78.23}{809.815}{nav.page.236}%
\Destination{79.36}{750.955}{nav.print.257}%
\Destination{79.36}{749.955}{nav.number.3.8.2}%
\Hotlink{102.194}{758.955}{25.078}{11.047}{nav.number.3.1.6}%
}

% Page 249 (body page 237)
\Page{%
\TextAt{2}{79.226}{796.943}{III.}%
\TextAt{2}{98.317}{796.943}{Riemann–Roch}%
\TextAt{2}{174.375}{796.943}{Theory}%
\Vector{q 0 G 0 g 0.299 w 7.226 18.729 m 460.769 18.729 l S Q}%
\TextAt{8}{79.226}{761.124}{is}%
\TextAt{8}{90.272}{761.124}{an}%
\TextAt{8}{105.813}{761.124}{isomorphism}%
\TextAt{8}{173.514}{761.124}{(see}%
\TextAt{8}{196.264}{761.124}{§1,}%
\TextAt{8}{214.375}{761.124}{exercise}%
\TextAt{8}{256.959}{761.124}{3),}%
\TextAt{8}{273.946}{761.124}{and}%
\TextAt{8}{295.991}{761.124}{we}%
\TextAt{8}{312.5}{761.124}{call}%
\TextAt{8}{333.241}{761.124}{the}%
\TextAt{8}{352.679}{761.124}{unique}%
\TextAt{8}{390.003}{761.124}{homomorphism}%
\TextAt{8}{471.954}{761.124}{arising}%
\TextAt{8}{509.027}{761.124}{from}%
\TextAt{8}{79.226}{744.078}{this,}%
\TextAt{13}{232.413}{727.032}{\Character{31}}%
\TextAt{9}{243.067}{727.032}{\Character{61}}%
\TextAt{8}{255.492}{727.032}{deg}%
\TextAt{16}{275.047}{727.032}{\Character{14}}%
\TextAt{8}{283.017}{727.032}{det}%
\TextAt{9}{302.597}{727.032}{\Character{58}}%
\TextAt{13}{309.17}{727.032}{K}%
\TextAt{17}{319.987}{731.968}{0}%
\TextAt{9}{324.719}{727.032}{\Character{40}}%
\Vector{q 0 G 0 g 0.478 w 257.272 -37.059 m 265.242 -37.059 l S Q}%
\TextAt{18}{329.272}{727.032}{Z}%
\TextAt{9}{337.242}{727.032}{\Character{41}}%
\NativeArrowAt{345.115}{727.032}{11.955}{19.26}{\longrightarrow}%
\TextAt{18}{367.697}{727.032}{R}%
\TextAt{13}{376.331}{727.032}{\Character{59}}%
\TextAt{19}{79.226}{705.136}{the}%
\TextAt{19}{98.653}{705.136}{Euler–Minkowski}%
\TextAt{19}{189.847}{705.136}{characteristic}%
\TextAt{8}{258.303}{705.136}{.}%
\TextAt{8}{266.751}{705.136}{It}%
\TextAt{8}{279.423}{705.136}{is}%
\TextAt{8}{291.187}{705.136}{computed}%
\TextAt{8}{345.165}{705.136}{explicitly}%
\TextAt{8}{395.891}{705.136}{as}%
\TextAt{8}{410.261}{705.136}{follows.}%
\TextAt{12}{79.226}{680.208}{Proposition}%
\TextAt{12}{151.937}{680.208}{(8.3).}%
\TextAt{8}{187.253}{680.208}{For}%
\TextAt{8}{208.222}{680.208}{a}%
\TextAt{8}{217.978}{680.208}{coherent}%
\TextAt{8}{265.129}{680.208}{metrized}%
\TextAt{18}{313.249}{680.208}{Z}%
\TextAt{8}{321.219}{680.208}{-module}%
\TextAt{13}{366.409}{680.208}{M}%
\TextAt{8}{378.983}{680.208}{,}%
\TextAt{8}{386.132}{680.208}{one}%
\TextAt{8}{407.592}{680.208}{has}%
\TextAt{13}{208.792}{655.001}{\Character{31}}%
\TextAt{9}{216.125}{655.001}{\Character{40}}%
\TextAt{13}{220.677}{655.001}{M}%
\TextAt{9}{233.251}{655.001}{\Character{41}}%
\TextAt{9}{241.106}{655.001}{\Character{61}}%
\TextAt{8}{253.55}{655.001}{log}%
\TextAt{8}{270.502}{655.001}{\char35\relax }%
\TextAt{13}{280.266}{655.001}{M}%
\TextAt{7}{291.59}{653.208}{tor}%
\TextAt{16}{305.569}{655.001}{\Character{0}}%
\TextAt{8}{317.524}{655.001}{log}%
\TextAt{8}{334.476}{655.001}{vol}%
\TextAt{9}{349.775}{655.001}{\Character{40}}%
\TextAt{13}{354.327}{655.001}{M}%
\TextAt{9}{366.901}{655.001}{\Character{47}}%
\TextAt{13}{372.754}{655.001}{M}%
\TextAt{7}{384.078}{653.208}{tor}%
\TextAt{9}{395.4}{655.001}{\Character{41}}%
\TextAt{13}{399.952}{655.001}{\Character{58}}%
\TextAt{8}{79.226}{629.794}{In}%
\TextAt{8}{95.884}{629.794}{this}%
\TextAt{8}{120.743}{629.794}{formula}%
\TextAt{13}{165.71}{629.794}{M}%
\TextAt{7}{177.035}{628}{tor}%
\TextAt{8}{194.29}{629.794}{denotes}%
\TextAt{8}{238.66}{629.794}{the}%
\TextAt{8}{260.853}{629.794}{torsion}%
\TextAt{8}{301.983}{629.794}{subgroup}%
\TextAt{8}{354.817}{629.794}{of}%
\TextAt{13}{370.183}{629.794}{M}%
\TextAt{8}{388.691}{629.794}{and}%
\TextAt{13}{413.49}{629.794}{M}%
\TextAt{9}{426.063}{629.794}{\Character{47}}%
\TextAt{13}{431.916}{629.794}{M}%
\TextAt{7}{443.241}{628}{tor}%
\TextAt{8}{460.496}{629.794}{the}%
\TextAt{8}{482.689}{629.794}{projective}%
\TextAt{8}{79.226}{612.748}{metrized}%
\TextAt{18}{127.346}{612.748}{Z}%
\TextAt{8}{135.316}{612.748}{-module}%
\TextAt{8}{180.506}{612.748}{which}%
\TextAt{8}{213.993}{612.748}{receives}%
\TextAt{8}{256.972}{612.748}{its}%
\TextAt{8}{273.29}{612.748}{metric}%
\TextAt{8}{309.706}{612.748}{from}%
\TextAt{13}{337.346}{612.748}{M}%
\TextAt{8}{353.817}{612.748}{via}%
\TextAt{13}{373.005}{612.748}{M}%
\TextAt{16}{388.236}{612.748}{\Character{10}}%
\TextAt{18}{400.191}{612.748}{R}%
\NativeMathAt{412.146}{612.748}{11.955}{\cong}%
\TextAt{13}{424.765}{612.748}{M}%
\TextAt{9}{437.339}{612.748}{\Character{47}}%
\TextAt{13}{443.192}{612.748}{M}%
\TextAt{7}{454.517}{610.955}{tor}%
\TextAt{16}{468.495}{612.748}{\Character{10}}%
\TextAt{18}{480.45}{612.748}{R}%
\TextAt{8}{489.084}{612.748}{.}%
\TextAt{19}{79.226}{579.937}{Proof}%
\TextAt{19}{110.62}{579.937}{of}%
\TextAt{19}{124.369}{579.937}{(8.3).}%
\TextAt{8}{159.039}{579.937}{If}%
\TextAt{13}{170.532}{579.937}{M}%
\TextAt{8}{186.804}{579.937}{is}%
\TextAt{8}{198.369}{579.937}{a}%
\TextAt{8}{207.926}{579.937}{finite}%
\TextAt{18}{237.639}{579.937}{Z}%
\TextAt{8}{245.609}{579.937}{-module,}%
\TextAt{8}{293.892}{579.937}{then}%
\TextAt{8}{320.354}{579.937}{the}%
\TextAt{8}{340.311}{579.937}{determinant}%
\TextAt{8}{406.129}{579.937}{of}%
\TextAt{8}{419.26}{579.937}{the}%
\TextAt{8}{439.218}{579.937}{class}%
\TextAt{9}{466.456}{579.937}{\Character{91}}%
\TextAt{13}{469.708}{579.937}{M}%
\TextAt{9}{482.281}{579.937}{\Character{93}}%
\NativeMathAt{488.854}{579.937}{11.955}{\in}%
\TextAt{13}{500.145}{579.937}{K}%
\TextAt{17}{510.962}{584.276}{0}%
\TextAt{9}{515.695}{579.937}{\Character{40}}%
\Vector{q 0 G 0 g 0.478 w 448.247 -184.153 m 456.217 -184.153 l S Q}%
\TextAt{18}{520.247}{579.937}{Z}%
\TextAt{9}{528.217}{579.937}{\Character{41}}%
\TextAt{8}{79.226}{562.891}{is}%
\TextAt{8}{90.99}{562.891}{computed}%
\TextAt{8}{144.968}{562.891}{from}%
\TextAt{8}{172.608}{562.891}{a}%
\TextAt{8}{182.363}{562.891}{free}%
\TextAt{8}{204.791}{562.891}{resolution}%
\TextAt{9}{234.759}{537.684}{0}%
\NativeArrowAt{243.933}{537.684}{11.955}{19.26}{\longrightarrow}%
\TextAt{13}{266.515}{537.684}{E}%
\NativeArrowAt{279.202}{537.684}{11.955}{19.26}{\longrightarrow}%
\TextAt{13}{301.784}{537.684}{F}%
\TextAt{21}{316.896}{544.461}{\Character{11}}%
\NativeArrowAt{314.309}{537.684}{11.955}{12.027}{\longrightarrow}%
\TextAt{13}{329.655}{537.684}{M}%
\NativeArrowAt{345.549}{537.684}{11.955}{19.26}{\longrightarrow}%
\TextAt{9}{368.132}{537.684}{0}%
\TextAt{13}{373.985}{537.684}{\Character{59}}%
\TextAt{8}{79.226}{512.477}{where}%
\TextAt{13}{114.36}{512.477}{F}%
\TextAt{9}{129.141}{512.477}{\Character{61}}%
\TextAt{18}{143.824}{512.477}{Z}%
\TextAt{21}{151.794}{516.816}{n}%
\TextAt{8}{162.652}{512.477}{and}%
\TextAt{13}{186.739}{512.477}{E}%
\TextAt{9}{201.683}{512.477}{\Character{61}}%
\TextAt{8}{216.365}{512.477}{ker}%
\TextAt{9}{232.301}{512.477}{\Character{40}}%
\TextAt{13}{236.854}{512.477}{\Character{11}}%
\TextAt{9}{244.375}{512.477}{\Character{41}}%
\NativeMathAt{254.506}{512.477}{11.955}{\cong}%
\TextAt{18}{269.382}{512.477}{Z}%
\TextAt{21}{277.352}{516.816}{n}%
\TextAt{8}{282.988}{512.477}{.}%
\TextAt{8}{295.41}{512.477}{If}%
\TextAt{8}{308.426}{512.477}{we}%
\TextAt{8}{326.978}{512.477}{equip}%
\TextAt{13}{359.84}{512.477}{F}%
\TextAt{16}{372.603}{512.477}{\Character{10}}%
\TextAt{18}{385.461}{512.477}{R}%
\TextAt{9}{399.673}{512.477}{\Character{61}}%
\TextAt{13}{414.356}{512.477}{E}%
\TextAt{16}{427.281}{512.477}{\Character{10}}%
\TextAt{18}{440.139}{512.477}{R}%
\TextAt{9}{454.351}{512.477}{\Character{61}}%
\TextAt{18}{469.034}{512.477}{R}%
\TextAt{21}{477.668}{516.816}{n}%
\TextAt{8}{488.526}{512.477}{with}%
\TextAt{8}{516.511}{512.477}{the}%
\TextAt{8}{79.226}{495.431}{standard}%
\TextAt{8}{129.151}{495.431}{metric,}%
\TextAt{8}{170.163}{495.431}{the}%
\TextAt{8}{191.395}{495.431}{sequence}%
\TextAt{8}{240.973}{495.431}{becomes}%
\TextAt{8}{288.602}{495.431}{a}%
\TextAt{8}{299.433}{495.431}{short}%
\TextAt{8}{330.493}{495.431}{exact}%
\TextAt{8}{362.461}{495.431}{sequence}%
\TextAt{8}{412.039}{495.431}{of}%
\TextAt{8}{426.444}{495.431}{metrized}%
\TextAt{18}{475.64}{495.431}{Z}%
\TextAt{8}{483.61}{495.431}{-modules,}%
\TextAt{8}{79.226}{478.385}{because}%
\TextAt{13}{122.528}{478.385}{M}%
\TextAt{16}{137.758}{478.385}{\Character{10}}%
\TextAt{18}{149.713}{478.385}{R}%
\TextAt{9}{161.668}{478.385}{\Character{61}}%
\TextAt{9}{174.078}{478.385}{0}%
\TextAt{8}{179.947}{478.385}{.}%
\TextAt{8}{188.395}{478.385}{We}%
\TextAt{8}{208.54}{478.385}{therefore}%
\TextAt{8}{257.639}{478.385}{have}%
\TextAt{8}{284.634}{478.385}{in}%
\TextAt{13}{298.287}{478.385}{K}%
\TextAt{17}{309.105}{482.724}{0}%
\TextAt{9}{313.837}{478.385}{\Character{40}}%
\Vector{q 0 G 0 g 0.478 w 246.389 -285.705 m 254.36 -285.705 l S Q}%
\TextAt{18}{318.389}{478.385}{Z}%
\TextAt{9}{326.36}{478.385}{\Character{41}}%
\TextAt{8}{330.912}{478.385}{:}%
\TextAt{9}{263.866}{453.178}{\Character{91}}%
\TextAt{13}{267.118}{453.178}{M}%
\TextAt{9}{279.692}{453.178}{\Character{93}}%
\TextAt{9}{286.243}{453.178}{\Character{61}}%
\TextAt{9}{298.665}{453.178}{\Character{91}}%
\TextAt{13}{301.941}{453.178}{F}%
\TextAt{9}{311.145}{453.178}{\Character{93}}%
\TextAt{16}{317.053}{453.178}{\Character{0}}%
\TextAt{9}{329.008}{453.178}{\Character{91}}%
\TextAt{13}{332.26}{453.178}{E}%
\TextAt{9}{341.626}{453.178}{\Character{93}}%
\TextAt{13}{344.878}{453.178}{\Character{58}}%
\ErrataAt{79.19}{427.971}{455.81}{Let $A$ be the matrix whose $j$-th column gives the coordinates of $e_j^{\prime}$ in the standard basis}%
\ErrataAt{79.19}{410.925}{455.81}{$e_1,\ldots,e_n$ of $F$, where $e_1^{\prime},\ldots,e_n^{\prime}$ is a $\mathbb Z$-basis of $E$. Then $x=e_1\wedge\cdots\wedge e_n$, resp.}%
\ErrataAt{79.19}{393.879}{455.81}{$x^{\prime}=e_1^{\prime}\wedge\cdots\wedge e_n^{\prime}$, is a generator of $\det F$, resp. $\det E$, and}%
\ErrataAt{199}{368.672}{245}{$x^{\prime}=\det(A)x,\qquad |\det(A)|=(F:E)=\#M.$}%
\TextAt{8}{79.226}{343.465}{The}%
\TextAt{8}{103.28}{343.465}{metric}%
\TextAt{16}{139.695}{343.465}{\Character{107}}%
\TextAt{16}{149.57}{343.465}{\Character{107}}%
\TextAt{8}{159.445}{343.465}{on}%
\TextAt{8}{175.704}{343.465}{det}%
\TextAt{13}{193.956}{343.465}{E}%
\TextAt{8}{207.22}{343.465}{is}%
\TextAt{8}{218.983}{343.465}{the}%
\TextAt{8}{239.14}{343.465}{same}%
\TextAt{8}{268.466}{343.465}{as}%
\TextAt{8}{282.836}{343.465}{that}%
\TextAt{8}{308.205}{343.465}{on}%
\TextAt{8}{324.464}{343.465}{det}%
\TextAt{13}{342.715}{343.465}{F}%
\TextAt{8}{351.919}{343.465}{,}%
\TextAt{8}{359.068}{343.465}{so}%
\TextAt{8}{373.438}{343.465}{that}%
\TextAt{13}{124.538}{318.258}{\Character{31}}%
\TextAt{9}{131.872}{318.258}{\Character{40}}%
\TextAt{13}{136.424}{318.258}{E}%
\TextAt{9}{145.79}{318.258}{\Character{41}}%
\TextAt{9}{153.645}{318.258}{\Character{61}}%
\TextAt{8}{166.089}{318.258}{deg}%
\TextAt{9}{183.651}{318.258}{\Character{40}}%
\TextAt{8}{188.203}{318.258}{det}%
\TextAt{13}{206.455}{318.258}{E}%
\TextAt{9}{215.821}{318.258}{\Character{41}}%
\TextAt{9}{223.676}{318.258}{\Character{61}}%
\TextAt{16}{236.12}{318.258}{\Character{0}}%
\TextAt{8}{247.411}{318.258}{log}%
\TextAt{16}{264.371}{318.258}{\Character{107}}%
\TextAt{13}{270.349}{318.258}{x}%
\TextAt{3}{277.001}{323.194}{\Character{48}}%
\TextAt{16}{279.796}{318.258}{\Character{107}}%
\TextAt{9}{289.094}{318.258}{\Character{61}}%
\TextAt{16}{301.52}{318.258}{\Character{0}}%
\TextAt{8}{312.811}{318.258}{log}%
\TextAt{9}{327.779}{318.258}{\Character{40}}%
\TextAt{8}{332.331}{318.258}{\char35\relax }%
\TextAt{13}{342.086}{318.258}{M}%
\TextAt{16}{354.66}{318.258}{\Character{107}}%
\TextAt{13}{360.638}{318.258}{x}%
\TextAt{16}{367.29}{318.258}{\Character{107}}%
\TextAt{9}{373.267}{318.258}{\Character{41}}%
\TextAt{9}{381.122}{318.258}{\Character{61}}%
\TextAt{16}{393.566}{318.258}{\Character{0}}%
\TextAt{8}{404.857}{318.258}{log}%
\TextAt{8}{421.81}{318.258}{\char35\relax }%
\TextAt{13}{431.573}{318.258}{M}%
\TextAt{9}{446.803}{318.258}{\Character{43}}%
\TextAt{13}{458.564}{318.258}{\Character{31}}%
\TextAt{9}{465.898}{318.258}{\Character{40}}%
\TextAt{13}{470.45}{318.258}{F}%
\TextAt{9}{479.654}{318.258}{\Character{41}}%
\TextAt{13}{484.206}{318.258}{\Character{59}}%
\TextAt{8}{79.226}{293.051}{and}%
\TextAt{8}{101.989}{293.051}{then}%
\TextAt{13}{182.585}{276.005}{\Character{31}}%
\TextAt{9}{189.918}{276.005}{\Character{40}}%
\TextAt{13}{194.471}{276.005}{M}%
\TextAt{9}{207.044}{276.005}{\Character{41}}%
\TextAt{9}{214.899}{276.005}{\Character{61}}%
\TextAt{13}{227.343}{276.005}{\Character{31}}%
\TextAt{9}{234.677}{276.005}{\Character{40}\Character{91}}%
\TextAt{13}{242.481}{276.005}{F}%
\TextAt{9}{251.684}{276.005}{\Character{93}}%
\TextAt{16}{257.592}{276.005}{\Character{0}}%
\TextAt{9}{269.548}{276.005}{\Character{91}}%
\TextAt{13}{272.799}{276.005}{E}%
\TextAt{9}{282.165}{276.005}{\Character{93}\Character{41}}%
\TextAt{9}{293.259}{276.005}{\Character{61}}%
\TextAt{13}{305.716}{276.005}{\Character{31}}%
\TextAt{9}{313.049}{276.005}{\Character{40}}%
\TextAt{13}{317.602}{276.005}{F}%
\TextAt{9}{326.805}{276.005}{\Character{41}}%
\TextAt{16}{334.014}{276.005}{\Character{0}}%
\TextAt{13}{345.969}{276.005}{\Character{31}}%
\TextAt{9}{353.303}{276.005}{\Character{40}}%
\TextAt{13}{357.855}{276.005}{E}%
\TextAt{9}{367.221}{276.005}{\Character{41}}%
\TextAt{9}{375.076}{276.005}{\Character{61}}%
\TextAt{8}{387.52}{276.005}{log}%
\TextAt{8}{404.473}{276.005}{\char35\relax }%
\TextAt{13}{414.236}{276.005}{M}%
\TextAt{13}{426.155}{276.005}{\Character{58}}%
\TextAt{8}{79.226}{254.11}{For}%
\TextAt{8}{100.196}{254.11}{an}%
\TextAt{8}{116.455}{254.11}{arbitrary}%
\TextAt{8}{166.224}{254.11}{coherent}%
\TextAt{8}{213.375}{254.11}{metrized}%
\TextAt{18}{261.495}{254.11}{Z}%
\TextAt{8}{269.465}{254.11}{-module}%
\TextAt{13}{314.655}{254.11}{M}%
\TextAt{8}{331.126}{254.11}{we}%
\TextAt{8}{348.354}{254.11}{have}%
\TextAt{8}{375.349}{254.11}{the}%
\TextAt{8}{395.505}{254.11}{direct}%
\TextAt{8}{428.669}{254.11}{sum}%
\TextAt{8}{453.44}{254.11}{decomposition}%
\TextAt{13}{252.672}{228.903}{M}%
\TextAt{9}{268.567}{228.903}{\Character{61}}%
\TextAt{13}{280.992}{228.903}{M}%
\TextAt{7}{292.317}{227.109}{tor}%
\TextAt{16}{306.295}{228.903}{\Character{8}}%
\TextAt{13}{318.25}{228.903}{M}%
\TextAt{9}{330.824}{228.903}{\Character{47}}%
\TextAt{13}{336.677}{228.903}{M}%
\TextAt{7}{348.002}{227.109}{tor}%
\TextAt{8}{79.226}{203.695}{into}%
\TextAt{8}{103.625}{203.695}{metrized}%
\TextAt{18}{152.4}{203.695}{Z}%
\TextAt{8}{160.37}{203.695}{-modules.}%
\TextAt{8}{216.694}{203.695}{If}%
\TextAt{13}{229.042}{203.695}{\Character{11}}%
\TextAt{17}{236.488}{201.902}{1}%
\TextAt{13}{241.22}{203.695}{\Character{59}}%
\TextAt{13}{246.444}{203.695}{\Character{58}}%
\TextAt{13}{251.681}{203.695}{\Character{58}}%
\TextAt{13}{256.905}{203.695}{\Character{58}}%
\TextAt{13}{262.142}{203.695}{\Character{59}}%
\TextAt{13}{267.378}{203.695}{\Character{11}}%
\TextAt{21}{274.886}{201.902}{n}%
\TextAt{8}{285.076}{203.695}{is}%
\TextAt{8}{297.495}{203.695}{a}%
\TextAt{8}{307.907}{203.695}{basis}%
\TextAt{8}{337.303}{203.695}{of}%
\TextAt{8}{351.289}{203.695}{the}%
\TextAt{8}{372.101}{203.695}{lattice}%
\TextAt{13}{408.526}{203.695}{M}%
\TextAt{9}{421.1}{203.695}{\Character{47}}%
\TextAt{13}{426.953}{203.695}{M}%
\TextAt{7}{438.278}{201.902}{tor}%
\TextAt{8}{449.599}{203.695}{,}%
\TextAt{8}{457.568}{203.695}{then}%
\TextAt{13}{484.884}{203.695}{x}%
\TextAt{9}{495.975}{203.695}{\Character{61}}%
\TextAt{13}{509.518}{203.695}{\Character{11}}%
\TextAt{17}{516.963}{201.902}{1}%
\TextAt{16}{524.799}{203.695}{\Character{94}}%
\TextAt{16}{79.226}{186.65}{\Character{1}}%
\TextAt{16}{84.522}{186.65}{\Character{1}}%
\TextAt{16}{89.83}{186.65}{\Character{1}}%
\TextAt{16}{94.851}{186.65}{\Character{94}}%
\TextAt{13}{104.567}{186.65}{\Character{11}}%
\TextAt{21}{112.013}{184.856}{n}%
\TextAt{8}{121.085}{186.65}{is}%
\TextAt{8}{132.387}{186.65}{a}%
\TextAt{8}{141.68}{186.65}{generator}%
\TextAt{8}{193.259}{186.65}{of}%
\TextAt{8}{206.127}{186.65}{det}%
\TextAt{13}{224.378}{186.65}{M}%
\TextAt{9}{236.952}{186.65}{\Character{47}}%
\TextAt{13}{242.805}{186.65}{M}%
\TextAt{7}{254.13}{184.856}{tor}%
\TextAt{8}{265.451}{186.65}{;}%
\TextAt{8}{272.292}{186.65}{then}%
\TextAt{13}{298.49}{186.65}{\Character{31}}%
\TextAt{9}{305.824}{186.65}{\Character{40}}%
\TextAt{13}{310.376}{186.65}{M}%
\TextAt{9}{322.95}{186.65}{\Character{47}}%
\TextAt{13}{328.803}{186.65}{M}%
\TextAt{7}{340.128}{184.856}{tor}%
\TextAt{9}{351.449}{186.65}{\Character{41}}%
\TextAt{9}{359.304}{186.65}{\Character{61}}%
\TextAt{8}{371.748}{186.65}{deg}%
\TextAt{9}{389.31}{186.65}{\Character{40}}%
\TextAt{8}{393.862}{186.65}{det}%
\TextAt{13}{412.114}{186.65}{M}%
\TextAt{9}{424.687}{186.65}{\Character{47}}%
\TextAt{13}{430.54}{186.65}{M}%
\TextAt{7}{441.865}{184.856}{tor}%
\TextAt{9}{453.187}{186.65}{\Character{41}}%
\TextAt{9}{461.042}{186.65}{\Character{61}}%
\TextAt{16}{473.485}{186.65}{\Character{0}}%
\TextAt{8}{484.776}{186.65}{log}%
\TextAt{16}{501.737}{186.65}{\Character{107}}%
\TextAt{13}{507.714}{186.65}{x}%
\TextAt{16}{514.366}{186.65}{\Character{107}}%
\TextAt{9}{523.665}{186.65}{\Character{61}}%
\TextAt{16}{79.226}{169.604}{\Character{0}}%
\TextAt{8}{90.517}{169.604}{log}%
\TextAt{8}{107.47}{169.604}{vol}%
\TextAt{9}{122.768}{169.604}{\Character{40}}%
\TextAt{13}{127.321}{169.604}{M}%
\TextAt{9}{139.894}{169.604}{\Character{47}}%
\TextAt{13}{145.747}{169.604}{M}%
\TextAt{7}{157.072}{167.81}{tor}%
\TextAt{9}{168.393}{169.604}{\Character{41}}%
\TextAt{8}{172.946}{169.604}{.}%
\TextAt{8}{181.394}{169.604}{We}%
\TextAt{8}{201.539}{169.604}{therefore}%
\TextAt{8}{250.638}{169.604}{conclude}%
\TextAt{8}{298.758}{169.604}{that}%
\TextAt{13}{145.412}{144.397}{\Character{31}}%
\TextAt{9}{152.745}{144.397}{\Character{40}}%
\TextAt{13}{157.297}{144.397}{M}%
\TextAt{9}{169.871}{144.397}{\Character{41}}%
\TextAt{9}{177.726}{144.397}{\Character{61}}%
\TextAt{13}{190.17}{144.397}{\Character{31}}%
\TextAt{9}{197.503}{144.397}{\Character{40}}%
\TextAt{13}{202.056}{144.397}{M}%
\TextAt{7}{213.38}{142.603}{tor}%
\TextAt{9}{224.702}{144.397}{\Character{41}}%
\TextAt{9}{231.899}{144.397}{\Character{43}}%
\TextAt{13}{243.672}{144.397}{\Character{31}}%
\TextAt{9}{251.006}{144.397}{\Character{40}}%
\TextAt{13}{255.558}{144.397}{M}%
\TextAt{9}{268.132}{144.397}{\Character{47}}%
\TextAt{13}{273.985}{144.397}{M}%
\TextAt{7}{285.31}{142.603}{tor}%
\TextAt{9}{296.631}{144.397}{\Character{41}}%
\TextAt{9}{304.486}{144.397}{\Character{61}}%
\TextAt{8}{316.93}{144.397}{log}%
\TextAt{8}{333.882}{144.397}{\char35\relax }%
\TextAt{13}{343.646}{144.397}{M}%
\TextAt{7}{354.97}{142.603}{tor}%
\TextAt{16}{368.949}{144.397}{\Character{0}}%
\TextAt{8}{380.904}{144.397}{log}%
\TextAt{8}{397.857}{144.397}{vol}%
\TextAt{9}{413.155}{144.397}{\Character{40}}%
\TextAt{13}{417.707}{144.397}{M}%
\TextAt{9}{430.281}{144.397}{\Character{47}}%
\TextAt{13}{436.134}{144.397}{M}%
\TextAt{7}{447.458}{142.603}{tor}%
\TextAt{9}{458.78}{144.397}{\Character{41}}%
\TextAt{13}{463.332}{144.397}{\Character{58}}%
\Vector{q 0 G 0 g 0.399 w 452.334 -654.81 m 452.334 -646.741 l S Q}%
\Vector{q 0 G 0 g 0.399 w 452.534 -646.94 m 459.707 -646.94 l S Q}%
\Vector{q 0 G 0 g 0.399 w 459.906 -654.81 m 459.906 -646.741 l S Q}%
\TextAt{8}{297.211}{43.033}{237}%
\Destination{78.23}{809.815}{nav.page.237}%
\Destination{79.36}{776.155}{nav.print.258}%
\Destination{79.36}{694.255}{nav.number.3.8.3}%
\Hotlink{201.443}{758.955}{6.858}{11.047}{nav.section.3.1}%
\Hotlink{123.869}{577.755}{26.106}{11.047}{nav.number.3.8.3}%
\Destination{79.226}{444.971}{errata.87}%
}

% Page 250 (body page 238)
\Page{%
\TextAt{2}{79.226}{796.943}{III.}%
\TextAt{2}{98.317}{796.943}{Riemann–Roch}%
\TextAt{2}{174.375}{796.943}{Theory}%
\Vector{q 0 G 0 g 0.299 w 7.226 18.729 m 460.769 18.729 l S Q}%
\TextAt{8}{97.159}{761.124}{The}%
\TextAt{8}{121.213}{761.124}{Euler–Minkowski}%
\TextAt{8}{212.813}{761.124}{characteristic}%
\TextAt{8}{284.747}{761.124}{of}%
\TextAt{8}{298.077}{761.124}{a}%
\TextAt{8}{307.833}{761.124}{replete}%
\TextAt{8}{346.197}{761.124}{ideal}%
\TextAt{14}{374.16}{761.124}{a}%
\TextAt{8}{380.143}{761.124}{,}%
\TextAt{13}{255.973}{737.481}{\Character{31}}%
\TextAt{9}{263.307}{737.481}{\Character{40}}%
\TextAt{14}{267.859}{737.481}{a}%
\TextAt{9}{273.842}{737.481}{\Character{41}}%
\TextAt{9}{281.697}{737.481}{\Character{61}}%
\TextAt{16}{294.141}{737.481}{\Character{0}}%
\TextAt{8}{305.432}{737.481}{log}%
\TextAt{8}{322.385}{737.481}{vol}%
\TextAt{9}{337.683}{737.481}{\Character{40}}%
\TextAt{14}{342.235}{737.481}{a}%
\TextAt{9}{348.218}{737.481}{\Character{41}}%
\TextAt{13}{352.771}{737.481}{\Character{59}}%
\TextAt{8}{79.226}{713.839}{which}%
\TextAt{8}{112.379}{713.839}{we}%
\TextAt{8}{129.274}{713.839}{defined}%
\TextAt{19}{169.253}{713.839}{ad}%
\TextAt{19}{185.035}{713.839}{hoc}%
\TextAt{8}{205.468}{713.839}{in}%
\TextAt{8}{218.787}{713.839}{§3}%
\TextAt{8}{233.888}{713.839}{via}%
\TextAt{8}{252.743}{713.839}{the}%
\TextAt{8}{272.566}{713.839}{Minkowski}%
\TextAt{8}{330.503}{713.839}{measure}%
\TextAt{8}{375.754}{713.839}{vol}%
\TextAt{9}{391.045}{713.839}{\Character{40}}%
\TextAt{14}{395.597}{713.839}{a}%
\TextAt{9}{401.58}{713.839}{\Character{41}}%
\TextAt{8}{409.697}{713.839}{now}%
\TextAt{8}{433.752}{713.839}{appears}%
\TextAt{8}{476.733}{713.839}{as}%
\TextAt{8}{490.769}{713.839}{a}%
\TextAt{8}{500.192}{713.839}{simple}%
\TextAt{8}{79.226}{696.793}{special}%
\TextAt{8}{117.856}{696.793}{case}%
\TextAt{8}{143.156}{696.793}{of}%
\TextAt{8}{157.014}{696.793}{the}%
\TextAt{8}{177.699}{696.793}{Euler–Minkowski}%
\TextAt{8}{269.828}{696.793}{characteristic}%
\TextAt{8}{342.291}{696.793}{for}%
\TextAt{8}{360.704}{696.793}{metrized}%
\TextAt{18}{409.353}{696.793}{Z}%
\TextAt{8}{417.323}{696.793}{-modules}%
\TextAt{8}{467.657}{696.793}{to}%
\TextAt{8}{482.495}{696.793}{which}%
\TextAt{8}{516.51}{696.793}{the}%
\TextAt{8}{79.226}{679.747}{detailed}%
\TextAt{8}{124.328}{679.747}{development}%
\TextAt{8}{193.173}{679.747}{of}%
\TextAt{8}{207.382}{679.747}{the}%
\TextAt{8}{228.418}{679.747}{theory}%
\TextAt{8}{266.048}{679.747}{has}%
\TextAt{8}{287.802}{679.747}{led}%
\TextAt{8}{307.535}{679.747}{us.}%
\TextAt{8}{329.74}{679.747}{Indeed,}%
\TextAt{8}{372.121}{679.747}{viewing}%
\TextAt{8}{415.597}{679.747}{the}%
\TextAt{8}{436.633}{679.747}{metrized}%
\TextAt{14}{485.632}{679.747}{o}%
\TextAt{8}{491.476}{679.747}{-module}%
\TextAt{13}{79.226}{662.701}{L}%
\TextAt{9}{87.191}{662.701}{\Character{40}}%
\TextAt{14}{91.743}{662.701}{a}%
\TextAt{9}{97.726}{662.701}{\Character{41}}%
\TextAt{8}{106.597}{662.701}{of}%
\TextAt{8}{120.348}{662.701}{rank}%
\TextAt{9}{147.765}{662.701}{1}%
\TextAt{8}{157.936}{662.701}{associated}%
\TextAt{8}{214.093}{662.701}{to}%
\TextAt{14}{228.824}{662.701}{a}%
\TextAt{8}{239.126}{662.701}{as}%
\TextAt{8}{253.918}{662.701}{the}%
\TextAt{8}{274.495}{662.701}{metrized}%
\TextAt{18}{323.036}{662.701}{Z}%
\TextAt{8}{331.006}{662.701}{-module}%
\TextAt{13}{376.618}{662.701}{i}%
\TextAt{3}{380.612}{660.908}{\Character{3}}%
\TextAt{13}{385.344}{662.701}{L}%
\TextAt{9}{393.309}{662.701}{\Character{40}}%
\TextAt{14}{397.861}{662.701}{a}%
\TextAt{9}{403.844}{662.701}{\Character{41}}%
\TextAt{8}{412.715}{662.701}{of}%
\TextAt{8}{426.466}{662.701}{rank}%
\TextAt{9}{453.883}{662.701}{\Character{91}}%
\TextAt{13}{457.134}{662.701}{K}%
\TextAt{9}{471.991}{662.701}{\Character{58}}%
\TextAt{18}{479.281}{662.701}{Q}%
\TextAt{9}{488.58}{662.701}{\Character{93}}%
\TextAt{8}{491.831}{662.701}{,}%
\TextAt{8}{499.507}{662.701}{we}%
\TextAt{8}{517.156}{662.701}{get}%
\TextAt{8}{79.226}{645.655}{the}%
\TextAt{12}{79.226}{621.174}{Proposition}%
\TextAt{12}{151.937}{621.174}{(8.4).}%
\TextAt{13}{187.253}{621.174}{\Character{31}}%
\TextAt{9}{194.587}{621.174}{\Character{40}}%
\TextAt{14}{199.139}{621.174}{a}%
\TextAt{9}{205.122}{621.174}{\Character{41}}%
\TextAt{9}{212.977}{621.174}{\Character{61}}%
\TextAt{13}{225.421}{621.174}{\Character{31}}%
\TextAt{9}{232.754}{621.174}{\Character{40}}%
\TextAt{13}{237.307}{621.174}{i}%
\TextAt{3}{241.3}{619.381}{\Character{3}}%
\TextAt{13}{246.032}{621.174}{L}%
\TextAt{9}{253.997}{621.174}{\Character{40}}%
\TextAt{14}{258.549}{621.174}{a}%
\TextAt{9}{264.532}{621.174}{\Character{41}\Character{41}}%
\TextAt{8}{273.637}{621.174}{.}%
\TextAt{19}{79.226}{589.257}{Proof.}%
\TextAt{8}{116}{589.257}{Let}%
\TextAt{14}{137.529}{589.257}{a}%
\TextAt{9}{147.786}{589.257}{\Character{61}}%
\TextAt{14}{161.165}{589.257}{a}%
\TextAt{21}{167.148}{587.464}{f}%
\TextAt{14}{172.593}{589.257}{a}%
\NativeMathAt{178.576}{587.464}{7.97}{\infty}%
\TextAt{9}{191.817}{589.257}{\Character{61}}%
\TextAt{14}{205.195}{589.257}{a}%
\TextAt{21}{211.178}{587.464}{f}%
\TextAt{20}{218.616}{598.224}{\Character{81}}%
\TextAt{27}{229.907}{585.77}{p}%
\TextAt{3}{233.896}{585.77}{\Character{106}}%
\NativeMathAt{236.247}{585.77}{7.97}{\infty}%
\TextAt{14}{247.207}{589.257}{p}%
\TextAt{21}{253.19}{593.596}{\Character{23}}%
\TextAt{28}{257.353}{592.599}{p}%
\TextAt{8}{261.934}{589.257}{.}%
\TextAt{8}{272.06}{589.257}{The}%
\TextAt{8}{296.673}{589.257}{metric}%
\TextAt{16}{333.648}{589.257}{\Character{104}}%
\TextAt{13}{342.754}{589.257}{\Character{59}}%
\TextAt{16}{352.455}{589.257}{\Character{105}}%
\TextAt{21}{357.104}{587.398}{i}%
\TextAt{23}{359.987}{586.401}{\Character{3}}%
\TextAt{21}{364.305}{587.398}{L}%
\TextAt{17}{370.064}{587.398}{\Character{40}}%
\TextAt{27}{373.357}{587.398}{a}%
\TextAt{17}{377.346}{587.398}{\Character{41}}%
\TextAt{8}{385.594}{589.257}{on}%
\TextAt{8}{402.412}{589.257}{the}%
\TextAt{18}{423.128}{589.257}{C}%
\TextAt{8}{431.762}{589.257}{-vector}%
\TextAt{8}{471.343}{589.257}{space}%
\TextAt{13}{503.177}{589.257}{K}%
\TextAt{25}{513.136}{587.464}{C}%
\TextAt{9}{523.665}{589.257}{\Character{61}}%
\TextAt{20}{79.226}{581.178}{\Character{81}}%
\TextAt{21}{90.517}{568.724}{\Character{28}}%
\NativeMathAt{95.196}{568.724}{7.97}{\in}%
\TextAt{21}{100.841}{568.724}{X}%
\TextAt{17}{108.428}{568.724}{\Character{40}}%
\TextAt{25}{111.721}{568.724}{C}%
\TextAt{17}{117.477}{568.724}{\Character{41}}%
\TextAt{18}{123.261}{572.211}{C}%
\TextAt{8}{135.793}{572.211}{is}%
\TextAt{8}{147.557}{572.211}{then}%
\TextAt{8}{174.217}{572.211}{given}%
\TextAt{8}{204.786}{572.211}{by}%
\TextAt{16}{238.079}{546.586}{\Character{104}}%
\TextAt{13}{242.728}{546.586}{x\Character{59}}%
\TextAt{13}{254.6}{546.586}{y}%
\TextAt{16}{260.761}{546.586}{\Character{105}}%
\TextAt{21}{265.411}{544.726}{i}%
\TextAt{23}{268.294}{543.73}{\Character{3}}%
\TextAt{21}{272.611}{544.726}{L}%
\TextAt{17}{278.37}{544.726}{\Character{40}}%
\TextAt{27}{281.663}{544.726}{a}%
\TextAt{17}{285.652}{544.726}{\Character{41}}%
\TextAt{9}{292.764}{546.586}{\Character{61}}%
\TextAt{20}{305.19}{557.944}{\Character{88}}%
\TextAt{21}{311.485}{532.838}{\Character{28}}%
\TextAt{13}{324.451}{546.586}{e}%
\TextAt{17}{329.876}{551.522}{2}%
\TextAt{21}{334.11}{551.522}{\Character{23}}%
\TextAt{28}{338.273}{550.526}{p}%
\TextAt{22}{341.858}{549.53}{\Character{28}}%
\TextAt{13}{347.523}{546.586}{x}%
\TextAt{21}{354.175}{544.793}{\Character{28}}%
\Vector{q 0 G 0 g 0.478 w 287.352 -220.593 m 293.488 -220.593 l S Q}%
\TextAt{13}{359.352}{546.586}{y}%
\TextAt{21}{365.488}{543.763}{\Character{28}}%
\TextAt{13}{370.665}{546.586}{\Character{59}}%
\TextAt{8}{79.226}{515.633}{where}%
\TextAt{14}{112.85}{515.633}{p}%
\TextAt{21}{118.833}{513.84}{\Character{28}}%
\TextAt{8}{127.721}{515.633}{is}%
\TextAt{8}{139.299}{515.633}{the}%
\TextAt{8}{159.27}{515.633}{infinite}%
\TextAt{8}{198.752}{515.633}{place}%
\TextAt{8}{228.478}{515.633}{of}%
\TextAt{13}{241.622}{515.633}{K}%
\TextAt{8}{256.152}{515.633}{corresponding}%
\TextAt{8}{331.152}{515.633}{to}%
\TextAt{8}{345.277}{515.633}{the}%
\TextAt{8}{365.248}{515.633}{embedding}%
\TextAt{13}{424.24}{515.633}{\Character{28}}%
\TextAt{9}{433.98}{515.633}{\Character{58}}%
\TextAt{13}{440.552}{515.633}{K}%
\NativeMathAt{454.691}{515.633}{11.955}{\rightarrow}%
\TextAt{18}{469.967}{515.633}{C}%
\TextAt{8}{478.601}{515.633}{.}%
\TextAt{8}{486.988}{515.633}{It}%
\TextAt{8}{499.474}{515.633}{results}%
\TextAt{8}{79.226}{498.587}{from}%
\TextAt{8}{106.866}{498.587}{the}%
\TextAt{8}{127.023}{498.587}{standard}%
\TextAt{8}{175.872}{498.587}{metric}%
\TextAt{16}{212.287}{498.587}{\Character{104}}%
\TextAt{13}{220.834}{498.587}{\Character{59}}%
\TextAt{16}{229.975}{498.587}{\Character{105}}%
\TextAt{8}{238.522}{498.587}{via}%
\TextAt{8}{257.71}{498.587}{the}%
\TextAt{13}{277.866}{498.587}{F}%
\TextAt{8}{287.07}{498.587}{-invariant}%
\TextAt{8}{340.091}{498.587}{transformation}%
\TextAt{13}{177.384}{474.945}{T}%
\TextAt{9}{189.192}{474.945}{\Character{58}}%
\TextAt{13}{195.764}{474.945}{K}%
\TextAt{25}{205.724}{473.151}{C}%
\NativeArrowAt{215.299}{474.945}{11.955}{19.26}{\longrightarrow}%
\TextAt{13}{237.881}{474.945}{K}%
\TextAt{25}{247.84}{473.151}{C}%
\TextAt{13}{254.094}{474.945}{\Character{59}}%
\TextAt{9}{283.249}{474.945}{\Character{40}}%
\TextAt{13}{287.801}{474.945}{x}%
\TextAt{21}{294.453}{473.151}{\Character{28}}%
\TextAt{9}{299.63}{474.945}{\Character{41}}%
\TextAt{21}{304.182}{473.085}{\Character{28}}%
\NativeMathAt{308.86}{473.085}{7.97}{\in}%
\TextAt{21}{314.506}{473.085}{X}%
\TextAt{17}{322.093}{473.085}{\Character{40}}%
\TextAt{25}{325.386}{473.085}{C}%
\TextAt{17}{331.142}{473.085}{\Character{41}}%
\NativeArrowAt{338.254}{474.945}{11.955}{19.26}{\longmapsto}%
\TextAt{9}{360.836}{474.945}{\Character{40}}%
\TextAt{13}{365.389}{474.945}{e}%
\TextAt{21}{370.814}{479.881}{\Character{23}}%
\TextAt{28}{374.977}{478.885}{p}%
\TextAt{22}{378.562}{477.888}{\Character{28}}%
\TextAt{13}{384.227}{474.945}{x}%
\TextAt{21}{390.879}{473.151}{\Character{28}}%
\TextAt{9}{396.056}{474.945}{\Character{41}}%
\TextAt{21}{400.608}{473.085}{\Character{28}}%
\NativeMathAt{405.286}{473.085}{7.97}{\in}%
\TextAt{21}{410.932}{473.085}{X}%
\TextAt{17}{418.519}{473.085}{\Character{40}}%
\TextAt{25}{421.812}{473.085}{C}%
\TextAt{17}{427.568}{473.085}{\Character{41}}%
\TextAt{13}{431.36}{474.945}{\Character{58}}%
\TextAt{8}{79.226}{451.302}{Equivalently,}%
\TextAt{16}{248.664}{434.256}{\Character{104}}%
\TextAt{13}{253.314}{434.256}{x\Character{59}}%
\TextAt{13}{265.185}{434.256}{y}%
\TextAt{16}{271.346}{434.256}{\Character{105}}%
\TextAt{21}{275.996}{432.397}{i}%
\TextAt{23}{278.879}{431.4}{\Character{3}}%
\TextAt{21}{283.196}{432.397}{L}%
\TextAt{17}{288.955}{432.397}{\Character{40}}%
\TextAt{27}{292.248}{432.397}{a}%
\TextAt{17}{296.237}{432.397}{\Character{41}}%
\TextAt{9}{303.349}{434.256}{\Character{61}}%
\TextAt{16}{315.775}{434.256}{\Character{104}}%
\TextAt{13}{320.424}{434.256}{T}%
\TextAt{13}{328.888}{434.256}{x\Character{59}}%
\TextAt{13}{340.772}{434.256}{T}%
\TextAt{13}{349.248}{434.256}{y}%
\TextAt{16}{355.431}{434.256}{\Character{105}}%
\TextAt{13}{360.08}{434.256}{\Character{58}}%
\TextAt{8}{79.226}{413.031}{The}%
\TextAt{8}{102.828}{413.031}{volume}%
\TextAt{8}{142.701}{413.031}{vol}%
\TextAt{9}{157.992}{413.031}{\Character{40}}%
\TextAt{13}{162.544}{413.031}{i}%
\TextAt{3}{166.537}{411.238}{\Character{3}}%
\TextAt{13}{171.27}{413.031}{L}%
\TextAt{9}{179.234}{413.031}{\Character{40}}%
\TextAt{14}{183.787}{413.031}{a}%
\TextAt{9}{189.77}{413.031}{\Character{41}\Character{41}}%
\TextAt{8}{202.32}{413.031}{of}%
\TextAt{8}{215.198}{413.031}{a}%
\TextAt{8}{224.502}{413.031}{fundamental}%
\TextAt{8}{291.692}{413.031}{mesh}%
\TextAt{8}{321.212}{413.031}{of}%
\TextAt{8}{334.09}{413.031}{the}%
\TextAt{8}{353.794}{413.031}{lattice}%
\TextAt{14}{389.112}{413.031}{a}%
\TextAt{21}{395.096}{411.238}{f}%
\TextAt{8}{403.986}{413.031}{in}%
\TextAt{13}{417.187}{413.031}{K}%
\TextAt{25}{427.146}{411.238}{R}%
\TextAt{8}{436.846}{413.031}{with}%
\TextAt{8}{463.054}{413.031}{respect}%
\TextAt{8}{502.652}{413.031}{to}%
\TextAt{8}{516.51}{413.031}{the}%
\TextAt{8}{79.226}{395.985}{Haar}%
\TextAt{8}{107.663}{395.985}{measure}%
\TextAt{8}{152.741}{395.985}{defined}%
\TextAt{8}{192.546}{395.985}{by}%
\TextAt{8}{208.298}{395.985}{the}%
\TextAt{8}{227.948}{395.985}{euclidean}%
\TextAt{8}{278.812}{395.985}{metric}%
\TextAt{8}{314.72}{395.985}{on}%
\TextAt{13}{330.472}{395.985}{K}%
\TextAt{25}{340.432}{394.192}{R}%
\TextAt{8}{350.076}{395.985}{is}%
\TextAt{8}{361.333}{395.985}{then}%
\TextAt{8}{387.486}{395.985}{the}%
\TextAt{8}{407.135}{395.985}{volume}%
\TextAt{8}{446.953}{395.985}{of}%
\TextAt{8}{459.776}{395.985}{a}%
\TextAt{8}{469.025}{395.985}{fundamental}%
\TextAt{8}{79.226}{378.939}{mesh}%
\TextAt{8}{109.198}{378.939}{of}%
\TextAt{8}{122.528}{378.939}{the}%
\TextAt{8}{142.684}{378.939}{lattice}%
\TextAt{13}{178.454}{378.939}{T}%
\TextAt{14}{186.941}{378.939}{a}%
\TextAt{21}{192.924}{377.146}{f}%
\TextAt{8}{202.266}{378.939}{with}%
\TextAt{8}{228.926}{378.939}{respect}%
\TextAt{8}{268.976}{378.939}{to}%
\TextAt{8}{283.286}{378.939}{the}%
\TextAt{8}{303.443}{378.939}{canonical}%
\TextAt{8}{354.826}{378.939}{measure}%
\TextAt{8}{400.411}{378.939}{defined}%
\TextAt{8}{440.724}{378.939}{by}%
\TextAt{16}{456.983}{378.939}{\Character{104}}%
\TextAt{13}{465.53}{378.939}{\Character{59}}%
\TextAt{16}{474.671}{378.939}{\Character{105}}%
\TextAt{8}{479.321}{378.939}{.}%
\TextAt{8}{487.769}{378.939}{Thus}%
\TextAt{8}{246.257}{355.297}{vol}%
\TextAt{9}{261.548}{355.297}{\Character{40}}%
\TextAt{13}{266.1}{355.297}{i}%
\TextAt{3}{270.094}{353.504}{\Character{3}}%
\TextAt{13}{274.826}{355.297}{L}%
\TextAt{9}{282.79}{355.297}{\Character{40}}%
\TextAt{14}{287.343}{355.297}{a}%
\TextAt{9}{293.326}{355.297}{\Character{41}\Character{41}}%
\TextAt{9}{305.724}{355.297}{\Character{61}}%
\TextAt{8}{318.177}{355.297}{vol}%
\TextAt{9}{333.467}{355.297}{\Character{40}}%
\TextAt{13}{338.02}{355.297}{T}%
\TextAt{14}{346.507}{355.297}{a}%
\TextAt{21}{352.49}{353.504}{f}%
\TextAt{9}{357.935}{355.297}{\Character{41}}%
\TextAt{13}{362.487}{355.297}{\Character{58}}%
\TextAt{8}{79.226}{331.655}{In}%
\TextAt{8}{93.847}{331.655}{the}%
\TextAt{8}{114.004}{331.655}{representation}%
\TextAt{13}{190.493}{331.655}{K}%
\TextAt{25}{200.452}{329.861}{R}%
\TextAt{9}{210.027}{331.655}{\Character{61}}%
\TextAt{20}{222.453}{340.621}{\Character{81}}%
\TextAt{27}{233.744}{328.168}{p}%
\TextAt{3}{237.732}{328.168}{\Character{106}}%
\NativeMathAt{240.083}{328.168}{7.97}{\infty}%
\TextAt{13}{251.044}{331.655}{K}%
\TextAt{27}{261.003}{329.861}{p}%
\TextAt{8}{265.49}{331.655}{,}%
\TextAt{8}{272.639}{331.655}{the}%
\TextAt{8}{292.795}{331.655}{canonical}%
\TextAt{8}{344.179}{331.655}{embedding}%
\TextAt{13}{220.326}{308.012}{K}%
\TextAt{25}{230.285}{306.219}{R}%
\TextAt{9}{239.86}{308.012}{\Character{61}}%
\TextAt{13}{252.285}{308.012}{K}%
\TextAt{16}{265.76}{308.012}{\Character{10}}%
\TextAt{25}{275.058}{306.219}{Q}%
\TextAt{18}{284.412}{308.012}{R}%
\NativeArrowAt{296.367}{308.012}{11.955}{19.26}{\longrightarrow}%
\TextAt{13}{318.949}{308.012}{K}%
\TextAt{25}{328.908}{306.219}{C}%
\TextAt{9}{338.484}{308.012}{\Character{61}}%
\TextAt{13}{350.909}{308.012}{K}%
\TextAt{16}{364.383}{308.012}{\Character{10}}%
\TextAt{25}{373.682}{306.219}{Q}%
\TextAt{18}{383.036}{308.012}{C}%
\TextAt{8}{79.226}{284.37}{maps}%
\TextAt{8}{110.749}{284.37}{an}%
\TextAt{8}{127.901}{284.37}{element}%
\TextAt{9}{172.037}{284.37}{\Character{40}}%
\TextAt{13}{176.589}{284.37}{x}%
\TextAt{27}{183.241}{282.577}{p}%
\TextAt{9}{187.728}{284.37}{\Character{41}}%
\TextAt{27}{192.28}{282.51}{p}%
\TextAt{3}{196.269}{282.51}{\Character{106}}%
\NativeMathAt{198.62}{282.51}{7.97}{\infty}%
\TextAt{8}{212.379}{284.37}{to}%
\TextAt{8}{227.582}{284.37}{the}%
\TextAt{8}{248.632}{284.37}{element}%
\TextAt{9}{292.768}{284.37}{\Character{40}}%
\TextAt{13}{297.32}{284.37}{x}%
\TextAt{21}{303.972}{282.577}{\Character{28}}%
\TextAt{9}{309.149}{284.37}{\Character{41}}%
\TextAt{21}{313.701}{282.51}{\Character{28}}%
\NativeMathAt{318.379}{282.51}{7.97}{\in}%
\TextAt{21}{324.025}{282.51}{X}%
\TextAt{17}{331.612}{282.51}{\Character{40}}%
\TextAt{25}{334.905}{282.51}{C}%
\TextAt{17}{340.661}{282.51}{\Character{41}}%
\TextAt{8}{349.243}{284.37}{with}%
\TextAt{13}{376.797}{284.37}{x}%
\TextAt{21}{383.449}{282.577}{\Character{28}}%
\TextAt{9}{393.469}{284.37}{\Character{61}}%
\TextAt{13}{407.417}{284.37}{\Character{28}}%
\TextAt{13}{413.825}{284.37}{x}%
\TextAt{27}{420.488}{282.577}{p}%
\TextAt{22}{424.477}{281.58}{\Character{28}}%
\TextAt{8}{429.644}{284.37}{.}%
\TextAt{8}{440.773}{284.37}{Here}%
\TextAt{8}{469.295}{284.37}{we}%
\TextAt{8}{487.416}{284.37}{extend}%
\TextAt{13}{526.351}{284.37}{\Character{28}}%
\TextAt{8}{79.226}{267.324}{to}%
\TextAt{13}{94.244}{267.324}{K}%
\TextAt{27}{104.203}{265.531}{p}%
\TextAt{22}{108.192}{264.534}{\Character{28}}%
\TextAt{8}{113.359}{267.324}{.}%
\TextAt{8}{123.931}{267.324}{The}%
\TextAt{8}{148.692}{267.324}{restriction}%
\TextAt{8}{205.398}{267.324}{of}%
\TextAt{8}{219.436}{267.324}{the}%
\TextAt{8}{240.3}{267.324}{transformation}%
\TextAt{13}{320.76}{267.324}{T}%
\TextAt{9}{333.774}{267.324}{\Character{58}}%
\TextAt{9}{341.533}{267.324}{\Character{40}}%
\TextAt{13}{346.105}{267.324}{x}%
\TextAt{21}{352.757}{265.531}{\Character{28}}%
\TextAt{9}{357.933}{267.324}{\Character{41}}%
\NativeArrowAt{367.013}{267.324}{11.955}{11.955}{\mapsto}%
\TextAt{9}{383.495}{267.324}{\Character{40}}%
\TextAt{13}{388.047}{267.324}{e}%
\TextAt{21}{393.472}{271.662}{\Character{23}}%
\TextAt{28}{397.635}{270.666}{p}%
\TextAt{22}{401.22}{269.67}{\Character{28}}%
\TextAt{13}{406.885}{267.324}{x}%
\TextAt{21}{413.537}{265.531}{\Character{28}}%
\TextAt{9}{418.714}{267.324}{\Character{41}}%
\TextAt{8}{427.871}{267.324}{to}%
\TextAt{13}{442.889}{267.324}{K}%
\TextAt{25}{452.848}{265.531}{R}%
\TextAt{9}{463.629}{267.324}{\Character{61}}%
\TextAt{20}{477.261}{276.29}{\Character{81}}%
\TextAt{27}{488.552}{263.837}{p}%
\TextAt{3}{492.541}{263.837}{\Character{106}}%
\NativeMathAt{494.892}{263.837}{7.97}{\infty}%
\TextAt{13}{505.852}{267.324}{K}%
\TextAt{27}{515.811}{265.531}{p}%
\TextAt{8}{524.903}{267.324}{is}%
\TextAt{8}{79.226}{250.278}{therefore}%
\TextAt{8}{129.223}{250.278}{given}%
\TextAt{8}{160.689}{250.278}{by}%
\TextAt{9}{177.846}{250.278}{\Character{40}}%
\TextAt{13}{182.398}{250.278}{x}%
\TextAt{27}{189.05}{248.485}{p}%
\TextAt{9}{193.537}{250.278}{\Character{41}}%
\NativeArrowAt{202.939}{250.278}{11.955}{11.955}{\mapsto}%
\TextAt{9}{219.743}{250.278}{\Character{40}}%
\TextAt{13}{224.296}{250.278}{e}%
\TextAt{21}{229.721}{254.616}{\Character{23}}%
\TextAt{28}{233.884}{253.62}{p}%
\TextAt{13}{238.465}{250.278}{x}%
\TextAt{27}{245.117}{248.485}{p}%
\TextAt{9}{249.604}{250.278}{\Character{41}}%
\TextAt{8}{254.156}{250.278}{.}%
\TextAt{8}{265.295}{250.278}{The}%
\TextAt{8}{290.246}{250.278}{lattice}%
\TextAt{13}{326.913}{250.278}{T}%
\TextAt{14}{335.4}{250.278}{a}%
\TextAt{21}{341.383}{248.485}{f}%
\TextAt{8}{351.622}{250.278}{is}%
\TextAt{8}{364.283}{250.278}{then}%
\TextAt{8}{391.84}{250.278}{the}%
\TextAt{8}{412.894}{250.278}{same}%
\TextAt{8}{443.117}{250.278}{lattice}%
\TextAt{8}{479.784}{250.278}{which}%
\TextAt{8}{514.167}{250.278}{was}%
\TextAt{8}{79.226}{233.232}{denoted}%
\TextAt{14}{123.448}{233.232}{a}%
\TextAt{8}{133.329}{233.232}{in}%
\TextAt{8}{146.982}{233.232}{§3.}%
\TextAt{8}{166.967}{233.232}{So}%
\TextAt{8}{183.226}{233.232}{we}%
\TextAt{8}{200.453}{233.232}{obtain}%
\TextAt{8}{253.223}{209.59}{vol}%
\TextAt{9}{268.514}{209.59}{\Character{40}}%
\TextAt{13}{273.066}{209.59}{i}%
\TextAt{3}{277.059}{207.796}{\Character{3}}%
\TextAt{13}{281.792}{209.59}{L}%
\TextAt{9}{289.756}{209.59}{\Character{40}}%
\TextAt{14}{294.309}{209.59}{a}%
\TextAt{9}{300.292}{209.59}{\Character{41}\Character{41}}%
\TextAt{9}{312.69}{209.59}{\Character{61}}%
\TextAt{8}{325.143}{209.59}{vol}%
\TextAt{9}{340.433}{209.59}{\Character{40}}%
\TextAt{14}{344.986}{209.59}{a}%
\TextAt{9}{350.969}{209.59}{\Character{41}}%
\TextAt{13}{355.521}{209.59}{\Character{59}}%
\TextAt{8}{79.226}{185.947}{i.e.,}%
\TextAt{13}{101.331}{185.947}{\Character{31}}%
\TextAt{9}{108.665}{185.947}{\Character{40}}%
\TextAt{13}{113.217}{185.947}{i}%
\TextAt{3}{117.211}{184.154}{\Character{3}}%
\TextAt{13}{121.943}{185.947}{L}%
\TextAt{9}{129.907}{185.947}{\Character{40}}%
\TextAt{14}{134.46}{185.947}{a}%
\TextAt{9}{140.443}{185.947}{\Character{41}\Character{41}}%
\TextAt{9}{152.841}{185.947}{\Character{61}}%
\TextAt{13}{165.294}{185.947}{\Character{31}}%
\TextAt{9}{172.627}{185.947}{\Character{40}}%
\TextAt{14}{177.18}{185.947}{a}%
\TextAt{9}{183.163}{185.947}{\Character{41}}%
\TextAt{8}{187.715}{185.947}{.}%
\Vector{q 0 G 0 g 0.399 w 452.334 -588.053 m 452.334 -579.983 l S Q}%
\Vector{q 0 G 0 g 0.399 w 452.534 -580.182 m 459.707 -580.182 l S Q}%
\Vector{q 0 G 0 g 0.399 w 452.534 -587.854 m 459.707 -587.854 l S Q}%
\Vector{q 0 G 0 g 0.399 w 459.906 -588.053 m 459.906 -579.983 l S Q}%
\TextAt{8}{97.159}{158.477}{Given}%
\TextAt{8}{131.967}{158.477}{this}%
\TextAt{8}{155.704}{158.477}{identification,}%
\TextAt{8}{230.671}{158.477}{the}%
\TextAt{8}{251.742}{158.477}{Riemann–Roch}%
\TextAt{8}{334.287}{158.477}{theorem}%
\TextAt{8}{380.728}{158.477}{(3.4)}%
\TextAt{8}{409.618}{158.477}{proven}%
\TextAt{8}{448.586}{158.477}{in}%
\TextAt{8}{463.154}{158.477}{§3}%
\TextAt{8}{479.503}{158.477}{for}%
\TextAt{8}{498.303}{158.477}{replete}%
\TextAt{8}{79.226}{141.431}{ideals}%
\TextAt{14}{111.804}{141.431}{a}%
\TextAt{8}{117.787}{141.431}{,}%
\TextAt{13}{250.613}{124.385}{\Character{31}}%
\TextAt{9}{257.947}{124.385}{\Character{40}}%
\TextAt{14}{262.499}{124.385}{a}%
\TextAt{9}{268.482}{124.385}{\Character{41}}%
\TextAt{9}{276.337}{124.385}{\Character{61}}%
\TextAt{8}{288.781}{124.385}{deg}%
\TextAt{9}{306.343}{124.385}{\Character{40}}%
\TextAt{14}{310.895}{124.385}{a}%
\TextAt{9}{316.878}{124.385}{\Character{41}}%
\TextAt{9}{324.075}{124.385}{\Character{43}}%
\TextAt{13}{335.849}{124.385}{\Character{31}}%
\TextAt{9}{343.182}{124.385}{\Character{40}}%
\TextAt{14}{347.734}{124.385}{o}%
\TextAt{9}{353.579}{124.385}{\Character{41}}%
\TextAt{13}{358.131}{124.385}{\Character{59}}%
\TextAt{8}{79.226}{103.16}{now}%
\TextAt{8}{103.615}{103.16}{appears}%
\TextAt{8}{146.928}{103.16}{as}%
\TextAt{8}{161.298}{103.16}{a}%
\TextAt{8}{171.054}{103.16}{special}%
\TextAt{8}{209.155}{103.16}{case}%
\TextAt{8}{233.926}{103.16}{of}%
\TextAt{8}{247.256}{103.16}{theorem}%
\TextAt{8}{292.781}{103.16}{(8.2),}%
\TextAt{8}{324.008}{103.16}{which}%
\TextAt{8}{357.495}{103.16}{says}%
\TextAt{8}{382.338}{103.16}{that}%
\TextAt{13}{224.818}{79.518}{\Character{31}}%
\TextAt{9}{232.152}{79.518}{\Character{40}}%
\TextAt{13}{236.704}{79.518}{i}%
\TextAt{3}{240.697}{77.724}{\Character{3}}%
\TextAt{13}{245.43}{79.518}{L}%
\TextAt{9}{253.394}{79.518}{\Character{40}}%
\TextAt{14}{257.947}{79.518}{a}%
\TextAt{9}{263.93}{79.518}{\Character{41}\Character{41}}%
\TextAt{9}{276.327}{79.518}{\Character{61}}%
\TextAt{8}{288.781}{79.518}{deg}%
\TextAt{9}{306.343}{79.518}{\Character{40}}%
\TextAt{13}{310.895}{79.518}{L}%
\TextAt{9}{318.86}{79.518}{\Character{40}}%
\TextAt{14}{323.412}{79.518}{a}%
\TextAt{9}{329.395}{79.518}{\Character{41}\Character{41}}%
\TextAt{9}{341.135}{79.518}{\Character{43}}%
\TextAt{13}{352.918}{79.518}{\Character{31}}%
\TextAt{9}{360.251}{79.518}{\Character{40}}%
\TextAt{13}{364.804}{79.518}{i}%
\TextAt{3}{368.797}{77.724}{\Character{3}}%
\TextAt{14}{373.529}{79.518}{o}%
\TextAt{9}{379.374}{79.518}{\Character{41}}%
\TextAt{13}{383.926}{79.518}{\Character{58}}%
\TextAt{8}{297.211}{43.033}{238}%
\Destination{78.23}{809.815}{nav.page.238}%
\Destination{79.36}{776.155}{nav.print.259}%
\Destination{79.36}{224.555}{nav.print.260}%
\Destination{79.36}{635.155}{nav.number.3.8.4}%
\Hotlink{223.966}{711.655}{6.858}{11.047}{nav.section.3.3}%
\Hotlink{152.161}{231.055}{6.858}{11.047}{nav.section.3.3}%
\Hotlink{468.333}{156.255}{6.858}{11.047}{nav.section.3.3}%
\Hotlink{380.228}{156.255}{25.078}{11.047}{nav.number.3.3.4}%
\Hotlink{292.281}{100.955}{25.078}{11.047}{nav.number.3.8.2}%
}


% Chapter 4
% Page 251 (body page 239)
\Page{%
\TextAt{1}{79.226}{673.054}{Chapter}%
\TextAt{1}{188.19}{673.054}{IV}%
\TextAt{1}{79.226}{614.873}{Abstract}%
\TextAt{1}{195.155}{614.873}{Class}%
\TextAt{1}{267.93}{614.873}{Field}%
\TextAt{1}{338.474}{614.873}{Theory}%
\TextAt{15}{79.226}{546.17}{§}%
\TextAt{15}{93.785}{546.17}{1}%
\TextAt{15}{120.693}{546.17}{Infinite}%
\TextAt{15}{188.401}{546.17}{Galois}%
\TextAt{15}{247.605}{546.17}{Theory}%
\TextAt{8}{79.226}{514.284}{Every}%
\TextAt{8}{114.101}{514.284}{field}%
\TextAt{13}{140.679}{514.284}{k}%
\TextAt{8}{152.286}{514.284}{is}%
\TextAt{8}{165.271}{514.284}{equipped}%
\TextAt{8}{216.56}{514.284}{with}%
\TextAt{8}{244.441}{514.284}{a}%
\TextAt{8}{255.417}{514.284}{distinguished}%
\TextAt{8}{327.652}{514.284}{Galois}%
\TextAt{8}{364.786}{514.284}{extension:}%
\TextAt{8}{423.544}{514.284}{the}%
\TextAt{8}{444.921}{514.284}{separable}%
\TextAt{8}{497.585}{514.284}{closure}%
\TextAt{9}{79.544}{500.393}{\Character{22}}%
\TextAt{13}{79.226}{497.238}{k}%
\TextAt{16}{90.382}{497.238}{\Character{106}}%
\TextAt{13}{98.368}{497.238}{k}%
\TextAt{8}{104.858}{497.238}{.}%
\TextAt{8}{115.674}{497.238}{Its}%
\TextAt{8}{133.751}{497.238}{Galois}%
\TextAt{8}{170.453}{497.238}{group}%
\TextAt{13}{204.418}{497.238}{G}%
\TextAt{21}{213.652}{495.445}{k}%
\TextAt{9}{223.437}{497.238}{\Character{61}}%
\TextAt{13}{237.208}{497.238}{G}%
\TextAt{9}{246.442}{497.238}{\Character{40}}%
\TextAt{9}{251.312}{500.393}{\Character{22}}%
\TextAt{13}{250.994}{497.238}{k}%
\TextAt{16}{262.149}{497.238}{\Character{106}}%
\TextAt{13}{270.136}{497.238}{k}%
\TextAt{9}{276.626}{497.238}{\Character{41}}%
\TextAt{8}{285.865}{497.238}{is}%
\TextAt{8}{298.418}{497.238}{called}%
\TextAt{8}{332.371}{497.238}{the}%
\TextAt{12}{353.317}{497.238}{absolute}%
\TextAt{12}{407.341}{497.238}{Galois}%
\TextAt{12}{449.363}{497.238}{group}%
\TextAt{8}{487.955}{497.238}{of}%
\TextAt{13}{502.074}{497.238}{k}%
\TextAt{8}{508.563}{497.238}{.}%
\TextAt{8}{519.38}{497.238}{As}%
\TextAt{8}{79.226}{480.192}{a}%
\TextAt{8}{89.606}{480.192}{rule,}%
\TextAt{8}{117.047}{480.192}{this}%
\TextAt{8}{140.494}{480.192}{extension}%
\TextAt{8}{192.885}{480.192}{will}%
\TextAt{8}{215.615}{480.192}{have}%
\TextAt{8}{243.235}{480.192}{infinite}%
\TextAt{8}{283.527}{480.192}{degree.}%
\TextAt{8}{326.367}{480.192}{It}%
\TextAt{8}{339.664}{480.192}{does,}%
\TextAt{8}{370.094}{480.192}{however,}%
\TextAt{8}{419.007}{480.192}{have}%
\TextAt{8}{446.626}{480.192}{the}%
\TextAt{8}{467.407}{480.192}{advantage}%
\TextAt{8}{523.337}{480.192}{of}%
\TextAt{8}{79.226}{463.147}{collecting}%
\TextAt{8}{131.324}{463.147}{all}%
\TextAt{8}{147.652}{463.147}{finite}%
\TextAt{8}{177.633}{463.147}{Galois}%
\TextAt{8}{213.615}{463.147}{extensions}%
\TextAt{8}{270.065}{463.147}{of}%
\TextAt{13}{283.464}{463.147}{k}%
\TextAt{8}{289.954}{463.147}{.}%
\TextAt{8}{298.609}{463.147}{This}%
\TextAt{8}{325.398}{463.147}{is}%
\TextAt{8}{337.231}{463.147}{why}%
\TextAt{8}{362.011}{463.147}{it}%
\TextAt{8}{373.784}{463.147}{is}%
\TextAt{8}{385.617}{463.147}{reasonable}%
\TextAt{8}{442.988}{463.147}{to}%
\TextAt{8}{457.367}{463.147}{try}%
\TextAt{8}{476.624}{463.147}{to}%
\TextAt{8}{491.003}{463.147}{give}%
\TextAt{8}{515.138}{463.147}{it}%
\TextAt{8}{526.911}{463.147}{a}%
\TextAt{8}{79.226}{446.101}{prominent}%
\TextAt{8}{135.377}{446.101}{place}%
\TextAt{8}{165.179}{446.101}{in}%
\TextAt{8}{178.722}{446.101}{Galois}%
\TextAt{8}{214.525}{446.101}{theory.}%
\TextAt{8}{254.809}{446.101}{But}%
\TextAt{8}{277.94}{446.101}{such}%
\TextAt{8}{304.227}{446.101}{an}%
\TextAt{8}{320.376}{446.101}{attempt}%
\TextAt{8}{365.146}{446.101}{faces}%
\TextAt{8}{393.382}{446.101}{the}%
\TextAt{8}{413.428}{446.101}{difficulty}%
\TextAt{8}{462.095}{446.101}{that}%
\TextAt{8}{487.354}{446.101}{the}%
\TextAt{8}{507.401}{446.101}{main}%
\TextAt{8}{79.226}{429.055}{theorem}%
\TextAt{8}{124.529}{429.055}{of}%
\TextAt{8}{137.637}{429.055}{Galois}%
\TextAt{8}{173.329}{429.055}{theory}%
\TextAt{8}{209.857}{429.055}{does}%
\TextAt{8}{236.032}{429.055}{not}%
\TextAt{8}{256.624}{429.055}{remain}%
\TextAt{8}{295.423}{429.055}{true}%
\TextAt{8}{319.912}{429.055}{for}%
\TextAt{8}{337.575}{429.055}{infinite}%
\TextAt{8}{377.02}{429.055}{extensions.}%
\TextAt{8}{437.878}{429.055}{Let}%
\TextAt{8}{458.625}{429.055}{us}%
\TextAt{8}{473.419}{429.055}{explain}%
\TextAt{8}{513.844}{429.055}{this}%
\TextAt{8}{79.226}{412.009}{in}%
\TextAt{8}{92.879}{412.009}{the}%
\TextAt{8}{113.035}{412.009}{following}%
\TextAt{12}{97.159}{391.974}{Example:}%
\TextAt{8}{156.689}{391.974}{The}%
\TextAt{8}{179.877}{391.974}{absolute}%
\TextAt{8}{225.254}{391.974}{Galois}%
\TextAt{8}{260.302}{391.974}{group}%
\TextAt{13}{292.611}{391.974}{G}%
\TextAt{25}{301.845}{390.181}{F}%
\TextAt{22}{306.716}{389.184}{p}%
\TextAt{9}{314.873}{391.974}{\Character{61}}%
\TextAt{13}{327.299}{391.974}{G}%
\TextAt{9}{336.532}{391.974}{\Character{40}}%
\Vector{q 0 G 0 g 0.478 w 269.085 -372.117 m 276.391 -372.117 l S Q}%
\TextAt{18}{341.085}{391.974}{F}%
\TextAt{21}{348.391}{390.181}{p}%
\TextAt{16}{356.472}{391.974}{\Character{106}}%
\TextAt{18}{363.114}{391.974}{F}%
\TextAt{21}{370.42}{390.181}{p}%
\TextAt{9}{375.181}{391.974}{\Character{41}}%
\TextAt{8}{382.765}{391.974}{of}%
\TextAt{8}{395.229}{391.974}{the}%
\TextAt{8}{414.52}{391.974}{field}%
\TextAt{18}{439.011}{391.974}{F}%
\TextAt{21}{446.317}{390.181}{p}%
\TextAt{8}{454.109}{391.974}{with}%
\TextAt{13}{479.904}{391.974}{p}%
\TextAt{8}{488.81}{391.974}{elements}%
\TextAt{8}{79.226}{374.928}{contains}%
\TextAt{8}{125.146}{374.928}{the}%
\TextAt{8}{145.302}{374.928}{Frobenius}%
\TextAt{8}{199.172}{374.928}{automorphism}%
\TextAt{13}{276.642}{374.928}{\Character{39}}%
\TextAt{8}{288.208}{374.928}{which}%
\TextAt{8}{321.694}{374.928}{is}%
\TextAt{8}{333.458}{374.928}{given}%
\TextAt{8}{364.027}{374.928}{by}%
\TextAt{13}{244.743}{349.962}{x}%
\TextAt{21}{251.395}{354.898}{\Character{39}}%
\TextAt{9}{260.732}{349.962}{\Character{61}}%
\TextAt{13}{273.158}{349.962}{x}%
\TextAt{21}{279.81}{354.898}{p}%
\TextAt{8}{296.526}{349.962}{for}%
\TextAt{8}{314.411}{349.962}{all}%
\TextAt{13}{330.67}{349.962}{x}%
\NativeMathAt{340.643}{349.962}{11.955}{\in}%
\Vector{q 0 G 0 g 0.478 w 279.934 -414.129 m 287.24 -414.129 l S Q}%
\TextAt{18}{351.934}{349.962}{F}%
\TextAt{21}{359.24}{348.168}{p}%
\TextAt{13}{364.001}{349.962}{\Character{58}}%
\TextAt{8}{79.226}{324.995}{The}%
\TextAt{8}{103.645}{324.995}{subgroup}%
\TextAt{16}{154.808}{324.995}{\Character{104}}%
\TextAt{13}{159.457}{324.995}{\Character{39}}%
\TextAt{16}{167.125}{324.995}{\Character{105}}%
\TextAt{9}{175.717}{324.995}{\Character{61}}%
\TextAt{16}{188.765}{324.995}{\Character{102}}%
\TextAt{13}{194.743}{324.995}{\Character{39}}%
\TextAt{21}{202.411}{329.334}{n}%
\TextAt{16}{211.99}{324.995}{\Character{106}}%
\TextAt{13}{219.254}{324.995}{n}%
\NativeMathAt{230.185}{324.995}{11.955}{\in}%
\TextAt{18}{242.098}{324.995}{Z}%
\TextAt{16}{250.068}{324.995}{\Character{103}}%
\TextAt{8}{260.308}{324.995}{has}%
\TextAt{8}{281.547}{324.995}{the}%
\TextAt{8}{302.068}{324.995}{same}%
\TextAt{8}{331.759}{324.995}{fixed}%
\TextAt{8}{360.41}{324.995}{field}%
\TextAt{18}{386.132}{324.995}{F}%
\TextAt{21}{393.438}{323.202}{p}%
\TextAt{8}{402.461}{324.995}{as}%
\TextAt{8}{417.196}{324.995}{the}%
\TextAt{8}{437.718}{324.995}{whole}%
\TextAt{8}{471.247}{324.995}{of}%
\TextAt{13}{484.942}{324.995}{G}%
\TextAt{25}{494.175}{323.202}{F}%
\TextAt{22}{499.046}{322.206}{p}%
\TextAt{8}{503.883}{324.995}{.}%
\TextAt{8}{513.426}{324.995}{But}%
\TextAt{8}{79.226}{307.949}{contrary}%
\TextAt{8}{126.552}{307.949}{to}%
\TextAt{8}{141.348}{307.949}{what}%
\TextAt{8}{171.1}{307.949}{we}%
\TextAt{8}{188.813}{307.949}{are}%
\TextAt{8}{208.809}{307.949}{used}%
\TextAt{8}{236.015}{307.949}{to}%
\TextAt{8}{250.811}{307.949}{in}%
\TextAt{8}{264.949}{307.949}{finite}%
\TextAt{8}{295.346}{307.949}{Galois}%
\TextAt{8}{331.745}{307.949}{theory,}%
\TextAt{8}{371.374}{307.949}{we}%
\TextAt{8}{389.087}{307.949}{find}%
\TextAt{16}{412.981}{307.949}{\Character{104}}%
\TextAt{13}{417.63}{307.949}{\Character{39}}%
\TextAt{16}{425.298}{307.949}{\Character{105}}%
\NativeMathAt{434.085}{307.949}{11.955}{\neq}%
\TextAt{13}{447.349}{307.949}{G}%
\TextAt{25}{456.582}{306.156}{F}%
\TextAt{22}{461.453}{305.16}{p}%
\TextAt{8}{466.29}{307.949}{.}%
\TextAt{8}{476.195}{307.949}{In}%
\TextAt{8}{491.302}{307.949}{order}%
\TextAt{8}{522.357}{307.949}{to}%
\TextAt{8}{79.226}{290.903}{check}%
\TextAt{8}{110.508}{290.903}{this,}%
\TextAt{8}{136.377}{290.903}{let}%
\TextAt{8}{153.025}{290.903}{us}%
\TextAt{8}{167.785}{290.903}{construct}%
\TextAt{8}{218.972}{290.903}{an}%
\TextAt{8}{234.975}{290.903}{element}%
\TextAt{13}{277.96}{290.903}{\Character{32}}%
\NativeMathAt{289.32}{290.903}{11.955}{\in}%
\TextAt{13}{300.611}{290.903}{G}%
\TextAt{25}{309.845}{289.11}{F}%
\TextAt{22}{314.716}{288.114}{p}%
\TextAt{8}{323.194}{290.903}{which}%
\TextAt{8}{356.424}{290.903}{does}%
\TextAt{8}{382.565}{290.903}{not}%
\TextAt{8}{403.122}{290.903}{belong}%
\TextAt{8}{440.262}{290.903}{to}%
\TextAt{16}{454.316}{290.903}{\Character{104}}%
\TextAt{13}{458.965}{290.903}{\Character{39}}%
\TextAt{16}{466.634}{290.903}{\Character{105}}%
\TextAt{8}{471.283}{290.903}{.}%
\TextAt{8}{479.646}{290.903}{We}%
\TextAt{8}{499.534}{290.903}{choose}%
\TextAt{8}{79.226}{273.857}{a}%
\TextAt{8}{88.982}{273.857}{sequence}%
\TextAt{16}{137.484}{273.857}{\Character{102}}%
\TextAt{13}{143.461}{273.857}{a}%
\TextAt{21}{149.606}{272.064}{n}%
\TextAt{16}{155.243}{273.857}{\Character{103}}%
\TextAt{21}{161.22}{272.064}{n}%
\NativeMathAt{166.358}{272.064}{7.97}{\in}%
\TextAt{25}{172.004}{272.064}{N}%
\TextAt{8}{182.156}{273.857}{of}%
\TextAt{8}{195.486}{273.857}{integers}%
\TextAt{8}{238.799}{273.857}{satisfying}%
\TextAt{13}{256.37}{248.891}{a}%
\TextAt{21}{262.515}{247.098}{n}%
\NativeMathAt{271.472}{248.891}{11.955}{\equiv}%
\TextAt{13}{284.091}{248.891}{a}%
\TextAt{21}{290.236}{247.098}{m}%
\TextAt{9}{310.18}{248.891}{\Character{40}}%
\TextAt{8}{314.732}{248.891}{mod}%
\TextAt{13}{340.834}{248.891}{m}%
\TextAt{9}{351.073}{248.891}{\Character{41}}%
\TextAt{8}{79.226}{223.925}{whenever}%
\TextAt{13}{131.607}{223.925}{m}%
\TextAt{16}{146.887}{223.925}{\Character{106}}%
\TextAt{13}{155.249}{223.925}{n}%
\TextAt{8}{162.236}{223.925}{,}%
\TextAt{8}{170.647}{223.925}{but}%
\TextAt{8}{193.116}{223.925}{such}%
\TextAt{8}{220.522}{223.925}{that}%
\TextAt{8}{246.901}{223.925}{there}%
\TextAt{8}{277.822}{223.925}{is}%
\TextAt{8}{290.595}{223.925}{no}%
\TextAt{8}{307.863}{223.925}{integer}%
\TextAt{13}{347.571}{223.925}{a}%
\TextAt{8}{358.623}{223.925}{satisfying}%
\TextAt{13}{411.793}{223.925}{a}%
\TextAt{21}{417.938}{222.131}{n}%
\NativeMathAt{428.615}{223.925}{11.955}{\equiv}%
\TextAt{13}{442.954}{223.925}{a}%
\TextAt{9}{454.413}{223.925}{\Character{40}}%
\TextAt{8}{458.965}{223.925}{mod}%
\TextAt{13}{485.067}{223.925}{n}%
\TextAt{9}{492.055}{223.925}{\Character{41}}%
\TextAt{8}{501.514}{223.925}{for}%
\TextAt{8}{520.408}{223.925}{all}%
\TextAt{13}{79.226}{206.879}{n}%
\NativeMathAt{89.535}{206.879}{11.955}{\in}%
\TextAt{18}{100.826}{206.879}{N}%
\TextAt{8}{109.46}{206.879}{.}%
\TextAt{8}{117.737}{206.879}{An}%
\TextAt{8}{136.398}{206.879}{example}%
\TextAt{8}{181.731}{206.879}{of}%
\TextAt{8}{194.545}{206.879}{such}%
\TextAt{8}{220.427}{206.879}{a}%
\TextAt{8}{229.668}{206.879}{sequence}%
\TextAt{8}{277.655}{206.879}{is}%
\TextAt{8}{288.904}{206.879}{given}%
\TextAt{8}{318.958}{206.879}{by}%
\TextAt{13}{334.702}{206.879}{a}%
\TextAt{21}{340.847}{205.085}{n}%
\TextAt{9}{349.804}{206.879}{\Character{61}}%
\TextAt{13}{362.23}{206.879}{n}%
\TextAt{3}{369.217}{211.217}{\Character{48}}%
\TextAt{13}{372.012}{206.879}{x}%
\TextAt{21}{378.664}{205.085}{n}%
\TextAt{8}{384.301}{206.879}{,}%
\TextAt{8}{391.038}{206.879}{where}%
\TextAt{8}{424.332}{206.879}{we}%
\TextAt{8}{441.044}{206.879}{write}%
\TextAt{13}{470.441}{206.879}{n}%
\TextAt{9}{480.749}{206.879}{\Character{61}}%
\TextAt{13}{493.175}{206.879}{n}%
\TextAt{3}{500.163}{211.217}{\Character{48}}%
\TextAt{13}{502.958}{206.879}{p}%
\TextAt{21}{508.833}{211.217}{v}%
\TextAt{22}{512.956}{210.221}{p}%
\TextAt{17}{517.295}{211.217}{\Character{40}}%
\TextAt{21}{520.588}{211.217}{n}%
\TextAt{17}{525.726}{211.217}{\Character{41}}%
\TextAt{8}{529.518}{206.879}{,}%
\TextAt{9}{79.226}{189.833}{\Character{40}}%
\TextAt{13}{83.778}{189.833}{n}%
\TextAt{3}{90.766}{194.171}{\Character{48}}%
\TextAt{13}{93.561}{189.833}{\Character{59}}%
\TextAt{13}{98.785}{189.833}{p}%
\TextAt{9}{104.68}{189.833}{\Character{41}}%
\TextAt{9}{112.535}{189.833}{\Character{61}}%
\TextAt{9}{124.956}{189.833}{1}%
\TextAt{8}{130.832}{189.833}{,}%
\TextAt{8}{137.981}{189.833}{and}%
\TextAt{9}{160.744}{189.833}{1}%
\TextAt{9}{169.902}{189.833}{\Character{61}}%
\TextAt{13}{182.343}{189.833}{n}%
\TextAt{3}{189.331}{194.171}{\Character{48}}%
\TextAt{13}{192.126}{189.833}{x}%
\TextAt{21}{198.778}{188.039}{n}%
\TextAt{9}{207.071}{189.833}{\Character{43}}%
\TextAt{13}{218.832}{189.833}{p}%
\TextAt{21}{224.707}{194.171}{v}%
\TextAt{22}{228.831}{193.175}{p}%
\TextAt{17}{233.17}{194.171}{\Character{40}}%
\TextAt{21}{236.463}{194.171}{n}%
\TextAt{17}{241.601}{194.171}{\Character{41}}%
\TextAt{13}{245.392}{189.833}{y}%
\TextAt{21}{251.1}{188.039}{n}%
\TextAt{8}{256.736}{189.833}{.}%
\TextAt{8}{265.185}{189.833}{Now}%
\TextAt{8}{291.845}{189.833}{put}%
\TextAt{13}{235.858}{164.866}{\Character{32}}%
\TextAt{21}{243.468}{163.073}{n}%
\TextAt{9}{252.425}{164.866}{\Character{61}}%
\TextAt{13}{264.851}{164.866}{\Character{39}}%
\TextAt{21}{272.519}{169.803}{a}%
\TextAt{22}{277.017}{168.806}{n}%
\TextAt{16}{282.614}{164.866}{\Character{106}}%
\TextAt{25}{285.935}{163.073}{F}%
\TextAt{22}{290.806}{161.529}{p}%
\TextAt{22}{294.646}{163.245}{n}%
\NativeMathAt{304.062}{164.866}{11.955}{\in}%
\TextAt{13}{315.353}{164.866}{G}%
\TextAt{9}{324.587}{164.866}{\Character{40}}%
\TextAt{18}{329.139}{164.866}{F}%
\TextAt{21}{336.445}{163.073}{p}%
\TextAt{22}{340.708}{165.342}{n}%
\TextAt{16}{349.625}{164.866}{\Character{106}}%
\TextAt{18}{356.267}{164.866}{F}%
\TextAt{21}{363.573}{163.073}{p}%
\TextAt{9}{368.334}{164.866}{\Character{41}}%
\TextAt{13}{372.886}{164.866}{\Character{58}}%
\TextAt{8}{79.226}{139.9}{If}%
\TextAt{18}{90.918}{139.9}{F}%
\TextAt{21}{98.224}{138.107}{p}%
\TextAt{22}{102.487}{140.376}{m}%
\NativeMathAt{113.342}{139.9}{11.955}{\subseteq}%
\TextAt{18}{125.961}{139.9}{F}%
\TextAt{21}{133.267}{138.107}{p}%
\TextAt{22}{137.53}{140.376}{n}%
\TextAt{8}{143.127}{139.9}{,}%
\TextAt{8}{150.276}{139.9}{then}%
\TextAt{13}{176.936}{139.9}{m}%
\TextAt{16}{190.497}{139.9}{\Character{106}}%
\TextAt{13}{197.138}{139.9}{n}%
\TextAt{8}{204.126}{139.9}{,}%
\TextAt{8}{211.275}{139.9}{so}%
\TextAt{8}{225.645}{139.9}{that}%
\TextAt{13}{251.014}{139.9}{a}%
\TextAt{21}{257.159}{138.107}{n}%
\NativeMathAt{266.116}{139.9}{11.955}{\equiv}%
\TextAt{13}{278.735}{139.9}{a}%
\TextAt{21}{284.88}{138.107}{m}%
\TextAt{9}{298.182}{139.9}{\Character{40}}%
\TextAt{8}{302.735}{139.9}{mod}%
\TextAt{13}{328.837}{139.9}{m}%
\TextAt{9}{339.076}{139.9}{\Character{41}}%
\TextAt{8}{343.628}{139.9}{,}%
\TextAt{8}{350.778}{139.9}{and}%
\TextAt{8}{373.54}{139.9}{therefore}%
\TextAt{13}{217.501}{114.934}{\Character{32}}%
\TextAt{21}{225.112}{113.14}{n}%
\TextAt{16}{230.748}{114.934}{\Character{106}}%
\TextAt{25}{234.069}{113.14}{F}%
\TextAt{22}{238.939}{111.596}{p}%
\TextAt{22}{242.78}{113.312}{m}%
\TextAt{9}{254.133}{114.934}{\Character{61}}%
\TextAt{13}{266.559}{114.934}{\Character{39}}%
\TextAt{21}{274.227}{119.87}{a}%
\TextAt{22}{278.725}{118.874}{n}%
\TextAt{16}{284.322}{114.934}{\Character{106}}%
\TextAt{25}{287.643}{113.14}{F}%
\TextAt{22}{292.513}{111.596}{p}%
\TextAt{22}{296.354}{113.312}{m}%
\TextAt{9}{307.707}{114.934}{\Character{61}}%
\TextAt{13}{320.133}{114.934}{\Character{39}}%
\TextAt{21}{327.801}{119.87}{a}%
\TextAt{22}{332.299}{118.874}{m}%
\TextAt{16}{339.833}{114.934}{\Character{106}}%
\TextAt{25}{343.154}{113.14}{F}%
\TextAt{22}{348.025}{111.596}{p}%
\TextAt{22}{351.865}{113.312}{m}%
\TextAt{9}{363.218}{114.934}{\Character{61}}%
\TextAt{13}{375.644}{114.934}{\Character{32}}%
\TextAt{21}{383.254}{113.14}{m}%
\TextAt{13}{391.243}{114.934}{\Character{58}}%
\TextAt{8}{79.226}{89.967}{Observe}%
\TextAt{8}{125.711}{89.967}{that}%
\TextAt{13}{152.626}{89.967}{\Character{39}}%
\TextAt{16}{160.294}{89.967}{\Character{106}}%
\TextAt{25}{163.615}{88.174}{F}%
\TextAt{22}{168.485}{86.63}{p}%
\TextAt{22}{172.326}{88.346}{m}%
\TextAt{8}{185.801}{89.967}{has}%
\TextAt{8}{208.22}{89.967}{order}%
\TextAt{13}{240.335}{89.967}{m}%
\TextAt{8}{250.574}{89.967}{.}%
\TextAt{8}{263.659}{89.967}{Therefore}%
\TextAt{8}{318.202}{89.967}{the}%
\TextAt{13}{339.904}{89.967}{\Character{32}}%
\TextAt{21}{347.514}{88.174}{n}%
\TextAt{8}{358.594}{89.967}{define}%
\TextAt{8}{393.948}{89.967}{an}%
\TextAt{8}{411.753}{89.967}{automorphism}%
\TextAt{8}{490.768}{89.967}{of}%
\Vector{q 0 G 0 g 0.478 w 433.643 -674.123 m 440.949 -674.123 l S Q}%
\TextAt{18}{505.643}{89.967}{F}%
\TextAt{21}{512.949}{88.174}{p}%
\TextAt{9}{523.665}{89.967}{\Character{61}}%
\TextAt{20}{79.226}{81.888}{\Character{83}}%
\NativeMathAt{89.189}{78.733}{7.97}{\infty}%
\TextAt{21}{89.189}{69.434}{n}%
\TextAt{17}{94.327}{69.434}{\Character{61}1}%
\TextAt{18}{107.638}{72.921}{F}%
\TextAt{21}{114.944}{71.128}{p}%
\TextAt{22}{119.207}{73.397}{n}%
\TextAt{8}{124.804}{72.921}{.}%
\TextAt{8}{133.719}{72.921}{Now}%
\TextAt{13}{160.535}{72.921}{\Character{32}}%
\TextAt{8}{172.627}{72.921}{cannot}%
\TextAt{8}{211.158}{72.921}{belong}%
\TextAt{8}{248.71}{72.921}{to}%
\TextAt{16}{263.175}{72.921}{\Character{104}}%
\TextAt{13}{267.825}{72.921}{\Character{39}}%
\TextAt{16}{275.493}{72.921}{\Character{105}}%
\TextAt{8}{284.195}{72.921}{because}%
\TextAt{13}{327.652}{72.921}{\Character{32}}%
\TextAt{9}{339.277}{72.921}{\Character{61}}%
\TextAt{13}{351.968}{72.921}{\Character{39}}%
\TextAt{21}{359.636}{77.26}{a}%
\TextAt{8}{364.632}{72.921}{,}%
\TextAt{8}{371.976}{72.921}{for}%
\TextAt{13}{390.016}{72.921}{a}%
\NativeMathAt{399.747}{72.921}{11.955}{\in}%
\TextAt{18}{411.303}{72.921}{Z}%
\TextAt{8}{419.273}{72.921}{,}%
\TextAt{8}{426.617}{72.921}{would}%
\TextAt{8}{460.916}{72.921}{imply}%
\TextAt{13}{493.913}{72.921}{\Character{32}}%
\TextAt{16}{501.952}{72.921}{\Character{106}}%
\TextAt{25}{505.273}{71.128}{F}%
\TextAt{22}{510.144}{69.584}{p}%
\TextAt{22}{513.984}{71.3}{n}%
\TextAt{9}{523.665}{72.921}{\Character{61}}%
\TextAt{8}{297.211}{43.033}{239}%
\Destination{79.23}{773.075}{nav.chapter.4}%
\Destination{78.23}{809.815}{nav.page.239}%
\Destination{79.23}{570.065}{nav.section.4.1}%
\Destination{79.36}{629.855}{nav.print.261}%
\Destination{79.23}{773.075}{outline.40}%
\Destination{79.23}{570.065}{outline.41}%
}

% Page 252 (body page 240)
\Page{%
\TextAt{2}{79.226}{796.943}{IV.}%
\TextAt{2}{98.012}{796.943}{Abstract}%
\TextAt{2}{143.252}{796.943}{Class}%
\TextAt{2}{171.844}{796.943}{Field}%
\TextAt{2}{199.575}{796.943}{Theory}%
\Vector{q 0 G 0 g 0.299 w 7.226 18.729 m 460.769 18.729 l S Q}%
\TextAt{13}{79.226}{761.124}{\Character{39}}%
\TextAt{21}{86.895}{765.462}{a}%
\TextAt{22}{91.393}{764.466}{n}%
\TextAt{16}{96.99}{761.124}{\Character{106}}%
\TextAt{25}{100.31}{759.33}{F}%
\TextAt{22}{105.181}{757.786}{p}%
\TextAt{22}{109.022}{759.502}{n}%
\TextAt{9}{120.439}{761.124}{\Character{61}}%
\TextAt{13}{134.866}{761.124}{\Character{39}}%
\TextAt{21}{142.534}{765.462}{a}%
\TextAt{16}{147.53}{761.124}{\Character{106}}%
\TextAt{25}{150.851}{759.33}{F}%
\TextAt{22}{155.722}{757.786}{p}%
\TextAt{22}{159.562}{759.502}{n}%
\TextAt{8}{170.729}{761.124}{and}%
\TextAt{8}{194.666}{761.124}{hence}%
\TextAt{13}{228.346}{761.124}{a}%
\TextAt{21}{234.491}{759.33}{n}%
\NativeMathAt{245.45}{761.124}{11.955}{\equiv}%
\TextAt{13}{260.07}{761.124}{a}%
\TextAt{9}{271.529}{761.124}{\Character{40}}%
\TextAt{8}{276.081}{761.124}{mod}%
\TextAt{13}{302.183}{761.124}{n}%
\TextAt{9}{309.171}{761.124}{\Character{41}}%
\TextAt{8}{318.795}{761.124}{for}%
\TextAt{8}{337.854}{761.124}{all}%
\TextAt{13}{355.287}{761.124}{n}%
\TextAt{8}{362.275}{761.124}{,}%
\TextAt{8}{370.892}{761.124}{which}%
\TextAt{8}{405.553}{761.124}{is}%
\TextAt{8}{418.491}{761.124}{what}%
\TextAt{8}{448.932}{761.124}{we}%
\TextAt{8}{467.333}{761.124}{ruled}%
\TextAt{8}{498.42}{761.124}{out}%
\TextAt{8}{520.408}{761.124}{by}%
\TextAt{8}{79.226}{744.078}{construction.}%
\TextAt{8}{97.159}{724.043}{The}%
\TextAt{8}{122.254}{724.043}{example}%
\TextAt{8}{169.143}{724.043}{does}%
\TextAt{8}{196.582}{724.043}{not}%
\TextAt{8}{218.437}{724.043}{mean,}%
\TextAt{8}{254.206}{724.043}{however,}%
\TextAt{8}{303.639}{724.043}{that}%
\TextAt{8}{330.049}{724.043}{we}%
\TextAt{8}{348.318}{724.043}{have}%
\TextAt{8}{376.354}{724.043}{to}%
\TextAt{8}{391.705}{724.043}{chuck}%
\TextAt{8}{425.265}{724.043}{the}%
\TextAt{8}{446.463}{724.043}{main}%
\TextAt{8}{476.77}{724.043}{theorem}%
\TextAt{8}{523.337}{724.043}{of}%
\TextAt{8}{79.226}{706.997}{Galois}%
\TextAt{8}{114.604}{706.997}{theory}%
\TextAt{8}{150.819}{706.997}{altogether}%
\TextAt{8}{205.577}{706.997}{in}%
\TextAt{8}{218.694}{706.997}{the}%
\TextAt{8}{238.316}{706.997}{case}%
\TextAt{8}{262.551}{706.997}{of}%
\TextAt{8}{275.346}{706.997}{infinite}%
\TextAt{8}{314.478}{706.997}{extensions.}%
\TextAt{8}{375.231}{706.997}{We}%
\TextAt{8}{394.841}{706.997}{just}%
\TextAt{8}{417.451}{706.997}{have}%
\TextAt{8}{443.91}{706.997}{to}%
\TextAt{8}{457.685}{706.997}{amend}%
\TextAt{8}{494.869}{706.997}{it}%
\TextAt{8}{506.038}{706.997}{using}%
\TextAt{8}{79.226}{689.951}{the}%
\TextAt{8}{99.324}{689.951}{observation}%
\TextAt{8}{161.456}{689.951}{that}%
\TextAt{8}{186.766}{689.951}{the}%
\TextAt{8}{206.863}{689.951}{Galois}%
\TextAt{8}{242.718}{689.951}{group}%
\TextAt{13}{275.835}{689.951}{G}%
\TextAt{9}{288.389}{689.951}{\Character{61}}%
\TextAt{13}{300.815}{689.951}{G}%
\TextAt{9}{310.048}{689.951}{\Character{40}\Character{10}}%
\TextAt{16}{326.376}{689.951}{\Character{106}}%
\TextAt{13}{333.018}{689.951}{k}%
\TextAt{9}{339.507}{689.951}{\Character{41}}%
\TextAt{8}{347.898}{689.951}{of}%
\TextAt{8}{361.169}{689.951}{any}%
\TextAt{8}{383.227}{689.951}{Galois}%
\TextAt{8}{419.082}{689.951}{extension}%
\TextAt{9}{470.789}{689.951}{\Character{10}}%
\TextAt{16}{482.564}{689.951}{\Character{106}}%
\TextAt{13}{489.206}{689.951}{k}%
\TextAt{8}{499.534}{689.951}{carries}%
\TextAt{8}{79.226}{672.905}{a}%
\TextAt{8}{90.13}{672.905}{canonical}%
\TextAt{8}{142.662}{672.905}{topology.}%
\TextAt{8}{197.821}{672.905}{This}%
\TextAt{8}{225.689}{672.905}{topology}%
\TextAt{8}{274.981}{672.905}{is}%
\TextAt{8}{287.893}{672.905}{called}%
\TextAt{8}{322.205}{672.905}{the}%
\TextAt{12}{343.51}{672.905}{Krull}%
\TextAt{12}{380.302}{672.905}{topology}%
\TextAt{8}{436.193}{672.905}{and}%
\TextAt{8}{460.105}{672.905}{is}%
\TextAt{8}{473.017}{672.905}{obtained}%
\TextAt{8}{522.297}{672.905}{as}%
\TextAt{8}{79.226}{655.859}{follows.}%
\TextAt{8}{122.213}{655.859}{For}%
\TextAt{8}{143.182}{655.859}{every}%
\TextAt{13}{174.074}{655.859}{\Character{27}}%
\NativeMathAt{184.478}{655.859}{11.955}{\in}%
\TextAt{13}{195.769}{655.859}{G}%
\TextAt{8}{208.9}{655.859}{we}%
\TextAt{8}{226.127}{655.859}{take}%
\TextAt{8}{251.496}{655.859}{the}%
\TextAt{8}{271.652}{655.859}{cosets}%
\TextAt{13}{278.67}{628.343}{\Character{27}}%
\TextAt{13}{285.736}{628.343}{G}%
\TextAt{9}{294.986}{628.343}{\Character{40}\Character{10}}%
\TextAt{16}{311.314}{628.343}{\Character{106}}%
\TextAt{13}{317.955}{628.343}{K}%
\TextAt{9}{328.773}{628.343}{\Character{41}}%
\TextAt{8}{79.226}{600.826}{as}%
\TextAt{8}{93.62}{600.826}{a}%
\TextAt{8}{103.4}{600.826}{basis}%
\TextAt{8}{132.164}{600.826}{of}%
\TextAt{8}{145.518}{600.826}{neighbourhoods}%
\TextAt{8}{230.161}{600.826}{of}%
\TextAt{13}{243.515}{600.826}{\Character{27}}%
\TextAt{8}{250.598}{600.826}{,}%
\TextAt{8}{257.777}{600.826}{with}%
\TextAt{13}{284.461}{600.826}{K}%
\TextAt{16}{298.641}{600.826}{\Character{106}}%
\TextAt{13}{305.323}{600.826}{k}%
\TextAt{8}{315.734}{600.826}{ranging}%
\TextAt{8}{358.044}{600.826}{over}%
\TextAt{8}{383.114}{600.826}{finite}%
\TextAt{8}{413.05}{600.826}{Galois}%
\TextAt{8}{448.987}{600.826}{subextensions}%
\TextAt{8}{523.337}{600.826}{of}%
\TextAt{9}{79.226}{583.78}{\Character{10}}%
\TextAt{16}{91.001}{583.78}{\Character{106}}%
\TextAt{13}{97.643}{583.78}{k}%
\TextAt{8}{104.133}{583.78}{.}%
\TextAt{8}{112.581}{583.78}{The}%
\TextAt{8}{136.635}{583.78}{multiplication}%
\TextAt{8}{211.761}{583.78}{and}%
\TextAt{8}{234.524}{583.78}{the}%
\TextAt{8}{254.68}{583.78}{inverse}%
\TextAt{8}{293.439}{583.78}{map}%
\TextAt{13}{167.029}{556.264}{G}%
\TextAt{16}{178.92}{556.264}{\Character{2}}%
\TextAt{13}{190.875}{556.264}{G}%
\NativeArrowAt{203.429}{556.264}{11.955}{19.26}{\longrightarrow}%
\TextAt{13}{226.011}{556.264}{G\Character{59}}%
\TextAt{9}{252.444}{556.264}{\Character{40}}%
\TextAt{13}{256.997}{556.264}{\Character{27}}%
\TextAt{13}{263.429}{556.264}{\Character{59}}%
\TextAt{13}{268.653}{556.264}{\Character{28}}%
\TextAt{9}{275.092}{556.264}{\Character{41}}%
\NativeArrowAt{282.965}{556.264}{11.955}{19.26}{\longmapsto}%
\TextAt{13}{305.547}{556.264}{\Character{27}}%
\TextAt{13}{312.613}{556.264}{\Character{28}}%
\TextAt{13}{318.387}{556.264}{\Character{59}}%
\TextAt{8}{347.552}{556.264}{and}%
\TextAt{13}{390.328}{556.264}{\Character{27}}%
\NativeArrowAt{400.731}{556.264}{11.955}{19.26}{\longmapsto}%
\TextAt{13}{423.313}{556.264}{\Character{27}}%
\TextAt{3}{430.396}{561.2}{\Character{0}}%
\TextAt{17}{436.982}{561.2}{1}%
\TextAt{13}{441.715}{556.264}{\Character{59}}%
\TextAt{8}{79.226}{528.748}{are}%
\TextAt{8}{98.833}{528.748}{continuous}%
\TextAt{8}{157.533}{528.748}{maps,}%
\TextAt{8}{191.534}{528.748}{since}%
\TextAt{8}{220.298}{528.748}{the}%
\TextAt{8}{240.55}{528.748}{preimage}%
\TextAt{8}{290.726}{528.748}{of}%
\TextAt{8}{304.152}{528.748}{a}%
\TextAt{8}{314.003}{528.748}{fundamental}%
\TextAt{8}{381.741}{528.748}{open}%
\TextAt{8}{410.123}{528.748}{neighbourhood}%
\TextAt{13}{490.223}{528.748}{\Character{27}}%
\TextAt{13}{497.288}{528.748}{\Character{28}}%
\TextAt{13}{503.709}{528.748}{G}%
\TextAt{9}{512.958}{528.748}{\Character{40}\Character{10}}%
\TextAt{16}{529.449}{528.748}{\Character{106}}%
\TextAt{13}{79.226}{511.702}{K}%
\TextAt{9}{90.044}{511.702}{\Character{41}}%
\TextAt{8}{94.596}{511.702}{,}%
\TextAt{8}{101.722}{511.702}{resp.}%
\TextAt{13}{131.034}{511.702}{\Character{27}}%
\TextAt{3}{138.116}{516.04}{\Character{0}}%
\TextAt{17}{144.703}{516.04}{1}%
\TextAt{13}{149.435}{511.702}{G}%
\TextAt{9}{158.669}{511.702}{\Character{40}\Character{10}}%
\TextAt{16}{174.996}{511.702}{\Character{106}}%
\TextAt{13}{181.638}{511.702}{K}%
\TextAt{9}{192.456}{511.702}{\Character{41}}%
\TextAt{8}{197.008}{511.702}{,}%
\TextAt{8}{204.133}{511.702}{contains}%
\TextAt{8}{250.024}{511.702}{the}%
\TextAt{8}{270.15}{511.702}{open}%
\TextAt{8}{298.407}{511.702}{neighbourhood}%
\TextAt{13}{378.381}{511.702}{\Character{27}}%
\TextAt{13}{385.447}{511.702}{G}%
\TextAt{9}{394.697}{511.702}{\Character{40}\Character{10}}%
\TextAt{16}{411.024}{511.702}{\Character{106}}%
\TextAt{13}{417.666}{511.702}{K}%
\TextAt{9}{428.484}{511.702}{\Character{41}}%
\TextAt{16}{435.632}{511.702}{\Character{2}}%
\TextAt{13}{447.527}{511.702}{\Character{28}}%
\TextAt{13}{453.935}{511.702}{G}%
\TextAt{9}{463.179}{511.702}{\Character{40}\Character{10}}%
\TextAt{16}{479.507}{511.702}{\Character{106}}%
\TextAt{13}{486.148}{511.702}{K}%
\TextAt{9}{496.966}{511.702}{\Character{41}}%
\TextAt{8}{501.518}{511.702}{,}%
\TextAt{8}{508.644}{511.702}{resp.}%
\TextAt{13}{79.226}{494.656}{\Character{27}}%
\TextAt{13}{86.292}{494.656}{G}%
\TextAt{9}{95.542}{494.656}{\Character{40}\Character{10}}%
\TextAt{16}{111.87}{494.656}{\Character{106}}%
\TextAt{13}{118.511}{494.656}{K}%
\TextAt{9}{129.329}{494.656}{\Character{41}}%
\TextAt{8}{133.881}{494.656}{.}%
\TextAt{8}{142.33}{494.656}{Thus}%
\TextAt{13}{171.979}{494.656}{G}%
\TextAt{8}{185.11}{494.656}{is}%
\TextAt{8}{196.873}{494.656}{a}%
\TextAt{8}{206.629}{494.656}{topological}%
\TextAt{8}{266.154}{494.656}{group}%
\TextAt{8}{299.329}{494.656}{which}%
\TextAt{8}{332.816}{494.656}{satisfies}%
\TextAt{8}{375.926}{494.656}{the}%
\TextAt{8}{396.082}{494.656}{following}%
\TextAt{12}{79.226}{469.068}{Proposition}%
\TextAt{12}{151.826}{469.068}{(1.1).}%
\TextAt{8}{187.045}{469.068}{For}%
\TextAt{8}{207.918}{469.068}{every}%
\TextAt{8}{238.713}{469.068}{(finite}%
\TextAt{8}{273.083}{469.068}{or}%
\TextAt{8}{287.297}{469.068}{infinite)}%
\TextAt{8}{331.422}{469.068}{Galois}%
\TextAt{8}{367.239}{469.068}{extension}%
\TextAt{9}{418.908}{469.068}{\Character{10}}%
\TextAt{16}{430.683}{469.068}{\Character{106}}%
\TextAt{13}{437.325}{469.068}{k}%
\TextAt{8}{447.615}{469.068}{the}%
\TextAt{8}{467.675}{469.068}{Galois}%
\TextAt{8}{503.491}{469.068}{group}%
\TextAt{13}{79.226}{452.022}{G}%
\TextAt{9}{91.781}{452.022}{\Character{61}}%
\TextAt{13}{104.206}{452.022}{G}%
\TextAt{9}{113.44}{452.022}{\Character{40}\Character{10}}%
\TextAt{16}{129.767}{452.022}{\Character{106}}%
\TextAt{13}{136.409}{452.022}{k}%
\TextAt{9}{142.898}{452.022}{\Character{41}}%
\TextAt{8}{151.348}{452.022}{is}%
\TextAt{8}{163.112}{452.022}{compact}%
\TextAt{8}{209.94}{452.022}{Hausdorff}%
\TextAt{8}{263.332}{452.022}{with}%
\TextAt{8}{289.992}{452.022}{respect}%
\TextAt{8}{330.042}{452.022}{to}%
\TextAt{8}{344.352}{452.022}{the}%
\TextAt{8}{364.509}{452.022}{Krull}%
\TextAt{8}{395.066}{452.022}{topology.}%
\TextAt{19}{79.226}{417.891}{Proof.}%
\TextAt{8}{116}{417.891}{If}%
\TextAt{13}{127.419}{417.891}{\Character{27}}%
\TextAt{13}{133.851}{417.891}{\Character{59}}%
\TextAt{13}{139.075}{417.891}{\Character{28}}%
\NativeMathAt{148.835}{417.891}{11.955}{\in}%
\TextAt{13}{160.126}{417.891}{G}%
\TextAt{8}{172.983}{417.891}{and}%
\TextAt{13}{195.472}{417.891}{\Character{27}}%
\NativeMathAt{205.876}{417.891}{11.955}{\neq}%
\TextAt{13}{218.301}{417.891}{\Character{28}}%
\TextAt{8}{224.72}{417.891}{,}%
\TextAt{8}{231.65}{417.891}{then}%
\TextAt{8}{258.037}{417.891}{there}%
\TextAt{8}{287.675}{417.891}{exists}%
\TextAt{8}{319.716}{417.891}{a}%
\TextAt{8}{329.198}{417.891}{finite}%
\TextAt{8}{358.836}{417.891}{Galois}%
\TextAt{8}{394.476}{417.891}{subextension}%
\TextAt{13}{463.913}{417.891}{K}%
\TextAt{16}{478.052}{417.891}{\Character{106}}%
\TextAt{13}{484.693}{417.891}{k}%
\TextAt{8}{494.807}{417.891}{of}%
\TextAt{9}{507.863}{417.891}{\Character{10}}%
\TextAt{16}{519.638}{417.891}{\Character{106}}%
\TextAt{13}{526.28}{417.891}{k}%
\TextAt{8}{79.226}{400.845}{such}%
\TextAt{8}{106.418}{400.845}{that}%
\TextAt{13}{132.582}{400.845}{\Character{27}}%
\TextAt{16}{139.665}{400.845}{\Character{106}}%
\TextAt{21}{142.986}{399.052}{K}%
\NativeMathAt{155.864}{400.845}{11.955}{\neq}%
\TextAt{13}{169.645}{400.845}{\Character{28}}%
\TextAt{16}{176.064}{400.845}{\Character{106}}%
\TextAt{21}{179.385}{399.052}{K}%
\TextAt{8}{187.587}{400.845}{,}%
\TextAt{8}{195.73}{400.845}{so}%
\TextAt{8}{210.896}{400.845}{that}%
\TextAt{13}{237.06}{400.845}{\Character{27}}%
\TextAt{13}{244.126}{400.845}{G}%
\TextAt{9}{253.376}{400.845}{\Character{40}\Character{10}}%
\TextAt{16}{271.058}{400.845}{\Character{106}}%
\TextAt{13}{279.055}{400.845}{K}%
\TextAt{9}{289.873}{400.845}{\Character{41}}%
\NativeMathAt{299.101}{400.845}{11.955}{\neq}%
\TextAt{13}{312.882}{400.845}{\Character{28}}%
\TextAt{13}{319.29}{400.845}{G}%
\TextAt{9}{328.534}{400.845}{\Character{40}\Character{10}}%
\TextAt{16}{346.217}{400.845}{\Character{106}}%
\TextAt{13}{354.214}{400.845}{K}%
\TextAt{9}{365.031}{400.845}{\Character{41}}%
\TextAt{8}{374.276}{400.845}{and}%
\TextAt{8}{397.834}{400.845}{thus}%
\TextAt{13}{424.381}{400.845}{\Character{27}}%
\TextAt{13}{431.447}{400.845}{G}%
\TextAt{9}{440.697}{400.845}{\Character{40}\Character{10}}%
\TextAt{16}{458.379}{400.845}{\Character{106}}%
\TextAt{13}{466.376}{400.845}{K}%
\TextAt{9}{477.194}{400.845}{\Character{41}}%
\TextAt{16}{484.945}{400.845}{\Character{92}}%
\TextAt{13}{496.113}{400.845}{\Character{28}}%
\TextAt{13}{502.521}{400.845}{G}%
\TextAt{9}{511.766}{400.845}{\Character{40}\Character{10}}%
\TextAt{16}{529.449}{400.845}{\Character{106}}%
\TextAt{13}{79.226}{383.799}{K}%
\TextAt{9}{90.044}{383.799}{\Character{41}}%
\TextAt{9}{98.915}{383.799}{\Character{61}}%
\TextAt{18}{112.36}{383.799}{\Character{63}}%
\TextAt{8}{121.658}{383.799}{.}%
\TextAt{8}{131.882}{383.799}{This}%
\TextAt{8}{159.194}{383.799}{shows}%
\TextAt{8}{193.404}{383.799}{that}%
\TextAt{13}{219.365}{383.799}{G}%
\TextAt{8}{233.088}{383.799}{is}%
\TextAt{8}{245.443}{383.799}{Hausdorff.}%
\TextAt{8}{305.162}{383.799}{In}%
\TextAt{8}{320.375}{383.799}{order}%
\TextAt{8}{351.536}{383.799}{to}%
\TextAt{8}{366.438}{383.799}{prove}%
\TextAt{8}{398.58}{383.799}{compactness,}%
\TextAt{8}{470.333}{383.799}{consider}%
\TextAt{8}{516.51}{383.799}{the}%
\TextAt{8}{79.226}{366.754}{mapping}%
\TextAt{13}{198.023}{349.708}{h}%
\TextAt{9}{208.082}{349.708}{\Character{58}}%
\TextAt{13}{214.655}{349.708}{G}%
\NativeArrowAt{227.209}{349.708}{11.955}{19.26}{\longrightarrow}%
\TextAt{20}{249.791}{361.065}{\Character{89}}%
\TextAt{21}{253.577}{335.693}{K}%
\TextAt{13}{267.06}{349.708}{G}%
\TextAt{9}{276.294}{349.708}{\Character{40}}%
\TextAt{13}{280.846}{349.708}{K}%
\TextAt{16}{294.984}{349.708}{\Character{106}}%
\TextAt{13}{301.626}{349.708}{k}%
\TextAt{9}{308.116}{349.708}{\Character{41}}%
\TextAt{13}{312.668}{349.708}{\Character{59}}%
\TextAt{13}{341.803}{349.708}{\Character{27}}%
\NativeArrowAt{352.226}{349.708}{11.955}{19.26}{\longmapsto}%
\TextAt{9}{374.808}{349.708}{\Character{40}}%
\TextAt{13}{379.36}{349.708}{\Character{27}}%
\TextAt{16}{386.442}{349.708}{\Character{106}}%
\TextAt{21}{389.763}{347.914}{K}%
\TextAt{9}{397.966}{349.708}{\Character{41}}%
\TextAt{21}{402.518}{347.914}{K}%
\TextAt{13}{410.721}{349.708}{\Character{59}}%
\TextAt{8}{79.226}{318.696}{where}%
\TextAt{13}{111.832}{318.696}{K}%
\TextAt{16}{125.97}{318.696}{\Character{106}}%
\TextAt{13}{132.612}{318.696}{k}%
\TextAt{8}{141.795}{318.696}{varies}%
\TextAt{8}{173.505}{318.696}{over}%
\TextAt{8}{197.347}{318.696}{the}%
\TextAt{8}{216.3}{318.696}{finite}%
\TextAt{8}{245.008}{318.696}{Galois}%
\TextAt{8}{279.718}{318.696}{subextensions.}%
\TextAt{8}{358.193}{318.696}{We}%
\TextAt{8}{377.134}{318.696}{view}%
\TextAt{8}{402.913}{318.696}{the}%
\TextAt{8}{421.866}{318.696}{finite}%
\TextAt{8}{450.575}{318.696}{groups}%
\TextAt{13}{487.161}{318.696}{G}%
\TextAt{9}{496.395}{318.696}{\Character{40}}%
\TextAt{13}{500.947}{318.696}{K}%
\TextAt{16}{515.086}{318.696}{\Character{106}}%
\TextAt{13}{521.728}{318.696}{k}%
\TextAt{9}{528.217}{318.696}{\Character{41}}%
\TextAt{8}{79.226}{301.65}{as}%
\TextAt{8}{93.467}{301.65}{discrete}%
\TextAt{8}{136.317}{301.65}{compact}%
\TextAt{8}{183.016}{301.65}{topological}%
\TextAt{8}{242.412}{301.65}{groups.}%
\TextAt{8}{284.71}{301.65}{Their}%
\TextAt{8}{316.441}{301.65}{product}%
\TextAt{8}{360.211}{301.65}{is}%
\TextAt{8}{371.846}{301.65}{therefore}%
\TextAt{8}{420.817}{301.65}{a}%
\TextAt{8}{430.443}{301.65}{compact}%
\TextAt{8}{477.142}{301.65}{topological}%
\TextAt{8}{79.226}{284.604}{space,}%
\TextAt{8}{115.737}{284.604}{by}%
\TextAt{8}{133.584}{284.604}{Tykhonov’s}%
\TextAt{8}{198.008}{284.604}{theorem}%
\TextAt{8}{245.12}{284.604}{(see}%
\TextAt{8}{270.176}{284.604}{[98]).}%
\TextAt{8}{306.162}{284.604}{The}%
\TextAt{8}{331.803}{284.604}{homomorphism}%
\TextAt{13}{416.06}{284.604}{h}%
\TextAt{8}{428.284}{284.604}{is}%
\TextAt{8}{441.635}{284.604}{injective,}%
\TextAt{8}{493.365}{284.604}{because}%
\TextAt{13}{79.226}{267.558}{\Character{27}}%
\TextAt{16}{86.309}{267.558}{\Character{106}}%
\TextAt{21}{89.629}{265.765}{K}%
\TextAt{9}{101.153}{267.558}{\Character{61}}%
\TextAt{9}{113.563}{267.558}{1}%
\TextAt{8}{123.329}{267.558}{for}%
\TextAt{8}{141.214}{267.558}{all}%
\TextAt{13}{157.473}{267.558}{K}%
\TextAt{8}{172.188}{267.558}{is}%
\TextAt{8}{183.952}{267.558}{equivalent}%
\TextAt{8}{239.567}{267.558}{to}%
\TextAt{13}{253.877}{267.558}{\Character{27}}%
\TextAt{9}{264.281}{267.558}{\Character{61}}%
\TextAt{9}{276.69}{267.558}{1}%
\TextAt{8}{282.559}{267.558}{.}%
\TextAt{8}{291.008}{267.558}{The}%
\TextAt{8}{315.061}{267.558}{sets}%
\TextAt{13}{240.108}{240.042}{U}%
\TextAt{9}{252.628}{240.042}{\Character{61}}%
\TextAt{20}{270.207}{251.399}{\Character{89}}%
\TextAt{21}{265.054}{225.939}{K}%
\NativeMathAt{272.758}{225.939}{7.97}{\neq}%
\TextAt{21}{279.345}{225.939}{K}%
\TextAt{24}{286.486}{224.832}{0}%
\TextAt{13}{292.63}{240.042}{G}%
\TextAt{9}{301.863}{240.042}{\Character{40}}%
\TextAt{13}{306.416}{240.042}{K}%
\TextAt{16}{320.554}{240.042}{\Character{106}}%
\TextAt{13}{327.196}{240.042}{k}%
\TextAt{9}{333.685}{240.042}{\Character{41}}%
\TextAt{16}{340.894}{240.042}{\Character{2}}%
\TextAt{16}{352.837}{240.042}{\Character{102}}%
\TextAt{9}{359.442}{240.042}{\Character{22}}%
\TextAt{13}{358.827}{240.042}{\Character{27}}%
\TextAt{16}{365.91}{240.042}{\Character{103}}%
\TextAt{8}{79.226}{202.76}{form}%
\TextAt{8}{107.369}{202.76}{a}%
\TextAt{8}{117.627}{202.76}{subbasis}%
\TextAt{8}{164.492}{202.76}{of}%
\TextAt{8}{178.325}{202.76}{open}%
\TextAt{8}{207.114}{202.76}{sets}%
\TextAt{8}{230.498}{202.76}{of}%
\TextAt{8}{244.331}{202.76}{the}%
\TextAt{8}{264.99}{202.76}{product}%
\TextAt{20}{309.392}{211.727}{\Character{81}}%
\TextAt{21}{320.683}{199.273}{K}%
\TextAt{13}{330.879}{202.76}{G}%
\TextAt{9}{340.112}{202.76}{\Character{40}}%
\TextAt{13}{344.665}{202.76}{K}%
\TextAt{16}{359.66}{202.76}{\Character{106}}%
\TextAt{13}{367.158}{202.76}{k}%
\TextAt{9}{373.648}{202.76}{\Character{41}}%
\TextAt{8}{378.2}{202.76}{,}%
\TextAt{8}{385.977}{202.76}{where}%
\TextAt{13}{420.289}{202.76}{K}%
\TextAt{17}{430.249}{200.967}{0}%
\TextAt{16}{439.158}{202.76}{\Character{106}}%
\TextAt{13}{446.657}{202.76}{k}%
\TextAt{8}{457.546}{202.76}{varies}%
\TextAt{8}{490.962}{202.76}{over}%
\TextAt{8}{516.51}{202.76}{the}%
\ErrataAt{79.19}{185.715}{455.81}{finite Galois subextensions of $\Omega\mid k$ and $\bar\sigma\in G(K_0\mid k)$. If $\sigma\in G$ is a preimage of $\bar\sigma$, then}%
\TextAt{13}{79.226}{168.669}{h}%
\TextAt{3}{85.965}{173.007}{\Character{0}}%
\TextAt{17}{92.551}{173.007}{1}%
\TextAt{9}{97.284}{168.669}{\Character{40}}%
\TextAt{13}{101.836}{168.669}{U}%
\TextAt{9}{111.035}{168.669}{\Character{41}}%
\TextAt{9}{119.715}{168.669}{\Character{61}}%
\TextAt{13}{132.966}{168.669}{\Character{27}}%
\TextAt{13}{140.032}{168.669}{G}%
\TextAt{9}{149.282}{168.669}{\Character{40}\Character{10}}%
\TextAt{16}{166.426}{168.669}{\Character{106}}%
\TextAt{13}{173.884}{168.669}{K}%
\TextAt{17}{183.843}{166.875}{0}%
\TextAt{9}{188.575}{168.669}{\Character{41}}%
\TextAt{8}{193.128}{168.669}{.}%
\TextAt{8}{203.013}{168.669}{Thus}%
\TextAt{13}{233.141}{168.669}{h}%
\TextAt{8}{244.256}{168.669}{is}%
\TextAt{8}{256.498}{168.669}{continuous.}%
\TextAt{8}{321.091}{168.669}{Moreover}%
\TextAt{13}{372.953}{168.669}{h}%
\TextAt{9}{379.691}{168.669}{\Character{40}}%
\TextAt{13}{384.244}{168.669}{\Character{27}}%
\TextAt{13}{391.31}{168.669}{G}%
\TextAt{9}{400.56}{168.669}{\Character{40}\Character{10}}%
\TextAt{16}{417.704}{168.669}{\Character{106}}%
\TextAt{13}{425.162}{168.669}{K}%
\TextAt{17}{435.121}{166.875}{0}%
\TextAt{9}{439.853}{168.669}{\Character{41}\Character{41}}%
\TextAt{9}{453.076}{168.669}{\Character{61}}%
\TextAt{13}{466.337}{168.669}{h}%
\TextAt{9}{473.075}{168.669}{\Character{40}}%
\TextAt{13}{477.627}{168.669}{G}%
\TextAt{9}{486.861}{168.669}{\Character{41}}%
\TextAt{16}{494.397}{168.669}{\Character{92}}%
\TextAt{13}{505.35}{168.669}{U}%
\TextAt{8}{514.549}{168.669}{,}%
\TextAt{8}{522.297}{168.669}{so}%
\TextAt{13}{79.226}{151.623}{h}%
\TextAt{9}{90.294}{151.623}{\Character{58}}%
\TextAt{13}{97.875}{151.623}{G}%
\NativeArrowAt{111.437}{151.623}{11.955}{19.26}{\longrightarrow}%
\TextAt{13}{135.027}{151.623}{h}%
\TextAt{9}{141.766}{151.623}{\Character{40}}%
\TextAt{13}{146.318}{151.623}{G}%
\TextAt{9}{155.552}{151.623}{\Character{41}}%
\TextAt{8}{164.593}{151.623}{is}%
\TextAt{8}{176.949}{151.623}{open,}%
\TextAt{8}{209.226}{151.623}{and}%
\TextAt{8}{232.58}{151.623}{thus}%
\TextAt{8}{258.924}{151.623}{a}%
\TextAt{8}{269.271}{151.623}{homeomorphism.}%
\TextAt{8}{363.467}{151.623}{It}%
\TextAt{8}{376.731}{151.623}{therefore}%
\TextAt{8}{426.423}{151.623}{suffices}%
\TextAt{8}{466.801}{151.623}{to}%
\TextAt{8}{481.703}{151.623}{show}%
\TextAt{8}{511.298}{151.623}{that}%
\TextAt{13}{79.226}{134.577}{h}%
\TextAt{9}{85.965}{134.577}{\Character{40}}%
\TextAt{13}{90.517}{134.577}{G}%
\TextAt{9}{99.751}{134.577}{\Character{41}}%
\TextAt{8}{108.898}{134.577}{is}%
\TextAt{8}{121.36}{134.577}{closed}%
\TextAt{8}{156.584}{134.577}{in}%
\TextAt{8}{170.934}{134.577}{the}%
\TextAt{8}{191.789}{134.577}{compact}%
\TextAt{8}{239.315}{134.577}{set}%
\TextAt{20}{258.28}{143.543}{\Character{81}}%
\TextAt{21}{269.571}{131.09}{K}%
\TextAt{13}{279.766}{134.577}{G}%
\TextAt{9}{289}{134.577}{\Character{40}}%
\TextAt{13}{293.552}{134.577}{K}%
\TextAt{16}{308.879}{134.577}{\Character{106}}%
\TextAt{13}{316.71}{134.577}{k}%
\TextAt{9}{323.2}{134.577}{\Character{41}}%
\TextAt{8}{327.752}{134.577}{.}%
\TextAt{8}{338.293}{134.577}{To}%
\TextAt{8}{356.219}{134.577}{see}%
\TextAt{8}{375.829}{134.577}{this}%
\TextAt{8}{399.349}{134.577}{we}%
\TextAt{8}{417.275}{134.577}{consider,}%
\TextAt{8}{466.984}{134.577}{for}%
\TextAt{8}{485.566}{134.577}{each}%
\TextAt{8}{512.601}{134.577}{pair}%
\TextAt{13}{79.226}{117.531}{L}%
\TextAt{3}{87.191}{121.869}{\Character{48}}%
\NativeMathAt{93.307}{117.531}{11.955}{\supseteq}%
\TextAt{13}{105.926}{117.531}{L}%
\TextAt{8}{117.788}{117.531}{of}%
\TextAt{8}{131.118}{117.531}{finite}%
\TextAt{8}{161.03}{117.531}{Galois}%
\TextAt{8}{196.943}{117.531}{subextensions}%
\TextAt{8}{271.268}{117.531}{of}%
\TextAt{9}{284.598}{117.531}{\Character{10}}%
\TextAt{16}{296.373}{117.531}{\Character{106}}%
\TextAt{13}{303.015}{117.531}{k}%
\TextAt{8}{309.505}{117.531}{,}%
\TextAt{8}{316.654}{117.531}{the}%
\TextAt{8}{336.81}{117.531}{set}%
\TextAt{13}{184.629}{82.153}{M}%
\TextAt{21}{195.954}{80.294}{L}%
\TextAt{23}{201.713}{82.563}{\Character{48}}%
\TextAt{3}{204.397}{80.294}{\Character{106}}%
\TextAt{21}{206.749}{80.294}{L}%
\TextAt{9}{216.327}{82.153}{\Character{61}}%
\TextAt{20}{228.753}{102.597}{\Character{40}}%
\TextAt{9}{238.383}{82.153}{\Character{40}}%
\TextAt{13}{242.936}{82.153}{\Character{27}}%
\TextAt{21}{249.589}{80.36}{K}%
\TextAt{9}{257.792}{82.153}{\Character{41}}%
\TextAt{21}{262.344}{80.36}{K}%
\NativeMathAt{273.868}{82.153}{11.955}{\in}%
\TextAt{20}{285.159}{93.511}{\Character{89}}%
\TextAt{21}{288.945}{68.139}{K}%
\TextAt{13}{302.427}{82.153}{G}%
\TextAt{9}{311.661}{82.153}{\Character{40}}%
\TextAt{13}{316.213}{82.153}{K}%
\TextAt{16}{330.352}{82.153}{\Character{106}}%
\TextAt{13}{336.993}{82.153}{k}%
\TextAt{9}{343.483}{82.153}{\Character{41}}%
\TextAt{20}{351.356}{103.075}{\Character{12}}%
\TextAt{20}{351.356}{95.902}{\Character{12}}%
\TextAt{20}{351.356}{88.729}{\Character{12}}%
\TextAt{20}{351.356}{81.556}{\Character{12}}%
\TextAt{20}{351.356}{74.383}{\Character{12}}%
\TextAt{13}{358.662}{82.154}{\Character{27}}%
\TextAt{21}{365.315}{80.36}{L}%
\TextAt{23}{371.074}{82.63}{\Character{48}}%
\TextAt{16}{374.257}{82.154}{\Character{106}}%
\TextAt{21}{377.578}{80.36}{L}%
\TextAt{9}{387.156}{82.153}{\Character{61}}%
\TextAt{13}{399.581}{82.153}{\Character{27}}%
\TextAt{21}{406.235}{80.36}{L}%
\TextAt{20}{412.492}{102.597}{\Character{41}}%
\TextAt{13}{424.115}{82.153}{\Character{58}}%
\TextAt{8}{297.211}{43.033}{240}%
\Destination{78.23}{809.815}{nav.page.240}%
\Destination{79.36}{739.055}{nav.print.262}%
\Destination{79.36}{483.055}{nav.number.4.1.1}%
\Hotlink{269.676}{282.355}{19.22}{11.047}{nav.cite.N098}%
\Destination{79.226}{202.715}{errata.113}%
}

% Page 253 (body page 241)
\Page{%
\TextAt{2}{79.226}{796.943}{IV.}%
\TextAt{2}{98.012}{796.943}{Abstract}%
\TextAt{2}{143.252}{796.943}{Class}%
\TextAt{2}{171.844}{796.943}{Field}%
\TextAt{2}{199.575}{796.943}{Theory}%
\Vector{q 0 G 0 g 0.299 w 7.226 18.729 m 460.769 18.729 l S Q}%
\TextAt{8}{79.226}{761.124}{One}%
\TextAt{8}{104.407}{761.124}{clearly}%
\TextAt{8}{142.273}{761.124}{has}%
\TextAt{13}{163.616}{761.124}{h}%
\TextAt{9}{170.355}{761.124}{\Character{40}}%
\TextAt{13}{174.907}{761.124}{G}%
\TextAt{9}{184.141}{761.124}{\Character{41}}%
\TextAt{9}{192.797}{761.124}{\Character{61}}%
\TextAt{20}{206.04}{770.09}{\Character{84}}%
\TextAt{21}{216.003}{757.637}{L}%
\TextAt{23}{221.762}{759.906}{\Character{48}}%
\NativeMathAt{224.446}{757.637}{7.97}{\supseteq}%
\TextAt{21}{231.033}{757.637}{L}%
\TextAt{13}{239.282}{761.124}{M}%
\TextAt{21}{250.607}{759.264}{L}%
\TextAt{23}{256.366}{761.533}{\Character{48}}%
\TextAt{3}{259.05}{759.264}{\Character{106}}%
\TextAt{21}{261.403}{759.264}{L}%
\TextAt{8}{267.66}{761.124}{.}%
\TextAt{8}{277.517}{761.124}{So}%
\TextAt{8}{294.246}{761.124}{it}%
\TextAt{8}{306.42}{761.124}{suffices}%
\TextAt{8}{346.676}{761.124}{to}%
\TextAt{8}{361.456}{761.124}{show}%
\TextAt{8}{390.929}{761.124}{that}%
\TextAt{13}{416.768}{761.124}{M}%
\TextAt{21}{428.093}{759.264}{L}%
\TextAt{23}{433.852}{761.533}{\Character{48}}%
\TextAt{3}{436.536}{759.264}{\Character{106}}%
\TextAt{21}{438.888}{759.264}{L}%
\TextAt{8}{449.512}{761.124}{is}%
\TextAt{8}{461.746}{761.124}{closed.}%
\TextAt{8}{502.233}{761.124}{But}%
\TextAt{8}{525.943}{761.124}{if}%
\TextAt{13}{79.226}{744.078}{G}%
\TextAt{9}{88.46}{744.078}{\Character{40}}%
\TextAt{13}{93.012}{744.078}{L}%
\TextAt{16}{104.297}{744.078}{\Character{106}}%
\TextAt{13}{110.939}{744.078}{k}%
\TextAt{9}{117.429}{744.078}{\Character{41}}%
\TextAt{9}{125.284}{744.078}{\Character{61}}%
\TextAt{16}{137.727}{744.078}{\Character{102}}%
\TextAt{13}{143.705}{744.078}{\Character{27}}%
\TextAt{17}{150.358}{742.284}{1}%
\TextAt{13}{155.091}{744.078}{\Character{59}}%
\TextAt{13}{160.315}{744.078}{\Character{58}}%
\TextAt{13}{165.552}{744.078}{\Character{58}}%
\TextAt{13}{170.776}{744.078}{\Character{58}}%
\TextAt{13}{176.013}{744.078}{\Character{59}}%
\TextAt{13}{181.249}{744.078}{\Character{27}}%
\TextAt{21}{187.965}{742.284}{n}%
\TextAt{16}{193.601}{744.078}{\Character{103}}%
\TextAt{8}{199.579}{744.078}{,}%
\TextAt{8}{206.728}{744.078}{and}%
\TextAt{13}{229.491}{744.078}{S}%
\TextAt{21}{236.69}{742.284}{i}%
\NativeMathAt{243.392}{744.078}{11.955}{\subseteq}%
\TextAt{13}{256.011}{744.078}{G}%
\TextAt{9}{265.245}{744.078}{\Character{40}}%
\TextAt{13}{269.797}{744.078}{L}%
\TextAt{3}{277.762}{748.416}{\Character{48}}%
\TextAt{16}{283.878}{744.078}{\Character{106}}%
\TextAt{13}{290.52}{744.078}{k}%
\TextAt{9}{297.009}{744.078}{\Character{41}}%
\TextAt{8}{305.459}{744.078}{is}%
\TextAt{8}{317.223}{744.078}{the}%
\TextAt{8}{337.379}{744.078}{set}%
\TextAt{8}{355.647}{744.078}{of}%
\TextAt{8}{368.977}{744.078}{extensions}%
\TextAt{8}{425.357}{744.078}{of}%
\TextAt{13}{438.687}{744.078}{\Character{27}}%
\TextAt{21}{445.341}{742.284}{i}%
\TextAt{8}{452.619}{744.078}{to}%
\TextAt{13}{466.93}{744.078}{L}%
\TextAt{3}{474.894}{748.416}{\Character{48}}%
\TextAt{8}{477.689}{744.078}{,}%
\TextAt{8}{484.838}{744.078}{then}%
\TextAt{13}{193.881}{707.912}{M}%
\TextAt{21}{205.206}{706.052}{L}%
\TextAt{23}{210.965}{708.321}{\Character{48}}%
\TextAt{3}{213.649}{706.052}{\Character{106}}%
\TextAt{21}{216.001}{706.052}{L}%
\TextAt{9}{225.579}{707.912}{\Character{61}}%
\TextAt{21}{242.288}{722.856}{n}%
\TextAt{20}{238.215}{719.269}{\Character{91}}%
\TextAt{21}{238.005}{694.075}{i}%
\TextAt{17}{240.888}{694.075}{\Character{61}1}%
\TextAt{20}{253.701}{728.355}{\Character{32}}%
\TextAt{20}{270.951}{719.269}{\Character{89}}%
\TextAt{21}{263.166}{693.753}{K}%
\NativeMathAt{270.87}{693.753}{7.97}{\neq}%
\TextAt{21}{277.457}{693.753}{L\Character{59}L}%
\TextAt{23}{291.327}{696.023}{\Character{48}}%
\TextAt{13}{296.004}{707.911}{G}%
\TextAt{9}{305.238}{707.911}{\Character{40}}%
\TextAt{13}{309.79}{707.911}{K}%
\TextAt{16}{323.928}{707.911}{\Character{106}}%
\TextAt{13}{330.57}{707.911}{k}%
\TextAt{9}{337.06}{707.911}{\Character{41}}%
\TextAt{16}{344.269}{707.911}{\Character{2}}%
\TextAt{13}{356.224}{707.911}{S}%
\TextAt{21}{363.423}{706.118}{i}%
\TextAt{16}{369.461}{707.911}{\Character{2}}%
\TextAt{16}{381.404}{707.911}{\Character{102}}%
\TextAt{13}{387.394}{707.911}{\Character{27}}%
\TextAt{21}{394.047}{706.118}{i}%
\TextAt{16}{397.429}{707.911}{\Character{103}}%
\TextAt{20}{403.406}{728.355}{\Character{33}}%
\TextAt{13}{414.863}{707.911}{\Character{59}}%
\TextAt{8}{79.226}{670.456}{i.e.,}%
\TextAt{13}{101.331}{670.456}{M}%
\TextAt{21}{112.656}{668.596}{L}%
\TextAt{23}{118.415}{670.865}{\Character{48}}%
\TextAt{3}{121.099}{668.596}{\Character{106}}%
\TextAt{21}{123.452}{668.596}{L}%
\TextAt{8}{133.606}{670.456}{is}%
\TextAt{8}{145.37}{670.456}{indeed}%
\TextAt{8}{182.431}{670.456}{closed.}%
\Vector{q 0 G 0 g 0.399 w 452.334 -103.544 m 452.334 -95.475 l S Q}%
\Vector{q 0 G 0 g 0.399 w 452.534 -95.674 m 459.707 -95.674 l S Q}%
\Vector{q 0 G 0 g 0.399 w 452.534 -103.345 m 459.707 -103.345 l S Q}%
\Vector{q 0 G 0 g 0.399 w 459.906 -103.544 m 459.906 -95.475 l S Q}%
\TextAt{8}{97.159}{641.653}{The}%
\TextAt{8}{121.628}{641.653}{main}%
\TextAt{8}{151.31}{641.653}{theorem}%
\TextAt{8}{197.251}{641.653}{of}%
\TextAt{8}{210.997}{641.653}{Galois}%
\TextAt{8}{247.326}{641.653}{theory}%
\TextAt{8}{284.492}{641.653}{for}%
\TextAt{8}{302.792}{641.653}{infinite}%
\TextAt{8}{342.875}{641.653}{extensions}%
\TextAt{8}{399.672}{641.653}{can}%
\TextAt{8}{421.547}{641.653}{now}%
\TextAt{8}{446.351}{641.653}{be}%
\TextAt{8}{462.691}{641.653}{formulated}%
\TextAt{8}{522.297}{641.653}{as}%
\TextAt{8}{79.226}{624.607}{follows.}%
\TextAt{12}{79.226}{598.794}{Theorem}%
\TextAt{12}{135.634}{598.794}{(1.2).}%
\TextAt{8}{170.423}{598.794}{Let}%
\TextAt{9}{190.866}{598.794}{\Character{10}}%
\TextAt{16}{202.641}{598.794}{\Character{106}}%
\TextAt{13}{209.283}{598.794}{k}%
\TextAt{8}{219.143}{598.794}{be}%
\TextAt{8}{234.541}{598.794}{a}%
\TextAt{8}{243.769}{598.794}{(finite}%
\TextAt{8}{277.709}{598.794}{or}%
\TextAt{8}{291.493}{598.794}{infinite)}%
\TextAt{8}{335.189}{598.794}{Galois}%
\TextAt{8}{370.575}{598.794}{extension.}%
\TextAt{8}{426.716}{598.794}{Then}%
\TextAt{8}{456.747}{598.794}{the}%
\TextAt{8}{476.377}{598.794}{assignment}%
\TextAt{13}{263.851}{570.488}{K}%
\NativeArrowAt{277.99}{570.488}{11.955}{19.26}{\longmapsto}%
\TextAt{13}{300.572}{570.488}{G}%
\TextAt{9}{309.805}{570.488}{\Character{40}\Character{10}}%
\TextAt{16}{326.133}{570.488}{\Character{106}}%
\TextAt{13}{332.774}{570.488}{K}%
\TextAt{9}{343.592}{570.488}{\Character{41}}%
\TextAt{8}{79.226}{542.183}{is}%
\TextAt{8}{91.995}{542.183}{a}%
\TextAt{9}{102.755}{542.183}{1}%
\TextAt{8}{108.608}{542.183}{-}%
\TextAt{9}{112.505}{542.183}{1}%
\TextAt{8}{118.358}{542.183}{-correspondence}%
\TextAt{8}{204.938}{542.183}{between}%
\TextAt{8}{251.134}{542.183}{the}%
\TextAt{8}{272.295}{542.183}{subextensions}%
\TextAt{13}{347.625}{542.183}{K}%
\TextAt{16}{363.476}{542.183}{\Character{106}}%
\TextAt{13}{371.83}{542.183}{k}%
\TextAt{8}{383.222}{542.183}{of}%
\TextAt{9}{397.557}{542.183}{\Character{10}}%
\TextAt{16}{411.044}{542.183}{\Character{106}}%
\TextAt{13}{419.398}{542.183}{k}%
\TextAt{8}{430.79}{542.183}{and}%
\TextAt{8}{454.558}{542.183}{the}%
\TextAt{8}{475.719}{542.183}{closed}%
\TextAt{8}{511.25}{542.183}{sub-}%
\TextAt{8}{79.226}{525.137}{groups}%
\TextAt{8}{117.718}{525.137}{of}%
\TextAt{13}{131.749}{525.137}{G}%
\TextAt{9}{140.982}{525.137}{\Character{40}\Character{10}}%
\TextAt{16}{158.505}{525.137}{\Character{106}}%
\TextAt{13}{166.341}{525.137}{k}%
\TextAt{9}{172.831}{525.137}{\Character{41}}%
\TextAt{8}{177.383}{525.137}{.}%
\TextAt{8}{187.935}{525.137}{The}%
\TextAt{8}{212.69}{525.137}{open}%
\TextAt{8}{241.677}{525.137}{subgroups}%
\TextAt{8}{297.79}{525.137}{of}%
\TextAt{13}{311.822}{525.137}{G}%
\TextAt{9}{321.055}{525.137}{\Character{40}\Character{10}}%
\TextAt{16}{338.578}{525.137}{\Character{106}}%
\TextAt{13}{346.414}{525.137}{k}%
\TextAt{9}{352.904}{525.137}{\Character{41}}%
\TextAt{8}{362.054}{525.137}{correspond}%
\TextAt{8}{422.328}{525.137}{precisely}%
\TextAt{8}{470.886}{525.137}{to}%
\TextAt{8}{485.897}{525.137}{the}%
\TextAt{8}{506.755}{525.137}{finite}%
\TextAt{8}{79.226}{508.091}{subextensions}%
\TextAt{8}{153.551}{508.091}{of}%
\TextAt{9}{166.881}{508.091}{\Character{10}}%
\TextAt{16}{178.657}{508.091}{\Character{106}}%
\TextAt{13}{185.298}{508.091}{k}%
\TextAt{8}{191.788}{508.091}{.}%
\TextAt{19}{79.226}{473.51}{Proof.}%
\TextAt{8}{116}{473.51}{Every}%
\TextAt{8}{150.184}{473.51}{open}%
\TextAt{8}{178.999}{473.51}{subgroup}%
\TextAt{8}{230.326}{473.51}{of}%
\TextAt{13}{244.185}{473.51}{G}%
\TextAt{9}{253.419}{473.51}{\Character{40}\Character{10}}%
\TextAt{16}{270.649}{473.51}{\Character{106}}%
\TextAt{13}{278.193}{473.51}{k}%
\TextAt{9}{284.682}{473.51}{\Character{41}}%
\TextAt{8}{293.662}{473.51}{is}%
\TextAt{8}{305.955}{473.51}{also}%
\TextAt{8}{329.964}{473.51}{closed,}%
\TextAt{8}{368.404}{473.51}{because}%
\TextAt{8}{412.236}{473.51}{it}%
\TextAt{8}{424.469}{473.51}{is}%
\TextAt{8}{436.763}{473.51}{the}%
\TextAt{8}{457.448}{473.51}{complement}%
\TextAt{8}{523.337}{473.51}{of}%
\TextAt{8}{79.226}{456.464}{the}%
\TextAt{8}{100.199}{456.464}{union}%
\TextAt{8}{133.534}{456.464}{of}%
\TextAt{8}{147.681}{456.464}{its}%
\TextAt{8}{164.816}{456.464}{open}%
\TextAt{8}{193.919}{456.464}{cosets.}%
\TextAt{8}{234.861}{456.464}{If}%
\TextAt{13}{247.369}{456.464}{K}%
\TextAt{16}{262.9}{456.464}{\Character{106}}%
\TextAt{13}{270.933}{456.464}{k}%
\TextAt{8}{282.137}{456.464}{is}%
\TextAt{8}{294.717}{456.464}{a}%
\TextAt{8}{305.29}{456.464}{finite}%
\TextAt{8}{336.018}{456.464}{subextension,}%
\TextAt{8}{410.001}{456.464}{then}%
\TextAt{13}{437.478}{456.464}{G}%
\TextAt{9}{446.712}{456.464}{\Character{40}\Character{10}}%
\TextAt{16}{464.431}{456.464}{\Character{106}}%
\TextAt{13}{472.464}{456.464}{K}%
\TextAt{9}{483.282}{456.464}{\Character{41}}%
\TextAt{8}{492.549}{456.464}{is}%
\TextAt{8}{505.129}{456.464}{open,}%
\TextAt{8}{79.226}{439.418}{because}%
\TextAt{8}{123.469}{439.418}{each}%
\TextAt{13}{150.748}{439.418}{\Character{27}}%
\NativeMathAt{162.755}{439.418}{11.955}{\in}%
\TextAt{13}{175.651}{439.418}{G}%
\TextAt{9}{184.884}{439.418}{\Character{40}\Character{10}}%
\TextAt{16}{202.816}{439.418}{\Character{106}}%
\TextAt{13}{211.062}{439.418}{K}%
\TextAt{9}{221.88}{439.418}{\Character{41}}%
\TextAt{8}{231.271}{439.418}{admits}%
\TextAt{8}{270.648}{439.418}{the}%
\TextAt{8}{291.746}{439.418}{open}%
\TextAt{8}{320.973}{439.418}{neighbourhood}%
\TextAt{13}{401.919}{439.418}{\Character{27}}%
\TextAt{13}{408.985}{439.418}{G}%
\TextAt{9}{418.235}{439.418}{\Character{40}\Character{10}}%
\TextAt{16}{436.167}{439.418}{\Character{106}}%
\TextAt{13}{444.413}{439.418}{N}%
\TextAt{9}{455.035}{439.418}{\Character{41}}%
\NativeMathAt{464.513}{439.418}{11.955}{\subseteq}%
\TextAt{13}{478.736}{439.418}{G}%
\TextAt{9}{487.97}{439.418}{\Character{40}\Character{10}}%
\TextAt{16}{505.902}{439.418}{\Character{106}}%
\TextAt{13}{514.148}{439.418}{K}%
\TextAt{9}{524.965}{439.418}{\Character{41}}%
\TextAt{8}{529.518}{439.418}{,}%
\TextAt{8}{79.226}{422.372}{where}%
\TextAt{13}{113.035}{422.372}{N}%
\TextAt{16}{126.979}{422.372}{\Character{106}}%
\TextAt{13}{133.621}{422.372}{k}%
\TextAt{8}{144.007}{422.372}{is}%
\TextAt{8}{155.771}{422.372}{the}%
\TextAt{8}{175.928}{422.372}{normal}%
\TextAt{8}{215.607}{422.372}{closure}%
\TextAt{8}{254.688}{422.372}{of}%
\TextAt{13}{268.018}{422.372}{K}%
\TextAt{16}{282.157}{422.372}{\Character{106}}%
\TextAt{13}{288.799}{422.372}{k}%
\TextAt{8}{295.288}{422.372}{.}%
\TextAt{8}{303.736}{422.372}{If}%
\TextAt{13}{315.429}{422.372}{K}%
\TextAt{16}{329.567}{422.372}{\Character{106}}%
\TextAt{13}{336.209}{422.372}{k}%
\TextAt{8}{346.596}{422.372}{is}%
\TextAt{8}{358.36}{422.372}{an}%
\TextAt{8}{374.619}{422.372}{arbitrary}%
\TextAt{8}{424.388}{422.372}{subextension,}%
\TextAt{8}{497.35}{422.372}{then}%
\TextAt{13}{240.027}{394.066}{G}%
\TextAt{9}{249.26}{394.067}{\Character{40}\Character{10}}%
\TextAt{16}{265.588}{394.066}{\Character{106}}%
\TextAt{13}{272.229}{394.067}{K}%
\TextAt{9}{283.047}{394.066}{\Character{41}}%
\TextAt{9}{290.902}{394.066}{\Character{61}}%
\TextAt{20}{303.346}{405.424}{\Character{92}}%
\TextAt{21}{308.546}{380.23}{i}%
\TextAt{13}{318.622}{394.067}{G}%
\TextAt{9}{327.855}{394.067}{\Character{40}\Character{10}}%
\TextAt{16}{344.183}{394.067}{\Character{106}}%
\TextAt{13}{350.825}{394.067}{K}%
\TextAt{21}{360.784}{392.273}{i}%
\TextAt{9}{364.165}{394.067}{\Character{41}}%
\TextAt{13}{368.717}{394.067}{\Character{59}}%
\TextAt{8}{79.226}{357.812}{where}%
\TextAt{13}{113.035}{357.812}{K}%
\TextAt{21}{122.994}{356.019}{i}%
\TextAt{16}{129.697}{357.812}{\Character{106}}%
\TextAt{13}{136.338}{357.812}{k}%
\TextAt{8}{146.725}{357.812}{varies}%
\TextAt{8}{179.638}{357.812}{over}%
\TextAt{8}{204.684}{357.812}{the}%
\TextAt{8}{224.84}{357.812}{finite}%
\TextAt{8}{254.752}{357.812}{subextensions}%
\TextAt{8}{329.077}{357.812}{of}%
\TextAt{13}{342.407}{357.812}{K}%
\TextAt{16}{356.546}{357.812}{\Character{106}}%
\TextAt{13}{363.188}{357.812}{k}%
\TextAt{8}{369.677}{357.812}{.}%
\TextAt{8}{378.125}{357.812}{Therefore}%
\TextAt{13}{431.123}{357.812}{G}%
\TextAt{9}{440.356}{357.812}{\Character{40}\Character{10}}%
\TextAt{16}{456.684}{357.812}{\Character{106}}%
\TextAt{13}{463.326}{357.812}{K}%
\TextAt{9}{474.143}{357.812}{\Character{41}}%
\TextAt{8}{482.593}{357.812}{is}%
\TextAt{8}{494.357}{357.812}{closed.}%
\TextAt{8}{97.159}{337.777}{The}%
\TextAt{8}{121.589}{337.777}{assignment}%
\TextAt{13}{182.256}{337.777}{K}%
\NativeArrowAt{197.036}{337.777}{11.955}{19.26}{\longmapsto}%
\TextAt{13}{220.26}{337.777}{G}%
\TextAt{9}{229.494}{337.777}{\Character{40}\Character{10}}%
\TextAt{16}{246.463}{337.777}{\Character{106}}%
\TextAt{13}{253.747}{337.777}{K}%
\TextAt{9}{264.564}{337.777}{\Character{41}}%
\TextAt{8}{273.391}{337.777}{is}%
\TextAt{8}{285.531}{337.777}{injective,}%
\TextAt{8}{335.748}{337.777}{since}%
\TextAt{13}{364.793}{337.777}{K}%
\TextAt{8}{379.885}{337.777}{is}%
\TextAt{8}{392.025}{337.777}{the}%
\TextAt{8}{412.558}{337.777}{fixed}%
\TextAt{8}{441.221}{337.777}{field}%
\TextAt{8}{466.954}{337.777}{of}%
\TextAt{13}{480.661}{337.777}{G}%
\TextAt{9}{489.895}{337.777}{\Character{40}\Character{10}}%
\TextAt{16}{506.864}{337.777}{\Character{106}}%
\TextAt{13}{514.148}{337.777}{K}%
\TextAt{9}{524.965}{337.777}{\Character{41}}%
\TextAt{8}{529.518}{337.777}{.}%
\TextAt{8}{79.226}{320.731}{To}%
\TextAt{8}{97.782}{320.731}{prove}%
\TextAt{8}{130.661}{320.731}{surjectivity,}%
\TextAt{8}{195.791}{320.731}{we}%
\TextAt{8}{214.347}{320.731}{have}%
\TextAt{8}{242.671}{320.731}{to}%
\TextAt{8}{258.31}{320.731}{show}%
\TextAt{8}{288.642}{320.731}{that,}%
\TextAt{8}{318.923}{320.731}{given}%
\TextAt{8}{350.821}{320.731}{an}%
\TextAt{8}{368.409}{320.731}{arbitrary}%
\TextAt{8}{419.507}{320.731}{closed}%
\TextAt{8}{455.362}{320.731}{subgroup}%
\TextAt{13}{507.488}{320.731}{H}%
\TextAt{8}{523.337}{320.731}{of}%
\TextAt{13}{79.226}{303.685}{G}%
\TextAt{9}{88.46}{303.685}{\Character{40}\Character{10}}%
\TextAt{16}{104.787}{303.685}{\Character{106}}%
\TextAt{13}{111.429}{303.685}{k}%
\TextAt{9}{117.918}{303.685}{\Character{41}}%
\TextAt{8}{122.471}{303.685}{,}%
\TextAt{8}{129.62}{303.685}{we}%
\TextAt{8}{146.847}{303.685}{always}%
\TextAt{8}{184.315}{303.685}{have}%
\TextAt{13}{267.401}{286.639}{H}%
\TextAt{9}{281.345}{286.639}{\Character{61}}%
\TextAt{13}{293.77}{286.639}{G}%
\TextAt{9}{303.004}{286.639}{\Character{40}\Character{10}}%
\TextAt{16}{319.331}{286.639}{\Character{106}}%
\TextAt{13}{325.973}{286.639}{K}%
\TextAt{9}{336.791}{286.639}{\Character{41}}%
\TextAt{13}{341.343}{286.639}{\Character{59}}%
\TextAt{8}{79.226}{263.416}{where}%
\TextAt{13}{113.454}{263.416}{K}%
\TextAt{8}{128.587}{263.416}{is}%
\TextAt{8}{140.769}{263.416}{the}%
\TextAt{8}{161.344}{263.416}{fixed}%
\TextAt{8}{190.048}{263.416}{field}%
\TextAt{8}{215.824}{263.416}{of}%
\TextAt{13}{229.572}{263.416}{H}%
\TextAt{8}{240.195}{263.416}{.}%
\TextAt{8}{249.898}{263.416}{The}%
\TextAt{8}{274.37}{263.416}{inclusion}%
\TextAt{13}{323.625}{263.416}{H}%
\NativeMathAt{338.282}{263.416}{11.955}{\subseteq}%
\TextAt{13}{351.614}{263.416}{G}%
\TextAt{9}{360.848}{263.416}{\Character{40}\Character{10}}%
\TextAt{16}{377.888}{263.416}{\Character{106}}%
\TextAt{13}{385.243}{263.416}{K}%
\TextAt{9}{396.06}{263.416}{\Character{41}}%
\TextAt{8}{404.928}{263.416}{is}%
\TextAt{8}{417.111}{263.416}{trivial.}%
\TextAt{8}{457.718}{263.416}{Conversely,}%
\TextAt{8}{519.762}{263.416}{let}%
\TextAt{13}{79.226}{246.37}{\Character{27}}%
\NativeMathAt{90.935}{246.37}{11.955}{\in}%
\TextAt{13}{103.532}{246.37}{G}%
\TextAt{9}{112.766}{246.37}{\Character{40}\Character{10}}%
\TextAt{16}{130.399}{246.37}{\Character{106}}%
\TextAt{13}{138.346}{246.37}{K}%
\TextAt{9}{149.164}{246.37}{\Character{41}}%
\TextAt{8}{153.716}{246.37}{.}%
\TextAt{8}{164.463}{246.37}{If}%
\TextAt{13}{176.922}{246.37}{L}%
\TextAt{16}{189.513}{246.37}{\Character{106}}%
\TextAt{13}{197.461}{246.37}{K}%
\TextAt{8}{212.942}{246.37}{is}%
\TextAt{8}{225.472}{246.37}{a}%
\TextAt{8}{235.994}{246.37}{finite}%
\TextAt{8}{266.672}{246.37}{Galois}%
\TextAt{8}{303.351}{246.37}{subextension}%
\TextAt{8}{373.828}{246.37}{of}%
\TextAt{9}{387.925}{246.37}{\Character{10}}%
\TextAt{16}{401.005}{246.37}{\Character{106}}%
\TextAt{13}{408.953}{246.37}{K}%
\TextAt{8}{419.771}{246.37}{,}%
\TextAt{8}{427.878}{246.37}{then}%
\TextAt{13}{455.304}{246.37}{\Character{27}}%
\TextAt{13}{462.37}{246.37}{G}%
\TextAt{9}{471.62}{246.37}{\Character{40}\Character{10}}%
\TextAt{16}{489.253}{246.37}{\Character{106}}%
\TextAt{13}{497.201}{246.37}{L}%
\TextAt{9}{505.165}{246.37}{\Character{41}}%
\TextAt{8}{514.381}{246.37}{is}%
\TextAt{8}{526.911}{246.37}{a}%
\TextAt{8}{79.226}{229.324}{fundamental}%
\TextAt{8}{145.875}{229.324}{open}%
\TextAt{8}{173.167}{229.324}{neighbourhood}%
\TextAt{8}{252.177}{229.324}{of}%
\TextAt{13}{264.514}{229.324}{\Character{27}}%
\TextAt{8}{274.5}{229.324}{in}%
\TextAt{13}{287.159}{229.324}{G}%
\TextAt{9}{296.393}{229.324}{\Character{40}\Character{10}}%
\TextAt{16}{312.72}{229.324}{\Character{106}}%
\TextAt{13}{319.362}{229.324}{K}%
\TextAt{9}{330.18}{229.324}{\Character{41}}%
\TextAt{8}{334.732}{229.324}{.}%
\TextAt{8}{342.849}{229.324}{The}%
\TextAt{8}{365.909}{229.324}{map}%
\TextAt{13}{390.93}{229.324}{H}%
\NativeArrowAt{404.874}{229.324}{11.955}{19.26}{\longrightarrow}%
\TextAt{13}{427.456}{229.324}{G}%
\TextAt{9}{436.689}{229.324}{\Character{40}}%
\TextAt{13}{441.242}{229.324}{L}%
\TextAt{16}{452.527}{229.324}{\Character{106}}%
\TextAt{13}{459.169}{229.324}{K}%
\TextAt{9}{469.986}{229.324}{\Character{41}}%
\TextAt{8}{477.442}{229.324}{is}%
\TextAt{8}{488.213}{229.324}{certainly}%
\TextAt{8}{79.226}{212.278}{surjective,}%
\TextAt{8}{134.984}{212.278}{because}%
\TextAt{8}{178.361}{212.278}{the}%
\TextAt{8}{198.593}{212.278}{image}%
\TextAt{9}{235.524}{215.3}{\Character{22}}%
\TextAt{13}{232.489}{212.278}{H}%
\TextAt{8}{247.085}{212.278}{has}%
\TextAt{8}{268.034}{212.278}{fixed}%
\TextAt{8}{296.395}{212.278}{field}%
\TextAt{13}{321.828}{212.278}{K}%
\TextAt{8}{336.618}{212.278}{and}%
\TextAt{8}{359.456}{212.278}{is}%
\TextAt{8}{371.295}{212.278}{therefore}%
\TextAt{8}{420.471}{212.278}{equal}%
\TextAt{8}{451.438}{212.278}{to}%
\TextAt{13}{465.824}{212.278}{G}%
\TextAt{9}{475.058}{212.278}{\Character{40}}%
\TextAt{13}{479.61}{212.278}{L}%
\TextAt{16}{491.024}{212.278}{\Character{106}}%
\TextAt{13}{497.794}{212.278}{K}%
\TextAt{9}{508.612}{212.278}{\Character{41}}%
\TextAt{8}{513.164}{212.278}{,}%
\TextAt{8}{520.408}{212.278}{by}%
\TextAt{8}{79.226}{195.232}{the}%
\TextAt{8}{100.057}{195.232}{main}%
\TextAt{8}{129.997}{195.232}{theorem}%
\TextAt{8}{176.197}{195.232}{of}%
\TextAt{8}{190.201}{195.232}{Galois}%
\TextAt{8}{226.789}{195.232}{theory}%
\TextAt{8}{264.213}{195.232}{for}%
\TextAt{8}{282.773}{195.232}{finite}%
\TextAt{8}{313.359}{195.232}{extensions.}%
\TextAt{8}{376.313}{195.232}{Thus}%
\TextAt{8}{406.636}{195.232}{we}%
\TextAt{8}{424.538}{195.232}{may}%
\TextAt{8}{450.581}{195.232}{choose}%
\TextAt{8}{488.388}{195.232}{a}%
\TextAt{13}{498.818}{195.232}{\Character{28}}%
\NativeMathAt{509.707}{195.232}{11.955}{\in}%
\TextAt{13}{522.147}{195.232}{H}%
\TextAt{8}{79.226}{178.186}{such}%
\TextAt{8}{105.452}{178.186}{that}%
\TextAt{13}{130.649}{178.186}{\Character{28}}%
\TextAt{16}{137.068}{178.186}{\Character{106}}%
\TextAt{21}{140.389}{176.393}{L}%
\TextAt{9}{149.967}{178.186}{\Character{61}}%
\TextAt{13}{162.393}{178.186}{\Character{27}}%
\TextAt{16}{169.475}{178.186}{\Character{106}}%
\TextAt{21}{172.796}{176.393}{L}%
\TextAt{8}{179.053}{178.186}{,}%
\TextAt{8}{186.065}{178.186}{i.e.,}%
\TextAt{13}{208.033}{178.186}{\Character{28}}%
\NativeMathAt{217.773}{178.186}{11.955}{\in}%
\TextAt{13}{229.064}{178.186}{H}%
\TextAt{16}{241.993}{178.186}{\Character{92}}%
\TextAt{13}{252.269}{178.186}{\Character{27}}%
\TextAt{13}{259.335}{178.186}{G}%
\TextAt{9}{268.585}{178.186}{\Character{40}\Character{10}}%
\TextAt{16}{284.913}{178.186}{\Character{106}}%
\TextAt{13}{291.555}{178.186}{L}%
\TextAt{9}{299.519}{178.186}{\Character{41}}%
\TextAt{8}{304.071}{178.186}{.}%
\TextAt{8}{312.463}{178.186}{This}%
\TextAt{8}{339.011}{178.186}{shows}%
\TextAt{8}{372.458}{178.186}{that}%
\TextAt{13}{397.655}{178.186}{\Character{27}}%
\TextAt{8}{408.464}{178.186}{belongs}%
\TextAt{8}{450.303}{178.186}{to}%
\TextAt{8}{464.442}{178.186}{the}%
\TextAt{8}{484.427}{178.186}{closure}%
\TextAt{8}{523.337}{178.186}{of}%
\TextAt{13}{79.226}{161.14}{H}%
\TextAt{8}{93.746}{161.14}{in}%
\TextAt{13}{107.399}{161.14}{G}%
\TextAt{9}{116.633}{161.14}{\Character{40}\Character{10}}%
\TextAt{16}{132.96}{161.14}{\Character{106}}%
\TextAt{13}{139.602}{161.14}{K}%
\TextAt{9}{150.419}{161.14}{\Character{41}}%
\TextAt{8}{154.972}{161.14}{,}%
\TextAt{8}{162.121}{161.14}{and}%
\TextAt{8}{184.884}{161.14}{thus}%
\TextAt{8}{210.635}{161.14}{to}%
\TextAt{13}{224.945}{161.14}{H}%
\TextAt{8}{239.465}{161.14}{itself,}%
\TextAt{8}{271.063}{161.14}{so}%
\TextAt{8}{285.433}{161.14}{that}%
\TextAt{13}{310.802}{161.14}{H}%
\TextAt{9}{324.745}{161.14}{\Character{61}}%
\TextAt{13}{337.171}{161.14}{G}%
\TextAt{9}{346.404}{161.14}{\Character{40}\Character{10}}%
\TextAt{16}{362.732}{161.14}{\Character{106}}%
\TextAt{13}{369.374}{161.14}{K}%
\TextAt{9}{380.191}{161.14}{\Character{41}}%
\TextAt{8}{384.744}{161.14}{.}%
\TextAt{8}{97.159}{141.105}{If}%
\TextAt{13}{108.736}{141.105}{H}%
\TextAt{8}{123.141}{141.105}{is}%
\TextAt{8}{134.79}{141.105}{an}%
\TextAt{8}{150.934}{141.105}{open}%
\TextAt{8}{179.105}{141.105}{subgroup}%
\TextAt{8}{229.787}{141.105}{of}%
\TextAt{13}{243.002}{141.105}{G}%
\TextAt{9}{252.236}{141.105}{\Character{40}\Character{10}}%
\TextAt{16}{268.563}{141.105}{\Character{106}}%
\TextAt{13}{275.205}{141.105}{k}%
\TextAt{9}{281.694}{141.105}{\Character{41}}%
\TextAt{8}{286.247}{141.105}{,}%
\TextAt{8}{293.304}{141.105}{then}%
\TextAt{8}{319.849}{141.105}{it}%
\TextAt{8}{331.438}{141.105}{is}%
\TextAt{8}{343.087}{141.105}{also}%
\TextAt{8}{366.452}{141.105}{closed,}%
\TextAt{8}{404.138}{141.105}{and}%
\TextAt{8}{426.785}{141.105}{therefore}%
\TextAt{8}{475.77}{141.105}{of}%
\TextAt{8}{488.985}{141.105}{the}%
\TextAt{8}{509.027}{141.105}{form}%
\TextAt{13}{79.226}{124.059}{H}%
\TextAt{9}{93.17}{124.059}{\Character{61}}%
\TextAt{13}{105.595}{124.059}{G}%
\TextAt{9}{114.829}{124.059}{\Character{40}\Character{10}}%
\TextAt{16}{131.156}{124.059}{\Character{106}}%
\TextAt{13}{137.798}{124.059}{K}%
\TextAt{9}{148.616}{124.059}{\Character{41}}%
\TextAt{8}{153.168}{124.059}{.}%
\TextAt{8}{161.425}{124.059}{But}%
\TextAt{13}{184.092}{124.059}{G}%
\TextAt{9}{193.325}{124.059}{\Character{40}\Character{10}}%
\TextAt{16}{209.653}{124.059}{\Character{106}}%
\TextAt{13}{216.294}{124.059}{k}%
\TextAt{9}{222.784}{124.059}{\Character{41}}%
\TextAt{8}{230.66}{124.059}{is}%
\TextAt{8}{241.849}{124.059}{the}%
\TextAt{8}{261.432}{124.059}{disjoint}%
\TextAt{8}{302.545}{124.059}{union}%
\TextAt{8}{334.489}{124.059}{of}%
\TextAt{8}{347.245}{124.059}{the}%
\TextAt{8}{366.828}{124.059}{open}%
\TextAt{8}{394.539}{124.059}{cosets}%
\TextAt{8}{427.906}{124.059}{of}%
\TextAt{13}{440.662}{124.059}{H}%
\TextAt{8}{451.285}{124.059}{.}%
\TextAt{8}{459.542}{124.059}{Since}%
\TextAt{13}{489.525}{124.059}{G}%
\TextAt{9}{498.758}{124.059}{\Character{40}\Character{10}}%
\TextAt{16}{515.086}{124.059}{\Character{106}}%
\TextAt{13}{521.728}{124.059}{k}%
\TextAt{9}{528.217}{124.059}{\Character{41}}%
\TextAt{8}{79.226}{107.013}{is}%
\TextAt{8}{91.925}{107.013}{compact,}%
\TextAt{8}{143.174}{107.013}{a}%
\TextAt{8}{153.864}{107.013}{finite}%
\TextAt{8}{184.711}{107.013}{number}%
\TextAt{8}{228.243}{107.013}{of}%
\TextAt{8}{242.508}{107.013}{cosets}%
\TextAt{8}{277.383}{107.013}{suffices}%
\TextAt{8}{318.105}{107.013}{to}%
\TextAt{8}{333.351}{107.013}{cover}%
\TextAt{8}{364.532}{107.013}{the}%
\TextAt{8}{385.624}{107.013}{group.}%
\TextAt{8}{426.155}{107.013}{Thus}%
\TextAt{8}{456.739}{107.013}{there}%
\TextAt{8}{487.586}{107.013}{is}%
\TextAt{8}{500.285}{107.013}{only}%
\TextAt{8}{526.911}{107.013}{a}%
\TextAt{8}{79.226}{89.967}{finite}%
\TextAt{8}{109.286}{89.967}{number}%
\TextAt{8}{152.03}{89.967}{of}%
\TextAt{8}{165.508}{89.967}{them;}%
\TextAt{13}{198.894}{89.967}{H}%
\TextAt{9}{213.09}{89.967}{\Character{61}}%
\TextAt{13}{225.767}{89.967}{G}%
\TextAt{9}{235.001}{89.967}{\Character{40}\Character{10}}%
\TextAt{16}{251.581}{89.967}{\Character{106}}%
\TextAt{13}{258.475}{89.967}{K}%
\TextAt{9}{269.292}{89.967}{\Character{41}}%
\TextAt{8}{277.89}{89.967}{has}%
\TextAt{8}{298.912}{89.967}{finite}%
\TextAt{8}{328.972}{89.967}{index}%
\TextAt{8}{360.658}{89.967}{in}%
\TextAt{13}{374.458}{89.967}{G}%
\TextAt{9}{383.692}{89.967}{\Character{40}\Character{10}}%
\TextAt{16}{400.272}{89.967}{\Character{106}}%
\TextAt{13}{407.166}{89.967}{k}%
\TextAt{9}{413.655}{89.967}{\Character{41}}%
\TextAt{8}{418.208}{89.967}{,}%
\TextAt{8}{425.542}{89.967}{and}%
\TextAt{8}{448.452}{89.967}{this}%
\TextAt{8}{471.423}{89.967}{implies}%
\TextAt{8}{511.298}{89.967}{that}%
\TextAt{13}{79.226}{72.921}{K}%
\TextAt{16}{93.365}{72.921}{\Character{106}}%
\TextAt{13}{100.006}{72.921}{k}%
\TextAt{8}{110.393}{72.921}{has}%
\TextAt{8}{131.267}{72.921}{finite}%
\TextAt{8}{161.179}{72.921}{degree.}%
\Vector{q 0 G 0 g 0.399 w 452.334 -701.079 m 452.334 -693.009 l S Q}%
\Vector{q 0 G 0 g 0.399 w 452.534 -693.208 m 459.707 -693.208 l S Q}%
\Vector{q 0 G 0 g 0.399 w 452.534 -700.879 m 459.707 -700.879 l S Q}%
\Vector{q 0 G 0 g 0.399 w 459.906 -701.079 m 459.906 -693.009 l S Q}%
\TextAt{8}{297.211}{43.033}{241}%
\Destination{78.23}{809.815}{nav.page.241}%
\Destination{79.36}{776.155}{nav.print.263}%
\Destination{79.36}{193.155}{nav.print.264}%
\Destination{79.36}{612.755}{nav.number.4.1.2}%
}

% Page 254 (body page 242)
\Page{%
\TextAt{2}{79.226}{796.943}{IV.}%
\TextAt{2}{98.012}{796.943}{Abstract}%
\TextAt{2}{143.252}{796.943}{Class}%
\TextAt{2}{171.844}{796.943}{Field}%
\TextAt{2}{199.575}{796.943}{Theory}%
\Vector{q 0 G 0 g 0.299 w 7.226 18.729 m 460.769 18.729 l S Q}%
\TextAt{8}{97.159}{761.124}{The}%
\TextAt{8}{121.774}{761.124}{topological}%
\TextAt{8}{181.86}{761.124}{Galois}%
\TextAt{8}{218.334}{761.124}{groups}%
\TextAt{13}{256.686}{761.124}{G}%
\TextAt{9}{270.197}{761.124}{\Character{61}}%
\TextAt{13}{283.578}{761.124}{G}%
\TextAt{9}{292.812}{761.124}{\Character{40}\Character{10}}%
\TextAt{16}{310.096}{761.124}{\Character{106}}%
\TextAt{13}{317.694}{761.124}{k}%
\TextAt{9}{324.183}{761.124}{\Character{41}}%
\TextAt{8}{333.194}{761.124}{have}%
\TextAt{8}{360.75}{761.124}{the}%
\TextAt{8}{381.468}{761.124}{special}%
\TextAt{8}{420.13}{761.124}{property}%
\TextAt{8}{468.5}{761.124}{that}%
\TextAt{8}{494.43}{761.124}{there}%
\TextAt{8}{524.903}{761.124}{is}%
\TextAt{8}{79.226}{744.078}{a}%
\TextAt{8}{89.157}{744.078}{fundamental}%
\TextAt{8}{156.974}{744.078}{system}%
\TextAt{8}{195.968}{744.078}{of}%
\TextAt{8}{209.473}{744.078}{neighbourhoods}%
\TextAt{8}{294.267}{744.078}{of}%
\TextAt{8}{307.773}{744.078}{the}%
\TextAt{8}{328.104}{744.078}{neutral}%
\TextAt{8}{368.604}{744.078}{element}%
\TextAt{9}{412.021}{744.078}{1}%
\NativeMathAt{421.494}{744.078}{11.955}{\in}%
\TextAt{13}{433.083}{744.078}{G}%
\TextAt{8}{446.39}{744.078}{which}%
\TextAt{8}{480.051}{744.078}{consists}%
\TextAt{8}{523.337}{744.078}{of}%
\TextAt{8}{79.226}{727.032}{normal}%
\TextAt{8}{118.119}{727.032}{subgroups.}%
\TextAt{8}{177.82}{727.032}{This}%
\TextAt{8}{203.753}{727.032}{property}%
\TextAt{8}{250.775}{727.032}{leads}%
\TextAt{8}{279.315}{727.032}{us}%
\TextAt{8}{293.544}{727.032}{to}%
\TextAt{8}{307.068}{727.032}{the}%
\TextAt{8}{326.438}{727.032}{abstract,}%
\TextAt{8}{374.657}{727.032}{purely}%
\TextAt{8}{409.963}{727.032}{group-theoretical}%
\TextAt{8}{500.239}{727.032}{notion}%
\TextAt{8}{79.226}{709.986}{of}%
\TextAt{8}{92.556}{709.986}{a}%
\TextAt{8}{102.312}{709.986}{profinite}%
\TextAt{8}{149.14}{709.986}{group.}%
\TextAt{12}{79.226}{684.887}{Definition}%
\TextAt{12}{142.522}{684.887}{(1.3).}%
\TextAt{8}{178.279}{684.887}{A}%
\TextAt{12}{191.392}{684.887}{profinite}%
\TextAt{12}{246.14}{684.887}{group}%
\TextAt{8}{284.383}{684.887}{is}%
\TextAt{8}{296.588}{684.887}{a}%
\TextAt{8}{306.784}{684.887}{topological}%
\TextAt{8}{366.75}{684.887}{group}%
\TextAt{13}{400.367}{684.887}{G}%
\TextAt{8}{413.939}{684.887}{which}%
\TextAt{8}{447.866}{684.887}{is}%
\TextAt{8}{460.071}{684.887}{Hausdorff}%
\TextAt{8}{513.904}{684.887}{and}%
\TextAt{8}{79.226}{667.841}{compact,}%
\TextAt{8}{131.078}{667.841}{and}%
\TextAt{8}{155.257}{667.841}{which}%
\TextAt{8}{190.161}{667.841}{admits}%
\TextAt{8}{230.013}{667.841}{a}%
\TextAt{8}{241.186}{667.841}{basis}%
\TextAt{8}{271.343}{667.841}{of}%
\TextAt{8}{286.09}{667.841}{neighbourhoods}%
\TextAt{8}{372.126}{667.841}{of}%
\TextAt{9}{386.873}{667.841}{1}%
\NativeMathAt{398.461}{667.841}{11.955}{\in}%
\TextAt{13}{412.167}{667.841}{G}%
\TextAt{8}{426.715}{667.841}{consisting}%
\TextAt{8}{482.241}{667.841}{of}%
\TextAt{8}{496.988}{667.841}{normal}%
\TextAt{8}{79.226}{650.795}{subgroups.}%
\TextAt{8}{97.159}{625.696}{It}%
\TextAt{8}{109.168}{625.696}{can}%
\TextAt{8}{129.963}{625.696}{be}%
\TextAt{8}{145.224}{625.696}{shown}%
\TextAt{8}{180.067}{625.696}{that}%
\TextAt{8}{204.772}{625.696}{the}%
\TextAt{8}{224.264}{625.696}{last}%
\TextAt{8}{245.778}{625.696}{condition}%
\TextAt{8}{296.497}{625.696}{is}%
\TextAt{8}{307.597}{625.696}{tantamount}%
\TextAt{8}{370.69}{625.696}{to}%
\TextAt{13}{384.337}{625.696}{G}%
\TextAt{8}{396.804}{625.696}{being}%
\TextAt{8}{427.678}{625.696}{totally}%
\TextAt{8}{464.422}{625.696}{disconnected,}%
\TextAt{8}{79.226}{608.65}{i.e.,}%
\TextAt{8}{102.372}{608.65}{to}%
\TextAt{8}{117.514}{608.65}{the}%
\TextAt{8}{138.503}{608.65}{condition}%
\TextAt{8}{190.718}{608.65}{that}%
\TextAt{8}{216.92}{608.65}{each}%
\TextAt{8}{244.089}{608.65}{element}%
\TextAt{8}{288.163}{608.65}{of}%
\TextAt{13}{302.326}{608.65}{G}%
\TextAt{8}{316.289}{608.65}{is}%
\TextAt{8}{328.885}{608.65}{equal}%
\TextAt{8}{360.609}{608.65}{to}%
\TextAt{8}{375.752}{608.65}{its}%
\TextAt{8}{392.903}{608.65}{own}%
\TextAt{8}{418.124}{608.65}{connected}%
\TextAt{8}{473.579}{608.65}{component.}%
\TextAt{8}{79.226}{591.604}{Every}%
\TextAt{8}{113.713}{591.604}{closed}%
\TextAt{8}{149.073}{591.604}{subgroup}%
\TextAt{13}{200.704}{591.604}{H}%
\TextAt{8}{216.058}{591.604}{of}%
\TextAt{13}{230.221}{591.604}{G}%
\TextAt{8}{244.186}{591.604}{is}%
\TextAt{8}{256.783}{591.604}{obviously}%
\TextAt{8}{309.394}{591.604}{again}%
\TextAt{8}{341.455}{591.604}{a}%
\TextAt{8}{352.044}{591.604}{profinite}%
\TextAt{8}{399.706}{591.604}{group.}%
\TextAt{8}{439.932}{591.604}{The}%
\TextAt{8}{464.82}{591.604}{disjoint}%
\TextAt{8}{507.341}{591.604}{coset}%
\TextAt{8}{79.226}{574.558}{decomposition}%
\TextAt{13}{275.541}{557.513}{G}%
\TextAt{9}{288.096}{557.513}{\Character{61}}%
\TextAt{20}{300.521}{568.87}{\Character{91}}%
\TextAt{21}{305.721}{543.676}{i}%
\TextAt{13}{315.797}{557.513}{\Character{27}}%
\TextAt{21}{322.451}{555.719}{i}%
\TextAt{13}{325.832}{557.513}{H}%
\TextAt{8}{79.226}{527.412}{shows}%
\TextAt{8}{112.844}{527.412}{immediately}%
\TextAt{8}{179.506}{527.412}{that}%
\TextAt{13}{204.875}{527.412}{H}%
\TextAt{8}{219.395}{527.412}{is}%
\TextAt{8}{231.159}{527.412}{open}%
\TextAt{8}{259.445}{527.412}{if}%
\TextAt{8}{270.169}{527.412}{and}%
\TextAt{8}{292.931}{527.412}{only}%
\TextAt{8}{318.623}{527.412}{if}%
\TextAt{8}{329.347}{527.412}{the}%
\TextAt{8}{349.503}{527.412}{index}%
\TextAt{9}{381.041}{527.412}{\Character{40}}%
\TextAt{13}{385.593}{527.412}{G}%
\TextAt{9}{398.148}{527.412}{\Character{58}}%
\TextAt{13}{404.72}{527.412}{H}%
\TextAt{9}{415.343}{527.412}{\Character{41}}%
\TextAt{8}{423.792}{527.412}{is}%
\TextAt{8}{435.556}{527.412}{finite.}%
\TextAt{8}{97.159}{507.377}{Profinite}%
\TextAt{8}{145.798}{507.377}{groups}%
\TextAt{8}{183.94}{507.377}{are}%
\TextAt{8}{203.803}{507.377}{fairly}%
\TextAt{8}{234.725}{507.377}{close}%
\TextAt{8}{263.1}{507.377}{relatives}%
\TextAt{8}{309.694}{507.377}{of}%
\TextAt{8}{323.377}{507.377}{finite}%
\TextAt{8}{353.641}{507.377}{groups.}%
\TextAt{8}{397.038}{507.377}{They}%
\TextAt{8}{427.625}{507.377}{can}%
\TextAt{8}{449.437}{507.377}{be}%
\TextAt{8}{465.713}{507.377}{reconstituted}%
\TextAt{8}{79.226}{490.331}{rather}%
\TextAt{8}{114.907}{490.331}{easily}%
\TextAt{8}{147.719}{490.331}{from}%
\TextAt{8}{175.917}{490.331}{their}%
\TextAt{8}{204.437}{490.331}{finite}%
\TextAt{8}{234.905}{490.331}{quotients.}%
\TextAt{8}{291.925}{490.331}{For}%
\TextAt{8}{313.451}{490.331}{the}%
\TextAt{8}{334.164}{490.331}{precise}%
\TextAt{8}{373.145}{490.331}{description}%
\TextAt{8}{433.597}{490.331}{of}%
\TextAt{8}{447.484}{490.331}{this}%
\TextAt{8}{470.864}{490.331}{we}%
\TextAt{8}{488.648}{490.331}{need}%
\TextAt{8}{516.51}{490.331}{the}%
\TextAt{8}{79.226}{473.286}{notion}%
\TextAt{8}{115.53}{473.286}{of}%
\TextAt{8}{128.737}{473.286}{projective}%
\TextAt{8}{182.591}{473.286}{limit,}%
\TextAt{8}{213.708}{473.286}{which}%
\TextAt{8}{247.071}{473.286}{naturally}%
\TextAt{8}{297.362}{473.286}{occurs}%
\TextAt{8}{333.392}{473.286}{in}%
\TextAt{8}{346.921}{473.286}{various}%
\TextAt{8}{386.872}{473.286}{places}%
\TextAt{8}{421.275}{473.286}{in}%
\TextAt{8}{434.804}{473.286}{number}%
\TextAt{8}{477.277}{473.286}{theory}%
\TextAt{8}{513.904}{473.286}{and}%
\TextAt{8}{79.226}{456.24}{which}%
\TextAt{8}{112.713}{456.24}{we}%
\TextAt{8}{129.94}{456.24}{will}%
\TextAt{8}{152.045}{456.24}{introduce}%
\TextAt{8}{204.074}{456.24}{next.}%
\TextAt{12}{79.226}{431.141}{Exercise}%
\TextAt{12}{132.003}{431.141}{IV.1.1.}%
\TextAt{8}{175.39}{431.141}{Let}%
\TextAt{13}{195.991}{431.141}{L}%
\TextAt{16}{207.276}{431.141}{\Character{106}}%
\TextAt{13}{213.918}{431.141}{k}%
\TextAt{8}{223.936}{431.141}{be}%
\TextAt{8}{239.491}{431.141}{a}%
\TextAt{8}{248.878}{431.141}{Galois}%
\TextAt{8}{284.423}{431.141}{extension}%
\TextAt{8}{335.82}{431.141}{and}%
\TextAt{13}{358.214}{431.141}{K}%
\TextAt{16}{372.353}{431.141}{\Character{106}}%
\TextAt{13}{378.994}{431.141}{k}%
\TextAt{8}{389.013}{431.141}{an}%
\TextAt{8}{404.903}{431.141}{arbitrary}%
\TextAt{8}{454.304}{431.141}{extension,}%
\TextAt{8}{509.027}{431.141}{both}%
\TextAt{8}{79.226}{414.095}{contained}%
\TextAt{8}{132.235}{414.095}{in}%
\TextAt{8}{145.888}{414.095}{a}%
\TextAt{8}{155.644}{414.095}{common}%
\TextAt{8}{202.472}{414.095}{extension}%
\TextAt{9}{254.238}{414.095}{\Character{10}}%
\TextAt{16}{266.013}{414.095}{\Character{106}}%
\TextAt{13}{272.655}{414.095}{k}%
\TextAt{8}{279.144}{414.095}{.}%
\TextAt{8}{287.593}{414.095}{If}%
\TextAt{13}{299.285}{414.095}{L}%
\TextAt{16}{309.906}{414.095}{\Character{92}}%
\TextAt{13}{320.533}{414.095}{K}%
\TextAt{9}{334.671}{414.095}{\Character{61}}%
\TextAt{13}{347.097}{414.095}{k}%
\TextAt{8}{353.586}{414.095}{,}%
\TextAt{8}{360.735}{414.095}{then}%
\TextAt{8}{387.395}{414.095}{the}%
\TextAt{8}{407.552}{414.095}{mapping}%
\TextAt{13}{201.693}{388.291}{G}%
\TextAt{9}{210.926}{388.291}{\Character{40}}%
\TextAt{13}{215.479}{388.291}{LK}%
\TextAt{16}{237.582}{388.291}{\Character{106}}%
\TextAt{13}{244.223}{388.291}{K}%
\TextAt{9}{255.041}{388.291}{\Character{41}}%
\NativeArrowAt{262.914}{388.291}{11.955}{19.26}{\longrightarrow}%
\TextAt{13}{285.496}{388.291}{G}%
\TextAt{9}{294.73}{388.291}{\Character{40}}%
\TextAt{13}{299.282}{388.291}{L}%
\TextAt{16}{310.567}{388.291}{\Character{106}}%
\TextAt{13}{317.209}{388.291}{k}%
\TextAt{9}{323.699}{388.291}{\Character{41}}%
\TextAt{13}{328.251}{388.291}{\Character{59}}%
\TextAt{13}{357.386}{388.291}{\Character{27}}%
\NativeArrowAt{367.809}{388.291}{11.955}{19.26}{\longmapsto}%
\TextAt{13}{390.391}{388.291}{\Character{27}}%
\TextAt{16}{397.473}{388.291}{\Character{106}}%
\TextAt{21}{400.794}{386.498}{L}%
\TextAt{13}{407.051}{388.291}{\Character{59}}%
\TextAt{8}{79.226}{362.488}{is}%
\TextAt{8}{90.579}{362.488}{a}%
\TextAt{8}{99.924}{362.488}{topological}%
\TextAt{8}{159.038}{362.488}{isomorphism,}%
\TextAt{8}{230.38}{362.488}{that}%
\TextAt{8}{255.338}{362.488}{is,}%
\TextAt{8}{270.025}{362.488}{an}%
\TextAt{8}{285.874}{362.488}{isomorphism}%
\TextAt{8}{353.882}{362.488}{of}%
\TextAt{8}{366.801}{362.488}{groups}%
\TextAt{8}{404.181}{362.488}{and}%
\TextAt{8}{426.533}{362.488}{a}%
\TextAt{8}{435.877}{362.488}{homeomorphism}%
\TextAt{8}{523.337}{362.488}{of}%
\TextAt{8}{79.226}{345.442}{topological}%
\TextAt{8}{138.751}{345.442}{spaces.}%
\TextAt{12}{79.226}{315.279}{Exercise}%
\TextAt{12}{132.982}{315.279}{IV.1.2.}%
\TextAt{8}{177.221}{315.279}{Given}%
\TextAt{8}{211.597}{315.279}{a}%
\TextAt{8}{221.835}{315.279}{family}%
\TextAt{8}{258.088}{315.279}{of}%
\TextAt{8}{271.901}{315.279}{Galois}%
\TextAt{8}{308.297}{315.279}{extensions}%
\TextAt{13}{365.16}{315.279}{K}%
\TextAt{21}{375.12}{313.486}{i}%
\TextAt{16}{382.645}{315.279}{\Character{106}}%
\TextAt{13}{390.109}{315.279}{k}%
\TextAt{8}{400.979}{315.279}{in}%
\TextAt{9}{415.115}{315.279}{\Character{10}}%
\TextAt{16}{427.713}{315.279}{\Character{106}}%
\TextAt{13}{435.178}{315.279}{k}%
\TextAt{8}{441.667}{315.279}{,}%
\TextAt{8}{449.42}{315.279}{let}%
\TextAt{13}{466.807}{315.279}{K}%
\TextAt{16}{481.769}{315.279}{\Character{106}}%
\TextAt{13}{489.233}{315.279}{k}%
\TextAt{8}{500.103}{315.279}{be}%
\TextAt{8}{516.51}{315.279}{the}%
\TextAt{8}{79.226}{298.233}{composite}%
\TextAt{8}{133.901}{298.233}{of}%
\TextAt{8}{146.888}{298.233}{all}%
\TextAt{13}{162.804}{298.233}{K}%
\TextAt{21}{172.763}{296.44}{i}%
\TextAt{16}{179.465}{298.233}{\Character{106}}%
\TextAt{13}{186.107}{298.233}{k}%
\TextAt{8}{192.596}{298.233}{,}%
\TextAt{8}{199.471}{298.233}{and}%
\TextAt{13}{221.89}{298.233}{K}%
\TextAt{3}{232.708}{302.572}{\Character{48}}%
\TextAt{21}{231.849}{295.278}{i}%
\TextAt{16}{238.824}{298.233}{\Character{106}}%
\TextAt{13}{245.466}{298.233}{k}%
\TextAt{8}{255.509}{298.233}{the}%
\TextAt{8}{275.323}{298.233}{composite}%
\TextAt{8}{329.997}{298.233}{of}%
\TextAt{8}{342.984}{298.233}{the}%
\TextAt{8}{362.798}{298.233}{extensions}%
\TextAt{13}{418.835}{298.233}{K}%
\TextAt{21}{428.794}{296.44}{j}%
\TextAt{16}{436.497}{298.233}{\Character{106}}%
\TextAt{13}{443.139}{298.233}{k}%
\TextAt{8}{453.183}{298.233}{such}%
\TextAt{8}{479.237}{298.233}{that}%
\TextAt{13}{504.262}{298.233}{j}%
\NativeMathAt{513.099}{298.233}{11.955}{\neq}%
\TextAt{13}{525.524}{298.233}{i}%
\TextAt{8}{529.518}{298.233}{.}%
\TextAt{8}{79.226}{281.187}{If}%
\TextAt{13}{90.918}{281.187}{K}%
\TextAt{3}{101.736}{285.526}{\Character{48}}%
\TextAt{21}{100.877}{278.232}{i}%
\TextAt{16}{107.188}{281.187}{\Character{92}}%
\TextAt{13}{117.815}{281.187}{K}%
\TextAt{21}{127.774}{279.394}{i}%
\TextAt{9}{134.476}{281.187}{\Character{61}}%
\TextAt{13}{146.901}{281.187}{k}%
\TextAt{8}{157.288}{281.187}{for}%
\TextAt{8}{175.173}{281.187}{all}%
\TextAt{13}{191.432}{281.187}{i}%
\TextAt{8}{195.426}{281.187}{,}%
\TextAt{8}{202.575}{281.187}{then}%
\TextAt{8}{229.235}{281.187}{one}%
\TextAt{8}{250.694}{281.187}{has}%
\TextAt{8}{271.568}{281.187}{a}%
\TextAt{8}{281.323}{281.187}{topological}%
\TextAt{8}{340.848}{281.187}{isomorphism}%
\TextAt{13}{240.898}{255.384}{G}%
\TextAt{9}{250.132}{255.384}{\Character{40}}%
\TextAt{13}{254.684}{255.384}{K}%
\TextAt{16}{268.823}{255.384}{\Character{106}}%
\TextAt{13}{275.464}{255.384}{k}%
\TextAt{9}{281.954}{255.384}{\Character{41}}%
\NativeMathAt{289.827}{255.384}{11.955}{\cong}%
\TextAt{20}{302.446}{266.742}{\Character{89}}%
\TextAt{21}{308.643}{241.547}{i}%
\TextAt{13}{319.715}{255.384}{G}%
\TextAt{9}{328.949}{255.384}{\Character{40}}%
\TextAt{13}{333.501}{255.384}{K}%
\TextAt{21}{343.46}{253.591}{i}%
\TextAt{16}{350.162}{255.384}{\Character{106}}%
\TextAt{13}{356.804}{255.384}{k}%
\TextAt{9}{363.293}{255.384}{\Character{41}}%
\TextAt{13}{367.846}{255.384}{\Character{58}}%
\TextAt{12}{79.226}{208.945}{Exercise}%
\TextAt{12}{133.304}{208.945}{IV.1.3.}%
\TextAt{8}{177.823}{208.945}{A}%
\TextAt{8}{191.257}{208.945}{compact}%
\TextAt{8}{238.849}{208.945}{Hausdorff}%
\TextAt{8}{293.003}{208.945}{group}%
\TextAt{8}{326.942}{208.945}{is}%
\TextAt{8}{339.468}{208.945}{totally}%
\TextAt{8}{377.639}{208.945}{disconnected}%
\TextAt{8}{447.395}{208.945}{if}%
\TextAt{8}{458.882}{208.945}{and}%
\TextAt{8}{482.407}{208.945}{only}%
\TextAt{8}{508.861}{208.945}{if}%
\TextAt{8}{520.348}{208.945}{its}%
\TextAt{8}{79.226}{191.899}{neutral}%
\TextAt{8}{119.551}{191.899}{element}%
\TextAt{8}{162.793}{191.899}{admits}%
\TextAt{8}{201.229}{191.899}{a}%
\TextAt{8}{210.984}{191.899}{basis}%
\TextAt{8}{239.724}{191.899}{of}%
\TextAt{8}{253.054}{191.899}{neighbourhoods}%
\TextAt{8}{337.673}{191.899}{consisting}%
\TextAt{8}{391.782}{191.899}{only}%
\TextAt{8}{417.474}{191.899}{of}%
\TextAt{8}{430.804}{191.899}{normal}%
\TextAt{8}{470.483}{191.899}{subgroups.}%
\TextAt{12}{79.226}{161.737}{Exercise}%
\TextAt{12}{131.11}{161.737}{IV.1.4.}%
\TextAt{8}{173.721}{161.737}{Every}%
\TextAt{8}{206.23}{161.737}{quotient}%
\TextAt{13}{251.269}{161.737}{G}%
\TextAt{9}{260.502}{161.737}{\Character{47}}%
\TextAt{13}{266.355}{161.737}{H}%
\TextAt{8}{279.731}{161.737}{of}%
\TextAt{8}{291.916}{161.737}{a}%
\TextAt{8}{300.527}{161.737}{profinite}%
\TextAt{8}{346.21}{161.737}{group}%
\TextAt{13}{378.241}{161.737}{G}%
\TextAt{8}{390.228}{161.737}{by}%
\TextAt{8}{405.342}{161.737}{a}%
\TextAt{8}{413.953}{161.737}{closed}%
\TextAt{8}{447.335}{161.737}{normal}%
\TextAt{8}{485.869}{161.737}{subgroup}%
\TextAt{13}{79.226}{144.691}{H}%
\TextAt{8}{93.746}{144.691}{is}%
\TextAt{8}{105.51}{144.691}{a}%
\TextAt{8}{115.265}{144.691}{profinite}%
\TextAt{8}{162.094}{144.691}{group.}%
\TextAt{12}{79.226}{114.528}{Exercise}%
\TextAt{12}{131.758}{114.528}{IV.1.5.}%
\TextAt{8}{174.932}{114.528}{Let}%
\TextAt{13}{195.32}{114.528}{G}%
\TextAt{3}{204.554}{118.866}{\Character{48}}%
\TextAt{8}{210.665}{114.528}{be}%
\TextAt{8}{226.008}{114.528}{the}%
\TextAt{8}{245.583}{114.528}{closure}%
\TextAt{8}{284.083}{114.528}{of}%
\TextAt{8}{296.832}{114.528}{the}%
\TextAt{8}{316.407}{114.528}{commutator}%
\TextAt{8}{381.854}{114.528}{subgroup}%
\TextAt{8}{432.07}{114.528}{of}%
\TextAt{8}{444.818}{114.528}{a}%
\TextAt{8}{453.992}{114.528}{profinite}%
\TextAt{8}{500.239}{114.528}{group,}%
\TextAt{8}{79.226}{97.482}{and}%
\TextAt{13}{102.404}{97.482}{G}%
\TextAt{7}{111.638}{101.821}{ab}%
\TextAt{9}{125.099}{97.482}{\Character{61}}%
\TextAt{13}{138.232}{97.482}{G}%
\TextAt{9}{147.466}{97.482}{\Character{47}}%
\TextAt{13}{153.319}{97.482}{G}%
\TextAt{3}{162.553}{101.821}{\Character{48}}%
\TextAt{8}{165.348}{97.482}{.}%
\TextAt{8}{175.042}{97.482}{Show}%
\TextAt{8}{206.35}{97.482}{that}%
\TextAt{8}{232.134}{97.482}{every}%
\TextAt{8}{263.441}{97.482}{continuous}%
\TextAt{8}{322.461}{97.482}{homomorphism}%
\TextAt{13}{405.546}{97.482}{G}%
\NativeArrowAt{418.809}{97.482}{11.955}{19.26}{\longrightarrow}%
\TextAt{13}{442.098}{97.482}{A}%
\TextAt{8}{455.187}{97.482}{into}%
\TextAt{8}{479.345}{97.482}{an}%
\TextAt{8}{496.019}{97.482}{abelian}%
\TextAt{8}{79.226}{80.436}{profinite}%
\TextAt{8}{126.055}{80.436}{group}%
\TextAt{8}{159.23}{80.436}{factorizes}%
\TextAt{8}{210.996}{80.436}{through}%
\TextAt{13}{255.23}{80.436}{G}%
\TextAt{7}{264.464}{84.775}{ab}%
\TextAt{8}{273.896}{80.436}{.}%
\TextAt{8}{297.211}{43.033}{242}%
\Destination{78.23}{809.815}{nav.page.242}%
\Destination{79.226}{445.155}{nav.ex.4.1.1}%
\Destination{79.226}{329.255}{nav.ex.4.1.2}%
\Destination{79.226}{222.955}{nav.ex.4.1.3}%
\Destination{79.226}{175.755}{nav.ex.4.1.4}%
\Destination{79.226}{128.555}{nav.ex.4.1.5}%
\Destination{79.36}{682.855}{nav.print.265}%
\Destination{79.36}{698.855}{nav.number.4.1.3}%
}

% Page 255 (body page 243)
\Page{%
\TextAt{2}{79.226}{796.943}{IV.}%
\TextAt{2}{98.012}{796.943}{Abstract}%
\TextAt{2}{143.252}{796.943}{Class}%
\TextAt{2}{171.844}{796.943}{Field}%
\TextAt{2}{199.575}{796.943}{Theory}%
\Vector{q 0 G 0 g 0.299 w 7.226 18.729 m 460.769 18.729 l S Q}%
\TextAt{15}{79.226}{760.804}{§}%
\TextAt{15}{93.785}{760.804}{2}%
\TextAt{15}{120.693}{760.804}{Projective}%
\TextAt{15}{213.346}{760.804}{and}%
\TextAt{15}{250.738}{760.804}{Inductive}%
\TextAt{15}{336.747}{760.804}{Limits}%
\TextAt{8}{79.226}{728.918}{The}%
\TextAt{8}{103.526}{728.918}{notions}%
\TextAt{8}{144.813}{728.918}{of}%
\TextAt{8}{158.389}{728.918}{projective,}%
\TextAt{8}{215.926}{728.918}{resp.}%
\TextAt{8}{245.985}{728.918}{inductive}%
\TextAt{8}{296.957}{728.918}{limit}%
\TextAt{8}{325.165}{728.918}{generalize}%
\TextAt{8}{379.389}{728.918}{the}%
\TextAt{8}{399.791}{728.918}{operations}%
\TextAt{8}{457.015}{728.918}{of}%
\TextAt{8}{470.591}{728.918}{intersection,}%
\TextAt{8}{79.226}{711.872}{resp.}%
\TextAt{8}{112.582}{711.872}{union.}%
\TextAt{8}{153.684}{711.872}{If}%
\TextAt{9}{166.721}{711.872}{\Character{40}}%
\TextAt{13}{171.273}{711.872}{X}%
\TextAt{21}{180.988}{710.079}{i}%
\TextAt{9}{184.369}{711.872}{\Character{41}}%
\TextAt{21}{188.922}{710.079}{i}%
\NativeMathAt{191.805}{710.079}{7.97}{\in}%
\TextAt{21}{197.45}{710.079}{I}%
\TextAt{8}{207.484}{711.872}{is}%
\TextAt{8}{220.592}{711.872}{a}%
\TextAt{8}{231.692}{711.872}{family}%
\TextAt{8}{268.807}{711.872}{of}%
\TextAt{8}{283.481}{711.872}{subsets}%
\TextAt{8}{325.33}{711.872}{of}%
\TextAt{8}{340.004}{711.872}{a}%
\TextAt{8}{351.104}{711.872}{topological}%
\TextAt{8}{411.973}{711.872}{space}%
\TextAt{13}{444.593}{711.872}{X}%
\TextAt{8}{460.49}{711.872}{which}%
\TextAt{8}{495.32}{711.872}{for}%
\TextAt{8}{514.55}{711.872}{any}%
\TextAt{8}{79.226}{694.826}{two}%
\TextAt{8}{101.592}{694.826}{sets}%
\TextAt{13}{124.724}{694.826}{X}%
\TextAt{21}{134.439}{693.033}{i}%
\TextAt{13}{137.82}{694.826}{\Character{59}}%
\TextAt{13}{143.044}{694.826}{X}%
\TextAt{21}{152.779}{693.033}{j}%
\TextAt{8}{161.308}{694.826}{also}%
\TextAt{8}{185.037}{694.826}{contains}%
\TextAt{8}{231.206}{694.826}{the}%
\TextAt{8}{251.611}{694.826}{set}%
\TextAt{13}{270.128}{694.826}{X}%
\TextAt{21}{279.843}{693.033}{i}%
\TextAt{16}{286.051}{694.826}{\Character{92}}%
\TextAt{13}{296.848}{694.826}{X}%
\TextAt{21}{306.563}{693.033}{j}%
\TextAt{8}{315.091}{694.826}{(resp.}%
\TextAt{13}{349.716}{694.826}{X}%
\TextAt{21}{359.431}{693.033}{i}%
\TextAt{16}{365.638}{694.826}{\Character{91}}%
\TextAt{13}{376.435}{694.826}{X}%
\TextAt{21}{386.15}{693.033}{j}%
\TextAt{8}{390.532}{694.826}{),}%
\TextAt{8}{402.548}{694.826}{then}%
\TextAt{8}{429.457}{694.826}{the}%
\TextAt{8}{449.862}{694.826}{projective}%
\TextAt{8}{504.089}{694.826}{(resp.}%
\TextAt{8}{79.226}{677.78}{inductive)}%
\TextAt{8}{134.507}{677.78}{limit}%
\TextAt{8}{162.47}{677.78}{of}%
\TextAt{8}{175.8}{677.78}{this}%
\TextAt{8}{198.622}{677.78}{family}%
\TextAt{8}{234.392}{677.78}{is}%
\TextAt{8}{246.156}{677.78}{simply}%
\TextAt{8}{283.612}{677.78}{defined}%
\TextAt{8}{323.925}{677.78}{by}%
\NativeMathAt{189.836}{643.394}{11.955}{\varprojlim}%
\TextAt{21}{191.554}{631.014}{i}%
\NativeMathAt{194.437}{631.014}{7.97}{\in}%
\TextAt{21}{200.083}{631.014}{I}%
\TextAt{13}{208.087}{643.394}{X}%
\TextAt{21}{217.802}{641.6}{i}%
\TextAt{9}{224.504}{643.394}{\Character{61}}%
\TextAt{20}{236.93}{654.751}{\Character{92}}%
\TextAt{21}{237.161}{629.38}{i}%
\NativeMathAt{240.044}{629.38}{7.97}{\in}%
\TextAt{21}{245.689}{629.38}{I}%
\TextAt{13}{252.206}{643.394}{X}%
\TextAt{21}{261.921}{641.6}{i}%
\TextAt{20}{291.205}{663.837}{\Character{32}}%
\TextAt{8}{300.67}{643.394}{resp.}%
\NativeMathAt{331.984}{643.394}{11.955}{\varinjlim}%
\TextAt{21}{333.703}{631.014}{i}%
\NativeMathAt{336.586}{631.014}{7.97}{\in}%
\TextAt{21}{342.231}{631.014}{I}%
\TextAt{13}{350.236}{643.394}{X}%
\TextAt{21}{359.951}{641.6}{i}%
\TextAt{9}{366.653}{643.394}{\Character{61}}%
\TextAt{20}{379.079}{654.751}{\Character{91}}%
\TextAt{21}{379.309}{629.38}{i}%
\NativeMathAt{382.193}{629.38}{7.97}{\in}%
\TextAt{21}{387.838}{629.38}{I}%
\TextAt{13}{394.355}{643.394}{X}%
\TextAt{21}{404.07}{641.6}{i}%
\TextAt{20}{407.451}{663.837}{\Character{33}}%
\TextAt{13}{418.908}{643.394}{\Character{58}}%
\TextAt{8}{79.226}{608.25}{Writing}%
\TextAt{13}{123.181}{608.25}{i}%
\NativeMathAt{132.261}{608.25}{11.955}{\leq}%
\TextAt{13}{146.646}{608.25}{j}%
\TextAt{8}{157.095}{608.25}{if}%
\TextAt{13}{168.855}{608.25}{X}%
\TextAt{21}{178.571}{606.457}{j}%
\NativeMathAt{188.039}{608.25}{11.955}{\subseteq}%
\TextAt{13}{202.424}{608.25}{X}%
\TextAt{21}{212.14}{606.457}{i}%
\TextAt{8}{220.454}{608.25}{(resp.}%
\TextAt{13}{257.44}{608.25}{X}%
\TextAt{21}{267.155}{606.457}{i}%
\NativeMathAt{275.623}{608.25}{11.955}{\subseteq}%
\TextAt{13}{290.008}{608.25}{X}%
\TextAt{21}{299.723}{606.457}{j}%
\TextAt{8}{304.105}{608.25}{)}%
\TextAt{8}{313.594}{608.25}{makes}%
\TextAt{8}{349.814}{608.25}{the}%
\TextAt{8}{371.006}{608.25}{indexing}%
\TextAt{8}{419.194}{608.25}{set}%
\TextAt{13}{438.497}{608.25}{I}%
\TextAt{8}{449.534}{608.25}{into}%
\TextAt{8}{474.313}{608.25}{a}%
\TextAt{12}{485.104}{608.25}{directed}%
\TextAt{12}{79.226}{591.204}{system}%
\TextAt{8}{119.527}{591.204}{,}%
\TextAt{8}{126.763}{591.204}{i.e.,}%
\TextAt{8}{148.956}{591.204}{an}%
\TextAt{8}{165.285}{591.204}{ordered}%
\TextAt{8}{207.628}{591.204}{set}%
\TextAt{8}{225.965}{591.204}{in}%
\TextAt{8}{239.688}{591.204}{which,}%
\TextAt{8}{276.513}{591.204}{for}%
\TextAt{8}{294.468}{591.204}{every}%
\TextAt{8}{325.43}{591.204}{pair}%
\TextAt{13}{349.565}{591.204}{i\Character{59}}%
\TextAt{13}{358.783}{591.204}{j}%
\TextAt{8}{364.319}{591.204}{,}%
\TextAt{8}{371.555}{591.204}{there}%
\TextAt{8}{401.537}{591.204}{exists}%
\TextAt{8}{433.921}{591.204}{a}%
\TextAt{13}{443.746}{591.204}{k}%
\TextAt{8}{454.203}{591.204}{such}%
\TextAt{8}{480.67}{591.204}{that}%
\TextAt{13}{506.109}{591.204}{i}%
\NativeMathAt{513.542}{591.204}{11.955}{\leq}%
\TextAt{13}{526.28}{591.204}{k}%
\TextAt{8}{79.226}{574.158}{and}%
\TextAt{13}{101.766}{574.158}{j}%
\NativeMathAt{110.602}{574.158}{11.955}{\leq}%
\TextAt{13}{123.222}{574.158}{k}%
\TextAt{8}{129.711}{574.158}{.}%
\TextAt{8}{138.085}{574.158}{In}%
\TextAt{8}{152.484}{574.158}{the}%
\TextAt{8}{172.418}{574.158}{case}%
\TextAt{8}{196.966}{574.158}{at}%
\TextAt{8}{211.054}{574.158}{hand,}%
\TextAt{8}{243.394}{574.158}{such}%
\TextAt{8}{269.568}{574.158}{a}%
\TextAt{13}{279.101}{574.158}{k}%
\TextAt{8}{289.265}{574.158}{is}%
\TextAt{8}{300.806}{574.158}{given}%
\TextAt{8}{331.153}{574.158}{by}%
\TextAt{13}{347.189}{574.158}{X}%
\TextAt{21}{356.905}{572.365}{k}%
\TextAt{9}{365.345}{574.158}{\Character{61}}%
\TextAt{13}{377.771}{574.158}{X}%
\TextAt{21}{387.486}{572.365}{i}%
\TextAt{16}{393.069}{574.158}{\Character{92}}%
\TextAt{13}{403.24}{574.158}{X}%
\TextAt{21}{412.955}{572.365}{j}%
\TextAt{8}{421.012}{574.158}{(resp.}%
\TextAt{13}{454.815}{574.158}{X}%
\TextAt{21}{464.53}{572.365}{k}%
\TextAt{9}{472.971}{574.158}{\Character{61}}%
\TextAt{13}{485.396}{574.158}{X}%
\TextAt{21}{495.111}{572.365}{i}%
\TextAt{16}{500.694}{574.158}{\Character{91}}%
\TextAt{13}{510.865}{574.158}{X}%
\TextAt{21}{520.581}{572.365}{j}%
\TextAt{8}{524.963}{574.158}{).}%
\TextAt{8}{79.226}{557.112}{For}%
\TextAt{13}{100.268}{557.112}{i}%
\NativeMathAt{107.706}{557.112}{11.955}{\leq}%
\TextAt{13}{120.449}{557.112}{j}%
\TextAt{8}{129.935}{557.112}{we}%
\TextAt{8}{147.235}{557.112}{denote}%
\TextAt{8}{185.026}{557.112}{the}%
\TextAt{8}{205.255}{557.112}{inclusion}%
\TextAt{13}{254.164}{557.112}{X}%
\TextAt{21}{263.879}{555.319}{j}%
\NativeArrowAt{271.706}{557.112}{11.955}{13.214}{\hookrightarrow}%
\TextAt{13}{288.365}{557.112}{X}%
\TextAt{21}{298.08}{555.319}{i}%
\TextAt{8}{305.431}{557.112}{(resp.}%
\TextAt{13}{339.526}{557.112}{X}%
\TextAt{21}{349.241}{555.319}{i}%
\NativeArrowAt{356.067}{557.112}{11.955}{13.214}{\hookrightarrow}%
\TextAt{13}{372.726}{557.112}{X}%
\TextAt{21}{382.441}{555.319}{j}%
\TextAt{8}{386.823}{557.112}{)}%
\TextAt{8}{395.348}{557.112}{by}%
\TextAt{13}{411.68}{557.112}{f}%
\TextAt{21}{417.45}{555.319}{ij}%
\TextAt{8}{428.685}{557.112}{and}%
\TextAt{8}{451.52}{557.112}{obtain}%
\TextAt{8}{488.02}{557.112}{a}%
\TextAt{8}{497.848}{557.112}{system}%
\TextAt{16}{79.226}{540.066}{\Character{102}}%
\TextAt{13}{85.204}{540.066}{X}%
\TextAt{21}{94.919}{538.273}{i}%
\TextAt{13}{98.3}{540.066}{\Character{59}}%
\TextAt{13}{103.524}{540.066}{f}%
\TextAt{21}{109.314}{538.273}{ij}%
\TextAt{16}{116.58}{540.066}{\Character{103}}%
\TextAt{8}{126.391}{540.066}{of}%
\TextAt{8}{139.657}{540.066}{sets}%
\TextAt{8}{162.475}{540.066}{and}%
\TextAt{8}{185.173}{540.066}{maps.}%
\TextAt{8}{220.332}{540.066}{The}%
\TextAt{8}{244.322}{540.066}{operations}%
\TextAt{8}{301.236}{540.066}{of}%
\TextAt{8}{314.502}{540.066}{intersection}%
\TextAt{8}{377.262}{540.066}{and}%
\TextAt{8}{399.961}{540.066}{union}%
\TextAt{8}{432.415}{540.066}{are}%
\TextAt{8}{451.861}{540.066}{now}%
\TextAt{8}{476.186}{540.066}{generalized}%
\TextAt{8}{79.226}{523.02}{by}%
\TextAt{8}{95.485}{523.02}{replacing}%
\TextAt{8}{145.565}{523.02}{the}%
\TextAt{8}{165.722}{523.02}{inclusions}%
\TextAt{13}{219.173}{523.02}{f}%
\TextAt{21}{224.943}{521.227}{ij}%
\TextAt{8}{236.106}{523.02}{with}%
\TextAt{8}{262.766}{523.02}{arbitrary}%
\TextAt{8}{312.535}{523.02}{maps.}%
\TextAt{12}{79.226}{497.719}{Definition}%
\TextAt{12}{142.328}{497.719}{(2.1).}%
\TextAt{8}{177.916}{497.719}{Let}%
\TextAt{13}{199.158}{497.719}{I}%
\TextAt{8}{209.431}{497.719}{be}%
\TextAt{8}{225.627}{497.719}{a}%
\TextAt{8}{235.655}{497.719}{directed}%
\TextAt{8}{280.796}{497.719}{system.}%
\TextAt{8}{324.982}{497.719}{A}%
\TextAt{12}{337.927}{497.719}{projective}%
\TextAt{8}{396.065}{497.719}{,}%
\TextAt{8}{403.555}{497.719}{resp.}%
\TextAt{12}{433.694}{497.719}{inductive}%
\TextAt{12}{492.469}{497.719}{system}%
\TextAt{8}{79.226}{480.673}{over}%
\TextAt{13}{104.272}{480.673}{I}%
\TextAt{8}{114.272}{480.673}{is}%
\TextAt{8}{126.036}{480.673}{a}%
\TextAt{8}{135.792}{480.673}{family}%
\TextAt{16}{171.562}{480.673}{\Character{102}}%
\TextAt{13}{177.539}{480.673}{X}%
\TextAt{21}{187.254}{478.879}{i}%
\TextAt{13}{190.636}{480.673}{\Character{59}}%
\TextAt{13}{195.86}{480.673}{f}%
\TextAt{21}{201.65}{478.879}{ij}%
\TextAt{16}{212.236}{480.673}{\Character{106}}%
\TextAt{13}{218.878}{480.673}{i\Character{59}}%
\TextAt{13}{228.096}{480.673}{j}%
\NativeMathAt{236.952}{480.673}{11.955}{\in}%
\TextAt{13}{248.243}{480.673}{I}%
\TextAt{13}{254.328}{480.673}{\Character{59}}%
\TextAt{13}{263.462}{480.673}{i}%
\NativeMathAt{270.801}{480.673}{11.955}{\leq}%
\TextAt{13}{283.42}{480.673}{j}%
\TextAt{16}{288.936}{480.673}{\Character{103}}%
\TextAt{8}{298.811}{480.673}{of}%
\TextAt{8}{312.141}{480.673}{topological}%
\TextAt{8}{371.666}{480.673}{spaces}%
\TextAt{13}{407.555}{480.673}{X}%
\TextAt{21}{417.27}{478.879}{i}%
\TextAt{8}{424.549}{480.673}{and}%
\TextAt{8}{447.312}{480.673}{continuous}%
\TextAt{8}{505.916}{480.673}{maps}%
\TextAt{13}{199.108}{454.159}{f}%
\TextAt{21}{204.878}{452.365}{ij}%
\TextAt{9}{215.465}{454.159}{\Character{58}}%
\TextAt{13}{222.037}{454.159}{X}%
\TextAt{21}{231.752}{452.365}{j}%
\NativeArrowAt{239.455}{454.159}{11.955}{19.26}{\longrightarrow}%
\TextAt{13}{262.037}{454.159}{X}%
\TextAt{21}{271.752}{452.365}{i}%
\TextAt{13}{275.134}{454.159}{\Character{59}}%
\TextAt{8}{304.288}{454.159}{resp.}%
\TextAt{13}{333.61}{454.159}{f}%
\TextAt{21}{339.38}{452.365}{ij}%
\TextAt{9}{349.966}{454.159}{\Character{58}}%
\TextAt{13}{356.539}{454.159}{X}%
\TextAt{21}{366.254}{452.365}{i}%
\NativeArrowAt{372.956}{454.159}{11.955}{19.26}{\longrightarrow}%
\TextAt{13}{395.538}{454.159}{X}%
\TextAt{21}{405.253}{452.365}{j}%
\TextAt{13}{409.636}{454.159}{\Character{59}}%
\TextAt{8}{79.226}{427.645}{such}%
\TextAt{8}{105.623}{427.645}{that}%
\TextAt{8}{130.992}{427.645}{one}%
\TextAt{8}{152.452}{427.645}{has}%
\TextAt{13}{173.325}{427.645}{f}%
\TextAt{21}{179.095}{425.851}{ii}%
\TextAt{9}{188.681}{427.645}{\Character{61}}%
\TextAt{8}{201.106}{427.645}{id}%
\TextAt{21}{210.861}{425.851}{X}%
\TextAt{22}{217.826}{424.636}{i}%
\TextAt{8}{225.384}{427.645}{and}%
\TextAt{13}{206.515}{401.131}{f}%
\TextAt{21}{212.285}{399.338}{ik}%
\TextAt{9}{223.608}{401.131}{\Character{61}}%
\TextAt{13}{236.034}{401.131}{f}%
\TextAt{21}{241.804}{399.338}{ij}%
\TextAt{16}{251.726}{401.131}{\Character{14}}%
\TextAt{13}{260.36}{401.131}{f}%
\TextAt{21}{266.13}{399.338}{j}%
\TextAt{21}{270.004}{399.338}{k}%
\TextAt{13}{275.134}{401.131}{\Character{59}}%
\TextAt{8}{304.288}{401.131}{resp.}%
\TextAt{13}{333.61}{401.131}{f}%
\TextAt{21}{339.38}{399.338}{ik}%
\TextAt{9}{350.704}{401.131}{\Character{61}}%
\TextAt{13}{363.129}{401.131}{f}%
\TextAt{21}{368.899}{399.338}{j}%
\TextAt{21}{372.773}{399.338}{k}%
\TextAt{16}{380.56}{401.131}{\Character{14}}%
\TextAt{13}{389.194}{401.131}{f}%
\TextAt{21}{394.964}{399.338}{ij}%
\TextAt{13}{402.229}{401.131}{\Character{59}}%
\TextAt{8}{79.226}{374.617}{when}%
\TextAt{13}{109.784}{374.617}{i}%
\NativeMathAt{117.098}{374.617}{11.955}{\leq}%
\TextAt{13}{129.717}{374.617}{j}%
\NativeMathAt{138.553}{374.617}{11.955}{\leq}%
\TextAt{13}{151.173}{374.617}{k}%
\TextAt{8}{157.662}{374.617}{.}%
\TextAt{8}{97.159}{349.315}{In}%
\TextAt{8}{111.805}{349.315}{order}%
\TextAt{8}{142.4}{349.315}{to}%
\TextAt{8}{156.735}{349.315}{define}%
\TextAt{8}{190.569}{349.315}{the}%
\TextAt{8}{210.751}{349.315}{projective,}%
\TextAt{8}{268.011}{349.315}{resp.}%
\TextAt{8}{297.409}{349.315}{inductive}%
\TextAt{8}{348.159}{349.315}{limit}%
\TextAt{8}{376.148}{349.315}{of}%
\TextAt{8}{389.503}{349.315}{a}%
\TextAt{8}{399.283}{349.315}{projective,}%
\TextAt{8}{456.544}{349.315}{resp.}%
\TextAt{8}{485.941}{349.315}{inductive}%
\TextAt{8}{79.226}{332.269}{system}%
\TextAt{16}{119.273}{332.269}{\Character{102}}%
\TextAt{13}{125.251}{332.269}{X}%
\TextAt{21}{134.966}{330.476}{i}%
\TextAt{13}{138.347}{332.269}{\Character{59}}%
\TextAt{13}{143.571}{332.269}{f}%
\TextAt{21}{149.362}{330.476}{ij}%
\TextAt{16}{156.627}{332.269}{\Character{103}}%
\TextAt{8}{162.605}{332.269}{,}%
\TextAt{8}{171.29}{332.269}{we}%
\TextAt{8}{189.746}{332.269}{make}%
\TextAt{8}{221.544}{332.269}{use}%
\TextAt{8}{242.989}{332.269}{of}%
\TextAt{8}{257.548}{332.269}{the}%
\TextAt{8}{278.933}{332.269}{direct}%
\TextAt{8}{313.325}{332.269}{product}%
\TextAt{20}{358.453}{341.236}{\Character{81}}%
\TextAt{21}{369.744}{328.782}{i}%
\NativeMathAt{372.627}{328.782}{7.97}{\in}%
\TextAt{21}{378.273}{328.782}{I}%
\TextAt{13}{385.057}{332.269}{X}%
\TextAt{21}{394.773}{330.476}{i}%
\TextAt{8}{398.154}{332.269}{,}%
\TextAt{8}{406.839}{332.269}{resp.}%
\TextAt{8}{439.847}{332.269}{the}%
\TextAt{8}{461.232}{332.269}{disjoint}%
\TextAt{8}{504.149}{332.269}{union}%
\TextAt{20}{79.226}{324.19}{\Character{96}}%
\TextAt{21}{90.517}{311.736}{i}%
\NativeMathAt{93.4}{311.736}{7.97}{\in}%
\TextAt{21}{99.046}{311.736}{I}%
\TextAt{13}{105.83}{315.223}{X}%
\TextAt{21}{115.545}{313.43}{i}%
\TextAt{8}{118.927}{315.223}{.}%
\TextAt{12}{79.226}{289.921}{Definition}%
\TextAt{12}{142.015}{289.921}{(2.2).}%
\TextAt{8}{177.33}{289.921}{The}%
\TextAt{12}{201.384}{289.921}{projective}%
\TextAt{12}{264.005}{289.921}{limit}%
\TextAt{13}{277.123}{263.407}{X}%
\TextAt{9}{291.099}{263.407}{\Character{61}}%
\NativeMathAt{303.525}{263.407}{11.955}{\varprojlim}%
\TextAt{21}{305.243}{251.027}{i}%
\NativeMathAt{308.126}{251.027}{7.97}{\in}%
\TextAt{21}{313.772}{251.027}{I}%
\TextAt{13}{321.776}{263.407}{X}%
\TextAt{21}{331.491}{261.614}{i}%
\TextAt{8}{79.226}{229.898}{of}%
\TextAt{8}{92.556}{229.898}{the}%
\TextAt{8}{112.713}{229.898}{projective}%
\TextAt{8}{166.69}{229.898}{system}%
\TextAt{16}{205.509}{229.898}{\Character{102}}%
\TextAt{13}{211.486}{229.898}{X}%
\TextAt{21}{221.201}{228.105}{i}%
\TextAt{13}{224.583}{229.898}{\Character{59}}%
\TextAt{13}{229.807}{229.898}{f}%
\TextAt{21}{235.597}{228.105}{ij}%
\TextAt{16}{242.862}{229.898}{\Character{103}}%
\TextAt{8}{252.737}{229.898}{is}%
\TextAt{8}{264.501}{229.898}{defined}%
\TextAt{8}{304.814}{229.898}{to}%
\TextAt{8}{319.124}{229.898}{be}%
\TextAt{8}{335.049}{229.898}{the}%
\TextAt{8}{355.205}{229.898}{subset}%
\TextAt{13}{186.043}{195.169}{X}%
\TextAt{9}{200.019}{195.169}{\Character{61}}%
\TextAt{20}{212.444}{215.613}{\Character{40}}%
\TextAt{9}{222.075}{195.169}{\Character{40}}%
\TextAt{13}{226.627}{195.169}{x}%
\TextAt{21}{233.279}{193.376}{i}%
\TextAt{9}{236.661}{195.169}{\Character{41}}%
\TextAt{21}{241.213}{193.376}{i}%
\NativeMathAt{244.096}{193.376}{7.97}{\in}%
\TextAt{21}{249.742}{193.376}{I}%
\NativeMathAt{257.854}{195.169}{11.955}{\in}%
\TextAt{20}{269.145}{206.527}{\Character{89}}%
\TextAt{21}{270.372}{181.155}{i}%
\NativeMathAt{273.255}{181.155}{7.97}{\in}%
\TextAt{21}{278.901}{181.155}{I}%
\TextAt{13}{286.414}{195.169}{X}%
\TextAt{21}{296.129}{193.376}{i}%
\TextAt{20}{302.831}{216.091}{\Character{12}}%
\TextAt{20}{302.831}{208.918}{\Character{12}}%
\TextAt{20}{302.831}{201.744}{\Character{12}}%
\TextAt{20}{302.831}{194.571}{\Character{12}}%
\TextAt{20}{302.831}{187.398}{\Character{12}}%
\TextAt{13}{310.137}{195.169}{f}%
\TextAt{21}{315.907}{193.376}{ij}%
\TextAt{9}{323.172}{195.169}{\Character{40}}%
\TextAt{13}{327.725}{195.169}{x}%
\TextAt{21}{334.377}{193.376}{j}%
\TextAt{9}{338.759}{195.169}{\Character{41}}%
\TextAt{9}{346.614}{195.169}{\Character{61}}%
\TextAt{13}{359.057}{195.169}{x}%
\TextAt{21}{365.71}{193.376}{i}%
\TextAt{8}{372.988}{195.169}{for}%
\TextAt{13}{390.873}{195.169}{i}%
\NativeMathAt{398.187}{195.169}{11.955}{\leq}%
\TextAt{13}{410.807}{195.169}{j}%
\TextAt{20}{416.322}{215.613}{\Character{41}}%
\TextAt{8}{79.226}{160.026}{of}%
\TextAt{8}{92.556}{160.026}{the}%
\TextAt{8}{112.713}{160.026}{product}%
\TextAt{20}{156.612}{168.992}{\Character{81}}%
\TextAt{21}{167.903}{156.539}{i}%
\NativeMathAt{170.786}{156.539}{7.97}{\in}%
\TextAt{21}{176.432}{156.539}{I}%
\TextAt{13}{183.216}{160.026}{X}%
\TextAt{21}{192.931}{158.232}{i}%
\TextAt{8}{196.313}{160.026}{.}%
\TextAt{8}{97.159}{134.724}{The}%
\TextAt{8}{121.107}{134.724}{product}%
\TextAt{20}{164.901}{143.69}{\Character{81}}%
\TextAt{21}{176.192}{131.237}{i}%
\NativeMathAt{179.075}{131.237}{7.97}{\in}%
\TextAt{21}{184.721}{131.237}{I}%
\TextAt{13}{191.505}{134.724}{X}%
\TextAt{21}{201.22}{132.931}{i}%
\TextAt{8}{208.393}{134.724}{is}%
\TextAt{8}{220.052}{134.724}{equipped}%
\TextAt{8}{270.014}{134.724}{with}%
\TextAt{8}{296.569}{134.724}{the}%
\TextAt{8}{316.62}{134.724}{product}%
\TextAt{8}{360.413}{134.724}{topology.}%
\TextAt{8}{412.092}{134.724}{If}%
\TextAt{8}{423.679}{134.724}{the}%
\TextAt{13}{443.73}{134.724}{X}%
\TextAt{21}{453.445}{132.931}{i}%
\TextAt{8}{460.618}{134.724}{are}%
\TextAt{8}{480.023}{134.724}{Hausdorff,}%
\TextAt{8}{79.226}{117.678}{then}%
\TextAt{8}{106.21}{117.678}{so}%
\TextAt{8}{120.904}{117.678}{is}%
\TextAt{8}{132.991}{117.678}{the}%
\TextAt{8}{153.472}{117.678}{product,}%
\TextAt{8}{201.028}{117.678}{and}%
\TextAt{8}{224.114}{117.678}{it}%
\TextAt{8}{236.142}{117.678}{contains}%
\TextAt{8}{282.385}{117.678}{in}%
\TextAt{8}{296.362}{117.678}{this}%
\TextAt{8}{319.508}{117.678}{case}%
\TextAt{13}{344.603}{117.678}{X}%
\TextAt{8}{359.479}{117.678}{as}%
\TextAt{8}{374.173}{117.678}{a}%
\TextAt{8}{384.252}{117.678}{closed}%
\TextAt{8}{419.103}{117.678}{subspace.}%
\TextAt{8}{473.521}{117.678}{Indeed,}%
\TextAt{8}{515.207}{117.678}{one}%
\TextAt{8}{79.226}{100.632}{has}%
\TextAt{13}{275.008}{83.586}{X}%
\TextAt{9}{288.984}{83.586}{\Character{61}}%
\TextAt{20}{301.444}{94.943}{\Character{92}}%
\TextAt{21}{301.409}{69.689}{i}%
\NativeMathAt{304.292}{69.689}{7.97}{\leq}%
\TextAt{21}{310.879}{69.689}{j}%
\TextAt{13}{316.756}{83.586}{X}%
\TextAt{21}{326.471}{81.793}{ij}%
\TextAt{13}{333.736}{83.586}{\Character{59}}%
\TextAt{8}{297.211}{43.033}{243}%
\Destination{78.23}{809.815}{nav.page.243}%
\Destination{79.23}{773.075}{nav.section.4.2}%
\Destination{79.36}{149.755}{nav.print.266}%
\Destination{79.36}{511.755}{nav.number.4.2.1}%
\Destination{79.36}{303.955}{nav.number.4.2.2}%
\Destination{79.23}{773.075}{outline.42}%
}

% Page 256 (body page 244)
\Page{%
\TextAt{2}{79.226}{796.943}{IV.}%
\TextAt{2}{98.012}{796.943}{Abstract}%
\TextAt{2}{143.252}{796.943}{Class}%
\TextAt{2}{171.844}{796.943}{Field}%
\TextAt{2}{199.575}{796.943}{Theory}%
\Vector{q 0 G 0 g 0.299 w 7.226 18.729 m 460.769 18.729 l S Q}%
\TextAt{8}{79.226}{761.124}{where}%
\TextAt{13}{113.426}{761.124}{X}%
\TextAt{21}{123.141}{759.33}{ij}%
\TextAt{9}{134.393}{761.124}{\Character{61}}%
\TextAt{16}{147.484}{761.124}{\Character{102}}%
\TextAt{9}{153.462}{761.124}{\Character{40}}%
\TextAt{13}{158.014}{761.124}{x}%
\TextAt{21}{164.666}{759.33}{k}%
\TextAt{9}{169.786}{761.124}{\Character{41}}%
\TextAt{21}{174.338}{759.33}{k}%
\NativeMathAt{178.96}{759.33}{7.97}{\in}%
\TextAt{21}{184.605}{759.33}{I}%
\NativeMathAt{193.384}{761.124}{11.955}{\in}%
\TextAt{20}{205.34}{770.09}{\Character{81}}%
\TextAt{21}{216.631}{757.637}{k}%
\TextAt{13}{223.744}{761.124}{X}%
\TextAt{21}{233.459}{759.33}{k}%
\TextAt{16}{242.565}{761.124}{\Character{106}}%
\TextAt{13}{249.872}{761.124}{f}%
\TextAt{21}{255.642}{759.33}{ij}%
\TextAt{9}{262.908}{761.124}{\Character{40}}%
\TextAt{13}{267.46}{761.124}{x}%
\TextAt{21}{274.112}{759.33}{j}%
\TextAt{9}{278.494}{761.124}{\Character{41}}%
\TextAt{9}{287.018}{761.124}{\Character{61}}%
\TextAt{13}{300.124}{761.124}{x}%
\TextAt{21}{306.776}{759.33}{i}%
\TextAt{16}{310.158}{761.124}{\Character{103}}%
\TextAt{8}{316.135}{761.124}{,}%
\TextAt{8}{323.773}{761.124}{so}%
\TextAt{8}{338.533}{761.124}{that}%
\TextAt{8}{364.293}{761.124}{it}%
\TextAt{8}{376.388}{761.124}{suffices}%
\TextAt{8}{416.565}{761.124}{to}%
\TextAt{8}{431.266}{761.124}{show}%
\TextAt{8}{460.66}{761.124}{the}%
\TextAt{8}{481.207}{761.124}{closedness}%
\TextAt{8}{79.226}{744.078}{of}%
\TextAt{8}{92.725}{744.078}{the}%
\TextAt{8}{113.05}{744.078}{sets}%
\TextAt{13}{136.101}{744.078}{X}%
\TextAt{21}{145.816}{742.284}{ij}%
\TextAt{8}{153.081}{744.078}{.}%
\TextAt{8}{162.036}{744.078}{Writing}%
\TextAt{13}{205.123}{744.078}{p}%
\TextAt{21}{210.999}{742.284}{i}%
\TextAt{9}{217.988}{744.078}{\Character{58}}%
\TextAt{20}{224.848}{753.044}{\Character{81}}%
\TextAt{21}{236.139}{740.591}{k}%
\NativeMathAt{240.761}{740.591}{7.97}{\in}%
\TextAt{21}{246.406}{740.591}{I}%
\TextAt{13}{253.19}{744.078}{X}%
\TextAt{21}{262.906}{742.284}{k}%
\NativeMathAt{271.634}{744.078}{11.955}{\rightarrow}%
\TextAt{13}{287.197}{744.078}{X}%
\TextAt{21}{296.912}{742.284}{i}%
\TextAt{8}{304.36}{744.078}{for}%
\TextAt{8}{322.413}{744.078}{the}%
\TextAt{13}{342.738}{744.078}{i}%
\TextAt{8}{346.732}{744.078}{-th}%
\TextAt{8}{365.754}{744.078}{projection,}%
\TextAt{8}{424.497}{744.078}{the}%
\TextAt{8}{444.822}{744.078}{two}%
\TextAt{8}{467.108}{744.078}{maps}%
\TextAt{13}{497.906}{744.078}{g}%
\TextAt{9}{507.548}{744.078}{\Character{61}}%
\TextAt{13}{520.261}{744.078}{p}%
\TextAt{21}{526.136}{742.284}{i}%
\TextAt{8}{529.518}{744.078}{,}%
\TextAt{13}{79.226}{727.032}{f}%
\TextAt{9}{89.593}{727.032}{\Character{61}}%
\TextAt{13}{102.019}{727.032}{f}%
\TextAt{21}{107.789}{725.238}{ij}%
\TextAt{16}{116.487}{727.032}{\Character{14}}%
\TextAt{13}{123.897}{727.032}{p}%
\TextAt{21}{129.772}{725.238}{j}%
\TextAt{9}{137.475}{727.032}{\Character{58}}%
\TextAt{20}{144.048}{735.998}{\Character{81}}%
\TextAt{21}{155.339}{723.545}{k}%
\NativeMathAt{159.961}{723.545}{7.97}{\in}%
\TextAt{21}{165.606}{723.545}{I}%
\TextAt{13}{172.39}{727.032}{X}%
\TextAt{21}{182.106}{725.238}{k}%
\NativeMathAt{190.546}{727.032}{11.955}{\rightarrow}%
\TextAt{13}{205.822}{727.032}{X}%
\TextAt{21}{215.537}{725.238}{i}%
\TextAt{8}{222.218}{727.032}{are}%
\TextAt{8}{241.13}{727.032}{continuous,}%
\TextAt{8}{302.507}{727.032}{and}%
\TextAt{8}{324.671}{727.032}{we}%
\TextAt{8}{341.3}{727.032}{may}%
\TextAt{8}{366.071}{727.032}{write}%
\TextAt{13}{395.384}{727.032}{X}%
\TextAt{21}{405.099}{725.238}{ij}%
\TextAt{9}{415.685}{727.032}{\Character{61}}%
\TextAt{16}{428.111}{727.032}{\Character{102}}%
\TextAt{13}{434.088}{727.032}{x}%
\NativeMathAt{444.061}{727.032}{11.955}{\in}%
\TextAt{20}{455.352}{735.998}{\Character{81}}%
\TextAt{21}{466.643}{723.545}{k}%
\TextAt{13}{473.756}{727.032}{X}%
\TextAt{21}{483.471}{725.238}{k}%
\TextAt{16}{491.911}{727.032}{\Character{106}}%
\TextAt{13}{498.553}{727.032}{g}%
\TextAt{9}{504.587}{727.032}{\Character{40}}%
\TextAt{13}{509.14}{727.032}{x}%
\TextAt{9}{515.792}{727.032}{\Character{41}}%
\TextAt{9}{523.647}{727.032}{\Character{61}}%
\TextAt{13}{79.226}{709.986}{f}%
\TextAt{9}{86.273}{709.986}{\Character{40}}%
\TextAt{13}{90.825}{709.986}{x}%
\TextAt{9}{97.477}{709.986}{\Character{41}}%
\TextAt{16}{102.029}{709.986}{\Character{103}}%
\TextAt{8}{108.007}{709.986}{.}%
\TextAt{8}{116.886}{709.986}{But}%
\TextAt{8}{140.27}{709.986}{in}%
\TextAt{8}{154.067}{709.986}{the}%
\TextAt{8}{174.367}{709.986}{Hausdorff}%
\TextAt{8}{227.902}{709.986}{case}%
\TextAt{8}{252.817}{709.986}{the}%
\TextAt{8}{273.117}{709.986}{equation}%
\TextAt{13}{321.069}{709.986}{g}%
\TextAt{9}{327.103}{709.986}{\Character{40}}%
\TextAt{13}{331.655}{709.986}{x}%
\TextAt{9}{338.308}{709.986}{\Character{41}}%
\TextAt{9}{346.414}{709.986}{\Character{61}}%
\TextAt{13}{359.095}{709.986}{f}%
\TextAt{9}{366.142}{709.986}{\Character{40}}%
\TextAt{13}{370.694}{709.986}{x}%
\TextAt{9}{377.346}{709.986}{\Character{41}}%
\TextAt{8}{385.94}{709.986}{defines}%
\TextAt{8}{424.507}{709.986}{a}%
\TextAt{8}{434.406}{709.986}{closed}%
\TextAt{8}{469.076}{709.986}{subset.}%
\TextAt{8}{509.947}{709.986}{This}%
\TextAt{8}{79.226}{692.94}{representation}%
\TextAt{13}{155.715}{692.94}{X}%
\TextAt{9}{169.691}{692.94}{\Character{61}}%
\TextAt{20}{182.117}{701.906}{\Character{84}}%
\TextAt{21}{192.079}{689.453}{i}%
\NativeMathAt{194.963}{689.453}{7.97}{\leq}%
\TextAt{21}{201.549}{689.453}{j}%
\TextAt{13}{207.924}{692.94}{X}%
\TextAt{21}{217.639}{691.147}{ij}%
\TextAt{8}{228.802}{692.94}{also}%
\TextAt{8}{252.282}{692.94}{gives}%
\TextAt{8}{280.962}{692.94}{the}%
\TextAt{8}{301.118}{692.94}{following}%
\TextAt{12}{79.226}{667.599}{Proposition}%
\TextAt{12}{151.829}{667.599}{(2.3).}%
\ErrataAt{186.95}{667.599}{348.05}{The projective limit $X=\varprojlim_i X_i$ of nonempty compact Hausdorff}%
\ErrataAt{79.19}{650.553}{455.81}{spaces $X_i$ is itself nonempty and compact.}%
\TextAt{19}{79.226}{616.916}{Proof.}%
\TextAt{8}{116}{616.916}{If}%
\TextAt{8}{127.281}{616.916}{all}%
\TextAt{8}{143.128}{616.916}{the}%
\TextAt{13}{162.873}{616.916}{X}%
\TextAt{21}{172.588}{615.123}{i}%
\TextAt{8}{179.455}{616.916}{are}%
\TextAt{8}{198.555}{616.916}{compact,}%
\TextAt{8}{248.305}{616.916}{then}%
\TextAt{8}{274.554}{616.916}{so}%
\TextAt{8}{288.512}{616.916}{is}%
\TextAt{8}{299.865}{616.916}{the}%
\TextAt{8}{319.609}{616.916}{product}%
\TextAt{20}{363.097}{625.882}{\Character{81}}%
\TextAt{21}{374.388}{613.429}{i}%
\NativeMathAt{377.271}{613.429}{7.97}{\in}%
\TextAt{21}{382.917}{613.429}{I}%
\TextAt{13}{389.701}{616.916}{X}%
\TextAt{21}{399.416}{615.123}{i}%
\TextAt{8}{402.798}{616.916}{,}%
\TextAt{8}{409.618}{616.916}{by}%
\TextAt{8}{425.465}{616.916}{Tykhonov’s}%
\TextAt{8}{487.89}{616.916}{theorem,}%
\TextAt{8}{79.226}{599.87}{and}%
\TextAt{8}{102.157}{599.87}{thus}%
\TextAt{8}{128.077}{599.87}{also}%
\TextAt{8}{151.725}{599.87}{the}%
\TextAt{8}{172.05}{599.87}{closed}%
\TextAt{8}{206.745}{599.87}{subset}%
\TextAt{13}{242.803}{599.87}{X}%
\TextAt{8}{253.458}{599.87}{.}%
\TextAt{8}{262.411}{599.87}{Furthermore,}%
\TextAt{13}{333.671}{599.87}{X}%
\TextAt{9}{347.934}{599.87}{\Character{61}}%
\TextAt{20}{360.646}{608.836}{\Character{84}}%
\TextAt{21}{370.609}{596.383}{i}%
\NativeMathAt{373.492}{596.383}{7.97}{\leq}%
\TextAt{21}{380.078}{596.383}{j}%
\TextAt{13}{386.453}{599.87}{X}%
\TextAt{21}{396.168}{598.077}{ij}%
\TextAt{8}{407.499}{599.87}{cannot}%
\TextAt{8}{446.044}{599.87}{be}%
\TextAt{8}{462.136}{599.87}{the}%
\TextAt{8}{482.461}{599.87}{empty}%
\TextAt{8}{518.399}{599.87}{set}%
\TextAt{8}{79.226}{582.824}{if}%
\TextAt{8}{89.679}{582.824}{the}%
\TextAt{13}{109.564}{582.824}{X}%
\TextAt{21}{119.279}{581.031}{i}%
\TextAt{8}{126.286}{582.824}{are}%
\TextAt{8}{145.526}{582.824}{nonempty.}%
\TextAt{8}{203.641}{582.824}{In}%
\TextAt{8}{217.991}{582.824}{fact,}%
\TextAt{8}{244.111}{582.824}{as}%
\TextAt{8}{258.21}{582.824}{the}%
\TextAt{8}{278.095}{582.824}{product}%
\TextAt{20}{321.724}{591.791}{\Character{81}}%
\TextAt{21}{333.015}{579.337}{i}%
\TextAt{13}{338.388}{582.824}{X}%
\TextAt{21}{348.104}{581.031}{i}%
\TextAt{8}{355.111}{582.824}{is}%
\TextAt{8}{366.604}{582.824}{compact,}%
\TextAt{8}{416.467}{582.824}{there}%
\TextAt{8}{446.107}{582.824}{would}%
\TextAt{8}{479.98}{582.824}{have}%
\TextAt{8}{506.704}{582.824}{to}%
\TextAt{8}{520.743}{582.824}{be}%
\TextAt{8}{79.226}{565.778}{an}%
\TextAt{8}{95.681}{565.778}{intersection}%
\TextAt{8}{158.702}{565.778}{of}%
\TextAt{8}{172.228}{565.778}{finitely}%
\TextAt{8}{211.769}{565.778}{many}%
\TextAt{13}{243.837}{565.778}{X}%
\TextAt{21}{253.553}{563.985}{ij}%
\TextAt{8}{264.911}{565.778}{which}%
\TextAt{8}{298.594}{565.778}{is}%
\TextAt{8}{310.554}{565.778}{empty.}%
\TextAt{8}{350.483}{565.778}{But}%
\TextAt{8}{373.92}{565.778}{this}%
\TextAt{8}{396.938}{565.778}{is}%
\TextAt{8}{408.898}{565.778}{impossible:}%
\TextAt{8}{470.868}{565.778}{if}%
\TextAt{8}{481.788}{565.778}{all}%
\TextAt{8}{498.243}{565.778}{indices}%
\TextAt{8}{79.226}{548.732}{entering}%
\TextAt{8}{125.392}{548.732}{into}%
\TextAt{8}{150.099}{548.732}{this}%
\TextAt{8}{173.885}{548.732}{finite}%
\TextAt{8}{204.76}{548.732}{intersection}%
\TextAt{8}{268.548}{548.732}{satisfy}%
\TextAt{13}{306.059}{548.732}{i\Character{59}}%
\TextAt{13}{315.276}{548.732}{j}%
\NativeMathAt{325.775}{548.732}{11.955}{\leq}%
\TextAt{13}{340.037}{548.732}{n}%
\TextAt{8}{347.024}{548.732}{,}%
\TextAt{8}{355.378}{548.732}{and}%
\TextAt{8}{379.104}{548.732}{if}%
\TextAt{13}{390.792}{548.732}{x}%
\TextAt{21}{397.444}{546.939}{n}%
\NativeMathAt{408.043}{548.732}{11.955}{\in}%
\TextAt{13}{420.976}{548.732}{X}%
\TextAt{21}{430.691}{546.939}{n}%
\TextAt{8}{436.328}{548.732}{,}%
\TextAt{8}{444.681}{548.732}{then}%
\TextAt{8}{472.305}{548.732}{the}%
\TextAt{8}{493.425}{548.732}{element}%
\TextAt{9}{79.226}{531.686}{\Character{40}}%
\TextAt{13}{83.778}{531.686}{x}%
\TextAt{21}{90.431}{529.893}{i}%
\TextAt{9}{93.812}{531.686}{\Character{41}}%
\TextAt{21}{98.364}{529.893}{i}%
\NativeMathAt{101.247}{529.893}{7.97}{\in}%
\TextAt{21}{106.893}{529.893}{I}%
\TextAt{8}{114.836}{531.686}{belongs}%
\TextAt{8}{156.1}{531.686}{to}%
\TextAt{8}{169.665}{531.686}{this}%
\TextAt{8}{191.741}{531.686}{intersection,}%
\TextAt{8}{257.22}{531.686}{where}%
\TextAt{8}{290.283}{531.686}{we}%
\TextAt{8}{306.765}{531.686}{choose}%
\TextAt{13}{343.151}{531.686}{x}%
\TextAt{21}{349.804}{529.893}{i}%
\TextAt{9}{356.506}{531.686}{\Character{61}}%
\TextAt{13}{368.931}{531.686}{f}%
\TextAt{21}{374.701}{529.893}{in}%
\TextAt{9}{383.221}{531.686}{\Character{40}}%
\TextAt{13}{387.773}{531.686}{x}%
\TextAt{21}{394.425}{529.893}{n}%
\TextAt{9}{400.061}{531.686}{\Character{41}}%
\TextAt{8}{407.765}{531.686}{for}%
\TextAt{13}{424.904}{531.686}{i}%
\NativeMathAt{432.218}{531.686}{11.955}{\leq}%
\TextAt{13}{444.837}{531.686}{n}%
\TextAt{8}{451.825}{531.686}{,}%
\TextAt{8}{458.377}{531.686}{and}%
\TextAt{8}{480.394}{531.686}{arbitrarily}%
\TextAt{8}{79.226}{514.64}{for}%
\TextAt{8}{97.111}{514.64}{all}%
\TextAt{8}{113.37}{514.64}{other}%
\TextAt{13}{143.939}{514.64}{i}%
\TextAt{8}{147.933}{514.64}{.}%
\Vector{q 0 G 0 g 0.399 w 452.334 -259.36 m 452.334 -251.29 l S Q}%
\Vector{q 0 G 0 g 0.399 w 452.534 -251.489 m 459.707 -251.489 l S Q}%
\Vector{q 0 G 0 g 0.399 w 452.534 -259.161 m 459.707 -259.161 l S Q}%
\Vector{q 0 G 0 g 0.399 w 459.906 -259.36 m 459.906 -251.29 l S Q}%
\TextAt{12}{79.226}{481.004}{Definition}%
\TextAt{12}{142.015}{481.004}{(2.4).}%
\TextAt{8}{177.33}{481.004}{The}%
\TextAt{12}{201.384}{481.004}{inductive}%
\TextAt{12}{259.845}{481.004}{limit}%
\TextAt{13}{277.123}{454.351}{X}%
\TextAt{9}{291.099}{454.351}{\Character{61}}%
\NativeMathAt{303.525}{454.351}{11.955}{\varinjlim}%
\TextAt{21}{305.243}{441.971}{i}%
\NativeMathAt{308.126}{441.971}{7.97}{\in}%
\TextAt{21}{313.772}{441.971}{I}%
\TextAt{13}{321.776}{454.351}{X}%
\TextAt{21}{331.491}{452.558}{i}%
\TextAt{8}{79.226}{420.704}{of}%
\TextAt{8}{92.556}{420.704}{an}%
\TextAt{8}{108.815}{420.704}{inductive}%
\TextAt{8}{159.541}{420.704}{system}%
\TextAt{16}{198.359}{420.704}{\Character{102}}%
\TextAt{13}{204.337}{420.704}{X}%
\TextAt{21}{214.052}{418.91}{i}%
\TextAt{13}{217.434}{420.704}{\Character{59}}%
\TextAt{13}{222.658}{420.704}{f}%
\TextAt{21}{228.448}{418.91}{ij}%
\TextAt{16}{235.713}{420.704}{\Character{103}}%
\TextAt{8}{245.588}{420.704}{is}%
\TextAt{8}{257.352}{420.704}{defined}%
\TextAt{8}{297.665}{420.704}{to}%
\TextAt{8}{311.975}{420.704}{be}%
\TextAt{8}{327.899}{420.704}{the}%
\TextAt{8}{348.056}{420.704}{quotient}%
\TextAt{13}{257.39}{385.836}{X}%
\TextAt{9}{271.366}{385.836}{\Character{61}}%
\TextAt{20}{283.792}{406.28}{\Character{32}}%
\TextAt{20}{293.256}{397.194}{\Character{97}}%
\TextAt{21}{294.483}{371.822}{i}%
\NativeMathAt{297.366}{371.822}{7.97}{\in}%
\TextAt{21}{303.012}{371.822}{I}%
\TextAt{13}{310.525}{385.836}{X}%
\TextAt{21}{320.24}{384.043}{i}%
\TextAt{20}{323.621}{406.28}{\Character{33}}%
\TextAt{20}{335.078}{395.52}{\Character{14}}%
\NativeMathAt{345.307}{385.836}{11.955}{\sim}%
\TextAt{8}{79.226}{350.554}{of}%
\TextAt{8}{94.038}{350.554}{the}%
\TextAt{8}{115.677}{350.554}{disjoint}%
\TextAt{8}{158.847}{350.554}{union}%
\TextAt{20}{192.847}{359.521}{\Character{96}}%
\TextAt{21}{204.138}{347.067}{i}%
\NativeMathAt{207.021}{347.067}{7.97}{\in}%
\TextAt{21}{212.667}{347.067}{I}%
\TextAt{13}{219.451}{350.554}{X}%
\TextAt{21}{229.167}{348.761}{i}%
\TextAt{8}{232.548}{350.554}{,}%
\TextAt{8}{241.55}{350.554}{where}%
\TextAt{8}{276.841}{350.554}{we}%
\TextAt{8}{295.551}{350.554}{consider}%
\TextAt{8}{342.618}{350.554}{two}%
\TextAt{8}{366.217}{350.554}{elements}%
\TextAt{13}{415.556}{350.554}{x}%
\TextAt{21}{422.208}{348.761}{i}%
\NativeMathAt{431.436}{350.554}{11.955}{\in}%
\TextAt{13}{445.253}{350.554}{X}%
\TextAt{21}{454.968}{348.761}{i}%
\TextAt{8}{463.729}{350.554}{and}%
\TextAt{13}{487.974}{350.554}{x}%
\TextAt{21}{494.626}{348.761}{j}%
\NativeMathAt{504.855}{350.554}{11.955}{\in}%
\TextAt{13}{518.672}{350.554}{X}%
\TextAt{21}{528.387}{348.761}{j}%
\TextAt{8}{79.226}{333.508}{equivalent}%
\TextAt{8}{134.842}{333.508}{if}%
\TextAt{8}{145.565}{333.508}{there}%
\TextAt{8}{175.477}{333.508}{exists}%
\TextAt{8}{207.792}{333.508}{a}%
\TextAt{13}{217.547}{333.508}{k}%
\NativeMathAt{227.358}{333.508}{11.955}{\geq}%
\TextAt{13}{239.977}{333.508}{i\Character{59}}%
\TextAt{13}{249.195}{333.508}{j}%
\TextAt{8}{258.628}{333.508}{such}%
\TextAt{8}{285.025}{333.508}{that}%
\TextAt{13}{262.587}{306.856}{f}%
\TextAt{21}{268.357}{305.063}{ik}%
\TextAt{9}{276.36}{306.856}{\Character{40}}%
\TextAt{13}{280.912}{306.856}{x}%
\TextAt{21}{287.564}{305.063}{i}%
\TextAt{9}{290.946}{306.856}{\Character{41}}%
\TextAt{9}{298.801}{306.856}{\Character{61}}%
\TextAt{13}{311.244}{306.856}{f}%
\TextAt{21}{317.014}{305.063}{j}%
\TextAt{21}{320.887}{305.063}{k}%
\TextAt{9}{326.018}{306.856}{\Character{40}}%
\TextAt{13}{330.57}{306.856}{x}%
\TextAt{21}{337.222}{305.063}{j}%
\TextAt{9}{341.605}{306.856}{\Character{41}}%
\TextAt{13}{346.157}{306.856}{\Character{58}}%
\TextAt{8}{97.159}{271.908}{In}%
\TextAt{8}{113.476}{271.908}{the}%
\TextAt{8}{135.328}{271.908}{applications,}%
\TextAt{8}{205.807}{271.908}{the}%
\TextAt{8}{227.659}{271.908}{projective}%
\TextAt{8}{283.332}{271.908}{and}%
\TextAt{8}{307.79}{271.908}{inductive}%
\TextAt{8}{360.211}{271.908}{systems}%
\TextAt{16}{405.34}{271.908}{\Character{102}}%
\TextAt{13}{411.318}{271.908}{X}%
\TextAt{21}{421.033}{270.115}{i}%
\TextAt{13}{424.414}{271.908}{\Character{59}}%
\TextAt{13}{429.638}{271.908}{f}%
\TextAt{21}{435.428}{270.115}{ij}%
\TextAt{16}{442.694}{271.908}{\Character{103}}%
\TextAt{8}{454.264}{271.908}{that}%
\TextAt{8}{481.329}{271.908}{occur}%
\TextAt{8}{514.562}{271.908}{will}%
\TextAt{8}{79.226}{254.863}{not}%
\TextAt{8}{99.596}{254.863}{just}%
\TextAt{8}{122.298}{254.863}{be}%
\TextAt{8}{137.778}{254.863}{systems}%
\TextAt{8}{180.767}{254.863}{of}%
\TextAt{8}{193.654}{254.863}{topological}%
\TextAt{8}{252.734}{254.863}{spaces}%
\TextAt{8}{288.18}{254.863}{and}%
\TextAt{8}{310.499}{254.863}{continuous}%
\TextAt{8}{368.659}{254.863}{maps,}%
\TextAt{8}{402.185}{254.863}{but}%
\TextAt{8}{423.201}{254.863}{the}%
\TextAt{13}{442.913}{254.863}{X}%
\TextAt{21}{452.629}{253.069}{i}%
\TextAt{8}{459.463}{254.863}{will}%
\TextAt{8}{481.125}{254.863}{usually}%
\TextAt{8}{520.743}{254.863}{be}%
\TextAt{8}{79.226}{237.817}{topological}%
\TextAt{8}{138.079}{237.817}{groups,}%
\TextAt{8}{178.584}{237.817}{rings}%
\TextAt{8}{206.593}{237.817}{or}%
\TextAt{8}{220.232}{237.817}{modules,}%
\TextAt{8}{268.854}{237.817}{etc.,}%
\TextAt{8}{293.674}{237.817}{and}%
\TextAt{8}{315.765}{237.817}{the}%
\TextAt{13}{335.25}{237.817}{f}%
\TextAt{21}{341.02}{236.023}{ij}%
\TextAt{8}{351.511}{237.817}{will}%
\TextAt{8}{372.945}{237.817}{be}%
\TextAt{8}{388.198}{237.817}{continuous}%
\TextAt{8}{446.13}{237.817}{homomorphisms.}%
\TextAt{8}{79.226}{220.771}{In}%
\TextAt{8}{92.852}{220.771}{what}%
\TextAt{8}{121.123}{220.771}{follows,}%
\TextAt{8}{162.015}{220.771}{we}%
\TextAt{8}{178.247}{220.771}{will}%
\TextAt{8}{199.357}{220.771}{deal}%
\TextAt{8}{223.074}{220.771}{explicitly}%
\TextAt{8}{272.804}{220.771}{only}%
\TextAt{8}{297.501}{220.771}{with}%
\TextAt{8}{323.166}{220.771}{projective}%
\TextAt{8}{376.148}{220.771}{and}%
\TextAt{8}{397.916}{220.771}{inductive}%
\TextAt{8}{447.647}{220.771}{systems}%
\TextAt{16}{490.085}{220.771}{\Character{102}}%
\TextAt{13}{496.062}{220.771}{G}%
\TextAt{21}{505.296}{218.977}{i}%
\TextAt{13}{508.677}{220.771}{\Character{59}}%
\TextAt{13}{513.901}{220.771}{g}%
\TextAt{21}{519.527}{218.977}{ij}%
\TextAt{16}{526.792}{220.771}{\Character{103}}%
\TextAt{8}{79.226}{203.725}{of}%
\TextAt{8}{93.312}{203.725}{topological}%
\TextAt{8}{153.593}{203.725}{groups.}%
\TextAt{8}{198.201}{203.725}{But}%
\TextAt{8}{222.198}{203.725}{since}%
\TextAt{8}{251.622}{203.725}{everything}%
\TextAt{8}{309.942}{203.725}{works}%
\TextAt{8}{343.934}{203.725}{exactly}%
\TextAt{8}{385.014}{203.725}{the}%
\TextAt{8}{405.926}{203.725}{same}%
\TextAt{8}{436.008}{203.725}{way}%
\TextAt{8}{460.507}{203.725}{for}%
\TextAt{8}{479.148}{203.725}{systems}%
\TextAt{8}{523.337}{203.725}{of}%
\TextAt{8}{79.226}{186.679}{rings}%
\TextAt{8}{107.907}{186.679}{or}%
\TextAt{8}{122.217}{186.679}{modules,}%
\TextAt{8}{171.377}{186.679}{these}%
\TextAt{8}{201.348}{186.679}{cases}%
\TextAt{8}{230.734}{186.679}{may}%
\TextAt{8}{256.103}{186.679}{be}%
\TextAt{8}{272.027}{186.679}{thought}%
\TextAt{8}{315.939}{186.679}{of}%
\TextAt{8}{329.269}{186.679}{tacitly}%
\TextAt{8}{366.019}{186.679}{as}%
\TextAt{8}{380.389}{186.679}{being}%
\TextAt{8}{411.927}{186.679}{treated}%
\TextAt{8}{452.251}{186.679}{as}%
\TextAt{8}{466.621}{186.679}{well.}%
\TextAt{8}{97.159}{166.644}{Let}%
\TextAt{16}{118.391}{166.644}{\Character{102}}%
\TextAt{13}{124.369}{166.644}{G}%
\TextAt{21}{133.602}{164.851}{i}%
\TextAt{13}{136.984}{166.644}{\Character{59}}%
\TextAt{13}{142.208}{166.644}{g}%
\TextAt{21}{147.833}{164.851}{ij}%
\TextAt{16}{155.098}{166.644}{\Character{103}}%
\TextAt{8}{165.236}{166.644}{be}%
\TextAt{8}{181.423}{166.644}{a}%
\TextAt{8}{191.442}{166.644}{projective,}%
\TextAt{8}{249}{166.644}{resp.}%
\TextAt{8}{279.11}{166.644}{inductive}%
\TextAt{8}{330.099}{166.644}{system}%
\TextAt{8}{369.18}{166.644}{of}%
\TextAt{8}{382.773}{166.644}{topological}%
\TextAt{8}{442.56}{166.644}{groups.}%
\TextAt{8}{485.69}{166.644}{Then}%
\TextAt{8}{516.51}{166.644}{the}%
\TextAt{8}{79.226}{149.598}{projective,}%
\TextAt{8}{136.456}{149.598}{resp.}%
\TextAt{8}{165.778}{149.598}{inductive}%
\TextAt{8}{216.503}{149.598}{limit}%
\TextAt{13}{220.913}{122.946}{G}%
\TextAt{9}{233.468}{122.946}{\Character{61}}%
\NativeMathAt{245.893}{122.946}{11.955}{\varprojlim}%
\TextAt{21}{247.611}{110.565}{i}%
\NativeMathAt{250.495}{110.565}{7.97}{\in}%
\TextAt{21}{256.14}{110.565}{I}%
\TextAt{13}{264.145}{122.946}{G}%
\TextAt{21}{273.378}{121.152}{i}%
\TextAt{13}{276.76}{122.946}{\Character{59}}%
\TextAt{8}{305.914}{122.946}{resp.}%
\TextAt{13}{335.236}{122.946}{G}%
\TextAt{9}{347.791}{122.946}{\Character{61}}%
\NativeMathAt{360.216}{122.946}{11.955}{\varinjlim}%
\TextAt{21}{361.934}{110.565}{i}%
\NativeMathAt{364.817}{110.565}{7.97}{\in}%
\TextAt{21}{370.463}{110.565}{I}%
\TextAt{13}{378.468}{122.946}{G}%
\TextAt{21}{387.701}{121.152}{i}%
\TextAt{8}{79.226}{89.967}{is}%
\TextAt{8}{91.035}{89.967}{a}%
\TextAt{12}{100.835}{89.967}{topological}%
\TextAt{12}{169.114}{89.967}{group}%
\TextAt{8}{206.961}{89.967}{as}%
\TextAt{8}{221.376}{89.967}{well.}%
\TextAt{8}{249.792}{89.967}{The}%
\TextAt{8}{273.89}{89.967}{multiplication}%
\TextAt{8}{349.061}{89.967}{in}%
\TextAt{8}{362.759}{89.967}{the}%
\TextAt{8}{382.96}{89.967}{projective}%
\TextAt{8}{436.982}{89.967}{limit}%
\TextAt{8}{464.99}{89.967}{is}%
\TextAt{8}{476.799}{89.967}{induced}%
\TextAt{8}{520.408}{89.967}{by}%
\TextAt{8}{79.226}{72.921}{the}%
\TextAt{8}{100.239}{72.921}{componentwise}%
\TextAt{8}{182.127}{72.921}{multiplication}%
\TextAt{8}{258.11}{72.921}{in}%
\TextAt{8}{272.619}{72.921}{the}%
\TextAt{8}{293.632}{72.921}{product}%
\TextAt{20}{338.388}{81.888}{\Character{81}}%
\TextAt{21}{349.679}{69.434}{i}%
\NativeMathAt{352.562}{69.434}{7.97}{\in}%
\TextAt{21}{358.208}{69.434}{I}%
\TextAt{13}{364.992}{72.921}{G}%
\TextAt{21}{374.225}{71.128}{i}%
\TextAt{8}{377.607}{72.921}{.}%
\TextAt{8}{388.624}{72.921}{In}%
\TextAt{8}{404.102}{72.921}{the}%
\TextAt{8}{425.114}{72.921}{case}%
\TextAt{8}{450.742}{72.921}{of}%
\TextAt{8}{464.928}{72.921}{the}%
\TextAt{8}{485.941}{72.921}{inductive}%
\TextAt{8}{297.211}{43.033}{244}%
\Destination{78.23}{809.815}{nav.page.244}%
\Destination{79.36}{269.855}{nav.print.267}%
\Destination{79.36}{681.555}{nav.number.4.2.3}%
\Destination{79.36}{494.955}{nav.number.4.2.4}%
\Destination{79.226}{684.599}{errata.102}%
}

% Page 257 (body page 245)
\Page{%
\TextAt{2}{79.226}{796.943}{IV.}%
\TextAt{2}{98.012}{796.943}{Abstract}%
\TextAt{2}{143.252}{796.943}{Class}%
\TextAt{2}{171.844}{796.943}{Field}%
\TextAt{2}{199.575}{796.943}{Theory}%
\Vector{q 0 G 0 g 0.299 w 7.226 18.729 m 460.769 18.729 l S Q}%
\TextAt{8}{79.226}{761.124}{limit,}%
\TextAt{8}{109.771}{761.124}{given}%
\TextAt{8}{139.504}{761.124}{two}%
\TextAt{8}{160.784}{761.124}{equivalence}%
\TextAt{8}{221.731}{761.124}{classes}%
\TextAt{13}{258.147}{761.124}{x\Character{59}}%
\TextAt{13}{270.019}{761.124}{y}%
\NativeMathAt{279.5}{761.124}{11.955}{\in}%
\TextAt{13}{290.791}{761.124}{G}%
\TextAt{9}{303.346}{761.124}{\Character{61}}%
\NativeMathAt{315.771}{761.124}{11.955}{\varinjlim}%
\TextAt{21}{332.03}{755.684}{i}%
\NativeMathAt{334.913}{755.684}{7.97}{\in}%
\TextAt{21}{340.559}{755.684}{I}%
\TextAt{13}{347.343}{761.124}{G}%
\TextAt{21}{356.577}{759.33}{i}%
\TextAt{8}{359.958}{761.124}{,}%
\TextAt{8}{366.438}{761.124}{one}%
\TextAt{8}{387.06}{761.124}{has}%
\TextAt{8}{407.097}{761.124}{to}%
\TextAt{8}{420.57}{761.124}{choose}%
\TextAt{8}{456.866}{761.124}{representatives}%
\TextAt{13}{79.226}{744.078}{x}%
\TextAt{21}{85.878}{742.284}{k}%
\TextAt{8}{94.895}{744.078}{and}%
\TextAt{13}{117.658}{744.078}{y}%
\TextAt{21}{123.366}{742.284}{k}%
\TextAt{8}{132.383}{744.078}{in}%
\TextAt{8}{146.036}{744.078}{the}%
\TextAt{8}{166.192}{744.078}{same}%
\TextAt{13}{195.518}{744.078}{G}%
\TextAt{21}{204.752}{742.284}{k}%
\TextAt{8}{213.769}{744.078}{in}%
\TextAt{8}{227.422}{744.078}{order}%
\TextAt{8}{257.991}{744.078}{to}%
\TextAt{8}{272.301}{744.078}{define}%
\TextAt{13}{227.529}{715.077}{xy}%
\TextAt{9}{243.639}{715.077}{\Character{61}}%
\TextAt{8}{256.064}{715.077}{equivalence}%
\TextAt{8}{317.848}{715.077}{class}%
\TextAt{8}{345.286}{715.077}{of}%
\TextAt{13}{358.616}{715.077}{x}%
\TextAt{21}{365.268}{713.283}{k}%
\TextAt{13}{370.387}{715.077}{y}%
\TextAt{21}{376.095}{713.283}{k}%
\TextAt{13}{381.215}{715.077}{\Character{58}}%
\ErrataAt{79.19}{686.075}{455.81}{The multiplication is independent of the representatives and defines a group. Equip it with}%
\ErrataAt{79.19}{669.029}{455.81}{the finest group topology for which the canonical maps $G_i\to G$ are continuous.}%
\TextAt{8}{97.159}{648.995}{The}%
\TextAt{8}{121.797}{648.995}{projections}%
\TextAt{13}{182.276}{648.995}{p}%
\TextAt{21}{188.151}{647.201}{i}%
\TextAt{9}{195.848}{648.995}{\Character{58}}%
\TextAt{20}{203.416}{657.961}{\Character{81}}%
\TextAt{21}{214.707}{645.508}{i}%
\NativeMathAt{217.59}{645.508}{7.97}{\in}%
\TextAt{21}{223.235}{645.508}{I}%
\TextAt{13}{230.02}{648.995}{G}%
\TextAt{21}{239.253}{647.201}{i}%
\NativeMathAt{246.95}{648.995}{11.955}{\rightarrow}%
\TextAt{13}{263.221}{648.995}{G}%
\TextAt{21}{272.455}{647.201}{i}%
\TextAt{8}{275.836}{648.995}{,}%
\TextAt{8}{283.715}{648.995}{resp.}%
\TextAt{8}{314.789}{648.995}{the}%
\TextAt{8}{335.529}{648.995}{inclusions}%
\TextAt{13}{389.565}{648.995}{\Character{19}}%
\TextAt{21}{393.702}{647.201}{i}%
\TextAt{9}{401.399}{648.995}{\Character{58}}%
\TextAt{13}{408.967}{648.995}{G}%
\TextAt{21}{418.2}{647.201}{i}%
\NativeMathAt{425.897}{648.995}{11.955}{\rightarrow}%
\TextAt{20}{442.169}{657.961}{\Character{96}}%
\TextAt{21}{453.46}{645.508}{i}%
\NativeMathAt{456.343}{645.508}{7.97}{\in}%
\TextAt{21}{461.988}{645.508}{I}%
\TextAt{13}{468.773}{648.995}{G}%
\TextAt{21}{478.006}{647.201}{i}%
\TextAt{8}{481.387}{648.995}{,}%
\TextAt{8}{489.267}{648.995}{induce}%
\TextAt{8}{526.911}{648.995}{a}%
\TextAt{8}{79.226}{631.949}{family}%
\TextAt{8}{114.996}{631.949}{of}%
\TextAt{8}{128.326}{631.949}{continuous}%
\TextAt{8}{186.93}{631.949}{homomorphisms}%
\TextAt{13}{210.128}{602.948}{g}%
\TextAt{21}{215.734}{601.154}{i}%
\TextAt{9}{222.436}{602.948}{\Character{58}}%
\TextAt{13}{229.008}{602.948}{G}%
\NativeArrowAt{241.563}{602.948}{11.955}{19.26}{\longrightarrow}%
\TextAt{13}{264.145}{602.948}{G}%
\TextAt{21}{273.378}{601.154}{i}%
\TextAt{13}{276.76}{602.948}{\Character{59}}%
\TextAt{8}{305.914}{602.948}{resp.}%
\TextAt{13}{335.236}{602.948}{g}%
\TextAt{21}{340.841}{601.154}{i}%
\TextAt{9}{347.543}{602.948}{\Character{58}}%
\TextAt{13}{354.116}{602.948}{G}%
\TextAt{21}{363.35}{601.154}{i}%
\NativeArrowAt{370.052}{602.948}{11.955}{19.26}{\longrightarrow}%
\TextAt{13}{392.634}{602.948}{G}%
\TextAt{8}{79.226}{573.946}{such}%
\TextAt{8}{105.084}{573.946}{that}%
\TextAt{13}{129.914}{573.946}{g}%
\TextAt{21}{135.52}{572.153}{i}%
\TextAt{9}{142.222}{573.946}{\Character{61}}%
\TextAt{13}{154.647}{573.946}{g}%
\TextAt{21}{160.252}{572.153}{ij}%
\TextAt{16}{169.072}{573.946}{\Character{14}}%
\TextAt{13}{176.605}{573.946}{g}%
\TextAt{21}{182.21}{572.153}{j}%
\TextAt{8}{186.592}{573.946}{,}%
\TextAt{8}{193.31}{573.946}{resp.}%
\TextAt{13}{222.453}{573.946}{g}%
\TextAt{21}{228.058}{572.153}{i}%
\TextAt{9}{234.76}{573.946}{\Character{61}}%
\TextAt{13}{247.186}{573.946}{g}%
\TextAt{21}{252.791}{572.153}{j}%
\TextAt{16}{258.728}{573.946}{\Character{14}}%
\TextAt{13}{266.26}{573.946}{g}%
\TextAt{21}{271.865}{572.153}{ij}%
\TextAt{8}{279.131}{573.946}{,}%
\TextAt{8}{285.849}{573.946}{for}%
\TextAt{13}{303.195}{573.946}{i}%
\NativeMathAt{310.509}{573.946}{11.955}{\leq}%
\TextAt{13}{323.128}{573.946}{j}%
\TextAt{8}{328.644}{573.946}{.}%
\TextAt{8}{336.912}{573.946}{This}%
\TextAt{8}{363.093}{573.946}{family}%
\TextAt{8}{398.324}{573.946}{has}%
\TextAt{8}{418.659}{573.946}{the}%
\TextAt{8}{438.277}{573.946}{following}%
\TextAt{8}{487.172}{573.946}{universal}%
\TextAt{8}{79.226}{556.9}{property.}%
\TextAt{12}{79.226}{530.888}{Proposition}%
\TextAt{12}{151.937}{530.888}{(2.5).}%
\TextAt{8}{187.253}{530.888}{If}%
\TextAt{13}{198.945}{530.888}{H}%
\TextAt{8}{213.465}{530.888}{is}%
\TextAt{8}{225.229}{530.888}{a}%
\TextAt{8}{234.985}{530.888}{topological}%
\TextAt{8}{294.509}{530.888}{group}%
\TextAt{8}{327.685}{530.888}{and}%
\TextAt{13}{207.606}{501.887}{h}%
\TextAt{21}{214.345}{500.094}{i}%
\TextAt{9}{221.047}{501.887}{\Character{58}}%
\TextAt{13}{227.619}{501.887}{H}%
\NativeArrowAt{241.563}{501.887}{11.955}{19.26}{\longrightarrow}%
\TextAt{13}{264.145}{501.887}{G}%
\TextAt{21}{273.378}{500.094}{i}%
\TextAt{13}{276.76}{501.887}{\Character{59}}%
\TextAt{8}{305.914}{501.887}{resp.}%
\TextAt{13}{335.236}{501.887}{h}%
\TextAt{21}{341.975}{500.094}{i}%
\TextAt{9}{348.677}{501.887}{\Character{58}}%
\TextAt{13}{355.249}{501.887}{G}%
\TextAt{21}{364.483}{500.094}{i}%
\NativeArrowAt{371.185}{501.887}{11.955}{19.26}{\longrightarrow}%
\TextAt{13}{393.767}{501.887}{H}%
\TextAt{8}{79.226}{472.886}{is}%
\TextAt{8}{90.99}{472.886}{a}%
\TextAt{8}{100.745}{472.886}{family}%
\TextAt{8}{136.515}{472.886}{of}%
\TextAt{8}{149.845}{472.886}{continuous}%
\TextAt{8}{208.45}{472.886}{homomorphisms}%
\TextAt{8}{295.734}{472.886}{such}%
\TextAt{8}{322.131}{472.886}{that}%
\TextAt{13}{215.611}{443.885}{h}%
\TextAt{21}{222.35}{442.091}{i}%
\TextAt{9}{229.052}{443.885}{\Character{61}}%
\TextAt{13}{241.477}{443.885}{g}%
\TextAt{21}{247.083}{442.091}{ij}%
\TextAt{16}{257.005}{443.885}{\Character{14}}%
\TextAt{13}{265.639}{443.885}{h}%
\TextAt{21}{272.377}{442.091}{j}%
\TextAt{13}{276.76}{443.885}{\Character{59}}%
\TextAt{8}{305.914}{443.885}{resp.}%
\TextAt{13}{335.236}{443.885}{h}%
\TextAt{21}{341.975}{442.091}{i}%
\TextAt{9}{348.677}{443.885}{\Character{61}}%
\TextAt{13}{361.102}{443.885}{h}%
\TextAt{21}{367.841}{442.091}{j}%
\TextAt{16}{374.88}{443.885}{\Character{14}}%
\TextAt{13}{383.514}{443.885}{g}%
\TextAt{21}{389.119}{442.091}{ij}%
\TextAt{8}{79.226}{414.884}{for}%
\TextAt{13}{97.111}{414.884}{i}%
\NativeMathAt{104.425}{414.884}{11.955}{\leq}%
\TextAt{13}{117.045}{414.884}{j}%
\TextAt{8}{122.56}{414.884}{,}%
\TextAt{8}{129.709}{414.884}{then}%
\TextAt{8}{156.369}{414.884}{there}%
\TextAt{8}{186.281}{414.884}{exists}%
\TextAt{8}{218.596}{414.884}{a}%
\TextAt{8}{228.351}{414.884}{unique}%
\TextAt{8}{266.393}{414.884}{continuous}%
\TextAt{8}{324.997}{414.884}{homomorphism}%
\TextAt{13}{167.756}{385.882}{h}%
\TextAt{9}{177.815}{385.882}{\Character{58}}%
\TextAt{13}{184.388}{385.882}{H}%
\NativeArrowAt{198.331}{385.882}{11.955}{19.26}{\longrightarrow}%
\TextAt{13}{220.913}{385.882}{G}%
\TextAt{9}{233.468}{385.882}{\Character{61}}%
\NativeMathAt{245.893}{385.882}{11.955}{\varprojlim}%
\TextAt{21}{252.581}{373.68}{i}%
\TextAt{13}{264.145}{385.882}{G}%
\TextAt{21}{273.378}{384.089}{i}%
\TextAt{13}{276.76}{385.882}{\Character{59}}%
\TextAt{8}{305.914}{385.882}{resp.}%
\TextAt{13}{335.236}{385.882}{h}%
\TextAt{9}{345.295}{385.882}{\Character{58}}%
\TextAt{13}{351.868}{385.882}{G}%
\TextAt{9}{364.422}{385.882}{\Character{61}}%
\NativeMathAt{376.848}{385.882}{11.955}{\varinjlim}%
\TextAt{21}{383.536}{373.68}{i}%
\TextAt{13}{395.099}{385.882}{G}%
\TextAt{21}{404.333}{384.089}{i}%
\NativeArrowAt{411.035}{385.882}{11.955}{19.26}{\longrightarrow}%
\TextAt{13}{433.617}{385.882}{H}%
\TextAt{8}{79.226}{351.116}{satisfying}%
\TextAt{13}{131.387}{351.116}{h}%
\TextAt{21}{138.125}{349.323}{i}%
\TextAt{9}{144.827}{351.116}{\Character{61}}%
\TextAt{13}{157.253}{351.116}{g}%
\TextAt{21}{162.858}{349.323}{i}%
\TextAt{16}{168.896}{351.116}{\Character{14}}%
\TextAt{13}{177.53}{351.116}{h}%
\TextAt{8}{184.269}{351.116}{,}%
\TextAt{8}{191.418}{351.116}{resp.}%
\TextAt{13}{220.74}{351.116}{h}%
\TextAt{21}{227.479}{349.323}{i}%
\TextAt{9}{234.181}{351.116}{\Character{61}}%
\TextAt{13}{246.606}{351.116}{h}%
\TextAt{16}{256.001}{351.116}{\Character{14}}%
\TextAt{13}{264.636}{351.116}{g}%
\TextAt{21}{270.241}{349.323}{i}%
\TextAt{8}{273.622}{351.116}{,}%
\TextAt{8}{280.771}{351.116}{for}%
\TextAt{8}{298.656}{351.116}{all}%
\TextAt{13}{314.915}{351.116}{i}%
\NativeMathAt{322.23}{351.116}{11.955}{\in}%
\TextAt{13}{333.521}{351.116}{I}%
\TextAt{8}{339.623}{351.116}{.}%
\TextAt{8}{97.159}{325.104}{The}%
\TextAt{8}{122.808}{325.104}{easy}%
\TextAt{8}{150.154}{325.104}{proof}%
\TextAt{8}{182.318}{325.104}{is}%
\TextAt{8}{195.677}{325.104}{left}%
\TextAt{8}{217.751}{325.104}{to}%
\TextAt{8}{233.656}{325.104}{the}%
\TextAt{8}{255.408}{325.104}{reader.}%
\TextAt{8}{300.513}{325.104}{A}%
\TextAt{12}{314.779}{325.104}{morphism}%
\TextAt{8}{378.923}{325.104}{between}%
\TextAt{8}{425.709}{325.104}{two}%
\TextAt{8}{449.421}{325.104}{projective,}%
\TextAt{8}{508.644}{325.104}{resp.}%
\TextAt{8}{79.226}{308.058}{inductive}%
\TextAt{8}{130.576}{308.058}{systems}%
\TextAt{16}{174.633}{308.058}{\Character{102}}%
\TextAt{13}{180.61}{308.058}{G}%
\TextAt{21}{189.844}{306.265}{i}%
\TextAt{13}{193.225}{308.058}{\Character{59}}%
\TextAt{13}{198.449}{308.058}{g}%
\TextAt{21}{204.075}{306.265}{ij}%
\TextAt{16}{211.34}{308.058}{\Character{103}}%
\TextAt{8}{221.839}{308.058}{and}%
\TextAt{16}{245.225}{308.058}{\Character{102}}%
\TextAt{13}{251.203}{308.058}{G}%
\TextAt{3}{260.437}{312.396}{\Character{48}}%
\TextAt{21}{260.437}{305.102}{i}%
\TextAt{13}{263.818}{308.058}{\Character{59}}%
\TextAt{13}{269.042}{308.058}{g}%
\TextAt{3}{275.096}{312.396}{\Character{48}}%
\TextAt{21}{274.667}{305.102}{ij}%
\TextAt{16}{281.933}{308.058}{\Character{103}}%
\TextAt{8}{292.431}{308.058}{of}%
\TextAt{8}{306.385}{308.058}{topological}%
\TextAt{8}{366.534}{308.058}{groups}%
\TextAt{8}{404.948}{308.058}{is}%
\TextAt{8}{417.336}{308.058}{a}%
\TextAt{8}{427.715}{308.058}{family}%
\TextAt{8}{464.109}{308.058}{of}%
\TextAt{8}{478.063}{308.058}{continuous}%
\TextAt{8}{79.226}{291.012}{homomorphisms}%
\TextAt{13}{166.511}{291.012}{f}%
\TextAt{21}{172.281}{289.219}{i}%
\TextAt{9}{178.983}{291.012}{\Character{58}}%
\TextAt{13}{185.555}{291.012}{G}%
\TextAt{21}{194.789}{289.219}{i}%
\NativeMathAt{201.491}{291.012}{11.955}{\rightarrow}%
\TextAt{13}{216.767}{291.012}{G}%
\TextAt{3}{226.001}{295.35}{\Character{48}}%
\TextAt{21}{226.001}{288.056}{i}%
\TextAt{8}{229.382}{291.012}{,}%
\TextAt{13}{236.531}{291.012}{i}%
\NativeMathAt{243.846}{291.012}{11.955}{\in}%
\TextAt{13}{255.136}{291.012}{I}%
\TextAt{8}{261.239}{291.012}{,}%
\TextAt{8}{268.388}{291.012}{such}%
\TextAt{8}{294.785}{291.012}{that}%
\TextAt{8}{320.154}{291.012}{the}%
\TextAt{8}{340.311}{291.012}{diagrams}%
\TextAt{13}{194.615}{258.812}{G}%
\TextAt{21}{203.848}{257.019}{j}%
\TextAt{13}{247.228}{258.812}{G}%
\TextAt{3}{256.462}{263.15}{\Character{48}}%
\TextAt{21}{256.462}{255.856}{j}%
\TextAt{13}{195.115}{215.261}{G}%
\TextAt{21}{204.348}{213.468}{i}%
\TextAt{13}{247.728}{215.261}{G}%
\TextAt{3}{256.962}{219.6}{\Character{48}}%
\TextAt{21}{256.962}{212.306}{i}%
\NativeDiagram{%
\path[draw=black,line width=0.478bp] (141.583,-512.199) -- (169.399,-512.199);%
\path[draw=black,line width=0.478bp,line cap=round,line join=round] (167.151,-509.33) .. controls (167.606,-511.051) and (168.618,-511.864) .. (169.638,-512.199) .. controls (168.618,-512.534) and (167.606,-513.347) .. (167.151,-515.068);%
}%
\TextAt{21}{223.745}{266.994}{f}%
\TextAt{22}{227.813}{265.779}{j}%
\NativeDiagram{%
\path[draw=black,line width=0.478bp] (129.423,-523.11) -- (129.423,-545.512);%
\path[draw=black,line width=0.478bp,line cap=round,line join=round] (132.292,-543.265) .. controls (130.571,-543.719) and (129.758,-544.731) .. (129.423,-545.751) .. controls (129.088,-544.731) and (128.275,-543.719) .. (126.554,-543.265);%
}%
\TextAt{21}{188.011}{238.923}{g}%
\TextAt{22}{192.041}{237.708}{ij}%
\NativeDiagram{%
\path[draw=black,line width=0.478bp] (182.037,-524.272) -- (182.037,-544.915);%
\path[draw=black,line width=0.478bp,line cap=round,line join=round] (184.906,-542.668) .. controls (183.185,-543.122) and (182.372,-544.134) .. (182.037,-545.154) .. controls (181.702,-544.134) and (180.889,-543.122) .. (179.168,-542.668);%
}%
\TextAt{21}{256.853}{237.972}{g}%
\TextAt{23}{261.169}{240.785}{\Character{48}}%
\TextAt{22}{260.883}{235.39}{ij}%
\NativeDiagram{%
\path[draw=black,line width=0.478bp] (141.082,-555.75) -- (169.899,-555.75);%
\path[draw=black,line width=0.478bp,line cap=round,line join=round] (167.652,-552.881) .. controls (168.106,-554.602) and (169.118,-555.415) .. (170.138,-555.75) .. controls (169.118,-556.085) and (168.106,-556.898) .. (167.652,-558.619);%
}%
\TextAt{21}{224.115}{209.9}{f}%
\TextAt{22}{228.183}{208.685}{i}%
\TextAt{8}{293.935}{239.942}{resp.}%
\TextAt{13}{351.15}{258.812}{G}%
\TextAt{21}{360.384}{257.019}{j}%
\TextAt{13}{403.764}{258.812}{G}%
\TextAt{3}{412.997}{263.15}{\Character{48}}%
\TextAt{21}{412.997}{255.856}{j}%
\TextAt{13}{351.65}{215.261}{G}%
\TextAt{21}{360.884}{213.468}{i}%
\TextAt{13}{404.264}{215.261}{G}%
\TextAt{3}{413.497}{219.6}{\Character{48}}%
\TextAt{21}{413.497}{212.306}{i}%
\NativeDiagram{%
\path[draw=black,line width=0.478bp] (298.119,-512.199) -- (325.935,-512.199);%
\path[draw=black,line width=0.478bp,line cap=round,line join=round] (323.687,-509.33) .. controls (324.142,-511.051) and (325.154,-511.864) .. (326.174,-512.199) .. controls (325.154,-512.534) and (324.142,-513.347) .. (323.687,-515.068);%
}%
\TextAt{21}{380.281}{266.994}{f}%
\TextAt{22}{384.349}{265.779}{j}%
\NativeDiagram{%
\path[draw=black,line width=0.478bp] (285.959,-545.99) -- (285.959,-523.588);%
\path[draw=black,line width=0.478bp,line cap=round,line join=round] (283.09,-525.836) .. controls (284.811,-525.381) and (285.624,-524.369) .. (285.959,-523.349) .. controls (286.294,-524.369) and (287.107,-525.381) .. (288.828,-525.836);%
}%
\TextAt{21}{344.547}{238.923}{g}%
\TextAt{22}{348.577}{237.708}{ij}%
\NativeDiagram{%
\path[draw=black,line width=0.478bp] (297.618,-555.75) -- (326.435,-555.75);%
\path[draw=black,line width=0.478bp,line cap=round,line join=round] (324.188,-552.881) .. controls (324.642,-554.602) and (325.655,-555.415) .. (326.675,-555.75) .. controls (325.655,-556.085) and (324.642,-556.898) .. (324.188,-558.619);%
}%
\TextAt{21}{380.651}{209.9}{f}%
\TextAt{22}{384.719}{208.685}{i}%
\NativeDiagram{%
\path[draw=black,line width=0.478bp] (338.573,-545.393) -- (338.573,-524.751);%
\path[draw=black,line width=0.478bp,line cap=round,line join=round] (335.704,-526.998) .. controls (337.425,-526.544) and (338.238,-525.531) .. (338.573,-524.511) .. controls (338.908,-525.531) and (339.721,-526.544) .. (341.442,-526.998);%
}%
\TextAt{21}{413.389}{237.972}{g}%
\TextAt{23}{417.704}{240.785}{\Character{48}}%
\TextAt{22}{417.419}{235.39}{ij}%
\TextAt{8}{79.226}{183.856}{commute}%
\TextAt{8}{129.629}{183.856}{for}%
\TextAt{13}{147.514}{183.856}{i}%
\NativeMathAt{154.828}{183.856}{11.955}{\leq}%
\TextAt{13}{167.448}{183.856}{j}%
\TextAt{8}{172.963}{183.856}{.}%
\TextAt{8}{181.411}{183.856}{Such}%
\TextAt{8}{209.697}{183.856}{a}%
\TextAt{8}{219.453}{183.856}{family}%
\TextAt{9}{255.223}{183.856}{\Character{40}}%
\TextAt{13}{259.775}{183.856}{f}%
\TextAt{21}{265.545}{182.063}{i}%
\TextAt{9}{268.926}{183.856}{\Character{41}}%
\TextAt{21}{273.479}{182.063}{i}%
\NativeMathAt{276.362}{182.063}{7.97}{\in}%
\TextAt{21}{282.007}{182.063}{I}%
\TextAt{8}{290.697}{183.856}{defines}%
\TextAt{8}{329.12}{183.856}{a}%
\TextAt{8}{338.876}{183.856}{mapping}%
\TextAt{13}{172.524}{154.855}{f}%
\TextAt{9}{182.891}{154.855}{\Character{58}}%
\TextAt{20}{189.464}{166.213}{\Character{89}}%
\TextAt{21}{190.691}{140.841}{i}%
\NativeMathAt{193.574}{140.841}{7.97}{\in}%
\TextAt{21}{199.219}{140.841}{I}%
\TextAt{13}{206.732}{154.855}{G}%
\TextAt{21}{215.966}{153.062}{i}%
\NativeArrowAt{222.668}{154.855}{11.955}{19.26}{\longrightarrow}%
\TextAt{20}{245.25}{166.213}{\Character{89}}%
\TextAt{21}{246.477}{140.841}{i}%
\NativeMathAt{249.36}{140.841}{7.97}{\in}%
\TextAt{21}{255.006}{140.841}{I}%
\TextAt{13}{262.519}{154.855}{G}%
\TextAt{3}{271.752}{159.791}{\Character{48}}%
\TextAt{21}{271.752}{151.9}{i}%
\TextAt{13}{275.134}{154.855}{\Character{59}}%
\TextAt{8}{304.288}{154.855}{resp.}%
\TextAt{13}{333.61}{154.855}{f}%
\TextAt{9}{343.978}{154.855}{\Character{58}}%
\TextAt{20}{350.55}{166.213}{\Character{97}}%
\TextAt{21}{351.777}{140.841}{i}%
\NativeMathAt{354.66}{140.841}{7.97}{\in}%
\TextAt{21}{360.306}{140.841}{I}%
\TextAt{13}{367.819}{154.855}{G}%
\TextAt{21}{377.052}{153.062}{i}%
\NativeArrowAt{383.754}{154.855}{11.955}{19.26}{\longrightarrow}%
\TextAt{20}{406.336}{166.213}{\Character{97}}%
\TextAt{21}{407.563}{140.841}{i}%
\NativeMathAt{410.446}{140.841}{7.97}{\in}%
\TextAt{21}{416.092}{140.841}{I}%
\TextAt{13}{423.605}{154.855}{G}%
\TextAt{3}{432.839}{159.791}{\Character{48}}%
\TextAt{21}{432.839}{151.9}{i}%
\TextAt{13}{436.22}{154.855}{\Character{59}}%
\TextAt{8}{79.226}{117.894}{which}%
\TextAt{8}{112.713}{117.894}{induces}%
\TextAt{8}{154.388}{117.894}{a}%
\TextAt{8}{164.144}{117.894}{homomorphism}%
\TextAt{13}{170.558}{88.893}{f}%
\TextAt{9}{180.925}{88.893}{\Character{58}}%
\NativeMathAt{187.498}{88.893}{11.955}{\varprojlim}%
\TextAt{21}{189.216}{76.513}{i}%
\NativeMathAt{192.099}{76.513}{7.97}{\in}%
\TextAt{21}{197.745}{76.513}{I}%
\TextAt{13}{205.749}{88.893}{G}%
\TextAt{21}{214.983}{87.1}{i}%
\NativeArrowAt{221.685}{88.893}{11.955}{19.26}{\longrightarrow}%
\NativeMathAt{244.267}{88.893}{11.955}{\varprojlim}%
\TextAt{21}{245.986}{76.513}{i}%
\NativeMathAt{248.869}{76.513}{7.97}{\in}%
\TextAt{21}{254.514}{76.513}{I}%
\TextAt{13}{262.519}{88.893}{G}%
\TextAt{3}{271.752}{93.829}{\Character{48}}%
\TextAt{21}{271.752}{85.938}{i}%
\TextAt{13}{275.134}{88.893}{\Character{59}}%
\TextAt{8}{304.288}{88.893}{resp.}%
\TextAt{13}{333.61}{88.893}{f}%
\TextAt{9}{343.978}{88.893}{\Character{58}}%
\NativeMathAt{350.55}{88.893}{11.955}{\varinjlim}%
\TextAt{21}{352.268}{76.513}{i}%
\NativeMathAt{355.151}{76.513}{7.97}{\in}%
\TextAt{21}{360.797}{76.513}{I}%
\TextAt{13}{368.802}{88.893}{G}%
\TextAt{21}{378.035}{87.1}{i}%
\NativeArrowAt{384.737}{88.893}{11.955}{19.26}{\longrightarrow}%
\NativeMathAt{407.319}{88.893}{11.955}{\varinjlim}%
\TextAt{21}{409.038}{76.513}{i}%
\NativeMathAt{411.921}{76.513}{7.97}{\in}%
\TextAt{21}{417.566}{76.513}{I}%
\TextAt{13}{425.571}{88.893}{G}%
\TextAt{3}{434.804}{93.829}{\Character{48}}%
\TextAt{21}{434.804}{85.938}{i}%
\TextAt{13}{438.186}{88.893}{\Character{58}}%
\TextAt{8}{297.211}{43.033}{245}%
\Destination{78.23}{809.815}{nav.page.245}%
\Destination{79.36}{429.855}{nav.print.268}%
\Destination{79.36}{544.855}{nav.number.4.2.5}%
\Destination{79.226}{703.075}{errata.107}%
}

% Page 258 (body page 246)
\Page{%
\TextAt{2}{79.226}{796.943}{IV.}%
\TextAt{2}{98.012}{796.943}{Abstract}%
\TextAt{2}{143.252}{796.943}{Class}%
\TextAt{2}{171.844}{796.943}{Field}%
\TextAt{2}{199.575}{796.943}{Theory}%
\Vector{q 0 G 0 g 0.299 w 7.226 18.729 m 460.769 18.729 l S Q}%
\TextAt{8}{79.226}{761.124}{In}%
\TextAt{8}{94.024}{761.124}{this}%
\TextAt{8}{117.022}{761.124}{way}%
\NativeMathAt{140.941}{761.124}{11.955}{\varprojlim}%
\TextAt{8}{157.2}{761.124}{,}%
\TextAt{8}{164.57}{761.124}{resp.}%
\NativeMathAt{194.421}{761.124}{11.955}{\varinjlim}%
\TextAt{8}{210.68}{761.124}{,}%
\TextAt{8}{218.049}{761.124}{becomes}%
\TextAt{8}{264.779}{761.124}{a}%
\TextAt{8}{274.71}{761.124}{functor.}%
\TextAt{8}{320.438}{761.124}{A}%
\TextAt{8}{333.286}{761.124}{particularly}%
\TextAt{8}{396.884}{761.124}{important}%
\TextAt{8}{452.353}{761.124}{property}%
\TextAt{8}{500.338}{761.124}{of}%
\TextAt{8}{513.844}{761.124}{this}%
\TextAt{8}{79.226}{744.078}{functor}%
\TextAt{8}{120.304}{744.078}{is}%
\TextAt{8}{132.497}{744.078}{its}%
\TextAt{8}{149.246}{744.078}{so-called}%
\TextAt{8}{197.209}{744.078}{“exactness”.}%
\TextAt{8}{264.296}{744.078}{For}%
\TextAt{8}{285.695}{744.078}{the}%
\TextAt{8}{306.282}{744.078}{inductive}%
\TextAt{8}{357.437}{744.078}{limit}%
\NativeMathAt{385.83}{744.078}{11.955}{\varinjlim}%
\TextAt{8}{402.089}{744.078}{,}%
\TextAt{8}{409.776}{744.078}{exactness}%
\TextAt{8}{462.031}{744.078}{holds}%
\TextAt{8}{493.09}{744.078}{without}%
\TextAt{8}{79.226}{727.032}{restrictions.}%
\TextAt{8}{144.39}{727.032}{In}%
\TextAt{8}{159.011}{727.032}{other}%
\TextAt{8}{189.58}{727.032}{words,}%
\TextAt{8}{226.39}{727.032}{one}%
\TextAt{8}{247.85}{727.032}{has}%
\TextAt{8}{268.724}{727.032}{the}%
\TextAt{12}{79.226}{702.041}{Proposition}%
\TextAt{12}{154.132}{702.041}{(2.6).}%
\TextAt{8}{191.355}{702.041}{Let}%
\TextAt{13}{214.232}{702.041}{\Character{11}}%
\TextAt{9}{228.325}{702.041}{\Character{58}}%
\TextAt{16}{238.149}{702.041}{\Character{102}}%
\TextAt{13}{244.126}{702.041}{G}%
\TextAt{3}{253.36}{706.38}{\Character{48}}%
\TextAt{21}{253.36}{699.086}{i}%
\TextAt{13}{256.741}{702.041}{\Character{59}}%
\TextAt{13}{261.965}{702.041}{g}%
\TextAt{3}{268.02}{706.38}{\Character{48}}%
\TextAt{21}{267.591}{699.086}{ij}%
\TextAt{16}{274.856}{702.041}{\Character{103}}%
\NativeMathAt{287.397}{702.041}{11.955}{\rightarrow}%
\TextAt{16}{305.928}{702.041}{\Character{102}}%
\TextAt{13}{311.91}{702.041}{G}%
\TextAt{21}{321.143}{700.248}{i}%
\TextAt{13}{324.525}{702.041}{\Character{59}}%
\TextAt{13}{329.749}{702.041}{g}%
\TextAt{21}{335.374}{700.248}{ij}%
\TextAt{16}{342.639}{702.041}{\Character{103}}%
\TextAt{8}{354.422}{702.041}{and}%
\TextAt{13}{379.092}{702.041}{\Character{12}}%
\TextAt{9}{392.935}{702.041}{\Character{58}}%
\TextAt{16}{402.759}{702.041}{\Character{102}}%
\TextAt{13}{408.736}{702.041}{G}%
\TextAt{21}{417.97}{700.248}{i}%
\TextAt{13}{421.351}{702.041}{\Character{59}}%
\TextAt{13}{426.575}{702.041}{g}%
\TextAt{21}{432.201}{700.248}{ij}%
\TextAt{16}{439.466}{702.041}{\Character{103}}%
\NativeMathAt{452.007}{702.041}{11.955}{\rightarrow}%
\TextAt{16}{470.538}{702.041}{\Character{102}}%
\TextAt{13}{476.52}{702.041}{G}%
\TextAt{3}{485.753}{706.38}{\Character{48}\Character{48}}%
\TextAt{21}{485.753}{699.086}{i}%
\TextAt{13}{490.845}{702.041}{\Character{59}}%
\TextAt{13}{496.069}{702.041}{g}%
\TextAt{3}{502.124}{706.38}{\Character{48}\Character{48}}%
\TextAt{21}{501.695}{699.086}{ij}%
\TextAt{16}{508.96}{702.041}{\Character{103}}%
\TextAt{8}{520.743}{702.041}{be}%
\TextAt{8}{79.226}{684.996}{morphisms}%
\TextAt{8}{138.536}{684.996}{between}%
\TextAt{8}{183.726}{684.996}{inductive}%
\TextAt{8}{234.452}{684.996}{systems}%
\TextAt{8}{277.885}{684.996}{of}%
\TextAt{8}{291.215}{684.996}{topological}%
\TextAt{8}{350.74}{684.996}{groups}%
\TextAt{8}{388.53}{684.996}{such}%
\TextAt{8}{414.927}{684.996}{that}%
\TextAt{8}{440.296}{684.996}{the}%
\TextAt{8}{460.453}{684.996}{sequence}%
\TextAt{13}{264.735}{659.572}{G}%
\TextAt{3}{273.968}{664.508}{\Character{48}}%
\TextAt{21}{273.968}{656.616}{i}%
\TextAt{21}{283.258}{666.501}{\Character{11}}%
\TextAt{22}{288.697}{665.286}{i}%
\NativeArrowAt{280.671}{659.572}{11.955}{15.183}{\longrightarrow}%
\TextAt{13}{299.179}{659.572}{G}%
\TextAt{21}{308.412}{657.778}{i}%
\TextAt{21}{317.702}{666.836}{\Character{12}}%
\TextAt{22}{322.454}{665.621}{i}%
\NativeArrowAt{315.114}{659.572}{11.955}{14.502}{\longrightarrow}%
\TextAt{13}{332.935}{659.572}{G}%
\TextAt{3}{342.169}{664.508}{\Character{48}\Character{48}}%
\TextAt{21}{342.169}{656.616}{i}%
\TextAt{8}{79.226}{634.148}{is}%
\TextAt{8}{90.99}{634.148}{exact}%
\TextAt{8}{121.882}{634.148}{for}%
\TextAt{8}{139.767}{634.148}{every}%
\TextAt{13}{170.659}{634.148}{i}%
\NativeMathAt{177.974}{634.148}{11.955}{\in}%
\TextAt{13}{189.265}{634.148}{I}%
\TextAt{8}{195.367}{634.148}{.}%
\TextAt{8}{203.816}{634.148}{Then}%
\TextAt{8}{234.373}{634.148}{the}%
\TextAt{8}{254.529}{634.148}{induced}%
\TextAt{8}{298.094}{634.148}{sequence}%
\NativeMathAt{240.211}{608.724}{11.955}{\varinjlim}%
\TextAt{21}{241.929}{596.344}{i}%
\NativeMathAt{244.812}{596.344}{7.97}{\in}%
\TextAt{21}{250.458}{596.344}{I}%
\TextAt{13}{258.462}{608.724}{G}%
\TextAt{3}{267.696}{613.66}{\Character{48}}%
\TextAt{21}{267.696}{605.768}{i}%
\TextAt{21}{276.986}{615.501}{\Character{11}}%
\NativeArrowAt{274.398}{608.724}{11.955}{12.027}{\longrightarrow}%
\NativeMathAt{289.744}{608.724}{11.955}{\varinjlim}%
\TextAt{21}{291.463}{596.344}{i}%
\NativeMathAt{294.346}{596.344}{7.97}{\in}%
\TextAt{21}{299.991}{596.344}{I}%
\TextAt{13}{307.996}{608.724}{G}%
\TextAt{21}{317.229}{606.931}{i}%
\TextAt{21}{326.623}{615.988}{\Character{12}}%
\NativeArrowAt{323.932}{608.724}{11.955}{11.955}{\longrightarrow}%
\NativeMathAt{339.208}{608.724}{11.955}{\varinjlim}%
\TextAt{21}{340.926}{596.344}{i}%
\NativeMathAt{343.809}{596.344}{7.97}{\in}%
\TextAt{21}{349.455}{596.344}{I}%
\TextAt{13}{357.459}{608.724}{G}%
\TextAt{3}{366.693}{613.66}{\Character{48}\Character{48}}%
\TextAt{21}{366.693}{605.768}{i}%
\TextAt{8}{79.226}{576.974}{is}%
\TextAt{8}{90.99}{576.974}{also}%
\TextAt{8}{114.47}{576.974}{exact.}%
\TextAt{19}{79.226}{544.04}{Proof.}%
\TextAt{8}{116}{544.04}{Let}%
\TextAt{13}{138.686}{544.04}{G}%
\TextAt{3}{147.92}{548.378}{\Character{48}}%
\TextAt{9}{156.961}{544.04}{\Character{61}}%
\NativeMathAt{172.311}{544.04}{11.955}{\varinjlim}%
\TextAt{21}{188.571}{538.6}{i}%
\TextAt{13}{193.944}{544.04}{G}%
\TextAt{3}{203.178}{548.378}{\Character{48}}%
\TextAt{21}{203.178}{541.084}{i}%
\TextAt{8}{206.559}{544.04}{,}%
\TextAt{13}{215.854}{544.04}{G}%
\TextAt{9}{231.334}{544.04}{\Character{61}}%
\NativeMathAt{246.684}{544.04}{11.955}{\varinjlim}%
\TextAt{21}{262.943}{538.6}{i}%
\TextAt{13}{268.317}{544.04}{G}%
\TextAt{21}{277.551}{542.246}{i}%
\TextAt{8}{280.932}{544.04}{,}%
\TextAt{13}{290.227}{544.04}{G}%
\TextAt{3}{299.461}{548.378}{\Character{48}\Character{48}}%
\TextAt{9}{310.799}{544.04}{\Character{61}}%
\NativeMathAt{326.149}{544.04}{11.955}{\varinjlim}%
\TextAt{21}{342.408}{538.6}{i}%
\TextAt{13}{347.782}{544.04}{G}%
\TextAt{3}{357.016}{548.378}{\Character{48}\Character{48}}%
\TextAt{21}{357.016}{541.084}{i}%
\TextAt{8}{362.108}{544.04}{.}%
\TextAt{8}{375.706}{544.04}{We}%
\TextAt{8}{397.567}{544.04}{consider}%
\TextAt{8}{444.868}{544.04}{the}%
\TextAt{8}{466.741}{544.04}{commutative}%
\TextAt{8}{79.226}{526.994}{diagram}%
\TextAt{13}{245.838}{508.081}{G}%
\TextAt{3}{255.072}{512.419}{\Character{48}}%
\TextAt{21}{255.072}{505.125}{i}%
\TextAt{13}{297.451}{508.081}{G}%
\TextAt{21}{306.684}{506.288}{i}%
\TextAt{13}{350.689}{508.081}{G}%
\TextAt{3}{359.923}{512.419}{\Character{48}\Character{48}}%
\TextAt{21}{359.923}{505.125}{i}%
\TextAt{13}{246.132}{466.08}{G}%
\TextAt{3}{255.366}{470.419}{\Character{48}}%
\TextAt{13}{299.142}{466.08}{G}%
\TextAt{13}{349.063}{466.08}{G}%
\TextAt{3}{358.297}{470.419}{\Character{48}\Character{48}}%
\TextAt{13}{363.389}{466.08}{\Character{58}}%
\NativeDiagram{%
\path[draw=black,line width=0.478bp] (191.806,-262.93) -- (219.622,-262.93);%
\path[draw=black,line width=0.478bp,line cap=round,line join=round] (217.375,-260.061) .. controls (217.829,-261.782) and (218.841,-262.595) .. (219.861,-262.93) .. controls (218.841,-263.264) and (217.829,-264.077) .. (217.375,-265.799);%
}%
\TextAt{21}{273.653}{515.101}{\Character{11}}%
\TextAt{22}{279.092}{513.886}{i}%
\NativeDiagram{%
\path[draw=black,line width=0.478bp] (180.147,-273.453) -- (180.147,-294.096);%
\path[draw=black,line width=0.478bp,line cap=round,line join=round] (183.016,-291.849) .. controls (181.294,-292.303) and (180.481,-293.315) .. (180.147,-294.335) .. controls (179.812,-293.315) and (178.999,-292.303) .. (177.277,-291.849);%
}%
\TextAt{21}{242.139}{488.21}{g}%
\TextAt{23}{246.455}{491.023}{\Character{48}}%
\TextAt{22}{246.169}{485.628}{i}%
\NativeDiagram{%
\path[draw=black,line width=0.478bp] (243.419,-262.93) -- (272.861,-262.93);%
\path[draw=black,line width=0.478bp,line cap=round,line join=round] (270.614,-260.061) .. controls (271.068,-261.782) and (272.08,-262.595) .. (273.1,-262.93) .. controls (272.08,-263.264) and (271.068,-264.077) .. (270.614,-265.799);%
}%
\TextAt{21}{326.423}{515.435}{\Character{12}}%
\TextAt{22}{331.175}{514.22}{i}%
\NativeDiagram{%
\path[draw=black,line width=0.478bp] (231.76,-272.291) -- (231.76,-294.693);%
\path[draw=black,line width=0.478bp,line cap=round,line join=round] (234.629,-292.446) .. controls (232.907,-292.9) and (232.095,-293.912) .. (231.76,-294.932) .. controls (231.425,-293.912) and (230.612,-292.9) .. (228.891,-292.446);%
}%
\TextAt{21}{306.575}{489.328}{g}%
\TextAt{22}{310.605}{488.113}{i}%
\NativeDiagram{%
\path[draw=black,line width=0.478bp] (285.854,-273.453) -- (285.854,-294.096);%
\path[draw=black,line width=0.478bp,line cap=round,line join=round] (288.723,-291.849) .. controls (287.002,-292.303) and (286.189,-293.315) .. (285.854,-294.335) .. controls (285.519,-293.315) and (284.707,-292.303) .. (282.985,-291.849);%
}%
\TextAt{21}{360.67}{488.21}{g}%
\TextAt{23}{364.986}{491.023}{\Character{48}\Character{48}}%
\TextAt{22}{364.7}{485.628}{i}%
\NativeDiagram{%
\path[draw=black,line width=0.478bp] (191.513,-304.931) -- (221.313,-304.931);%
\path[draw=black,line width=0.478bp,line cap=round,line join=round] (219.065,-302.062) .. controls (219.519,-303.783) and (220.532,-304.596) .. (221.552,-304.931) .. controls (220.532,-305.266) and (219.519,-306.079) .. (219.065,-307.8);%
}%
\TextAt{21}{275.932}{462.822}{\Character{11}}%
\NativeDiagram{%
\path[draw=black,line width=0.478bp] (241.729,-304.931) -- (271.235,-304.931);%
\path[draw=black,line width=0.478bp,line cap=round,line join=round] (268.988,-302.062) .. controls (269.442,-303.783) and (270.454,-304.596) .. (271.474,-304.931) .. controls (270.454,-305.266) and (269.442,-306.079) .. (268.988,-307.8);%
}%
\TextAt{21}{326.14}{460.719}{\Character{12}}%
\TextAt{8}{79.226}{441.688}{Let}%
\TextAt{13}{101.387}{441.688}{x}%
\NativeMathAt{113.39}{441.688}{11.955}{\in}%
\TextAt{13}{126.711}{441.688}{G}%
\TextAt{8}{141.033}{441.688}{be}%
\TextAt{8}{158.149}{441.688}{such}%
\TextAt{8}{185.737}{441.688}{that}%
\TextAt{13}{212.297}{441.688}{\Character{12}}%
\TextAt{9}{219.568}{441.688}{\Character{40}}%
\TextAt{13}{224.121}{441.688}{x}%
\TextAt{9}{230.773}{441.688}{\Character{41}}%
\TextAt{9}{240.66}{441.688}{\Character{61}}%
\TextAt{9}{255.114}{441.688}{1}%
\TextAt{8}{260.984}{441.688}{.}%
\TextAt{8}{273.007}{441.688}{Then}%
\TextAt{8}{304.755}{441.688}{there}%
\TextAt{8}{335.858}{441.688}{exists}%
\TextAt{8}{369.364}{441.688}{an}%
\TextAt{13}{386.815}{441.688}{i}%
\TextAt{8}{395.897}{441.688}{and}%
\TextAt{8}{419.851}{441.688}{an}%
\TextAt{13}{437.301}{441.688}{x}%
\TextAt{21}{443.953}{439.895}{i}%
\NativeMathAt{452.685}{441.688}{11.955}{\in}%
\TextAt{13}{466.006}{441.688}{G}%
\TextAt{21}{475.24}{439.895}{i}%
\TextAt{8}{483.71}{441.688}{such}%
\TextAt{8}{511.298}{441.688}{that}%
\TextAt{13}{79.226}{424.642}{g}%
\TextAt{21}{84.831}{422.849}{i}%
\TextAt{9}{88.213}{424.642}{\Character{40}}%
\TextAt{13}{92.765}{424.642}{x}%
\TextAt{21}{99.417}{422.849}{i}%
\TextAt{9}{102.798}{424.642}{\Character{41}}%
\TextAt{9}{110.653}{424.642}{\Character{61}}%
\TextAt{13}{123.097}{424.642}{x}%
\TextAt{8}{129.749}{424.642}{.}%
\TextAt{8}{138.197}{424.642}{As}%
\TextAt{13}{228.487}{407.596}{g}%
\TextAt{3}{234.521}{412.532}{\Character{48}\Character{48}}%
\TextAt{21}{234.092}{404.641}{i}%
\TextAt{13}{239.613}{407.596}{\Character{12}}%
\TextAt{21}{246.222}{405.803}{i}%
\TextAt{9}{249.603}{407.596}{\Character{40}}%
\TextAt{13}{254.155}{407.596}{x}%
\TextAt{21}{260.807}{405.803}{i}%
\TextAt{9}{264.189}{407.596}{\Character{41}}%
\TextAt{9}{272.044}{407.596}{\Character{61}}%
\TextAt{13}{284.487}{407.596}{\Character{12}}%
\TextAt{13}{291.756}{407.596}{g}%
\TextAt{21}{297.364}{405.803}{i}%
\TextAt{9}{300.745}{407.596}{\Character{40}}%
\TextAt{13}{305.298}{407.596}{x}%
\TextAt{21}{311.95}{405.803}{i}%
\TextAt{9}{315.331}{407.596}{\Character{41}}%
\TextAt{9}{323.186}{407.596}{\Character{61}}%
\TextAt{13}{335.63}{407.596}{\Character{12}}%
\TextAt{9}{342.901}{407.596}{\Character{40}}%
\TextAt{13}{347.453}{407.596}{x}%
\TextAt{9}{354.105}{407.596}{\Character{41}}%
\TextAt{9}{361.96}{407.596}{\Character{61}}%
\TextAt{9}{374.381}{407.596}{1}%
\TextAt{13}{380.257}{407.596}{\Character{59}}%
\TextAt{8}{79.226}{385.607}{there}%
\TextAt{8}{109.388}{385.607}{exists}%
\TextAt{13}{141.953}{385.607}{j}%
\NativeMathAt{151.216}{385.608}{11.955}{\geq}%
\TextAt{13}{164.261}{385.607}{i}%
\TextAt{8}{172.402}{385.607}{such}%
\TextAt{8}{199.05}{385.607}{that}%
\TextAt{13}{224.669}{385.607}{\Character{12}}%
\TextAt{21}{231.277}{383.814}{i}%
\TextAt{9}{234.658}{385.608}{\Character{40}}%
\TextAt{13}{239.211}{385.608}{x}%
\TextAt{21}{245.863}{383.814}{i}%
\TextAt{9}{249.244}{385.608}{\Character{41}}%
\TextAt{8}{257.944}{385.607}{equals}%
\TextAt{9}{293.701}{385.607}{1}%
\TextAt{8}{303.701}{385.607}{in}%
\TextAt{13}{317.604}{385.607}{G}%
\TextAt{3}{326.838}{389.946}{\Character{48}\Character{48}}%
\TextAt{21}{326.838}{382.652}{j}%
\TextAt{8}{331.93}{385.607}{.}%
\TextAt{8}{341.129}{385.607}{Changing}%
\TextAt{8}{394.066}{385.607}{notation,}%
\TextAt{8}{444.47}{385.607}{we}%
\TextAt{8}{461.948}{385.607}{may}%
\TextAt{8}{487.567}{385.607}{therefore}%
\TextAt{8}{79.226}{368.562}{assume}%
\TextAt{8}{121.096}{368.562}{that}%
\TextAt{13}{147.891}{368.562}{\Character{12}}%
\TextAt{21}{154.5}{366.768}{i}%
\TextAt{9}{157.881}{368.562}{\Character{40}}%
\TextAt{13}{162.433}{368.562}{x}%
\TextAt{21}{169.085}{366.768}{i}%
\TextAt{9}{172.467}{368.562}{\Character{41}}%
\TextAt{9}{182.749}{368.562}{\Character{61}}%
\TextAt{9}{197.609}{368.562}{1}%
\TextAt{8}{203.478}{368.562}{,}%
\TextAt{8}{212.409}{368.562}{so}%
\TextAt{8}{228.205}{368.562}{that}%
\TextAt{8}{255}{368.562}{there}%
\TextAt{8}{286.338}{368.562}{exists}%
\TextAt{13}{320.078}{368.562}{y}%
\TextAt{21}{325.786}{366.768}{i}%
\NativeMathAt{334.918}{368.562}{11.955}{\in}%
\TextAt{13}{348.639}{368.562}{G}%
\TextAt{3}{357.873}{372.9}{\Character{48}}%
\TextAt{21}{357.873}{365.606}{i}%
\TextAt{8}{366.577}{368.562}{such}%
\TextAt{8}{394.4}{368.562}{that}%
\TextAt{13}{421.195}{368.562}{\Character{11}}%
\TextAt{21}{428.64}{366.768}{i}%
\TextAt{9}{432.022}{368.562}{\Character{40}}%
\TextAt{13}{436.574}{368.562}{y}%
\TextAt{21}{442.282}{366.768}{i}%
\TextAt{9}{445.663}{368.562}{\Character{41}}%
\TextAt{9}{455.956}{368.562}{\Character{61}}%
\TextAt{13}{470.821}{368.562}{x}%
\TextAt{21}{477.473}{366.768}{i}%
\TextAt{8}{480.855}{368.562}{.}%
\TextAt{8}{493.58}{368.562}{Putting}%
\TextAt{13}{79.226}{351.516}{y}%
\TextAt{9}{88.684}{351.516}{\Character{61}}%
\TextAt{13}{101.109}{351.516}{g}%
\TextAt{3}{107.143}{355.854}{\Character{48}}%
\TextAt{21}{106.714}{348.56}{i}%
\TextAt{9}{110.096}{351.516}{\Character{40}}%
\TextAt{13}{114.648}{351.516}{y}%
\TextAt{21}{120.356}{349.722}{i}%
\TextAt{9}{123.737}{351.516}{\Character{41}}%
\TextAt{8}{128.289}{351.516}{,}%
\TextAt{8}{135.439}{351.516}{we}%
\TextAt{8}{152.666}{351.516}{have}%
\TextAt{13}{179.661}{351.516}{\Character{11}}%
\TextAt{9}{187.182}{351.516}{\Character{40}}%
\TextAt{13}{191.735}{351.516}{y}%
\TextAt{9}{197.871}{351.516}{\Character{41}}%
\TextAt{9}{205.726}{351.516}{\Character{61}}%
\TextAt{13}{218.17}{351.516}{x}%
\TextAt{8}{224.822}{351.516}{.}%
\Vector{q 0 G 0 g 0.399 w 452.334 -422.484 m 452.334 -414.415 l S Q}%
\Vector{q 0 G 0 g 0.399 w 452.534 -414.614 m 459.707 -414.614 l S Q}%
\Vector{q 0 G 0 g 0.399 w 452.534 -422.285 m 459.707 -422.285 l S Q}%
\Vector{q 0 G 0 g 0.399 w 459.906 -422.484 m 459.906 -414.415 l S Q}%
\ErrataAt{97.1}{323.536}{437.9}{The projective limit need not preserve exactness in general. It does preserve exactness}%
\ErrataAt{79.19}{306.49}{455.81}{for projective systems of compact Hausdorff groups, as expressed by the following.}%
\TextAt{12}{79.226}{281.5}{Proposition}%
\TextAt{12}{152.479}{281.5}{(2.7).}%
\TextAt{8}{188.265}{281.5}{Let}%
\TextAt{13}{209.704}{281.5}{\Character{11}}%
\TextAt{9}{221.348}{281.5}{\Character{58}}%
\TextAt{16}{228.722}{281.5}{\Character{102}}%
\TextAt{13}{234.7}{281.5}{G}%
\TextAt{3}{243.934}{285.839}{\Character{48}}%
\TextAt{21}{243.934}{278.545}{i}%
\TextAt{13}{247.315}{281.5}{\Character{59}}%
\TextAt{13}{252.539}{281.5}{g}%
\TextAt{3}{258.593}{285.839}{\Character{48}}%
\TextAt{21}{258.164}{278.545}{ij}%
\TextAt{16}{265.43}{281.5}{\Character{103}}%
\NativeMathAt{275.52}{281.5}{11.955}{\rightarrow}%
\TextAt{16}{291.6}{281.5}{\Character{102}}%
\TextAt{13}{297.585}{281.5}{G}%
\TextAt{21}{306.819}{279.707}{i}%
\TextAt{13}{310.2}{281.5}{\Character{59}}%
\TextAt{13}{315.424}{281.5}{g}%
\TextAt{21}{321.049}{279.707}{ij}%
\TextAt{16}{328.315}{281.5}{\Character{103}}%
\TextAt{8}{338.66}{281.5}{and}%
\TextAt{13}{361.893}{281.5}{\Character{12}}%
\TextAt{9}{373.287}{281.5}{\Character{58}}%
\TextAt{16}{380.661}{281.5}{\Character{102}}%
\TextAt{13}{386.638}{281.5}{G}%
\TextAt{21}{395.872}{279.707}{i}%
\TextAt{13}{399.253}{281.5}{\Character{59}}%
\TextAt{13}{404.477}{281.5}{g}%
\TextAt{21}{410.103}{279.707}{ij}%
\TextAt{16}{417.368}{281.5}{\Character{103}}%
\NativeMathAt{427.458}{281.5}{11.955}{\rightarrow}%
\TextAt{16}{443.538}{281.5}{\Character{102}}%
\TextAt{13}{449.523}{281.5}{G}%
\TextAt{3}{458.757}{285.839}{\Character{48}\Character{48}}%
\TextAt{21}{458.757}{278.545}{i}%
\TextAt{13}{463.849}{281.5}{\Character{59}}%
\TextAt{13}{469.073}{281.5}{g}%
\TextAt{3}{475.127}{285.839}{\Character{48}\Character{48}}%
\TextAt{21}{474.698}{278.545}{ij}%
\TextAt{16}{481.964}{281.5}{\Character{103}}%
\TextAt{8}{492.309}{281.5}{be}%
\TextAt{8}{508.704}{281.5}{mor-}%
\ErrataAt{79.19}{264.454}{455.81}{phisms between projective systems of compact Hausdorff groups such that the sequence}%
\TextAt{13}{264.735}{239.03}{G}%
\TextAt{3}{273.968}{243.967}{\Character{48}}%
\TextAt{21}{273.968}{236.075}{i}%
\TextAt{21}{283.258}{245.96}{\Character{11}}%
\TextAt{22}{288.697}{244.745}{i}%
\NativeArrowAt{280.671}{239.03}{11.955}{15.183}{\longrightarrow}%
\TextAt{13}{299.179}{239.03}{G}%
\TextAt{21}{308.412}{237.237}{i}%
\TextAt{21}{317.702}{246.295}{\Character{12}}%
\TextAt{22}{322.454}{245.08}{i}%
\NativeArrowAt{315.114}{239.03}{11.955}{14.502}{\longrightarrow}%
\TextAt{13}{332.935}{239.03}{G}%
\TextAt{3}{342.169}{243.967}{\Character{48}\Character{48}}%
\TextAt{21}{342.169}{236.075}{i}%
\TextAt{8}{79.226}{213.607}{is}%
\TextAt{8}{90.99}{213.607}{exact}%
\TextAt{8}{121.882}{213.607}{for}%
\TextAt{8}{139.767}{213.607}{every}%
\TextAt{13}{170.659}{213.607}{i}%
\NativeMathAt{177.974}{213.607}{11.955}{\in}%
\TextAt{13}{189.265}{213.607}{I}%
\TextAt{8}{195.367}{213.607}{.}%
\TextAt{8}{203.816}{213.607}{Then}%
\NativeMathAt{240.211}{188.183}{11.955}{\varprojlim}%
\TextAt{21}{246.899}{175.98}{i}%
\TextAt{13}{258.462}{188.183}{G}%
\TextAt{3}{267.696}{193.119}{\Character{48}}%
\TextAt{21}{267.696}{185.227}{i}%
\TextAt{21}{276.986}{194.96}{\Character{11}}%
\NativeArrowAt{274.398}{188.183}{11.955}{12.027}{\longrightarrow}%
\NativeMathAt{289.744}{188.183}{11.955}{\varprojlim}%
\TextAt{21}{296.432}{175.98}{i}%
\TextAt{13}{307.996}{188.183}{G}%
\TextAt{21}{317.229}{186.389}{i}%
\TextAt{21}{326.623}{195.447}{\Character{12}}%
\NativeArrowAt{323.932}{188.183}{11.955}{11.955}{\longrightarrow}%
\NativeMathAt{339.208}{188.183}{11.955}{\varprojlim}%
\TextAt{21}{345.896}{175.98}{i}%
\TextAt{13}{357.459}{188.183}{G}%
\TextAt{3}{366.693}{193.119}{\Character{48}\Character{48}}%
\TextAt{21}{366.693}{185.227}{i}%
\ErrataAt{79.19}{156.994}{455.81}{is again an exact sequence of compact Hausdorff groups.}%
\TextAt{19}{79.226}{124.059}{Proof.}%
\TextAt{8}{116}{124.059}{Let}%
\TextAt{13}{138.056}{124.059}{x}%
\TextAt{9}{149.881}{124.059}{\Character{61}}%
\TextAt{9}{164.144}{124.059}{\Character{40}}%
\TextAt{13}{168.71}{124.059}{x}%
\TextAt{21}{175.362}{122.266}{i}%
\TextAt{9}{178.744}{124.059}{\Character{41}}%
\TextAt{21}{183.296}{122.266}{i}%
\NativeMathAt{186.179}{122.266}{7.97}{\in}%
\TextAt{21}{191.825}{122.266}{I}%
\NativeMathAt{201.789}{124.059}{11.955}{\in}%
\NativeMathAt{214.931}{124.059}{11.955}{\varprojlim}%
\TextAt{21}{231.191}{118.62}{i}%
\TextAt{13}{236.564}{124.059}{G}%
\TextAt{21}{245.798}{122.266}{i}%
\TextAt{8}{254.163}{124.059}{and}%
\TextAt{13}{278.012}{124.059}{\Character{12}}%
\TextAt{9}{285.284}{124.059}{\Character{40}}%
\TextAt{13}{289.836}{124.059}{x}%
\TextAt{9}{296.488}{124.059}{\Character{41}}%
\TextAt{9}{306.196}{124.059}{\Character{61}}%
\TextAt{9}{320.47}{124.059}{1}%
\TextAt{8}{326.343}{124.059}{,}%
\TextAt{8}{334.85}{124.059}{so}%
\TextAt{8}{350.307}{124.059}{that}%
\TextAt{13}{376.762}{124.059}{\Character{12}}%
\TextAt{21}{383.371}{122.266}{i}%
\TextAt{9}{386.752}{124.059}{\Character{40}}%
\TextAt{13}{391.305}{124.059}{x}%
\TextAt{21}{397.957}{122.266}{i}%
\TextAt{9}{401.338}{124.059}{\Character{41}}%
\TextAt{9}{411.046}{124.059}{\Character{61}}%
\TextAt{9}{425.32}{124.059}{1}%
\TextAt{8}{436.177}{124.059}{for}%
\TextAt{8}{455.148}{124.059}{all}%
\TextAt{13}{472.494}{124.059}{i}%
\NativeMathAt{481.66}{124.059}{11.955}{\in}%
\TextAt{13}{494.802}{124.059}{I}%
\TextAt{8}{500.905}{124.059}{.}%
\TextAt{8}{512.613}{124.059}{The}%
\TextAt{8}{79.226}{107.013}{preimages}%
\TextAt{13}{135.154}{107.013}{Y}%
\TextAt{21}{141.946}{105.22}{i}%
\TextAt{9}{150.748}{107.013}{\Character{61}}%
\TextAt{13}{165.274}{107.013}{\Character{11}}%
\TextAt{3}{172.796}{111.961}{\Character{0}}%
\TextAt{17}{179.383}{111.961}{1}%
\TextAt{21}{172.72}{103.949}{i}%
\TextAt{9}{184.115}{107.013}{\Character{40}}%
\TextAt{13}{188.667}{107.013}{x}%
\TextAt{21}{195.319}{105.22}{i}%
\TextAt{9}{198.701}{107.013}{\Character{41}}%
\NativeMathAt{208.674}{107.013}{11.955}{\subseteq}%
\TextAt{13}{223.394}{107.013}{G}%
\TextAt{3}{232.628}{111.352}{\Character{48}}%
\TextAt{21}{232.628}{104.058}{i}%
\TextAt{8}{241.139}{107.013}{then}%
\TextAt{8}{269.032}{107.013}{form}%
\TextAt{8}{297.905}{107.013}{a}%
\TextAt{8}{308.893}{107.013}{projective}%
\TextAt{8}{364.103}{107.013}{system}%
\TextAt{8}{404.155}{107.013}{of}%
\TextAt{8}{418.717}{107.013}{nonempty}%
\TextAt{8}{474.585}{107.013}{closed,}%
\TextAt{8}{513.904}{107.013}{and}%
\TextAt{8}{79.226}{89.967}{hence}%
\TextAt{8}{112.937}{89.967}{compact}%
\TextAt{8}{160.97}{89.967}{subsets}%
\TextAt{8}{202.679}{89.967}{of}%
\TextAt{8}{217.213}{89.967}{the}%
\TextAt{13}{238.574}{89.967}{G}%
\TextAt{3}{247.808}{94.306}{\Character{48}}%
\TextAt{21}{247.808}{87.012}{i}%
\TextAt{8}{251.189}{89.967}{.}%
\TextAt{8}{263.252}{89.967}{By}%
\TextAt{8}{282.819}{89.967}{(2.3),}%
\TextAt{8}{315.552}{89.967}{this}%
\TextAt{8}{339.579}{89.967}{means}%
\TextAt{8}{376.613}{89.967}{that}%
\TextAt{8}{403.187}{89.967}{the}%
\TextAt{8}{424.548}{89.967}{projective}%
\TextAt{8}{479.73}{89.967}{limit}%
\TextAt{13}{508.898}{89.967}{Y}%
\TextAt{9}{523.665}{89.967}{\Character{61}}%
\NativeMathAt{79.226}{72.921}{11.955}{\varprojlim}%
\TextAt{21}{95.485}{67.482}{i}%
\TextAt{13}{100.859}{72.921}{Y}%
\TextAt{21}{107.651}{71.128}{i}%
\NativeMathAt{114.353}{72.921}{11.955}{\subseteq}%
\NativeMathAt{126.972}{72.921}{11.955}{\varprojlim}%
\TextAt{21}{143.231}{67.482}{i}%
\TextAt{13}{148.605}{72.921}{G}%
\TextAt{3}{157.839}{77.26}{\Character{48}}%
\TextAt{21}{157.839}{69.966}{i}%
\TextAt{8}{165.117}{72.921}{is}%
\TextAt{8}{176.881}{72.921}{nonempty,}%
\TextAt{8}{233.788}{72.921}{and}%
\TextAt{13}{256.55}{72.921}{\Character{11}}%
\TextAt{8}{267.97}{72.921}{maps}%
\TextAt{8}{298.599}{72.921}{every}%
\TextAt{8}{329.491}{72.921}{element}%
\TextAt{13}{372.733}{72.921}{y}%
\NativeMathAt{382.19}{72.921}{11.955}{\in}%
\TextAt{13}{393.481}{72.921}{Y}%
\TextAt{8}{406.772}{72.921}{to}%
\TextAt{13}{421.082}{72.921}{x}%
\TextAt{8}{427.734}{72.921}{.}%
\Vector{q 0 G 0 g 0.399 w 452.334 -701.079 m 452.334 -693.009 l S Q}%
\Vector{q 0 G 0 g 0.399 w 452.534 -693.208 m 459.707 -693.208 l S Q}%
\Vector{q 0 G 0 g 0.399 w 452.534 -700.879 m 459.707 -700.879 l S Q}%
\Vector{q 0 G 0 g 0.399 w 459.906 -701.079 m 459.906 -693.009 l S Q}%
\TextAt{8}{297.211}{43.033}{246}%
\Destination{78.23}{809.815}{nav.page.246}%
\Destination{79.36}{699.955}{nav.print.269}%
\Destination{79.36}{716.055}{nav.number.4.2.6}%
\Destination{79.36}{295.455}{nav.number.4.2.7}%
\Hotlink{282.319}{87.755}{25.078}{11.047}{nav.number.4.2.3}%
\Destination{79.226}{340.536}{errata.103}%
}

% Page 259 (body page 247)
\Page{%
\TextAt{2}{79.226}{796.943}{IV.}%
\TextAt{2}{98.012}{796.943}{Abstract}%
\TextAt{2}{143.252}{796.943}{Class}%
\TextAt{2}{171.844}{796.943}{Field}%
\TextAt{2}{199.575}{796.943}{Theory}%
\Vector{q 0 G 0 g 0.299 w 7.226 18.729 m 460.769 18.729 l S Q}%
\TextAt{8}{97.159}{761.124}{Now}%
\TextAt{8}{125.145}{761.124}{that}%
\TextAt{8}{151.841}{761.124}{we}%
\TextAt{8}{170.395}{761.124}{have}%
\TextAt{8}{198.716}{761.124}{at}%
\TextAt{8}{214.353}{761.124}{our}%
\TextAt{8}{236.493}{761.124}{disposal}%
\TextAt{8}{282.496}{761.124}{the}%
\TextAt{8}{303.979}{761.124}{notion}%
\TextAt{8}{341.733}{761.124}{of}%
\TextAt{8}{356.389}{761.124}{projective}%
\TextAt{8}{411.693}{761.124}{limit,}%
\TextAt{8}{444.566}{761.124}{we}%
\TextAt{8}{463.12}{761.124}{return}%
\TextAt{8}{500.216}{761.124}{to}%
\TextAt{8}{515.853}{761.124}{our}%
\TextAt{8}{79.226}{744.078}{starting}%
\TextAt{8}{122.054}{744.078}{point,}%
\TextAt{8}{155.219}{744.078}{the}%
\TextAt{8}{174.555}{744.078}{profinite}%
\TextAt{8}{220.563}{744.078}{groups.}%
\TextAt{8}{262.631}{744.078}{Recall}%
\TextAt{8}{297.09}{744.078}{that}%
\TextAt{8}{321.638}{744.078}{these}%
\TextAt{8}{350.789}{744.078}{are}%
\TextAt{8}{369.48}{744.078}{the}%
\TextAt{8}{388.816}{744.078}{topological}%
\TextAt{8}{447.52}{744.078}{groups}%
\TextAt{8}{484.49}{744.078}{which}%
\TextAt{8}{517.156}{744.078}{are}%
\TextAt{8}{79.226}{727.032}{Hausdorff,}%
\TextAt{8}{135.134}{727.032}{compact}%
\TextAt{8}{181.043}{727.032}{and}%
\TextAt{8}{202.886}{727.032}{totally}%
\TextAt{8}{239.374}{727.032}{disconnected,}%
\TextAt{8}{310.883}{727.032}{i.e.,}%
\TextAt{8}{332.252}{727.032}{they}%
\TextAt{8}{357.67}{727.032}{admit}%
\TextAt{8}{390.571}{727.032}{a}%
\TextAt{8}{399.407}{727.032}{basis}%
\TextAt{8}{427.228}{727.032}{of}%
\TextAt{8}{439.638}{727.032}{neighbourhoods}%
\TextAt{8}{523.337}{727.032}{of}%
\TextAt{8}{79.226}{709.986}{the}%
\TextAt{8}{98.446}{709.986}{neutral}%
\TextAt{8}{137.835}{709.986}{element}%
\TextAt{8}{180.14}{709.986}{consisting}%
\TextAt{8}{233.313}{709.986}{of}%
\TextAt{8}{245.706}{709.986}{normal}%
\TextAt{8}{284.449}{709.986}{subgroups.}%
\TextAt{8}{344.1}{709.986}{The}%
\TextAt{8}{367.218}{709.986}{next}%
\TextAt{8}{392.619}{709.986}{proposition}%
\TextAt{8}{453.216}{709.986}{shows}%
\TextAt{8}{485.897}{709.986}{that}%
\TextAt{8}{510.33}{709.986}{they}%
\TextAt{8}{79.226}{692.94}{are}%
\TextAt{8}{98.687}{692.94}{precisely}%
\TextAt{8}{146.494}{692.94}{the}%
\TextAt{8}{166.601}{692.94}{projective}%
\TextAt{8}{220.529}{692.94}{limits}%
\TextAt{8}{253.057}{692.94}{of}%
\TextAt{8}{266.337}{692.94}{finite}%
\TextAt{8}{296.2}{692.94}{groups}%
\TextAt{8}{333.94}{692.94}{(which}%
\TextAt{8}{371.932}{692.94}{we}%
\TextAt{8}{389.11}{692.94}{view}%
\TextAt{8}{416.043}{692.94}{as}%
\TextAt{8}{430.363}{692.94}{compact}%
\TextAt{8}{477.142}{692.94}{topological}%
\TextAt{8}{79.226}{675.894}{groups}%
\TextAt{8}{117.016}{675.894}{with}%
\TextAt{8}{143.676}{675.894}{respect}%
\TextAt{8}{183.726}{675.894}{to}%
\TextAt{8}{198.037}{675.894}{the}%
\TextAt{8}{218.193}{675.894}{discrete}%
\TextAt{8}{261.172}{675.894}{topology).}%
\TextAt{12}{79.226}{650.108}{Proposition}%
\TextAt{12}{153.859}{650.108}{(2.8).}%
\TextAt{8}{190.846}{650.108}{If}%
\TextAt{13}{204.208}{650.108}{G}%
\TextAt{8}{219.01}{650.108}{is}%
\TextAt{8}{232.445}{650.108}{a}%
\TextAt{8}{243.871}{650.108}{profinite}%
\TextAt{8}{292.37}{650.108}{group,}%
\TextAt{8}{330.885}{650.108}{and}%
\TextAt{8}{355.319}{650.108}{if}%
\TextAt{13}{367.713}{650.108}{N}%
\TextAt{8}{383.904}{650.108}{varies}%
\TextAt{8}{418.487}{650.108}{over}%
\TextAt{8}{445.204}{650.108}{the}%
\TextAt{8}{467.031}{650.108}{open}%
\TextAt{8}{496.988}{650.108}{normal}%
\TextAt{8}{79.226}{633.062}{subgroups}%
\TextAt{8}{134.638}{633.062}{of}%
\TextAt{13}{147.968}{633.062}{G}%
\TextAt{8}{157.202}{633.062}{,}%
\TextAt{8}{164.351}{633.062}{then}%
\TextAt{8}{191.011}{633.062}{one}%
\TextAt{8}{212.471}{633.062}{has,}%
\TextAt{8}{236.596}{633.062}{algebraically}%
\TextAt{8}{304.573}{633.062}{as}%
\TextAt{8}{318.944}{633.062}{well}%
\TextAt{8}{342.675}{633.062}{as}%
\TextAt{8}{357.045}{633.062}{topologically,}%
\TextAt{8}{428.274}{633.062}{that}%
\TextAt{13}{270.13}{604.852}{G}%
\NativeMathAt{282.684}{604.852}{11.955}{\cong}%
\NativeMathAt{295.304}{604.852}{11.955}{\varprojlim}%
\TextAt{21}{299.648}{592.472}{N}%
\TextAt{13}{313.555}{604.852}{G}%
\TextAt{9}{322.789}{604.852}{\Character{47}}%
\TextAt{13}{328.642}{604.852}{N}%
\TextAt{13}{338.613}{604.852}{\Character{58}}%
\TextAt{8}{79.226}{570.15}{If}%
\TextAt{8}{90.918}{570.15}{conversely}%
\TextAt{16}{146.916}{570.15}{\Character{102}}%
\TextAt{13}{152.894}{570.15}{G}%
\TextAt{21}{162.128}{568.356}{i}%
\TextAt{13}{165.509}{570.15}{\Character{59}}%
\TextAt{13}{170.733}{570.15}{g}%
\TextAt{21}{176.358}{568.356}{ij}%
\TextAt{16}{183.624}{570.15}{\Character{103}}%
\TextAt{8}{193.499}{570.15}{is}%
\TextAt{8}{205.262}{570.15}{a}%
\TextAt{8}{215.018}{570.15}{projective}%
\TextAt{8}{268.995}{570.15}{system}%
\TextAt{8}{307.814}{570.15}{of}%
\TextAt{8}{321.144}{570.15}{finite}%
\TextAt{8}{351.056}{570.15}{(or}%
\TextAt{8}{369.921}{570.15}{even}%
\TextAt{8}{396.581}{570.15}{profinite)}%
\TextAt{8}{447.964}{570.15}{groups,}%
\TextAt{8}{489.006}{570.15}{then}%
\TextAt{13}{278.075}{541.94}{G}%
\TextAt{9}{290.629}{541.94}{\Character{61}}%
\NativeMathAt{303.055}{541.94}{11.955}{\varprojlim}%
\TextAt{21}{309.743}{529.737}{i}%
\TextAt{13}{321.306}{541.94}{G}%
\TextAt{21}{330.54}{540.147}{i}%
\TextAt{8}{79.226}{507.965}{is}%
\TextAt{8}{90.99}{507.965}{a}%
\TextAt{8}{100.745}{507.965}{profinite}%
\TextAt{8}{147.574}{507.965}{group.}%
\TextAt{19}{79.226}{473.438}{Proof.}%
\TextAt{8}{116}{473.438}{Let}%
\TextAt{13}{137.358}{473.438}{G}%
\TextAt{8}{150.877}{473.438}{be}%
\TextAt{8}{167.19}{473.438}{a}%
\TextAt{8}{177.334}{473.438}{profinite}%
\TextAt{8}{224.551}{473.438}{group}%
\TextAt{8}{258.114}{473.438}{and}%
\TextAt{8}{281.265}{473.438}{let}%
\TextAt{16}{298.558}{473.438}{\Character{102}}%
\TextAt{13}{304.536}{473.438}{N}%
\TextAt{21}{313.91}{471.645}{i}%
\TextAt{16}{321.274}{473.438}{\Character{106}}%
\TextAt{13}{328.577}{473.438}{i}%
\NativeMathAt{336.553}{473.438}{11.955}{\in}%
\TextAt{13}{348.506}{473.438}{I}%
\TextAt{16}{354.609}{473.438}{\Character{103}}%
\TextAt{8}{364.872}{473.438}{be}%
\TextAt{8}{381.185}{473.438}{the}%
\TextAt{8}{401.73}{473.438}{family}%
\TextAt{8}{437.888}{473.438}{of}%
\TextAt{8}{451.606}{473.438}{its}%
\TextAt{8}{468.313}{473.438}{open}%
\TextAt{8}{496.988}{473.438}{normal}%
\TextAt{8}{79.226}{456.392}{subgroups.}%
\TextAt{8}{139.949}{456.392}{We}%
\TextAt{8}{160.347}{456.392}{make}%
\TextAt{13}{191.169}{456.392}{I}%
\TextAt{8}{201.423}{456.392}{into}%
\TextAt{8}{225.419}{456.392}{a}%
\TextAt{8}{235.427}{456.392}{directed}%
\TextAt{8}{280.548}{456.392}{system}%
\TextAt{8}{319.62}{456.392}{by}%
\TextAt{8}{336.132}{456.392}{defining}%
\TextAt{13}{380.608}{456.392}{i}%
\NativeMathAt{388.353}{456.392}{11.955}{\leq}%
\TextAt{13}{401.404}{456.392}{j}%
\TextAt{8}{411.07}{456.392}{if}%
\TextAt{13}{422.047}{456.392}{N}%
\TextAt{21}{431.421}{454.599}{i}%
\NativeMathAt{438.555}{456.392}{11.955}{\supseteq}%
\TextAt{13}{451.606}{456.392}{N}%
\TextAt{21}{460.979}{454.599}{j}%
\TextAt{8}{465.362}{456.392}{.}%
\TextAt{8}{474.57}{456.392}{The}%
\TextAt{8}{498.877}{456.392}{groups}%
\TextAt{13}{79.226}{439.346}{G}%
\TextAt{21}{88.46}{437.553}{i}%
\TextAt{9}{95.162}{439.346}{\Character{61}}%
\TextAt{13}{107.587}{439.346}{G}%
\TextAt{9}{116.821}{439.346}{\Character{47}}%
\TextAt{13}{122.674}{439.346}{N}%
\TextAt{21}{132.048}{437.553}{i}%
\TextAt{8}{138.848}{439.346}{are}%
\TextAt{8}{157.88}{439.346}{finite}%
\TextAt{8}{187.313}{439.346}{since}%
\TextAt{8}{215.503}{439.346}{the}%
\TextAt{8}{235.18}{439.346}{cosets}%
\TextAt{8}{268.642}{439.346}{of}%
\TextAt{13}{281.494}{439.346}{N}%
\TextAt{21}{290.867}{437.553}{i}%
\TextAt{8}{297.667}{439.346}{in}%
\TextAt{13}{310.841}{439.346}{G}%
\TextAt{8}{323.494}{439.346}{form}%
\TextAt{8}{350.655}{439.346}{a}%
\TextAt{8}{359.932}{439.346}{disjoint}%
\TextAt{8}{401.141}{439.346}{open}%
\TextAt{8}{428.948}{439.346}{covering}%
\TextAt{8}{474.329}{439.346}{of}%
\TextAt{13}{487.181}{439.346}{G}%
\TextAt{8}{496.414}{439.346}{,}%
\TextAt{8}{503.18}{439.346}{which}%
\TextAt{8}{79.226}{422.3}{must}%
\TextAt{8}{108.901}{422.3}{be}%
\TextAt{8}{125.497}{422.3}{finite}%
\TextAt{8}{156.081}{422.3}{because}%
\TextAt{13}{200.054}{422.3}{G}%
\TextAt{8}{213.857}{422.3}{is}%
\TextAt{8}{226.293}{422.3}{compact.}%
\TextAt{8}{279.688}{422.3}{For}%
\TextAt{13}{301.329}{422.3}{i}%
\NativeMathAt{309.788}{422.3}{11.955}{\leq}%
\TextAt{13}{323.552}{422.3}{j}%
\TextAt{8}{333.637}{422.3}{we}%
\TextAt{8}{351.536}{422.3}{have}%
\TextAt{8}{379.203}{422.3}{the}%
\TextAt{8}{400.031}{422.3}{projections}%
\TextAt{13}{460.598}{422.3}{g}%
\TextAt{21}{466.204}{420.507}{ij}%
\TextAt{9}{477.935}{422.3}{\Character{58}}%
\TextAt{13}{485.652}{422.3}{G}%
\TextAt{21}{494.886}{420.507}{j}%
\NativeMathAt{503.734}{422.3}{11.955}{\rightarrow}%
\TextAt{13}{520.155}{422.3}{G}%
\TextAt{21}{529.388}{420.507}{i}%
\TextAt{8}{79.226}{405.254}{and}%
\TextAt{8}{101.302}{405.254}{obtain}%
\TextAt{8}{137.044}{405.254}{a}%
\TextAt{8}{146.113}{405.254}{projective}%
\TextAt{8}{199.404}{405.254}{system}%
\TextAt{16}{237.536}{405.254}{\Character{102}}%
\TextAt{13}{243.514}{405.254}{G}%
\TextAt{21}{252.747}{403.461}{i}%
\TextAt{13}{256.129}{405.254}{\Character{59}}%
\TextAt{13}{261.353}{405.254}{g}%
\TextAt{21}{266.978}{403.461}{ij}%
\TextAt{16}{274.243}{405.254}{\Character{103}}%
\TextAt{8}{283.432}{405.254}{of}%
\TextAt{8}{296.076}{405.254}{finite,}%
\TextAt{8}{328.69}{405.254}{and}%
\TextAt{8}{350.767}{405.254}{hence}%
\TextAt{8}{382.586}{405.254}{discrete,}%
\TextAt{8}{428.268}{405.254}{compact}%
\TextAt{8}{474.41}{405.254}{groups.}%
\TextAt{8}{516.522}{405.254}{We}%
\TextAt{8}{79.226}{388.209}{show}%
\TextAt{8}{108.229}{388.209}{that}%
\TextAt{8}{133.598}{388.209}{the}%
\TextAt{8}{153.755}{388.209}{homomorphism}%
\TextAt{13}{163.835}{359.999}{f}%
\TextAt{9}{174.202}{359.999}{\Character{58}}%
\TextAt{13}{180.775}{359.999}{G}%
\NativeArrowAt{193.329}{359.999}{11.955}{19.26}{\longrightarrow}%
\NativeMathAt{215.911}{359.999}{11.955}{\varprojlim}%
\TextAt{21}{217.63}{347.619}{i}%
\NativeMathAt{220.513}{347.619}{7.97}{\in}%
\TextAt{21}{226.158}{347.619}{I}%
\TextAt{13}{234.163}{359.999}{G}%
\TextAt{21}{243.396}{358.205}{i}%
\TextAt{13}{246.778}{359.999}{\Character{59}}%
\TextAt{13}{275.913}{359.999}{\Character{27}}%
\NativeArrowAt{286.335}{359.999}{11.955}{19.26}{\longmapsto}%
\TextAt{9}{308.917}{359.999}{\Character{40}}%
\TextAt{13}{313.47}{359.999}{\Character{27}}%
\TextAt{21}{320.123}{358.205}{i}%
\TextAt{9}{323.504}{359.999}{\Character{41}}%
\TextAt{21}{328.057}{358.205}{i}%
\NativeMathAt{330.94}{358.205}{7.97}{\in}%
\TextAt{21}{336.585}{358.205}{I}%
\TextAt{13}{341.377}{359.999}{\Character{59}}%
\TextAt{13}{370.512}{359.999}{\Character{27}}%
\TextAt{21}{377.185}{358.205}{i}%
\TextAt{9}{383.887}{359.999}{\Character{61}}%
\TextAt{13}{396.313}{359.999}{\Character{27}}%
\TextAt{8}{406.716}{359.999}{mod}%
\TextAt{13}{432.154}{359.999}{N}%
\TextAt{21}{441.528}{358.205}{i}%
\TextAt{13}{444.909}{359.999}{\Character{59}}%
\TextAt{8}{79.226}{325.458}{is}%
\TextAt{8}{91.572}{325.458}{an}%
\TextAt{8}{108.413}{325.458}{isomorphism}%
\TextAt{8}{177.415}{325.458}{and}%
\TextAt{8}{200.759}{325.458}{a}%
\TextAt{8}{211.097}{325.458}{homeomorphism.}%
\TextAt{13}{305.264}{325.458}{f}%
\TextAt{8}{316.79}{325.458}{is}%
\TextAt{8}{329.136}{325.458}{injective}%
\TextAt{8}{376.212}{325.458}{because}%
\TextAt{8}{420.095}{325.458}{its}%
\TextAt{8}{436.996}{325.458}{kernel}%
\TextAt{8}{472.045}{325.458}{is}%
\TextAt{8}{484.391}{325.458}{the}%
\TextAt{8}{505.129}{325.458}{inter-}%
\TextAt{8}{79.226}{308.412}{section}%
\TextAt{20}{119.741}{317.378}{\Character{84}}%
\TextAt{21}{129.703}{304.925}{i}%
\NativeMathAt{132.587}{304.925}{7.97}{\in}%
\TextAt{21}{138.232}{304.925}{I}%
\TextAt{13}{145.016}{308.412}{N}%
\TextAt{21}{154.39}{306.618}{i}%
\TextAt{8}{157.771}{308.412}{,}%
\TextAt{8}{166.712}{308.412}{which}%
\TextAt{8}{201.632}{308.412}{equals}%
\TextAt{16}{238.571}{308.412}{\Character{102}}%
\TextAt{9}{244.549}{308.412}{1}%
\TextAt{16}{250.402}{308.412}{\Character{103}}%
\TextAt{8}{261.71}{308.412}{because}%
\TextAt{13}{306.445}{308.412}{G}%
\TextAt{8}{321.009}{308.412}{is}%
\TextAt{8}{334.206}{308.412}{Hausdorff}%
\TextAt{8}{389.031}{308.412}{and}%
\TextAt{8}{413.227}{308.412}{the}%
\TextAt{13}{434.816}{308.412}{N}%
\TextAt{21}{444.19}{306.618}{i}%
\TextAt{8}{452.902}{308.412}{form}%
\TextAt{8}{481.975}{308.412}{a}%
\TextAt{8}{493.164}{308.412}{basis}%
\TextAt{8}{523.337}{308.412}{of}%
\TextAt{8}{79.226}{291.366}{neighbourhoods}%
\TextAt{8}{163.845}{291.366}{of}%
\TextAt{9}{177.175}{291.366}{1}%
\TextAt{8}{183.028}{291.366}{.}%
\TextAt{8}{191.476}{291.366}{The}%
\TextAt{8}{215.53}{291.366}{groups}%
\TextAt{13}{245.558}{263.156}{U}%
\TextAt{21}{253.508}{261.363}{S}%
\TextAt{9}{262.923}{263.156}{\Character{61}}%
\TextAt{20}{275.348}{274.513}{\Character{89}}%
\TextAt{21}{275.924}{248.61}{i}%
\TextAt{17}{279.748}{248.61}{\Character{47}}%
\NativeMathAt{278.807}{248.61}{7.97}{\in}%
\TextAt{21}{284.453}{248.61}{S}%
\TextAt{13}{292.617}{263.156}{G}%
\TextAt{21}{301.85}{261.363}{i}%
\TextAt{16}{307.888}{263.156}{\Character{2}}%
\TextAt{20}{319.843}{274.513}{\Character{89}}%
\TextAt{21}{320.419}{249.142}{i}%
\NativeMathAt{323.302}{249.142}{7.97}{\in}%
\TextAt{21}{328.948}{249.142}{S}%
\TextAt{16}{335.12}{263.156}{\Character{102}}%
\TextAt{9}{341.097}{263.156}{1}%
\TextAt{21}{346.95}{261.363}{G}%
\TextAt{22}{353.549}{260.148}{i}%
\TextAt{16}{357.209}{263.156}{\Character{103}}%
\TextAt{13}{363.186}{263.156}{\Character{59}}%
\TextAt{8}{79.226}{226.335}{with}%
\TextAt{13}{105.712}{226.335}{S}%
\TextAt{8}{117.331}{226.335}{varying}%
\TextAt{8}{158.797}{226.335}{over}%
\TextAt{8}{183.669}{226.335}{the}%
\TextAt{8}{203.652}{226.335}{finite}%
\TextAt{8}{233.389}{226.335}{subsets}%
\TextAt{8}{273.72}{226.335}{of}%
\TextAt{13}{286.876}{226.335}{I}%
\TextAt{8}{292.978}{226.335}{,}%
\TextAt{8}{299.989}{226.335}{form}%
\TextAt{8}{327.455}{226.335}{a}%
\TextAt{8}{337.036}{226.335}{basis}%
\TextAt{8}{365.603}{226.335}{of}%
\TextAt{8}{378.759}{226.335}{neighbourhoods}%
\TextAt{8}{463.204}{226.335}{of}%
\TextAt{8}{476.36}{226.335}{the}%
\TextAt{8}{496.342}{226.335}{neutral}%
\TextAt{8}{79.226}{209.289}{element}%
\TextAt{8}{122.609}{209.289}{in}%
\TextAt{20}{136.402}{218.256}{\Character{81}}%
\TextAt{21}{147.693}{205.802}{i}%
\NativeMathAt{150.576}{205.802}{7.97}{\in}%
\TextAt{21}{156.222}{205.802}{I}%
\TextAt{13}{163.006}{209.289}{G}%
\TextAt{21}{172.24}{207.496}{i}%
\TextAt{8}{175.621}{209.289}{.}%
\TextAt{8}{184.491}{209.289}{As}%
\TextAt{13}{201.918}{209.289}{f}%
\TextAt{3}{208.965}{213.628}{\Character{0}}%
\TextAt{17}{215.551}{213.628}{1}%
\TextAt{9}{220.284}{209.289}{\Character{40}}%
\TextAt{13}{224.836}{209.289}{U}%
\TextAt{21}{232.786}{207.496}{S}%
\TextAt{16}{241.633}{209.289}{\Character{92}}%
\NativeMathAt{252.355}{209.289}{11.955}{\varprojlim}%
\TextAt{21}{268.614}{203.85}{i}%
\TextAt{13}{273.988}{209.289}{G}%
\TextAt{21}{283.222}{207.496}{i}%
\TextAt{9}{286.603}{209.289}{\Character{41}}%
\TextAt{9}{294.697}{209.289}{\Character{61}}%
\TextAt{20}{307.381}{218.256}{\Character{84}}%
\TextAt{21}{317.343}{205.802}{i}%
\NativeMathAt{320.226}{205.802}{7.97}{\in}%
\TextAt{21}{325.872}{205.802}{S}%
\TextAt{13}{333.958}{209.289}{N}%
\TextAt{21}{343.332}{207.496}{i}%
\TextAt{8}{346.713}{209.289}{,}%
\TextAt{8}{354.038}{209.289}{we}%
\TextAt{8}{371.406}{209.289}{see}%
\TextAt{8}{390.46}{209.289}{that}%
\TextAt{13}{415.969}{209.289}{f}%
\TextAt{8}{427.054}{209.289}{is}%
\TextAt{8}{438.958}{209.289}{continuous.}%
\TextAt{8}{502.535}{209.289}{More-}%
\TextAt{8}{79.226}{192.243}{over,}%
\TextAt{8}{107.978}{192.243}{as}%
\TextAt{13}{122.712}{192.243}{G}%
\TextAt{8}{136.206}{192.243}{is}%
\TextAt{8}{148.333}{192.243}{compact,}%
\TextAt{8}{198.867}{192.243}{the}%
\TextAt{8}{219.387}{192.243}{image}%
\TextAt{13}{253.572}{192.243}{f}%
\TextAt{9}{260.618}{192.243}{\Character{40}}%
\TextAt{13}{265.17}{192.243}{G}%
\TextAt{9}{274.404}{192.243}{\Character{41}}%
\TextAt{8}{283.217}{192.243}{is}%
\TextAt{8}{295.344}{192.243}{closed}%
\TextAt{8}{330.234}{192.243}{in}%
\NativeMathAt{344.25}{192.243}{11.955}{\varprojlim}%
\TextAt{21}{360.509}{186.804}{i}%
\TextAt{13}{365.883}{192.243}{G}%
\TextAt{21}{375.117}{190.45}{i}%
\TextAt{8}{378.498}{192.243}{.}%
\TextAt{8}{388.036}{192.243}{On}%
\TextAt{8}{407.91}{192.243}{the}%
\TextAt{8}{428.43}{192.243}{other}%
\TextAt{8}{459.363}{192.243}{hand}%
\TextAt{8}{488.992}{192.243}{it}%
\TextAt{8}{501.06}{192.243}{is}%
\TextAt{8}{513.187}{192.243}{also}%
\TextAt{8}{79.226}{175.197}{dense.}%
\TextAt{8}{115.578}{175.197}{For}%
\TextAt{8}{136.187}{175.197}{if}%
\TextAt{13}{146.552}{175.197}{x}%
\TextAt{9}{156.525}{175.197}{\Character{61}}%
\TextAt{9}{168.935}{175.197}{\Character{40}}%
\TextAt{13}{173.503}{175.197}{x}%
\TextAt{21}{180.155}{173.404}{i}%
\TextAt{9}{183.536}{175.197}{\Character{41}}%
\TextAt{21}{188.088}{173.404}{i}%
\NativeMathAt{190.971}{173.404}{7.97}{\in}%
\TextAt{21}{196.617}{173.404}{I}%
\NativeMathAt{204.73}{175.197}{11.955}{\in}%
\NativeMathAt{216.021}{175.197}{11.955}{\varprojlim}%
\TextAt{21}{232.28}{169.758}{i}%
\TextAt{13}{237.653}{175.197}{G}%
\TextAt{21}{246.887}{173.404}{i}%
\TextAt{8}{253.806}{175.197}{and}%
\TextAt{13}{276.209}{175.197}{x}%
\TextAt{9}{282.861}{175.197}{\Character{40}}%
\TextAt{13}{287.414}{175.197}{U}%
\TextAt{21}{295.364}{173.404}{S}%
\TextAt{16}{303.379}{175.197}{\Character{92}}%
\NativeMathAt{313.271}{175.197}{11.955}{\varprojlim}%
\TextAt{21}{329.53}{169.758}{i}%
\TextAt{13}{334.904}{175.197}{G}%
\TextAt{21}{344.138}{173.404}{i}%
\TextAt{9}{347.519}{175.197}{\Character{41}}%
\TextAt{8}{355.609}{175.197}{is}%
\TextAt{8}{367.014}{175.197}{a}%
\TextAt{8}{376.409}{175.197}{fundamental}%
\TextAt{8}{443.692}{175.197}{neighbourhood}%
\TextAt{8}{523.337}{175.197}{of}%
\TextAt{13}{79.226}{158.151}{x}%
\TextAt{8}{85.878}{158.151}{,}%
\TextAt{8}{93.066}{158.151}{then}%
\TextAt{8}{119.757}{158.151}{we}%
\TextAt{8}{137.015}{158.151}{may}%
\TextAt{8}{162.415}{158.151}{choose}%
\TextAt{8}{199.578}{158.151}{a}%
\TextAt{13}{209.365}{158.151}{y}%
\NativeMathAt{218.875}{158.151}{11.955}{\in}%
\TextAt{13}{230.218}{158.151}{G}%
\TextAt{8}{243.38}{158.151}{which}%
\TextAt{8}{276.897}{158.151}{is}%
\TextAt{8}{288.692}{158.151}{mapped}%
\TextAt{8}{333.268}{158.151}{to}%
\TextAt{13}{347.609}{158.151}{x}%
\TextAt{21}{354.261}{156.358}{k}%
\TextAt{8}{363.309}{158.151}{under}%
\TextAt{8}{396.503}{158.151}{the}%
\TextAt{8}{416.69}{158.151}{projection}%
\TextAt{13}{472.002}{158.151}{G}%
\NativeMathAt{484.609}{158.151}{11.955}{\rightarrow}%
\TextAt{13}{499.938}{158.151}{G}%
\TextAt{9}{509.171}{158.151}{\Character{47}}%
\TextAt{13}{515.024}{158.151}{N}%
\TextAt{21}{524.398}{156.358}{k}%
\TextAt{8}{529.518}{158.151}{,}%
\TextAt{8}{79.226}{141.105}{where}%
\TextAt{8}{113.988}{141.105}{we}%
\TextAt{8}{132.168}{141.105}{put}%
\TextAt{13}{154.58}{141.105}{N}%
\TextAt{21}{163.953}{139.312}{k}%
\TextAt{9}{174.017}{141.105}{\Character{61}}%
\TextAt{20}{188.066}{150.072}{\Character{84}}%
\TextAt{21}{198.028}{137.618}{i}%
\NativeMathAt{200.911}{137.618}{7.97}{\in}%
\TextAt{21}{206.557}{137.618}{S}%
\TextAt{13}{214.643}{141.105}{N}%
\TextAt{21}{224.017}{139.312}{i}%
\TextAt{8}{227.398}{141.105}{.}%
\TextAt{8}{238.704}{141.105}{Then}%
\TextAt{13}{270.214}{141.105}{y}%
\TextAt{8}{279.671}{141.105}{mod}%
\TextAt{13}{305.109}{141.105}{N}%
\TextAt{21}{314.483}{139.312}{i}%
\TextAt{9}{322.808}{141.105}{\Character{61}}%
\TextAt{13}{336.857}{141.105}{x}%
\TextAt{21}{343.509}{139.312}{i}%
\TextAt{8}{351.74}{141.105}{for}%
\TextAt{8}{370.577}{141.105}{all}%
\TextAt{13}{387.789}{141.105}{i}%
\NativeMathAt{396.726}{141.105}{11.955}{\in}%
\TextAt{13}{409.64}{141.105}{S}%
\TextAt{8}{417.535}{141.105}{,}%
\TextAt{8}{425.875}{141.105}{so}%
\TextAt{8}{441.198}{141.105}{that}%
\TextAt{13}{467.519}{141.105}{f}%
\TextAt{9}{474.565}{141.105}{\Character{40}}%
\TextAt{13}{479.118}{141.105}{y}%
\TextAt{9}{485.254}{141.105}{\Character{41}}%
\TextAt{8}{494.656}{141.105}{belongs}%
\TextAt{8}{79.226}{124.059}{to}%
\TextAt{8}{94.029}{124.059}{the}%
\TextAt{8}{114.677}{124.059}{neighbourhood}%
\TextAt{13}{195.173}{124.059}{x}%
\TextAt{9}{201.826}{124.059}{\Character{40}}%
\TextAt{13}{206.378}{124.059}{U}%
\TextAt{21}{214.328}{122.266}{S}%
\TextAt{16}{223.414}{124.059}{\Character{92}}%
\NativeMathAt{234.377}{124.059}{11.955}{\varprojlim}%
\TextAt{21}{250.636}{118.62}{i}%
\TextAt{13}{256.009}{124.059}{G}%
\TextAt{21}{265.243}{122.266}{i}%
\TextAt{9}{268.624}{124.059}{\Character{41}}%
\TextAt{8}{273.177}{124.059}{.}%
\TextAt{8}{283.101}{124.059}{Therefore}%
\TextAt{8}{336.591}{124.059}{the}%
\TextAt{8}{357.24}{124.059}{closed}%
\TextAt{8}{392.258}{124.059}{set}%
\TextAt{13}{411.018}{124.059}{f}%
\TextAt{9}{418.064}{124.059}{\Character{40}}%
\TextAt{13}{422.617}{124.059}{G}%
\TextAt{9}{431.85}{124.059}{\Character{41}}%
\TextAt{8}{440.792}{124.059}{is}%
\TextAt{8}{453.048}{124.059}{indeed}%
\TextAt{8}{490.602}{124.059}{dense}%
\TextAt{8}{523.014}{124.059}{in}%
\NativeMathAt{79.226}{107.013}{11.955}{\varprojlim}%
\TextAt{21}{95.485}{101.574}{i}%
\TextAt{13}{100.859}{107.013}{G}%
\TextAt{21}{110.093}{105.22}{i}%
\TextAt{8}{113.474}{107.013}{,}%
\TextAt{8}{121.01}{107.013}{and}%
\TextAt{8}{144.083}{107.013}{so}%
\TextAt{13}{158.762}{107.013}{f}%
\TextAt{9}{165.809}{107.013}{\Character{40}}%
\TextAt{13}{170.361}{107.013}{G}%
\TextAt{9}{179.595}{107.013}{\Character{41}}%
\TextAt{9}{187.976}{107.013}{\Character{61}}%
\NativeMathAt{200.949}{107.013}{11.955}{\varprojlim}%
\TextAt{21}{217.208}{101.574}{i}%
\TextAt{13}{222.582}{107.013}{G}%
\TextAt{21}{231.816}{105.22}{i}%
\TextAt{8}{235.197}{107.013}{.}%
\TextAt{8}{244.574}{107.013}{Since}%
\TextAt{13}{275.442}{107.013}{G}%
\TextAt{8}{288.882}{107.013}{is}%
\TextAt{8}{300.956}{107.013}{compact,}%
\TextAt{13}{351.423}{107.013}{f}%
\TextAt{8}{362.677}{107.013}{maps}%
\TextAt{8}{393.616}{107.013}{closed}%
\TextAt{8}{428.452}{107.013}{sets}%
\TextAt{8}{451.644}{107.013}{into}%
\TextAt{8}{475.697}{107.013}{closed}%
\TextAt{8}{510.533}{107.013}{sets,}%
\TextAt{8}{79.226}{89.967}{and}%
\TextAt{8}{101.22}{89.967}{thus}%
\TextAt{8}{126.203}{89.967}{also}%
\TextAt{8}{148.914}{89.967}{open}%
\TextAt{8}{176.432}{89.967}{sets}%
\TextAt{8}{198.545}{89.967}{into}%
\TextAt{8}{221.52}{89.967}{open}%
\TextAt{8}{249.037}{89.967}{sets.}%
\TextAt{8}{276.214}{89.967}{This}%
\TextAt{8}{302.165}{89.967}{shows}%
\TextAt{8}{335.015}{89.967}{that}%
\TextAt{13}{359.615}{89.967}{f}%
\TextAt{9}{369.982}{89.967}{\Character{58}}%
\TextAt{13}{376.555}{89.967}{G}%
\NativeMathAt{389.109}{89.967}{11.955}{\rightarrow}%
\NativeMathAt{404.385}{89.967}{11.955}{\varprojlim}%
\TextAt{21}{420.644}{84.528}{i}%
\TextAt{13}{426.018}{89.967}{G}%
\TextAt{21}{435.252}{88.174}{i}%
\TextAt{8}{441.762}{89.967}{is}%
\TextAt{8}{452.757}{89.967}{an}%
\TextAt{8}{468.247}{89.967}{isomorphism}%
\TextAt{8}{79.226}{72.921}{and}%
\TextAt{8}{101.989}{72.921}{a}%
\TextAt{8}{111.744}{72.921}{homeomorphism.}%
\TextAt{8}{297.211}{43.033}{247}%
\Destination{78.23}{809.815}{nav.page.247}%
\Destination{79.36}{665.155}{nav.print.270}%
\Destination{79.36}{664.155}{nav.number.4.2.8}%
}

% Page 260 (body page 248)
\Page{%
\TextAt{2}{79.226}{796.943}{IV.}%
\TextAt{2}{98.012}{796.943}{Abstract}%
\TextAt{2}{143.252}{796.943}{Class}%
\TextAt{2}{171.844}{796.943}{Field}%
\TextAt{2}{199.575}{796.943}{Theory}%
\Vector{q 0 G 0 g 0.299 w 7.226 18.729 m 460.769 18.729 l S Q}%
\TextAt{8}{97.159}{761.124}{Conversely,}%
\TextAt{8}{160.324}{761.124}{let}%
\TextAt{16}{178.544}{761.124}{\Character{102}}%
\TextAt{13}{184.522}{761.124}{G}%
\TextAt{21}{193.755}{759.33}{i}%
\TextAt{13}{197.137}{761.124}{\Character{59}}%
\TextAt{13}{202.361}{761.124}{g}%
\TextAt{21}{207.986}{759.33}{ij}%
\TextAt{16}{215.251}{761.124}{\Character{103}}%
\TextAt{8}{226.442}{761.124}{be}%
\TextAt{8}{243.681}{761.124}{a}%
\TextAt{8}{254.752}{761.124}{projective}%
\TextAt{8}{310.045}{761.124}{system}%
\TextAt{8}{350.178}{761.124}{of}%
\TextAt{8}{364.824}{761.124}{profinite}%
\TextAt{8}{412.967}{761.124}{groups.}%
\TextAt{8}{459.254}{761.124}{As}%
\TextAt{8}{477.857}{761.124}{the}%
\TextAt{13}{499.329}{761.124}{G}%
\TextAt{21}{508.562}{759.33}{i}%
\TextAt{8}{517.156}{761.124}{are}%
\TextAt{8}{79.226}{744.078}{Hausdorff}%
\TextAt{8}{132.026}{744.078}{and}%
\TextAt{8}{154.197}{744.078}{compact,}%
\TextAt{8}{203.804}{744.078}{so}%
\TextAt{8}{217.582}{744.078}{is}%
\TextAt{8}{228.754}{744.078}{the}%
\TextAt{8}{248.319}{744.078}{projective}%
\TextAt{8}{301.705}{744.078}{limit}%
\TextAt{13}{329.076}{744.078}{G}%
\TextAt{9}{341.63}{744.078}{\Character{61}}%
\NativeMathAt{354.056}{744.078}{11.955}{\varprojlim}%
\TextAt{21}{370.315}{738.638}{i}%
\TextAt{13}{375.689}{744.078}{G}%
\TextAt{21}{384.922}{742.284}{i}%
\TextAt{8}{388.304}{744.078}{,}%
\TextAt{8}{394.979}{744.078}{by}%
\TextAt{8}{410.647}{744.078}{(2.3).}%
\TextAt{8}{442.975}{744.078}{If}%
\TextAt{13}{454.076}{744.078}{N}%
\TextAt{21}{463.45}{742.284}{i}%
\TextAt{8}{470.136}{744.078}{varies}%
\TextAt{8}{502.457}{744.078}{over}%
\TextAt{8}{526.911}{744.078}{a}%
\TextAt{8}{79.226}{727.032}{basis}%
\TextAt{8}{107.602}{727.032}{of}%
\TextAt{8}{120.568}{727.032}{neighbourhoods}%
\TextAt{8}{204.822}{727.032}{of}%
\TextAt{8}{217.788}{727.032}{the}%
\TextAt{8}{237.58}{727.032}{neutral}%
\TextAt{8}{277.54}{727.032}{element}%
\TextAt{8}{320.418}{727.032}{in}%
\TextAt{13}{333.706}{727.032}{G}%
\TextAt{21}{342.94}{725.238}{i}%
\TextAt{8}{349.854}{727.032}{which}%
\TextAt{8}{382.976}{727.032}{consists}%
\TextAt{8}{425.722}{727.032}{of}%
\TextAt{8}{438.688}{727.032}{normal}%
\TextAt{8}{478.003}{727.032}{subgroups,}%
\TextAt{8}{79.226}{709.986}{then}%
\TextAt{8}{105.886}{709.986}{the}%
\TextAt{8}{126.043}{709.986}{groups}%
\TextAt{13}{252.217}{692.94}{U}%
\TextAt{21}{260.168}{691.147}{S}%
\TextAt{9}{269.582}{692.94}{\Character{61}}%
\TextAt{20}{282.008}{704.297}{\Character{89}}%
\TextAt{21}{282.584}{678.394}{i}%
\TextAt{17}{286.408}{678.394}{\Character{47}}%
\NativeMathAt{285.467}{678.394}{7.97}{\in}%
\TextAt{21}{291.112}{678.394}{S}%
\TextAt{13}{299.276}{692.94}{G}%
\TextAt{21}{308.51}{691.147}{i}%
\TextAt{16}{314.548}{692.94}{\Character{2}}%
\TextAt{20}{326.503}{704.297}{\Character{89}}%
\TextAt{21}{327.079}{678.926}{i}%
\NativeMathAt{329.962}{678.926}{7.97}{\in}%
\TextAt{21}{335.608}{678.926}{S}%
\TextAt{13}{343.772}{692.94}{N}%
\TextAt{21}{353.145}{691.147}{i}%
\TextAt{13}{356.527}{692.94}{\Character{59}}%
\TextAt{8}{79.226}{660.807}{with}%
\TextAt{13}{106.978}{660.807}{S}%
\TextAt{8}{119.863}{660.807}{varying}%
\TextAt{8}{162.594}{660.807}{over}%
\TextAt{8}{188.732}{660.807}{the}%
\TextAt{8}{209.981}{660.807}{finite}%
\TextAt{8}{240.984}{660.807}{subsets}%
\TextAt{8}{282.58}{660.807}{of}%
\TextAt{13}{297.002}{660.807}{I}%
\TextAt{8}{303.105}{660.807}{,}%
\TextAt{8}{311.619}{660.807}{make}%
\TextAt{8}{343.28}{660.807}{up}%
\TextAt{8}{361.277}{660.807}{a}%
\TextAt{8}{372.124}{660.807}{basis}%
\TextAt{8}{401.956}{660.807}{of}%
\TextAt{8}{416.378}{660.807}{neighbourhoods}%
\TextAt{8}{502.089}{660.807}{of}%
\TextAt{8}{516.51}{660.807}{the}%
\TextAt{8}{79.226}{643.761}{neutral}%
\TextAt{8}{119.72}{643.761}{element}%
\TextAt{8}{163.13}{643.761}{in}%
\TextAt{20}{176.952}{652.728}{\Character{81}}%
\TextAt{21}{188.243}{640.274}{i}%
\NativeMathAt{191.126}{640.274}{7.97}{\in}%
\TextAt{21}{196.772}{640.274}{I}%
\TextAt{13}{203.556}{643.761}{G}%
\TextAt{21}{212.79}{641.968}{i}%
\TextAt{8}{220.237}{643.761}{consisting}%
\TextAt{8}{274.515}{643.761}{of}%
\TextAt{8}{288.014}{643.761}{normal}%
\TextAt{8}{327.862}{643.761}{subgroups.}%
\TextAt{8}{388.332}{643.761}{The}%
\TextAt{8}{412.554}{643.761}{normal}%
\TextAt{8}{452.402}{643.761}{subgroups}%
\TextAt{13}{507.983}{643.761}{U}%
\TextAt{21}{515.934}{641.968}{S}%
\TextAt{16}{524.799}{643.761}{\Character{92}}%
\NativeMathAt{79.226}{626.715}{11.955}{\varprojlim}%
\TextAt{21}{95.485}{621.276}{i}%
\TextAt{13}{100.859}{626.715}{G}%
\TextAt{21}{110.093}{624.922}{i}%
\TextAt{8}{117.983}{626.715}{therefore}%
\TextAt{8}{167.694}{626.715}{form}%
\TextAt{8}{195.946}{626.715}{a}%
\TextAt{8}{206.313}{626.715}{basis}%
\TextAt{8}{235.665}{626.715}{of}%
\TextAt{8}{249.607}{626.715}{neighbourhoods}%
\TextAt{8}{334.837}{626.715}{of}%
\TextAt{8}{348.779}{626.715}{the}%
\TextAt{8}{369.547}{626.715}{neutral}%
\TextAt{8}{410.483}{626.715}{element}%
\TextAt{8}{454.337}{626.715}{in}%
\NativeMathAt{468.601}{626.715}{11.955}{\varprojlim}%
\TextAt{21}{484.86}{621.276}{i}%
\TextAt{13}{490.234}{626.715}{G}%
\TextAt{21}{499.468}{624.922}{i}%
\TextAt{8}{502.849}{626.715}{;}%
\TextAt{8}{510.915}{626.715}{thus}%
\NativeMathAt{79.226}{609.669}{11.955}{\varprojlim}%
\TextAt{21}{95.485}{604.23}{i}%
\TextAt{13}{100.859}{609.669}{G}%
\TextAt{21}{110.093}{607.876}{i}%
\TextAt{8}{117.371}{609.669}{is}%
\TextAt{8}{129.135}{609.669}{a}%
\TextAt{8}{138.891}{609.669}{profinite}%
\TextAt{8}{185.719}{609.669}{group.}%
\Vector{q 0 G 0 g 0.399 w 452.334 -164.331 m 452.334 -156.261 l S Q}%
\Vector{q 0 G 0 g 0.399 w 452.534 -156.46 m 459.707 -156.46 l S Q}%
\Vector{q 0 G 0 g 0.399 w 452.534 -164.131 m 459.707 -164.131 l S Q}%
\Vector{q 0 G 0 g 0.399 w 459.906 -164.331 m 459.906 -156.261 l S Q}%
\TextAt{8}{97.159}{580.668}{Let}%
\TextAt{8}{119.79}{580.668}{us}%
\TextAt{8}{136.467}{580.668}{now}%
\TextAt{8}{162.518}{580.668}{illustrate}%
\TextAt{8}{213.674}{580.668}{the}%
\TextAt{8}{235.492}{580.668}{notions}%
\TextAt{8}{278.196}{580.668}{of}%
\TextAt{8}{293.188}{580.668}{profinite}%
\TextAt{8}{341.678}{580.668}{group}%
\TextAt{8}{376.515}{580.668}{and}%
\TextAt{8}{400.94}{580.668}{projective}%
\TextAt{8}{456.579}{580.668}{limit}%
\TextAt{8}{486.204}{580.668}{by}%
\TextAt{8}{504.125}{580.668}{a}%
\TextAt{8}{515.542}{580.668}{few}%
\TextAt{8}{79.226}{563.622}{concrete}%
\TextAt{8}{125.397}{563.622}{examples.}%
\TextAt{12}{97.159}{543.587}{Example}%
\TextAt{12}{154.733}{543.587}{1:}%
\TextAt{8}{173.993}{543.587}{The}%
\TextAt{8}{199.843}{543.587}{Galois}%
\TextAt{8}{237.552}{543.587}{group}%
\TextAt{13}{272.524}{543.587}{G}%
\TextAt{9}{288.139}{543.587}{\Character{61}}%
\TextAt{13}{303.625}{543.587}{G}%
\TextAt{9}{312.858}{543.587}{\Character{40}\Character{10}}%
\TextAt{16}{332.246}{543.587}{\Character{106}}%
\TextAt{13}{341.948}{543.587}{k}%
\TextAt{9}{348.438}{543.587}{\Character{41}}%
\TextAt{8}{358.683}{543.587}{of}%
\TextAt{8}{373.809}{543.587}{a}%
\TextAt{8}{385.36}{543.587}{Galois}%
\TextAt{8}{423.07}{543.587}{extension}%
\TextAt{9}{476.631}{543.587}{\Character{10}}%
\TextAt{16}{491.467}{543.587}{\Character{106}}%
\TextAt{13}{501.169}{543.587}{k}%
\TextAt{8}{513.352}{543.587}{is}%
\TextAt{8}{526.911}{543.587}{a}%
\TextAt{8}{79.226}{526.541}{profinite}%
\TextAt{8}{127.33}{526.541}{group}%
\TextAt{8}{161.78}{526.541}{with}%
\TextAt{8}{189.715}{526.541}{respect}%
\TextAt{8}{231.04}{526.541}{to}%
\TextAt{8}{246.625}{526.541}{the}%
\TextAt{8}{268.057}{526.541}{Krull}%
\TextAt{8}{299.889}{526.541}{topology.}%
\TextAt{8}{355.428}{526.541}{This}%
\TextAt{8}{383.423}{526.541}{was}%
\TextAt{8}{407.198}{526.541}{already}%
\TextAt{8}{449.778}{526.541}{stated}%
\TextAt{8}{486.237}{526.541}{in}%
\TextAt{8}{501.165}{526.541}{§1.}%
\TextAt{8}{524.975}{526.541}{If}%
\TextAt{13}{79.226}{509.495}{K}%
\TextAt{16}{93.395}{509.495}{\Character{106}}%
\TextAt{13}{100.067}{509.495}{k}%
\TextAt{8}{110.472}{509.495}{varies}%
\TextAt{8}{143.402}{509.495}{over}%
\TextAt{8}{168.466}{509.495}{the}%
\TextAt{8}{188.64}{509.495}{finite}%
\TextAt{8}{218.57}{509.495}{Galois}%
\TextAt{8}{254.501}{509.495}{subextensions}%
\TextAt{8}{328.844}{509.495}{of}%
\TextAt{9}{342.192}{509.495}{\Character{10}}%
\TextAt{16}{353.997}{509.495}{\Character{106}}%
\TextAt{13}{360.669}{509.495}{k}%
\TextAt{8}{367.159}{509.495}{,}%
\TextAt{8}{374.33}{509.495}{then,}%
\TextAt{8}{404.264}{509.495}{by}%
\TextAt{8}{420.541}{509.495}{definition}%
\TextAt{8}{472.587}{509.495}{of}%
\TextAt{8}{485.935}{509.495}{the}%
\TextAt{8}{506.109}{509.495}{Krull}%
\TextAt{8}{79.226}{492.45}{topology,}%
\TextAt{13}{129.94}{492.45}{G}%
\TextAt{9}{139.174}{492.45}{\Character{40}\Character{10}}%
\TextAt{16}{155.91}{492.45}{\Character{106}}%
\TextAt{13}{162.959}{492.45}{K}%
\TextAt{9}{173.777}{492.45}{\Character{41}}%
\TextAt{8}{182.466}{492.45}{varies}%
\TextAt{8}{215.618}{492.45}{over}%
\TextAt{8}{240.904}{492.45}{the}%
\TextAt{8}{261.3}{492.45}{open}%
\TextAt{8}{289.825}{492.45}{normal}%
\TextAt{8}{329.744}{492.45}{subgroups}%
\TextAt{8}{385.395}{492.45}{of}%
\TextAt{13}{398.965}{492.45}{G}%
\TextAt{8}{408.198}{492.45}{.}%
\TextAt{8}{417.365}{492.45}{In}%
\TextAt{8}{432.225}{492.45}{view}%
\TextAt{8}{459.448}{492.45}{of}%
\TextAt{8}{473.017}{492.45}{the}%
\TextAt{8}{493.413}{492.45}{identity}%
\TextAt{13}{79.226}{475.404}{G}%
\TextAt{9}{88.46}{475.404}{\Character{40}}%
\TextAt{13}{93.012}{475.404}{K}%
\TextAt{16}{109.846}{475.404}{\Character{106}}%
\TextAt{13}{119.184}{475.404}{k}%
\TextAt{9}{125.673}{475.404}{\Character{41}}%
\TextAt{9}{136.229}{475.404}{\Character{61}}%
\TextAt{13}{151.364}{475.404}{G}%
\TextAt{9}{160.597}{475.404}{\Character{40}\Character{10}}%
\TextAt{16}{179.62}{475.404}{\Character{106}}%
\TextAt{13}{188.958}{475.404}{k}%
\TextAt{9}{195.448}{475.404}{\Character{41}\Character{47}}%
\TextAt{13}{205.853}{475.404}{G}%
\TextAt{9}{215.086}{475.404}{\Character{40}\Character{10}}%
\TextAt{16}{234.11}{475.404}{\Character{106}}%
\TextAt{13}{243.447}{475.404}{K}%
\TextAt{9}{254.265}{475.404}{\Character{41}}%
\TextAt{8}{264.297}{475.404}{and}%
\TextAt{8}{288.641}{475.404}{of}%
\TextAt{8}{303.553}{475.404}{(2.8),}%
\TextAt{8}{336.757}{475.404}{we}%
\TextAt{8}{355.567}{475.404}{therefore}%
\TextAt{8}{406.248}{475.404}{obtain}%
\TextAt{8}{444.258}{475.404}{the}%
\TextAt{8}{465.996}{475.404}{Galois}%
\TextAt{8}{503.491}{475.404}{group}%
\TextAt{13}{79.226}{458.358}{G}%
\TextAt{9}{88.46}{458.358}{\Character{40}\Character{10}}%
\TextAt{16}{104.787}{458.358}{\Character{106}}%
\TextAt{13}{111.429}{458.358}{k}%
\TextAt{9}{117.918}{458.358}{\Character{41}}%
\TextAt{8}{126.368}{458.358}{as}%
\TextAt{8}{140.738}{458.358}{the}%
\TextAt{8}{160.895}{458.358}{projective}%
\TextAt{8}{214.872}{458.358}{limit}%
\TextAt{13}{244.476}{429.356}{G}%
\TextAt{9}{253.709}{429.356}{\Character{40}\Character{10}}%
\TextAt{16}{270.037}{429.356}{\Character{106}}%
\TextAt{13}{276.678}{429.356}{k}%
\TextAt{9}{283.168}{429.356}{\Character{41}}%
\NativeMathAt{291.041}{429.356}{11.955}{\cong}%
\NativeMathAt{303.66}{429.356}{11.955}{\varprojlim}%
\TextAt{21}{307.938}{416.976}{K}%
\TextAt{13}{321.912}{429.356}{G}%
\TextAt{9}{331.146}{429.356}{\Character{40}}%
\TextAt{13}{335.698}{429.356}{K}%
\TextAt{16}{349.836}{429.356}{\Character{106}}%
\TextAt{13}{356.478}{429.356}{k}%
\TextAt{9}{362.968}{429.356}{\Character{41}}%
\TextAt{8}{79.226}{393.863}{of}%
\TextAt{8}{92.556}{393.863}{the}%
\TextAt{8}{112.713}{393.863}{finite}%
\TextAt{8}{142.624}{393.863}{Galois}%
\TextAt{8}{178.538}{393.863}{groups}%
\TextAt{13}{216.328}{393.863}{G}%
\TextAt{9}{225.562}{393.863}{\Character{40}}%
\TextAt{13}{230.114}{393.863}{K}%
\TextAt{16}{244.252}{393.863}{\Character{106}}%
\TextAt{13}{250.894}{393.863}{k}%
\TextAt{9}{257.384}{393.863}{\Character{41}}%
\TextAt{8}{261.936}{393.863}{.}%
\TextAt{12}{97.159}{373.828}{Example}%
\TextAt{12}{152.759}{373.828}{2:}%
\TextAt{8}{168.589}{373.828}{If}%
\TextAt{13}{180.361}{373.828}{p}%
\TextAt{8}{190.214}{373.828}{is}%
\TextAt{8}{202.058}{373.828}{a}%
\TextAt{8}{211.893}{373.828}{prime}%
\TextAt{8}{245.137}{373.828}{number,}%
\TextAt{8}{291.085}{373.828}{then}%
\TextAt{8}{317.825}{373.828}{the}%
\TextAt{8}{338.062}{373.828}{rings}%
\TextAt{18}{366.823}{373.828}{Z}%
\TextAt{9}{374.793}{373.828}{\Character{47}}%
\TextAt{13}{380.646}{373.828}{p}%
\TextAt{21}{386.521}{378.167}{n}%
\TextAt{18}{392.157}{373.828}{Z}%
\TextAt{8}{400.127}{373.828}{,}%
\TextAt{13}{407.377}{373.828}{n}%
\NativeMathAt{417.822}{373.828}{11.955}{\in}%
\TextAt{18}{429.249}{373.828}{N}%
\TextAt{8}{437.884}{373.828}{,}%
\TextAt{8}{445.133}{373.828}{form}%
\TextAt{8}{472.854}{373.828}{a}%
\TextAt{8}{482.689}{373.828}{projective}%
\TextAt{8}{79.226}{356.782}{system}%
\TextAt{8}{117.783}{356.782}{with}%
\TextAt{8}{144.181}{356.782}{respect}%
\TextAt{8}{183.969}{356.782}{to}%
\TextAt{8}{198.018}{356.782}{the}%
\TextAt{8}{217.912}{356.782}{projections}%
\TextAt{18}{277.546}{356.782}{Z}%
\TextAt{9}{285.516}{356.782}{\Character{47}}%
\TextAt{13}{291.369}{356.782}{p}%
\TextAt{21}{297.244}{361.121}{n}%
\TextAt{18}{302.881}{356.782}{Z}%
\NativeMathAt{314.172}{356.782}{11.955}{\rightarrow}%
\TextAt{18}{329.448}{356.782}{Z}%
\TextAt{9}{337.418}{356.782}{\Character{47}}%
\TextAt{13}{343.271}{356.782}{p}%
\TextAt{21}{349.146}{361.121}{m}%
\TextAt{18}{357.135}{356.782}{Z}%
\TextAt{8}{365.105}{356.782}{,}%
\TextAt{8}{372.045}{356.782}{for}%
\TextAt{13}{389.668}{356.782}{n}%
\NativeMathAt{399.976}{356.782}{11.955}{\geq}%
\TextAt{13}{412.596}{356.782}{m}%
\TextAt{8}{422.835}{356.782}{.}%
\TextAt{8}{431.196}{356.782}{The}%
\TextAt{8}{454.988}{356.782}{projective}%
\TextAt{8}{508.704}{356.782}{limit}%
\TextAt{18}{265.981}{327.781}{Z}%
\TextAt{21}{273.951}{325.988}{p}%
\TextAt{9}{282.033}{327.781}{\Character{61}}%
\NativeMathAt{294.458}{327.781}{11.955}{\varprojlim}%
\TextAt{21}{300.019}{315.667}{n}%
\TextAt{18}{312.71}{327.781}{Z}%
\TextAt{9}{320.68}{327.781}{\Character{47}}%
\TextAt{13}{326.533}{327.781}{p}%
\TextAt{21}{332.408}{332.717}{n}%
\TextAt{18}{338.045}{327.781}{Z}%
\TextAt{8}{79.226}{292.553}{is}%
\TextAt{8}{90.99}{292.553}{the}%
\TextAt{8}{111.146}{292.553}{ring}%
\TextAt{8}{135.212}{292.553}{of}%
\TextAt{13}{148.542}{292.553}{p}%
\TextAt{8}{154.417}{292.553}{-adic}%
\TextAt{8}{183.026}{292.553}{integers}%
\TextAt{8}{226.34}{292.553}{(see}%
\TextAt{8}{249.808}{292.553}{chap.}%
\TextAt{8}{281.999}{292.553}{II,}%
\TextAt{8}{297.911}{292.553}{§1).}%
\TextAt{12}{97.159}{272.519}{Example}%
\TextAt{12}{151.912}{272.519}{3:}%
\TextAt{8}{167.253}{272.519}{Let}%
\TextAt{14}{187.566}{272.519}{o}%
\TextAt{8}{196.651}{272.519}{be}%
\TextAt{8}{211.919}{272.519}{the}%
\TextAt{8}{231.419}{272.519}{valuation}%
\TextAt{8}{281.835}{272.519}{ring}%
\TextAt{8}{305.244}{272.519}{in}%
\TextAt{8}{318.24}{272.519}{a}%
\TextAt{13}{327.339}{272.519}{p}%
\TextAt{8}{333.214}{272.519}{-adic}%
\TextAt{8}{361.167}{272.519}{number}%
\TextAt{8}{403.106}{272.519}{field}%
\TextAt{13}{427.807}{272.519}{K}%
\TextAt{8}{441.866}{272.519}{and}%
\TextAt{14}{463.972}{272.519}{p}%
\TextAt{8}{473.196}{272.519}{its}%
\TextAt{8}{488.858}{272.519}{maximal}%
\TextAt{8}{79.226}{255.473}{ideal.}%
\TextAt{8}{114.184}{255.473}{The}%
\TextAt{8}{139.053}{255.473}{ideals}%
\TextAt{14}{172.446}{255.473}{p}%
\TextAt{21}{178.429}{259.811}{n}%
\TextAt{8}{184.065}{255.473}{,}%
\TextAt{13}{192.233}{255.473}{n}%
\NativeMathAt{203.93}{255.473}{11.955}{\in}%
\TextAt{18}{216.609}{255.473}{N}%
\TextAt{8}{225.243}{255.473}{,}%
\TextAt{8}{233.411}{255.473}{make}%
\TextAt{8}{264.795}{255.473}{up}%
\TextAt{8}{282.514}{255.473}{a}%
\TextAt{8}{293.085}{255.473}{basis}%
\TextAt{8}{322.64}{255.473}{of}%
\TextAt{8}{336.784}{255.473}{neighbourhoods}%
\TextAt{8}{422.218}{255.473}{of}%
\TextAt{8}{436.363}{255.473}{the}%
\TextAt{8}{457.334}{255.473}{zero}%
\TextAt{8}{482.86}{255.473}{element}%
\TextAt{9}{526.917}{255.473}{0}%
\TextAt{8}{79.226}{238.427}{in}%
\TextAt{14}{93.18}{238.427}{o}%
\TextAt{8}{99.024}{238.427}{.}%
\TextAt{14}{108.374}{238.427}{o}%
\TextAt{8}{118.416}{238.427}{is}%
\TextAt{8}{130.481}{238.427}{Hausdorff}%
\TextAt{8}{184.173}{238.427}{and}%
\TextAt{8}{207.237}{238.427}{compact,}%
\TextAt{8}{257.693}{238.427}{and}%
\TextAt{8}{280.756}{238.427}{so}%
\TextAt{8}{295.427}{238.427}{is}%
\TextAt{8}{307.491}{238.427}{a}%
\TextAt{8}{317.548}{238.427}{profinite}%
\TextAt{8}{364.677}{238.427}{ring.}%
\TextAt{8}{394.195}{238.427}{The}%
\TextAt{8}{418.55}{238.427}{rings}%
\TextAt{14}{447.531}{238.427}{o}%
\TextAt{9}{453.375}{238.427}{\Character{47}}%
\TextAt{14}{459.228}{238.427}{p}%
\TextAt{21}{465.211}{242.765}{n}%
\TextAt{8}{470.848}{238.427}{,}%
\TextAt{13}{478.373}{238.427}{n}%
\NativeMathAt{489.193}{238.427}{11.955}{\in}%
\TextAt{18}{500.997}{238.427}{N}%
\TextAt{8}{509.631}{238.427}{,}%
\TextAt{8}{517.156}{238.427}{are}%
\TextAt{8}{79.226}{221.381}{finite}%
\TextAt{8}{109.138}{221.381}{and}%
\TextAt{8}{131.901}{221.381}{we}%
\TextAt{8}{149.128}{221.381}{have}%
\TextAt{8}{176.123}{221.381}{a}%
\TextAt{8}{185.878}{221.381}{topological}%
\TextAt{8}{245.403}{221.381}{isomorphism}%
\TextAt{14}{205.762}{192.38}{o}%
\NativeMathAt{214.927}{192.379}{11.955}{\cong}%
\NativeMathAt{227.546}{192.38}{11.955}{\varprojlim}%
\TextAt{21}{233.106}{180.265}{n}%
\TextAt{14}{245.798}{192.38}{o}%
\TextAt{9}{251.642}{192.38}{\Character{47}}%
\TextAt{14}{257.495}{192.38}{p}%
\TextAt{21}{263.478}{197.316}{n}%
\TextAt{13}{269.114}{192.38}{\Character{59}}%
\TextAt{13}{298.249}{192.38}{a}%
\NativeArrowAt{307.734}{192.38}{11.955}{19.26}{\longmapsto}%
\TextAt{9}{330.316}{192.38}{\Character{40}}%
\TextAt{13}{334.869}{192.38}{a}%
\TextAt{8}{344.335}{192.38}{mod}%
\TextAt{14}{369.772}{192.38}{p}%
\TextAt{21}{375.756}{197.316}{n}%
\TextAt{9}{381.392}{192.38}{\Character{41}}%
\TextAt{21}{385.944}{190.586}{n}%
\NativeMathAt{391.082}{190.586}{7.97}{\in}%
\TextAt{25}{396.728}{190.586}{N}%
\TextAt{13}{402.982}{192.38}{\Character{58}}%
\TextAt{8}{79.226}{157.702}{The}%
\TextAt{8}{102.467}{157.702}{group}%
\TextAt{8}{134.829}{157.702}{of}%
\TextAt{8}{147.346}{157.702}{units}%
\TextAt{13}{175.859}{157.702}{U}%
\TextAt{9}{188.379}{157.702}{\Character{61}}%
\TextAt{14}{200.804}{157.702}{o}%
\TextAt{3}{206.648}{162.04}{\Character{3}}%
\TextAt{8}{214.465}{157.702}{is}%
\TextAt{8}{225.416}{157.702}{closed}%
\TextAt{8}{259.129}{157.702}{in}%
\TextAt{14}{271.969}{157.702}{o}%
\TextAt{8}{277.813}{157.702}{,}%
\TextAt{8}{284.311}{157.702}{hence}%
\TextAt{8}{316.004}{157.702}{Hausdorff}%
\TextAt{8}{368.583}{157.702}{and}%
\TextAt{8}{390.532}{157.702}{compact,}%
\TextAt{8}{439.962}{157.702}{and}%
\TextAt{8}{461.911}{157.702}{the}%
\TextAt{8}{481.255}{157.702}{subgroups}%
\TextAt{13}{79.226}{140.656}{U}%
\TextAt{17}{88.425}{144.994}{\Character{40}}%
\TextAt{21}{91.718}{144.994}{n}%
\TextAt{17}{96.857}{144.994}{\Character{41}}%
\TextAt{9}{103.969}{140.656}{\Character{61}}%
\TextAt{9}{116.379}{140.656}{1}%
\TextAt{9}{124.879}{140.656}{\Character{43}}%
\TextAt{14}{136.665}{140.656}{p}%
\TextAt{21}{142.648}{144.994}{n}%
\TextAt{8}{152.182}{140.656}{form}%
\TextAt{8}{179.822}{140.656}{a}%
\TextAt{8}{189.578}{140.656}{basis}%
\TextAt{8}{218.318}{140.656}{of}%
\TextAt{8}{231.648}{140.656}{neighbourhoods}%
\TextAt{8}{316.267}{140.656}{of}%
\TextAt{9}{329.597}{140.656}{1}%
\NativeMathAt{338.771}{140.656}{11.955}{\in}%
\TextAt{13}{350.062}{140.656}{U}%
\TextAt{8}{359.261}{140.656}{.}%
\TextAt{8}{367.709}{140.656}{Thus}%
\TextAt{13}{266.065}{111.655}{U}%
\NativeMathAt{278.585}{111.655}{11.955}{\cong}%
\NativeMathAt{291.205}{111.655}{11.955}{\varprojlim}%
\TextAt{21}{296.765}{99.54}{n}%
\TextAt{13}{309.456}{111.655}{U}%
\TextAt{9}{318.655}{111.655}{\Character{47}}%
\TextAt{13}{324.508}{111.655}{U}%
\TextAt{17}{333.707}{116.591}{\Character{40}}%
\TextAt{21}{337.001}{116.591}{n}%
\TextAt{17}{342.139}{116.591}{\Character{41}}%
\TextAt{8}{79.226}{76.869}{is}%
\TextAt{8}{90.99}{76.869}{also}%
\TextAt{8}{114.47}{76.869}{a}%
\TextAt{8}{124.225}{76.869}{profinite}%
\TextAt{8}{171.054}{76.869}{group.}%
\TextAt{8}{208.78}{76.869}{In}%
\TextAt{8}{223.402}{76.869}{fact,}%
\TextAt{8}{249.739}{76.869}{we}%
\TextAt{8}{266.966}{76.869}{have}%
\TextAt{8}{293.961}{76.869}{seen}%
\TextAt{8}{319.378}{76.869}{all}%
\TextAt{8}{335.637}{76.869}{this}%
\TextAt{8}{358.459}{76.869}{already}%
\TextAt{8}{399.764}{76.869}{in}%
\TextAt{8}{413.417}{76.869}{chap.}%
\TextAt{8}{445.608}{76.869}{II,}%
\TextAt{8}{461.521}{76.869}{§4.}%
\TextAt{8}{297.211}{43.033}{248}%
\Destination{78.23}{809.815}{nav.page.248}%
\Destination{79.36}{742.055}{nav.print.271}%
\Hotlink{506.344}{524.355}{6.858}{11.047}{nav.section.4.1}%
\Hotlink{281.499}{290.355}{9.751}{11.047}{nav.chapter.2}%
\Hotlink{303.09}{290.355}{6.858}{11.047}{nav.section.2.1}%
\Hotlink{445.108}{74.655}{9.763}{11.047}{nav.chapter.2}%
\Hotlink{466.7}{74.655}{6.858}{11.047}{nav.section.2.4}%
\Hotlink{410.147}{741.855}{25.078}{11.047}{nav.number.4.2.3}%
\Hotlink{303.053}{473.155}{25.078}{11.047}{nav.number.4.2.8}%
}

% Page 261 (body page 249)
\Page{%
\TextAt{2}{79.226}{796.943}{IV.}%
\TextAt{2}{98.012}{796.943}{Abstract}%
\TextAt{2}{143.252}{796.943}{Class}%
\TextAt{2}{171.844}{796.943}{Field}%
\TextAt{2}{199.575}{796.943}{Theory}%
\Vector{q 0 G 0 g 0.299 w 7.226 18.729 m 460.769 18.729 l S Q}%
\TextAt{12}{97.159}{761.124}{Example}%
\TextAt{12}{153.362}{761.124}{4:}%
\TextAt{8}{170.24}{761.124}{The}%
\TextAt{8}{194.898}{761.124}{rings}%
\TextAt{18}{224.183}{761.124}{Z}%
\TextAt{9}{232.153}{761.124}{\Character{47}}%
\TextAt{13}{238.006}{761.124}{n}%
\TextAt{18}{244.994}{761.124}{Z}%
\TextAt{8}{252.964}{761.124}{,}%
\TextAt{13}{260.869}{761.124}{n}%
\NativeMathAt{272.207}{761.124}{11.955}{\in}%
\TextAt{18}{284.528}{761.124}{N}%
\TextAt{8}{293.163}{761.124}{,}%
\TextAt{8}{301.067}{761.124}{form}%
\TextAt{8}{329.312}{761.124}{a}%
\TextAt{8}{339.672}{761.124}{projective}%
\TextAt{8}{394.254}{761.124}{system}%
\TextAt{8}{433.677}{761.124}{with}%
\TextAt{8}{460.941}{761.124}{respect}%
\TextAt{8}{501.596}{761.124}{to}%
\TextAt{8}{516.51}{761.124}{the}%
\ErrataAt{79.22}{744.055}{453.64}{projections $\mathbb Z/n\mathbb Z\to\mathbb Z/m\mathbb Z$, $m\mid n$, where the ordering on $\mathbb N$ is given by divisibility,}%
\TextAt{13}{79.226}{727.032}{n}%
\TextAt{16}{89.535}{727.032}{\Character{106}}%
\TextAt{13}{96.176}{727.032}{m}%
\TextAt{8}{106.416}{727.032}{.}%
\TextAt{8}{130.059}{727.032}{The}%
\TextAt{8}{154.113}{727.032}{projective}%
\TextAt{8}{208.09}{727.032}{limit}%
\TextAt{20}{271.288}{713.074}{\Character{98}}%
\TextAt{18}{270.623}{709.986}{Z}%
\TextAt{9}{281.914}{709.986}{\Character{61}}%
\NativeMathAt{294.34}{709.986}{11.955}{\varprojlim}%
\TextAt{21}{299.9}{697.871}{n}%
\TextAt{18}{312.591}{709.986}{Z}%
\TextAt{9}{320.562}{709.986}{\Character{47}}%
\TextAt{13}{326.415}{709.986}{n}%
\TextAt{18}{333.402}{709.986}{Z}%
\TextAt{8}{79.226}{681.134}{was}%
\TextAt{8}{102.861}{681.134}{originally}%
\TextAt{8}{155.713}{681.134}{called}%
\TextAt{8}{190.012}{681.134}{the}%
\TextAt{12}{211.303}{681.134}{Prüfer}%
\TextAt{12}{255.001}{681.134}{ring}%
\TextAt{8}{278.445}{681.134}{,}%
\TextAt{8}{287.013}{681.134}{whereas}%
\TextAt{8}{332.43}{681.134}{nowadays}%
\TextAt{8}{386.323}{681.134}{it}%
\TextAt{8}{399.161}{681.134}{has}%
\TextAt{8}{421.17}{681.134}{become}%
\TextAt{8}{464.243}{681.134}{customary}%
\TextAt{8}{522.357}{681.134}{to}%
\TextAt{8}{79.226}{664.088}{refer}%
\TextAt{8}{105.825}{664.088}{to}%
\TextAt{8}{119.751}{664.088}{it}%
\TextAt{8}{131.07}{664.088}{by}%
\TextAt{8}{146.945}{664.088}{the}%
\TextAt{8}{166.717}{664.088}{somewhat}%
\TextAt{8}{221.028}{664.088}{curt}%
\TextAt{8}{245.355}{664.088}{abbreviation}%
\TextAt{8}{312.947}{664.088}{“zed-hat”}%
\TextAt{8}{365.226}{664.088}{(or}%
\TextAt{8}{383.706}{664.088}{“zee-hat”).}%
\TextAt{8}{444.044}{664.088}{This}%
\TextAt{8}{470.379}{664.088}{ring}%
\TextAt{8}{494.06}{664.088}{is}%
\TextAt{8}{505.44}{664.088}{going}%
\TextAt{8}{79.226}{647.042}{to}%
\TextAt{8}{93.576}{647.042}{occupy}%
\TextAt{8}{132.961}{647.042}{quite}%
\TextAt{8}{162.59}{647.042}{an}%
\TextAt{8}{178.888}{647.042}{important}%
\TextAt{8}{234.221}{647.042}{position}%
\TextAt{8}{278.877}{647.042}{in}%
\TextAt{8}{292.57}{647.042}{what}%
\TextAt{8}{321.876}{647.042}{follows.}%
\TextAt{8}{364.983}{647.042}{It}%
\TextAt{8}{377.695}{647.042}{contains}%
\TextAt{18}{423.655}{647.042}{Z}%
\TextAt{8}{435.562}{647.042}{as}%
\TextAt{8}{449.972}{647.042}{a}%
\TextAt{8}{459.767}{647.042}{dense}%
\TextAt{8}{491.727}{647.042}{subring.}%
\TextAt{8}{79.226}{629.996}{The}%
\TextAt{8}{102.808}{629.996}{groups}%
\TextAt{13}{140.126}{629.996}{n}%
\TextAt{20}{147.778}{633.084}{\Character{98}}%
\TextAt{18}{147.113}{629.996}{Z}%
\TextAt{8}{155.084}{629.996}{,}%
\TextAt{13}{161.855}{629.996}{n}%
\NativeMathAt{172.163}{629.996}{11.955}{\in}%
\TextAt{18}{183.454}{629.996}{N}%
\TextAt{8}{192.089}{629.996}{,}%
\TextAt{8}{198.86}{629.996}{are}%
\TextAt{8}{217.899}{629.996}{precisely}%
\TextAt{8}{265.283}{629.996}{the}%
\TextAt{8}{284.967}{629.996}{open}%
\TextAt{8}{312.781}{629.996}{subgroups}%
\TextAt{8}{367.721}{629.996}{of}%
\TextAt{20}{381.243}{633.084}{\Character{98}}%
\TextAt{18}{380.579}{629.996}{Z}%
\TextAt{8}{388.549}{629.996}{,}%
\TextAt{8}{395.321}{629.996}{and}%
\TextAt{8}{417.611}{629.996}{it}%
\TextAt{8}{428.843}{629.996}{is}%
\TextAt{8}{440.135}{629.996}{easy}%
\TextAt{8}{465.414}{629.996}{to}%
\TextAt{8}{479.252}{629.996}{verify}%
\TextAt{8}{511.298}{629.996}{that}%
\TextAt{20}{268.285}{604.779}{\Character{98}}%
\TextAt{18}{267.621}{601.691}{Z}%
\TextAt{9}{275.591}{601.691}{\Character{47}}%
\TextAt{13}{281.444}{601.691}{n}%
\TextAt{20}{289.096}{604.779}{\Character{98}}%
\TextAt{18}{288.432}{601.691}{Z}%
\NativeMathAt{299.723}{601.691}{11.955}{\cong}%
\TextAt{18}{312.342}{601.691}{Z}%
\TextAt{9}{320.312}{601.691}{\Character{47}}%
\TextAt{13}{326.165}{601.691}{n}%
\TextAt{18}{333.153}{601.691}{Z}%
\TextAt{13}{341.123}{601.691}{\Character{58}}%
\TextAt{8}{79.226}{573.386}{Taking,}%
\TextAt{8}{123.593}{573.386}{for}%
\TextAt{8}{143.153}{573.386}{each}%
\TextAt{8}{171.165}{573.386}{natural}%
\TextAt{8}{213.822}{573.386}{number}%
\TextAt{13}{258.094}{573.386}{n}%
\TextAt{8}{265.081}{573.386}{,}%
\TextAt{8}{274.324}{573.386}{the}%
\TextAt{8}{296.155}{573.386}{prime}%
\TextAt{8}{330.994}{573.386}{factorization}%
\TextAt{13}{400.646}{573.386}{n}%
\TextAt{9}{413.809}{573.386}{\Character{61}}%
\TextAt{20}{429.089}{582.352}{\Character{81}}%
\TextAt{21}{440.38}{569.899}{p}%
\TextAt{13}{447.133}{573.386}{p}%
\TextAt{21}{453.008}{577.724}{v}%
\TextAt{22}{457.132}{576.728}{p}%
\TextAt{8}{461.968}{573.386}{,}%
\TextAt{8}{471.211}{573.386}{the}%
\TextAt{8}{493.042}{573.386}{Chinese}%
\TextAt{8}{79.226}{556.34}{remainder}%
\TextAt{8}{134.507}{556.34}{theorem}%
\TextAt{8}{180.032}{556.34}{implies}%
\TextAt{8}{219.759}{556.34}{the}%
\TextAt{8}{239.916}{556.34}{decomposition}%
\TextAt{18}{255.063}{528.035}{Z}%
\TextAt{9}{263.033}{528.035}{\Character{47}}%
\TextAt{13}{268.886}{528.035}{n}%
\TextAt{18}{275.874}{528.035}{Z}%
\NativeMathAt{287.165}{528.034}{11.955}{\cong}%
\TextAt{20}{299.784}{539.392}{\Character{89}}%
\TextAt{21}{305.291}{514.286}{p}%
\TextAt{18}{317.052}{528.035}{Z}%
\TextAt{9}{325.023}{528.035}{\Character{47}}%
\TextAt{13}{330.876}{528.035}{p}%
\TextAt{21}{336.751}{532.971}{v}%
\TextAt{22}{340.874}{531.974}{p}%
\TextAt{18}{345.711}{528.035}{Z}%
\TextAt{13}{353.681}{528.035}{\Character{59}}%
\TextAt{8}{79.226}{490.989}{and}%
\TextAt{8}{101.989}{490.989}{passing}%
\TextAt{8}{143.091}{490.989}{to}%
\TextAt{8}{157.401}{490.989}{the}%
\TextAt{8}{177.557}{490.989}{projective}%
\TextAt{8}{231.535}{490.989}{limit,}%
\TextAt{20}{278.081}{465.772}{\Character{98}}%
\TextAt{18}{277.417}{462.684}{Z}%
\NativeMathAt{288.708}{462.683}{11.955}{\cong}%
\TextAt{20}{301.327}{474.041}{\Character{89}}%
\TextAt{21}{306.834}{448.935}{p}%
\TextAt{18}{318.596}{462.684}{Z}%
\TextAt{21}{326.566}{460.89}{p}%
\TextAt{13}{331.327}{462.684}{\Character{58}}%
\TextAt{8}{79.226}{422.212}{This}%
\TextAt{8}{106.838}{422.212}{takes}%
\TextAt{8}{137.713}{422.212}{the}%
\TextAt{8}{158.762}{422.212}{natural}%
\TextAt{8}{200.636}{422.212}{embedding}%
\TextAt{8}{260.706}{422.212}{of}%
\TextAt{18}{274.928}{422.212}{Z}%
\TextAt{8}{287.688}{422.212}{into}%
\TextAt{20}{312.987}{425.3}{\Character{98}}%
\TextAt{18}{312.323}{422.212}{Z}%
\TextAt{8}{325.082}{422.212}{to}%
\TextAt{8}{340.284}{422.212}{the}%
\TextAt{8}{361.333}{422.212}{diagonal}%
\TextAt{8}{409.065}{422.212}{embedding}%
\TextAt{18}{469.135}{422.212}{Z}%
\NativeMathAt{481.946}{422.212}{11.955}{\rightarrow}%
\TextAt{20}{498.742}{431.178}{\Character{81}}%
\TextAt{21}{510.033}{418.725}{p}%
\TextAt{18}{516.787}{422.212}{Z}%
\TextAt{21}{524.757}{420.418}{p}%
\TextAt{8}{529.518}{422.212}{,}%
\TextAt{13}{79.226}{405.166}{a}%
\NativeArrowAt{88.692}{405.166}{11.955}{11.955}{\mapsto}%
\TextAt{9}{103.968}{405.166}{\Character{40}}%
\TextAt{13}{108.52}{405.166}{a\Character{59}}%
\TextAt{13}{119.889}{405.166}{a\Character{59}}%
\TextAt{13}{131.271}{405.166}{a\Character{59}}%
\TextAt{13}{142.652}{405.166}{\Character{58}}%
\TextAt{13}{147.877}{405.166}{\Character{58}}%
\TextAt{13}{153.113}{405.166}{\Character{58}}%
\TextAt{9}{156.428}{405.166}{\Character{41}}%
\TextAt{8}{160.98}{405.166}{.}%
\TextAt{12}{97.159}{385.131}{Example}%
\TextAt{12}{152.667}{385.131}{5:}%
\TextAt{8}{168.336}{385.131}{For}%
\TextAt{8}{189.305}{385.131}{the}%
\TextAt{8}{209.462}{385.131}{field}%
\TextAt{18}{234.819}{385.131}{F}%
\TextAt{21}{242.125}{383.338}{q}%
\TextAt{8}{250.587}{385.131}{with}%
\TextAt{13}{277.247}{385.131}{q}%
\TextAt{8}{286.764}{385.131}{elements,}%
\TextAt{8}{337.872}{385.131}{we}%
\TextAt{8}{355.099}{385.131}{get}%
\TextAt{8}{374.61}{385.131}{isomorphisms}%
\TextAt{13}{253.441}{356.826}{G}%
\TextAt{9}{262.674}{356.826}{\Character{40}}%
\TextAt{18}{267.227}{356.826}{F}%
\TextAt{21}{274.533}{355.033}{q}%
\TextAt{22}{278.599}{357.302}{n}%
\TextAt{16}{287.517}{356.826}{\Character{106}}%
\TextAt{18}{294.159}{356.826}{F}%
\TextAt{21}{301.465}{355.033}{q}%
\TextAt{9}{306.03}{356.826}{\Character{41}}%
\NativeMathAt{313.903}{356.826}{11.955}{\cong}%
\TextAt{18}{326.522}{356.826}{Z}%
\TextAt{9}{334.492}{356.826}{\Character{47}}%
\TextAt{13}{340.345}{356.826}{n}%
\TextAt{18}{347.333}{356.826}{Z}%
\TextAt{13}{355.303}{356.826}{\Character{59}}%
\TextAt{8}{79.226}{328.521}{one}%
\TextAt{8}{100.808}{328.521}{for}%
\TextAt{8}{118.816}{328.521}{every}%
\TextAt{13}{149.831}{328.521}{n}%
\NativeMathAt{160.348}{328.521}{11.955}{\in}%
\TextAt{18}{171.849}{328.521}{N}%
\TextAt{8}{180.483}{328.521}{,}%
\TextAt{8}{187.785}{328.521}{by}%
\TextAt{8}{204.167}{328.521}{mapping}%
\TextAt{8}{252.421}{328.521}{the}%
\TextAt{8}{272.701}{328.521}{Frobenius}%
\TextAt{8}{326.693}{328.521}{automorphism}%
\TextAt{13}{404.285}{328.521}{\Character{39}}%
\TextAt{21}{411.954}{326.728}{n}%
\TextAt{8}{421.61}{328.521}{to}%
\TextAt{9}{436.043}{328.521}{1}%
\TextAt{8}{445.217}{328.521}{mod}%
\TextAt{13}{470.655}{328.521}{n}%
\TextAt{18}{477.643}{328.521}{Z}%
\TextAt{8}{485.613}{328.521}{.}%
\TextAt{8}{494.429}{328.521}{Passing}%
\TextAt{8}{79.226}{311.475}{to}%
\TextAt{8}{93.536}{311.475}{the}%
\TextAt{8}{113.693}{311.475}{projective}%
\TextAt{8}{167.67}{311.475}{limit}%
\TextAt{8}{195.634}{311.475}{gives}%
\TextAt{8}{224.314}{311.475}{an}%
\TextAt{8}{240.573}{311.475}{isomorphism}%
\TextAt{13}{268.021}{283.17}{G}%
\TextAt{9}{277.255}{283.17}{\Character{40}}%
\Vector{q 0 G 0 g 0.478 w 209.807 -480.921 m 217.113 -480.921 l S Q}%
\TextAt{18}{281.807}{283.17}{F}%
\TextAt{21}{289.113}{281.377}{q}%
\TextAt{16}{296.999}{283.17}{\Character{106}}%
\TextAt{18}{303.641}{283.17}{F}%
\TextAt{21}{310.947}{281.377}{q}%
\TextAt{9}{315.512}{283.17}{\Character{41}}%
\NativeMathAt{323.385}{283.17}{11.955}{\cong}%
\TextAt{20}{336.668}{286.258}{\Character{98}}%
\TextAt{18}{336.004}{283.17}{Z}%
\TextAt{8}{79.226}{254.865}{which}%
\TextAt{8}{113.655}{254.865}{sends}%
\TextAt{8}{145.932}{254.865}{the}%
\TextAt{8}{167.03}{254.865}{Frobenius}%
\TextAt{8}{221.843}{254.865}{automorphism}%
\TextAt{13}{300.254}{254.865}{\Character{39}}%
\NativeMathAt{312.849}{254.865}{11.955}{\in}%
\TextAt{13}{325.746}{254.865}{G}%
\TextAt{9}{334.98}{254.865}{\Character{40}}%
\Vector{q 0 G 0 g 0.478 w 267.532 -509.226 m 274.838 -509.226 l S Q}%
\TextAt{18}{339.532}{254.865}{F}%
\TextAt{21}{346.838}{253.072}{q}%
\TextAt{16}{356.329}{254.865}{\Character{106}}%
\TextAt{18}{364.577}{254.865}{F}%
\TextAt{21}{371.883}{253.072}{q}%
\TextAt{9}{376.448}{254.865}{\Character{41}}%
\TextAt{8}{385.84}{254.865}{to}%
\TextAt{9}{401.092}{254.865}{1}%
\NativeMathAt{411.872}{254.865}{11.955}{\in}%
\TextAt{20}{425.433}{257.953}{\Character{98}}%
\TextAt{18}{424.769}{254.865}{Z}%
\TextAt{8}{432.739}{254.865}{,}%
\TextAt{8}{441.066}{254.865}{and}%
\TextAt{8}{464.771}{254.865}{the}%
\TextAt{8}{485.869}{254.865}{subgroup}%
\TextAt{16}{79.226}{237.819}{\Character{104}}%
\TextAt{13}{83.875}{237.819}{\Character{39}}%
\TextAt{16}{91.544}{237.819}{\Character{105}}%
\TextAt{9}{100.534}{237.819}{\Character{61}}%
\TextAt{16}{113.98}{237.819}{\Character{102}}%
\TextAt{13}{119.958}{237.819}{\Character{39}}%
\TextAt{21}{127.626}{242.157}{n}%
\TextAt{16}{137.604}{237.819}{\Character{106}}%
\TextAt{13}{145.266}{237.819}{n}%
\NativeMathAt{156.595}{237.819}{11.955}{\in}%
\TextAt{18}{168.907}{237.819}{Z}%
\TextAt{16}{176.877}{237.819}{\Character{103}}%
\TextAt{8}{187.35}{237.819}{onto}%
\TextAt{8}{214.298}{237.819}{the}%
\TextAt{8}{235.054}{237.819}{dense}%
\TextAt{8}{267.573}{237.819}{(but}%
\TextAt{8}{294.186}{237.819}{not}%
\TextAt{8}{315.599}{237.819}{closed)}%
\TextAt{8}{355.279}{237.819}{subgroup}%
\TextAt{18}{406.676}{237.819}{Z}%
\TextAt{8}{419.142}{237.819}{of}%
\TextAt{20}{433.735}{240.907}{\Character{98}}%
\TextAt{18}{433.071}{237.819}{Z}%
\TextAt{8}{441.041}{237.819}{.}%
\TextAt{8}{451.286}{237.819}{Given}%
\TextAt{8}{485.777}{237.819}{this,}%
\TextAt{8}{512.6}{237.819}{it}%
\TextAt{8}{524.903}{237.819}{is}%
\TextAt{8}{79.226}{220.773}{now}%
\TextAt{8}{103.207}{220.773}{clear,}%
\TextAt{8}{134.096}{220.773}{in}%
\TextAt{8}{147.341}{220.773}{the}%
\TextAt{8}{167.09}{220.773}{example}%
\TextAt{8}{212.53}{220.773}{at}%
\TextAt{8}{226.433}{220.773}{the}%
\TextAt{8}{246.181}{220.773}{beginning}%
\TextAt{8}{299.428}{220.773}{of}%
\TextAt{8}{312.351}{220.773}{this}%
\TextAt{8}{334.765}{220.773}{chapter,}%
\TextAt{8}{379.642}{220.773}{how}%
\TextAt{8}{403.623}{220.773}{we}%
\TextAt{8}{420.442}{220.773}{were}%
\TextAt{8}{447.017}{220.773}{able}%
\TextAt{8}{471.321}{220.773}{to}%
\TextAt{8}{485.224}{220.773}{construct}%
\TextAt{8}{79.226}{203.727}{an}%
\TextAt{8}{96.336}{203.727}{element}%
\TextAt{13}{140.428}{203.727}{\Character{32}}%
\NativeMathAt{153.238}{203.727}{11.955}{\in}%
\TextAt{13}{165.978}{203.727}{G}%
\TextAt{9}{175.212}{203.727}{\Character{40}}%
\Vector{q 0 G 0 g 0.478 w 107.764 -560.364 m 115.07 -560.364 l S Q}%
\TextAt{18}{179.764}{203.727}{F}%
\TextAt{21}{187.07}{201.934}{q}%
\TextAt{16}{196.406}{203.727}{\Character{106}}%
\TextAt{18}{204.497}{203.727}{F}%
\TextAt{21}{211.803}{201.934}{q}%
\TextAt{9}{216.368}{203.727}{\Character{41}}%
\TextAt{8}{225.668}{203.727}{which}%
\TextAt{8}{260.005}{203.727}{did}%
\TextAt{8}{281.012}{203.727}{not}%
\TextAt{8}{302.677}{203.727}{belong}%
\TextAt{8}{340.923}{203.727}{to}%
\TextAt{16}{356.084}{203.727}{\Character{104}}%
\TextAt{13}{360.733}{203.727}{\Character{39}}%
\TextAt{16}{368.402}{203.727}{\Character{105}}%
\TextAt{8}{373.051}{203.727}{.}%
\TextAt{8}{384.051}{203.727}{In}%
\TextAt{8}{399.523}{203.727}{fact,}%
\TextAt{8}{426.923}{203.727}{looking}%
\TextAt{8}{468.756}{203.727}{at}%
\TextAt{8}{483.917}{203.727}{it}%
\TextAt{8}{496.472}{203.727}{via}%
\TextAt{8}{516.511}{203.727}{the}%
\TextAt{8}{79.226}{186.681}{isomorphism}%
\TextAt{13}{147.646}{186.681}{G}%
\TextAt{9}{156.879}{186.681}{\Character{40}}%
\Vector{q 0 G 0 g 0.478 w 89.432 -577.41 m 96.737 -577.41 l S Q}%
\TextAt{18}{161.432}{186.681}{F}%
\TextAt{21}{168.737}{184.888}{q}%
\TextAt{16}{176.623}{186.681}{\Character{106}}%
\TextAt{18}{183.265}{186.681}{F}%
\TextAt{21}{190.571}{184.888}{q}%
\TextAt{9}{195.136}{186.681}{\Character{41}}%
\NativeMathAt{203.009}{186.681}{11.955}{\cong}%
\TextAt{20}{216.292}{189.769}{\Character{98}}%
\TextAt{18}{215.628}{186.681}{Z}%
\TextAt{8}{223.598}{186.681}{,}%
\TextAt{8}{230.748}{186.681}{what}%
\TextAt{8}{260.014}{186.681}{we}%
\TextAt{8}{277.241}{186.681}{did}%
\TextAt{8}{297.398}{186.681}{amounted}%
\TextAt{8}{351.71}{186.681}{to}%
\TextAt{8}{366.02}{186.681}{writing}%
\TextAt{8}{406.345}{186.681}{down}%
\TextAt{8}{437.237}{186.681}{the}%
\TextAt{8}{457.394}{186.681}{element}%
\TextAt{9}{218.576}{158.376}{\Character{40}}%
\TextAt{13}{223.128}{158.376}{\Character{58}}%
\TextAt{13}{228.352}{158.376}{\Character{58}}%
\TextAt{13}{233.589}{158.376}{\Character{58}}%
\TextAt{13}{238.825}{158.376}{\Character{59}}%
\TextAt{9}{244.105}{158.376}{0}%
\TextAt{13}{249.958}{158.376}{\Character{59}}%
\TextAt{9}{255.202}{158.376}{0}%
\TextAt{13}{261.055}{158.376}{\Character{59}}%
\TextAt{9}{266.299}{158.376}{1}%
\TextAt{21}{272.152}{156.583}{p}%
\TextAt{13}{276.913}{158.376}{\Character{59}}%
\TextAt{9}{282.157}{158.376}{0}%
\TextAt{13}{288.01}{158.376}{\Character{59}}%
\TextAt{9}{293.254}{158.376}{0}%
\TextAt{13}{299.107}{158.376}{\Character{59}}%
\TextAt{13}{304.331}{158.376}{\Character{58}}%
\TextAt{13}{309.568}{158.376}{\Character{58}}%
\TextAt{13}{314.804}{158.376}{\Character{58}}%
\TextAt{9}{318.091}{158.376}{\Character{41}}%
\NativeMathAt{325.965}{158.376}{11.955}{\in}%
\TextAt{20}{337.255}{169.733}{\Character{89}}%
\TextAt{21}{343.164}{144.273}{\Character{96}}%
\TextAt{18}{354.524}{158.376}{Z}%
\TextAt{21}{362.494}{156.583}{\Character{96}}%
\TextAt{9}{369.772}{158.376}{\Character{61}}%
\TextAt{20}{382.862}{161.464}{\Character{98}}%
\TextAt{18}{382.198}{158.376}{Z}%
\TextAt{13}{390.168}{158.376}{\Character{59}}%
\TextAt{8}{79.226}{122.525}{which}%
\TextAt{8}{112.713}{122.525}{does}%
\TextAt{8}{139.11}{122.525}{not}%
\TextAt{8}{159.924}{122.525}{belong}%
\TextAt{8}{197.319}{122.525}{to}%
\TextAt{18}{211.63}{122.525}{Z}%
\TextAt{8}{219.6}{122.525}{.}%
\TextAt{8}{297.211}{43.033}{249}%
\Destination{78.23}{809.815}{nav.page.249}%
\Destination{79.36}{776.155}{nav.print.272}%
%
\Destination{79.226}{758.055}{errata.25}%
\Destination{97.16}{774}{nav.example.4.2.4}%
\Destination{97.16}{398}{nav.example.4.2.5}%
}

% Page 262 (body page 250)
\Page{%
\TextAt{2}{79.226}{796.943}{IV.}%
\TextAt{2}{98.012}{796.943}{Abstract}%
\TextAt{2}{143.252}{796.943}{Class}%
\TextAt{2}{171.844}{796.943}{Field}%
\TextAt{2}{199.575}{796.943}{Theory}%
\Vector{q 0 G 0 g 0.299 w 7.226 18.729 m 460.769 18.729 l S Q}%
\TextAt{12}{97.159}{761.124}{Example}%
\TextAt{12}{153.945}{761.124}{6:}%
\TextAt{8}{171.838}{761.124}{Let}%
\TextAt{20}{195.247}{764.212}{\Character{101}}%
\TextAt{18}{193.918}{761.124}{Q}%
\TextAt{16}{208.432}{761.124}{\Character{106}}%
\TextAt{18}{216.968}{761.124}{Q}%
\TextAt{8}{231.275}{761.124}{be}%
\TextAt{8}{248.311}{761.124}{the}%
\TextAt{8}{269.579}{761.124}{extension}%
\TextAt{8}{322.456}{761.124}{obtained}%
\TextAt{8}{371.699}{761.124}{by}%
\TextAt{8}{389.07}{761.124}{adjoining}%
\TextAt{8}{441.242}{761.124}{all}%
\TextAt{8}{458.612}{761.124}{roots}%
\TextAt{8}{489.385}{761.124}{of}%
\TextAt{8}{503.826}{761.124}{unity.}%
\TextAt{8}{79.226}{744.078}{Its}%
\TextAt{8}{96.597}{744.078}{Galois}%
\TextAt{8}{132.595}{744.078}{group}%
\TextAt{13}{165.854}{744.078}{G}%
\TextAt{9}{175.088}{744.078}{\Character{40}}%
\TextAt{20}{180.968}{747.166}{\Character{101}}%
\TextAt{18}{179.64}{744.078}{Q}%
\TextAt{16}{192.402}{744.078}{\Character{106}}%
\TextAt{18}{199.187}{744.078}{Q}%
\TextAt{9}{208.486}{744.078}{\Character{41}}%
\TextAt{8}{217.019}{744.078}{is}%
\TextAt{8}{228.867}{744.078}{then}%
\TextAt{8}{255.611}{744.078}{canonically}%
\TextAt{8}{316.511}{744.078}{isomorphic}%
\TextAt{8}{375.845}{744.078}{(as}%
\TextAt{8}{394.854}{744.078}{a}%
\TextAt{8}{404.693}{744.078}{topological}%
\TextAt{8}{464.302}{744.078}{group)}%
\TextAt{8}{502.116}{744.078}{to}%
\TextAt{8}{516.51}{744.078}{the}%
\TextAt{8}{79.226}{727.032}{group}%
\TextAt{8}{112.402}{727.032}{of}%
\TextAt{8}{125.732}{727.032}{units}%
\TextAt{20}{155.722}{730.12}{\Character{98}}%
\TextAt{18}{155.058}{727.032}{Z}%
\TextAt{3}{163.028}{731.37}{\Character{3}}%
\NativeMathAt{171.081}{727.032}{11.955}{\cong}%
\TextAt{20}{183.7}{735.998}{\Character{81}}%
\TextAt{21}{194.991}{723.545}{p}%
\TextAt{18}{201.745}{727.032}{Z}%
\TextAt{3}{209.715}{731.37}{\Character{3}}%
\TextAt{21}{209.715}{724.076}{p}%
\TextAt{8}{218.373}{727.032}{of}%
\TextAt{20}{232.367}{730.12}{\Character{98}}%
\TextAt{18}{231.703}{727.032}{Z}%
\TextAt{8}{239.673}{727.032}{,}%
\TextAt{13}{266.602}{698.031}{G}%
\TextAt{9}{275.835}{698.031}{\Character{40}}%
\TextAt{20}{281.716}{701.119}{\Character{101}}%
\TextAt{18}{280.388}{698.031}{Q}%
\TextAt{16}{293.007}{698.031}{\Character{106}}%
\TextAt{18}{299.649}{698.031}{Q}%
\TextAt{9}{308.947}{698.031}{\Character{41}}%
\NativeMathAt{316.82}{698.03}{11.955}{\cong}%
\TextAt{20}{330.104}{701.119}{\Character{98}}%
\TextAt{18}{329.44}{698.031}{Z}%
\TextAt{3}{337.41}{702.967}{\Character{3}}%
\TextAt{13}{342.142}{698.031}{\Character{58}}%
\TextAt{8}{79.226}{669.029}{This}%
\TextAt{8}{106.557}{669.029}{isomorphism}%
\TextAt{8}{175.587}{669.029}{is}%
\TextAt{8}{187.961}{669.029}{obtained}%
\TextAt{8}{236.703}{669.029}{by}%
\TextAt{8}{253.573}{669.029}{passing}%
\TextAt{8}{295.286}{669.029}{to}%
\TextAt{8}{310.206}{669.029}{the}%
\TextAt{8}{330.974}{669.029}{projective}%
\TextAt{8}{385.562}{669.029}{limit}%
\TextAt{8}{414.136}{669.029}{from}%
\TextAt{8}{442.387}{669.029}{the}%
\TextAt{8}{463.154}{669.029}{canonical}%
\TextAt{8}{515.148}{669.029}{iso-}%
\TextAt{8}{79.226}{651.983}{morphisms}%
\TextAt{13}{240.752}{634.938}{G}%
\TextAt{9}{249.986}{634.938}{\Character{40}}%
\TextAt{18}{254.538}{634.938}{Q}%
\TextAt{9}{263.837}{634.938}{\Character{40}}%
\TextAt{13}{268.389}{634.938}{\Character{22}}%
\TextAt{21}{275.432}{633.144}{n}%
\TextAt{9}{281.068}{634.938}{\Character{41}}%
\TextAt{16}{288.941}{634.938}{\Character{106}}%
\TextAt{18}{295.583}{634.938}{Q}%
\TextAt{9}{304.882}{634.938}{\Character{41}}%
\NativeMathAt{312.755}{634.937}{11.955}{\cong}%
\TextAt{9}{325.374}{634.938}{\Character{40}}%
\TextAt{18}{329.926}{634.938}{Z}%
\TextAt{9}{337.896}{634.938}{\Character{47}}%
\TextAt{13}{343.749}{634.938}{n}%
\TextAt{18}{350.737}{634.938}{Z}%
\TextAt{9}{358.707}{634.938}{\Character{41}}%
\TextAt{3}{363.26}{639.874}{\Character{3}}%
\TextAt{13}{367.992}{634.938}{\Character{59}}%
\TextAt{8}{79.226}{611.416}{where}%
\TextAt{13}{113.035}{611.416}{\Character{22}}%
\TextAt{21}{120.078}{609.623}{n}%
\TextAt{8}{129.612}{611.416}{denotes}%
\TextAt{8}{171.945}{611.416}{the}%
\TextAt{8}{192.102}{611.416}{group}%
\TextAt{8}{225.277}{611.416}{of}%
\TextAt{13}{238.607}{611.416}{n}%
\TextAt{8}{245.595}{611.416}{-th}%
\TextAt{8}{264.448}{611.416}{roots}%
\TextAt{8}{294.109}{611.416}{of}%
\TextAt{8}{307.439}{611.416}{unity.}%
\TextAt{12}{97.159}{591.381}{Example}%
\TextAt{12}{151.331}{591.381}{7:}%
\TextAt{8}{166.419}{591.381}{The}%
\TextAt{8}{189.312}{591.381}{groups}%
\TextAt{18}{225.941}{591.381}{Z}%
\TextAt{21}{233.911}{589.588}{p}%
\TextAt{8}{241.408}{591.381}{and}%
\TextAt{20}{263.673}{594.469}{\Character{98}}%
\TextAt{18}{263.009}{591.381}{Z}%
\TextAt{8}{273.715}{591.381}{are}%
\TextAt{8}{292.064}{591.381}{(additive)}%
\TextAt{8}{344.893}{591.381}{special}%
\TextAt{8}{381.832}{591.381}{cases}%
\TextAt{8}{410.057}{591.381}{of}%
\TextAt{8}{422.225}{591.381}{the}%
\TextAt{8}{441.22}{591.381}{class}%
\TextAt{8}{467.496}{591.381}{of}%
\TextAt{12}{480.179}{591.381}{procyclic}%
\TextAt{12}{79.226}{574.335}{groups}%
\TextAt{8}{118.439}{574.335}{.}%
\TextAt{8}{132.195}{574.335}{These}%
\TextAt{8}{167.834}{574.335}{are}%
\TextAt{8}{189.114}{574.335}{profinite}%
\TextAt{8}{237.712}{574.335}{groups}%
\TextAt{13}{277.271}{574.335}{G}%
\TextAt{8}{292.172}{574.335}{which}%
\TextAt{8}{327.428}{574.335}{are}%
\TextAt{8}{348.708}{574.335}{topologically}%
\TextAt{8}{419.435}{574.335}{generated}%
\TextAt{8}{474.536}{574.335}{by}%
\TextAt{8}{492.564}{574.335}{a}%
\TextAt{8}{504.089}{574.335}{single}%
\TextAt{8}{79.226}{557.289}{element}%
\TextAt{13}{123.702}{557.289}{\Character{27}}%
\TextAt{8}{130.784}{557.289}{;}%
\TextAt{8}{139.784}{557.289}{i.e.,}%
\TextAt{13}{163.431}{557.289}{G}%
\TextAt{8}{177.796}{557.289}{is}%
\TextAt{8}{190.794}{557.289}{the}%
\TextAt{8}{212.184}{557.289}{closure}%
\Vector{q 0 G 0 g 0.478 w 180.499 -206.071 m 196.88 -206.071 l S Q}%
\TextAt{16}{252.499}{557.289}{\Character{104}}%
\TextAt{13}{257.148}{557.289}{\Character{27}}%
\TextAt{16}{264.231}{557.289}{\Character{105}}%
\TextAt{8}{274.011}{557.289}{of}%
\TextAt{8}{288.575}{557.289}{the}%
\TextAt{8}{309.965}{557.289}{subgroup}%
\TextAt{16}{361.996}{557.289}{\Character{104}}%
\TextAt{13}{366.645}{557.289}{\Character{27}}%
\TextAt{16}{373.728}{557.289}{\Character{105}}%
\TextAt{9}{383.8}{557.289}{\Character{61}}%
\TextAt{16}{398.328}{557.289}{\Character{102}}%
\TextAt{13}{404.306}{557.289}{\Character{27}}%
\TextAt{21}{411.388}{561.628}{n}%
\TextAt{16}{422.448}{557.289}{\Character{106}}%
\TextAt{13}{431.192}{557.289}{n}%
\NativeMathAt{443.603}{557.289}{11.955}{\in}%
\TextAt{18}{456.996}{557.289}{Z}%
\TextAt{16}{464.966}{557.289}{\Character{103}}%
\TextAt{8}{470.944}{557.289}{.}%
\TextAt{8}{483.093}{557.289}{The}%
\TextAt{8}{508.381}{557.289}{open}%
\TextAt{8}{79.226}{540.243}{subgroups}%
\TextAt{8}{135.668}{540.243}{of}%
\TextAt{8}{150.029}{540.243}{a}%
\TextAt{8}{160.814}{540.243}{procyclic}%
\TextAt{8}{211.267}{540.243}{group}%
\TextAt{13}{245.473}{540.243}{G}%
\TextAt{9}{259.783}{540.243}{\Character{61}}%
\Vector{q 0 G 0 g 0.478 w 201.964 -223.117 m 218.345 -223.117 l S Q}%
\TextAt{16}{273.964}{540.243}{\Character{104}}%
\TextAt{13}{278.613}{540.243}{\Character{27}}%
\TextAt{16}{285.695}{540.243}{\Character{105}}%
\TextAt{8}{295.272}{540.243}{are}%
\TextAt{8}{315.813}{540.243}{all}%
\TextAt{8}{333.102}{540.243}{of}%
\TextAt{8}{347.462}{540.243}{the}%
\TextAt{8}{368.649}{540.243}{form}%
\TextAt{13}{397.319}{540.243}{G}%
\TextAt{21}{406.553}{544.582}{n}%
\TextAt{8}{412.189}{540.243}{.}%
\TextAt{8}{423.728}{540.243}{Indeed,}%
\TextAt{13}{466.297}{540.243}{G}%
\TextAt{21}{475.531}{544.582}{n}%
\TextAt{8}{486.095}{540.243}{is}%
\TextAt{8}{498.889}{540.243}{closed,}%
\TextAt{8}{79.226}{523.197}{being}%
\TextAt{8}{111.196}{523.197}{the}%
\TextAt{8}{131.786}{523.197}{image}%
\TextAt{8}{166.039}{523.197}{of}%
\TextAt{8}{179.802}{523.197}{the}%
\TextAt{8}{200.391}{523.197}{continuous}%
\TextAt{8}{259.428}{523.197}{map}%
\TextAt{13}{285.875}{523.197}{G}%
\NativeMathAt{299.166}{523.197}{11.955}{\rightarrow}%
\TextAt{13}{315.18}{523.197}{G}%
\TextAt{8}{324.413}{523.197}{,}%
\TextAt{13}{332.103}{523.197}{x}%
\NativeArrowAt{342.813}{523.197}{11.955}{11.955}{\mapsto}%
\TextAt{13}{358.827}{523.197}{x}%
\TextAt{21}{365.479}{527.536}{n}%
\TextAt{8}{371.115}{523.197}{,}%
\TextAt{8}{378.805}{523.197}{and}%
\TextAt{8}{402}{523.197}{the}%
\TextAt{8}{422.589}{523.197}{quotient}%
\TextAt{8}{469.205}{523.197}{group}%
\TextAt{13}{502.813}{523.197}{G}%
\TextAt{9}{512.047}{523.197}{\Character{47}}%
\TextAt{13}{517.9}{523.197}{G}%
\TextAt{21}{527.133}{527.536}{n}%
\TextAt{8}{79.226}{506.151}{is}%
\TextAt{8}{92.831}{506.151}{finite,}%
\TextAt{8}{128.295}{506.151}{because}%
\TextAt{8}{173.437}{506.151}{it}%
\TextAt{8}{186.982}{506.151}{contains}%
\TextAt{8}{234.742}{506.151}{the}%
\TextAt{8}{256.739}{506.151}{finite}%
\TextAt{8}{288.491}{506.151}{group}%
\TextAt{16}{323.507}{506.151}{\Character{102}}%
\TextAt{13}{329.485}{506.151}{\Character{27}}%
\TextAt{21}{336.567}{510.49}{\Character{23}}%
\TextAt{8}{345.07}{506.151}{mod}%
\TextAt{13}{370.508}{506.151}{G}%
\TextAt{21}{379.742}{510.49}{n}%
\TextAt{16}{391.835}{506.151}{\Character{106}}%
\TextAt{9}{401.614}{506.151}{0}%
\NativeMathAt{413.924}{506.151}{11.955}{\leq}%
\TextAt{13}{429.68}{506.151}{\Character{23}}%
\TextAt{13}{442.663}{506.151}{\Character{60}}%
\TextAt{13}{458.217}{506.151}{n}%
\TextAt{16}{465.225}{506.151}{\Character{103}}%
\TextAt{8}{476.94}{506.151}{as}%
\TextAt{8}{493.151}{506.151}{a}%
\TextAt{8}{504.747}{506.151}{dense}%
\TextAt{8}{79.226}{489.105}{subgroup,}%
\TextAt{8}{134.158}{489.105}{and}%
\TextAt{8}{157.626}{489.105}{is}%
\TextAt{8}{170.095}{489.105}{therefore}%
\TextAt{8}{219.901}{489.105}{equal}%
\TextAt{8}{251.499}{489.105}{to}%
\TextAt{8}{266.515}{489.105}{it.}%
\TextAt{8}{284.887}{489.105}{Conversely,}%
\TextAt{8}{347.29}{489.105}{if}%
\TextAt{13}{358.719}{489.105}{H}%
\TextAt{8}{373.945}{489.105}{is}%
\TextAt{8}{386.415}{489.105}{a}%
\TextAt{8}{396.876}{489.105}{subgroup}%
\TextAt{8}{448.379}{489.105}{of}%
\TextAt{13}{462.414}{489.105}{G}%
\TextAt{8}{476.251}{489.105}{of}%
\TextAt{8}{490.287}{489.105}{index}%
\TextAt{13}{522.53}{489.105}{n}%
\TextAt{8}{529.518}{489.105}{,}%
\TextAt{8}{79.226}{472.059}{then}%
\TextAt{13}{105.886}{472.059}{G}%
\TextAt{21}{115.12}{476.398}{n}%
\NativeMathAt{124.077}{472.059}{11.955}{\subseteq}%
\TextAt{13}{136.696}{472.059}{H}%
\NativeMathAt{150.64}{472.059}{11.955}{\subseteq}%
\TextAt{13}{163.259}{472.059}{G}%
\TextAt{8}{176.39}{472.059}{and}%
\TextAt{13}{199.153}{472.059}{n}%
\TextAt{9}{209.461}{472.059}{\Character{61}}%
\TextAt{9}{221.871}{472.059}{\Character{40}}%
\TextAt{13}{226.439}{472.059}{G}%
\TextAt{9}{238.993}{472.059}{\Character{58}}%
\TextAt{13}{245.566}{472.059}{H}%
\TextAt{9}{256.189}{472.059}{\Character{41}}%
\NativeMathAt{264.062}{472.059}{11.955}{\leq}%
\TextAt{9}{276.681}{472.059}{\Character{40}}%
\TextAt{13}{281.233}{472.059}{G}%
\TextAt{9}{293.788}{472.059}{\Character{58}}%
\TextAt{13}{300.36}{472.059}{G}%
\TextAt{21}{309.594}{476.398}{n}%
\TextAt{9}{315.23}{472.059}{\Character{41}}%
\NativeMathAt{323.103}{472.059}{11.955}{\leq}%
\TextAt{13}{335.723}{472.059}{n}%
\TextAt{8}{342.71}{472.059}{,}%
\TextAt{8}{349.86}{472.059}{so}%
\TextAt{8}{364.23}{472.059}{that}%
\TextAt{13}{389.599}{472.059}{H}%
\TextAt{9}{403.542}{472.059}{\Character{61}}%
\TextAt{13}{415.967}{472.059}{G}%
\TextAt{21}{425.201}{476.398}{n}%
\TextAt{8}{430.837}{472.059}{.}%
\TextAt{8}{97.159}{452.025}{Every}%
\TextAt{8}{131.509}{452.025}{procyclic}%
\TextAt{8}{181.629}{452.025}{group}%
\TextAt{13}{215.501}{452.025}{G}%
\TextAt{8}{229.329}{452.025}{is}%
\TextAt{8}{241.789}{452.025}{a}%
\TextAt{8}{252.241}{452.025}{quotient}%
\TextAt{8}{299.121}{452.025}{of}%
\TextAt{8}{313.147}{452.025}{the}%
\TextAt{8}{334}{452.025}{group}%
\TextAt{20}{368.537}{455.113}{\Character{98}}%
\TextAt{18}{367.873}{452.025}{Z}%
\TextAt{8}{375.843}{452.025}{.}%
\TextAt{8}{386.381}{452.025}{In}%
\TextAt{8}{401.699}{452.025}{fact,}%
\TextAt{8}{428.907}{452.025}{if}%
\TextAt{13}{440.328}{452.025}{G}%
\TextAt{9}{454.069}{452.025}{\Character{61}}%
\Vector{q 0 G 0 g 0.478 w 395.682 -311.335 m 412.063 -311.335 l S Q}%
\TextAt{16}{467.682}{452.025}{\Character{104}}%
\TextAt{13}{472.331}{452.025}{\Character{27}}%
\TextAt{16}{479.414}{452.025}{\Character{105}}%
\TextAt{8}{484.063}{452.025}{,}%
\TextAt{8}{492.083}{452.025}{then}%
\TextAt{8}{519.439}{452.025}{we}%
\TextAt{8}{79.226}{434.979}{have}%
\TextAt{8}{106.221}{434.979}{for}%
\TextAt{8}{124.106}{434.979}{every}%
\TextAt{13}{154.998}{434.979}{n}%
\TextAt{8}{165.883}{434.979}{the}%
\TextAt{8}{186.039}{434.979}{surjective}%
\TextAt{8}{238.451}{434.979}{homomorphism}%
\TextAt{18}{184.383}{405.977}{Z}%
\TextAt{9}{192.353}{405.977}{\Character{47}}%
\TextAt{13}{198.206}{405.977}{n}%
\TextAt{18}{205.194}{405.977}{Z}%
\NativeArrowAt{216.485}{405.977}{11.955}{19.26}{\longrightarrow}%
\TextAt{13}{239.067}{405.977}{G}%
\TextAt{9}{248.3}{405.977}{\Character{47}}%
\TextAt{13}{254.153}{405.977}{G}%
\TextAt{21}{263.387}{410.914}{n}%
\TextAt{13}{269.023}{405.977}{\Character{59}}%
\TextAt{9}{298.178}{405.977}{1}%
\TextAt{8}{307.351}{405.977}{mod}%
\TextAt{13}{332.789}{405.977}{n}%
\TextAt{18}{339.777}{405.977}{Z}%
\NativeArrowAt{351.068}{405.977}{11.955}{19.26}{\longmapsto}%
\TextAt{13}{373.65}{405.977}{\Character{27}}%
\TextAt{8}{384.053}{405.977}{mod}%
\TextAt{13}{409.491}{405.977}{G}%
\TextAt{21}{418.725}{410.914}{n}%
\TextAt{13}{424.361}{405.977}{\Character{59}}%
\TextAt{8}{79.226}{376.976}{and}%
\TextAt{8}{101.558}{376.976}{in}%
\TextAt{8}{114.779}{376.976}{view}%
\TextAt{8}{141.331}{376.976}{of}%
\TextAt{8}{154.23}{376.976}{(2.7),}%
\TextAt{8}{185.112}{376.976}{passing}%
\TextAt{8}{225.783}{376.976}{to}%
\TextAt{8}{239.662}{376.976}{the}%
\TextAt{8}{259.388}{376.976}{projective}%
\TextAt{8}{312.934}{376.976}{limit}%
\TextAt{8}{340.466}{376.976}{yields}%
\TextAt{8}{372.936}{376.976}{a}%
\TextAt{8}{382.26}{376.976}{continuous}%
\TextAt{8}{440.434}{376.976}{surjection}%
\TextAt{20}{494.381}{380.065}{\Character{98}}%
\TextAt{18}{493.717}{376.976}{Z}%
\NativeMathAt{505.008}{376.976}{11.955}{\rightarrow}%
\TextAt{13}{520.284}{376.976}{G}%
\TextAt{8}{529.518}{376.976}{.}%
\TextAt{12}{97.159}{356.941}{Example}%
\TextAt{12}{152.667}{356.941}{8:}%
\TextAt{8}{168.336}{356.941}{Let}%
\TextAt{13}{189.305}{356.941}{A}%
\TextAt{8}{201.978}{356.941}{be}%
\TextAt{8}{217.902}{356.941}{an}%
\TextAt{8}{234.161}{356.941}{abelian}%
\TextAt{8}{274.809}{356.941}{torsion}%
\TextAt{8}{313.902}{356.941}{group.}%
\TextAt{8}{351.629}{356.941}{Then}%
\TextAt{8}{382.186}{356.941}{the}%
\TextAt{12}{402.343}{356.941}{Pontryagin}%
\TextAt{12}{471.432}{356.941}{dual}%
\TextAt{13}{250.201}{327.94}{\Character{31}}%
\TextAt{9}{257.534}{327.94}{\Character{40}}%
\TextAt{13}{262.087}{327.94}{A}%
\TextAt{9}{270.862}{327.94}{\Character{41}}%
\TextAt{9}{278.717}{327.94}{\Character{61}}%
\TextAt{8}{291.161}{327.94}{Hom}%
\TextAt{9}{315.549}{327.94}{\Character{40}}%
\TextAt{13}{320.101}{327.94}{A\Character{59}}%
\TextAt{18}{334.121}{327.94}{Q}%
\TextAt{9}{343.419}{327.94}{\Character{47}}%
\TextAt{18}{349.272}{327.94}{Z}%
\TextAt{9}{357.243}{327.94}{\Character{41}}%
\TextAt{8}{79.226}{298.939}{is}%
\TextAt{8}{90.99}{298.939}{a}%
\TextAt{8}{100.745}{298.939}{profinite}%
\TextAt{8}{147.574}{298.939}{group.}%
\TextAt{8}{185.3}{298.939}{For}%
\TextAt{8}{206.27}{298.939}{one}%
\TextAt{8}{227.729}{298.939}{has}%
\TextAt{13}{278.395}{281.893}{A}%
\TextAt{9}{290.491}{281.893}{\Character{61}}%
\TextAt{20}{302.916}{293.251}{\Character{91}}%
\TextAt{21}{308.117}{268.056}{i}%
\TextAt{13}{318.193}{281.893}{A}%
\TextAt{21}{326.968}{280.1}{i}%
\TextAt{13}{330.349}{281.893}{\Character{59}}%
\TextAt{8}{79.226}{250.972}{where}%
\TextAt{13}{113.035}{250.972}{A}%
\TextAt{21}{121.811}{249.179}{i}%
\TextAt{8}{129.089}{250.972}{varies}%
\TextAt{8}{162.002}{250.972}{over}%
\TextAt{8}{187.048}{250.972}{the}%
\TextAt{8}{207.204}{250.972}{finite}%
\TextAt{8}{237.116}{250.972}{subgroups}%
\TextAt{8}{292.529}{250.972}{of}%
\TextAt{13}{305.859}{250.972}{A}%
\TextAt{8}{314.634}{250.972}{,}%
\TextAt{8}{321.783}{250.972}{and}%
\TextAt{8}{344.546}{250.972}{thus}%
\TextAt{13}{261.998}{221.971}{\Character{31}}%
\TextAt{9}{269.331}{221.971}{\Character{40}}%
\TextAt{13}{273.884}{221.971}{A}%
\TextAt{9}{282.659}{221.971}{\Character{41}}%
\NativeMathAt{290.532}{221.971}{11.955}{\cong}%
\NativeMathAt{303.151}{221.971}{11.955}{\varprojlim}%
\TextAt{21}{309.839}{209.768}{i}%
\TextAt{13}{321.403}{221.971}{\Character{31}}%
\TextAt{9}{328.737}{221.971}{\Character{40}}%
\TextAt{13}{333.289}{221.971}{A}%
\TextAt{21}{342.064}{220.178}{i}%
\TextAt{9}{345.446}{221.971}{\Character{41}}%
\TextAt{8}{79.226}{186.655}{with}%
\TextAt{8}{105.886}{186.655}{finite}%
\TextAt{8}{135.798}{186.655}{groups}%
\TextAt{13}{173.588}{186.655}{\Character{31}}%
\TextAt{9}{180.922}{186.655}{\Character{40}}%
\TextAt{13}{185.474}{186.655}{A}%
\TextAt{21}{194.25}{184.862}{i}%
\TextAt{9}{197.631}{186.655}{\Character{41}}%
\TextAt{8}{202.183}{186.655}{.}%
\TextAt{8}{210.631}{186.655}{If}%
\TextAt{8}{222.324}{186.655}{for}%
\TextAt{8}{240.209}{186.655}{instance,}%
\TextAt{13}{248.75}{157.654}{A}%
\TextAt{9}{260.846}{157.654}{\Character{61}}%
\TextAt{18}{273.271}{157.654}{Q}%
\TextAt{9}{282.57}{157.654}{\Character{47}}%
\TextAt{18}{288.423}{157.654}{Z}%
\TextAt{9}{299.714}{157.654}{\Character{61}}%
\TextAt{20}{313.768}{169.011}{\Character{91}}%
\TextAt{21}{312.139}{143.596}{n}%
\NativeMathAt{317.278}{143.596}{7.97}{\in}%
\TextAt{25}{322.923}{143.596}{N}%
\TextAt{17}{332.319}{162.361}{1}%
\Vector{q 0 G 0 g 0.478 w 259.867 -613.357 m 265.006 -613.357 l S Q}%
\TextAt{21}{331.867}{153.531}{n}%
\TextAt{18}{338.201}{157.654}{Z}%
\TextAt{9}{346.171}{157.654}{\Character{47}}%
\TextAt{18}{352.024}{157.654}{Z}%
\TextAt{13}{359.994}{157.654}{\Character{59}}%
\TextAt{8}{79.226}{119.102}{then}%
\TextAt{13}{105.886}{119.102}{\Character{31}}%
\TextAt{9}{113.22}{119.102}{\Character{40}}%
\TextAt{17}{119.42}{123.809}{1}%
\Vector{q 0 G 0 g 0.478 w 46.968 -651.909 m 52.106 -651.909 l S Q}%
\TextAt{21}{118.968}{114.979}{n}%
\TextAt{18}{125.301}{119.102}{Z}%
\TextAt{9}{133.271}{119.102}{\Character{47}}%
\TextAt{18}{139.124}{119.102}{Z}%
\TextAt{9}{147.095}{119.102}{\Character{41}}%
\TextAt{9}{154.95}{119.102}{\Character{61}}%
\TextAt{18}{167.393}{119.102}{Z}%
\TextAt{9}{175.363}{119.102}{\Character{47}}%
\TextAt{13}{181.216}{119.102}{n}%
\TextAt{18}{188.204}{119.102}{Z}%
\TextAt{8}{196.174}{119.102}{,}%
\TextAt{8}{203.323}{119.102}{so}%
\TextAt{8}{217.693}{119.102}{that}%
\TextAt{13}{241.248}{90.101}{\Character{31}}%
\TextAt{9}{248.581}{90.101}{\Character{40}}%
\TextAt{18}{253.133}{90.101}{Q}%
\TextAt{9}{262.432}{90.101}{\Character{47}}%
\TextAt{18}{268.285}{90.101}{Z}%
\TextAt{9}{276.255}{90.101}{\Character{41}}%
\NativeMathAt{284.128}{90.101}{11.955}{\cong}%
\NativeMathAt{296.748}{90.101}{11.955}{\varprojlim}%
\TextAt{21}{302.308}{77.986}{n}%
\TextAt{18}{314.999}{90.101}{Z}%
\TextAt{9}{322.969}{90.101}{\Character{47}}%
\TextAt{13}{328.822}{90.101}{n}%
\TextAt{18}{335.81}{90.101}{Z}%
\TextAt{9}{347.101}{90.101}{\Character{61}}%
\TextAt{20}{360.19}{93.189}{\Character{98}}%
\TextAt{18}{359.526}{90.101}{Z}%
\TextAt{13}{367.496}{90.101}{\Character{58}}%
\TextAt{8}{297.211}{43.033}{250}%
\Destination{78.23}{809.815}{nav.page.250}%
\Destination{79.36}{776.155}{nav.print.273}%
\Hotlink{153.73}{374.755}{25.078}{11.047}{nav.number.4.2.7}%
}

% Page 263 (body page 251)
\Page{%
\TextAt{2}{79.226}{796.943}{IV.}%
\TextAt{2}{98.012}{796.943}{Abstract}%
\TextAt{2}{143.252}{796.943}{Class}%
\TextAt{2}{171.844}{796.943}{Field}%
\TextAt{2}{199.575}{796.943}{Theory}%
\Vector{q 0 G 0 g 0.299 w 7.226 18.729 m 460.769 18.729 l S Q}%
\TextAt{12}{97.159}{761.124}{Example}%
\TextAt{12}{152.066}{761.124}{9:}%
\TextAt{8}{167.474}{761.124}{If}%
\TextAt{13}{178.644}{761.124}{G}%
\TextAt{8}{191.253}{761.124}{is}%
\TextAt{8}{202.495}{761.124}{any}%
\TextAt{8}{224.09}{761.124}{group}%
\TextAt{8}{256.743}{761.124}{and}%
\TextAt{13}{278.984}{761.124}{N}%
\TextAt{8}{292.981}{761.124}{varies}%
\TextAt{8}{325.372}{761.124}{over}%
\TextAt{8}{349.896}{761.124}{all}%
\TextAt{8}{365.633}{761.124}{normal}%
\TextAt{8}{404.79}{761.124}{subgroups}%
\TextAt{8}{459.68}{761.124}{of}%
\TextAt{8}{472.488}{761.124}{finite}%
\TextAt{8}{501.877}{761.124}{index,}%
\TextAt{8}{79.226}{744.078}{then}%
\TextAt{8}{105.886}{744.078}{the}%
\TextAt{8}{126.043}{744.078}{profinite}%
\TextAt{8}{172.871}{744.078}{group}%
\TextAt{20}{273.799}{730.054}{\Character{98}}%
\TextAt{13}{271.527}{727.032}{G}%
\TextAt{9}{284.082}{727.032}{\Character{61}}%
\NativeMathAt{296.507}{727.032}{11.955}{\varprojlim}%
\TextAt{21}{300.852}{714.652}{N}%
\TextAt{13}{314.759}{727.032}{G}%
\TextAt{9}{323.993}{727.032}{\Character{47}}%
\TextAt{13}{329.846}{727.032}{N}%
\TextAt{8}{79.226}{697.568}{is}%
\TextAt{8}{91.041}{697.568}{called}%
\TextAt{8}{124.255}{697.568}{the}%
\TextAt{12}{144.462}{697.568}{profinite}%
\TextAt{12}{198.76}{697.568}{completion}%
\TextAt{8}{267.158}{697.568}{of}%
\TextAt{13}{280.539}{697.568}{G}%
\TextAt{8}{289.773}{697.568}{.}%
\TextAt{8}{298.372}{697.568}{The}%
\TextAt{8}{322.477}{697.568}{profinite}%
\TextAt{8}{369.356}{697.568}{completion}%
\TextAt{8}{429.242}{697.568}{of}%
\TextAt{18}{442.622}{697.568}{Z}%
\TextAt{8}{450.592}{697.568}{,}%
\TextAt{8}{457.805}{697.568}{for}%
\TextAt{8}{475.74}{697.568}{example,}%
\TextAt{8}{524.903}{697.568}{is}%
\TextAt{8}{79.226}{680.522}{the}%
\TextAt{8}{99.383}{680.522}{group}%
\TextAt{20}{133.222}{683.61}{\Character{98}}%
\TextAt{18}{132.558}{680.522}{Z}%
\TextAt{9}{143.849}{680.522}{\Character{61}}%
\NativeMathAt{156.275}{680.522}{11.955}{\varprojlim}%
\TextAt{21}{172.534}{675.082}{n}%
\TextAt{18}{180.162}{680.522}{Z}%
\TextAt{9}{188.133}{680.522}{\Character{47}}%
\TextAt{13}{193.986}{680.522}{n}%
\TextAt{18}{200.973}{680.522}{Z}%
\TextAt{8}{208.943}{680.522}{.}%
\TextAt{12}{79.226}{654.509}{Exercise}%
\TextAt{12}{134.397}{654.509}{IV.2.1.}%
\TextAt{8}{179.865}{654.509}{Show}%
\TextAt{8}{212.469}{654.509}{that,}%
\TextAt{8}{243.231}{654.509}{for}%
\TextAt{8}{262.828}{654.509}{a}%
\TextAt{8}{274.296}{654.509}{profinite}%
\TextAt{8}{322.837}{654.509}{group}%
\TextAt{13}{357.725}{654.509}{G}%
\TextAt{8}{366.959}{654.509}{,}%
\TextAt{8}{376.248}{654.509}{the}%
\TextAt{8}{398.117}{654.509}{power}%
\TextAt{8}{433.974}{654.509}{map}%
\TextAt{13}{461.7}{654.509}{G}%
\TextAt{16}{474.758}{654.509}{\Character{2}}%
\TextAt{18}{487.881}{654.509}{Z}%
\NativeMathAt{502.09}{654.509}{11.955}{\rightarrow}%
\TextAt{13}{520.284}{654.509}{G}%
\TextAt{8}{529.518}{654.509}{,}%
\TextAt{9}{79.226}{637.463}{\Character{40}}%
\TextAt{13}{83.778}{637.463}{\Character{27}}%
\TextAt{13}{90.21}{637.463}{\Character{59}}%
\TextAt{13}{95.434}{637.463}{n}%
\TextAt{9}{102.442}{637.463}{\Character{41}}%
\NativeArrowAt{110.315}{637.463}{11.955}{11.955}{\mapsto}%
\TextAt{13}{125.591}{637.463}{\Character{27}}%
\TextAt{21}{132.674}{641.802}{n}%
\TextAt{8}{138.31}{637.463}{,}%
\TextAt{8}{145.459}{637.463}{extends}%
\TextAt{8}{188.115}{637.463}{to}%
\TextAt{8}{202.426}{637.463}{a}%
\TextAt{8}{212.181}{637.463}{continuous}%
\TextAt{8}{270.785}{637.463}{map}%
\TextAt{13}{223.865}{608.462}{G}%
\TextAt{16}{235.755}{608.462}{\Character{2}}%
\TextAt{20}{248.375}{611.551}{\Character{98}}%
\TextAt{18}{247.711}{608.462}{Z}%
\NativeArrowAt{259.002}{608.462}{11.955}{19.26}{\longrightarrow}%
\TextAt{13}{281.584}{608.462}{G\Character{59}}%
\TextAt{9}{319.972}{608.462}{\Character{40}}%
\TextAt{13}{324.524}{608.462}{\Character{27}}%
\TextAt{13}{330.956}{608.462}{\Character{59}}%
\TextAt{13}{336.18}{608.462}{a}%
\TextAt{9}{342.345}{608.462}{\Character{41}}%
\NativeArrowAt{350.218}{608.462}{11.955}{19.26}{\longmapsto}%
\TextAt{13}{372.8}{608.462}{\Character{27}}%
\TextAt{21}{379.883}{613.398}{a}%
\TextAt{13}{384.879}{608.462}{\Character{59}}%
\TextAt{8}{79.226}{579.461}{and}%
\TextAt{8}{101.989}{579.461}{that}%
\TextAt{8}{127.358}{579.461}{one}%
\TextAt{8}{148.817}{579.461}{has}%
\TextAt{9}{169.691}{579.461}{\Character{40}}%
\TextAt{13}{174.243}{579.461}{\Character{27}}%
\TextAt{21}{181.326}{583.8}{a}%
\TextAt{9}{186.322}{579.461}{\Character{41}}%
\TextAt{21}{190.874}{583.8}{b}%
\TextAt{9}{198.316}{579.461}{\Character{61}}%
\TextAt{13}{210.741}{579.461}{\Character{27}}%
\TextAt{21}{217.824}{583.8}{ab}%
\TextAt{8}{230.34}{579.461}{and}%
\TextAt{13}{253.102}{579.461}{\Character{27}}%
\TextAt{21}{260.185}{583.8}{a}%
\TextAt{17}{264.683}{583.8}{\Character{43}}%
\TextAt{21}{271.269}{583.8}{b}%
\TextAt{9}{278.711}{579.461}{\Character{61}}%
\TextAt{13}{291.136}{579.461}{\Character{27}}%
\TextAt{21}{298.219}{583.8}{a}%
\TextAt{13}{303.215}{579.461}{\Character{27}}%
\TextAt{21}{310.297}{583.8}{b}%
\TextAt{8}{318.315}{579.461}{if}%
\TextAt{13}{329.039}{579.461}{G}%
\TextAt{8}{342.17}{579.461}{is}%
\TextAt{8}{353.934}{579.461}{abelian.}%
\TextAt{12}{79.226}{547.471}{Exercise}%
\TextAt{12}{132.551}{547.471}{IV.2.2.}%
\TextAt{8}{176.416}{547.471}{If}%
\TextAt{13}{188.216}{547.471}{\Character{27}}%
\NativeMathAt{198.804}{547.471}{11.955}{\in}%
\TextAt{13}{210.28}{547.471}{G}%
\TextAt{8}{223.52}{547.471}{and}%
\TextAt{13}{246.391}{547.471}{a}%
\TextAt{9}{256.041}{547.471}{\Character{61}}%
\TextAt{8}{268.651}{547.471}{lim}%
\TextAt{21}{284.911}{545.678}{i}%
\NativeMathAt{287.794}{545.678}{7.97}{\rightarrow}%
\NativeMathAt{296.258}{545.678}{7.97}{\infty}%
\TextAt{13}{307.221}{547.471}{a}%
\TextAt{21}{313.366}{545.678}{i}%
\NativeMathAt{320.253}{547.471}{11.955}{\in}%
\TextAt{20}{332.393}{550.56}{\Character{98}}%
\TextAt{18}{331.729}{547.471}{Z}%
\TextAt{8}{343.705}{547.471}{with}%
\TextAt{13}{370.473}{547.471}{a}%
\TextAt{21}{376.618}{545.678}{i}%
\NativeMathAt{383.505}{547.471}{11.955}{\in}%
\TextAt{18}{394.981}{547.471}{Z}%
\TextAt{8}{402.951}{547.471}{,}%
\TextAt{8}{410.236}{547.471}{then}%
\TextAt{13}{437.004}{547.471}{\Character{27}}%
\TextAt{21}{444.086}{551.81}{a}%
\TextAt{9}{452.588}{547.471}{\Character{61}}%
\TextAt{8}{465.198}{547.471}{lim}%
\TextAt{21}{481.457}{545.678}{i}%
\NativeMathAt{484.341}{545.678}{7.97}{\rightarrow}%
\NativeMathAt{492.805}{545.678}{7.97}{\infty}%
\TextAt{13}{503.768}{547.471}{\Character{27}}%
\TextAt{21}{510.85}{551.81}{a}%
\TextAt{22}{515.348}{550.595}{i}%
\TextAt{8}{523.014}{547.471}{in}%
\TextAt{13}{79.226}{530.425}{G}%
\TextAt{8}{88.46}{530.425}{.}%
\TextAt{12}{79.226}{498.435}{Exercise}%
\TextAt{12}{132.49}{498.435}{IV.2.3.}%
\TextAt{8}{176.3}{498.435}{A}%
\TextAt{12}{189.027}{498.435}{pro-}%
\TextAt{13}{213.213}{498.435}{p}%
\TextAt{12}{219.088}{498.435}{-group}%
\TextAt{8}{261.428}{498.435}{is}%
\TextAt{8}{273.247}{498.435}{a}%
\TextAt{8}{283.057}{498.435}{profinite}%
\TextAt{8}{329.94}{498.435}{group}%
\TextAt{13}{363.17}{498.435}{G}%
\TextAt{8}{376.356}{498.435}{whose}%
\TextAt{8}{410.937}{498.435}{quotients}%
\TextAt{13}{461.789}{498.435}{G}%
\TextAt{9}{471.023}{498.435}{\Character{47}}%
\TextAt{13}{476.876}{498.435}{N}%
\TextAt{8}{487.499}{498.435}{,}%
\TextAt{8}{494.716}{498.435}{modulo}%
\TextAt{8}{79.226}{481.389}{all}%
\TextAt{8}{95.206}{481.389}{open}%
\TextAt{8}{123.213}{481.389}{normal}%
\TextAt{8}{162.614}{481.389}{subgroups}%
\TextAt{13}{217.747}{481.389}{N}%
\TextAt{8}{228.37}{481.389}{,}%
\TextAt{8}{235.296}{481.389}{are}%
\TextAt{8}{254.528}{481.389}{finite}%
\TextAt{13}{284.161}{481.389}{p}%
\TextAt{8}{290.036}{481.389}{-groups.}%
\TextAt{8}{336.182}{481.389}{Imitating}%
\TextAt{8}{387.609}{481.389}{exercise}%
\TextAt{8}{430.632}{481.389}{1,}%
\TextAt{8}{443.416}{481.389}{make}%
\TextAt{8}{473.707}{481.389}{sense}%
\TextAt{8}{503.459}{481.389}{of}%
\TextAt{8}{516.51}{481.389}{the}%
\TextAt{8}{79.226}{464.343}{powers}%
\TextAt{13}{117.985}{464.343}{\Character{27}}%
\TextAt{21}{125.067}{468.682}{a}%
\TextAt{8}{130.063}{464.343}{,}%
\TextAt{8}{137.213}{464.343}{for}%
\TextAt{8}{155.097}{464.343}{all}%
\TextAt{13}{171.357}{464.343}{\Character{27}}%
\NativeMathAt{181.76}{464.343}{11.955}{\in}%
\TextAt{13}{193.051}{464.343}{G}%
\TextAt{8}{206.182}{464.343}{and}%
\TextAt{13}{228.944}{464.343}{a}%
\NativeMathAt{238.41}{464.343}{11.955}{\in}%
\TextAt{18}{249.701}{464.343}{Z}%
\TextAt{21}{257.671}{462.55}{p}%
\TextAt{8}{262.432}{464.343}{.}%
\TextAt{12}{79.226}{432.353}{Exercise}%
\TextAt{12}{132.366}{432.353}{IV.2.4.}%
\TextAt{8}{176.069}{432.353}{A}%
\TextAt{8}{188.689}{432.353}{closed}%
\TextAt{8}{223.163}{432.353}{subgroup}%
\TextAt{13}{273.908}{432.353}{H}%
\TextAt{8}{288.376}{432.353}{of}%
\TextAt{8}{301.653}{432.353}{a}%
\TextAt{8}{311.356}{432.353}{profinite}%
\TextAt{8}{358.132}{432.353}{group}%
\TextAt{13}{391.255}{432.353}{G}%
\TextAt{8}{404.333}{432.353}{is}%
\TextAt{8}{416.044}{432.353}{called}%
\TextAt{8}{449.155}{432.353}{a}%
\TextAt{13}{458.858}{432.353}{p}%
\TextAt{12}{464.733}{432.353}{-Sylow}%
\TextAt{12}{508.034}{432.353}{sub-}%
\TextAt{12}{79.226}{415.307}{group}%
\TextAt{8}{117.449}{415.307}{of}%
\TextAt{13}{131.2}{415.307}{G}%
\TextAt{8}{144.752}{415.307}{if,}%
\TextAt{8}{159.254}{415.307}{for}%
\TextAt{8}{177.56}{415.307}{every}%
\TextAt{8}{208.873}{415.307}{open}%
\TextAt{8}{237.58}{415.307}{normal}%
\TextAt{8}{277.68}{415.307}{subgroup}%
\TextAt{13}{328.898}{415.307}{N}%
\TextAt{8}{343.839}{415.307}{of}%
\TextAt{13}{357.59}{415.307}{G}%
\TextAt{8}{366.824}{415.307}{,}%
\TextAt{8}{374.499}{415.307}{the}%
\TextAt{8}{395.076}{415.307}{group}%
\TextAt{13}{428.673}{415.307}{H}%
\TextAt{13}{439.289}{415.307}{N}%
\TextAt{9}{449.918}{415.307}{\Character{47}}%
\TextAt{13}{455.771}{415.307}{N}%
\TextAt{8}{470.712}{415.307}{is}%
\TextAt{8}{482.897}{415.307}{a}%
\TextAt{13}{493.073}{415.307}{p}%
\TextAt{8}{498.948}{415.307}{-Sylow}%
\TextAt{8}{79.226}{398.262}{subgroup}%
\TextAt{8}{130.024}{398.262}{of}%
\TextAt{13}{143.354}{398.262}{G}%
\TextAt{9}{152.587}{398.261}{\Character{47}}%
\TextAt{13}{158.44}{398.262}{N}%
\TextAt{8}{169.063}{398.262}{.}%
\TextAt{8}{177.511}{398.262}{Show:}%
\TextAt{8}{89.579}{378.227}{(i)}%
\TextAt{8}{107.918}{378.227}{For}%
\TextAt{8}{128.888}{378.227}{every}%
\TextAt{8}{159.78}{378.227}{prime}%
\TextAt{8}{192.944}{378.227}{number}%
\TextAt{13}{235.54}{378.227}{p}%
\TextAt{8}{241.415}{378.227}{,}%
\TextAt{8}{248.564}{378.227}{there}%
\TextAt{8}{278.476}{378.227}{exists}%
\TextAt{8}{310.791}{378.227}{a}%
\TextAt{13}{320.546}{378.227}{p}%
\TextAt{8}{326.421}{378.227}{-Sylow}%
\TextAt{8}{364.14}{378.227}{subgroup}%
\TextAt{8}{414.938}{378.227}{of}%
\TextAt{13}{428.268}{378.227}{G}%
\TextAt{8}{437.501}{378.227}{.}%
\TextAt{8}{86.327}{361.181}{(ii)}%
\TextAt{8}{107.918}{361.181}{Every}%
\TextAt{8}{141.572}{361.181}{pro-}%
\TextAt{13}{162.386}{361.181}{p}%
\TextAt{8}{168.261}{361.181}{-subgroup}%
\TextAt{8}{222.956}{361.181}{of}%
\TextAt{13}{236.286}{361.181}{G}%
\TextAt{8}{249.417}{361.181}{is}%
\TextAt{8}{261.181}{361.181}{contained}%
\TextAt{8}{314.19}{361.181}{in}%
\TextAt{8}{327.843}{361.181}{a}%
\TextAt{13}{337.599}{361.181}{p}%
\TextAt{8}{343.474}{361.181}{-Sylow}%
\TextAt{8}{381.192}{361.181}{subgroup.}%
\TextAt{8}{83.076}{344.135}{(iii)}%
\TextAt{8}{107.918}{344.135}{Every}%
\TextAt{8}{141.572}{344.135}{two}%
\TextAt{13}{163.689}{344.135}{p}%
\TextAt{8}{169.564}{344.135}{-Sylow}%
\TextAt{8}{207.283}{344.135}{subgroups}%
\TextAt{8}{262.695}{344.135}{of}%
\TextAt{13}{276.025}{344.135}{G}%
\TextAt{8}{289.156}{344.135}{are}%
\TextAt{8}{308.667}{344.135}{conjugate.}%
\TextAt{12}{79.226}{312.145}{Exercise}%
\TextAt{12}{132.427}{312.145}{IV.2.5.}%
\TextAt{8}{176.183}{312.145}{What}%
\TextAt{8}{209.023}{312.145}{is}%
\TextAt{8}{220.787}{312.145}{the}%
\TextAt{13}{240.944}{312.145}{p}%
\TextAt{8}{246.819}{312.145}{-Sylow}%
\TextAt{8}{284.537}{312.145}{subgroup}%
\TextAt{8}{335.335}{312.145}{of}%
\TextAt{20}{349.329}{315.233}{\Character{98}}%
\TextAt{18}{348.665}{312.145}{Z}%
\TextAt{8}{360.532}{312.145}{and}%
\TextAt{8}{383.295}{312.145}{of}%
\TextAt{18}{396.625}{312.145}{Z}%
\TextAt{3}{404.595}{316.483}{\Character{3}}%
\TextAt{21}{404.595}{309.189}{p}%
\TextAt{8}{409.356}{312.145}{?}%
\TextAt{12}{79.226}{280.155}{Exercise}%
\TextAt{12}{131.739}{280.155}{IV.2.6.}%
\TextAt{8}{174.898}{280.155}{If}%
\TextAt{16}{185.992}{280.155}{\Character{102}}%
\TextAt{13}{191.97}{280.155}{G}%
\TextAt{21}{201.203}{278.362}{i}%
\TextAt{16}{204.585}{280.155}{\Character{103}}%
\TextAt{8}{213.862}{280.155}{is}%
\TextAt{8}{225.028}{280.155}{a}%
\TextAt{8}{234.186}{280.155}{projective}%
\TextAt{8}{287.566}{280.155}{system}%
\TextAt{8}{325.787}{280.155}{of}%
\TextAt{8}{338.519}{280.155}{profinite}%
\TextAt{8}{384.75}{280.155}{groups}%
\TextAt{8}{421.943}{280.155}{and}%
\TextAt{13}{444.108}{280.155}{G}%
\TextAt{9}{456.663}{280.155}{\Character{61}}%
\NativeMathAt{469.088}{280.155}{11.955}{\varprojlim}%
\TextAt{21}{485.347}{274.715}{i}%
\TextAt{13}{490.721}{280.155}{G}%
\TextAt{21}{499.954}{278.362}{i}%
\TextAt{8}{503.336}{280.155}{,}%
\TextAt{8}{510.007}{280.155}{then}%
\TextAt{13}{79.226}{263.109}{G}%
\TextAt{7}{88.46}{267.447}{ab}%
\TextAt{9}{101.213}{263.109}{\Character{61}}%
\NativeMathAt{113.639}{263.109}{11.955}{\varprojlim}%
\TextAt{21}{129.898}{257.67}{i}%
\TextAt{13}{135.272}{263.109}{G}%
\TextAt{7}{144.505}{267.447}{ab}%
\TextAt{21}{144.505}{260.153}{i}%
\TextAt{8}{157.835}{263.109}{(see}%
\TextAt{8}{181.303}{263.109}{§1,}%
\TextAt{8}{199.989}{263.109}{exercise}%
\TextAt{8}{243.291}{263.109}{5).}%
\TextAt{15}{79.226}{210.245}{§}%
\TextAt{15}{93.785}{210.245}{3}%
\TextAt{15}{120.693}{210.245}{Abstract}%
\TextAt{15}{201.21}{210.245}{Galois}%
\TextAt{15}{260.414}{210.245}{Theory}%
\TextAt{8}{79.226}{178.359}{Class}%
\TextAt{8}{110.012}{178.359}{field}%
\TextAt{8}{135.465}{178.359}{theory}%
\TextAt{8}{172.312}{178.359}{is}%
\TextAt{8}{184.172}{178.359}{the}%
\TextAt{8}{204.425}{178.359}{final}%
\TextAt{8}{230.536}{178.359}{outcome}%
\TextAt{8}{277.461}{178.359}{of}%
\TextAt{8}{290.887}{178.359}{a}%
\TextAt{8}{300.739}{178.359}{long}%
\TextAt{8}{326.205}{178.359}{development}%
\TextAt{8}{394.266}{178.359}{of}%
\TextAt{8}{407.693}{178.359}{algebraic}%
\TextAt{8}{457.224}{178.359}{number}%
\TextAt{8}{499.917}{178.359}{theory}%
\TextAt{8}{79.226}{161.313}{the}%
\TextAt{8}{99.383}{161.313}{beginning}%
\TextAt{8}{153.037}{161.313}{of}%
\TextAt{8}{166.367}{161.313}{which}%
\TextAt{8}{199.854}{161.313}{was}%
\TextAt{8}{222.353}{161.313}{Gauss’s}%
\TextAt{8}{264.89}{161.313}{reciprocity}%
\TextAt{8}{323.1}{161.313}{law}%
\TextAt{20}{243.195}{141.846}{\Character{16}}%
\TextAt{13}{251.531}{136.664}{a}%
\Vector{q 0 G 0 g 0.478 w 179.531 -642.435 m 185.675 -642.435 l S Q}%
\TextAt{13}{252.114}{120.375}{b}%
\TextAt{20}{258.871}{141.846}{\Character{17}}%
\TextAt{20}{268.003}{145.433}{\Character{18}}%
\TextAt{13}{278.583}{136.664}{b}%
\Vector{q 0 G 0 g 0.478 w 205.999 -642.435 m 212.144 -642.435 l S Q}%
\TextAt{13}{277.999}{120.375}{a}%
\TextAt{20}{285.34}{145.433}{\Character{19}}%
\TextAt{9}{297.461}{128.576}{\Character{61}}%
\TextAt{9}{309.87}{128.576}{\Character{40}}%
\TextAt{16}{314.439}{128.576}{\Character{0}}%
\TextAt{9}{323.737}{128.576}{1\Character{41}}%
\TextAt{22}{335.338}{136.873}{a}%
\TextAt{23}{339.372}{136.873}{\Character{0}}%
\TextAt{24}{345.128}{136.873}{1}%
\Vector{q 0 G 0 g 0.359 w 263.338 -638.495 m 276.781 -638.495 l S Q}%
\TextAt{24}{340.233}{130.387}{2}%
\TextAt{22}{351.172}{136.873}{b}%
\TextAt{23}{354.446}{136.873}{\Character{0}}%
\TextAt{24}{360.202}{136.873}{1}%
\Vector{q 0 G 0 g 0.359 w 279.172 -638.495 m 291.855 -638.495 l S Q}%
\TextAt{24}{355.687}{130.387}{2}%
\TextAt{13}{365.549}{128.576}{\Character{58}}%
\TextAt{8}{79.226}{95.851}{The}%
\TextAt{8}{103.079}{95.851}{endeavours}%
\TextAt{8}{163.109}{95.851}{to}%
\TextAt{8}{177.219}{95.851}{generalize}%
\TextAt{8}{230.996}{95.851}{this}%
\TextAt{8}{253.618}{95.851}{law}%
\TextAt{8}{274.554}{95.851}{finally}%
\TextAt{8}{309.801}{95.851}{produced}%
\TextAt{8}{360.649}{95.851}{a}%
\TextAt{8}{370.204}{95.851}{theory}%
\TextAt{8}{406.754}{95.851}{of}%
\TextAt{8}{419.883}{95.851}{the}%
\TextAt{8}{439.839}{95.851}{abelian}%
\TextAt{8}{480.286}{95.851}{extensions}%
\TextAt{8}{79.226}{78.805}{of}%
\TextAt{8}{93.345}{78.805}{algebraic}%
\TextAt{8}{143.568}{78.805}{and}%
\TextAt{13}{167.12}{78.805}{p}%
\TextAt{8}{172.995}{78.805}{-adic}%
\TextAt{8}{202.392}{78.805}{number}%
\TextAt{8}{245.777}{78.805}{fields.}%
\TextAt{8}{282.666}{78.805}{These}%
\TextAt{8}{317.324}{78.805}{extensions}%
\TextAt{13}{374.493}{78.805}{L}%
\TextAt{16}{387.123}{78.805}{\Character{106}}%
\TextAt{13}{395.109}{78.805}{K}%
\TextAt{8}{410.612}{78.805}{are}%
\TextAt{8}{430.912}{78.805}{classified}%
\TextAt{8}{480.597}{78.805}{by}%
\TextAt{8}{497.645}{78.805}{certain}%
\TextAt{8}{297.211}{43.033}{251}%
\Destination{79.23}{392.285}{nav.Item.81}%
\Destination{79.23}{375.225}{nav.Item.82}%
\Destination{79.23}{358.175}{nav.Item.83}%
\Destination{78.23}{809.815}{nav.page.251}%
\Destination{79.23}{233.655}{nav.section.4.3}%
\Destination{79.226}{668.555}{nav.ex.4.2.1}%
\Destination{79.226}{561.455}{nav.ex.4.2.2}%
\Destination{79.226}{512.455}{nav.ex.4.2.3}%
\Destination{79.226}{446.355}{nav.ex.4.2.4}%
\Destination{79.226}{326.155}{nav.ex.4.2.5}%
\Destination{79.226}{294.155}{nav.ex.4.2.6}%
\Destination{79.36}{776.155}{nav.print.274}%
\Destination{79.36}{193.355}{nav.print.275}%
\Destination{79.23}{233.655}{outline.43}%
\Hotlink{430.132}{479.155}{6.858}{11.047}{nav.ex.4.2.1}%
\Hotlink{242.791}{260.955}{6.858}{11.047}{nav.ex.4.1.5}%
\Hotlink{186.482}{260.955}{6.858}{11.047}{nav.section.4.1}%
}

% Page 264 (body page 252)
\Page{%
\TextAt{2}{79.226}{796.943}{IV.}%
\TextAt{2}{98.012}{796.943}{Abstract}%
\TextAt{2}{143.252}{796.943}{Class}%
\TextAt{2}{171.844}{796.943}{Field}%
\TextAt{2}{199.575}{796.943}{Theory}%
\Vector{q 0 G 0 g 0.299 w 7.226 18.729 m 460.769 18.729 l S Q}%
\TextAt{8}{79.226}{761.124}{subgroups}%
\TextAt{16}{134.56}{761.124}{\Character{78}}%
\TextAt{21}{144.369}{759.33}{L}%
\TextAt{9}{153.947}{761.124}{\Character{61}}%
\TextAt{13}{166.372}{761.124}{N}%
\TextAt{21}{175.746}{759.264}{L}%
\TextAt{3}{181.505}{759.264}{\Character{106}}%
\TextAt{21}{183.857}{759.264}{K}%
\TextAt{13}{192.06}{761.124}{A}%
\TextAt{21}{200.835}{759.33}{L}%
\TextAt{8}{210.911}{761.124}{of}%
\TextAt{8}{224.163}{761.124}{a}%
\TextAt{8}{233.839}{761.124}{group}%
\TextAt{13}{266.936}{761.124}{A}%
\TextAt{21}{275.712}{759.33}{K}%
\TextAt{8}{287.733}{761.124}{attached}%
\TextAt{8}{335.463}{761.124}{to}%
\TextAt{8}{349.695}{761.124}{the}%
\TextAt{8}{369.772}{761.124}{base}%
\TextAt{8}{395.768}{761.124}{field.}%
\TextAt{8}{425.65}{761.124}{In}%
\TextAt{8}{440.192}{761.124}{the}%
\TextAt{8}{460.27}{761.124}{local}%
\TextAt{8}{487.832}{761.124}{case,}%
\TextAt{13}{515.791}{761.124}{A}%
\TextAt{21}{524.567}{759.33}{K}%
\TextAt{8}{79.226}{744.078}{is}%
\TextAt{8}{90.627}{744.078}{the}%
\TextAt{8}{110.42}{744.078}{multiplicative}%
\TextAt{8}{183.88}{744.078}{group}%
\TextAt{13}{216.692}{744.078}{K}%
\TextAt{3}{227.509}{748.416}{\Character{3}}%
\TextAt{8}{235.776}{744.078}{and}%
\TextAt{8}{258.175}{744.078}{in}%
\TextAt{8}{271.465}{744.078}{the}%
\TextAt{8}{291.258}{744.078}{global}%
\TextAt{8}{325.373}{744.078}{case}%
\TextAt{8}{349.781}{744.078}{it}%
\TextAt{8}{361.121}{744.078}{is}%
\TextAt{8}{372.522}{744.078}{a}%
\TextAt{8}{381.914}{744.078}{modification}%
\TextAt{8}{448.87}{744.078}{of}%
\TextAt{8}{461.837}{744.078}{the}%
\TextAt{8}{481.63}{744.078}{ideal}%
\TextAt{8}{509.23}{744.078}{class}%
\TextAt{8}{79.226}{727.032}{group.}%
\TextAt{8}{116.953}{727.032}{At}%
\TextAt{8}{133.857}{727.032}{the}%
\TextAt{8}{154.014}{727.032}{heart}%
\TextAt{8}{184.583}{727.032}{of}%
\TextAt{8}{197.913}{727.032}{this}%
\TextAt{8}{220.736}{727.032}{theory}%
\TextAt{8}{257.486}{727.032}{there}%
\TextAt{8}{287.398}{727.032}{is}%
\TextAt{8}{299.161}{727.032}{a}%
\TextAt{8}{308.917}{727.032}{mysterious}%
\TextAt{8}{367.581}{727.032}{canonical}%
\TextAt{8}{418.964}{727.032}{isomorphism}%
\TextAt{13}{241.085}{702.326}{G}%
\TextAt{9}{250.318}{702.326}{\Character{40}}%
\TextAt{13}{254.871}{702.326}{L}%
\TextAt{16}{266.156}{702.326}{\Character{106}}%
\TextAt{13}{272.798}{702.326}{K}%
\TextAt{9}{283.615}{702.326}{\Character{41}}%
\NativeMathAt{291.489}{702.326}{11.955}{\cong}%
\TextAt{13}{304.108}{702.326}{A}%
\TextAt{21}{312.883}{700.533}{K}%
\TextAt{9}{321.086}{702.326}{\Character{47}}%
\TextAt{13}{326.939}{702.326}{N}%
\TextAt{21}{336.313}{700.466}{L}%
\TextAt{3}{342.072}{700.466}{\Character{106}}%
\TextAt{21}{344.424}{700.466}{K}%
\TextAt{13}{352.627}{702.326}{A}%
\TextAt{21}{361.402}{700.533}{L}%
\TextAt{13}{367.659}{702.326}{\Character{59}}%
\TextAt{8}{79.226}{677.62}{which—if}%
\TextAt{8}{131.74}{677.62}{we}%
\TextAt{8}{149.464}{677.62}{view}%
\TextAt{8}{176.943}{677.62}{things}%
\TextAt{8}{212.624}{677.62}{in}%
\TextAt{8}{226.773}{677.62}{the}%
\TextAt{8}{247.426}{677.62}{right}%
\TextAt{8}{276.221}{677.62}{way—encapsulates}%
\TextAt{8}{376.029}{677.62}{the}%
\TextAt{8}{396.682}{677.62}{reciprocity}%
\TextAt{8}{455.388}{677.62}{law}%
\TextAt{8}{477.022}{677.62}{in}%
\TextAt{8}{491.171}{677.62}{its}%
\TextAt{8}{507.986}{677.62}{most}%
\TextAt{8}{79.226}{660.574}{general}%
\TextAt{8}{119.831}{660.574}{form.}%
\TextAt{8}{152.861}{660.574}{Now,}%
\TextAt{8}{183.122}{660.574}{this}%
\TextAt{8}{206.224}{660.574}{map}%
\TextAt{8}{232.518}{660.574}{can}%
\TextAt{8}{254.258}{660.574}{be}%
\TextAt{8}{270.462}{660.574}{abstracted}%
\TextAt{8}{328.042}{660.574}{completely}%
\TextAt{8}{387.177}{660.574}{from}%
\TextAt{8}{415.097}{660.574}{the}%
\TextAt{8}{435.533}{660.574}{field-theoretic}%
\TextAt{8}{509.947}{660.574}{situ-}%
\TextAt{8}{79.226}{643.528}{ation}%
\TextAt{8}{109.353}{643.528}{and}%
\TextAt{8}{132.319}{643.528}{treated}%
\TextAt{8}{172.847}{643.528}{on}%
\TextAt{8}{189.31}{643.528}{a}%
\TextAt{8}{199.268}{643.528}{purely}%
\TextAt{8}{235.564}{643.528}{group}%
\TextAt{8}{268.943}{643.528}{theoretical}%
\TextAt{8}{327.033}{643.528}{basis.}%
\TextAt{8}{360.934}{643.528}{In}%
\TextAt{8}{375.759}{643.528}{this}%
\TextAt{8}{398.785}{643.528}{way,}%
\TextAt{8}{425.053}{643.528}{class}%
\TextAt{8}{452.694}{643.528}{field}%
\TextAt{8}{478.254}{643.528}{theory}%
\TextAt{8}{515.207}{643.528}{can}%
\TextAt{8}{79.226}{626.482}{be}%
\TextAt{8}{95.15}{626.482}{given}%
\TextAt{8}{125.72}{626.482}{an}%
\TextAt{8}{141.979}{626.482}{abstract,}%
\TextAt{8}{190.828}{626.482}{but}%
\TextAt{8}{212.287}{626.482}{elementary}%
\TextAt{8}{272.123}{626.482}{foundation,}%
\TextAt{8}{334.242}{626.482}{to}%
\TextAt{8}{348.552}{626.482}{which}%
\TextAt{8}{382.039}{626.482}{we}%
\TextAt{8}{399.266}{626.482}{will}%
\TextAt{8}{421.371}{626.482}{now}%
\TextAt{8}{445.76}{626.482}{turn.}%
\TextAt{8}{97.159}{606.448}{We}%
\TextAt{8}{117.617}{606.448}{begin}%
\TextAt{8}{149.468}{606.448}{our}%
\TextAt{8}{170.595}{606.448}{considerations}%
\TextAt{8}{247.134}{606.448}{by}%
\TextAt{8}{263.706}{606.448}{giving}%
\TextAt{8}{298.821}{606.448}{ourselves}%
\TextAt{8}{348.688}{606.448}{a}%
\TextAt{8}{358.757}{606.448}{profinite}%
\TextAt{8}{405.899}{606.448}{group}%
\TextAt{13}{439.387}{606.448}{G}%
\TextAt{8}{448.621}{606.448}{.}%
\TextAt{8}{458.009}{606.448}{The}%
\TextAt{8}{482.376}{606.448}{theory}%
\TextAt{8}{519.439}{606.448}{we}%
\TextAt{8}{79.226}{589.402}{are}%
\TextAt{8}{98.102}{589.402}{about}%
\TextAt{8}{130.965}{589.402}{to}%
\TextAt{8}{144.64}{589.402}{develop}%
\TextAt{8}{186.278}{589.402}{is}%
\TextAt{8}{197.407}{589.402}{purely}%
\TextAt{8}{232.865}{589.402}{group}%
\TextAt{8}{265.405}{589.402}{theoretical}%
\TextAt{8}{322.657}{589.402}{in}%
\TextAt{8}{335.674}{589.402}{nature.}%
\TextAt{8}{377.086}{589.402}{However,}%
\TextAt{8}{426.981}{589.402}{the}%
\TextAt{8}{446.503}{589.402}{only}%
\TextAt{8}{471.559}{589.402}{applications}%
\TextAt{8}{79.226}{572.356}{we}%
\TextAt{8}{96.763}{572.356}{have}%
\TextAt{8}{124.067}{572.356}{in}%
\TextAt{8}{138.03}{572.356}{mind}%
\TextAt{8}{168.251}{572.356}{are}%
\TextAt{8}{188.072}{572.356}{field}%
\TextAt{8}{213.738}{572.356}{theoretical,}%
\TextAt{8}{275.264}{572.356}{and}%
\TextAt{8}{298.336}{572.356}{the}%
\TextAt{8}{318.802}{572.356}{language}%
\TextAt{8}{367.901}{572.356}{of}%
\TextAt{8}{381.54}{572.356}{field}%
\TextAt{8}{407.207}{572.356}{theory}%
\TextAt{8}{444.267}{572.356}{allows}%
\TextAt{8}{479.437}{572.356}{immediate}%
\TextAt{8}{79.226}{555.31}{insights}%
\TextAt{8}{123.609}{555.31}{into}%
\TextAt{8}{149.006}{555.31}{the}%
\TextAt{8}{170.817}{555.31}{group}%
\TextAt{8}{205.648}{555.31}{theoretical}%
\TextAt{8}{265.189}{555.31}{relations.}%
\TextAt{8}{322.25}{555.31}{We}%
\TextAt{8}{344.049}{555.31}{will}%
\TextAt{8}{367.809}{555.31}{therefore}%
\TextAt{8}{418.564}{555.31}{formally}%
\TextAt{8}{466.401}{555.31}{interpret}%
\TextAt{8}{516.51}{555.31}{the}%
\TextAt{8}{79.226}{538.264}{profinite}%
\TextAt{8}{126.046}{538.264}{group}%
\TextAt{13}{159.213}{538.264}{G}%
\TextAt{8}{172.335}{538.264}{as}%
\TextAt{8}{186.697}{538.264}{a}%
\TextAt{8}{196.444}{538.264}{Galois}%
\TextAt{8}{232.348}{538.264}{group}%
\TextAt{8}{265.515}{538.264}{in}%
\TextAt{8}{279.16}{538.264}{the}%
\TextAt{8}{299.308}{538.264}{following}%
\TextAt{8}{348.734}{538.264}{way.}%
\TextAt{8}{376.044}{538.264}{(Let}%
\TextAt{8}{401.56}{538.264}{us}%
\TextAt{8}{416.567}{538.264}{remark}%
\TextAt{8}{456.56}{538.264}{in}%
\TextAt{8}{470.205}{538.264}{passing}%
\TextAt{8}{511.298}{538.264}{that}%
\TextAt{8}{79.226}{521.218}{every}%
\TextAt{8}{110.229}{521.218}{profinite}%
\TextAt{8}{157.168}{521.218}{group}%
\TextAt{8}{190.455}{521.218}{is}%
\TextAt{8}{202.329}{521.218}{indeed}%
\TextAt{8}{239.501}{521.218}{the}%
\TextAt{8}{259.768}{521.218}{Galois}%
\TextAt{8}{295.793}{521.218}{group}%
\TextAt{13}{329.079}{521.218}{G}%
\TextAt{9}{341.822}{521.218}{\Character{61}}%
\TextAt{13}{354.436}{521.218}{G}%
\TextAt{9}{363.67}{521.218}{\Character{40}}%
\TextAt{9}{368.541}{524.373}{\Character{22}}%
\TextAt{13}{368.222}{521.218}{k}%
\TextAt{16}{378.222}{521.218}{\Character{106}}%
\TextAt{13}{385.052}{521.218}{k}%
\TextAt{9}{391.542}{521.218}{\Character{41}}%
\TextAt{8}{400.102}{521.218}{of}%
\TextAt{8}{413.543}{521.218}{a}%
\TextAt{8}{423.409}{521.218}{Galois}%
\TextAt{8}{459.433}{521.218}{field}%
\TextAt{8}{484.901}{521.218}{extension}%
\TextAt{9}{79.544}{507.327}{\Character{22}}%
\TextAt{13}{79.226}{504.172}{k}%
\TextAt{16}{92.167}{504.172}{\Character{106}}%
\TextAt{13}{101.94}{504.172}{k}%
\TextAt{8}{108.429}{504.172}{;}%
\TextAt{8}{118.334}{504.172}{this}%
\TextAt{8}{142.994}{504.172}{will}%
\TextAt{8}{166.936}{504.172}{allow}%
\TextAt{8}{199.02}{504.172}{the}%
\TextAt{8}{221.014}{504.172}{reader}%
\TextAt{8}{258.621}{504.172}{to}%
\TextAt{8}{274.768}{504.172}{rely}%
\TextAt{8}{299.691}{504.172}{on}%
\TextAt{8}{317.787}{504.172}{his}%
\TextAt{8}{337.892}{504.172}{standard}%
\TextAt{8}{388.578}{504.172}{knowledge}%
\TextAt{8}{446.999}{504.172}{of}%
\TextAt{8}{462.166}{504.172}{Galois}%
\TextAt{8}{499.917}{504.172}{theory}%
\TextAt{8}{79.226}{487.126}{whenever}%
\TextAt{8}{130.598}{487.126}{the}%
\TextAt{8}{150.754}{487.126}{formal}%
\TextAt{8}{187.504}{487.126}{development}%
\TextAt{8}{255.469}{487.126}{in}%
\TextAt{8}{269.122}{487.126}{terms}%
\TextAt{8}{301.7}{487.126}{of}%
\TextAt{8}{315.03}{487.126}{group}%
\TextAt{8}{348.206}{487.126}{theory}%
\TextAt{8}{384.956}{487.126}{alone}%
\TextAt{8}{415.525}{487.126}{would}%
\TextAt{8}{449.669}{487.126}{seem}%
\TextAt{8}{478.338}{487.126}{odd.)}%
\TextAt{8}{97.159}{467.091}{We}%
\TextAt{8}{117.033}{467.091}{denote}%
\TextAt{8}{154.482}{467.091}{the}%
\TextAt{8}{174.368}{467.091}{closed}%
\TextAt{8}{208.624}{467.091}{subgroups}%
\TextAt{8}{263.766}{467.091}{of}%
\TextAt{13}{276.826}{467.091}{G}%
\TextAt{8}{289.687}{467.091}{by}%
\TextAt{13}{305.676}{467.091}{G}%
\TextAt{21}{314.909}{465.298}{K}%
\TextAt{8}{323.112}{467.091}{,}%
\TextAt{8}{330.045}{467.091}{and}%
\TextAt{8}{352.538}{467.091}{call}%
\TextAt{8}{373.727}{467.091}{these}%
\TextAt{8}{403.428}{467.091}{indices}%
\TextAt{13}{441.582}{467.091}{K}%
\TextAt{8}{456.027}{467.091}{“fields”;}%
\TextAt{13}{500.117}{467.091}{K}%
\TextAt{8}{514.562}{467.091}{will}%
\TextAt{8}{79.226}{450.045}{be}%
\TextAt{8}{96.078}{450.045}{called}%
\TextAt{8}{130.169}{450.045}{the}%
\TextAt{8}{151.252}{450.045}{fixed}%
\TextAt{8}{180.465}{450.045}{field}%
\TextAt{8}{206.75}{450.045}{of}%
\TextAt{13}{221.007}{450.045}{G}%
\TextAt{21}{230.24}{448.252}{K}%
\TextAt{8}{238.443}{450.045}{.}%
\TextAt{8}{249.673}{450.045}{The}%
\TextAt{8}{274.654}{450.045}{field}%
\TextAt{13}{300.938}{450.045}{k}%
\TextAt{8}{312.252}{450.045}{such}%
\TextAt{8}{339.577}{450.045}{that}%
\TextAt{13}{365.873}{450.045}{G}%
\TextAt{21}{375.106}{448.252}{k}%
\TextAt{9}{385.127}{450.045}{\Character{61}}%
\TextAt{13}{399.133}{450.045}{G}%
\TextAt{8}{413.191}{450.045}{is}%
\TextAt{8}{425.882}{450.045}{called}%
\TextAt{8}{459.973}{450.045}{the}%
\TextAt{8}{481.057}{450.045}{base}%
\TextAt{8}{508.058}{450.045}{field,}%
\TextAt{8}{79.226}{432.999}{and}%
\TextAt{9}{102.181}{436.154}{\Character{22}}%
\TextAt{13}{101.863}{432.999}{k}%
\TextAt{8}{112.124}{432.999}{denotes}%
\TextAt{8}{154.331}{432.999}{the}%
\TextAt{8}{174.362}{432.999}{field}%
\TextAt{8}{199.593}{432.999}{satisfying}%
\TextAt{13}{251.627}{432.999}{G}%
\TextAt{17}{261.054}{432.545}{\Character{22}}%
\TextAt{21}{260.861}{430.442}{k}%
\TextAt{9}{269.301}{432.999}{\Character{61}}%
\TextAt{16}{281.727}{432.999}{\Character{102}}%
\TextAt{9}{287.704}{432.999}{1}%
\TextAt{16}{293.557}{432.999}{\Character{103}}%
\TextAt{8}{299.535}{432.999}{.}%
\TextAt{8}{307.941}{432.999}{The}%
\TextAt{8}{331.869}{432.999}{field}%
\TextAt{8}{357.1}{432.999}{belonging}%
\TextAt{8}{409.984}{432.999}{to}%
\TextAt{8}{424.168}{432.999}{the}%
\TextAt{8}{444.198}{432.999}{closure}%
\Vector{q 0 G 0 g 0.478 w 411.154 -330.361 m 427.535 -330.361 l S Q}%
\TextAt{16}{483.154}{432.999}{\Character{104}}%
\TextAt{13}{487.803}{432.999}{\Character{27}}%
\TextAt{16}{494.886}{432.999}{\Character{105}}%
\TextAt{8}{503.306}{432.999}{of}%
\TextAt{8}{516.51}{432.999}{the}%
\TextAt{8}{79.226}{415.953}{cyclic}%
\TextAt{8}{111.253}{415.953}{group}%
\TextAt{16}{144.272}{415.953}{\Character{104}}%
\TextAt{13}{148.922}{415.953}{\Character{27}}%
\TextAt{16}{156.004}{415.953}{\Character{105}}%
\TextAt{9}{163.974}{415.953}{\Character{61}}%
\TextAt{16}{176.399}{415.953}{\Character{102}}%
\TextAt{13}{182.377}{415.953}{\Character{27}}%
\TextAt{21}{189.46}{420.292}{k}%
\TextAt{16}{197.9}{415.953}{\Character{106}}%
\TextAt{13}{204.542}{415.953}{k}%
\NativeMathAt{214.352}{415.953}{11.955}{\in}%
\TextAt{18}{225.643}{415.953}{Z}%
\TextAt{16}{233.613}{415.953}{\Character{103}}%
\TextAt{8}{243.332}{415.953}{generated}%
\TextAt{8}{296.507}{415.953}{by}%
\TextAt{8}{312.61}{415.953}{an}%
\TextAt{8}{328.713}{415.953}{element}%
\TextAt{13}{371.798}{415.953}{\Character{27}}%
\NativeMathAt{382.201}{415.953}{11.955}{\in}%
\TextAt{13}{393.492}{415.953}{G}%
\TextAt{8}{406.467}{415.953}{is}%
\TextAt{8}{418.075}{415.953}{simply}%
\TextAt{8}{455.374}{415.953}{called}%
\TextAt{8}{488.381}{415.953}{the}%
\TextAt{8}{508.381}{415.953}{fixed}%
\TextAt{8}{79.226}{398.907}{field}%
\TextAt{8}{104.583}{398.907}{of}%
\TextAt{13}{117.913}{398.907}{\Character{27}}%
\TextAt{8}{124.995}{398.907}{.}%
\TextAt{8}{97.159}{378.873}{We}%
\TextAt{8}{117.91}{378.873}{write}%
\TextAt{8}{148.429}{378.873}{formally}%
\TextAt{13}{195.219}{378.873}{K}%
\NativeMathAt{210.392}{378.873}{11.955}{\subseteq}%
\TextAt{13}{224.045}{378.873}{L}%
\TextAt{8}{236.514}{378.873}{or}%
\TextAt{13}{251.431}{378.873}{L}%
\TextAt{16}{263.751}{378.873}{\Character{106}}%
\TextAt{13}{271.427}{378.873}{K}%
\TextAt{8}{286.749}{378.873}{if}%
\TextAt{13}{298.08}{378.873}{G}%
\TextAt{21}{307.313}{377.079}{L}%
\NativeMathAt{317.926}{378.873}{11.955}{\subseteq}%
\TextAt{13}{331.579}{378.873}{G}%
\TextAt{21}{340.813}{377.079}{K}%
\TextAt{8}{349.016}{378.873}{,}%
\TextAt{8}{356.923}{378.873}{and}%
\TextAt{8}{380.293}{378.873}{we}%
\TextAt{8}{398.127}{378.873}{call}%
\TextAt{8}{420.194}{378.873}{the}%
\TextAt{8}{440.957}{378.873}{pair}%
\TextAt{13}{465.63}{378.873}{L}%
\TextAt{16}{477.95}{378.873}{\Character{106}}%
\TextAt{13}{485.626}{378.873}{K}%
\TextAt{8}{500.948}{378.873}{a}%
\TextAt{8}{511.31}{378.873}{field}%
\TextAt{8}{79.226}{361.827}{extension.}%
\TextAt{13}{135.491}{361.827}{L}%
\TextAt{16}{146.776}{361.827}{\Character{106}}%
\TextAt{13}{153.418}{361.827}{K}%
\TextAt{8}{167.976}{361.827}{is}%
\TextAt{8}{179.584}{361.827}{called}%
\TextAt{8}{212.591}{361.827}{a}%
\TextAt{8}{222.19}{361.827}{finite}%
\TextAt{8}{251.946}{361.827}{extension,}%
\TextAt{8}{306.838}{361.827}{if}%
\TextAt{13}{317.405}{361.827}{G}%
\TextAt{21}{326.639}{360.033}{L}%
\TextAt{8}{336.637}{361.827}{is}%
\TextAt{8}{348.245}{361.827}{open,}%
\TextAt{8}{379.657}{361.827}{i.e.,}%
\TextAt{8}{401.637}{361.827}{of}%
\TextAt{8}{414.811}{361.827}{finite}%
\TextAt{8}{444.566}{361.827}{index}%
\TextAt{8}{475.947}{361.827}{in}%
\TextAt{13}{489.444}{361.827}{G}%
\TextAt{21}{498.677}{360.033}{K}%
\TextAt{8}{506.88}{361.827}{,}%
\TextAt{8}{513.904}{361.827}{and}%
\TextAt{8}{79.226}{344.781}{this}%
\TextAt{8}{102.049}{344.781}{index}%
\TextAt{9}{252.947}{327.735}{\Character{91}}%
\TextAt{13}{256.199}{327.735}{L}%
\TextAt{9}{267.484}{327.735}{\Character{58}}%
\TextAt{13}{274.056}{327.735}{K}%
\TextAt{9}{284.874}{327.735}{\Character{93}}%
\TextAt{9}{291.425}{327.735}{\Character{58}\Character{61}}%
\TextAt{9}{307.087}{327.735}{\Character{40}}%
\TextAt{13}{311.676}{327.735}{G}%
\TextAt{21}{320.91}{325.941}{K}%
\TextAt{9}{332.433}{327.735}{\Character{58}}%
\TextAt{13}{339.006}{327.735}{G}%
\TextAt{21}{348.239}{325.941}{L}%
\TextAt{9}{354.496}{327.735}{\Character{41}}%
\TextAt{8}{79.226}{306.054}{will}%
\TextAt{8}{101.286}{306.054}{be}%
\TextAt{8}{117.165}{306.054}{called}%
\TextAt{8}{150.284}{306.054}{the}%
\TextAt{8}{170.395}{306.054}{degree}%
\TextAt{8}{206.766}{306.054}{of}%
\TextAt{13}{220.05}{306.054}{L}%
\TextAt{16}{231.336}{306.054}{\Character{106}}%
\TextAt{13}{237.978}{306.054}{K}%
\TextAt{8}{248.795}{306.054}{.}%
\TextAt{13}{257.229}{306.054}{L}%
\TextAt{16}{268.514}{306.054}{\Character{106}}%
\TextAt{13}{275.156}{306.054}{K}%
\TextAt{8}{289.826}{306.054}{is}%
\TextAt{8}{301.544}{306.054}{said}%
\TextAt{8}{325.625}{306.054}{to}%
\TextAt{8}{339.89}{306.054}{be}%
\TextAt{8}{355.769}{306.054}{normal}%
\TextAt{8}{395.403}{306.054}{or}%
\TextAt{8}{409.669}{306.054}{Galois}%
\TextAt{8}{445.537}{306.054}{if}%
\TextAt{13}{456.215}{306.054}{G}%
\TextAt{21}{465.449}{304.261}{L}%
\TextAt{8}{475.559}{306.054}{is}%
\TextAt{8}{487.277}{306.054}{a}%
\TextAt{8}{496.988}{306.054}{normal}%
\TextAt{8}{79.226}{289.008}{subgroup}%
\TextAt{8}{130.024}{289.008}{of}%
\TextAt{13}{143.354}{289.008}{G}%
\TextAt{21}{152.587}{287.215}{K}%
\TextAt{8}{160.79}{289.008}{.}%
\TextAt{8}{169.238}{289.008}{If}%
\TextAt{8}{180.93}{289.008}{this}%
\TextAt{8}{203.753}{289.008}{is}%
\TextAt{8}{215.517}{289.008}{the}%
\TextAt{8}{235.673}{289.008}{case,}%
\TextAt{8}{263.696}{289.008}{we}%
\TextAt{8}{280.924}{289.008}{define}%
\TextAt{8}{314.733}{289.008}{the}%
\TextAt{8}{334.889}{289.008}{Galois}%
\TextAt{8}{370.802}{289.008}{group}%
\TextAt{8}{403.978}{289.008}{of}%
\TextAt{13}{417.308}{289.008}{L}%
\TextAt{16}{428.593}{289.008}{\Character{106}}%
\TextAt{13}{435.235}{289.008}{K}%
\TextAt{8}{449.95}{289.008}{by}%
\TextAt{13}{253.567}{264.302}{G}%
\TextAt{9}{262.801}{264.302}{\Character{40}}%
\TextAt{13}{267.353}{264.302}{L}%
\TextAt{16}{278.639}{264.302}{\Character{106}}%
\TextAt{13}{285.28}{264.302}{K}%
\TextAt{9}{296.098}{264.302}{\Character{41}}%
\TextAt{9}{303.953}{264.302}{\Character{61}}%
\TextAt{13}{316.397}{264.302}{G}%
\TextAt{21}{325.63}{262.509}{K}%
\TextAt{9}{333.833}{264.302}{\Character{47}}%
\TextAt{13}{339.686}{264.302}{G}%
\TextAt{21}{348.919}{262.509}{L}%
\TextAt{13}{355.177}{264.302}{\Character{58}}%
\TextAt{8}{79.226}{239.597}{If}%
\TextAt{13}{92.784}{239.597}{N}%
\NativeMathAt{109.907}{239.596}{11.955}{\supseteq}%
\TextAt{13}{125.707}{239.597}{L}%
\NativeMathAt{140.172}{239.596}{11.955}{\supseteq}%
\TextAt{13}{155.971}{239.597}{K}%
\TextAt{8}{172.552}{239.597}{are}%
\TextAt{8}{193.928}{239.597}{Galois}%
\TextAt{8}{231.708}{239.597}{extensions}%
\TextAt{8}{289.954}{239.597}{of}%
\TextAt{13}{305.15}{239.597}{K}%
\TextAt{8}{315.968}{239.597}{,}%
\TextAt{8}{325.449}{239.597}{we}%
\TextAt{8}{344.543}{239.597}{define}%
\TextAt{8}{380.218}{239.597}{the}%
\TextAt{8}{402.24}{239.597}{restriction}%
\TextAt{8}{460.104}{239.597}{of}%
\TextAt{8}{475.3}{239.597}{an}%
\TextAt{8}{493.425}{239.597}{element}%
\TextAt{13}{79.226}{222.551}{\Character{27}}%
\NativeMathAt{89.629}{222.551}{11.955}{\in}%
\TextAt{13}{100.92}{222.551}{G}%
\TextAt{9}{110.154}{222.551}{\Character{40}}%
\TextAt{13}{114.706}{222.551}{N}%
\TextAt{16}{128.65}{222.551}{\Character{106}}%
\TextAt{13}{135.291}{222.551}{K}%
\TextAt{9}{146.109}{222.551}{\Character{41}}%
\TextAt{8}{154.559}{222.551}{to}%
\TextAt{13}{168.869}{222.551}{L}%
\TextAt{8}{180.731}{222.551}{by}%
\TextAt{13}{215.957}{197.845}{\Character{27}}%
\TextAt{16}{223.039}{197.845}{\Character{106}}%
\TextAt{21}{226.36}{196.052}{L}%
\TextAt{9}{235.938}{197.845}{\Character{61}}%
\TextAt{13}{248.363}{197.845}{\Character{27}}%
\TextAt{8}{258.767}{197.845}{mod}%
\TextAt{13}{284.205}{197.845}{G}%
\TextAt{9}{293.438}{197.845}{\Character{40}}%
\TextAt{13}{297.991}{197.845}{N}%
\TextAt{16}{311.934}{197.845}{\Character{106}}%
\TextAt{13}{318.576}{197.845}{L}%
\TextAt{9}{326.54}{197.845}{\Character{41}}%
\NativeMathAt{334.413}{197.845}{11.955}{\in}%
\TextAt{13}{345.704}{197.845}{G}%
\TextAt{9}{354.938}{197.845}{\Character{40}}%
\TextAt{13}{359.49}{197.845}{L}%
\TextAt{16}{370.776}{197.845}{\Character{106}}%
\TextAt{13}{377.417}{197.845}{K}%
\TextAt{9}{388.235}{197.845}{\Character{41}}%
\TextAt{13}{392.787}{197.845}{\Character{58}}%
\TextAt{8}{79.226}{173.139}{This}%
\TextAt{8}{105.946}{173.139}{gives}%
\TextAt{8}{134.626}{173.139}{a}%
\TextAt{8}{144.382}{173.139}{homomorphism}%
\TextAt{13}{203.608}{148.433}{G}%
\TextAt{9}{212.842}{148.434}{\Character{40}}%
\TextAt{13}{217.394}{148.433}{N}%
\TextAt{16}{231.338}{148.434}{\Character{106}}%
\TextAt{13}{237.979}{148.433}{K}%
\TextAt{9}{248.797}{148.434}{\Character{41}}%
\NativeArrowAt{256.67}{148.433}{11.955}{19.26}{\longrightarrow}%
\TextAt{13}{279.252}{148.434}{G}%
\TextAt{9}{288.486}{148.433}{\Character{40}}%
\TextAt{13}{293.038}{148.434}{L}%
\TextAt{16}{304.324}{148.433}{\Character{106}}%
\TextAt{13}{310.965}{148.434}{K}%
\TextAt{9}{321.783}{148.433}{\Character{41}}%
\TextAt{13}{326.335}{148.434}{\Character{59}}%
\TextAt{13}{355.47}{148.434}{\Character{27}}%
\NativeArrowAt{365.893}{148.433}{11.955}{19.26}{\longmapsto}%
\TextAt{13}{388.475}{148.434}{\Character{27}}%
\TextAt{16}{395.557}{148.433}{\Character{106}}%
\TextAt{21}{398.878}{146.64}{L}%
\TextAt{13}{405.136}{148.434}{\Character{59}}%
\TextAt{8}{79.226}{123.728}{with}%
\TextAt{8}{106.096}{123.728}{kernel}%
\TextAt{13}{140.773}{123.728}{G}%
\TextAt{9}{150.007}{123.728}{\Character{40}}%
\TextAt{13}{154.559}{123.728}{N}%
\TextAt{16}{168.861}{123.728}{\Character{106}}%
\TextAt{13}{175.86}{123.728}{L}%
\TextAt{9}{183.825}{123.728}{\Character{41}}%
\TextAt{8}{188.377}{123.728}{.}%
\TextAt{8}{197.456}{123.728}{The}%
\TextAt{8}{221.72}{123.728}{extension}%
\TextAt{13}{273.696}{123.728}{L}%
\TextAt{16}{285.339}{123.728}{\Character{106}}%
\TextAt{13}{292.339}{123.728}{K}%
\TextAt{8}{307.264}{123.728}{is}%
\TextAt{8}{319.238}{123.728}{called}%
\TextAt{8}{352.612}{123.728}{cyclic,}%
\TextAt{8}{388.309}{123.728}{abelian,}%
\TextAt{8}{432.471}{123.728}{solvable,}%
\TextAt{8}{479.957}{123.728}{etc.,}%
\TextAt{8}{505.577}{123.728}{if}%
\TextAt{8}{516.51}{123.728}{the}%
\TextAt{8}{79.226}{106.682}{Galois}%
\TextAt{8}{115.139}{106.682}{group}%
\TextAt{13}{148.315}{106.682}{G}%
\TextAt{9}{157.549}{106.682}{\Character{40}}%
\TextAt{13}{162.101}{106.682}{L}%
\TextAt{16}{173.386}{106.682}{\Character{106}}%
\TextAt{13}{180.028}{106.682}{K}%
\TextAt{9}{190.846}{106.682}{\Character{41}}%
\TextAt{8}{199.295}{106.682}{has}%
\TextAt{8}{220.169}{106.682}{these}%
\TextAt{8}{250.141}{106.682}{properties.}%
\TextAt{8}{309.709}{106.682}{We}%
\TextAt{8}{329.854}{106.682}{put}%
\TextAt{13}{226.911}{81.976}{K}%
\TextAt{9}{241.049}{81.976}{\Character{61}}%
\TextAt{20}{253.475}{93.334}{\Character{92}}%
\TextAt{21}{258.675}{68.139}{i}%
\TextAt{13}{268.751}{81.976}{K}%
\TextAt{21}{278.71}{80.183}{i}%
\TextAt{8}{306.001}{81.976}{(“intersection”)}%
\TextAt{8}{297.211}{43.033}{252}%
\Destination{78.23}{809.815}{nav.page.252}%
\Destination{79.36}{321.055}{nav.print.276}%
}

% Page 265 (body page 253)
\Page{%
\TextAt{2}{79.226}{796.943}{IV.}%
\TextAt{2}{98.012}{796.943}{Abstract}%
\TextAt{2}{143.252}{796.943}{Class}%
\TextAt{2}{171.844}{796.943}{Field}%
\TextAt{2}{199.575}{796.943}{Theory}%
\Vector{q 0 G 0 g 0.299 w 7.226 18.729 m 460.769 18.729 l S Q}%
\TextAt{8}{79.226}{761.124}{if}%
\TextAt{13}{89.95}{761.124}{G}%
\TextAt{21}{99.184}{759.33}{K}%
\TextAt{8}{111.284}{761.124}{is}%
\TextAt{8}{123.048}{761.124}{topologically}%
\TextAt{8}{192.005}{761.124}{generated}%
\TextAt{8}{245.337}{761.124}{by}%
\TextAt{8}{261.596}{761.124}{the}%
\TextAt{8}{281.752}{761.124}{subgroups}%
\TextAt{13}{337.165}{761.124}{G}%
\TextAt{21}{346.398}{759.33}{K}%
\TextAt{22}{353.54}{758.115}{i}%
\TextAt{8}{357.199}{761.124}{,}%
\TextAt{8}{364.349}{761.124}{and}%
\TextAt{13}{229.818}{735.222}{K}%
\TextAt{9}{243.956}{735.222}{\Character{61}}%
\TextAt{20}{256.382}{746.579}{\Character{89}}%
\TextAt{21}{262.578}{721.385}{i}%
\TextAt{13}{273.65}{735.222}{K}%
\TextAt{21}{283.609}{733.429}{i}%
\TextAt{8}{310.901}{735.222}{(“composite”)}%
\TextAt{8}{79.226}{700.863}{if}%
\TextAt{13}{89.95}{700.863}{G}%
\TextAt{21}{99.184}{699.07}{K}%
\TextAt{9}{110.707}{700.863}{\Character{61}}%
\TextAt{20}{123.133}{709.829}{\Character{84}}%
\TextAt{21}{133.095}{697.376}{i}%
\TextAt{13}{138.469}{700.863}{G}%
\TextAt{21}{147.703}{699.07}{K}%
\TextAt{22}{154.844}{697.854}{i}%
\TextAt{8}{158.504}{700.863}{.}%
\TextAt{8}{166.952}{700.863}{If}%
\TextAt{13}{178.644}{700.863}{G}%
\TextAt{21}{187.878}{699.07}{K}%
\TextAt{23}{195.583}{701.339}{\Character{48}}%
\TextAt{9}{202.086}{700.863}{\Character{61}}%
\TextAt{13}{214.511}{700.863}{\Character{27}}%
\TextAt{3}{221.594}{705.201}{\Character{0}}%
\TextAt{17}{228.18}{705.201}{1}%
\TextAt{13}{232.913}{700.863}{G}%
\TextAt{21}{242.146}{699.07}{K}%
\TextAt{13}{250.349}{700.863}{\Character{27}}%
\TextAt{8}{261.329}{700.863}{for}%
\TextAt{13}{279.214}{700.863}{\Character{27}}%
\NativeMathAt{289.617}{700.863}{11.955}{\in}%
\TextAt{13}{300.908}{700.863}{G}%
\TextAt{8}{310.141}{700.863}{,}%
\TextAt{8}{317.291}{700.863}{we}%
\TextAt{8}{334.518}{700.863}{write}%
\TextAt{13}{364.43}{700.863}{K}%
\TextAt{3}{375.248}{705.201}{\Character{48}}%
\TextAt{9}{381.363}{700.863}{\Character{61}}%
\TextAt{13}{393.789}{700.863}{K}%
\TextAt{21}{404.607}{705.201}{\Character{27}}%
\TextAt{8}{410.23}{700.863}{.}%
\TextAt{8}{97.159}{680.828}{Now}%
\TextAt{8}{123.502}{680.828}{let}%
\TextAt{13}{140.089}{680.828}{A}%
\TextAt{8}{152.445}{680.828}{be}%
\TextAt{8}{168.052}{680.828}{a}%
\TextAt{19}{177.49}{680.828}{continuous}%
\TextAt{19}{236.029}{680.828}{multiplicative}%
\TextAt{13}{307.719}{680.828}{G}%
\TextAt{19}{316.952}{680.828}{-module}%
\TextAt{8}{356.703}{680.828}{.}%
\TextAt{8}{365.046}{680.828}{By}%
\TextAt{8}{383.092}{680.828}{this}%
\TextAt{8}{405.597}{680.828}{we}%
\TextAt{8}{422.508}{680.828}{mean}%
\TextAt{8}{453.405}{680.828}{a}%
\TextAt{8}{462.844}{680.828}{multiplicative}%
\TextAt{8}{79.226}{663.782}{abelian}%
\TextAt{8}{120.605}{663.782}{group}%
\TextAt{13}{154.511}{663.782}{A}%
\TextAt{8}{167.914}{663.782}{on}%
\TextAt{8}{184.904}{663.782}{which}%
\TextAt{8}{219.121}{663.782}{the}%
\TextAt{8}{240.009}{663.782}{elements}%
\TextAt{13}{288.596}{663.782}{\Character{27}}%
\NativeMathAt{300.245}{663.782}{11.955}{\in}%
\TextAt{13}{312.781}{663.782}{G}%
\TextAt{8}{326.643}{663.782}{operate}%
\TextAt{8}{369.324}{663.782}{as}%
\TextAt{8}{384.425}{663.782}{automorphisms}%
\TextAt{8}{467.24}{663.782}{on}%
\TextAt{8}{484.23}{663.782}{the}%
\TextAt{8}{505.117}{663.782}{right,}%
\TextAt{13}{79.226}{646.736}{\Character{27}}%
\TextAt{9}{89.629}{646.736}{\Character{58}}%
\TextAt{13}{96.202}{646.736}{A}%
\NativeMathAt{108.298}{646.736}{11.955}{\rightarrow}%
\TextAt{13}{123.574}{646.736}{A}%
\TextAt{8}{132.349}{646.736}{,}%
\TextAt{13}{139.499}{646.736}{a}%
\NativeArrowAt{148.964}{646.736}{11.955}{11.955}{\mapsto}%
\TextAt{13}{164.24}{646.736}{a}%
\TextAt{21}{170.385}{651.074}{\Character{27}}%
\TextAt{8}{176.009}{646.736}{.}%
\TextAt{8}{184.457}{646.736}{This}%
\TextAt{8}{211.177}{646.736}{action}%
\TextAt{8}{246.301}{646.736}{must}%
\TextAt{8}{275.304}{646.736}{satisfy}%
\TextAt{8}{89.579}{626.701}{(i)}%
\TextAt{13}{107.918}{626.701}{a}%
\TextAt{17}{114.063}{631.04}{1}%
\TextAt{9}{122.117}{626.701}{\Character{61}}%
\TextAt{13}{134.542}{626.701}{a}%
\TextAt{8}{140.687}{626.701}{,}%
\TextAt{8}{86.327}{609.655}{(ii)}%
\TextAt{9}{107.918}{609.655}{\Character{40}}%
\TextAt{13}{112.471}{609.655}{ab}%
\TextAt{9}{123.593}{609.655}{\Character{41}}%
\TextAt{21}{128.145}{613.994}{\Character{27}}%
\TextAt{9}{137.089}{609.655}{\Character{61}}%
\TextAt{13}{149.515}{609.655}{a}%
\TextAt{21}{155.66}{613.994}{\Character{27}}%
\TextAt{13}{161.283}{609.655}{b}%
\TextAt{21}{166.26}{613.994}{\Character{27}}%
\TextAt{8}{171.884}{609.655}{,}%
\TextAt{8}{83.076}{592.609}{(iii)}%
\TextAt{13}{107.918}{592.609}{a}%
\TextAt{21}{114.063}{596.948}{\Character{27}}%
\TextAt{21}{119.18}{596.948}{\Character{28}}%
\TextAt{9}{127.686}{592.609}{\Character{61}}%
\TextAt{9}{140.096}{592.609}{\Character{40}}%
\TextAt{13}{144.664}{592.609}{a}%
\TextAt{21}{150.809}{596.948}{\Character{27}}%
\TextAt{9}{156.432}{592.609}{\Character{41}}%
\TextAt{21}{160.985}{596.948}{\Character{28}}%
\TextAt{8}{166.161}{592.609}{,}%
\TextAt{8}{83.398}{575.563}{(iv)}%
\TextAt{13}{107.918}{575.563}{A}%
\TextAt{9}{120.015}{575.563}{\Character{61}}%
\TextAt{20}{132.44}{584.53}{\Character{83}}%
\TextAt{17}{142.403}{572.076}{\Character{91}}%
\TextAt{21}{144.755}{572.076}{K}%
\TextAt{17}{152.46}{572.076}{\Character{58}}%
\TextAt{21}{154.812}{572.076}{k}%
\TextAt{17}{159.434}{572.076}{\Character{93}}%
\TextAt{21}{161.786}{572.076}{\Character{60}}%
\NativeMathAt{168.372}{572.076}{7.97}{\infty}%
\TextAt{13}{179.331}{575.563}{A}%
\TextAt{21}{188.107}{573.77}{K}%
\TextAt{8}{196.309}{575.563}{,}%
\TextAt{8}{79.226}{555.529}{where}%
\TextAt{13}{113.035}{555.529}{A}%
\TextAt{21}{121.811}{553.735}{K}%
\TextAt{8}{133.911}{555.529}{in}%
\TextAt{8}{147.564}{555.529}{the}%
\TextAt{8}{167.72}{555.529}{last}%
\TextAt{8}{189.897}{555.529}{condition}%
\TextAt{8}{241.28}{555.529}{denotes}%
\TextAt{8}{283.613}{555.529}{the}%
\TextAt{8}{303.77}{555.529}{fixed}%
\TextAt{8}{332.056}{555.529}{module}%
\TextAt{13}{373.349}{555.529}{A}%
\TextAt{21}{382.124}{559.867}{G}%
\TextAt{22}{388.723}{558.528}{K}%
\TextAt{8}{400.161}{555.529}{under}%
\TextAt{13}{433.325}{555.529}{G}%
\TextAt{21}{442.559}{553.735}{K}%
\TextAt{8}{450.761}{555.529}{,}%
\TextAt{8}{457.911}{555.529}{so}%
\TextAt{8}{472.281}{555.529}{that}%
\TextAt{13}{206.869}{529.627}{A}%
\TextAt{21}{215.645}{527.834}{K}%
\TextAt{9}{227.168}{529.627}{\Character{61}}%
\TextAt{16}{239.594}{529.627}{\Character{102}}%
\TextAt{13}{245.571}{529.627}{a}%
\NativeMathAt{255.037}{529.627}{11.955}{\in}%
\TextAt{13}{266.328}{529.627}{A}%
\TextAt{16}{278.424}{529.627}{\Character{106}}%
\TextAt{13}{285.066}{529.627}{a}%
\TextAt{21}{291.211}{534.563}{\Character{27}}%
\TextAt{9}{300.155}{529.627}{\Character{61}}%
\TextAt{13}{312.58}{529.627}{a}%
\TextAt{8}{322.623}{529.627}{for}%
\TextAt{8}{340.508}{529.627}{all}%
\TextAt{13}{356.767}{529.627}{\Character{27}}%
\NativeMathAt{367.17}{529.627}{11.955}{\in}%
\TextAt{13}{378.461}{529.627}{G}%
\TextAt{21}{387.695}{527.834}{K}%
\TextAt{16}{395.897}{529.627}{\Character{103}}%
\TextAt{13}{401.875}{529.627}{\Character{59}}%
\TextAt{8}{79.226}{503.725}{and}%
\TextAt{8}{101.512}{503.725}{where}%
\TextAt{13}{134.845}{503.725}{K}%
\TextAt{8}{149.083}{503.725}{varies}%
\TextAt{8}{181.519}{503.725}{over}%
\TextAt{8}{206.088}{503.725}{all}%
\TextAt{8}{221.871}{503.725}{extensions}%
\TextAt{8}{277.775}{503.725}{that}%
\TextAt{8}{302.667}{503.725}{are}%
\TextAt{8}{321.701}{503.725}{finite}%
\TextAt{8}{351.136}{503.725}{over}%
\TextAt{13}{375.706}{503.725}{k}%
\TextAt{8}{382.195}{503.725}{.}%
\TextAt{8}{390.485}{503.725}{The}%
\TextAt{8}{414.062}{503.725}{condition}%
\TextAt{8}{464.969}{503.725}{(iv)}%
\TextAt{8}{486.932}{503.725}{says}%
\TextAt{8}{511.298}{503.725}{that}%
\TextAt{13}{79.226}{486.679}{G}%
\TextAt{8}{92.357}{486.679}{operates}%
\TextAt{8}{138.923}{486.679}{continuously}%
\TextAt{8}{206.959}{486.679}{on}%
\TextAt{13}{223.218}{486.679}{A}%
\TextAt{8}{231.994}{486.679}{,}%
\TextAt{8}{239.143}{486.679}{i.e.,}%
\TextAt{8}{261.248}{486.679}{the}%
\TextAt{8}{281.405}{486.679}{map}%
\TextAt{13}{225.473}{460.777}{G}%
\TextAt{16}{237.363}{460.777}{\Character{2}}%
\TextAt{13}{249.318}{460.777}{A}%
\NativeArrowAt{261.414}{460.777}{11.955}{19.26}{\longrightarrow}%
\TextAt{13}{283.996}{460.777}{A\Character{59}}%
\TextAt{9}{321.926}{460.777}{\Character{40}}%
\TextAt{13}{326.478}{460.777}{\Character{27}}%
\TextAt{13}{332.91}{460.777}{\Character{59}}%
\TextAt{13}{338.134}{460.777}{a}%
\TextAt{9}{344.3}{460.777}{\Character{41}}%
\NativeArrowAt{352.173}{460.777}{11.955}{19.26}{\longmapsto}%
\TextAt{13}{374.755}{460.777}{a}%
\TextAt{21}{380.9}{465.714}{\Character{27}}%
\TextAt{8}{79.226}{434.876}{is}%
\TextAt{8}{91.173}{434.876}{continuous,}%
\TextAt{8}{153.258}{434.876}{where}%
\TextAt{13}{187.25}{434.876}{A}%
\TextAt{8}{200.106}{434.876}{is}%
\TextAt{8}{212.053}{434.876}{equipped}%
\TextAt{8}{262.304}{434.876}{with}%
\TextAt{8}{289.147}{434.876}{the}%
\TextAt{8}{309.487}{434.876}{discrete}%
\TextAt{8}{352.649}{434.876}{topology.}%
\TextAt{8}{404.912}{434.876}{Indeed,}%
\TextAt{8}{446.422}{434.876}{this}%
\TextAt{8}{469.427}{434.876}{continuity}%
\TextAt{8}{524.903}{434.876}{is}%
\TextAt{8}{79.226}{417.83}{equivalent}%
\TextAt{8}{133.928}{417.83}{to}%
\TextAt{8}{147.325}{417.83}{the}%
\TextAt{8}{166.568}{417.83}{fact}%
\TextAt{8}{188.74}{417.83}{that,}%
\TextAt{8}{216.631}{417.83}{for}%
\TextAt{8}{233.602}{417.83}{every}%
\TextAt{8}{263.581}{417.83}{element}%
\TextAt{9}{305.909}{417.83}{\Character{40}}%
\TextAt{13}{310.462}{417.83}{\Character{27}}%
\TextAt{13}{316.894}{417.83}{\Character{59}}%
\TextAt{13}{322.118}{417.83}{a}%
\TextAt{9}{328.283}{417.83}{\Character{41}}%
\NativeMathAt{336.156}{417.83}{11.955}{\in}%
\TextAt{13}{347.447}{417.83}{G}%
\TextAt{16}{357.47}{417.83}{\Character{2}}%
\TextAt{13}{367.557}{417.83}{A}%
\TextAt{8}{376.333}{417.83}{,}%
\TextAt{8}{382.751}{417.83}{there}%
\TextAt{8}{411.75}{417.83}{exists}%
\TextAt{8}{443.151}{417.83}{an}%
\TextAt{8}{458.497}{417.83}{open}%
\TextAt{8}{485.869}{417.83}{subgroup}%
\TextAt{13}{79.226}{400.784}{U}%
\TextAt{9}{91.746}{400.784}{\Character{61}}%
\TextAt{13}{104.172}{400.784}{G}%
\TextAt{21}{113.405}{398.991}{K}%
\TextAt{8}{125.421}{400.784}{of}%
\TextAt{13}{138.668}{400.784}{G}%
\TextAt{8}{151.715}{400.784}{such}%
\TextAt{8}{178.028}{400.784}{that}%
\TextAt{8}{203.313}{400.784}{the}%
\TextAt{8}{223.386}{400.784}{neighbourhood}%
\TextAt{13}{303.306}{400.784}{\Character{27}}%
\TextAt{13}{310.372}{400.784}{U}%
\TextAt{16}{322.073}{400.784}{\Character{2}}%
\TextAt{16}{333.837}{400.784}{\Character{102}}%
\TextAt{13}{339.835}{400.784}{a}%
\TextAt{16}{345.98}{400.784}{\Character{103}}%
\TextAt{8}{355.771}{400.784}{of}%
\TextAt{9}{369.017}{400.784}{\Character{40}}%
\TextAt{13}{373.569}{400.784}{\Character{27}}%
\TextAt{13}{380.001}{400.784}{\Character{59}}%
\TextAt{13}{385.225}{400.784}{a}%
\TextAt{9}{391.391}{400.784}{\Character{41}}%
\TextAt{8}{399.757}{400.784}{is}%
\TextAt{8}{411.437}{400.784}{mapped}%
\TextAt{8}{455.898}{400.784}{to}%
\TextAt{8}{470.125}{400.784}{the}%
\TextAt{8}{490.197}{400.784}{open}%
\TextAt{8}{518.399}{400.784}{set}%
\TextAt{16}{79.226}{383.738}{\Character{102}}%
\TextAt{13}{85.204}{383.738}{a}%
\TextAt{21}{91.349}{388.076}{\Character{27}}%
\TextAt{16}{96.972}{383.738}{\Character{103}}%
\TextAt{8}{102.95}{383.738}{,}%
\TextAt{8}{110.099}{383.738}{and}%
\TextAt{8}{132.862}{383.738}{this}%
\TextAt{8}{155.684}{383.738}{means}%
\TextAt{8}{191.514}{383.738}{simply}%
\TextAt{8}{228.969}{383.738}{that}%
\TextAt{13}{254.338}{383.738}{a}%
\TextAt{21}{260.483}{388.076}{\Character{27}}%
\NativeMathAt{269.427}{383.738}{11.955}{\in}%
\TextAt{13}{280.718}{383.738}{A}%
\TextAt{21}{289.494}{388.076}{U}%
\TextAt{9}{299.89}{383.738}{\Character{61}}%
\TextAt{13}{312.316}{383.738}{A}%
\TextAt{21}{321.091}{381.945}{K}%
\TextAt{8}{329.294}{383.738}{.}%
\TextAt{12}{79.226}{361.157}{Remark.}%
\TextAt{8}{132.544}{361.157}{In}%
\TextAt{8}{146.326}{361.157}{the}%
\TextAt{8}{165.644}{361.157}{exponential}%
\TextAt{8}{227.57}{361.157}{notation}%
\TextAt{13}{273.571}{361.157}{a}%
\TextAt{21}{279.716}{365.495}{\Character{27}}%
\TextAt{8}{285.339}{361.157}{,}%
\TextAt{8}{291.817}{361.157}{the}%
\TextAt{8}{311.135}{361.157}{operation}%
\TextAt{8}{362.66}{361.157}{of}%
\TextAt{13}{375.151}{361.157}{G}%
\TextAt{8}{387.443}{361.157}{on}%
\TextAt{13}{402.863}{361.157}{A}%
\TextAt{8}{414.697}{361.157}{appears}%
\TextAt{8}{457.171}{361.157}{as}%
\TextAt{8}{470.702}{361.157}{an}%
\TextAt{8}{486.122}{361.157}{action}%
\TextAt{8}{520.408}{361.157}{on}%
\TextAt{8}{79.226}{344.111}{the}%
\TextAt{8}{99.446}{344.111}{right.}%
\TextAt{8}{132.484}{344.111}{This}%
\TextAt{8}{159.267}{344.111}{notation}%
\TextAt{8}{206.171}{344.111}{is}%
\TextAt{8}{217.998}{344.111}{adequate}%
\TextAt{8}{267.819}{344.111}{for}%
\TextAt{8}{285.767}{344.111}{many}%
\TextAt{8}{317.703}{344.111}{computations}%
\TextAt{8}{390.68}{344.111}{in}%
\TextAt{8}{404.396}{344.111}{the}%
\TextAt{8}{424.616}{344.111}{case}%
\TextAt{8}{449.45}{344.111}{of}%
\TextAt{8}{462.844}{344.111}{multiplicative}%
\TextAt{13}{79.226}{327.065}{G}%
\TextAt{8}{88.46}{327.065}{-modules}%
\TextAt{13}{139.583}{327.065}{A}%
\TextAt{8}{148.359}{327.065}{.}%
\TextAt{8}{160.762}{327.065}{For}%
\TextAt{8}{183.05}{327.065}{instance,}%
\TextAt{8}{233.535}{327.065}{the}%
\TextAt{8}{255.009}{327.065}{notation}%
\TextAt{13}{303.168}{327.065}{a}%
\TextAt{21}{309.313}{331.404}{\Character{27}}%
\TextAt{3}{314.438}{331.404}{\Character{0}}%
\TextAt{17}{321.025}{331.404}{1}%
\TextAt{9}{331.324}{327.065}{\Character{58}\Character{61}}%
\TextAt{13}{349.248}{327.065}{a}%
\TextAt{21}{355.393}{331.404}{\Character{27}}%
\TextAt{13}{361.017}{327.065}{a}%
\TextAt{3}{367.162}{331.404}{\Character{0}}%
\TextAt{17}{373.748}{331.404}{1}%
\TextAt{8}{383.696}{327.065}{is}%
\TextAt{8}{396.778}{327.065}{to}%
\TextAt{8}{412.407}{327.065}{be}%
\TextAt{8}{429.65}{327.065}{preferred}%
\TextAt{8}{480.713}{327.065}{to}%
\TextAt{8}{496.342}{327.065}{writing}%
\TextAt{9}{79.226}{310.019}{\Character{40}}%
\TextAt{13}{83.778}{310.019}{\Character{27}}%
\TextAt{16}{93.953}{310.019}{\Character{0}}%
\TextAt{9}{106.343}{310.019}{1\Character{41}}%
\TextAt{13}{116.748}{310.019}{a}%
\TextAt{9}{127.301}{310.019}{\Character{61}}%
\TextAt{13}{140.814}{310.019}{\Character{27}}%
\TextAt{13}{147.88}{310.019}{a}%
\TextAt{16}{157.133}{310.019}{\Character{1}}%
\TextAt{13}{163.546}{310.019}{a}%
\TextAt{3}{169.691}{314.358}{\Character{0}}%
\TextAt{17}{176.277}{314.358}{1}%
\TextAt{8}{181.009}{310.019}{.}%
\TextAt{8}{191.372}{310.019}{On}%
\TextAt{8}{211.521}{310.019}{the}%
\TextAt{8}{232.315}{310.019}{other}%
\TextAt{8}{263.523}{310.019}{hand,}%
\TextAt{8}{296.839}{310.019}{classical}%
\TextAt{8}{342.476}{310.019}{usage}%
\TextAt{8}{375.046}{310.019}{often}%
\TextAt{8}{405.274}{310.019}{calls}%
\TextAt{8}{431.986}{310.019}{for}%
\TextAt{8}{450.509}{310.019}{an}%
\TextAt{8}{467.406}{310.019}{operation}%
\TextAt{8}{520.408}{310.019}{on}%
\TextAt{8}{79.226}{292.973}{the}%
\TextAt{8}{99.83}{292.973}{left.}%
\TextAt{8}{126.203}{292.973}{Thus}%
\TextAt{8}{156.3}{292.973}{in}%
\TextAt{8}{170.4}{292.973}{the}%
\TextAt{8}{191.004}{292.973}{case}%
\TextAt{8}{216.223}{292.973}{of}%
\TextAt{8}{230.001}{292.973}{a}%
\TextAt{8}{240.204}{292.973}{Galois}%
\TextAt{8}{276.565}{292.973}{extension}%
\TextAt{13}{328.778}{292.973}{L}%
\TextAt{16}{340.826}{292.973}{\Character{106}}%
\TextAt{13}{348.231}{292.973}{K}%
\TextAt{8}{363.394}{292.973}{of}%
\TextAt{8}{377.171}{292.973}{actual}%
\TextAt{8}{412.743}{292.973}{fields,}%
\TextAt{8}{446.526}{292.973}{the}%
\TextAt{8}{467.13}{292.973}{Galois}%
\TextAt{8}{503.491}{292.973}{group}%
\TextAt{13}{79.226}{275.927}{G}%
\TextAt{9}{88.46}{275.927}{\Character{40}}%
\TextAt{13}{93.012}{275.927}{L}%
\TextAt{16}{105.061}{275.927}{\Character{106}}%
\TextAt{13}{112.467}{275.927}{K}%
\TextAt{9}{123.284}{275.927}{\Character{41}}%
\TextAt{8}{132.182}{275.927}{acts}%
\TextAt{8}{156.756}{275.927}{as}%
\TextAt{8}{171.574}{275.927}{the}%
\TextAt{8}{192.179}{275.927}{automorphism}%
\TextAt{8}{270.097}{275.927}{group}%
\TextAt{8}{303.721}{275.927}{on}%
\TextAt{13}{320.428}{275.927}{L}%
\TextAt{8}{332.738}{275.927}{from}%
\TextAt{8}{360.826}{275.927}{the}%
\TextAt{8}{381.431}{275.927}{left,}%
\TextAt{8}{405.722}{275.927}{and}%
\TextAt{8}{428.933}{275.927}{therefore}%
\TextAt{8}{478.481}{275.927}{also}%
\TextAt{8}{502.409}{275.927}{in}%
\TextAt{8}{516.51}{275.927}{the}%
\TextAt{8}{79.226}{258.881}{same}%
\TextAt{8}{107.87}{258.881}{way}%
\TextAt{8}{130.931}{258.881}{on}%
\TextAt{8}{146.508}{258.881}{the}%
\TextAt{8}{165.982}{258.881}{multiplicative}%
\TextAt{8}{239.124}{258.881}{group}%
\TextAt{13}{271.617}{258.881}{L}%
\TextAt{3}{279.582}{263.22}{\Character{3}}%
\TextAt{8}{284.314}{258.881}{.}%
\TextAt{8}{292.535}{258.881}{This}%
\TextAt{8}{318.573}{258.881}{occasional}%
\TextAt{8}{373.566}{258.881}{switch}%
\TextAt{8}{409.036}{258.881}{from}%
\TextAt{8}{435.995}{258.881}{the}%
\TextAt{8}{455.469}{258.881}{left}%
\TextAt{8}{475.266}{258.881}{to}%
\TextAt{8}{488.895}{258.881}{the}%
\TextAt{8}{508.369}{258.881}{right}%
\TextAt{8}{79.226}{241.835}{should}%
\TextAt{8}{116.359}{241.835}{not}%
\TextAt{8}{137.173}{241.835}{confuse}%
\TextAt{8}{178.526}{241.835}{the}%
\TextAt{8}{198.682}{241.835}{reader.}%
\TextAt{8}{97.159}{219.254}{For}%
\TextAt{8}{118.199}{219.254}{every}%
\TextAt{8}{149.162}{219.254}{extension}%
\TextAt{13}{200.999}{219.254}{L}%
\TextAt{16}{212.405}{219.254}{\Character{106}}%
\TextAt{13}{219.167}{219.254}{K}%
\TextAt{8}{233.953}{219.254}{we}%
\TextAt{8}{251.251}{219.254}{have}%
\TextAt{13}{278.317}{219.254}{A}%
\TextAt{21}{287.092}{217.461}{K}%
\NativeMathAt{298.736}{219.254}{11.955}{\subseteq}%
\TextAt{13}{311.476}{219.254}{A}%
\TextAt{21}{320.252}{217.461}{L}%
\TextAt{8}{326.509}{219.254}{,}%
\TextAt{8}{333.747}{219.254}{and}%
\TextAt{8}{356.58}{219.254}{if}%
\TextAt{13}{367.375}{219.254}{L}%
\TextAt{16}{378.781}{219.254}{\Character{106}}%
\TextAt{13}{385.543}{219.254}{K}%
\TextAt{8}{400.329}{219.254}{is}%
\TextAt{8}{412.164}{219.254}{finite,}%
\TextAt{8}{445.416}{219.254}{then}%
\TextAt{8}{472.147}{219.254}{we}%
\TextAt{8}{489.445}{219.254}{have}%
\TextAt{8}{516.51}{219.254}{the}%
\TextAt{8}{79.226}{202.208}{norm}%
\TextAt{8}{109.796}{202.208}{map}%
\TextAt{13}{200.187}{185.162}{N}%
\TextAt{21}{209.561}{183.303}{L}%
\TextAt{3}{215.32}{183.303}{\Character{106}}%
\TextAt{21}{217.672}{183.303}{K}%
\TextAt{9}{229.196}{185.162}{\Character{58}}%
\TextAt{13}{235.768}{185.162}{A}%
\TextAt{21}{244.544}{183.369}{L}%
\NativeArrowAt{254.122}{185.162}{11.955}{19.26}{\longrightarrow}%
\TextAt{13}{276.704}{185.162}{A}%
\TextAt{21}{285.479}{183.369}{K}%
\TextAt{13}{293.682}{185.162}{\Character{59}}%
\TextAt{13}{322.817}{185.162}{N}%
\TextAt{21}{332.21}{183.303}{L}%
\TextAt{3}{337.969}{183.303}{\Character{106}}%
\TextAt{21}{340.321}{183.303}{K}%
\TextAt{9}{348.524}{185.162}{\Character{40}}%
\TextAt{13}{353.076}{185.162}{a}%
\TextAt{9}{359.221}{185.162}{\Character{41}}%
\TextAt{9}{367.076}{185.162}{\Character{61}}%
\TextAt{20}{379.52}{196.52}{\Character{89}}%
\TextAt{21}{384.595}{171.414}{\Character{27}}%
\TextAt{13}{396.788}{185.162}{a}%
\TextAt{21}{402.933}{190.099}{\Character{27}}%
\TextAt{13}{408.557}{185.162}{\Character{59}}%
\TextAt{8}{79.226}{155.108}{where}%
\TextAt{13}{113.015}{155.108}{\Character{27}}%
\TextAt{8}{123.974}{155.108}{varies}%
\TextAt{8}{156.866}{155.108}{over}%
\TextAt{8}{181.891}{155.108}{a}%
\TextAt{8}{191.626}{155.108}{system}%
\TextAt{8}{230.424}{155.108}{of}%
\TextAt{8}{243.733}{155.108}{representatives}%
\TextAt{8}{323.514}{155.108}{of}%
\TextAt{13}{336.823}{155.108}{G}%
\TextAt{21}{346.057}{153.315}{L}%
\TextAt{16}{352.314}{155.108}{\Character{110}}%
\TextAt{13}{358.291}{155.108}{G}%
\TextAt{21}{367.525}{153.315}{K}%
\TextAt{8}{375.728}{155.108}{.}%
\TextAt{8}{384.169}{155.108}{If}%
\TextAt{13}{395.841}{155.108}{L}%
\TextAt{16}{407.126}{155.108}{\Character{106}}%
\TextAt{13}{413.768}{155.108}{K}%
\TextAt{8}{428.462}{155.108}{is}%
\TextAt{8}{440.206}{155.108}{Galois,}%
\TextAt{8}{479.354}{155.108}{then}%
\TextAt{13}{505.994}{155.108}{A}%
\TextAt{21}{514.769}{153.315}{L}%
\TextAt{8}{524.903}{155.108}{is}%
\TextAt{8}{79.226}{138.062}{a}%
\TextAt{13}{88.982}{138.062}{G}%
\TextAt{9}{98.215}{138.062}{\Character{40}}%
\TextAt{13}{102.768}{138.062}{L}%
\TextAt{16}{114.053}{138.062}{\Character{106}}%
\TextAt{13}{120.695}{138.062}{K}%
\TextAt{9}{131.512}{138.062}{\Character{41}}%
\TextAt{8}{136.065}{138.062}{-module}%
\TextAt{8}{181.255}{138.062}{and}%
\TextAt{8}{204.018}{138.062}{one}%
\TextAt{8}{225.477}{138.062}{has}%
\TextAt{13}{268.873}{121.016}{A}%
\TextAt{21}{277.648}{127.127}{G}%
\TextAt{17}{284.247}{127.127}{\Character{40}}%
\TextAt{21}{287.54}{127.127}{L}%
\TextAt{3}{293.299}{127.127}{\Character{106}}%
\TextAt{21}{295.651}{127.127}{K}%
\TextAt{17}{303.356}{127.127}{\Character{41}}%
\TextAt{21}{277.648}{117.775}{L}%
\TextAt{9}{310.468}{121.016}{\Character{61}}%
\TextAt{13}{322.893}{121.016}{A}%
\TextAt{21}{331.669}{119.223}{K}%
\TextAt{13}{339.871}{121.016}{\Character{58}}%
\TextAt{8}{79.226}{98.823}{At}%
\TextAt{8}{96.131}{98.823}{the}%
\TextAt{8}{116.287}{98.823}{center}%
\TextAt{8}{151.077}{98.823}{of}%
\TextAt{8}{164.407}{98.823}{class}%
\TextAt{8}{191.844}{98.823}{field}%
\TextAt{8}{217.201}{98.823}{theory}%
\TextAt{8}{253.951}{98.823}{there}%
\TextAt{8}{283.863}{98.823}{is}%
\TextAt{8}{295.627}{98.823}{the}%
\TextAt{12}{315.783}{98.823}{norm}%
\TextAt{12}{351.182}{98.823}{residue}%
\TextAt{12}{397.425}{98.823}{group}%
\TextAt{13}{218.814}{72.921}{H}%
\TextAt{17}{229.436}{77.858}{0}%
\TextAt{9}{234.168}{72.921}{\Character{40}}%
\TextAt{13}{238.721}{72.921}{G}%
\TextAt{9}{247.954}{72.921}{\Character{40}}%
\TextAt{13}{252.507}{72.921}{L}%
\TextAt{16}{263.792}{72.921}{\Character{106}}%
\TextAt{13}{270.434}{72.921}{K}%
\TextAt{9}{281.251}{72.921}{\Character{41}}%
\TextAt{13}{285.804}{72.921}{\Character{59}}%
\TextAt{13}{291.028}{72.921}{A}%
\TextAt{21}{299.823}{71.128}{L}%
\TextAt{9}{306.08}{72.921}{\Character{41}}%
\TextAt{9}{313.935}{72.921}{\Character{61}}%
\TextAt{13}{326.379}{72.921}{A}%
\TextAt{21}{335.154}{71.128}{K}%
\TextAt{9}{343.357}{72.921}{\Character{47}}%
\TextAt{13}{349.21}{72.921}{N}%
\TextAt{21}{358.584}{71.062}{L}%
\TextAt{3}{364.343}{71.062}{\Character{106}}%
\TextAt{21}{366.695}{71.062}{K}%
\TextAt{13}{374.898}{72.921}{A}%
\TextAt{21}{383.673}{71.128}{L}%
\TextAt{13}{389.931}{72.921}{\Character{58}}%
\TextAt{8}{297.211}{43.033}{253}%
\Destination{79.23}{641.305}{nav.Item.84}%
\Destination{79.23}{623.705}{nav.Item.85}%
\Destination{79.23}{606.655}{nav.Item.86}%
\Destination{79.23}{589.605}{nav.Item.87}%
\Destination{78.23}{809.815}{nav.page.253}%
\Destination{79.36}{376.155}{nav.print.277}%
}

% Page 266 (body page 254)
\Page{%
\TextAt{2}{79.226}{796.943}{IV.}%
\TextAt{2}{98.012}{796.943}{Abstract}%
\TextAt{2}{143.252}{796.943}{Class}%
\TextAt{2}{171.844}{796.943}{Field}%
\TextAt{2}{199.575}{796.943}{Theory}%
\Vector{q 0 G 0 g 0.299 w 7.226 18.729 m 460.769 18.729 l S Q}%
\TextAt{8}{79.226}{761.124}{We}%
\TextAt{8}{99.371}{761.124}{also}%
\TextAt{8}{122.851}{761.124}{consider}%
\TextAt{8}{168.436}{761.124}{the}%
\TextAt{8}{188.592}{761.124}{group}%
\TextAt{13}{201.211}{735.377}{H}%
\TextAt{3}{211.833}{740.314}{\Character{0}}%
\TextAt{17}{218.42}{740.314}{1}%
\TextAt{9}{223.152}{735.377}{\Character{40}}%
\TextAt{13}{227.704}{735.377}{G}%
\TextAt{9}{236.938}{735.377}{\Character{40}}%
\TextAt{13}{241.49}{735.377}{L}%
\TextAt{16}{252.776}{735.377}{\Character{106}}%
\TextAt{13}{259.417}{735.377}{K}%
\TextAt{9}{270.235}{735.377}{\Character{41}}%
\TextAt{13}{274.787}{735.377}{\Character{59}}%
\TextAt{13}{280.011}{735.377}{A}%
\TextAt{21}{288.807}{733.584}{L}%
\TextAt{9}{295.064}{735.377}{\Character{41}}%
\TextAt{9}{302.919}{735.377}{\Character{61}}%
\TextAt{21}{315.363}{733.584}{N}%
\TextAt{22}{322.084}{731.846}{L}%
\TextAt{23}{327.143}{731.846}{\Character{106}}%
\TextAt{22}{329.412}{731.846}{K}%
\TextAt{13}{336.953}{735.377}{A}%
\TextAt{21}{345.729}{733.584}{L}%
\TextAt{9}{351.986}{735.377}{\Character{47}}%
\TextAt{13}{357.839}{735.377}{I}%
\TextAt{21}{363.002}{733.518}{G}%
\TextAt{17}{369.6}{733.518}{\Character{40}}%
\TextAt{21}{372.894}{733.518}{L}%
\TextAt{3}{378.653}{733.518}{\Character{106}}%
\TextAt{21}{381.005}{733.518}{K}%
\TextAt{17}{388.71}{733.518}{\Character{41}}%
\TextAt{13}{392.501}{735.377}{A}%
\TextAt{21}{401.276}{733.584}{L}%
\TextAt{13}{407.533}{735.377}{\Character{59}}%
\TextAt{8}{79.226}{709.631}{where}%
\TextAt{21}{219.691}{690.792}{N}%
\TextAt{22}{226.413}{689.054}{L}%
\TextAt{23}{231.471}{689.054}{\Character{106}}%
\TextAt{22}{233.741}{689.054}{K}%
\TextAt{13}{241.282}{692.585}{A}%
\TextAt{21}{250.057}{690.792}{L}%
\TextAt{9}{259.635}{692.585}{\Character{61}}%
\TextAt{16}{272.061}{692.585}{\Character{102}}%
\TextAt{13}{278.038}{692.585}{a}%
\NativeMathAt{287.504}{692.585}{11.955}{\in}%
\TextAt{13}{298.795}{692.585}{A}%
\TextAt{21}{307.57}{690.792}{L}%
\TextAt{16}{317.148}{692.585}{\Character{106}}%
\TextAt{13}{323.79}{692.585}{N}%
\TextAt{21}{333.164}{690.725}{L}%
\TextAt{3}{338.923}{690.725}{\Character{106}}%
\TextAt{21}{341.275}{690.725}{K}%
\TextAt{9}{349.478}{692.585}{\Character{40}}%
\TextAt{13}{354.03}{692.585}{a}%
\TextAt{9}{360.175}{692.585}{\Character{41}}%
\TextAt{9}{368.03}{692.585}{\Character{61}}%
\TextAt{9}{380.451}{692.585}{1}%
\TextAt{16}{386.327}{692.585}{\Character{103}}%
\TextAt{8}{79.226}{670.458}{is}%
\TextAt{8}{91.098}{670.458}{the}%
\TextAt{8}{111.361}{670.458}{“norm-one}%
\TextAt{8}{169.021}{670.458}{group”}%
\TextAt{8}{207.828}{670.458}{and}%
\TextAt{13}{230.698}{670.458}{I}%
\TextAt{21}{235.861}{668.599}{G}%
\TextAt{17}{242.459}{668.599}{\Character{40}}%
\TextAt{21}{245.753}{668.599}{L}%
\TextAt{3}{251.512}{668.599}{\Character{106}}%
\TextAt{21}{253.864}{668.599}{K}%
\TextAt{17}{261.568}{668.599}{\Character{41}}%
\TextAt{13}{265.36}{670.458}{A}%
\TextAt{21}{274.135}{668.665}{L}%
\TextAt{8}{284.397}{670.458}{is}%
\TextAt{8}{296.269}{670.458}{the}%
\TextAt{8}{316.533}{670.458}{subgroup}%
\TextAt{8}{367.438}{670.458}{of}%
\TextAt{21}{380.875}{668.665}{N}%
\TextAt{22}{387.596}{666.927}{L}%
\TextAt{23}{392.655}{666.927}{\Character{106}}%
\TextAt{22}{394.925}{666.927}{K}%
\TextAt{13}{402.466}{670.458}{A}%
\TextAt{21}{411.241}{668.665}{L}%
\TextAt{8}{421.503}{670.458}{which}%
\TextAt{8}{455.097}{670.458}{is}%
\TextAt{8}{466.968}{670.458}{generated}%
\TextAt{8}{520.408}{670.458}{by}%
\TextAt{8}{79.226}{653.412}{all}%
\TextAt{8}{95.485}{653.412}{elements}%
\TextAt{13}{268.962}{636.366}{a}%
\TextAt{21}{275.107}{641.303}{\Character{27}}%
\TextAt{3}{280.233}{641.303}{\Character{0}}%
\TextAt{17}{286.819}{641.303}{1}%
\TextAt{9}{294.872}{636.366}{\Character{58}\Character{61}}%
\TextAt{13}{310.549}{636.366}{a}%
\TextAt{21}{316.694}{641.303}{\Character{27}}%
\TextAt{13}{322.318}{636.366}{a}%
\TextAt{3}{328.463}{641.303}{\Character{0}}%
\TextAt{17}{335.049}{641.303}{1}%
\TextAt{13}{339.782}{636.366}{\Character{59}}%
\TextAt{8}{79.226}{614.24}{with}%
\TextAt{13}{105.146}{614.24}{a}%
\NativeMathAt{114.612}{614.24}{11.955}{\in}%
\TextAt{13}{125.903}{614.24}{A}%
\TextAt{21}{134.678}{612.446}{L}%
\TextAt{8}{140.935}{614.24}{,}%
\TextAt{8}{147.492}{614.24}{and}%
\TextAt{13}{169.515}{614.24}{\Character{27}}%
\NativeMathAt{179.918}{614.24}{11.955}{\in}%
\TextAt{13}{191.209}{614.24}{G}%
\TextAt{9}{200.443}{614.24}{\Character{40}}%
\TextAt{13}{204.995}{614.24}{L}%
\TextAt{16}{216.28}{614.24}{\Character{106}}%
\TextAt{13}{222.922}{614.24}{K}%
\TextAt{9}{233.74}{614.24}{\Character{41}}%
\TextAt{8}{238.292}{614.24}{.}%
\TextAt{8}{246.494}{614.24}{If}%
\TextAt{13}{257.446}{614.24}{G}%
\TextAt{9}{266.679}{614.24}{\Character{40}}%
\TextAt{13}{271.232}{614.24}{L}%
\TextAt{16}{282.517}{614.24}{\Character{106}}%
\TextAt{13}{289.159}{614.24}{K}%
\TextAt{9}{299.977}{614.24}{\Character{41}}%
\TextAt{8}{307.686}{614.24}{is}%
\TextAt{8}{318.71}{614.24}{cyclic}%
\TextAt{8}{350.153}{614.24}{and}%
\TextAt{13}{372.176}{614.24}{\Character{27}}%
\TextAt{8}{382.415}{614.24}{is}%
\TextAt{8}{393.439}{614.24}{a}%
\TextAt{8}{402.454}{614.24}{generator,}%
\TextAt{8}{457.155}{614.24}{then}%
\TextAt{13}{483.075}{614.24}{I}%
\TextAt{21}{488.238}{612.38}{G}%
\TextAt{17}{494.836}{612.38}{\Character{40}}%
\TextAt{21}{498.13}{612.38}{L}%
\TextAt{3}{503.889}{612.38}{\Character{106}}%
\TextAt{21}{506.241}{612.38}{K}%
\TextAt{17}{513.946}{612.38}{\Character{41}}%
\TextAt{13}{517.737}{614.24}{A}%
\TextAt{21}{526.512}{612.446}{L}%
\TextAt{8}{79.226}{597.194}{is}%
\TextAt{8}{90.99}{597.194}{simply}%
\TextAt{8}{128.446}{597.194}{the}%
\TextAt{8}{148.602}{597.194}{group}%
\TextAt{13}{243.741}{580.148}{A}%
\TextAt{21}{252.516}{585.096}{\Character{27}}%
\TextAt{3}{257.642}{585.096}{\Character{0}}%
\TextAt{17}{264.228}{585.096}{1}%
\TextAt{21}{252.516}{576.907}{L}%
\TextAt{9}{272.281}{580.148}{\Character{61}}%
\TextAt{16}{284.707}{580.148}{\Character{102}}%
\TextAt{13}{290.684}{580.148}{a}%
\TextAt{21}{296.829}{585.084}{\Character{27}}%
\TextAt{3}{301.955}{585.084}{\Character{0}}%
\TextAt{17}{308.541}{585.084}{1}%
\TextAt{16}{316.594}{580.148}{\Character{106}}%
\TextAt{13}{323.236}{580.148}{a}%
\NativeMathAt{332.702}{580.148}{11.955}{\in}%
\TextAt{13}{343.993}{580.148}{A}%
\TextAt{21}{352.768}{578.355}{L}%
\TextAt{16}{359.025}{580.148}{\Character{103}}%
\TextAt{13}{365.003}{580.148}{\Character{58}}%
\TextAt{8}{79.226}{558.021}{In}%
\TextAt{8}{93.51}{558.021}{fact,}%
\TextAt{8}{119.578}{558.021}{the}%
\TextAt{8}{139.397}{558.021}{formal}%
\TextAt{8}{175.811}{558.021}{identity}%
\TextAt{13}{218.728}{558.021}{\Character{27}}%
\TextAt{21}{225.81}{562.359}{k}%
\TextAt{16}{232.897}{558.021}{\Character{0}}%
\TextAt{9}{244.163}{558.021}{1}%
\TextAt{9}{253.321}{558.021}{\Character{61}}%
\TextAt{9}{265.742}{558.021}{\Character{40}1}%
\TextAt{9}{278.104}{558.021}{\Character{43}}%
\TextAt{13}{289.208}{558.021}{\Character{27}}%
\TextAt{9}{298.258}{558.021}{\Character{43}}%
\TextAt{16}{309.33}{558.021}{\Character{1}}%
\TextAt{16}{314.626}{558.021}{\Character{1}}%
\TextAt{16}{319.934}{558.021}{\Character{1}}%
\TextAt{9}{325.245}{558.021}{\Character{43}}%
\TextAt{13}{336.317}{558.021}{\Character{27}}%
\TextAt{21}{343.4}{562.359}{k}%
\TextAt{3}{348.021}{562.36}{\Character{0}}%
\TextAt{17}{354.608}{562.359}{1}%
\TextAt{9}{359.34}{558.021}{\Character{41}\Character{40}}%
\TextAt{13}{368.445}{558.021}{\Character{27}}%
\TextAt{16}{377.495}{558.021}{\Character{0}}%
\TextAt{9}{388.761}{558.021}{1\Character{41}}%
\TextAt{8}{402.727}{558.021}{implies}%
\TextAt{13}{442.117}{558.021}{a}%
\TextAt{21}{448.262}{562.359}{\Character{27}}%
\TextAt{22}{453.387}{565.172}{k}%
\TextAt{3}{457.96}{562.36}{\Character{0}}%
\TextAt{17}{464.547}{562.359}{1}%
\TextAt{9}{472.6}{558.021}{\Character{61}}%
\TextAt{13}{485.025}{558.021}{b}%
\TextAt{21}{490.002}{562.359}{\Character{27}}%
\TextAt{3}{495.128}{562.36}{\Character{0}}%
\TextAt{17}{501.714}{562.359}{1}%
\TextAt{8}{510.007}{558.021}{with}%
\TextAt{13}{79.226}{540.975}{b}%
\TextAt{9}{87.524}{540.975}{\Character{61}}%
\TextAt{20}{99.95}{549.942}{\Character{81}}%
\TextAt{21}{111.241}{546.787}{k}%
\TextAt{3}{115.862}{546.787}{\Character{0}}%
\TextAt{17}{122.449}{546.787}{1}%
\TextAt{21}{111.241}{537.488}{i}%
\TextAt{17}{114.124}{537.488}{\Character{61}0}%
\TextAt{13}{129.174}{540.975}{a}%
\TextAt{21}{135.318}{545.314}{\Character{27}}%
\TextAt{22}{140.444}{548.126}{i}%
\TextAt{8}{144.104}{540.975}{.}%
\TextAt{8}{97.159}{520.94}{Let}%
\TextAt{8}{118.052}{520.94}{us}%
\TextAt{8}{132.992}{520.94}{now}%
\TextAt{8}{157.304}{520.94}{apply}%
\TextAt{8}{189.424}{520.94}{the}%
\TextAt{8}{209.504}{520.94}{notions}%
\TextAt{8}{250.47}{520.94}{introduced}%
\TextAt{8}{308.927}{520.94}{so}%
\TextAt{8}{323.221}{520.94}{far}%
\TextAt{8}{341.03}{520.94}{to}%
\TextAt{8}{355.264}{520.94}{the}%
\TextAt{8}{375.344}{520.94}{example}%
\TextAt{8}{421.116}{520.94}{of}%
\TextAt{12}{434.884}{520.94}{Kummer}%
\TextAt{12}{491.345}{520.94}{theory}%
\TextAt{8}{529.518}{520.94}{.}%
\TextAt{8}{79.226}{503.894}{For}%
\TextAt{8}{100.196}{503.894}{this,}%
\TextAt{8}{126.27}{503.894}{we}%
\TextAt{8}{143.497}{503.894}{impose}%
\TextAt{8}{182.901}{503.894}{on}%
\TextAt{8}{199.16}{503.894}{the}%
\TextAt{13}{219.317}{503.894}{G}%
\TextAt{8}{228.55}{503.894}{-module}%
\TextAt{13}{273.741}{503.894}{A}%
\TextAt{8}{286.414}{503.894}{the}%
\TextAt{8}{306.57}{503.894}{following}%
\TextAt{8}{356.005}{503.894}{axiomatic}%
\TextAt{8}{409.672}{503.894}{condition:}%
\TextAt{12}{79.226}{478.812}{Axiom}%
\TextAt{12}{122.659}{478.812}{(3.1).}%
\TextAt{8}{157.975}{478.812}{One}%
\TextAt{8}{182.686}{478.812}{has}%
\TextAt{13}{203.56}{478.812}{H}%
\TextAt{3}{214.183}{483.15}{\Character{0}}%
\TextAt{17}{220.769}{483.15}{1}%
\TextAt{9}{225.501}{478.812}{\Character{40}}%
\TextAt{13}{230.054}{478.812}{G}%
\TextAt{9}{239.287}{478.812}{\Character{40}}%
\TextAt{13}{243.84}{478.812}{L}%
\TextAt{16}{255.125}{478.812}{\Character{106}}%
\TextAt{13}{261.767}{478.812}{K}%
\TextAt{9}{272.584}{478.812}{\Character{41}}%
\TextAt{13}{277.137}{478.812}{\Character{59}}%
\TextAt{13}{282.361}{478.812}{A}%
\TextAt{21}{291.156}{477.019}{L}%
\TextAt{9}{297.413}{478.812}{\Character{41}}%
\TextAt{9}{305.268}{478.812}{\Character{61}}%
\TextAt{9}{317.689}{478.812}{1}%
\TextAt{8}{327.462}{478.812}{for}%
\TextAt{8}{345.347}{478.812}{all}%
\TextAt{8}{361.606}{478.812}{finite}%
\TextAt{8}{391.518}{478.812}{cyclic}%
\TextAt{8}{423.701}{478.812}{extensions}%
\TextAt{13}{480.082}{478.812}{L}%
\TextAt{16}{491.367}{478.812}{\Character{106}}%
\TextAt{13}{498.009}{478.812}{K}%
\TextAt{8}{508.827}{478.812}{.}%
\TextAt{8}{97.159}{453.73}{The}%
\TextAt{8}{120.824}{453.73}{theory}%
\TextAt{8}{157.185}{453.73}{we}%
\TextAt{8}{174.024}{453.73}{are}%
\TextAt{8}{193.146}{453.73}{about}%
\TextAt{8}{226.255}{453.73}{to}%
\TextAt{8}{240.176}{453.73}{develop}%
\TextAt{8}{282.061}{453.73}{makes}%
\TextAt{8}{316.856}{453.73}{reference}%
\TextAt{8}{365.555}{453.73}{to}%
\TextAt{8}{379.477}{453.73}{a}%
\TextAt{8}{388.843}{453.73}{surjective}%
\TextAt{13}{440.866}{453.73}{G}%
\TextAt{8}{450.099}{453.73}{-homomorphism}%
\TextAt{13}{237.366}{427.983}{\Character{125}}%
\TextAt{9}{248.139}{427.983}{\Character{58}}%
\TextAt{13}{254.712}{427.983}{A}%
\NativeArrowAt{266.808}{427.983}{11.955}{19.26}{\longrightarrow}%
\TextAt{13}{289.39}{427.983}{A\Character{59}}%
\TextAt{13}{327.3}{427.983}{a}%
\NativeArrowAt{336.786}{427.983}{11.955}{19.26}{\longmapsto}%
\TextAt{13}{359.368}{427.983}{a}%
\TextAt{21}{365.513}{432.92}{\Character{125}}%
\TextAt{13}{371.378}{427.983}{\Character{59}}%
\TextAt{8}{79.226}{402.237}{with}%
\TextAt{8}{106.546}{402.237}{finite}%
\TextAt{8}{137.118}{402.237}{cyclic}%
\TextAt{8}{169.961}{402.237}{kernel}%
\TextAt{13}{205.088}{402.237}{\Character{22}}%
\TextAt{21}{212.131}{400.444}{\Character{125}}%
\TextAt{8}{217.996}{402.237}{.}%
\TextAt{8}{228.424}{402.237}{The}%
\TextAt{8}{253.138}{402.237}{order}%
\TextAt{13}{284.367}{402.237}{n}%
\TextAt{9}{295.8}{402.237}{\Character{61}}%
\TextAt{8}{309.35}{402.237}{\char35\relax }%
\TextAt{13}{319.105}{402.237}{\Character{22}}%
\TextAt{21}{326.148}{400.444}{\Character{125}}%
\TextAt{8}{336.571}{402.237}{is}%
\TextAt{8}{348.995}{402.237}{called}%
\TextAt{8}{382.818}{402.237}{the}%
\TextAt{8}{403.635}{402.237}{exponent}%
\TextAt{8}{454.698}{402.237}{of}%
\TextAt{8}{468.687}{402.237}{the}%
\TextAt{8}{489.504}{402.237}{operator}%
\TextAt{13}{79.226}{385.191}{\Character{125}}%
\TextAt{8}{86.679}{385.191}{.}%
\TextAt{8}{100.007}{385.191}{The}%
\TextAt{8}{125.688}{385.191}{case}%
\TextAt{8}{152.086}{385.191}{of}%
\TextAt{8}{167.042}{385.191}{prime}%
\TextAt{8}{201.833}{385.191}{interest}%
\TextAt{8}{245.47}{385.191}{to}%
\TextAt{8}{261.407}{385.191}{us}%
\TextAt{8}{278.049}{385.191}{is}%
\TextAt{8}{291.44}{385.191}{when}%
\TextAt{13}{323.624}{385.191}{\Character{125}}%
\TextAt{8}{336.601}{385.191}{is}%
\TextAt{8}{349.991}{385.191}{the}%
\TextAt{13}{371.775}{385.191}{n}%
\TextAt{8}{378.762}{385.191}{-th}%
\TextAt{8}{399.242}{385.191}{power}%
\TextAt{8}{435.013}{385.191}{map}%
\TextAt{13}{462.654}{385.191}{a}%
\NativeArrowAt{474.892}{385.191}{11.955}{11.955}{\mapsto}%
\TextAt{13}{492.94}{385.191}{a}%
\TextAt{21}{499.085}{389.53}{n}%
\TextAt{8}{504.722}{385.191}{,}%
\TextAt{8}{513.904}{385.191}{and}%
\TextAt{13}{79.226}{368.145}{\Character{22}}%
\TextAt{21}{86.269}{366.352}{\Character{125}}%
\TextAt{9}{95.456}{368.145}{\Character{61}}%
\TextAt{13}{107.881}{368.145}{\Character{22}}%
\TextAt{21}{114.924}{366.352}{n}%
\TextAt{9}{123.881}{368.145}{\Character{61}}%
\TextAt{16}{136.307}{368.145}{\Character{102}}%
\TextAt{13}{142.284}{368.145}{\Character{24}}%
\NativeMathAt{151.296}{368.145}{11.955}{\in}%
\TextAt{13}{162.586}{368.145}{A}%
\TextAt{16}{174.683}{368.145}{\Character{106}}%
\TextAt{13}{181.324}{368.145}{\Character{24}}%
\TextAt{21}{187.015}{372.484}{n}%
\TextAt{9}{195.972}{368.145}{\Character{61}}%
\TextAt{9}{208.382}{368.145}{1}%
\TextAt{16}{214.25}{368.145}{\Character{103}}%
\TextAt{8}{224.125}{368.145}{is}%
\TextAt{8}{235.889}{368.145}{the}%
\TextAt{8}{256.046}{368.145}{group}%
\TextAt{8}{289.221}{368.145}{of}%
\TextAt{8}{302.551}{368.145}{“}%
\TextAt{13}{308.075}{368.145}{n}%
\TextAt{8}{315.062}{368.145}{-th}%
\TextAt{8}{333.916}{368.145}{roots}%
\TextAt{8}{363.576}{368.145}{of}%
\TextAt{8}{376.906}{368.145}{unity”}%
\TextAt{8}{412.999}{368.145}{in}%
\TextAt{13}{426.652}{368.145}{A}%
\TextAt{8}{435.427}{368.145}{.}%
\TextAt{8}{97.159}{348.11}{We}%
\TextAt{8}{117.531}{348.11}{now}%
\TextAt{8}{142.147}{348.11}{fix}%
\TextAt{8}{158.956}{348.11}{a}%
\TextAt{8}{168.939}{348.11}{field}%
\TextAt{13}{194.523}{348.11}{K}%
\TextAt{8}{209.466}{348.11}{such}%
\TextAt{8}{236.09}{348.11}{that}%
\TextAt{13}{261.687}{348.11}{\Character{22}}%
\TextAt{21}{268.73}{346.317}{\Character{125}}%
\NativeMathAt{278.304}{348.11}{11.955}{\subseteq}%
\TextAt{13}{291.31}{348.11}{A}%
\TextAt{21}{300.086}{346.317}{K}%
\TextAt{8}{308.288}{348.11}{.}%
\TextAt{8}{317.419}{348.11}{For}%
\TextAt{8}{338.616}{348.11}{every}%
\TextAt{8}{369.736}{348.11}{subset}%
\TextAt{13}{405.852}{348.11}{B}%
\NativeMathAt{419.057}{348.11}{11.955}{\subseteq}%
\TextAt{13}{432.064}{348.11}{A}%
\TextAt{8}{440.839}{348.11}{,}%
\TextAt{8}{448.273}{348.11}{let}%
\TextAt{13}{465.405}{348.11}{K}%
\TextAt{9}{476.223}{348.11}{\Character{40}}%
\TextAt{13}{480.775}{348.11}{B}%
\TextAt{9}{490.271}{348.11}{\Character{41}}%
\TextAt{8}{498.948}{348.11}{denote}%
\TextAt{8}{79.226}{331.064}{the}%
\TextAt{8}{99.383}{331.064}{fixed}%
\TextAt{8}{127.668}{331.064}{field}%
\TextAt{8}{153.025}{331.064}{of}%
\TextAt{8}{166.355}{331.064}{the}%
\TextAt{8}{186.512}{331.064}{closed}%
\TextAt{8}{221.038}{331.064}{subgroup}%
\TextAt{13}{213.064}{305.318}{H}%
\TextAt{9}{227.007}{305.318}{\Character{61}}%
\TextAt{16}{239.433}{305.318}{\Character{102}}%
\TextAt{13}{245.411}{305.318}{\Character{27}}%
\NativeMathAt{255.814}{305.318}{11.955}{\in}%
\TextAt{13}{267.105}{305.318}{G}%
\TextAt{21}{276.338}{303.525}{K}%
\TextAt{16}{287.862}{305.318}{\Character{106}}%
\TextAt{13}{294.504}{305.318}{b}%
\TextAt{21}{299.481}{310.254}{\Character{27}}%
\TextAt{9}{308.425}{305.318}{\Character{61}}%
\TextAt{13}{320.85}{305.318}{b}%
\TextAt{8}{329.725}{305.318}{for}%
\TextAt{8}{347.61}{305.318}{all}%
\TextAt{13}{363.869}{305.318}{b}%
\NativeMathAt{372.167}{305.318}{11.955}{\in}%
\TextAt{13}{383.458}{305.318}{B}%
\TextAt{16}{392.954}{305.318}{\Character{103}}%
\TextAt{8}{79.226}{279.572}{of}%
\TextAt{13}{92.079}{279.572}{G}%
\TextAt{21}{101.313}{277.779}{K}%
\TextAt{8}{109.515}{279.572}{.}%
\TextAt{8}{117.805}{279.572}{If}%
\TextAt{13}{129.02}{279.572}{B}%
\TextAt{8}{141.937}{279.572}{is}%
\TextAt{13}{153.224}{279.572}{G}%
\TextAt{21}{162.457}{277.779}{K}%
\TextAt{8}{170.66}{279.572}{-invariant,}%
\TextAt{8}{226.551}{279.572}{then}%
\TextAt{13}{252.734}{279.572}{K}%
\TextAt{9}{263.552}{279.572}{\Character{40}}%
\TextAt{13}{268.104}{279.572}{B}%
\TextAt{9}{277.601}{279.572}{\Character{41}}%
\TextAt{16}{285.474}{279.572}{\Character{106}}%
\TextAt{13}{292.116}{279.572}{K}%
\TextAt{8}{306.354}{279.572}{is}%
\TextAt{8}{317.641}{279.572}{obviously}%
\TextAt{8}{368.941}{279.572}{Galois.}%
\TextAt{8}{409.247}{279.572}{A}%
\TextAt{12}{421.442}{279.572}{Kummer}%
\TextAt{12}{477.441}{279.572}{extension}%
\TextAt{8}{79.226}{262.526}{(with}%
\TextAt{8}{110.441}{262.526}{respect}%
\TextAt{8}{150.491}{262.526}{to}%
\TextAt{13}{164.801}{262.526}{\Character{125}}%
\TextAt{8}{172.254}{262.526}{)}%
\TextAt{8}{180.706}{262.526}{is}%
\TextAt{8}{192.47}{262.526}{by}%
\TextAt{8}{208.729}{262.526}{definition}%
\TextAt{8}{260.758}{262.526}{an}%
\TextAt{8}{277.017}{262.526}{extension}%
\TextAt{8}{328.783}{262.526}{of}%
\TextAt{8}{342.113}{262.526}{the}%
\TextAt{8}{362.269}{262.526}{form}%
\TextAt{13}{265.53}{236.78}{K}%
\TextAt{9}{276.348}{236.78}{\Character{40}}%
\TextAt{13}{280.9}{236.78}{\Character{125}}%
\TextAt{3}{288.353}{241.716}{\Character{0}}%
\TextAt{17}{294.939}{241.716}{1}%
\TextAt{9}{299.672}{236.78}{\Character{40}\Character{1}\Character{41}\Character{41}}%
\TextAt{16}{326.405}{236.78}{\Character{106}}%
\TextAt{13}{333.046}{236.78}{K}%
\TextAt{13}{343.208}{236.78}{\Character{59}}%
\TextAt{8}{79.226}{211.033}{where}%
\TextAt{9}{113.717}{211.033}{\Character{1}}%
\NativeMathAt{127.955}{211.033}{11.955}{\subseteq}%
\TextAt{13}{141.737}{211.033}{A}%
\TextAt{21}{150.512}{209.24}{K}%
\TextAt{8}{158.715}{211.033}{.}%
\TextAt{8}{169.209}{211.033}{A}%
\TextAt{8}{182.563}{211.033}{Kummer}%
\TextAt{8}{232.01}{211.033}{extension}%
\TextAt{13}{284.458}{211.033}{K}%
\TextAt{9}{295.276}{211.033}{\Character{40}}%
\TextAt{13}{299.828}{211.033}{\Character{125}}%
\TextAt{3}{307.28}{215.372}{\Character{0}}%
\TextAt{17}{313.867}{215.372}{1}%
\TextAt{9}{318.599}{211.033}{\Character{40}\Character{1}\Character{41}\Character{41}}%
\TextAt{16}{346.494}{211.033}{\Character{106}}%
\TextAt{13}{354.298}{211.033}{K}%
\TextAt{8}{369.695}{211.033}{is}%
\TextAt{8}{382.141}{211.033}{always}%
\TextAt{8}{420.291}{211.033}{Galois,}%
\TextAt{8}{460.308}{211.033}{and}%
\TextAt{8}{483.753}{211.033}{its}%
\TextAt{8}{500.754}{211.033}{Galois}%
\TextAt{8}{79.226}{193.987}{group}%
\TextAt{8}{113.419}{193.987}{is}%
\TextAt{8}{126.201}{193.987}{abelian}%
\TextAt{8}{167.866}{193.987}{of}%
\TextAt{8}{182.214}{193.987}{exponent}%
\TextAt{13}{233.635}{193.987}{n}%
\TextAt{8}{240.622}{193.987}{.}%
\TextAt{8}{252.123}{193.987}{Indeed,}%
\TextAt{8}{294.677}{193.987}{for}%
\TextAt{8}{313.579}{193.987}{an}%
\TextAt{8}{330.856}{193.987}{extension}%
\TextAt{13}{383.64}{193.987}{K}%
\TextAt{9}{394.457}{193.987}{\Character{40}}%
\TextAt{13}{399.01}{193.987}{\Character{125}}%
\TextAt{3}{406.462}{198.326}{\Character{0}}%
\TextAt{17}{413.049}{198.326}{1}%
\TextAt{9}{417.781}{193.987}{\Character{40}}%
\TextAt{13}{422.333}{193.987}{a}%
\TextAt{9}{428.478}{193.987}{\Character{41}\Character{41}}%
\TextAt{16}{442.638}{193.987}{\Character{106}}%
\TextAt{13}{451.014}{193.987}{K}%
\TextAt{8}{461.832}{193.987}{,}%
\TextAt{8}{470.253}{193.987}{we}%
\TextAt{8}{488.498}{193.987}{have}%
\TextAt{8}{516.51}{193.987}{the}%
\TextAt{8}{79.226}{176.941}{injective}%
\TextAt{8}{125.72}{176.941}{homomorphism}%
\TextAt{13}{198.406}{151.195}{G}%
\TextAt{9}{207.64}{151.195}{\Character{40}}%
\TextAt{13}{212.192}{151.195}{K}%
\TextAt{9}{223.01}{151.195}{\Character{40}}%
\TextAt{13}{227.562}{151.195}{\Character{125}}%
\TextAt{3}{235.014}{156.131}{\Character{0}}%
\TextAt{17}{241.601}{156.131}{1}%
\TextAt{9}{246.333}{151.195}{\Character{40}}%
\TextAt{13}{250.886}{151.195}{a}%
\TextAt{9}{257.031}{151.195}{\Character{41}\Character{41}}%
\TextAt{16}{269.456}{151.195}{\Character{106}}%
\TextAt{13}{276.098}{151.195}{K}%
\TextAt{9}{286.915}{151.195}{\Character{41}}%
\NativeArrowAt{294.789}{151.195}{11.955}{13.214}{\hookrightarrow}%
\TextAt{13}{311.324}{151.195}{\Character{22}}%
\TextAt{21}{318.367}{149.402}{\Character{125}}%
\TextAt{13}{324.232}{151.195}{\Character{59}}%
\TextAt{13}{353.367}{151.195}{\Character{27}}%
\NativeArrowAt{363.79}{151.195}{11.955}{19.26}{\longmapsto}%
\TextAt{13}{386.372}{151.195}{\Character{11}}%
\TextAt{21}{393.894}{156.131}{\Character{27}}%
\TextAt{3}{399.019}{156.131}{\Character{0}}%
\TextAt{17}{405.606}{156.131}{1}%
\TextAt{13}{410.338}{151.195}{\Character{59}}%
\TextAt{8}{79.226}{125.449}{where}%
\TextAt{13}{113.651}{125.449}{\Character{11}}%
\NativeMathAt{125.543}{125.449}{11.955}{\in}%
\TextAt{13}{137.884}{125.449}{\Character{125}}%
\TextAt{3}{145.337}{129.787}{\Character{0}}%
\TextAt{17}{151.923}{129.787}{1}%
\TextAt{9}{156.655}{125.449}{\Character{40}}%
\TextAt{13}{161.208}{125.449}{a}%
\TextAt{9}{167.353}{125.449}{\Character{41}}%
\TextAt{8}{171.905}{125.449}{.}%
\TextAt{8}{182.201}{125.449}{Since}%
\TextAt{13}{213.375}{125.449}{\Character{22}}%
\TextAt{21}{220.417}{123.656}{\Character{125}}%
\NativeMathAt{230.654}{125.449}{11.955}{\subseteq}%
\TextAt{13}{244.323}{125.449}{A}%
\TextAt{21}{253.098}{123.656}{K}%
\TextAt{8}{261.301}{125.449}{,}%
\TextAt{8}{269.22}{125.449}{this}%
\TextAt{8}{292.658}{125.449}{definition}%
\TextAt{8}{345.303}{125.449}{does}%
\TextAt{8}{372.316}{125.449}{not}%
\TextAt{8}{393.746}{125.449}{depend}%
\TextAt{8}{434.997}{125.449}{on}%
\TextAt{8}{451.872}{125.449}{the}%
\TextAt{8}{472.645}{125.449}{choice}%
\TextAt{8}{508.05}{125.449}{of}%
\TextAt{13}{521.996}{125.449}{\Character{11}}%
\TextAt{8}{529.518}{125.449}{.}%
\TextAt{8}{79.226}{108.403}{Thus,}%
\TextAt{8}{112.127}{108.403}{for}%
\TextAt{8}{130.012}{108.403}{a}%
\TextAt{8}{139.767}{108.403}{Kummer}%
\TextAt{8}{188.532}{108.403}{extension}%
\TextAt{13}{240.298}{108.403}{L}%
\TextAt{9}{251.583}{108.403}{\Character{61}}%
\TextAt{13}{264.009}{108.403}{K}%
\TextAt{9}{274.827}{108.403}{\Character{40}}%
\TextAt{13}{279.379}{108.403}{\Character{125}}%
\TextAt{3}{286.831}{112.741}{\Character{0}}%
\TextAt{17}{293.418}{112.741}{1}%
\TextAt{9}{298.15}{108.403}{\Character{40}\Character{1}\Character{41}\Character{41}}%
\TextAt{9}{324.834}{108.403}{\Character{61}}%
\TextAt{20}{337.309}{117.369}{\Character{81}}%
\TextAt{21}{348.6}{104.916}{a}%
\NativeMathAt{353.098}{104.916}{7.97}{\in}%
\TextAt{17}{358.743}{104.916}{\Character{1}}%
\TextAt{13}{368.291}{108.403}{K}%
\TextAt{9}{379.108}{108.403}{\Character{40}}%
\TextAt{13}{383.661}{108.403}{\Character{125}}%
\TextAt{3}{391.113}{112.741}{\Character{0}}%
\TextAt{17}{397.7}{112.741}{1}%
\TextAt{9}{402.432}{108.403}{\Character{40}}%
\TextAt{13}{406.985}{108.403}{a}%
\TextAt{9}{413.129}{108.403}{\Character{41}\Character{41}}%
\TextAt{8}{422.234}{108.403}{,}%
\TextAt{8}{429.383}{108.403}{the}%
\TextAt{8}{449.54}{108.403}{composite}%
\TextAt{8}{504.557}{108.403}{map}%
\TextAt{13}{193.127}{82.657}{G}%
\TextAt{9}{202.361}{82.657}{\Character{40}}%
\TextAt{13}{206.913}{82.657}{L}%
\TextAt{16}{218.198}{82.657}{\Character{106}}%
\TextAt{13}{224.84}{82.657}{K}%
\TextAt{9}{235.658}{82.657}{\Character{41}}%
\NativeArrowAt{243.531}{82.657}{11.955}{19.26}{\longrightarrow}%
\TextAt{20}{267.075}{94.014}{\Character{89}}%
\TextAt{21}{266.113}{68.642}{a}%
\NativeMathAt{270.611}{68.642}{7.97}{\in}%
\TextAt{17}{276.256}{68.642}{\Character{1}}%
\TextAt{13}{285.306}{82.657}{G}%
\TextAt{9}{294.54}{82.657}{\Character{40}}%
\TextAt{13}{299.092}{82.657}{K}%
\TextAt{9}{309.91}{82.657}{\Character{40}}%
\TextAt{13}{314.462}{82.657}{\Character{125}}%
\TextAt{3}{321.914}{87.593}{\Character{0}}%
\TextAt{17}{328.501}{87.593}{1}%
\TextAt{9}{333.233}{82.657}{\Character{40}}%
\TextAt{13}{337.786}{82.657}{a}%
\TextAt{9}{343.93}{82.657}{\Character{41}\Character{41}}%
\TextAt{16}{356.356}{82.657}{\Character{106}}%
\TextAt{13}{362.998}{82.657}{K}%
\TextAt{9}{373.815}{82.657}{\Character{41}}%
\NativeArrowAt{381.689}{82.657}{11.955}{19.26}{\longrightarrow}%
\TextAt{13}{404.271}{82.657}{\Character{22}}%
\TextAt{17}{411.314}{87.593}{\Character{1}}%
\TextAt{21}{411.314}{79.701}{\Character{125}}%
\TextAt{8}{297.211}{43.033}{254}%
\Destination{78.23}{809.815}{nav.page.254}%
\Destination{79.36}{493.855}{nav.print.278}%
\Destination{79.23}{486}{nav.number.4.3.1}%
}

% Page 267 (body page 255)
\Page{%
\TextAt{2}{79.226}{796.943}{IV.}%
\TextAt{2}{98.012}{796.943}{Abstract}%
\TextAt{2}{143.252}{796.943}{Class}%
\TextAt{2}{171.844}{796.943}{Field}%
\TextAt{2}{199.575}{796.943}{Theory}%
\Vector{q 0 G 0 g 0.299 w 7.226 18.729 m 460.769 18.729 l S Q}%
\TextAt{8}{79.226}{761.124}{is}%
\TextAt{8}{90.99}{761.124}{an}%
\TextAt{8}{107.249}{761.124}{injective}%
\TextAt{8}{153.743}{761.124}{homomorphism.}%
\TextAt{8}{97.159}{741.089}{The}%
\TextAt{8}{121.461}{741.089}{following}%
\TextAt{8}{171.144}{741.089}{proposition}%
\TextAt{8}{232.926}{741.089}{says}%
\TextAt{8}{258.017}{741.089}{that}%
\TextAt{8}{283.635}{741.089}{conversely,}%
\TextAt{8}{342.215}{741.089}{any}%
\TextAt{8}{364.58}{741.089}{abelian}%
\TextAt{8}{405.476}{741.089}{extension}%
\TextAt{13}{457.49}{741.089}{L}%
\TextAt{16}{469.199}{741.089}{\Character{106}}%
\TextAt{13}{476.264}{741.089}{K}%
\TextAt{8}{491.228}{741.089}{of}%
\TextAt{8}{504.806}{741.089}{expo-}%
\TextAt{8}{79.226}{724.043}{nent}%
\TextAt{13}{105.563}{724.043}{n}%
\TextAt{8}{116.448}{724.043}{is}%
\TextAt{8}{128.212}{724.043}{a}%
\TextAt{8}{137.968}{724.043}{Kummer}%
\TextAt{8}{186.733}{724.043}{extension.}%
\TextAt{12}{79.226}{698.34}{Proposition}%
\TextAt{12}{151.937}{698.34}{(3.2).}%
\TextAt{8}{187.253}{698.34}{If}%
\TextAt{13}{198.945}{698.34}{L}%
\TextAt{16}{210.231}{698.34}{\Character{106}}%
\TextAt{13}{216.872}{698.34}{K}%
\TextAt{8}{231.587}{698.34}{is}%
\TextAt{8}{243.351}{698.34}{an}%
\TextAt{8}{259.61}{698.34}{abelian}%
\TextAt{8}{300.258}{698.34}{extension}%
\TextAt{8}{352.024}{698.34}{of}%
\TextAt{8}{365.354}{698.34}{exponent}%
\TextAt{13}{415.757}{698.34}{n}%
\TextAt{8}{422.744}{698.34}{,}%
\TextAt{8}{429.894}{698.34}{then}%
\TextAt{13}{203.057}{670.423}{L}%
\TextAt{9}{214.342}{670.423}{\Character{61}}%
\TextAt{13}{226.768}{670.423}{K}%
\TextAt{9}{237.586}{670.423}{\Character{40}}%
\TextAt{13}{242.138}{670.423}{\Character{125}}%
\TextAt{3}{249.59}{675.359}{\Character{0}}%
\TextAt{17}{256.177}{675.359}{1}%
\TextAt{9}{260.909}{670.423}{\Character{40}\Character{1}\Character{41}\Character{41}}%
\TextAt{8}{308.232}{670.423}{with}%
\TextAt{9}{334.892}{670.423}{\Character{1}}%
\TextAt{9}{347.947}{670.423}{\Character{61}}%
\TextAt{13}{360.393}{670.423}{A}%
\TextAt{21}{369.168}{676.09}{\Character{125}}%
\TextAt{21}{369.168}{667.182}{L}%
\TextAt{16}{378.082}{670.423}{\Character{92}}%
\TextAt{13}{388.709}{670.423}{A}%
\TextAt{21}{397.484}{668.63}{K}%
\TextAt{13}{405.687}{670.423}{\Character{58}}%
\TextAt{8}{79.226}{642.505}{If}%
\TextAt{8}{90.918}{642.505}{in}%
\TextAt{8}{104.571}{642.505}{particular,}%
\TextAt{13}{161.812}{642.505}{L}%
\TextAt{16}{173.098}{642.505}{\Character{106}}%
\TextAt{13}{179.74}{642.505}{K}%
\TextAt{8}{194.455}{642.505}{is}%
\TextAt{8}{206.218}{642.505}{cyclic,}%
\TextAt{8}{241.654}{642.505}{then}%
\TextAt{8}{268.314}{642.505}{we}%
\TextAt{8}{285.541}{642.505}{find}%
\TextAt{13}{308.949}{642.505}{L}%
\TextAt{9}{320.235}{642.505}{\Character{61}}%
\TextAt{13}{332.66}{642.505}{K}%
\TextAt{9}{343.478}{642.505}{\Character{40}}%
\TextAt{13}{348.03}{642.505}{\Character{11}}%
\TextAt{9}{355.552}{642.505}{\Character{41}}%
\TextAt{8}{364.002}{642.505}{with}%
\TextAt{13}{390.662}{642.505}{\Character{11}}%
\TextAt{21}{398.183}{646.844}{\Character{125}}%
\TextAt{9}{407.37}{642.505}{\Character{61}}%
\TextAt{13}{419.795}{642.505}{a}%
\NativeMathAt{429.261}{642.505}{11.955}{\in}%
\TextAt{13}{440.552}{642.505}{A}%
\TextAt{21}{449.327}{640.712}{K}%
\TextAt{8}{457.53}{642.505}{.}%
\TextAt{19}{79.226}{608.146}{Proof.}%
\TextAt{8}{116}{608.146}{We}%
\TextAt{8}{137.443}{608.146}{have}%
\TextAt{13}{165.736}{608.146}{\Character{125}}%
\TextAt{3}{173.189}{612.485}{\Character{0}}%
\TextAt{17}{179.775}{612.485}{1}%
\TextAt{9}{184.508}{608.146}{\Character{40}\Character{1}\Character{41}}%
\NativeMathAt{208.901}{608.146}{11.955}{\subseteq}%
\TextAt{13}{223.732}{608.146}{A}%
\TextAt{21}{232.508}{606.353}{L}%
\TextAt{8}{238.765}{608.146}{,}%
\TextAt{8}{247.537}{608.146}{for}%
\TextAt{8}{266.72}{608.146}{if}%
\TextAt{13}{278.743}{608.146}{x}%
\NativeMathAt{290.928}{608.146}{11.955}{\in}%
\TextAt{13}{304.432}{608.146}{A}%
\TextAt{8}{318.403}{608.146}{and}%
\TextAt{13}{342.464}{608.146}{x}%
\TextAt{21}{349.116}{612.485}{\Character{125}}%
\TextAt{9}{360.515}{608.146}{\Character{61}}%
\TextAt{13}{375.153}{608.146}{\Character{11}}%
\TextAt{21}{382.675}{612.485}{\Character{125}}%
\TextAt{9}{394.074}{608.146}{\Character{61}}%
\TextAt{13}{408.712}{608.146}{a}%
\NativeMathAt{420.39}{608.146}{11.955}{\in}%
\TextAt{13}{433.894}{608.146}{A}%
\TextAt{21}{442.669}{606.353}{K}%
\TextAt{8}{450.871}{608.146}{,}%
\TextAt{13}{459.644}{608.146}{\Character{11}}%
\NativeMathAt{472.699}{608.146}{11.955}{\in}%
\TextAt{13}{486.202}{608.146}{A}%
\TextAt{21}{494.978}{606.353}{L}%
\TextAt{8}{501.235}{608.146}{,}%
\TextAt{8}{510.007}{608.146}{then}%
\TextAt{13}{79.226}{591.1}{x}%
\TextAt{9}{90.381}{591.1}{\Character{61}}%
\TextAt{13}{103.989}{591.1}{\Character{24}}%
\TextAt{13}{109.668}{591.1}{\Character{11}}%
\NativeMathAt{121.704}{591.1}{11.955}{\in}%
\TextAt{13}{134.178}{591.1}{A}%
\TextAt{21}{142.953}{589.307}{L}%
\TextAt{8}{153.801}{591.1}{for}%
\TextAt{8}{172.38}{591.1}{some}%
\TextAt{13}{202.4}{591.1}{\Character{24}}%
\NativeMathAt{212.593}{591.1}{11.955}{\in}%
\TextAt{13}{225.067}{591.1}{\Character{22}}%
\TextAt{21}{232.11}{589.307}{\Character{125}}%
\NativeMathAt{242.478}{591.1}{11.955}{\subseteq}%
\TextAt{13}{256.28}{591.1}{A}%
\TextAt{21}{265.055}{589.307}{K}%
\TextAt{8}{273.258}{591.1}{.}%
\TextAt{8}{283.788}{591.1}{Therefore}%
\TextAt{13}{337.479}{591.1}{K}%
\TextAt{9}{348.296}{591.1}{\Character{40}}%
\TextAt{13}{352.849}{591.1}{\Character{125}}%
\TextAt{3}{360.301}{595.439}{\Character{0}}%
\TextAt{17}{366.888}{595.439}{1}%
\TextAt{9}{371.62}{591.1}{\Character{40}\Character{1}\Character{41}\Character{41}}%
\NativeMathAt{399.535}{591.1}{11.955}{\subseteq}%
\TextAt{13}{413.337}{591.1}{L}%
\TextAt{8}{421.301}{591.1}{.}%
\TextAt{8}{431.831}{591.1}{On}%
\TextAt{8}{452.036}{591.1}{the}%
\TextAt{8}{472.886}{591.1}{other}%
\TextAt{8}{504.149}{591.1}{hand,}%
\TextAt{8}{79.226}{574.054}{the}%
\TextAt{8}{99.672}{574.054}{extension}%
\TextAt{13}{151.727}{574.054}{L}%
\TextAt{16}{163.506}{574.054}{\Character{106}}%
\TextAt{13}{170.641}{574.054}{K}%
\TextAt{8}{185.646}{574.054}{is}%
\TextAt{8}{197.699}{574.054}{the}%
\TextAt{8}{218.145}{574.054}{composite}%
\TextAt{8}{273.452}{574.054}{of}%
\TextAt{8}{287.072}{574.054}{its}%
\TextAt{8}{303.68}{574.054}{cyclic}%
\TextAt{8}{336.153}{574.054}{subextensions.}%
\TextAt{8}{415.897}{574.054}{For}%
\TextAt{8}{437.156}{574.054}{it}%
\TextAt{8}{449.15}{574.054}{is}%
\TextAt{8}{461.203}{574.054}{the}%
\TextAt{8}{481.649}{574.054}{composite}%
\TextAt{8}{79.226}{557.008}{of}%
\TextAt{8}{93.242}{557.008}{its}%
\TextAt{8}{110.246}{557.008}{finite}%
\TextAt{8}{140.843}{557.008}{subextensions,}%
\TextAt{8}{219.277}{557.008}{and}%
\TextAt{8}{242.725}{557.008}{the}%
\TextAt{8}{263.567}{557.008}{Galois}%
\TextAt{8}{300.166}{557.008}{group}%
\TextAt{8}{334.027}{557.008}{of}%
\TextAt{8}{348.042}{557.008}{a}%
\TextAt{8}{358.483}{557.008}{finite}%
\TextAt{8}{389.08}{557.008}{subextension}%
\TextAt{8}{459.476}{557.008}{is}%
\TextAt{8}{471.926}{557.008}{the}%
\TextAt{8}{492.767}{557.008}{product}%
\TextAt{8}{79.226}{539.962}{of}%
\TextAt{8}{92.721}{539.962}{cyclic}%
\TextAt{8}{125.069}{539.962}{groups,}%
\TextAt{8}{166.317}{539.962}{which}%
\TextAt{8}{199.969}{539.962}{may}%
\TextAt{8}{225.503}{539.962}{be}%
\TextAt{8}{241.592}{539.962}{interpreted}%
\TextAt{8}{301.915}{539.962}{as}%
\TextAt{8}{316.45}{539.962}{Galois}%
\TextAt{8}{352.528}{539.962}{groups}%
\TextAt{8}{390.483}{539.962}{of}%
\TextAt{8}{403.978}{539.962}{cyclic}%
\TextAt{8}{436.327}{539.962}{subextensions.}%
\TextAt{8}{515.697}{539.962}{Let}%
\TextAt{8}{79.226}{522.916}{now}%
\TextAt{13}{103.747}{522.916}{M}%
\TextAt{16}{119.868}{522.916}{\Character{106}}%
\TextAt{13}{126.736}{522.916}{K}%
\TextAt{8}{141.583}{522.916}{be}%
\TextAt{8}{157.64}{522.916}{a}%
\TextAt{8}{167.528}{522.916}{cyclic}%
\TextAt{8}{199.844}{522.916}{subextension}%
\TextAt{8}{269.688}{522.916}{of}%
\TextAt{13}{283.15}{522.916}{L}%
\TextAt{16}{294.662}{522.916}{\Character{106}}%
\TextAt{13}{301.529}{522.916}{K}%
\TextAt{8}{312.347}{522.916}{.}%
\TextAt{8}{321.193}{522.916}{It}%
\TextAt{8}{333.999}{522.916}{suffices}%
\TextAt{8}{373.918}{522.916}{to}%
\TextAt{8}{388.361}{522.916}{show}%
\TextAt{8}{417.497}{522.916}{that}%
\TextAt{13}{442.998}{522.916}{M}%
\NativeMathAt{459.119}{522.916}{11.955}{\subseteq}%
\TextAt{13}{471.964}{522.916}{K}%
\TextAt{9}{482.782}{522.916}{\Character{40}}%
\TextAt{13}{487.334}{522.916}{\Character{125}}%
\TextAt{3}{494.787}{527.255}{\Character{0}}%
\TextAt{17}{501.373}{527.255}{1}%
\TextAt{9}{506.106}{522.916}{\Character{40}\Character{1}\Character{41}\Character{41}}%
\TextAt{8}{529.518}{522.916}{.}%
\TextAt{8}{79.226}{505.87}{Let}%
\TextAt{13}{100.601}{505.87}{\Character{27}}%
\TextAt{8}{111.986}{505.87}{be}%
\TextAt{8}{128.315}{505.87}{a}%
\TextAt{8}{138.476}{505.87}{generator}%
\TextAt{8}{190.922}{505.87}{of}%
\TextAt{13}{204.657}{505.87}{G}%
\TextAt{9}{213.89}{505.87}{\Character{40}}%
\TextAt{13}{218.443}{505.87}{M}%
\TextAt{16}{235.028}{505.87}{\Character{106}}%
\TextAt{13}{242.36}{505.87}{K}%
\TextAt{9}{253.177}{505.87}{\Character{41}}%
\TextAt{8}{262.032}{505.87}{and}%
\TextAt{13}{285.2}{505.87}{\Character{16}}%
\TextAt{8}{295.518}{505.87}{a}%
\TextAt{8}{305.679}{505.87}{generator}%
\TextAt{8}{358.125}{505.87}{of}%
\TextAt{13}{371.86}{505.87}{\Character{22}}%
\TextAt{21}{378.903}{504.077}{\Character{125}}%
\TextAt{8}{384.769}{505.87}{.}%
\TextAt{8}{394.432}{505.87}{Let}%
\TextAt{13}{415.807}{505.87}{d}%
\TextAt{9}{425.901}{505.87}{\Character{61}}%
\TextAt{9}{439.004}{505.87}{\Character{91}}%
\TextAt{13}{442.268}{505.87}{M}%
\TextAt{9}{458.853}{505.87}{\Character{58}}%
\TextAt{13}{466.116}{505.87}{K}%
\TextAt{9}{476.934}{505.87}{\Character{93}}%
\TextAt{8}{480.186}{505.87}{,}%
\TextAt{13}{487.841}{505.87}{d}%
\TextAt{3}{493.924}{510.209}{\Character{48}}%
\TextAt{9}{500.73}{505.87}{\Character{61}}%
\TextAt{13}{513.846}{505.87}{n}%
\TextAt{9}{520.834}{505.87}{\Character{47}}%
\TextAt{13}{526.687}{505.87}{d}%
\TextAt{8}{79.226}{488.824}{and}%
\TextAt{13}{101.986}{488.824}{\Character{24}}%
\TextAt{9}{110.997}{488.824}{\Character{61}}%
\TextAt{13}{123.423}{488.824}{\Character{16}}%
\TextAt{21}{129.439}{493.163}{d}%
\TextAt{23}{133.796}{495.975}{\Character{48}}%
\TextAt{8}{136.978}{488.824}{.}%
\TextAt{8}{145.426}{488.824}{Since}%
\TextAt{13}{175.981}{488.824}{N}%
\TextAt{21}{185.354}{486.965}{M}%
\TextAt{3}{194.336}{486.965}{\Character{106}}%
\TextAt{21}{196.688}{486.965}{K}%
\TextAt{9}{204.891}{488.824}{\Character{40}}%
\TextAt{13}{209.443}{488.824}{\Character{24}}%
\TextAt{9}{215.134}{488.824}{\Character{41}}%
\TextAt{9}{222.989}{488.824}{\Character{61}}%
\TextAt{13}{235.432}{488.824}{\Character{24}}%
\TextAt{21}{241.123}{493.163}{d}%
\TextAt{9}{249.299}{488.824}{\Character{61}}%
\TextAt{9}{261.709}{488.824}{1}%
\TextAt{8}{267.578}{488.824}{,}%
\TextAt{8}{274.725}{488.824}{(3.1)}%
\TextAt{8}{302.697}{488.824}{shows}%
\TextAt{8}{336.313}{488.824}{that}%
\TextAt{13}{361.679}{488.824}{\Character{24}}%
\TextAt{9}{370.69}{488.824}{\Character{61}}%
\TextAt{13}{383.116}{488.824}{\Character{11}}%
\TextAt{21}{390.637}{493.163}{\Character{27}}%
\TextAt{3}{395.763}{493.163}{\Character{0}}%
\TextAt{17}{402.349}{493.163}{1}%
\TextAt{8}{410.976}{488.824}{for}%
\TextAt{8}{428.858}{488.824}{some}%
\TextAt{13}{458.182}{488.824}{\Character{11}}%
\NativeMathAt{469.024}{488.824}{11.955}{\in}%
\TextAt{13}{480.315}{488.824}{A}%
\TextAt{21}{489.091}{487.031}{M}%
\TextAt{8}{498.571}{488.824}{.}%
\TextAt{8}{507.018}{488.824}{Thus}%
\TextAt{13}{79.226}{471.778}{K}%
\NativeMathAt{94.01}{471.778}{11.955}{\subseteq}%
\TextAt{13}{107.275}{471.778}{K}%
\TextAt{9}{118.093}{471.778}{\Character{40}}%
\TextAt{13}{122.645}{471.778}{\Character{11}}%
\TextAt{9}{130.167}{471.778}{\Character{41}}%
\NativeMathAt{138.685}{471.778}{11.955}{\subseteq}%
\TextAt{13}{151.95}{471.778}{M}%
\TextAt{8}{164.524}{471.778}{.}%
\TextAt{8}{174.108}{471.778}{But}%
\TextAt{13}{197.728}{471.778}{\Character{11}}%
\TextAt{21}{205.25}{476.117}{\Character{27}}%
\TextAt{22}{210.375}{478.93}{i}%
\TextAt{9}{218.001}{471.778}{\Character{61}}%
\TextAt{13}{231.072}{471.778}{\Character{24}}%
\TextAt{21}{236.763}{476.117}{i}%
\TextAt{13}{240.144}{471.778}{\Character{11}}%
\TextAt{8}{247.666}{471.778}{.}%
\TextAt{8}{257.25}{471.778}{Thus}%
\TextAt{13}{287.278}{471.778}{\Character{11}}%
\TextAt{21}{294.8}{476.117}{\Character{27}}%
\TextAt{22}{299.925}{478.93}{i}%
\TextAt{9}{307.551}{471.778}{\Character{61}}%
\TextAt{13}{320.622}{471.778}{\Character{11}}%
\TextAt{8}{332.42}{471.778}{is}%
\TextAt{8}{344.563}{471.778}{equivalent}%
\TextAt{8}{400.557}{471.778}{to}%
\TextAt{13}{415.246}{471.778}{i}%
\NativeMathAt{423.206}{471.778}{11.955}{\equiv}%
\TextAt{9}{436.471}{471.778}{0}%
\TextAt{9}{447.625}{471.778}{\Character{40}}%
\TextAt{8}{452.189}{471.778}{mod}%
\TextAt{13}{478.291}{471.778}{d}%
\TextAt{9}{484.374}{471.778}{\Character{41}}%
\TextAt{8}{488.926}{471.778}{,}%
\TextAt{8}{496.549}{471.778}{so}%
\TextAt{8}{511.298}{471.778}{that}%
\TextAt{13}{79.226}{454.732}{K}%
\TextAt{9}{90.044}{454.732}{\Character{40}}%
\TextAt{13}{94.596}{454.732}{\Character{11}}%
\TextAt{9}{102.118}{454.732}{\Character{41}}%
\TextAt{9}{110.63}{454.732}{\Character{61}}%
\TextAt{13}{123.734}{454.732}{M}%
\TextAt{8}{136.307}{454.732}{.}%
\TextAt{8}{145.915}{454.732}{But}%
\TextAt{9}{169.542}{454.732}{\Character{40}}%
\TextAt{13}{174.095}{454.732}{\Character{11}}%
\TextAt{21}{181.616}{459.071}{\Character{125}}%
\TextAt{9}{187.482}{454.732}{\Character{41}}%
\TextAt{21}{192.034}{459.071}{\Character{27}}%
\TextAt{3}{197.16}{459.071}{\Character{0}}%
\TextAt{17}{203.746}{459.071}{1}%
\TextAt{9}{212.458}{454.732}{\Character{61}}%
\TextAt{9}{225.525}{454.732}{\Character{40}}%
\TextAt{13}{230.094}{454.732}{\Character{11}}%
\TextAt{21}{237.616}{459.071}{\Character{27}}%
\TextAt{3}{242.741}{459.071}{\Character{0}}%
\TextAt{17}{249.328}{459.071}{1}%
\TextAt{9}{254.06}{454.732}{\Character{41}}%
\TextAt{21}{258.612}{459.071}{\Character{125}}%
\TextAt{9}{268.457}{454.732}{\Character{61}}%
\TextAt{13}{281.541}{454.732}{\Character{24}}%
\TextAt{21}{287.232}{459.071}{\Character{125}}%
\TextAt{9}{297.077}{454.732}{\Character{61}}%
\TextAt{9}{310.144}{454.732}{1}%
\TextAt{8}{316.014}{454.732}{,}%
\TextAt{8}{323.646}{454.732}{so}%
\TextAt{8}{338.403}{454.732}{that}%
\TextAt{13}{364.158}{454.732}{a}%
\TextAt{9}{374.282}{454.732}{\Character{61}}%
\TextAt{13}{387.366}{454.732}{\Character{11}}%
\TextAt{21}{394.888}{459.071}{\Character{125}}%
\NativeMathAt{404.733}{454.732}{11.955}{\in}%
\TextAt{13}{416.683}{454.732}{A}%
\TextAt{21}{425.458}{452.939}{K}%
\TextAt{8}{433.661}{454.732}{;}%
\TextAt{8}{441.39}{454.732}{then}%
\TextAt{13}{468.436}{454.732}{\Character{11}}%
\NativeMathAt{479.937}{454.732}{11.955}{\in}%
\TextAt{13}{491.887}{454.732}{\Character{125}}%
\TextAt{3}{499.339}{459.071}{\Character{0}}%
\TextAt{17}{505.926}{459.071}{1}%
\TextAt{9}{510.658}{454.732}{\Character{40}\Character{1}\Character{41}}%
\TextAt{8}{529.518}{454.732}{,}%
\TextAt{8}{79.226}{437.686}{and}%
\TextAt{8}{101.989}{437.686}{therefore}%
\TextAt{13}{151.089}{437.686}{M}%
\NativeMathAt{166.983}{437.686}{11.955}{\subseteq}%
\TextAt{13}{179.602}{437.686}{K}%
\TextAt{9}{190.42}{437.686}{\Character{40}}%
\TextAt{13}{194.972}{437.686}{\Character{125}}%
\TextAt{3}{202.425}{442.025}{\Character{0}}%
\TextAt{17}{209.012}{442.025}{1}%
\TextAt{9}{213.744}{437.686}{\Character{40}\Character{1}\Character{41}\Character{41}}%
\TextAt{8}{237.156}{437.686}{.}%
\Vector{q 0 G 0 g 0.399 w 452.334 -336.314 m 452.334 -328.244 l S Q}%
\Vector{q 0 G 0 g 0.399 w 452.534 -328.443 m 459.707 -328.443 l S Q}%
\Vector{q 0 G 0 g 0.399 w 452.534 -336.114 m 459.707 -336.114 l S Q}%
\Vector{q 0 G 0 g 0.399 w 459.906 -336.314 m 459.906 -328.244 l S Q}%
\TextAt{8}{97.159}{408.995}{As}%
\TextAt{8}{114.446}{408.995}{the}%
\TextAt{8}{134.602}{408.995}{main}%
\TextAt{8}{163.869}{408.995}{result}%
\TextAt{8}{196.447}{408.995}{of}%
\TextAt{8}{209.777}{408.995}{general}%
\TextAt{8}{250.101}{408.995}{Kummer}%
\TextAt{8}{298.867}{408.995}{theory,}%
\TextAt{8}{337.888}{408.995}{we}%
\TextAt{8}{355.116}{408.995}{now}%
\TextAt{8}{379.504}{408.995}{obtain}%
\TextAt{8}{415.932}{408.995}{the}%
\TextAt{8}{436.088}{408.995}{following}%
\TextAt{12}{79.226}{383.292}{Theorem}%
\TextAt{12}{136.24}{383.292}{(3.3).}%
\TextAt{8}{171.556}{383.292}{The}%
\TextAt{8}{195.61}{383.292}{correspondence}%
\TextAt{9}{247.537}{355.375}{\Character{1}}%
\NativeArrowAt{260.613}{355.375}{11.955}{19.26}{\longmapsto}%
\TextAt{13}{283.195}{355.375}{L}%
\TextAt{9}{294.48}{355.375}{\Character{61}}%
\TextAt{13}{306.905}{355.375}{K}%
\TextAt{9}{317.723}{355.375}{\Character{40}}%
\TextAt{13}{322.275}{355.375}{\Character{125}}%
\TextAt{3}{329.728}{360.311}{\Character{0}}%
\TextAt{17}{336.315}{360.311}{1}%
\TextAt{9}{341.047}{355.375}{\Character{40}\Character{1}\Character{41}\Character{41}}%
\TextAt{8}{79.226}{327.457}{is}%
\TextAt{8}{91.156}{327.457}{a}%
\TextAt{9}{101.078}{327.457}{1}%
\TextAt{8}{106.931}{327.457}{-}%
\TextAt{9}{110.829}{327.457}{1}%
\TextAt{8}{116.682}{327.457}{-correspondence}%
\TextAt{8}{202.423}{327.457}{between}%
\TextAt{8}{247.78}{327.457}{the}%
\TextAt{8}{268.103}{327.457}{groups}%
\TextAt{9}{306.06}{327.457}{\Character{1}}%
\TextAt{8}{319.879}{327.457}{such}%
\TextAt{8}{346.442}{327.457}{that}%
\TextAt{13}{371.978}{327.457}{A}%
\TextAt{21}{380.753}{333.125}{\Character{125}}%
\TextAt{21}{380.753}{324.216}{K}%
\NativeMathAt{392.56}{327.457}{11.955}{\subseteq}%
\TextAt{9}{405.463}{327.457}{\Character{1}}%
\NativeMathAt{418.823}{327.457}{11.955}{\subseteq}%
\TextAt{13}{431.726}{327.457}{A}%
\TextAt{21}{440.501}{325.664}{K}%
\TextAt{8}{452.767}{327.457}{and}%
\TextAt{8}{475.696}{327.457}{the}%
\TextAt{8}{496.019}{327.457}{abelian}%
\TextAt{8}{79.226}{310.411}{extensions}%
\TextAt{13}{135.607}{310.411}{L}%
\TextAt{16}{146.892}{310.411}{\Character{106}}%
\TextAt{13}{153.534}{310.411}{K}%
\TextAt{8}{168.249}{310.411}{of}%
\TextAt{8}{181.579}{310.411}{exponent}%
\TextAt{13}{231.982}{310.411}{n}%
\TextAt{8}{238.969}{310.411}{.}%
\TextAt{8}{97.159}{290.377}{If}%
\TextAt{9}{109.514}{290.377}{\Character{1}}%
\TextAt{8}{123.829}{290.377}{and}%
\TextAt{13}{147.254}{290.377}{L}%
\TextAt{8}{159.778}{290.377}{correspond}%
\TextAt{8}{220.014}{290.377}{to}%
\TextAt{8}{234.987}{290.377}{each}%
\TextAt{8}{261.986}{290.377}{other,}%
\TextAt{8}{296.636}{290.377}{then}%
\TextAt{13}{323.958}{290.377}{A}%
\TextAt{21}{332.734}{296.044}{\Character{125}}%
\TextAt{21}{332.734}{287.136}{L}%
\TextAt{16}{342.099}{290.377}{\Character{92}}%
\TextAt{13}{353.178}{290.377}{A}%
\TextAt{21}{361.953}{288.583}{K}%
\TextAt{9}{374.606}{290.377}{\Character{61}}%
\TextAt{9}{388.151}{290.377}{\Character{1}}%
\TextAt{8}{397.915}{290.377}{,}%
\TextAt{8}{405.893}{290.377}{and}%
\TextAt{8}{429.318}{290.377}{we}%
\TextAt{8}{447.208}{290.377}{have}%
\TextAt{8}{474.866}{290.377}{a}%
\TextAt{8}{485.284}{290.377}{canonical}%
\TextAt{8}{79.226}{273.331}{isomorphism}%
\TextAt{9}{171.11}{256.285}{\Character{1}\Character{47}}%
\TextAt{13}{186.718}{256.285}{A}%
\TextAt{21}{195.493}{261.952}{\Character{125}}%
\TextAt{21}{195.493}{253.044}{K}%
\NativeMathAt{207.017}{256.285}{11.955}{\cong}%
\TextAt{8}{219.636}{256.285}{Hom}%
\TextAt{9}{244.025}{256.285}{\Character{40}}%
\TextAt{13}{248.577}{256.285}{G}%
\TextAt{9}{257.811}{256.285}{\Character{40}}%
\TextAt{13}{262.363}{256.285}{L}%
\TextAt{16}{273.648}{256.285}{\Character{106}}%
\TextAt{13}{280.29}{256.285}{K}%
\TextAt{9}{291.108}{256.285}{\Character{41}}%
\TextAt{13}{295.66}{256.285}{\Character{59}}%
\TextAt{13}{300.884}{256.285}{\Character{22}}%
\TextAt{21}{307.947}{254.491}{\Character{125}}%
\TextAt{9}{313.813}{256.285}{\Character{41}}%
\TextAt{13}{318.365}{256.285}{\Character{59}}%
\TextAt{13}{347.5}{256.285}{a}%
\TextAt{8}{356.985}{256.285}{mod}%
\TextAt{13}{382.423}{256.285}{A}%
\TextAt{21}{391.199}{261.952}{\Character{125}}%
\TextAt{21}{391.199}{253.044}{K}%
\NativeArrowAt{402.722}{256.285}{11.955}{19.26}{\longmapsto}%
\TextAt{13}{425.304}{256.285}{\Character{31}}%
\TextAt{21}{432.638}{254.491}{a}%
\TextAt{13}{437.634}{256.285}{\Character{59}}%
\TextAt{8}{79.226}{233.228}{where}%
\TextAt{8}{113.035}{233.228}{the}%
\TextAt{8}{133.192}{233.228}{character}%
\TextAt{13}{184.252}{233.228}{\Character{31}}%
\TextAt{21}{191.586}{231.434}{a}%
\TextAt{9}{199.903}{233.227}{\Character{58}}%
\TextAt{13}{206.475}{233.227}{G}%
\TextAt{9}{215.709}{233.227}{\Character{40}}%
\TextAt{13}{220.261}{233.227}{L}%
\TextAt{16}{231.547}{233.227}{\Character{106}}%
\TextAt{13}{238.188}{233.227}{K}%
\TextAt{9}{249.006}{233.227}{\Character{41}}%
\NativeMathAt{256.879}{233.227}{11.955}{\rightarrow}%
\TextAt{13}{272.155}{233.227}{\Character{22}}%
\TextAt{21}{279.198}{231.434}{\Character{125}}%
\TextAt{8}{288.961}{233.228}{is}%
\TextAt{8}{300.725}{233.228}{given}%
\TextAt{8}{331.294}{233.228}{by}%
\TextAt{13}{347.553}{233.228}{\Character{31}}%
\TextAt{21}{354.887}{231.434}{a}%
\TextAt{9}{359.883}{233.227}{\Character{40}}%
\TextAt{13}{364.435}{233.227}{\Character{27}}%
\TextAt{9}{371.518}{233.227}{\Character{41}}%
\TextAt{9}{379.373}{233.227}{\Character{61}}%
\TextAt{13}{391.816}{233.227}{\Character{11}}%
\TextAt{21}{399.338}{237.566}{\Character{27}}%
\TextAt{3}{404.463}{237.566}{\Character{0}}%
\TextAt{17}{411.05}{237.566}{1}%
\TextAt{8}{415.782}{233.228}{,}%
\TextAt{8}{422.932}{233.228}{for}%
\TextAt{13}{440.816}{233.228}{\Character{11}}%
\NativeMathAt{451.659}{233.227}{11.955}{\in}%
\TextAt{13}{462.95}{233.228}{\Character{125}}%
\TextAt{3}{470.402}{237.566}{\Character{0}}%
\TextAt{17}{476.989}{237.566}{1}%
\TextAt{9}{481.721}{233.228}{\Character{40}}%
\TextAt{13}{486.274}{233.227}{a}%
\TextAt{9}{492.419}{233.228}{\Character{41}}%
\TextAt{8}{496.971}{233.228}{.}%
\TextAt{19}{79.226}{198.868}{Proof.}%
\TextAt{8}{116}{198.868}{Let}%
\TextAt{13}{138.005}{198.868}{L}%
\TextAt{16}{151.054}{198.868}{\Character{106}}%
\TextAt{13}{159.461}{198.868}{K}%
\TextAt{8}{175.211}{198.868}{be}%
\TextAt{8}{192.17}{198.868}{an}%
\TextAt{8}{209.465}{198.868}{abelian}%
\TextAt{8}{251.148}{198.868}{extension}%
\TextAt{8}{303.949}{198.868}{of}%
\TextAt{8}{318.314}{198.868}{exponent}%
\TextAt{13}{369.752}{198.868}{n}%
\TextAt{8}{376.74}{198.868}{.}%
\TextAt{8}{388.294}{198.868}{By}%
\TextAt{8}{407.693}{198.868}{(3.2),}%
\TextAt{8}{440.214}{198.868}{we}%
\TextAt{8}{458.476}{198.868}{then}%
\TextAt{8}{486.172}{198.868}{find}%
\TextAt{13}{510.615}{198.868}{L}%
\TextAt{9}{523.665}{198.868}{\Character{61}}%
\TextAt{13}{79.226}{181.822}{K}%
\TextAt{9}{90.044}{181.822}{\Character{40}}%
\TextAt{13}{94.596}{181.822}{\Character{125}}%
\TextAt{3}{102.049}{186.161}{\Character{0}}%
\TextAt{17}{108.635}{186.161}{1}%
\TextAt{9}{113.368}{181.822}{\Character{40}\Character{1}\Character{41}\Character{41}}%
\TextAt{8}{140.677}{181.822}{with}%
\TextAt{9}{167.337}{181.822}{\Character{1}}%
\TextAt{9}{180.392}{181.822}{\Character{61}}%
\TextAt{13}{192.838}{181.822}{A}%
\TextAt{21}{201.614}{187.49}{\Character{125}}%
\TextAt{21}{201.614}{178.581}{L}%
\TextAt{16}{210.527}{181.822}{\Character{92}}%
\TextAt{13}{221.154}{181.822}{A}%
\TextAt{21}{229.93}{180.029}{K}%
\TextAt{8}{238.132}{181.822}{.}%
\TextAt{8}{246.581}{181.822}{We}%
\TextAt{8}{266.725}{181.822}{consider}%
\TextAt{8}{312.31}{181.822}{the}%
\TextAt{8}{332.466}{181.822}{homomorphism}%
\TextAt{9}{200.413}{153.905}{\Character{1}}%
\NativeArrowAt{213.488}{153.905}{11.955}{19.26}{\longrightarrow}%
\TextAt{8}{236.07}{153.905}{Hom}%
\TextAt{9}{260.459}{153.905}{\Character{40}}%
\TextAt{13}{265.011}{153.905}{G}%
\TextAt{9}{274.245}{153.905}{\Character{40}}%
\TextAt{13}{278.797}{153.905}{L}%
\TextAt{16}{290.083}{153.905}{\Character{106}}%
\TextAt{13}{296.724}{153.905}{K}%
\TextAt{9}{307.542}{153.905}{\Character{41}}%
\TextAt{13}{312.094}{153.905}{\Character{59}}%
\TextAt{13}{317.318}{153.905}{\Character{22}}%
\TextAt{21}{324.381}{152.111}{\Character{125}}%
\TextAt{9}{330.247}{153.905}{\Character{41}}%
\TextAt{13}{334.799}{153.905}{\Character{59}}%
\TextAt{13}{363.934}{153.905}{a}%
\NativeArrowAt{373.42}{153.905}{11.955}{19.26}{\longmapsto}%
\TextAt{13}{396.002}{153.905}{\Character{31}}%
\TextAt{21}{403.335}{152.111}{a}%
\TextAt{13}{408.331}{153.905}{\Character{59}}%
\TextAt{8}{79.226}{125.987}{where}%
\TextAt{13}{113.035}{125.987}{\Character{31}}%
\TextAt{21}{120.369}{124.194}{a}%
\TextAt{9}{125.365}{125.987}{\Character{40}}%
\TextAt{13}{129.917}{125.987}{\Character{27}}%
\TextAt{9}{137}{125.987}{\Character{41}}%
\TextAt{9}{144.855}{125.987}{\Character{61}}%
\TextAt{13}{157.298}{125.987}{\Character{11}}%
\TextAt{21}{164.82}{130.326}{\Character{27}}%
\TextAt{3}{169.945}{130.326}{\Character{0}}%
\TextAt{17}{176.532}{130.326}{1}%
\TextAt{8}{181.264}{125.987}{,}%
\TextAt{13}{188.413}{125.987}{\Character{11}}%
\NativeMathAt{199.256}{125.987}{11.955}{\in}%
\TextAt{13}{210.547}{125.987}{\Character{125}}%
\TextAt{3}{218}{130.326}{\Character{0}}%
\TextAt{17}{224.586}{130.326}{1}%
\TextAt{9}{229.318}{125.987}{\Character{40}}%
\TextAt{13}{233.871}{125.987}{a}%
\TextAt{9}{240.016}{125.987}{\Character{41}}%
\TextAt{8}{244.568}{125.987}{.}%
\TextAt{8}{253.016}{125.987}{Since}%
\TextAt{13}{195.241}{98.07}{\Character{31}}%
\TextAt{21}{202.574}{96.277}{a}%
\TextAt{9}{210.891}{98.07}{\Character{61}}%
\TextAt{9}{223.3}{98.07}{1}%
\NativeArrowAt{235.812}{98.07}{11.955}{21.926}{\Longleftrightarrow}%
\TextAt{13}{264.371}{98.07}{\Character{11}}%
\TextAt{21}{271.893}{103.006}{\Character{27}}%
\TextAt{3}{277.018}{103.006}{\Character{0}}%
\TextAt{17}{283.605}{103.006}{1}%
\TextAt{9}{291.658}{98.07}{\Character{61}}%
\TextAt{9}{304.067}{98.07}{1}%
\TextAt{8}{313.834}{98.07}{for}%
\TextAt{8}{331.719}{98.07}{all}%
\TextAt{13}{347.978}{98.07}{\Character{27}}%
\NativeMathAt{358.381}{98.07}{11.955}{\in}%
\TextAt{13}{369.672}{98.07}{G}%
\TextAt{9}{378.905}{98.07}{\Character{40}}%
\TextAt{13}{383.458}{98.07}{L}%
\TextAt{16}{394.743}{98.07}{\Character{106}}%
\TextAt{13}{401.385}{98.07}{K}%
\TextAt{9}{412.202}{98.07}{\Character{41}}%
\NativeArrowAt{235.812}{78.035}{11.955}{21.926}{\Longleftrightarrow}%
\TextAt{13}{264.371}{78.035}{\Character{11}}%
\NativeMathAt{275.214}{78.035}{11.955}{\in}%
\TextAt{13}{286.505}{78.035}{A}%
\TextAt{21}{295.28}{76.242}{K}%
\NativeArrowAt{310.124}{78.035}{11.955}{21.926}{\Longleftrightarrow}%
\TextAt{13}{338.684}{78.035}{a}%
\TextAt{9}{348.15}{78.035}{\Character{61}}%
\TextAt{13}{360.575}{78.035}{\Character{11}}%
\TextAt{21}{368.097}{82.971}{\Character{125}}%
\NativeMathAt{377.283}{78.035}{11.955}{\in}%
\TextAt{13}{388.574}{78.035}{A}%
\TextAt{21}{397.35}{83.703}{\Character{125}}%
\TextAt{21}{397.35}{74.794}{K}%
\TextAt{13}{405.552}{78.035}{\Character{59}}%
\TextAt{8}{297.211}{43.033}{255}%
\Destination{78.23}{809.815}{nav.page.255}%
\Destination{79.36}{623.155}{nav.print.279}%
\Destination{79.36}{712.355}{nav.number.4.3.2}%
\Destination{79.36}{397.255}{nav.number.4.3.3}%
\Hotlink{407.193}{196.655}{25.078}{11.047}{nav.number.4.3.2}%
\Hotlink{275.884}{486.6}{21.722}{11.047}{nav.number.4.3.1}%
}

% Page 268 (body page 256)
\Page{%
\TextAt{2}{79.226}{796.943}{IV.}%
\TextAt{2}{98.012}{796.943}{Abstract}%
\TextAt{2}{143.252}{796.943}{Class}%
\TextAt{2}{171.844}{796.943}{Field}%
\TextAt{2}{199.575}{796.943}{Theory}%
\Vector{q 0 G 0 g 0.299 w 7.226 18.729 m 460.769 18.729 l S Q}%
\TextAt{8}{79.226}{761.124}{it}%
\TextAt{8}{90.887}{761.124}{has}%
\TextAt{8}{111.717}{761.124}{the}%
\TextAt{8}{131.83}{761.124}{kernel}%
\TextAt{13}{166.253}{761.124}{A}%
\TextAt{21}{175.028}{766.791}{\Character{125}}%
\TextAt{21}{175.028}{757.883}{K}%
\TextAt{8}{183.231}{761.124}{.}%
\TextAt{8}{191.665}{761.124}{To}%
\TextAt{8}{208.849}{761.124}{prove}%
\TextAt{8}{240.355}{761.124}{the}%
\TextAt{8}{260.468}{761.124}{surjectivity,}%
\TextAt{8}{323.903}{761.124}{we}%
\TextAt{8}{341.086}{761.124}{let}%
\TextAt{13}{357.948}{761.124}{\Character{31}}%
\NativeMathAt{368.602}{761.124}{11.955}{\in}%
\TextAt{8}{379.893}{761.124}{Hom}%
\TextAt{9}{404.281}{761.124}{\Character{40}}%
\TextAt{13}{408.834}{761.124}{G}%
\TextAt{9}{418.067}{761.124}{\Character{40}}%
\TextAt{13}{422.62}{761.124}{L}%
\TextAt{16}{433.905}{761.124}{\Character{106}}%
\TextAt{13}{440.547}{761.124}{K}%
\TextAt{9}{451.364}{761.124}{\Character{41}}%
\TextAt{13}{455.917}{761.124}{\Character{59}}%
\TextAt{13}{461.141}{761.124}{\Character{22}}%
\TextAt{21}{468.204}{759.33}{\Character{125}}%
\TextAt{9}{474.07}{761.124}{\Character{41}}%
\TextAt{8}{478.622}{761.124}{.}%
\TextAt{13}{487.056}{761.124}{\Character{31}}%
\TextAt{8}{498.243}{761.124}{defines}%
\TextAt{8}{79.226}{744.078}{a}%
\TextAt{8}{88.982}{744.078}{cyclic}%
\TextAt{8}{121.165}{744.078}{extension}%
\TextAt{13}{172.931}{744.078}{M}%
\TextAt{16}{188.825}{744.078}{\Character{106}}%
\TextAt{13}{195.467}{744.078}{K}%
\TextAt{8}{210.182}{744.078}{and}%
\TextAt{8}{232.945}{744.078}{is}%
\TextAt{8}{244.709}{744.078}{the}%
\TextAt{8}{264.865}{744.078}{composite}%
\TextAt{8}{319.883}{744.078}{of}%
\TextAt{8}{333.213}{744.078}{homomorphisms}%
\TextAt{13}{226.28}{719.602}{G}%
\TextAt{9}{235.514}{719.602}{\Character{40}}%
\TextAt{13}{240.066}{719.602}{L}%
\TextAt{16}{251.351}{719.602}{\Character{106}}%
\TextAt{13}{257.993}{719.602}{K}%
\TextAt{9}{268.811}{719.602}{\Character{41}}%
\NativeArrowAt{276.684}{719.602}{11.955}{19.26}{\longrightarrow}%
\TextAt{13}{299.266}{719.602}{G}%
\TextAt{9}{308.5}{719.602}{\Character{40}}%
\TextAt{13}{313.052}{719.602}{M}%
\TextAt{16}{328.946}{719.602}{\Character{106}}%
\TextAt{13}{335.588}{719.602}{K}%
\TextAt{9}{346.406}{719.602}{\Character{41}}%
\TextAt{17}{357.904}{726.866}{\Character{22}}%
\TextAt{21}{356.922}{726.866}{\Character{31}}%
\NativeArrowAt{354.279}{719.602}{11.955}{11.955}{\longrightarrow}%
\TextAt{13}{369.555}{719.602}{\Character{22}}%
\TextAt{21}{376.598}{717.808}{\Character{125}}%
\TextAt{13}{382.464}{719.602}{\Character{58}}%
\TextAt{8}{79.226}{695.125}{Let}%
\TextAt{13}{100.454}{695.125}{\Character{27}}%
\TextAt{8}{111.693}{695.125}{be}%
\TextAt{8}{127.876}{695.125}{a}%
\TextAt{8}{137.89}{695.125}{generator}%
\TextAt{8}{190.189}{695.125}{of}%
\TextAt{13}{203.778}{695.125}{G}%
\TextAt{9}{213.012}{695.125}{\Character{40}}%
\TextAt{13}{217.564}{695.125}{M}%
\TextAt{16}{233.899}{695.125}{\Character{106}}%
\TextAt{13}{240.982}{695.125}{K}%
\TextAt{9}{251.8}{695.125}{\Character{41}}%
\TextAt{8}{256.352}{695.125}{.}%
\TextAt{8}{265.576}{695.125}{Since}%
\TextAt{13}{296.393}{695.125}{N}%
\TextAt{21}{305.766}{693.266}{M}%
\TextAt{3}{314.748}{693.266}{\Character{106}}%
\TextAt{21}{317.1}{693.266}{K}%
\TextAt{9}{325.303}{695.125}{\Character{40}}%
\TextAt{9}{331.233}{695.125}{\Character{22}}%
\TextAt{13}{329.855}{695.125}{\Character{31}}%
\TextAt{9}{337.189}{695.125}{\Character{40}}%
\TextAt{13}{341.741}{695.125}{\Character{27}}%
\TextAt{9}{348.824}{695.125}{\Character{41}\Character{41}}%
\TextAt{9}{361.664}{695.125}{\Character{61}}%
\TextAt{9}{375.915}{695.125}{\Character{22}}%
\TextAt{13}{374.556}{695.125}{\Character{31}}%
\TextAt{9}{381.89}{695.125}{\Character{40}}%
\TextAt{13}{386.442}{695.125}{\Character{27}}%
\TextAt{9}{393.525}{695.125}{\Character{41}}%
\TextAt{17}{398.077}{699.464}{\Character{91}}%
\TextAt{21}{400.429}{699.464}{M}%
\TextAt{17}{409.411}{699.464}{\Character{58}}%
\TextAt{21}{411.763}{699.464}{K}%
\TextAt{17}{419.468}{699.464}{\Character{93}}%
\TextAt{9}{426.08}{695.125}{\Character{61}}%
\TextAt{9}{438.932}{695.125}{1}%
\TextAt{8}{444.799}{695.125}{,}%
\TextAt{8}{452.272}{695.125}{we}%
\TextAt{8}{469.758}{695.125}{deduce}%
\TextAt{8}{509.027}{695.125}{from}%
\TextAt{8}{79.226}{678.079}{(3.1)}%
\TextAt{8}{107.333}{678.079}{that}%
\TextAt{9}{134.225}{678.079}{\Character{22}}%
\TextAt{13}{132.834}{678.079}{\Character{31}}%
\TextAt{9}{140.168}{678.079}{\Character{40}}%
\TextAt{13}{144.72}{678.079}{\Character{27}}%
\TextAt{9}{151.802}{678.079}{\Character{41}}%
\TextAt{9}{159.884}{678.079}{\Character{61}}%
\TextAt{13}{172.551}{678.079}{\Character{11}}%
\TextAt{21}{180.073}{682.418}{\Character{27}}%
\TextAt{3}{185.198}{682.418}{\Character{0}}%
\TextAt{17}{191.785}{682.418}{1}%
\TextAt{8}{200.546}{678.079}{for}%
\TextAt{8}{218.563}{678.079}{some}%
\TextAt{13}{248.021}{678.079}{\Character{11}}%
\NativeMathAt{259.089}{678.079}{11.955}{\in}%
\TextAt{13}{270.605}{678.079}{A}%
\TextAt{21}{279.38}{676.286}{M}%
\TextAt{8}{288.86}{678.079}{.}%
\TextAt{8}{297.705}{678.079}{Now,}%
\TextAt{9}{327.781}{678.079}{\Character{40}}%
\TextAt{13}{332.334}{678.079}{\Character{11}}%
\TextAt{21}{339.855}{682.418}{\Character{125}}%
\TextAt{9}{345.721}{678.079}{\Character{41}}%
\TextAt{21}{350.273}{682.418}{\Character{27}}%
\TextAt{3}{355.399}{682.418}{\Character{0}}%
\TextAt{17}{361.985}{682.418}{1}%
\TextAt{9}{370.263}{678.079}{\Character{61}}%
\TextAt{9}{382.9}{678.079}{\Character{40}}%
\TextAt{13}{387.466}{678.079}{\Character{11}}%
\TextAt{21}{394.988}{682.418}{\Character{27}}%
\TextAt{3}{400.113}{682.418}{\Character{0}}%
\TextAt{17}{406.7}{682.418}{1}%
\TextAt{9}{411.432}{678.079}{\Character{41}}%
\TextAt{21}{415.984}{682.418}{\Character{125}}%
\TextAt{9}{425.396}{678.079}{\Character{61}}%
\TextAt{9}{439.419}{678.079}{\Character{22}}%
\TextAt{13}{438.046}{678.079}{\Character{31}}%
\TextAt{9}{445.38}{678.079}{\Character{40}}%
\TextAt{13}{449.932}{678.079}{\Character{27}}%
\TextAt{9}{457.014}{678.079}{\Character{41}}%
\TextAt{21}{461.567}{682.418}{\Character{125}}%
\TextAt{9}{470.978}{678.079}{\Character{61}}%
\TextAt{9}{483.615}{678.079}{1}%
\TextAt{8}{489.482}{678.079}{,}%
\TextAt{8}{496.796}{678.079}{so}%
\TextAt{8}{511.298}{678.079}{that}%
\TextAt{13}{79.226}{661.034}{a}%
\TextAt{9}{89.201}{661.034}{\Character{61}}%
\TextAt{13}{102.136}{661.034}{\Character{11}}%
\TextAt{21}{109.658}{665.372}{\Character{125}}%
\NativeMathAt{119.353}{661.034}{11.955}{\in}%
\TextAt{13}{131.154}{661.034}{A}%
\TextAt{21}{139.929}{666.701}{\Character{125}}%
\TextAt{21}{139.929}{657.792}{L}%
\TextAt{16}{149.046}{661.034}{\Character{92}}%
\TextAt{13}{159.877}{661.034}{A}%
\TextAt{21}{168.652}{659.24}{K}%
\TextAt{9}{180.685}{661.034}{\Character{61}}%
\TextAt{9}{193.609}{661.034}{\Character{1}}%
\TextAt{8}{203.375}{661.034}{.}%
\TextAt{8}{212.72}{661.034}{For}%
\TextAt{13}{233.988}{661.034}{\Character{28}}%
\NativeMathAt{244.237}{661.034}{11.955}{\in}%
\TextAt{13}{256.037}{661.034}{G}%
\TextAt{9}{265.271}{661.034}{\Character{40}}%
\TextAt{13}{269.823}{661.034}{L}%
\TextAt{16}{281.618}{661.034}{\Character{106}}%
\TextAt{13}{288.769}{661.034}{K}%
\TextAt{9}{299.586}{661.034}{\Character{41}}%
\TextAt{8}{304.139}{661.034}{,}%
\TextAt{8}{311.661}{661.034}{one}%
\TextAt{8}{333.42}{661.034}{has}%
\TextAt{13}{354.592}{661.034}{\Character{31}}%
\TextAt{9}{361.926}{661.034}{\Character{40}}%
\TextAt{13}{366.478}{661.034}{\Character{28}}%
\TextAt{9}{372.897}{661.034}{\Character{41}}%
\TextAt{9}{381.266}{661.034}{\Character{61}}%
\TextAt{9}{395.588}{661.034}{\Character{22}}%
\TextAt{13}{394.214}{661.034}{\Character{31}}%
\TextAt{9}{401.548}{661.034}{\Character{40}}%
\TextAt{13}{406.1}{661.034}{\Character{28}}%
\TextAt{16}{412.519}{661.034}{\Character{106}}%
\TextAt{21}{415.84}{659.24}{M}%
\TextAt{9}{425.32}{661.034}{\Character{41}}%
\TextAt{9}{433.689}{661.034}{\Character{61}}%
\TextAt{13}{446.637}{661.034}{\Character{11}}%
\TextAt{21}{454.159}{665.372}{\Character{28}}%
\TextAt{3}{458.837}{665.372}{\Character{0}}%
\TextAt{17}{465.424}{665.372}{1}%
\TextAt{9}{473.986}{661.034}{\Character{61}}%
\TextAt{13}{486.921}{661.034}{\Character{31}}%
\TextAt{21}{494.254}{659.24}{a}%
\TextAt{9}{499.25}{661.034}{\Character{40}}%
\TextAt{13}{503.803}{661.034}{\Character{28}}%
\TextAt{9}{510.222}{661.034}{\Character{41}}%
\TextAt{8}{514.774}{661.034}{,}%
\TextAt{8}{522.297}{661.034}{so}%
\TextAt{8}{79.226}{643.988}{that}%
\TextAt{13}{104.595}{643.988}{\Character{31}}%
\TextAt{9}{115.249}{643.988}{\Character{61}}%
\TextAt{13}{127.675}{643.988}{\Character{31}}%
\TextAt{21}{135.008}{642.194}{a}%
\TextAt{8}{140.005}{643.988}{.}%
\TextAt{8}{148.453}{643.988}{This}%
\TextAt{8}{175.173}{643.988}{proves}%
\TextAt{8}{211.337}{643.988}{the}%
\TextAt{8}{231.494}{643.988}{surjectivity,}%
\TextAt{8}{294.964}{643.988}{and}%
\TextAt{8}{317.726}{643.988}{we}%
\TextAt{8}{334.954}{643.988}{obtain}%
\TextAt{8}{371.381}{643.988}{an}%
\TextAt{8}{387.64}{643.988}{isomorphism}%
\TextAt{9}{230.744}{619.511}{\Character{1}\Character{47}}%
\TextAt{13}{246.352}{619.511}{A}%
\TextAt{21}{255.128}{625.179}{\Character{125}}%
\TextAt{21}{255.128}{616.27}{K}%
\NativeMathAt{266.651}{619.511}{11.955}{\cong}%
\TextAt{8}{279.271}{619.511}{Hom}%
\TextAt{9}{303.659}{619.511}{\Character{40}}%
\TextAt{13}{308.211}{619.511}{G}%
\TextAt{9}{317.445}{619.511}{\Character{40}}%
\TextAt{13}{321.997}{619.511}{L}%
\TextAt{16}{333.283}{619.511}{\Character{106}}%
\TextAt{13}{339.924}{619.511}{K}%
\TextAt{9}{350.742}{619.511}{\Character{41}}%
\TextAt{13}{355.294}{619.511}{\Character{59}}%
\TextAt{13}{360.518}{619.511}{\Character{22}}%
\TextAt{21}{367.582}{617.718}{\Character{125}}%
\TextAt{9}{373.447}{619.511}{\Character{41}}%
\TextAt{13}{377.999}{619.511}{\Character{58}}%
\TextAt{8}{79.226}{595.035}{If}%
\TextAt{9}{91.447}{595.035}{\Character{1}}%
\TextAt{8}{105.628}{595.035}{is}%
\TextAt{8}{117.92}{595.035}{any}%
\TextAt{8}{140.566}{595.035}{group}%
\TextAt{8}{174.27}{595.035}{between}%
\TextAt{13}{219.989}{595.035}{A}%
\TextAt{21}{228.764}{600.703}{\Character{125}}%
\TextAt{21}{228.764}{591.794}{K}%
\TextAt{8}{241.393}{595.035}{and}%
\TextAt{13}{264.684}{595.035}{A}%
\TextAt{21}{273.459}{593.242}{K}%
\TextAt{8}{286.088}{595.035}{and}%
\TextAt{8}{309.379}{595.035}{if}%
\TextAt{13}{320.631}{595.035}{L}%
\TextAt{9}{332.817}{595.035}{\Character{61}}%
\TextAt{13}{346.143}{595.035}{K}%
\TextAt{9}{356.961}{595.035}{\Character{40}}%
\TextAt{13}{361.513}{595.035}{\Character{125}}%
\TextAt{3}{368.966}{599.374}{\Character{0}}%
\TextAt{17}{375.552}{599.374}{1}%
\TextAt{9}{380.285}{595.035}{\Character{40}\Character{1}\Character{41}\Character{41}}%
\TextAt{8}{403.697}{595.035}{,}%
\TextAt{8}{411.506}{595.035}{then}%
\TextAt{9}{438.695}{595.035}{\Character{1}}%
\TextAt{9}{452.659}{595.035}{\Character{61}}%
\TextAt{13}{465.997}{595.035}{A}%
\TextAt{21}{474.773}{600.703}{\Character{125}}%
\TextAt{21}{474.773}{591.794}{L}%
\TextAt{16}{484.047}{595.035}{\Character{92}}%
\TextAt{13}{495.034}{595.035}{A}%
\TextAt{21}{503.809}{593.242}{K}%
\TextAt{8}{512.012}{595.035}{.}%
\TextAt{8}{522.046}{595.035}{In}%
\TextAt{8}{79.226}{577.989}{fact,}%
\TextAt{8}{105.563}{577.989}{putting}%
\TextAt{9}{147.191}{577.989}{\Character{1}}%
\TextAt{3}{156.946}{582.328}{\Character{48}}%
\TextAt{9}{163.062}{577.989}{\Character{61}}%
\TextAt{13}{175.488}{577.989}{A}%
\TextAt{21}{184.263}{583.657}{\Character{125}}%
\TextAt{21}{184.263}{574.748}{L}%
\TextAt{16}{193.177}{577.989}{\Character{92}}%
\TextAt{13}{203.804}{577.989}{A}%
\TextAt{21}{212.579}{576.196}{K}%
\TextAt{8}{220.782}{577.989}{,}%
\TextAt{8}{227.931}{577.989}{we}%
\TextAt{8}{245.158}{577.989}{have}%
\TextAt{8}{272.153}{577.989}{just}%
\TextAt{8}{295.298}{577.989}{seen}%
\TextAt{8}{320.715}{577.989}{that}%
\TextAt{8}{346.084}{577.989}{one}%
\TextAt{8}{367.543}{577.989}{has}%
\TextAt{9}{229.347}{553.513}{\Character{1}}%
\TextAt{3}{239.102}{558.449}{\Character{48}}%
\TextAt{9}{241.897}{553.513}{\Character{47}}%
\TextAt{13}{247.75}{553.513}{A}%
\TextAt{21}{256.525}{559.181}{\Character{125}}%
\TextAt{21}{256.525}{550.272}{K}%
\NativeMathAt{268.049}{553.513}{11.955}{\cong}%
\TextAt{8}{280.668}{553.513}{Hom}%
\TextAt{9}{305.057}{553.513}{\Character{40}}%
\TextAt{13}{309.609}{553.513}{G}%
\TextAt{9}{318.843}{553.513}{\Character{40}}%
\TextAt{13}{323.395}{553.513}{L}%
\TextAt{16}{334.68}{553.513}{\Character{106}}%
\TextAt{13}{341.322}{553.513}{K}%
\TextAt{9}{352.14}{553.513}{\Character{41}}%
\TextAt{13}{356.692}{553.513}{\Character{59}}%
\TextAt{13}{361.916}{553.513}{\Character{22}}%
\TextAt{21}{368.979}{551.72}{\Character{125}}%
\TextAt{9}{374.845}{553.513}{\Character{41}}%
\TextAt{13}{379.397}{553.513}{\Character{58}}%
\TextAt{8}{79.226}{529.037}{The}%
\TextAt{8}{103.499}{529.037}{subgroup}%
\TextAt{9}{154.516}{529.037}{\Character{1}\Character{47}}%
\TextAt{13}{170.124}{529.037}{A}%
\TextAt{21}{178.899}{534.705}{\Character{125}}%
\TextAt{21}{178.899}{525.796}{K}%
\TextAt{8}{191.219}{529.037}{corresponds}%
\TextAt{8}{255.625}{529.037}{under}%
\TextAt{8}{289.008}{529.037}{Pontryagin}%
\TextAt{8}{349.242}{529.037}{duality}%
\TextAt{8}{389.141}{529.037}{to}%
\TextAt{8}{403.67}{529.037}{the}%
\TextAt{8}{424.046}{529.037}{subgroup}%
\TextAt{8}{475.063}{529.037}{Hom}%
\TextAt{9}{499.451}{529.037}{\Character{40}}%
\TextAt{13}{504.004}{529.037}{G}%
\TextAt{9}{513.237}{529.037}{\Character{40}}%
\TextAt{13}{517.79}{529.037}{L}%
\TextAt{16}{529.449}{529.037}{\Character{106}}%
\TextAt{13}{79.226}{511.991}{K}%
\TextAt{9}{90.044}{511.991}{\Character{41}\Character{47}}%
\TextAt{13}{100.449}{511.991}{H}%
\TextAt{13}{110.408}{511.991}{\Character{59}}%
\TextAt{13}{115.644}{511.991}{\Character{22}}%
\TextAt{21}{122.709}{510.198}{\Character{125}}%
\TextAt{9}{128.574}{511.991}{\Character{41}}%
\TextAt{8}{133.126}{511.991}{,}%
\TextAt{8}{140.276}{511.991}{where}%
\TextAt{13}{186.505}{487.515}{H}%
\TextAt{9}{200.449}{487.515}{\Character{61}}%
\TextAt{16}{212.874}{487.515}{\Character{102}}%
\TextAt{13}{218.852}{487.515}{\Character{27}}%
\NativeMathAt{229.255}{487.515}{11.955}{\in}%
\TextAt{13}{240.546}{487.515}{G}%
\TextAt{9}{249.78}{487.515}{\Character{40}}%
\TextAt{13}{254.332}{487.515}{L}%
\TextAt{16}{265.617}{487.515}{\Character{106}}%
\TextAt{13}{272.259}{487.515}{K}%
\TextAt{9}{283.077}{487.515}{\Character{41}}%
\TextAt{16}{290.95}{487.515}{\Character{106}}%
\TextAt{13}{297.592}{487.515}{\Character{31}}%
\TextAt{21}{304.925}{485.722}{a}%
\TextAt{9}{309.921}{487.515}{\Character{40}}%
\TextAt{13}{314.474}{487.515}{\Character{27}}%
\TextAt{9}{321.556}{487.515}{\Character{41}}%
\TextAt{9}{329.411}{487.515}{\Character{61}}%
\TextAt{9}{341.832}{487.515}{1}%
\TextAt{8}{351.605}{487.515}{for}%
\TextAt{8}{369.49}{487.515}{all}%
\TextAt{13}{385.749}{487.515}{a}%
\NativeMathAt{395.215}{487.515}{11.955}{\in}%
\TextAt{9}{406.506}{487.515}{\Character{1}}%
\TextAt{16}{416.261}{487.515}{\Character{103}}%
\TextAt{13}{422.239}{487.515}{\Character{58}}%
\TextAt{8}{79.226}{463.039}{As}%
\TextAt{13}{98.958}{463.039}{\Character{11}}%
\TextAt{21}{106.48}{467.377}{\Character{27}}%
\TextAt{3}{111.605}{467.377}{\Character{0}}%
\TextAt{17}{118.191}{467.377}{1}%
\TextAt{9}{130.41}{463.039}{\Character{61}}%
\TextAt{13}{147.001}{463.039}{\Character{31}}%
\TextAt{21}{154.335}{461.246}{a}%
\TextAt{9}{159.331}{463.039}{\Character{40}}%
\TextAt{13}{163.883}{463.039}{\Character{27}}%
\TextAt{9}{170.966}{463.039}{\Character{41}}%
\TextAt{8}{181.86}{463.039}{for}%
\TextAt{13}{202.189}{463.039}{\Character{11}}%
\NativeMathAt{217.198}{463.039}{11.955}{\in}%
\TextAt{13}{232.654}{463.039}{\Character{125}}%
\TextAt{3}{240.107}{467.377}{\Character{0}}%
\TextAt{17}{246.693}{467.377}{1}%
\TextAt{9}{251.426}{463.039}{\Character{40}}%
\TextAt{13}{255.978}{463.039}{a}%
\TextAt{9}{262.123}{463.039}{\Character{41}}%
\TextAt{8}{266.675}{463.039}{,}%
\TextAt{13}{276.88}{463.039}{H}%
\TextAt{8}{293.844}{463.039}{leaves}%
\TextAt{8}{329.847}{463.039}{fixed}%
\TextAt{8}{360.577}{463.039}{the}%
\TextAt{8}{383.178}{463.039}{elements}%
\TextAt{8}{433.479}{463.039}{of}%
\TextAt{13}{449.254}{463.039}{\Character{125}}%
\TextAt{3}{456.706}{467.377}{\Character{0}}%
\TextAt{17}{463.293}{467.377}{1}%
\TextAt{9}{468.025}{463.039}{\Character{40}\Character{1}\Character{41}}%
\TextAt{8}{486.885}{463.039}{,}%
\TextAt{8}{497.09}{463.039}{and}%
\TextAt{8}{522.297}{463.039}{as}%
\TextAt{13}{79.226}{445.993}{K}%
\TextAt{9}{90.044}{445.993}{\Character{40}}%
\TextAt{13}{94.596}{445.993}{\Character{125}}%
\TextAt{3}{102.049}{450.331}{\Character{0}}%
\TextAt{17}{108.635}{450.331}{1}%
\TextAt{9}{113.368}{445.993}{\Character{40}\Character{1}\Character{41}\Character{41}}%
\TextAt{9}{142.467}{445.993}{\Character{61}}%
\TextAt{13}{157.359}{445.993}{L}%
\TextAt{8}{165.323}{445.993}{,}%
\TextAt{8}{174.245}{445.993}{we}%
\TextAt{8}{192.89}{445.993}{find}%
\TextAt{8}{217.716}{445.993}{that}%
\TextAt{13}{244.503}{445.993}{H}%
\TextAt{9}{260.863}{445.993}{\Character{61}}%
\TextAt{9}{275.687}{445.993}{1}%
\TextAt{8}{281.558}{445.993}{,}%
\TextAt{8}{290.479}{445.993}{so}%
\TextAt{8}{306.267}{445.993}{that}%
\TextAt{8}{333.054}{445.993}{Hom}%
\TextAt{9}{357.443}{445.993}{\Character{40}}%
\TextAt{13}{361.995}{445.993}{G}%
\TextAt{9}{371.229}{445.993}{\Character{40}}%
\TextAt{13}{375.781}{445.993}{L}%
\TextAt{16}{389.483}{445.993}{\Character{106}}%
\TextAt{13}{398.541}{445.993}{K}%
\TextAt{9}{409.358}{445.993}{\Character{41}\Character{47}}%
\TextAt{13}{419.764}{445.993}{H}%
\TextAt{13}{429.723}{445.993}{\Character{59}}%
\TextAt{13}{434.959}{445.993}{\Character{22}}%
\TextAt{21}{442.023}{444.2}{\Character{125}}%
\TextAt{9}{447.889}{445.993}{\Character{41}}%
\TextAt{9}{458.159}{445.993}{\Character{61}}%
\TextAt{8}{473.02}{445.993}{Hom}%
\TextAt{9}{497.409}{445.993}{\Character{40}}%
\TextAt{13}{501.961}{445.993}{G}%
\TextAt{9}{511.195}{445.993}{\Character{40}}%
\TextAt{13}{515.747}{445.993}{L}%
\TextAt{16}{529.449}{445.993}{\Character{106}}%
\TextAt{13}{79.226}{428.947}{K}%
\TextAt{9}{90.044}{428.947}{\Character{41}}%
\TextAt{13}{94.596}{428.947}{\Character{59}}%
\TextAt{13}{99.82}{428.947}{\Character{22}}%
\TextAt{21}{106.883}{427.154}{\Character{125}}%
\TextAt{9}{112.749}{428.947}{\Character{41}}%
\TextAt{8}{117.301}{428.947}{.}%
\TextAt{8}{125.75}{428.947}{It}%
\TextAt{8}{138.422}{428.947}{follows}%
\TextAt{8}{176.858}{428.947}{that}%
\TextAt{9}{202.227}{428.947}{\Character{1}\Character{47}}%
\TextAt{13}{217.835}{428.947}{A}%
\TextAt{21}{226.61}{434.614}{\Character{125}}%
\TextAt{21}{226.61}{425.706}{K}%
\TextAt{9}{238.134}{428.947}{\Character{61}}%
\TextAt{9}{250.544}{428.947}{\Character{1}}%
\TextAt{3}{260.314}{433.285}{\Character{48}}%
\TextAt{9}{263.109}{428.947}{\Character{47}}%
\TextAt{13}{268.962}{428.947}{A}%
\TextAt{21}{277.737}{434.614}{\Character{125}}%
\TextAt{21}{277.737}{425.706}{K}%
\TextAt{8}{285.94}{428.947}{,}%
\TextAt{8}{293.089}{428.947}{i.e.,}%
\TextAt{9}{315.194}{428.947}{\Character{1}}%
\TextAt{9}{328.249}{428.947}{\Character{61}}%
\TextAt{9}{340.671}{428.947}{\Character{1}}%
\TextAt{3}{350.451}{433.285}{\Character{48}}%
\TextAt{8}{353.246}{428.947}{.}%
\TextAt{8}{97.159}{408.912}{It}%
\TextAt{8}{112.511}{408.912}{is}%
\TextAt{8}{126.955}{408.912}{therefore}%
\TextAt{8}{178.735}{408.912}{clear}%
\TextAt{8}{209.378}{408.912}{that}%
\TextAt{8}{237.427}{408.912}{the}%
\TextAt{8}{260.264}{408.912}{correspondence}%
\TextAt{9}{344.621}{408.912}{\Character{1}}%
\NativeArrowAt{362.264}{408.912}{11.955}{19.26}{\longmapsto}%
\TextAt{13}{389.414}{408.912}{L}%
\TextAt{9}{405.266}{408.912}{\Character{61}}%
\TextAt{13}{422.259}{408.912}{K}%
\TextAt{9}{433.076}{408.912}{\Character{40}}%
\TextAt{13}{437.629}{408.912}{\Character{125}}%
\TextAt{3}{445.081}{413.251}{\Character{0}}%
\TextAt{17}{451.668}{413.251}{1}%
\TextAt{9}{456.4}{408.912}{\Character{40}\Character{1}\Character{41}\Character{41}}%
\TextAt{8}{486.389}{408.912}{is}%
\TextAt{8}{500.833}{408.912}{a}%
\TextAt{9}{513.269}{408.912}{1}%
\TextAt{8}{519.122}{408.912}{-}%
\TextAt{9}{523.019}{408.912}{1}%
\TextAt{8}{528.872}{408.912}{-}%
\TextAt{8}{79.226}{391.866}{correspondence,}%
\TextAt{8}{164.156}{391.866}{as}%
\TextAt{8}{178.526}{391.866}{claimed.}%
\TextAt{8}{225.996}{391.866}{This}%
\TextAt{8}{252.716}{391.866}{finishes}%
\TextAt{8}{293.805}{391.866}{the}%
\TextAt{8}{313.962}{391.866}{proof}%
\TextAt{8}{344.531}{391.866}{of}%
\TextAt{8}{357.861}{391.866}{the}%
\TextAt{8}{378.018}{391.866}{theorem.}%
\Vector{q 0 G 0 g 0.399 w 452.334 -382.134 m 452.334 -374.064 l S Q}%
\Vector{q 0 G 0 g 0.399 w 452.534 -374.263 m 459.707 -374.263 l S Q}%
\Vector{q 0 G 0 g 0.399 w 452.534 -381.935 m 459.707 -381.935 l S Q}%
\Vector{q 0 G 0 g 0.399 w 459.906 -382.134 m 459.906 -374.064 l S Q}%
\TextAt{12}{97.159}{364.158}{Remarks}%
\TextAt{12}{153.293}{364.158}{and}%
\TextAt{12}{179.085}{364.158}{Examples:}%
\TextAt{8}{244.28}{364.158}{1)}%
\TextAt{8}{258.439}{364.158}{If}%
\TextAt{13}{269.979}{364.158}{L}%
\TextAt{16}{281.265}{364.158}{\Character{106}}%
\TextAt{13}{287.906}{364.158}{K}%
\TextAt{8}{302.47}{364.158}{is}%
\TextAt{8}{314.082}{364.158}{infinite,}%
\TextAt{8}{356.879}{364.158}{then}%
\TextAt{8}{383.388}{364.158}{Hom}%
\TextAt{9}{407.776}{364.158}{\Character{40}}%
\TextAt{13}{412.328}{364.158}{G}%
\TextAt{9}{421.562}{364.158}{\Character{40}}%
\TextAt{13}{426.114}{364.158}{L}%
\TextAt{16}{437.4}{364.158}{\Character{106}}%
\TextAt{13}{444.041}{364.158}{K}%
\TextAt{9}{454.859}{364.158}{\Character{41}}%
\TextAt{13}{459.411}{364.158}{\Character{59}}%
\TextAt{13}{464.635}{364.158}{\Character{22}}%
\TextAt{21}{471.699}{362.365}{\Character{125}}%
\TextAt{9}{477.564}{364.158}{\Character{41}}%
\TextAt{8}{485.862}{364.158}{has}%
\TextAt{8}{506.584}{364.158}{to}%
\TextAt{8}{520.743}{364.158}{be}%
\TextAt{8}{79.226}{347.112}{interpreted}%
\TextAt{8}{140.086}{347.112}{as}%
\TextAt{8}{155.158}{347.112}{the}%
\TextAt{8}{176.016}{347.112}{group}%
\TextAt{8}{209.894}{347.112}{of}%
\TextAt{8}{223.925}{347.112}{all}%
\TextAt{19}{240.886}{347.112}{continuous}%
\TextAt{8}{300.312}{347.112}{homomorphisms}%
\TextAt{13}{388.298}{347.112}{\Character{31}}%
\TextAt{9}{400.148}{347.112}{\Character{58}}%
\TextAt{13}{407.916}{347.112}{G}%
\TextAt{9}{417.15}{347.112}{\Character{40}}%
\TextAt{13}{421.702}{347.112}{L}%
\TextAt{16}{434.184}{347.112}{\Character{106}}%
\TextAt{13}{442.021}{347.112}{K}%
\TextAt{9}{452.839}{347.112}{\Character{41}}%
\NativeMathAt{461.908}{347.112}{11.955}{\rightarrow}%
\TextAt{13}{478.38}{347.112}{\Character{22}}%
\TextAt{21}{485.423}{345.319}{\Character{125}}%
\TextAt{8}{491.288}{347.112}{,}%
\TextAt{8}{499.314}{347.112}{i.e.,}%
\TextAt{8}{522.297}{347.112}{as}%
\TextAt{8}{79.226}{330.066}{the}%
\TextAt{8}{99.383}{330.066}{character}%
\TextAt{8}{150.443}{330.066}{group}%
\TextAt{8}{183.619}{330.066}{of}%
\TextAt{8}{196.949}{330.066}{the}%
\TextAt{8}{217.105}{330.066}{topological}%
\TextAt{8}{276.63}{330.066}{group}%
\TextAt{13}{309.806}{330.066}{G}%
\TextAt{9}{319.039}{330.066}{\Character{40}}%
\TextAt{13}{323.591}{330.066}{L}%
\TextAt{16}{334.877}{330.066}{\Character{106}}%
\TextAt{13}{341.518}{330.066}{K}%
\TextAt{9}{352.336}{330.066}{\Character{41}}%
\TextAt{8}{356.889}{330.066}{.}%
\TextAt{8}{97.159}{310.031}{2)}%
\TextAt{8}{111.793}{310.031}{The}%
\TextAt{8}{136.171}{310.031}{composite}%
\TextAt{8}{191.513}{310.031}{of}%
\TextAt{8}{205.167}{310.031}{two}%
\TextAt{8}{227.608}{310.031}{abelian}%
\TextAt{8}{268.58}{310.031}{extensions}%
\TextAt{8}{325.284}{310.031}{of}%
\TextAt{13}{338.938}{310.031}{K}%
\TextAt{8}{353.977}{310.031}{of}%
\TextAt{8}{367.631}{310.031}{exponent}%
\TextAt{13}{418.358}{310.031}{n}%
\TextAt{8}{429.567}{310.031}{is}%
\TextAt{8}{441.655}{310.031}{again}%
\TextAt{8}{473.206}{310.031}{of}%
\TextAt{8}{486.86}{310.031}{the}%
\TextAt{8}{507.341}{310.031}{same}%
\TextAt{8}{79.226}{292.985}{type,}%
\TextAt{8}{109.174}{292.985}{and}%
\TextAt{8}{132.223}{292.985}{all}%
\TextAt{8}{148.769}{292.985}{of}%
\TextAt{8}{162.386}{292.985}{them}%
\TextAt{8}{192.585}{292.985}{lie}%
\TextAt{8}{208.473}{292.985}{in}%
\TextAt{8}{222.413}{292.985}{the}%
\TextAt{8}{242.856}{292.985}{maximal}%
\TextAt{8}{290.951}{292.985}{abelian}%
\TextAt{8}{331.886}{292.985}{extension}%
\TextAt{8}{383.938}{292.985}{of}%
\TextAt{8}{397.555}{292.985}{exponent}%
\TextAt{13}{448.245}{292.985}{n}%
\TextAt{8}{455.233}{292.985}{.}%
\TextAt{8}{464.542}{292.985}{It}%
\TextAt{8}{477.501}{292.985}{is}%
\TextAt{8}{489.552}{292.985}{given}%
\TextAt{8}{520.408}{292.985}{by}%
\TextAt{20}{81.964}{278.961}{\Character{98}}%
\TextAt{13}{79.226}{275.939}{K}%
\TextAt{9}{93.365}{275.939}{\Character{61}}%
\TextAt{13}{105.79}{275.939}{K}%
\TextAt{9}{116.608}{275.939}{\Character{40}}%
\TextAt{13}{121.16}{275.939}{\Character{125}}%
\TextAt{3}{128.613}{280.278}{\Character{0}}%
\TextAt{17}{135.199}{280.278}{1}%
\TextAt{9}{139.931}{275.939}{\Character{40}}%
\TextAt{13}{144.484}{275.939}{A}%
\TextAt{21}{153.259}{274.146}{K}%
\TextAt{9}{161.462}{275.939}{\Character{41}\Character{41}}%
\TextAt{8}{170.566}{275.939}{,}%
\TextAt{8}{177.716}{275.939}{and}%
\TextAt{8}{200.478}{275.939}{for}%
\TextAt{8}{218.363}{275.939}{the}%
\TextAt{8}{238.52}{275.939}{Pontryagin}%
\TextAt{8}{298.535}{275.939}{dual}%
\TextAt{13}{156.229}{251.463}{G}%
\TextAt{9}{165.462}{251.463}{\Character{40}}%
\TextAt{20}{172.753}{254.485}{\Character{98}}%
\TextAt{13}{170.015}{251.463}{K}%
\TextAt{16}{184.153}{251.463}{\Character{106}}%
\TextAt{13}{190.795}{251.463}{K}%
\TextAt{9}{201.612}{251.463}{\Character{41}}%
\TextAt{3}{206.165}{256.399}{\Character{3}}%
\TextAt{9}{214.218}{251.463}{\Character{61}}%
\TextAt{8}{226.643}{251.463}{Hom}%
\TextAt{9}{251.032}{251.463}{\Character{40}}%
\TextAt{13}{255.584}{251.463}{G}%
\TextAt{9}{264.818}{251.463}{\Character{40}}%
\TextAt{20}{272.109}{254.485}{\Character{98}}%
\TextAt{13}{269.37}{251.463}{K}%
\TextAt{16}{283.509}{251.463}{\Character{106}}%
\TextAt{13}{290.15}{251.463}{K}%
\TextAt{9}{300.968}{251.463}{\Character{41}}%
\TextAt{13}{305.52}{251.463}{\Character{59}}%
\TextAt{18}{310.765}{251.463}{Q}%
\TextAt{9}{320.063}{251.463}{\Character{47}}%
\TextAt{18}{325.916}{251.463}{Z}%
\TextAt{9}{333.886}{251.463}{\Character{41}}%
\TextAt{9}{341.741}{251.463}{\Character{61}}%
\TextAt{8}{354.185}{251.463}{Hom}%
\TextAt{9}{378.573}{251.463}{\Character{40}}%
\TextAt{13}{383.126}{251.463}{G}%
\TextAt{9}{392.359}{251.463}{\Character{40}}%
\TextAt{20}{399.65}{254.485}{\Character{98}}%
\TextAt{13}{396.912}{251.463}{K}%
\TextAt{16}{411.05}{251.463}{\Character{106}}%
\TextAt{13}{417.692}{251.463}{K}%
\TextAt{9}{428.51}{251.463}{\Character{41}}%
\TextAt{13}{433.062}{251.463}{\Character{59}}%
\TextAt{13}{438.286}{251.463}{\Character{22}}%
\TextAt{21}{445.349}{249.67}{\Character{125}}%
\TextAt{9}{451.215}{251.463}{\Character{41}}%
\TextAt{8}{79.226}{226.987}{we}%
\TextAt{8}{96.454}{226.987}{have}%
\TextAt{8}{123.448}{226.987}{by}%
\TextAt{8}{139.707}{226.987}{(3.3)}%
\TextAt{8}{167.682}{226.987}{that}%
\TextAt{13}{249.163}{209.941}{G}%
\TextAt{9}{258.397}{209.941}{\Character{40}}%
\TextAt{20}{265.687}{212.963}{\Character{98}}%
\TextAt{13}{262.949}{209.941}{K}%
\TextAt{16}{277.088}{209.941}{\Character{106}}%
\TextAt{13}{283.729}{209.941}{K}%
\TextAt{9}{294.547}{209.941}{\Character{41}}%
\TextAt{3}{299.099}{214.877}{\Character{3}}%
\NativeMathAt{307.152}{209.941}{11.955}{\cong}%
\TextAt{13}{319.772}{209.941}{A}%
\TextAt{21}{328.547}{208.148}{K}%
\TextAt{9}{336.75}{209.941}{\Character{47}}%
\TextAt{13}{342.603}{209.941}{A}%
\TextAt{21}{351.378}{215.609}{\Character{125}}%
\TextAt{21}{351.378}{206.7}{K}%
\TextAt{13}{359.581}{209.941}{\Character{58}}%
\TextAt{8}{79.226}{188.359}{3)}%
\TextAt{8}{94.183}{188.359}{If}%
\TextAt{13}{106.521}{188.359}{k}%
\TextAt{8}{117.554}{188.359}{is}%
\TextAt{8}{129.965}{188.359}{an}%
\TextAt{8}{146.87}{188.359}{actual}%
\TextAt{8}{182.641}{188.359}{field}%
\TextAt{8}{208.644}{188.359}{of}%
\TextAt{8}{222.62}{188.359}{positive}%
\TextAt{8}{266.58}{188.359}{characteristic}%
\TextAt{13}{339.161}{188.359}{p}%
\TextAt{8}{349.58}{188.359}{and}%
\TextAt{9}{373.307}{191.513}{\Character{22}}%
\TextAt{13}{372.988}{188.359}{k}%
\TextAt{8}{384.022}{188.359}{is}%
\TextAt{8}{396.432}{188.359}{the}%
\TextAt{8}{417.235}{188.359}{separable}%
\TextAt{8}{469.324}{188.359}{closure}%
\TextAt{8}{509.052}{188.359}{of}%
\TextAt{13}{523.028}{188.359}{k}%
\TextAt{8}{529.518}{188.359}{,}%
\TextAt{8}{79.226}{171.313}{then}%
\TextAt{13}{105.886}{171.313}{A}%
\TextAt{8}{118.559}{171.313}{may}%
\TextAt{8}{143.928}{171.313}{be}%
\TextAt{8}{159.852}{171.313}{chosen}%
\TextAt{8}{197.308}{171.313}{to}%
\TextAt{8}{211.618}{171.313}{be}%
\TextAt{8}{227.542}{171.313}{the}%
\TextAt{8}{247.699}{171.313}{additive}%
\TextAt{8}{292.578}{171.313}{group}%
\TextAt{9}{326.072}{174.467}{\Character{22}}%
\TextAt{13}{325.754}{171.313}{k}%
\TextAt{8}{336.141}{171.313}{and}%
\TextAt{13}{358.904}{171.313}{\Character{125}}%
\TextAt{8}{370.253}{171.313}{to}%
\TextAt{8}{384.564}{171.313}{be}%
\TextAt{8}{400.488}{171.313}{the}%
\TextAt{8}{420.644}{171.313}{operator}%
\TextAt{13}{215.154}{146.836}{\Character{125}}%
\TextAt{9}{225.927}{146.836}{\Character{58}}%
\TextAt{9}{232.818}{149.991}{\Character{22}}%
\TextAt{13}{232.5}{146.836}{k}%
\NativeArrowAt{242.31}{146.836}{11.955}{19.26}{\longrightarrow}%
\TextAt{9}{265.21}{149.991}{\Character{22}}%
\TextAt{13}{264.892}{146.836}{k}%
\TextAt{13}{271.372}{146.836}{\Character{59}}%
\TextAt{13}{300.519}{146.836}{a}%
\NativeArrowAt{310.002}{146.836}{11.955}{19.26}{\longmapsto}%
\TextAt{13}{332.584}{146.836}{\Character{125}a}%
\TextAt{9}{349.502}{146.836}{\Character{61}}%
\TextAt{13}{361.928}{146.836}{a}%
\TextAt{21}{368.073}{151.773}{p}%
\TextAt{16}{375.49}{146.836}{\Character{0}}%
\TextAt{13}{387.445}{146.836}{a\Character{58}}%
\TextAt{8}{79.226}{122.36}{Then}%
\TextAt{8}{109.784}{122.36}{axiom}%
\TextAt{8}{144.585}{122.36}{(3.1)}%
\TextAt{8}{172.56}{122.36}{is}%
\TextAt{8}{184.324}{122.36}{indeed}%
\TextAt{8}{221.385}{122.36}{satisfied,}%
\TextAt{8}{269.636}{122.36}{for}%
\TextAt{8}{287.521}{122.36}{we}%
\TextAt{8}{304.748}{122.36}{have,}%
\TextAt{8}{334.995}{122.36}{in}%
\TextAt{8}{348.648}{122.36}{complete}%
\TextAt{8}{398.071}{122.36}{generality:}%
\TextAt{12}{79.226}{97.641}{Proposition}%
\TextAt{12}{153.423}{97.641}{(3.4).}%
\TextAt{8}{190.03}{97.641}{For}%
\TextAt{8}{212.291}{97.641}{every}%
\TextAt{8}{244.474}{97.641}{cyclic}%
\TextAt{8}{277.949}{97.641}{finite}%
\TextAt{8}{309.152}{97.641}{field}%
\TextAt{8}{335.8}{97.641}{extension}%
\TextAt{13}{388.857}{97.641}{L}%
\TextAt{16}{402.343}{97.641}{\Character{106}}%
\TextAt{13}{411.186}{97.641}{K}%
\TextAt{8}{422.004}{97.641}{,}%
\TextAt{8}{430.767}{97.641}{one}%
\TextAt{8}{453.518}{97.641}{has}%
\TextAt{13}{475.683}{97.641}{H}%
\TextAt{3}{486.306}{101.979}{\Character{0}}%
\TextAt{17}{492.892}{101.979}{1}%
\TextAt{9}{497.624}{97.641}{\Character{40}}%
\TextAt{13}{502.177}{97.641}{G}%
\TextAt{9}{511.41}{97.641}{\Character{40}}%
\TextAt{13}{515.963}{97.641}{L}%
\TextAt{16}{529.449}{97.641}{\Character{106}}%
\TextAt{13}{79.226}{80.595}{K}%
\TextAt{9}{90.044}{80.595}{\Character{41}}%
\TextAt{13}{94.596}{80.595}{\Character{59}}%
\TextAt{13}{99.82}{80.595}{L}%
\TextAt{9}{107.805}{80.595}{\Character{41}}%
\TextAt{9}{115.66}{80.595}{\Character{61}}%
\TextAt{9}{128.081}{80.595}{1}%
\TextAt{8}{133.956}{80.595}{.}%
\TextAt{8}{297.211}{43.033}{256}%
\Destination{78.23}{809.815}{nav.page.256}%
\Destination{79.36}{693.055}{nav.print.280}%
\Destination{79.36}{111.655}{nav.number.4.3.4}%
\Hotlink{139.207}{224.755}{25.078}{11.047}{nav.number.4.3.3}%
\Hotlink{80.397}{675.9}{21.722}{11.047}{nav.number.4.3.1}%
\Hotlink{145.755}{120.2}{21.722}{11.047}{nav.number.4.3.1}%
}

% Page 269 (body page 257)
\Page{%
\TextAt{2}{79.226}{796.943}{IV.}%
\TextAt{2}{98.012}{796.943}{Abstract}%
\TextAt{2}{143.252}{796.943}{Class}%
\TextAt{2}{171.844}{796.943}{Field}%
\TextAt{2}{199.575}{796.943}{Theory}%
\Vector{q 0 G 0 g 0.299 w 7.226 18.729 m 460.769 18.729 l S Q}%
\TextAt{19}{79.226}{761.124}{Proof.}%
\TextAt{8}{116}{761.124}{The}%
\TextAt{8}{140.407}{761.124}{extension}%
\TextAt{13}{192.525}{761.124}{L}%
\TextAt{16}{204.412}{761.124}{\Character{106}}%
\TextAt{13}{211.654}{761.124}{K}%
\TextAt{8}{226.722}{761.124}{always}%
\TextAt{8}{264.542}{761.124}{admits}%
\TextAt{8}{303.331}{761.124}{a}%
\TextAt{8}{313.439}{761.124}{normal}%
\TextAt{8}{353.471}{761.124}{basis}%
\TextAt{16}{382.564}{761.124}{\Character{102}}%
\TextAt{13}{388.541}{761.124}{\Character{27}}%
\TextAt{13}{395.607}{761.124}{c}%
\TextAt{16}{404.584}{761.124}{\Character{106}}%
\TextAt{13}{411.826}{761.124}{\Character{27}}%
\NativeMathAt{422.831}{761.124}{11.955}{\in}%
\TextAt{13}{434.723}{761.124}{G}%
\TextAt{9}{443.956}{761.124}{\Character{40}}%
\TextAt{13}{448.508}{761.124}{L}%
\TextAt{16}{460.395}{761.124}{\Character{106}}%
\TextAt{13}{467.638}{761.124}{K}%
\TextAt{9}{478.455}{761.124}{\Character{41}}%
\TextAt{16}{483.008}{761.124}{\Character{103}}%
\TextAt{8}{488.985}{761.124}{,}%
\TextAt{8}{496.575}{761.124}{so}%
\TextAt{8}{511.298}{761.124}{that}%
\TextAt{13}{79.226}{744.078}{L}%
\TextAt{9}{90.932}{744.078}{\Character{61}}%
\TextAt{20}{103.777}{753.044}{\Character{76}}%
\TextAt{21}{117.061}{740.591}{\Character{27}}%
\TextAt{13}{124.677}{744.078}{K}%
\TextAt{13}{135.485}{744.078}{\Character{27}}%
\TextAt{13}{142.562}{744.078}{c}%
\TextAt{8}{147.615}{744.078}{.}%
\TextAt{8}{156.803}{744.078}{This}%
\TextAt{8}{183.769}{744.078}{means}%
\TextAt{8}{219.845}{744.078}{that}%
\TextAt{13}{245.461}{744.078}{L}%
\TextAt{8}{257.569}{744.078}{is}%
\TextAt{8}{269.58}{744.078}{a}%
\TextAt{13}{279.582}{744.078}{G}%
\TextAt{9}{288.815}{744.078}{\Character{40}}%
\TextAt{13}{293.368}{744.078}{L}%
\TextAt{16}{305.073}{744.078}{\Character{106}}%
\TextAt{13}{312.135}{744.078}{K}%
\TextAt{9}{322.953}{744.078}{\Character{41}}%
\TextAt{8}{327.505}{744.078}{-induced}%
\TextAt{8}{375.214}{744.078}{module}%
\TextAt{8}{416.753}{744.078}{in}%
\TextAt{8}{430.653}{744.078}{the}%
\TextAt{8}{451.056}{744.078}{sense}%
\TextAt{8}{481.334}{744.078}{of}%
\TextAt{8}{494.91}{744.078}{§7,}%
\TextAt{8}{513.904}{744.078}{and}%
\TextAt{8}{79.226}{727.032}{then}%
\TextAt{13}{105.886}{727.032}{H}%
\TextAt{3}{116.509}{731.37}{\Character{0}}%
\TextAt{17}{123.095}{731.37}{1}%
\TextAt{9}{127.828}{727.032}{\Character{40}}%
\TextAt{13}{132.38}{727.032}{G}%
\TextAt{9}{141.614}{727.032}{\Character{40}}%
\TextAt{13}{146.166}{727.032}{L}%
\TextAt{16}{157.451}{727.032}{\Character{106}}%
\TextAt{13}{164.093}{727.032}{K}%
\TextAt{9}{174.911}{727.032}{\Character{41}}%
\TextAt{13}{179.463}{727.032}{\Character{59}}%
\TextAt{13}{184.687}{727.032}{L}%
\TextAt{9}{192.672}{727.032}{\Character{41}}%
\TextAt{9}{200.527}{727.032}{\Character{61}}%
\TextAt{9}{212.948}{727.032}{1}%
\TextAt{8}{218.823}{727.032}{,}%
\TextAt{8}{225.972}{727.032}{by}%
\TextAt{8}{242.231}{727.032}{(7.4).}%
\Vector{q 0 G 0 g 0.399 w 452.334 -46.968 m 452.334 -38.898 l S Q}%
\Vector{q 0 G 0 g 0.399 w 452.534 -39.098 m 459.707 -39.098 l S Q}%
\Vector{q 0 G 0 g 0.399 w 452.534 -46.769 m 459.707 -46.769 l S Q}%
\Vector{q 0 G 0 g 0.399 w 459.906 -46.968 m 459.906 -38.898 l S Q}%
\TextAt{8}{97.159}{698.478}{The}%
\TextAt{8}{120.527}{698.478}{Kummer}%
\TextAt{8}{168.606}{698.478}{theory}%
\TextAt{8}{204.671}{698.478}{with}%
\TextAt{8}{230.645}{698.478}{respect}%
\TextAt{8}{270.009}{698.478}{to}%
\TextAt{8}{283.634}{698.478}{the}%
\TextAt{8}{303.104}{698.478}{operator}%
\TextAt{13}{349.582}{698.478}{\Character{125}a}%
\TextAt{9}{366.5}{698.478}{\Character{61}}%
\TextAt{13}{378.926}{698.478}{a}%
\TextAt{21}{385.07}{702.816}{p}%
\TextAt{16}{391.086}{698.478}{\Character{0}}%
\TextAt{13}{401.639}{698.478}{a}%
\TextAt{8}{410.995}{698.478}{is}%
\TextAt{8}{422.073}{698.478}{usually}%
\TextAt{8}{461.449}{698.478}{called}%
\TextAt{12}{493.927}{698.478}{Artin–}%
\TextAt{12}{79.226}{681.432}{Schreier}%
\TextAt{12}{131.267}{681.432}{theory}%
\TextAt{8}{169.44}{681.432}{.}%
\TextAt{8}{97.159}{661.397}{4)}%
\TextAt{8}{111.83}{661.397}{The}%
\TextAt{8}{136.244}{661.397}{chief}%
\TextAt{8}{163.911}{661.397}{application}%
\TextAt{8}{224.764}{661.397}{of}%
\TextAt{8}{238.455}{661.397}{the}%
\TextAt{8}{258.972}{661.397}{theory}%
\TextAt{8}{296.083}{661.397}{developed}%
\TextAt{8}{350.744}{661.397}{above}%
\TextAt{8}{384.28}{661.397}{is}%
\TextAt{8}{396.405}{661.397}{to}%
\TextAt{8}{411.076}{661.397}{the}%
\TextAt{8}{431.593}{661.397}{case}%
\TextAt{8}{456.724}{661.397}{where}%
\TextAt{13}{490.894}{661.397}{G}%
\TextAt{8}{504.386}{661.397}{is}%
\TextAt{8}{516.51}{661.397}{the}%
\TextAt{8}{79.226}{644.351}{absolute}%
\TextAt{8}{125.027}{644.351}{Galois}%
\TextAt{8}{160.499}{644.351}{group}%
\TextAt{13}{193.233}{644.351}{G}%
\TextAt{9}{202.467}{644.351}{\Character{40}}%
\TextAt{9}{207.338}{647.506}{\Character{22}}%
\TextAt{13}{207.019}{644.351}{k}%
\TextAt{16}{216.83}{644.351}{\Character{106}}%
\TextAt{13}{223.471}{644.351}{k}%
\TextAt{9}{229.961}{644.351}{\Character{41}}%
\TextAt{8}{237.969}{644.351}{of}%
\TextAt{8}{250.858}{644.351}{an}%
\TextAt{8}{266.675}{644.351}{actual}%
\TextAt{8}{301.358}{644.351}{field}%
\TextAt{13}{326.274}{644.351}{k}%
\TextAt{8}{332.763}{644.351}{,}%
\TextAt{13}{339.559}{644.351}{A}%
\TextAt{8}{351.791}{644.351}{is}%
\TextAt{8}{363.113}{644.351}{the}%
\TextAt{8}{382.828}{644.351}{multiplicative}%
\TextAt{8}{456.21}{644.351}{group}%
\TextAt{9}{489.262}{647.506}{\Character{22}}%
\TextAt{13}{488.944}{644.351}{k}%
\TextAt{3}{495.434}{648.689}{\Character{3}}%
\TextAt{8}{503.622}{644.351}{of}%
\TextAt{8}{516.51}{644.351}{the}%
\TextAt{8}{79.226}{627.305}{algebraic}%
\TextAt{8}{128.305}{627.305}{closure,}%
\TextAt{8}{170.354}{627.305}{and}%
\TextAt{13}{192.761}{627.305}{\Character{125}}%
\TextAt{8}{203.755}{627.305}{is}%
\TextAt{8}{215.163}{627.305}{the}%
\TextAt{13}{234.964}{627.305}{n}%
\TextAt{8}{241.952}{627.305}{-th}%
\TextAt{8}{260.449}{627.305}{power}%
\TextAt{8}{294.238}{627.305}{map}%
\TextAt{13}{319.896}{627.305}{a}%
\NativeArrowAt{329.362}{627.305}{11.955}{11.955}{\mapsto}%
\TextAt{13}{344.638}{627.305}{a}%
\TextAt{21}{350.783}{631.643}{n}%
\TextAt{8}{356.42}{627.305}{,}%
\TextAt{8}{363.284}{627.305}{for}%
\TextAt{8}{380.813}{627.305}{some}%
\TextAt{8}{409.784}{627.305}{natural}%
\TextAt{8}{450.41}{627.305}{number}%
\TextAt{13}{492.651}{627.305}{n}%
\TextAt{8}{503.18}{627.305}{which}%
\TextAt{8}{79.226}{610.259}{is}%
\TextAt{8}{91.088}{610.259}{relatively}%
\TextAt{8}{142.246}{610.259}{prime}%
\TextAt{8}{175.508}{610.259}{to}%
\TextAt{8}{189.916}{610.259}{the}%
\TextAt{8}{210.171}{610.259}{characteristic}%
\TextAt{8}{282.203}{610.259}{of}%
\TextAt{13}{295.631}{610.259}{k}%
\TextAt{8}{306.115}{610.259}{(in}%
\TextAt{8}{324.421}{610.259}{particular,}%
\TextAt{13}{381.785}{610.259}{n}%
\TextAt{8}{392.768}{610.259}{is}%
\TextAt{8}{404.63}{610.259}{arbitrary}%
\TextAt{8}{454.497}{610.259}{if}%
\TextAt{8}{465.319}{610.259}{char}%
\TextAt{9}{487.436}{610.259}{\Character{40}}%
\TextAt{13}{491.988}{610.259}{k}%
\TextAt{9}{498.477}{610.259}{\Character{41}}%
\TextAt{9}{506.499}{610.259}{\Character{61}}%
\TextAt{9}{519.088}{610.259}{0}%
\TextAt{8}{524.963}{610.259}{).}%
\TextAt{8}{79.226}{593.213}{Axiom}%
\TextAt{8}{115.857}{593.213}{(3.1)}%
\TextAt{8}{142.745}{593.213}{is}%
\TextAt{8}{153.422}{593.213}{always}%
\TextAt{8}{189.802}{593.213}{satisfied}%
\TextAt{8}{233.714}{593.213}{in}%
\TextAt{8}{246.279}{593.213}{this}%
\TextAt{8}{268.014}{593.213}{case}%
\TextAt{8}{291.698}{593.213}{and}%
\TextAt{8}{313.373}{593.213}{is}%
\TextAt{8}{324.049}{593.213}{called}%
\TextAt{12}{356.126}{593.213}{Hilbert}%
\TextAt{12}{402.074}{593.213}{90}%
\TextAt{8}{418.345}{593.213}{because}%
\TextAt{8}{460.56}{593.213}{this}%
\TextAt{8}{482.295}{593.213}{statement}%
\TextAt{8}{79.226}{576.167}{occurs}%
\TextAt{8}{114.841}{576.167}{as}%
\TextAt{19}{128.673}{576.167}{Satz}%
\TextAt{8}{153.922}{576.167}{number}%
\TextAt{8}{195.98}{576.167}{90}%
\TextAt{8}{211.056}{576.167}{among}%
\TextAt{8}{248.248}{576.167}{the}%
\TextAt{8}{267.867}{576.167}{169}%
\TextAt{8}{288.8}{576.167}{theorems}%
\TextAt{8}{338.402}{576.167}{in}%
\TextAt{8}{351.517}{576.167}{Hilbert’s}%
\TextAt{8}{399.158}{576.167}{famous}%
\TextAt{8}{438.682}{576.167}{“Zahlbericht”}%
\TextAt{8}{511.298}{576.167}{[72].}%
\TextAt{8}{79.226}{559.121}{Thus}%
\TextAt{8}{108.875}{559.121}{we}%
\TextAt{8}{126.102}{559.121}{have}%
\TextAt{8}{153.097}{559.121}{the}%
\TextAt{12}{79.226}{533.556}{Theorem}%
\TextAt{12}{136.24}{533.556}{(3.5)}%
\TextAt{12}{168.4}{533.556}{(Hilbert}%
\TextAt{12}{220.835}{533.556}{90).}%
\TextAt{8}{247.172}{533.556}{For}%
\TextAt{8}{268.142}{533.556}{a}%
\TextAt{8}{277.897}{533.556}{cyclic}%
\TextAt{8}{310.08}{533.556}{field}%
\TextAt{8}{335.437}{533.556}{extension}%
\TextAt{13}{387.203}{533.556}{L}%
\TextAt{16}{398.489}{533.556}{\Character{106}}%
\TextAt{13}{405.13}{533.556}{K}%
\TextAt{8}{419.845}{533.556}{one}%
\TextAt{8}{441.305}{533.556}{always}%
\TextAt{8}{478.772}{533.556}{has}%
\TextAt{13}{245.537}{506.119}{H}%
\TextAt{3}{256.16}{511.055}{\Character{0}}%
\TextAt{17}{262.746}{511.055}{1}%
\TextAt{9}{267.479}{506.119}{\Character{40}}%
\TextAt{13}{272.031}{506.119}{G}%
\TextAt{9}{281.265}{506.119}{\Character{40}}%
\TextAt{13}{285.817}{506.119}{L}%
\TextAt{16}{297.102}{506.119}{\Character{106}}%
\TextAt{13}{303.744}{506.119}{K}%
\TextAt{9}{314.562}{506.119}{\Character{41}}%
\TextAt{13}{319.114}{506.119}{\Character{59}}%
\TextAt{13}{324.338}{506.119}{L}%
\TextAt{3}{332.323}{511.055}{\Character{3}}%
\TextAt{9}{337.055}{506.119}{\Character{41}}%
\TextAt{9}{344.91}{506.119}{\Character{61}}%
\TextAt{9}{357.331}{506.119}{1}%
\TextAt{13}{363.207}{506.119}{\Character{58}}%
\TextAt{8}{79.226}{478.682}{In}%
\TextAt{8}{93.847}{478.682}{other}%
\TextAt{8}{124.417}{478.682}{words:}%
\TextAt{8}{97.159}{458.647}{An}%
\TextAt{8}{116.505}{458.647}{element}%
\TextAt{13}{159.918}{458.647}{\Character{11}}%
\NativeMathAt{171.05}{458.647}{11.955}{\in}%
\TextAt{13}{182.632}{458.647}{L}%
\TextAt{3}{190.596}{462.985}{\Character{3}}%
\TextAt{8}{199.396}{458.647}{of}%
\TextAt{8}{212.897}{458.647}{norm}%
\TextAt{13}{243.636}{458.647}{N}%
\TextAt{21}{253.01}{456.787}{L}%
\TextAt{3}{258.769}{456.787}{\Character{106}}%
\TextAt{21}{261.121}{456.787}{K}%
\TextAt{9}{269.324}{458.647}{\Character{40}}%
\TextAt{13}{273.876}{458.647}{\Character{11}}%
\TextAt{9}{281.398}{458.647}{\Character{41}}%
\TextAt{9}{289.551}{458.647}{\Character{61}}%
\TextAt{9}{302.26}{458.647}{1}%
\TextAt{8}{312.198}{458.647}{is}%
\TextAt{8}{324.132}{458.647}{of}%
\TextAt{8}{337.633}{458.647}{the}%
\TextAt{8}{357.96}{458.647}{form}%
\TextAt{13}{385.77}{458.647}{\Character{11}}%
\TextAt{9}{396.903}{458.647}{\Character{61}}%
\TextAt{13}{409.619}{458.647}{\Character{12}}%
\TextAt{21}{416.89}{462.985}{\Character{27}}%
\TextAt{3}{422.015}{462.985}{\Character{0}}%
\TextAt{17}{428.602}{462.985}{1}%
\TextAt{8}{433.334}{458.647}{,}%
\TextAt{8}{440.696}{458.647}{where}%
\TextAt{13}{474.676}{458.647}{\Character{12}}%
\NativeMathAt{485.558}{458.647}{11.955}{\in}%
\TextAt{13}{497.14}{458.647}{L}%
\TextAt{3}{505.104}{462.985}{\Character{3}}%
\TextAt{8}{513.904}{458.647}{and}%
\TextAt{13}{79.226}{441.601}{\Character{27}}%
\TextAt{8}{90.206}{441.601}{is}%
\TextAt{8}{101.97}{441.601}{a}%
\TextAt{8}{111.725}{441.601}{generator}%
\TextAt{8}{163.766}{441.601}{of}%
\TextAt{13}{177.096}{441.601}{G}%
\TextAt{9}{186.33}{441.601}{\Character{40}}%
\TextAt{13}{190.882}{441.601}{L}%
\TextAt{16}{202.167}{441.601}{\Character{106}}%
\TextAt{13}{208.809}{441.601}{K}%
\TextAt{9}{219.627}{441.601}{\Character{41}}%
\TextAt{8}{224.179}{441.601}{.}%
\TextAt{19}{79.226}{407.516}{Proof.}%
\TextAt{8}{116}{407.516}{Let}%
\TextAt{13}{138.392}{407.516}{n}%
\TextAt{9}{151.123}{407.516}{\Character{61}}%
\TextAt{9}{165.959}{407.516}{\Character{91}}%
\TextAt{13}{169.224}{407.516}{L}%
\TextAt{9}{182.933}{407.516}{\Character{58}}%
\TextAt{13}{191.928}{407.516}{K}%
\TextAt{9}{202.746}{407.516}{\Character{93}}%
\TextAt{8}{205.998}{407.516}{.}%
\TextAt{8}{218.712}{407.516}{By}%
\TextAt{8}{238.497}{407.516}{virtue}%
\TextAt{8}{274.063}{407.516}{of}%
\TextAt{8}{288.815}{407.516}{the}%
\TextAt{8}{310.394}{407.516}{linear}%
\TextAt{8}{344.334}{407.516}{independence}%
\TextAt{8}{418.252}{407.516}{of}%
\TextAt{8}{433.004}{407.516}{the}%
\TextAt{8}{454.583}{407.516}{automorphisms}%
\TextAt{9}{79.226}{390.47}{1}%
\TextAt{13}{85.079}{390.47}{\Character{59}}%
\TextAt{13}{90.303}{390.47}{\Character{27}}%
\TextAt{13}{96.747}{390.47}{\Character{59}}%
\TextAt{13}{101.96}{390.47}{\Character{58}}%
\TextAt{13}{107.196}{390.47}{\Character{58}}%
\TextAt{13}{112.433}{390.47}{\Character{58}}%
\TextAt{13}{117.657}{390.47}{\Character{59}}%
\TextAt{13}{122.893}{390.47}{\Character{27}}%
\TextAt{21}{130.059}{394.809}{n}%
\TextAt{3}{135.197}{394.809}{\Character{0}}%
\TextAt{17}{141.783}{394.809}{1}%
\TextAt{8}{150.413}{390.47}{(see}%
\TextAt{8}{173.881}{390.47}{[15],}%
\TextAt{8}{199.25}{390.47}{chap.}%
\TextAt{8}{231.441}{390.47}{5,}%
\TextAt{8}{244.448}{390.47}{§7,}%
\TextAt{8}{263.134}{390.47}{no.}%
\TextAt{8}{283.944}{390.47}{5),}%
\TextAt{8}{301.506}{390.47}{there}%
\TextAt{8}{331.418}{390.47}{exists}%
\TextAt{8}{363.733}{390.47}{an}%
\TextAt{8}{379.992}{390.47}{element}%
\TextAt{13}{423.234}{390.47}{\Character{13}}%
\NativeMathAt{433.277}{390.47}{11.955}{\in}%
\TextAt{13}{444.568}{390.47}{L}%
\TextAt{3}{452.532}{394.809}{\Character{3}}%
\TextAt{8}{461.162}{390.47}{such}%
\TextAt{8}{487.559}{390.47}{that}%
\TextAt{13}{167.386}{363.033}{\Character{12}}%
\TextAt{9}{177.978}{363.033}{\Character{61}}%
\TextAt{13}{190.404}{363.033}{\Character{13}}%
\TextAt{9}{199.783}{363.033}{\Character{43}}%
\TextAt{13}{211.544}{363.033}{\Character{11}\Character{13}}%
\TextAt{21}{225.788}{367.969}{\Character{27}}%
\TextAt{9}{234.068}{363.033}{\Character{43}}%
\TextAt{13}{245.829}{363.033}{\Character{11}}%
\TextAt{17}{253.351}{367.969}{1\Character{43}}%
\TextAt{21}{264.172}{367.969}{\Character{27}}%
\TextAt{13}{269.795}{363.033}{\Character{13}}%
\TextAt{21}{276.518}{367.969}{\Character{27}}%
\TextAt{24}{281.643}{370.782}{2}%
\TextAt{9}{288.949}{363.033}{\Character{43}}%
\TextAt{16}{300.71}{363.033}{\Character{1}}%
\TextAt{16}{306.006}{363.033}{\Character{1}}%
\TextAt{16}{311.314}{363.033}{\Character{1}}%
\TextAt{9}{317.314}{363.033}{\Character{43}}%
\TextAt{13}{329.076}{363.033}{\Character{11}}%
\TextAt{17}{336.597}{367.969}{1\Character{43}}%
\TextAt{21}{347.418}{367.969}{\Character{27}}%
\TextAt{17}{352.543}{367.969}{\Character{43}}%
\TextAt{3}{359.13}{367.969}{\Character{1}\Character{1}\Character{1}}%
\TextAt{17}{366.187}{367.969}{\Character{43}}%
\TextAt{21}{372.773}{367.969}{\Character{27}}%
\TextAt{22}{377.899}{370.782}{n}%
\TextAt{23}{382.499}{370.782}{\Character{0}}%
\TextAt{24}{388.256}{370.782}{2}%
\TextAt{13}{392.905}{363.033}{\Character{13}}%
\TextAt{21}{399.627}{367.969}{\Character{27}}%
\TextAt{22}{404.752}{370.782}{n}%
\TextAt{23}{409.353}{370.782}{\Character{0}}%
\TextAt{24}{415.109}{370.782}{1}%
\NativeMathAt{423.079}{363.033}{11.955}{\neq}%
\TextAt{9}{435.488}{363.033}{0}%
\TextAt{13}{441.358}{363.033}{\Character{58}}%
\TextAt{8}{79.226}{335.596}{As}%
\TextAt{13}{96.513}{335.596}{N}%
\TextAt{21}{105.887}{333.737}{L}%
\TextAt{3}{111.646}{333.737}{\Character{106}}%
\TextAt{21}{113.999}{333.737}{K}%
\TextAt{9}{122.201}{335.596}{\Character{40}}%
\TextAt{13}{126.754}{335.596}{\Character{11}}%
\TextAt{9}{134.275}{335.596}{\Character{41}}%
\TextAt{9}{142.13}{335.596}{\Character{61}}%
\TextAt{9}{154.551}{335.596}{1}%
\TextAt{8}{160.427}{335.596}{,}%
\TextAt{8}{167.576}{335.596}{one}%
\TextAt{8}{189.036}{335.596}{gets}%
\TextAt{13}{213.161}{335.596}{\Character{11}\Character{12}}%
\TextAt{21}{227.954}{339.935}{\Character{27}}%
\TextAt{9}{236.898}{335.596}{\Character{61}}%
\TextAt{13}{249.324}{335.596}{\Character{12}}%
\TextAt{8}{256.595}{335.596}{,}%
\TextAt{8}{263.744}{335.596}{and}%
\TextAt{8}{286.507}{335.596}{thus}%
\TextAt{13}{312.258}{335.596}{\Character{11}}%
\TextAt{9}{323.101}{335.596}{\Character{61}}%
\TextAt{13}{335.526}{335.596}{\Character{12}}%
\TextAt{17}{342.798}{339.935}{1}%
\TextAt{3}{347.032}{339.935}{\Character{0}}%
\TextAt{21}{353.618}{339.935}{\Character{27}}%
\TextAt{8}{359.242}{335.596}{.}%
\Vector{q 0 G 0 g 0.399 w 452.334 -438.404 m 452.334 -430.334 l S Q}%
\Vector{q 0 G 0 g 0.399 w 452.534 -430.533 m 459.707 -430.533 l S Q}%
\Vector{q 0 G 0 g 0.399 w 452.534 -438.205 m 459.707 -438.205 l S Q}%
\Vector{q 0 G 0 g 0.399 w 459.906 -438.404 m 459.906 -430.334 l S Q}%
\TextAt{8}{97.159}{307.042}{If}%
\TextAt{8}{108.276}{307.042}{now}%
\TextAt{8}{132.09}{307.042}{the}%
\TextAt{8}{151.672}{307.042}{field}%
\TextAt{13}{176.454}{307.042}{K}%
\TextAt{8}{190.595}{307.042}{contains}%
\TextAt{8}{235.94}{307.042}{the}%
\TextAt{8}{255.522}{307.042}{group}%
\TextAt{13}{288.123}{307.042}{\Character{22}}%
\TextAt{21}{295.166}{305.249}{n}%
\TextAt{8}{304.125}{307.042}{of}%
\TextAt{13}{316.88}{307.042}{n}%
\TextAt{8}{323.868}{307.042}{-th}%
\TextAt{8}{342.147}{307.042}{roots}%
\TextAt{8}{371.233}{307.042}{of}%
\TextAt{8}{383.988}{307.042}{unity,}%
\TextAt{8}{416.369}{307.042}{the}%
\TextAt{8}{435.951}{307.042}{operator}%
\TextAt{13}{482.54}{307.042}{\Character{125}}%
\TextAt{9}{489.992}{307.042}{\Character{40}}%
\TextAt{13}{494.545}{307.042}{a}%
\TextAt{9}{500.69}{307.042}{\Character{41}}%
\TextAt{9}{508.545}{307.042}{\Character{61}}%
\TextAt{13}{520.988}{307.042}{a}%
\TextAt{21}{527.133}{311.38}{n}%
\TextAt{8}{79.226}{289.996}{has}%
\TextAt{8}{99.203}{289.996}{exponent}%
\TextAt{13}{148.709}{289.996}{n}%
\TextAt{8}{155.696}{289.996}{,}%
\TextAt{8}{162.128}{289.996}{and}%
\TextAt{8}{183.993}{289.996}{we}%
\TextAt{8}{200.324}{289.996}{obtain}%
\TextAt{8}{235.854}{289.996}{the}%
\TextAt{8}{255.114}{289.996}{following}%
\TextAt{8}{303.651}{289.996}{corollary,}%
\TextAt{8}{353.671}{289.996}{which}%
\TextAt{8}{386.261}{289.996}{is}%
\TextAt{8}{397.127}{289.996}{the}%
\TextAt{8}{416.387}{289.996}{most}%
\TextAt{8}{444.17}{289.996}{important}%
\TextAt{8}{498.566}{289.996}{special}%
\TextAt{8}{79.226}{272.95}{case}%
\TextAt{8}{103.997}{272.95}{of}%
\TextAt{8}{117.327}{272.95}{theorem}%
\TextAt{8}{162.853}{272.95}{(3.3).}%
\TextAt{12}{79.226}{247.385}{Corollary}%
\TextAt{12}{137.823}{247.385}{(3.6).}%
\TextAt{8}{172.103}{247.385}{Let}%
\TextAt{13}{192.037}{247.385}{n}%
\TextAt{8}{201.887}{247.385}{be}%
\TextAt{8}{216.776}{247.385}{a}%
\TextAt{8}{225.496}{247.385}{natural}%
\TextAt{8}{265.443}{247.385}{number}%
\TextAt{8}{307.004}{247.385}{which}%
\TextAt{8}{339.455}{247.385}{is}%
\TextAt{8}{350.183}{247.385}{relatively}%
\TextAt{8}{400.208}{247.385}{prime}%
\TextAt{8}{432.337}{247.385}{to}%
\TextAt{8}{445.612}{247.385}{the}%
\TextAt{8}{464.733}{247.385}{characteristic}%
\TextAt{8}{79.226}{230.339}{of}%
\TextAt{8}{92.556}{230.339}{the}%
\TextAt{8}{112.713}{230.339}{field}%
\TextAt{13}{138.07}{230.339}{K}%
\TextAt{8}{148.887}{230.339}{,}%
\TextAt{8}{156.036}{230.339}{and}%
\TextAt{8}{178.799}{230.339}{assume}%
\TextAt{8}{219.243}{230.339}{that}%
\TextAt{13}{244.612}{230.339}{\Character{22}}%
\TextAt{21}{251.655}{228.545}{n}%
\NativeMathAt{260.612}{230.339}{11.955}{\subseteq}%
\TextAt{13}{273.232}{230.339}{K}%
\TextAt{8}{284.049}{230.339}{.}%
\TextAt{8}{97.159}{210.304}{Then}%
\TextAt{8}{128.726}{210.304}{the}%
\TextAt{8}{149.892}{210.304}{abelian}%
\TextAt{8}{191.549}{210.304}{extensions}%
\TextAt{13}{248.939}{210.304}{L}%
\TextAt{16}{261.944}{210.304}{\Character{106}}%
\TextAt{13}{270.306}{210.304}{K}%
\TextAt{8}{286.031}{210.304}{of}%
\TextAt{8}{300.37}{210.304}{exponent}%
\TextAt{13}{351.782}{210.304}{n}%
\TextAt{8}{363.677}{210.304}{correspond}%
\TextAt{9}{424.259}{210.304}{1}%
\TextAt{8}{430.112}{210.304}{-}%
\TextAt{9}{434.009}{210.304}{1}%
\TextAt{8}{444.769}{210.304}{to}%
\TextAt{8}{460.089}{210.304}{the}%
\TextAt{8}{481.255}{210.304}{subgroups}%
\TextAt{9}{79.226}{193.258}{\Character{1}}%
\NativeMathAt{92.302}{193.258}{11.955}{\subseteq}%
\TextAt{13}{104.921}{193.258}{K}%
\TextAt{3}{115.739}{197.596}{\Character{3}}%
\TextAt{8}{124.369}{193.258}{which}%
\TextAt{8}{157.855}{193.258}{contain}%
\TextAt{13}{199.16}{193.258}{K}%
\TextAt{3}{209.978}{197.596}{\Character{3}}%
\TextAt{21}{214.212}{197.596}{n}%
\TextAt{8}{219.848}{193.258}{,}%
\TextAt{8}{226.998}{193.258}{via}%
\TextAt{8}{246.186}{193.258}{the}%
\TextAt{8}{266.342}{193.258}{rule}%
\TextAt{9}{254.228}{165.821}{\Character{1}}%
\NativeArrowAt{267.304}{165.821}{11.955}{19.26}{\longmapsto}%
\TextAt{13}{289.886}{165.821}{L}%
\TextAt{9}{301.171}{165.821}{\Character{61}}%
\TextAt{13}{313.596}{165.821}{K}%
\TextAt{9}{324.414}{165.821}{\Character{40}}%
\TextAt{22}{332.287}{172.356}{n}%
\TextAt{16}{330.246}{176.527}{\Character{112}}%
\Vector{q 0 G 0 g 0.478 w 268.209 -597.234 m 277.964 -597.234 l S Q}%
\TextAt{9}{340.209}{165.821}{\Character{1}\Character{41}}%
\TextAt{13}{354.516}{165.821}{\Character{59}}%
\TextAt{8}{79.226}{138.384}{and}%
\TextAt{8}{101.989}{138.384}{we}%
\TextAt{8}{119.216}{138.384}{have}%
\TextAt{13}{229.004}{121.338}{G}%
\TextAt{9}{238.238}{121.338}{\Character{40}}%
\TextAt{13}{242.79}{121.338}{L}%
\TextAt{16}{254.075}{121.338}{\Character{106}}%
\TextAt{13}{260.717}{121.338}{K}%
\TextAt{9}{271.535}{121.338}{\Character{41}}%
\NativeMathAt{279.408}{121.338}{11.955}{\cong}%
\TextAt{8}{292.027}{121.338}{Hom}%
\TextAt{9}{316.416}{121.338}{\Character{40}\Character{1}\Character{47}}%
\TextAt{13}{336.576}{121.338}{K}%
\TextAt{3}{347.394}{126.274}{\Character{3}}%
\TextAt{21}{351.628}{126.274}{n}%
\TextAt{13}{357.264}{121.338}{\Character{59}}%
\TextAt{13}{362.488}{121.338}{\Character{22}}%
\TextAt{21}{369.551}{119.545}{n}%
\TextAt{9}{375.188}{121.338}{\Character{41}}%
\TextAt{13}{379.74}{121.338}{\Character{58}}%
\TextAt{8}{97.159}{89.967}{Hilbert’s}%
\TextAt{8}{145.976}{89.967}{theorem}%
\TextAt{8}{192.139}{89.967}{90,}%
\TextAt{8}{211.801}{89.967}{which}%
\TextAt{8}{245.925}{89.967}{is}%
\TextAt{8}{258.327}{89.967}{the}%
\TextAt{8}{279.121}{89.967}{main}%
\TextAt{8}{309.024}{89.967}{basis}%
\TextAt{8}{338.402}{89.967}{of}%
\TextAt{8}{352.37}{89.967}{this}%
\TextAt{8}{375.83}{89.967}{corollary,}%
\TextAt{8}{427.365}{89.967}{admits}%
\TextAt{8}{466.438}{89.967}{the}%
\TextAt{8}{487.232}{89.967}{following}%
\TextAt{8}{79.226}{72.921}{generalization}%
\TextAt{8}{153.043}{72.921}{to}%
\TextAt{8}{166.367}{72.921}{arbitrary}%
\TextAt{8}{215.149}{72.921}{Galois}%
\TextAt{8}{250.076}{72.921}{extensions}%
\TextAt{13}{305.47}{72.921}{L}%
\TextAt{16}{316.755}{72.921}{\Character{106}}%
\TextAt{13}{323.397}{72.921}{K}%
\TextAt{8}{334.214}{72.921}{,}%
\TextAt{8}{340.574}{72.921}{which}%
\TextAt{8}{373.074}{72.921}{goes}%
\TextAt{8}{397.839}{72.921}{back}%
\TextAt{8}{424.17}{72.921}{to}%
\TextAt{8}{437.493}{72.921}{the}%
\TextAt{8}{456.663}{72.921}{mathematician}%
\TextAt{8}{297.211}{43.033}{257}%
\Destination{78.23}{809.815}{nav.page.257}%
\Destination{79.36}{776.155}{nav.print.281}%
\Destination{79.36}{625.255}{nav.print.282}%
\Destination{79.36}{547.555}{nav.number.4.3.5}%
\Destination{79.36}{261.355}{nav.number.4.3.6}%
\Hotlink{500.089}{741.855}{6.858}{11.047}{nav.section.4.7}%
\Hotlink{510.798}{573.955}{19.22}{11.047}{nav.cite.N072}%
%
\Hotlink{173.381}{388.255}{19.22}{11.047}{nav.cite.N015}%
\Hotlink{241.731}{724.855}{25.078}{11.047}{nav.number.4.7.4}%
\Hotlink{162.353}{270.755}{25.078}{11.047}{nav.number.4.3.3}%
\Hotlink{117.027}{591}{21.722}{11.047}{nav.number.4.3.1}%
}

% Page 270 (body page 258)
\Page{%
\TextAt{2}{79.226}{796.943}{IV.}%
\TextAt{2}{98.012}{796.943}{Abstract}%
\TextAt{2}{143.252}{796.943}{Class}%
\TextAt{2}{171.844}{796.943}{Field}%
\TextAt{2}{199.575}{796.943}{Theory}%
\Vector{q 0 G 0 g 0.299 w 7.226 18.729 m 460.769 18.729 l S Q}%
\TextAt{26}{79.226}{761.124}{Emmy}%
\TextAt{26}{117.143}{761.124}{Noether}%
\TextAt{8}{172.736}{761.124}{(1882–1935).}%
\TextAt{8}{242.826}{761.124}{Let}%
\TextAt{13}{263.224}{761.124}{G}%
\TextAt{8}{275.783}{761.124}{be}%
\TextAt{8}{291.136}{761.124}{a}%
\TextAt{8}{300.32}{761.124}{finite}%
\TextAt{8}{329.66}{761.124}{group}%
\TextAt{8}{362.264}{761.124}{and}%
\TextAt{13}{384.455}{761.124}{A}%
\TextAt{8}{396.556}{761.124}{a}%
\TextAt{8}{405.739}{761.124}{multiplicative}%
\TextAt{13}{478.991}{761.124}{G}%
\TextAt{8}{488.225}{761.124}{-module.}%
\TextAt{8}{79.226}{744.078}{A}%
\TextAt{12}{90.887}{744.078}{1-cocycle}%
\TextAt{8}{144.112}{744.078}{,}%
\TextAt{8}{150.452}{744.078}{or}%
\TextAt{12}{163.752}{744.078}{crossed}%
\TextAt{12}{209.502}{744.078}{homomorphism}%
\TextAt{8}{300.302}{744.078}{,}%
\TextAt{8}{306.642}{744.078}{of}%
\TextAt{13}{318.962}{744.078}{G}%
\TextAt{8}{331.083}{744.078}{with}%
\TextAt{8}{356.732}{744.078}{values}%
\TextAt{8}{390.583}{744.078}{in}%
\TextAt{13}{403.226}{744.078}{A}%
\TextAt{8}{414.888}{744.078}{is}%
\TextAt{8}{425.641}{744.078}{a}%
\TextAt{8}{434.386}{744.078}{function}%
\TextAt{13}{479.224}{744.078}{f}%
\TextAt{9}{489.591}{744.078}{\Character{58}}%
\TextAt{13}{496.164}{744.078}{G}%
\NativeMathAt{508.718}{744.078}{11.955}{\rightarrow}%
\TextAt{13}{523.994}{744.078}{A}%
\TextAt{8}{79.226}{727.032}{satisfying}%
\TextAt{13}{257.808}{709.986}{f}%
\TextAt{9}{264.855}{709.986}{\Character{40}}%
\TextAt{13}{269.407}{709.986}{\Character{27}}%
\TextAt{13}{276.473}{709.986}{\Character{28}}%
\TextAt{9}{282.909}{709.986}{\Character{41}}%
\TextAt{9}{290.764}{709.986}{\Character{61}}%
\TextAt{13}{303.207}{709.986}{f}%
\TextAt{9}{310.254}{709.986}{\Character{40}}%
\TextAt{13}{314.806}{709.986}{\Character{27}}%
\TextAt{9}{321.888}{709.986}{\Character{41}}%
\TextAt{21}{326.441}{714.922}{\Character{28}}%
\TextAt{13}{331.617}{709.986}{f}%
\TextAt{9}{338.664}{709.986}{\Character{40}}%
\TextAt{13}{343.216}{709.986}{\Character{28}}%
\TextAt{9}{349.635}{709.986}{\Character{41}}%
\TextAt{8}{79.226}{686.536}{for}%
\TextAt{8}{97.665}{686.536}{all}%
\TextAt{13}{114.479}{686.536}{\Character{27}}%
\TextAt{13}{120.911}{686.536}{\Character{59}}%
\TextAt{13}{126.135}{686.536}{\Character{28}}%
\NativeMathAt{136.839}{686.536}{11.955}{\in}%
\TextAt{13}{149.075}{686.536}{G}%
\TextAt{8}{158.309}{686.536}{.}%
\TextAt{8}{168.42}{686.536}{The}%
\TextAt{8}{193.028}{686.536}{1-cocycles}%
\TextAt{8}{248.265}{686.536}{form}%
\TextAt{8}{276.46}{686.536}{an}%
\TextAt{8}{293.273}{686.536}{abelian}%
\TextAt{8}{334.475}{686.536}{group}%
\TextAt{13}{368.205}{686.536}{Z}%
\TextAt{17}{377.076}{690.875}{1}%
\TextAt{9}{381.808}{686.536}{\Character{40}}%
\TextAt{13}{386.36}{686.536}{G\Character{59}}%
\TextAt{13}{400.814}{686.536}{A}%
\TextAt{9}{409.613}{686.536}{\Character{41}}%
\TextAt{8}{414.166}{686.536}{.}%
\TextAt{8}{424.277}{686.536}{For}%
\TextAt{8}{445.801}{686.536}{every}%
\TextAt{13}{477.247}{686.536}{a}%
\NativeMathAt{487.657}{686.536}{11.955}{\in}%
\TextAt{13}{499.893}{686.536}{A}%
\TextAt{8}{508.668}{686.536}{,}%
\TextAt{8}{516.51}{686.536}{the}%
\TextAt{8}{79.226}{669.49}{function}%
\TextAt{13}{273.354}{652.444}{f}%
\TextAt{21}{279.124}{650.651}{a}%
\TextAt{9}{284.12}{652.444}{\Character{40}}%
\TextAt{13}{288.672}{652.444}{\Character{27}}%
\TextAt{9}{295.754}{652.444}{\Character{41}}%
\TextAt{9}{303.609}{652.444}{\Character{61}}%
\TextAt{13}{316.053}{652.444}{a}%
\TextAt{21}{322.198}{657.381}{\Character{27}}%
\TextAt{3}{327.323}{657.381}{\Character{0}}%
\TextAt{17}{333.91}{657.381}{1}%
\TextAt{8}{79.226}{628.995}{is}%
\TextAt{8}{90.99}{628.995}{a}%
\TextAt{8}{100.745}{628.995}{1-cocycle,}%
\TextAt{8}{154.066}{628.995}{for}%
\TextAt{8}{171.95}{628.995}{one}%
\TextAt{8}{193.41}{628.995}{has}%
\TextAt{13}{191.717}{600.163}{f}%
\TextAt{21}{197.487}{598.369}{a}%
\TextAt{9}{202.483}{600.163}{\Character{40}}%
\TextAt{13}{207.035}{600.163}{\Character{27}}%
\TextAt{13}{214.101}{600.163}{\Character{28}}%
\TextAt{9}{220.537}{600.163}{\Character{41}}%
\TextAt{9}{228.392}{600.163}{\Character{61}}%
\TextAt{13}{240.835}{600.163}{a}%
\TextAt{21}{246.98}{605.099}{\Character{27}}%
\TextAt{21}{252.097}{605.099}{\Character{28}}%
\TextAt{3}{256.784}{605.099}{\Character{0}}%
\TextAt{17}{263.37}{605.099}{1}%
\TextAt{9}{271.424}{600.163}{\Character{61}}%
\TextAt{9}{283.834}{600.163}{\Character{40}}%
\TextAt{13}{288.401}{600.163}{a}%
\TextAt{21}{294.546}{605.099}{\Character{27}}%
\TextAt{3}{299.672}{605.099}{\Character{0}}%
\TextAt{17}{306.258}{605.099}{1}%
\TextAt{9}{310.99}{600.163}{\Character{41}}%
\TextAt{21}{315.543}{605.099}{\Character{28}}%
\TextAt{13}{320.719}{600.163}{a}%
\TextAt{21}{326.864}{605.099}{\Character{28}}%
\TextAt{3}{331.543}{605.099}{\Character{0}}%
\TextAt{17}{338.129}{605.099}{1}%
\TextAt{9}{346.182}{600.163}{\Character{61}}%
\TextAt{13}{358.608}{600.163}{f}%
\TextAt{21}{364.378}{598.369}{a}%
\TextAt{9}{369.374}{600.163}{\Character{40}}%
\TextAt{13}{373.926}{600.163}{\Character{27}}%
\TextAt{9}{381.009}{600.163}{\Character{41}}%
\TextAt{21}{385.561}{605.099}{\Character{28}}%
\TextAt{13}{390.737}{600.163}{f}%
\TextAt{21}{396.507}{598.369}{a}%
\TextAt{9}{401.503}{600.163}{\Character{40}}%
\TextAt{13}{406.056}{600.163}{\Character{28}}%
\TextAt{9}{412.475}{600.163}{\Character{41}}%
\TextAt{13}{417.027}{600.163}{\Character{58}}%
\TextAt{8}{79.226}{571.33}{The}%
\TextAt{8}{102.92}{571.33}{functions}%
\TextAt{13}{153.022}{571.33}{f}%
\TextAt{21}{158.792}{569.537}{a}%
\TextAt{8}{167.325}{571.33}{are}%
\TextAt{8}{186.476}{571.33}{called}%
\TextAt{12}{219.279}{571.33}{1-coboundaries}%
\TextAt{8}{310.949}{571.33}{and}%
\TextAt{8}{333.352}{571.33}{form}%
\TextAt{8}{360.632}{571.33}{a}%
\TextAt{8}{370.027}{571.33}{subgroup}%
\TextAt{13}{420.464}{571.33}{B}%
\TextAt{17}{429.96}{575.669}{1}%
\TextAt{9}{434.692}{571.33}{\Character{40}}%
\TextAt{13}{439.245}{571.33}{G\Character{59}}%
\TextAt{13}{453.699}{571.33}{A}%
\TextAt{9}{462.498}{571.33}{\Character{41}}%
\TextAt{8}{470.587}{571.33}{of}%
\TextAt{13}{483.557}{571.33}{Z}%
\TextAt{17}{492.428}{575.669}{1}%
\TextAt{9}{497.16}{571.33}{\Character{40}}%
\TextAt{13}{501.712}{571.33}{G\Character{59}}%
\TextAt{13}{516.166}{571.33}{A}%
\TextAt{9}{524.965}{571.33}{\Character{41}}%
\TextAt{8}{529.518}{571.33}{.}%
\TextAt{8}{79.226}{554.284}{We}%
\TextAt{8}{99.371}{554.284}{define}%
\TextAt{13}{225.068}{537.238}{H}%
\TextAt{17}{235.691}{542.174}{1}%
\TextAt{9}{240.423}{537.238}{\Character{40}}%
\TextAt{13}{244.975}{537.238}{G\Character{59}}%
\TextAt{13}{259.429}{537.238}{A}%
\TextAt{9}{268.229}{537.238}{\Character{41}}%
\TextAt{9}{276.084}{537.238}{\Character{61}}%
\TextAt{13}{288.527}{537.238}{Z}%
\TextAt{17}{297.398}{542.174}{1}%
\TextAt{9}{302.13}{537.238}{\Character{40}}%
\TextAt{13}{306.683}{537.238}{G\Character{59}}%
\TextAt{13}{321.137}{537.238}{A}%
\TextAt{9}{329.936}{537.238}{\Character{41}\Character{47}}%
\TextAt{13}{340.341}{537.238}{B}%
\TextAt{17}{349.837}{542.174}{1}%
\TextAt{9}{354.57}{537.238}{\Character{40}}%
\TextAt{13}{359.122}{537.238}{G\Character{59}}%
\TextAt{13}{373.576}{537.238}{A}%
\TextAt{9}{382.375}{537.238}{\Character{41}}%
\TextAt{8}{79.226}{513.789}{and}%
\TextAt{8}{101.989}{513.789}{obtain}%
\TextAt{8}{138.416}{513.789}{as}%
\TextAt{8}{152.786}{513.789}{a}%
\TextAt{8}{162.542}{513.789}{first}%
\TextAt{8}{186.667}{513.789}{result}%
\TextAt{8}{219.245}{513.789}{about}%
\TextAt{8}{252.743}{513.789}{this}%
\TextAt{8}{275.566}{513.789}{group}%
\TextAt{8}{308.741}{513.789}{the}%
\TextAt{12}{79.226}{487.825}{Proposition}%
\TextAt{12}{151.937}{487.825}{(3.7).}%
\TextAt{8}{187.253}{487.825}{If}%
\TextAt{13}{198.945}{487.825}{G}%
\TextAt{8}{212.076}{487.825}{is}%
\TextAt{8}{223.84}{487.825}{cyclic,}%
\TextAt{8}{259.275}{487.825}{then}%
\TextAt{13}{285.935}{487.825}{H}%
\TextAt{17}{296.558}{492.163}{1}%
\TextAt{9}{301.29}{487.825}{\Character{40}}%
\TextAt{13}{305.843}{487.825}{G\Character{59}}%
\TextAt{13}{320.297}{487.825}{A}%
\TextAt{9}{329.096}{487.825}{\Character{41}}%
\NativeMathAt{336.969}{487.825}{11.955}{\cong}%
\TextAt{13}{349.588}{487.825}{H}%
\TextAt{3}{360.211}{492.163}{\Character{0}}%
\TextAt{17}{366.797}{492.163}{1}%
\TextAt{9}{371.53}{487.825}{\Character{40}}%
\TextAt{13}{376.082}{487.825}{G\Character{59}}%
\TextAt{13}{390.536}{487.825}{A}%
\TextAt{9}{399.335}{487.825}{\Character{41}}%
\TextAt{8}{403.887}{487.825}{.}%
\TextAt{19}{79.226}{452.942}{Proof.}%
\TextAt{8}{116}{452.942}{Let}%
\TextAt{13}{136.97}{452.942}{G}%
\TextAt{9}{149.524}{452.942}{\Character{61}}%
\TextAt{16}{161.95}{452.942}{\Character{104}}%
\TextAt{13}{166.599}{452.942}{\Character{27}}%
\TextAt{16}{173.681}{452.942}{\Character{105}}%
\TextAt{8}{178.33}{452.942}{.}%
\TextAt{8}{186.779}{452.942}{If}%
\TextAt{13}{198.471}{452.942}{f}%
\NativeMathAt{208.838}{452.942}{11.955}{\in}%
\TextAt{13}{220.129}{452.942}{Z}%
\TextAt{17}{229}{457.281}{1}%
\TextAt{9}{233.732}{452.942}{\Character{40}}%
\TextAt{13}{238.285}{452.942}{G\Character{59}}%
\TextAt{13}{252.739}{452.942}{A}%
\TextAt{9}{261.538}{452.942}{\Character{41}}%
\TextAt{8}{266.09}{452.942}{,}%
\TextAt{8}{273.239}{452.942}{then}%
\TextAt{8}{299.899}{452.942}{for}%
\TextAt{13}{317.784}{452.942}{k}%
\NativeMathAt{327.595}{452.942}{11.955}{\geq}%
\TextAt{9}{340.214}{452.942}{1}%
\TextAt{13}{142.467}{415.496}{f}%
\TextAt{9}{149.514}{415.496}{\Character{40}}%
\TextAt{13}{154.066}{415.496}{\Character{27}}%
\TextAt{21}{161.148}{420.433}{k}%
\TextAt{9}{166.268}{415.496}{\Character{41}}%
\TextAt{9}{174.123}{415.496}{\Character{61}}%
\TextAt{13}{186.567}{415.496}{f}%
\TextAt{9}{193.613}{415.496}{\Character{40}}%
\TextAt{13}{198.166}{415.496}{\Character{27}}%
\TextAt{21}{205.248}{420.433}{k}%
\TextAt{3}{209.87}{420.433}{\Character{0}}%
\TextAt{17}{216.456}{420.433}{1}%
\TextAt{9}{221.188}{415.496}{\Character{41}}%
\TextAt{21}{225.741}{420.433}{\Character{27}}%
\TextAt{13}{231.364}{415.496}{f}%
\TextAt{9}{238.411}{415.496}{\Character{40}}%
\TextAt{13}{242.963}{415.496}{\Character{27}}%
\TextAt{9}{250.045}{415.496}{\Character{41}}%
\TextAt{9}{257.9}{415.496}{\Character{61}}%
\TextAt{13}{270.344}{415.496}{f}%
\TextAt{9}{277.39}{415.496}{\Character{40}}%
\TextAt{13}{281.943}{415.496}{\Character{27}}%
\TextAt{21}{289.025}{420.433}{k}%
\TextAt{3}{293.647}{420.433}{\Character{0}}%
\TextAt{17}{300.233}{420.433}{2}%
\TextAt{9}{304.966}{415.496}{\Character{41}}%
\TextAt{21}{309.518}{420.433}{\Character{27}}%
\TextAt{24}{314.643}{423.245}{2}%
\TextAt{13}{319.292}{415.496}{f}%
\TextAt{9}{326.339}{415.496}{\Character{40}}%
\TextAt{13}{330.891}{415.496}{\Character{27}}%
\TextAt{9}{337.974}{415.496}{\Character{41}}%
\TextAt{21}{342.526}{420.433}{\Character{27}}%
\TextAt{13}{348.149}{415.496}{f}%
\TextAt{9}{355.196}{415.496}{\Character{40}}%
\TextAt{13}{359.748}{415.496}{\Character{27}}%
\TextAt{9}{366.831}{415.496}{\Character{41}}%
\TextAt{9}{374.686}{415.496}{\Character{61}}%
\TextAt{16}{387.129}{415.496}{\Character{1}}%
\TextAt{16}{392.425}{415.496}{\Character{1}}%
\TextAt{16}{397.733}{415.496}{\Character{1}}%
\TextAt{9}{404.398}{415.496}{\Character{61}}%
\TextAt{21}{416.823}{430.44}{k}%
\TextAt{3}{421.445}{430.44}{\Character{0}}%
\TextAt{17}{428.031}{430.44}{1}%
\TextAt{20}{416.906}{426.854}{\Character{89}}%
\TextAt{21}{417.692}{401.66}{i}%
\TextAt{17}{420.576}{401.66}{\Character{61}0}%
\TextAt{13}{434.258}{415.496}{f}%
\TextAt{9}{441.304}{415.496}{\Character{40}}%
\TextAt{13}{445.857}{415.496}{\Character{27}}%
\TextAt{9}{452.939}{415.496}{\Character{41}}%
\TextAt{21}{457.491}{420.433}{\Character{27}}%
\TextAt{22}{462.617}{423.245}{i}%
\TextAt{13}{466.277}{415.496}{\Character{59}}%
\TextAt{8}{79.226}{378.715}{and}%
\TextAt{13}{101.989}{378.715}{f}%
\TextAt{9}{109.035}{378.715}{\Character{40}1\Character{41}}%
\TextAt{9}{127.279}{378.715}{\Character{61}}%
\TextAt{9}{139.7}{378.715}{1}%
\TextAt{8}{149.49}{378.715}{because}%
\TextAt{13}{192.791}{378.715}{f}%
\TextAt{9}{199.838}{378.715}{\Character{40}1\Character{41}}%
\TextAt{9}{218.082}{378.715}{\Character{61}}%
\TextAt{13}{230.542}{378.715}{f}%
\TextAt{9}{237.588}{378.715}{\Character{40}1\Character{41}}%
\TextAt{13}{252.546}{378.715}{f}%
\TextAt{9}{259.592}{378.715}{\Character{40}1\Character{41}}%
\TextAt{8}{274.55}{378.715}{.}%
\TextAt{8}{282.998}{378.715}{If}%
\TextAt{13}{294.69}{378.715}{n}%
\TextAt{9}{304.999}{378.715}{\Character{61}}%
\TextAt{8}{317.424}{378.715}{\char35\relax }%
\TextAt{13}{327.18}{378.715}{G}%
\TextAt{8}{336.413}{378.715}{,}%
\TextAt{8}{343.562}{378.715}{then}%
\TextAt{13}{199.679}{341.668}{N}%
\TextAt{21}{209.053}{339.874}{G}%
\TextAt{13}{216.149}{341.668}{f}%
\TextAt{9}{223.196}{341.668}{\Character{40}}%
\TextAt{13}{227.748}{341.668}{\Character{27}}%
\TextAt{9}{234.831}{341.668}{\Character{41}}%
\TextAt{9}{242.686}{341.668}{\Character{61}}%
\TextAt{21}{255.129}{356.612}{n}%
\TextAt{3}{260.267}{356.612}{\Character{0}}%
\TextAt{17}{266.854}{356.612}{1}%
\TextAt{20}{255.471}{353.025}{\Character{89}}%
\TextAt{21}{256.257}{327.831}{i}%
\TextAt{17}{259.14}{327.831}{\Character{61}0}%
\TextAt{13}{273.081}{341.668}{f}%
\TextAt{9}{280.127}{341.668}{\Character{40}}%
\TextAt{13}{284.679}{341.668}{\Character{27}}%
\TextAt{9}{291.762}{341.668}{\Character{41}}%
\TextAt{21}{296.314}{346.604}{\Character{27}}%
\TextAt{22}{301.439}{349.417}{i}%
\TextAt{9}{308.42}{341.668}{\Character{61}}%
\TextAt{13}{320.846}{341.668}{f}%
\TextAt{9}{327.892}{341.668}{\Character{40}}%
\TextAt{13}{332.444}{341.668}{\Character{27}}%
\TextAt{21}{339.527}{346.604}{n}%
\TextAt{9}{345.163}{341.668}{\Character{41}}%
\TextAt{9}{353.018}{341.668}{\Character{61}}%
\TextAt{13}{365.462}{341.668}{f}%
\TextAt{9}{372.508}{341.668}{\Character{40}1\Character{41}}%
\TextAt{9}{390.752}{341.668}{\Character{61}}%
\TextAt{9}{403.173}{341.668}{1}%
\TextAt{13}{409.065}{341.668}{\Character{59}}%
\TextAt{8}{79.226}{302.741}{so}%
\TextAt{8}{94.016}{302.741}{that}%
\TextAt{13}{119.804}{302.741}{f}%
\TextAt{9}{126.85}{302.741}{\Character{40}}%
\TextAt{13}{131.403}{302.741}{\Character{27}}%
\TextAt{9}{138.485}{302.741}{\Character{41}}%
\NativeMathAt{147.073}{302.741}{11.955}{\in}%
\TextAt{21}{159.079}{300.948}{N}%
\TextAt{22}{165.8}{299.608}{G}%
\TextAt{13}{172.493}{302.741}{A}%
\TextAt{9}{185.304}{302.741}{\Character{61}}%
\TextAt{16}{198.444}{302.741}{\Character{102}}%
\TextAt{13}{204.422}{302.741}{a}%
\NativeMathAt{214.602}{302.741}{11.955}{\in}%
\TextAt{13}{226.608}{302.741}{A}%
\TextAt{16}{239.419}{302.741}{\Character{106}}%
\TextAt{13}{246.775}{302.741}{N}%
\TextAt{21}{256.149}{300.948}{G}%
\TextAt{13}{263.246}{302.741}{a}%
\TextAt{9}{273.427}{302.741}{\Character{61}}%
\TextAt{20}{286.567}{311.707}{\Character{81}}%
\TextAt{21}{297.858}{308.552}{n}%
\TextAt{3}{302.996}{308.552}{\Character{0}}%
\TextAt{17}{309.583}{308.552}{1}%
\TextAt{21}{297.858}{299.254}{i}%
\TextAt{17}{300.741}{299.254}{\Character{61}0}%
\TextAt{13}{316.307}{302.741}{a}%
\TextAt{21}{322.452}{307.079}{\Character{27}}%
\TextAt{22}{327.578}{309.892}{i}%
\TextAt{9}{335.273}{302.741}{\Character{61}}%
\TextAt{9}{348.4}{302.741}{1}%
\TextAt{16}{354.266}{302.741}{\Character{103}}%
\TextAt{8}{360.244}{302.741}{.}%
\TextAt{8}{369.951}{302.741}{Conversely}%
\TextAt{8}{429.62}{302.741}{we}%
\TextAt{8}{447.267}{302.741}{obtain,}%
\TextAt{8}{487.47}{302.741}{for}%
\TextAt{8}{505.775}{302.741}{every}%
\TextAt{13}{79.226}{285.695}{a}%
\NativeMathAt{88.692}{285.695}{11.955}{\in}%
\TextAt{13}{99.983}{285.695}{A}%
\TextAt{8}{112.656}{285.695}{such}%
\TextAt{8}{139.053}{285.695}{that}%
\TextAt{13}{164.422}{285.695}{N}%
\TextAt{21}{173.795}{283.902}{G}%
\TextAt{13}{180.892}{285.695}{a}%
\TextAt{9}{190.358}{285.695}{\Character{61}}%
\TextAt{9}{202.768}{285.695}{1}%
\TextAt{8}{208.636}{285.695}{,}%
\TextAt{8}{215.786}{285.695}{a}%
\TextAt{8}{225.541}{285.695}{1-cocycle}%
\TextAt{8}{275.609}{285.695}{by}%
\TextAt{8}{291.868}{285.695}{putting}%
\TextAt{13}{333.496}{285.695}{f}%
\TextAt{9}{340.543}{285.695}{\Character{40}}%
\TextAt{13}{345.095}{285.695}{\Character{27}}%
\TextAt{9}{352.177}{285.695}{\Character{41}}%
\TextAt{9}{360.032}{285.695}{\Character{61}}%
\TextAt{13}{372.476}{285.695}{a}%
\TextAt{8}{382.518}{285.695}{and}%
\TextAt{13}{266.14}{248.249}{f}%
\TextAt{9}{273.186}{248.249}{\Character{40}}%
\TextAt{13}{277.739}{248.249}{\Character{27}}%
\TextAt{21}{284.821}{253.185}{k}%
\TextAt{9}{289.941}{248.249}{\Character{41}}%
\TextAt{9}{297.796}{248.249}{\Character{61}}%
\TextAt{21}{310.239}{263.193}{k}%
\TextAt{3}{314.861}{263.193}{\Character{0}}%
\TextAt{17}{321.447}{263.193}{1}%
\TextAt{20}{310.322}{259.606}{\Character{89}}%
\TextAt{21}{311.109}{234.412}{i}%
\TextAt{17}{313.992}{234.412}{\Character{61}0}%
\TextAt{13}{327.674}{248.249}{a}%
\TextAt{21}{333.819}{253.185}{\Character{27}}%
\TextAt{22}{338.944}{255.998}{i}%
\TextAt{13}{342.604}{248.249}{\Character{58}}%
\TextAt{8}{79.226}{211.468}{The}%
\TextAt{8}{104.65}{211.468}{reader}%
\TextAt{8}{141.79}{211.468}{is}%
\TextAt{8}{154.925}{211.468}{invited}%
\TextAt{8}{195.317}{211.468}{to}%
\TextAt{8}{210.997}{211.468}{check}%
\TextAt{8}{243.905}{211.468}{this.}%
\TextAt{8}{275.39}{211.468}{The}%
\TextAt{8}{300.814}{211.468}{map}%
\TextAt{13}{328.199}{211.468}{f}%
\NativeArrowAt{340.901}{211.467}{11.955}{11.955}{\mapsto}%
\TextAt{13}{358.512}{211.468}{f}%
\TextAt{9}{365.559}{211.467}{\Character{40}}%
\TextAt{13}{370.111}{211.468}{\Character{27}}%
\TextAt{9}{377.194}{211.467}{\Character{41}}%
\TextAt{8}{387.014}{211.468}{therefore}%
\TextAt{8}{437.484}{211.468}{is}%
\TextAt{8}{450.618}{211.468}{an}%
\TextAt{8}{468.247}{211.468}{isomorphism}%
\TextAt{8}{79.226}{194.422}{between}%
\TextAt{13}{125.264}{194.422}{Z}%
\TextAt{17}{134.135}{198.76}{1}%
\TextAt{9}{138.867}{194.422}{\Character{40}}%
\TextAt{13}{143.419}{194.422}{G\Character{59}}%
\TextAt{13}{157.873}{194.422}{A}%
\TextAt{9}{166.673}{194.422}{\Character{41}}%
\TextAt{8}{175.969}{194.422}{and}%
\TextAt{21}{199.579}{192.628}{N}%
\TextAt{22}{206.301}{191.289}{G}%
\TextAt{13}{212.993}{194.422}{A}%
\TextAt{8}{221.768}{194.422}{.}%
\TextAt{8}{232.759}{194.422}{This}%
\TextAt{8}{260.326}{194.422}{isomorphism}%
\TextAt{8}{329.592}{194.422}{maps}%
\TextAt{13}{361.069}{194.422}{B}%
\TextAt{17}{370.565}{198.76}{1}%
\TextAt{9}{375.297}{194.422}{\Character{40}}%
\TextAt{13}{379.85}{194.422}{G\Character{59}}%
\TextAt{13}{394.304}{194.422}{A}%
\TextAt{9}{403.103}{194.422}{\Character{41}}%
\TextAt{8}{412.4}{194.422}{onto}%
\TextAt{13}{439.596}{194.422}{I}%
\TextAt{21}{444.759}{192.628}{G}%
\TextAt{13}{451.856}{194.422}{A}%
\TextAt{8}{460.631}{194.422}{,}%
\TextAt{8}{468.839}{194.422}{because}%
\TextAt{13}{512.988}{194.422}{f}%
\NativeMathAt{524.799}{194.422}{11.955}{\in}%
\TextAt{13}{79.226}{177.376}{B}%
\TextAt{17}{88.722}{181.714}{1}%
\TextAt{9}{93.455}{177.376}{\Character{40}}%
\TextAt{13}{98.007}{177.376}{G\Character{59}}%
\TextAt{13}{112.461}{177.376}{A}%
\TextAt{9}{121.26}{177.376}{\Character{41}}%
\NativeArrowAt{132.454}{177.376}{11.955}{21.926}{\Longleftrightarrow}%
\TextAt{13}{161.014}{177.376}{f}%
\TextAt{9}{168.06}{177.376}{\Character{40}}%
\TextAt{13}{172.612}{177.376}{\Character{27}}%
\TextAt{21}{179.695}{181.714}{k}%
\TextAt{9}{184.815}{177.376}{\Character{41}}%
\TextAt{9}{192.67}{177.376}{\Character{61}}%
\TextAt{13}{205.113}{177.376}{a}%
\TextAt{21}{211.258}{181.714}{\Character{27}}%
\TextAt{22}{216.384}{184.527}{k}%
\TextAt{3}{220.957}{181.714}{\Character{0}}%
\TextAt{17}{227.543}{181.714}{1}%
\TextAt{8}{236.072}{177.376}{for}%
\TextAt{8}{253.857}{177.376}{some}%
\TextAt{8}{283.082}{177.376}{fixed}%
\TextAt{13}{311.267}{177.376}{a}%
\NativeArrowAt{324.054}{177.376}{11.955}{21.926}{\Longleftrightarrow}%
\TextAt{13}{352.614}{177.376}{f}%
\TextAt{9}{359.66}{177.376}{\Character{40}}%
\TextAt{13}{364.212}{177.376}{\Character{27}}%
\TextAt{9}{371.295}{177.376}{\Character{41}}%
\TextAt{9}{379.15}{177.376}{\Character{61}}%
\TextAt{13}{391.593}{177.376}{a}%
\TextAt{21}{397.738}{181.714}{\Character{27}}%
\TextAt{3}{402.864}{181.714}{\Character{0}}%
\TextAt{17}{409.45}{181.714}{1}%
\NativeArrowAt{420.824}{177.376}{11.955}{21.926}{\Longleftrightarrow}%
\TextAt{13}{449.384}{177.376}{f}%
\TextAt{9}{456.43}{177.376}{\Character{40}}%
\TextAt{13}{460.982}{177.376}{\Character{27}}%
\TextAt{9}{468.065}{177.376}{\Character{41}}%
\NativeMathAt{475.938}{177.376}{11.955}{\in}%
\TextAt{13}{487.229}{177.376}{I}%
\TextAt{21}{492.392}{175.582}{G}%
\TextAt{13}{499.489}{177.376}{A}%
\TextAt{8}{508.264}{177.376}{.}%
\Vector{q 0 G 0 g 0.399 w 452.334 -596.624 m 452.334 -588.555 l S Q}%
\Vector{q 0 G 0 g 0.399 w 452.534 -588.754 m 459.707 -588.754 l S Q}%
\Vector{q 0 G 0 g 0.399 w 452.534 -596.425 m 459.707 -596.425 l S Q}%
\Vector{q 0 G 0 g 0.399 w 459.906 -596.624 m 459.906 -588.555 l S Q}%
\TextAt{8}{97.159}{148.423}{Noether’s}%
\TextAt{8}{149.893}{148.423}{generalization}%
\TextAt{8}{224.697}{148.423}{of}%
\TextAt{8}{238.027}{148.423}{Hilbert’s}%
\TextAt{8}{286.206}{148.423}{theorem}%
\TextAt{8}{331.731}{148.423}{90}%
\TextAt{8}{347.345}{148.423}{now}%
\TextAt{8}{371.733}{148.423}{reads:}%
\TextAt{12}{79.226}{122.458}{Proposition}%
\TextAt{12}{151.937}{122.458}{(3.8).}%
\TextAt{8}{187.253}{122.458}{For}%
\TextAt{8}{208.222}{122.458}{a}%
\TextAt{8}{217.978}{122.458}{finite}%
\TextAt{8}{247.89}{122.458}{Galois}%
\TextAt{8}{283.803}{122.458}{field}%
\TextAt{8}{309.16}{122.458}{extension}%
\TextAt{13}{360.926}{122.458}{L}%
\TextAt{16}{372.211}{122.458}{\Character{106}}%
\TextAt{13}{378.853}{122.458}{K}%
\TextAt{8}{389.671}{122.458}{,}%
\TextAt{8}{396.82}{122.458}{one}%
\TextAt{8}{418.279}{122.458}{has}%
\TextAt{8}{439.153}{122.458}{that}%
\TextAt{13}{248.831}{93.626}{H}%
\TextAt{17}{259.453}{98.562}{1}%
\TextAt{9}{264.185}{93.626}{\Character{40}}%
\TextAt{13}{268.738}{93.626}{G}%
\TextAt{9}{277.971}{93.626}{\Character{40}}%
\TextAt{13}{282.524}{93.626}{L}%
\TextAt{16}{293.809}{93.626}{\Character{106}}%
\TextAt{13}{300.451}{93.626}{K}%
\TextAt{9}{311.268}{93.626}{\Character{41}}%
\TextAt{13}{315.821}{93.626}{\Character{59}}%
\TextAt{13}{321.045}{93.626}{L}%
\TextAt{3}{329.029}{98.562}{\Character{3}}%
\TextAt{9}{333.762}{93.626}{\Character{41}}%
\TextAt{9}{341.617}{93.626}{\Character{61}}%
\TextAt{9}{354.038}{93.626}{1}%
\TextAt{13}{359.913}{93.626}{\Character{58}}%
\TextAt{8}{297.211}{43.033}{258}%
\Destination{78.23}{809.815}{nav.page.258}%
\Destination{79.36}{300.655}{nav.print.283}%
\Destination{79.36}{501.855}{nav.number.4.3.7}%
\Destination{79.36}{136.455}{nav.number.4.3.8}%
}

% Page 271 (body page 259)
\Page{%
\TextAt{2}{79.226}{796.943}{IV.}%
\TextAt{2}{98.012}{796.943}{Abstract}%
\TextAt{2}{143.252}{796.943}{Class}%
\TextAt{2}{171.844}{796.943}{Field}%
\TextAt{2}{199.575}{796.943}{Theory}%
\Vector{q 0 G 0 g 0.299 w 7.226 18.729 m 460.769 18.729 l S Q}%
\TextAt{19}{79.226}{761.124}{Proof.}%
\TextAt{8}{116}{761.124}{Let}%
\TextAt{13}{136.97}{761.124}{f}%
\TextAt{9}{147.337}{761.124}{\Character{58}}%
\TextAt{13}{153.909}{761.124}{G}%
\NativeMathAt{166.464}{761.124}{11.955}{\rightarrow}%
\TextAt{13}{181.74}{761.124}{L}%
\TextAt{3}{189.704}{765.462}{\Character{3}}%
\TextAt{8}{198.334}{761.124}{be}%
\TextAt{8}{214.258}{761.124}{a}%
\TextAt{8}{224.014}{761.124}{1-cocycle.}%
\TextAt{8}{278.633}{761.124}{For}%
\TextAt{13}{299.602}{761.124}{c}%
\NativeMathAt{307.961}{761.124}{11.955}{\in}%
\TextAt{13}{319.252}{761.124}{L}%
\TextAt{3}{327.217}{765.462}{\Character{3}}%
\TextAt{8}{331.949}{761.124}{,}%
\TextAt{8}{339.098}{761.124}{we}%
\TextAt{8}{356.326}{761.124}{put}%
\TextAt{13}{254.908}{732.123}{\Character{11}}%
\TextAt{9}{265.751}{732.123}{\Character{61}}%
\TextAt{20}{289.428}{743.48}{\Character{88}}%
\TextAt{21}{278.176}{717.577}{\Character{27}}%
\NativeMathAt{283.302}{717.577}{7.97}{\in}%
\TextAt{21}{288.947}{717.577}{G}%
\TextAt{17}{295.546}{717.577}{\Character{40}}%
\TextAt{21}{298.839}{717.577}{L}%
\TextAt{3}{304.598}{717.577}{\Character{106}}%
\TextAt{21}{306.95}{717.577}{K}%
\TextAt{17}{314.655}{717.577}{\Character{41}}%
\TextAt{13}{319.941}{732.123}{f}%
\TextAt{9}{326.987}{732.123}{\Character{40}}%
\TextAt{13}{331.54}{732.123}{\Character{27}}%
\TextAt{9}{338.622}{732.123}{\Character{41}}%
\TextAt{13}{343.174}{732.123}{c}%
\TextAt{21}{348.212}{737.059}{\Character{27}}%
\TextAt{13}{353.836}{732.123}{\Character{58}}%
\TextAt{8}{79.226}{692.471}{Since}%
\TextAt{8}{109.48}{692.471}{the}%
\TextAt{8}{129.333}{692.471}{automorphisms}%
\TextAt{13}{211.113}{692.471}{\Character{27}}%
\TextAt{8}{221.79}{692.471}{are}%
\TextAt{8}{240.997}{692.471}{linearly}%
\TextAt{8}{282.644}{692.471}{independent}%
\TextAt{8}{348.668}{692.471}{(see}%
\TextAt{8}{371.832}{692.471}{[15],}%
\TextAt{8}{396.958}{692.471}{chap.}%
\TextAt{8}{429.048}{692.471}{5,}%
\TextAt{8}{441.812}{692.471}{§7,}%
\TextAt{8}{460.255}{692.471}{no.}%
\TextAt{8}{480.964}{692.471}{5),}%
\TextAt{8}{498.284}{692.471}{we}%
\TextAt{8}{515.207}{692.471}{can}%
\TextAt{8}{79.226}{675.425}{choose}%
\TextAt{13}{116.359}{675.425}{c}%
\NativeMathAt{124.718}{675.425}{11.955}{\in}%
\TextAt{13}{136.009}{675.425}{L}%
\TextAt{3}{143.973}{679.764}{\Character{3}}%
\TextAt{8}{152.603}{675.425}{such}%
\TextAt{8}{179}{675.425}{that}%
\TextAt{13}{204.369}{675.425}{\Character{11}}%
\NativeMathAt{215.211}{675.425}{11.955}{\neq}%
\TextAt{9}{227.621}{675.425}{0}%
\TextAt{8}{233.49}{675.425}{.}%
\TextAt{8}{241.938}{675.425}{For}%
\TextAt{13}{262.908}{675.425}{\Character{28}}%
\NativeMathAt{272.647}{675.425}{11.955}{\in}%
\TextAt{13}{283.938}{675.425}{G}%
\TextAt{9}{293.172}{675.425}{\Character{40}}%
\TextAt{13}{297.724}{675.425}{L}%
\TextAt{16}{309.009}{675.425}{\Character{106}}%
\TextAt{13}{315.651}{675.425}{K}%
\TextAt{9}{326.469}{675.425}{\Character{41}}%
\TextAt{8}{331.021}{675.425}{,}%
\TextAt{8}{338.17}{675.425}{we}%
\TextAt{8}{355.398}{675.425}{obtain}%
\TextAt{13}{173.122}{646.424}{\Character{11}}%
\TextAt{21}{180.643}{651.36}{\Character{28}}%
\TextAt{9}{189.141}{646.424}{\Character{61}}%
\TextAt{20}{201.566}{657.782}{\Character{88}}%
\TextAt{21}{207.638}{632.676}{\Character{27}}%
\TextAt{13}{220.827}{646.424}{f}%
\TextAt{9}{227.874}{646.424}{\Character{40}}%
\TextAt{13}{232.426}{646.424}{\Character{27}}%
\TextAt{9}{239.508}{646.424}{\Character{41}}%
\TextAt{21}{244.061}{651.36}{\Character{28}}%
\TextAt{13}{249.237}{646.424}{c}%
\TextAt{21}{254.275}{651.36}{\Character{27}}%
\TextAt{21}{259.392}{651.36}{\Character{28}}%
\TextAt{9}{267.898}{646.424}{\Character{61}}%
\TextAt{20}{280.323}{657.782}{\Character{88}}%
\TextAt{21}{286.395}{632.676}{\Character{27}}%
\TextAt{13}{299.584}{646.424}{f}%
\TextAt{9}{306.631}{646.424}{\Character{40}}%
\TextAt{13}{311.183}{646.424}{\Character{28}}%
\TextAt{9}{317.602}{646.424}{\Character{41}}%
\TextAt{3}{322.154}{651.36}{\Character{0}}%
\TextAt{17}{328.741}{651.36}{1}%
\TextAt{13}{333.473}{646.424}{f}%
\TextAt{9}{340.52}{646.424}{\Character{40}}%
\TextAt{13}{345.072}{646.424}{\Character{27}}%
\TextAt{13}{352.137}{646.424}{\Character{28}}%
\TextAt{9}{358.573}{646.424}{\Character{41}}%
\TextAt{13}{363.126}{646.424}{c}%
\TextAt{21}{368.164}{651.36}{\Character{27}}%
\TextAt{21}{373.281}{651.36}{\Character{28}}%
\TextAt{9}{381.786}{646.424}{\Character{61}}%
\TextAt{13}{394.212}{646.424}{f}%
\TextAt{9}{401.258}{646.424}{\Character{40}}%
\TextAt{13}{405.811}{646.424}{\Character{28}}%
\TextAt{9}{412.229}{646.424}{\Character{41}}%
\TextAt{3}{416.782}{651.36}{\Character{0}}%
\TextAt{17}{423.368}{651.36}{1}%
\TextAt{13}{428.101}{646.424}{\Character{11}\Character{59}}%
\TextAt{8}{79.226}{609.054}{i.e.,}%
\TextAt{13}{101.331}{609.054}{f}%
\TextAt{9}{108.378}{609.054}{\Character{40}}%
\TextAt{13}{112.93}{609.054}{\Character{28}}%
\TextAt{9}{119.349}{609.054}{\Character{41}}%
\TextAt{9}{127.204}{609.054}{\Character{61}}%
\TextAt{13}{139.648}{609.054}{\Character{12}}%
\TextAt{21}{146.919}{613.392}{\Character{28}}%
\TextAt{3}{151.597}{613.392}{\Character{0}}%
\TextAt{17}{158.184}{613.392}{1}%
\TextAt{8}{166.813}{609.054}{with}%
\TextAt{13}{193.473}{609.054}{\Character{12}}%
\TextAt{9}{204.066}{609.054}{\Character{61}}%
\TextAt{13}{216.491}{609.054}{\Character{11}}%
\TextAt{3}{224.013}{613.392}{\Character{0}}%
\TextAt{17}{230.599}{613.392}{1}%
\TextAt{8}{235.332}{609.054}{.}%
\Vector{q 0 G 0 g 0.399 w 452.334 -164.946 m 452.334 -156.876 l S Q}%
\Vector{q 0 G 0 g 0.399 w 452.534 -157.076 m 459.707 -157.076 l S Q}%
\Vector{q 0 G 0 g 0.399 w 452.534 -164.747 m 459.707 -164.747 l S Q}%
\Vector{q 0 G 0 g 0.399 w 459.906 -164.946 m 459.906 -156.876 l S Q}%
\TextAt{8}{97.159}{580.053}{This}%
\TextAt{8}{123.879}{580.053}{proposition}%
\TextAt{8}{185.412}{580.053}{will}%
\TextAt{8}{207.517}{580.053}{only}%
\TextAt{8}{233.209}{580.053}{be}%
\TextAt{8}{249.133}{580.053}{applied}%
\TextAt{8}{290.103}{580.053}{once}%
\TextAt{8}{316.763}{580.053}{in}%
\TextAt{8}{330.416}{580.053}{this}%
\TextAt{8}{353.239}{580.053}{book}%
\TextAt{8}{382.182}{580.053}{(see}%
\TextAt{8}{405.65}{580.053}{chap.}%
\TextAt{8}{437.841}{580.053}{VI,}%
\TextAt{8}{457.986}{580.053}{(2.5)).}%
\TextAt{12}{79.226}{554.04}{Exercise}%
\TextAt{12}{132.142}{554.04}{IV.3.1.}%
\TextAt{8}{175.651}{554.04}{Show}%
\TextAt{8}{206.296}{554.04}{that}%
\TextAt{8}{231.417}{554.04}{Hilbert}%
\TextAt{8}{271.483}{554.04}{90}%
\TextAt{8}{286.849}{554.04}{in}%
\TextAt{8}{300.255}{554.04}{Noether’s}%
\TextAt{8}{352.742}{554.04}{formulation}%
\TextAt{8}{415.594}{554.04}{also}%
\TextAt{8}{438.826}{554.04}{holds}%
\TextAt{8}{469.208}{554.04}{for}%
\TextAt{8}{486.846}{554.04}{the}%
\TextAt{8}{506.755}{554.04}{addi-}%
\TextAt{8}{79.226}{536.994}{tive}%
\TextAt{8}{101.989}{536.994}{group}%
\TextAt{13}{135.164}{536.994}{L}%
\TextAt{8}{147.026}{536.994}{of}%
\TextAt{8}{160.356}{536.994}{a}%
\TextAt{8}{170.112}{536.994}{Galois}%
\TextAt{8}{206.025}{536.994}{extension}%
\TextAt{13}{257.791}{536.994}{L}%
\TextAt{16}{269.076}{536.994}{\Character{106}}%
\TextAt{13}{275.718}{536.994}{K}%
\TextAt{8}{286.536}{536.994}{.}%
\TextAt{12}{97.159}{516.96}{Hint:}%
\TextAt{8}{132.698}{516.96}{Use}%
\TextAt{8}{155.185}{516.96}{the}%
\TextAt{8}{175.342}{516.96}{normal}%
\TextAt{8}{215.021}{516.96}{basis}%
\TextAt{8}{243.761}{516.96}{theorem.}%
\TextAt{12}{79.226}{484.97}{Exercise}%
\TextAt{12}{133.603}{484.97}{IV.3.2.}%
\TextAt{8}{178.382}{484.97}{Let}%
\TextAt{13}{200.374}{484.97}{k}%
\TextAt{8}{211.784}{484.97}{be}%
\TextAt{8}{228.731}{484.97}{a}%
\TextAt{8}{239.509}{484.97}{field}%
\TextAt{8}{265.889}{484.97}{of}%
\TextAt{8}{280.242}{484.97}{characteristic}%
\TextAt{13}{353.199}{484.97}{p}%
\TextAt{8}{363.994}{484.97}{and}%
\TextAt{9}{388.098}{488.125}{\Character{22}}%
\TextAt{13}{387.779}{484.97}{k}%
\TextAt{8}{399.189}{484.97}{its}%
\TextAt{8}{416.531}{484.97}{separable}%
\TextAt{8}{468.997}{484.97}{closure.}%
\TextAt{8}{515.697}{484.97}{For}%
\TextAt{8}{79.226}{467.924}{fixed}%
\TextAt{13}{107.278}{467.924}{n}%
\NativeMathAt{117.586}{467.924}{11.955}{\geq}%
\TextAt{9}{130.205}{467.924}{1}%
\TextAt{8}{136.058}{467.924}{,}%
\TextAt{8}{143.02}{467.924}{consider}%
\TextAt{8}{188.371}{467.924}{in}%
\TextAt{8}{201.789}{467.924}{the}%
\TextAt{8}{221.711}{467.924}{ring}%
\TextAt{8}{245.542}{467.924}{of}%
\TextAt{8}{258.638}{467.924}{Witt}%
\TextAt{8}{286.689}{467.924}{vectors}%
\TextAt{13}{326.194}{467.924}{W}%
\TextAt{9}{338.871}{467.924}{\Character{40}}%
\TextAt{9}{343.742}{471.079}{\Character{22}}%
\TextAt{13}{343.424}{467.924}{k}%
\TextAt{9}{349.913}{467.924}{\Character{41}}%
\TextAt{8}{358.128}{467.924}{(see}%
\TextAt{8}{381.362}{467.924}{chap.}%
\TextAt{8}{413.475}{467.924}{II,}%
\TextAt{8}{429.153}{467.924}{§4,}%
\TextAt{8}{447.651}{467.924}{exercise}%
\TextAt{8}{490.718}{467.924}{2–6)}%
\TextAt{8}{516.51}{467.924}{the}%
\TextAt{8}{79.226}{450.878}{additive}%
\TextAt{8}{124.248}{450.878}{group}%
\TextAt{13}{157.566}{450.878}{W}%
\TextAt{21}{168.617}{449.084}{n}%
\TextAt{9}{174.254}{450.878}{\Character{40}}%
\TextAt{9}{179.124}{454.033}{\Character{22}}%
\TextAt{13}{178.806}{450.878}{k}%
\TextAt{9}{185.296}{450.878}{\Character{41}}%
\TextAt{8}{193.887}{450.878}{of}%
\TextAt{8}{207.36}{450.878}{truncated}%
\TextAt{8}{260.834}{450.878}{Witt}%
\TextAt{8}{289.262}{450.878}{vectors}%
\TextAt{13}{329.143}{450.878}{a}%
\TextAt{9}{338.851}{450.878}{\Character{61}}%
\TextAt{9}{351.512}{450.878}{\Character{40}}%
\TextAt{13}{356.071}{450.878}{a}%
\TextAt{17}{362.216}{449.084}{0}%
\TextAt{13}{366.949}{450.878}{\Character{59}}%
\TextAt{13}{372.173}{450.878}{a}%
\TextAt{17}{378.338}{449.084}{1}%
\TextAt{13}{383.07}{450.878}{\Character{59}}%
\TextAt{13}{388.294}{450.878}{\Character{58}}%
\TextAt{13}{393.531}{450.878}{\Character{58}}%
\TextAt{13}{398.755}{450.878}{\Character{58}}%
\TextAt{13}{403.992}{450.878}{\Character{59}}%
\TextAt{13}{409.228}{450.878}{a}%
\TextAt{21}{415.436}{449.084}{n}%
\TextAt{3}{420.574}{449.084}{\Character{0}}%
\TextAt{17}{427.16}{449.084}{1}%
\TextAt{9}{431.893}{450.878}{\Character{41}}%
\TextAt{8}{436.445}{450.878}{.}%
\TextAt{8}{445.32}{450.878}{Show}%
\TextAt{8}{476.354}{450.878}{that}%
\TextAt{8}{501.865}{450.878}{axiom}%
\TextAt{8}{79.226}{433.832}{(3.1)}%
\TextAt{8}{107.201}{433.832}{holds}%
\TextAt{8}{137.83}{433.832}{for}%
\TextAt{8}{155.715}{433.832}{the}%
\TextAt{13}{175.872}{433.832}{G}%
\TextAt{9}{185.105}{433.832}{\Character{40}}%
\TextAt{9}{189.976}{436.987}{\Character{22}}%
\TextAt{13}{189.658}{433.832}{k}%
\TextAt{16}{199.468}{433.832}{\Character{106}}%
\TextAt{13}{206.11}{433.832}{k}%
\TextAt{9}{212.599}{433.832}{\Character{41}}%
\TextAt{8}{217.152}{433.832}{-module}%
\TextAt{13}{262.342}{433.832}{A}%
\TextAt{9}{274.438}{433.832}{\Character{61}}%
\TextAt{13}{286.864}{433.832}{W}%
\TextAt{21}{297.915}{432.039}{n}%
\TextAt{9}{303.552}{433.832}{\Character{40}}%
\TextAt{9}{308.422}{436.987}{\Character{22}}%
\TextAt{13}{308.104}{433.832}{k}%
\TextAt{9}{314.593}{433.832}{\Character{41}}%
\TextAt{8}{319.146}{433.832}{.}%
\TextAt{12}{79.226}{401.842}{Exercise}%
\TextAt{12}{132.427}{401.842}{IV.3.3.}%
\TextAt{8}{176.183}{401.842}{Show}%
\TextAt{8}{207.075}{401.842}{that}%
\TextAt{8}{232.444}{401.842}{the}%
\TextAt{8}{252.6}{401.842}{operator}%
\TextAt{13}{203.164}{372.841}{\Character{125}}%
\TextAt{9}{213.937}{372.841}{\Character{58}}%
\TextAt{13}{220.51}{372.841}{W}%
\TextAt{21}{231.561}{371.047}{n}%
\TextAt{9}{237.198}{372.841}{\Character{40}}%
\TextAt{9}{242.068}{375.996}{\Character{22}}%
\TextAt{13}{241.75}{372.841}{k}%
\TextAt{9}{248.239}{372.841}{\Character{41}}%
\NativeArrowAt{256.113}{372.841}{11.955}{19.26}{\longrightarrow}%
\TextAt{13}{278.695}{372.841}{W}%
\TextAt{21}{289.746}{371.047}{n}%
\TextAt{9}{295.382}{372.841}{\Character{40}}%
\TextAt{9}{300.253}{375.996}{\Character{22}}%
\TextAt{13}{299.935}{372.841}{k}%
\TextAt{9}{306.424}{372.841}{\Character{41}}%
\TextAt{13}{310.977}{372.841}{\Character{59}}%
\TextAt{13}{340.112}{372.841}{\Character{125}a}%
\TextAt{9}{357.049}{372.841}{\Character{61}}%
\TextAt{13}{369.475}{372.841}{F}%
\TextAt{13}{378.657}{372.841}{a}%
\TextAt{16}{387.48}{372.841}{\Character{0}}%
\TextAt{13}{399.435}{372.841}{a\Character{59}}%
\TextAt{8}{79.226}{343.84}{is}%
\TextAt{8}{90.465}{343.84}{a}%
\TextAt{8}{99.696}{343.84}{homomorphism}%
\TextAt{8}{181.841}{343.84}{with}%
\TextAt{8}{207.977}{343.84}{cyclic}%
\TextAt{8}{239.635}{343.84}{kernel}%
\TextAt{13}{273.577}{343.84}{\Character{22}}%
\TextAt{21}{280.62}{342.046}{\Character{125}}%
\TextAt{8}{289.859}{343.84}{of}%
\TextAt{8}{302.664}{343.84}{order}%
\TextAt{13}{332.709}{343.84}{p}%
\TextAt{21}{338.584}{348.178}{n}%
\TextAt{8}{344.22}{343.84}{.}%
\TextAt{8}{352.493}{343.84}{Discuss}%
\TextAt{8}{393.609}{343.84}{the}%
\TextAt{8}{413.24}{343.84}{corresponding}%
\TextAt{8}{487.902}{343.84}{Kummer}%
\TextAt{8}{79.226}{326.794}{theory}%
\TextAt{8}{115.976}{326.794}{for}%
\TextAt{8}{133.861}{326.794}{the}%
\TextAt{8}{154.018}{326.794}{abelian}%
\TextAt{8}{194.665}{326.794}{extensions}%
\TextAt{8}{251.046}{326.794}{of}%
\TextAt{8}{264.376}{326.794}{exponent}%
\TextAt{13}{314.779}{326.794}{p}%
\TextAt{21}{320.654}{331.132}{n}%
\TextAt{8}{326.29}{326.794}{.}%
\TextAt{12}{79.226}{294.804}{Exercise}%
\TextAt{12}{132.427}{294.804}{IV.3.4.}%
\TextAt{8}{176.183}{294.804}{Let}%
\TextAt{13}{197.152}{294.804}{G}%
\TextAt{8}{210.283}{294.804}{be}%
\TextAt{8}{226.207}{294.804}{a}%
\TextAt{8}{235.963}{294.804}{profinite}%
\TextAt{8}{282.791}{294.804}{group}%
\TextAt{8}{315.967}{294.804}{and}%
\TextAt{13}{338.729}{294.804}{A}%
\TextAt{8}{351.402}{294.804}{a}%
\TextAt{8}{361.157}{294.804}{continuous}%
\TextAt{13}{419.762}{294.804}{G}%
\TextAt{8}{428.995}{294.804}{-module.}%
\TextAt{8}{478.737}{294.804}{Put}%
\TextAt{13}{223.442}{265.803}{H}%
\TextAt{17}{234.065}{270.739}{1}%
\TextAt{9}{238.797}{265.803}{\Character{40}}%
\TextAt{13}{243.35}{265.803}{G\Character{59}}%
\TextAt{13}{257.804}{265.803}{A}%
\TextAt{9}{266.603}{265.803}{\Character{41}}%
\TextAt{9}{274.458}{265.803}{\Character{61}}%
\TextAt{13}{286.901}{265.803}{Z}%
\TextAt{17}{295.772}{270.739}{1}%
\TextAt{9}{300.505}{265.803}{\Character{40}}%
\TextAt{13}{305.057}{265.803}{G\Character{59}}%
\TextAt{13}{319.511}{265.803}{A}%
\TextAt{9}{328.31}{265.803}{\Character{41}\Character{47}}%
\TextAt{13}{338.715}{265.803}{B}%
\TextAt{17}{348.212}{270.739}{1}%
\TextAt{9}{352.944}{265.803}{\Character{40}}%
\TextAt{13}{357.496}{265.803}{G\Character{59}}%
\TextAt{13}{371.95}{265.803}{A}%
\TextAt{9}{380.749}{265.803}{\Character{41}}%
\TextAt{13}{385.302}{265.803}{\Character{59}}%
\TextAt{8}{79.226}{236.801}{where}%
\TextAt{13}{113.338}{236.801}{Z}%
\TextAt{17}{122.208}{241.14}{1}%
\TextAt{9}{126.941}{236.801}{\Character{40}}%
\TextAt{13}{131.493}{236.801}{G\Character{59}}%
\TextAt{13}{145.947}{236.801}{A}%
\TextAt{9}{154.746}{236.801}{\Character{41}}%
\TextAt{8}{163.498}{236.801}{consists}%
\TextAt{8}{206.911}{236.801}{of}%
\TextAt{8}{220.543}{236.801}{all}%
\TextAt{8}{237.104}{236.801}{continuous}%
\TextAt{8}{296.011}{236.801}{maps}%
\TextAt{13}{326.942}{236.801}{f}%
\TextAt{9}{337.824}{236.801}{\Character{58}}%
\TextAt{13}{344.912}{236.801}{G}%
\NativeMathAt{357.982}{236.801}{11.955}{\rightarrow}%
\TextAt{13}{373.773}{236.801}{A}%
\TextAt{8}{386.748}{236.801}{(with}%
\TextAt{8}{418.265}{236.801}{respect}%
\TextAt{8}{458.617}{236.801}{to}%
\TextAt{8}{473.229}{236.801}{the}%
\TextAt{8}{493.688}{236.801}{discrete}%
\TextAt{8}{79.226}{219.755}{topology}%
\TextAt{8}{128.058}{219.755}{on}%
\TextAt{13}{145.005}{219.755}{A}%
\TextAt{8}{153.781}{219.755}{)}%
\TextAt{8}{162.921}{219.755}{such}%
\TextAt{8}{190.007}{219.755}{that}%
\TextAt{13}{216.064}{219.755}{f}%
\TextAt{9}{223.11}{219.755}{\Character{40}}%
\TextAt{13}{227.663}{219.755}{\Character{27}}%
\TextAt{13}{234.729}{219.755}{\Character{28}}%
\TextAt{9}{241.164}{219.755}{\Character{41}}%
\TextAt{9}{250.19}{219.755}{\Character{61}}%
\TextAt{13}{263.809}{219.755}{f}%
\TextAt{9}{270.855}{219.755}{\Character{40}}%
\TextAt{13}{275.407}{219.755}{\Character{27}}%
\TextAt{9}{282.49}{219.755}{\Character{41}}%
\TextAt{21}{287.042}{224.094}{\Character{28}}%
\TextAt{13}{292.219}{219.755}{f}%
\TextAt{9}{299.265}{219.755}{\Character{40}}%
\TextAt{13}{303.817}{219.755}{\Character{28}}%
\TextAt{9}{310.236}{219.755}{\Character{41}}%
\TextAt{8}{314.789}{219.755}{,}%
\TextAt{8}{322.798}{219.755}{and}%
\TextAt{13}{346.249}{219.755}{B}%
\TextAt{17}{355.746}{224.094}{1}%
\TextAt{9}{360.478}{219.755}{\Character{40}}%
\TextAt{13}{365.03}{219.755}{G\Character{59}}%
\TextAt{13}{379.484}{219.755}{A}%
\TextAt{9}{388.283}{219.755}{\Character{41}}%
\TextAt{8}{397.421}{219.755}{consists}%
\TextAt{8}{441.22}{219.755}{of}%
\TextAt{8}{455.238}{219.755}{all}%
\TextAt{8}{472.186}{219.755}{functions}%
\TextAt{8}{523.337}{219.755}{of}%
\TextAt{8}{79.226}{202.709}{the}%
\TextAt{8}{99.246}{202.709}{form}%
\TextAt{13}{126.751}{202.709}{f}%
\TextAt{21}{132.521}{200.916}{a}%
\TextAt{9}{137.517}{202.709}{\Character{40}}%
\TextAt{13}{142.069}{202.709}{\Character{27}}%
\TextAt{9}{149.151}{202.709}{\Character{41}}%
\TextAt{9}{157.006}{202.709}{\Character{61}}%
\TextAt{13}{169.45}{202.709}{a}%
\TextAt{21}{175.595}{207.048}{\Character{27}}%
\TextAt{3}{180.72}{207.048}{\Character{0}}%
\TextAt{17}{187.307}{207.048}{1}%
\TextAt{8}{192.039}{202.709}{,}%
\TextAt{13}{199.079}{202.709}{a}%
\NativeMathAt{208.545}{202.709}{11.955}{\in}%
\TextAt{13}{219.836}{202.709}{A}%
\TextAt{8}{228.612}{202.709}{.}%
\TextAt{8}{237.014}{202.709}{Show}%
\TextAt{8}{267.77}{202.709}{that}%
\TextAt{8}{293.003}{202.709}{if}%
\TextAt{14}{303.591}{202.709}{g}%
\TextAt{8}{313.374}{202.709}{is}%
\TextAt{8}{325.002}{202.709}{a}%
\TextAt{8}{334.621}{202.709}{closed}%
\TextAt{8}{369.011}{202.709}{normal}%
\TextAt{8}{408.554}{202.709}{subgroup}%
\TextAt{8}{459.216}{202.709}{of}%
\TextAt{13}{472.41}{202.709}{G}%
\TextAt{8}{481.643}{202.709}{,}%
\TextAt{8}{488.683}{202.709}{then}%
\TextAt{8}{515.207}{202.709}{one}%
\TextAt{8}{79.226}{185.663}{has}%
\TextAt{8}{100.1}{185.663}{an}%
\TextAt{8}{116.359}{185.663}{exact}%
\TextAt{8}{147.251}{185.663}{sequence}%
\TextAt{9}{184.434}{156.662}{1}%
\NativeArrowAt{193.608}{156.662}{11.955}{19.26}{\longrightarrow}%
\TextAt{13}{216.19}{156.662}{H}%
\TextAt{17}{226.813}{161.599}{1}%
\TextAt{9}{231.545}{156.662}{\Character{40}}%
\TextAt{13}{236.097}{156.662}{G}%
\TextAt{9}{245.331}{156.662}{\Character{47}}%
\TextAt{14}{251.184}{156.662}{g}%
\TextAt{13}{257.206}{156.662}{\Character{59}}%
\TextAt{13}{262.43}{156.662}{A}%
\TextAt{27}{271.225}{161.599}{g}%
\TextAt{9}{275.738}{156.662}{\Character{41}}%
\NativeArrowAt{283.611}{156.662}{11.955}{19.26}{\longrightarrow}%
\TextAt{13}{306.193}{156.662}{H}%
\TextAt{17}{316.816}{161.599}{1}%
\TextAt{9}{321.548}{156.662}{\Character{40}}%
\TextAt{13}{326.1}{156.662}{G\Character{59}}%
\TextAt{13}{340.554}{156.662}{A}%
\TextAt{9}{349.354}{156.662}{\Character{41}}%
\NativeArrowAt{357.227}{156.662}{11.955}{19.26}{\longrightarrow}%
\TextAt{13}{379.809}{156.662}{H}%
\TextAt{17}{390.431}{161.599}{1}%
\TextAt{9}{395.164}{156.662}{\Character{40}}%
\TextAt{14}{399.716}{156.662}{g}%
\TextAt{13}{405.738}{156.662}{\Character{59}}%
\TextAt{13}{410.962}{156.662}{A}%
\TextAt{9}{419.757}{156.662}{\Character{41}}%
\TextAt{13}{424.31}{156.662}{\Character{58}}%
\TextAt{12}{79.226}{112.717}{Exercise}%
\TextAt{12}{133.607}{112.717}{IV.3.5.}%
\TextAt{8}{178.389}{112.717}{Show}%
\TextAt{8}{210.307}{112.717}{that}%
\TextAt{13}{236.702}{112.717}{H}%
\TextAt{17}{247.324}{117.056}{1}%
\TextAt{9}{252.057}{112.717}{\Character{40}}%
\TextAt{13}{256.609}{112.717}{G\Character{59}}%
\TextAt{13}{271.063}{112.717}{A}%
\TextAt{9}{279.862}{112.717}{\Character{41}}%
\TextAt{9}{289.474}{112.717}{\Character{61}}%
\NativeMathAt{303.658}{112.717}{11.955}{\varinjlim}%
\TextAt{21}{319.917}{107.278}{N}%
\TextAt{13}{329.978}{112.717}{H}%
\TextAt{17}{340.6}{117.056}{1}%
\TextAt{9}{345.333}{112.717}{\Character{40}}%
\TextAt{13}{349.885}{112.717}{G}%
\TextAt{9}{359.119}{112.717}{\Character{47}}%
\TextAt{13}{364.972}{112.717}{N}%
\TextAt{13}{374.943}{112.717}{\Character{59}}%
\TextAt{13}{380.167}{112.717}{A}%
\TextAt{21}{388.963}{117.056}{N}%
\TextAt{9}{397.032}{112.717}{\Character{41}}%
\TextAt{8}{401.584}{112.717}{,}%
\TextAt{8}{410.016}{112.717}{where}%
\TextAt{13}{444.851}{112.717}{N}%
\TextAt{8}{460.397}{112.717}{varies}%
\TextAt{8}{494.336}{112.717}{over}%
\TextAt{8}{520.408}{112.717}{all}%
\TextAt{8}{79.226}{95.671}{the}%
\TextAt{8}{99.383}{95.671}{open}%
\TextAt{8}{127.668}{95.671}{normal}%
\TextAt{8}{167.348}{95.671}{subgroups}%
\TextAt{8}{222.76}{95.671}{of}%
\TextAt{13}{236.09}{95.671}{G}%
\TextAt{8}{245.324}{95.671}{.}%
\TextAt{8}{297.211}{43.033}{259}%
\Destination{78.23}{809.815}{nav.page.259}%
\Destination{79.226}{568.055}{nav.ex.4.3.1}%
\Destination{79.226}{498.955}{nav.ex.4.3.2}%
\Destination{79.226}{415.855}{nav.ex.4.3.3}%
\Destination{79.226}{308.755}{nav.ex.4.3.4}%
\Destination{79.226}{126.755}{nav.ex.4.3.5}%
\Destination{79.36}{309.755}{nav.print.284}%
%
\Hotlink{371.332}{690.255}{19.22}{11.047}{nav.cite.N015}%
\Hotlink{437.341}{577.855}{13.995}{11.047}{nav.chapter.6}%
\Hotlink{490.218}{465.755}{6.858}{11.047}{nav.ex.2.4.2}%
\Hotlink{412.975}{465.755}{9.763}{11.047}{nav.chapter.2}%
\Hotlink{434.332}{465.755}{6.858}{11.047}{nav.section.2.4}%
\Hotlink{457.486}{577.855}{25.078}{11.047}{nav.number.6.2.5}%
\Hotlink{80.397}{431.6}{21.722}{11.047}{nav.number.4.3.1}%
\Hotlink{502.911}{465.7}{4.866}{11.047}{nav.ex.2.4.6}%
}

% Page 272 (body page 260)
\Page{%
\TextAt{2}{79.226}{796.943}{IV.}%
\TextAt{2}{98.012}{796.943}{Abstract}%
\TextAt{2}{143.252}{796.943}{Class}%
\TextAt{2}{171.844}{796.943}{Field}%
\TextAt{2}{199.575}{796.943}{Theory}%
\Vector{q 0 G 0 g 0.299 w 7.226 18.729 m 460.769 18.729 l S Q}%
\TextAt{12}{79.226}{761.124}{Exercise}%
\TextAt{12}{131.254}{761.124}{IV.3.6.}%
\TextAt{8}{173.991}{761.124}{If}%
\TextAt{9}{184.664}{761.124}{1}%
\NativeMathAt{193.838}{761.124}{11.955}{\rightarrow}%
\TextAt{13}{209.114}{761.124}{A}%
\NativeMathAt{221.21}{761.124}{11.955}{\rightarrow}%
\TextAt{13}{236.486}{761.124}{B}%
\NativeMathAt{249.303}{761.124}{11.955}{\rightarrow}%
\TextAt{13}{264.579}{761.124}{C}%
\NativeMathAt{277.133}{761.124}{11.955}{\rightarrow}%
\TextAt{9}{292.409}{761.124}{1}%
\TextAt{8}{301.14}{761.124}{is}%
\TextAt{8}{311.885}{761.124}{an}%
\TextAt{8}{327.125}{761.124}{exact}%
\TextAt{8}{356.998}{761.124}{sequence}%
\TextAt{8}{404.481}{761.124}{of}%
\TextAt{8}{416.791}{761.124}{continuous}%
\TextAt{13}{474.376}{761.124}{G}%
\TextAt{8}{483.61}{761.124}{-modules,}%
\TextAt{8}{79.226}{744.078}{then}%
\TextAt{8}{105.886}{744.078}{one}%
\TextAt{8}{127.346}{744.078}{has}%
\TextAt{8}{148.219}{744.078}{an}%
\TextAt{8}{164.478}{744.078}{exact}%
\TextAt{8}{195.371}{744.078}{sequence}%
\TextAt{9}{127.181}{717.082}{1}%
\NativeArrowAt{136.354}{717.082}{11.955}{19.26}{\longrightarrow}%
\TextAt{13}{158.936}{717.082}{A}%
\TextAt{21}{167.712}{722.018}{G}%
\NativeArrowAt{178.129}{717.082}{11.955}{19.26}{\longrightarrow}%
\TextAt{13}{200.711}{717.082}{B}%
\TextAt{21}{210.208}{722.018}{G}%
\NativeArrowAt{220.625}{717.082}{11.955}{19.26}{\longrightarrow}%
\TextAt{13}{243.207}{717.082}{C}%
\TextAt{21}{252.441}{722.018}{G}%
\NativeArrowAt{262.858}{717.082}{11.955}{19.26}{\longrightarrow}%
\TextAt{13}{285.44}{717.082}{H}%
\TextAt{17}{296.063}{722.018}{1}%
\TextAt{9}{300.795}{717.082}{\Character{40}}%
\TextAt{13}{305.348}{717.082}{G\Character{59}}%
\TextAt{13}{319.802}{717.082}{A}%
\TextAt{9}{328.601}{717.082}{\Character{41}}%
\NativeArrowAt{336.474}{717.082}{11.955}{19.26}{\longrightarrow}%
\TextAt{13}{359.056}{717.082}{H}%
\TextAt{17}{369.679}{722.018}{1}%
\TextAt{9}{374.411}{717.082}{\Character{40}}%
\TextAt{13}{378.963}{717.082}{G\Character{59}}%
\TextAt{13}{393.417}{717.082}{B}%
\TextAt{9}{402.937}{717.082}{\Character{41}}%
\NativeArrowAt{410.81}{717.082}{11.955}{19.26}{\longrightarrow}%
\TextAt{13}{433.392}{717.082}{H}%
\TextAt{17}{444.015}{722.018}{1}%
\TextAt{9}{448.747}{717.082}{\Character{40}}%
\TextAt{13}{453.3}{717.082}{G\Character{59}}%
\TextAt{13}{467.754}{717.082}{C}%
\TextAt{9}{477.011}{717.082}{\Character{41}}%
\TextAt{13}{481.563}{717.082}{\Character{58}}%
\TextAt{12}{97.159}{687.097}{Remark:}%
\TextAt{8}{156.189}{687.097}{The}%
\TextAt{8}{182.03}{687.097}{group}%
\TextAt{13}{216.993}{687.097}{H}%
\TextAt{17}{227.616}{691.435}{1}%
\TextAt{9}{232.348}{687.097}{\Character{40}}%
\TextAt{13}{236.9}{687.097}{G\Character{59}}%
\TextAt{13}{251.354}{687.097}{A}%
\TextAt{9}{260.153}{687.097}{\Character{41}}%
\TextAt{8}{270.39}{687.097}{is}%
\TextAt{8}{283.941}{687.097}{only}%
\TextAt{8}{311.42}{687.097}{the}%
\TextAt{8}{333.364}{687.097}{first}%
\TextAt{8}{359.277}{687.097}{term}%
\TextAt{8}{389.027}{687.097}{of}%
\TextAt{8}{404.144}{687.097}{a}%
\TextAt{8}{415.687}{687.097}{whole}%
\TextAt{8}{450.638}{687.097}{series}%
\TextAt{8}{483.759}{687.097}{of}%
\TextAt{8}{498.877}{687.097}{groups}%
\TextAt{13}{79.226}{670.051}{H}%
\TextAt{21}{89.849}{674.389}{i}%
\TextAt{9}{93.23}{670.051}{\Character{40}}%
\TextAt{13}{97.782}{670.051}{G\Character{59}}%
\TextAt{13}{112.236}{670.051}{A}%
\TextAt{9}{121.035}{670.051}{\Character{41}}%
\TextAt{8}{125.588}{670.051}{,}%
\TextAt{13}{132.781}{670.051}{i}%
\TextAt{9}{140.155}{670.051}{\Character{61}}%
\TextAt{9}{152.624}{670.051}{1}%
\TextAt{13}{158.494}{670.051}{\Character{59}}%
\TextAt{9}{163.738}{670.051}{2}%
\TextAt{13}{169.591}{670.051}{\Character{59}}%
\TextAt{9}{174.835}{670.051}{3}%
\TextAt{13}{180.688}{670.051}{\Character{59}}%
\TextAt{13}{185.912}{670.051}{\Character{58}}%
\TextAt{13}{191.149}{670.051}{\Character{58}}%
\TextAt{13}{196.373}{670.051}{\Character{58}}%
\TextAt{8}{199.672}{670.051}{,}%
\TextAt{8}{206.865}{670.051}{which}%
\TextAt{8}{240.387}{670.051}{are}%
\TextAt{8}{259.933}{670.051}{the}%
\TextAt{8}{280.125}{670.051}{objects}%
\TextAt{8}{319.887}{670.051}{of}%
\TextAt{12}{333.766}{670.051}{group}%
\TextAt{12}{372.195}{670.051}{cohomology}%
\TextAt{8}{445.288}{670.051}{(see}%
\TextAt{8}{468.791}{670.051}{[145]).}%
\TextAt{8}{505.978}{670.051}{Class}%
\TextAt{8}{79.226}{653.005}{field}%
\TextAt{8}{104.583}{653.005}{theory}%
\TextAt{8}{141.333}{653.005}{can}%
\TextAt{8}{162.793}{653.005}{also}%
\TextAt{8}{186.273}{653.005}{be}%
\TextAt{8}{202.197}{653.005}{built}%
\TextAt{8}{230.16}{653.005}{upon}%
\TextAt{8}{259.749}{653.005}{this}%
\TextAt{8}{282.572}{653.005}{theory}%
\TextAt{8}{319.322}{653.005}{(see}%
\TextAt{8}{342.79}{653.005}{[10],}%
\TextAt{8}{368.159}{653.005}{[108]).}%
\TextAt{12}{79.226}{622.161}{Exercise}%
\TextAt{12}{133.148}{622.161}{IV.3.7.}%
\TextAt{8}{177.531}{622.161}{Even}%
\TextAt{8}{207.579}{622.161}{for}%
\TextAt{8}{226.091}{622.161}{infinite}%
\TextAt{8}{266.386}{622.161}{Galois}%
\TextAt{8}{302.926}{622.161}{extensions}%
\TextAt{13}{359.933}{622.161}{L}%
\TextAt{16}{372.287}{622.161}{\Character{106}}%
\TextAt{13}{379.997}{622.161}{K}%
\TextAt{8}{390.815}{622.161}{,}%
\TextAt{8}{398.748}{622.161}{one}%
\TextAt{8}{420.835}{622.161}{has}%
\TextAt{8}{442.335}{622.161}{Hilbert’s}%
\TextAt{8}{491.142}{622.161}{theorem}%
\TextAt{8}{79.226}{605.115}{90:}%
\TextAt{13}{248.831}{588.069}{H}%
\TextAt{17}{259.453}{593.005}{1}%
\TextAt{9}{264.185}{588.069}{\Character{40}}%
\TextAt{13}{268.738}{588.069}{G}%
\TextAt{9}{277.971}{588.069}{\Character{40}}%
\TextAt{13}{282.524}{588.069}{L}%
\TextAt{16}{293.809}{588.069}{\Character{106}}%
\TextAt{13}{300.451}{588.069}{K}%
\TextAt{9}{311.268}{588.069}{\Character{41}}%
\TextAt{13}{315.821}{588.069}{\Character{59}}%
\TextAt{13}{321.045}{588.069}{L}%
\TextAt{3}{329.029}{593.005}{\Character{3}}%
\TextAt{9}{333.762}{588.069}{\Character{41}}%
\TextAt{9}{341.617}{588.069}{\Character{61}}%
\TextAt{9}{354.038}{588.069}{1}%
\TextAt{13}{359.913}{588.069}{\Character{58}}%
\TextAt{12}{79.226}{551.608}{Exercise}%
\TextAt{12}{131.428}{551.608}{IV.3.8.}%
\TextAt{8}{174.317}{551.608}{If}%
\TextAt{13}{185.141}{551.608}{n}%
\TextAt{8}{195.158}{551.608}{is}%
\TextAt{8}{206.054}{551.608}{not}%
\TextAt{8}{226}{551.608}{divisible}%
\TextAt{8}{271.04}{551.608}{by}%
\TextAt{8}{286.432}{551.608}{the}%
\TextAt{8}{305.72}{551.608}{characteristic}%
\TextAt{8}{376.787}{551.608}{of}%
\TextAt{8}{389.249}{551.608}{the}%
\TextAt{8}{408.537}{551.608}{field}%
\TextAt{13}{433.027}{551.608}{K}%
\TextAt{8}{446.874}{551.608}{and}%
\TextAt{8}{468.769}{551.608}{if}%
\TextAt{13}{478.625}{551.608}{\Character{22}}%
\TextAt{21}{485.668}{549.815}{n}%
\TextAt{8}{494.334}{551.608}{denotes}%
\TextAt{8}{79.226}{534.562}{the}%
\TextAt{8}{99.383}{534.562}{group}%
\TextAt{8}{132.558}{534.562}{of}%
\TextAt{13}{145.888}{534.562}{n}%
\TextAt{8}{152.876}{534.562}{-th}%
\TextAt{8}{171.729}{534.562}{roots}%
\TextAt{8}{201.39}{534.562}{of}%
\TextAt{8}{214.72}{534.562}{unity}%
\TextAt{8}{245.289}{534.562}{in}%
\TextAt{8}{258.942}{534.562}{the}%
\TextAt{8}{279.098}{534.562}{separable}%
\TextAt{8}{330.542}{534.562}{closure}%
\TextAt{9}{372.756}{537.584}{\Character{22}}%
\TextAt{13}{369.623}{534.562}{K}%
\TextAt{8}{380.441}{534.562}{,}%
\TextAt{8}{387.59}{534.562}{then}%
\TextAt{13}{245.447}{507.566}{H}%
\TextAt{17}{256.069}{512.502}{1}%
\TextAt{9}{260.802}{507.566}{\Character{40}}%
\TextAt{13}{265.354}{507.566}{G}%
\TextAt{21}{274.588}{505.773}{K}%
\TextAt{13}{282.79}{507.566}{\Character{59}}%
\TextAt{13}{288.014}{507.566}{\Character{22}}%
\TextAt{21}{295.077}{505.773}{n}%
\TextAt{9}{300.714}{507.566}{\Character{41}}%
\NativeMathAt{308.587}{507.566}{11.955}{\cong}%
\TextAt{13}{321.206}{507.566}{K}%
\TextAt{3}{332.024}{512.502}{\Character{3}}%
\TextAt{9}{336.756}{507.566}{\Character{47}}%
\TextAt{13}{342.609}{507.566}{K}%
\TextAt{3}{353.427}{512.502}{\Character{3}}%
\TextAt{21}{357.661}{512.502}{n}%
\TextAt{13}{363.297}{507.566}{\Character{58}}%
\TextAt{15}{79.226}{445.622}{§}%
\TextAt{15}{93.785}{445.622}{4}%
\TextAt{15}{120.693}{445.622}{Abstract}%
\TextAt{15}{201.21}{445.622}{Valuation}%
\TextAt{15}{289.043}{445.622}{Theory}%
\TextAt{8}{79.226}{413.736}{The}%
\TextAt{8}{103.844}{413.736}{further}%
\TextAt{8}{143.752}{413.736}{development}%
\TextAt{8}{212.28}{413.736}{will}%
\TextAt{8}{234.949}{413.736}{now}%
\TextAt{8}{259.901}{413.736}{be}%
\TextAt{8}{276.389}{413.736}{based}%
\TextAt{8}{309.531}{413.736}{on}%
\TextAt{8}{326.353}{413.736}{a}%
\TextAt{8}{336.672}{413.736}{fixed}%
\TextAt{8}{365.522}{413.736}{choice}%
\TextAt{8}{400.875}{413.736}{of}%
\TextAt{8}{414.769}{413.736}{a}%
\TextAt{8}{425.088}{413.736}{surjective}%
\TextAt{8}{478.063}{413.736}{continuous}%
\TextAt{8}{79.226}{396.69}{homomorphism}%
\TextAt{13}{276.457}{379.644}{d}%
\TextAt{9}{285.86}{379.644}{\Character{58}}%
\TextAt{13}{292.433}{379.644}{G}%
\NativeArrowAt{304.987}{379.644}{11.955}{19.26}{\longrightarrow}%
\TextAt{20}{328.233}{382.732}{\Character{98}}%
\TextAt{18}{327.569}{379.644}{Z}%
\TextAt{8}{79.226}{356.982}{from}%
\TextAt{8}{108.088}{356.982}{the}%
\TextAt{8}{129.465}{356.982}{profinite}%
\TextAt{8}{177.515}{356.982}{group}%
\TextAt{13}{211.911}{356.982}{G}%
\TextAt{8}{226.263}{356.982}{onto}%
\TextAt{8}{253.834}{356.982}{the}%
\TextAt{8}{275.211}{356.982}{procyclic}%
\TextAt{8}{325.855}{356.982}{group}%
\TextAt{20}{360.916}{360.07}{\Character{98}}%
\TextAt{18}{360.252}{356.982}{Z}%
\TextAt{9}{373.624}{356.982}{\Character{61}}%
\NativeMathAt{388.13}{356.982}{11.955}{\varprojlim}%
\TextAt{21}{404.389}{351.542}{n}%
\TextAt{18}{412.018}{356.982}{Z}%
\TextAt{9}{419.988}{356.982}{\Character{47}}%
\TextAt{13}{425.841}{356.982}{n}%
\TextAt{18}{432.829}{356.982}{Z}%
\TextAt{8}{445.917}{356.982}{(see}%
\TextAt{8}{470.606}{356.982}{§2,}%
\TextAt{8}{490.819}{356.982}{example}%
\TextAt{8}{79.226}{339.936}{4).}%
\TextAt{8}{103.03}{339.936}{This}%
\TextAt{8}{131.398}{339.936}{homomorphism}%
\TextAt{8}{215.716}{339.936}{will}%
\TextAt{8}{239.468}{339.936}{produce}%
\TextAt{8}{285.661}{339.936}{a}%
\TextAt{8}{297.064}{339.936}{theory}%
\TextAt{8}{335.462}{339.936}{which}%
\TextAt{8}{370.596}{339.936}{is}%
\TextAt{8}{384.008}{339.936}{an}%
\TextAt{8}{401.915}{339.936}{abstract}%
\TextAt{8}{449.159}{339.936}{reflection}%
\TextAt{8}{501.533}{339.936}{of}%
\TextAt{8}{516.51}{339.936}{the}%
\TextAt{8}{79.226}{322.89}{ramification}%
\TextAt{8}{144.628}{322.89}{theory}%
\TextAt{8}{181.731}{322.89}{of}%
\TextAt{13}{195.415}{322.89}{p}%
\TextAt{8}{201.29}{322.89}{-adic}%
\TextAt{8}{230.252}{322.89}{number}%
\TextAt{8}{273.202}{322.89}{fields.}%
\TextAt{8}{308.785}{322.89}{Indeed,}%
\TextAt{8}{350.509}{322.89}{in}%
\TextAt{8}{364.515}{322.89}{the}%
\TextAt{8}{385.025}{322.89}{case}%
\TextAt{8}{410.15}{322.89}{where}%
\TextAt{13}{444.312}{322.89}{G}%
\TextAt{8}{457.797}{322.89}{is}%
\TextAt{8}{469.914}{322.89}{the}%
\TextAt{8}{490.424}{322.89}{absolute}%
\TextAt{8}{79.226}{305.844}{Galois}%
\TextAt{8}{115.28}{305.844}{group}%
\TextAt{13}{148.596}{305.844}{G}%
\TextAt{21}{157.83}{304.05}{k}%
\TextAt{9}{166.51}{305.844}{\Character{61}}%
\TextAt{13}{179.175}{305.844}{G}%
\TextAt{9}{188.409}{305.844}{\Character{40}}%
\TextAt{9}{193.279}{308.998}{\Character{22}}%
\TextAt{13}{192.961}{305.844}{k}%
\TextAt{16}{203.011}{305.844}{\Character{106}}%
\TextAt{13}{209.892}{305.844}{k}%
\TextAt{9}{216.381}{305.844}{\Character{41}}%
\TextAt{8}{224.972}{305.844}{of}%
\TextAt{8}{238.442}{305.844}{a}%
\TextAt{13}{248.338}{305.844}{p}%
\TextAt{8}{254.213}{305.844}{-adic}%
\TextAt{8}{282.963}{305.844}{number}%
\TextAt{8}{325.699}{305.844}{field}%
\TextAt{13}{351.197}{305.844}{k}%
\TextAt{8}{357.686}{305.844}{,}%
\TextAt{8}{365.011}{305.844}{such}%
\TextAt{8}{391.549}{305.844}{a}%
\TextAt{8}{401.445}{305.844}{surjective}%
\TextAt{8}{453.997}{305.844}{homomorphism}%
\TextAt{13}{79.226}{288.798}{d}%
\TextAt{9}{88.63}{288.798}{\Character{58}}%
\TextAt{13}{95.202}{288.798}{G}%
\NativeMathAt{107.757}{288.798}{11.955}{\rightarrow}%
\TextAt{20}{123.697}{291.886}{\Character{98}}%
\TextAt{18}{123.033}{288.798}{Z}%
\TextAt{8}{134.159}{288.798}{arises}%
\TextAt{8}{165.411}{288.798}{via}%
\TextAt{8}{183.858}{288.798}{the}%
\TextAt{8}{203.274}{288.798}{maximal}%
\TextAt{8}{250.342}{288.798}{unramified}%
\TextAt{8}{308.133}{288.798}{extension}%
\TextAt{20}{359.082}{291.953}{\Character{101}}%
\TextAt{13}{359.158}{288.798}{k}%
\TextAt{16}{368.969}{288.798}{\Character{106}}%
\TextAt{13}{375.611}{288.798}{k}%
\TextAt{8}{382.1}{288.798}{:}%
\TextAt{8}{390.178}{288.798}{if}%
\TextAt{18}{400.161}{288.798}{F}%
\TextAt{21}{407.467}{287.004}{q}%
\TextAt{8}{415.189}{288.798}{is}%
\TextAt{8}{426.212}{288.798}{the}%
\TextAt{8}{445.627}{288.798}{residue}%
\TextAt{8}{484.614}{288.798}{class}%
\TextAt{8}{511.31}{288.798}{field}%
\TextAt{8}{79.226}{271.752}{of}%
\TextAt{13}{92.556}{271.752}{k}%
\TextAt{8}{99.046}{271.752}{,}%
\TextAt{8}{106.195}{271.752}{then,}%
\TextAt{8}{136.107}{271.752}{by}%
\TextAt{8}{152.366}{271.752}{chap.}%
\TextAt{8}{184.557}{271.752}{II,}%
\TextAt{8}{200.469}{271.752}{§9,}%
\TextAt{8}{219.155}{271.752}{p.173}%
\TextAt{8}{250.382}{271.752}{and}%
\TextAt{8}{273.145}{271.752}{example}%
\TextAt{8}{318.993}{271.752}{5}%
\TextAt{8}{328.748}{271.752}{in}%
\TextAt{8}{342.401}{271.752}{§2,}%
\TextAt{8}{361.087}{271.752}{we}%
\TextAt{8}{378.314}{271.752}{have}%
\TextAt{8}{405.309}{271.752}{canonical}%
\TextAt{8}{456.693}{271.752}{isomorphisms}%
\TextAt{13}{239.411}{244.756}{G}%
\TextAt{9}{248.645}{244.756}{\Character{40}}%
\TextAt{20}{253.121}{247.911}{\Character{101}}%
\TextAt{13}{253.197}{244.756}{k}%
\TextAt{16}{263.008}{244.756}{\Character{106}}%
\TextAt{13}{269.649}{244.756}{k}%
\TextAt{9}{276.139}{244.756}{\Character{41}}%
\NativeMathAt{284.012}{244.756}{11.955}{\cong}%
\TextAt{13}{296.631}{244.756}{G}%
\TextAt{9}{305.865}{244.756}{\Character{40}}%
\Vector{q 0 G 0 g 0.478 w 238.417 -519.335 m 245.723 -519.335 l S Q}%
\TextAt{18}{310.417}{244.756}{F}%
\TextAt{21}{317.723}{242.962}{q}%
\TextAt{16}{325.609}{244.756}{\Character{106}}%
\TextAt{18}{332.251}{244.756}{F}%
\TextAt{21}{339.557}{242.962}{q}%
\TextAt{9}{344.122}{244.756}{\Character{41}}%
\NativeMathAt{351.995}{244.756}{11.955}{\cong}%
\TextAt{20}{365.278}{247.844}{\Character{98}}%
\TextAt{18}{364.614}{244.756}{Z}%
\TextAt{8}{79.226}{217.76}{which}%
\TextAt{8}{113.433}{217.76}{associate}%
\TextAt{8}{163.386}{217.76}{to}%
\TextAt{8}{178.417}{217.76}{the}%
\TextAt{8}{199.294}{217.76}{element}%
\TextAt{9}{243.257}{217.76}{1}%
\NativeMathAt{253.659}{217.76}{11.955}{\in}%
\TextAt{20}{266.843}{220.848}{\Character{98}}%
\TextAt{18}{266.178}{217.76}{Z}%
\TextAt{8}{278.767}{217.76}{the}%
\TextAt{8}{299.644}{217.76}{Frobenius}%
\TextAt{8}{354.235}{217.76}{automorphism}%
\TextAt{13}{432.425}{217.76}{\Character{39}}%
\NativeMathAt{444.643}{217.76}{11.955}{\in}%
\TextAt{13}{457.162}{217.76}{G}%
\TextAt{9}{466.396}{217.76}{\Character{40}}%
\TextAt{20}{470.872}{220.915}{\Character{101}}%
\TextAt{13}{470.948}{217.76}{k}%
\TextAt{16}{481.987}{217.76}{\Character{106}}%
\TextAt{13}{489.857}{217.76}{k}%
\TextAt{9}{496.347}{217.76}{\Character{41}}%
\TextAt{8}{500.899}{217.76}{.}%
\TextAt{8}{511.51}{217.76}{It}%
\TextAt{8}{524.903}{217.76}{is}%
\TextAt{8}{79.226}{200.714}{defined}%
\TextAt{8}{119.539}{200.714}{by}%
\TextAt{13}{224.196}{183.668}{a}%
\TextAt{21}{230.341}{188.604}{\Character{39}}%
\NativeMathAt{239.678}{183.668}{11.955}{\equiv}%
\TextAt{13}{252.297}{183.668}{a}%
\TextAt{21}{258.442}{188.604}{q}%
\TextAt{9}{274.962}{183.668}{\Character{40}}%
\TextAt{8}{279.514}{183.668}{mod}%
\TextAt{20}{305.287}{184.784}{\Character{101}}%
\TextAt{14}{305.617}{183.668}{p}%
\TextAt{9}{311.6}{183.668}{\Character{41}}%
\TextAt{8}{340.062}{183.668}{for}%
\TextAt{13}{357.947}{183.668}{a}%
\NativeMathAt{367.413}{183.668}{11.955}{\in}%
\TextAt{20}{378.305}{184.203}{\Character{101}}%
\TextAt{14}{378.704}{183.668}{o}%
\TextAt{13}{384.548}{183.668}{\Character{59}}%
\TextAt{8}{79.226}{161.005}{where}%
\TextAt{20}{112.74}{161.541}{\Character{101}}%
\TextAt{14}{113.139}{161.005}{o}%
\TextAt{8}{118.983}{161.005}{,}%
\TextAt{8}{126.261}{161.005}{resp.}%
\TextAt{20}{155.564}{162.122}{\Character{101}}%
\TextAt{14}{155.893}{161.005}{p}%
\TextAt{8}{161.876}{161.005}{,}%
\TextAt{8}{169.154}{161.005}{denote}%
\TextAt{8}{206.976}{161.005}{the}%
\TextAt{8}{227.236}{161.005}{valuation}%
\TextAt{8}{278.412}{161.005}{ring}%
\TextAt{8}{302.581}{161.005}{of}%
\TextAt{20}{315.938}{164.16}{\Character{101}}%
\TextAt{13}{316.014}{161.005}{k}%
\TextAt{8}{322.503}{161.005}{,}%
\TextAt{8}{329.782}{161.005}{resp.}%
\TextAt{8}{359.414}{161.005}{its}%
\TextAt{8}{375.836}{161.005}{maximal}%
\TextAt{8}{423.748}{161.005}{ideal.}%
\TextAt{8}{456.572}{161.005}{The}%
\TextAt{8}{480.729}{161.005}{homomor-}%
\TextAt{8}{79.226}{143.959}{phism}%
\TextAt{13}{113.753}{143.959}{d}%
\TextAt{9}{123.156}{143.959}{\Character{58}}%
\TextAt{13}{129.729}{143.959}{G}%
\NativeMathAt{142.283}{143.959}{11.955}{\rightarrow}%
\TextAt{20}{158.223}{147.048}{\Character{98}}%
\TextAt{18}{157.559}{143.959}{Z}%
\TextAt{8}{169.427}{143.959}{in}%
\TextAt{8}{183.079}{143.959}{question}%
\TextAt{8}{229.645}{143.959}{is}%
\TextAt{8}{241.409}{143.959}{then}%
\TextAt{8}{268.069}{143.959}{given,}%
\TextAt{8}{301.89}{143.959}{in}%
\TextAt{8}{315.543}{143.959}{this}%
\TextAt{8}{338.365}{143.959}{concrete}%
\TextAt{8}{384.536}{143.959}{case,}%
\TextAt{8}{412.559}{143.959}{as}%
\TextAt{8}{426.929}{143.959}{the}%
\TextAt{8}{447.086}{143.959}{composite}%
\TextAt{13}{252.271}{116.963}{G}%
\NativeArrowAt{264.826}{116.963}{11.955}{19.26}{\longrightarrow}%
\TextAt{13}{287.408}{116.963}{G}%
\TextAt{9}{296.642}{116.963}{\Character{40}}%
\TextAt{20}{301.118}{120.118}{\Character{101}}%
\TextAt{13}{301.194}{116.963}{k}%
\TextAt{16}{311.004}{116.963}{\Character{106}}%
\TextAt{13}{317.646}{116.963}{k}%
\TextAt{9}{324.135}{116.963}{\Character{41}}%
\NativeMathAt{334.596}{123.74}{7.97}{\sim}%
\NativeArrowAt{332.009}{116.963}{11.955}{13.175}{\longrightarrow}%
\TextAt{20}{349.167}{120.052}{\Character{98}}%
\TextAt{18}{348.502}{116.963}{Z}%
\TextAt{13}{356.473}{116.963}{\Character{58}}%
\TextAt{8}{79.226}{89.967}{In}%
\TextAt{8}{95.085}{89.967}{the}%
\TextAt{8}{116.479}{89.967}{abstract}%
\TextAt{8}{163.313}{89.967}{situation,}%
\TextAt{8}{216.961}{89.967}{the}%
\TextAt{8}{238.355}{89.967}{initial}%
\TextAt{8}{273.413}{89.967}{choice}%
\TextAt{8}{309.44}{89.967}{of}%
\TextAt{8}{324.008}{89.967}{a}%
\TextAt{8}{335.001}{89.967}{surjective}%
\TextAt{8}{388.65}{89.967}{homomorphism}%
\TextAt{13}{472.557}{89.967}{d}%
\TextAt{9}{484.07}{89.967}{\Character{58}}%
\TextAt{13}{492.751}{89.967}{G}%
\NativeMathAt{507.414}{89.967}{11.955}{\rightarrow}%
\TextAt{20}{525.464}{93.056}{\Character{98}}%
\TextAt{18}{524.799}{89.967}{Z}%
\TextAt{8}{79.226}{72.921}{mimics}%
\TextAt{8}{119.593}{72.921}{the}%
\TextAt{13}{140.388}{72.921}{p}%
\TextAt{8}{146.263}{72.921}{-adic}%
\TextAt{8}{175.512}{72.921}{case,}%
\TextAt{8}{204.334}{72.921}{but}%
\TextAt{8}{226.433}{72.921}{the}%
\TextAt{8}{247.228}{72.921}{applications}%
\TextAt{8}{312.976}{72.921}{of}%
\TextAt{8}{326.945}{72.921}{the}%
\TextAt{8}{347.741}{72.921}{theory}%
\TextAt{8}{385.13}{72.921}{are}%
\TextAt{8}{405.28}{72.921}{by}%
\TextAt{8}{422.179}{72.921}{no}%
\TextAt{8}{439.077}{72.921}{means}%
\TextAt{8}{475.546}{72.921}{confined}%
\TextAt{8}{522.357}{72.921}{to}%
\TextAt{8}{297.211}{43.033}{260}%
\Destination{78.23}{809.815}{nav.page.260}%
\Destination{79.23}{471.485}{nav.section.4.4}%
\Destination{79.226}{775.155}{nav.ex.4.3.6}%
\Destination{79.226}{636.155}{nav.ex.4.3.7}%
\Destination{79.226}{565.655}{nav.ex.4.3.8}%
\Destination{79.36}{232.755}{nav.print.285}%
\Destination{79.23}{471.485}{outline.44}%
\Hotlink{468.291}{667.855}{25.078}{11.047}{nav.cite.N145}%
\Hotlink{342.29}{650.855}{19.22}{11.047}{nav.cite.N010}%
\Hotlink{367.659}{650.855}{25.078}{11.047}{nav.cite.N108}%
\Hotlink{475.785}{354.755}{6.858}{11.047}{nav.section.4.2}%
\Hotlink{184.057}{269.555}{9.763}{11.047}{nav.chapter.2}%
\Hotlink{205.648}{269.555}{6.858}{11.047}{nav.section.2.9}%
\Hotlink{347.58}{269.555}{6.858}{11.047}{nav.section.4.2}%
\Hotlink{228.41}{269.555}{18.574}{11.047}{nav.print.173}%
\Hotlink{79.548}{337.7}{5.2}{11.047}{nav.example.4.2.4}%
\Hotlink{319.567}{269.5}{4.674}{11.047}{nav.example.4.2.5}%
}

% Page 273 (body page 261)
\Page{%
\TextAt{2}{79.226}{796.943}{IV.}%
\TextAt{2}{98.012}{796.943}{Abstract}%
\TextAt{2}{143.252}{796.943}{Class}%
\TextAt{2}{171.844}{796.943}{Field}%
\TextAt{2}{199.575}{796.943}{Theory}%
\Vector{q 0 G 0 g 0.299 w 7.226 18.729 m 460.769 18.729 l S Q}%
\TextAt{13}{79.226}{761.124}{p}%
\TextAt{8}{85.101}{761.124}{-adic}%
\TextAt{8}{114.149}{761.124}{number}%
\TextAt{8}{157.184}{761.124}{fields.}%
\TextAt{8}{193.023}{761.124}{The}%
\TextAt{8}{217.516}{761.124}{kernel}%
\TextAt{13}{252.422}{761.124}{I}%
\TextAt{8}{262.861}{761.124}{of}%
\TextAt{13}{276.63}{761.124}{d}%
\TextAt{8}{287.049}{761.124}{has}%
\TextAt{8}{308.362}{761.124}{a}%
\TextAt{8}{318.556}{761.124}{certain}%
\TextAt{8}{358.017}{761.124}{fixed}%
\TextAt{8}{386.741}{761.124}{field}%
\TextAt{20}{412.461}{764.279}{\Character{101}}%
\TextAt{13}{412.537}{761.124}{k}%
\TextAt{16}{423.096}{761.124}{\Character{106}}%
\TextAt{13}{430.485}{761.124}{k}%
\TextAt{8}{436.975}{761.124}{,}%
\TextAt{8}{444.673}{761.124}{and}%
\TextAt{13}{467.874}{761.124}{d}%
\TextAt{8}{478.293}{761.124}{induces}%
\TextAt{8}{520.408}{761.124}{an}%
\TextAt{8}{79.226}{744.078}{isomorphism}%
\TextAt{13}{147.646}{744.078}{G}%
\TextAt{9}{156.879}{744.078}{\Character{40}}%
\TextAt{20}{161.355}{747.233}{\Character{101}}%
\TextAt{13}{161.432}{744.078}{k}%
\TextAt{16}{171.242}{744.078}{\Character{106}}%
\TextAt{13}{177.884}{744.078}{k}%
\TextAt{9}{184.373}{744.078}{\Character{41}}%
\NativeMathAt{192.246}{744.078}{11.955}{\cong}%
\TextAt{20}{205.53}{747.166}{\Character{98}}%
\TextAt{18}{204.866}{744.078}{Z}%
\TextAt{8}{212.836}{744.078}{.}%
\TextAt{8}{97.159}{724.043}{More}%
\TextAt{8}{126.854}{724.043}{generally,}%
\TextAt{8}{178.451}{724.043}{for}%
\TextAt{8}{195.796}{724.043}{any}%
\TextAt{8}{217.373}{724.043}{field}%
\TextAt{13}{242.19}{724.043}{K}%
\TextAt{8}{256.365}{724.043}{we}%
\TextAt{8}{273.052}{724.043}{denote}%
\TextAt{8}{310.231}{724.043}{by}%
\TextAt{13}{325.95}{724.043}{I}%
\TextAt{21}{331.113}{722.25}{K}%
\TextAt{8}{342.673}{724.043}{the}%
\TextAt{8}{362.29}{724.043}{kernel}%
\TextAt{8}{396.216}{724.043}{of}%
\TextAt{8}{409.007}{724.043}{the}%
\TextAt{8}{428.623}{724.043}{restriction}%
\TextAt{13}{484.081}{724.043}{d}%
\TextAt{9}{493.485}{724.043}{\Character{58}}%
\TextAt{13}{500.057}{724.043}{G}%
\TextAt{21}{509.291}{722.25}{K}%
\NativeMathAt{520.814}{724.043}{11.955}{\rightarrow}%
\TextAt{20}{79.89}{710.085}{\Character{98}}%
\TextAt{18}{79.226}{706.997}{Z}%
\TextAt{8}{87.196}{706.997}{,}%
\TextAt{8}{94.345}{706.997}{and}%
\TextAt{8}{117.108}{706.997}{call}%
\TextAt{8}{138.568}{706.997}{it}%
\TextAt{8}{150.272}{706.997}{the}%
\TextAt{12}{170.428}{706.997}{inertia}%
\TextAt{12}{213.276}{706.997}{group}%
\TextAt{8}{251.078}{706.997}{over}%
\TextAt{13}{276.124}{706.997}{K}%
\TextAt{8}{286.942}{706.997}{.}%
\TextAt{8}{295.39}{706.997}{Since}%
\TextAt{13}{222.093}{682.951}{I}%
\TextAt{21}{227.256}{681.158}{K}%
\TextAt{9}{238.779}{682.951}{\Character{61}}%
\TextAt{13}{251.205}{682.951}{G}%
\TextAt{21}{260.438}{681.158}{K}%
\TextAt{16}{271.298}{682.951}{\Character{92}}%
\TextAt{13}{281.925}{682.951}{I}%
\TextAt{9}{291.348}{682.951}{\Character{61}}%
\TextAt{13}{303.774}{682.951}{G}%
\TextAt{21}{313.007}{681.158}{K}%
\TextAt{16}{323.867}{682.951}{\Character{92}}%
\TextAt{13}{334.493}{682.951}{G}%
\TextAt{30}{343.686}{681.313}{\Character{101}}%
\TextAt{21}{343.727}{679.21}{k}%
\TextAt{9}{352.168}{682.951}{\Character{61}}%
\TextAt{13}{364.593}{682.951}{G}%
\TextAt{21}{373.827}{679.21}{K}%
\TextAt{30}{381.49}{681.313}{\Character{101}}%
\TextAt{21}{381.531}{679.21}{k}%
\TextAt{13}{386.651}{682.951}{\Character{59}}%
\TextAt{8}{79.226}{658.905}{the}%
\TextAt{8}{99.383}{658.905}{fixed}%
\TextAt{8}{127.668}{658.905}{field}%
\TextAt{20}{155.764}{661.927}{\Character{101}}%
\TextAt{13}{153.025}{658.905}{K}%
\TextAt{8}{167.74}{658.905}{of}%
\TextAt{13}{181.071}{658.905}{I}%
\TextAt{21}{186.233}{657.112}{K}%
\TextAt{8}{198.333}{658.905}{is}%
\TextAt{8}{210.097}{658.905}{the}%
\TextAt{8}{230.254}{658.905}{composite}%
\TextAt{20}{285.175}{637.881}{\Character{101}}%
\TextAt{13}{282.436}{634.859}{K}%
\TextAt{9}{296.575}{634.859}{\Character{61}}%
\TextAt{13}{309}{634.859}{K}%
\TextAt{20}{319.742}{638.014}{\Character{101}}%
\TextAt{13}{319.818}{634.859}{k}%
\TextAt{13}{326.298}{634.859}{\Character{58}}%
\TextAt{8}{79.226}{610.813}{We}%
\TextAt{8}{99.371}{610.813}{call}%
\TextAt{20}{123.568}{613.835}{\Character{101}}%
\TextAt{13}{120.83}{610.813}{K}%
\TextAt{16}{134.969}{610.813}{\Character{106}}%
\TextAt{13}{141.61}{610.813}{K}%
\TextAt{8}{156.325}{610.813}{the}%
\TextAt{12}{176.482}{610.813}{maximal}%
\TextAt{12}{231.057}{610.813}{unramified}%
\TextAt{12}{298.556}{610.813}{extension}%
\TextAt{8}{357.782}{610.813}{of}%
\TextAt{13}{371.112}{610.813}{K}%
\TextAt{8}{381.93}{610.813}{.}%
\TextAt{8}{390.378}{610.813}{We}%
\TextAt{8}{410.522}{610.813}{put}%
\TextAt{13}{212.845}{586.767}{f}%
\TextAt{21}{218.615}{584.974}{K}%
\TextAt{9}{230.138}{586.767}{\Character{61}}%
\TextAt{9}{242.548}{586.767}{\Character{40}}%
\TextAt{20}{247.78}{589.856}{\Character{98}}%
\TextAt{18}{247.116}{586.767}{Z}%
\TextAt{9}{258.407}{586.767}{\Character{58}}%
\TextAt{13}{264.98}{586.767}{d}%
\TextAt{9}{271.062}{586.767}{\Character{40}}%
\TextAt{13}{275.615}{586.767}{G}%
\TextAt{21}{284.848}{584.974}{K}%
\TextAt{9}{293.051}{586.767}{\Character{41}\Character{41}}%
\TextAt{13}{302.156}{586.767}{\Character{59}}%
\TextAt{13}{331.291}{586.767}{e}%
\TextAt{21}{336.735}{584.974}{K}%
\TextAt{9}{348.259}{586.767}{\Character{61}}%
\TextAt{9}{360.668}{586.767}{\Character{40}}%
\TextAt{13}{365.237}{586.767}{I}%
\TextAt{9}{374.66}{586.767}{\Character{58}}%
\TextAt{13}{381.233}{586.767}{I}%
\TextAt{21}{386.396}{584.974}{K}%
\TextAt{9}{394.598}{586.767}{\Character{41}}%
\TextAt{8}{79.226}{562.721}{and}%
\TextAt{8}{101.989}{562.721}{obtain,}%
\TextAt{8}{141.668}{562.721}{when}%
\TextAt{13}{172.225}{562.721}{f}%
\TextAt{21}{177.995}{560.928}{K}%
\TextAt{8}{190.095}{562.721}{is}%
\TextAt{8}{201.859}{562.721}{finite,}%
\TextAt{8}{235.023}{562.721}{a}%
\TextAt{8}{244.778}{562.721}{surjective}%
\TextAt{8}{297.19}{562.721}{homomorphism}%
\TextAt{13}{249.157}{536.494}{d}%
\TextAt{21}{255.24}{534.701}{K}%
\TextAt{9}{266.764}{536.494}{\Character{61}}%
\TextAt{9}{284.445}{544.582}{1}%
\Vector{q 0 G 0 g 0.478 w 208.385 -234.517 m 222.357 -234.517 l S Q}%
\TextAt{13}{280.385}{528.294}{f}%
\TextAt{21}{286.155}{526.5}{K}%
\TextAt{13}{295.553}{536.494}{d}%
\TextAt{9}{304.956}{536.494}{\Character{58}}%
\TextAt{13}{311.529}{536.494}{G}%
\TextAt{21}{320.762}{534.701}{K}%
\NativeArrowAt{332.286}{536.494}{11.955}{19.26}{\longrightarrow}%
\TextAt{20}{355.532}{539.583}{\Character{98}}%
\TextAt{18}{354.868}{536.494}{Z}%
\TextAt{8}{79.226}{509.676}{with}%
\TextAt{8}{105.886}{509.676}{kernel}%
\TextAt{13}{140.353}{509.676}{I}%
\TextAt{21}{145.516}{507.883}{K}%
\TextAt{8}{153.719}{509.676}{,}%
\TextAt{8}{160.868}{509.676}{and}%
\TextAt{8}{183.63}{509.676}{an}%
\TextAt{8}{199.889}{509.676}{isomorphism}%
\TextAt{13}{253.422}{485.63}{d}%
\TextAt{21}{259.505}{483.837}{K}%
\TextAt{9}{271.028}{485.63}{\Character{58}}%
\TextAt{13}{277.601}{485.63}{G}%
\TextAt{9}{286.834}{485.63}{\Character{40}}%
\TextAt{20}{294.125}{488.652}{\Character{101}}%
\TextAt{13}{291.387}{485.63}{K}%
\TextAt{16}{305.525}{485.63}{\Character{106}}%
\TextAt{13}{312.167}{485.63}{K}%
\TextAt{9}{322.985}{485.63}{\Character{41}}%
\NativeMathAt{333.445}{492.407}{7.97}{\sim}%
\NativeArrowAt{330.858}{485.63}{11.955}{13.175}{\longrightarrow}%
\TextAt{20}{348.016}{488.718}{\Character{98}}%
\TextAt{18}{347.352}{485.63}{Z}%
\TextAt{13}{355.322}{485.63}{\Character{58}}%
\TextAt{12}{79.226}{454.033}{Definition}%
\TextAt{12}{143.489}{454.033}{(4.1).}%
\TextAt{8}{180.087}{454.033}{The}%
\TextAt{8}{205.422}{454.033}{element}%
\TextAt{13}{249.946}{454.033}{\Character{39}}%
\TextAt{21}{257.614}{452.24}{K}%
\NativeMathAt{271.322}{454.033}{11.955}{\in}%
\TextAt{13}{284.797}{454.033}{G}%
\TextAt{9}{294.031}{454.033}{\Character{40}}%
\TextAt{20}{301.322}{457.055}{\Character{101}}%
\TextAt{13}{298.583}{454.033}{K}%
\TextAt{16}{314.906}{454.033}{\Character{106}}%
\TextAt{13}{323.732}{454.033}{K}%
\TextAt{9}{334.55}{454.033}{\Character{41}}%
\TextAt{8}{344.282}{454.033}{such}%
\TextAt{8}{371.96}{454.033}{that}%
\TextAt{13}{398.611}{454.033}{d}%
\TextAt{21}{404.694}{452.24}{K}%
\TextAt{9}{412.896}{454.033}{\Character{40}}%
\TextAt{13}{417.449}{454.033}{\Character{39}}%
\TextAt{21}{425.117}{452.24}{K}%
\TextAt{9}{433.32}{454.033}{\Character{41}}%
\TextAt{9}{443.362}{454.033}{\Character{61}}%
\TextAt{9}{457.96}{454.033}{1}%
\TextAt{8}{469.019}{454.033}{is}%
\TextAt{8}{482.065}{454.033}{called}%
\TextAt{8}{516.51}{454.033}{the}%
\TextAt{12}{79.226}{436.987}{Frobenius}%
\TextAt{8}{140.652}{436.987}{over}%
\TextAt{13}{165.698}{436.987}{K}%
\TextAt{8}{176.516}{436.987}{.}%
\TextAt{8}{97.159}{412.391}{For}%
\TextAt{8}{118.129}{412.391}{a}%
\TextAt{8}{127.886}{412.391}{field}%
\TextAt{8}{153.244}{412.391}{extension}%
\TextAt{13}{205.011}{412.391}{L}%
\TextAt{16}{216.298}{412.391}{\Character{106}}%
\TextAt{13}{222.941}{412.391}{K}%
\TextAt{8}{237.657}{412.391}{we}%
\TextAt{8}{254.885}{412.391}{define}%
\TextAt{8}{288.696}{412.391}{the}%
\TextAt{12}{308.853}{412.391}{inertia}%
\TextAt{12}{351.702}{412.391}{degree}%
\TextAt{13}{393.701}{412.391}{f}%
\TextAt{21}{399.471}{410.531}{L}%
\TextAt{3}{405.23}{410.531}{\Character{106}}%
\TextAt{21}{407.582}{410.531}{K}%
\TextAt{8}{419.683}{412.391}{and}%
\TextAt{8}{442.447}{412.391}{the}%
\TextAt{12}{462.605}{412.391}{ramification}%
\TextAt{12}{79.226}{395.345}{index}%
\TextAt{13}{115.044}{395.345}{e}%
\TextAt{21}{120.469}{393.485}{L}%
\TextAt{3}{126.228}{393.485}{\Character{106}}%
\TextAt{21}{128.581}{393.485}{K}%
\TextAt{8}{140.681}{395.345}{by}%
\TextAt{13}{170.329}{371.299}{f}%
\TextAt{21}{176.099}{369.439}{L}%
\TextAt{3}{181.858}{369.439}{\Character{106}}%
\TextAt{21}{184.211}{369.439}{K}%
\TextAt{9}{195.734}{371.299}{\Character{61}}%
\TextAt{9}{208.144}{371.299}{\Character{40}}%
\TextAt{13}{212.712}{371.299}{d}%
\TextAt{9}{218.795}{371.299}{\Character{40}}%
\TextAt{13}{223.347}{371.299}{G}%
\TextAt{21}{232.581}{369.506}{K}%
\TextAt{9}{240.783}{371.299}{\Character{41}}%
\TextAt{9}{248.638}{371.299}{\Character{58}}%
\TextAt{13}{255.229}{371.299}{d}%
\TextAt{9}{261.312}{371.299}{\Character{40}}%
\TextAt{13}{265.864}{371.299}{G}%
\TextAt{21}{275.098}{369.506}{L}%
\TextAt{9}{281.355}{371.299}{\Character{41}\Character{41}}%
\TextAt{8}{314.37}{371.299}{and}%
\TextAt{13}{357.145}{371.299}{e}%
\TextAt{21}{362.571}{369.439}{L}%
\TextAt{3}{368.33}{369.439}{\Character{106}}%
\TextAt{21}{370.682}{369.439}{K}%
\TextAt{9}{382.206}{371.299}{\Character{61}}%
\TextAt{9}{394.615}{371.299}{\Character{40}}%
\TextAt{13}{399.183}{371.299}{I}%
\TextAt{21}{404.346}{369.506}{K}%
\TextAt{9}{415.87}{371.299}{\Character{58}}%
\TextAt{13}{422.442}{371.299}{I}%
\TextAt{21}{427.605}{369.506}{L}%
\TextAt{9}{433.862}{371.299}{\Character{41}}%
\TextAt{13}{438.415}{371.299}{\Character{58}}%
\TextAt{8}{79.226}{347.253}{For}%
\TextAt{8}{100.196}{347.253}{a}%
\TextAt{8}{109.951}{347.253}{tower}%
\TextAt{8}{141.823}{347.253}{of}%
\TextAt{8}{155.153}{347.253}{fields}%
\TextAt{13}{185.125}{347.253}{K}%
\NativeMathAt{199.264}{347.253}{11.955}{\subseteq}%
\TextAt{13}{211.883}{347.253}{L}%
\NativeMathAt{223.168}{347.253}{11.955}{\subseteq}%
\TextAt{13}{235.788}{347.253}{M}%
\TextAt{8}{252.259}{347.253}{this}%
\TextAt{8}{275.081}{347.253}{definition}%
\TextAt{8}{327.11}{347.253}{obviously}%
\TextAt{8}{378.888}{347.253}{implies}%
\TextAt{8}{418.615}{347.253}{that}%
\TextAt{13}{185.048}{323.207}{f}%
\TextAt{21}{190.818}{321.347}{M}%
\TextAt{3}{199.799}{321.347}{\Character{106}}%
\TextAt{21}{202.152}{321.347}{K}%
\TextAt{9}{213.675}{323.207}{\Character{61}}%
\TextAt{13}{226.101}{323.207}{f}%
\TextAt{21}{231.871}{321.347}{L}%
\TextAt{3}{237.63}{321.347}{\Character{106}}%
\TextAt{21}{239.982}{321.347}{K}%
\TextAt{13}{248.185}{323.207}{f}%
\TextAt{21}{253.955}{321.347}{M}%
\TextAt{3}{262.936}{321.347}{\Character{106}}%
\TextAt{21}{265.289}{321.347}{L}%
\TextAt{8}{295.456}{323.207}{and}%
\TextAt{13}{338.232}{323.207}{e}%
\TextAt{21}{343.657}{321.347}{M}%
\TextAt{3}{352.639}{321.347}{\Character{106}}%
\TextAt{21}{354.991}{321.347}{K}%
\TextAt{9}{366.515}{323.207}{\Character{61}}%
\TextAt{13}{378.94}{323.207}{e}%
\TextAt{21}{384.366}{321.347}{L}%
\TextAt{3}{390.125}{321.347}{\Character{106}}%
\TextAt{21}{392.477}{321.347}{K}%
\TextAt{13}{400.68}{323.207}{e}%
\TextAt{21}{406.105}{321.347}{M}%
\TextAt{3}{415.087}{321.347}{\Character{106}}%
\TextAt{21}{417.439}{321.347}{L}%
\TextAt{13}{423.696}{323.207}{\Character{58}}%
\ErrataAt{79.19}{291.589}{455.81}{\textbf{Proposition (4.2).} \emph{For every extension $L\mid K$, the subgroup index satisfies}}%
\ErrataAt{224}{267.6}{196}{$[G_K:G_L]=f_{L\mid K}e_{L\mid K}.$}%
\TextAt{19}{79.226}{228.418}{Proof.}%
\TextAt{8}{116}{228.418}{The}%
\TextAt{8}{140.054}{228.418}{exact}%
\TextAt{8}{170.946}{228.418}{commutative}%
\TextAt{8}{240.872}{228.418}{diagram}%
\TextAt{9}{190.435}{204.624}{1}%
\TextAt{13}{236.259}{204.624}{I}%
\TextAt{21}{241.422}{202.831}{L}%
\TextAt{13}{288.623}{204.624}{G}%
\TextAt{21}{297.856}{202.831}{L}%
\TextAt{13}{345.056}{204.624}{d}%
\TextAt{9}{351.139}{204.624}{\Character{40}}%
\TextAt{13}{355.691}{204.624}{G}%
\TextAt{21}{364.925}{202.831}{L}%
\TextAt{9}{371.182}{204.624}{\Character{41}}%
\TextAt{9}{415.704}{204.624}{1}%
\TextAt{9}{190.435}{162.39}{1}%
\TextAt{13}{235.286}{162.39}{I}%
\TextAt{21}{240.449}{160.597}{K}%
\TextAt{13}{287.65}{162.39}{G}%
\TextAt{21}{296.883}{160.597}{K}%
\TextAt{13}{344.083}{162.39}{d}%
\TextAt{9}{350.166}{162.39}{\Character{40}}%
\TextAt{13}{354.718}{162.39}{G}%
\TextAt{21}{363.952}{160.597}{K}%
\TextAt{9}{372.154}{162.39}{\Character{41}}%
\TextAt{9}{415.704}{162.39}{1}%
\NativeDiagram{%
\path[draw=black,line width=0.478bp] (129.64,-566.387) -- (158.429,-566.387);%
\path[draw=black,line width=0.478bp,line cap=round,line join=round] (156.182,-563.518) .. controls (156.636,-565.239) and (157.648,-566.052) .. (158.668,-566.387) .. controls (157.648,-566.721) and (156.636,-567.534) .. (156.182,-569.256);%
\path[draw=black,line width=0.478bp] (181.032,-566.387) -- (210.793,-566.387);%
\path[draw=black,line width=0.478bp,line cap=round,line join=round] (208.546,-563.518) .. controls (209,-565.239) and (210.012,-566.052) .. (211.032,-566.387) .. controls (210.012,-566.721) and (209,-567.534) .. (208.546,-569.256);%
\path[draw=black,line width=0.478bp] (169.97,-575.748) -- (169.97,-598.383);%
\path[draw=black,line width=0.478bp,line cap=round,line join=round] (172.839,-596.136) .. controls (171.117,-596.59) and (170.304,-597.603) .. (169.97,-598.622) .. controls (169.635,-597.603) and (168.822,-596.59) .. (167.1,-596.136);%
\path[draw=black,line width=0.478bp] (237.466,-566.387) -- (267.228,-566.387);%
\path[draw=black,line width=0.478bp,line cap=round,line join=round] (264.98,-563.518) .. controls (265.435,-565.239) and (266.447,-566.052) .. (267.467,-566.387) .. controls (266.447,-566.721) and (265.435,-567.534) .. (264.98,-569.256);%
\path[draw=black,line width=0.478bp] (224.369,-575.748) -- (224.369,-598.383);%
\path[draw=black,line width=0.478bp,line cap=round,line join=round] (227.238,-596.136) .. controls (225.517,-596.59) and (224.704,-597.603) .. (224.369,-598.622) .. controls (224.034,-597.603) and (223.221,-596.59) .. (221.5,-596.136);%
\path[draw=black,line width=0.478bp] (309.089,-566.387) -- (337.877,-566.387);%
\path[draw=black,line width=0.478bp,line cap=round,line join=round] (335.63,-563.518) .. controls (336.084,-565.239) and (337.096,-566.052) .. (338.116,-566.387) .. controls (337.096,-566.721) and (336.084,-567.534) .. (335.63,-569.256);%
\path[draw=black,line width=0.478bp] (288.397,-576.943) -- (288.397,-597.586);%
\path[draw=black,line width=0.478bp,line cap=round,line join=round] (291.266,-595.339) .. controls (289.545,-595.793) and (288.732,-596.806) .. (288.397,-597.826) .. controls (288.063,-596.806) and (287.25,-595.793) .. (285.528,-595.339);%
\path[draw=black,line width=0.478bp] (129.64,-608.621) -- (157.456,-608.621);%
\path[draw=black,line width=0.478bp,line cap=round,line join=round] (155.209,-605.752) .. controls (155.663,-607.474) and (156.676,-608.287) .. (157.696,-608.621) .. controls (156.676,-608.956) and (155.663,-609.769) .. (155.209,-611.49);%
\path[draw=black,line width=0.478bp] (182.004,-608.621) -- (209.82,-608.621);%
\path[draw=black,line width=0.478bp,line cap=round,line join=round] (207.573,-605.752) .. controls (208.027,-607.474) and (209.04,-608.287) .. (210.06,-608.621) .. controls (209.04,-608.956) and (208.027,-609.769) .. (207.573,-611.49);%
\path[draw=black,line width=0.478bp] (238.439,-608.621) -- (266.255,-608.621);%
\path[draw=black,line width=0.478bp,line cap=round,line join=round] (264.008,-605.752) .. controls (264.462,-607.474) and (265.474,-608.287) .. (266.494,-608.621) .. controls (265.474,-608.956) and (264.462,-609.769) .. (264.008,-611.49);%
\path[draw=black,line width=0.478bp] (310.061,-608.621) -- (337.877,-608.621);%
\path[draw=black,line width=0.478bp,line cap=round,line join=round] (335.63,-605.752) .. controls (336.084,-607.474) and (337.096,-608.287) .. (338.116,-608.621) .. controls (337.096,-608.956) and (336.084,-609.769) .. (335.63,-611.49);%
}%
\TextAt{8}{79.226}{138.059}{immediately}%
\TextAt{8}{145.888}{138.059}{yields,}%
\TextAt{8}{182.041}{138.059}{if}%
\TextAt{13}{192.764}{138.059}{L}%
\TextAt{16}{204.05}{138.059}{\Character{106}}%
\TextAt{13}{210.691}{138.059}{K}%
\TextAt{8}{225.407}{138.059}{is}%
\TextAt{8}{237.17}{138.059}{Galois,}%
\TextAt{8}{276.336}{138.059}{the}%
\TextAt{8}{296.492}{138.059}{exact}%
\TextAt{8}{327.384}{138.059}{sequence}%
\TextAt{9}{173.275}{114.013}{1}%
\NativeArrowAt{182.449}{114.013}{11.955}{19.26}{\longrightarrow}%
\TextAt{13}{205.031}{114.013}{I}%
\TextAt{21}{210.194}{112.22}{K}%
\TextAt{9}{218.397}{114.013}{\Character{47}}%
\TextAt{13}{224.25}{114.013}{I}%
\TextAt{21}{229.412}{112.22}{L}%
\NativeArrowAt{238.99}{114.013}{11.955}{19.26}{\longrightarrow}%
\TextAt{13}{261.572}{114.013}{G}%
\TextAt{9}{270.806}{114.013}{\Character{40}}%
\TextAt{13}{275.358}{114.013}{L}%
\TextAt{16}{286.644}{114.013}{\Character{106}}%
\TextAt{13}{293.285}{114.013}{K}%
\TextAt{9}{304.103}{114.013}{\Character{41}}%
\NativeArrowAt{311.976}{114.013}{11.955}{19.26}{\longrightarrow}%
\TextAt{13}{334.558}{114.013}{d}%
\TextAt{9}{340.641}{114.013}{\Character{40}}%
\TextAt{13}{345.193}{114.013}{G}%
\TextAt{21}{354.427}{112.22}{K}%
\TextAt{9}{362.63}{114.013}{\Character{41}\Character{47}}%
\TextAt{13}{373.035}{114.013}{d}%
\TextAt{9}{379.118}{114.013}{\Character{40}}%
\TextAt{13}{383.67}{114.013}{G}%
\TextAt{21}{392.904}{112.22}{L}%
\TextAt{9}{399.161}{114.013}{\Character{41}}%
\NativeArrowAt{407.034}{114.013}{11.955}{19.26}{\longrightarrow}%
\TextAt{9}{429.616}{114.013}{1}%
\TextAt{13}{435.469}{114.013}{\Character{58}}%
\ErrataAt{79.19}{90.002}{455.81}{In general, factor $[G_K:G_L]$ through $I_KG_L$. The two indices are $[d(G_K):d(G_L)]$}%
\ErrataAt{79.19}{72.956}{455.81}{and $[I_K:I_L]$. For finite extensions, $[G_K:G_L]$ is the ordinary degree.}%
\Vector{q 0 G 0 g 0.399 w 452.334 -701.079 m 452.334 -693.009 l S Q}%
\Vector{q 0 G 0 g 0.399 w 452.534 -693.208 m 459.707 -693.208 l S Q}%
\Vector{q 0 G 0 g 0.399 w 452.534 -700.879 m 459.707 -700.879 l S Q}%
\Vector{q 0 G 0 g 0.399 w 459.906 -701.079 m 459.906 -693.009 l S Q}%
\TextAt{8}{297.211}{43.033}{261}%
\Destination{78.23}{809.815}{nav.page.261}%
\Destination{79.36}{243.455}{nav.print.286}%
\Destination{79.36}{468.055}{nav.number.4.4.1}%
\Destination{79.36}{305.655}{nav.number.4.4.2}%
\Destination{79.226}{305}{errata.195}%
}

% Page 274 (body page 262)
\Page{%
\TextAt{2}{79.226}{796.943}{IV.}%
\TextAt{2}{98.012}{796.943}{Abstract}%
\TextAt{2}{143.252}{796.943}{Class}%
\TextAt{2}{171.844}{796.943}{Field}%
\TextAt{2}{199.575}{796.943}{Theory}%
\Vector{q 0 G 0 g 0.299 w 7.226 18.729 m 460.769 18.729 l S Q}%
\TextAt{13}{97.159}{761.124}{L}%
\TextAt{16}{110.747}{761.124}{\Character{106}}%
\TextAt{13}{119.692}{761.124}{K}%
\TextAt{8}{135.758}{761.124}{is}%
\TextAt{8}{148.874}{761.124}{called}%
\TextAt{12}{183.389}{761.124}{unramified}%
\TextAt{8}{251.653}{761.124}{if}%
\TextAt{13}{263.728}{761.124}{e}%
\TextAt{21}{269.154}{759.264}{L}%
\TextAt{3}{274.913}{759.264}{\Character{106}}%
\TextAt{21}{277.265}{759.264}{K}%
\TextAt{9}{291.091}{761.124}{\Character{61}}%
\TextAt{9}{305.808}{761.124}{1}%
\TextAt{8}{311.673}{761.124}{,}%
\TextAt{8}{320.511}{761.124}{i.e.,}%
\TextAt{8}{344.306}{761.124}{if}%
\TextAt{13}{356.381}{761.124}{L}%
\NativeMathAt{369.969}{761.124}{11.955}{\subseteq}%
\TextAt{20}{387.63}{764.146}{\Character{101}}%
\TextAt{13}{384.891}{761.124}{K}%
\TextAt{8}{395.709}{761.124}{.}%
\TextAt{13}{408.212}{761.124}{L}%
\TextAt{16}{421.8}{761.124}{\Character{106}}%
\TextAt{13}{430.744}{761.124}{K}%
\TextAt{8}{446.811}{761.124}{is}%
\TextAt{8}{459.926}{761.124}{called}%
\TextAt{12}{494.441}{761.124}{totally}%
\TextAt{12}{79.226}{744.078}{ramified}%
\TextAt{8}{131.441}{744.078}{if}%
\TextAt{13}{142.411}{744.078}{f}%
\TextAt{21}{148.181}{742.218}{L}%
\TextAt{3}{153.94}{742.218}{\Character{106}}%
\TextAt{21}{156.292}{742.218}{K}%
\TextAt{9}{168.235}{744.078}{\Character{61}}%
\TextAt{9}{181.063}{744.078}{1}%
\TextAt{8}{186.932}{744.078}{,}%
\TextAt{8}{194.389}{744.078}{i.e.,}%
\TextAt{8}{216.802}{744.078}{if}%
\TextAt{13}{227.771}{744.078}{L}%
\TextAt{16}{238.56}{744.078}{\Character{92}}%
\TextAt{20}{252.093}{747.1}{\Character{101}}%
\TextAt{13}{249.355}{744.078}{K}%
\TextAt{9}{263.912}{744.078}{\Character{61}}%
\TextAt{13}{276.757}{744.078}{K}%
\TextAt{8}{287.575}{744.078}{.}%
\TextAt{8}{296.761}{744.078}{In}%
\TextAt{8}{311.628}{744.078}{the}%
\TextAt{8}{332.03}{744.078}{unramified}%
\TextAt{8}{390.809}{744.078}{case,}%
\TextAt{8}{419.139}{744.078}{we}%
\TextAt{8}{436.612}{744.078}{have}%
\TextAt{8}{463.853}{744.078}{the}%
\TextAt{8}{484.255}{744.078}{surjective}%
\TextAt{8}{79.226}{727.032}{homomorphism}%
\TextAt{13}{244.537}{709.986}{G}%
\TextAt{9}{253.77}{709.986}{\Character{40}}%
\TextAt{20}{261.061}{713.008}{\Character{101}}%
\TextAt{13}{258.323}{709.986}{K}%
\TextAt{16}{272.461}{709.986}{\Character{106}}%
\TextAt{13}{279.103}{709.986}{K}%
\TextAt{9}{289.921}{709.986}{\Character{41}}%
\NativeArrowAt{297.794}{709.986}{11.955}{19.26}{\longrightarrow}%
\TextAt{13}{320.376}{709.986}{G}%
\TextAt{9}{329.609}{709.986}{\Character{40}}%
\TextAt{13}{334.162}{709.986}{L}%
\TextAt{16}{345.447}{709.986}{\Character{106}}%
\TextAt{13}{352.089}{709.986}{K}%
\TextAt{9}{362.906}{709.986}{\Character{41}}%
\TextAt{8}{79.226}{686.464}{and,}%
\TextAt{8}{105.143}{686.464}{if}%
\TextAt{13}{115.744}{686.464}{f}%
\TextAt{21}{121.514}{684.671}{K}%
\TextAt{13}{133.037}{686.464}{\Character{60}}%
\NativeMathAt{145.463}{686.464}{11.955}{\infty}%
\TextAt{8}{157.418}{686.464}{,}%
\TextAt{8}{164.469}{686.464}{we}%
\TextAt{8}{181.574}{686.464}{call}%
\TextAt{8}{202.911}{686.464}{the}%
\TextAt{8}{222.945}{686.464}{image}%
\TextAt{13}{256.644}{686.464}{\Character{39}}%
\TextAt{21}{264.313}{684.604}{L}%
\TextAt{3}{270.072}{684.604}{\Character{106}}%
\TextAt{21}{272.424}{684.604}{K}%
\TextAt{8}{284.402}{686.464}{of}%
\TextAt{13}{297.609}{686.464}{\Character{39}}%
\TextAt{21}{305.278}{684.671}{K}%
\TextAt{8}{317.255}{686.464}{the}%
\TextAt{12}{337.289}{686.464}{Frobenius}%
\TextAt{12}{399.16}{686.464}{automorphism}%
\TextAt{8}{487.565}{686.464}{of}%
\TextAt{13}{500.773}{686.464}{L}%
\TextAt{16}{512.058}{686.464}{\Character{106}}%
\TextAt{13}{518.7}{686.464}{K}%
\TextAt{8}{529.518}{686.464}{.}%
\TextAt{8}{97.159}{666.429}{For}%
\TextAt{8}{118.128}{666.429}{an}%
\TextAt{8}{134.387}{666.429}{arbitrary}%
\TextAt{8}{184.157}{666.429}{extension}%
\TextAt{13}{235.923}{666.429}{L}%
\TextAt{16}{247.208}{666.429}{\Character{106}}%
\TextAt{13}{253.85}{666.429}{K}%
\TextAt{8}{268.565}{666.429}{one}%
\TextAt{8}{290.024}{666.429}{has}%
\TextAt{20}{284.112}{640.45}{\Character{101}}%
\TextAt{13}{283.125}{637.428}{L}%
\TextAt{9}{294.411}{637.428}{\Character{61}}%
\TextAt{13}{306.836}{637.428}{L}%
\TextAt{20}{317.539}{640.45}{\Character{101}}%
\TextAt{13}{314.801}{637.428}{K}%
\TextAt{13}{325.608}{637.428}{\Character{59}}%
\TextAt{8}{79.226}{608.427}{since}%
\TextAt{13}{108.188}{608.427}{L}%
\TextAt{20}{118.89}{611.449}{\Character{101}}%
\TextAt{13}{116.152}{608.427}{K}%
\TextAt{9}{130.79}{608.427}{\Character{61}}%
\TextAt{13}{143.715}{608.427}{LK}%
\TextAt{20}{162.421}{611.582}{\Character{101}}%
\TextAt{13}{162.497}{608.427}{k}%
\TextAt{9}{172.806}{608.427}{\Character{61}}%
\TextAt{13}{185.731}{608.427}{L}%
\TextAt{20}{193.62}{611.582}{\Character{101}}%
\TextAt{13}{193.696}{608.427}{k}%
\TextAt{9}{204.005}{608.427}{\Character{61}}%
\TextAt{20}{217.917}{611.449}{\Character{101}}%
\TextAt{13}{216.93}{608.427}{L}%
\TextAt{8}{224.895}{608.427}{,}%
\TextAt{8}{232.41}{608.427}{and}%
\TextAt{13}{255.466}{608.427}{L}%
\TextAt{16}{266.287}{608.427}{\Character{92}}%
\TextAt{20}{279.851}{611.449}{\Character{101}}%
\TextAt{13}{277.113}{608.427}{K}%
\TextAt{16}{291.751}{608.427}{\Character{106}}%
\TextAt{13}{298.892}{608.427}{K}%
\TextAt{8}{313.9}{608.427}{is}%
\TextAt{8}{325.957}{608.427}{the}%
\TextAt{8}{346.406}{608.427}{maximal}%
\TextAt{8}{394.508}{608.427}{unramified}%
\TextAt{8}{453.333}{608.427}{subextension}%
\TextAt{8}{523.337}{608.427}{of}%
\ErrataAt{79.19}{591.366}{455.81}{$L\mid K$. Its subgroup index is the inertia degree. With $M=L\cap\widetilde K$ one has}%
\ErrataAt{140}{574.32}{395}{$f_{L\mid K}=[G_K:G_M],\qquad [G_K:G_M]=[M:K]\quad\text{if }[M:K]\text{ is finite}.$}%
\TextAt{8}{79.226}{550.813}{Equally}%
\TextAt{8}{122.313}{550.813}{obvious}%
\TextAt{8}{164.658}{550.813}{is}%
\TextAt{8}{176.422}{550.813}{the}%
\TextAt{12}{79.226}{524.801}{Proposition}%
\TextAt{12}{152.527}{524.801}{(4.3).}%
\TextAt{8}{188.354}{524.801}{If}%
\TextAt{13}{200.558}{524.801}{f}%
\TextAt{21}{206.328}{523.008}{K}%
\TextAt{8}{218.941}{524.801}{and}%
\TextAt{13}{242.215}{524.801}{f}%
\TextAt{21}{247.985}{523.008}{L}%
\TextAt{8}{258.652}{524.801}{are}%
\TextAt{8}{278.675}{524.801}{finite,}%
\TextAt{8}{312.479}{524.801}{then}%
\TextAt{13}{339.651}{524.801}{f}%
\TextAt{21}{345.421}{522.941}{L}%
\TextAt{3}{351.18}{522.941}{\Character{106}}%
\TextAt{21}{353.532}{522.941}{K}%
\TextAt{9}{365.928}{524.801}{\Character{61}}%
\TextAt{13}{379.226}{524.801}{f}%
\TextAt{21}{384.996}{523.008}{L}%
\TextAt{9}{391.254}{524.801}{\Character{47}}%
\TextAt{13}{397.107}{524.801}{f}%
\TextAt{21}{402.877}{523.008}{K}%
\TextAt{8}{411.079}{524.801}{,}%
\TextAt{8}{418.869}{524.801}{and}%
\TextAt{8}{442.143}{524.801}{we}%
\TextAt{8}{459.883}{524.801}{have}%
\TextAt{8}{487.39}{524.801}{the}%
\TextAt{8}{508.058}{524.801}{com-}%
\TextAt{8}{79.226}{507.755}{mutative}%
\TextAt{8}{128.338}{507.755}{diagram}%
\TextAt{13}{266.609}{488.205}{G}%
\TextAt{21}{275.843}{486.411}{L}%
\TextAt{20}{324.36}{491.293}{\Character{98}}%
\TextAt{18}{323.696}{488.205}{Z}%
\TextAt{13}{265.636}{444.41}{G}%
\TextAt{21}{274.87}{442.617}{K}%
\TextAt{20}{322.734}{447.498}{\Character{98}}%
\TextAt{18}{322.07}{444.41}{Z}%
\TextAt{13}{330.04}{444.41}{\Character{58}}%
\NativeDiagram{%
\path[draw=black,line width=0.478bp] (215.451,-282.805) -- (245.866,-282.805);%
\path[draw=black,line width=0.478bp,line cap=round,line join=round] (243.618,-279.936) .. controls (244.073,-281.658) and (245.085,-282.47) .. (246.105,-282.805) .. controls (245.085,-283.14) and (244.073,-283.953) .. (243.618,-285.674);%
}%
\TextAt{21}{297.941}{495.35}{d}%
\TextAt{22}{302.298}{494.01}{L}%
\NativeDiagram{%
\path[draw=black,line width=0.478bp] (202.354,-292.166) -- (202.354,-316.363);%
\path[draw=black,line width=0.478bp,line cap=round,line join=round] (205.223,-314.115) .. controls (203.502,-314.57) and (202.689,-315.582) .. (202.354,-316.602) .. controls (202.019,-315.582) and (201.206,-314.57) .. (199.485,-314.115);%
\path[draw=black,line width=0.478bp] (255.681,-290.373) -- (255.681,-312.809);%
\path[draw=black,line width=0.478bp,line cap=round,line join=round] (258.551,-310.562) .. controls (256.829,-311.016) and (256.016,-312.029) .. (255.681,-313.048) .. controls (255.347,-312.029) and (254.534,-311.016) .. (252.812,-310.562);%
}%
\TextAt{21}{330.497}{471.018}{f}%
\TextAt{22}{334.565}{469.28}{L}%
\TextAt{23}{339.624}{469.28}{\Character{106}}%
\TextAt{22}{341.893}{469.28}{K}%
\NativeDiagram{%
\path[draw=black,line width=0.478bp] (216.424,-326.601) -- (244.24,-326.601);%
\path[draw=black,line width=0.478bp,line cap=round,line join=round] (241.993,-323.732) .. controls (242.447,-325.453) and (243.459,-326.266) .. (244.479,-326.601) .. controls (243.459,-326.935) and (242.447,-327.748) .. (241.993,-329.47);%
}%
\TextAt{21}{296.871}{439.049}{d}%
\TextAt{22}{301.228}{437.709}{K}%
\TextAt{8}{79.226}{414.776}{In}%
\TextAt{8}{93.847}{414.776}{particular,}%
\TextAt{8}{151.089}{414.776}{one}%
\TextAt{8}{172.548}{414.776}{has}%
\TextAt{13}{193.422}{414.776}{\Character{39}}%
\TextAt{21}{201.09}{412.983}{L}%
\TextAt{16}{207.347}{414.776}{\Character{106}}%
\TextAt{30}{212.639}{413.138}{\Character{101}}%
\TextAt{21}{210.668}{411.123}{K}%
\TextAt{9}{222.192}{414.776}{\Character{61}}%
\TextAt{13}{234.617}{414.776}{\Character{39}}%
\TextAt{21}{242.286}{422.127}{f}%
\TextAt{22}{246.354}{420.389}{L}%
\TextAt{23}{251.413}{420.389}{\Character{106}}%
\TextAt{22}{253.682}{420.389}{K}%
\TextAt{21}{242.286}{411.535}{K}%
\TextAt{8}{261.223}{414.776}{.}%
\TextAt{8}{97.159}{388.764}{The}%
\TextAt{8}{122.123}{388.764}{Frobenius}%
\TextAt{8}{176.904}{388.764}{automorphism}%
\TextAt{8}{255.284}{388.764}{governs}%
\TextAt{8}{298.216}{388.764}{the}%
\TextAt{8}{319.283}{388.764}{entire}%
\TextAt{8}{353.035}{388.764}{class}%
\TextAt{8}{381.382}{388.764}{field}%
\TextAt{8}{407.65}{388.764}{theory}%
\TextAt{8}{445.31}{388.764}{like}%
\TextAt{8}{467.68}{388.764}{a}%
\TextAt{8}{478.346}{388.764}{king.}%
\TextAt{8}{511.32}{388.764}{It}%
\TextAt{8}{524.903}{388.764}{is}%
\TextAt{8}{79.226}{371.718}{therefore}%
\TextAt{8}{129.973}{371.718}{most}%
\TextAt{8}{160.3}{371.718}{remarkable}%
\TextAt{8}{222.763}{371.718}{that}%
\TextAt{8}{249.778}{371.718}{in}%
\TextAt{8}{265.078}{371.718}{the}%
\TextAt{8}{286.881}{371.718}{case}%
\TextAt{8}{313.299}{371.718}{of}%
\TextAt{8}{328.276}{371.718}{a}%
\TextAt{8}{339.678}{371.718}{finite}%
\TextAt{8}{371.237}{371.718}{Galois}%
\TextAt{8}{408.797}{371.718}{extension}%
\TextAt{13}{462.21}{371.718}{L}%
\TextAt{16}{476.301}{371.718}{\Character{106}}%
\TextAt{13}{485.749}{371.718}{K}%
\TextAt{8}{496.567}{371.718}{,}%
\TextAt{8}{505.775}{371.718}{every}%
\TextAt{13}{79.226}{354.672}{\Character{27}}%
\NativeMathAt{89.971}{354.672}{11.955}{\in}%
\TextAt{13}{101.604}{354.672}{G}%
\TextAt{9}{110.838}{354.672}{\Character{40}}%
\TextAt{13}{115.39}{354.672}{L}%
\TextAt{16}{127.018}{354.672}{\Character{106}}%
\TextAt{13}{134.001}{354.672}{K}%
\TextAt{9}{144.819}{354.672}{\Character{41}}%
\TextAt{8}{153.469}{354.672}{becomes}%
\TextAt{8}{200.224}{354.672}{a}%
\TextAt{8}{210.18}{354.672}{Frobenius}%
\TextAt{8}{264.25}{354.672}{automorphism}%
\TextAt{8}{341.92}{354.672}{once}%
\TextAt{8}{368.781}{354.672}{it}%
\TextAt{8}{380.686}{354.672}{is}%
\TextAt{8}{392.651}{354.672}{manoeuvered}%
\TextAt{8}{464.068}{354.672}{into}%
\TextAt{8}{488.012}{354.672}{the}%
\TextAt{8}{508.369}{354.672}{right}%
\TextAt{8}{79.226}{337.626}{position.}%
\TextAt{8}{133.312}{337.626}{This}%
\TextAt{8}{161.672}{337.626}{is}%
\TextAt{8}{175.075}{337.626}{achieved}%
\TextAt{8}{223.866}{337.626}{in}%
\TextAt{8}{239.158}{337.626}{the}%
\TextAt{8}{260.954}{337.626}{following}%
\TextAt{8}{312.028}{337.626}{manner.}%
\TextAt{8}{363.771}{337.626}{For}%
\TextAt{8}{386.38}{337.626}{what}%
\TextAt{8}{417.286}{337.626}{follows,}%
\TextAt{8}{461.023}{337.626}{let}%
\TextAt{8}{479.567}{337.626}{us}%
\TextAt{8}{496.223}{337.626}{assume}%
\TextAt{8}{79.226}{320.58}{systematically}%
\TextAt{8}{155.728}{320.58}{that}%
\TextAt{13}{181.373}{320.58}{f}%
\TextAt{21}{187.143}{318.787}{K}%
\TextAt{13}{199.136}{320.58}{\Character{60}}%
\NativeMathAt{212.031}{320.58}{11.955}{\infty}%
\TextAt{8}{223.987}{320.58}{.}%
\TextAt{8}{233.262}{320.58}{We}%
\TextAt{8}{253.682}{320.58}{pass}%
\TextAt{8}{279.446}{320.58}{from}%
\TextAt{8}{307.362}{320.58}{the}%
\TextAt{8}{327.795}{320.58}{Galois}%
\TextAt{8}{363.984}{320.58}{extension}%
\TextAt{13}{416.025}{320.58}{L}%
\TextAt{16}{427.78}{320.58}{\Character{106}}%
\TextAt{13}{434.892}{320.58}{K}%
\TextAt{8}{449.883}{320.58}{to}%
\TextAt{8}{464.469}{320.58}{the}%
\TextAt{8}{484.901}{320.58}{extension}%
\TextAt{20}{80.213}{306.556}{\Character{101}}%
\TextAt{13}{79.226}{303.534}{L}%
\TextAt{16}{90.511}{303.534}{\Character{106}}%
\TextAt{13}{97.153}{303.534}{K}%
\TextAt{8}{111.868}{303.534}{and}%
\TextAt{8}{134.631}{303.534}{consider}%
\TextAt{8}{180.216}{303.534}{in}%
\TextAt{8}{193.869}{303.534}{the}%
\TextAt{8}{214.025}{303.534}{Galois}%
\TextAt{8}{249.939}{303.534}{group}%
\TextAt{13}{283.114}{303.534}{G}%
\TextAt{9}{292.348}{303.534}{\Character{40}}%
\TextAt{20}{297.887}{306.556}{\Character{101}}%
\TextAt{13}{296.9}{303.534}{L}%
\TextAt{16}{308.185}{303.534}{\Character{106}}%
\TextAt{13}{314.827}{303.534}{K}%
\TextAt{9}{325.645}{303.534}{\Character{41}}%
\TextAt{8}{334.095}{303.534}{the}%
\TextAt{8}{354.251}{303.534}{semigroup}%
\TextAt{8}{193.089}{274.533}{Frob}%
\TextAt{9}{217.645}{274.533}{\Character{40}}%
\TextAt{20}{223.184}{277.555}{\Character{101}}%
\TextAt{13}{222.198}{274.533}{L}%
\TextAt{16}{233.483}{274.533}{\Character{106}}%
\TextAt{13}{240.125}{274.533}{K}%
\TextAt{9}{250.942}{274.533}{\Character{41}}%
\TextAt{9}{258.797}{274.533}{\Character{61}}%
\TextAt{16}{271.241}{274.533}{\Character{102}}%
\TextAt{13}{277.219}{274.533}{\Character{27}}%
\NativeMathAt{287.622}{274.533}{11.955}{\in}%
\TextAt{13}{298.913}{274.533}{G}%
\TextAt{9}{308.147}{274.533}{\Character{40}}%
\TextAt{20}{313.685}{277.555}{\Character{101}}%
\TextAt{13}{312.699}{274.533}{L}%
\TextAt{16}{323.984}{274.533}{\Character{106}}%
\TextAt{13}{330.626}{274.533}{K}%
\TextAt{9}{341.444}{274.533}{\Character{41}}%
\TextAt{16}{349.317}{274.533}{\Character{106}}%
\TextAt{13}{355.958}{274.533}{d}%
\TextAt{21}{362.041}{272.74}{K}%
\TextAt{9}{370.244}{274.533}{\Character{40}}%
\TextAt{13}{374.796}{274.533}{\Character{27}}%
\TextAt{9}{381.879}{274.533}{\Character{41}}%
\NativeMathAt{389.752}{274.533}{11.955}{\in}%
\TextAt{18}{401.043}{274.533}{N}%
\TextAt{16}{409.677}{274.533}{\Character{103}}%
\TextAt{13}{415.655}{274.533}{\Character{58}}%
\TextAt{8}{79.226}{245.532}{Observe}%
\TextAt{8}{124.976}{245.532}{here}%
\TextAt{8}{151.144}{245.532}{that}%
\TextAt{13}{177.324}{245.532}{d}%
\TextAt{21}{183.406}{243.739}{K}%
\TextAt{9}{196.311}{245.532}{\Character{58}}%
\TextAt{13}{204.265}{245.532}{G}%
\TextAt{21}{213.499}{243.739}{K}%
\NativeMathAt{226.404}{245.532}{11.955}{\rightarrow}%
\TextAt{20}{243.726}{248.62}{\Character{98}}%
\TextAt{18}{243.062}{245.532}{Z}%
\TextAt{8}{255.74}{245.532}{factorizes}%
\TextAt{8}{308.317}{245.532}{through}%
\TextAt{13}{353.362}{245.532}{G}%
\TextAt{9}{362.595}{245.532}{\Character{40}}%
\TextAt{20}{368.134}{248.554}{\Character{101}}%
\TextAt{13}{367.148}{245.532}{L}%
\TextAt{16}{379.815}{245.532}{\Character{106}}%
\TextAt{13}{387.838}{245.532}{K}%
\TextAt{9}{398.656}{245.532}{\Character{41}}%
\TextAt{8}{407.916}{245.532}{because}%
\TextAt{13}{452.028}{245.532}{G}%
\TextAt{30}{462.024}{243.894}{\Character{101}}%
\TextAt{21}{461.262}{241.879}{L}%
\TextAt{9}{472.222}{245.532}{\Character{61}}%
\TextAt{13}{486.029}{245.532}{I}%
\TextAt{21}{491.192}{243.739}{L}%
\NativeMathAt{502.151}{245.532}{11.955}{\subseteq}%
\TextAt{13}{516.152}{245.532}{I}%
\TextAt{21}{521.315}{243.739}{K}%
\TextAt{8}{529.518}{245.532}{;}%
\TextAt{8}{79.226}{228.486}{recall}%
\TextAt{8}{110.441}{228.486}{also}%
\TextAt{8}{133.921}{228.486}{that}%
\TextAt{9}{159.29}{228.486}{0}%
\NativeMathAt{168.464}{228.486}{11.955}{\notin}%
\TextAt{18}{179.755}{228.486}{N}%
\TextAt{8}{188.389}{228.486}{.}%
\TextAt{8}{196.837}{228.486}{Firstly,}%
\TextAt{8}{237.055}{228.486}{we}%
\TextAt{8}{254.282}{228.486}{have}%
\TextAt{8}{281.277}{228.486}{the}%
\TextAt{12}{79.226}{202.474}{Proposition}%
\TextAt{12}{151.937}{202.474}{(4.4).}%
\TextAt{8}{187.253}{202.474}{For}%
\TextAt{8}{208.222}{202.474}{a}%
\TextAt{8}{217.978}{202.474}{finite}%
\TextAt{8}{247.89}{202.474}{Galois}%
\TextAt{8}{283.803}{202.474}{extension}%
\TextAt{13}{335.569}{202.474}{L}%
\TextAt{16}{346.854}{202.474}{\Character{106}}%
\TextAt{13}{353.496}{202.474}{K}%
\TextAt{8}{368.211}{202.474}{the}%
\TextAt{8}{388.367}{202.474}{mapping}%
\TextAt{8}{197.276}{173.473}{Frob}%
\TextAt{9}{221.832}{173.472}{\Character{40}}%
\TextAt{20}{227.371}{176.495}{\Character{101}}%
\TextAt{13}{226.385}{173.473}{L}%
\TextAt{16}{237.67}{173.472}{\Character{106}}%
\TextAt{13}{244.312}{173.473}{K}%
\TextAt{9}{255.129}{173.472}{\Character{41}}%
\NativeArrowAt{263.002}{173.473}{11.955}{19.26}{\longrightarrow}%
\TextAt{13}{285.584}{173.472}{G}%
\TextAt{9}{294.818}{173.473}{\Character{40}}%
\TextAt{13}{299.37}{173.472}{L}%
\TextAt{16}{310.656}{173.473}{\Character{106}}%
\TextAt{13}{317.297}{173.472}{K}%
\TextAt{9}{328.115}{173.473}{\Character{41}}%
\TextAt{13}{332.667}{173.472}{\Character{59}}%
\TextAt{13}{361.802}{173.472}{\Character{27}}%
\NativeArrowAt{372.225}{173.473}{11.955}{19.26}{\longmapsto}%
\TextAt{13}{394.807}{173.472}{\Character{27}}%
\TextAt{16}{401.89}{173.473}{\Character{106}}%
\TextAt{21}{405.21}{171.679}{L}%
\TextAt{13}{411.468}{173.472}{\Character{59}}%
\TextAt{8}{79.226}{144.471}{is}%
\TextAt{8}{90.99}{144.471}{surjective.}%
\TextAt{19}{79.226}{109.493}{Proof.}%
\TextAt{8}{116}{109.493}{Let}%
\TextAt{13}{136.509}{109.493}{\Character{27}}%
\NativeMathAt{146.912}{109.493}{11.955}{\in}%
\TextAt{13}{158.203}{109.493}{G}%
\TextAt{9}{167.437}{109.493}{\Character{40}}%
\TextAt{13}{171.989}{109.493}{L}%
\TextAt{16}{183.274}{109.493}{\Character{106}}%
\TextAt{13}{189.916}{109.493}{K}%
\TextAt{9}{200.734}{109.493}{\Character{41}}%
\TextAt{8}{208.723}{109.493}{and}%
\TextAt{8}{231.025}{109.493}{let}%
\TextAt{13}{247.469}{109.493}{\Character{39}}%
\NativeMathAt{258.458}{109.493}{11.955}{\in}%
\TextAt{13}{269.749}{109.493}{G}%
\TextAt{9}{278.983}{109.493}{\Character{40}}%
\TextAt{20}{284.522}{112.515}{\Character{101}}%
\TextAt{13}{283.535}{109.493}{L}%
\TextAt{16}{294.82}{109.493}{\Character{106}}%
\TextAt{13}{301.462}{109.493}{K}%
\TextAt{9}{312.28}{109.493}{\Character{41}}%
\TextAt{8}{320.269}{109.493}{be}%
\TextAt{8}{335.733}{109.493}{an}%
\TextAt{8}{351.531}{109.493}{element}%
\TextAt{8}{394.312}{109.493}{such}%
\TextAt{8}{420.249}{109.493}{that}%
\TextAt{13}{445.157}{109.493}{d}%
\TextAt{21}{451.24}{107.699}{K}%
\TextAt{9}{459.442}{109.493}{\Character{40}}%
\TextAt{13}{463.995}{109.493}{\Character{39}}%
\TextAt{9}{471.663}{109.493}{\Character{41}}%
\TextAt{9}{479.518}{109.493}{\Character{61}}%
\TextAt{9}{491.939}{109.493}{1}%
\TextAt{8}{497.815}{109.493}{.}%
\TextAt{8}{506.11}{109.493}{Then}%
\TextAt{13}{79.226}{92.447}{\Character{39}}%
\TextAt{16}{86.895}{92.447}{\Character{106}}%
\TextAt{30}{92.186}{90.808}{\Character{101}}%
\TextAt{21}{90.215}{88.794}{K}%
\TextAt{9}{102.34}{92.447}{\Character{61}}%
\TextAt{13}{115.367}{92.447}{\Character{39}}%
\TextAt{21}{123.035}{90.653}{K}%
\TextAt{8}{135.488}{92.447}{and}%
\TextAt{13}{158.604}{92.447}{\Character{39}}%
\TextAt{16}{166.272}{92.447}{\Character{106}}%
\TextAt{21}{169.593}{88.794}{L}%
\TextAt{3}{175.352}{88.794}{\Character{92}}%
\TextAt{30}{182.968}{90.808}{\Character{101}}%
\TextAt{21}{180.998}{88.794}{K}%
\TextAt{9}{193.122}{92.447}{\Character{61}}%
\TextAt{13}{206.149}{92.447}{\Character{39}}%
\TextAt{21}{213.818}{88.794}{L}%
\TextAt{3}{219.577}{88.794}{\Character{92}}%
\TextAt{30}{227.193}{90.808}{\Character{101}}%
\TextAt{21}{225.222}{88.794}{K}%
\TextAt{3}{232.927}{88.794}{\Character{106}}%
\TextAt{21}{235.279}{88.794}{K}%
\TextAt{8}{243.482}{92.447}{.}%
\TextAt{8}{252.989}{92.447}{Restricting}%
\TextAt{13}{313.404}{92.447}{\Character{27}}%
\TextAt{8}{324.737}{92.447}{to}%
\TextAt{8}{339.4}{92.447}{the}%
\TextAt{8}{359.909}{92.447}{maximal}%
\TextAt{8}{408.071}{92.447}{unramified}%
\TextAt{8}{466.956}{92.447}{subextension}%
\TextAt{8}{297.211}{43.033}{262}%
\Destination{78.23}{809.815}{nav.page.262}%
\Destination{79.36}{403.755}{nav.print.287}%
\Destination{79.36}{538.755}{nav.number.4.4.3}%
\Destination{79.36}{216.455}{nav.number.4.4.4}%
}

% Page 275 (body page 263)
\Page{%
\TextAt{2}{79.226}{796.943}{IV.}%
\TextAt{2}{98.012}{796.943}{Abstract}%
\TextAt{2}{143.252}{796.943}{Class}%
\TextAt{2}{171.844}{796.943}{Field}%
\TextAt{2}{199.575}{796.943}{Theory}%
\Vector{q 0 G 0 g 0.299 w 7.226 18.729 m 460.769 18.729 l S Q}%
\TextAt{13}{79.226}{761.124}{L}%
\TextAt{16}{90.031}{761.124}{\Character{92}}%
\TextAt{20}{103.58}{764.146}{\Character{101}}%
\TextAt{13}{100.842}{761.124}{K}%
\TextAt{16}{115.44}{761.124}{\Character{106}}%
\TextAt{13}{122.541}{761.124}{K}%
\TextAt{8}{133.359}{761.124}{,}%
\TextAt{8}{140.845}{761.124}{it}%
\TextAt{8}{152.819}{761.124}{becomes}%
\TextAt{8}{199.642}{761.124}{a}%
\TextAt{8}{209.667}{761.124}{power}%
\TextAt{8}{244.081}{761.124}{of}%
\TextAt{8}{257.681}{761.124}{the}%
\TextAt{8}{278.107}{761.124}{Frobenius}%
\TextAt{8}{332.246}{761.124}{automorphism,}%
\TextAt{13}{413.305}{761.124}{\Character{27}}%
\TextAt{16}{420.387}{761.124}{\Character{106}}%
\TextAt{21}{423.708}{757.471}{L}%
\TextAt{3}{429.467}{757.471}{\Character{92}}%
\TextAt{30}{437.083}{759.485}{\Character{101}}%
\TextAt{21}{435.113}{757.471}{K}%
\TextAt{9}{447.096}{761.124}{\Character{61}}%
\TextAt{13}{459.981}{761.124}{\Character{39}}%
\TextAt{21}{467.649}{765.462}{n}%
\TextAt{21}{467.649}{755.779}{L}%
\TextAt{3}{473.408}{755.779}{\Character{92}}%
\TextAt{30}{481.024}{757.793}{\Character{101}}%
\TextAt{21}{479.054}{755.779}{K}%
\TextAt{3}{486.758}{755.779}{\Character{106}}%
\TextAt{21}{489.111}{755.779}{K}%
\TextAt{8}{497.313}{761.124}{,}%
\TextAt{8}{504.8}{761.124}{so}%
\TextAt{8}{519.439}{761.124}{we}%
\TextAt{8}{79.226}{741.133}{may}%
\TextAt{8}{104.595}{741.133}{choose}%
\TextAt{13}{141.728}{741.133}{n}%
\TextAt{8}{152.613}{741.133}{in}%
\TextAt{18}{166.266}{741.133}{N}%
\TextAt{8}{174.9}{741.133}{.}%
\TextAt{8}{183.348}{741.133}{As}%
\TextAt{20}{201.622}{744.155}{\Character{101}}%
\TextAt{13}{200.635}{741.133}{L}%
\TextAt{9}{211.921}{741.133}{\Character{61}}%
\TextAt{13}{224.346}{741.133}{L}%
\TextAt{20}{235.049}{744.155}{\Character{101}}%
\TextAt{13}{232.311}{741.133}{K}%
\TextAt{8}{243.128}{741.133}{,}%
\TextAt{8}{250.278}{741.133}{we}%
\TextAt{8}{267.505}{741.133}{have}%
\TextAt{13}{238.695}{713.398}{G}%
\TextAt{9}{247.929}{713.398}{\Character{40}}%
\TextAt{20}{253.467}{716.42}{\Character{101}}%
\TextAt{13}{252.481}{713.398}{L}%
\TextAt{16}{263.766}{713.398}{\Character{106}}%
\TextAt{20}{273.146}{716.42}{\Character{101}}%
\TextAt{13}{270.408}{713.398}{K}%
\TextAt{9}{281.226}{713.398}{\Character{41}}%
\NativeMathAt{289.099}{713.398}{11.955}{\cong}%
\TextAt{13}{301.718}{713.398}{G}%
\TextAt{9}{310.952}{713.398}{\Character{40}}%
\TextAt{13}{315.504}{713.398}{L}%
\TextAt{16}{326.789}{713.398}{\Character{106}}%
\TextAt{13}{333.431}{713.398}{L}%
\TextAt{16}{344.052}{713.398}{\Character{92}}%
\TextAt{20}{357.417}{716.42}{\Character{101}}%
\TextAt{13}{354.679}{713.398}{K}%
\TextAt{9}{365.497}{713.398}{\Character{41}}%
\TextAt{13}{370.049}{713.398}{\Character{58}}%
\TextAt{8}{79.226}{685.662}{If}%
\TextAt{8}{91.323}{685.662}{now}%
\TextAt{13}{116.115}{685.662}{\Character{28}}%
\NativeMathAt{126.544}{685.662}{11.955}{\in}%
\TextAt{13}{138.524}{685.662}{G}%
\TextAt{9}{147.758}{685.662}{\Character{40}}%
\TextAt{20}{153.297}{688.684}{\Character{101}}%
\TextAt{13}{152.31}{685.662}{L}%
\TextAt{16}{164.284}{685.662}{\Character{106}}%
\TextAt{20}{174.353}{688.684}{\Character{101}}%
\TextAt{13}{171.615}{685.662}{K}%
\TextAt{9}{182.433}{685.662}{\Character{41}}%
\TextAt{8}{191.287}{685.662}{is}%
\TextAt{8}{203.455}{685.662}{mapped}%
\TextAt{8}{248.404}{685.662}{to}%
\TextAt{13}{263.119}{685.662}{\Character{27}}%
\TextAt{13}{270.185}{685.662}{\Character{39}}%
\TextAt{3}{277.87}{690}{\Character{0}}%
\TextAt{21}{284.456}{690}{n}%
\TextAt{16}{290.092}{685.662}{\Character{106}}%
\TextAt{21}{293.413}{683.869}{L}%
\TextAt{8}{303.972}{685.662}{under}%
\TextAt{8}{337.54}{685.662}{this}%
\TextAt{8}{360.767}{685.662}{isomorphism,}%
\TextAt{8}{432.943}{685.662}{then}%
\TextAt{9}{460.622}{685.662}{\Character{22}}%
\TextAt{13}{460.008}{685.662}{\Character{27}}%
\TextAt{9}{471.1}{685.662}{\Character{61}}%
\TextAt{13}{484.214}{685.662}{\Character{28}}%
\TextAt{13}{490.622}{685.662}{\Character{39}}%
\TextAt{21}{498.302}{690}{n}%
\TextAt{8}{508.24}{685.662}{is}%
\TextAt{8}{520.408}{685.662}{an}%
\TextAt{8}{79.226}{668.616}{element}%
\TextAt{8}{122.798}{668.616}{satisfying}%
\TextAt{9}{175.903}{668.616}{\Character{22}}%
\TextAt{13}{175.288}{668.616}{\Character{27}}%
\TextAt{16}{182.37}{668.616}{\Character{106}}%
\TextAt{21}{185.691}{666.823}{L}%
\TextAt{9}{195.831}{668.616}{\Character{61}}%
\TextAt{13}{208.819}{668.616}{\Character{28}}%
\TextAt{13}{215.227}{668.616}{\Character{39}}%
\TextAt{21}{222.906}{672.954}{n}%
\TextAt{16}{228.543}{668.616}{\Character{106}}%
\TextAt{21}{231.863}{666.823}{L}%
\TextAt{9}{242.003}{668.616}{\Character{61}}%
\TextAt{13}{254.991}{668.616}{\Character{27}}%
\TextAt{13}{262.057}{668.616}{\Character{39}}%
\TextAt{3}{269.742}{672.954}{\Character{0}}%
\TextAt{21}{276.328}{672.954}{n}%
\TextAt{13}{281.965}{668.616}{\Character{39}}%
\TextAt{21}{289.633}{672.954}{n}%
\TextAt{16}{295.269}{668.616}{\Character{106}}%
\TextAt{21}{298.59}{666.823}{L}%
\TextAt{9}{308.73}{668.616}{\Character{61}}%
\TextAt{13}{321.718}{668.616}{\Character{27}}%
\TextAt{8}{333.027}{668.616}{and}%
\TextAt{9}{356.735}{668.616}{\Character{22}}%
\TextAt{13}{356.12}{668.616}{\Character{27}}%
\TextAt{16}{363.202}{668.616}{\Character{106}}%
\TextAt{30}{368.494}{666.978}{\Character{101}}%
\TextAt{21}{366.523}{664.963}{K}%
\TextAt{9}{378.609}{668.616}{\Character{61}}%
\TextAt{13}{391.596}{668.616}{\Character{39}}%
\TextAt{21}{399.265}{672.954}{n}%
\TextAt{21}{399.265}{665.595}{K}%
\TextAt{8}{407.467}{668.616}{.}%
\TextAt{8}{416.905}{668.616}{Hence}%
\TextAt{13}{452.012}{668.616}{d}%
\TextAt{21}{458.095}{666.823}{K}%
\TextAt{9}{466.298}{668.616}{\Character{40}}%
\TextAt{9}{471.451}{668.616}{\Character{22}}%
\TextAt{13}{470.85}{668.616}{\Character{27}}%
\TextAt{9}{477.932}{668.616}{\Character{41}}%
\TextAt{9}{486.349}{668.616}{\Character{61}}%
\TextAt{13}{499.355}{668.616}{n}%
\TextAt{8}{506.343}{668.616}{,}%
\TextAt{8}{513.904}{668.616}{and}%
\TextAt{8}{79.226}{651.57}{so}%
\TextAt{9}{94.211}{651.57}{\Character{22}}%
\TextAt{13}{93.596}{651.57}{\Character{27}}%
\NativeMathAt{103.999}{651.57}{11.955}{\in}%
\TextAt{8}{115.29}{651.57}{Frob}%
\TextAt{9}{139.846}{651.57}{\Character{40}}%
\TextAt{20}{145.385}{654.592}{\Character{101}}%
\TextAt{13}{144.399}{651.57}{L}%
\TextAt{16}{155.684}{651.57}{\Character{106}}%
\TextAt{13}{162.326}{651.57}{K}%
\TextAt{9}{173.143}{651.57}{\Character{41}}%
\TextAt{8}{177.696}{651.57}{.}%
\Vector{q 0 G 0 g 0.399 w 452.334 -122.43 m 452.334 -114.36 l S Q}%
\Vector{q 0 G 0 g 0.399 w 452.534 -114.559 m 459.707 -114.559 l S Q}%
\Vector{q 0 G 0 g 0.399 w 452.534 -122.231 m 459.707 -122.231 l S Q}%
\Vector{q 0 G 0 g 0.399 w 459.906 -122.43 m 459.906 -114.36 l S Q}%
\TextAt{8}{97.159}{622.93}{Thus}%
\TextAt{8}{126.849}{622.93}{every}%
\TextAt{8}{157.783}{622.93}{element}%
\TextAt{13}{201.067}{622.93}{\Character{27}}%
\NativeMathAt{211.541}{622.93}{11.955}{\in}%
\TextAt{13}{222.903}{622.93}{G}%
\TextAt{9}{232.137}{622.93}{\Character{40}}%
\TextAt{13}{236.689}{622.93}{L}%
\TextAt{16}{248.045}{622.93}{\Character{106}}%
\TextAt{13}{254.758}{622.93}{K}%
\TextAt{9}{265.576}{622.93}{\Character{41}}%
\TextAt{8}{274.067}{622.93}{may}%
\TextAt{8}{299.478}{622.93}{be}%
\TextAt{8}{315.444}{622.93}{lifted}%
\TextAt{8}{345.72}{622.93}{to}%
\TextAt{8}{360.072}{622.93}{an}%
\TextAt{8}{376.373}{622.93}{element}%
\TextAt{9}{420.271}{622.93}{\Character{22}}%
\TextAt{13}{419.656}{622.93}{\Character{27}}%
\NativeMathAt{430.13}{622.93}{11.955}{\in}%
\TextAt{8}{441.492}{622.93}{Frob}%
\TextAt{9}{466.048}{622.93}{\Character{40}}%
\TextAt{20}{471.587}{625.952}{\Character{101}}%
\TextAt{13}{470.601}{622.93}{L}%
\TextAt{16}{481.957}{622.93}{\Character{106}}%
\TextAt{13}{488.67}{622.93}{K}%
\TextAt{9}{499.487}{622.93}{\Character{41}}%
\TextAt{8}{504.04}{622.93}{.}%
\TextAt{8}{512.613}{622.93}{The}%
\TextAt{8}{79.226}{605.884}{following}%
\TextAt{8}{128.283}{605.884}{proposition}%
\TextAt{8}{189.439}{605.884}{shows}%
\TextAt{8}{222.679}{605.884}{that}%
\TextAt{8}{247.67}{605.884}{this}%
\TextAt{8}{270.115}{605.884}{lifting,}%
\TextAt{8}{307.209}{605.884}{considered}%
\TextAt{8}{364.12}{605.884}{over}%
\TextAt{8}{388.789}{605.884}{its}%
\TextAt{8}{404.73}{605.884}{fixed}%
\TextAt{8}{432.638}{605.884}{field,}%
\TextAt{8}{460.945}{605.884}{is}%
\TextAt{8}{472.331}{605.884}{actually}%
\TextAt{8}{516.51}{605.884}{the}%
\TextAt{8}{79.226}{588.839}{Frobenius}%
\TextAt{8}{133.096}{588.839}{automorphism.}%
\TextAt{12}{79.226}{563.188}{Proposition}%
\TextAt{12}{152.869}{563.188}{(4.5).}%
\TextAt{8}{188.995}{563.188}{Let}%
\TextAt{9}{211.389}{563.188}{\Character{22}}%
\TextAt{13}{210.774}{563.188}{\Character{27}}%
\NativeMathAt{222.558}{563.188}{11.955}{\in}%
\TextAt{8}{235.23}{563.188}{Frob}%
\TextAt{9}{259.786}{563.188}{\Character{40}}%
\TextAt{20}{265.324}{566.21}{\Character{101}}%
\TextAt{13}{264.338}{563.188}{L}%
\TextAt{16}{277.004}{563.188}{\Character{106}}%
\TextAt{13}{285.026}{563.188}{K}%
\TextAt{9}{295.844}{563.188}{\Character{41}}%
\TextAt{8}{300.396}{563.188}{,}%
\TextAt{8}{308.558}{563.188}{and}%
\TextAt{8}{332.13}{563.188}{let}%
\TextAt{9}{349.845}{563.188}{\Character{6}}%
\TextAt{8}{363.007}{563.188}{be}%
\TextAt{8}{379.741}{563.188}{the}%
\TextAt{8}{400.708}{563.188}{fixed}%
\TextAt{8}{429.804}{563.188}{field}%
\TextAt{8}{455.971}{563.188}{of}%
\TextAt{9}{470.726}{563.188}{\Character{22}}%
\TextAt{13}{470.111}{563.188}{\Character{27}}%
\TextAt{8}{477.193}{563.188}{.}%
\TextAt{8}{488.072}{563.188}{Then}%
\TextAt{8}{519.439}{563.188}{we}%
\TextAt{8}{79.226}{546.142}{have:}%
\TextAt{8}{137.434}{529.096}{(i)}%
\TextAt{13}{153.693}{529.096}{f}%
\TextAt{17}{159.463}{527.236}{\Character{6}}%
\TextAt{3}{165.579}{527.236}{\Character{106}}%
\TextAt{21}{167.931}{527.236}{K}%
\TextAt{9}{179.455}{529.096}{\Character{61}}%
\TextAt{13}{191.88}{529.096}{d}%
\TextAt{21}{197.963}{527.303}{K}%
\TextAt{9}{206.165}{529.096}{\Character{40}}%
\TextAt{9}{211.318}{529.096}{\Character{22}}%
\TextAt{13}{210.718}{529.096}{\Character{27}}%
\TextAt{9}{217.8}{529.096}{\Character{41}}%
\TextAt{13}{222.352}{529.096}{\Character{59}}%
\TextAt{8}{239.552}{529.096}{(ii)}%
\TextAt{9}{259.063}{529.096}{\Character{91}\Character{6}}%
\TextAt{9}{274.067}{529.096}{\Character{58}}%
\TextAt{13}{280.662}{529.096}{K}%
\TextAt{9}{291.48}{529.096}{\Character{93}}%
\TextAt{13}{298.052}{529.096}{\Character{60}}%
\NativeMathAt{310.478}{529.096}{11.955}{\infty}%
\TextAt{13}{322.433}{529.096}{\Character{59}}%
\TextAt{8}{339.632}{529.096}{(iii)}%
\TextAt{20}{363.301}{532.118}{\Character{101}}%
\TextAt{9}{362.395}{529.096}{\Character{6}}%
\TextAt{9}{374.159}{529.096}{\Character{61}}%
\TextAt{20}{387.582}{532.118}{\Character{101}}%
\TextAt{13}{386.595}{529.096}{L\Character{59}}%
\TextAt{8}{411.759}{529.096}{(iv)}%
\TextAt{9}{434.814}{529.096}{\Character{22}}%
\TextAt{13}{434.199}{529.096}{\Character{27}}%
\TextAt{9}{444.602}{529.096}{\Character{61}}%
\TextAt{13}{457.028}{529.096}{\Character{39}}%
\TextAt{17}{464.696}{527.303}{\Character{6}}%
\TextAt{13}{471.31}{529.096}{\Character{58}}%
\TextAt{19}{79.226}{488.907}{Proof.}%
\TextAt{8}{116}{488.907}{(i)}%
\TextAt{9}{132.259}{488.907}{\Character{6}}%
\TextAt{16}{143.37}{488.907}{\Character{92}}%
\TextAt{20}{156.735}{491.929}{\Character{101}}%
\TextAt{13}{153.997}{488.907}{K}%
\TextAt{8}{168.712}{488.907}{is}%
\TextAt{8}{180.476}{488.907}{the}%
\TextAt{8}{200.632}{488.907}{fixed}%
\TextAt{8}{228.918}{488.907}{field}%
\TextAt{8}{254.275}{488.907}{of}%
\TextAt{9}{268.22}{488.907}{\Character{22}}%
\TextAt{13}{267.605}{488.907}{\Character{27}}%
\TextAt{16}{274.688}{488.907}{\Character{106}}%
\TextAt{30}{279.979}{487.268}{\Character{101}}%
\TextAt{21}{278.009}{485.254}{K}%
\TextAt{9}{289.532}{488.907}{\Character{61}}%
\TextAt{13}{301.958}{488.907}{\Character{39}}%
\TextAt{21}{309.626}{495.017}{d}%
\TextAt{22}{313.983}{493.678}{K}%
\TextAt{17}{321.026}{495.017}{\Character{40}}%
\TextAt{17}{324.756}{495.017}{\Character{22}}%
\TextAt{21}{324.32}{495.017}{\Character{27}}%
\TextAt{17}{329.445}{495.017}{\Character{41}}%
\TextAt{21}{309.626}{485.666}{K}%
\TextAt{8}{333.236}{488.907}{,}%
\TextAt{8}{340.385}{488.907}{so}%
\TextAt{8}{354.756}{488.907}{that}%
\TextAt{13}{232.284}{461.171}{f}%
\TextAt{17}{238.054}{459.311}{\Character{6}}%
\TextAt{3}{244.17}{459.311}{\Character{106}}%
\TextAt{21}{246.522}{459.311}{K}%
\TextAt{9}{258.046}{461.171}{\Character{61}}%
\TextAt{9}{270.456}{461.171}{\Character{91}\Character{6}}%
\TextAt{16}{284.834}{461.171}{\Character{92}}%
\TextAt{20}{298.199}{464.193}{\Character{101}}%
\TextAt{13}{295.461}{461.171}{K}%
\TextAt{9}{309.599}{461.171}{\Character{58}}%
\TextAt{13}{316.172}{461.171}{K}%
\TextAt{9}{326.989}{461.171}{\Character{93}}%
\TextAt{9}{333.54}{461.171}{\Character{61}}%
\TextAt{13}{345.987}{461.171}{d}%
\TextAt{21}{352.07}{459.378}{K}%
\TextAt{9}{360.273}{461.171}{\Character{40}}%
\TextAt{9}{365.426}{461.171}{\Character{22}}%
\TextAt{13}{364.825}{461.171}{\Character{27}}%
\TextAt{9}{371.908}{461.171}{\Character{41}}%
\TextAt{13}{376.46}{461.171}{\Character{58}}%
\TextAt{8}{79.226}{433.435}{(ii)}%
\TextAt{8}{98.737}{433.435}{One}%
\TextAt{8}{123.448}{433.435}{has}%
\TextAt{20}{147.06}{436.457}{\Character{101}}%
\TextAt{13}{144.322}{433.435}{K}%
\NativeMathAt{158.461}{433.435}{11.955}{\subseteq}%
\TextAt{9}{171.08}{433.435}{\Character{6}}%
\TextAt{20}{182.272}{436.457}{\Character{101}}%
\TextAt{13}{179.534}{433.435}{K}%
\TextAt{9}{193.673}{433.435}{\Character{61}}%
\TextAt{20}{207.004}{436.457}{\Character{101}}%
\TextAt{9}{206.098}{433.435}{\Character{6}}%
\NativeMathAt{217.873}{433.435}{11.955}{\subseteq}%
\TextAt{20}{231.479}{436.457}{\Character{101}}%
\TextAt{13}{230.493}{433.435}{L}%
\TextAt{8}{238.457}{433.435}{;}%
\TextAt{8}{245.606}{433.435}{thus}%
\TextAt{13}{192.08}{405.7}{e}%
\TextAt{17}{197.505}{403.84}{\Character{6}}%
\TextAt{3}{203.621}{403.84}{\Character{106}}%
\TextAt{21}{205.973}{403.84}{K}%
\TextAt{9}{217.497}{405.7}{\Character{61}}%
\TextAt{9}{229.907}{405.7}{\Character{40}}%
\TextAt{13}{234.475}{405.7}{I}%
\TextAt{21}{239.638}{403.906}{K}%
\TextAt{9}{251.161}{405.7}{\Character{58}}%
\TextAt{13}{257.734}{405.7}{I}%
\TextAt{17}{262.896}{403.906}{\Character{6}}%
\TextAt{9}{269.511}{405.7}{\Character{41}}%
\TextAt{9}{277.366}{405.7}{\Character{61}}%
\TextAt{8}{289.809}{405.7}{\char35\relax }%
\TextAt{13}{299.565}{405.7}{G}%
\TextAt{9}{308.798}{405.7}{\Character{40}}%
\TextAt{20}{314.257}{408.722}{\Character{101}}%
\TextAt{9}{313.351}{405.7}{\Character{6}}%
\TextAt{16}{325.126}{405.7}{\Character{106}}%
\TextAt{20}{334.506}{408.722}{\Character{101}}%
\TextAt{13}{331.768}{405.7}{K}%
\TextAt{9}{342.585}{405.7}{\Character{41}}%
\NativeMathAt{350.458}{405.7}{11.955}{\leq}%
\TextAt{8}{363.078}{405.7}{\char35\relax }%
\TextAt{13}{372.833}{405.7}{G}%
\TextAt{9}{382.067}{405.7}{\Character{40}}%
\TextAt{20}{387.606}{408.722}{\Character{101}}%
\TextAt{13}{386.619}{405.7}{L}%
\TextAt{16}{397.904}{405.7}{\Character{106}}%
\TextAt{20}{407.284}{408.722}{\Character{101}}%
\TextAt{13}{404.546}{405.7}{K}%
\TextAt{9}{415.364}{405.7}{\Character{41}}%
\TextAt{8}{79.226}{377.964}{is}%
\TextAt{8}{90.99}{377.964}{finite.}%
\TextAt{8}{125.453}{377.964}{Therefore}%
\TextAt{9}{178.45}{377.964}{\Character{91}\Character{6}}%
\TextAt{9}{193.454}{377.964}{\Character{58}}%
\TextAt{13}{200.049}{377.964}{K}%
\TextAt{9}{210.867}{377.964}{\Character{93}}%
\TextAt{9}{217.418}{377.964}{\Character{61}}%
\TextAt{13}{229.865}{377.964}{f}%
\TextAt{17}{235.635}{376.104}{\Character{6}}%
\TextAt{3}{241.751}{376.104}{\Character{106}}%
\TextAt{21}{244.103}{376.104}{K}%
\TextAt{13}{252.306}{377.964}{e}%
\TextAt{17}{257.731}{376.104}{\Character{6}}%
\TextAt{3}{263.848}{376.104}{\Character{106}}%
\TextAt{21}{266.2}{376.104}{K}%
\TextAt{8}{278.3}{377.964}{is}%
\TextAt{8}{290.064}{377.964}{finite}%
\TextAt{8}{319.976}{377.964}{as}%
\TextAt{8}{334.346}{377.964}{well.}%
\TextAt{8}{97.159}{357.929}{(iii)}%
\TextAt{8}{120.103}{357.929}{The}%
\TextAt{8}{144.337}{357.929}{canonical}%
\TextAt{8}{195.902}{357.929}{surjection}%
\TextAt{9}{249.797}{357.929}{\Character{0}}%
\TextAt{9}{260.724}{357.929}{\Character{61}}%
\TextAt{13}{273.475}{357.929}{G}%
\TextAt{9}{282.709}{357.929}{\Character{40}}%
\TextAt{20}{288.248}{360.951}{\Character{101}}%
\TextAt{13}{287.261}{357.929}{L}%
\TextAt{16}{298.855}{357.929}{\Character{106}}%
\TextAt{9}{305.805}{357.929}{\Character{6}\Character{41}}%
\NativeMathAt{322.441}{357.929}{11.955}{\rightarrow}%
\TextAt{13}{338.026}{357.929}{G}%
\TextAt{9}{347.259}{357.929}{\Character{40}}%
\TextAt{20}{352.718}{360.951}{\Character{101}}%
\TextAt{9}{351.812}{357.929}{\Character{6}}%
\TextAt{16}{363.895}{357.929}{\Character{106}}%
\TextAt{9}{370.846}{357.929}{\Character{6}\Character{41}}%
\NativeMathAt{387.482}{357.929}{11.955}{\cong}%
\TextAt{20}{401.074}{361.018}{\Character{98}}%
\TextAt{18}{400.41}{357.929}{Z}%
\TextAt{8}{412.458}{357.929}{has}%
\TextAt{8}{433.513}{357.929}{to}%
\TextAt{8}{448.004}{357.929}{be}%
\TextAt{8}{464.11}{357.929}{bijective.}%
\TextAt{8}{515.697}{357.929}{For}%
\TextAt{8}{79.226}{340.883}{since}%
\TextAt{9}{108.251}{340.883}{\Character{0}}%
\TextAt{9}{119.477}{340.883}{\Character{61}}%
\Vector{q 0 G 0 g 0.478 w 60.526 -422.477 m 76.907 -422.477 l S Q}%
\TextAt{16}{132.526}{340.883}{\Character{104}}%
\TextAt{9}{137.79}{340.883}{\Character{22}}%
\TextAt{13}{137.175}{340.883}{\Character{27}}%
\TextAt{16}{144.258}{340.883}{\Character{105}}%
\TextAt{8}{153.161}{340.883}{is}%
\TextAt{8}{165.281}{340.883}{procyclic,}%
\TextAt{8}{218.401}{340.883}{one}%
\TextAt{8}{240.217}{340.883}{finds}%
\TextAt{9}{268.596}{340.883}{\Character{40}\Character{0}}%
\TextAt{9}{284.365}{340.883}{\Character{58}}%
\TextAt{9}{291.538}{340.883}{\Character{0}}%
\TextAt{21}{298.884}{345.222}{n}%
\TextAt{9}{304.521}{340.883}{\Character{41}}%
\NativeMathAt{313.001}{340.883}{11.955}{\leq}%
\TextAt{13}{326.228}{340.883}{n}%
\TextAt{8}{337.469}{340.883}{for}%
\TextAt{8}{355.71}{340.883}{every}%
\TextAt{13}{386.959}{340.883}{n}%
\NativeMathAt{397.874}{340.883}{11.955}{\in}%
\TextAt{18}{409.773}{340.883}{N}%
\TextAt{8}{422.661}{340.883}{(see}%
\TextAt{8}{446.485}{340.883}{§2,}%
\TextAt{8}{465.616}{340.883}{p.273).}%
\TextAt{8}{507.018}{340.883}{Thus}%
\TextAt{8}{79.226}{323.837}{the}%
\TextAt{8}{98.962}{323.837}{induced}%
\TextAt{8}{142.106}{323.837}{maps}%
\TextAt{9}{172.314}{323.837}{\Character{0}\Character{47}\Character{0}}%
\TextAt{21}{192.796}{328.176}{n}%
\NativeMathAt{201.753}{323.837}{11.955}{\rightarrow}%
\TextAt{20}{217.693}{326.926}{\Character{98}}%
\TextAt{18}{217.029}{323.837}{Z}%
\TextAt{9}{224.999}{323.837}{\Character{47}}%
\TextAt{13}{230.852}{323.837}{n}%
\TextAt{20}{238.504}{326.926}{\Character{98}}%
\TextAt{18}{237.84}{323.837}{Z}%
\TextAt{8}{249.287}{323.837}{are}%
\TextAt{8}{268.377}{323.837}{bijective}%
\TextAt{8}{314.45}{323.837}{and}%
\TextAt{8}{336.792}{323.837}{so}%
\TextAt{8}{350.741}{323.837}{is}%
\TextAt{9}{362.085}{323.837}{\Character{0}}%
\NativeMathAt{372.72}{323.837}{11.955}{\rightarrow}%
\TextAt{20}{388.66}{326.926}{\Character{98}}%
\TextAt{18}{387.996}{323.837}{Z}%
\TextAt{8}{395.966}{323.837}{.}%
\TextAt{8}{404.274}{323.837}{But}%
\TextAt{13}{427.094}{323.837}{G}%
\TextAt{9}{436.328}{323.837}{\Character{40}}%
\TextAt{20}{441.867}{326.859}{\Character{101}}%
\TextAt{13}{440.88}{323.837}{L}%
\TextAt{16}{452.165}{323.837}{\Character{106}}%
\TextAt{9}{458.807}{323.837}{\Character{6}\Character{41}}%
\TextAt{9}{475.114}{323.837}{\Character{61}}%
\TextAt{13}{487.56}{323.837}{G}%
\TextAt{9}{496.794}{323.837}{\Character{40}}%
\TextAt{20}{502.252}{326.859}{\Character{101}}%
\TextAt{9}{501.346}{323.837}{\Character{6}}%
\TextAt{16}{513.121}{323.837}{\Character{106}}%
\TextAt{9}{519.763}{323.837}{\Character{6}\Character{41}}%
\TextAt{8}{79.226}{306.791}{implies}%
\TextAt{8}{118.953}{306.791}{that}%
\TextAt{20}{145.309}{309.813}{\Character{101}}%
\TextAt{13}{144.322}{306.791}{L}%
\TextAt{9}{155.607}{306.791}{\Character{61}}%
\TextAt{20}{168.939}{309.813}{\Character{101}}%
\TextAt{9}{168.033}{306.791}{\Character{6}}%
\TextAt{8}{176.487}{306.791}{.}%
\TextAt{8}{97.159}{286.756}{(iv)}%
\TextAt{13}{119.599}{286.756}{f}%
\TextAt{17}{125.369}{284.897}{\Character{6}}%
\TextAt{3}{131.485}{284.897}{\Character{106}}%
\TextAt{21}{133.837}{284.897}{K}%
\TextAt{13}{142.04}{286.756}{d}%
\TextAt{17}{148.122}{284.963}{\Character{6}}%
\TextAt{9}{154.737}{286.756}{\Character{40}}%
\TextAt{9}{159.89}{286.756}{\Character{22}}%
\TextAt{13}{159.289}{286.756}{\Character{27}}%
\TextAt{9}{166.371}{286.756}{\Character{41}}%
\TextAt{9}{174.226}{286.756}{\Character{61}}%
\TextAt{13}{186.67}{286.756}{d}%
\TextAt{21}{192.753}{284.963}{K}%
\TextAt{9}{200.955}{286.756}{\Character{40}}%
\TextAt{9}{206.108}{286.756}{\Character{22}}%
\TextAt{13}{205.508}{286.756}{\Character{27}}%
\TextAt{9}{212.59}{286.756}{\Character{41}}%
\TextAt{9}{220.445}{286.756}{\Character{61}}%
\TextAt{13}{232.889}{286.756}{f}%
\TextAt{17}{238.659}{284.897}{\Character{6}}%
\TextAt{3}{244.775}{284.897}{\Character{106}}%
\TextAt{21}{247.127}{284.897}{K}%
\TextAt{8}{255.33}{286.756}{;}%
\TextAt{8}{262.479}{286.756}{thus}%
\TextAt{13}{288.23}{286.756}{d}%
\TextAt{17}{294.313}{284.963}{\Character{6}}%
\TextAt{9}{300.927}{286.756}{\Character{40}}%
\TextAt{9}{306.08}{286.756}{\Character{22}}%
\TextAt{13}{305.48}{286.756}{\Character{27}}%
\TextAt{9}{312.562}{286.756}{\Character{41}}%
\TextAt{9}{320.417}{286.756}{\Character{61}}%
\TextAt{9}{332.838}{286.756}{1}%
\TextAt{8}{338.714}{286.756}{,}%
\TextAt{8}{345.863}{286.756}{and}%
\TextAt{8}{368.625}{286.756}{so}%
\TextAt{9}{383.61}{286.756}{\Character{22}}%
\TextAt{13}{382.996}{286.756}{\Character{27}}%
\TextAt{9}{393.399}{286.756}{\Character{61}}%
\TextAt{13}{405.824}{286.756}{\Character{39}}%
\TextAt{17}{413.493}{284.963}{\Character{6}}%
\TextAt{8}{420.107}{286.756}{.}%
\Vector{q 0 G 0 g 0.399 w 452.334 -487.244 m 452.334 -479.174 l S Q}%
\Vector{q 0 G 0 g 0.399 w 452.534 -479.373 m 459.707 -479.373 l S Q}%
\Vector{q 0 G 0 g 0.399 w 452.534 -487.044 m 459.707 -487.044 l S Q}%
\Vector{q 0 G 0 g 0.399 w 459.906 -487.244 m 459.906 -479.174 l S Q}%
\TextAt{8}{97.159}{258.117}{Let}%
\TextAt{8}{118.691}{258.117}{us}%
\TextAt{8}{134.27}{258.117}{illustrate}%
\TextAt{8}{184.327}{258.117}{the}%
\TextAt{8}{205.047}{258.117}{situation}%
\TextAt{8}{254.459}{258.117}{described}%
\TextAt{8}{306.776}{258.117}{in}%
\TextAt{8}{320.992}{258.117}{the}%
\TextAt{8}{341.711}{258.117}{last}%
\TextAt{8}{364.451}{258.117}{proposition}%
\TextAt{8}{426.547}{258.117}{by}%
\TextAt{8}{443.369}{258.117}{a}%
\TextAt{8}{453.688}{258.117}{diagram,}%
\TextAt{8}{503.18}{258.117}{which}%
\TextAt{8}{79.226}{241.071}{one}%
\TextAt{8}{100.686}{241.071}{should}%
\TextAt{8}{137.818}{241.071}{keep}%
\TextAt{8}{164.478}{241.071}{in}%
\TextAt{8}{178.131}{241.071}{mind}%
\TextAt{8}{208.043}{241.071}{for}%
\TextAt{8}{225.928}{241.071}{the}%
\TextAt{8}{246.084}{241.071}{sequel.}%
\TextAt{13}{231.472}{150.647}{K}%
\TextAt{13}{275.419}{207.341}{L}%
\TextAt{9}{293.599}{177.577}{\Character{6}}%
\TextAt{20}{361.772}{151.926}{\Character{101}}%
\TextAt{13}{359.033}{148.904}{K}%
\TextAt{20}{349.346}{208.619}{\Character{101}}%
\TextAt{13}{348.359}{205.597}{L}%
\TextAt{9}{359.645}{205.597}{\Character{61}}%
\TextAt{20}{372.977}{208.619}{\Character{101}}%
\TextAt{9}{372.07}{205.597}{\Character{6}}%
\NativeDiagram{%
\path[draw=black,line width=0.399bp] (174.474,-619.268) -- (282.849,-619.268);%
}%
\TextAt{13}{304.845}{161.241}{\Character{39}}%
\TextAt{21}{312.514}{159.448}{K}%
\NativeDiagram{%
\path[draw=black,line width=0.399bp] (215.567,-562.574) -- (272.175,-562.574);%
}%
\TextAt{13}{308.908}{217.934}{\Character{39}}%
\TextAt{21}{316.577}{216.141}{L}%
\NativeDiagram{%
\path[draw=black,line width=0.399bp] (207.401,-619.268) -- (207.401,-570.843);%
\path[draw=black,line width=0.399bp] (225.826,-619.268) -- (225.826,-600.607);%
\path[draw=black,line width=0.399bp] (292.442,-609.256) -- (292.442,-572.587);%
\path[draw=black,line width=0.399bp] (234.238,-588.582) -- (272.175,-571.628);%
}%
\TextAt{9}{317.205}{182.958}{\Character{22}}%
\TextAt{13}{316.59}{182.958}{\Character{27}}%
\TextAt{9}{326.994}{182.958}{\Character{61}}%
\TextAt{13}{339.419}{182.958}{\Character{39}}%
\TextAt{17}{347.088}{181.165}{\Character{6}}%
\TextAt{13}{248.516}{159.594}{f}%
\TextAt{21}{254.286}{157.735}{L}%
\TextAt{3}{260.045}{157.735}{\Character{106}}%
\TextAt{21}{262.398}{157.735}{K}%
\TextAt{13}{256.842}{142.585}{f}%
\TextAt{17}{262.612}{140.726}{\Character{6}}%
\TextAt{3}{268.728}{140.726}{\Character{106}}%
\TextAt{21}{271.081}{140.726}{K}%
\TextAt{8}{79.226}{107.013}{All}%
\TextAt{8}{98.311}{107.013}{the}%
\TextAt{8}{118.375}{107.013}{preceding}%
\TextAt{8}{170.958}{107.013}{discussions}%
\TextAt{8}{230.296}{107.013}{arose}%
\TextAt{8}{260.188}{107.013}{entirely}%
\TextAt{8}{302.369}{107.013}{from}%
\TextAt{8}{329.918}{107.013}{the}%
\TextAt{8}{349.983}{107.013}{initial}%
\TextAt{8}{383.712}{107.013}{datum}%
\TextAt{8}{420.694}{107.013}{of}%
\TextAt{8}{433.932}{107.013}{the}%
\TextAt{8}{453.997}{107.013}{homomorphism}%
\TextAt{13}{79.226}{89.967}{d}%
\TextAt{9}{89.069}{89.967}{\Character{58}}%
\TextAt{13}{96.08}{89.967}{G}%
\NativeMathAt{109.074}{89.967}{11.955}{\rightarrow}%
\TextAt{20}{125.453}{93.056}{\Character{98}}%
\TextAt{18}{124.789}{89.967}{Z}%
\TextAt{8}{132.759}{89.967}{.}%
\TextAt{8}{141.98}{89.967}{We}%
\TextAt{8}{162.382}{89.967}{now}%
\TextAt{8}{187.028}{89.967}{add}%
\TextAt{8}{210.049}{89.967}{to}%
\TextAt{8}{224.617}{89.967}{the}%
\TextAt{8}{245.031}{89.967}{data}%
\TextAt{8}{271.96}{89.967}{a}%
\TextAt{8}{281.973}{89.967}{multiplicative}%
\TextAt{13}{356.054}{89.967}{G}%
\TextAt{8}{365.288}{89.967}{-module}%
\TextAt{13}{410.736}{89.967}{A}%
\TextAt{8}{419.511}{89.967}{,}%
\TextAt{8}{426.982}{89.967}{which}%
\TextAt{8}{460.726}{89.967}{we}%
\TextAt{8}{478.211}{89.967}{equip}%
\TextAt{8}{510.007}{89.967}{with}%
\TextAt{8}{79.226}{72.921}{a}%
\TextAt{8}{88.982}{72.921}{homomorphism}%
\TextAt{8}{171.652}{72.921}{that}%
\TextAt{8}{197.02}{72.921}{is}%
\TextAt{8}{208.784}{72.921}{to}%
\TextAt{8}{223.095}{72.921}{play}%
\TextAt{8}{248.464}{72.921}{the}%
\TextAt{8}{268.62}{72.921}{role}%
\TextAt{8}{291.383}{72.921}{of}%
\TextAt{8}{304.713}{72.921}{a}%
\TextAt{8}{314.468}{72.921}{henselian}%
\TextAt{8}{365.254}{72.921}{valuation.}%
\TextAt{8}{297.211}{43.033}{263}%
\Destination{78.23}{809.815}{nav.page.263}%
\Destination{79.36}{448.455}{nav.print.288}%
\Destination{79.36}{577.155}{nav.number.4.4.5}%
\Hotlink{451.664}{338.655}{6.858}{11.047}{nav.section.4.2}%
\Hotlink{474.871}{338.655}{18.574}{11.047}{nav.print.273}%
}

% Page 276 (body page 264)
\Page{%
\TextAt{2}{79.226}{796.943}{IV.}%
\TextAt{2}{98.012}{796.943}{Abstract}%
\TextAt{2}{143.252}{796.943}{Class}%
\TextAt{2}{171.844}{796.943}{Field}%
\TextAt{2}{199.575}{796.943}{Theory}%
\Vector{q 0 G 0 g 0.299 w 7.226 18.729 m 460.769 18.729 l S Q}%
\TextAt{12}{79.226}{761.124}{Definition}%
\TextAt{12}{142.32}{761.124}{(4.6).}%
\TextAt{8}{177.901}{761.124}{A}%
\TextAt{12}{190.839}{761.124}{henselian}%
\TextAt{12}{249.641}{761.124}{valuation}%
\TextAt{8}{307.626}{761.124}{of}%
\TextAt{13}{321.222}{761.124}{A}%
\TextAt{21}{329.997}{759.33}{k}%
\TextAt{8}{339.28}{761.124}{with}%
\TextAt{8}{366.206}{761.124}{respect}%
\TextAt{8}{406.521}{761.124}{to}%
\TextAt{13}{421.097}{761.124}{d}%
\TextAt{9}{430.953}{761.124}{\Character{58}}%
\TextAt{13}{437.978}{761.124}{G}%
\NativeMathAt{450.985}{761.124}{11.955}{\rightarrow}%
\TextAt{20}{467.378}{764.212}{\Character{98}}%
\TextAt{18}{466.713}{761.124}{Z}%
\TextAt{8}{478.847}{761.124}{is}%
\TextAt{8}{490.876}{761.124}{a}%
\TextAt{8}{500.897}{761.124}{homo-}%
\TextAt{8}{79.226}{744.078}{morphism}%
\TextAt{13}{274.123}{727.032}{v}%
\TextAt{9}{283.532}{727.032}{\Character{58}}%
\TextAt{13}{290.105}{727.032}{A}%
\TextAt{21}{298.88}{725.238}{k}%
\NativeArrowAt{307.32}{727.032}{11.955}{19.26}{\longrightarrow}%
\TextAt{20}{330.567}{730.12}{\Character{98}}%
\TextAt{18}{329.902}{727.032}{Z}%
\TextAt{8}{79.226}{703.51}{satisfying}%
\TextAt{8}{131.387}{703.51}{the}%
\TextAt{8}{151.543}{703.51}{following}%
\TextAt{8}{200.978}{703.51}{properties:}%
\TextAt{8}{89.579}{683.475}{(i)}%
\TextAt{13}{107.918}{683.475}{v}%
\TextAt{9}{114.007}{683.475}{\Character{40}}%
\TextAt{13}{118.559}{683.475}{A}%
\TextAt{21}{127.334}{681.682}{k}%
\TextAt{9}{132.454}{683.475}{\Character{41}}%
\TextAt{9}{140.309}{683.475}{\Character{61}}%
\TextAt{13}{152.753}{683.475}{Z}%
\NativeMathAt{164.944}{683.475}{11.955}{\supseteq}%
\TextAt{18}{177.564}{683.475}{Z}%
\TextAt{8}{189.431}{683.475}{and}%
\TextAt{13}{212.194}{683.475}{Z}%
\TextAt{9}{221.065}{683.475}{\Character{47}}%
\TextAt{13}{226.918}{683.475}{nZ}%
\NativeMathAt{246.097}{683.475}{11.955}{\cong}%
\TextAt{18}{258.716}{683.475}{Z}%
\TextAt{9}{266.686}{683.475}{\Character{47}}%
\TextAt{13}{272.539}{683.475}{n}%
\TextAt{18}{279.527}{683.475}{Z}%
\TextAt{8}{291.395}{683.475}{for}%
\TextAt{8}{309.279}{683.475}{all}%
\TextAt{13}{325.539}{683.475}{n}%
\NativeMathAt{335.847}{683.475}{11.955}{\in}%
\TextAt{18}{347.138}{683.475}{N}%
\TextAt{8}{355.772}{683.475}{,}%
\TextAt{8}{86.327}{666.429}{(ii)}%
\TextAt{13}{107.918}{666.429}{v}%
\TextAt{9}{114.007}{666.429}{\Character{40}}%
\TextAt{13}{118.559}{666.429}{N}%
\TextAt{21}{127.933}{664.57}{K}%
\TextAt{3}{135.637}{664.57}{\Character{106}}%
\TextAt{21}{137.99}{664.57}{k}%
\TextAt{13}{143.109}{666.429}{A}%
\TextAt{21}{151.885}{664.636}{K}%
\TextAt{9}{160.087}{666.429}{\Character{41}}%
\TextAt{9}{167.942}{666.429}{\Character{61}}%
\TextAt{13}{180.386}{666.429}{f}%
\TextAt{21}{186.156}{664.636}{K}%
\TextAt{13}{194.359}{666.429}{Z}%
\TextAt{8}{207.127}{666.429}{for}%
\TextAt{8}{225.012}{666.429}{all}%
\TextAt{8}{241.271}{666.429}{finite}%
\TextAt{8}{271.183}{666.429}{extensions}%
\TextAt{13}{327.563}{666.429}{K}%
\TextAt{16}{341.702}{666.429}{\Character{106}}%
\TextAt{13}{348.344}{666.429}{k}%
\TextAt{8}{354.833}{666.429}{.}%
\TextAt{8}{97.159}{640.417}{Exactly}%
\TextAt{8}{139.648}{640.417}{like}%
\TextAt{8}{160.509}{640.417}{the}%
\TextAt{8}{180.068}{640.417}{original}%
\TextAt{8}{221.756}{640.417}{homomorphism}%
\TextAt{13}{303.828}{640.417}{d}%
\TextAt{9}{313.232}{640.417}{\Character{58}}%
\TextAt{13}{319.804}{640.417}{G}%
\TextAt{21}{329.038}{638.624}{k}%
\NativeMathAt{337.478}{640.417}{11.955}{\rightarrow}%
\TextAt{20}{353.419}{643.505}{\Character{98}}%
\TextAt{18}{352.754}{640.417}{Z}%
\TextAt{8}{360.724}{640.417}{,}%
\TextAt{8}{367.396}{640.417}{the}%
\TextAt{8}{386.954}{640.417}{henselian}%
\TextAt{8}{437.142}{640.417}{valuation}%
\TextAt{13}{487.617}{640.417}{v}%
\TextAt{9}{497.026}{640.417}{\Character{58}}%
\TextAt{13}{503.598}{640.417}{A}%
\TextAt{21}{512.374}{638.624}{k}%
\NativeMathAt{520.814}{640.417}{11.955}{\rightarrow}%
\TextAt{20}{79.89}{626.459}{\Character{98}}%
\TextAt{18}{79.226}{623.371}{Z}%
\TextAt{8}{91.094}{623.371}{has}%
\TextAt{8}{111.967}{623.371}{the}%
\TextAt{8}{132.124}{623.371}{property}%
\TextAt{8}{179.933}{623.371}{of}%
\TextAt{8}{193.263}{623.371}{reproducing}%
\TextAt{8}{257.976}{623.371}{itself}%
\TextAt{8}{286.322}{623.371}{over}%
\TextAt{8}{311.368}{623.371}{every}%
\TextAt{8}{342.26}{623.371}{finite}%
\TextAt{8}{372.172}{623.371}{extension}%
\TextAt{13}{423.938}{623.371}{K}%
\TextAt{8}{438.653}{623.371}{of}%
\TextAt{13}{451.983}{623.371}{k}%
\TextAt{8}{458.472}{623.371}{.}%
\TextAt{12}{79.226}{597.359}{Proposition}%
\TextAt{12}{151.937}{597.359}{(4.7).}%
\TextAt{8}{187.253}{597.359}{For}%
\TextAt{8}{208.222}{597.359}{every}%
\TextAt{8}{239.115}{597.359}{field}%
\TextAt{13}{264.472}{597.359}{K}%
\TextAt{8}{279.187}{597.359}{which}%
\TextAt{8}{312.673}{597.359}{is}%
\TextAt{8}{324.437}{597.359}{finite}%
\TextAt{8}{354.349}{597.359}{over}%
\TextAt{13}{379.395}{597.359}{k}%
\TextAt{8}{385.884}{597.359}{,}%
\TextAt{8}{393.033}{597.359}{the}%
\TextAt{8}{413.19}{597.359}{formula}%
\TextAt{13}{231.225}{565.614}{v}%
\TextAt{21}{236.884}{563.821}{K}%
\TextAt{9}{248.407}{565.614}{\Character{61}}%
\TextAt{9}{266.088}{573.702}{1}%
\Vector{q 0 G 0 g 0.478 w 190.028 -205.397 m 204.001 -205.397 l S Q}%
\TextAt{13}{262.028}{557.414}{f}%
\TextAt{21}{267.798}{555.62}{K}%
\TextAt{13}{277.197}{565.614}{v}%
\TextAt{16}{285.941}{565.614}{\Character{14}}%
\TextAt{13}{294.576}{565.614}{N}%
\TextAt{21}{303.949}{563.755}{K}%
\TextAt{3}{311.654}{563.755}{\Character{106}}%
\TextAt{21}{314.006}{563.755}{k}%
\TextAt{9}{322.447}{565.614}{\Character{58}}%
\TextAt{13}{329.019}{565.614}{A}%
\TextAt{21}{337.795}{563.821}{K}%
\NativeArrowAt{349.318}{565.614}{11.955}{19.26}{\longrightarrow}%
\TextAt{13}{371.9}{565.614}{Z}%
\TextAt{8}{79.226}{533.721}{defines}%
\TextAt{8}{117.65}{533.721}{a}%
\TextAt{8}{127.405}{533.721}{surjective}%
\TextAt{8}{179.817}{533.721}{homomorphism}%
\TextAt{8}{262.487}{533.721}{satisfying}%
\TextAt{8}{314.647}{533.721}{the}%
\TextAt{8}{334.804}{533.721}{following}%
\TextAt{8}{384.238}{533.721}{properties:}%
\TextAt{8}{89.579}{513.686}{(i)}%
\TextAt{13}{107.918}{513.686}{v}%
\TextAt{21}{113.578}{511.893}{K}%
\TextAt{9}{125.101}{513.686}{\Character{61}}%
\TextAt{13}{137.527}{513.686}{v}%
\TextAt{21}{143.186}{511.893}{K}%
\TextAt{22}{150.891}{514.162}{\Character{27}}%
\TextAt{16}{159.041}{513.686}{\Character{14}}%
\TextAt{13}{167.675}{513.686}{\Character{27}}%
\TextAt{8}{178.655}{513.686}{for}%
\TextAt{8}{196.54}{513.686}{all}%
\TextAt{13}{212.799}{513.686}{\Character{27}}%
\NativeMathAt{223.202}{513.686}{11.955}{\in}%
\TextAt{13}{234.493}{513.686}{G}%
\TextAt{8}{243.727}{513.686}{.}%
\TextAt{8}{86.327}{496.64}{(ii)}%
\TextAt{8}{107.918}{496.64}{For}%
\TextAt{8}{128.888}{496.64}{every}%
\TextAt{8}{159.78}{496.64}{finite}%
\TextAt{8}{189.692}{496.64}{extension}%
\TextAt{13}{241.458}{496.64}{L}%
\TextAt{16}{252.743}{496.64}{\Character{106}}%
\TextAt{13}{259.385}{496.64}{K}%
\TextAt{8}{270.202}{496.64}{,}%
\TextAt{8}{277.352}{496.64}{one}%
\TextAt{8}{298.811}{496.64}{has}%
\TextAt{8}{319.685}{496.64}{the}%
\TextAt{8}{339.841}{496.64}{commutative}%
\TextAt{8}{409.767}{496.64}{diagram}%
\TextAt{13}{287.605}{464.585}{A}%
\TextAt{21}{296.381}{462.792}{L}%
\TextAt{20}{344.898}{467.673}{\Character{98}}%
\TextAt{18}{344.234}{464.585}{Z}%
\TextAt{13}{286.632}{420.79}{A}%
\TextAt{21}{295.408}{418.997}{K}%
\TextAt{20}{343.272}{423.878}{\Character{98}}%
\TextAt{18}{342.608}{420.79}{Z}%
\TextAt{13}{350.578}{420.79}{\Character{58}}%
\NativeDiagram{%
\path[draw=black,line width=0.478bp] (235.99,-306.425) -- (266.405,-306.425);%
\path[draw=black,line width=0.478bp,line cap=round,line join=round] (264.157,-303.556) .. controls (264.612,-305.278) and (265.624,-306.09) .. (266.644,-306.425) .. controls (265.624,-306.76) and (264.612,-307.573) .. (264.157,-309.294);%
}%
\TextAt{21}{318.596}{471.73}{v}%
\TextAt{22}{322.719}{470.391}{L}%
\NativeDiagram{%
\path[draw=black,line width=0.478bp] (223.122,-315.786) -- (223.122,-339.983);%
\path[draw=black,line width=0.478bp,line cap=round,line join=round] (225.991,-337.735) .. controls (224.27,-338.189) and (223.457,-339.202) .. (223.122,-340.222) .. controls (222.787,-339.202) and (221.974,-338.189) .. (220.253,-337.735);%
}%
\TextAt{21}{271.214}{444.769}{N}%
\TextAt{22}{277.935}{443.031}{L}%
\TextAt{23}{282.994}{443.031}{\Character{106}}%
\TextAt{22}{285.263}{443.031}{K}%
\NativeDiagram{%
\path[draw=black,line width=0.478bp] (276.22,-313.993) -- (276.22,-336.429);%
\path[draw=black,line width=0.478bp,line cap=round,line join=round] (279.089,-334.182) .. controls (277.368,-334.636) and (276.555,-335.649) .. (276.22,-336.668) .. controls (275.885,-335.649) and (275.073,-334.636) .. (273.351,-334.182);%
}%
\TextAt{21}{351.036}{447.398}{f}%
\TextAt{22}{355.104}{445.66}{L}%
\TextAt{23}{360.162}{445.66}{\Character{106}}%
\TextAt{22}{362.432}{445.66}{K}%
\NativeDiagram{%
\path[draw=black,line width=0.478bp] (236.963,-350.221) -- (264.779,-350.221);%
\path[draw=black,line width=0.478bp,line cap=round,line join=round] (262.531,-347.352) .. controls (262.986,-349.073) and (263.998,-349.886) .. (265.018,-350.221) .. controls (263.998,-350.555) and (262.986,-351.368) .. (262.531,-353.09);%
}%
\TextAt{21}{317.527}{417.532}{v}%
\TextAt{22}{321.65}{416.193}{K}%
\TextAt{19}{79.226}{376.247}{Proof.}%
\TextAt{8}{116}{376.247}{(i)}%
\TextAt{8}{133.138}{376.247}{If}%
\TextAt{13}{145.709}{376.247}{\Character{28}}%
\TextAt{8}{156.904}{376.247}{runs}%
\TextAt{8}{183.857}{376.247}{through}%
\TextAt{8}{228.97}{376.247}{a}%
\TextAt{8}{239.604}{376.247}{system}%
\TextAt{8}{279.301}{376.247}{of}%
\TextAt{8}{293.51}{376.247}{representatives}%
\TextAt{8}{374.189}{376.247}{of}%
\TextAt{13}{388.398}{376.247}{G}%
\TextAt{21}{397.632}{374.454}{k}%
\TextAt{9}{402.752}{376.247}{\Character{47}}%
\TextAt{13}{408.605}{376.247}{G}%
\TextAt{21}{417.838}{374.454}{K}%
\TextAt{8}{426.041}{376.247}{,}%
\TextAt{8}{434.288}{376.247}{then}%
\TextAt{13}{461.827}{376.247}{\Character{27}}%
\TextAt{3}{468.91}{380.585}{\Character{0}}%
\TextAt{17}{475.496}{380.585}{1}%
\TextAt{13}{480.228}{376.247}{\Character{28}}%
\TextAt{13}{486.636}{376.247}{\Character{27}}%
\TextAt{8}{498.506}{376.247}{sweeps}%
\TextAt{8}{79.226}{359.201}{across}%
\TextAt{8}{113.117}{359.201}{a}%
\TextAt{8}{122.164}{359.201}{system}%
\TextAt{8}{160.275}{359.201}{of}%
\TextAt{8}{172.897}{359.201}{representatives}%
\TextAt{8}{251.99}{359.201}{of}%
\TextAt{13}{264.612}{359.201}{G}%
\TextAt{21}{273.846}{357.408}{k}%
\TextAt{9}{278.966}{359.201}{\Character{47}}%
\TextAt{13}{284.819}{359.201}{\Character{27}}%
\TextAt{3}{291.901}{363.539}{\Character{0}}%
\TextAt{17}{298.488}{363.539}{1}%
\TextAt{13}{303.22}{359.201}{G}%
\TextAt{21}{312.454}{357.408}{K}%
\TextAt{13}{320.656}{359.201}{\Character{27}}%
\TextAt{9}{331.06}{359.201}{\Character{61}}%
\TextAt{13}{343.485}{359.201}{G}%
\TextAt{21}{352.719}{357.408}{k}%
\TextAt{9}{357.838}{359.201}{\Character{47}}%
\TextAt{13}{363.691}{359.201}{G}%
\TextAt{21}{372.925}{357.408}{K}%
\TextAt{22}{380.63}{359.677}{\Character{27}}%
\TextAt{8}{386.124}{359.201}{.}%
\TextAt{8}{394.336}{359.201}{Hence}%
\TextAt{8}{428.406}{359.201}{we}%
\TextAt{8}{444.925}{359.201}{have,}%
\TextAt{8}{474.606}{359.201}{for}%
\TextAt{13}{491.783}{359.201}{a}%
\NativeMathAt{501.249}{359.201}{11.955}{\in}%
\TextAt{13}{512.54}{359.201}{A}%
\TextAt{21}{521.315}{357.408}{K}%
\TextAt{8}{529.518}{359.201}{,}%
\TextAt{13}{171.007}{322.062}{v}%
\TextAt{21}{176.666}{320.269}{K}%
\TextAt{22}{184.371}{322.538}{\Character{27}}%
\TextAt{9}{189.865}{322.062}{\Character{40}}%
\TextAt{13}{194.417}{322.062}{a}%
\TextAt{21}{200.562}{326.998}{\Character{27}}%
\TextAt{9}{206.185}{322.062}{\Character{41}}%
\TextAt{9}{214.04}{322.062}{\Character{61}}%
\TextAt{9}{234.237}{330.15}{1}%
\Vector{q 0 G 0 g 0.478 w 155.68 -448.949 m 174.648 -448.949 l S Q}%
\TextAt{13}{227.68}{313.861}{f}%
\TextAt{21}{233.45}{312.068}{K}%
\TextAt{22}{241.154}{314.337}{\Character{27}}%
\TextAt{13}{247.844}{322.062}{v}%
\TextAt{20}{255.924}{342.506}{\Character{32}}%
\TextAt{20}{265.389}{333.42}{\Character{89}}%
\TextAt{21}{270.688}{308.314}{\Character{28}}%
\TextAt{13}{282.657}{322.062}{a}%
\TextAt{21}{288.802}{326.998}{\Character{27}}%
\TextAt{21}{293.919}{326.998}{\Character{27}}%
\TextAt{23}{299.053}{329.811}{\Character{0}}%
\TextAt{24}{304.809}{329.811}{1}%
\TextAt{21}{308.96}{326.998}{\Character{28}}%
\TextAt{21}{313.631}{326.998}{\Character{27}}%
\TextAt{20}{319.262}{342.506}{\Character{33}}%
\TextAt{9}{332.047}{322.062}{\Character{61}}%
\TextAt{9}{349.728}{330.15}{1}%
\Vector{q 0 G 0 g 0.478 w 273.668 -448.949 m 287.641 -448.949 l S Q}%
\TextAt{13}{345.668}{313.861}{f}%
\TextAt{21}{351.438}{312.068}{K}%
\TextAt{13}{360.836}{322.062}{v}%
\TextAt{20}{368.917}{342.506}{\Character{32}\Character{32}}%
\TextAt{20}{387.846}{333.42}{\Character{89}}%
\TextAt{21}{393.145}{308.314}{\Character{28}}%
\TextAt{13}{405.115}{322.062}{a}%
\TextAt{21}{411.26}{326.998}{\Character{28}}%
\TextAt{20}{416.436}{342.506}{\Character{33}}%
\TextAt{21}{425.901}{339.829}{\Character{27}}%
\TextAt{20}{431.524}{342.506}{\Character{33}}%
\TextAt{9}{214.059}{287.341}{\Character{61}}%
\TextAt{9}{231.739}{295.429}{1}%
\Vector{q 0 G 0 g 0.478 w 155.68 -483.67 m 169.652 -483.67 l S Q}%
\TextAt{13}{227.68}{279.14}{f}%
\TextAt{21}{233.45}{277.347}{K}%
\TextAt{13}{242.848}{287.341}{v}%
\TextAt{9}{248.936}{287.341}{\Character{40}}%
\TextAt{13}{253.488}{287.341}{N}%
\TextAt{21}{262.862}{285.481}{K}%
\TextAt{3}{270.567}{285.481}{\Character{106}}%
\TextAt{21}{272.919}{285.481}{k}%
\TextAt{9}{278.039}{287.341}{\Character{40}}%
\TextAt{13}{282.591}{287.341}{a}%
\TextAt{9}{288.736}{287.341}{\Character{41}\Character{41}}%
\TextAt{9}{301.134}{287.341}{\Character{61}}%
\TextAt{13}{313.587}{287.341}{v}%
\TextAt{21}{319.246}{285.547}{K}%
\TextAt{9}{327.449}{287.341}{\Character{40}}%
\TextAt{13}{332.001}{287.341}{a}%
\TextAt{9}{338.146}{287.341}{\Character{41}}%
\TextAt{13}{342.698}{287.341}{\Character{58}}%
\TextAt{8}{79.226}{254.898}{(ii)}%
\TextAt{8}{98.737}{254.898}{For}%
\TextAt{13}{119.706}{254.898}{a}%
\NativeMathAt{129.172}{254.898}{11.955}{\in}%
\TextAt{13}{140.463}{254.898}{A}%
\TextAt{21}{149.238}{253.104}{L}%
\TextAt{8}{159.393}{254.898}{one}%
\TextAt{8}{180.853}{254.898}{has:}%
\TextAt{13}{178.974}{223.153}{f}%
\TextAt{21}{184.744}{221.294}{L}%
\TextAt{3}{190.503}{221.294}{\Character{106}}%
\TextAt{21}{192.855}{221.294}{K}%
\TextAt{13}{201.058}{223.153}{v}%
\TextAt{21}{206.717}{221.36}{L}%
\TextAt{9}{212.975}{223.153}{\Character{40}}%
\TextAt{13}{217.527}{223.153}{a}%
\TextAt{9}{223.672}{223.153}{\Character{41}}%
\TextAt{9}{231.527}{223.153}{\Character{61}}%
\TextAt{13}{243.97}{223.153}{f}%
\TextAt{21}{249.74}{221.294}{L}%
\TextAt{3}{255.499}{221.294}{\Character{106}}%
\TextAt{21}{257.852}{221.294}{K}%
\TextAt{9}{270.337}{231.241}{1}%
\Vector{q 0 G 0 g 0.478 w 195.25 -547.858 m 207.277 -547.858 l S Q}%
\TextAt{13}{267.25}{214.953}{f}%
\TextAt{21}{273.02}{213.159}{L}%
\TextAt{13}{280.473}{223.153}{v}%
\TextAt{9}{286.561}{223.153}{\Character{40}}%
\TextAt{13}{291.113}{223.153}{N}%
\TextAt{21}{300.487}{221.294}{L}%
\TextAt{3}{306.246}{221.294}{\Character{106}}%
\TextAt{21}{308.598}{221.294}{k}%
\TextAt{9}{313.718}{223.153}{\Character{40}}%
\TextAt{13}{318.27}{223.153}{a}%
\TextAt{9}{324.415}{223.153}{\Character{41}\Character{41}}%
\TextAt{9}{231.545}{192.851}{\Character{61}}%
\TextAt{9}{249.226}{200.939}{1}%
\Vector{q 0 G 0 g 0.478 w 173.166 -578.16 m 187.139 -578.16 l S Q}%
\TextAt{13}{245.166}{184.65}{f}%
\TextAt{21}{250.936}{182.857}{K}%
\TextAt{13}{260.334}{192.851}{v}%
\TextAt{9}{266.422}{192.851}{\Character{40}}%
\TextAt{13}{270.975}{192.851}{N}%
\TextAt{21}{280.349}{190.991}{K}%
\TextAt{3}{288.053}{190.991}{\Character{106}}%
\TextAt{21}{290.405}{190.991}{k}%
\TextAt{9}{295.525}{192.851}{\Character{40}}%
\TextAt{13}{300.077}{192.851}{N}%
\TextAt{21}{309.451}{190.991}{L}%
\TextAt{3}{315.21}{190.991}{\Character{106}}%
\TextAt{21}{317.563}{190.991}{K}%
\TextAt{9}{325.765}{192.851}{\Character{40}}%
\TextAt{13}{330.318}{192.851}{a}%
\TextAt{9}{336.463}{192.851}{\Character{41}\Character{41}\Character{41}}%
\TextAt{9}{353.404}{192.851}{\Character{61}}%
\TextAt{13}{365.866}{192.851}{v}%
\TextAt{21}{371.525}{191.058}{K}%
\TextAt{9}{379.728}{192.851}{\Character{40}}%
\TextAt{13}{384.28}{192.851}{N}%
\TextAt{21}{393.654}{190.991}{L}%
\TextAt{3}{399.413}{190.991}{\Character{106}}%
\TextAt{21}{401.765}{190.991}{K}%
\TextAt{9}{409.968}{192.851}{\Character{40}}%
\TextAt{13}{414.52}{192.851}{a}%
\TextAt{9}{420.665}{192.851}{\Character{41}\Character{41}}%
\TextAt{13}{429.77}{192.851}{\Character{58}}%
\Vector{q 0 G 0 g 0.399 w 452.334 -612.696 m 452.334 -604.626 l S Q}%
\Vector{q 0 G 0 g 0.399 w 452.534 -604.825 m 459.707 -604.825 l S Q}%
\Vector{q 0 G 0 g 0.399 w 452.534 -612.496 m 459.707 -612.496 l S Q}%
\Vector{q 0 G 0 g 0.399 w 459.906 -612.696 m 459.906 -604.626 l S Q}%
\TextAt{12}{79.226}{126.326}{Definition}%
\TextAt{12}{140.666}{126.326}{(4.8).}%
\TextAt{8}{174.809}{126.326}{A}%
\TextAt{12}{186.308}{126.326}{prime}%
\TextAt{12}{223.503}{126.326}{element}%
\TextAt{8}{271.921}{126.326}{of}%
\TextAt{13}{284.079}{126.326}{A}%
\TextAt{21}{292.854}{124.532}{K}%
\TextAt{8}{303.782}{126.326}{is}%
\TextAt{8}{314.373}{126.326}{an}%
\TextAt{8}{329.46}{126.326}{element}%
\TextAt{13}{371.53}{126.326}{\Character{25}}%
\TextAt{21}{378.17}{124.532}{K}%
\NativeMathAt{389.694}{126.326}{11.955}{\in}%
\TextAt{13}{400.984}{126.326}{A}%
\TextAt{21}{409.76}{124.532}{K}%
\TextAt{8}{420.688}{126.326}{such}%
\TextAt{8}{445.912}{126.326}{that}%
\TextAt{13}{470.109}{126.326}{v}%
\TextAt{21}{475.768}{124.532}{K}%
\TextAt{9}{483.971}{126.326}{\Character{40}}%
\TextAt{13}{488.523}{126.326}{\Character{25}}%
\TextAt{21}{495.163}{124.532}{K}%
\TextAt{9}{503.366}{126.326}{\Character{41}}%
\TextAt{9}{511.221}{126.326}{\Character{61}}%
\TextAt{9}{523.642}{126.326}{1}%
\TextAt{8}{529.518}{126.326}{.}%
\TextAt{8}{79.226}{109.28}{We}%
\TextAt{8}{99.371}{109.28}{put}%
\TextAt{13}{232.723}{92.234}{U}%
\TextAt{21}{240.673}{90.44}{K}%
\TextAt{9}{252.197}{92.234}{\Character{61}}%
\TextAt{16}{264.622}{92.234}{\Character{102}}%
\TextAt{13}{270.6}{92.234}{u}%
\NativeMathAt{280.583}{92.234}{11.955}{\in}%
\TextAt{13}{291.874}{92.234}{A}%
\TextAt{21}{300.65}{90.44}{K}%
\TextAt{16}{312.173}{92.234}{\Character{106}}%
\TextAt{13}{318.815}{92.234}{v}%
\TextAt{21}{324.474}{90.44}{K}%
\TextAt{9}{332.677}{92.234}{\Character{40}}%
\TextAt{13}{337.229}{92.234}{u}%
\TextAt{9}{343.892}{92.234}{\Character{41}}%
\TextAt{9}{351.747}{92.234}{\Character{61}}%
\TextAt{9}{364.168}{92.234}{0}%
\TextAt{16}{370.043}{92.234}{\Character{103}}%
\TextAt{13}{376.021}{92.234}{\Character{58}}%
\TextAt{8}{297.211}{43.033}{264}%
\Destination{79.23}{698.075}{nav.Item.88}%
\Destination{79.23}{680.475}{nav.Item.89}%
\Destination{79.23}{528.285}{nav.Item.90}%
\Destination{79.23}{510.685}{nav.Item.91}%
\Destination{78.23}{809.815}{nav.page.264}%
\Destination{79.36}{698.455}{nav.print.289}%
\Destination{79.36}{775.155}{nav.number.4.4.6}%
\Destination{79.36}{611.355}{nav.number.4.4.7}%
\Destination{79.36}{140.355}{nav.number.4.4.8}%
}

% Page 277 (body page 265)
\Page{%
\TextAt{2}{79.226}{796.943}{IV.}%
\TextAt{2}{98.012}{796.943}{Abstract}%
\TextAt{2}{143.252}{796.943}{Class}%
\TextAt{2}{171.844}{796.943}{Field}%
\TextAt{2}{199.575}{796.943}{Theory}%
\Vector{q 0 G 0 g 0.299 w 7.226 18.729 m 460.769 18.729 l S Q}%
\TextAt{8}{97.159}{761.124}{For}%
\TextAt{8}{117.999}{761.124}{an}%
\TextAt{8}{134.129}{761.124}{unramified}%
\TextAt{8}{192.532}{761.124}{extension}%
\TextAt{13}{244.169}{761.124}{L}%
\TextAt{16}{255.454}{761.124}{\Character{106}}%
\TextAt{13}{262.096}{761.124}{K}%
\TextAt{8}{272.913}{761.124}{,}%
\TextAt{8}{279.959}{761.124}{that}%
\TextAt{8}{305.199}{761.124}{is,}%
\TextAt{8}{320.111}{761.124}{an}%
\TextAt{8}{336.241}{761.124}{extension}%
\TextAt{8}{387.877}{761.124}{such}%
\TextAt{8}{414.145}{761.124}{that}%
\TextAt{13}{439.385}{761.124}{f}%
\TextAt{21}{445.155}{759.264}{L}%
\TextAt{3}{450.914}{759.264}{\Character{106}}%
\TextAt{21}{453.266}{759.264}{K}%
\TextAt{9}{464.789}{761.124}{\Character{61}}%
\TextAt{9}{477.199}{761.124}{\Character{91}}%
\TextAt{13}{480.467}{761.124}{L}%
\TextAt{9}{491.752}{761.124}{\Character{58}}%
\TextAt{13}{498.324}{761.124}{K}%
\TextAt{9}{509.142}{761.124}{\Character{93}}%
\TextAt{8}{512.394}{761.124}{,}%
\TextAt{8}{519.439}{761.124}{we}%
\TextAt{8}{79.226}{744.078}{have}%
\TextAt{8}{106.398}{744.078}{from}%
\TextAt{8}{134.215}{744.078}{(4.7),}%
\TextAt{8}{165.664}{744.078}{(ii)}%
\TextAt{8}{185.352}{744.078}{that}%
\TextAt{13}{210.898}{744.078}{v}%
\TextAt{21}{216.557}{742.284}{L}%
\TextAt{16}{222.814}{744.078}{\Character{106}}%
\TextAt{21}{226.135}{742.284}{A}%
\TextAt{22}{232.478}{740.945}{K}%
\TextAt{9}{243.642}{744.078}{\Character{61}}%
\TextAt{13}{256.369}{744.078}{v}%
\TextAt{21}{262.028}{742.284}{K}%
\TextAt{8}{270.231}{744.078}{.}%
\TextAt{8}{279.21}{744.078}{In}%
\TextAt{8}{294.009}{744.078}{particular,}%
\TextAt{8}{351.471}{744.078}{a}%
\TextAt{8}{361.404}{744.078}{prime}%
\TextAt{8}{394.745}{744.078}{element}%
\TextAt{8}{438.164}{744.078}{of}%
\TextAt{13}{451.671}{744.078}{A}%
\TextAt{21}{460.446}{742.284}{K}%
\TextAt{8}{472.723}{744.078}{is}%
\TextAt{8}{484.664}{744.078}{itself}%
\TextAt{8}{513.187}{744.078}{also}%
\TextAt{8}{79.226}{727.032}{a}%
\TextAt{8}{88.774}{727.032}{prime}%
\TextAt{8}{121.73}{727.032}{element}%
\TextAt{8}{164.764}{727.032}{of}%
\TextAt{13}{177.887}{727.032}{A}%
\TextAt{21}{186.662}{725.238}{L}%
\TextAt{8}{192.919}{727.032}{.}%
\TextAt{8}{201.298}{727.032}{If}%
\TextAt{8}{212.783}{727.032}{on}%
\TextAt{8}{228.834}{727.032}{the}%
\TextAt{8}{248.783}{727.032}{other}%
\TextAt{8}{279.145}{727.032}{hand,}%
\TextAt{13}{311.497}{727.032}{L}%
\TextAt{16}{322.782}{727.032}{\Character{106}}%
\TextAt{13}{329.424}{727.032}{K}%
\TextAt{8}{343.931}{727.032}{is}%
\TextAt{8}{355.488}{727.032}{totally}%
\TextAt{8}{392.688}{727.032}{ramified,}%
\TextAt{8}{441.299}{727.032}{i.e.,}%
\TextAt{13}{463.238}{727.032}{f}%
\TextAt{21}{469.008}{725.172}{L}%
\TextAt{3}{474.767}{725.172}{\Character{106}}%
\TextAt{21}{477.119}{725.172}{K}%
\TextAt{9}{488.643}{727.032}{\Character{61}}%
\TextAt{9}{501.053}{727.032}{1}%
\TextAt{8}{506.921}{727.032}{,}%
\TextAt{8}{513.904}{727.032}{and}%
\TextAt{8}{79.226}{709.986}{if}%
\TextAt{13}{89.95}{709.986}{\Character{25}}%
\TextAt{21}{96.59}{708.193}{L}%
\TextAt{8}{106.745}{709.986}{is}%
\TextAt{8}{118.509}{709.986}{a}%
\TextAt{8}{128.264}{709.986}{prime}%
\TextAt{8}{161.428}{709.986}{element}%
\TextAt{8}{204.67}{709.986}{of}%
\TextAt{13}{218}{709.986}{A}%
\TextAt{21}{226.775}{708.193}{L}%
\TextAt{8}{233.032}{709.986}{,}%
\TextAt{8}{240.181}{709.986}{then}%
\TextAt{13}{266.841}{709.986}{\Character{25}}%
\TextAt{21}{273.482}{708.193}{K}%
\TextAt{9}{285.005}{709.986}{\Character{61}}%
\TextAt{13}{297.431}{709.986}{N}%
\TextAt{21}{306.805}{708.126}{L}%
\TextAt{3}{312.564}{708.126}{\Character{106}}%
\TextAt{21}{314.916}{708.126}{K}%
\TextAt{9}{323.119}{709.986}{\Character{40}}%
\TextAt{13}{327.671}{709.986}{\Character{25}}%
\TextAt{21}{334.311}{708.193}{L}%
\TextAt{9}{340.568}{709.986}{\Character{41}}%
\TextAt{8}{349.018}{709.986}{is}%
\TextAt{8}{360.782}{709.986}{a}%
\TextAt{8}{370.537}{709.986}{prime}%
\TextAt{8}{403.701}{709.986}{element}%
\TextAt{8}{446.943}{709.986}{of}%
\TextAt{13}{460.273}{709.986}{A}%
\TextAt{21}{469.048}{708.193}{K}%
\TextAt{8}{477.251}{709.986}{.}%
\TextAt{12}{79.226}{683.973}{Exercise}%
\TextAt{12}{131.556}{683.973}{IV.4.1.}%
\TextAt{8}{174.556}{683.973}{Assume}%
\TextAt{8}{217.161}{683.973}{that}%
\TextAt{8}{241.773}{683.973}{every}%
\TextAt{8}{271.908}{683.973}{closed}%
\TextAt{8}{305.678}{683.973}{abelian}%
\TextAt{8}{345.569}{683.973}{subgroup}%
\TextAt{8}{395.61}{683.973}{of}%
\TextAt{13}{408.184}{683.973}{G}%
\TextAt{8}{420.558}{683.973}{is}%
\TextAt{8}{431.566}{683.973}{procyclic.}%
\TextAt{8}{485.287}{683.973}{Let}%
\TextAt{13}{505.5}{683.973}{K}%
\TextAt{16}{519.638}{683.973}{\Character{106}}%
\TextAt{13}{526.28}{683.973}{k}%
\TextAt{8}{79.226}{666.927}{be}%
\TextAt{8}{95.03}{666.927}{a}%
\TextAt{8}{104.666}{666.927}{finite}%
\TextAt{8}{134.457}{666.927}{extension.}%
\TextAt{8}{190.734}{666.927}{A}%
\TextAt{12}{203.287}{666.927}{microprime}%
\TextAt{14}{274.288}{666.927}{p}%
\TextAt{8}{284.048}{666.927}{of}%
\TextAt{13}{297.258}{666.927}{K}%
\TextAt{8}{311.853}{666.927}{is}%
\TextAt{8}{323.497}{666.927}{by}%
\TextAt{8}{339.636}{666.927}{definition}%
\TextAt{8}{391.545}{666.927}{a}%
\TextAt{8}{401.18}{666.927}{conjugacy}%
\TextAt{8}{455.695}{666.927}{class}%
\TextAt{16}{483.012}{666.927}{\Character{104}}%
\TextAt{13}{487.661}{666.927}{\Character{27}}%
\TextAt{16}{494.744}{666.927}{\Character{105}}%
\NativeMathAt{502.694}{666.927}{11.955}{\subseteq}%
\TextAt{13}{515.333}{666.927}{G}%
\TextAt{21}{524.567}{665.134}{K}%
\TextAt{8}{79.226}{649.881}{of}%
\TextAt{8}{93.086}{649.881}{some}%
\TextAt{8}{122.941}{649.881}{Frobenius}%
\TextAt{8}{177.34}{649.881}{element}%
\TextAt{13}{221.112}{649.881}{\Character{27}}%
\NativeMathAt{232.417}{649.881}{11.955}{\in}%
\TextAt{8}{244.61}{649.881}{Frob}%
\TextAt{9}{269.166}{649.881}{\Character{40}}%
\TextAt{9}{274.037}{653.036}{\Character{22}}%
\TextAt{13}{273.719}{649.881}{k}%
\TextAt{16}{284.431}{649.881}{\Character{106}}%
\TextAt{13}{291.975}{649.881}{K}%
\TextAt{9}{302.793}{649.881}{\Character{41}}%
\TextAt{8}{311.772}{649.881}{which}%
\TextAt{8}{345.788}{649.881}{is}%
\TextAt{8}{358.081}{649.881}{not}%
\TextAt{8}{379.424}{649.881}{a}%
\TextAt{8}{389.709}{649.881}{proper}%
\TextAt{8}{427.634}{649.881}{power}%
\TextAt{13}{462.308}{649.881}{\Character{27}}%
\TextAt{3}{469.39}{654.22}{\Character{48}}%
\TextAt{21}{471.687}{654.22}{n}%
\TextAt{8}{477.324}{649.881}{,}%
\TextAt{13}{485.135}{649.881}{n}%
\TextAt{13}{496.337}{649.881}{\Character{62}}%
\TextAt{9}{509.673}{649.881}{1}%
\TextAt{8}{515.526}{649.881}{,}%
\TextAt{8}{523.337}{649.881}{of}%
\TextAt{8}{79.226}{632.836}{some}%
\TextAt{8}{107.663}{632.836}{other}%
\TextAt{8}{137.343}{632.836}{Frobenius}%
\TextAt{8}{190.324}{632.836}{element}%
\TextAt{13}{232.676}{632.836}{\Character{27}}%
\TextAt{3}{239.758}{637.174}{\Character{48}}%
\NativeMathAt{245.874}{632.836}{11.955}{\in}%
\TextAt{8}{257.165}{632.836}{Frob}%
\TextAt{9}{281.721}{632.836}{\Character{40}}%
\TextAt{9}{286.592}{635.99}{\Character{22}}%
\TextAt{13}{286.274}{632.836}{k}%
\TextAt{16}{296.084}{632.836}{\Character{106}}%
\TextAt{13}{302.726}{632.836}{K}%
\TextAt{9}{313.543}{632.836}{\Character{41}}%
\TextAt{8}{318.096}{632.836}{.}%
\TextAt{8}{326.248}{632.836}{Let}%
\TextAt{8}{346.328}{632.836}{spec}%
\TextAt{9}{367.847}{632.836}{\Character{40}}%
\TextAt{13}{372.399}{632.836}{K}%
\TextAt{9}{383.217}{632.836}{\Character{41}}%
\TextAt{8}{390.777}{632.836}{be}%
\TextAt{8}{405.812}{632.836}{the}%
\TextAt{8}{425.079}{632.836}{set}%
\TextAt{8}{442.457}{632.836}{of}%
\TextAt{8}{454.898}{632.836}{all}%
\TextAt{8}{470.268}{632.836}{microprimes}%
\TextAt{8}{79.226}{615.79}{of}%
\TextAt{13}{92.556}{615.79}{K}%
\TextAt{8}{103.374}{615.79}{.}%
\TextAt{8}{111.822}{615.79}{Show}%
\TextAt{8}{142.714}{615.79}{that}%
\TextAt{8}{168.083}{615.79}{if}%
\TextAt{13}{178.807}{615.79}{L}%
\TextAt{16}{190.092}{615.79}{\Character{106}}%
\TextAt{13}{196.734}{615.79}{K}%
\TextAt{8}{211.449}{615.79}{is}%
\TextAt{8}{223.213}{615.79}{a}%
\TextAt{8}{232.968}{615.79}{finite}%
\TextAt{8}{262.88}{615.79}{extension,}%
\TextAt{8}{317.898}{615.79}{then}%
\TextAt{8}{344.558}{615.79}{there}%
\TextAt{8}{374.47}{615.79}{is}%
\TextAt{8}{386.234}{615.79}{a}%
\TextAt{8}{395.989}{615.79}{canonical}%
\TextAt{8}{447.372}{615.79}{mapping}%
\TextAt{13}{242.924}{586.788}{\Character{25}}%
\TextAt{9}{253.314}{586.788}{\Character{58}}%
\TextAt{8}{259.887}{586.788}{spec}%
\TextAt{9}{281.406}{586.788}{\Character{40}}%
\TextAt{13}{285.958}{586.788}{L}%
\TextAt{9}{293.923}{586.788}{\Character{41}}%
\NativeArrowAt{301.796}{586.788}{11.955}{19.26}{\longrightarrow}%
\TextAt{8}{324.378}{586.788}{spec}%
\TextAt{9}{345.897}{586.788}{\Character{40}}%
\TextAt{13}{350.45}{586.788}{K}%
\TextAt{9}{361.267}{586.788}{\Character{41}}%
\TextAt{13}{365.82}{586.788}{\Character{58}}%
\TextAt{8}{79.226}{557.787}{Above}%
\TextAt{8}{116.098}{557.787}{any}%
\TextAt{8}{138.994}{557.787}{microprime}%
\TextAt{14}{201.557}{557.787}{p}%
\TextAt{8}{212.216}{557.787}{there}%
\TextAt{8}{242.907}{557.787}{are}%
\TextAt{8}{263.197}{557.787}{only}%
\TextAt{8}{289.667}{557.787}{finitely}%
\TextAt{8}{329.79}{557.787}{many}%
\TextAt{8}{362.442}{557.787}{microprimes}%
\TextAt{14}{429.62}{557.787}{P}%
\TextAt{8}{444.189}{557.787}{of}%
\TextAt{13}{458.298}{557.787}{L}%
\TextAt{8}{466.263}{557.787}{,}%
\TextAt{8}{474.385}{557.787}{i.e.,}%
\TextAt{8}{497.464}{557.787}{the}%
\TextAt{8}{518.399}{557.787}{set}%
\TextAt{13}{79.226}{540.741}{\Character{25}}%
\TextAt{3}{86.295}{545.08}{\Character{0}}%
\TextAt{17}{92.882}{545.08}{1}%
\TextAt{9}{97.614}{540.741}{\Character{40}}%
\TextAt{14}{102.167}{540.741}{p}%
\TextAt{9}{108.15}{540.741}{\Character{41}}%
\TextAt{8}{116.599}{540.741}{is}%
\TextAt{8}{128.363}{540.741}{finite.}%
\TextAt{8}{162.826}{540.741}{We}%
\TextAt{8}{182.97}{540.741}{write}%
\TextAt{14}{212.882}{540.741}{P}%
\TextAt{16}{226.096}{540.741}{\Character{106}}%
\TextAt{14}{232.738}{540.741}{p}%
\TextAt{8}{242.619}{540.741}{to}%
\TextAt{8}{256.929}{540.741}{mean}%
\TextAt{14}{288.144}{540.741}{P}%
\NativeMathAt{301.358}{540.741}{11.955}{\in}%
\TextAt{13}{312.649}{540.741}{\Character{25}}%
\TextAt{3}{319.718}{545.08}{\Character{0}}%
\TextAt{17}{326.305}{545.08}{1}%
\TextAt{9}{331.037}{540.741}{\Character{40}}%
\TextAt{14}{335.589}{540.741}{p}%
\TextAt{9}{341.572}{540.741}{\Character{41}}%
\TextAt{8}{346.125}{540.741}{.}%
\TextAt{12}{79.226}{508.751}{Exercise}%
\TextAt{12}{132.955}{508.751}{IV.4.2.}%
\TextAt{8}{177.17}{508.751}{For}%
\TextAt{8}{198.599}{508.751}{a}%
\TextAt{8}{208.813}{508.751}{finite}%
\TextAt{8}{239.184}{508.751}{extension}%
\TextAt{13}{291.409}{508.751}{L}%
\TextAt{16}{303.477}{508.751}{\Character{106}}%
\TextAt{13}{310.901}{508.751}{K}%
\TextAt{8}{326.076}{508.751}{and}%
\TextAt{8}{349.297}{508.751}{a}%
\TextAt{8}{359.512}{508.751}{microprime}%
\TextAt{14}{421.756}{508.751}{P}%
\TextAt{16}{435.752}{508.751}{\Character{106}}%
\TextAt{14}{443.176}{508.751}{p}%
\TextAt{8}{453.516}{508.751}{of}%
\TextAt{13}{467.305}{508.751}{L}%
\TextAt{8}{475.27}{508.751}{,}%
\TextAt{8}{482.993}{508.751}{let}%
\TextAt{13}{500.357}{508.751}{f}%
\TextAt{27}{506.127}{506.892}{P}%
\TextAt{3}{512.722}{506.892}{\Character{106}}%
\TextAt{27}{515.075}{506.892}{p}%
\TextAt{9}{523.665}{508.751}{\Character{61}}%
\TextAt{13}{79.226}{491.705}{d}%
\TextAt{9}{85.309}{491.705}{\Character{40}}%
\TextAt{14}{89.861}{491.705}{P}%
\TextAt{9}{99.754}{491.705}{\Character{41}\Character{47}}%
\TextAt{13}{110.16}{491.705}{d}%
\TextAt{9}{116.242}{491.705}{\Character{40}}%
\TextAt{14}{120.795}{491.705}{p}%
\TextAt{9}{126.778}{491.705}{\Character{41}}%
\TextAt{8}{131.33}{491.705}{.}%
\TextAt{8}{139.778}{491.705}{Show}%
\TextAt{8}{170.671}{491.705}{that}%
\TextAt{20}{259.677}{486.017}{\Character{88}}%
\TextAt{27}{261.843}{460.114}{P}%
\TextAt{3}{268.438}{460.114}{\Character{106}}%
\TextAt{27}{270.79}{460.114}{p}%
\TextAt{13}{278.938}{474.659}{f}%
\TextAt{27}{284.708}{472.8}{P}%
\TextAt{3}{291.303}{472.8}{\Character{106}}%
\TextAt{27}{293.655}{472.8}{p}%
\TextAt{9}{301.463}{474.659}{\Character{61}}%
\TextAt{9}{313.873}{474.659}{\Character{91}}%
\TextAt{13}{317.14}{474.659}{L}%
\TextAt{9}{328.426}{474.659}{\Character{58}}%
\TextAt{13}{334.998}{474.659}{K}%
\TextAt{9}{345.816}{474.659}{\Character{93}}%
\TextAt{13}{349.067}{474.659}{\Character{58}}%
\TextAt{12}{79.226}{425.544}{Exercise}%
\TextAt{12}{132.427}{425.544}{IV.4.3.}%
\TextAt{8}{176.183}{425.544}{For}%
\TextAt{8}{197.152}{425.544}{an}%
\TextAt{8}{213.411}{425.544}{infinite}%
\TextAt{8}{253.078}{425.544}{extension}%
\TextAt{13}{304.844}{425.544}{L}%
\TextAt{16}{316.129}{425.544}{\Character{106}}%
\TextAt{13}{322.771}{425.544}{K}%
\TextAt{8}{333.589}{425.544}{,}%
\TextAt{8}{340.738}{425.544}{let}%
\TextAt{8}{245.816}{396.542}{spec}%
\TextAt{9}{267.335}{396.542}{\Character{40}}%
\TextAt{13}{271.888}{396.542}{L}%
\TextAt{9}{279.852}{396.542}{\Character{41}}%
\TextAt{9}{287.707}{396.542}{\Character{61}}%
\NativeMathAt{300.151}{396.542}{11.955}{\varprojlim}%
\TextAt{21}{305.561}{384.428}{\Character{11}}%
\TextAt{8}{318.402}{396.542}{spec}%
\TextAt{9}{339.922}{396.542}{\Character{40}}%
\TextAt{13}{344.474}{396.542}{L}%
\TextAt{21}{352.439}{394.749}{\Character{11}}%
\TextAt{9}{358.376}{396.542}{\Character{41}}%
\TextAt{13}{362.928}{396.542}{\Character{59}}%
\TextAt{8}{79.226}{361.315}{where}%
\TextAt{13}{113.03}{361.315}{L}%
\TextAt{21}{120.995}{359.521}{\Character{11}}%
\TextAt{16}{130.253}{361.315}{\Character{106}}%
\TextAt{13}{136.895}{361.315}{K}%
\TextAt{8}{151.605}{361.315}{varies}%
\TextAt{8}{184.513}{361.315}{over}%
\TextAt{8}{209.554}{361.315}{the}%
\TextAt{8}{229.705}{361.315}{finite}%
\TextAt{8}{259.612}{361.315}{subextensions}%
\TextAt{8}{333.933}{361.315}{of}%
\TextAt{13}{347.258}{361.315}{L}%
\TextAt{16}{358.543}{361.315}{\Character{106}}%
\TextAt{13}{365.185}{361.315}{K}%
\TextAt{8}{376.003}{361.315}{.}%
\TextAt{8}{384.449}{361.315}{What}%
\TextAt{8}{417.285}{361.315}{are}%
\TextAt{8}{436.791}{361.315}{the}%
\TextAt{8}{456.943}{361.315}{microprimes}%
\TextAt{8}{523.337}{361.315}{of}%
\TextAt{9}{79.544}{347.423}{\Character{22}}%
\TextAt{13}{79.226}{344.269}{k}%
\TextAt{8}{85.716}{344.269}{?}%
\TextAt{12}{79.226}{312.279}{Exercise}%
\TextAt{12}{132.114}{312.279}{IV.4.4.}%
\ErrataAt{175}{312.279}{360}{Let $L\mid K$ be finite Galois. Show that $G(L\mid K)$ operates}%
\ErrataAt{79.19}{295.233}{455.81}{transitively on each fibre of $\operatorname{spec}(L)\to\operatorname{spec}(K)$. The decomposition group}%
\TextAt{13}{205.111}{266.232}{G}%
\TextAt{27}{214.344}{264.438}{P}%
\TextAt{9}{221.438}{266.232}{\Character{40}}%
\TextAt{13}{225.99}{266.232}{L}%
\TextAt{16}{237.276}{266.232}{\Character{106}}%
\TextAt{13}{243.917}{266.232}{K}%
\TextAt{9}{254.735}{266.232}{\Character{41}}%
\TextAt{9}{262.59}{266.232}{\Character{61}}%
\TextAt{16}{275.034}{266.232}{\Character{102}}%
\TextAt{13}{281.011}{266.232}{\Character{27}}%
\NativeMathAt{291.415}{266.232}{11.955}{\in}%
\TextAt{13}{302.706}{266.232}{G}%
\TextAt{9}{311.939}{266.232}{\Character{40}}%
\TextAt{13}{316.491}{266.232}{L}%
\TextAt{16}{327.777}{266.232}{\Character{106}}%
\TextAt{13}{334.419}{266.232}{K}%
\TextAt{9}{345.236}{266.232}{\Character{41}}%
\TextAt{16}{353.109}{266.232}{\Character{106}}%
\TextAt{14}{359.751}{266.232}{P}%
\TextAt{21}{369.644}{271.168}{\Character{27}}%
\TextAt{9}{378.589}{266.232}{\Character{61}}%
\TextAt{14}{391.014}{266.232}{P}%
\TextAt{16}{400.907}{266.232}{\Character{103}}%
\TextAt{8}{79.226}{237.23}{is}%
\TextAt{8}{91.242}{237.23}{cyclic,}%
\TextAt{8}{126.992}{237.23}{and}%
\TextAt{8}{150.006}{237.23}{if}%
\TextAt{13}{160.982}{237.23}{Z}%
\TextAt{27}{168.994}{235.437}{P}%
\TextAt{9}{179.838}{237.23}{\Character{61}}%
\TextAt{13}{192.693}{237.23}{L}%
\TextAt{21}{200.657}{241.569}{G}%
\TextAt{28}{207.256}{240.178}{P}%
\TextAt{17}{213.453}{241.569}{\Character{40}}%
\TextAt{21}{216.746}{241.569}{L}%
\TextAt{3}{222.505}{241.569}{\Character{106}}%
\TextAt{21}{224.857}{241.569}{K}%
\TextAt{17}{232.562}{241.569}{\Character{41}}%
\TextAt{8}{240.503}{237.23}{is}%
\TextAt{8}{252.518}{237.23}{the}%
\TextAt{8}{272.927}{237.23}{“decomposition}%
\TextAt{8}{355.837}{237.23}{field”}%
\TextAt{8}{386.969}{237.23}{of}%
\TextAt{14}{400.55}{237.23}{P}%
\NativeMathAt{414.194}{237.23}{11.955}{\in}%
\TextAt{8}{425.914}{237.23}{spec}%
\TextAt{9}{447.433}{237.23}{\Character{40}}%
\TextAt{13}{451.985}{237.23}{L}%
\TextAt{9}{459.95}{237.23}{\Character{41}}%
\TextAt{8}{464.502}{237.23}{,}%
\TextAt{8}{471.966}{237.23}{then}%
\TextAt{13}{498.878}{237.23}{L}%
\TextAt{16}{510.593}{237.23}{\Character{106}}%
\TextAt{13}{517.664}{237.23}{Z}%
\TextAt{27}{525.676}{235.437}{P}%
\TextAt{8}{79.226}{220.185}{is}%
\TextAt{8}{90.99}{220.185}{unramified.}%
\TextAt{15}{79.226}{167.321}{§}%
\TextAt{15}{93.785}{167.321}{5}%
\TextAt{15}{120.693}{167.321}{The}%
\TextAt{15}{160.202}{167.321}{Reciprocity}%
\TextAt{15}{264.063}{167.321}{Map}%
\TextAt{8}{79.226}{135.435}{Continuing}%
\TextAt{8}{139.493}{135.435}{with}%
\TextAt{8}{165.927}{135.435}{the}%
\TextAt{8}{185.858}{135.435}{notation}%
\TextAt{8}{232.472}{135.435}{of}%
\TextAt{8}{245.576}{135.435}{the}%
\TextAt{8}{265.506}{135.435}{previous}%
\TextAt{8}{311.846}{135.435}{section,}%
\TextAt{8}{353.998}{135.435}{we}%
\TextAt{8}{371}{135.435}{consider}%
\TextAt{8}{416.359}{135.435}{again}%
\TextAt{8}{447.359}{135.435}{a}%
\TextAt{8}{456.889}{135.435}{profinite}%
\TextAt{8}{503.491}{135.435}{group}%
\TextAt{13}{79.226}{118.389}{G}%
\TextAt{8}{88.46}{118.389}{,}%
\TextAt{8}{95.609}{118.389}{a}%
\TextAt{8}{105.364}{118.389}{continuous}%
\TextAt{13}{163.969}{118.389}{G}%
\TextAt{8}{173.202}{118.389}{-module}%
\TextAt{13}{218.393}{118.389}{A}%
\TextAt{8}{227.168}{118.389}{,}%
\TextAt{8}{234.317}{118.389}{and}%
\TextAt{8}{257.08}{118.389}{a}%
\TextAt{8}{266.835}{118.389}{pair}%
\TextAt{8}{290.901}{118.389}{of}%
\TextAt{8}{304.231}{118.389}{homomorphisms}%
\TextAt{13}{228.379}{89.388}{d}%
\TextAt{9}{237.782}{89.388}{\Character{58}}%
\TextAt{13}{244.355}{89.388}{G}%
\NativeArrowAt{256.909}{89.388}{11.955}{19.26}{\longrightarrow}%
\TextAt{20}{280.155}{92.476}{\Character{98}}%
\TextAt{18}{279.491}{89.388}{Z}%
\TextAt{13}{287.461}{89.388}{\Character{59}}%
\TextAt{13}{316.596}{89.388}{v}%
\TextAt{9}{326.025}{89.388}{\Character{58}}%
\TextAt{13}{332.597}{89.388}{A}%
\TextAt{21}{341.373}{87.595}{k}%
\NativeArrowAt{349.813}{89.388}{11.955}{19.26}{\longrightarrow}%
\TextAt{20}{373.059}{92.476}{\Character{98}}%
\TextAt{18}{372.395}{89.388}{Z}%
\TextAt{13}{380.365}{89.388}{\Character{59}}%
\TextAt{8}{297.211}{43.033}{265}%
\Destination{78.23}{809.815}{nav.page.265}%
\Destination{79.23}{196.055}{nav.section.4.5}%
\Destination{79.226}{697.955}{nav.ex.4.4.1}%
\Destination{79.226}{522.755}{nav.ex.4.4.2}%
\Destination{79.226}{439.555}{nav.ex.4.4.3}%
\Destination{79.226}{326.255}{nav.ex.4.4.4}%
\Destination{79.36}{698.955}{nav.print.290}%
\Destination{79.23}{196.055}{outline.45}%
\Hotlink{133.715}{741.855}{25.078}{11.047}{nav.number.4.4.7}%
\Destination{79.226}{329.279}{errata.104}%
}

% Page 278 (body page 266)
\Page{%
\TextAt{2}{79.226}{796.943}{IV.}%
\TextAt{2}{98.012}{796.943}{Abstract}%
\TextAt{2}{143.252}{796.943}{Class}%
\TextAt{2}{171.844}{796.943}{Field}%
\TextAt{2}{199.575}{796.943}{Theory}%
\Vector{q 0 G 0 g 0.299 w 7.226 18.729 m 460.769 18.729 l S Q}%
\TextAt{8}{79.226}{761.124}{such}%
\TextAt{8}{106.49}{761.124}{that}%
\TextAt{13}{132.726}{761.124}{d}%
\TextAt{8}{143.573}{761.124}{is}%
\TextAt{8}{156.204}{761.124}{continuous}%
\TextAt{8}{215.676}{761.124}{and}%
\TextAt{8}{239.306}{761.124}{surjective}%
\TextAt{8}{292.584}{761.124}{and}%
\TextAt{13}{316.214}{761.124}{v}%
\TextAt{8}{327.067}{761.124}{is}%
\TextAt{8}{339.698}{761.124}{a}%
\TextAt{8}{350.32}{761.124}{henselian}%
\TextAt{8}{401.973}{761.124}{valuation}%
\TextAt{8}{453.912}{761.124}{with}%
\TextAt{8}{481.44}{761.124}{respect}%
\TextAt{8}{522.357}{761.124}{to}%
\TextAt{13}{79.226}{744.078}{d}%
\TextAt{8}{85.309}{744.078}{.}%
\TextAt{8}{96.86}{744.078}{In}%
\TextAt{8}{112.515}{744.078}{the}%
\TextAt{8}{133.706}{744.078}{following}%
\TextAt{8}{184.174}{744.078}{we}%
\TextAt{8}{202.436}{744.078}{introduce}%
\TextAt{8}{255.499}{744.078}{the}%
\TextAt{8}{276.69}{744.078}{convention}%
\TextAt{8}{336.268}{744.078}{that}%
\TextAt{8}{362.671}{744.078}{the}%
\TextAt{8}{383.862}{744.078}{letter}%
\TextAt{13}{416.111}{744.078}{K}%
\TextAt{8}{426.929}{744.078}{,}%
\TextAt{8}{435.371}{744.078}{whenever}%
\TextAt{8}{487.776}{744.078}{it}%
\TextAt{8}{500.514}{744.078}{occurs}%
\TextAt{8}{79.226}{727.032}{without}%
\TextAt{8}{124.558}{727.032}{embellishments}%
\TextAt{8}{208.052}{727.032}{or}%
\TextAt{8}{224.118}{727.032}{commentary}%
\TextAt{8}{292.87}{727.032}{to}%
\TextAt{8}{308.936}{727.032}{the}%
\TextAt{8}{330.848}{727.032}{contrary,}%
\TextAt{8}{382.155}{727.032}{will}%
\TextAt{8}{406.016}{727.032}{always}%
\TextAt{8}{445.239}{727.032}{denote}%
\TextAt{8}{484.713}{727.032}{a}%
\TextAt{8}{496.224}{727.032}{field}%
\TextAt{8}{523.337}{727.032}{of}%
\TextAt{8}{79.226}{709.986}{finite}%
\TextAt{8}{108.544}{709.986}{degree}%
\TextAt{8}{144.365}{709.986}{over}%
\TextAt{13}{168.816}{709.986}{k}%
\TextAt{8}{175.306}{709.986}{.}%
\TextAt{8}{183.556}{709.986}{We}%
\TextAt{8}{203.106}{709.986}{furthermore}%
\TextAt{8}{267.225}{709.986}{impose}%
\TextAt{8}{306.035}{709.986}{the}%
\TextAt{8}{325.597}{709.986}{following}%
\TextAt{8}{374.438}{709.986}{axiomatic}%
\TextAt{8}{427.51}{709.986}{condition,}%
\TextAt{8}{481.67}{709.986}{which}%
\TextAt{8}{514.562}{709.986}{will}%
\TextAt{8}{79.226}{692.94}{be}%
\TextAt{8}{95.15}{692.94}{systematically}%
\TextAt{8}{171.377}{692.94}{assumed}%
\TextAt{8}{218.325}{692.94}{in}%
\TextAt{8}{231.977}{692.94}{the}%
\TextAt{8}{252.134}{692.94}{sequel.}%
\TextAt{12}{79.226}{668.161}{Axiom}%
\TextAt{12}{122.659}{668.161}{(5.1).}%
\TextAt{8}{157.975}{668.161}{For}%
\TextAt{8}{178.944}{668.161}{every}%
\TextAt{8}{209.836}{668.161}{unramified}%
\TextAt{8}{268.369}{668.161}{finite}%
\TextAt{8}{298.281}{668.161}{extension}%
\TextAt{13}{350.047}{668.161}{L}%
\TextAt{16}{361.332}{668.161}{\Character{106}}%
\TextAt{13}{367.974}{668.161}{K}%
\TextAt{8}{382.689}{668.161}{one}%
\TextAt{8}{404.148}{668.161}{has}%
\TextAt{13}{204.859}{643.478}{H}%
\TextAt{21}{215.482}{648.414}{i}%
\TextAt{9}{218.863}{643.478}{\Character{40}}%
\TextAt{13}{223.415}{643.478}{G}%
\TextAt{9}{232.649}{643.478}{\Character{40}}%
\TextAt{13}{237.201}{643.478}{L}%
\TextAt{16}{248.486}{643.478}{\Character{106}}%
\TextAt{13}{255.128}{643.478}{K}%
\TextAt{9}{265.946}{643.478}{\Character{41}}%
\TextAt{13}{270.498}{643.478}{\Character{59}}%
\TextAt{13}{275.722}{643.478}{U}%
\TextAt{21}{283.693}{641.684}{L}%
\TextAt{9}{289.95}{643.478}{\Character{41}}%
\TextAt{9}{297.805}{643.478}{\Character{61}}%
\TextAt{9}{310.226}{643.478}{1}%
\TextAt{8}{340.012}{643.478}{for}%
\TextAt{13}{357.897}{643.478}{i}%
\TextAt{9}{365.211}{643.478}{\Character{61}}%
\TextAt{9}{377.62}{643.478}{0}%
\TextAt{13}{383.489}{643.478}{\Character{59}}%
\TextAt{16}{388.734}{643.478}{\Character{0}}%
\TextAt{9}{398.032}{643.478}{1}%
\TextAt{13}{403.885}{643.478}{\Character{58}}%
\TextAt{8}{97.159}{611.062}{For}%
\TextAt{8}{118.128}{611.062}{an}%
\TextAt{8}{134.387}{611.062}{infinite}%
\TextAt{8}{174.055}{611.062}{extension}%
\TextAt{13}{225.82}{611.062}{L}%
\TextAt{16}{237.106}{611.062}{\Character{106}}%
\TextAt{13}{243.747}{611.062}{K}%
\TextAt{8}{258.463}{611.062}{we}%
\TextAt{8}{275.69}{611.062}{set}%
\TextAt{13}{244.918}{586.378}{N}%
\TextAt{21}{254.292}{584.518}{L}%
\TextAt{3}{260.051}{584.518}{\Character{106}}%
\TextAt{21}{262.403}{584.518}{K}%
\TextAt{13}{270.606}{586.378}{A}%
\TextAt{21}{279.381}{584.585}{L}%
\TextAt{9}{288.959}{586.378}{\Character{61}}%
\TextAt{20}{301.384}{597.736}{\Character{92}}%
\TextAt{21}{303.535}{572.364}{M}%
\TextAt{13}{316.661}{586.378}{N}%
\TextAt{21}{326.034}{584.518}{M}%
\TextAt{3}{335.016}{584.518}{\Character{106}}%
\TextAt{21}{337.368}{584.518}{K}%
\TextAt{13}{345.571}{586.378}{A}%
\TextAt{21}{354.346}{584.585}{M}%
\TextAt{13}{363.826}{586.378}{\Character{59}}%
\TextAt{8}{79.226}{553.568}{with}%
\TextAt{13}{105.886}{553.568}{M}%
\TextAt{16}{121.781}{553.568}{\Character{106}}%
\TextAt{13}{128.422}{553.568}{K}%
\TextAt{8}{143.137}{553.568}{varying}%
\TextAt{8}{184.777}{553.568}{over}%
\TextAt{8}{209.823}{553.568}{the}%
\TextAt{8}{229.98}{553.568}{finite}%
\TextAt{8}{259.892}{553.568}{subextensions}%
\TextAt{8}{334.217}{553.568}{of}%
\TextAt{13}{347.547}{553.568}{L}%
\TextAt{16}{358.832}{553.568}{\Character{106}}%
\TextAt{13}{365.474}{553.568}{K}%
\TextAt{8}{376.292}{553.568}{.}%
\TextAt{8}{97.159}{533.534}{Our}%
\TextAt{8}{121.225}{533.534}{goal}%
\TextAt{8}{145.948}{533.534}{is}%
\TextAt{8}{157.712}{533.534}{to}%
\TextAt{8}{172.022}{533.534}{define}%
\TextAt{8}{205.831}{533.534}{a}%
\TextAt{8}{215.587}{533.534}{canonical}%
\TextAt{8}{266.97}{533.534}{homomorphism}%
\TextAt{13}{221.988}{508.85}{r}%
\TextAt{21}{227.263}{506.991}{L}%
\TextAt{3}{233.022}{506.991}{\Character{106}}%
\TextAt{21}{235.375}{506.991}{K}%
\TextAt{9}{246.898}{508.85}{\Character{58}}%
\TextAt{13}{253.471}{508.85}{G}%
\TextAt{9}{262.704}{508.85}{\Character{40}}%
\TextAt{13}{267.256}{508.85}{L}%
\TextAt{16}{278.542}{508.85}{\Character{106}}%
\TextAt{13}{285.184}{508.85}{K}%
\TextAt{9}{296.001}{508.85}{\Character{41}}%
\NativeArrowAt{303.874}{508.85}{11.955}{19.26}{\longrightarrow}%
\TextAt{13}{326.456}{508.85}{A}%
\TextAt{21}{335.232}{507.057}{K}%
\TextAt{9}{343.434}{508.85}{\Character{47}}%
\TextAt{13}{349.287}{508.85}{N}%
\TextAt{21}{358.661}{506.991}{L}%
\TextAt{3}{364.42}{506.991}{\Character{106}}%
\TextAt{21}{366.773}{506.991}{K}%
\TextAt{13}{374.975}{508.85}{A}%
\TextAt{21}{383.751}{507.057}{L}%
\TextAt{8}{79.226}{484.167}{for}%
\TextAt{8}{97.057}{484.167}{every}%
\TextAt{8}{127.896}{484.167}{finite}%
\TextAt{8}{157.754}{484.167}{Galois}%
\TextAt{8}{193.614}{484.167}{extension}%
\TextAt{13}{245.326}{484.167}{L}%
\TextAt{16}{256.611}{484.167}{\Character{106}}%
\TextAt{13}{263.253}{484.167}{K}%
\TextAt{8}{274.071}{484.167}{.}%
\TextAt{8}{282.501}{484.167}{To}%
\TextAt{8}{299.675}{484.167}{this}%
\TextAt{8}{322.444}{484.167}{end,}%
\TextAt{8}{347.758}{484.167}{we}%
\TextAt{8}{364.932}{484.167}{pass}%
\TextAt{8}{390.366}{484.167}{from}%
\TextAt{13}{417.953}{484.167}{L}%
\TextAt{16}{429.238}{484.167}{\Character{106}}%
\TextAt{13}{435.88}{484.167}{K}%
\TextAt{8}{450.542}{484.167}{to}%
\TextAt{8}{464.798}{484.167}{the}%
\TextAt{8}{484.901}{484.167}{extension}%
\TextAt{20}{80.213}{470.143}{\Character{101}}%
\TextAt{13}{79.226}{467.121}{L}%
\TextAt{16}{90.511}{467.121}{\Character{106}}%
\TextAt{13}{97.153}{467.121}{K}%
\TextAt{8}{111.868}{467.121}{and}%
\TextAt{8}{134.631}{467.121}{define}%
\TextAt{8}{168.44}{467.121}{first}%
\TextAt{8}{192.566}{467.121}{a}%
\TextAt{8}{202.321}{467.121}{mapping}%
\TextAt{8}{250.453}{467.121}{on}%
\TextAt{8}{266.712}{467.121}{the}%
\TextAt{8}{286.868}{467.121}{semigroup}%
\TextAt{8}{193.089}{442.438}{Frob}%
\TextAt{9}{217.645}{442.438}{\Character{40}}%
\TextAt{20}{223.184}{445.46}{\Character{101}}%
\TextAt{13}{222.198}{442.438}{L}%
\TextAt{16}{233.483}{442.438}{\Character{106}}%
\TextAt{13}{240.125}{442.438}{K}%
\TextAt{9}{250.942}{442.438}{\Character{41}}%
\TextAt{9}{258.797}{442.438}{\Character{61}}%
\TextAt{16}{271.241}{442.438}{\Character{102}}%
\TextAt{13}{277.219}{442.438}{\Character{27}}%
\NativeMathAt{287.622}{442.438}{11.955}{\in}%
\TextAt{13}{298.913}{442.438}{G}%
\TextAt{9}{308.147}{442.438}{\Character{40}}%
\TextAt{20}{313.685}{445.46}{\Character{101}}%
\TextAt{13}{312.699}{442.438}{L}%
\TextAt{16}{323.984}{442.438}{\Character{106}}%
\TextAt{13}{330.626}{442.438}{K}%
\TextAt{9}{341.444}{442.438}{\Character{41}}%
\TextAt{16}{349.317}{442.438}{\Character{106}}%
\TextAt{13}{355.958}{442.438}{d}%
\TextAt{21}{362.041}{440.644}{K}%
\TextAt{9}{370.244}{442.438}{\Character{40}}%
\TextAt{13}{374.796}{442.438}{\Character{27}}%
\TextAt{9}{381.879}{442.438}{\Character{41}}%
\NativeMathAt{389.752}{442.438}{11.955}{\in}%
\TextAt{18}{401.043}{442.438}{N}%
\TextAt{16}{409.677}{442.438}{\Character{103}}%
\TextAt{13}{415.655}{442.438}{\Character{58}}%
\TextAt{12}{79.226}{410.021}{Definition}%
\TextAt{12}{142.015}{410.021}{(5.2).}%
\TextAt{8}{177.33}{410.021}{The}%
\TextAt{12}{201.384}{410.021}{reciprocity}%
\TextAt{12}{268.979}{410.021}{map}%
\TextAt{13}{214.327}{385.338}{r}%
\TextAt{30}{220.364}{383.7}{\Character{101}}%
\TextAt{21}{219.602}{381.685}{L}%
\TextAt{3}{225.361}{381.685}{\Character{106}}%
\TextAt{21}{227.713}{381.685}{K}%
\TextAt{9}{239.237}{385.338}{\Character{58}}%
\TextAt{8}{245.809}{385.338}{Frob}%
\TextAt{9}{270.365}{385.338}{\Character{40}}%
\TextAt{20}{275.904}{388.36}{\Character{101}}%
\TextAt{13}{274.918}{385.338}{L}%
\TextAt{16}{286.203}{385.338}{\Character{106}}%
\TextAt{13}{292.845}{385.338}{K}%
\TextAt{9}{303.662}{385.338}{\Character{41}}%
\NativeArrowAt{311.536}{385.338}{11.955}{19.26}{\longrightarrow}%
\TextAt{13}{334.118}{385.338}{A}%
\TextAt{21}{342.893}{383.545}{K}%
\TextAt{9}{351.096}{385.338}{\Character{47}}%
\TextAt{13}{356.949}{385.338}{N}%
\TextAt{30}{367.085}{383.7}{\Character{101}}%
\TextAt{21}{366.322}{381.685}{L}%
\TextAt{3}{372.081}{381.685}{\Character{106}}%
\TextAt{21}{374.434}{381.685}{K}%
\TextAt{13}{382.636}{385.338}{A}%
\TextAt{30}{392.174}{383.7}{\Character{101}}%
\TextAt{21}{391.412}{381.685}{L}%
\TextAt{8}{79.226}{360.655}{is}%
\TextAt{8}{90.99}{360.655}{defined}%
\TextAt{8}{131.303}{360.655}{by}%
\TextAt{13}{218.669}{343.609}{r}%
\TextAt{30}{224.707}{341.97}{\Character{101}}%
\TextAt{21}{223.944}{339.956}{L}%
\TextAt{3}{229.703}{339.956}{\Character{106}}%
\TextAt{21}{232.056}{339.956}{K}%
\TextAt{9}{240.258}{343.609}{\Character{40}}%
\TextAt{13}{244.811}{343.609}{\Character{27}}%
\TextAt{9}{251.893}{343.609}{\Character{41}}%
\TextAt{9}{259.748}{343.609}{\Character{61}}%
\TextAt{13}{272.192}{343.609}{N}%
\TextAt{17}{281.566}{341.749}{\Character{6}}%
\TextAt{3}{287.682}{341.749}{\Character{106}}%
\TextAt{21}{290.034}{341.749}{K}%
\TextAt{9}{298.237}{343.609}{\Character{40}}%
\TextAt{13}{302.789}{343.609}{\Character{25}}%
\TextAt{17}{309.429}{341.815}{\Character{6}}%
\TextAt{9}{316.043}{343.609}{\Character{41}}%
\TextAt{8}{323.917}{343.609}{mod}%
\TextAt{13}{349.355}{343.609}{N}%
\TextAt{30}{359.491}{341.97}{\Character{101}}%
\TextAt{21}{358.728}{339.956}{L}%
\TextAt{3}{364.487}{339.956}{\Character{106}}%
\TextAt{21}{366.84}{339.956}{K}%
\TextAt{13}{375.042}{343.609}{A}%
\TextAt{30}{384.58}{341.97}{\Character{101}}%
\TextAt{21}{383.818}{339.956}{L}%
\TextAt{13}{390.075}{343.609}{\Character{59}}%
\TextAt{8}{79.226}{321.937}{where}%
\TextAt{9}{113.035}{321.937}{\Character{6}}%
\TextAt{8}{125.387}{321.937}{is}%
\TextAt{8}{137.151}{321.937}{the}%
\TextAt{8}{157.307}{321.937}{fixed}%
\TextAt{8}{185.593}{321.937}{field}%
\TextAt{8}{210.95}{321.937}{of}%
\TextAt{13}{224.28}{321.937}{\Character{27}}%
\TextAt{8}{235.26}{321.937}{and}%
\TextAt{13}{258.023}{321.937}{\Character{25}}%
\TextAt{17}{264.663}{320.144}{\Character{6}}%
\NativeMathAt{274.598}{321.937}{11.955}{\in}%
\TextAt{13}{285.889}{321.937}{A}%
\TextAt{17}{294.664}{320.144}{\Character{6}}%
\TextAt{8}{305.176}{321.937}{is}%
\TextAt{8}{316.94}{321.937}{a}%
\TextAt{8}{326.695}{321.937}{prime}%
\TextAt{8}{359.859}{321.937}{element.}%
\TextAt{8}{97.159}{297.159}{Observe}%
\TextAt{8}{143.663}{297.159}{that}%
\TextAt{9}{170.596}{297.159}{\Character{6}}%
\TextAt{8}{184.512}{297.159}{is}%
\TextAt{8}{197.84}{297.159}{of}%
\TextAt{8}{212.735}{297.159}{finite}%
\TextAt{8}{244.211}{297.159}{degree}%
\TextAt{8}{282.191}{297.159}{over}%
\TextAt{13}{308.802}{297.159}{K}%
\TextAt{8}{325.081}{297.159}{by}%
\TextAt{8}{342.904}{297.159}{(4.5),}%
\TextAt{8}{376.087}{297.159}{and}%
\TextAt{13}{400.414}{297.159}{\Character{27}}%
\TextAt{8}{412.958}{297.159}{becomes}%
\TextAt{8}{461.076}{297.159}{the}%
\TextAt{8}{482.797}{297.159}{Frobenius}%
\TextAt{8}{79.226}{280.113}{automorphism}%
\TextAt{13}{156.432}{280.113}{\Character{39}}%
\TextAt{17}{164.1}{278.319}{\Character{6}}%
\TextAt{8}{174.348}{280.113}{over}%
\TextAt{9}{199.13}{280.113}{\Character{6}}%
\TextAt{8}{207.584}{280.113}{.}%
\TextAt{8}{215.945}{280.113}{The}%
\TextAt{8}{239.735}{280.113}{definition}%
\TextAt{8}{291.5}{280.113}{of}%
\TextAt{13}{304.566}{280.113}{r}%
\TextAt{30}{310.604}{278.474}{\Character{101}}%
\TextAt{21}{309.841}{276.46}{L}%
\TextAt{3}{315.6}{276.46}{\Character{106}}%
\TextAt{21}{317.953}{276.46}{K}%
\TextAt{9}{326.155}{280.113}{\Character{40}}%
\TextAt{13}{330.708}{280.113}{\Character{27}}%
\TextAt{9}{337.79}{280.113}{\Character{41}}%
\TextAt{8}{345.976}{280.113}{does}%
\TextAt{8}{372.109}{280.113}{not}%
\TextAt{8}{392.659}{280.113}{depend}%
\TextAt{8}{433.031}{280.113}{on}%
\TextAt{8}{449.026}{280.113}{the}%
\TextAt{8}{468.919}{280.113}{choice}%
\TextAt{8}{503.444}{280.113}{of}%
\TextAt{8}{516.51}{280.113}{the}%
\TextAt{8}{79.226}{263.067}{element}%
\TextAt{13}{122.577}{263.067}{\Character{25}}%
\TextAt{17}{129.217}{261.274}{\Character{6}}%
\TextAt{8}{135.831}{263.067}{.}%
\TextAt{8}{144.606}{263.067}{For}%
\TextAt{8}{165.684}{263.067}{another}%
\TextAt{8}{208.724}{263.067}{one}%
\TextAt{8}{230.293}{263.067}{differs}%
\TextAt{8}{265.251}{263.067}{from}%
\TextAt{13}{293}{263.067}{\Character{25}}%
\TextAt{17}{299.64}{261.274}{\Character{6}}%
\TextAt{8}{310.261}{263.067}{only}%
\TextAt{8}{336.061}{263.067}{by}%
\TextAt{8}{352.429}{263.067}{an}%
\TextAt{8}{368.797}{263.067}{element}%
\TextAt{13}{412.148}{263.067}{u}%
\NativeMathAt{422.317}{263.067}{11.955}{\in}%
\TextAt{13}{433.793}{263.067}{U}%
\TextAt{17}{441.743}{261.274}{\Character{6}}%
\TextAt{8}{448.358}{263.067}{,}%
\TextAt{8}{455.643}{263.067}{and}%
\TextAt{8}{478.514}{263.067}{for}%
\TextAt{8}{496.508}{263.067}{this}%
\TextAt{8}{519.439}{263.067}{we}%
\TextAt{8}{79.226}{246.021}{have}%
\TextAt{13}{105.609}{246.021}{N}%
\TextAt{17}{114.983}{244.161}{\Character{6}}%
\TextAt{3}{121.099}{244.161}{\Character{106}}%
\TextAt{21}{123.451}{244.161}{K}%
\TextAt{9}{131.654}{246.021}{\Character{40}}%
\TextAt{13}{136.206}{246.021}{u}%
\TextAt{9}{142.868}{246.021}{\Character{41}}%
\NativeMathAt{150.742}{246.021}{11.955}{\in}%
\TextAt{13}{162.033}{246.021}{N}%
\TextAt{30}{172.169}{244.382}{\Character{101}}%
\TextAt{21}{171.406}{242.368}{L}%
\TextAt{3}{177.165}{242.368}{\Character{106}}%
\TextAt{21}{179.518}{242.368}{K}%
\TextAt{13}{187.72}{246.021}{A}%
\TextAt{30}{197.258}{244.382}{\Character{101}}%
\TextAt{21}{196.496}{242.368}{L}%
\TextAt{8}{202.753}{246.021}{,}%
\TextAt{8}{209.413}{246.021}{so}%
\TextAt{8}{223.171}{246.021}{that}%
\TextAt{13}{247.927}{246.021}{N}%
\TextAt{17}{257.301}{244.161}{\Character{6}}%
\TextAt{3}{263.417}{244.161}{\Character{106}}%
\TextAt{21}{265.77}{244.161}{K}%
\TextAt{9}{273.972}{246.021}{\Character{40}}%
\TextAt{13}{278.525}{246.021}{u}%
\TextAt{9}{285.187}{246.021}{\Character{41}}%
\NativeMathAt{293.06}{246.021}{11.955}{\in}%
\TextAt{13}{304.351}{246.021}{N}%
\TextAt{21}{313.725}{244.161}{M}%
\TextAt{3}{322.707}{244.161}{\Character{106}}%
\TextAt{21}{325.059}{244.161}{K}%
\TextAt{13}{333.262}{246.021}{A}%
\TextAt{21}{342.037}{244.228}{M}%
\TextAt{8}{354.802}{246.021}{for}%
\TextAt{8}{372.075}{246.021}{every}%
\TextAt{8}{402.355}{246.021}{finite}%
\TextAt{8}{431.655}{246.021}{Galois}%
\TextAt{8}{466.956}{246.021}{subextension}%
\TextAt{13}{79.226}{227.723}{M}%
\TextAt{16}{95.298}{227.723}{\Character{106}}%
\TextAt{13}{102.116}{227.723}{K}%
\TextAt{8}{116.935}{227.723}{of}%
\TextAt{20}{131.356}{230.745}{\Character{101}}%
\TextAt{13}{130.369}{227.723}{L}%
\TextAt{16}{141.832}{227.723}{\Character{106}}%
\TextAt{13}{148.65}{227.723}{K}%
\TextAt{8}{159.468}{227.723}{.}%
\TextAt{8}{168.228}{227.723}{To}%
\TextAt{8}{185.559}{227.723}{see}%
\TextAt{8}{204.576}{227.723}{this,}%
\TextAt{8}{230.781}{227.723}{we}%
\TextAt{8}{248.112}{227.723}{may}%
\TextAt{8}{273.585}{227.723}{clearly}%
\TextAt{8}{311.084}{227.723}{assume}%
\TextAt{8}{351.633}{227.723}{that}%
\TextAt{9}{377.105}{227.723}{\Character{6}}%
\NativeMathAt{389.058}{227.723}{11.955}{\subseteq}%
\TextAt{13}{401.854}{227.723}{M}%
\TextAt{8}{414.428}{227.723}{.}%
\TextAt{8}{423.188}{227.723}{Applying}%
\TextAt{8}{474.017}{227.723}{(5.1)}%
\TextAt{8}{502.096}{227.723}{to}%
\TextAt{8}{516.51}{227.723}{the}%
\TextAt{8}{79.226}{210.677}{unramified}%
\TextAt{8}{137.759}{210.677}{extension}%
\TextAt{13}{189.525}{210.677}{M}%
\TextAt{16}{205.419}{210.677}{\Character{106}}%
\TextAt{9}{212.061}{210.677}{\Character{6}}%
\TextAt{8}{220.515}{210.677}{,}%
\TextAt{8}{227.664}{210.677}{one}%
\TextAt{8}{249.124}{210.677}{finds}%
\TextAt{13}{277.147}{210.677}{u}%
\TextAt{9}{287.13}{210.677}{\Character{61}}%
\TextAt{13}{299.555}{210.677}{N}%
\TextAt{21}{308.929}{208.817}{M}%
\TextAt{3}{317.911}{208.817}{\Character{106}}%
\TextAt{17}{320.263}{208.817}{\Character{6}}%
\TextAt{9}{326.877}{210.677}{\Character{40}}%
\TextAt{13}{331.43}{210.677}{\Character{34}}%
\TextAt{9}{336.908}{210.677}{\Character{41}}%
\TextAt{8}{341.46}{210.677}{,}%
\TextAt{13}{348.609}{210.677}{\Character{34}}%
\NativeMathAt{357.408}{210.677}{11.955}{\in}%
\TextAt{13}{368.699}{210.677}{U}%
\TextAt{21}{376.649}{208.883}{M}%
\TextAt{8}{386.129}{210.677}{,}%
\TextAt{8}{393.278}{210.677}{and}%
\TextAt{8}{416.041}{210.677}{thus}%
\TextAt{13}{176.557}{185.993}{N}%
\TextAt{17}{185.931}{184.134}{\Character{6}}%
\TextAt{3}{192.047}{184.134}{\Character{106}}%
\TextAt{21}{194.399}{184.134}{K}%
\TextAt{9}{202.602}{185.993}{\Character{40}}%
\TextAt{13}{207.154}{185.993}{u}%
\TextAt{9}{213.817}{185.993}{\Character{41}}%
\TextAt{9}{221.672}{185.993}{\Character{61}}%
\TextAt{13}{234.116}{185.993}{N}%
\TextAt{17}{243.489}{184.134}{\Character{6}}%
\TextAt{3}{249.605}{184.134}{\Character{106}}%
\TextAt{21}{251.958}{184.134}{K}%
\TextAt{9}{260.16}{185.993}{\Character{40}}%
\TextAt{13}{264.713}{185.993}{N}%
\TextAt{21}{274.087}{184.134}{M}%
\TextAt{3}{283.068}{184.134}{\Character{106}}%
\TextAt{17}{285.421}{184.134}{\Character{6}}%
\TextAt{9}{292.035}{185.993}{\Character{40}}%
\TextAt{13}{296.587}{185.993}{\Character{34}}%
\TextAt{9}{302.065}{185.993}{\Character{41}\Character{41}}%
\TextAt{9}{314.463}{185.993}{\Character{61}}%
\TextAt{13}{326.916}{185.993}{N}%
\TextAt{21}{336.29}{184.134}{M}%
\TextAt{3}{345.272}{184.134}{\Character{106}}%
\TextAt{21}{347.624}{184.134}{K}%
\TextAt{9}{355.827}{185.993}{\Character{40}}%
\TextAt{13}{360.379}{185.993}{\Character{34}}%
\TextAt{9}{365.857}{185.993}{\Character{41}}%
\NativeMathAt{373.73}{185.993}{11.955}{\in}%
\TextAt{13}{385.021}{185.993}{N}%
\TextAt{21}{394.395}{184.134}{M}%
\TextAt{3}{403.377}{184.134}{\Character{106}}%
\TextAt{21}{405.729}{184.134}{K}%
\TextAt{13}{413.932}{185.993}{A}%
\TextAt{21}{422.707}{184.2}{M}%
\TextAt{13}{432.187}{185.993}{\Character{58}}%
\TextAt{8}{79.226}{161.31}{Next}%
\TextAt{8}{108.85}{161.31}{we}%
\TextAt{8}{127.092}{161.31}{want}%
\TextAt{8}{156.728}{161.31}{to}%
\TextAt{8}{172.053}{161.31}{show}%
\TextAt{8}{202.071}{161.31}{that}%
\TextAt{8}{228.455}{161.31}{the}%
\TextAt{8}{249.627}{161.31}{reciprocity}%
\TextAt{8}{308.851}{161.31}{map}%
\TextAt{13}{335.881}{161.31}{r}%
\TextAt{30}{341.918}{159.672}{\Character{101}}%
\TextAt{21}{341.156}{157.657}{L}%
\TextAt{3}{346.915}{157.657}{\Character{106}}%
\TextAt{21}{349.267}{157.657}{K}%
\TextAt{8}{362.382}{161.31}{is}%
\TextAt{8}{375.161}{161.31}{multiplicative.}%
\TextAt{8}{456.58}{161.31}{To}%
\TextAt{8}{474.823}{161.31}{do}%
\TextAt{8}{492.097}{161.31}{this,}%
\TextAt{8}{519.439}{161.31}{we}%
\TextAt{8}{79.226}{143.012}{consider}%
\TextAt{8}{124.811}{143.012}{for}%
\TextAt{8}{142.696}{143.012}{every}%
\TextAt{13}{173.588}{143.012}{\Character{27}}%
\NativeMathAt{183.992}{143.012}{11.955}{\in}%
\TextAt{13}{195.283}{143.012}{G}%
\TextAt{9}{204.516}{143.012}{\Character{40}}%
\TextAt{20}{210.055}{146.034}{\Character{101}}%
\TextAt{13}{209.068}{143.012}{L}%
\TextAt{16}{220.354}{143.012}{\Character{106}}%
\TextAt{13}{226.996}{143.012}{K}%
\TextAt{9}{237.813}{143.012}{\Character{41}}%
\TextAt{8}{246.263}{143.012}{and}%
\TextAt{8}{269.026}{143.012}{every}%
\TextAt{13}{299.918}{143.012}{n}%
\NativeMathAt{310.226}{143.012}{11.955}{\in}%
\TextAt{18}{321.517}{143.012}{N}%
\TextAt{8}{334.049}{143.012}{the}%
\TextAt{8}{354.205}{143.012}{endomorphisms}%
\TextAt{13}{164.684}{118.328}{\Character{27}}%
\TextAt{16}{174.423}{118.329}{\Character{0}}%
\TextAt{9}{186.379}{118.328}{1}%
\TextAt{9}{195.537}{118.328}{\Character{58}}%
\TextAt{13}{202.125}{118.329}{A}%
\TextAt{30}{211.663}{116.69}{\Character{101}}%
\TextAt{21}{210.9}{114.676}{L}%
\NativeArrowAt{220.478}{118.328}{11.955}{19.26}{\longrightarrow}%
\TextAt{13}{243.06}{118.328}{A}%
\TextAt{30}{252.598}{116.69}{\Character{101}}%
\TextAt{21}{251.836}{114.676}{L}%
\TextAt{13}{258.093}{118.328}{\Character{59}}%
\TextAt{13}{346.803}{118.328}{a}%
\NativeArrowAt{356.268}{118.328}{11.955}{19.26}{\longmapsto}%
\TextAt{13}{378.85}{118.328}{a}%
\TextAt{21}{384.995}{123.265}{\Character{27}}%
\TextAt{3}{390.121}{123.265}{\Character{0}}%
\TextAt{17}{396.707}{123.265}{1}%
\TextAt{9}{404.76}{118.328}{\Character{61}}%
\TextAt{13}{417.186}{118.328}{a}%
\TextAt{21}{423.331}{123.265}{\Character{27}}%
\TextAt{9}{428.954}{118.328}{\Character{47}}%
\TextAt{13}{434.807}{118.328}{a\Character{59}}%
\TextAt{13}{179.942}{87.954}{\Character{27}}%
\TextAt{21}{186.595}{86.16}{n}%
\TextAt{9}{195.552}{87.954}{\Character{58}}%
\TextAt{13}{202.125}{87.954}{A}%
\TextAt{30}{211.663}{86.315}{\Character{101}}%
\TextAt{21}{210.9}{84.301}{L}%
\NativeArrowAt{220.478}{87.954}{11.955}{19.26}{\longrightarrow}%
\TextAt{13}{243.06}{87.954}{A}%
\TextAt{30}{252.598}{86.315}{\Character{101}}%
\TextAt{21}{251.836}{84.301}{L}%
\TextAt{13}{258.093}{87.954}{\Character{59}}%
\TextAt{13}{346.803}{87.954}{a}%
\NativeArrowAt{356.268}{87.954}{11.955}{19.26}{\longmapsto}%
\TextAt{13}{378.85}{87.954}{a}%
\TextAt{21}{384.995}{92.89}{\Character{27}}%
\TextAt{22}{389.835}{91.894}{n}%
\TextAt{9}{398.753}{87.954}{\Character{61}}%
\TextAt{21}{411.178}{102.898}{n}%
\TextAt{3}{416.316}{102.898}{\Character{0}}%
\TextAt{17}{422.903}{102.898}{1}%
\TextAt{20}{411.519}{99.311}{\Character{89}}%
\TextAt{21}{412.306}{74.117}{i}%
\TextAt{17}{415.189}{74.117}{\Character{61}0}%
\TextAt{13}{429.129}{87.954}{a}%
\TextAt{21}{435.274}{92.89}{\Character{27}}%
\TextAt{22}{440.4}{95.703}{i}%
\TextAt{13}{444.06}{87.954}{\Character{58}}%
\TextAt{8}{297.211}{43.033}{266}%
\Destination{78.23}{809.815}{nav.page.266}%
\Destination{79.36}{568.555}{nav.print.291}%
\Destination{79.36}{424.055}{nav.number.4.5.2}%
\Hotlink{342.404}{294.955}{25.078}{11.047}{nav.number.4.4.5}%
\Destination{79.226}{680}{nav.number.4.5.1}%
}

% Page 279 (body page 267)
\Page{%
\TextAt{2}{79.226}{796.943}{IV.}%
\TextAt{2}{98.012}{796.943}{Abstract}%
\TextAt{2}{143.252}{796.943}{Class}%
\TextAt{2}{171.844}{796.943}{Field}%
\TextAt{2}{199.575}{796.943}{Theory}%
\Vector{q 0 G 0 g 0.299 w 7.226 18.729 m 460.769 18.729 l S Q}%
\TextAt{8}{79.226}{761.124}{In}%
\TextAt{8}{93.847}{761.124}{formal}%
\TextAt{8}{130.598}{761.124}{notation,}%
\TextAt{8}{180.69}{761.124}{this}%
\TextAt{8}{203.512}{761.124}{gives}%
\TextAt{13}{232.193}{761.124}{\Character{27}}%
\TextAt{21}{238.846}{759.33}{n}%
\TextAt{9}{247.803}{761.124}{\Character{61}}%
\TextAt{9}{260.212}{761.124}{\Character{40}}%
\TextAt{13}{264.781}{761.124}{\Character{27}}%
\TextAt{21}{271.863}{765.462}{n}%
\TextAt{16}{280.156}{761.124}{\Character{0}}%
\TextAt{9}{292.112}{761.124}{1\Character{41}\Character{47}\Character{40}}%
\TextAt{13}{312.922}{761.124}{\Character{27}}%
\TextAt{16}{322.661}{761.124}{\Character{0}}%
\TextAt{9}{334.616}{761.124}{1\Character{41}}%
\TextAt{8}{345.022}{761.124}{,}%
\TextAt{8}{352.171}{761.124}{and}%
\TextAt{8}{374.934}{761.124}{we}%
\TextAt{8}{392.161}{761.124}{find}%
\TextAt{8}{415.569}{761.124}{that}%
\TextAt{9}{211.801}{732.732}{\Character{40}}%
\TextAt{13}{216.354}{732.732}{\Character{27}}%
\TextAt{16}{226.093}{732.732}{\Character{0}}%
\TextAt{9}{238.048}{732.732}{1\Character{41}}%
\TextAt{16}{251.11}{732.732}{\Character{14}}%
\TextAt{13}{259.744}{732.732}{\Character{27}}%
\TextAt{21}{266.398}{730.938}{n}%
\TextAt{9}{275.355}{732.732}{\Character{61}}%
\TextAt{13}{287.78}{732.732}{\Character{27}}%
\TextAt{21}{294.434}{730.938}{n}%
\TextAt{16}{302.727}{732.732}{\Character{14}}%
\TextAt{9}{311.361}{732.732}{\Character{40}}%
\TextAt{13}{315.913}{732.732}{\Character{27}}%
\TextAt{16}{325.652}{732.732}{\Character{0}}%
\TextAt{9}{337.608}{732.732}{1\Character{41}}%
\TextAt{9}{351.309}{732.732}{\Character{61}}%
\TextAt{13}{363.759}{732.732}{\Character{27}}%
\TextAt{21}{370.842}{737.668}{n}%
\TextAt{16}{379.135}{732.732}{\Character{0}}%
\TextAt{9}{391.09}{732.732}{1}%
\TextAt{13}{396.943}{732.732}{\Character{58}}%
\TextAt{8}{79.226}{704.34}{Now}%
\TextAt{8}{105.886}{704.34}{we}%
\TextAt{8}{123.114}{704.34}{introduce}%
\TextAt{8}{175.142}{704.34}{the}%
\TextAt{8}{195.299}{704.34}{homomorphism}%
\TextAt{13}{246.066}{675.947}{N}%
\TextAt{9}{260.01}{675.947}{\Character{61}}%
\TextAt{13}{272.435}{675.947}{N}%
\TextAt{30}{282.571}{674.309}{\Character{101}}%
\TextAt{21}{281.809}{672.295}{L}%
\TextAt{3}{287.568}{672.294}{\Character{106}}%
\TextAt{30}{291.891}{674.309}{\Character{101}}%
\TextAt{21}{289.92}{672.295}{K}%
\TextAt{9}{301.444}{675.947}{\Character{58}}%
\TextAt{13}{308.016}{675.947}{A}%
\TextAt{30}{317.554}{674.309}{\Character{101}}%
\TextAt{21}{316.792}{672.295}{L}%
\NativeArrowAt{326.37}{675.947}{11.955}{19.26}{\longrightarrow}%
\TextAt{13}{348.952}{675.947}{A}%
\TextAt{30}{359.697}{674.309}{\Character{101}}%
\TextAt{21}{357.727}{672.295}{K}%
\TextAt{8}{79.226}{647.555}{and}%
\TextAt{8}{101.989}{647.555}{prove}%
\TextAt{8}{133.538}{647.555}{two}%
\TextAt{8}{155.656}{647.555}{lemmas}%
\TextAt{8}{197.989}{647.555}{for}%
\TextAt{8}{215.874}{647.555}{it.}%
\TextAt{12}{79.226}{621.717}{Lemma}%
\TextAt{12}{126.367}{621.717}{(5.3).}%
\TextAt{8}{161.226}{621.717}{Let}%
\TextAt{13}{181.739}{621.717}{\Character{39}\Character{59}}%
\TextAt{13}{194.627}{621.717}{\Character{27}}%
\NativeMathAt{205.055}{621.717}{11.955}{\in}%
\TextAt{8}{216.346}{621.717}{Frob}%
\TextAt{9}{240.902}{621.717}{\Character{40}}%
\TextAt{20}{246.441}{624.739}{\Character{101}}%
\TextAt{13}{245.454}{621.717}{L}%
\TextAt{16}{256.74}{621.717}{\Character{106}}%
\TextAt{13}{263.381}{621.717}{K}%
\TextAt{9}{274.199}{621.717}{\Character{41}}%
\TextAt{8}{282.192}{621.717}{with}%
\TextAt{13}{308.396}{621.717}{d}%
\TextAt{21}{314.479}{619.924}{K}%
\TextAt{9}{322.682}{621.717}{\Character{40}}%
\TextAt{13}{327.234}{621.717}{\Character{39}}%
\TextAt{9}{334.902}{621.717}{\Character{41}}%
\TextAt{9}{342.757}{621.717}{\Character{61}}%
\TextAt{9}{355.178}{621.717}{1}%
\TextAt{8}{361.054}{621.717}{,}%
\TextAt{13}{367.838}{621.717}{d}%
\TextAt{21}{373.921}{619.924}{K}%
\TextAt{9}{382.124}{621.717}{\Character{40}}%
\TextAt{13}{386.676}{621.717}{\Character{27}}%
\TextAt{9}{393.758}{621.717}{\Character{41}}%
\TextAt{9}{401.613}{621.717}{\Character{61}}%
\TextAt{13}{414.057}{621.717}{n}%
\TextAt{8}{421.045}{621.717}{.}%
\TextAt{8}{429.341}{621.717}{If}%
\TextAt{9}{440.577}{621.717}{\Character{6}}%
\TextAt{8}{452.472}{621.717}{is}%
\TextAt{8}{463.78}{621.717}{the}%
\TextAt{8}{483.48}{621.717}{fixed}%
\TextAt{8}{511.31}{621.717}{field}%
\TextAt{8}{79.226}{604.671}{of}%
\TextAt{13}{92.556}{604.671}{\Character{27}}%
\TextAt{8}{103.536}{604.671}{and}%
\TextAt{13}{126.299}{604.671}{a}%
\NativeMathAt{135.764}{604.671}{11.955}{\in}%
\TextAt{13}{147.055}{604.671}{A}%
\TextAt{17}{155.831}{602.878}{\Character{6}}%
\TextAt{8}{162.445}{604.671}{,}%
\TextAt{8}{169.594}{604.671}{then}%
\TextAt{13}{208.406}{576.279}{N}%
\TextAt{17}{217.78}{574.419}{\Character{6}}%
\TextAt{3}{223.896}{574.419}{\Character{106}}%
\TextAt{21}{226.248}{574.419}{K}%
\TextAt{9}{234.451}{576.279}{\Character{40}}%
\TextAt{13}{239.003}{576.279}{a}%
\TextAt{9}{245.148}{576.279}{\Character{41}}%
\TextAt{9}{253.003}{576.279}{\Character{61}}%
\TextAt{9}{265.424}{576.279}{\Character{40}}%
\TextAt{13}{269.999}{576.279}{N}%
\TextAt{16}{283.278}{576.279}{\Character{14}}%
\TextAt{13}{291.913}{576.279}{\Character{39}}%
\TextAt{21}{299.581}{574.486}{n}%
\TextAt{9}{305.217}{576.279}{\Character{41}\Character{40}}%
\TextAt{13}{314.322}{576.279}{a}%
\TextAt{9}{320.467}{576.279}{\Character{41}}%
\TextAt{9}{328.322}{576.279}{\Character{61}}%
\TextAt{9}{340.743}{576.279}{\Character{40}}%
\TextAt{13}{345.318}{576.279}{\Character{39}}%
\TextAt{21}{352.986}{574.486}{n}%
\TextAt{16}{361.279}{576.279}{\Character{14}}%
\TextAt{13}{369.913}{576.279}{N}%
\TextAt{9}{380.536}{576.279}{\Character{41}\Character{40}}%
\TextAt{13}{389.641}{576.279}{a}%
\TextAt{9}{395.786}{576.279}{\Character{41}}%
\TextAt{13}{400.338}{576.279}{\Character{58}}%
\TextAt{19}{79.226}{530.302}{Proof.}%
\TextAt{8}{116}{530.302}{The}%
\TextAt{8}{141.121}{530.302}{maximal}%
\TextAt{8}{189.997}{530.302}{unramified}%
\TextAt{8}{249.597}{530.302}{subextension}%
\TextAt{9}{320.374}{530.302}{\Character{6}}%
\TextAt{17}{328.829}{534.641}{0}%
\TextAt{9}{338.7}{530.302}{\Character{61}}%
\TextAt{9}{352.927}{530.302}{\Character{6}}%
\TextAt{16}{364.783}{530.302}{\Character{92}}%
\TextAt{20}{378.875}{533.324}{\Character{101}}%
\TextAt{13}{376.137}{530.302}{K}%
\TextAt{16}{392.094}{530.302}{\Character{106}}%
\TextAt{13}{400.554}{530.302}{K}%
\TextAt{8}{416.337}{530.302}{is}%
\TextAt{8}{429.168}{530.302}{of}%
\TextAt{8}{443.565}{530.302}{degree}%
\TextAt{13}{481.048}{530.302}{n}%
\TextAt{8}{488.035}{530.302}{,}%
\TextAt{8}{496.518}{530.302}{and}%
\TextAt{8}{520.348}{530.302}{its}%
\TextAt{8}{79.226}{513.256}{Galois}%
\TextAt{8}{115.641}{513.256}{group}%
\TextAt{13}{149.318}{513.256}{G}%
\TextAt{9}{158.551}{513.256}{\Character{40}\Character{6}}%
\TextAt{17}{171.558}{517.595}{0}%
\TextAt{16}{180.465}{513.256}{\Character{106}}%
\TextAt{13}{187.961}{513.256}{K}%
\TextAt{9}{198.779}{513.256}{\Character{41}}%
\TextAt{8}{207.73}{513.256}{is}%
\TextAt{8}{219.995}{513.256}{generated}%
\TextAt{8}{273.828}{513.256}{by}%
\TextAt{8}{290.588}{513.256}{the}%
\TextAt{8}{311.246}{513.256}{Frobenius}%
\TextAt{8}{365.617}{513.256}{automorphism}%
\TextAt{13}{443.588}{513.256}{\Character{39}}%
\TextAt{17}{451.256}{511.253}{\Character{6}}%
\TextAt{24}{457.372}{513.522}{0}%
\TextAt{3}{461.523}{511.253}{\Character{106}}%
\TextAt{21}{463.875}{511.253}{K}%
\TextAt{9}{476.253}{513.256}{\Character{61}}%
\TextAt{13}{489.533}{513.256}{\Character{39}}%
\TextAt{21}{497.201}{511.463}{K}%
\TextAt{16}{505.404}{513.256}{\Character{106}}%
\TextAt{17}{508.725}{511.253}{\Character{6}}%
\TextAt{24}{514.841}{513.522}{0}%
\TextAt{9}{523.665}{513.256}{\Character{61}}%
\TextAt{13}{79.226}{496.21}{\Character{39}}%
\TextAt{16}{86.895}{496.21}{\Character{106}}%
\TextAt{30}{92.186}{494.572}{\Character{101}}%
\TextAt{21}{90.215}{492.557}{K}%
\TextAt{16}{98.418}{496.21}{\Character{106}}%
\TextAt{17}{101.739}{494.207}{\Character{6}}%
\TextAt{24}{107.855}{496.476}{0}%
\TextAt{9}{115.889}{496.21}{\Character{61}}%
\TextAt{13}{128.378}{496.21}{\Character{39}}%
\TextAt{16}{136.046}{496.21}{\Character{106}}%
\TextAt{17}{139.367}{494.207}{\Character{6}}%
\TextAt{24}{145.483}{496.476}{0}%
\TextAt{8}{150.132}{496.21}{.}%
\TextAt{8}{158.692}{496.21}{Consequently,}%
\TextAt{13}{233.59}{496.21}{N}%
\TextAt{17}{242.964}{494.207}{\Character{6}}%
\TextAt{24}{249.08}{496.476}{0}%
\TextAt{3}{253.231}{494.207}{\Character{106}}%
\TextAt{21}{255.584}{494.207}{K}%
\TextAt{9}{267.171}{496.21}{\Character{61}}%
\TextAt{13}{279.66}{496.21}{\Character{39}}%
\TextAt{21}{287.328}{494.417}{n}%
\TextAt{16}{292.964}{496.21}{\Character{106}}%
\TextAt{21}{296.285}{494.417}{A}%
\TextAt{24}{302.628}{491.594}{\Character{6}}%
\TextAt{24}{307.831}{493.31}{0}%
\TextAt{8}{312.978}{496.21}{.}%
\TextAt{8}{321.538}{496.21}{On}%
\TextAt{8}{341.087}{496.21}{the}%
\TextAt{8}{361.28}{496.21}{other}%
\TextAt{8}{391.887}{496.21}{hand,}%
\TextAt{8}{424.452}{496.21}{one}%
\TextAt{8}{445.949}{496.21}{has}%
\TextAt{9}{466.86}{496.21}{\Character{6}}%
\TextAt{20}{478.052}{499.232}{\Character{101}}%
\TextAt{13}{475.314}{496.21}{K}%
\TextAt{9}{489.516}{496.21}{\Character{61}}%
\TextAt{20}{502.992}{499.232}{\Character{101}}%
\TextAt{13}{502.005}{496.21}{L}%
\TextAt{8}{513.904}{496.21}{and}%
\TextAt{9}{79.226}{479.164}{\Character{6}}%
\TextAt{16}{90.337}{479.164}{\Character{92}}%
\TextAt{20}{103.702}{482.186}{\Character{101}}%
\TextAt{13}{100.964}{479.164}{K}%
\TextAt{9}{115.102}{479.164}{\Character{61}}%
\TextAt{9}{127.512}{479.164}{\Character{6}}%
\TextAt{17}{135.982}{483.503}{0}%
\TextAt{8}{140.715}{479.164}{,}%
\TextAt{8}{147.864}{479.164}{and}%
\TextAt{8}{170.626}{479.164}{therefore}%
\TextAt{13}{219.726}{479.164}{N}%
\TextAt{17}{229.1}{477.161}{\Character{6}}%
\TextAt{3}{235.216}{477.161}{\Character{106}}%
\TextAt{17}{237.568}{477.161}{\Character{6}}%
\TextAt{24}{243.685}{479.43}{0}%
\TextAt{9}{251.655}{479.164}{\Character{61}}%
\TextAt{13}{264.08}{479.164}{N}%
\TextAt{16}{274.703}{479.164}{\Character{106}}%
\TextAt{21}{278.023}{477.371}{A}%
\TextAt{24}{284.366}{476.032}{\Character{6}}%
\TextAt{8}{290.565}{479.164}{.}%
\TextAt{8}{299.014}{479.164}{For}%
\TextAt{13}{319.983}{479.164}{a}%
\NativeMathAt{329.449}{479.164}{11.955}{\in}%
\TextAt{13}{340.74}{479.164}{A}%
\TextAt{17}{349.515}{477.371}{\Character{6}}%
\TextAt{8}{360.027}{479.164}{we}%
\TextAt{8}{377.254}{479.164}{thus}%
\TextAt{8}{403.006}{479.164}{get}%
\TextAt{13}{181.539}{450.772}{N}%
\TextAt{17}{190.913}{448.913}{\Character{6}}%
\TextAt{3}{197.029}{448.913}{\Character{106}}%
\TextAt{21}{199.381}{448.913}{K}%
\TextAt{9}{207.584}{450.772}{\Character{40}}%
\TextAt{13}{212.136}{450.772}{a}%
\TextAt{9}{218.281}{450.772}{\Character{41}}%
\TextAt{9}{226.136}{450.772}{\Character{61}}%
\TextAt{13}{238.58}{450.772}{N}%
\TextAt{17}{247.954}{448.769}{\Character{6}}%
\TextAt{24}{254.07}{451.038}{0}%
\TextAt{3}{258.221}{448.769}{\Character{106}}%
\TextAt{21}{260.573}{448.769}{K}%
\TextAt{9}{268.776}{450.772}{\Character{40}}%
\TextAt{13}{273.328}{450.772}{N}%
\TextAt{17}{282.702}{448.769}{\Character{6}}%
\TextAt{3}{288.818}{448.769}{\Character{106}}%
\TextAt{17}{291.17}{448.769}{\Character{6}}%
\TextAt{24}{297.286}{451.038}{0}%
\TextAt{9}{301.935}{450.772}{\Character{40}}%
\TextAt{13}{306.488}{450.772}{a}%
\TextAt{9}{312.633}{450.772}{\Character{41}\Character{41}}%
\TextAt{9}{325.031}{450.772}{\Character{61}}%
\TextAt{13}{337.484}{450.772}{N}%
\TextAt{9}{348.106}{450.772}{\Character{40}}%
\TextAt{13}{352.659}{450.772}{a}%
\TextAt{9}{358.803}{450.772}{\Character{41}}%
\TextAt{21}{363.356}{455.709}{\Character{39}}%
\TextAt{22}{368.874}{454.712}{n}%
\TextAt{9}{377.792}{450.772}{\Character{61}}%
\TextAt{13}{390.217}{450.772}{N}%
\TextAt{9}{400.84}{450.772}{\Character{40}}%
\TextAt{13}{405.392}{450.772}{a}%
\TextAt{21}{411.537}{455.708}{\Character{39}}%
\TextAt{22}{417.056}{454.712}{n}%
\TextAt{9}{422.653}{450.772}{\Character{41}}%
\TextAt{13}{427.205}{450.772}{\Character{58}}%
\TextAt{8}{79.226}{422.38}{The}%
\TextAt{8}{103.28}{422.38}{last}%
\TextAt{8}{125.457}{422.38}{equation}%
\TextAt{8}{173.266}{422.38}{follows}%
\TextAt{8}{211.701}{422.38}{from}%
\TextAt{13}{239.342}{422.38}{\Character{39}G}%
\TextAt{9}{256.244}{422.38}{\Character{40}}%
\TextAt{20}{261.783}{425.402}{\Character{101}}%
\TextAt{13}{260.796}{422.38}{L}%
\TextAt{16}{272.081}{422.38}{\Character{106}}%
\TextAt{20}{281.461}{425.402}{\Character{101}}%
\TextAt{13}{278.723}{422.38}{K}%
\TextAt{9}{289.541}{422.38}{\Character{41}}%
\TextAt{9}{297.396}{422.38}{\Character{61}}%
\TextAt{13}{309.839}{422.38}{G}%
\TextAt{9}{319.073}{422.38}{\Character{40}}%
\TextAt{20}{324.612}{425.402}{\Character{101}}%
\TextAt{13}{323.625}{422.38}{L}%
\TextAt{16}{334.911}{422.38}{\Character{106}}%
\TextAt{20}{344.291}{425.402}{\Character{101}}%
\TextAt{13}{341.552}{422.38}{K}%
\TextAt{9}{352.37}{422.38}{\Character{41}}%
\TextAt{13}{356.922}{422.38}{\Character{39}}%
\TextAt{8}{364.591}{422.38}{.}%
\Vector{q 0 G 0 g 0.399 w 452.334 -351.62 m 452.334 -343.55 l S Q}%
\Vector{q 0 G 0 g 0.399 w 452.534 -343.749 m 459.707 -343.749 l S Q}%
\Vector{q 0 G 0 g 0.399 w 452.534 -351.42 m 459.707 -351.42 l S Q}%
\Vector{q 0 G 0 g 0.399 w 459.906 -351.62 m 459.906 -343.55 l S Q}%
\TextAt{8}{97.159}{393.553}{The}%
\TextAt{8}{120.904}{393.553}{subgroup}%
\TextAt{13}{171.393}{393.553}{I}%
\TextAt{21}{176.556}{389.9}{G}%
\TextAt{17}{183.155}{389.9}{\Character{40}}%
\TextAt{30}{187.211}{391.915}{\Character{101}}%
\TextAt{21}{186.448}{389.9}{L}%
\TextAt{3}{192.207}{389.9}{\Character{106}}%
\TextAt{30}{196.53}{391.915}{\Character{101}}%
\TextAt{21}{194.56}{389.9}{K}%
\TextAt{17}{202.264}{389.9}{\Character{41}}%
\TextAt{13}{206.056}{393.553}{U}%
\TextAt{30}{214.768}{391.915}{\Character{101}}%
\TextAt{21}{214.006}{389.9}{L}%
\TextAt{8}{220.263}{393.553}{,}%
\TextAt{8}{227.166}{393.553}{which}%
\TextAt{8}{260.344}{393.553}{is}%
\TextAt{8}{271.799}{393.553}{generated}%
\TextAt{8}{324.823}{393.553}{by}%
\TextAt{8}{340.773}{393.553}{all}%
\TextAt{8}{356.724}{393.553}{elements}%
\TextAt{8}{404.272}{393.553}{of}%
\TextAt{8}{417.294}{393.553}{the}%
\TextAt{8}{437.142}{393.553}{form}%
\TextAt{13}{464.474}{393.553}{u}%
\TextAt{21}{471.136}{397.892}{\Character{28}}%
\TextAt{3}{475.815}{397.892}{\Character{0}}%
\TextAt{17}{482.401}{397.892}{1}%
\TextAt{8}{487.134}{393.553}{,}%
\TextAt{13}{494.036}{393.553}{u}%
\NativeMathAt{504.019}{393.553}{11.955}{\in}%
\TextAt{13}{515.31}{393.553}{U}%
\TextAt{30}{524.023}{391.915}{\Character{101}}%
\TextAt{21}{523.261}{389.9}{L}%
\TextAt{8}{529.518}{393.553}{,}%
\TextAt{13}{79.226}{375.255}{\Character{28}}%
\NativeMathAt{88.966}{375.255}{11.955}{\in}%
\TextAt{13}{100.257}{375.255}{G}%
\TextAt{9}{109.491}{375.255}{\Character{40}}%
\TextAt{20}{115.029}{378.277}{\Character{101}}%
\TextAt{13}{114.043}{375.255}{L}%
\TextAt{16}{125.328}{375.255}{\Character{106}}%
\TextAt{20}{134.708}{378.277}{\Character{101}}%
\TextAt{13}{131.97}{375.255}{K}%
\TextAt{9}{142.788}{375.255}{\Character{41}}%
\TextAt{8}{147.34}{375.255}{,}%
\TextAt{8}{153.974}{375.255}{is}%
\TextAt{8}{165.095}{375.255}{mapped}%
\TextAt{8}{208.996}{375.255}{to}%
\TextAt{9}{222.663}{375.255}{1}%
\TextAt{8}{231.77}{375.255}{by}%
\TextAt{8}{247.385}{375.255}{the}%
\TextAt{8}{266.898}{375.255}{homomorphism}%
\TextAt{13}{348.925}{375.255}{N}%
\TextAt{9}{362.868}{375.255}{\Character{61}}%
\TextAt{13}{375.294}{375.255}{N}%
\TextAt{30}{385.43}{373.617}{\Character{101}}%
\TextAt{21}{384.668}{371.602}{L}%
\TextAt{3}{390.427}{371.602}{\Character{106}}%
\TextAt{30}{394.749}{373.617}{\Character{101}}%
\TextAt{21}{392.779}{371.602}{K}%
\TextAt{9}{404.303}{375.255}{\Character{58}}%
\TextAt{13}{410.875}{375.255}{U}%
\TextAt{30}{419.588}{373.617}{\Character{101}}%
\TextAt{21}{418.825}{371.602}{L}%
\NativeMathAt{428.403}{375.255}{11.955}{\rightarrow}%
\TextAt{13}{443.679}{375.255}{U}%
\TextAt{30}{453.6}{373.617}{\Character{101}}%
\TextAt{21}{451.63}{371.602}{K}%
\TextAt{8}{459.832}{375.255}{.}%
\TextAt{8}{468.066}{375.255}{We}%
\TextAt{8}{487.567}{375.255}{therefore}%
\TextAt{8}{79.226}{358.209}{obtain}%
\TextAt{8}{115.654}{358.209}{an}%
\TextAt{8}{131.913}{358.209}{induced}%
\TextAt{8}{175.477}{358.209}{homomorphism}%
\TextAt{13}{229.677}{329.817}{N}%
\TextAt{9}{243.62}{329.817}{\Character{58}}%
\TextAt{13}{250.193}{329.817}{H}%
\TextAt{17}{259.892}{328.024}{0}%
\TextAt{9}{264.624}{329.817}{\Character{40}}%
\TextAt{13}{269.176}{329.817}{G}%
\TextAt{9}{278.41}{329.817}{\Character{40}}%
\TextAt{20}{283.949}{332.839}{\Character{101}}%
\TextAt{13}{282.962}{329.817}{L}%
\TextAt{16}{294.248}{329.817}{\Character{106}}%
\TextAt{20}{303.628}{332.839}{\Character{101}}%
\TextAt{13}{300.889}{329.817}{K}%
\TextAt{9}{311.707}{329.817}{\Character{41}}%
\TextAt{13}{316.259}{329.817}{\Character{59}}%
\TextAt{13}{321.483}{329.817}{U}%
\TextAt{30}{330.216}{328.179}{\Character{101}}%
\TextAt{21}{329.454}{326.164}{L}%
\TextAt{9}{335.711}{329.817}{\Character{41}}%
\NativeArrowAt{343.584}{329.817}{11.955}{19.26}{\longrightarrow}%
\TextAt{13}{366.166}{329.817}{U}%
\TextAt{30}{376.087}{328.179}{\Character{101}}%
\TextAt{21}{374.116}{326.164}{K}%
\TextAt{8}{79.226}{301.425}{on}%
\TextAt{8}{96.213}{301.425}{the}%
\TextAt{8}{117.097}{301.425}{quotient}%
\TextAt{8}{164.007}{301.425}{group}%
\TextAt{13}{197.911}{301.425}{H}%
\TextAt{17}{207.61}{299.632}{0}%
\TextAt{9}{212.342}{301.425}{\Character{40}}%
\TextAt{13}{216.894}{301.425}{G}%
\TextAt{9}{226.128}{301.425}{\Character{40}}%
\TextAt{20}{231.667}{304.447}{\Character{101}}%
\TextAt{13}{230.68}{301.425}{L}%
\TextAt{16}{243.206}{301.425}{\Character{106}}%
\TextAt{20}{253.826}{304.447}{\Character{101}}%
\TextAt{13}{251.087}{301.425}{K}%
\TextAt{9}{261.905}{301.425}{\Character{41}}%
\TextAt{13}{266.457}{301.425}{\Character{59}}%
\TextAt{13}{271.681}{301.425}{U}%
\TextAt{30}{280.414}{299.787}{\Character{101}}%
\TextAt{21}{279.652}{297.772}{L}%
\TextAt{9}{285.909}{301.425}{\Character{41}}%
\TextAt{9}{295.007}{301.425}{\Character{61}}%
\TextAt{13}{308.688}{301.425}{U}%
\TextAt{30}{317.4}{299.787}{\Character{101}}%
\TextAt{21}{316.638}{297.772}{L}%
\TextAt{9}{322.895}{301.425}{\Character{47}}%
\TextAt{13}{328.748}{301.425}{I}%
\TextAt{21}{333.911}{297.772}{G}%
\TextAt{17}{340.51}{297.772}{\Character{40}}%
\TextAt{30}{344.565}{299.787}{\Character{101}}%
\TextAt{21}{343.803}{297.772}{L}%
\TextAt{3}{349.562}{297.772}{\Character{106}}%
\TextAt{30}{353.885}{299.787}{\Character{101}}%
\TextAt{21}{351.914}{297.772}{K}%
\TextAt{17}{359.619}{297.772}{\Character{41}}%
\TextAt{13}{363.41}{301.425}{U}%
\TextAt{30}{372.123}{299.787}{\Character{101}}%
\TextAt{21}{371.36}{297.772}{L}%
\TextAt{8}{377.618}{301.425}{.}%
\TextAt{8}{388.249}{301.425}{For}%
\TextAt{8}{409.946}{301.425}{this}%
\TextAt{8}{433.496}{301.425}{group,}%
\TextAt{8}{470.833}{301.425}{we}%
\TextAt{8}{488.788}{301.425}{have}%
\TextAt{8}{516.51}{301.425}{the}%
\TextAt{8}{79.226}{284.379}{following}%
\TextAt{8}{128.661}{284.379}{lemma.}%
\TextAt{12}{79.226}{258.541}{Lemma}%
\TextAt{12}{127.202}{258.541}{(5.4).}%
\TextAt{8}{162.788}{258.541}{If}%
\TextAt{13}{174.751}{258.541}{x}%
\NativeMathAt{185.184}{258.541}{11.955}{\in}%
\TextAt{13}{196.936}{258.541}{H}%
\TextAt{17}{206.635}{256.747}{0}%
\TextAt{9}{211.367}{258.541}{\Character{40}}%
\TextAt{13}{215.92}{258.541}{G}%
\TextAt{9}{225.153}{258.541}{\Character{40}}%
\TextAt{20}{230.692}{261.563}{\Character{101}}%
\TextAt{13}{229.706}{258.541}{L}%
\TextAt{16}{241.452}{258.541}{\Character{106}}%
\TextAt{20}{251.292}{261.563}{\Character{101}}%
\TextAt{13}{248.554}{258.541}{K}%
\TextAt{9}{259.372}{258.541}{\Character{41}}%
\TextAt{13}{263.924}{258.541}{\Character{59}}%
\TextAt{13}{269.148}{258.541}{U}%
\TextAt{30}{277.881}{256.902}{\Character{101}}%
\TextAt{21}{277.118}{254.888}{L}%
\TextAt{9}{283.376}{258.541}{\Character{41}}%
\TextAt{8}{292.096}{258.541}{is}%
\TextAt{8}{304.13}{258.541}{fixed}%
\TextAt{8}{332.686}{258.541}{by}%
\TextAt{8}{349.216}{258.541}{an}%
\TextAt{8}{365.745}{258.541}{element}%
\TextAt{13}{409.257}{258.541}{\Character{39}}%
\NativeMathAt{420.707}{258.541}{11.955}{\in}%
\TextAt{13}{432.459}{258.541}{G}%
\TextAt{9}{441.692}{258.541}{\Character{40}}%
\TextAt{20}{447.231}{261.563}{\Character{101}}%
\TextAt{13}{446.245}{258.541}{L}%
\TextAt{16}{457.991}{258.541}{\Character{106}}%
\TextAt{13}{465.093}{258.541}{K}%
\TextAt{9}{475.911}{258.541}{\Character{41}}%
\TextAt{8}{484.631}{258.541}{such}%
\TextAt{8}{511.298}{258.541}{that}%
\TextAt{13}{79.226}{241.495}{d}%
\TextAt{21}{85.309}{239.701}{K}%
\TextAt{9}{93.512}{241.495}{\Character{40}}%
\TextAt{13}{98.064}{241.495}{\Character{39}}%
\TextAt{9}{105.732}{241.495}{\Character{41}}%
\TextAt{9}{113.587}{241.495}{\Character{61}}%
\TextAt{9}{126.008}{241.495}{1}%
\TextAt{8}{131.884}{241.495}{,}%
\TextAt{8}{139.033}{241.495}{i.e.,}%
\TextAt{13}{161.138}{241.495}{x}%
\TextAt{21}{167.79}{245.833}{\Character{39}}%
\TextAt{9}{177.127}{241.495}{\Character{61}}%
\TextAt{13}{189.553}{241.495}{x}%
\TextAt{8}{196.205}{241.495}{,}%
\TextAt{8}{203.354}{241.495}{then}%
\TextAt{13}{263.929}{224.449}{N}%
\TextAt{9}{274.551}{224.449}{\Character{40}}%
\TextAt{13}{279.104}{224.449}{x}%
\TextAt{9}{285.756}{224.449}{\Character{41}}%
\NativeMathAt{293.629}{224.449}{11.955}{\in}%
\TextAt{13}{304.92}{224.449}{N}%
\TextAt{30}{315.056}{222.81}{\Character{101}}%
\TextAt{21}{314.294}{220.796}{L}%
\TextAt{3}{320.053}{220.796}{\Character{106}}%
\TextAt{21}{322.405}{220.796}{K}%
\TextAt{13}{330.608}{224.449}{U}%
\TextAt{30}{339.32}{222.81}{\Character{101}}%
\TextAt{21}{338.558}{220.796}{L}%
\TextAt{13}{344.815}{224.449}{\Character{58}}%
\TextAt{19}{79.226}{183.603}{Proof.}%
\TextAt{8}{116}{183.603}{Let}%
\TextAt{13}{136.97}{183.603}{x}%
\TextAt{9}{146.943}{183.603}{\Character{61}}%
\TextAt{13}{159.368}{183.603}{u}%
\TextAt{8}{169.351}{183.603}{mod}%
\TextAt{13}{194.789}{183.603}{I}%
\TextAt{21}{199.952}{179.95}{G}%
\TextAt{17}{206.551}{179.95}{\Character{40}}%
\TextAt{30}{210.606}{181.965}{\Character{101}}%
\TextAt{21}{209.844}{179.95}{L}%
\TextAt{3}{215.603}{179.95}{\Character{106}}%
\TextAt{30}{219.926}{181.965}{\Character{101}}%
\TextAt{21}{217.955}{179.95}{K}%
\TextAt{17}{225.66}{179.95}{\Character{41}}%
\TextAt{13}{229.451}{183.603}{U}%
\TextAt{30}{238.164}{181.965}{\Character{101}}%
\TextAt{21}{237.402}{179.95}{L}%
\TextAt{8}{243.659}{183.603}{,}%
\TextAt{8}{250.808}{183.603}{with}%
\TextAt{13}{277.468}{183.603}{x}%
\TextAt{21}{284.12}{187.942}{\Character{39}}%
\TextAt{3}{289.638}{187.942}{\Character{0}}%
\TextAt{17}{296.225}{187.942}{1}%
\TextAt{9}{304.278}{183.603}{\Character{61}}%
\TextAt{9}{316.688}{183.603}{1}%
\TextAt{8}{322.556}{183.603}{,}%
\TextAt{8}{329.706}{183.603}{so}%
\TextAt{8}{344.076}{183.603}{that}%
\TextAt{13}{185.976}{146.044}{u}%
\TextAt{21}{192.638}{150.981}{\Character{39}}%
\TextAt{3}{198.157}{150.981}{\Character{0}}%
\TextAt{17}{204.743}{150.981}{1}%
\TextAt{9}{212.796}{146.044}{\Character{61}}%
\TextAt{21}{230.832}{160.988}{r}%
\TextAt{20}{225.222}{157.402}{\Character{89}}%
\TextAt{21}{226.008}{132.207}{i}%
\TextAt{17}{228.891}{132.207}{\Character{61}1}%
\TextAt{13}{242.49}{146.044}{u}%
\TextAt{21}{249.153}{151.377}{\Character{28}}%
\TextAt{22}{252.911}{150.162}{i}%
\TextAt{3}{256.073}{151.377}{\Character{0}}%
\TextAt{17}{262.659}{151.377}{1}%
\TextAt{21}{249.153}{142.981}{i}%
\TextAt{13}{267.392}{146.044}{\Character{59}}%
\TextAt{13}{296.527}{146.044}{u}%
\TextAt{21}{303.208}{144.251}{i}%
\NativeMathAt{309.911}{146.044}{11.955}{\in}%
\TextAt{13}{321.202}{146.044}{U}%
\TextAt{30}{329.914}{144.406}{\Character{101}}%
\TextAt{21}{329.152}{142.391}{L}%
\TextAt{13}{335.409}{146.044}{\Character{59}}%
\TextAt{13}{352.589}{146.044}{\Character{28}}%
\TextAt{21}{357.692}{144.251}{i}%
\NativeMathAt{364.394}{146.044}{11.955}{\in}%
\TextAt{13}{375.685}{146.044}{G}%
\TextAt{9}{384.919}{146.044}{\Character{40}}%
\TextAt{20}{390.458}{149.066}{\Character{101}}%
\TextAt{13}{389.471}{146.044}{L}%
\TextAt{16}{400.756}{146.044}{\Character{106}}%
\TextAt{20}{410.136}{149.066}{\Character{101}}%
\TextAt{13}{407.398}{146.044}{K}%
\TextAt{9}{418.216}{146.044}{\Character{41}}%
\TextAt{13}{422.768}{146.044}{\Character{58}}%
\TextAt{8}{517.802}{146.044}{(*)}%
\TextAt{8}{79.226}{107.013}{Let}%
\TextAt{13}{101.398}{107.013}{M}%
\TextAt{16}{119.341}{107.013}{\Character{106}}%
\TextAt{13}{128.032}{107.013}{K}%
\TextAt{8}{143.95}{107.013}{be}%
\TextAt{8}{161.076}{107.013}{a}%
\TextAt{8}{172.034}{107.013}{finite}%
\TextAt{8}{203.149}{107.013}{Galois}%
\TextAt{8}{240.264}{107.013}{subextension}%
\TextAt{8}{311.177}{107.013}{of}%
\TextAt{20}{326.696}{110.035}{\Character{101}}%
\TextAt{13}{325.71}{107.013}{L}%
\TextAt{16}{339.044}{107.013}{\Character{106}}%
\TextAt{13}{347.735}{107.013}{K}%
\TextAt{8}{358.553}{107.013}{.}%
\TextAt{8}{370.608}{107.013}{In}%
\TextAt{8}{386.432}{107.013}{order}%
\TextAt{8}{418.204}{107.013}{to}%
\TextAt{8}{433.716}{107.013}{prove}%
\TextAt{8}{466.468}{107.013}{that}%
\TextAt{13}{493.04}{107.013}{N}%
\TextAt{9}{503.662}{107.013}{\Character{40}}%
\TextAt{13}{508.215}{107.013}{u}%
\TextAt{9}{514.877}{107.013}{\Character{41}}%
\NativeMathAt{524.799}{107.013}{11.955}{\in}%
\TextAt{13}{79.226}{89.967}{N}%
\TextAt{21}{88.6}{88.108}{M}%
\TextAt{3}{97.582}{88.108}{\Character{106}}%
\TextAt{21}{99.934}{88.108}{K}%
\TextAt{13}{108.137}{89.967}{U}%
\TextAt{21}{116.087}{88.174}{M}%
\TextAt{8}{125.567}{89.967}{,}%
\TextAt{8}{132.688}{89.967}{we}%
\TextAt{8}{149.881}{89.967}{may}%
\TextAt{8}{175.216}{89.967}{assume}%
\TextAt{8}{215.626}{89.967}{that}%
\TextAt{13}{240.96}{89.967}{u\Character{59}}%
\TextAt{13}{252.844}{89.967}{u}%
\TextAt{21}{259.529}{88.174}{i}%
\NativeMathAt{266.232}{89.967}{11.955}{\in}%
\TextAt{13}{277.523}{89.967}{U}%
\TextAt{21}{285.473}{88.174}{M}%
\TextAt{8}{298.816}{89.967}{and}%
\TextAt{13}{321.544}{89.967}{L}%
\NativeMathAt{332.829}{89.967}{11.955}{\subseteq}%
\TextAt{13}{345.449}{89.967}{M}%
\TextAt{8}{358.022}{89.967}{.}%
\TextAt{8}{366.459}{89.967}{Let}%
\TextAt{13}{387.394}{89.967}{n}%
\TextAt{9}{397.702}{89.967}{\Character{61}}%
\TextAt{9}{410.112}{89.967}{\Character{91}}%
\TextAt{13}{413.38}{89.967}{M}%
\TextAt{9}{429.274}{89.967}{\Character{58}}%
\TextAt{13}{435.847}{89.967}{K}%
\TextAt{9}{446.664}{89.967}{\Character{93}}%
\TextAt{8}{449.916}{89.967}{,}%
\TextAt{13}{457.038}{89.967}{\Character{27}}%
\TextAt{9}{467.441}{89.967}{\Character{61}}%
\TextAt{13}{479.866}{89.967}{\Character{39}}%
\TextAt{21}{487.535}{94.306}{n}%
\TextAt{8}{497.034}{89.967}{and}%
\TextAt{8}{519.762}{89.967}{let}%
\TextAt{9}{79.226}{72.921}{\Character{6}}%
\NativeMathAt{91.205}{72.921}{11.955}{\supseteq}%
\TextAt{13}{104.028}{72.921}{M}%
\TextAt{8}{120.618}{72.921}{be}%
\TextAt{8}{136.662}{72.921}{the}%
\TextAt{8}{156.938}{72.921}{fixed}%
\TextAt{8}{185.344}{72.921}{field}%
\TextAt{8}{210.82}{72.921}{of}%
\TextAt{13}{224.27}{72.921}{\Character{27}}%
\TextAt{8}{231.352}{72.921}{.}%
\TextAt{8}{240.159}{72.921}{Further,}%
\TextAt{8}{285.989}{72.921}{let}%
\TextAt{9}{303.013}{72.921}{\Character{6}}%
\TextAt{21}{311.467}{71.128}{n}%
\TextAt{16}{320.628}{72.921}{\Character{106}}%
\TextAt{9}{327.474}{72.921}{\Character{6}}%
\TextAt{8}{339.945}{72.921}{be}%
\TextAt{8}{355.989}{72.921}{the}%
\TextAt{8}{376.265}{72.921}{unramified}%
\TextAt{8}{434.917}{72.921}{extension}%
\TextAt{8}{486.802}{72.921}{of}%
\TextAt{8}{500.251}{72.921}{degree}%
\TextAt{8}{297.211}{43.033}{267}%
\Destination{522.36}{163.095}{nav.AMS.52}%
\Destination{78.23}{809.815}{nav.page.267}%
\Destination{79.36}{619.655}{nav.print.292}%
\Destination{79.36}{635.755}{nav.number.4.5.3}%
\Destination{79.36}{272.555}{nav.number.4.5.4}%
}

% Page 280 (body page 268)
\Page{%
\TextAt{2}{79.226}{796.943}{IV.}%
\TextAt{2}{98.012}{796.943}{Abstract}%
\TextAt{2}{143.252}{796.943}{Class}%
\TextAt{2}{171.844}{796.943}{Field}%
\TextAt{2}{199.575}{796.943}{Theory}%
\Vector{q 0 G 0 g 0.299 w 7.226 18.729 m 460.769 18.729 l S Q}%
\TextAt{13}{79.226}{761.124}{n}%
\TextAt{8}{86.214}{761.124}{,}%
\TextAt{8}{93.88}{761.124}{i.e.,}%
\TextAt{8}{116.501}{761.124}{the}%
\TextAt{8}{137.071}{761.124}{fixed}%
\TextAt{8}{165.77}{761.124}{field}%
\TextAt{8}{191.54}{761.124}{of}%
\TextAt{13}{205.283}{761.124}{\Character{27}}%
\TextAt{21}{212.366}{765.462}{n}%
\TextAt{9}{222.027}{761.124}{\Character{61}}%
\TextAt{13}{235.157}{761.124}{\Character{39}}%
\TextAt{21}{242.825}{765.462}{n}%
\TextAt{17}{242.825}{758.103}{\Character{6}}%
\TextAt{8}{249.439}{761.124}{.}%
\TextAt{8}{259.128}{761.124}{By}%
\TextAt{8}{277.904}{761.124}{(5.1),}%
\TextAt{8}{309.647}{761.124}{we}%
\TextAt{8}{327.288}{761.124}{can}%
\TextAt{8}{349.161}{761.124}{then}%
\TextAt{8}{376.234}{761.124}{find}%
\TextAt{8}{400.056}{761.124}{elements}%
\TextAt{20}{448.661}{761.124}{\Character{101}}%
\TextAt{13}{448.325}{761.124}{u\Character{59}}%
\TextAt{20}{460.567}{761.124}{\Character{101}}%
\TextAt{13}{460.232}{761.124}{u}%
\TextAt{21}{466.894}{759.33}{i}%
\NativeMathAt{474.301}{761.124}{11.955}{\in}%
\TextAt{13}{486.296}{761.124}{U}%
\TextAt{17}{494.246}{759.33}{\Character{6}}%
\TextAt{22}{500.362}{758.334}{n}%
\TextAt{8}{510.27}{761.124}{such}%
\TextAt{8}{79.226}{744.078}{that}%
\TextAt{13}{185.837}{727.032}{u}%
\TextAt{9}{195.82}{727.032}{\Character{61}}%
\TextAt{13}{208.246}{727.032}{N}%
\TextAt{17}{217.62}{725.172}{\Character{6}}%
\TextAt{22}{223.736}{724.176}{n}%
\TextAt{3}{228.835}{725.172}{\Character{106}}%
\TextAt{17}{231.187}{725.172}{\Character{6}}%
\TextAt{9}{237.801}{727.032}{\Character{40}}%
\TextAt{20}{242.689}{727.032}{\Character{101}}%
\TextAt{13}{242.354}{727.032}{u}%
\TextAt{9}{249.016}{727.032}{\Character{41}}%
\TextAt{9}{256.871}{727.032}{\Character{61}}%
\TextAt{20}{269.65}{727.032}{\Character{101}}%
\TextAt{13}{269.315}{727.032}{u}%
\TextAt{21}{275.977}{731.968}{\Character{27}}%
\TextAt{22}{280.816}{730.972}{n}%
\TextAt{13}{286.413}{727.032}{\Character{59}}%
\TextAt{13}{315.548}{727.032}{u}%
\TextAt{21}{322.23}{725.238}{i}%
\TextAt{9}{328.932}{727.032}{\Character{61}}%
\TextAt{13}{341.358}{727.032}{N}%
\TextAt{17}{350.732}{725.172}{\Character{6}}%
\TextAt{22}{356.848}{724.176}{n}%
\TextAt{3}{361.947}{725.172}{\Character{106}}%
\TextAt{17}{364.299}{725.172}{\Character{6}}%
\TextAt{9}{370.913}{727.032}{\Character{40}}%
\TextAt{20}{375.801}{727.032}{\Character{101}}%
\TextAt{13}{375.466}{727.032}{u}%
\TextAt{21}{382.128}{725.238}{i}%
\TextAt{9}{385.509}{727.032}{\Character{41}}%
\TextAt{9}{393.364}{727.032}{\Character{61}}%
\TextAt{20}{406.143}{727.032}{\Character{101}}%
\TextAt{13}{405.808}{727.032}{u}%
\TextAt{21}{412.47}{732.146}{\Character{27}}%
\TextAt{22}{417.31}{731.15}{n}%
\TextAt{21}{412.47}{723.968}{i}%
\TextAt{13}{422.907}{727.032}{\Character{58}}%
\TextAt{8}{79.226}{703.676}{By}%
\TextAt{8}{99.008}{703.676}{(*),}%
\TextAt{8}{122.898}{703.676}{the}%
\TextAt{8}{144.473}{703.676}{elements}%
\TextAt{20}{194.083}{703.676}{\Character{101}}%
\TextAt{13}{193.748}{703.676}{u}%
\TextAt{21}{200.41}{708.014}{\Character{39}}%
\TextAt{3}{205.929}{708.014}{\Character{0}}%
\TextAt{17}{212.515}{708.014}{1}%
\TextAt{8}{222.563}{703.676}{and}%
\TextAt{20}{246.744}{712.642}{\Character{81}}%
\TextAt{21}{258.035}{700.189}{i}%
\TextAt{20}{263.745}{703.676}{\Character{101}}%
\TextAt{13}{263.409}{703.676}{u}%
\TextAt{21}{270.072}{709.009}{\Character{28}}%
\TextAt{22}{273.83}{707.794}{i}%
\TextAt{3}{276.992}{709.009}{\Character{0}}%
\TextAt{17}{283.578}{709.009}{1}%
\TextAt{21}{270.072}{700.612}{i}%
\TextAt{8}{293.626}{703.676}{only}%
\TextAt{8}{320.736}{703.676}{differ}%
\TextAt{8}{352.389}{703.676}{by}%
\TextAt{8}{370.067}{703.676}{an}%
\TextAt{8}{387.745}{703.676}{element}%
\TextAt{20}{432.735}{703.676}{\Character{101}}%
\TextAt{13}{432.405}{703.676}{x}%
\NativeMathAt{444.795}{703.676}{11.955}{\in}%
\TextAt{13}{458.503}{703.676}{U}%
\TextAt{17}{466.454}{701.882}{\Character{6}}%
\TextAt{22}{472.57}{700.886}{n}%
\TextAt{8}{483.483}{703.676}{such}%
\TextAt{8}{511.298}{703.676}{that}%
\TextAt{13}{79.226}{686.63}{N}%
\TextAt{17}{88.6}{684.77}{\Character{6}}%
\TextAt{22}{94.716}{683.774}{n}%
\TextAt{3}{99.815}{684.77}{\Character{106}}%
\TextAt{17}{102.167}{684.77}{\Character{6}}%
\TextAt{9}{108.781}{686.63}{\Character{40}}%
\TextAt{20}{113.664}{686.63}{\Character{101}}%
\TextAt{13}{113.334}{686.63}{x}%
\TextAt{9}{119.986}{686.63}{\Character{41}}%
\TextAt{9}{128.462}{686.63}{\Character{61}}%
\TextAt{9}{141.493}{686.63}{1}%
\TextAt{8}{147.369}{686.63}{.}%
\TextAt{8}{156.901}{686.63}{Hence—again}%
\TextAt{8}{231.074}{686.63}{by}%
\TextAt{8}{247.694}{686.63}{(5.1)—they}%
\TextAt{8}{310.175}{686.63}{differ}%
\TextAt{8}{340.771}{686.63}{by}%
\TextAt{8}{357.391}{686.63}{an}%
\TextAt{8}{374.011}{686.63}{element}%
\TextAt{8}{417.614}{686.63}{of}%
\TextAt{8}{431.306}{686.63}{the}%
\TextAt{8}{451.824}{686.63}{form}%
\TextAt{20}{480.223}{686.63}{\Character{101}}%
\TextAt{13}{479.825}{686.63}{y}%
\TextAt{21}{485.962}{690.968}{\Character{27}}%
\TextAt{3}{491.087}{690.968}{\Character{0}}%
\TextAt{17}{497.674}{690.968}{1}%
\TextAt{8}{502.406}{686.63}{,}%
\TextAt{8}{510.007}{686.63}{with}%
\TextAt{20}{79.624}{669.584}{\Character{101}}%
\TextAt{13}{79.226}{669.584}{y}%
\NativeMathAt{88.684}{669.584}{11.955}{\in}%
\TextAt{13}{99.975}{669.584}{U}%
\TextAt{17}{107.925}{667.791}{\Character{6}}%
\TextAt{22}{114.041}{666.794}{n}%
\TextAt{8}{119.638}{669.584}{.}%
\TextAt{8}{128.086}{669.584}{We}%
\TextAt{8}{148.231}{669.584}{may}%
\TextAt{8}{173.6}{669.584}{thus}%
\TextAt{8}{199.351}{669.584}{write}%
\TextAt{20}{197.416}{640.969}{\Character{101}}%
\TextAt{13}{197.08}{640.969}{u}%
\TextAt{21}{203.743}{645.905}{\Character{39}}%
\TextAt{3}{209.261}{645.905}{\Character{0}}%
\TextAt{17}{215.848}{645.905}{1}%
\TextAt{9}{223.901}{640.969}{\Character{61}}%
\TextAt{20}{236.724}{640.969}{\Character{101}}%
\TextAt{13}{236.326}{640.969}{y}%
\TextAt{21}{242.463}{645.905}{\Character{39}}%
\TextAt{22}{247.981}{648.718}{n}%
\TextAt{3}{253.08}{645.905}{\Character{0}}%
\TextAt{17}{259.667}{645.905}{1}%
\TextAt{20}{266.391}{652.327}{\Character{89}}%
\TextAt{21}{272.588}{627.132}{i}%
\TextAt{20}{283.995}{640.969}{\Character{101}}%
\TextAt{13}{283.66}{640.969}{u}%
\TextAt{21}{290.322}{646.302}{\Character{28}}%
\TextAt{22}{294.081}{645.087}{i}%
\TextAt{3}{297.242}{646.302}{\Character{0}}%
\TextAt{17}{303.829}{646.302}{1}%
\TextAt{21}{290.322}{637.905}{i}%
\TextAt{9}{311.882}{640.969}{\Character{61}}%
\TextAt{9}{324.292}{640.969}{\Character{40}}%
\TextAt{20}{329.258}{640.969}{\Character{101}}%
\TextAt{13}{328.86}{640.969}{y}%
\TextAt{21}{334.996}{645.905}{\Character{39}}%
\TextAt{22}{340.515}{644.909}{n}%
\TextAt{9}{346.112}{640.969}{\Character{41}}%
\TextAt{21}{350.664}{645.905}{\Character{39}}%
\TextAt{3}{356.182}{645.905}{\Character{0}}%
\TextAt{17}{362.769}{645.905}{1}%
\TextAt{20}{369.494}{652.327}{\Character{89}}%
\TextAt{21}{375.69}{627.132}{i}%
\TextAt{20}{387.098}{640.969}{\Character{101}}%
\TextAt{13}{386.762}{640.969}{u}%
\TextAt{21}{393.425}{646.302}{\Character{28}}%
\TextAt{22}{397.183}{645.087}{i}%
\TextAt{3}{400.345}{646.302}{\Character{0}}%
\TextAt{17}{406.931}{646.302}{1}%
\TextAt{21}{393.425}{637.905}{i}%
\TextAt{13}{411.663}{640.969}{\Character{58}}%
\TextAt{8}{79.226}{603.897}{Applying}%
\TextAt{13}{129.952}{603.897}{N}%
\TextAt{8}{144.472}{603.897}{gives}%
\TextAt{13}{173.152}{603.897}{N}%
\TextAt{9}{183.775}{603.897}{\Character{40}}%
\TextAt{20}{188.663}{603.897}{\Character{101}}%
\TextAt{13}{188.327}{603.897}{u}%
\TextAt{9}{194.99}{603.897}{\Character{41}}%
\TextAt{21}{199.542}{608.236}{\Character{39}}%
\TextAt{3}{205.06}{608.236}{\Character{0}}%
\TextAt{17}{211.647}{608.236}{1}%
\TextAt{9}{219.7}{603.897}{\Character{61}}%
\TextAt{13}{232.125}{603.897}{N}%
\TextAt{9}{242.748}{603.897}{\Character{40}}%
\TextAt{20}{247.698}{603.897}{\Character{101}}%
\TextAt{13}{247.3}{603.897}{y}%
\TextAt{21}{253.437}{608.236}{\Character{39}}%
\TextAt{22}{258.955}{607.239}{n}%
\TextAt{9}{264.552}{603.897}{\Character{41}}%
\TextAt{21}{269.105}{608.236}{\Character{39}}%
\TextAt{3}{274.623}{608.236}{\Character{0}}%
\TextAt{17}{281.209}{608.236}{1}%
\TextAt{8}{285.942}{603.897}{,}%
\TextAt{8}{293.091}{603.897}{so}%
\TextAt{8}{307.461}{603.897}{that}%
\TextAt{13}{257.512}{575.283}{N}%
\TextAt{9}{268.135}{575.283}{\Character{40}}%
\TextAt{20}{273.022}{575.283}{\Character{101}}%
\TextAt{13}{272.687}{575.283}{u}%
\TextAt{9}{279.349}{575.283}{\Character{41}}%
\TextAt{9}{287.204}{575.283}{\Character{61}}%
\TextAt{13}{299.648}{575.283}{N}%
\TextAt{9}{310.271}{575.283}{\Character{40}}%
\TextAt{20}{315.221}{575.283}{\Character{101}}%
\TextAt{13}{314.823}{575.283}{y}%
\TextAt{21}{320.96}{580.219}{\Character{39}}%
\TextAt{22}{326.478}{579.223}{n}%
\TextAt{9}{332.075}{575.283}{\Character{41}}%
\TextAt{16}{339.284}{575.283}{\Character{1}}%
\TextAt{13}{345.261}{575.283}{z}%
\TextAt{13}{351.215}{575.283}{\Character{59}}%
\TextAt{8}{79.226}{546.668}{for}%
\TextAt{8}{97.19}{546.668}{some}%
\TextAt{13}{126.596}{546.668}{z}%
\NativeMathAt{136.023}{546.668}{11.955}{\in}%
\TextAt{13}{147.449}{546.668}{U}%
\TextAt{30}{157.37}{545.03}{\Character{101}}%
\TextAt{21}{155.399}{543.015}{K}%
\TextAt{8}{167.579}{546.668}{such}%
\TextAt{8}{194.055}{546.668}{that}%
\TextAt{13}{219.503}{546.668}{z}%
\TextAt{21}{225.474}{551.007}{\Character{39}}%
\TextAt{3}{230.992}{551.007}{\Character{0}}%
\TextAt{17}{237.579}{551.007}{1}%
\TextAt{9}{245.767}{546.668}{\Character{61}}%
\TextAt{9}{258.32}{546.668}{1}%
\TextAt{8}{264.181}{546.668}{;}%
\TextAt{8}{271.449}{546.668}{therefore}%
\TextAt{13}{320.629}{546.668}{z}%
\TextAt{21}{326.599}{551.007}{\Character{39}}%
\TextAt{9}{336.072}{546.668}{\Character{61}}%
\TextAt{13}{348.633}{546.668}{z}%
\TextAt{8}{354.603}{546.668}{,}%
\TextAt{8}{361.852}{546.668}{and}%
\TextAt{13}{384.694}{546.668}{z}%
\NativeMathAt{394.12}{546.668}{11.955}{\in}%
\TextAt{13}{405.547}{546.668}{U}%
\TextAt{21}{413.497}{544.875}{K}%
\TextAt{8}{421.7}{546.668}{.}%
\TextAt{8}{430.386}{546.668}{Finally,}%
\TextAt{8}{472.592}{546.668}{applying}%
\TextAt{13}{520.48}{546.668}{\Character{27}}%
\TextAt{21}{527.133}{544.875}{n}%
\TextAt{8}{79.226}{529.622}{and}%
\TextAt{8}{101.682}{529.622}{putting}%
\TextAt{13}{143.002}{529.622}{y}%
\TextAt{9}{152.46}{529.622}{\Character{61}}%
\TextAt{20}{165.283}{529.622}{\Character{101}}%
\TextAt{13}{164.885}{529.622}{y}%
\TextAt{21}{171.022}{533.961}{\Character{27}}%
\TextAt{22}{175.861}{532.964}{n}%
\TextAt{9}{184.779}{529.622}{\Character{61}}%
\TextAt{13}{197.205}{529.622}{N}%
\TextAt{17}{206.578}{527.762}{\Character{6}}%
\TextAt{22}{212.694}{526.766}{n}%
\TextAt{3}{217.793}{527.762}{\Character{106}}%
\TextAt{17}{220.146}{527.762}{\Character{6}}%
\TextAt{9}{226.76}{529.622}{\Character{40}}%
\TextAt{20}{231.71}{529.622}{\Character{101}}%
\TextAt{13}{231.312}{529.622}{y}%
\TextAt{9}{237.449}{529.622}{\Character{41}}%
\NativeMathAt{245.322}{529.622}{11.955}{\in}%
\TextAt{13}{256.613}{529.622}{U}%
\TextAt{17}{264.563}{527.829}{\Character{6}}%
\TextAt{8}{271.177}{529.622}{,}%
\TextAt{8}{278.081}{529.622}{we}%
\TextAt{8}{295.001}{529.622}{obtain,}%
\TextAt{8}{334.434}{529.622}{observing}%
\TextAt{13}{386.551}{529.622}{n}%
\TextAt{9}{396.859}{529.622}{\Character{61}}%
\TextAt{9}{409.269}{529.622}{\Character{91}}%
\TextAt{13}{412.536}{529.622}{M}%
\TextAt{9}{428.431}{529.622}{\Character{58}}%
\TextAt{13}{435.003}{529.622}{K}%
\TextAt{9}{445.821}{529.622}{\Character{93}}%
\TextAt{8}{452.663}{529.622}{and}%
\TextAt{8}{475.118}{529.622}{using}%
\TextAt{8}{505.44}{529.622}{(5.3),}%
\TextAt{8}{79.226}{512.576}{that}%
\TextAt{13}{196.966}{483.962}{N}%
\TextAt{9}{207.589}{483.962}{\Character{40}}%
\TextAt{13}{212.141}{483.962}{u}%
\TextAt{9}{218.803}{483.962}{\Character{41}}%
\TextAt{9}{226.658}{483.962}{\Character{61}}%
\TextAt{13}{239.102}{483.962}{N}%
\TextAt{9}{249.725}{483.962}{\Character{40}}%
\TextAt{20}{254.613}{483.962}{\Character{101}}%
\TextAt{13}{254.277}{483.962}{u}%
\TextAt{9}{260.939}{483.962}{\Character{41}}%
\TextAt{21}{265.492}{488.898}{\Character{27}}%
\TextAt{22}{270.331}{487.901}{n}%
\TextAt{9}{279.249}{483.962}{\Character{61}}%
\TextAt{13}{291.674}{483.962}{N}%
\TextAt{9}{302.297}{483.962}{\Character{40}}%
\TextAt{20}{307.247}{483.962}{\Character{101}}%
\TextAt{13}{306.849}{483.962}{y}%
\TextAt{21}{312.986}{488.898}{\Character{39}}%
\TextAt{22}{318.504}{487.901}{n}%
\TextAt{9}{324.101}{483.962}{\Character{41}}%
\TextAt{21}{328.654}{488.898}{\Character{27}}%
\TextAt{22}{333.493}{487.901}{n}%
\TextAt{13}{339.09}{483.962}{z}%
\TextAt{21}{345.061}{488.898}{n}%
\TextAt{9}{354.018}{483.962}{\Character{61}}%
\TextAt{13}{366.443}{483.962}{N}%
\TextAt{9}{377.066}{483.962}{\Character{40}}%
\TextAt{13}{381.618}{483.962}{y}%
\TextAt{21}{387.755}{488.898}{\Character{39}}%
\TextAt{22}{393.273}{487.901}{n}%
\TextAt{9}{398.87}{483.962}{\Character{41}}%
\TextAt{13}{403.423}{483.961}{z}%
\TextAt{21}{409.393}{488.898}{n}%
\TextAt{9}{226.677}{463.927}{\Character{61}}%
\TextAt{13}{239.102}{463.927}{N}%
\TextAt{17}{248.476}{462.067}{\Character{6}}%
\TextAt{3}{254.592}{462.067}{\Character{106}}%
\TextAt{21}{256.944}{462.067}{K}%
\TextAt{9}{265.147}{463.927}{\Character{40}}%
\TextAt{13}{269.699}{463.927}{y}%
\TextAt{9}{275.836}{463.927}{\Character{41}}%
\TextAt{13}{280.388}{463.927}{N}%
\TextAt{21}{289.762}{462.067}{M}%
\TextAt{3}{298.744}{462.067}{\Character{106}}%
\TextAt{21}{301.096}{462.067}{K}%
\TextAt{9}{309.299}{463.927}{\Character{40}}%
\TextAt{13}{313.851}{463.927}{z}%
\TextAt{9}{319.822}{463.927}{\Character{41}}%
\NativeMathAt{327.695}{463.927}{11.955}{\in}%
\TextAt{13}{338.986}{463.927}{N}%
\TextAt{21}{348.36}{462.067}{M}%
\TextAt{3}{357.341}{462.067}{\Character{106}}%
\TextAt{21}{359.694}{462.067}{K}%
\TextAt{13}{367.896}{463.927}{U}%
\TextAt{21}{375.847}{462.133}{M}%
\TextAt{13}{385.326}{463.927}{\Character{58}}%
\Vector{q 0 G 0 g 0.399 w 452.334 -338.688 m 452.334 -330.618 l S Q}%
\Vector{q 0 G 0 g 0.399 w 452.534 -330.817 m 459.707 -330.817 l S Q}%
\Vector{q 0 G 0 g 0.399 w 452.534 -338.489 m 459.707 -338.489 l S Q}%
\Vector{q 0 G 0 g 0.399 w 459.906 -338.688 m 459.906 -330.618 l S Q}%
\TextAt{12}{79.226}{400.554}{Proposition}%
\TextAt{12}{151.937}{400.554}{(5.5).}%
\TextAt{8}{187.253}{400.554}{The}%
\TextAt{8}{211.307}{400.554}{reciprocity}%
\TextAt{8}{269.517}{400.554}{map}%
\TextAt{13}{214.327}{371.94}{r}%
\TextAt{30}{220.364}{370.301}{\Character{101}}%
\TextAt{21}{219.602}{368.287}{L}%
\TextAt{3}{225.361}{368.287}{\Character{106}}%
\TextAt{21}{227.713}{368.287}{K}%
\TextAt{9}{239.237}{371.94}{\Character{58}}%
\TextAt{8}{245.809}{371.94}{Frob}%
\TextAt{9}{270.365}{371.94}{\Character{40}}%
\TextAt{20}{275.904}{374.962}{\Character{101}}%
\TextAt{13}{274.918}{371.94}{L}%
\TextAt{16}{286.203}{371.94}{\Character{106}}%
\TextAt{13}{292.845}{371.94}{K}%
\TextAt{9}{303.662}{371.94}{\Character{41}}%
\NativeArrowAt{311.536}{371.94}{11.955}{19.26}{\longrightarrow}%
\TextAt{13}{334.118}{371.94}{A}%
\TextAt{21}{342.893}{370.147}{K}%
\TextAt{9}{351.096}{371.94}{\Character{47}}%
\TextAt{13}{356.949}{371.94}{N}%
\TextAt{30}{367.085}{370.301}{\Character{101}}%
\TextAt{21}{366.322}{368.287}{L}%
\TextAt{3}{372.081}{368.287}{\Character{106}}%
\TextAt{21}{374.434}{368.287}{K}%
\TextAt{13}{382.636}{371.94}{A}%
\TextAt{30}{392.174}{370.301}{\Character{101}}%
\TextAt{21}{391.412}{368.287}{L}%
\TextAt{8}{79.226}{343.325}{is}%
\TextAt{8}{90.99}{343.325}{multiplicative.}%
\TextAt{19}{79.226}{308.567}{Proof.}%
\TextAt{8}{116}{308.567}{Let}%
\TextAt{13}{136.907}{308.567}{\Character{27}}%
\TextAt{17}{143.561}{306.774}{1}%
\TextAt{13}{148.293}{308.567}{\Character{27}}%
\TextAt{17}{154.946}{306.774}{2}%
\TextAt{9}{163}{308.567}{\Character{61}}%
\TextAt{13}{175.425}{308.567}{\Character{27}}%
\TextAt{17}{182.079}{306.774}{3}%
\TextAt{8}{190.646}{308.567}{be}%
\TextAt{8}{206.508}{308.567}{an}%
\TextAt{8}{222.704}{308.567}{equation}%
\TextAt{8}{270.451}{308.567}{in}%
\TextAt{8}{284.041}{308.567}{Frob}%
\TextAt{9}{308.597}{308.567}{\Character{40}}%
\TextAt{20}{314.136}{311.589}{\Character{101}}%
\TextAt{13}{313.149}{308.567}{L}%
\TextAt{16}{324.434}{308.567}{\Character{106}}%
\TextAt{13}{331.076}{308.567}{K}%
\TextAt{9}{341.894}{308.567}{\Character{41}}%
\TextAt{8}{346.446}{308.567}{,}%
\TextAt{13}{353.545}{308.567}{n}%
\TextAt{21}{360.533}{306.774}{i}%
\TextAt{9}{367.235}{308.567}{\Character{61}}%
\TextAt{13}{379.661}{308.567}{d}%
\TextAt{21}{385.743}{306.774}{K}%
\TextAt{9}{393.946}{308.567}{\Character{40}}%
\TextAt{13}{398.498}{308.567}{\Character{27}}%
\TextAt{21}{405.152}{306.774}{i}%
\TextAt{9}{408.533}{308.567}{\Character{41}}%
\TextAt{8}{413.085}{308.567}{,}%
\TextAt{9}{420.185}{308.567}{\Character{6}}%
\TextAt{21}{428.639}{306.774}{i}%
\TextAt{8}{435.855}{308.567}{the}%
\TextAt{8}{455.949}{308.567}{fixed}%
\TextAt{8}{484.173}{308.567}{field}%
\TextAt{8}{509.467}{308.567}{of}%
\TextAt{13}{522.735}{308.567}{\Character{27}}%
\TextAt{21}{529.388}{306.774}{i}%
\TextAt{8}{79.226}{291.521}{and}%
\TextAt{13}{101.989}{291.521}{\Character{25}}%
\TextAt{21}{108.629}{289.728}{i}%
\NativeMathAt{115.331}{291.521}{11.955}{\in}%
\TextAt{13}{126.622}{291.521}{A}%
\TextAt{17}{135.398}{289.728}{\Character{6}}%
\TextAt{22}{141.514}{288.513}{i}%
\TextAt{8}{149.071}{291.521}{a}%
\TextAt{8}{158.826}{291.521}{prime}%
\TextAt{8}{191.99}{291.521}{element,}%
\TextAt{8}{238.484}{291.521}{for}%
\TextAt{13}{256.369}{291.521}{i}%
\TextAt{9}{263.683}{291.521}{\Character{61}}%
\TextAt{9}{276.093}{291.521}{1}%
\TextAt{13}{281.961}{291.521}{\Character{59}}%
\TextAt{9}{287.205}{291.521}{2}%
\TextAt{13}{293.058}{291.521}{\Character{59}}%
\TextAt{9}{298.303}{291.521}{3}%
\TextAt{8}{304.156}{291.521}{.}%
\TextAt{8}{312.604}{291.521}{We}%
\TextAt{8}{332.748}{291.521}{have}%
\TextAt{8}{359.743}{291.521}{to}%
\TextAt{8}{374.053}{291.521}{show}%
\TextAt{8}{403.057}{291.521}{that}%
\TextAt{13}{176.451}{262.907}{N}%
\TextAt{17}{185.825}{261.047}{\Character{6}}%
\TextAt{24}{191.941}{259.94}{1}%
\TextAt{3}{196.092}{261.047}{\Character{106}}%
\TextAt{21}{198.444}{261.047}{K}%
\TextAt{9}{206.647}{262.907}{\Character{40}}%
\TextAt{13}{211.199}{262.907}{\Character{25}}%
\TextAt{17}{217.84}{261.114}{1}%
\TextAt{9}{222.572}{262.907}{\Character{41}}%
\TextAt{13}{227.124}{262.907}{N}%
\TextAt{17}{236.498}{261.047}{\Character{6}}%
\TextAt{24}{242.614}{259.94}{2}%
\TextAt{3}{246.765}{261.047}{\Character{106}}%
\TextAt{21}{249.118}{261.047}{K}%
\TextAt{9}{257.32}{262.907}{\Character{40}}%
\TextAt{13}{261.873}{262.907}{\Character{25}}%
\TextAt{17}{268.513}{261.114}{2}%
\TextAt{9}{273.245}{262.907}{\Character{41}}%
\NativeMathAt{281.118}{262.907}{11.955}{\equiv}%
\TextAt{13}{293.738}{262.907}{N}%
\TextAt{17}{303.111}{261.047}{\Character{6}}%
\TextAt{24}{309.227}{259.94}{3}%
\TextAt{3}{313.379}{261.047}{\Character{106}}%
\TextAt{21}{315.731}{261.047}{K}%
\TextAt{9}{323.934}{262.907}{\Character{40}}%
\TextAt{13}{328.486}{262.907}{\Character{25}}%
\TextAt{17}{335.126}{261.114}{3}%
\TextAt{9}{339.858}{262.907}{\Character{41}}%
\TextAt{9}{356.356}{262.907}{\Character{40}}%
\TextAt{8}{360.918}{262.907}{mod}%
\TextAt{13}{387.02}{262.907}{N}%
\TextAt{30}{397.156}{261.269}{\Character{101}}%
\TextAt{21}{396.394}{259.254}{L}%
\TextAt{3}{402.153}{259.254}{\Character{106}}%
\TextAt{21}{404.505}{259.254}{K}%
\TextAt{13}{412.708}{262.907}{A}%
\TextAt{30}{422.246}{261.269}{\Character{101}}%
\TextAt{21}{421.483}{259.254}{L}%
\TextAt{9}{427.741}{262.907}{\Character{41}}%
\TextAt{13}{432.293}{262.907}{\Character{58}}%
\TextAt{8}{79.226}{233.04}{Choose}%
\TextAt{8}{119.933}{233.04}{a}%
\TextAt{8}{129.689}{233.04}{fixed}%
\TextAt{13}{157.975}{233.04}{\Character{39}}%
\NativeMathAt{168.964}{233.04}{11.955}{\in}%
\TextAt{13}{180.255}{233.04}{G}%
\TextAt{9}{189.489}{233.04}{\Character{40}}%
\TextAt{20}{195.027}{236.062}{\Character{101}}%
\TextAt{13}{194.041}{233.04}{L}%
\TextAt{16}{205.326}{233.04}{\Character{106}}%
\TextAt{13}{211.968}{233.04}{K}%
\TextAt{9}{222.786}{233.04}{\Character{41}}%
\TextAt{8}{231.235}{233.04}{such}%
\TextAt{8}{257.632}{233.04}{that}%
\TextAt{13}{283.001}{233.04}{d}%
\TextAt{21}{289.084}{231.247}{K}%
\TextAt{9}{297.287}{233.04}{\Character{40}}%
\TextAt{13}{301.839}{233.04}{\Character{39}}%
\TextAt{9}{309.507}{233.04}{\Character{41}}%
\TextAt{9}{317.362}{233.04}{\Character{61}}%
\TextAt{9}{329.783}{233.04}{1}%
\TextAt{8}{339.556}{233.04}{and}%
\TextAt{8}{362.319}{233.04}{put}%
\TextAt{13}{243.985}{204.426}{\Character{28}}%
\TextAt{21}{249.069}{202.632}{i}%
\TextAt{9}{255.771}{204.426}{\Character{61}}%
\TextAt{13}{268.196}{204.426}{\Character{27}}%
\TextAt{3}{275.279}{209.374}{\Character{0}}%
\TextAt{17}{281.865}{209.374}{1}%
\TextAt{21}{274.85}{201.362}{i}%
\TextAt{13}{286.598}{204.426}{\Character{39}}%
\TextAt{21}{294.266}{209.362}{n}%
\TextAt{22}{299.404}{208.147}{i}%
\NativeMathAt{306.385}{204.426}{11.955}{\in}%
\TextAt{13}{317.676}{204.426}{G}%
\TextAt{9}{326.91}{204.426}{\Character{40}}%
\TextAt{20}{332.448}{207.448}{\Character{101}}%
\TextAt{13}{331.462}{204.426}{L}%
\TextAt{16}{342.747}{204.426}{\Character{106}}%
\TextAt{20}{352.127}{207.448}{\Character{101}}%
\TextAt{13}{349.389}{204.426}{K}%
\TextAt{9}{360.207}{204.426}{\Character{41}}%
\TextAt{13}{364.759}{204.426}{\Character{58}}%
\TextAt{8}{79.226}{175.811}{From}%
\TextAt{13}{109.951}{175.811}{\Character{27}}%
\TextAt{17}{116.604}{174.018}{1}%
\TextAt{13}{121.337}{175.811}{\Character{27}}%
\TextAt{17}{127.99}{174.018}{2}%
\TextAt{9}{136.043}{175.811}{\Character{61}}%
\TextAt{13}{148.469}{175.811}{\Character{27}}%
\TextAt{17}{155.122}{174.018}{3}%
\TextAt{8}{163.752}{175.811}{and}%
\TextAt{13}{186.515}{175.811}{n}%
\TextAt{17}{193.502}{174.018}{1}%
\TextAt{9}{200.891}{175.811}{\Character{43}}%
\TextAt{13}{212.652}{175.811}{n}%
\TextAt{17}{219.64}{174.018}{2}%
\TextAt{9}{227.693}{175.811}{\Character{61}}%
\TextAt{13}{240.119}{175.811}{n}%
\TextAt{17}{247.106}{174.018}{3}%
\TextAt{8}{251.839}{175.811}{,}%
\TextAt{8}{258.988}{175.811}{we}%
\TextAt{8}{276.215}{175.811}{then}%
\TextAt{8}{302.875}{175.811}{deduce}%
\TextAt{8}{341.885}{175.811}{that}%
\TextAt{13}{185.341}{147.196}{\Character{28}}%
\TextAt{17}{190.424}{145.403}{3}%
\TextAt{9}{198.478}{147.196}{\Character{61}}%
\TextAt{13}{210.903}{147.196}{\Character{27}}%
\TextAt{3}{217.985}{152.133}{\Character{0}}%
\TextAt{17}{224.572}{152.133}{1}%
\TextAt{17}{217.557}{144.241}{2}%
\TextAt{13}{229.304}{147.196}{\Character{27}}%
\TextAt{3}{236.387}{152.133}{\Character{0}}%
\TextAt{17}{242.973}{152.133}{1}%
\TextAt{17}{235.958}{144.241}{1}%
\TextAt{13}{247.706}{147.196}{\Character{39}}%
\TextAt{21}{255.374}{152.133}{n}%
\TextAt{24}{260.512}{151.026}{2}%
\TextAt{17}{264.663}{152.133}{\Character{43}}%
\TextAt{21}{271.25}{152.133}{n}%
\TextAt{24}{276.388}{151.026}{1}%
\TextAt{9}{284.358}{147.196}{\Character{61}}%
\TextAt{13}{296.783}{147.196}{\Character{27}}%
\TextAt{3}{303.866}{152.133}{\Character{0}}%
\TextAt{17}{310.452}{152.133}{1}%
\TextAt{17}{303.437}{144.241}{2}%
\TextAt{13}{315.185}{147.196}{\Character{39}}%
\TextAt{21}{322.853}{152.133}{n}%
\TextAt{24}{327.991}{151.026}{2}%
\TextAt{9}{332.64}{147.196}{\Character{40}}%
\TextAt{13}{337.193}{147.196}{\Character{39}}%
\TextAt{3}{344.861}{152.133}{\Character{0}}%
\TextAt{21}{351.447}{152.133}{n}%
\TextAt{24}{356.586}{151.026}{2}%
\TextAt{13}{361.235}{147.196}{\Character{27}}%
\TextAt{17}{367.888}{145.403}{1}%
\TextAt{13}{372.621}{147.196}{\Character{39}}%
\TextAt{21}{380.289}{152.133}{n}%
\TextAt{24}{385.427}{151.026}{2}%
\TextAt{9}{390.076}{147.196}{\Character{41}}%
\TextAt{3}{394.629}{152.133}{\Character{0}}%
\TextAt{17}{401.215}{152.133}{1}%
\TextAt{13}{405.948}{147.196}{\Character{39}}%
\TextAt{21}{413.616}{152.133}{n}%
\TextAt{24}{418.754}{151.026}{1}%
\TextAt{13}{423.403}{147.196}{\Character{58}}%
\TextAt{8}{79.226}{118.582}{Putting}%
\TextAt{13}{123.493}{118.582}{\Character{27}}%
\TextAt{17}{130.147}{116.789}{4}%
\TextAt{9}{140.212}{118.582}{\Character{61}}%
\TextAt{13}{154.649}{118.582}{\Character{39}}%
\TextAt{3}{162.317}{122.92}{\Character{0}}%
\TextAt{21}{168.904}{122.92}{n}%
\TextAt{24}{174.042}{121.813}{2}%
\TextAt{13}{178.691}{118.582}{\Character{27}}%
\TextAt{17}{185.345}{116.789}{1}%
\TextAt{13}{190.077}{118.582}{\Character{39}}%
\TextAt{21}{197.745}{122.92}{n}%
\TextAt{24}{202.884}{121.813}{2}%
\TextAt{8}{207.533}{118.582}{,}%
\TextAt{13}{216.158}{118.582}{n}%
\TextAt{17}{223.145}{116.789}{4}%
\TextAt{9}{233.21}{118.582}{\Character{61}}%
\TextAt{13}{247.648}{118.582}{d}%
\TextAt{21}{253.73}{116.789}{K}%
\TextAt{9}{261.933}{118.582}{\Character{40}}%
\TextAt{13}{266.485}{118.582}{\Character{27}}%
\TextAt{17}{273.139}{116.789}{4}%
\TextAt{9}{277.871}{118.582}{\Character{41}}%
\TextAt{9}{287.746}{118.582}{\Character{61}}%
\TextAt{13}{302.193}{118.582}{n}%
\TextAt{17}{309.181}{116.789}{1}%
\TextAt{8}{313.913}{118.582}{,}%
\TextAt{9}{322.538}{118.582}{\Character{6}}%
\TextAt{17}{330.993}{116.789}{4}%
\TextAt{9}{341.058}{118.582}{\Character{61}}%
\TextAt{9}{355.488}{118.582}{\Character{6}}%
\TextAt{21}{363.949}{124.25}{\Character{39}}%
\TextAt{22}{369.468}{127.062}{n}%
\TextAt{24}{374.068}{125.269}{2}%
\TextAt{17}{363.949}{115.651}{1}%
\TextAt{8}{379.216}{118.582}{,}%
\TextAt{13}{387.841}{118.582}{\Character{25}}%
\TextAt{17}{394.481}{116.789}{4}%
\TextAt{9}{404.546}{118.582}{\Character{61}}%
\TextAt{13}{418.983}{118.582}{\Character{25}}%
\TextAt{21}{426.053}{124.25}{\Character{39}}%
\TextAt{22}{431.571}{127.062}{n}%
\TextAt{24}{436.172}{125.269}{2}%
\TextAt{17}{425.624}{115.651}{1}%
\NativeMathAt{446.652}{118.582}{11.955}{\in}%
\TextAt{13}{459.954}{118.582}{A}%
\TextAt{17}{468.73}{116.789}{\Character{6}}%
\TextAt{24}{474.846}{115.682}{4}%
\TextAt{8}{484.573}{118.582}{and}%
\TextAt{13}{508.516}{118.582}{\Character{28}}%
\TextAt{17}{513.6}{116.789}{4}%
\TextAt{9}{523.665}{118.582}{\Character{61}}%
\TextAt{13}{79.226}{101.536}{\Character{27}}%
\TextAt{3}{86.309}{106.484}{\Character{0}}%
\TextAt{17}{92.895}{106.484}{1}%
\TextAt{17}{85.88}{98.605}{4}%
\TextAt{13}{97.627}{101.536}{\Character{39}}%
\TextAt{21}{105.296}{105.874}{n}%
\TextAt{24}{110.434}{104.767}{4}%
\TextAt{8}{115.083}{101.536}{,}%
\TextAt{8}{122.232}{101.536}{we}%
\TextAt{8}{139.46}{101.536}{find}%
\TextAt{13}{162.868}{101.536}{\Character{28}}%
\TextAt{17}{167.952}{99.743}{3}%
\TextAt{9}{176.005}{101.536}{\Character{61}}%
\TextAt{13}{188.43}{101.536}{\Character{28}}%
\TextAt{17}{193.514}{99.743}{2}%
\TextAt{13}{198.246}{101.536}{\Character{28}}%
\TextAt{17}{203.33}{99.743}{4}%
\TextAt{8}{211.96}{101.536}{and}%
\TextAt{13}{245.826}{72.921}{N}%
\TextAt{17}{255.199}{71.062}{\Character{6}}%
\TextAt{24}{261.315}{69.955}{4}%
\TextAt{3}{265.467}{71.062}{\Character{106}}%
\TextAt{21}{267.819}{71.062}{K}%
\TextAt{9}{276.022}{72.921}{\Character{40}}%
\TextAt{13}{280.574}{72.921}{\Character{25}}%
\TextAt{17}{287.214}{71.128}{4}%
\TextAt{9}{291.946}{72.921}{\Character{41}}%
\TextAt{9}{299.801}{72.921}{\Character{61}}%
\TextAt{13}{312.245}{72.921}{N}%
\TextAt{17}{321.619}{71.062}{\Character{6}}%
\TextAt{24}{327.735}{69.955}{1}%
\TextAt{3}{331.886}{71.062}{\Character{106}}%
\TextAt{21}{334.238}{71.062}{K}%
\TextAt{9}{342.441}{72.921}{\Character{40}}%
\TextAt{13}{346.993}{72.921}{\Character{25}}%
\TextAt{17}{353.634}{71.128}{1}%
\TextAt{9}{358.366}{72.921}{\Character{41}}%
\TextAt{13}{362.918}{72.921}{\Character{58}}%
\TextAt{8}{297.211}{43.033}{268}%
\Destination{78.23}{809.815}{nav.page.268}%
\Destination{79.36}{701.655}{nav.print.293}%
\Destination{79.36}{414.555}{nav.number.4.5.5}%
\Hotlink{504.94}{527.455}{25.078}{11.047}{nav.number.4.5.3}%
\Hotlink{279.061}{758.126}{21.722}{11.047}{nav.number.4.5.1}%
\Hotlink{248.853}{683.633}{21.722}{11.047}{nav.number.4.5.1}%
}

% Page 281 (body page 269)
\Page{%
\TextAt{2}{79.226}{796.943}{IV.}%
\TextAt{2}{98.012}{796.943}{Abstract}%
\TextAt{2}{143.252}{796.943}{Class}%
\TextAt{2}{171.844}{796.943}{Field}%
\TextAt{2}{199.575}{796.943}{Theory}%
\Vector{q 0 G 0 g 0.299 w 7.226 18.729 m 460.769 18.729 l S Q}%
\TextAt{8}{79.226}{761.124}{We}%
\TextAt{8}{99.371}{761.124}{may}%
\TextAt{8}{124.739}{761.124}{therefore}%
\TextAt{8}{173.839}{761.124}{pass}%
\TextAt{8}{199.328}{761.124}{to}%
\TextAt{8}{213.638}{761.124}{the}%
\TextAt{8}{233.795}{761.124}{congruence}%
\TextAt{13}{176.451}{732.463}{N}%
\TextAt{17}{185.825}{730.603}{\Character{6}}%
\TextAt{24}{191.941}{729.496}{3}%
\TextAt{3}{196.092}{730.603}{\Character{106}}%
\TextAt{21}{198.444}{730.603}{K}%
\TextAt{9}{206.647}{732.463}{\Character{40}}%
\TextAt{13}{211.199}{732.463}{\Character{25}}%
\TextAt{17}{217.84}{730.67}{3}%
\TextAt{9}{222.572}{732.463}{\Character{41}}%
\NativeMathAt{230.445}{732.463}{11.955}{\equiv}%
\TextAt{13}{243.064}{732.463}{N}%
\TextAt{17}{252.438}{730.603}{\Character{6}}%
\TextAt{24}{258.554}{729.496}{2}%
\TextAt{3}{262.705}{730.603}{\Character{106}}%
\TextAt{21}{265.058}{730.603}{K}%
\TextAt{9}{273.26}{732.463}{\Character{40}}%
\TextAt{13}{277.813}{732.463}{\Character{25}}%
\TextAt{17}{284.453}{730.67}{2}%
\TextAt{9}{289.185}{732.463}{\Character{41}}%
\TextAt{13}{293.738}{732.463}{N}%
\TextAt{17}{303.111}{730.603}{\Character{6}}%
\TextAt{24}{309.228}{729.496}{4}%
\TextAt{3}{313.379}{730.603}{\Character{106}}%
\TextAt{21}{315.731}{730.603}{K}%
\TextAt{9}{323.934}{732.463}{\Character{40}}%
\TextAt{13}{328.486}{732.463}{\Character{25}}%
\TextAt{17}{335.126}{730.67}{4}%
\TextAt{9}{339.858}{732.463}{\Character{41}}%
\TextAt{9}{356.356}{732.463}{\Character{40}}%
\TextAt{8}{360.918}{732.463}{mod}%
\TextAt{13}{387.02}{732.463}{N}%
\TextAt{30}{397.156}{730.825}{\Character{101}}%
\TextAt{21}{396.394}{728.81}{L}%
\TextAt{3}{402.153}{728.81}{\Character{106}}%
\TextAt{21}{404.505}{728.81}{K}%
\TextAt{13}{412.708}{732.463}{A}%
\TextAt{30}{422.246}{730.825}{\Character{101}}%
\TextAt{21}{421.483}{728.81}{L}%
\TextAt{9}{427.741}{732.463}{\Character{41}}%
\TextAt{13}{432.293}{732.463}{\Character{59}}%
\TextAt{8}{79.226}{703.802}{the}%
\TextAt{8}{99.341}{703.802}{proof}%
\TextAt{8}{129.869}{703.802}{of}%
\TextAt{8}{143.158}{703.802}{which}%
\TextAt{8}{176.603}{703.802}{uses}%
\TextAt{8}{201.392}{703.802}{the}%
\TextAt{8}{221.507}{703.802}{identity}%
\TextAt{13}{264.72}{703.802}{\Character{28}}%
\TextAt{17}{269.803}{702.009}{3}%
\TextAt{9}{277.856}{703.802}{\Character{61}}%
\TextAt{13}{290.282}{703.802}{\Character{28}}%
\TextAt{17}{295.366}{702.009}{2}%
\TextAt{13}{300.098}{703.802}{\Character{28}}%
\TextAt{17}{305.182}{702.009}{4}%
\TextAt{8}{309.914}{703.802}{.}%
\TextAt{8}{318.348}{703.802}{From}%
\TextAt{8}{349.032}{703.802}{(5.3),}%
\TextAt{8}{380.226}{703.802}{we}%
\TextAt{8}{397.412}{703.802}{have}%
\TextAt{13}{424.365}{703.802}{N}%
\TextAt{17}{433.739}{701.942}{\Character{6}}%
\TextAt{22}{439.855}{700.727}{i}%
\TextAt{3}{443.017}{701.942}{\Character{106}}%
\TextAt{21}{445.369}{701.942}{K}%
\TextAt{9}{453.572}{703.802}{\Character{40}}%
\TextAt{13}{458.124}{703.802}{\Character{25}}%
\TextAt{21}{464.764}{702.009}{i}%
\TextAt{9}{468.145}{703.802}{\Character{41}}%
\TextAt{9}{476}{703.802}{\Character{61}}%
\TextAt{13}{488.444}{703.802}{N}%
\TextAt{9}{499.067}{703.802}{\Character{40}}%
\TextAt{13}{503.619}{703.802}{\Character{25}}%
\TextAt{21}{510.688}{710.818}{\Character{39}}%
\TextAt{22}{516.207}{709.821}{n}%
\TextAt{22}{520.807}{707.92}{i}%
\TextAt{21}{510.259}{700.738}{i}%
\TextAt{9}{524.965}{703.802}{\Character{41}}%
\TextAt{8}{529.518}{703.802}{.}%
\TextAt{8}{79.226}{686.756}{Thus,}%
\TextAt{8}{112.127}{686.756}{if}%
\TextAt{8}{122.851}{686.756}{we}%
\TextAt{8}{140.078}{686.756}{put}%
\TextAt{13}{253.078}{669.71}{u}%
\TextAt{9}{263.061}{669.71}{\Character{61}}%
\TextAt{13}{275.486}{669.71}{\Character{25}}%
\TextAt{21}{282.556}{676.618}{\Character{39}}%
\TextAt{22}{288.074}{675.621}{n}%
\TextAt{24}{292.675}{673.828}{3}%
\TextAt{17}{282.127}{666.779}{3}%
\TextAt{13}{297.822}{669.71}{\Character{25}}%
\TextAt{3}{304.891}{676.618}{\Character{0}}%
\TextAt{21}{311.478}{676.618}{\Character{39}}%
\TextAt{22}{316.996}{675.621}{n}%
\TextAt{24}{321.597}{673.828}{4}%
\TextAt{17}{304.462}{666.779}{4}%
\TextAt{13}{326.744}{669.71}{\Character{25}}%
\TextAt{3}{333.813}{676.618}{\Character{0}}%
\TextAt{21}{340.4}{676.618}{\Character{39}}%
\TextAt{22}{345.918}{675.621}{n}%
\TextAt{24}{350.519}{673.828}{2}%
\TextAt{17}{333.384}{666.779}{2}%
\TextAt{13}{355.666}{669.71}{\Character{59}}%
\TextAt{8}{79.226}{646.334}{then}%
\TextAt{8}{106.05}{646.334}{the}%
\TextAt{8}{126.369}{646.334}{congruence}%
\TextAt{8}{187.014}{646.334}{amounts}%
\TextAt{8}{234.4}{646.334}{simply}%
\TextAt{8}{272.019}{646.334}{to}%
\TextAt{8}{286.493}{646.334}{the}%
\TextAt{8}{306.813}{646.334}{relation}%
\TextAt{13}{349.907}{646.334}{N}%
\TextAt{9}{360.53}{646.334}{\Character{40}}%
\TextAt{13}{365.082}{646.334}{u}%
\TextAt{9}{371.744}{646.334}{\Character{41}}%
\NativeMathAt{379.896}{646.334}{11.955}{\in}%
\TextAt{13}{391.465}{646.334}{N}%
\TextAt{30}{401.602}{644.696}{\Character{101}}%
\TextAt{21}{400.839}{642.681}{L}%
\TextAt{3}{406.598}{642.681}{\Character{106}}%
\TextAt{21}{408.951}{642.681}{K}%
\TextAt{13}{417.153}{646.334}{A}%
\TextAt{30}{426.691}{644.696}{\Character{101}}%
\TextAt{21}{425.929}{642.681}{L}%
\TextAt{8}{432.186}{646.334}{.}%
\TextAt{8}{441.124}{646.334}{For}%
\TextAt{8}{462.257}{646.334}{this,}%
\TextAt{8}{488.535}{646.334}{however,}%
\TextAt{8}{79.226}{629.288}{lemma}%
\TextAt{8}{116.945}{629.288}{(5.4)}%
\TextAt{8}{144.92}{629.288}{gives}%
\TextAt{8}{173.6}{629.288}{us}%
\TextAt{8}{188.616}{629.288}{all}%
\TextAt{8}{204.875}{629.288}{that}%
\TextAt{8}{230.244}{629.288}{we}%
\TextAt{8}{247.471}{629.288}{need.}%
\TextAt{8}{97.159}{608.107}{Since}%
\TextAt{13}{127.716}{608.107}{\Character{39}}%
\TextAt{21}{135.385}{606.314}{n}%
\TextAt{22}{140.523}{605.099}{i}%
\TextAt{16}{146.839}{608.107}{\Character{14}}%
\TextAt{9}{155.474}{608.107}{\Character{40}}%
\TextAt{13}{160.026}{608.107}{\Character{39}}%
\TextAt{16}{170.351}{608.107}{\Character{0}}%
\TextAt{9}{182.306}{608.107}{1\Character{41}}%
\TextAt{9}{196.007}{608.107}{\Character{61}}%
\TextAt{13}{208.458}{608.107}{\Character{39}}%
\TextAt{21}{216.126}{612.446}{n}%
\TextAt{22}{221.264}{611.231}{i}%
\TextAt{16}{227.581}{608.107}{\Character{0}}%
\TextAt{9}{239.536}{608.107}{1}%
\TextAt{8}{249.287}{608.107}{and}%
\TextAt{13}{272.049}{608.107}{\Character{25}}%
\TextAt{21}{279.118}{613.775}{\Character{39}}%
\TextAt{22}{284.637}{616.588}{n}%
\TextAt{22}{289.238}{614.686}{i}%
\TextAt{3}{292.897}{613.775}{\Character{0}}%
\TextAt{17}{299.484}{613.775}{1}%
\TextAt{21}{278.69}{605.043}{i}%
\TextAt{9}{307.537}{608.107}{\Character{61}}%
\TextAt{13}{319.963}{608.107}{\Character{25}}%
\TextAt{21}{327.032}{614.875}{\Character{27}}%
\TextAt{23}{332.157}{618.45}{\Character{0}}%
\TextAt{24}{337.913}{618.45}{1}%
\TextAt{22}{331.871}{612.225}{i}%
\TextAt{21}{342.064}{614.875}{\Character{39}}%
\TextAt{22}{347.583}{617.688}{n}%
\TextAt{22}{352.183}{615.786}{i}%
\TextAt{3}{355.843}{614.875}{\Character{0}}%
\TextAt{17}{362.43}{614.875}{1}%
\TextAt{21}{326.603}{605.043}{i}%
\TextAt{9}{370.483}{608.107}{\Character{61}}%
\TextAt{13}{382.908}{608.107}{\Character{25}}%
\TextAt{21}{389.978}{613.44}{\Character{28}}%
\TextAt{22}{393.736}{612.225}{i}%
\TextAt{3}{396.898}{613.44}{\Character{0}}%
\TextAt{17}{403.484}{613.44}{1}%
\TextAt{21}{389.549}{605.043}{i}%
\TextAt{8}{408.216}{608.107}{,}%
\TextAt{8}{415.366}{608.107}{we}%
\TextAt{8}{432.593}{608.107}{have}%
\TextAt{13}{245.303}{579.447}{u}%
\TextAt{21}{251.966}{584.383}{\Character{39}}%
\TextAt{3}{257.484}{584.383}{\Character{0}}%
\TextAt{17}{264.07}{584.383}{1}%
\TextAt{9}{272.123}{579.446}{\Character{61}}%
\TextAt{13}{284.549}{579.447}{\Character{25}}%
\TextAt{21}{291.618}{584.671}{\Character{28}}%
\TextAt{24}{295.376}{583.564}{3}%
\TextAt{3}{299.527}{584.671}{\Character{0}}%
\TextAt{17}{306.114}{584.671}{1}%
\TextAt{17}{291.189}{576.515}{3}%
\TextAt{13}{310.846}{579.446}{\Character{25}}%
\TextAt{17}{317.916}{584.671}{1}%
\TextAt{3}{322.15}{584.671}{\Character{0}}%
\TextAt{21}{328.736}{584.671}{\Character{28}}%
\TextAt{24}{332.494}{583.564}{4}%
\TextAt{17}{317.487}{576.515}{4}%
\TextAt{13}{337.144}{579.446}{\Character{25}}%
\TextAt{17}{344.213}{584.671}{1}%
\TextAt{3}{348.447}{584.671}{\Character{0}}%
\TextAt{21}{355.034}{584.671}{\Character{28}}%
\TextAt{24}{358.792}{583.564}{2}%
\TextAt{17}{343.784}{576.515}{2}%
\TextAt{13}{363.441}{579.446}{\Character{58}}%
\TextAt{8}{79.226}{550.786}{From}%
\TextAt{8}{109.578}{550.786}{the}%
\TextAt{8}{129.362}{550.786}{identity}%
\TextAt{13}{172.243}{550.786}{\Character{28}}%
\TextAt{17}{177.327}{548.992}{3}%
\TextAt{9}{185.38}{550.786}{\Character{61}}%
\TextAt{13}{197.805}{550.786}{\Character{28}}%
\TextAt{17}{202.889}{548.992}{2}%
\TextAt{13}{207.621}{550.786}{\Character{28}}%
\TextAt{17}{212.705}{548.992}{4}%
\TextAt{8}{217.437}{550.786}{,}%
\TextAt{8}{224.288}{550.786}{it}%
\TextAt{8}{235.62}{550.786}{follows}%
\TextAt{8}{273.683}{550.786}{that}%
\TextAt{9}{298.679}{550.786}{\Character{40}}%
\TextAt{13}{303.231}{550.786}{\Character{28}}%
\TextAt{17}{308.315}{548.992}{3}%
\TextAt{16}{314.942}{550.786}{\Character{0}}%
\TextAt{9}{326.135}{550.786}{1\Character{41}}%
\TextAt{9}{338.413}{550.786}{\Character{43}}%
\TextAt{9}{349.4}{550.786}{\Character{40}1}%
\TextAt{16}{361.733}{550.786}{\Character{0}}%
\TextAt{13}{372.926}{550.786}{\Character{28}}%
\TextAt{17}{378.01}{548.992}{2}%
\TextAt{9}{382.742}{550.786}{\Character{41}}%
\TextAt{9}{389.174}{550.786}{\Character{43}}%
\TextAt{9}{400.161}{550.786}{\Character{40}1}%
\TextAt{16}{412.488}{550.786}{\Character{0}}%
\TextAt{13}{423.681}{550.786}{\Character{28}}%
\TextAt{17}{428.765}{548.992}{4}%
\TextAt{9}{433.497}{550.786}{\Character{41}}%
\TextAt{9}{441.352}{550.786}{\Character{61}}%
\TextAt{9}{453.773}{550.786}{\Character{40}1}%
\TextAt{16}{466.095}{550.786}{\Character{0}}%
\TextAt{13}{477.288}{550.786}{\Character{28}}%
\TextAt{17}{482.372}{548.992}{2}%
\TextAt{9}{487.104}{550.786}{\Character{41}\Character{40}1}%
\TextAt{16}{503.956}{550.786}{\Character{0}}%
\TextAt{13}{515.149}{550.786}{\Character{28}}%
\TextAt{17}{520.233}{548.992}{4}%
\TextAt{9}{524.965}{550.786}{\Character{41}}%
\TextAt{8}{529.518}{550.786}{.}%
\TextAt{8}{79.226}{533.74}{Putting}%
\TextAt{8}{122.313}{533.74}{now}%
\TextAt{13}{158.982}{505.079}{\Character{25}}%
\TextAt{17}{165.622}{503.286}{3}%
\TextAt{9}{173.675}{505.079}{\Character{61}}%
\TextAt{13}{186.101}{505.079}{u}%
\TextAt{17}{192.763}{503.286}{3}%
\TextAt{13}{197.496}{505.079}{\Character{25}}%
\TextAt{17}{204.136}{503.286}{4}%
\TextAt{13}{208.868}{505.079}{\Character{59}}%
\TextAt{13}{238.003}{505.079}{\Character{25}}%
\TextAt{17}{244.663}{503.286}{2}%
\TextAt{9}{252.716}{505.079}{\Character{61}}%
\TextAt{13}{265.142}{505.079}{u}%
\TextAt{3}{271.804}{510.015}{\Character{0}}%
\TextAt{17}{278.391}{510.015}{1}%
\TextAt{17}{271.804}{502.123}{2}%
\TextAt{13}{283.123}{505.079}{\Character{25}}%


\TextAt{17}{289.763}{502.148}{4}%
\TextAt{13}{298.6}{505.079}{\Character{59}}%
\TextAt{13}{327.735}{505.079}{\Character{25}}%
\TextAt{21}{334.823}{510.304}{\Character{28}}%
\TextAt{24}{338.582}{509.197}{2}%
\TextAt{17}{334.394}{502.148}{4}%
\TextAt{9}{346.552}{505.079}{\Character{61}}%
\TextAt{13}{358.977}{505.079}{u}%
\TextAt{17}{365.64}{503.286}{4}%
\TextAt{13}{370.372}{505.079}{\Character{25}}%
\TextAt{17}{377.012}{503.286}{4}%
\TextAt{13}{381.744}{505.079}{\Character{59}}%
\TextAt{13}{410.879}{505.079}{u}%
\TextAt{21}{417.561}{503.286}{i}%
\NativeMathAt{424.264}{505.079}{11.955}{\in}%
\TextAt{13}{435.555}{505.079}{U}%
\TextAt{30}{444.267}{503.441}{\Character{101}}%
\TextAt{21}{443.505}{501.426}{L}%
\TextAt{13}{449.762}{505.079}{\Character{59}}%
\TextAt{8}{79.226}{476.418}{we}%
\TextAt{8}{96.454}{476.418}{obtain}%
\TextAt{13}{263.664}{454.026}{u}%
\TextAt{21}{270.327}{458.963}{\Character{39}}%
\TextAt{3}{275.845}{458.963}{\Character{0}}%
\TextAt{17}{282.431}{458.963}{1}%
\TextAt{9}{290.485}{454.026}{\Character{61}}%
\TextAt{17}{308.431}{468.97}{4}%
\TextAt{20}{302.91}{465.384}{\Character{89}}%
\TextAt{21}{303.696}{440.19}{i}%
\TextAt{17}{306.579}{440.19}{\Character{61}2}%
\TextAt{13}{320.179}{454.026}{u}%
\TextAt{21}{326.841}{459.359}{\Character{28}}%
\TextAt{22}{330.599}{458.144}{i}%
\TextAt{3}{333.761}{459.359}{\Character{0}}%
\TextAt{17}{340.347}{459.359}{1}%
\TextAt{21}{326.841}{450.963}{i}%
\TextAt{13}{345.08}{454.026}{\Character{58}}%
\TextAt{8}{79.226}{420.012}{For}%
\TextAt{8}{99.608}{420.012}{the}%
\TextAt{8}{119.176}{420.012}{element}%
\TextAt{13}{161.83}{420.012}{x}%
\TextAt{9}{171.803}{420.012}{\Character{61}}%
\TextAt{13}{184.228}{420.012}{u}%
\TextAt{8}{194.211}{420.012}{mod}%
\TextAt{13}{219.649}{420.012}{I}%
\TextAt{21}{224.812}{416.359}{G}%
\TextAt{17}{231.411}{416.359}{\Character{40}}%
\TextAt{30}{235.467}{418.373}{\Character{101}}%
\TextAt{21}{234.704}{416.359}{L}%
\TextAt{3}{240.463}{416.359}{\Character{106}}%
\TextAt{30}{244.786}{418.373}{\Character{101}}%
\TextAt{21}{242.816}{416.359}{K}%
\TextAt{17}{250.52}{416.359}{\Character{41}}%
\TextAt{13}{254.311}{420.012}{U}%
\TextAt{30}{263.024}{418.373}{\Character{101}}%
\TextAt{21}{262.262}{416.359}{L}%
\NativeMathAt{271.84}{420.012}{11.955}{\in}%
\TextAt{13}{283.131}{420.012}{H}%
\TextAt{17}{292.83}{418.218}{0}%
\TextAt{9}{297.562}{420.012}{\Character{40}}%
\TextAt{13}{302.114}{420.012}{G}%
\TextAt{9}{311.348}{420.012}{\Character{40}}%
\TextAt{20}{316.887}{423.034}{\Character{101}}%
\TextAt{13}{315.9}{420.012}{L}%
\TextAt{16}{327.186}{420.012}{\Character{106}}%
\TextAt{20}{336.566}{423.034}{\Character{101}}%
\TextAt{13}{333.827}{420.012}{K}%
\TextAt{9}{344.645}{420.012}{\Character{41}}%
\TextAt{13}{349.197}{420.012}{\Character{59}}%
\TextAt{13}{354.421}{420.012}{U}%
\TextAt{30}{363.154}{418.373}{\Character{101}}%
\TextAt{21}{362.392}{416.359}{L}%
\TextAt{9}{368.649}{420.012}{\Character{41}}%
\TextAt{8}{373.201}{420.012}{,}%
\TextAt{8}{379.88}{420.012}{this}%
\TextAt{8}{402.114}{420.012}{means}%
\TextAt{8}{437.356}{420.012}{that}%
\TextAt{13}{462.137}{420.012}{x}%
\TextAt{21}{468.789}{424.35}{\Character{39}}%
\TextAt{3}{474.307}{424.35}{\Character{0}}%
\TextAt{17}{480.894}{424.35}{1}%
\TextAt{9}{488.947}{420.012}{\Character{61}}%
\TextAt{9}{501.357}{420.012}{1}%
\TextAt{8}{507.225}{420.012}{,}%
\TextAt{8}{513.904}{420.012}{and}%
\TextAt{8}{79.226}{402.966}{so}%
\TextAt{13}{93.596}{402.966}{x}%
\TextAt{21}{100.248}{407.304}{\Character{39}}%
\TextAt{9}{109.586}{402.966}{\Character{61}}%
\TextAt{13}{122.011}{402.966}{x}%
\TextAt{8}{128.663}{402.966}{;}%
\TextAt{8}{135.812}{402.966}{then}%
\TextAt{8}{162.472}{402.966}{by}%
\TextAt{8}{178.731}{402.966}{(5.4),}%
\TextAt{8}{209.958}{402.966}{we}%
\TextAt{8}{227.186}{402.966}{do}%
\TextAt{8}{243.445}{402.966}{get}%
\TextAt{13}{262.956}{402.966}{N}%
\TextAt{9}{273.578}{402.966}{\Character{40}}%
\TextAt{13}{278.131}{402.966}{u}%
\TextAt{9}{284.793}{402.966}{\Character{41}}%
\TextAt{9}{292.648}{402.966}{\Character{61}}%
\TextAt{13}{305.092}{402.966}{N}%
\TextAt{9}{315.714}{402.966}{\Character{40}}%
\TextAt{13}{320.267}{402.966}{x}%
\TextAt{9}{326.919}{402.966}{\Character{41}}%
\NativeMathAt{334.792}{402.966}{11.955}{\in}%
\TextAt{13}{346.083}{402.966}{N}%
\TextAt{30}{356.219}{401.327}{\Character{101}}%
\TextAt{21}{355.457}{399.313}{L}%
\TextAt{3}{361.216}{399.313}{\Character{106}}%
\TextAt{21}{363.568}{399.313}{K}%
\TextAt{13}{371.771}{402.966}{A}%
\TextAt{30}{381.308}{401.327}{\Character{101}}%
\TextAt{21}{380.546}{399.313}{L}%
\TextAt{8}{386.803}{402.966}{.}%
\Vector{q 0 G 0 g 0.399 w 452.334 -371.034 m 452.334 -362.965 l S Q}%
\Vector{q 0 G 0 g 0.399 w 452.534 -363.164 m 459.707 -363.164 l S Q}%
\Vector{q 0 G 0 g 0.399 w 452.534 -370.835 m 459.707 -370.835 l S Q}%
\Vector{q 0 G 0 g 0.399 w 459.906 -371.034 m 459.906 -362.965 l S Q}%
\TextAt{8}{97.159}{374.062}{From}%
\TextAt{8}{127.884}{374.062}{the}%
\TextAt{8}{148.04}{374.062}{surjectivity}%
\TextAt{8}{209.239}{374.062}{of}%
\TextAt{8}{222.569}{374.062}{the}%
\TextAt{8}{242.725}{374.062}{mapping}%
\TextAt{8}{238.302}{345.401}{Frob}%
\TextAt{9}{262.858}{345.401}{\Character{40}}%
\TextAt{20}{268.397}{348.423}{\Character{101}}%
\TextAt{13}{267.41}{345.401}{L}%
\TextAt{16}{278.696}{345.401}{\Character{106}}%
\TextAt{13}{285.338}{345.401}{K}%
\TextAt{9}{296.155}{345.401}{\Character{41}}%
\NativeArrowAt{304.028}{345.401}{11.955}{19.26}{\longrightarrow}%
\TextAt{13}{326.61}{345.401}{G}%
\TextAt{9}{335.844}{345.401}{\Character{40}}%
\TextAt{13}{340.396}{345.401}{L}%
\TextAt{16}{351.682}{345.401}{\Character{106}}%
\TextAt{13}{358.323}{345.401}{K}%
\TextAt{9}{369.141}{345.401}{\Character{41}}%
\TextAt{8}{79.226}{316.74}{and}%
\TextAt{8}{101.989}{316.74}{the}%
\TextAt{8}{122.145}{316.74}{fact}%
\TextAt{8}{145.231}{316.74}{that}%
\TextAt{13}{170.6}{316.74}{N}%
\TextAt{30}{180.736}{315.102}{\Character{101}}%
\TextAt{21}{179.973}{313.087}{L}%
\TextAt{3}{185.732}{313.087}{\Character{106}}%
\TextAt{21}{188.085}{313.087}{K}%
\TextAt{13}{196.287}{316.74}{A}%
\TextAt{30}{205.825}{315.102}{\Character{101}}%
\TextAt{21}{205.063}{313.087}{L}%
\NativeMathAt{214.641}{316.74}{11.955}{\subseteq}%
\TextAt{13}{227.26}{316.74}{N}%
\TextAt{21}{236.634}{314.88}{L}%
\TextAt{3}{242.393}{314.88}{\Character{106}}%
\TextAt{21}{244.745}{314.88}{K}%
\TextAt{13}{252.948}{316.74}{A}%
\TextAt{21}{261.723}{314.947}{L}%
\TextAt{8}{267.98}{316.74}{,}%
\TextAt{8}{275.13}{316.74}{we}%
\TextAt{8}{292.357}{316.74}{now}%
\TextAt{8}{316.746}{316.74}{have}%
\TextAt{8}{343.74}{316.74}{the}%
\TextAt{12}{79.226}{290.825}{Proposition}%
\TextAt{12}{152.336}{290.825}{(5.6).}%
\TextAt{8}{187.999}{290.825}{For}%
\TextAt{8}{209.315}{290.825}{every}%
\TextAt{8}{240.554}{290.825}{finite}%
\TextAt{8}{270.813}{290.825}{Galois}%
\TextAt{8}{307.073}{290.825}{extension}%
\TextAt{13}{359.186}{290.825}{L}%
\TextAt{16}{371.062}{290.825}{\Character{106}}%
\TextAt{13}{378.295}{290.825}{K}%
\TextAt{8}{389.112}{290.825}{,}%
\TextAt{8}{396.695}{290.825}{there}%
\TextAt{8}{426.954}{290.825}{is}%
\TextAt{8}{439.065}{290.825}{a}%
\TextAt{8}{449.167}{290.825}{canonical}%
\TextAt{8}{500.897}{290.825}{homo-}%
\TextAt{8}{79.226}{273.779}{morphism}%
\TextAt{13}{221.988}{256.733}{r}%
\TextAt{21}{227.263}{254.873}{L}%
\TextAt{3}{233.022}{254.873}{\Character{106}}%
\TextAt{21}{235.375}{254.873}{K}%
\TextAt{9}{246.898}{256.733}{\Character{58}}%
\TextAt{13}{253.471}{256.733}{G}%
\TextAt{9}{262.704}{256.733}{\Character{40}}%
\TextAt{13}{267.256}{256.733}{L}%
\TextAt{16}{278.542}{256.733}{\Character{106}}%
\TextAt{13}{285.184}{256.733}{K}%
\TextAt{9}{296.001}{256.733}{\Character{41}}%
\NativeArrowAt{303.874}{256.733}{11.955}{19.26}{\longrightarrow}%
\TextAt{13}{326.456}{256.733}{A}%
\TextAt{21}{335.232}{254.94}{K}%
\TextAt{9}{343.434}{256.733}{\Character{47}}%
\TextAt{13}{349.287}{256.733}{N}%
\TextAt{21}{358.661}{254.873}{L}%
\TextAt{3}{364.42}{254.873}{\Character{106}}%
\TextAt{21}{366.773}{254.873}{K}%
\TextAt{13}{374.975}{256.733}{A}%
\TextAt{21}{383.751}{254.94}{L}%
\TextAt{8}{79.226}{233.357}{given}%
\TextAt{8}{109.796}{233.357}{by}%
\TextAt{13}{218.669}{216.311}{r}%
\TextAt{21}{223.944}{214.452}{L}%
\TextAt{3}{229.703}{214.452}{\Character{106}}%
\TextAt{21}{232.056}{214.452}{K}%
\TextAt{9}{240.258}{216.311}{\Character{40}}%
\TextAt{13}{244.811}{216.311}{\Character{27}}%
\TextAt{9}{251.893}{216.311}{\Character{41}}%
\TextAt{9}{259.748}{216.311}{\Character{61}}%
\TextAt{13}{272.192}{216.311}{N}%
\TextAt{17}{281.566}{214.452}{\Character{6}}%
\TextAt{3}{287.682}{214.452}{\Character{106}}%
\TextAt{21}{290.034}{214.452}{K}%
\TextAt{9}{298.237}{216.311}{\Character{40}}%
\TextAt{13}{302.789}{216.311}{\Character{25}}%
\TextAt{17}{309.429}{214.518}{\Character{6}}%
\TextAt{9}{316.043}{216.311}{\Character{41}}%
\TextAt{8}{323.917}{216.311}{mod}%
\TextAt{13}{349.355}{216.311}{N}%
\TextAt{21}{358.728}{214.452}{L}%
\TextAt{3}{364.487}{214.452}{\Character{106}}%
\TextAt{21}{366.84}{214.452}{K}%
\TextAt{13}{375.042}{216.311}{A}%
\TextAt{21}{383.818}{214.518}{L}%
\TextAt{13}{390.075}{216.311}{\Character{59}}%
\TextAt{8}{79.226}{192.935}{where}%
\TextAt{9}{113.323}{192.935}{\Character{6}}%
\TextAt{8}{125.962}{192.935}{is}%
\TextAt{8}{138.013}{192.935}{the}%
\TextAt{8}{158.457}{192.935}{fixed}%
\TextAt{8}{187.03}{192.935}{field}%
\TextAt{8}{212.675}{192.935}{of}%
\TextAt{8}{226.292}{192.935}{a}%
\TextAt{8}{236.335}{192.935}{preimage}%
\TextAt{9}{287.317}{192.935}{\Character{22}}%
\TextAt{13}{286.702}{192.935}{\Character{27}}%
\NativeMathAt{297.595}{192.935}{11.955}{\in}%
\TextAt{8}{309.376}{192.935}{Frob}%
\TextAt{9}{333.932}{192.935}{\Character{40}}%
\TextAt{20}{339.471}{195.957}{\Character{101}}%
\TextAt{13}{338.484}{192.935}{L}%
\TextAt{16}{350.259}{192.935}{\Character{106}}%
\TextAt{13}{357.391}{192.935}{K}%
\TextAt{9}{368.209}{192.935}{\Character{41}}%
\TextAt{8}{376.946}{192.935}{of}%
\TextAt{13}{390.563}{192.935}{\Character{27}}%
\NativeMathAt{401.456}{192.935}{11.955}{\in}%
\TextAt{13}{413.237}{192.935}{G}%
\TextAt{9}{422.47}{192.935}{\Character{40}}%
\TextAt{13}{427.023}{192.935}{L}%
\TextAt{16}{438.798}{192.935}{\Character{106}}%
\TextAt{13}{445.929}{192.935}{K}%
\TextAt{9}{456.747}{192.935}{\Character{41}}%
\TextAt{8}{465.484}{192.935}{and}%
\TextAt{13}{488.534}{192.935}{\Character{25}}%
\TextAt{17}{495.175}{191.142}{\Character{6}}%
\NativeMathAt{505.599}{192.935}{11.955}{\in}%
\TextAt{13}{517.38}{192.935}{A}%
\TextAt{17}{526.155}{191.142}{\Character{6}}%
\TextAt{8}{79.226}{175.889}{is}%
\TextAt{8}{90.99}{175.889}{a}%
\TextAt{8}{100.745}{175.889}{prime}%
\TextAt{8}{133.909}{175.889}{element.}%
\TextAt{8}{181.702}{175.889}{It}%
\TextAt{8}{194.374}{175.889}{is}%
\TextAt{8}{206.138}{175.889}{called}%
\TextAt{8}{239.302}{175.889}{the}%
\TextAt{12}{259.458}{175.889}{reciprocity}%
\TextAt{12}{327.053}{175.889}{homomorphism}%
\TextAt{8}{421.75}{175.889}{of}%
\TextAt{13}{435.08}{175.889}{L}%
\TextAt{16}{446.365}{175.89}{\Character{106}}%
\TextAt{13}{453.007}{175.889}{K}%
\TextAt{8}{463.824}{175.889}{.}%
\TextAt{19}{79.226}{141.105}{Proof.}%
\TextAt{8}{116}{141.105}{We}%
\TextAt{8}{137.215}{141.105}{first}%
\TextAt{8}{162.41}{141.105}{show}%
\TextAt{8}{192.484}{141.105}{that}%
\TextAt{8}{218.923}{141.105}{the}%
\TextAt{8}{240.149}{141.105}{definition}%
\TextAt{8}{293.248}{141.105}{of}%
\TextAt{13}{307.649}{141.105}{r}%
\TextAt{21}{312.924}{139.246}{L}%
\TextAt{3}{318.683}{139.246}{\Character{106}}%
\TextAt{21}{321.035}{139.246}{K}%
\TextAt{9}{329.238}{141.105}{\Character{40}}%
\TextAt{13}{333.79}{141.105}{\Character{27}}%
\TextAt{9}{340.873}{141.105}{\Character{41}}%
\TextAt{8}{350.392}{141.105}{is}%
\TextAt{8}{363.227}{141.105}{independent}%
\TextAt{8}{430.624}{141.105}{of}%
\TextAt{8}{445.024}{141.105}{the}%
\TextAt{8}{466.251}{141.105}{choice}%
\TextAt{8}{502.11}{141.105}{of}%
\TextAt{8}{516.51}{141.105}{the}%
\TextAt{8}{79.226}{124.059}{preimage}%
\TextAt{9}{130.832}{124.059}{\Character{22}}%
\TextAt{13}{130.217}{124.059}{\Character{27}}%
\NativeMathAt{142.173}{124.059}{11.955}{\in}%
\TextAt{8}{155.016}{124.059}{Frob}%
\TextAt{9}{179.572}{124.059}{\Character{40}}%
\TextAt{20}{185.111}{127.081}{\Character{101}}%
\TextAt{13}{184.124}{124.059}{L}%
\TextAt{16}{196.962}{124.059}{\Character{106}}%
\TextAt{13}{205.156}{124.059}{K}%
\TextAt{9}{215.973}{124.059}{\Character{41}}%
\TextAt{8}{225.334}{124.059}{of}%
\TextAt{13}{239.575}{124.059}{\Character{27}}%
\TextAt{8}{246.657}{124.059}{.}%
\TextAt{8}{257.838}{124.059}{For}%
\TextAt{8}{279.718}{124.059}{this,}%
\TextAt{8}{306.931}{124.059}{let}%
\TextAt{9}{325.361}{124.059}{\Character{22}}%
\TextAt{13}{324.746}{124.059}{\Character{27}}%
\TextAt{3}{331.829}{128.398}{\Character{48}}%
\NativeMathAt{339.497}{124.059}{11.955}{\in}%
\TextAt{8}{352.34}{124.059}{Frob}%
\TextAt{9}{376.896}{124.059}{\Character{40}}%
\TextAt{20}{382.435}{127.081}{\Character{101}}%
\TextAt{13}{381.448}{124.059}{L}%
\TextAt{16}{394.286}{124.059}{\Character{106}}%
\TextAt{13}{402.48}{124.059}{K}%
\TextAt{9}{413.297}{124.059}{\Character{41}}%
\TextAt{8}{422.658}{124.059}{be}%
\TextAt{8}{439.493}{124.059}{another}%
\TextAt{8}{483.335}{124.059}{preimage,}%
\TextAt{9}{79.226}{107.013}{\Character{6}}%
\TextAt{3}{87.68}{111.352}{\Character{48}}%
\TextAt{8}{95.259}{107.013}{its}%
\TextAt{8}{112.465}{107.013}{fixed}%
\TextAt{8}{141.637}{107.013}{field}%
\TextAt{8}{167.881}{107.013}{and}%
\TextAt{13}{191.53}{107.013}{\Character{25}}%
\TextAt{17}{198.17}{105.22}{\Character{6}}%
\TextAt{23}{204.286}{107.489}{\Character{48}}%
\NativeMathAt{212.3}{107.013}{11.955}{\in}%
\TextAt{13}{225.102}{107.013}{A}%
\TextAt{17}{233.877}{105.22}{\Character{6}}%
\TextAt{23}{239.993}{107.489}{\Character{48}}%
\TextAt{8}{247.96}{107.013}{a}%
\TextAt{8}{258.602}{107.013}{prime}%
\TextAt{8}{292.652}{107.013}{element.}%
\TextAt{8}{343.104}{107.013}{If}%
\TextAt{13}{355.683}{107.013}{d}%
\TextAt{21}{361.766}{105.22}{K}%
\TextAt{9}{369.968}{107.013}{\Character{40}}%
\TextAt{9}{375.121}{107.013}{\Character{22}}%
\TextAt{13}{374.521}{107.013}{\Character{27}}%
\TextAt{9}{381.603}{107.013}{\Character{41}}%
\TextAt{9}{390.976}{107.013}{\Character{61}}%
\TextAt{13}{404.923}{107.013}{d}%
\TextAt{21}{411.006}{105.22}{K}%
\TextAt{9}{419.209}{107.013}{\Character{40}}%
\TextAt{9}{424.362}{107.013}{\Character{22}}%
\TextAt{13}{423.761}{107.013}{\Character{27}}%
\TextAt{3}{430.843}{111.352}{\Character{48}}%
\TextAt{9}{433.638}{107.013}{\Character{41}}%
\TextAt{8}{438.191}{107.013}{,}%
\TextAt{8}{446.448}{107.013}{then}%
\TextAt{9}{474.609}{107.013}{\Character{22}}%
\TextAt{13}{473.995}{107.013}{\Character{27}}%
\TextAt{16}{481.077}{107.013}{\Character{106}}%
\TextAt{30}{486.368}{105.375}{\Character{101}}%
\TextAt{21}{484.398}{103.36}{K}%
\TextAt{9}{497.432}{107.013}{\Character{61}}%
\TextAt{9}{511.97}{107.013}{\Character{22}}%
\TextAt{13}{511.368}{107.013}{\Character{27}}%
\TextAt{3}{518.451}{111.352}{\Character{48}}%
\TextAt{16}{521.246}{107.013}{\Character{106}}%
\TextAt{30}{526.537}{105.375}{\Character{101}}%
\TextAt{21}{524.567}{103.36}{K}%
\TextAt{8}{79.226}{89.967}{and}%
\TextAt{9}{102.996}{89.967}{\Character{22}}%
\TextAt{13}{102.381}{89.967}{\Character{27}}%
\TextAt{16}{109.464}{89.967}{\Character{106}}%
\TextAt{21}{112.784}{88.174}{L}%
\TextAt{9}{123.031}{89.967}{\Character{61}}%
\TextAt{9}{136.732}{89.967}{\Character{22}}%
\TextAt{13}{136.125}{89.967}{\Character{27}}%
\TextAt{3}{143.208}{94.306}{\Character{48}}%
\TextAt{16}{146.003}{89.967}{\Character{106}}%
\TextAt{21}{149.324}{88.174}{L}%
\TextAt{8}{155.581}{89.967}{,}%
\TextAt{8}{163.221}{89.967}{so}%
\TextAt{8}{177.983}{89.967}{that}%
\TextAt{9}{204.359}{89.967}{\Character{22}}%
\TextAt{13}{203.744}{89.967}{\Character{27}}%
\TextAt{9}{214.816}{89.967}{\Character{61}}%
\TextAt{9}{228.517}{89.967}{\Character{22}}%
\TextAt{13}{227.91}{89.967}{\Character{27}}%
\TextAt{3}{234.993}{94.306}{\Character{48}}%
\TextAt{8}{237.788}{89.967}{,}%
\TextAt{8}{245.428}{89.967}{and}%
\TextAt{8}{268.583}{89.967}{there}%
\TextAt{8}{298.887}{89.967}{is}%
\TextAt{8}{311.043}{89.967}{nothing}%
\TextAt{8}{354.367}{89.967}{to}%
\TextAt{8}{369.069}{89.967}{show.}%
\TextAt{8}{403.801}{89.967}{However,}%
\TextAt{8}{454.694}{89.967}{if}%
\TextAt{8}{465.81}{89.967}{we}%
\TextAt{8}{483.43}{89.967}{have,}%
\TextAt{8}{514.167}{89.967}{say,}%
\TextAt{13}{79.226}{72.921}{d}%
\TextAt{21}{85.309}{71.128}{K}%
\TextAt{9}{93.512}{72.921}{\Character{40}}%
\TextAt{9}{98.665}{72.921}{\Character{22}}%
\TextAt{13}{98.064}{72.921}{\Character{27}}%
\TextAt{9}{105.146}{72.921}{\Character{41}}%
\TextAt{13}{113.019}{72.921}{\Character{60}}%
\TextAt{13}{125.429}{72.921}{d}%
\TextAt{21}{131.528}{71.128}{K}%
\TextAt{9}{139.73}{72.921}{\Character{40}}%
\TextAt{9}{144.883}{72.921}{\Character{22}}%
\TextAt{13}{144.283}{72.921}{\Character{27}}%
\TextAt{3}{151.365}{77.26}{\Character{48}}%
\TextAt{9}{154.16}{72.921}{\Character{41}}%
\TextAt{8}{158.712}{72.921}{,}%
\TextAt{8}{165.812}{72.921}{then}%
\TextAt{9}{193.024}{72.921}{\Character{22}}%
\TextAt{13}{192.409}{72.921}{\Character{27}}%
\TextAt{3}{199.491}{77.26}{\Character{48}}%
\TextAt{9}{205.607}{72.921}{\Character{61}}%
\TextAt{9}{218.638}{72.921}{\Character{22}}%
\TextAt{13}{218.033}{72.921}{\Character{27}}%
\TextAt{9}{225.723}{72.921}{\Character{22}}%
\TextAt{13}{225.115}{72.921}{\Character{28}}%
\TextAt{8}{235.369}{72.921}{for}%
\TextAt{8}{253.191}{72.921}{some}%
\TextAt{9}{283.063}{72.921}{\Character{22}}%
\TextAt{13}{282.455}{72.921}{\Character{28}}%
\NativeMathAt{292.195}{72.921}{11.955}{\in}%
\TextAt{8}{303.486}{72.921}{Frob}%
\TextAt{9}{328.042}{72.921}{\Character{40}}%
\TextAt{20}{333.58}{75.943}{\Character{101}}%
\TextAt{13}{332.594}{72.921}{L}%
\TextAt{16}{343.879}{72.921}{\Character{106}}%
\TextAt{13}{350.521}{72.921}{K}%
\TextAt{9}{361.339}{72.921}{\Character{41}}%
\TextAt{8}{365.891}{72.921}{,}%
\TextAt{8}{372.99}{72.921}{and}%
\TextAt{9}{396.298}{72.921}{\Character{22}}%
\TextAt{13}{395.69}{72.921}{\Character{28}}%
\TextAt{16}{402.109}{72.921}{\Character{106}}%
\TextAt{21}{405.43}{71.128}{L}%
\TextAt{9}{415.008}{72.921}{\Character{61}}%
\TextAt{9}{427.418}{72.921}{1}%
\TextAt{8}{433.287}{72.921}{.}%
\TextAt{8}{441.714}{72.921}{The}%
\TextAt{8}{465.705}{72.921}{fixed}%
\TextAt{8}{493.929}{72.921}{field}%
\TextAt{9}{519.223}{72.921}{\Character{6}}%
\TextAt{3}{527.677}{77.26}{\Character{48}\Character{48}}%
\TextAt{8}{297.211}{43.033}{269}%
\Destination{78.23}{809.815}{nav.page.269}%
\Destination{79.36}{718.755}{nav.print.294}%
\Destination{79.36}{122.055}{nav.print.295}%
\Destination{79.36}{304.855}{nav.number.4.5.6}%
\Hotlink{348.532}{701.555}{25.078}{11.047}{nav.number.4.5.3}%
\Hotlink{116.445}{627.055}{25.078}{11.047}{nav.number.4.5.4}%
\Hotlink{178.231}{400.755}{25.078}{11.047}{nav.number.4.5.4}%
\Destination{79.226}{527.304}{errata.109}%
}

% Page 282 (body page 270)
\Page{%
\TextAt{2}{79.226}{796.943}{IV.}%
\TextAt{2}{98.012}{796.943}{Abstract}%
\TextAt{2}{143.252}{796.943}{Class}%
\TextAt{2}{171.844}{796.943}{Field}%
\TextAt{2}{199.575}{796.943}{Theory}%
\Vector{q 0 G 0 g 0.299 w 7.226 18.729 m 460.769 18.729 l S Q}%
\TextAt{8}{79.226}{761.124}{of}%
\TextAt{9}{93.725}{761.124}{\Character{22}}%
\TextAt{13}{93.117}{761.124}{\Character{28}}%
\TextAt{8}{103.994}{761.124}{contains}%
\TextAt{13}{150.475}{761.124}{L}%
\TextAt{8}{158.439}{761.124}{,}%
\TextAt{8}{166.29}{761.124}{so}%
\TextAt{13}{181.221}{761.124}{r}%
\TextAt{30}{187.259}{759.485}{\Character{101}}%
\TextAt{21}{186.496}{757.471}{L}%
\TextAt{3}{192.255}{757.471}{\Character{106}}%
\TextAt{21}{194.607}{757.471}{K}%
\TextAt{9}{202.81}{761.124}{\Character{40}}%
\TextAt{9}{207.951}{761.124}{\Character{22}}%
\TextAt{13}{207.362}{761.124}{\Character{28}}%
\TextAt{9}{213.781}{761.124}{\Character{41}}%
\NativeMathAt{222.61}{761.124}{11.955}{\equiv}%
\TextAt{13}{236.186}{761.124}{N}%
\TextAt{17}{245.559}{759.264}{\Character{6}}%
\TextAt{23}{251.675}{761.533}{\Character{48}\Character{48}}%
\TextAt{3}{256.546}{759.264}{\Character{106}}%
\TextAt{21}{258.898}{759.264}{K}%
\TextAt{9}{267.101}{761.124}{\Character{40}}%
\TextAt{13}{271.653}{761.124}{\Character{25}}%
\TextAt{17}{278.294}{759.33}{\Character{6}}%
\TextAt{23}{284.41}{761.6}{\Character{48}\Character{48}}%
\TextAt{9}{289.778}{761.124}{\Character{41}}%
\NativeMathAt{298.607}{761.124}{11.955}{\equiv}%
\TextAt{9}{312.183}{761.124}{1}%
\TextAt{9}{323.337}{761.124}{\Character{40}}%
\TextAt{8}{327.901}{761.124}{mod}%
\TextAt{13}{354.003}{761.124}{N}%
\TextAt{21}{363.377}{759.264}{L}%
\TextAt{3}{369.136}{759.264}{\Character{106}}%
\TextAt{21}{371.489}{759.264}{K}%
\TextAt{13}{379.691}{761.124}{A}%
\TextAt{21}{388.467}{759.33}{L}%
\TextAt{9}{394.724}{761.124}{\Character{41}}%
\TextAt{8}{399.276}{761.124}{.}%
\TextAt{8}{409.407}{761.124}{It}%
\TextAt{8}{422.64}{761.124}{follows}%
\TextAt{8}{461.637}{761.124}{therefore}%
\TextAt{8}{511.298}{761.124}{that}%
\ErrataAt{79.19}{744.078}{455.81}{$r_{\widetilde L\mid K}(\bar\sigma^{\prime})=r_{\widetilde L\mid K}(\bar\sigma)r_{\widetilde L\mid K}(\bar\tau)\equiv r_{\widetilde L\mid K}(\bar\sigma)\pmod{N_{L\mid K}A_L}.$}%
\TextAt{8}{97.159}{724.043}{The}%
\TextAt{8}{122.514}{724.043}{fact}%
\TextAt{8}{146.9}{724.043}{that}%
\TextAt{8}{173.57}{724.043}{the}%
\TextAt{8}{195.027}{724.043}{mapping}%
\TextAt{8}{244.459}{724.043}{is}%
\TextAt{8}{257.524}{724.043}{a}%
\TextAt{8}{268.58}{724.043}{homomorphism}%
\TextAt{8}{352.551}{724.043}{now}%
\TextAt{8}{378.24}{724.043}{follows}%
\TextAt{8}{417.977}{724.043}{directly}%
\TextAt{8}{461.874}{724.043}{from}%
\TextAt{8}{490.815}{724.043}{(5.5):}%
\TextAt{8}{525.943}{724.043}{if}%
\TextAt{9}{79.841}{706.997}{\Character{22}}%
\TextAt{13}{79.226}{706.997}{\Character{27}}%
\TextAt{17}{85.88}{705.204}{1}%
\TextAt{13}{90.612}{706.997}{\Character{59}}%
\TextAt{9}{96.471}{706.997}{\Character{22}}%
\TextAt{13}{95.856}{706.997}{\Character{27}}%
\TextAt{17}{102.51}{705.204}{2}%
\NativeMathAt{110.563}{706.997}{11.955}{\in}%
\TextAt{8}{121.854}{706.997}{Frob}%
\TextAt{9}{146.41}{706.997}{\Character{40}}%
\TextAt{20}{151.948}{710.019}{\Character{101}}%
\TextAt{13}{150.962}{706.997}{L}%
\TextAt{16}{162.247}{706.997}{\Character{106}}%
\TextAt{13}{168.889}{706.997}{K}%
\TextAt{9}{179.707}{706.997}{\Character{41}}%
\TextAt{8}{188.026}{706.997}{are}%
\TextAt{8}{207.407}{706.997}{preimages}%
\TextAt{8}{261.972}{706.997}{of}%
\TextAt{13}{275.171}{706.997}{\Character{27}}%
\TextAt{17}{281.825}{705.204}{1}%
\TextAt{13}{286.557}{706.997}{\Character{59}}%
\TextAt{13}{291.781}{706.997}{\Character{27}}%
\TextAt{17}{298.455}{705.204}{2}%
\NativeMathAt{306.508}{706.997}{11.955}{\in}%
\TextAt{13}{317.799}{706.997}{G}%
\TextAt{9}{327.033}{706.997}{\Character{40}}%
\TextAt{13}{331.585}{706.997}{L}%
\TextAt{16}{342.87}{706.997}{\Character{106}}%
\TextAt{13}{349.512}{706.997}{K}%
\TextAt{9}{360.33}{706.997}{\Character{41}}%
\TextAt{8}{364.882}{706.997}{,}%
\TextAt{8}{371.927}{706.997}{then}%
\TextAt{9}{399.072}{706.997}{\Character{22}}%
\TextAt{13}{398.457}{706.997}{\Character{27}}%
\TextAt{17}{405.11}{705.204}{3}%
\TextAt{9}{413.164}{706.997}{\Character{61}}%
\TextAt{9}{426.195}{706.997}{\Character{22}}%
\TextAt{13}{425.589}{706.997}{\Character{27}}%
\TextAt{17}{432.243}{705.204}{1}%
\TextAt{9}{437.59}{706.997}{\Character{22}}%
\TextAt{13}{436.975}{706.997}{\Character{27}}%
\TextAt{17}{443.628}{705.204}{2}%
\TextAt{8}{452.128}{706.997}{is}%
\TextAt{8}{463.762}{706.997}{a}%
\TextAt{8}{473.387}{706.997}{preimage}%
\TextAt{8}{523.337}{706.997}{of}%
\TextAt{13}{79.226}{689.951}{\Character{27}}%
\TextAt{17}{85.88}{688.158}{3}%
\TextAt{9}{93.933}{689.951}{\Character{61}}%
\TextAt{13}{106.358}{689.951}{\Character{27}}%
\TextAt{17}{113.012}{688.158}{1}%
\TextAt{13}{117.744}{689.951}{\Character{27}}%
\TextAt{17}{124.397}{688.158}{2}%
\TextAt{8}{129.13}{689.951}{.}%
\Vector{q 0 G 0 g 0.399 w 452.334 -84.049 m 452.334 -75.979 l S Q}%
\Vector{q 0 G 0 g 0.399 w 452.534 -76.178 m 459.707 -76.178 l S Q}%
\Vector{q 0 G 0 g 0.399 w 452.534 -83.85 m 459.707 -83.85 l S Q}%
\Vector{q 0 G 0 g 0.399 w 459.906 -84.049 m 459.906 -75.979 l S Q}%
\TextAt{8}{97.159}{662.79}{The}%
\TextAt{8}{120.637}{662.79}{definition}%
\TextAt{8}{172.09}{662.79}{of}%
\TextAt{8}{184.844}{662.79}{the}%
\TextAt{8}{204.425}{662.79}{reciprocity}%
\TextAt{8}{262.059}{662.79}{map}%
\TextAt{8}{287.497}{662.79}{expresses}%
\TextAt{8}{337.504}{662.79}{the}%
\TextAt{8}{357.084}{662.79}{fundamental}%
\TextAt{8}{424.151}{662.79}{principle}%
\TextAt{8}{471.694}{662.79}{of}%
\TextAt{8}{484.449}{662.79}{class}%
\TextAt{8}{511.31}{662.79}{field}%
\TextAt{8}{79.226}{645.744}{theory}%
\TextAt{8}{117.177}{645.744}{to}%
\TextAt{8}{132.688}{645.744}{the}%
\TextAt{8}{154.046}{645.744}{effect}%
\TextAt{8}{186.127}{645.744}{that}%
\TextAt{8}{212.696}{645.744}{Frobenius}%
\TextAt{8}{267.767}{645.744}{automorphisms}%
\TextAt{8}{351.052}{645.744}{correspond}%
\TextAt{8}{411.826}{645.744}{to}%
\TextAt{8}{427.337}{645.744}{prime}%
\TextAt{8}{461.701}{645.744}{elements:}%
\TextAt{8}{516.51}{645.744}{the}%
\TextAt{8}{79.226}{628.698}{element}%
\TextAt{13}{123.756}{628.698}{\Character{27}}%
\TextAt{9}{136.355}{628.698}{\Character{61}}%
\TextAt{13}{150.976}{628.698}{\Character{39}}%
\TextAt{17}{158.645}{626.904}{\Character{6}}%
\NativeMathAt{170.776}{628.698}{11.955}{\in}%
\TextAt{13}{184.262}{628.698}{G}%
\TextAt{9}{193.496}{628.698}{\Character{40}}%
\TextAt{20}{199.035}{631.72}{\Character{101}}%
\TextAt{13}{198.048}{628.698}{L}%
\TextAt{16}{211.529}{628.698}{\Character{106}}%
\TextAt{9}{220.366}{628.698}{\Character{6}\Character{41}}%
\TextAt{8}{238.559}{628.698}{is}%
\TextAt{8}{251.611}{628.698}{mapped}%
\TextAt{8}{297.445}{628.698}{to}%
\TextAt{13}{313.043}{628.698}{\Character{25}}%
\TextAt{17}{319.684}{626.904}{\Character{6}}%
\NativeMathAt{331.814}{628.698}{11.955}{\in}%
\TextAt{13}{345.301}{628.698}{A}%
\TextAt{17}{354.076}{626.904}{\Character{6}}%
\TextAt{8}{360.691}{628.698}{;}%
\TextAt{8}{369.772}{628.698}{for}%
\TextAt{8}{388.946}{628.698}{reasons}%
\TextAt{8}{431.336}{628.698}{of}%
\TextAt{8}{445.954}{628.698}{functoriality,}%
\TextAt{8}{516.51}{628.698}{the}%
\TextAt{8}{79.226}{611.652}{inclusion}%
\TextAt{13}{129.032}{611.652}{G}%
\TextAt{9}{138.266}{611.652}{\Character{40}}%
\TextAt{20}{143.804}{614.674}{\Character{101}}%
\TextAt{13}{142.818}{611.652}{L}%
\TextAt{16}{155.754}{611.652}{\Character{106}}%
\TextAt{9}{164.047}{611.652}{\Character{6}\Character{41}}%
\NativeArrowAt{182.026}{611.652}{11.955}{13.214}{\hookrightarrow}%
\TextAt{13}{200.212}{611.652}{G}%
\TextAt{9}{209.446}{611.652}{\Character{40}}%
\TextAt{20}{214.985}{614.674}{\Character{101}}%
\TextAt{13}{213.998}{611.652}{L}%
\TextAt{16}{226.934}{611.652}{\Character{106}}%
\TextAt{13}{235.227}{611.652}{K}%
\TextAt{9}{246.045}{611.652}{\Character{41}}%
\TextAt{8}{255.464}{611.652}{corresponds}%
\TextAt{8}{320.62}{611.652}{to}%
\TextAt{8}{335.899}{611.652}{the}%
\TextAt{8}{357.024}{611.652}{norm}%
\TextAt{8}{388.563}{611.652}{map}%
\TextAt{13}{415.546}{611.652}{N}%
\TextAt{17}{424.92}{609.792}{\Character{6}}%
\TextAt{3}{431.036}{609.792}{\Character{106}}%
\TextAt{21}{433.388}{609.792}{K}%
\TextAt{9}{446.563}{611.652}{\Character{58}}%
\TextAt{13}{454.786}{611.652}{A}%
\TextAt{17}{463.562}{609.858}{\Character{6}}%
\NativeMathAt{475.148}{611.652}{11.955}{\rightarrow}%
\TextAt{13}{492.075}{611.652}{A}%
\TextAt{21}{500.85}{609.858}{K}%
\TextAt{8}{509.053}{611.652}{.}%
\TextAt{8}{520.408}{611.652}{So}%
\TextAt{8}{79.226}{594.606}{the}%
\TextAt{8}{99.55}{594.606}{definition}%
\TextAt{8}{151.747}{594.606}{of}%
\TextAt{13}{165.245}{594.606}{r}%
\TextAt{21}{170.52}{592.746}{L}%
\TextAt{3}{176.279}{592.746}{\Character{106}}%
\TextAt{21}{178.631}{592.746}{K}%
\TextAt{9}{186.834}{594.606}{\Character{40}}%
\TextAt{13}{191.386}{594.606}{\Character{27}}%
\TextAt{9}{198.469}{594.606}{\Character{41}}%
\TextAt{8}{207.086}{594.606}{is}%
\TextAt{8}{219.018}{594.606}{already}%
\TextAt{8}{260.491}{594.606}{forced}%
\TextAt{8}{295.448}{594.606}{upon}%
\TextAt{8}{325.205}{594.606}{us}%
\TextAt{8}{340.388}{594.606}{by}%
\TextAt{8}{356.815}{594.606}{these}%
\TextAt{8}{386.955}{594.606}{requirements.}%
\TextAt{8}{461.66}{594.606}{This}%
\TextAt{8}{488.547}{594.606}{principle}%
\TextAt{8}{79.226}{577.56}{appears}%
\TextAt{8}{122.54}{577.56}{at}%
\TextAt{8}{136.85}{577.56}{its}%
\TextAt{8}{153.169}{577.56}{purest}%
\TextAt{8}{188.999}{577.56}{in}%
\TextAt{8}{202.651}{577.56}{the}%
\TextAt{12}{79.226}{553.387}{Proposition}%
\TextAt{12}{151.937}{553.387}{(5.7).}%
\TextAt{8}{187.253}{553.387}{If}%
\TextAt{13}{198.945}{553.387}{L}%
\TextAt{16}{210.231}{553.387}{\Character{106}}%
\TextAt{13}{216.872}{553.387}{K}%
\TextAt{8}{231.587}{553.387}{is}%
\TextAt{8}{243.351}{553.387}{an}%
\TextAt{8}{259.61}{553.387}{unramified}%
\TextAt{8}{318.143}{553.387}{extension,}%
\TextAt{8}{373.16}{553.387}{then}%
\TextAt{8}{399.82}{553.387}{the}%
\TextAt{8}{419.977}{553.387}{reciprocity}%
\TextAt{8}{478.187}{553.387}{map}%
\TextAt{13}{221.988}{530.825}{r}%
\TextAt{21}{227.263}{528.965}{L}%
\TextAt{3}{233.022}{528.965}{\Character{106}}%
\TextAt{21}{235.375}{528.965}{K}%
\TextAt{9}{246.898}{530.825}{\Character{58}}%
\TextAt{13}{253.471}{530.825}{G}%
\TextAt{9}{262.704}{530.825}{\Character{40}}%
\TextAt{13}{267.256}{530.825}{L}%
\TextAt{16}{278.542}{530.825}{\Character{106}}%
\TextAt{13}{285.184}{530.825}{K}%
\TextAt{9}{296.001}{530.825}{\Character{41}}%
\NativeArrowAt{303.874}{530.825}{11.955}{19.26}{\longrightarrow}%
\TextAt{13}{326.456}{530.825}{A}%
\TextAt{21}{335.232}{529.031}{K}%
\TextAt{9}{343.434}{530.825}{\Character{47}}%
\TextAt{13}{349.287}{530.825}{N}%
\TextAt{21}{358.661}{528.965}{L}%
\TextAt{3}{364.42}{528.965}{\Character{106}}%
\TextAt{21}{366.773}{528.965}{K}%
\TextAt{13}{374.975}{530.825}{A}%
\TextAt{21}{383.751}{529.031}{L}%
\TextAt{8}{79.226}{508.262}{is}%
\TextAt{8}{90.99}{508.262}{given}%
\TextAt{8}{121.559}{508.262}{by}%
\TextAt{13}{227}{491.216}{r}%
\TextAt{21}{232.275}{489.356}{L}%
\TextAt{3}{238.034}{489.356}{\Character{106}}%
\TextAt{21}{240.386}{489.356}{K}%
\TextAt{9}{248.589}{491.216}{\Character{40}}%
\TextAt{13}{253.141}{491.216}{\Character{39}}%
\TextAt{21}{260.81}{489.356}{L}%
\TextAt{3}{266.569}{489.356}{\Character{106}}%
\TextAt{21}{268.921}{489.356}{K}%
\TextAt{9}{277.124}{491.216}{\Character{41}}%
\TextAt{9}{284.979}{491.216}{\Character{61}}%
\TextAt{13}{297.422}{491.216}{\Character{25}}%
\TextAt{21}{304.063}{489.423}{K}%
\TextAt{8}{315.586}{491.216}{mod}%
\TextAt{13}{341.024}{491.216}{N}%
\TextAt{21}{350.398}{489.356}{L}%
\TextAt{3}{356.157}{489.356}{\Character{106}}%
\TextAt{21}{358.509}{489.356}{K}%
\TextAt{13}{366.712}{491.216}{A}%
\TextAt{21}{375.487}{489.423}{L}%
\TextAt{13}{381.744}{491.216}{\Character{59}}%
\TextAt{8}{79.226}{470.454}{and}%
\TextAt{8}{101.989}{470.454}{is}%
\TextAt{8}{113.753}{470.454}{an}%
\TextAt{8}{130.012}{470.454}{isomorphism.}%
\TextAt{19}{79.226}{439.155}{Proof.}%
\TextAt{8}{116}{439.155}{In}%
\TextAt{8}{130.739}{439.155}{this}%
\TextAt{8}{153.68}{439.155}{case}%
\TextAt{8}{178.569}{439.155}{one}%
\TextAt{8}{200.146}{439.155}{has}%
\TextAt{20}{222.125}{442.177}{\Character{101}}%
\TextAt{13}{221.138}{439.155}{L}%
\TextAt{9}{232.625}{439.155}{\Character{61}}%
\TextAt{20}{247.989}{442.177}{\Character{101}}%
\TextAt{13}{245.251}{439.155}{K}%
\TextAt{8}{260.084}{439.155}{and}%
\TextAt{13}{282.965}{439.155}{\Character{39}}%
\TextAt{21}{290.633}{437.361}{K}%
\NativeMathAt{302.358}{439.155}{11.955}{\in}%
\TextAt{13}{313.85}{439.155}{G}%
\TextAt{9}{323.084}{439.155}{\Character{40}}%
\TextAt{20}{330.374}{442.177}{\Character{101}}%
\TextAt{13}{327.636}{439.155}{K}%
\TextAt{16}{341.975}{439.155}{\Character{106}}%
\TextAt{13}{348.818}{439.155}{K}%
\TextAt{9}{359.636}{439.155}{\Character{41}}%
\TextAt{8}{368.204}{439.155}{is}%
\TextAt{8}{380.086}{439.155}{a}%
\TextAt{8}{389.959}{439.155}{preimage}%
\TextAt{8}{440.157}{439.155}{of}%
\TextAt{13}{453.605}{439.155}{\Character{39}}%
\TextAt{21}{461.274}{437.295}{L}%
\TextAt{3}{467.033}{437.295}{\Character{106}}%
\TextAt{21}{469.385}{437.295}{K}%
\TextAt{8}{481.603}{439.155}{with}%
\TextAt{8}{508.381}{439.155}{fixed}%
\TextAt{8}{79.226}{422.109}{field}%
\TextAt{13}{105.508}{422.109}{K}%
\TextAt{8}{116.326}{422.109}{,}%
\TextAt{8}{124.632}{422.109}{i.e.,}%
\TextAt{13}{147.894}{422.109}{r}%
\TextAt{21}{153.169}{420.249}{L}%
\TextAt{3}{158.928}{420.249}{\Character{106}}%
\TextAt{21}{161.28}{420.249}{K}%
\TextAt{9}{169.483}{422.109}{\Character{40}}%
\TextAt{13}{174.035}{422.109}{\Character{39}}%
\TextAt{21}{181.704}{420.249}{L}%
\TextAt{3}{187.463}{420.249}{\Character{106}}%
\TextAt{21}{189.815}{420.249}{K}%
\TextAt{9}{198.018}{422.109}{\Character{41}}%
\TextAt{9}{207.451}{422.109}{\Character{61}}%
\TextAt{13}{221.47}{422.109}{\Character{25}}%
\TextAt{21}{228.111}{420.315}{K}%
\TextAt{8}{239.634}{422.109}{mod}%
\TextAt{13}{265.072}{422.109}{N}%
\TextAt{21}{274.446}{420.249}{L}%
\TextAt{3}{280.205}{420.249}{\Character{106}}%
\TextAt{21}{282.557}{420.249}{K}%
\TextAt{13}{290.76}{422.109}{A}%
\TextAt{21}{299.535}{420.315}{L}%
\TextAt{8}{305.792}{422.109}{.}%
\TextAt{8}{317.017}{422.109}{The}%
\TextAt{8}{341.996}{422.109}{fact}%
\TextAt{8}{366.007}{422.109}{that}%
\TextAt{8}{392.301}{422.109}{we}%
\TextAt{8}{410.454}{422.109}{have}%
\TextAt{8}{438.374}{422.109}{an}%
\TextAt{8}{455.558}{422.109}{isomorphism}%
\TextAt{8}{524.903}{422.109}{is}%
\TextAt{8}{79.226}{405.063}{seen}%
\TextAt{8}{104.643}{405.063}{from}%
\TextAt{8}{132.283}{405.063}{the}%
\TextAt{8}{152.44}{405.063}{composite}%
\TextAt{13}{185.5}{382.5}{G}%
\TextAt{9}{194.734}{382.5}{\Character{40}}%
\TextAt{13}{199.286}{382.5}{L}%
\TextAt{16}{210.572}{382.5}{\Character{106}}%
\TextAt{13}{217.213}{382.5}{K}%
\TextAt{9}{228.031}{382.5}{\Character{41}}%
\NativeArrowAt{235.904}{382.5}{11.955}{19.26}{\longrightarrow}%
\TextAt{13}{258.486}{382.5}{A}%
\TextAt{21}{267.262}{380.707}{K}%
\TextAt{9}{275.464}{382.5}{\Character{47}}%
\TextAt{13}{281.317}{382.5}{N}%
\TextAt{21}{290.691}{380.641}{L}%
\TextAt{3}{296.45}{380.64}{\Character{106}}%
\TextAt{21}{298.802}{380.641}{K}%
\TextAt{13}{307.005}{382.5}{A}%
\TextAt{21}{315.78}{380.707}{L}%
\NativeArrowAt{325.358}{382.5}{11.955}{19.26}{\longrightarrow}%
\TextAt{13}{347.94}{382.5}{Z}%
\TextAt{9}{356.811}{382.5}{\Character{47}}%
\TextAt{13}{362.664}{382.5}{nZ}%
\NativeMathAt{381.843}{382.5}{11.955}{\cong}%
\TextAt{18}{394.463}{382.5}{Z}%
\TextAt{9}{402.433}{382.5}{\Character{47}}%
\TextAt{13}{408.286}{382.5}{n}%
\TextAt{18}{415.273}{382.5}{Z}%
\TextAt{13}{423.244}{382.5}{\Character{59}}%
\TextAt{8}{79.226}{359.938}{with}%
\TextAt{13}{105.017}{359.938}{n}%
\TextAt{9}{115.325}{359.938}{\Character{61}}%
\TextAt{9}{127.735}{359.938}{\Character{91}}%
\TextAt{13}{131.003}{359.938}{L}%
\TextAt{9}{142.288}{359.938}{\Character{58}}%
\TextAt{13}{148.86}{359.938}{K}%
\TextAt{9}{159.678}{359.938}{\Character{93}}%
\TextAt{8}{162.93}{359.938}{,}%
\TextAt{8}{169.384}{359.938}{where}%
\TextAt{8}{202.324}{359.938}{the}%
\TextAt{8}{221.611}{359.938}{second}%
\TextAt{8}{258.52}{359.938}{map}%
\TextAt{8}{283.665}{359.938}{is}%
\TextAt{8}{294.56}{359.938}{induced}%
\TextAt{8}{337.256}{359.938}{by}%
\TextAt{8}{352.645}{359.938}{the}%
\TextAt{8}{371.933}{359.938}{valuation}%
\TextAt{13}{422.136}{359.938}{v}%
\TextAt{21}{427.795}{358.145}{K}%
\TextAt{9}{439.319}{359.938}{\Character{58}}%
\TextAt{13}{445.891}{359.938}{A}%
\TextAt{21}{454.667}{358.145}{K}%
\NativeMathAt{466.19}{359.938}{11.955}{\rightarrow}%
\TextAt{13}{481.466}{359.938}{Z}%
\TextAt{8}{493.365}{359.938}{because}%
\ErrataAt{79.19}{342.894}{455.81}{$v_K(N_{L\mid K}A_L)\subseteq nZ$. If $v_K(a)\in nZ$, choose $b\in A_K$ with $v_K(b)=v_K(a)/n$. Then}%
\ErrataAt{79.19}{325.848}{455.81}{$u=ab^{-n}\in U_K$. By \hyperlink{nav.number.4.5.1}{(5.1)}, $u=N_{L\mid K}(\epsilon)$ for some $\epsilon\in U_L$, so $a=N_{L\mid K}(\epsilon b)$.}%
\ErrataAt{79.19}{308.802}{455.81}{Thus the valuation map is injective. It is surjective since $v_K(\pi_K)=1$. The elements}%
\ErrataAt{79.19}{291.756}{455.81}{$\varphi_{L\mid K}$, $\pi_K\bmod N_{L\mid K}A_L$ and $1\bmod nZ$ correspond and generate the respective groups.}%
\Vector{q 0 G 0 g 0.399 w 452.334 -482.246 m 452.334 -474.176 l S Q}%
\Vector{q 0 G 0 g 0.399 w 452.534 -474.376 m 459.707 -474.376 l S Q}%
\Vector{q 0 G 0 g 0.399 w 452.534 -482.047 m 459.707 -482.047 l S Q}%
\Vector{q 0 G 0 g 0.399 w 459.906 -482.246 m 459.906 -474.176 l S Q}%
\TextAt{8}{97.159}{264.592}{The}%
\TextAt{8}{121.213}{264.592}{reciprocity}%
\TextAt{8}{179.422}{264.592}{homomorphism}%
\TextAt{13}{262.092}{264.592}{r}%
\TextAt{21}{267.368}{262.733}{L}%
\TextAt{3}{273.127}{262.733}{\Character{106}}%
\TextAt{21}{275.479}{262.733}{K}%
\TextAt{8}{287.579}{264.592}{exhibits}%
\TextAt{8}{331.538}{264.592}{the}%
\TextAt{8}{351.695}{264.592}{following}%
\TextAt{8}{401.129}{264.592}{functorial}%
\TextAt{8}{454.139}{264.592}{behaviour.}%
\TextAt{12}{79.226}{240.42}{Proposition}%
\TextAt{12}{152.416}{240.42}{(5.8).}%
\TextAt{8}{188.148}{240.42}{Let}%
\TextAt{13}{209.534}{240.42}{L}%
\TextAt{16}{221.529}{240.42}{\Character{106}}%
\TextAt{13}{228.88}{240.42}{K}%
\TextAt{8}{244.011}{240.42}{and}%
\TextAt{13}{267.19}{240.42}{L}%
\TextAt{3}{275.155}{244.758}{\Character{48}}%
\TextAt{16}{281.98}{240.42}{\Character{106}}%
\TextAt{13}{289.331}{240.42}{K}%
\TextAt{3}{300.149}{244.758}{\Character{48}}%
\TextAt{8}{307.257}{240.42}{be}%
\TextAt{8}{323.598}{240.42}{finite}%
\TextAt{8}{353.926}{240.42}{Galois}%
\TextAt{8}{390.256}{240.42}{extensions,}%
\TextAt{8}{450.408}{240.42}{so}%
\TextAt{8}{465.195}{240.42}{that}%
\TextAt{13}{490.98}{240.42}{K}%
\NativeMathAt{505.828}{240.42}{11.955}{\subseteq}%
\TextAt{13}{519.157}{240.42}{K}%
\TextAt{3}{529.974}{244.758}{\Character{48}}%
\TextAt{8}{79.226}{223.374}{and}%
\TextAt{13}{101.989}{223.374}{L}%
\NativeMathAt{113.274}{223.374}{11.955}{\subseteq}%
\TextAt{13}{125.893}{223.374}{L}%
\TextAt{3}{133.858}{227.712}{\Character{48}}%
\TextAt{8}{136.653}{223.374}{,}%
\TextAt{8}{143.802}{223.374}{and}%
\TextAt{8}{166.565}{223.374}{let}%
\TextAt{13}{183.469}{223.374}{\Character{27}}%
\NativeMathAt{193.873}{223.374}{11.955}{\in}%
\TextAt{13}{205.164}{223.374}{G}%
\TextAt{8}{214.397}{223.374}{.}%
\TextAt{8}{222.846}{223.374}{Then}%
\TextAt{8}{253.403}{223.374}{we}%
\TextAt{8}{270.63}{223.374}{have}%
\TextAt{8}{297.625}{223.374}{the}%
\TextAt{8}{317.782}{223.374}{commutative}%
\TextAt{8}{387.707}{223.374}{diagrams}%
\TextAt{13}{230.19}{198.812}{G}%
\TextAt{9}{239.424}{198.812}{\Character{40}}%
\TextAt{13}{243.976}{198.812}{L}%
\TextAt{3}{251.94}{203.151}{\Character{48}}%
\TextAt{16}{258.056}{198.812}{\Character{106}}%
\TextAt{13}{264.698}{198.812}{K}%
\TextAt{3}{275.516}{203.151}{\Character{48}}%
\TextAt{9}{278.311}{198.812}{\Character{41}}%
\TextAt{13}{307.515}{198.812}{A}%
\TextAt{21}{316.29}{197.019}{K}%
\TextAt{23}{323.994}{199.288}{\Character{48}}%
\TextAt{9}{327.177}{198.812}{\Character{47}}%
\TextAt{13}{333.03}{198.812}{N}%
\TextAt{21}{342.404}{196.952}{L}%
\TextAt{23}{348.163}{199.222}{\Character{48}}%
\TextAt{3}{350.847}{196.952}{\Character{106}}%
\TextAt{21}{353.2}{196.952}{K}%
\TextAt{23}{360.904}{199.222}{\Character{48}}%
\TextAt{13}{364.087}{198.812}{A}%
\TextAt{21}{372.862}{197.019}{L}%
\TextAt{23}{378.621}{199.288}{\Character{48}}%
\TextAt{13}{232.985}{155.715}{G}%
\TextAt{9}{242.219}{155.715}{\Character{40}}%
\TextAt{13}{246.771}{155.715}{L}%
\TextAt{16}{258.056}{155.715}{\Character{106}}%
\TextAt{13}{264.698}{155.715}{K}%
\TextAt{9}{275.516}{155.715}{\Character{41}}%
\TextAt{13}{312.884}{155.715}{A}%
\TextAt{21}{321.659}{153.921}{K}%
\TextAt{9}{329.862}{155.715}{\Character{47}}%
\TextAt{13}{335.715}{155.715}{N}%
\TextAt{21}{345.088}{153.855}{L}%
\TextAt{3}{350.847}{153.855}{\Character{106}}%
\TextAt{21}{353.2}{153.855}{K}%
\TextAt{13}{361.402}{155.715}{A}%
\TextAt{21}{370.178}{153.921}{L}%
\NativeDiagram{%
\path[draw=black,line width=0.478bp] (216.216,-572.198) -- (229.686,-572.198);%
\path[draw=black,line width=0.478bp,line cap=round,line join=round] (227.438,-569.329) .. controls (227.892,-571.051) and (228.905,-571.863) .. (229.925,-572.198) .. controls (228.905,-572.533) and (227.892,-573.346) .. (227.438,-575.067);%
}%
\TextAt{21}{283.409}{208.403}{r}%
\TextAt{22}{287.23}{206.111}{L}%
\TextAt{23}{292.289}{207.827}{\Character{48}}%
\TextAt{23}{294.973}{206.111}{\Character{106}}%
\TextAt{22}{297.243}{206.111}{K}%
\TextAt{23}{303.787}{207.827}{\Character{48}}%
\NativeDiagram{%
\path[draw=black,line width=0.478bp] (184.527,-582.755) -- (184.527,-604.261);%
\path[draw=black,line width=0.478bp,line cap=round,line join=round] (187.396,-602.014) .. controls (185.675,-602.468) and (184.862,-603.481) .. (184.527,-604.5) .. controls (184.192,-603.481) and (183.38,-602.468) .. (181.658,-602.014);%
\path[draw=black,line width=0.478bp] (272.661,-583.618) -- (272.661,-604.261);%
\path[draw=black,line width=0.478bp,line cap=round,line join=round] (275.53,-602.014) .. controls (273.808,-602.468) and (272.996,-603.481) .. (272.661,-604.5) .. controls (272.326,-603.481) and (271.513,-602.468) .. (269.792,-602.014);%
}%
\TextAt{21}{347.476}{178.991}{N}%
\TextAt{22}{354.197}{176.699}{K}%
\TextAt{23}{360.742}{178.415}{\Character{48}}%
\TextAt{23}{363.426}{176.699}{\Character{106}}%
\TextAt{22}{365.696}{176.699}{K}%
\NativeDiagram{%
\path[draw=black,line width=0.478bp] (213.421,-615.296) -- (235.055,-615.296);%
\path[draw=black,line width=0.478bp,line cap=round,line join=round] (232.807,-612.427) .. controls (233.261,-614.149) and (234.274,-614.962) .. (235.294,-615.296) .. controls (234.274,-615.631) and (233.261,-616.444) .. (232.807,-618.165);%
}%
\TextAt{21}{287.381}{152.457}{r}%
\TextAt{22}{291.202}{150.719}{L}%
\TextAt{23}{296.261}{150.719}{\Character{106}}%
\TextAt{22}{298.53}{150.719}{K}%
\TextAt{8}{297.211}{43.033}{270}%
\Destination{78.23}{809.815}{nav.page.270}%
\Destination{79.36}{567.355}{nav.number.4.5.7}%
\Destination{79.36}{254.455}{nav.number.4.5.8}%
\Hotlink{490.315}{721.855}{25.078}{11.047}{nav.number.4.5.5}%
\Hotlink{274.353}{127.155}{21.483}{9.572}{nav.number.4.4.6}%
\Destination{79.226}{359.894}{errata.100}%
\Destination{79.226}{364.23}{errata.108}%
\Destination{79.226}{761.078}{errata.110}%
\Hotlink{279.35}{322.853}{21.722}{11.047}{nav.number.4.5.1}%
}

% Page 283 (body page 271)
\Page{%
\TextAt{2}{79.226}{796.943}{IV.}%
\TextAt{2}{98.012}{796.943}{Abstract}%
\TextAt{2}{143.252}{796.943}{Class}%
\TextAt{2}{171.844}{796.943}{Field}%
\TextAt{2}{199.575}{796.943}{Theory}%
\Vector{q 0 G 0 g 0.299 w 7.226 18.729 m 460.769 18.729 l S Q}%
\TextAt{13}{226.735}{758.034}{G}%
\TextAt{9}{235.969}{758.034}{\Character{40}}%
\TextAt{13}{240.521}{758.034}{L}%
\TextAt{16}{251.807}{758.034}{\Character{106}}%
\TextAt{13}{258.448}{758.034}{K}%
\TextAt{9}{269.266}{758.034}{\Character{41}}%
\TextAt{13}{315.711}{758.034}{A}%
\TextAt{21}{324.486}{756.241}{K}%
\TextAt{9}{332.689}{758.034}{\Character{47}}%
\TextAt{13}{338.542}{758.034}{N}%
\TextAt{21}{347.916}{756.175}{L}%
\TextAt{3}{353.675}{756.175}{\Character{106}}%
\TextAt{21}{356.027}{756.175}{K}%
\TextAt{13}{364.23}{758.034}{A}%
\TextAt{21}{373.005}{756.241}{L}%
\TextAt{13}{221.112}{714.937}{G}%
\TextAt{9}{230.346}{714.937}{\Character{40}}%
\TextAt{13}{234.898}{714.937}{L}%
\TextAt{21}{242.863}{719.275}{\Character{27}}%
\TextAt{16}{251.807}{714.937}{\Character{106}}%
\TextAt{13}{258.449}{714.937}{K}%
\TextAt{21}{269.267}{719.275}{\Character{27}}%
\TextAt{9}{274.89}{714.937}{\Character{41}}%
\TextAt{13}{304.093}{714.937}{A}%
\TextAt{21}{312.868}{713.143}{K}%
\TextAt{22}{320.573}{715.413}{\Character{27}}%
\TextAt{9}{326.067}{714.937}{\Character{47}}%
\TextAt{13}{331.92}{714.937}{N}%
\TextAt{21}{341.294}{713.077}{L}%
\TextAt{22}{347.053}{715.346}{\Character{27}}%
\TextAt{3}{352.049}{713.077}{\Character{106}}%
\TextAt{21}{354.401}{713.077}{K}%
\TextAt{22}{362.105}{715.346}{\Character{27}}%
\TextAt{13}{367.599}{714.937}{A}%
\TextAt{21}{376.375}{713.143}{L}%
\TextAt{22}{382.134}{715.413}{\Character{27}}%
\TextAt{13}{387.628}{714.937}{\Character{59}}%
\NativeDiagram{%
\path[draw=black,line width=0.478bp] (207.172,-12.976) -- (237.883,-12.976);%
\path[draw=black,line width=0.478bp,line cap=round,line join=round] (235.635,-10.107) .. controls (236.09,-11.829) and (237.102,-12.641) .. (238.122,-12.976) .. controls (237.102,-13.311) and (236.09,-14.124) .. (235.635,-15.845);%
}%
\TextAt{21}{285.67}{767.072}{r}%
\TextAt{22}{289.491}{765.334}{L}%
\TextAt{23}{294.55}{765.334}{\Character{106}}%
\TextAt{22}{296.819}{765.334}{K}%
\NativeDiagram{%
\path[draw=black,line width=0.478bp] (178.278,-23.533) -- (178.278,-45.039);%
\path[draw=black,line width=0.478bp,line cap=round,line join=round] (181.148,-42.792) .. controls (179.426,-43.246) and (178.613,-44.259) .. (178.278,-45.278) .. controls (177.944,-44.259) and (177.131,-43.246) .. (175.409,-42.792);%
}%
\TextAt{21}{238.02}{736.678}{\Character{27}}%
\TextAt{23}{243.145}{739.49}{\Character{3}}%
\NativeDiagram{%
\path[draw=black,line width=0.478bp] (275.489,-24.396) -- (275.489,-45.039);%
\path[draw=black,line width=0.478bp,line cap=round,line join=round] (278.358,-42.792) .. controls (276.637,-43.246) and (275.824,-44.259) .. (275.489,-45.278) .. controls (275.154,-44.259) and (274.342,-43.246) .. (272.62,-42.792);%
}%
\TextAt{21}{350.305}{737.327}{\Character{27}}%
\NativeDiagram{%
\path[draw=black,line width=0.478bp] (212.796,-56.074) -- (226.265,-56.074);%
\path[draw=black,line width=0.478bp,line cap=round,line join=round] (224.018,-53.205) .. controls (224.472,-54.927) and (225.484,-55.74) .. (226.504,-56.074) .. controls (225.484,-56.409) and (224.472,-57.222) .. (224.018,-58.943);%
}%
\TextAt{21}{277.677}{705.179}{r}%
\TextAt{22}{281.498}{703.441}{L}%
\TextAt{22}{286.557}{705.157}{\Character{27}}%
\TextAt{23}{291.553}{703.441}{\Character{106}}%
\TextAt{22}{293.822}{703.441}{K}%
\TextAt{22}{300.367}{705.157}{\Character{27}}%
\TextAt{8}{79.226}{690.605}{where}%
\TextAt{8}{113.895}{690.605}{the}%
\TextAt{8}{134.912}{690.605}{vertical}%
\TextAt{8}{177.4}{690.605}{arrows}%
\TextAt{8}{215.727}{690.605}{on}%
\TextAt{8}{232.846}{690.605}{the}%
\TextAt{8}{253.863}{690.605}{left}%
\TextAt{8}{275.202}{690.605}{are}%
\TextAt{8}{295.573}{690.605}{given}%
\TextAt{8}{327.002}{690.605}{by}%
\TextAt{13}{344.122}{690.605}{\Character{27}}%
\TextAt{3}{351.204}{694.943}{\Character{48}}%
\NativeArrowAt{358.785}{690.605}{11.955}{11.955}{\mapsto}%
\TextAt{13}{375.527}{690.605}{\Character{27}}%
\TextAt{3}{382.61}{694.943}{\Character{48}}%
\TextAt{16}{385.405}{690.605}{\Character{106}}%
\TextAt{21}{388.726}{688.811}{L}%
\TextAt{8}{394.983}{690.605}{,}%
\TextAt{8}{403.207}{690.605}{resp.}%
\TextAt{8}{435.109}{690.605}{by}%
\TextAt{8}{452.228}{690.605}{the}%
\TextAt{8}{473.245}{690.605}{conjugation}%
\TextAt{13}{79.226}{673.559}{\Character{28}}%
\NativeArrowAt{88.966}{673.559}{11.955}{11.955}{\mapsto}%
\TextAt{13}{104.242}{673.559}{\Character{27}}%
\TextAt{3}{111.324}{677.897}{\Character{0}}%
\TextAt{17}{117.911}{677.897}{1}%
\TextAt{13}{122.643}{673.559}{\Character{28}}%
\TextAt{13}{129.051}{673.559}{\Character{27}}%
\TextAt{8}{136.145}{673.559}{.}%
\TextAt{19}{79.226}{640.144}{Proof.}%
\TextAt{8}{116}{640.144}{Let}%
\TextAt{13}{137.081}{640.144}{\Character{27}}%
\TextAt{3}{144.163}{644.482}{\Character{48}}%
\NativeMathAt{150.469}{640.144}{11.955}{\in}%
\TextAt{13}{161.95}{640.144}{G}%
\TextAt{9}{171.184}{640.144}{\Character{40}}%
\TextAt{13}{175.736}{640.144}{L}%
\TextAt{3}{183.701}{644.482}{\Character{48}}%
\TextAt{16}{190.007}{640.144}{\Character{106}}%
\TextAt{13}{196.838}{640.144}{K}%
\TextAt{3}{207.656}{644.482}{\Character{48}}%
\TextAt{9}{210.451}{640.144}{\Character{41}}%
\TextAt{8}{219.012}{640.144}{and}%
\TextAt{13}{241.886}{640.144}{\Character{27}}%
\TextAt{9}{252.479}{640.144}{\Character{61}}%
\TextAt{13}{265.095}{640.144}{\Character{27}}%
\TextAt{3}{272.177}{644.482}{\Character{48}}%
\TextAt{16}{274.972}{640.144}{\Character{106}}%
\TextAt{21}{278.293}{638.351}{L}%
\NativeMathAt{288.061}{640.144}{11.955}{\in}%
\TextAt{13}{299.542}{640.144}{G}%
\TextAt{9}{308.776}{640.144}{\Character{40}}%
\TextAt{13}{313.328}{640.144}{L}%
\TextAt{16}{324.803}{640.144}{\Character{106}}%
\TextAt{13}{331.635}{640.144}{K}%
\TextAt{9}{342.453}{640.144}{\Character{41}}%
\TextAt{8}{347.005}{640.144}{.}%
\TextAt{8}{355.788}{640.144}{If}%
\TextAt{9}{368.206}{640.144}{\Character{22}}%
\TextAt{13}{367.591}{640.144}{\Character{27}}%
\TextAt{3}{374.674}{644.482}{\Character{48}}%
\NativeMathAt{380.979}{640.144}{11.955}{\in}%
\TextAt{8}{392.46}{640.144}{Frob}%
\TextAt{9}{417.016}{640.144}{\Character{40}}%
\TextAt{20}{422.555}{643.166}{\Character{101}}%
\TextAt{13}{421.569}{640.144}{L}%
\TextAt{3}{429.533}{644.482}{\Character{48}}%
\TextAt{16}{435.839}{640.144}{\Character{106}}%
\TextAt{13}{442.671}{640.144}{K}%
\TextAt{3}{453.488}{644.482}{\Character{48}}%
\TextAt{9}{456.283}{640.144}{\Character{41}}%
\TextAt{8}{464.844}{640.144}{is}%
\TextAt{8}{476.72}{640.144}{a}%
\TextAt{8}{486.587}{640.144}{preimage}%
\TextAt{8}{79.226}{623.098}{of}%
\TextAt{13}{91.655}{623.098}{\Character{27}}%
\TextAt{3}{98.738}{627.436}{\Character{48}}%
\TextAt{8}{101.533}{623.098}{,}%
\TextAt{8}{107.961}{623.098}{then}%
\TextAt{9}{134.335}{623.098}{\Character{22}}%
\TextAt{13}{133.72}{623.098}{\Character{27}}%
\TextAt{9}{144.124}{623.098}{\Character{61}}%
\TextAt{9}{157.155}{623.098}{\Character{22}}%
\TextAt{13}{156.549}{623.098}{\Character{27}}%
\TextAt{3}{163.631}{627.436}{\Character{48}}%
\TextAt{16}{166.427}{623.098}{\Character{106}}%
\TextAt{30}{170.51}{621.46}{\Character{101}}%
\TextAt{21}{169.747}{619.445}{L}%
\NativeMathAt{179.325}{623.098}{11.955}{\in}%
\TextAt{8}{190.616}{623.098}{Frob}%
\TextAt{9}{215.172}{623.098}{\Character{40}}%
\TextAt{20}{220.711}{626.12}{\Character{101}}%
\TextAt{13}{219.725}{623.098}{L}%
\TextAt{16}{231.01}{623.098}{\Character{106}}%
\TextAt{13}{237.652}{623.098}{K}%
\TextAt{9}{248.469}{623.098}{\Character{41}}%
\TextAt{8}{256.018}{623.098}{is}%
\TextAt{8}{266.881}{623.098}{a}%
\TextAt{8}{275.736}{623.098}{preimage}%
\TextAt{8}{324.915}{623.098}{of}%
\TextAt{13}{337.344}{623.098}{\Character{27}}%
\TextAt{8}{347.423}{623.098}{such}%
\TextAt{8}{372.919}{623.098}{that}%
\TextAt{13}{397.387}{623.098}{d}%
\TextAt{21}{403.47}{621.305}{K}%
\TextAt{9}{411.672}{623.098}{\Character{40}}%
\TextAt{9}{416.825}{623.098}{\Character{22}}%
\TextAt{13}{416.225}{623.098}{\Character{27}}%
\TextAt{9}{423.307}{623.098}{\Character{41}}%
\TextAt{9}{431.162}{623.098}{\Character{61}}%
\TextAt{13}{443.606}{623.098}{f}%
\TextAt{21}{449.376}{621.238}{K}%
\TextAt{23}{457.08}{623.507}{\Character{48}}%
\TextAt{3}{459.765}{621.238}{\Character{106}}%
\TextAt{21}{462.117}{621.238}{K}%
\TextAt{13}{470.32}{623.098}{d}%
\TextAt{21}{476.402}{621.305}{K}%
\TextAt{23}{484.107}{623.574}{\Character{48}}%
\TextAt{9}{487.289}{623.098}{\Character{40}}%
\TextAt{9}{492.442}{623.098}{\Character{22}}%
\TextAt{13}{491.842}{623.098}{\Character{27}}%
\TextAt{3}{498.924}{627.436}{\Character{48}}%
\TextAt{9}{501.719}{623.098}{\Character{41}}%
\NativeMathAt{509.592}{623.098}{11.955}{\in}%
\TextAt{18}{520.883}{623.098}{N}%
\TextAt{8}{529.518}{623.098}{.}%
\TextAt{8}{79.226}{606.052}{Let}%
\TextAt{9}{99.181}{606.052}{\Character{6}}%
\TextAt{3}{107.636}{610.39}{\Character{48}}%
\TextAt{8}{113.314}{606.052}{be}%
\TextAt{8}{128.224}{606.052}{the}%
\TextAt{8}{147.366}{606.052}{fixed}%
\TextAt{8}{174.638}{606.052}{field}%
\TextAt{8}{198.98}{606.052}{of}%
\TextAt{9}{211.911}{606.052}{\Character{22}}%
\TextAt{13}{211.296}{606.052}{\Character{27}}%
\TextAt{3}{218.379}{610.39}{\Character{48}}%
\TextAt{8}{221.174}{606.052}{.}%
\TextAt{8}{229.284}{606.052}{Then}%
\TextAt{9}{258.827}{606.052}{\Character{6}}%
\TextAt{9}{270.591}{606.052}{\Character{61}}%
\TextAt{9}{283.012}{606.052}{\Character{6}}%
\TextAt{3}{291.482}{610.39}{\Character{48}}%
\TextAt{16}{294.86}{606.052}{\Character{92}}%
\TextAt{20}{304.399}{609.074}{\Character{101}}%
\TextAt{13}{303.413}{606.052}{L}%
\TextAt{9}{314.698}{606.052}{\Character{61}}%
\TextAt{9}{327.108}{606.052}{\Character{6}}%
\TextAt{3}{335.578}{610.39}{\Character{48}}%
\TextAt{16}{338.955}{606.052}{\Character{92}}%
\TextAt{20}{348.414}{609.074}{\Character{101}}%
\TextAt{9}{347.508}{606.052}{\Character{6}}%
\TextAt{8}{358.846}{606.052}{is}%
\TextAt{8}{369.595}{606.052}{the}%
\TextAt{8}{388.737}{606.052}{fixed}%
\TextAt{8}{416.009}{606.052}{field}%
\TextAt{8}{440.352}{606.052}{of}%
\TextAt{9}{453.282}{606.052}{\Character{22}}%
\TextAt{13}{452.667}{606.052}{\Character{27}}%
\TextAt{8}{462.633}{606.052}{and}%
\TextAt{13}{484.381}{606.052}{f}%
\TextAt{17}{490.151}{604.192}{\Character{6}}%
\TextAt{23}{496.268}{606.462}{\Character{48}}%
\TextAt{3}{498.952}{604.192}{\Character{106}}%
\TextAt{17}{501.304}{604.192}{\Character{6}}%
\TextAt{9}{511.239}{606.052}{\Character{61}}%
\TextAt{9}{523.649}{606.052}{1}%
\TextAt{8}{529.518}{606.052}{.}%
\TextAt{8}{79.226}{589.006}{If}%
\TextAt{8}{90.883}{589.006}{now}%
\TextAt{13}{115.237}{589.006}{\Character{25}}%
\TextAt{17}{121.877}{587.213}{\Character{6}}%
\TextAt{23}{127.993}{589.482}{\Character{48}}%
\NativeMathAt{134.496}{589.006}{11.955}{\in}%
\TextAt{13}{145.787}{589.006}{A}%
\TextAt{17}{154.563}{587.213}{\Character{6}}%
\TextAt{23}{160.679}{589.482}{\Character{48}}%
\TextAt{8}{167.724}{589.006}{is}%
\TextAt{8}{179.452}{589.006}{a}%
\TextAt{8}{189.173}{589.006}{prime}%
\TextAt{8}{222.301}{589.006}{element}%
\TextAt{8}{265.508}{589.006}{of}%
\TextAt{9}{278.803}{589.006}{\Character{6}}%
\TextAt{3}{287.257}{593.344}{\Character{48}}%
\TextAt{8}{290.053}{589.006}{,}%
\TextAt{8}{297.174}{589.006}{then}%
\TextAt{13}{323.799}{589.006}{\Character{25}}%
\TextAt{17}{330.439}{587.213}{\Character{6}}%
\TextAt{9}{340.374}{589.006}{\Character{61}}%
\TextAt{13}{352.799}{589.006}{N}%
\TextAt{17}{362.173}{587.146}{\Character{6}}%
\TextAt{23}{368.289}{589.416}{\Character{48}}%
\TextAt{3}{370.974}{587.146}{\Character{106}}%
\TextAt{17}{373.326}{587.146}{\Character{6}}%
\TextAt{9}{379.94}{589.006}{\Character{40}}%
\TextAt{13}{384.492}{589.006}{\Character{25}}%
\TextAt{17}{391.133}{587.213}{\Character{6}}%
\TextAt{23}{397.249}{589.482}{\Character{48}}%
\TextAt{9}{400.431}{589.006}{\Character{41}}%
\NativeMathAt{408.304}{589.006}{11.955}{\in}%
\TextAt{13}{419.595}{589.006}{A}%
\TextAt{17}{428.371}{587.213}{\Character{6}}%
\TextAt{8}{438.847}{589.006}{is}%
\TextAt{8}{450.576}{589.006}{a}%
\TextAt{8}{460.296}{589.006}{prime}%
\TextAt{8}{493.425}{589.006}{element}%
\TextAt{8}{79.226}{571.96}{of}%
\TextAt{9}{92.229}{571.96}{\Character{6}}%
\TextAt{8}{100.683}{571.96}{.}%
\TextAt{8}{109.022}{571.96}{The}%
\TextAt{8}{132.748}{571.96}{commutativity}%
\TextAt{8}{211.133}{571.96}{of}%
\TextAt{8}{224.136}{571.96}{the}%
\TextAt{8}{243.964}{571.96}{diagram}%
\TextAt{8}{289.174}{571.96}{on}%
\TextAt{8}{305.105}{571.96}{the}%
\TextAt{8}{324.934}{571.96}{left}%
\TextAt{8}{345.086}{571.96}{therefore}%
\TextAt{8}{393.858}{571.96}{follows}%
\TextAt{8}{431.966}{571.96}{from}%
\TextAt{8}{459.279}{571.96}{the}%
\TextAt{8}{479.108}{571.96}{equality}%
\TextAt{8}{523.337}{571.96}{of}%
\TextAt{8}{79.226}{554.914}{norms}%
\TextAt{13}{189.598}{528.65}{N}%
\TextAt{17}{198.971}{526.79}{\Character{6}}%
\TextAt{3}{205.087}{526.79}{\Character{106}}%
\TextAt{21}{207.44}{526.79}{K}%
\TextAt{9}{215.642}{528.65}{\Character{40}}%
\TextAt{13}{220.195}{528.65}{\Character{25}}%
\TextAt{17}{226.835}{526.857}{\Character{6}}%
\TextAt{9}{233.449}{528.65}{\Character{41}}%
\TextAt{9}{241.304}{528.65}{\Character{61}}%
\TextAt{13}{253.748}{528.65}{N}%
\TextAt{17}{263.122}{526.79}{\Character{6}}%
\TextAt{3}{269.238}{526.79}{\Character{106}}%
\TextAt{21}{271.59}{526.79}{K}%
\TextAt{9}{279.793}{528.65}{\Character{40}}%
\TextAt{13}{284.345}{528.65}{N}%
\TextAt{17}{293.719}{526.79}{\Character{6}}%
\TextAt{23}{299.835}{529.059}{\Character{48}}%
\TextAt{3}{302.519}{526.79}{\Character{106}}%
\TextAt{17}{304.872}{526.79}{\Character{6}}%
\TextAt{9}{311.486}{528.65}{\Character{40}}%
\TextAt{13}{316.038}{528.65}{\Character{25}}%
\TextAt{17}{322.678}{526.857}{\Character{6}}%
\TextAt{23}{328.794}{529.126}{\Character{48}}%
\TextAt{9}{331.977}{528.65}{\Character{41}\Character{41}}%
\TextAt{9}{241.322}{508.615}{\Character{61}}%
\TextAt{13}{253.748}{508.615}{N}%
\TextAt{17}{263.122}{506.755}{\Character{6}}%
\TextAt{23}{269.238}{509.025}{\Character{48}}%
\TextAt{3}{271.922}{506.755}{\Character{106}}%
\TextAt{21}{274.274}{506.755}{K}%
\TextAt{9}{282.477}{508.615}{\Character{40}}%
\TextAt{13}{287.029}{508.615}{\Character{25}}%
\TextAt{17}{293.67}{506.822}{\Character{6}}%
\TextAt{23}{299.786}{509.091}{\Character{48}}%
\TextAt{9}{302.968}{508.615}{\Character{41}}%
\TextAt{9}{310.823}{508.615}{\Character{61}}%
\TextAt{13}{323.267}{508.615}{N}%
\TextAt{21}{332.641}{506.755}{K}%
\TextAt{23}{340.345}{509.025}{\Character{48}}%
\TextAt{3}{343.03}{506.755}{\Character{106}}%
\TextAt{21}{345.382}{506.755}{K}%
\TextAt{9}{353.585}{508.615}{\Character{40}}%
\TextAt{13}{358.137}{508.615}{N}%
\TextAt{17}{367.511}{506.755}{\Character{6}}%
\TextAt{23}{373.627}{509.025}{\Character{48}}%
\TextAt{3}{376.311}{506.755}{\Character{106}}%
\TextAt{21}{378.664}{506.755}{K}%
\TextAt{23}{386.368}{509.025}{\Character{48}}%
\TextAt{9}{389.551}{508.615}{\Character{40}}%
\TextAt{13}{394.103}{508.615}{\Character{25}}%
\TextAt{17}{400.743}{506.822}{\Character{6}}%
\TextAt{23}{406.859}{509.091}{\Character{48}}%
\TextAt{9}{410.042}{508.615}{\Character{41}\Character{41}}%
\TextAt{13}{419.146}{508.615}{\Character{58}}%
\TextAt{8}{79.226}{482.351}{On}%
\TextAt{8}{98.81}{482.351}{the}%
\TextAt{8}{119.039}{482.351}{other}%
\TextAt{8}{149.681}{482.351}{hand,}%
\TextAt{8}{182.29}{482.351}{let}%
\TextAt{13}{199.268}{482.351}{\Character{28}}%
\NativeMathAt{209.132}{482.351}{11.955}{\in}%
\TextAt{13}{220.547}{482.351}{G}%
\TextAt{9}{229.78}{482.351}{\Character{40}}%
\TextAt{13}{234.333}{482.351}{L}%
\TextAt{16}{245.742}{482.351}{\Character{106}}%
\TextAt{13}{252.508}{482.351}{K}%
\TextAt{9}{263.326}{482.351}{\Character{41}}%
\TextAt{8}{267.878}{482.351}{,}%
\TextAt{8}{275.118}{482.351}{and}%
\TextAt{8}{297.954}{482.351}{let}%
\TextAt{9}{315.539}{482.351}{\Character{22}}%
\TextAt{13}{314.931}{482.351}{\Character{28}}%
\TextAt{8}{325.32}{482.351}{be}%
\TextAt{8}{341.317}{482.351}{a}%
\TextAt{8}{351.146}{482.351}{preimage}%
\TextAt{8}{401.299}{482.351}{in}%
\TextAt{8}{415.024}{482.351}{Frob}%
\TextAt{9}{439.58}{482.351}{\Character{40}}%
\TextAt{20}{445.119}{485.373}{\Character{101}}%
\TextAt{13}{444.133}{482.351}{L}%
\TextAt{16}{455.542}{482.351}{\Character{106}}%
\TextAt{13}{462.308}{482.351}{K}%
\TextAt{9}{473.126}{482.351}{\Character{41}}%
\TextAt{8}{481.648}{482.351}{with}%
\TextAt{8}{508.381}{482.351}{fixed}%
\TextAt{8}{79.226}{465.305}{field}%
\TextAt{9}{103.91}{465.305}{\Character{6}}%
\TextAt{8}{112.364}{465.305}{,}%
\TextAt{8}{118.975}{465.305}{and}%
\TextAt{20}{141.278}{465.305}{\Character{98}}%
\TextAt{13}{141.065}{465.305}{\Character{28}}%
\NativeMathAt{150.804}{465.305}{11.955}{\in}%
\TextAt{13}{162.095}{465.305}{G}%
\TextAt{8}{174.553}{465.305}{a}%
\TextAt{8}{183.636}{465.305}{lifting}%
\TextAt{8}{217.106}{465.305}{of}%
\TextAt{9}{230.372}{465.305}{\Character{22}}%
\TextAt{13}{229.763}{465.305}{\Character{28}}%
\TextAt{8}{239.407}{465.305}{to}%
\TextAt{9}{253.362}{468.46}{\Character{22}}%
\TextAt{13}{253.044}{465.305}{k}%
\TextAt{8}{259.533}{465.305}{.}%
\TextAt{8}{267.757}{465.305}{Then}%
\TextAt{9}{297.642}{465.305}{\Character{6}}%
\TextAt{21}{306.096}{469.643}{\Character{27}}%
\TextAt{8}{314.944}{465.305}{is}%
\TextAt{8}{326.035}{465.305}{the}%
\TextAt{8}{345.518}{465.305}{fixed}%
\TextAt{8}{373.131}{465.305}{field}%
\TextAt{8}{397.815}{465.305}{of}%
\TextAt{13}{410.471}{465.305}{\Character{27}}%
\TextAt{3}{417.554}{469.643}{\Character{0}}%
\TextAt{17}{424.14}{469.643}{1}%
\TextAt{20}{429.086}{465.305}{\Character{98}}%
\TextAt{13}{428.873}{465.305}{\Character{28}}%
\TextAt{13}{435.281}{465.305}{\Character{27}}%
\TextAt{16}{442.374}{465.305}{\Character{106}}%
\TextAt{30}{446.838}{463.445}{\Character{102}}%
\TextAt{21}{445.695}{461.431}{L}%
\TextAt{22}{451.454}{463.7}{\Character{27}}%
\TextAt{8}{456.948}{465.305}{,}%
\TextAt{8}{463.559}{465.305}{and}%
\TextAt{8}{485.648}{465.305}{if}%
\TextAt{13}{495.699}{465.305}{\Character{25}}%
\NativeMathAt{506.089}{465.305}{11.955}{\in}%
\TextAt{13}{517.38}{465.305}{A}%
\TextAt{17}{526.155}{463.512}{\Character{6}}%
\TextAt{8}{79.226}{448.259}{is}%
\TextAt{8}{91.103}{448.259}{a}%
\TextAt{8}{100.972}{448.259}{prime}%
\TextAt{8}{134.25}{448.259}{element}%
\TextAt{8}{177.605}{448.259}{of}%
\TextAt{9}{191.048}{448.259}{\Character{6}}%
\TextAt{8}{199.503}{448.259}{,}%
\TextAt{8}{206.794}{448.259}{then}%
\TextAt{13}{233.567}{448.259}{\Character{25}}%
\TextAt{21}{240.636}{452.597}{\Character{27}}%
\NativeMathAt{249.774}{448.259}{11.955}{\in}%
\TextAt{13}{261.258}{448.259}{A}%
\TextAt{17}{270.034}{446.466}{\Character{6}}%
\TextAt{22}{276.15}{448.735}{\Character{27}}%
\TextAt{8}{285.655}{448.259}{is}%
\TextAt{8}{297.532}{448.259}{a}%
\TextAt{8}{307.401}{448.259}{prime}%
\TextAt{8}{340.678}{448.259}{element}%
\TextAt{8}{384.033}{448.259}{of}%
\TextAt{9}{397.477}{448.259}{\Character{6}}%
\TextAt{21}{405.931}{452.597}{\Character{27}}%
\TextAt{8}{411.555}{448.259}{.}%
\TextAt{8}{420.343}{448.259}{The}%
\TextAt{8}{444.511}{448.259}{commutativity}%
\TextAt{8}{523.337}{448.259}{of}%
\TextAt{8}{79.226}{431.213}{the}%
\TextAt{8}{99.383}{431.213}{diagram}%
\TextAt{8}{144.92}{431.213}{on}%
\TextAt{8}{161.179}{431.213}{the}%
\TextAt{8}{181.335}{431.213}{right}%
\TextAt{8}{209.633}{431.213}{therefore}%
\TextAt{8}{258.733}{431.213}{follows}%
\TextAt{8}{297.169}{431.213}{from}%
\TextAt{8}{324.809}{431.213}{the}%
\TextAt{8}{344.966}{431.213}{equality}%
\TextAt{8}{389.523}{431.213}{of}%
\TextAt{8}{402.853}{431.213}{norms}%
\TextAt{13}{243.661}{404.949}{N}%
\TextAt{17}{253.035}{403.089}{\Character{6}}%
\TextAt{22}{259.151}{405.358}{\Character{27}}%
\TextAt{3}{264.146}{403.089}{\Character{106}}%
\TextAt{21}{266.499}{403.089}{K}%
\TextAt{22}{274.203}{405.358}{\Character{27}}%
\TextAt{9}{279.697}{404.949}{\Character{40}}%
\TextAt{13}{284.25}{404.949}{\Character{25}}%
\TextAt{21}{291.319}{409.885}{\Character{27}}%
\TextAt{9}{296.942}{404.949}{\Character{41}}%
\TextAt{9}{304.797}{404.949}{\Character{61}}%
\TextAt{13}{317.241}{404.949}{N}%
\TextAt{17}{326.615}{403.089}{\Character{6}}%
\TextAt{3}{332.731}{403.089}{\Character{106}}%
\TextAt{21}{335.083}{403.089}{K}%
\TextAt{9}{343.286}{404.949}{\Character{40}}%
\TextAt{13}{347.838}{404.949}{\Character{25}}%
\TextAt{9}{354.907}{404.949}{\Character{41}}%
\TextAt{21}{359.46}{409.885}{\Character{27}}%
\TextAt{13}{365.083}{404.949}{\Character{58}}%
\Vector{q 0 G 0 g 0.399 w 452.334 -395.315 m 452.334 -387.246 l S Q}%
\Vector{q 0 G 0 g 0.399 w 452.534 -387.445 m 459.707 -387.445 l S Q}%
\Vector{q 0 G 0 g 0.399 w 452.534 -395.116 m 459.707 -395.116 l S Q}%
\Vector{q 0 G 0 g 0.399 w 459.906 -395.315 m 459.906 -387.246 l S Q}%
\TextAt{8}{97.159}{350.465}{Another}%
\TextAt{8}{144.164}{350.465}{very}%
\TextAt{8}{171.014}{350.465}{interesting}%
\TextAt{8}{229.795}{350.465}{functorial}%
\TextAt{8}{283.962}{350.465}{property}%
\TextAt{8}{332.928}{350.465}{of}%
\TextAt{8}{347.415}{350.465}{the}%
\TextAt{8}{368.729}{350.465}{reciprocity}%
\TextAt{8}{428.097}{350.465}{map}%
\TextAt{8}{455.268}{350.465}{is}%
\TextAt{8}{468.19}{350.465}{obtained}%
\TextAt{8}{517.479}{350.465}{via}%
\TextAt{8}{79.226}{333.419}{the}%
\TextAt{19}{101.241}{333.419}{transfer}%
\TextAt{8}{147.274}{333.419}{(}%
\TextAt{19}{151.829}{333.419}{Verlagerung}%
\TextAt{8}{218.927}{333.419}{in}%
\TextAt{8}{234.438}{333.419}{German).}%
\TextAt{8}{294.072}{333.419}{For}%
\TextAt{8}{316.899}{333.419}{an}%
\TextAt{8}{335.017}{333.419}{arbitrary}%
\TextAt{8}{386.645}{333.419}{group}%
\TextAt{13}{421.679}{333.419}{G}%
\TextAt{8}{430.913}{333.419}{,}%
\TextAt{8}{440.385}{333.419}{let}%
\TextAt{13}{459.149}{333.419}{G}%
\TextAt{3}{468.382}{337.758}{\Character{48}}%
\TextAt{8}{476.933}{333.419}{denote}%
\TextAt{8}{516.51}{333.419}{the}%
\TextAt{8}{79.226}{316.373}{commutator}%
\TextAt{8}{145.255}{316.373}{subgroup}%
\TextAt{8}{196.052}{316.373}{and}%
\TextAt{8}{218.815}{316.373}{write}%
\TextAt{13}{275.234}{299.327}{G}%
\TextAt{7}{284.467}{304.264}{ab}%
\TextAt{9}{297.221}{299.327}{\Character{61}}%
\TextAt{13}{309.646}{299.327}{G}%
\TextAt{9}{318.88}{299.327}{\Character{47}}%
\TextAt{13}{324.733}{299.327}{G}%
\TextAt{3}{333.967}{304.264}{\Character{48}}%
\TextAt{8}{79.226}{276.979}{for}%
\TextAt{8}{97.179}{276.979}{the}%
\TextAt{8}{117.404}{276.979}{maximal}%
\TextAt{8}{165.281}{276.979}{abelian}%
\TextAt{8}{205.997}{276.979}{quotient}%
\TextAt{8}{252.248}{276.979}{group.}%
\TextAt{8}{290.18}{276.979}{If}%
\TextAt{8}{301.941}{276.979}{then}%
\TextAt{13}{328.669}{276.979}{H}%
\NativeMathAt{342.729}{276.979}{11.955}{\subseteq}%
\TextAt{13}{355.465}{276.979}{G}%
\TextAt{8}{368.664}{276.979}{is}%
\TextAt{8}{380.496}{276.979}{a}%
\TextAt{8}{390.32}{276.979}{subgroup}%
\TextAt{8}{441.186}{276.979}{of}%
\TextAt{8}{454.584}{276.979}{finite}%
\TextAt{8}{484.564}{276.979}{index,}%
\TextAt{8}{519.439}{276.979}{we}%
\TextAt{8}{79.226}{259.933}{have}%
\TextAt{8}{106.221}{259.933}{a}%
\TextAt{8}{115.976}{259.933}{canonical}%
\TextAt{8}{167.36}{259.933}{homomorphism}%
\TextAt{8}{259.474}{233.669}{Ver}%
\TextAt{9}{281.325}{233.669}{\Character{58}}%
\TextAt{13}{287.898}{233.669}{G}%
\TextAt{7}{297.131}{238.605}{ab}%
\NativeArrowAt{309.885}{233.669}{11.955}{19.26}{\longrightarrow}%
\TextAt{13}{332.467}{233.669}{H}%
\TextAt{7}{343.089}{238.605}{ab}%
\TextAt{8}{79.226}{207.404}{which}%
\TextAt{8}{112.159}{207.404}{is}%
\TextAt{8}{123.37}{207.404}{called}%
\TextAt{12}{155.98}{207.404}{transfer}%
\TextAt{8}{205.124}{207.404}{from}%
\TextAt{13}{232.211}{207.404}{G}%
\TextAt{8}{244.788}{207.404}{to}%
\TextAt{13}{258.545}{207.404}{H}%
\TextAt{8}{269.168}{207.404}{.}%
\TextAt{8}{277.432}{207.404}{This}%
\TextAt{8}{303.598}{207.404}{homomorphism}%
\TextAt{8}{385.715}{207.404}{is}%
\TextAt{8}{396.925}{207.404}{defined}%
\TextAt{8}{436.684}{207.404}{as}%
\TextAt{8}{450.501}{207.404}{follows}%
\TextAt{8}{488.383}{207.404}{(see}%
\TextAt{8}{511.298}{207.404}{[75],}%
\TextAt{8}{79.226}{190.358}{chap.}%
\TextAt{8}{111.417}{190.358}{IV,}%
\TextAt{8}{131.562}{190.358}{§1).}%
\TextAt{8}{97.159}{170.324}{Let}%
\TextAt{13}{117.895}{170.324}{R}%
\TextAt{8}{130.568}{170.324}{be}%
\TextAt{8}{146.259}{170.324}{a}%
\TextAt{8}{155.781}{170.324}{system}%
\TextAt{8}{194.366}{170.324}{of}%
\TextAt{8}{207.463}{170.324}{representatives}%
\TextAt{8}{287.03}{170.324}{for}%
\TextAt{8}{304.682}{170.324}{the}%
\TextAt{8}{324.605}{170.324}{left}%
\TextAt{8}{344.851}{170.324}{cosets}%
\TextAt{8}{378.558}{170.324}{of}%
\TextAt{13}{391.655}{170.324}{H}%
\TextAt{8}{405.942}{170.324}{in}%
\TextAt{13}{419.361}{170.324}{G}%
\TextAt{8}{428.595}{170.324}{,}%
\TextAt{13}{435.557}{170.324}{G}%
\TextAt{9}{448.112}{170.324}{\Character{61}}%
\TextAt{13}{460.537}{170.324}{RH}%
\TextAt{8}{480.168}{170.324}{,}%
\TextAt{9}{487.131}{170.324}{1}%
\NativeMathAt{496.305}{170.324}{11.955}{\in}%
\TextAt{13}{507.596}{170.324}{R}%
\TextAt{8}{516.604}{170.324}{.}%
\TextAt{8}{524.975}{170.324}{If}%
\TextAt{13}{79.226}{153.278}{\Character{27}}%
\NativeMathAt{89.629}{153.278}{11.955}{\in}%
\TextAt{13}{100.92}{153.278}{G}%
\TextAt{8}{114.051}{153.278}{we}%
\TextAt{8}{131.279}{153.278}{write,}%
\TextAt{8}{164.442}{153.278}{for}%
\TextAt{8}{182.327}{153.278}{every}%
\TextAt{13}{213.22}{153.278}{\Character{26}}%
\NativeMathAt{222.578}{153.278}{11.955}{\in}%
\TextAt{13}{233.869}{153.278}{R}%
\TextAt{8}{242.878}{153.278}{,}%
\TextAt{13}{222.339}{127.013}{\Character{27}}%
\TextAt{13}{229.405}{127.013}{\Character{26}}%
\TextAt{9}{238.78}{127.013}{\Character{61}}%
\TextAt{13}{251.205}{127.013}{\Character{26}}%
\TextAt{3}{257.243}{131.95}{\Character{48}}%
\TextAt{13}{260.038}{127.013}{\Character{27}}%
\TextAt{21}{266.691}{125.22}{\Character{26}}%
\TextAt{13}{271.526}{127.013}{\Character{59}}%
\TextAt{13}{300.661}{127.013}{\Character{27}}%
\TextAt{21}{307.334}{125.22}{\Character{26}}%
\NativeMathAt{315.49}{127.013}{11.955}{\in}%
\TextAt{13}{326.781}{127.013}{H}%
\TextAt{13}{336.74}{127.013}{\Character{59}}%
\TextAt{13}{353.931}{127.013}{\Character{26}}%
\TextAt{3}{359.99}{131.95}{\Character{48}}%
\NativeMathAt{366.106}{127.013}{11.955}{\in}%
\TextAt{13}{377.397}{127.013}{R\Character{59}}%
\TextAt{8}{79.226}{100.749}{and}%
\TextAt{8}{101.989}{100.749}{we}%
\TextAt{8}{119.216}{100.749}{define}%
\TextAt{8}{222.713}{83.703}{Ver}%
\TextAt{9}{241.244}{83.703}{\Character{40}}%
\TextAt{13}{245.796}{83.703}{\Character{27}}%
\TextAt{8}{256.199}{83.703}{mod}%
\TextAt{13}{281.637}{83.703}{G}%
\TextAt{3}{290.871}{88.639}{\Character{48}}%
\TextAt{9}{293.666}{83.703}{\Character{41}}%
\TextAt{9}{301.521}{83.703}{\Character{61}}%
\TextAt{20}{314.531}{95.061}{\Character{89}}%
\TextAt{21}{313.964}{69.689}{\Character{26}}%
\NativeMathAt{318.301}{69.689}{7.97}{\in}%
\TextAt{21}{323.946}{69.689}{R}%
\TextAt{13}{332.366}{83.703}{\Character{27}}%
\TextAt{21}{339.02}{81.91}{\Character{26}}%
\TextAt{8}{347.175}{83.703}{mod}%
\TextAt{13}{372.613}{83.703}{H}%
\TextAt{3}{383.236}{88.639}{\Character{48}}%
\TextAt{13}{386.031}{83.703}{\Character{58}}%
\TextAt{8}{297.211}{43.033}{271}%
\Destination{78.23}{809.815}{nav.page.271}%
\Destination{79.36}{705.555}{nav.print.296}%
\Hotlink{510.798}{205.155}{19.22}{11.047}{nav.cite.N075}%
\Hotlink{110.917}{188.155}{13.995}{11.047}{nav.chapter.4}%
\Hotlink{136.741}{188.155}{6.858}{11.047}{nav.section.4.1}%
}

% Page 284 (body page 272)
\Page{%
\TextAt{2}{79.226}{796.943}{IV.}%
\TextAt{2}{98.012}{796.943}{Abstract}%
\TextAt{2}{143.252}{796.943}{Class}%
\TextAt{2}{171.844}{796.943}{Field}%
\TextAt{2}{199.575}{796.943}{Theory}%
\Vector{q 0 G 0 g 0.299 w 7.226 18.729 m 460.769 18.729 l S Q}%
\TextAt{8}{79.226}{761.124}{Another}%
\TextAt{8}{125.074}{761.124}{description}%
\TextAt{8}{184.97}{761.124}{of}%
\TextAt{8}{198.3}{761.124}{the}%
\TextAt{8}{218.456}{761.124}{transfer}%
\TextAt{8}{261.77}{761.124}{results}%
\TextAt{8}{298.962}{761.124}{from}%
\TextAt{8}{326.603}{761.124}{the}%
\TextAt{8}{346.759}{761.124}{double}%
\TextAt{8}{384.477}{761.124}{coset}%
\TextAt{8}{413.804}{761.124}{decomposition}%
\TextAt{13}{270.155}{732.123}{G}%
\TextAt{9}{282.709}{732.123}{\Character{61}}%
\TextAt{20}{295.135}{743.48}{\Character{91}}%
\TextAt{21}{299.437}{718.374}{\Character{28}}%
\TextAt{16}{308.418}{732.123}{\Character{104}}%
\TextAt{13}{313.068}{732.123}{\Character{27}}%
\TextAt{16}{320.15}{732.123}{\Character{105}}%
\TextAt{13}{324.799}{732.123}{\Character{28}}%
\TextAt{13}{331.207}{732.123}{H}%
\TextAt{8}{79.226}{695.261}{of}%
\TextAt{13}{93.867}{695.261}{G}%
\TextAt{8}{108.308}{695.261}{in}%
\TextAt{8}{123.271}{695.261}{terms}%
\TextAt{8}{157.159}{695.261}{of}%
\TextAt{8}{171.8}{695.261}{the}%
\TextAt{8}{193.267}{695.261}{subgroups}%
\TextAt{16}{249.989}{695.261}{\Character{104}}%
\TextAt{13}{254.639}{695.261}{\Character{27}}%
\TextAt{16}{261.721}{695.261}{\Character{105}}%
\TextAt{8}{271.578}{695.261}{and}%
\TextAt{13}{295.651}{695.261}{H}%
\TextAt{8}{306.274}{695.261}{.}%
\TextAt{8}{318.653}{695.261}{Letting}%
\TextAt{13}{361.101}{695.261}{f}%
\TextAt{9}{368.148}{695.261}{\Character{40}}%
\TextAt{13}{372.7}{695.261}{\Character{28}}%
\TextAt{9}{379.119}{695.261}{\Character{41}}%
\TextAt{8}{388.879}{695.261}{denote}%
\TextAt{8}{427.908}{695.261}{the}%
\TextAt{8}{449.375}{695.261}{smallest}%
\TextAt{8}{495.685}{695.261}{natural}%
\TextAt{8}{79.226}{678.215}{number}%
\TextAt{8}{121.822}{678.215}{such}%
\TextAt{8}{148.219}{678.215}{that}%
\TextAt{13}{173.588}{678.215}{\Character{27}}%
\TextAt{21}{180.242}{676.421}{\Character{28}}%
\TextAt{9}{188.739}{678.215}{\Character{61}}%
\TextAt{13}{201.165}{678.215}{\Character{28}}%
\TextAt{3}{207.583}{682.553}{\Character{0}}%
\TextAt{17}{214.17}{682.553}{1}%
\TextAt{13}{218.902}{678.215}{\Character{27}}%
\TextAt{21}{225.985}{682.553}{f}%
\TextAt{17}{230.931}{682.553}{\Character{40}}%
\TextAt{21}{234.225}{682.553}{\Character{28}}%
\TextAt{17}{238.903}{682.553}{\Character{41}}%
\TextAt{13}{242.694}{678.215}{\Character{28}}%
\NativeMathAt{252.434}{678.215}{11.955}{\in}%
\TextAt{13}{263.725}{678.215}{H}%
\TextAt{8}{274.348}{678.215}{,}%
\TextAt{8}{281.497}{678.215}{one}%
\TextAt{8}{302.957}{678.215}{has}%
\TextAt{13}{323.83}{678.215}{H}%
\TextAt{16}{337.11}{678.215}{\Character{92}}%
\TextAt{16}{347.726}{678.215}{\Character{104}}%
\TextAt{13}{352.386}{678.215}{\Character{28}}%
\TextAt{3}{358.805}{682.553}{\Character{0}}%
\TextAt{17}{365.391}{682.553}{1}%
\TextAt{13}{370.123}{678.215}{\Character{27}}%
\TextAt{13}{377.189}{678.215}{\Character{28}}%
\TextAt{16}{383.625}{678.215}{\Character{105}}%
\TextAt{9}{391.595}{678.215}{\Character{61}}%
\TextAt{16}{404.02}{678.215}{\Character{104}}%
\TextAt{13}{408.67}{678.215}{\Character{27}}%
\TextAt{21}{415.323}{676.421}{\Character{28}}%
\TextAt{16}{420.499}{678.215}{\Character{105}}%
\TextAt{8}{425.149}{678.215}{,}%
\TextAt{8}{432.298}{678.215}{and}%
\TextAt{8}{455.061}{678.215}{we}%
\TextAt{8}{472.288}{678.215}{find}%
\TextAt{8}{495.696}{678.215}{that}%
\TextAt{8}{223.109}{649.214}{Ver}%
\TextAt{9}{241.639}{649.214}{\Character{40}}%
\TextAt{13}{246.192}{649.214}{\Character{27}}%
\TextAt{8}{256.595}{649.214}{mod}%
\TextAt{13}{282.033}{649.214}{G}%
\TextAt{3}{291.266}{654.15}{\Character{48}}%
\TextAt{9}{294.062}{649.214}{\Character{41}}%
\TextAt{9}{301.917}{649.214}{\Character{61}}%
\TextAt{20}{314.36}{660.571}{\Character{89}}%
\TextAt{21}{319.659}{635.465}{\Character{28}}%
\TextAt{13}{331.629}{649.214}{\Character{27}}%
\TextAt{21}{338.282}{647.42}{\Character{28}}%
\TextAt{8}{346.78}{649.214}{mod}%
\TextAt{13}{372.217}{649.214}{H}%
\TextAt{3}{382.84}{654.15}{\Character{48}}%
\TextAt{13}{385.635}{649.214}{\Character{58}}%
\TextAt{8}{79.226}{611.711}{This}%
\TextAt{8}{108.002}{611.711}{formula}%
\TextAt{8}{152.99}{611.711}{is}%
\TextAt{8}{166.81}{611.711}{obtained}%
\TextAt{8}{216.998}{611.711}{from}%
\TextAt{8}{246.695}{611.711}{the}%
\TextAt{8}{268.908}{611.711}{one}%
\TextAt{8}{292.424}{611.711}{above}%
\TextAt{8}{327.656}{611.711}{by}%
\TextAt{8}{345.972}{611.711}{choosing}%
\TextAt{8}{395.574}{611.711}{for}%
\TextAt{13}{415.515}{611.711}{R}%
\TextAt{8}{430.477}{611.711}{the}%
\TextAt{8}{452.69}{611.711}{set}%
\TextAt{16}{473.014}{611.711}{\Character{102}}%
\TextAt{13}{478.992}{611.711}{\Character{27}}%
\TextAt{21}{486.074}{616.049}{i}%
\TextAt{13}{489.456}{611.711}{\Character{28}}%
\TextAt{16}{502.7}{611.711}{\Character{106}}%
\TextAt{13}{512.846}{611.711}{i}%
\TextAt{9}{523.665}{611.711}{\Character{61}}%
\TextAt{9}{79.226}{594.665}{1}%
\TextAt{13}{85.079}{594.665}{\Character{59}}%
\TextAt{13}{90.303}{594.665}{\Character{58}}%
\TextAt{13}{95.54}{594.665}{\Character{58}}%
\TextAt{13}{100.776}{594.665}{\Character{58}}%
\TextAt{13}{106.001}{594.665}{\Character{59}}%
\TextAt{13}{111.237}{594.665}{f}%
\TextAt{9}{118.346}{594.665}{\Character{40}}%
\TextAt{13}{122.899}{594.665}{\Character{28}}%
\TextAt{9}{129.318}{594.665}{\Character{41}}%
\TextAt{16}{133.87}{594.665}{\Character{103}}%
\TextAt{8}{139.848}{594.665}{.}%
\TextAt{8}{148.296}{594.665}{Applying}%
\TextAt{8}{199.022}{594.665}{this}%
\TextAt{8}{221.844}{594.665}{to}%
\TextAt{8}{236.154}{594.665}{the}%
\TextAt{8}{256.311}{594.665}{reciprocity}%
\TextAt{8}{314.521}{594.665}{homomorphism}%
\TextAt{13}{217.272}{565.664}{r}%
\TextAt{21}{222.547}{563.804}{L}%
\TextAt{3}{228.306}{563.804}{\Character{106}}%
\TextAt{21}{230.658}{563.804}{K}%
\TextAt{9}{242.182}{565.664}{\Character{58}}%
\TextAt{13}{248.754}{565.664}{G}%
\TextAt{9}{257.988}{565.664}{\Character{40}}%
\TextAt{13}{262.54}{565.664}{L}%
\TextAt{16}{273.825}{565.664}{\Character{106}}%
\TextAt{13}{280.467}{565.664}{K}%
\TextAt{9}{291.285}{565.664}{\Character{41}}%
\TextAt{7}{295.837}{570.6}{ab}%
\NativeArrowAt{308.591}{565.664}{11.955}{19.26}{\longrightarrow}%
\TextAt{13}{331.173}{565.664}{A}%
\TextAt{21}{339.948}{563.87}{K}%
\TextAt{9}{348.151}{565.664}{\Character{47}}%
\TextAt{13}{354.004}{565.664}{N}%
\TextAt{21}{363.377}{563.804}{L}%
\TextAt{3}{369.137}{563.804}{\Character{106}}%
\TextAt{21}{371.489}{563.804}{K}%
\TextAt{13}{379.692}{565.664}{A}%
\TextAt{21}{388.467}{563.87}{L}%
\TextAt{8}{79.226}{536.662}{we}%
\TextAt{8}{96.454}{536.662}{get}%
\TextAt{8}{115.964}{536.662}{the}%
\TextAt{12}{79.226}{510.65}{Proposition}%
\TextAt{12}{151.887}{510.65}{(5.9).}%
\TextAt{8}{187.159}{510.65}{Let}%
\TextAt{13}{208.085}{510.65}{L}%
\TextAt{16}{219.37}{510.65}{\Character{106}}%
\TextAt{13}{226.012}{510.65}{K}%
\TextAt{8}{240.683}{510.65}{be}%
\TextAt{8}{256.564}{510.65}{a}%
\TextAt{8}{266.276}{510.65}{finite}%
\TextAt{8}{296.144}{510.65}{Galois}%
\TextAt{8}{332.014}{510.65}{extension}%
\TextAt{8}{383.736}{510.65}{and}%
\TextAt{13}{406.455}{510.65}{K}%
\TextAt{3}{417.272}{514.989}{\Character{48}}%
\TextAt{8}{423.921}{510.65}{an}%
\TextAt{8}{440.137}{510.65}{intermediate}%
\TextAt{8}{508.058}{510.65}{field.}%
\TextAt{8}{79.226}{493.604}{Then}%
\TextAt{8}{109.784}{493.604}{we}%
\TextAt{8}{127.011}{493.604}{have}%
\TextAt{8}{154.006}{493.604}{the}%
\TextAt{8}{174.162}{493.604}{commutative}%
\TextAt{8}{244.088}{493.604}{diagram}%
\TextAt{13}{222.382}{462.604}{G}%
\TextAt{9}{231.615}{462.604}{\Character{40}}%
\TextAt{13}{236.168}{462.604}{L}%
\TextAt{16}{247.453}{462.604}{\Character{106}}%
\TextAt{13}{254.095}{462.604}{K}%
\TextAt{3}{264.913}{466.943}{\Character{48}}%
\TextAt{9}{267.708}{462.604}{\Character{41}}%
\TextAt{7}{272.26}{466.943}{ab}%
\TextAt{13}{320.69}{462.604}{A}%
\TextAt{21}{329.465}{460.811}{K}%
\TextAt{23}{337.17}{463.08}{\Character{48}}%
\TextAt{9}{340.352}{462.604}{\Character{47}}%
\TextAt{13}{346.205}{462.604}{N}%
\TextAt{21}{355.579}{460.744}{L}%
\TextAt{3}{361.338}{460.744}{\Character{106}}%
\TextAt{21}{363.69}{460.744}{K}%
\TextAt{23}{371.395}{463.014}{\Character{48}}%
\TextAt{13}{374.578}{462.604}{A}%
\TextAt{21}{383.353}{460.811}{L}%
\TextAt{13}{223.78}{418.603}{G}%
\TextAt{9}{233.013}{418.603}{\Character{40}}%
\TextAt{13}{237.566}{418.603}{L}%
\TextAt{16}{248.851}{418.603}{\Character{106}}%
\TextAt{13}{255.493}{418.603}{K}%
\TextAt{9}{266.311}{418.603}{\Character{41}}%
\TextAt{7}{270.863}{422.942}{ab}%
\TextAt{13}{321.749}{418.603}{A}%
\TextAt{21}{330.524}{416.81}{K}%
\TextAt{9}{338.727}{418.603}{\Character{47}}%
\TextAt{13}{344.58}{418.603}{N}%
\TextAt{21}{353.954}{416.744}{L}%
\TextAt{3}{359.713}{416.744}{\Character{106}}%
\TextAt{21}{362.065}{416.744}{K}%
\TextAt{13}{370.268}{418.603}{A}%
\TextAt{21}{379.043}{416.81}{L}%
\TextAt{13}{385.3}{418.603}{\Character{59}}%
\NativeDiagram{%
\path[draw=black,line width=0.478bp] (215.046,-308.407) -- (242.862,-308.407);%
\path[draw=black,line width=0.478bp,line cap=round,line join=round] (240.615,-305.538) .. controls (241.069,-307.26) and (242.081,-308.072) .. (243.101,-308.407) .. controls (242.081,-308.742) and (241.069,-309.555) .. (240.615,-311.276);%
}%
\TextAt{21}{290.755}{472.194}{r}%
\TextAt{22}{294.576}{469.903}{L}%
\TextAt{23}{299.635}{469.903}{\Character{106}}%
\TextAt{22}{301.904}{469.903}{K}%
\TextAt{23}{308.449}{471.619}{\Character{48}}%
\NativeDiagram{%
\path[draw=black,line width=0.478bp] (180.038,-340.949) -- (180.038,-319.442);%
\path[draw=black,line width=0.478bp,line cap=round,line join=round] (177.169,-321.69) .. controls (178.891,-321.235) and (179.704,-320.223) .. (180.038,-319.203) .. controls (180.373,-320.223) and (181.186,-321.235) .. (182.908,-321.69);%
}%
\TextAt{7}{235.817}{441.406}{Ver}%
\NativeDiagram{%
\path[draw=black,line width=0.478bp] (213.649,-352.409) -- (243.921,-352.409);%
\path[draw=black,line width=0.478bp,line cap=round,line join=round] (241.673,-349.54) .. controls (242.127,-351.261) and (243.14,-352.074) .. (244.16,-352.409) .. controls (243.14,-352.743) and (242.127,-353.556) .. (241.673,-355.278);%
}%
\TextAt{21}{291.927}{415.344}{r}%
\TextAt{22}{295.748}{413.606}{L}%
\TextAt{23}{300.807}{413.606}{\Character{106}}%
\TextAt{22}{303.076}{413.606}{K}%
\NativeDiagram{%
\path[draw=black,line width=0.478bp] (283.153,-341.852) -- (283.153,-320.306);%
\path[draw=black,line width=0.478bp,line cap=round,line join=round] (280.284,-322.553) .. controls (282.005,-322.099) and (282.818,-321.086) .. (283.153,-320.067) .. controls (283.488,-321.086) and (284.3,-322.099) .. (286.022,-322.553);%
}%
\TextAt{8}{79.226}{388.288}{where}%
\TextAt{8}{113.035}{388.288}{the}%
\TextAt{8}{133.192}{388.288}{arrow}%
\TextAt{8}{166.045}{388.288}{on}%
\TextAt{8}{182.304}{388.288}{the}%
\TextAt{8}{202.46}{388.288}{right}%
\TextAt{8}{230.758}{388.288}{is}%
\TextAt{8}{242.522}{388.288}{induced}%
\TextAt{8}{286.086}{388.288}{by}%
\TextAt{8}{302.346}{388.288}{inclusion.}%
\TextAt{19}{79.226}{353.309}{Proof.}%
\TextAt{8}{116}{353.309}{Let}%
\TextAt{8}{137.321}{353.309}{us}%
\TextAt{8}{152.689}{353.309}{write}%
\TextAt{8}{182.952}{353.309}{temporarily}%
\TextAt{13}{247.049}{353.309}{G}%
\TextAt{9}{260.203}{353.309}{\Character{61}}%
\TextAt{13}{273.228}{353.309}{G}%
\TextAt{9}{282.462}{353.309}{\Character{40}}%
\TextAt{20}{288}{356.331}{\Character{101}}%
\TextAt{13}{287.014}{353.309}{L}%
\TextAt{16}{298.899}{353.309}{\Character{106}}%
\TextAt{13}{306.14}{353.309}{K}%
\TextAt{9}{316.957}{353.309}{\Character{41}}%
\TextAt{8}{325.759}{353.309}{and}%
\TextAt{13}{348.873}{353.309}{H}%
\TextAt{9}{363.416}{353.309}{\Character{61}}%
\TextAt{13}{376.441}{353.309}{G}%
\TextAt{9}{385.675}{353.309}{\Character{40}}%
\TextAt{20}{391.214}{356.331}{\Character{101}}%
\TextAt{13}{390.227}{353.309}{L}%
\TextAt{16}{402.112}{353.309}{\Character{106}}%
\TextAt{13}{409.353}{353.309}{K}%
\TextAt{3}{420.171}{357.648}{\Character{48}}%
\TextAt{9}{422.966}{353.309}{\Character{41}}%
\TextAt{8}{427.518}{353.309}{.}%
\TextAt{8}{437.022}{353.309}{Let}%
\TextAt{13}{458.343}{353.309}{\Character{27}}%
\NativeMathAt{469.345}{353.309}{11.955}{\in}%
\TextAt{13}{481.236}{353.309}{G}%
\TextAt{9}{490.469}{353.309}{\Character{40}}%
\TextAt{13}{495.022}{353.309}{L}%
\TextAt{16}{506.907}{353.309}{\Character{106}}%
\TextAt{13}{514.148}{353.309}{K}%
\TextAt{9}{524.965}{353.309}{\Character{41}}%
\TextAt{8}{529.518}{353.309}{,}%
\TextAt{8}{79.226}{336.263}{and}%
\TextAt{8}{103.041}{336.263}{let}%
\TextAt{9}{121.613}{336.263}{\Character{22}}%
\TextAt{13}{120.998}{336.263}{\Character{27}}%
\TextAt{8}{133.03}{336.263}{be}%
\TextAt{8}{150.006}{336.263}{a}%
\TextAt{8}{160.814}{336.263}{preimage}%
\TextAt{8}{211.946}{336.263}{in}%
\TextAt{8}{226.651}{336.263}{Frob}%
\TextAt{9}{251.207}{336.263}{\Character{40}}%
\TextAt{20}{256.746}{339.285}{\Character{101}}%
\TextAt{13}{255.76}{336.263}{L}%
\TextAt{16}{268.838}{336.263}{\Character{106}}%
\TextAt{13}{277.273}{336.263}{K}%
\TextAt{9}{288.09}{336.263}{\Character{41}}%
\TextAt{8}{297.592}{336.263}{with}%
\TextAt{8}{325.305}{336.263}{fixed}%
\TextAt{8}{354.643}{336.263}{field}%
\TextAt{9}{381.052}{336.263}{\Character{6}}%
\TextAt{8}{394.456}{336.263}{and}%
\TextAt{13}{418.271}{336.263}{S}%
\TextAt{9}{431.28}{336.263}{\Character{61}}%
\TextAt{13}{445.498}{336.263}{G}%
\TextAt{9}{454.732}{336.263}{\Character{40}}%
\TextAt{20}{460.271}{339.285}{\Character{101}}%
\TextAt{13}{459.284}{336.263}{L}%
\TextAt{16}{472.363}{336.263}{\Character{106}}%
\TextAt{9}{480.798}{336.263}{\Character{6}\Character{41}}%
\TextAt{9}{498.898}{336.263}{\Character{61}}%
\Vector{q 0 G 0 g 0.478 w 441.137 -427.097 m 457.518 -427.097 l S Q}%
\TextAt{16}{513.137}{336.263}{\Character{104}}%
\TextAt{9}{518.401}{336.263}{\Character{22}}%
\TextAt{13}{517.786}{336.263}{\Character{27}}%
\TextAt{16}{524.868}{336.263}{\Character{105}}%
\TextAt{8}{529.518}{336.263}{.}%
\TextAt{8}{79.226}{319.217}{We}%
\TextAt{8}{99.451}{319.217}{consider}%
\TextAt{8}{145.117}{319.217}{the}%
\TextAt{8}{165.354}{319.217}{double}%
\TextAt{8}{203.153}{319.217}{coset}%
\TextAt{8}{232.56}{319.217}{decomposition}%
\TextAt{13}{309.776}{319.217}{G}%
\TextAt{9}{322.468}{319.217}{\Character{61}}%
\TextAt{20}{335.031}{328.184}{\Character{83}}%
\TextAt{21}{344.993}{315.73}{\Character{28}}%
\TextAt{13}{352.162}{319.217}{S}%
\TextAt{13}{360.052}{319.217}{\Character{28}}%
\TextAt{13}{366.46}{319.217}{H}%
\TextAt{8}{381.077}{319.217}{and}%
\TextAt{8}{403.921}{319.217}{put}%
\TextAt{13}{425.461}{319.217}{S}%
\TextAt{21}{432.66}{317.424}{\Character{28}}%
\TextAt{9}{441.295}{319.217}{\Character{61}}%
\TextAt{13}{453.858}{319.217}{\Character{28}}%
\TextAt{3}{460.277}{323.556}{\Character{0}}%
\TextAt{17}{466.863}{323.556}{1}%
\TextAt{13}{471.596}{319.217}{S}%
\TextAt{13}{479.486}{319.217}{\Character{28}}%
\TextAt{16}{488.622}{319.217}{\Character{92}}%
\TextAt{13}{499.304}{319.217}{H}%
\TextAt{8}{513.904}{319.217}{and}%
\TextAt{9}{79.841}{302.171}{\Character{22}}%
\TextAt{13}{79.226}{302.171}{\Character{27}}%
\TextAt{21}{85.88}{300.378}{\Character{28}}%
\TextAt{9}{95.871}{302.171}{\Character{61}}%
\TextAt{13}{109.791}{302.171}{\Character{28}}%
\TextAt{3}{116.21}{306.51}{\Character{0}}%
\TextAt{17}{122.797}{306.51}{1}%
\TextAt{9}{128.144}{302.171}{\Character{22}}%
\TextAt{13}{127.529}{302.171}{\Character{27}}%
\TextAt{21}{134.611}{306.51}{f}%
\TextAt{17}{139.558}{306.51}{\Character{40}}%
\TextAt{21}{142.852}{306.51}{\Character{28}}%
\TextAt{17}{147.53}{306.51}{\Character{41}}%
\TextAt{13}{151.321}{302.171}{\Character{28}}%
\TextAt{8}{162.515}{302.171}{as}%
\TextAt{8}{177.762}{302.171}{above.}%
\TextAt{8}{218.119}{302.171}{Let}%
\TextAt{9}{242.631}{305.193}{\Character{22}}%
\TextAt{13}{239.965}{302.171}{G}%
\TextAt{9}{254.014}{302.171}{\Character{61}}%
\TextAt{13}{267.934}{302.171}{G}%
\TextAt{9}{277.168}{302.171}{\Character{40}}%
\TextAt{13}{281.72}{302.171}{L}%
\TextAt{16}{294.5}{302.171}{\Character{106}}%
\TextAt{13}{302.636}{302.171}{K}%
\TextAt{9}{313.454}{302.171}{\Character{41}}%
\TextAt{8}{318.006}{302.171}{,}%
\TextAt{9}{329.287}{305.193}{\Character{22}}%
\TextAt{13}{326.252}{302.171}{H}%
\TextAt{9}{341.689}{302.171}{\Character{61}}%
\TextAt{13}{355.609}{302.171}{G}%
\TextAt{9}{364.843}{302.171}{\Character{40}}%
\TextAt{13}{369.395}{302.171}{L}%
\TextAt{16}{382.175}{302.171}{\Character{106}}%
\TextAt{13}{390.311}{302.171}{K}%
\TextAt{3}{401.129}{306.51}{\Character{48}}%
\TextAt{9}{403.924}{302.171}{\Character{41}}%
\TextAt{8}{408.476}{302.171}{,}%
\TextAt{9}{418.718}{305.193}{\Character{22}}%
\TextAt{13}{416.722}{302.171}{S}%
\TextAt{9}{429.432}{302.171}{\Character{61}}%
\TextAt{16}{443.352}{302.171}{\Character{104}}%
\TextAt{13}{448.002}{302.171}{\Character{27}}%
\TextAt{16}{455.084}{302.171}{\Character{105}}%
\TextAt{8}{459.733}{302.171}{,}%
\TextAt{9}{468.587}{302.171}{\Character{22}}%
\TextAt{13}{467.979}{302.171}{\Character{28}}%
\TextAt{9}{479.213}{302.171}{\Character{61}}%
\TextAt{13}{493.133}{302.171}{\Character{28}}%
\TextAt{16}{499.552}{302.171}{\Character{106}}%
\TextAt{21}{502.873}{300.378}{L}%
\TextAt{8}{513.904}{302.171}{and}%
\TextAt{13}{79.226}{285.125}{\Character{27}}%
\TextAt{21}{85.88}{283.332}{\Character{28}}%
\TextAt{9}{94.377}{285.125}{\Character{61}}%
\TextAt{9}{107.408}{285.125}{\Character{22}}%
\TextAt{13}{106.802}{285.125}{\Character{27}}%
\TextAt{21}{113.456}{283.332}{\Character{28}}%
\TextAt{16}{118.632}{285.125}{\Character{106}}%
\TextAt{21}{121.953}{283.332}{L}%
\TextAt{8}{128.21}{285.125}{.}%
\TextAt{8}{136.659}{285.125}{Then}%
\TextAt{8}{167.216}{285.125}{we}%
\TextAt{8}{184.444}{285.125}{obviously}%
\TextAt{8}{236.221}{285.125}{have}%
\TextAt{9}{274.441}{259.146}{\Character{22}}%
\TextAt{13}{271.776}{256.124}{G}%
\TextAt{9}{284.33}{256.124}{\Character{61}}%
\TextAt{20}{296.755}{267.482}{\Character{91}}%
\TextAt{17}{301.515}{242.376}{\Character{22}}%
\TextAt{21}{301.058}{242.376}{\Character{28}}%
\TextAt{9}{314.028}{259.146}{\Character{22}}%
\TextAt{13}{312.032}{256.124}{S}%
\TextAt{9}{320.535}{256.124}{\Character{22}}%
\TextAt{13}{319.927}{256.124}{\Character{28}}%
\TextAt{9}{329.381}{259.146}{\Character{22}}%
\TextAt{13}{326.346}{256.124}{H}%
\TextAt{13}{336.962}{256.124}{\Character{59}}%
\TextAt{8}{79.226}{219.812}{and}%
\TextAt{8}{101.989}{219.812}{therefore}%
\TextAt{8}{184.556}{202.766}{Ver}%
\TextAt{9}{203.087}{202.766}{\Character{40}}%
\TextAt{13}{207.639}{202.766}{\Character{27}}%
\TextAt{8}{218.042}{202.766}{mod}%
\TextAt{13}{243.48}{202.766}{G}%
\TextAt{9}{252.714}{202.766}{\Character{40}}%
\TextAt{13}{257.266}{202.766}{L}%
\TextAt{16}{268.552}{202.766}{\Character{106}}%
\TextAt{13}{275.193}{202.766}{K}%
\TextAt{9}{286.011}{202.766}{\Character{41}}%
\TextAt{3}{290.563}{207.703}{\Character{48}}%
\TextAt{9}{293.358}{202.766}{\Character{41}}%
\TextAt{9}{301.213}{202.766}{\Character{61}}%
\TextAt{20}{313.657}{214.124}{\Character{89}}%
\TextAt{21}{318.956}{189.018}{\Character{28}}%
\TextAt{13}{330.926}{202.766}{\Character{27}}%
\TextAt{21}{337.579}{200.973}{\Character{28}}%
\TextAt{8}{346.076}{202.766}{mod}%
\TextAt{13}{371.514}{202.766}{G}%
\TextAt{9}{380.748}{202.766}{\Character{40}}%
\TextAt{13}{385.3}{202.766}{L}%
\TextAt{16}{396.586}{202.766}{\Character{106}}%
\TextAt{13}{403.227}{202.766}{K}%
\TextAt{3}{414.045}{207.703}{\Character{48}}%
\TextAt{9}{416.84}{202.766}{\Character{41}}%
\TextAt{3}{421.392}{207.703}{\Character{48}}%
\TextAt{13}{424.187}{202.766}{\Character{58}}%
\TextAt{8}{79.226}{171.384}{For}%
\TextAt{8}{100.508}{171.384}{every}%
\TextAt{13}{131.712}{171.384}{\Character{28}}%
\TextAt{8}{138.131}{171.384}{,}%
\TextAt{8}{145.67}{171.384}{let}%
\TextAt{13}{162.887}{171.384}{\Character{33}}%
\TextAt{21}{170.176}{169.591}{\Character{28}}%
\TextAt{8}{179.562}{171.384}{vary}%
\TextAt{8}{205.9}{171.384}{over}%
\TextAt{8}{231.259}{171.384}{a}%
\TextAt{8}{241.326}{171.384}{system}%
\TextAt{8}{280.457}{171.384}{of}%
\TextAt{8}{294.099}{171.384}{right}%
\TextAt{8}{322.709}{171.384}{coset}%
\TextAt{8}{352.347}{171.384}{representatives}%
\TextAt{8}{432.46}{171.384}{of}%
\TextAt{13}{446.102}{171.384}{H}%
\TextAt{9}{456.724}{171.384}{\Character{47}}%
\TextAt{13}{462.577}{171.384}{S}%
\TextAt{21}{469.777}{169.591}{\Character{28}}%
\TextAt{8}{474.953}{171.384}{.}%
\TextAt{8}{484.338}{171.384}{Then}%
\TextAt{8}{515.207}{171.384}{one}%
\TextAt{8}{79.226}{154.338}{has}%
\TextAt{13}{202.662}{137.292}{H}%
\TextAt{9}{216.605}{137.292}{\Character{61}}%
\TextAt{20}{229.031}{148.649}{\Character{91}}%
\TextAt{21}{230.705}{123.543}{\Character{33}}%
\TextAt{22}{235.971}{122.547}{\Character{28}}%
\TextAt{13}{244.307}{137.292}{S}%
\TextAt{21}{251.506}{135.499}{\Character{28}}%
\TextAt{13}{256.683}{137.292}{\Character{33}}%
\TextAt{21}{263.972}{135.499}{\Character{28}}%
\TextAt{8}{293.059}{137.292}{and}%
\TextAt{13}{335.835}{137.292}{G}%
\TextAt{9}{348.389}{137.292}{\Character{61}}%
\TextAt{20}{362.42}{148.649}{\Character{91}}%
\TextAt{21}{360.814}{123.543}{\Character{28}}%
\TextAt{21}{365.014}{123.543}{\Character{59}\Character{33}}%
\TextAt{22}{372.64}{122.547}{\Character{28}}%
\TextAt{13}{379.302}{137.292}{S}%
\TextAt{13}{387.192}{137.292}{\Character{28}}%
\TextAt{13}{393.6}{137.292}{\Character{33}}%
\TextAt{21}{400.905}{135.499}{\Character{28}}%
\TextAt{13}{406.082}{137.292}{\Character{58}}%
\TextAt{8}{79.226}{103.852}{Let}%
\TextAt{9}{100.235}{103.852}{\Character{6}}%
\TextAt{21}{108.689}{102.058}{\Character{28}}%
\TextAt{8}{117.802}{103.852}{be}%
\TextAt{8}{133.766}{103.852}{the}%
\TextAt{8}{153.962}{103.852}{fixed}%
\TextAt{8}{182.287}{103.852}{field}%
\TextAt{8}{207.683}{103.852}{of}%
\TextAt{9}{221.667}{103.852}{\Character{22}}%
\TextAt{13}{221.052}{103.851}{\Character{27}}%
\TextAt{21}{227.706}{102.058}{\Character{28}}%
\TextAt{8}{232.882}{103.852}{,}%
\TextAt{8}{240.081}{103.852}{i.e.,}%
\TextAt{8}{262.235}{103.852}{the}%
\TextAt{8}{282.431}{103.852}{fixed}%
\TextAt{8}{310.756}{103.852}{field}%
\TextAt{8}{336.152}{103.852}{of}%
\TextAt{13}{349.522}{103.852}{S}%
\TextAt{21}{356.721}{102.058}{\Character{28}}%
\TextAt{8}{361.897}{103.852}{.}%
\TextAt{9}{370.464}{103.852}{\Character{6}}%
\TextAt{21}{378.918}{108.19}{\Character{28}}%
\TextAt{8}{388.031}{103.852}{is}%
\TextAt{8}{399.834}{103.852}{the}%
\TextAt{8}{420.03}{103.852}{fixed}%
\TextAt{8}{448.355}{103.852}{field}%
\TextAt{8}{473.752}{103.852}{of}%
\TextAt{13}{487.121}{103.852}{\Character{28}}%
\TextAt{3}{493.54}{108.19}{\Character{0}}%
\TextAt{17}{500.126}{108.19}{1}%
\TextAt{9}{505.473}{103.852}{\Character{22}}%
\TextAt{13}{504.859}{103.851}{\Character{27}}%
\TextAt{13}{511.925}{103.851}{\Character{28}}%
\TextAt{8}{522.297}{103.852}{so}%
\TextAt{8}{79.226}{86.806}{that}%
\TextAt{9}{105.057}{86.806}{\Character{6}}%
\TextAt{21}{113.511}{85.012}{\Character{28}}%
\TextAt{16}{122.795}{86.806}{\Character{106}}%
\TextAt{9}{130.223}{86.806}{\Character{6}}%
\TextAt{21}{138.677}{91.144}{\Character{28}}%
\TextAt{8}{148.212}{86.806}{is}%
\TextAt{8}{160.438}{86.806}{the}%
\TextAt{8}{181.056}{86.806}{unramified}%
\TextAt{8}{240.05}{86.806}{subextension}%
\TextAt{8}{310.222}{86.806}{of}%
\TextAt{8}{324.013}{86.806}{degree}%
\TextAt{13}{360.89}{86.806}{f}%
\TextAt{9}{367.937}{86.806}{\Character{40}}%
\TextAt{13}{372.489}{86.806}{\Character{28}}%
\TextAt{9}{378.908}{86.806}{\Character{41}}%
\TextAt{8}{387.819}{86.806}{in}%
\TextAt{20}{402.92}{89.828}{\Character{101}}%
\TextAt{13}{401.933}{86.806}{L}%
\TextAt{16}{414.005}{86.806}{\Character{106}}%
\TextAt{9}{421.433}{86.806}{\Character{6}}%
\TextAt{21}{429.888}{91.144}{\Character{28}}%
\TextAt{8}{435.064}{86.806}{.}%
\TextAt{8}{444.897}{86.806}{If}%
\TextAt{8}{457.051}{86.806}{now}%
\TextAt{13}{481.901}{86.806}{\Character{25}}%
\NativeMathAt{493.077}{86.806}{11.955}{\in}%
\TextAt{13}{505.155}{86.806}{A}%
\TextAt{17}{513.93}{85.012}{\Character{6}}%
\TextAt{8}{524.903}{86.806}{is}%
\TextAt{8}{297.211}{43.033}{272}%
\Destination{78.23}{809.815}{nav.page.272}%
\Destination{79.36}{776.155}{nav.print.297}%
\Destination{79.36}{186.355}{nav.print.298}%
\Destination{79.36}{524.655}{nav.number.4.5.9}%
}

% Page 285 (body page 273)
\Page{%
\TextAt{2}{79.226}{796.943}{IV.}%
\TextAt{2}{98.012}{796.943}{Abstract}%
\TextAt{2}{143.252}{796.943}{Class}%
\TextAt{2}{171.844}{796.943}{Field}%
\TextAt{2}{199.575}{796.943}{Theory}%
\Vector{q 0 G 0 g 0.299 w 7.226 18.729 m 460.769 18.729 l S Q}%
\TextAt{8}{79.226}{761.124}{a}%
\TextAt{8}{89.434}{761.124}{prime}%
\TextAt{8}{123.05}{761.124}{element}%
\TextAt{8}{166.744}{761.124}{of}%
\TextAt{9}{180.526}{761.124}{\Character{6}}%
\TextAt{8}{188.981}{761.124}{,}%
\TextAt{8}{196.695}{761.124}{then}%
\TextAt{13}{223.807}{761.124}{\Character{25}}%
\TextAt{21}{230.877}{765.462}{\Character{28}}%
\NativeMathAt{240.145}{761.124}{11.955}{\in}%
\TextAt{13}{252.206}{761.124}{A}%
\TextAt{17}{260.982}{759.33}{\Character{6}}%
\TextAt{22}{267.098}{761.6}{\Character{28}}%
\TextAt{8}{276.615}{761.124}{is}%
\TextAt{8}{288.831}{761.124}{a}%
\TextAt{8}{299.039}{761.124}{prime}%
\TextAt{8}{332.655}{761.124}{element}%
\TextAt{8}{376.349}{761.124}{of}%
\TextAt{9}{390.131}{761.124}{\Character{6}}%
\TextAt{21}{398.585}{765.462}{\Character{28}}%
\TextAt{8}{403.762}{761.124}{,}%
\TextAt{8}{411.477}{761.124}{and}%
\TextAt{8}{434.691}{761.124}{thus}%
\TextAt{8}{460.895}{761.124}{also}%
\TextAt{8}{484.827}{761.124}{of}%
\TextAt{9}{498.61}{761.124}{\Character{6}}%
\TextAt{21}{507.064}{759.33}{\Character{28}}%
\TextAt{8}{512.241}{761.124}{.}%
\TextAt{8}{522.046}{761.124}{In}%
\TextAt{8}{79.226}{744.078}{view}%
\TextAt{8}{106.209}{744.078}{of}%
\TextAt{8}{119.539}{744.078}{the}%
\TextAt{8}{139.695}{744.078}{above}%
\TextAt{8}{172.871}{744.078}{double}%
\TextAt{8}{210.59}{744.078}{coset}%
\TextAt{8}{239.916}{744.078}{decomposition,}%
\TextAt{8}{320.302}{744.078}{we}%
\TextAt{8}{337.53}{744.078}{therefore}%
\TextAt{8}{386.629}{744.078}{find}%
\TextAt{13}{161.672}{707.885}{N}%
\TextAt{17}{171.046}{706.026}{\Character{6}}%
\TextAt{3}{177.162}{706.026}{\Character{106}}%
\TextAt{21}{179.514}{706.026}{K}%
\TextAt{9}{187.717}{707.885}{\Character{40}}%
\TextAt{13}{192.269}{707.885}{\Character{25}}%
\TextAt{9}{199.338}{707.885}{\Character{41}}%
\TextAt{9}{207.193}{707.885}{\Character{61}}%
\TextAt{20}{220.246}{719.243}{\Character{89}}%
\TextAt{21}{219.637}{694.137}{\Character{28}}%
\TextAt{21}{223.837}{694.137}{\Character{59}\Character{33}}%
\TextAt{22}{231.463}{693.141}{\Character{28}}%
\TextAt{13}{238.124}{707.885}{\Character{25}}%
\TextAt{21}{245.194}{712.821}{\Character{28}}%
\TextAt{21}{249.865}{712.821}{\Character{33}}%
\TextAt{22}{255.138}{711.825}{\Character{28}}%
\TextAt{9}{263.626}{707.885}{\Character{61}}%
\TextAt{20}{276.051}{719.243}{\Character{89}}%
\TextAt{21}{281.35}{694.137}{\Character{28}}%
\TextAt{20}{293.32}{728.329}{\Character{32}}%
\TextAt{20}{302.784}{719.243}{\Character{89}}%
\TextAt{21}{305.455}{694.137}{\Character{33}}%
\TextAt{22}{310.721}{693.141}{\Character{28}}%
\TextAt{9}{318.06}{707.885}{\Character{40}}%
\TextAt{13}{322.613}{707.885}{\Character{25}}%
\TextAt{21}{329.682}{712.822}{\Character{28}}%
\TextAt{9}{334.859}{707.885}{\Character{41}}%
\TextAt{21}{339.411}{712.822}{\Character{33}}%
\TextAt{22}{344.676}{711.825}{\Character{28}}%
\TextAt{20}{349.844}{728.329}{\Character{33}}%
\TextAt{9}{362.629}{707.885}{\Character{61}}%
\TextAt{20}{375.055}{719.243}{\Character{89}}%
\TextAt{21}{380.354}{694.137}{\Character{28}}%
\TextAt{13}{392.323}{707.885}{N}%
\TextAt{17}{401.697}{706.026}{\Character{6}}%
\TextAt{22}{407.813}{705.029}{\Character{28}}%
\TextAt{3}{412.482}{706.026}{\Character{106}}%
\TextAt{21}{414.835}{706.026}{K}%
\TextAt{23}{422.539}{708.295}{\Character{48}}%
\TextAt{9}{425.722}{707.885}{\Character{40}}%
\TextAt{13}{430.274}{707.885}{\Character{25}}%
\TextAt{21}{437.343}{712.821}{\Character{28}}%
\TextAt{9}{442.52}{707.885}{\Character{41}}%
\TextAt{13}{447.072}{707.885}{\Character{59}}%
\TextAt{8}{79.226}{666.784}{and}%
\TextAt{8}{101.989}{666.784}{since}%
\TextAt{9}{131.272}{666.784}{\Character{22}}%
\TextAt{13}{130.657}{666.784}{\Character{27}}%
\TextAt{21}{137.311}{664.991}{\Character{28}}%
\NativeMathAt{145.808}{666.784}{11.955}{\in}%
\TextAt{8}{157.099}{666.784}{Frob}%
\TextAt{9}{181.655}{666.784}{\Character{40}}%
\TextAt{20}{187.194}{669.806}{\Character{101}}%
\TextAt{13}{186.207}{666.784}{L}%
\TextAt{16}{197.493}{666.784}{\Character{106}}%
\TextAt{13}{204.134}{666.784}{K}%
\TextAt{3}{214.952}{671.122}{\Character{48}}%
\TextAt{9}{217.747}{666.784}{\Character{41}}%
\TextAt{8}{226.197}{666.784}{is}%
\TextAt{8}{237.961}{666.784}{a}%
\TextAt{8}{247.716}{666.784}{preimage}%
\TextAt{8}{297.796}{666.784}{of}%
\TextAt{13}{311.126}{666.784}{\Character{27}}%
\TextAt{21}{317.78}{664.991}{\Character{28}}%
\NativeMathAt{326.277}{666.784}{11.955}{\in}%
\TextAt{13}{337.568}{666.784}{G}%
\TextAt{9}{346.802}{666.784}{\Character{40}}%
\TextAt{13}{351.354}{666.784}{L}%
\TextAt{16}{362.639}{666.784}{\Character{106}}%
\TextAt{13}{369.281}{666.784}{K}%
\TextAt{3}{380.099}{671.122}{\Character{48}}%
\TextAt{9}{382.894}{666.784}{\Character{41}}%
\TextAt{8}{387.446}{666.784}{,}%
\TextAt{8}{394.595}{666.784}{it}%
\TextAt{8}{406.299}{666.784}{follows}%
\TextAt{8}{444.735}{666.784}{that}%
\TextAt{13}{119.822}{629.922}{r}%
\TextAt{21}{125.097}{628.062}{L}%
\TextAt{3}{130.856}{628.062}{\Character{106}}%
\TextAt{21}{133.208}{628.062}{K}%
\TextAt{9}{141.411}{629.922}{\Character{40}}%
\TextAt{13}{145.963}{629.922}{\Character{27}}%
\TextAt{9}{153.046}{629.922}{\Character{41}}%
\NativeMathAt{160.919}{629.922}{11.955}{\equiv}%
\TextAt{20}{173.538}{641.279}{\Character{89}}%
\TextAt{21}{178.837}{616.173}{\Character{28}}%
\TextAt{13}{190.807}{629.922}{r}%
\TextAt{21}{196.082}{628.062}{L}%
\TextAt{3}{201.841}{628.062}{\Character{106}}%
\TextAt{21}{204.194}{628.062}{K}%
\TextAt{23}{211.898}{630.331}{\Character{48}}%
\TextAt{9}{215.081}{629.922}{\Character{40}}%
\TextAt{13}{219.633}{629.922}{\Character{27}}%
\TextAt{21}{226.286}{628.129}{\Character{28}}%
\TextAt{9}{231.463}{629.922}{\Character{41}}%
\NativeMathAt{239.336}{629.922}{11.955}{\equiv}%
\TextAt{13}{251.955}{629.922}{r}%
\TextAt{21}{257.231}{628.062}{L}%
\TextAt{3}{262.99}{628.062}{\Character{106}}%
\TextAt{21}{265.342}{628.062}{K}%
\TextAt{23}{273.047}{630.331}{\Character{48}}%
\TextAt{20}{278.222}{650.365}{\Character{32}}%
\TextAt{20}{287.686}{641.279}{\Character{89}}%
\TextAt{21}{292.985}{616.173}{\Character{28}}%
\TextAt{13}{304.955}{629.922}{\Character{27}}%
\TextAt{21}{311.608}{628.129}{\Character{28}}%
\TextAt{20}{316.785}{650.365}{\Character{33}}%
\NativeMathAt{329.57}{629.922}{11.955}{\equiv}%
\TextAt{13}{342.189}{629.922}{r}%
\TextAt{21}{347.465}{628.062}{L}%
\TextAt{3}{353.224}{628.062}{\Character{106}}%
\TextAt{21}{355.576}{628.062}{K}%
\TextAt{23}{363.281}{630.331}{\Character{48}}%
\TextAt{9}{366.463}{629.922}{\Character{40}}%
\TextAt{8}{371.015}{629.922}{Ver}%
\TextAt{9}{389.546}{629.922}{\Character{40}}%
\TextAt{13}{394.098}{629.922}{\Character{27}}%
\TextAt{8}{404.501}{629.922}{mod}%
\TextAt{13}{429.939}{629.922}{G}%
\TextAt{9}{439.173}{629.922}{\Character{40}}%
\TextAt{13}{443.725}{629.922}{L}%
\TextAt{16}{455.011}{629.922}{\Character{106}}%
\TextAt{13}{461.652}{629.922}{K}%
\TextAt{9}{472.47}{629.922}{\Character{41}}%
\TextAt{3}{477.022}{634.858}{\Character{48}}%
\TextAt{9}{479.817}{629.922}{\Character{41}\Character{41}}%
\TextAt{13}{488.922}{629.922}{\Character{58}}%
\Vector{q 0 G 0 g 0.399 w 452.334 -180.044 m 452.334 -171.974 l S Q}%
\Vector{q 0 G 0 g 0.399 w 452.534 -172.173 m 459.707 -172.173 l S Q}%
\Vector{q 0 G 0 g 0.399 w 452.534 -179.844 m 459.707 -179.844 l S Q}%
\Vector{q 0 G 0 g 0.399 w 459.906 -180.044 m 459.906 -171.974 l S Q}%
\TextAt{12}{79.226}{558.978}{Exercise}%
\TextAt{12}{131.552}{558.978}{IV.5.1.}%
\TextAt{8}{174.548}{558.978}{Let}%
\TextAt{13}{194.757}{558.978}{L}%
\TextAt{16}{206.042}{558.978}{\Character{106}}%
\TextAt{13}{212.684}{558.978}{K}%
\TextAt{8}{226.639}{558.978}{be}%
\TextAt{8}{241.803}{558.978}{abelian}%
\TextAt{8}{281.69}{558.978}{and}%
\TextAt{8}{303.693}{558.978}{totally}%
\TextAt{8}{340.34}{558.978}{ramified,}%
\TextAt{8}{388.509}{558.978}{and}%
\TextAt{8}{410.512}{558.978}{let}%
\TextAt{13}{426.656}{558.978}{\Character{25}}%
\TextAt{8}{436.862}{558.978}{be}%
\TextAt{8}{452.026}{558.978}{a}%
\TextAt{8}{461.022}{558.978}{prime}%
\TextAt{8}{493.425}{558.978}{element}%
\TextAt{8}{79.226}{541.932}{of}%
\TextAt{13}{92.556}{541.932}{A}%
\TextAt{21}{101.332}{540.139}{L}%
\TextAt{8}{107.589}{541.932}{.}%
\TextAt{8}{116.037}{541.932}{If}%
\TextAt{8}{127.729}{541.932}{then}%
\TextAt{13}{154.389}{541.932}{\Character{27}}%
\NativeMathAt{164.792}{541.932}{11.955}{\in}%
\TextAt{13}{176.083}{541.932}{G}%
\TextAt{9}{185.317}{541.932}{\Character{40}}%
\TextAt{13}{189.869}{541.932}{L}%
\TextAt{16}{201.155}{541.932}{\Character{106}}%
\TextAt{13}{207.796}{541.932}{K}%
\TextAt{9}{218.614}{541.932}{\Character{41}}%
\TextAt{8}{227.064}{541.932}{and}%
\TextAt{13}{274.881}{524.886}{y}%
\TextAt{21}{281.018}{529.822}{\Character{39}}%
\TextAt{3}{286.536}{529.822}{\Character{0}}%
\TextAt{17}{293.123}{529.822}{1}%
\TextAt{9}{301.176}{524.886}{\Character{61}}%
\TextAt{13}{313.601}{524.886}{\Character{25}}%
\TextAt{21}{320.67}{529.822}{\Character{27}}%
\TextAt{3}{325.796}{529.822}{\Character{0}}%
\TextAt{17}{332.382}{529.822}{1}%
\ErrataAt{79.19}{501.389}{455.81}{with $y\in U_{\widetilde L}$, then $r_{L\mid K}(\sigma^{-1})=N(y)\bmod N_{L\mid K}A_L$, where $N=N_{\widetilde L\mid\widetilde K}$ (B.~Dwork, see \hyperlink{nav.cite.N122}{[122]},}%
\TextAt{8}{79.226}{484.318}{chap.}%
\TextAt{8}{111.417}{484.318}{XIII,}%
\TextAt{8}{140.648}{484.318}{§5).}%
\TextAt{12}{79.226}{452.328}{Exercise}%
\TextAt{12}{132.359}{452.328}{IV.5.2.}%
\TextAt{8}{176.055}{452.328}{Generalize}%
\TextAt{8}{233.297}{452.328}{the}%
\TextAt{8}{253.394}{452.328}{theory}%
\TextAt{8}{290.085}{452.328}{developed}%
\TextAt{8}{344.326}{452.328}{so}%
\TextAt{8}{358.637}{452.328}{far}%
\TextAt{8}{376.463}{452.328}{in}%
\TextAt{8}{390.057}{452.328}{the}%
\TextAt{8}{410.154}{452.328}{following}%
\TextAt{8}{459.529}{452.328}{way.}%
\TextAt{8}{486.823}{452.328}{Let}%
\TextAt{13}{507.733}{452.328}{P}%
\TextAt{8}{520.743}{452.328}{be}%
\TextAt{8}{79.226}{435.282}{a}%
\TextAt{8}{89.607}{435.282}{set}%
\TextAt{8}{108.5}{435.282}{of}%
\TextAt{8}{122.455}{435.282}{prime}%
\TextAt{8}{156.245}{435.282}{numbers}%
\TextAt{8}{204.081}{435.282}{and}%
\TextAt{8}{227.469}{435.282}{let}%
\TextAt{13}{244.999}{435.282}{G}%
\TextAt{8}{258.756}{435.282}{be}%
\TextAt{8}{275.305}{435.282}{a}%
\TextAt{8}{285.686}{435.282}{pro-}%
\TextAt{13}{306.5}{435.282}{P}%
\TextAt{8}{315.671}{435.282}{-group,}%
\TextAt{8}{356.778}{435.282}{i.e.,}%
\TextAt{8}{379.665}{435.282}{a}%
\TextAt{8}{390.046}{435.282}{profinite}%
\TextAt{8}{437.499}{435.282}{group}%
\TextAt{8}{471.3}{435.282}{all}%
\TextAt{8}{488.185}{435.282}{of}%
\TextAt{8}{502.14}{435.282}{whose}%
\TextAt{8}{79.226}{418.236}{quotients}%
\TextAt{13}{130.024}{418.236}{G}%
\TextAt{9}{139.257}{418.236}{\Character{47}}%
\TextAt{13}{145.11}{418.236}{N}%
\TextAt{8}{159.63}{418.236}{by}%
\TextAt{8}{175.889}{418.236}{open}%
\TextAt{8}{204.175}{418.236}{normal}%
\TextAt{8}{243.854}{418.236}{subgroups}%
\TextAt{13}{299.267}{418.236}{N}%
\TextAt{8}{313.787}{418.236}{have}%
\TextAt{8}{340.781}{418.236}{order}%
\TextAt{8}{371.351}{418.236}{divisible}%
\TextAt{8}{417.259}{418.236}{only}%
\TextAt{8}{442.95}{418.236}{by}%
\TextAt{8}{459.209}{418.236}{primes}%
\TextAt{8}{496.988}{418.236}{in}%
\TextAt{13}{510.64}{418.236}{P}%
\TextAt{8}{519.812}{418.236}{.}%
\TextAt{8}{97.159}{398.202}{Let}%
\TextAt{13}{118.702}{398.202}{d}%
\TextAt{9}{129.083}{398.202}{\Character{58}}%
\TextAt{13}{136.632}{398.202}{G}%
\NativeMathAt{150.164}{398.202}{11.955}{\rightarrow}%
\TextAt{18}{166.417}{398.202}{Z}%
\TextAt{21}{174.387}{396.408}{P}%
\TextAt{8}{185.901}{398.202}{be}%
\TextAt{8}{202.399}{398.202}{a}%
\TextAt{8}{212.728}{398.202}{surjective}%
\TextAt{8}{265.713}{398.202}{homomorphism}%
\TextAt{8}{348.956}{398.202}{onto}%
\TextAt{8}{375.879}{398.202}{the}%
\TextAt{8}{396.609}{398.202}{group}%
\TextAt{18}{430.358}{398.202}{Z}%
\TextAt{21}{438.328}{396.408}{P}%
\TextAt{9}{449.669}{398.202}{\Character{61}}%
\TextAt{20}{463.072}{407.168}{\Character{81}}%
\TextAt{21}{474.363}{394.715}{p}%
\NativeMathAt{478.626}{394.715}{7.97}{\in}%
\TextAt{21}{484.272}{394.715}{P}%
\TextAt{18}{493.307}{398.202}{Z}%
\TextAt{21}{501.277}{396.408}{p}%
\TextAt{8}{506.038}{398.202}{,}%
\TextAt{8}{513.904}{398.202}{and}%
\TextAt{8}{79.226}{381.156}{let}%
\TextAt{13}{97.289}{381.156}{A}%
\TextAt{8}{111.119}{381.156}{be}%
\TextAt{8}{128.202}{381.156}{a}%
\TextAt{13}{139.115}{381.156}{G}%
\TextAt{8}{148.349}{381.156}{-module.}%
\TextAt{8}{201.564}{381.156}{A}%
\TextAt{12}{215.393}{381.156}{henselian}%
\TextAt{13}{275.222}{381.156}{P}%
\TextAt{12}{284.393}{381.156}{-valuation}%
\TextAt{8}{347.754}{381.156}{with}%
\TextAt{8}{375.572}{381.156}{respect}%
\TextAt{8}{416.779}{381.156}{to}%
\TextAt{13}{432.248}{381.156}{d}%
\TextAt{8}{443.386}{381.156}{is}%
\TextAt{8}{456.308}{381.156}{by}%
\TextAt{8}{473.725}{381.156}{definition}%
\TextAt{8}{526.911}{381.156}{a}%
\TextAt{8}{79.226}{364.11}{homomorphism}%
\TextAt{13}{270.601}{347.064}{v}%
\TextAt{9}{280.011}{347.064}{\Character{58}}%
\TextAt{13}{286.583}{347.064}{A}%
\TextAt{21}{295.358}{345.27}{k}%
\NativeArrowAt{303.799}{347.064}{11.955}{19.26}{\longrightarrow}%
\TextAt{18}{326.381}{347.064}{Z}%
\TextAt{21}{334.351}{345.27}{P}%
\TextAt{8}{79.226}{323.542}{which}%
\TextAt{8}{112.713}{323.542}{satisfies}%
\TextAt{8}{155.823}{323.542}{the}%
\TextAt{8}{175.979}{323.542}{following}%
\TextAt{8}{225.414}{323.542}{properties:}%
\ErrataAt{89.5}{303.533}{445.5}{(i) $v(A_k)=Z\supseteq\mathbb Z$ and $Z/nZ\cong\mathbb Z/n\mathbb Z$ for all natural numbers $n$ which are divisible}%
\TextAt{8}{107.918}{286.461}{only}%
\TextAt{8}{133.61}{286.461}{by}%
\TextAt{8}{149.869}{286.461}{primes}%
\TextAt{8}{187.648}{286.461}{in}%
\TextAt{13}{201.3}{286.461}{P}%
\TextAt{8}{210.471}{286.461}{,}%
\TextAt{8}{86.327}{269.415}{(ii)}%
\TextAt{13}{107.918}{269.415}{v}%
\TextAt{9}{114.007}{269.415}{\Character{40}}%
\TextAt{13}{118.559}{269.415}{N}%
\TextAt{21}{127.933}{267.556}{K}%
\TextAt{3}{135.637}{267.555}{\Character{106}}%
\TextAt{21}{137.99}{267.556}{k}%
\TextAt{13}{143.109}{269.415}{A}%
\TextAt{21}{151.885}{267.622}{K}%
\TextAt{9}{160.087}{269.415}{\Character{41}}%
\TextAt{9}{167.942}{269.415}{\Character{61}}%
\TextAt{13}{180.386}{269.415}{f}%
\TextAt{21}{186.156}{267.622}{K}%
\TextAt{13}{194.359}{269.415}{Z}%
\TextAt{8}{207.127}{269.415}{for}%
\TextAt{8}{225.012}{269.415}{all}%
\TextAt{8}{241.271}{269.415}{finite}%
\TextAt{8}{271.183}{269.415}{extensions}%
\TextAt{13}{327.563}{269.415}{K}%
\TextAt{16}{341.702}{269.415}{\Character{106}}%
\TextAt{13}{348.344}{269.415}{k}%
\TextAt{8}{354.833}{269.415}{,}%
\TextAt{8}{361.982}{269.415}{where}%
\TextAt{13}{395.791}{269.415}{f}%
\TextAt{21}{401.561}{267.622}{K}%
\TextAt{9}{413.085}{269.415}{\Character{61}}%
\TextAt{9}{425.495}{269.415}{\Character{40}}%
\TextAt{13}{430.063}{269.415}{d}%
\TextAt{9}{436.145}{269.415}{\Character{40}}%
\TextAt{13}{440.698}{269.415}{G}%
\TextAt{9}{449.931}{269.415}{\Character{41}}%
\TextAt{9}{457.786}{269.415}{\Character{58}}%
\TextAt{13}{464.377}{269.415}{d}%
\TextAt{9}{470.46}{269.415}{\Character{40}}%
\TextAt{13}{475.012}{269.415}{G}%
\TextAt{21}{484.246}{267.622}{K}%
\TextAt{9}{492.448}{269.415}{\Character{41}\Character{41}}%
\TextAt{8}{501.553}{269.415}{.}%
\TextAt{8}{79.226}{249.38}{Under}%
\TextAt{8}{114.882}{249.38}{the}%
\TextAt{8}{135.259}{249.38}{hypothesis}%
\TextAt{8}{193.164}{249.38}{that}%
\TextAt{13}{218.754}{249.38}{H}%
\TextAt{21}{229.376}{253.719}{i}%
\TextAt{9}{232.758}{249.38}{\Character{40}}%
\TextAt{13}{237.31}{249.38}{G}%
\TextAt{9}{246.544}{249.38}{\Character{40}}%
\TextAt{13}{251.096}{249.38}{L}%
\TextAt{16}{262.758}{249.38}{\Character{106}}%
\TextAt{13}{269.776}{249.38}{K}%
\TextAt{9}{280.594}{249.38}{\Character{41}}%
\TextAt{13}{285.146}{249.38}{\Character{59}}%
\TextAt{13}{290.37}{249.38}{U}%
\TextAt{21}{298.34}{247.587}{L}%
\TextAt{9}{304.597}{249.38}{\Character{41}}%
\TextAt{9}{312.834}{249.38}{\Character{61}}%
\TextAt{9}{325.626}{249.38}{1}%
\TextAt{8}{335.62}{249.38}{for}%
\TextAt{13}{353.726}{249.38}{i}%
\TextAt{9}{361.417}{249.38}{\Character{61}}%
\TextAt{9}{374.209}{249.38}{0}%
\TextAt{13}{380.072}{249.38}{\Character{59}}%
\TextAt{16}{385.316}{249.38}{\Character{0}}%
\TextAt{9}{394.614}{249.38}{1}%
\TextAt{8}{400.467}{249.38}{,}%
\TextAt{8}{407.893}{249.38}{for}%
\TextAt{8}{425.998}{249.38}{all}%
\TextAt{8}{442.478}{249.38}{unramified}%
\TextAt{8}{501.232}{249.38}{exten-}%
\TextAt{8}{79.226}{232.334}{sions}%
\TextAt{13}{108.615}{232.334}{L}%
\TextAt{16}{121.006}{232.334}{\Character{106}}%
\TextAt{13}{128.753}{232.334}{K}%
\TextAt{8}{139.57}{232.334}{,}%
\TextAt{8}{147.53}{232.334}{prove}%
\TextAt{8}{179.728}{232.334}{the}%
\TextAt{8}{200.534}{232.334}{existence}%
\TextAt{8}{250.987}{232.334}{of}%
\TextAt{8}{264.966}{232.334}{a}%
\TextAt{8}{275.37}{232.334}{canonical}%
\TextAt{8}{327.402}{232.334}{reciprocity}%
\TextAt{8}{386.26}{232.334}{homomorphism}%
\TextAt{13}{469.579}{232.334}{r}%
\TextAt{21}{474.854}{230.475}{L}%
\TextAt{3}{480.613}{230.475}{\Character{106}}%
\TextAt{21}{482.965}{230.475}{K}%
\TextAt{9}{495.594}{232.334}{\Character{58}}%
\TextAt{13}{503.272}{232.334}{G}%
\TextAt{9}{512.506}{232.334}{\Character{40}}%
\TextAt{13}{517.058}{232.334}{L}%
\TextAt{16}{529.449}{232.334}{\Character{106}}%
\TextAt{13}{79.226}{215.288}{K}%
\TextAt{9}{90.044}{215.289}{\Character{41}}%
\TextAt{7}{94.596}{219.627}{ab}%
\NativeMathAt{107.35}{215.289}{11.955}{\rightarrow}%
\TextAt{13}{122.626}{215.288}{A}%
\TextAt{21}{131.401}{213.495}{K}%
\TextAt{9}{139.604}{215.289}{\Character{47}}%
\TextAt{13}{145.457}{215.289}{N}%
\TextAt{21}{154.83}{213.429}{L}%
\TextAt{3}{160.589}{213.429}{\Character{106}}%
\TextAt{21}{162.942}{213.429}{K}%
\TextAt{13}{171.144}{215.289}{A}%
\TextAt{21}{179.92}{213.495}{L}%
\TextAt{8}{190.074}{215.289}{for}%
\TextAt{8}{207.959}{215.289}{every}%
\TextAt{8}{238.852}{215.289}{finite}%
\TextAt{8}{268.763}{215.289}{Galois}%
\TextAt{8}{304.677}{215.289}{extension}%
\TextAt{13}{356.443}{215.289}{L}%
\TextAt{16}{367.728}{215.288}{\Character{106}}%
\TextAt{13}{374.37}{215.289}{K}%
\TextAt{8}{385.187}{215.289}{.}%
\TextAt{12}{79.226}{183.299}{Exercise}%
\TextAt{12}{133.026}{183.299}{IV.5.3.}%
\TextAt{8}{177.303}{183.299}{Let}%
\TextAt{13}{198.793}{183.299}{d}%
\TextAt{9}{209.084}{183.299}{\Character{58}}%
\TextAt{13}{216.544}{183.299}{G}%
\NativeMathAt{229.986}{183.299}{11.955}{\rightarrow}%
\TextAt{20}{246.814}{186.387}{\Character{98}}%
\TextAt{18}{246.15}{183.299}{Z}%
\TextAt{8}{258.539}{183.299}{be}%
\TextAt{8}{274.984}{183.299}{a}%
\TextAt{8}{285.26}{183.299}{surjective}%
\TextAt{8}{338.192}{183.299}{homomorphism,}%
\TextAt{13}{424.765}{183.299}{A}%
\TextAt{8}{437.959}{183.299}{a}%
\TextAt{13}{448.235}{183.299}{G}%
\TextAt{8}{457.469}{183.299}{-module}%
\TextAt{8}{503.18}{183.299}{which}%
\TextAt{8}{79.226}{166.253}{satisfies}%
\TextAt{8}{122.52}{166.253}{axiom}%
\TextAt{8}{157.505}{166.253}{(5.1),}%
\TextAt{8}{188.962}{166.253}{and}%
\TextAt{8}{211.908}{166.253}{let}%
\TextAt{13}{228.996}{166.253}{v}%
\TextAt{9}{238.718}{166.253}{\Character{58}}%
\TextAt{13}{245.603}{166.253}{A}%
\TextAt{21}{254.379}{164.459}{k}%
\NativeMathAt{263.132}{166.253}{11.955}{\rightarrow}%
\TextAt{20}{279.385}{169.341}{\Character{98}}%
\TextAt{18}{278.721}{166.253}{Z}%
\TextAt{8}{290.772}{166.253}{be}%
\TextAt{8}{306.88}{166.253}{a}%
\TextAt{8}{316.819}{166.253}{henselian}%
\TextAt{8}{367.789}{166.253}{valuation}%
\TextAt{8}{419.045}{166.253}{with}%
\TextAt{8}{445.888}{166.253}{respect}%
\TextAt{8}{486.122}{166.253}{to}%
\TextAt{13}{500.616}{166.253}{d}%
\TextAt{8}{506.698}{166.253}{.}%
\TextAt{8}{515.698}{166.253}{Let}%
\TextAt{13}{79.226}{149.207}{K}%
\TextAt{16}{93.365}{149.207}{\Character{106}}%
\TextAt{13}{100.006}{149.207}{k}%
\TextAt{8}{109.58}{149.207}{be}%
\TextAt{8}{124.69}{149.207}{a}%
\TextAt{8}{133.632}{149.207}{finite}%
\TextAt{8}{162.731}{149.207}{extension}%
\TextAt{8}{213.683}{149.207}{and}%
\TextAt{8}{235.632}{149.207}{let}%
\TextAt{8}{251.723}{149.207}{spec}%
\TextAt{9}{273.242}{149.207}{\Character{40}}%
\TextAt{13}{277.795}{149.207}{K}%
\TextAt{9}{288.612}{149.207}{\Character{41}}%
\TextAt{8}{296.248}{149.207}{be}%
\TextAt{8}{311.359}{149.207}{the}%
\TextAt{8}{330.702}{149.207}{set}%
\TextAt{8}{348.156}{149.207}{of}%
\TextAt{8}{360.672}{149.207}{microprimes}%
\TextAt{8}{426.258}{149.207}{of}%
\TextAt{13}{438.774}{149.207}{K}%
\TextAt{8}{452.676}{149.207}{(see}%
\TextAt{8}{475.33}{149.207}{§4,}%
\TextAt{8}{493.365}{149.207}{exercise}%
\ErrataAt{79.19}{132.205}{455.81}{\hyperlink{nav.ex.4.4.1}{1}--\hyperlink{nav.ex.4.4.4}{4}). Define a canonical mapping}%
\TextAt{13}{228.376}{103.16}{r}%
\TextAt{21}{233.651}{101.366}{K}%
\TextAt{9}{245.175}{103.16}{\Character{58}}%
\TextAt{8}{251.747}{103.16}{spec}%
\TextAt{9}{273.266}{103.16}{\Character{40}}%
\TextAt{13}{277.819}{103.16}{K}%
\TextAt{9}{288.636}{103.16}{\Character{41}}%
\NativeArrowAt{296.51}{103.16}{11.955}{19.26}{\longrightarrow}%
\TextAt{13}{319.092}{103.16}{A}%
\TextAt{21}{327.867}{101.366}{K}%
\TextAt{9}{336.07}{103.16}{\Character{47}}%
\TextAt{13}{341.923}{103.16}{N}%
\TextAt{17}{351.49}{102.706}{\Character{22}}%
\TextAt{21}{351.296}{100.602}{k}%
\TextAt{3}{355.918}{100.602}{\Character{106}}%
\TextAt{21}{358.27}{100.602}{K}%
\TextAt{13}{366.473}{103.16}{A}%
\TextAt{17}{375.442}{102.706}{\Character{22}}%
\TextAt{21}{375.248}{100.602}{k}%
\TextAt{13}{380.368}{103.16}{\Character{59}}%
\TextAt{8}{297.211}{43.033}{273}%
\Destination{79.23}{318.105}{nav.Item.92}%
\Destination{79.23}{284.025}{nav.Item.93}%
\Destination{78.23}{809.815}{nav.page.273}%
\Destination{79.226}{572.955}{nav.ex.4.5.1}%
\Destination{79.226}{466.355}{nav.ex.4.5.2}%
\Destination{79.226}{197.255}{nav.ex.4.5.3}%
\Destination{79.36}{181.255}{nav.print.299}%
\Hotlink{145.827}{482.155}{6.858}{11.047}{nav.section.4.5}%
\Hotlink{480.509}{147.055}{6.858}{11.047}{nav.section.4.4}%
\Destination{79.226}{320.533}{errata.105}%
\Destination{79.226}{518.389}{errata.117}%
\Destination{79.226}{165}{errata.199}%
}

% Page 286 (body page 274)
\Page{%
\TextAt{2}{79.226}{796.943}{IV.}%
\TextAt{2}{98.012}{796.943}{Abstract}%
\TextAt{2}{143.252}{796.943}{Class}%
\TextAt{2}{171.844}{796.943}{Field}%
\TextAt{2}{199.575}{796.943}{Theory}%
\Vector{q 0 G 0 g 0.299 w 7.226 18.729 m 460.769 18.729 l S Q}%
\TextAt{8}{79.226}{761.124}{and}%
\TextAt{8}{101.989}{761.124}{show}%
\TextAt{8}{130.992}{761.124}{that,}%
\TextAt{8}{159.613}{761.124}{for}%
\TextAt{8}{177.498}{761.124}{a}%
\TextAt{8}{187.253}{761.124}{finite}%
\TextAt{8}{217.165}{761.124}{extension,}%
\TextAt{8}{272.183}{761.124}{the}%
\TextAt{8}{292.339}{761.124}{diagram}%
\TextAt{8}{236.565}{733.493}{spec}%
\TextAt{9}{258.085}{733.493}{\Character{40}}%
\TextAt{13}{262.637}{733.493}{L}%
\TextAt{9}{270.601}{733.493}{\Character{41}}%
\TextAt{13}{317.523}{733.493}{A}%
\TextAt{21}{326.298}{731.7}{L}%
\TextAt{9}{332.556}{733.493}{\Character{47}}%
\TextAt{13}{338.409}{733.493}{N}%
\TextAt{17}{347.976}{733.039}{\Character{22}}%
\TextAt{21}{347.782}{730.936}{k}%
\TextAt{3}{352.404}{730.936}{\Character{106}}%
\TextAt{21}{354.756}{730.936}{L}%
\TextAt{13}{361.013}{733.493}{A}%
\TextAt{17}{369.983}{733.039}{\Character{22}}%
\TextAt{21}{369.789}{730.936}{k}%
\TextAt{8}{235.138}{689.698}{spec}%
\TextAt{9}{256.658}{689.698}{\Character{40}}%
\TextAt{13}{261.21}{689.698}{K}%
\TextAt{9}{272.028}{689.698}{\Character{41}}%
\TextAt{13}{315.577}{689.698}{A}%
\TextAt{21}{324.352}{687.905}{K}%
\TextAt{9}{332.555}{689.698}{\Character{47}}%
\TextAt{13}{338.408}{689.698}{N}%
\TextAt{17}{347.976}{689.245}{\Character{22}}%
\TextAt{21}{347.782}{687.141}{k}%
\TextAt{3}{352.403}{687.141}{\Character{106}}%
\TextAt{21}{354.756}{687.141}{K}%
\TextAt{13}{362.958}{689.698}{A}%
\TextAt{17}{371.928}{689.245}{\Character{22}}%
\TextAt{21}{371.734}{687.141}{k}%
\NativeDiagram{%
\path[draw=black,line width=0.478bp] (208.507,-37.518) -- (239.695,-37.518);%
\path[draw=black,line width=0.478bp,line cap=round,line join=round] (237.447,-34.649) .. controls (237.902,-36.37) and (238.914,-37.183) .. (239.934,-37.518) .. controls (238.914,-37.853) and (237.902,-38.665) .. (237.447,-40.387);%
}%
\TextAt{21}{291.651}{740.637}{r}%
\TextAt{22}{295.472}{739.298}{L}%
\NativeDiagram{%
\path[draw=black,line width=0.478bp] (183.86,-48.075) -- (183.86,-70.278);%
\path[draw=black,line width=0.478bp,line cap=round,line join=round] (186.729,-68.031) .. controls (185.008,-68.485) and (184.195,-69.498) .. (183.86,-70.517) .. controls (183.526,-69.498) and (182.713,-68.485) .. (180.991,-68.031);%
}%
\TextAt{21}{247.896}{712.869}{\Character{25}}%
\NativeDiagram{%
\path[draw=black,line width=0.478bp] (274.218,-49.635) -- (274.218,-70.278);%
\path[draw=black,line width=0.478bp,line cap=round,line join=round] (277.087,-68.031) .. controls (275.366,-68.485) and (274.553,-69.498) .. (274.218,-70.517) .. controls (273.883,-69.498) and (273.071,-68.485) .. (271.349,-68.031);%
}%
\TextAt{21}{349.034}{712.697}{N}%
\TextAt{22}{355.755}{710.959}{L}%
\TextAt{23}{360.814}{710.959}{\Character{106}}%
\TextAt{22}{363.083}{710.959}{K}%
\NativeDiagram{%
\path[draw=black,line width=0.478bp] (209.933,-81.313) -- (237.749,-81.313);%
\path[draw=black,line width=0.478bp,line cap=round,line join=round] (235.502,-78.444) .. controls (235.956,-80.166) and (236.968,-80.979) .. (237.988,-81.313) .. controls (236.968,-81.648) and (235.956,-82.461) .. (235.502,-84.182);%
}%
\TextAt{21}{290.648}{686.44}{r}%
\TextAt{22}{294.469}{685.101}{K}%
\TextAt{8}{79.226}{659.97}{commutes.}%
\TextAt{8}{140.983}{659.97}{Show}%
\TextAt{8}{172.605}{659.97}{furthermore}%
\TextAt{8}{238.047}{659.97}{that,}%
\TextAt{8}{267.58}{659.97}{for}%
\TextAt{8}{286.194}{659.97}{every}%
\TextAt{8}{317.816}{659.97}{finite}%
\TextAt{8}{348.457}{659.97}{Galois}%
\TextAt{8}{385.099}{659.97}{extension}%
\TextAt{13}{437.595}{659.97}{L}%
\TextAt{16}{450.123}{659.97}{\Character{106}}%
\TextAt{13}{458.008}{659.97}{K}%
\TextAt{8}{468.825}{659.97}{,}%
\TextAt{13}{476.886}{659.97}{r}%
\TextAt{21}{482.162}{658.176}{K}%
\TextAt{8}{494.991}{659.97}{induces}%
\ErrataAt{79.19}{642.938}{455.81}{the reciprocity homomorphism}%
\TextAt{13}{220.362}{615.041}{r}%
\TextAt{21}{225.637}{613.181}{L}%
\TextAt{3}{231.396}{613.181}{\Character{106}}%
\TextAt{21}{233.749}{613.181}{K}%
\TextAt{9}{245.272}{615.041}{\Character{58}}%
\TextAt{13}{251.845}{615.041}{G}%
\TextAt{9}{261.078}{615.041}{\Character{40}}%
\TextAt{13}{265.631}{615.041}{L}%
\TextAt{16}{276.916}{615.041}{\Character{106}}%
\TextAt{13}{283.558}{615.041}{K}%
\TextAt{9}{294.375}{615.041}{\Character{41}}%
\NativeArrowAt{302.249}{615.041}{11.955}{19.26}{\longrightarrow}%
\TextAt{13}{324.831}{615.041}{A}%
\TextAt{21}{333.606}{613.247}{K}%
\TextAt{9}{341.809}{615.041}{\Character{47}}%
\TextAt{13}{347.662}{615.041}{N}%
\TextAt{21}{357.035}{613.181}{L}%
\TextAt{3}{362.794}{613.181}{\Character{106}}%
\TextAt{21}{365.147}{613.181}{K}%
\TextAt{13}{373.349}{615.041}{A}%
\TextAt{21}{382.125}{613.247}{L}%
\TextAt{13}{388.382}{615.041}{\Character{58}}%
\TextAt{12}{97.159}{584.169}{Hint:}%
\TextAt{8}{132.814}{584.169}{Let}%
\TextAt{13}{153.842}{584.169}{\Character{39}}%
\NativeMathAt{164.93}{584.169}{11.955}{\in}%
\TextAt{13}{176.32}{584.169}{G}%
\TextAt{21}{185.554}{582.376}{K}%
\TextAt{8}{197.712}{584.169}{be}%
\TextAt{8}{213.694}{584.169}{an}%
\TextAt{8}{230.011}{584.169}{element}%
\TextAt{8}{273.311}{584.169}{such}%
\TextAt{8}{299.767}{584.169}{that}%
\TextAt{13}{325.194}{584.169}{d}%
\TextAt{21}{331.276}{582.376}{K}%
\TextAt{9}{339.479}{584.169}{\Character{40}}%
\TextAt{13}{344.031}{584.169}{\Character{39}}%
\TextAt{9}{351.7}{584.169}{\Character{41}}%
\NativeMathAt{359.672}{584.169}{11.955}{\in}%
\TextAt{18}{371.062}{584.169}{N}%
\TextAt{8}{379.696}{584.169}{.}%
\TextAt{8}{388.319}{584.169}{Let}%
\TextAt{9}{409.347}{584.169}{\Character{6}}%
\TextAt{8}{421.757}{584.169}{be}%
\TextAt{8}{437.739}{584.169}{the}%
\TextAt{8}{457.954}{584.169}{fixed}%
\TextAt{8}{486.298}{584.169}{field}%
\TextAt{8}{511.713}{584.169}{of}%
\TextAt{13}{525.101}{584.169}{\Character{39}}%
\TextAt{8}{79.226}{567.123}{and}%
\TextAt{20}{271.51}{553.099}{\Character{98}}%
\TextAt{13}{268.817}{550.077}{A}%
\TextAt{17}{277.593}{548.284}{\Character{6}}%
\TextAt{9}{287.528}{550.077}{\Character{61}}%
\NativeMathAt{299.953}{550.077}{11.955}{\varprojlim}%
\TextAt{21}{305.363}{537.963}{\Character{11}}%
\TextAt{13}{318.205}{550.077}{A}%
\TextAt{21}{326.98}{548.284}{K}%
\TextAt{22}{334.121}{547.287}{\Character{11}}%
\TextAt{13}{339.927}{550.077}{\Character{59}}%
\TextAt{8}{79.226}{520.808}{where}%
\TextAt{13}{112.436}{520.808}{K}%
\TextAt{21}{122.395}{519.015}{\Character{11}}%
\TextAt{16}{131.653}{520.808}{\Character{106}}%
\TextAt{13}{138.294}{520.808}{K}%
\TextAt{8}{152.41}{520.808}{varies}%
\TextAt{8}{184.722}{520.808}{over}%
\TextAt{8}{209.169}{520.808}{the}%
\TextAt{8}{228.725}{520.808}{finite}%
\TextAt{8}{258.037}{520.808}{subextensions}%
\TextAt{8}{331.763}{520.808}{of}%
\TextAt{9}{344.493}{520.808}{\Character{6}}%
\TextAt{16}{356.268}{520.808}{\Character{106}}%
\TextAt{13}{362.91}{520.808}{K}%
\TextAt{8}{373.728}{520.808}{,}%
\TextAt{8}{380.397}{520.808}{and}%
\TextAt{8}{402.56}{520.808}{where}%
\TextAt{8}{435.769}{520.808}{the}%
\TextAt{8}{455.326}{520.808}{projective}%
\TextAt{8}{508.704}{520.808}{limit}%
\TextAt{8}{79.226}{503.762}{is}%
\TextAt{8}{90.998}{503.762}{taken}%
\TextAt{8}{122.879}{503.762}{with}%
\TextAt{8}{149.548}{503.762}{respect}%
\TextAt{8}{189.606}{503.762}{to}%
\TextAt{8}{203.925}{503.762}{the}%
\TextAt{8}{224.09}{503.762}{norm}%
\TextAt{8}{254.668}{503.762}{maps}%
\TextAt{13}{285.305}{503.762}{N}%
\TextAt{21}{294.679}{501.902}{K}%
\TextAt{22}{301.82}{500.496}{\Character{12}}%
\TextAt{3}{306.798}{501.902}{\Character{106}}%
\TextAt{21}{309.15}{501.902}{K}%
\TextAt{22}{316.292}{500.906}{\Character{11}}%
\TextAt{9}{325.432}{503.762}{\Character{58}}%
\TextAt{13}{332.019}{503.762}{A}%
\TextAt{21}{340.795}{501.969}{K}%
\TextAt{22}{347.936}{500.563}{\Character{12}}%
\NativeMathAt{356.747}{503.762}{11.955}{\rightarrow}%
\TextAt{13}{372.037}{503.762}{A}%
\TextAt{21}{380.813}{501.969}{K}%
\TextAt{22}{387.954}{500.972}{\Character{11}}%
\TextAt{8}{393.76}{503.762}{.}%
\TextAt{8}{402.233}{503.762}{Then}%
\TextAt{8}{432.799}{503.762}{there}%
\TextAt{8}{462.719}{503.762}{is}%
\TextAt{8}{474.492}{503.762}{a}%
\TextAt{8}{484.255}{503.762}{surjective}%
\TextAt{8}{79.226}{486.716}{homomorphism}%
\TextAt{13}{161.896}{486.716}{v}%
\TextAt{17}{167.555}{484.923}{\Character{6}}%
\TextAt{9}{177.49}{486.716}{\Character{58}}%
\TextAt{20}{186.756}{489.738}{\Character{98}}%
\TextAt{13}{184.063}{486.716}{A}%
\TextAt{17}{192.838}{484.923}{\Character{6}}%
\NativeMathAt{202.773}{486.716}{11.955}{\rightarrow}%
\TextAt{13}{218.049}{486.716}{Z}%
\TextAt{8}{230.817}{486.716}{and}%
\TextAt{8}{253.58}{486.716}{a}%
\TextAt{8}{263.336}{486.716}{homomorphism}%
\TextAt{13}{346.006}{486.716}{N}%
\TextAt{17}{355.379}{484.856}{\Character{6}}%
\TextAt{3}{361.495}{484.856}{\Character{106}}%
\TextAt{21}{363.848}{484.856}{K}%
\TextAt{9}{375.371}{486.716}{\Character{58}}%
\TextAt{20}{384.636}{489.738}{\Character{98}}%
\TextAt{13}{381.944}{486.716}{A}%
\TextAt{17}{390.719}{484.923}{\Character{6}}%
\NativeMathAt{400.654}{486.716}{11.955}{\rightarrow}%
\TextAt{13}{415.93}{486.716}{A}%
\TextAt{21}{424.705}{484.923}{K}%
\TextAt{8}{432.908}{486.716}{.}%
\TextAt{15}{79.226}{434.337}{§}%
\TextAt{15}{93.785}{434.337}{6}%
\TextAt{15}{120.693}{434.337}{The}%
\TextAt{15}{160.202}{434.337}{General}%
\TextAt{15}{233.024}{434.337}{Reciprocity}%
\TextAt{15}{336.884}{434.337}{Law}%
\TextAt{8}{79.226}{402.451}{We}%
\TextAt{8}{99.371}{402.451}{now}%
\TextAt{8}{123.759}{402.451}{impose}%
\TextAt{8}{163.163}{402.451}{on}%
\TextAt{8}{179.422}{402.451}{the}%
\TextAt{8}{199.579}{402.451}{continuous}%
\TextAt{13}{258.183}{402.451}{G}%
\TextAt{8}{267.417}{402.451}{-module}%
\TextAt{13}{312.607}{402.451}{A}%
\TextAt{8}{325.28}{402.451}{the}%
\TextAt{8}{345.436}{402.451}{following}%
\TextAt{8}{394.871}{402.451}{condition.}%
\TextAt{12}{79.226}{376.758}{Class}%
\TextAt{12}{114.327}{376.758}{Field}%
\TextAt{12}{148.351}{376.758}{Axiom}%
\TextAt{12}{191.784}{376.758}{(6.1).}%
\TextAt{8}{227.1}{376.758}{For}%
\TextAt{8}{248.069}{376.758}{every}%
\TextAt{8}{278.961}{376.758}{cyclic}%
\TextAt{8}{311.144}{376.758}{extension}%
\TextAt{13}{362.91}{376.758}{L}%
\TextAt{16}{374.196}{376.758}{\Character{106}}%
\TextAt{13}{380.837}{376.758}{K}%
\TextAt{8}{395.552}{376.758}{one}%
\TextAt{8}{417.012}{376.758}{has}%
\TextAt{8}{190.521}{337.416}{\char35\relax }%
\TextAt{13}{200.276}{337.416}{H}%
\TextAt{21}{210.899}{342.352}{i}%
\TextAt{9}{214.28}{337.416}{\Character{40}}%
\TextAt{13}{218.832}{337.416}{G}%
\TextAt{9}{228.066}{337.416}{\Character{40}}%
\TextAt{13}{232.618}{337.416}{L}%
\TextAt{16}{243.904}{337.416}{\Character{106}}%
\TextAt{13}{250.545}{337.416}{K}%
\TextAt{9}{261.363}{337.416}{\Character{41}}%
\TextAt{13}{265.915}{337.416}{\Character{59}}%
\TextAt{13}{271.139}{337.416}{A}%
\TextAt{21}{279.935}{335.623}{L}%
\TextAt{9}{286.192}{337.416}{\Character{41}}%
\TextAt{9}{294.047}{337.416}{\Character{61}}%
\TextAt{20}{306.491}{361.924}{\Character{56}}%
\TextAt{20}{306.491}{351.165}{\Character{60}}%
\TextAt{20}{306.491}{329.645}{\Character{58}}%
\TextAt{9}{317.118}{346.541}{\Character{91}}%
\TextAt{13}{320.369}{346.541}{L}%
\TextAt{9}{331.655}{346.541}{\Character{58}}%
\TextAt{13}{338.227}{346.541}{K}%
\TextAt{9}{349.045}{346.541}{\Character{93}}%
\TextAt{8}{364.252}{346.541}{for}%
\TextAt{13}{382.136}{346.541}{i}%
\TextAt{9}{389.451}{346.541}{\Character{61}}%
\TextAt{9}{401.861}{346.541}{0}%
\TextAt{13}{407.729}{346.541}{\Character{59}}%
\TextAt{9}{317.118}{326.086}{1}%
\TextAt{8}{364.252}{326.086}{for}%
\TextAt{13}{382.136}{326.086}{i}%
\TextAt{9}{389.451}{326.086}{\Character{61}}%
\TextAt{16}{401.876}{326.086}{\Character{0}}%
\TextAt{9}{411.175}{326.086}{1}%
\TextAt{13}{417.028}{326.086}{\Character{58}}%
\TextAt{8}{97.159}{289.869}{Among}%
\TextAt{8}{137.216}{289.869}{the}%
\TextAt{8}{156.783}{289.869}{cyclic}%
\TextAt{8}{188.376}{289.869}{extensions}%
\TextAt{8}{244.167}{289.869}{there}%
\TextAt{8}{273.489}{289.869}{are}%
\TextAt{8}{292.409}{289.869}{in}%
\TextAt{8}{305.472}{289.869}{particular}%
\TextAt{8}{358.872}{289.869}{the}%
\TextAt{8}{378.438}{289.869}{unramified}%
\TextAt{8}{436.381}{289.869}{ones.}%
\TextAt{8}{466.809}{289.869}{For}%
\TextAt{8}{487.189}{289.869}{them}%
\TextAt{8}{516.51}{289.869}{the}%
\TextAt{8}{79.226}{272.823}{above}%
\TextAt{8}{112.402}{272.823}{condition}%
\TextAt{8}{163.785}{272.823}{amounts}%
\TextAt{8}{211.008}{272.823}{precisely}%
\TextAt{8}{258.865}{272.823}{to}%
\TextAt{8}{273.175}{272.823}{requiring}%
\TextAt{8}{322.932}{272.823}{axiom}%
\TextAt{8}{357.734}{272.823}{(5.1),}%
\TextAt{8}{388.961}{272.823}{so}%
\TextAt{8}{403.331}{272.823}{that}%
\TextAt{8}{428.7}{272.823}{one}%
\TextAt{8}{450.159}{272.823}{has}%
\TextAt{12}{79.226}{247.13}{Proposition}%
\TextAt{12}{151.937}{247.13}{(6.2).}%
\TextAt{8}{187.253}{247.13}{For}%
\TextAt{8}{208.222}{247.13}{a}%
\TextAt{8}{217.978}{247.13}{finite}%
\TextAt{8}{247.89}{247.13}{unramified}%
\TextAt{8}{306.422}{247.13}{extension}%
\TextAt{13}{358.188}{247.13}{L}%
\TextAt{16}{369.473}{247.13}{\Character{106}}%
\TextAt{13}{376.115}{247.13}{K}%
\TextAt{8}{386.933}{247.13}{,}%
\TextAt{8}{394.082}{247.13}{one}%
\TextAt{8}{415.542}{247.13}{has}%
\TextAt{13}{204.859}{219.248}{H}%
\TextAt{21}{215.482}{224.184}{i}%
\TextAt{9}{218.863}{219.248}{\Character{40}}%
\TextAt{13}{223.415}{219.248}{G}%
\TextAt{9}{232.649}{219.248}{\Character{40}}%
\TextAt{13}{237.201}{219.248}{L}%
\TextAt{16}{248.486}{219.248}{\Character{106}}%
\TextAt{13}{255.128}{219.248}{K}%
\TextAt{9}{265.946}{219.248}{\Character{41}}%
\TextAt{13}{270.498}{219.248}{\Character{59}}%
\TextAt{13}{275.722}{219.248}{U}%
\TextAt{21}{283.693}{217.454}{L}%
\TextAt{9}{289.95}{219.248}{\Character{41}}%
\TextAt{9}{297.805}{219.248}{\Character{61}}%
\TextAt{9}{310.226}{219.248}{1}%
\TextAt{8}{340.012}{219.248}{for}%
\TextAt{13}{357.897}{219.248}{i}%
\TextAt{9}{365.211}{219.248}{\Character{61}}%
\TextAt{9}{377.62}{219.248}{0}%
\TextAt{13}{383.489}{219.248}{\Character{59}}%
\TextAt{16}{388.734}{219.248}{\Character{0}}%
\TextAt{9}{398.032}{219.248}{1}%
\TextAt{13}{403.885}{219.248}{\Character{58}}%
\TextAt{19}{79.226}{174.071}{Proof.}%
\TextAt{8}{116}{174.071}{Since}%
\TextAt{13}{147.098}{174.071}{L}%
\TextAt{16}{159.303}{174.071}{\Character{106}}%
\TextAt{13}{166.866}{174.071}{K}%
\TextAt{8}{182.121}{174.071}{is}%
\TextAt{8}{194.425}{174.071}{unramified,}%
\TextAt{8}{256.884}{174.071}{a}%
\TextAt{8}{267.18}{174.071}{prime}%
\TextAt{8}{300.884}{174.071}{element}%
\TextAt{13}{344.665}{174.071}{\Character{25}}%
\TextAt{21}{351.306}{172.278}{K}%
\TextAt{8}{363.946}{174.071}{of}%
\TextAt{13}{377.816}{174.071}{A}%
\TextAt{21}{386.591}{172.278}{K}%
\TextAt{8}{399.232}{174.071}{is}%
\TextAt{8}{411.536}{174.071}{also}%
\TextAt{8}{435.556}{174.071}{a}%
\TextAt{8}{445.851}{174.071}{prime}%
\TextAt{8}{479.555}{174.071}{element}%
\TextAt{8}{523.337}{174.071}{of}%
\TextAt{13}{79.226}{157.025}{A}%
\TextAt{21}{88.002}{155.232}{L}%
\TextAt{8}{94.259}{157.025}{.}%
\TextAt{8}{104.05}{157.025}{As}%
\TextAt{13}{121.785}{157.025}{H}%
\TextAt{3}{132.408}{161.363}{\Character{0}}%
\TextAt{17}{138.994}{161.363}{1}%
\TextAt{9}{143.727}{157.025}{\Character{40}}%
\TextAt{13}{148.279}{157.025}{G}%
\TextAt{9}{157.513}{157.025}{\Character{40}}%
\TextAt{13}{162.065}{157.025}{L}%
\TextAt{16}{174.113}{157.025}{\Character{106}}%
\TextAt{13}{181.518}{157.025}{K}%
\TextAt{9}{192.336}{157.025}{\Character{41}}%
\TextAt{13}{196.888}{157.025}{\Character{59}}%
\TextAt{13}{202.112}{157.025}{A}%
\TextAt{21}{210.908}{155.232}{L}%
\TextAt{9}{217.165}{157.025}{\Character{41}}%
\TextAt{9}{225.785}{157.025}{\Character{61}}%
\TextAt{9}{238.971}{157.025}{1}%
\TextAt{8}{244.843}{157.025}{,}%
\TextAt{8}{252.552}{157.025}{every}%
\TextAt{8}{283.891}{157.025}{element}%
\TextAt{13}{327.581}{157.025}{u}%
\NativeMathAt{338.327}{157.025}{11.955}{\in}%
\TextAt{13}{350.381}{157.025}{U}%
\TextAt{21}{358.332}{155.232}{L}%
\TextAt{8}{368.934}{157.025}{such}%
\TextAt{8}{395.779}{157.025}{that}%
\TextAt{13}{421.596}{157.025}{N}%
\TextAt{21}{430.969}{155.165}{L}%
\TextAt{3}{436.728}{155.165}{\Character{106}}%
\TextAt{21}{439.081}{155.165}{K}%
\TextAt{9}{447.283}{157.025}{\Character{40}}%
\TextAt{13}{451.836}{157.025}{u}%
\TextAt{9}{458.498}{157.025}{\Character{41}}%
\TextAt{9}{467.118}{157.025}{\Character{61}}%
\TextAt{9}{480.304}{157.025}{1}%
\TextAt{8}{490.521}{157.025}{is}%
\TextAt{8}{502.733}{157.025}{of}%
\TextAt{8}{516.51}{157.025}{the}%
\ErrataAt{79.19}{139.96}{455.81}{form $u=a^{\sigma-1}$, with $a\in A_L$ and $\sigma=\varphi_{L\mid K}$. Choose $b\in A_K$ with $v_K(b)=v_L(a)$.}%
\ErrataAt{79.19}{122.914}{455.81}{As $a/b\in U_L$ and $(a/b)^{\sigma-1}=a^{\sigma-1}$, it follows that $H^{-1}(G(L\mid K),U_L)=1$.}%
%
%
%
%
%
%
%
%
%
%
%
%
%
%
%
%
%
%
%
%
%
%
%
%
%
%
%
%
%
%
%
%
%
%
%
%
%
%
%
%
%
%
%
%
%
%
%
%
%
%
%
%
%
%
%
%
%
%
%
%
%
%
%
%
%
%
%
%
%
%
%
%
%
%
%
%
%
%
%
%
%
%
\TextAt{8}{297.211}{43.033}{274}%
%
\Destination{78.23}{809.815}{nav.page.274}%
\Destination{79.23}{459.335}{nav.section.4.6}%
\Destination{79.36}{261.155}{nav.number.4.6.2}%
\Destination{79.23}{459.335}{outline.46}%
%
\Destination{79.226}{659.938}{errata.106}%
\Destination{79.226}{171}{errata.200}%
\Destination{79.226}{400}{nav.number.4.6.1}%
}

% Page 287 (body page 275)
\Page{%
\TextAt{2}{79.226}{796.943}{IV.}%
\TextAt{2}{98.012}{796.943}{Abstract}%
\TextAt{2}{143.252}{796.943}{Class}%
\TextAt{2}{171.844}{796.943}{Field}%
\TextAt{2}{199.575}{796.943}{Theory}%
\Vector{q 0 G 0 g 0.299 w 7.226 18.729 m 460.769 18.729 l S Q}%
\TextAt{8}{97.159}{761.124}{On}%
\TextAt{8}{116.67}{761.124}{the}%
\TextAt{8}{136.826}{761.124}{other}%
\TextAt{8}{167.396}{761.124}{hand,}%
\TextAt{8}{199.914}{761.124}{the}%
\TextAt{8}{220.07}{761.124}{homomorphism}%
\TextAt{13}{302.74}{761.124}{v}%
\TextAt{21}{308.399}{759.33}{K}%
\TextAt{9}{319.923}{761.124}{\Character{58}}%
\TextAt{13}{326.495}{761.124}{A}%
\TextAt{21}{335.271}{759.33}{K}%
\NativeMathAt{346.794}{761.124}{11.955}{\rightarrow}%
\TextAt{13}{362.07}{761.124}{Z}%
\TextAt{8}{374.838}{761.124}{gives}%
\TextAt{8}{403.519}{761.124}{rise}%
\TextAt{8}{425.038}{761.124}{to}%
\TextAt{8}{439.348}{761.124}{a}%
\TextAt{8}{449.104}{761.124}{homomorphism}%
\TextAt{13}{210.213}{733.274}{v}%
\TextAt{21}{215.872}{731.481}{K}%
\TextAt{9}{227.395}{733.274}{\Character{58}}%
\TextAt{13}{233.968}{733.274}{A}%
\TextAt{21}{242.743}{731.481}{K}%
\TextAt{9}{250.946}{733.274}{\Character{47}}%
\TextAt{13}{256.799}{733.274}{N}%
\TextAt{21}{266.173}{731.415}{L}%
\TextAt{3}{271.932}{731.415}{\Character{106}}%
\TextAt{21}{274.284}{731.415}{K}%
\TextAt{13}{282.487}{733.274}{A}%
\TextAt{21}{291.262}{731.481}{L}%
\NativeArrowAt{300.84}{733.274}{11.955}{19.26}{\longrightarrow}%
\TextAt{13}{323.422}{733.274}{Z}%
\TextAt{9}{332.293}{733.274}{\Character{47}}%
\TextAt{13}{338.146}{733.274}{nZ}%
\TextAt{9}{357.325}{733.274}{\Character{61}}%
\TextAt{18}{369.751}{733.274}{Z}%
\TextAt{9}{377.721}{733.274}{\Character{47}}%
\TextAt{13}{383.574}{733.274}{n}%
\TextAt{18}{390.561}{733.274}{Z}%
\TextAt{13}{398.531}{733.274}{\Character{59}}%
\TextAt{8}{79.226}{705.425}{where}%
\TextAt{13}{113.266}{705.425}{n}%
\TextAt{9}{123.968}{705.425}{\Character{61}}%
\TextAt{9}{136.772}{705.425}{\Character{91}}%
\TextAt{13}{140.039}{705.425}{L}%
\TextAt{9}{151.718}{705.425}{\Character{58}}%
\TextAt{13}{158.685}{705.425}{K}%
\TextAt{9}{169.502}{705.425}{\Character{93}}%
\TextAt{9}{176.448}{705.425}{\Character{61}}%
\TextAt{13}{189.288}{705.425}{f}%
\TextAt{21}{195.058}{703.566}{L}%
\TextAt{3}{200.817}{703.566}{\Character{106}}%
\TextAt{21}{203.169}{703.566}{K}%
\TextAt{8}{211.372}{705.425}{,}%
\TextAt{8}{218.81}{705.425}{because}%
\TextAt{13}{262.342}{705.425}{v}%
\TextAt{21}{268.001}{703.632}{K}%
\TextAt{9}{276.204}{705.425}{\Character{40}}%
\TextAt{13}{280.756}{705.425}{N}%
\TextAt{21}{290.13}{703.566}{L}%
\TextAt{3}{295.889}{703.566}{\Character{106}}%
\TextAt{21}{298.242}{703.566}{K}%
\TextAt{13}{306.444}{705.425}{A}%
\TextAt{21}{315.22}{703.632}{L}%
\TextAt{9}{321.477}{705.425}{\Character{41}}%
\TextAt{9}{329.726}{705.425}{\Character{61}}%
\TextAt{13}{342.563}{705.425}{f}%
\TextAt{21}{348.333}{703.566}{L}%
\TextAt{3}{354.092}{703.566}{\Character{106}}%
\TextAt{21}{356.444}{703.566}{K}%
\TextAt{13}{364.647}{705.425}{Z}%
\TextAt{9}{377.232}{705.425}{\Character{61}}%
\TextAt{13}{390.051}{705.425}{nZ}%
\TextAt{8}{405.91}{705.425}{.}%
\TextAt{8}{415.051}{705.425}{This}%
\TextAt{8}{442.002}{705.425}{homomorphism}%
\TextAt{8}{524.903}{705.425}{is}%
\TextAt{8}{79.226}{688.379}{surjective}%
\TextAt{8}{131.602}{688.379}{as}%
\TextAt{13}{145.936}{688.379}{v}%
\TextAt{21}{151.595}{686.586}{K}%
\TextAt{9}{159.798}{688.379}{\Character{40}}%
\TextAt{13}{164.35}{688.379}{\Character{25}}%
\TextAt{21}{170.991}{686.586}{K}%
\TextAt{8}{182.514}{688.379}{mod}%
\TextAt{13}{207.952}{688.379}{N}%
\TextAt{21}{217.326}{686.52}{L}%
\TextAt{3}{223.085}{686.52}{\Character{106}}%
\TextAt{21}{225.437}{686.52}{K}%
\TextAt{13}{233.64}{688.379}{A}%
\TextAt{21}{242.415}{686.586}{L}%
\TextAt{9}{248.672}{688.379}{\Character{41}}%
\TextAt{9}{256.527}{688.379}{\Character{61}}%
\TextAt{9}{268.948}{688.379}{1}%
\TextAt{8}{278.145}{688.379}{mod}%
\TextAt{13}{303.583}{688.379}{nZ}%
\TextAt{8}{319.441}{688.379}{,}%
\TextAt{8}{326.562}{688.379}{and}%
\TextAt{8}{349.289}{688.379}{it}%
\TextAt{8}{360.957}{688.379}{is}%
\TextAt{8}{372.685}{688.379}{bijective}%
\TextAt{8}{419.143}{688.379}{as}%
\TextAt{8}{433.477}{688.379}{\char35\relax }%
\TextAt{13}{443.232}{688.379}{A}%
\TextAt{21}{452.008}{686.586}{K}%
\TextAt{9}{460.21}{688.379}{\Character{47}}%
\TextAt{13}{466.063}{688.379}{N}%
\TextAt{21}{475.437}{686.52}{L}%
\TextAt{3}{481.196}{686.52}{\Character{106}}%
\TextAt{21}{483.549}{686.52}{K}%
\TextAt{13}{491.751}{688.379}{A}%
\TextAt{21}{500.527}{686.586}{L}%
\TextAt{9}{510.105}{688.379}{\Character{61}}%
\TextAt{13}{522.53}{688.379}{n}%
\TextAt{8}{529.518}{688.379}{.}%
\TextAt{8}{79.226}{671.333}{If}%
\TextAt{8}{91.983}{671.333}{now}%
\TextAt{13}{117.437}{671.333}{u}%
\NativeMathAt{129.235}{671.333}{11.955}{\in}%
\TextAt{13}{142.341}{671.333}{U}%
\TextAt{21}{150.292}{669.54}{K}%
\TextAt{8}{158.494}{671.333}{,}%
\TextAt{8}{166.975}{671.333}{then}%
\TextAt{8}{194.7}{671.333}{we}%
\TextAt{8}{212.993}{671.333}{have}%
\TextAt{13}{241.052}{671.333}{u}%
\TextAt{9}{252.851}{671.333}{\Character{61}}%
\TextAt{13}{267.091}{671.333}{N}%
\TextAt{21}{276.465}{669.474}{L}%
\TextAt{3}{282.224}{669.474}{\Character{106}}%
\TextAt{21}{284.576}{669.474}{K}%
\TextAt{9}{292.779}{671.333}{\Character{40}}%
\TextAt{13}{297.331}{671.333}{a}%
\TextAt{9}{303.476}{671.333}{\Character{41}}%
\TextAt{8}{308.029}{671.333}{,}%
\TextAt{8}{316.509}{671.333}{with}%
\TextAt{13}{344.234}{671.333}{a}%
\NativeMathAt{355.515}{671.333}{11.955}{\in}%
\TextAt{13}{368.621}{671.333}{A}%
\TextAt{21}{377.397}{669.54}{L}%
\TextAt{8}{383.654}{671.333}{,}%
\TextAt{8}{392.134}{671.333}{since}%
\TextAt{13}{421.868}{671.333}{v}%
\TextAt{21}{427.527}{669.54}{K}%
\TextAt{9}{435.73}{671.333}{\Character{40}}%
\TextAt{13}{440.282}{671.333}{u}%
\TextAt{9}{446.945}{671.333}{\Character{41}}%
\TextAt{9}{456.617}{671.333}{\Character{61}}%
\TextAt{9}{470.855}{671.333}{0}%
\TextAt{8}{476.726}{671.333}{.}%
\TextAt{8}{488.37}{671.333}{But}%
\TextAt{9}{512.676}{671.333}{0}%
\TextAt{9}{523.651}{671.333}{\Character{61}}%
\TextAt{13}{79.226}{654.287}{v}%
\TextAt{21}{84.885}{652.494}{K}%
\TextAt{9}{93.088}{654.287}{\Character{40}}%
\TextAt{13}{97.64}{654.287}{u}%
\TextAt{9}{104.303}{654.287}{\Character{41}}%
\TextAt{9}{113.532}{654.287}{\Character{61}}%
\TextAt{13}{127.337}{654.287}{v}%
\TextAt{21}{132.996}{652.494}{K}%
\TextAt{9}{141.199}{654.287}{\Character{40}}%
\TextAt{13}{145.751}{654.287}{N}%
\TextAt{21}{155.125}{652.428}{L}%
\TextAt{3}{160.884}{652.428}{\Character{106}}%
\TextAt{21}{163.236}{652.428}{K}%
\TextAt{9}{171.439}{654.287}{\Character{40}}%
\TextAt{13}{175.991}{654.287}{a}%
\TextAt{9}{182.136}{654.287}{\Character{41}\Character{41}}%
\TextAt{9}{195.908}{654.287}{\Character{61}}%
\TextAt{13}{209.722}{654.287}{nv}%
\TextAt{21}{222.369}{652.494}{L}%
\TextAt{9}{228.626}{654.287}{\Character{40}}%
\TextAt{13}{233.178}{654.287}{a}%
\TextAt{9}{239.323}{654.287}{\Character{41}}%
\TextAt{8}{243.876}{654.287}{,}%
\TextAt{8}{252.028}{654.287}{so}%
\TextAt{8}{267.201}{654.287}{we}%
\TextAt{8}{285.23}{654.287}{get}%
\TextAt{8}{305.544}{654.287}{in}%
\TextAt{8}{319.999}{654.287}{fact}%
\TextAt{13}{343.887}{654.287}{a}%
\NativeMathAt{354.72}{654.287}{11.955}{\in}%
\TextAt{13}{367.379}{654.287}{U}%
\TextAt{21}{375.329}{652.494}{L}%
\TextAt{8}{381.586}{654.287}{.}%
\TextAt{8}{392.442}{654.287}{This}%
\TextAt{8}{419.964}{654.287}{proves}%
\TextAt{8}{456.931}{654.287}{that}%
\TextAt{13}{483.103}{654.287}{H}%
\TextAt{17}{493.725}{658.626}{0}%
\TextAt{9}{498.457}{654.287}{\Character{40}}%
\TextAt{13}{503.01}{654.287}{G}%
\TextAt{9}{512.243}{654.287}{\Character{40}}%
\TextAt{13}{516.796}{654.287}{L}%
\TextAt{16}{529.449}{654.287}{\Character{106}}%
\TextAt{13}{79.226}{637.241}{K}%
\TextAt{9}{90.044}{637.241}{\Character{41}}%
\TextAt{13}{94.596}{637.241}{\Character{59}}%
\TextAt{13}{99.82}{637.241}{U}%
\TextAt{21}{107.791}{635.448}{L}%
\TextAt{9}{114.048}{637.241}{\Character{41}}%
\TextAt{9}{121.903}{637.241}{\Character{61}}%
\TextAt{9}{134.324}{637.241}{1}%
\TextAt{8}{140.199}{637.241}{.}%
\Vector{q 0 G 0 g 0.399 w 452.334 -136.759 m 452.334 -128.689 l S Q}%
\Vector{q 0 G 0 g 0.399 w 452.534 -128.888 m 459.707 -128.888 l S Q}%
\Vector{q 0 G 0 g 0.399 w 452.534 -136.559 m 459.707 -136.559 l S Q}%
\Vector{q 0 G 0 g 0.399 w 459.906 -136.759 m 459.906 -128.689 l S Q}%
\TextAt{8}{97.159}{608.569}{By}%
\TextAt{8}{115.522}{608.569}{definition,}%
\TextAt{8}{170.803}{608.569}{a}%
\TextAt{12}{180.558}{608.569}{class}%
\TextAt{12}{211.917}{608.569}{field}%
\TextAt{12}{241.219}{608.569}{theory}%
\TextAt{8}{283.289}{608.569}{is}%
\TextAt{8}{295.053}{608.569}{a}%
\TextAt{8}{304.808}{608.569}{pair}%
\TextAt{8}{328.874}{608.569}{of}%
\TextAt{8}{342.204}{608.569}{homomorphisms}%
\TextAt{9}{237.11}{580.72}{\Character{40}}%
\TextAt{13}{241.662}{580.72}{d}%
\TextAt{9}{251.066}{580.72}{\Character{58}}%
\TextAt{13}{257.638}{580.72}{G}%
\NativeMathAt{270.193}{580.72}{11.955}{\rightarrow}%
\TextAt{20}{286.133}{583.809}{\Character{98}}%
\TextAt{18}{285.469}{580.72}{Z}%
\TextAt{13}{293.439}{580.72}{\Character{59}}%
\TextAt{13}{310.619}{580.72}{v}%
\TextAt{9}{320.047}{580.72}{\Character{58}}%
\TextAt{13}{326.62}{580.72}{A}%
\TextAt{21}{335.395}{578.927}{k}%
\NativeMathAt{343.836}{580.72}{11.955}{\rightarrow}%
\TextAt{20}{359.776}{583.809}{\Character{98}}%
\TextAt{18}{359.112}{580.72}{Z}%
\TextAt{9}{367.082}{580.72}{\Character{41}}%
\TextAt{13}{371.634}{580.72}{\Character{59}}%
\TextAt{8}{79.226}{552.871}{where}%
\TextAt{13}{112.596}{552.871}{A}%
\TextAt{8}{124.829}{552.871}{is}%
\TextAt{8}{136.154}{552.871}{a}%
\TextAt{13}{145.47}{552.871}{G}%
\TextAt{8}{154.704}{552.871}{-module}%
\TextAt{8}{199.455}{552.871}{which}%
\TextAt{8}{232.502}{552.871}{satisfies}%
\TextAt{8}{275.173}{552.871}{axiom}%
\TextAt{8}{309.535}{552.871}{(6.1),}%
\TextAt{13}{340.41}{552.871}{d}%
\TextAt{8}{349.951}{552.871}{is}%
\TextAt{8}{361.276}{552.871}{a}%
\TextAt{8}{370.592}{552.871}{surjective}%
\TextAt{8}{422.564}{552.871}{continuous}%
\TextAt{8}{480.729}{552.871}{homomor-}%
\TextAt{8}{79.226}{535.825}{phism,}%
\TextAt{8}{116.71}{535.825}{and}%
\TextAt{13}{139.105}{535.825}{v}%
\TextAt{8}{148.722}{535.825}{is}%
\TextAt{8}{160.118}{535.825}{a}%
\TextAt{8}{169.506}{535.825}{henselian}%
\TextAt{8}{219.923}{535.825}{valuation.}%
\TextAt{8}{275.424}{535.825}{From}%
\TextAt{8}{305.781}{535.825}{proposition}%
\TextAt{8}{366.946}{535.825}{(6.2)}%
\TextAt{8}{394.553}{535.825}{and}%
\TextAt{8}{416.948}{535.825}{§5,}%
\TextAt{8}{435.339}{535.825}{we}%
\TextAt{8}{452.198}{535.825}{obtain}%
\TextAt{8}{488.258}{535.825}{for}%
\TextAt{8}{505.775}{535.825}{every}%
\TextAt{8}{79.226}{518.779}{finite}%
\TextAt{8}{109.138}{518.779}{Galois}%
\TextAt{8}{145.051}{518.779}{extension}%
\TextAt{13}{196.817}{518.779}{L}%
\TextAt{16}{208.103}{518.779}{\Character{106}}%
\TextAt{13}{214.744}{518.779}{K}%
\TextAt{8}{225.562}{518.779}{,}%
\TextAt{8}{232.711}{518.779}{the}%
\TextAt{8}{252.868}{518.779}{reciprocity}%
\TextAt{8}{311.077}{518.779}{homomorphism}%
\TextAt{13}{215.646}{490.93}{r}%
\TextAt{21}{220.921}{489.07}{L}%
\TextAt{3}{226.68}{489.07}{\Character{106}}%
\TextAt{21}{229.032}{489.07}{K}%
\TextAt{9}{240.556}{490.93}{\Character{58}}%
\TextAt{13}{247.128}{490.93}{G}%
\TextAt{9}{256.362}{490.93}{\Character{40}}%
\TextAt{13}{260.914}{490.93}{L}%
\TextAt{16}{272.2}{490.93}{\Character{106}}%
\TextAt{13}{278.841}{490.93}{K}%
\TextAt{9}{289.659}{490.93}{\Character{41}}%
\TextAt{7}{294.211}{495.866}{ab}%
\NativeArrowAt{306.965}{490.93}{11.955}{19.26}{\longrightarrow}%
\TextAt{13}{329.547}{490.93}{A}%
\TextAt{21}{338.322}{489.136}{K}%
\TextAt{9}{346.525}{490.93}{\Character{47}}%
\TextAt{13}{352.378}{490.93}{N}%
\TextAt{21}{361.752}{489.07}{L}%
\TextAt{3}{367.511}{489.07}{\Character{106}}%
\TextAt{21}{369.863}{489.07}{K}%
\TextAt{13}{378.066}{490.93}{A}%
\TextAt{21}{386.841}{489.136}{L}%
\TextAt{13}{393.098}{490.93}{\Character{58}}%
\TextAt{8}{79.226}{463.081}{But}%
\TextAt{8}{103.777}{463.081}{the}%
\TextAt{8}{125.243}{463.081}{class}%
\TextAt{8}{153.99}{463.081}{field}%
\TextAt{8}{180.656}{463.081}{axiom}%
\TextAt{8}{216.767}{463.081}{yields}%
\TextAt{8}{250.978}{463.081}{moreover}%
\TextAt{8}{302.702}{463.081}{the}%
\TextAt{8}{324.169}{463.081}{following}%
\TextAt{8}{374.913}{463.081}{theorem,}%
\TextAt{8}{425.327}{463.081}{which}%
\TextAt{8}{460.123}{463.081}{represents}%
\TextAt{8}{516.51}{463.081}{the}%
\TextAt{8}{79.226}{446.035}{main}%
\TextAt{8}{108.436}{446.035}{theorem}%
\TextAt{8}{153.904}{446.035}{of}%
\TextAt{8}{167.177}{446.035}{class}%
\TextAt{8}{194.558}{446.035}{field}%
\TextAt{8}{219.858}{446.035}{theory,}%
\TextAt{8}{258.834}{446.035}{and}%
\TextAt{8}{281.54}{446.035}{which}%
\TextAt{8}{314.969}{446.035}{we}%
\TextAt{8}{332.14}{446.035}{will}%
\TextAt{8}{354.188}{446.035}{call}%
\TextAt{8}{375.591}{446.035}{the}%
\TextAt{12}{395.691}{446.035}{general}%
\TextAt{12}{442.358}{446.035}{reciprocity}%
\TextAt{12}{509.887}{446.035}{law}%
\TextAt{8}{529.518}{446.035}{.}%
\TextAt{12}{79.226}{420.351}{Theorem}%
\TextAt{12}{136.24}{420.351}{(6.3).}%
\TextAt{8}{171.556}{420.351}{For}%
\TextAt{8}{192.525}{420.351}{every}%
\TextAt{8}{223.417}{420.351}{finite}%
\TextAt{8}{253.329}{420.351}{Galois}%
\TextAt{8}{289.243}{420.351}{extension}%
\TextAt{13}{341.008}{420.351}{L}%
\TextAt{16}{352.294}{420.351}{\Character{106}}%
\TextAt{13}{358.936}{420.351}{K}%
\TextAt{8}{369.753}{420.351}{,}%
\TextAt{8}{376.902}{420.351}{the}%
\TextAt{8}{397.059}{420.351}{homomorphism}%
\TextAt{13}{217.272}{392.502}{r}%
\TextAt{21}{222.547}{390.642}{L}%
\TextAt{3}{228.306}{390.642}{\Character{106}}%
\TextAt{21}{230.658}{390.642}{K}%
\TextAt{9}{242.182}{392.502}{\Character{58}}%
\TextAt{13}{248.754}{392.502}{G}%
\TextAt{9}{257.988}{392.502}{\Character{40}}%
\TextAt{13}{262.54}{392.502}{L}%
\TextAt{16}{273.825}{392.502}{\Character{106}}%
\TextAt{13}{280.467}{392.502}{K}%
\TextAt{9}{291.285}{392.502}{\Character{41}}%
\TextAt{7}{295.837}{397.438}{ab}%
\NativeArrowAt{308.591}{392.502}{11.955}{19.26}{\longrightarrow}%
\TextAt{13}{331.173}{392.502}{A}%
\TextAt{21}{339.948}{390.709}{K}%
\TextAt{9}{348.151}{392.502}{\Character{47}}%
\TextAt{13}{354.004}{392.502}{N}%
\TextAt{21}{363.377}{390.642}{L}%
\TextAt{3}{369.137}{390.642}{\Character{106}}%
\TextAt{21}{371.489}{390.642}{K}%
\TextAt{13}{379.692}{392.502}{A}%
\TextAt{21}{388.467}{390.709}{L}%
\TextAt{8}{79.226}{364.653}{is}%
\TextAt{8}{90.99}{364.653}{an}%
\TextAt{8}{107.249}{364.653}{isomorphism.}%
\TextAt{19}{79.226}{330.332}{Proof.}%
\TextAt{8}{116}{330.332}{If}%
\TextAt{13}{128.336}{330.332}{M}%
\TextAt{16}{145.328}{330.332}{\Character{106}}%
\TextAt{13}{153.067}{330.332}{K}%
\TextAt{8}{168.426}{330.332}{is}%
\TextAt{8}{180.834}{330.332}{a}%
\TextAt{8}{191.233}{330.332}{Galois}%
\TextAt{8}{227.79}{330.332}{subextension}%
\TextAt{8}{298.145}{330.332}{of}%
\TextAt{13}{312.118}{330.332}{L}%
\TextAt{16}{324.501}{330.332}{\Character{106}}%
\TextAt{13}{332.24}{330.332}{K}%
\TextAt{8}{343.058}{330.332}{,}%
\TextAt{8}{351.012}{330.332}{we}%
\TextAt{8}{368.883}{330.332}{get}%
\TextAt{8}{389.038}{330.332}{from}%
\TextAt{8}{417.322}{330.332}{(5.8)}%
\TextAt{8}{445.941}{330.332}{the}%
\TextAt{8}{466.741}{330.332}{commutative}%
\TextAt{8}{79.226}{313.286}{exact}%
\TextAt{8}{110.118}{313.286}{diagram}%
\TextAt{9}{149.387}{285.69}{1}%
\TextAt{13}{188.525}{285.69}{G}%
\TextAt{9}{197.758}{285.69}{\Character{40}}%
\TextAt{13}{202.311}{285.69}{L}%
\TextAt{16}{213.596}{285.69}{\Character{106}}%
\TextAt{13}{220.238}{285.69}{M}%
\TextAt{9}{232.811}{285.69}{\Character{41}}%
\TextAt{13}{278.883}{285.69}{G}%
\TextAt{9}{288.116}{285.69}{\Character{40}}%
\TextAt{13}{292.669}{285.69}{L}%
\TextAt{16}{303.954}{285.69}{\Character{106}}%
\TextAt{13}{310.596}{285.69}{K}%
\TextAt{9}{321.414}{285.69}{\Character{41}}%
\TextAt{13}{368.004}{285.69}{G}%
\TextAt{9}{377.238}{285.69}{\Character{40}}%
\TextAt{13}{381.79}{285.69}{M}%
\TextAt{16}{397.685}{285.69}{\Character{106}}%
\TextAt{13}{404.326}{285.69}{K}%
\TextAt{9}{415.144}{285.69}{\Character{41}}%
\TextAt{9}{455.126}{285.69}{1}%
\TextAt{13}{179.892}{243.456}{A}%
\TextAt{21}{188.667}{241.663}{M}%
\TextAt{9}{198.147}{243.456}{\Character{47}}%
\TextAt{13}{204}{243.456}{N}%
\TextAt{21}{213.374}{241.596}{L}%
\TextAt{3}{219.133}{241.596}{\Character{106}}%
\TextAt{21}{221.485}{241.596}{M}%
\TextAt{13}{230.965}{243.456}{A}%
\TextAt{21}{239.74}{241.663}{L}%
\TextAt{13}{270.649}{243.456}{A}%
\TextAt{21}{279.424}{241.663}{K}%
\TextAt{9}{287.627}{243.456}{\Character{47}}%
\TextAt{13}{293.48}{243.456}{N}%
\TextAt{21}{302.854}{241.596}{L}%
\TextAt{3}{308.613}{241.596}{\Character{106}}%
\TextAt{21}{310.965}{241.596}{K}%
\TextAt{13}{319.168}{243.456}{A}%
\TextAt{21}{327.943}{241.663}{L}%
\TextAt{13}{358.851}{243.456}{A}%
\TextAt{21}{367.627}{241.663}{K}%
\TextAt{9}{375.829}{243.456}{\Character{47}}%
\TextAt{13}{381.682}{243.456}{N}%
\TextAt{21}{391.056}{241.596}{M}%
\TextAt{3}{400.038}{241.596}{\Character{106}}%
\TextAt{21}{402.39}{241.596}{K}%
\TextAt{13}{410.593}{243.456}{A}%
\TextAt{21}{419.368}{241.663}{M}%
\TextAt{9}{453.501}{243.456}{1}%
\TextAt{13}{459.354}{243.456}{\Character{58}}%
\NativeDiagram{%
\path[draw=black,line width=0.478bp] (88.592,-485.321) -- (110.695,-485.321);%
\path[draw=black,line width=0.478bp,line cap=round,line join=round] (108.448,-482.452) .. controls (108.902,-484.173) and (109.915,-484.986) .. (110.935,-485.321) .. controls (109.915,-485.656) and (108.902,-486.469) .. (108.448,-488.19);%
\path[draw=black,line width=0.478bp] (170.717,-485.321) -- (201.055,-485.321);%
\path[draw=black,line width=0.478bp,line cap=round,line join=round] (198.807,-482.452) .. controls (199.262,-484.173) and (200.274,-484.986) .. (201.294,-485.321) .. controls (200.274,-485.656) and (199.262,-486.469) .. (198.807,-488.19);%
\path[draw=black,line width=0.478bp] (140.945,-495.878) -- (140.945,-516.521);%
\path[draw=black,line width=0.478bp,line cap=round,line join=round] (143.815,-514.273) .. controls (142.093,-514.727) and (141.28,-515.74) .. (140.946,-516.76) .. controls (140.611,-515.74) and (139.798,-514.727) .. (138.076,-514.273);%
}%
\TextAt{21}{190.894}{267.462}{r}%
\TextAt{22}{194.715}{265.724}{L}%
\TextAt{23}{199.774}{265.724}{\Character{106}}%
\TextAt{22}{202.043}{265.724}{M}%
\NativeDiagram{%
\path[draw=black,line width=0.478bp] (259.32,-485.321) -- (290.177,-485.321);%
\path[draw=black,line width=0.478bp,line cap=round,line join=round] (287.929,-482.452) .. controls (288.384,-484.173) and (289.396,-484.986) .. (290.416,-485.321) .. controls (289.396,-485.656) and (288.384,-486.469) .. (287.929,-488.19);%
\path[draw=black,line width=0.478bp] (230.427,-495.878) -- (230.427,-516.521);%
\path[draw=black,line width=0.478bp,line cap=round,line join=round] (233.296,-514.273) .. controls (231.574,-514.727) and (230.761,-515.74) .. (230.427,-516.76) .. controls (230.092,-515.74) and (229.279,-514.727) .. (227.557,-514.273);%
}%
\TextAt{21}{305.242}{267.462}{r}%
\TextAt{22}{309.063}{265.724}{L}%
\TextAt{23}{314.122}{265.724}{\Character{106}}%
\TextAt{22}{316.391}{265.724}{K}%
\NativeDiagram{%
\path[draw=black,line width=0.478bp] (353.052,-485.321) -- (377.3,-485.321);%
\path[draw=black,line width=0.478bp,line cap=round,line join=round] (375.052,-482.452) .. controls (375.506,-484.173) and (376.519,-484.986) .. (377.539,-485.321) .. controls (376.519,-485.656) and (375.506,-486.469) .. (375.052,-488.19);%
\path[draw=black,line width=0.478bp] (321.853,-495.878) -- (321.853,-516.521);%
\path[draw=black,line width=0.478bp,line cap=round,line join=round] (324.722,-514.273) .. controls (323.001,-514.727) and (322.188,-515.74) .. (321.853,-516.76) .. controls (321.519,-515.74) and (320.706,-514.727) .. (318.984,-514.273);%
}%
\TextAt{21}{396.669}{267.462}{r}%
\TextAt{22}{400.49}{265.724}{M}%
\TextAt{23}{408.078}{265.724}{\Character{106}}%
\TextAt{22}{410.348}{265.724}{K}%
\NativeDiagram{%
\path[draw=black,line width=0.478bp] (179.351,-527.556) -- (192.82,-527.556);%
\path[draw=black,line width=0.478bp,line cap=round,line join=round] (190.573,-524.686) .. controls (191.027,-526.408) and (192.04,-527.221) .. (193.06,-527.556) .. controls (192.04,-527.89) and (191.027,-528.703) .. (190.573,-530.425);%
}%
\TextAt{21}{246.514}{238.183}{N}%
\TextAt{22}{253.235}{236.445}{M}%
\TextAt{23}{260.823}{236.445}{\Character{106}}%
\TextAt{22}{263.093}{236.445}{K}%
\NativeDiagram{%
\path[draw=black,line width=0.478bp] (267.555,-527.556) -- (281.024,-527.556);%
\path[draw=black,line width=0.478bp,line cap=round,line join=round] (278.777,-524.686) .. controls (279.231,-526.408) and (280.244,-527.221) .. (281.264,-527.556) .. controls (280.244,-527.89) and (279.231,-528.703) .. (278.777,-530.425);%
\path[draw=black,line width=0.478bp] (362.204,-527.556) -- (375.674,-527.556);%
\path[draw=black,line width=0.478bp,line cap=round,line join=round] (373.426,-524.686) .. controls (373.881,-526.408) and (374.893,-527.221) .. (375.913,-527.556) .. controls (374.893,-527.89) and (373.881,-528.703) .. (373.426,-530.425);%
}%
\TextAt{8}{79.226}{212.278}{We}%
\TextAt{8}{99.371}{212.278}{use}%
\TextAt{8}{119.587}{212.278}{this}%
\TextAt{8}{142.409}{212.278}{diagram}%
\TextAt{8}{187.946}{212.278}{to}%
\TextAt{8}{202.257}{212.278}{make}%
\TextAt{8}{232.826}{212.278}{three}%
\TextAt{8}{262.738}{212.278}{reductions.}%
\TextAt{19}{97.159}{192.243}{First}%
\TextAt{19}{127.553}{192.243}{reduction.}%
\TextAt{8}{187.003}{192.243}{We}%
\TextAt{8}{208.443}{192.243}{may}%
\TextAt{8}{235.107}{192.243}{assume}%
\TextAt{8}{276.847}{192.243}{that}%
\TextAt{13}{303.511}{192.243}{G}%
\TextAt{9}{312.745}{192.243}{\Character{40}}%
\TextAt{13}{317.297}{192.243}{L}%
\TextAt{16}{330.79}{192.243}{\Character{106}}%
\TextAt{13}{339.639}{192.243}{K}%
\TextAt{9}{350.457}{192.243}{\Character{41}}%
\TextAt{8}{360.202}{192.243}{is}%
\TextAt{8}{373.262}{192.243}{abelian.}%
\TextAt{8}{422.346}{192.243}{For}%
\TextAt{8}{444.611}{192.243}{if}%
\TextAt{8}{456.63}{192.243}{the}%
\TextAt{8}{478.082}{192.243}{theorem}%
\TextAt{8}{524.903}{192.243}{is}%
\TextAt{8}{79.226}{175.197}{proved}%
\TextAt{8}{117.6}{175.197}{in}%
\TextAt{8}{131.573}{175.197}{this}%
\TextAt{8}{154.716}{175.197}{case,}%
\TextAt{8}{183.139}{175.197}{then,}%
\TextAt{8}{213.451}{175.197}{putting}%
\TextAt{13}{255.4}{175.197}{M}%
\TextAt{9}{271.84}{175.197}{\Character{61}}%
\TextAt{13}{284.811}{175.197}{L}%
\TextAt{7}{292.776}{179.536}{ab}%
\TextAt{8}{306.426}{175.197}{the}%
\TextAt{8}{326.903}{175.197}{maximal}%
\TextAt{8}{375.032}{175.197}{abelian}%
\TextAt{8}{416}{175.197}{subextension}%
\TextAt{8}{486.031}{175.197}{of}%
\TextAt{13}{499.681}{175.197}{L}%
\TextAt{16}{511.512}{175.197}{\Character{106}}%
\TextAt{13}{518.7}{175.197}{K}%
\TextAt{8}{529.518}{175.197}{,}%
\TextAt{8}{79.226}{158.151}{we}%
\TextAt{8}{96.162}{158.151}{find}%
\TextAt{13}{119.279}{158.151}{G}%
\TextAt{9}{128.513}{158.151}{\Character{40}}%
\TextAt{13}{133.065}{158.151}{L}%
\TextAt{16}{144.351}{158.151}{\Character{106}}%
\TextAt{13}{150.992}{158.151}{K}%
\TextAt{9}{161.81}{158.151}{\Character{41}}%
\TextAt{7}{166.362}{162.49}{ab}%
\TextAt{9}{179.116}{158.151}{\Character{61}}%
\TextAt{13}{191.541}{158.151}{G}%
\TextAt{9}{200.775}{158.151}{\Character{40}}%
\TextAt{13}{205.327}{158.151}{M}%
\TextAt{16}{221.222}{158.151}{\Character{106}}%
\TextAt{13}{227.863}{158.151}{K}%
\TextAt{9}{238.681}{158.151}{\Character{41}}%
\TextAt{8}{243.233}{158.151}{,}%
\TextAt{8}{250.15}{158.151}{and}%
\TextAt{8}{272.621}{158.151}{the}%
\TextAt{8}{292.486}{158.151}{commutator}%
\TextAt{8}{358.224}{158.151}{subgroup}%
\TextAt{13}{408.73}{158.151}{G}%
\TextAt{9}{417.963}{158.151}{\Character{40}}%
\TextAt{13}{422.516}{158.151}{L}%
\TextAt{16}{433.801}{158.151}{\Character{106}}%
\TextAt{13}{440.443}{158.151}{M}%
\TextAt{9}{453.016}{158.151}{\Character{41}}%
\TextAt{8}{461.175}{158.151}{of}%
\TextAt{13}{474.214}{158.151}{G}%
\TextAt{9}{483.447}{158.151}{\Character{40}}%
\TextAt{13}{488}{158.151}{L}%
\TextAt{16}{499.285}{158.151}{\Character{106}}%
\TextAt{13}{505.927}{158.151}{K}%
\TextAt{9}{516.744}{158.151}{\Character{41}}%
\TextAt{8}{524.903}{158.151}{is}%
\TextAt{8}{79.226}{141.105}{precisely}%
\TextAt{8}{127.091}{141.105}{the}%
\TextAt{8}{147.256}{141.105}{kernel}%
\TextAt{8}{181.732}{141.105}{of}%
\TextAt{13}{195.07}{141.105}{r}%
\TextAt{21}{200.345}{139.246}{L}%
\TextAt{3}{206.104}{139.246}{\Character{106}}%
\TextAt{21}{208.457}{139.246}{K}%
\TextAt{8}{216.659}{141.105}{,}%
\TextAt{8}{223.819}{141.105}{i.e.,}%
\TextAt{13}{245.935}{141.105}{G}%
\TextAt{9}{255.169}{141.105}{\Character{40}}%
\TextAt{13}{259.721}{141.105}{L}%
\TextAt{16}{271.021}{141.105}{\Character{106}}%
\TextAt{13}{277.677}{141.105}{K}%
\TextAt{9}{288.495}{141.105}{\Character{41}}%
\TextAt{7}{293.047}{145.444}{ab}%
\NativeMathAt{305.815}{141.105}{11.955}{\rightarrow}%
\TextAt{13}{321.106}{141.105}{A}%
\TextAt{21}{329.881}{139.312}{K}%
\TextAt{9}{338.084}{141.105}{\Character{47}}%
\TextAt{13}{343.937}{141.105}{N}%
\TextAt{21}{353.311}{139.246}{L}%
\TextAt{3}{359.07}{139.246}{\Character{106}}%
\TextAt{21}{361.422}{139.246}{K}%
\TextAt{13}{369.625}{141.105}{A}%
\TextAt{21}{378.4}{139.312}{L}%
\TextAt{8}{388.563}{141.105}{is}%
\TextAt{8}{400.336}{141.105}{injective.}%
\TextAt{8}{451.406}{141.105}{The}%
\TextAt{8}{475.468}{141.105}{surjectivity}%
\TextAt{8}{79.226}{124.059}{follows}%
\TextAt{8}{118.076}{124.059}{by}%
\TextAt{8}{134.749}{124.059}{induction}%
\TextAt{8}{187.191}{124.059}{on}%
\TextAt{8}{203.864}{124.059}{the}%
\TextAt{8}{224.434}{124.059}{degree.}%
\TextAt{8}{266.642}{124.059}{Indeed,}%
\TextAt{8}{308.44}{124.059}{in}%
\TextAt{8}{322.507}{124.059}{the}%
\TextAt{8}{343.077}{124.059}{case}%
\TextAt{8}{368.262}{124.059}{where}%
\TextAt{13}{402.485}{124.059}{G}%
\TextAt{9}{411.719}{124.059}{\Character{40}}%
\TextAt{13}{416.271}{124.059}{L}%
\TextAt{16}{428.261}{124.059}{\Character{106}}%
\TextAt{13}{435.608}{124.059}{K}%
\TextAt{9}{446.426}{124.059}{\Character{41}}%
\TextAt{8}{455.29}{124.059}{is}%
\TextAt{8}{467.467}{124.059}{solvable,}%
\TextAt{8}{515.207}{124.059}{one}%
\TextAt{8}{79.226}{107.013}{has}%
\TextAt{8}{100.605}{107.013}{either}%
\TextAt{13}{134.274}{107.013}{M}%
\TextAt{9}{151.029}{107.013}{\Character{61}}%
\TextAt{13}{164.315}{107.013}{L}%
\TextAt{8}{176.683}{107.013}{or}%
\TextAt{9}{191.498}{107.013}{\Character{91}}%
\TextAt{13}{194.75}{107.013}{L}%
\TextAt{9}{206.896}{107.013}{\Character{58}}%
\TextAt{13}{214.329}{107.013}{M}%
\TextAt{9}{226.903}{107.013}{\Character{93}}%
\TextAt{13}{234.336}{107.013}{\Character{60}}%
\TextAt{9}{247.622}{107.013}{\Character{91}}%
\TextAt{13}{250.874}{107.013}{L}%
\TextAt{9}{263.02}{107.013}{\Character{58}}%
\TextAt{13}{270.454}{107.013}{K}%
\TextAt{9}{281.271}{107.013}{\Character{93}}%
\TextAt{8}{284.523}{107.013}{,}%
\TextAt{8}{292.304}{107.013}{and}%
\TextAt{8}{315.571}{107.013}{if}%
\TextAt{13}{326.8}{107.013}{r}%
\TextAt{21}{332.076}{105.154}{M}%
\TextAt{3}{341.057}{105.154}{\Character{106}}%
\TextAt{21}{343.41}{105.154}{K}%
\TextAt{8}{356.015}{107.013}{and}%
\TextAt{13}{379.283}{107.013}{r}%
\TextAt{21}{384.558}{105.154}{L}%
\TextAt{3}{390.317}{105.154}{\Character{106}}%
\TextAt{21}{392.669}{105.154}{M}%
\TextAt{8}{406.552}{107.013}{are}%
\TextAt{8}{426.568}{107.013}{surjective,}%
\TextAt{8}{482.863}{107.013}{then}%
\TextAt{8}{510.028}{107.013}{so}%
\TextAt{8}{524.903}{107.013}{is}%
\TextAt{13}{79.226}{89.967}{r}%
\TextAt{21}{84.501}{88.108}{L}%
\TextAt{3}{90.26}{88.108}{\Character{106}}%
\TextAt{21}{92.613}{88.108}{K}%
\TextAt{8}{100.815}{89.967}{.}%
\TextAt{8}{111.114}{89.967}{In}%
\TextAt{8}{126.352}{89.967}{the}%
\TextAt{8}{147.126}{89.967}{general}%
\TextAt{8}{188.067}{89.967}{case,}%
\TextAt{8}{216.861}{89.967}{let}%
\TextAt{13}{234.383}{89.967}{M}%
\TextAt{8}{251.471}{89.967}{be}%
\TextAt{8}{268.012}{89.967}{the}%
\TextAt{8}{288.785}{89.967}{fixed}%
\TextAt{8}{317.688}{89.967}{field}%
\TextAt{8}{343.662}{89.967}{of}%
\TextAt{8}{357.609}{89.967}{a}%
\TextAt{13}{367.981}{89.967}{p}%
\TextAt{8}{373.856}{89.967}{-Sylow}%
\TextAt{8}{412.192}{89.967}{subgroup.}%
\TextAt{13}{469.391}{89.967}{M}%
\TextAt{16}{486.336}{89.967}{\Character{106}}%
\TextAt{13}{494.029}{89.967}{K}%
\TextAt{8}{509.361}{89.967}{need}%
\TextAt{8}{79.226}{72.921}{not}%
\TextAt{8}{100.4}{72.921}{be}%
\TextAt{8}{116.684}{72.921}{Galois,}%
\TextAt{8}{156.298}{72.921}{but}%
\TextAt{8}{178.118}{72.921}{we}%
\TextAt{8}{195.705}{72.921}{may}%
\TextAt{8}{221.433}{72.921}{use}%
\TextAt{8}{242.009}{72.921}{the}%
\TextAt{8}{262.525}{72.921}{left}%
\TextAt{8}{283.364}{72.921}{part}%
\TextAt{8}{309.093}{72.921}{of}%
\TextAt{8}{322.782}{72.921}{the}%
\TextAt{8}{343.298}{72.921}{diagram,}%
\TextAt{8}{392.537}{72.921}{where}%
\TextAt{13}{426.706}{72.921}{r}%
\TextAt{21}{431.981}{71.062}{L}%
\TextAt{3}{437.74}{71.062}{\Character{106}}%
\TextAt{21}{440.093}{71.062}{M}%
\TextAt{8}{453.829}{72.921}{is}%
\TextAt{8}{465.953}{72.921}{surjective.}%
\TextAt{8}{523.994}{72.921}{It}%
\TextAt{8}{297.211}{43.033}{275}%
\Destination{78.23}{809.815}{nav.page.275}%
\Destination{79.36}{776.155}{nav.print.300}%
\Destination{79.36}{190.155}{nav.print.301}%
\Destination{79.36}{434.355}{nav.number.4.6.3}%
\Hotlink{422.127}{533.655}{6.858}{11.047}{nav.section.4.5}%
\Hotlink{366.446}{533.655}{25.078}{11.047}{nav.number.4.6.2}%
\Hotlink{416.822}{328.155}{25.078}{11.047}{nav.number.4.5.8}%
}

% Page 288 (body page 276)
\Page{%
\TextAt{2}{79.226}{796.943}{IV.}%
\TextAt{2}{98.012}{796.943}{Abstract}%
\TextAt{2}{143.252}{796.943}{Class}%
\TextAt{2}{171.844}{796.943}{Field}%
\TextAt{2}{199.575}{796.943}{Theory}%
\Vector{q 0 G 0 g 0.299 w 7.226 18.729 m 460.769 18.729 l S Q}%
\ErrataAt{79.19}{761.124}{455.81}{then suffices to show that the image of $N_{M\mid K}$ contains the $p$-primary subgroup $S_p$ of}%
\ErrataAt{79.19}{744.078}{455.81}{$A_K/N_{L\mid K}A_L$ for every $p$. The quotient has exponent dividing $[L:K]$, so the primary}%
\ErrataAt{79.19}{727.032}{455.81}{subgroups generate it. Now the inclusion $A_K\subseteq A_M$ induces a homomorphism}%
\TextAt{13}{219.649}{698.031}{i}%
\TextAt{9}{226.963}{698.031}{\Character{58}}%
\TextAt{13}{233.535}{698.031}{A}%
\TextAt{21}{242.311}{696.237}{K}%
\TextAt{9}{250.513}{698.031}{\Character{47}}%
\TextAt{13}{256.366}{698.031}{N}%
\TextAt{21}{265.74}{696.171}{L}%
\TextAt{3}{271.499}{696.171}{\Character{106}}%
\TextAt{21}{273.852}{696.171}{K}%
\TextAt{13}{282.054}{698.031}{A}%
\TextAt{21}{290.83}{696.237}{L}%
\NativeArrowAt{300.408}{698.031}{11.955}{19.26}{\longrightarrow}%
\TextAt{13}{322.99}{698.031}{A}%
\TextAt{21}{331.765}{696.237}{M}%
\TextAt{9}{341.245}{698.031}{\Character{47}}%
\TextAt{13}{347.098}{698.031}{N}%
\TextAt{21}{356.472}{696.171}{L}%
\TextAt{3}{362.231}{696.171}{\Character{106}}%
\TextAt{21}{364.583}{696.171}{M}%
\TextAt{13}{374.063}{698.031}{A}%
\TextAt{21}{382.838}{696.237}{L}%
\TextAt{13}{389.095}{698.031}{\Character{59}}%
\TextAt{8}{79.226}{666.347}{such}%
\TextAt{8}{105.908}{666.347}{that}%
\TextAt{13}{131.561}{666.347}{N}%
\TextAt{21}{140.935}{664.487}{M}%
\TextAt{3}{149.916}{664.487}{\Character{106}}%
\TextAt{21}{152.269}{664.487}{K}%
\TextAt{16}{163.322}{666.347}{\Character{14}}%
\TextAt{13}{172.15}{666.347}{i}%
\TextAt{9}{179.949}{666.347}{\Character{61}}%
\TextAt{9}{192.849}{666.347}{\Character{91}}%
\TextAt{13}{196.111}{666.347}{M}%
\TextAt{9}{212.49}{666.347}{\Character{58}}%
\TextAt{13}{219.547}{666.347}{K}%
\TextAt{9}{230.365}{666.347}{\Character{93}}%
\TextAt{8}{233.616}{666.347}{.}%
\TextAt{8}{242.918}{666.347}{As}%
\TextAt{9}{260.489}{666.347}{\Character{40}\Character{91}}%
\TextAt{13}{268.293}{666.347}{M}%
\TextAt{9}{284.673}{666.347}{\Character{58}}%
\TextAt{13}{291.73}{666.347}{K}%
\TextAt{9}{302.547}{666.347}{\Character{93}}%
\TextAt{13}{305.799}{666.347}{\Character{59}}%
\TextAt{13}{311.023}{666.347}{p}%
\TextAt{9}{316.918}{666.347}{\Character{41}}%
\TextAt{9}{325.263}{666.347}{\Character{61}}%
\TextAt{9}{338.162}{666.347}{1}%
\TextAt{8}{344.039}{666.347}{,}%
\TextAt{13}{351.544}{666.347}{S}%
\TextAt{21}{358.743}{664.554}{p}%
\TextAt{17}{369.897}{674.054}{\Character{91}}%
\TextAt{21}{372.25}{674.054}{M}%
\TextAt{17}{381.231}{674.054}{\Character{58}}%
\TextAt{21}{383.584}{674.054}{K}%
\TextAt{17}{391.288}{674.054}{\Character{93}}%
\NativeArrowAt{367.31}{666.347}{11.955}{30.306}{\longrightarrow}%
\TextAt{13}{401.445}{666.347}{S}%
\TextAt{21}{408.644}{664.554}{p}%
\TextAt{8}{417.587}{666.347}{is}%
\TextAt{8}{429.635}{666.347}{surjective,}%
\TextAt{8}{485.654}{666.347}{so}%
\TextAt{13}{500.309}{666.347}{S}%
\TextAt{21}{507.508}{664.554}{p}%
\TextAt{8}{516.451}{666.347}{lies}%
\TextAt{8}{79.226}{649.301}{in}%
\TextAt{8}{92.879}{649.301}{the}%
\TextAt{8}{113.035}{649.301}{image}%
\TextAt{8}{146.857}{649.301}{of}%
\TextAt{13}{160.187}{649.301}{N}%
\TextAt{21}{169.56}{647.441}{M}%
\TextAt{3}{178.542}{647.441}{\Character{106}}%
\TextAt{21}{180.894}{647.441}{K}%
\TextAt{8}{189.097}{649.301}{,}%
\TextAt{8}{196.246}{649.301}{and}%
\TextAt{8}{219.009}{649.301}{therefore}%
\TextAt{8}{268.109}{649.301}{of}%
\TextAt{13}{281.439}{649.301}{r}%
\TextAt{21}{286.714}{647.441}{L}%
\TextAt{3}{292.473}{647.441}{\Character{106}}%
\TextAt{21}{294.825}{647.441}{K}%
\TextAt{8}{303.028}{649.301}{.}%
\TextAt{19}{97.159}{629.266}{Second}%
\TextAt{19}{137.169}{629.266}{reduction.}%
\TextAt{8}{195.962}{629.266}{We}%
\TextAt{8}{217.184}{629.266}{may}%
\TextAt{8}{243.629}{629.266}{assume}%
\TextAt{8}{285.15}{629.266}{that}%
\TextAt{13}{311.595}{629.266}{L}%
\TextAt{16}{324.715}{629.266}{\Character{106}}%
\TextAt{13}{333.192}{629.266}{K}%
\TextAt{8}{348.984}{629.266}{is}%
\TextAt{8}{361.824}{629.266}{cyclic.}%
\TextAt{8}{401.788}{629.266}{For}%
\TextAt{8}{423.834}{629.266}{if}%
\TextAt{13}{435.634}{629.266}{M}%
\TextAt{16}{453.364}{629.266}{\Character{106}}%
\TextAt{13}{461.84}{629.266}{K}%
\TextAt{8}{477.632}{629.266}{varies}%
\TextAt{8}{511.621}{629.266}{over}%
\TextAt{8}{79.226}{612.22}{all}%
\TextAt{8}{95.85}{612.22}{cyclic}%
\TextAt{8}{128.397}{612.22}{subextensions}%
\TextAt{8}{203.087}{612.22}{of}%
\TextAt{13}{216.782}{612.22}{L}%
\TextAt{16}{228.688}{612.22}{\Character{106}}%
\TextAt{13}{235.951}{612.22}{K}%
\TextAt{8}{246.768}{612.22}{,}%
\TextAt{8}{254.373}{612.22}{then}%
\TextAt{8}{281.398}{612.22}{the}%
\TextAt{8}{301.918}{612.22}{diagram}%
\TextAt{8}{347.82}{612.22}{shows}%
\TextAt{8}{381.802}{612.22}{that}%
\TextAt{8}{407.536}{612.22}{the}%
\TextAt{8}{428.057}{612.22}{kernel}%
\TextAt{8}{462.888}{612.22}{of}%
\TextAt{13}{476.582}{612.22}{r}%
\TextAt{21}{481.858}{610.361}{L}%
\TextAt{3}{487.617}{610.36}{\Character{106}}%
\TextAt{21}{489.969}{610.361}{K}%
\TextAt{8}{502.433}{612.22}{lies}%
\TextAt{8}{523.014}{612.22}{in}%
\TextAt{8}{79.226}{595.174}{the}%
\TextAt{8}{99.706}{595.174}{kernel}%
\TextAt{8}{134.496}{595.174}{of}%
\TextAt{8}{148.15}{595.174}{the}%
\TextAt{8}{168.63}{595.174}{map}%
\TextAt{13}{194.967}{595.174}{G}%
\TextAt{9}{204.201}{595.174}{\Character{40}}%
\TextAt{13}{208.753}{595.174}{L}%
\TextAt{16}{220.59}{595.174}{\Character{106}}%
\TextAt{13}{227.783}{595.174}{K}%
\TextAt{9}{238.6}{595.174}{\Character{41}}%
\NativeMathAt{247.025}{595.174}{11.955}{\rightarrow}%
\TextAt{20}{262.852}{604.141}{\Character{81}}%
\TextAt{21}{274.143}{591.687}{M}%
\TextAt{13}{285.615}{595.174}{G}%
\TextAt{9}{294.849}{595.174}{\Character{40}}%
\TextAt{13}{299.401}{595.174}{M}%
\TextAt{16}{315.847}{595.174}{\Character{106}}%
\TextAt{13}{323.04}{595.174}{K}%
\TextAt{9}{333.858}{595.174}{\Character{41}}%
\TextAt{8}{338.41}{595.174}{.}%
\TextAt{8}{347.828}{595.174}{Since}%
\TextAt{13}{378.709}{595.174}{G}%
\TextAt{9}{387.943}{595.174}{\Character{40}}%
\TextAt{13}{392.495}{595.174}{L}%
\TextAt{16}{404.332}{595.174}{\Character{106}}%
\TextAt{13}{411.525}{595.174}{K}%
\TextAt{9}{422.342}{595.174}{\Character{41}}%
\TextAt{8}{431.116}{595.174}{is}%
\TextAt{8}{443.203}{595.174}{abelian,}%
\TextAt{8}{487.507}{595.174}{this}%
\TextAt{8}{510.652}{595.174}{map}%
\TextAt{8}{79.226}{578.128}{is}%
\TextAt{8}{91.892}{578.128}{injective}%
\TextAt{8}{139.287}{578.128}{and}%
\TextAt{8}{162.951}{578.128}{hence}%
\TextAt{8}{196.358}{578.128}{the}%
\TextAt{8}{217.416}{578.128}{same}%
\TextAt{8}{247.644}{578.128}{is}%
\TextAt{8}{260.309}{578.128}{true}%
\TextAt{8}{285.922}{578.128}{of}%
\TextAt{13}{300.153}{578.128}{r}%
\TextAt{21}{305.429}{576.269}{L}%
\TextAt{3}{311.188}{576.269}{\Character{106}}%
\TextAt{21}{313.54}{576.269}{K}%
\TextAt{8}{321.743}{578.128}{.}%
\TextAt{8}{332.896}{578.128}{Choosing}%
\TextAt{8}{384.917}{578.128}{a}%
\TextAt{8}{395.574}{578.128}{proper}%
\TextAt{8}{433.871}{578.128}{cyclic}%
\TextAt{8}{466.956}{578.128}{subextension}%
\TextAt{13}{79.226}{561.082}{M}%
\TextAt{16}{95.596}{561.082}{\Character{106}}%
\TextAt{13}{102.713}{561.082}{K}%
\TextAt{8}{117.707}{561.082}{of}%
\TextAt{13}{131.317}{561.082}{L}%
\TextAt{16}{143.077}{561.082}{\Character{106}}%
\TextAt{13}{150.195}{561.082}{K}%
\TextAt{8}{161.012}{561.082}{,}%
\TextAt{8}{168.51}{561.082}{surjectivity}%
\TextAt{8}{229.988}{561.082}{follows}%
\TextAt{8}{268.703}{561.082}{by}%
\TextAt{8}{285.241}{561.082}{induction}%
\TextAt{8}{337.549}{561.082}{on}%
\TextAt{8}{354.087}{561.082}{the}%
\TextAt{8}{374.522}{561.082}{degree}%
\TextAt{8}{411.217}{561.082}{as}%
\TextAt{8}{425.866}{561.082}{in}%
\TextAt{8}{439.798}{561.082}{the}%
\TextAt{8}{460.233}{561.082}{first}%
\TextAt{8}{484.638}{561.082}{reduction}%
\TextAt{8}{79.226}{544.036}{for}%
\TextAt{8}{97.111}{544.036}{solvable}%
\TextAt{8}{141.082}{544.036}{extensions.}%
\TextAt{19}{97.159}{524.002}{Third}%
\TextAt{19}{129.699}{524.002}{reduction.}%
\TextAt{8}{185.557}{524.002}{Let}%
\TextAt{13}{206.624}{524.002}{L}%
\TextAt{16}{218.076}{524.002}{\Character{106}}%
\TextAt{13}{224.885}{524.002}{K}%
\TextAt{8}{239.699}{524.002}{be}%
\TextAt{8}{255.721}{524.002}{cyclic.}%
\TextAt{8}{292.749}{524.002}{We}%
\TextAt{8}{312.992}{524.002}{may}%
\TextAt{8}{338.459}{524.002}{assume}%
\TextAt{8}{379.001}{524.002}{that}%
\TextAt{13}{404.468}{524.002}{f}%
\TextAt{21}{410.238}{522.142}{L}%
\TextAt{3}{415.997}{522.142}{\Character{106}}%
\TextAt{21}{418.35}{522.142}{K}%
\TextAt{9}{430.04}{524.002}{\Character{61}}%
\TextAt{9}{442.617}{524.002}{1}%
\TextAt{8}{448.486}{524.002}{.}%
\TextAt{8}{457.229}{524.002}{To}%
\TextAt{8}{474.554}{524.002}{see}%
\TextAt{8}{493.565}{524.002}{this,}%
\TextAt{8}{519.762}{524.002}{let}%
\TextAt{13}{79.226}{506.956}{M}%
\TextAt{9}{95.121}{506.956}{\Character{61}}%
\TextAt{13}{107.546}{506.956}{L}%
\TextAt{16}{118.113}{506.956}{\Character{92}}%
\TextAt{20}{131.423}{509.978}{\Character{101}}%
\TextAt{13}{128.685}{506.956}{K}%
\TextAt{8}{143.373}{506.956}{be}%
\TextAt{8}{159.271}{506.956}{the}%
\TextAt{8}{179.4}{506.956}{maximal}%
\TextAt{8}{227.183}{506.956}{unramified}%
\TextAt{8}{285.688}{506.956}{subextension}%
\TextAt{8}{355.372}{506.956}{of}%
\TextAt{13}{368.675}{506.956}{L}%
\TextAt{16}{379.961}{506.956}{\Character{106}}%
\TextAt{13}{386.603}{506.956}{K}%
\TextAt{8}{397.42}{506.956}{.}%
\TextAt{8}{405.86}{506.956}{Then}%
\TextAt{13}{436.39}{506.956}{f}%
\TextAt{21}{442.16}{505.096}{L}%
\TextAt{3}{447.919}{505.096}{\Character{106}}%
\TextAt{21}{450.272}{505.096}{M}%
\TextAt{9}{463.072}{506.956}{\Character{61}}%
\TextAt{9}{475.482}{506.956}{1}%
\TextAt{8}{485.222}{506.956}{and}%
\TextAt{13}{507.957}{506.956}{r}%
\TextAt{21}{513.233}{505.096}{M}%
\TextAt{3}{522.214}{505.096}{\Character{106}}%
\TextAt{21}{524.567}{505.096}{K}%
\TextAt{8}{79.226}{489.91}{is}%
\TextAt{8}{91.799}{489.91}{an}%
\TextAt{8}{108.867}{489.91}{isomorphism}%
\TextAt{8}{178.095}{489.91}{by}%
\TextAt{8}{195.163}{489.91}{(5.7).}%
\TextAt{8}{230.116}{489.91}{In}%
\TextAt{8}{245.546}{489.91}{the}%
\TextAt{8}{266.511}{489.91}{bottom}%
\TextAt{8}{308.625}{489.91}{sequence}%
\TextAt{8}{357.936}{489.91}{of}%
\TextAt{8}{372.075}{489.91}{our}%
\TextAt{8}{393.698}{489.91}{diagram,}%
\TextAt{8}{443.498}{489.91}{the}%
\TextAt{8}{464.463}{489.91}{map}%
\TextAt{13}{491.286}{489.91}{N}%
\TextAt{21}{500.66}{488.05}{M}%
\TextAt{3}{509.642}{488.05}{\Character{106}}%
\TextAt{21}{511.994}{488.05}{K}%
\TextAt{8}{524.903}{489.91}{is}%
\TextAt{8}{79.226}{472.864}{injective}%
\TextAt{8}{125.804}{472.864}{because}%
\TextAt{8}{169.191}{472.864}{the}%
\TextAt{8}{189.432}{472.864}{groups}%
\TextAt{8}{227.307}{472.864}{in}%
\TextAt{8}{241.044}{472.864}{this}%
\TextAt{8}{263.951}{472.864}{sequence}%
\TextAt{8}{312.538}{472.864}{have}%
\TextAt{8}{339.617}{472.864}{the}%
\TextAt{8}{359.859}{472.864}{respective}%
\TextAt{8}{414.303}{472.864}{orders}%
\TextAt{9}{449.572}{472.864}{\Character{91}}%
\TextAt{13}{452.824}{472.864}{L}%
\TextAt{9}{464.253}{472.864}{\Character{58}}%
\TextAt{13}{470.97}{472.864}{M}%
\TextAt{9}{483.544}{472.864}{\Character{93}}%
\TextAt{8}{486.795}{472.864}{,}%
\TextAt{9}{494.05}{472.864}{\Character{91}}%
\TextAt{13}{497.302}{472.864}{L}%
\TextAt{9}{508.732}{472.864}{\Character{58}}%
\TextAt{13}{515.448}{472.864}{K}%
\TextAt{9}{526.266}{472.864}{\Character{93}}%
\TextAt{8}{529.518}{472.864}{,}%
\TextAt{9}{79.226}{455.818}{\Character{91}}%
\TextAt{13}{82.478}{455.818}{M}%
\TextAt{9}{98.372}{455.818}{\Character{58}}%
\TextAt{13}{104.945}{455.818}{K}%
\TextAt{9}{115.762}{455.818}{\Character{93}}%
\TextAt{8}{122.911}{455.818}{by}%
\TextAt{8}{139.17}{455.818}{axiom}%
\TextAt{8}{173.972}{455.818}{(6.1).}%
\TextAt{8}{206.498}{455.818}{Therefore}%
\TextAt{13}{259.495}{455.818}{r}%
\TextAt{21}{264.771}{453.958}{L}%
\TextAt{3}{270.53}{453.958}{\Character{106}}%
\TextAt{21}{272.882}{453.958}{K}%
\TextAt{8}{284.982}{455.818}{is}%
\TextAt{8}{296.746}{455.818}{an}%
\TextAt{8}{313.005}{455.818}{isomorphism}%
\TextAt{8}{381.424}{455.818}{if}%
\TextAt{13}{392.148}{455.818}{r}%
\TextAt{21}{397.423}{453.958}{L}%
\TextAt{3}{403.183}{453.958}{\Character{106}}%
\TextAt{21}{405.535}{453.958}{M}%
\TextAt{8}{418.912}{455.818}{is.}%
\TextAt{8}{97.159}{435.783}{Now}%
\TextAt{8}{124.777}{435.783}{let}%
\TextAt{13}{142.64}{435.783}{L}%
\TextAt{16}{155.558}{435.783}{\Character{106}}%
\TextAt{13}{163.832}{435.783}{K}%
\TextAt{8}{179.505}{435.783}{be}%
\TextAt{8}{196.388}{435.783}{cyclic}%
\TextAt{8}{229.529}{435.783}{and}%
\TextAt{8}{253.25}{435.783}{totally}%
\TextAt{8}{291.616}{435.783}{ramified,}%
\TextAt{8}{341.591}{435.783}{i.e.,}%
\TextAt{13}{364.893}{435.783}{f}%
\TextAt{21}{370.663}{433.923}{L}%
\TextAt{3}{376.422}{433.923}{\Character{106}}%
\TextAt{21}{378.775}{433.923}{K}%
\TextAt{9}{391.931}{435.783}{\Character{61}}%
\TextAt{9}{405.978}{435.783}{1}%
\TextAt{8}{411.842}{435.783}{.}%
\TextAt{8}{423.165}{435.783}{Let}%
\TextAt{13}{445.092}{435.783}{\Character{27}}%
\TextAt{8}{457.03}{435.783}{be}%
\TextAt{8}{473.912}{435.783}{a}%
\TextAt{8}{484.626}{435.783}{generator}%
\TextAt{8}{79.226}{418.737}{of}%
\TextAt{13}{93.864}{418.737}{G}%
\TextAt{9}{103.098}{418.737}{\Character{40}}%
\TextAt{13}{107.65}{418.737}{L}%
\TextAt{16}{121.164}{418.737}{\Character{106}}%
\TextAt{13}{130.035}{418.737}{K}%
\TextAt{9}{140.852}{418.737}{\Character{41}}%
\TextAt{8}{145.405}{418.737}{.}%
\TextAt{8}{157.777}{418.737}{We}%
\TextAt{8}{179.229}{418.737}{view}%
\TextAt{13}{207.52}{418.737}{\Character{27}}%
\TextAt{8}{219.808}{418.737}{via}%
\TextAt{8}{240.304}{418.737}{the}%
\TextAt{8}{261.768}{418.737}{isomorphism}%
\TextAt{13}{331.495}{418.737}{G}%
\TextAt{9}{340.729}{418.737}{\Character{40}}%
\TextAt{13}{345.281}{418.737}{L}%
\TextAt{16}{358.796}{418.737}{\Character{106}}%
\TextAt{13}{367.666}{418.737}{K}%
\TextAt{9}{378.484}{418.737}{\Character{41}}%
\NativeMathAt{388.586}{418.737}{11.955}{\cong}%
\TextAt{13}{403.434}{418.737}{G}%
\TextAt{9}{412.668}{418.737}{\Character{40}}%
\TextAt{20}{418.206}{421.759}{\Character{101}}%
\TextAt{13}{417.22}{418.737}{L}%
\TextAt{16}{430.734}{418.737}{\Character{106}}%
\TextAt{20}{442.343}{421.759}{\Character{101}}%
\TextAt{13}{439.605}{418.737}{K}%
\TextAt{9}{450.422}{418.737}{\Character{41}}%
\TextAt{8}{460.18}{418.737}{as}%
\TextAt{8}{475.858}{418.737}{an}%
\TextAt{8}{493.425}{418.737}{element}%
\TextAt{8}{79.226}{401.691}{of}%
\TextAt{13}{93.387}{401.691}{G}%
\TextAt{9}{102.62}{401.691}{\Character{40}}%
\TextAt{20}{108.159}{404.713}{\Character{101}}%
\TextAt{13}{107.173}{401.691}{L}%
\TextAt{16}{119.873}{401.691}{\Character{106}}%
\TextAt{20}{130.668}{404.713}{\Character{101}}%
\TextAt{13}{127.93}{401.691}{K}%
\TextAt{9}{138.748}{401.691}{\Character{41}}%
\TextAt{8}{143.3}{401.691}{,}%
\TextAt{8}{151.487}{401.691}{and}%
\TextAt{8}{175.08}{401.691}{obtain}%
\TextAt{8}{212.338}{401.691}{the}%
\TextAt{8}{233.325}{401.691}{element}%
\TextAt{20}{277.618}{401.691}{\Character{101}}%
\TextAt{13}{277.397}{401.691}{\Character{27}}%
\TextAt{9}{289.216}{401.691}{\Character{61}}%
\TextAt{13}{303.056}{401.691}{\Character{27}}%
\TextAt{13}{310.122}{401.691}{\Character{39}}%
\TextAt{21}{317.807}{399.898}{L}%
\NativeMathAt{328.8}{401.691}{11.955}{\in}%
\TextAt{8}{341.506}{401.691}{Frob}%
\TextAt{9}{366.062}{401.691}{\Character{40}}%
\TextAt{20}{371.601}{404.713}{\Character{101}}%
\TextAt{13}{370.615}{401.691}{L}%
\TextAt{16}{383.315}{401.691}{\Character{106}}%
\TextAt{13}{391.372}{401.691}{K}%
\TextAt{9}{402.19}{401.691}{\Character{41}}%
\TextAt{8}{406.742}{401.691}{,}%
\TextAt{8}{414.929}{401.691}{which}%
\TextAt{8}{449.246}{401.691}{is}%
\TextAt{8}{461.84}{401.691}{a}%
\TextAt{8}{472.426}{401.691}{preimage}%
\TextAt{8}{523.337}{401.691}{of}%
\TextAt{13}{79.226}{384.645}{\Character{27}}%
\NativeMathAt{90.838}{384.645}{11.955}{\in}%
\TextAt{13}{103.338}{384.645}{G}%
\TextAt{9}{112.572}{384.645}{\Character{40}}%
\TextAt{13}{117.124}{384.645}{L}%
\TextAt{16}{129.618}{384.645}{\Character{106}}%
\TextAt{13}{137.468}{384.645}{K}%
\TextAt{9}{148.286}{384.645}{\Character{41}}%
\TextAt{8}{157.445}{384.645}{such}%
\TextAt{8}{184.552}{384.645}{that}%
\TextAt{13}{210.63}{384.645}{d}%
\TextAt{21}{216.712}{382.852}{K}%
\TextAt{9}{224.915}{384.645}{\Character{40}}%
\TextAt{20}{229.688}{384.645}{\Character{101}}%
\TextAt{13}{229.467}{384.645}{\Character{27}}%
\TextAt{9}{236.55}{384.645}{\Character{41}}%
\TextAt{9}{245.612}{384.645}{\Character{61}}%
\TextAt{13}{259.266}{384.645}{d}%
\TextAt{21}{265.349}{382.852}{K}%
\TextAt{9}{273.551}{384.645}{\Character{40}}%
\TextAt{13}{278.104}{384.645}{\Character{39}}%
\TextAt{21}{285.772}{382.852}{L}%
\TextAt{9}{292.029}{384.645}{\Character{41}}%
\TextAt{9}{301.091}{384.645}{\Character{61}}%
\TextAt{13}{314.745}{384.645}{f}%
\TextAt{21}{320.515}{382.785}{L}%
\TextAt{3}{326.274}{382.785}{\Character{106}}%
\TextAt{21}{328.627}{382.785}{K}%
\TextAt{9}{341.359}{384.645}{\Character{61}}%
\TextAt{9}{354.976}{384.645}{1}%
\TextAt{8}{360.846}{384.645}{.}%
\TextAt{8}{371.423}{384.645}{We}%
\TextAt{8}{392.276}{384.645}{thus}%
\TextAt{8}{418.737}{384.645}{find}%
\TextAt{8}{442.855}{384.645}{for}%
\TextAt{8}{461.449}{384.645}{the}%
\TextAt{8}{482.315}{384.645}{fixed}%
\TextAt{8}{511.31}{384.645}{field}%
\TextAt{9}{79.226}{367.599}{\Character{6}}%
\TextAt{16}{91.001}{367.599}{\Character{106}}%
\TextAt{13}{97.643}{367.599}{K}%
\TextAt{8}{112.318}{367.599}{of}%
\TextAt{20}{125.828}{367.599}{\Character{101}}%
\TextAt{13}{125.607}{367.599}{\Character{27}}%
\TextAt{8}{136.547}{367.599}{that}%
\TextAt{13}{161.875}{367.599}{f}%
\TextAt{17}{167.645}{365.739}{\Character{6}}%
\TextAt{3}{173.761}{365.739}{\Character{106}}%
\TextAt{21}{176.113}{365.739}{K}%
\TextAt{9}{187.637}{367.599}{\Character{61}}%
\TextAt{9}{200.047}{367.599}{1}%
\TextAt{8}{205.915}{367.599}{,}%
\TextAt{8}{213.032}{367.599}{and}%
\TextAt{8}{235.755}{367.599}{so}%
\TextAt{9}{250.084}{367.599}{\Character{6}}%
\TextAt{16}{261.113}{367.599}{\Character{92}}%
\TextAt{20}{274.395}{370.621}{\Character{101}}%
\TextAt{13}{271.657}{367.599}{K}%
\TextAt{9}{285.795}{367.599}{\Character{61}}%
\TextAt{13}{298.221}{367.599}{K}%
\TextAt{8}{309.038}{367.599}{.}%
\TextAt{8}{317.473}{367.599}{Let}%
\TextAt{13}{338.402}{367.599}{M}%
\TextAt{16}{354.297}{367.599}{\Character{106}}%
\TextAt{13}{360.938}{367.599}{K}%
\TextAt{8}{375.613}{367.599}{be}%
\TextAt{8}{391.497}{367.599}{a}%
\TextAt{8}{401.212}{367.599}{finite}%
\TextAt{8}{431.083}{367.599}{Galois}%
\TextAt{8}{466.956}{367.599}{subextension}%
\TextAt{8}{79.226}{350.553}{of}%
\TextAt{20}{92.852}{353.575}{\Character{101}}%
\TextAt{13}{91.865}{350.553}{L}%
\TextAt{16}{103.151}{350.553}{\Character{106}}%
\TextAt{13}{109.792}{350.553}{K}%
\TextAt{8}{123.816}{350.553}{containing}%
\TextAt{9}{180.044}{350.553}{\Character{6}}%
\TextAt{8}{191.705}{350.553}{and}%
\TextAt{13}{213.777}{350.553}{L}%
\TextAt{8}{221.741}{350.553}{,}%
\TextAt{8}{228.338}{350.553}{let}%
\TextAt{13}{244.551}{350.553}{M}%
\TextAt{17}{257.125}{354.891}{0}%
\TextAt{9}{265.178}{350.553}{\Character{61}}%
\TextAt{13}{277.603}{350.553}{M}%
\TextAt{16}{291.421}{350.553}{\Character{92}}%
\TextAt{20}{303.373}{353.575}{\Character{101}}%
\TextAt{13}{300.635}{350.553}{K}%
\TextAt{8}{314.659}{350.553}{be}%
\TextAt{8}{329.892}{350.553}{the}%
\TextAt{8}{349.358}{350.553}{maximal}%
\TextAt{8}{396.476}{350.553}{unramified}%
\TextAt{8}{454.317}{350.553}{subextension}%
\TextAt{8}{523.337}{350.553}{of}%
\TextAt{13}{79.226}{333.507}{M}%
\TextAt{16}{95.156}{333.507}{\Character{106}}%
\TextAt{13}{101.833}{333.507}{K}%
\TextAt{8}{112.651}{333.507}{,}%
\TextAt{8}{119.826}{333.507}{and}%
\TextAt{8}{142.61}{333.507}{put}%
\TextAt{13}{164.09}{333.507}{N}%
\TextAt{9}{178.069}{333.507}{\Character{61}}%
\TextAt{13}{190.53}{333.507}{N}%
\TextAt{21}{199.904}{331.503}{M}%
\TextAt{3}{208.886}{331.503}{\Character{106}}%
\TextAt{21}{211.238}{331.503}{M}%
\TextAt{24}{220.22}{333.773}{0}%
\TextAt{8}{224.869}{333.507}{.}%
\TextAt{8}{233.38}{333.507}{As}%
\TextAt{13}{250.688}{333.507}{f}%
\TextAt{17}{256.458}{331.647}{\Character{6}}%
\TextAt{3}{262.574}{331.647}{\Character{106}}%
\TextAt{21}{264.926}{331.647}{K}%
\TextAt{9}{276.485}{333.507}{\Character{61}}%
\TextAt{13}{288.946}{333.507}{f}%
\TextAt{21}{294.716}{331.647}{L}%
\TextAt{3}{300.475}{331.647}{\Character{106}}%
\TextAt{21}{302.827}{331.647}{K}%
\TextAt{9}{314.386}{333.507}{\Character{61}}%
\TextAt{9}{326.831}{333.507}{1}%
\TextAt{8}{332.7}{333.507}{,}%
\TextAt{8}{339.876}{333.507}{one}%
\TextAt{8}{361.356}{333.507}{finds}%
\TextAt{13}{389.4}{333.507}{N}%
\TextAt{16}{400.023}{333.507}{\Character{106}}%
\TextAt{21}{403.343}{331.714}{A}%
\TextAt{24}{409.686}{330.374}{\Character{6}}%
\TextAt{9}{419.242}{333.507}{\Character{61}}%
\TextAt{13}{431.703}{333.507}{N}%
\TextAt{17}{441.076}{331.647}{\Character{6}}%
\TextAt{3}{447.192}{331.647}{\Character{106}}%
\TextAt{21}{449.545}{331.647}{K}%
\TextAt{8}{457.747}{333.507}{,}%
\TextAt{13}{464.923}{333.507}{N}%
\TextAt{16}{475.545}{333.507}{\Character{106}}%
\TextAt{21}{478.866}{331.714}{A}%
\TextAt{22}{485.209}{330.374}{L}%
\TextAt{9}{494.621}{333.507}{\Character{61}}%
\TextAt{13}{507.082}{333.507}{N}%
\TextAt{21}{516.455}{331.647}{L}%
\TextAt{3}{522.214}{331.647}{\Character{106}}%
\TextAt{21}{524.567}{331.647}{K}%
\TextAt{8}{79.226}{316.461}{(see}%
\TextAt{8}{102.694}{316.461}{the}%
\TextAt{8}{122.851}{316.461}{proof}%
\TextAt{8}{153.42}{316.461}{of}%
\TextAt{8}{166.75}{316.461}{(5.3)).}%
\TextAt{8}{97.159}{296.426}{For}%
\TextAt{8}{118.131}{296.426}{the}%
\TextAt{8}{138.289}{296.426}{injectivity}%
\TextAt{8}{193.572}{296.426}{of}%
\TextAt{13}{206.905}{296.426}{r}%
\TextAt{21}{212.18}{294.567}{L}%
\TextAt{3}{217.939}{294.567}{\Character{106}}%
\TextAt{21}{220.291}{294.567}{K}%
\TextAt{8}{228.494}{296.426}{,}%
\TextAt{8}{235.646}{296.426}{we}%
\TextAt{8}{252.876}{296.426}{have}%
\TextAt{8}{279.873}{296.426}{to}%
\TextAt{8}{294.185}{296.426}{prove}%
\TextAt{8}{325.737}{296.426}{this:}%
\TextAt{8}{353.115}{296.426}{if}%
\TextAt{13}{363.841}{296.426}{r}%
\TextAt{21}{369.117}{294.567}{L}%
\TextAt{3}{374.876}{294.567}{\Character{106}}%
\TextAt{21}{377.228}{294.567}{K}%
\TextAt{9}{385.431}{296.426}{\Character{40}}%
\TextAt{13}{389.983}{296.426}{\Character{27}}%
\TextAt{21}{397.065}{300.765}{k}%
\TextAt{9}{402.185}{296.426}{\Character{41}}%
\TextAt{9}{410.052}{296.426}{\Character{61}}%
\TextAt{9}{422.473}{296.426}{1}%
\TextAt{8}{428.344}{296.426}{,}%
\TextAt{8}{435.496}{296.426}{where}%
\TextAt{9}{469.308}{296.426}{0}%
\NativeMathAt{478.486}{296.426}{11.955}{\leq}%
\TextAt{13}{491.109}{296.426}{k}%
\TextAt{13}{500.912}{296.426}{\Character{60}}%
\TextAt{13}{513.334}{296.426}{n}%
\TextAt{9}{523.665}{296.426}{\Character{61}}%
\TextAt{9}{79.226}{279.38}{\Character{91}}%
\TextAt{13}{82.478}{279.38}{L}%
\TextAt{9}{93.763}{279.38}{\Character{58}}%
\TextAt{13}{100.336}{279.38}{K}%
\TextAt{9}{111.153}{279.38}{\Character{93}}%
\TextAt{8}{114.405}{279.38}{,}%
\TextAt{8}{121.554}{279.38}{then}%
\TextAt{13}{148.214}{279.38}{k}%
\TextAt{9}{158.025}{279.38}{\Character{61}}%
\TextAt{9}{170.435}{279.38}{0}%
\TextAt{8}{176.303}{279.38}{.}%
\TextAt{8}{97.159}{259.346}{In}%
\TextAt{8}{111.358}{259.346}{order}%
\TextAt{8}{141.506}{259.346}{to}%
\TextAt{8}{155.395}{259.346}{do}%
\TextAt{8}{171.232}{259.346}{this,}%
\TextAt{8}{196.969}{259.346}{let}%
\TextAt{13}{213.452}{259.346}{\Character{25}}%
\TextAt{17}{220.093}{257.552}{\Character{6}}%
\NativeMathAt{230.028}{259.346}{11.955}{\in}%
\TextAt{13}{241.319}{259.346}{A}%
\TextAt{17}{250.094}{257.552}{\Character{6}}%
\TextAt{8}{256.708}{259.346}{,}%
\TextAt{13}{263.52}{259.346}{\Character{25}}%
\TextAt{21}{270.16}{257.552}{L}%
\NativeMathAt{279.738}{259.346}{11.955}{\in}%
\TextAt{13}{291.029}{259.346}{A}%
\TextAt{21}{299.805}{257.552}{L}%
\TextAt{8}{309.538}{259.346}{be}%
\TextAt{8}{325.04}{259.346}{prime}%
\TextAt{8}{357.783}{259.346}{elements.}%
\TextAt{8}{410.049}{259.346}{Since}%
\TextAt{9}{440.185}{259.346}{\Character{6}}%
\TextAt{13}{448.64}{259.346}{\Character{59}}%
\TextAt{13}{453.864}{259.346}{L}%
\NativeMathAt{465.169}{259.346}{11.955}{\subseteq}%
\TextAt{13}{477.788}{259.346}{M}%
\TextAt{8}{490.362}{259.346}{,}%
\TextAt{13}{497.174}{259.346}{\Character{25}}%
\TextAt{17}{503.814}{257.552}{\Character{6}}%
\TextAt{8}{513.904}{259.346}{and}%
\TextAt{13}{79.226}{242.3}{\Character{25}}%
\TextAt{21}{85.866}{240.506}{L}%
\TextAt{8}{96.021}{242.3}{are}%
\TextAt{8}{115.532}{242.3}{both}%
\TextAt{8}{143.172}{242.3}{prime}%
\TextAt{8}{176.336}{242.3}{elements}%
\TextAt{8}{224.192}{242.3}{of}%
\TextAt{13}{237.522}{242.3}{M}%
\TextAt{8}{250.096}{242.3}{.}%
\TextAt{8}{258.544}{242.3}{Putting}%
\TextAt{13}{301.631}{242.3}{\Character{25}}%
\TextAt{21}{308.7}{246.638}{k}%
\TextAt{17}{308.271}{239.279}{\Character{6}}%
\TextAt{9}{318.206}{242.3}{\Character{61}}%
\TextAt{13}{330.632}{242.3}{u\Character{25}}%
\TextAt{21}{344.363}{246.638}{k}%
\TextAt{21}{343.934}{239.279}{L}%
\TextAt{8}{350.192}{242.3}{,}%
\TextAt{13}{357.341}{242.3}{u}%
\NativeMathAt{367.324}{242.3}{11.955}{\in}%
\TextAt{13}{378.615}{242.3}{U}%
\TextAt{21}{386.565}{240.506}{M}%
\TextAt{8}{396.045}{242.3}{,}%
\TextAt{8}{403.194}{242.3}{we}%
\TextAt{8}{420.422}{242.3}{obtain}%
\TextAt{13}{161.055}{213.298}{r}%
\TextAt{21}{166.33}{211.439}{L}%
\TextAt{3}{172.089}{211.439}{\Character{106}}%
\TextAt{21}{174.442}{211.439}{K}%
\TextAt{9}{182.644}{213.298}{\Character{40}}%
\TextAt{13}{187.197}{213.299}{\Character{27}}%
\TextAt{21}{194.279}{218.235}{k}%
\TextAt{9}{199.399}{213.298}{\Character{41}}%
\TextAt{9}{207.254}{213.298}{\Character{61}}%
\TextAt{13}{219.697}{213.299}{N}%
\TextAt{9}{230.32}{213.298}{\Character{40}}%
\TextAt{13}{234.872}{213.299}{\Character{25}}%
\TextAt{21}{241.941}{218.235}{k}%
\TextAt{17}{241.513}{210.343}{\Character{6}}%
\TextAt{9}{248.127}{213.299}{\Character{41}}%
\TextAt{9}{255.982}{213.299}{\Character{61}}%
\TextAt{13}{268.425}{213.298}{N}%
\TextAt{9}{279.048}{213.299}{\Character{40}}%
\TextAt{13}{283.6}{213.298}{u}%
\TextAt{9}{290.263}{213.299}{\Character{41}}%
\TextAt{16}{297.472}{213.298}{\Character{1}}%
\TextAt{13}{303.449}{213.299}{N}%
\TextAt{9}{314.072}{213.298}{\Character{40}}%
\TextAt{13}{318.624}{213.299}{\Character{25}}%
\TextAt{21}{325.693}{218.235}{k}%
\TextAt{21}{325.265}{210.343}{L}%
\TextAt{9}{331.522}{213.299}{\Character{41}}%
\TextAt{9}{339.377}{213.299}{\Character{61}}%
\TextAt{13}{351.82}{213.298}{N}%
\TextAt{9}{362.443}{213.299}{\Character{40}}%
\TextAt{13}{366.995}{213.298}{u}%
\TextAt{9}{373.658}{213.299}{\Character{41}}%
\TextAt{8}{381.531}{213.298}{mod}%
\TextAt{13}{406.969}{213.299}{N}%
\TextAt{21}{416.343}{211.439}{L}%
\TextAt{3}{422.102}{211.439}{\Character{106}}%
\TextAt{21}{424.454}{211.439}{K}%
\TextAt{13}{432.657}{213.298}{A}%
\TextAt{21}{441.432}{211.505}{L}%
\TextAt{13}{447.689}{213.299}{\Character{58}}%
\TextAt{8}{79.226}{184.297}{From}%
\TextAt{13}{109.333}{184.297}{r}%
\TextAt{21}{114.609}{182.438}{L}%
\TextAt{3}{120.368}{182.438}{\Character{106}}%
\TextAt{21}{122.72}{182.438}{K}%
\TextAt{9}{130.923}{184.297}{\Character{40}}%
\TextAt{13}{135.475}{184.297}{\Character{27}}%
\TextAt{21}{142.557}{188.636}{k}%
\TextAt{9}{147.677}{184.297}{\Character{41}}%
\TextAt{9}{155.532}{184.297}{\Character{61}}%
\TextAt{9}{167.953}{184.297}{1}%
\TextAt{8}{173.829}{184.297}{,}%
\TextAt{8}{180.484}{184.297}{it}%
\TextAt{8}{191.57}{184.297}{thus}%
\TextAt{8}{216.704}{184.297}{follows}%
\TextAt{8}{254.523}{184.297}{that}%
\TextAt{13}{279.274}{184.297}{N}%
\TextAt{9}{289.896}{184.297}{\Character{40}}%
\TextAt{13}{294.449}{184.297}{u}%
\TextAt{9}{301.111}{184.297}{\Character{41}}%
\TextAt{9}{308.966}{184.297}{\Character{61}}%
\TextAt{13}{321.41}{184.297}{N}%
\TextAt{9}{332.032}{184.297}{\Character{40}}%
\TextAt{13}{336.585}{184.297}{v}%
\TextAt{9}{342.673}{184.297}{\Character{41}}%
\TextAt{8}{350.505}{184.297}{for}%
\TextAt{8}{367.772}{184.297}{some}%
\TextAt{13}{396.481}{184.297}{v}%
\NativeMathAt{405.89}{184.297}{11.955}{\in}%
\TextAt{13}{417.181}{184.297}{U}%
\TextAt{21}{425.131}{182.504}{L}%
\TextAt{8}{431.388}{184.297}{,}%
\TextAt{8}{438.044}{184.297}{so}%
\TextAt{8}{451.796}{184.297}{that}%
\TextAt{13}{476.547}{184.297}{N}%
\TextAt{9}{487.17}{184.297}{\Character{40}}%
\TextAt{13}{491.722}{184.297}{u}%
\TextAt{3}{498.385}{188.636}{\Character{0}}%
\TextAt{17}{504.971}{188.636}{1}%
\TextAt{13}{509.703}{184.297}{v}%
\TextAt{9}{515.792}{184.297}{\Character{41}}%
\TextAt{9}{523.647}{184.297}{\Character{61}}%
\TextAt{9}{79.226}{167.251}{1}%
\TextAt{8}{85.079}{167.251}{.}%
\TextAt{8}{95.017}{167.251}{From}%
\TextAt{8}{126.238}{167.251}{axiom}%
\TextAt{8}{161.536}{167.251}{(6.1),}%
\TextAt{8}{193.383}{167.251}{we}%
\TextAt{8}{211.107}{167.251}{may}%
\TextAt{8}{236.972}{167.251}{write}%
\TextAt{13}{267.381}{167.251}{u}%
\TextAt{3}{274.043}{171.59}{\Character{0}}%
\TextAt{17}{280.63}{171.59}{1}%
\TextAt{13}{285.362}{167.251}{v}%
\TextAt{9}{295.617}{167.251}{\Character{61}}%
\TextAt{13}{308.888}{167.251}{a}%
\TextAt{21}{315.033}{171.59}{\Character{27}}%
\TextAt{3}{320.159}{171.59}{\Character{0}}%
\TextAt{17}{326.745}{171.59}{1}%
\TextAt{8}{335.871}{167.251}{for}%
\TextAt{8}{354.253}{167.251}{some}%
\TextAt{13}{384.075}{167.251}{a}%
\NativeMathAt{394.387}{167.251}{11.955}{\in}%
\TextAt{13}{406.524}{167.251}{A}%
\TextAt{21}{415.299}{165.458}{M}%
\TextAt{8}{424.779}{167.251}{,}%
\TextAt{8}{432.549}{167.251}{and}%
\TextAt{8}{455.808}{167.251}{find}%
\TextAt{8}{479.712}{167.251}{in}%
\TextAt{13}{493.862}{167.251}{A}%
\TextAt{21}{502.637}{165.458}{M}%
\TextAt{8}{516.51}{167.251}{the}%
\TextAt{8}{79.226}{150.205}{equation}%
\TextAt{9}{148.771}{133.159}{\Character{40}}%
\TextAt{13}{153.323}{133.159}{\Character{25}}%
\TextAt{21}{160.392}{138.096}{k}%
\TextAt{21}{159.963}{130.204}{L}%
\TextAt{13}{166.221}{133.159}{v}%
\TextAt{9}{172.309}{133.159}{\Character{41}}%
\TextAt{21}{176.861}{138.096}{\Character{27}}%
\TextAt{3}{181.986}{138.096}{\Character{0}}%
\TextAt{17}{188.573}{138.096}{1}%
\TextAt{9}{196.626}{133.159}{\Character{61}}%
\TextAt{9}{209.036}{133.159}{\Character{40}}%
\TextAt{13}{213.604}{133.159}{\Character{25}}%
\TextAt{21}{220.673}{138.096}{k}%
\TextAt{21}{220.244}{130.204}{L}%
\TextAt{13}{226.501}{133.159}{v}%
\TextAt{9}{232.59}{133.159}{\Character{41}}%
\TextAt{30}{237.352}{138.096}{\Character{101}}%
\TextAt{21}{237.142}{138.096}{\Character{27}}%
\TextAt{3}{242.267}{138.096}{\Character{0}}%
\TextAt{17}{248.854}{138.096}{1}%
\TextAt{9}{256.907}{133.159}{\Character{61}}%
\TextAt{9}{269.317}{133.159}{\Character{40}}%
\TextAt{13}{273.885}{133.159}{\Character{25}}%
\TextAt{21}{280.954}{138.096}{k}%
\TextAt{17}{280.525}{130.204}{\Character{6}}%
\TextAt{13}{287.139}{133.159}{u}%
\TextAt{3}{293.802}{138.096}{\Character{0}}%
\TextAt{17}{300.388}{138.096}{1}%
\TextAt{13}{305.12}{133.159}{v}%
\TextAt{9}{311.209}{133.159}{\Character{41}}%
\TextAt{30}{315.971}{138.096}{\Character{101}}%
\TextAt{21}{315.761}{138.096}{\Character{27}}%
\TextAt{3}{320.886}{138.096}{\Character{0}}%
\TextAt{17}{327.473}{138.096}{1}%
\TextAt{9}{335.526}{133.159}{\Character{61}}%
\TextAt{9}{347.935}{133.159}{\Character{40}}%
\TextAt{13}{352.504}{133.159}{a}%
\TextAt{21}{358.649}{138.096}{\Character{27}}%
\TextAt{3}{363.774}{138.096}{\Character{0}}%
\TextAt{17}{370.36}{138.096}{1}%
\TextAt{9}{375.093}{133.159}{\Character{41}}%
\TextAt{30}{379.855}{138.096}{\Character{101}}%
\TextAt{21}{379.645}{138.096}{\Character{27}}%
\TextAt{3}{384.77}{138.096}{\Character{0}}%
\TextAt{17}{391.357}{138.096}{1}%
\TextAt{9}{399.41}{133.159}{\Character{61}}%
\TextAt{9}{411.82}{133.159}{\Character{40}}%
\TextAt{13}{416.388}{133.159}{a}%
\TextAt{30}{422.743}{138.096}{\Character{101}}%
\TextAt{21}{422.533}{138.096}{\Character{27}}%
\TextAt{3}{427.658}{138.096}{\Character{0}}%
\TextAt{17}{434.244}{138.096}{1}%
\TextAt{9}{438.977}{133.159}{\Character{41}}%
\TextAt{21}{443.529}{138.096}{\Character{27}}%
\TextAt{3}{448.654}{138.096}{\Character{0}}%
\TextAt{17}{455.241}{138.096}{1}%
\TextAt{13}{459.973}{133.159}{\Character{59}}%
\TextAt{8}{79.226}{109.638}{and}%
\TextAt{8}{102.183}{109.638}{so}%
\TextAt{13}{116.747}{109.638}{x}%
\TextAt{9}{127.05}{109.638}{\Character{61}}%
\TextAt{13}{139.807}{109.638}{\Character{25}}%
\TextAt{21}{146.876}{113.976}{k}%
\TextAt{21}{146.447}{106.617}{L}%
\TextAt{13}{152.704}{109.638}{v}%
\TextAt{13}{158.777}{109.638}{a}%
\TextAt{17}{164.937}{113.976}{1}%
\TextAt{3}{169.171}{113.976}{\Character{0}}%
\TextAt{30}{175.968}{113.976}{\Character{101}}%
\TextAt{21}{175.758}{113.976}{\Character{27}}%
\NativeMathAt{185.033}{109.638}{11.955}{\in}%
\TextAt{13}{196.654}{109.638}{A}%
\TextAt{21}{205.43}{107.634}{M}%
\TextAt{24}{214.411}{109.903}{0}%
\TextAt{8}{219.061}{109.638}{.}%
\TextAt{8}{228.091}{109.638}{Now}%
\TextAt{13}{254.945}{109.638}{v}%
\TextAt{21}{260.604}{107.634}{M}%
\TextAt{24}{269.586}{109.903}{0}%
\TextAt{9}{274.235}{109.638}{\Character{40}}%
\TextAt{13}{278.788}{109.638}{x}%
\TextAt{9}{285.44}{109.638}{\Character{41}}%
\NativeMathAt{293.643}{109.638}{11.955}{\in}%
\TextAt{20}{305.929}{112.726}{\Character{98}}%
\TextAt{18}{305.265}{109.638}{Z}%
\TextAt{8}{317.327}{109.638}{and}%
\TextAt{13}{340.283}{109.638}{nv}%
\TextAt{21}{352.93}{107.634}{M}%
\TextAt{24}{361.912}{109.903}{0}%
\TextAt{9}{366.561}{109.638}{\Character{40}}%
\TextAt{13}{371.113}{109.638}{x}%
\TextAt{9}{377.765}{109.638}{\Character{41}}%
\TextAt{9}{385.954}{109.638}{\Character{61}}%
\TextAt{13}{398.725}{109.638}{v}%
\TextAt{21}{404.385}{107.844}{M}%
\TextAt{9}{413.864}{109.638}{\Character{40}}%
\TextAt{13}{418.417}{109.638}{x}%
\TextAt{9}{425.069}{109.638}{\Character{41}}%
\TextAt{9}{433.258}{109.638}{\Character{61}}%
\TextAt{13}{446.029}{109.638}{k}%
\TextAt{8}{456.61}{109.638}{imply}%
\TextAt{8}{489.645}{109.638}{that}%
\TextAt{8}{515.207}{109.638}{one}%
\TextAt{8}{79.226}{92.592}{has}%
\TextAt{13}{100.1}{92.592}{k}%
\TextAt{9}{109.91}{92.592}{\Character{61}}%
\TextAt{9}{122.32}{92.592}{0}%
\TextAt{8}{128.189}{92.592}{,}%
\TextAt{8}{135.338}{92.592}{and}%
\TextAt{8}{158.101}{92.592}{so}%
\TextAt{13}{172.471}{92.592}{r}%
\TextAt{21}{177.746}{90.732}{L}%
\TextAt{3}{183.505}{90.732}{\Character{106}}%
\TextAt{21}{185.857}{90.732}{K}%
\TextAt{8}{197.957}{92.592}{is}%
\TextAt{8}{209.721}{92.592}{injective.}%
\TextAt{8}{260.766}{92.592}{The}%
\TextAt{8}{284.82}{92.592}{surjectivity}%
\TextAt{8}{346.018}{92.592}{then}%
\TextAt{8}{372.678}{92.592}{follows}%
\TextAt{8}{411.114}{92.592}{from}%
\TextAt{8}{438.754}{92.592}{(6.1).}%
\Vector{q 0 G 0 g 0.399 w 452.334 -681.408 m 452.334 -673.338 l S Q}%
\Vector{q 0 G 0 g 0.399 w 452.534 -673.538 m 459.707 -673.538 l S Q}%
\Vector{q 0 G 0 g 0.399 w 452.534 -681.209 m 459.707 -681.209 l S Q}%
\Vector{q 0 G 0 g 0.399 w 459.906 -681.408 m 459.906 -673.338 l S Q}%
\TextAt{8}{297.211}{43.033}{276}%
\Destination{78.23}{809.815}{nav.page.276}%
\Destination{79.36}{274.355}{nav.print.302}%
\Hotlink{194.663}{487.755}{25.078}{11.047}{nav.number.4.5.7}%
\Hotlink{166.25}{314.255}{25.078}{11.047}{nav.number.4.5.3}%
\Destination{79.226}{778.124}{errata.101}%
\Hotlink{439.912}{89.602}{21.722}{11.047}{nav.number.4.6.1}%
}

% Page 289 (body page 277)
\Page{%
\TextAt{2}{79.226}{796.943}{IV.}%
\TextAt{2}{98.012}{796.943}{Abstract}%
\TextAt{2}{143.252}{796.943}{Class}%
\TextAt{2}{171.844}{796.943}{Field}%
\TextAt{2}{199.575}{796.943}{Theory}%
\Vector{q 0 G 0 g 0.299 w 7.226 18.729 m 460.769 18.729 l S Q}%
\TextAt{8}{97.159}{761.124}{The}%
\TextAt{8}{121.747}{761.124}{inverse}%
\TextAt{8}{161.041}{761.124}{of}%
\TextAt{8}{174.905}{761.124}{the}%
\TextAt{8}{195.596}{761.124}{mapping}%
\TextAt{13}{244.263}{761.124}{r}%
\TextAt{21}{249.538}{759.264}{L}%
\TextAt{3}{255.297}{759.264}{\Character{106}}%
\TextAt{21}{257.649}{759.264}{K}%
\TextAt{9}{270.084}{761.124}{\Character{58}}%
\TextAt{13}{277.568}{761.124}{G}%
\TextAt{9}{286.801}{761.124}{\Character{40}}%
\TextAt{13}{291.354}{761.124}{L}%
\TextAt{16}{303.55}{761.124}{\Character{106}}%
\TextAt{13}{311.103}{761.124}{K}%
\TextAt{9}{321.921}{761.124}{\Character{41}}%
\TextAt{7}{326.473}{765.462}{ab}%
\NativeMathAt{340.137}{761.124}{11.955}{\rightarrow}%
\TextAt{13}{356.325}{761.124}{A}%
\TextAt{21}{365.1}{759.33}{K}%
\TextAt{9}{373.303}{761.124}{\Character{47}}%
\TextAt{13}{379.156}{761.124}{N}%
\TextAt{21}{388.529}{759.264}{L}%
\TextAt{3}{394.288}{759.264}{\Character{106}}%
\TextAt{21}{396.641}{759.264}{K}%
\TextAt{13}{404.843}{761.124}{A}%
\TextAt{21}{413.619}{759.33}{L}%
\TextAt{8}{424.308}{761.124}{gives,}%
\TextAt{8}{456.909}{761.124}{for}%
\TextAt{8}{475.328}{761.124}{every}%
\TextAt{8}{506.755}{761.124}{finite}%
\TextAt{8}{79.226}{744.078}{Galois}%
\TextAt{8}{115.139}{744.078}{extension}%
\TextAt{13}{166.905}{744.078}{L}%
\TextAt{16}{178.191}{744.078}{\Character{106}}%
\TextAt{13}{184.832}{744.078}{K}%
\TextAt{8}{195.65}{744.078}{,}%
\TextAt{8}{202.799}{744.078}{a}%
\TextAt{8}{212.555}{744.078}{surjective}%
\TextAt{8}{264.966}{744.078}{homomorphism}%
\TextAt{9}{227.857}{718.945}{\Character{40}}%
\TextAt{13}{236.307}{718.945}{\Character{59}}%
\TextAt{13}{241.531}{718.945}{L}%
\TextAt{16}{252.837}{718.945}{\Character{106}}%
\TextAt{13}{259.478}{718.945}{K}%
\TextAt{9}{270.296}{718.945}{\Character{41}}%
\TextAt{9}{278.151}{718.945}{\Character{58}}%
\TextAt{13}{284.742}{718.945}{A}%
\TextAt{21}{293.517}{717.151}{K}%
\NativeArrowAt{305.041}{718.945}{11.955}{19.26}{\longrightarrow}%
\TextAt{13}{327.623}{718.945}{G}%
\TextAt{9}{336.856}{718.945}{\Character{40}}%
\TextAt{13}{341.409}{718.945}{L}%
\TextAt{16}{352.694}{718.945}{\Character{106}}%
\TextAt{13}{359.336}{718.945}{K}%
\TextAt{9}{370.153}{718.945}{\Character{41}}%
\TextAt{7}{374.706}{723.881}{ab}%
\TextAt{8}{79.226}{693.812}{with}%
\TextAt{8}{105.346}{693.812}{kernel}%
\TextAt{13}{139.273}{693.812}{N}%
\TextAt{21}{148.647}{691.952}{L}%
\TextAt{3}{154.406}{691.952}{\Character{106}}%
\TextAt{21}{156.758}{691.952}{K}%
\TextAt{13}{164.961}{693.812}{A}%
\TextAt{21}{173.736}{692.018}{L}%
\TextAt{8}{179.993}{693.812}{.}%
\TextAt{8}{188.261}{693.812}{This}%
\TextAt{8}{214.441}{693.812}{map}%
\TextAt{8}{239.916}{693.812}{is}%
\TextAt{8}{251.139}{693.812}{called}%
\TextAt{8}{283.763}{693.812}{the}%
\TextAt{12}{303.379}{693.812}{norm}%
\TextAt{12}{338.157}{693.812}{residue}%
\TextAt{12}{383.779}{693.812}{symbol}%
\TextAt{8}{428.704}{693.812}{of}%
\TextAt{13}{441.494}{693.812}{L}%
\TextAt{16}{452.779}{693.812}{\Character{106}}%
\TextAt{13}{459.421}{693.812}{K}%
\TextAt{8}{470.239}{693.812}{.}%
\TextAt{8}{478.507}{693.812}{From}%
\TextAt{8}{508.692}{693.812}{(5.8)}%
\TextAt{8}{79.226}{676.766}{and}%
\TextAt{8}{101.989}{676.766}{(5.9)}%
\TextAt{8}{129.964}{676.766}{we}%
\TextAt{8}{147.191}{676.766}{have}%
\TextAt{8}{174.186}{676.766}{the}%
\TextAt{12}{79.226}{651.858}{Proposition}%
\TextAt{12}{151.971}{651.858}{(6.4).}%
\TextAt{8}{187.315}{651.858}{Let}%
\TextAt{13}{208.313}{651.858}{L}%
\TextAt{16}{219.647}{651.858}{\Character{106}}%
\TextAt{13}{226.338}{651.858}{K}%
\TextAt{8}{241.082}{651.858}{and}%
\TextAt{13}{263.873}{651.858}{L}%
\TextAt{3}{271.838}{656.197}{\Character{48}}%
\TextAt{16}{278.003}{651.858}{\Character{106}}%
\TextAt{13}{284.693}{651.858}{K}%
\TextAt{3}{295.511}{656.197}{\Character{48}}%
\TextAt{8}{302.232}{651.858}{be}%
\TextAt{8}{318.185}{651.858}{finite}%
\TextAt{8}{348.126}{651.858}{Galois}%
\TextAt{8}{384.068}{651.858}{extensions}%
\TextAt{8}{440.477}{651.858}{such}%
\TextAt{8}{466.903}{651.858}{that}%
\TextAt{13}{492.301}{651.858}{K}%
\NativeMathAt{506.488}{651.858}{11.955}{\subseteq}%
\TextAt{13}{519.157}{651.858}{K}%
\TextAt{3}{529.974}{656.197}{\Character{48}}%
\TextAt{8}{79.226}{634.813}{and}%
\TextAt{13}{101.989}{634.813}{L}%
\NativeMathAt{113.274}{634.813}{11.955}{\subseteq}%
\TextAt{13}{125.893}{634.813}{L}%
\TextAt{3}{133.858}{639.151}{\Character{48}}%
\TextAt{8}{136.653}{634.813}{,}%
\TextAt{8}{143.802}{634.813}{and}%
\TextAt{8}{166.565}{634.813}{let}%
\TextAt{13}{183.469}{634.813}{\Character{27}}%
\NativeMathAt{193.873}{634.813}{11.955}{\in}%
\TextAt{13}{205.164}{634.813}{G}%
\TextAt{8}{214.397}{634.813}{.}%
\TextAt{8}{222.846}{634.813}{Then}%
\TextAt{8}{253.403}{634.813}{we}%
\TextAt{8}{270.63}{634.813}{have}%
\TextAt{8}{297.625}{634.813}{the}%
\TextAt{8}{317.782}{634.813}{commutative}%
\TextAt{8}{387.707}{634.813}{diagrams}%
\TextAt{13}{171.975}{606.772}{A}%
\TextAt{21}{180.751}{604.979}{K}%
\TextAt{23}{188.455}{607.248}{\Character{48}}%
\TextAt{13}{230.635}{606.772}{G}%
\TextAt{9}{239.869}{606.772}{\Character{40}}%
\TextAt{13}{244.421}{606.772}{L}%
\TextAt{3}{252.385}{611.11}{\Character{48}}%
\TextAt{16}{258.501}{606.772}{\Character{106}}%
\TextAt{13}{265.143}{606.772}{K}%
\TextAt{3}{275.961}{611.11}{\Character{48}}%
\TextAt{9}{278.756}{606.772}{\Character{41}}%
\TextAt{7}{283.308}{611.11}{ab}%
\TextAt{13}{173.317}{563.634}{A}%
\TextAt{21}{182.093}{561.841}{K}%
\TextAt{13}{233.43}{563.634}{G}%
\TextAt{9}{242.664}{563.634}{\Character{40}}%
\TextAt{13}{247.216}{563.634}{L}%
\TextAt{16}{258.501}{563.634}{\Character{106}}%
\TextAt{13}{265.143}{563.634}{K}%
\TextAt{9}{275.961}{563.634}{\Character{41}}%
\TextAt{7}{280.513}{567.973}{ab}%
\NativeDiagram{%
\path[draw=black,line width=0.478bp] (124.989,-164.24) -- (152.805,-164.24);%
\path[draw=black,line width=0.478bp,line cap=round,line join=round] (150.557,-161.371) .. controls (151.012,-163.092) and (152.024,-163.905) .. (153.044,-164.24) .. controls (152.024,-164.574) and (151.012,-165.387) .. (150.557,-167.109);%
}%
\TextAt{17}{194.125}{614.568}{\Character{40}}%
\TextAt{21}{201.316}{614.568}{\Character{59}L}%
\TextAt{23}{209.427}{617.381}{\Character{48}}%
\TextAt{3}{212.112}{614.568}{\Character{106}}%
\TextAt{21}{214.464}{614.568}{K}%
\TextAt{23}{222.169}{617.381}{\Character{48}}%
\TextAt{17}{224.853}{614.568}{\Character{41}}%
\NativeDiagram{%
\path[draw=black,line width=0.478bp] (109.806,-173.601) -- (109.806,-197.14);%
\path[draw=black,line width=0.478bp,line cap=round,line join=round] (112.675,-194.892) .. controls (110.953,-195.347) and (110.14,-196.359) .. (109.806,-197.379) .. controls (109.471,-196.359) and (108.658,-195.347) .. (106.936,-194.892);%
}%
\TextAt{21}{153.727}{587.56}{N}%
\TextAt{22}{160.449}{585.269}{K}%
\TextAt{23}{166.994}{586.985}{\Character{48}}%
\TextAt{23}{169.678}{585.269}{\Character{106}}%
\TextAt{22}{171.947}{585.269}{K}%
\NativeDiagram{%
\path[draw=black,line width=0.478bp] (189.688,-174.796) -- (189.688,-195.439);%
\path[draw=black,line width=0.478bp,line cap=round,line join=round] (192.558,-193.192) .. controls (190.836,-193.646) and (190.023,-194.659) .. (189.688,-195.678) .. controls (189.354,-194.659) and (188.541,-193.646) .. (186.819,-193.192);%
\path[draw=black,line width=0.478bp] (123.647,-207.378) -- (155.6,-207.378);%
\path[draw=black,line width=0.478bp,line cap=round,line join=round] (153.352,-204.509) .. controls (153.807,-206.23) and (154.819,-207.043) .. (155.839,-207.378) .. controls (154.819,-207.712) and (153.807,-208.525) .. (153.352,-210.247);%
}%
\TextAt{17}{197.536}{557.829}{\Character{40}}%
\TextAt{21}{204.727}{557.829}{\Character{59}L}%
\TextAt{3}{212.838}{557.829}{\Character{106}}%
\TextAt{21}{215.191}{557.829}{K}%
\TextAt{17}{222.895}{557.829}{\Character{41}}%
\TextAt{13}{329.454}{606.772}{A}%
\TextAt{21}{338.229}{604.979}{K}%
\TextAt{13}{393.55}{606.772}{G}%
\TextAt{9}{402.784}{606.772}{\Character{40}}%
\TextAt{13}{407.336}{606.772}{L}%
\TextAt{16}{418.622}{606.772}{\Character{106}}%
\TextAt{13}{425.263}{606.772}{K}%
\TextAt{9}{436.081}{606.772}{\Character{41}}%
\TextAt{7}{440.633}{611.11}{ab}%
\TextAt{13}{326.956}{563.634}{A}%
\TextAt{21}{335.731}{561.841}{K}%
\TextAt{22}{343.435}{564.11}{\Character{27}}%
\TextAt{13}{387.926}{563.634}{G}%
\TextAt{9}{397.16}{563.634}{\Character{40}}%
\TextAt{13}{401.712}{563.634}{L}%
\TextAt{21}{409.677}{567.973}{\Character{27}}%
\TextAt{16}{418.621}{563.634}{\Character{106}}%
\TextAt{13}{425.263}{563.634}{K}%
\TextAt{21}{436.081}{567.973}{\Character{27}}%
\TextAt{9}{441.704}{563.634}{\Character{41}}%
\TextAt{7}{446.256}{567.973}{ab}%
\NativeDiagram{%
\path[draw=black,line width=0.478bp] (279.784,-164.24) -- (315.722,-164.24);%
\path[draw=black,line width=0.478bp,line cap=round,line join=round] (313.474,-161.371) .. controls (313.929,-163.092) and (314.941,-163.905) .. (315.961,-164.24) .. controls (314.941,-164.574) and (313.929,-165.387) .. (313.474,-167.109);%
}%
\TextAt{17}{355.666}{614.568}{\Character{40}}%
\TextAt{21}{362.857}{614.568}{\Character{59}L}%
\TextAt{3}{370.968}{614.568}{\Character{106}}%
\TextAt{21}{373.32}{614.568}{K}%
\TextAt{17}{381.025}{614.568}{\Character{41}}%
\NativeDiagram{%
\path[draw=black,line width=0.478bp] (265.943,-173.601) -- (265.943,-197.14);%
\path[draw=black,line width=0.478bp,line cap=round,line join=round] (268.812,-194.892) .. controls (267.091,-195.347) and (266.278,-196.359) .. (265.943,-197.379) .. controls (265.608,-196.359) and (264.796,-195.347) .. (263.074,-194.892);%
}%
\TextAt{21}{330.002}{586.675}{\Character{27}}%
\NativeDiagram{%
\path[draw=black,line width=0.478bp] (349.81,-174.796) -- (349.81,-195.439);%
\path[draw=black,line width=0.478bp,line cap=round,line join=round] (352.679,-193.192) .. controls (350.958,-193.646) and (350.145,-194.659) .. (349.81,-195.678) .. controls (349.475,-194.659) and (348.663,-193.646) .. (346.941,-193.192);%
}%
\TextAt{21}{424.626}{585.846}{\Character{27}}%
\TextAt{23}{429.751}{588.659}{\Character{3}}%
\NativeDiagram{%
\path[draw=black,line width=0.478bp] (282.282,-207.378) -- (310.098,-207.378);%
\path[draw=black,line width=0.478bp,line cap=round,line join=round] (307.851,-204.509) .. controls (308.305,-206.23) and (309.317,-207.043) .. (310.337,-207.378) .. controls (309.317,-207.712) and (308.305,-208.525) .. (307.851,-210.247);%
}%
\TextAt{17}{349.107}{557.829}{\Character{40}}%
\TextAt{21}{356.298}{557.829}{\Character{59}L}%
\TextAt{22}{364.409}{560.642}{\Character{27}}%
\TextAt{3}{369.405}{557.829}{\Character{106}}%
\TextAt{21}{371.757}{557.829}{K}%
\TextAt{22}{379.462}{560.642}{\Character{27}}%
\TextAt{17}{384.458}{557.829}{\Character{41}}%
\TextAt{8}{79.226}{535.411}{and}%
\TextAt{8}{101.989}{535.411}{if}%
\TextAt{13}{112.713}{535.411}{K}%
\TextAt{3}{123.53}{539.75}{\Character{48}}%
\NativeMathAt{129.646}{535.411}{11.955}{\subseteq}%
\TextAt{13}{142.266}{535.411}{L}%
\TextAt{8}{150.23}{535.411}{,}%
\TextAt{8}{157.379}{535.411}{we}%
\TextAt{8}{174.607}{535.411}{have}%
\TextAt{8}{201.601}{535.411}{the}%
\TextAt{8}{221.758}{535.411}{commutative}%
\TextAt{8}{291.684}{535.411}{diagram}%
\TextAt{13}{246.784}{507.371}{A}%
\TextAt{21}{255.559}{505.577}{K}%
\TextAt{23}{263.264}{507.847}{\Character{48}}%
\TextAt{13}{305.671}{507.371}{G}%
\TextAt{9}{314.905}{507.371}{\Character{40}}%
\TextAt{13}{319.457}{507.371}{L}%
\TextAt{16}{330.742}{507.371}{\Character{106}}%
\TextAt{13}{337.384}{507.371}{K}%
\TextAt{3}{348.202}{511.709}{\Character{48}}%
\TextAt{9}{350.997}{507.371}{\Character{41}}%
\TextAt{7}{355.549}{511.709}{ab}%
\TextAt{13}{248.126}{464.233}{A}%
\TextAt{21}{256.901}{462.44}{K}%
\TextAt{13}{305.443}{464.233}{G}%
\TextAt{9}{314.677}{464.233}{\Character{40}}%
\TextAt{13}{319.229}{464.233}{L}%
\TextAt{16}{330.514}{464.233}{\Character{106}}%
\TextAt{13}{337.156}{464.233}{K}%
\TextAt{9}{347.974}{464.233}{\Character{41}}%
\TextAt{7}{352.526}{468.572}{ab}%
\TextAt{13}{361.959}{464.233}{\Character{58}}%
\NativeDiagram{%
\path[draw=black,line width=0.478bp] (199.798,-263.641) -- (227.842,-263.641);%
\path[draw=black,line width=0.478bp,line cap=round,line join=round] (225.595,-260.772) .. controls (226.049,-262.493) and (227.062,-263.306) .. (228.082,-263.641) .. controls (227.062,-263.975) and (226.049,-264.788) .. (225.595,-266.51);%
}%
\TextAt{17}{270.391}{515.167}{\Character{40}}%
\TextAt{21}{277.581}{515.167}{\Character{59}L}%
\TextAt{3}{285.693}{515.167}{\Character{106}}%
\TextAt{21}{288.045}{515.167}{K}%
\TextAt{23}{295.75}{517.98}{\Character{48}}%
\TextAt{17}{298.434}{515.167}{\Character{41}}%
\NativeDiagram{%
\path[draw=black,line width=0.478bp] (184.615,-297.019) -- (184.615,-273.48);%
\path[draw=black,line width=0.478bp,line cap=round,line join=round] (181.746,-275.728) .. controls (183.467,-275.273) and (184.28,-274.261) .. (184.615,-273.241) .. controls (184.95,-274.261) and (185.763,-275.273) .. (187.484,-275.728);%
\path[draw=black,line width=0.478bp] (198.456,-306.779) -- (227.614,-306.779);%
\path[draw=black,line width=0.478bp,line cap=round,line join=round] (225.367,-303.909) .. controls (225.821,-305.631) and (226.833,-306.444) .. (227.853,-306.779) .. controls (226.833,-307.113) and (225.821,-307.926) .. (225.367,-309.648);%
}%
\TextAt{17}{270.948}{458.428}{\Character{40}}%
\TextAt{21}{278.138}{458.428}{\Character{59}L}%
\TextAt{3}{286.25}{458.428}{\Character{106}}%
\TextAt{21}{288.602}{458.428}{K}%
\TextAt{17}{296.307}{458.428}{\Character{41}}%
\NativeDiagram{%
\path[draw=black,line width=0.478bp] (263.328,-295.319) -- (263.328,-274.676);%
\path[draw=black,line width=0.478bp,line cap=round,line join=round] (260.459,-276.923) .. controls (262.181,-276.469) and (262.994,-275.456) .. (263.328,-274.436) .. controls (263.663,-275.456) and (264.476,-276.469) .. (266.198,-276.923);%
}%
\TextAt{7}{338.144}{486.604}{Ver}%
\TextAt{8}{97.159}{428.499}{The}%
\TextAt{8}{122.413}{428.499}{definition}%
\TextAt{8}{175.641}{428.499}{of}%
\TextAt{8}{190.171}{428.499}{the}%
\TextAt{8}{211.527}{428.499}{norm}%
\TextAt{8}{243.297}{428.499}{residue}%
\TextAt{8}{284.224}{428.499}{symbol}%
\TextAt{8}{325.485}{428.499}{automatically}%
\TextAt{8}{400.52}{428.499}{extends}%
\TextAt{8}{444.376}{428.499}{to}%
\TextAt{8}{459.886}{428.499}{infinite}%
\TextAt{8}{500.754}{428.499}{Galois}%
\TextAt{8}{79.226}{411.453}{extensions}%
\TextAt{13}{135.607}{411.453}{L}%
\TextAt{16}{146.892}{411.453}{\Character{106}}%
\TextAt{13}{153.534}{411.453}{K}%
\TextAt{8}{164.351}{411.453}{.}%
\TextAt{8}{172.8}{411.453}{For}%
\TextAt{8}{193.769}{411.453}{if}%
\TextAt{13}{204.493}{411.453}{L}%
\TextAt{21}{212.457}{409.66}{i}%
\TextAt{16}{219.16}{411.453}{\Character{106}}%
\TextAt{13}{225.801}{411.453}{K}%
\TextAt{8}{240.516}{411.453}{varies}%
\TextAt{8}{273.429}{411.453}{over}%
\TextAt{8}{298.475}{411.453}{the}%
\TextAt{8}{318.631}{411.453}{finite}%
\TextAt{8}{348.543}{411.453}{Galois}%
\TextAt{8}{384.457}{411.453}{subextensions,}%
\TextAt{8}{462.034}{411.453}{then}%
\TextAt{13}{230.793}{386.32}{G}%
\TextAt{9}{240.026}{386.32}{\Character{40}}%
\TextAt{13}{244.579}{386.32}{L}%
\TextAt{16}{255.864}{386.32}{\Character{106}}%
\TextAt{13}{262.506}{386.32}{K}%
\TextAt{9}{273.323}{386.32}{\Character{41}}%
\TextAt{7}{277.876}{391.256}{ab}%
\TextAt{9}{290.629}{386.32}{\Character{61}}%
\NativeMathAt{303.055}{386.32}{11.955}{\varprojlim}%
\TextAt{21}{309.743}{374.117}{i}%
\TextAt{13}{321.306}{386.32}{G}%
\TextAt{9}{330.54}{386.32}{\Character{40}}%
\TextAt{13}{335.092}{386.32}{L}%
\TextAt{21}{343.057}{384.527}{i}%
\TextAt{16}{349.759}{386.32}{\Character{106}}%
\TextAt{13}{356.4}{386.32}{K}%
\TextAt{9}{367.218}{386.32}{\Character{41}}%
\TextAt{7}{371.77}{391.256}{ab}%
\TextAt{8}{79.226}{354.872}{(see}%
\TextAt{8}{103.401}{354.872}{§2,}%
\TextAt{8}{122.97}{354.872}{exercise}%
\TextAt{8}{166.978}{354.872}{6).}%
\TextAt{8}{187.959}{354.872}{As}%
\TextAt{9}{205.953}{354.872}{\Character{40}}%
\TextAt{13}{210.505}{354.872}{a\Character{59}}%
\TextAt{13}{221.874}{354.872}{L}%
\TextAt{21}{229.859}{353.079}{i}%
\TextAt{23}{232.742}{355.348}{\Character{48}}%
\TextAt{16}{240.45}{354.872}{\Character{106}}%
\TextAt{13}{248.295}{354.872}{K}%
\TextAt{9}{259.113}{354.872}{\Character{41}}%
\TextAt{16}{263.665}{354.872}{\Character{106}}%
\TextAt{21}{266.986}{352.036}{L}%
\TextAt{40}{272.745}{354.841}{ab}%
\TextAt{22}{272.745}{349.387}{i}%
\TextAt{9}{285.96}{354.872}{\Character{61}}%
\TextAt{9}{299.577}{354.872}{\Character{40}}%
\TextAt{13}{304.142}{354.872}{a\Character{59}}%
\TextAt{13}{315.511}{354.872}{L}%
\TextAt{21}{323.495}{353.079}{i}%
\TextAt{16}{331.402}{354.872}{\Character{106}}%
\TextAt{13}{339.248}{354.872}{K}%
\TextAt{9}{350.065}{354.872}{\Character{41}}%
\TextAt{8}{359.222}{354.872}{for}%
\TextAt{13}{377.813}{354.872}{L}%
\TextAt{21}{385.778}{353.079}{i}%
\TextAt{23}{388.661}{355.348}{\Character{48}}%
\NativeMathAt{396.368}{354.872}{11.955}{\supseteq}%
\TextAt{13}{410.192}{354.872}{L}%
\TextAt{21}{418.156}{353.079}{i}%
\TextAt{8}{421.537}{354.872}{,}%
\TextAt{8}{429.57}{354.872}{the}%
\TextAt{8}{450.433}{354.872}{individual}%
\TextAt{8}{506.097}{354.872}{norm}%
\TextAt{8}{79.226}{337.826}{residue}%
\TextAt{8}{118.953}{337.826}{symbols}%
\TextAt{9}{163.63}{337.826}{\Character{40}}%
\TextAt{13}{168.182}{337.826}{a\Character{59}}%
\TextAt{13}{179.551}{337.826}{L}%
\TextAt{21}{187.536}{336.033}{i}%
\TextAt{16}{194.238}{337.826}{\Character{106}}%
\TextAt{13}{200.879}{337.826}{K}%
\TextAt{9}{211.697}{337.826}{\Character{41}}%
\TextAt{8}{216.249}{337.826}{,}%
\TextAt{13}{223.399}{337.826}{a}%
\NativeMathAt{232.864}{337.826}{11.955}{\in}%
\TextAt{13}{244.155}{337.826}{A}%
\TextAt{21}{252.931}{336.033}{K}%
\TextAt{8}{261.133}{337.826}{,}%
\TextAt{8}{268.283}{337.826}{determine}%
\TextAt{8}{322.906}{337.826}{an}%
\TextAt{8}{339.165}{337.826}{element}%
\TextAt{9}{244.189}{312.693}{\Character{40}}%
\TextAt{13}{248.741}{312.693}{a\Character{59}}%
\TextAt{13}{260.11}{312.693}{L}%
\TextAt{16}{271.416}{312.693}{\Character{106}}%
\TextAt{13}{278.058}{312.693}{K}%
\TextAt{9}{288.875}{312.693}{\Character{41}}%
\NativeMathAt{296.748}{312.693}{11.955}{\in}%
\TextAt{13}{308.039}{312.693}{G}%
\TextAt{9}{317.273}{312.693}{\Character{40}}%
\TextAt{13}{321.825}{312.693}{L}%
\TextAt{16}{333.111}{312.693}{\Character{106}}%
\TextAt{13}{339.752}{312.693}{K}%
\TextAt{9}{350.57}{312.693}{\Character{41}}%
\TextAt{7}{355.122}{317.629}{ab}%
\TextAt{13}{364.555}{312.693}{\Character{58}}%
\TextAt{8}{79.226}{287.56}{In}%
\TextAt{8}{95.08}{287.56}{the}%
\TextAt{8}{116.468}{287.56}{special}%
\TextAt{8}{155.802}{287.56}{case}%
\TextAt{8}{181.805}{287.56}{of}%
\TextAt{8}{196.367}{287.56}{the}%
\TextAt{8}{217.756}{287.56}{extension}%
\TextAt{20}{273.492}{290.582}{\Character{101}}%
\TextAt{13}{270.754}{287.56}{K}%
\TextAt{16}{286.992}{287.56}{\Character{106}}%
\TextAt{13}{295.734}{287.56}{K}%
\TextAt{8}{306.551}{287.56}{,}%
\TextAt{8}{315.241}{287.56}{we}%
\TextAt{8}{333.701}{287.56}{find}%
\TextAt{8}{358.341}{287.56}{the}%
\TextAt{8}{379.73}{287.56}{following}%
\TextAt{8}{430.397}{287.56}{intimate}%
\TextAt{8}{478.134}{287.56}{connection}%
\TextAt{8}{79.226}{270.514}{between}%
\TextAt{8}{124.417}{270.514}{the}%
\TextAt{8}{144.573}{270.514}{maps}%
\TextAt{13}{175.202}{270.514}{d}%
\TextAt{21}{181.285}{268.721}{K}%
\TextAt{8}{189.488}{270.514}{,}%
\TextAt{13}{196.637}{270.514}{v}%
\TextAt{21}{202.296}{268.721}{K}%
\TextAt{8}{210.499}{270.514}{,}%
\TextAt{8}{217.648}{270.514}{and}%
\TextAt{9}{240.411}{270.514}{\Character{40}}%
\TextAt{13}{248.86}{270.514}{\Character{59}}%
\TextAt{20}{256.843}{273.536}{\Character{101}}%
\TextAt{13}{254.104}{270.514}{K}%
\TextAt{16}{268.243}{270.514}{\Character{106}}%
\TextAt{13}{274.885}{270.514}{K}%
\TextAt{9}{285.702}{270.514}{\Character{41}}%
\TextAt{8}{290.255}{270.514}{.}%
\TextAt{12}{79.226}{245.607}{Proposition}%
\TextAt{12}{151.937}{245.607}{(6.5).}%
\TextAt{8}{187.253}{245.607}{One}%
\TextAt{8}{211.964}{245.607}{has}%
\TextAt{9}{163.896}{220.474}{\Character{40}}%
\TextAt{13}{168.448}{220.474}{a\Character{59}}%
\TextAt{20}{182.575}{223.496}{\Character{101}}%
\TextAt{13}{179.837}{220.474}{K}%
\TextAt{16}{193.976}{220.474}{\Character{106}}%
\TextAt{13}{200.617}{220.474}{K}%
\TextAt{9}{211.435}{220.474}{\Character{41}}%
\TextAt{9}{219.29}{220.474}{\Character{61}}%
\TextAt{13}{231.734}{220.474}{\Character{39}}%
\TextAt{21}{239.402}{226.584}{v}%
\TextAt{22}{243.525}{225.245}{K}%
\TextAt{17}{250.568}{226.584}{\Character{40}}%
\TextAt{21}{253.862}{226.584}{a}%
\TextAt{17}{258.36}{226.584}{\Character{41}}%
\TextAt{21}{239.402}{217.233}{K}%
\TextAt{13}{262.151}{220.474}{\Character{59}}%
\TextAt{8}{291.306}{220.474}{and}%
\TextAt{8}{314.068}{220.474}{thus}%
\TextAt{13}{339.82}{220.474}{d}%
\TextAt{21}{345.902}{218.681}{K}%
\TextAt{16}{356.762}{220.474}{\Character{14}}%
\TextAt{9}{365.396}{220.474}{\Character{40}}%
\TextAt{13}{373.846}{220.474}{\Character{59}}%
\TextAt{20}{381.828}{223.496}{\Character{101}}%
\TextAt{13}{379.09}{220.474}{K}%
\TextAt{16}{393.228}{220.474}{\Character{106}}%
\TextAt{13}{399.87}{220.474}{K}%
\TextAt{9}{410.688}{220.474}{\Character{41}}%
\TextAt{9}{418.543}{220.474}{\Character{61}}%
\TextAt{13}{430.986}{220.474}{v}%
\TextAt{21}{436.646}{218.681}{K}%
\TextAt{13}{444.848}{220.474}{\Character{58}}%
\TextAt{19}{79.226}{179.618}{Proof.}%
\TextAt{8}{116}{179.618}{Let}%
\TextAt{13}{137.1}{179.618}{L}%
\TextAt{16}{148.607}{179.618}{\Character{106}}%
\TextAt{13}{155.47}{179.618}{K}%
\TextAt{8}{170.315}{179.618}{be}%
\TextAt{8}{186.369}{179.618}{the}%
\TextAt{8}{206.656}{179.618}{subextension}%
\TextAt{8}{276.496}{179.618}{of}%
\TextAt{20}{292.695}{182.64}{\Character{101}}%
\TextAt{13}{289.956}{179.618}{K}%
\TextAt{16}{304.317}{179.618}{\Character{106}}%
\TextAt{13}{311.18}{179.618}{K}%
\TextAt{8}{326.025}{179.618}{of}%
\TextAt{8}{339.485}{179.618}{degree}%
\TextAt{13}{376.031}{179.618}{f}%
\TextAt{8}{383.077}{179.618}{.}%
\TextAt{8}{391.915}{179.618}{As}%
\TextAt{13}{409.333}{179.618}{Z}%
\TextAt{9}{418.203}{179.618}{\Character{47}}%
\TextAt{13}{424.056}{179.618}{f}%
\TextAt{13}{431.086}{179.618}{Z}%
\TextAt{9}{443.516}{179.618}{\Character{61}}%
\TextAt{18}{456.163}{179.618}{Z}%
\TextAt{9}{464.133}{179.618}{\Character{47}}%
\TextAt{13}{469.986}{179.618}{f}%
\TextAt{18}{477.033}{179.618}{Z}%
\TextAt{8}{485.003}{179.618}{,}%
\TextAt{8}{492.315}{179.618}{we}%
\TextAt{8}{509.672}{179.618}{have}%
\TextAt{13}{79.226}{162.572}{v}%
\TextAt{21}{84.885}{160.779}{K}%
\TextAt{9}{93.088}{162.572}{\Character{40}}%
\TextAt{13}{97.64}{162.572}{a}%
\TextAt{9}{103.785}{162.572}{\Character{41}}%
\TextAt{9}{112.728}{162.572}{\Character{61}}%
\TextAt{13}{126.26}{162.572}{n}%
\TextAt{9}{136.339}{162.572}{\Character{43}}%
\TextAt{13}{148.535}{162.572}{f}%
\TextAt{13}{155.565}{162.572}{z}%
\TextAt{8}{161.552}{162.572}{,}%
\TextAt{8}{169.5}{162.572}{with}%
\TextAt{13}{196.798}{162.572}{n}%
\NativeMathAt{208.194}{162.572}{11.955}{\in}%
\TextAt{18}{220.573}{162.572}{Z}%
\TextAt{8}{228.543}{162.572}{,}%
\TextAt{13}{236.49}{162.572}{z}%
\NativeMathAt{246.87}{162.572}{11.955}{\in}%
\TextAt{13}{259.248}{162.572}{Z}%
\TextAt{8}{268.119}{162.572}{;}%
\TextAt{8}{276.226}{162.572}{that}%
\TextAt{8}{302.233}{162.572}{is,}%
\TextAt{13}{318.047}{162.572}{a}%
\TextAt{9}{328.6}{162.572}{\Character{61}}%
\TextAt{13}{342.114}{162.572}{u\Character{25}}%
\TextAt{21}{355.845}{166.911}{n}%
\TextAt{21}{355.416}{159.552}{K}%
\TextAt{13}{363.619}{162.572}{b}%
\TextAt{21}{368.596}{166.911}{f}%
\TextAt{8}{374.041}{162.572}{,}%
\TextAt{8}{381.988}{162.572}{with}%
\TextAt{13}{409.286}{162.572}{u}%
\NativeMathAt{420.358}{162.572}{11.955}{\in}%
\TextAt{13}{432.736}{162.572}{U}%
\TextAt{21}{440.687}{160.779}{K}%
\TextAt{8}{448.889}{162.572}{,}%
\TextAt{13}{456.836}{162.572}{b}%
\NativeMathAt{466.222}{162.572}{11.955}{\in}%
\TextAt{13}{478.601}{162.572}{A}%
\TextAt{21}{487.376}{160.779}{K}%
\TextAt{8}{495.579}{162.572}{.}%
\TextAt{8}{505.942}{162.572}{From}%
\TextAt{8}{79.226}{145.526}{(5.7),}%
\TextAt{8}{110.453}{145.526}{we}%
\TextAt{8}{127.68}{145.526}{obtain}%
\TextAt{9}{176.253}{120.393}{\Character{40}}%
\TextAt{13}{180.806}{120.393}{a\Character{59}}%
\TextAt{20}{194.933}{123.415}{\Character{101}}%
\TextAt{13}{192.195}{120.393}{K}%
\TextAt{16}{206.333}{120.393}{\Character{106}}%
\TextAt{13}{212.975}{120.393}{K}%
\TextAt{9}{223.793}{120.393}{\Character{41}}%
\TextAt{16}{228.345}{120.393}{\Character{106}}%
\TextAt{21}{231.666}{118.6}{L}%
\TextAt{9}{241.244}{120.393}{\Character{61}}%
\TextAt{9}{253.654}{120.393}{\Character{40}}%
\TextAt{13}{258.222}{120.393}{a\Character{59}}%
\TextAt{13}{269.591}{120.393}{L}%
\TextAt{16}{280.896}{120.393}{\Character{106}}%
\TextAt{13}{287.538}{120.393}{K}%
\TextAt{9}{298.355}{120.393}{\Character{41}}%
\TextAt{9}{306.21}{120.393}{\Character{61}}%
\TextAt{9}{318.631}{120.393}{\Character{40}}%
\TextAt{13}{323.206}{120.393}{\Character{25}}%
\TextAt{21}{329.847}{118.6}{K}%
\TextAt{13}{338.049}{120.393}{\Character{59}}%
\TextAt{13}{343.273}{120.393}{L}%
\TextAt{16}{354.579}{120.393}{\Character{106}}%
\TextAt{13}{361.221}{120.393}{K}%
\TextAt{9}{372.038}{120.393}{\Character{41}}%
\TextAt{21}{376.591}{125.33}{n}%
\TextAt{9}{382.227}{120.393}{\Character{40}}%
\TextAt{13}{386.779}{120.393}{b\Character{59}}%
\TextAt{13}{396.977}{120.393}{L}%
\TextAt{16}{408.286}{120.393}{\Character{106}}%
\TextAt{13}{414.928}{120.393}{K}%
\TextAt{9}{425.745}{120.393}{\Character{41}}%
\TextAt{21}{430.298}{125.33}{f}%
\TextAt{9}{241.244}{99.207}{\Character{61}}%
\TextAt{13}{253.669}{99.207}{\Character{39}}%
\TextAt{21}{261.338}{104.143}{n}%
\TextAt{21}{261.338}{96.251}{L}%
\TextAt{3}{267.097}{96.251}{\Character{106}}%
\TextAt{21}{269.449}{96.251}{K}%
\TextAt{9}{280.972}{99.207}{\Character{61}}%
\TextAt{13}{293.398}{99.207}{\Character{39}}%
\TextAt{21}{301.066}{105.317}{v}%
\TextAt{22}{305.19}{103.977}{K}%
\TextAt{17}{312.233}{105.317}{\Character{40}}%
\TextAt{21}{315.526}{105.317}{a}%
\TextAt{17}{320.024}{105.317}{\Character{41}}%
\TextAt{21}{301.066}{95.965}{K}%
\TextAt{16}{323.815}{99.207}{\Character{106}}%
\TextAt{21}{327.136}{97.413}{L}%
\TextAt{13}{333.394}{99.207}{\Character{58}}%
\TextAt{8}{79.226}{72.921}{Thus}%
\TextAt{8}{108.875}{72.921}{we}%
\TextAt{8}{126.102}{72.921}{must}%
\TextAt{8}{155.106}{72.921}{have}%
\TextAt{9}{182.1}{72.921}{\Character{40}}%
\TextAt{13}{186.653}{72.921}{a\Character{59}}%
\TextAt{20}{200.78}{75.943}{\Character{101}}%
\TextAt{13}{198.042}{72.921}{K}%
\TextAt{16}{212.18}{72.921}{\Character{106}}%
\TextAt{13}{218.822}{72.921}{K}%
\TextAt{9}{229.64}{72.921}{\Character{41}}%
\TextAt{9}{237.495}{72.921}{\Character{61}}%
\TextAt{13}{249.938}{72.921}{\Character{39}}%
\TextAt{21}{257.607}{79.032}{v}%
\TextAt{22}{261.73}{77.692}{K}%
\TextAt{17}{268.773}{79.032}{\Character{40}}%
\TextAt{21}{272.066}{79.032}{a}%
\TextAt{17}{276.564}{79.032}{\Character{41}}%
\TextAt{21}{257.607}{69.68}{K}%
\TextAt{8}{280.356}{72.921}{.}%
\Vector{q 0 G 0 g 0.399 w 452.334 -701.079 m 452.334 -693.009 l S Q}%
\Vector{q 0 G 0 g 0.399 w 452.534 -693.208 m 459.707 -693.208 l S Q}%
\Vector{q 0 G 0 g 0.399 w 452.534 -700.879 m 459.707 -700.879 l S Q}%
\Vector{q 0 G 0 g 0.399 w 459.906 -701.079 m 459.906 -693.009 l S Q}%
\TextAt{8}{297.211}{43.033}{277}%
\Destination{78.23}{809.815}{nav.page.277}%
\Destination{79.36}{426.455}{nav.print.303}%
\Destination{79.36}{665.855}{nav.number.4.6.4}%
\Destination{79.36}{259.655}{nav.number.4.6.5}%
\Hotlink{166.478}{352.655}{6.858}{11.047}{nav.ex.4.2.6}%
\Hotlink{108.58}{352.655}{6.858}{11.047}{nav.section.4.2}%
\Hotlink{508.192}{691.655}{25.078}{11.047}{nav.number.4.5.8}%
\Hotlink{101.489}{674.555}{25.078}{11.047}{nav.number.4.5.9}%
\Hotlink{78.726}{143.355}{25.078}{11.047}{nav.number.4.5.7}%
}

% Page 290 (body page 278)
\Page{%
\TextAt{2}{79.226}{796.943}{IV.}%
\TextAt{2}{98.012}{796.943}{Abstract}%
\TextAt{2}{143.252}{796.943}{Class}%
\TextAt{2}{171.844}{796.943}{Field}%
\TextAt{2}{199.575}{796.943}{Theory}%
\Vector{q 0 G 0 g 0.299 w 7.226 18.729 m 460.769 18.729 l S Q}%
\TextAt{8}{97.159}{761.124}{The}%
\TextAt{8}{121.241}{761.124}{main}%
\TextAt{8}{150.536}{761.124}{goal}%
\TextAt{8}{175.288}{761.124}{of}%
\TextAt{8}{188.646}{761.124}{field}%
\TextAt{8}{214.032}{761.124}{theory}%
\TextAt{8}{250.81}{761.124}{is}%
\TextAt{8}{262.603}{761.124}{to}%
\TextAt{8}{276.942}{761.124}{classify}%
\TextAt{8}{317.414}{761.124}{all}%
\TextAt{8}{333.702}{761.124}{algebraic}%
\TextAt{8}{383.165}{761.124}{extensions}%
\TextAt{8}{439.574}{761.124}{of}%
\TextAt{8}{452.933}{761.124}{a}%
\TextAt{8}{462.717}{761.124}{given}%
\TextAt{8}{493.315}{761.124}{field}%
\TextAt{13}{518.7}{761.124}{K}%
\TextAt{8}{529.518}{761.124}{.}%
\TextAt{8}{79.226}{744.078}{The}%
\TextAt{8}{104.153}{744.078}{law}%
\TextAt{8}{126.162}{744.078}{governing}%
\TextAt{8}{180.056}{744.078}{the}%
\TextAt{8}{201.086}{744.078}{constitution}%
\TextAt{8}{267.066}{744.078}{of}%
\TextAt{8}{281.269}{744.078}{extensions}%
\TextAt{8}{338.523}{744.078}{of}%
\TextAt{13}{352.726}{744.078}{K}%
\TextAt{8}{368.313}{744.078}{is}%
\TextAt{8}{380.95}{744.078}{hidden}%
\TextAt{8}{420.187}{744.078}{in}%
\TextAt{8}{434.713}{744.078}{the}%
\TextAt{8}{455.742}{744.078}{inner}%
\TextAt{8}{486.527}{744.078}{structure}%
\TextAt{8}{79.226}{727.032}{of}%
\TextAt{8}{92.947}{727.032}{the}%
\TextAt{8}{113.494}{727.032}{base}%
\TextAt{8}{139.958}{727.032}{field}%
\TextAt{13}{165.706}{727.032}{K}%
\TextAt{8}{180.811}{727.032}{itself,}%
\TextAt{8}{212.897}{727.032}{and}%
\TextAt{8}{236.05}{727.032}{should}%
\TextAt{8}{273.574}{727.032}{therefore}%
\TextAt{8}{323.064}{727.032}{be}%
\TextAt{8}{339.379}{727.032}{expressed}%
\TextAt{8}{392.24}{727.032}{in}%
\TextAt{8}{406.284}{727.032}{terms}%
\TextAt{8}{439.252}{727.032}{of}%
\TextAt{8}{452.973}{727.032}{entities}%
\TextAt{8}{494.071}{727.032}{directly}%
\TextAt{8}{79.226}{709.986}{associated}%
\TextAt{8}{134.995}{709.986}{with}%
\TextAt{8}{161.689}{709.986}{it.}%
\TextAt{8}{178.046}{709.986}{Class}%
\TextAt{8}{208.768}{709.986}{field}%
\TextAt{8}{234.159}{709.986}{theory}%
\TextAt{8}{270.943}{709.986}{solves}%
\TextAt{8}{304.272}{709.986}{this}%
\TextAt{8}{327.129}{709.986}{problem}%
\TextAt{8}{372.688}{709.986}{as}%
\TextAt{8}{387.092}{709.986}{far}%
\TextAt{8}{405.011}{709.986}{as}%
\TextAt{8}{419.415}{709.986}{the}%
\TextAt{8}{439.605}{709.986}{abelian}%
\TextAt{8}{480.286}{709.986}{extensions}%
\TextAt{8}{79.226}{692.94}{of}%
\TextAt{13}{93.689}{692.94}{K}%
\TextAt{8}{109.536}{692.94}{are}%
\TextAt{8}{130.18}{692.94}{concerned.}%
\TextAt{8}{192.751}{692.94}{It}%
\TextAt{8}{206.556}{692.94}{establishes}%
\TextAt{8}{265.755}{692.94}{a}%
\TextAt{8}{276.643}{692.94}{1-1-correspondence}%
\TextAt{8}{378.964}{692.94}{between}%
\TextAt{8}{425.287}{692.94}{these}%
\TextAt{8}{456.391}{692.94}{extensions}%
\TextAt{8}{513.904}{692.94}{and}%
\TextAt{8}{79.226}{675.894}{certain}%
\TextAt{8}{118.248}{675.894}{subgroups}%
\TextAt{8}{173.66}{675.894}{of}%
\TextAt{13}{186.99}{675.894}{A}%
\TextAt{21}{195.765}{674.101}{K}%
\TextAt{8}{203.968}{675.894}{.}%
\TextAt{8}{212.416}{675.894}{More}%
\TextAt{8}{242.651}{675.894}{precisely,}%
\TextAt{8}{292.779}{675.894}{this}%
\TextAt{8}{315.601}{675.894}{is}%
\TextAt{8}{327.365}{675.894}{done}%
\TextAt{8}{355.329}{675.894}{as}%
\TextAt{8}{369.699}{675.894}{follows.}%
\TextAt{8}{97.159}{655.859}{For}%
\TextAt{8}{119.201}{655.859}{every}%
\TextAt{8}{151.165}{655.859}{field}%
\TextAt{13}{177.595}{655.859}{K}%
\TextAt{8}{188.412}{655.859}{,}%
\TextAt{8}{196.902}{655.859}{we}%
\TextAt{8}{215.202}{655.859}{equip}%
\TextAt{8}{247.812}{655.859}{the}%
\TextAt{8}{269.041}{655.859}{group}%
\TextAt{13}{303.289}{655.859}{A}%
\TextAt{21}{312.064}{654.066}{K}%
\TextAt{8}{325.237}{655.859}{with}%
\TextAt{8}{352.969}{655.859}{a}%
\TextAt{8}{363.797}{655.859}{topology}%
\TextAt{8}{413.013}{655.859}{by}%
\TextAt{8}{430.345}{655.859}{declaring}%
\TextAt{8}{481.497}{655.859}{the}%
\TextAt{8}{502.726}{655.859}{cosets}%
\TextAt{13}{79.226}{638.813}{aN}%
\TextAt{21}{94.745}{636.953}{L}%
\TextAt{3}{100.504}{636.953}{\Character{106}}%
\TextAt{21}{102.856}{636.953}{K}%
\TextAt{13}{111.059}{638.813}{A}%
\TextAt{21}{119.834}{637.02}{L}%
\TextAt{8}{130.521}{638.813}{to}%
\TextAt{8}{145.363}{638.813}{be}%
\TextAt{8}{161.819}{638.813}{a}%
\TextAt{8}{172.106}{638.813}{basis}%
\TextAt{8}{201.378}{638.813}{of}%
\TextAt{8}{215.24}{638.813}{neighbourhoods}%
\TextAt{8}{300.391}{638.813}{of}%
\TextAt{13}{314.253}{638.813}{a}%
\NativeMathAt{324.625}{638.813}{11.955}{\in}%
\TextAt{13}{336.822}{638.813}{A}%
\TextAt{21}{345.597}{637.02}{K}%
\TextAt{8}{353.8}{638.813}{,}%
\TextAt{8}{361.614}{638.813}{where}%
\TextAt{13}{395.955}{638.813}{L}%
\TextAt{16}{408.147}{638.813}{\Character{106}}%
\TextAt{13}{415.695}{638.813}{K}%
\TextAt{8}{430.942}{638.813}{varies}%
\TextAt{8}{464.386}{638.813}{over}%
\TextAt{8}{489.964}{638.813}{all}%
\TextAt{8}{506.755}{638.813}{finite}%
\TextAt{8}{79.226}{621.767}{Galois}%
\TextAt{8}{115.139}{621.767}{extensions}%
\TextAt{8}{171.52}{621.767}{of}%
\TextAt{13}{184.85}{621.767}{K}%
\TextAt{8}{195.668}{621.767}{.}%
\TextAt{8}{204.116}{621.767}{We}%
\TextAt{8}{224.261}{621.767}{call}%
\TextAt{8}{245.72}{621.767}{this}%
\TextAt{8}{268.542}{621.767}{topology}%
\TextAt{8}{316.686}{621.767}{the}%
\TextAt{12}{336.842}{621.767}{norm}%
\TextAt{12}{372.242}{621.767}{topology}%
\TextAt{8}{426.984}{621.767}{of}%
\TextAt{13}{440.314}{621.767}{A}%
\TextAt{21}{449.09}{619.974}{K}%
\TextAt{8}{457.292}{621.767}{.}%
\TextAt{12}{79.226}{595.755}{Proposition}%
\TextAt{12}{151.973}{595.755}{(6.6).}%
\TextAt{8}{187.319}{595.755}{(i)}%
\TextAt{8}{203.609}{595.755}{The}%
\TextAt{8}{227.694}{595.755}{open}%
\TextAt{8}{256.01}{595.755}{subgroups}%
\TextAt{8}{311.453}{595.755}{of}%
\TextAt{13}{324.814}{595.755}{A}%
\TextAt{21}{333.59}{593.962}{K}%
\TextAt{8}{345.72}{595.755}{are}%
\TextAt{8}{365.262}{595.755}{precisely}%
\TextAt{8}{413.149}{595.755}{the}%
\TextAt{8}{433.337}{595.755}{closed}%
\TextAt{8}{467.894}{595.755}{subgroups}%
\TextAt{8}{523.337}{595.755}{of}%
\TextAt{8}{79.226}{578.709}{finite}%
\TextAt{8}{109.138}{578.709}{index.}%
\TextAt{8}{97.159}{558.674}{(ii)}%
\TextAt{8}{116.67}{558.674}{The}%
\TextAt{8}{140.724}{558.674}{valuation}%
\TextAt{13}{191.796}{558.674}{v}%
\TextAt{21}{197.455}{556.881}{K}%
\TextAt{9}{208.979}{558.674}{\Character{58}}%
\TextAt{13}{215.551}{558.674}{A}%
\TextAt{21}{224.327}{556.881}{K}%
\NativeMathAt{235.85}{558.674}{11.955}{\rightarrow}%
\TextAt{20}{251.79}{561.762}{\Character{98}}%
\TextAt{18}{251.126}{558.674}{Z}%
\TextAt{8}{262.994}{558.674}{is}%
\TextAt{8}{274.758}{558.674}{continuous.}%
\TextAt{8}{97.159}{538.639}{(iii)}%
\TextAt{8}{119.922}{538.639}{If}%
\TextAt{13}{131.614}{538.639}{L}%
\TextAt{16}{142.899}{538.639}{\Character{106}}%
\TextAt{13}{149.541}{538.639}{K}%
\TextAt{8}{164.256}{538.639}{is}%
\TextAt{8}{176.02}{538.639}{a}%
\TextAt{8}{185.775}{538.639}{finite}%
\TextAt{8}{215.687}{538.639}{extension,}%
\TextAt{8}{270.705}{538.639}{then}%
\TextAt{13}{297.365}{538.639}{N}%
\TextAt{21}{306.738}{536.78}{L}%
\TextAt{3}{312.498}{536.78}{\Character{106}}%
\TextAt{21}{314.85}{536.78}{K}%
\TextAt{9}{326.373}{538.639}{\Character{58}}%
\TextAt{13}{332.946}{538.639}{A}%
\TextAt{21}{341.721}{536.846}{L}%
\NativeMathAt{351.299}{538.639}{11.955}{\rightarrow}%
\TextAt{13}{366.575}{538.639}{A}%
\TextAt{21}{375.351}{536.846}{K}%
\TextAt{8}{387.451}{538.639}{is}%
\TextAt{8}{399.215}{538.639}{continuous.}%
\TextAt{8}{97.159}{518.605}{(iv)}%
\TextAt{13}{119.599}{518.605}{A}%
\TextAt{21}{128.374}{516.811}{K}%
\TextAt{8}{140.474}{518.605}{is}%
\TextAt{8}{152.238}{518.605}{Hausdorff}%
\TextAt{8}{205.63}{518.605}{if}%
\TextAt{8}{216.354}{518.605}{and}%
\TextAt{8}{239.116}{518.605}{only}%
\TextAt{8}{264.808}{518.605}{if}%
\TextAt{8}{275.532}{518.605}{the}%
\TextAt{8}{295.688}{518.605}{group}%
\TextAt{13}{261.637}{489.603}{A}%
\TextAt{17}{270.413}{494.54}{0}%
\TextAt{21}{270.413}{486.648}{K}%
\TextAt{9}{281.936}{489.603}{\Character{61}}%
\TextAt{20}{294.362}{500.961}{\Character{92}}%
\TextAt{21}{298.124}{475.589}{L}%
\TextAt{13}{309.638}{489.603}{N}%
\TextAt{21}{319.012}{487.744}{L}%
\TextAt{3}{324.771}{487.744}{\Character{106}}%
\TextAt{21}{327.123}{487.744}{K}%
\TextAt{13}{335.326}{489.603}{A}%
\TextAt{21}{344.101}{487.81}{L}%
\TextAt{8}{79.226}{453.026}{of}%
\TextAt{8}{92.556}{453.026}{universal}%
\TextAt{8}{142.051}{453.026}{norms}%
\TextAt{8}{177.235}{453.026}{is}%
\TextAt{8}{188.999}{453.026}{trivial.}%
\TextAt{19}{79.226}{418.047}{Proof.}%
\TextAt{8}{116}{418.047}{(i)}%
\TextAt{8}{132.259}{418.047}{If}%
\TextAt{16}{143.951}{418.047}{\Character{78}}%
\TextAt{8}{159.42}{418.047}{is}%
\TextAt{8}{171.184}{418.047}{a}%
\TextAt{8}{180.939}{418.047}{subgroup}%
\TextAt{8}{231.736}{418.047}{of}%
\TextAt{13}{245.066}{418.047}{A}%
\TextAt{21}{253.842}{416.254}{K}%
\TextAt{8}{262.044}{418.047}{,}%
\TextAt{8}{269.194}{418.047}{then}%
\TextAt{16}{253.067}{389.046}{\Character{78}}%
\TextAt{9}{267.959}{389.046}{\Character{61}}%
\TextAt{13}{280.384}{389.046}{A}%
\TextAt{21}{289.16}{387.253}{K}%
\TextAt{16}{300.019}{389.046}{\Character{110}}%
\TextAt{20}{315.669}{400.403}{\Character{91}}%
\TextAt{21}{308.653}{374.943}{a}%
\TextAt{3}{313.151}{374.943}{\Character{78}}%
\NativeMathAt{321.265}{374.943}{7.97}{\neq}%
\TextAt{3}{327.853}{374.943}{\Character{78}}%
\TextAt{13}{337.961}{389.046}{a}%
\TextAt{16}{344.106}{389.046}{\Character{78}}%
\TextAt{13}{355.677}{389.046}{\Character{58}}%
\TextAt{8}{79.226}{350.83}{Now,}%
\TextAt{8}{109.195}{350.83}{if}%
\TextAt{16}{119.964}{350.83}{\Character{78}}%
\TextAt{8}{135.477}{350.83}{is}%
\TextAt{8}{147.286}{350.83}{open,}%
\TextAt{8}{178.881}{350.83}{so}%
\TextAt{8}{193.296}{350.83}{are}%
\TextAt{8}{212.852}{350.83}{all}%
\TextAt{8}{229.157}{350.83}{cosets}%
\TextAt{13}{263.143}{350.83}{a}%
\TextAt{16}{269.288}{350.83}{\Character{78}}%
\TextAt{8}{280.859}{350.83}{,}%
\TextAt{8}{288.064}{350.83}{so}%
\TextAt{8}{302.48}{350.83}{that}%
\TextAt{16}{327.894}{350.83}{\Character{78}}%
\TextAt{8}{343.408}{350.83}{is}%
\TextAt{8}{355.217}{350.83}{closed,}%
\TextAt{8}{393.052}{350.83}{and}%
\TextAt{8}{415.86}{350.83}{since}%
\TextAt{16}{444.574}{350.83}{\Character{78}}%
\TextAt{8}{460.087}{350.83}{has}%
\TextAt{8}{481.006}{350.83}{to}%
\TextAt{8}{495.362}{350.83}{contain}%
\TextAt{8}{79.226}{333.784}{one}%
\TextAt{8}{100.23}{333.784}{of}%
\TextAt{8}{113.104}{333.784}{the}%
\TextAt{8}{132.804}{333.784}{neighbourhoods}%
\TextAt{13}{216.967}{333.784}{N}%
\TextAt{21}{226.341}{331.924}{L}%
\TextAt{3}{232.1}{331.924}{\Character{106}}%
\TextAt{21}{234.452}{331.924}{K}%
\TextAt{13}{242.655}{333.784}{A}%
\TextAt{21}{251.43}{331.991}{L}%
\TextAt{8}{261.129}{333.784}{of}%
\TextAt{8}{274.003}{333.784}{the}%
\TextAt{8}{293.704}{333.784}{basis}%
\TextAt{8}{321.988}{333.784}{of}%
\TextAt{8}{334.862}{333.784}{neighbourhoods}%
\TextAt{8}{419.025}{333.784}{of}%
\TextAt{9}{431.899}{333.784}{1}%
\TextAt{8}{437.752}{333.784}{,}%
\TextAt{16}{444.537}{333.784}{\Character{78}}%
\TextAt{8}{459.549}{333.784}{is}%
\TextAt{8}{470.857}{333.784}{also}%
\TextAt{8}{493.881}{333.784}{of}%
\TextAt{8}{506.755}{333.784}{finite}%
\TextAt{8}{79.226}{316.738}{index.}%
\TextAt{8}{115.562}{316.738}{If,}%
\TextAt{8}{130.61}{316.738}{conversely,}%
\TextAt{16}{188.982}{316.738}{\Character{78}}%
\TextAt{8}{204.533}{316.738}{is}%
\TextAt{8}{216.38}{316.738}{closed}%
\TextAt{8}{250.989}{316.738}{and}%
\TextAt{8}{273.834}{316.738}{of}%
\TextAt{8}{287.246}{316.738}{finite}%
\TextAt{8}{317.241}{316.738}{index,}%
\TextAt{8}{352.134}{316.738}{then}%
\TextAt{8}{378.876}{316.738}{the}%
\TextAt{8}{399.115}{316.738}{union}%
\TextAt{8}{431.716}{316.738}{of}%
\TextAt{8}{445.128}{316.738}{the}%
\TextAt{8}{465.367}{316.738}{finitely}%
\TextAt{8}{504.794}{316.738}{many}%
\TextAt{8}{79.226}{299.692}{cosets}%
\TextAt{13}{113.167}{299.692}{a}%
\TextAt{16}{119.312}{299.692}{\Character{78}}%
\NativeMathAt{134.196}{299.692}{11.955}{\neq}%
\TextAt{16}{146.629}{299.692}{\Character{78}}%
\TextAt{8}{162.097}{299.692}{is}%
\TextAt{8}{173.861}{299.692}{closed,}%
\TextAt{8}{211.639}{299.692}{and}%
\TextAt{8}{234.402}{299.692}{so}%
\TextAt{16}{248.772}{299.692}{\Character{78}}%
\TextAt{8}{264.24}{299.692}{is}%
\TextAt{8}{276.004}{299.692}{open.}%
\TextAt{8}{97.159}{279.657}{(ii)}%
\TextAt{8}{117.051}{279.657}{The}%
\TextAt{8}{141.486}{279.657}{groups}%
\TextAt{13}{179.658}{279.657}{f}%
\TextAt{20}{187.368}{282.746}{\Character{98}}%
\TextAt{18}{186.704}{279.657}{Z}%
\TextAt{8}{194.674}{279.657}{,}%
\TextAt{13}{202.3}{279.657}{f}%
\NativeMathAt{213.317}{279.657}{11.955}{\in}%
\TextAt{18}{225.257}{279.657}{N}%
\TextAt{8}{233.892}{279.657}{,}%
\TextAt{8}{241.518}{279.657}{form}%
\TextAt{8}{269.539}{279.657}{a}%
\TextAt{8}{279.676}{279.657}{basis}%
\TextAt{8}{308.797}{279.657}{of}%
\TextAt{8}{322.508}{279.657}{neighbourhoods}%
\TextAt{8}{407.508}{279.657}{of}%
\TextAt{9}{421.22}{279.657}{0}%
\NativeMathAt{431.043}{279.657}{11.955}{\in}%
\TextAt{20}{443.648}{282.746}{\Character{98}}%
\TextAt{18}{442.984}{279.657}{Z}%
\TextAt{8}{455.233}{279.657}{(see}%
\TextAt{8}{479.082}{279.657}{§2),}%
\TextAt{8}{502.799}{279.657}{and}%
\TextAt{8}{525.943}{279.657}{if}%
\TextAt{13}{79.226}{262.611}{L}%
\TextAt{16}{90.511}{262.611}{\Character{106}}%
\TextAt{13}{97.153}{262.611}{K}%
\TextAt{8}{111.868}{262.611}{is}%
\TextAt{8}{123.632}{262.611}{the}%
\TextAt{8}{143.789}{262.611}{unramified}%
\TextAt{8}{202.321}{262.611}{extension}%
\TextAt{8}{254.087}{262.611}{of}%
\TextAt{8}{267.417}{262.611}{degree}%
\TextAt{13}{303.832}{262.611}{f}%
\TextAt{8}{310.879}{262.611}{,}%
\TextAt{8}{318.028}{262.611}{then}%
\TextAt{8}{344.688}{262.611}{it}%
\TextAt{8}{356.392}{262.611}{follows}%
\TextAt{8}{394.828}{262.611}{from}%
\TextAt{8}{422.468}{262.611}{(4.7)}%
\TextAt{8}{450.444}{262.611}{that}%
\TextAt{13}{227.627}{233.61}{v}%
\TextAt{21}{233.286}{231.817}{K}%
\TextAt{9}{241.489}{233.61}{\Character{40}}%
\TextAt{13}{246.041}{233.61}{N}%
\TextAt{21}{255.415}{231.751}{L}%
\TextAt{3}{261.174}{231.751}{\Character{106}}%
\TextAt{21}{263.526}{231.751}{K}%
\TextAt{13}{271.729}{233.61}{A}%
\TextAt{21}{280.504}{231.817}{L}%
\TextAt{9}{286.762}{233.61}{\Character{41}}%
\TextAt{9}{294.617}{233.61}{\Character{61}}%
\TextAt{13}{307.06}{233.61}{f}%
\TextAt{13}{314.09}{233.61}{v}%
\TextAt{21}{319.766}{231.817}{L}%
\TextAt{9}{326.023}{233.61}{\Character{40}}%
\TextAt{13}{330.575}{233.61}{A}%
\TextAt{21}{339.351}{231.817}{L}%
\TextAt{9}{345.608}{233.61}{\Character{41}}%
\NativeMathAt{353.481}{233.61}{11.955}{\subseteq}%
\TextAt{13}{366.1}{233.61}{f}%
\TextAt{20}{373.811}{236.699}{\Character{98}}%
\TextAt{18}{373.147}{233.61}{Z}%
\TextAt{13}{381.117}{233.61}{\Character{59}}%
\TextAt{8}{79.226}{204.609}{which}%
\TextAt{8}{112.713}{204.609}{shows}%
\TextAt{8}{146.331}{204.609}{the}%
\TextAt{8}{166.487}{204.609}{continuity}%
\TextAt{8}{221.78}{204.609}{of}%
\TextAt{13}{235.11}{204.609}{v}%
\TextAt{21}{240.769}{202.816}{K}%
\TextAt{8}{248.972}{204.609}{.}%
\TextAt{8}{97.159}{184.574}{(iii)}%
\TextAt{8}{119.922}{184.574}{Let}%
\TextAt{13}{140.891}{184.574}{N}%
\TextAt{21}{150.265}{182.715}{M}%
\TextAt{3}{159.246}{182.715}{\Character{106}}%
\TextAt{21}{161.599}{182.715}{K}%
\TextAt{13}{169.801}{184.574}{A}%
\TextAt{21}{178.577}{182.781}{M}%
\TextAt{8}{191.954}{184.574}{be}%
\TextAt{8}{207.878}{184.574}{an}%
\TextAt{8}{224.137}{184.574}{open}%
\TextAt{8}{252.423}{184.574}{neighbourhood}%
\TextAt{8}{332.427}{184.574}{of}%
\TextAt{9}{345.757}{184.574}{1}%
\NativeMathAt{354.931}{184.574}{11.955}{\in}%
\TextAt{13}{366.222}{184.574}{A}%
\TextAt{21}{374.997}{182.781}{K}%
\TextAt{8}{383.2}{184.574}{.}%
\TextAt{8}{391.648}{184.574}{Then}%
\TextAt{13}{189.839}{155.573}{N}%
\TextAt{21}{199.212}{153.714}{L}%
\TextAt{3}{204.972}{153.714}{\Character{106}}%
\TextAt{21}{207.324}{153.714}{K}%
\TextAt{9}{215.527}{155.573}{\Character{40}}%
\TextAt{13}{220.079}{155.573}{N}%
\TextAt{21}{229.453}{153.714}{M}%
\TextAt{21}{238.427}{153.714}{L}%
\TextAt{3}{244.193}{153.714}{\Character{106}}%
\TextAt{21}{246.546}{153.714}{L}%
\TextAt{13}{252.803}{155.573}{A}%
\TextAt{21}{261.578}{153.78}{M}%
\TextAt{21}{270.552}{153.78}{L}%
\TextAt{9}{276.817}{155.573}{\Character{41}}%
\TextAt{9}{284.672}{155.573}{\Character{61}}%
\TextAt{13}{297.116}{155.573}{N}%
\TextAt{21}{306.49}{153.714}{M}%
\TextAt{21}{315.464}{153.714}{L}%
\TextAt{3}{321.23}{153.714}{\Character{106}}%
\TextAt{21}{323.583}{153.714}{K}%
\TextAt{13}{331.785}{155.573}{A}%
\TextAt{21}{340.561}{153.78}{M}%
\TextAt{21}{349.535}{153.78}{L}%
\NativeMathAt{359.12}{155.573}{11.955}{\subseteq}%
\TextAt{13}{371.74}{155.573}{N}%
\TextAt{21}{381.114}{153.714}{M}%
\TextAt{3}{390.095}{153.714}{\Character{106}}%
\TextAt{21}{392.448}{153.714}{K}%
\TextAt{13}{400.65}{155.573}{A}%
\TextAt{21}{409.426}{153.78}{M}%
\TextAt{13}{418.905}{155.573}{\Character{59}}%
\TextAt{8}{79.226}{126.572}{which}%
\TextAt{8}{112.713}{126.572}{shows}%
\TextAt{8}{146.331}{126.572}{the}%
\TextAt{8}{166.487}{126.572}{continuity}%
\TextAt{8}{221.78}{126.572}{of}%
\TextAt{13}{235.11}{126.572}{N}%
\TextAt{21}{244.483}{124.712}{L}%
\TextAt{3}{250.242}{124.712}{\Character{106}}%
\TextAt{21}{252.595}{124.712}{K}%
\TextAt{8}{260.797}{126.572}{.}%
\TextAt{8}{97.159}{106.537}{(iv)}%
\TextAt{8}{119.599}{106.537}{is}%
\TextAt{8}{131.363}{106.537}{self-evident.}%
\Vector{q 0 G 0 g 0.399 w 452.334 -667.463 m 452.334 -659.393 l S Q}%
\Vector{q 0 G 0 g 0.399 w 452.534 -659.592 m 459.707 -659.592 l S Q}%
\Vector{q 0 G 0 g 0.399 w 452.534 -667.263 m 459.707 -667.263 l S Q}%
\Vector{q 0 G 0 g 0.399 w 459.906 -667.463 m 459.906 -659.393 l S Q}%
\TextAt{8}{97.159}{77.536}{The}%
\TextAt{8}{121.213}{77.536}{finite}%
\TextAt{8}{151.125}{77.536}{abelian}%
\TextAt{8}{191.772}{77.536}{extensions}%
\TextAt{13}{248.153}{77.536}{L}%
\TextAt{16}{259.438}{77.536}{\Character{106}}%
\TextAt{13}{266.08}{77.536}{K}%
\TextAt{8}{280.795}{77.536}{are}%
\TextAt{8}{300.306}{77.536}{now}%
\TextAt{8}{324.694}{77.536}{classified}%
\TextAt{8}{373.591}{77.536}{as}%
\TextAt{8}{387.961}{77.536}{follows.}%
\TextAt{8}{297.211}{43.033}{278}%
\Destination{78.23}{809.815}{nav.page.278}%
\Destination{79.36}{533.555}{nav.print.304}%
\Destination{79.36}{609.755}{nav.number.4.6.6}%
\Hotlink{484.261}{277.455}{6.858}{11.047}{nav.section.4.2}%
\Hotlink{421.968}{260.455}{25.078}{11.047}{nav.number.4.4.7}%
}

% Page 291 (body page 279)
\Page{%
\TextAt{2}{79.226}{796.943}{IV.}%
\TextAt{2}{98.012}{796.943}{Abstract}%
\TextAt{2}{143.252}{796.943}{Class}%
\TextAt{2}{171.844}{796.943}{Field}%
\TextAt{2}{199.575}{796.943}{Theory}%
\Vector{q 0 G 0 g 0.299 w 7.226 18.729 m 460.769 18.729 l S Q}%
\TextAt{12}{79.226}{761.124}{Theorem}%
\TextAt{12}{136.24}{761.124}{(6.7).}%
\TextAt{8}{171.556}{761.124}{Associating}%
\TextAt{13}{252.798}{735.136}{L}%
\NativeArrowAt{264.083}{735.136}{11.955}{19.26}{\longmapsto}%
\TextAt{16}{286.654}{735.136}{\Character{78}}%
\TextAt{21}{296.474}{733.343}{L}%
\TextAt{9}{306.052}{735.136}{\Character{61}}%
\TextAt{13}{318.478}{735.136}{N}%
\TextAt{21}{327.851}{733.277}{L}%
\TextAt{3}{333.61}{733.277}{\Character{106}}%
\TextAt{21}{335.963}{733.277}{K}%
\TextAt{13}{344.165}{735.136}{A}%
\TextAt{21}{352.941}{733.343}{L}%
\TextAt{8}{79.226}{709.149}{sets}%
\TextAt{8}{102.795}{709.149}{up}%
\TextAt{8}{120.386}{709.149}{a}%
\TextAt{8}{130.828}{709.149}{1-1-correspondence}%
\TextAt{8}{232.702}{709.149}{between}%
\TextAt{8}{278.579}{709.149}{the}%
\TextAt{8}{299.422}{709.149}{finite}%
\TextAt{8}{330.02}{709.149}{abelian}%
\TextAt{8}{371.354}{709.149}{extensions}%
\TextAt{13}{428.421}{709.149}{L}%
\TextAt{16}{440.876}{709.149}{\Character{106}}%
\TextAt{13}{448.688}{709.149}{K}%
\TextAt{8}{464.089}{709.149}{and}%
\TextAt{8}{487.538}{709.149}{the}%
\TextAt{8}{508.381}{709.149}{open}%
\TextAt{8}{79.226}{692.103}{subgroups}%
\TextAt{16}{134.638}{692.103}{\Character{78}}%
\TextAt{8}{150.107}{692.103}{of}%
\TextAt{13}{163.437}{692.103}{A}%
\TextAt{21}{172.212}{690.31}{K}%
\TextAt{8}{180.415}{692.103}{.}%
\TextAt{8}{188.863}{692.103}{Furthermore,}%
\TextAt{8}{259.912}{692.103}{one}%
\TextAt{8}{281.372}{692.103}{has}%
\TextAt{13}{116.311}{666.116}{L}%
\TextAt{17}{124.275}{664.322}{1}%
\NativeMathAt{132.328}{666.116}{11.955}{\subseteq}%
\TextAt{13}{144.948}{666.116}{L}%
\TextAt{17}{152.912}{664.322}{2}%
\NativeArrowAt{160.965}{666.116}{11.955}{21.926}{\Longleftrightarrow}%
\TextAt{16}{186.202}{666.116}{\Character{78}}%
\TextAt{21}{196.013}{664.322}{L}%
\TextAt{24}{201.772}{663.216}{1}%
\NativeMathAt{209.742}{666.116}{11.955}{\supseteq}%
\TextAt{16}{222.343}{666.116}{\Character{78}}%
\TextAt{21}{232.17}{664.322}{L}%
\TextAt{24}{237.93}{663.216}{2}%
\TextAt{13}{242.579}{666.116}{\Character{59}}%
\TextAt{16}{271.733}{666.116}{\Character{78}}%
\TextAt{21}{281.542}{664.322}{L}%
\TextAt{24}{287.301}{663.216}{1}%
\TextAt{21}{291.452}{664.322}{L}%
\TextAt{24}{297.211}{663.216}{2}%
\TextAt{9}{305.181}{666.116}{\Character{61}}%
\TextAt{16}{317.607}{666.116}{\Character{78}}%
\TextAt{21}{327.416}{664.322}{L}%
\TextAt{24}{333.175}{663.216}{1}%
\TextAt{16}{340.481}{666.116}{\Character{92}}%
\TextAt{16}{351.097}{666.116}{\Character{78}}%
\TextAt{21}{360.917}{664.322}{L}%
\TextAt{24}{366.676}{663.216}{2}%
\TextAt{13}{371.325}{666.116}{\Character{59}}%
\TextAt{16}{400.48}{666.116}{\Character{78}}%
\TextAt{21}{410.289}{664.322}{L}%
\TextAt{24}{416.048}{663.216}{1}%
\TextAt{3}{420.199}{664.322}{\Character{92}}%
\TextAt{21}{425.844}{664.322}{L}%
\TextAt{24}{431.603}{663.216}{2}%
\TextAt{9}{439.573}{666.116}{\Character{61}}%
\TextAt{16}{451.999}{666.116}{\Character{78}}%
\TextAt{21}{461.808}{664.322}{L}%
\TextAt{24}{467.567}{663.216}{1}%
\TextAt{16}{472.216}{666.116}{\Character{78}}%
\TextAt{21}{482.025}{664.322}{L}%
\TextAt{24}{487.784}{663.216}{2}%
\TextAt{13}{492.434}{666.116}{\Character{58}}%
\TextAt{8}{97.159}{632.023}{The}%
\TextAt{8}{120.369}{632.023}{field}%
\TextAt{13}{144.882}{632.023}{L}%
\TextAt{8}{155.9}{632.023}{corresponding}%
\TextAt{8}{230.243}{632.023}{to}%
\TextAt{8}{243.709}{632.023}{the}%
\TextAt{8}{263.022}{632.023}{subgroup}%
\TextAt{16}{312.976}{632.023}{\Character{78}}%
\TextAt{8}{327.6}{632.023}{of}%
\TextAt{13}{340.087}{632.023}{A}%
\TextAt{21}{348.862}{630.23}{K}%
\TextAt{8}{360.118}{632.023}{is}%
\TextAt{8}{371.039}{632.023}{called}%
\TextAt{8}{403.359}{632.023}{the}%
\TextAt{12}{422.671}{632.023}{class}%
\TextAt{12}{453.059}{632.023}{field}%
\TextAt{8}{480.932}{632.023}{associated}%
\TextAt{8}{79.226}{614.977}{with}%
\TextAt{16}{105.886}{614.977}{\Character{78}}%
\TextAt{8}{117.457}{614.977}{.}%
\TextAt{8}{125.905}{614.977}{By}%
\TextAt{8}{144.268}{614.977}{(6.3),}%
\TextAt{8}{175.495}{614.977}{it}%
\TextAt{8}{187.199}{614.977}{satisfies}%
\TextAt{13}{255.66}{597.931}{G}%
\TextAt{9}{264.893}{597.931}{\Character{40}}%
\TextAt{13}{269.445}{597.931}{L}%
\TextAt{16}{280.731}{597.931}{\Character{106}}%
\TextAt{13}{287.372}{597.931}{K}%
\TextAt{9}{298.19}{597.931}{\Character{41}}%
\NativeMathAt{306.063}{597.931}{11.955}{\cong}%
\TextAt{13}{318.683}{597.931}{A}%
\TextAt{21}{327.458}{596.138}{K}%
\TextAt{9}{335.661}{597.931}{\Character{47}}%
\TextAt{16}{341.514}{597.931}{\Character{78}}%
\TextAt{13}{353.084}{597.931}{\Character{58}}%
\TextAt{19}{79.226}{569.791}{Proof}%
\TextAt{19}{110.62}{569.791}{of}%
\TextAt{19}{124.369}{569.791}{(6.7).}%
\TextAt{8}{159.039}{569.791}{If}%
\TextAt{13}{172.313}{569.791}{L}%
\TextAt{17}{180.277}{567.998}{1}%
\TextAt{8}{190.489}{569.791}{and}%
\TextAt{13}{214.833}{569.791}{L}%
\TextAt{17}{222.798}{567.998}{2}%
\TextAt{8}{233.009}{569.791}{are}%
\TextAt{8}{254.102}{569.791}{abelian}%
\TextAt{8}{296.331}{569.791}{extensions}%
\TextAt{8}{354.293}{569.791}{of}%
\TextAt{13}{369.205}{569.791}{K}%
\TextAt{8}{380.023}{569.791}{,}%
\TextAt{8}{389.149}{569.791}{then}%
\TextAt{8}{417.391}{569.791}{the}%
\TextAt{8}{439.129}{569.791}{transitivity}%
\TextAt{8}{501.599}{569.791}{of}%
\TextAt{8}{516.51}{569.791}{the}%
\TextAt{8}{79.226}{552.745}{norm}%
\TextAt{8}{111.472}{552.745}{implies}%
\TextAt{16}{152.875}{552.745}{\Character{78}}%
\TextAt{21}{162.684}{550.952}{L}%
\TextAt{24}{168.443}{549.845}{1}%
\TextAt{21}{172.594}{550.952}{L}%
\TextAt{24}{178.353}{549.845}{2}%
\NativeMathAt{189.18}{552.745}{11.955}{\subseteq}%
\TextAt{16}{204.638}{552.745}{\Character{78}}%
\TextAt{21}{214.464}{550.952}{L}%
\TextAt{24}{220.223}{549.845}{1}%
\TextAt{16}{228.672}{552.745}{\Character{92}}%
\TextAt{16}{240.424}{552.745}{\Character{78}}%
\TextAt{21}{250.25}{550.952}{L}%
\TextAt{24}{256.009}{549.845}{2}%
\TextAt{8}{260.658}{552.745}{.}%
\TextAt{8}{274.135}{552.745}{If,}%
\TextAt{8}{291.174}{552.745}{conversely,}%
\TextAt{13}{351.539}{552.745}{a}%
\NativeMathAt{363.861}{552.745}{11.955}{\in}%
\TextAt{16}{377.992}{552.745}{\Character{78}}%
\TextAt{21}{387.818}{550.952}{L}%
\TextAt{24}{393.577}{549.845}{1}%
\TextAt{16}{402.025}{552.745}{\Character{92}}%
\TextAt{16}{413.777}{552.745}{\Character{78}}%
\TextAt{21}{423.604}{550.952}{L}%
\TextAt{24}{429.363}{549.845}{2}%
\TextAt{8}{434.012}{552.745}{,}%
\TextAt{8}{443.256}{552.745}{then}%
\TextAt{8}{471.592}{552.745}{the}%
\TextAt{8}{493.425}{552.745}{element}%
\TextAt{9}{79.226}{535.699}{\Character{40}}%
\TextAt{13}{83.778}{535.699}{a\Character{59}}%
\TextAt{13}{95.147}{535.699}{L}%
\TextAt{17}{103.132}{533.906}{1}%
\TextAt{13}{107.864}{535.699}{L}%
\TextAt{17}{115.829}{533.906}{2}%
\TextAt{16}{123.882}{535.699}{\Character{106}}%
\TextAt{13}{130.524}{535.699}{K}%
\TextAt{9}{141.341}{535.699}{\Character{41}}%
\NativeMathAt{149.215}{535.699}{11.955}{\in}%
\TextAt{13}{160.506}{535.699}{G}%
\TextAt{9}{169.739}{535.699}{\Character{40}}%
\TextAt{13}{174.292}{535.699}{L}%
\TextAt{17}{182.256}{533.906}{1}%
\TextAt{13}{186.988}{535.699}{L}%
\TextAt{17}{194.953}{533.906}{2}%
\TextAt{16}{203.006}{535.699}{\Character{106}}%
\TextAt{13}{209.648}{535.699}{K}%
\TextAt{9}{220.465}{535.699}{\Character{41}}%
\TextAt{8}{228.769}{535.699}{projects}%
\TextAt{8}{272.904}{535.699}{trivially}%
\TextAt{8}{316.992}{535.699}{onto}%
\TextAt{13}{343.195}{535.699}{G}%
\TextAt{9}{352.428}{535.699}{\Character{40}}%
\TextAt{13}{356.981}{535.699}{L}%
\TextAt{21}{364.945}{533.906}{i}%
\TextAt{16}{371.647}{535.699}{\Character{106}}%
\TextAt{13}{378.289}{535.699}{K}%
\TextAt{9}{389.107}{535.699}{\Character{41}}%
\TextAt{8}{393.659}{535.699}{,}%
\TextAt{8}{400.691}{535.699}{that}%
\TextAt{8}{425.913}{535.699}{is,}%
\TextAt{9}{440.812}{535.699}{\Character{40}}%
\TextAt{13}{445.364}{535.699}{a\Character{59}}%
\TextAt{13}{456.733}{535.699}{L}%
\TextAt{21}{464.718}{533.906}{i}%
\TextAt{16}{471.42}{535.699}{\Character{106}}%
\TextAt{13}{478.062}{535.699}{K}%
\TextAt{9}{488.879}{535.699}{\Character{41}}%
\TextAt{9}{496.734}{535.699}{\Character{61}}%
\TextAt{9}{509.155}{535.699}{1}%
\TextAt{8}{518.782}{535.699}{for}%
\TextAt{13}{79.226}{518.653}{i}%
\TextAt{9}{86.696}{518.653}{\Character{61}}%
\TextAt{9}{99.261}{518.653}{1}%
\TextAt{13}{105.13}{518.653}{\Character{59}}%
\TextAt{9}{110.374}{518.653}{2}%
\TextAt{8}{116.227}{518.653}{.}%
\TextAt{8}{124.949}{518.653}{Thus}%
\TextAt{9}{154.689}{518.653}{\Character{40}}%
\TextAt{13}{159.241}{518.653}{a\Character{59}}%
\TextAt{13}{170.61}{518.653}{L}%
\TextAt{17}{178.595}{516.86}{1}%
\TextAt{13}{183.327}{518.653}{L}%
\TextAt{17}{191.292}{516.86}{2}%
\TextAt{16}{199.501}{518.653}{\Character{106}}%
\TextAt{13}{206.298}{518.653}{K}%
\TextAt{9}{217.115}{518.653}{\Character{41}}%
\TextAt{9}{225.125}{518.653}{\Character{61}}%
\TextAt{9}{237.702}{518.653}{1}%
\TextAt{8}{243.578}{518.653}{,}%
\TextAt{8}{250.841}{518.653}{i.e.,}%
\TextAt{13}{273.06}{518.653}{a}%
\NativeMathAt{282.682}{518.653}{11.955}{\in}%
\TextAt{16}{294.111}{518.653}{\Character{78}}%
\TextAt{21}{303.937}{516.86}{L}%
\TextAt{24}{309.696}{515.753}{1}%
\TextAt{21}{313.847}{516.86}{L}%
\TextAt{24}{319.606}{515.753}{2}%
\TextAt{8}{324.256}{518.653}{.}%
\TextAt{8}{332.978}{518.653}{We}%
\TextAt{8}{353.213}{518.653}{therefore}%
\TextAt{8}{402.404}{518.653}{have}%
\TextAt{16}{429.49}{518.653}{\Character{78}}%
\TextAt{21}{439.3}{516.86}{L}%
\TextAt{24}{445.059}{515.753}{1}%
\TextAt{21}{449.21}{516.86}{L}%
\TextAt{24}{454.969}{515.753}{2}%
\TextAt{9}{463.094}{518.653}{\Character{61}}%
\TextAt{16}{475.675}{518.653}{\Character{78}}%
\TextAt{21}{485.484}{516.86}{L}%
\TextAt{24}{491.243}{515.753}{1}%
\TextAt{16}{498.611}{518.653}{\Character{92}}%
\TextAt{16}{509.287}{518.653}{\Character{78}}%
\TextAt{21}{519.109}{516.86}{L}%
\TextAt{24}{524.869}{515.753}{2}%
\TextAt{8}{529.518}{518.653}{,}%
\TextAt{8}{79.226}{501.607}{and}%
\TextAt{8}{101.989}{501.607}{so}%
\TextAt{16}{172.825}{475.62}{\Character{78}}%
\TextAt{21}{182.635}{473.827}{L}%
\TextAt{24}{188.394}{472.72}{1}%
\NativeMathAt{196.364}{475.62}{11.955}{\supseteq}%
\TextAt{16}{208.965}{475.62}{\Character{78}}%
\TextAt{21}{218.792}{473.827}{L}%
\TextAt{24}{224.551}{472.72}{2}%
\NativeArrowAt{232.521}{475.62}{11.955}{21.926}{\Longleftrightarrow}%
\TextAt{16}{257.758}{475.62}{\Character{78}}%
\TextAt{21}{267.569}{473.827}{L}%
\TextAt{24}{273.328}{472.72}{1}%
\TextAt{16}{280.634}{475.62}{\Character{92}}%
\TextAt{16}{291.25}{475.62}{\Character{78}}%
\TextAt{21}{301.07}{473.827}{L}%
\TextAt{24}{306.829}{472.72}{2}%
\TextAt{9}{314.799}{475.62}{\Character{61}}%
\TextAt{16}{327.224}{475.62}{\Character{78}}%
\TextAt{21}{337.033}{473.827}{L}%
\TextAt{24}{342.792}{472.72}{1}%
\TextAt{21}{346.943}{473.827}{L}%
\TextAt{24}{352.702}{472.72}{2}%
\TextAt{9}{360.672}{475.62}{\Character{61}}%
\TextAt{16}{373.098}{475.62}{\Character{78}}%
\TextAt{21}{382.907}{473.827}{L}%
\TextAt{24}{388.666}{472.72}{2}%
\NativeArrowAt{232.521}{455.585}{11.955}{21.926}{\Longleftrightarrow}%
\TextAt{9}{257.76}{455.585}{\Character{91}}%
\TextAt{13}{261.011}{455.585}{L}%
\TextAt{17}{268.976}{453.792}{1}%
\TextAt{13}{273.708}{455.585}{L}%
\TextAt{17}{281.673}{453.792}{2}%
\TextAt{9}{289.726}{455.585}{\Character{58}}%
\TextAt{13}{296.298}{455.585}{K}%
\TextAt{9}{307.116}{455.585}{\Character{93}}%
\TextAt{9}{313.667}{455.585}{\Character{61}}%
\TextAt{9}{326.089}{455.585}{\Character{91}}%
\TextAt{13}{329.366}{455.585}{L}%
\TextAt{17}{337.33}{453.792}{2}%
\TextAt{9}{345.383}{455.585}{\Character{58}}%
\TextAt{13}{351.956}{455.585}{K}%
\TextAt{9}{362.773}{455.585}{\Character{93}}%
\NativeArrowAt{369.346}{455.585}{11.955}{21.926}{\Longleftrightarrow}%
\TextAt{13}{394.585}{455.585}{L}%
\TextAt{17}{402.549}{453.792}{1}%
\NativeMathAt{410.602}{455.585}{11.955}{\subseteq}%
\TextAt{13}{423.222}{455.585}{L}%
\TextAt{17}{431.186}{453.792}{2}%
\TextAt{13}{435.918}{455.585}{\Character{58}}%
\TextAt{8}{79.226}{429.598}{This}%
\TextAt{8}{105.946}{429.598}{shows}%
\TextAt{8}{139.564}{429.598}{the}%
\TextAt{8}{159.72}{429.598}{injectivity}%
\TextAt{8}{215.001}{429.598}{of}%
\TextAt{8}{228.331}{429.598}{the}%
\TextAt{8}{248.487}{429.598}{correspondence}%
\TextAt{13}{330.165}{429.598}{L}%
\NativeArrowAt{341.45}{429.598}{11.955}{11.955}{\mapsto}%
\TextAt{16}{356.717}{429.598}{\Character{78}}%
\TextAt{21}{366.536}{427.805}{L}%
\TextAt{8}{372.793}{429.598}{.}%
\TextAt{8}{97.159}{409.563}{If}%
\TextAt{16}{109.353}{409.563}{\Character{78}}%
\TextAt{8}{125.323}{409.563}{is}%
\TextAt{8}{137.588}{409.563}{any}%
\TextAt{8}{160.207}{409.563}{open}%
\TextAt{8}{188.995}{409.563}{subgroup,}%
\TextAt{8}{243.671}{409.563}{then}%
\TextAt{8}{270.833}{409.563}{it}%
\TextAt{8}{283.039}{409.563}{contains}%
\TextAt{8}{329.46}{409.563}{the}%
\TextAt{8}{350.119}{409.563}{norm}%
\TextAt{8}{381.19}{409.563}{group}%
\TextAt{16}{414.867}{409.563}{\Character{78}}%
\TextAt{21}{424.676}{407.77}{L}%
\TextAt{9}{435.109}{409.563}{\Character{61}}%
\TextAt{13}{448.39}{409.563}{N}%
\TextAt{21}{457.763}{407.703}{L}%
\TextAt{3}{463.522}{407.703}{\Character{106}}%
\TextAt{21}{465.875}{407.703}{K}%
\TextAt{13}{474.077}{409.563}{A}%
\TextAt{21}{482.853}{407.77}{L}%
\TextAt{8}{493.509}{409.563}{of}%
\TextAt{8}{507.341}{409.563}{some}%
\TextAt{8}{79.226}{392.517}{Galois}%
\TextAt{8}{115.767}{392.517}{extension}%
\TextAt{13}{168.161}{392.517}{L}%
\TextAt{16}{180.516}{392.517}{\Character{106}}%
\TextAt{13}{188.227}{392.517}{K}%
\TextAt{8}{199.045}{392.517}{.}%
\TextAt{8}{209.376}{392.517}{(6.3)}%
\TextAt{8}{237.979}{392.517}{implies}%
\TextAt{8}{278.334}{392.517}{that}%
\TextAt{16}{304.33}{392.517}{\Character{78}}%
\TextAt{21}{314.14}{390.724}{L}%
\TextAt{9}{324.787}{392.517}{\Character{61}}%
\TextAt{16}{338.282}{392.517}{\Character{78}}%
\TextAt{21}{348.091}{390.217}{L}%
\TextAt{40}{353.851}{392.487}{ab}%
\TextAt{8}{362.54}{392.517}{,}%
\TextAt{8}{370.474}{392.517}{so}%
\TextAt{8}{385.472}{392.517}{we}%
\TextAt{8}{403.327}{392.517}{may}%
\TextAt{8}{429.323}{392.517}{assume}%
\TextAt{13}{470.395}{392.517}{L}%
\TextAt{16}{482.75}{392.517}{\Character{106}}%
\TextAt{13}{490.462}{392.517}{K}%
\TextAt{8}{505.805}{392.517}{to}%
\TextAt{8}{520.743}{392.517}{be}%
\TextAt{8}{79.226}{375.471}{abelian.}%
\TextAt{8}{127.833}{375.471}{But}%
\TextAt{9}{152.21}{375.471}{\Character{40}}%
\TextAt{16}{156.763}{375.471}{\Character{78}}%
\TextAt{13}{168.334}{375.471}{\Character{59}}%
\TextAt{13}{173.558}{375.471}{L}%
\TextAt{16}{186.799}{375.471}{\Character{106}}%
\TextAt{13}{195.377}{375.471}{K}%
\TextAt{9}{206.195}{375.471}{\Character{41}}%
\TextAt{9}{215.986}{375.471}{\Character{61}}%
\TextAt{13}{230.366}{375.471}{G}%
\TextAt{9}{239.6}{375.471}{\Character{40}}%
\TextAt{13}{244.152}{375.471}{L}%
\TextAt{16}{257.374}{375.471}{\Character{106}}%
\TextAt{13}{265.952}{375.471}{L}%
\TextAt{3}{273.917}{379.81}{\Character{48}}%
\TextAt{9}{276.712}{375.471}{\Character{41}}%
\TextAt{8}{286.298}{375.471}{for}%
\TextAt{8}{305.319}{375.471}{some}%
\TextAt{8}{335.781}{375.471}{intermediate}%
\TextAt{8}{404.882}{375.471}{field}%
\TextAt{13}{431.376}{375.471}{L}%
\TextAt{3}{439.34}{379.81}{\Character{48}}%
\TextAt{8}{447.169}{375.471}{of}%
\TextAt{13}{461.635}{375.471}{L}%
\TextAt{16}{474.857}{375.471}{\Character{106}}%
\TextAt{13}{483.435}{375.471}{K}%
\TextAt{8}{494.252}{375.471}{.}%
\TextAt{8}{506.109}{375.471}{Since}%
\TextAt{16}{79.226}{358.425}{\Character{78}}%
\NativeMathAt{94.11}{358.425}{11.955}{\supseteq}%
\TextAt{16}{106.711}{358.425}{\Character{78}}%
\TextAt{21}{116.546}{356.632}{L}%
\TextAt{8}{122.803}{358.425}{,}%
\TextAt{8}{129.884}{358.425}{the}%
\TextAt{8}{149.954}{358.425}{group}%
\TextAt{16}{183.044}{358.425}{\Character{78}}%
\TextAt{8}{198.426}{358.425}{is}%
\TextAt{8}{210.103}{358.425}{the}%
\TextAt{8}{230.174}{358.425}{full}%
\TextAt{8}{250.567}{358.425}{preimage}%
\TextAt{8}{300.561}{358.425}{of}%
\TextAt{13}{313.805}{358.425}{G}%
\TextAt{9}{323.039}{358.425}{\Character{40}}%
\TextAt{13}{327.591}{358.425}{L}%
\TextAt{16}{338.876}{358.425}{\Character{106}}%
\TextAt{13}{345.518}{358.425}{L}%
\TextAt{3}{353.482}{362.764}{\Character{48}}%
\TextAt{9}{356.277}{358.425}{\Character{41}}%
\TextAt{8}{364.641}{358.425}{under}%
\TextAt{8}{397.719}{358.425}{the}%
\TextAt{8}{417.789}{358.425}{map}%
\TextAt{9}{443.717}{358.425}{\Character{40}}%
\TextAt{13}{452.081}{358.425}{\Character{59}}%
\TextAt{13}{457.305}{358.425}{L}%
\TextAt{16}{468.61}{358.425}{\Character{106}}%
\TextAt{13}{475.252}{358.425}{K}%
\TextAt{9}{486.07}{358.425}{\Character{41}}%
\TextAt{9}{493.925}{358.425}{\Character{58}}%
\TextAt{13}{500.515}{358.425}{A}%
\TextAt{21}{509.291}{356.632}{K}%
\NativeMathAt{520.814}{358.425}{11.955}{\rightarrow}%
\TextAt{13}{79.226}{341.379}{G}%
\TextAt{9}{88.46}{341.379}{\Character{40}}%
\TextAt{13}{93.012}{341.379}{L}%
\TextAt{16}{104.609}{341.379}{\Character{106}}%
\TextAt{13}{111.563}{341.379}{K}%
\TextAt{9}{122.38}{341.379}{\Character{41}}%
\TextAt{8}{126.933}{341.379}{,}%
\TextAt{8}{134.31}{341.379}{and}%
\TextAt{8}{157.256}{341.379}{thus}%
\TextAt{8}{183.19}{341.379}{it}%
\TextAt{8}{195.077}{341.379}{is}%
\TextAt{8}{207.024}{341.379}{the}%
\TextAt{8}{227.364}{341.379}{full}%
\TextAt{8}{248.026}{341.379}{kernel}%
\TextAt{8}{282.675}{341.379}{of}%
\TextAt{9}{296.188}{341.379}{\Character{40}}%
\TextAt{13}{304.821}{341.379}{\Character{59}}%
\TextAt{13}{310.045}{341.379}{L}%
\TextAt{3}{318.03}{345.718}{\Character{48}}%
\TextAt{16}{324.457}{341.379}{\Character{106}}%
\TextAt{13}{331.411}{341.379}{K}%
\TextAt{9}{342.229}{341.379}{\Character{41}}%
\TextAt{9}{350.394}{341.379}{\Character{58}}%
\TextAt{13}{357.298}{341.379}{A}%
\TextAt{21}{366.073}{339.586}{K}%
\NativeMathAt{377.908}{341.379}{11.955}{\rightarrow}%
\TextAt{13}{393.496}{341.379}{G}%
\TextAt{9}{402.73}{341.379}{\Character{40}}%
\TextAt{13}{407.282}{341.379}{L}%
\TextAt{3}{415.246}{345.718}{\Character{48}}%
\TextAt{16}{421.674}{341.379}{\Character{106}}%
\TextAt{13}{428.628}{341.379}{K}%
\TextAt{9}{439.445}{341.379}{\Character{41}}%
\TextAt{8}{443.998}{341.379}{.}%
\TextAt{8}{452.995}{341.379}{Thus}%
\TextAt{16}{482.826}{341.379}{\Character{78}}%
\TextAt{9}{498.03}{341.379}{\Character{61}}%
\TextAt{16}{510.767}{341.379}{\Character{78}}%
\TextAt{21}{520.576}{339.586}{L}%
\TextAt{23}{526.335}{341.855}{\Character{48}}%
\TextAt{8}{529.518}{341.379}{.}%
\TextAt{8}{79.226}{324.333}{This}%
\TextAt{8}{105.946}{324.333}{shows}%
\TextAt{8}{139.564}{324.333}{that}%
\TextAt{8}{164.933}{324.333}{the}%
\TextAt{8}{185.089}{324.333}{correspondence}%
\TextAt{13}{266.767}{324.333}{L}%
\NativeArrowAt{278.052}{324.333}{11.955}{11.955}{\mapsto}%
\TextAt{16}{293.319}{324.333}{\Character{78}}%
\TextAt{21}{303.137}{322.54}{L}%
\TextAt{8}{313.292}{324.333}{is}%
\TextAt{8}{325.056}{324.333}{surjective.}%
\TextAt{8}{97.159}{304.298}{Finally,}%
\TextAt{8}{139.908}{304.298}{the}%
\TextAt{8}{160.579}{304.298}{equality}%
\TextAt{16}{205.651}{304.298}{\Character{78}}%
\TextAt{21}{215.46}{302.505}{L}%
\TextAt{24}{221.219}{301.398}{1}%
\TextAt{3}{225.37}{302.505}{\Character{92}}%
\TextAt{21}{231.015}{302.505}{L}%
\TextAt{24}{236.774}{301.398}{2}%
\TextAt{9}{245.621}{304.298}{\Character{61}}%
\TextAt{16}{258.924}{304.298}{\Character{78}}%
\TextAt{21}{268.733}{302.505}{L}%
\TextAt{24}{274.492}{301.398}{1}%
\TextAt{16}{279.141}{304.298}{\Character{78}}%
\TextAt{21}{288.95}{302.505}{L}%
\TextAt{24}{294.709}{301.398}{2}%
\TextAt{8}{303.77}{304.298}{is}%
\TextAt{8}{316.049}{304.298}{obtained}%
\TextAt{8}{364.695}{304.298}{as}%
\TextAt{8}{379.579}{304.298}{follows.}%
\TextAt{13}{424.11}{304.298}{L}%
\TextAt{17}{432.074}{302.505}{1}%
\TextAt{16}{439.814}{304.298}{\Character{92}}%
\TextAt{13}{450.791}{304.298}{L}%
\TextAt{17}{458.756}{302.505}{2}%
\NativeMathAt{467.686}{304.298}{11.955}{\subseteq}%
\TextAt{13}{481.182}{304.298}{L}%
\TextAt{21}{489.147}{302.505}{i}%
\TextAt{8}{496.94}{304.298}{implies}%
\TextAt{8}{79.226}{287.252}{that}%
\TextAt{16}{104.595}{287.252}{\Character{78}}%
\TextAt{21}{114.404}{285.459}{L}%
\TextAt{24}{120.163}{284.352}{1}%
\TextAt{3}{124.314}{285.459}{\Character{92}}%
\TextAt{21}{129.96}{285.459}{L}%
\TextAt{24}{135.719}{284.352}{2}%
\NativeMathAt{143.689}{287.252}{11.955}{\supseteq}%
\TextAt{16}{156.29}{287.252}{\Character{78}}%
\TextAt{21}{166.117}{285.459}{L}%
\TextAt{22}{171.876}{284.244}{i}%
\TextAt{8}{175.536}{287.252}{,}%
\TextAt{13}{182.685}{287.252}{i}%
\TextAt{9}{190}{287.252}{\Character{61}}%
\TextAt{9}{202.41}{287.252}{1}%
\TextAt{13}{208.278}{287.252}{\Character{59}}%
\TextAt{9}{213.522}{287.252}{2}%
\TextAt{8}{219.375}{287.252}{,}%
\TextAt{8}{226.524}{287.252}{and}%
\TextAt{8}{249.287}{287.252}{thus}%
\TextAt{16}{258.298}{261.265}{\Character{78}}%
\TextAt{21}{268.107}{259.472}{L}%
\TextAt{24}{273.866}{258.365}{1}%
\TextAt{3}{278.017}{259.472}{\Character{92}}%
\TextAt{21}{283.663}{259.472}{L}%
\TextAt{24}{289.422}{258.365}{2}%
\NativeMathAt{297.392}{261.265}{11.955}{\supseteq}%
\TextAt{16}{309.993}{261.265}{\Character{78}}%
\TextAt{21}{319.82}{259.472}{L}%
\TextAt{24}{325.579}{258.365}{1}%
\TextAt{16}{330.229}{261.265}{\Character{78}}%
\TextAt{21}{340.038}{259.472}{L}%
\TextAt{24}{345.797}{258.365}{2}%
\TextAt{13}{350.446}{261.265}{\Character{58}}%
\TextAt{8}{79.226}{235.278}{As}%
\TextAt{16}{98.336}{235.278}{\Character{78}}%
\TextAt{21}{108.145}{233.485}{L}%
\TextAt{24}{113.904}{232.378}{1}%
\TextAt{16}{118.553}{235.278}{\Character{78}}%
\TextAt{21}{128.362}{233.485}{L}%
\TextAt{24}{134.121}{232.378}{2}%
\TextAt{8}{144.491}{235.278}{is}%
\TextAt{8}{158.077}{235.278}{open,}%
\TextAt{8}{191.893}{235.278}{we}%
\TextAt{8}{210.943}{235.278}{have}%
\TextAt{8}{239.761}{235.278}{just}%
\TextAt{8}{264.728}{235.278}{shown}%
\TextAt{8}{302.058}{235.278}{that}%
\TextAt{16}{329.249}{235.278}{\Character{78}}%
\TextAt{21}{339.058}{233.485}{L}%
\TextAt{24}{344.818}{232.378}{1}%
\TextAt{16}{349.467}{235.278}{\Character{78}}%
\TextAt{21}{359.276}{233.485}{L}%
\TextAt{24}{365.035}{232.378}{2}%
\TextAt{9}{376.111}{235.278}{\Character{61}}%
\TextAt{16}{391.642}{235.278}{\Character{78}}%
\TextAt{21}{401.451}{233.485}{L}%
\TextAt{8}{413.429}{235.278}{for}%
\TextAt{8}{433.136}{235.278}{some}%
\TextAt{8}{464.285}{235.278}{finite}%
\TextAt{8}{496.019}{235.278}{abelian}%
\TextAt{8}{79.226}{218.232}{extension}%
\TextAt{13}{130.992}{218.232}{L}%
\TextAt{16}{142.277}{218.232}{\Character{106}}%
\TextAt{13}{148.919}{218.232}{K}%
\TextAt{8}{159.737}{218.232}{.}%
\TextAt{8}{168.185}{218.232}{But}%
\TextAt{16}{191.426}{218.232}{\Character{78}}%
\TextAt{21}{201.235}{216.439}{L}%
\TextAt{22}{206.994}{215.224}{i}%
\NativeMathAt{213.975}{218.232}{11.955}{\subseteq}%
\TextAt{16}{226.576}{218.232}{\Character{78}}%
\TextAt{21}{236.403}{216.439}{L}%
\TextAt{8}{246.558}{218.232}{implies}%
\TextAt{13}{286.285}{218.232}{L}%
\NativeMathAt{297.57}{218.232}{11.955}{\subseteq}%
\TextAt{13}{310.19}{218.232}{L}%
\TextAt{17}{318.154}{216.439}{1}%
\TextAt{16}{325.543}{218.232}{\Character{92}}%
\TextAt{13}{336.17}{218.232}{L}%
\TextAt{17}{344.134}{216.439}{2}%
\TextAt{8}{348.867}{218.232}{,}%
\TextAt{8}{356.016}{218.232}{so}%
\TextAt{8}{370.386}{218.232}{that}%
\TextAt{16}{242.392}{192.245}{\Character{78}}%
\TextAt{21}{252.201}{190.451}{L}%
\TextAt{24}{257.96}{189.344}{1}%
\TextAt{16}{262.609}{192.244}{\Character{78}}%
\TextAt{21}{272.418}{190.451}{L}%
\TextAt{24}{278.177}{189.344}{2}%
\TextAt{9}{286.147}{192.244}{\Character{61}}%
\TextAt{16}{298.573}{192.244}{\Character{78}}%
\TextAt{21}{308.382}{190.451}{L}%
\NativeMathAt{317.96}{192.244}{11.955}{\supseteq}%
\TextAt{16}{330.561}{192.244}{\Character{78}}%
\TextAt{21}{340.388}{190.451}{L}%
\TextAt{24}{346.147}{189.344}{1}%
\TextAt{3}{350.298}{190.451}{\Character{92}}%
\TextAt{21}{355.944}{190.451}{L}%
\TextAt{24}{361.703}{189.344}{2}%
\TextAt{13}{366.352}{192.244}{\Character{58}}%
\Vector{q 0 G 0 g 0.399 w 452.334 -607.743 m 452.334 -599.673 l S Q}%
\Vector{q 0 G 0 g 0.399 w 452.534 -599.872 m 459.707 -599.872 l S Q}%
\Vector{q 0 G 0 g 0.399 w 452.534 -607.544 m 459.707 -607.544 l S Q}%
\Vector{q 0 G 0 g 0.399 w 459.906 -607.743 m 459.906 -599.673 l S Q}%
\TextAt{12}{79.226}{133.001}{Exercise}%
\TextAt{12}{132.281}{133.001}{IV.6.1.}%
\TextAt{8}{175.91}{133.001}{Let}%
\TextAt{13}{196.753}{133.001}{n}%
\TextAt{8}{207.512}{133.001}{be}%
\TextAt{8}{223.31}{133.001}{a}%
\TextAt{8}{232.938}{133.001}{natural}%
\TextAt{8}{273.794}{133.001}{number,}%
\TextAt{8}{319.541}{133.001}{and}%
\TextAt{8}{342.177}{133.001}{assume}%
\TextAt{8}{382.495}{133.001}{the}%
\TextAt{8}{402.525}{133.001}{group}%
\TextAt{13}{435.574}{133.001}{\Character{22}}%
\TextAt{21}{442.617}{131.207}{n}%
\TextAt{9}{451.574}{133.001}{\Character{61}}%
\TextAt{16}{464}{133.001}{\Character{102}}%
\TextAt{13}{469.977}{133.001}{\Character{24}}%
\NativeMathAt{478.988}{133.001}{11.955}{\in}%
\TextAt{13}{490.279}{133.001}{A}%
\TextAt{16}{502.376}{133.001}{\Character{106}}%
\TextAt{13}{509.017}{133.001}{\Character{24}}%
\TextAt{21}{514.708}{137.339}{n}%
\TextAt{9}{523.665}{133.001}{\Character{61}}%
\TextAt{9}{79.226}{115.955}{1}%
\TextAt{16}{85.079}{115.955}{\Character{103}}%
\TextAt{8}{95.253}{115.955}{is}%
\TextAt{8}{107.315}{115.955}{cyclic}%
\TextAt{8}{139.797}{115.955}{of}%
\TextAt{8}{153.426}{115.955}{order}%
\TextAt{13}{184.294}{115.955}{n}%
\TextAt{8}{191.281}{115.955}{,}%
\TextAt{8}{198.804}{115.955}{and}%
\TextAt{13}{221.865}{115.955}{A}%
\NativeMathAt{234.47}{115.955}{11.955}{\subseteq}%
\TextAt{13}{247.598}{115.955}{A}%
\TextAt{21}{256.374}{120.293}{n}%
\TextAt{8}{262.01}{115.955}{.}%
\TextAt{8}{271.354}{115.955}{Let}%
\TextAt{13}{292.622}{115.955}{K}%
\TextAt{8}{307.636}{115.955}{be}%
\TextAt{8}{323.859}{115.955}{a}%
\TextAt{8}{333.913}{115.955}{field}%
\TextAt{8}{359.568}{115.955}{such}%
\TextAt{8}{386.264}{115.955}{that}%
\TextAt{13}{411.931}{115.955}{\Character{22}}%
\TextAt{21}{418.974}{114.161}{n}%
\NativeMathAt{428.44}{115.955}{11.955}{\subseteq}%
\TextAt{13}{441.569}{115.955}{A}%
\TextAt{21}{450.344}{114.161}{K}%
\TextAt{8}{462.743}{115.955}{and}%
\TextAt{8}{485.804}{115.955}{let}%
\TextAt{13}{503.007}{115.955}{L}%
\TextAt{16}{514.801}{115.955}{\Character{106}}%
\TextAt{13}{521.952}{115.955}{K}%
\TextAt{8}{79.226}{98.909}{be}%
\TextAt{8}{95.15}{98.909}{the}%
\TextAt{8}{115.307}{98.909}{maximal}%
\TextAt{8}{163.116}{98.909}{abelian}%
\TextAt{8}{203.763}{98.909}{extension}%
\TextAt{8}{255.529}{98.909}{of}%
\TextAt{8}{268.859}{98.909}{exponent}%
\TextAt{13}{319.262}{98.909}{n}%
\TextAt{8}{326.25}{98.909}{.}%
\TextAt{8}{334.698}{98.909}{If}%
\TextAt{13}{346.39}{98.909}{L}%
\TextAt{16}{357.675}{98.909}{\Character{106}}%
\TextAt{13}{364.317}{98.909}{K}%
\TextAt{8}{379.032}{98.909}{is}%
\TextAt{8}{390.796}{98.909}{finite,}%
\TextAt{8}{423.96}{98.909}{then}%
\TextAt{8}{450.62}{98.909}{one}%
\TextAt{8}{472.079}{98.909}{has}%
\TextAt{13}{267.65}{72.921}{N}%
\TextAt{21}{277.023}{71.062}{L}%
\TextAt{3}{282.782}{71.062}{\Character{106}}%
\TextAt{21}{285.135}{71.062}{K}%
\TextAt{13}{293.338}{72.921}{A}%
\TextAt{21}{302.113}{71.128}{L}%
\TextAt{9}{311.691}{72.921}{\Character{61}}%
\TextAt{13}{324.116}{72.921}{A}%
\TextAt{21}{332.892}{77.858}{n}%
\TextAt{21}{332.892}{69.966}{K}%
\TextAt{13}{341.094}{72.921}{\Character{58}}%
\TextAt{8}{297.211}{43.033}{279}%
\Destination{78.23}{809.815}{nav.page.279}%
\Destination{79.226}{146.955}{nav.ex.4.6.1}%
\Destination{79.36}{584.755}{nav.print.305}%
\Destination{79.36}{775.155}{nav.number.4.6.7}%
\Hotlink{143.768}{612.755}{25.078}{11.047}{nav.number.4.6.3}%
\Hotlink{123.869}{567.555}{26.106}{11.047}{nav.number.4.6.7}%
\Hotlink{208.876}{390.355}{25.078}{11.047}{nav.number.4.6.3}%
}

% Page 292 (body page 280)
\Page{%
\TextAt{2}{79.226}{796.943}{IV.}%
\TextAt{2}{98.012}{796.943}{Abstract}%
\TextAt{2}{143.252}{796.943}{Class}%
\TextAt{2}{171.844}{796.943}{Field}%
\TextAt{2}{199.575}{796.943}{Theory}%
\Vector{q 0 G 0 g 0.299 w 7.226 18.729 m 460.769 18.729 l S Q}%
\TextAt{12}{79.226}{761.124}{Exercise}%
\TextAt{12}{131.582}{761.124}{IV.6.2.}%
\TextAt{8}{174.604}{761.124}{Under}%
\TextAt{8}{209.306}{761.124}{the}%
\TextAt{8}{228.728}{761.124}{hypotheses}%
\TextAt{8}{287.627}{761.124}{of}%
\TextAt{8}{300.223}{761.124}{exercise}%
\TextAt{8}{342.79}{761.124}{1,}%
\TextAt{8}{355.211}{761.124}{Kummer}%
\TextAt{8}{403.242}{761.124}{theory}%
\TextAt{8}{439.258}{761.124}{and}%
\TextAt{8}{461.287}{761.124}{class}%
\TextAt{8}{487.99}{761.124}{field}%
\TextAt{8}{512.613}{761.124}{the-}%
\TextAt{8}{79.226}{744.078}{ory}%
\TextAt{8}{99.757}{744.078}{yield,}%
\TextAt{8}{131.344}{744.078}{via}%
\TextAt{8}{150.572}{744.078}{Pontryagin}%
\TextAt{8}{210.627}{744.078}{duality}%
\TextAt{13}{250.345}{744.078}{G}%
\TextAt{9}{259.579}{744.078}{\Character{40}}%
\TextAt{13}{264.131}{744.078}{L}%
\TextAt{16}{275.484}{744.078}{\Character{106}}%
\TextAt{13}{282.193}{744.078}{K}%
\TextAt{9}{293.011}{744.078}{\Character{41}}%
\TextAt{16}{300.247}{744.078}{\Character{2}}%
\TextAt{8}{312.229}{744.078}{Hom}%
\TextAt{9}{336.618}{744.078}{\Character{40}}%
\TextAt{13}{341.17}{744.078}{G}%
\TextAt{9}{350.404}{744.078}{\Character{40}}%
\TextAt{13}{354.956}{744.078}{L}%
\TextAt{16}{366.309}{744.078}{\Character{106}}%
\TextAt{13}{373.018}{744.078}{K}%
\TextAt{9}{383.836}{744.078}{\Character{41}}%
\TextAt{13}{388.388}{744.078}{\Character{59}}%
\TextAt{13}{393.612}{744.078}{\Character{22}}%
\TextAt{21}{400.675}{742.284}{n}%
\TextAt{9}{406.312}{744.078}{\Character{41}}%
\NativeMathAt{414.252}{744.078}{11.955}{\rightarrow}%
\TextAt{13}{429.596}{744.078}{\Character{22}}%
\TextAt{21}{436.639}{742.284}{n}%
\TextAt{8}{442.275}{744.078}{,}%
\TextAt{8}{449.474}{744.078}{a}%
\TextAt{8}{459.269}{744.078}{nondegenerate}%
\TextAt{8}{79.226}{727.032}{bilinear}%
\TextAt{8}{121.5}{727.032}{mapping}%
\TextAt{8}{169.631}{727.032}{(the}%
\TextAt{8}{194.342}{727.032}{abstract}%
\TextAt{8}{239.939}{727.032}{“Hilbert}%
\TextAt{8}{285.776}{727.032}{symbol”)}%
\TextAt{9}{221.948}{700.248}{\Character{40}}%
\TextAt{13}{230.397}{700.248}{\Character{59}}%
\TextAt{9}{239.539}{700.248}{\Character{41}}%
\TextAt{9}{247.394}{700.248}{\Character{58}}%
\TextAt{13}{253.984}{700.248}{A}%
\TextAt{21}{262.76}{698.455}{K}%
\TextAt{9}{270.962}{700.248}{\Character{47}}%
\TextAt{13}{276.815}{700.248}{A}%
\TextAt{21}{285.591}{705.184}{n}%
\TextAt{21}{285.591}{697.292}{K}%
\TextAt{16}{296.45}{700.248}{\Character{2}}%
\TextAt{13}{308.405}{700.248}{A}%
\TextAt{21}{317.181}{698.455}{K}%
\TextAt{9}{325.383}{700.248}{\Character{47}}%
\TextAt{13}{331.236}{700.248}{A}%
\TextAt{21}{340.012}{705.184}{n}%
\TextAt{21}{340.012}{697.292}{K}%
\NativeArrowAt{351.535}{700.248}{11.955}{19.26}{\longrightarrow}%
\TextAt{13}{374.117}{700.248}{\Character{22}}%
\TextAt{21}{381.16}{698.455}{n}%
\TextAt{13}{386.796}{700.248}{\Character{58}}%
\TextAt{12}{79.226}{659.787}{Exercise}%
\TextAt{12}{132.164}{659.787}{IV.6.3.}%
\TextAt{8}{175.692}{659.787}{Let}%
\TextAt{13}{196.434}{659.787}{p}%
\TextAt{8}{205.978}{659.787}{be}%
\TextAt{8}{221.675}{659.787}{a}%
\TextAt{8}{231.202}{659.787}{prime}%
\TextAt{8}{264.138}{659.787}{number}%
\TextAt{8}{306.506}{659.787}{and}%
\TextAt{9}{329.041}{659.787}{\Character{40}}%
\TextAt{13}{333.593}{659.787}{d}%
\TextAt{9}{342.996}{659.787}{\Character{58}}%
\TextAt{13}{349.569}{659.787}{G}%
\NativeMathAt{362.123}{659.787}{11.955}{\rightarrow}%
\TextAt{18}{377.399}{659.787}{Z}%
\TextAt{21}{385.37}{657.994}{p}%
\TextAt{13}{390.13}{659.787}{\Character{59}}%
\TextAt{13}{395.354}{659.787}{v}%
\TextAt{9}{404.784}{659.787}{\Character{58}}%
\TextAt{13}{411.356}{659.787}{A}%
\TextAt{21}{420.132}{657.994}{k}%
\NativeMathAt{428.572}{659.787}{11.955}{\rightarrow}%
\TextAt{18}{443.848}{659.787}{Z}%
\TextAt{21}{451.818}{657.994}{p}%
\TextAt{9}{456.579}{659.787}{\Character{41}}%
\TextAt{8}{464.801}{659.787}{a}%
\TextAt{13}{474.328}{659.787}{p}%
\TextAt{8}{480.203}{659.787}{-class}%
\TextAt{8}{511.31}{659.787}{field}%
\TextAt{8}{79.226}{642.741}{theory}%
\TextAt{8}{114.809}{642.741}{in}%
\TextAt{8}{127.295}{642.741}{the}%
\TextAt{8}{146.284}{642.741}{sense}%
\TextAt{8}{175.148}{642.741}{of}%
\TextAt{8}{187.311}{642.741}{§5,}%
\TextAt{8}{205.063}{642.741}{exercise}%
\TextAt{8}{247.198}{642.741}{2.}%
\TextAt{8}{261.115}{642.741}{Let}%
\TextAt{13}{280.917}{642.741}{d}%
\TextAt{3}{287}{647.08}{\Character{48}}%
\TextAt{9}{293.116}{642.741}{\Character{58}}%
\TextAt{13}{299.688}{642.741}{G}%
\NativeMathAt{312.243}{642.741}{11.955}{\rightarrow}%
\TextAt{18}{327.519}{642.741}{Z}%
\TextAt{21}{335.489}{640.948}{p}%
\TextAt{8}{342.98}{642.741}{be}%
\TextAt{8}{357.737}{642.741}{another}%
\TextAt{8}{399.501}{642.741}{surjective}%
\TextAt{8}{450.745}{642.741}{homomorphism,}%
\TextAt{8}{79.226}{625.695}{and}%
\TextAt{20}{104.727}{628.717}{\Character{101}}%
\TextAt{13}{101.989}{625.695}{K}%
\TextAt{3}{112.806}{630.034}{\Character{48}}%
\TextAt{16}{118.922}{625.695}{\Character{106}}%
\TextAt{13}{125.564}{625.695}{K}%
\TextAt{8}{140.279}{625.695}{the}%
\TextAt{18}{160.436}{625.695}{Z}%
\TextAt{21}{168.406}{623.902}{p}%
\TextAt{8}{173.167}{625.695}{-extension}%
\TextAt{8}{228.83}{625.695}{defined}%
\TextAt{8}{269.143}{625.695}{by}%
\TextAt{13}{285.402}{625.695}{d}%
\TextAt{3}{291.484}{630.034}{\Character{48}}%
\TextAt{8}{294.28}{625.695}{.}%
\TextAt{8}{302.728}{625.695}{Let}%
\TextAt{13}{323.697}{625.695}{v}%
\TextAt{3}{329.785}{630.034}{\Character{48}}%
\TextAt{9}{335.901}{625.695}{\Character{58}}%
\TextAt{13}{342.474}{625.695}{A}%
\TextAt{21}{351.249}{623.902}{k}%
\NativeMathAt{359.69}{625.695}{11.955}{\rightarrow}%
\TextAt{18}{374.966}{625.695}{Z}%
\TextAt{21}{382.936}{623.902}{p}%
\TextAt{8}{391.594}{625.695}{be}%
\TextAt{8}{407.519}{625.695}{the}%
\TextAt{8}{427.675}{625.695}{composite}%
\TextAt{8}{482.693}{625.695}{of}%
\TextAt{13}{231.303}{594.945}{A}%
\TextAt{21}{240.079}{593.152}{k}%
\TextAt{17}{251.107}{602.652}{\Character{40}}%
\TextAt{21}{258.297}{602.652}{\Character{59}}%
\TextAt{30}{262.62}{604.666}{\Character{101}}%
\TextAt{21}{260.65}{602.652}{K}%
\TextAt{23}{268.354}{605.465}{\Character{48}}%
\TextAt{3}{271.039}{602.652}{\Character{106}}%
\TextAt{21}{273.391}{602.652}{K}%
\TextAt{17}{281.095}{602.652}{\Character{41}}%
\NativeArrowAt{248.519}{594.945}{11.955}{39.823}{\longrightarrow}%
\TextAt{13}{291.709}{594.945}{G}%
\TextAt{9}{300.942}{594.945}{\Character{40}}%
\TextAt{20}{308.233}{597.967}{\Character{101}}%
\TextAt{13}{305.494}{594.945}{K}%
\TextAt{3}{316.312}{599.881}{\Character{48}}%
\TextAt{16}{322.428}{594.945}{\Character{106}}%
\TextAt{13}{329.07}{594.945}{K}%
\TextAt{9}{339.887}{594.945}{\Character{41}}%
\TextAt{21}{350.348}{601.722}{d}%
\TextAt{23}{354.705}{604.535}{\Character{48}}%
\NativeArrowAt{347.761}{594.945}{11.955}{13.629}{\longrightarrow}%
\TextAt{18}{364.71}{594.945}{Z}%
\TextAt{21}{372.68}{593.152}{p}%
\TextAt{13}{377.441}{594.945}{\Character{58}}%
\TextAt{8}{79.226}{568.161}{Then}%
\TextAt{9}{110.002}{568.161}{\Character{40}}%
\TextAt{13}{114.555}{568.161}{d}%
\TextAt{3}{120.637}{572.499}{\Character{48}}%
\TextAt{13}{123.432}{568.161}{\Character{59}}%
\TextAt{13}{128.656}{568.161}{v}%
\TextAt{3}{134.765}{572.499}{\Character{48}}%
\TextAt{9}{137.56}{568.161}{\Character{41}}%
\TextAt{8}{146.228}{568.161}{is}%
\TextAt{8}{158.211}{568.161}{also}%
\TextAt{8}{181.91}{568.161}{a}%
\TextAt{13}{191.884}{568.161}{p}%
\TextAt{8}{197.759}{568.161}{-class}%
\TextAt{8}{229.312}{568.161}{field}%
\TextAt{8}{254.888}{568.161}{theory.}%
\TextAt{8}{295.865}{568.161}{The}%
\TextAt{8}{320.138}{568.161}{norm}%
\TextAt{8}{350.926}{568.161}{residue}%
\TextAt{8}{390.872}{568.161}{symbols}%
\TextAt{8}{435.767}{568.161}{associated}%
\TextAt{8}{491.721}{568.161}{to}%
\TextAt{9}{506.25}{568.161}{\Character{40}}%
\TextAt{13}{510.802}{568.161}{d\Character{59}}%
\TextAt{13}{522.1}{568.161}{v}%
\TextAt{9}{528.217}{568.161}{\Character{41}}%
\TextAt{8}{79.226}{551.115}{and}%
\TextAt{9}{102.449}{551.115}{\Character{40}}%
\TextAt{13}{107.001}{551.115}{d}%
\TextAt{3}{113.084}{555.453}{\Character{48}}%
\TextAt{13}{115.879}{551.115}{\Character{59}}%
\TextAt{13}{121.103}{551.115}{v}%
\TextAt{3}{127.211}{555.453}{\Character{48}}%
\TextAt{9}{130.006}{551.115}{\Character{41}}%
\TextAt{8}{138.916}{551.115}{coincide.}%
\TextAt{8}{189.713}{551.115}{(No}%
\TextAt{8}{213.259}{551.115}{analogous}%
\TextAt{8}{267.78}{551.115}{statement}%
\TextAt{8}{322.611}{551.115}{holds}%
\TextAt{8}{353.7}{551.115}{in}%
\TextAt{8}{367.813}{551.115}{the}%
\TextAt{8}{388.429}{551.115}{case}%
\TextAt{8}{413.66}{551.115}{of}%
\TextAt{20}{428.114}{554.203}{\Character{98}}%
\TextAt{18}{427.449}{551.115}{Z}%
\TextAt{8}{435.42}{551.115}{-class}%
\TextAt{8}{467.214}{551.115}{field}%
\TextAt{8}{493.031}{551.115}{theories}%
\TextAt{9}{79.226}{534.069}{\Character{40}}%
\TextAt{13}{83.778}{534.069}{d}%
\TextAt{9}{93.182}{534.069}{\Character{58}}%
\TextAt{13}{99.754}{534.069}{G}%
\NativeMathAt{112.309}{534.069}{11.955}{\rightarrow}%
\TextAt{20}{128.249}{537.158}{\Character{98}}%
\TextAt{18}{127.585}{534.069}{Z}%
\TextAt{13}{135.555}{534.069}{\Character{59}}%
\TextAt{13}{140.779}{534.069}{v}%
\TextAt{9}{150.208}{534.069}{\Character{58}}%
\TextAt{13}{156.781}{534.069}{A}%
\TextAt{21}{165.556}{532.276}{k}%
\NativeMathAt{173.997}{534.069}{11.955}{\rightarrow}%
\TextAt{20}{189.937}{537.158}{\Character{98}}%
\TextAt{18}{189.273}{534.069}{Z}%
\TextAt{9}{197.243}{534.069}{\Character{41}}%
\TextAt{8}{201.795}{534.069}{.)}%
\TextAt{12}{79.226}{503.346}{Exercise}%
\TextAt{12}{132.977}{503.346}{IV.6.4.}%
\TextAt{8}{177.212}{503.346}{(Generalization}%
\TextAt{8}{260.372}{503.346}{to}%
\TextAt{8}{275.161}{503.346}{infinite}%
\TextAt{8}{315.307}{503.346}{extensions.)}%
\TextAt{8}{382.229}{503.346}{Let}%
\TextAt{9}{403.677}{503.346}{\Character{40}}%
\TextAt{13}{408.229}{503.346}{d}%
\TextAt{9}{418.448}{503.346}{\Character{58}}%
\TextAt{13}{425.836}{503.346}{G}%
\NativeMathAt{439.206}{503.346}{11.955}{\rightarrow}%
\TextAt{20}{455.962}{506.435}{\Character{98}}%
\TextAt{18}{455.297}{503.346}{Z}%
\TextAt{13}{463.267}{503.346}{\Character{59}}%
\TextAt{13}{468.491}{503.346}{v}%
\TextAt{9}{478.736}{503.346}{\Character{58}}%
\TextAt{13}{486.124}{503.346}{A}%
\TextAt{21}{494.899}{501.553}{k}%
\NativeMathAt{504.156}{503.346}{11.955}{\rightarrow}%
\TextAt{20}{520.911}{506.435}{\Character{98}}%
\TextAt{18}{520.247}{503.346}{Z}%
\TextAt{9}{528.217}{503.346}{\Character{41}}%
\TextAt{8}{79.226}{486.3}{be}%
\TextAt{8}{95.865}{486.3}{a}%
\TextAt{8}{106.334}{486.3}{class}%
\TextAt{8}{134.486}{486.3}{field}%
\TextAt{8}{160.557}{486.3}{theory.}%
\TextAt{8}{203.02}{486.3}{We}%
\TextAt{8}{223.879}{486.3}{assume}%
\TextAt{8}{265.037}{486.3}{that}%
\TextAt{8}{291.121}{486.3}{the}%
\TextAt{8}{311.991}{486.3}{kernel}%
\TextAt{13}{347.172}{486.3}{U}%
\TextAt{21}{355.122}{484.507}{K}%
\TextAt{8}{367.937}{486.3}{of}%
\TextAt{13}{381.981}{486.3}{v}%
\TextAt{21}{387.64}{484.507}{K}%
\TextAt{9}{400.381}{486.3}{\Character{58}}%
\TextAt{13}{408.17}{486.3}{A}%
\TextAt{21}{416.946}{484.507}{K}%
\NativeMathAt{429.686}{486.3}{11.955}{\rightarrow}%
\TextAt{20}{446.844}{489.389}{\Character{98}}%
\TextAt{18}{446.18}{486.3}{Z}%
\TextAt{8}{458.761}{486.3}{is}%
\TextAt{8}{471.239}{486.3}{compact}%
\TextAt{8}{518.782}{486.3}{for}%
\TextAt{8}{79.226}{469.254}{every}%
\TextAt{8}{110.118}{469.254}{finite}%
\TextAt{8}{140.03}{469.254}{extension}%
\TextAt{13}{191.796}{469.254}{K}%
\TextAt{16}{205.935}{469.254}{\Character{106}}%
\TextAt{13}{212.576}{469.254}{k}%
\TextAt{8}{219.066}{469.254}{.}%
\TextAt{8}{227.514}{469.254}{For}%
\TextAt{8}{248.483}{469.254}{an}%
\TextAt{8}{264.742}{469.254}{infinite}%
\TextAt{8}{304.41}{469.254}{extension}%
\TextAt{13}{356.176}{469.254}{K}%
\TextAt{16}{370.314}{469.254}{\Character{106}}%
\TextAt{13}{376.956}{469.254}{k}%
\TextAt{8}{383.445}{469.254}{,}%
\TextAt{8}{390.595}{469.254}{put}%
\TextAt{20}{270.716}{445.492}{\Character{98}}%
\TextAt{13}{268.023}{442.47}{A}%
\TextAt{21}{276.798}{440.677}{K}%
\TextAt{9}{288.322}{442.47}{\Character{61}}%
\NativeMathAt{300.747}{442.47}{11.955}{\varprojlim}%
\TextAt{21}{306.157}{430.356}{\Character{11}}%
\TextAt{13}{318.999}{442.47}{A}%
\TextAt{21}{327.774}{440.677}{K}%
\TextAt{22}{334.916}{439.681}{\Character{11}}%
\TextAt{13}{340.721}{442.47}{\Character{59}}%
\TextAt{8}{79.226}{409.46}{where}%
\TextAt{13}{112.467}{409.46}{K}%
\TextAt{21}{122.426}{407.667}{\Character{11}}%
\TextAt{16}{131.684}{409.46}{\Character{106}}%
\TextAt{13}{138.326}{409.46}{k}%
\TextAt{8}{148.144}{409.46}{varies}%
\TextAt{8}{180.488}{409.46}{over}%
\TextAt{8}{204.965}{409.46}{the}%
\TextAt{8}{224.553}{409.46}{finite}%
\TextAt{8}{253.897}{409.46}{subextensions}%
\TextAt{8}{327.653}{409.46}{of}%
\TextAt{13}{340.415}{409.46}{K}%
\TextAt{16}{354.553}{409.46}{\Character{106}}%
\TextAt{13}{361.195}{409.46}{k}%
\TextAt{8}{371.013}{409.46}{and}%
\TextAt{8}{393.208}{409.46}{the}%
\TextAt{8}{412.795}{409.46}{projective}%
\TextAt{8}{466.204}{409.46}{limit}%
\TextAt{8}{493.599}{409.46}{is}%
\TextAt{8}{504.794}{409.46}{taken}%
\TextAt{8}{79.226}{392.414}{with}%
\TextAt{8}{105.886}{392.414}{respect}%
\TextAt{8}{145.936}{392.414}{to}%
\TextAt{8}{160.246}{392.414}{the}%
\TextAt{8}{180.403}{392.414}{norm}%
\TextAt{8}{210.972}{392.414}{maps}%
\TextAt{13}{241.601}{392.414}{N}%
\TextAt{21}{250.975}{390.554}{K}%
\TextAt{22}{258.116}{389.148}{\Character{12}}%
\TextAt{3}{263.094}{390.554}{\Character{106}}%
\TextAt{21}{265.447}{390.554}{K}%
\TextAt{22}{272.588}{389.558}{\Character{11}}%
\TextAt{9}{281.714}{392.414}{\Character{58}}%
\TextAt{13}{288.287}{392.414}{A}%
\TextAt{21}{297.062}{390.621}{K}%
\TextAt{22}{304.203}{389.215}{\Character{12}}%
\NativeMathAt{313}{392.414}{11.955}{\rightarrow}%
\TextAt{13}{328.276}{392.414}{A}%
\TextAt{21}{337.051}{390.621}{K}%
\TextAt{22}{344.193}{389.625}{\Character{11}}%
\TextAt{8}{349.998}{392.414}{.}%
\TextAt{8}{358.447}{392.414}{Show:}%
\TextAt{8}{91.528}{372.379}{1)}%
\TextAt{8}{107.918}{372.379}{For}%
\TextAt{8}{128.888}{372.379}{every}%
\TextAt{8}{159.78}{372.379}{(finite}%
\TextAt{8}{194.247}{372.379}{or}%
\TextAt{8}{208.557}{372.379}{infinite)}%
\TextAt{8}{252.779}{372.379}{extension}%
\TextAt{13}{304.545}{372.379}{L}%
\TextAt{16}{315.83}{372.379}{\Character{106}}%
\TextAt{13}{322.472}{372.379}{K}%
\TextAt{8}{333.29}{372.379}{,}%
\TextAt{8}{340.439}{372.379}{one}%
\TextAt{8}{361.899}{372.379}{has}%
\TextAt{8}{382.772}{372.379}{a}%
\TextAt{8}{392.528}{372.379}{norm}%
\TextAt{8}{423.097}{372.379}{map}%
\TextAt{13}{271.971}{345.595}{N}%
\TextAt{21}{281.345}{343.736}{L}%
\TextAt{3}{287.104}{343.736}{\Character{106}}%
\TextAt{21}{289.456}{343.736}{K}%
\TextAt{9}{300.98}{345.595}{\Character{58}}%
\TextAt{20}{310.245}{348.617}{\Character{98}}%
\TextAt{13}{307.552}{345.595}{A}%
\TextAt{21}{316.327}{343.802}{L}%
\NativeArrowAt{325.905}{345.595}{11.955}{19.26}{\longrightarrow}%
\TextAt{20}{351.18}{348.617}{\Character{98}}%
\TextAt{13}{348.487}{345.595}{A}%
\TextAt{21}{357.263}{343.802}{K}%
\TextAt{13}{365.465}{345.595}{\Character{59}}%
\TextAt{8}{107.918}{318.812}{and}%
\TextAt{8}{130.703}{318.812}{if}%
\TextAt{13}{141.449}{318.812}{L}%
\TextAt{16}{152.772}{318.812}{\Character{106}}%
\TextAt{13}{159.451}{318.812}{K}%
\TextAt{8}{174.188}{318.812}{is}%
\TextAt{8}{185.973}{318.812}{finite,}%
\TextAt{8}{219.165}{318.812}{there}%
\TextAt{8}{249.098}{318.812}{is}%
\TextAt{8}{260.884}{318.812}{an}%
\TextAt{8}{277.165}{318.812}{injection}%
\TextAt{13}{324.984}{318.812}{i}%
\TextAt{21}{328.977}{316.952}{L}%
\TextAt{3}{334.736}{316.952}{\Character{106}}%
\TextAt{21}{337.089}{316.952}{K}%
\TextAt{9}{348.649}{318.812}{\Character{58}}%
\TextAt{20}{357.952}{321.834}{\Character{98}}%
\TextAt{13}{355.259}{318.812}{A}%
\TextAt{21}{364.035}{317.018}{K}%
\NativeArrowAt{375.596}{318.812}{11.955}{13.214}{\hookrightarrow}%
\TextAt{20}{394.861}{321.834}{\Character{98}}%
\TextAt{13}{392.168}{318.812}{A}%
\TextAt{21}{400.943}{317.018}{L}%
\TextAt{8}{407.201}{318.812}{.}%
\TextAt{8}{415.715}{318.812}{If}%
\TextAt{8}{427.429}{318.812}{furthermore}%
\TextAt{13}{492.164}{318.812}{L}%
\TextAt{16}{503.487}{318.812}{\Character{106}}%
\TextAt{13}{510.166}{318.812}{K}%
\TextAt{8}{524.903}{318.812}{is}%
\TextAt{8}{107.918}{301.766}{Galois,}%
\TextAt{8}{147.084}{301.766}{then}%
\TextAt{8}{173.744}{301.766}{one}%
\TextAt{8}{195.203}{301.766}{has}%
\TextAt{20}{218.77}{304.788}{\Character{98}}%
\TextAt{13}{216.077}{301.766}{A}%
\TextAt{21}{224.852}{299.972}{K}%
\TextAt{9}{236.376}{301.766}{\Character{61}}%
\TextAt{20}{251.494}{304.788}{\Character{98}}%
\TextAt{13}{248.801}{301.766}{A}%
\TextAt{21}{257.577}{307.876}{G}%
\TextAt{17}{264.175}{307.876}{\Character{40}}%
\TextAt{21}{267.468}{307.876}{L}%
\TextAt{3}{273.228}{307.876}{\Character{106}}%
\TextAt{21}{275.58}{307.876}{K}%
\TextAt{17}{283.284}{307.876}{\Character{41}}%
\TextAt{21}{257.577}{298.525}{L}%
\TextAt{8}{287.076}{301.766}{.}%
\TextAt{8}{91.528}{284.72}{2)}%
\TextAt{8}{107.918}{284.72}{For}%
\TextAt{8}{128.955}{284.72}{every}%
\TextAt{8}{159.913}{284.72}{extension}%
\TextAt{13}{211.746}{284.72}{K}%
\TextAt{16}{225.998}{284.72}{\Character{106}}%
\TextAt{13}{232.753}{284.72}{k}%
\TextAt{8}{243.207}{284.72}{with}%
\TextAt{8}{269.934}{284.72}{finite}%
\TextAt{8}{299.912}{284.72}{inertia}%
\TextAt{8}{337.052}{284.72}{degree}%
\TextAt{13}{373.534}{284.72}{f}%
\TextAt{21}{379.304}{282.926}{K}%
\TextAt{9}{390.941}{284.72}{\Character{61}}%
\TextAt{9}{403.47}{284.72}{\Character{91}}%
\TextAt{13}{406.732}{284.72}{K}%
\TextAt{16}{420.251}{284.72}{\Character{92}}%
\TextAt{20}{430.848}{287.875}{\Character{101}}%
\TextAt{13}{430.924}{284.72}{k}%
\TextAt{9}{440.848}{284.72}{\Character{58}}%
\TextAt{13}{447.534}{284.72}{k}%
\TextAt{9}{454.023}{284.72}{\Character{93}}%
\TextAt{8}{457.275}{284.72}{,}%
\TextAt{9}{464.507}{284.72}{\Character{40}}%
\TextAt{13}{469.06}{284.72}{d\Character{59}}%
\TextAt{13}{480.358}{284.72}{v}%
\TextAt{9}{486.475}{284.72}{\Character{41}}%
\TextAt{8}{494.991}{284.72}{induces}%
\TextAt{8}{107.918}{267.674}{a}%
\TextAt{8}{117.674}{267.674}{class}%
\TextAt{8}{145.111}{267.674}{field}%
\TextAt{8}{170.468}{267.674}{theory}%
\TextAt{9}{207.218}{267.674}{\Character{40}}%
\TextAt{13}{211.77}{267.674}{d}%
\TextAt{21}{217.853}{265.88}{K}%
\TextAt{9}{229.377}{267.674}{\Character{58}}%
\TextAt{13}{235.949}{267.674}{G}%
\TextAt{21}{245.183}{265.88}{K}%
\NativeMathAt{256.706}{267.674}{11.955}{\rightarrow}%
\TextAt{20}{272.646}{270.762}{\Character{98}}%
\TextAt{18}{271.982}{267.674}{Z}%
\TextAt{13}{279.952}{267.674}{\Character{59}}%
\TextAt{13}{285.176}{267.674}{v}%
\TextAt{21}{290.856}{265.88}{K}%
\TextAt{9}{302.379}{267.674}{\Character{58}}%
\TextAt{20}{311.644}{270.696}{\Character{98}}%
\TextAt{13}{308.952}{267.674}{A}%
\TextAt{21}{317.727}{265.88}{K}%
\NativeMathAt{329.251}{267.674}{11.955}{\rightarrow}%
\TextAt{20}{345.191}{270.762}{\Character{98}}%
\TextAt{18}{344.527}{267.674}{Z}%
\TextAt{9}{352.497}{267.674}{\Character{41}}%
\TextAt{8}{357.049}{267.674}{.}%
\TextAt{8}{91.528}{250.628}{3)}%
\TextAt{8}{107.918}{250.628}{If}%
\TextAt{13}{119.812}{250.628}{K}%
\NativeMathAt{134.294}{250.628}{11.955}{\subseteq}%
\TextAt{13}{147.257}{250.628}{K}%
\TextAt{3}{158.075}{254.966}{\Character{48}}%
\TextAt{8}{164.969}{250.628}{are}%
\TextAt{8}{184.681}{250.628}{extensions}%
\TextAt{8}{241.263}{250.628}{of}%
\TextAt{13}{254.795}{250.628}{k}%
\TextAt{8}{265.383}{250.628}{with}%
\TextAt{13}{292.245}{250.628}{f}%
\TextAt{21}{298.015}{248.835}{K}%
\TextAt{13}{306.217}{250.628}{\Character{59}}%
\TextAt{13}{311.441}{250.628}{f}%
\TextAt{21}{317.231}{248.835}{K}%
\TextAt{23}{324.936}{251.104}{\Character{48}}%
\TextAt{13}{331.783}{250.628}{\Character{60}}%
\NativeMathAt{344.552}{250.628}{11.955}{\infty}%
\TextAt{8}{356.507}{250.628}{,}%
\TextAt{8}{363.908}{250.628}{and}%
\TextAt{13}{386.872}{250.628}{L}%
\TextAt{16}{398.501}{250.628}{\Character{106}}%
\TextAt{13}{405.486}{250.628}{K}%
\TextAt{8}{420.403}{250.628}{and}%
\TextAt{13}{443.367}{250.628}{L}%
\TextAt{3}{451.332}{254.966}{\Character{48}}%
\TextAt{16}{457.791}{250.628}{\Character{106}}%
\TextAt{13}{464.776}{250.628}{K}%
\TextAt{3}{475.594}{254.966}{\Character{48}}%
\TextAt{8}{482.488}{250.628}{are}%
\TextAt{8}{502.2}{250.628}{(finite}%
\TextAt{8}{107.918}{233.582}{or}%
\TextAt{8}{122.229}{233.582}{infinite)}%
\TextAt{8}{166.451}{233.582}{Galois}%
\TextAt{8}{202.364}{233.582}{extensions}%
\TextAt{8}{258.745}{233.582}{with}%
\TextAt{13}{285.405}{233.582}{L}%
\NativeMathAt{296.69}{233.582}{11.955}{\subseteq}%
\TextAt{13}{309.31}{233.582}{L}%
\TextAt{3}{317.274}{237.92}{\Character{48}}%
\TextAt{8}{320.069}{233.582}{,}%
\TextAt{8}{327.218}{233.582}{then}%
\TextAt{8}{353.878}{233.582}{one}%
\TextAt{8}{375.338}{233.582}{has}%
\TextAt{8}{396.212}{233.582}{a}%
\TextAt{8}{405.967}{233.582}{commutative}%
\TextAt{8}{475.893}{233.582}{diagram}%
\TextAt{20}{270.489}{206.362}{\Character{98}}%
\TextAt{13}{267.796}{203.34}{A}%
\TextAt{21}{276.572}{201.547}{K}%
\TextAt{23}{284.276}{203.816}{\Character{48}}%
\TextAt{13}{326.456}{203.34}{G}%
\TextAt{9}{335.69}{203.34}{\Character{40}}%
\TextAt{13}{340.242}{203.34}{L}%
\TextAt{3}{348.206}{207.678}{\Character{48}}%
\TextAt{16}{354.322}{203.34}{\Character{106}}%
\TextAt{13}{360.964}{203.34}{K}%
\TextAt{3}{371.782}{207.678}{\Character{48}}%
\TextAt{9}{374.577}{203.34}{\Character{41}}%
\TextAt{7}{379.129}{207.678}{ab}%
\TextAt{20}{271.831}{161.438}{\Character{98}}%
\TextAt{13}{269.138}{158.416}{A}%
\TextAt{21}{277.914}{156.623}{K}%
\TextAt{13}{327.625}{158.416}{G}%
\TextAt{9}{336.859}{158.416}{\Character{40}}%
\TextAt{13}{341.411}{158.416}{L}%
\TextAt{16}{352.696}{158.416}{\Character{106}}%
\TextAt{13}{359.338}{158.416}{K}%
\TextAt{9}{370.156}{158.416}{\Character{41}}%
\TextAt{7}{374.708}{162.754}{ab}%
\TextAt{13}{384.141}{158.416}{\Character{58}}%
\NativeDiagram{%
\path[draw=black,line width=0.478bp] (220.81,-567.671) -- (248.626,-567.671);%
\path[draw=black,line width=0.478bp,line cap=round,line join=round] (246.378,-564.802) .. controls (246.833,-566.523) and (247.845,-567.336) .. (248.865,-567.671) .. controls (247.845,-568.005) and (246.833,-568.818) .. (246.378,-570.54);%
}%
\TextAt{17}{289.946}{211.137}{\Character{40}}%
\TextAt{21}{297.137}{211.137}{\Character{59}L}%
\TextAt{23}{305.248}{213.95}{\Character{48}}%
\TextAt{3}{307.933}{211.137}{\Character{106}}%
\TextAt{21}{310.285}{211.137}{K}%
\TextAt{23}{317.99}{213.95}{\Character{48}}%
\TextAt{17}{320.674}{211.137}{\Character{41}}%
\NativeDiagram{%
\path[draw=black,line width=0.478bp] (205.627,-577.032) -- (205.627,-598.87);%
\path[draw=black,line width=0.478bp,line cap=round,line join=round] (208.496,-596.623) .. controls (206.774,-597.077) and (205.961,-598.09) .. (205.627,-599.11) .. controls (205.292,-598.09) and (204.479,-597.077) .. (202.757,-596.623);%
}%
\TextAt{21}{249.548}{184.98}{N}%
\TextAt{22}{256.27}{182.689}{K}%
\TextAt{23}{262.815}{184.404}{\Character{48}}%
\TextAt{23}{265.499}{182.689}{\Character{106}}%
\TextAt{22}{267.768}{182.689}{K}%
\NativeDiagram{%
\path[draw=black,line width=0.478bp] (285.509,-578.227) -- (285.509,-600.657);%
\path[draw=black,line width=0.478bp,line cap=round,line join=round] (288.379,-598.409) .. controls (286.657,-598.864) and (285.844,-599.876) .. (285.509,-600.896) .. controls (285.175,-599.876) and (284.362,-598.864) .. (282.64,-598.409);%
\path[draw=black,line width=0.478bp] (219.468,-612.595) -- (249.795,-612.595);%
\path[draw=black,line width=0.478bp,line cap=round,line join=round] (247.548,-609.726) .. controls (248.002,-611.448) and (249.014,-612.261) .. (250.034,-612.595) .. controls (249.014,-612.93) and (248.002,-613.743) .. (247.548,-615.464);%
}%
\TextAt{17}{292.544}{152.612}{\Character{40}}%
\TextAt{21}{299.735}{152.612}{\Character{59}L}%
\TextAt{3}{307.846}{152.612}{\Character{106}}%
\TextAt{21}{310.199}{152.612}{K}%
\TextAt{17}{317.903}{152.612}{\Character{41}}%
\TextAt{12}{79.226}{116.751}{Exercise}%
\TextAt{12}{132.189}{116.751}{IV.6.5.}%
\TextAt{8}{175.738}{116.751}{If}%
\TextAt{13}{187.224}{116.751}{L}%
\TextAt{16}{198.509}{116.751}{\Character{106}}%
\TextAt{13}{205.151}{116.751}{K}%
\TextAt{8}{219.66}{116.751}{is}%
\TextAt{8}{231.217}{116.751}{a}%
\TextAt{8}{240.766}{116.751}{finite}%
\TextAt{8}{270.471}{116.751}{Galois}%
\TextAt{8}{306.178}{116.751}{extension,}%
\TextAt{8}{361.03}{116.751}{then}%
\TextAt{13}{387.484}{116.751}{G}%
\TextAt{7}{396.717}{121.09}{ab}%
\TextAt{21}{396.717}{113.651}{L}%
\TextAt{8}{409.841}{116.751}{is}%
\TextAt{8}{421.398}{116.751}{a}%
\TextAt{13}{430.947}{116.751}{G}%
\TextAt{9}{440.181}{116.751}{\Character{40}}%
\TextAt{13}{444.733}{116.751}{L}%
\TextAt{16}{456.018}{116.751}{\Character{106}}%
\TextAt{13}{462.66}{116.751}{K}%
\TextAt{9}{473.478}{116.751}{\Character{41}}%
\TextAt{8}{478.03}{116.751}{-module}%
\TextAt{8}{523.014}{116.751}{in}%
\TextAt{8}{79.226}{99.705}{a}%
\TextAt{8}{88.982}{99.705}{canonical}%
\TextAt{8}{140.365}{99.705}{way,}%
\TextAt{8}{166.379}{99.705}{and}%
\TextAt{8}{189.142}{99.705}{the}%
\TextAt{8}{209.298}{99.705}{transfer}%
\TextAt{8}{252.612}{99.705}{from}%
\TextAt{13}{280.252}{99.705}{G}%
\TextAt{21}{289.486}{97.912}{K}%
\TextAt{8}{301.586}{99.705}{to}%
\TextAt{13}{315.896}{99.705}{G}%
\TextAt{21}{325.13}{97.912}{L}%
\TextAt{8}{335.285}{99.705}{is}%
\TextAt{8}{347.048}{99.705}{a}%
\TextAt{8}{356.804}{99.705}{homomorphism}%
\TextAt{8}{239.24}{72.921}{Ver}%
\TextAt{9}{261.092}{72.921}{\Character{58}}%
\TextAt{13}{267.664}{72.921}{G}%
\TextAt{7}{276.898}{77.858}{ab}%
\TextAt{21}{276.898}{69.966}{K}%
\NativeArrowAt{289.651}{72.921}{11.955}{19.26}{\longrightarrow}%
\TextAt{9}{312.233}{72.921}{\Character{40}}%
\TextAt{13}{316.786}{72.921}{G}%
\TextAt{7}{326.019}{77.858}{ab}%
\TextAt{21}{326.019}{69.966}{L}%
\TextAt{9}{335.452}{72.921}{\Character{41}}%
\TextAt{21}{340.004}{77.858}{G}%
\TextAt{17}{346.603}{77.858}{\Character{40}}%
\TextAt{21}{349.896}{77.858}{L}%
\TextAt{3}{355.655}{77.858}{\Character{106}}%
\TextAt{21}{358.008}{77.858}{K}%
\TextAt{17}{365.712}{77.858}{\Character{41}}%
\TextAt{13}{369.504}{72.921}{\Character{58}}%
\TextAt{8}{297.211}{43.033}{280}%
\Destination{79.23}{384.995}{nav.Item.94}%
\Destination{79.23}{298.525}{nav.Item.95}%
\Destination{79.23}{264.685}{nav.Item.96}%
\Destination{78.23}{809.815}{nav.page.280}%
\Destination{79.226}{775.155}{nav.ex.4.6.2}%
\Destination{79.226}{673.755}{nav.ex.4.6.3}%
\Destination{79.226}{517.355}{nav.ex.4.6.4}%
\Destination{79.226}{130.755}{nav.ex.4.6.5}%
\Destination{79.36}{518.355}{nav.print.306}%
\Hotlink{342.29}{758.955}{6.858}{11.047}{nav.ex.4.6.1}%
\Hotlink{192.49}{640.555}{6.858}{11.047}{nav.section.4.5}%
}

% Page 293 (body page 281)
\Page{%
\TextAt{2}{79.226}{796.943}{IV.}%
\TextAt{2}{98.012}{796.943}{Abstract}%
\TextAt{2}{143.252}{796.943}{Class}%
\TextAt{2}{171.844}{796.943}{Field}%
\TextAt{2}{199.575}{796.943}{Theory}%
\Vector{q 0 G 0 g 0.299 w 7.226 18.729 m 460.769 18.729 l S Q}%
\TextAt{12}{79.226}{761.124}{Exercise}%
\TextAt{12}{133.263}{761.124}{IV.6.6.}%
\TextAt{8}{177.746}{761.124}{(Tautological}%
\TextAt{8}{249.703}{761.124}{class}%
\TextAt{8}{277.867}{761.124}{field}%
\TextAt{8}{303.951}{761.124}{theory.)}%
\TextAt{8}{351.009}{761.124}{Assume}%
\TextAt{8}{395.097}{761.124}{that}%
\TextAt{8}{421.194}{761.124}{the}%
\TextAt{8}{442.077}{761.124}{profinite}%
\TextAt{8}{489.633}{761.124}{group}%
\TextAt{13}{523.536}{761.124}{G}%
\TextAt{8}{79.226}{744.078}{satisfies}%
\TextAt{8}{122.336}{744.078}{the}%
\TextAt{8}{142.493}{744.078}{condition:}%
\TextAt{8}{198.427}{744.078}{for}%
\TextAt{8}{216.312}{744.078}{every}%
\TextAt{8}{247.204}{744.078}{finite}%
\TextAt{8}{277.116}{744.078}{Galois}%
\TextAt{8}{313.029}{744.078}{extension,}%
\TextAt{8}{240.866}{715.077}{Ver}%
\TextAt{9}{262.718}{715.077}{\Character{58}}%
\TextAt{13}{269.29}{715.077}{G}%
\TextAt{7}{278.524}{720.013}{ab}%
\TextAt{21}{278.524}{712.121}{K}%
\NativeArrowAt{291.277}{715.077}{11.955}{19.26}{\longrightarrow}%
\TextAt{9}{313.859}{715.077}{\Character{40}}%
\TextAt{13}{318.412}{715.077}{G}%
\TextAt{7}{327.645}{720.013}{ab}%
\TextAt{21}{327.645}{712.121}{L}%
\TextAt{9}{337.078}{715.077}{\Character{41}}%
\TextAt{21}{341.63}{720.013}{G}%
\TextAt{17}{348.229}{720.013}{\Character{40}}%
\TextAt{21}{351.522}{720.013}{L}%
\TextAt{3}{357.281}{720.013}{\Character{106}}%
\TextAt{21}{359.633}{720.013}{K}%
\TextAt{17}{367.338}{720.013}{\Character{41}}%
\TextAt{8}{79.226}{686.075}{is}%
\TextAt{8}{90.829}{686.075}{an}%
\TextAt{8}{106.927}{686.075}{isomorphism.}%
\TextAt{8}{179.844}{686.075}{(These}%
\TextAt{8}{218.107}{686.075}{are}%
\TextAt{8}{237.456}{686.075}{the}%
\TextAt{8}{257.452}{686.075}{profinite}%
\TextAt{8}{304.119}{686.075}{groups}%
\TextAt{8}{341.748}{686.075}{of}%
\TextAt{8}{354.917}{686.075}{“strict}%
\TextAt{8}{390.909}{686.075}{cohomological}%
\TextAt{8}{466.209}{686.075}{dimension}%
\TextAt{8}{521.388}{686.075}{2”}%
\TextAt{8}{79.226}{669.029}{(see}%
\TextAt{8}{102.726}{669.029}{[145],}%
\TextAt{8}{133.993}{669.029}{chap.}%
\TextAt{8}{164.917}{669.029}{III,}%
\TextAt{8}{185.404}{669.029}{th.}%
\TextAt{8}{203.643}{669.029}{(3.6.4)).)}%
\TextAt{8}{254.484}{669.029}{Put}%
\TextAt{13}{277.434}{669.029}{A}%
\TextAt{21}{286.209}{667.236}{K}%
\TextAt{9}{297.787}{669.029}{\Character{61}}%
\TextAt{13}{310.267}{669.029}{G}%
\TextAt{7}{319.5}{673.368}{ab}%
\TextAt{21}{319.5}{665.929}{K}%
\TextAt{8}{332.862}{669.029}{and}%
\TextAt{8}{355.657}{669.029}{form}%
\TextAt{8}{383.329}{669.029}{the}%
\TextAt{8}{403.517}{669.029}{direct}%
\TextAt{8}{436.712}{669.029}{limit}%
\TextAt{13}{464.707}{669.029}{A}%
\TextAt{9}{476.858}{669.029}{\Character{61}}%
\NativeMathAt{489.337}{669.029}{11.955}{\varinjlim}%
\TextAt{21}{505.596}{663.59}{K}%
\TextAt{13}{515.791}{669.029}{A}%
\TextAt{21}{524.567}{667.236}{K}%
\TextAt{8}{79.226}{651.983}{via}%
\TextAt{8}{98.414}{651.983}{the}%
\TextAt{8}{118.571}{651.983}{transfer.}%
\TextAt{8}{166.435}{651.983}{Then}%
\TextAt{13}{196.993}{651.983}{A}%
\TextAt{21}{205.768}{650.19}{K}%
\TextAt{8}{217.868}{651.983}{is}%
\TextAt{8}{229.632}{651.983}{identified}%
\TextAt{8}{280.68}{651.983}{with}%
\TextAt{13}{307.34}{651.983}{A}%
\TextAt{21}{316.116}{656.322}{G}%
\TextAt{22}{322.714}{654.982}{K}%
\TextAt{8}{330.256}{651.983}{.}%
\TextAt{8}{97.159}{631.949}{Show}%
\TextAt{8}{128.051}{631.949}{that}%
\TextAt{8}{153.42}{631.949}{for}%
\TextAt{8}{171.305}{631.949}{every}%
\TextAt{8}{202.197}{631.949}{cyclic}%
\TextAt{8}{234.38}{631.949}{extension}%
\TextAt{13}{286.146}{631.949}{L}%
\TextAt{16}{297.432}{631.949}{\Character{106}}%
\TextAt{13}{304.073}{631.949}{K}%
\TextAt{8}{318.788}{631.949}{one}%
\TextAt{8}{340.248}{631.949}{has}%
\TextAt{8}{190.521}{591.5}{\char35\relax }%
\TextAt{13}{200.276}{591.5}{H}%
\TextAt{21}{210.899}{596.436}{i}%
\TextAt{9}{214.28}{591.5}{\Character{40}}%
\TextAt{13}{218.832}{591.5}{G}%
\TextAt{9}{228.066}{591.5}{\Character{40}}%
\TextAt{13}{232.618}{591.5}{L}%
\TextAt{16}{243.904}{591.5}{\Character{106}}%
\TextAt{13}{250.545}{591.5}{K}%
\TextAt{9}{261.363}{591.5}{\Character{41}}%
\TextAt{13}{265.915}{591.5}{\Character{59}}%
\TextAt{13}{271.139}{591.5}{A}%
\TextAt{21}{279.935}{589.707}{L}%
\TextAt{9}{286.192}{591.5}{\Character{41}}%
\TextAt{9}{294.047}{591.5}{\Character{61}}%
\TextAt{20}{306.491}{616.009}{\Character{56}}%
\TextAt{20}{306.491}{605.249}{\Character{60}}%
\TextAt{20}{306.491}{583.729}{\Character{58}}%
\TextAt{9}{317.118}{600.626}{\Character{91}}%
\TextAt{13}{320.369}{600.626}{L}%
\TextAt{9}{331.655}{600.626}{\Character{58}}%
\TextAt{13}{338.227}{600.626}{K}%
\TextAt{9}{349.045}{600.626}{\Character{93}}%
\TextAt{8}{364.252}{600.626}{for}%
\TextAt{13}{382.136}{600.626}{i}%
\TextAt{9}{389.451}{600.626}{\Character{61}}%
\TextAt{9}{401.861}{600.626}{0}%
\TextAt{13}{407.729}{600.626}{\Character{59}}%
\TextAt{9}{317.118}{580.17}{1}%
\TextAt{8}{364.252}{580.17}{for}%
\TextAt{13}{382.136}{580.17}{i}%
\TextAt{9}{389.451}{580.17}{\Character{61}}%
\TextAt{16}{401.876}{580.17}{\Character{0}}%
\TextAt{9}{411.175}{580.17}{1}%
\TextAt{13}{417.028}{580.17}{\Character{59}}%
\TextAt{8}{79.226}{548.295}{and}%
\TextAt{8}{103.368}{548.295}{that}%
\TextAt{8}{130.116}{548.295}{for}%
\TextAt{8}{149.38}{548.295}{every}%
\TextAt{8}{181.652}{548.295}{surjective}%
\TextAt{8}{235.442}{548.295}{homomorphism}%
\TextAt{13}{319.492}{548.295}{d}%
\TextAt{9}{331.245}{548.295}{\Character{58}}%
\TextAt{13}{340.168}{548.295}{G}%
\NativeMathAt{355.073}{548.295}{11.955}{\rightarrow}%
\TextAt{20}{373.364}{551.384}{\Character{98}}%
\TextAt{18}{372.7}{548.295}{Z}%
\TextAt{8}{380.67}{548.295}{,}%
\TextAt{8}{389.543}{548.295}{the}%
\TextAt{8}{411.079}{548.295}{induced}%
\TextAt{8}{456.023}{548.295}{map}%
\TextAt{13}{483.416}{548.295}{v}%
\TextAt{9}{495.176}{548.295}{\Character{58}}%
\TextAt{13}{504.099}{548.295}{A}%
\TextAt{21}{512.874}{546.502}{k}%
\TextAt{9}{523.665}{548.295}{\Character{61}}%
\TextAt{13}{79.226}{531.249}{G}%
\TextAt{7}{88.46}{535.588}{ab}%
\NativeMathAt{101.793}{531.249}{11.955}{\rightarrow}%
\TextAt{20}{118.312}{534.338}{\Character{98}}%
\TextAt{18}{117.648}{531.249}{Z}%
\TextAt{8}{129.855}{531.249}{is}%
\TextAt{8}{141.959}{531.249}{a}%
\TextAt{8}{152.054}{531.249}{henselian}%
\TextAt{8}{203.18}{531.249}{valuation}%
\TextAt{8}{254.592}{531.249}{with}%
\TextAt{8}{281.592}{531.249}{respect}%
\TextAt{8}{321.982}{531.249}{to}%
\TextAt{13}{336.632}{531.249}{d}%
\TextAt{8}{342.715}{531.249}{.}%
\TextAt{8}{352.183}{531.249}{The}%
\TextAt{8}{376.577}{531.249}{corresponding}%
\TextAt{8}{452.103}{531.249}{reciprocity}%
\TextAt{8}{510.652}{531.249}{map}%
\TextAt{13}{79.226}{514.203}{r}%
\TextAt{21}{84.501}{512.344}{L}%
\TextAt{3}{90.26}{512.344}{\Character{106}}%
\TextAt{21}{92.613}{512.344}{K}%
\TextAt{9}{104.136}{514.203}{\Character{58}}%
\TextAt{13}{110.709}{514.203}{G}%
\TextAt{9}{119.942}{514.203}{\Character{40}}%
\TextAt{13}{124.495}{514.203}{L}%
\TextAt{16}{135.78}{514.203}{\Character{106}}%
\TextAt{13}{142.422}{514.203}{K}%
\TextAt{9}{153.24}{514.203}{\Character{41}}%
\NativeMathAt{161.113}{514.203}{11.955}{\rightarrow}%
\TextAt{13}{176.389}{514.203}{A}%
\TextAt{21}{185.164}{512.41}{K}%
\TextAt{9}{193.367}{514.203}{\Character{47}}%
\TextAt{13}{199.22}{514.203}{N}%
\TextAt{21}{208.594}{512.344}{L}%
\TextAt{3}{214.353}{512.344}{\Character{106}}%
\TextAt{21}{216.705}{512.344}{K}%
\TextAt{13}{224.908}{514.203}{A}%
\TextAt{21}{233.683}{512.41}{L}%
\TextAt{8}{243.837}{514.203}{is}%
\TextAt{8}{255.601}{514.203}{essentially}%
\TextAt{8}{311.659}{514.203}{the}%
\TextAt{8}{331.816}{514.203}{identity.}%
\TextAt{8}{97.159}{488.191}{Abstract}%
\TextAt{8}{145.846}{488.191}{class}%
\TextAt{8}{173.457}{488.191}{field}%
\TextAt{8}{198.987}{488.191}{theory}%
\TextAt{8}{235.911}{488.191}{acquires}%
\TextAt{8}{281.346}{488.191}{a}%
\TextAt{8}{291.275}{488.191}{much}%
\TextAt{8}{322.663}{488.191}{broader}%
\TextAt{8}{365.768}{488.191}{range}%
\TextAt{8}{397.814}{488.191}{of}%
\TextAt{8}{411.317}{488.191}{applications}%
\TextAt{8}{476.598}{488.191}{if}%
\TextAt{8}{487.495}{488.191}{it}%
\TextAt{8}{499.373}{488.191}{is}%
\TextAt{8}{511.31}{488.191}{gen-}%
\TextAt{8}{79.226}{471.145}{eralized}%
\TextAt{8}{122.145}{471.145}{as}%
\TextAt{8}{136.515}{471.145}{follows.}%
\TextAt{12}{79.226}{445.133}{Exercise}%
\TextAt{12}{131.986}{445.133}{IV.6.7.}%
\TextAt{8}{175.359}{445.133}{Let}%
\TextAt{13}{195.945}{445.133}{G}%
\TextAt{8}{208.693}{445.133}{be}%
\TextAt{8}{224.235}{445.133}{a}%
\TextAt{8}{233.607}{445.133}{profinite}%
\TextAt{8}{280.053}{445.133}{group}%
\TextAt{8}{312.845}{445.133}{and}%
\TextAt{13}{335.225}{445.133}{B}%
\TextAt{9}{344.721}{445.133}{\Character{40}}%
\TextAt{13}{349.273}{445.133}{G}%
\TextAt{9}{358.507}{445.133}{\Character{41}}%
\TextAt{8}{366.574}{445.133}{the}%
\TextAt{8}{386.347}{445.133}{category}%
\TextAt{8}{433.127}{445.133}{of}%
\TextAt{8}{446.074}{445.133}{finite}%
\TextAt{13}{475.603}{445.133}{G}%
\TextAt{8}{484.837}{445.133}{-sets,}%
\TextAt{8}{514.562}{445.133}{i.e.,}%
\TextAt{8}{79.226}{428.087}{of}%
\TextAt{8}{92.763}{428.087}{finite}%
\TextAt{8}{122.881}{428.087}{sets}%
\TextAt{13}{145.97}{428.087}{X}%
\TextAt{8}{160.729}{428.087}{with}%
\TextAt{8}{187.595}{428.087}{a}%
\TextAt{8}{197.557}{428.087}{continuous}%
\TextAt{13}{256.368}{428.087}{G}%
\TextAt{8}{265.602}{428.087}{-operation.}%
\TextAt{8}{327.033}{428.087}{Show}%
\TextAt{8}{358.132}{428.087}{that}%
\TextAt{8}{383.707}{428.087}{the}%
\TextAt{8}{404.07}{428.087}{connected,}%
\TextAt{8}{462.203}{428.087}{i.e.,}%
\TextAt{8}{484.566}{428.087}{transitive}%
\TextAt{13}{79.226}{411.041}{G}%
\TextAt{8}{88.46}{411.041}{-sets}%
\TextAt{8}{115.584}{411.041}{in}%
\TextAt{13}{129.582}{411.041}{B}%
\TextAt{9}{139.078}{411.041}{\Character{40}}%
\TextAt{13}{143.631}{411.041}{G}%
\TextAt{9}{152.864}{411.041}{\Character{41}}%
\TextAt{8}{161.659}{411.041}{are,}%
\TextAt{8}{184.853}{411.041}{up}%
\TextAt{8}{202.103}{411.041}{to}%
\TextAt{8}{216.758}{411.041}{isomorphism,}%
\TextAt{8}{288.861}{411.041}{the}%
\TextAt{8}{309.362}{411.041}{sets}%
\TextAt{13}{332.589}{411.041}{G}%
\TextAt{9}{341.823}{411.041}{\Character{47}}%
\TextAt{13}{347.676}{411.041}{G}%
\TextAt{21}{356.909}{409.247}{K}%
\TextAt{8}{365.112}{411.041}{,}%
\TextAt{8}{372.693}{411.041}{where}%
\TextAt{13}{406.847}{411.041}{G}%
\TextAt{21}{416.08}{409.247}{K}%
\TextAt{8}{428.525}{411.041}{is}%
\TextAt{8}{440.634}{411.041}{an}%
\TextAt{8}{457.238}{411.041}{open}%
\TextAt{8}{485.869}{411.041}{subgroup}%
\TextAt{8}{79.226}{393.995}{of}%
\TextAt{13}{92.556}{393.995}{G}%
\TextAt{8}{101.79}{393.995}{,}%
\TextAt{8}{108.939}{393.995}{and}%
\TextAt{13}{131.702}{393.995}{G}%
\TextAt{8}{144.833}{393.995}{operates}%
\TextAt{8}{191.398}{393.995}{via}%
\TextAt{8}{210.586}{393.995}{multiplication}%
\TextAt{8}{285.712}{393.995}{on}%
\TextAt{8}{301.971}{393.995}{the}%
\TextAt{8}{322.128}{393.995}{left.}%
\TextAt{8}{97.159}{367.982}{If}%
\TextAt{13}{108.851}{367.982}{X}%
\TextAt{8}{123.404}{367.982}{is}%
\TextAt{8}{135.167}{367.982}{a}%
\TextAt{8}{144.923}{367.982}{finite}%
\TextAt{13}{174.835}{367.982}{G}%
\TextAt{8}{184.068}{367.982}{-set}%
\TextAt{8}{206.233}{367.982}{and}%
\TextAt{13}{228.996}{367.982}{x}%
\NativeMathAt{238.969}{367.982}{11.955}{\in}%
\TextAt{13}{250.26}{367.982}{X}%
\TextAt{8}{260.915}{367.982}{,}%
\TextAt{8}{268.064}{367.982}{then}%
\TextAt{13}{217.324}{338.981}{\Character{25}}%
\TextAt{17}{223.964}{337.188}{1}%
\TextAt{9}{228.696}{338.981}{\Character{40}}%
\TextAt{13}{233.249}{338.981}{X}%
\TextAt{13}{243.244}{338.981}{\Character{59}}%
\TextAt{13}{248.468}{338.981}{x}%
\TextAt{9}{255.15}{338.981}{\Character{41}}%
\TextAt{9}{263.005}{338.981}{\Character{61}}%
\TextAt{13}{275.448}{338.981}{G}%
\TextAt{21}{284.682}{337.188}{x}%
\TextAt{9}{293.268}{338.981}{\Character{61}}%
\TextAt{16}{305.693}{338.981}{\Character{102}}%
\TextAt{13}{311.671}{338.981}{\Character{27}}%
\NativeMathAt{322.074}{338.981}{11.955}{\in}%
\TextAt{13}{333.365}{338.981}{G}%
\TextAt{16}{345.92}{338.981}{\Character{106}}%
\TextAt{13}{352.561}{338.981}{\Character{27}}%
\TextAt{13}{359.627}{338.981}{x}%
\TextAt{9}{369.617}{338.981}{\Character{61}}%
\TextAt{13}{382.042}{338.981}{x}%
\TextAt{16}{388.694}{338.981}{\Character{103}}%
\TextAt{8}{79.226}{309.98}{is}%
\TextAt{8}{91.861}{309.98}{called}%
\TextAt{8}{125.896}{309.98}{the}%
\TextAt{12}{146.923}{309.98}{fundamental}%
\TextAt{12}{225.443}{309.98}{group}%
\TextAt{8}{264.116}{309.98}{of}%
\TextAt{13}{278.317}{309.98}{X}%
\TextAt{8}{293.74}{309.98}{with}%
\TextAt{8}{321.272}{309.98}{base}%
\TextAt{8}{348.217}{309.98}{point}%
\TextAt{13}{379.657}{309.98}{x}%
\TextAt{8}{386.309}{309.98}{.}%
\TextAt{8}{397.371}{309.98}{For}%
\TextAt{8}{419.211}{309.98}{a}%
\TextAt{8}{429.838}{309.98}{map}%
\TextAt{13}{456.723}{309.98}{f}%
\TextAt{9}{468.575}{309.98}{\Character{58}}%
\TextAt{13}{476.632}{309.98}{X}%
\NativeMathAt{492.092}{309.98}{11.955}{\rightarrow}%
\TextAt{13}{508.852}{309.98}{Y}%
\TextAt{8}{523.014}{309.98}{in}%
\TextAt{13}{79.226}{292.934}{B}%
\TextAt{9}{88.722}{292.934}{\Character{40}}%
\TextAt{13}{93.275}{292.934}{G}%
\TextAt{9}{102.508}{292.934}{\Character{41}}%
\TextAt{8}{107.061}{292.934}{,}%
\TextAt{8}{114.21}{292.934}{we}%
\TextAt{8}{131.437}{292.934}{put}%
\TextAt{13}{248.884}{275.888}{G}%
\TextAt{9}{258.117}{275.888}{\Character{40}}%
\TextAt{13}{262.67}{275.888}{X}%
\TextAt{16}{276.645}{275.888}{\Character{106}}%
\TextAt{13}{283.287}{275.888}{Y}%
\TextAt{9}{292.68}{275.888}{\Character{41}}%
\TextAt{9}{300.535}{275.888}{\Character{61}}%
\TextAt{8}{312.979}{275.888}{Aut}%
\TextAt{21}{332.813}{274.095}{Y}%
\TextAt{9}{340.101}{275.888}{\Character{40}}%
\TextAt{13}{344.653}{275.888}{X}%
\TextAt{9}{355.308}{275.888}{\Character{41}}%
\TextAt{13}{359.86}{275.888}{\Character{58}}%
\TextAt{13}{79.226}{252.367}{f}%
\TextAt{8}{90.499}{252.367}{is}%
\TextAt{8}{102.592}{252.367}{called}%
\TextAt{12}{136.084}{252.367}{Galois}%
\TextAt{8}{176.941}{252.367}{if}%
\TextAt{13}{187.994}{252.367}{X}%
\TextAt{8}{202.875}{252.367}{and}%
\TextAt{13}{225.967}{252.367}{Y}%
\TextAt{8}{239.586}{252.367}{are}%
\TextAt{8}{259.426}{252.367}{connected}%
\TextAt{8}{314.378}{252.367}{and}%
\TextAt{13}{337.47}{252.367}{G}%
\TextAt{9}{346.703}{252.367}{\Character{40}}%
\TextAt{13}{351.256}{252.367}{X}%
\TextAt{16}{365.792}{252.367}{\Character{106}}%
\TextAt{13}{372.994}{252.367}{Y}%
\TextAt{9}{382.387}{252.367}{\Character{41}}%
\TextAt{8}{391.166}{252.367}{operates}%
\TextAt{8}{438.06}{252.367}{transitively}%
\TextAt{8}{499.923}{252.367}{on}%
\TextAt{8}{516.51}{252.367}{the}%
\TextAt{8}{79.226}{235.321}{fibres}%
\TextAt{13}{110.501}{235.321}{f}%
\TextAt{3}{117.547}{239.659}{\Character{0}}%
\TextAt{17}{124.134}{239.659}{1}%
\TextAt{9}{128.866}{235.321}{\Character{40}}%
\TextAt{13}{133.418}{235.321}{y}%
\TextAt{9}{139.555}{235.321}{\Character{41}}%
\TextAt{8}{144.107}{235.321}{.}%
\TextAt{12}{79.226}{209.308}{Exercise}%
\TextAt{12}{133.034}{209.308}{IV.6.8.}%
\TextAt{8}{177.318}{209.308}{Let}%
\TextAt{13}{198.815}{209.308}{f}%
\TextAt{9}{210.082}{209.308}{\Character{58}}%
\TextAt{13}{217.554}{209.308}{X}%
\NativeMathAt{232.43}{209.308}{11.955}{\rightarrow}%
\TextAt{13}{248.606}{209.308}{Y}%
\TextAt{8}{262.424}{209.308}{be}%
\TextAt{8}{278.877}{209.308}{a}%
\TextAt{8}{289.16}{209.308}{map}%
\TextAt{8}{315.702}{209.308}{of}%
\TextAt{8}{329.56}{209.308}{connected}%
\TextAt{8}{384.711}{209.308}{finite}%
\TextAt{13}{415.151}{209.308}{G}%
\TextAt{8}{424.385}{209.308}{-sets,}%
\TextAt{8}{455.076}{209.308}{and}%
\TextAt{8}{478.367}{209.308}{let}%
\TextAt{13}{495.799}{209.308}{x}%
\NativeMathAt{506.672}{209.308}{11.955}{\in}%
\TextAt{13}{518.863}{209.308}{X}%
\TextAt{8}{529.518}{209.308}{,}%
\TextAt{13}{79.226}{192.262}{y}%
\TextAt{9}{90.395}{192.262}{\Character{61}}%
\TextAt{13}{104.532}{192.262}{f}%
\TextAt{9}{111.578}{192.262}{\Character{40}}%
\TextAt{13}{116.13}{192.262}{x}%
\TextAt{9}{122.782}{192.262}{\Character{41}}%
\NativeMathAt{132.367}{192.262}{11.955}{\in}%
\TextAt{13}{145.369}{192.262}{Y}%
\TextAt{8}{154.762}{192.262}{.}%
\TextAt{8}{166.223}{192.262}{Show}%
\TextAt{8}{198.119}{192.262}{that}%
\TextAt{13}{224.492}{192.262}{f}%
\TextAt{8}{236.44}{192.262}{is}%
\TextAt{8}{249.208}{192.262}{Galois}%
\TextAt{8}{286.126}{192.262}{if}%
\TextAt{8}{297.854}{192.262}{and}%
\TextAt{8}{321.621}{192.262}{only}%
\TextAt{8}{348.316}{192.262}{if}%
\TextAt{13}{360.044}{192.262}{\Character{25}}%
\TextAt{17}{366.685}{190.469}{1}%
\TextAt{9}{371.417}{192.262}{\Character{40}}%
\TextAt{13}{375.969}{192.262}{X}%
\TextAt{13}{385.964}{192.262}{\Character{59}}%
\TextAt{13}{391.188}{192.262}{x}%
\TextAt{9}{397.87}{192.262}{\Character{41}}%
\TextAt{8}{407.324}{192.262}{is}%
\TextAt{8}{420.092}{192.262}{a}%
\TextAt{8}{430.852}{192.262}{normal}%
\TextAt{8}{471.535}{192.262}{subgroup}%
\TextAt{8}{523.337}{192.262}{of}%
\TextAt{13}{79.226}{175.216}{\Character{25}}%
\TextAt{17}{85.866}{173.423}{1}%
\TextAt{9}{90.599}{175.216}{\Character{40}}%
\TextAt{13}{95.151}{175.216}{Y}%
\TextAt{13}{102.587}{175.216}{\Character{59}}%
\TextAt{13}{107.824}{175.216}{y}%
\TextAt{9}{113.974}{175.216}{\Character{41}}%
\TextAt{8}{118.526}{175.216}{.}%
\TextAt{8}{126.975}{175.216}{In}%
\TextAt{8}{141.596}{175.216}{this}%
\TextAt{8}{164.418}{175.216}{case,}%
\TextAt{8}{192.441}{175.216}{one}%
\TextAt{8}{213.901}{175.216}{has}%
\TextAt{8}{234.775}{175.216}{a}%
\TextAt{8}{244.53}{175.216}{canonical}%
\TextAt{8}{295.913}{175.216}{isomorphism}%
\TextAt{13}{228.462}{146.215}{G}%
\TextAt{9}{237.695}{146.215}{\Character{40}}%
\TextAt{13}{242.248}{146.215}{X}%
\TextAt{16}{256.224}{146.215}{\Character{106}}%
\TextAt{13}{262.865}{146.215}{Y}%
\TextAt{9}{272.258}{146.215}{\Character{41}}%
\NativeMathAt{280.132}{146.215}{11.955}{\cong}%
\TextAt{13}{292.751}{146.215}{\Character{25}}%
\TextAt{17}{299.391}{144.422}{1}%
\TextAt{9}{304.123}{146.215}{\Character{40}}%
\TextAt{13}{308.676}{146.215}{Y}%
\TextAt{13}{316.112}{146.215}{\Character{59}}%
\TextAt{13}{321.349}{146.215}{y}%
\TextAt{9}{327.499}{146.215}{\Character{41}\Character{47}}%
\TextAt{13}{337.904}{146.215}{\Character{25}}%
\TextAt{17}{344.544}{144.422}{1}%
\TextAt{9}{349.277}{146.215}{\Character{40}}%
\TextAt{13}{353.829}{146.215}{X}%
\TextAt{13}{363.824}{146.215}{\Character{59}}%
\TextAt{13}{369.048}{146.215}{x}%
\TextAt{9}{375.73}{146.215}{\Character{41}}%
\TextAt{13}{380.282}{146.215}{\Character{58}}%
\TextAt{8}{97.159}{108.248}{A}%
\TextAt{8}{109.831}{108.248}{pair}%
\TextAt{8}{133.897}{108.248}{of}%
\TextAt{8}{147.227}{108.248}{functors}%
\TextAt{13}{229.5}{91.202}{A}%
\TextAt{9}{241.597}{91.202}{\Character{61}}%
\TextAt{9}{254.007}{91.202}{\Character{40}}%
\TextAt{13}{258.574}{91.202}{A}%
\TextAt{3}{267.35}{96.138}{\Character{3}}%
\TextAt{13}{272.082}{91.202}{\Character{59}}%
\TextAt{13}{277.306}{91.202}{A}%
\TextAt{3}{286.102}{89.408}{\Character{3}}%
\TextAt{9}{290.834}{91.202}{\Character{41}}%
\TextAt{9}{298.689}{91.202}{\Character{58}}%
\TextAt{13}{305.28}{91.202}{B}%
\TextAt{9}{314.776}{91.202}{\Character{40}}%
\TextAt{13}{319.328}{91.202}{G}%
\TextAt{9}{328.562}{91.202}{\Character{41}}%
\NativeArrowAt{336.435}{91.202}{11.955}{19.26}{\longrightarrow}%
\TextAt{9}{359.017}{91.202}{\Character{40}}%
\TextAt{13}{363.569}{91.202}{ab}%
\TextAt{9}{374.691}{91.202}{\Character{41}}%
\TextAt{13}{379.244}{91.202}{\Character{59}}%
\TextAt{8}{297.211}{43.033}{281}%
\Destination{78.23}{809.815}{nav.page.281}%
\Destination{79.226}{775.155}{nav.ex.4.6.6}%
\Destination{79.226}{459.155}{nav.ex.4.6.7}%
\Destination{79.226}{223.355}{nav.ex.4.6.8}%
\Destination{79.36}{324.955}{nav.print.307}%
\Hotlink{164.417}{666.855}{14.294}{11.047}{nav.chapter.3}%
\Hotlink{102.226}{666.855}{25.078}{11.047}{nav.cite.N145}%
}

% Page 294 (body page 282)
\Page{%
\TextAt{2}{79.226}{796.943}{IV.}%
\TextAt{2}{98.012}{796.943}{Abstract}%
\TextAt{2}{143.252}{796.943}{Class}%
\TextAt{2}{171.844}{796.943}{Field}%
\TextAt{2}{199.575}{796.943}{Theory}%
\Vector{q 0 G 0 g 0.299 w 7.226 18.729 m 460.769 18.729 l S Q}%
\TextAt{8}{79.226}{761.124}{consisting}%
\TextAt{8}{134.462}{761.124}{of}%
\TextAt{8}{148.92}{761.124}{a}%
\TextAt{8}{159.802}{761.124}{contravariant}%
\TextAt{8}{232.505}{761.124}{functor}%
\TextAt{13}{274.28}{761.124}{A}%
\TextAt{3}{283.055}{765.462}{\Character{3}}%
\TextAt{8}{292.812}{761.124}{and}%
\TextAt{8}{316.702}{761.124}{a}%
\TextAt{8}{327.585}{761.124}{covariant}%
\TextAt{8}{379.139}{761.124}{functor}%
\TextAt{13}{420.914}{761.124}{A}%
\TextAt{3}{429.689}{759.33}{\Character{3}}%
\TextAt{8}{439.446}{761.124}{from}%
\TextAt{13}{468.214}{761.124}{B}%
\TextAt{9}{477.71}{761.124}{\Character{40}}%
\TextAt{13}{482.262}{761.124}{G}%
\TextAt{9}{491.496}{761.124}{\Character{41}}%
\TextAt{8}{501.073}{761.124}{to}%
\TextAt{8}{516.51}{761.124}{the}%
\TextAt{8}{79.226}{744.078}{category}%
\TextAt{9}{126.389}{744.078}{\Character{40}}%
\TextAt{13}{130.942}{744.078}{ab}%
\TextAt{9}{142.064}{744.078}{\Character{41}}%
\TextAt{8}{150.513}{744.078}{of}%
\TextAt{8}{163.843}{744.078}{abelian}%
\TextAt{8}{204.491}{744.078}{groups}%
\TextAt{8}{242.281}{744.078}{is}%
\TextAt{8}{254.045}{744.078}{called}%
\TextAt{8}{287.209}{744.078}{a}%
\TextAt{12}{296.964}{744.078}{double}%
\TextAt{12}{340.469}{744.078}{functor}%
\TextAt{8}{386.867}{744.078}{if}%
\TextAt{13}{241.091}{715.46}{A}%
\TextAt{3}{249.866}{720.396}{\Character{3}}%
\TextAt{9}{254.598}{715.46}{\Character{40}}%
\TextAt{13}{259.151}{715.46}{X}%
\TextAt{9}{269.806}{715.46}{\Character{41}}%
\TextAt{9}{277.661}{715.46}{\Character{61}}%
\TextAt{13}{290.104}{715.46}{A}%
\TextAt{3}{298.88}{713.667}{\Character{3}}%
\TextAt{9}{303.612}{715.46}{\Character{40}}%
\TextAt{13}{308.164}{715.46}{X}%
\TextAt{9}{318.82}{715.46}{\Character{41}}%
\TextAt{9}{326.675}{715.46}{\Character{61}\Character{58}}%
\TextAt{13}{342.37}{715.46}{A}%
\TextAt{9}{351.145}{715.46}{\Character{40}}%
\TextAt{13}{355.698}{715.46}{X}%
\TextAt{9}{366.353}{715.46}{\Character{41}}%
\TextAt{8}{79.226}{686.843}{for}%
\TextAt{8}{97.111}{686.843}{all}%
\TextAt{13}{113.37}{686.843}{X}%
\NativeMathAt{127.346}{686.843}{11.955}{\in}%
\TextAt{13}{138.637}{686.843}{B}%
\TextAt{9}{148.133}{686.843}{\Character{40}}%
\TextAt{13}{152.686}{686.843}{G}%
\TextAt{9}{161.919}{686.843}{\Character{41}}%
\TextAt{8}{166.472}{686.843}{.}%
\TextAt{8}{174.92}{686.843}{We}%
\TextAt{8}{195.064}{686.843}{define}%
\TextAt{13}{262.808}{669.797}{A}%
\TextAt{21}{271.584}{668.003}{K}%
\TextAt{9}{283.107}{669.797}{\Character{61}}%
\TextAt{13}{295.533}{669.797}{A}%
\TextAt{9}{304.308}{669.797}{\Character{40}}%
\TextAt{13}{308.86}{669.797}{G}%
\TextAt{9}{318.094}{669.797}{\Character{47}}%
\TextAt{13}{323.947}{669.797}{G}%
\TextAt{21}{333.181}{668.003}{K}%
\TextAt{9}{341.383}{669.797}{\Character{41}}%
\TextAt{13}{345.936}{669.797}{\Character{58}}%
\TextAt{8}{79.226}{646.439}{If}%
\TextAt{13}{90.918}{646.439}{f}%
\TextAt{9}{101.286}{646.439}{\Character{58}}%
\TextAt{13}{107.858}{646.439}{X}%
\NativeMathAt{121.834}{646.439}{11.955}{\rightarrow}%
\TextAt{13}{137.11}{646.439}{Y}%
\TextAt{8}{150.401}{646.439}{is}%
\TextAt{8}{162.164}{646.439}{a}%
\TextAt{8}{171.92}{646.439}{morphism}%
\TextAt{8}{226.615}{646.439}{in}%
\TextAt{13}{240.268}{646.439}{B}%
\TextAt{9}{249.764}{646.439}{\Character{40}}%
\TextAt{13}{254.316}{646.439}{G}%
\TextAt{9}{263.55}{646.439}{\Character{41}}%
\TextAt{8}{268.102}{646.439}{,}%
\TextAt{8}{275.251}{646.439}{then}%
\TextAt{8}{301.911}{646.439}{we}%
\TextAt{8}{319.139}{646.439}{put}%
\TextAt{13}{214.483}{617.822}{A}%
\TextAt{3}{223.259}{622.758}{\Character{3}}%
\TextAt{9}{227.991}{617.822}{\Character{40}}%
\TextAt{13}{232.543}{617.822}{f}%
\TextAt{9}{239.59}{617.822}{\Character{41}}%
\TextAt{9}{247.445}{617.822}{\Character{61}}%
\TextAt{13}{259.888}{617.822}{f}%
\TextAt{3}{266.935}{622.758}{\Character{3}}%
\TextAt{8}{295.578}{617.822}{and}%
\TextAt{13}{338.353}{617.822}{A}%
\TextAt{3}{347.129}{616.029}{\Character{3}}%
\TextAt{9}{351.861}{617.822}{\Character{40}}%
\TextAt{13}{356.413}{617.822}{f}%
\TextAt{9}{363.46}{617.822}{\Character{41}}%
\TextAt{9}{371.315}{617.822}{\Character{61}}%
\TextAt{13}{383.758}{617.822}{f}%
\TextAt{3}{389.528}{616.029}{\Character{3}}%
\TextAt{13}{394.261}{617.822}{\Character{58}}%
\TextAt{8}{79.226}{589.204}{A}%
\TextAt{8}{92.59}{589.204}{homomorphism}%
\TextAt{13}{175.953}{589.204}{h}%
\TextAt{9}{187.192}{589.204}{\Character{58}}%
\TextAt{13}{194.945}{589.204}{A}%
\NativeMathAt{208.221}{589.204}{11.955}{\rightarrow}%
\TextAt{13}{224.678}{589.204}{B}%
\TextAt{8}{238.764}{589.204}{of}%
\TextAt{8}{252.786}{589.204}{double}%
\TextAt{8}{291.197}{589.204}{functors}%
\TextAt{8}{337.152}{589.204}{is}%
\TextAt{8}{349.609}{589.204}{a}%
\TextAt{8}{360.056}{589.204}{family}%
\TextAt{8}{396.519}{589.204}{of}%
\TextAt{8}{410.541}{589.204}{homomorphisms}%
\TextAt{13}{498.519}{589.204}{h}%
\TextAt{9}{505.257}{589.204}{\Character{40}}%
\TextAt{13}{509.809}{589.204}{X}%
\TextAt{9}{520.465}{589.204}{\Character{41}}%
\TextAt{9}{529.503}{589.204}{\Character{58}}%
\TextAt{13}{79.226}{572.158}{A}%
\TextAt{9}{88.002}{572.158}{\Character{40}}%
\TextAt{13}{92.554}{572.158}{X}%
\TextAt{9}{103.209}{572.158}{\Character{41}}%
\NativeMathAt{113.464}{572.158}{11.955}{\rightarrow}%
\TextAt{13}{131.123}{572.158}{B}%
\TextAt{9}{140.619}{572.158}{\Character{40}}%
\TextAt{13}{145.171}{572.158}{X}%
\TextAt{9}{155.826}{572.158}{\Character{41}}%
\TextAt{8}{165.674}{572.158}{representing}%
\TextAt{8}{233.148}{572.158}{natural}%
\TextAt{8}{275.528}{572.158}{transformations}%
\TextAt{13}{361.294}{572.158}{A}%
\TextAt{3}{370.069}{576.497}{\Character{3}}%
\NativeMathAt{380.505}{572.158}{11.955}{\rightarrow}%
\TextAt{13}{398.163}{572.158}{B}%
\TextAt{3}{407.659}{576.497}{\Character{3}}%
\TextAt{8}{417.687}{572.158}{and}%
\TextAt{13}{441.847}{572.158}{A}%
\TextAt{3}{450.623}{570.365}{\Character{3}}%
\NativeMathAt{461.058}{572.158}{11.955}{\rightarrow}%
\TextAt{13}{478.716}{572.158}{B}%
\TextAt{3}{487.609}{570.365}{\Character{3}}%
\TextAt{8}{492.342}{572.158}{.}%
\TextAt{8}{504.984}{572.158}{A}%
\TextAt{13}{519.053}{572.158}{G}%
\TextAt{12}{528.286}{572.158}{-}%
\TextAt{12}{79.226}{555.112}{modulation}%
\TextAt{8}{149.845}{555.112}{is}%
\TextAt{8}{161.609}{555.112}{defined}%
\TextAt{8}{201.922}{555.112}{to}%
\TextAt{8}{216.232}{555.112}{be}%
\TextAt{8}{232.157}{555.112}{a}%
\TextAt{8}{241.912}{555.112}{double}%
\TextAt{8}{279.631}{555.112}{functor}%
\TextAt{13}{320.278}{555.112}{A}%
\TextAt{8}{332.951}{555.112}{such}%
\TextAt{8}{359.348}{555.112}{that}%
\TextAt{8}{89.579}{535.078}{(i)}%
\TextAt{13}{107.918}{535.078}{A}%
\TextAt{9}{116.694}{535.078}{\Character{40}}%
\TextAt{13}{121.246}{535.078}{X}%
\TextAt{16}{134.558}{535.078}{\Character{116}}%
\TextAt{13}{145.185}{535.078}{Y}%
\TextAt{9}{154.578}{535.078}{\Character{41}}%
\TextAt{9}{162.433}{535.078}{\Character{61}}%
\TextAt{13}{174.877}{535.078}{A}%
\TextAt{9}{183.652}{535.078}{\Character{40}}%
\TextAt{13}{188.204}{535.078}{X}%
\TextAt{9}{198.859}{535.078}{\Character{41}}%
\TextAt{16}{206.068}{535.078}{\Character{2}}%
\TextAt{13}{218.023}{535.078}{A}%
\TextAt{9}{226.799}{535.078}{\Character{40}}%
\TextAt{13}{231.351}{535.078}{Y}%
\TextAt{9}{240.744}{535.078}{\Character{41}}%
\TextAt{8}{245.297}{535.078}{.}%
\TextAt{8}{86.327}{518.032}{(ii)}%
\TextAt{8}{107.918}{518.032}{If}%
\TextAt{8}{119.611}{518.032}{among}%
\TextAt{8}{157.341}{518.032}{the}%
\TextAt{8}{177.498}{518.032}{two}%
\TextAt{8}{199.615}{518.032}{diagrams}%
\TextAt{13}{200.292}{483.41}{X}%
\TextAt{13}{249.945}{483.41}{X}%
\TextAt{3}{260.6}{487.748}{\Character{48}}%
\TextAt{13}{200.923}{444.365}{Y}%
\TextAt{13}{250.576}{444.365}{Y}%
\TextAt{3}{259.969}{448.703}{\Character{48}}%
\NativeDiagram{%
\path[draw=black,line width=0.478bp] (133.62,-295.168) -- (133.62,-316.408);%
\path[draw=black,line width=0.478bp,line cap=round,line join=round] (136.489,-314.161) .. controls (134.767,-314.615) and (133.954,-315.627) .. (133.62,-316.647) .. controls (133.285,-315.627) and (132.472,-314.615) .. (130.751,-314.161);%
}%
\TextAt{21}{197.857}{465.98}{f}%
\NativeDiagram{%
\path[draw=black,line width=0.478bp] (172.594,-287.601) -- (144.778,-287.601);%
\path[draw=black,line width=0.478bp,line cap=round,line join=round] (147.025,-290.47) .. controls (146.571,-288.748) and (145.558,-287.935) .. (144.538,-287.601) .. controls (145.558,-287.266) and (146.571,-286.453) .. (147.025,-284.731);%
}%
\TextAt{21}{226.946}{490.765}{g}%
\TextAt{23}{231.262}{493.578}{\Character{48}}%
\NativeDiagram{%
\path[draw=black,line width=0.478bp] (184.671,-295.168) -- (184.671,-315.811);%
\path[draw=black,line width=0.478bp,line cap=round,line join=round] (187.54,-313.564) .. controls (185.818,-314.018) and (185.005,-315.031) .. (184.671,-316.05) .. controls (184.336,-315.031) and (183.523,-314.018) .. (181.802,-313.564);%
}%
\TextAt{21}{259.486}{465.979}{f}%
\TextAt{23}{264.433}{468.792}{\Character{48}}%
\NativeDiagram{%
\path[draw=black,line width=0.478bp] (173.225,-326.646) -- (144.147,-326.646);%
\path[draw=black,line width=0.478bp,line cap=round,line join=round] (146.394,-329.515) .. controls (145.94,-327.794) and (144.927,-326.981) .. (143.908,-326.646) .. controls (144.927,-326.311) and (145.94,-325.499) .. (146.394,-323.777);%
}%
\TextAt{21}{228.289}{441.107}{g}%
\TextAt{8}{293.604}{468.271}{and}%
\TextAt{13}{341.531}{486.399}{A}%
\TextAt{9}{350.306}{486.399}{\Character{40}}%
\TextAt{13}{354.859}{486.399}{X}%
\TextAt{9}{365.514}{486.399}{\Character{41}}%
\TextAt{13}{409.064}{486.399}{A}%
\TextAt{9}{417.84}{486.399}{\Character{40}}%
\TextAt{13}{422.392}{486.399}{X}%
\TextAt{3}{433.047}{490.737}{\Character{48}}%
\TextAt{9}{435.842}{486.399}{\Character{41}}%
\TextAt{13}{342.162}{444.165}{A}%
\TextAt{9}{350.937}{444.165}{\Character{40}}%
\TextAt{13}{355.49}{444.165}{Y}%
\TextAt{9}{364.883}{444.165}{\Character{41}}%
\TextAt{13}{409.695}{444.165}{A}%
\TextAt{9}{418.471}{444.165}{\Character{40}}%
\TextAt{13}{423.023}{444.165}{Y}%
\TextAt{3}{432.416}{448.503}{\Character{48}}%
\TextAt{9}{435.211}{444.165}{\Character{41}}%
\NativeDiagram{%
\path[draw=black,line width=0.478bp] (303.419,-284.612) -- (331.235,-284.612);%
\path[draw=black,line width=0.478bp,line cap=round,line join=round] (328.988,-281.743) .. controls (329.442,-283.464) and (330.454,-284.277) .. (331.474,-284.612) .. controls (330.454,-284.946) and (329.442,-285.759) .. (328.988,-287.481);%
}%
\TextAt{21}{384.157}{493.754}{g}%
\TextAt{23}{388.472}{496.567}{\Character{48}\Character{3}}%
\NativeDiagram{%
\path[draw=black,line width=0.478bp] (283.8,-295.168) -- (283.8,-315.811);%
\path[draw=black,line width=0.478bp,line cap=round,line join=round] (286.669,-313.564) .. controls (284.947,-314.018) and (284.134,-315.031) .. (283.8,-316.05) .. controls (283.465,-315.031) and (282.652,-314.018) .. (280.93,-313.564);%
}%
\TextAt{21}{344.599}{466.278}{f}%
\TextAt{23}{348.667}{465.282}{\Character{3}}%
\NativeDiagram{%
\path[draw=black,line width=0.478bp] (352.731,-295.168) -- (352.731,-315.811);%
\path[draw=black,line width=0.478bp,line cap=round,line join=round] (355.6,-313.564) .. controls (353.878,-314.018) and (353.066,-315.031) .. (352.731,-316.05) .. controls (352.396,-315.031) and (351.583,-314.018) .. (349.862,-313.564);%
}%
\TextAt{21}{427.546}{466.2}{f}%
\TextAt{23}{432.492}{469.013}{\Character{48}}%
\TextAt{23}{431.614}{464.207}{\Character{3}}%
\NativeDiagram{%
\path[draw=black,line width=0.478bp] (302.788,-326.846) -- (331.866,-326.846);%
\path[draw=black,line width=0.478bp,line cap=round,line join=round] (329.619,-323.977) .. controls (330.073,-325.699) and (331.085,-326.512) .. (332.105,-326.846) .. controls (331.085,-327.181) and (330.073,-327.994) .. (329.619,-329.715);%
}%
\TextAt{21}{385.25}{438.744}{g}%
\TextAt{23}{389.565}{441.557}{\Character{3}}%
\TextAt{8}{107.918}{413.084}{in}%
\TextAt{13}{122.363}{413.084}{B}%
\TextAt{9}{131.859}{413.084}{\Character{40}}%
\TextAt{13}{136.412}{413.084}{G}%
\TextAt{9}{145.645}{413.084}{\Character{41}}%
\TextAt{8}{150.198}{413.084}{,}%
\TextAt{8}{158.337}{413.084}{resp.}%
\TextAt{9}{190.035}{413.084}{\Character{40}}%
\TextAt{13}{194.587}{413.084}{ab}%
\TextAt{9}{205.709}{413.084}{\Character{41}}%
\TextAt{8}{210.261}{413.084}{,}%
\TextAt{8}{218.4}{413.084}{the}%
\TextAt{8}{239.349}{413.084}{one}%
\TextAt{8}{261.6}{413.084}{on}%
\TextAt{8}{278.651}{413.084}{the}%
\TextAt{8}{299.599}{413.084}{left}%
\TextAt{8}{320.87}{413.084}{is}%
\TextAt{8}{333.426}{413.084}{cartesian,}%
\TextAt{8}{387.162}{413.084}{then}%
\TextAt{8}{414.614}{413.084}{the}%
\TextAt{8}{435.563}{413.084}{one}%
\TextAt{8}{457.814}{413.084}{on}%
\TextAt{8}{474.865}{413.084}{the}%
\TextAt{8}{495.813}{413.084}{right}%
\TextAt{8}{524.903}{413.084}{is}%
\TextAt{8}{107.918}{396.038}{commutative.}%
\TextAt{12}{97.159}{376.003}{Remark:}%
\TextAt{13}{154.171}{376.003}{G}%
\TextAt{8}{163.405}{376.003}{-modulations}%
\TextAt{8}{234.814}{376.003}{were}%
\TextAt{8}{262.575}{376.003}{introduced}%
\TextAt{8}{321.885}{376.003}{in}%
\TextAt{8}{336.316}{376.003}{a}%
\TextAt{8}{346.85}{376.003}{general}%
\TextAt{8}{387.953}{376.003}{context}%
\TextAt{8}{430.359}{376.003}{by}%
\TextAt{8}{447.396}{376.003}{A.}%
\TextAt{26}{464.098}{376.003}{Dress}%
\TextAt{8}{503.503}{376.003}{under}%
\TextAt{8}{79.226}{358.957}{the}%
\TextAt{8}{99.383}{358.957}{name}%
\TextAt{8}{130.598}{358.957}{of}%
\TextAt{12}{144.442}{358.957}{Mackey}%
\TextAt{12}{193.793}{358.957}{functors}%
\TextAt{8}{245.499}{358.957}{(see}%
\TextAt{8}{268.967}{358.957}{[32]).}%
\TextAt{12}{79.226}{333.055}{Exercise}%
\TextAt{12}{132.427}{333.055}{IV.6.9.}%
\TextAt{13}{176.183}{333.055}{G}%
\TextAt{8}{185.416}{333.055}{-modulations}%
\TextAt{8}{256.047}{333.055}{form}%
\TextAt{8}{283.688}{333.055}{an}%
\TextAt{8}{299.947}{333.055}{abelian}%
\TextAt{8}{340.594}{333.055}{category.}%
\TextAt{12}{79.226}{301.284}{Exercise}%
\TextAt{12}{132.98}{301.284}{IV.6.10.}%
\TextAt{8}{183.947}{301.284}{If}%
\TextAt{13}{196.12}{301.284}{A}%
\TextAt{8}{209.274}{301.284}{is}%
\TextAt{8}{221.518}{301.284}{a}%
\TextAt{13}{231.755}{301.284}{G}%
\TextAt{8}{240.988}{301.284}{-module,}%
\TextAt{8}{290.032}{301.284}{then}%
\TextAt{8}{317.172}{301.284}{the}%
\TextAt{8}{337.81}{301.284}{function}%
\TextAt{13}{384.138}{301.284}{A}%
\TextAt{9}{392.914}{301.284}{\Character{40}}%
\TextAt{13}{397.466}{301.284}{G}%
\TextAt{9}{406.7}{301.284}{\Character{47}}%
\TextAt{13}{412.553}{301.284}{G}%
\TextAt{21}{421.786}{299.491}{K}%
\TextAt{9}{429.989}{301.284}{\Character{41}}%
\TextAt{9}{438.668}{301.284}{\Character{61}}%
\TextAt{13}{451.926}{301.284}{A}%
\TextAt{21}{460.702}{305.622}{G}%
\TextAt{22}{467.3}{304.283}{K}%
\TextAt{8}{479.22}{301.284}{extends}%
\TextAt{8}{522.357}{301.284}{to}%
\TextAt{8}{79.226}{284.238}{a}%
\TextAt{13}{89.394}{284.238}{G}%
\TextAt{8}{98.627}{284.238}{-modulation}%
\TextAt{13}{165.056}{284.238}{A}%
\Vector{q 0 G 0 g 0.478 w 93.056 -491.436 m 101.831 -491.436 l S Q}%
\TextAt{8}{178.141}{284.238}{in}%
\TextAt{8}{192.206}{284.238}{such}%
\TextAt{8}{219.015}{284.238}{a}%
\TextAt{8}{229.183}{284.238}{way}%
\TextAt{8}{253.338}{284.238}{that,}%
\TextAt{8}{282.473}{284.238}{for}%
\TextAt{8}{300.771}{284.238}{an}%
\TextAt{8}{317.442}{284.238}{extension}%
\TextAt{13}{369.62}{284.238}{L}%
\TextAt{16}{381.607}{284.238}{\Character{106}}%
\TextAt{13}{388.952}{284.238}{K}%
\TextAt{8}{399.769}{284.238}{,}%
\TextAt{8}{407.434}{284.238}{the}%
\TextAt{8}{428.002}{284.238}{map}%
\TextAt{13}{454.429}{284.238}{f}%
\TextAt{3}{461.475}{288.576}{\Character{3}}%
\TextAt{9}{470.231}{284.238}{\Character{58}}%
\TextAt{13}{477.506}{284.238}{A}%
\TextAt{21}{486.281}{282.445}{K}%
\NativeMathAt{498.507}{284.238}{11.955}{\rightarrow}%
\TextAt{13}{514.485}{284.238}{A}%
\TextAt{21}{523.261}{282.445}{L}%
\TextAt{8}{529.518}{284.238}{,}%
\TextAt{8}{79.226}{267.192}{resp.}%
\TextAt{13}{110.694}{267.192}{f}%
\TextAt{3}{116.464}{265.399}{\Character{3}}%
\TextAt{9}{125.736}{267.192}{\Character{58}}%
\TextAt{13}{133.527}{267.192}{A}%
\TextAt{21}{142.302}{265.399}{L}%
\NativeMathAt{153.099}{267.192}{11.955}{\rightarrow}%
\TextAt{13}{169.594}{267.192}{A}%
\TextAt{21}{178.369}{265.399}{K}%
\TextAt{8}{186.572}{267.192}{,}%
\TextAt{8}{194.615}{267.192}{induced}%
\TextAt{8}{238.895}{267.192}{by}%
\TextAt{13}{255.869}{267.192}{f}%
\TextAt{9}{267.455}{267.192}{\Character{58}}%
\TextAt{13}{275.247}{267.192}{G}%
\TextAt{9}{284.48}{267.192}{\Character{47}}%
\TextAt{13}{290.333}{267.192}{G}%
\TextAt{21}{299.567}{265.399}{L}%
\NativeMathAt{310.364}{267.192}{11.955}{\rightarrow}%
\TextAt{13}{326.859}{267.192}{G}%
\TextAt{9}{336.092}{267.192}{\Character{47}}%
\TextAt{13}{341.945}{267.192}{G}%
\TextAt{21}{351.179}{265.399}{K}%
\TextAt{8}{359.382}{267.192}{,}%
\TextAt{8}{367.425}{267.192}{is}%
\TextAt{8}{379.904}{267.192}{the}%
\TextAt{8}{400.776}{267.192}{inclusion,}%
\TextAt{8}{453.758}{267.192}{resp.}%
\TextAt{8}{485.226}{267.192}{the}%
\TextAt{8}{506.098}{267.192}{norm}%
\TextAt{13}{79.226}{250.146}{N}%
\TextAt{21}{88.6}{248.286}{L}%
\TextAt{3}{94.359}{248.286}{\Character{106}}%
\TextAt{21}{96.711}{248.286}{K}%
\TextAt{8}{104.914}{250.146}{.}%
\TextAt{8}{97.159}{230.111}{The}%
\TextAt{8}{120.231}{230.111}{rule}%
\TextAt{13}{142.657}{230.111}{A}%
\NativeArrowAt{154.753}{230.111}{11.955}{11.955}{\mapsto}%
\TextAt{13}{170.029}{230.111}{A}%
\Vector{q 0 G 0 g 0.478 w 98.029 -545.563 m 106.804 -545.563 l S Q}%
\TextAt{8}{181.72}{230.111}{is}%
\TextAt{8}{192.501}{230.111}{an}%
\TextAt{8}{207.778}{230.111}{equivalence}%
\TextAt{8}{268.58}{230.111}{between}%
\TextAt{8}{312.789}{230.111}{the}%
\TextAt{8}{331.963}{230.111}{category}%
\TextAt{8}{378.144}{230.111}{of}%
\TextAt{13}{390.492}{230.111}{G}%
\TextAt{8}{399.726}{230.111}{-modules}%
\TextAt{8}{448.549}{230.111}{and}%
\TextAt{8}{470.329}{230.111}{the}%
\TextAt{8}{489.504}{230.111}{category}%
\TextAt{8}{79.226}{213.065}{of}%
\TextAt{13}{92.556}{213.065}{G}%
\TextAt{8}{101.79}{213.065}{-modulations}%
\TextAt{8}{172.421}{213.065}{with}%
\TextAt{8}{199.081}{213.065}{“Galois}%
\TextAt{8}{240.518}{213.065}{descent”,}%
\TextAt{8}{290.646}{213.065}{i.e.,}%
\TextAt{8}{312.751}{213.065}{of}%
\TextAt{8}{326.081}{213.065}{those}%
\TextAt{13}{356.71}{213.065}{G}%
\TextAt{8}{365.944}{213.065}{-modulations}%
\TextAt{13}{436.575}{213.065}{A}%
\TextAt{8}{449.247}{213.065}{such}%
\TextAt{8}{475.644}{213.065}{that}%
\TextAt{13}{237.474}{184.448}{f}%
\TextAt{3}{244.521}{189.384}{\Character{3}}%
\TextAt{9}{252.574}{184.448}{\Character{58}}%
\TextAt{13}{259.146}{184.448}{A}%
\TextAt{9}{267.922}{184.448}{\Character{40}}%
\TextAt{13}{272.474}{184.448}{Y}%
\TextAt{9}{281.867}{184.448}{\Character{41}}%
\NativeArrowAt{289.74}{184.448}{11.955}{19.26}{\longrightarrow}%
\TextAt{13}{312.322}{184.448}{A}%
\TextAt{9}{321.098}{184.448}{\Character{40}}%
\TextAt{13}{325.65}{184.448}{X}%
\TextAt{9}{336.305}{184.448}{\Character{41}}%
\TextAt{21}{340.857}{189.384}{G}%
\TextAt{17}{347.456}{189.384}{\Character{40}}%
\TextAt{21}{350.749}{189.384}{X}%
\TextAt{3}{358.336}{189.384}{\Character{106}}%
\TextAt{21}{360.689}{189.384}{Y}%
\TextAt{17}{367.478}{189.384}{\Character{41}}%
\TextAt{13}{371.27}{184.448}{\Character{59}}%
\TextAt{8}{79.226}{155.83}{for}%
\TextAt{8}{97.111}{155.83}{every}%
\TextAt{8}{128.003}{155.83}{Galois}%
\TextAt{8}{163.917}{155.83}{mapping}%
\TextAt{13}{212.048}{155.83}{f}%
\TextAt{9}{222.415}{155.83}{\Character{58}}%
\TextAt{13}{228.988}{155.83}{X}%
\NativeMathAt{242.964}{155.83}{11.955}{\rightarrow}%
\TextAt{13}{258.24}{155.83}{Y}%
\TextAt{8}{267.633}{155.83}{,}%
\TextAt{8}{274.782}{155.83}{is}%
\TextAt{8}{286.546}{155.83}{an}%
\TextAt{8}{302.805}{155.83}{isomorphism.}%
\TextAt{12}{79.226}{124.059}{Exercise}%
\TextAt{12}{132.337}{124.059}{IV.6.11.}%
\TextAt{13}{182.745}{124.059}{G}%
\TextAt{8}{191.979}{124.059}{-modulations}%
\TextAt{8}{262.532}{124.059}{are}%
\TextAt{8}{281.964}{124.059}{explicitly}%
\TextAt{8}{332.612}{124.059}{given}%
\TextAt{8}{363.103}{124.059}{by}%
\TextAt{8}{379.283}{124.059}{the}%
\TextAt{8}{399.362}{124.059}{following}%
\TextAt{8}{448.718}{124.059}{data.}%
\TextAt{8}{479.915}{124.059}{Let}%
\TextAt{13}{500.806}{124.059}{B}%
\TextAt{17}{509.699}{122.266}{0}%
\TextAt{9}{514.431}{124.059}{\Character{40}}%
\TextAt{13}{518.984}{124.059}{G}%
\TextAt{9}{528.217}{124.059}{\Character{41}}%
\TextAt{8}{79.226}{107.013}{be}%
\TextAt{8}{94.077}{107.013}{the}%
\TextAt{8}{113.159}{107.013}{category}%
\TextAt{8}{159.249}{107.013}{whose}%
\TextAt{8}{192.702}{107.013}{objects}%
\TextAt{8}{231.355}{107.013}{are}%
\TextAt{8}{249.792}{107.013}{the}%
\TextAt{13}{268.875}{107.013}{G}%
\TextAt{8}{278.108}{107.013}{-sets}%
\TextAt{13}{303.814}{107.013}{G}%
\TextAt{9}{313.048}{107.013}{\Character{47}}%
\TextAt{13}{318.901}{107.013}{U}%
\TextAt{8}{328.1}{107.013}{,}%
\TextAt{8}{334.39}{107.013}{where}%
\TextAt{13}{367.126}{107.013}{U}%
\TextAt{8}{379.148}{107.013}{varies}%
\TextAt{8}{410.987}{107.013}{over}%
\TextAt{8}{434.96}{107.013}{the}%
\TextAt{8}{454.042}{107.013}{open}%
\TextAt{8}{481.255}{107.013}{subgroups}%
\TextAt{8}{79.226}{89.967}{of}%
\TextAt{13}{92.66}{89.967}{G}%
\TextAt{8}{101.894}{89.967}{,}%
\TextAt{8}{109.173}{89.967}{and}%
\TextAt{8}{132.039}{89.967}{whose}%
\TextAt{8}{166.67}{89.967}{morphisms}%
\TextAt{8}{226.084}{89.967}{are}%
\TextAt{8}{245.698}{89.967}{just}%
\TextAt{8}{268.947}{89.967}{the}%
\TextAt{8}{289.208}{89.967}{projections}%
\TextAt{13}{349.207}{89.967}{\Character{25}}%
\TextAt{9}{359.774}{89.967}{\Character{58}}%
\TextAt{13}{366.524}{89.967}{G}%
\TextAt{9}{375.758}{89.967}{\Character{47}}%
\TextAt{13}{381.611}{89.967}{U}%
\NativeMathAt{394.308}{89.967}{11.955}{\rightarrow}%
\TextAt{13}{409.761}{89.967}{G}%
\TextAt{9}{418.995}{89.967}{\Character{47}}%
\TextAt{13}{424.848}{89.967}{V}%
\TextAt{8}{438.275}{89.967}{for}%
\TextAt{13}{456.263}{89.967}{U}%
\NativeMathAt{468.961}{89.967}{11.955}{\subseteq}%
\TextAt{13}{481.757}{89.967}{V}%
\TextAt{8}{491.183}{89.967}{,}%
\TextAt{8}{498.462}{89.967}{as}%
\TextAt{8}{512.936}{89.967}{well}%
\TextAt{8}{79.226}{72.921}{as}%
\TextAt{8}{93.596}{72.921}{the}%
\TextAt{8}{113.753}{72.921}{maps}%
\TextAt{13}{144.382}{72.921}{c}%
\TextAt{9}{149.42}{72.921}{\Character{40}}%
\TextAt{13}{153.972}{72.921}{\Character{27}}%
\TextAt{9}{161.055}{72.921}{\Character{41}}%
\TextAt{9}{168.91}{72.921}{\Character{58}}%
\TextAt{13}{175.5}{72.921}{G}%
\TextAt{9}{184.734}{72.921}{\Character{47}}%
\TextAt{13}{190.587}{72.921}{U}%
\NativeMathAt{203.107}{72.921}{11.955}{\rightarrow}%
\TextAt{13}{218.383}{72.921}{G}%
\TextAt{9}{227.616}{72.921}{\Character{47}}%
\TextAt{13}{233.469}{72.921}{\Character{27}}%
\TextAt{13}{240.535}{72.921}{U}%
\TextAt{13}{249.74}{72.921}{\Character{27}}%
\TextAt{3}{256.833}{77.26}{\Character{0}}%
\TextAt{17}{263.42}{77.26}{1}%
\TextAt{8}{268.152}{72.921}{,}%
\TextAt{13}{275.301}{72.921}{\Character{28}}%
\TextAt{13}{281.709}{72.921}{U}%
\NativeArrowAt{294.24}{72.921}{11.955}{11.955}{\mapsto}%
\TextAt{13}{309.516}{72.921}{\Character{28}}%
\TextAt{13}{315.924}{72.921}{U}%
\TextAt{13}{325.13}{72.921}{\Character{27}}%
\TextAt{3}{332.216}{77.26}{\Character{0}}%
\TextAt{17}{338.803}{77.26}{1}%
\TextAt{9}{346.856}{72.921}{\Character{61}}%
\TextAt{13}{359.282}{72.921}{\Character{28}}%
\TextAt{13}{365.69}{72.921}{\Character{27}}%
\TextAt{3}{372.783}{77.26}{\Character{0}}%
\TextAt{17}{379.369}{77.26}{1}%
\TextAt{9}{384.102}{72.921}{\Character{40}}%
\TextAt{13}{388.654}{72.921}{\Character{27}}%
\TextAt{13}{395.72}{72.921}{U}%
\TextAt{13}{404.925}{72.921}{\Character{27}}%
\TextAt{3}{412.018}{77.26}{\Character{0}}%
\TextAt{17}{418.604}{77.26}{1}%
\TextAt{9}{423.337}{72.921}{\Character{41}}%
\TextAt{8}{427.889}{72.921}{,}%
\TextAt{8}{435.038}{72.921}{for}%
\TextAt{13}{452.923}{72.921}{\Character{27}}%
\NativeMathAt{463.326}{72.921}{11.955}{\in}%
\TextAt{13}{474.617}{72.921}{G}%
\TextAt{8}{483.851}{72.921}{.}%
\TextAt{8}{297.211}{43.033}{282}%
\Destination{79.23}{552.005}{nav.Item.97}%
\Destination{79.23}{532.075}{nav.Item.98}%
\Destination{78.23}{809.815}{nav.page.282}%
\Destination{79.226}{347.055}{nav.ex.4.6.9}%
\Destination{79.226}{315.255}{nav.ex.4.6.10}%
\Destination{79.226}{138.055}{nav.ex.4.6.11}%
\Destination{79.36}{245.155}{nav.print.308}%
\Hotlink{268.467}{356.755}{19.22}{11.047}{nav.cite.N032}%
}

% Page 295 (body page 283)
\Page{%
\TextAt{2}{79.226}{796.943}{IV.}%
\TextAt{2}{98.012}{796.943}{Abstract}%
\TextAt{2}{143.252}{796.943}{Class}%
\TextAt{2}{171.844}{796.943}{Field}%
\TextAt{2}{199.575}{796.943}{Theory}%
\Vector{q 0 G 0 g 0.299 w 7.226 18.729 m 460.769 18.729 l S Q}%
\TextAt{8}{97.159}{761.124}{Let}%
\TextAt{13}{119.531}{761.124}{A}%
\TextAt{9}{134.018}{761.124}{\Character{61}}%
\TextAt{9}{148.819}{761.124}{\Character{40}}%
\TextAt{13}{153.386}{761.124}{A}%
\TextAt{3}{162.161}{765.462}{\Character{3}}%
\TextAt{13}{166.894}{761.124}{\Character{59}}%
\TextAt{13}{172.118}{761.124}{A}%
\TextAt{3}{180.913}{759.33}{\Character{3}}%
\TextAt{9}{185.645}{761.124}{\Character{41}}%
\TextAt{9}{195.891}{761.124}{\Character{58}}%
\TextAt{13}{204.872}{761.124}{B}%
\TextAt{17}{213.765}{759.33}{0}%
\TextAt{9}{218.497}{761.124}{\Character{40}}%
\TextAt{13}{223.05}{761.124}{G}%
\TextAt{9}{232.283}{761.124}{\Character{41}}%
\NativeMathAt{242.547}{761.124}{11.955}{\rightarrow}%
\TextAt{9}{260.213}{761.124}{\Character{40}}%
\TextAt{13}{264.766}{761.124}{ab}%
\TextAt{9}{275.888}{761.124}{\Character{41}}%
\TextAt{8}{285.74}{761.124}{be}%
\TextAt{8}{303.067}{761.124}{a}%
\TextAt{8}{314.226}{761.124}{double}%
\TextAt{8}{353.347}{761.124}{functor}%
\TextAt{8}{395.397}{761.124}{and}%
\TextAt{8}{419.563}{761.124}{for}%
\TextAt{13}{438.85}{761.124}{\Character{25}}%
\TextAt{9}{451.631}{761.124}{\Character{58}}%
\TextAt{13}{460.594}{761.124}{G}%
\TextAt{9}{469.827}{761.124}{\Character{47}}%
\TextAt{13}{475.68}{761.124}{U}%
\NativeMathAt{490.591}{761.124}{11.955}{\rightarrow}%
\TextAt{13}{508.257}{761.124}{G}%
\TextAt{9}{517.491}{761.124}{\Character{47}}%
\TextAt{13}{523.344}{761.124}{V}%
\TextAt{9}{79.226}{744.078}{\Character{40}}%
\TextAt{13}{83.778}{744.078}{U}%
\NativeMathAt{96.298}{744.078}{11.955}{\subseteq}%
\TextAt{13}{108.918}{744.078}{V}%
\TextAt{9}{118.343}{744.078}{\Character{41}}%
\TextAt{8}{122.896}{744.078}{,}%
\TextAt{8}{130.045}{744.078}{resp.}%
\TextAt{13}{159.367}{744.078}{c}%
\TextAt{9}{164.405}{744.078}{\Character{40}}%
\TextAt{13}{168.957}{744.078}{\Character{27}}%
\TextAt{9}{176.04}{744.078}{\Character{41}}%
\TextAt{9}{183.895}{744.078}{\Character{58}}%
\TextAt{13}{190.485}{744.078}{G}%
\TextAt{9}{199.719}{744.078}{\Character{47}}%
\TextAt{13}{205.572}{744.078}{U}%
\NativeMathAt{218.092}{744.078}{11.955}{\rightarrow}%
\TextAt{13}{233.368}{744.078}{G}%
\TextAt{9}{242.601}{744.078}{\Character{47}}%
\TextAt{13}{248.454}{744.078}{\Character{27}}%
\TextAt{13}{255.52}{744.078}{U}%
\TextAt{13}{264.725}{744.078}{\Character{27}}%
\TextAt{3}{271.818}{748.416}{\Character{0}}%
\TextAt{17}{278.405}{748.416}{1}%
\TextAt{9}{287.034}{744.078}{\Character{40}}%
\TextAt{13}{291.587}{744.078}{\Character{27}}%
\NativeMathAt{301.99}{744.078}{11.955}{\in}%
\TextAt{13}{313.281}{744.078}{G}%
\TextAt{9}{322.515}{744.078}{\Character{41}}%
\TextAt{8}{327.067}{744.078}{,}%
\TextAt{8}{334.216}{744.078}{define}%
\TextAt{8}{190.214}{715.077}{Ind}%
\TextAt{21}{207.442}{720.219}{U}%
\TextAt{21}{207.442}{712.121}{V}%
\TextAt{9}{218.074}{715.077}{\Character{61}}%
\TextAt{13}{230.5}{715.077}{A}%
\TextAt{3}{239.275}{713.283}{\Character{3}}%
\TextAt{9}{244.007}{715.077}{\Character{40}}%
\TextAt{13}{248.56}{715.077}{\Character{25}}%
\TextAt{9}{255.629}{715.077}{\Character{41}}%
\TextAt{9}{263.484}{715.077}{\Character{58}}%
\TextAt{13}{270.075}{715.077}{A}%
\TextAt{9}{278.85}{715.077}{\Character{40}}%
\TextAt{13}{283.402}{715.077}{G}%
\TextAt{9}{292.636}{715.077}{\Character{47}}%
\TextAt{13}{298.489}{715.077}{U}%
\TextAt{9}{307.688}{715.077}{\Character{41}}%
\NativeArrowAt{315.561}{715.077}{11.955}{19.26}{\longrightarrow}%
\TextAt{13}{338.143}{715.077}{A}%
\TextAt{9}{346.918}{715.077}{\Character{40}}%
\TextAt{13}{351.471}{715.077}{G}%
\TextAt{9}{360.704}{715.077}{\Character{47}}%
\TextAt{13}{366.557}{715.077}{V}%
\TextAt{9}{375.983}{715.077}{\Character{41}}%
\TextAt{13}{380.535}{715.077}{\Character{59}}%
\TextAt{8}{189.007}{695.042}{Res}%
\TextAt{21}{207.442}{700.052}{V}%
\TextAt{21}{207.442}{692.086}{U}%
\TextAt{9}{218.074}{695.042}{\Character{61}}%
\TextAt{13}{230.5}{695.042}{A}%
\TextAt{3}{239.275}{699.978}{\Character{3}}%
\TextAt{9}{244.007}{695.042}{\Character{40}}%
\TextAt{13}{248.56}{695.042}{\Character{25}}%
\TextAt{9}{255.629}{695.042}{\Character{41}}%
\TextAt{9}{263.484}{695.042}{\Character{58}}%
\TextAt{13}{270.075}{695.042}{A}%
\TextAt{9}{278.85}{695.042}{\Character{40}}%
\TextAt{13}{283.402}{695.042}{G}%
\TextAt{9}{292.636}{695.042}{\Character{47}}%
\TextAt{13}{298.489}{695.042}{V}%
\TextAt{9}{307.914}{695.042}{\Character{41}}%
\NativeArrowAt{315.788}{695.042}{11.955}{19.26}{\longrightarrow}%
\TextAt{13}{338.37}{695.042}{A}%
\TextAt{9}{347.145}{695.042}{\Character{40}}%
\TextAt{13}{351.697}{695.042}{G}%
\TextAt{9}{360.931}{695.042}{\Character{47}}%
\TextAt{13}{366.784}{695.042}{U}%
\TextAt{9}{375.983}{695.042}{\Character{41}}%
\TextAt{13}{380.535}{695.042}{\Character{59}}%
\TextAt{13}{188.796}{675.007}{c}%
\TextAt{9}{193.834}{675.007}{\Character{40}}%
\TextAt{13}{198.386}{675.007}{\Character{27}}%
\TextAt{9}{205.469}{675.007}{\Character{41}}%
\TextAt{3}{210.021}{673.214}{\Character{3}}%
\TextAt{9}{218.074}{675.007}{\Character{61}}%
\TextAt{13}{230.5}{675.007}{A}%
\TextAt{3}{239.275}{673.214}{\Character{3}}%
\TextAt{9}{244.007}{675.007}{\Character{40}}%
\TextAt{13}{248.56}{675.007}{c}%
\TextAt{9}{253.598}{675.007}{\Character{40}}%
\TextAt{13}{258.15}{675.007}{\Character{27}}%
\TextAt{9}{265.232}{675.007}{\Character{41}\Character{41}}%
\TextAt{9}{277.63}{675.007}{\Character{58}}%
\TextAt{13}{284.23}{675.007}{A}%
\TextAt{9}{293.006}{675.007}{\Character{40}}%
\TextAt{13}{297.558}{675.007}{G}%
\TextAt{9}{306.792}{675.007}{\Character{47}}%
\TextAt{13}{312.645}{675.007}{U}%
\TextAt{9}{321.844}{675.007}{\Character{41}}%
\NativeArrowAt{329.717}{675.007}{11.955}{19.26}{\longrightarrow}%
\TextAt{13}{352.299}{675.007}{A}%
\TextAt{9}{361.074}{675.007}{\Character{40}}%
\TextAt{13}{365.626}{675.007}{G}%
\TextAt{9}{374.86}{675.007}{\Character{47}}%
\TextAt{13}{380.713}{675.007}{\Character{27}}%
\TextAt{13}{387.779}{675.007}{U}%
\TextAt{13}{396.984}{675.007}{\Character{27}}%
\TextAt{3}{404.077}{679.943}{\Character{0}}%
\TextAt{17}{410.663}{679.943}{1}%
\TextAt{9}{415.396}{675.007}{\Character{41}}%
\TextAt{13}{419.948}{675.007}{\Character{58}}%
\TextAt{8}{79.226}{646.006}{If}%
\TextAt{8}{90.918}{646.006}{for}%
\TextAt{8}{108.803}{646.006}{any}%
\TextAt{8}{130.92}{646.006}{three}%
\TextAt{8}{160.832}{646.006}{open}%
\TextAt{8}{189.118}{646.006}{subgroups}%
\TextAt{13}{244.53}{646.006}{U\Character{59}}%
\TextAt{13}{257.657}{646.006}{V}%
\NativeMathAt{270.419}{646.006}{11.955}{\subseteq}%
\TextAt{13}{283.039}{646.006}{W}%
\TextAt{8}{299.613}{646.006}{of}%
\TextAt{13}{312.943}{646.006}{G}%
\TextAt{8}{322.177}{646.006}{,}%
\TextAt{8}{329.326}{646.006}{one}%
\TextAt{8}{350.786}{646.006}{has}%
\TextAt{8}{371.66}{646.006}{the}%
\TextAt{12}{391.816}{646.006}{induction}%
\TextAt{12}{451.616}{646.006}{formula}%
\TextAt{8}{167.38}{617.005}{Res}%
\TextAt{21}{185.815}{622.015}{W}%
\TextAt{21}{185.815}{614.049}{U}%
\TextAt{16}{197.472}{617.005}{\Character{14}}%
\TextAt{8}{205.442}{617.005}{Ind}%
\TextAt{21}{222.669}{622.147}{V}%
\TextAt{21}{222.669}{614.049}{W}%
\TextAt{9}{235.654}{617.005}{\Character{61}}%
\TextAt{20}{254.958}{628.362}{\Character{88}}%
\TextAt{21}{248.079}{602.459}{U}%
\TextAt{3}{254.657}{602.459}{\Character{110}}%
\TextAt{21}{258.891}{602.459}{W}%
\TextAt{17}{268.057}{602.459}{\Character{47}}%
\TextAt{21}{272.291}{602.459}{V}%
\TextAt{8}{281.097}{617.005}{Ind}%
\TextAt{21}{298.325}{622.147}{U}%
\TextAt{3}{304.903}{622.147}{\Character{92}}%
\TextAt{21}{310.548}{622.147}{\Character{27}}%
\TextAt{21}{315.665}{622.147}{V}%
\TextAt{21}{322.479}{622.147}{\Character{27}}%
\TextAt{23}{327.612}{624.96}{\Character{0}}%
\TextAt{24}{333.368}{624.96}{1}%
\TextAt{21}{298.325}{614.049}{U}%
\TextAt{16}{340.01}{617.005}{\Character{14}}%
\TextAt{13}{345.988}{617.005}{c}%
\TextAt{9}{351.026}{617.005}{\Character{40}}%
\TextAt{13}{355.578}{617.005}{\Character{27}}%
\TextAt{9}{362.66}{617.005}{\Character{41}}%
\TextAt{3}{367.213}{615.212}{\Character{3}}%
\TextAt{16}{374.602}{617.005}{\Character{14}}%
\TextAt{8}{383.236}{617.005}{Res}%
\TextAt{21}{401.671}{622.015}{V}%
\TextAt{21}{401.671}{613.981}{V}%
\TextAt{3}{408.484}{613.981}{\Character{92}}%
\TextAt{21}{414.13}{613.981}{\Character{27}}%
\TextAt{23}{419.255}{616.25}{\Character{0}}%
\TextAt{24}{425.011}{616.25}{1}%
\TextAt{21}{429.162}{613.981}{U}%
\TextAt{21}{435.729}{613.981}{\Character{27}}%
\TextAt{13}{441.364}{617.005}{\Character{59}}%
\TextAt{8}{79.226}{577.353}{then}%
\TextAt{13}{105.886}{577.353}{A}%
\TextAt{8}{118.559}{577.353}{extends}%
\TextAt{8}{161.215}{577.353}{uniquely}%
\TextAt{8}{208.689}{577.353}{to}%
\TextAt{8}{222.999}{577.353}{a}%
\TextAt{13}{232.755}{577.353}{G}%
\TextAt{8}{241.988}{577.353}{-modulation}%
\TextAt{13}{308.005}{577.353}{A}%
\TextAt{9}{320.101}{577.353}{\Character{58}}%
\TextAt{13}{326.673}{577.353}{B}%
\TextAt{9}{336.17}{577.353}{\Character{40}}%
\TextAt{13}{340.722}{577.353}{G}%
\TextAt{9}{349.956}{577.353}{\Character{41}}%
\NativeMathAt{357.829}{577.353}{11.955}{\rightarrow}%
\TextAt{9}{373.105}{577.353}{\Character{40}}%
\TextAt{13}{377.657}{577.353}{ab}%
\TextAt{9}{388.779}{577.353}{\Character{41}}%
\TextAt{8}{393.332}{577.353}{.}%
\TextAt{12}{97.159}{557.319}{Hint:}%
\TextAt{8}{132.698}{557.319}{If}%
\TextAt{13}{144.39}{557.319}{X}%
\TextAt{8}{158.942}{557.319}{is}%
\TextAt{8}{170.706}{557.319}{an}%
\TextAt{8}{186.965}{557.319}{arbitrary}%
\TextAt{8}{236.735}{557.319}{finite}%
\TextAt{13}{266.646}{557.319}{G}%
\TextAt{8}{275.88}{557.319}{-set,}%
\TextAt{8}{301.297}{557.319}{then}%
\TextAt{8}{327.957}{557.319}{the}%
\TextAt{8}{348.113}{557.319}{disjoint}%
\TextAt{8}{389.801}{557.319}{union}%
\TextAt{13}{255.966}{528.317}{A}%
\TextAt{21}{264.741}{526.524}{X}%
\TextAt{9}{276.147}{528.317}{\Character{61}}%
\TextAt{20}{289.934}{539.675}{\Character{97}}%
\TextAt{21}{288.573}{514.303}{x}%
\NativeMathAt{293.339}{514.303}{7.97}{\in}%
\TextAt{21}{298.985}{514.303}{X}%
\TextAt{13}{308.564}{528.317}{A}%
\TextAt{9}{317.34}{528.317}{\Character{40}}%
\TextAt{13}{321.892}{528.317}{G}%
\TextAt{9}{331.126}{528.317}{\Character{47}}%
\TextAt{13}{336.979}{528.317}{G}%
\TextAt{21}{346.212}{526.524}{x}%
\TextAt{9}{351.477}{528.317}{\Character{41}}%
\TextAt{8}{79.226}{490.687}{is}%
\TextAt{8}{90.99}{490.687}{again}%
\TextAt{8}{122.217}{490.687}{a}%
\TextAt{13}{131.972}{490.687}{G}%
\TextAt{8}{141.206}{490.687}{-set,}%
\TextAt{8}{166.623}{490.687}{because}%
\TextAt{13}{209.924}{490.687}{c}%
\TextAt{9}{214.962}{490.687}{\Character{40}}%
\TextAt{13}{219.515}{490.687}{\Character{27}}%
\TextAt{9}{226.597}{490.687}{\Character{41}}%
\TextAt{3}{231.149}{488.894}{\Character{3}}%
\TextAt{13}{235.882}{490.687}{A}%
\TextAt{9}{244.657}{490.687}{\Character{40}}%
\TextAt{13}{249.209}{490.687}{G}%
\TextAt{9}{258.443}{490.687}{\Character{47}}%
\TextAt{13}{264.296}{490.687}{G}%
\TextAt{21}{273.53}{488.894}{x}%
\TextAt{9}{278.795}{490.687}{\Character{41}}%
\TextAt{9}{286.65}{490.687}{\Character{61}}%
\TextAt{13}{299.093}{490.687}{A}%
\TextAt{9}{307.869}{490.687}{\Character{40}}%
\TextAt{13}{312.421}{490.687}{G}%
\TextAt{9}{321.655}{490.687}{\Character{47}}%
\TextAt{13}{327.508}{490.687}{G}%
\TextAt{21}{336.741}{488.894}{\Character{27}}%
\TextAt{21}{341.858}{488.894}{x}%
\TextAt{9}{347.131}{490.687}{\Character{41}}%
\TextAt{8}{351.684}{490.687}{.}%
\TextAt{8}{360.132}{490.687}{Define}%
\TextAt{13}{396.38}{490.687}{A}%
\TextAt{9}{405.156}{490.687}{\Character{40}}%
\TextAt{13}{409.708}{490.687}{X}%
\TextAt{9}{420.363}{490.687}{\Character{41}}%
\TextAt{8}{428.813}{490.687}{to}%
\TextAt{8}{443.123}{490.687}{be}%
\TextAt{8}{459.047}{490.687}{the}%
\TextAt{8}{479.204}{490.687}{group}%
\TextAt{13}{247.013}{461.686}{A}%
\TextAt{9}{255.789}{461.686}{\Character{40}}%
\TextAt{13}{260.341}{461.686}{X}%
\TextAt{9}{270.996}{461.686}{\Character{41}}%
\TextAt{9}{278.851}{461.686}{\Character{61}}%
\TextAt{8}{291.295}{461.686}{Hom}%
\TextAt{21}{315.683}{459.892}{X}%
\TextAt{9}{323.768}{461.686}{\Character{40}}%
\TextAt{13}{328.321}{461.686}{X}%
\TextAt{13}{338.315}{461.686}{\Character{59}}%
\TextAt{13}{343.54}{461.686}{A}%
\TextAt{21}{352.345}{459.892}{X}%
\TextAt{9}{360.43}{461.686}{\Character{41}}%
\TextAt{8}{79.226}{432.685}{of}%
\TextAt{8}{92.556}{432.685}{all}%
\TextAt{13}{108.815}{432.685}{G}%
\TextAt{8}{118.049}{432.685}{-equivariant}%
\TextAt{8}{182.774}{432.685}{sections}%
\TextAt{13}{226.47}{432.685}{X}%
\NativeMathAt{240.446}{432.685}{11.955}{\rightarrow}%
\TextAt{13}{255.722}{432.685}{A}%
\TextAt{21}{264.498}{430.891}{X}%
\TextAt{8}{276.48}{432.685}{of}%
\TextAt{8}{289.81}{432.685}{the}%
\TextAt{8}{309.966}{432.685}{projection}%
\TextAt{13}{365.247}{432.685}{A}%
\TextAt{21}{374.023}{430.891}{X}%
\NativeMathAt{385.428}{432.685}{11.955}{\rightarrow}%
\TextAt{13}{400.704}{432.685}{X}%
\TextAt{8}{411.36}{432.685}{.}%
\TextAt{12}{79.226}{400.695}{Exercise}%
\TextAt{12}{132.427}{400.695}{IV.6.12.}%
\TextAt{8}{182.913}{400.695}{The}%
\TextAt{8}{206.967}{400.695}{function}%
\TextAt{13}{252.815}{400.695}{\Character{25}}%
\TextAt{7}{259.884}{405.033}{ab}%
\TextAt{9}{269.317}{400.695}{\Character{40}}%
\TextAt{13}{273.869}{400.695}{G}%
\TextAt{9}{283.103}{400.695}{\Character{47}}%
\TextAt{13}{288.956}{400.695}{G}%
\TextAt{21}{298.19}{398.901}{K}%
\TextAt{9}{306.392}{400.695}{\Character{41}}%
\TextAt{9}{314.247}{400.695}{\Character{61}}%
\TextAt{13}{326.691}{400.695}{G}%
\TextAt{7}{335.925}{405.033}{ab}%
\TextAt{21}{335.925}{397.594}{K}%
\TextAt{8}{349.255}{400.695}{extends}%
\TextAt{8}{391.911}{400.695}{to}%
\TextAt{8}{406.221}{400.695}{a}%
\TextAt{13}{415.976}{400.695}{G}%
\TextAt{8}{425.21}{400.695}{-modulation}%
\TextAt{13}{245.318}{371.694}{\Character{25}}%
\TextAt{7}{252.387}{376.63}{ab}%
\TextAt{9}{265.141}{371.693}{\Character{58}}%
\TextAt{13}{271.713}{371.694}{B}%
\TextAt{9}{281.209}{371.693}{\Character{40}}%
\TextAt{13}{285.762}{371.694}{G}%
\TextAt{9}{294.995}{371.693}{\Character{41}}%
\NativeArrowAt{302.869}{371.694}{11.955}{19.26}{\longrightarrow}%
\TextAt{9}{325.451}{371.693}{\Character{40}}%
\TextAt{13}{330.003}{371.694}{pr}%
\TextAt{13}{341.468}{371.694}{o}%
\TextAt{8}{347.106}{371.694}{-}%
\TextAt{13}{351.003}{371.694}{ab}%
\TextAt{9}{362.125}{371.693}{\Character{41}}%
\TextAt{8}{79.226}{342.692}{into}%
\TextAt{8}{104.012}{342.692}{the}%
\TextAt{8}{125.212}{342.692}{category}%
\TextAt{8}{173.419}{342.692}{of}%
\TextAt{8}{187.792}{342.692}{pro-abelian}%
\TextAt{8}{250.297}{342.692}{groups.}%
\TextAt{8}{295.769}{342.692}{Thus,}%
\TextAt{8}{329.973}{342.692}{for}%
\TextAt{8}{348.902}{342.692}{an}%
\TextAt{8}{366.204}{342.692}{extension}%
\TextAt{13}{419.013}{342.692}{L}%
\TextAt{16}{432.077}{342.692}{\Character{106}}%
\TextAt{13}{440.497}{342.692}{K}%
\TextAt{8}{451.314}{342.692}{,}%
\TextAt{8}{459.768}{342.692}{the}%
\TextAt{8}{480.968}{342.692}{maps}%
\TextAt{13}{512.64}{342.692}{f}%
\TextAt{3}{519.687}{347.031}{\Character{3}}%
\TextAt{9}{529.518}{342.692}{\Character{58}}%
\TextAt{13}{79.226}{325.646}{G}%
\TextAt{7}{88.46}{329.985}{ab}%
\TextAt{21}{88.46}{322.546}{K}%
\NativeMathAt{103.166}{325.646}{11.955}{\rightarrow}%
\TextAt{13}{120.396}{325.646}{G}%
\TextAt{7}{129.629}{329.985}{ab}%
\TextAt{21}{129.629}{322.546}{L}%
\TextAt{8}{139.062}{325.646}{,}%
\TextAt{8}{147.644}{325.646}{resp.}%
\TextAt{13}{180.404}{325.646}{f}%
\TextAt{3}{186.174}{323.853}{\Character{3}}%
\TextAt{9}{196.18}{325.646}{\Character{58}}%
\TextAt{13}{204.706}{325.646}{G}%
\TextAt{7}{213.94}{329.985}{ab}%
\TextAt{21}{213.94}{322.546}{L}%
\NativeMathAt{228.646}{325.646}{11.955}{\rightarrow}%
\TextAt{13}{245.876}{325.646}{G}%
\TextAt{7}{255.109}{329.985}{ab}%
\TextAt{21}{255.109}{322.546}{K}%
\TextAt{8}{264.542}{325.646}{,}%
\TextAt{8}{273.124}{325.646}{induced}%
\TextAt{8}{317.834}{325.646}{by}%
\TextAt{13}{335.24}{325.646}{f}%
\TextAt{9}{347.56}{325.646}{\Character{58}}%
\TextAt{13}{356.086}{325.646}{G}%
\TextAt{9}{365.319}{325.646}{\Character{47}}%
\TextAt{13}{371.172}{325.646}{G}%
\TextAt{21}{380.406}{323.853}{L}%
\NativeMathAt{391.937}{325.646}{11.955}{\rightarrow}%
\TextAt{13}{409.166}{325.646}{G}%
\TextAt{9}{418.4}{325.646}{\Character{47}}%
\TextAt{13}{424.253}{325.646}{G}%
\TextAt{21}{433.487}{323.853}{K}%
\TextAt{8}{446.733}{325.646}{are}%
\TextAt{8}{467.39}{325.646}{given}%
\TextAt{8}{499.105}{325.646}{by}%
\TextAt{8}{516.51}{325.646}{the}%
\TextAt{8}{79.226}{308.6}{transfer,}%
\TextAt{8}{125.792}{308.6}{resp.}%
\TextAt{8}{155.114}{308.6}{the}%
\TextAt{8}{175.27}{308.6}{inclusion}%
\TextAt{13}{224.107}{308.6}{G}%
\TextAt{21}{233.34}{306.807}{L}%
\NativeArrowAt{242.918}{308.6}{11.955}{13.215}{\hookrightarrow}%
\TextAt{13}{259.454}{308.6}{G}%
\TextAt{21}{268.687}{306.807}{K}%
\TextAt{8}{276.89}{308.6}{.}%
\TextAt{12}{79.226}{276.61}{Exercise}%
\TextAt{12}{132.427}{276.61}{IV.6.13.}%
\TextAt{8}{182.913}{276.61}{Let}%
\TextAt{13}{203.883}{276.61}{A}%
\TextAt{8}{216.555}{276.61}{be}%
\TextAt{8}{232.48}{276.61}{a}%
\TextAt{13}{242.235}{276.61}{G}%
\TextAt{8}{251.469}{276.61}{-modulation.}%
\TextAt{8}{322.036}{276.61}{For}%
\TextAt{8}{343.006}{276.61}{every}%
\TextAt{8}{373.898}{276.61}{connected}%
\TextAt{8}{428.521}{276.61}{finite}%
\TextAt{13}{458.433}{276.61}{G}%
\TextAt{8}{467.666}{276.61}{-set}%
\TextAt{13}{489.831}{276.61}{X}%
\TextAt{8}{500.486}{276.61}{,}%
\TextAt{8}{507.636}{276.61}{let}%
\TextAt{13}{250.394}{247.609}{N}%
\TextAt{13}{261.01}{247.609}{A}%
\TextAt{9}{269.792}{247.609}{\Character{40}}%
\TextAt{13}{274.345}{247.609}{X}%
\TextAt{9}{285}{247.609}{\Character{41}}%
\TextAt{9}{292.855}{247.609}{\Character{61}}%
\TextAt{20}{305.298}{258.967}{\Character{92}}%
\TextAt{21}{309.467}{233.507}{f}%
\TextAt{13}{320.574}{247.609}{f}%
\TextAt{3}{326.344}{245.816}{\Character{3}}%
\TextAt{13}{331.077}{247.609}{A}%
\TextAt{9}{339.852}{247.609}{\Character{40}}%
\TextAt{13}{344.404}{247.609}{Y}%
\TextAt{9}{353.797}{247.609}{\Character{41}}%
\TextAt{13}{358.35}{247.609}{\Character{59}}%
\TextAt{8}{79.226}{209.393}{where}%
\TextAt{8}{113.167}{209.393}{the}%
\TextAt{8}{133.455}{209.393}{intersection}%
\TextAt{8}{196.411}{209.393}{is}%
\TextAt{8}{208.306}{209.393}{taken}%
\TextAt{8}{240.31}{209.393}{over}%
\TextAt{8}{265.487}{209.393}{all}%
\TextAt{8}{281.878}{209.393}{Galois}%
\TextAt{8}{317.923}{209.393}{maps}%
\TextAt{13}{348.683}{209.393}{f}%
\TextAt{9}{359.275}{209.393}{\Character{58}}%
\TextAt{13}{366.071}{209.393}{Y}%
\NativeMathAt{379.009}{209.393}{11.955}{\rightarrow}%
\TextAt{13}{394.509}{209.393}{X}%
\TextAt{8}{405.164}{209.393}{.}%
\TextAt{8}{414.007}{209.393}{Show}%
\TextAt{8}{445.031}{209.393}{that}%
\TextAt{8}{470.531}{209.393}{the}%
\TextAt{8}{490.819}{209.393}{function}%
\TextAt{13}{79.226}{192.347}{N}%
\TextAt{13}{89.842}{192.347}{A}%
\TextAt{9}{98.624}{192.348}{\Character{40}}%
\TextAt{13}{103.176}{192.347}{X}%
\TextAt{9}{113.832}{192.348}{\Character{41}}%
\TextAt{8}{122.281}{192.347}{defines}%
\TextAt{8}{160.705}{192.347}{a}%
\TextAt{13}{170.461}{192.347}{G}%
\TextAt{8}{179.694}{192.347}{-submodulation}%
\TextAt{13}{263.333}{192.347}{N}%
\TextAt{13}{273.949}{192.347}{A}%
\TextAt{8}{286.628}{192.347}{of}%
\TextAt{13}{299.958}{192.347}{A}%
\TextAt{8}{308.733}{192.347}{,}%
\TextAt{8}{315.883}{192.347}{the}%
\TextAt{8}{336.039}{192.347}{modulation}%
\TextAt{8}{398.158}{192.347}{of}%
\TextAt{8}{411.488}{192.347}{universal}%
\TextAt{8}{460.982}{192.347}{norms.}%
\TextAt{12}{79.226}{160.358}{Exercise}%
\TextAt{12}{130.946}{160.358}{IV.6.14.}%
\TextAt{8}{180.146}{160.358}{If}%
\TextAt{13}{190.551}{160.358}{A}%
\TextAt{8}{201.937}{160.358}{is}%
\TextAt{8}{212.414}{160.358}{a}%
\TextAt{13}{220.882}{160.358}{G}%
\TextAt{8}{230.116}{160.358}{-modulation,}%
\TextAt{8}{298.355}{160.358}{then}%
\TextAt{8}{323.728}{160.358}{the}%
\TextAt{8}{342.597}{160.358}{completion}%
\TextAt{20}{403.838}{163.38}{\Character{98}}%
\TextAt{13}{401.146}{160.358}{A}%
\TextAt{8}{412.531}{160.358}{is}%
\TextAt{8}{423.008}{160.358}{again}%
\TextAt{8}{452.948}{160.358}{a}%
\TextAt{13}{461.417}{160.358}{G}%
\TextAt{8}{470.65}{160.358}{-modulation}%
\TextAt{8}{79.226}{143.312}{which,}%
\TextAt{8}{115.964}{143.312}{for}%
\TextAt{8}{133.849}{143.312}{connected}%
\TextAt{13}{188.472}{143.312}{X}%
\TextAt{8}{199.128}{143.312}{,}%
\TextAt{8}{206.277}{143.312}{is}%
\TextAt{8}{218.041}{143.312}{given}%
\TextAt{8}{248.61}{143.312}{by}%
\TextAt{20}{239.716}{117.332}{\Character{98}}%
\TextAt{13}{237.024}{114.31}{A}%
\TextAt{9}{245.799}{114.31}{\Character{40}}%
\TextAt{13}{250.351}{114.31}{X}%
\TextAt{9}{261.007}{114.31}{\Character{41}}%
\TextAt{9}{268.862}{114.31}{\Character{61}}%
\NativeMathAt{281.305}{114.31}{11.955}{\varprojlim}%
\TextAt{21}{286.961}{101.842}{f}%
\TextAt{13}{299.557}{114.31}{A}%
\TextAt{9}{308.332}{114.31}{\Character{40}}%
\TextAt{13}{312.884}{114.31}{X}%
\TextAt{9}{323.539}{114.31}{\Character{41}\Character{47}}%
\TextAt{13}{333.945}{114.31}{f}%
\TextAt{3}{339.715}{112.517}{\Character{3}}%
\TextAt{13}{344.447}{114.31}{A}%
\TextAt{9}{353.222}{114.31}{\Character{40}}%
\TextAt{13}{357.775}{114.31}{Y}%
\TextAt{9}{367.168}{114.31}{\Character{41}}%
\TextAt{13}{371.72}{114.31}{\Character{59}}%
\TextAt{8}{79.226}{77.729}{where}%
\TextAt{8}{113.035}{77.729}{the}%
\TextAt{8}{133.192}{77.729}{projective}%
\TextAt{8}{187.169}{77.729}{limit}%
\TextAt{8}{215.133}{77.729}{is}%
\TextAt{8}{226.896}{77.729}{taken}%
\TextAt{8}{258.769}{77.729}{over}%
\TextAt{8}{283.815}{77.729}{all}%
\TextAt{8}{300.074}{77.729}{Galois}%
\TextAt{8}{335.987}{77.729}{maps}%
\TextAt{13}{366.616}{77.729}{f}%
\TextAt{9}{376.984}{77.729}{\Character{58}}%
\TextAt{13}{383.556}{77.729}{Y}%
\NativeMathAt{396.27}{77.729}{11.955}{\rightarrow}%
\TextAt{13}{411.546}{77.729}{X}%
\TextAt{8}{422.201}{77.729}{.}%
\TextAt{8}{297.211}{43.033}{283}%
\Destination{78.23}{809.815}{nav.page.283}%
\Destination{79.226}{414.655}{nav.ex.4.6.12}%
\Destination{79.226}{290.655}{nav.ex.4.6.13}%
\Destination{79.226}{174.355}{nav.ex.4.6.14}%
}

% Page 296 (body page 284)
\Page{%
\TextAt{2}{79.226}{796.943}{IV.}%
\TextAt{2}{98.012}{796.943}{Abstract}%
\TextAt{2}{143.252}{796.943}{Class}%
\TextAt{2}{171.844}{796.943}{Field}%
\TextAt{2}{199.575}{796.943}{Theory}%
\Vector{q 0 G 0 g 0.299 w 7.226 18.729 m 460.769 18.729 l S Q}%
\TextAt{8}{97.159}{761.124}{For}%
\TextAt{8}{118.036}{761.124}{the}%
\TextAt{8}{138.101}{761.124}{following,}%
\TextAt{8}{190.713}{761.124}{let}%
\TextAt{13}{207.526}{761.124}{d}%
\TextAt{9}{216.93}{761.124}{\Character{58}}%
\TextAt{13}{223.502}{761.124}{G}%
\NativeMathAt{236.056}{761.124}{11.955}{\rightarrow}%
\TextAt{20}{251.997}{764.212}{\Character{98}}%
\TextAt{18}{251.333}{761.124}{Z}%
\TextAt{8}{263.108}{761.124}{be}%
\TextAt{8}{278.94}{761.124}{a}%
\TextAt{8}{288.604}{761.124}{fixed}%
\TextAt{8}{316.798}{761.124}{surjective}%
\TextAt{8}{369.117}{761.124}{homomorphism.}%
\TextAt{8}{456.307}{761.124}{Let}%
\TextAt{13}{477.185}{761.124}{f}%
\TextAt{9}{487.552}{761.124}{\Character{58}}%
\TextAt{13}{494.124}{761.124}{X}%
\NativeMathAt{508.1}{761.124}{11.955}{\rightarrow}%
\TextAt{13}{523.376}{761.124}{Y}%
\TextAt{8}{79.226}{744.078}{be}%
\TextAt{8}{95.993}{744.078}{a}%
\TextAt{8}{106.591}{744.078}{map}%
\TextAt{8}{133.448}{744.078}{of}%
\TextAt{8}{147.62}{744.078}{connected}%
\TextAt{8}{203.086}{744.078}{finite}%
\TextAt{13}{233.84}{744.078}{G}%
\TextAt{8}{243.074}{744.078}{-sets}%
\TextAt{8}{270.696}{744.078}{and}%
\TextAt{13}{294.301}{744.078}{x}%
\NativeMathAt{305.71}{744.078}{11.955}{\in}%
\TextAt{13}{318.437}{744.078}{X}%
\TextAt{8}{329.092}{744.078}{,}%
\TextAt{13}{337.294}{744.078}{y}%
\TextAt{9}{348.187}{744.078}{\Character{61}}%
\TextAt{13}{362.048}{744.078}{f}%
\TextAt{9}{369.095}{744.078}{\Character{40}}%
\TextAt{13}{373.647}{744.078}{x}%
\TextAt{9}{380.299}{744.078}{\Character{41}}%
\NativeMathAt{389.608}{744.078}{11.955}{\in}%
\TextAt{13}{402.335}{744.078}{Y}%
\TextAt{8}{411.728}{744.078}{.}%
\TextAt{8}{422.704}{744.078}{The}%
\TextAt{12}{447.6}{744.078}{inertia}%
\TextAt{12}{491.417}{744.078}{degree}%
\TextAt{8}{529.518}{744.078}{,}%
\TextAt{8}{79.226}{727.032}{resp.}%
\TextAt{8}{108.548}{727.032}{the}%
\TextAt{12}{128.705}{727.032}{ramification}%
\TextAt{12}{203.353}{727.032}{index}%
\TextAt{8}{235.273}{727.032}{,}%
\TextAt{8}{242.422}{727.032}{of}%
\TextAt{13}{255.752}{727.032}{f}%
\TextAt{8}{266.696}{727.032}{is}%
\TextAt{8}{278.46}{727.032}{defined}%
\TextAt{8}{318.773}{727.032}{by}%
\TextAt{13}{168.389}{698.031}{f}%
\TextAt{21}{174.159}{696.171}{X}%
\TextAt{3}{181.746}{696.171}{\Character{106}}%
\TextAt{21}{184.098}{696.171}{Y}%
\TextAt{9}{194.707}{698.031}{\Character{61}}%
\TextAt{9}{207.117}{698.031}{\Character{40}}%
\TextAt{13}{211.685}{698.031}{d}%
\TextAt{9}{217.767}{698.031}{\Character{40}}%
\TextAt{13}{222.32}{698.031}{G}%
\TextAt{21}{231.553}{696.237}{y}%
\TextAt{9}{236.523}{698.031}{\Character{41}}%
\TextAt{9}{244.378}{698.031}{\Character{58}}%
\TextAt{13}{250.969}{698.031}{d}%
\TextAt{9}{257.052}{698.031}{\Character{40}}%
\TextAt{13}{261.604}{698.031}{G}%
\TextAt{21}{270.837}{696.237}{x}%
\TextAt{9}{276.102}{698.031}{\Character{41}\Character{41}}%
\TextAt{13}{285.207}{698.031}{\Character{59}}%
\TextAt{8}{314.362}{698.031}{resp.}%
\TextAt{13}{362.398}{698.031}{e}%
\TextAt{21}{367.823}{696.171}{X}%
\TextAt{3}{375.41}{696.171}{\Character{106}}%
\TextAt{21}{377.762}{696.171}{Y}%
\TextAt{9}{388.371}{698.031}{\Character{61}}%
\TextAt{9}{400.781}{698.031}{\Character{40}}%
\TextAt{13}{405.349}{698.031}{I}%
\TextAt{21}{410.512}{696.237}{y}%
\TextAt{9}{418.802}{698.031}{\Character{58}}%
\TextAt{13}{425.375}{698.031}{I}%
\TextAt{21}{430.538}{696.237}{x}%
\TextAt{9}{435.803}{698.031}{\Character{41}}%
\TextAt{13}{440.355}{698.031}{\Character{59}}%
\TextAt{8}{79.226}{669.029}{where}%
\TextAt{13}{112.881}{669.029}{I}%
\TextAt{21}{118.044}{667.236}{y}%
\TextAt{8}{123.014}{669.029}{,}%
\TextAt{8}{130.04}{669.029}{resp.}%
\TextAt{13}{159.31}{669.029}{I}%
\TextAt{21}{164.473}{667.236}{x}%
\TextAt{8}{169.738}{669.029}{,}%
\TextAt{8}{176.764}{669.029}{is}%
\TextAt{8}{188.374}{669.029}{the}%
\TextAt{8}{208.376}{669.029}{kernel}%
\TextAt{8}{242.688}{669.029}{of}%
\TextAt{13}{255.864}{669.029}{d}%
\TextAt{9}{265.267}{669.029}{\Character{58}}%
\TextAt{13}{271.84}{669.029}{G}%
\TextAt{21}{281.074}{667.236}{y}%
\NativeMathAt{289.364}{669.029}{11.955}{\rightarrow}%
\TextAt{20}{305.304}{672.118}{\Character{98}}%
\TextAt{18}{304.64}{669.029}{Z}%
\TextAt{8}{312.61}{669.029}{,}%
\TextAt{8}{319.636}{669.029}{resp.}%
\TextAt{13}{348.907}{669.029}{d}%
\TextAt{9}{358.31}{669.029}{\Character{58}}%
\TextAt{13}{364.883}{669.029}{G}%
\TextAt{21}{374.116}{667.236}{x}%
\NativeMathAt{382.702}{669.029}{11.955}{\rightarrow}%
\TextAt{20}{398.642}{672.118}{\Character{98}}%
\TextAt{18}{397.978}{669.029}{Z}%
\TextAt{8}{405.948}{669.029}{.}%
\TextAt{13}{414.345}{669.029}{f}%
\TextAt{8}{425.135}{669.029}{is}%
\TextAt{8}{436.744}{669.029}{called}%
\TextAt{12}{469.754}{669.029}{unramified}%
\TextAt{8}{79.226}{651.983}{if}%
\TextAt{13}{89.95}{651.983}{e}%
\TextAt{21}{95.375}{650.124}{X}%
\TextAt{3}{102.962}{650.124}{\Character{106}}%
\TextAt{21}{105.315}{650.124}{Y}%
\TextAt{9}{115.924}{651.983}{\Character{61}}%
\TextAt{9}{128.334}{651.983}{1}%
\TextAt{8}{134.202}{651.983}{.}%
\TextAt{12}{79.226}{625.971}{Exercise}%
\TextAt{12}{132.048}{625.971}{IV.6.15.}%
\TextAt{13}{182.205}{625.971}{d}%
\TextAt{8}{191.856}{625.971}{defines}%
\TextAt{8}{229.951}{625.971}{a}%
\TextAt{13}{239.377}{625.971}{G}%
\TextAt{8}{248.61}{625.971}{-modulation}%
\TextAt{20}{314.962}{629.06}{\Character{98}}%
\TextAt{18}{314.297}{625.971}{Z}%
\TextAt{8}{325.836}{625.971}{such}%
\TextAt{8}{351.903}{625.971}{that}%
\TextAt{8}{376.943}{625.971}{the}%
\TextAt{8}{396.77}{625.971}{maps}%
\TextAt{13}{427.07}{625.971}{f}%
\TextAt{3}{434.116}{630.31}{\Character{3}}%
\TextAt{13}{438.849}{625.971}{\Character{59}}%
\TextAt{13}{444.073}{625.971}{f}%
\TextAt{3}{449.863}{624.178}{\Character{3}}%
\TextAt{8}{454.595}{625.971}{,}%
\TextAt{8}{461.481}{625.971}{corresponding}%
\TextAt{8}{79.226}{608.925}{to}%
\TextAt{8}{93.536}{608.925}{a}%
\TextAt{8}{103.292}{608.925}{mapping}%
\TextAt{13}{151.423}{608.925}{f}%
\TextAt{9}{161.791}{608.925}{\Character{58}}%
\TextAt{13}{168.363}{608.925}{X}%
\NativeMathAt{182.339}{608.925}{11.955}{\rightarrow}%
\TextAt{13}{197.615}{608.925}{Y}%
\TextAt{8}{210.906}{608.925}{of}%
\TextAt{8}{224.236}{608.925}{connected}%
\TextAt{13}{278.859}{608.925}{G}%
\TextAt{8}{288.093}{608.925}{-sets,}%
\TextAt{8}{318.124}{608.925}{are}%
\TextAt{8}{337.635}{608.925}{given}%
\TextAt{8}{368.204}{608.925}{by}%
\TextAt{20}{241.312}{579.354}{\Character{98}}%
\TextAt{18}{240.648}{576.265}{Z}%
\TextAt{9}{248.618}{576.265}{\Character{40}}%
\TextAt{13}{253.171}{576.265}{Y}%
\TextAt{9}{262.564}{576.265}{\Character{41}}%
\TextAt{9}{270.419}{576.265}{\Character{61}}%
\TextAt{20}{283.527}{579.354}{\Character{98}}%
\TextAt{18}{282.862}{576.265}{Z}%
\TextAt{21}{294.223}{588.896}{e}%
\TextAt{22}{298.152}{587.158}{X}%
\TextAt{23}{304.6}{587.158}{\Character{106}}%
\TextAt{22}{306.869}{587.158}{Y}%
\NativeMathAt{297.763}{576.265}{11.955}{\rightleftarrows}%
\TextAt{21}{294.153}{566.577}{f}%
\TextAt{22}{298.222}{564.839}{X}%
\TextAt{23}{304.67}{564.839}{\Character{106}}%
\TextAt{22}{306.939}{564.839}{Y}%
\TextAt{20}{317.314}{579.354}{\Character{98}}%
\TextAt{18}{316.649}{576.265}{Z}%
\TextAt{9}{327.94}{576.265}{\Character{61}}%
\TextAt{20}{341.03}{579.354}{\Character{98}}%
\TextAt{18}{340.366}{576.265}{Z}%
\TextAt{9}{348.336}{576.265}{\Character{40}}%
\TextAt{13}{352.888}{576.265}{X}%
\TextAt{9}{363.543}{576.265}{\Character{41}}%
\TextAt{13}{368.096}{576.265}{\Character{58}}%
\TextAt{8}{79.226}{540.781}{This}%
\TextAt{8}{105.946}{540.781}{gives}%
\TextAt{8}{134.626}{540.781}{a}%
\TextAt{8}{144.382}{540.781}{homomorphism}%
\TextAt{8}{227.052}{540.781}{of}%
\TextAt{13}{240.382}{540.781}{G}%
\TextAt{8}{249.615}{540.781}{-modulations}%
\TextAt{13}{271.197}{511.78}{d}%
\TextAt{9}{280.6}{511.78}{\Character{58}}%
\TextAt{13}{287.173}{511.78}{\Character{25}}%
\TextAt{7}{294.242}{516.716}{ab}%
\NativeArrowAt{306.995}{511.78}{11.955}{19.26}{\longrightarrow}%
\TextAt{20}{330.241}{514.868}{\Character{98}}%
\TextAt{18}{329.577}{511.78}{Z}%
\TextAt{13}{337.547}{511.78}{\Character{58}}%
\TextAt{12}{79.226}{467.835}{Exercise}%
\TextAt{12}{132.794}{467.835}{IV.6.16.}%
\TextAt{8}{183.6}{467.835}{An}%
\TextAt{8}{203.096}{467.835}{unramified}%
\TextAt{8}{261.948}{467.835}{map}%
\TextAt{13}{288.282}{467.835}{f}%
\TextAt{9}{299.194}{467.835}{\Character{58}}%
\TextAt{13}{306.311}{467.835}{X}%
\NativeMathAt{320.831}{467.835}{11.955}{\rightarrow}%
\TextAt{13}{336.652}{467.835}{Y}%
\TextAt{8}{350.262}{467.835}{of}%
\TextAt{8}{363.912}{467.835}{connected}%
\TextAt{8}{418.854}{467.835}{finite}%
\TextAt{13}{449.086}{467.835}{G}%
\TextAt{8}{458.319}{467.835}{-sets}%
\TextAt{8}{485.418}{467.835}{is}%
\TextAt{8}{497.502}{467.835}{Galois,}%
\TextAt{8}{79.226}{450.789}{and}%
\TextAt{13}{101.989}{450.789}{d}%
\TextAt{8}{111.969}{450.789}{induces}%
\TextAt{8}{153.645}{450.789}{an}%
\TextAt{8}{169.904}{450.789}{isomorphism}%
\TextAt{13}{249.832}{421.788}{G}%
\TextAt{9}{259.066}{421.788}{\Character{40}}%
\TextAt{13}{263.618}{421.788}{X}%
\TextAt{16}{277.594}{421.788}{\Character{106}}%
\TextAt{13}{284.236}{421.788}{Y}%
\TextAt{9}{293.629}{421.788}{\Character{41}}%
\NativeMathAt{301.502}{421.787}{11.955}{\cong}%
\TextAt{18}{314.121}{421.788}{Z}%
\TextAt{9}{322.091}{421.788}{\Character{47}}%
\TextAt{13}{327.944}{421.788}{f}%
\TextAt{21}{333.714}{419.928}{X}%
\TextAt{3}{341.301}{419.928}{\Character{106}}%
\TextAt{21}{343.654}{419.928}{Y}%
\TextAt{18}{350.942}{421.788}{Z}%
\TextAt{13}{358.912}{421.788}{\Character{58}}%
\TextAt{8}{79.226}{392.786}{Let}%
\TextAt{13}{100.196}{392.786}{\Character{39}}%
\TextAt{21}{107.864}{390.927}{X}%
\TextAt{3}{115.451}{390.927}{\Character{106}}%
\TextAt{21}{117.803}{390.927}{Y}%
\NativeMathAt{128.412}{392.786}{11.955}{\in}%
\TextAt{13}{139.703}{392.786}{G}%
\TextAt{9}{148.937}{392.786}{\Character{40}}%
\TextAt{13}{153.489}{392.786}{X}%
\TextAt{16}{167.465}{392.786}{\Character{106}}%
\TextAt{13}{174.107}{392.786}{Y}%
\TextAt{9}{183.5}{392.786}{\Character{41}}%
\TextAt{8}{191.949}{392.786}{be}%
\TextAt{8}{207.874}{392.786}{the}%
\TextAt{8}{228.03}{392.786}{element}%
\TextAt{8}{271.272}{392.786}{which}%
\TextAt{8}{304.758}{392.786}{is}%
\TextAt{8}{316.522}{392.786}{mapped}%
\TextAt{8}{361.067}{392.786}{to}%
\TextAt{9}{375.378}{392.786}{1}%
\TextAt{8}{384.551}{392.786}{mod}%
\TextAt{13}{409.989}{392.786}{f}%
\TextAt{21}{415.759}{390.927}{X}%
\TextAt{3}{423.346}{390.927}{\Character{106}}%
\TextAt{21}{425.699}{390.927}{Y}%
\TextAt{18}{432.987}{392.786}{Z}%
\TextAt{8}{440.957}{392.786}{.}%
\TextAt{8}{97.159}{366.774}{Let}%
\TextAt{13}{117.143}{366.774}{A}%
\TextAt{8}{128.831}{366.774}{be}%
\TextAt{8}{143.77}{366.774}{a}%
\TextAt{13}{152.541}{366.774}{G}%
\TextAt{8}{161.774}{366.774}{-modulation.}%
\TextAt{8}{232.013}{366.774}{We}%
\TextAt{8}{251.173}{366.774}{define}%
\TextAt{8}{283.997}{366.774}{a}%
\TextAt{19}{292.767}{366.774}{henselian}%
\TextAt{19}{343.117}{366.774}{valuation}%
\TextAt{8}{392.929}{366.774}{of}%
\TextAt{13}{405.274}{366.774}{A}%
\TextAt{8}{416.962}{366.774}{to}%
\TextAt{8}{430.287}{366.774}{be}%
\TextAt{8}{445.227}{366.774}{a}%
\TextAt{8}{453.997}{366.774}{homomorphism}%
\TextAt{13}{276.683}{337.773}{v}%
\TextAt{9}{286.092}{337.773}{\Character{58}}%
\TextAt{13}{292.664}{337.773}{A}%
\NativeArrowAt{304.761}{337.773}{11.955}{19.26}{\longrightarrow}%
\TextAt{20}{328.007}{340.861}{\Character{98}}%
\TextAt{18}{327.343}{337.773}{Z}%
\TextAt{8}{79.226}{308.772}{such}%
\TextAt{8}{105.74}{308.772}{that}%
\TextAt{8}{131.227}{308.772}{the}%
\TextAt{8}{151.5}{308.772}{submodulation}%
\TextAt{13}{231.359}{308.772}{v}%
\TextAt{9}{237.447}{308.772}{\Character{40}}%
\TextAt{13}{241.999}{308.772}{A}%
\TextAt{9}{250.775}{308.772}{\Character{41}}%
\TextAt{8}{259.342}{308.772}{of}%
\TextAt{20}{273.453}{311.86}{\Character{98}}%
\TextAt{18}{272.789}{308.772}{Z}%
\TextAt{8}{284.774}{308.772}{comes}%
\TextAt{8}{319.418}{308.772}{from}%
\TextAt{8}{347.175}{308.772}{a}%
\TextAt{8}{357.048}{308.772}{subgroup}%
\TextAt{13}{407.963}{308.772}{Z}%
\NativeMathAt{420.354}{308.772}{11.955}{\subseteq}%
\TextAt{20}{433.838}{311.86}{\Character{98}}%
\TextAt{18}{433.174}{308.772}{Z}%
\TextAt{8}{445.159}{308.772}{which}%
\TextAt{8}{478.762}{308.772}{contains}%
\TextAt{18}{524.799}{308.772}{Z}%
\TextAt{8}{79.226}{291.726}{and}%
\TextAt{8}{101.989}{291.726}{satisfies}%
\TextAt{13}{145.099}{291.726}{Z}%
\TextAt{9}{153.97}{291.726}{\Character{47}}%
\TextAt{13}{159.823}{291.726}{nZ}%
\TextAt{9}{179.002}{291.726}{\Character{61}}%
\TextAt{18}{191.428}{291.726}{Z}%
\TextAt{9}{199.398}{291.726}{\Character{47}}%
\TextAt{13}{205.251}{291.726}{n}%
\TextAt{18}{212.238}{291.726}{Z}%
\TextAt{8}{224.106}{291.726}{for}%
\TextAt{8}{241.991}{291.726}{all}%
\TextAt{13}{258.25}{291.726}{n}%
\NativeMathAt{268.558}{291.726}{11.955}{\in}%
\TextAt{18}{279.849}{291.726}{N}%
\TextAt{8}{288.484}{291.726}{.}%
\TextAt{8}{296.932}{291.726}{Let}%
\TextAt{13}{317.901}{291.726}{U}%
\TextAt{8}{330.998}{291.726}{denote}%
\TextAt{8}{368.716}{291.726}{the}%
\TextAt{8}{388.873}{291.726}{kernel}%
\TextAt{8}{423.339}{291.726}{of}%
\TextAt{13}{436.67}{291.726}{v}%
\TextAt{8}{445.445}{291.726}{.}%
\TextAt{12}{79.226}{265.713}{Exercise}%
\TextAt{12}{132.35}{265.713}{IV.6.17.}%
\TextAt{8}{182.769}{265.713}{Compare}%
\TextAt{8}{232.783}{265.713}{this}%
\TextAt{8}{255.538}{265.713}{definition}%
\TextAt{8}{307.5}{265.713}{with}%
\TextAt{8}{334.093}{265.713}{the}%
\TextAt{8}{354.183}{265.713}{definition}%
\TextAt{8}{406.145}{265.713}{(4.6)}%
\TextAt{8}{434.053}{265.713}{of}%
\TextAt{8}{447.316}{265.713}{a}%
\TextAt{8}{457.005}{265.713}{henselian}%
\TextAt{8}{507.723}{265.713}{valu-}%
\TextAt{8}{79.226}{248.667}{ation}%
\TextAt{8}{109.15}{248.667}{of}%
\TextAt{8}{122.48}{248.667}{a}%
\TextAt{13}{132.235}{248.667}{G}%
\TextAt{8}{141.469}{248.667}{-module}%
\TextAt{13}{186.66}{248.667}{A}%
\TextAt{8}{195.435}{248.667}{.}%
\TextAt{12}{79.226}{216.678}{Exercise}%
\TextAt{12}{131.844}{216.678}{IV.6.18.}%
\TextAt{8}{181.824}{216.678}{Assume}%
\TextAt{8}{224.679}{216.678}{that}%
\TextAt{8}{249.541}{216.678}{for}%
\TextAt{8}{266.92}{216.678}{every}%
\TextAt{8}{297.305}{216.678}{unramified}%
\TextAt{8}{355.331}{216.678}{map}%
\TextAt{13}{380.839}{216.678}{f}%
\TextAt{9}{391.206}{216.678}{\Character{58}}%
\TextAt{13}{397.779}{216.678}{X}%
\NativeMathAt{411.755}{216.678}{11.955}{\rightarrow}%
\TextAt{13}{427.031}{216.678}{Y}%
\TextAt{8}{439.815}{216.678}{of}%
\TextAt{8}{452.638}{216.678}{connected}%
\TextAt{8}{506.755}{216.678}{finite}%
\TextAt{13}{79.226}{199.632}{G}%
\TextAt{8}{88.46}{199.632}{-sets,}%
\TextAt{8}{118.491}{199.632}{the}%
\TextAt{8}{138.648}{199.632}{sequence}%
\TextAt{9}{173.196}{167.258}{0}%
\NativeArrowAt{182.37}{167.258}{11.955}{19.26}{\longrightarrow}%
\TextAt{13}{204.952}{167.258}{U}%
\TextAt{9}{214.151}{167.258}{\Character{40}}%
\TextAt{13}{218.704}{167.258}{Y}%
\TextAt{9}{228.097}{167.258}{\Character{41}}%
\TextAt{21}{238.558}{174.522}{f}%
\TextAt{23}{243.504}{177.334}{\Character{3}}%
\NativeArrowAt{235.97}{167.258}{11.955}{15.853}{\longrightarrow}%
\TextAt{13}{255.141}{167.258}{U}%
\TextAt{9}{264.34}{167.258}{\Character{40}}%
\TextAt{13}{268.893}{167.258}{X}%
\TextAt{9}{279.548}{167.258}{\Character{41}}%
\TextAt{21}{290.008}{177.571}{\Character{39}}%
\TextAt{23}{295.527}{180.384}{\Character{3}}%
\TextAt{22}{295.527}{174.466}{X}%
\TextAt{23}{301.975}{174.466}{\Character{106}}%
\TextAt{22}{304.244}{174.466}{Y}%
\TextAt{3}{310.634}{177.571}{\Character{0}}%
\TextAt{17}{317.22}{177.571}{1}%
\NativeArrowAt{287.421}{167.258}{11.955}{38.006}{\longrightarrow}%
\TextAt{13}{328.774}{167.258}{U}%
\TextAt{9}{337.973}{167.258}{\Character{40}}%
\TextAt{13}{342.526}{167.258}{X}%
\TextAt{9}{353.181}{167.258}{\Character{41}}%
\TextAt{21}{363.642}{174.522}{f}%
\TextAt{23}{367.71}{173.525}{\Character{3}}%
\NativeArrowAt{361.054}{167.258}{11.955}{14.968}{\longrightarrow}%
\TextAt{13}{379.347}{167.257}{U}%
\TextAt{9}{388.546}{167.258}{\Character{40}}%
\TextAt{13}{393.098}{167.257}{Y}%
\TextAt{9}{402.491}{167.258}{\Character{41}}%
\NativeArrowAt{410.364}{167.257}{11.955}{19.26}{\longrightarrow}%
\TextAt{9}{432.946}{167.258}{0}%
\TextAt{8}{79.226}{138.256}{is}%
\TextAt{8}{90.538}{138.256}{exact,}%
\TextAt{8}{124.321}{138.256}{and}%
\TextAt{8}{146.632}{138.256}{that}%
\TextAt{13}{171.549}{138.256}{A}%
\TextAt{9}{180.324}{138.256}{\Character{40}}%
\TextAt{13}{184.877}{138.256}{Y}%
\TextAt{9}{194.27}{138.256}{\Character{41}}%
\TextAt{17}{198.822}{142.595}{\Character{91}}%
\TextAt{21}{201.174}{142.595}{X}%
\TextAt{17}{208.761}{142.595}{\Character{58}}%
\TextAt{21}{211.114}{142.595}{Y}%
\TextAt{17}{217.904}{142.595}{\Character{93}}%
\NativeMathAt{224.075}{138.256}{11.955}{\subseteq}%
\TextAt{13}{236.694}{138.256}{f}%
\TextAt{3}{242.464}{136.463}{\Character{3}}%
\TextAt{13}{247.196}{138.256}{A}%
\TextAt{9}{255.972}{138.256}{\Character{40}}%
\TextAt{13}{260.524}{138.256}{X}%
\TextAt{9}{271.179}{138.256}{\Character{41}}%
\TextAt{8}{279.177}{138.256}{for}%
\TextAt{8}{296.61}{138.256}{every}%
\TextAt{8}{327.051}{138.256}{Galois}%
\TextAt{8}{362.512}{138.256}{mapping}%
\TextAt{13}{410.192}{138.256}{f}%
\TextAt{9}{420.559}{138.256}{\Character{58}}%
\TextAt{13}{427.132}{138.256}{X}%
\NativeMathAt{441.108}{138.256}{11.955}{\rightarrow}%
\TextAt{13}{456.384}{138.256}{Y}%
\TextAt{8}{469.223}{138.256}{(the}%
\TextAt{8}{493.482}{138.256}{latter}%
\TextAt{8}{524.903}{138.256}{is}%
\TextAt{8}{79.226}{121.21}{a}%
\TextAt{8}{88.455}{121.21}{consequence}%
\TextAt{8}{153.993}{121.21}{of}%
\TextAt{8}{166.797}{121.21}{the}%
\TextAt{8}{186.427}{121.21}{condition}%
\TextAt{8}{237.284}{121.21}{which}%
\TextAt{8}{270.245}{121.21}{will}%
\TextAt{8}{291.823}{121.21}{be}%
\TextAt{8}{307.221}{121.21}{imposed}%
\TextAt{8}{352.603}{121.21}{in}%
\TextAt{8}{365.73}{121.21}{exercise}%
\TextAt{8}{408.505}{121.21}{19).}%
\TextAt{8}{433.049}{121.21}{Then}%
\TextAt{8}{463.08}{121.21}{the}%
\TextAt{8}{482.71}{121.21}{pair}%
\TextAt{9}{506.25}{121.21}{\Character{40}}%
\TextAt{13}{510.802}{121.21}{d\Character{59}}%
\TextAt{13}{522.1}{121.21}{v}%
\TextAt{9}{528.217}{121.21}{\Character{41}}%
\TextAt{8}{79.226}{104.164}{gives,}%
\TextAt{8}{111.158}{104.164}{for}%
\TextAt{8}{129.043}{104.164}{every}%
\TextAt{8}{159.935}{104.164}{Galois}%
\TextAt{8}{195.849}{104.164}{mapping}%
\TextAt{13}{243.98}{104.164}{f}%
\TextAt{9}{254.348}{104.164}{\Character{58}}%
\TextAt{13}{260.92}{104.164}{X}%
\NativeMathAt{274.896}{104.164}{11.955}{\rightarrow}%
\TextAt{13}{290.172}{104.164}{Y}%
\TextAt{8}{299.565}{104.164}{,}%
\TextAt{8}{306.714}{104.164}{a}%
\TextAt{8}{316.47}{104.164}{canonical}%
\TextAt{8}{367.853}{104.164}{“reciprocity}%
\TextAt{8}{431.586}{104.164}{homomorphism”}%
\TextAt{13}{214.966}{75.163}{r}%
\TextAt{21}{220.242}{73.304}{X}%
\TextAt{3}{227.829}{73.304}{\Character{106}}%
\TextAt{21}{230.181}{73.304}{Y}%
\TextAt{9}{240.79}{75.163}{\Character{58}}%
\TextAt{13}{247.362}{75.163}{G}%
\TextAt{9}{256.596}{75.163}{\Character{40}}%
\TextAt{13}{261.148}{75.163}{X}%
\TextAt{16}{275.124}{75.163}{\Character{106}}%
\TextAt{13}{281.766}{75.163}{Y}%
\TextAt{9}{291.159}{75.163}{\Character{41}}%
\NativeArrowAt{299.032}{75.163}{11.955}{19.26}{\longrightarrow}%
\TextAt{13}{321.614}{75.163}{A}%
\TextAt{9}{330.389}{75.163}{\Character{40}}%
\TextAt{13}{334.942}{75.163}{Y}%
\TextAt{9}{344.335}{75.163}{\Character{41}\Character{47}}%
\TextAt{13}{354.74}{75.163}{f}%
\TextAt{3}{360.51}{73.37}{\Character{3}}%
\TextAt{13}{365.243}{75.163}{A}%
\TextAt{9}{374.018}{75.163}{\Character{40}}%
\TextAt{13}{378.57}{75.163}{X}%
\TextAt{9}{389.225}{75.163}{\Character{41}}%
\TextAt{13}{393.778}{75.163}{\Character{58}}%
\TextAt{8}{297.211}{43.033}{284}%
\Destination{78.23}{809.815}{nav.page.284}%
\Destination{79.226}{639.955}{nav.ex.4.6.15}%
\Destination{79.226}{481.855}{nav.ex.4.6.16}%
\Destination{79.226}{279.755}{nav.ex.4.6.17}%
\Destination{79.226}{230.655}{nav.ex.4.6.18}%
\Destination{79.36}{776.155}{nav.print.309}%
\Hotlink{408.005}{119.055}{12.716}{11.047}{nav.ex.4.6.19}%
\Hotlink{405.645}{263.555}{25.078}{11.047}{nav.number.4.4.6}%
\Destination{79.226}{308.726}{errata.111}%
}

% Page 297 (body page 285)
\Page{%
\TextAt{2}{79.226}{796.943}{IV.}%
\TextAt{2}{98.012}{796.943}{Abstract}%
\TextAt{2}{143.252}{796.943}{Class}%
\TextAt{2}{171.844}{796.943}{Field}%
\TextAt{2}{199.575}{796.943}{Theory}%
\Vector{q 0 G 0 g 0.299 w 7.226 18.729 m 460.769 18.729 l S Q}%
\TextAt{12}{79.226}{761.124}{Exercise}%
\TextAt{12}{132.342}{761.124}{IV.6.19.}%
\TextAt{8}{182.756}{761.124}{Assume,}%
\TextAt{8}{229.31}{761.124}{beyond}%
\TextAt{8}{269.885}{761.124}{the}%
\TextAt{8}{289.968}{761.124}{condition}%
\TextAt{8}{341.278}{761.124}{required}%
\TextAt{8}{387.053}{761.124}{in}%
\TextAt{8}{400.632}{761.124}{exercise}%
\TextAt{8}{443.861}{761.124}{18,}%
\TextAt{8}{462.667}{761.124}{that}%
\TextAt{8}{487.963}{761.124}{for}%
\TextAt{8}{505.775}{761.124}{every}%
\TextAt{8}{79.226}{744.078}{Galois}%
\TextAt{8}{115.139}{744.078}{mapping}%
\TextAt{13}{163.271}{744.078}{f}%
\TextAt{9}{173.638}{744.078}{\Character{58}}%
\TextAt{13}{180.211}{744.078}{X}%
\NativeMathAt{194.187}{744.078}{11.955}{\rightarrow}%
\TextAt{13}{209.463}{744.078}{Y}%
\TextAt{8}{222.753}{744.078}{with}%
\TextAt{8}{249.413}{744.078}{cyclic}%
\TextAt{8}{281.597}{744.078}{Galois}%
\TextAt{8}{317.51}{744.078}{group}%
\TextAt{13}{350.686}{744.078}{G}%
\TextAt{9}{359.919}{744.078}{\Character{40}}%
\TextAt{13}{364.471}{744.078}{X}%
\TextAt{16}{378.447}{744.078}{\Character{106}}%
\TextAt{13}{385.089}{744.078}{Y}%
\TextAt{9}{394.482}{744.078}{\Character{41}}%
\TextAt{8}{399.035}{744.078}{,}%
\TextAt{8}{406.184}{744.078}{one}%
\TextAt{8}{427.643}{744.078}{has}%
\TextAt{9}{151.999}{715.077}{\Character{40}}%
\TextAt{13}{156.551}{715.077}{A}%
\TextAt{9}{165.326}{715.077}{\Character{40}}%
\TextAt{13}{169.879}{715.077}{Y}%
\TextAt{9}{179.272}{715.077}{\Character{41}}%
\TextAt{9}{187.127}{715.077}{\Character{58}}%
\TextAt{13}{193.717}{715.077}{f}%
\TextAt{3}{199.487}{713.283}{\Character{3}}%
\TextAt{13}{204.22}{715.077}{A}%
\TextAt{9}{212.995}{715.077}{\Character{40}}%
\TextAt{13}{217.547}{715.077}{X}%
\TextAt{9}{228.202}{715.077}{\Character{41}\Character{41}}%
\TextAt{9}{240.6}{715.077}{\Character{61}}%
\TextAt{9}{253.021}{715.077}{\Character{91}}%
\TextAt{13}{256.305}{715.077}{X}%
\TextAt{9}{270.281}{715.077}{\Character{58}}%
\TextAt{13}{276.853}{715.077}{Y}%
\TextAt{9}{286.247}{715.077}{\Character{93}}%
\TextAt{8}{313.409}{715.077}{and}%
\TextAt{8}{358.177}{715.077}{ker}%
\TextAt{13}{376.105}{715.077}{f}%
\TextAt{3}{381.875}{713.283}{\Character{3}}%
\TextAt{9}{389.929}{715.077}{\Character{61}}%
\TextAt{8}{402.354}{715.077}{im}%
\TextAt{9}{415.361}{715.077}{\Character{40}}%
\TextAt{13}{419.914}{715.077}{\Character{27}}%
\TextAt{3}{426.996}{720.013}{\Character{3}}%
\TextAt{16}{434.385}{715.077}{\Character{0}}%
\TextAt{9}{446.34}{715.077}{1\Character{41}}%
\TextAt{13}{456.746}{715.077}{\Character{59}}%
\TextAt{8}{79.226}{686.075}{where}%
\TextAt{9}{113.039}{686.075}{\Character{91}}%
\TextAt{13}{116.29}{686.075}{X}%
\TextAt{9}{130.272}{686.075}{\Character{58}}%
\TextAt{13}{136.851}{686.075}{Y}%
\TextAt{9}{146.244}{686.075}{\Character{93}}%
\TextAt{9}{152.807}{686.075}{\Character{61}}%
\TextAt{8}{165.253}{686.075}{\char35\relax }%
\TextAt{13}{175.009}{686.075}{f}%
\TextAt{3}{182.055}{690.414}{\Character{0}}%
\TextAt{17}{188.642}{690.414}{1}%
\TextAt{9}{193.374}{686.075}{\Character{40}}%
\TextAt{13}{197.926}{686.075}{y}%
\TextAt{9}{204.063}{686.075}{\Character{41}}%
\TextAt{8}{208.615}{686.075}{,}%
\TextAt{8}{215.769}{686.075}{with}%
\TextAt{13}{242.432}{686.075}{y}%
\NativeMathAt{251.896}{686.075}{11.955}{\in}%
\TextAt{13}{263.192}{686.075}{Y}%
\TextAt{8}{272.586}{686.075}{,}%
\TextAt{8}{279.739}{686.075}{and}%
\TextAt{13}{302.505}{686.075}{\Character{27}}%
\TextAt{8}{313.488}{686.075}{is}%
\TextAt{8}{325.256}{686.075}{a}%
\TextAt{8}{335.015}{686.075}{generator}%
\TextAt{8}{387.059}{686.075}{of}%
\TextAt{13}{400.392}{686.075}{G}%
\TextAt{9}{409.626}{686.075}{\Character{40}}%
\TextAt{13}{414.178}{686.075}{X}%
\TextAt{16}{428.16}{686.075}{\Character{106}}%
\TextAt{13}{434.808}{686.075}{Y}%
\TextAt{9}{444.201}{686.075}{\Character{41}}%
\TextAt{8}{448.753}{686.075}{.}%
\TextAt{8}{457.212}{686.075}{Then}%
\TextAt{8}{487.772}{686.075}{if}%
\TextAt{13}{498.5}{686.075}{r}%
\TextAt{21}{503.775}{684.216}{X}%
\TextAt{3}{511.362}{684.216}{\Character{106}}%
\TextAt{21}{513.714}{684.216}{Y}%
\TextAt{8}{524.903}{686.075}{is}%
\TextAt{8}{79.226}{669.029}{an}%
\TextAt{8}{95.485}{669.029}{isomorphism}%
\TextAt{8}{163.905}{669.029}{for}%
\TextAt{8}{181.79}{669.029}{every}%
\TextAt{8}{212.682}{669.029}{Galois}%
\TextAt{8}{248.595}{669.029}{mapping}%
\TextAt{13}{296.727}{669.029}{f}%
\TextAt{9}{307.094}{669.029}{\Character{58}}%
\TextAt{13}{313.666}{669.029}{X}%
\NativeMathAt{327.642}{669.029}{11.955}{\rightarrow}%
\TextAt{13}{342.918}{669.029}{Y}%
\TextAt{8}{356.209}{669.029}{of}%
\TextAt{8}{369.539}{669.029}{prime}%
\TextAt{8}{402.703}{669.029}{degree}%
\TextAt{9}{439.118}{669.029}{\Character{91}}%
\TextAt{13}{442.37}{669.029}{X}%
\TextAt{9}{456.346}{669.029}{\Character{58}}%
\TextAt{13}{462.918}{669.029}{Y}%
\TextAt{9}{472.311}{669.029}{\Character{93}}%
\TextAt{8}{475.563}{669.029}{,}%
\TextAt{8}{482.712}{669.029}{so}%
\TextAt{8}{497.082}{669.029}{is}%
\TextAt{13}{210.25}{640.028}{r}%
\TextAt{21}{215.525}{638.169}{X}%
\TextAt{3}{223.112}{638.169}{\Character{106}}%
\TextAt{21}{225.465}{638.169}{Y}%
\TextAt{9}{236.073}{640.028}{\Character{58}}%
\TextAt{13}{242.646}{640.028}{G}%
\TextAt{9}{251.88}{640.028}{\Character{40}}%
\TextAt{13}{256.432}{640.028}{X}%
\TextAt{16}{270.408}{640.028}{\Character{106}}%
\TextAt{13}{277.05}{640.028}{Y}%
\TextAt{9}{286.443}{640.028}{\Character{41}}%
\TextAt{7}{290.995}{644.965}{ab}%
\NativeArrowAt{303.748}{640.028}{11.955}{19.26}{\longrightarrow}%
\TextAt{13}{326.33}{640.028}{A}%
\TextAt{9}{335.106}{640.028}{\Character{40}}%
\TextAt{13}{339.658}{640.028}{Y}%
\TextAt{9}{349.051}{640.028}{\Character{41}\Character{47}}%
\TextAt{13}{359.457}{640.028}{f}%
\TextAt{3}{365.227}{638.235}{\Character{3}}%
\TextAt{13}{369.959}{640.028}{A}%
\TextAt{9}{378.734}{640.028}{\Character{40}}%
\TextAt{13}{383.287}{640.028}{X}%
\TextAt{9}{393.942}{640.028}{\Character{41}}%
\TextAt{13}{398.494}{640.028}{\Character{59}}%
\TextAt{8}{79.226}{611.027}{for}%
\TextAt{8}{97.111}{611.027}{every}%
\TextAt{8}{128.003}{611.027}{Galois}%
\TextAt{8}{163.917}{611.027}{mapping}%
\TextAt{13}{212.048}{611.027}{f}%
\TextAt{9}{222.415}{611.027}{\Character{58}}%
\TextAt{13}{228.988}{611.027}{X}%
\NativeMathAt{242.964}{611.027}{11.955}{\rightarrow}%
\TextAt{13}{258.24}{611.027}{Y}%
\TextAt{8}{267.633}{611.027}{.}%
\TextAt{12}{79.226}{579.037}{Exercise}%
\TextAt{12}{132.8}{579.037}{IV.6.20.}%
\TextAt{8}{183.611}{579.037}{Under}%
\TextAt{8}{219.371}{579.037}{the}%
\TextAt{8}{239.852}{579.037}{hypotheses}%
\TextAt{8}{299.809}{579.037}{of}%
\TextAt{8}{313.463}{579.037}{exercise}%
\TextAt{8}{357.09}{579.037}{18}%
\TextAt{8}{373.028}{579.037}{and}%
\TextAt{8}{396.115}{579.037}{19}%
\TextAt{8}{412.053}{579.037}{one}%
\TextAt{8}{433.837}{579.037}{obtains}%
\TextAt{8}{475.204}{579.037}{a}%
\TextAt{8}{485.284}{579.037}{canonical}%
\TextAt{8}{79.226}{561.991}{homomorphism}%
\TextAt{8}{161.896}{561.991}{of}%
\TextAt{13}{175.226}{561.991}{G}%
\TextAt{8}{184.46}{561.991}{-modulations}%
\TextAt{13}{280.408}{544.945}{A}%
\NativeArrowAt{292.504}{544.945}{11.955}{19.26}{\longrightarrow}%
\TextAt{13}{315.086}{544.945}{\Character{25}}%
\TextAt{7}{322.155}{549.882}{ab}%
\TextAt{8}{79.226}{521.424}{whose}%
\TextAt{8}{113.593}{521.424}{kernel}%
\TextAt{8}{147.901}{521.424}{is}%
\TextAt{8}{159.506}{521.424}{the}%
\TextAt{13}{179.503}{521.424}{G}%
\TextAt{8}{188.736}{521.424}{-modulation}%
\TextAt{13}{254.594}{521.424}{N}%
\TextAt{13}{265.21}{521.424}{A}%
\TextAt{8}{277.73}{521.424}{of}%
\TextAt{8}{290.901}{521.424}{universal}%
\TextAt{8}{340.236}{521.424}{norms}%
\TextAt{8}{375.261}{521.424}{(see}%
\TextAt{8}{398.569}{521.424}{exercise}%
\TextAt{8}{441.712}{521.424}{13).}%
\TextAt{8}{466.378}{521.424}{It}%
\TextAt{8}{478.891}{521.424}{induces}%
\TextAt{8}{520.408}{521.424}{an}%
\TextAt{8}{79.226}{504.378}{isomorphism}%
\TextAt{20}{286.144}{490.354}{\Character{98}}%
\TextAt{13}{283.452}{487.332}{A}%
\NativeMathAt{298.136}{494.109}{7.97}{\sim}%
\NativeArrowAt{295.548}{487.332}{11.955}{13.175}{\longrightarrow}%
\TextAt{13}{312.042}{487.332}{\Character{25}}%
\TextAt{7}{319.111}{492.268}{ab}%
\TextAt{8}{79.226}{463.81}{of}%
\TextAt{8}{92.556}{463.81}{the}%
\TextAt{8}{112.713}{463.81}{completion}%
\TextAt{20}{175.241}{466.832}{\Character{98}}%
\TextAt{13}{172.548}{463.81}{A}%
\TextAt{8}{185.221}{463.81}{of}%
\TextAt{13}{198.551}{463.81}{A}%
\TextAt{8}{211.224}{463.81}{(see}%
\TextAt{8}{234.692}{463.81}{exercise}%
\TextAt{8}{277.993}{463.81}{14).}%
\TextAt{12}{97.159}{437.798}{Remark:}%
\TextAt{8}{152.515}{437.798}{The}%
\TextAt{8}{176.369}{437.798}{theory}%
\TextAt{8}{212.92}{437.798}{sketched}%
\TextAt{8}{259.931}{437.798}{above}%
\TextAt{8}{292.907}{437.798}{and}%
\TextAt{8}{315.47}{437.798}{contained}%
\TextAt{8}{368.28}{437.798}{in}%
\TextAt{8}{381.733}{437.798}{the}%
\TextAt{8}{401.69}{437.798}{exercises}%
\TextAt{8}{449.407}{437.798}{has}%
\TextAt{8}{470.081}{437.798}{a}%
\TextAt{8}{479.637}{437.798}{very}%
\TextAt{8}{505.129}{437.798}{inter-}%
\TextAt{8}{79.226}{420.752}{esting}%
\TextAt{8}{112.682}{420.752}{application}%
\TextAt{8}{172.751}{420.752}{to}%
\TextAt{8}{186.636}{420.752}{higher}%
\TextAt{8}{221.981}{420.752}{dimensional}%
\TextAt{8}{286.007}{420.752}{class}%
\TextAt{8}{313.019}{420.752}{field}%
\TextAt{8}{337.951}{420.752}{theory.}%
\TextAt{8}{378.13}{420.752}{In}%
\TextAt{8}{392.327}{420.752}{chap.}%
\TextAt{8}{422.794}{420.752}{V,}%
\TextAt{8}{438.294}{420.752}{(1.3),}%
\TextAt{8}{469.181}{420.752}{we}%
\TextAt{8}{485.983}{420.752}{will}%
\TextAt{8}{507.664}{420.752}{show}%
\TextAt{8}{79.226}{403.706}{that,}%
\TextAt{8}{107.847}{403.706}{for}%
\TextAt{8}{125.732}{403.706}{a}%
\TextAt{8}{135.487}{403.706}{Galois}%
\TextAt{8}{171.4}{403.706}{extension}%
\TextAt{13}{223.166}{403.706}{L}%
\TextAt{16}{234.452}{403.706}{\Character{106}}%
\TextAt{13}{241.093}{403.706}{K}%
\TextAt{8}{255.808}{403.706}{of}%
\TextAt{8}{269.139}{403.706}{local}%
\TextAt{8}{296.779}{403.706}{fields,}%
\TextAt{8}{330.002}{403.706}{there}%
\TextAt{8}{359.914}{403.706}{is}%
\TextAt{8}{371.678}{403.706}{a}%
\TextAt{8}{381.433}{403.706}{reciprocity}%
\TextAt{8}{439.643}{403.706}{isomorphism}%
\TextAt{13}{238.25}{374.705}{G}%
\TextAt{9}{247.484}{374.705}{\Character{40}}%
\TextAt{13}{252.036}{374.705}{L}%
\TextAt{16}{263.322}{374.705}{\Character{106}}%
\TextAt{13}{269.963}{374.705}{K}%
\TextAt{9}{280.781}{374.705}{\Character{41}}%
\TextAt{7}{285.333}{379.641}{ab}%
\NativeMathAt{298.087}{374.704}{11.955}{\cong}%
\TextAt{13}{310.706}{374.705}{K}%
\TextAt{3}{321.524}{379.641}{\Character{3}}%
\TextAt{9}{326.256}{374.705}{\Character{47}}%
\TextAt{13}{332.109}{374.705}{N}%
\TextAt{21}{341.483}{372.845}{L}%
\TextAt{3}{347.242}{372.845}{\Character{106}}%
\TextAt{21}{349.594}{372.845}{K}%
\TextAt{13}{357.797}{374.705}{L}%
\TextAt{3}{365.761}{379.641}{\Character{3}}%
\TextAt{13}{370.494}{374.705}{\Character{58}}%
\TextAt{8}{79.226}{345.703}{The}%
\TextAt{8}{103.328}{345.703}{multiplicative}%
\TextAt{8}{177.199}{345.703}{group}%
\TextAt{13}{210.423}{345.703}{K}%
\TextAt{3}{221.24}{350.042}{\Character{3}}%
\TextAt{8}{229.918}{345.703}{may}%
\TextAt{8}{255.335}{345.703}{be}%
\TextAt{8}{271.307}{345.703}{interpreted}%
\TextAt{8}{331.513}{345.703}{in}%
\TextAt{13}{345.214}{345.703}{K}%
\TextAt{8}{356.032}{345.703}{-theory}%
\TextAt{8}{396.727}{345.703}{as}%
\TextAt{8}{411.145}{345.703}{the}%
\TextAt{8}{431.35}{345.703}{group}%
\TextAt{13}{464.573}{345.703}{K}%
\TextAt{17}{474.532}{343.91}{1}%
\TextAt{9}{479.265}{345.703}{\Character{40}}%
\TextAt{13}{483.817}{345.703}{K}%
\TextAt{9}{494.635}{345.703}{\Character{41}}%
\TextAt{8}{503.132}{345.703}{of}%
\TextAt{8}{516.51}{345.703}{the}%
\TextAt{8}{79.226}{328.657}{field}%
\TextAt{13}{104.583}{328.657}{K}%
\TextAt{8}{115.401}{328.657}{.}%
\TextAt{8}{123.849}{328.657}{The}%
\TextAt{8}{147.903}{328.657}{group}%
\TextAt{13}{181.078}{328.657}{K}%
\TextAt{17}{191.038}{326.864}{2}%
\TextAt{9}{195.77}{328.657}{\Character{40}}%
\TextAt{13}{200.322}{328.657}{K}%
\TextAt{9}{211.14}{328.657}{\Character{41}}%
\TextAt{8}{219.59}{328.657}{is}%
\TextAt{8}{231.353}{328.657}{defined}%
\TextAt{8}{271.666}{328.657}{to}%
\TextAt{8}{285.977}{328.657}{be}%
\TextAt{8}{301.901}{328.657}{the}%
\TextAt{8}{322.057}{328.657}{quotient}%
\TextAt{8}{368.24}{328.657}{group}%
\TextAt{13}{244.353}{299.656}{K}%
\TextAt{17}{254.312}{297.863}{2}%
\TextAt{9}{259.044}{299.656}{\Character{40}}%
\TextAt{13}{263.597}{299.656}{K}%
\TextAt{9}{274.414}{299.656}{\Character{41}}%
\TextAt{9}{282.269}{299.656}{\Character{61}}%
\TextAt{9}{294.69}{299.656}{\Character{40}}%
\TextAt{13}{299.265}{299.656}{K}%
\TextAt{3}{310.083}{304.593}{\Character{3}}%
\TextAt{16}{317.472}{299.656}{\Character{10}}%
\TextAt{13}{329.427}{299.656}{K}%
\TextAt{3}{340.245}{304.593}{\Character{3}}%
\TextAt{9}{344.977}{299.656}{\Character{41}\Character{47}}%
\TextAt{13}{355.383}{299.656}{R\Character{59}}%
\TextAt{8}{79.226}{270.655}{where}%
\TextAt{13}{112.974}{270.655}{R}%
\TextAt{8}{125.819}{270.655}{is}%
\TextAt{8}{137.522}{270.655}{generated}%
\TextAt{8}{190.793}{270.655}{by}%
\TextAt{8}{206.991}{270.655}{all}%
\TextAt{8}{223.189}{270.655}{elements}%
\TextAt{8}{270.985}{270.655}{of}%
\TextAt{8}{284.254}{270.655}{the}%
\TextAt{8}{304.35}{270.655}{form}%
\TextAt{13}{331.929}{270.655}{x}%
\TextAt{16}{341.113}{270.655}{\Character{10}}%
\TextAt{9}{352.943}{270.655}{\Character{40}1}%
\TextAt{16}{365.881}{270.655}{\Character{0}}%
\TextAt{13}{377.711}{270.655}{x}%
\TextAt{9}{384.363}{270.655}{\Character{41}}%
\TextAt{8}{388.916}{270.655}{.}%
\TextAt{8}{397.344}{270.655}{Treating}%
\TextAt{8}{444.434}{270.655}{Galois}%
\TextAt{8}{480.286}{270.655}{extensions}%
\TextAt{13}{79.226}{253.609}{L}%
\TextAt{16}{90.511}{253.609}{\Character{106}}%
\TextAt{13}{97.153}{253.609}{K}%
\TextAt{8}{110.871}{253.609}{of}%
\TextAt{8}{123.203}{253.609}{“2-local}%
\TextAt{8}{165.124}{253.609}{fields”—these}%
\TextAt{8}{237.4}{253.609}{are}%
\TextAt{8}{255.913}{253.609}{discretely}%
\TextAt{8}{307.327}{253.609}{valued}%
\TextAt{8}{343.08}{253.609}{complete}%
\TextAt{8}{391.505}{253.609}{fields}%
\TextAt{8}{420.479}{253.609}{with}%
\TextAt{8}{446.141}{253.609}{residue}%
\TextAt{8}{484.87}{253.609}{class}%
\TextAt{8}{511.31}{253.609}{field}%
\TextAt{8}{79.226}{236.563}{a}%
\TextAt{8}{89.781}{236.563}{local}%
\TextAt{8}{118.221}{236.563}{field}%
\TextAt{8}{144.378}{236.563}{(e.g.,}%
\TextAt{18}{174.644}{236.563}{Q}%
\TextAt{21}{183.942}{234.77}{p}%
\TextAt{9}{188.703}{236.563}{\Character{40}\Character{40}}%
\TextAt{13}{197.808}{236.563}{x}%
\TextAt{9}{204.46}{236.563}{\Character{41}\Character{41}}%
\TextAt{8}{213.564}{236.563}{,}%
\TextAt{18}{221.713}{236.563}{F}%
\TextAt{21}{229.019}{234.77}{p}%
\TextAt{9}{233.78}{236.563}{\Character{40}\Character{40}}%
\TextAt{13}{242.885}{236.563}{x}%
\TextAt{9}{249.537}{236.563}{\Character{41}\Character{41}\Character{40}\Character{40}}%
\TextAt{13}{267.746}{236.563}{y}%
\TextAt{9}{273.883}{236.563}{\Character{41}\Character{41}}%
\TextAt{8}{282.988}{236.563}{)—the}%
\TextAt{8}{320.203}{236.563}{Japanese}%
\TextAt{8}{370.652}{236.563}{mathematician}%
\TextAt{26}{451.456}{236.563}{Kazuya}%
\TextAt{26}{501.495}{236.563}{Kato}%
\TextAt{8}{79.226}{219.517}{(see}%
\TextAt{8}{102.694}{219.517}{[83])}%
\TextAt{8}{129.366}{219.517}{has}%
\TextAt{8}{150.24}{219.517}{established}%
\TextAt{8}{210.195}{219.517}{a}%
\TextAt{8}{219.95}{219.517}{canonical}%
\TextAt{8}{271.334}{219.517}{isomorphism}%
\TextAt{13}{219.186}{190.516}{G}%
\TextAt{9}{228.42}{190.516}{\Character{40}}%
\TextAt{13}{232.972}{190.516}{L}%
\TextAt{16}{244.258}{190.516}{\Character{106}}%
\TextAt{13}{250.899}{190.516}{K}%
\TextAt{9}{261.717}{190.516}{\Character{41}}%
\TextAt{7}{266.269}{195.452}{ab}%
\NativeMathAt{279.023}{190.516}{11.955}{\cong}%
\TextAt{13}{291.642}{190.516}{K}%
\TextAt{17}{301.601}{188.723}{2}%
\TextAt{9}{306.334}{190.516}{\Character{40}}%
\TextAt{13}{310.886}{190.516}{K}%
\TextAt{9}{321.704}{190.516}{\Character{41}\Character{47}}%
\TextAt{13}{332.109}{190.516}{N}%
\TextAt{21}{341.483}{188.656}{L}%
\TextAt{3}{347.242}{188.656}{\Character{106}}%
\TextAt{21}{349.594}{188.656}{K}%
\TextAt{13}{357.797}{190.516}{K}%
\TextAt{17}{367.756}{188.723}{2}%
\TextAt{9}{372.488}{190.516}{\Character{40}}%
\TextAt{13}{377.041}{190.516}{L}%
\TextAt{9}{385.005}{190.516}{\Character{41}}%
\TextAt{13}{389.557}{190.516}{\Character{58}}%
\TextAt{26}{79.226}{161.515}{Kato}%
\TextAt{8}{110.501}{161.515}{’s}%
\TextAt{8}{122.889}{161.515}{proof}%
\TextAt{8}{154.083}{161.515}{is}%
\TextAt{8}{166.471}{161.515}{intricate}%
\TextAt{8}{213.601}{161.515}{and}%
\TextAt{8}{236.988}{161.515}{needs}%
\TextAt{8}{269.533}{161.515}{heavy}%
\TextAt{8}{303.655}{161.515}{machinery.}%
\TextAt{8}{365.683}{161.515}{It}%
\TextAt{8}{378.979}{161.515}{was}%
\TextAt{8}{402.103}{161.515}{simplified}%
\TextAt{8}{455.462}{161.515}{by}%
\TextAt{8}{472.345}{161.515}{the}%
\TextAt{8}{493.126}{161.515}{Russian}%
\TextAt{8}{79.226}{144.469}{mathematician}%
\TextAt{8}{160.087}{144.469}{I.}%
\TextAt{26}{172.313}{144.469}{Fesenko}%
\TextAt{8}{226.681}{144.469}{(see}%
\TextAt{8}{251.006}{144.469}{[36],}%
\TextAt{8}{277.445}{144.469}{[37],}%
\TextAt{8}{303.885}{144.469}{[38]).}%
\TextAt{8}{337.678}{144.469}{His}%
\TextAt{8}{359.074}{144.469}{proof}%
\TextAt{8}{390.5}{144.469}{may}%
\TextAt{8}{416.725}{144.469}{be}%
\TextAt{8}{433.506}{144.469}{viewed}%
\TextAt{8}{472.727}{144.469}{as}%
\TextAt{8}{487.954}{144.469}{a}%
\TextAt{8}{498.566}{144.469}{special}%
\TextAt{8}{79.226}{127.423}{case}%
\TextAt{8}{104.454}{127.423}{of}%
\TextAt{8}{118.241}{127.423}{the}%
\TextAt{8}{138.854}{127.423}{theory}%
\TextAt{8}{176.061}{127.423}{sketched}%
\TextAt{8}{223.729}{127.423}{above.}%
\TextAt{8}{262.826}{127.423}{The}%
\TextAt{8}{287.337}{127.423}{basic}%
\TextAt{8}{317.12}{127.423}{idea}%
\TextAt{8}{342.288}{127.423}{is}%
\TextAt{8}{354.509}{127.423}{the}%
\TextAt{8}{375.122}{127.423}{following.}%
\TextAt{8}{430.478}{127.423}{The}%
\TextAt{8}{454.989}{127.423}{correspondence}%
\TextAt{13}{79.226}{110.377}{K}%
\NativeArrowAt{93.365}{110.377}{11.955}{11.955}{\mapsto}%
\TextAt{13}{108.641}{110.377}{K}%
\TextAt{17}{118.6}{108.584}{2}%
\TextAt{9}{123.332}{110.377}{\Character{40}}%
\TextAt{13}{127.884}{110.377}{K}%
\TextAt{9}{138.702}{110.377}{\Character{41}}%
\TextAt{8}{146.283}{110.377}{may}%
\TextAt{8}{170.784}{110.377}{be}%
\TextAt{8}{185.84}{110.377}{extended}%
\TextAt{8}{234.717}{110.377}{to}%
\TextAt{8}{248.158}{110.377}{a}%
\TextAt{13}{257.045}{110.377}{G}%
\TextAt{8}{266.279}{110.377}{-modulation}%
\TextAt{13}{331.427}{110.377}{K}%
\TextAt{17}{341.386}{108.584}{2}%
\TextAt{8}{346.118}{110.377}{.}%
\TextAt{8}{354.277}{110.377}{It}%
\TextAt{8}{366.081}{110.377}{does}%
\TextAt{8}{391.61}{110.377}{not}%
\TextAt{8}{411.555}{110.377}{satisfy}%
\TextAt{8}{447.234}{110.377}{the}%
\TextAt{8}{466.522}{110.377}{hypothesis}%
\TextAt{8}{523.337}{110.377}{of}%
\TextAt{8}{79.226}{93.331}{exercise}%
\TextAt{8}{122.273}{93.331}{15,}%
\TextAt{8}{140.934}{93.331}{so}%
\TextAt{8}{155.049}{93.331}{that}%
\TextAt{8}{180.163}{93.331}{one}%
\TextAt{8}{201.367}{93.331}{may}%
\TextAt{8}{226.481}{93.331}{not}%
\TextAt{8}{247.04}{93.331}{apply}%
\TextAt{8}{278.98}{93.331}{the}%
\TextAt{8}{298.881}{93.331}{abstract}%
\TextAt{8}{344.223}{93.331}{theory}%
\TextAt{8}{380.719}{93.331}{directly}%
\TextAt{8}{423.06}{93.331}{to}%
\TextAt{13}{437.115}{93.331}{K}%
\TextAt{17}{447.074}{91.538}{2}%
\TextAt{8}{451.806}{93.331}{.}%
\TextAt{8}{460.17}{93.331}{But}%
\TextAt{26}{483.156}{93.331}{Fesenko}%
\TextAt{8}{79.226}{76.285}{considers}%
\TextAt{8}{128.784}{76.285}{on}%
\TextAt{13}{144.402}{76.285}{K}%
\TextAt{17}{154.361}{74.492}{2}%
\TextAt{8}{162.349}{76.285}{the}%
\TextAt{8}{181.864}{76.285}{finest}%
\TextAt{8}{212.497}{76.285}{topology}%
\TextAt{8}{259.999}{76.285}{for}%
\TextAt{8}{277.243}{76.285}{which}%
\TextAt{8}{310.087}{76.285}{the}%
\TextAt{8}{329.602}{76.285}{canonical}%
\TextAt{8}{380.344}{76.285}{map}%
\TextAt{9}{405.717}{76.285}{\Character{40}}%
\TextAt{13}{413.525}{76.285}{\Character{59}}%
\TextAt{9}{422.025}{76.285}{\Character{41}}%
\TextAt{9}{429.88}{76.285}{\Character{58}}%
\TextAt{13}{436.471}{76.285}{K}%
\TextAt{3}{447.289}{80.624}{\Character{3}}%
\TextAt{16}{453.366}{76.285}{\Character{2}}%
\TextAt{13}{464.009}{76.285}{K}%
\TextAt{3}{474.827}{80.624}{\Character{3}}%
\NativeMathAt{482.88}{76.285}{11.955}{\rightarrow}%
\TextAt{13}{498.156}{76.285}{K}%
\TextAt{17}{508.115}{74.492}{2}%
\TextAt{9}{512.847}{76.285}{\Character{40}}%
\TextAt{13}{517.399}{76.285}{K}%
\TextAt{9}{528.217}{76.285}{\Character{41}}%
\TextAt{8}{297.211}{43.033}{285}%
\Destination{78.23}{809.815}{nav.page.285}%
\Destination{79.226}{775.155}{nav.ex.4.6.19}%
\Destination{79.226}{593.055}{nav.ex.4.6.20}%
\Destination{79.36}{594.055}{nav.print.310}%
\Hotlink{443.361}{758.955}{12.716}{11.047}{nav.ex.4.6.18}%
\Hotlink{356.59}{576.855}{12.716}{11.047}{nav.ex.4.6.18}%
\Hotlink{441.212}{519.255}{12.716}{11.047}{nav.ex.4.6.13}%
\Hotlink{277.493}{461.655}{12.716}{11.047}{nav.ex.4.6.14}%
\Hotlink{422.294}{418.555}{9.775}{11.047}{nav.chapter.5}%
\Hotlink{102.194}{217.355}{19.22}{11.047}{nav.cite.N083}%
\Hotlink{250.506}{142.255}{19.22}{11.047}{nav.cite.N036}%
\Hotlink{276.945}{142.255}{19.22}{11.047}{nav.cite.N037}%
\Hotlink{303.385}{142.255}{19.22}{11.047}{nav.cite.N038}%
\Hotlink{121.773}{91.155}{12.716}{11.047}{nav.ex.4.6.15}%
\Hotlink{437.794}{418.555}{25.078}{11.047}{nav.number.5.1.3}%
}

% Page 298 (body page 286)
\Page{%
\TextAt{2}{79.226}{796.943}{IV.}%
\TextAt{2}{98.012}{796.943}{Abstract}%
\TextAt{2}{143.252}{796.943}{Class}%
\TextAt{2}{171.844}{796.943}{Field}%
\TextAt{2}{199.575}{796.943}{Theory}%
\Vector{q 0 G 0 g 0.299 w 7.226 18.729 m 460.769 18.729 l S Q}%
\TextAt{8}{79.226}{761.124}{is}%
\TextAt{8}{90.199}{761.124}{sequentially}%
\TextAt{8}{153.536}{761.124}{continuous,}%
\TextAt{8}{214.759}{761.124}{and}%
\TextAt{8}{236.731}{761.124}{for}%
\TextAt{8}{253.825}{761.124}{which}%
\TextAt{8}{286.521}{761.124}{one}%
\TextAt{8}{307.19}{761.124}{has}%
\TextAt{13}{327.273}{761.124}{x}%
\TextAt{21}{333.925}{759.33}{n}%
\TextAt{9}{340.601}{761.124}{\Character{43}}%
\TextAt{13}{350.745}{761.124}{y}%
\TextAt{21}{356.453}{759.33}{n}%
\NativeMathAt{365.41}{761.124}{11.955}{\rightarrow}%
\TextAt{9}{380.686}{761.124}{0}%
\TextAt{8}{386.539}{761.124}{,}%
\TextAt{16}{393.055}{761.124}{\Character{0}}%
\TextAt{13}{402.354}{761.124}{x}%
\TextAt{21}{409.006}{759.33}{n}%
\NativeMathAt{417.963}{761.124}{11.955}{\rightarrow}%
\TextAt{9}{433.239}{761.124}{0}%
\TextAt{8}{442.199}{761.124}{whenever}%
\TextAt{13}{492.779}{761.124}{x}%
\TextAt{21}{499.431}{759.33}{n}%
\NativeMathAt{508.389}{761.124}{11.955}{\rightarrow}%
\TextAt{9}{523.665}{761.124}{0}%
\TextAt{8}{529.518}{761.124}{,}%
\TextAt{13}{79.226}{744.078}{y}%
\TextAt{21}{84.934}{742.284}{n}%
\NativeMathAt{93.891}{744.078}{11.955}{\rightarrow}%
\TextAt{9}{109.167}{744.078}{0}%
\TextAt{8}{115.02}{744.078}{.}%
\TextAt{8}{123.468}{744.078}{He}%
\TextAt{8}{141.341}{744.078}{puts}%
\TextAt{13}{238.143}{727.032}{K}%
\TextAt{7}{248.961}{732.696}{top}%
\TextAt{17}{248.103}{724.101}{2}%
\TextAt{9}{261.685}{727.032}{\Character{40}}%
\TextAt{13}{266.238}{727.032}{K}%
\TextAt{9}{277.055}{727.032}{\Character{41}}%
\TextAt{9}{284.91}{727.032}{\Character{61}}%
\TextAt{13}{297.354}{727.032}{K}%
\TextAt{17}{307.313}{725.238}{2}%
\TextAt{9}{312.046}{727.032}{\Character{40}}%
\TextAt{13}{316.598}{727.032}{K}%
\TextAt{9}{327.416}{727.032}{\Character{41}\Character{47}\Character{3}}%
\TextAt{17}{345.946}{725.238}{2}%
\TextAt{9}{350.678}{727.032}{\Character{40}}%
\TextAt{13}{355.231}{727.032}{K}%
\TextAt{9}{366.048}{727.032}{\Character{41}}%
\TextAt{13}{370.601}{727.032}{\Character{59}}%
\TextAt{8}{79.226}{703.51}{where}%
\TextAt{9}{112.984}{703.51}{\Character{3}}%
\TextAt{17}{121.109}{701.717}{2}%
\TextAt{9}{125.841}{703.51}{\Character{40}}%
\TextAt{13}{130.394}{703.51}{K}%
\TextAt{9}{141.211}{703.51}{\Character{41}}%
\TextAt{8}{149.609}{703.51}{is}%
\TextAt{8}{161.322}{703.51}{the}%
\TextAt{8}{181.427}{703.51}{intersection}%
\TextAt{8}{244.2}{703.51}{of}%
\TextAt{8}{257.478}{703.51}{all}%
\TextAt{8}{273.686}{703.51}{open}%
\TextAt{8}{301.92}{703.51}{neighbourhoods}%
\TextAt{8}{386.487}{703.51}{of}%
\TextAt{9}{399.766}{703.51}{0}%
\TextAt{8}{409.465}{703.51}{in}%
\TextAt{13}{423.066}{703.51}{K}%
\TextAt{17}{433.025}{701.717}{2}%
\TextAt{9}{437.757}{703.51}{\Character{40}}%
\TextAt{13}{442.31}{703.51}{K}%
\TextAt{9}{453.127}{703.51}{\Character{41}}%
\TextAt{8}{457.68}{703.51}{,}%
\TextAt{8}{464.788}{703.51}{and}%
\TextAt{8}{487.499}{703.51}{he}%
\TextAt{8}{503.049}{703.51}{shows}%
\TextAt{8}{79.226}{686.464}{that}%
\TextAt{13}{191.262}{669.418}{K}%
\TextAt{7}{202.08}{675.082}{top}%
\TextAt{17}{201.221}{666.487}{2}%
\TextAt{9}{214.804}{669.418}{\Character{40}}%
\TextAt{13}{219.356}{669.418}{K}%
\TextAt{9}{230.174}{669.418}{\Character{41}\Character{47}}%
\TextAt{13}{240.579}{669.418}{N}%
\TextAt{21}{249.953}{667.558}{L}%
\TextAt{3}{255.712}{667.558}{\Character{106}}%
\TextAt{21}{258.064}{667.558}{K}%
\TextAt{13}{266.267}{669.418}{K}%
\TextAt{7}{277.085}{675.082}{top}%
\TextAt{17}{276.226}{666.487}{2}%
\TextAt{9}{289.809}{669.418}{\Character{40}}%
\TextAt{13}{294.361}{669.418}{L}%
\TextAt{9}{302.326}{669.418}{\Character{41}}%
\NativeMathAt{310.199}{669.418}{11.955}{\cong}%
\TextAt{13}{322.818}{669.418}{K}%
\TextAt{17}{332.778}{667.625}{2}%
\TextAt{9}{337.51}{669.418}{\Character{40}}%
\TextAt{13}{342.062}{669.418}{K}%
\TextAt{9}{352.88}{669.418}{\Character{41}\Character{47}}%
\TextAt{13}{363.285}{669.418}{N}%
\TextAt{21}{372.659}{667.558}{L}%
\TextAt{3}{378.418}{667.558}{\Character{106}}%
\TextAt{21}{380.77}{667.558}{K}%
\TextAt{13}{388.973}{669.418}{K}%
\TextAt{17}{398.932}{667.625}{2}%
\TextAt{9}{403.664}{669.418}{\Character{40}}%
\TextAt{13}{408.217}{669.418}{L}%
\TextAt{9}{416.181}{669.418}{\Character{41}}%
\TextAt{8}{79.226}{645.896}{for}%
\TextAt{8}{97.429}{645.896}{every}%
\TextAt{8}{128.64}{645.896}{Galois}%
\TextAt{8}{164.872}{645.896}{extension}%
\TextAt{13}{216.956}{645.896}{L}%
\TextAt{16}{228.784}{645.896}{\Character{106}}%
\TextAt{13}{235.968}{645.896}{K}%
\TextAt{8}{246.786}{645.896}{,}%
\TextAt{8}{254.333}{645.896}{and}%
\TextAt{8}{277.414}{645.896}{that}%
\TextAt{13}{303.101}{645.896}{K}%
\TextAt{7}{313.919}{651.561}{top}%
\TextAt{17}{313.061}{642.965}{2}%
\TextAt{9}{326.643}{645.896}{\Character{40}}%
\TextAt{13}{331.196}{645.896}{K}%
\TextAt{9}{342.013}{645.896}{\Character{41}}%
\TextAt{8}{350.782}{645.896}{satisfies}%
\TextAt{8}{394.21}{645.896}{properties}%
\TextAt{8}{449.546}{645.896}{which}%
\TextAt{8}{483.351}{645.896}{imply}%
\TextAt{8}{516.51}{645.896}{the}%
\TextAt{8}{79.226}{628.85}{hypothesis}%
\TextAt{8}{138.424}{628.85}{of}%
\TextAt{8}{153.268}{628.85}{exercise}%
\TextAt{8}{198.084}{628.85}{18}%
\TextAt{8}{215.212}{628.85}{and}%
\TextAt{8}{239.488}{628.85}{19}%
\TextAt{8}{256.616}{628.85}{when}%
\TextAt{8}{288.688}{628.85}{viewing}%
\TextAt{13}{332.798}{628.85}{K}%
\TextAt{7}{343.616}{634.515}{top}%
\TextAt{17}{342.757}{625.919}{2}%
\TextAt{8}{361.751}{628.85}{as}%
\TextAt{8}{377.636}{628.85}{a}%
\TextAt{13}{388.905}{628.85}{G}%
\TextAt{8}{398.139}{628.85}{-modulation.}%
\TextAt{8}{473.249}{628.85}{This}%
\TextAt{8}{501.483}{628.85}{makes}%
\TextAt{26}{79.226}{611.804}{Kato}%
\TextAt{8}{110.501}{611.804}{’s}%
\TextAt{8}{122.265}{611.804}{theorem}%
\TextAt{8}{167.79}{611.804}{into}%
\TextAt{8}{191.533}{611.804}{a}%
\TextAt{8}{201.288}{611.804}{special}%
\TextAt{8}{239.39}{611.804}{case}%
\TextAt{8}{264.161}{611.804}{of}%
\TextAt{8}{277.491}{611.804}{the}%
\TextAt{8}{297.647}{611.804}{theory}%
\TextAt{8}{334.397}{611.804}{developed}%
\TextAt{8}{388.698}{611.804}{above.}%
\TextAt{15}{79.226}{564.918}{§}%
\TextAt{15}{93.785}{564.918}{7}%
\TextAt{15}{120.693}{564.918}{The}%
\TextAt{15}{160.202}{564.918}{Herbrand}%
\TextAt{15}{248.173}{564.918}{Quotient}%
\TextAt{8}{79.226}{533.033}{The}%
\TextAt{8}{103.383}{533.033}{preceding}%
\TextAt{8}{156.161}{533.033}{section}%
\TextAt{8}{195.346}{533.033}{concluded}%
\TextAt{8}{250.073}{533.033}{abstract}%
\TextAt{8}{295.773}{533.033}{class}%
\TextAt{8}{323.313}{533.033}{field}%
\TextAt{8}{348.774}{533.033}{theory.}%
\TextAt{8}{389.405}{533.033}{In}%
\TextAt{8}{404.129}{533.033}{order}%
\TextAt{8}{434.802}{533.033}{to}%
\TextAt{8}{449.216}{533.033}{be}%
\TextAt{8}{465.243}{533.033}{able}%
\TextAt{8}{490.058}{533.033}{to}%
\TextAt{8}{504.472}{533.033}{apply}%
\TextAt{8}{79.226}{515.987}{it}%
\TextAt{8}{91.573}{515.987}{to}%
\TextAt{8}{106.526}{515.987}{the}%
\TextAt{8}{127.325}{515.987}{concrete}%
\TextAt{8}{174.138}{515.987}{situations}%
\TextAt{8}{228.245}{515.987}{encountered}%
\TextAt{8}{294.246}{515.987}{in}%
\TextAt{8}{308.542}{515.987}{number}%
\TextAt{8}{351.781}{515.987}{theory,}%
\TextAt{8}{391.606}{515.987}{it}%
\TextAt{8}{403.952}{515.987}{is}%
\TextAt{8}{416.359}{515.987}{all}%
\TextAt{8}{433.261}{515.987}{important}%
\TextAt{8}{489.196}{515.987}{to}%
\TextAt{8}{504.149}{515.987}{verify}%
\TextAt{8}{79.226}{498.941}{the}%
\TextAt{19}{99.549}{498.941}{class}%
\TextAt{19}{127.823}{498.941}{field}%
\TextAt{19}{153.108}{498.941}{axiom}%
\TextAt{8}{187.933}{498.941}{(6.1)}%
\TextAt{8}{216.075}{498.941}{in}%
\TextAt{8}{229.895}{498.941}{these}%
\TextAt{8}{260.033}{498.941}{contexts.}%
\TextAt{8}{311.327}{498.941}{An}%
\TextAt{8}{330.67}{498.941}{excellent}%
\TextAt{8}{378.957}{498.941}{tool}%
\TextAt{8}{402.867}{498.941}{for}%
\TextAt{8}{420.918}{498.941}{this}%
\TextAt{8}{443.908}{498.941}{is}%
\TextAt{8}{455.838}{498.941}{the}%
\TextAt{12}{476.162}{498.941}{Herbrand}%
\TextAt{12}{79.226}{481.895}{quotient}%
\TextAt{8}{127.979}{481.895}{.}%
\TextAt{8}{136.428}{481.895}{It}%
\TextAt{8}{149.1}{481.895}{is}%
\TextAt{8}{160.864}{481.895}{a}%
\TextAt{8}{170.619}{481.895}{group-theoretic}%
\TextAt{8}{252.572}{481.895}{formalism,}%
\TextAt{8}{310.196}{481.895}{which}%
\TextAt{8}{343.682}{481.895}{we}%
\TextAt{8}{360.91}{481.895}{develop}%
\TextAt{8}{403.183}{481.895}{here}%
\TextAt{8}{428.54}{481.895}{for}%
\TextAt{8}{446.425}{481.895}{future}%
\TextAt{8}{481.215}{481.895}{use.}%
\TextAt{8}{97.159}{461.86}{Let}%
\TextAt{13}{117.911}{461.86}{G}%
\TextAt{8}{130.826}{461.86}{be}%
\TextAt{8}{146.533}{461.86}{a}%
\TextAt{8}{156.072}{461.86}{finite}%
\TextAt{8}{185.767}{461.86}{cyclic}%
\TextAt{8}{217.733}{461.86}{group}%
\TextAt{8}{250.692}{461.86}{of}%
\TextAt{8}{263.805}{461.86}{order}%
\TextAt{13}{294.157}{461.86}{n}%
\TextAt{8}{301.145}{461.86}{,}%
\TextAt{8}{308.121}{461.86}{let}%
\TextAt{13}{324.809}{461.86}{\Character{27}}%
\TextAt{8}{335.571}{461.86}{be}%
\TextAt{8}{351.279}{461.86}{a}%
\TextAt{8}{360.817}{461.86}{generator,}%
\TextAt{8}{415.937}{461.86}{and}%
\TextAt{13}{438.482}{461.86}{A}%
\TextAt{8}{450.938}{461.86}{a}%
\TextAt{13}{460.477}{461.86}{G}%
\TextAt{8}{469.71}{461.86}{-module.}%
\TextAt{8}{519.38}{461.86}{As}%
\TextAt{8}{79.226}{444.814}{before,}%
\TextAt{8}{117.59}{444.814}{we}%
\TextAt{8}{134.818}{444.814}{form}%
\TextAt{8}{162.458}{444.814}{the}%
\TextAt{8}{182.614}{444.814}{two}%
\TextAt{8}{204.732}{444.814}{groups}%
\TextAt{13}{156.253}{415.813}{H}%
\TextAt{17}{166.875}{420.749}{0}%
\TextAt{9}{171.608}{415.813}{\Character{40}}%
\TextAt{13}{176.16}{415.813}{G\Character{59}}%
\TextAt{13}{190.614}{415.813}{A}%
\TextAt{9}{199.413}{415.813}{\Character{41}}%
\TextAt{9}{207.268}{415.813}{\Character{61}}%
\TextAt{13}{219.712}{415.813}{A}%
\TextAt{21}{228.487}{420.749}{G}%
\TextAt{9}{235.584}{415.813}{\Character{47}}%
\TextAt{13}{241.437}{415.813}{N}%
\TextAt{21}{250.811}{414.02}{G}%
\TextAt{13}{257.907}{415.813}{A}%
\TextAt{8}{290.593}{415.813}{and}%
\TextAt{13}{333.369}{415.813}{H}%
\TextAt{3}{343.991}{420.749}{\Character{0}}%
\TextAt{17}{350.578}{420.749}{1}%
\TextAt{9}{355.31}{415.813}{\Character{40}}%
\TextAt{13}{359.862}{415.813}{G\Character{59}}%
\TextAt{13}{374.316}{415.813}{A}%
\TextAt{9}{383.116}{415.813}{\Character{41}}%
\TextAt{9}{390.971}{415.813}{\Character{61}}%
\TextAt{21}{403.414}{414.02}{N}%
\TextAt{22}{410.136}{412.68}{G}%
\TextAt{13}{416.828}{415.813}{A}%
\TextAt{9}{425.603}{415.813}{\Character{47}}%
\TextAt{13}{431.456}{415.813}{I}%
\TextAt{21}{436.619}{414.02}{G}%
\TextAt{13}{443.716}{415.813}{A\Character{59}}%
\TextAt{8}{79.226}{386.812}{where}%
\TextAt{13}{132.14}{352.465}{A}%
\TextAt{21}{140.916}{357.401}{G}%
\TextAt{9}{151.333}{352.465}{\Character{61}}%
\TextAt{16}{163.759}{352.465}{\Character{102}}%
\TextAt{13}{169.736}{352.465}{a}%
\NativeMathAt{179.202}{352.465}{11.955}{\in}%
\TextAt{13}{190.493}{352.465}{A}%
\TextAt{16}{202.589}{352.465}{\Character{106}}%
\TextAt{13}{209.231}{352.465}{a}%
\TextAt{21}{215.376}{357.401}{\Character{27}}%
\TextAt{9}{224.32}{352.465}{\Character{61}}%
\TextAt{13}{236.746}{352.465}{a}%
\TextAt{16}{242.891}{352.465}{\Character{103}}%
\TextAt{13}{248.868}{352.465}{\Character{59}}%
\TextAt{13}{309.272}{352.465}{N}%
\TextAt{21}{318.646}{350.672}{G}%
\TextAt{13}{325.743}{352.465}{A}%
\TextAt{9}{337.839}{352.465}{\Character{61}}%
\TextAt{20}{350.265}{372.908}{\Character{40}}%
\TextAt{13}{359.895}{352.465}{N}%
\TextAt{21}{369.269}{350.672}{G}%
\TextAt{13}{376.366}{352.465}{a}%
\TextAt{9}{385.832}{352.465}{\Character{61}}%
\TextAt{21}{398.257}{367.409}{n}%
\TextAt{3}{403.395}{367.409}{\Character{0}}%
\TextAt{17}{409.982}{367.409}{1}%
\TextAt{20}{398.599}{363.822}{\Character{89}}%
\TextAt{21}{399.385}{338.628}{i}%
\TextAt{17}{402.268}{338.628}{\Character{61}0}%
\TextAt{13}{416.209}{352.465}{a}%
\TextAt{21}{422.353}{357.401}{\Character{27}}%
\TextAt{22}{427.479}{360.214}{i}%
\TextAt{20}{434.46}{366.213}{\Character{12}}%
\TextAt{20}{434.46}{359.04}{\Character{12}}%
\TextAt{20}{434.46}{351.867}{\Character{12}}%
\TextAt{13}{441.765}{352.465}{a}%
\NativeMathAt{451.231}{352.465}{11.955}{\in}%
\TextAt{13}{462.522}{352.465}{A}%
\TextAt{20}{471.298}{372.908}{\Character{41}}%
\TextAt{13}{482.921}{352.465}{\Character{59}}%
\TextAt{21}{125.823}{319.722}{N}%
\TextAt{22}{132.545}{318.383}{G}%
\TextAt{13}{139.237}{321.515}{A}%
\TextAt{9}{151.333}{321.515}{\Character{61}}%
\TextAt{16}{163.759}{321.515}{\Character{102}}%
\TextAt{13}{169.736}{321.515}{a}%
\NativeMathAt{179.202}{321.515}{11.955}{\in}%
\TextAt{13}{190.493}{321.515}{A}%
\TextAt{16}{202.589}{321.515}{\Character{106}}%
\TextAt{13}{209.231}{321.515}{N}%
\TextAt{21}{218.605}{319.722}{G}%
\TextAt{13}{225.702}{321.515}{a}%
\TextAt{9}{235.168}{321.515}{\Character{61}}%
\TextAt{9}{247.578}{321.515}{1}%
\TextAt{16}{253.446}{321.515}{\Character{103}}%
\TextAt{13}{259.424}{321.515}{\Character{59}}%
\TextAt{13}{313.483}{321.515}{I}%
\TextAt{21}{318.646}{319.722}{G}%
\TextAt{13}{325.743}{321.515}{A}%
\TextAt{9}{337.839}{321.515}{\Character{61}}%
\TextAt{16}{350.265}{321.515}{\Character{102}}%
\TextAt{13}{356.242}{321.515}{a}%
\TextAt{21}{362.387}{326.451}{\Character{27}}%
\TextAt{3}{367.513}{326.451}{\Character{0}}%
\TextAt{17}{374.099}{326.451}{1}%
\TextAt{16}{382.152}{321.515}{\Character{106}}%
\TextAt{13}{388.794}{321.515}{a}%
\NativeMathAt{398.26}{321.515}{11.955}{\in}%
\TextAt{13}{409.551}{321.515}{A}%
\TextAt{16}{418.326}{321.515}{\Character{103}}%
\TextAt{13}{424.304}{321.515}{\Character{58}}%
\TextAt{12}{79.226}{280.755}{Proposition}%
\TextAt{12}{152.441}{280.755}{(7.1).}%
\TextAt{8}{188.195}{280.755}{If}%
\TextAt{9}{200.325}{280.755}{1}%
\NativeMathAt{210.245}{280.755}{11.955}{\rightarrow}%
\TextAt{13}{226.268}{280.755}{A}%
\TextAt{21}{242.941}{287.532}{i}%
\NativeArrowAt{239.111}{280.755}{11.955}{11.955}{\longrightarrow}%
\TextAt{13}{255.133}{280.755}{B}%
\TextAt{21}{272.027}{288.019}{j}%
\NativeArrowAt{268.697}{280.755}{11.955}{11.955}{\longrightarrow}%
\TextAt{13}{284.719}{280.755}{C}%
\NativeMathAt{298.02}{280.755}{11.955}{\rightarrow}%
\TextAt{9}{314.043}{280.755}{1}%
\TextAt{8}{324.231}{280.755}{is}%
\TextAt{8}{336.433}{280.755}{an}%
\TextAt{8}{353.13}{280.755}{exact}%
\TextAt{8}{384.46}{280.755}{sequence}%
\TextAt{8}{433.401}{280.755}{of}%
\TextAt{13}{447.169}{280.755}{G}%
\TextAt{8}{456.402}{280.755}{-modules,}%
\TextAt{8}{510.007}{280.755}{then}%
\TextAt{8}{79.226}{263.709}{we}%
\TextAt{8}{96.454}{263.709}{obtain}%
\TextAt{8}{132.881}{263.709}{an}%
\TextAt{8}{149.14}{263.709}{exact}%
\TextAt{8}{180.032}{263.709}{hexagon}%
\TextAt{13}{231.61}{232.841}{H}%
\TextAt{17}{242.233}{237.179}{0}%
\TextAt{9}{246.965}{232.841}{\Character{40}}%
\TextAt{13}{251.517}{232.841}{G\Character{59}}%
\TextAt{13}{265.971}{232.841}{A}%
\TextAt{9}{274.77}{232.841}{\Character{41}}%
\TextAt{13}{339.254}{232.841}{H}%
\TextAt{17}{349.876}{237.179}{0}%
\TextAt{9}{354.609}{232.841}{\Character{40}}%
\TextAt{13}{359.161}{232.841}{G\Character{59}}%
\TextAt{13}{373.615}{232.841}{B}%
\TextAt{9}{383.135}{232.841}{\Character{41}}%
\TextAt{13}{119.855}{179.339}{H}%
\TextAt{3}{130.478}{183.677}{\Character{0}}%
\TextAt{17}{137.064}{183.677}{1}%
\TextAt{9}{141.797}{179.339}{\Character{40}}%
\TextAt{13}{146.349}{179.339}{G\Character{59}}%
\TextAt{13}{160.803}{179.339}{C}%
\TextAt{9}{170.06}{179.339}{\Character{41}}%
\TextAt{13}{443.964}{179.339}{H}%
\TextAt{17}{454.586}{183.677}{0}%
\TextAt{9}{459.319}{179.339}{\Character{40}}%
\TextAt{13}{463.871}{179.339}{G\Character{59}}%
\TextAt{13}{478.325}{179.339}{C}%
\TextAt{9}{487.582}{179.339}{\Character{41}}%
\TextAt{13}{227.957}{125.837}{H}%
\TextAt{3}{238.58}{130.175}{\Character{0}}%
\TextAt{17}{245.166}{130.175}{1}%
\TextAt{9}{249.898}{125.837}{\Character{40}}%
\TextAt{13}{254.451}{125.837}{G\Character{59}}%
\TextAt{13}{268.905}{125.837}{B}%
\TextAt{9}{278.425}{125.837}{\Character{41}}%
\TextAt{13}{336.321}{125.837}{H}%
\TextAt{3}{346.943}{130.175}{\Character{0}}%
\TextAt{17}{353.53}{130.175}{1}%
\TextAt{9}{358.262}{125.837}{\Character{40}}%
\TextAt{13}{362.814}{125.837}{G\Character{59}}%
\TextAt{13}{377.268}{125.837}{A}%
\TextAt{9}{386.068}{125.837}{\Character{41}}%
\TextAt{13}{466.424}{125.837}{\Character{58}}%
\NativeDiagram{%
\path[draw=black,line width=0.478bp] (212.678,-538.168) -- (261.427,-538.168);%
\path[draw=black,line width=0.478bp,line cap=round,line join=round] (259.179,-535.299) .. controls (259.634,-537.021) and (260.646,-537.834) .. (261.666,-538.168) .. controls (260.646,-538.503) and (259.634,-539.316) .. (259.179,-541.038);%
}%
\TextAt{21}{305.182}{240.197}{f}%
\TextAt{24}{309.25}{239.09}{1}%
\NativeDiagram{%
\path[draw=black,line width=0.478bp] (312.103,-548.725) -- (374.006,-580.388);%
\path[draw=black,line width=0.478bp,line cap=round,line join=round] (373.312,-576.811) .. controls (372.932,-578.55) and (373.464,-579.735) .. (374.219,-580.497) .. controls (373.159,-580.331) and (371.887,-580.594) .. (370.699,-581.919);%
}%
\TextAt{21}{418.083}{213.7}{f}%
\TextAt{24}{422.151}{212.593}{2}%
\NativeDiagram{%
\path[draw=black,line width=0.478bp] (97.616,-580.606) -- (161.69,-548.937);%
\path[draw=black,line width=0.478bp,line cap=round,line join=round] (158.405,-547.361) .. controls (159.575,-548.703) and (160.842,-548.983) .. (161.905,-548.831) .. controls (161.139,-549.583) and (160.591,-550.76) .. (160.947,-552.505);%
}%
\TextAt{21}{190.833}{213.7}{f}%
\TextAt{24}{194.901}{212.593}{6}%
\NativeDiagram{%
\path[draw=black,line width=0.478bp] (375.425,-602.228) -- (313.522,-633.891);%
\path[draw=black,line width=0.478bp,line cap=round,line join=round] (316.829,-635.422) .. controls (315.641,-634.096) and (314.37,-633.834) .. (313.309,-634) .. controls (314.065,-633.238) and (314.596,-632.053) .. (314.217,-630.313);%
}%
\TextAt{21}{419.076}{147.481}{f}%
\TextAt{24}{423.144}{146.374}{3}%
\NativeDiagram{%
\path[draw=black,line width=0.478bp] (161.088,-634.109) -- (97.013,-602.44);%
\path[draw=black,line width=0.478bp,line cap=round,line join=round] (97.757,-606.007) .. controls (98.112,-604.263) and (97.565,-603.086) .. (96.799,-602.334) .. controls (97.862,-602.486) and (99.129,-602.205) .. (100.299,-600.864);%
}%
\TextAt{21}{189.802}{147.481}{f}%
\TextAt{24}{193.87}{146.374}{5}%
\NativeDiagram{%
\path[draw=black,line width=0.478bp] (258.972,-645.174) -- (216.81,-645.174);%
\path[draw=black,line width=0.478bp,line cap=round,line join=round] (219.057,-648.043) .. controls (218.603,-646.322) and (217.59,-645.509) .. (216.571,-645.174) .. controls (217.59,-644.839) and (218.603,-644.026) .. (219.057,-642.305);%
}%
\TextAt{21}{305.542}{120.476}{f}%
\TextAt{24}{309.61}{119.369}{4}%
\TextAt{8}{297.211}{43.033}{286}%
\Destination{78.23}{809.815}{nav.page.286}%
\Destination{79.23}{591.335}{nav.section.4.7}%
\Destination{79.36}{513.955}{nav.print.311}%
\Destination{79.36}{294.755}{nav.number.4.7.1}%
\Destination{79.23}{591.335}{outline.47}%
\Hotlink{197.584}{626.655}{12.716}{11.047}{nav.ex.4.6.18}%
\Destination{79.226}{720.51}{errata.112}%
}

% Page 299 (body page 287)
\Page{%
\TextAt{2}{79.226}{796.943}{IV.}%
\TextAt{2}{98.012}{796.943}{Abstract}%
\TextAt{2}{143.252}{796.943}{Class}%
\TextAt{2}{171.844}{796.943}{Field}%
\TextAt{2}{199.575}{796.943}{Theory}%
\Vector{q 0 G 0 g 0.299 w 7.226 18.729 m 460.769 18.729 l S Q}%
\TextAt{19}{79.226}{759.35}{Proof.}%
\TextAt{8}{116}{759.35}{The}%
\TextAt{8}{140.937}{759.35}{homomorphisms}%
\TextAt{13}{229.104}{759.35}{f}%
\TextAt{17}{234.874}{757.557}{1}%
\TextAt{13}{239.606}{759.35}{\Character{59}}%
\TextAt{13}{244.83}{759.35}{f}%
\TextAt{17}{250.62}{757.557}{4}%
\TextAt{8}{260.133}{759.35}{and}%
\TextAt{13}{283.778}{759.35}{f}%
\TextAt{17}{289.548}{757.557}{2}%
\TextAt{13}{294.28}{759.35}{\Character{59}}%
\TextAt{13}{299.504}{759.35}{f}%
\TextAt{17}{305.294}{757.557}{5}%
\TextAt{8}{314.807}{759.35}{are}%
\TextAt{8}{335.2}{759.35}{induced}%
\TextAt{8}{379.647}{759.35}{by}%
\TextAt{13}{396.789}{759.35}{A}%
\TextAt{21}{414.22}{766.127}{i}%
\NativeArrowAt{410.389}{759.35}{11.955}{11.955}{\longrightarrow}%
\TextAt{13}{427.17}{759.35}{B}%
\TextAt{8}{441.446}{759.35}{and}%
\TextAt{13}{465.091}{759.35}{B}%
\TextAt{21}{482.742}{766.614}{j}%
\NativeArrowAt{479.412}{759.35}{11.955}{11.955}{\longrightarrow}%
\TextAt{13}{496.193}{759.35}{C}%
\TextAt{8}{505.426}{759.35}{.}%
\TextAt{8}{516.522}{759.35}{We}%
\TextAt{8}{79.226}{742.304}{identify}%
\TextAt{13}{122.833}{742.304}{A}%
\TextAt{8}{136.516}{742.304}{with}%
\TextAt{8}{164.186}{742.304}{its}%
\TextAt{8}{181.515}{742.304}{image}%
\TextAt{8}{216.346}{742.304}{in}%
\TextAt{13}{231.009}{742.304}{B}%
\TextAt{8}{245.413}{742.304}{so}%
\TextAt{8}{260.793}{742.304}{that}%
\TextAt{13}{287.172}{742.304}{i}%
\TextAt{8}{296.073}{742.304}{becomes}%
\TextAt{8}{343.637}{742.304}{the}%
\TextAt{8}{364.804}{742.304}{inclusion.}%
\TextAt{8}{421.222}{742.304}{Then}%
\TextAt{13}{452.79}{742.304}{f}%
\TextAt{17}{458.56}{740.511}{3}%
\TextAt{8}{468.2}{742.304}{is}%
\TextAt{8}{480.974}{742.304}{defined}%
\TextAt{8}{522.297}{742.304}{as}%
\TextAt{8}{79.226}{725.258}{follows.}%
\TextAt{8}{125.418}{725.258}{Let}%
\TextAt{13}{147.456}{725.258}{c}%
\NativeMathAt{157.636}{725.258}{11.955}{\in}%
\TextAt{13}{170.747}{725.258}{C}%
\TextAt{21}{179.981}{729.597}{G}%
\TextAt{8}{192.044}{725.258}{and}%
\TextAt{8}{215.875}{725.258}{let}%
\TextAt{13}{233.848}{725.258}{b}%
\NativeMathAt{243.967}{725.258}{11.955}{\in}%
\TextAt{13}{257.078}{725.258}{B}%
\TextAt{8}{271.54}{725.258}{be}%
\TextAt{8}{288.533}{725.258}{an}%
\TextAt{8}{305.861}{725.258}{element}%
\TextAt{8}{350.171}{725.258}{such}%
\TextAt{8}{377.636}{725.258}{that}%
\TextAt{13}{404.074}{725.258}{j}%
\TextAt{9}{409.589}{725.258}{\Character{40}}%
\TextAt{13}{414.142}{725.258}{b}%
\TextAt{9}{419.119}{725.258}{\Character{41}}%
\TextAt{9}{428.803}{725.258}{\Character{61}}%
\TextAt{13}{443.059}{725.258}{c}%
\TextAt{8}{448.097}{725.258}{.}%
\TextAt{8}{459.75}{725.258}{Then}%
\TextAt{8}{491.376}{725.258}{we}%
\TextAt{8}{509.672}{725.258}{have}%
\TextAt{13}{79.226}{708.212}{j}%
\TextAt{9}{84.742}{708.212}{\Character{40}}%
\TextAt{13}{89.294}{708.212}{b}%
\TextAt{21}{94.271}{712.551}{\Character{27}}%
\TextAt{3}{99.396}{712.551}{\Character{0}}%
\TextAt{17}{105.983}{712.551}{1}%
\TextAt{9}{110.715}{708.212}{\Character{41}}%
\TextAt{9}{119.418}{708.212}{\Character{61}}%
\TextAt{13}{132.702}{708.212}{c}%
\TextAt{21}{137.74}{712.551}{\Character{27}}%
\TextAt{3}{142.865}{712.551}{\Character{0}}%
\TextAt{17}{149.451}{712.551}{1}%
\TextAt{9}{158.348}{708.212}{\Character{61}}%
\TextAt{9}{171.606}{708.212}{1}%
\TextAt{8}{181.863}{708.212}{and}%
\TextAt{13}{205.121}{708.212}{N}%
\TextAt{21}{214.495}{706.419}{G}%
\TextAt{9}{221.592}{708.212}{\Character{40}}%
\TextAt{13}{226.144}{708.212}{b}%
\TextAt{21}{231.121}{712.551}{\Character{27}}%
\TextAt{3}{236.247}{712.551}{\Character{0}}%
\TextAt{17}{242.833}{712.551}{1}%
\TextAt{9}{247.565}{708.212}{\Character{41}}%
\TextAt{9}{256.268}{708.212}{\Character{61}}%
\TextAt{13}{269.552}{708.212}{N}%
\TextAt{21}{278.926}{706.419}{G}%
\TextAt{9}{286.022}{708.212}{\Character{40}}%
\TextAt{13}{290.575}{708.212}{b}%
\TextAt{21}{295.552}{712.551}{\Character{27}}%
\TextAt{9}{301.175}{708.212}{\Character{41}\Character{47}}%
\TextAt{13}{311.581}{708.212}{N}%
\TextAt{21}{320.954}{706.419}{G}%
\TextAt{9}{328.051}{708.212}{\Character{40}}%
\TextAt{13}{332.604}{708.212}{b}%
\TextAt{9}{337.581}{708.212}{\Character{41}}%
\TextAt{9}{346.284}{708.212}{\Character{61}}%
\TextAt{9}{359.543}{708.212}{1}%
\TextAt{8}{365.42}{708.212}{,}%
\TextAt{8}{373.188}{708.212}{so}%
\TextAt{8}{388.054}{708.212}{that}%
\TextAt{13}{413.918}{708.212}{b}%
\TextAt{21}{418.895}{712.551}{\Character{27}}%
\TextAt{3}{424.02}{712.551}{\Character{0}}%
\TextAt{17}{430.607}{712.551}{1}%
\NativeMathAt{439.504}{708.212}{11.955}{\in}%
\TextAt{21}{451.639}{706.419}{N}%
\TextAt{22}{458.36}{705.08}{G}%
\TextAt{13}{465.052}{708.212}{A}%
\TextAt{8}{473.828}{708.212}{.}%
\TextAt{13}{483.762}{708.212}{f}%
\TextAt{17}{489.532}{706.419}{3}%
\TextAt{8}{498.656}{708.212}{is}%
\TextAt{8}{510.916}{708.212}{thus}%
\TextAt{8}{79.226}{691.166}{defined}%
\TextAt{8}{118.912}{691.166}{by}%
\TextAt{13}{134.544}{691.166}{c}%
\TextAt{8}{142.903}{691.166}{mod}%
\TextAt{13}{168.341}{691.166}{N}%
\TextAt{21}{177.715}{689.373}{G}%
\TextAt{13}{184.811}{691.166}{C}%
\NativeArrowAt{197.366}{691.166}{11.955}{11.955}{\mapsto}%
\TextAt{13}{212.642}{691.166}{b}%
\TextAt{21}{217.619}{695.505}{\Character{27}}%
\TextAt{3}{222.744}{695.505}{\Character{0}}%
\TextAt{17}{229.331}{695.505}{1}%
\TextAt{8}{237.384}{691.166}{mod}%
\TextAt{13}{262.822}{691.166}{I}%
\TextAt{21}{267.985}{689.373}{G}%
\TextAt{13}{275.082}{691.166}{A}%
\TextAt{8}{283.857}{691.166}{.}%
\TextAt{8}{292.096}{691.166}{In}%
\TextAt{8}{306.091}{691.166}{order}%
\TextAt{8}{336.033}{691.166}{to}%
\TextAt{8}{349.716}{691.166}{define}%
\TextAt{13}{382.899}{691.166}{f}%
\TextAt{17}{388.669}{689.373}{6}%
\TextAt{8}{393.401}{691.166}{,}%
\TextAt{8}{400.049}{691.166}{let}%
\TextAt{13}{416.326}{691.166}{c}%
\NativeMathAt{424.685}{691.166}{11.955}{\in}%
\TextAt{21}{435.976}{689.373}{N}%
\TextAt{22}{442.697}{688.034}{G}%
\TextAt{13}{449.39}{691.166}{C}%
\TextAt{8}{458.623}{691.166}{,}%
\TextAt{8}{465.271}{691.166}{and}%
\TextAt{8}{487.407}{691.166}{let}%
\TextAt{13}{503.684}{691.166}{b}%
\NativeMathAt{511.982}{691.166}{11.955}{\in}%
\TextAt{13}{523.273}{691.166}{B}%
\TextAt{8}{79.226}{674.12}{be}%
\TextAt{8}{95.341}{674.12}{an}%
\TextAt{8}{111.791}{674.12}{element}%
\TextAt{8}{155.224}{674.12}{such}%
\TextAt{8}{181.811}{674.12}{that}%
\TextAt{13}{207.371}{674.12}{j}%
\TextAt{9}{212.887}{674.12}{\Character{40}}%
\TextAt{13}{217.439}{674.12}{b}%
\TextAt{9}{222.416}{674.12}{\Character{41}}%
\TextAt{9}{230.593}{674.12}{\Character{61}}%
\TextAt{13}{243.365}{674.12}{c}%
\TextAt{8}{248.403}{674.12}{.}%
\TextAt{8}{257.423}{674.12}{Then}%
\TextAt{13}{288.172}{674.12}{j}%
\TextAt{9}{293.687}{674.12}{\Character{40}}%
\TextAt{13}{298.239}{674.12}{N}%
\TextAt{21}{307.613}{672.327}{G}%
\TextAt{13}{314.71}{674.12}{b}%
\TextAt{9}{319.687}{674.12}{\Character{41}}%
\TextAt{9}{327.864}{674.12}{\Character{61}}%
\TextAt{13}{340.636}{674.12}{N}%
\TextAt{21}{350.01}{672.327}{G}%
\TextAt{13}{357.107}{674.12}{c}%
\TextAt{9}{365.791}{674.12}{\Character{61}}%
\TextAt{9}{378.523}{674.12}{1}%
\TextAt{8}{384.394}{674.12}{,}%
\TextAt{8}{391.782}{674.12}{so}%
\TextAt{8}{406.343}{674.12}{that}%
\TextAt{13}{431.902}{674.12}{N}%
\TextAt{21}{441.276}{672.327}{G}%
\TextAt{13}{448.373}{674.12}{b}%
\NativeMathAt{456.996}{674.12}{11.955}{\in}%
\TextAt{13}{468.612}{674.12}{A}%
\TextAt{8}{477.387}{674.12}{.}%
\TextAt{8}{486.408}{674.12}{The}%
\TextAt{8}{510.653}{674.12}{map}%
\ErrataAt{79.19}{657.1}{455.81}{$f_6$ is now given by $c\bmod I_GC\longmapsto N_Gb\bmod N_GA$.}%
\TextAt{8}{97.159}{637.04}{We}%
\TextAt{8}{118.869}{637.04}{now}%
\TextAt{8}{144.824}{637.04}{prove}%
\TextAt{8}{177.94}{637.04}{exactness}%
\TextAt{8}{231.332}{637.04}{at}%
\TextAt{8}{247.208}{637.04}{the}%
\TextAt{8}{268.931}{637.04}{place}%
\TextAt{13}{300.409}{637.04}{H}%
\TextAt{17}{311.031}{641.378}{0}%
\TextAt{9}{315.764}{637.04}{\Character{40}}%
\TextAt{13}{320.316}{637.04}{G\Character{59}}%
\TextAt{13}{334.77}{637.04}{A}%
\TextAt{9}{343.569}{637.04}{\Character{41}}%
\TextAt{8}{348.121}{637.04}{.}%
\TextAt{8}{361.268}{637.04}{Let}%
\TextAt{13}{383.804}{637.04}{a}%
\NativeMathAt{395.938}{637.04}{11.955}{\in}%
\TextAt{13}{409.898}{637.04}{A}%
\TextAt{21}{418.674}{641.378}{G}%
\TextAt{8}{431.234}{637.04}{such}%
\TextAt{8}{459.197}{637.04}{that}%
\TextAt{13}{486.132}{637.04}{f}%
\TextAt{17}{491.902}{635.246}{1}%
\TextAt{9}{496.634}{637.04}{\Character{40}}%
\TextAt{13}{501.187}{637.04}{a}%
\TextAt{8}{510.652}{637.04}{mod}%
\TextAt{13}{79.226}{619.994}{N}%
\TextAt{21}{88.6}{618.2}{G}%
\TextAt{13}{95.697}{619.994}{A}%
\TextAt{9}{104.472}{619.994}{\Character{41}}%
\TextAt{9}{112.649}{619.994}{\Character{61}}%
\TextAt{9}{125.394}{619.994}{1}%
\TextAt{8}{131.263}{619.994}{;}%
\TextAt{8}{138.693}{619.994}{in}%
\TextAt{8}{152.533}{619.994}{other}%
\TextAt{8}{183.29}{619.994}{words,}%
\TextAt{13}{220.334}{619.994}{a}%
\TextAt{9}{230.12}{619.994}{\Character{61}}%
\TextAt{13}{242.865}{619.994}{N}%
\TextAt{21}{252.238}{618.2}{G}%
\TextAt{13}{259.335}{619.994}{b}%
\TextAt{8}{268.397}{619.994}{for}%
\TextAt{8}{286.47}{619.994}{some}%
\TextAt{13}{315.983}{619.994}{b}%
\NativeMathAt{324.601}{619.994}{11.955}{\in}%
\TextAt{13}{336.211}{619.994}{B}%
\TextAt{8}{345.707}{619.994}{.}%
\TextAt{8}{354.718}{619.994}{Writing}%
\TextAt{13}{397.824}{619.994}{c}%
\TextAt{9}{406.503}{619.994}{\Character{61}}%
\TextAt{13}{419.248}{619.994}{j}%
\TextAt{9}{424.763}{619.994}{\Character{40}}%
\TextAt{13}{429.316}{619.994}{b}%
\TextAt{9}{434.293}{619.994}{\Character{41}}%
\TextAt{8}{438.845}{619.994}{,}%
\TextAt{8}{446.229}{619.994}{we}%
\TextAt{8}{463.643}{619.994}{find}%
\TextAt{13}{487.239}{619.994}{f}%
\TextAt{17}{493.009}{618.2}{6}%
\TextAt{9}{497.741}{619.994}{\Character{40}}%
\TextAt{13}{502.294}{619.994}{c}%
\TextAt{8}{510.653}{619.994}{mod}%
\TextAt{13}{79.226}{602.948}{I}%
\TextAt{21}{84.389}{601.154}{G}%
\TextAt{13}{91.486}{602.948}{C}%
\TextAt{9}{100.719}{602.948}{\Character{41}}%
\TextAt{9}{110.2}{602.948}{\Character{61}}%
\TextAt{13}{124.254}{602.948}{a}%
\TextAt{8}{133.719}{602.948}{mod}%
\TextAt{13}{159.157}{602.948}{N}%
\TextAt{21}{168.531}{601.154}{G}%
\TextAt{13}{175.628}{602.948}{A}%
\TextAt{8}{184.403}{602.948}{.}%
\TextAt{8}{195.699}{602.948}{Exactness}%
\TextAt{8}{251.236}{602.948}{at}%
\TextAt{13}{266.496}{602.948}{H}%
\TextAt{3}{277.118}{607.286}{\Character{0}}%
\TextAt{17}{283.705}{607.286}{1}%
\TextAt{9}{288.437}{602.948}{\Character{40}}%
\TextAt{13}{292.989}{602.948}{G\Character{59}}%
\TextAt{13}{307.443}{602.948}{A}%
\TextAt{9}{316.243}{602.948}{\Character{41}}%
\TextAt{8}{325.642}{602.948}{is}%
\TextAt{8}{338.355}{602.948}{deduced}%
\TextAt{8}{384.817}{602.948}{as}%
\TextAt{8}{400.137}{602.948}{follows:}%
\TextAt{8}{445.022}{602.948}{let}%
\TextAt{13}{462.876}{602.948}{a}%
\NativeMathAt{473.96}{602.948}{11.955}{\in}%
\TextAt{21}{486.868}{601.154}{N}%
\TextAt{22}{493.59}{599.815}{G}%
\TextAt{13}{500.282}{602.948}{A}%
\TextAt{8}{513.904}{602.948}{and}%
\TextAt{13}{79.226}{585.902}{f}%
\TextAt{17}{84.996}{584.108}{4}%
\TextAt{9}{89.728}{585.902}{\Character{40}}%
\TextAt{13}{94.281}{585.902}{a}%
\TextAt{8}{103.747}{585.902}{mod}%
\TextAt{13}{129.184}{585.902}{I}%
\TextAt{21}{134.347}{584.108}{G}%
\TextAt{13}{141.444}{585.902}{A}%
\TextAt{9}{150.22}{585.902}{\Character{41}}%
\TextAt{9}{160.466}{585.902}{\Character{61}}%
\TextAt{9}{175.266}{585.902}{1}%
\TextAt{8}{181.141}{585.902}{,}%
\TextAt{8}{190.039}{585.902}{i.e.,}%
\TextAt{13}{213.894}{585.902}{a}%
\TextAt{9}{225.745}{585.902}{\Character{61}}%
\TextAt{13}{240.555}{585.902}{b}%
\TextAt{21}{245.532}{590.24}{\Character{27}}%
\TextAt{3}{250.657}{590.24}{\Character{0}}%
\TextAt{17}{257.244}{590.24}{1}%
\TextAt{8}{261.976}{585.902}{,}%
\TextAt{8}{270.875}{585.902}{with}%
\TextAt{13}{298.934}{585.902}{b}%
\NativeMathAt{309.617}{585.902}{11.955}{\in}%
\TextAt{13}{323.293}{585.902}{B}%
\TextAt{8}{332.789}{585.902}{.}%
\TextAt{8}{345.436}{585.902}{Writing}%
\TextAt{13}{389.755}{585.902}{c}%
\TextAt{9}{400.498}{585.902}{\Character{61}}%
\TextAt{13}{415.309}{585.902}{j}%
\TextAt{9}{420.824}{585.902}{\Character{40}}%
\TextAt{13}{425.377}{585.902}{b}%
\TextAt{9}{430.354}{585.902}{\Character{41}}%
\TextAt{8}{434.906}{585.902}{,}%
\TextAt{8}{443.805}{585.902}{we}%
\TextAt{8}{462.431}{585.902}{find}%
\TextAt{13}{487.239}{585.902}{f}%
\TextAt{17}{493.009}{584.108}{3}%
\TextAt{9}{497.741}{585.902}{\Character{40}}%
\TextAt{13}{502.294}{585.902}{c}%
\TextAt{8}{510.653}{585.902}{mod}%
\TextAt{13}{79.226}{568.856}{N}%
\TextAt{21}{88.6}{567.063}{G}%
\TextAt{13}{95.697}{568.856}{C}%
\TextAt{9}{104.93}{568.856}{\Character{41}}%
\TextAt{9}{112.785}{568.856}{\Character{61}}%
\TextAt{13}{125.229}{568.856}{a}%
\TextAt{8}{134.695}{568.856}{mod}%
\TextAt{13}{160.133}{568.856}{I}%
\TextAt{21}{165.296}{567.063}{G}%
\TextAt{13}{172.392}{568.856}{A}%
\TextAt{8}{181.168}{568.856}{.}%
\TextAt{8}{189.616}{568.856}{The}%
\TextAt{8}{213.67}{568.856}{exactness}%
\TextAt{8}{265.496}{568.856}{at}%
\TextAt{8}{279.806}{568.856}{all}%
\TextAt{8}{296.065}{568.856}{other}%
\TextAt{8}{326.634}{568.856}{places}%
\TextAt{8}{361.161}{568.856}{is}%
\TextAt{8}{372.925}{568.856}{seen}%
\TextAt{8}{398.341}{568.856}{even}%
\TextAt{8}{425.001}{568.856}{more}%
\TextAt{8}{454.268}{568.856}{easily.}%
\Vector{q 0 G 0 g 0.399 w 452.334 -205.144 m 452.334 -197.074 l S Q}%
\Vector{q 0 G 0 g 0.399 w 452.534 -197.274 m 459.707 -197.274 l S Q}%
\Vector{q 0 G 0 g 0.399 w 452.534 -204.945 m 459.707 -204.945 l S Q}%
\Vector{q 0 G 0 g 0.399 w 459.906 -205.144 m 459.906 -197.074 l S Q}%
\TextAt{12}{79.226}{536.622}{Definition}%
\TextAt{12}{142.015}{536.622}{(7.2).}%
\TextAt{8}{177.33}{536.622}{The}%
\TextAt{12}{201.384}{536.622}{Herbrand}%
\TextAt{12}{262.475}{536.622}{quotient}%
\TextAt{8}{315.126}{536.622}{of}%
\TextAt{8}{328.456}{536.622}{the}%
\TextAt{13}{348.612}{536.622}{G}%
\TextAt{8}{357.846}{536.622}{-module}%
\TextAt{13}{403.036}{536.622}{A}%
\TextAt{8}{415.709}{536.622}{is}%
\TextAt{8}{427.473}{536.622}{defined}%
\TextAt{8}{467.786}{536.622}{to}%
\TextAt{8}{482.096}{536.622}{be}%
\TextAt{13}{243.728}{507.91}{h}%
\TextAt{9}{250.466}{507.91}{\Character{40}}%
\TextAt{13}{255.019}{507.91}{G\Character{59}}%
\TextAt{13}{269.473}{507.91}{A}%
\TextAt{9}{278.272}{507.91}{\Character{41}}%
\TextAt{9}{286.127}{507.91}{\Character{61}}%
\TextAt{8}{303.059}{515.998}{\char35\relax }%
\TextAt{13}{312.815}{515.998}{H}%
\TextAt{17}{323.437}{520.337}{0}%
\TextAt{9}{328.17}{515.998}{\Character{40}}%
\TextAt{13}{332.722}{515.998}{G\Character{59}}%
\TextAt{13}{347.176}{515.998}{A}%
\TextAt{9}{355.975}{515.998}{\Character{41}}%
\Vector{q 0 G 0 g 0.478 w 227.766 -263.101 m 291.821 -263.101 l S Q}%
\TextAt{8}{299.766}{499.71}{\char35\relax }%
\TextAt{13}{309.521}{499.71}{H}%
\TextAt{3}{320.144}{503.163}{\Character{0}}%
\TextAt{17}{326.731}{503.163}{1}%
\TextAt{9}{331.463}{499.71}{\Character{40}}%
\TextAt{13}{336.015}{499.71}{G\Character{59}}%
\TextAt{13}{350.469}{499.71}{A}%
\TextAt{9}{359.268}{499.71}{\Character{41}}%
\TextAt{13}{365.016}{507.91}{\Character{59}}%
\TextAt{8}{79.226}{480.156}{provided}%
\TextAt{8}{127.358}{480.156}{that}%
\TextAt{8}{152.727}{480.156}{both}%
\TextAt{8}{180.367}{480.156}{orders}%
\TextAt{8}{215.551}{480.156}{are}%
\TextAt{8}{235.062}{480.156}{finite.}%
\TextAt{8}{97.159}{455.516}{The}%
\TextAt{8}{121.213}{455.516}{salient}%
\TextAt{8}{158.023}{455.516}{property}%
\TextAt{8}{205.831}{455.516}{of}%
\TextAt{8}{219.161}{455.516}{the}%
\TextAt{8}{239.318}{455.516}{Herbrand}%
\TextAt{8}{291.67}{455.516}{quotient}%
\TextAt{8}{337.852}{455.516}{is}%
\TextAt{8}{349.616}{455.516}{its}%
\TextAt{19}{365.935}{455.516}{multiplicativity}%
\TextAt{8}{441.551}{455.516}{.}%
\TextAt{12}{79.226}{430.876}{Proposition}%
\TextAt{12}{152.441}{430.876}{(7.3).}%
\TextAt{8}{188.195}{430.876}{If}%
\TextAt{9}{200.325}{430.876}{1}%
\NativeMathAt{210.246}{430.876}{11.955}{\rightarrow}%
\TextAt{13}{226.268}{430.876}{A}%
\NativeMathAt{239.111}{430.876}{11.955}{\rightarrow}%
\TextAt{13}{255.133}{430.876}{B}%
\NativeMathAt{268.697}{430.876}{11.955}{\rightarrow}%
\TextAt{13}{284.719}{430.876}{C}%
\NativeMathAt{298.02}{430.876}{11.955}{\rightarrow}%
\TextAt{9}{314.043}{430.876}{1}%
\TextAt{8}{324.231}{430.876}{is}%
\TextAt{8}{336.433}{430.876}{an}%
\TextAt{8}{353.13}{430.876}{exact}%
\TextAt{8}{384.46}{430.876}{sequence}%
\TextAt{8}{433.401}{430.876}{of}%
\TextAt{13}{447.169}{430.876}{G}%
\TextAt{8}{456.402}{430.876}{-modules,}%
\TextAt{8}{510.007}{430.876}{then}%
\TextAt{8}{79.226}{413.83}{one}%
\TextAt{8}{100.686}{413.83}{has}%
\TextAt{13}{238.891}{396.784}{h}%
\TextAt{9}{245.629}{396.785}{\Character{40}}%
\TextAt{13}{250.181}{396.784}{G\Character{59}}%
\TextAt{13}{264.635}{396.784}{B}%
\TextAt{9}{274.156}{396.785}{\Character{41}}%
\TextAt{9}{282.011}{396.785}{\Character{61}}%
\TextAt{13}{294.454}{396.784}{h}%
\TextAt{9}{301.193}{396.785}{\Character{40}}%
\TextAt{13}{305.745}{396.784}{G\Character{59}}%
\TextAt{13}{320.199}{396.784}{A}%
\TextAt{9}{328.998}{396.785}{\Character{41}}%
\TextAt{13}{333.55}{396.784}{h}%
\TextAt{9}{340.289}{396.785}{\Character{40}}%
\TextAt{13}{344.841}{396.784}{G\Character{59}}%
\TextAt{13}{359.295}{396.784}{C}%
\TextAt{9}{368.553}{396.785}{\Character{41}}%
\TextAt{8}{79.226}{375.321}{in}%
\TextAt{8}{93.839}{375.321}{the}%
\TextAt{8}{114.955}{375.321}{sense}%
\TextAt{8}{145.947}{375.321}{that,}%
\TextAt{8}{175.768}{375.321}{whenever}%
\TextAt{8}{228.099}{375.321}{two}%
\TextAt{8}{251.176}{375.321}{of}%
\TextAt{8}{265.466}{375.321}{these}%
\TextAt{8}{296.398}{375.321}{quotients}%
\TextAt{8}{348.155}{375.321}{are}%
\TextAt{8}{368.626}{375.321}{defined,}%
\TextAt{8}{413.391}{375.321}{so}%
\TextAt{8}{428.721}{375.321}{is}%
\TextAt{8}{441.445}{375.321}{the}%
\TextAt{8}{462.561}{375.321}{third}%
\TextAt{8}{492.788}{375.321}{and}%
\TextAt{8}{516.51}{375.321}{the}%
\TextAt{8}{79.226}{358.275}{identity}%
\TextAt{8}{122.48}{358.275}{holds.}%
\TextAt{8}{97.159}{338.241}{For}%
\TextAt{8}{118.128}{338.241}{a}%
\TextAt{8}{127.884}{338.241}{finite}%
\TextAt{13}{157.796}{338.241}{G}%
\TextAt{8}{167.029}{338.241}{-module}%
\TextAt{13}{212.22}{338.241}{A}%
\TextAt{8}{220.995}{338.241}{,}%
\TextAt{8}{228.144}{338.241}{one}%
\TextAt{8}{249.604}{338.241}{has}%
\TextAt{13}{270.478}{338.241}{h}%
\TextAt{9}{277.216}{338.241}{\Character{40}}%
\TextAt{13}{281.768}{338.241}{G\Character{59}}%
\TextAt{13}{296.222}{338.241}{A}%
\TextAt{9}{305.022}{338.241}{\Character{41}}%
\TextAt{9}{312.877}{338.241}{\Character{61}}%
\TextAt{9}{325.298}{338.241}{1}%
\TextAt{8}{331.173}{338.241}{.}%
\TextAt{19}{79.226}{306.007}{Proof.}%
\TextAt{8}{116}{306.007}{We}%
\TextAt{8}{136.494}{306.007}{consider}%
\TextAt{8}{182.428}{306.007}{the}%
\TextAt{8}{202.933}{306.007}{exact}%
\TextAt{8}{234.174}{306.007}{hexagon}%
\TextAt{8}{280.383}{306.007}{(7.1).}%
\TextAt{8}{313.956}{306.007}{Calling}%
\TextAt{13}{354.63}{306.007}{n}%
\TextAt{21}{361.618}{304.213}{i}%
\TextAt{8}{369.245}{306.007}{the}%
\TextAt{8}{389.751}{306.007}{order}%
\TextAt{8}{420.669}{306.007}{of}%
\TextAt{8}{434.348}{306.007}{the}%
\TextAt{8}{454.854}{306.007}{image}%
\TextAt{8}{489.024}{306.007}{of}%
\TextAt{13}{502.703}{306.007}{f}%
\TextAt{21}{508.473}{304.213}{i}%
\TextAt{8}{511.854}{306.007}{,}%
\TextAt{8}{519.439}{306.007}{we}%
\TextAt{8}{79.226}{288.961}{find}%
\TextAt{8}{119.034}{264.763}{\char35\relax }%
\TextAt{13}{128.79}{264.763}{H}%
\TextAt{17}{139.412}{269.699}{0}%
\TextAt{9}{144.145}{264.763}{\Character{40}}%
\TextAt{13}{148.697}{264.763}{G\Character{59}}%
\TextAt{13}{163.151}{264.763}{A}%
\TextAt{9}{171.95}{264.763}{\Character{41}}%
\TextAt{9}{179.805}{264.763}{\Character{61}}%
\TextAt{13}{192.249}{264.763}{n}%
\TextAt{17}{199.236}{262.97}{6}%
\TextAt{13}{203.969}{264.763}{n}%
\TextAt{17}{210.956}{262.97}{1}%
\TextAt{13}{215.689}{264.763}{\Character{59}}%
\TextAt{8}{258.748}{264.763}{\char35\relax }%
\TextAt{13}{268.504}{264.763}{H}%
\TextAt{17}{279.127}{269.699}{0}%
\TextAt{9}{283.859}{264.763}{\Character{40}}%
\TextAt{13}{288.411}{264.763}{G\Character{59}}%
\TextAt{13}{302.865}{264.763}{B}%
\TextAt{9}{312.385}{264.763}{\Character{41}}%
\TextAt{9}{320.24}{264.763}{\Character{61}}%
\TextAt{13}{332.684}{264.763}{n}%
\TextAt{17}{339.671}{262.97}{1}%
\TextAt{13}{344.404}{264.763}{n}%
\TextAt{17}{351.391}{262.97}{2}%
\TextAt{13}{356.124}{264.763}{\Character{59}}%
\TextAt{8}{399.184}{264.763}{\char35\relax }%
\TextAt{13}{408.939}{264.763}{H}%
\TextAt{17}{419.562}{269.699}{0}%
\TextAt{9}{424.294}{264.763}{\Character{40}}%
\TextAt{13}{428.846}{264.763}{G\Character{59}}%
\TextAt{13}{443.3}{264.763}{C}%
\TextAt{9}{452.558}{264.763}{\Character{41}}%
\TextAt{9}{460.413}{264.763}{\Character{61}}%
\TextAt{13}{472.856}{264.763}{n}%
\TextAt{17}{479.844}{262.97}{2}%
\TextAt{13}{484.576}{264.763}{n}%
\TextAt{17}{491.564}{262.97}{3}%
\TextAt{13}{496.296}{264.763}{\Character{59}}%
\TextAt{8}{112.448}{244.728}{\char35\relax }%
\TextAt{13}{122.203}{244.728}{H}%
\TextAt{3}{132.826}{249.664}{\Character{0}}%
\TextAt{17}{139.412}{249.664}{1}%
\TextAt{9}{144.145}{244.728}{\Character{40}}%
\TextAt{13}{148.697}{244.728}{G\Character{59}}%
\TextAt{13}{163.151}{244.728}{A}%
\TextAt{9}{171.95}{244.728}{\Character{41}}%
\TextAt{9}{179.805}{244.728}{\Character{61}}%
\TextAt{13}{192.249}{244.728}{n}%
\TextAt{17}{199.236}{242.935}{3}%
\TextAt{13}{203.969}{244.728}{n}%
\TextAt{17}{210.956}{242.935}{4}%
\TextAt{13}{215.689}{244.728}{\Character{59}}%
\TextAt{8}{252.162}{244.728}{\char35\relax }%
\TextAt{13}{261.917}{244.728}{H}%
\TextAt{3}{272.54}{249.664}{\Character{0}}%
\TextAt{17}{279.127}{249.664}{1}%
\TextAt{9}{283.859}{244.728}{\Character{40}}%
\TextAt{13}{288.411}{244.728}{G\Character{59}}%
\TextAt{13}{302.865}{244.728}{B}%
\TextAt{9}{312.385}{244.728}{\Character{41}}%
\TextAt{9}{320.24}{244.728}{\Character{61}}%
\TextAt{13}{332.684}{244.728}{n}%
\TextAt{17}{339.671}{242.935}{4}%
\TextAt{13}{344.404}{244.728}{n}%
\TextAt{17}{351.391}{242.935}{5}%
\TextAt{13}{356.124}{244.728}{\Character{59}}%
\TextAt{8}{392.597}{244.728}{\char35\relax }%
\TextAt{13}{402.352}{244.728}{H}%
\TextAt{3}{412.975}{249.664}{\Character{0}}%
\TextAt{17}{419.562}{249.664}{1}%
\TextAt{9}{424.294}{244.728}{\Character{40}}%
\TextAt{13}{428.846}{244.728}{G\Character{59}}%
\TextAt{13}{443.3}{244.728}{C}%
\TextAt{9}{452.558}{244.728}{\Character{41}}%
\TextAt{9}{460.413}{244.728}{\Character{61}}%
\TextAt{13}{472.856}{244.728}{n}%
\TextAt{17}{479.844}{242.935}{5}%
\TextAt{13}{484.576}{244.728}{n}%
\TextAt{17}{491.564}{242.935}{6}%
\TextAt{13}{496.296}{244.728}{\Character{59}}%
\TextAt{8}{79.226}{220.53}{and}%
\TextAt{8}{101.989}{220.53}{thus}%
\TextAt{8}{218.41}{196.332}{\char35\relax }%
\TextAt{13}{228.166}{196.332}{H}%
\TextAt{17}{238.788}{201.269}{0}%
\TextAt{9}{243.521}{196.332}{\Character{40}}%
\TextAt{13}{248.073}{196.332}{G\Character{59}}%
\TextAt{13}{262.527}{196.332}{A}%
\TextAt{9}{271.326}{196.332}{\Character{41}}%
\TextAt{16}{278.535}{196.332}{\Character{1}}%
\TextAt{8}{284.513}{196.332}{\char35\relax }%
\TextAt{13}{294.268}{196.332}{H}%
\TextAt{17}{304.891}{201.269}{0}%
\TextAt{9}{309.623}{196.332}{\Character{40}}%
\TextAt{13}{314.175}{196.332}{G\Character{59}}%
\TextAt{13}{328.629}{196.332}{C}%
\TextAt{9}{337.887}{196.332}{\Character{41}}%
\TextAt{16}{345.096}{196.332}{\Character{1}}%
\TextAt{8}{351.073}{196.332}{\char35\relax }%
\TextAt{13}{360.829}{196.332}{H}%
\TextAt{3}{371.451}{201.269}{\Character{0}}%
\TextAt{17}{378.038}{201.269}{1}%
\TextAt{9}{382.77}{196.332}{\Character{40}}%
\TextAt{13}{387.323}{196.332}{G\Character{59}}%
\TextAt{13}{401.777}{196.332}{B}%
\TextAt{9}{411.297}{196.332}{\Character{41}}%
\TextAt{9}{196.147}{176.298}{\Character{61}}%
\TextAt{8}{208.572}{176.298}{\char35\relax }%
\TextAt{13}{218.328}{176.298}{H}%
\TextAt{17}{228.95}{181.234}{0}%
\TextAt{9}{233.683}{176.298}{\Character{40}}%
\TextAt{13}{238.235}{176.298}{G\Character{59}}%
\TextAt{13}{252.689}{176.298}{B}%
\TextAt{9}{262.209}{176.298}{\Character{41}}%
\TextAt{16}{269.418}{176.298}{\Character{1}}%
\TextAt{8}{275.395}{176.298}{\char35\relax }%
\TextAt{13}{285.151}{176.298}{H}%
\TextAt{3}{295.773}{181.234}{\Character{0}}%
\TextAt{17}{302.36}{181.234}{1}%
\TextAt{9}{307.092}{176.298}{\Character{40}}%
\TextAt{13}{311.645}{176.298}{G\Character{59}}%
\TextAt{13}{326.099}{176.298}{A}%
\TextAt{9}{334.898}{176.298}{\Character{41}}%
\TextAt{16}{342.107}{176.298}{\Character{1}}%
\TextAt{8}{348.084}{176.298}{\char35\relax }%
\TextAt{13}{357.84}{176.298}{H}%
\TextAt{3}{368.462}{181.234}{\Character{0}}%
\TextAt{17}{375.049}{181.234}{1}%
\TextAt{9}{379.781}{176.298}{\Character{40}}%
\TextAt{13}{384.334}{176.298}{G\Character{59}}%
\TextAt{13}{398.788}{176.298}{C}%
\TextAt{9}{408.045}{176.298}{\Character{41}}%
\TextAt{13}{412.597}{176.298}{\Character{58}}%
\TextAt{8}{79.226}{152.1}{At}%
\TextAt{8}{96.165}{152.1}{the}%
\TextAt{8}{116.356}{152.1}{same}%
\TextAt{8}{145.716}{152.1}{time,}%
\TextAt{8}{175.671}{152.1}{we}%
\TextAt{8}{192.932}{152.1}{see}%
\TextAt{8}{211.88}{152.1}{that}%
\TextAt{8}{237.283}{152.1}{if}%
\TextAt{8}{248.041}{152.1}{any}%
\TextAt{8}{270.192}{152.1}{two}%
\TextAt{8}{292.343}{152.1}{of}%
\TextAt{8}{305.708}{152.1}{the}%
\TextAt{8}{325.898}{152.1}{quotients}%
\TextAt{8}{376.73}{152.1}{are}%
\TextAt{8}{396.275}{152.1}{well-defined,}%
\TextAt{8}{463.614}{152.1}{then}%
\TextAt{8}{490.308}{152.1}{so}%
\TextAt{8}{504.712}{152.1}{is}%
\TextAt{8}{516.51}{152.1}{the}%
\TextAt{8}{79.226}{135.054}{third.}%
\TextAt{8}{113.544}{135.054}{And}%
\TextAt{8}{139.391}{135.054}{from}%
\TextAt{8}{167.199}{135.054}{the}%
\TextAt{8}{187.522}{135.054}{last}%
\TextAt{8}{209.866}{135.054}{equation,}%
\TextAt{8}{261.135}{135.054}{we}%
\TextAt{8}{278.53}{135.054}{obtain}%
\TextAt{13}{315.124}{135.054}{h}%
\TextAt{9}{321.863}{135.054}{\Character{40}}%
\TextAt{13}{326.415}{135.054}{G\Character{59}}%
\TextAt{13}{340.869}{135.054}{B}%
\TextAt{9}{350.389}{135.054}{\Character{41}}%
\TextAt{9}{358.531}{135.054}{\Character{61}}%
\TextAt{13}{371.257}{135.054}{h}%
\TextAt{9}{377.996}{135.054}{\Character{40}}%
\TextAt{13}{382.548}{135.054}{G\Character{59}}%
\TextAt{13}{397.002}{135.054}{A}%
\TextAt{9}{405.801}{135.054}{\Character{41}}%
\TextAt{13}{410.353}{135.054}{h}%
\TextAt{9}{417.092}{135.054}{\Character{40}}%
\TextAt{13}{421.644}{135.054}{G\Character{59}}%
\TextAt{13}{436.098}{135.054}{C}%
\TextAt{9}{445.356}{135.054}{\Character{41}}%
\TextAt{8}{449.908}{135.054}{.}%
\TextAt{8}{458.857}{135.054}{Finally,}%
\TextAt{8}{501.172}{135.054}{if}%
\TextAt{13}{512.063}{135.054}{A}%
\TextAt{8}{524.903}{135.054}{is}%
\TextAt{8}{79.226}{118.008}{a}%
\TextAt{8}{88.982}{118.008}{finite}%
\TextAt{13}{118.893}{118.008}{G}%
\TextAt{8}{128.127}{118.008}{-module,}%
\TextAt{8}{176.569}{118.008}{then}%
\TextAt{8}{203.229}{118.008}{the}%
\TextAt{8}{223.386}{118.008}{exact}%
\TextAt{8}{254.278}{118.008}{sequences}%
\TextAt{9}{121.775}{93.81}{1}%
\NativeArrowAt{130.949}{93.81}{11.955}{19.26}{\longrightarrow}%
\TextAt{13}{153.531}{93.81}{A}%
\TextAt{21}{162.306}{98.746}{G}%
\NativeArrowAt{172.724}{93.81}{11.955}{19.26}{\longrightarrow}%
\TextAt{13}{195.306}{93.81}{A}%
\TextAt{21}{209.989}{100.587}{\Character{27}}%
\TextAt{3}{215.115}{100.587}{\Character{0}}%
\TextAt{17}{221.701}{100.587}{1}%
\NativeArrowAt{207.402}{93.81}{11.955}{22.524}{\longrightarrow}%
\TextAt{13}{233.255}{93.81}{I}%
\TextAt{21}{238.418}{92.017}{G}%
\TextAt{13}{245.515}{93.81}{A}%
\NativeArrowAt{257.611}{93.81}{11.955}{19.26}{\longrightarrow}%
\TextAt{9}{280.193}{93.81}{1}%
\TextAt{13}{286.046}{93.81}{\Character{59}}%
\TextAt{9}{315.2}{93.81}{1}%
\NativeArrowAt{324.374}{93.81}{11.955}{19.26}{\longrightarrow}%
\TextAt{21}{346.956}{92.017}{N}%
\TextAt{22}{353.678}{90.677}{G}%
\TextAt{13}{360.37}{93.81}{A}%
\NativeArrowAt{372.466}{93.81}{11.955}{19.26}{\longrightarrow}%
\TextAt{13}{395.048}{93.81}{A}%
\TextAt{21}{409.732}{100.864}{N}%
\TextAt{22}{416.453}{99.524}{G}%
\NativeArrowAt{407.144}{93.81}{11.955}{19.487}{\longrightarrow}%
\TextAt{13}{429.967}{93.81}{N}%
\TextAt{21}{439.341}{92.017}{G}%
\TextAt{13}{446.438}{93.81}{A}%
\NativeArrowAt{458.534}{93.81}{11.955}{19.26}{\longrightarrow}%
\TextAt{9}{481.116}{93.81}{1}%
\TextAt{13}{486.969}{93.81}{\Character{59}}%
%
%
%
%
%
%
%
%
%
%
%
%
%
%
%
%
%
%
%
%
%
%
%
\TextAt{8}{297.211}{43.033}{287}%
%
\Destination{78.23}{809.815}{nav.page.287}%
\Destination{79.36}{470.555}{nav.print.312}%
\Destination{79.36}{550.655}{nav.number.4.7.2}%
\Destination{79.36}{444.855}{nav.number.4.7.3}%
%
\Hotlink{279.883}{303.855}{25.078}{11.047}{nav.number.4.7.1}%
\Destination{79.226}{670}{errata.201}%
}

% Page 300 (body page 288)
\Page{%
\TextAt{2}{79.226}{796.943}{IV.}%
\TextAt{2}{98.012}{796.943}{Abstract}%
\TextAt{2}{143.252}{796.943}{Class}%
\TextAt{2}{171.844}{796.943}{Field}%
\TextAt{2}{199.575}{796.943}{Theory}%
\Vector{q 0 G 0 g 0.299 w 7.226 18.729 m 460.769 18.729 l S Q}%
\TextAt{8}{79.226}{761.124}{show}%
\TextAt{8}{108.229}{761.124}{that}%
\TextAt{8}{133.598}{761.124}{\char35\relax }%
\TextAt{13}{143.354}{761.124}{A}%
\TextAt{9}{155.45}{761.124}{\Character{61}}%
\TextAt{8}{167.875}{761.124}{\char35\relax }%
\TextAt{13}{177.631}{761.124}{A}%
\TextAt{21}{186.406}{765.462}{G}%
\TextAt{16}{196.16}{761.124}{\Character{1}}%
\TextAt{8}{202.137}{761.124}{\char35\relax }%
\TextAt{13}{211.893}{761.124}{I}%
\TextAt{21}{217.055}{759.33}{G}%
\TextAt{13}{224.152}{761.124}{A}%
\TextAt{9}{236.248}{761.124}{\Character{61}}%
\TextAt{8}{248.674}{761.124}{\char35\relax }%
\TextAt{21}{258.429}{759.33}{N}%
\TextAt{22}{265.151}{757.991}{G}%
\TextAt{13}{271.843}{761.124}{A}%
\TextAt{16}{283.275}{761.124}{\Character{1}}%
\TextAt{8}{289.253}{761.124}{\char35\relax }%
\TextAt{13}{299.008}{761.124}{N}%
\TextAt{21}{308.382}{759.33}{G}%
\TextAt{13}{315.479}{761.124}{A}%
\TextAt{8}{324.254}{761.124}{,}%
\TextAt{8}{331.403}{761.124}{and}%
\TextAt{13}{354.166}{761.124}{h}%
\TextAt{9}{360.904}{761.124}{\Character{40}}%
\TextAt{13}{365.457}{761.124}{G\Character{59}}%
\TextAt{13}{379.911}{761.124}{A}%
\TextAt{9}{388.71}{761.124}{\Character{41}}%
\TextAt{9}{396.565}{761.124}{\Character{61}}%
\TextAt{9}{408.986}{761.124}{1}%
\TextAt{8}{414.862}{761.124}{.}%
\Vector{q 0 G 0 g 0.399 w 452.334 -12.876 m 452.334 -4.807 l S Q}%
\Vector{q 0 G 0 g 0.399 w 452.534 -5.006 m 459.707 -5.006 l S Q}%
\Vector{q 0 G 0 g 0.399 w 452.534 -12.677 m 459.707 -12.677 l S Q}%
\Vector{q 0 G 0 g 0.399 w 459.906 -12.876 m 459.906 -4.807 l S Q}%
\TextAt{8}{97.159}{732.123}{If}%
\TextAt{13}{107.913}{732.123}{G}%
\TextAt{8}{120.106}{732.123}{is}%
\TextAt{8}{130.933}{732.123}{an}%
\TextAt{8}{146.254}{732.123}{arbitrary}%
\TextAt{8}{195.085}{732.123}{group}%
\TextAt{8}{227.323}{732.123}{and}%
\TextAt{13}{249.148}{732.123}{g}%
\TextAt{8}{258.142}{732.123}{a}%
\TextAt{8}{266.959}{732.123}{subgroup,}%
\TextAt{8}{320.259}{732.123}{then}%
\TextAt{8}{345.981}{732.123}{to}%
\TextAt{8}{359.353}{732.123}{any}%
\TextAt{13}{380.533}{732.123}{g}%
\TextAt{8}{386.567}{732.123}{-module}%
\TextAt{13}{430.82}{732.123}{B}%
\TextAt{8}{440.316}{732.123}{,}%
\TextAt{8}{446.715}{732.123}{we}%
\TextAt{8}{463.004}{732.123}{may}%
\TextAt{8}{487.435}{732.123}{associate}%
\TextAt{8}{79.226}{715.077}{the}%
\TextAt{8}{99.383}{715.077}{so-called}%
\TextAt{12}{146.916}{715.077}{induced}%
\TextAt{13}{197.14}{715.077}{G}%
\TextAt{12}{206.374}{715.077}{-module}%
\TextAt{13}{270.649}{686.075}{A}%
\TextAt{9}{282.745}{686.075}{\Character{61}}%
\TextAt{8}{295.17}{686.075}{Ind}%
\TextAt{21}{312.398}{691.743}{g}%
\TextAt{21}{312.398}{682.834}{G}%
\TextAt{9}{319.494}{686.075}{\Character{40}}%
\TextAt{13}{324.047}{686.075}{B}%
\TextAt{9}{333.543}{686.075}{\Character{41}}%
\TextAt{13}{338.095}{686.075}{\Character{58}}%
\TextAt{8}{79.226}{657.074}{It}%
\TextAt{8}{90.987}{657.074}{consists}%
\TextAt{8}{133.186}{657.074}{of}%
\TextAt{8}{145.604}{657.074}{all}%
\TextAt{8}{160.951}{657.074}{functions}%
\TextAt{13}{210.502}{657.074}{f}%
\TextAt{9}{220.87}{657.074}{\Character{58}}%
\TextAt{13}{227.442}{657.074}{G}%
\NativeMathAt{239.997}{657.074}{11.955}{\rightarrow}%
\TextAt{13}{255.273}{657.074}{B}%
\TextAt{8}{267.755}{657.074}{such}%
\TextAt{8}{293.24}{657.074}{that}%
\TextAt{13}{317.697}{657.074}{f}%
\TextAt{9}{324.744}{657.074}{\Character{40}}%
\TextAt{13}{329.296}{657.074}{x\Character{28}}%
\TextAt{9}{342.367}{657.074}{\Character{41}}%
\TextAt{9}{350.222}{657.074}{\Character{61}}%
\TextAt{13}{362.666}{657.074}{f}%
\TextAt{9}{369.712}{657.074}{\Character{40}}%
\TextAt{13}{374.264}{657.074}{x}%
\TextAt{9}{380.917}{657.074}{\Character{41}}%
\TextAt{21}{385.469}{661.413}{\Character{28}}%
\TextAt{8}{393.631}{657.074}{for}%
\TextAt{8}{410.604}{657.074}{all}%
\TextAt{13}{425.952}{657.074}{\Character{28}}%
\NativeMathAt{435.691}{657.074}{11.955}{\in}%
\TextAt{13}{446.982}{657.074}{g}%
\TextAt{8}{453.017}{657.074}{.}%
\TextAt{8}{461.161}{657.074}{The}%
\TextAt{8}{484.303}{657.074}{operation}%
\TextAt{8}{79.226}{640.028}{of}%
\TextAt{13}{92.556}{640.028}{\Character{27}}%
\NativeMathAt{102.959}{640.028}{11.955}{\in}%
\TextAt{13}{114.25}{640.028}{G}%
\TextAt{8}{127.381}{640.028}{is}%
\TextAt{8}{139.145}{640.028}{given}%
\TextAt{8}{169.715}{640.028}{by}%
\TextAt{13}{267.343}{622.982}{f}%
\TextAt{21}{274.389}{627.919}{\Character{27}}%
\TextAt{9}{280.013}{622.982}{\Character{40}}%
\TextAt{13}{284.565}{622.982}{x}%
\TextAt{9}{291.217}{622.982}{\Character{41}}%
\TextAt{9}{299.072}{622.982}{\Character{61}}%
\TextAt{13}{311.516}{622.982}{f}%
\TextAt{9}{318.562}{622.982}{\Character{40}}%
\TextAt{13}{323.114}{622.982}{\Character{27}}%
\TextAt{13}{330.18}{622.982}{x}%
\TextAt{9}{336.849}{622.982}{\Character{41}}%
\TextAt{13}{341.401}{622.982}{\Character{58}}%
\TextAt{8}{79.226}{599.461}{If}%
\TextAt{13}{90.836}{599.461}{g}%
\TextAt{9}{100.191}{599.461}{\Character{61}}%
\TextAt{16}{112.616}{599.461}{\Character{102}}%
\TextAt{9}{118.594}{599.461}{1}%
\TextAt{16}{124.447}{599.461}{\Character{103}}%
\TextAt{8}{130.425}{599.461}{,}%
\TextAt{8}{137.508}{599.461}{we}%
\TextAt{8}{154.653}{599.461}{write}%
\TextAt{8}{184.482}{599.461}{Ind}%
\TextAt{21}{201.709}{597.667}{G}%
\TextAt{9}{208.806}{599.461}{\Character{40}}%
\TextAt{13}{213.358}{599.461}{B}%
\TextAt{9}{222.855}{599.461}{\Character{41}}%
\TextAt{8}{231.222}{599.461}{instead}%
\TextAt{8}{271.524}{599.461}{of}%
\TextAt{8}{284.771}{599.461}{Ind}%
\TextAt{21}{301.999}{605.128}{g}%
\TextAt{21}{301.999}{596.22}{G}%
\TextAt{9}{309.096}{599.461}{\Character{40}}%
\TextAt{13}{313.648}{599.461}{B}%
\TextAt{9}{323.144}{599.461}{\Character{41}}%
\TextAt{8}{327.697}{599.461}{.}%
\TextAt{8}{336.117}{599.461}{We}%
\TextAt{8}{356.179}{599.461}{have}%
\TextAt{8}{383.092}{599.461}{a}%
\TextAt{8}{392.764}{599.461}{canonical}%
\TextAt{13}{444.065}{599.461}{g}%
\TextAt{8}{450.1}{599.461}{-homomorphism}%
\TextAt{13}{214.675}{570.46}{\Character{25}}%
\TextAt{9}{225.065}{570.46}{\Character{58}}%
\TextAt{8}{231.637}{570.46}{Ind}%
\TextAt{21}{248.865}{576.127}{g}%
\TextAt{21}{248.865}{567.218}{G}%
\TextAt{9}{255.961}{570.46}{\Character{40}}%
\TextAt{13}{260.514}{570.46}{B}%
\TextAt{9}{270.01}{570.46}{\Character{41}}%
\NativeArrowAt{277.883}{570.46}{11.955}{19.26}{\longrightarrow}%
\TextAt{13}{300.465}{570.46}{B}%
\TextAt{13}{309.946}{570.46}{\Character{59}}%
\TextAt{13}{339.092}{570.46}{f}%
\NativeArrowAt{349.483}{570.46}{11.955}{19.26}{\longmapsto}%
\TextAt{13}{372.065}{570.46}{f}%
\TextAt{9}{379.112}{570.46}{\Character{40}1\Character{41}}%
\TextAt{13}{394.069}{570.46}{\Character{59}}%
\TextAt{8}{79.226}{541.458}{which}%
\TextAt{8}{112.713}{541.458}{maps}%
\TextAt{8}{143.342}{541.458}{the}%
\TextAt{13}{163.498}{541.458}{g}%
\TextAt{8}{169.532}{541.458}{-submodule}%
\TextAt{13}{201.987}{512.457}{B}%
\TextAt{3}{211.483}{517.393}{\Character{48}}%
\TextAt{9}{217.599}{512.457}{\Character{61}}%
\TextAt{16}{230.025}{512.457}{\Character{102}}%
\TextAt{13}{236.002}{512.457}{f}%
\NativeMathAt{246.37}{512.457}{11.955}{\in}%
\TextAt{8}{257.661}{512.457}{Ind}%
\TextAt{21}{274.888}{518.125}{g}%
\TextAt{21}{274.888}{509.216}{G}%
\TextAt{9}{281.985}{512.457}{\Character{40}}%
\TextAt{13}{286.537}{512.457}{B}%
\TextAt{9}{296.033}{512.457}{\Character{41}}%
\TextAt{16}{303.906}{512.457}{\Character{106}}%
\TextAt{13}{310.548}{512.457}{f}%
\TextAt{9}{317.595}{512.457}{\Character{40}}%
\TextAt{13}{322.147}{512.457}{x}%
\TextAt{9}{328.799}{512.457}{\Character{41}}%
\TextAt{9}{336.654}{512.457}{\Character{61}}%
\TextAt{9}{349.075}{512.457}{1}%
\TextAt{8}{358.848}{512.457}{for}%
\TextAt{13}{376.733}{512.457}{x}%
\NativeMathAt{386.706}{512.457}{11.955}{\notin}%
\TextAt{13}{397.997}{512.457}{g}%
\TextAt{16}{404.031}{512.457}{\Character{103}}%
\TextAt{8}{79.226}{483.456}{isomorphically}%
\TextAt{8}{157.018}{483.456}{onto}%
\TextAt{13}{183.368}{483.456}{B}%
\TextAt{8}{192.864}{483.456}{.}%
\TextAt{8}{201.312}{483.456}{We}%
\TextAt{8}{221.457}{483.456}{identify}%
\TextAt{13}{264.053}{483.456}{B}%
\TextAt{3}{273.549}{487.795}{\Character{48}}%
\TextAt{8}{280.242}{483.456}{with}%
\TextAt{13}{306.902}{483.456}{B}%
\TextAt{8}{316.398}{483.456}{.}%
\TextAt{8}{324.846}{483.456}{If}%
\TextAt{13}{336.538}{483.456}{g}%
\TextAt{8}{346.47}{483.456}{is}%
\TextAt{8}{358.234}{483.456}{of}%
\TextAt{8}{371.564}{483.456}{finite}%
\TextAt{8}{401.476}{483.456}{index,}%
\TextAt{8}{436.265}{483.456}{we}%
\TextAt{8}{453.493}{483.456}{find}%
\ErrataAt{247}{454.455}{148}{$\displaystyle\operatorname{Ind}_G^g(B)=\prod_{\sigma\in R}B^\sigma,$}%
\ErrataAt{79.19}{414.804}{455.81}{where $R$ is a system of right coset representatives for $g\backslash G$. The supports are the left cosets}%
\ErrataAt{79.19}{397.758}{455.81}{$\sigma^{-1}g$ for $\sigma\in R$.}%
\TextAt{8}{97.159}{377.723}{Indeed,}%
\TextAt{8}{139.746}{377.723}{for}%
\TextAt{8}{158.675}{377.723}{any}%
\TextAt{13}{181.836}{377.723}{f}%
\NativeMathAt{193.983}{377.723}{11.955}{\in}%
\TextAt{8}{207.054}{377.723}{Ind}%
\TextAt{21}{224.281}{383.39}{g}%
\TextAt{21}{224.281}{374.482}{G}%
\TextAt{9}{231.378}{377.723}{\Character{40}}%
\TextAt{13}{235.931}{377.723}{B}%
\TextAt{9}{245.427}{377.723}{\Character{41}}%
\TextAt{8}{254.921}{377.723}{we}%
\TextAt{8}{273.193}{377.723}{have}%
\TextAt{8}{301.232}{377.723}{a}%
\TextAt{8}{312.032}{377.723}{unique}%
\TextAt{8}{351.117}{377.723}{factorization}%
\TextAt{13}{420.139}{377.723}{f}%
\TextAt{9}{432.286}{377.723}{\Character{61}}%
\TextAt{20}{446.491}{386.689}{\Character{81}}%
\TextAt{21}{457.782}{374.236}{\Character{27}}%
\TextAt{13}{465.398}{377.723}{f}%
\TextAt{21}{472.444}{382.061}{\Character{27}}%
\TextAt{21}{471.168}{374.767}{\Character{27}}%
\TextAt{8}{478.068}{377.723}{,}%
\TextAt{8}{486.522}{377.723}{where}%
\TextAt{13}{521.376}{377.723}{f}%
\TextAt{21}{527.146}{375.93}{\Character{27}}%
\TextAt{8}{79.226}{360.677}{denotes}%
\TextAt{8}{122.056}{360.677}{the}%
\TextAt{8}{142.71}{360.677}{function}%
\TextAt{8}{189.055}{360.677}{in}%
\TextAt{13}{203.205}{360.677}{B}%
\TextAt{3}{212.701}{365.015}{\Character{48}}%
\TextAt{8}{219.89}{360.677}{which}%
\TextAt{8}{253.874}{360.677}{is}%
\TextAt{8}{266.134}{360.677}{determined}%
\TextAt{8}{327.758}{360.677}{by}%
\TextAt{13}{344.514}{360.677}{f}%
\TextAt{21}{350.284}{358.884}{\Character{27}}%
\TextAt{9}{355.908}{360.677}{\Character{40}1\Character{41}}%
\TextAt{9}{375}{360.677}{\Character{61}}%
\TextAt{13}{388.305}{360.677}{f}%
\TextAt{9}{395.352}{360.677}{\Character{40}}%
\TextAt{13}{399.904}{360.677}{\Character{27}}%
\TextAt{3}{406.986}{365.015}{\Character{0}}%
\TextAt{17}{413.573}{365.015}{1}%
\TextAt{9}{418.305}{360.677}{\Character{41}}%
\TextAt{8}{422.858}{360.677}{.}%
\TextAt{8}{432.797}{360.677}{If}%
\TextAt{8}{444.986}{360.677}{conversely}%
\TextAt{13}{501.481}{360.677}{A}%
\TextAt{8}{514.651}{360.677}{is}%
\TextAt{8}{526.911}{360.677}{a}%
\TextAt{13}{79.226}{343.631}{G}%
\TextAt{8}{88.46}{343.631}{-module}%
\TextAt{8}{133.65}{343.631}{with}%
\TextAt{8}{160.31}{343.631}{a}%
\TextAt{13}{170.066}{343.631}{g}%
\TextAt{8}{176.1}{343.631}{-submodule}%
\TextAt{13}{238.912}{343.631}{B}%
\TextAt{8}{252.306}{343.631}{such}%
\TextAt{8}{278.703}{343.631}{that}%
\TextAt{13}{304.072}{343.631}{A}%
\TextAt{8}{316.745}{343.631}{is}%
\TextAt{8}{328.509}{343.631}{the}%
\TextAt{8}{348.665}{343.631}{direct}%
\TextAt{8}{381.829}{343.631}{product}%
\ErrataAt{255}{314.63}{130}{$\displaystyle A=\prod_{\sigma\in R}B^\sigma,$}%
\TextAt{8}{79.226}{274.846}{then}%
\TextAt{13}{105.886}{274.846}{A}%
\NativeMathAt{117.982}{274.845}{11.955}{\cong}%
\TextAt{8}{130.602}{274.846}{Ind}%
\TextAt{21}{147.829}{280.513}{g}%
\TextAt{21}{147.829}{271.604}{G}%
\TextAt{9}{154.926}{274.846}{\Character{40}}%
\TextAt{13}{159.478}{274.846}{B}%
\TextAt{9}{168.974}{274.846}{\Character{41}}%
\TextAt{8}{177.424}{274.846}{via}%
\TextAt{13}{196.612}{274.846}{B}%
\NativeMathAt{209.429}{274.845}{11.955}{\cong}%
\TextAt{13}{222.049}{274.846}{B}%
\TextAt{3}{231.545}{279.184}{\Character{48}}%
\TextAt{8}{234.34}{274.846}{.}%
\TextAt{12}{79.226}{248.833}{Proposition}%
\TextAt{12}{151.063}{248.833}{(7.4).}%
\TextAt{8}{185.618}{248.833}{Let}%
\TextAt{13}{205.828}{248.833}{G}%
\TextAt{8}{218.198}{248.833}{be}%
\TextAt{8}{233.363}{248.833}{a}%
\TextAt{8}{242.358}{248.833}{finite}%
\TextAt{8}{271.509}{248.833}{cyclic}%
\TextAt{8}{302.933}{248.833}{group,}%
\TextAt{13}{338.752}{248.833}{g}%
\TextAt{8}{347.923}{248.833}{a}%
\TextAt{8}{356.919}{248.833}{subgroup}%
\TextAt{8}{406.956}{248.833}{and}%
\TextAt{13}{428.958}{248.833}{B}%
\TextAt{8}{441.592}{248.833}{a}%
\TextAt{13}{450.587}{248.833}{g}%
\TextAt{8}{456.621}{248.833}{-module.}%
\TextAt{8}{506.109}{248.833}{Then}%
\TextAt{8}{79.226}{231.787}{we}%
\TextAt{8}{96.454}{231.787}{have}%
\TextAt{8}{123.448}{231.787}{canonically}%
\TextAt{13}{190.313}{202.786}{H}%
\TextAt{21}{200.935}{207.722}{i}%
\TextAt{9}{204.317}{202.786}{\Character{40}}%
\TextAt{13}{208.869}{202.786}{G\Character{59}}%
\TextAt{8}{223.347}{202.786}{Ind}%
\TextAt{21}{240.574}{208.454}{g}%
\TextAt{21}{240.574}{199.545}{G}%
\TextAt{9}{247.671}{202.786}{\Character{40}}%
\TextAt{13}{252.223}{202.786}{B}%
\TextAt{9}{261.72}{202.786}{\Character{41}\Character{41}}%
\NativeMathAt{274.145}{202.786}{11.955}{\cong}%
\TextAt{13}{286.764}{202.786}{H}%
\TextAt{21}{297.387}{207.722}{i}%
\TextAt{9}{300.768}{202.786}{\Character{40}}%
\TextAt{13}{305.321}{202.786}{g}%
\TextAt{13}{311.334}{202.786}{\Character{59}}%
\TextAt{13}{316.571}{202.786}{B}%
\TextAt{9}{326.095}{202.786}{\Character{41}}%
\TextAt{8}{354.558}{202.786}{for}%
\TextAt{13}{372.443}{202.786}{i}%
\TextAt{9}{379.757}{202.786}{\Character{61}}%
\TextAt{9}{392.167}{202.786}{0}%
\TextAt{13}{398.036}{202.786}{\Character{59}}%
\TextAt{16}{403.28}{202.786}{\Character{0}}%
\TextAt{9}{412.578}{202.786}{1}%
\TextAt{13}{418.431}{202.786}{\Character{58}}%
\TextAt{19}{79.226}{155.852}{Proof.}%
\TextAt{8}{116}{155.852}{Let}%
\TextAt{13}{136.324}{155.852}{A}%
\TextAt{9}{148.42}{155.852}{\Character{61}}%
\TextAt{8}{160.846}{155.852}{Ind}%
\TextAt{21}{178.073}{161.52}{g}%
\TextAt{21}{178.073}{152.611}{G}%
\TextAt{9}{185.17}{155.852}{\Character{40}}%
\TextAt{13}{189.722}{155.852}{B}%
\TextAt{9}{199.219}{155.852}{\Character{41}}%
\TextAt{8}{207.023}{155.852}{and}%
\TextAt{8}{229.14}{155.852}{let}%
\TextAt{13}{245.399}{155.852}{R}%
\TextAt{8}{257.66}{155.852}{be}%
\TextAt{8}{272.939}{155.852}{a}%
\TextAt{8}{282.049}{155.852}{system}%
\TextAt{8}{320.222}{155.852}{of}%
\TextAt{8}{332.906}{155.852}{right}%
\TextAt{8}{360.559}{155.852}{coset}%
\TextAt{8}{389.239}{155.852}{representatives}%
\TextAt{8}{468.395}{155.852}{for}%
\TextAt{13}{485.634}{155.852}{G}%
\TextAt{9}{494.868}{155.852}{\Character{47}}%
\TextAt{13}{500.721}{155.852}{g}%
\TextAt{8}{510.007}{155.852}{with}%
\TextAt{9}{79.226}{138.806}{1}%
\NativeMathAt{88.4}{138.806}{11.955}{\in}%
\TextAt{13}{99.691}{138.806}{R}%
\TextAt{8}{108.699}{138.806}{.}%
\TextAt{8}{117.148}{138.806}{We}%
\TextAt{8}{137.292}{138.806}{consider}%
\TextAt{8}{182.877}{138.806}{the}%
\TextAt{13}{203.034}{138.806}{g}%
\TextAt{8}{209.068}{138.806}{-homomorphisms}%
\TextAt{13}{147.477}{109.805}{\Character{25}}%
\TextAt{9}{157.867}{109.805}{\Character{58}}%
\TextAt{13}{164.44}{109.805}{A}%
\NativeArrowAt{176.536}{109.805}{11.955}{19.26}{\longrightarrow}%
\TextAt{13}{199.118}{109.805}{B}%
\TextAt{13}{208.599}{109.805}{\Character{59}}%
\TextAt{13}{225.79}{109.805}{f}%
\NativeArrowAt{236.181}{109.805}{11.955}{19.26}{\longmapsto}%
\TextAt{13}{258.763}{109.805}{f}%
\TextAt{9}{265.809}{109.805}{\Character{40}1\Character{41}\Character{59}}%
\TextAt{13}{309.922}{109.805}{\Character{23}}%
\TextAt{9}{319.78}{109.805}{\Character{58}}%
\TextAt{13}{326.353}{109.805}{A}%
\NativeArrowAt{338.449}{109.805}{11.955}{19.26}{\longrightarrow}%
\TextAt{13}{361.031}{109.805}{B}%
\TextAt{13}{370.512}{109.805}{\Character{59}}%
\TextAt{13}{387.703}{109.805}{f}%
\NativeArrowAt{398.094}{109.805}{11.955}{19.26}{\longmapsto}%
\TextAt{20}{421.243}{121.163}{\Character{89}}%
\TextAt{21}{420.676}{95.791}{\Character{26}}%
\NativeMathAt{425.012}{95.791}{7.97}{\in}%
\TextAt{21}{430.658}{95.791}{R}%
\TextAt{13}{439.078}{109.805}{f}%
\TextAt{9}{446.124}{109.805}{\Character{40}}%
\TextAt{13}{450.677}{109.805}{\Character{26}}%
\TextAt{9}{456.714}{109.805}{\Character{41}}%
\TextAt{13}{461.267}{109.805}{\Character{58}}%
\TextAt{8}{297.211}{43.033}{288}%
\Destination{78.23}{809.815}{nav.page.288}%
\Destination{79.36}{672.055}{nav.print.313}%
\Destination{79.36}{262.855}{nav.number.4.7.4}%
\Destination{79.226}{471.455}{errata.114}%
}

% Page 301 (body page 289)
\Page{%
\TextAt{2}{79.226}{796.943}{IV.}%
\TextAt{2}{98.012}{796.943}{Abstract}%
\TextAt{2}{143.252}{796.943}{Class}%
\TextAt{2}{171.844}{796.943}{Field}%
\TextAt{2}{199.575}{796.943}{Theory}%
\Vector{q 0 G 0 g 0.299 w 7.226 18.729 m 460.769 18.729 l S Q}%
\TextAt{8}{79.226}{761.124}{Both}%
\TextAt{8}{108.325}{761.124}{admit}%
\TextAt{8}{142.146}{761.124}{the}%
\TextAt{13}{162.303}{761.124}{g}%
\TextAt{8}{168.337}{761.124}{-homomorphism}%
\TextAt{13}{183.656}{724.775}{s}%
\TextAt{9}{192.491}{724.775}{\Character{58}}%
\TextAt{13}{199.063}{724.775}{B}%
\NativeArrowAt{211.881}{724.775}{11.955}{19.26}{\longrightarrow}%
\TextAt{13}{234.463}{724.775}{A\Character{59}}%
\TextAt{13}{272.373}{724.775}{b}%
\NativeArrowAt{280.69}{724.775}{11.955}{19.26}{\longmapsto}%
\TextAt{13}{303.272}{724.775}{f}%
\TextAt{21}{309.042}{722.982}{b}%
\TextAt{9}{313.163}{724.775}{\Character{40}}%
\TextAt{13}{317.715}{724.775}{\Character{27}}%
\TextAt{9}{324.798}{724.775}{\Character{41}}%
\TextAt{9}{332.653}{724.775}{\Character{61}}%
\TextAt{20}{345.096}{749.284}{\Character{56}}%
\TextAt{20}{345.096}{738.524}{\Character{60}}%
\TextAt{20}{345.096}{717.004}{\Character{58}}%
\TextAt{13}{355.723}{733.901}{b}%
\TextAt{21}{360.7}{738.239}{\Character{27}}%
\TextAt{8}{378.279}{733.901}{for}%
\TextAt{13}{396.164}{733.901}{\Character{27}}%
\NativeMathAt{406.567}{733.901}{11.955}{\in}%
\TextAt{13}{417.858}{733.901}{g}%
\TextAt{13}{423.872}{733.901}{\Character{59}}%
\TextAt{9}{355.723}{713.446}{1}%
\TextAt{8}{378.279}{713.446}{for}%
\TextAt{13}{396.164}{713.446}{\Character{27}}%
\NativeMathAt{406.567}{713.446}{11.955}{\notin}%
\TextAt{13}{417.858}{713.446}{g}%
\TextAt{13}{423.872}{713.446}{\Character{59}}%
\TextAt{8}{79.226}{688.432}{as}%
\TextAt{8}{93.596}{688.432}{a}%
\TextAt{8}{103.352}{688.432}{section,}%
\TextAt{8}{145.685}{688.432}{i.e.,}%
\TextAt{13}{167.79}{688.432}{\Character{25}}%
\TextAt{16}{177.516}{688.432}{\Character{14}}%
\TextAt{13}{186.15}{688.432}{s}%
\TextAt{9}{194.985}{688.432}{\Character{61}}%
\TextAt{13}{207.411}{688.432}{\Character{23}}%
\TextAt{16}{216.605}{688.432}{\Character{14}}%
\TextAt{13}{225.239}{688.432}{s}%
\TextAt{9}{234.074}{688.432}{\Character{61}}%
\TextAt{8}{246.5}{688.432}{id,}%
\TextAt{8}{263.404}{688.432}{and}%
\TextAt{8}{286.167}{688.432}{we}%
\TextAt{8}{303.394}{688.432}{have}%
\TextAt{13}{263.4}{662.867}{\Character{25}}%
\TextAt{16}{273.126}{662.867}{\Character{14}}%
\TextAt{13}{281.761}{662.867}{N}%
\TextAt{21}{291.134}{661.074}{G}%
\TextAt{9}{301.552}{662.867}{\Character{61}}%
\TextAt{13}{313.977}{662.867}{N}%
\TextAt{21}{323.351}{661.074}{g}%
\TextAt{16}{330.822}{662.867}{\Character{14}}%
\TextAt{13}{339.456}{662.867}{\Character{23}\Character{59}}%
\TextAt{8}{79.226}{637.302}{because}%
\TextAt{8}{122.528}{637.302}{one}%
\TextAt{8}{143.987}{637.302}{finds}%
\TextAt{8}{172.01}{637.302}{that,}%
\TextAt{8}{200.631}{637.302}{for}%
\TextAt{13}{218.516}{637.302}{f}%
\NativeMathAt{228.883}{637.302}{11.955}{\in}%
\TextAt{13}{240.174}{637.302}{A}%
\TextAt{8}{248.949}{637.302}{,}%
\TextAt{9}{125.069}{604.264}{\Character{40}}%
\TextAt{13}{129.622}{604.263}{N}%
\TextAt{21}{138.995}{602.47}{G}%
\TextAt{13}{146.092}{604.264}{f}%
\TextAt{9}{153.139}{604.264}{\Character{41}\Character{40}1\Character{41}}%
\TextAt{9}{175.926}{604.264}{\Character{61}}%
\TextAt{20}{188.395}{615.621}{\Character{89}}%
\TextAt{21}{188.713}{590.515}{\Character{28}}%
\NativeMathAt{193.391}{590.515}{7.97}{\in}%
\TextAt{21}{199.037}{590.515}{g}%
\TextAt{20}{206.23}{615.621}{\Character{89}}%
\TextAt{21}{205.663}{590.249}{\Character{26}}%
\NativeMathAt{210}{590.249}{7.97}{\in}%
\TextAt{21}{215.646}{590.249}{R}%
\TextAt{13}{224.065}{604.264}{f}%
\TextAt{21}{231.112}{609.2}{\Character{26}\Character{28}}%
\TextAt{9}{240.625}{604.264}{\Character{40}1\Character{41}}%
\TextAt{9}{258.869}{604.264}{\Character{61}}%
\TextAt{20}{271.329}{615.621}{\Character{89}}%
\TextAt{21}{276.628}{590.515}{\Character{28}}%
\TextAt{20}{288.597}{615.621}{\Character{89}}%
\TextAt{21}{294.067}{590.515}{\Character{26}}%
\TextAt{13}{305.866}{604.264}{f}%
\TextAt{9}{312.912}{604.264}{\Character{40}}%
\TextAt{13}{317.465}{604.264}{\Character{26}\Character{28}}%
\TextAt{9}{329.921}{604.264}{\Character{41}}%
\TextAt{9}{337.776}{604.264}{\Character{61}}%
\TextAt{20}{350.22}{615.621}{\Character{89}}%
\TextAt{21}{355.519}{590.515}{\Character{28}}%
\TextAt{20}{367.489}{624.707}{\Character{32}}%
\TextAt{20}{376.953}{615.621}{\Character{89}}%
\TextAt{21}{382.423}{590.515}{\Character{26}}%
\TextAt{13}{394.222}{604.264}{f}%
\TextAt{9}{401.268}{604.264}{\Character{40}}%
\TextAt{13}{405.821}{604.264}{\Character{26}}%
\TextAt{9}{411.858}{604.264}{\Character{41}}%
\TextAt{20}{416.411}{624.707}{\Character{33}}%
\TextAt{21}{425.875}{622.03}{\Character{28}}%
\TextAt{9}{434.372}{604.264}{\Character{61}}%
\TextAt{13}{446.798}{604.264}{N}%
\TextAt{21}{456.172}{602.47}{g}%
\TextAt{9}{460.986}{604.264}{\Character{40}}%
\TextAt{13}{465.538}{604.264}{\Character{23}}%
\TextAt{13}{472.065}{604.264}{f}%
\TextAt{9}{479.122}{604.264}{\Character{41}}%
\TextAt{13}{483.675}{604.264}{\Character{58}}%
\TextAt{8}{79.226}{568.204}{If}%
\TextAt{13}{91.31}{568.204}{f}%
\NativeMathAt{102.345}{568.204}{11.955}{\in}%
\TextAt{13}{114.303}{568.204}{A}%
\TextAt{21}{123.078}{572.542}{G}%
\TextAt{8}{130.175}{568.204}{,}%
\TextAt{8}{137.814}{568.204}{then}%
\TextAt{13}{164.866}{568.204}{f}%
\TextAt{9}{171.912}{568.204}{\Character{40}}%
\TextAt{13}{176.464}{568.204}{\Character{27}}%
\TextAt{9}{183.547}{568.204}{\Character{41}}%
\TextAt{9}{192.071}{568.204}{\Character{61}}%
\TextAt{13}{205.18}{568.204}{f}%
\TextAt{9}{212.227}{568.204}{\Character{40}1\Character{41}}%
\TextAt{8}{231.474}{568.204}{for}%
\TextAt{8}{249.75}{568.204}{all}%
\TextAt{13}{266.401}{568.204}{\Character{27}}%
\NativeMathAt{277.471}{568.204}{11.955}{\in}%
\TextAt{13}{289.43}{568.204}{G}%
\TextAt{8}{298.663}{568.204}{,}%
\TextAt{8}{306.302}{568.204}{and}%
\TextAt{13}{329.457}{568.204}{f}%
\TextAt{9}{336.503}{568.204}{\Character{40}1\Character{41}}%
\TextAt{9}{355.416}{568.204}{\Character{61}}%
\TextAt{13}{368.542}{568.204}{f}%
\TextAt{9}{375.588}{568.204}{\Character{40}}%
\TextAt{13}{380.141}{568.204}{\Character{28}}%
\TextAt{9}{386.56}{568.204}{\Character{41}}%
\TextAt{9}{395.084}{568.204}{\Character{61}}%
\TextAt{13}{408.193}{568.204}{f}%
\TextAt{9}{415.239}{568.204}{\Character{40}1\Character{41}}%
\TextAt{21}{430.197}{572.542}{\Character{28}}%
\TextAt{8}{439.663}{568.204}{for}%
\TextAt{8}{457.939}{568.204}{all}%
\TextAt{13}{474.59}{568.204}{\Character{28}}%
\NativeMathAt{484.997}{568.204}{11.955}{\in}%
\TextAt{13}{496.956}{568.204}{g}%
\TextAt{8}{502.99}{568.204}{.}%
\TextAt{8}{512.613}{568.204}{The}%
\TextAt{8}{79.226}{551.158}{map}%
\TextAt{13}{105.241}{551.158}{\Character{25}}%
\TextAt{8}{116.207}{551.158}{therefore}%
\TextAt{8}{165.307}{551.158}{induces}%
\TextAt{8}{206.983}{551.158}{an}%
\TextAt{8}{223.242}{551.158}{isomorphism}%
\TextAt{13}{267.848}{525.593}{\Character{25}}%
\TextAt{9}{278.238}{525.593}{\Character{58}}%
\TextAt{13}{284.811}{525.593}{A}%
\TextAt{21}{293.586}{530.529}{G}%
\NativeArrowAt{304.004}{525.593}{11.955}{19.26}{\longrightarrow}%
\TextAt{13}{326.586}{525.593}{B}%
\TextAt{21}{336.082}{530.529}{g}%
\TextAt{13}{340.896}{525.593}{\Character{58}}%
\TextAt{8}{79.226}{500.028}{It}%
\TextAt{8}{91.765}{500.028}{sends}%
\TextAt{13}{122.966}{500.028}{N}%
\TextAt{21}{132.34}{498.234}{G}%
\TextAt{13}{139.436}{500.028}{A}%
\TextAt{8}{151.975}{500.028}{onto}%
\TextAt{13}{178.191}{500.028}{N}%
\TextAt{21}{187.565}{498.234}{g}%
\TextAt{13}{192.379}{500.028}{B}%
\TextAt{8}{201.875}{500.028}{,}%
\TextAt{8}{208.917}{500.028}{for}%
\TextAt{8}{226.668}{500.028}{one}%
\TextAt{8}{247.994}{500.028}{has}%
\TextAt{13}{268.734}{500.028}{\Character{25}}%
\TextAt{9}{275.803}{500.028}{\Character{40}}%
\TextAt{13}{280.356}{500.028}{N}%
\TextAt{21}{289.73}{498.234}{G}%
\TextAt{13}{296.826}{500.028}{A}%
\TextAt{9}{305.602}{500.028}{\Character{41}}%
\TextAt{9}{313.457}{500.028}{\Character{61}}%
\TextAt{13}{325.9}{500.028}{N}%
\TextAt{21}{335.274}{498.234}{g}%
\TextAt{9}{340.088}{500.028}{\Character{40}}%
\TextAt{13}{344.641}{500.028}{\Character{23}}%
\TextAt{13}{351.168}{500.028}{A}%
\TextAt{9}{359.954}{500.028}{\Character{41}}%
\NativeMathAt{367.827}{500.028}{11.955}{\subseteq}%
\TextAt{13}{380.446}{500.028}{N}%
\TextAt{21}{389.82}{498.234}{g}%
\TextAt{13}{394.634}{500.028}{B}%
\TextAt{8}{407.894}{500.028}{on}%
\TextAt{8}{424.019}{500.028}{the}%
\TextAt{8}{444.042}{500.028}{one}%
\TextAt{8}{465.368}{500.028}{hand,}%
\TextAt{8}{497.779}{500.028}{and}%
\TextAt{8}{520.408}{500.028}{on}%
\TextAt{8}{79.226}{482.982}{the}%
\TextAt{8}{99.383}{482.982}{other,}%
\TextAt{13}{133.204}{482.982}{N}%
\TextAt{21}{142.578}{481.188}{g}%
\TextAt{9}{147.392}{482.982}{\Character{40}}%
\TextAt{13}{151.944}{482.982}{B}%
\TextAt{9}{161.44}{482.982}{\Character{41}}%
\TextAt{9}{169.295}{482.982}{\Character{61}}%
\TextAt{13}{181.739}{482.982}{N}%
\TextAt{21}{191.113}{481.188}{g}%
\TextAt{9}{195.926}{482.982}{\Character{40}}%
\TextAt{13}{200.479}{482.982}{\Character{23}}%
\TextAt{13}{207.007}{482.982}{sB}%
\TextAt{9}{222.027}{482.982}{\Character{41}}%
\TextAt{9}{229.882}{482.982}{\Character{61}}%
\TextAt{13}{242.326}{482.982}{\Character{25}}%
\TextAt{9}{249.395}{482.982}{\Character{40}}%
\TextAt{13}{253.947}{482.982}{N}%
\TextAt{21}{263.321}{481.188}{G}%
\TextAt{9}{270.418}{482.982}{\Character{40}}%
\TextAt{13}{274.97}{482.982}{sB}%
\TextAt{9}{289.98}{482.982}{\Character{41}\Character{41}}%
\NativeMathAt{302.406}{482.982}{11.955}{\subseteq}%
\TextAt{13}{315.025}{482.982}{\Character{25}}%
\TextAt{9}{322.095}{482.982}{\Character{40}}%
\TextAt{13}{326.647}{482.982}{N}%
\TextAt{21}{336.021}{481.188}{G}%
\TextAt{13}{343.117}{482.982}{A}%
\TextAt{9}{351.893}{482.982}{\Character{41}}%
\TextAt{8}{356.445}{482.982}{.}%
\TextAt{8}{364.893}{482.982}{Therefore}%
\TextAt{13}{417.891}{482.982}{H}%
\TextAt{17}{428.513}{487.32}{0}%
\TextAt{9}{433.246}{482.982}{\Character{40}}%
\TextAt{13}{437.798}{482.982}{G\Character{59}}%
\TextAt{13}{452.252}{482.982}{A}%
\TextAt{9}{461.051}{482.982}{\Character{41}}%
\TextAt{9}{468.906}{482.982}{\Character{61}}%
\TextAt{13}{481.35}{482.982}{H}%
\TextAt{17}{491.972}{487.32}{0}%
\TextAt{9}{496.705}{482.982}{\Character{40}}%
\TextAt{13}{501.257}{482.982}{g}%
\TextAt{13}{507.271}{482.982}{\Character{59}}%
\TextAt{13}{512.507}{482.982}{B}%
\TextAt{9}{522.032}{482.982}{\Character{41}}%
\TextAt{8}{526.584}{482.982}{.}%
\TextAt{8}{97.159}{462.947}{As}%
\TextAt{13}{114.446}{462.947}{N}%
\TextAt{21}{123.82}{461.154}{g}%
\TextAt{16}{131.291}{462.947}{\Character{14}}%
\TextAt{13}{139.925}{462.947}{\Character{23}}%
\TextAt{9}{149.784}{462.947}{\Character{61}}%
\TextAt{13}{162.209}{462.947}{\Character{25}}%
\TextAt{16}{171.935}{462.947}{\Character{14}}%
\TextAt{13}{180.569}{462.947}{N}%
\TextAt{21}{189.943}{461.154}{G}%
\TextAt{8}{197.04}{462.947}{,}%
\TextAt{8}{204.189}{462.947}{the}%
\TextAt{13}{224.345}{462.947}{g}%
\TextAt{8}{230.38}{462.947}{-homomorphism}%
\TextAt{13}{316.947}{462.947}{\Character{23}}%
\TextAt{9}{326.806}{462.947}{\Character{58}}%
\TextAt{13}{333.378}{462.947}{A}%
\NativeMathAt{345.475}{462.947}{11.955}{\rightarrow}%
\TextAt{13}{360.751}{462.947}{B}%
\TextAt{8}{374.144}{462.947}{induces}%
\TextAt{8}{415.82}{462.947}{a}%
\TextAt{13}{425.575}{462.947}{g}%
\TextAt{8}{431.61}{462.947}{-homomorphism}%
\TextAt{13}{261.58}{437.382}{\Character{23}}%
\TextAt{9}{271.439}{437.382}{\Character{58}}%
\TextAt{21}{278.011}{435.589}{N}%
\TextAt{22}{284.732}{434.249}{G}%
\TextAt{13}{291.425}{437.382}{A}%
\NativeArrowAt{303.521}{437.382}{11.955}{19.26}{\longrightarrow}%
\TextAt{21}{326.103}{435.589}{N}%
\TextAt{22}{332.824}{434.592}{g}%
\TextAt{13}{337.668}{437.382}{B}%
\TextAt{13}{347.149}{437.382}{\Character{58}}%
\TextAt{8}{79.226}{411.817}{It}%
\TextAt{8}{93.963}{411.817}{is}%
\TextAt{8}{107.79}{411.817}{surjective}%
\TextAt{8}{162.266}{411.817}{since}%
\TextAt{13}{192.998}{411.817}{\Character{23}}%
\TextAt{16}{203.599}{411.817}{\Character{14}}%
\TextAt{13}{213.641}{411.817}{s}%
\TextAt{9}{225.993}{411.817}{\Character{61}}%
\TextAt{8}{241.935}{411.817}{id.}%
\TextAt{8}{266.331}{411.817}{We}%
\TextAt{8}{288.539}{411.817}{show}%
\TextAt{8}{319.606}{411.817}{that}%
\TextAt{13}{347.039}{411.817}{I}%
\TextAt{21}{352.202}{410.023}{G}%
\TextAt{13}{359.299}{411.817}{A}%
\TextAt{8}{374.035}{411.817}{is}%
\TextAt{8}{387.863}{411.817}{the}%
\TextAt{8}{410.083}{411.817}{preimage}%
\TextAt{8}{462.227}{411.817}{of}%
\TextAt{13}{477.621}{411.817}{I}%
\TextAt{21}{482.784}{410.023}{g}%
\TextAt{13}{487.598}{411.817}{B}%
\TextAt{8}{497.095}{411.817}{.}%
\TextAt{13}{511.735}{411.817}{I}%
\TextAt{21}{516.897}{410.023}{G}%
\TextAt{13}{523.994}{411.817}{A}%
\TextAt{8}{79.226}{394.771}{consists}%
\TextAt{8}{124.353}{394.771}{of}%
\TextAt{8}{139.7}{394.771}{all}%
\TextAt{8}{157.976}{394.771}{elements}%
\TextAt{13}{207.85}{394.771}{f}%
\TextAt{21}{214.896}{399.109}{\Character{27}}%
\TextAt{3}{220.021}{399.109}{\Character{0}}%
\TextAt{17}{226.608}{399.109}{1}%
\TextAt{8}{231.34}{394.771}{,}%
\TextAt{13}{241.01}{394.771}{f}%
\NativeMathAt{254.815}{394.771}{11.955}{\in}%
\TextAt{13}{269.543}{394.771}{A}%
\TextAt{8}{278.318}{394.771}{,}%
\TextAt{13}{287.988}{394.771}{\Character{27}}%
\NativeMathAt{301.829}{394.771}{11.955}{\in}%
\TextAt{13}{316.557}{394.771}{G}%
\TextAt{8}{325.79}{394.771}{.}%
\TextAt{8}{340.289}{394.771}{For}%
\TextAt{8}{363.275}{394.771}{if}%
\TextAt{13}{376.016}{394.771}{G}%
\TextAt{9}{392.008}{394.771}{\Character{61}}%
\TextAt{16}{407.87}{394.771}{\Character{104}}%
\TextAt{13}{412.519}{394.771}{\Character{27}}%
\TextAt{17}{419.173}{392.977}{0}%
\TextAt{16}{423.905}{394.771}{\Character{105}}%
\TextAt{8}{434.469}{394.771}{and}%
\TextAt{13}{459.248}{394.771}{\Character{27}}%
\TextAt{9}{473.088}{394.771}{\Character{61}}%
\TextAt{13}{488.951}{394.771}{\Character{27}}%
\TextAt{21}{496.033}{399.109}{i}%
\TextAt{17}{495.604}{391.815}{0}%
\TextAt{8}{500.337}{394.771}{,}%
\TextAt{8}{510.007}{394.771}{then}%
\TextAt{13}{79.226}{377.725}{f}%
\TextAt{21}{86.273}{382.063}{\Character{27}}%
\TextAt{3}{91.398}{382.063}{\Character{0}}%
\TextAt{17}{97.984}{382.063}{1}%
\TextAt{9}{106.038}{377.725}{\Character{61}}%
\TextAt{13}{118.463}{377.725}{f}%
\TextAt{17}{125.509}{382.063}{\Character{40}1\Character{43}}%
\TextAt{21}{139.623}{382.063}{\Character{27}}%
\TextAt{24}{144.463}{380.956}{0}%
\TextAt{17}{148.614}{382.063}{\Character{43}}%
\TextAt{3}{155.2}{382.063}{\Character{1}\Character{1}\Character{1}}%
\TextAt{17}{162.257}{382.063}{\Character{43}}%
\TextAt{21}{168.844}{382.063}{\Character{27}}%
\TextAt{22}{173.969}{385.639}{i}%
\TextAt{23}{176.633}{385.639}{\Character{0}}%
\TextAt{24}{182.389}{385.639}{1}%
\TextAt{24}{173.683}{379.522}{0}%
\TextAt{17}{186.54}{382.063}{\Character{41}\Character{40}}%
\TextAt{21}{193.126}{382.063}{\Character{27}}%
\TextAt{24}{197.966}{380.956}{0}%
\TextAt{3}{202.117}{382.063}{\Character{0}}%
\TextAt{17}{208.703}{382.063}{1\Character{41}}%
\NativeMathAt{220.05}{377.725}{11.955}{\in}%
\TextAt{13}{231.341}{377.725}{I}%
\TextAt{21}{236.503}{375.931}{G}%
\TextAt{13}{243.6}{377.725}{A}%
\TextAt{8}{252.376}{377.725}{.}%
\TextAt{8}{260.647}{377.725}{In}%
\TextAt{8}{274.736}{377.725}{the}%
\TextAt{8}{294.36}{377.725}{same}%
\TextAt{8}{323.154}{377.725}{way,}%
\TextAt{8}{348.743}{377.725}{one}%
\TextAt{8}{369.67}{377.725}{has}%
\TextAt{13}{390.012}{377.725}{I}%
\TextAt{21}{395.175}{375.931}{g}%
\TextAt{13}{399.989}{377.725}{B}%
\TextAt{9}{412.806}{377.725}{\Character{61}}%
\TextAt{16}{425.231}{377.725}{\Character{102}}%
\TextAt{13}{431.209}{377.725}{b}%
\TextAt{21}{436.186}{382.063}{\Character{28}}%
\TextAt{3}{440.865}{382.063}{\Character{0}}%
\TextAt{17}{447.451}{382.063}{1}%
\TextAt{16}{455.504}{377.725}{\Character{106}}%
\TextAt{13}{462.146}{377.725}{b}%
\NativeMathAt{470.444}{377.725}{11.955}{\in}%
\TextAt{13}{481.735}{377.725}{B}%
\TextAt{13}{491.216}{377.725}{\Character{59}}%
\TextAt{13}{496.452}{377.725}{\Character{28}}%
\NativeMathAt{506.215}{377.725}{11.955}{\in}%
\TextAt{13}{517.506}{377.725}{g}%
\TextAt{16}{523.54}{377.725}{\Character{103}}%
\TextAt{8}{529.518}{377.725}{.}%
\TextAt{8}{79.226}{360.679}{Writing}%
\TextAt{8}{122.145}{360.679}{now}%
\TextAt{13}{146.534}{360.679}{\Character{27}}%
\TextAt{13}{153.6}{360.679}{\Character{26}}%
\TextAt{9}{162.975}{360.679}{\Character{61}}%
\TextAt{13}{175.4}{360.679}{\Character{26}}%
\TextAt{3}{181.438}{365.017}{\Character{48}}%
\TextAt{13}{184.233}{360.679}{\Character{28}}%
\TextAt{21}{189.317}{358.885}{\Character{26}}%
\TextAt{8}{194.151}{360.679}{,}%
\TextAt{8}{201.301}{360.679}{with}%
\TextAt{13}{227.961}{360.679}{\Character{26}\Character{59}}%
\TextAt{13}{239.223}{360.679}{\Character{26}}%
\TextAt{3}{245.28}{365.017}{\Character{48}}%
\NativeMathAt{251.396}{360.679}{11.955}{\in}%
\TextAt{13}{262.687}{360.679}{R}%
\TextAt{8}{271.696}{360.679}{,}%
\TextAt{13}{278.845}{360.679}{\Character{28}}%
\TextAt{21}{283.928}{358.885}{\Character{26}}%
\NativeMathAt{292.084}{360.679}{11.955}{\in}%
\TextAt{13}{303.375}{360.679}{g}%
\TextAt{8}{309.409}{360.679}{,}%
\TextAt{8}{316.558}{360.679}{we}%
\TextAt{8}{333.786}{360.679}{obtain}%
\TextAt{13}{171.682}{330.766}{\Character{23}}%
\TextAt{9}{178.22}{330.766}{\Character{40}}%
\TextAt{13}{182.773}{330.766}{f}%
\TextAt{21}{189.819}{335.702}{\Character{27}}%
\TextAt{3}{194.944}{335.702}{\Character{0}}%
\TextAt{17}{201.531}{335.702}{1}%
\TextAt{9}{206.263}{330.766}{\Character{41}}%
\TextAt{9}{214.118}{330.766}{\Character{61}}%
\TextAt{20}{227.128}{342.124}{\Character{89}}%
\TextAt{21}{226.562}{316.752}{\Character{26}}%
\NativeMathAt{230.898}{316.752}{7.97}{\in}%
\TextAt{21}{236.544}{316.752}{R}%
\TextAt{13}{246.159}{338.854}{f}%
\TextAt{9}{253.206}{338.854}{\Character{40}}%
\TextAt{13}{257.758}{338.854}{\Character{27}}%
\TextAt{13}{264.824}{338.854}{\Character{26}}%
\TextAt{9}{270.878}{338.854}{\Character{41}}%
\Vector{q 0 G 0 g 0.478 w 174.159 -440.245 m 203.43 -440.245 l S Q}%
\TextAt{13}{249.7}{322.565}{f}%
\TextAt{9}{256.747}{322.565}{\Character{40}}%
\TextAt{13}{261.299}{322.565}{\Character{26}}%
\TextAt{9}{267.337}{322.565}{\Character{41}}%
\TextAt{9}{279.947}{330.766}{\Character{61}}%
\TextAt{20}{292.372}{342.124}{\Character{89}}%
\TextAt{21}{296.5}{316.608}{\Character{26}}%
\TextAt{23}{300.836}{318.877}{\Character{48}}%
\TextAt{13}{310.836}{338.854}{f}%
\TextAt{9}{317.883}{338.854}{\Character{40}}%
\TextAt{13}{322.435}{338.854}{\Character{26}}%
\TextAt{3}{328.473}{343.192}{\Character{48}}%
\TextAt{9}{331.268}{338.854}{\Character{41}}%
\TextAt{21}{335.82}{343.192}{\Character{28}}%
\TextAt{22}{339.578}{342.196}{\Character{26}}%
\Vector{q 0 G 0 g 0.478 w 238.836 -440.245 m 272.394 -440.245 l S Q}%
\TextAt{13}{315.123}{322.565}{f}%
\TextAt{9}{322.17}{322.565}{\Character{40}}%
\TextAt{13}{326.722}{322.565}{\Character{26}}%
\TextAt{3}{332.76}{326.019}{\Character{48}}%
\TextAt{9}{335.555}{322.565}{\Character{41}}%
\TextAt{9}{348.911}{330.766}{\Character{61}}%
\TextAt{20}{361.336}{342.124}{\Character{89}}%
\TextAt{21}{366.806}{317.018}{\Character{26}}%
\TextAt{13}{378.605}{330.766}{b}%
\TextAt{21}{383.582}{335.702}{\Character{28}}%
\TextAt{22}{387.34}{334.706}{\Character{26}}%
\TextAt{3}{391.658}{335.702}{\Character{0}}%
\TextAt{17}{398.245}{335.702}{1}%
\TextAt{21}{383.582}{327.811}{\Character{26}}%
\NativeMathAt{406.298}{330.766}{11.955}{\in}%
\TextAt{13}{417.589}{330.766}{I}%
\TextAt{21}{422.752}{328.973}{g}%
\TextAt{13}{427.566}{330.766}{B}%
\TextAt{13}{437.047}{330.766}{\Character{58}}%
\TextAt{8}{79.226}{294.873}{On}%
\TextAt{8}{99.898}{294.873}{the}%
\TextAt{8}{121.216}{294.873}{other}%
\TextAt{8}{152.946}{294.873}{hand,}%
\TextAt{8}{186.916}{294.873}{for}%
\TextAt{13}{205.962}{294.873}{b}%
\TextAt{21}{210.939}{299.211}{\Character{28}}%
\TextAt{3}{215.617}{299.211}{\Character{0}}%
\TextAt{17}{222.204}{299.211}{1}%
\NativeMathAt{232.236}{294.873}{11.955}{\in}%
\TextAt{13}{245.506}{294.873}{I}%
\TextAt{21}{250.668}{293.079}{g}%
\TextAt{13}{255.482}{294.873}{B}%
\TextAt{8}{264.979}{294.873}{,}%
\TextAt{8}{273.579}{294.873}{the}%
\TextAt{8}{294.897}{294.873}{function}%
\TextAt{13}{341.906}{294.873}{f}%
\TextAt{21}{348.953}{299.211}{\Character{28}}%
\TextAt{3}{353.631}{299.211}{\Character{0}}%
\TextAt{17}{360.217}{299.211}{1}%
\TextAt{8}{364.95}{294.873}{,}%
\TextAt{8}{373.55}{294.873}{with}%
\TextAt{13}{401.372}{294.873}{f}%
\TextAt{9}{413.718}{294.873}{\Character{61}}%
\TextAt{13}{428.122}{294.873}{sb}%
\TextAt{8}{438.613}{294.873}{,}%
\TextAt{8}{447.214}{294.873}{is}%
\TextAt{8}{460.139}{294.873}{a}%
\TextAt{8}{471.055}{294.873}{preimage}%
\TextAt{8}{522.297}{294.873}{as}%
\TextAt{13}{79.226}{277.827}{\Character{23}}%
\TextAt{9}{85.764}{277.827}{\Character{40}}%
\TextAt{13}{90.316}{277.827}{f}%
\TextAt{21}{97.363}{282.165}{\Character{28}}%
\TextAt{3}{102.041}{282.165}{\Character{0}}%
\TextAt{17}{108.628}{282.165}{1}%
\TextAt{9}{113.36}{277.827}{\Character{41}}%
\TextAt{9}{121.215}{277.827}{\Character{61}}%
\TextAt{13}{133.659}{277.827}{\Character{23}}%
\TextAt{13}{140.187}{277.827}{s}%
\TextAt{9}{145.711}{277.827}{\Character{40}}%
\TextAt{13}{150.263}{277.827}{b}%
\TextAt{9}{155.24}{277.827}{\Character{41}}%
\TextAt{21}{159.792}{282.165}{\Character{28}}%
\TextAt{3}{164.471}{282.165}{\Character{0}}%
\TextAt{17}{171.057}{282.165}{1}%
\TextAt{9}{179.11}{277.827}{\Character{61}}%
\TextAt{13}{191.536}{277.827}{b}%
\TextAt{21}{196.513}{282.165}{\Character{28}}%
\TextAt{3}{201.191}{282.165}{\Character{0}}%
\TextAt{17}{207.778}{282.165}{1}%
\TextAt{8}{212.51}{277.827}{.}%
\TextAt{8}{220.695}{277.827}{After}%
\TextAt{8}{250.46}{277.827}{this}%
\TextAt{8}{272.491}{277.827}{it}%
\TextAt{8}{283.403}{277.827}{remains}%
\TextAt{8}{326.248}{277.827}{to}%
\TextAt{8}{339.766}{277.827}{show}%
\TextAt{8}{367.978}{277.827}{ker}%
\TextAt{9}{383.914}{277.827}{\Character{40}}%
\TextAt{13}{388.466}{277.827}{\Character{23}}%
\TextAt{9}{395.004}{277.827}{\Character{41}}%
\NativeMathAt{402.878}{277.827}{11.955}{\subseteq}%
\TextAt{13}{415.497}{277.827}{I}%
\TextAt{21}{420.66}{276.033}{G}%
\TextAt{13}{427.757}{277.827}{A}%
\TextAt{8}{436.532}{277.827}{.}%
\TextAt{8}{444.716}{277.827}{Let}%
\TextAt{13}{464.894}{277.827}{G}%
\TextAt{9}{477.448}{277.827}{\Character{61}}%
\TextAt{16}{489.874}{277.827}{\Character{104}}%
\TextAt{13}{494.523}{277.827}{\Character{39}}%
\TextAt{16}{502.192}{277.827}{\Character{105}}%
\TextAt{8}{506.841}{277.827}{,}%
\TextAt{13}{513.357}{277.827}{n}%
\TextAt{9}{523.665}{277.827}{\Character{61}}%
\ErrataAt{79.19}{260.781}{455.81}{$(G:g)$, $R=\{1,\varphi,\ldots,\varphi^{n-1}\}$. If $\nu(f)=1$, define $h\in A$ on these representatives by}%
\ErrataAt{79.19}{243.735}{455.81}{$h(\varphi^k)=\prod_{i=0}^k f(\varphi^i)$ for $0\le k<n$. Then $h(\varphi^{n-1})=\nu(f)=1$, so equivariance gives}%
\ErrataAt{79.19}{226.689}{455.81}{$h(\varphi^{-1})=1$. For $0\le k<n$ the quotient $h(\varphi^k)/h(\varphi^{k-1})$ is $f(\varphi^k)$. Thus}%
\ErrataAt{79.19}{209.643}{455.81}{$f=h^{1-\varphi^{-1}}\in I_GA$. Therefore $H^{-1}(G,A)=H^{-1}(g,B)$.}%
\Vector{q 0 G 0 g 0.399 w 452.334 -564.357 m 452.334 -556.287 l S Q}%
\Vector{q 0 G 0 g 0.399 w 452.534 -556.487 m 459.707 -556.487 l S Q}%
\Vector{q 0 G 0 g 0.399 w 452.534 -564.158 m 459.707 -564.158 l S Q}%
\Vector{q 0 G 0 g 0.399 w 459.906 -564.357 m 459.906 -556.287 l S Q}%
\TextAt{12}{79.226}{176.627}{Exercise}%
\TextAt{12}{133.143}{176.627}{IV.7.1.}%
\TextAt{8}{177.521}{176.627}{Let}%
\TextAt{13}{199.113}{176.627}{f}%
\TextAt{13}{205.497}{176.627}{\Character{59}}%
\TextAt{13}{210.734}{176.627}{g}%
\TextAt{8}{221.308}{176.627}{be}%
\TextAt{8}{237.855}{176.627}{endomorphisms}%
\TextAt{8}{321.853}{176.627}{of}%
\TextAt{8}{335.806}{176.627}{an}%
\TextAt{8}{352.687}{176.627}{abelian}%
\TextAt{8}{393.958}{176.627}{group}%
\TextAt{13}{427.756}{176.627}{A}%
\TextAt{8}{441.051}{176.627}{such}%
\TextAt{8}{468.071}{176.627}{that}%
\TextAt{13}{494.062}{176.627}{f}%
\TextAt{16}{504.19}{176.627}{\Character{14}}%
\TextAt{13}{513.249}{176.627}{g}%
\TextAt{9}{523.665}{176.627}{\Character{61}}%
\TextAt{13}{79.226}{159.581}{g}%
\TextAt{16}{87.917}{159.581}{\Character{14}}%
\TextAt{13}{96.551}{159.581}{f}%
\TextAt{9}{106.919}{159.581}{\Character{61}}%
\TextAt{9}{119.329}{159.581}{0}%
\TextAt{8}{125.197}{159.581}{.}%
\TextAt{8}{133.645}{159.581}{Make}%
\TextAt{8}{165.183}{159.581}{sense}%
\TextAt{8}{195.215}{159.581}{of}%
\TextAt{8}{208.545}{159.581}{the}%
\TextAt{8}{228.701}{159.581}{following}%
\TextAt{8}{278.136}{159.581}{statement.}%
\TextAt{8}{337.059}{159.581}{The}%
\TextAt{8}{361.112}{159.581}{quotient}%
\TextAt{13}{247.069}{130.561}{q}%
\TextAt{21}{252.259}{128.768}{f}%
\TextAt{21}{256.73}{128.768}{\Character{59}g}%
\TextAt{9}{263.902}{130.561}{\Character{40}}%
\TextAt{13}{268.454}{130.561}{A}%
\TextAt{9}{277.23}{130.561}{\Character{41}}%
\TextAt{9}{285.085}{130.561}{\Character{61}}%
\TextAt{9}{298.724}{138.649}{\Character{40}}%
\TextAt{8}{303.276}{138.649}{ker}%
\TextAt{13}{321.205}{138.649}{f}%
\TextAt{9}{331.572}{138.649}{\Character{58}}%
\TextAt{8}{338.145}{138.649}{im}%
\TextAt{13}{353.144}{138.649}{g}%
\TextAt{9}{359.179}{138.649}{\Character{41}}%
\Vector{q 0 G 0 g 0.478 w 226.724 -640.45 m 291.731 -640.45 l S Q}%
\TextAt{9}{298.724}{122.36}{\Character{40}}%
\TextAt{8}{303.276}{122.36}{ker}%
\TextAt{13}{321.205}{122.36}{g}%
\TextAt{9}{330.56}{122.36}{\Character{58}}%
\TextAt{8}{337.132}{122.36}{im}%
\TextAt{13}{352.132}{122.36}{f}%
\TextAt{9}{359.179}{122.36}{\Character{41}}%
%
%
%
%
%
%
%
%
%
%
%
%
%
%
%
%
%
%
%
%
%
%
%
%
%
%
%
%
%
%
%
%
%
%
%
%
%
%
%
%
%
%
%
%
%
%
%
%
%
%
%
%
%
%
%
%
%
%
%
%
%
%
%
%
%
%
%
%
%
%
%
%
%
%
%
%
%
%
%
%
%
%
%
%
%
%
%
%
%
%
%
%
%
%
%
%
%
%
\TextAt{8}{297.211}{43.033}{289}%
%
\Destination{78.23}{809.815}{nav.page.289}%
\Destination{79.226}{190.655}{nav.ex.4.7.1}%
\Destination{79.36}{652.255}{nav.print.314}%
\Destination{79.226}{277.781}{errata.115}%
}

% Page 302 (body page 290)
\Page{%
\TextAt{2}{79.226}{796.943}{IV.}%
\TextAt{2}{98.012}{796.943}{Abstract}%
\TextAt{2}{143.252}{796.943}{Class}%
\TextAt{2}{171.844}{796.943}{Field}%
\TextAt{2}{199.575}{796.943}{Theory}%
\Vector{q 0 G 0 g 0.299 w 7.226 18.729 m 460.769 18.729 l S Q}%
\TextAt{8}{79.226}{761.124}{is}%
\TextAt{8}{90.99}{761.124}{multiplicative.}%
\TextAt{12}{79.226}{729.134}{Exercise}%
\TextAt{12}{133.977}{729.134}{IV.7.2.}%
\TextAt{8}{179.08}{729.134}{Let}%
\TextAt{13}{201.397}{729.134}{f}%
\TextAt{13}{207.781}{729.134}{\Character{59}}%
\TextAt{13}{213.018}{729.134}{g}%
\TextAt{8}{224.317}{729.134}{be}%
\TextAt{8}{241.589}{729.134}{two}%
\TextAt{8}{265.053}{729.134}{commuting}%
\TextAt{8}{327.217}{729.134}{endomorphisms}%
\TextAt{8}{411.94}{729.134}{of}%
\TextAt{8}{426.617}{729.134}{an}%
\TextAt{8}{444.224}{729.134}{abelian}%
\TextAt{8}{486.219}{729.134}{group}%
\TextAt{13}{520.742}{729.134}{A}%
\TextAt{8}{529.518}{729.134}{.}%
\TextAt{8}{79.226}{712.088}{Show}%
\TextAt{8}{110.118}{712.088}{that}%
\TextAt{13}{242.004}{695.042}{q}%
\TextAt{17}{247.194}{693.249}{0}%
\TextAt{21}{251.428}{693.249}{\Character{59}g}%
\TextAt{21}{258.083}{693.249}{f}%
\TextAt{9}{263.541}{695.042}{\Character{40}}%
\TextAt{13}{268.094}{695.042}{A}%
\TextAt{9}{276.869}{695.042}{\Character{41}}%
\TextAt{9}{284.724}{695.042}{\Character{61}}%
\TextAt{13}{297.168}{695.042}{q}%
\TextAt{17}{302.358}{693.249}{0}%
\TextAt{21}{306.592}{693.249}{\Character{59}g}%
\TextAt{9}{313.758}{695.042}{\Character{40}}%
\TextAt{13}{318.311}{695.042}{A}%
\TextAt{9}{327.086}{695.042}{\Character{41}}%
\TextAt{13}{331.638}{695.042}{q}%
\TextAt{17}{336.829}{693.249}{0}%
\TextAt{21}{341.063}{693.249}{\Character{59}f}%
\TextAt{9}{348.86}{695.042}{\Character{40}}%
\TextAt{13}{353.412}{695.042}{A}%
\TextAt{9}{362.188}{695.042}{\Character{41}}%
\TextAt{13}{366.74}{695.042}{\Character{59}}%
\TextAt{8}{79.226}{671.52}{provided}%
\TextAt{8}{127.358}{671.52}{all}%
\TextAt{8}{143.617}{671.52}{quotients}%
\TextAt{8}{194.414}{671.52}{are}%
\TextAt{8}{213.925}{671.52}{defined.}%
\TextAt{12}{79.226}{639.53}{Exercise}%
\TextAt{12}{133.274}{639.53}{IV.7.3.}%
\TextAt{8}{177.766}{639.53}{Let}%
\TextAt{13}{199.472}{639.53}{G}%
\TextAt{8}{213.339}{639.53}{be}%
\TextAt{8}{230}{639.53}{a}%
\TextAt{8}{240.492}{639.53}{cyclic}%
\TextAt{8}{273.412}{639.53}{group}%
\TextAt{8}{307.324}{639.53}{of}%
\TextAt{8}{321.39}{639.53}{prime}%
\TextAt{8}{355.29}{639.53}{order}%
\TextAt{13}{386.596}{639.53}{p}%
\TextAt{8}{392.471}{639.53}{,}%
\TextAt{8}{400.541}{639.53}{and}%
\TextAt{8}{424.04}{639.53}{let}%
\TextAt{13}{441.681}{639.53}{A}%
\TextAt{8}{455.09}{639.53}{be}%
\TextAt{8}{471.751}{639.53}{a}%
\TextAt{13}{482.243}{639.53}{G}%
\TextAt{8}{491.476}{639.53}{-module}%
\TextAt{8}{79.226}{622.484}{such}%
\TextAt{8}{105.623}{622.484}{that}%
\TextAt{13}{130.992}{622.484}{q}%
\TextAt{17}{136.182}{620.691}{0}%
\TextAt{21}{140.416}{620.691}{\Character{59}p}%
\TextAt{9}{147.53}{622.484}{\Character{40}}%
\TextAt{13}{152.082}{622.484}{A}%
\TextAt{9}{160.857}{622.484}{\Character{41}}%
\TextAt{8}{169.307}{622.484}{is}%
\TextAt{8}{181.071}{622.484}{defined.}%
\TextAt{8}{225.935}{622.484}{Show}%
\TextAt{8}{256.827}{622.484}{that}%
\TextAt{13}{225.887}{593.483}{h}%
\TextAt{9}{232.625}{593.483}{\Character{40}}%
\TextAt{13}{237.178}{593.483}{G\Character{59}}%
\TextAt{13}{251.632}{593.483}{A}%
\TextAt{9}{260.431}{593.483}{\Character{41}}%
\TextAt{21}{264.983}{598.419}{p}%
\TextAt{3}{269.246}{598.419}{\Character{0}}%
\TextAt{17}{275.833}{598.419}{1}%
\TextAt{9}{283.886}{593.483}{\Character{61}}%
\TextAt{13}{296.311}{593.483}{q}%
\TextAt{17}{301.501}{591.69}{0}%
\TextAt{21}{305.736}{591.69}{\Character{59}p}%
\TextAt{9}{312.849}{593.483}{\Character{40}}%
\TextAt{13}{317.401}{593.483}{A}%
\TextAt{21}{326.176}{598.419}{G}%
\TextAt{9}{333.273}{593.483}{\Character{41}}%
\TextAt{21}{337.826}{598.419}{p}%
\TextAt{9}{342.586}{593.483}{\Character{47}}%
\TextAt{13}{348.439}{593.483}{q}%
\TextAt{17}{353.63}{591.69}{0}%
\TextAt{21}{357.864}{591.69}{\Character{59}p}%
\TextAt{9}{364.977}{593.483}{\Character{40}}%
\TextAt{13}{369.529}{593.483}{A}%
\TextAt{9}{378.305}{593.483}{\Character{41}}%
\TextAt{13}{382.857}{593.483}{\Character{58}}%
\TextAt{12}{97.159}{561.493}{Hint:}%
\TextAt{8}{132.698}{561.493}{Use}%
\TextAt{8}{155.185}{561.493}{the}%
\TextAt{8}{175.342}{561.493}{exact}%
\TextAt{8}{206.234}{561.493}{sequence}%
\TextAt{9}{220.144}{532.492}{0}%
\NativeArrowAt{229.318}{532.492}{11.955}{19.26}{\longrightarrow}%
\TextAt{13}{251.9}{532.492}{A}%
\TextAt{21}{260.675}{537.428}{G}%
\NativeArrowAt{271.093}{532.492}{11.955}{19.26}{\longrightarrow}%
\TextAt{13}{293.675}{532.492}{A}%
\TextAt{21}{308.359}{539.269}{\Character{27}}%
\TextAt{3}{313.484}{539.269}{\Character{0}}%
\TextAt{17}{320.071}{539.269}{1}%
\NativeArrowAt{305.771}{532.492}{11.955}{22.524}{\longrightarrow}%
\TextAt{13}{331.625}{532.492}{A}%
\TextAt{21}{340.4}{537.428}{\Character{27}}%
\TextAt{3}{345.525}{537.428}{\Character{0}}%
\TextAt{17}{352.112}{537.428}{1}%
\NativeArrowAt{360.165}{532.492}{11.955}{19.26}{\longrightarrow}%
\TextAt{9}{382.747}{532.492}{0}%
\TextAt{13}{388.6}{532.492}{\Character{58}}%
\TextAt{8}{79.226}{503.491}{Let}%
\TextAt{13}{101.4}{503.491}{N}%
\TextAt{9}{117.397}{503.491}{\Character{61}}%
\TextAt{9}{131.863}{503.491}{1}%
\TextAt{9}{141.188}{503.491}{\Character{43}}%
\TextAt{13}{153.789}{503.491}{\Character{27}}%
\TextAt{9}{164.349}{503.491}{\Character{43}}%
\TextAt{16}{176.932}{503.491}{\Character{1}}%
\TextAt{16}{182.228}{503.491}{\Character{1}}%
\TextAt{16}{187.536}{503.491}{\Character{1}}%
\TextAt{9}{194.357}{503.491}{\Character{43}}%
\TextAt{13}{206.94}{503.491}{\Character{27}}%
\TextAt{21}{214.022}{507.829}{p}%
\TextAt{3}{218.285}{507.829}{\Character{0}}%
\TextAt{17}{224.872}{507.829}{1}%
\TextAt{8}{234.706}{503.491}{in}%
\TextAt{8}{249.564}{503.491}{the}%
\TextAt{8}{270.925}{503.491}{group}%
\TextAt{8}{305.305}{503.491}{ring}%
\TextAt{18}{330.576}{503.491}{Z}%
\TextAt{9}{338.546}{503.491}{\Character{91}}%
\TextAt{13}{341.798}{503.491}{G}%
\TextAt{9}{351.031}{503.491}{\Character{93}}%
\TextAt{8}{354.283}{503.491}{.}%
\TextAt{8}{366.346}{503.491}{Show}%
\TextAt{8}{398.443}{503.491}{that}%
\TextAt{8}{425.016}{503.491}{the}%
\TextAt{8}{446.377}{503.491}{ring}%
\TextAt{18}{471.648}{503.491}{Z}%
\TextAt{9}{479.618}{503.491}{\Character{91}}%
\TextAt{13}{482.87}{503.491}{G}%
\TextAt{9}{492.103}{503.491}{\Character{93}\Character{47}}%
\TextAt{18}{501.208}{503.491}{Z}%
\TextAt{13}{509.178}{503.491}{N}%
\TextAt{8}{524.903}{503.491}{is}%
\TextAt{8}{79.226}{486.445}{isomorphic}%
\TextAt{8}{138.476}{486.445}{to}%
\TextAt{18}{152.786}{486.445}{Z}%
\TextAt{9}{160.756}{486.445}{\Character{91}}%
\TextAt{13}{164.008}{486.445}{\Character{16}}%
\TextAt{9}{170.024}{486.445}{\Character{93}}%
\TextAt{8}{173.275}{486.445}{,}%
\TextAt{8}{180.425}{486.445}{for}%
\TextAt{13}{198.309}{486.445}{\Character{16}}%
\TextAt{8}{208.222}{486.445}{a}%
\TextAt{8}{217.978}{486.445}{primitive}%
\TextAt{13}{268.058}{486.445}{p}%
\TextAt{8}{273.933}{486.445}{-th}%
\TextAt{8}{292.787}{486.445}{root}%
\TextAt{8}{317.833}{486.445}{of}%
\TextAt{8}{331.163}{486.445}{unity,}%
\TextAt{8}{364.003}{486.445}{and}%
\TextAt{8}{386.766}{486.445}{that}%
\TextAt{8}{412.135}{486.445}{in}%
\TextAt{8}{425.788}{486.445}{this}%
\TextAt{8}{448.61}{486.445}{ring}%
\TextAt{8}{472.676}{486.445}{one}%
\TextAt{8}{494.136}{486.445}{has}%
\TextAt{13}{264.706}{457.444}{p}%
\TextAt{9}{273.901}{457.444}{\Character{61}}%
\TextAt{9}{286.31}{457.444}{\Character{40}}%
\TextAt{13}{290.879}{457.444}{\Character{27}}%
\TextAt{16}{300.618}{457.444}{\Character{0}}%
\TextAt{9}{312.574}{457.444}{1\Character{41}}%
\TextAt{21}{322.979}{462.38}{p}%
\TextAt{3}{327.242}{462.38}{\Character{0}}%
\TextAt{17}{333.828}{462.38}{1}%
\TextAt{13}{338.56}{457.444}{\Character{34}\Character{59}}%
\TextAt{8}{79.226}{428.443}{where}%
\TextAt{13}{113.035}{428.443}{\Character{34}}%
\TextAt{8}{122.411}{428.443}{is}%
\TextAt{8}{134.175}{428.443}{a}%
\TextAt{8}{143.93}{428.443}{unit}%
\TextAt{8}{168.641}{428.443}{in}%
\TextAt{18}{182.294}{428.443}{Z}%
\TextAt{9}{190.264}{428.443}{\Character{91}}%
\TextAt{13}{193.516}{428.443}{G}%
\TextAt{9}{202.75}{428.443}{\Character{93}\Character{47}}%
\TextAt{18}{211.854}{428.443}{Z}%
\TextAt{13}{219.824}{428.443}{N}%
\TextAt{8}{230.447}{428.443}{.}%
\TextAt{12}{79.226}{396.453}{Exercise}%
\TextAt{12}{133.702}{396.453}{IV.7.4.}%
\TextAt{8}{178.566}{396.453}{Let}%
\TextAt{13}{200.644}{396.453}{L}%
\TextAt{16}{213.819}{396.453}{\Character{106}}%
\TextAt{13}{222.35}{396.453}{K}%
\TextAt{8}{238.174}{396.453}{be}%
\TextAt{8}{255.207}{396.453}{a}%
\TextAt{8}{266.071}{396.453}{cyclic}%
\TextAt{8}{299.363}{396.453}{extension}%
\TextAt{8}{352.237}{396.453}{of}%
\TextAt{8}{366.676}{396.453}{prime}%
\TextAt{8}{400.948}{396.453}{degree.}%
\TextAt{8}{445.24}{396.453}{Using}%
\TextAt{8}{479.249}{396.453}{exercise}%
\TextAt{8}{523.66}{396.453}{3,}%
\TextAt{8}{79.226}{379.407}{compute}%
\TextAt{8}{127.78}{379.407}{the}%
\TextAt{8}{149.016}{379.407}{Herbrand}%
\TextAt{8}{202.448}{379.407}{quotient}%
\TextAt{8}{249.71}{379.407}{of}%
\TextAt{8}{264.12}{379.407}{the}%
\TextAt{8}{285.356}{379.407}{group}%
\TextAt{8}{319.612}{379.407}{of}%
\TextAt{8}{334.022}{379.407}{units}%
\TextAt{14}{364.428}{379.407}{o}%
\TextAt{3}{370.272}{383.745}{\Character{3}}%
\TextAt{21}{370.272}{376.386}{L}%
\TextAt{8}{381.506}{379.407}{of}%
\TextAt{13}{395.916}{379.407}{L}%
\TextAt{8}{403.881}{379.407}{,}%
\TextAt{8}{412.38}{379.407}{viewed}%
\TextAt{8}{451.824}{379.407}{as}%
\TextAt{8}{467.273}{379.407}{a}%
\TextAt{13}{478.109}{379.407}{G}%
\TextAt{9}{487.342}{379.407}{\Character{40}}%
\TextAt{13}{491.895}{379.407}{L}%
\TextAt{16}{505.02}{379.407}{\Character{106}}%
\TextAt{13}{513.502}{379.407}{K}%
\TextAt{9}{524.32}{379.407}{\Character{41}}%
\TextAt{8}{528.872}{379.407}{-}%
\TextAt{8}{79.226}{362.361}{module.}%
\TextAt{12}{79.226}{330.371}{Exercise}%
\TextAt{12}{132.427}{330.371}{IV.7.5.}%
\TextAt{8}{176.183}{330.371}{If}%
\TextAt{13}{187.875}{330.371}{G}%
\TextAt{8}{201.006}{330.371}{is}%
\TextAt{8}{212.77}{330.371}{a}%
\TextAt{8}{222.525}{330.371}{group,}%
\TextAt{13}{258.952}{330.371}{g}%
\TextAt{8}{268.884}{330.371}{a}%
\TextAt{8}{278.64}{330.371}{normal}%
\TextAt{8}{318.319}{330.371}{subgroup}%
\TextAt{8}{369.116}{330.371}{and}%
\TextAt{13}{391.879}{330.371}{A}%
\TextAt{8}{404.552}{330.371}{a}%
\TextAt{13}{414.307}{330.371}{g}%
\TextAt{8}{420.341}{330.371}{-module,}%
\TextAt{8}{468.784}{330.371}{then}%
\TextAt{13}{233.574}{301.37}{H}%
\TextAt{17}{244.197}{306.306}{1}%
\TextAt{9}{248.929}{301.37}{\Character{40}}%
\TextAt{13}{253.482}{301.37}{G\Character{59}}%
\TextAt{8}{267.959}{301.37}{Ind}%
\TextAt{21}{285.187}{307.037}{g}%
\TextAt{21}{285.187}{298.129}{G}%
\TextAt{9}{292.284}{301.37}{\Character{40}}%
\TextAt{13}{296.836}{301.37}{A}%
\TextAt{9}{305.611}{301.37}{\Character{41}\Character{41}}%
\NativeMathAt{318.037}{301.37}{11.955}{\cong}%
\TextAt{13}{330.656}{301.37}{H}%
\TextAt{17}{341.279}{306.306}{1}%
\TextAt{9}{346.011}{301.37}{\Character{40}}%
\TextAt{13}{350.563}{301.37}{g}%
\TextAt{13}{356.577}{301.37}{\Character{59}}%
\TextAt{13}{361.813}{301.37}{A}%
\TextAt{9}{370.617}{301.37}{\Character{41}}%
\TextAt{13}{375.17}{301.37}{\Character{58}}%
\TextAt{8}{297.211}{43.033}{290}%
\Destination{78.23}{809.815}{nav.page.290}%
\Destination{79.226}{743.155}{nav.ex.4.7.2}%
\Destination{79.226}{653.555}{nav.ex.4.7.3}%
\Destination{79.226}{410.455}{nav.ex.4.7.4}%
\Destination{79.226}{344.355}{nav.ex.4.7.5}%
\Destination{79.36}{654.555}{nav.print.315}%
\Hotlink{523.16}{394.255}{6.858}{11.047}{nav.ex.4.7.3}%
}


% Chapter 5
% Page 303 (body page 291)
\Page{%
\TextAt{1}{79.226}{673.054}{Chapter}%
\TextAt{1}{188.19}{673.054}{V}%
\TextAt{1}{79.226}{614.873}{Local}%
\TextAt{1}{153.711}{614.873}{Class}%
\TextAt{1}{226.486}{614.873}{Field}%
\TextAt{1}{297.03}{614.873}{Theory}%
\TextAt{15}{79.226}{546.17}{§}%
\TextAt{15}{93.785}{546.17}{1}%
\TextAt{15}{120.693}{546.17}{The}%
\TextAt{15}{160.202}{546.17}{Local}%
\TextAt{15}{211.935}{546.17}{Reciprocity}%
\TextAt{15}{315.796}{546.17}{Law}%
\TextAt{8}{79.226}{514.284}{The}%
\TextAt{8}{103.882}{514.284}{abstract}%
\TextAt{8}{150.082}{514.284}{class}%
\TextAt{8}{178.122}{514.284}{field}%
\TextAt{8}{204.081}{514.284}{theory}%
\TextAt{8}{241.434}{514.284}{that}%
\TextAt{8}{267.405}{514.284}{we}%
\TextAt{8}{285.235}{514.284}{have}%
\TextAt{8}{312.833}{514.284}{developed}%
\TextAt{8}{367.735}{514.284}{in}%
\TextAt{8}{381.991}{514.284}{the}%
\TextAt{8}{402.75}{514.284}{last}%
\TextAt{8}{425.529}{514.284}{chapter}%
\TextAt{8}{468.082}{514.284}{is}%
\TextAt{8}{480.449}{514.284}{now}%
\TextAt{8}{505.44}{514.284}{going}%
\TextAt{8}{79.226}{497.238}{to}%
\TextAt{8}{93.655}{497.238}{be}%
\TextAt{8}{109.698}{497.238}{applied}%
\TextAt{8}{150.787}{497.238}{to}%
\TextAt{8}{165.215}{497.238}{the}%
\TextAt{8}{185.49}{497.238}{case}%
\TextAt{8}{210.38}{497.238}{of}%
\TextAt{8}{223.828}{497.238}{a}%
\TextAt{19}{233.702}{497.238}{local}%
\TextAt{19}{260.131}{497.238}{field}%
\TextAt{8}{281.053}{497.238}{,}%
\TextAt{8}{288.35}{497.238}{i.e.,}%
\TextAt{8}{310.603}{497.238}{to}%
\TextAt{8}{325.032}{497.238}{a}%
\TextAt{8}{334.906}{497.238}{field}%
\TextAt{8}{360.381}{497.238}{which}%
\TextAt{8}{393.986}{497.238}{is}%
\TextAt{8}{405.869}{497.238}{complete}%
\TextAt{8}{455.41}{497.238}{with}%
\TextAt{8}{482.188}{497.238}{respect}%
\TextAt{8}{522.357}{497.238}{to}%
\TextAt{8}{79.226}{480.192}{a}%
\TextAt{8}{89.436}{480.192}{discrete}%
\TextAt{8}{132.869}{480.192}{valuation,}%
\TextAt{8}{187.76}{480.192}{and}%
\TextAt{8}{210.977}{480.192}{which}%
\TextAt{8}{244.918}{480.192}{has}%
\TextAt{8}{266.246}{480.192}{a}%
\TextAt{8}{276.455}{480.192}{finite}%
\TextAt{8}{306.821}{480.192}{residue}%
\TextAt{8}{347.002}{480.192}{class}%
\TextAt{8}{374.893}{480.192}{field.}%
\TextAt{8}{406.163}{480.192}{By}%
\TextAt{8}{424.981}{480.192}{chap.}%
\TextAt{8}{458.534}{480.192}{II,}%
\TextAt{8}{474.901}{480.192}{(5.2),}%
\TextAt{8}{506.695}{480.192}{these}%
\TextAt{8}{79.226}{463.147}{are}%
\TextAt{8}{98.618}{463.147}{precisely}%
\TextAt{8}{146.356}{463.147}{the}%
\TextAt{8}{166.393}{463.147}{finite}%
\TextAt{8}{196.186}{463.147}{extensions}%
\TextAt{13}{252.448}{463.147}{K}%
\TextAt{8}{267.044}{463.147}{of}%
\TextAt{8}{280.255}{463.147}{the}%
\TextAt{8}{300.292}{463.147}{fields}%
\TextAt{18}{330.145}{463.147}{Q}%
\TextAt{21}{339.443}{461.353}{p}%
\TextAt{8}{347.983}{463.147}{or}%
\TextAt{18}{362.174}{463.147}{F}%
\TextAt{21}{369.48}{461.353}{p}%
\TextAt{9}{374.241}{463.147}{\Character{40}\Character{40}}%
\TextAt{13}{383.345}{463.147}{t}%
\TextAt{9}{387.573}{463.147}{\Character{41}\Character{41}}%
\TextAt{8}{396.677}{463.147}{.}%
\TextAt{8}{405.086}{463.147}{We}%
\TextAt{8}{425.111}{463.147}{will}%
\TextAt{8}{447.098}{463.147}{use}%
\TextAt{8}{467.195}{463.147}{the}%
\TextAt{8}{487.232}{463.147}{following}%
\TextAt{8}{79.226}{446.101}{notation.}%
\TextAt{8}{130.617}{446.101}{Let}%
\TextAt{13}{79.226}{426.066}{v}%
\TextAt{21}{84.885}{424.273}{K}%
\TextAt{8}{127.415}{426.066}{be}%
\TextAt{8}{143.339}{426.066}{the}%
\TextAt{8}{163.496}{426.066}{discrete}%
\TextAt{8}{206.475}{426.066}{valuation}%
\TextAt{8}{257.547}{426.066}{normalized}%
\TextAt{8}{317.383}{426.066}{by}%
\TextAt{13}{333.642}{426.066}{v}%
\TextAt{21}{339.301}{424.273}{K}%
\TextAt{9}{347.504}{426.066}{\Character{40}}%
\TextAt{13}{352.056}{426.066}{K}%
\TextAt{3}{362.874}{430.404}{\Character{3}}%
\TextAt{9}{367.606}{426.066}{\Character{41}}%
\TextAt{9}{375.461}{426.066}{\Character{61}}%
\TextAt{18}{387.905}{426.066}{Z}%
\TextAt{8}{395.875}{426.066}{,}%
\TextAt{14}{79.226}{409.02}{o}%
\TextAt{21}{85.07}{407.227}{K}%
\TextAt{9}{127.415}{409.02}{\Character{61}}%
\TextAt{16}{139.841}{409.02}{\Character{102}}%
\TextAt{13}{145.818}{409.02}{a}%
\NativeMathAt{155.284}{409.02}{11.955}{\in}%
\TextAt{13}{166.575}{409.02}{K}%
\TextAt{16}{180.713}{409.02}{\Character{106}}%
\TextAt{13}{187.355}{409.02}{v}%
\TextAt{21}{193.014}{407.227}{K}%
\TextAt{9}{201.217}{409.02}{\Character{40}}%
\TextAt{13}{205.769}{409.02}{a}%
\TextAt{9}{211.914}{409.02}{\Character{41}}%
\NativeMathAt{219.788}{409.02}{11.955}{\geq}%
\TextAt{9}{232.407}{409.02}{0}%
\TextAt{16}{238.26}{409.02}{\Character{103}}%
\TextAt{8}{248.135}{409.02}{the}%
\TextAt{8}{268.291}{409.02}{valuation}%
\TextAt{8}{319.364}{409.02}{ring,}%
\TextAt{14}{79.226}{391.974}{p}%
\TextAt{21}{85.209}{390.181}{K}%
\TextAt{9}{127.415}{391.974}{\Character{61}}%
\TextAt{16}{139.841}{391.974}{\Character{102}}%
\TextAt{13}{145.818}{391.974}{a}%
\NativeMathAt{155.284}{391.974}{11.955}{\in}%
\TextAt{13}{166.575}{391.974}{K}%
\TextAt{16}{180.713}{391.974}{\Character{106}}%
\TextAt{13}{187.355}{391.974}{v}%
\TextAt{21}{193.014}{390.181}{K}%
\TextAt{9}{201.217}{391.974}{\Character{40}}%
\TextAt{13}{205.769}{391.974}{a}%
\TextAt{9}{211.914}{391.974}{\Character{41}}%
\TextAt{13}{219.788}{391.974}{\Character{62}}%
\TextAt{9}{232.213}{391.974}{0}%
\TextAt{16}{238.066}{391.974}{\Character{103}}%
\TextAt{8}{247.941}{391.974}{the}%
\TextAt{8}{268.097}{391.974}{maximal}%
\TextAt{8}{315.906}{391.974}{ideal,}%
\TextAt{13}{79.226}{374.928}{\Character{20}}%
\TextAt{9}{127.415}{374.928}{\Character{61}}%
\TextAt{14}{139.841}{374.928}{o}%
\TextAt{21}{145.685}{373.135}{K}%
\TextAt{9}{153.887}{374.928}{\Character{47}}%
\TextAt{14}{159.74}{374.928}{p}%
\TextAt{21}{165.724}{373.135}{K}%
\TextAt{8}{177.824}{374.928}{the}%
\TextAt{8}{197.98}{374.928}{residue}%
\TextAt{8}{237.707}{374.928}{class}%
\TextAt{8}{265.144}{374.928}{field,}%
\TextAt{13}{79.226}{357.882}{U}%
\TextAt{21}{87.176}{356.089}{K}%
\TextAt{9}{127.415}{357.882}{\Character{61}}%
\TextAt{16}{139.841}{357.882}{\Character{102}}%
\TextAt{13}{145.818}{357.882}{a}%
\NativeMathAt{155.284}{357.882}{11.955}{\in}%
\TextAt{13}{166.575}{357.882}{K}%
\TextAt{3}{177.393}{362.22}{\Character{3}}%
\TextAt{16}{185.446}{357.882}{\Character{106}}%
\TextAt{13}{192.087}{357.882}{v}%
\TextAt{21}{197.747}{356.089}{K}%
\TextAt{9}{205.949}{357.882}{\Character{40}}%
\TextAt{13}{210.502}{357.882}{a}%
\TextAt{9}{216.647}{357.882}{\Character{41}}%
\TextAt{9}{224.502}{357.882}{\Character{61}}%
\TextAt{9}{236.923}{357.882}{0}%
\TextAt{16}{242.798}{357.882}{\Character{103}}%
\TextAt{8}{252.673}{357.882}{the}%
\TextAt{8}{272.83}{357.882}{unit}%
\TextAt{8}{297.541}{357.882}{group,}%
\TextAt{13}{79.226}{340.836}{U}%
\TextAt{17}{88.425}{346.946}{\Character{40}}%
\TextAt{21}{91.718}{346.946}{n}%
\TextAt{17}{96.857}{346.946}{\Character{41}}%
\TextAt{21}{87.176}{337.595}{K}%
\TextAt{9}{127.415}{340.836}{\Character{61}}%
\TextAt{9}{139.825}{340.836}{1}%
\TextAt{9}{148.337}{340.836}{\Character{43}}%
\TextAt{14}{160.112}{340.836}{p}%
\TextAt{21}{166.095}{345.174}{n}%
\TextAt{21}{166.095}{337.815}{K}%
\TextAt{8}{178.195}{340.836}{the}%
\TextAt{8}{198.351}{340.836}{group}%
\TextAt{8}{231.527}{340.836}{of}%
\TextAt{13}{244.857}{340.836}{n}%
\TextAt{8}{251.844}{340.836}{-th}%
\TextAt{8}{270.698}{340.836}{higher}%
\TextAt{8}{306.468}{340.836}{units,}%
\TextAt{13}{339.045}{340.836}{n}%
\TextAt{9}{349.354}{340.836}{\Character{61}}%
\TextAt{9}{361.764}{340.836}{1}%
\TextAt{13}{367.632}{340.836}{\Character{59}}%
\TextAt{9}{372.876}{340.836}{2}%
\TextAt{13}{378.729}{340.836}{\Character{59}}%
\TextAt{13}{383.953}{340.836}{\Character{58}}%
\TextAt{13}{389.19}{340.836}{\Character{58}}%
\TextAt{13}{394.426}{340.836}{\Character{58}}%
\TextAt{8}{397.714}{340.836}{,}%
\TextAt{13}{79.226}{323.79}{q}%
\TextAt{9}{127.415}{323.79}{\Character{61}}%
\TextAt{13}{139.841}{323.79}{q}%
\TextAt{21}{145.031}{321.997}{K}%
\TextAt{9}{156.554}{323.79}{\Character{61}}%
\TextAt{8}{168.98}{323.79}{\char35\relax }%
\TextAt{13}{178.735}{323.79}{\Character{20}}%
\TextAt{8}{185.474}{323.79}{,}%
\TextAt{16}{79.226}{306.744}{\Character{106}}%
\TextAt{13}{82.547}{306.744}{a}%
\TextAt{16}{88.692}{306.744}{\Character{106}}%
\TextAt{27}{92.013}{304.951}{p}%
\TextAt{9}{127.415}{306.744}{\Character{61}}%
\TextAt{13}{139.841}{306.744}{q}%
\TextAt{3}{145.46}{311.082}{\Character{0}}%
\TextAt{21}{152.046}{311.082}{v}%
\TextAt{22}{156.17}{309.743}{K}%
\TextAt{17}{163.213}{311.082}{\Character{40}}%
\TextAt{21}{166.506}{311.082}{a}%
\TextAt{17}{171.004}{311.082}{\Character{41}}%
\TextAt{8}{178.693}{306.744}{the}%
\TextAt{8}{198.849}{306.744}{normalized}%
\TextAt{14}{258.685}{306.744}{p}%
\TextAt{8}{264.668}{306.744}{-adic}%
\TextAt{8}{293.277}{306.744}{absolute}%
\TextAt{8}{339.519}{306.744}{value,}%
\TextAt{13}{79.226}{289.698}{\Character{22}}%
\TextAt{21}{86.269}{287.905}{n}%
\TextAt{8}{127.415}{289.698}{the}%
\TextAt{8}{147.572}{289.698}{group}%
\TextAt{8}{180.747}{289.698}{of}%
\TextAt{13}{194.077}{289.698}{n}%
\TextAt{8}{201.065}{289.698}{-th}%
\TextAt{8}{219.918}{289.698}{roots}%
\TextAt{8}{249.579}{289.698}{of}%
\TextAt{8}{262.909}{289.698}{unity,}%
\TextAt{8}{295.75}{289.698}{and}%
\TextAt{13}{318.512}{289.698}{\Character{22}}%
\TextAt{21}{325.555}{287.905}{n}%
\TextAt{9}{331.192}{289.698}{\Character{40}}%
\TextAt{13}{335.744}{289.698}{K}%
\TextAt{9}{346.562}{289.698}{\Character{41}}%
\TextAt{9}{354.417}{289.698}{\Character{61}}%
\TextAt{13}{366.86}{289.698}{\Character{22}}%
\TextAt{21}{373.903}{287.905}{n}%
\TextAt{16}{382.196}{289.698}{\Character{92}}%
\TextAt{13}{392.823}{289.698}{K}%
\TextAt{3}{403.641}{294.036}{\Character{3}}%
\TextAt{8}{408.373}{289.698}{.}%
\TextAt{13}{79.226}{269.663}{\Character{25}}%
\TextAt{21}{85.866}{267.87}{K}%
\TextAt{8}{94.069}{269.663}{,}%
\TextAt{8}{101.218}{269.663}{or}%
\TextAt{8}{115.529}{269.663}{simply}%
\TextAt{13}{152.984}{269.663}{\Character{25}}%
\TextAt{8}{160.053}{269.663}{,}%
\TextAt{8}{167.203}{269.663}{denotes}%
\TextAt{8}{209.536}{269.663}{a}%
\TextAt{8}{219.291}{269.663}{prime}%
\TextAt{8}{252.455}{269.663}{element}%
\TextAt{8}{295.697}{269.663}{of}%
\TextAt{13}{309.027}{269.663}{K}%
\TextAt{8}{319.845}{269.663}{,}%
\TextAt{8}{326.994}{269.663}{i.e.,}%
\TextAt{14}{349.099}{269.663}{p}%
\TextAt{21}{355.082}{267.87}{K}%
\TextAt{9}{366.605}{269.663}{\Character{61}}%
\TextAt{13}{379.031}{269.663}{\Character{25}}%
\TextAt{14}{386.1}{269.663}{o}%
\TextAt{21}{391.944}{267.87}{K}%
\TextAt{8}{400.147}{269.663}{.}%
\TextAt{8}{97.159}{249.628}{In}%
\TextAt{8}{113.365}{249.628}{local}%
\TextAt{8}{142.59}{249.628}{class}%
\TextAt{8}{171.612}{249.628}{field}%
\TextAt{8}{198.554}{249.628}{theory,}%
\TextAt{8}{239.557}{249.628}{the}%
\TextAt{8}{261.298}{249.628}{rôle}%
\TextAt{8}{285.645}{249.628}{of}%
\TextAt{8}{300.56}{249.628}{the}%
\TextAt{8}{322.302}{249.628}{profinite}%
\TextAt{8}{370.715}{249.628}{group}%
\TextAt{13}{405.475}{249.628}{G}%
\TextAt{8}{420.191}{249.628}{of}%
\TextAt{8}{435.106}{249.628}{abstract}%
\TextAt{8}{482.288}{249.628}{class}%
\TextAt{8}{511.31}{249.628}{field}%
\TextAt{8}{79.226}{232.583}{theory}%
\TextAt{8}{116.4}{232.583}{is}%
\TextAt{8}{128.587}{232.583}{taken}%
\TextAt{8}{160.883}{232.583}{by}%
\TextAt{8}{177.565}{232.583}{the}%
\TextAt{8}{198.145}{232.583}{absolute}%
\TextAt{8}{244.811}{232.583}{Galois}%
\TextAt{8}{281.148}{232.583}{group}%
\TextAt{13}{314.746}{232.583}{G}%
\TextAt{9}{323.98}{232.583}{\Character{40}}%
\TextAt{9}{328.851}{235.737}{\Character{22}}%
\TextAt{13}{328.532}{232.583}{k}%
\TextAt{16}{335.022}{232.583}{\Character{106}}%
\TextAt{13}{338.343}{232.583}{k}%
\TextAt{9}{344.832}{232.583}{\Character{41}}%
\TextAt{8}{353.705}{232.583}{of}%
\TextAt{8}{367.459}{232.583}{a}%
\TextAt{8}{377.637}{232.583}{fixed}%
\TextAt{8}{406.347}{232.583}{local}%
\TextAt{8}{434.41}{232.583}{field}%
\TextAt{13}{460.191}{232.583}{k}%
\TextAt{8}{466.68}{232.583}{,}%
\TextAt{8}{474.359}{232.583}{and}%
\TextAt{8}{497.545}{232.583}{that}%
\TextAt{8}{523.337}{232.583}{of}%
\TextAt{8}{79.226}{215.537}{the}%
\TextAt{13}{100.397}{215.537}{G}%
\TextAt{8}{109.631}{215.537}{-module}%
\TextAt{13}{155.836}{215.537}{A}%
\TextAt{8}{169.523}{215.537}{by}%
\TextAt{8}{186.797}{215.537}{the}%
\TextAt{8}{207.968}{215.537}{multiplicative}%
\TextAt{8}{282.806}{215.537}{group}%
\TextAt{9}{317.314}{218.691}{\Character{22}}%
\TextAt{13}{316.996}{215.537}{k}%
\TextAt{3}{323.486}{219.875}{\Character{3}}%
\TextAt{8}{333.13}{215.537}{of}%
\TextAt{8}{347.475}{215.537}{the}%
\TextAt{8}{368.646}{215.537}{separable}%
\TextAt{8}{421.103}{215.537}{closure}%
\TextAt{9}{461.518}{218.691}{\Character{22}}%
\TextAt{13}{461.2}{215.537}{k}%
\TextAt{8}{472.601}{215.537}{of}%
\TextAt{13}{486.946}{215.537}{k}%
\TextAt{8}{493.435}{215.537}{.}%
\TextAt{8}{504.927}{215.537}{For}%
\TextAt{8}{526.911}{215.537}{a}%
\TextAt{8}{79.226}{198.491}{finite}%
\TextAt{8}{110.103}{198.491}{extension}%
\TextAt{13}{162.833}{198.491}{K}%
\TextAt{16}{173.651}{198.491}{\Character{106}}%
\TextAt{13}{176.972}{198.491}{k}%
\TextAt{8}{188.323}{198.491}{we}%
\TextAt{8}{206.516}{198.491}{thus}%
\TextAt{8}{233.232}{198.491}{have}%
\TextAt{13}{261.191}{198.491}{A}%
\TextAt{21}{269.967}{196.697}{K}%
\TextAt{9}{283.134}{198.491}{\Character{61}}%
\TextAt{13}{297.204}{198.491}{K}%
\TextAt{3}{308.021}{202.829}{\Character{3}}%
\TextAt{8}{312.754}{198.491}{,}%
\TextAt{8}{321.109}{198.491}{and}%
\TextAt{8}{344.836}{198.491}{the}%
\TextAt{8}{365.957}{198.491}{crucial}%
\TextAt{8}{404.64}{198.491}{point}%
\TextAt{8}{436.174}{198.491}{is}%
\TextAt{8}{448.903}{198.491}{to}%
\TextAt{8}{464.178}{198.491}{verify}%
\TextAt{8}{497.661}{198.491}{for}%
\TextAt{8}{516.51}{198.491}{the}%
\TextAt{8}{79.226}{181.445}{multiplicative}%
\TextAt{8}{153.049}{181.445}{group}%
\TextAt{8}{186.225}{181.445}{of}%
\TextAt{8}{199.555}{181.445}{a}%
\TextAt{8}{209.31}{181.445}{local}%
\TextAt{8}{236.951}{181.445}{field}%
\TextAt{8}{262.308}{181.445}{the}%
\TextAt{8}{282.464}{181.445}{class}%
\TextAt{8}{309.901}{181.445}{field}%
\TextAt{8}{335.258}{181.445}{axiom:}%
\TextAt{12}{79.226}{155.432}{Theorem}%
\TextAt{12}{136.24}{155.432}{(1.1).}%
\TextAt{8}{171.556}{155.432}{For}%
\TextAt{8}{192.525}{155.432}{a}%
\TextAt{8}{202.281}{155.432}{cyclic}%
\TextAt{8}{234.464}{155.432}{extension}%
\TextAt{13}{286.23}{155.432}{L}%
\TextAt{16}{294.194}{155.432}{\Character{106}}%
\TextAt{13}{297.515}{155.432}{K}%
\TextAt{8}{312.23}{155.432}{of}%
\TextAt{8}{325.56}{155.432}{local}%
\TextAt{8}{353.201}{155.432}{fields,}%
\TextAt{8}{386.424}{155.432}{one}%
\TextAt{8}{407.884}{155.432}{has}%
\TextAt{8}{195.009}{114.972}{\char35\relax }%
\TextAt{13}{204.765}{114.972}{H}%
\TextAt{21}{215.387}{119.908}{i}%
\TextAt{9}{218.769}{114.972}{\Character{40}}%
\TextAt{13}{223.321}{114.972}{G}%
\TextAt{9}{232.555}{114.972}{\Character{40}}%
\TextAt{13}{237.107}{114.972}{L}%
\TextAt{16}{245.071}{114.972}{\Character{106}}%
\TextAt{13}{248.392}{114.972}{K}%
\TextAt{9}{259.21}{114.972}{\Character{41}}%
\TextAt{13}{263.762}{114.972}{\Character{59}}%
\TextAt{13}{268.986}{114.972}{L}%
\TextAt{3}{276.971}{119.908}{\Character{3}}%
\TextAt{9}{281.703}{114.972}{\Character{41}}%
\TextAt{9}{289.558}{114.972}{\Character{61}}%
\TextAt{20}{302.002}{139.48}{\Character{56}}%
\TextAt{20}{302.002}{128.72}{\Character{60}}%
\TextAt{20}{302.002}{107.201}{\Character{58}}%
\TextAt{9}{312.629}{124.097}{\Character{91}}%
\TextAt{13}{315.881}{124.097}{L}%
\TextAt{9}{327.166}{124.097}{\Character{58}}%
\TextAt{13}{333.738}{124.097}{K}%
\TextAt{9}{344.556}{124.097}{\Character{93}}%
\TextAt{8}{359.763}{124.097}{for}%
\TextAt{13}{377.648}{124.097}{i}%
\TextAt{9}{384.962}{124.097}{\Character{61}}%
\TextAt{9}{397.372}{124.097}{0}%
\TextAt{13}{403.241}{124.097}{\Character{59}}%
\TextAt{9}{312.629}{103.642}{1}%
\TextAt{8}{359.763}{103.642}{for}%
\TextAt{13}{377.648}{103.642}{i}%
\TextAt{9}{384.962}{103.642}{\Character{61}}%
\TextAt{16}{397.388}{103.642}{\Character{0}}%
\TextAt{9}{406.686}{103.642}{1}%
\TextAt{13}{412.539}{103.642}{\Character{58}}%
\TextAt{8}{297.211}{43.033}{291}%
\Destination{79.23}{773.075}{nav.chapter.5}%
\Destination{78.23}{809.815}{nav.page.291}%
\Destination{79.23}{570.065}{nav.section.5.1}%
\Destination{79.36}{629.855}{nav.print.317}%
\Destination{79.36}{169.455}{nav.number.5.1.1}%
\Destination{79.23}{773.075}{outline.48}%
\Destination{79.23}{570.065}{outline.49}%
\Hotlink{458.034}{477.955}{9.763}{11.047}{nav.chapter.2}%
\Hotlink{474.401}{477.955}{25.078}{11.047}{nav.number.2.5.2}%
}

% Page 304 (body page 292)
\Page{%
\TextAt{2}{79.226}{796.943}{V.}%
\TextAt{2}{94.073}{796.943}{Local}%
\TextAt{2}{123.615}{796.943}{Class}%
\TextAt{2}{152.208}{796.943}{Field}%
\TextAt{2}{179.939}{796.943}{Theory}%
\Vector{q 0 G 0 g 0.299 w 7.226 18.729 m 460.769 18.729 l S Q}%
\TextAt{19}{79.226}{761.124}{Proof.}%
\TextAt{8}{116}{761.124}{For}%
\TextAt{13}{137.68}{761.124}{i}%
\TextAt{9}{146.204}{761.124}{\Character{61}}%
\TextAt{16}{159.84}{761.124}{\Character{0}}%
\TextAt{9}{169.138}{761.124}{1}%
\TextAt{8}{179.599}{761.124}{this}%
\TextAt{8}{203.131}{761.124}{is}%
\TextAt{8}{215.605}{761.124}{the}%
\TextAt{8}{236.472}{761.124}{claim}%
\TextAt{8}{268.397}{761.124}{of}%
\TextAt{8}{282.437}{761.124}{proposition}%
\TextAt{8}{344.68}{761.124}{(3.5)}%
\TextAt{8}{373.366}{761.124}{(“Hilbert}%
\TextAt{8}{424.467}{761.124}{90”)}%
\TextAt{8}{450.869}{761.124}{in}%
\TextAt{8}{465.232}{761.124}{chap.}%
\TextAt{8}{499.553}{761.124}{IV.}%
\TextAt{8}{520.408}{761.124}{So}%
\TextAt{8}{79.226}{744.078}{all}%
\TextAt{8}{95.854}{744.078}{we}%
\TextAt{8}{113.45}{744.078}{have}%
\TextAt{8}{140.814}{744.078}{to}%
\TextAt{8}{155.493}{744.078}{show}%
\TextAt{8}{184.866}{744.078}{is}%
\TextAt{8}{196.998}{744.078}{that}%
\TextAt{8}{222.736}{744.078}{the}%
\TextAt{8}{243.261}{744.078}{Herbrand}%
\TextAt{8}{295.982}{744.078}{quotient}%
\TextAt{8}{342.534}{744.078}{is}%
\TextAt{13}{354.667}{744.078}{h}%
\TextAt{9}{361.405}{744.078}{\Character{40}}%
\TextAt{13}{365.958}{744.078}{G\Character{59}}%
\TextAt{13}{380.412}{744.078}{L}%
\TextAt{3}{388.4}{748.416}{\Character{3}}%
\TextAt{9}{393.132}{744.078}{\Character{41}}%
\TextAt{9}{401.62}{744.078}{\Character{61}}%
\TextAt{8}{414.688}{744.078}{\char35\relax }%
\TextAt{13}{424.444}{744.078}{H}%
\TextAt{17}{435.066}{748.416}{0}%
\TextAt{9}{439.799}{744.078}{\Character{40}}%
\TextAt{13}{444.351}{744.078}{G\Character{59}}%
\TextAt{13}{458.805}{744.078}{L}%
\TextAt{3}{466.793}{748.416}{\Character{3}}%
\TextAt{9}{471.525}{744.078}{\Character{41}}%
\TextAt{9}{480.013}{744.078}{\Character{61}}%
\TextAt{9}{493.056}{744.078}{\Character{91}}%
\TextAt{13}{496.333}{744.078}{L}%
\TextAt{9}{508.247}{744.078}{\Character{58}}%
\TextAt{13}{515.448}{744.078}{K}%
\TextAt{9}{526.266}{744.078}{\Character{93}}%
\TextAt{8}{529.518}{744.078}{,}%
\TextAt{8}{79.226}{727.032}{where}%
\TextAt{8}{113.035}{727.032}{we}%
\TextAt{8}{130.263}{727.032}{have}%
\TextAt{8}{157.258}{727.032}{put}%
\TextAt{13}{178.717}{727.032}{G}%
\TextAt{9}{191.272}{727.032}{\Character{61}}%
\TextAt{13}{203.697}{727.032}{G}%
\TextAt{9}{212.931}{727.032}{\Character{40}}%
\TextAt{13}{217.483}{727.032}{L}%
\TextAt{16}{225.447}{727.032}{\Character{106}}%
\TextAt{13}{228.768}{727.032}{K}%
\TextAt{9}{239.586}{727.032}{\Character{41}}%
\TextAt{8}{244.138}{727.032}{.}%
\TextAt{8}{252.587}{727.032}{The}%
\TextAt{8}{276.64}{727.032}{exact}%
\TextAt{8}{307.533}{727.032}{sequence}%
\TextAt{9}{230.773}{702.074}{1}%
\NativeArrowAt{239.947}{702.074}{11.955}{19.26}{\longrightarrow}%
\TextAt{13}{262.529}{702.074}{U}%
\TextAt{21}{270.479}{700.281}{L}%
\NativeArrowAt{280.057}{702.074}{11.955}{19.26}{\longrightarrow}%
\TextAt{13}{302.639}{702.074}{L}%
\TextAt{3}{310.604}{707.011}{\Character{3}}%
\TextAt{21}{321.245}{709.128}{v}%
\TextAt{22}{325.368}{707.789}{L}%
\NativeArrowAt{318.657}{702.074}{11.955}{16.259}{\longrightarrow}%
\TextAt{18}{338.245}{702.074}{Z}%
\NativeArrowAt{349.536}{702.074}{11.955}{19.26}{\longrightarrow}%
\TextAt{9}{372.118}{702.074}{0}%
\TextAt{13}{377.971}{702.074}{\Character{59}}%
\TextAt{8}{79.226}{677.117}{in}%
\TextAt{8}{92.879}{677.117}{which}%
\TextAt{18}{126.365}{677.117}{Z}%
\TextAt{8}{138.233}{677.117}{has}%
\TextAt{8}{159.107}{677.117}{to}%
\TextAt{8}{173.417}{677.117}{be}%
\TextAt{8}{189.341}{677.117}{viewed}%
\TextAt{8}{227.705}{677.117}{as}%
\TextAt{8}{242.076}{677.117}{the}%
\TextAt{8}{262.232}{677.117}{trivial}%
\TextAt{13}{297.033}{677.117}{G}%
\TextAt{8}{306.267}{677.117}{-module,}%
\TextAt{8}{354.709}{677.117}{yields,}%
\TextAt{8}{390.862}{677.117}{by}%
\TextAt{8}{407.121}{677.117}{chap.}%
\TextAt{8}{439.312}{677.117}{IV,}%
\TextAt{8}{459.457}{677.117}{(7.3),}%
\TextAt{13}{185.853}{652.16}{h}%
\TextAt{9}{192.592}{652.16}{\Character{40}}%
\TextAt{13}{197.144}{652.16}{G\Character{59}}%
\TextAt{13}{211.598}{652.16}{L}%
\TextAt{3}{219.587}{657.096}{\Character{3}}%
\TextAt{9}{224.319}{652.16}{\Character{41}}%
\TextAt{9}{232.174}{652.16}{\Character{61}}%
\TextAt{13}{244.618}{652.16}{h}%
\TextAt{9}{251.356}{652.16}{\Character{40}}%
\TextAt{13}{255.908}{652.16}{G\Character{59}}%
\TextAt{18}{270.386}{652.16}{Z}%
\TextAt{9}{278.356}{652.16}{\Character{41}}%
\TextAt{13}{282.909}{652.16}{h}%
\TextAt{9}{289.647}{652.16}{\Character{40}}%
\TextAt{13}{294.2}{652.16}{G\Character{59}}%
\TextAt{13}{308.654}{652.16}{U}%
\TextAt{21}{316.628}{650.366}{L}%
\TextAt{9}{322.885}{652.16}{\Character{41}}%
\TextAt{9}{330.74}{652.16}{\Character{61}}%
\TextAt{9}{343.161}{652.16}{\Character{91}}%
\TextAt{13}{346.435}{652.16}{L}%
\TextAt{9}{357.72}{652.16}{\Character{58}}%
\TextAt{13}{364.293}{652.16}{K}%
\TextAt{9}{375.111}{652.16}{\Character{93}}%
\TextAt{13}{378.362}{652.16}{h}%
\TextAt{9}{385.101}{652.16}{\Character{40}}%
\TextAt{13}{389.653}{652.16}{G\Character{59}}%
\TextAt{13}{404.107}{652.16}{U}%
\TextAt{21}{412.081}{650.366}{L}%
\TextAt{9}{418.338}{652.16}{\Character{41}}%
\TextAt{13}{422.891}{652.16}{\Character{58}}%
\TextAt{8}{79.226}{627.202}{Hence}%
\TextAt{8}{113.176}{627.202}{we}%
\TextAt{8}{129.576}{627.202}{have}%
\TextAt{8}{155.744}{627.202}{to}%
\TextAt{8}{169.227}{627.202}{show}%
\TextAt{8}{197.402}{627.202}{that}%
\TextAt{13}{221.944}{627.202}{h}%
\TextAt{9}{228.682}{627.202}{\Character{40}}%
\TextAt{13}{233.235}{627.202}{G\Character{59}}%
\TextAt{13}{247.689}{627.202}{U}%
\TextAt{21}{255.663}{625.409}{L}%
\TextAt{9}{261.92}{627.202}{\Character{41}}%
\TextAt{9}{269.775}{627.202}{\Character{61}}%
\TextAt{9}{282.196}{627.202}{1}%
\TextAt{8}{288.071}{627.202}{.}%
\TextAt{8}{296.244}{627.202}{For}%
\TextAt{8}{316.386}{627.202}{this}%
\TextAt{8}{338.381}{627.202}{we}%
\TextAt{8}{354.781}{627.202}{choose}%
\TextAt{8}{391.086}{627.202}{a}%
\TextAt{8}{400.014}{627.202}{normal}%
\TextAt{8}{438.866}{627.202}{basis}%
\TextAt{16}{466.779}{627.202}{\Character{102}}%
\TextAt{13}{472.756}{627.202}{\Character{11}}%
\TextAt{21}{480.278}{631.541}{\Character{27}}%
\TextAt{16}{489.222}{627.202}{\Character{106}}%
\TextAt{13}{495.864}{627.202}{\Character{27}}%
\NativeMathAt{506.267}{627.202}{11.955}{\in}%
\TextAt{13}{517.558}{627.202}{G}%
\TextAt{16}{526.792}{627.202}{\Character{103}}%
\TextAt{8}{79.226}{610.156}{of}%
\TextAt{13}{91.961}{610.156}{L}%
\TextAt{16}{99.926}{610.156}{\Character{106}}%
\TextAt{13}{103.247}{610.156}{K}%
\TextAt{8}{117.367}{610.156}{(see}%
\TextAt{8}{140.24}{610.156}{[93],}%
\TextAt{8}{165.132}{610.156}{chap.}%
\TextAt{8}{197.125}{610.156}{VIII,}%
\TextAt{8}{225.761}{610.156}{§12,}%
\TextAt{8}{249.829}{610.156}{th.20),}%
\TextAt{13}{287.083}{610.156}{\Character{11}}%
\NativeMathAt{297.926}{610.156}{11.955}{\in}%
\TextAt{14}{309.217}{610.156}{o}%
\TextAt{21}{315.061}{608.363}{L}%
\TextAt{8}{321.318}{610.156}{,}%
\TextAt{8}{327.991}{610.156}{and}%
\TextAt{8}{350.159}{610.156}{consider}%
\TextAt{8}{395.149}{610.156}{in}%
\TextAt{14}{408.207}{610.156}{o}%
\TextAt{21}{414.051}{608.363}{L}%
\TextAt{8}{423.611}{610.156}{the}%
\TextAt{8}{443.172}{610.156}{open}%
\TextAt{8}{470.863}{610.156}{(and}%
\TextAt{8}{497.585}{610.156}{closed)}%
\TextAt{13}{79.226}{593.11}{G}%
\TextAt{8}{88.46}{593.11}{-module}%
\TextAt{13}{133.65}{593.11}{M}%
\TextAt{9}{149.545}{593.11}{\Character{61}}%
\TextAt{20}{161.97}{602.077}{\Character{80}}%
\TextAt{21}{174.59}{589.623}{\Character{27}}%
\NativeMathAt{179.715}{589.623}{7.97}{\in}%
\TextAt{21}{185.36}{589.623}{G}%
\TextAt{14}{194.45}{593.11}{o}%
\TextAt{21}{200.294}{591.317}{K}%
\TextAt{13}{208.497}{593.11}{\Character{11}}%
\TextAt{21}{216.018}{597.449}{\Character{27}}%
\TextAt{8}{221.642}{593.11}{.}%
\TextAt{8}{230.09}{593.11}{Then}%
\TextAt{8}{260.648}{593.11}{the}%
\TextAt{8}{280.804}{593.11}{open}%
\TextAt{8}{309.09}{593.11}{sets}%
\TextAt{13}{220.542}{568.153}{V}%
\TextAt{21}{229.967}{573.089}{n}%
\TextAt{9}{238.924}{568.153}{\Character{61}}%
\TextAt{9}{251.334}{568.153}{1}%
\TextAt{9}{259.846}{568.153}{\Character{43}}%
\TextAt{13}{271.621}{568.153}{\Character{25}}%
\TextAt{21}{278.69}{573.089}{n}%
\TextAt{21}{278.261}{565.198}{K}%
\TextAt{13}{286.464}{568.153}{M}%
\TextAt{13}{298.383}{568.153}{\Character{59}}%
\TextAt{13}{327.518}{568.153}{n}%
\TextAt{9}{337.85}{568.153}{\Character{61}}%
\TextAt{9}{350.26}{568.153}{1}%
\TextAt{13}{356.129}{568.153}{\Character{59}}%
\TextAt{9}{361.373}{568.153}{2}%
\TextAt{13}{367.226}{568.153}{\Character{59}}%
\TextAt{13}{372.45}{568.153}{\Character{58}}%
\TextAt{13}{377.687}{568.153}{\Character{58}}%
\TextAt{13}{382.911}{568.153}{\Character{58}}%
\TextAt{13}{388.148}{568.153}{\Character{59}}%
\TextAt{8}{79.226}{543.196}{form}%
\TextAt{8}{107.396}{543.196}{a}%
\TextAt{8}{117.682}{543.196}{basis}%
\TextAt{8}{146.952}{543.196}{of}%
\TextAt{8}{160.812}{543.196}{open}%
\TextAt{8}{189.628}{543.196}{neighbourhoods}%
\TextAt{8}{274.777}{543.196}{of}%
\TextAt{9}{288.637}{543.196}{1}%
\TextAt{8}{298.917}{543.196}{in}%
\TextAt{13}{313.1}{543.196}{U}%
\TextAt{21}{321.05}{541.403}{L}%
\TextAt{8}{327.307}{543.196}{.}%
\TextAt{8}{337.346}{543.196}{Since}%
\TextAt{13}{368.433}{543.196}{M}%
\TextAt{8}{385.434}{543.196}{is}%
\TextAt{8}{397.728}{543.196}{open,}%
\TextAt{8}{429.928}{543.196}{we}%
\TextAt{8}{447.686}{543.196}{have}%
\TextAt{13}{475.21}{543.196}{\Character{25}}%
\TextAt{21}{482.28}{547.534}{N}%
\TextAt{21}{481.851}{540.175}{K}%
\TextAt{14}{490.348}{543.196}{o}%
\TextAt{21}{496.192}{541.403}{L}%
\NativeMathAt{506.673}{543.196}{11.955}{\subseteq}%
\TextAt{13}{520.196}{543.196}{M}%
\TextAt{8}{79.226}{526.15}{for}%
\TextAt{8}{96.949}{526.15}{suitable}%
\TextAt{13}{140.424}{526.15}{N}%
\TextAt{8}{151.046}{526.15}{,}%
\TextAt{8}{158.066}{526.15}{and}%
\TextAt{8}{180.667}{526.15}{for}%
\TextAt{13}{198.39}{526.15}{n}%
\NativeMathAt{208.699}{526.15}{11.955}{\geq}%
\TextAt{13}{221.318}{526.15}{N}%
\TextAt{8}{235.676}{526.15}{the}%
\TextAt{13}{255.671}{526.15}{V}%
\TextAt{21}{265.096}{530.488}{n}%
\TextAt{8}{274.468}{526.15}{are}%
\TextAt{8}{293.817}{526.15}{even}%
\TextAt{8}{320.316}{526.15}{subgroups}%
\TextAt{8}{375.566}{526.15}{(of}%
\TextAt{8}{393.289}{526.15}{finite}%
\TextAt{8}{423.039}{526.15}{index)}%
\TextAt{8}{458.97}{526.15}{of}%
\TextAt{13}{472.138}{526.15}{U}%
\TextAt{21}{480.088}{524.357}{L}%
\TextAt{8}{486.346}{526.15}{,}%
\TextAt{8}{493.365}{526.15}{because}%
\TextAt{8}{79.226}{509.104}{we}%
\TextAt{8}{96.454}{509.104}{have}%
\TextAt{9}{165.285}{492.058}{\Character{40}}%
\TextAt{13}{169.837}{492.058}{\Character{25}}%
\TextAt{21}{176.907}{496.994}{n}%
\TextAt{21}{176.478}{489.102}{K}%
\TextAt{13}{184.68}{492.058}{M}%
\TextAt{9}{197.254}{492.058}{\Character{41}\Character{40}}%
\TextAt{13}{206.359}{492.058}{\Character{25}}%
\TextAt{21}{213.428}{496.994}{n}%
\TextAt{21}{212.999}{489.102}{K}%
\TextAt{13}{221.202}{492.058}{M}%
\TextAt{9}{233.775}{492.058}{\Character{41}}%
\TextAt{9}{241.63}{492.058}{\Character{61}}%
\TextAt{13}{254.074}{492.058}{\Character{25}}%
\TextAt{17}{261.143}{496.994}{2}%
\TextAt{21}{265.377}{496.994}{n}%
\TextAt{21}{260.714}{489.102}{K}%
\TextAt{13}{271.014}{492.058}{M}%
\TextAt{13}{283.579}{492.058}{M}%
\NativeMathAt{299.482}{492.058}{11.955}{\subseteq}%
\TextAt{13}{312.101}{492.058}{\Character{25}}%
\TextAt{17}{319.17}{496.994}{2}%
\TextAt{21}{323.404}{496.994}{n}%
\TextAt{21}{318.741}{489.102}{K}%
\TextAt{14}{329.041}{492.058}{o}%
\TextAt{21}{334.885}{490.265}{L}%
\NativeMathAt{344.463}{492.058}{11.955}{\subseteq}%
\TextAt{13}{357.082}{492.058}{\Character{25}}%
\TextAt{17}{364.152}{497.006}{2}%
\TextAt{21}{368.386}{497.006}{n}%
\TextAt{3}{373.524}{497.006}{\Character{0}}%
\TextAt{21}{380.11}{497.006}{N}%
\TextAt{21}{363.723}{488.817}{K}%
\TextAt{13}{388.179}{492.058}{M}%
\NativeMathAt{404.073}{492.058}{11.955}{\subseteq}%
\TextAt{13}{416.693}{492.058}{\Character{25}}%
\TextAt{21}{423.762}{496.994}{n}%
\TextAt{21}{423.333}{489.102}{K}%
\TextAt{13}{431.536}{492.058}{M}%
\TextAt{13}{443.455}{492.058}{\Character{58}}%
\TextAt{8}{79.226}{470.269}{Hence}%
\TextAt{13}{114.756}{470.269}{V}%
\TextAt{21}{124.182}{474.608}{n}%
\TextAt{13}{129.818}{470.269}{V}%
\TextAt{21}{139.244}{474.608}{n}%
\NativeMathAt{149.484}{470.269}{11.955}{\subseteq}%
\TextAt{13}{163.386}{470.269}{V}%
\TextAt{21}{172.812}{474.608}{n}%
\TextAt{8}{178.448}{470.269}{,}%
\TextAt{8}{186.538}{470.269}{and}%
\TextAt{8}{210.053}{470.269}{since}%
\TextAt{9}{239.474}{470.269}{1}%
\TextAt{16}{248.497}{470.269}{\Character{0}}%
\TextAt{13}{260.965}{470.269}{\Character{25}}%
\TextAt{21}{268.034}{474.608}{n}%
\TextAt{21}{267.606}{467.249}{K}%
\TextAt{13}{275.808}{470.269}{\Character{22}}%
\TextAt{8}{282.851}{470.269}{,}%
\TextAt{8}{290.941}{470.269}{for}%
\TextAt{13}{309.579}{470.269}{\Character{22}}%
\NativeMathAt{321.225}{470.269}{11.955}{\in}%
\TextAt{13}{333.799}{470.269}{M}%
\TextAt{8}{346.372}{470.269}{,}%
\TextAt{8}{354.462}{470.269}{lies}%
\TextAt{8}{375.431}{470.269}{in}%
\TextAt{13}{389.837}{470.269}{V}%
\TextAt{21}{399.262}{474.608}{n}%
\TextAt{8}{404.899}{470.269}{,}%
\TextAt{8}{412.989}{470.269}{so}%
\TextAt{8}{428.111}{470.269}{does}%
\TextAt{9}{455.261}{470.269}{\Character{40}1}%
\TextAt{16}{468.836}{470.269}{\Character{0}}%
\TextAt{13}{481.304}{470.269}{\Character{25}}%
\TextAt{21}{488.374}{474.608}{n}%
\TextAt{21}{487.945}{467.249}{K}%
\TextAt{13}{496.147}{470.269}{\Character{22}}%
\TextAt{9}{503.19}{470.269}{\Character{41}}%
\TextAt{3}{507.743}{474.608}{\Character{0}}%
\TextAt{17}{514.329}{474.608}{1}%
\TextAt{9}{523.665}{470.269}{\Character{61}}%
\TextAt{9}{79.226}{453.223}{1}%
\TextAt{9}{88.826}{453.223}{\Character{43}}%
\TextAt{13}{101.695}{453.223}{\Character{25}}%
\TextAt{21}{108.764}{457.562}{n}%
\TextAt{21}{108.335}{450.203}{K}%
\TextAt{9}{116.538}{453.223}{\Character{40}}%
\TextAt{20}{121.09}{462.19}{\Character{80}}%
\NativeMathAt{133.71}{459.035}{7.97}{\infty}%
\TextAt{21}{133.71}{449.736}{i}%
\TextAt{17}{136.593}{449.736}{\Character{61}1}%
\TextAt{13}{149.904}{453.223}{\Character{22}}%
\TextAt{21}{156.947}{457.562}{i}%
\TextAt{13}{160.328}{453.223}{\Character{25}}%
\TextAt{21}{167.398}{459.334}{n}%
\TextAt{17}{172.536}{459.334}{\Character{40}}%
\TextAt{21}{175.829}{459.334}{i}%
\TextAt{3}{178.712}{459.334}{\Character{0}}%
\TextAt{17}{185.299}{459.334}{1\Character{41}}%
\TextAt{21}{166.969}{449.982}{K}%
\TextAt{9}{193.324}{453.223}{\Character{41}}%
\TextAt{8}{197.877}{453.223}{.}%
\TextAt{8}{211.161}{453.223}{Via}%
\TextAt{8}{234.556}{453.223}{the}%
\TextAt{8}{256.324}{453.223}{correspondence}%
\TextAt{9}{339.614}{453.223}{1}%
\TextAt{9}{349.214}{453.223}{\Character{43}}%
\TextAt{13}{362.083}{453.223}{\Character{25}}%
\TextAt{21}{369.152}{457.562}{n}%
\TextAt{21}{368.723}{450.203}{K}%
\TextAt{13}{376.926}{453.223}{\Character{11}}%
\NativeArrowAt{390.515}{453.223}{11.955}{11.955}{\mapsto}%
\TextAt{13}{408.539}{453.223}{\Character{11}}%
\TextAt{8}{419.381}{453.223}{mod}%
\TextAt{13}{444.819}{453.223}{\Character{25}}%
\TextAt{21}{451.459}{451.43}{K}%
\TextAt{13}{459.662}{453.223}{M}%
\TextAt{8}{472.236}{453.223}{,}%
\TextAt{8}{481.4}{453.223}{we}%
\TextAt{8}{500.239}{453.223}{obtain}%
\ErrataAt{79.19}{436.177}{455.81}{$G$-isomorphisms for $n>N$ as in II, \hyperlink{nav.number.2.3.10}{(3.10)},}%
\TextAt{13}{163.377}{411.22}{V}%
\TextAt{21}{172.803}{416.156}{n}%
\TextAt{9}{178.439}{411.22}{\Character{47}}%
\TextAt{13}{184.292}{411.22}{V}%
\TextAt{21}{193.718}{416.156}{n}%
\TextAt{17}{198.856}{416.156}{\Character{43}1}%
\NativeMathAt{213.496}{411.22}{11.955}{\cong}%
\TextAt{13}{226.115}{411.22}{M}%
\TextAt{9}{238.689}{411.22}{\Character{47}}%
\TextAt{13}{244.542}{411.22}{\Character{25}}%
\TextAt{21}{251.182}{409.427}{K}%
\TextAt{13}{259.385}{411.22}{M}%
\TextAt{9}{275.279}{411.22}{\Character{61}}%
\TextAt{20}{287.705}{422.577}{\Character{77}}%
\TextAt{21}{288.053}{397.206}{\Character{27}}%
\NativeMathAt{293.178}{397.206}{7.97}{\in}%
\TextAt{21}{298.824}{397.206}{G}%
\TextAt{9}{305.77}{411.22}{\Character{40}}%
\TextAt{14}{310.323}{411.22}{o}%
\TextAt{21}{316.167}{409.427}{K}%
\TextAt{9}{324.37}{411.22}{\Character{47}}%
\TextAt{14}{330.223}{411.22}{p}%
\TextAt{21}{336.206}{409.427}{K}%
\TextAt{9}{344.408}{411.22}{\Character{41}}%
\TextAt{13}{348.961}{411.22}{\Character{11}}%
\TextAt{21}{356.482}{416.156}{\Character{27}}%
\TextAt{9}{365.427}{411.22}{\Character{61}}%
\TextAt{8}{377.852}{411.22}{Ind}%
\TextAt{21}{395.08}{409.427}{G}%
\TextAt{9}{402.176}{411.22}{\Character{40}}%
\TextAt{14}{406.729}{411.22}{o}%
\TextAt{21}{412.573}{409.427}{K}%
\TextAt{9}{420.776}{411.22}{\Character{47}}%
\TextAt{14}{426.629}{411.22}{p}%
\TextAt{21}{432.612}{409.427}{K}%
\TextAt{9}{440.814}{411.22}{\Character{41}}%
\TextAt{13}{445.367}{411.22}{\Character{58}}%
\TextAt{8}{79.226}{376.992}{So}%
\TextAt{8}{96.008}{376.992}{by}%
\TextAt{8}{112.789}{376.992}{chap.}%
\TextAt{8}{146.547}{376.992}{IV,}%
\TextAt{8}{167.214}{376.992}{(7.4),}%
\TextAt{8}{199.094}{376.992}{we}%
\TextAt{8}{216.844}{376.992}{have}%
\TextAt{13}{244.361}{376.992}{H}%
\TextAt{21}{254.984}{381.331}{i}%
\TextAt{9}{258.365}{376.992}{\Character{40}}%
\TextAt{13}{262.917}{376.992}{G\Character{59}}%
\TextAt{13}{277.371}{376.992}{V}%
\TextAt{21}{286.821}{381.331}{n}%
\TextAt{9}{292.457}{376.992}{\Character{47}}%
\TextAt{13}{298.31}{376.992}{V}%
\TextAt{21}{307.736}{381.331}{n}%
\TextAt{17}{312.874}{381.331}{\Character{43}1}%
\TextAt{9}{324.193}{376.992}{\Character{41}}%
\TextAt{9}{332.944}{376.992}{\Character{61}}%
\TextAt{9}{346.25}{376.992}{1}%
\TextAt{8}{356.545}{376.992}{for}%
\TextAt{13}{374.952}{376.992}{i}%
\TextAt{9}{383.156}{376.992}{\Character{61}}%
\TextAt{9}{396.462}{376.992}{0}%
\TextAt{13}{402.325}{376.992}{\Character{59}}%
\TextAt{16}{407.569}{376.992}{\Character{0}}%
\TextAt{9}{416.868}{376.992}{1}%
\TextAt{8}{427.141}{376.992}{and}%
\TextAt{13}{450.426}{376.992}{n}%
\NativeMathAt{461.624}{376.992}{11.955}{>}%
\TextAt{13}{475.134}{376.992}{N}%
\TextAt{8}{485.756}{376.992}{.}%
\TextAt{8}{495.772}{376.992}{This}%
\TextAt{8}{523.014}{376.992}{in}%
\TextAt{8}{79.226}{359.946}{turn}%
\TextAt{8}{105.624}{359.946}{implies}%
\TextAt{8}{145.734}{359.946}{that}%
\TextAt{13}{171.486}{359.946}{H}%
\TextAt{21}{182.109}{364.285}{i}%
\TextAt{9}{185.49}{359.946}{\Character{40}}%
\TextAt{13}{190.043}{359.946}{G\Character{59}}%
\TextAt{13}{204.497}{359.946}{V}%
\TextAt{21}{213.946}{364.285}{n}%
\TextAt{9}{219.582}{359.946}{\Character{41}}%
\TextAt{9}{228.094}{359.946}{\Character{61}}%
\TextAt{9}{241.161}{359.946}{1}%
\TextAt{8}{251.321}{359.946}{for}%
\TextAt{13}{269.59}{359.946}{i}%
\TextAt{9}{277.557}{359.946}{\Character{61}}%
\TextAt{9}{290.624}{359.946}{0}%
\TextAt{13}{296.489}{359.946}{\Character{59}}%
\TextAt{16}{301.733}{359.946}{\Character{0}}%
\TextAt{9}{311.031}{359.946}{1}%
\TextAt{8}{321.165}{359.946}{and}%
\TextAt{13}{344.311}{359.946}{n}%
\NativeMathAt{355.273}{359.946}{11.955}{>}%
\TextAt{13}{368.545}{359.946}{N}%
\TextAt{8}{379.168}{359.946}{.}%
\TextAt{8}{388.766}{359.946}{Indeed,}%
\TextAt{8}{430.527}{359.946}{if}%
\TextAt{8}{441.634}{359.946}{for}%
\TextAt{8}{459.902}{359.946}{instance}%
\TextAt{13}{505.87}{359.946}{i}%
\TextAt{9}{513.838}{359.946}{\Character{61}}%
\TextAt{9}{526.905}{359.946}{0}%
\TextAt{8}{79.226}{342.9}{and}%
\TextAt{13}{102.009}{342.9}{a}%
\NativeMathAt{111.509}{342.9}{11.955}{\in}%
\TextAt{9}{122.833}{342.9}{\Character{40}}%
\TextAt{13}{127.386}{342.9}{V}%
\TextAt{21}{136.811}{347.239}{n}%
\TextAt{9}{142.448}{342.9}{\Character{41}}%
\TextAt{21}{147}{347.239}{G}%
\TextAt{8}{154.097}{342.9}{,}%
\TextAt{8}{161.271}{342.9}{then}%
\TextAt{13}{187.951}{342.9}{a}%
\TextAt{9}{197.451}{342.9}{\Character{61}}%
\TextAt{9}{209.896}{342.9}{\Character{40}}%
\TextAt{13}{214.463}{342.9}{N}%
\TextAt{21}{223.836}{341.107}{G}%
\TextAt{13}{230.933}{342.9}{b}%
\TextAt{17}{235.91}{341.107}{0}%
\TextAt{9}{240.643}{342.9}{\Character{41}}%
\TextAt{13}{245.195}{342.9}{a}%
\TextAt{17}{251.34}{341.107}{1}%
\TextAt{8}{256.072}{342.9}{,}%
\TextAt{8}{263.246}{342.9}{with}%
\TextAt{13}{289.926}{342.9}{b}%
\TextAt{17}{294.903}{341.107}{0}%
\NativeMathAt{302.991}{342.9}{11.955}{\in}%
\TextAt{13}{314.316}{342.9}{V}%
\TextAt{21}{323.741}{347.239}{n}%
\TextAt{8}{329.378}{342.9}{,}%
\TextAt{13}{336.552}{342.9}{a}%
\TextAt{17}{342.697}{341.107}{1}%
\NativeMathAt{350.784}{342.9}{11.955}{\in}%
\TextAt{9}{362.109}{342.9}{\Character{40}}%
\TextAt{13}{366.661}{342.9}{V}%
\TextAt{21}{376.087}{347.239}{n}%
\TextAt{17}{381.225}{347.239}{\Character{43}1}%
\TextAt{9}{392.544}{342.9}{\Character{41}}%
\TextAt{21}{397.096}{347.239}{G}%
\TextAt{8}{408.11}{342.9}{and}%
\TextAt{8}{430.893}{342.9}{thus}%
\TextAt{13}{456.664}{342.9}{a}%
\TextAt{17}{462.809}{341.107}{1}%
\TextAt{9}{470.896}{342.9}{\Character{61}}%
\TextAt{9}{483.341}{342.9}{\Character{40}}%
\TextAt{13}{487.908}{342.9}{N}%
\TextAt{21}{497.282}{341.107}{G}%
\TextAt{13}{504.379}{342.9}{b}%
\TextAt{17}{509.356}{341.107}{1}%
\TextAt{9}{514.088}{342.9}{\Character{41}}%
\TextAt{13}{518.641}{342.9}{a}%
\TextAt{17}{524.785}{341.107}{2}%
\TextAt{8}{529.518}{342.9}{,}%
\TextAt{8}{79.226}{325.854}{for}%
\TextAt{8}{97.111}{325.854}{some}%
\TextAt{13}{126.437}{325.854}{b}%
\TextAt{17}{131.414}{324.061}{1}%
\NativeMathAt{139.467}{325.854}{11.955}{\in}%
\TextAt{13}{150.758}{325.854}{V}%
\TextAt{21}{160.184}{330.193}{n}%
\TextAt{17}{165.322}{330.193}{\Character{43}1}%
\TextAt{8}{176.641}{325.854}{,}%
\TextAt{13}{183.79}{325.854}{a}%
\TextAt{17}{189.935}{324.061}{2}%
\NativeMathAt{197.988}{325.854}{11.955}{\in}%
\TextAt{9}{209.279}{325.854}{\Character{40}}%
\TextAt{13}{213.832}{325.854}{V}%
\TextAt{21}{223.257}{330.193}{n}%
\TextAt{17}{228.396}{330.193}{\Character{43}2}%
\TextAt{9}{239.714}{325.854}{\Character{41}}%
\TextAt{21}{244.267}{330.193}{G}%
\TextAt{8}{251.363}{325.854}{,}%
\TextAt{8}{258.513}{325.854}{etc.;}%
\TextAt{8}{283.87}{325.854}{in}%
\TextAt{8}{297.522}{325.854}{general,}%
\TextAt{13}{174.411}{300.897}{a}%
\TextAt{21}{180.556}{299.104}{i}%
\TextAt{9}{187.258}{300.897}{\Character{61}}%
\TextAt{9}{199.668}{300.897}{\Character{40}}%
\TextAt{13}{204.236}{300.897}{N}%
\TextAt{21}{213.61}{299.104}{G}%
\TextAt{13}{220.707}{300.897}{b}%
\TextAt{21}{225.684}{299.104}{i}%
\TextAt{9}{229.065}{300.897}{\Character{41}}%
\TextAt{13}{233.618}{300.897}{a}%
\TextAt{21}{239.763}{299.104}{i}%
\TextAt{17}{242.646}{299.104}{\Character{43}1}%
\TextAt{13}{253.964}{300.897}{\Character{59}}%
\TextAt{13}{283.099}{300.897}{b}%
\TextAt{21}{288.096}{299.104}{i}%
\NativeMathAt{294.798}{300.897}{11.955}{\in}%
\TextAt{13}{306.089}{300.897}{V}%
\TextAt{21}{315.515}{305.833}{n}%
\TextAt{17}{320.653}{305.833}{\Character{43}}%
\TextAt{21}{327.24}{305.833}{i}%
\TextAt{13}{330.621}{300.897}{\Character{59}}%
\TextAt{13}{347.801}{300.897}{a}%
\TextAt{21}{353.965}{299.104}{i}%
\TextAt{17}{356.848}{299.104}{\Character{43}1}%
\NativeMathAt{371.488}{300.897}{11.955}{\in}%
\TextAt{9}{382.779}{300.897}{\Character{40}}%
\TextAt{13}{387.331}{300.897}{V}%
\TextAt{21}{396.757}{305.833}{n}%
\TextAt{17}{401.895}{305.833}{\Character{43}}%
\TextAt{21}{408.482}{305.833}{i}%
\TextAt{17}{411.365}{305.833}{\Character{43}1}%
\TextAt{9}{422.684}{300.897}{\Character{41}}%
\TextAt{21}{427.236}{305.833}{G}%
\TextAt{13}{434.333}{300.897}{\Character{58}}%
\TextAt{8}{79.226}{275.94}{This}%
\TextAt{8}{105.386}{275.94}{yields}%
\TextAt{13}{137.726}{275.94}{a}%
\TextAt{9}{147.192}{275.94}{\Character{61}}%
\TextAt{13}{159.617}{275.94}{N}%
\TextAt{21}{168.991}{274.146}{G}%
\TextAt{13}{176.088}{275.94}{b}%
\TextAt{8}{181.065}{275.94}{,}%
\TextAt{8}{187.766}{275.94}{with}%
\TextAt{8}{213.865}{275.94}{the}%
\TextAt{8}{233.461}{275.94}{convergent}%
\TextAt{8}{291.445}{275.94}{product}%
\TextAt{13}{334.784}{275.94}{b}%
\TextAt{9}{343.082}{275.94}{\Character{61}}%
\TextAt{20}{355.508}{284.906}{\Character{81}}%
\NativeMathAt{366.799}{281.751}{7.97}{\infty}%
\TextAt{21}{366.799}{272.453}{i}%
\TextAt{17}{369.682}{272.453}{\Character{61}0}%
\TextAt{13}{382.993}{275.94}{b}%
\TextAt{21}{387.97}{274.146}{i}%
\NativeMathAt{394.673}{275.94}{11.955}{\in}%
\TextAt{13}{405.964}{275.94}{V}%
\TextAt{21}{415.389}{280.278}{n}%
\TextAt{8}{421.026}{275.94}{,}%
\TextAt{8}{427.726}{275.94}{so}%
\TextAt{8}{441.536}{275.94}{that}%
\TextAt{13}{466.345}{275.94}{H}%
\TextAt{17}{476.967}{280.278}{0}%
\TextAt{9}{481.7}{275.94}{\Character{40}}%
\TextAt{13}{486.252}{275.94}{G\Character{59}}%
\TextAt{13}{500.706}{275.94}{V}%
\TextAt{21}{510.155}{280.278}{n}%
\TextAt{9}{515.792}{275.94}{\Character{41}}%
\TextAt{9}{523.647}{275.94}{\Character{61}}%
\TextAt{9}{79.226}{258.894}{1}%
\TextAt{8}{85.079}{258.894}{.}%
\TextAt{8}{96.481}{258.894}{In}%
\TextAt{8}{112.087}{258.894}{the}%
\TextAt{8}{133.228}{258.894}{same}%
\TextAt{8}{163.539}{258.894}{way}%
\TextAt{8}{188.266}{258.894}{we}%
\TextAt{8}{206.478}{258.894}{have}%
\TextAt{8}{234.458}{258.894}{for}%
\TextAt{13}{253.327}{258.894}{a}%
\NativeMathAt{264.471}{258.894}{11.955}{\in}%
\TextAt{13}{277.44}{258.894}{V}%
\TextAt{21}{286.866}{263.232}{n}%
\TextAt{8}{297.384}{258.894}{such}%
\TextAt{8}{324.765}{258.894}{that}%
\TextAt{13}{351.119}{258.894}{N}%
\TextAt{21}{360.493}{257.1}{G}%
\TextAt{13}{367.59}{258.894}{a}%
\TextAt{9}{378.733}{258.894}{\Character{61}}%
\TextAt{9}{392.828}{258.894}{1}%
\TextAt{8}{398.69}{258.894}{,}%
\TextAt{8}{407.07}{258.894}{that}%
\TextAt{13}{433.423}{258.894}{a}%
\TextAt{9}{444.567}{258.894}{\Character{61}}%
\TextAt{13}{458.67}{258.894}{b}%
\TextAt{21}{463.647}{263.232}{\Character{27}}%
\TextAt{3}{468.773}{263.232}{\Character{0}}%
\TextAt{17}{475.359}{263.232}{1}%
\TextAt{8}{480.091}{258.894}{,}%
\TextAt{8}{488.471}{258.894}{for}%
\TextAt{8}{507.341}{258.894}{some}%
\TextAt{13}{79.226}{241.848}{b}%
\NativeMathAt{87.524}{241.848}{11.955}{\in}%
\TextAt{13}{98.815}{241.848}{V}%
\TextAt{21}{108.241}{246.186}{n}%
\TextAt{8}{113.877}{241.848}{,}%
\TextAt{8}{121.026}{241.848}{where}%
\TextAt{13}{154.835}{241.848}{\Character{27}}%
\TextAt{8}{165.815}{241.848}{is}%
\TextAt{8}{177.579}{241.848}{a}%
\TextAt{8}{187.335}{241.848}{generator}%
\TextAt{8}{239.375}{241.848}{of}%
\TextAt{13}{252.705}{241.848}{G}%
\TextAt{8}{261.939}{241.848}{.}%
\TextAt{8}{270.387}{241.848}{Thus}%
\TextAt{13}{300.036}{241.848}{H}%
\TextAt{3}{310.659}{246.186}{\Character{0}}%
\TextAt{17}{317.245}{246.186}{1}%
\TextAt{9}{321.978}{241.848}{\Character{40}}%
\TextAt{13}{326.53}{241.848}{G\Character{59}}%
\TextAt{13}{340.984}{241.848}{V}%
\TextAt{21}{350.433}{246.186}{n}%
\TextAt{9}{356.07}{241.848}{\Character{41}}%
\TextAt{9}{363.925}{241.848}{\Character{61}}%
\TextAt{9}{376.346}{241.848}{1}%
\TextAt{8}{382.221}{241.848}{.}%
\TextAt{8}{390.67}{241.848}{We}%
\TextAt{8}{410.814}{241.848}{now}%
\TextAt{8}{435.203}{241.848}{obtain}%
\TextAt{13}{209.648}{216.89}{h}%
\TextAt{9}{216.386}{216.89}{\Character{40}}%
\TextAt{13}{220.938}{216.89}{G\Character{59}}%
\TextAt{13}{235.392}{216.89}{U}%
\TextAt{21}{243.367}{215.097}{L}%
\TextAt{9}{249.624}{216.89}{\Character{41}}%
\TextAt{9}{257.479}{216.89}{\Character{61}}%
\TextAt{13}{269.922}{216.89}{h}%
\TextAt{9}{276.661}{216.89}{\Character{40}}%
\TextAt{13}{281.213}{216.89}{G\Character{59}}%
\TextAt{13}{295.667}{216.89}{U}%
\TextAt{21}{303.641}{215.097}{L}%
\TextAt{9}{309.898}{216.89}{\Character{47}}%
\TextAt{13}{315.751}{216.89}{V}%
\TextAt{21}{325.177}{221.827}{n}%
\TextAt{9}{330.813}{216.89}{\Character{41}}%
\TextAt{13}{335.366}{216.89}{h}%
\TextAt{9}{342.104}{216.89}{\Character{40}}%
\TextAt{13}{346.657}{216.89}{G\Character{59}}%
\TextAt{13}{361.111}{216.89}{V}%
\TextAt{21}{370.56}{221.827}{n}%
\TextAt{9}{376.196}{216.89}{\Character{41}}%
\TextAt{9}{384.051}{216.89}{\Character{61}}%
\TextAt{9}{396.472}{216.89}{1}%
\TextAt{8}{79.226}{191.933}{because}%
\TextAt{13}{122.528}{191.933}{U}%
\TextAt{21}{130.478}{190.14}{L}%
\TextAt{9}{136.735}{191.933}{\Character{47}}%
\TextAt{13}{142.588}{191.933}{V}%
\TextAt{21}{152.014}{196.271}{n}%
\TextAt{8}{161.548}{191.933}{is}%
\TextAt{8}{173.311}{191.933}{finite.}%
\Vector{q 0 G 0 g 0.399 w 452.334 -582.067 m 452.334 -573.997 l S Q}%
\Vector{q 0 G 0 g 0.399 w 452.534 -574.196 m 459.707 -574.196 l S Q}%
\Vector{q 0 G 0 g 0.399 w 452.534 -581.868 m 459.707 -581.868 l S Q}%
\Vector{q 0 G 0 g 0.399 w 459.906 -582.067 m 459.906 -573.997 l S Q}%
\TextAt{12}{79.226}{159.265}{Corollary}%
\TextAt{12}{138.427}{159.265}{(1.2).}%
\TextAt{8}{173.233}{159.265}{If}%
\TextAt{13}{184.415}{159.265}{L}%
\TextAt{16}{192.38}{159.265}{\Character{106}}%
\TextAt{13}{195.7}{159.265}{K}%
\TextAt{8}{209.906}{159.265}{is}%
\TextAt{8}{221.159}{159.265}{an}%
\TextAt{8}{236.908}{159.265}{unramified}%
\TextAt{8}{294.931}{159.265}{extension}%
\TextAt{8}{346.187}{159.265}{of}%
\TextAt{8}{359.007}{159.265}{local}%
\TextAt{8}{386.137}{159.265}{fields,}%
\TextAt{8}{418.953}{159.265}{then}%
\TextAt{8}{445.103}{159.265}{for}%
\TextAt{13}{462.478}{159.265}{i}%
\TextAt{9}{469.792}{159.265}{\Character{61}}%
\TextAt{9}{482.202}{159.265}{0}%
\TextAt{13}{488.07}{159.265}{\Character{59}}%
\TextAt{16}{493.315}{159.265}{\Character{0}}%
\TextAt{9}{502.613}{159.265}{1}%
\TextAt{8}{508.466}{159.265}{,}%
\TextAt{8}{515.207}{159.265}{one}%
\TextAt{8}{79.226}{142.219}{has}%
\TextAt{13}{132.148}{125.173}{H}%
\TextAt{21}{142.77}{130.109}{i}%
\TextAt{9}{146.152}{125.173}{\Character{40}}%
\TextAt{13}{150.704}{125.173}{G}%
\TextAt{9}{159.938}{125.173}{\Character{40}}%
\TextAt{13}{164.49}{125.173}{L}%
\TextAt{16}{172.455}{125.173}{\Character{106}}%
\TextAt{13}{175.775}{125.173}{K}%
\TextAt{9}{186.593}{125.173}{\Character{41}}%
\TextAt{13}{191.145}{125.173}{\Character{59}}%
\TextAt{13}{196.369}{125.173}{U}%
\TextAt{21}{204.34}{123.38}{L}%
\TextAt{9}{210.597}{125.173}{\Character{41}}%
\TextAt{9}{218.452}{125.173}{\Character{61}}%
\TextAt{9}{230.873}{125.173}{1}%
\TextAt{8}{248.704}{125.173}{and}%
\TextAt{13}{279.524}{125.173}{H}%
\TextAt{21}{290.147}{130.109}{i}%
\TextAt{9}{293.528}{125.173}{\Character{40}}%
\TextAt{13}{298.08}{125.173}{G}%
\TextAt{9}{307.314}{125.173}{\Character{40}}%
\TextAt{13}{311.866}{125.173}{L}%
\TextAt{16}{319.831}{125.173}{\Character{106}}%
\TextAt{13}{323.152}{125.173}{K}%
\TextAt{9}{333.969}{125.173}{\Character{41}}%
\TextAt{13}{338.522}{125.173}{\Character{59}}%
\TextAt{13}{343.746}{125.173}{U}%
\TextAt{17}{352.965}{131.283}{\Character{40}}%
\TextAt{21}{356.258}{131.283}{n}%
\TextAt{17}{361.396}{131.283}{\Character{41}}%
\TextAt{21}{351.716}{121.932}{L}%
\TextAt{9}{365.188}{125.173}{\Character{41}}%
\TextAt{9}{373.043}{125.173}{\Character{61}}%
\TextAt{9}{385.464}{125.173}{1}%
\TextAt{8}{403.295}{125.173}{for}%
\TextAt{13}{421.18}{125.173}{n}%
\TextAt{9}{431.488}{125.173}{\Character{61}}%
\TextAt{9}{443.898}{125.173}{1}%
\TextAt{13}{449.767}{125.173}{\Character{59}}%
\TextAt{9}{455.011}{125.173}{2}%
\TextAt{13}{460.864}{125.173}{\Character{59}}%
\TextAt{13}{466.088}{125.173}{\Character{58}}%
\TextAt{13}{471.325}{125.173}{\Character{58}}%
\TextAt{13}{476.549}{125.173}{\Character{58}}%
\TextAt{8}{79.226}{103.384}{In}%
\TextAt{8}{93.847}{103.384}{particular,}%
\TextAt{13}{204.948}{86.338}{N}%
\TextAt{21}{214.322}{84.479}{L}%
\TextAt{3}{220.081}{84.479}{\Character{106}}%
\TextAt{21}{222.433}{84.479}{K}%
\TextAt{13}{230.636}{86.338}{U}%
\TextAt{21}{238.586}{84.545}{L}%
\TextAt{9}{248.164}{86.338}{\Character{61}}%
\TextAt{13}{260.59}{86.338}{U}%
\TextAt{21}{268.54}{84.545}{K}%
\TextAt{8}{288.698}{86.338}{and}%
\TextAt{13}{319.518}{86.338}{N}%
\TextAt{21}{328.892}{84.479}{L}%
\TextAt{3}{334.651}{84.479}{\Character{106}}%
\TextAt{21}{337.003}{84.479}{K}%
\TextAt{13}{345.206}{86.338}{U}%
\TextAt{17}{354.405}{92.449}{\Character{40}}%
\TextAt{21}{357.698}{92.449}{n}%
\TextAt{17}{362.836}{92.449}{\Character{41}}%
\TextAt{21}{353.156}{83.097}{L}%
\TextAt{9}{369.949}{86.338}{\Character{61}}%
\TextAt{13}{382.374}{86.338}{U}%
\TextAt{17}{391.573}{92.449}{\Character{40}}%
\TextAt{21}{394.866}{92.449}{n}%
\TextAt{17}{400.005}{92.449}{\Character{41}}%
\TextAt{21}{390.324}{83.097}{K}%
\TextAt{13}{403.796}{86.338}{\Character{58}}%
\TextAt{8}{297.211}{43.033}{292}%
\Destination{78.23}{809.815}{nav.page.292}%
\Destination{79.36}{776.155}{nav.print.318}%
\Destination{79.36}{173.255}{nav.number.5.1.2}%
\Hotlink{499.053}{758.955}{13.995}{11.047}{nav.chapter.4}%
\Hotlink{438.812}{674.955}{13.995}{11.047}{nav.chapter.4}%
\Hotlink{139.74}{607.955}{19.22}{11.047}{nav.cite.N093}%
\Hotlink{146.047}{374.755}{13.995}{11.047}{nav.chapter.4}%
\Hotlink{344.18}{758.955}{25.078}{11.047}{nav.number.5.3.5}%
\Hotlink{458.957}{674.955}{25.078}{11.047}{nav.number.4.7.3}%
\Hotlink{166.714}{374.755}{25.078}{11.047}{nav.number.4.7.4}%
\Destination{79.226}{453.177}{errata.130}%
\Hotlink{210.496}{433.183}{27.58}{11.047}{nav.number.2.3.10}%
}

% Page 305 (body page 293)
\Page{%
\TextAt{2}{79.226}{796.943}{V.}%
\TextAt{2}{94.073}{796.943}{Local}%
\TextAt{2}{123.615}{796.943}{Class}%
\TextAt{2}{152.208}{796.943}{Field}%
\TextAt{2}{179.939}{796.943}{Theory}%
\Vector{q 0 G 0 g 0.299 w 7.226 18.729 m 460.769 18.729 l S Q}%
\TextAt{19}{79.226}{761.124}{Proof.}%
\TextAt{8}{116}{761.124}{Let}%
\TextAt{13}{137.524}{761.124}{G}%
\TextAt{9}{151.023}{761.124}{\Character{61}}%
\TextAt{13}{164.393}{761.124}{G}%
\TextAt{9}{173.626}{761.124}{\Character{40}}%
\TextAt{13}{178.179}{761.124}{L}%
\TextAt{16}{186.143}{761.124}{\Character{106}}%
\TextAt{13}{189.464}{761.124}{K}%
\TextAt{9}{200.282}{761.124}{\Character{41}}%
\TextAt{8}{204.834}{761.124}{.}%
\TextAt{8}{214.945}{761.124}{We}%
\TextAt{8}{235.644}{761.124}{have}%
\TextAt{8}{263.193}{761.124}{already}%
\TextAt{8}{305.052}{761.124}{seen}%
\TextAt{8}{331.023}{761.124}{that}%
\TextAt{13}{356.946}{761.124}{H}%
\TextAt{21}{367.569}{765.462}{i}%
\TextAt{9}{370.95}{761.124}{\Character{40}}%
\TextAt{13}{375.503}{761.124}{G\Character{59}}%
\TextAt{13}{389.957}{761.124}{U}%
\TextAt{21}{397.931}{759.33}{L}%
\TextAt{9}{404.188}{761.124}{\Character{41}}%
\TextAt{9}{412.987}{761.124}{\Character{61}}%
\TextAt{9}{426.353}{761.124}{1}%
\TextAt{8}{436.68}{761.124}{in}%
\TextAt{8}{450.887}{761.124}{chap.}%
\TextAt{8}{484.741}{761.124}{IV,}%
\TextAt{8}{505.44}{761.124}{(6.2).}%
\TextAt{8}{79.226}{744.078}{In}%
\TextAt{8}{93.847}{744.078}{order}%
\TextAt{8}{124.417}{744.078}{to}%
\TextAt{8}{138.727}{744.078}{prove}%
\TextAt{13}{170.277}{744.078}{H}%
\TextAt{21}{180.899}{748.416}{i}%
\TextAt{9}{184.281}{744.078}{\Character{40}}%
\TextAt{13}{188.833}{744.078}{G\Character{59}}%
\TextAt{13}{203.287}{744.078}{U}%
\TextAt{17}{212.51}{750.188}{\Character{40}}%
\TextAt{21}{215.803}{750.188}{n}%
\TextAt{17}{220.941}{750.188}{\Character{41}}%
\TextAt{21}{211.261}{740.837}{L}%
\TextAt{9}{224.733}{744.078}{\Character{41}}%
\TextAt{9}{232.588}{744.078}{\Character{61}}%
\TextAt{9}{245.009}{744.078}{1}%
\TextAt{8}{250.884}{744.078}{,}%
\TextAt{8}{258.033}{744.078}{we}%
\TextAt{8}{275.261}{744.078}{first}%
\TextAt{8}{299.386}{744.078}{show}%
\TextAt{8}{328.39}{744.078}{that}%
\TextAt{13}{214.604}{715.845}{H}%
\TextAt{21}{225.227}{720.781}{i}%
\TextAt{9}{228.608}{715.845}{\Character{40}}%
\TextAt{13}{233.16}{715.845}{G\Character{59}}%
\TextAt{13}{247.614}{715.845}{\Character{21}}%
\TextAt{3}{254.466}{720.781}{\Character{3}}%
\TextAt{9}{259.199}{715.845}{\Character{41}}%
\TextAt{9}{267.054}{715.845}{\Character{61}}%
\TextAt{9}{279.475}{715.845}{1}%
\TextAt{8}{297.306}{715.845}{and}%
\TextAt{13}{328.126}{715.845}{H}%
\TextAt{21}{338.749}{720.781}{i}%
\TextAt{9}{342.13}{715.845}{\Character{40}}%
\TextAt{13}{346.682}{715.845}{G\Character{59}}%
\TextAt{13}{361.136}{715.845}{\Character{21}}%
\TextAt{9}{367.988}{715.845}{\Character{41}}%
\TextAt{9}{375.843}{715.845}{\Character{61}}%
\TextAt{9}{388.264}{715.845}{1}%
\TextAt{13}{394.14}{715.845}{\Character{59}}%
\TextAt{8}{79.226}{687.612}{for}%
\TextAt{8}{97.121}{687.612}{the}%
\TextAt{8}{117.287}{687.612}{residue}%
\TextAt{8}{157.024}{687.612}{class}%
\TextAt{8}{184.471}{687.612}{field}%
\TextAt{13}{209.838}{687.612}{\Character{21}}%
\TextAt{8}{220.573}{687.612}{of}%
\TextAt{13}{233.913}{687.612}{L}%
\TextAt{8}{241.878}{687.612}{.}%
\TextAt{8}{250.356}{687.612}{It}%
\TextAt{8}{263.038}{687.612}{is}%
\TextAt{8}{274.812}{687.612}{enough}%
\TextAt{8}{315.146}{687.612}{to}%
\TextAt{8}{329.466}{687.612}{prove}%
\TextAt{8}{361.026}{687.612}{this}%
\TextAt{8}{383.858}{687.612}{for}%
\TextAt{13}{401.753}{687.612}{i}%
\TextAt{9}{409.084}{687.612}{\Character{61}}%
\TextAt{16}{421.526}{687.612}{\Character{0}}%
\TextAt{9}{430.825}{687.612}{1}%
\TextAt{8}{436.677}{687.612}{,}%
\TextAt{8}{443.839}{687.612}{as}%
\TextAt{13}{458.219}{687.612}{\Character{21}}%
\TextAt{8}{468.955}{687.612}{is}%
\TextAt{8}{480.728}{687.612}{finite,}%
\TextAt{8}{513.904}{687.612}{and}%
\TextAt{8}{79.226}{670.566}{so}%
\TextAt{13}{93.785}{670.566}{h}%
\TextAt{9}{100.523}{670.566}{\Character{40}}%
\TextAt{13}{105.076}{670.566}{G\Character{59}}%
\TextAt{13}{119.53}{670.566}{\Character{21}}%
\TextAt{3}{126.382}{674.904}{\Character{3}}%
\TextAt{9}{131.114}{670.566}{\Character{41}}%
\TextAt{9}{139.291}{670.566}{\Character{61}}%
\TextAt{13}{152.055}{670.566}{h}%
\TextAt{9}{158.794}{670.566}{\Character{40}}%
\TextAt{13}{163.346}{670.566}{G\Character{59}}%
\TextAt{13}{177.8}{670.566}{\Character{21}}%
\TextAt{9}{184.652}{670.566}{\Character{41}}%
\TextAt{9}{192.829}{670.566}{\Character{61}}%
\TextAt{9}{205.574}{670.566}{1}%
\TextAt{8}{211.446}{670.566}{.}%
\TextAt{8}{220.46}{670.566}{We}%
\TextAt{8}{240.793}{670.566}{have}%
\TextAt{13}{267.976}{670.566}{H}%
\TextAt{3}{278.599}{674.904}{\Character{0}}%
\TextAt{17}{285.185}{674.904}{1}%
\TextAt{9}{289.918}{670.566}{\Character{40}}%
\TextAt{13}{294.47}{670.566}{G\Character{59}}%
\TextAt{13}{308.924}{670.566}{\Character{21}}%
\TextAt{3}{315.776}{674.904}{\Character{3}}%
\TextAt{9}{320.508}{670.566}{\Character{41}}%
\TextAt{9}{328.685}{670.566}{\Character{61}}%
\TextAt{9}{341.43}{670.566}{1}%
\TextAt{8}{351.388}{670.566}{by}%
\TextAt{8}{367.836}{670.566}{Hilbert}%
\TextAt{8}{408.337}{670.566}{90}%
\TextAt{8}{424.139}{670.566}{(see}%
\TextAt{8}{447.795}{670.566}{chap.}%
\TextAt{8}{480.552}{670.566}{IV,}%
\TextAt{8}{500.885}{670.566}{(3.5)).}%
\TextAt{8}{79.226}{653.52}{Let}%
\TextAt{13}{100.821}{653.52}{f}%
\TextAt{9}{112.253}{653.52}{\Character{61}}%
\TextAt{9}{125.727}{653.52}{\Character{91}}%
\TextAt{13}{128.996}{653.52}{\Character{21}}%
\TextAt{9}{140.21}{653.52}{\Character{58}}%
\TextAt{13}{147.848}{653.52}{\Character{20}}%
\TextAt{9}{154.587}{653.52}{\Character{93}}%
\TextAt{8}{162.361}{653.52}{be}%
\TextAt{8}{178.91}{653.52}{the}%
\TextAt{8}{199.692}{653.52}{degree}%
\TextAt{8}{236.733}{653.52}{of}%
\TextAt{13}{250.688}{653.52}{\Character{21}}%
\TextAt{8}{262.039}{653.52}{over}%
\TextAt{8}{287.71}{653.52}{the}%
\TextAt{8}{308.492}{653.52}{residue}%
\TextAt{8}{348.844}{653.52}{class}%
\TextAt{8}{376.906}{653.52}{field}%
\TextAt{13}{402.888}{653.52}{\Character{20}}%
\TextAt{8}{414.149}{653.52}{of}%
\TextAt{13}{428.104}{653.52}{K}%
\TextAt{8}{438.922}{653.52}{,}%
\TextAt{8}{446.853}{653.52}{and}%
\TextAt{8}{470.24}{653.52}{let}%
\TextAt{13}{487.77}{653.52}{\Character{39}}%
\TextAt{8}{499.961}{653.52}{be}%
\TextAt{8}{516.51}{653.52}{the}%
\TextAt{8}{79.226}{636.474}{Frobenius}%
\TextAt{8}{133.096}{636.474}{automorphism}%
\TextAt{8}{210.566}{636.474}{of}%
\TextAt{13}{223.896}{636.474}{\Character{21}}%
\TextAt{16}{230.724}{636.474}{\Character{106}}%
\TextAt{13}{234.045}{636.474}{\Character{20}}%
\TextAt{8}{240.784}{636.474}{.}%
\TextAt{8}{249.232}{636.474}{Then}%
\TextAt{8}{279.789}{636.474}{we}%
\TextAt{8}{297.017}{636.474}{have}%
\TextAt{8}{175.734}{599.14}{\char35\relax }%
\TextAt{21}{185.49}{597.347}{N}%
\TextAt{22}{192.211}{596.008}{G}%
\TextAt{13}{198.903}{599.14}{\Character{21}}%
\TextAt{9}{209.053}{599.14}{\Character{61}}%
\TextAt{8}{221.478}{599.14}{\char35\relax }%
\TextAt{20}{233.226}{619.584}{\Character{40}}%
\TextAt{13}{242.857}{599.14}{x}%
\NativeMathAt{252.83}{599.14}{11.955}{\in}%
\TextAt{13}{264.121}{599.14}{\Character{21}}%
\TextAt{20}{274.847}{620.062}{\Character{12}}%
\TextAt{20}{274.847}{612.889}{\Character{12}}%
\TextAt{20}{274.847}{605.716}{\Character{12}}%
\TextAt{20}{274.847}{598.543}{\Character{12}}%
\TextAt{20}{274.847}{591.369}{\Character{12}}%
\TextAt{21}{283.48}{614.571}{f}%
\TextAt{3}{288.426}{614.572}{\Character{0}}%
\TextAt{17}{295.013}{614.571}{1}%
\TextAt{20}{282.729}{610.498}{\Character{88}}%
\TextAt{21}{284.511}{585.304}{i}%
\TextAt{17}{287.395}{585.304}{\Character{61}0}%
\TextAt{13}{301.99}{599.14}{x}%
\TextAt{21}{308.642}{604.077}{\Character{39}}%
\TextAt{22}{314.161}{606.889}{i}%
\TextAt{9}{321.141}{599.14}{\Character{61}}%
\TextAt{21}{334.317}{614.572}{f}%
\TextAt{3}{339.264}{614.572}{\Character{0}}%
\TextAt{17}{345.851}{614.572}{1}%
\TextAt{20}{333.567}{610.498}{\Character{88}}%
\TextAt{21}{335.349}{585.304}{i}%
\TextAt{17}{338.232}{585.304}{\Character{61}0}%
\TextAt{13}{352.828}{599.14}{x}%
\TextAt{21}{359.48}{604.077}{q}%
\TextAt{22}{363.547}{606.889}{i}%
\TextAt{9}{370.527}{599.14}{\Character{61}}%
\TextAt{9}{382.937}{599.14}{0}%
\TextAt{20}{388.806}{619.584}{\Character{41}}%
\NativeMathAt{401.757}{599.14}{11.955}{\leq}%
\TextAt{13}{414.377}{599.14}{q}%
\TextAt{21}{419.996}{604.077}{f}%
\TextAt{3}{424.943}{604.077}{\Character{0}}%
\TextAt{17}{431.529}{604.077}{1}%
\TextAt{8}{79.226}{563.628}{and}%
\TextAt{8}{258.646}{546.582}{\char35\relax }%
\TextAt{9}{268.401}{546.582}{\Character{40}}%
\TextAt{13}{272.953}{546.582}{\Character{39}}%
\TextAt{16}{283.278}{546.582}{\Character{0}}%
\TextAt{9}{295.234}{546.582}{1\Character{41}}%
\TextAt{13}{305.639}{546.582}{\Character{21}}%
\TextAt{9}{315.788}{546.582}{\Character{61}}%
\TextAt{13}{328.214}{546.582}{q}%
\TextAt{21}{333.833}{551.518}{f}%
\TextAt{3}{338.78}{551.518}{\Character{0}}%
\TextAt{17}{345.366}{551.518}{1}%
\TextAt{13}{350.098}{546.582}{\Character{59}}%
\TextAt{8}{79.226}{523.389}{since}%
\TextAt{8}{107.895}{523.389}{the}%
\TextAt{8}{128.051}{523.389}{map}%
\TextAt{13}{154.066}{523.389}{\Character{21}}%
\TextAt{21}{166.802}{530.654}{\Character{39}}%
\TextAt{3}{172.321}{530.654}{\Character{0}}%
\TextAt{17}{178.907}{530.654}{1}%
\NativeArrowAt{164.215}{523.389}{11.955}{22.918}{\longrightarrow}%
\TextAt{13}{190.461}{523.389}{\Character{21}}%
\TextAt{8}{201.187}{523.389}{has}%
\TextAt{8}{222.061}{523.389}{kernel}%
\TextAt{13}{256.527}{523.389}{\Character{20}}%
\TextAt{8}{263.266}{523.389}{.}%
\TextAt{8}{271.714}{523.389}{Therefore}%
\TextAt{13}{324.712}{523.389}{H}%
\TextAt{3}{335.334}{527.728}{\Character{0}}%
\TextAt{17}{341.921}{527.728}{1}%
\TextAt{9}{346.653}{523.389}{\Character{40}}%
\TextAt{13}{351.205}{523.389}{G\Character{59}}%
\TextAt{13}{365.659}{523.389}{\Character{21}}%
\TextAt{9}{372.512}{523.389}{\Character{41}}%
\TextAt{9}{380.367}{523.389}{\Character{61}}%
\TextAt{21}{392.81}{521.596}{N}%
\TextAt{22}{399.532}{520.257}{G}%
\TextAt{13}{406.224}{523.389}{\Character{21}}%
\TextAt{9}{413.053}{523.389}{\Character{47}\Character{40}}%
\TextAt{13}{423.458}{523.389}{\Character{39}}%
\TextAt{16}{433.783}{523.389}{\Character{0}}%
\TextAt{9}{445.738}{523.389}{1\Character{41}}%
\TextAt{13}{456.143}{523.389}{\Character{21}}%
\TextAt{9}{466.293}{523.389}{\Character{61}}%
\TextAt{9}{478.702}{523.389}{1}%
\TextAt{8}{484.571}{523.389}{.}%
\TextAt{8}{97.159}{503.355}{Applying}%
\TextAt{8}{146.694}{503.355}{now}%
\TextAt{8}{169.892}{503.355}{the}%
\TextAt{8}{188.858}{503.355}{exact}%
\TextAt{8}{218.559}{503.355}{hexagon}%
\TextAt{8}{263.228}{503.355}{of}%
\TextAt{8}{275.368}{503.355}{chap.}%
\TextAt{8}{307.162}{503.355}{IV,}%
\TextAt{8}{326.116}{503.355}{(7.1),}%
\TextAt{8}{356.39}{503.355}{to}%
\TextAt{8}{369.51}{503.355}{the}%
\TextAt{8}{388.476}{503.355}{exact}%
\TextAt{8}{418.177}{503.355}{sequence}%
\TextAt{8}{465.489}{503.355}{of}%
\TextAt{13}{477.628}{503.355}{G}%
\TextAt{8}{486.862}{503.355}{-modules}%
\TextAt{9}{223.57}{475.122}{1}%
\NativeArrowAt{232.744}{475.122}{11.955}{19.26}{\longrightarrow}%
\TextAt{13}{255.326}{475.122}{U}%
\TextAt{17}{264.525}{481.232}{\Character{40}1\Character{41}}%
\TextAt{21}{263.276}{471.881}{L}%
\NativeArrowAt{279.165}{475.122}{11.955}{19.26}{\longrightarrow}%
\TextAt{13}{301.747}{475.122}{U}%
\TextAt{21}{309.697}{473.328}{L}%
\NativeArrowAt{319.275}{475.122}{11.955}{19.26}{\longrightarrow}%
\TextAt{13}{341.857}{475.122}{\Character{21}}%
\TextAt{3}{348.686}{480.058}{\Character{3}}%
\NativeArrowAt{356.739}{475.122}{11.955}{19.26}{\longrightarrow}%
\TextAt{9}{379.321}{475.122}{1}%
\TextAt{13}{385.174}{475.122}{\Character{59}}%
\TextAt{8}{79.226}{446.889}{we}%
\TextAt{8}{95.761}{446.889}{obtain}%
\TextAt{13}{131.496}{446.889}{H}%
\TextAt{21}{142.119}{451.227}{i}%
\TextAt{9}{145.5}{446.889}{\Character{40}}%
\TextAt{13}{150.052}{446.889}{G\Character{59}}%
\TextAt{13}{164.506}{446.889}{U}%
\TextAt{17}{173.729}{452.999}{\Character{40}1\Character{41}}%
\TextAt{21}{172.48}{443.648}{L}%
\TextAt{9}{185.048}{446.889}{\Character{41}}%
\TextAt{9}{192.903}{446.889}{\Character{61}}%
\TextAt{13}{205.347}{446.889}{H}%
\TextAt{21}{215.969}{451.227}{i}%
\TextAt{9}{219.35}{446.889}{\Character{40}}%
\TextAt{13}{223.903}{446.889}{G\Character{59}}%
\TextAt{13}{238.357}{446.889}{U}%
\TextAt{21}{246.331}{445.095}{L}%
\TextAt{9}{252.588}{446.889}{\Character{41}}%
\TextAt{9}{260.443}{446.889}{\Character{61}}%
\TextAt{9}{272.864}{446.889}{1}%
\TextAt{8}{278.74}{446.889}{,}%
\TextAt{8}{285.335}{446.889}{because}%
\TextAt{13}{327.944}{446.889}{H}%
\TextAt{21}{338.567}{451.227}{i}%
\TextAt{9}{341.948}{446.889}{\Character{40}}%
\TextAt{13}{346.5}{446.889}{G\Character{59}}%
\TextAt{13}{360.954}{446.889}{\Character{21}}%
\TextAt{3}{367.806}{451.227}{\Character{3}}%
\TextAt{9}{372.539}{446.889}{\Character{41}}%
\TextAt{9}{380.394}{446.889}{\Character{61}}%
\TextAt{9}{392.815}{446.889}{1}%
\TextAt{8}{398.69}{446.889}{.}%
\TextAt{8}{406.908}{446.889}{If}%
\TextAt{13}{417.908}{446.889}{\Character{25}}%
\TextAt{8}{428.182}{446.889}{is}%
\TextAt{8}{439.253}{446.889}{a}%
\TextAt{8}{448.316}{446.889}{prime}%
\TextAt{8}{480.787}{446.889}{element}%
\TextAt{8}{523.337}{446.889}{of}%
\TextAt{13}{79.226}{429.843}{K}%
\TextAt{8}{90.044}{429.843}{,}%
\TextAt{8}{96.589}{429.843}{then}%
\TextAt{13}{122.493}{429.843}{\Character{25}}%
\TextAt{8}{132.704}{429.843}{is}%
\TextAt{8}{143.713}{429.843}{also}%
\TextAt{8}{166.437}{429.843}{a}%
\TextAt{8}{175.437}{429.843}{prime}%
\TextAt{8}{207.845}{429.843}{element}%
\TextAt{8}{250.332}{429.843}{of}%
\TextAt{13}{262.906}{429.843}{L}%
\TextAt{8}{270.871}{429.843}{,}%
\TextAt{8}{277.416}{429.843}{so}%
\TextAt{8}{291.03}{429.843}{the}%
\TextAt{8}{310.431}{429.843}{map}%
\TextAt{13}{335.69}{429.843}{U}%
\TextAt{17}{344.889}{435.953}{\Character{40}}%
\TextAt{21}{348.183}{435.953}{n}%
\TextAt{17}{353.321}{435.953}{\Character{41}}%
\TextAt{21}{343.641}{426.602}{L}%
\NativeMathAt{360.433}{429.843}{11.955}{\rightarrow}%
\TextAt{13}{375.709}{429.843}{\Character{21}}%
\TextAt{8}{385.679}{429.843}{given}%
\TextAt{8}{415.493}{429.843}{by}%
\TextAt{9}{430.997}{429.843}{1}%
\TextAt{9}{437.943}{429.843}{\Character{43}}%
\TextAt{13}{448.178}{429.843}{a\Character{25}}%
\TextAt{21}{461.392}{434.181}{n}%
\NativeArrowAt{470.35}{429.843}{11.955}{11.955}{\mapsto}%
\TextAt{13}{485.626}{429.843}{a}%
\TextAt{8}{495.091}{429.843}{mod}%
\TextAt{14}{520.529}{429.843}{p}%
\TextAt{21}{526.512}{428.049}{L}%
\TextAt{8}{79.226}{412.797}{is}%
\TextAt{8}{90.99}{412.797}{a}%
\TextAt{13}{100.745}{412.797}{G}%
\TextAt{8}{109.979}{412.797}{-homomorphism.}%
\TextAt{8}{201.097}{412.797}{From}%
\TextAt{8}{231.822}{412.797}{the}%
\TextAt{8}{251.979}{412.797}{exact}%
\TextAt{8}{282.871}{412.797}{sequence}%
\TextAt{9}{216.467}{384.564}{1}%
\NativeArrowAt{225.641}{384.564}{11.955}{19.26}{\longrightarrow}%
\TextAt{13}{248.223}{384.564}{U}%
\TextAt{17}{257.422}{390.674}{\Character{40}}%
\TextAt{21}{260.715}{390.674}{n}%
\TextAt{17}{265.853}{390.674}{\Character{43}1\Character{41}}%
\TextAt{21}{256.173}{381.323}{L}%
\NativeArrowAt{283.786}{384.564}{11.955}{19.26}{\longrightarrow}%
\TextAt{13}{306.368}{384.564}{U}%
\TextAt{17}{315.567}{390.674}{\Character{40}}%
\TextAt{21}{318.86}{390.674}{n}%
\TextAt{17}{323.999}{390.674}{\Character{41}}%
\TextAt{21}{314.318}{381.323}{L}%
\NativeArrowAt{331.111}{384.564}{11.955}{19.26}{\longrightarrow}%
\TextAt{13}{353.693}{384.564}{\Character{21}}%
\NativeArrowAt{363.842}{384.564}{11.955}{19.26}{\longrightarrow}%
\TextAt{9}{386.424}{384.564}{1}%
\TextAt{13}{392.277}{384.564}{\Character{59}}%
\TextAt{8}{79.226}{356.331}{we}%
\TextAt{8}{96.454}{356.331}{now}%
\TextAt{8}{120.842}{356.331}{deduce}%
\TextAt{8}{159.852}{356.331}{by}%
\TextAt{8}{176.111}{356.331}{induction}%
\TextAt{8}{228.14}{356.331}{just}%
\TextAt{8}{251.285}{356.331}{as}%
\TextAt{8}{265.655}{356.331}{above,}%
\TextAt{8}{302.082}{356.331}{because}%
\TextAt{13}{345.384}{356.331}{H}%
\TextAt{21}{356.007}{360.669}{i}%
\TextAt{9}{359.388}{356.331}{\Character{40}}%
\TextAt{13}{363.94}{356.331}{G\Character{59}}%
\TextAt{13}{378.394}{356.331}{\Character{21}}%
\TextAt{9}{385.247}{356.331}{\Character{41}}%
\TextAt{9}{393.102}{356.331}{\Character{61}}%
\TextAt{9}{405.523}{356.331}{0}%
\TextAt{8}{411.398}{356.331}{,}%
\TextAt{8}{418.547}{356.331}{that}%
\TextAt{13}{221.281}{328.098}{H}%
\TextAt{21}{231.903}{333.034}{i}%
\TextAt{9}{235.285}{328.098}{\Character{40}}%
\TextAt{13}{239.837}{328.098}{G\Character{59}}%
\TextAt{13}{254.291}{328.098}{U}%
\TextAt{17}{263.514}{334.208}{\Character{40}}%
\TextAt{21}{266.807}{334.208}{n}%
\TextAt{17}{271.945}{334.208}{\Character{43}1\Character{41}}%
\TextAt{21}{262.265}{324.857}{L}%
\TextAt{9}{286.557}{328.098}{\Character{41}}%
\TextAt{9}{294.412}{328.098}{\Character{61}}%
\TextAt{13}{306.856}{328.098}{H}%
\TextAt{21}{317.478}{333.034}{i}%
\TextAt{9}{320.86}{328.098}{\Character{40}}%
\TextAt{13}{325.412}{328.098}{G\Character{59}}%
\TextAt{13}{339.866}{328.098}{U}%
\TextAt{17}{349.089}{334.208}{\Character{40}}%
\TextAt{21}{352.382}{334.208}{n}%
\TextAt{17}{357.52}{334.208}{\Character{41}}%
\TextAt{21}{347.84}{324.857}{L}%
\TextAt{9}{361.312}{328.098}{\Character{41}}%
\TextAt{9}{369.167}{328.098}{\Character{61}}%
\TextAt{9}{381.588}{328.098}{1}%
\TextAt{13}{387.463}{328.098}{\Character{59}}%
\TextAt{8}{79.226}{299.865}{since}%
\TextAt{13}{107.895}{299.865}{H}%
\TextAt{21}{118.517}{304.203}{i}%
\TextAt{9}{121.899}{299.865}{\Character{40}}%
\TextAt{13}{126.451}{299.865}{G\Character{59}}%
\TextAt{13}{140.905}{299.865}{U}%
\TextAt{17}{150.128}{305.975}{\Character{40}1\Character{41}}%
\TextAt{21}{148.879}{296.624}{L}%
\TextAt{9}{161.446}{299.865}{\Character{41}}%
\TextAt{9}{169.301}{299.865}{\Character{61}}%
\TextAt{9}{181.722}{299.865}{1}%
\TextAt{8}{187.598}{299.865}{.}%
\Vector{q 0 G 0 g 0.399 w 452.334 -474.135 m 452.334 -466.065 l S Q}%
\Vector{q 0 G 0 g 0.399 w 452.534 -466.265 m 459.707 -466.265 l S Q}%
\Vector{q 0 G 0 g 0.399 w 452.534 -473.936 m 459.707 -473.936 l S Q}%
\Vector{q 0 G 0 g 0.399 w 459.906 -474.135 m 459.906 -466.065 l S Q}%
\TextAt{8}{97.159}{271.083}{We}%
\TextAt{8}{117.796}{271.083}{now}%
\TextAt{8}{142.677}{271.083}{consider}%
\TextAt{8}{188.755}{271.083}{the}%
\TextAt{8}{209.404}{271.083}{maximal}%
\TextAt{8}{257.705}{271.083}{unramified}%
\TextAt{8}{316.73}{271.083}{extension}%
\TextAt{20}{368.913}{274.238}{\Character{101}}%
\TextAt{13}{368.989}{271.083}{k}%
\TextAt{16}{375.478}{271.083}{\Character{106}}%
\TextAt{13}{378.799}{271.083}{k}%
\TextAt{8}{389.679}{271.083}{over}%
\TextAt{8}{415.218}{271.083}{the}%
\TextAt{8}{435.867}{271.083}{ground}%
\TextAt{8}{476.038}{271.083}{field}%
\TextAt{13}{501.888}{271.083}{k}%
\TextAt{8}{508.377}{271.083}{.}%
\TextAt{8}{518.304}{271.083}{By}%
\TextAt{8}{79.226}{254.037}{chap.}%
\TextAt{8}{111.363}{254.037}{II,}%
\TextAt{8}{127.11}{254.037}{§9,}%
\TextAt{8}{145.664}{254.037}{the}%
\TextAt{8}{165.656}{254.037}{residue}%
\TextAt{8}{205.218}{254.037}{class}%
\TextAt{8}{232.49}{254.037}{field}%
\TextAt{8}{257.683}{254.037}{of}%
\TextAt{20}{270.772}{257.192}{\Character{101}}%
\TextAt{13}{270.848}{254.037}{k}%
\TextAt{8}{281.07}{254.037}{is}%
\TextAt{8}{292.669}{254.037}{the}%
\TextAt{8}{312.661}{254.037}{algebraic}%
\TextAt{8}{361.931}{254.037}{closure}%
\TextAt{9}{401.29}{254.037}{\Character{22}}%
\TextAt{13}{400.847}{254.037}{\Character{20}}%
\TextAt{8}{411.318}{254.037}{of}%
\TextAt{8}{424.484}{254.037}{the}%
\TextAt{8}{444.475}{254.037}{residue}%
\TextAt{8}{484.038}{254.037}{class}%
\TextAt{8}{511.31}{254.037}{field}%
\TextAt{13}{79.226}{236.991}{\Character{20}}%
\TextAt{8}{89.862}{236.991}{of}%
\TextAt{13}{103.192}{236.991}{k}%
\TextAt{8}{109.682}{236.991}{.}%
\TextAt{8}{118.13}{236.991}{By}%
\TextAt{8}{136.493}{236.991}{chap.}%
\TextAt{8}{168.684}{236.991}{II,}%
\TextAt{8}{184.597}{236.991}{(9.9),}%
\TextAt{8}{215.824}{236.991}{we}%
\TextAt{8}{233.051}{236.991}{get}%
\TextAt{8}{252.562}{236.991}{a}%
\TextAt{8}{262.317}{236.991}{canonical}%
\TextAt{8}{313.701}{236.991}{isomorphism}%
\TextAt{13}{249.56}{208.758}{G}%
\TextAt{9}{258.793}{208.758}{\Character{40}}%
\TextAt{20}{263.269}{211.913}{\Character{101}}%
\TextAt{13}{263.345}{208.758}{k}%
\TextAt{16}{269.835}{208.758}{\Character{106}}%
\TextAt{13}{273.156}{208.758}{k}%
\TextAt{9}{279.645}{208.758}{\Character{41}}%
\NativeMathAt{287.519}{208.758}{11.955}{\cong}%
\TextAt{13}{300.138}{208.758}{G}%
\TextAt{9}{309.371}{208.758}{\Character{40}}%
\TextAt{9}{314.356}{208.758}{\Character{22}}%
\TextAt{13}{313.924}{208.758}{\Character{20}}%
\TextAt{16}{320.662}{208.758}{\Character{106}}%
\TextAt{13}{323.983}{208.758}{\Character{20}}%
\TextAt{9}{330.722}{208.758}{\Character{41}}%
\NativeMathAt{338.595}{208.758}{11.955}{\cong}%
\TextAt{20}{351.878}{211.847}{\Character{98}}%
\TextAt{18}{351.214}{208.758}{Z}%
\TextAt{13}{359.184}{208.758}{\Character{58}}%
\TextAt{8}{79.226}{180.525}{It}%
\TextAt{8}{92.44}{180.525}{associates}%
\TextAt{8}{146.827}{180.525}{to}%
\TextAt{8}{161.679}{180.525}{the}%
\TextAt{8}{182.377}{180.525}{element}%
\TextAt{9}{226.16}{180.525}{1}%
\NativeMathAt{236.256}{180.525}{11.955}{\in}%
\TextAt{20}{249.134}{183.614}{\Character{98}}%
\TextAt{18}{248.47}{180.525}{Z}%
\TextAt{8}{260.879}{180.525}{the}%
\TextAt{8}{281.576}{180.525}{Frobenius}%
\TextAt{8}{335.988}{180.525}{automorphism}%
\TextAt{13}{413.998}{180.525}{x}%
\NativeArrowAt{424.894}{180.525}{11.955}{11.955}{\mapsto}%
\TextAt{13}{441.092}{180.525}{x}%
\TextAt{21}{447.744}{184.864}{q}%
\TextAt{8}{456.748}{180.525}{in}%
\TextAt{13}{470.942}{180.525}{G}%
\TextAt{9}{480.176}{180.525}{\Character{40}}%
\TextAt{9}{485.161}{180.525}{\Character{22}}%
\TextAt{13}{484.728}{180.525}{\Character{20}}%
\TextAt{16}{491.467}{180.525}{\Character{106}}%
\TextAt{13}{494.787}{180.525}{\Character{20}}%
\TextAt{9}{501.526}{180.525}{\Character{41}}%
\TextAt{8}{506.078}{180.525}{,}%
\TextAt{8}{513.904}{180.525}{and}%
\TextAt{8}{79.226}{163.479}{the}%
\TextAt{8}{99.383}{163.479}{Frobenius}%
\TextAt{8}{153.253}{163.479}{automorphism}%
\TextAt{13}{230.722}{163.479}{\Character{39}}%
\TextAt{21}{238.39}{161.686}{k}%
\TextAt{8}{247.408}{163.479}{in}%
\TextAt{13}{261.06}{163.479}{G}%
\TextAt{9}{270.294}{163.479}{\Character{40}}%
\TextAt{20}{274.77}{166.634}{\Character{101}}%
\TextAt{13}{274.846}{163.479}{k}%
\TextAt{16}{281.336}{163.479}{\Character{106}}%
\TextAt{13}{284.657}{163.479}{k}%
\TextAt{9}{291.146}{163.479}{\Character{41}}%
\TextAt{8}{299.596}{163.479}{which}%
\TextAt{8}{333.082}{163.479}{is}%
\TextAt{8}{344.846}{163.479}{given}%
\TextAt{8}{375.416}{163.479}{by}%
\TextAt{13}{232.311}{135.246}{a}%
\TextAt{21}{238.456}{140.182}{\Character{39}}%
\TextAt{22}{243.975}{138.777}{k}%
\NativeMathAt{252.367}{135.246}{11.955}{\equiv}%
\TextAt{13}{264.986}{135.246}{a}%
\TextAt{21}{271.131}{140.182}{q}%
\TextAt{8}{279.017}{135.246}{mod}%
\TextAt{14}{304.455}{135.246}{p}%
\TextAt{30}{310.396}{133.608}{\Character{101}}%
\TextAt{21}{310.438}{131.505}{k}%
\TextAt{13}{315.557}{135.246}{\Character{59}}%
\TextAt{13}{344.692}{135.246}{a}%
\NativeMathAt{354.178}{135.246}{11.955}{\in}%
\TextAt{14}{365.469}{135.246}{o}%
\TextAt{30}{371.271}{133.608}{\Character{101}}%
\TextAt{21}{371.313}{131.505}{k}%
\TextAt{13}{376.433}{135.246}{\Character{58}}%
\TextAt{8}{79.226}{107.013}{For}%
\TextAt{8}{99.841}{107.013}{the}%
\TextAt{8}{119.643}{107.013}{absolute}%
\TextAt{8}{165.531}{107.013}{Galois}%
\TextAt{8}{201.09}{107.013}{group}%
\TextAt{13}{233.911}{107.013}{G}%
\TextAt{9}{246.466}{107.013}{\Character{61}}%
\TextAt{13}{258.891}{107.013}{G}%
\TextAt{9}{268.125}{107.013}{\Character{40}}%
\TextAt{9}{272.995}{110.168}{\Character{22}}%
\TextAt{13}{272.677}{107.013}{k}%
\TextAt{16}{279.166}{107.013}{\Character{106}}%
\TextAt{13}{282.487}{107.013}{k}%
\TextAt{9}{288.977}{107.013}{\Character{41}}%
\TextAt{8}{297.072}{107.013}{we}%
\TextAt{8}{313.945}{107.013}{therefore}%
\TextAt{8}{362.69}{107.013}{obtain}%
\TextAt{8}{398.763}{107.013}{the}%
\TextAt{8}{418.565}{107.013}{continuous}%
\TextAt{8}{476.815}{107.013}{and}%
\TextAt{8}{499.223}{107.013}{surjec-}%
\TextAt{8}{79.226}{89.967}{tive}%
\TextAt{8}{101.989}{89.967}{homomorphism}%
\TextAt{13}{274.831}{72.921}{d}%
\TextAt{9}{284.234}{72.921}{\Character{58}}%
\TextAt{13}{290.807}{72.921}{G}%
\NativeArrowAt{303.361}{72.921}{11.955}{19.26}{\longrightarrow}%
\TextAt{20}{326.607}{76.01}{\Character{98}}%
\TextAt{18}{325.943}{72.921}{Z}%
\TextAt{13}{333.913}{72.921}{\Character{58}}%
\TextAt{8}{297.211}{43.033}{293}%
\Destination{78.23}{809.815}{nav.page.293}%
\Destination{79.36}{776.155}{nav.print.319}%
\Destination{79.36}{286.055}{nav.print.320}%
\Hotlink{484.241}{758.955}{13.995}{11.047}{nav.chapter.4}%
\Hotlink{480.052}{668.355}{13.995}{11.047}{nav.chapter.4}%
\Hotlink{306.662}{501.155}{13.995}{11.047}{nav.chapter.4}%
\Hotlink{110.863}{251.855}{9.751}{11.047}{nav.chapter.2}%
\Hotlink{132.289}{251.855}{6.858}{11.047}{nav.section.2.9}%
\Hotlink{168.184}{234.755}{9.763}{11.047}{nav.chapter.2}%
\Hotlink{504.94}{758.955}{25.078}{11.047}{nav.number.4.6.2}%
\Hotlink{500.385}{668.355}{25.078}{11.047}{nav.number.4.3.5}%
\Hotlink{325.616}{501.155}{25.078}{11.047}{nav.number.4.7.1}%
\Hotlink{184.097}{234.755}{25.078}{11.047}{nav.number.2.9.9}%
}

% Page 306 (body page 294)
\Page{%
\TextAt{2}{79.226}{796.943}{V.}%
\TextAt{2}{94.073}{796.943}{Local}%
\TextAt{2}{123.615}{796.943}{Class}%
\TextAt{2}{152.208}{796.943}{Field}%
\TextAt{2}{179.939}{796.943}{Theory}%
\Vector{q 0 G 0 g 0.299 w 7.226 18.729 m 460.769 18.729 l S Q}%
\TextAt{8}{79.226}{761.124}{Thus}%
\TextAt{8}{109.382}{761.124}{the}%
\TextAt{8}{130.045}{761.124}{abstract}%
\TextAt{8}{176.149}{761.124}{notions}%
\TextAt{8}{217.699}{761.124}{of}%
\TextAt{8}{231.536}{761.124}{chap.}%
\TextAt{8}{265.248}{761.124}{IV,}%
\TextAt{8}{285.899}{761.124}{§4,}%
\TextAt{8}{305.219}{761.124}{based}%
\TextAt{8}{338.304}{761.124}{on}%
\TextAt{8}{355.07}{761.124}{this}%
\TextAt{8}{378.399}{761.124}{homomorphism,}%
\TextAt{8}{464.955}{761.124}{like}%
\TextAt{8}{486.921}{761.124}{“unrami-}%
\TextAt{8}{79.226}{744.078}{fied”,}%
\TextAt{8}{110.125}{744.078}{“ramification}%
\TextAt{8}{180.711}{744.078}{index”,}%
\TextAt{8}{221.042}{744.078}{“inertia}%
\TextAt{8}{263.653}{744.078}{degree”,}%
\TextAt{8}{308.862}{744.078}{etc.,}%
\TextAt{8}{334.237}{744.078}{do}%
\TextAt{8}{350.511}{744.078}{agree,}%
\TextAt{8}{384.35}{744.078}{in}%
\TextAt{8}{398.018}{744.078}{the}%
\TextAt{8}{418.189}{744.078}{case}%
\TextAt{8}{442.974}{744.078}{at}%
\TextAt{8}{457.299}{744.078}{hand,}%
\TextAt{8}{489.836}{744.078}{with}%
\TextAt{8}{516.51}{744.078}{the}%
\TextAt{8}{79.226}{727.032}{corresponding}%
\TextAt{8}{154.412}{727.032}{concrete}%
\TextAt{8}{200.583}{727.032}{notions}%
\TextAt{8}{241.625}{727.032}{defined}%
\TextAt{8}{281.938}{727.032}{in}%
\TextAt{8}{295.591}{727.032}{chap.}%
\TextAt{8}{327.782}{727.032}{II.}%
\TextAt{8}{97.159}{706.997}{As}%
\TextAt{8}{114.938}{706.997}{stated}%
\TextAt{8}{150.614}{706.997}{above}%
\TextAt{8}{184.282}{706.997}{we}%
\TextAt{8}{202.001}{706.997}{choose}%
\TextAt{13}{239.626}{706.997}{A}%
\TextAt{9}{252.561}{706.997}{\Character{61}}%
\TextAt{9}{266.143}{710.152}{\Character{22}}%
\TextAt{13}{265.825}{706.997}{k}%
\TextAt{3}{272.315}{711.335}{\Character{3}}%
\TextAt{8}{281.436}{706.997}{to}%
\TextAt{8}{296.239}{706.997}{be}%
\TextAt{8}{312.655}{706.997}{our}%
\TextAt{13}{333.961}{706.997}{G}%
\TextAt{8}{343.195}{706.997}{-module.}%
\TextAt{8}{394.412}{706.997}{Hence}%
\TextAt{13}{429.682}{706.997}{A}%
\TextAt{21}{438.457}{705.204}{K}%
\TextAt{9}{450.819}{706.997}{\Character{61}}%
\TextAt{13}{464.083}{706.997}{K}%
\TextAt{3}{474.901}{711.335}{\Character{3}}%
\TextAt{8}{479.633}{706.997}{,}%
\TextAt{8}{487.398}{706.997}{for}%
\TextAt{8}{505.775}{706.997}{every}%
\TextAt{8}{79.226}{689.951}{finite}%
\TextAt{8}{110.071}{689.951}{extension}%
\TextAt{13}{162.769}{689.951}{K}%
\TextAt{16}{173.587}{689.951}{\Character{106}}%
\TextAt{13}{176.908}{689.951}{k}%
\TextAt{8}{183.397}{689.951}{.}%
\TextAt{8}{194.643}{689.951}{The}%
\TextAt{8}{219.63}{689.951}{usual}%
\TextAt{8}{251.191}{689.951}{normalized}%
\TextAt{8}{311.96}{689.951}{exponential}%
\TextAt{8}{375.657}{689.951}{valuation}%
\TextAt{13}{427.662}{689.951}{v}%
\TextAt{21}{433.321}{688.158}{k}%
\TextAt{9}{443.351}{689.951}{\Character{58}}%
\TextAt{13}{451.513}{689.951}{k}%
\TextAt{3}{458.002}{694.289}{\Character{3}}%
\NativeMathAt{467.645}{689.951}{11.955}{\rightarrow}%
\TextAt{18}{484.51}{689.951}{Z}%
\TextAt{8}{497.31}{689.951}{is}%
\TextAt{8}{510.007}{689.951}{then}%
\TextAt{8}{79.226}{672.905}{henselian}%
\TextAt{8}{129.436}{672.905}{with}%
\TextAt{8}{155.52}{672.905}{respect}%
\TextAt{8}{194.994}{672.905}{to}%
\TextAt{13}{208.729}{672.905}{d}%
\TextAt{8}{214.811}{672.905}{,}%
\TextAt{8}{221.5}{672.905}{in}%
\TextAt{8}{234.577}{672.905}{the}%
\TextAt{8}{254.158}{672.905}{sense}%
\TextAt{8}{283.613}{672.905}{of}%
\TextAt{8}{296.367}{672.905}{chap.}%
\TextAt{8}{328.367}{672.905}{IV,}%
\TextAt{8}{347.935}{672.905}{(4.6).}%
\TextAt{8}{380.27}{672.905}{For,}%
\TextAt{8}{404.03}{672.905}{given}%
\TextAt{8}{434.024}{672.905}{any}%
\TextAt{8}{455.565}{672.905}{finite}%
\TextAt{8}{484.901}{672.905}{extension}%
\TextAt{13}{79.226}{655.859}{K}%
\TextAt{16}{90.044}{655.859}{\Character{106}}%
\TextAt{13}{93.365}{655.859}{k}%
\TextAt{8}{99.854}{655.859}{,}%
\TextAt{17}{111.568}{660.566}{1}%
\Vector{q 0 G 0 g 0.478 w 36.199 -115.152 m 47.17 -115.152 l S Q}%
\TextAt{21}{108.199}{651.736}{e}%
\TextAt{22}{112.127}{650.397}{K}%
\TextAt{13}{120.366}{655.859}{v}%
\TextAt{21}{126.025}{654.066}{K}%
\TextAt{8}{138.125}{655.859}{is}%
\TextAt{8}{149.889}{655.859}{the}%
\TextAt{8}{170.045}{655.859}{extension}%
\TextAt{8}{221.811}{655.859}{of}%
\TextAt{13}{235.141}{655.859}{v}%
\TextAt{21}{240.801}{654.066}{k}%
\TextAt{8}{249.818}{655.859}{to}%
\TextAt{13}{264.128}{655.859}{K}%
\TextAt{3}{274.946}{660.197}{\Character{3}}%
\TextAt{8}{279.678}{655.859}{,}%
\TextAt{8}{286.827}{655.859}{and}%
\TextAt{8}{309.59}{655.859}{by}%
\TextAt{8}{325.849}{655.859}{chap.}%
\TextAt{8}{358.04}{655.859}{II,}%
\TextAt{8}{373.953}{655.859}{(4.8),}%
\TextAt{9}{173.413}{630.261}{1}%
\Vector{q 0 G 0 g 0.478 w 97.526 -148.837 m 111.154 -148.837 l S Q}%
\TextAt{13}{169.526}{613.973}{e}%
\TextAt{21}{174.951}{612.18}{K}%
\TextAt{13}{184.35}{622.174}{v}%
\TextAt{21}{190.009}{620.38}{K}%
\TextAt{9}{198.212}{622.174}{\Character{40}}%
\TextAt{13}{202.764}{622.174}{K}%
\TextAt{3}{213.581}{627.11}{\Character{3}}%
\TextAt{9}{218.314}{622.174}{\Character{41}}%
\TextAt{9}{226.169}{622.174}{\Character{61}}%
\TextAt{9}{253.733}{630.262}{1}%
\Vector{q 0 G 0 g 0.478 w 167.808 -148.837 m 201.512 -148.837 l S Q}%
\TextAt{9}{239.808}{613.973}{\Character{91}}%
\TextAt{13}{243.06}{613.973}{K}%
\TextAt{9}{257.198}{613.973}{\Character{58}}%
\TextAt{13}{263.771}{613.973}{k}%
\TextAt{9}{270.26}{613.973}{\Character{93}}%
\TextAt{13}{274.707}{622.174}{v}%
\TextAt{21}{280.367}{620.38}{k}%
\TextAt{9}{285.486}{622.174}{\Character{40}}%
\TextAt{13}{290.039}{622.174}{N}%
\TextAt{21}{299.412}{620.314}{K}%
\TextAt{3}{307.117}{620.314}{\Character{106}}%
\TextAt{21}{309.469}{620.314}{k}%
\TextAt{13}{314.589}{622.174}{K}%
\TextAt{3}{325.407}{627.11}{\Character{3}}%
\TextAt{9}{330.139}{622.174}{\Character{41}}%
\TextAt{9}{337.994}{622.174}{\Character{61}}%
\TextAt{9}{362.507}{630.262}{1}%
\Vector{q 0 G 0 g 0.478 w 279.633 -148.837 m 307.234 -148.837 l S Q}%
\TextAt{13}{351.633}{613.973}{e}%
\TextAt{21}{357.059}{612.18}{K}%
\TextAt{13}{365.261}{613.973}{f}%
\TextAt{21}{371.031}{612.18}{K}%
\TextAt{13}{380.43}{622.174}{v}%
\TextAt{21}{386.089}{620.38}{k}%
\TextAt{9}{391.209}{622.174}{\Character{40}}%
\TextAt{13}{395.761}{622.174}{N}%
\TextAt{21}{405.135}{620.314}{K}%
\TextAt{3}{412.839}{620.314}{\Character{106}}%
\TextAt{21}{415.192}{620.314}{k}%
\TextAt{13}{420.311}{622.174}{K}%
\TextAt{3}{431.129}{627.11}{\Character{3}}%
\TextAt{9}{435.861}{622.174}{\Character{41}}%
\TextAt{13}{440.414}{622.174}{\Character{59}}%
\TextAt{8}{79.226}{589.587}{i.e.,}%
\TextAt{13}{101.331}{589.587}{v}%
\TextAt{21}{106.991}{587.794}{k}%
\TextAt{9}{112.11}{589.587}{\Character{40}}%
\TextAt{13}{116.663}{589.587}{N}%
\TextAt{21}{126.036}{587.727}{K}%
\TextAt{3}{133.741}{587.727}{\Character{106}}%
\TextAt{21}{136.093}{587.727}{k}%
\TextAt{13}{141.213}{589.587}{K}%
\TextAt{3}{152.031}{593.925}{\Character{3}}%
\TextAt{9}{156.763}{589.587}{\Character{41}}%
\TextAt{9}{164.618}{589.587}{\Character{61}}%
\TextAt{13}{177.062}{589.587}{f}%
\TextAt{21}{182.832}{587.794}{K}%
\TextAt{13}{191.034}{589.587}{v}%
\TextAt{21}{196.694}{587.794}{K}%
\TextAt{9}{204.896}{589.587}{\Character{40}}%
\TextAt{13}{209.449}{589.587}{K}%
\TextAt{3}{220.266}{593.925}{\Character{3}}%
\TextAt{9}{224.999}{589.587}{\Character{41}}%
\TextAt{9}{232.854}{589.587}{\Character{61}}%
\TextAt{13}{245.297}{589.587}{f}%
\TextAt{21}{251.067}{587.794}{K}%
\TextAt{18}{259.27}{589.587}{Z}%
\TextAt{8}{267.24}{589.587}{.}%
\TextAt{8}{275.688}{589.587}{The}%
\TextAt{8}{299.742}{589.587}{pair}%
\TextAt{8}{323.808}{589.587}{of}%
\TextAt{8}{337.138}{589.587}{homomorphisms}%
\TextAt{9}{230.421}{561.106}{\Character{40}}%
\TextAt{13}{234.973}{561.106}{d}%
\TextAt{9}{244.377}{561.106}{\Character{58}}%
\TextAt{13}{250.949}{561.106}{G}%
\NativeArrowAt{263.504}{561.106}{11.955}{19.26}{\longrightarrow}%
\TextAt{20}{286.75}{564.194}{\Character{98}}%
\TextAt{18}{286.086}{561.106}{Z}%
\TextAt{13}{294.056}{561.106}{\Character{59}}%
\TextAt{13}{311.236}{561.106}{v}%
\TextAt{21}{316.914}{559.313}{k}%
\TextAt{9}{325.355}{561.106}{\Character{58}}%
\TextAt{13}{331.928}{561.106}{k}%
\TextAt{3}{338.417}{566.042}{\Character{3}}%
\NativeArrowAt{346.47}{561.106}{11.955}{19.26}{\longrightarrow}%
\TextAt{18}{369.052}{561.106}{Z}%
\TextAt{9}{377.022}{561.106}{\Character{41}}%
\TextAt{8}{79.226}{532.625}{therefore}%
\TextAt{8}{130.22}{532.625}{satisfies}%
\TextAt{8}{175.223}{532.625}{all}%
\TextAt{8}{193.376}{532.625}{the}%
\TextAt{8}{215.426}{532.625}{properties}%
\TextAt{8}{272.337}{532.625}{of}%
\TextAt{8}{287.561}{532.625}{a}%
\TextAt{8}{299.21}{532.625}{class}%
\TextAt{8}{328.54}{532.625}{field}%
\TextAt{8}{355.791}{532.625}{theory,}%
\TextAt{8}{397.179}{532.625}{and}%
\TextAt{8}{421.836}{532.625}{we}%
\TextAt{8}{440.956}{532.625}{obtain}%
\TextAt{8}{479.277}{532.625}{the}%
\TextAt{12}{501.327}{532.625}{Local}%
\TextAt{12}{79.226}{515.579}{Reciprocity}%
\TextAt{12}{151.352}{515.579}{Law}%
\TextAt{8}{175.322}{515.579}{:}%
\TextAt{12}{79.226}{489.716}{Theorem}%
\TextAt{12}{135.7}{489.716}{(1.3).}%
\TextAt{8}{170.546}{489.716}{For}%
\TextAt{8}{191.046}{489.716}{every}%
\TextAt{8}{221.468}{489.716}{finite}%
\TextAt{8}{250.911}{489.716}{Galois}%
\TextAt{8}{286.354}{489.716}{extension}%
\TextAt{13}{337.651}{489.716}{L}%
\TextAt{16}{345.615}{489.716}{\Character{106}}%
\TextAt{13}{348.936}{489.716}{K}%
\TextAt{8}{363.182}{489.716}{of}%
\TextAt{8}{376.042}{489.716}{local}%
\TextAt{8}{403.213}{489.716}{fields}%
\TextAt{8}{432.715}{489.716}{we}%
\TextAt{8}{449.473}{489.716}{have}%
\TextAt{8}{475.998}{489.716}{a}%
\TextAt{8}{485.284}{489.716}{canonical}%
\TextAt{8}{79.226}{472.67}{isomorphism}%
\TextAt{13}{223.892}{455.624}{r}%
\TextAt{21}{229.168}{453.764}{L}%
\TextAt{3}{234.927}{453.764}{\Character{106}}%
\TextAt{21}{237.279}{453.764}{K}%
\TextAt{9}{248.803}{455.624}{\Character{58}}%
\TextAt{13}{255.375}{455.624}{G}%
\TextAt{9}{264.609}{455.624}{\Character{40}}%
\TextAt{13}{269.161}{455.624}{L}%
\TextAt{16}{277.126}{455.624}{\Character{106}}%
\TextAt{13}{280.446}{455.624}{K}%
\TextAt{9}{291.264}{455.624}{\Character{41}}%
\TextAt{7}{295.816}{460.56}{ab}%
\NativeMathAt{311.157}{462.401}{7.97}{\sim}%
\NativeArrowAt{308.57}{455.624}{11.955}{13.175}{\longrightarrow}%
\TextAt{13}{325.064}{455.624}{K}%
\TextAt{3}{335.881}{460.56}{\Character{3}}%
\TextAt{9}{340.614}{455.624}{\Character{47}}%
\TextAt{13}{346.467}{455.624}{N}%
\TextAt{21}{355.841}{453.764}{L}%
\TextAt{3}{361.6}{453.764}{\Character{106}}%
\TextAt{21}{363.952}{453.764}{K}%
\TextAt{13}{372.155}{455.624}{L}%
\TextAt{3}{380.119}{460.56}{\Character{3}}%
\TextAt{13}{384.851}{455.624}{\Character{58}}%
\TextAt{8}{97.159}{423.507}{The}%
\TextAt{8}{120.952}{423.507}{general}%
\TextAt{8}{161.016}{423.507}{definition}%
\TextAt{8}{212.784}{423.507}{of}%
\TextAt{8}{225.853}{423.507}{the}%
\TextAt{8}{245.749}{423.507}{reciprocity}%
\TextAt{8}{303.698}{423.507}{map}%
\TextAt{8}{329.452}{423.507}{in}%
\TextAt{8}{342.844}{423.507}{chap.}%
\TextAt{8}{374.948}{423.507}{IV,}%
\TextAt{8}{394.832}{423.507}{(5.6),}%
\TextAt{8}{425.85}{423.507}{was}%
\TextAt{8}{448.089}{423.507}{actually}%
\TextAt{8}{492.385}{423.507}{inspired}%
\TextAt{8}{79.226}{406.461}{by}%
\TextAt{8}{95.206}{406.461}{the}%
\TextAt{8}{115.083}{406.461}{case}%
\TextAt{8}{139.575}{406.461}{of}%
\TextAt{8}{152.625}{406.461}{local}%
\TextAt{8}{179.986}{406.461}{class}%
\TextAt{8}{207.144}{406.461}{field}%
\TextAt{8}{232.222}{406.461}{theory.}%
\TextAt{8}{272.449}{406.461}{This}%
\TextAt{8}{298.89}{406.461}{is}%
\TextAt{8}{310.374}{406.461}{why}%
\TextAt{8}{334.806}{406.461}{it}%
\TextAt{8}{346.231}{406.461}{is}%
\TextAt{8}{357.716}{406.461}{especially}%
\TextAt{8}{410.17}{406.461}{transparent}%
\TextAt{8}{472.727}{406.461}{in}%
\TextAt{8}{486.101}{406.461}{this}%
\TextAt{8}{508.644}{406.461}{case:}%
\TextAt{8}{79.226}{389.415}{let}%
\TextAt{13}{95.71}{389.415}{\Character{27}}%
\NativeMathAt{106.113}{389.415}{11.955}{\in}%
\TextAt{13}{117.404}{389.415}{G}%
\TextAt{9}{126.638}{389.415}{\Character{40}}%
\TextAt{13}{131.19}{389.415}{L}%
\TextAt{16}{139.155}{389.415}{\Character{106}}%
\TextAt{13}{142.475}{389.415}{K}%
\TextAt{9}{153.293}{389.415}{\Character{41}}%
\TextAt{8}{157.845}{389.415}{,}%
\TextAt{8}{164.658}{389.415}{and}%
\TextAt{8}{187}{389.415}{let}%
\TextAt{20}{203.704}{389.415}{\Character{101}}%
\TextAt{13}{203.483}{389.415}{\Character{27}}%
\TextAt{8}{214.042}{389.415}{be}%
\TextAt{8}{229.546}{389.415}{an}%
\TextAt{8}{245.384}{389.415}{extension}%
\TextAt{8}{296.729}{389.415}{of}%
\TextAt{13}{309.638}{389.415}{\Character{27}}%
\TextAt{8}{320.197}{389.415}{to}%
\TextAt{8}{334.086}{389.415}{the}%
\TextAt{8}{353.822}{389.415}{maximal}%
\TextAt{8}{401.21}{389.415}{unramified}%
\TextAt{8}{459.321}{389.415}{extension}%
\TextAt{20}{511.653}{392.437}{\Character{101}}%
\TextAt{13}{510.666}{389.415}{L}%
\TextAt{16}{518.631}{389.415}{\Character{106}}%
\TextAt{13}{521.952}{389.415}{K}%
\TextAt{8}{79.226}{372.369}{of}%
\TextAt{13}{92.184}{372.369}{L}%
\TextAt{8}{103.673}{372.369}{such}%
\TextAt{8}{129.698}{372.369}{that}%
\TextAt{13}{154.695}{372.369}{d}%
\TextAt{21}{160.777}{370.576}{K}%
\TextAt{9}{168.98}{372.369}{\Character{40}}%
\TextAt{20}{173.753}{372.369}{\Character{101}}%
\TextAt{13}{173.532}{372.369}{\Character{27}}%
\TextAt{9}{180.615}{372.369}{\Character{41}}%
\NativeMathAt{188.488}{372.369}{11.955}{\in}%
\TextAt{18}{199.779}{372.369}{N}%
\TextAt{8}{211.938}{372.369}{or,}%
\TextAt{8}{229.203}{372.369}{in}%
\TextAt{8}{242.483}{372.369}{other}%
\TextAt{8}{272.68}{372.369}{words,}%
\TextAt{20}{309.412}{372.369}{\Character{101}}%
\TextAt{13}{309.192}{372.369}{\Character{27}}%
\TextAt{16}{316.275}{372.369}{\Character{106}}%
\TextAt{20}{322.334}{375.391}{\Character{101}}%
\TextAt{13}{319.595}{372.369}{K}%
\TextAt{9}{333.734}{372.369}{\Character{61}}%
\TextAt{13}{346.159}{372.369}{\Character{39}}%
\TextAt{21}{353.828}{376.708}{n}%
\TextAt{21}{353.828}{369.349}{K}%
\TextAt{8}{362.03}{372.369}{,}%
\TextAt{8}{368.882}{372.369}{for}%
\TextAt{8}{386.394}{372.369}{some}%
\TextAt{13}{415.348}{372.369}{n}%
\NativeMathAt{425.657}{372.369}{11.955}{\in}%
\TextAt{18}{436.948}{372.369}{N}%
\TextAt{8}{445.582}{372.369}{.}%
\TextAt{8}{453.906}{372.369}{If}%
\TextAt{9}{465.226}{372.369}{\Character{6}}%
\TextAt{8}{477.205}{372.369}{is}%
\TextAt{8}{488.597}{372.369}{the}%
\TextAt{8}{508.381}{372.369}{fixed}%
\TextAt{8}{79.226}{355.323}{field}%
\TextAt{8}{104.583}{355.323}{of}%
\TextAt{20}{118.133}{355.323}{\Character{101}}%
\TextAt{13}{117.913}{355.323}{\Character{27}}%
\TextAt{8}{128.893}{355.323}{and}%
\TextAt{13}{151.656}{355.323}{\Character{25}}%
\TextAt{17}{158.296}{353.53}{\Character{6}}%
\NativeMathAt{168.231}{355.323}{11.955}{\in}%
\TextAt{9}{179.522}{355.323}{\Character{6}}%
\TextAt{8}{191.874}{355.323}{is}%
\TextAt{8}{203.637}{355.323}{a}%
\TextAt{8}{213.393}{355.323}{prime}%
\TextAt{8}{246.556}{355.323}{element,}%
\TextAt{8}{293.05}{355.323}{then}%
\TextAt{13}{219.837}{326.843}{r}%
\TextAt{21}{225.112}{324.983}{L}%
\TextAt{3}{230.871}{324.983}{\Character{106}}%
\TextAt{21}{233.224}{324.983}{K}%
\TextAt{9}{241.426}{326.843}{\Character{40}}%
\TextAt{13}{245.979}{326.843}{\Character{27}}%
\TextAt{9}{253.061}{326.843}{\Character{41}}%
\TextAt{9}{260.916}{326.843}{\Character{61}}%
\TextAt{13}{273.36}{326.843}{N}%
\TextAt{17}{282.733}{324.983}{\Character{6}}%
\TextAt{3}{288.85}{324.983}{\Character{106}}%
\TextAt{21}{291.202}{324.983}{K}%
\TextAt{9}{299.405}{326.843}{\Character{40}}%
\TextAt{13}{303.957}{326.843}{\Character{25}}%
\TextAt{17}{310.597}{325.049}{\Character{6}}%
\TextAt{9}{317.211}{326.843}{\Character{41}}%
\TextAt{8}{325.084}{326.843}{mod}%
\TextAt{13}{350.522}{326.843}{N}%
\TextAt{21}{359.896}{324.983}{L}%
\TextAt{3}{365.655}{324.983}{\Character{106}}%
\TextAt{21}{368.008}{324.983}{K}%
\TextAt{13}{376.21}{326.843}{L}%
\TextAt{3}{384.175}{331.779}{\Character{3}}%
\TextAt{13}{388.907}{326.843}{\Character{58}}%
\TextAt{8}{79.226}{298.362}{Inverting}%
\TextAt{13}{129.306}{298.362}{r}%
\TextAt{21}{134.582}{296.502}{L}%
\TextAt{3}{140.341}{296.502}{\Character{106}}%
\TextAt{21}{142.693}{296.502}{K}%
\TextAt{8}{154.793}{298.362}{gives}%
\TextAt{8}{183.474}{298.362}{us}%
\TextAt{8}{198.489}{298.362}{the}%
\TextAt{8}{218.646}{298.362}{local}%
\TextAt{12}{246.286}{298.362}{norm}%
\TextAt{12}{281.685}{298.362}{residue}%
\TextAt{12}{327.928}{298.362}{symbol}%
\TextAt{9}{233.587}{269.881}{\Character{40}}%
\TextAt{13}{242.037}{269.881}{\Character{59}}%
\TextAt{13}{247.261}{269.881}{L}%
\TextAt{16}{255.246}{269.881}{\Character{106}}%
\TextAt{13}{258.567}{269.881}{K}%
\TextAt{9}{269.384}{269.881}{\Character{41}}%
\TextAt{9}{277.239}{269.881}{\Character{58}}%
\TextAt{13}{283.83}{269.881}{K}%
\TextAt{3}{294.648}{274.817}{\Character{3}}%
\NativeArrowAt{302.701}{269.881}{11.955}{19.26}{\longrightarrow}%
\TextAt{13}{325.283}{269.881}{G}%
\TextAt{9}{334.516}{269.881}{\Character{40}}%
\TextAt{13}{339.069}{269.881}{L}%
\TextAt{16}{347.033}{269.881}{\Character{106}}%
\TextAt{13}{350.354}{269.881}{K}%
\TextAt{9}{361.172}{269.881}{\Character{41}}%
\TextAt{7}{365.724}{274.817}{ab}%
\TextAt{13}{375.157}{269.881}{\Character{58}}%
\TextAt{8}{79.226}{241.4}{It}%
\TextAt{8}{91.899}{241.4}{is}%
\TextAt{8}{103.663}{241.4}{surjective}%
\TextAt{8}{156.074}{241.4}{and}%
\TextAt{8}{178.837}{241.4}{has}%
\TextAt{8}{199.71}{241.4}{kernel}%
\TextAt{13}{234.177}{241.4}{N}%
\TextAt{21}{243.551}{239.541}{L}%
\TextAt{3}{249.31}{239.541}{\Character{106}}%
\TextAt{21}{251.662}{239.541}{K}%
\TextAt{13}{259.865}{241.4}{L}%
\TextAt{3}{267.829}{245.739}{\Character{3}}%
\TextAt{8}{272.562}{241.4}{.}%
\TextAt{8}{97.159}{221.365}{In}%
\TextAt{8}{112.029}{221.365}{global}%
\TextAt{8}{146.756}{221.365}{class}%
\TextAt{8}{174.442}{221.365}{field}%
\TextAt{8}{200.047}{221.365}{theory}%
\TextAt{8}{237.046}{221.365}{we}%
\TextAt{8}{254.522}{221.365}{will}%
\TextAt{8}{276.876}{221.365}{have}%
\TextAt{8}{304.12}{221.365}{to}%
\TextAt{8}{318.679}{221.365}{take}%
\TextAt{8}{344.296}{221.365}{into}%
\TextAt{8}{368.288}{221.365}{account}%
\TextAt{8}{411.79}{221.365}{the}%
\TextAt{8}{432.195}{221.365}{field}%
\TextAt{18}{457.801}{221.365}{R}%
\TextAt{9}{470.18}{221.365}{\Character{61}}%
\TextAt{18}{483.029}{221.365}{Q}%
\NativeMathAt{492.327}{219.572}{7.97}{\infty}%
\TextAt{8}{505.44}{221.365}{along}%
\TextAt{8}{79.226}{204.32}{with}%
\TextAt{8}{106.44}{204.32}{the}%
\TextAt{14}{127.151}{204.32}{p}%
\TextAt{8}{133.134}{204.32}{-adic}%
\TextAt{8}{162.297}{204.32}{number}%
\TextAt{8}{205.447}{204.32}{fields}%
\TextAt{18}{235.973}{204.32}{Q}%
\TextAt{21}{245.271}{202.526}{p}%
\TextAt{8}{250.032}{204.32}{.}%
\TextAt{8}{260.142}{204.32}{It}%
\TextAt{8}{273.369}{204.32}{also}%
\TextAt{8}{297.403}{204.32}{admits}%
\TextAt{8}{336.393}{204.32}{a}%
\TextAt{8}{346.702}{204.32}{reciprocity}%
\TextAt{8}{405.466}{204.32}{law:}%
\TextAt{8}{432.262}{204.32}{for}%
\TextAt{8}{450.701}{204.32}{the}%
\TextAt{8}{471.411}{204.32}{unique}%
\TextAt{8}{510.007}{204.32}{non-}%
\TextAt{8}{79.226}{187.274}{trivial}%
\TextAt{8}{114.028}{187.274}{Galois}%
\TextAt{8}{149.941}{187.274}{extension}%
\TextAt{18}{201.707}{187.274}{C}%
\TextAt{16}{210.341}{187.274}{\Character{106}}%
\TextAt{18}{213.662}{187.274}{R}%
\TextAt{8}{222.296}{187.274}{,}%
\TextAt{8}{229.446}{187.274}{we}%
\TextAt{8}{246.673}{187.274}{define}%
\TextAt{8}{280.482}{187.274}{the}%
\TextAt{8}{300.639}{187.274}{norm}%
\TextAt{8}{331.208}{187.274}{residue}%
\TextAt{8}{370.935}{187.274}{symbol}%
\TextAt{9}{242.535}{158.793}{\Character{40}}%
\TextAt{13}{250.984}{158.793}{\Character{59}}%
\TextAt{18}{256.229}{158.793}{C}%
\TextAt{16}{264.863}{158.793}{\Character{106}}%
\TextAt{18}{268.184}{158.793}{R}%
\TextAt{9}{276.818}{158.793}{\Character{41}}%
\TextAt{9}{284.673}{158.793}{\Character{58}}%
\TextAt{18}{291.264}{158.793}{R}%
\TextAt{3}{299.898}{163.729}{\Character{3}}%
\NativeArrowAt{307.951}{158.793}{11.955}{19.26}{\longrightarrow}%
\TextAt{13}{330.533}{158.793}{G}%
\TextAt{9}{339.767}{158.793}{\Character{40}}%
\TextAt{18}{344.319}{158.793}{C}%
\TextAt{16}{352.953}{158.793}{\Character{106}}%
\TextAt{18}{356.274}{158.793}{R}%
\TextAt{9}{364.909}{158.793}{\Character{41}}%
\TextAt{8}{79.226}{130.312}{by}%
\TextAt{9}{238.209}{113.266}{\Character{40}}%
\TextAt{13}{242.761}{113.266}{a\Character{59}}%
\TextAt{18}{254.15}{113.266}{C}%
\TextAt{16}{262.785}{113.266}{\Character{106}}%
\TextAt{18}{266.106}{113.266}{R}%
\TextAt{9}{274.74}{113.266}{\Character{41}}%
\TextAt{16}{279.292}{123.241}{\Character{112}}%
\Vector{q 0 G 0 g 0.478 w 217.255 -650.52 m 232.406 -650.52 l S Q}%
\TextAt{16}{289.255}{113.266}{\Character{0}}%
\TextAt{9}{298.553}{113.266}{1}%
\TextAt{9}{307.711}{113.266}{\Character{61}}%
\TextAt{16}{320.153}{123.241}{\Character{112}}%
\Vector{q 0 G 0 g 0.478 w 258.115 -650.52 m 273.267 -650.52 l S Q}%
\TextAt{16}{330.115}{113.266}{\Character{0}}%
\TextAt{9}{339.414}{113.266}{1}%
\TextAt{7}{346.678}{121.043}{sgn}%
\TextAt{17}{358.952}{121.043}{\Character{40}}%
\TextAt{21}{362.246}{121.043}{a}%
\TextAt{17}{366.744}{121.043}{\Character{41}}%
\TextAt{13}{370.535}{113.266}{\Character{58}}%
\TextAt{8}{79.226}{89.967}{The}%
\TextAt{8}{104.172}{89.967}{kernel}%
\TextAt{8}{139.53}{89.967}{of}%
\TextAt{9}{153.752}{89.967}{\Character{40}}%
\TextAt{13}{163.094}{89.967}{\Character{59}}%
\TextAt{18}{168.338}{89.967}{C}%
\TextAt{16}{176.972}{89.967}{\Character{106}}%
\TextAt{18}{180.293}{89.967}{R}%
\TextAt{9}{188.927}{89.967}{\Character{41}}%
\TextAt{8}{198.269}{89.967}{is}%
\TextAt{8}{210.924}{89.967}{the}%
\TextAt{8}{231.973}{89.967}{group}%
\TextAt{18}{266.04}{89.967}{R}%
\TextAt{3}{274.674}{94.306}{\Character{3}}%
\TextAt{17}{274.674}{87.012}{\Character{43}}%
\TextAt{8}{286.548}{89.967}{of}%
\TextAt{8}{300.77}{89.967}{all}%
\TextAt{8}{317.921}{89.967}{positive}%
\TextAt{8}{362.126}{89.967}{real}%
\TextAt{8}{385.78}{89.967}{numbers,}%
\TextAt{8}{437.358}{89.967}{which}%
\TextAt{8}{471.736}{89.967}{is}%
\TextAt{8}{484.392}{89.967}{again}%
\TextAt{8}{516.51}{89.967}{the}%
\TextAt{8}{79.226}{72.921}{group}%
\TextAt{8}{112.402}{72.921}{of}%
\TextAt{8}{125.732}{72.921}{norms}%
\TextAt{13}{160.916}{72.921}{N}%
\TextAt{25}{170.29}{71.062}{C}%
\TextAt{3}{176.046}{71.062}{\Character{106}}%
\TextAt{25}{178.398}{71.062}{R}%
\TextAt{18}{184.652}{72.921}{C}%
\TextAt{3}{193.287}{77.26}{\Character{3}}%
\TextAt{9}{201.34}{72.921}{\Character{61}}%
\TextAt{16}{213.765}{72.921}{\Character{102}}%
\TextAt{13}{219.743}{72.921}{z}%
\TextAt{9}{226.423}{72.921}{\Character{22}}%
\TextAt{13}{225.714}{72.921}{z}%
\TextAt{16}{235.005}{72.921}{\Character{106}}%
\TextAt{13}{241.647}{72.921}{z}%
\NativeMathAt{250.938}{72.921}{11.955}{\in}%
\TextAt{18}{262.229}{72.921}{C}%
\TextAt{3}{270.864}{77.26}{\Character{3}}%
\TextAt{16}{275.596}{72.921}{\Character{103}}%
\TextAt{8}{281.573}{72.921}{.}%
\TextAt{8}{297.211}{43.033}{294}%
\Destination{78.23}{809.815}{nav.page.294}%
\Destination{79.36}{370.355}{nav.print.321}%
\Destination{79.36}{503.755}{nav.number.5.1.3}%
\Hotlink{264.748}{758.955}{13.995}{11.047}{nav.chapter.4}%
\Hotlink{291.078}{758.955}{6.858}{11.047}{nav.section.4.4}%
\Hotlink{327.282}{724.855}{9.763}{11.047}{nav.chapter.2}%
\Hotlink{327.867}{670.755}{13.995}{11.047}{nav.chapter.4}%
\Hotlink{357.54}{653.655}{9.763}{11.047}{nav.chapter.2}%
\Hotlink{374.448}{421.355}{13.995}{11.047}{nav.chapter.4}%
\Hotlink{347.435}{670.755}{25.078}{11.047}{nav.number.4.4.6}%
\Hotlink{373.453}{653.655}{25.078}{11.047}{nav.number.2.4.8}%
\Hotlink{394.332}{421.355}{25.078}{11.047}{nav.number.4.5.6}%
}

% Page 307 (body page 295)
\Page{%
\TextAt{2}{79.226}{796.943}{V.}%
\TextAt{2}{94.073}{796.943}{Local}%
\TextAt{2}{123.615}{796.943}{Class}%
\TextAt{2}{152.208}{796.943}{Field}%
\TextAt{2}{179.939}{796.943}{Theory}%
\Vector{q 0 G 0 g 0.299 w 7.226 18.729 m 460.769 18.729 l S Q}%
\TextAt{8}{97.159}{761.124}{The}%
\TextAt{8}{121.103}{761.124}{reciprocity}%
\TextAt{8}{179.204}{761.124}{law}%
\TextAt{8}{200.231}{761.124}{gives}%
\TextAt{8}{228.802}{761.124}{us}%
\TextAt{8}{243.709}{761.124}{a}%
\TextAt{8}{253.355}{761.124}{very}%
\TextAt{8}{278.937}{761.124}{simple}%
\TextAt{8}{315.303}{761.124}{classification}%
\TextAt{8}{383.613}{761.124}{of}%
\TextAt{8}{396.834}{761.124}{the}%
\TextAt{8}{416.881}{761.124}{abelian}%
\TextAt{8}{457.419}{761.124}{extensions}%
\TextAt{8}{513.691}{761.124}{of}%
\TextAt{8}{526.911}{761.124}{a}%
\TextAt{8}{79.226}{744.078}{local}%
\TextAt{8}{106.866}{744.078}{field}%
\TextAt{13}{132.223}{744.078}{K}%
\TextAt{8}{143.041}{744.078}{.}%
\TextAt{8}{151.489}{744.078}{It}%
\TextAt{8}{164.162}{744.078}{is}%
\TextAt{8}{175.926}{744.078}{formulated}%
\TextAt{8}{235.116}{744.078}{in}%
\TextAt{8}{248.769}{744.078}{the}%
\TextAt{8}{268.925}{744.078}{following}%
\TextAt{12}{79.226}{718.065}{Theorem}%
\TextAt{12}{136.24}{718.065}{(1.4).}%
\TextAt{8}{171.556}{718.065}{The}%
\TextAt{8}{195.61}{718.065}{rule}%
\TextAt{13}{253.966}{701.019}{L}%
\NativeArrowAt{265.251}{701.019}{11.955}{19.26}{\longmapsto}%
\TextAt{16}{287.822}{701.019}{\Character{78}}%
\TextAt{21}{297.642}{699.226}{L}%
\TextAt{9}{307.22}{701.019}{\Character{61}}%
\TextAt{13}{319.645}{701.019}{N}%
\TextAt{21}{329.019}{699.16}{L}%
\TextAt{3}{334.778}{699.16}{\Character{106}}%
\TextAt{21}{337.131}{699.16}{K}%
\TextAt{13}{345.333}{701.019}{L}%
\TextAt{3}{353.298}{705.956}{\Character{3}}%
\TextAt{8}{79.226}{677.498}{gives}%
\TextAt{8}{108.677}{677.498}{a}%
\TextAt{9}{119.203}{677.498}{1}%
\TextAt{8}{125.056}{677.498}{–}%
\TextAt{9}{130.914}{677.498}{1}%
\TextAt{8}{136.767}{677.498}{-correspondence}%
\TextAt{8}{223.113}{677.498}{between}%
\TextAt{8}{269.074}{677.498}{the}%
\TextAt{8}{290.001}{677.498}{finite}%
\TextAt{8}{320.684}{677.498}{abelian}%
\TextAt{8}{362.102}{677.498}{extensions}%
\TextAt{8}{419.253}{677.498}{of}%
\TextAt{8}{433.354}{677.498}{a}%
\TextAt{8}{443.88}{677.498}{local}%
\TextAt{8}{472.291}{677.498}{field}%
\TextAt{13}{498.419}{677.498}{K}%
\TextAt{8}{513.904}{677.498}{and}%
\TextAt{8}{79.226}{660.452}{the}%
\TextAt{8}{99.383}{660.452}{open}%
\TextAt{8}{127.668}{660.452}{subgroups}%
\TextAt{16}{183.081}{660.452}{\Character{78}}%
\TextAt{8}{198.549}{660.452}{of}%
\TextAt{8}{211.879}{660.452}{finite}%
\TextAt{8}{241.791}{660.452}{index}%
\TextAt{8}{273.328}{660.452}{in}%
\TextAt{13}{286.981}{660.452}{K}%
\TextAt{3}{297.799}{664.79}{\Character{3}}%
\TextAt{8}{302.531}{660.452}{.}%
\TextAt{8}{310.98}{660.452}{Furthermore,}%
\TextAt{13}{116.311}{631.451}{L}%
\TextAt{17}{124.275}{629.657}{1}%
\NativeMathAt{132.328}{631.451}{11.955}{\subseteq}%
\TextAt{13}{144.948}{631.451}{L}%
\TextAt{17}{152.912}{629.657}{2}%
\NativeArrowAt{160.965}{631.451}{11.955}{21.926}{\Longleftrightarrow}%
\TextAt{16}{186.202}{631.451}{\Character{78}}%
\TextAt{21}{196.013}{629.657}{L}%
\TextAt{24}{201.772}{628.55}{1}%
\NativeMathAt{209.742}{631.451}{11.955}{\supseteq}%
\TextAt{16}{222.343}{631.451}{\Character{78}}%
\TextAt{21}{232.17}{629.657}{L}%
\TextAt{24}{237.93}{628.55}{2}%
\TextAt{13}{242.579}{631.451}{\Character{59}}%
\TextAt{16}{271.733}{631.451}{\Character{78}}%
\TextAt{21}{281.542}{629.657}{L}%
\TextAt{24}{287.301}{628.55}{1}%
\TextAt{21}{291.452}{629.657}{L}%
\TextAt{24}{297.211}{628.55}{2}%
\TextAt{9}{305.181}{631.451}{\Character{61}}%
\TextAt{16}{317.607}{631.451}{\Character{78}}%
\TextAt{21}{327.416}{629.657}{L}%
\TextAt{24}{333.175}{628.55}{1}%
\TextAt{16}{340.481}{631.451}{\Character{92}}%
\TextAt{16}{351.097}{631.451}{\Character{78}}%
\TextAt{21}{360.917}{629.657}{L}%
\TextAt{24}{366.676}{628.55}{2}%
\TextAt{13}{371.325}{631.451}{\Character{59}}%
\TextAt{16}{400.48}{631.451}{\Character{78}}%
\TextAt{21}{410.289}{629.657}{L}%
\TextAt{24}{416.048}{628.55}{1}%
\TextAt{3}{420.199}{629.657}{\Character{92}}%
\TextAt{21}{425.844}{629.657}{L}%
\TextAt{24}{431.603}{628.55}{2}%
\TextAt{9}{439.573}{631.451}{\Character{61}}%
\TextAt{16}{451.999}{631.451}{\Character{78}}%
\TextAt{21}{461.808}{629.657}{L}%
\TextAt{24}{467.567}{628.55}{1}%
\TextAt{16}{472.216}{631.451}{\Character{78}}%
\TextAt{21}{482.025}{629.657}{L}%
\TextAt{24}{487.784}{628.55}{2}%
\TextAt{13}{492.434}{631.45}{\Character{58}}%
\TextAt{19}{79.226}{584.517}{Proof.}%
\TextAt{8}{116}{584.517}{By}%
\TextAt{8}{134.415}{584.517}{chap.}%
\TextAt{8}{166.763}{584.517}{IV,}%
\TextAt{8}{186.959}{584.517}{(6.7),}%
\TextAt{8}{218.251}{584.517}{all}%
\TextAt{8}{234.562}{584.517}{we}%
\TextAt{8}{251.842}{584.517}{have}%
\TextAt{8}{278.889}{584.517}{to}%
\TextAt{8}{293.251}{584.517}{show}%
\TextAt{8}{322.306}{584.517}{is}%
\TextAt{8}{334.122}{584.517}{that}%
\TextAt{8}{359.543}{584.517}{the}%
\TextAt{8}{379.752}{584.517}{subgroups}%
\TextAt{16}{435.216}{584.517}{\Character{78}}%
\TextAt{8}{450.736}{584.517}{of}%
\TextAt{13}{464.118}{584.517}{K}%
\TextAt{3}{474.936}{588.855}{\Character{3}}%
\TextAt{8}{483.618}{584.517}{which}%
\TextAt{8}{517.156}{584.517}{are}%
\TextAt{8}{79.226}{567.471}{open}%
\TextAt{8}{107.777}{567.471}{in}%
\TextAt{8}{121.696}{567.471}{the}%
\TextAt{8}{142.118}{567.471}{norm}%
\TextAt{8}{172.952}{567.471}{topology}%
\TextAt{8}{221.361}{567.471}{are}%
\TextAt{8}{241.137}{567.471}{precisely}%
\TextAt{8}{289.259}{567.471}{the}%
\TextAt{8}{309.681}{567.471}{subgroups}%
\TextAt{8}{365.359}{567.471}{of}%
\TextAt{8}{378.954}{567.471}{finite}%
\TextAt{8}{409.131}{567.471}{index}%
\TextAt{8}{440.935}{567.471}{which}%
\TextAt{8}{474.686}{567.471}{are}%
\TextAt{8}{494.463}{567.471}{open}%
\TextAt{8}{523.014}{567.471}{in}%
\TextAt{8}{79.226}{550.425}{the}%
\TextAt{8}{99.84}{550.425}{valuation}%
\TextAt{8}{151.37}{550.425}{topology.}%
\TextAt{8}{204.458}{550.425}{A}%
\TextAt{8}{217.587}{550.425}{subgroup}%
\TextAt{16}{268.842}{550.425}{\Character{78}}%
\TextAt{8}{284.768}{550.425}{which}%
\TextAt{8}{318.712}{550.425}{is}%
\TextAt{8}{330.934}{550.425}{open}%
\TextAt{8}{359.678}{550.425}{in}%
\TextAt{8}{373.788}{550.425}{the}%
\TextAt{8}{394.402}{550.425}{norm}%
\TextAt{8}{425.429}{550.425}{topology}%
\TextAt{8}{474.03}{550.425}{contains}%
\TextAt{8}{520.408}{550.425}{by}%
\TextAt{8}{79.226}{533.379}{definition}%
\TextAt{8}{131.161}{533.379}{a}%
\TextAt{8}{140.822}{533.379}{group}%
\TextAt{8}{173.904}{533.379}{of}%
\TextAt{8}{187.14}{533.379}{norms}%
\TextAt{13}{222.23}{533.379}{N}%
\TextAt{21}{231.603}{531.519}{L}%
\TextAt{3}{237.362}{531.519}{\Character{106}}%
\TextAt{21}{239.715}{531.519}{K}%
\TextAt{13}{247.918}{533.379}{L}%
\TextAt{3}{255.882}{537.717}{\Character{3}}%
\TextAt{8}{260.614}{533.379}{.}%
\TextAt{8}{269.031}{533.379}{By}%
\TextAt{8}{287.3}{533.379}{(1.3),}%
\TextAt{8}{318.452}{533.379}{this}%
\TextAt{8}{341.18}{533.379}{has}%
\TextAt{8}{361.96}{533.379}{finite}%
\TextAt{8}{391.778}{533.379}{index}%
\TextAt{8}{423.221}{533.379}{in}%
\TextAt{13}{436.78}{533.379}{K}%
\TextAt{3}{447.598}{537.717}{\Character{3}}%
\TextAt{8}{452.33}{533.379}{.}%
\TextAt{8}{460.747}{533.379}{It}%
\TextAt{8}{473.325}{533.379}{is}%
\TextAt{8}{484.995}{533.379}{also}%
\TextAt{8}{508.381}{533.379}{open}%
\TextAt{8}{79.226}{516.333}{because}%
\TextAt{8}{122.903}{516.333}{it}%
\TextAt{8}{134.981}{516.333}{contains}%
\TextAt{8}{181.276}{516.333}{the}%
\TextAt{8}{201.807}{516.333}{subgroup}%
\TextAt{13}{252.98}{516.333}{N}%
\TextAt{21}{262.353}{514.473}{L}%
\TextAt{3}{268.112}{514.473}{\Character{106}}%
\TextAt{21}{270.465}{514.473}{K}%
\TextAt{13}{278.668}{516.333}{U}%
\TextAt{21}{286.618}{514.54}{L}%
\TextAt{8}{297.147}{516.333}{which}%
\TextAt{8}{331.008}{516.333}{itself}%
\TextAt{8}{359.729}{516.333}{is}%
\TextAt{8}{371.868}{516.333}{open,}%
\TextAt{8}{403.874}{516.333}{for}%
\TextAt{8}{422.133}{516.333}{it}%
\TextAt{8}{434.212}{516.333}{is}%
\TextAt{8}{446.351}{516.333}{closed,}%
\TextAt{8}{484.598}{516.333}{being}%
\TextAt{8}{516.51}{516.333}{the}%
\TextAt{8}{79.226}{499.287}{image}%
\TextAt{8}{113.206}{499.287}{of}%
\TextAt{8}{126.694}{499.287}{the}%
\TextAt{8}{147.009}{499.287}{compact}%
\TextAt{8}{193.996}{499.287}{group}%
\TextAt{13}{227.33}{499.287}{U}%
\TextAt{21}{235.28}{497.494}{L}%
\TextAt{8}{241.537}{499.287}{,}%
\TextAt{8}{248.884}{499.287}{and}%
\TextAt{8}{271.806}{499.287}{has}%
\TextAt{8}{292.838}{499.287}{finite}%
\TextAt{8}{322.908}{499.287}{index}%
\TextAt{8}{354.604}{499.287}{in}%
\TextAt{13}{368.415}{499.287}{U}%
\TextAt{21}{376.366}{497.494}{K}%
\TextAt{8}{384.568}{499.287}{.}%
\TextAt{8}{393.492}{499.287}{We}%
\TextAt{8}{413.795}{499.287}{prove}%
\TextAt{8}{445.503}{499.287}{the}%
\TextAt{8}{465.818}{499.287}{converse}%
\TextAt{8}{512.541}{499.287}{first}%
\TextAt{8}{79.226}{482.241}{in}%
\TextAt{19}{97.159}{462.206}{The}%
\TextAt{19}{122.175}{462.206}{case}%
\TextAt{8}{148.387}{462.206}{char}%
\TextAt{9}{170.504}{462.206}{\Character{40}}%
\TextAt{13}{175.056}{462.206}{K}%
\TextAt{9}{185.874}{462.206}{\Character{41}}%
\TextAt{18}{195.503}{462.206}{\Character{45}}%
\TextAt{13}{203.9}{462.206}{n}%
\TextAt{19}{210.887}{462.206}{.}%
\TextAt{8}{222.761}{462.206}{Let}%
\TextAt{16}{244.76}{462.206}{\Character{78}}%
\TextAt{8}{261.259}{462.206}{be}%
\TextAt{8}{278.213}{462.206}{a}%
\TextAt{8}{288.999}{462.206}{subgroup}%
\TextAt{8}{340.826}{462.206}{of}%
\TextAt{8}{355.186}{462.206}{finite}%
\TextAt{8}{386.128}{462.206}{index}%
\TextAt{13}{418.696}{462.206}{n}%
\TextAt{9}{430.76}{462.206}{\Character{61}}%
\TextAt{9}{444.927}{462.206}{\Character{40}}%
\TextAt{13}{449.493}{462.206}{K}%
\TextAt{3}{460.311}{466.545}{\Character{3}}%
\TextAt{9}{470.12}{462.206}{\Character{58}}%
\TextAt{16}{478.448}{462.206}{\Character{78}}%
\TextAt{9}{490.018}{462.206}{\Character{41}}%
\TextAt{8}{494.571}{462.206}{.}%
\TextAt{8}{506.109}{462.206}{Then}%
\TextAt{13}{79.226}{445.16}{K}%
\TextAt{3}{90.044}{449.499}{\Character{3}}%
\TextAt{21}{94.278}{449.499}{n}%
\NativeMathAt{103.235}{445.16}{11.955}{\subseteq}%
\TextAt{16}{115.836}{445.16}{\Character{78}}%
\TextAt{8}{127.425}{445.16}{,}%
\TextAt{8}{134.531}{445.16}{and}%
\TextAt{8}{157.239}{445.16}{it}%
\TextAt{8}{168.888}{445.16}{is}%
\TextAt{8}{180.597}{445.16}{enough}%
\TextAt{8}{220.867}{445.16}{to}%
\TextAt{8}{235.123}{445.16}{show}%
\TextAt{8}{264.072}{445.16}{that}%
\TextAt{13}{289.386}{445.16}{K}%
\TextAt{3}{300.203}{449.499}{\Character{3}}%
\TextAt{21}{304.438}{449.499}{n}%
\TextAt{8}{313.917}{445.16}{contains}%
\TextAt{8}{359.782}{445.16}{a}%
\TextAt{8}{369.483}{445.16}{group}%
\TextAt{8}{402.603}{445.16}{of}%
\TextAt{8}{415.879}{445.16}{norms.}%
\TextAt{8}{455.596}{445.16}{For}%
\TextAt{8}{476.51}{445.16}{this}%
\TextAt{8}{499.278}{445.16}{we}%
\TextAt{8}{516.451}{445.16}{use}%
\TextAt{8}{79.226}{428.114}{Kummer}%
\TextAt{8}{128.604}{428.114}{theory}%
\TextAt{8}{165.968}{428.114}{(see}%
\TextAt{8}{190.049}{428.114}{chap.}%
\TextAt{8}{224.079}{428.114}{IV,}%
\TextAt{8}{244.837}{428.114}{§3).}%
\TextAt{8}{271.216}{428.114}{We}%
\TextAt{8}{291.974}{428.114}{may}%
\TextAt{8}{317.956}{428.114}{assume}%
\TextAt{8}{359.013}{428.114}{that}%
\TextAt{13}{384.995}{428.114}{K}%
\TextAt{3}{395.813}{432.453}{\Character{3}}%
\TextAt{8}{405.056}{428.114}{contains}%
\TextAt{8}{451.589}{428.114}{the}%
\TextAt{8}{472.358}{428.114}{group}%
\TextAt{13}{506.147}{428.114}{\Character{22}}%
\TextAt{21}{513.19}{426.321}{n}%
\TextAt{8}{523.337}{428.114}{of}%
\TextAt{13}{79.226}{411.068}{n}%
\TextAt{8}{86.214}{411.068}{-th}%
\TextAt{8}{105.761}{411.068}{roots}%
\TextAt{8}{136.116}{411.068}{of}%
\TextAt{8}{150.14}{411.068}{unity.}%
\TextAt{8}{186.362}{411.068}{For}%
\TextAt{8}{208.026}{411.068}{if}%
\TextAt{8}{219.443}{411.068}{it}%
\TextAt{8}{231.842}{411.068}{does}%
\TextAt{8}{258.933}{411.068}{not,}%
\TextAt{8}{283.866}{411.068}{we}%
\TextAt{8}{301.787}{411.068}{put}%
\TextAt{13}{323.941}{411.068}{K}%
\TextAt{17}{333.9}{409.275}{1}%
\TextAt{9}{343.136}{411.068}{\Character{61}}%
\TextAt{13}{356.744}{411.068}{K}%
\TextAt{9}{367.562}{411.068}{\Character{40}}%
\TextAt{13}{372.114}{411.068}{\Character{22}}%
\TextAt{21}{379.157}{409.275}{n}%
\TextAt{9}{384.793}{411.068}{\Character{41}}%
\TextAt{8}{389.346}{411.068}{.}%
\ErrataAt{399.85}{411.068}{135.15}{If $K_1^{*n}$ contains a group of}%
\TextAt{8}{79.226}{394.022}{norms}%
\TextAt{13}{114.41}{394.022}{N}%
\TextAt{21}{123.784}{392.163}{L}%
\TextAt{24}{129.543}{391.056}{1}%
\TextAt{3}{133.694}{392.163}{\Character{106}}%
\TextAt{21}{136.046}{392.163}{K}%
\TextAt{24}{143.188}{391.056}{1}%
\TextAt{13}{147.837}{394.022}{L}%
\TextAt{3}{155.802}{398.361}{\Character{3}}%
\TextAt{17}{155.801}{391.067}{1}%
\TextAt{8}{160.534}{394.022}{,}%
\TextAt{8}{167.683}{394.022}{and}%
\TextAt{13}{190.446}{394.022}{L}%
\TextAt{16}{198.41}{394.022}{\Character{106}}%
\TextAt{13}{201.731}{394.022}{K}%
\TextAt{8}{216.446}{394.022}{is}%
\TextAt{8}{228.21}{394.022}{a}%
\TextAt{8}{237.965}{394.022}{Galois}%
\TextAt{8}{273.879}{394.022}{extension}%
\TextAt{8}{325.645}{394.022}{containing}%
\TextAt{13}{382.563}{394.022}{L}%
\TextAt{17}{390.528}{392.229}{1}%
\TextAt{8}{395.26}{394.022}{,}%
\TextAt{8}{402.409}{394.022}{then}%
\TextAt{13}{186.588}{365.286}{N}%
\TextAt{21}{195.962}{363.427}{L}%
\TextAt{3}{201.721}{363.427}{\Character{106}}%
\TextAt{21}{204.073}{363.427}{K}%
\TextAt{13}{212.276}{365.286}{L}%
\TextAt{3}{220.241}{370.223}{\Character{3}}%
\TextAt{9}{228.294}{365.286}{\Character{61}}%
\TextAt{13}{240.719}{365.286}{N}%
\TextAt{21}{250.093}{363.427}{K}%
\TextAt{24}{257.234}{362.32}{1}%
\TextAt{3}{261.386}{363.427}{\Character{106}}%
\TextAt{21}{263.738}{363.427}{K}%
\TextAt{9}{271.94}{365.286}{\Character{40}}%
\TextAt{13}{276.493}{365.286}{N}%
\TextAt{21}{285.867}{363.427}{L}%
\TextAt{3}{291.626}{363.427}{\Character{106}}%
\TextAt{21}{293.978}{363.427}{K}%
\TextAt{24}{301.119}{362.32}{1}%
\TextAt{13}{305.769}{365.286}{L}%
\TextAt{3}{313.733}{370.223}{\Character{3}}%
\TextAt{9}{318.465}{365.286}{\Character{41}}%
\NativeMathAt{326.339}{365.286}{11.955}{\subseteq}%
\TextAt{13}{338.958}{365.286}{N}%
\TextAt{21}{348.332}{363.427}{K}%
\TextAt{24}{355.473}{362.32}{1}%
\TextAt{3}{359.624}{363.427}{\Character{106}}%
\TextAt{21}{361.976}{363.427}{K}%
\TextAt{9}{370.179}{365.286}{\Character{40}}%
\TextAt{13}{374.731}{365.286}{N}%
\TextAt{21}{384.105}{363.427}{L}%
\TextAt{24}{389.864}{362.32}{1}%
\TextAt{3}{394.015}{363.427}{\Character{106}}%
\TextAt{21}{396.368}{363.427}{K}%
\TextAt{24}{403.509}{362.32}{1}%
\TextAt{13}{408.158}{365.286}{L}%
\TextAt{3}{416.123}{370.223}{\Character{3}}%
\TextAt{17}{416.123}{362.331}{1}%
\TextAt{9}{420.855}{365.286}{\Character{41}}%
\NativeMathAt{228.294}{345.252}{11.955}{\subseteq}%
\TextAt{13}{240.913}{345.252}{N}%
\TextAt{21}{250.287}{343.392}{K}%
\TextAt{24}{257.428}{342.285}{1}%
\TextAt{3}{261.579}{343.392}{\Character{106}}%
\TextAt{21}{263.932}{343.392}{K}%
\TextAt{9}{272.134}{345.252}{\Character{40}}%
\TextAt{13}{276.687}{345.252}{K}%
\TextAt{3}{287.504}{350.188}{\Character{3}}%
\TextAt{21}{291.739}{350.188}{n}%
\TextAt{17}{286.646}{342.296}{1}%
\TextAt{9}{297.375}{345.252}{\Character{41}}%
\NativeMathAt{305.248}{345.252}{11.955}{\subseteq}%
\TextAt{13}{317.867}{345.252}{K}%
\TextAt{3}{328.685}{350.188}{\Character{3}}%
\TextAt{21}{332.919}{350.188}{n}%
\TextAt{13}{338.556}{345.252}{\Character{58}}%
\TextAt{8}{79.226}{316.108}{So}%
\TextAt{8}{95.96}{316.108}{let}%
\TextAt{13}{113.339}{316.108}{\Character{22}}%
\TextAt{21}{120.382}{314.315}{n}%
\NativeMathAt{130.149}{316.108}{11.955}{\subseteq}%
\TextAt{13}{143.577}{316.108}{K}%
\TextAt{8}{154.395}{316.108}{,}%
\TextAt{8}{162.137}{316.108}{and}%
\TextAt{8}{185.375}{316.108}{let}%
\TextAt{13}{202.754}{316.108}{L}%
\TextAt{9}{214.849}{316.108}{\Character{61}}%
\TextAt{13}{228.083}{316.108}{K}%
\TextAt{9}{238.901}{316.108}{\Character{40}}%
\TextAt{22}{246.774}{321.942}{n}%
\TextAt{16}{244.733}{326.23}{\Character{112}}%
\Vector{q 0 G 0 g 0.478 w 182.696 -447.531 m 198.246 -447.531 l S Q}%
\TextAt{13}{254.696}{316.108}{K}%
\TextAt{3}{265.514}{319.562}{\Character{3}}%
\TextAt{9}{270.246}{316.108}{\Character{41}}%
\TextAt{8}{279.171}{316.108}{be}%
\TextAt{8}{295.57}{316.108}{the}%
\TextAt{8}{316.201}{316.108}{maximal}%
\TextAt{8}{364.484}{316.108}{abelian}%
\TextAt{8}{405.607}{316.108}{extension}%
\TextAt{8}{457.847}{316.108}{of}%
\TextAt{8}{471.652}{316.108}{exponent}%
\TextAt{13}{522.53}{316.108}{n}%
\TextAt{8}{529.518}{316.108}{.}%
\TextAt{8}{79.226}{299.062}{Then}%
\TextAt{8}{109.784}{299.062}{by}%
\TextAt{8}{126.043}{299.062}{chap.}%
\TextAt{8}{158.234}{299.062}{IV,}%
\TextAt{8}{178.378}{299.062}{§3,}%
\TextAt{8}{197.064}{299.062}{we}%
\TextAt{8}{214.292}{299.062}{have}%
\TextAt{8}{229.427}{270.061}{Hom}%
\TextAt{9}{253.816}{270.061}{\Character{40}}%
\TextAt{13}{258.368}{270.061}{G}%
\TextAt{9}{267.602}{270.061}{\Character{40}}%
\TextAt{13}{272.154}{270.061}{L}%
\TextAt{16}{280.119}{270.061}{\Character{106}}%
\TextAt{13}{283.44}{270.061}{K}%
\TextAt{9}{294.257}{270.061}{\Character{41}}%
\TextAt{13}{298.81}{270.061}{\Character{59}}%
\TextAt{13}{304.034}{270.061}{\Character{22}}%
\TextAt{21}{311.097}{268.268}{n}%
\TextAt{9}{316.733}{270.061}{\Character{41}}%
\NativeMathAt{324.606}{270.061}{11.955}{\cong}%
\TextAt{13}{337.225}{270.061}{K}%
\TextAt{3}{348.043}{274.997}{\Character{3}}%
\TextAt{9}{352.775}{270.061}{\Character{47}}%
\TextAt{13}{358.628}{270.061}{K}%
\TextAt{3}{369.446}{274.997}{\Character{3}}%
\TextAt{21}{373.68}{274.997}{n}%
\TextAt{13}{379.317}{270.061}{\Character{58}}%
\TextAt{8}{517.802}{270.061}{(*)}%
\TextAt{8}{79.226}{241.06}{By}%
\TextAt{8}{97.34}{241.06}{chap.}%
\TextAt{8}{129.448}{241.06}{II,}%
\TextAt{8}{145.111}{241.06}{(5.8),}%
\TextAt{13}{176.138}{241.06}{K}%
\TextAt{3}{186.956}{245.398}{\Character{3}}%
\TextAt{9}{191.688}{241.06}{\Character{47}}%
\TextAt{13}{197.541}{241.06}{K}%
\TextAt{3}{208.359}{245.398}{\Character{3}}%
\TextAt{21}{212.593}{245.398}{n}%
\TextAt{8}{221.877}{241.06}{is}%
\TextAt{8}{233.392}{241.06}{finite,}%
\TextAt{8}{266.356}{241.06}{and}%
\TextAt{8}{288.869}{241.06}{then}%
\TextAt{8}{315.28}{241.06}{so}%
\TextAt{8}{329.4}{241.06}{is}%
\TextAt{13}{340.915}{241.06}{G}%
\TextAt{9}{350.148}{241.06}{\Character{40}}%
\TextAt{13}{354.701}{241.06}{L}%
\TextAt{16}{362.665}{241.06}{\Character{106}}%
\TextAt{13}{365.986}{241.06}{K}%
\TextAt{9}{376.804}{241.06}{\Character{41}}%
\TextAt{8}{381.356}{241.06}{.}%
\TextAt{8}{389.721}{241.06}{Since}%
\TextAt{13}{420.029}{241.06}{K}%
\TextAt{3}{430.847}{245.398}{\Character{3}}%
\TextAt{9}{435.579}{241.06}{\Character{47}}%
\TextAt{13}{441.432}{241.06}{N}%
\TextAt{21}{450.806}{239.2}{L}%
\TextAt{3}{456.565}{239.2}{\Character{106}}%
\TextAt{21}{458.917}{239.2}{K}%
\TextAt{13}{467.12}{241.06}{L}%
\TextAt{3}{475.084}{245.398}{\Character{3}}%
\TextAt{8}{483.465}{241.06}{is}%
\TextAt{8}{494.979}{241.06}{isomor-}%
\TextAt{8}{79.226}{224.014}{phic}%
\TextAt{8}{104.583}{224.014}{to}%
\TextAt{13}{118.893}{224.014}{G}%
\TextAt{9}{128.127}{224.014}{\Character{40}}%
\TextAt{13}{132.679}{224.014}{L}%
\TextAt{16}{140.644}{224.014}{\Character{106}}%
\TextAt{13}{143.965}{224.014}{K}%
\TextAt{9}{154.782}{224.014}{\Character{41}}%
\TextAt{8}{163.232}{224.014}{and}%
\TextAt{8}{185.995}{224.014}{has}%
\TextAt{8}{206.869}{224.014}{exponent}%
\TextAt{13}{257.272}{224.014}{n}%
\TextAt{8}{264.259}{224.014}{,}%
\TextAt{8}{271.408}{224.014}{we}%
\TextAt{8}{288.636}{224.014}{have}%
\TextAt{8}{315.631}{224.014}{that}%
\TextAt{13}{340.999}{224.014}{K}%
\TextAt{3}{351.817}{228.352}{\Character{3}}%
\TextAt{21}{356.051}{228.352}{n}%
\NativeMathAt{365.008}{224.014}{11.955}{\subseteq}%
\TextAt{13}{377.628}{224.014}{N}%
\TextAt{21}{387.002}{222.154}{L}%
\TextAt{3}{392.761}{222.154}{\Character{106}}%
\TextAt{21}{395.113}{222.154}{K}%
\TextAt{13}{403.316}{224.014}{L}%
\TextAt{3}{411.28}{228.352}{\Character{3}}%
\TextAt{8}{416.012}{224.014}{,}%
\TextAt{8}{423.162}{224.014}{and}%
\TextAt{8}{445.924}{224.014}{(*)}%
\TextAt{8}{464.79}{224.014}{yields}%
\TextAt{8}{202.832}{195.013}{\char35\relax }%
\TextAt{13}{212.588}{195.013}{K}%
\TextAt{3}{223.406}{199.949}{\Character{3}}%
\TextAt{9}{228.138}{195.013}{\Character{47}}%
\TextAt{13}{233.991}{195.013}{K}%
\TextAt{3}{244.809}{199.949}{\Character{3}}%
\TextAt{21}{249.043}{199.949}{n}%
\TextAt{9}{258}{195.013}{\Character{61}}%
\TextAt{8}{270.425}{195.013}{\char35\relax }%
\TextAt{13}{280.181}{195.013}{G}%
\TextAt{9}{289.414}{195.013}{\Character{40}}%
\TextAt{13}{293.967}{195.013}{L}%
\TextAt{16}{301.931}{195.013}{\Character{106}}%
\TextAt{13}{305.252}{195.013}{K}%
\TextAt{9}{316.07}{195.013}{\Character{41}}%
\TextAt{9}{323.925}{195.013}{\Character{61}}%
\TextAt{8}{336.368}{195.013}{\char35\relax }%
\TextAt{13}{346.124}{195.013}{K}%
\TextAt{3}{356.942}{199.949}{\Character{3}}%
\TextAt{9}{361.674}{195.013}{\Character{47}}%
\TextAt{13}{367.527}{195.013}{N}%
\TextAt{21}{376.901}{193.153}{L}%
\TextAt{3}{382.66}{193.153}{\Character{106}}%
\TextAt{21}{385.012}{193.153}{K}%
\TextAt{13}{393.215}{195.013}{L}%
\TextAt{3}{401.179}{199.949}{\Character{3}}%
\TextAt{13}{405.912}{195.013}{\Character{59}}%
\TextAt{8}{79.226}{166.012}{and}%
\TextAt{8}{101.989}{166.012}{therefore}%
\TextAt{13}{151.089}{166.012}{K}%
\TextAt{3}{161.906}{170.35}{\Character{3}}%
\TextAt{21}{166.141}{170.35}{n}%
\TextAt{9}{175.098}{166.012}{\Character{61}}%
\TextAt{13}{187.523}{166.011}{N}%
\TextAt{21}{196.897}{164.152}{L}%
\TextAt{3}{202.656}{164.152}{\Character{106}}%
\TextAt{21}{205.008}{164.152}{K}%
\TextAt{13}{213.211}{166.011}{L}%
\TextAt{3}{221.176}{170.35}{\Character{3}}%
\TextAt{8}{225.908}{166.012}{.}%
\TextAt{19}{97.159}{145.977}{The}%
\TextAt{19}{121.776}{145.977}{case}%
\TextAt{8}{147.589}{145.977}{char}%
\TextAt{9}{169.706}{145.977}{\Character{40}}%
\TextAt{13}{174.259}{145.977}{K}%
\TextAt{9}{185.076}{145.977}{\Character{41}}%
\TextAt{9}{194.054}{145.977}{\Character{61}}%
\TextAt{13}{207.62}{145.977}{p}%
\TextAt{16}{213.495}{145.977}{\Character{106}}%
\TextAt{13}{216.816}{145.977}{n}%
\TextAt{19}{223.803}{145.977}{.}%
\TextAt{8}{234.562}{145.977}{In}%
\TextAt{8}{249.842}{145.977}{this}%
\TextAt{8}{273.323}{145.977}{case}%
\TextAt{8}{298.753}{145.977}{the}%
\TextAt{8}{319.568}{145.977}{proof}%
\TextAt{8}{350.796}{145.977}{will}%
\TextAt{8}{373.559}{145.977}{follow}%
\TextAt{8}{408.039}{145.977}{from}%
\TextAt{8}{436.338}{145.977}{Lubin–Tate}%
\TextAt{8}{499.917}{145.977}{theory}%
\TextAt{8}{79.226}{128.931}{which}%
\TextAt{8}{113.962}{128.931}{we}%
\TextAt{8}{132.438}{128.931}{will}%
\TextAt{8}{155.792}{128.931}{develop}%
\TextAt{8}{199.315}{128.931}{in}%
\TextAt{8}{214.217}{128.931}{§4.}%
\TextAt{8}{237.95}{128.931}{But}%
\TextAt{8}{262.44}{128.931}{it}%
\TextAt{8}{275.393}{128.931}{is}%
\TextAt{8}{288.406}{128.931}{also}%
\TextAt{8}{313.135}{128.931}{possible}%
\TextAt{8}{358.403}{128.931}{to}%
\TextAt{8}{373.962}{128.931}{do}%
\TextAt{8}{391.471}{128.931}{without}%
\TextAt{8}{436.296}{128.931}{this}%
\TextAt{8}{460.368}{128.931}{theory,}%
\TextAt{8}{500.951}{128.931}{at}%
\TextAt{8}{516.51}{128.931}{the}%
\TextAt{8}{79.226}{111.885}{expense}%
\TextAt{8}{123.003}{111.885}{of}%
\TextAt{19}{136.486}{111.885}{ad}%
\TextAt{19}{152.789}{111.885}{hoc}%
\TextAt{8}{173.708}{111.885}{arguments}%
\TextAt{8}{230.839}{111.885}{which}%
\TextAt{8}{264.479}{111.885}{turn}%
\TextAt{8}{290.646}{111.885}{out}%
\TextAt{8}{311.612}{111.885}{to}%
\TextAt{8}{326.075}{111.885}{be}%
\TextAt{8}{342.152}{111.885}{somewhat}%
\TextAt{8}{397}{111.885}{elaborate.}%
\TextAt{8}{453.07}{111.885}{Since}%
\TextAt{8}{483.78}{111.885}{the}%
\TextAt{8}{504.089}{111.885}{result}%
\TextAt{8}{79.226}{94.839}{has}%
\TextAt{8}{99.018}{94.839}{no}%
\TextAt{8}{114.196}{94.839}{further}%
\TextAt{8}{152.458}{94.839}{use}%
\TextAt{8}{171.593}{94.839}{in}%
\TextAt{8}{184.164}{94.839}{the}%
\TextAt{8}{203.239}{94.839}{remainder}%
\TextAt{8}{257.438}{94.839}{of}%
\TextAt{8}{269.686}{94.839}{this}%
\TextAt{8}{291.427}{94.839}{book,}%
\TextAt{8}{322.757}{94.839}{we}%
\TextAt{8}{338.903}{94.839}{simply}%
\TextAt{8}{375.277}{94.839}{refer}%
\TextAt{8}{401.178}{94.839}{the}%
\TextAt{8}{420.253}{94.839}{reader}%
\TextAt{8}{454.941}{94.839}{to}%
\TextAt{8}{468.169}{94.839}{the}%
\TextAt{8}{487.244}{94.839}{beautiful}%
\TextAt{8}{79.226}{77.793}{treatment}%
\TextAt{8}{133.538}{77.793}{in}%
\TextAt{8}{147.191}{77.793}{[122],}%
\TextAt{8}{178.418}{77.793}{chap.}%
\TextAt{8}{210.609}{77.793}{XI,}%
\TextAt{8}{230.754}{77.793}{§5,}%
\TextAt{8}{249.44}{77.793}{and}%
\TextAt{8}{272.203}{77.793}{chap.}%
\TextAt{8}{304.394}{77.793}{XIV,}%
\TextAt{8}{333.313}{77.793}{§6.}%
\Vector{q 0 G 0 g 0.399 w 452.334 -696.207 m 452.334 -688.137 l S Q}%
\Vector{q 0 G 0 g 0.399 w 452.534 -688.337 m 459.707 -688.337 l S Q}%
\Vector{q 0 G 0 g 0.399 w 452.534 -696.008 m 459.707 -696.008 l S Q}%
\Vector{q 0 G 0 g 0.399 w 459.906 -696.207 m 459.906 -688.137 l S Q}%
\TextAt{8}{297.211}{43.033}{295}%
\Destination{522.36}{287.105}{nav.AMS.53}%
\Destination{78.23}{809.815}{nav.page.295}%
\Destination{79.36}{443.155}{nav.print.322}%
\Destination{79.36}{732.055}{nav.number.5.1.4}%
\Hotlink{166.263}{582.355}{13.995}{11.047}{nav.chapter.4}%
\Hotlink{223.579}{425.955}{13.995}{11.047}{nav.chapter.4}%
\Hotlink{250.016}{425.955}{6.858}{11.047}{nav.section.4.3}%
\Hotlink{157.734}{296.855}{13.995}{11.047}{nav.chapter.4}%
\Hotlink{183.557}{296.855}{6.858}{11.047}{nav.section.4.3}%
\Hotlink{128.948}{238.855}{9.751}{11.047}{nav.chapter.2}%
\Hotlink{219.396}{126.755}{6.858}{11.047}{nav.section.5.4}%
\Hotlink{235.933}{75.555}{6.858}{11.047}{nav.section.5.5}%
\Hotlink{338.492}{75.555}{6.858}{11.047}{nav.section.5.6}%
\Hotlink{146.691}{75.555}{25.078}{11.047}{nav.cite.N122}%
\Hotlink{186.459}{582.355}{25.078}{11.047}{nav.number.4.6.7}%
\Hotlink{286.8}{531.155}{25.078}{11.047}{nav.number.5.1.3}%
\Hotlink{144.611}{238.855}{25.078}{11.047}{nav.number.2.5.8}%
\Destination{79.226}{428.068}{errata.129}%
}

% Page 308 (body page 296)
\Page{%
\TextAt{2}{79.226}{796.943}{V.}%
\TextAt{2}{94.073}{796.943}{Local}%
\TextAt{2}{123.615}{796.943}{Class}%
\TextAt{2}{152.208}{796.943}{Field}%
\TextAt{2}{179.939}{796.943}{Theory}%
\Vector{q 0 G 0 g 0.299 w 7.226 18.729 m 460.769 18.729 l S Q}%
\TextAt{8}{97.159}{761.124}{The}%
\TextAt{8}{121.213}{761.124}{proof}%
\TextAt{8}{151.782}{761.124}{also}%
\TextAt{8}{175.262}{761.124}{shows}%
\TextAt{8}{208.88}{761.124}{the}%
\TextAt{8}{229.036}{761.124}{following}%
\TextAt{12}{79.226}{736.259}{Proposition}%
\TextAt{12}{152.06}{736.259}{(1.5).}%
\TextAt{8}{187.481}{736.259}{If}%
\TextAt{13}{199.28}{736.259}{K}%
\TextAt{8}{214.101}{736.259}{contains}%
\TextAt{8}{260.127}{736.259}{the}%
\TextAt{13}{280.39}{736.259}{n}%
\TextAt{8}{287.377}{736.259}{-th}%
\TextAt{8}{306.337}{736.259}{roots}%
\TextAt{8}{336.104}{736.259}{of}%
\TextAt{8}{349.54}{736.259}{unity,}%
\TextAt{8}{382.514}{736.259}{and}%
\TextAt{8}{405.382}{736.259}{if}%
\TextAt{8}{416.212}{736.259}{the}%
\TextAt{8}{436.475}{736.259}{characteristic}%
\TextAt{8}{508.516}{736.259}{of}%
\TextAt{13}{521.952}{736.259}{K}%
\TextAt{8}{79.226}{719.213}{does}%
\TextAt{8}{105.623}{719.213}{not}%
\TextAt{8}{126.437}{719.213}{divide}%
\TextAt{13}{161.227}{719.213}{n}%
\TextAt{8}{168.214}{719.213}{,}%
\TextAt{8}{175.363}{719.213}{then}%
\TextAt{8}{202.023}{719.213}{the}%
\TextAt{8}{222.18}{719.213}{extension}%
\TextAt{13}{273.946}{719.213}{L}%
\TextAt{9}{285.231}{719.213}{\Character{61}}%
\TextAt{13}{297.657}{719.213}{K}%
\TextAt{9}{308.474}{719.213}{\Character{40}}%
\TextAt{22}{316.347}{725.048}{n}%
\TextAt{16}{314.307}{729.335}{\Character{112}}%
\Vector{q 0 G 0 g 0.478 w 252.269 -44.425 m 267.819 -44.425 l S Q}%
\TextAt{13}{324.269}{719.213}{K}%
\TextAt{3}{335.087}{722.667}{\Character{3}}%
\TextAt{9}{339.819}{719.213}{\Character{41}}%
\TextAt{16}{344.372}{719.213}{\Character{106}}%
\TextAt{13}{347.692}{719.213}{K}%
\TextAt{8}{362.408}{719.213}{is}%
\TextAt{8}{374.171}{719.213}{finite,}%
\TextAt{8}{407.335}{719.213}{and}%
\TextAt{8}{430.098}{719.213}{one}%
\TextAt{8}{451.557}{719.213}{has}%
\TextAt{13}{196.338}{694.231}{N}%
\TextAt{21}{205.712}{692.371}{L}%
\TextAt{3}{211.471}{692.371}{\Character{106}}%
\TextAt{21}{213.823}{692.371}{K}%
\TextAt{13}{222.026}{694.231}{L}%
\TextAt{3}{229.991}{699.167}{\Character{3}}%
\TextAt{9}{238.044}{694.231}{\Character{61}}%
\TextAt{13}{250.469}{694.231}{K}%
\TextAt{3}{261.287}{699.167}{\Character{3}}%
\TextAt{21}{265.521}{699.167}{n}%
\TextAt{8}{283.113}{694.231}{and}%
\TextAt{13}{313.933}{694.231}{G}%
\TextAt{9}{323.167}{694.231}{\Character{40}}%
\TextAt{13}{327.719}{694.231}{L}%
\TextAt{16}{335.683}{694.231}{\Character{106}}%
\TextAt{13}{339.004}{694.231}{K}%
\TextAt{9}{349.822}{694.231}{\Character{41}}%
\NativeMathAt{357.695}{694.231}{11.955}{\cong}%
\TextAt{13}{370.315}{694.231}{K}%
\TextAt{3}{381.132}{699.167}{\Character{3}}%
\TextAt{9}{385.865}{694.231}{\Character{47}}%
\TextAt{13}{391.718}{694.231}{K}%
\TextAt{3}{402.535}{699.167}{\Character{3}}%
\TextAt{21}{406.769}{699.167}{n}%
\TextAt{13}{412.406}{694.231}{\Character{58}}%
\TextAt{8}{97.159}{661.43}{Theorem}%
\TextAt{8}{147.578}{661.43}{(1.4)}%
\TextAt{8}{176.549}{661.43}{is}%
\TextAt{8}{189.309}{661.43}{called}%
\TextAt{8}{223.469}{661.43}{the}%
\TextAt{12}{244.622}{661.43}{existence}%
\TextAt{12}{303.487}{661.43}{theorem}%
\TextAt{8}{351.906}{661.43}{,}%
\TextAt{8}{360.3}{661.43}{because}%
\TextAt{8}{404.598}{661.43}{its}%
\TextAt{8}{421.913}{661.43}{essential}%
\TextAt{8}{469.535}{661.43}{statement}%
\TextAt{8}{524.903}{661.43}{is}%
\TextAt{8}{79.226}{644.384}{that,}%
\TextAt{8}{108.126}{644.384}{for}%
\TextAt{8}{126.234}{644.384}{every}%
\TextAt{8}{157.35}{644.384}{open}%
\TextAt{8}{185.859}{644.384}{subgroup}%
\TextAt{16}{236.879}{644.384}{\Character{78}}%
\TextAt{8}{252.571}{644.384}{of}%
\TextAt{8}{266.124}{644.384}{finite}%
\TextAt{8}{296.259}{644.384}{index}%
\TextAt{8}{328.02}{644.384}{in}%
\TextAt{13}{341.896}{644.384}{K}%
\TextAt{3}{352.714}{648.722}{\Character{3}}%
\TextAt{8}{357.446}{644.384}{,}%
\TextAt{8}{364.875}{644.384}{there}%
\TextAt{8}{395.01}{644.384}{exists}%
\TextAt{8}{427.548}{644.384}{an}%
\TextAt{8}{444.03}{644.384}{abelian}%
\TextAt{8}{484.901}{644.384}{extension}%
\TextAt{13}{79.226}{627.338}{L}%
\TextAt{16}{87.191}{627.338}{\Character{106}}%
\TextAt{13}{90.512}{627.338}{K}%
\TextAt{8}{104.345}{627.338}{such}%
\TextAt{8}{129.861}{627.338}{that}%
\TextAt{13}{154.348}{627.338}{N}%
\TextAt{21}{163.722}{625.478}{L}%
\TextAt{3}{169.481}{625.478}{\Character{106}}%
\TextAt{21}{171.833}{625.478}{K}%
\TextAt{13}{180.036}{627.338}{L}%
\TextAt{3}{188}{631.676}{\Character{3}}%
\TextAt{9}{196.054}{627.338}{\Character{61}}%
\TextAt{16}{208.479}{627.338}{\Character{78}}%
\TextAt{8}{220.05}{627.338}{.}%
\TextAt{8}{228.204}{627.338}{This}%
\TextAt{8}{254.043}{627.338}{is}%
\TextAt{8}{264.925}{627.338}{the}%
\TextAt{8}{284.2}{627.338}{“class}%
\TextAt{8}{316.279}{627.338}{field”}%
\TextAt{8}{346.278}{627.338}{of}%
\TextAt{16}{358.726}{627.338}{\Character{78}}%
\TextAt{8}{370.297}{627.338}{.}%
\TextAt{8}{378.452}{627.338}{(Incidentally,}%
\TextAt{8}{448.629}{627.338}{when}%
\TextAt{8}{478.305}{627.338}{char}%
\TextAt{9}{500.422}{627.338}{\Character{40}}%
\TextAt{13}{504.974}{627.338}{K}%
\TextAt{9}{515.792}{627.338}{\Character{41}}%
\TextAt{9}{523.647}{627.338}{\Character{61}}%
\TextAt{9}{79.226}{610.292}{0}%
\TextAt{8}{85.079}{610.292}{,}%
\TextAt{8}{94.423}{610.292}{every}%
\TextAt{8}{127.07}{610.292}{subgroup}%
\TextAt{8}{179.623}{610.292}{of}%
\TextAt{8}{194.709}{610.292}{finite}%
\TextAt{8}{226.376}{610.292}{index}%
\TextAt{8}{259.669}{610.292}{is}%
\TextAt{8}{273.189}{610.292}{automatically}%
\TextAt{8}{348.779}{610.292}{open—see}%
\TextAt{8}{405.54}{610.292}{chap.}%
\TextAt{8}{442.998}{610.292}{II,}%
\TextAt{8}{460.666}{610.292}{(5.7).)}%
\TextAt{8}{503.013}{610.292}{Every}%
\TextAt{8}{79.226}{593.246}{open}%
\TextAt{8}{109.349}{593.246}{subgroup}%
\TextAt{8}{161.983}{593.246}{of}%
\TextAt{13}{177.15}{593.246}{K}%
\TextAt{3}{187.967}{597.584}{\Character{3}}%
\TextAt{8}{198.434}{593.246}{contains}%
\TextAt{8}{246.19}{593.246}{some}%
\TextAt{8}{277.353}{593.246}{higher}%
\TextAt{8}{314.959}{593.246}{unit}%
\TextAt{8}{341.507}{593.246}{group}%
\TextAt{13}{376.52}{593.246}{U}%
\TextAt{17}{385.719}{599.356}{\Character{40}}%
\TextAt{21}{389.012}{599.356}{n}%
\TextAt{17}{394.15}{599.356}{\Character{41}}%
\TextAt{21}{384.47}{590.005}{K}%
\TextAt{8}{403.676}{593.246}{as}%
\TextAt{8}{419.883}{593.246}{these}%
\TextAt{8}{451.691}{593.246}{form}%
\TextAt{8}{481.168}{593.246}{a}%
\TextAt{8}{492.76}{593.246}{basis}%
\TextAt{8}{523.337}{593.246}{of}%
\TextAt{8}{79.226}{576.2}{neighbourhoods}%
\TextAt{8}{163.845}{576.2}{of}%
\TextAt{9}{177.175}{576.2}{1}%
\TextAt{8}{186.925}{576.2}{in}%
\TextAt{13}{200.578}{576.2}{K}%
\TextAt{3}{211.396}{580.538}{\Character{3}}%
\TextAt{8}{216.128}{576.2}{.}%
\TextAt{8}{224.576}{576.2}{We}%
\TextAt{8}{244.721}{576.2}{put}%
\TextAt{13}{266.18}{576.2}{U}%
\TextAt{17}{275.379}{582.31}{\Character{40}0\Character{41}}%
\TextAt{21}{274.131}{572.959}{K}%
\TextAt{9}{290.019}{576.2}{\Character{61}}%
\TextAt{13}{302.445}{576.2}{U}%
\TextAt{21}{310.395}{574.406}{K}%
\TextAt{8}{322.495}{576.2}{and}%
\TextAt{8}{345.258}{576.2}{define:}%
\TextAt{12}{79.226}{551.335}{Definition}%
\TextAt{12}{141.795}{551.335}{(1.6).}%
\TextAt{8}{176.92}{551.335}{Let}%
\TextAt{13}{197.699}{551.335}{L}%
\TextAt{16}{205.664}{551.335}{\Character{106}}%
\TextAt{13}{208.984}{551.335}{K}%
\TextAt{8}{223.509}{551.335}{be}%
\TextAt{8}{239.243}{551.335}{a}%
\TextAt{8}{248.807}{551.335}{finite}%
\TextAt{8}{278.529}{551.335}{abelian}%
\TextAt{8}{318.986}{551.335}{extension,}%
\TextAt{8}{373.851}{551.335}{and}%
\TextAt{13}{396.423}{551.335}{n}%
\TextAt{8}{407.117}{551.335}{the}%
\TextAt{8}{427.083}{551.335}{smallest}%
\TextAt{8}{471.892}{551.335}{number}%
\NativeMathAt{514.297}{551.335}{11.955}{\geq}%
\TextAt{9}{526.916}{551.335}{0}%
\TextAt{8}{79.226}{534.289}{such}%
\TextAt{8}{105.623}{534.289}{that}%
\TextAt{13}{130.992}{534.289}{U}%
\TextAt{17}{140.191}{540.4}{\Character{40}}%
\TextAt{21}{143.484}{540.4}{n}%
\TextAt{17}{148.623}{540.4}{\Character{41}}%
\TextAt{21}{138.942}{531.048}{K}%
\NativeMathAt{155.735}{534.289}{11.955}{\subseteq}%
\TextAt{13}{168.354}{534.289}{N}%
\TextAt{21}{177.728}{532.43}{L}%
\TextAt{3}{183.487}{532.43}{\Character{106}}%
\TextAt{21}{185.839}{532.43}{K}%
\TextAt{13}{194.042}{534.289}{L}%
\TextAt{3}{202.006}{538.628}{\Character{3}}%
\TextAt{8}{206.739}{534.289}{.}%
\TextAt{8}{215.187}{534.289}{Then}%
\TextAt{8}{245.744}{534.289}{the}%
\TextAt{8}{265.901}{534.289}{ideal}%
\TextAt{14}{289.082}{509.307}{f}%
\TextAt{9}{296.303}{509.307}{\Character{61}}%
\TextAt{14}{308.728}{509.307}{p}%
\TextAt{21}{314.711}{514.243}{n}%
\TextAt{21}{314.711}{506.351}{K}%
\TextAt{8}{79.226}{484.324}{is}%
\TextAt{8}{90.99}{484.324}{called}%
\TextAt{8}{124.154}{484.324}{the}%
\TextAt{12}{144.31}{484.324}{conductor}%
\TextAt{8}{206.776}{484.324}{of}%
\TextAt{13}{220.106}{484.324}{L}%
\TextAt{16}{228.07}{484.324}{\Character{106}}%
\TextAt{13}{231.391}{484.324}{K}%
\TextAt{8}{242.209}{484.324}{.}%
\TextAt{12}{79.226}{454.63}{Proposition}%
\TextAt{12}{152.469}{454.63}{(1.7).}%
\TextAt{8}{188.247}{454.63}{A}%
\TextAt{8}{201.382}{454.63}{finite}%
\TextAt{8}{231.756}{454.63}{abelian}%
\TextAt{8}{272.866}{454.63}{extension}%
\TextAt{13}{325.094}{454.63}{L}%
\TextAt{16}{333.058}{454.63}{\Character{106}}%
\TextAt{13}{336.379}{454.63}{K}%
\TextAt{8}{351.557}{454.63}{is}%
\TextAt{8}{363.783}{454.63}{unramified}%
\TextAt{8}{422.778}{454.63}{if}%
\TextAt{8}{433.964}{454.63}{and}%
\TextAt{8}{457.189}{454.63}{only}%
\TextAt{8}{483.343}{454.63}{if}%
\TextAt{8}{494.529}{454.63}{its}%
\TextAt{8}{511.31}{454.63}{con-}%
\TextAt{8}{79.226}{437.584}{ductor}%
\TextAt{8}{116.299}{437.584}{is}%
\TextAt{14}{128.063}{437.584}{f}%
\TextAt{9}{135.284}{437.584}{\Character{61}}%
\TextAt{9}{147.694}{437.584}{1}%
\TextAt{8}{153.563}{437.584}{.}%
\TextAt{19}{79.226}{404.902}{Proof.}%
\TextAt{8}{116}{404.902}{If}%
\TextAt{13}{128.547}{404.902}{L}%
\TextAt{16}{136.511}{404.902}{\Character{106}}%
\TextAt{13}{139.832}{404.902}{K}%
\TextAt{8}{155.401}{404.902}{is}%
\TextAt{8}{168.02}{404.902}{unramified,}%
\TextAt{8}{230.872}{404.902}{then}%
\TextAt{13}{258.386}{404.902}{U}%
\TextAt{21}{266.336}{403.108}{K}%
\TextAt{9}{279.316}{404.902}{\Character{61}}%
\TextAt{13}{293.197}{404.902}{N}%
\TextAt{21}{302.571}{403.042}{L}%
\TextAt{3}{308.33}{403.042}{\Character{106}}%
\TextAt{21}{310.682}{403.042}{K}%
\TextAt{13}{318.885}{404.902}{U}%
\TextAt{21}{326.835}{403.108}{L}%
\TextAt{8}{337.844}{404.902}{by}%
\TextAt{8}{354.957}{404.902}{(1.2),}%
\TextAt{8}{387.252}{404.902}{so}%
\TextAt{8}{402.477}{404.902}{that}%
\TextAt{14}{428.7}{404.902}{f}%
\TextAt{9}{437.377}{404.902}{\Character{61}}%
\TextAt{9}{451.245}{404.902}{1}%
\TextAt{8}{457.111}{404.902}{.}%
\TextAt{8}{468.122}{404.902}{If}%
\TextAt{8}{480.669}{404.902}{conversely}%
\TextAt{14}{79.226}{387.856}{f}%
\TextAt{9}{89.339}{387.856}{\Character{61}}%
\TextAt{9}{104.642}{387.856}{1}%
\TextAt{8}{110.509}{387.856}{,}%
\TextAt{8}{119.779}{387.856}{then}%
\TextAt{13}{148.136}{387.856}{U}%
\TextAt{21}{156.086}{386.062}{K}%
\NativeMathAt{170.502}{387.856}{11.955}{\subseteq}%
\TextAt{13}{186.012}{387.856}{N}%
\TextAt{21}{195.386}{385.996}{L}%
\TextAt{3}{201.145}{385.996}{\Character{106}}%
\TextAt{21}{203.498}{385.996}{K}%
\TextAt{13}{211.7}{387.856}{U}%
\TextAt{21}{219.651}{386.062}{L}%
\TextAt{8}{231.502}{387.856}{and}%
\TextAt{13}{255.961}{387.856}{\Character{25}}%
\TextAt{21}{263.031}{392.194}{n}%
\TextAt{21}{262.602}{384.835}{K}%
\NativeMathAt{277.017}{387.856}{11.955}{\in}%
\TextAt{13}{291.199}{387.856}{N}%
\TextAt{21}{300.573}{385.996}{L}%
\TextAt{3}{306.332}{385.996}{\Character{106}}%
\TextAt{21}{308.685}{385.996}{K}%
\TextAt{13}{316.887}{387.856}{L}%
\TextAt{3}{324.852}{392.194}{\Character{3}}%
\TextAt{8}{329.584}{387.856}{,}%
\TextAt{8}{338.854}{387.856}{for}%
\TextAt{13}{358.436}{387.856}{n}%
\TextAt{9}{371.636}{387.856}{\Character{61}}%
\TextAt{9}{386.939}{387.856}{\Character{40}}%
\TextAt{13}{391.505}{387.856}{K}%
\TextAt{3}{402.323}{392.194}{\Character{3}}%
\TextAt{9}{413.268}{387.856}{\Character{58}}%
\TextAt{13}{422.732}{387.856}{N}%
\TextAt{21}{432.106}{385.996}{L}%
\TextAt{3}{437.865}{385.996}{\Character{106}}%
\TextAt{21}{440.217}{385.996}{K}%
\TextAt{13}{448.42}{387.856}{L}%
\TextAt{3}{456.384}{392.194}{\Character{3}}%
\TextAt{9}{461.117}{387.856}{\Character{41}}%
\TextAt{8}{465.669}{387.856}{.}%
\TextAt{8}{479.208}{387.856}{If}%
\TextAt{13}{492.597}{387.856}{M}%
\TextAt{16}{505.17}{387.856}{\Character{106}}%
\TextAt{13}{508.491}{387.856}{K}%
\TextAt{8}{524.903}{387.856}{is}%
\TextAt{8}{79.226}{370.81}{the}%
\TextAt{8}{99.884}{370.81}{unramified}%
\TextAt{8}{158.919}{370.81}{extension}%
\TextAt{8}{211.186}{370.81}{of}%
\TextAt{8}{225.018}{370.81}{degree}%
\TextAt{13}{261.935}{370.81}{n}%
\TextAt{8}{268.923}{370.81}{,}%
\TextAt{8}{276.699}{370.81}{then}%
\TextAt{13}{303.861}{370.81}{N}%
\TextAt{21}{313.235}{368.95}{M}%
\TextAt{3}{322.217}{368.95}{\Character{106}}%
\TextAt{21}{324.569}{368.95}{K}%
\TextAt{13}{332.772}{370.81}{M}%
\TextAt{3}{345.345}{375.148}{\Character{3}}%
\TextAt{9}{354.254}{370.81}{\Character{61}}%
\TextAt{9}{367.524}{370.81}{\Character{40}}%
\TextAt{13}{372.087}{370.81}{\Character{25}}%
\TextAt{21}{379.156}{375.148}{n}%
\TextAt{21}{378.727}{367.789}{K}%
\TextAt{9}{386.93}{370.81}{\Character{41}}%
\TextAt{16}{394.481}{370.81}{\Character{2}}%
\TextAt{13}{406.778}{370.81}{U}%
\TextAt{21}{414.728}{369.016}{K}%
\NativeMathAt{427.107}{370.81}{11.955}{\subseteq}%
\TextAt{13}{440.581}{370.81}{N}%
\TextAt{21}{449.955}{368.95}{L}%
\TextAt{3}{455.714}{368.95}{\Character{106}}%
\TextAt{21}{458.066}{368.95}{K}%
\TextAt{13}{466.269}{370.81}{L}%
\TextAt{3}{474.234}{375.148}{\Character{3}}%
\TextAt{8}{478.966}{370.81}{,}%
\TextAt{8}{486.742}{370.81}{and}%
\TextAt{8}{510.007}{370.81}{then}%
\TextAt{13}{79.226}{353.764}{M}%
\NativeMathAt{95.121}{353.764}{11.955}{\supseteq}%
\TextAt{13}{107.74}{353.764}{L}%
\TextAt{8}{115.704}{353.764}{,}%
\TextAt{8}{122.854}{353.764}{i.e.,}%
\TextAt{13}{144.959}{353.764}{L}%
\TextAt{16}{152.923}{353.764}{\Character{106}}%
\TextAt{13}{156.244}{353.764}{K}%
\TextAt{8}{170.959}{353.764}{is}%
\TextAt{8}{182.723}{353.764}{unramified.}%
\Vector{q 0 G 0 g 0.399 w 452.334 -420.236 m 452.334 -412.167 l S Q}%
\Vector{q 0 G 0 g 0.399 w 452.534 -412.366 m 459.707 -412.366 l S Q}%
\Vector{q 0 G 0 g 0.399 w 452.534 -420.037 m 459.707 -420.037 l S Q}%
\Vector{q 0 G 0 g 0.399 w 459.906 -420.236 m 459.906 -412.167 l S Q}%
\TextAt{8}{97.159}{325.911}{Every}%
\TextAt{8}{131.831}{325.911}{open}%
\TextAt{8}{161.136}{325.911}{subgroup}%
\TextAt{16}{212.952}{325.911}{\Character{78}}%
\TextAt{8}{229.439}{325.911}{of}%
\TextAt{8}{243.787}{325.911}{finite}%
\TextAt{8}{274.718}{325.911}{index}%
\TextAt{8}{307.274}{325.911}{in}%
\TextAt{13}{321.946}{325.911}{K}%
\TextAt{3}{332.763}{330.249}{\Character{3}}%
\TextAt{8}{342.412}{325.911}{contains}%
\TextAt{8}{389.35}{325.911}{a}%
\TextAt{8}{400.124}{325.911}{group}%
\TextAt{8}{434.318}{325.911}{of}%
\TextAt{8}{448.667}{325.911}{the}%
\TextAt{8}{469.842}{325.911}{form}%
\TextAt{9}{498.501}{325.911}{\Character{40}}%
\TextAt{13}{503.053}{325.911}{\Character{25}}%
\TextAt{21}{510.123}{330.249}{f}%
\TextAt{9}{515.568}{325.911}{\Character{41}}%
\TextAt{16}{523.471}{325.911}{\Character{2}}%
\TextAt{13}{79.226}{308.865}{U}%
\TextAt{17}{88.425}{314.975}{\Character{40}}%
\TextAt{21}{91.718}{314.975}{n}%
\TextAt{17}{96.857}{314.975}{\Character{41}}%
\TextAt{21}{87.176}{305.624}{K}%
\TextAt{8}{100.648}{308.865}{.}%
\TextAt{8}{109.368}{308.865}{This}%
\TextAt{8}{136.179}{308.865}{is}%
\TextAt{8}{148.033}{308.865}{again}%
\TextAt{8}{179.351}{308.865}{open}%
\TextAt{8}{207.727}{308.865}{and}%
\TextAt{8}{230.58}{308.865}{of}%
\TextAt{8}{244.001}{308.865}{finite}%
\TextAt{8}{274.003}{308.865}{index.}%
\TextAt{8}{310.364}{308.865}{Hence}%
\TextAt{8}{345.232}{308.865}{every}%
\TextAt{8}{376.215}{308.865}{finite}%
\TextAt{8}{406.217}{308.865}{abelian}%
\TextAt{8}{446.955}{308.865}{extension}%
\TextAt{13}{498.812}{308.865}{L}%
\TextAt{16}{506.776}{308.865}{\Character{106}}%
\TextAt{13}{510.097}{308.865}{K}%
\TextAt{8}{524.903}{308.865}{is}%
\TextAt{8}{79.226}{291.819}{contained}%
\TextAt{8}{132.261}{291.819}{in}%
\TextAt{8}{145.939}{291.819}{the}%
\TextAt{8}{166.121}{291.819}{class}%
\TextAt{8}{193.583}{291.819}{field}%
\TextAt{8}{218.965}{291.819}{of}%
\TextAt{8}{232.321}{291.819}{such}%
\TextAt{8}{258.743}{291.819}{a}%
\TextAt{8}{268.524}{291.819}{group}%
\TextAt{9}{301.725}{291.819}{\Character{40}}%
\TextAt{13}{306.277}{291.819}{\Character{25}}%
\TextAt{21}{313.346}{296.157}{f}%
\TextAt{9}{318.791}{291.819}{\Character{41}}%
\TextAt{16}{326.018}{291.819}{\Character{2}}%
\TextAt{13}{337.99}{291.819}{U}%
\TextAt{17}{347.189}{297.929}{\Character{40}}%
\TextAt{21}{350.482}{297.929}{n}%
\TextAt{17}{355.621}{297.929}{\Character{41}}%
\TextAt{21}{345.94}{288.578}{K}%
\TextAt{8}{359.412}{291.819}{.}%
\TextAt{8}{367.936}{291.819}{Therefore}%
\TextAt{8}{420.959}{291.819}{the}%
\TextAt{8}{441.141}{291.819}{class}%
\TextAt{8}{468.603}{291.819}{fields}%
\TextAt{8}{498.6}{291.819}{for}%
\TextAt{8}{516.51}{291.819}{the}%
\TextAt{8}{79.226}{274.773}{groups}%
\TextAt{9}{116.696}{274.773}{\Character{40}}%
\TextAt{13}{121.249}{274.773}{\Character{25}}%
\TextAt{21}{128.318}{279.111}{f}%
\TextAt{9}{133.763}{274.773}{\Character{41}}%
\TextAt{16}{140.318}{274.773}{\Character{2}}%
\TextAt{13}{151.618}{274.773}{U}%
\TextAt{17}{160.817}{280.883}{\Character{40}}%
\TextAt{21}{164.111}{280.883}{n}%
\TextAt{17}{169.249}{280.883}{\Character{41}}%
\TextAt{21}{159.569}{271.532}{K}%
\TextAt{8}{176.618}{274.773}{are}%
\TextAt{8}{195.808}{274.773}{particularly}%
\TextAt{8}{258.911}{274.773}{important.}%
\TextAt{8}{318.648}{274.773}{We}%
\TextAt{8}{338.472}{274.773}{will}%
\TextAt{8}{360.257}{274.773}{characterize}%
\TextAt{8}{424.651}{274.773}{them}%
\TextAt{8}{454.242}{274.773}{explicitly}%
\TextAt{8}{504.648}{274.773}{in}%
\TextAt{8}{517.981}{274.773}{§5,}%
\TextAt{8}{79.226}{257.727}{as}%
\TextAt{8}{93.899}{257.727}{immediate}%
\TextAt{8}{151.431}{257.727}{analogues}%
\TextAt{8}{205.137}{257.727}{of}%
\TextAt{8}{218.77}{257.727}{the}%
\TextAt{8}{239.229}{257.727}{cyclotomic}%
\TextAt{8}{297.741}{257.727}{fields}%
\TextAt{8}{328.016}{257.727}{over}%
\TextAt{18}{353.364}{257.727}{Q}%
\TextAt{21}{362.663}{255.934}{p}%
\TextAt{8}{367.424}{257.727}{.}%
\TextAt{8}{376.78}{257.727}{In}%
\TextAt{8}{391.704}{257.727}{the}%
\TextAt{8}{412.163}{257.727}{case}%
\TextAt{8}{437.236}{257.727}{of}%
\TextAt{8}{450.869}{257.727}{the}%
\TextAt{8}{471.328}{257.727}{ground}%
\TextAt{8}{511.31}{257.727}{field}%
\TextAt{13}{79.226}{240.681}{K}%
\TextAt{9}{93.365}{240.681}{\Character{61}}%
\TextAt{18}{105.79}{240.681}{Q}%
\TextAt{21}{115.089}{238.888}{p}%
\TextAt{8}{119.85}{240.681}{,}%
\TextAt{8}{126.689}{240.681}{the}%
\TextAt{8}{146.459}{240.681}{class}%
\TextAt{8}{173.509}{240.681}{field}%
\TextAt{8}{198.479}{240.681}{of}%
\TextAt{8}{211.422}{240.681}{the}%
\TextAt{8}{231.192}{240.681}{group}%
\TextAt{9}{263.98}{240.681}{\Character{40}}%
\TextAt{13}{268.533}{240.681}{p}%
\TextAt{9}{274.408}{240.681}{\Character{41}}%
\TextAt{16}{280.825}{240.681}{\Character{2}}%
\TextAt{13}{291.989}{240.681}{U}%
\TextAt{17}{301.188}{246.791}{\Character{40}}%
\TextAt{21}{304.482}{246.791}{n}%
\TextAt{17}{309.62}{246.791}{\Character{41}}%
\TextAt{21}{299.94}{237.44}{K}%
\TextAt{8}{316.922}{240.681}{is}%
\TextAt{8}{328.299}{240.681}{precisely}%
\TextAt{8}{375.768}{240.681}{the}%
\TextAt{8}{395.538}{240.681}{field}%
\TextAt{18}{420.508}{240.681}{Q}%
\TextAt{21}{429.806}{238.888}{p}%
\TextAt{9}{434.567}{240.681}{\Character{40}}%
\TextAt{13}{439.12}{240.681}{\Character{22}}%
\TextAt{21}{446.163}{238.888}{p}%
\TextAt{22}{450.425}{241.157}{n}%
\TextAt{9}{456.022}{240.681}{\Character{41}}%
\TextAt{8}{464.085}{240.681}{of}%
\TextAt{13}{477.028}{240.681}{p}%
\TextAt{21}{482.903}{245.019}{n}%
\TextAt{8}{488.54}{240.681}{-th}%
\TextAt{8}{507.006}{240.681}{roots}%
\TextAt{8}{79.226}{223.635}{of}%
\TextAt{8}{92.556}{223.635}{unity:}%
\TextAt{12}{79.226}{198.771}{Proposition}%
\TextAt{12}{152.033}{198.771}{(1.8).}%
\TextAt{8}{187.431}{198.771}{The}%
\TextAt{8}{211.568}{198.771}{group}%
\TextAt{8}{244.826}{198.771}{of}%
\TextAt{8}{258.239}{198.771}{norms}%
\TextAt{8}{293.506}{198.771}{of}%
\TextAt{8}{306.919}{198.771}{the}%
\TextAt{8}{327.158}{198.771}{extension}%
\TextAt{18}{379.006}{198.771}{Q}%
\TextAt{21}{388.305}{196.977}{p}%
\TextAt{9}{393.066}{198.771}{\Character{40}}%
\TextAt{13}{397.618}{198.771}{\Character{22}}%
\TextAt{21}{404.661}{196.977}{p}%
\TextAt{22}{408.924}{199.247}{n}%
\TextAt{9}{414.521}{198.771}{\Character{41}}%
\TextAt{16}{419.073}{198.771}{\Character{106}}%
\TextAt{18}{422.394}{198.771}{Q}%
\TextAt{21}{431.693}{196.977}{p}%
\TextAt{8}{440.434}{198.771}{is}%
\TextAt{8}{452.28}{198.771}{the}%
\TextAt{8}{472.52}{198.771}{group}%
\TextAt{9}{505.778}{198.771}{\Character{40}}%
\TextAt{13}{510.33}{198.771}{p}%
\TextAt{9}{516.206}{198.771}{\Character{41}}%
\TextAt{16}{523.471}{198.771}{\Character{2}}%
\TextAt{13}{79.226}{181.725}{U}%
\TextAt{17}{88.425}{187.835}{\Character{40}}%
\TextAt{21}{91.718}{187.835}{n}%
\TextAt{17}{96.857}{187.835}{\Character{41}}%
\TextAt{25}{87.176}{178.439}{Q}%
\TextAt{22}{93.375}{177.443}{p}%
\TextAt{8}{100.648}{181.725}{.}%
\TextAt{19}{79.226}{149.042}{Proof.}%
\TextAt{8}{116}{149.042}{Let}%
\TextAt{13}{137.383}{149.042}{K}%
\TextAt{9}{152.227}{149.042}{\Character{61}}%
\TextAt{18}{165.357}{149.042}{Q}%
\TextAt{21}{174.656}{147.249}{p}%
\TextAt{8}{183.728}{149.042}{and}%
\TextAt{13}{206.904}{149.042}{L}%
\TextAt{9}{218.894}{149.042}{\Character{61}}%
\TextAt{18}{232.025}{149.042}{Q}%
\TextAt{21}{241.323}{147.249}{p}%
\TextAt{9}{246.084}{149.042}{\Character{40}}%
\TextAt{13}{250.636}{149.042}{\Character{22}}%
\TextAt{21}{257.679}{147.249}{p}%
\TextAt{22}{261.942}{149.518}{n}%
\TextAt{9}{267.539}{149.042}{\Character{41}}%
\TextAt{8}{272.091}{149.042}{.}%
\TextAt{8}{281.781}{149.042}{By}%
\TextAt{8}{300.558}{149.042}{chap.}%
\TextAt{8}{333.99}{149.042}{II,}%
\TextAt{8}{350.316}{149.042}{(7.13),}%
\TextAt{8}{387.918}{149.042}{the}%
\TextAt{8}{408.488}{149.042}{extension}%
\TextAt{13}{460.667}{149.042}{L}%
\TextAt{16}{468.632}{149.042}{\Character{106}}%
\TextAt{13}{471.953}{149.042}{K}%
\TextAt{8}{487.082}{149.042}{is}%
\TextAt{8}{499.259}{149.042}{totally}%
\TextAt{8}{79.226}{131.996}{ramified}%
\TextAt{8}{124.856}{131.996}{of}%
\TextAt{8}{138.29}{131.996}{degree}%
\TextAt{13}{174.809}{131.996}{p}%
\TextAt{21}{180.684}{136.335}{n}%
\TextAt{3}{185.823}{136.335}{\Character{0}}%
\TextAt{17}{192.409}{136.335}{1}%
\TextAt{9}{197.141}{131.996}{\Character{40}}%
\TextAt{13}{201.694}{131.996}{p}%
\TextAt{16}{210.297}{131.996}{\Character{0}}%
\TextAt{9}{222.323}{131.996}{1\Character{41}}%
\TextAt{8}{232.728}{131.996}{,}%
\TextAt{8}{240.007}{131.996}{and}%
\TextAt{8}{262.874}{131.996}{if}%
\TextAt{13}{273.702}{131.996}{\Character{16}}%
\TextAt{8}{283.719}{131.996}{is}%
\TextAt{8}{295.587}{131.996}{a}%
\TextAt{8}{305.447}{131.996}{primitive}%
\TextAt{13}{355.631}{131.996}{p}%
\TextAt{21}{361.506}{136.335}{n}%
\TextAt{8}{367.143}{131.996}{-th}%
\TextAt{8}{386.1}{131.996}{root}%
\TextAt{8}{411.25}{131.996}{of}%
\TextAt{8}{424.684}{131.996}{unity,}%
\TextAt{8}{457.655}{131.996}{then}%
\TextAt{9}{484.42}{131.996}{1}%
\TextAt{16}{493}{131.996}{\Character{0}}%
\TextAt{13}{505.026}{131.996}{\Character{16}}%
\TextAt{8}{515.043}{131.996}{is}%
\TextAt{8}{526.911}{131.996}{a}%
\TextAt{8}{79.226}{114.95}{prime}%
\TextAt{8}{112.805}{114.95}{element}%
\TextAt{8}{156.463}{114.95}{of}%
\TextAt{13}{170.209}{114.95}{L}%
\TextAt{8}{182.486}{114.95}{of}%
\TextAt{8}{196.232}{114.95}{norm}%
\TextAt{13}{227.217}{114.95}{N}%
\TextAt{21}{236.591}{113.09}{L}%
\TextAt{3}{242.35}{113.09}{\Character{106}}%
\TextAt{21}{244.702}{113.09}{K}%
\TextAt{9}{252.905}{114.95}{\Character{40}1}%
\TextAt{16}{266.25}{114.95}{\Character{0}}%
\TextAt{13}{278.488}{114.95}{\Character{16}}%
\TextAt{9}{284.504}{114.95}{\Character{41}}%
\TextAt{9}{293.076}{114.95}{\Character{61}}%
\TextAt{13}{306.219}{114.95}{p}%
\TextAt{8}{312.094}{114.95}{.}%
\TextAt{8}{321.789}{114.95}{We}%
\TextAt{8}{342.35}{114.95}{now}%
\TextAt{8}{367.154}{114.95}{consider}%
\TextAt{8}{413.154}{114.95}{the}%
\TextAt{8}{433.727}{114.95}{exponential}%
\TextAt{8}{496.907}{114.95}{map}%
\TextAt{8}{523.337}{114.95}{of}%
\TextAt{18}{79.226}{97.904}{Q}%
\TextAt{21}{88.525}{96.111}{p}%
\TextAt{8}{93.286}{97.904}{.}%
\TextAt{8}{101.734}{97.904}{By}%
\TextAt{8}{120.097}{97.904}{chap.}%
\TextAt{8}{152.288}{97.904}{II,}%
\TextAt{8}{168.201}{97.904}{(5.5),}%
\TextAt{8}{199.428}{97.904}{it}%
\TextAt{8}{211.132}{97.904}{gives}%
\TextAt{8}{239.812}{97.904}{an}%
\TextAt{8}{256.071}{97.904}{isomorphism}%
\TextAt{8}{261.581}{72.921}{exp}%
\TextAt{9}{282.786}{72.921}{\Character{58}}%
\TextAt{14}{289.359}{72.921}{p}%
\TextAt{21}{295.342}{77.858}{\Character{23}}%
\TextAt{21}{295.342}{69.966}{K}%
\NativeArrowAt{306.865}{72.921}{11.955}{19.26}{\longrightarrow}%
\TextAt{13}{329.447}{72.921}{U}%
\TextAt{17}{338.647}{79.032}{\Character{40}}%
\TextAt{21}{341.94}{79.032}{\Character{23}}%
\TextAt{17}{346.624}{79.032}{\Character{41}}%
\TextAt{21}{337.398}{69.68}{K}%
\TextAt{8}{297.211}{43.033}{296}%
\Destination{78.23}{809.815}{nav.page.296}%
\Destination{79.36}{659.355}{nav.print.323}%
\Destination{79.36}{750.255}{nav.number.5.1.5}%
\Destination{79.36}{565.355}{nav.number.5.1.6}%
\Destination{79.36}{468.655}{nav.number.5.1.7}%
\Destination{79.36}{212.755}{nav.number.5.1.8}%
\Hotlink{442.498}{608.055}{9.751}{11.047}{nav.chapter.2}%
\Hotlink{523.16}{272.555}{6.858}{11.047}{nav.section.5.5}%
\Hotlink{333.49}{146.855}{9.751}{11.047}{nav.chapter.2}%
\Hotlink{151.788}{95.655}{9.763}{11.047}{nav.chapter.2}%
\Hotlink{147.078}{659.255}{25.078}{11.047}{nav.number.5.1.4}%
\Hotlink{460.166}{608.055}{25.078}{11.047}{nav.number.2.5.7}%
\Hotlink{354.457}{402.655}{25.078}{11.047}{nav.number.5.1.2}%
\Hotlink{349.816}{146.855}{30.936}{11.047}{nav.number.2.7.13}%
\Hotlink{167.701}{95.655}{25.078}{11.047}{nav.number.2.5.5}%
}

% Page 309 (body page 297)
\Page{%
\TextAt{2}{79.226}{796.943}{V.}%
\TextAt{2}{94.073}{796.943}{Local}%
\TextAt{2}{123.615}{796.943}{Class}%
\TextAt{2}{152.208}{796.943}{Field}%
\TextAt{2}{179.939}{796.943}{Theory}%
\Vector{q 0 G 0 g 0.299 w 7.226 18.729 m 460.769 18.729 l S Q}%
\TextAt{8}{79.226}{761.124}{for}%
\TextAt{13}{97.53}{761.124}{\Character{23}}%
\NativeMathAt{108.101}{761.124}{11.955}{\geq}%
\TextAt{9}{121.434}{761.124}{1}%
\TextAt{8}{127.287}{761.124}{,}%
\TextAt{8}{134.959}{761.124}{provided}%
\TextAt{13}{183.509}{761.124}{p}%
\NativeMathAt{193.418}{761.124}{11.955}{\neq}%
\TextAt{9}{206.545}{761.124}{2}%
\TextAt{8}{212.41}{761.124}{,}%
\TextAt{8}{220.082}{761.124}{and}%
\TextAt{8}{243.263}{761.124}{for}%
\TextAt{13}{261.567}{761.124}{\Character{23}}%
\NativeMathAt{272.139}{761.124}{11.955}{\geq}%
\TextAt{9}{285.471}{761.124}{2}%
\TextAt{8}{291.324}{761.124}{,}%
\TextAt{8}{298.997}{761.124}{even}%
\TextAt{8}{326.075}{761.124}{if}%
\TextAt{13}{337.217}{761.124}{p}%
\TextAt{9}{347.126}{761.124}{\Character{61}}%
\TextAt{9}{360.253}{761.124}{2}%
\TextAt{8}{366.118}{761.124}{.}%
\TextAt{8}{375.822}{761.124}{It}%
\TextAt{8}{388.913}{761.124}{transforms}%
\TextAt{8}{447.673}{761.124}{the}%
\TextAt{8}{468.247}{761.124}{isomorphism}%
\TextAt{14}{79.226}{744.078}{p}%
\TextAt{21}{85.209}{748.416}{\Character{23}}%
\TextAt{21}{85.209}{741.057}{K}%
\NativeMathAt{96.893}{744.078}{11.955}{\rightarrow}%
\TextAt{14}{112.328}{744.078}{p}%
\TextAt{21}{118.311}{749.026}{\Character{23}}%
\TextAt{17}{122.995}{749.026}{\Character{43}}%
\TextAt{21}{129.582}{749.026}{s}%
\TextAt{3}{133.498}{749.026}{\Character{0}}%
\TextAt{17}{140.084}{749.026}{1}%
\TextAt{21}{118.311}{740.837}{K}%
\TextAt{8}{148.808}{744.078}{given}%
\TextAt{8}{179.471}{744.078}{by}%
\TextAt{13}{195.824}{744.078}{a}%
\NativeArrowAt{205.449}{744.078}{11.955}{11.955}{\mapsto}%
\TextAt{13}{220.885}{744.078}{p}%
\TextAt{21}{226.76}{748.416}{s}%
\TextAt{3}{230.676}{748.416}{\Character{0}}%
\TextAt{17}{237.263}{748.416}{1}%
\TextAt{9}{241.995}{744.078}{\Character{40}}%
\TextAt{13}{246.547}{744.078}{p}%
\TextAt{16}{255.143}{744.078}{\Character{0}}%
\TextAt{9}{267.162}{744.078}{1\Character{41}}%
\TextAt{13}{277.567}{744.078}{a}%
\TextAt{8}{283.712}{744.078}{,}%
\TextAt{8}{290.979}{744.078}{into}%
\TextAt{8}{314.815}{744.078}{the}%
\TextAt{8}{335.066}{744.078}{isomorphism}%
\TextAt{13}{403.579}{744.078}{U}%
\TextAt{17}{412.778}{750.188}{\Character{40}}%
\TextAt{21}{416.071}{750.188}{\Character{23}}%
\TextAt{17}{420.755}{750.188}{\Character{41}}%
\TextAt{21}{411.529}{740.837}{K}%
\NativeMathAt{428.027}{744.078}{11.955}{\rightarrow}%
\TextAt{13}{443.463}{744.078}{U}%
\TextAt{17}{452.662}{750.188}{\Character{40}}%
\TextAt{21}{455.955}{750.188}{\Character{23}}%
\TextAt{17}{460.639}{750.188}{\Character{43}}%
\TextAt{21}{467.226}{750.188}{s}%
\TextAt{3}{471.142}{750.188}{\Character{0}}%
\TextAt{17}{477.728}{750.188}{1\Character{41}}%
\TextAt{21}{451.413}{740.837}{K}%
\TextAt{8}{489.745}{744.078}{given}%
\TextAt{8}{520.408}{744.078}{by}%
\TextAt{13}{79.226}{727.032}{x}%
\NativeArrowAt{90.728}{727.032}{11.955}{11.955}{\mapsto}%
\TextAt{13}{107.532}{727.032}{x}%
\TextAt{21}{114.184}{731.37}{p}%
\TextAt{22}{118.447}{734.183}{s}%
\TextAt{23}{121.941}{734.183}{\Character{0}}%
\TextAt{24}{127.697}{734.183}{1}%
\TextAt{17}{131.848}{731.37}{\Character{40}}%
\TextAt{21}{135.141}{731.37}{p}%
\TextAt{3}{139.404}{731.37}{\Character{0}}%
\TextAt{17}{145.991}{731.37}{1\Character{41}}%
\TextAt{8}{154.016}{727.032}{,}%
\TextAt{8}{162.287}{727.032}{so}%
\TextAt{8}{177.554}{727.032}{that}%
\TextAt{9}{203.82}{727.032}{\Character{40}}%
\TextAt{13}{208.372}{727.032}{U}%
\TextAt{17}{217.571}{733.142}{\Character{40}1\Character{41}}%
\TextAt{21}{216.322}{723.791}{K}%
\TextAt{9}{228.89}{727.032}{\Character{41}}%
\TextAt{21}{233.442}{731.37}{p}%
\TextAt{22}{237.705}{734.183}{n}%
\TextAt{23}{242.306}{734.183}{\Character{0}}%
\TextAt{24}{248.062}{734.183}{1}%
\TextAt{17}{252.213}{731.37}{\Character{40}}%
\TextAt{21}{255.506}{731.37}{p}%
\TextAt{3}{259.769}{731.37}{\Character{0}}%
\TextAt{17}{266.355}{731.37}{1\Character{41}}%
\TextAt{9}{279.23}{727.032}{\Character{61}}%
\TextAt{13}{293.184}{727.032}{U}%
\TextAt{17}{302.383}{733.142}{\Character{40}}%
\TextAt{21}{305.677}{733.142}{n}%
\TextAt{17}{310.815}{733.142}{\Character{41}}%
\TextAt{21}{301.135}{723.791}{K}%
\TextAt{8}{319.401}{727.032}{if}%
\TextAt{13}{331.021}{727.032}{p}%
\NativeMathAt{341.746}{727.032}{11.955}{\neq}%
\TextAt{9}{355.686}{727.032}{2}%
\TextAt{8}{361.553}{727.032}{,}%
\TextAt{8}{369.823}{727.032}{and}%
\TextAt{9}{393.483}{727.032}{\Character{40}}%
\TextAt{13}{398.035}{727.032}{U}%
\TextAt{17}{407.234}{733.142}{\Character{40}2\Character{41}}%
\TextAt{21}{405.986}{723.791}{K}%
\TextAt{9}{418.553}{727.032}{\Character{41}}%
\TextAt{17}{423.105}{731.37}{2}%
\TextAt{22}{427.34}{734.183}{n}%
\TextAt{23}{431.94}{734.183}{\Character{0}}%
\TextAt{24}{437.697}{734.183}{2}%
\TextAt{9}{447.195}{727.032}{\Character{61}}%
\TextAt{13}{461.149}{727.032}{U}%
\TextAt{17}{470.348}{733.142}{\Character{40}}%
\TextAt{21}{473.641}{733.142}{n}%
\TextAt{17}{478.78}{733.142}{\Character{41}}%
\TextAt{21}{469.099}{723.791}{K}%
\TextAt{8}{487.365}{727.032}{if}%
\TextAt{13}{498.986}{727.032}{p}%
\TextAt{9}{509.711}{727.032}{\Character{61}}%
\TextAt{9}{523.651}{727.032}{2}%
\TextAt{8}{529.518}{727.032}{,}%
\TextAt{13}{79.226}{709.986}{n}%
\TextAt{13}{89.986}{709.986}{\Character{62}}%
\TextAt{9}{102.893}{709.986}{1}%
\TextAt{8}{112.917}{709.986}{(the}%
\TextAt{8}{137.901}{709.986}{case}%
\TextAt{13}{162.946}{709.986}{p}%
\TextAt{9}{172.608}{709.986}{\Character{61}}%
\TextAt{9}{185.484}{709.986}{2}%
\TextAt{8}{191.353}{709.986}{,}%
\TextAt{13}{198.844}{709.986}{n}%
\TextAt{9}{209.619}{709.986}{\Character{61}}%
\TextAt{9}{222.495}{709.986}{1}%
\TextAt{8}{232.535}{709.986}{is}%
\TextAt{8}{244.572}{709.986}{trivial).}%
\TextAt{8}{289.3}{709.986}{Consequently,}%
\TextAt{8}{364.494}{709.986}{we}%
\TextAt{8}{381.995}{709.986}{have}%
\TextAt{13}{409.263}{709.986}{U}%
\TextAt{17}{418.462}{716.096}{\Character{40}}%
\TextAt{21}{421.755}{716.096}{n}%
\TextAt{17}{426.894}{716.096}{\Character{41}}%
\TextAt{21}{417.213}{706.745}{K}%
\NativeMathAt{434.472}{709.986}{11.955}{\subseteq}%
\TextAt{13}{447.558}{709.986}{N}%
\TextAt{21}{456.932}{708.126}{L}%
\TextAt{3}{462.691}{708.126}{\Character{106}}%
\TextAt{21}{465.043}{708.126}{K}%
\TextAt{13}{473.246}{709.986}{L}%
\TextAt{3}{481.21}{714.324}{\Character{3}}%
\TextAt{8}{490.113}{709.986}{if}%
\TextAt{13}{501.111}{709.986}{p}%
\NativeMathAt{510.773}{709.986}{11.955}{\neq}%
\TextAt{9}{523.649}{709.986}{2}%
\TextAt{8}{529.518}{709.986}{.}%
\TextAt{8}{79.226}{692.94}{For}%
\TextAt{13}{100.196}{692.94}{p}%
\TextAt{9}{109.391}{692.94}{\Character{61}}%
\TextAt{9}{121.801}{692.94}{2}%
\TextAt{8}{131.567}{692.94}{we}%
\TextAt{8}{148.795}{692.94}{note}%
\TextAt{8}{174.809}{692.94}{that}%
\TextAt{13}{204.358}{664.353}{U}%
\TextAt{17}{213.557}{670.463}{\Character{40}2\Character{41}}%
\TextAt{21}{212.308}{661.112}{K}%
\TextAt{9}{228.196}{664.353}{\Character{61}}%
\TextAt{13}{240.622}{664.353}{U}%
\TextAt{17}{249.821}{670.463}{\Character{40}3\Character{41}}%
\TextAt{21}{248.572}{661.112}{K}%
\TextAt{16}{263.796}{664.353}{\Character{91}}%
\TextAt{9}{274.423}{664.353}{5}%
\TextAt{13}{280.276}{664.353}{U}%
\TextAt{17}{289.475}{670.463}{\Character{40}3\Character{41}}%
\TextAt{21}{288.226}{661.112}{K}%
\TextAt{9}{304.115}{664.353}{\Character{61}}%
\TextAt{9}{316.524}{664.353}{\Character{40}}%
\TextAt{13}{321.093}{664.353}{U}%
\TextAt{17}{330.292}{670.463}{\Character{40}2\Character{41}}%
\TextAt{21}{329.043}{661.112}{K}%
\TextAt{9}{341.61}{664.353}{\Character{41}}%
\TextAt{17}{346.163}{669.289}{2}%
\TextAt{16}{353.552}{664.353}{\Character{91}}%
\TextAt{9}{364.179}{664.353}{5\Character{40}}%
\TextAt{13}{374.584}{664.353}{U}%
\TextAt{17}{383.783}{670.463}{\Character{40}2\Character{41}}%
\TextAt{21}{382.534}{661.112}{K}%
\TextAt{9}{395.102}{664.353}{\Character{41}}%
\TextAt{17}{399.654}{669.289}{2}%
\TextAt{13}{404.386}{664.353}{\Character{59}}%
\TextAt{8}{79.226}{635.766}{because}%
\TextAt{8}{122.528}{635.766}{a}%
\TextAt{8}{132.283}{635.766}{number}%
\TextAt{8}{174.879}{635.766}{that}%
\TextAt{8}{200.248}{635.766}{is}%
\TextAt{8}{212.012}{635.766}{congruent}%
\TextAt{8}{266.325}{635.766}{to}%
\TextAt{9}{280.635}{635.766}{1}%
\TextAt{8}{289.809}{635.766}{mod}%
\TextAt{9}{315.247}{635.766}{4}%
\TextAt{8}{324.997}{635.766}{is}%
\TextAt{8}{336.761}{635.766}{congruent}%
\TextAt{8}{391.073}{635.766}{to}%
\TextAt{9}{405.384}{635.766}{1}%
\TextAt{8}{415.134}{635.766}{or}%
\TextAt{9}{429.444}{635.766}{5}%
\TextAt{8}{438.618}{635.766}{mod}%
\TextAt{9}{464.056}{635.766}{8}%
\TextAt{8}{469.909}{635.766}{.}%
\TextAt{8}{478.357}{635.766}{Hence}%
\TextAt{13}{217.737}{607.179}{U}%
\TextAt{17}{226.936}{613.29}{\Character{40}}%
\TextAt{21}{230.229}{613.29}{n}%
\TextAt{17}{235.367}{613.29}{\Character{41}}%
\TextAt{21}{225.687}{603.938}{K}%
\TextAt{9}{242.479}{607.179}{\Character{61}}%
\TextAt{9}{254.889}{607.179}{\Character{40}}%
\TextAt{13}{259.457}{607.179}{U}%
\TextAt{17}{268.656}{613.29}{\Character{40}2\Character{41}}%
\TextAt{21}{267.408}{603.938}{K}%
\TextAt{9}{279.975}{607.179}{\Character{41}}%
\TextAt{17}{284.527}{612.116}{2}%
\TextAt{22}{288.762}{614.928}{n}%
\TextAt{23}{293.362}{614.928}{\Character{0}}%
\TextAt{24}{299.119}{614.928}{1}%
\TextAt{16}{306.424}{607.179}{\Character{91}}%
\TextAt{9}{317.051}{607.179}{5}%
\TextAt{17}{322.904}{612.116}{2}%
\TextAt{22}{327.138}{614.928}{n}%
\TextAt{23}{331.739}{614.928}{\Character{0}}%
\TextAt{24}{337.495}{614.928}{2}%
\TextAt{9}{342.144}{607.179}{\Character{40}}%
\TextAt{13}{346.697}{607.179}{U}%
\TextAt{17}{355.896}{613.29}{\Character{40}2\Character{41}}%
\TextAt{21}{354.647}{603.938}{K}%
\TextAt{9}{367.215}{607.179}{\Character{41}}%
\TextAt{17}{371.767}{612.116}{2}%
\TextAt{22}{376.001}{614.928}{n}%
\TextAt{23}{380.602}{614.928}{\Character{0}}%
\TextAt{24}{386.358}{614.928}{1}%
\TextAt{13}{391.007}{607.179}{\Character{58}}%
\TextAt{8}{79.226}{578.593}{It}%
\TextAt{8}{92.315}{578.593}{is}%
\TextAt{8}{104.496}{578.593}{easy}%
\TextAt{8}{130.665}{578.593}{to}%
\TextAt{8}{145.392}{578.593}{show}%
\TextAt{8}{174.812}{578.593}{that}%
\TextAt{9}{200.598}{578.593}{5}%
\TextAt{17}{206.451}{582.931}{2}%
\TextAt{22}{210.685}{585.744}{n}%
\TextAt{23}{215.285}{585.744}{\Character{0}}%
\TextAt{24}{221.042}{585.744}{2}%
\TextAt{9}{229.722}{578.593}{\Character{61}}%
\TextAt{13}{242.858}{578.593}{N}%
\TextAt{21}{252.232}{576.733}{L}%
\TextAt{3}{257.991}{576.733}{\Character{106}}%
\TextAt{21}{260.343}{576.733}{K}%
\TextAt{9}{268.546}{578.593}{\Character{40}2}%
\TextAt{9}{281.864}{578.593}{\Character{43}}%
\TextAt{13}{293.937}{578.593}{i}%
\TextAt{9}{297.931}{578.593}{\Character{41}}%
\TextAt{8}{302.483}{578.593}{,}%
\TextAt{8}{310.153}{578.593}{so}%
\TextAt{13}{324.94}{578.593}{U}%
\TextAt{17}{334.139}{584.703}{\Character{40}}%
\TextAt{21}{337.433}{584.703}{n}%
\TextAt{17}{342.571}{584.703}{\Character{41}}%
\TextAt{21}{332.891}{575.351}{K}%
\NativeMathAt{350.393}{578.593}{11.955}{\subseteq}%
\TextAt{13}{363.723}{578.593}{N}%
\TextAt{21}{373.097}{576.733}{L}%
\TextAt{3}{378.856}{576.733}{\Character{106}}%
\TextAt{21}{381.208}{576.733}{K}%
\TextAt{13}{389.411}{578.593}{L}%
\TextAt{3}{397.375}{582.931}{\Character{3}}%
\TextAt{8}{406.422}{578.593}{holds}%
\TextAt{8}{437.468}{578.593}{also}%
\TextAt{8}{461.365}{578.593}{in}%
\TextAt{8}{475.435}{578.593}{case}%
\TextAt{13}{500.622}{578.593}{p}%
\TextAt{9}{510.529}{578.593}{\Character{61}}%
\TextAt{9}{523.656}{578.593}{2}%
\TextAt{8}{529.518}{578.593}{.}%
\TextAt{8}{79.226}{561.547}{Since}%
\TextAt{13}{110.02}{561.547}{p}%
\TextAt{9}{119.618}{561.547}{\Character{61}}%
\TextAt{13}{132.447}{561.547}{N}%
\TextAt{21}{141.82}{559.687}{L}%
\TextAt{3}{147.579}{559.687}{\Character{106}}%
\TextAt{21}{149.932}{559.687}{K}%
\TextAt{9}{158.134}{561.547}{\Character{40}1}%
\TextAt{16}{171.357}{561.547}{\Character{0}}%
\TextAt{13}{183.474}{561.547}{\Character{16}}%
\TextAt{9}{189.489}{561.547}{\Character{41}}%
\TextAt{8}{194.042}{561.547}{,}%
\TextAt{8}{201.486}{561.547}{we}%
\TextAt{8}{218.95}{561.547}{have}%
\TextAt{9}{246.181}{561.547}{\Character{40}}%
\TextAt{13}{250.733}{561.547}{p}%
\TextAt{9}{256.608}{561.547}{\Character{41}}%
\TextAt{16}{263.978}{561.547}{\Character{2}}%
\TextAt{13}{276.095}{561.547}{U}%
\TextAt{17}{285.294}{567.657}{\Character{40}}%
\TextAt{21}{288.587}{567.657}{n}%
\TextAt{17}{293.725}{567.657}{\Character{41}}%
\TextAt{21}{284.045}{558.305}{K}%
\NativeMathAt{301.24}{561.547}{11.955}{\subseteq}%
\TextAt{13}{314.262}{561.547}{N}%
\TextAt{21}{323.636}{559.687}{L}%
\TextAt{3}{329.395}{559.687}{\Character{106}}%
\TextAt{21}{331.747}{559.687}{K}%
\TextAt{13}{339.95}{561.547}{L}%
\TextAt{3}{347.914}{565.885}{\Character{3}}%
\TextAt{8}{352.647}{561.547}{,}%
\TextAt{8}{360.091}{561.547}{and}%
\TextAt{8}{383.09}{561.547}{since}%
\TextAt{8}{411.995}{561.547}{both}%
\TextAt{8}{439.872}{561.547}{groups}%
\TextAt{8}{477.898}{561.547}{have}%
\TextAt{8}{505.129}{561.547}{index}%
\TextAt{13}{79.226}{544.501}{p}%
\TextAt{21}{85.101}{548.839}{n}%
\TextAt{3}{90.239}{548.839}{\Character{0}}%
\TextAt{17}{96.826}{548.839}{1}%
\TextAt{9}{101.558}{544.501}{\Character{40}}%
\TextAt{13}{106.111}{544.501}{p}%
\TextAt{16}{114.642}{544.501}{\Character{0}}%
\TextAt{9}{126.598}{544.501}{1\Character{41}}%
\TextAt{8}{140.9}{544.501}{in}%
\TextAt{13}{154.553}{544.501}{K}%
\TextAt{3}{165.371}{548.839}{\Character{3}}%
\TextAt{8}{170.103}{544.501}{,}%
\TextAt{8}{177.252}{544.501}{we}%
\TextAt{8}{194.48}{544.501}{do}%
\TextAt{8}{210.739}{544.501}{find}%
\TextAt{8}{234.147}{544.501}{that}%
\TextAt{13}{259.516}{544.501}{N}%
\TextAt{21}{268.89}{542.641}{L}%
\TextAt{3}{274.649}{542.641}{\Character{106}}%
\TextAt{21}{277.001}{542.641}{K}%
\TextAt{13}{285.204}{544.501}{L}%
\TextAt{3}{293.168}{548.839}{\Character{3}}%
\TextAt{9}{301.221}{544.501}{\Character{61}}%
\TextAt{9}{313.631}{544.501}{\Character{40}}%
\TextAt{13}{318.199}{544.501}{p}%
\TextAt{9}{324.074}{544.501}{\Character{41}}%
\TextAt{16}{331.283}{544.501}{\Character{2}}%
\TextAt{13}{343.238}{544.501}{U}%
\TextAt{17}{352.437}{550.611}{\Character{40}}%
\TextAt{21}{355.731}{550.611}{n}%
\TextAt{17}{360.869}{550.611}{\Character{41}}%
\TextAt{21}{351.189}{541.26}{K}%
\TextAt{8}{368.558}{544.501}{as}%
\TextAt{8}{382.928}{544.501}{claimed.}%
\Vector{q 0 G 0 g 0.399 w 452.334 -229.499 m 452.334 -221.43 l S Q}%
\Vector{q 0 G 0 g 0.399 w 452.534 -221.629 m 459.707 -221.629 l S Q}%
\Vector{q 0 G 0 g 0.399 w 452.534 -229.3 m 459.707 -229.3 l S Q}%
\Vector{q 0 G 0 g 0.399 w 459.906 -229.499 m 459.906 -221.43 l S Q}%
\TextAt{8}{97.159}{515.618}{As}%
\TextAt{8}{114.104}{515.618}{an}%
\TextAt{8}{130.021}{515.618}{immediate}%
\TextAt{8}{186.908}{515.618}{consequence}%
\TextAt{8}{252.63}{515.618}{of}%
\TextAt{8}{265.618}{515.618}{this}%
\TextAt{8}{288.098}{515.618}{last}%
\TextAt{8}{309.933}{515.618}{proposition,}%
\TextAt{8}{374.444}{515.618}{we}%
\TextAt{8}{391.329}{515.618}{obtain}%
\TextAt{8}{427.414}{515.618}{a}%
\TextAt{8}{436.828}{515.618}{local}%
\TextAt{8}{464.126}{515.618}{version}%
\TextAt{8}{503.523}{515.618}{of}%
\TextAt{8}{516.51}{515.618}{the}%
\TextAt{8}{79.226}{498.572}{famous}%
\TextAt{8}{119.591}{498.572}{theorem}%
\TextAt{8}{165.42}{498.572}{of}%
\TextAt{8}{179.054}{498.572}{Kronecker–Weber,}%
\TextAt{8}{276.975}{498.572}{to}%
\TextAt{8}{291.589}{498.572}{the}%
\TextAt{8}{312.049}{498.572}{effect}%
\TextAt{8}{343.233}{498.572}{that}%
\TextAt{8}{368.905}{498.572}{every}%
\TextAt{8}{400.101}{498.572}{finite}%
\TextAt{8}{430.316}{498.572}{abelian}%
\TextAt{8}{471.267}{498.572}{extension}%
\TextAt{8}{523.337}{498.572}{of}%
\TextAt{18}{79.226}{481.526}{Q}%
\TextAt{8}{92.422}{481.526}{is}%
\TextAt{8}{104.186}{481.526}{contained}%
\TextAt{8}{157.195}{481.526}{in}%
\TextAt{8}{170.848}{481.526}{a}%
\TextAt{8}{180.603}{481.526}{cyclotomic}%
\TextAt{8}{238.813}{481.526}{field.}%
\TextAt{12}{79.226}{455.632}{Corollary}%
\TextAt{12}{140.231}{455.632}{(1.9).}%
\TextAt{8}{176.606}{455.632}{Every}%
\TextAt{8}{211.318}{455.632}{finite}%
\TextAt{8}{242.288}{455.632}{abelian}%
\TextAt{8}{283.994}{455.632}{extension}%
\TextAt{8}{336.818}{455.632}{of}%
\TextAt{13}{351.207}{455.632}{L}%
\TextAt{16}{359.171}{455.632}{\Character{106}}%
\TextAt{18}{362.492}{455.632}{Q}%
\TextAt{21}{371.791}{453.839}{p}%
\TextAt{8}{381.508}{455.632}{is}%
\TextAt{8}{394.33}{455.632}{contained}%
\TextAt{8}{448.398}{455.632}{in}%
\TextAt{8}{463.109}{455.632}{a}%
\TextAt{8}{473.923}{455.632}{field}%
\TextAt{18}{500.338}{455.632}{Q}%
\TextAt{21}{509.637}{453.839}{p}%
\TextAt{9}{514.397}{455.632}{\Character{40}}%
\TextAt{13}{518.95}{455.632}{\Character{16}}%
\TextAt{9}{524.965}{455.632}{\Character{41}}%
\TextAt{8}{529.518}{455.632}{,}%
\TextAt{8}{79.226}{438.586}{where}%
\TextAt{13}{113.035}{438.586}{\Character{16}}%
\TextAt{8}{122.948}{438.586}{is}%
\TextAt{8}{134.712}{438.586}{a}%
\TextAt{8}{144.468}{438.586}{root}%
\TextAt{8}{169.514}{438.586}{of}%
\TextAt{8}{182.844}{438.586}{unity.}%
\TextAt{8}{216.984}{438.586}{In}%
\TextAt{8}{231.605}{438.586}{other}%
\TextAt{8}{262.174}{438.586}{words:}%
\TextAt{8}{97.159}{418.551}{The}%
\TextAt{8}{121.213}{418.551}{maximal}%
\TextAt{8}{169.021}{418.551}{abelian}%
\TextAt{8}{209.669}{418.551}{extension}%
\TextAt{18}{261.435}{418.551}{Q}%
\TextAt{7}{270.733}{422.89}{ab}%
\TextAt{21}{270.733}{415.596}{p}%
\TextAt{16}{280.166}{418.551}{\Character{106}}%
\TextAt{18}{283.487}{418.551}{Q}%
\TextAt{21}{292.785}{416.758}{p}%
\TextAt{8}{301.444}{418.551}{is}%
\TextAt{8}{313.208}{418.551}{generated}%
\TextAt{8}{366.54}{418.551}{by}%
\TextAt{8}{382.799}{418.551}{adjoining}%
\TextAt{8}{433.859}{418.551}{all}%
\TextAt{8}{450.118}{418.551}{roots}%
\TextAt{8}{479.779}{418.551}{of}%
\TextAt{8}{493.109}{418.551}{unity.}%
\TextAt{19}{79.226}{383.809}{Proof.}%
\TextAt{8}{116}{383.809}{For}%
\TextAt{8}{137.047}{383.809}{suitable}%
\TextAt{13}{180.761}{383.809}{f}%
\TextAt{8}{191.783}{383.809}{and}%
\TextAt{13}{214.623}{383.809}{n}%
\TextAt{8}{221.611}{383.809}{,}%
\TextAt{8}{228.857}{383.809}{we}%
\TextAt{8}{246.163}{383.809}{have}%
\TextAt{9}{273.235}{383.809}{\Character{40}}%
\TextAt{13}{277.787}{383.809}{p}%
\TextAt{21}{283.663}{388.148}{f}%
\TextAt{9}{289.107}{383.809}{\Character{41}}%
\TextAt{16}{296.369}{383.809}{\Character{2}}%
\TextAt{13}{308.378}{383.809}{U}%
\TextAt{17}{317.577}{389.92}{\Character{40}}%
\TextAt{21}{320.87}{389.92}{n}%
\TextAt{17}{326.008}{389.92}{\Character{41}}%
\TextAt{25}{316.328}{380.524}{Q}%
\TextAt{22}{322.527}{379.528}{p}%
\NativeMathAt{333.253}{383.809}{11.955}{\subseteq}%
\TextAt{13}{346.005}{383.809}{N}%
\TextAt{21}{355.378}{381.95}{L}%
\TextAt{3}{361.137}{381.95}{\Character{106}}%
\TextAt{21}{363.49}{381.95}{K}%
\TextAt{13}{371.692}{383.809}{L}%
\TextAt{3}{379.657}{388.148}{\Character{3}}%
\TextAt{8}{384.389}{383.809}{.}%
\TextAt{8}{393.071}{383.809}{Therefore}%
\TextAt{13}{446.146}{383.809}{L}%
\TextAt{8}{458.085}{383.809}{is}%
\TextAt{8}{469.927}{383.809}{contained}%
\TextAt{8}{523.014}{383.809}{in}%
\TextAt{8}{79.226}{366.763}{the}%
\TextAt{8}{99.383}{366.763}{class}%
\TextAt{8}{126.82}{366.763}{field}%
\TextAt{13}{152.177}{366.763}{M}%
\TextAt{8}{168.648}{366.763}{of}%
\TextAt{8}{181.978}{366.763}{the}%
\TextAt{8}{202.134}{366.763}{group}%
\TextAt{9}{200.005}{338.176}{\Character{40}}%
\TextAt{13}{204.558}{338.176}{p}%
\TextAt{21}{210.433}{343.113}{f}%
\TextAt{9}{215.878}{338.176}{\Character{41}}%
\TextAt{16}{223.087}{338.176}{\Character{2}}%
\TextAt{13}{235.042}{338.176}{U}%
\TextAt{17}{244.241}{344.287}{\Character{40}}%
\TextAt{21}{247.534}{344.287}{n}%
\TextAt{17}{252.672}{344.287}{\Character{41}}%
\TextAt{25}{242.992}{334.891}{Q}%
\TextAt{22}{249.191}{333.895}{p}%
\TextAt{9}{259.785}{338.176}{\Character{61}}%
\TextAt{9}{272.195}{338.176}{\Character{40}\Character{40}}%
\TextAt{13}{281.315}{338.176}{p}%
\TextAt{21}{287.19}{343.113}{f}%
\TextAt{9}{292.635}{338.176}{\Character{41}}%
\TextAt{16}{299.844}{338.176}{\Character{2}}%
\TextAt{13}{311.799}{338.176}{U}%
\TextAt{25}{319.749}{336.383}{Q}%
\TextAt{22}{325.948}{335.387}{p}%
\TextAt{9}{330.785}{338.176}{\Character{41}}%
\TextAt{16}{337.994}{338.177}{\Character{92}}%
\TextAt{9}{348.621}{338.176}{\Character{40}\Character{40}}%
\TextAt{13}{357.725}{338.177}{p}%
\TextAt{9}{363.6}{338.176}{\Character{41}}%
\TextAt{16}{370.809}{338.177}{\Character{2}}%
\TextAt{13}{382.765}{338.176}{U}%
\TextAt{17}{391.964}{344.287}{\Character{40}}%
\TextAt{21}{395.257}{344.287}{n}%
\TextAt{17}{400.395}{344.287}{\Character{41}}%
\TextAt{25}{390.715}{334.891}{Q}%
\TextAt{22}{396.914}{333.895}{p}%
\TextAt{9}{404.186}{338.177}{\Character{41}}%
\TextAt{13}{408.739}{338.177}{\Character{58}}%
\TextAt{8}{79.226}{309.59}{By}%
\TextAt{8}{97.685}{309.59}{(1.4),}%
\TextAt{13}{129.032}{309.59}{M}%
\TextAt{8}{145.6}{309.59}{is}%
\TextAt{8}{157.46}{309.59}{the}%
\TextAt{8}{177.712}{309.59}{composite}%
\TextAt{8}{232.826}{309.59}{of}%
\TextAt{8}{246.252}{309.59}{the}%
\TextAt{8}{266.505}{309.59}{class}%
\TextAt{8}{294.038}{309.59}{field}%
\TextAt{8}{319.491}{309.59}{for}%
\TextAt{9}{337.472}{309.59}{\Character{40}}%
\TextAt{13}{342.024}{309.59}{p}%
\TextAt{21}{347.899}{313.928}{f}%
\TextAt{9}{353.344}{309.59}{\Character{41}}%
\TextAt{16}{360.619}{309.59}{\Character{2}}%
\TextAt{13}{372.639}{309.59}{U}%
\TextAt{25}{380.59}{307.796}{Q}%
\TextAt{22}{386.789}{306.8}{p}%
\TextAt{8}{391.625}{309.59}{—this}%
\TextAt{8}{426.248}{309.59}{being}%
\TextAt{8}{457.882}{309.59}{the}%
\TextAt{8}{478.134}{309.59}{unramified}%
\TextAt{8}{79.226}{292.544}{extension}%
\TextAt{8}{130.822}{292.544}{of}%
\TextAt{8}{143.982}{292.544}{degree}%
\TextAt{13}{180.227}{292.544}{f}%
\TextAt{8}{187.274}{292.544}{—and}%
\TextAt{8}{221.57}{292.544}{the}%
\TextAt{8}{241.557}{292.544}{class}%
\TextAt{8}{268.824}{292.544}{field}%
\TextAt{8}{294.011}{292.544}{for}%
\TextAt{9}{311.725}{292.544}{\Character{40}}%
\TextAt{13}{316.278}{292.544}{p}%
\TextAt{9}{322.153}{292.544}{\Character{41}}%
\TextAt{16}{329.014}{292.544}{\Character{2}}%
\TextAt{13}{340.622}{292.544}{U}%
\TextAt{17}{349.821}{298.654}{\Character{40}}%
\TextAt{21}{353.114}{298.654}{n}%
\TextAt{17}{358.252}{298.654}{\Character{41}}%
\TextAt{25}{348.572}{289.258}{Q}%
\TextAt{22}{354.771}{288.262}{p}%
\TextAt{8}{362.043}{292.544}{.}%
\TextAt{13}{370.435}{292.544}{M}%
\TextAt{8}{386.736}{292.544}{is}%
\TextAt{8}{398.33}{292.544}{therefore}%
\TextAt{8}{447.26}{292.544}{generated}%
\TextAt{8}{500.422}{292.544}{by}%
\TextAt{8}{516.51}{292.544}{the}%
\TextAt{9}{79.226}{275.498}{\Character{40}}%
\TextAt{13}{83.778}{275.498}{p}%
\TextAt{21}{89.654}{279.836}{f}%
\TextAt{16}{97.755}{275.498}{\Character{0}}%
\TextAt{9}{109.71}{275.498}{1\Character{41}}%
\TextAt{13}{120.116}{275.498}{p}%
\TextAt{21}{125.991}{279.836}{n}%
\TextAt{8}{131.627}{275.498}{-th}%
\TextAt{8}{150.48}{275.498}{roots}%
\TextAt{8}{180.141}{275.498}{of}%
\TextAt{8}{193.471}{275.498}{unity.}%
\Vector{q 0 G 0 g 0.399 w 452.334 -498.502 m 452.334 -490.432 l S Q}%
\Vector{q 0 G 0 g 0.399 w 452.534 -490.632 m 459.707 -490.632 l S Q}%
\Vector{q 0 G 0 g 0.399 w 452.534 -498.303 m 459.707 -498.303 l S Q}%
\Vector{q 0 G 0 g 0.399 w 459.906 -498.502 m 459.906 -490.432 l S Q}%
\TextAt{8}{97.159}{246.615}{From}%
\TextAt{8}{126.987}{246.615}{the}%
\TextAt{8}{146.247}{246.615}{local}%
\TextAt{8}{172.991}{246.615}{Kronecker–Weber}%
\TextAt{8}{266.384}{246.615}{theorem,}%
\TextAt{8}{314.444}{246.615}{one}%
\TextAt{8}{335.007}{246.615}{may}%
\TextAt{8}{359.479}{246.615}{readily}%
\TextAt{8}{397.282}{246.615}{deduce}%
\TextAt{8}{435.395}{246.615}{the}%
\TextAt{8}{454.654}{246.615}{global,}%
\TextAt{8}{491.668}{246.615}{classical}%
\TextAt{12}{79.226}{229.569}{Theorem}%
\TextAt{12}{136.24}{229.569}{of}%
\TextAt{12}{151.567}{229.569}{Kronecker–Weber}%
\TextAt{8}{257.011}{229.569}{.}%
\TextAt{12}{79.226}{203.675}{Theorem}%
\TextAt{12}{136.324}{203.675}{(1.10).}%
\TextAt{8}{178.443}{203.675}{Every}%
\TextAt{8}{212.169}{203.675}{finite}%
\TextAt{8}{242.153}{203.675}{abelian}%
\TextAt{8}{282.873}{203.675}{extension}%
\TextAt{13}{334.712}{203.675}{L}%
\TextAt{16}{342.676}{203.675}{\Character{106}}%
\TextAt{18}{345.997}{203.675}{Q}%
\TextAt{8}{359.265}{203.675}{is}%
\TextAt{8}{371.102}{203.675}{contained}%
\TextAt{8}{424.183}{203.675}{in}%
\TextAt{8}{437.909}{203.675}{a}%
\TextAt{8}{447.737}{203.675}{field}%
\TextAt{18}{473.166}{203.675}{Q}%
\TextAt{9}{482.464}{203.675}{\Character{40}}%
\TextAt{13}{487.017}{203.675}{\Character{16}}%
\TextAt{9}{493.032}{203.675}{\Character{41}}%
\TextAt{8}{501.555}{203.675}{gener-}%
\TextAt{8}{79.226}{186.629}{ated}%
\TextAt{8}{105.241}{186.629}{by}%
\TextAt{8}{121.5}{186.629}{a}%
\TextAt{8}{131.255}{186.629}{root}%
\TextAt{8}{156.301}{186.629}{of}%
\TextAt{8}{169.631}{186.629}{unity}%
\TextAt{13}{200.201}{186.629}{\Character{16}}%
\TextAt{8}{206.216}{186.629}{.}%
\TextAt{19}{79.226}{151.887}{Proof.}%
\TextAt{8}{116}{151.887}{Let}%
\TextAt{13}{137.088}{151.887}{S}%
\TextAt{8}{148.999}{151.887}{be}%
\TextAt{8}{165.042}{151.887}{the}%
\TextAt{8}{185.316}{151.887}{set}%
\TextAt{8}{203.702}{151.887}{of}%
\TextAt{8}{217.151}{151.887}{all}%
\TextAt{8}{233.528}{151.887}{prime}%
\TextAt{8}{266.81}{151.887}{numbers}%
\TextAt{13}{314.139}{151.887}{p}%
\TextAt{8}{324.03}{151.887}{that}%
\TextAt{8}{349.517}{151.887}{are}%
\TextAt{8}{369.146}{151.887}{ramified}%
\TextAt{8}{414.79}{151.887}{in}%
\TextAt{13}{428.561}{151.887}{L}%
\TextAt{8}{436.526}{151.887}{,}%
\TextAt{8}{443.823}{151.887}{and}%
\TextAt{8}{466.704}{151.887}{let}%
\TextAt{13}{483.727}{151.887}{L}%
\TextAt{21}{491.691}{150.094}{p}%
\TextAt{8}{500.468}{151.887}{be}%
\TextAt{8}{516.51}{151.887}{the}%
\TextAt{8}{79.226}{134.841}{completion}%
\TextAt{8}{139.611}{134.841}{of}%
\TextAt{13}{153.49}{134.841}{L}%
\TextAt{8}{165.901}{134.841}{with}%
\TextAt{8}{193.111}{134.841}{respect}%
\TextAt{8}{233.71}{134.841}{to}%
\TextAt{8}{248.569}{134.841}{some}%
\TextAt{8}{278.445}{134.841}{prime}%
\TextAt{8}{312.157}{134.841}{lying}%
\TextAt{8}{341.65}{134.841}{above}%
\TextAt{13}{375.375}{134.841}{p}%
\TextAt{8}{381.25}{134.841}{.}%
\TextAt{8}{391.346}{134.841}{Then}%
\TextAt{13}{422.453}{134.841}{L}%
\TextAt{21}{430.417}{133.048}{p}%
\TextAt{16}{435.178}{134.841}{\Character{106}}%
\TextAt{18}{438.499}{134.841}{Q}%
\TextAt{21}{447.798}{133.048}{p}%
\TextAt{8}{457.005}{134.841}{is}%
\TextAt{8}{469.318}{134.841}{abelian,}%
\TextAt{8}{513.904}{134.841}{and}%
\TextAt{8}{79.226}{117.795}{therefore}%
\TextAt{13}{128.593}{117.795}{L}%
\TextAt{21}{136.557}{116.002}{p}%
\NativeMathAt{145.094}{117.795}{11.955}{\subseteq}%
\TextAt{18}{158.168}{117.795}{Q}%
\TextAt{21}{167.466}{116.002}{p}%
\TextAt{9}{172.227}{117.795}{\Character{40}}%
\TextAt{13}{176.78}{117.795}{\Character{22}}%
\TextAt{21}{183.823}{116.002}{n}%
\TextAt{22}{188.961}{115.006}{p}%
\TextAt{9}{193.797}{117.795}{\Character{41}}%
\TextAt{8}{198.35}{117.795}{,}%
\TextAt{8}{205.833}{117.795}{for}%
\TextAt{8}{223.984}{117.795}{a}%
\TextAt{8}{234.007}{117.795}{suitable}%
\TextAt{13}{277.91}{117.795}{n}%
\TextAt{21}{284.897}{116.002}{p}%
\TextAt{8}{289.658}{117.795}{.}%
\TextAt{8}{298.907}{117.795}{Let}%
\TextAt{13}{320.143}{117.795}{p}%
\TextAt{21}{326.018}{122.134}{e}%
\TextAt{22}{329.947}{121.137}{p}%
\TextAt{8}{338.948}{117.795}{be}%
\TextAt{8}{355.139}{117.795}{the}%
\TextAt{8}{375.562}{117.795}{precise}%
\TextAt{8}{414.253}{117.795}{power}%
\TextAt{8}{448.664}{117.795}{of}%
\TextAt{13}{462.26}{117.795}{p}%
\TextAt{8}{472.3}{117.795}{dividing}%
\TextAt{13}{517.769}{117.795}{n}%
\TextAt{21}{524.757}{116.002}{p}%
\TextAt{8}{529.518}{117.795}{,}%
\TextAt{8}{79.226}{100.749}{and}%
\TextAt{8}{101.989}{100.749}{let}%
\TextAt{13}{276.937}{83.703}{n}%
\TextAt{9}{287.245}{83.703}{\Character{61}}%
\TextAt{20}{299.785}{95.061}{\Character{89}}%
\TextAt{21}{299.671}{69.689}{p}%
\NativeMathAt{303.933}{69.689}{7.97}{\in}%
\TextAt{21}{309.579}{69.689}{S}%
\TextAt{13}{317.167}{83.703}{p}%
\TextAt{21}{323.042}{88.639}{e}%
\TextAt{22}{326.971}{87.643}{p}%
\TextAt{13}{331.807}{83.703}{\Character{58}}%
\TextAt{8}{297.211}{43.033}{297}%
\Destination{78.23}{809.815}{nav.page.297}%
\Destination{79.36}{724.955}{nav.print.324}%
\Destination{79.36}{166.855}{nav.print.325}%
\Destination{79.36}{469.655}{nav.number.5.1.9}%
\Destination{79.36}{217.655}{nav.number.5.1.10}%
\Hotlink{97.185}{307.355}{25.078}{11.047}{nav.number.5.1.4}%
}

% Page 310 (body page 298)
\Page{%
\TextAt{2}{79.226}{796.943}{V.}%
\TextAt{2}{94.073}{796.943}{Local}%
\TextAt{2}{123.615}{796.943}{Class}%
\TextAt{2}{152.208}{796.943}{Field}%
\TextAt{2}{179.939}{796.943}{Theory}%
\Vector{q 0 G 0 g 0.299 w 7.226 18.729 m 460.769 18.729 l S Q}%
\TextAt{8}{79.226}{761.124}{We}%
\TextAt{8}{99.359}{761.124}{will}%
\TextAt{8}{121.453}{761.124}{show}%
\TextAt{8}{150.445}{761.124}{that}%
\TextAt{13}{175.803}{761.124}{L}%
\NativeMathAt{187.088}{761.124}{11.955}{\subseteq}%
\TextAt{18}{199.707}{761.124}{Q}%
\TextAt{9}{209.006}{761.124}{\Character{40}}%
\TextAt{13}{213.558}{761.124}{\Character{22}}%
\TextAt{21}{220.601}{759.33}{n}%
\TextAt{9}{226.237}{761.124}{\Character{41}}%
\TextAt{8}{230.79}{761.124}{.}%
\TextAt{8}{239.234}{761.124}{For}%
\TextAt{8}{260.192}{761.124}{this}%
\TextAt{8}{283.004}{761.124}{let}%
\TextAt{13}{299.897}{761.124}{M}%
\TextAt{9}{315.791}{761.124}{\Character{61}}%
\TextAt{13}{328.217}{761.124}{L}%
\TextAt{9}{336.181}{761.124}{\Character{40}}%
\TextAt{13}{340.734}{761.124}{\Character{22}}%
\TextAt{21}{347.777}{759.33}{n}%
\TextAt{9}{353.413}{761.124}{\Character{41}}%
\TextAt{8}{357.965}{761.124}{.}%
\TextAt{8}{366.41}{761.124}{Then}%
\TextAt{13}{396.956}{761.124}{M}%
\TextAt{16}{409.53}{761.124}{\Character{106}}%
\TextAt{18}{412.85}{761.124}{Q}%
\TextAt{8}{426.035}{761.124}{is}%
\TextAt{8}{437.788}{761.124}{abelian,}%
\TextAt{8}{481.678}{761.124}{and}%
\TextAt{8}{504.429}{761.124}{if}%
\TextAt{13}{515.142}{761.124}{p}%
\TextAt{8}{524.903}{761.124}{is}%
\TextAt{8}{79.226}{744.078}{ramified}%
\TextAt{8}{124.431}{744.078}{in}%
\TextAt{13}{137.763}{744.078}{M}%
\TextAt{16}{150.337}{744.078}{\Character{106}}%
\TextAt{18}{153.658}{744.078}{Q}%
\TextAt{8}{162.956}{744.078}{,}%
\TextAt{8}{169.849}{744.078}{then}%
\TextAt{13}{196.188}{744.078}{p}%
\TextAt{8}{205.64}{744.078}{must}%
\TextAt{8}{234.323}{744.078}{lie}%
\TextAt{8}{249.604}{744.078}{in}%
\TextAt{13}{262.936}{744.078}{S}%
\TextAt{8}{270.831}{744.078}{.}%
\TextAt{8}{279.173}{744.078}{If}%
\TextAt{13}{290.544}{744.078}{M}%
\TextAt{21}{301.869}{742.284}{p}%
\TextAt{8}{310.207}{744.078}{is}%
\TextAt{8}{321.65}{744.078}{the}%
\TextAt{8}{341.486}{744.078}{completion}%
\TextAt{8}{401.001}{744.078}{with}%
\TextAt{8}{427.34}{744.078}{respect}%
\TextAt{8}{467.069}{744.078}{to}%
\TextAt{8}{481.059}{744.078}{a}%
\TextAt{8}{490.494}{744.078}{prime}%
\TextAt{8}{523.337}{744.078}{of}%
\TextAt{13}{79.226}{727.032}{M}%
\TextAt{8}{95.697}{727.032}{above}%
\TextAt{13}{128.873}{727.032}{p}%
\TextAt{8}{138.645}{727.032}{whose}%
\TextAt{8}{173.172}{727.032}{restriction}%
\TextAt{8}{229.17}{727.032}{to}%
\TextAt{13}{243.48}{727.032}{L}%
\TextAt{8}{255.342}{727.032}{gives}%
\TextAt{8}{284.022}{727.032}{the}%
\TextAt{8}{304.179}{727.032}{completion}%
\TextAt{13}{364.015}{727.032}{L}%
\TextAt{21}{371.979}{725.238}{p}%
\TextAt{8}{376.74}{727.032}{,}%
\TextAt{8}{383.889}{727.032}{then}%
\TextAt{13}{188.954}{702.477}{M}%
\TextAt{21}{200.279}{700.684}{p}%
\TextAt{9}{208.361}{702.477}{\Character{61}}%
\TextAt{13}{220.786}{702.477}{L}%
\TextAt{21}{228.751}{700.684}{p}%
\TextAt{9}{233.512}{702.477}{\Character{40}}%
\TextAt{13}{238.064}{702.477}{\Character{22}}%
\TextAt{21}{245.107}{700.684}{n}%
\TextAt{9}{250.743}{702.477}{\Character{41}}%
\TextAt{9}{258.598}{702.477}{\Character{61}}%
\TextAt{18}{271.042}{702.477}{Q}%
\TextAt{21}{280.341}{700.684}{p}%
\TextAt{9}{285.101}{702.477}{\Character{40}}%
\TextAt{13}{289.654}{702.477}{\Character{22}}%
\TextAt{21}{296.697}{700.684}{p}%
\TextAt{22}{300.96}{703.7}{e}%
\TextAt{22}{304.48}{702.704}{p}%
\TextAt{21}{309.317}{700.684}{n}%
\TextAt{23}{314.455}{702.953}{\Character{48}}%
\TextAt{9}{317.638}{702.477}{\Character{41}}%
\TextAt{9}{325.493}{702.477}{\Character{61}}%
\TextAt{18}{337.936}{702.477}{Q}%
\TextAt{21}{347.235}{700.684}{p}%
\TextAt{9}{351.996}{702.477}{\Character{40}}%
\TextAt{13}{356.548}{702.477}{\Character{22}}%
\TextAt{21}{363.591}{700.684}{p}%
\TextAt{22}{367.854}{703.7}{e}%
\TextAt{22}{371.375}{702.704}{p}%
\TextAt{9}{376.71}{702.477}{\Character{41}}%
\TextAt{18}{381.262}{702.477}{Q}%
\TextAt{21}{390.56}{700.684}{p}%
\TextAt{9}{395.321}{702.477}{\Character{40}}%
\TextAt{13}{399.874}{702.477}{\Character{22}}%
\TextAt{21}{406.917}{700.684}{n}%
\TextAt{23}{412.055}{702.953}{\Character{48}}%
\TextAt{9}{415.237}{702.477}{\Character{41}}%
\TextAt{13}{419.79}{702.477}{\Character{59}}%
\TextAt{8}{79.226}{677.923}{with}%
\TextAt{9}{106.854}{677.923}{\Character{40}}%
\TextAt{13}{111.406}{677.923}{n}%
\TextAt{3}{118.394}{682.261}{\Character{48}}%
\TextAt{13}{121.189}{677.923}{\Character{59}}%
\TextAt{13}{126.413}{677.923}{p}%
\TextAt{9}{132.308}{677.923}{\Character{41}}%
\TextAt{9}{141.812}{677.923}{\Character{61}}%
\TextAt{9}{155.884}{677.923}{1}%
\TextAt{8}{161.758}{677.923}{.}%
\TextAt{18}{173.109}{677.923}{Q}%
\TextAt{21}{182.407}{676.129}{p}%
\TextAt{9}{187.168}{677.923}{\Character{40}}%
\TextAt{13}{191.72}{677.923}{\Character{22}}%
\TextAt{21}{198.763}{676.129}{n}%
\TextAt{23}{203.902}{678.399}{\Character{48}}%
\TextAt{9}{207.084}{677.923}{\Character{41}}%
\TextAt{16}{211.636}{677.923}{\Character{106}}%
\TextAt{18}{214.957}{677.923}{Q}%
\TextAt{21}{224.256}{676.129}{p}%
\TextAt{8}{233.882}{677.923}{is}%
\TextAt{8}{246.613}{677.923}{the}%
\TextAt{8}{267.737}{677.923}{maximal}%
\TextAt{8}{316.513}{677.923}{unramified}%
\TextAt{8}{376.014}{677.923}{subextension}%
\TextAt{8}{446.692}{677.923}{of}%
\TextAt{18}{460.989}{677.923}{Q}%
\TextAt{21}{470.288}{676.129}{p}%
\TextAt{9}{475.049}{677.923}{\Character{40}}%
\TextAt{13}{479.601}{677.923}{\Character{22}}%
\TextAt{21}{486.644}{676.129}{p}%
\TextAt{22}{490.907}{679.146}{e}%
\TextAt{22}{494.428}{678.15}{p}%
\TextAt{21}{499.264}{676.129}{n}%
\TextAt{23}{504.403}{678.399}{\Character{48}}%
\TextAt{9}{507.585}{677.923}{\Character{41}}%
\TextAt{16}{512.137}{677.923}{\Character{106}}%
\TextAt{18}{515.458}{677.923}{Q}%
\TextAt{21}{524.757}{676.129}{p}%
\TextAt{8}{529.518}{677.923}{.}%
\TextAt{8}{79.226}{660.877}{The}%
\TextAt{8}{104.049}{660.877}{inertia}%
\TextAt{8}{141.89}{660.877}{group}%
\TextAt{13}{175.835}{660.877}{I}%
\TextAt{21}{180.997}{659.083}{p}%
\TextAt{8}{190.424}{660.877}{of}%
\TextAt{13}{204.523}{660.877}{M}%
\TextAt{16}{217.097}{660.877}{\Character{106}}%
\TextAt{18}{220.418}{660.877}{Q}%
\TextAt{8}{234.382}{660.877}{is}%
\TextAt{8}{246.915}{660.877}{therefore}%
\TextAt{8}{296.783}{660.877}{isomorphic}%
\TextAt{8}{356.802}{660.877}{to}%
\TextAt{8}{371.881}{660.877}{the}%
\TextAt{8}{392.806}{660.877}{group}%
\TextAt{13}{426.75}{660.877}{G}%
\TextAt{9}{435.984}{660.877}{\Character{40}}%
\TextAt{18}{440.536}{660.877}{Q}%
\TextAt{21}{449.835}{659.083}{p}%
\TextAt{9}{454.596}{660.877}{\Character{40}}%
\TextAt{13}{459.148}{660.877}{\Character{22}}%
\TextAt{21}{466.191}{659.083}{p}%
\TextAt{22}{470.454}{662.1}{e}%
\TextAt{22}{473.974}{661.104}{p}%
\TextAt{9}{479.309}{660.877}{\Character{41}}%
\TextAt{16}{483.862}{660.877}{\Character{106}}%
\TextAt{18}{487.182}{660.877}{Q}%
\TextAt{21}{496.481}{659.083}{p}%
\TextAt{9}{501.242}{660.877}{\Character{41}}%
\TextAt{8}{505.794}{660.877}{,}%
\TextAt{8}{513.904}{660.877}{and}%
\TextAt{8}{79.226}{643.831}{consequently}%
\TextAt{8}{149.794}{643.831}{has}%
\TextAt{8}{171.908}{643.831}{order}%
\TextAt{13}{203.718}{643.831}{\Character{39}}%
\TextAt{9}{211.387}{643.831}{\Character{40}}%
\TextAt{13}{215.939}{643.831}{p}%
\TextAt{21}{221.814}{648.169}{e}%
\TextAt{22}{225.742}{647.173}{p}%
\TextAt{9}{230.579}{643.831}{\Character{41}}%
\TextAt{8}{235.131}{643.831}{,}%
\TextAt{8}{243.831}{643.831}{where}%
\TextAt{13}{278.88}{643.831}{\Character{39}}%
\TextAt{8}{291.686}{643.831}{is}%
\TextAt{8}{304.691}{643.831}{Euler’s}%
\TextAt{8}{345.168}{643.831}{function.}%
\TextAt{8}{399.288}{643.831}{Let}%
\TextAt{13}{421.497}{643.831}{I}%
\TextAt{8}{432.738}{643.831}{be}%
\TextAt{8}{449.902}{643.831}{the}%
\TextAt{8}{471.299}{643.831}{subgroup}%
\TextAt{8}{523.337}{643.831}{of}%
\TextAt{13}{79.226}{626.785}{G}%
\TextAt{9}{88.46}{626.785}{\Character{40}}%
\TextAt{13}{93.012}{626.785}{M}%
\TextAt{16}{105.586}{626.785}{\Character{106}}%
\TextAt{18}{108.907}{626.785}{Q}%
\TextAt{9}{118.205}{626.785}{\Character{41}}%
\TextAt{8}{126.487}{626.785}{generated}%
\TextAt{8}{179.651}{626.785}{by}%
\TextAt{8}{195.743}{626.785}{all}%
\TextAt{13}{211.834}{626.785}{I}%
\TextAt{21}{216.997}{624.991}{p}%
\TextAt{8}{221.758}{626.785}{,}%
\TextAt{13}{228.773}{626.785}{p}%
\NativeMathAt{237.969}{626.785}{11.955}{\in}%
\TextAt{13}{249.26}{626.785}{S}%
\TextAt{8}{257.155}{626.785}{.}%
\TextAt{8}{265.548}{626.785}{The}%
\TextAt{8}{289.434}{626.785}{fixed}%
\TextAt{8}{317.552}{626.785}{field}%
\TextAt{8}{342.741}{626.785}{of}%
\TextAt{13}{355.903}{626.785}{I}%
\TextAt{8}{365.736}{626.785}{is}%
\TextAt{8}{377.332}{626.785}{then}%
\TextAt{8}{403.824}{626.785}{unramified,}%
\TextAt{8}{465.475}{626.785}{and}%
\TextAt{8}{488.069}{626.785}{hence}%
\TextAt{8}{520.408}{626.785}{by}%
\TextAt{8}{79.226}{609.739}{Minkowski’s}%
\TextAt{8}{145.234}{609.739}{theorem}%
\TextAt{8}{190.631}{609.739}{from}%
\TextAt{8}{218.142}{609.739}{chap.}%
\TextAt{8}{250.291}{609.739}{III,}%
\TextAt{8}{270.618}{609.739}{(2.18),}%
\TextAt{8}{307.6}{609.739}{it}%
\TextAt{8}{319.176}{609.739}{equals}%
\TextAt{18}{354.554}{609.739}{Q}%
\TextAt{8}{363.852}{609.739}{,}%
\TextAt{8}{370.899}{609.739}{i.e.,}%
\TextAt{13}{392.901}{609.739}{I}%
\TextAt{9}{402.325}{609.739}{\Character{61}}%
\TextAt{13}{414.75}{609.739}{G}%
\TextAt{9}{423.984}{609.739}{\Character{40}}%
\TextAt{13}{428.536}{609.739}{M}%
\TextAt{16}{441.11}{609.739}{\Character{106}}%
\TextAt{18}{444.431}{609.739}{Q}%
\TextAt{9}{453.729}{609.739}{\Character{41}}%
\TextAt{8}{458.282}{609.739}{.}%
\TextAt{8}{466.687}{609.739}{On}%
\TextAt{8}{486.07}{609.739}{the}%
\TextAt{8}{506.098}{609.739}{other}%
\TextAt{8}{79.226}{592.693}{hand}%
\TextAt{8}{108.492}{592.693}{we}%
\TextAt{8}{125.72}{592.693}{have}%
\TextAt{8}{181.541}{575.647}{\char35\relax }%
\TextAt{13}{191.297}{575.647}{I}%
\NativeMathAt{200.72}{575.647}{11.955}{\leq}%
\TextAt{20}{213.454}{587.004}{\Character{89}}%
\TextAt{21}{213.34}{561.633}{p}%
\NativeMathAt{217.603}{561.633}{7.97}{\in}%
\TextAt{21}{223.248}{561.633}{S}%
\TextAt{8}{230.836}{575.647}{\char35\relax }%
\TextAt{13}{240.592}{575.647}{I}%
\TextAt{21}{245.755}{573.854}{p}%
\TextAt{9}{253.836}{575.647}{\Character{61}}%
\TextAt{20}{266.376}{587.004}{\Character{89}}%
\TextAt{21}{266.262}{561.633}{p}%
\NativeMathAt{270.525}{561.633}{7.97}{\in}%
\TextAt{21}{276.17}{561.633}{S}%
\TextAt{13}{283.759}{575.647}{\Character{39}}%
\TextAt{9}{291.427}{575.647}{\Character{40}}%
\TextAt{13}{295.979}{575.647}{p}%
\TextAt{21}{301.854}{580.583}{e}%
\TextAt{22}{305.783}{579.587}{p}%
\TextAt{9}{310.619}{575.647}{\Character{41}}%
\TextAt{9}{318.474}{575.647}{\Character{61}}%
\TextAt{13}{330.918}{575.647}{\Character{39}}%
\TextAt{9}{338.586}{575.647}{\Character{40}}%
\TextAt{13}{343.139}{575.647}{n}%
\TextAt{9}{350.126}{575.647}{\Character{41}}%
\TextAt{9}{357.981}{575.647}{\Character{61}}%
\TextAt{9}{370.402}{575.647}{\Character{91}}%
\TextAt{18}{373.677}{575.647}{Q}%
\TextAt{9}{382.975}{575.647}{\Character{40}}%
\TextAt{13}{387.527}{575.647}{\Character{22}}%
\TextAt{21}{394.57}{573.854}{n}%
\TextAt{9}{400.207}{575.647}{\Character{41}}%
\TextAt{9}{408.062}{575.647}{\Character{58}}%
\TextAt{18}{414.652}{575.647}{Q}%
\TextAt{9}{423.951}{575.647}{\Character{93}}%
\TextAt{13}{427.203}{575.647}{\Character{59}}%
\TextAt{8}{79.226}{544.355}{and}%
\TextAt{8}{100.876}{544.355}{therefore}%
\TextAt{9}{148.863}{544.355}{\Character{91}}%
\TextAt{13}{152.115}{544.355}{M}%
\TextAt{9}{168.009}{544.355}{\Character{58}}%
\TextAt{18}{174.582}{544.355}{Q}%
\TextAt{9}{183.88}{544.355}{\Character{93}}%
\TextAt{9}{190.431}{544.355}{\Character{61}}%
\TextAt{9}{202.853}{544.355}{\Character{91}}%
\TextAt{18}{206.13}{544.355}{Q}%
\TextAt{9}{215.429}{544.355}{\Character{40}}%
\TextAt{13}{219.981}{544.355}{\Character{22}}%
\TextAt{21}{227.024}{542.561}{n}%
\TextAt{9}{232.66}{544.355}{\Character{41}}%
\TextAt{9}{240.515}{544.355}{\Character{58}}%
\TextAt{18}{247.106}{544.355}{Q}%
\TextAt{9}{256.404}{544.355}{\Character{93}}%
\TextAt{8}{259.656}{544.355}{,}%
\TextAt{8}{265.915}{544.355}{so}%
\TextAt{8}{279.172}{544.355}{that}%
\TextAt{13}{303.429}{544.355}{M}%
\TextAt{9}{319.323}{544.355}{\Character{61}}%
\TextAt{18}{331.748}{544.355}{Q}%
\TextAt{9}{341.047}{544.355}{\Character{40}}%
\TextAt{13}{345.599}{544.355}{\Character{22}}%
\TextAt{21}{352.642}{542.561}{n}%
\TextAt{9}{358.279}{544.355}{\Character{41}}%
\TextAt{8}{362.831}{544.355}{.}%
\TextAt{8}{370.908}{544.355}{This}%
\TextAt{8}{396.515}{544.355}{shows}%
\TextAt{8}{429.021}{544.355}{that}%
\TextAt{13}{453.277}{544.355}{L}%
\NativeMathAt{464.562}{544.355}{11.955}{\subseteq}%
\TextAt{18}{477.182}{544.355}{Q}%
\TextAt{9}{486.48}{544.355}{\Character{40}}%
\TextAt{13}{491.032}{544.355}{\Character{22}}%
\TextAt{21}{498.075}{542.561}{n}%
\TextAt{9}{503.712}{544.355}{\Character{41}}%
\TextAt{8}{508.264}{544.355}{.}%
\Vector{q 0 G 0 g 0.399 w 452.334 -229.645 m 452.334 -221.576 l S Q}%
\Vector{q 0 G 0 g 0.399 w 452.534 -221.775 m 459.707 -221.775 l S Q}%
\Vector{q 0 G 0 g 0.399 w 452.534 -229.446 m 459.707 -229.446 l S Q}%
\Vector{q 0 G 0 g 0.399 w 459.906 -229.645 m 459.906 -221.576 l S Q}%
\TextAt{8}{97.159}{516.624}{The}%
\TextAt{8}{121.213}{516.624}{following}%
\TextAt{8}{170.647}{516.624}{exercises}%
\TextAt{8}{218.564}{516.624}{1–3}%
\TextAt{8}{240.035}{516.624}{presuppose}%
\TextAt{8}{300.313}{516.624}{exercises}%
\TextAt{8}{348.229}{516.624}{4–8}%
\TextAt{8}{369.701}{516.624}{of}%
\TextAt{8}{383.031}{516.624}{chap.}%
\TextAt{8}{415.222}{516.624}{IV,}%
\TextAt{8}{435.367}{516.624}{§3.}%
\TextAt{12}{79.226}{491.882}{Exercise}%
\TextAt{12}{132.427}{491.882}{V.1.1.}%
\TextAt{8}{171.173}{491.882}{For}%
\TextAt{8}{192.143}{491.882}{the}%
\TextAt{8}{212.299}{491.882}{Galois}%
\TextAt{8}{248.212}{491.882}{group}%
\TextAt{9}{281.388}{491.882}{\Character{0}}%
\TextAt{9}{292.004}{491.882}{\Character{61}}%
\TextAt{13}{304.449}{491.882}{G}%
\TextAt{9}{313.682}{491.882}{\Character{40}}%
\TextAt{20}{320.973}{494.904}{\Character{101}}%
\TextAt{13}{318.234}{491.882}{K}%
\TextAt{16}{329.052}{491.882}{\Character{106}}%
\TextAt{13}{332.373}{491.882}{K}%
\TextAt{9}{343.191}{491.882}{\Character{41}}%
\TextAt{8}{347.743}{491.882}{,}%
\TextAt{8}{354.892}{491.882}{one}%
\TextAt{8}{376.352}{491.882}{has}%
\TextAt{8}{397.226}{491.882}{canonically}%
\TextAt{13}{163.214}{467.327}{H}%
\TextAt{17}{173.837}{472.264}{1}%
\TextAt{9}{178.569}{467.327}{\Character{40}\Character{0}}%
\TextAt{13}{190.436}{467.327}{\Character{59}}%
\TextAt{18}{195.68}{467.327}{Z}%
\TextAt{9}{203.65}{467.327}{\Character{47}}%
\TextAt{13}{209.503}{467.327}{n}%
\TextAt{18}{216.491}{467.327}{Z}%
\TextAt{9}{224.461}{467.327}{\Character{41}}%
\NativeMathAt{232.334}{467.327}{11.955}{\cong}%
\TextAt{18}{244.953}{467.327}{Z}%
\TextAt{9}{252.924}{467.327}{\Character{47}}%
\TextAt{13}{258.777}{467.327}{n}%
\TextAt{18}{265.764}{467.327}{Z}%
\TextAt{8}{285.689}{467.327}{and}%
\TextAt{13}{316.51}{467.327}{H}%
\TextAt{17}{327.132}{472.264}{1}%
\TextAt{9}{331.865}{467.327}{\Character{40}\Character{0}}%
\TextAt{13}{343.731}{467.327}{\Character{59}}%
\TextAt{13}{348.955}{467.327}{\Character{22}}%
\TextAt{21}{356.018}{465.534}{n}%
\TextAt{9}{361.655}{467.327}{\Character{41}}%
\NativeMathAt{369.528}{467.327}{11.955}{\cong}%
\TextAt{13}{382.147}{467.327}{U}%
\TextAt{21}{390.098}{465.534}{K}%
\TextAt{13}{398.3}{467.327}{K}%
\TextAt{3}{409.118}{472.264}{\Character{3}}%
\TextAt{21}{413.352}{472.264}{n}%
\TextAt{9}{418.988}{467.327}{\Character{47}}%
\TextAt{13}{424.841}{467.327}{K}%
\TextAt{3}{435.659}{472.264}{\Character{3}}%
\TextAt{21}{439.893}{472.264}{n}%
\TextAt{13}{445.53}{467.327}{\Character{59}}%
\TextAt{8}{79.226}{442.773}{the}%
\TextAt{8}{99.383}{442.773}{latter}%
\TextAt{8}{131.255}{442.773}{provided}%
\TextAt{8}{179.387}{442.773}{that}%
\TextAt{13}{204.755}{442.773}{n}%
\TextAt{8}{215.64}{442.773}{is}%
\TextAt{8}{227.404}{442.773}{not}%
\TextAt{8}{248.218}{442.773}{divisible}%
\TextAt{8}{294.126}{442.773}{by}%
\TextAt{8}{310.385}{442.773}{the}%
\TextAt{8}{330.542}{442.773}{residue}%
\TextAt{8}{370.269}{442.773}{characteristic.}%
\TextAt{12}{79.226}{413.324}{Exercise}%
\TextAt{12}{132.427}{413.324}{V.1.2.}%
\TextAt{8}{171.173}{413.324}{For}%
\TextAt{8}{192.143}{413.324}{an}%
\TextAt{8}{208.402}{413.324}{arbitrary}%
\TextAt{8}{258.171}{413.324}{field}%
\TextAt{13}{283.528}{413.324}{K}%
\TextAt{8}{298.243}{413.324}{and}%
\TextAt{8}{321.006}{413.324}{a}%
\TextAt{13}{330.761}{413.324}{G}%
\TextAt{21}{339.995}{411.531}{K}%
\TextAt{8}{348.197}{413.324}{-module}%
\TextAt{13}{393.388}{413.324}{A}%
\TextAt{8}{402.163}{413.324}{,}%
\TextAt{8}{409.313}{413.324}{put}%
\TextAt{13}{244.218}{388.769}{H}%
\TextAt{17}{254.841}{393.705}{1}%
\TextAt{9}{259.573}{388.769}{\Character{40}}%
\TextAt{13}{264.125}{388.769}{K}%
\TextAt{13}{274.287}{388.769}{\Character{59}}%
\TextAt{13}{279.523}{388.769}{A}%
\TextAt{9}{288.312}{388.769}{\Character{41}}%
\TextAt{9}{296.167}{388.769}{\Character{61}}%
\TextAt{13}{308.611}{388.769}{H}%
\TextAt{17}{319.233}{393.705}{1}%
\TextAt{9}{323.966}{388.769}{\Character{40}}%
\TextAt{13}{328.518}{388.769}{G}%
\TextAt{21}{337.752}{386.976}{K}%
\TextAt{13}{345.954}{388.769}{\Character{59}}%
\TextAt{13}{351.178}{388.769}{A}%
\TextAt{9}{359.974}{388.769}{\Character{41}}%
\TextAt{13}{364.526}{388.769}{\Character{58}}%
\TextAt{8}{79.226}{364.215}{If}%
\TextAt{13}{90.961}{364.215}{K}%
\TextAt{8}{105.719}{364.215}{is}%
\TextAt{8}{117.526}{364.215}{a}%
\TextAt{14}{127.324}{364.215}{p}%
\TextAt{8}{133.307}{364.215}{-adic}%
\TextAt{8}{161.959}{364.215}{number}%
\TextAt{8}{204.598}{364.215}{field}%
\TextAt{8}{229.997}{364.215}{and}%
\TextAt{13}{252.803}{364.215}{n}%
\TextAt{8}{263.731}{364.215}{a}%
\TextAt{8}{273.529}{364.215}{natural}%
\TextAt{8}{314.554}{364.215}{number,}%
\TextAt{8}{360.456}{364.215}{then}%
\TextAt{8}{387.159}{364.215}{there}%
\TextAt{8}{417.113}{364.215}{exists}%
\TextAt{8}{449.471}{364.215}{a}%
\TextAt{8}{459.269}{364.215}{nondegenerate}%
\TextAt{8}{79.226}{347.169}{pairing}%
\TextAt{13}{210.749}{330.123}{H}%
\TextAt{17}{221.372}{335.059}{1}%
\TextAt{9}{226.104}{330.123}{\Character{40}}%
\TextAt{13}{230.656}{330.123}{K}%
\TextAt{13}{240.818}{330.123}{\Character{59}}%
\TextAt{18}{246.068}{330.123}{Z}%
\TextAt{9}{254.038}{330.123}{\Character{47}}%
\TextAt{13}{259.891}{330.123}{n}%
\TextAt{18}{266.878}{330.123}{Z}%
\TextAt{9}{274.849}{330.123}{\Character{41}}%
\TextAt{16}{282.058}{330.123}{\Character{2}}%
\TextAt{13}{294.013}{330.123}{H}%
\TextAt{17}{304.635}{335.059}{1}%
\TextAt{9}{309.368}{330.123}{\Character{40}}%
\TextAt{13}{313.92}{330.123}{K}%
\TextAt{13}{324.082}{330.123}{\Character{59}}%
\TextAt{13}{329.318}{330.123}{\Character{22}}%
\TextAt{21}{336.374}{328.33}{n}%
\TextAt{9}{342.011}{330.123}{\Character{41}}%
\NativeArrowAt{349.884}{330.123}{11.955}{19.26}{\longrightarrow}%
\TextAt{18}{372.466}{330.123}{Z}%
\TextAt{9}{380.436}{330.123}{\Character{47}}%
\TextAt{13}{386.289}{330.123}{n}%
\TextAt{18}{393.277}{330.123}{Z}%
\TextAt{8}{79.226}{308.507}{of}%
\TextAt{8}{92.556}{308.507}{finite}%
\TextAt{8}{122.468}{308.507}{groups}%
\TextAt{8}{160.258}{308.507}{given}%
\TextAt{8}{190.828}{308.507}{by}%
\TextAt{9}{246.563}{291.461}{\Character{40}}%
\TextAt{13}{251.115}{291.461}{\Character{31}\Character{59}}%
\TextAt{13}{263.668}{291.461}{a}%
\TextAt{9}{269.838}{291.461}{\Character{41}}%
\NativeArrowAt{277.711}{291.461}{11.955}{19.26}{\longmapsto}%
\TextAt{13}{300.293}{291.461}{\Character{31}}%
\TextAt{9}{307.626}{291.461}{\Character{40}\Character{40}}%
\TextAt{13}{316.731}{291.461}{a\Character{59}}%
\TextAt{9}{331.253}{294.483}{\Character{22}}%
\TextAt{13}{328.12}{291.461}{K}%
\TextAt{16}{338.938}{291.461}{\Character{106}}%
\TextAt{13}{342.259}{291.461}{K}%
\TextAt{9}{353.076}{291.461}{\Character{41}\Character{41}}%
\TextAt{13}{362.181}{291.461}{\Character{58}}%
\TextAt{8}{79.226}{269.845}{If}%
\TextAt{13}{90.918}{269.845}{n}%
\TextAt{8}{101.803}{269.845}{is}%
\TextAt{8}{113.567}{269.845}{not}%
\TextAt{8}{134.381}{269.845}{divisible}%
\TextAt{8}{180.289}{269.845}{by}%
\TextAt{8}{196.548}{269.845}{the}%
\TextAt{8}{216.704}{269.845}{residue}%
\TextAt{8}{256.431}{269.845}{characteristic}%
\TextAt{13}{328.366}{269.845}{p}%
\TextAt{8}{334.241}{269.845}{,}%
\TextAt{8}{341.39}{269.845}{then}%
\TextAt{8}{368.05}{269.845}{the}%
\TextAt{8}{388.207}{269.845}{orthogonal}%
\TextAt{8}{446.763}{269.845}{complement}%
\TextAt{8}{512.122}{269.845}{of}%
\TextAt{13}{171.55}{245.29}{H}%
\TextAt{17}{182.173}{250.226}{1}%
\TextAt{7}{181.249}{242.335}{nr}%
\TextAt{9}{189.749}{245.29}{\Character{40}}%
\TextAt{13}{194.302}{245.29}{K}%
\TextAt{13}{204.464}{245.29}{\Character{59}}%
\TextAt{18}{209.713}{245.29}{Z}%
\TextAt{9}{217.683}{245.29}{\Character{47}}%
\TextAt{13}{223.536}{245.29}{n}%
\TextAt{9}{230.524}{245.29}{\Character{41}}%
\TextAt{9}{238.379}{245.29}{\Character{58}\Character{61}}%
\TextAt{13}{254.074}{245.29}{H}%
\TextAt{17}{264.697}{250.226}{1}%
\TextAt{9}{269.429}{245.29}{\Character{40}}%
\TextAt{13}{273.981}{245.29}{G}%
\TextAt{9}{283.215}{245.29}{\Character{40}}%
\TextAt{20}{290.506}{248.312}{\Character{101}}%
\TextAt{13}{287.767}{245.29}{K}%
\TextAt{16}{298.585}{245.29}{\Character{106}}%
\TextAt{13}{301.906}{245.29}{K}%
\TextAt{9}{312.724}{245.29}{\Character{41}}%
\TextAt{13}{317.276}{245.29}{\Character{59}}%
\TextAt{18}{322.52}{245.29}{Z}%
\TextAt{9}{330.49}{245.29}{\Character{47}}%
\TextAt{13}{336.343}{245.29}{n}%
\TextAt{18}{343.331}{245.29}{Z}%
\TextAt{9}{351.301}{245.29}{\Character{41}}%
\NativeMathAt{359.174}{245.29}{11.955}{\subseteq}%
\TextAt{13}{371.793}{245.29}{H}%
\TextAt{17}{382.416}{250.226}{1}%
\TextAt{9}{387.148}{245.29}{\Character{40}}%
\TextAt{13}{391.701}{245.29}{K}%
\TextAt{13}{401.863}{245.29}{\Character{59}}%
\TextAt{18}{407.112}{245.29}{Z}%
\TextAt{9}{415.082}{245.29}{\Character{47}}%
\TextAt{13}{420.935}{245.29}{n}%
\TextAt{18}{427.923}{245.29}{Z}%
\TextAt{9}{435.893}{245.29}{\Character{41}}%
\TextAt{8}{79.226}{220.736}{is}%
\TextAt{8}{90.99}{220.736}{the}%
\TextAt{8}{111.146}{220.736}{group}%
\TextAt{13}{190.092}{203.69}{H}%
\TextAt{17}{200.714}{208.626}{1}%
\TextAt{7}{199.791}{200.734}{nr}%
\TextAt{9}{208.291}{203.69}{\Character{40}}%
\TextAt{13}{212.843}{203.69}{K}%
\TextAt{13}{223.005}{203.69}{\Character{59}}%
\TextAt{13}{228.241}{203.69}{\Character{22}}%
\TextAt{21}{235.298}{201.896}{n}%
\TextAt{9}{240.934}{203.69}{\Character{41}}%
\TextAt{9}{248.789}{203.69}{\Character{58}\Character{61}}%
\TextAt{13}{264.484}{203.69}{H}%
\TextAt{17}{275.107}{208.626}{1}%
\TextAt{9}{279.839}{203.69}{\Character{40}}%
\TextAt{13}{284.391}{203.69}{G}%
\TextAt{9}{293.625}{203.69}{\Character{40}}%
\TextAt{20}{300.916}{206.712}{\Character{101}}%
\TextAt{13}{298.177}{203.69}{K}%
\TextAt{16}{308.995}{203.69}{\Character{106}}%
\TextAt{13}{312.316}{203.69}{K}%
\TextAt{9}{323.134}{203.69}{\Character{41}}%
\TextAt{13}{327.686}{203.69}{\Character{59}}%
\TextAt{13}{332.91}{203.69}{\Character{22}}%
\TextAt{21}{339.973}{201.896}{n}%
\TextAt{9}{345.609}{203.69}{\Character{41}}%
\NativeMathAt{353.483}{203.69}{11.955}{\subseteq}%
\TextAt{13}{366.102}{203.69}{H}%
\TextAt{17}{376.725}{208.626}{1}%
\TextAt{9}{381.457}{203.69}{\Character{40}}%
\TextAt{13}{386.009}{203.69}{K}%
\TextAt{13}{396.171}{203.69}{\Character{59}}%
\TextAt{13}{401.407}{203.69}{\Character{22}}%
\TextAt{21}{408.464}{201.896}{n}%
\TextAt{9}{414.1}{203.69}{\Character{41}}%
\TextAt{13}{418.652}{203.69}{\Character{58}}%
\TextAt{12}{79.226}{169.671}{Exercise}%
\TextAt{12}{134.078}{169.671}{V.1.3.}%
\TextAt{8}{174.261}{169.671}{If}%
\TextAt{13}{187.389}{169.671}{L}%
\TextAt{16}{195.354}{169.671}{\Character{106}}%
\TextAt{13}{198.674}{169.671}{K}%
\TextAt{8}{214.825}{169.671}{is}%
\TextAt{8}{228.025}{169.671}{a}%
\TextAt{8}{239.217}{169.671}{finite}%
\TextAt{8}{270.564}{169.671}{extension}%
\TextAt{8}{323.766}{169.671}{of}%
\TextAt{14}{338.532}{169.671}{p}%
\TextAt{8}{344.515}{169.671}{-adic}%
\TextAt{8}{374.56}{169.671}{number}%
\TextAt{8}{418.592}{169.671}{fields,}%
\TextAt{8}{453.61}{169.671}{then}%
\TextAt{8}{481.706}{169.671}{one}%
\TextAt{8}{504.602}{169.671}{has}%
\TextAt{8}{526.911}{169.671}{a}%
\TextAt{8}{79.226}{152.625}{commutative}%
\TextAt{8}{149.152}{152.625}{diagram}%
\TextAt{13}{178.761}{127.814}{H}%
\TextAt{17}{189.384}{132.153}{1}%
\TextAt{9}{194.116}{127.814}{\Character{40}}%
\TextAt{13}{198.669}{127.814}{L\Character{59}}%
\TextAt{18}{211.877}{127.814}{Z}%
\TextAt{9}{219.848}{127.814}{\Character{47}}%
\TextAt{13}{225.701}{127.814}{n}%
\TextAt{18}{232.688}{127.814}{Z}%
\TextAt{9}{240.658}{127.814}{\Character{41}}%
\TextAt{13}{300.758}{127.814}{H}%
\TextAt{17}{311.38}{132.153}{1}%
\TextAt{9}{316.113}{127.814}{\Character{40}}%
\TextAt{13}{320.665}{127.814}{L\Character{59}}%
\TextAt{13}{333.852}{127.814}{\Character{22}}%
\TextAt{21}{340.916}{126.021}{n}%
\TextAt{9}{346.553}{127.814}{\Character{41}}%
\TextAt{18}{405.55}{127.814}{Z}%
\TextAt{9}{413.52}{127.814}{\Character{47}}%
\TextAt{13}{419.373}{127.814}{n}%
\TextAt{18}{426.361}{127.814}{Z}%
\TextAt{13}{177.66}{74.312}{H}%
\TextAt{17}{188.283}{78.651}{1}%
\TextAt{9}{193.015}{74.312}{\Character{40}}%
\TextAt{13}{197.568}{74.312}{K}%
\TextAt{13}{207.73}{74.312}{\Character{59}}%
\TextAt{18}{212.979}{74.312}{Z}%
\TextAt{9}{220.949}{74.312}{\Character{47}}%
\TextAt{13}{226.802}{74.312}{n}%
\TextAt{18}{233.79}{74.312}{Z}%
\TextAt{9}{241.76}{74.312}{\Character{41}}%
\TextAt{13}{299.656}{74.312}{H}%
\TextAt{17}{310.278}{78.651}{1}%
\TextAt{9}{315.011}{74.312}{\Character{40}}%
\TextAt{13}{319.563}{74.312}{K}%
\TextAt{13}{329.725}{74.312}{\Character{59}}%
\TextAt{13}{334.961}{74.312}{\Character{22}}%
\TextAt{21}{342.017}{72.519}{n}%
\TextAt{9}{347.654}{74.312}{\Character{41}}%
\TextAt{18}{405.55}{74.312}{Z}%
\TextAt{9}{413.52}{74.312}{\Character{47}}%
\TextAt{13}{419.373}{74.312}{n}%
\TextAt{18}{426.361}{74.312}{Z}%
\TextAt{16}{268.337}{127.814}{\Character{2}}%
\NativeDiagram{%
\path[draw=black,line width=0.478bp] (284.46,-643.197) -- (327.724,-643.197);%
\path[draw=black,line width=0.478bp,line cap=round,line join=round] (325.476,-640.328) .. controls (325.931,-642.049) and (326.943,-642.862) .. (327.963,-643.197) .. controls (326.943,-643.532) and (325.931,-644.344) .. (325.476,-646.066);%
\path[draw=black,line width=0.478bp] (253.934,-653.754) -- (253.934,-685.156);%
\path[draw=black,line width=0.478bp,line cap=round,line join=round] (256.803,-682.909) .. controls (255.081,-683.363) and (254.269,-684.375) .. (253.934,-685.395) .. controls (253.599,-684.375) and (252.786,-683.363) .. (251.065,-682.909);%
}%
\TextAt{21}{328.749}{103.199}{N}%
\TextAt{22}{335.47}{101.461}{L}%
\TextAt{23}{340.529}{101.461}{\Character{106}}%
\TextAt{22}{342.798}{101.461}{K}%
\NativeDiagram{%
\path[draw=black,line width=2.797bp] (347.945,-653.754) -- (347.945,-686.143);%
\path[draw=black!0,line width=1.841bp] (347.945,-653.754) -- (347.945,-686.143);%
\path[draw=black,line width=0.478bp] (139.987,-685.634) -- (139.987,-654.232);%
\path[draw=black,line width=0.478bp,line cap=round,line join=round] (137.118,-656.479) .. controls (138.84,-656.025) and (139.652,-655.013) .. (139.987,-653.993) .. controls (140.322,-655.013) and (141.135,-656.025) .. (142.856,-656.479);%
}%
\TextAt{7}{198.771}{102.568}{res}%
\TextAt{16}{268.337}{74.312}{\Character{2}}%
\NativeDiagram{%
\path[draw=black,line width=0.478bp] (285.561,-696.7) -- (327.724,-696.7);%
\path[draw=black,line width=0.478bp,line cap=round,line join=round] (325.476,-693.83) .. controls (325.931,-695.552) and (326.943,-696.365) .. (327.963,-696.699) .. controls (326.943,-697.034) and (325.931,-697.847) .. (325.476,-699.569);%
}%
\TextAt{8}{297.211}{43.033}{298}%
\Destination{78.23}{809.815}{nav.page.298}%
\Destination{79.226}{505.855}{nav.ex.5.1.1}%
\Destination{79.226}{427.355}{nav.ex.5.1.2}%
\Destination{79.226}{183.655}{nav.ex.5.1.3}%
\Hotlink{249.791}{607.555}{14.294}{11.047}{nav.chapter.3}%
\Hotlink{218.064}{514.455}{6.858}{11.047}{nav.ex.5.3.1}%
\Hotlink{347.729}{514.455}{6.858}{11.047}{nav.ex.5.3.4}%
\Hotlink{414.722}{514.455}{13.995}{11.047}{nav.chapter.4}%
\Hotlink{440.546}{514.455}{6.858}{11.047}{nav.section.4.3}%
\Hotlink{270.118}{607.555}{30.936}{11.047}{nav.number.3.2.18}%
}

% Page 311 (body page 299)
\Page{%
\TextAt{2}{79.226}{796.943}{V.}%
\TextAt{2}{94.073}{796.943}{Local}%
\TextAt{2}{123.615}{796.943}{Class}%
\TextAt{2}{152.208}{796.943}{Field}%
\TextAt{2}{179.939}{796.943}{Theory}%
\Vector{q 0 G 0 g 0.299 w 7.226 18.729 m 460.769 18.729 l S Q}%
\TextAt{12}{79.226}{761.124}{Exercise}%
\TextAt{12}{132.952}{761.124}{V.1.4.}%
\TextAt{12}{172.156}{761.124}{(Local}%
\TextAt{12}{213.843}{761.124}{Tate}%
\TextAt{12}{244.986}{761.124}{Duality).}%
\TextAt{8}{304.3}{761.124}{Show}%
\TextAt{8}{335.649}{761.124}{that}%
\TextAt{8}{361.474}{761.124}{the}%
\TextAt{8}{382.088}{761.124}{statements}%
\TextAt{8}{441.532}{761.124}{of}%
\TextAt{8}{455.319}{761.124}{exercises}%
\TextAt{8}{503.692}{761.124}{2}%
\TextAt{8}{513.904}{761.124}{and}%
\TextAt{8}{79.226}{744.078}{3}%
\TextAt{8}{89.076}{744.078}{generalize}%
\TextAt{8}{143.147}{744.078}{to}%
\TextAt{8}{157.552}{744.078}{an}%
\TextAt{8}{173.905}{744.078}{arbitrary}%
\TextAt{8}{223.768}{744.078}{finite}%
\TextAt{13}{253.774}{744.078}{G}%
\TextAt{21}{263.008}{742.284}{K}%
\TextAt{8}{271.211}{744.078}{-module}%
\TextAt{13}{316.495}{744.078}{A}%
\TextAt{8}{329.262}{744.078}{instead}%
\TextAt{8}{369.741}{744.078}{of}%
\TextAt{18}{383.165}{744.078}{Z}%
\TextAt{9}{391.135}{744.078}{\Character{47}}%
\TextAt{13}{396.988}{744.078}{n}%
\TextAt{18}{403.976}{744.078}{Z}%
\TextAt{8}{411.946}{744.078}{,}%
\TextAt{8}{419.213}{744.078}{and}%
\TextAt{13}{442.069}{744.078}{A}%
\TextAt{3}{450.845}{748.416}{\Character{48}}%
\TextAt{9}{457.121}{744.078}{\Character{61}}%
\TextAt{8}{469.707}{744.078}{Hom}%
\TextAt{9}{494.095}{744.078}{\Character{40}}%
\TextAt{13}{498.648}{744.078}{A\Character{59}}%
\TextAt{9}{515.8}{747.1}{\Character{22}}%
\TextAt{13}{512.667}{744.078}{K}%
\TextAt{3}{523.485}{748.416}{\Character{3}}%
\TextAt{9}{528.217}{744.078}{\Character{41}}%
\TextAt{8}{79.226}{727.032}{instead}%
\TextAt{8}{119.611}{727.032}{of}%
\TextAt{13}{132.941}{727.032}{\Character{22}}%
\TextAt{21}{139.984}{725.238}{n}%
\TextAt{8}{145.62}{727.032}{.}%
\TextAt{12}{97.159}{706.997}{Hint:}%
\TextAt{8}{132.698}{706.997}{Use}%
\TextAt{8}{155.185}{706.997}{exercises}%
\TextAt{8}{203.102}{706.997}{4–8}%
\TextAt{8}{224.573}{706.997}{of}%
\TextAt{8}{237.903}{706.997}{chap.}%
\TextAt{8}{270.094}{706.997}{IV,}%
\TextAt{8}{290.239}{706.997}{§3.}%
\TextAt{12}{79.226}{678.008}{Exercise}%
\TextAt{12}{132.361}{678.008}{V.1.5.}%
\TextAt{8}{171.051}{678.008}{Let}%
\TextAt{13}{191.964}{678.008}{L}%
\TextAt{16}{199.928}{678.007}{\Character{106}}%
\TextAt{13}{203.249}{678.008}{K}%
\TextAt{8}{217.908}{678.008}{be}%
\TextAt{8}{233.775}{678.008}{the}%
\TextAt{8}{253.875}{678.008}{composite}%
\TextAt{8}{308.836}{678.008}{of}%
\TextAt{8}{322.109}{678.008}{all}%
\TextAt{18}{338.311}{678.008}{Z}%
\TextAt{21}{346.281}{676.214}{p}%
\TextAt{8}{351.042}{678.008}{-extensions}%
\TextAt{8}{411.263}{678.008}{of}%
\TextAt{8}{424.537}{678.008}{a}%
\TextAt{14}{434.235}{678.008}{p}%
\TextAt{8}{440.218}{678.008}{-adic}%
\TextAt{8}{468.77}{678.008}{number}%
\TextAt{8}{511.31}{678.008}{field}%
\TextAt{13}{79.226}{660.962}{K}%
\TextAt{8}{95.196}{660.962}{(i.e.,}%
\TextAt{8}{123.424}{660.962}{extensions}%
\TextAt{8}{181.06}{660.962}{with}%
\TextAt{8}{208.975}{660.962}{Galois}%
\TextAt{8}{246.143}{660.962}{group}%
\TextAt{8}{280.573}{660.962}{isomorphic}%
\TextAt{8}{341.078}{660.962}{to}%
\TextAt{18}{356.643}{660.962}{Z}%
\TextAt{21}{364.613}{659.168}{p}%
\TextAt{8}{369.374}{660.962}{).}%
\TextAt{8}{386.141}{660.962}{Show}%
\TextAt{8}{418.288}{660.962}{that}%
\TextAt{8}{444.912}{660.962}{the}%
\TextAt{8}{466.323}{660.962}{Galois}%
\TextAt{8}{503.491}{660.962}{group}%
\TextAt{13}{79.226}{643.916}{G}%
\TextAt{9}{88.46}{643.916}{\Character{40}}%
\TextAt{13}{93.012}{643.916}{L}%
\TextAt{16}{100.977}{643.916}{\Character{106}}%
\TextAt{13}{104.297}{643.916}{K}%
\TextAt{9}{115.115}{643.916}{\Character{41}}%
\TextAt{8}{123.565}{643.916}{is}%
\TextAt{8}{135.329}{643.916}{a}%
\TextAt{8}{145.084}{643.916}{free,}%
\TextAt{8}{170.764}{643.916}{finitely}%
\TextAt{8}{210.108}{643.916}{generated}%
\TextAt{18}{263.44}{643.916}{Z}%
\TextAt{21}{271.41}{642.122}{p}%
\TextAt{8}{276.171}{643.916}{-module}%
\TextAt{8}{321.362}{643.916}{and}%
\TextAt{8}{344.125}{643.916}{determine}%
\TextAt{8}{398.748}{643.916}{its}%
\TextAt{8}{415.067}{643.916}{rank.}%
\TextAt{12}{97.159}{623.881}{Hint:}%
\TextAt{8}{132.698}{623.881}{Use}%
\TextAt{8}{155.185}{623.881}{chap.}%
\TextAt{8}{187.377}{623.881}{II,}%
\TextAt{8}{203.289}{623.881}{(5.7).}%
\TextAt{12}{79.226}{594.891}{Exercise}%
\TextAt{12}{132.842}{594.891}{V.1.6.}%
\TextAt{8}{171.95}{594.891}{There}%
\TextAt{8}{206.12}{594.891}{is}%
\TextAt{8}{218.245}{594.891}{only}%
\TextAt{8}{244.298}{594.891}{one}%
\TextAt{8}{266.118}{594.891}{unramified}%
\TextAt{18}{325.012}{594.891}{Z}%
\TextAt{21}{332.982}{593.098}{p}%
\TextAt{8}{337.743}{594.891}{-extension}%
\TextAt{8}{393.767}{594.891}{of}%
\TextAt{13}{407.459}{594.891}{K}%
\TextAt{8}{418.276}{594.891}{.}%
\TextAt{8}{427.808}{594.891}{Generate}%
\TextAt{8}{478.321}{594.891}{it}%
\TextAt{8}{490.386}{594.891}{by}%
\TextAt{8}{507.006}{594.891}{roots}%
\TextAt{8}{79.226}{577.845}{of}%
\TextAt{8}{92.556}{577.845}{unity.}%
\TextAt{12}{79.226}{548.856}{Exercise}%
\TextAt{12}{131.663}{548.856}{V.1.7.}%
\ErrataAt{169.72}{548.856}{365.28}{Let $K$ be a $p$-adic number field, and let $L$ be the field generated}%
\TextAt{8}{79.226}{531.81}{by}%
\TextAt{8}{95.499}{531.81}{all}%
\TextAt{8}{111.772}{531.81}{roots}%
\TextAt{8}{141.447}{531.81}{of}%
\TextAt{8}{154.791}{531.81}{unity}%
\TextAt{8}{185.375}{531.81}{of}%
\TextAt{13}{198.719}{531.81}{p}%
\TextAt{8}{204.594}{531.81}{-power}%
\TextAt{8}{242.649}{531.81}{order.}%
\TextAt{8}{277.812}{531.81}{The}%
\TextAt{8}{301.88}{531.81}{fixed}%
\TextAt{8}{330.18}{531.81}{field}%
\TextAt{8}{355.55}{531.81}{of}%
\TextAt{8}{368.895}{531.81}{the}%
\TextAt{8}{389.065}{531.81}{torsion}%
\TextAt{8}{428.172}{531.81}{subgroup}%
\TextAt{8}{478.984}{531.81}{of}%
\TextAt{13}{492.328}{531.81}{G}%
\TextAt{9}{501.562}{531.81}{\Character{40}}%
\TextAt{13}{506.114}{531.81}{L}%
\TextAt{16}{514.079}{531.81}{\Character{106}}%
\TextAt{13}{517.399}{531.81}{K}%
\TextAt{9}{528.217}{531.81}{\Character{41}}%
\TextAt{8}{79.226}{514.764}{is}%
\TextAt{8}{90.99}{514.764}{a}%
\TextAt{18}{100.745}{514.764}{Z}%
\TextAt{21}{108.716}{512.971}{p}%
\TextAt{8}{113.477}{514.764}{-extension.}%
\TextAt{8}{173.691}{514.764}{It}%
\TextAt{8}{186.363}{514.764}{is}%
\TextAt{8}{198.127}{514.764}{called}%
\TextAt{8}{231.291}{514.764}{the}%
\TextAt{12}{251.447}{514.764}{cyclotomic}%
\TextAt{18}{318.36}{514.764}{Z}%
\TextAt{21}{326.33}{512.971}{p}%
\TextAt{12}{331.091}{514.764}{-extension}%
\TextAt{8}{390.903}{514.764}{.}%
\TextAt{12}{79.226}{485.774}{Exercise}%
\TextAt{12}{133.414}{485.774}{V.1.8.}%
\TextAt{8}{173.019}{485.774}{Let}%
\TextAt{20}{196.175}{488.863}{\Character{98}}%
\TextAt{18}{194.847}{485.774}{Q}%
\TextAt{21}{204.146}{483.981}{p}%
\TextAt{16}{208.907}{485.774}{\Character{106}}%
\TextAt{18}{212.227}{485.774}{Q}%
\TextAt{21}{221.526}{483.981}{p}%
\TextAt{8}{231.043}{485.774}{be}%
\TextAt{8}{247.825}{485.774}{the}%
\TextAt{8}{268.84}{485.774}{cyclotomic}%
\TextAt{18}{327.908}{485.774}{Z}%
\TextAt{21}{335.879}{483.981}{p}%
\TextAt{8}{340.64}{485.774}{-extension}%
\TextAt{8}{397.161}{485.774}{of}%
\TextAt{18}{411.35}{485.774}{Q}%
\TextAt{21}{420.648}{483.981}{p}%
\TextAt{8}{425.409}{485.774}{,}%
\TextAt{8}{433.631}{485.774}{let}%
\TextAt{13}{451.395}{485.774}{G}%
\TextAt{9}{460.628}{485.774}{\Character{40}}%
\TextAt{20}{466.509}{488.863}{\Character{98}}%
\TextAt{18}{465.18}{485.774}{Q}%
\TextAt{21}{474.479}{483.981}{p}%
\TextAt{16}{479.24}{485.774}{\Character{106}}%
\TextAt{18}{482.561}{485.774}{Q}%
\TextAt{21}{491.859}{483.981}{p}%
\TextAt{9}{496.62}{485.774}{\Character{41}}%
\NativeMathAt{505.956}{485.774}{11.955}{\cong}%
\TextAt{18}{520.038}{485.774}{Z}%
\TextAt{21}{528.009}{483.981}{p}%
\TextAt{8}{79.226}{468.728}{be}%
\TextAt{8}{95.974}{468.728}{a}%
\TextAt{8}{106.552}{468.728}{chosen}%
\TextAt{8}{144.831}{468.728}{isomorphism,}%
\TextAt{8}{217.531}{468.728}{and}%
\TextAt{8}{241.117}{468.728}{let}%
\TextAt{20}{260.516}{471.883}{\Character{98}}%
\TextAt{13}{258.845}{468.728}{d}%
\TextAt{9}{269.651}{468.728}{\Character{58}}%
\TextAt{13}{277.626}{468.728}{G}%
\TextAt{25}{286.86}{466.935}{Q}%
\TextAt{22}{293.059}{465.939}{p}%
\NativeMathAt{302.619}{468.728}{11.955}{\rightarrow}%
\TextAt{18}{319.298}{468.728}{Z}%
\TextAt{21}{327.268}{466.935}{p}%
\TextAt{8}{336.749}{468.728}{be}%
\TextAt{8}{353.497}{468.728}{the}%
\TextAt{8}{374.476}{468.728}{induced}%
\TextAt{8}{418.864}{468.728}{homomorphism}%
\TextAt{8}{502.357}{468.728}{of}%
\TextAt{8}{516.51}{468.728}{the}%
\TextAt{8}{79.226}{451.683}{absolute}%
\TextAt{8}{125.469}{451.683}{Galois}%
\TextAt{8}{161.382}{451.683}{group.}%
\TextAt{8}{199.109}{451.683}{Show:}%
\ErrataAt{97.13}{431.578}{437.87}{For a suitable topological generator $u$ of $1+p\mathbb Z_p$ for odd $p$, or of $1+4\mathbb Z_2$ for $p=2$,}%
\TextAt{20}{223.957}{404.219}{\Character{98}}%
\TextAt{13}{223.908}{404.219}{v}%
\TextAt{9}{233.317}{404.219}{\Character{58}}%
\TextAt{18}{239.89}{404.219}{Q}%
\TextAt{3}{249.188}{409.156}{\Character{3}}%
\TextAt{21}{249.188}{401.264}{p}%
\NativeArrowAt{257.27}{404.219}{11.955}{19.26}{\longrightarrow}%
\TextAt{18}{279.852}{404.219}{Z}%
\TextAt{21}{287.822}{402.426}{p}%
\TextAt{13}{292.583}{404.219}{\Character{59}}%
\TextAt{20}{321.786}{404.219}{\Character{98}}%
\TextAt{13}{321.738}{404.219}{v}%
\TextAt{9}{327.826}{404.219}{\Character{40}}%
\TextAt{13}{332.378}{404.219}{a}%
\TextAt{9}{338.523}{404.219}{\Character{41}}%
\TextAt{9}{346.378}{404.219}{\Character{61}}%
\TextAt{8}{360.276}{412.307}{log}%
\TextAt{13}{377.236}{412.307}{a}%
\Vector{q 0 G 0 g 0.478 w 288.017 -366.792 m 311.64 -366.792 l S Q}%
\TextAt{8}{360.017}{396.019}{log}%
\TextAt{13}{376.978}{396.019}{u}%
\TextAt{13}{384.836}{404.219}{\Character{59}}%
\TextAt{8}{79.226}{374.078}{defines}%
\TextAt{8}{117.046}{374.078}{a}%
\TextAt{8}{126.198}{374.078}{henselian}%
\TextAt{8}{176.38}{374.078}{valuation}%
\TextAt{8}{226.849}{374.078}{with}%
\TextAt{8}{252.906}{374.078}{respect}%
\TextAt{8}{292.352}{374.078}{to}%
\TextAt{20}{307.73}{377.233}{\Character{98}}%
\TextAt{13}{306.059}{374.078}{d}%
\TextAt{8}{312.141}{374.078}{,}%
\TextAt{8}{318.808}{374.078}{in}%
\TextAt{8}{331.857}{374.078}{the}%
\TextAt{8}{351.41}{374.078}{sense}%
\TextAt{8}{380.837}{374.078}{of}%
\TextAt{8}{393.564}{374.078}{abstract}%
\TextAt{13}{438.557}{374.078}{p}%
\TextAt{8}{444.432}{374.078}{-class}%
\TextAt{8}{475.163}{374.078}{field}%
\TextAt{8}{499.917}{374.078}{theory}%
\TextAt{8}{79.226}{357.032}{(see}%
\TextAt{8}{102.694}{357.032}{chap.}%
\TextAt{8}{134.885}{357.032}{IV,}%
\TextAt{8}{155.03}{357.032}{§5,}%
\TextAt{8}{173.716}{357.032}{exercise}%
\TextAt{8}{217.017}{357.032}{2).}%
%
\TextAt{12}{79.226}{328.043}{Exercise}%
\TextAt{12}{132.496}{328.043}{V.1.9.}%
\TextAt{8}{171.302}{328.043}{Determine}%
\TextAt{8}{228.424}{328.043}{all}%
\TextAt{13}{244.743}{328.043}{p}%
\TextAt{8}{250.618}{328.043}{-class}%
\TextAt{8}{282.013}{328.043}{field}%
\TextAt{8}{307.429}{328.043}{theories}%
\TextAt{9}{351.126}{328.043}{\Character{40}}%
\TextAt{13}{355.678}{328.043}{d}%
\TextAt{9}{365.184}{328.043}{\Character{58}}%
\TextAt{13}{371.858}{328.043}{G}%
\TextAt{21}{381.092}{326.25}{K}%
\NativeMathAt{392.717}{328.043}{11.955}{\rightarrow}%
\TextAt{18}{408.096}{328.043}{Z}%
\TextAt{21}{416.066}{326.25}{p}%
\TextAt{13}{420.827}{328.043}{\Character{59}}%
\TextAt{13}{430.009}{328.043}{v}%
\TextAt{9}{439.539}{328.043}{\Character{58}}%
\TextAt{13}{446.214}{328.043}{K}%
\TextAt{3}{457.031}{332.381}{\Character{3}}%
\NativeMathAt{465.187}{328.043}{11.955}{\rightarrow}%
\TextAt{18}{480.565}{328.043}{Z}%
\TextAt{21}{488.535}{326.25}{p}%
\TextAt{9}{493.296}{328.043}{\Character{41}}%
\TextAt{8}{501.805}{328.043}{over}%
\TextAt{8}{526.911}{328.043}{a}%
\TextAt{14}{79.226}{310.997}{p}%
\TextAt{8}{85.209}{310.997}{-adic}%
\TextAt{8}{113.818}{310.997}{number}%
\TextAt{8}{156.414}{310.997}{field}%
\TextAt{13}{181.771}{310.997}{K}%
\TextAt{8}{192.589}{310.997}{.}%
\TextAt{12}{79.226}{282.007}{Exercise}%
\TextAt{12}{133.103}{282.007}{V.1.10.}%
\TextAt{8}{179.169}{282.007}{Determine}%
\TextAt{8}{236.819}{282.007}{all}%
\TextAt{8}{253.666}{282.007}{class}%
\TextAt{8}{281.691}{282.007}{field}%
\TextAt{8}{307.637}{282.007}{theories}%
\TextAt{9}{351.861}{282.007}{\Character{40}}%
\TextAt{13}{356.413}{282.007}{d}%
\TextAt{9}{366.819}{282.007}{\Character{58}}%
\TextAt{13}{374.394}{282.007}{G}%
\TextAt{21}{383.628}{280.214}{K}%
\NativeMathAt{396.153}{282.007}{11.955}{\rightarrow}%
\TextAt{20}{413.096}{285.096}{\Character{98}}%
\TextAt{18}{412.432}{282.007}{Z}%
\TextAt{13}{420.402}{282.007}{\Character{59}}%
\TextAt{13}{430.11}{282.007}{v}%
\TextAt{9}{440.543}{282.007}{\Character{58}}%
\TextAt{13}{448.118}{282.007}{K}%
\TextAt{3}{458.935}{286.346}{\Character{3}}%
\NativeMathAt{467.991}{282.007}{11.955}{\rightarrow}%
\TextAt{20}{484.933}{285.096}{\Character{98}}%
\TextAt{18}{484.269}{282.007}{Z}%
\TextAt{9}{492.239}{282.007}{\Character{41}}%
\TextAt{8}{501.277}{282.007}{over}%
\TextAt{8}{526.911}{282.007}{a}%
\TextAt{14}{79.226}{264.961}{p}%
\TextAt{8}{85.209}{264.961}{-adic}%
\TextAt{8}{113.818}{264.961}{number}%
\TextAt{8}{156.414}{264.961}{field}%
\TextAt{13}{181.771}{264.961}{K}%
\TextAt{8}{192.589}{264.961}{.}%
\TextAt{12}{79.226}{235.972}{Exercise}%
\TextAt{12}{131.866}{235.972}{V.1.11.}%
\TextAt{8}{176.856}{235.972}{The}%
\TextAt{12}{200.423}{235.972}{Weil}%
\TextAt{12}{230.731}{235.972}{group}%
\TextAt{8}{268.046}{235.972}{of}%
\TextAt{8}{280.889}{235.972}{a}%
\TextAt{8}{290.157}{235.972}{local}%
\TextAt{8}{317.31}{235.972}{field}%
\TextAt{13}{342.18}{235.972}{K}%
\TextAt{8}{356.408}{235.972}{is}%
\TextAt{8}{367.684}{235.972}{the}%
\TextAt{8}{387.354}{235.972}{preimage}%
\TextAt{13}{436.947}{235.972}{W}%
\TextAt{21}{447.998}{234.179}{K}%
\TextAt{8}{459.611}{235.972}{of}%
\TextAt{18}{472.454}{235.972}{Z}%
\TextAt{8}{483.834}{235.972}{under}%
\TextAt{8}{516.51}{235.972}{the}%
\TextAt{8}{79.226}{218.926}{mapping}%
\TextAt{13}{127.358}{218.926}{d}%
\TextAt{21}{133.44}{217.133}{K}%
\TextAt{9}{144.964}{218.926}{\Character{58}}%
\TextAt{13}{151.536}{218.926}{G}%
\TextAt{21}{160.77}{217.133}{K}%
\NativeMathAt{172.293}{218.926}{11.955}{\rightarrow}%
\TextAt{20}{188.234}{222.014}{\Character{98}}%
\TextAt{18}{187.569}{218.926}{Z}%
\TextAt{8}{195.54}{218.926}{.}%
\TextAt{8}{203.988}{218.926}{Show:}%
\TextAt{8}{97.159}{198.891}{The}%
\TextAt{8}{121.518}{198.891}{norm}%
\TextAt{8}{152.393}{198.891}{residue}%
\TextAt{8}{192.426}{198.891}{symbol}%
\TextAt{9}{232.794}{198.891}{\Character{40}}%
\TextAt{13}{241.549}{198.891}{\Character{59}}%
\TextAt{13}{246.773}{198.891}{K}%
\TextAt{7}{257.611}{203.23}{ab}%
\TextAt{16}{267.043}{198.891}{\Character{106}}%
\TextAt{13}{270.364}{198.891}{K}%
\TextAt{9}{281.182}{198.891}{\Character{41}}%
\TextAt{8}{289.937}{198.891}{of}%
\TextAt{8}{303.573}{198.891}{the}%
\TextAt{8}{324.035}{198.891}{maximal}%
\TextAt{8}{372.15}{198.891}{abelian}%
\TextAt{8}{413.103}{198.891}{extension}%
\TextAt{13}{465.174}{198.891}{K}%
\TextAt{7}{475.992}{203.23}{ab}%
\TextAt{16}{485.425}{198.891}{\Character{106}}%
\TextAt{13}{488.746}{198.891}{K}%
\TextAt{8}{503.766}{198.891}{yields}%
\TextAt{8}{79.226}{181.845}{an}%
\TextAt{8}{95.485}{181.845}{isomorphism}%
\TextAt{9}{244.37}{164.799}{\Character{40}}%
\TextAt{13}{252.82}{164.799}{\Character{59}}%
\TextAt{13}{258.044}{164.799}{K}%
\TextAt{7}{268.882}{169.735}{ab}%
\TextAt{16}{278.315}{164.799}{\Character{106}}%
\TextAt{13}{281.636}{164.799}{K}%
\TextAt{9}{292.453}{164.799}{\Character{41}}%
\TextAt{9}{300.308}{164.799}{\Character{58}}%
\TextAt{13}{306.899}{164.799}{K}%
\TextAt{3}{317.717}{169.735}{\Character{3}}%
\NativeMathAt{328.357}{171.576}{7.97}{\sim}%
\NativeArrowAt{325.77}{164.799}{11.955}{13.175}{\longrightarrow}%
\TextAt{13}{342.264}{164.799}{W}%
\TextAt{7}{354.941}{169.735}{ab}%
\TextAt{21}{353.315}{161.844}{K}%
\TextAt{13}{364.374}{164.799}{\Character{59}}%
\TextAt{8}{79.226}{143.528}{which}%
\TextAt{8}{111.6}{143.528}{maps}%
\TextAt{8}{141.116}{143.528}{the}%
\TextAt{8}{160.16}{143.528}{unit}%
\TextAt{8}{183.758}{143.528}{group}%
\TextAt{13}{215.821}{143.528}{U}%
\TextAt{21}{223.771}{141.735}{K}%
\TextAt{8}{234.758}{143.528}{onto}%
\TextAt{8}{259.995}{143.528}{the}%
\TextAt{8}{279.038}{143.528}{inertia}%
\TextAt{8}{314.998}{143.528}{group}%
\TextAt{13}{347.061}{143.528}{I}%
\TextAt{9}{353.164}{143.528}{\Character{40}}%
\TextAt{13}{357.716}{143.528}{K}%
\TextAt{7}{368.534}{147.866}{ab}%
\TextAt{16}{377.967}{143.528}{\Character{106}}%
\TextAt{13}{381.288}{143.528}{K}%
\TextAt{9}{392.105}{143.528}{\Character{41}}%
\TextAt{8}{396.658}{143.528}{,}%
\TextAt{8}{402.916}{143.528}{and}%
\TextAt{8}{424.566}{143.528}{the}%
\TextAt{8}{443.61}{143.528}{group}%
\TextAt{8}{475.673}{143.528}{of}%
\TextAt{8}{487.89}{143.528}{principal}%
\TextAt{8}{79.226}{126.482}{units}%
\TextAt{13}{108.552}{126.482}{U}%
\TextAt{17}{117.751}{132.592}{\Character{40}1\Character{41}}%
\TextAt{21}{116.502}{123.241}{K}%
\TextAt{8}{132.967}{126.482}{onto}%
\TextAt{8}{159.317}{126.482}{the}%
\TextAt{8}{179.473}{126.482}{ramification}%
\TextAt{8}{244.521}{126.482}{group}%
\TextAt{13}{277.697}{126.482}{R}%
\TextAt{9}{286.705}{126.482}{\Character{40}}%
\TextAt{13}{291.258}{126.482}{K}%
\TextAt{7}{302.075}{130.82}{ab}%
\TextAt{16}{311.508}{126.482}{\Character{106}}%
\TextAt{13}{314.829}{126.482}{K}%
\TextAt{9}{325.646}{126.482}{\Character{41}}%
\TextAt{8}{330.199}{126.482}{.}%
%
%
%
%
%
%
%
%
%
%
%
%
%
%
%
%
%
%
%
%
%
%
%
%
%
%
%
%
%
%
%
%
%
%
%
%
%
%
%
%
%
%
%
%
%
%
%
%
%
%
%
%
%
%
%
%
%
%
%
%
%
%
\TextAt{8}{297.211}{43.033}{299}%
%
\Destination{78.23}{809.815}{nav.page.299}%
\Destination{79.226}{775.155}{nav.ex.5.1.4}%
\Destination{79.226}{692.055}{nav.ex.5.1.5}%
\Destination{79.226}{608.855}{nav.ex.5.1.6}%
\Destination{79.226}{562.855}{nav.ex.5.1.7}%
\Destination{79.226}{499.755}{nav.ex.5.1.8}%
\Destination{79.226}{342.055}{nav.ex.5.1.9}%
\Destination{79.226}{296.055}{nav.ex.5.1.10}%
\Destination{79.226}{249.955}{nav.ex.5.1.11}%
\Destination{79.36}{776.155}{nav.print.326}%
%
\Hotlink{503.192}{758.955}{6.858}{11.047}{nav.ex.5.1.2}%
\Hotlink{202.602}{704.755}{6.858}{11.047}{nav.ex.4.3.4}%
\Hotlink{269.594}{704.755}{13.995}{11.047}{nav.chapter.4}%
\Hotlink{295.418}{704.755}{6.858}{11.047}{nav.section.4.3}%
\Hotlink{186.877}{621.655}{9.751}{11.047}{nav.chapter.2}%
\Hotlink{216.517}{354.855}{6.858}{11.047}{nav.ex.4.5.2}%
\Hotlink{134.385}{354.855}{13.995}{11.047}{nav.chapter.4}%
\Hotlink{160.209}{354.855}{6.858}{11.047}{nav.section.4.5}%
\Hotlink{202.789}{621.655}{25.078}{11.047}{nav.number.2.5.7}%
\Destination{79.226}{565.856}{errata.131}%
\Destination{79.226}{446}{errata.202}%
}

% Page 312 (body page 300)
\Page{%
\TextAt{2}{79.226}{796.943}{V.}%
\TextAt{2}{94.073}{796.943}{Local}%
\TextAt{2}{123.615}{796.943}{Class}%
\TextAt{2}{152.208}{796.943}{Field}%
\TextAt{2}{179.939}{796.943}{Theory}%
\Vector{q 0 G 0 g 0.299 w 7.226 18.729 m 460.769 18.729 l S Q}%
\TextAt{15}{79.226}{760.804}{§}%
\TextAt{15}{93.785}{760.804}{2}%
\TextAt{15}{120.693}{760.804}{The}%
\TextAt{15}{160.202}{760.804}{Norm}%
\TextAt{15}{215.567}{760.804}{Residue}%
\TextAt{15}{288.681}{760.804}{Symbol}%
\TextAt{15}{358.111}{760.804}{over}%
\TextAt{46}{400.168}{760.804}{Q}%
\TextAt{13}{413.558}{758.222}{p}%
\TextAt{8}{79.226}{728.918}{If}%
\TextAt{13}{91.16}{728.918}{\Character{16}}%
\TextAt{8}{101.314}{728.918}{is}%
\TextAt{8}{113.32}{728.918}{a}%
\TextAt{8}{123.317}{728.918}{primitive}%
\TextAt{13}{173.638}{728.918}{m}%
\TextAt{8}{183.878}{728.918}{-th}%
\TextAt{8}{202.972}{728.918}{root}%
\TextAt{8}{228.26}{728.918}{of}%
\TextAt{8}{241.831}{728.918}{unity,}%
\TextAt{8}{274.974}{728.918}{with}%
\TextAt{9}{301.876}{728.918}{\Character{40}}%
\TextAt{13}{306.428}{728.918}{m\Character{59}}%
\TextAt{13}{321.886}{728.918}{p}%
\TextAt{9}{327.787}{728.918}{\Character{41}}%
\TextAt{9}{336.06}{728.918}{\Character{61}}%
\TextAt{9}{348.888}{728.918}{1}%
\TextAt{8}{354.761}{728.918}{,}%
\TextAt{8}{362.212}{728.918}{then}%
\TextAt{18}{389.114}{728.918}{Q}%
\TextAt{21}{398.412}{727.125}{p}%
\TextAt{9}{403.173}{728.918}{\Character{40}}%
\TextAt{13}{407.726}{728.918}{\Character{16}}%
\TextAt{9}{413.741}{728.918}{\Character{41}}%
\TextAt{16}{418.293}{728.918}{\Character{106}}%
\TextAt{18}{421.614}{728.918}{Q}%
\TextAt{21}{430.913}{727.125}{p}%
\TextAt{8}{439.813}{728.918}{is}%
\TextAt{8}{451.818}{728.918}{unramified,}%
\TextAt{8}{513.904}{728.918}{and}%
\TextAt{8}{79.226}{711.872}{the}%
\TextAt{8}{99.383}{711.872}{norm}%
\TextAt{8}{129.952}{711.872}{residue}%
\TextAt{8}{169.679}{711.872}{symbol}%
\TextAt{8}{209.741}{711.872}{is}%
\TextAt{8}{221.505}{711.872}{obviously}%
\TextAt{8}{273.282}{711.872}{given}%
\TextAt{8}{303.852}{711.872}{by}%
\TextAt{9}{245.394}{683.75}{\Character{40}}%
\TextAt{13}{249.947}{683.75}{a\Character{59}}%
\TextAt{18}{261.336}{683.75}{Q}%
\TextAt{21}{270.634}{681.956}{p}%
\TextAt{9}{275.395}{683.75}{\Character{40}}%
\TextAt{13}{279.948}{683.75}{\Character{16}}%
\TextAt{9}{285.963}{683.75}{\Character{41}}%
\TextAt{16}{290.516}{683.75}{\Character{106}}%
\TextAt{18}{293.836}{683.75}{Q}%
\TextAt{21}{303.135}{681.956}{p}%
\TextAt{9}{307.896}{683.75}{\Character{41}}%
\TextAt{13}{312.448}{683.75}{\Character{16}}%
\TextAt{9}{321.785}{683.75}{\Character{61}}%
\TextAt{13}{334.21}{683.75}{\Character{16}}%
\TextAt{21}{340.226}{688.686}{p}%
\TextAt{22}{344.488}{691.702}{v}%
\TextAt{22}{348.224}{690.706}{p}%
\TextAt{24}{352.563}{691.702}{\Character{40}}%
\TextAt{22}{355.441}{691.702}{a}%
\TextAt{24}{359.475}{691.702}{\Character{41}}%
\TextAt{13}{363.35}{683.75}{\Character{58}}%
\TextAt{8}{79.226}{655.627}{But}%
\TextAt{8}{103.143}{655.627}{if}%
\TextAt{13}{114.544}{655.627}{\Character{16}}%
\TextAt{8}{125.133}{655.627}{is}%
\TextAt{8}{137.573}{655.627}{a}%
\TextAt{8}{148.005}{655.627}{primitive}%
\TextAt{13}{198.762}{655.627}{p}%
\TextAt{21}{204.637}{659.965}{n}%
\TextAt{8}{210.273}{655.627}{-th}%
\TextAt{8}{229.803}{655.627}{root}%
\TextAt{8}{255.525}{655.627}{of}%
\TextAt{8}{269.532}{655.627}{unity,}%
\TextAt{8}{303.218}{655.627}{then}%
\TextAt{8}{330.554}{655.627}{we}%
\TextAt{8}{348.458}{655.627}{obtain}%
\TextAt{8}{385.562}{655.627}{the}%
\TextAt{8}{406.395}{655.627}{norm}%
\TextAt{8}{437.64}{655.627}{residue}%
\TextAt{8}{478.044}{655.627}{symbol}%
\TextAt{8}{518.782}{655.627}{for}%
\TextAt{8}{79.226}{638.581}{the}%
\TextAt{8}{99.383}{638.581}{extension}%
\TextAt{18}{151.148}{638.581}{Q}%
\TextAt{21}{160.447}{636.788}{p}%
\TextAt{9}{165.208}{638.581}{\Character{40}}%
\TextAt{13}{169.76}{638.581}{\Character{16}}%
\TextAt{9}{175.776}{638.581}{\Character{41}}%
\TextAt{16}{180.328}{638.581}{\Character{106}}%
\TextAt{18}{183.649}{638.581}{Q}%
\TextAt{21}{192.948}{636.788}{p}%
\TextAt{8}{201.606}{638.581}{explicitly}%
\TextAt{8}{252.332}{638.581}{in}%
\TextAt{8}{265.984}{638.581}{the}%
\TextAt{8}{286.141}{638.581}{simple}%
\TextAt{8}{322.616}{638.581}{form}%
\TextAt{9}{249.302}{610.458}{\Character{40}}%
\TextAt{13}{253.855}{610.458}{a\Character{59}}%
\TextAt{18}{265.244}{610.458}{Q}%
\TextAt{21}{274.542}{608.665}{p}%
\TextAt{9}{279.303}{610.458}{\Character{40}}%
\TextAt{13}{283.855}{610.458}{\Character{16}}%
\TextAt{9}{289.871}{610.458}{\Character{41}}%
\TextAt{16}{294.423}{610.458}{\Character{106}}%
\TextAt{18}{297.744}{610.458}{Q}%
\TextAt{21}{307.043}{608.665}{p}%
\TextAt{9}{311.804}{610.458}{\Character{41}}%
\TextAt{13}{316.356}{610.458}{\Character{16}}%
\TextAt{9}{325.692}{610.458}{\Character{61}}%
\TextAt{13}{338.118}{610.458}{\Character{16}}%
\TextAt{21}{344.133}{615.394}{u}%
\TextAt{23}{349.036}{618.207}{\Character{0}}%
\TextAt{24}{354.793}{618.207}{1}%
\TextAt{13}{359.442}{610.458}{\Character{59}}%
\TextAt{8}{79.226}{582.335}{where}%
\TextAt{13}{113.848}{582.335}{a}%
\TextAt{9}{124.698}{582.335}{\Character{61}}%
\TextAt{13}{138.507}{582.335}{up}%
\TextAt{21}{151.045}{586.674}{v}%
\TextAt{22}{155.168}{585.677}{p}%
\TextAt{17}{159.507}{586.674}{\Character{40}}%
\TextAt{21}{162.8}{586.674}{a}%
\TextAt{17}{167.298}{586.674}{\Character{41}}%
\TextAt{8}{171.09}{582.335}{,}%
\TextAt{8}{179.254}{582.335}{and}%
\TextAt{13}{202.829}{582.335}{\Character{16}}%
\TextAt{21}{208.845}{586.674}{u}%
\TextAt{23}{213.748}{589.486}{\Character{0}}%
\TextAt{24}{219.504}{589.486}{1}%
\TextAt{8}{228.863}{582.335}{is}%
\TextAt{8}{241.439}{582.335}{the}%
\TextAt{8}{262.407}{582.335}{power}%
\TextAt{13}{297.364}{582.335}{\Character{16}}%
\TextAt{21}{303.379}{586.674}{r}%
\TextAt{8}{312.644}{582.335}{with}%
\TextAt{8}{340.116}{582.335}{any}%
\TextAt{8}{363.045}{582.335}{rational}%
\TextAt{8}{407.446}{582.335}{integer}%
\TextAt{13}{446.957}{582.335}{r}%
\NativeMathAt{457.263}{582.335}{11.955}{\equiv}%
\TextAt{13}{471.266}{582.335}{u}%
\TextAt{3}{477.929}{586.674}{\Character{0}}%
\TextAt{17}{484.515}{586.674}{1}%
\TextAt{8}{492.568}{582.335}{mod}%
\TextAt{13}{518.006}{582.335}{p}%
\TextAt{21}{523.881}{586.674}{n}%
\TextAt{8}{529.518}{582.335}{.}%
\TextAt{8}{79.226}{565.289}{This}%
\TextAt{8}{106.037}{565.289}{result}%
\TextAt{8}{138.705}{565.289}{is}%
\TextAt{8}{150.559}{565.289}{important,}%
\TextAt{8}{209.217}{565.289}{not}%
\TextAt{8}{230.122}{565.289}{only}%
\TextAt{8}{255.904}{565.289}{in}%
\TextAt{8}{269.647}{565.289}{the}%
\TextAt{8}{289.894}{565.289}{local}%
\TextAt{8}{317.625}{565.289}{situation,}%
\TextAt{8}{369.839}{565.289}{but}%
\TextAt{8}{391.389}{565.289}{it}%
\TextAt{8}{403.184}{565.289}{will}%
\TextAt{8}{425.379}{565.289}{play}%
\TextAt{8}{450.839}{565.289}{an}%
\TextAt{8}{467.188}{565.289}{essential}%
\TextAt{8}{513.904}{565.289}{rôle}%
\TextAt{8}{79.226}{548.243}{when}%
\TextAt{8}{110.169}{548.243}{we}%
\TextAt{8}{127.782}{548.243}{develop}%
\TextAt{8}{170.441}{548.243}{global}%
\TextAt{8}{205.306}{548.243}{class}%
\TextAt{8}{233.128}{548.243}{field}%
\TextAt{8}{258.871}{548.243}{theory}%
\TextAt{8}{296.007}{548.243}{(see}%
\TextAt{8}{319.86}{548.243}{chap.}%
\TextAt{8}{353.209}{548.243}{VI,}%
\TextAt{8}{373.739}{548.243}{§5).}%
\TextAt{8}{399.435}{548.243}{Unfortunately,}%
\TextAt{8}{477.961}{548.243}{there}%
\TextAt{8}{508.258}{548.243}{is}%
\TextAt{8}{520.408}{548.243}{no}%
\TextAt{8}{79.226}{531.197}{direct}%
\TextAt{8}{113.473}{531.197}{algebraic}%
\TextAt{8}{163.992}{531.197}{proof}%
\TextAt{8}{195.645}{531.197}{of}%
\TextAt{8}{210.059}{531.197}{this}%
\TextAt{8}{233.965}{531.197}{fact}%
\TextAt{8}{258.134}{531.197}{known}%
\TextAt{8}{296.29}{531.197}{to}%
\TextAt{8}{311.684}{531.197}{date.}%
\TextAt{8}{345.501}{531.197}{We}%
\TextAt{8}{366.729}{531.197}{have}%
\TextAt{8}{394.808}{531.197}{to}%
\TextAt{8}{410.202}{531.197}{invoke}%
\TextAt{8}{447.39}{531.197}{a}%
\TextAt{8}{458.229}{531.197}{transcendental}%
\TextAt{8}{79.226}{514.151}{method}%
\TextAt{8}{121.769}{514.151}{which}%
\TextAt{8}{155.202}{514.151}{makes}%
\TextAt{8}{190.332}{514.151}{use}%
\TextAt{8}{210.495}{514.151}{of}%
\TextAt{8}{223.771}{514.151}{the}%
\TextAt{19}{243.874}{514.151}{completion}%
\TextAt{20}{305.426}{517.173}{\Character{98}}%
\TextAt{13}{302.688}{514.151}{K}%
\TextAt{8}{317.35}{514.151}{of}%
\TextAt{8}{330.626}{514.151}{the}%
\TextAt{8}{350.729}{514.151}{maximal}%
\TextAt{8}{398.484}{514.151}{unramified}%
\TextAt{8}{456.963}{514.151}{extension}%
\TextAt{20}{511.414}{517.173}{\Character{101}}%
\TextAt{13}{508.675}{514.151}{K}%
\TextAt{8}{523.337}{514.151}{of}%
\TextAt{8}{79.226}{497.105}{a}%
\TextAt{8}{88.418}{497.105}{local}%
\TextAt{8}{115.495}{497.105}{field}%
\TextAt{13}{140.289}{497.105}{K}%
\TextAt{8}{151.106}{497.105}{.}%
\TextAt{8}{159.367}{497.105}{We}%
\TextAt{8}{178.948}{497.105}{extend}%
\TextAt{8}{216.426}{497.105}{the}%
\TextAt{8}{236.019}{497.105}{Frobenius}%
\TextAt{13}{289.325}{497.105}{\Character{39}}%
\NativeMathAt{300.314}{497.105}{11.955}{\in}%
\TextAt{13}{311.605}{497.105}{G}%
\TextAt{9}{320.839}{497.105}{\Character{40}}%
\TextAt{20}{328.13}{500.127}{\Character{101}}%
\TextAt{13}{325.391}{497.105}{K}%
\TextAt{16}{336.209}{497.105}{\Character{106}}%
\TextAt{13}{339.53}{497.105}{K}%
\TextAt{9}{350.348}{497.105}{\Character{41}}%
\TextAt{8}{358.234}{497.105}{to}%
\TextAt{20}{374.719}{500.127}{\Character{98}}%
\TextAt{13}{371.981}{497.105}{K}%
\TextAt{8}{386.132}{497.105}{by}%
\TextAt{8}{401.828}{497.105}{continuity.}%
\TextAt{8}{460.504}{497.105}{First}%
\TextAt{8}{488.453}{497.105}{we}%
\TextAt{8}{505.117}{497.105}{prove}%
\TextAt{8}{79.226}{480.059}{the}%
\TextAt{12}{79.226}{454.298}{Lemma}%
\TextAt{12}{126.891}{454.298}{(2.1).}%
\TextAt{8}{162.207}{454.298}{For}%
\TextAt{8}{183.176}{454.298}{every}%
\TextAt{13}{214.069}{454.298}{c}%
\NativeMathAt{222.427}{454.298}{11.955}{\in}%
\TextAt{14}{233.718}{454.298}{o}%
\TextAt{30}{241.533}{452.66}{\Character{98}}%
\TextAt{21}{239.563}{450.645}{K}%
\TextAt{8}{247.765}{454.298}{,}%
\TextAt{8}{254.914}{454.298}{resp.}%
\TextAt{8}{284.236}{454.298}{every}%
\TextAt{13}{315.129}{454.298}{c}%
\NativeMathAt{323.487}{454.298}{11.955}{\in}%
\TextAt{13}{334.778}{454.298}{U}%
\TextAt{30}{344.699}{452.66}{\Character{98}}%
\TextAt{21}{342.729}{450.645}{K}%
\TextAt{8}{350.931}{454.298}{,}%
\TextAt{8}{358.081}{454.298}{the}%
\TextAt{8}{378.237}{454.298}{equation}%
\TextAt{13}{216.282}{426.175}{x}%
\TextAt{21}{222.934}{431.111}{\Character{39}}%
\TextAt{16}{231.607}{426.175}{\Character{0}}%
\TextAt{13}{243.562}{426.175}{x}%
\TextAt{9}{253.535}{426.175}{\Character{61}}%
\TextAt{13}{265.96}{426.175}{c\Character{59}}%
\TextAt{8}{300.153}{426.175}{resp.}%
\TextAt{13}{348.189}{426.175}{x}%
\TextAt{21}{354.841}{431.111}{\Character{39}}%
\TextAt{3}{360.359}{431.111}{\Character{0}}%
\TextAt{17}{366.946}{431.111}{1}%
\TextAt{9}{374.999}{426.175}{\Character{61}}%
\TextAt{13}{387.424}{426.175}{c\Character{59}}%
\TextAt{8}{79.226}{398.052}{admits}%
\TextAt{8}{117.662}{398.052}{a}%
\TextAt{8}{127.417}{398.052}{solution}%
\TextAt{8}{171.711}{398.052}{in}%
\TextAt{14}{185.364}{398.052}{o}%
\TextAt{30}{193.179}{396.414}{\Character{98}}%
\TextAt{21}{191.208}{394.399}{K}%
\TextAt{8}{199.411}{398.052}{,}%
\TextAt{8}{206.56}{398.052}{resp.}%
\TextAt{8}{235.882}{398.052}{in}%
\TextAt{13}{249.535}{398.052}{U}%
\TextAt{30}{259.456}{396.414}{\Character{98}}%
\TextAt{21}{257.485}{394.399}{K}%
\TextAt{8}{265.688}{398.052}{.}%
\TextAt{8}{274.136}{398.052}{If}%
\TextAt{13}{285.829}{398.052}{x}%
\TextAt{21}{292.481}{402.391}{\Character{39}}%
\TextAt{9}{301.818}{398.052}{\Character{61}}%
\TextAt{13}{314.243}{398.052}{x}%
\TextAt{8}{324.793}{398.052}{for}%
\TextAt{13}{342.678}{398.052}{x}%
\NativeMathAt{352.651}{398.052}{11.955}{\in}%
\TextAt{14}{363.942}{398.052}{o}%
\TextAt{30}{371.756}{396.414}{\Character{98}}%
\TextAt{21}{369.786}{394.399}{K}%
\TextAt{8}{377.989}{398.052}{,}%
\TextAt{8}{385.138}{398.052}{then}%
\TextAt{13}{411.798}{398.052}{x}%
\NativeMathAt{421.771}{398.052}{11.955}{\in}%
\TextAt{14}{433.062}{398.052}{o}%
\TextAt{21}{438.906}{396.259}{K}%
\TextAt{8}{447.109}{398.052}{.}%
\TextAt{19}{79.226}{363.576}{Proof.}%
\TextAt{8}{116}{363.576}{Let}%
\TextAt{13}{136.575}{363.576}{\Character{25}}%
\TextAt{8}{147.146}{363.576}{be}%
\TextAt{8}{162.676}{363.576}{a}%
\TextAt{8}{172.036}{363.576}{prime}%
\TextAt{8}{204.805}{363.576}{element}%
\TextAt{8}{247.652}{363.576}{of}%
\TextAt{13}{260.587}{363.576}{K}%
\TextAt{8}{271.405}{363.576}{.}%
\TextAt{8}{279.721}{363.576}{Then}%
\TextAt{13}{309.884}{363.576}{\Character{25}}%
\TextAt{8}{320.456}{363.576}{is}%
\TextAt{8}{331.825}{363.576}{also}%
\TextAt{8}{354.91}{363.576}{a}%
\TextAt{8}{364.27}{363.576}{prime}%
\TextAt{8}{397.039}{363.576}{element}%
\TextAt{8}{439.886}{363.576}{of}%
\TextAt{20}{455.559}{366.598}{\Character{98}}%
\TextAt{13}{452.821}{363.576}{K}%
\TextAt{8}{463.639}{363.576}{,}%
\TextAt{8}{470.472}{363.576}{and}%
\TextAt{8}{492.84}{363.576}{we}%
\TextAt{8}{509.672}{363.576}{have}%
\TextAt{8}{79.226}{346.53}{the}%
\TextAt{13}{99.383}{346.53}{\Character{39}}%
\TextAt{8}{107.051}{346.53}{-invariant}%
\TextAt{8}{160.072}{346.53}{isomorphisms}%
\TextAt{13}{215.355}{318.407}{U}%
\TextAt{30}{225.276}{316.768}{\Character{98}}%
\TextAt{21}{223.305}{314.754}{K}%
\TextAt{9}{231.508}{318.407}{\Character{47}}%
\TextAt{13}{237.361}{318.407}{U}%
\TextAt{17}{246.56}{324.517}{\Character{40}1\Character{41}}%
\TextAt{30}{247.282}{314.856}{\Character{98}}%
\TextAt{21}{245.311}{312.841}{K}%
\NativeMathAt{261.2}{318.407}{11.955}{\cong}%
\TextAt{9}{274.262}{318.407}{\Character{22}}%
\TextAt{13}{273.819}{318.407}{\Character{20}}%
\TextAt{3}{280.558}{323.343}{\Character{3}}%
\TextAt{13}{285.29}{318.407}{\Character{59}}%
\TextAt{13}{314.425}{318.407}{U}%
\TextAt{17}{323.643}{324.517}{\Character{40}}%
\TextAt{21}{326.937}{324.517}{n}%
\TextAt{17}{332.075}{324.517}{\Character{41}}%
\TextAt{30}{324.365}{314.856}{\Character{98}}%
\TextAt{21}{322.395}{312.841}{K}%
\TextAt{9}{335.866}{318.407}{\Character{47}}%
\TextAt{13}{341.719}{318.407}{U}%
\TextAt{17}{350.918}{324.517}{\Character{40}}%
\TextAt{21}{354.212}{324.517}{n}%
\TextAt{17}{359.35}{324.517}{\Character{43}1\Character{41}}%
\TextAt{30}{351.64}{314.856}{\Character{98}}%
\TextAt{21}{349.67}{312.841}{K}%
\NativeMathAt{377.283}{318.407}{11.955}{\cong}%
\TextAt{9}{390.345}{318.407}{\Character{22}}%
\TextAt{13}{389.902}{318.407}{\Character{20}}%
\TextAt{8}{79.226}{290.284}{(see}%
\TextAt{8}{103.257}{290.284}{chap.}%
\TextAt{8}{137.135}{290.284}{II,}%
\TextAt{8}{153.61}{290.284}{(3.10)).}%
\TextAt{8}{198.236}{290.284}{Let}%
\TextAt{13}{219.768}{290.284}{c}%
\NativeMathAt{229.085}{290.284}{11.955}{\in}%
\TextAt{13}{241.335}{290.284}{U}%
\TextAt{30}{251.255}{288.646}{\Character{98}}%
\TextAt{21}{249.285}{286.631}{K}%
\TextAt{8}{261.947}{290.284}{and}%
\TextAt{9}{285.515}{290.284}{\Character{22}}%
\TextAt{13}{285.272}{290.284}{c}%
\TextAt{9}{294.59}{290.284}{\Character{61}}%
\TextAt{13}{307.974}{290.284}{c}%
\TextAt{8}{316.332}{290.284}{mod}%
\TextAt{14}{341.77}{290.284}{p}%
\TextAt{30}{349.724}{288.646}{\Character{98}}%
\TextAt{21}{347.753}{286.631}{K}%
\TextAt{8}{355.956}{290.284}{.}%
\TextAt{8}{366.092}{290.284}{Since}%
\TextAt{8}{397.211}{290.284}{the}%
\TextAt{8}{417.93}{290.284}{residue}%
\TextAt{8}{458.22}{290.284}{class}%
\TextAt{8}{486.219}{290.284}{field}%
\TextAt{9}{512.581}{290.284}{\Character{22}}%
\TextAt{13}{512.139}{290.284}{\Character{20}}%
\TextAt{8}{523.337}{290.284}{of}%
\TextAt{20}{81.964}{276.26}{\Character{98}}%
\TextAt{13}{79.226}{273.238}{K}%
\TextAt{8}{94.69}{273.238}{is}%
\TextAt{8}{107.203}{273.238}{algebraically}%
\TextAt{8}{175.928}{273.238}{closed,}%
\TextAt{8}{214.643}{273.238}{the}%
\TextAt{8}{235.548}{273.238}{equation}%
\TextAt{9}{284.83}{273.238}{\Character{22}}%
\TextAt{13}{284.105}{273.238}{x}%
\TextAt{21}{290.757}{277.576}{\Character{39}}%
\TextAt{9}{301.37}{273.238}{\Character{61}}%
\TextAt{9}{315.788}{273.238}{\Character{22}}%
\TextAt{13}{315.072}{273.238}{x}%
\TextAt{21}{321.724}{277.576}{q}%
\TextAt{9}{330.886}{273.238}{\Character{61}}%
\TextAt{9}{345.304}{273.238}{\Character{22}}%
\TextAt{13}{344.587}{273.238}{x}%
\TextAt{16}{354.406}{273.238}{\Character{1}}%
\TextAt{9}{361.137}{273.238}{\Character{22}}%
\TextAt{13}{360.894}{273.238}{c}%
\TextAt{8}{370.578}{273.238}{(}%
\TextAt{13}{375.133}{273.238}{q}%
\TextAt{9}{385.349}{273.238}{\Character{61}}%
\TextAt{13}{399.05}{273.238}{q}%
\TextAt{21}{404.241}{271.445}{K}%
\TextAt{8}{412.443}{273.238}{)}%
\TextAt{8}{421.644}{273.238}{has}%
\TextAt{8}{443.267}{273.238}{a}%
\TextAt{8}{453.771}{273.238}{solution}%
\NativeMathAt{498.814}{273.238}{11.955}{\neq}%
\TextAt{9}{512.503}{273.238}{0}%
\TextAt{8}{523.014}{273.238}{in}%
\TextAt{9}{79.669}{256.192}{\Character{22}}%
\TextAt{13}{79.226}{256.192}{\Character{20}}%
\TextAt{9}{89.286}{256.192}{\Character{61}}%
\TextAt{14}{101.711}{256.192}{o}%
\TextAt{30}{109.526}{254.554}{\Character{98}}%
\TextAt{21}{107.555}{252.539}{K}%
\TextAt{9}{115.758}{256.192}{\Character{47}}%
\TextAt{14}{121.611}{256.192}{p}%
\TextAt{30}{129.564}{254.554}{\Character{98}}%
\TextAt{21}{127.594}{252.539}{K}%
\TextAt{8}{135.797}{256.192}{,}%
\TextAt{8}{142.946}{256.192}{i.e.,}%
\TextAt{13}{209.542}{239.146}{c}%
\TextAt{9}{217.901}{239.146}{\Character{61}}%
\TextAt{13}{230.326}{239.146}{x}%
\TextAt{21}{236.978}{244.814}{\Character{39}}%
\TextAt{3}{242.496}{244.814}{\Character{0}}%
\TextAt{17}{249.083}{244.814}{1}%
\TextAt{17}{236.978}{236.215}{1}%
\TextAt{13}{253.815}{239.146}{a}%
\TextAt{17}{259.96}{237.353}{1}%
\TextAt{13}{264.692}{239.146}{\Character{59}}%
\TextAt{13}{293.827}{239.146}{x}%
\TextAt{17}{300.499}{237.353}{1}%
\NativeMathAt{308.552}{239.146}{11.955}{\in}%
\TextAt{13}{319.843}{239.146}{U}%
\TextAt{30}{329.764}{237.508}{\Character{98}}%
\TextAt{21}{327.793}{235.493}{K}%
\TextAt{13}{335.996}{239.146}{\Character{59}}%
\TextAt{13}{353.176}{239.146}{a}%
\TextAt{17}{359.34}{237.353}{1}%
\NativeMathAt{367.394}{239.146}{11.955}{\in}%
\TextAt{13}{378.684}{239.146}{U}%
\TextAt{17}{387.884}{245.256}{\Character{40}1\Character{41}}%
\TextAt{30}{388.605}{235.595}{\Character{98}}%
\TextAt{21}{386.635}{233.58}{K}%
\TextAt{13}{399.202}{239.146}{\Character{58}}%
\TextAt{8}{79.226}{214.397}{For}%
\TextAt{8}{101.572}{214.397}{similar}%
\TextAt{8}{141.385}{214.397}{reasons,}%
\TextAt{8}{187.459}{214.397}{we}%
\TextAt{8}{206.063}{214.397}{find}%
\TextAt{8}{230.848}{214.397}{that}%
\TextAt{13}{257.593}{214.397}{a}%
\TextAt{17}{263.738}{212.604}{1}%
\TextAt{9}{274.137}{214.397}{\Character{61}}%
\TextAt{13}{288.909}{214.397}{x}%
\TextAt{21}{295.561}{220.064}{\Character{39}}%
\TextAt{3}{301.079}{220.064}{\Character{0}}%
\TextAt{17}{307.665}{220.064}{1}%
\TextAt{17}{295.561}{211.466}{2}%
\TextAt{13}{312.398}{214.397}{a}%
\TextAt{17}{318.543}{212.604}{2}%
\TextAt{8}{323.275}{214.397}{,}%
\TextAt{8}{332.145}{214.397}{for}%
\TextAt{8}{351.406}{214.397}{some}%
\TextAt{13}{382.109}{214.397}{x}%
\TextAt{17}{388.761}{212.604}{2}%
\NativeMathAt{399.16}{214.397}{11.955}{\in}%
\TextAt{13}{412.797}{214.397}{U}%
\TextAt{17}{421.996}{220.507}{\Character{40}1\Character{41}}%
\TextAt{30}{422.718}{210.846}{\Character{98}}%
\TextAt{21}{420.747}{208.831}{K}%
\TextAt{8}{438.589}{214.397}{and}%
\TextAt{13}{462.728}{214.397}{a}%
\TextAt{17}{468.873}{212.604}{2}%
\NativeMathAt{479.272}{214.397}{11.955}{\in}%
\TextAt{13}{492.909}{214.397}{U}%
\TextAt{17}{502.108}{220.507}{\Character{40}2\Character{41}}%
\TextAt{30}{502.83}{210.846}{\Character{98}}%
\TextAt{21}{500.859}{208.831}{K}%
\TextAt{8}{513.427}{214.397}{,}%
\TextAt{8}{522.297}{214.397}{so}%
\TextAt{8}{79.226}{197.351}{that}%
\TextAt{13}{106.148}{197.351}{c}%
\TextAt{9}{117.153}{197.351}{\Character{61}}%
\TextAt{9}{132.217}{197.351}{\Character{40}}%
\TextAt{13}{136.777}{197.351}{x}%
\TextAt{17}{143.429}{195.558}{1}%
\TextAt{13}{148.161}{197.351}{x}%
\TextAt{17}{154.813}{195.558}{2}%
\TextAt{9}{159.545}{197.351}{\Character{41}}%
\TextAt{21}{164.098}{201.689}{\Character{39}}%
\TextAt{3}{169.616}{201.689}{\Character{0}}%
\TextAt{17}{176.203}{201.689}{1}%
\TextAt{13}{180.935}{197.351}{a}%
\TextAt{17}{187.08}{195.558}{2}%
\TextAt{8}{191.812}{197.351}{.}%
\TextAt{8}{204.919}{197.351}{Indeed,}%
\TextAt{8}{248.141}{197.351}{putting}%
\TextAt{13}{291.322}{197.351}{a}%
\TextAt{17}{297.467}{195.558}{1}%
\TextAt{9}{308.166}{197.351}{\Character{61}}%
\TextAt{9}{323.23}{197.351}{1}%
\TextAt{9}{332.782}{197.351}{\Character{43}}%
\TextAt{13}{345.625}{197.351}{b}%
\TextAt{17}{350.602}{195.558}{1}%
\TextAt{13}{355.335}{197.351}{\Character{25}}%
\TextAt{8}{362.404}{197.351}{,}%
\TextAt{13}{371.494}{197.351}{x}%
\TextAt{17}{378.146}{195.558}{2}%
\TextAt{9}{388.845}{197.351}{\Character{61}}%
\TextAt{9}{403.909}{197.351}{1}%
\TextAt{9}{413.461}{197.351}{\Character{43}}%
\TextAt{13}{426.305}{197.351}{y}%
\TextAt{17}{432.013}{195.558}{2}%
\TextAt{13}{436.745}{197.351}{\Character{25}}%
\TextAt{8}{443.814}{197.351}{,}%
\TextAt{8}{452.904}{197.351}{gives}%
\TextAt{13}{483.138}{197.351}{a}%
\TextAt{17}{489.283}{195.558}{1}%
\TextAt{13}{494.015}{197.351}{x}%
\TextAt{17}{500.667}{203.018}{1}%
\TextAt{3}{504.901}{203.018}{\Character{0}}%
\TextAt{21}{511.488}{203.018}{\Character{39}}%
\TextAt{17}{500.667}{194.42}{2}%
\NativeMathAt{523.471}{197.351}{11.955}{\equiv}%
\TextAt{9}{79.226}{180.305}{1}%
\TextAt{16}{87.394}{180.305}{\Character{0}}%
\TextAt{9}{99.008}{180.305}{\Character{40}}%
\TextAt{13}{103.56}{180.305}{y}%
\TextAt{21}{109.697}{185.973}{\Character{39}}%
\TextAt{17}{109.268}{177.374}{2}%
\TextAt{16}{118.028}{180.305}{\Character{0}}%
\TextAt{13}{129.642}{180.305}{y}%
\TextAt{17}{135.349}{178.512}{2}%
\TextAt{16}{142.397}{180.305}{\Character{0}}%
\TextAt{13}{154.01}{180.305}{b}%
\TextAt{17}{158.987}{178.512}{1}%
\TextAt{9}{163.72}{180.305}{\Character{41}}%
\TextAt{13}{168.272}{180.305}{\Character{25}}%
\TextAt{8}{178.662}{180.305}{mod}%
\TextAt{13}{204.1}{180.305}{\Character{25}}%
\TextAt{17}{211.169}{184.643}{2}%
\TextAt{8}{215.902}{180.305}{,}%
\TextAt{8}{222.917}{180.305}{i.e.,}%
\TextAt{8}{244.889}{180.305}{we}%
\TextAt{8}{261.949}{180.305}{have}%
\TextAt{8}{288.777}{180.305}{to}%
\TextAt{8}{302.92}{180.305}{solve}%
\TextAt{8}{331.433}{180.305}{the}%
\TextAt{8}{351.423}{180.305}{congruence}%
\TextAt{13}{411.737}{180.305}{y}%
\TextAt{21}{417.874}{185.972}{\Character{39}}%
\TextAt{17}{417.445}{177.374}{2}%
\TextAt{16}{426.205}{180.305}{\Character{0}}%
\TextAt{13}{437.819}{180.305}{y}%
\TextAt{17}{443.526}{178.512}{2}%
\TextAt{16}{450.574}{180.305}{\Character{0}}%
\TextAt{13}{462.187}{180.305}{b}%
\TextAt{17}{467.164}{178.512}{1}%
\NativeMathAt{475.217}{180.305}{11.955}{\equiv}%
\TextAt{9}{487.837}{180.305}{0}%
\TextAt{8}{497.011}{180.305}{mod}%
\TextAt{13}{522.449}{180.305}{\Character{25}}%
\TextAt{8}{529.518}{180.305}{,}%
\TextAt{8}{79.226}{163.259}{or}%
\TextAt{8}{92.538}{163.259}{equivalently}%
\TextAt{8}{156.588}{163.259}{the}%
\TextAt{8}{175.747}{163.259}{equation}%
\TextAt{9}{223.349}{163.259}{\Character{22}}%
\TextAt{13}{222.557}{163.259}{y}%
\TextAt{21}{228.694}{168.927}{q}%
\TextAt{17}{228.265}{160.328}{2}%
\TextAt{16}{233.874}{163.259}{\Character{0}}%
\TextAt{9}{244.58}{163.259}{\Character{22}}%
\TextAt{13}{243.788}{163.259}{y}%
\TextAt{17}{249.496}{161.466}{2}%
\TextAt{16}{254.844}{163.259}{\Character{0}}%
\TextAt{9}{264.32}{166.414}{\Character{22}}%
\TextAt{13}{264.758}{163.259}{b}%
\TextAt{17}{269.735}{161.466}{1}%
\TextAt{9}{277.788}{163.259}{\Character{61}}%
\TextAt{9}{290.197}{163.259}{0}%
\TextAt{8}{298.966}{163.259}{in}%
\TextAt{9}{312.063}{163.259}{\Character{22}}%
\TextAt{13}{311.62}{163.259}{\Character{20}}%
\TextAt{8}{318.359}{163.259}{.}%
\TextAt{8}{326.475}{163.259}{This}%
\TextAt{8}{352.196}{163.259}{is}%
\TextAt{8}{362.962}{163.259}{possible}%
\TextAt{8}{405.983}{163.259}{because}%
\TextAt{9}{448.729}{163.259}{\Character{22}}%
\TextAt{13}{448.286}{163.259}{\Character{20}}%
\TextAt{8}{457.924}{163.259}{is}%
\TextAt{8}{468.69}{163.259}{algebraically}%
\TextAt{8}{79.226}{146.213}{closed.}%
\TextAt{8}{118.304}{146.213}{Continuing}%
\TextAt{8}{178.797}{146.213}{in}%
\TextAt{8}{192.45}{146.213}{this}%
\TextAt{8}{215.272}{146.213}{way,}%
\TextAt{8}{241.286}{146.213}{we}%
\TextAt{8}{258.514}{146.213}{get}%
\TextAt{13}{171.968}{118.09}{c}%
\TextAt{9}{180.326}{118.09}{\Character{61}}%
\TextAt{9}{192.736}{118.09}{\Character{40}}%
\TextAt{13}{197.304}{118.09}{x}%
\TextAt{17}{203.956}{116.297}{1}%
\TextAt{13}{208.689}{118.09}{x}%
\TextAt{17}{215.341}{116.297}{2}%
\TextAt{16}{222.066}{118.09}{\Character{1}}%
\TextAt{16}{227.362}{118.09}{\Character{1}}%
\TextAt{16}{232.67}{118.09}{\Character{1}}%
\TextAt{13}{238.006}{118.09}{x}%
\TextAt{21}{244.658}{116.297}{n}%
\TextAt{9}{250.294}{118.09}{\Character{41}}%
\TextAt{21}{254.846}{123.026}{\Character{39}}%
\TextAt{3}{260.365}{123.026}{\Character{0}}%
\TextAt{17}{266.951}{123.026}{1}%
\TextAt{13}{271.684}{118.09}{a}%
\TextAt{21}{277.828}{116.297}{n}%
\TextAt{13}{283.465}{118.09}{\Character{59}}%
\TextAt{13}{312.6}{118.09}{x}%
\TextAt{21}{319.271}{116.297}{n}%
\NativeMathAt{328.229}{118.09}{11.955}{\in}%
\TextAt{13}{339.52}{118.09}{U}%
\TextAt{17}{348.719}{124.201}{\Character{40}}%
\TextAt{21}{352.012}{124.201}{n}%
\TextAt{3}{357.15}{124.201}{\Character{0}}%
\TextAt{17}{363.737}{124.201}{1\Character{41}}%
\TextAt{30}{349.44}{114.539}{\Character{98}}%
\TextAt{21}{347.47}{112.524}{K}%
\TextAt{13}{371.762}{118.09}{\Character{59}}%
\TextAt{13}{388.942}{118.09}{a}%
\TextAt{21}{395.106}{116.297}{n}%
\NativeMathAt{404.064}{118.09}{11.955}{\in}%
\TextAt{13}{415.355}{118.09}{U}%
\TextAt{17}{424.554}{124.201}{\Character{40}}%
\TextAt{21}{427.847}{124.201}{n}%
\TextAt{17}{432.985}{124.201}{\Character{41}}%
\TextAt{30}{425.275}{114.539}{\Character{98}}%
\TextAt{21}{423.305}{112.524}{K}%
\TextAt{13}{436.776}{118.09}{\Character{59}}%
\TextAt{8}{79.226}{89.967}{and}%
\TextAt{8}{101.831}{89.967}{passing}%
\TextAt{8}{142.776}{89.967}{to}%
\TextAt{8}{156.929}{89.967}{the}%
\TextAt{8}{176.928}{89.967}{limit}%
\TextAt{8}{204.734}{89.967}{finally}%
\TextAt{8}{240.024}{89.967}{gives}%
\TextAt{13}{268.547}{89.967}{c}%
\TextAt{9}{276.905}{89.967}{\Character{61}}%
\TextAt{13}{289.331}{89.967}{x}%
\TextAt{21}{295.983}{94.306}{\Character{39}}%
\TextAt{3}{301.501}{94.306}{\Character{0}}%
\TextAt{17}{308.088}{94.306}{1}%
\TextAt{8}{312.82}{89.967}{,}%
\TextAt{8}{319.843}{89.967}{where}%
\TextAt{13}{353.495}{89.967}{x}%
\TextAt{9}{363.468}{89.967}{\Character{61}}%
\TextAt{20}{375.894}{98.934}{\Character{81}}%
\NativeMathAt{387.185}{95.779}{7.97}{\infty}%
\TextAt{21}{387.185}{86.48}{n}%
\TextAt{17}{392.323}{86.48}{\Character{61}1}%
\TextAt{13}{405.634}{89.967}{x}%
\TextAt{21}{412.286}{88.174}{n}%
\NativeMathAt{421.243}{89.967}{11.955}{\in}%
\TextAt{13}{432.534}{89.967}{U}%
\TextAt{30}{442.455}{88.329}{\Character{98}}%
\TextAt{21}{440.485}{86.314}{K}%
\TextAt{8}{448.687}{89.967}{.}%
\TextAt{8}{457.083}{89.967}{The}%
\TextAt{8}{480.98}{89.967}{solvability}%
\TextAt{8}{79.226}{72.921}{of}%
\TextAt{8}{92.556}{72.921}{the}%
\TextAt{8}{112.713}{72.921}{equation}%
\TextAt{13}{160.521}{72.921}{x}%
\TextAt{21}{167.173}{77.26}{\Character{39}}%
\TextAt{16}{175.846}{72.921}{\Character{0}}%
\TextAt{13}{187.802}{72.921}{x}%
\TextAt{9}{197.775}{72.921}{\Character{61}}%
\TextAt{13}{210.2}{72.921}{c}%
\TextAt{8}{219.135}{72.921}{follows}%
\TextAt{8}{257.571}{72.921}{analogously,}%
\TextAt{8}{323.337}{72.921}{using}%
\TextAt{8}{353.966}{72.921}{the}%
\TextAt{8}{374.122}{72.921}{isomorphisms}%
\TextAt{14}{447.156}{72.921}{p}%
\TextAt{21}{453.139}{77.26}{n}%
\TextAt{30}{455.11}{69.591}{\Character{98}}%
\TextAt{21}{453.139}{67.576}{K}%
\TextAt{9}{461.342}{72.921}{\Character{47}}%
\TextAt{14}{467.195}{72.921}{p}%
\TextAt{21}{473.178}{77.869}{n}%
\TextAt{17}{478.316}{77.869}{\Character{43}1}%
\TextAt{30}{475.149}{69.37}{\Character{98}}%
\TextAt{21}{473.178}{67.356}{K}%
\NativeMathAt{492.956}{72.921}{11.955}{\cong}%
\TextAt{9}{506.018}{72.921}{\Character{22}}%
\TextAt{13}{505.575}{72.921}{\Character{20}}%
\TextAt{8}{512.314}{72.921}{.}%
\TextAt{8}{297.211}{43.033}{300}%
\Destination{78.23}{809.815}{nav.page.300}%
\Destination{79.23}{773.075}{nav.section.5.2}%
\Destination{79.36}{743.955}{nav.print.327}%
\Destination{79.36}{104.955}{nav.print.328}%
\Destination{79.36}{468.255}{nav.number.5.2.1}%
\Destination{79.23}{773.075}{outline.50}%
\Hotlink{352.709}{546.055}{13.995}{11.047}{nav.chapter.6}%
\Hotlink{378.918}{546.055}{6.858}{11.047}{nav.section.6.5}%
\Hotlink{136.635}{288.055}{9.763}{11.047}{nav.chapter.2}%
\Hotlink{153.11}{288.055}{30.936}{11.047}{nav.number.2.3.10}%
}

% Page 313 (body page 301)
\Page{%
\TextAt{2}{79.226}{796.943}{V.}%
\TextAt{2}{94.073}{796.943}{Local}%
\TextAt{2}{123.615}{796.943}{Class}%
\TextAt{2}{152.208}{796.943}{Field}%
\TextAt{2}{179.939}{796.943}{Theory}%
\Vector{q 0 G 0 g 0.299 w 7.226 18.729 m 460.769 18.729 l S Q}%
\TextAt{8}{97.159}{761.124}{Now}%
\TextAt{8}{123.819}{761.124}{let}%
\TextAt{13}{140.724}{761.124}{x}%
\NativeMathAt{150.696}{761.124}{11.955}{\in}%
\TextAt{14}{161.987}{761.124}{o}%
\TextAt{30}{169.802}{759.485}{\Character{98}}%
\TextAt{21}{167.832}{757.471}{K}%
\TextAt{8}{179.932}{761.124}{and}%
\TextAt{13}{202.694}{761.124}{x}%
\TextAt{21}{209.346}{765.462}{\Character{39}}%
\TextAt{9}{218.684}{761.124}{\Character{61}}%
\TextAt{13}{231.109}{761.124}{x}%
\TextAt{8}{237.761}{761.124}{.}%
\TextAt{8}{246.21}{761.124}{Then,}%
\TextAt{8}{280.019}{761.124}{for}%
\TextAt{8}{297.904}{761.124}{every}%
\TextAt{13}{328.796}{761.124}{n}%
\NativeMathAt{339.104}{761.124}{11.955}{\geq}%
\TextAt{9}{351.724}{761.124}{1}%
\TextAt{8}{357.577}{761.124}{,}%
\TextAt{8}{364.726}{761.124}{one}%
\TextAt{8}{386.185}{761.124}{has}%
\TextAt{13}{194.682}{732.123}{x}%
\TextAt{9}{204.655}{732.123}{\Character{61}}%
\TextAt{13}{217.081}{732.123}{x}%
\TextAt{21}{223.733}{730.329}{n}%
\TextAt{9}{232.026}{732.123}{\Character{43}}%
\TextAt{13}{243.787}{732.123}{\Character{25}}%
\TextAt{21}{250.856}{737.059}{n}%
\TextAt{13}{256.493}{732.123}{y}%
\TextAt{21}{262.2}{730.329}{n}%
\TextAt{8}{279.792}{732.123}{with}%
\TextAt{13}{306.452}{732.123}{x}%
\TextAt{21}{313.104}{730.329}{n}%
\NativeMathAt{322.061}{732.123}{11.955}{\in}%
\TextAt{14}{333.352}{732.123}{o}%
\TextAt{21}{339.196}{730.329}{K}%
\TextAt{8}{351.296}{732.123}{and}%
\TextAt{13}{374.059}{732.123}{y}%
\TextAt{21}{379.767}{730.329}{n}%
\NativeMathAt{388.724}{732.123}{11.955}{\in}%
\TextAt{14}{400.015}{732.123}{o}%
\TextAt{30}{407.83}{730.484}{\Character{98}}%
\TextAt{21}{405.859}{728.47}{K}%
\TextAt{13}{414.062}{732.123}{\Character{58}}%
\TextAt{8}{517.802}{732.123}{(*)}%
\TextAt{8}{79.226}{703.121}{Indeed,}%
\TextAt{8}{122.504}{703.121}{for}%
\TextAt{13}{141.986}{703.121}{n}%
\TextAt{9}{155.016}{703.121}{\Character{61}}%
\TextAt{9}{170.151}{703.121}{1}%
\TextAt{8}{181.511}{703.121}{we}%
\TextAt{8}{200.335}{703.121}{have}%
\TextAt{13}{228.927}{703.121}{x}%
\TextAt{9}{241.622}{703.121}{\Character{61}}%
\TextAt{13}{256.769}{703.121}{a}%
\TextAt{9}{266.659}{703.121}{\Character{43}}%
\TextAt{13}{279.509}{703.121}{\Character{25}}%
\TextAt{13}{286.563}{703.121}{b}%
\TextAt{8}{291.555}{703.121}{,}%
\TextAt{8}{300.701}{703.121}{with}%
\TextAt{13}{328.958}{703.121}{a}%
\NativeMathAt{341.145}{703.121}{11.955}{\in}%
\TextAt{14}{355.158}{703.121}{o}%
\TextAt{30}{362.973}{701.483}{\Character{101}}%
\TextAt{21}{361.002}{699.468}{K}%
\TextAt{8}{369.205}{703.121}{,}%
\TextAt{13}{378.351}{703.121}{b}%
\NativeMathAt{389.37}{703.121}{11.955}{\in}%
\TextAt{14}{403.383}{703.121}{o}%
\TextAt{30}{411.198}{701.483}{\Character{98}}%
\TextAt{21}{409.227}{699.468}{K}%
\TextAt{8}{417.43}{703.121}{,}%
\TextAt{8}{426.575}{703.121}{and}%
\TextAt{13}{450.935}{703.121}{x}%
\TextAt{21}{457.587}{707.46}{\Character{39}}%
\TextAt{9}{469.646}{703.121}{\Character{61}}%
\TextAt{13}{484.793}{703.121}{x}%
\TextAt{8}{496.94}{703.121}{implies}%
\TextAt{13}{79.226}{686.075}{a}%
\TextAt{21}{85.371}{690.414}{\Character{39}}%
\NativeMathAt{97.061}{686.075}{11.955}{\equiv}%
\TextAt{13}{112.033}{686.075}{a}%
\TextAt{8}{121.498}{686.075}{mod}%
\TextAt{13}{146.936}{686.075}{\Character{25}}%
\TextAt{8}{154.006}{686.075}{.}%
\TextAt{8}{166.595}{686.075}{Hence}%
\TextAt{13}{202.753}{686.075}{a}%
\TextAt{9}{214.572}{686.075}{\Character{61}}%
\TextAt{13}{229.35}{686.075}{x}%
\TextAt{17}{236.002}{684.282}{1}%
\TextAt{9}{244.332}{686.075}{\Character{43}}%
\TextAt{13}{257.034}{686.075}{\Character{25}}%
\TextAt{13}{264.088}{686.075}{c}%
\TextAt{8}{269.141}{686.075}{,}%
\TextAt{8}{278.016}{686.075}{with}%
\TextAt{13}{306.057}{686.075}{x}%
\TextAt{17}{312.709}{684.282}{1}%
\NativeMathAt{323.114}{686.075}{11.955}{\in}%
\TextAt{14}{336.758}{686.075}{o}%
\TextAt{21}{342.602}{684.282}{K}%
\TextAt{8}{350.805}{686.075}{,}%
\TextAt{13}{359.679}{686.075}{c}%
\NativeMathAt{370.391}{686.075}{11.955}{\in}%
\TextAt{14}{384.034}{686.075}{o}%
\TextAt{30}{391.849}{684.437}{\Character{101}}%
\TextAt{21}{389.879}{682.422}{K}%
\TextAt{8}{398.081}{686.075}{,}%
\TextAt{8}{406.956}{686.075}{and}%
\TextAt{8}{431.099}{686.075}{therefore}%
\TextAt{13}{481.579}{686.075}{x}%
\TextAt{9}{493.905}{686.075}{\Character{61}}%
\TextAt{13}{508.683}{686.075}{x}%
\TextAt{17}{515.335}{684.282}{1}%
\TextAt{9}{523.665}{686.075}{\Character{43}}%
\TextAt{13}{79.226}{669.029}{\Character{25}}%
\TextAt{9}{86.295}{669.029}{\Character{40}}%
\TextAt{13}{90.848}{669.029}{b}%
\TextAt{9}{99.186}{669.029}{\Character{43}}%
\TextAt{13}{111.652}{669.029}{c}%
\TextAt{9}{116.69}{669.029}{\Character{41}}%
\TextAt{9}{126.314}{669.029}{\Character{61}}%
\TextAt{13}{140.511}{669.029}{x}%
\TextAt{17}{147.163}{667.236}{1}%
\TextAt{9}{155.257}{669.029}{\Character{43}}%
\TextAt{13}{167.723}{669.029}{\Character{25}}%
\TextAt{13}{174.777}{669.029}{y}%
\TextAt{17}{180.5}{667.236}{1}%
\TextAt{8}{185.232}{669.029}{,}%
\TextAt{13}{193.673}{669.029}{y}%
\TextAt{17}{199.381}{667.236}{1}%
\NativeMathAt{209.195}{669.029}{11.955}{\in}%
\TextAt{14}{222.248}{669.029}{o}%
\TextAt{30}{230.062}{667.391}{\Character{98}}%
\TextAt{21}{228.092}{665.376}{K}%
\TextAt{8}{236.295}{669.029}{.}%
\TextAt{8}{247.844}{669.029}{The}%
\TextAt{8}{272.931}{669.029}{equation}%
\TextAt{13}{321.773}{669.029}{x}%
\TextAt{9}{333.508}{669.029}{\Character{61}}%
\TextAt{13}{347.694}{669.029}{x}%
\TextAt{21}{354.347}{667.236}{n}%
\TextAt{9}{363.344}{669.029}{\Character{43}}%
\TextAt{13}{375.81}{669.029}{\Character{25}}%
\TextAt{21}{382.879}{673.368}{n}%
\TextAt{13}{388.516}{669.029}{y}%
\TextAt{21}{394.223}{667.236}{n}%
\TextAt{8}{404.791}{669.029}{implies}%
\TextAt{8}{445.551}{669.029}{furthermore}%
\TextAt{8}{511.298}{669.029}{that}%
\TextAt{13}{79.226}{651.983}{y}%
\TextAt{21}{85.363}{656.322}{\Character{39}}%
\TextAt{21}{84.934}{649.028}{n}%
\TextAt{9}{95.697}{651.983}{\Character{61}}%
\TextAt{13}{109.119}{651.983}{y}%
\TextAt{21}{114.827}{650.19}{n}%
\TextAt{8}{120.463}{651.983}{,}%
\TextAt{8}{128.343}{651.983}{so}%
\TextAt{8}{143.298}{651.983}{that}%
\TextAt{8}{169.252}{651.983}{we}%
\TextAt{8}{187.064}{651.983}{get}%
\TextAt{8}{207.16}{651.983}{as}%
\TextAt{8}{222.115}{651.983}{above}%
\TextAt{13}{255.875}{651.983}{y}%
\TextAt{21}{261.583}{650.19}{n}%
\TextAt{9}{271.537}{651.983}{\Character{61}}%
\TextAt{13}{284.959}{651.983}{c}%
\TextAt{21}{289.997}{650.19}{n}%
\TextAt{9}{298.688}{651.983}{\Character{43}}%
\TextAt{13}{310.848}{651.983}{\Character{25}}%
\TextAt{13}{317.902}{651.983}{d}%
\TextAt{21}{324}{650.19}{n}%
\TextAt{8}{329.637}{651.983}{,}%
\TextAt{8}{337.517}{651.983}{with}%
\TextAt{13}{364.762}{651.983}{c}%
\TextAt{21}{369.8}{650.19}{n}%
\NativeMathAt{379.754}{651.983}{11.955}{\in}%
\TextAt{14}{392.041}{651.983}{o}%
\TextAt{21}{397.885}{650.19}{K}%
\TextAt{8}{406.088}{651.983}{,}%
\TextAt{13}{413.968}{651.983}{d}%
\TextAt{21}{420.051}{650.19}{n}%
\NativeMathAt{430.005}{651.983}{11.955}{\in}%
\TextAt{14}{442.292}{651.983}{o}%
\TextAt{30}{450.107}{650.345}{\Character{98}}%
\TextAt{21}{448.137}{648.331}{K}%
\TextAt{8}{456.339}{651.983}{,}%
\TextAt{8}{464.22}{651.983}{and}%
\TextAt{8}{487.567}{651.983}{therefore}%
\TextAt{13}{79.226}{634.938}{x}%
\TextAt{9}{89.199}{634.938}{\Character{61}}%
\TextAt{9}{101.609}{634.938}{\Character{40}}%
\TextAt{13}{106.177}{634.938}{x}%
\TextAt{21}{112.829}{633.144}{n}%
\TextAt{9}{119.566}{634.938}{\Character{43}}%
\TextAt{13}{129.772}{634.938}{c}%
\TextAt{21}{134.81}{633.144}{n}%
\TextAt{13}{140.446}{634.938}{\Character{25}}%
\TextAt{21}{147.515}{639.276}{n}%
\TextAt{9}{153.152}{634.938}{\Character{41}}%
\TextAt{9}{158.795}{634.938}{\Character{43}}%
\TextAt{13}{169.011}{634.938}{\Character{25}}%
\TextAt{21}{176.08}{639.276}{n}%
\TextAt{17}{181.218}{639.276}{\Character{43}1}%
\TextAt{13}{192.537}{634.938}{d}%
\TextAt{21}{198.62}{633.144}{n}%
\TextAt{9}{207.577}{634.938}{\Character{61}}%
\TextAt{13}{220.002}{634.938}{x}%
\TextAt{21}{226.654}{633.144}{n}%
\TextAt{17}{231.793}{633.144}{\Character{43}1}%
\TextAt{9}{244.212}{634.938}{\Character{43}}%
\TextAt{13}{254.418}{634.938}{\Character{25}}%
\TextAt{21}{261.487}{639.276}{n}%
\TextAt{17}{266.626}{639.276}{\Character{43}1}%
\TextAt{13}{277.944}{634.938}{y}%
\TextAt{21}{283.652}{633.144}{n}%
\TextAt{17}{288.79}{633.144}{\Character{43}1}%
\TextAt{8}{300.109}{634.938}{,}%
\TextAt{8}{306.65}{634.938}{for}%
\TextAt{8}{323.774}{634.938}{some}%
\TextAt{13}{352.339}{634.938}{x}%
\TextAt{21}{358.991}{633.144}{n}%
\TextAt{17}{364.129}{633.144}{\Character{43}1}%
\NativeMathAt{378.769}{634.938}{11.955}{\in}%
\TextAt{14}{390.06}{634.938}{o}%
\TextAt{21}{395.904}{633.144}{K}%
\TextAt{8}{404.107}{634.938}{,}%
\TextAt{13}{410.648}{634.938}{y}%
\TextAt{21}{416.355}{633.144}{n}%
\TextAt{17}{421.493}{633.144}{\Character{43}1}%
\NativeMathAt{436.133}{634.938}{11.955}{\in}%
\TextAt{14}{447.424}{634.938}{o}%
\TextAt{30}{455.239}{633.299}{\Character{98}}%
\TextAt{21}{453.268}{631.285}{K}%
\TextAt{8}{461.471}{634.938}{.}%
\TextAt{8}{469.666}{634.938}{Now}%
\TextAt{8}{495.565}{634.938}{passing}%
\TextAt{8}{79.226}{617.892}{to}%
\TextAt{8}{93.536}{617.892}{the}%
\TextAt{8}{113.693}{617.892}{limit}%
\TextAt{8}{141.656}{617.892}{in}%
\TextAt{8}{155.309}{617.892}{the}%
\TextAt{8}{175.465}{617.892}{equation}%
\TextAt{8}{223.274}{617.892}{(*)}%
\TextAt{8}{242.139}{617.892}{gives}%
\TextAt{13}{270.82}{617.892}{x}%
\TextAt{9}{280.793}{617.892}{\Character{61}}%
\TextAt{8}{293.218}{617.892}{lim}%
\TextAt{21}{309.477}{616.098}{n}%
\NativeMathAt{314.615}{616.098}{7.97}{\rightarrow}%
\NativeMathAt{323.079}{616.098}{7.97}{\infty}%
\TextAt{13}{334.043}{617.892}{x}%
\TextAt{21}{340.695}{616.098}{n}%
\NativeMathAt{349.652}{617.892}{11.955}{\in}%
\TextAt{14}{360.943}{617.892}{o}%
\TextAt{21}{366.787}{616.098}{K}%
\TextAt{8}{374.99}{617.892}{,}%
\TextAt{8}{382.139}{617.892}{because}%
\TextAt{13}{425.441}{617.892}{K}%
\TextAt{8}{440.156}{617.892}{is}%
\TextAt{8}{451.92}{617.892}{complete.}%
\Vector{q 0 G 0 g 0.399 w 452.334 -156.108 m 452.334 -148.039 l S Q}%
\Vector{q 0 G 0 g 0.399 w 452.534 -148.238 m 459.707 -148.238 l S Q}%
\Vector{q 0 G 0 g 0.399 w 452.534 -155.909 m 459.707 -155.909 l S Q}%
\Vector{q 0 G 0 g 0.399 w 459.906 -156.108 m 459.906 -148.039 l S Q}%
\TextAt{8}{97.159}{588.89}{For}%
\TextAt{8}{119.009}{588.89}{a}%
\TextAt{8}{129.645}{588.89}{power}%
\TextAt{8}{164.669}{588.89}{series}%
\TextAt{13}{196.885}{588.89}{F}%
\TextAt{9}{206.088}{588.89}{\Character{40}}%
\TextAt{13}{210.64}{588.89}{X}%
\TextAt{17}{220.356}{587.097}{1}%
\TextAt{13}{225.088}{588.89}{\Character{59}}%
\TextAt{13}{230.312}{588.89}{\Character{58}}%
\TextAt{13}{235.549}{588.89}{\Character{58}}%
\TextAt{13}{240.773}{588.89}{\Character{58}}%
\TextAt{13}{246.01}{588.89}{\Character{59}}%
\TextAt{13}{251.246}{588.89}{X}%
\TextAt{21}{261.024}{587.097}{n}%
\TextAt{9}{266.66}{588.89}{\Character{41}}%
\NativeMathAt{276.034}{588.89}{11.955}{\in}%
\TextAt{14}{288.826}{588.89}{o}%
\TextAt{30}{296.641}{587.252}{\Character{98}}%
\TextAt{21}{294.67}{585.237}{K}%
\TextAt{9}{302.873}{588.89}{\Character{91}\Character{91}}%
\TextAt{13}{309.376}{588.89}{X}%
\TextAt{17}{319.091}{587.097}{1}%
\TextAt{13}{323.824}{588.89}{\Character{59}}%
\TextAt{13}{329.048}{588.89}{\Character{58}}%
\TextAt{13}{334.285}{588.89}{\Character{58}}%
\TextAt{13}{339.509}{588.89}{\Character{58}}%
\TextAt{13}{344.746}{588.89}{\Character{59}}%
\TextAt{13}{349.982}{588.89}{X}%
\TextAt{21}{359.76}{587.097}{n}%
\TextAt{9}{365.396}{588.89}{\Character{93}\Character{93}}%
\TextAt{8}{371.899}{588.89}{,}%
\TextAt{8}{380.149}{588.89}{let}%
\TextAt{13}{397.935}{588.89}{F}%
\TextAt{21}{407.138}{593.229}{\Character{39}}%
\TextAt{8}{417.933}{588.89}{be}%
\TextAt{8}{434.737}{588.89}{the}%
\TextAt{8}{455.774}{588.89}{power}%
\TextAt{8}{490.799}{588.89}{series}%
\TextAt{8}{523.014}{588.89}{in}%
\TextAt{14}{79.226}{571.844}{o}%
\TextAt{30}{87.041}{570.206}{\Character{98}}%
\TextAt{21}{85.07}{568.191}{K}%
\TextAt{9}{93.273}{571.844}{\Character{91}\Character{91}}%
\TextAt{13}{99.776}{571.844}{X}%
\TextAt{17}{109.492}{570.051}{1}%
\TextAt{13}{114.224}{571.844}{\Character{59}}%
\TextAt{13}{119.448}{571.844}{\Character{58}}%
\TextAt{13}{124.685}{571.844}{\Character{58}}%
\TextAt{13}{129.909}{571.844}{\Character{58}}%
\TextAt{13}{135.146}{571.844}{\Character{59}}%
\TextAt{13}{140.382}{571.844}{X}%
\TextAt{21}{150.16}{570.051}{n}%
\TextAt{9}{155.796}{571.844}{\Character{93}\Character{93}}%
\TextAt{8}{166.737}{571.844}{which}%
\TextAt{8}{200.763}{571.844}{arises}%
\TextAt{8}{233.295}{571.844}{from}%
\TextAt{13}{261.475}{571.844}{F}%
\TextAt{8}{275.116}{571.844}{by}%
\TextAt{8}{291.915}{571.844}{applying}%
\TextAt{13}{340.264}{571.844}{\Character{39}}%
\TextAt{8}{352.369}{571.844}{to}%
\TextAt{8}{367.219}{571.844}{the}%
\TextAt{8}{387.916}{571.844}{coefficients}%
\TextAt{8}{447.693}{571.844}{of}%
\TextAt{13}{461.563}{571.844}{F}%
\TextAt{8}{470.767}{571.844}{.}%
\TextAt{8}{480.835}{571.844}{A}%
\TextAt{12}{494.047}{571.844}{Lubin-}%
\TextAt{12}{79.226}{554.798}{Tate}%
\TextAt{12}{110.529}{554.798}{series}%
\TextAt{8}{147.145}{554.798}{for}%
\TextAt{8}{165.626}{554.798}{a}%
\TextAt{8}{175.977}{554.798}{prime}%
\TextAt{8}{209.737}{554.798}{element}%
\TextAt{13}{253.574}{554.798}{\Character{25}}%
\TextAt{8}{265.137}{554.798}{of}%
\TextAt{20}{281.801}{557.82}{\Character{98}}%
\TextAt{13}{279.062}{554.798}{K}%
\TextAt{8}{294.373}{554.798}{is}%
\TextAt{8}{306.733}{554.798}{by}%
\TextAt{8}{323.588}{554.798}{definition}%
\TextAt{8}{376.212}{554.798}{a}%
\TextAt{8}{386.563}{554.798}{power}%
\TextAt{8}{421.303}{554.798}{series}%
\TextAt{13}{453.233}{554.798}{e}%
\TextAt{9}{458.659}{554.798}{\Character{40}}%
\TextAt{13}{463.211}{554.798}{X}%
\TextAt{9}{473.866}{554.798}{\Character{41}}%
\NativeMathAt{482.755}{554.798}{11.955}{\in}%
\TextAt{14}{495.061}{554.798}{o}%
\TextAt{30}{502.875}{553.16}{\Character{98}}%
\TextAt{21}{500.905}{551.145}{K}%
\TextAt{9}{509.108}{554.798}{\Character{91}\Character{91}}%
\TextAt{13}{515.611}{554.798}{X}%
\TextAt{9}{526.266}{554.798}{\Character{93}\Character{93}}%
\TextAt{8}{79.226}{537.752}{with}%
\TextAt{8}{105.886}{537.752}{the}%
\TextAt{8}{126.043}{537.752}{properties}%
\TextAt{13}{180.389}{508.751}{e}%
\TextAt{9}{185.815}{508.751}{\Character{40}}%
\TextAt{13}{190.367}{508.751}{X}%
\TextAt{9}{201.022}{508.751}{\Character{41}}%
\NativeMathAt{208.896}{508.751}{11.955}{\equiv}%
\TextAt{13}{221.515}{508.751}{\Character{25}}%
\TextAt{13}{228.569}{508.751}{X}%
\TextAt{8}{242.56}{508.751}{mod}%
\TextAt{8}{267.989}{508.751}{deg}%
\TextAt{9}{287.553}{508.751}{2}%
\TextAt{8}{305.361}{508.751}{and}%
\TextAt{13}{336.181}{508.751}{e}%
\TextAt{9}{341.607}{508.751}{\Character{40}}%
\TextAt{13}{346.159}{508.751}{X}%
\TextAt{9}{356.814}{508.751}{\Character{41}}%
\NativeMathAt{364.687}{508.751}{11.955}{\equiv}%
\TextAt{13}{377.307}{508.751}{X}%
\TextAt{21}{387.962}{513.688}{q}%
\TextAt{8}{395.847}{508.751}{mod}%
\TextAt{13}{421.285}{508.751}{\Character{25}}%
\TextAt{13}{428.339}{508.751}{\Character{59}}%
\TextAt{8}{79.226}{479.75}{where}%
\TextAt{13}{113.18}{479.75}{q}%
\TextAt{9}{122.366}{479.75}{\Character{61}}%
\TextAt{13}{135.038}{479.75}{q}%
\TextAt{21}{140.229}{477.957}{K}%
\TextAt{8}{152.473}{479.75}{denotes,}%
\TextAt{8}{198.239}{479.75}{as}%
\TextAt{8}{212.754}{479.75}{always,}%
\TextAt{8}{253.654}{479.75}{the}%
\TextAt{8}{273.955}{479.75}{number}%
\TextAt{8}{316.696}{479.75}{of}%
\TextAt{8}{330.171}{479.75}{elements}%
\TextAt{8}{378.172}{479.75}{in}%
\TextAt{8}{391.969}{479.75}{the}%
\TextAt{8}{412.27}{479.75}{residue}%
\TextAt{8}{452.142}{479.75}{class}%
\TextAt{8}{479.724}{479.75}{field}%
\TextAt{8}{505.225}{479.75}{of}%
\TextAt{13}{518.7}{479.75}{K}%
\TextAt{8}{529.518}{479.75}{.}%
\TextAt{8}{79.226}{462.704}{The}%
\TextAt{8}{103.809}{462.704}{totality}%
\TextAt{8}{145.977}{462.704}{of}%
\TextAt{8}{159.836}{462.704}{all}%
\TextAt{8}{176.624}{462.704}{Lubin-Tate}%
\TextAt{8}{238.112}{462.704}{series}%
\TextAt{8}{269.975}{462.704}{is}%
\TextAt{8}{282.268}{462.704}{denoted}%
\TextAt{8}{327.019}{462.704}{by}%
\TextAt{16}{343.807}{462.704}{\Character{69}}%
\TextAt{21}{350.116}{460.911}{\Character{25}}%
\TextAt{8}{355.763}{462.704}{.}%
\TextAt{8}{365.798}{462.704}{In}%
\TextAt{16}{380.948}{462.704}{\Character{69}}%
\TextAt{21}{387.257}{460.911}{\Character{25}}%
\TextAt{8}{397.33}{462.704}{there}%
\TextAt{8}{427.771}{462.704}{are}%
\TextAt{8}{447.811}{462.704}{in}%
\TextAt{8}{461.992}{462.704}{particular}%
\TextAt{8}{516.51}{462.704}{the}%
\TextAt{8}{79.226}{445.658}{polynomials}%
\TextAt{13}{183.613}{428.612}{e}%
\TextAt{9}{189.039}{428.612}{\Character{40}}%
\TextAt{13}{193.591}{428.612}{X}%
\TextAt{9}{204.246}{428.612}{\Character{41}}%
\TextAt{9}{212.101}{428.612}{\Character{61}}%
\TextAt{13}{224.545}{428.612}{uX}%
\TextAt{21}{241.862}{433.548}{q}%
\TextAt{9}{249.084}{428.612}{\Character{43}}%
\TextAt{13}{260.845}{428.612}{\Character{25}}%
\TextAt{9}{267.915}{428.612}{\Character{40}}%
\TextAt{13}{272.467}{428.612}{a}%
\TextAt{21}{278.612}{426.819}{q}%
\TextAt{3}{282.679}{426.819}{\Character{0}}%
\TextAt{17}{289.265}{426.819}{1}%
\TextAt{13}{293.997}{428.612}{X}%
\TextAt{21}{304.652}{433.548}{q}%
\TextAt{3}{308.719}{433.548}{\Character{0}}%
\TextAt{17}{315.306}{433.548}{1}%
\TextAt{9}{322.695}{428.612}{\Character{43}}%
\TextAt{16}{334.456}{428.612}{\Character{1}}%
\TextAt{16}{339.752}{428.612}{\Character{1}}%
\TextAt{16}{345.06}{428.612}{\Character{1}}%
\TextAt{9}{351.06}{428.612}{\Character{43}}%
\TextAt{13}{362.822}{428.612}{a}%
\TextAt{17}{368.967}{426.819}{2}%
\TextAt{13}{373.699}{428.612}{X}%
\TextAt{17}{384.354}{433.548}{2}%
\TextAt{9}{389.086}{428.612}{\Character{41}}%
\TextAt{9}{396.283}{428.612}{\Character{43}}%
\TextAt{13}{408.057}{428.612}{\Character{25}}%
\TextAt{13}{415.111}{428.612}{X}%
\TextAt{13}{425.117}{428.612}{\Character{59}}%
\TextAt{8}{79.226}{405.091}{where}%
\TextAt{13}{112.518}{405.091}{u\Character{59}}%
\TextAt{13}{124.402}{405.091}{a}%
\TextAt{21}{130.57}{403.297}{i}%
\NativeMathAt{137.272}{405.091}{11.955}{\in}%
\TextAt{14}{148.563}{405.091}{o}%
\TextAt{30}{156.378}{403.452}{\Character{98}}%
\TextAt{21}{154.407}{401.438}{K}%
\TextAt{8}{165.99}{405.091}{and}%
\TextAt{13}{188.236}{405.091}{u}%
\NativeMathAt{198.219}{405.091}{11.955}{\equiv}%
\TextAt{9}{210.838}{405.091}{1}%
\TextAt{8}{220.012}{405.091}{mod}%
\TextAt{13}{245.45}{405.091}{\Character{25}}%
\TextAt{8}{252.519}{405.091}{.}%
\TextAt{8}{260.795}{405.091}{These}%
\TextAt{8}{294.147}{405.091}{are}%
\TextAt{8}{313.141}{405.091}{called}%
\TextAt{8}{345.788}{405.091}{the}%
\TextAt{12}{365.427}{405.091}{Lubin-Tate}%
\TextAt{12}{434.172}{405.091}{polynomials}%
\TextAt{8}{504.337}{405.091}{.}%
\TextAt{8}{512.613}{405.091}{The}%
\TextAt{8}{79.226}{388.045}{simplest}%
\TextAt{8}{125.044}{388.045}{one}%
\TextAt{8}{146.676}{388.045}{among}%
\TextAt{8}{184.579}{388.045}{them}%
\TextAt{8}{214.664}{388.045}{is}%
\TextAt{8}{226.601}{388.045}{the}%
\TextAt{8}{246.93}{388.045}{polynomial}%
\TextAt{13}{307.596}{388.045}{X}%
\TextAt{21}{318.251}{392.383}{q}%
\TextAt{9}{325.59}{388.045}{\Character{43}}%
\TextAt{13}{337.469}{388.045}{\Character{25}}%
\TextAt{13}{344.523}{388.045}{X}%
\TextAt{8}{355.193}{388.045}{.}%
\TextAt{8}{364.16}{388.045}{In}%
\TextAt{8}{378.954}{388.045}{the}%
\TextAt{8}{399.283}{388.045}{case}%
\TextAt{13}{424.227}{388.045}{K}%
\TextAt{9}{438.66}{388.045}{\Character{61}}%
\TextAt{18}{451.38}{388.045}{Q}%
\TextAt{21}{460.678}{386.251}{p}%
\TextAt{8}{469.509}{388.045}{for}%
\TextAt{8}{487.567}{388.045}{example,}%
\TextAt{13}{79.226}{370.999}{e}%
\TextAt{9}{84.652}{370.999}{\Character{40}}%
\TextAt{13}{89.204}{370.999}{X}%
\TextAt{9}{99.859}{370.999}{\Character{41}}%
\TextAt{9}{107.714}{370.999}{\Character{61}}%
\TextAt{9}{120.135}{370.999}{\Character{40}1}%
\TextAt{9}{133.178}{370.999}{\Character{43}}%
\TextAt{13}{144.981}{370.999}{X}%
\TextAt{9}{155.636}{370.999}{\Character{41}}%
\TextAt{21}{160.188}{375.337}{p}%
\TextAt{16}{167.606}{370.999}{\Character{0}}%
\TextAt{9}{179.561}{370.999}{1}%
\TextAt{8}{189.312}{370.999}{is}%
\TextAt{8}{201.075}{370.999}{a}%
\TextAt{8}{210.831}{370.999}{Lubin-Tate}%
\TextAt{8}{271.79}{370.999}{polynomial}%
\TextAt{8}{332.283}{370.999}{for}%
\TextAt{8}{350.168}{370.999}{the}%
\TextAt{8}{370.325}{370.999}{prime}%
\TextAt{8}{403.488}{370.999}{element}%
\TextAt{13}{446.73}{370.999}{p}%
\TextAt{8}{452.605}{370.999}{.}%
\TextAt{12}{79.226}{344.986}{Proposition}%
\TextAt{12}{152.006}{344.986}{(2.2).}%
\TextAt{8}{187.382}{344.986}{Let}%
\TextAt{13}{208.411}{344.986}{\Character{25}}%
\TextAt{8}{219.438}{344.986}{and}%
\TextAt{9}{242.868}{344.986}{\Character{22}}%
\TextAt{13}{242.26}{344.986}{\Character{25}}%
\TextAt{8}{253.287}{344.986}{be}%
\TextAt{8}{269.271}{344.986}{prime}%
\TextAt{8}{302.495}{344.986}{elements}%
\TextAt{8}{350.411}{344.986}{of}%
\TextAt{20}{366.539}{348.008}{\Character{98}}%
\TextAt{13}{363.801}{344.986}{K}%
\TextAt{8}{374.619}{344.986}{,}%
\TextAt{8}{381.843}{344.986}{and}%
\TextAt{8}{404.665}{344.986}{let}%
\TextAt{13}{421.63}{344.986}{e}%
\TextAt{9}{427.055}{344.986}{\Character{40}}%
\TextAt{13}{431.608}{344.986}{X}%
\TextAt{9}{442.263}{344.986}{\Character{41}}%
\NativeMathAt{450.238}{344.986}{11.955}{\in}%
\TextAt{16}{461.619}{344.986}{\Character{69}}%
\TextAt{21}{467.941}{343.193}{\Character{25}}%
\TextAt{8}{473.588}{344.986}{,}%
\TextAt{9}{481.248}{344.986}{\Character{22}}%
\TextAt{13}{480.812}{344.986}{e}%
\TextAt{9}{486.237}{344.986}{\Character{40}}%
\TextAt{13}{490.789}{344.986}{X}%
\TextAt{9}{501.445}{344.986}{\Character{41}}%
\NativeMathAt{509.42}{344.986}{11.955}{\in}%
\TextAt{16}{520.801}{344.986}{\Character{69}}%
\TextAt{17}{527.58}{343.193}{\Character{22}}%
\TextAt{21}{527.123}{343.193}{\Character{25}}%
\TextAt{8}{79.226}{327.94}{be}%
\TextAt{8}{95.656}{327.94}{Lubin-Tate}%
\TextAt{8}{157.122}{327.94}{series.}%
\TextAt{8}{194.525}{327.94}{Let}%
\TextAt{13}{216.001}{327.94}{L}%
\TextAt{9}{223.965}{327.94}{\Character{40}}%
\TextAt{13}{228.518}{327.94}{X}%
\TextAt{17}{238.233}{326.147}{1}%
\TextAt{13}{242.965}{327.94}{\Character{59}}%
\TextAt{13}{248.189}{327.94}{\Character{58}}%
\TextAt{13}{253.426}{327.94}{\Character{58}}%
\TextAt{13}{258.662}{327.94}{\Character{58}}%
\TextAt{13}{263.887}{327.94}{\Character{59}}%
\TextAt{13}{269.123}{327.94}{X}%
\TextAt{21}{278.901}{326.147}{n}%
\TextAt{9}{284.538}{327.94}{\Character{41}}%
\TextAt{9}{293.253}{327.94}{\Character{61}}%
\TextAt{20}{306.561}{336.907}{\Character{80}}%
\TextAt{21}{319.18}{333.752}{n}%
\TextAt{21}{319.18}{324.453}{i}%
\TextAt{17}{322.063}{324.453}{\Character{61}1}%
\TextAt{13}{335.375}{327.94}{a}%
\TextAt{21}{341.52}{326.147}{i}%
\TextAt{13}{344.901}{327.94}{X}%
\TextAt{21}{354.616}{326.147}{i}%
\TextAt{8}{362.401}{327.94}{be}%
\TextAt{8}{378.831}{327.94}{a}%
\TextAt{8}{389.093}{327.94}{linear}%
\TextAt{8}{422.117}{327.94}{form}%
\TextAt{8}{450.263}{327.94}{with}%
\TextAt{8}{477.429}{327.94}{coefficients}%
\TextAt{13}{79.226}{310.894}{a}%
\TextAt{21}{85.371}{309.101}{i}%
\NativeMathAt{92.073}{310.894}{11.955}{\in}%
\TextAt{14}{103.364}{310.894}{o}%
\TextAt{30}{111.179}{309.256}{\Character{98}}%
\TextAt{21}{109.208}{307.241}{K}%
\TextAt{8}{121.308}{310.894}{such}%
\TextAt{8}{147.705}{310.894}{that}%
\TextAt{13}{213.332}{293.848}{\Character{25}}%
\TextAt{13}{220.386}{293.848}{L}%
\TextAt{9}{228.366}{293.848}{\Character{40}}%
\TextAt{13}{232.918}{293.848}{X}%
\TextAt{17}{242.634}{292.055}{1}%
\TextAt{13}{247.366}{293.848}{\Character{59}}%
\TextAt{13}{252.59}{293.848}{\Character{58}}%
\TextAt{13}{257.827}{293.848}{\Character{58}}%
\TextAt{13}{263.051}{293.848}{\Character{58}}%
\TextAt{13}{268.288}{293.848}{\Character{59}}%
\TextAt{13}{273.524}{293.848}{X}%
\TextAt{21}{283.302}{292.055}{n}%
\TextAt{9}{288.938}{293.848}{\Character{41}}%
\TextAt{9}{296.793}{293.848}{\Character{61}}%
\TextAt{9}{309.824}{293.848}{\Character{22}}%
\TextAt{13}{309.237}{293.848}{\Character{25}}%
\TextAt{13}{316.291}{293.848}{L}%
\TextAt{21}{324.271}{298.785}{\Character{39}}%
\TextAt{9}{330.287}{293.848}{\Character{40}}%
\TextAt{13}{334.839}{293.848}{X}%
\TextAt{17}{344.555}{292.055}{1}%
\TextAt{13}{349.287}{293.848}{\Character{59}}%
\TextAt{13}{354.511}{293.848}{\Character{58}}%
\TextAt{13}{359.748}{293.848}{\Character{58}}%
\TextAt{13}{364.972}{293.848}{\Character{58}}%
\TextAt{13}{370.209}{293.848}{\Character{59}}%
\TextAt{13}{375.445}{293.848}{X}%
\TextAt{21}{385.223}{292.055}{n}%
\TextAt{9}{390.859}{293.848}{\Character{41}}%
\TextAt{13}{395.412}{293.848}{\Character{58}}%
\TextAt{8}{79.226}{270.327}{Then}%
\TextAt{8}{109.996}{270.327}{there}%
\TextAt{8}{140.121}{270.327}{is}%
\TextAt{8}{152.097}{270.327}{a}%
\TextAt{8}{162.066}{270.327}{uniquely}%
\TextAt{8}{209.752}{270.327}{determined}%
\TextAt{8}{271.092}{270.327}{power}%
\TextAt{8}{305.449}{270.327}{series}%
\TextAt{13}{336.996}{270.327}{F}%
\TextAt{9}{346.199}{270.327}{\Character{40}}%
\TextAt{13}{350.752}{270.327}{X}%
\TextAt{17}{360.467}{268.533}{1}%
\TextAt{13}{365.199}{270.327}{\Character{59}}%
\TextAt{13}{370.423}{270.327}{\Character{58}}%
\TextAt{13}{375.66}{270.327}{\Character{58}}%
\TextAt{13}{380.896}{270.327}{\Character{58}}%
\TextAt{13}{386.121}{270.327}{\Character{59}}%
\TextAt{13}{391.357}{270.327}{X}%
\TextAt{21}{401.135}{268.533}{n}%
\TextAt{9}{406.772}{270.327}{\Character{41}}%
\NativeMathAt{415.007}{270.327}{11.955}{\in}%
\TextAt{14}{426.661}{270.327}{o}%
\TextAt{30}{434.476}{268.688}{\Character{98}}%
\TextAt{21}{432.505}{266.674}{K}%
\TextAt{9}{440.708}{270.327}{\Character{91}\Character{91}}%
\TextAt{13}{447.211}{270.327}{X}%
\TextAt{17}{456.926}{268.533}{1}%
\TextAt{13}{461.659}{270.327}{\Character{59}}%
\TextAt{13}{466.883}{270.327}{\Character{58}}%
\TextAt{13}{472.12}{270.327}{\Character{58}}%
\TextAt{13}{477.344}{270.327}{\Character{58}}%
\TextAt{13}{482.581}{270.327}{\Character{59}}%
\TextAt{13}{487.817}{270.327}{X}%
\TextAt{21}{497.595}{268.533}{n}%
\TextAt{9}{503.231}{270.327}{\Character{93}\Character{93}}%
\TextAt{8}{513.845}{270.327}{sat-}%
\TextAt{8}{79.226}{253.281}{isfying}%
\TextAt{13}{195.61}{236.235}{F}%
\TextAt{9}{204.814}{236.235}{\Character{40}}%
\TextAt{13}{209.366}{236.235}{X}%
\TextAt{17}{219.081}{234.442}{1}%
\TextAt{13}{223.814}{236.235}{\Character{59}}%
\TextAt{13}{229.038}{236.235}{\Character{58}}%
\TextAt{13}{234.275}{236.235}{\Character{58}}%
\TextAt{13}{239.499}{236.235}{\Character{58}}%
\TextAt{13}{244.736}{236.235}{\Character{59}}%
\TextAt{13}{249.972}{236.235}{X}%
\TextAt{21}{259.75}{234.442}{n}%
\TextAt{9}{265.386}{236.235}{\Character{41}}%
\NativeMathAt{273.259}{236.235}{11.955}{\equiv}%
\TextAt{13}{285.878}{236.235}{L}%
\TextAt{9}{293.843}{236.235}{\Character{40}}%
\TextAt{13}{298.395}{236.235}{X}%
\TextAt{17}{308.111}{234.442}{1}%
\TextAt{13}{312.843}{236.235}{\Character{59}}%
\TextAt{13}{318.067}{236.235}{\Character{58}}%
\TextAt{13}{323.304}{236.235}{\Character{58}}%
\TextAt{13}{328.528}{236.235}{\Character{58}}%
\TextAt{13}{333.765}{236.235}{\Character{59}}%
\TextAt{13}{339.001}{236.235}{X}%
\TextAt{21}{348.779}{234.442}{n}%
\TextAt{9}{354.415}{236.235}{\Character{41}}%
\TextAt{8}{362.288}{236.235}{mod}%
\TextAt{8}{387.717}{236.235}{deg}%
\TextAt{9}{407.281}{236.235}{2}%
\TextAt{13}{413.134}{236.235}{\Character{59}}%
\TextAt{13}{197.367}{212.713}{e}%
\TextAt{9}{202.793}{212.713}{\Character{40}}%
\TextAt{13}{207.345}{212.713}{F}%
\TextAt{9}{216.549}{212.713}{\Character{40}}%
\TextAt{13}{221.101}{212.713}{X}%
\TextAt{17}{230.816}{210.92}{1}%
\TextAt{13}{235.549}{212.713}{\Character{59}}%
\TextAt{13}{240.773}{212.713}{\Character{58}}%
\TextAt{13}{246.01}{212.713}{\Character{58}}%
\TextAt{13}{251.234}{212.713}{\Character{58}}%
\TextAt{13}{256.471}{212.713}{\Character{59}}%
\TextAt{13}{261.707}{212.713}{X}%
\TextAt{21}{271.485}{210.92}{n}%
\TextAt{9}{277.121}{212.713}{\Character{41}\Character{41}}%
\TextAt{9}{289.519}{212.713}{\Character{61}}%
\TextAt{13}{301.972}{212.713}{F}%
\TextAt{21}{311.175}{217.649}{\Character{39}}%
\TextAt{9}{317.192}{212.713}{\Character{40}}%
\TextAt{9}{322.165}{212.713}{\Character{22}}%
\TextAt{13}{321.744}{212.713}{e}%
\TextAt{9}{327.17}{212.713}{\Character{40}}%
\TextAt{13}{331.722}{212.713}{X}%
\TextAt{17}{341.437}{210.92}{1}%
\TextAt{9}{346.17}{212.713}{\Character{41}}%
\TextAt{13}{350.722}{212.713}{\Character{59}}%
\TextAt{13}{355.946}{212.713}{\Character{58}}%
\TextAt{13}{361.183}{212.713}{\Character{58}}%
\TextAt{13}{366.407}{212.713}{\Character{58}}%
\TextAt{13}{371.644}{212.713}{\Character{59}}%
\TextAt{9}{377.379}{212.713}{\Character{22}}%
\TextAt{13}{376.943}{212.713}{e}%
\TextAt{9}{382.368}{212.713}{\Character{40}}%
\TextAt{13}{386.92}{212.713}{X}%
\TextAt{21}{396.636}{210.92}{n}%
\TextAt{9}{402.272}{212.713}{\Character{41}\Character{41}}%
\TextAt{13}{411.377}{212.713}{\Character{58}}%
\ErrataAt{79.19}{189.191}{455.81}{If $e,\bar e,L$ have coefficients in a complete subring $\mathfrak o\subseteq\mathfrak o_{\widehat K}$ with $\mathfrak o^\varphi=\mathfrak o$}%
\ErrataAt{79.19}{172.145}{455.81}{and $\mathfrak o\cap\pi\mathfrak o_{\widehat K}=\pi\mathfrak o$, then $F$ has coefficients in $\mathfrak o$ as well.}%
\ErrataAt{79.19}{137.167}{455.81}{\textbf{Proof.} Let $\mathfrak o$ be such a subring containing the coefficients of $e,\bar e,L$ (possibly $\mathfrak o_{\widehat K}$).}%
\ErrataAt{79.19}{120.121}{455.81}{Put $X=(X_1,\ldots,X_n)$ and $\bar e(X)=(\bar e(X_1),\ldots,\bar e(X_n))$. Let}%
\TextAt{13}{235.15}{84.609}{F}%
\TextAt{9}{244.354}{84.609}{\Character{40}}%
\TextAt{13}{248.906}{84.609}{X}%
\TextAt{9}{259.561}{84.609}{\Character{41}}%
\TextAt{9}{267.416}{84.609}{\Character{61}}%
\NativeMathAt{284.26}{99.553}{7.97}{\infty}%
\TextAt{20}{279.86}{95.967}{\Character{88}}%
\TextAt{21}{280.742}{70.861}{\Character{23}}%
\TextAt{17}{285.426}{70.861}{\Character{61}1}%
\TextAt{13}{299.121}{84.609}{E}%
\TextAt{21}{307.786}{82.816}{\Character{23}}%
\TextAt{9}{312.968}{84.609}{\Character{40}}%
\TextAt{13}{317.52}{84.609}{X}%
\TextAt{9}{328.176}{84.609}{\Character{41}}%
\NativeMathAt{336.049}{84.609}{11.955}{\in}%
\TextAt{14}{347.34}{84.609}{o}%
\TextAt{9}{353.184}{84.609}{\Character{91}\Character{91}}%
\TextAt{13}{359.687}{84.609}{X}%
\TextAt{9}{370.342}{84.609}{\Character{93}\Character{93}}%
\TextAt{8}{297.211}{43.033}{301}%
\Destination{522.36}{749.165}{nav.AMS.54}%
\Destination{78.23}{809.815}{nav.page.301}%
\Destination{79.36}{358.955}{nav.number.5.2.2}%
\Destination{79.226}{206.191}{errata.126}%
}

% Page 314 (body page 302)
\Page{%
\TextAt{2}{79.226}{796.943}{V.}%
\TextAt{2}{94.073}{796.943}{Local}%
\TextAt{2}{123.615}{796.943}{Class}%
\TextAt{2}{152.208}{796.943}{Field}%
\TextAt{2}{179.939}{796.943}{Theory}%
\Vector{q 0 G 0 g 0.299 w 7.226 18.729 m 460.769 18.729 l S Q}%
\TextAt{8}{79.226}{761.124}{be}%
\TextAt{8}{95.15}{761.124}{a}%
\TextAt{8}{104.906}{761.124}{power}%
\TextAt{8}{139.05}{761.124}{series,}%
\TextAt{13}{173.636}{761.124}{E}%
\TextAt{21}{182.301}{759.33}{\Character{23}}%
\TextAt{9}{187.484}{761.124}{\Character{40}}%
\TextAt{13}{192.036}{761.124}{X}%
\TextAt{9}{202.691}{761.124}{\Character{41}}%
\TextAt{8}{211.141}{761.124}{its}%
\TextAt{8}{227.459}{761.124}{homogeneous}%
\TextAt{8}{299.071}{761.124}{part}%
\TextAt{8}{324.44}{761.124}{of}%
\TextAt{8}{337.77}{761.124}{degree}%
\TextAt{13}{374.185}{761.124}{\Character{23}}%
\TextAt{8}{380.723}{761.124}{,}%
\TextAt{8}{387.872}{761.124}{and}%
\TextAt{8}{410.635}{761.124}{let}%
\TextAt{13}{254.119}{729.585}{F}%
\TextAt{21}{261.696}{727.792}{r}%
\TextAt{9}{266.251}{729.585}{\Character{40}}%
\TextAt{13}{270.803}{729.585}{X}%
\TextAt{9}{281.459}{729.585}{\Character{41}}%
\TextAt{9}{289.314}{729.585}{\Character{61}}%
\TextAt{21}{308.363}{744.529}{r}%
\TextAt{20}{301.757}{740.942}{\Character{88}}%
\TextAt{21}{302.639}{715.836}{\Character{23}}%
\TextAt{17}{307.323}{715.836}{\Character{61}1}%
\TextAt{13}{321.018}{729.585}{E}%
\TextAt{21}{329.684}{727.792}{\Character{23}}%
\TextAt{9}{334.866}{729.585}{\Character{40}}%
\TextAt{13}{339.418}{729.585}{X}%
\TextAt{9}{350.073}{729.585}{\Character{41}}%
\TextAt{13}{354.625}{729.585}{\Character{58}}%
\TextAt{8}{79.226}{696.696}{Clearly,}%
\TextAt{13}{122.145}{696.696}{F}%
\TextAt{9}{131.349}{696.696}{\Character{40}}%
\TextAt{13}{135.901}{696.696}{X}%
\TextAt{9}{146.556}{696.696}{\Character{41}}%
\TextAt{8}{155.006}{696.696}{is}%
\TextAt{8}{166.77}{696.696}{a}%
\TextAt{8}{176.525}{696.696}{solution}%
\TextAt{8}{220.819}{696.696}{of}%
\TextAt{8}{234.149}{696.696}{the}%
\TextAt{8}{254.306}{696.696}{above}%
\TextAt{8}{287.481}{696.696}{problem}%
\TextAt{8}{333.006}{696.696}{if}%
\TextAt{8}{343.73}{696.696}{and}%
\TextAt{8}{366.493}{696.696}{only}%
\TextAt{8}{392.185}{696.696}{if}%
\TextAt{13}{402.908}{696.696}{F}%
\TextAt{17}{410.486}{694.903}{1}%
\TextAt{9}{415.218}{696.696}{\Character{40}}%
\TextAt{13}{419.771}{696.696}{X}%
\TextAt{9}{430.426}{696.696}{\Character{41}}%
\TextAt{9}{438.281}{696.696}{\Character{61}}%
\TextAt{13}{450.725}{696.696}{L}%
\TextAt{9}{458.689}{696.696}{\Character{40}}%
\TextAt{13}{463.241}{696.696}{X}%
\TextAt{9}{473.897}{696.696}{\Character{41}}%
\TextAt{8}{482.346}{696.696}{and}%
\TextAt{13}{209.413}{671.667}{e}%
\TextAt{9}{214.839}{671.667}{\Character{40}}%
\TextAt{13}{219.391}{671.667}{F}%
\TextAt{21}{226.969}{669.874}{r}%
\TextAt{9}{231.523}{671.667}{\Character{40}}%
\TextAt{13}{236.076}{671.667}{X}%
\TextAt{9}{246.731}{671.667}{\Character{41}\Character{41}}%
\NativeMathAt{259.156}{671.667}{11.955}{\equiv}%
\TextAt{13}{271.776}{671.667}{F}%
\TextAt{21}{280.979}{676.603}{\Character{39}}%
\TextAt{21}{279.353}{668.712}{r}%
\TextAt{9}{286.996}{671.667}{\Character{40}}%
\TextAt{9}{291.969}{671.667}{\Character{22}}%
\TextAt{13}{291.548}{671.667}{e}%
\TextAt{9}{296.973}{671.667}{\Character{40}}%
\TextAt{13}{301.526}{671.667}{X}%
\TextAt{9}{312.181}{671.667}{\Character{41}\Character{41}}%
\TextAt{8}{324.606}{671.667}{mod}%
\TextAt{8}{350.035}{671.667}{deg}%
\TextAt{9}{367.606}{671.667}{\Character{40}}%
\TextAt{13}{372.159}{671.667}{r}%
\TextAt{9}{380.416}{671.667}{\Character{43}}%
\TextAt{9}{392.168}{671.667}{1\Character{41}}%
\TextAt{8}{517.802}{671.667}{(1)}%
\TextAt{8}{79.226}{646.639}{for}%
\TextAt{8}{98.067}{646.639}{every}%
\TextAt{13}{129.915}{646.639}{r}%
\NativeMathAt{140.466}{646.639}{11.955}{\geq}%
\TextAt{9}{154.715}{646.639}{1}%
\TextAt{8}{160.568}{646.639}{.}%
\TextAt{8}{171.884}{646.639}{We}%
\TextAt{8}{192.985}{646.639}{determine}%
\TextAt{8}{248.564}{646.639}{the}%
\TextAt{8}{269.676}{646.639}{polynomials}%
\TextAt{13}{335.74}{646.639}{E}%
\TextAt{21}{344.406}{644.845}{\Character{23}}%
\TextAt{9}{349.588}{646.639}{\Character{40}}%
\TextAt{13}{354.14}{646.639}{X}%
\TextAt{9}{364.795}{646.639}{\Character{41}}%
\TextAt{8}{374.201}{646.639}{inductively.}%
\TextAt{8}{440.798}{646.639}{For}%
\TextAt{13}{462.723}{646.639}{\Character{23}}%
\TextAt{9}{474.211}{646.639}{\Character{61}}%
\TextAt{9}{488.258}{646.639}{1}%
\TextAt{8}{498.973}{646.639}{we}%
\TextAt{8}{517.156}{646.639}{are}%
\TextAt{8}{79.226}{629.593}{forced}%
\TextAt{8}{114.726}{629.593}{to}%
\TextAt{8}{129.748}{629.593}{take}%
\TextAt{13}{155.827}{629.593}{E}%
\TextAt{17}{164.493}{627.8}{1}%
\TextAt{9}{169.225}{629.593}{\Character{40}}%
\TextAt{13}{173.777}{629.593}{X}%
\TextAt{9}{184.432}{629.593}{\Character{41}}%
\TextAt{9}{193.506}{629.593}{\Character{61}}%
\TextAt{13}{207.153}{629.593}{L}%
\TextAt{9}{215.118}{629.593}{\Character{40}}%
\TextAt{13}{219.67}{629.593}{X}%
\TextAt{9}{230.325}{629.593}{\Character{41}}%
\TextAt{8}{234.878}{629.593}{.}%
\TextAt{8}{245.458}{629.593}{Condition}%
\TextAt{8}{300.804}{629.593}{(1)}%
\TextAt{8}{320.38}{629.593}{is}%
\TextAt{8}{332.855}{629.593}{then}%
\TextAt{8}{360.226}{629.593}{satisfied}%
\TextAt{8}{405.936}{629.593}{for}%
\TextAt{13}{424.531}{629.593}{r}%
\TextAt{9}{434.664}{629.593}{\Character{61}}%
\TextAt{9}{448.293}{629.593}{1}%
\TextAt{8}{458.762}{629.593}{by}%
\TextAt{8}{475.731}{629.593}{hypothesis.}%
\TextAt{8}{79.226}{612.547}{Assume}%
\TextAt{8}{123.567}{612.547}{that}%
\TextAt{8}{149.915}{612.547}{the}%
\TextAt{13}{171.051}{612.547}{E}%
\TextAt{21}{179.716}{610.754}{\Character{23}}%
\TextAt{9}{184.898}{612.547}{\Character{40}}%
\TextAt{13}{189.451}{612.547}{X}%
\TextAt{9}{200.106}{612.547}{\Character{41}}%
\TextAt{8}{204.658}{612.547}{,}%
\TextAt{8}{213.032}{612.547}{for}%
\TextAt{13}{231.896}{612.547}{\Character{23}}%
\TextAt{9}{243.424}{612.547}{\Character{61}}%
\TextAt{9}{257.507}{612.547}{1}%
\TextAt{13}{263.371}{612.547}{\Character{59}}%
\TextAt{13}{268.595}{612.547}{\Character{58}}%
\TextAt{13}{273.832}{612.547}{\Character{58}}%
\TextAt{13}{279.068}{612.547}{\Character{58}}%
\TextAt{13}{284.293}{612.547}{\Character{59}}%
\TextAt{13}{289.529}{612.547}{r}%
\TextAt{8}{295.192}{612.547}{,}%
\TextAt{8}{303.566}{612.547}{have}%
\TextAt{8}{331.54}{612.547}{already}%
\TextAt{8}{373.824}{612.547}{been}%
\TextAt{8}{402.432}{612.547}{found,}%
\TextAt{8}{439.749}{612.547}{and}%
\TextAt{8}{463.491}{612.547}{that}%
\TextAt{8}{489.839}{612.547}{they}%
\TextAt{8}{517.156}{612.547}{are}%
\TextAt{8}{79.226}{595.501}{uniquely}%
\TextAt{8}{126.738}{595.501}{determined}%
\TextAt{8}{187.903}{595.501}{by}%
\TextAt{8}{204.199}{595.501}{condition}%
\TextAt{8}{255.621}{595.501}{(1).}%
\TextAt{8}{279.15}{595.501}{We}%
\TextAt{8}{299.333}{595.501}{then}%
\TextAt{8}{326.03}{595.501}{put}%
\TextAt{13}{347.528}{595.501}{F}%
\TextAt{21}{355.106}{593.708}{r}%
\TextAt{17}{359.162}{593.708}{\Character{43}1}%
\TextAt{9}{370.481}{595.501}{\Character{40}}%
\TextAt{13}{375.033}{595.501}{X}%
\TextAt{9}{385.688}{595.501}{\Character{41}}%
\TextAt{9}{393.614}{595.501}{\Character{61}}%
\TextAt{13}{406.116}{595.501}{F}%
\TextAt{21}{413.694}{593.708}{r}%
\TextAt{9}{418.249}{595.501}{\Character{40}}%
\TextAt{13}{422.801}{595.501}{X}%
\TextAt{9}{433.456}{595.501}{\Character{41}}%
\TextAt{9}{440.677}{595.501}{\Character{43}}%
\TextAt{13}{452.478}{595.501}{E}%
\TextAt{21}{461.143}{593.708}{r}%
\TextAt{17}{465.2}{593.708}{\Character{43}1}%
\TextAt{9}{476.519}{595.501}{\Character{40}}%
\TextAt{13}{481.071}{595.501}{X}%
\TextAt{9}{491.726}{595.501}{\Character{41}}%
\TextAt{8}{500.214}{595.501}{with}%
\TextAt{8}{526.911}{595.501}{a}%
\TextAt{8}{79.226}{578.455}{homogeneous}%
\TextAt{8}{150.309}{578.455}{polynomial}%
\TextAt{13}{210.274}{578.455}{E}%
\TextAt{21}{218.939}{576.662}{r}%
\TextAt{17}{222.996}{576.662}{\Character{43}1}%
\TextAt{9}{234.315}{578.455}{\Character{40}}%
\TextAt{13}{238.867}{578.455}{X}%
\TextAt{9}{249.522}{578.455}{\Character{41}}%
\NativeMathAt{257.395}{578.455}{11.955}{\in}%
\TextAt{14}{268.686}{578.455}{o}%
\TextAt{9}{274.53}{578.455}{\Character{91}}%
\TextAt{13}{277.782}{578.455}{X}%
\TextAt{9}{288.437}{578.455}{\Character{93}}%
\TextAt{8}{295.058}{578.455}{of}%
\TextAt{8}{307.859}{578.455}{degree}%
\TextAt{13}{343.746}{578.455}{r}%
\TextAt{9}{350.923}{578.455}{\Character{43}}%
\TextAt{9}{361.587}{578.455}{1}%
\TextAt{8}{370.825}{578.455}{which}%
\TextAt{8}{403.783}{578.455}{has}%
\TextAt{8}{424.128}{578.455}{yet}%
\TextAt{8}{443.111}{578.455}{to}%
\TextAt{8}{456.893}{578.455}{be}%
\TextAt{8}{472.288}{578.455}{determined.}%
\TextAt{8}{79.226}{561.409}{The}%
\TextAt{8}{103.28}{561.409}{congruences}%
\TextAt{13}{163.671}{536.38}{e}%
\TextAt{9}{169.097}{536.38}{\Character{40}}%
\TextAt{13}{173.649}{536.38}{F}%
\TextAt{21}{181.227}{534.587}{r}%
\TextAt{17}{185.283}{534.587}{\Character{43}1}%
\TextAt{9}{196.602}{536.38}{\Character{40}}%
\TextAt{13}{201.154}{536.38}{X}%
\TextAt{9}{211.81}{536.38}{\Character{41}\Character{41}}%
\NativeMathAt{224.235}{536.38}{11.955}{\equiv}%
\TextAt{13}{236.854}{536.38}{e}%
\TextAt{9}{242.28}{536.38}{\Character{40}}%
\TextAt{13}{246.832}{536.38}{F}%
\TextAt{21}{254.41}{534.587}{r}%
\TextAt{9}{258.965}{536.38}{\Character{40}}%
\TextAt{13}{263.517}{536.38}{X}%
\TextAt{9}{274.172}{536.38}{\Character{41}\Character{41}}%
\TextAt{9}{285.912}{536.38}{\Character{43}}%
\TextAt{13}{297.695}{536.38}{\Character{25}}%
\TextAt{13}{304.749}{536.38}{E}%
\TextAt{21}{313.429}{534.587}{r}%
\TextAt{17}{317.486}{534.587}{\Character{43}1}%
\TextAt{9}{328.805}{536.38}{\Character{40}}%
\TextAt{13}{333.357}{536.38}{X}%
\TextAt{9}{344.012}{536.38}{\Character{41}}%
\TextAt{8}{351.885}{536.38}{mod}%
\TextAt{8}{377.314}{536.38}{deg}%
\TextAt{9}{394.885}{536.38}{\Character{40}}%
\TextAt{13}{399.438}{536.38}{r}%
\TextAt{9}{407.695}{536.38}{\Character{43}}%
\TextAt{9}{419.447}{536.38}{2\Character{41}}%
\TextAt{13}{429.861}{536.38}{\Character{59}}%
\TextAt{13}{163.671}{516.346}{F}%
\TextAt{21}{172.875}{522.013}{\Character{39}}%
\TextAt{21}{171.249}{513.414}{r}%
\TextAt{17}{175.306}{513.414}{\Character{43}1}%
\TextAt{9}{186.624}{516.346}{\Character{40}}%
\TextAt{9}{191.597}{516.346}{\Character{22}}%
\TextAt{13}{191.177}{516.346}{e}%
\TextAt{9}{196.602}{516.346}{\Character{40}}%
\TextAt{13}{201.154}{516.346}{X}%
\TextAt{9}{211.81}{516.346}{\Character{41}\Character{41}}%
\NativeMathAt{224.235}{516.346}{11.955}{\equiv}%
\TextAt{13}{236.854}{516.346}{F}%
\TextAt{21}{246.058}{521.282}{\Character{39}}%
\TextAt{21}{244.432}{513.39}{r}%
\TextAt{9}{252.074}{516.346}{\Character{40}}%
\TextAt{9}{257.047}{516.346}{\Character{22}}%
\TextAt{13}{256.627}{516.346}{e}%
\TextAt{9}{262.052}{516.346}{\Character{40}}%
\TextAt{13}{266.605}{516.346}{X}%
\TextAt{9}{277.26}{516.346}{\Character{41}\Character{41}}%
\TextAt{9}{289}{516.346}{\Character{43}}%
\TextAt{9}{301.362}{516.346}{\Character{22}}%
\TextAt{13}{300.782}{516.346}{\Character{25}}%
\TextAt{21}{307.852}{521.282}{r}%
\TextAt{17}{311.908}{521.282}{\Character{43}1}%
\TextAt{13}{323.227}{516.346}{E}%
\TextAt{21}{332.593}{522.013}{\Character{39}}%
\TextAt{21}{331.892}{513.414}{r}%
\TextAt{17}{335.949}{513.414}{\Character{43}1}%
\TextAt{9}{347.268}{516.346}{\Character{40}}%
\TextAt{13}{351.82}{516.346}{X}%
\TextAt{9}{362.475}{516.346}{\Character{41}}%
\TextAt{8}{370.348}{516.346}{mod}%
\TextAt{8}{395.777}{516.346}{deg}%
\TextAt{9}{413.348}{516.346}{\Character{40}}%
\TextAt{13}{417.901}{516.346}{r}%
\TextAt{9}{426.158}{516.346}{\Character{43}}%
\TextAt{9}{437.91}{516.346}{2\Character{41}}%
\TextAt{8}{79.226}{491.317}{show}%
\TextAt{8}{108.229}{491.317}{that}%
\TextAt{13}{133.598}{491.317}{E}%
\TextAt{21}{142.264}{489.524}{r}%
\TextAt{17}{146.32}{489.524}{\Character{43}1}%
\TextAt{9}{157.639}{491.317}{\Character{40}}%
\TextAt{13}{162.191}{491.317}{X}%
\TextAt{9}{172.846}{491.317}{\Character{41}}%
\TextAt{8}{181.296}{491.317}{has}%
\TextAt{8}{202.17}{491.317}{to}%
\TextAt{8}{216.48}{491.317}{satisfy}%
\TextAt{8}{253.027}{491.317}{the}%
\TextAt{8}{273.184}{491.317}{congruence}%
\TextAt{13}{159.196}{466.289}{G}%
\TextAt{21}{168.43}{464.495}{r}%
\TextAt{17}{172.486}{464.495}{\Character{43}1}%
\TextAt{9}{183.805}{466.289}{\Character{40}}%
\TextAt{13}{188.357}{466.289}{X}%
\TextAt{9}{199.012}{466.289}{\Character{41}}%
\TextAt{9}{206.209}{466.289}{\Character{43}}%
\TextAt{13}{217.983}{466.289}{\Character{25}}%
\TextAt{13}{225.037}{466.289}{E}%
\TextAt{21}{233.717}{464.495}{r}%
\TextAt{17}{237.774}{464.495}{\Character{43}1}%
\TextAt{9}{249.093}{466.289}{\Character{40}}%
\TextAt{13}{253.645}{466.289}{X}%
\TextAt{9}{264.3}{466.289}{\Character{41}}%
\TextAt{16}{271.509}{466.289}{\Character{0}}%
\TextAt{9}{284.072}{466.289}{\Character{22}}%
\TextAt{13}{283.464}{466.289}{\Character{25}}%
\TextAt{21}{290.534}{471.225}{r}%
\TextAt{17}{294.59}{471.225}{\Character{43}1}%
\TextAt{13}{305.909}{466.289}{E}%
\TextAt{21}{315.275}{471.956}{\Character{39}}%
\TextAt{21}{314.574}{463.357}{r}%
\TextAt{17}{318.631}{463.357}{\Character{43}1}%
\TextAt{9}{329.95}{466.289}{\Character{40}}%
\TextAt{13}{334.502}{466.289}{X}%
\TextAt{9}{345.157}{466.289}{\Character{41}}%
\NativeMathAt{353.03}{466.289}{11.955}{\equiv}%
\TextAt{9}{365.65}{466.289}{0}%
\TextAt{8}{374.824}{466.289}{mod}%
\TextAt{8}{400.253}{466.289}{deg}%
\TextAt{9}{417.824}{466.289}{\Character{40}}%
\TextAt{13}{422.376}{466.289}{r}%
\TextAt{9}{430.633}{466.289}{\Character{43}}%
\TextAt{9}{442.385}{466.289}{2\Character{41}}%
\TextAt{8}{517.802}{466.289}{(2)}%
\TextAt{8}{79.226}{441.26}{with}%
\TextAt{13}{106.364}{441.26}{G}%
\TextAt{21}{115.598}{439.467}{r}%
\TextAt{17}{119.654}{439.467}{\Character{43}1}%
\TextAt{9}{130.973}{441.26}{\Character{40}}%
\TextAt{13}{135.526}{441.26}{X}%
\TextAt{9}{146.181}{441.26}{\Character{41}}%
\TextAt{9}{154.849}{441.26}{\Character{61}}%
\TextAt{13}{168.109}{441.26}{e}%
\TextAt{9}{173.534}{441.26}{\Character{40}}%
\TextAt{13}{178.086}{441.26}{F}%
\TextAt{21}{185.664}{439.467}{r}%
\TextAt{9}{190.219}{441.26}{\Character{40}}%
\TextAt{13}{194.771}{441.26}{X}%
\TextAt{9}{205.426}{441.26}{\Character{41}\Character{41}}%
\TextAt{16}{217.514}{441.26}{\Character{0}}%
\TextAt{13}{229.795}{441.26}{F}%
\TextAt{21}{238.998}{445.599}{\Character{39}}%
\TextAt{21}{237.372}{438.305}{r}%
\TextAt{9}{245.015}{441.26}{\Character{40}}%
\TextAt{9}{249.988}{441.26}{\Character{22}}%
\TextAt{13}{249.567}{441.26}{e}%
\TextAt{9}{254.992}{441.26}{\Character{40}}%
\TextAt{13}{259.545}{441.26}{X}%
\TextAt{9}{270.2}{441.26}{\Character{41}\Character{41}}%
\NativeMathAt{283.44}{441.26}{11.955}{\in}%
\TextAt{14}{295.546}{441.26}{o}%
\TextAt{9}{301.39}{441.26}{\Character{91}\Character{91}}%
\TextAt{13}{307.893}{441.26}{X}%
\TextAt{9}{318.548}{441.26}{\Character{93}\Character{93}}%
\TextAt{8}{325.052}{441.26}{.}%
\TextAt{8}{334.934}{441.26}{We}%
\TextAt{8}{355.557}{441.26}{have}%
\TextAt{13}{383.03}{441.26}{G}%
\TextAt{21}{392.263}{439.467}{r}%
\TextAt{17}{396.32}{439.467}{\Character{43}1}%
\TextAt{9}{407.639}{441.26}{\Character{40}}%
\TextAt{13}{412.191}{441.26}{X}%
\TextAt{9}{422.846}{441.26}{\Character{41}}%
\NativeMathAt{431.534}{441.26}{11.955}{\equiv}%
\TextAt{9}{444.968}{441.26}{0}%
\TextAt{8}{454.142}{441.26}{mod}%
\TextAt{8}{479.571}{441.26}{deg}%
\TextAt{9}{497.142}{441.26}{\Character{40}}%
\TextAt{13}{501.694}{441.26}{r}%
\TextAt{9}{510.277}{441.26}{\Character{43}}%
\TextAt{9}{522.352}{441.26}{1\Character{41}}%
\TextAt{8}{79.226}{424.214}{and}%
\TextAt{13}{201.726}{407.168}{G}%
\TextAt{21}{210.96}{405.375}{r}%
\TextAt{17}{215.016}{405.375}{\Character{43}1}%
\TextAt{9}{226.335}{407.168}{\Character{40}}%
\TextAt{13}{230.887}{407.168}{X}%
\TextAt{9}{241.542}{407.168}{\Character{41}}%
\NativeMathAt{249.416}{407.168}{11.955}{\equiv}%
\TextAt{13}{262.035}{407.168}{F}%
\TextAt{21}{269.613}{405.375}{r}%
\TextAt{9}{274.167}{407.168}{\Character{40}}%
\TextAt{13}{278.72}{407.168}{X}%
\TextAt{9}{289.375}{407.168}{\Character{41}}%
\TextAt{21}{293.927}{412.104}{q}%
\TextAt{16}{301.149}{407.168}{\Character{0}}%
\TextAt{13}{313.104}{407.168}{F}%
\TextAt{21}{322.308}{412.104}{\Character{39}}%
\TextAt{21}{320.682}{404.213}{r}%
\TextAt{9}{328.324}{407.168}{\Character{40}}%
\TextAt{13}{332.876}{407.168}{X}%
\TextAt{21}{343.531}{412.104}{q}%
\TextAt{9}{348.096}{407.168}{\Character{41}}%
\NativeMathAt{355.969}{407.168}{11.955}{\equiv}%
\TextAt{9}{368.589}{407.168}{0}%
\TextAt{8}{377.763}{407.168}{mod}%
\TextAt{13}{403.2}{407.168}{\Character{25}}%
\TextAt{8}{517.802}{407.168}{(3)}%
\TextAt{8}{79.226}{385.349}{because}%
\TextAt{13}{122.528}{385.349}{e}%
\TextAt{9}{127.953}{385.349}{\Character{40}}%
\TextAt{13}{132.506}{385.349}{X}%
\TextAt{9}{143.161}{385.349}{\Character{41}}%
\NativeMathAt{151.034}{385.349}{11.955}{\equiv}%
\TextAt{9}{164.09}{385.349}{\Character{22}}%
\TextAt{13}{163.653}{385.349}{e}%
\TextAt{9}{169.079}{385.349}{\Character{40}}%
\TextAt{13}{173.631}{385.349}{X}%
\TextAt{9}{184.286}{385.349}{\Character{41}}%
\NativeMathAt{192.159}{385.349}{11.955}{\equiv}%
\TextAt{13}{204.778}{385.349}{X}%
\TextAt{21}{215.434}{389.687}{q}%
\TextAt{8}{223.319}{385.349}{mod}%
\TextAt{13}{248.757}{385.349}{\Character{25}}%
\TextAt{8}{259.724}{385.349}{and}%
\TextAt{13}{282.487}{385.349}{\Character{11}}%
\TextAt{21}{290.008}{389.687}{\Character{39}}%
\NativeMathAt{299.345}{385.349}{11.955}{\equiv}%
\TextAt{13}{311.965}{385.349}{\Character{11}}%
\TextAt{21}{319.487}{389.687}{q}%
\TextAt{8}{327.372}{385.349}{mod}%
\TextAt{13}{352.81}{385.349}{\Character{25}}%
\TextAt{8}{363.777}{385.349}{for}%
\TextAt{13}{381.662}{385.349}{\Character{11}}%
\NativeMathAt{392.504}{385.349}{11.955}{\in}%
\TextAt{14}{403.795}{385.349}{o}%
\TextAt{8}{409.639}{385.349}{.}%
\TextAt{8}{97.159}{365.314}{Now}%
\TextAt{8}{123.091}{365.314}{let}%
\TextAt{13}{139.267}{365.314}{X}%
\TextAt{21}{149.923}{369.653}{i}%
\TextAt{9}{156.625}{365.314}{\Character{61}}%
\TextAt{13}{169.05}{365.314}{X}%
\TextAt{21}{179.705}{370.539}{i}%
\TextAt{24}{182.588}{369.432}{1}%
\TextAt{17}{178.765}{362.383}{1}%
\TextAt{16}{189.23}{365.314}{\Character{1}}%
\TextAt{16}{194.526}{365.314}{\Character{1}}%
\TextAt{16}{199.834}{365.314}{\Character{1}}%
\TextAt{13}{205.17}{365.314}{X}%
\TextAt{21}{215.825}{369.653}{i}%
\TextAt{22}{218.709}{368.656}{n}%
\TextAt{21}{214.885}{362.359}{n}%
\TextAt{8}{227.475}{365.314}{be}%
\TextAt{8}{242.671}{365.314}{a}%
\TextAt{8}{251.699}{365.314}{monomial}%
\TextAt{8}{304.96}{365.314}{of}%
\TextAt{8}{317.562}{365.314}{degree}%
\TextAt{13}{353.249}{365.314}{r}%
\TextAt{9}{360.018}{365.314}{\Character{43}}%
\TextAt{9}{370.276}{365.314}{1}%
\TextAt{8}{379.313}{365.314}{in}%
\TextAt{14}{392.237}{365.314}{o}%
\TextAt{9}{398.082}{365.314}{\Character{91}}%
\TextAt{13}{401.333}{365.314}{X}%
\TextAt{9}{411.988}{365.314}{\Character{93}}%
\TextAt{8}{415.24}{365.314}{.}%
\TextAt{8}{423.446}{365.314}{By}%
\TextAt{8}{441.081}{365.314}{(3),}%
\TextAt{8}{462.615}{365.314}{the}%
\TextAt{8}{482.044}{365.314}{coefficient}%
\TextAt{8}{79.226}{348.268}{of}%
\TextAt{13}{93.892}{348.268}{X}%
\TextAt{21}{104.547}{352.607}{i}%
\TextAt{8}{113.161}{348.268}{in}%
\TextAt{13}{128.15}{348.268}{G}%
\TextAt{21}{137.383}{346.475}{r}%
\TextAt{17}{141.44}{346.475}{\Character{43}1}%
\TextAt{8}{157.992}{348.268}{is}%
\TextAt{8}{171.091}{348.268}{of}%
\TextAt{8}{185.757}{348.268}{the}%
\TextAt{8}{207.249}{348.268}{form}%
\TextAt{16}{236.225}{348.268}{\Character{0}}%
\TextAt{13}{245.523}{348.268}{\Character{25}}%
\TextAt{13}{252.577}{348.268}{\Character{12}}%
\TextAt{8}{259.864}{348.268}{,}%
\TextAt{8}{268.683}{348.268}{with}%
\TextAt{13}{296.678}{348.268}{\Character{12}}%
\NativeMathAt{309.546}{348.268}{11.955}{\in}%
\TextAt{14}{323.113}{348.268}{o}%
\TextAt{8}{328.958}{348.268}{.}%
\TextAt{8}{341.413}{348.268}{Let}%
\TextAt{13}{363.718}{348.268}{\Character{11}}%
\TextAt{8}{376.472}{348.268}{be}%
\TextAt{8}{393.732}{348.268}{the}%
\TextAt{8}{415.224}{348.268}{coefficient}%
\TextAt{8}{471.183}{348.268}{of}%
\TextAt{8}{485.849}{348.268}{the}%
\TextAt{8}{507.341}{348.268}{same}%
\TextAt{8}{79.226}{331.222}{monomial}%
\TextAt{13}{133.749}{331.222}{X}%
\TextAt{21}{144.404}{335.561}{i}%
\TextAt{8}{152.216}{331.222}{in}%
\TextAt{13}{166.402}{331.222}{E}%
\TextAt{21}{175.067}{329.429}{r}%
\TextAt{17}{179.124}{329.429}{\Character{43}1}%
\TextAt{8}{190.443}{331.222}{.}%
\TextAt{8}{200.491}{331.222}{Then}%
\TextAt{13}{231.581}{331.222}{\Character{25}}%
\TextAt{13}{238.635}{331.222}{\Character{11}}%
\TextAt{16}{249.192}{331.222}{\Character{0}}%
\TextAt{9}{262.119}{331.222}{\Character{22}}%
\TextAt{13}{261.511}{331.222}{\Character{25}}%
\TextAt{21}{268.58}{335.561}{r}%
\TextAt{17}{272.637}{335.561}{\Character{43}1}%
\TextAt{13}{283.956}{331.222}{\Character{11}}%
\TextAt{21}{291.477}{335.561}{\Character{39}}%
\TextAt{8}{301.924}{331.222}{is}%
\TextAt{8}{314.221}{331.222}{the}%
\TextAt{8}{334.911}{331.222}{coefficient}%
\TextAt{8}{390.067}{331.222}{of}%
\TextAt{13}{403.931}{331.222}{X}%
\TextAt{21}{414.586}{335.561}{i}%
\TextAt{8}{422.398}{331.222}{in}%
\TextAt{13}{436.584}{331.222}{\Character{25}}%
\TextAt{13}{443.638}{331.222}{E}%
\TextAt{21}{452.318}{329.429}{r}%
\TextAt{17}{456.375}{329.429}{\Character{43}1}%
\TextAt{16}{470.714}{331.222}{\Character{0}}%
\TextAt{9}{483.64}{331.222}{\Character{22}}%
\TextAt{13}{483.032}{331.222}{\Character{25}}%
\TextAt{21}{490.102}{335.561}{r}%
\TextAt{17}{494.158}{335.561}{\Character{43}1}%
\TextAt{13}{505.477}{331.222}{E}%
\TextAt{21}{514.843}{336.89}{\Character{39}}%
\TextAt{21}{514.142}{328.291}{r}%
\TextAt{17}{518.199}{328.291}{\Character{43}1}%
\TextAt{8}{529.518}{331.222}{.}%
\TextAt{8}{79.226}{314.176}{Since}%
\TextAt{13}{109.438}{314.176}{G}%
\TextAt{21}{118.672}{312.383}{r}%
\TextAt{17}{122.729}{312.383}{\Character{43}1}%
\TextAt{9}{134.047}{314.176}{\Character{40}}%
\TextAt{13}{138.6}{314.176}{X}%
\TextAt{9}{149.255}{314.176}{\Character{41}}%
\NativeMathAt{157.128}{314.176}{11.955}{\equiv}%
\TextAt{9}{169.747}{314.176}{0}%
\TextAt{8}{178.921}{314.176}{mod}%
\TextAt{8}{204.35}{314.176}{deg}%
\TextAt{9}{221.921}{314.176}{\Character{40}}%
\TextAt{13}{226.474}{314.176}{r}%
\TextAt{9}{234.025}{314.176}{\Character{43}}%
\TextAt{9}{245.072}{314.176}{1\Character{41}}%
\TextAt{8}{255.486}{314.176}{,}%
\TextAt{8}{262.359}{314.176}{(2)}%
\TextAt{8}{280.879}{314.176}{holds}%
\TextAt{8}{311.163}{314.176}{if}%
\TextAt{8}{321.542}{314.176}{and}%
\TextAt{8}{343.959}{314.176}{only}%
\TextAt{8}{369.306}{314.176}{if}%
\TextAt{8}{379.684}{314.176}{the}%
\TextAt{8}{399.496}{314.176}{coefficient}%
\TextAt{13}{453.774}{314.176}{\Character{11}}%
\TextAt{8}{464.848}{314.176}{of}%
\TextAt{13}{477.832}{314.176}{X}%
\TextAt{21}{488.488}{318.515}{i}%
\TextAt{8}{495.421}{314.176}{in}%
\TextAt{13}{508.729}{314.176}{E}%
\TextAt{21}{517.394}{312.383}{r}%
\TextAt{17}{521.451}{312.383}{\Character{43}1}%
\TextAt{8}{79.226}{297.13}{satisfies}%
\TextAt{8}{122.336}{297.13}{the}%
\TextAt{8}{142.493}{297.13}{equation}%
\TextAt{16}{243.577}{280.084}{\Character{0}}%
\TextAt{13}{252.875}{280.084}{\Character{25}}%
\TextAt{13}{259.929}{280.084}{\Character{12}}%
\TextAt{9}{269.873}{280.084}{\Character{43}}%
\TextAt{13}{281.634}{280.084}{\Character{25}}%
\TextAt{13}{288.688}{280.084}{\Character{11}}%
\TextAt{16}{298.882}{280.084}{\Character{0}}%
\TextAt{9}{311.445}{280.084}{\Character{22}}%
\TextAt{13}{310.837}{280.084}{\Character{25}}%
\TextAt{21}{317.906}{285.021}{r}%
\TextAt{17}{321.963}{285.021}{\Character{43}1}%
\TextAt{13}{333.281}{280.084}{\Character{11}}%
\TextAt{21}{340.803}{285.021}{\Character{39}}%
\TextAt{9}{350.14}{280.084}{\Character{61}}%
\TextAt{9}{362.55}{280.084}{0}%
\TextAt{8}{517.802}{280.084}{(4)}%
\TextAt{8}{79.226}{258.265}{for}%
\TextAt{8}{98.141}{258.265}{every}%
\TextAt{8}{130.063}{258.265}{monomial}%
\TextAt{13}{185.083}{258.265}{X}%
\TextAt{21}{195.738}{262.604}{i}%
\TextAt{8}{204.047}{258.265}{of}%
\TextAt{8}{218.407}{258.265}{degree}%
\TextAt{13}{255.852}{258.265}{r}%
\TextAt{9}{264.811}{258.265}{\Character{43}}%
\TextAt{9}{277.256}{258.265}{1}%
\TextAt{8}{283.128}{258.265}{.}%
\TextAt{8}{294.666}{258.265}{This}%
\TextAt{8}{322.416}{258.265}{equation}%
\TextAt{8}{371.254}{258.265}{has}%
\TextAt{8}{393.158}{258.265}{a}%
\TextAt{8}{403.944}{258.265}{unique}%
\TextAt{8}{443.015}{258.265}{solution}%
\TextAt{13}{488.339}{258.265}{\Character{11}}%
\TextAt{8}{500.788}{258.265}{in}%
\TextAt{14}{515.471}{258.265}{o}%
\TextAt{30}{523.285}{256.627}{\Character{98}}%
\TextAt{21}{521.315}{254.612}{K}%
\TextAt{8}{529.518}{258.265}{,}%
\TextAt{8}{79.226}{241.219}{which}%
\TextAt{8}{112.713}{241.219}{actually}%
\TextAt{8}{157.269}{241.219}{belongs}%
\TextAt{8}{199.28}{241.219}{to}%
\TextAt{14}{213.59}{241.219}{o}%
\TextAt{8}{219.435}{241.219}{.}%
\TextAt{8}{227.883}{241.219}{For}%
\TextAt{8}{248.852}{241.219}{if}%
\TextAt{8}{259.576}{241.219}{we}%
\TextAt{8}{276.803}{241.219}{put}%
\TextAt{13}{298.263}{241.219}{\Character{13}}%
\TextAt{9}{308.306}{241.219}{\Character{61}}%
\TextAt{13}{320.731}{241.219}{\Character{25}}%
\TextAt{3}{327.801}{245.558}{\Character{0}}%
\TextAt{17}{334.387}{245.558}{1}%
\TextAt{9}{339.728}{241.219}{\Character{22}}%
\TextAt{13}{339.12}{241.219}{\Character{25}}%
\TextAt{21}{346.189}{245.558}{r}%
\TextAt{17}{350.245}{245.558}{\Character{43}1}%
\TextAt{8}{361.564}{241.219}{,}%
\TextAt{8}{368.713}{241.219}{we}%
\TextAt{8}{385.941}{241.219}{obtain}%
\TextAt{8}{422.368}{241.219}{the}%
\TextAt{8}{442.525}{241.219}{equation}%
\TextAt{13}{271.666}{216.191}{\Character{11}}%
\TextAt{16}{281.845}{216.191}{\Character{0}}%
\TextAt{13}{293.8}{216.191}{\Character{13}}%
\TextAt{13}{300.507}{216.191}{\Character{11}}%
\TextAt{21}{308.044}{221.127}{\Character{39}}%
\TextAt{9}{317.381}{216.191}{\Character{61}}%
\TextAt{13}{329.806}{216.191}{\Character{12}}%
\TextAt{13}{337.075}{216.191}{\Character{59}}%
\TextAt{8}{79.226}{191.162}{which}%
\TextAt{8}{112.713}{191.162}{is}%
\TextAt{8}{124.476}{191.162}{clearly}%
\TextAt{8}{161.872}{191.162}{solved}%
\TextAt{8}{197.056}{191.162}{by}%
\TextAt{8}{213.315}{191.162}{the}%
\TextAt{8}{233.472}{191.162}{series}%
\TextAt{13}{221.395}{166.134}{\Character{11}}%
\TextAt{9}{232.238}{166.134}{\Character{61}}%
\TextAt{13}{244.663}{166.134}{\Character{12}}%
\TextAt{9}{254.591}{166.134}{\Character{43}}%
\TextAt{13}{266.353}{166.134}{\Character{13}}%
\TextAt{13}{273.06}{166.134}{\Character{12}}%
\TextAt{21}{280.346}{171.07}{\Character{39}}%
\TextAt{9}{289.019}{166.134}{\Character{43}}%
\TextAt{13}{300.781}{166.134}{\Character{13}}%
\TextAt{17}{307.503}{171.07}{1\Character{43}}%
\TextAt{21}{318.323}{171.07}{\Character{39}}%
\TextAt{13}{324.34}{166.134}{\Character{12}}%
\TextAt{21}{331.611}{171.07}{\Character{39}}%
\TextAt{24}{337.129}{173.883}{2}%
\TextAt{9}{344.435}{166.134}{\Character{43}}%
\TextAt{16}{356.197}{166.134}{\Character{1}}%
\TextAt{16}{361.493}{166.134}{\Character{1}}%
\TextAt{16}{366.801}{166.134}{\Character{1}}%
\NativeMathAt{373.436}{166.134}{11.955}{\in}%
\TextAt{14}{384.756}{166.134}{o}%
\TextAt{8}{79.226}{141.105}{(the}%
\TextAt{8}{105.844}{141.105}{series}%
\TextAt{8}{139.085}{141.105}{converges}%
\TextAt{8}{193.415}{141.105}{because}%
\TextAt{13}{238.622}{141.105}{v}%
\TextAt{30}{246.252}{139.467}{\Character{98}}%
\TextAt{21}{244.282}{137.452}{K}%
\TextAt{9}{252.484}{141.105}{\Character{40}}%
\TextAt{13}{257.037}{141.105}{\Character{13}}%
\TextAt{9}{263.759}{141.105}{\Character{41}}%
\NativeMathAt{274.881}{141.105}{11.955}{\geq}%
\TextAt{9}{290.749}{141.105}{1}%
\TextAt{8}{296.602}{141.105}{).}%
\TextAt{8}{315.324}{141.105}{If}%
\TextAt{13}{328.923}{141.105}{\Character{11}}%
\TextAt{3}{336.444}{145.444}{\Character{48}}%
\TextAt{8}{345.043}{141.105}{is}%
\TextAt{8}{358.714}{141.105}{another}%
\TextAt{8}{403.551}{141.105}{solution,}%
\TextAt{8}{453.48}{141.105}{then}%
\TextAt{13}{482.046}{141.105}{\Character{11}}%
\TextAt{16}{493.524}{141.105}{\Character{0}}%
\TextAt{13}{506.779}{141.105}{\Character{11}}%
\TextAt{3}{514.3}{145.444}{\Character{48}}%
\TextAt{9}{523.665}{141.105}{\Character{61}}%
\TextAt{13}{79.226}{124.059}{\Character{13}}%
\TextAt{9}{85.948}{124.059}{\Character{40}}%
\TextAt{13}{90.501}{124.059}{\Character{11}}%
\TextAt{21}{98.022}{128.398}{\Character{39}}%
\TextAt{16}{107.714}{124.059}{\Character{0}}%
\TextAt{13}{120.688}{124.059}{\Character{11}}%
\TextAt{3}{128.21}{128.398}{\Character{48}}%
\TextAt{21}{130.507}{128.398}{\Character{39}}%
\TextAt{9}{136.523}{124.059}{\Character{41}}%
\TextAt{8}{141.075}{124.059}{,}%
\TextAt{8}{150.093}{124.059}{hence}%
\TextAt{13}{184.093}{124.059}{v}%
\TextAt{30}{191.723}{122.421}{\Character{98}}%
\TextAt{21}{189.752}{120.406}{K}%
\TextAt{9}{197.955}{124.059}{\Character{40}}%
\TextAt{13}{202.507}{124.059}{\Character{11}}%
\TextAt{16}{213.704}{124.059}{\Character{0}}%
\TextAt{13}{226.678}{124.059}{\Character{11}}%
\TextAt{3}{234.2}{128.398}{\Character{48}}%
\TextAt{9}{236.995}{124.059}{\Character{41}}%
\TextAt{9}{247.396}{124.059}{\Character{61}}%
\TextAt{13}{262.387}{124.059}{v}%
\TextAt{30}{270.017}{122.421}{\Character{98}}%
\TextAt{21}{268.046}{120.406}{K}%
\TextAt{9}{276.249}{124.059}{\Character{40}}%
\TextAt{13}{280.801}{124.059}{\Character{13}}%
\TextAt{9}{287.524}{124.059}{\Character{41}}%
\TextAt{9}{295.737}{124.059}{\Character{43}}%
\TextAt{13}{308.531}{124.059}{v}%
\TextAt{30}{316.161}{122.421}{\Character{98}}%
\TextAt{21}{314.19}{120.406}{K}%
\TextAt{9}{322.393}{124.059}{\Character{40}\Character{40}}%
\TextAt{13}{331.498}{124.059}{\Character{11}}%
\TextAt{16}{342.695}{124.059}{\Character{0}}%
\TextAt{13}{355.669}{124.059}{\Character{11}}%
\TextAt{3}{363.19}{128.398}{\Character{48}}%
\TextAt{9}{365.986}{124.059}{\Character{41}}%
\TextAt{21}{370.538}{128.398}{\Character{39}}%
\TextAt{9}{376.554}{124.059}{\Character{41}}%
\TextAt{9}{386.955}{124.059}{\Character{61}}%
\TextAt{13}{401.946}{124.059}{v}%
\TextAt{30}{409.576}{122.421}{\Character{98}}%
\TextAt{21}{407.605}{120.406}{K}%
\TextAt{9}{415.808}{124.059}{\Character{40}}%
\TextAt{13}{420.36}{124.059}{\Character{13}}%
\TextAt{9}{427.083}{124.059}{\Character{41}}%
\TextAt{9}{435.296}{124.059}{\Character{43}}%
\TextAt{13}{448.09}{124.059}{v}%
\TextAt{30}{455.72}{122.421}{\Character{98}}%
\TextAt{21}{453.75}{120.406}{K}%
\TextAt{9}{461.952}{124.059}{\Character{40}}%
\TextAt{13}{466.505}{124.059}{\Character{11}}%
\TextAt{16}{477.702}{124.059}{\Character{0}}%
\TextAt{13}{490.676}{124.059}{\Character{11}}%
\TextAt{3}{498.197}{128.398}{\Character{48}}%
\TextAt{9}{500.992}{124.059}{\Character{41}}%
\TextAt{8}{505.545}{124.059}{,}%
\TextAt{8}{514.562}{124.059}{i.e.,}%
\TextAt{13}{79.226}{107.013}{v}%
\TextAt{30}{86.856}{105.375}{\Character{98}}%
\TextAt{21}{84.885}{103.36}{K}%
\TextAt{9}{93.088}{107.013}{\Character{40}}%
\TextAt{13}{97.64}{107.013}{\Character{11}}%
\TextAt{16}{108.087}{107.013}{\Character{0}}%
\TextAt{13}{120.309}{107.013}{\Character{11}}%
\TextAt{3}{127.831}{111.352}{\Character{48}}%
\TextAt{9}{130.626}{107.013}{\Character{41}}%
\TextAt{9}{139.15}{107.013}{\Character{61}}%
\NativeMathAt{152.263}{107.013}{11.955}{\infty}%
\TextAt{8}{168.509}{107.013}{because}%
\TextAt{13}{212.203}{107.013}{v}%
\TextAt{30}{219.833}{105.375}{\Character{98}}%
\TextAt{21}{217.862}{103.36}{K}%
\TextAt{9}{226.065}{107.013}{\Character{40}}%
\TextAt{13}{230.617}{107.013}{\Character{13}}%
\TextAt{9}{237.34}{107.013}{\Character{41}}%
\NativeMathAt{245.882}{107.013}{11.955}{\geq}%
\TextAt{9}{259.171}{107.013}{1}%
\TextAt{8}{265.024}{107.013}{,}%
\TextAt{8}{272.664}{107.013}{and}%
\TextAt{8}{295.819}{107.013}{therefore}%
\TextAt{13}{345.312}{107.013}{\Character{11}}%
\TextAt{9}{356.824}{107.013}{\Character{61}}%
\TextAt{13}{369.919}{107.013}{\Character{11}}%
\TextAt{3}{377.44}{111.352}{\Character{48}}%
\TextAt{8}{380.235}{107.013}{.}%
\TextAt{8}{389.862}{107.013}{As}%
\TextAt{8}{407.542}{107.013}{a}%
\TextAt{8}{417.69}{107.013}{consequence,}%
\TextAt{8}{487.497}{107.013}{for}%
\TextAt{8}{505.775}{107.013}{every}%
\TextAt{8}{79.226}{89.967}{monomial}%
\TextAt{13}{133.255}{89.967}{X}%
\TextAt{21}{143.91}{94.306}{i}%
\TextAt{8}{151.228}{89.967}{of}%
\TextAt{8}{164.598}{89.967}{degree}%
\TextAt{13}{201.053}{89.967}{r}%
\TextAt{9}{209.337}{89.967}{\Character{43}}%
\TextAt{9}{221.113}{89.967}{1}%
\TextAt{8}{226.978}{89.967}{,}%
\TextAt{8}{234.177}{89.967}{equation}%
\TextAt{8}{282.025}{89.967}{(4)}%
\TextAt{8}{300.93}{89.967}{has}%
\TextAt{8}{321.843}{89.967}{a}%
\TextAt{8}{331.638}{89.967}{unique}%
\TextAt{8}{369.719}{89.967}{solution}%
\TextAt{13}{414.052}{89.967}{\Character{11}}%
\TextAt{8}{425.511}{89.967}{in}%
\TextAt{14}{439.203}{89.967}{o}%
\TextAt{8}{445.048}{89.967}{,}%
\TextAt{8}{452.246}{89.967}{i.e.,}%
\TextAt{8}{474.401}{89.967}{there}%
\TextAt{8}{504.352}{89.967}{exists}%
\TextAt{8}{79.226}{72.921}{a}%
\TextAt{8}{88.982}{72.921}{unique}%
\TextAt{13}{127.023}{72.921}{E}%
\TextAt{21}{135.688}{71.128}{r}%
\TextAt{17}{139.745}{71.128}{\Character{43}1}%
\TextAt{9}{151.064}{72.921}{\Character{40}}%
\TextAt{13}{155.616}{72.921}{X}%
\TextAt{9}{166.271}{72.921}{\Character{41}}%
\NativeMathAt{174.144}{72.921}{11.955}{\in}%
\TextAt{14}{185.435}{72.921}{o}%
\TextAt{9}{191.279}{72.921}{\Character{91}}%
\TextAt{13}{194.531}{72.921}{X}%
\TextAt{9}{205.186}{72.921}{\Character{93}}%
\TextAt{8}{212.335}{72.921}{satisfying}%
\TextAt{8}{264.496}{72.921}{(2).}%
\TextAt{8}{287.912}{72.921}{This}%
\TextAt{8}{314.632}{72.921}{finishes}%
\TextAt{8}{355.722}{72.921}{the}%
\TextAt{8}{375.878}{72.921}{proof.}%
\Vector{q 0 G 0 g 0.399 w 452.334 -701.079 m 452.334 -693.009 l S Q}%
\Vector{q 0 G 0 g 0.399 w 452.534 -693.208 m 459.707 -693.208 l S Q}%
\Vector{q 0 G 0 g 0.399 w 452.534 -700.879 m 459.707 -700.879 l S Q}%
\Vector{q 0 G 0 g 0.399 w 459.906 -701.079 m 459.906 -693.009 l S Q}%
\TextAt{8}{297.211}{43.033}{302}%
\Destination{522.36}{688.715}{nav.AMS.55}%
\Destination{522.36}{483.335}{nav.AMS.56}%
\Destination{522.36}{424.215}{nav.AMS.57}%
\Destination{522.36}{297.135}{nav.AMS.58}%
\Destination{78.23}{809.815}{nav.page.302}%
\Destination{79.36}{776.155}{nav.print.329}%
\Destination{79.36}{363.255}{nav.print.330}%
}

% Page 315 (body page 303)
\Page{%
\TextAt{2}{79.226}{796.943}{V.}%
\TextAt{2}{94.073}{796.943}{Local}%
\TextAt{2}{123.615}{796.943}{Class}%
\TextAt{2}{152.208}{796.943}{Field}%
\TextAt{2}{179.939}{796.943}{Theory}%
\Vector{q 0 G 0 g 0.299 w 7.226 18.729 m 460.769 18.729 l S Q}%
\TextAt{12}{79.226}{761.124}{Corollary}%
\TextAt{12}{139.009}{761.124}{(2.3).}%
\TextAt{8}{174.32}{761.124}{Let}%
\TextAt{13}{195.285}{761.124}{\Character{25}}%
\TextAt{8}{206.247}{761.124}{and}%
\TextAt{9}{229.614}{761.124}{\Character{22}}%
\TextAt{13}{229.006}{761.124}{\Character{25}}%
\TextAt{8}{239.968}{761.124}{be}%
\TextAt{8}{255.888}{761.124}{prime}%
\TextAt{8}{289.047}{761.124}{elements}%
\TextAt{8}{336.899}{761.124}{of}%
\TextAt{13}{350.225}{761.124}{K}%
\TextAt{8}{361.042}{761.124}{,}%
\TextAt{8}{368.188}{761.124}{and}%
\TextAt{8}{390.946}{761.124}{let}%
\TextAt{13}{407.847}{761.124}{e}%
\NativeMathAt{416.593}{761.124}{11.955}{\in}%
\TextAt{16}{427.867}{761.124}{\Character{69}}%
\TextAt{21}{434.194}{759.33}{\Character{25}}%
\TextAt{8}{439.84}{761.124}{,}%
\TextAt{9}{447.423}{761.124}{\Character{22}}%
\TextAt{13}{446.986}{761.124}{e}%
\NativeMathAt{455.732}{761.124}{11.955}{\in}%
\TextAt{16}{467.006}{761.124}{\Character{69}}%
\TextAt{17}{473.79}{759.33}{\Character{22}}%
\TextAt{21}{473.333}{759.33}{\Character{25}}%
\TextAt{8}{482.873}{761.124}{be}%
\TextAt{8}{498.793}{761.124}{Lubin-}%
\TextAt{8}{79.226}{744.078}{Tate}%
\TextAt{8}{105.128}{744.078}{series}%
\TextAt{8}{135.382}{744.078}{with}%
\TextAt{8}{160.961}{744.078}{coefficients}%
\TextAt{8}{219.118}{744.078}{in}%
\TextAt{14}{231.69}{744.078}{o}%
\TextAt{21}{237.534}{742.284}{K}%
\TextAt{8}{245.737}{744.078}{.}%
\TextAt{8}{253.825}{744.078}{Let}%
\TextAt{13}{273.714}{744.078}{\Character{25}}%
\TextAt{9}{284.104}{744.078}{\Character{61}}%
\TextAt{13}{296.529}{744.078}{u}%
\TextAt{9}{303.8}{744.078}{\Character{22}}%
\TextAt{13}{303.192}{744.078}{\Character{25}}%
\TextAt{8}{310.261}{744.078}{,}%
\TextAt{13}{316.546}{744.078}{u}%
\NativeMathAt{326.529}{744.078}{11.955}{\in}%
\TextAt{13}{337.82}{744.078}{U}%
\TextAt{21}{345.77}{742.284}{K}%
\TextAt{8}{353.973}{744.078}{,}%
\TextAt{8}{360.257}{744.078}{and}%
\TextAt{13}{381.939}{744.078}{u}%
\TextAt{9}{391.922}{744.078}{\Character{61}}%
\TextAt{13}{404.348}{744.078}{\Character{34}}%
\TextAt{21}{409.826}{748.416}{\Character{39}}%
\TextAt{3}{415.344}{748.416}{\Character{0}}%
\TextAt{17}{421.931}{748.416}{1}%
\TextAt{8}{426.663}{744.078}{,}%
\TextAt{13}{432.948}{744.078}{\Character{34}}%
\NativeMathAt{441.747}{744.078}{11.955}{\in}%
\TextAt{13}{453.037}{744.078}{U}%
\TextAt{30}{462.958}{742.439}{\Character{98}}%
\TextAt{21}{460.988}{740.425}{K}%
\TextAt{8}{469.19}{744.078}{.}%
\TextAt{8}{477.278}{744.078}{Then}%
\TextAt{8}{506.755}{744.078}{there}%
\TextAt{8}{79.226}{727.032}{is}%
\TextAt{8}{91.789}{727.032}{a}%
\TextAt{8}{102.344}{727.032}{uniquely}%
\TextAt{8}{150.617}{727.032}{determined}%
\TextAt{8}{212.543}{727.032}{power}%
\TextAt{8}{247.486}{727.032}{series}%
\TextAt{13}{279.62}{727.032}{\Character{18}}%
\TextAt{9}{285.4}{727.032}{\Character{40}}%
\TextAt{13}{289.953}{727.032}{X}%
\TextAt{9}{300.608}{727.032}{\Character{41}}%
\NativeMathAt{309.843}{727.032}{11.955}{\in}%
\TextAt{14}{322.496}{727.032}{o}%
\TextAt{30}{330.311}{725.393}{\Character{98}}%
\TextAt{21}{328.34}{723.379}{K}%
\TextAt{9}{336.543}{727.032}{\Character{91}\Character{91}}%
\TextAt{13}{343.046}{727.032}{X}%
\TextAt{9}{353.701}{727.032}{\Character{93}\Character{93}}%
\TextAt{8}{364.901}{727.032}{such}%
\TextAt{8}{392.098}{727.032}{that}%
\TextAt{13}{418.266}{727.032}{\Character{18}}%
\TextAt{9}{424.046}{727.032}{\Character{40}}%
\TextAt{13}{428.598}{727.032}{X}%
\TextAt{9}{439.254}{727.032}{\Character{41}}%
\NativeMathAt{448.489}{727.032}{11.955}{\equiv}%
\TextAt{13}{462.47}{727.032}{\Character{34}X}%
\TextAt{8}{481.924}{727.032}{mod}%
\TextAt{8}{507.353}{727.032}{deg}%
\TextAt{9}{526.917}{727.032}{2}%
\TextAt{8}{79.226}{709.986}{and}%
\TextAt{13}{270.994}{692.94}{e}%
\TextAt{16}{279.076}{692.94}{\Character{14}}%
\TextAt{13}{287.71}{692.94}{\Character{18}}%
\TextAt{9}{296.811}{692.94}{\Character{61}}%
\TextAt{13}{309.237}{692.94}{\Character{18}}%
\TextAt{21}{315.017}{697.876}{\Character{39}}%
\TextAt{16}{323.69}{692.94}{\Character{14}}%
\TextAt{9}{332.761}{692.94}{\Character{22}}%
\TextAt{13}{332.325}{692.94}{e\Character{58}}%
\TextAt{8}{79.226}{669.418}{Furthermore,}%
\TextAt{8}{152.104}{669.418}{there}%
\TextAt{8}{183.478}{669.418}{is}%
\TextAt{8}{196.704}{669.418}{a}%
\TextAt{8}{207.921}{669.418}{uniquely}%
\TextAt{8}{256.858}{669.418}{determined}%
\TextAt{8}{319.447}{669.418}{power}%
\TextAt{8}{355.053}{669.418}{series}%
\TextAt{9}{387.849}{669.418}{\Character{91}}%
\TextAt{13}{391.101}{669.418}{u}%
\TextAt{9}{397.764}{669.418}{\Character{93}\Character{40}}%
\TextAt{13}{405.568}{669.418}{X}%
\TextAt{9}{416.223}{669.418}{\Character{41}}%
\NativeMathAt{426.588}{669.418}{11.955}{\in}%
\TextAt{14}{440.37}{669.418}{o}%
\TextAt{21}{446.215}{667.625}{K}%
\TextAt{9}{454.417}{669.418}{\Character{91}\Character{91}}%
\TextAt{13}{460.921}{669.418}{X}%
\TextAt{9}{471.576}{669.418}{\Character{93}\Character{93}}%
\TextAt{8}{483.439}{669.418}{such}%
\TextAt{8}{511.298}{669.418}{that}%
\TextAt{9}{79.226}{652.372}{\Character{91}}%
\TextAt{13}{82.478}{652.372}{u}%
\TextAt{9}{89.14}{652.372}{\Character{93}\Character{40}}%
\TextAt{13}{96.944}{652.372}{X}%
\TextAt{9}{107.599}{652.372}{\Character{41}}%
\NativeMathAt{115.473}{652.372}{11.955}{\equiv}%
\TextAt{13}{128.092}{652.372}{uX}%
\TextAt{8}{148.73}{652.372}{mod}%
\TextAt{8}{174.159}{652.372}{deg}%
\TextAt{9}{193.723}{652.372}{2}%
\TextAt{8}{203.473}{652.372}{and}%
\TextAt{9}{267.053}{635.326}{\Character{22}}%
\TextAt{13}{266.617}{635.326}{e}%
\TextAt{16}{274.699}{635.326}{\Character{14}}%
\TextAt{9}{283.333}{635.326}{\Character{91}}%
\TextAt{13}{286.585}{635.326}{u}%
\TextAt{9}{293.247}{635.326}{\Character{93}}%
\TextAt{9}{299.798}{635.326}{\Character{61}}%
\TextAt{9}{312.22}{635.326}{\Character{91}}%
\TextAt{13}{315.497}{635.326}{u}%
\TextAt{9}{322.159}{635.326}{\Character{93}}%
\TextAt{16}{328.068}{635.326}{\Character{14}}%
\TextAt{9}{337.138}{635.326}{\Character{22}}%
\TextAt{13}{336.702}{635.326}{e\Character{58}}%
\TextAt{8}{79.226}{611.804}{They}%
\TextAt{8}{109.461}{611.804}{satisfy}%
\TextAt{13}{275.482}{594.759}{\Character{18}}%
\TextAt{21}{281.262}{599.695}{\Character{39}}%
\TextAt{9}{290.599}{594.759}{\Character{61}}%
\TextAt{13}{303.025}{594.759}{\Character{18}}%
\TextAt{16}{311.462}{594.759}{\Character{14}}%
\TextAt{9}{320.096}{594.759}{\Character{91}}%
\TextAt{13}{323.348}{594.759}{u}%
\TextAt{9}{330.01}{594.759}{\Character{93}}%
\TextAt{13}{333.262}{594.759}{\Character{58}}%
\TextAt{19}{79.226}{553.304}{Proof.}%
\TextAt{8}{116}{553.304}{Putting}%
\TextAt{13}{159.899}{553.304}{L}%
\TextAt{9}{167.863}{553.304}{\Character{40}}%
\TextAt{13}{172.416}{553.304}{X}%
\TextAt{9}{183.071}{553.304}{\Character{41}}%
\TextAt{9}{192.312}{553.304}{\Character{61}}%
\TextAt{13}{206.137}{553.304}{\Character{34}X}%
\TextAt{8}{222.271}{553.304}{,}%
\TextAt{8}{230.435}{553.304}{we}%
\TextAt{8}{248.474}{553.304}{have}%
\TextAt{13}{276.281}{553.304}{\Character{25}}%
\TextAt{13}{283.335}{553.304}{L}%
\TextAt{9}{291.315}{553.304}{\Character{40}}%
\TextAt{13}{295.868}{553.304}{X}%
\TextAt{9}{306.523}{553.304}{\Character{41}}%
\TextAt{9}{315.764}{553.304}{\Character{61}}%
\TextAt{9}{330.17}{553.304}{\Character{22}}%
\TextAt{13}{329.589}{553.304}{\Character{25}}%
\TextAt{13}{336.643}{553.304}{L}%
\TextAt{21}{344.623}{557.643}{\Character{39}}%
\TextAt{9}{350.639}{553.304}{\Character{40}}%
\TextAt{13}{355.192}{553.304}{X}%
\TextAt{9}{365.847}{553.304}{\Character{41}}%
\TextAt{8}{375.109}{553.304}{and}%
\TextAt{8}{398.684}{553.304}{the}%
\TextAt{8}{419.652}{553.304}{first}%
\TextAt{8}{444.59}{553.304}{claim}%
\TextAt{8}{476.617}{553.304}{follows}%
\TextAt{8}{515.865}{553.304}{im-}%
\TextAt{8}{79.226}{536.258}{mediately}%
\TextAt{8}{133.732}{536.258}{from}%
\TextAt{8}{162.223}{536.258}{(2.2).}%
\TextAt{8}{197.301}{536.258}{In}%
\TextAt{8}{212.773}{536.258}{the}%
\TextAt{8}{233.781}{536.258}{same}%
\TextAt{8}{263.958}{536.258}{way,}%
\TextAt{8}{291.036}{536.258}{with}%
\TextAt{8}{318.546}{536.258}{the}%
\TextAt{8}{339.554}{536.258}{linear}%
\TextAt{8}{372.923}{536.258}{form}%
\TextAt{13}{401.414}{536.258}{L}%
\TextAt{9}{409.378}{536.258}{\Character{40}}%
\TextAt{13}{413.931}{536.258}{X}%
\TextAt{9}{424.586}{536.258}{\Character{41}}%
\TextAt{9}{433.899}{536.258}{\Character{61}}%
\TextAt{13}{447.784}{536.258}{uX}%
\TextAt{8}{465.102}{536.258}{,}%
\TextAt{8}{473.314}{536.258}{one}%
\TextAt{8}{495.625}{536.258}{obtains}%
\TextAt{8}{79.226}{519.212}{the}%
\TextAt{8}{101.091}{519.212}{existence}%
\TextAt{8}{152.605}{519.212}{and}%
\TextAt{8}{177.076}{519.212}{uniqueness}%
\TextAt{8}{237.76}{519.212}{of}%
\TextAt{8}{252.798}{519.212}{the}%
\TextAt{8}{274.663}{519.212}{power}%
\TextAt{8}{310.516}{519.212}{series}%
\TextAt{9}{343.559}{519.212}{\Character{91}}%
\TextAt{13}{346.811}{519.212}{u}%
\TextAt{9}{353.473}{519.212}{\Character{93}\Character{40}}%
\TextAt{13}{361.277}{519.212}{X}%
\TextAt{9}{371.932}{519.212}{\Character{41}}%
\NativeMathAt{382.717}{519.212}{11.955}{\in}%
\TextAt{14}{396.92}{519.212}{o}%
\TextAt{21}{402.764}{517.419}{K}%
\TextAt{9}{410.967}{519.212}{\Character{91}\Character{91}}%
\TextAt{13}{417.47}{519.212}{X}%
\TextAt{9}{428.125}{519.212}{\Character{93}\Character{93}}%
\TextAt{8}{434.629}{519.212}{.}%
\TextAt{8}{448.203}{519.212}{Finally,}%
\TextAt{8}{492.445}{519.212}{defining}%
\TextAt{13}{79.226}{502.166}{\Character{18}}%
\TextAt{17}{84.681}{500.373}{1}%
\TextAt{9}{92.734}{502.166}{\Character{61}}%
\TextAt{13}{105.16}{502.166}{\Character{18}}%
\TextAt{21}{110.94}{506.505}{\Character{39}}%
\TextAt{23}{116.459}{509.317}{\Character{0}}%
\TextAt{24}{122.215}{509.317}{1}%
\TextAt{16}{129.521}{502.166}{\Character{14}}%
\TextAt{9}{138.155}{502.166}{\Character{91}}%
\TextAt{13}{141.406}{502.166}{u}%
\TextAt{9}{148.069}{502.166}{\Character{93}}%
\TextAt{8}{151.321}{502.166}{,}%
\TextAt{8}{158.47}{502.166}{we}%
\TextAt{8}{175.697}{502.166}{get}%
\TextAt{13}{201.81}{473.165}{e}%
\TextAt{16}{209.892}{473.165}{\Character{14}}%
\TextAt{13}{218.526}{473.165}{\Character{18}}%
\TextAt{17}{223.981}{471.372}{1}%
\TextAt{9}{232.034}{473.165}{\Character{61}}%
\TextAt{9}{244.444}{473.165}{\Character{40}}%
\TextAt{13}{249.012}{473.165}{e}%
\TextAt{16}{257.094}{473.165}{\Character{14}}%
\TextAt{13}{265.729}{473.165}{\Character{18}}%
\TextAt{9}{271.509}{473.165}{\Character{41}}%
\TextAt{21}{276.061}{478.101}{\Character{39}}%
\TextAt{23}{281.58}{480.914}{\Character{0}}%
\TextAt{24}{287.336}{480.914}{1}%
\TextAt{16}{294.641}{473.165}{\Character{14}}%
\TextAt{9}{303.276}{473.165}{\Character{91}}%
\TextAt{13}{306.527}{473.165}{u}%
\TextAt{9}{313.19}{473.165}{\Character{93}}%
\TextAt{9}{319.741}{473.165}{\Character{61}}%
\TextAt{9}{332.163}{473.165}{\Character{40}}%
\TextAt{13}{336.74}{473.165}{\Character{18}}%
\TextAt{21}{342.52}{478.101}{\Character{39}}%
\TextAt{16}{351.194}{473.165}{\Character{14}}%
\TextAt{9}{360.264}{473.165}{\Character{22}}%
\TextAt{13}{359.828}{473.165}{e}%
\TextAt{9}{365.253}{473.165}{\Character{41}}%
\TextAt{21}{369.806}{478.101}{\Character{39}}%
\TextAt{23}{375.324}{480.914}{\Character{0}}%
\TextAt{24}{381.08}{480.914}{1}%
\TextAt{16}{388.386}{473.165}{\Character{14}}%
\TextAt{9}{397.02}{473.165}{\Character{91}}%
\TextAt{13}{400.272}{473.165}{u}%
\TextAt{9}{406.934}{473.165}{\Character{93}}%
\TextAt{9}{232.034}{453.13}{\Character{61}}%
\TextAt{9}{244.444}{453.13}{\Character{40}}%
\TextAt{13}{249.012}{453.13}{\Character{18}}%
\TextAt{21}{254.792}{458.066}{\Character{39}}%
\TextAt{23}{260.311}{460.879}{\Character{0}}%
\TextAt{24}{266.067}{460.879}{1}%
\TextAt{16}{273.373}{453.13}{\Character{14}}%
\TextAt{9}{282.007}{453.13}{\Character{91}}%
\TextAt{13}{285.259}{453.13}{u}%
\TextAt{9}{291.921}{453.13}{\Character{93}\Character{41}}%
\TextAt{21}{299.725}{458.066}{\Character{39}}%
\TextAt{16}{308.398}{453.13}{\Character{14}}%
\TextAt{9}{317.469}{453.13}{\Character{22}}%
\TextAt{13}{317.032}{453.13}{e}%
\TextAt{9}{325.779}{453.13}{\Character{61}}%
\TextAt{13}{338.204}{453.13}{\Character{18}}%
\TextAt{21}{343.985}{458.798}{\Character{39}}%
\TextAt{17}{343.659}{450.199}{1}%
\TextAt{16}{352.658}{453.13}{\Character{14}}%
\TextAt{9}{361.728}{453.13}{\Character{22}}%
\TextAt{13}{361.292}{453.13}{e\Character{59}}%
\TextAt{8}{79.226}{424.129}{and}%
\TextAt{8}{101.989}{424.129}{thus}%
\TextAt{13}{127.74}{424.129}{\Character{18}}%
\TextAt{17}{133.195}{422.336}{1}%
\TextAt{9}{141.249}{424.129}{\Character{61}}%
\TextAt{13}{153.674}{424.129}{\Character{18}}%
\TextAt{8}{163.352}{424.129}{because}%
\TextAt{8}{206.653}{424.129}{of}%
\TextAt{8}{219.983}{424.129}{uniqueness.}%
\TextAt{8}{283.509}{424.129}{Hence}%
\TextAt{13}{318.287}{424.129}{\Character{18}}%
\TextAt{21}{324.067}{428.468}{\Character{39}}%
\TextAt{9}{333.404}{424.129}{\Character{61}}%
\TextAt{13}{345.83}{424.129}{\Character{18}}%
\TextAt{16}{354.267}{424.129}{\Character{14}}%
\TextAt{9}{362.901}{424.129}{\Character{91}}%
\TextAt{13}{366.153}{424.129}{u}%
\TextAt{9}{372.815}{424.129}{\Character{93}}%
\TextAt{8}{376.067}{424.129}{.}%
\Vector{q 0 G 0 g 0.399 w 452.334 -349.871 m 452.334 -341.801 l S Q}%
\Vector{q 0 G 0 g 0.399 w 452.534 -342 m 459.707 -342 l S Q}%
\Vector{q 0 G 0 g 0.399 w 452.534 -349.672 m 459.707 -349.672 l S Q}%
\Vector{q 0 G 0 g 0.399 w 459.906 -349.871 m 459.906 -341.801 l S Q}%
\TextAt{12}{79.226}{389.15}{Theorem}%
\TextAt{12}{135.813}{389.15}{(2.4).}%
\TextAt{8}{170.758}{389.15}{Let}%
\TextAt{13}{191.356}{389.15}{a}%
\TextAt{9}{200.822}{389.15}{\Character{61}}%
\TextAt{13}{213.247}{389.15}{up}%
\TextAt{21}{225.785}{393.489}{v}%
\TextAt{22}{229.908}{392.493}{p}%
\TextAt{17}{234.247}{393.489}{\Character{40}}%
\TextAt{21}{237.54}{393.489}{a}%
\TextAt{17}{242.038}{393.489}{\Character{41}}%
\NativeMathAt{249.15}{389.15}{11.955}{\in}%
\TextAt{18}{260.441}{389.15}{Q}%
\TextAt{3}{269.74}{393.489}{\Character{3}}%
\TextAt{21}{269.74}{386.195}{p}%
\TextAt{8}{274.501}{389.15}{,}%
\TextAt{8}{281.353}{389.15}{and}%
\TextAt{8}{303.744}{389.15}{let}%
\TextAt{13}{320.278}{389.15}{\Character{16}}%
\TextAt{8}{329.82}{389.15}{be}%
\TextAt{8}{345.373}{389.15}{a}%
\TextAt{8}{354.757}{389.15}{primitive}%
\TextAt{13}{404.466}{389.15}{p}%
\TextAt{21}{410.341}{393.489}{n}%
\TextAt{8}{415.977}{389.15}{-th}%
\TextAt{8}{434.46}{389.15}{root}%
\TextAt{8}{459.134}{389.15}{of}%
\TextAt{8}{472.093}{389.15}{unity.}%
\TextAt{8}{506.109}{389.15}{Then}%
\TextAt{8}{79.226}{372.104}{one}%
\TextAt{8}{100.686}{372.104}{has}%
\TextAt{9}{249.302}{355.058}{\Character{40}}%
\TextAt{13}{253.855}{355.058}{a\Character{59}}%
\TextAt{18}{265.244}{355.058}{Q}%
\TextAt{21}{274.542}{353.265}{p}%
\TextAt{9}{279.303}{355.058}{\Character{40}}%
\TextAt{13}{283.855}{355.058}{\Character{16}}%
\TextAt{9}{289.871}{355.058}{\Character{41}}%
\TextAt{16}{294.423}{355.058}{\Character{106}}%
\TextAt{18}{297.744}{355.058}{Q}%
\TextAt{21}{307.043}{353.265}{p}%
\TextAt{9}{311.804}{355.058}{\Character{41}}%
\TextAt{13}{316.356}{355.058}{\Character{16}}%
\TextAt{9}{325.692}{355.058}{\Character{61}}%
\TextAt{13}{338.118}{355.058}{\Character{16}}%
\TextAt{21}{344.133}{359.995}{u}%
\TextAt{23}{349.036}{362.807}{\Character{0}}%
\TextAt{24}{354.793}{362.807}{1}%
\TextAt{13}{359.442}{355.058}{\Character{58}}%
\TextAt{19}{79.226}{313.604}{Proof.}%
\TextAt{8}{116}{313.604}{As}%
\TextAt{18}{134.939}{313.604}{N}%
\TextAt{8}{149.122}{313.604}{is}%
\TextAt{8}{162.538}{313.604}{dense}%
\TextAt{8}{196.11}{313.604}{in}%
\TextAt{18}{211.414}{313.604}{Z}%
\TextAt{21}{219.384}{311.811}{p}%
\TextAt{8}{224.145}{313.604}{,}%
\TextAt{8}{233.359}{313.604}{we}%
\TextAt{8}{252.238}{313.604}{may}%
\TextAt{8}{279.258}{313.604}{assume}%
\TextAt{8}{321.354}{313.604}{that}%
\TextAt{13}{348.375}{313.604}{u}%
\NativeMathAt{361.173}{313.604}{11.955}{\in}%
\TextAt{18}{375.278}{313.604}{N}%
\TextAt{8}{383.913}{313.604}{,}%
\TextAt{9}{393.126}{313.604}{\Character{40}}%
\TextAt{13}{397.679}{313.604}{u\Character{59}}%
\TextAt{13}{409.562}{313.604}{p}%
\TextAt{9}{415.46}{313.604}{\Character{41}}%
\TextAt{9}{426.136}{313.604}{\Character{61}}%
\TextAt{9}{441.367}{313.604}{1}%
\TextAt{8}{447.241}{313.604}{.}%
\TextAt{8}{460.644}{313.604}{Let}%
\TextAt{13}{483.265}{313.604}{K}%
\TextAt{9}{500.218}{313.604}{\Character{61}}%
\TextAt{18}{515.458}{313.604}{Q}%
\TextAt{21}{524.757}{311.811}{p}%
\TextAt{8}{529.518}{313.604}{,}%
\TextAt{13}{79.226}{296.558}{L}%
\TextAt{9}{90.511}{296.558}{\Character{61}}%
\TextAt{18}{102.937}{296.558}{Q}%
\TextAt{21}{112.235}{294.765}{p}%
\TextAt{9}{116.996}{296.558}{\Character{40}}%
\TextAt{13}{121.549}{296.558}{\Character{16}}%
\TextAt{9}{127.564}{296.558}{\Character{41}}%
\TextAt{8}{132.117}{296.558}{,}%
\TextAt{8}{139.266}{296.558}{and}%
\TextAt{8}{162.028}{296.558}{let}%
\TextAt{13}{178.933}{296.558}{\Character{27}}%
\NativeMathAt{189.336}{296.558}{11.955}{\in}%
\TextAt{13}{200.627}{296.558}{G}%
\TextAt{9}{209.861}{296.558}{\Character{40}}%
\TextAt{13}{214.413}{296.558}{L}%
\TextAt{16}{222.378}{296.558}{\Character{106}}%
\TextAt{13}{225.699}{296.558}{K}%
\TextAt{9}{236.516}{296.558}{\Character{41}}%
\TextAt{8}{244.966}{296.558}{be}%
\TextAt{8}{260.89}{296.558}{the}%
\TextAt{8}{281.047}{296.558}{automorphism}%
\TextAt{8}{358.516}{296.558}{defined}%
\TextAt{8}{398.829}{296.558}{by}%
\TextAt{13}{280.017}{267.557}{\Character{16}}%
\TextAt{21}{286.033}{272.493}{\Character{27}}%
\TextAt{9}{294.977}{267.557}{\Character{61}}%
\TextAt{13}{307.403}{267.557}{\Character{16}}%
\TextAt{21}{313.418}{272.493}{u}%
\TextAt{23}{318.321}{275.306}{\Character{0}}%
\TextAt{24}{324.077}{275.306}{1}%
\TextAt{13}{328.727}{267.557}{\Character{58}}%
\TextAt{8}{79.226}{238.556}{Since}%
\TextAt{18}{111.067}{238.556}{Q}%
\TextAt{21}{120.366}{236.762}{p}%
\TextAt{9}{125.126}{238.556}{\Character{40}}%
\TextAt{13}{129.679}{238.556}{\Character{16}}%
\TextAt{9}{135.694}{238.556}{\Character{41}}%
\TextAt{16}{140.247}{238.556}{\Character{106}}%
\TextAt{18}{143.568}{238.556}{Q}%
\TextAt{21}{152.866}{236.762}{p}%
\TextAt{8}{162.808}{238.556}{is}%
\TextAt{8}{175.855}{238.556}{totally}%
\TextAt{8}{214.546}{238.556}{ramified,}%
\TextAt{8}{264.928}{238.556}{we}%
\TextAt{8}{283.439}{238.556}{have}%
\TextAt{13}{311.717}{238.556}{G}%
\TextAt{9}{320.951}{238.556}{\Character{40}}%
\TextAt{13}{325.503}{238.556}{L}%
\TextAt{16}{333.467}{238.556}{\Character{106}}%
\TextAt{13}{336.788}{238.556}{K}%
\TextAt{9}{347.606}{238.556}{\Character{41}}%
\NativeMathAt{357.666}{238.556}{11.955}{\cong}%
\TextAt{13}{372.473}{238.556}{G}%
\TextAt{9}{381.706}{238.556}{\Character{40}}%
\TextAt{20}{387.245}{241.578}{\Character{101}}%
\TextAt{13}{386.259}{238.556}{L}%
\TextAt{16}{394.223}{238.556}{\Character{106}}%
\TextAt{20}{400.282}{241.578}{\Character{101}}%
\TextAt{13}{397.544}{238.556}{K}%
\TextAt{9}{408.362}{238.556}{\Character{41}}%
\TextAt{8}{412.914}{238.556}{,}%
\TextAt{8}{421.668}{238.556}{and}%
\TextAt{8}{445.714}{238.556}{we}%
\TextAt{8}{464.225}{238.556}{view}%
\TextAt{13}{492.491}{238.556}{\Character{27}}%
\TextAt{8}{504.754}{238.556}{as}%
\TextAt{8}{520.408}{238.556}{an}%
\TextAt{8}{79.226}{221.51}{element}%
\TextAt{8}{122.535}{221.51}{of}%
\TextAt{13}{135.931}{221.51}{G}%
\TextAt{9}{145.165}{221.51}{\Character{40}}%
\TextAt{20}{150.704}{224.532}{\Character{101}}%
\TextAt{13}{149.717}{221.51}{L}%
\TextAt{16}{157.682}{221.51}{\Character{106}}%
\TextAt{20}{163.741}{224.532}{\Character{101}}%
\TextAt{13}{161.003}{221.51}{K}%
\TextAt{9}{171.82}{221.51}{\Character{41}}%
\TextAt{8}{176.373}{221.51}{.}%
\TextAt{8}{185.021}{221.51}{Then}%
\TextAt{20}{215.866}{221.51}{\Character{101}}%
\TextAt{13}{215.645}{221.51}{\Character{27}}%
\TextAt{9}{226.162}{221.51}{\Character{61}}%
\TextAt{13}{238.702}{221.51}{\Character{27}}%
\TextAt{13}{245.768}{221.51}{\Character{39}}%
\TextAt{21}{253.452}{219.717}{L}%
\NativeMathAt{263.144}{221.51}{11.955}{\in}%
\TextAt{8}{274.549}{221.51}{Frob}%
\TextAt{9}{299.105}{221.51}{\Character{40}}%
\TextAt{20}{304.644}{224.532}{\Character{101}}%
\TextAt{13}{303.657}{221.51}{L}%
\TextAt{16}{311.622}{221.51}{\Character{106}}%
\TextAt{13}{314.943}{221.51}{K}%
\TextAt{9}{325.76}{221.51}{\Character{41}}%
\TextAt{8}{334.277}{221.51}{is}%
\TextAt{8}{346.107}{221.51}{an}%
\TextAt{8}{362.433}{221.51}{element}%
\TextAt{8}{405.742}{221.51}{such}%
\TextAt{8}{432.205}{221.51}{that}%
\TextAt{13}{457.641}{221.51}{d}%
\TextAt{21}{463.724}{219.717}{K}%
\TextAt{9}{471.926}{221.51}{\Character{40}}%
\TextAt{20}{476.699}{221.51}{\Character{101}}%
\TextAt{13}{476.479}{221.51}{\Character{27}}%
\TextAt{9}{483.561}{221.51}{\Character{41}}%
\TextAt{9}{491.535}{221.51}{\Character{61}}%
\TextAt{9}{504.064}{221.51}{1}%
\TextAt{8}{513.904}{221.51}{and}%
\TextAt{20}{79.446}{204.464}{\Character{101}}%
\TextAt{13}{79.226}{204.464}{\Character{27}}%
\TextAt{16}{86.309}{204.464}{\Character{106}}%
\TextAt{21}{89.629}{202.671}{L}%
\TextAt{9}{100.277}{204.464}{\Character{61}}%
\TextAt{13}{113.773}{204.464}{\Character{27}}%
\TextAt{8}{120.855}{204.464}{.}%
\TextAt{8}{131.187}{204.464}{The}%
\TextAt{8}{155.869}{204.464}{fixed}%
\TextAt{8}{184.783}{204.464}{field}%
\TextAt{9}{210.768}{204.464}{\Character{6}}%
\TextAt{8}{223.747}{204.464}{of}%
\TextAt{20}{237.926}{204.464}{\Character{101}}%
\TextAt{13}{237.705}{204.464}{\Character{27}}%
\TextAt{8}{249.313}{204.464}{is}%
\TextAt{8}{261.705}{204.464}{totally}%
\TextAt{8}{299.741}{204.464}{ramified}%
\TextAt{8}{345.894}{204.464}{because}%
\TextAt{13}{389.823}{204.464}{f}%
\TextAt{17}{395.593}{202.604}{\Character{6}}%
\TextAt{3}{401.709}{202.604}{\Character{106}}%
\TextAt{21}{404.062}{202.604}{K}%
\TextAt{9}{416.655}{204.464}{\Character{61}}%
\TextAt{13}{430.151}{204.464}{d}%
\TextAt{21}{436.233}{202.671}{K}%
\TextAt{9}{444.436}{204.464}{\Character{40}}%
\TextAt{20}{449.209}{204.464}{\Character{101}}%
\TextAt{13}{448.988}{204.464}{\Character{27}}%
\TextAt{9}{456.071}{204.464}{\Character{41}}%
\TextAt{9}{465.002}{204.464}{\Character{61}}%
\TextAt{9}{478.487}{204.464}{1}%
\TextAt{8}{488.888}{204.464}{by}%
\TextAt{8}{505.775}{204.464}{chap.}%
\TextAt{8}{79.226}{187.418}{IV,}%
\TextAt{8}{99.137}{187.418}{(4.5).}%
\TextAt{8}{131.585}{187.418}{The}%
\TextAt{8}{155.405}{187.418}{proof}%
\TextAt{8}{185.741}{187.418}{of}%
\TextAt{8}{198.837}{187.418}{the}%
\TextAt{8}{218.76}{187.418}{theorem}%
\TextAt{8}{264.051}{187.418}{is}%
\TextAt{8}{275.582}{187.418}{based}%
\TextAt{8}{307.926}{187.418}{on}%
\TextAt{8}{323.951}{187.418}{the}%
\TextAt{8}{343.874}{187.418}{fact}%
\TextAt{8}{366.726}{187.418}{that}%
\TextAt{8}{391.861}{187.418}{the}%
\TextAt{8}{411.783}{187.418}{field}%
\TextAt{9}{436.907}{187.418}{\Character{6}}%
\TextAt{8}{449.025}{187.418}{can}%
\TextAt{8}{470.25}{187.418}{be}%
\TextAt{8}{485.941}{187.418}{explicitly}%
\TextAt{8}{79.226}{170.372}{generated}%
\TextAt{8}{132.558}{170.372}{by}%
\TextAt{8}{148.817}{170.372}{a}%
\TextAt{8}{158.573}{170.372}{prime}%
\TextAt{8}{191.736}{170.372}{element}%
\TextAt{13}{234.978}{170.372}{\Character{25}}%
\TextAt{17}{241.618}{168.579}{\Character{6}}%
\TextAt{8}{252.13}{170.372}{which}%
\TextAt{8}{285.616}{170.372}{is}%
\TextAt{8}{297.38}{170.372}{given}%
\TextAt{8}{327.95}{170.372}{by}%
\TextAt{8}{344.209}{170.372}{the}%
\TextAt{8}{364.365}{170.372}{power}%
\TextAt{8}{398.509}{170.372}{series}%
\TextAt{13}{429.844}{170.372}{\Character{18}}%
\TextAt{8}{439.521}{170.372}{of}%
\TextAt{8}{452.851}{170.372}{(2.3).}%
\TextAt{8}{97.159}{150.337}{In}%
\TextAt{8}{112.651}{150.337}{order}%
\TextAt{8}{144.091}{150.337}{to}%
\TextAt{8}{159.272}{150.337}{do}%
\TextAt{8}{176.402}{150.337}{this,}%
\TextAt{8}{203.565}{150.337}{assume}%
\TextAt{20}{245.1}{150.337}{\Character{101}}%
\TextAt{13}{244.88}{150.337}{\Character{27}}%
\TextAt{8}{256.731}{150.337}{and}%
\TextAt{13}{280.364}{150.337}{\Character{39}}%
\TextAt{9}{292.837}{150.337}{\Character{61}}%
\TextAt{13}{306.747}{150.337}{\Character{39}}%
\TextAt{21}{314.415}{148.544}{L}%
\TextAt{8}{325.441}{150.337}{have}%
\TextAt{8}{353.306}{150.337}{been}%
\TextAt{8}{381.805}{150.337}{extended}%
\TextAt{8}{432.422}{150.337}{continuously}%
\TextAt{8}{501.329}{150.337}{to}%
\TextAt{8}{516.51}{150.337}{the}%
\TextAt{8}{79.226}{133.291}{completion}%
\TextAt{20}{140.048}{136.313}{\Character{98}}%
\TextAt{13}{139.062}{133.291}{L}%
\TextAt{8}{150.924}{133.291}{of}%
\TextAt{20}{165.24}{136.313}{\Character{101}}%
\TextAt{13}{164.254}{133.291}{L}%
\TextAt{8}{172.218}{133.291}{,}%
\TextAt{8}{179.367}{133.291}{and}%
\TextAt{8}{202.13}{133.291}{consider}%
\TextAt{8}{247.715}{133.291}{the}%
\TextAt{8}{267.872}{133.291}{two}%
\TextAt{8}{289.989}{133.291}{Lubin-Tate}%
\TextAt{8}{350.948}{133.291}{polynomials}%
\TextAt{13}{184.537}{104.29}{e}%
\TextAt{9}{189.962}{104.29}{\Character{40}}%
\TextAt{13}{194.515}{104.29}{X}%
\TextAt{9}{205.17}{104.29}{\Character{41}}%
\TextAt{9}{213.025}{104.29}{\Character{61}}%
\TextAt{13}{225.468}{104.29}{upX}%
\TextAt{9}{251.318}{104.29}{\Character{43}}%
\TextAt{13}{263.079}{104.29}{X}%
\TextAt{21}{273.734}{109.226}{p}%
\TextAt{8}{290.45}{104.29}{and}%
\TextAt{9}{321.707}{104.29}{\Character{22}}%
\TextAt{13}{321.271}{104.29}{e}%
\TextAt{9}{326.696}{104.29}{\Character{40}}%
\TextAt{13}{331.249}{104.29}{X}%
\TextAt{9}{341.904}{104.29}{\Character{41}}%
\TextAt{9}{349.759}{104.29}{\Character{61}}%
\TextAt{9}{362.18}{104.29}{\Character{40}1}%
\TextAt{9}{375.223}{104.29}{\Character{43}}%
\TextAt{13}{387.026}{104.29}{X}%
\TextAt{9}{397.681}{104.29}{\Character{41}}%
\TextAt{21}{402.233}{109.226}{p}%
\TextAt{16}{409.651}{104.29}{\Character{0}}%
\TextAt{9}{421.606}{104.29}{1}%
\TextAt{8}{297.211}{43.033}{303}%
\Destination{78.23}{809.815}{nav.page.303}%
\Destination{79.36}{404.155}{nav.print.331}%
\Destination{79.36}{775.155}{nav.number.5.2.3}%
\Destination{79.36}{403.155}{nav.number.5.2.4}%
\Hotlink{161.723}{534.055}{25.078}{11.047}{nav.number.5.2.2}%
\Hotlink{98.637}{185.255}{25.078}{11.047}{nav.number.4.4.5}%
\Hotlink{452.351}{168.155}{25.078}{11.047}{nav.number.5.2.3}%
}

% Page 316 (body page 304)
\Page{%
\TextAt{2}{79.226}{796.943}{V.}%
\TextAt{2}{94.073}{796.943}{Local}%
\TextAt{2}{123.615}{796.943}{Class}%
\TextAt{2}{152.208}{796.943}{Field}%
\TextAt{2}{179.939}{796.943}{Theory}%
\Vector{q 0 G 0 g 0.299 w 7.226 18.729 m 460.769 18.729 l S Q}%
\TextAt{8}{79.226}{761.124}{as}%
\TextAt{8}{94.513}{761.124}{well}%
\TextAt{8}{119.161}{761.124}{as}%
\TextAt{8}{134.448}{761.124}{the}%
\TextAt{8}{155.521}{761.124}{polynomial}%
\TextAt{9}{216.931}{761.124}{\Character{91}}%
\TextAt{13}{220.182}{761.124}{u}%
\TextAt{9}{226.845}{761.124}{\Character{93}\Character{40}}%
\TextAt{13}{234.649}{761.124}{X}%
\TextAt{9}{245.304}{761.124}{\Character{41}}%
\TextAt{9}{254.725}{761.124}{\Character{61}}%
\TextAt{9}{268.7}{761.124}{\Character{40}1}%
\TextAt{9}{282.377}{761.124}{\Character{43}}%
\TextAt{13}{294.8}{761.124}{X}%
\TextAt{9}{305.455}{761.124}{\Character{41}}%
\TextAt{21}{310.008}{765.462}{u}%
\TextAt{16}{318.69}{761.124}{\Character{0}}%
\TextAt{9}{331.27}{761.124}{1}%
\TextAt{8}{337.123}{761.124}{.}%
\TextAt{8}{348.322}{761.124}{Then}%
\TextAt{9}{380.233}{761.124}{\Character{22}}%
\TextAt{13}{379.796}{761.124}{e}%
\TextAt{9}{385.222}{761.124}{\Character{40}\Character{91}}%
\TextAt{13}{393.026}{761.124}{u}%
\TextAt{9}{399.688}{761.124}{\Character{93}\Character{40}}%
\TextAt{13}{407.492}{761.124}{X}%
\TextAt{9}{418.147}{761.124}{\Character{41}\Character{41}}%
\TextAt{9}{432.111}{761.124}{\Character{61}}%
\TextAt{9}{446.086}{761.124}{\Character{40}1}%
\TextAt{9}{459.763}{761.124}{\Character{43}}%
\TextAt{13}{472.196}{761.124}{X}%
\TextAt{9}{482.851}{761.124}{\Character{41}}%
\TextAt{21}{487.403}{765.462}{up}%
\TextAt{16}{500.349}{761.124}{\Character{0}}%
\TextAt{9}{512.929}{761.124}{1}%
\TextAt{9}{523.653}{761.124}{\Character{61}}%
\TextAt{9}{79.226}{744.078}{\Character{91}}%
\TextAt{13}{82.478}{744.078}{u}%
\TextAt{9}{89.14}{744.078}{\Character{93}\Character{40}}%
\TextAt{9}{97.353}{744.078}{\Character{22}}%
\TextAt{13}{96.944}{744.078}{e}%
\TextAt{9}{102.37}{744.078}{\Character{40}}%
\TextAt{13}{106.922}{744.078}{X}%
\TextAt{9}{117.577}{744.078}{\Character{41}\Character{41}}%
\TextAt{8}{126.682}{744.078}{.}%
\TextAt{8}{135.13}{744.078}{By}%
\TextAt{8}{153.493}{744.078}{(2.3),}%
\TextAt{8}{184.72}{744.078}{there}%
\TextAt{8}{214.632}{744.078}{is}%
\TextAt{8}{226.396}{744.078}{a}%
\TextAt{8}{236.151}{744.078}{power}%
\TextAt{8}{270.295}{744.078}{series}%
\TextAt{13}{301.63}{744.078}{\Character{18}}%
\TextAt{9}{307.41}{744.078}{\Character{40}}%
\TextAt{13}{311.962}{744.078}{X}%
\TextAt{9}{322.618}{744.078}{\Character{41}}%
\NativeMathAt{330.491}{744.078}{11.955}{\in}%
\TextAt{14}{341.782}{744.078}{o}%
\TextAt{30}{349.596}{742.439}{\Character{98}}%
\TextAt{21}{347.626}{740.425}{K}%
\TextAt{9}{355.829}{744.078}{\Character{91}\Character{91}}%
\TextAt{13}{362.332}{744.078}{X}%
\TextAt{9}{372.987}{744.078}{\Character{93}\Character{93}}%
\TextAt{8}{383.388}{744.078}{such}%
\TextAt{8}{409.785}{744.078}{that}%
\TextAt{13}{220.716}{715.077}{e}%
\TextAt{16}{228.798}{715.077}{\Character{14}}%
\TextAt{13}{237.432}{715.077}{\Character{18}}%
\TextAt{9}{246.534}{715.077}{\Character{61}}%
\TextAt{13}{258.959}{715.077}{\Character{18}}%
\TextAt{21}{264.739}{720.013}{\Character{39}}%
\TextAt{16}{273.412}{715.077}{\Character{14}}%
\TextAt{9}{282.483}{715.077}{\Character{22}}%
\TextAt{13}{282.047}{715.077}{e}%
\TextAt{8}{299.427}{715.077}{and}%
\TextAt{13}{330.248}{715.077}{\Character{18}}%
\TextAt{21}{336.028}{720.013}{\Character{39}}%
\TextAt{9}{345.365}{715.077}{\Character{61}}%
\TextAt{13}{357.791}{715.077}{\Character{18}}%
\TextAt{16}{366.228}{715.077}{\Character{14}}%
\TextAt{9}{374.862}{715.077}{\Character{91}}%
\TextAt{13}{378.114}{715.077}{u}%
\TextAt{9}{384.776}{715.077}{\Character{93}}%
\TextAt{13}{388.028}{715.077}{\Character{58}}%
\TextAt{8}{79.226}{686.075}{Substituting}%
\TextAt{8}{146.283}{686.075}{the}%
\TextAt{8}{166.439}{686.075}{prime}%
\TextAt{8}{199.603}{686.075}{element}%
\TextAt{13}{242.845}{686.075}{\Character{21}}%
\TextAt{9}{252.994}{686.075}{\Character{61}}%
\TextAt{13}{265.419}{686.075}{\Character{16}}%
\TextAt{16}{274.092}{686.075}{\Character{0}}%
\TextAt{9}{286.047}{686.075}{1}%
\TextAt{8}{295.797}{686.075}{of}%
\TextAt{13}{309.127}{686.075}{L}%
\TextAt{8}{317.092}{686.075}{,}%
\TextAt{8}{324.241}{686.075}{we}%
\TextAt{8}{341.468}{686.075}{obtain}%
\TextAt{8}{377.896}{686.075}{a}%
\TextAt{8}{387.651}{686.075}{prime}%
\TextAt{8}{420.815}{686.075}{element}%
\TextAt{8}{464.057}{686.075}{of}%
\TextAt{9}{477.387}{686.075}{\Character{6}}%
\TextAt{8}{489.738}{686.075}{by}%
\TextAt{13}{279.015}{657.074}{\Character{25}}%
\TextAt{17}{285.655}{655.281}{\Character{6}}%
\TextAt{9}{295.59}{657.074}{\Character{61}}%
\TextAt{13}{308.016}{657.074}{\Character{18}}%
\TextAt{9}{313.796}{657.074}{\Character{40}}%
\TextAt{13}{318.348}{657.074}{\Character{21}}%
\TextAt{9}{325.177}{657.074}{\Character{41}}%
\TextAt{13}{329.729}{657.074}{\Character{58}}%
\TextAt{8}{79.226}{628.073}{Indeed,}%
\TextAt{9}{120.507}{628.073}{\Character{91}}%
\TextAt{13}{123.759}{628.073}{u}%
\TextAt{9}{130.421}{628.073}{\Character{93}\Character{40}}%
\TextAt{13}{138.225}{628.073}{\Character{21}}%
\TextAt{21}{145.054}{632.412}{\Character{27}}%
\TextAt{9}{150.677}{628.073}{\Character{41}}%
\TextAt{9}{158.532}{628.073}{\Character{61}}%
\TextAt{9}{170.953}{628.073}{\Character{40}1}%
\TextAt{9}{183.996}{628.073}{\Character{43}}%
\TextAt{13}{195.799}{628.073}{\Character{21}}%
\TextAt{21}{202.628}{632.412}{\Character{27}}%
\TextAt{9}{208.251}{628.073}{\Character{41}}%
\TextAt{21}{212.804}{632.412}{u}%
\TextAt{16}{220.861}{628.073}{\Character{0}}%
\TextAt{9}{232.816}{628.073}{1}%
\TextAt{9}{241.974}{628.073}{\Character{61}}%
\TextAt{13}{254.416}{628.073}{\Character{16}}%
\TextAt{21}{260.431}{632.412}{\Character{27}}%
\TextAt{21}{265.548}{632.412}{u}%
\TextAt{16}{273.614}{628.073}{\Character{0}}%
\TextAt{9}{285.57}{628.073}{1}%
\TextAt{9}{294.728}{628.073}{\Character{61}}%
\TextAt{13}{307.169}{628.073}{\Character{16}}%
\TextAt{16}{315.841}{628.073}{\Character{0}}%
\TextAt{9}{327.796}{628.073}{1}%
\TextAt{9}{336.954}{628.073}{\Character{61}}%
\TextAt{13}{349.396}{628.073}{\Character{21}}%
\TextAt{8}{356.224}{628.073}{,}%
\TextAt{8}{363.373}{628.073}{and}%
\TextAt{8}{386.136}{628.073}{therefore}%
\TextAt{13}{207.288}{599.072}{\Character{25}}%
\TextAt{30}{214.567}{604.008}{\Character{101}}%
\TextAt{21}{214.357}{604.008}{\Character{27}}%
\TextAt{17}{213.928}{596.116}{\Character{6}}%
\TextAt{9}{223.863}{599.072}{\Character{61}}%
\TextAt{13}{236.289}{599.072}{\Character{18}}%
\TextAt{21}{242.069}{604.008}{\Character{39}}%
\TextAt{9}{248.085}{599.072}{\Character{40}}%
\TextAt{13}{252.638}{599.072}{\Character{21}}%
\TextAt{21}{259.466}{604.008}{\Character{27}}%
\TextAt{9}{265.09}{599.072}{\Character{41}}%
\TextAt{9}{272.945}{599.072}{\Character{61}}%
\TextAt{13}{285.388}{599.072}{\Character{18}}%
\TextAt{9}{291.169}{599.072}{\Character{40}\Character{91}}%
\TextAt{13}{298.973}{599.072}{u}%
\TextAt{9}{305.635}{599.072}{\Character{93}\Character{40}}%
\TextAt{13}{313.439}{599.072}{\Character{21}}%
\TextAt{21}{320.267}{604.008}{\Character{27}}%
\TextAt{9}{325.891}{599.072}{\Character{41}\Character{41}}%
\TextAt{9}{338.288}{599.072}{\Character{61}}%
\TextAt{13}{350.742}{599.072}{\Character{18}}%
\TextAt{9}{356.522}{599.072}{\Character{40}}%
\TextAt{13}{361.075}{599.072}{\Character{21}}%
\TextAt{9}{367.903}{599.072}{\Character{41}}%
\TextAt{9}{375.758}{599.072}{\Character{61}}%
\TextAt{13}{388.202}{599.072}{\Character{25}}%
\TextAt{17}{394.842}{597.279}{\Character{6}}%
\TextAt{13}{401.456}{599.072}{\Character{59}}%
\TextAt{8}{79.226}{570.071}{i.e.,}%
\TextAt{13}{101.331}{570.071}{\Character{25}}%
\TextAt{17}{107.972}{568.278}{\Character{6}}%
\NativeMathAt{117.907}{570.071}{11.955}{\in}%
\TextAt{9}{129.198}{570.071}{\Character{6}}%
\TextAt{8}{137.652}{570.071}{.}%
\TextAt{8}{146.1}{570.071}{We}%
\TextAt{8}{166.245}{570.071}{will}%
\TextAt{8}{188.35}{570.071}{show}%
\TextAt{8}{217.353}{570.071}{that}%
\TextAt{13}{219.319}{541.07}{P}%
\TextAt{9}{228.49}{541.07}{\Character{40}}%
\TextAt{13}{233.042}{541.07}{X}%
\TextAt{9}{243.697}{541.07}{\Character{41}}%
\TextAt{9}{251.552}{541.07}{\Character{61}}%
\TextAt{13}{263.996}{541.07}{e}%
\TextAt{21}{269.422}{546.006}{n}%
\TextAt{3}{274.56}{546.006}{\Character{0}}%
\TextAt{17}{281.146}{546.006}{1}%
\TextAt{9}{285.879}{541.07}{\Character{40}}%
\TextAt{13}{290.431}{541.07}{X}%
\TextAt{9}{301.086}{541.07}{\Character{41}}%
\TextAt{21}{305.638}{546.006}{p}%
\TextAt{3}{309.901}{546.006}{\Character{0}}%
\TextAt{17}{316.488}{546.006}{1}%
\TextAt{9}{323.877}{541.07}{\Character{43}}%
\TextAt{13}{335.638}{541.07}{up}%
\NativeMathAt{351.496}{541.07}{11.955}{\in}%
\TextAt{18}{362.787}{541.07}{Z}%
\TextAt{21}{370.757}{539.276}{p}%
\TextAt{9}{375.518}{541.07}{\Character{91}}%
\TextAt{13}{378.77}{541.07}{X}%
\TextAt{9}{389.425}{541.07}{\Character{93}}%
\TextAt{8}{79.226}{512.069}{is}%
\TextAt{8}{92.258}{512.069}{the}%
\TextAt{8}{113.683}{512.069}{minimal}%
\TextAt{8}{160.477}{512.069}{polynomial}%
\TextAt{8}{222.239}{512.069}{of}%
\TextAt{13}{236.837}{512.069}{\Character{25}}%
\TextAt{17}{243.477}{510.275}{\Character{6}}%
\TextAt{8}{250.091}{512.069}{,}%
\TextAt{8}{258.826}{512.069}{where}%
\TextAt{13}{293.904}{512.069}{e}%
\TextAt{21}{299.329}{516.407}{i}%
\TextAt{9}{302.711}{512.069}{\Character{40}}%
\TextAt{13}{307.263}{512.069}{X}%
\TextAt{9}{317.918}{512.069}{\Character{41}}%
\TextAt{8}{327.636}{512.069}{is}%
\TextAt{8}{340.668}{512.069}{defined}%
\TextAt{8}{382.25}{512.069}{by}%
\TextAt{13}{399.777}{512.069}{e}%
\TextAt{17}{405.203}{516.407}{0}%
\TextAt{9}{409.935}{512.069}{\Character{40}}%
\TextAt{13}{414.487}{512.069}{X}%
\TextAt{9}{425.142}{512.069}{\Character{41}}%
\TextAt{9}{435.161}{512.069}{\Character{61}}%
\TextAt{13}{449.764}{512.069}{X}%
\TextAt{8}{465.585}{512.069}{and}%
\TextAt{13}{489.616}{512.069}{e}%
\TextAt{21}{495.041}{516.407}{i}%
\TextAt{9}{498.423}{512.069}{\Character{40}}%
\TextAt{13}{502.975}{512.069}{X}%
\TextAt{9}{513.63}{512.069}{\Character{41}}%
\TextAt{9}{523.649}{512.069}{\Character{61}}%
\TextAt{13}{79.226}{495.023}{e}%
\TextAt{9}{84.652}{495.023}{\Character{40}}%
\TextAt{13}{89.204}{495.023}{e}%
\TextAt{21}{94.629}{499.361}{i}%
\TextAt{3}{97.512}{499.361}{\Character{0}}%
\TextAt{17}{104.099}{499.361}{1}%
\TextAt{9}{108.831}{495.023}{\Character{40}}%
\TextAt{13}{113.384}{495.023}{X}%
\TextAt{9}{124.039}{495.023}{\Character{41}\Character{41}}%
\TextAt{8}{133.143}{495.023}{.}%
\TextAt{13}{141.458}{495.023}{P}%
\TextAt{9}{150.629}{495.023}{\Character{40}}%
\TextAt{13}{155.181}{495.023}{X}%
\TextAt{9}{165.836}{495.023}{\Character{41}}%
\TextAt{8}{173.884}{495.023}{is}%
\TextAt{8}{185.245}{495.023}{monic}%
\TextAt{8}{219.31}{495.023}{of}%
\TextAt{8}{232.237}{495.023}{degree}%
\TextAt{13}{268.251}{495.023}{p}%
\TextAt{21}{274.126}{499.361}{n}%
\TextAt{3}{279.264}{499.361}{\Character{0}}%
\TextAt{17}{285.85}{499.361}{1}%
\TextAt{9}{290.583}{495.023}{\Character{40}}%
\TextAt{13}{295.135}{495.023}{p}%
\TextAt{16}{302.844}{495.023}{\Character{0}}%
\TextAt{9}{313.977}{495.023}{1\Character{41}}%
\TextAt{8}{327.877}{495.023}{and}%
\TextAt{8}{350.237}{495.023}{irreducible}%
\TextAt{8}{407.71}{495.023}{by}%
\TextAt{8}{423.567}{495.023}{Eisenstein’s}%
\TextAt{8}{486.587}{495.023}{criterion,}%
\TextAt{8}{79.226}{477.977}{as}%
\TextAt{13}{93.536}{477.977}{e}%
\TextAt{9}{98.962}{477.977}{\Character{40}}%
\TextAt{13}{103.514}{477.977}{X}%
\TextAt{9}{114.169}{477.977}{\Character{41}}%
\NativeMathAt{122.042}{477.977}{11.955}{\equiv}%
\TextAt{13}{134.662}{477.977}{X}%
\TextAt{21}{145.317}{482.315}{p}%
\TextAt{8}{153.399}{477.977}{mod}%
\TextAt{13}{178.836}{477.977}{p}%
\TextAt{8}{184.712}{477.977}{,}%
\TextAt{8}{191.813}{477.977}{and}%
\TextAt{8}{214.516}{477.977}{so}%
\TextAt{13}{228.826}{477.977}{e}%
\TextAt{21}{234.251}{482.315}{n}%
\TextAt{3}{239.39}{482.315}{\Character{0}}%
\TextAt{17}{245.976}{482.315}{1}%
\TextAt{9}{250.708}{477.977}{\Character{40}}%
\TextAt{13}{255.261}{477.977}{X}%
\TextAt{9}{265.916}{477.977}{\Character{41}}%
\TextAt{21}{270.468}{482.315}{p}%
\TextAt{3}{274.731}{482.315}{\Character{0}}%
\TextAt{17}{281.317}{482.315}{1}%
\NativeMathAt{289.371}{477.977}{11.955}{\equiv}%
\TextAt{13}{301.99}{477.977}{X}%
\TextAt{21}{312.645}{482.315}{p}%
\TextAt{22}{316.908}{485.128}{n}%
\TextAt{23}{321.509}{485.128}{\Character{0}}%
\TextAt{24}{327.265}{485.128}{1}%
\TextAt{17}{331.416}{482.315}{\Character{40}}%
\TextAt{21}{334.709}{482.315}{p}%
\TextAt{3}{338.972}{482.315}{\Character{0}}%
\TextAt{17}{345.558}{482.315}{1\Character{41}}%
\TextAt{8}{356.905}{477.977}{mod}%
\TextAt{13}{382.343}{477.977}{p}%
\TextAt{8}{388.218}{477.977}{.}%
\TextAt{8}{396.646}{477.977}{Finally,}%
\TextAt{13}{438.704}{477.977}{e}%
\TextAt{21}{444.13}{482.315}{n}%
\TextAt{9}{449.766}{477.977}{\Character{40}}%
\TextAt{13}{454.318}{477.977}{X}%
\TextAt{9}{464.974}{477.977}{\Character{41}}%
\TextAt{9}{472.829}{477.977}{\Character{61}}%
\TextAt{13}{485.272}{477.977}{e}%
\TextAt{21}{490.698}{482.315}{n}%
\TextAt{3}{495.836}{482.315}{\Character{0}}%
\TextAt{17}{502.422}{482.315}{1}%
\TextAt{9}{507.155}{477.977}{\Character{40}}%
\TextAt{13}{511.707}{477.977}{X}%
\TextAt{9}{522.362}{477.977}{\Character{41}}%
\TextAt{16}{529.449}{477.977}{\Character{1}}%
\TextAt{9}{79.226}{460.931}{\Character{40}}%
\TextAt{13}{83.778}{460.931}{up}%
\TextAt{9}{98.973}{460.931}{\Character{43}}%
\TextAt{13}{110.734}{460.931}{e}%
\TextAt{21}{116.159}{465.269}{n}%
\TextAt{3}{121.298}{465.269}{\Character{0}}%
\TextAt{17}{127.884}{465.269}{1}%
\TextAt{9}{132.616}{460.931}{\Character{40}}%
\TextAt{13}{137.169}{460.931}{X}%
\TextAt{9}{147.824}{460.931}{\Character{41}}%
\TextAt{21}{152.376}{465.269}{p}%
\TextAt{3}{156.639}{465.269}{\Character{0}}%
\TextAt{17}{163.226}{465.269}{1}%
\TextAt{9}{167.958}{460.931}{\Character{41}}%
\TextAt{9}{175.813}{460.931}{\Character{61}}%
\TextAt{13}{188.256}{460.931}{e}%
\TextAt{21}{193.682}{465.269}{n}%
\TextAt{3}{198.82}{465.269}{\Character{0}}%
\TextAt{17}{205.407}{465.269}{1}%
\TextAt{9}{210.139}{460.931}{\Character{40}}%
\TextAt{13}{214.691}{460.931}{X}%
\TextAt{9}{225.346}{460.931}{\Character{41}}%
\TextAt{13}{229.899}{460.931}{P}%
\TextAt{9}{239.07}{460.931}{\Character{40}}%
\TextAt{13}{243.622}{460.931}{X}%
\TextAt{9}{254.277}{460.931}{\Character{41}}%
\TextAt{8}{258.83}{460.931}{,}%
\TextAt{8}{265.979}{460.931}{so}%
\TextAt{8}{280.349}{460.931}{that}%
\TextAt{13}{241.902}{431.93}{P}%
\TextAt{9}{251.074}{431.93}{\Character{40}}%
\TextAt{13}{255.626}{431.93}{\Character{25}}%
\TextAt{17}{262.266}{430.136}{\Character{6}}%
\TextAt{9}{268.88}{431.93}{\Character{41}}%
\TextAt{13}{273.433}{431.93}{e}%
\TextAt{21}{278.858}{436.866}{n}%
\TextAt{3}{283.996}{436.866}{\Character{0}}%
\TextAt{17}{290.583}{436.866}{1}%
\TextAt{9}{295.315}{431.93}{\Character{40}}%
\TextAt{13}{299.867}{431.93}{\Character{25}}%
\TextAt{17}{306.508}{430.136}{\Character{6}}%
\TextAt{9}{313.122}{431.93}{\Character{41}}%
\TextAt{9}{320.977}{431.93}{\Character{61}}%
\TextAt{13}{333.421}{431.93}{e}%
\TextAt{21}{338.846}{436.866}{n}%
\TextAt{9}{344.482}{431.93}{\Character{40}}%
\TextAt{13}{349.035}{431.93}{\Character{25}}%
\TextAt{17}{355.675}{430.136}{\Character{6}}%
\TextAt{9}{362.289}{431.93}{\Character{41}}%
\TextAt{13}{366.842}{431.93}{\Character{58}}%
\TextAt{8}{79.226}{402.928}{Since}%
\TextAt{13}{108.766}{402.928}{e}%
\TextAt{21}{114.191}{407.267}{i}%
\TextAt{9}{117.572}{402.928}{\Character{40}}%
\TextAt{13}{122.125}{402.928}{\Character{25}}%
\TextAt{17}{128.765}{401.135}{\Character{6}}%
\TextAt{9}{135.379}{402.928}{\Character{41}}%
\TextAt{9}{143.234}{402.928}{\Character{61}}%
\TextAt{13}{155.678}{402.928}{e}%
\TextAt{21}{161.103}{407.267}{i}%
\TextAt{9}{164.485}{402.928}{\Character{40}}%
\TextAt{13}{169.037}{402.928}{\Character{18}}%
\TextAt{9}{174.817}{402.928}{\Character{40}}%
\TextAt{13}{179.37}{402.928}{\Character{21}}%
\TextAt{9}{186.198}{402.928}{\Character{41}\Character{41}}%
\TextAt{9}{198.596}{402.928}{\Character{61}}%
\TextAt{13}{211.049}{402.928}{\Character{18}}%
\TextAt{21}{216.829}{407.267}{\Character{39}}%
\TextAt{22}{222.348}{410.08}{i}%
\TextAt{9}{226.007}{402.928}{\Character{40}}%
\TextAt{9}{230.98}{402.928}{\Character{22}}%
\TextAt{13}{230.56}{402.928}{e}%
\TextAt{21}{235.985}{407.267}{i}%
\TextAt{9}{239.366}{402.928}{\Character{40}}%
\TextAt{13}{243.919}{402.928}{\Character{21}}%
\TextAt{9}{250.747}{402.928}{\Character{41}\Character{41}}%
\TextAt{9}{263.144}{402.928}{\Character{61}}%
\TextAt{13}{275.598}{402.928}{\Character{18}}%
\TextAt{21}{281.379}{407.267}{\Character{39}}%
\TextAt{22}{286.897}{410.08}{i}%
\TextAt{9}{290.557}{402.928}{\Character{40}\Character{40}1}%
\TextAt{9}{306.063}{402.928}{\Character{43}}%
\TextAt{13}{315.769}{402.928}{\Character{21}}%
\TextAt{9}{322.598}{402.928}{\Character{41}}%
\TextAt{21}{327.15}{407.267}{p}%
\TextAt{22}{331.413}{410.08}{i}%
\TextAt{16}{335.648}{402.928}{\Character{0}}%
\TextAt{9}{345.521}{402.928}{1\Character{41}}%
\TextAt{9}{359.222}{402.928}{\Character{61}}%
\TextAt{13}{371.673}{402.928}{\Character{18}}%
\TextAt{21}{377.453}{407.267}{\Character{39}}%
\TextAt{22}{382.972}{410.08}{i}%
\TextAt{9}{386.632}{402.928}{\Character{40}}%
\TextAt{13}{391.184}{402.928}{\Character{16}}%
\TextAt{21}{397.199}{407.267}{p}%
\TextAt{22}{401.462}{410.08}{i}%
\TextAt{16}{405.697}{402.928}{\Character{0}}%
\TextAt{9}{415.571}{402.928}{1\Character{41}}%
\TextAt{8}{425.976}{402.928}{,}%
\TextAt{8}{432.311}{402.928}{we}%
\TextAt{8}{448.521}{402.928}{have}%
\TextAt{13}{474.497}{402.928}{e}%
\TextAt{21}{479.923}{407.267}{n}%
\TextAt{9}{485.559}{402.928}{\Character{40}}%
\TextAt{13}{490.112}{402.928}{\Character{25}}%
\TextAt{17}{496.752}{401.135}{\Character{6}}%
\TextAt{9}{503.366}{402.928}{\Character{41}}%
\TextAt{9}{511.221}{402.928}{\Character{61}}%
\TextAt{9}{523.642}{402.928}{0}%
\TextAt{8}{529.518}{402.928}{,}%
\TextAt{13}{79.226}{385.882}{e}%
\TextAt{21}{84.652}{390.221}{n}%
\TextAt{3}{89.79}{390.221}{\Character{0}}%
\TextAt{17}{96.376}{390.221}{1}%
\TextAt{9}{101.109}{385.882}{\Character{40}}%
\TextAt{13}{105.661}{385.882}{\Character{25}}%
\TextAt{17}{112.301}{384.089}{\Character{6}}%
\TextAt{9}{118.915}{385.882}{\Character{41}}%
\NativeMathAt{126.789}{385.882}{11.955}{\neq}%
\TextAt{9}{139.199}{385.882}{0}%
\TextAt{8}{145.067}{385.882}{,}%
\TextAt{8}{152.216}{385.882}{and}%
\TextAt{8}{174.979}{385.882}{thus}%
\TextAt{13}{200.73}{385.882}{P}%
\TextAt{9}{209.901}{385.882}{\Character{40}}%
\TextAt{13}{214.454}{385.882}{\Character{25}}%
\TextAt{17}{221.094}{384.089}{\Character{6}}%
\TextAt{9}{227.708}{385.882}{\Character{41}}%
\TextAt{9}{235.563}{385.882}{\Character{61}}%
\TextAt{9}{247.984}{385.882}{0}%
\TextAt{8}{253.86}{385.882}{.}%
\TextAt{8}{97.159}{365.848}{Observing}%
\TextAt{8}{152.834}{365.848}{that}%
\TextAt{13}{178.203}{365.848}{N}%
\TextAt{21}{187.577}{363.988}{L}%
\TextAt{3}{193.336}{363.988}{\Character{106}}%
\TextAt{21}{195.688}{363.988}{K}%
\TextAt{9}{203.891}{365.848}{\Character{40}}%
\TextAt{13}{208.443}{365.848}{\Character{16}}%
\TextAt{16}{217.115}{365.848}{\Character{0}}%
\TextAt{9}{229.071}{365.848}{1\Character{41}}%
\TextAt{9}{242.772}{365.848}{\Character{61}}%
\TextAt{9}{255.193}{365.848}{\Character{40}}%
\TextAt{16}{259.775}{365.848}{\Character{0}}%
\TextAt{9}{269.073}{365.848}{1\Character{41}}%
\TextAt{21}{279.478}{370.186}{d}%
\TextAt{13}{284.334}{365.848}{p}%
\TextAt{8}{290.209}{365.848}{,}%
\TextAt{13}{297.358}{365.848}{d}%
\TextAt{9}{306.762}{365.848}{\Character{61}}%
\TextAt{9}{319.172}{365.848}{\Character{91}}%
\TextAt{13}{322.439}{365.848}{L}%
\TextAt{9}{333.724}{365.848}{\Character{58}}%
\TextAt{13}{340.297}{365.848}{K}%
\TextAt{9}{351.114}{365.848}{\Character{93}}%
\TextAt{8}{358.263}{365.848}{(see}%
\TextAt{8}{381.731}{365.848}{chap.}%
\TextAt{8}{413.923}{365.848}{II,}%
\TextAt{8}{429.835}{365.848}{(7.13)),}%
\TextAt{8}{471.475}{365.848}{we}%
\TextAt{8}{488.702}{365.848}{obtain}%
\TextAt{13}{173.732}{336.847}{N}%
\TextAt{17}{183.106}{334.987}{\Character{6}}%
\TextAt{3}{189.222}{334.987}{\Character{106}}%
\TextAt{21}{191.574}{334.987}{K}%
\TextAt{9}{199.777}{336.847}{\Character{40}}%
\TextAt{13}{204.329}{336.847}{\Character{25}}%
\TextAt{17}{210.969}{335.053}{\Character{6}}%
\TextAt{9}{217.583}{336.847}{\Character{41}}%
\TextAt{9}{225.438}{336.847}{\Character{61}}%
\TextAt{9}{237.859}{336.847}{\Character{40}}%
\TextAt{16}{242.434}{336.847}{\Character{0}}%
\TextAt{9}{251.733}{336.847}{1\Character{41}}%
\TextAt{21}{262.138}{341.783}{d}%
\TextAt{13}{266.994}{336.847}{P}%
\TextAt{9}{276.165}{336.847}{\Character{40}0\Character{41}}%
\TextAt{9}{294.409}{336.847}{\Character{61}}%
\TextAt{9}{306.83}{336.847}{\Character{40}}%
\TextAt{16}{311.421}{336.847}{\Character{0}}%
\TextAt{9}{320.72}{336.847}{1\Character{41}}%
\TextAt{21}{331.125}{341.783}{d}%
\TextAt{13}{335.98}{336.847}{pu}%
\NativeMathAt{351.839}{336.847}{11.955}{\equiv}%
\TextAt{13}{364.458}{336.847}{u}%
\TextAt{8}{374.441}{336.847}{mod}%
\TextAt{13}{399.879}{336.847}{N}%
\TextAt{21}{409.253}{334.987}{L}%
\TextAt{3}{415.012}{334.987}{\Character{106}}%
\TextAt{21}{417.364}{334.987}{K}%
\TextAt{13}{425.567}{336.847}{L}%
\TextAt{3}{433.532}{341.783}{\Character{3}}%
\TextAt{8}{79.226}{307.845}{and}%
\TextAt{8}{101.989}{307.845}{therefore}%
\TextAt{13}{151.089}{307.845}{r}%
\TextAt{21}{156.364}{305.986}{L}%
\TextAt{3}{162.123}{305.986}{\Character{106}}%
\TextAt{21}{164.475}{305.986}{K}%
\TextAt{9}{172.678}{307.845}{\Character{40}}%
\TextAt{13}{177.23}{307.845}{\Character{27}}%
\TextAt{9}{184.313}{307.845}{\Character{41}}%
\NativeMathAt{192.186}{307.845}{11.955}{\equiv}%
\TextAt{13}{204.805}{307.845}{u}%
\TextAt{8}{214.788}{307.845}{mod}%
\TextAt{13}{240.226}{307.845}{N}%
\TextAt{21}{249.6}{305.986}{L}%
\TextAt{3}{255.359}{305.986}{\Character{106}}%
\TextAt{21}{257.712}{305.986}{K}%
\TextAt{13}{265.914}{307.845}{L}%
\TextAt{3}{273.879}{312.184}{\Character{3}}%
\TextAt{8}{278.611}{307.845}{,}%
\TextAt{8}{285.76}{307.845}{i.e.,}%
\TextAt{9}{307.865}{307.845}{\Character{40}}%
\TextAt{13}{312.418}{307.845}{u\Character{59}}%
\TextAt{13}{324.302}{307.845}{L}%
\TextAt{16}{332.289}{307.845}{\Character{106}}%
\TextAt{13}{335.61}{307.845}{K}%
\TextAt{9}{346.427}{307.845}{\Character{41}}%
\TextAt{9}{354.282}{307.845}{\Character{61}}%
\TextAt{9}{366.703}{307.845}{\Character{40}}%
\TextAt{13}{371.278}{307.845}{a\Character{59}}%
\TextAt{13}{382.647}{307.845}{L}%
\TextAt{16}{390.632}{307.845}{\Character{106}}%
\TextAt{13}{393.953}{307.845}{K}%
\TextAt{9}{404.771}{307.845}{\Character{41}}%
\TextAt{9}{412.626}{307.845}{\Character{61}}%
\TextAt{13}{425.069}{307.845}{\Character{27}}%
\TextAt{8}{432.152}{307.845}{,}%
\TextAt{8}{439.301}{307.845}{as}%
\TextAt{8}{453.671}{307.845}{required.}%
\Vector{q 0 G 0 g 0.399 w 452.334 -466.155 m 452.334 -458.085 l S Q}%
\Vector{q 0 G 0 g 0.399 w 452.534 -458.284 m 459.707 -458.284 l S Q}%
\Vector{q 0 G 0 g 0.399 w 452.534 -465.955 m 459.707 -465.955 l S Q}%
\Vector{q 0 G 0 g 0.399 w 459.906 -466.155 m 459.906 -458.085 l S Q}%
\TextAt{8}{97.159}{278.844}{In}%
\TextAt{8}{112.027}{278.844}{order}%
\TextAt{8}{142.844}{278.844}{to}%
\TextAt{8}{157.401}{278.844}{really}%
\TextAt{8}{189.844}{278.844}{understand}%
\TextAt{8}{251.289}{278.844}{this}%
\TextAt{8}{274.359}{278.844}{proof}%
\TextAt{8}{305.176}{278.844}{of}%
\TextAt{8}{318.753}{278.844}{theorem}%
\TextAt{8}{364.525}{278.844}{(2.4),}%
\TextAt{8}{396.061}{278.844}{one}%
\TextAt{8}{417.768}{278.844}{has}%
\TextAt{8}{438.889}{278.844}{to}%
\TextAt{8}{453.446}{278.844}{read}%
\TextAt{8}{479.708}{278.844}{§4.}%
\TextAt{8}{500.435}{278.844}{Let}%
\TextAt{8}{521.651}{278.844}{us}%
\TextAt{8}{79.226}{261.798}{note}%
\TextAt{8}{105.16}{261.798}{that}%
\TextAt{8}{130.447}{261.798}{one}%
\TextAt{8}{151.826}{261.798}{would}%
\TextAt{8}{185.889}{261.798}{get}%
\TextAt{8}{205.318}{261.798}{a}%
\TextAt{8}{214.993}{261.798}{direct,}%
\TextAt{8}{251.343}{261.798}{purely}%
\TextAt{8}{287.355}{261.798}{algebraic}%
\TextAt{8}{336.709}{261.798}{proof,}%
\TextAt{8}{370.465}{261.798}{if}%
\TextAt{8}{381.108}{261.798}{one}%
\TextAt{8}{402.486}{261.798}{could}%
\TextAt{8}{433.62}{261.798}{show}%
\TextAt{8}{462.542}{261.798}{without}%
\TextAt{8}{506.038}{261.798}{using}%
\TextAt{8}{79.226}{244.752}{the}%
\TextAt{8}{99.141}{244.752}{power}%
\TextAt{8}{133.044}{244.752}{series}%
\TextAt{13}{164.138}{244.752}{\Character{18}}%
\TextAt{8}{173.575}{244.752}{that}%
\TextAt{8}{198.702}{244.752}{the}%
\TextAt{8}{218.618}{244.752}{splitting}%
\TextAt{8}{264.619}{244.752}{field}%
\TextAt{8}{289.735}{244.752}{of}%
\TextAt{8}{302.824}{244.752}{the}%
\TextAt{8}{322.739}{244.752}{polynomial}%
\TextAt{13}{382.991}{244.752}{e}%
\TextAt{21}{388.417}{249.091}{n}%
\TextAt{9}{394.053}{244.752}{\Character{40}}%
\TextAt{13}{398.606}{244.752}{X}%
\TextAt{9}{409.261}{244.752}{\Character{41}}%
\TextAt{8}{417.469}{244.752}{is}%
\TextAt{8}{428.992}{244.752}{abelian,}%
\TextAt{8}{472.699}{244.752}{and}%
\TextAt{8}{495.22}{244.752}{that}%
\TextAt{8}{520.348}{244.752}{its}%
\TextAt{8}{79.226}{227.706}{elements}%
\TextAt{8}{126.98}{227.706}{are}%
\TextAt{8}{146.388}{227.706}{all}%
\TextAt{8}{162.544}{227.706}{fixed}%
\TextAt{8}{190.727}{227.706}{under}%
\TextAt{20}{224.008}{227.706}{\Character{101}}%
\TextAt{13}{223.787}{227.706}{\Character{27}}%
\TextAt{9}{234.191}{227.706}{\Character{61}}%
\TextAt{13}{246.616}{227.706}{\Character{27}}%
\TextAt{13}{253.681}{227.706}{\Character{39}}%
\TextAt{21}{261.367}{225.913}{L}%
\TextAt{8}{267.624}{227.706}{.}%
\TextAt{8}{276.038}{227.706}{This}%
\TextAt{8}{302.655}{227.706}{splitting}%
\TextAt{8}{348.795}{227.706}{field}%
\TextAt{8}{374.049}{227.706}{would}%
\TextAt{8}{408.09}{227.706}{then}%
\TextAt{8}{434.647}{227.706}{have}%
\TextAt{8}{461.539}{227.706}{to}%
\TextAt{8}{475.746}{227.706}{be}%
\TextAt{8}{491.567}{227.706}{equal}%
\TextAt{8}{522.357}{227.706}{to}%
\TextAt{8}{79.226}{210.66}{the}%
\TextAt{8}{99.581}{210.66}{field}%
\TextAt{9}{125.136}{210.66}{\Character{6}}%
\TextAt{8}{137.686}{210.66}{and}%
\TextAt{8}{160.647}{210.66}{every}%
\TextAt{8}{191.738}{210.66}{zero}%
\TextAt{8}{216.648}{210.66}{of}%
\TextAt{13}{230.176}{210.66}{P}%
\TextAt{9}{239.347}{210.66}{\Character{40}}%
\TextAt{13}{243.9}{210.66}{X}%
\TextAt{9}{254.555}{210.66}{\Character{41}}%
\TextAt{9}{262.756}{210.66}{\Character{61}}%
\TextAt{13}{275.53}{210.66}{e}%
\TextAt{21}{280.955}{214.999}{n}%
\TextAt{9}{286.591}{210.66}{\Character{40}}%
\TextAt{13}{291.144}{210.66}{X}%
\TextAt{9}{301.799}{210.66}{\Character{41}\Character{47}}%
\TextAt{13}{312.204}{210.66}{e}%
\TextAt{21}{317.63}{214.999}{n}%
\TextAt{3}{322.768}{214.999}{\Character{0}}%
\TextAt{17}{329.354}{214.999}{1}%
\TextAt{9}{334.087}{210.66}{\Character{40}}%
\TextAt{13}{338.639}{210.66}{X}%
\TextAt{9}{349.294}{210.66}{\Character{41}}%
\TextAt{8}{357.942}{210.66}{would}%
\TextAt{8}{392.285}{210.66}{have}%
\TextAt{8}{419.478}{210.66}{to}%
\TextAt{8}{433.986}{210.66}{be}%
\TextAt{8}{450.109}{210.66}{a}%
\TextAt{8}{460.063}{210.66}{prime}%
\TextAt{8}{493.425}{210.66}{element}%
\TextAt{13}{79.226}{193.614}{\Character{25}}%
\TextAt{17}{85.866}{191.821}{\Character{6}}%
\NativeMathAt{97.125}{193.614}{11.955}{\in}%
\TextAt{9}{109.74}{193.614}{\Character{6}}%
\TextAt{8}{122.869}{193.614}{such}%
\TextAt{8}{150.043}{193.614}{that}%
\TextAt{13}{176.188}{193.614}{N}%
\TextAt{17}{185.562}{191.755}{\Character{6}}%
\TextAt{3}{191.678}{191.755}{\Character{106}}%
\TextAt{21}{194.031}{191.755}{K}%
\TextAt{9}{202.233}{193.614}{\Character{40}}%
\TextAt{13}{206.786}{193.614}{\Character{25}}%
\TextAt{17}{213.426}{191.821}{\Character{6}}%
\TextAt{9}{220.04}{193.614}{\Character{41}}%
\NativeMathAt{229.237}{193.614}{11.955}{\equiv}%
\TextAt{13}{243.18}{193.614}{u}%
\TextAt{8}{253.164}{193.614}{mod}%
\TextAt{13}{278.602}{193.614}{N}%
\TextAt{21}{287.975}{191.755}{L}%
\TextAt{3}{293.734}{191.755}{\Character{106}}%
\TextAt{21}{296.087}{191.755}{K}%
\TextAt{13}{304.289}{193.614}{L}%
\TextAt{3}{312.254}{197.953}{\Character{3}}%
\TextAt{8}{316.986}{193.614}{,}%
\TextAt{8}{325.107}{193.614}{in}%
\TextAt{8}{339.536}{193.614}{which}%
\TextAt{8}{373.8}{193.614}{case}%
\TextAt{13}{399.347}{193.614}{r}%
\TextAt{21}{404.623}{191.755}{L}%
\TextAt{3}{410.382}{191.755}{\Character{106}}%
\TextAt{21}{412.734}{191.755}{K}%
\TextAt{9}{420.937}{193.614}{\Character{40}}%
\TextAt{13}{425.489}{193.614}{\Character{27}}%
\TextAt{9}{432.572}{193.614}{\Character{41}}%
\NativeMathAt{441.769}{193.614}{11.955}{\equiv}%
\TextAt{13}{455.712}{193.614}{u}%
\TextAt{8}{465.695}{193.614}{mod}%
\TextAt{13}{491.133}{193.614}{N}%
\TextAt{21}{500.507}{191.755}{L}%
\TextAt{3}{506.266}{191.755}{\Character{106}}%
\TextAt{21}{508.618}{191.755}{K}%
\TextAt{13}{516.821}{193.614}{L}%
\TextAt{3}{524.785}{197.953}{\Character{3}}%
\TextAt{8}{529.518}{193.614}{,}%
\TextAt{8}{79.226}{176.568}{and}%
\TextAt{8}{101.989}{176.568}{so}%
\TextAt{9}{116.359}{176.568}{\Character{40}}%
\TextAt{13}{120.911}{176.568}{u\Character{59}}%
\TextAt{13}{132.795}{176.568}{L}%
\TextAt{16}{140.782}{176.568}{\Character{106}}%
\TextAt{13}{144.103}{176.568}{K}%
\TextAt{9}{154.921}{176.568}{\Character{41}}%
\TextAt{9}{162.776}{176.568}{\Character{61}}%
\TextAt{13}{175.22}{176.568}{\Character{27}}%
\TextAt{8}{182.302}{176.568}{.}%
\TextAt{12}{79.226}{150.556}{Exercise}%
\TextAt{12}{132.339}{150.556}{V.2.1.}%
\TextAt{8}{171.009}{150.556}{The}%
\TextAt{13}{194.986}{150.556}{p}%
\TextAt{8}{200.862}{150.556}{-class}%
\TextAt{8}{232.12}{150.556}{field}%
\TextAt{8}{257.4}{150.556}{theory}%
\TextAt{9}{294.074}{150.556}{\Character{40}}%
\TextAt{13}{298.626}{150.556}{d}%
\TextAt{9}{308.03}{150.556}{\Character{58}}%
\TextAt{13}{314.602}{150.556}{G}%
\TextAt{25}{323.836}{148.763}{Q}%
\TextAt{22}{330.035}{147.767}{p}%
\NativeMathAt{338.192}{150.556}{11.955}{\rightarrow}%
\TextAt{18}{353.468}{150.556}{Z}%
\TextAt{21}{361.439}{148.763}{p}%
\TextAt{13}{366.199}{150.556}{\Character{59}}%
\TextAt{13}{375.249}{150.556}{v}%
\TextAt{9}{384.674}{150.556}{\Character{58}}%
\TextAt{18}{391.246}{150.556}{Q}%
\TextAt{3}{400.545}{154.895}{\Character{3}}%
\TextAt{21}{400.545}{147.601}{p}%
\NativeMathAt{408.626}{150.556}{11.955}{\rightarrow}%
\TextAt{18}{423.902}{150.556}{Z}%
\TextAt{9}{431.873}{150.556}{\Character{41}}%
\TextAt{8}{440.246}{150.556}{for}%
\TextAt{8}{458.054}{150.556}{the}%
\TextAt{8}{478.134}{150.556}{unramified}%
\TextAt{18}{79.226}{133.51}{Z}%
\TextAt{21}{87.196}{131.717}{p}%
\TextAt{8}{91.957}{133.51}{-extension}%
\TextAt{8}{148.549}{133.51}{of}%
\TextAt{18}{162.808}{133.51}{Q}%
\TextAt{21}{172.107}{131.717}{p}%
\TextAt{8}{176.868}{133.51}{,}%
\TextAt{8}{185.178}{133.51}{and}%
\TextAt{8}{208.87}{133.51}{the}%
\TextAt{13}{229.955}{133.51}{p}%
\TextAt{8}{235.83}{133.51}{-class}%
\TextAt{8}{268.094}{133.51}{field}%
\TextAt{8}{294.379}{133.51}{theory}%
\TextAt{9}{332.059}{133.51}{\Character{40}}%
\TextAt{13}{336.611}{133.51}{d}%
\TextAt{9}{347.597}{133.51}{\Character{58}}%
\TextAt{13}{355.753}{133.51}{G}%
\TextAt{25}{364.987}{131.717}{Q}%
\TextAt{22}{371.186}{130.721}{p}%
\NativeMathAt{380.926}{133.51}{11.955}{\rightarrow}%
\TextAt{18}{397.785}{133.51}{Z}%
\TextAt{21}{405.755}{131.717}{p}%
\TextAt{13}{410.516}{133.51}{\Character{59}}%
\TextAt{14}{420.587}{133.51}{v}%
\TextAt{9}{431.609}{133.51}{\Character{58}}%
\TextAt{18}{439.765}{133.51}{Q}%
\TextAt{3}{449.063}{137.849}{\Character{3}}%
\TextAt{21}{449.063}{130.555}{p}%
\NativeMathAt{458.728}{133.51}{11.955}{\rightarrow}%
\TextAt{18}{475.587}{133.51}{Z}%
\TextAt{21}{483.557}{131.717}{p}%
\TextAt{9}{488.318}{133.51}{\Character{41}}%
\TextAt{8}{497.697}{133.51}{for}%
\TextAt{8}{516.511}{133.51}{the}%
\TextAt{8}{79.226}{116.464}{cyclotomic}%
\TextAt{18}{138.316}{116.464}{Z}%
\TextAt{21}{146.286}{114.671}{p}%
\TextAt{8}{151.047}{116.464}{-extension}%
\TextAt{8}{207.59}{116.464}{of}%
\TextAt{18}{221.8}{116.464}{Q}%
\TextAt{21}{231.098}{114.671}{p}%
\TextAt{8}{240.636}{116.464}{(see}%
\TextAt{8}{264.984}{116.464}{§1,}%
\TextAt{8}{284.77}{116.464}{exercise}%
\TextAt{8}{328.951}{116.464}{7)}%
\TextAt{8}{344.141}{116.464}{yield}%
\TextAt{8}{373.307}{116.464}{the}%
\TextAt{8}{394.343}{116.464}{same}%
\TextAt{8}{424.549}{116.464}{norm}%
\TextAt{8}{455.998}{116.464}{residue}%
\TextAt{8}{496.605}{116.464}{symbol}%
\TextAt{9}{79.226}{99.418}{\Character{40}}%
\TextAt{13}{87.676}{99.418}{\Character{59}}%
\TextAt{13}{92.9}{99.418}{L}%
\TextAt{16}{100.885}{99.418}{\Character{106}}%
\TextAt{13}{104.205}{99.418}{K}%
\TextAt{9}{115.023}{99.418}{\Character{41}}%
\TextAt{8}{119.575}{99.418}{.}%
\TextAt{19}{97.159}{79.383}{Hint:}%
\TextAt{8}{128.657}{79.383}{Show}%
\TextAt{8}{159.502}{79.383}{that}%
\TextAt{8}{184.824}{79.383}{this}%
\TextAt{8}{207.6}{79.383}{statement}%
\TextAt{8}{261.925}{79.383}{is}%
\TextAt{8}{273.642}{79.383}{equivalent}%
\TextAt{8}{329.211}{79.383}{to}%
\TextAt{8}{343.474}{79.383}{formula}%
\TextAt{8}{386.358}{79.383}{(2.4):}%
\TextAt{9}{418.861}{79.383}{\Character{40}}%
\TextAt{13}{423.413}{79.383}{u\Character{59}}%
\TextAt{18}{435.32}{79.383}{Q}%
\TextAt{21}{444.618}{77.59}{p}%
\TextAt{9}{449.379}{79.383}{\Character{40}}%
\TextAt{13}{453.931}{79.383}{\Character{16}}%
\TextAt{9}{459.947}{79.383}{\Character{41}}%
\TextAt{16}{464.499}{79.383}{\Character{106}}%
\TextAt{18}{467.82}{79.383}{Q}%
\TextAt{21}{477.119}{77.59}{p}%
\TextAt{9}{481.88}{79.383}{\Character{41}}%
\TextAt{13}{486.432}{79.383}{\Character{16}}%
\TextAt{9}{495.768}{79.383}{\Character{61}}%
\TextAt{13}{508.194}{79.383}{\Character{16}}%
\TextAt{21}{514.209}{83.722}{u}%
\TextAt{23}{519.112}{86.535}{\Character{0}}%
\TextAt{24}{524.868}{86.535}{1}%
\TextAt{8}{529.518}{79.383}{.}%
\TextAt{8}{297.211}{43.033}{304}%
\Destination{78.23}{809.815}{nav.page.304}%
\Destination{79.226}{164.555}{nav.ex.5.2.1}%
\Destination{79.36}{293.855}{nav.print.332}%
\Hotlink{413.423}{363.655}{9.751}{11.047}{nav.chapter.2}%
\Hotlink{484.887}{276.655}{6.858}{11.047}{nav.section.5.4}%
\Hotlink{270.163}{114.255}{6.858}{11.047}{nav.section.5.1}%
\Hotlink{152.993}{741.855}{25.078}{11.047}{nav.number.5.2.3}%
\Hotlink{429.335}{363.655}{30.936}{11.047}{nav.number.2.7.13}%
\Hotlink{364.025}{276.655}{25.078}{11.047}{nav.number.5.2.4}%
\Hotlink{385.858}{77.155}{25.078}{11.047}{nav.number.5.2.4}%
}

% Page 317 (body page 305)
\Page{%
\TextAt{2}{79.226}{796.943}{V.}%
\TextAt{2}{94.073}{796.943}{Local}%
\TextAt{2}{123.615}{796.943}{Class}%
\TextAt{2}{152.208}{796.943}{Field}%
\TextAt{2}{179.939}{796.943}{Theory}%
\Vector{q 0 G 0 g 0.299 w 7.226 18.729 m 460.769 18.729 l S Q}%
\TextAt{12}{79.226}{761.124}{Exercise}%
\TextAt{12}{132.296}{761.124}{V.2.2.}%
\TextAt{8}{170.929}{761.124}{Let}%
\TextAt{13}{191.785}{761.124}{L}%
\TextAt{16}{199.75}{761.124}{\Character{106}}%
\TextAt{13}{203.071}{761.124}{K}%
\TextAt{8}{217.672}{761.124}{be}%
\TextAt{8}{233.483}{761.124}{a}%
\TextAt{8}{243.125}{761.124}{totally}%
\TextAt{8}{280.419}{761.124}{ramified}%
\TextAt{8}{325.831}{761.124}{Galois}%
\TextAt{8}{361.631}{761.124}{extension,}%
\TextAt{8}{416.558}{761.124}{and}%
\TextAt{8}{439.207}{761.124}{let}%
\TextAt{20}{456.985}{764.146}{\Character{98}}%
\TextAt{13}{455.998}{761.124}{L}%
\TextAt{8}{467.747}{761.124}{(resp.}%
\TextAt{20}{504.324}{764.146}{\Character{98}}%
\TextAt{13}{501.586}{761.124}{K}%
\TextAt{8}{512.404}{761.124}{)}%
\TextAt{8}{520.743}{761.124}{be}%
\TextAt{8}{79.226}{744.078}{the}%
\TextAt{8}{99.736}{744.078}{completion}%
\TextAt{8}{159.925}{744.078}{of}%
\TextAt{8}{173.608}{744.078}{the}%
\TextAt{8}{194.117}{744.078}{maximal}%
\TextAt{8}{242.279}{744.078}{unramified}%
\TextAt{8}{301.165}{744.078}{extension}%
\TextAt{20}{354.271}{747.1}{\Character{101}}%
\TextAt{13}{353.284}{744.078}{L}%
\TextAt{8}{365.499}{744.078}{(resp.}%
\TextAt{20}{403.174}{747.1}{\Character{101}}%
\TextAt{13}{400.436}{744.078}{K}%
\TextAt{8}{411.254}{744.078}{)}%
\TextAt{8}{420.059}{744.078}{of}%
\TextAt{13}{433.742}{744.078}{L}%
\TextAt{8}{445.958}{744.078}{(resp.}%
\TextAt{13}{480.894}{744.078}{K}%
\TextAt{8}{491.712}{744.078}{).}%
\TextAt{8}{505.775}{744.078}{Show}%
\TextAt{8}{79.226}{727.032}{that}%
\TextAt{13}{104.669}{727.032}{N}%
\TextAt{30}{114.805}{725.393}{\Character{98}}%
\TextAt{21}{114.043}{723.379}{L}%
\TextAt{3}{119.802}{723.379}{\Character{106}}%
\TextAt{30}{124.124}{725.393}{\Character{98}}%
\TextAt{21}{122.154}{723.379}{K}%
\TextAt{20}{131.343}{730.054}{\Character{98}}%
\TextAt{13}{130.357}{727.032}{L}%
\TextAt{3}{138.321}{731.37}{\Character{3}}%
\TextAt{9}{146.5}{727.032}{\Character{61}}%
\TextAt{20}{161.79}{730.054}{\Character{98}}%
\TextAt{13}{159.052}{727.032}{K}%
\TextAt{3}{169.869}{731.37}{\Character{3}}%
\TextAt{8}{174.602}{727.032}{,}%
\TextAt{8}{181.843}{727.032}{and}%
\TextAt{8}{204.68}{727.032}{that}%
\TextAt{8}{230.122}{727.032}{every}%
\TextAt{13}{261.088}{727.032}{y}%
\NativeMathAt{270.672}{727.032}{11.955}{\in}%
\TextAt{20}{283.075}{730.054}{\Character{98}}%
\TextAt{13}{282.088}{727.032}{L}%
\TextAt{3}{290.053}{731.37}{\Character{3}}%
\TextAt{8}{298.756}{727.032}{with}%
\TextAt{13}{325.49}{727.032}{N}%
\TextAt{30}{335.627}{725.393}{\Character{98}}%
\TextAt{21}{334.864}{723.379}{L}%
\TextAt{3}{340.623}{723.379}{\Character{106}}%
\TextAt{30}{344.946}{725.393}{\Character{98}}%
\TextAt{21}{342.975}{723.379}{K}%
\TextAt{9}{351.178}{727.032}{\Character{40}}%
\TextAt{13}{355.73}{727.032}{y}%
\TextAt{9}{361.867}{727.032}{\Character{41}}%
\TextAt{9}{369.853}{727.032}{\Character{61}}%
\TextAt{9}{382.394}{727.032}{1}%
\TextAt{8}{392.242}{727.032}{is}%
\TextAt{8}{404.079}{727.032}{of}%
\TextAt{8}{417.483}{727.032}{the}%
\TextAt{8}{437.714}{727.032}{form}%
\TextAt{13}{465.428}{727.032}{y}%
\TextAt{9}{475.011}{727.032}{\Character{61}}%
\TextAt{20}{487.562}{735.998}{\Character{81}}%
\TextAt{21}{498.853}{723.545}{i}%
\TextAt{13}{504.227}{727.032}{z}%
\TextAt{21}{510.198}{732.365}{\Character{27}}%
\TextAt{22}{515.037}{731.15}{i}%
\TextAt{3}{518.199}{732.365}{\Character{0}}%
\TextAt{17}{524.785}{732.365}{1}%
\TextAt{21}{509.665}{723.968}{i}%
\TextAt{8}{529.518}{727.032}{,}%
\TextAt{13}{79.226}{709.986}{\Character{27}}%
\TextAt{21}{85.88}{708.193}{i}%
\NativeMathAt{92.582}{709.986}{11.955}{\in}%
\TextAt{13}{103.873}{709.986}{G}%
\TextAt{9}{113.106}{709.986}{\Character{40}}%
\TextAt{13}{117.659}{709.986}{L}%
\TextAt{16}{125.623}{709.986}{\Character{106}}%
\TextAt{13}{128.944}{709.986}{K}%
\TextAt{9}{139.762}{709.986}{\Character{41}}%
\TextAt{8}{144.314}{709.986}{.}%
\TextAt{12}{79.226}{677.996}{Exercise}%
\TextAt{12}{131.134}{677.996}{V.2.3.}%
\TextAt{12}{168.757}{677.996}{(Theorem}%
\TextAt{12}{229.716}{677.996}{of}%
\TextAt{12}{243.75}{677.996}{Dwork).}%
\TextAt{8}{296.536}{677.996}{Let}%
\TextAt{13}{316.382}{677.996}{L}%
\TextAt{16}{324.346}{677.996}{\Character{106}}%
\TextAt{13}{327.667}{677.996}{K}%
\TextAt{8}{341.258}{677.996}{be}%
\TextAt{8}{356.059}{677.996}{a}%
\TextAt{8}{364.691}{677.996}{totally}%
\TextAt{8}{400.975}{677.996}{ramified}%
\TextAt{8}{445.377}{677.996}{abelian}%
\TextAt{8}{484.901}{677.996}{extension}%
\TextAt{8}{79.226}{660.95}{of}%
\TextAt{13}{92.668}{660.95}{p}%
\TextAt{8}{98.543}{660.95}{-adic}%
\TextAt{8}{127.264}{660.95}{number}%
\TextAt{8}{169.972}{660.95}{fields.}%
\TextAt{8}{204.829}{660.95}{Let}%
\TextAt{13}{225.91}{660.95}{x}%
\NativeMathAt{236.074}{660.95}{11.955}{\in}%
\TextAt{13}{247.555}{660.95}{K}%
\TextAt{3}{258.373}{665.288}{\Character{3}}%
\TextAt{8}{267.114}{660.95}{and}%
\TextAt{13}{289.989}{660.95}{y}%
\NativeMathAt{299.637}{660.95}{11.955}{\in}%
\TextAt{20}{312.105}{663.972}{\Character{98}}%
\TextAt{13}{311.118}{660.95}{L}%
\TextAt{3}{319.083}{665.288}{\Character{3}}%
\TextAt{8}{327.824}{660.95}{such}%
\TextAt{8}{354.333}{660.95}{that}%
\TextAt{13}{379.813}{660.95}{N}%
\TextAt{30}{389.95}{659.312}{\Character{98}}%
\TextAt{21}{389.187}{657.297}{L}%
\TextAt{3}{394.946}{657.297}{\Character{106}}%
\TextAt{30}{399.269}{659.312}{\Character{98}}%
\TextAt{21}{397.299}{657.297}{K}%
\TextAt{9}{405.501}{660.95}{\Character{40}}%
\TextAt{13}{410.054}{660.95}{y}%
\TextAt{9}{416.19}{660.95}{\Character{41}}%
\TextAt{9}{424.236}{660.95}{\Character{61}}%
\TextAt{13}{436.87}{660.95}{x}%
\TextAt{8}{443.522}{660.95}{.}%
\TextAt{8}{452.305}{660.95}{Let}%
\TextAt{13}{473.387}{660.95}{z}%
\TextAt{21}{478.824}{659.157}{i}%
\NativeMathAt{485.717}{660.95}{11.955}{\in}%
\TextAt{20}{498.185}{663.972}{\Character{98}}%
\TextAt{13}{497.198}{660.95}{L}%
\TextAt{3}{505.163}{665.288}{\Character{3}}%
\TextAt{8}{513.904}{660.95}{and}%
\TextAt{8}{79.226}{643.904}{choose}%
\TextAt{13}{116.359}{643.904}{\Character{27}}%
\TextAt{21}{123.012}{642.111}{i}%
\NativeMathAt{129.714}{643.904}{11.955}{\in}%
\TextAt{13}{141.005}{643.904}{G}%
\TextAt{9}{150.239}{643.904}{\Character{40}}%
\TextAt{13}{154.791}{643.904}{L}%
\TextAt{16}{162.756}{643.904}{\Character{106}}%
\TextAt{13}{166.077}{643.904}{K}%
\TextAt{9}{176.894}{643.904}{\Character{41}}%
\TextAt{8}{185.344}{643.904}{such}%
\TextAt{8}{211.741}{643.904}{that}%
\TextAt{13}{260.211}{626.858}{y}%
\TextAt{21}{266.348}{631.794}{\Character{39}}%
\TextAt{22}{271.866}{630.455}{K}%
\TextAt{3}{278.909}{631.794}{\Character{0}}%
\TextAt{17}{285.495}{631.794}{1}%
\TextAt{9}{293.548}{626.858}{\Character{61}}%
\TextAt{20}{305.974}{638.215}{\Character{89}}%
\TextAt{21}{312.17}{613.021}{i}%
\TextAt{13}{323.243}{626.858}{z}%
\TextAt{21}{329.213}{632.191}{\Character{27}}%
\TextAt{22}{334.052}{630.976}{i}%
\TextAt{3}{337.214}{632.191}{\Character{0}}%
\TextAt{17}{343.801}{632.191}{1}%
\TextAt{21}{328.68}{623.794}{i}%
\TextAt{13}{348.533}{626.858}{\Character{58}}%
\ErrataAt{79.19}{594.879}{455.81}{Putting $\sigma=\prod_i\sigma_i^{v_{\widehat L}(z_i)}$, one has $(x,L\mid K)=\sigma^{-1}$.}%
\TextAt{19}{97.159}{574.844}{Hint:}%
\TextAt{8}{128.681}{574.844}{See}%
\TextAt{8}{149.483}{574.844}{chap.}%
\TextAt{8}{181.674}{574.844}{IV,}%
\TextAt{8}{201.818}{574.844}{§5,}%
\TextAt{8}{220.504}{574.844}{exercise}%
\TextAt{8}{263.806}{574.844}{1.}%
\TextAt{12}{79.226}{542.854}{Exercise}%
\TextAt{12}{133.004}{542.854}{V.2.4.}%
\TextAt{8}{172.253}{542.854}{Deduce}%
\TextAt{8}{214.204}{542.854}{from}%
\TextAt{8}{242.346}{542.854}{exercises}%
\TextAt{8}{290.765}{542.854}{2}%
\TextAt{8}{301.022}{542.854}{and}%
\TextAt{8}{324.287}{542.854}{3}%
\TextAt{8}{334.545}{542.854}{the}%
\TextAt{8}{355.203}{542.854}{formula}%
\TextAt{9}{398.637}{542.854}{\Character{40}}%
\TextAt{13}{403.189}{542.854}{u\Character{59}}%
\TextAt{18}{415.095}{542.854}{Q}%
\TextAt{21}{424.394}{541.061}{p}%
\TextAt{9}{429.155}{542.854}{\Character{40}}%
\TextAt{13}{433.707}{542.854}{\Character{16}}%
\TextAt{9}{439.723}{542.854}{\Character{41}}%
\TextAt{16}{444.275}{542.854}{\Character{106}}%
\TextAt{18}{447.596}{542.854}{Q}%
\TextAt{21}{456.895}{541.061}{p}%
\TextAt{9}{461.655}{542.854}{\Character{41}}%
\TextAt{13}{466.208}{542.854}{\Character{16}}%
\TextAt{9}{476.4}{542.854}{\Character{61}}%
\TextAt{13}{489.681}{542.854}{\Character{16}}%
\TextAt{21}{495.697}{547.193}{u}%
\TextAt{23}{500.6}{550.005}{\Character{0}}%
\TextAt{24}{506.356}{550.005}{1}%
\TextAt{8}{511.005}{542.854}{,}%
\TextAt{8}{518.782}{542.854}{for}%
\TextAt{8}{79.226}{525.808}{some}%
\TextAt{13}{108.552}{525.808}{p}%
\TextAt{21}{114.427}{530.147}{n}%
\TextAt{8}{120.064}{525.808}{-th}%
\TextAt{8}{138.917}{525.808}{root}%
\TextAt{8}{163.963}{525.808}{of}%
\TextAt{8}{177.293}{525.808}{unity}%
\TextAt{13}{207.862}{525.808}{\Character{16}}%
\TextAt{8}{213.878}{525.808}{.}%
\TextAt{15}{79.226}{472.945}{§}%
\TextAt{15}{93.785}{472.945}{3}%
\TextAt{15}{120.693}{472.945}{The}%
\TextAt{15}{160.202}{472.945}{Hilbert}%
\TextAt{15}{228.169}{472.945}{Symbol}%
\TextAt{8}{79.226}{441.059}{Let}%
\TextAt{13}{99.242}{441.059}{K}%
\TextAt{8}{113.003}{441.059}{be}%
\TextAt{8}{127.974}{441.059}{a}%
\TextAt{8}{136.776}{441.059}{local}%
\TextAt{8}{163.462}{441.059}{field,}%
\TextAt{8}{191.308}{441.059}{or}%
\TextAt{13}{204.665}{441.059}{K}%
\TextAt{9}{218.803}{441.059}{\Character{61}}%
\TextAt{18}{231.229}{441.059}{R}%
\TextAt{8}{239.863}{441.059}{,}%
\TextAt{13}{246.249}{441.059}{K}%
\TextAt{9}{260.388}{441.059}{\Character{61}}%
\TextAt{18}{272.813}{441.059}{C}%
\TextAt{8}{281.447}{441.059}{.}%
\TextAt{8}{289.578}{441.059}{We}%
\TextAt{8}{308.769}{441.059}{assume}%
\TextAt{8}{348.259}{441.059}{that}%
\TextAt{13}{372.674}{441.059}{K}%
\TextAt{8}{386.436}{441.059}{contains}%
\TextAt{8}{431.402}{441.059}{the}%
\TextAt{8}{450.605}{441.059}{group}%
\TextAt{13}{482.827}{441.059}{\Character{22}}%
\TextAt{21}{489.87}{439.265}{n}%
\TextAt{8}{498.45}{441.059}{of}%
\TextAt{13}{510.826}{441.059}{n}%
\TextAt{8}{517.814}{441.059}{-th}%
\TextAt{8}{79.226}{424.013}{roots}%
\TextAt{8}{108.613}{424.013}{of}%
\TextAt{8}{121.668}{424.013}{unity,}%
\TextAt{8}{154.29}{424.013}{where}%
\TextAt{13}{187.825}{424.013}{n}%
\TextAt{8}{198.436}{424.013}{is}%
\TextAt{8}{209.925}{424.013}{a}%
\TextAt{8}{219.407}{424.013}{natural}%
\TextAt{8}{260.115}{424.013}{number}%
\TextAt{8}{302.437}{424.013}{which}%
\TextAt{8}{335.649}{424.013}{is}%
\TextAt{8}{347.139}{424.013}{relatively}%
\TextAt{8}{397.925}{424.013}{prime}%
\TextAt{8}{430.814}{424.013}{to}%
\TextAt{8}{444.85}{424.013}{the}%
\TextAt{8}{464.733}{424.013}{characteristic}%
\TextAt{8}{79.226}{406.967}{of}%
\TextAt{13}{92.83}{406.967}{K}%
\TextAt{8}{107.819}{406.967}{(i.e.,}%
\TextAt{13}{134.821}{406.967}{n}%
\TextAt{8}{145.98}{406.967}{can}%
\TextAt{8}{167.714}{406.967}{be}%
\TextAt{8}{183.912}{406.967}{arbitrary}%
\TextAt{8}{233.955}{406.967}{if}%
\TextAt{8}{244.953}{406.967}{char}%
\TextAt{9}{267.07}{406.967}{\Character{40}}%
\TextAt{13}{271.622}{406.967}{K}%
\TextAt{9}{282.44}{406.967}{\Character{41}}%
\TextAt{9}{290.761}{406.967}{\Character{61}}%
\TextAt{9}{303.649}{406.967}{0}%
\TextAt{8}{309.525}{406.967}{).}%
\TextAt{8}{323.35}{406.967}{Over}%
\TextAt{8}{352.245}{406.967}{such}%
\TextAt{8}{378.916}{406.967}{a}%
\TextAt{8}{388.945}{406.967}{field}%
\TextAt{13}{414.576}{406.967}{K}%
\TextAt{8}{429.565}{406.967}{we}%
\TextAt{8}{447.066}{406.967}{then}%
\TextAt{8}{474}{406.967}{have}%
\TextAt{8}{501.269}{406.967}{at}%
\TextAt{8}{515.853}{406.967}{our}%
\TextAt{8}{79.226}{389.921}{disposal,}%
\TextAt{8}{128.135}{389.921}{on}%
\TextAt{8}{145.179}{389.921}{the}%
\TextAt{8}{166.12}{389.921}{one}%
\TextAt{8}{188.364}{389.921}{hand,}%
\TextAt{8}{221.863}{389.921}{Kummer}%
\TextAt{8}{271.413}{389.921}{theory}%
\TextAt{8}{308.947}{389.921}{(see}%
\TextAt{8}{333.2}{389.921}{chap.}%
\TextAt{8}{367.745}{389.921}{IV,}%
\TextAt{8}{388.674}{389.921}{§3),}%
\TextAt{8}{412.896}{389.921}{and}%
\TextAt{8}{436.443}{389.921}{on}%
\TextAt{8}{453.487}{389.921}{the}%
\TextAt{8}{474.428}{389.921}{other,}%
\TextAt{8}{509.23}{389.921}{class}%
\TextAt{8}{79.226}{372.875}{field}%
\TextAt{8}{105.466}{372.875}{theory.}%
\TextAt{8}{148.436}{372.875}{It}%
\TextAt{8}{161.992}{372.875}{is}%
\TextAt{8}{174.639}{372.875}{the}%
\TextAt{8}{195.678}{372.875}{interplay}%
\TextAt{8}{245.673}{372.875}{between}%
\TextAt{8}{291.747}{372.875}{both}%
\TextAt{8}{320.27}{372.875}{theories,}%
\TextAt{8}{368.262}{372.875}{which}%
\TextAt{8}{402.632}{372.875}{gives}%
\TextAt{8}{432.195}{372.875}{rise}%
\TextAt{8}{454.598}{372.875}{to}%
\TextAt{8}{469.791}{372.875}{the}%
\TextAt{8}{490.831}{372.875}{“Hilbert}%
\TextAt{8}{79.226}{355.829}{symbol”.}%
\TextAt{8}{128.127}{355.829}{This}%
\TextAt{8}{154.977}{355.829}{is}%
\TextAt{8}{166.871}{355.829}{a}%
\TextAt{8}{176.756}{355.829}{highly}%
\TextAt{8}{212.334}{355.829}{remarkable}%
\TextAt{8}{273.28}{355.829}{phenomenon}%
\TextAt{8}{341.698}{355.829}{which}%
\TextAt{8}{375.315}{355.829}{will}%
\TextAt{8}{397.55}{355.829}{lead}%
\TextAt{8}{422.391}{355.829}{us}%
\TextAt{8}{437.537}{355.829}{to}%
\TextAt{8}{451.978}{355.829}{a}%
\TextAt{8}{461.863}{355.829}{generalization}%
\TextAt{8}{79.226}{338.783}{of}%
\TextAt{8}{92.556}{338.783}{the}%
\TextAt{8}{112.713}{338.783}{classical}%
\TextAt{8}{157.712}{338.783}{reciprocity}%
\TextAt{8}{215.922}{338.783}{law}%
\TextAt{8}{237.058}{338.783}{of}%
\TextAt{8}{250.388}{338.783}{Gauss,}%
\TextAt{8}{288.31}{338.783}{to}%
\TextAt{13}{302.62}{338.783}{n}%
\TextAt{8}{309.608}{338.783}{-th}%
\TextAt{8}{328.461}{338.783}{power}%
\TextAt{8}{362.605}{338.783}{residues.}%
\TextAt{8}{97.159}{318.748}{Let}%
\TextAt{13}{117.738}{318.748}{L}%
\TextAt{9}{129.023}{318.748}{\Character{61}}%
\TextAt{13}{141.449}{318.748}{K}%
\TextAt{9}{152.267}{318.748}{\Character{40}}%
\TextAt{22}{160.14}{324.582}{n}%
\TextAt{16}{158.099}{328.87}{\Character{112}}%
\Vector{q 0 G 0 g 0.478 w 96.062 -444.891 m 111.612 -444.891 l S Q}%
\TextAt{13}{168.062}{318.748}{K}%
\TextAt{3}{178.879}{322.202}{\Character{3}}%
\TextAt{9}{183.612}{318.748}{\Character{41}}%
\TextAt{8}{191.671}{318.748}{be}%
\TextAt{8}{207.205}{318.748}{the}%
\TextAt{8}{226.971}{318.748}{maximal}%
\TextAt{8}{274.39}{318.748}{abelian}%
\TextAt{8}{314.647}{318.748}{extension}%
\TextAt{8}{366.023}{318.748}{of}%
\TextAt{8}{378.963}{318.748}{exponent}%
\TextAt{13}{428.976}{318.748}{n}%
\TextAt{8}{435.964}{318.748}{.}%
\TextAt{8}{444.282}{318.748}{By}%
\TextAt{8}{462.255}{318.748}{(1.5),}%
\TextAt{8}{493.17}{318.748}{we}%
\TextAt{8}{510.007}{318.748}{then}%
\TextAt{8}{79.226}{301.702}{have}%
\TextAt{13}{266.962}{284.656}{N}%
\TextAt{21}{276.336}{282.796}{L}%
\TextAt{3}{282.095}{282.796}{\Character{106}}%
\TextAt{21}{284.448}{282.796}{K}%
\TextAt{13}{292.65}{284.656}{L}%
\TextAt{3}{300.615}{289.592}{\Character{3}}%
\TextAt{9}{308.668}{284.656}{\Character{61}}%
\TextAt{13}{321.093}{284.656}{K}%
\TextAt{3}{331.911}{289.592}{\Character{3}}%
\TextAt{21}{336.145}{289.592}{n}%
\TextAt{13}{341.782}{284.656}{\Character{59}}%
\TextAt{8}{79.226}{261.134}{and}%
\TextAt{8}{101.989}{261.134}{class}%
\TextAt{8}{129.426}{261.134}{field}%
\TextAt{8}{154.783}{261.134}{theory}%
\TextAt{8}{191.533}{261.134}{gives}%
\TextAt{8}{220.213}{261.134}{us}%
\TextAt{8}{235.229}{261.134}{the}%
\TextAt{8}{255.386}{261.134}{canonical}%
\TextAt{8}{306.769}{261.134}{isomorphism}%
\TextAt{13}{255.136}{232.133}{G}%
\TextAt{9}{264.369}{232.133}{\Character{40}}%
\TextAt{13}{268.922}{232.133}{L}%
\TextAt{16}{276.886}{232.133}{\Character{106}}%
\TextAt{13}{280.207}{232.133}{K}%
\TextAt{9}{291.025}{232.133}{\Character{41}}%
\NativeMathAt{298.898}{232.133}{11.955}{\cong}%
\TextAt{13}{311.517}{232.133}{K}%
\TextAt{3}{322.335}{237.069}{\Character{3}}%
\TextAt{9}{327.067}{232.133}{\Character{47}}%
\TextAt{13}{332.92}{232.133}{K}%
\TextAt{3}{343.738}{237.069}{\Character{3}}%
\TextAt{21}{347.972}{237.069}{n}%
\TextAt{13}{353.608}{232.133}{\Character{58}}%
\TextAt{8}{79.226}{203.132}{On}%
\TextAt{8}{98.737}{203.132}{the}%
\TextAt{8}{118.893}{203.132}{other}%
\TextAt{8}{149.463}{203.132}{hand,}%
\TextAt{8}{181.981}{203.132}{Kummer}%
\TextAt{8}{230.746}{203.132}{theory}%
\TextAt{8}{267.496}{203.132}{gives}%
\TextAt{8}{296.177}{203.132}{the}%
\TextAt{8}{316.333}{203.132}{canonical}%
\TextAt{8}{367.716}{203.132}{isomorphism}%
\TextAt{8}{229.427}{174.131}{Hom}%
\TextAt{9}{253.816}{174.131}{\Character{40}}%
\TextAt{13}{258.368}{174.131}{G}%
\TextAt{9}{267.602}{174.131}{\Character{40}}%
\TextAt{13}{272.154}{174.131}{L}%
\TextAt{16}{280.119}{174.131}{\Character{106}}%
\TextAt{13}{283.44}{174.131}{K}%
\TextAt{9}{294.257}{174.131}{\Character{41}}%
\TextAt{13}{298.81}{174.131}{\Character{59}}%
\TextAt{13}{304.034}{174.131}{\Character{22}}%
\TextAt{21}{311.097}{172.338}{n}%
\TextAt{9}{316.733}{174.131}{\Character{41}}%
\NativeMathAt{324.606}{174.131}{11.955}{\cong}%
\TextAt{13}{337.225}{174.131}{K}%
\TextAt{3}{348.043}{179.067}{\Character{3}}%
\TextAt{9}{352.775}{174.131}{\Character{47}}%
\TextAt{13}{358.628}{174.131}{K}%
\TextAt{3}{369.446}{179.067}{\Character{3}}%
\TextAt{21}{373.68}{179.067}{n}%
\TextAt{13}{379.317}{174.131}{\Character{58}}%
\TextAt{8}{297.211}{43.033}{305}%
\Destination{78.23}{809.815}{nav.page.305}%
\Destination{79.23}{499.355}{nav.section.5.3}%
\Destination{79.226}{775.155}{nav.ex.5.2.2}%
\Destination{79.226}{691.955}{nav.ex.5.2.3}%
\Destination{79.226}{556.855}{nav.ex.5.2.4}%
\Destination{79.36}{456.055}{nav.print.333}%
\Destination{79.23}{499.355}{outline.51}%
\Hotlink{181.174}{572.655}{13.995}{11.047}{nav.chapter.4}%
\Hotlink{206.997}{572.655}{6.858}{11.047}{nav.section.4.5}%
\Hotlink{290.265}{540.655}{6.858}{11.047}{nav.ex.5.2.2}%
\Hotlink{367.245}{387.755}{13.995}{11.047}{nav.chapter.4}%
\Hotlink{393.853}{387.755}{6.858}{11.047}{nav.section.4.3}%
\Hotlink{187.492}{149.755}{5.981}{9.572}{nav.ex.5.3.3}%
\Hotlink{461.755}{316.555}{25.078}{11.047}{nav.number.5.1.5}%
\Destination{79.226}{611.879}{errata.116}%
}

% Page 318 (body page 306)
\Page{%
\TextAt{2}{79.226}{796.943}{V.}%
\TextAt{2}{94.073}{796.943}{Local}%
\TextAt{2}{123.615}{796.943}{Class}%
\TextAt{2}{152.208}{796.943}{Field}%
\TextAt{2}{179.939}{796.943}{Theory}%
\Vector{q 0 G 0 g 0.299 w 7.226 18.729 m 460.769 18.729 l S Q}%
\TextAt{8}{79.226}{761.124}{The}%
\TextAt{8}{103.28}{761.124}{bilinear}%
\TextAt{8}{145.553}{761.124}{map}%
\TextAt{13}{158.279}{734.554}{G}%
\TextAt{9}{167.513}{734.554}{\Character{40}}%
\TextAt{13}{172.065}{734.554}{L}%
\TextAt{16}{180.03}{734.554}{\Character{106}}%
\TextAt{13}{183.35}{734.554}{K}%
\TextAt{9}{194.168}{734.554}{\Character{41}}%
\TextAt{16}{201.377}{734.554}{\Character{2}}%
\TextAt{8}{213.332}{734.554}{Hom}%
\TextAt{9}{237.721}{734.554}{\Character{40}}%
\TextAt{13}{242.273}{734.554}{G}%
\TextAt{9}{251.507}{734.554}{\Character{40}}%
\TextAt{13}{256.059}{734.554}{L}%
\TextAt{16}{264.024}{734.554}{\Character{106}}%
\TextAt{13}{267.344}{734.554}{K}%
\TextAt{9}{278.162}{734.554}{\Character{41}}%
\TextAt{13}{282.715}{734.554}{\Character{59}}%
\TextAt{13}{287.939}{734.554}{\Character{22}}%
\TextAt{21}{295.002}{732.761}{n}%
\TextAt{9}{300.638}{734.554}{\Character{41}}%
\NativeArrowAt{308.511}{734.554}{11.955}{19.26}{\longrightarrow}%
\TextAt{13}{331.093}{734.554}{\Character{22}}%
\TextAt{21}{338.136}{732.761}{n}%
\TextAt{13}{343.772}{734.554}{\Character{59}}%
\TextAt{9}{372.927}{734.554}{\Character{40}}%
\TextAt{13}{377.479}{734.554}{\Character{27}}%
\TextAt{13}{383.911}{734.554}{\Character{59}}%
\TextAt{13}{389.135}{734.554}{\Character{31}}%
\TextAt{9}{396.489}{734.554}{\Character{41}}%
\NativeArrowAt{404.362}{734.554}{11.955}{19.26}{\longmapsto}%
\TextAt{13}{426.944}{734.554}{\Character{31}}%
\TextAt{9}{434.278}{734.554}{\Character{40}}%
\TextAt{13}{438.83}{734.554}{\Character{27}}%
\TextAt{9}{445.913}{734.554}{\Character{41}}%
\TextAt{13}{450.465}{734.554}{\Character{59}}%
\TextAt{8}{79.226}{707.984}{therefore}%
\TextAt{8}{128.326}{707.984}{defines}%
\TextAt{8}{166.75}{707.984}{a}%
\TextAt{8}{176.505}{707.984}{nondegenerate}%
\TextAt{8}{253.903}{707.984}{bilinear}%
\TextAt{8}{296.177}{707.984}{pairing}%
\TextAt{20}{216.844}{694.535}{\Character{18}}%
\TextAt{13}{230.737}{685.766}{\Character{59}}%
\Vector{q 0 G 0 g 0.478 w 154.84 -93.333 m 165.886 -93.333 l S Q}%
\TextAt{14}{229.371}{669.477}{p}%
\TextAt{20}{239.082}{694.535}{\Character{19}}%
\TextAt{9}{251.203}{677.678}{\Character{58}}%
\TextAt{13}{257.775}{677.678}{K}%
\TextAt{3}{268.593}{682.614}{\Character{3}}%
\TextAt{9}{273.325}{677.678}{\Character{47}}%
\TextAt{13}{279.178}{677.678}{K}%
\TextAt{3}{289.996}{682.614}{\Character{3}}%
\TextAt{21}{294.23}{682.614}{n}%
\TextAt{16}{302.523}{677.678}{\Character{2}}%
\TextAt{13}{314.478}{677.678}{K}%
\TextAt{3}{325.296}{682.614}{\Character{3}}%
\TextAt{9}{330.028}{677.678}{\Character{47}}%
\TextAt{13}{335.881}{677.678}{K}%
\TextAt{3}{346.699}{682.614}{\Character{3}}%
\TextAt{21}{350.933}{682.614}{n}%
\NativeArrowAt{359.89}{677.678}{11.955}{19.26}{\longrightarrow}%
\TextAt{13}{382.472}{677.678}{\Character{22}}%
\TextAt{21}{389.515}{675.885}{n}%
\TextAt{8}{79.226}{646.834}{(bilinear}%
\TextAt{8}{127.16}{646.834}{in}%
\TextAt{8}{141.919}{646.834}{the}%
\TextAt{8}{163.181}{646.834}{multiplicative}%
\TextAt{8}{238.11}{646.834}{sense).}%
\TextAt{8}{280.564}{646.834}{This}%
\TextAt{8}{308.389}{646.834}{pairing}%
\TextAt{8}{349.174}{646.834}{is}%
\TextAt{8}{362.044}{646.834}{called}%
\TextAt{8}{396.313}{646.834}{the}%
\TextAt{12}{417.575}{646.834}{Hilbert}%
\TextAt{12}{466.046}{646.834}{symbol}%
\TextAt{8}{507.614}{646.834}{.}%
\TextAt{8}{519.38}{646.834}{Its}%
\TextAt{8}{79.226}{629.788}{relation}%
\TextAt{8}{122.157}{629.788}{to}%
\TextAt{8}{136.467}{629.788}{the}%
\TextAt{8}{156.624}{629.788}{norm}%
\TextAt{8}{187.193}{629.788}{residue}%
\TextAt{8}{226.92}{629.788}{symbol}%
\TextAt{8}{266.982}{629.788}{is}%
\TextAt{8}{278.746}{629.788}{described}%
\TextAt{8}{330.5}{629.788}{explicitly}%
\TextAt{8}{381.226}{629.788}{in}%
\TextAt{8}{394.878}{629.788}{the}%
\TextAt{8}{415.035}{629.788}{following}%
\TextAt{8}{464.47}{629.788}{proposition.}%
\TextAt{12}{79.226}{604.47}{Proposition}%
\TextAt{12}{151.937}{604.47}{(3.1).}%
\TextAt{8}{187.253}{604.47}{For}%
\TextAt{13}{208.222}{604.47}{a\Character{59}}%
\TextAt{13}{219.591}{604.47}{b}%
\NativeMathAt{227.909}{604.47}{11.955}{\in}%
\TextAt{13}{239.2}{604.47}{K}%
\TextAt{3}{250.018}{608.809}{\Character{3}}%
\TextAt{8}{254.75}{604.47}{,}%
\TextAt{8}{261.9}{604.47}{the}%
\TextAt{8}{282.056}{604.47}{Hilbert}%
\TextAt{8}{322.369}{604.47}{symbol}%
\TextAt{20}{362.431}{614.154}{\Character{0}}%
\TextAt{21}{369.106}{609.726}{a\Character{59}b}%
\Vector{q 0 G 0 g 0.478 w 297.106 -166.541 m 307.579 -166.541 l S Q}%
\TextAt{27}{372.348}{600.348}{p}%
\TextAt{20}{380.774}{614.154}{\Character{1}}%
\NativeMathAt{389.574}{604.47}{11.955}{\in}%
\TextAt{13}{400.865}{604.47}{\Character{22}}%
\TextAt{21}{407.908}{602.677}{n}%
\TextAt{8}{417.442}{604.47}{is}%
\TextAt{8}{429.206}{604.47}{given}%
\TextAt{8}{459.775}{604.47}{by}%
\TextAt{9}{225.717}{570.986}{\Character{40}}%
\TextAt{13}{230.269}{570.986}{a\Character{59}}%
\TextAt{13}{241.638}{570.986}{K}%
\TextAt{9}{252.476}{570.986}{\Character{40}}%
\TextAt{22}{260.349}{577.6}{n}%
\TextAt{16}{258.308}{581.758}{\Character{112}}%
\Vector{q 0 G 0 g 0.478 w 196.271 -192.003 m 201.248 -192.003 l S Q}%
\TextAt{13}{268.271}{570.986}{b}%
\TextAt{9}{273.248}{570.986}{\Character{41}}%
\TextAt{16}{277.8}{570.986}{\Character{106}}%
\TextAt{13}{281.121}{570.986}{K}%
\TextAt{9}{291.939}{570.986}{\Character{41}}%
\TextAt{22}{299.812}{577.6}{n}%
\TextAt{16}{297.771}{581.758}{\Character{112}}%
\Vector{q 0 G 0 g 0.478 w 235.734 -192.003 m 240.711 -192.003 l S Q}%
\TextAt{13}{307.734}{570.986}{b}%
\TextAt{9}{316.032}{570.986}{\Character{61}}%
\TextAt{20}{328.457}{587.843}{\Character{18}}%
\TextAt{13}{338.453}{579.073}{a\Character{59}}%
\TextAt{13}{349.822}{579.073}{b}%
\Vector{q 0 G 0 g 0.478 w 266.453 -200.026 m 282.819 -200.026 l S Q}%
\TextAt{14}{343.645}{562.785}{p}%
\TextAt{20}{356.015}{587.843}{\Character{19}}%
\TextAt{22}{370.128}{577.6}{n}%
\TextAt{16}{368.087}{581.758}{\Character{112}}%
\Vector{q 0 G 0 g 0.478 w 306.05 -192.003 m 311.027 -192.003 l S Q}%
\TextAt{13}{378.05}{570.986}{b\Character{58}}%
\TextAt{19}{79.226}{523.598}{Proof.}%
\TextAt{8}{116}{523.598}{The}%
\TextAt{8}{140.098}{523.598}{image}%
\TextAt{8}{173.962}{523.598}{of}%
\TextAt{13}{187.336}{523.598}{a}%
\TextAt{8}{197.422}{523.598}{under}%
\TextAt{8}{230.629}{523.598}{the}%
\TextAt{8}{250.829}{523.598}{isomorphism}%
\TextAt{13}{319.292}{523.598}{K}%
\TextAt{3}{330.11}{527.937}{\Character{3}}%
\TextAt{9}{334.842}{523.598}{\Character{47}}%
\TextAt{13}{340.695}{523.598}{K}%
\TextAt{3}{351.513}{527.937}{\Character{3}}%
\TextAt{21}{355.747}{527.937}{n}%
\NativeMathAt{364.778}{523.598}{11.955}{\cong}%
\TextAt{13}{377.472}{523.598}{G}%
\TextAt{9}{386.706}{523.598}{\Character{40}}%
\TextAt{13}{391.258}{523.598}{L}%
\TextAt{16}{399.223}{523.598}{\Character{106}}%
\TextAt{13}{402.543}{523.598}{K}%
\TextAt{9}{413.361}{523.598}{\Character{41}}%
\TextAt{8}{421.854}{523.598}{of}%
\TextAt{8}{435.228}{523.598}{class}%
\TextAt{8}{462.709}{523.598}{field}%
\TextAt{8}{488.109}{523.598}{theory}%
\TextAt{8}{524.903}{523.598}{is}%
\TextAt{8}{79.226}{506.552}{the}%
\TextAt{8}{98.732}{506.552}{norm}%
\TextAt{8}{128.651}{506.552}{residue}%
\TextAt{8}{167.727}{506.552}{symbol}%
\TextAt{13}{207.138}{506.552}{\Character{27}}%
\TextAt{9}{217.541}{506.552}{\Character{61}}%
\TextAt{9}{229.951}{506.552}{\Character{40}}%
\TextAt{13}{234.519}{506.552}{a\Character{59}}%
\TextAt{13}{245.888}{506.552}{L}%
\TextAt{16}{253.873}{506.552}{\Character{106}}%
\TextAt{13}{257.193}{506.552}{K}%
\TextAt{9}{268.011}{506.552}{\Character{41}}%
\TextAt{8}{272.563}{506.552}{.}%
\TextAt{8}{280.795}{506.552}{The}%
\TextAt{8}{304.198}{506.552}{image}%
\TextAt{8}{337.368}{506.552}{of}%
\TextAt{13}{350.048}{506.552}{b}%
\TextAt{8}{358.272}{506.552}{under}%
\TextAt{8}{390.785}{506.552}{the}%
\TextAt{8}{410.29}{506.552}{isomorphism}%
\TextAt{13}{478.059}{506.552}{K}%
\TextAt{3}{488.877}{510.891}{\Character{3}}%
\TextAt{9}{493.609}{506.552}{\Character{47}}%
\TextAt{13}{499.462}{506.552}{K}%
\TextAt{3}{510.28}{510.891}{\Character{3}}%
\TextAt{21}{514.514}{510.891}{n}%
\NativeMathAt{523.471}{506.552}{11.955}{\cong}%
\TextAt{8}{79.226}{489.506}{Hom}%
\TextAt{9}{103.615}{489.506}{\Character{40}}%
\TextAt{13}{108.167}{489.506}{G}%
\TextAt{9}{117.401}{489.506}{\Character{40}}%
\TextAt{13}{121.953}{489.506}{L}%
\TextAt{16}{129.917}{489.506}{\Character{106}}%
\TextAt{13}{133.238}{489.506}{K}%
\TextAt{9}{144.056}{489.506}{\Character{41}}%
\TextAt{13}{148.608}{489.506}{\Character{59}}%
\TextAt{13}{153.832}{489.506}{\Character{22}}%
\TextAt{21}{160.895}{487.713}{n}%
\TextAt{9}{166.532}{489.506}{\Character{41}}%
\TextAt{8}{174.258}{489.506}{of}%
\TextAt{8}{186.865}{489.506}{Kummer}%
\TextAt{8}{234.907}{489.506}{theory}%
\TextAt{8}{270.934}{489.506}{is}%
\TextAt{8}{281.975}{489.506}{the}%
\TextAt{8}{301.408}{489.506}{character}%
\TextAt{13}{351.745}{489.506}{\Character{31}}%
\TextAt{21}{359.079}{487.713}{b}%
\TextAt{9}{366.52}{489.506}{\Character{58}}%
\TextAt{13}{373.093}{489.506}{G}%
\TextAt{9}{382.326}{489.506}{\Character{40}}%
\TextAt{13}{386.879}{489.506}{L}%
\TextAt{16}{394.843}{489.506}{\Character{106}}%
\TextAt{13}{398.164}{489.506}{K}%
\TextAt{9}{408.982}{489.506}{\Character{41}}%
\NativeMathAt{416.855}{489.506}{11.955}{\rightarrow}%
\TextAt{13}{432.131}{489.506}{\Character{22}}%
\TextAt{21}{439.174}{487.713}{n}%
\TextAt{8}{447.984}{489.506}{given}%
\TextAt{8}{477.83}{489.506}{by}%
\TextAt{13}{493.366}{489.506}{\Character{31}}%
\TextAt{21}{500.7}{487.713}{b}%
\TextAt{9}{504.82}{489.506}{\Character{40}}%
\TextAt{13}{509.373}{489.506}{\Character{28}}%
\TextAt{9}{515.792}{489.506}{\Character{41}}%
\TextAt{9}{523.647}{489.506}{\Character{61}}%
\TextAt{13}{79.226}{472.46}{\Character{28}}%
\TextAt{22}{88.966}{478.374}{n}%
\TextAt{16}{86.925}{482.649}{\Character{112}}%
\Vector{q 0 G 0 g 0.478 w 24.888 -291.112 m 29.865 -291.112 l S Q}%
\TextAt{13}{96.888}{472.46}{b}%
\TextAt{9}{101.865}{472.46}{\Character{47}}%
\TextAt{22}{111.039}{478.374}{n}%
\TextAt{16}{108.998}{482.649}{\Character{112}}%
\Vector{q 0 G 0 g 0.478 w 46.96 -291.112 m 51.937 -291.112 l S Q}%
\TextAt{13}{118.96}{472.46}{b}%
\TextAt{8}{123.937}{472.46}{.}%
\TextAt{8}{132.386}{472.46}{By}%
\TextAt{8}{150.749}{472.46}{definition}%
\TextAt{8}{202.778}{472.46}{of}%
\TextAt{8}{216.108}{472.46}{the}%
\TextAt{8}{236.264}{472.46}{Hilbert}%
\TextAt{8}{276.577}{472.46}{symbol,}%
\TextAt{8}{319.891}{472.46}{we}%
\TextAt{8}{337.118}{472.46}{have}%
\TextAt{20}{233.939}{458.473}{\Character{18}}%
\TextAt{13}{243.935}{449.704}{a\Character{59}}%
\TextAt{13}{255.304}{449.704}{b}%
\Vector{q 0 G 0 g 0.478 w 171.935 -329.395 m 188.301 -329.395 l S Q}%
\TextAt{14}{249.126}{433.415}{p}%
\TextAt{20}{261.496}{458.473}{\Character{19}}%
\TextAt{9}{273.617}{441.616}{\Character{61}}%
\TextAt{13}{286.043}{441.616}{\Character{31}}%
\TextAt{21}{293.376}{439.823}{b}%
\TextAt{9}{297.497}{441.616}{\Character{40}}%
\TextAt{13}{302.049}{441.616}{\Character{27}}%
\TextAt{9}{309.132}{441.616}{\Character{41}}%
\TextAt{9}{316.987}{441.616}{\Character{61}}%
\TextAt{13}{329.431}{441.616}{\Character{27}}%
\TextAt{22}{339.834}{448.23}{n}%
\TextAt{16}{337.793}{452.388}{\Character{112}}%
\Vector{q 0 G 0 g 0.478 w 275.756 -321.373 m 280.733 -321.373 l S Q}%
\TextAt{13}{347.756}{441.616}{b}%
\TextAt{9}{352.733}{441.616}{\Character{47}}%
\TextAt{22}{361.906}{448.23}{n}%
\TextAt{16}{359.866}{452.388}{\Character{112}}%
\Vector{q 0 G 0 g 0.478 w 297.828 -321.373 m 302.805 -321.373 l S Q}%
\TextAt{13}{369.828}{441.616}{b\Character{59}}%
\TextAt{8}{79.226}{408.594}{hence}%
\TextAt{9}{111.732}{408.594}{\Character{40}}%
\TextAt{13}{116.285}{408.594}{a\Character{59}}%
\TextAt{13}{127.654}{408.594}{K}%
\TextAt{9}{138.491}{408.594}{\Character{40}}%
\TextAt{22}{146.365}{414.507}{n}%
\TextAt{16}{144.324}{418.782}{\Character{112}}%
\Vector{q 0 G 0 g 0.478 w 82.286 -354.979 m 87.263 -354.979 l S Q}%
\TextAt{13}{154.286}{408.594}{b}%
\TextAt{9}{159.263}{408.594}{\Character{41}}%
\TextAt{16}{163.816}{408.594}{\Character{106}}%
\TextAt{13}{167.137}{408.594}{K}%
\TextAt{9}{177.954}{408.594}{\Character{41}}%
\TextAt{22}{185.827}{414.507}{n}%
\TextAt{16}{183.787}{418.782}{\Character{112}}%
\Vector{q 0 G 0 g 0.478 w 121.749 -354.979 m 126.726 -354.979 l S Q}%
\TextAt{13}{193.749}{408.594}{b}%
\TextAt{9}{202.047}{408.594}{\Character{61}}%
\TextAt{9}{214.457}{408.594}{\Character{40}}%
\TextAt{13}{219.025}{408.594}{a\Character{59}}%
\TextAt{13}{230.394}{408.594}{L}%
\TextAt{16}{238.379}{408.594}{\Character{106}}%
\TextAt{13}{241.699}{408.594}{K}%
\TextAt{9}{252.517}{408.594}{\Character{41}}%
\TextAt{22}{260.39}{414.507}{n}%
\TextAt{16}{258.349}{418.782}{\Character{112}}%
\Vector{q 0 G 0 g 0.478 w 196.312 -354.979 m 201.289 -354.979 l S Q}%
\TextAt{13}{268.312}{408.594}{b}%
\TextAt{9}{276.61}{408.594}{\Character{61}}%
\TextAt{20}{289.035}{418.277}{\Character{0}}%
\TextAt{21}{295.71}{413.849}{a\Character{59}b}%
\Vector{q 0 G 0 g 0.478 w 223.71 -362.418 m 234.183 -362.418 l S Q}%
\TextAt{27}{298.953}{404.471}{p}%
\TextAt{20}{307.379}{418.277}{\Character{1}}%
\TextAt{22}{316.179}{414.507}{n}%
\TextAt{16}{314.138}{418.782}{\Character{112}}%
\Vector{q 0 G 0 g 0.478 w 252.101 -354.979 m 257.078 -354.979 l S Q}%
\TextAt{13}{324.101}{408.594}{b}%
\TextAt{8}{329.078}{408.594}{.}%
\Vector{q 0 G 0 g 0.399 w 452.334 -365.406 m 452.334 -357.337 l S Q}%
\Vector{q 0 G 0 g 0.399 w 452.534 -357.536 m 459.707 -357.536 l S Q}%
\Vector{q 0 G 0 g 0.399 w 452.534 -365.207 m 459.707 -365.207 l S Q}%
\Vector{q 0 G 0 g 0.399 w 459.906 -365.406 m 459.906 -357.337 l S Q}%
\TextAt{8}{97.159}{380.287}{The}%
\TextAt{8}{121.213}{380.287}{Hilbert}%
\TextAt{8}{161.526}{380.287}{symbol}%
\TextAt{8}{201.587}{380.287}{has}%
\TextAt{8}{222.461}{380.287}{the}%
\TextAt{8}{242.617}{380.287}{following}%
\TextAt{8}{292.052}{380.287}{fundamental}%
\TextAt{8}{359.694}{380.287}{properties:}%
\TextAt{12}{79.226}{354.969}{Proposition}%
\TextAt{12}{151.937}{354.969}{(3.2).}%
\TextAt{8}{89.579}{330.649}{(i)}%
\TextAt{20}{107.918}{347.505}{\Character{18}}%
\TextAt{13}{117.914}{338.736}{aa}%
\TextAt{3}{130.204}{343.075}{\Character{48}}%
\TextAt{13}{132.999}{338.736}{\Character{59}}%
\TextAt{13}{138.223}{338.736}{b}%
\Vector{q 0 G 0 g 0.478 w 45.914 -440.363 m 71.221 -440.363 l S Q}%
\TextAt{14}{127.576}{322.448}{p}%
\TextAt{20}{144.416}{347.505}{\Character{19}}%
\TextAt{9}{156.537}{330.649}{\Character{61}}%
\TextAt{20}{168.963}{347.505}{\Character{18}}%
\TextAt{13}{178.959}{338.736}{a\Character{59}}%
\TextAt{13}{190.328}{338.736}{b}%
\Vector{q 0 G 0 g 0.478 w 106.959 -440.363 m 123.325 -440.363 l S Q}%
\TextAt{14}{184.15}{322.448}{p}%
\TextAt{20}{196.52}{347.505}{\Character{19}}%
\TextAt{20}{207.304}{347.505}{\Character{18}}%
\TextAt{13}{217.309}{338.736}{a}%
\TextAt{3}{223.454}{343.075}{\Character{48}}%
\TextAt{13}{226.249}{338.736}{\Character{59}}%
\TextAt{13}{231.473}{338.736}{b}%
\Vector{q 0 G 0 g 0.478 w 145.309 -440.363 m 164.47 -440.363 l S Q}%
\TextAt{14}{223.898}{322.448}{p}%
\TextAt{20}{237.666}{347.505}{\Character{19}}%
\TextAt{8}{246.466}{330.649}{,}%
\TextAt{8}{86.327}{300.96}{(ii)}%
\TextAt{20}{107.918}{317.817}{\Character{18}}%
\TextAt{13}{117.914}{309.047}{a\Character{59}}%
\TextAt{13}{129.283}{309.047}{bb}%
\TextAt{3}{139.258}{313.386}{\Character{48}}%
\Vector{q 0 G 0 g 0.478 w 45.914 -470.052 m 70.053 -470.052 l S Q}%
\TextAt{14}{126.992}{292.759}{p}%
\TextAt{20}{143.248}{317.817}{\Character{19}}%
\TextAt{9}{155.369}{300.96}{\Character{61}}%
\TextAt{20}{167.795}{317.817}{\Character{18}}%
\TextAt{13}{177.791}{309.047}{a\Character{59}}%
\TextAt{13}{189.16}{309.047}{b}%
\Vector{q 0 G 0 g 0.478 w 105.791 -470.052 m 122.157 -470.052 l S Q}%
\TextAt{14}{182.982}{292.759}{p}%
\TextAt{20}{195.353}{317.817}{\Character{19}}%
\TextAt{20}{206.137}{317.817}{\Character{18}}%
\TextAt{13}{216.141}{309.047}{a\Character{59}}%
\TextAt{13}{227.51}{309.047}{b}%
\TextAt{3}{232.507}{313.386}{\Character{48}}%
\Vector{q 0 G 0 g 0.478 w 144.141 -470.052 m 163.303 -470.052 l S Q}%
\TextAt{14}{222.73}{292.759}{p}%
\TextAt{20}{236.498}{317.817}{\Character{19}}%
\TextAt{8}{245.298}{300.96}{,}%
\TextAt{8}{83.076}{271.271}{(iii)}%
\TextAt{20}{107.918}{288.128}{\Character{18}}%
\TextAt{13}{117.914}{279.358}{a\Character{59}}%
\TextAt{13}{129.283}{279.358}{b}%
\Vector{q 0 G 0 g 0.478 w 45.914 -499.741 m 62.281 -499.741 l S Q}%
\TextAt{14}{123.106}{263.07}{p}%
\TextAt{20}{135.476}{288.128}{\Character{19}}%
\TextAt{9}{147.597}{271.271}{\Character{61}}%
\TextAt{9}{160.007}{271.271}{1}%
\NativeArrowAt{181.152}{271.271}{11.955}{21.926}{\Longleftrightarrow}%
\TextAt{13}{218.346}{271.271}{a}%
\TextAt{8}{228.388}{271.271}{is}%
\TextAt{8}{240.152}{271.271}{a}%
\TextAt{8}{249.907}{271.271}{norm}%
\TextAt{8}{280.477}{271.271}{from}%
\TextAt{8}{308.117}{271.271}{the}%
\TextAt{8}{328.273}{271.271}{extension}%
\TextAt{13}{380.039}{271.271}{K}%
\TextAt{9}{390.857}{271.271}{\Character{40}}%
\TextAt{22}{398.73}{277.885}{n}%
\TextAt{16}{396.689}{282.043}{\Character{112}}%
\Vector{q 0 G 0 g 0.478 w 334.652 -491.718 m 339.629 -491.718 l S Q}%
\TextAt{13}{406.652}{271.271}{b}%
\TextAt{9}{411.629}{271.271}{\Character{41}}%
\TextAt{16}{416.181}{271.271}{\Character{106}}%
\TextAt{13}{419.502}{271.271}{K}%
\TextAt{8}{430.32}{271.271}{,}%
\TextAt{8}{83.398}{239.6}{(iv)}%
\TextAt{20}{107.918}{256.457}{\Character{18}}%
\TextAt{13}{117.914}{247.688}{a\Character{59}}%
\TextAt{13}{129.283}{247.688}{b}%
\Vector{q 0 G 0 g 0.478 w 45.914 -531.411 m 62.281 -531.411 l S Q}%
\TextAt{14}{123.106}{231.4}{p}%
\TextAt{20}{135.476}{256.457}{\Character{19}}%
\TextAt{9}{147.597}{239.6}{\Character{61}}%
\TextAt{20}{160.023}{256.457}{\Character{18}}%
\TextAt{13}{170.019}{247.688}{b\Character{59}}%
\TextAt{13}{180.217}{247.688}{a}%
\Vector{q 0 G 0 g 0.478 w 98.019 -531.411 m 114.385 -531.411 l S Q}%
\TextAt{14}{175.21}{231.4}{p}%
\TextAt{20}{187.58}{256.457}{\Character{19}}%
\TextAt{3}{196.381}{253.781}{\Character{0}}%
\TextAt{17}{202.967}{253.781}{1}%
\TextAt{8}{207.7}{239.6}{,}%
\TextAt{8}{86.65}{209.911}{(v)}%
\TextAt{20}{107.918}{226.768}{\Character{18}}%
\TextAt{13}{117.914}{217.999}{a\Character{59}}%
\TextAt{9}{129.303}{217.999}{1}%
\TextAt{16}{137.813}{217.999}{\Character{0}}%
\TextAt{13}{149.768}{217.999}{a}%
\Vector{q 0 G 0 g 0.478 w 45.914 -561.1 m 83.913 -561.1 l S Q}%
\TextAt{14}{133.922}{201.711}{p}%
\TextAt{20}{157.109}{226.768}{\Character{19}}%
\TextAt{9}{169.23}{209.911}{\Character{61}}%
\TextAt{9}{181.64}{209.911}{1}%
\TextAt{8}{191.406}{209.911}{and}%
\TextAt{20}{216.161}{226.768}{\Character{18}}%
\TextAt{13}{226.157}{217.999}{a\Character{59}}%
\TextAt{16}{237.546}{217.999}{\Character{0}}%
\TextAt{13}{246.844}{217.999}{a}%
\Vector{q 0 G 0 g 0.478 w 154.157 -561.1 m 180.989 -561.1 l S Q}%
\TextAt{14}{236.582}{201.711}{p}%
\TextAt{20}{254.185}{226.768}{\Character{19}}%
\TextAt{9}{266.306}{209.911}{\Character{61}}%
\TextAt{9}{278.716}{209.911}{1}%
\TextAt{8}{284.585}{209.911}{,}%
\TextAt{8}{83.398}{186.767}{(vi)}%
\TextAt{8}{107.918}{186.767}{If}%
\TextAt{20}{119.611}{196.451}{\Character{0}}%
\TextAt{21}{126.286}{192.023}{a\Character{59}b}%
\Vector{q 0 G 0 g 0.478 w 54.286 -584.244 m 64.759 -584.244 l S Q}%
\TextAt{27}{129.528}{182.644}{p}%
\TextAt{20}{137.954}{196.451}{\Character{1}}%
\TextAt{9}{146.754}{186.767}{\Character{61}}%
\TextAt{9}{159.164}{186.767}{1}%
\TextAt{8}{168.93}{186.767}{for}%
\TextAt{8}{186.815}{186.767}{all}%
\TextAt{13}{203.074}{186.767}{b}%
\NativeMathAt{211.372}{186.767}{11.955}{\in}%
\TextAt{13}{222.663}{186.767}{K}%
\TextAt{3}{233.481}{191.105}{\Character{3}}%
\TextAt{8}{238.213}{186.767}{,}%
\TextAt{8}{245.362}{186.767}{then}%
\TextAt{13}{272.022}{186.767}{a}%
\NativeMathAt{281.488}{186.767}{11.955}{\in}%
\TextAt{13}{292.779}{186.767}{K}%
\TextAt{3}{303.597}{191.105}{\Character{3}}%
\TextAt{21}{307.831}{191.105}{n}%
\TextAt{8}{313.467}{186.767}{.}%
\TextAt{19}{79.226}{153.177}{Proof.}%
\TextAt{8}{116}{153.177}{(i)}%
\TextAt{8}{132.174}{153.177}{and}%
\TextAt{8}{154.851}{153.177}{(ii)}%
\TextAt{8}{174.277}{153.177}{are}%
\TextAt{8}{193.703}{153.177}{clear}%
\TextAt{8}{221.58}{153.177}{from}%
\TextAt{8}{249.135}{153.177}{the}%
\TextAt{8}{269.207}{153.177}{definition,}%
\TextAt{8}{324.419}{153.177}{(iii)}%
\TextAt{8}{347.097}{153.177}{follows}%
\TextAt{8}{385.447}{153.177}{from}%
\TextAt{8}{413.002}{153.177}{(3.1),}%
\TextAt{8}{444.161}{153.177}{and}%
\TextAt{8}{466.838}{153.177}{(vi)}%
\TextAt{8}{489.193}{153.177}{reformu-}%
\TextAt{8}{79.226}{136.132}{lates}%
\TextAt{8}{106.603}{136.132}{the}%
\TextAt{8}{126.76}{136.132}{nondegenerateness}%
\TextAt{8}{225.091}{136.132}{of}%
\TextAt{8}{238.421}{136.132}{the}%
\TextAt{8}{258.578}{136.132}{Hilbert}%
\TextAt{8}{298.89}{136.132}{symbol.}%
\TextAt{8}{97.159}{116.097}{If}%
\TextAt{13}{108.851}{116.097}{b}%
\NativeMathAt{117.149}{116.097}{11.955}{\in}%
\TextAt{13}{128.44}{116.097}{K}%
\TextAt{3}{139.258}{120.435}{\Character{3}}%
\TextAt{8}{147.887}{116.097}{and}%
\TextAt{13}{170.65}{116.097}{x}%
\NativeMathAt{180.623}{116.097}{11.955}{\in}%
\TextAt{13}{191.914}{116.097}{K}%
\TextAt{8}{206.629}{116.097}{such}%
\TextAt{8}{233.026}{116.097}{that}%
\TextAt{13}{258.395}{116.097}{x}%
\TextAt{21}{265.047}{120.435}{n}%
\TextAt{16}{273.34}{116.097}{\Character{0}}%
\TextAt{13}{285.295}{116.097}{b}%
\NativeMathAt{293.593}{116.097}{11.955}{\neq}%
\TextAt{9}{306.002}{116.097}{0}%
\TextAt{8}{311.872}{116.097}{,}%
\TextAt{8}{319.021}{116.097}{then}%
\TextAt{13}{217.67}{81.976}{x}%
\TextAt{21}{224.322}{86.912}{n}%
\TextAt{16}{232.615}{81.976}{\Character{0}}%
\TextAt{13}{244.57}{81.976}{b}%
\TextAt{9}{252.868}{81.976}{\Character{61}}%
\TextAt{21}{265.293}{96.92}{n}%
\TextAt{3}{270.432}{96.92}{\Character{0}}%
\TextAt{17}{277.018}{96.92}{1}%
\TextAt{20}{265.635}{93.334}{\Character{89}}%
\TextAt{21}{266.421}{68.139}{i}%
\TextAt{17}{269.304}{68.139}{\Character{61}0}%
\TextAt{9}{281.252}{81.976}{\Character{40}}%
\TextAt{13}{285.805}{81.976}{x}%
\TextAt{16}{295.113}{81.976}{\Character{0}}%
\TextAt{13}{307.068}{81.976}{\Character{16}}%
\TextAt{21}{313.084}{86.912}{i}%
\TextAt{13}{316.465}{81.976}{\Character{12}}%
\TextAt{9}{323.737}{81.976}{\Character{41}}%
\TextAt{13}{328.289}{81.976}{\Character{59}}%
\TextAt{13}{357.424}{81.976}{\Character{12}}%
\TextAt{21}{364.715}{86.912}{n}%
\TextAt{9}{373.672}{81.976}{\Character{61}}%
\TextAt{13}{386.097}{81.976}{b\Character{59}}%
\TextAt{8}{297.211}{43.033}{306}%
\Destination{79.23}{319.295}{nav.Item.100}%
\Destination{79.23}{289.605}{nav.Item.101}%
\Destination{79.23}{259.915}{nav.Item.102}%
\Destination{79.23}{228.245}{nav.Item.103}%
\Destination{79.23}{198.555}{nav.Item.104}%
\Destination{79.23}{348.975}{nav.Item.99}%
\Destination{78.23}{809.815}{nav.page.306}%
\Destination{79.36}{395.255}{nav.print.334}%
\Destination{79.36}{618.455}{nav.number.5.3.1}%
\Destination{79.36}{368.955}{nav.number.5.3.2}%
\Hotlink{412.502}{150.955}{25.078}{11.047}{nav.number.5.3.1}%
}

% Page 319 (body page 307)
\Page{%
\TextAt{2}{79.226}{796.943}{V.}%
\TextAt{2}{94.073}{796.943}{Local}%
\TextAt{2}{123.615}{796.943}{Class}%
\TextAt{2}{152.208}{796.943}{Field}%
\TextAt{2}{179.939}{796.943}{Theory}%
\Vector{q 0 G 0 g 0.299 w 7.226 18.729 m 460.769 18.729 l S Q}%
\TextAt{8}{79.226}{761.124}{for}%
\TextAt{8}{96.527}{761.124}{some}%
\TextAt{8}{125.268}{761.124}{primitive}%
\TextAt{13}{174.764}{761.124}{n}%
\TextAt{8}{181.751}{761.124}{-th}%
\TextAt{8}{200.02}{761.124}{root}%
\TextAt{8}{224.482}{761.124}{of}%
\TextAt{8}{237.227}{761.124}{unity}%
\TextAt{13}{267.212}{761.124}{\Character{16}}%
\TextAt{8}{273.228}{761.124}{.}%
\TextAt{8}{281.481}{761.124}{Let}%
\TextAt{13}{301.866}{761.124}{d}%
\TextAt{8}{311.262}{761.124}{be}%
\TextAt{8}{326.601}{761.124}{the}%
\TextAt{8}{346.173}{761.124}{greatest}%
\TextAt{8}{389.883}{761.124}{divisor}%
\TextAt{8}{427.411}{761.124}{of}%
\TextAt{13}{440.157}{761.124}{n}%
\TextAt{8}{450.457}{761.124}{such}%
\TextAt{8}{476.27}{761.124}{that}%
\TextAt{13}{501.054}{761.124}{y}%
\TextAt{21}{507.191}{765.462}{d}%
\TextAt{9}{515.367}{761.124}{\Character{61}}%
\TextAt{13}{527.792}{761.124}{b}%
\TextAt{8}{79.226}{744.078}{has}%
\TextAt{8}{100.328}{744.078}{a}%
\TextAt{8}{110.311}{744.078}{solution}%
\TextAt{8}{154.833}{744.078}{in}%
\TextAt{13}{168.714}{744.078}{K}%
\TextAt{8}{179.531}{744.078}{,}%
\TextAt{8}{186.965}{744.078}{and}%
\TextAt{8}{209.956}{744.078}{let}%
\TextAt{13}{227.088}{744.078}{n}%
\TextAt{9}{237.785}{744.078}{\Character{61}}%
\TextAt{13}{250.599}{744.078}{dm}%
\TextAt{8}{266.921}{744.078}{.}%
\TextAt{8}{276.053}{744.078}{Then}%
\TextAt{8}{306.839}{744.078}{the}%
\TextAt{8}{327.223}{744.078}{extension}%
\TextAt{13}{379.217}{744.078}{K}%
\TextAt{9}{390.034}{744.078}{\Character{40}}%
\TextAt{13}{394.587}{744.078}{\Character{12}}%
\TextAt{9}{401.858}{744.078}{\Character{41}}%
\TextAt{16}{406.41}{744.078}{\Character{106}}%
\TextAt{13}{409.731}{744.078}{K}%
\TextAt{8}{424.674}{744.078}{is}%
\TextAt{8}{436.666}{744.078}{cyclic}%
\TextAt{8}{469.077}{744.078}{of}%
\TextAt{8}{482.635}{744.078}{degree}%
\TextAt{13}{519.278}{744.078}{m}%
\TextAt{8}{529.518}{744.078}{,}%
\TextAt{8}{79.226}{727.032}{and}%
\TextAt{8}{102.243}{727.032}{the}%
\TextAt{8}{122.653}{727.032}{conjugates}%
\TextAt{8}{180.531}{727.032}{of}%
\TextAt{13}{194.115}{727.032}{x}%
\TextAt{16}{203.597}{727.032}{\Character{0}}%
\TextAt{13}{215.725}{727.032}{\Character{16}}%
\TextAt{21}{221.74}{731.37}{i}%
\TextAt{13}{225.122}{727.032}{\Character{12}}%
\TextAt{8}{236.544}{727.032}{are}%
\TextAt{8}{256.309}{727.032}{the}%
\TextAt{8}{276.719}{727.032}{elements}%
\TextAt{13}{324.83}{727.032}{x}%
\TextAt{16}{334.312}{727.032}{\Character{0}}%
\TextAt{13}{346.44}{727.032}{\Character{16}}%
\TextAt{21}{352.455}{731.37}{j}%
\TextAt{13}{356.838}{727.032}{\Character{12}}%
\TextAt{8}{368.26}{727.032}{such}%
\TextAt{8}{394.911}{727.032}{that}%
\TextAt{13}{420.534}{727.032}{j}%
\NativeMathAt{429.803}{727.032}{11.955}{\equiv}%
\TextAt{13}{442.855}{727.032}{i}%
\TextAt{8}{450.169}{727.032}{mod}%
\TextAt{13}{475.607}{727.032}{d}%
\TextAt{8}{481.69}{727.032}{.}%
\TextAt{8}{490.9}{727.032}{We}%
\TextAt{8}{511.298}{727.032}{may}%
\TextAt{8}{79.226}{709.986}{therefore}%
\TextAt{8}{128.326}{709.986}{write}%
\TextAt{13}{228.717}{687.196}{x}%
\TextAt{21}{235.369}{692.132}{n}%
\TextAt{16}{243.662}{687.196}{\Character{0}}%
\TextAt{13}{255.617}{687.196}{b}%
\TextAt{9}{263.915}{687.196}{\Character{61}}%
\TextAt{21}{276.39}{702.14}{d}%
\TextAt{3}{280.747}{702.14}{\Character{0}}%
\TextAt{17}{287.334}{702.14}{1}%
\TextAt{20}{276.341}{698.553}{\Character{89}}%
\TextAt{21}{277.127}{673.359}{i}%
\TextAt{17}{280.01}{673.359}{\Character{61}0}%
\TextAt{13}{293.609}{687.196}{N}%
\TextAt{21}{302.983}{685.336}{K}%
\TextAt{17}{310.688}{685.336}{\Character{40}}%
\TextAt{21}{313.981}{685.336}{\Character{12}}%
\TextAt{17}{319.142}{685.336}{\Character{41}}%
\TextAt{3}{322.435}{685.336}{\Character{106}}%
\TextAt{21}{324.788}{685.336}{K}%
\TextAt{9}{332.99}{687.196}{\Character{40}}%
\TextAt{13}{337.543}{687.196}{x}%
\TextAt{16}{346.851}{687.196}{\Character{0}}%
\TextAt{13}{358.807}{687.196}{\Character{16}}%
\TextAt{21}{364.822}{692.132}{i}%
\TextAt{13}{368.203}{687.196}{\Character{12}}%
\TextAt{9}{375.475}{687.196}{\Character{41}}%
\TextAt{13}{380.027}{687.196}{\Character{58}}%
\TextAt{8}{79.226}{653.857}{Hence}%
\TextAt{13}{114.004}{653.857}{x}%
\TextAt{21}{120.656}{658.195}{n}%
\TextAt{16}{128.949}{653.857}{\Character{0}}%
\TextAt{13}{140.904}{653.857}{b}%
\TextAt{8}{149.778}{653.857}{is}%
\TextAt{8}{161.542}{653.857}{a}%
\TextAt{8}{171.298}{653.857}{norm}%
\TextAt{8}{201.867}{653.857}{from}%
\TextAt{13}{229.508}{653.857}{K}%
\TextAt{9}{240.325}{653.857}{\Character{40}}%
\TextAt{22}{248.198}{659.771}{n}%
\TextAt{16}{246.157}{664.045}{\Character{112}}%
\Vector{q 0 G 0 g 0.478 w 184.12 -109.716 m 189.097 -109.716 l S Q}%
\TextAt{13}{256.12}{653.857}{b}%
\TextAt{9}{261.097}{653.857}{\Character{41}}%
\TextAt{16}{265.65}{653.857}{\Character{106}}%
\TextAt{13}{268.97}{653.857}{K}%
\TextAt{8}{279.788}{653.857}{,}%
\TextAt{8}{286.937}{653.857}{i.e.,}%
\TextAt{20}{262.527}{638.162}{\Character{18}}%
\TextAt{13}{272.523}{629.393}{x}%
\TextAt{21}{279.175}{633.732}{n}%
\TextAt{16}{287.468}{629.393}{\Character{0}}%
\TextAt{13}{299.423}{629.393}{b\Character{59}}%
\TextAt{13}{309.621}{629.393}{b}%
\Vector{q 0 G 0 g 0.478 w 200.523 -149.706 m 242.622 -149.706 l S Q}%
\TextAt{14}{290.581}{613.105}{p}%
\TextAt{20}{315.817}{638.162}{\Character{19}}%
\TextAt{9}{327.938}{621.306}{\Character{61}}%
\TextAt{9}{340.348}{621.306}{1}%
\TextAt{13}{346.217}{621.306}{\Character{58}}%
\TextAt{8}{79.226}{588.754}{Choosing}%
\TextAt{13}{130.346}{588.754}{x}%
\TextAt{9}{140.319}{588.754}{\Character{61}}%
\TextAt{9}{152.729}{588.754}{1}%
\TextAt{8}{158.598}{588.754}{,}%
\TextAt{13}{165.747}{588.754}{b}%
\TextAt{9}{174.045}{588.754}{\Character{61}}%
\TextAt{9}{186.455}{588.754}{1}%
\TextAt{16}{194.98}{588.754}{\Character{0}}%
\TextAt{13}{206.935}{588.754}{a}%
\TextAt{8}{213.08}{588.754}{,}%
\TextAt{8}{220.229}{588.754}{and}%
\TextAt{13}{242.992}{588.754}{x}%
\TextAt{9}{252.965}{588.754}{\Character{61}}%
\TextAt{9}{265.375}{588.754}{0}%
\TextAt{8}{271.243}{588.754}{,}%
\TextAt{13}{278.393}{588.754}{b}%
\TextAt{9}{286.691}{588.754}{\Character{61}}%
\TextAt{16}{299.116}{588.754}{\Character{0}}%
\TextAt{13}{308.415}{588.754}{a}%
\TextAt{8}{318.457}{588.754}{then}%
\TextAt{8}{345.117}{588.754}{yield}%
\TextAt{8}{373.403}{588.754}{(v).}%
\TextAt{8}{397.142}{588.754}{(iv)}%
\TextAt{8}{419.582}{588.754}{finally}%
\TextAt{8}{455.029}{588.754}{follows}%
\TextAt{8}{493.465}{588.754}{from}%
\TextAt{20}{159.869}{573.048}{\Character{18}}%
\TextAt{13}{169.865}{564.279}{a\Character{59}}%
\TextAt{13}{181.234}{564.279}{b}%
\Vector{q 0 G 0 g 0.478 w 97.865 -214.82 m 114.231 -214.82 l S Q}%
\TextAt{14}{175.056}{547.99}{p}%
\TextAt{20}{187.426}{573.048}{\Character{19}}%
\TextAt{20}{198.21}{573.048}{\Character{18}}%
\TextAt{13}{208.215}{564.279}{b\Character{59}}%
\TextAt{13}{218.413}{564.279}{a}%
\Vector{q 0 G 0 g 0.478 w 136.215 -214.82 m 152.581 -214.82 l S Q}%
\TextAt{14}{213.407}{547.99}{p}%
\TextAt{20}{225.777}{573.048}{\Character{19}}%
\TextAt{9}{237.898}{556.191}{\Character{61}}%
\TextAt{20}{250.324}{573.048}{\Character{18}}%
\TextAt{13}{260.319}{564.279}{a\Character{59}}%
\TextAt{16}{271.709}{564.279}{\Character{0}}%
\TextAt{13}{281.007}{564.279}{a}%
\Vector{q 0 G 0 g 0.478 w 188.319 -214.82 m 215.152 -214.82 l S Q}%
\TextAt{14}{270.744}{547.99}{p}%
\TextAt{20}{288.347}{573.048}{\Character{19}}%
\TextAt{20}{299.131}{573.048}{\Character{18}}%
\TextAt{13}{309.136}{564.279}{a\Character{59}}%
\TextAt{13}{320.505}{564.279}{b}%
\Vector{q 0 G 0 g 0.478 w 237.136 -214.82 m 253.502 -214.82 l S Q}%
\TextAt{14}{314.328}{547.99}{p}%
\TextAt{20}{326.698}{573.048}{\Character{19}}%
\TextAt{20}{337.482}{573.048}{\Character{18}}%
\TextAt{13}{347.487}{564.279}{b\Character{59}}%
\TextAt{13}{357.685}{564.279}{a}%
\Vector{q 0 G 0 g 0.478 w 275.487 -214.82 m 291.853 -214.82 l S Q}%
\TextAt{14}{352.678}{547.99}{p}%
\TextAt{20}{365.048}{573.048}{\Character{19}}%
\TextAt{20}{375.832}{573.048}{\Character{18}}%
\TextAt{13}{385.837}{564.279}{b\Character{59}}%
\TextAt{16}{396.058}{564.279}{\Character{0}}%
\TextAt{13}{405.357}{564.279}{b}%
\Vector{q 0 G 0 g 0.478 w 313.837 -214.82 m 338.334 -214.82 l S Q}%
\TextAt{14}{395.094}{547.99}{p}%
\TextAt{20}{411.53}{573.048}{\Character{19}}%
\TextAt{9}{237.898}{523.513}{\Character{61}}%
\TextAt{20}{250.324}{540.37}{\Character{18}}%
\TextAt{13}{260.319}{531.601}{a\Character{59}}%
\TextAt{16}{271.709}{531.601}{\Character{0}}%
\TextAt{13}{281.007}{531.601}{ab}%
\Vector{q 0 G 0 g 0.478 w 188.319 -247.498 m 220.129 -247.498 l S Q}%
\TextAt{14}{273.233}{515.313}{p}%
\TextAt{20}{293.325}{540.37}{\Character{19}}%
\TextAt{20}{304.109}{540.37}{\Character{18}}%
\TextAt{13}{314.113}{531.601}{b\Character{59}}%
\TextAt{16}{324.335}{531.601}{\Character{0}}%
\TextAt{13}{333.633}{531.601}{ab}%
\Vector{q 0 G 0 g 0.478 w 242.113 -247.498 m 272.755 -247.498 l S Q}%
\TextAt{14}{326.443}{515.313}{p}%
\TextAt{20}{345.951}{540.37}{\Character{19}}%
\TextAt{9}{358.072}{523.513}{\Character{61}}%
\TextAt{20}{370.497}{540.37}{\Character{18}}%
\TextAt{13}{380.493}{531.601}{ab\Character{59}}%
\TextAt{16}{396.859}{531.601}{\Character{0}}%
\TextAt{13}{406.158}{531.601}{ab}%
\Vector{q 0 G 0 g 0.478 w 308.493 -247.498 m 345.28 -247.498 l S Q}%
\TextAt{14}{395.895}{515.313}{p}%
\TextAt{20}{418.475}{540.37}{\Character{19}}%
\TextAt{9}{430.597}{523.513}{\Character{61}}%
\TextAt{9}{443.007}{523.513}{1}%
\TextAt{13}{448.875}{523.513}{\Character{58}}%
\Vector{q 0 G 0 g 0.399 w 452.334 -282.141 m 452.334 -274.072 l S Q}%
\Vector{q 0 G 0 g 0.399 w 452.534 -274.271 m 459.707 -274.271 l S Q}%
\Vector{q 0 G 0 g 0.399 w 452.534 -281.942 m 459.707 -281.942 l S Q}%
\Vector{q 0 G 0 g 0.399 w 459.906 -282.141 m 459.906 -274.072 l S Q}%
\TextAt{8}{97.159}{463.064}{In}%
\TextAt{8}{111.828}{463.064}{the}%
\TextAt{8}{132.031}{463.064}{case}%
\TextAt{13}{156.85}{463.064}{K}%
\TextAt{9}{171.07}{463.064}{\Character{61}}%
\TextAt{18}{183.576}{463.064}{R}%
\TextAt{8}{196.155}{463.064}{we}%
\TextAt{8}{213.43}{463.064}{have}%
\TextAt{13}{240.472}{463.064}{n}%
\TextAt{9}{250.862}{463.064}{\Character{61}}%
\TextAt{9}{263.355}{463.064}{1}%
\TextAt{8}{273.166}{463.064}{or}%
\TextAt{13}{287.524}{463.064}{n}%
\TextAt{9}{297.913}{463.064}{\Character{61}}%
\TextAt{9}{310.406}{463.064}{2}%
\TextAt{8}{316.273}{463.064}{.}%
\TextAt{8}{324.864}{463.064}{For}%
\TextAt{13}{345.88}{463.064}{n}%
\TextAt{9}{356.27}{463.064}{\Character{61}}%
\TextAt{9}{368.763}{463.064}{1}%
\TextAt{8}{378.574}{463.064}{one}%
\TextAt{8}{400.081}{463.064}{finds,}%
\TextAt{8}{431.415}{463.064}{of}%
\TextAt{8}{444.793}{463.064}{course,}%
\TextAt{20}{483.934}{472.748}{\Character{0}}%
\TextAt{21}{490.609}{468.32}{a\Character{59}b}%
\Vector{q 0 G 0 g 0.478 w 418.609 -307.947 m 429.081 -307.947 l S Q}%
\TextAt{27}{493.851}{458.942}{p}%
\TextAt{20}{502.277}{472.748}{\Character{1}}%
\TextAt{9}{511.158}{463.064}{\Character{61}}%
\TextAt{9}{523.651}{463.064}{1}%
\TextAt{8}{529.518}{463.064}{,}%
\TextAt{8}{79.226}{446.018}{and}%
\TextAt{8}{101.989}{446.018}{for}%
\TextAt{13}{119.874}{446.018}{n}%
\TextAt{9}{130.182}{446.018}{\Character{61}}%
\TextAt{9}{142.592}{446.018}{2}%
\TextAt{8}{152.358}{446.018}{we}%
\TextAt{8}{169.585}{446.018}{have}%
\TextAt{20}{237.429}{444.424}{\Character{18}}%
\TextAt{13}{247.425}{435.655}{a\Character{59}}%
\TextAt{13}{258.794}{435.655}{b}%
\Vector{q 0 G 0 g 0.478 w 175.425 -343.444 m 191.791 -343.444 l S Q}%
\TextAt{14}{252.617}{419.367}{p}%
\TextAt{20}{264.987}{444.424}{\Character{19}}%
\TextAt{9}{277.108}{427.568}{\Character{61}}%
\TextAt{9}{289.518}{427.568}{\Character{40}}%
\TextAt{16}{294.086}{427.568}{\Character{0}}%
\TextAt{9}{303.384}{427.568}{1\Character{41}}%
\TextAt{40}{314.985}{436.26}{sgn}%
\TextAt{22}{326.868}{436.26}{a}%
\TextAt{23}{330.903}{436.26}{\Character{0}}%
\TextAt{24}{336.659}{436.26}{1}%
\Vector{q 0 G 0 g 0.359 w 242.985 -339.504 m 268.312 -339.504 l S Q}%
\TextAt{24}{325.822}{429.378}{2}%
\TextAt{3}{341.507}{432.504}{\Character{1}}%
\TextAt{40}{345.055}{436.26}{sgn}%
\TextAt{22}{356.938}{436.26}{b}%
\TextAt{23}{360.212}{436.26}{\Character{0}}%
\TextAt{24}{365.968}{436.26}{1}%
\Vector{q 0 G 0 g 0.359 w 273.055 -339.504 m 297.621 -339.504 l S Q}%
\TextAt{24}{355.512}{429.378}{2}%
\TextAt{13}{371.315}{427.568}{\Character{59}}%
\TextAt{8}{79.226}{397.904}{because}%
\TextAt{9}{123.41}{397.904}{\Character{40}}%
\TextAt{13}{127.962}{397.904}{a\Character{59}}%
\TextAt{18}{139.351}{397.904}{R}%
\TextAt{9}{147.986}{397.904}{\Character{40}}%
\TextAt{16}{152.538}{408.092}{\Character{112}}%
\Vector{q 0 G 0 g 0.478 w 90.501 -365.669 m 95.478 -365.669 l S Q}%
\TextAt{13}{162.501}{397.904}{b}%
\TextAt{9}{167.478}{397.904}{\Character{41}}%
\TextAt{16}{172.03}{397.904}{\Character{106}}%
\TextAt{18}{175.351}{397.904}{R}%
\TextAt{9}{183.985}{397.904}{\Character{41}}%
\TextAt{9}{193.346}{397.904}{\Character{61}}%
\TextAt{9}{207.274}{397.904}{1}%
\TextAt{8}{217.924}{397.904}{for}%
\TextAt{13}{236.691}{397.904}{b}%
\TextAt{13}{246.482}{397.904}{\Character{62}}%
\TextAt{9}{260.422}{397.904}{0}%
\TextAt{8}{266.275}{397.904}{,}%
\TextAt{8}{274.527}{397.904}{and}%
\TextAt{9}{298.172}{397.904}{\Character{61}}%
\TextAt{9}{312.088}{397.904}{\Character{40}}%
\TextAt{16}{316.653}{397.904}{\Character{0}}%
\TextAt{9}{325.952}{397.904}{1\Character{41}}%
\TextAt{17}{336.357}{402.242}{\Character{40}}%
\TextAt{7}{339.65}{402.242}{sgn}%
\TextAt{21}{353.336}{402.242}{a}%
\TextAt{3}{357.834}{402.242}{\Character{0}}%
\TextAt{17}{364.42}{402.242}{1\Character{41}\Character{47}2}%
\TextAt{8}{385.694}{397.904}{for}%
\TextAt{13}{404.461}{397.904}{b}%
\TextAt{13}{414.252}{397.904}{\Character{60}}%
\TextAt{9}{428.191}{397.904}{0}%
\TextAt{8}{434.044}{397.904}{.}%
\TextAt{8}{445.14}{397.904}{Here}%
\TextAt{8}{473.65}{397.904}{the}%
\TextAt{8}{494.689}{397.904}{letter}%
\TextAt{14}{526.786}{397.904}{p}%
\TextAt{8}{79.226}{380.858}{symbolically}%
\TextAt{8}{146.283}{380.858}{stands}%
\TextAt{8}{182.83}{380.858}{for}%
\TextAt{8}{200.715}{380.858}{an}%
\TextAt{8}{216.974}{380.858}{infinite}%
\TextAt{8}{256.641}{380.858}{place.}%
\TextAt{8}{97.159}{360.823}{Next}%
\TextAt{8}{126.179}{360.823}{we}%
\TextAt{8}{143.818}{360.823}{determine}%
\TextAt{8}{198.852}{360.823}{the}%
\TextAt{8}{219.42}{360.823}{Hilbert}%
\TextAt{8}{260.144}{360.823}{symbol}%
\TextAt{8}{300.617}{360.823}{explicitly}%
\TextAt{8}{351.755}{360.823}{in}%
\TextAt{8}{365.819}{360.823}{the}%
\TextAt{8}{386.387}{360.823}{case}%
\TextAt{8}{411.569}{360.823}{where}%
\TextAt{13}{445.79}{360.823}{K}%
\TextAt{8}{460.916}{360.823}{is}%
\TextAt{8}{473.091}{360.823}{a}%
\TextAt{8}{483.258}{360.823}{local}%
\TextAt{8}{511.31}{360.823}{field}%
\TextAt{8}{79.226}{343.777}{(}%
\NativeMathAt{83.781}{343.777}{11.955}{\neq}%
\TextAt{18}{97.415}{343.777}{R}%
\TextAt{13}{106.049}{343.777}{\Character{59}}%
\TextAt{18}{111.293}{343.777}{C}%
\TextAt{8}{119.927}{343.777}{)}%
\TextAt{8}{129.088}{343.777}{whose}%
\TextAt{8}{164.324}{343.777}{residue}%
\TextAt{8}{204.76}{343.777}{characteristic}%
\TextAt{13}{277.403}{343.777}{p}%
\TextAt{8}{287.884}{343.777}{does}%
\TextAt{8}{314.99}{343.777}{not}%
\TextAt{8}{336.513}{343.777}{divide}%
\TextAt{13}{372.011}{343.777}{n}%
\TextAt{8}{378.999}{343.777}{.}%
\TextAt{8}{389.573}{343.777}{We}%
\TextAt{8}{410.427}{343.777}{call}%
\TextAt{8}{432.595}{343.777}{this}%
\TextAt{8}{456.126}{343.777}{the}%
\TextAt{8}{476.992}{343.777}{case}%
\TextAt{8}{502.472}{343.777}{of}%
\TextAt{8}{516.51}{343.777}{the}%
\TextAt{12}{79.226}{326.731}{tame}%
\TextAt{12}{112.081}{326.731}{Hilbert}%
\TextAt{12}{158.529}{326.731}{symbol}%
\TextAt{8}{200.097}{326.731}{.}%
\TextAt{8}{208.328}{326.731}{Since}%
\TextAt{13}{238.232}{326.731}{\Character{22}}%
\TextAt{21}{245.275}{324.938}{n}%
\NativeMathAt{254.233}{326.731}{11.955}{\subseteq}%
\TextAt{13}{266.852}{326.731}{\Character{22}}%
\TextAt{21}{273.895}{324.938}{q}%
\TextAt{3}{277.962}{324.938}{\Character{0}}%
\TextAt{17}{284.548}{324.938}{1}%
\TextAt{8}{292.525}{326.731}{one}%
\TextAt{8}{313.332}{326.731}{has}%
\TextAt{13}{333.552}{326.731}{n}%
\TextAt{16}{343.861}{326.731}{\Character{106}}%
\TextAt{13}{350.503}{326.731}{q}%
\TextAt{16}{357.443}{326.731}{\Character{0}}%
\TextAt{9}{368.063}{326.731}{1}%
\TextAt{8}{377.161}{326.731}{in}%
\TextAt{8}{390.161}{326.731}{that}%
\TextAt{8}{414.877}{326.731}{case.}%
\TextAt{8}{443.981}{326.731}{First}%
\TextAt{8}{471.841}{326.731}{we}%
\TextAt{8}{488.416}{326.731}{establish}%
\TextAt{8}{79.226}{309.685}{the}%
\TextAt{12}{79.226}{283.879}{Lemma}%
\TextAt{12}{126.621}{283.879}{(3.3).}%
\TextAt{8}{161.702}{283.879}{Let}%
\TextAt{9}{182.436}{283.879}{\Character{40}}%
\TextAt{13}{186.989}{283.879}{n\Character{59}}%
\TextAt{13}{199.195}{283.879}{p}%
\TextAt{9}{205.095}{283.879}{\Character{41}}%
\TextAt{9}{212.95}{283.879}{\Character{61}}%
\TextAt{9}{225.371}{283.879}{1}%
\TextAt{8}{234.909}{283.879}{and}%
\TextAt{13}{257.437}{283.879}{x}%
\NativeMathAt{267.41}{283.879}{11.955}{\in}%
\TextAt{13}{278.701}{283.879}{K}%
\TextAt{3}{289.519}{288.218}{\Character{3}}%
\TextAt{8}{294.251}{283.879}{.}%
\TextAt{8}{302.621}{283.879}{The}%
\TextAt{8}{326.44}{283.879}{extension}%
\TextAt{13}{377.971}{283.879}{K}%
\TextAt{9}{388.789}{283.879}{\Character{40}}%
\TextAt{22}{396.662}{287.9}{n}%
\TextAt{16}{394.621}{292.49}{\Character{112}}%
\Vector{q 0 G 0 g 0.478 w 332.583 -481.27 m 339.236 -481.27 l S Q}%
\TextAt{13}{404.583}{283.879}{x}%
\TextAt{9}{411.236}{283.879}{\Character{41}}%
\TextAt{16}{415.788}{283.879}{\Character{106}}%
\TextAt{13}{419.109}{283.879}{K}%
\TextAt{8}{433.589}{283.879}{is}%
\TextAt{8}{445.118}{283.879}{unramified}%
\TextAt{8}{503.415}{283.879}{if}%
\TextAt{8}{513.904}{283.879}{and}%
\TextAt{8}{79.226}{266.833}{only}%
\TextAt{8}{104.918}{266.833}{if}%
\TextAt{13}{115.642}{266.833}{x}%
\NativeMathAt{125.615}{266.833}{11.955}{\in}%
\TextAt{13}{136.906}{266.833}{U}%
\TextAt{21}{144.856}{265.04}{K}%
\TextAt{13}{153.058}{266.833}{K}%
\TextAt{3}{163.876}{271.172}{\Character{3}}%
\TextAt{21}{168.11}{271.172}{n}%
\TextAt{8}{173.747}{266.833}{.}%
\TextAt{19}{79.226}{232.268}{Proof.}%
\TextAt{8}{116}{232.268}{Let}%
\TextAt{13}{138.083}{232.268}{x}%
\TextAt{9}{149.955}{232.268}{\Character{61}}%
\TextAt{13}{164.278}{232.268}{uy}%
\TextAt{21}{177.077}{236.607}{n}%
\TextAt{8}{187.725}{232.268}{with}%
\TextAt{13}{215.499}{232.268}{u}%
\NativeMathAt{227.38}{232.268}{11.955}{\in}%
\TextAt{13}{240.57}{232.268}{U}%
\TextAt{21}{248.52}{230.475}{K}%
\TextAt{8}{256.722}{232.268}{,}%
\TextAt{13}{265.264}{232.268}{y}%
\NativeMathAt{276.62}{232.268}{11.955}{\in}%
\TextAt{13}{289.809}{232.268}{K}%
\TextAt{3}{300.627}{236.607}{\Character{3}}%
\TextAt{8}{305.359}{232.268}{,}%
\TextAt{8}{313.9}{232.268}{so}%
\TextAt{8}{329.384}{232.268}{that}%
\TextAt{13}{355.867}{232.268}{K}%
\TextAt{9}{366.685}{232.268}{\Character{40}}%
\TextAt{22}{374.558}{236.289}{n}%
\TextAt{16}{372.517}{240.88}{\Character{112}}%
\Vector{q 0 G 0 g 0.478 w 310.48 -532.881 m 317.132 -532.881 l S Q}%
\TextAt{13}{382.48}{232.268}{x}%
\TextAt{9}{389.132}{232.268}{\Character{41}}%
\TextAt{9}{398.887}{232.268}{\Character{61}}%
\TextAt{13}{413.227}{232.268}{K}%
\TextAt{9}{424.045}{232.268}{\Character{40}}%
\TextAt{22}{431.918}{236.289}{n}%
\TextAt{16}{429.877}{240.88}{\Character{112}}%
\Vector{q 0 G 0 g 0.478 w 367.839 -532.881 m 374.502 -532.881 l S Q}%
\TextAt{13}{439.839}{232.268}{u}%
\TextAt{9}{446.502}{232.268}{\Character{41}}%
\TextAt{8}{451.054}{232.268}{.}%
\TextAt{8}{462.844}{232.268}{Let}%
\TextAt{13}{484.927}{232.268}{\Character{20}}%
\TextAt{3}{491.666}{236.607}{\Character{48}}%
\TextAt{8}{499.472}{232.268}{be}%
\TextAt{8}{516.511}{232.268}{the}%
\TextAt{8}{79.226}{215.223}{splitting}%
\TextAt{8}{125.062}{215.223}{field}%
\TextAt{8}{150.013}{215.223}{of}%
\TextAt{8}{162.936}{215.223}{the}%
\TextAt{8}{182.686}{215.223}{polynomial}%
\TextAt{13}{242.773}{215.223}{X}%
\TextAt{21}{253.428}{219.561}{n}%
\TextAt{16}{260.89}{215.223}{\Character{0}}%
\TextAt{13}{272.014}{215.223}{u}%
\TextAt{8}{281.998}{215.223}{mod}%
\TextAt{14}{307.436}{215.223}{p}%
\TextAt{8}{316.91}{215.223}{over}%
\TextAt{8}{341.549}{215.223}{the}%
\TextAt{8}{361.299}{215.223}{residue}%
\TextAt{8}{400.62}{215.223}{class}%
\TextAt{8}{427.651}{215.223}{field}%
\TextAt{13}{452.601}{215.223}{\Character{20}}%
\TextAt{8}{459.34}{215.223}{,}%
\TextAt{8}{466.164}{215.223}{and}%
\TextAt{8}{488.52}{215.223}{let}%
\TextAt{13}{505.018}{215.223}{K}%
\TextAt{3}{515.836}{219.561}{\Character{48}}%
\TextAt{16}{518.631}{215.223}{\Character{106}}%
\TextAt{13}{521.952}{215.223}{K}%
\TextAt{8}{79.226}{198.177}{be}%
\TextAt{8}{96.294}{198.177}{the}%
\TextAt{8}{117.594}{198.177}{unramified}%
\TextAt{8}{177.27}{198.177}{extension}%
\TextAt{8}{230.179}{198.177}{with}%
\TextAt{8}{257.983}{198.177}{residue}%
\TextAt{8}{298.853}{198.177}{class}%
\TextAt{8}{327.434}{198.177}{field}%
\TextAt{13}{353.934}{198.177}{\Character{20}}%
\TextAt{3}{360.673}{202.515}{\Character{48}}%
\TextAt{8}{368.509}{198.177}{(see}%
\TextAt{8}{393.12}{198.177}{chap.}%
\TextAt{8}{428.742}{198.177}{II,}%
\TextAt{8}{445.798}{198.177}{§9,}%
\TextAt{8}{465.913}{198.177}{p.}%
\TextAt{8}{484.296}{198.177}{173).}%
\TextAt{8}{518.304}{198.177}{By}%
\TextAt{8}{79.226}{181.131}{Hensel’s}%
\TextAt{8}{124.513}{181.131}{lemma,}%
\TextAt{13}{165.464}{181.131}{X}%
\TextAt{21}{176.119}{185.469}{n}%
\TextAt{16}{184.364}{181.131}{\Character{0}}%
\TextAt{13}{196.271}{181.131}{u}%
\TextAt{8}{206.807}{181.131}{splits}%
\TextAt{8}{237.473}{181.131}{over}%
\TextAt{13}{262.495}{181.131}{K}%
\TextAt{3}{273.313}{185.469}{\Character{48}}%
\TextAt{8}{279.982}{181.131}{into}%
\TextAt{8}{303.701}{181.131}{linear}%
\TextAt{8}{336.196}{181.131}{factors,}%
\TextAt{8}{377.542}{181.131}{so}%
\TextAt{13}{391.888}{181.131}{K}%
\TextAt{9}{402.706}{181.131}{\Character{40}}%
\TextAt{22}{410.579}{185.151}{n}%
\TextAt{16}{408.538}{189.742}{\Character{112}}%
\Vector{q 0 G 0 g 0.478 w 346.501 -584.019 m 353.163 -584.019 l S Q}%
\TextAt{13}{418.501}{181.131}{u}%
\TextAt{9}{425.163}{181.131}{\Character{41}}%
\NativeMathAt{433.036}{181.131}{11.955}{\subseteq}%
\TextAt{13}{445.656}{181.131}{K}%
\TextAt{3}{456.473}{185.469}{\Character{48}}%
\TextAt{8}{463.142}{181.131}{is}%
\TextAt{8}{474.883}{181.131}{unramified.}%
\TextAt{8}{79.226}{164.085}{Assume}%
\TextAt{8}{121.841}{164.085}{conversely}%
\TextAt{8}{177.093}{164.085}{that}%
\TextAt{13}{201.715}{164.085}{L}%
\TextAt{9}{213.001}{164.085}{\Character{61}}%
\TextAt{13}{225.426}{164.085}{K}%
\TextAt{9}{236.244}{164.085}{\Character{40}}%
\TextAt{22}{244.117}{168.106}{n}%
\TextAt{16}{242.076}{172.696}{\Character{112}}%
\Vector{q 0 G 0 g 0.478 w 180.039 -601.065 m 186.691 -601.065 l S Q}%
\TextAt{13}{252.039}{164.085}{x}%
\TextAt{9}{258.691}{164.085}{\Character{41}}%
\TextAt{8}{266.394}{164.085}{is}%
\TextAt{8}{277.412}{164.085}{unramified}%
\TextAt{8}{335.198}{164.085}{over}%
\TextAt{13}{359.498}{164.085}{K}%
\TextAt{8}{370.316}{164.085}{,}%
\TextAt{8}{376.868}{164.085}{and}%
\TextAt{8}{398.884}{164.085}{let}%
\TextAt{13}{415.042}{164.085}{x}%
\TextAt{9}{425.015}{164.085}{\Character{61}}%
\TextAt{13}{437.441}{164.085}{u\Character{25}}%
\TextAt{21}{451.173}{168.423}{r}%
\TextAt{8}{455.727}{164.085}{,}%
\TextAt{8}{462.279}{164.085}{where}%
\TextAt{13}{495.342}{164.085}{u}%
\NativeMathAt{505.326}{164.085}{11.955}{\in}%
\TextAt{13}{516.617}{164.085}{U}%
\TextAt{21}{524.567}{162.291}{K}%
\TextAt{8}{79.226}{147.039}{and}%
\TextAt{13}{102.659}{147.039}{\Character{25}}%
\TextAt{8}{114.295}{147.039}{is}%
\TextAt{8}{126.729}{147.039}{a}%
\TextAt{8}{137.154}{147.039}{prime}%
\TextAt{8}{170.988}{147.039}{element}%
\TextAt{8}{214.899}{147.039}{of}%
\TextAt{13}{228.899}{147.039}{K}%
\TextAt{8}{239.717}{147.039}{.}%
\TextAt{8}{250.175}{147.039}{Then}%
\TextAt{13}{281.402}{147.039}{v}%
\TextAt{21}{287.061}{145.245}{L}%
\TextAt{9}{293.318}{147.039}{\Character{40}}%
\TextAt{22}{301.192}{152.102}{n}%
\TextAt{16}{299.151}{156.519}{\Character{112}}%
\Vector{q 0 G 0 g 0.478 w 237.113 -617.242 m 255.4 -617.242 l S Q}%
\TextAt{13}{309.113}{147.039}{u\Character{25}}%
\TextAt{21}{322.845}{150.492}{r}%
\TextAt{9}{327.4}{147.039}{\Character{41}}%
\TextAt{9}{336.402}{147.039}{\Character{61}}%
\TextAt{17}{351.629}{151.746}{1}%
\Vector{q 0 G 0 g 0.478 w 279.177 -623.973 m 284.315 -623.973 l S Q}%
\TextAt{21}{351.177}{142.916}{n}%
\TextAt{13}{357.511}{147.039}{v}%
\TextAt{21}{363.17}{145.245}{L}%
\TextAt{9}{369.427}{147.039}{\Character{40}}%
\TextAt{13}{373.979}{147.039}{\Character{25}}%
\TextAt{21}{381.049}{151.377}{r}%
\TextAt{9}{385.603}{147.039}{\Character{41}}%
\TextAt{9}{394.605}{147.039}{\Character{61}}%
\TextAt{21}{409.921}{151.746}{r}%
\Vector{q 0 G 0 g 0.478 w 337.381 -623.973 m 342.519 -623.973 l S Q}%
\TextAt{21}{409.381}{142.916}{n}%
\NativeMathAt{420.177}{147.039}{11.955}{\in}%
\TextAt{18}{432.609}{147.039}{Z}%
\TextAt{8}{440.579}{147.039}{,}%
\TextAt{8}{448.566}{147.039}{hence}%
\TextAt{13}{481.742}{147.039}{n}%
\TextAt{16}{493.192}{147.039}{\Character{106}}%
\TextAt{13}{500.975}{147.039}{r}%
\TextAt{8}{506.575}{147.039}{,}%
\TextAt{8}{514.562}{147.039}{i.e.,}%
\TextAt{13}{79.226}{129.993}{\Character{25}}%
\TextAt{21}{86.295}{134.331}{r}%
\NativeMathAt{94.171}{129.993}{11.955}{\in}%
\TextAt{13}{105.462}{129.993}{K}%
\TextAt{3}{116.28}{134.331}{\Character{3}}%
\TextAt{21}{120.514}{134.331}{n}%
\TextAt{8}{126.15}{129.993}{,}%
\TextAt{8}{133.299}{129.993}{and}%
\TextAt{8}{156.062}{129.993}{thus}%
\TextAt{13}{181.813}{129.993}{x}%
\NativeMathAt{191.786}{129.993}{11.955}{\in}%
\TextAt{13}{203.077}{129.993}{U}%
\TextAt{21}{211.028}{128.199}{K}%
\TextAt{13}{219.23}{129.993}{K}%
\TextAt{3}{230.048}{134.331}{\Character{3}}%
\TextAt{21}{234.282}{134.331}{n}%
\TextAt{8}{239.918}{129.993}{.}%
\Vector{q 0 G 0 g 0.399 w 452.334 -644.007 m 452.334 -635.938 l S Q}%
\Vector{q 0 G 0 g 0.399 w 452.534 -636.137 m 459.707 -636.137 l S Q}%
\Vector{q 0 G 0 g 0.399 w 452.534 -643.808 m 459.707 -643.808 l S Q}%
\Vector{q 0 G 0 g 0.399 w 459.906 -644.007 m 459.906 -635.938 l S Q}%
\TextAt{8}{97.159}{101.198}{Since}%
\TextAt{13}{127.716}{101.198}{U}%
\TextAt{21}{135.667}{99.405}{K}%
\TextAt{9}{147.19}{101.198}{\Character{61}}%
\TextAt{13}{159.616}{101.198}{\Character{22}}%
\TextAt{21}{166.659}{99.405}{q}%
\TextAt{3}{170.725}{99.405}{\Character{0}}%
\TextAt{17}{177.312}{99.405}{1}%
\TextAt{16}{184.701}{101.198}{\Character{2}}%
\TextAt{13}{196.656}{101.198}{U}%
\TextAt{17}{205.855}{107.309}{\Character{40}1\Character{41}}%
\TextAt{21}{204.606}{97.957}{K}%
\TextAt{8}{217.174}{101.198}{,}%
\TextAt{8}{224.323}{101.198}{every}%
\TextAt{8}{255.215}{101.198}{unit}%
\TextAt{13}{279.926}{101.198}{u}%
\NativeMathAt{289.91}{101.198}{11.955}{\in}%
\TextAt{13}{301.201}{101.198}{U}%
\TextAt{21}{309.151}{99.405}{K}%
\TextAt{8}{321.251}{101.198}{has}%
\TextAt{8}{342.125}{101.198}{a}%
\TextAt{8}{351.88}{101.198}{unique}%
\TextAt{8}{389.922}{101.198}{decomposition}%
\TextAt{13}{275.07}{72.921}{u}%
\TextAt{9}{285.054}{72.921}{\Character{61}}%
\TextAt{13}{297.479}{72.921}{\Character{33}}%
\TextAt{9}{305.197}{72.921}{\Character{40}}%
\TextAt{13}{309.75}{72.921}{u}%
\TextAt{9}{316.412}{72.921}{\Character{41}}%
\TextAt{16}{320.964}{72.921}{\Character{104}}%
\TextAt{13}{325.614}{72.921}{u}%
\TextAt{16}{332.276}{72.921}{\Character{105}}%
\TextAt{8}{297.211}{43.033}{307}%
\Destination{78.23}{809.815}{nav.page.307}%
\Destination{79.36}{461.055}{nav.print.335}%
\Destination{79.36}{297.855}{nav.number.5.3.3}%
\Hotlink{428.242}{195.955}{9.763}{11.047}{nav.chapter.2}%
\Hotlink{450.977}{195.955}{6.858}{11.047}{nav.section.2.9}%
\Hotlink{483.796}{195.955}{18.574}{11.047}{nav.print.173}%
}

% Page 320 (body page 308)
\Page{%
\TextAt{2}{79.226}{796.943}{V.}%
\TextAt{2}{94.073}{796.943}{Local}%
\TextAt{2}{123.615}{796.943}{Class}%
\TextAt{2}{152.208}{796.943}{Field}%
\TextAt{2}{179.939}{796.943}{Theory}%
\Vector{q 0 G 0 g 0.299 w 7.226 18.729 m 460.769 18.729 l S Q}%
\TextAt{8}{79.226}{760.991}{with}%
\TextAt{13}{105.925}{760.991}{\Character{33}}%
\TextAt{9}{113.643}{760.991}{\Character{40}}%
\TextAt{13}{118.195}{760.991}{u}%
\TextAt{9}{124.858}{760.991}{\Character{41}}%
\NativeMathAt{132.797}{760.991}{11.955}{\in}%
\TextAt{13}{144.154}{760.991}{\Character{22}}%
\TextAt{21}{151.197}{759.198}{q}%
\TextAt{3}{155.264}{759.198}{\Character{0}}%
\TextAt{17}{161.85}{759.198}{1}%
\TextAt{8}{170.519}{760.991}{and}%
\TextAt{16}{193.32}{760.991}{\Character{104}}%
\TextAt{13}{197.969}{760.991}{u}%
\TextAt{16}{204.632}{760.991}{\Character{105}}%
\NativeMathAt{212.654}{760.991}{11.955}{\in}%
\TextAt{13}{224.025}{760.991}{U}%
\TextAt{17}{233.224}{767.101}{\Character{40}1\Character{41}}%
\TextAt{21}{231.975}{757.75}{K}%
\TextAt{8}{244.543}{760.991}{,}%
\TextAt{13}{251.74}{760.991}{u}%
\NativeMathAt{261.79}{760.991}{11.955}{\equiv}%
\TextAt{13}{274.475}{760.991}{\Character{33}}%
\TextAt{9}{282.193}{760.991}{\Character{40}}%
\TextAt{13}{286.745}{760.991}{u}%
\TextAt{9}{293.408}{760.991}{\Character{41}}%
\TextAt{8}{301.281}{760.991}{mod}%
\TextAt{14}{326.719}{760.991}{p}%
\TextAt{8}{332.702}{760.991}{.}%
\TextAt{8}{341.266}{760.991}{With}%
\TextAt{8}{371.54}{760.991}{this}%
\TextAt{8}{394.401}{760.991}{notation}%
\TextAt{8}{441.28}{760.991}{we}%
\TextAt{8}{458.546}{760.991}{will}%
\TextAt{8}{480.69}{760.991}{now}%
\TextAt{8}{505.117}{760.991}{prove}%
\TextAt{8}{79.226}{743.945}{the}%
\TextAt{12}{79.226}{718.577}{Proposition}%
\TextAt{12}{151.937}{718.577}{(3.4).}%
\TextAt{8}{187.253}{718.577}{If}%
\TextAt{9}{198.945}{718.577}{\Character{40}}%
\TextAt{13}{203.498}{718.577}{n\Character{59}}%
\TextAt{13}{215.704}{718.577}{p}%
\TextAt{9}{221.604}{718.577}{\Character{41}}%
\TextAt{9}{229.459}{718.577}{\Character{61}}%
\TextAt{9}{241.88}{718.577}{1}%
\TextAt{8}{251.653}{718.577}{and}%
\TextAt{13}{274.416}{718.577}{a\Character{59}}%
\TextAt{13}{285.785}{718.577}{b}%
\NativeMathAt{294.103}{718.577}{11.955}{\in}%
\TextAt{13}{305.394}{718.577}{K}%
\TextAt{3}{316.212}{722.916}{\Character{3}}%
\TextAt{8}{320.944}{718.577}{,}%
\TextAt{8}{328.093}{718.577}{then}%
\TextAt{20}{223.221}{701.581}{\Character{18}}%
\TextAt{13}{233.217}{692.812}{a\Character{59}}%
\TextAt{13}{244.586}{692.812}{b}%
\Vector{q 0 G 0 g 0.478 w 161.217 -86.287 m 177.583 -86.287 l S Q}%
\TextAt{14}{238.408}{676.524}{p}%
\TextAt{20}{250.778}{701.581}{\Character{19}}%
\TextAt{9}{262.9}{684.724}{\Character{61}}%
\TextAt{13}{275.325}{684.724}{\Character{33}}%
\TextAt{20}{285.036}{701.581}{\Character{18}}%
\TextAt{9}{293.836}{684.724}{\Character{40}}%
\TextAt{16}{298.388}{684.724}{\Character{0}}%
\TextAt{9}{307.687}{684.724}{1\Character{41}}%
\TextAt{21}{318.092}{689.661}{\Character{11}\Character{12}}%
\TextAt{13}{330.831}{692.812}{b}%
\TextAt{21}{335.808}{697.151}{\Character{11}}%
\Vector{q 0 G 0 g 0.478 w 258.386 -86.287 m 270.19 -86.287 l S Q}%
\TextAt{13}{330.386}{676.524}{a}%
\TextAt{21}{336.531}{679.977}{\Character{12}}%
\TextAt{20}{343.386}{701.581}{\Character{19}}%
\TextAt{17}{352.186}{698.905}{\Character{40}}%
\TextAt{21}{355.48}{698.905}{q}%
\TextAt{3}{359.546}{698.905}{\Character{0}}%
\TextAt{17}{366.133}{698.905}{1\Character{41}\Character{47}}%
\TextAt{21}{377.894}{698.905}{n}%
\TextAt{13}{385.523}{684.724}{\Character{59}}%
\TextAt{8}{79.226}{653.706}{where}%
\TextAt{13}{113.035}{653.706}{\Character{11}}%
\TextAt{9}{123.878}{653.706}{\Character{61}}%
\TextAt{13}{136.303}{653.706}{v}%
\TextAt{21}{141.963}{651.913}{K}%
\TextAt{9}{150.165}{653.706}{\Character{40}}%
\TextAt{13}{154.718}{653.706}{a}%
\TextAt{9}{160.863}{653.706}{\Character{41}}%
\TextAt{8}{165.415}{653.706}{,}%
\TextAt{13}{172.564}{653.706}{\Character{12}}%
\TextAt{9}{183.156}{653.706}{\Character{61}}%
\TextAt{13}{195.582}{653.706}{v}%
\TextAt{21}{201.241}{651.913}{K}%
\TextAt{9}{209.444}{653.706}{\Character{40}}%
\TextAt{13}{213.996}{653.706}{b}%
\TextAt{9}{218.973}{653.706}{\Character{41}}%
\TextAt{8}{223.525}{653.706}{.}%
\TextAt{19}{79.226}{620.017}{Proof.}%
\TextAt{8}{116}{620.017}{The}%
\TextAt{8}{140.054}{620.017}{function}%
\TextAt{16}{229.564}{596.424}{\Character{104}}%
\TextAt{13}{234.213}{596.424}{a\Character{59}}%
\TextAt{13}{245.582}{596.424}{b}%
\TextAt{16}{250.579}{596.424}{\Character{105}}%
\TextAt{9}{258.549}{596.424}{\Character{58}\Character{61}}%
\TextAt{13}{274.226}{596.424}{\Character{33}}%
\TextAt{20}{283.937}{613.281}{\Character{18}}%
\TextAt{9}{292.737}{596.424}{\Character{40}}%
\TextAt{16}{297.29}{596.424}{\Character{0}}%
\TextAt{9}{306.588}{596.424}{1\Character{41}}%
\TextAt{21}{316.994}{601.36}{\Character{11}\Character{12}}%
\TextAt{13}{329.732}{604.512}{b}%
\TextAt{21}{334.71}{608.85}{\Character{11}}%
\Vector{q 0 G 0 g 0.478 w 257.287 -174.587 m 269.092 -174.587 l S Q}%
\TextAt{13}{329.287}{588.223}{a}%
\TextAt{21}{335.432}{591.677}{\Character{12}}%
\TextAt{20}{342.287}{613.281}{\Character{19}}%
\TextAt{17}{351.088}{610.604}{\Character{40}}%
\TextAt{21}{354.381}{610.604}{q}%
\TextAt{3}{358.448}{610.604}{\Character{0}}%
\TextAt{17}{365.034}{610.604}{1\Character{41}\Character{47}}%
\TextAt{21}{376.796}{610.604}{n}%
\TextAt{8}{79.226}{569.595}{is}%
\TextAt{8}{91.461}{569.595}{obviously}%
\TextAt{8}{143.71}{569.595}{bilinear}%
\TextAt{8}{186.455}{569.595}{(in}%
\TextAt{8}{205.134}{569.595}{the}%
\TextAt{8}{225.762}{569.595}{multiplicative}%
\TextAt{8}{300.056}{569.595}{sense).}%
\TextAt{8}{340.608}{569.595}{We}%
\TextAt{8}{361.223}{569.595}{may}%
\TextAt{8}{387.064}{569.595}{therefore}%
\TextAt{8}{436.635}{569.595}{assume}%
\TextAt{8}{477.55}{569.595}{that}%
\TextAt{13}{503.391}{569.595}{a}%
\TextAt{8}{513.904}{569.595}{and}%
\TextAt{13}{79.226}{552.55}{b}%
\TextAt{8}{88.437}{552.55}{are}%
\TextAt{8}{108.284}{552.55}{prime}%
\TextAt{8}{141.784}{552.55}{elements:}%
\TextAt{13}{194.864}{552.55}{a}%
\TextAt{9}{204.903}{552.549}{\Character{61}}%
\TextAt{13}{217.901}{552.55}{\Character{25}}%
\TextAt{8}{224.971}{552.55}{,}%
\TextAt{13}{232.54}{552.55}{b}%
\TextAt{9}{241.411}{552.549}{\Character{61}}%
\TextAt{16}{254.41}{552.55}{\Character{0}}%
\TextAt{13}{263.708}{552.549}{\Character{25}}%
\TextAt{13}{270.762}{552.549}{u}%
\TextAt{8}{277.44}{552.55}{,}%
\TextAt{13}{285.01}{552.55}{u}%
\NativeMathAt{295.566}{552.549}{11.955}{\in}%
\TextAt{13}{307.43}{552.55}{U}%
\TextAt{21}{315.38}{550.756}{K}%
\TextAt{8}{323.583}{552.55}{.}%
\TextAt{8}{333.04}{552.55}{Since}%
\TextAt{8}{363.934}{552.55}{clearly}%
\TextAt{16}{401.666}{552.55}{\Character{104}}%
\TextAt{13}{406.315}{552.549}{\Character{25}}%
\TextAt{13}{413.369}{552.549}{\Character{59}}%
\TextAt{16}{418.629}{552.55}{\Character{0}}%
\TextAt{13}{427.927}{552.549}{\Character{25}}%
\TextAt{16}{434.996}{552.55}{\Character{105}}%
\TextAt{9}{443.54}{552.549}{\Character{61}}%
\TextAt{20}{456.538}{562.233}{\Character{0}}%
\TextAt{21}{463.213}{557.805}{\Character{25}}%
\TextAt{21}{468.354}{557.805}{\Character{59}}%
\TextAt{3}{470.714}{557.805}{\Character{0}}%
\TextAt{21}{477.301}{557.805}{\Character{25}}%
\Vector{q 0 G 0 g 0.478 w 391.213 -218.462 m 410.45 -218.462 l S Q}%
\TextAt{27}{470.837}{548.427}{p}%
\TextAt{20}{483.645}{562.233}{\Character{1}}%
\TextAt{9}{493.018}{552.549}{\Character{61}}%
\TextAt{9}{506.001}{552.549}{1}%
\TextAt{8}{511.87}{552.55}{,}%
\TextAt{8}{519.44}{552.55}{we}%
\TextAt{8}{79.226}{535.504}{may}%
\TextAt{8}{104.595}{535.504}{restrict}%
\TextAt{8}{144.98}{535.504}{to}%
\TextAt{8}{159.29}{535.504}{the}%
\TextAt{8}{179.446}{535.504}{case}%
\TextAt{13}{204.217}{535.504}{a}%
\TextAt{9}{213.683}{535.504}{\Character{61}}%
\TextAt{13}{226.109}{535.504}{\Character{25}}%
\TextAt{8}{233.178}{535.504}{,}%
\TextAt{13}{240.327}{535.504}{b}%
\TextAt{9}{248.625}{535.504}{\Character{61}}%
\TextAt{13}{261.051}{535.504}{u}%
\TextAt{8}{267.713}{535.504}{.}%
\TextAt{8}{276.161}{535.504}{Let}%
\TextAt{13}{297.131}{535.504}{y}%
\TextAt{9}{306.588}{535.504}{\Character{61}}%
\TextAt{22}{322.335}{539.524}{n}%
\TextAt{16}{320.294}{544.115}{\Character{112}}%
\Vector{q 0 G 0 g 0.478 w 258.256 -229.646 m 264.919 -229.646 l S Q}%
\TextAt{13}{330.256}{535.504}{u}%
\TextAt{8}{340.816}{535.504}{and}%
\TextAt{13}{363.579}{535.504}{K}%
\TextAt{3}{374.396}{539.842}{\Character{48}}%
\TextAt{9}{380.512}{535.504}{\Character{61}}%
\TextAt{13}{392.938}{535.504}{K}%
\TextAt{9}{403.756}{535.504}{\Character{40}}%
\TextAt{13}{408.308}{535.504}{y}%
\TextAt{9}{414.445}{535.504}{\Character{41}}%
\TextAt{8}{418.997}{535.504}{.}%
\TextAt{8}{427.445}{535.504}{Then}%
\TextAt{8}{458.003}{535.504}{we}%
\TextAt{8}{475.23}{535.504}{have}%
\TextAt{16}{174.484}{504.485}{\Character{104}}%
\TextAt{13}{179.134}{504.485}{\Character{25}}%
\TextAt{13}{186.188}{504.485}{\Character{59}}%
\TextAt{13}{191.424}{504.485}{u}%
\TextAt{16}{198.11}{504.485}{\Character{105}}%
\TextAt{9}{206.08}{504.485}{\Character{61}}%
\TextAt{13}{218.505}{504.485}{\Character{33}}%
\TextAt{9}{226.223}{504.485}{\Character{40}}%
\TextAt{13}{230.776}{504.485}{u}%
\TextAt{9}{237.438}{504.485}{\Character{41}}%
\TextAt{17}{241.99}{509.421}{\Character{40}}%
\TextAt{21}{245.284}{509.421}{q}%
\TextAt{3}{249.35}{509.421}{\Character{0}}%
\TextAt{17}{255.937}{509.421}{1\Character{41}\Character{47}}%
\TextAt{21}{267.699}{509.421}{n}%
\TextAt{8}{285.29}{504.485}{and}%
\TextAt{9}{316.11}{504.485}{\Character{40}}%
\TextAt{13}{320.663}{504.485}{\Character{25}}%
\TextAt{13}{327.717}{504.485}{\Character{59}}%
\TextAt{13}{332.953}{504.485}{K}%
\TextAt{3}{343.794}{509.421}{\Character{48}}%
\TextAt{16}{346.589}{504.485}{\Character{106}}%
\TextAt{13}{349.91}{504.485}{K}%
\TextAt{9}{360.728}{504.485}{\Character{41}}%
\TextAt{13}{365.28}{504.485}{y}%
\TextAt{9}{374.737}{504.485}{\Character{61}}%
\TextAt{20}{387.163}{521.342}{\Character{18}}%
\TextAt{13}{397.159}{512.573}{\Character{25}}%
\TextAt{13}{404.213}{512.573}{\Character{59}}%
\TextAt{13}{409.449}{512.573}{u}%
\Vector{q 0 G 0 g 0.478 w 325.159 -266.526 m 344.135 -266.526 l S Q}%
\TextAt{14}{403.655}{496.285}{p}%
\TextAt{20}{417.33}{521.342}{\Character{19}}%
\TextAt{13}{428.123}{504.485}{y}%
\TextAt{13}{434.244}{504.485}{\Character{58}}%
\TextAt{8}{79.226}{473.467}{By}%
\TextAt{8}{96.448}{473.467}{(3.3),}%
\TextAt{8}{126.762}{473.467}{we}%
\TextAt{8}{142.848}{473.467}{see}%
\TextAt{8}{160.619}{473.467}{that}%
\TextAt{13}{184.847}{473.467}{K}%
\TextAt{3}{195.664}{477.805}{\Character{48}}%
\TextAt{16}{198.459}{473.467}{\Character{106}}%
\TextAt{13}{201.78}{473.467}{K}%
\TextAt{8}{215.354}{473.467}{is}%
\TextAt{8}{225.976}{473.467}{unramified}%
\TextAt{8}{283.368}{473.467}{and}%
\TextAt{8}{304.989}{473.467}{by}%
\TextAt{8}{320.106}{473.467}{chap.}%
\TextAt{8}{351.917}{473.467}{IV,}%
\TextAt{8}{370.92}{473.467}{(5.7),}%
\TextAt{9}{401.234}{473.467}{\Character{40}}%
\TextAt{13}{405.786}{473.467}{\Character{25}}%
\TextAt{13}{412.84}{473.467}{\Character{59}}%
\TextAt{13}{418.076}{473.467}{K}%
\TextAt{3}{428.917}{477.805}{\Character{48}}%
\TextAt{16}{431.713}{473.467}{\Character{106}}%
\TextAt{13}{435.033}{473.467}{K}%
\TextAt{9}{445.851}{473.467}{\Character{41}}%
\TextAt{8}{453.159}{473.467}{is}%
\TextAt{8}{463.782}{473.467}{the}%
\TextAt{8}{482.797}{473.467}{Frobenius}%
\TextAt{8}{79.226}{456.421}{automorphism}%
\TextAt{13}{156.696}{456.421}{\Character{39}}%
\TextAt{9}{167.685}{456.421}{\Character{61}}%
\TextAt{13}{180.11}{456.421}{\Character{39}}%
\TextAt{21}{187.779}{454.561}{K}%
\TextAt{23}{195.483}{456.83}{\Character{48}}%
\TextAt{3}{198.168}{454.561}{\Character{106}}%
\TextAt{21}{200.52}{454.561}{K}%
\TextAt{8}{208.723}{456.421}{.}%
\TextAt{8}{217.171}{456.421}{Consequently,}%
\TextAt{20}{148.543}{441.396}{\Character{18}}%
\TextAt{13}{158.539}{432.627}{\Character{25}}%
\TextAt{13}{165.593}{432.627}{\Character{59}}%
\TextAt{13}{170.829}{432.627}{u}%
\Vector{q 0 G 0 g 0.478 w 86.539 -346.472 m 105.515 -346.472 l S Q}%
\TextAt{14}{165.035}{416.338}{p}%
\TextAt{20}{178.71}{441.396}{\Character{19}}%
\TextAt{9}{190.831}{424.539}{\Character{61}}%
\TextAt{13}{204.452}{432.627}{\Character{39}y}%
\Vector{q 0 G 0 g 0.478 w 132.452 -346.472 m 146.257 -346.472 l S Q}%
\TextAt{13}{208.287}{416.338}{y}%
\NativeMathAt{222.774}{424.539}{11.955}{\equiv}%
\TextAt{13}{235.393}{424.539}{y}%
\TextAt{21}{241.53}{429.475}{q}%
\TextAt{3}{245.597}{429.475}{\Character{0}}%
\TextAt{17}{252.183}{429.475}{1}%
\TextAt{9}{260.236}{424.539}{\Character{61}}%
\TextAt{13}{272.662}{424.539}{u}%
\TextAt{17}{279.324}{429.475}{\Character{40}}%
\TextAt{21}{282.617}{429.475}{q}%
\TextAt{3}{286.684}{429.475}{\Character{0}}%
\TextAt{17}{293.271}{429.475}{1\Character{41}\Character{47}}%
\TextAt{21}{305.032}{429.475}{n}%
\NativeMathAt{313.989}{424.539}{11.955}{\equiv}%
\TextAt{13}{326.609}{424.539}{\Character{33}}%
\TextAt{9}{334.327}{424.539}{\Character{40}}%
\TextAt{13}{338.879}{424.539}{u}%
\TextAt{9}{345.542}{424.539}{\Character{41}}%
\TextAt{17}{350.094}{429.475}{\Character{40}}%
\TextAt{21}{353.387}{429.475}{q}%
\TextAt{3}{357.454}{429.475}{\Character{0}}%
\TextAt{17}{364.041}{429.475}{1\Character{41}\Character{47}}%
\TextAt{21}{375.802}{429.475}{n}%
\TextAt{9}{384.759}{424.539}{\Character{61}}%
\TextAt{16}{397.185}{424.539}{\Character{104}}%
\TextAt{13}{401.834}{424.539}{\Character{25}}%
\TextAt{13}{408.888}{424.539}{\Character{59}}%
\TextAt{13}{414.124}{424.539}{u}%
\TextAt{16}{420.81}{424.539}{\Character{105}}%
\TextAt{8}{428.78}{424.539}{mod}%
\TextAt{14}{454.218}{424.539}{p}%
\TextAt{13}{460.201}{424.539}{\Character{59}}%
\TextAt{8}{79.226}{392.325}{hence}%
\TextAt{20}{111.732}{402.009}{\Character{0}}%
\TextAt{21}{118.407}{397.581}{\Character{25}}%
\TextAt{21}{123.548}{397.581}{\Character{59}u}%
\Vector{q 0 G 0 g 0.478 w 46.407 -378.686 m 58.811 -378.686 l S Q}%
\TextAt{27}{122.615}{388.203}{p}%
\TextAt{20}{132.007}{402.009}{\Character{1}}%
\TextAt{9}{140.807}{392.325}{\Character{61}}%
\TextAt{16}{153.233}{392.325}{\Character{104}}%
\TextAt{13}{157.882}{392.325}{\Character{25}}%
\TextAt{13}{164.936}{392.325}{\Character{59}}%
\TextAt{13}{170.172}{392.325}{u}%
\TextAt{16}{176.858}{392.325}{\Character{105}}%
\TextAt{8}{181.507}{392.325}{,}%
\TextAt{8}{188.656}{392.325}{because}%
\TextAt{13}{231.958}{392.325}{\Character{22}}%
\TextAt{21}{239.001}{390.532}{q}%
\TextAt{3}{243.068}{390.532}{\Character{0}}%
\TextAt{17}{249.654}{390.532}{1}%
\TextAt{8}{258.284}{392.325}{is}%
\TextAt{8}{270.048}{392.325}{mapped}%
\TextAt{8}{314.593}{392.325}{isomorphically}%
\TextAt{8}{392.385}{392.325}{onto}%
\TextAt{13}{418.734}{392.325}{\Character{20}}%
\TextAt{3}{425.473}{396.664}{\Character{3}}%
\TextAt{8}{434.102}{392.325}{by}%
\TextAt{13}{450.361}{392.325}{U}%
\TextAt{21}{458.312}{390.532}{K}%
\NativeMathAt{469.835}{392.325}{11.955}{\rightarrow}%
\TextAt{13}{485.111}{392.325}{\Character{20}}%
\TextAt{3}{491.85}{396.664}{\Character{3}}%
\TextAt{8}{496.582}{392.325}{.}%
\Vector{q 0 G 0 g 0.399 w 452.334 -381.675 m 452.334 -373.605 l S Q}%
\Vector{q 0 G 0 g 0.399 w 452.534 -373.804 m 459.707 -373.804 l S Q}%
\Vector{q 0 G 0 g 0.399 w 452.534 -381.476 m 459.707 -381.476 l S Q}%
\Vector{q 0 G 0 g 0.399 w 459.906 -381.675 m 459.906 -373.605 l S Q}%
\TextAt{8}{97.159}{363.969}{The}%
\TextAt{8}{121.213}{363.969}{proposition}%
\TextAt{8}{182.746}{363.969}{shows}%
\TextAt{8}{216.364}{363.969}{in}%
\TextAt{8}{230.017}{363.969}{particular}%
\TextAt{8}{284.006}{363.969}{that}%
\TextAt{8}{309.375}{363.969}{the}%
\TextAt{8}{329.532}{363.969}{Hilbert}%
\TextAt{8}{369.844}{363.969}{symbol}%
\TextAt{20}{251.226}{350.357}{\Character{18}}%
\TextAt{13}{261.222}{341.588}{\Character{25}}%
\TextAt{13}{268.276}{341.588}{\Character{59}}%
\TextAt{13}{273.512}{341.588}{u}%
\Vector{q 0 G 0 g 0.478 w 189.222 -437.511 m 208.198 -437.511 l S Q}%
\TextAt{14}{267.718}{325.3}{p}%
\TextAt{20}{281.393}{350.357}{\Character{19}}%
\TextAt{9}{293.514}{333.501}{\Character{61}}%
\TextAt{13}{305.94}{333.501}{\Character{33}}%
\TextAt{9}{313.658}{333.501}{\Character{40}}%
\TextAt{13}{318.21}{333.501}{u}%
\TextAt{9}{324.873}{333.501}{\Character{41}}%
\TextAt{17}{329.425}{338.437}{\Character{40}}%
\TextAt{21}{332.718}{338.437}{q}%
\TextAt{3}{336.785}{338.437}{\Character{0}}%
\TextAt{17}{343.372}{338.437}{1\Character{41}\Character{47}}%
\TextAt{21}{355.133}{338.437}{n}%
\TextAt{8}{79.226}{302.482}{(in}%
\TextAt{8}{98.402}{302.482}{the}%
\TextAt{8}{119.526}{302.482}{case}%
\TextAt{9}{145.265}{302.482}{\Character{40}}%
\TextAt{13}{149.817}{302.482}{n\Character{59}}%
\TextAt{13}{162.023}{302.482}{p}%
\TextAt{9}{167.924}{302.482}{\Character{41}}%
\TextAt{9}{177.428}{302.482}{\Character{61}}%
\TextAt{9}{191.5}{302.482}{1}%
\TextAt{8}{197.375}{302.482}{)}%
\TextAt{8}{206.795}{302.482}{is}%
\TextAt{8}{219.526}{302.482}{independent}%
\TextAt{8}{286.822}{302.482}{of}%
\TextAt{8}{301.119}{302.482}{the}%
\TextAt{8}{322.244}{302.482}{choice}%
\TextAt{8}{358.001}{302.482}{of}%
\TextAt{8}{372.299}{302.482}{the}%
\TextAt{8}{393.423}{302.482}{prime}%
\TextAt{8}{427.555}{302.482}{element}%
\TextAt{13}{471.765}{302.482}{\Character{25}}%
\TextAt{8}{478.834}{302.482}{.}%
\TextAt{8}{490.186}{302.482}{We}%
\TextAt{8}{511.298}{302.482}{may}%
\TextAt{8}{79.226}{285.436}{therefore}%
\TextAt{8}{128.326}{285.436}{put}%
\TextAt{20}{227.432}{281.642}{\Character{18}}%
\TextAt{13}{237.428}{272.873}{u}%
\Vector{q 0 G 0 g 0.478 w 165.428 -506.226 m 172.091 -506.226 l S Q}%
\TextAt{14}{237.768}{256.585}{p}%
\TextAt{20}{245.286}{281.642}{\Character{19}}%
\TextAt{9}{257.407}{264.786}{\Character{58}\Character{61}}%
\TextAt{20}{273.084}{281.642}{\Character{18}}%
\TextAt{13}{283.08}{272.873}{\Character{25}}%
\TextAt{13}{290.134}{272.873}{\Character{59}}%
\TextAt{13}{295.37}{272.873}{u}%
\Vector{q 0 G 0 g 0.478 w 211.08 -506.226 m 230.056 -506.226 l S Q}%
\TextAt{14}{289.577}{256.585}{p}%
\TextAt{20}{303.252}{281.642}{\Character{19}}%
\TextAt{8}{326}{264.786}{for}%
\TextAt{13}{343.885}{264.786}{u}%
\NativeMathAt{353.868}{264.786}{11.955}{\in}%
\TextAt{13}{365.159}{264.786}{U}%
\TextAt{21}{373.109}{262.992}{K}%
\TextAt{13}{381.312}{264.786}{\Character{58}}%
\TextAt{20}{79.226}{246.445}{\Character{0}}%
\TextAt{21}{85.901}{241.468}{u}%
\Vector{q 0 G 0 g 0.478 w 13.901 -534.25 m 18.804 -534.25 l S Q}%
\TextAt{27}{86.358}{232.639}{p}%
\TextAt{20}{92}{246.445}{\Character{1}}%
\TextAt{8}{101.376}{236.761}{is}%
\TextAt{8}{113.14}{236.761}{the}%
\TextAt{8}{133.297}{236.761}{root}%
\TextAt{8}{158.343}{236.761}{of}%
\TextAt{8}{171.673}{236.761}{unity}%
\TextAt{8}{202.242}{236.761}{determined}%
\TextAt{8}{263.369}{236.761}{by}%
\TextAt{20}{242.599}{219.959}{\Character{18}}%
\TextAt{13}{252.595}{211.19}{u}%
\Vector{q 0 G 0 g 0.478 w 180.595 -567.909 m 187.257 -567.909 l S Q}%
\TextAt{14}{252.935}{194.902}{p}%
\TextAt{20}{260.453}{219.959}{\Character{19}}%
\NativeMathAt{272.574}{203.102}{11.955}{\equiv}%
\TextAt{13}{285.193}{203.102}{u}%
\TextAt{17}{291.856}{208.039}{\Character{40}}%
\TextAt{21}{295.149}{208.039}{q}%
\TextAt{3}{299.216}{208.039}{\Character{0}}%
\TextAt{17}{305.802}{208.039}{1\Character{41}\Character{47}}%
\TextAt{21}{317.564}{208.039}{n}%
\TextAt{8}{326.521}{203.102}{mod}%
\TextAt{14}{351.959}{203.102}{p}%
\TextAt{21}{357.942}{201.309}{K}%
\TextAt{13}{366.145}{203.102}{\Character{58}}%
\TextAt{8}{79.226}{172.754}{We}%
\TextAt{8}{98.907}{172.754}{call}%
\TextAt{8}{119.902}{172.754}{it}%
\TextAt{8}{131.142}{172.754}{the}%
\TextAt{12}{150.834}{172.754}{Legendre}%
\TextAt{12}{208.439}{172.754}{symbol}%
\TextAt{8}{250.007}{172.754}{,}%
\TextAt{8}{256.785}{172.754}{or}%
\TextAt{8}{270.631}{172.754}{the}%
\TextAt{13}{290.323}{172.754}{n}%
\TextAt{12}{297.311}{172.754}{-th}%
\TextAt{12}{318.452}{172.754}{power}%
\TextAt{12}{357.585}{172.754}{residue}%
\TextAt{12}{403.294}{172.754}{symbol}%
\TextAt{8}{444.862}{172.754}{.}%
\TextAt{8}{453.156}{172.754}{Both}%
\TextAt{8}{481.79}{172.754}{names}%
\TextAt{8}{517.156}{172.754}{are}%
\TextAt{8}{79.226}{155.708}{justified}%
\TextAt{8}{123.831}{155.708}{by}%
\TextAt{8}{140.09}{155.708}{the}%
\TextAt{12}{79.226}{130.34}{Proposition}%
\TextAt{12}{151.937}{130.34}{(3.5).}%
\TextAt{8}{187.253}{130.34}{Let}%
\TextAt{9}{208.222}{130.34}{\Character{40}}%
\TextAt{13}{212.775}{130.34}{n\Character{59}}%
\TextAt{13}{224.981}{130.34}{p}%
\TextAt{9}{230.882}{130.34}{\Character{41}}%
\TextAt{9}{238.737}{130.34}{\Character{61}}%
\TextAt{9}{251.158}{130.34}{1}%
\TextAt{8}{260.931}{130.34}{and}%
\TextAt{13}{283.693}{130.34}{u}%
\NativeMathAt{293.677}{130.34}{11.955}{\in}%
\TextAt{13}{304.968}{130.34}{U}%
\TextAt{21}{312.918}{128.547}{K}%
\TextAt{8}{321.12}{130.34}{.}%
\TextAt{8}{329.569}{130.34}{Then}%
\TextAt{8}{360.126}{130.34}{one}%
\TextAt{8}{381.586}{130.34}{has}%
\TextAt{20}{184.465}{116.167}{\Character{18}}%
\TextAt{13}{194.461}{107.398}{u}%
\Vector{q 0 G 0 g 0.478 w 122.461 -671.701 m 129.123 -671.701 l S Q}%
\TextAt{14}{194.801}{91.109}{p}%
\TextAt{20}{202.319}{116.167}{\Character{19}}%
\TextAt{9}{214.44}{99.31}{\Character{61}}%
\TextAt{9}{226.85}{99.31}{1}%
\NativeArrowAt{247.994}{99.31}{11.955}{21.926}{\Longleftrightarrow}%
\TextAt{13}{285.188}{99.31}{u}%
\TextAt{8}{295.748}{99.31}{is}%
\TextAt{8}{307.512}{99.31}{an}%
\TextAt{13}{323.771}{99.31}{n}%
\TextAt{8}{330.759}{99.31}{-th}%
\TextAt{8}{349.612}{99.31}{power}%
\TextAt{8}{383.756}{99.31}{mod}%
\TextAt{14}{410.093}{99.31}{p}%
\TextAt{21}{416.076}{97.517}{K}%
\TextAt{13}{424.279}{99.31}{\Character{58}}%
\TextAt{8}{297.211}{43.033}{308}%
\Destination{78.23}{809.815}{nav.page.308}%
\Destination{79.36}{584.555}{nav.print.336}%
\Destination{79.36}{732.555}{nav.number.5.3.4}%
\Destination{79.36}{144.355}{nav.number.5.3.5}%
\Hotlink{351.417}{471.255}{13.995}{11.047}{nav.chapter.4}%
\Hotlink{95.948}{471.255}{25.078}{11.047}{nav.number.5.3.3}%
\Hotlink{370.42}{471.255}{25.078}{11.047}{nav.number.4.5.7}%
}

% Page 321 (body page 309)
\Page{%
\TextAt{2}{79.226}{796.943}{V.}%
\TextAt{2}{94.073}{796.943}{Local}%
\TextAt{2}{123.615}{796.943}{Class}%
\TextAt{2}{152.208}{796.943}{Field}%
\TextAt{2}{179.939}{796.943}{Theory}%
\Vector{q 0 G 0 g 0.299 w 7.226 18.729 m 460.769 18.729 l S Q}%
\TextAt{19}{79.226}{761.124}{Proof.}%
\TextAt{8}{116}{761.124}{Let}%
\TextAt{13}{138.321}{761.124}{\Character{16}}%
\TextAt{8}{149.585}{761.124}{be}%
\TextAt{8}{166.861}{761.124}{a}%
\TextAt{8}{177.968}{761.124}{primitive}%
\TextAt{9}{229.399}{761.124}{\Character{40}}%
\TextAt{13}{233.952}{761.124}{q}%
\TextAt{16}{243.149}{761.124}{\Character{0}}%
\TextAt{9}{256.025}{761.124}{1\Character{41}}%
\TextAt{8}{266.43}{761.124}{-th}%
\TextAt{8}{286.635}{761.124}{root}%
\TextAt{8}{313.033}{761.124}{of}%
\TextAt{8}{327.714}{761.124}{unity,}%
\TextAt{8}{362.244}{761.124}{and}%
\TextAt{8}{386.358}{761.124}{let}%
\TextAt{13}{404.614}{761.124}{m}%
\TextAt{9}{420.477}{761.124}{\Character{61}}%
\TextAt{21}{436.401}{766.379}{q}%
\TextAt{3}{440.468}{766.38}{\Character{0}}%
\TextAt{17}{447.054}{766.379}{1}%
\Vector{q 0 G 0 g 0.478 w 364.401 -9.888 m 379.289 -9.888 l S Q}%
\TextAt{21}{441.276}{757.001}{n}%
\TextAt{8}{452.484}{761.124}{.}%
\TextAt{8}{464.987}{761.124}{Then}%
\TextAt{13}{496.895}{761.124}{\Character{16}}%
\TextAt{21}{502.911}{765.462}{n}%
\TextAt{8}{513.796}{761.124}{is}%
\TextAt{8}{526.911}{761.124}{a}%
\TextAt{8}{79.226}{744.078}{primitive}%
\TextAt{13}{129.306}{744.078}{m}%
\TextAt{8}{139.546}{744.078}{-th}%
\TextAt{8}{158.399}{744.078}{root}%
\TextAt{8}{183.445}{744.078}{of}%
\TextAt{8}{196.775}{744.078}{unity,}%
\TextAt{8}{229.616}{744.078}{and}%
\TextAt{20}{147.542}{728.209}{\Character{18}}%
\TextAt{13}{157.538}{719.44}{u}%
\Vector{q 0 G 0 g 0.478 w 85.538 -59.659 m 92.2 -59.659 l S Q}%
\TextAt{14}{157.878}{703.152}{p}%
\TextAt{20}{165.396}{728.209}{\Character{19}}%
\TextAt{9}{177.517}{711.352}{\Character{61}}%
\TextAt{13}{189.942}{711.352}{\Character{33}}%
\TextAt{9}{197.661}{711.352}{\Character{40}}%
\TextAt{13}{202.213}{711.352}{u}%
\TextAt{9}{208.875}{711.352}{\Character{41}}%
\TextAt{21}{213.428}{716.289}{m}%
\TextAt{9}{224.737}{711.352}{\Character{61}}%
\TextAt{9}{237.147}{711.352}{1}%
\NativeArrowAt{258.292}{711.352}{11.955}{21.926}{\Longleftrightarrow}%
\TextAt{13}{295.486}{711.352}{\Character{33}}%
\TextAt{9}{303.204}{711.352}{\Character{40}}%
\TextAt{13}{307.756}{711.352}{u}%
\TextAt{9}{314.419}{711.352}{\Character{41}}%
\NativeMathAt{322.292}{711.352}{11.955}{\in}%
\TextAt{13}{333.583}{711.352}{\Character{22}}%
\TextAt{21}{340.626}{709.559}{m}%
\NativeArrowAt{363.89}{711.352}{11.955}{21.926}{\Longleftrightarrow}%
\TextAt{13}{401.084}{711.352}{\Character{33}}%
\TextAt{9}{408.802}{711.352}{\Character{40}}%
\TextAt{13}{413.355}{711.352}{u}%
\TextAt{9}{420.017}{711.352}{\Character{41}}%
\TextAt{9}{427.872}{711.352}{\Character{61}}%
\TextAt{9}{440.293}{711.352}{\Character{40}}%
\TextAt{13}{444.868}{711.352}{\Character{16}}%
\TextAt{21}{450.884}{716.289}{n}%
\TextAt{9}{456.52}{711.352}{\Character{41}}%
\TextAt{21}{461.072}{716.289}{i}%
\NativeArrowAt{258.292}{684.078}{11.955}{21.926}{\Longleftrightarrow}%
\TextAt{13}{295.486}{684.078}{u}%
\NativeMathAt{305.469}{684.078}{11.955}{\equiv}%
\TextAt{13}{318.088}{684.078}{\Character{33}}%
\TextAt{9}{325.806}{684.078}{\Character{40}}%
\TextAt{13}{330.359}{684.078}{u}%
\TextAt{9}{337.021}{684.078}{\Character{41}}%
\TextAt{9}{344.876}{684.078}{\Character{61}}%
\TextAt{9}{357.297}{684.078}{\Character{40}}%
\TextAt{13}{361.872}{684.078}{\Character{16}}%
\TextAt{21}{367.888}{689.014}{i}%
\TextAt{9}{371.269}{684.078}{\Character{41}}%
\TextAt{21}{375.821}{689.014}{n}%
\TextAt{8}{384.778}{684.078}{mod}%
\TextAt{14}{410.216}{684.078}{p}%
\TextAt{21}{416.199}{682.284}{K}%
\TextAt{13}{424.402}{684.078}{\Character{58}}%
\Vector{q 0 G 0 g 0.399 w 452.334 -118.924 m 452.334 -110.854 l S Q}%
\Vector{q 0 G 0 g 0.399 w 452.534 -111.053 m 459.707 -111.053 l S Q}%
\Vector{q 0 G 0 g 0.399 w 452.534 -118.724 m 459.707 -118.724 l S Q}%
\Vector{q 0 G 0 g 0.399 w 459.906 -118.924 m 459.906 -110.854 l S Q}%
\TextAt{8}{97.159}{626.075}{It}%
\TextAt{8}{110.38}{626.075}{is}%
\TextAt{8}{122.692}{626.075}{an}%
\TextAt{8}{139.499}{626.075}{important,}%
\TextAt{8}{198.729}{626.075}{but}%
\TextAt{8}{220.737}{626.075}{in}%
\TextAt{8}{234.938}{626.075}{general}%
\TextAt{8}{275.811}{626.075}{difficult,}%
\TextAt{8}{322.667}{626.075}{problem}%
\TextAt{8}{368.741}{626.075}{to}%
\TextAt{8}{383.599}{626.075}{find}%
\TextAt{8}{407.556}{626.075}{explicit}%
\TextAt{8}{449.397}{626.075}{formulae}%
\TextAt{8}{498.077}{626.075}{for}%
\TextAt{8}{516.51}{626.075}{the}%
\TextAt{8}{79.226}{609.029}{Hilbert}%
\TextAt{8}{120.601}{609.029}{symbol}%
\TextAt{20}{161.725}{618.713}{\Character{0}}%
\TextAt{21}{168.4}{614.285}{a\Character{59}b}%
\Vector{q 0 G 0 g 0.478 w 96.4 -161.982 m 106.873 -161.982 l S Q}%
\TextAt{27}{171.643}{604.907}{p}%
\TextAt{20}{180.069}{618.713}{\Character{1}}%
\TextAt{8}{190.508}{609.029}{also}%
\TextAt{8}{215.05}{609.029}{in}%
\TextAt{8}{229.766}{609.029}{the}%
\TextAt{8}{250.984}{609.029}{case}%
\TextAt{13}{276.818}{609.029}{p}%
\TextAt{16}{287.824}{609.029}{\Character{106}}%
\TextAt{13}{296.276}{609.029}{n}%
\TextAt{8}{303.264}{609.029}{.}%
\TextAt{8}{314.899}{609.029}{Let}%
\TextAt{8}{336.931}{609.029}{us}%
\TextAt{8}{353.009}{609.029}{look}%
\TextAt{8}{379.44}{609.029}{at}%
\TextAt{8}{394.813}{609.029}{the}%
\TextAt{8}{416.032}{609.029}{case}%
\TextAt{8}{441.865}{609.029}{where}%
\TextAt{13}{476.737}{609.029}{n}%
\TextAt{9}{488.856}{609.029}{\Character{61}}%
\TextAt{9}{503.083}{609.029}{2}%
\TextAt{8}{513.904}{609.029}{and}%
\TextAt{13}{79.226}{591.983}{K}%
\TextAt{9}{93.365}{591.983}{\Character{61}}%
\TextAt{18}{105.79}{591.983}{Q}%
\TextAt{21}{115.089}{590.19}{p}%
\TextAt{8}{119.85}{591.983}{.}%
\TextAt{8}{128.298}{591.983}{If}%
\TextAt{13}{139.99}{591.983}{a}%
\NativeMathAt{149.456}{591.983}{11.955}{\in}%
\TextAt{18}{160.747}{591.983}{Z}%
\TextAt{17}{168.717}{590.19}{2}%
\TextAt{8}{173.449}{591.983}{,}%
\TextAt{8}{180.598}{591.983}{then}%
\TextAt{9}{207.258}{591.983}{\Character{40}}%
\TextAt{16}{211.811}{591.983}{\Character{0}}%
\TextAt{9}{221.109}{591.983}{1\Character{41}}%
\TextAt{21}{231.515}{596.322}{a}%
\TextAt{8}{240.408}{591.983}{means}%
\TextAt{9}{267.467}{562.982}{\Character{40}}%
\TextAt{16}{272.02}{562.982}{\Character{0}}%
\TextAt{9}{281.318}{562.982}{1\Character{41}}%
\TextAt{21}{291.723}{567.918}{a}%
\TextAt{9}{300.04}{562.982}{\Character{61}}%
\TextAt{9}{312.45}{562.982}{\Character{40}}%
\TextAt{16}{317.018}{562.982}{\Character{0}}%
\TextAt{9}{326.317}{562.982}{1\Character{41}}%
\TextAt{21}{336.722}{567.918}{r}%
\TextAt{13}{341.277}{562.982}{\Character{59}}%
\TextAt{8}{79.226}{533.981}{where}%
\TextAt{13}{113.035}{533.981}{r}%
\TextAt{8}{122.533}{533.981}{is}%
\TextAt{8}{134.297}{533.981}{a}%
\TextAt{8}{144.053}{533.981}{rational}%
\TextAt{8}{187.641}{533.981}{integer}%
\NativeMathAt{226.34}{533.981}{11.955}{\equiv}%
\TextAt{13}{238.959}{533.981}{a}%
\TextAt{8}{248.425}{533.981}{mod}%
\TextAt{9}{273.863}{533.981}{2}%
\TextAt{8}{279.716}{533.981}{.}%
\TextAt{12}{79.226}{507.969}{Theorem}%
\TextAt{12}{136.24}{507.969}{(3.6).}%
\TextAt{8}{171.556}{507.969}{Let}%
\TextAt{13}{192.525}{507.969}{n}%
\TextAt{9}{202.834}{507.969}{\Character{61}}%
\TextAt{9}{215.244}{507.969}{2}%
\TextAt{8}{221.112}{507.969}{.}%
\TextAt{8}{229.561}{507.969}{For}%
\TextAt{13}{250.53}{507.969}{a\Character{59}}%
\TextAt{13}{261.899}{507.969}{b}%
\NativeMathAt{270.217}{507.969}{11.955}{\in}%
\TextAt{18}{281.508}{507.969}{Q}%
\TextAt{3}{290.806}{512.307}{\Character{3}}%
\TextAt{21}{290.806}{505.013}{p}%
\TextAt{8}{295.567}{507.969}{,}%
\TextAt{8}{302.717}{507.969}{we}%
\TextAt{8}{319.944}{507.969}{write}%
\TextAt{13}{206.104}{478.968}{a}%
\TextAt{9}{215.569}{478.968}{\Character{61}}%
\TextAt{13}{227.995}{478.968}{p}%
\TextAt{21}{233.87}{483.904}{\Character{11}}%
\TextAt{13}{239.807}{478.968}{a}%
\TextAt{3}{245.952}{483.904}{\Character{48}}%
\TextAt{13}{248.747}{478.968}{\Character{59}}%
\TextAt{13}{277.882}{478.968}{b}%
\TextAt{9}{286.2}{478.968}{\Character{61}}%
\TextAt{13}{298.625}{478.968}{p}%
\TextAt{21}{304.5}{483.904}{\Character{12}}%
\TextAt{13}{310.16}{478.968}{b}%
\TextAt{3}{315.137}{483.904}{\Character{48}}%
\TextAt{13}{317.932}{478.968}{\Character{59}}%
\TextAt{13}{347.067}{478.968}{a}%
\TextAt{3}{353.231}{483.904}{\Character{48}}%
\TextAt{13}{356.026}{478.968}{\Character{59}}%
\TextAt{13}{361.25}{478.968}{b}%
\TextAt{3}{366.248}{483.904}{\Character{48}}%
\NativeMathAt{372.363}{478.968}{11.955}{\in}%
\TextAt{13}{383.654}{478.968}{U}%
\TextAt{25}{391.605}{477.174}{Q}%
\TextAt{22}{397.804}{476.178}{p}%
\TextAt{13}{402.64}{478.968}{\Character{58}}%
\TextAt{8}{79.226}{449.966}{If}%
\TextAt{13}{90.918}{449.966}{p}%
\NativeMathAt{100.114}{449.966}{11.955}{\neq}%
\TextAt{9}{112.524}{449.966}{2}%
\TextAt{8}{118.393}{449.966}{,}%
\TextAt{8}{125.542}{449.966}{then}%
\TextAt{20}{215.688}{443.787}{\Character{18}}%
\TextAt{13}{225.683}{435.018}{a\Character{59}}%
\TextAt{13}{237.052}{435.018}{b}%
\Vector{q 0 G 0 g 0.478 w 153.683 -344.081 m 170.05 -344.081 l S Q}%
\TextAt{13}{230.929}{418.73}{p}%
\TextAt{20}{243.245}{443.787}{\Character{19}}%
\TextAt{9}{255.366}{426.931}{\Character{61}}%
\TextAt{9}{267.776}{426.931}{\Character{40}}%
\TextAt{16}{272.344}{426.93}{\Character{0}}%
\TextAt{9}{281.643}{426.931}{1\Character{41}}%
\TextAt{22}{293.243}{435.56}{p}%
\TextAt{23}{297.084}{435.56}{\Character{0}}%
\TextAt{24}{302.84}{435.56}{1}%
\Vector{q 0 G 0 g 0.359 w 221.243 -340.141 m 234.493 -340.141 l S Q}%
\TextAt{24}{298.042}{428.741}{2}%
\TextAt{21}{307.688}{431.867}{\Character{11}\Character{12}}%
\TextAt{20}{320.779}{443.787}{\Character{18}}%
\TextAt{13}{330.775}{435.018}{a}%
\TextAt{3}{336.92}{439.357}{\Character{48}}%
\Vector{q 0 G 0 g 0.478 w 258.775 -344.081 m 267.715 -344.081 l S Q}%
\TextAt{13}{332.308}{418.73}{p}%
\TextAt{20}{340.911}{443.787}{\Character{19}}%
\TextAt{21}{349.711}{441.111}{\Character{12}}%
\TextAt{20}{357.363}{443.787}{\Character{18}}%
\TextAt{13}{367.359}{435.018}{b}%
\TextAt{3}{372.336}{439.357}{\Character{48}}%
\Vector{q 0 G 0 g 0.478 w 295.359 -344.081 m 303.131 -344.081 l S Q}%
\TextAt{13}{368.307}{418.73}{p}%
\TextAt{20}{376.327}{443.787}{\Character{19}}%
\TextAt{21}{385.127}{441.111}{\Character{11}}%
\TextAt{13}{393.056}{426.931}{\Character{58}}%
\TextAt{8}{79.226}{397.785}{In}%
\TextAt{8}{94.514}{397.785}{particular,}%
\TextAt{8}{152.588}{397.785}{one}%
\TextAt{8}{174.714}{397.785}{has}%
\TextAt{20}{196.254}{407.469}{\Character{0}}%
\TextAt{21}{202.929}{403.041}{p\Character{59}p}%
\Vector{q 0 G 0 g 0.478 w 130.929 -373.226 m 141.807 -373.226 l S Q}%
\TextAt{21}{206.237}{393.662}{p}%
\TextAt{20}{215.003}{407.469}{\Character{1}}%
\TextAt{9}{224.939}{397.785}{\Character{61}}%
\TextAt{9}{238.484}{397.785}{\Character{40}}%
\TextAt{16}{243.052}{397.785}{\Character{0}}%
\TextAt{9}{252.351}{397.785}{1\Character{41}}%
\TextAt{17}{262.756}{402.123}{\Character{40}}%
\TextAt{21}{266.049}{402.123}{p}%
\TextAt{3}{270.312}{402.123}{\Character{0}}%
\TextAt{17}{276.898}{402.123}{1\Character{41}\Character{47}2}%
\TextAt{8}{297.956}{397.785}{and}%
\TextAt{20}{321.385}{407.469}{\Character{0}}%
\TextAt{21}{328.06}{403.041}{p\Character{59}u}%
\Vector{q 0 G 0 g 0.478 w 256.06 -373.226 m 267.578 -373.226 l S Q}%
\TextAt{21}{331.688}{393.662}{p}%
\TextAt{20}{340.774}{407.469}{\Character{1}}%
\TextAt{9}{350.71}{397.785}{\Character{61}}%
\TextAt{20}{364.271}{407.469}{\Character{0}}%
\TextAt{21}{370.946}{402.492}{u}%
\Vector{q 0 G 0 g 0.478 w 298.946 -373.226 m 303.849 -373.226 l S Q}%
\TextAt{21}{371.266}{393.662}{p}%
\TextAt{20}{377.044}{407.469}{\Character{1}}%
\TextAt{8}{382.524}{397.785}{,}%
\TextAt{8}{390.506}{397.785}{if}%
\TextAt{13}{401.896}{397.785}{u}%
\TextAt{8}{413.122}{397.785}{is}%
\TextAt{8}{425.553}{397.785}{a}%
\TextAt{8}{435.975}{397.785}{unit.}%
\TextAt{8}{467.236}{397.785}{If}%
\TextAt{13}{479.595}{397.785}{p}%
\TextAt{9}{489.926}{397.785}{\Character{61}}%
\TextAt{9}{503.471}{397.785}{2}%
\TextAt{8}{513.904}{397.785}{and}%
\TextAt{13}{79.226}{380.739}{a\Character{59}}%
\TextAt{13}{90.595}{380.739}{b}%
\NativeMathAt{98.913}{380.739}{11.955}{\in}%
\TextAt{13}{110.204}{380.739}{U}%
\TextAt{25}{118.154}{378.946}{Q}%
\TextAt{24}{124.353}{377.839}{2}%
\TextAt{8}{129.003}{380.739}{,}%
\TextAt{8}{136.152}{380.739}{then}%
\TextAt{20}{248.242}{376.143}{\Character{18}}%
\TextAt{9}{258.238}{367.374}{2}%
\TextAt{13}{264.091}{367.374}{\Character{59}}%
\TextAt{13}{269.315}{367.374}{a}%
\Vector{q 0 G 0 g 0.478 w 186.238 -411.725 m 203.48 -411.725 l S Q}%
\TextAt{9}{263.933}{351.085}{2}%
\TextAt{20}{276.676}{376.143}{\Character{19}}%
\TextAt{9}{288.797}{359.286}{\Character{61}}%
\TextAt{9}{301.207}{359.286}{\Character{40}}%
\TextAt{16}{305.775}{359.286}{\Character{0}}%
\TextAt{9}{315.073}{359.286}{1\Character{41}}%
\TextAt{17}{325.479}{364.222}{\Character{40}}%
\TextAt{21}{328.772}{364.222}{a}%
\TextAt{24}{333.27}{367.035}{2}%
\TextAt{3}{337.421}{364.222}{\Character{0}}%
\TextAt{17}{344.008}{364.222}{1\Character{41}\Character{47}8}%
\TextAt{13}{360.502}{359.286}{\Character{59}}%
\TextAt{20}{224.437}{339.978}{\Character{18}}%
\TextAt{13}{234.432}{331.209}{a\Character{59}}%
\TextAt{13}{245.801}{331.209}{b}%
\Vector{q 0 G 0 g 0.478 w 162.432 -447.89 m 178.799 -447.89 l S Q}%
\TextAt{9}{239.689}{314.921}{2}%
\TextAt{20}{251.994}{339.978}{\Character{19}}%
\TextAt{9}{264.115}{323.121}{\Character{61}}%
\TextAt{20}{276.541}{339.978}{\Character{18}}%
\TextAt{13}{286.537}{331.209}{b\Character{59}}%
\TextAt{13}{296.735}{331.209}{a}%
\Vector{q 0 G 0 g 0.478 w 214.537 -447.89 m 230.903 -447.89 l S Q}%
\TextAt{9}{291.793}{314.921}{2}%
\TextAt{20}{304.098}{339.978}{\Character{19}}%
\TextAt{9}{316.22}{323.121}{\Character{61}}%
\TextAt{9}{328.63}{323.121}{\Character{40}}%
\TextAt{16}{333.197}{323.121}{\Character{0}}%
\TextAt{9}{342.496}{323.121}{1\Character{41}}%
\TextAt{22}{354.097}{331.418}{a}%
\TextAt{23}{358.131}{331.418}{\Character{0}}%
\TextAt{24}{363.887}{331.418}{1}%
\Vector{q 0 G 0 g 0.359 w 282.097 -443.95 m 295.54 -443.95 l S Q}%
\TextAt{24}{358.992}{324.932}{2}%
\TextAt{22}{369.931}{331.418}{b}%
\TextAt{23}{373.205}{331.418}{\Character{0}}%
\TextAt{24}{378.961}{331.418}{1}%
\Vector{q 0 G 0 g 0.359 w 297.931 -443.95 m 310.614 -443.95 l S Q}%
\TextAt{24}{374.446}{324.932}{2}%
\TextAt{13}{384.307}{323.121}{\Character{58}}%
\TextAt{19}{79.226}{277.393}{Proof.}%
\TextAt{8}{116}{277.393}{The}%
\TextAt{8}{140.878}{277.393}{claim}%
\TextAt{8}{172.917}{277.393}{for}%
\TextAt{8}{191.626}{277.393}{the}%
\TextAt{8}{212.606}{277.393}{case}%
\TextAt{13}{238.201}{277.393}{p}%
\NativeMathAt{248.801}{277.393}{11.955}{\neq}%
\TextAt{9}{262.621}{277.393}{2}%
\TextAt{8}{273.204}{277.393}{is}%
\TextAt{8}{285.792}{277.393}{an}%
\TextAt{8}{302.875}{277.393}{immediate}%
\TextAt{8}{360.928}{277.393}{consequence}%
\TextAt{8}{427.816}{277.393}{of}%
\TextAt{8}{441.97}{277.393}{(3.4),}%
\TextAt{8}{474.227}{277.393}{and}%
\TextAt{8}{497.814}{277.393}{will}%
\TextAt{8}{520.743}{277.393}{be}%
\TextAt{8}{79.226}{260.347}{left}%
\TextAt{8}{100.531}{260.347}{to}%
\TextAt{8}{115.666}{260.347}{the}%
\TextAt{8}{136.648}{260.347}{reader.}%
\TextAt{8}{179.445}{260.347}{So}%
\TextAt{8}{196.529}{260.347}{let}%
\TextAt{13}{214.259}{260.347}{p}%
\TextAt{9}{224.861}{260.347}{\Character{61}}%
\TextAt{9}{238.681}{260.347}{2}%
\TextAt{8}{244.546}{260.347}{.}%
\TextAt{8}{255.47}{260.347}{We}%
\TextAt{8}{276.44}{260.347}{put}%
\TextAt{13}{298.725}{260.347}{\Character{17}}%
\TextAt{9}{304.929}{260.347}{\Character{40}}%
\TextAt{13}{309.482}{260.347}{a}%
\TextAt{9}{315.626}{260.347}{\Character{41}}%
\TextAt{9}{324.891}{260.347}{\Character{61}}%
\TextAt{21}{339.933}{265.054}{a}%
\TextAt{24}{344.431}{267.867}{2}%
\TextAt{3}{348.582}{265.054}{\Character{0}}%
\TextAt{17}{355.169}{265.054}{1}%
\Vector{q 0 G 0 g 0.478 w 267.933 -510.665 m 287.403 -510.665 l S Q}%
\TextAt{17}{347.551}{256.224}{8}%
\TextAt{8}{365.321}{260.347}{and}%
\TextAt{13}{388.909}{260.347}{\Character{34}}%
\TextAt{9}{394.387}{260.347}{\Character{40}}%
\TextAt{13}{398.94}{260.347}{a}%
\TextAt{9}{405.085}{260.347}{\Character{41}}%
\TextAt{9}{414.35}{260.347}{\Character{61}}%
\TextAt{21}{429.391}{265.054}{a}%
\TextAt{3}{433.89}{265.054}{\Character{0}}%
\TextAt{17}{440.476}{265.054}{1}%
\Vector{q 0 G 0 g 0.478 w 357.391 -510.665 m 372.71 -510.665 l S Q}%
\TextAt{17}{434.934}{256.224}{2}%
\TextAt{8}{445.906}{260.347}{.}%
\TextAt{8}{456.83}{260.347}{An}%
\TextAt{8}{476.831}{260.347}{elementary}%
\TextAt{8}{79.226}{243.301}{computation}%
\TextAt{8}{147.526}{243.301}{shows}%
\TextAt{8}{181.144}{243.301}{that}%
\TextAt{13}{129.668}{214.3}{\Character{17}}%
\TextAt{9}{135.872}{214.3}{\Character{40}}%
\TextAt{13}{140.424}{214.3}{a}%
\TextAt{17}{146.569}{212.506}{1}%
\TextAt{13}{151.302}{214.3}{a}%
\TextAt{17}{157.447}{212.506}{2}%
\TextAt{9}{162.179}{214.3}{\Character{41}}%
\NativeMathAt{170.052}{214.3}{11.955}{\equiv}%
\TextAt{13}{182.671}{214.3}{\Character{17}}%
\TextAt{9}{188.876}{214.3}{\Character{40}}%
\TextAt{13}{193.428}{214.3}{a}%
\TextAt{17}{199.573}{212.506}{1}%
\TextAt{9}{204.305}{214.3}{\Character{41}}%
\TextAt{9}{211.502}{214.3}{\Character{43}}%
\TextAt{13}{223.276}{214.3}{\Character{17}}%
\TextAt{9}{229.48}{214.3}{\Character{40}}%
\TextAt{13}{234.033}{214.3}{a}%
\TextAt{17}{240.177}{212.506}{2}%
\TextAt{9}{244.91}{214.3}{\Character{41}}%
\TextAt{8}{252.783}{214.3}{mod}%
\TextAt{9}{278.221}{214.3}{2}%
\TextAt{8}{296.029}{214.3}{and}%
\TextAt{13}{326.849}{214.3}{\Character{34}}%
\TextAt{9}{332.327}{214.3}{\Character{40}}%
\TextAt{13}{336.88}{214.3}{a}%
\TextAt{17}{343.025}{212.506}{1}%
\TextAt{13}{347.757}{214.3}{a}%
\TextAt{17}{353.902}{212.506}{2}%
\TextAt{9}{358.634}{214.3}{\Character{41}}%
\NativeMathAt{366.507}{214.3}{11.955}{\equiv}%
\TextAt{13}{379.127}{214.3}{\Character{34}}%
\TextAt{9}{384.605}{214.3}{\Character{40}}%
\TextAt{13}{389.157}{214.3}{a}%
\TextAt{17}{395.302}{212.506}{1}%
\TextAt{9}{400.034}{214.3}{\Character{41}}%
\TextAt{9}{407.231}{214.3}{\Character{43}}%
\TextAt{13}{419.005}{214.3}{\Character{34}}%
\TextAt{9}{424.483}{214.3}{\Character{40}}%
\TextAt{13}{429.035}{214.3}{a}%
\TextAt{17}{435.18}{212.506}{2}%
\TextAt{9}{439.912}{214.3}{\Character{41}}%
\TextAt{8}{447.785}{214.3}{mod}%
\TextAt{9}{473.223}{214.3}{2}%
\TextAt{13}{479.076}{214.3}{\Character{58}}%
\TextAt{8}{79.226}{185.298}{Thus}%
\TextAt{8}{108.992}{185.298}{both}%
\TextAt{8}{136.75}{185.298}{sides}%
\TextAt{8}{164.95}{185.298}{of}%
\TextAt{8}{178.397}{185.298}{the}%
\TextAt{8}{198.671}{185.298}{equations}%
\TextAt{8}{251.212}{185.298}{we}%
\TextAt{8}{268.556}{185.298}{have}%
\TextAt{8}{295.669}{185.298}{to}%
\TextAt{8}{310.096}{185.298}{prove}%
\TextAt{8}{341.763}{185.298}{are}%
\TextAt{8}{361.391}{185.298}{multiplicative}%
\TextAt{8}{435.332}{185.298}{and}%
\TextAt{8}{458.212}{185.298}{it}%
\TextAt{8}{470.033}{185.298}{is}%
\TextAt{8}{481.914}{185.298}{enough}%
\TextAt{8}{522.357}{185.298}{to}%
\TextAt{8}{79.226}{168.252}{check}%
\TextAt{8}{111.278}{168.252}{the}%
\TextAt{8}{131.949}{168.252}{claim}%
\TextAt{8}{163.679}{168.252}{for}%
\TextAt{8}{182.078}{168.252}{a}%
\TextAt{8}{192.348}{168.252}{set}%
\TextAt{8}{211.13}{168.252}{of}%
\TextAt{8}{224.974}{168.252}{generators}%
\TextAt{8}{282.145}{168.252}{of}%
\TextAt{13}{295.989}{168.252}{U}%
\TextAt{25}{303.939}{166.459}{Q}%
\TextAt{24}{310.138}{165.352}{2}%
\TextAt{9}{314.787}{168.252}{\Character{47}}%
\TextAt{13}{320.64}{168.252}{U}%
\TextAt{17}{329.839}{172.591}{2}%
\TextAt{25}{328.591}{165.188}{Q}%
\TextAt{24}{334.79}{164.081}{2}%
\TextAt{8}{339.439}{168.252}{.}%
\TextAt{16}{349.431}{168.252}{\Character{102}}%
\TextAt{9}{355.408}{168.252}{5}%
\TextAt{13}{361.261}{168.252}{\Character{59}}%
\TextAt{16}{366.505}{168.252}{\Character{0}}%
\TextAt{9}{375.804}{168.252}{1}%
\TextAt{16}{381.657}{168.252}{\Character{103}}%
\TextAt{8}{392.046}{168.252}{is}%
\TextAt{8}{404.325}{168.252}{such}%
\TextAt{8}{431.236}{168.252}{a}%
\TextAt{8}{441.506}{168.252}{set.}%
\TextAt{8}{465.868}{168.252}{We}%
\TextAt{8}{486.527}{168.252}{postpone}%
\TextAt{8}{79.226}{151.207}{this}%
\TextAt{8}{102.049}{151.207}{for}%
\TextAt{8}{119.933}{151.207}{the}%
\TextAt{8}{140.09}{151.207}{moment}%
\TextAt{8}{185.292}{151.207}{and}%
\TextAt{8}{208.055}{151.207}{define}%
\TextAt{9}{241.864}{151.207}{\Character{40}}%
\TextAt{13}{246.417}{151.206}{a\Character{59}}%
\TextAt{13}{257.786}{151.206}{b}%
\TextAt{9}{262.783}{151.207}{\Character{41}}%
\TextAt{9}{270.638}{151.207}{\Character{61}}%
\TextAt{20}{283.081}{160.89}{\Character{0}}%
\TextAt{21}{289.756}{156.462}{a\Character{59}b}%
\Vector{q 0 G 0 g 0.478 w 217.756 -619.805 m 228.229 -619.805 l S Q}%
\TextAt{17}{292.876}{147.084}{2}%
\TextAt{20}{301.425}{160.89}{\Character{1}}%
\TextAt{8}{306.904}{151.207}{.}%
\TextAt{8}{97.159}{131.172}{We}%
\TextAt{8}{117.751}{131.172}{have}%
\TextAt{9}{145.193}{131.172}{\Character{40}}%
\TextAt{16}{149.745}{131.172}{\Character{0}}%
\TextAt{9}{159.044}{131.172}{1}%
\TextAt{13}{164.897}{131.172}{\Character{59}}%
\TextAt{13}{170.121}{131.172}{x}%
\TextAt{9}{176.793}{131.172}{\Character{41}}%
\TextAt{9}{185.413}{131.172}{\Character{61}}%
\TextAt{9}{198.599}{131.172}{1}%
\TextAt{8}{208.814}{131.172}{if}%
\TextAt{8}{219.985}{131.172}{and}%
\TextAt{8}{243.195}{131.172}{only}%
\TextAt{8}{269.334}{131.172}{if}%
\TextAt{13}{280.505}{131.172}{x}%
\TextAt{8}{291.502}{131.172}{is}%
\TextAt{8}{303.714}{131.172}{a}%
\TextAt{8}{313.916}{131.172}{norm}%
\TextAt{8}{344.933}{131.172}{from}%
\TextAt{18}{373.021}{131.172}{Q}%
\TextAt{17}{382.319}{129.378}{2}%
\TextAt{9}{387.052}{131.172}{\Character{40}}%
\TextAt{16}{391.604}{140.563}{\Character{112}}%
\Vector{q 0 G 0 g 0.478 w 329.567 -633.198 m 344.718 -633.198 l S Q}%
\TextAt{16}{401.567}{131.172}{\Character{0}}%
\TextAt{9}{410.865}{131.172}{1\Character{41}}%
\TextAt{16}{421.27}{131.172}{\Character{106}}%
\TextAt{18}{424.591}{131.172}{Q}%
\TextAt{17}{433.89}{129.378}{2}%
\TextAt{8}{438.622}{131.172}{,}%
\TextAt{8}{446.33}{131.172}{i.e.,}%
\TextAt{13}{468.995}{131.172}{x}%
\TextAt{9}{479.73}{131.172}{\Character{61}}%
\TextAt{13}{492.918}{131.172}{y}%
\TextAt{17}{499.055}{135.51}{2}%
\TextAt{9}{506.749}{131.172}{\Character{43}}%
\TextAt{13}{518.815}{131.172}{z}%
\TextAt{17}{524.785}{135.51}{2}%
\TextAt{8}{529.518}{131.172}{,}%
\TextAt{13}{79.226}{114.126}{y}%
\TextAt{13}{85.347}{114.126}{\Character{59}}%
\TextAt{13}{90.583}{114.126}{z}%
\NativeMathAt{100.336}{114.126}{11.955}{\in}%
\TextAt{18}{112.065}{114.126}{Q}%
\TextAt{17}{121.364}{112.333}{2}%
\TextAt{8}{126.096}{114.126}{.}%
\TextAt{8}{135.315}{114.126}{Since}%
\TextAt{9}{166.13}{114.126}{5}%
\TextAt{9}{175.73}{114.126}{\Character{61}}%
\TextAt{9}{188.582}{114.126}{4}%
\TextAt{9}{197.261}{114.126}{\Character{43}}%
\TextAt{9}{209.193}{114.126}{1}%
\TextAt{8}{219.233}{114.126}{and}%
\TextAt{9}{242.253}{114.126}{2}%
\TextAt{9}{251.853}{114.126}{\Character{61}}%
\TextAt{9}{264.705}{114.126}{1}%
\TextAt{9}{273.384}{114.126}{\Character{43}}%
\TextAt{9}{285.316}{114.126}{1}%
\TextAt{8}{291.202}{114.126}{,}%
\TextAt{8}{298.673}{114.126}{we}%
\TextAt{8}{316.157}{114.126}{find}%
\TextAt{8}{339.822}{114.126}{that}%
\TextAt{9}{365.448}{114.126}{\Character{40}}%
\TextAt{16}{370}{114.126}{\Character{0}}%
\TextAt{9}{379.299}{114.126}{1}%
\TextAt{13}{385.152}{114.126}{\Character{59}}%
\TextAt{9}{390.396}{114.126}{2\Character{41}}%
\TextAt{9}{404.539}{114.126}{\Character{61}}%
\TextAt{9}{417.391}{114.126}{\Character{40}}%
\TextAt{16}{421.976}{114.126}{\Character{0}}%
\TextAt{9}{431.274}{114.126}{1}%
\TextAt{13}{437.127}{114.126}{\Character{59}}%
\TextAt{9}{442.371}{114.126}{5\Character{41}}%
\TextAt{9}{456.514}{114.126}{\Character{61}}%
\TextAt{9}{469.366}{114.126}{1}%
\TextAt{8}{475.252}{114.126}{.}%
\TextAt{8}{484.471}{114.126}{If}%
\TextAt{8}{496.42}{114.126}{we}%
\TextAt{8}{513.904}{114.126}{had}%
\TextAt{9}{79.226}{97.08}{\Character{40}}%
\TextAt{16}{83.778}{97.08}{\Character{0}}%
\TextAt{9}{93.077}{97.08}{1}%
\TextAt{13}{98.93}{97.08}{\Character{59}}%
\TextAt{16}{104.174}{97.08}{\Character{0}}%
\TextAt{9}{113.473}{97.08}{1\Character{41}}%
\TextAt{9}{127.556}{97.08}{\Character{61}}%
\TextAt{9}{140.348}{97.08}{1}%
\TextAt{8}{146.226}{97.08}{,}%
\TextAt{8}{153.65}{97.08}{then}%
\TextAt{8}{180.53}{97.08}{it}%
\TextAt{8}{192.454}{97.08}{would}%
\TextAt{8}{226.818}{97.08}{follow}%
\TextAt{8}{260.859}{97.08}{that}%
\TextAt{9}{286.448}{97.08}{\Character{40}}%
\TextAt{16}{291}{97.08}{\Character{0}}%
\TextAt{9}{300.298}{97.08}{1}%
\TextAt{13}{306.151}{97.08}{\Character{59}}%
\TextAt{13}{311.375}{97.08}{x}%
\TextAt{9}{318.048}{97.08}{\Character{41}}%
\TextAt{9}{326.285}{97.08}{\Character{61}}%
\TextAt{9}{339.077}{97.08}{1}%
\TextAt{8}{349.066}{97.08}{for}%
\TextAt{8}{367.17}{97.08}{all}%
\TextAt{13}{383.649}{97.08}{x}%
\TextAt{8}{390.301}{97.08}{,}%
\TextAt{8}{397.725}{97.08}{i.e.,}%
\TextAt{16}{420.105}{97.08}{\Character{0}}%
\TextAt{9}{429.404}{97.08}{1}%
\TextAt{8}{439.374}{97.08}{would}%
\TextAt{8}{473.738}{97.08}{be}%
\TextAt{8}{489.882}{97.08}{a}%
\TextAt{8}{499.857}{97.08}{square}%
\TextAt{8}{79.226}{80.034}{in}%
\TextAt{18}{92.879}{80.034}{Q}%
\TextAt{3}{102.177}{84.372}{\Character{3}}%
\TextAt{17}{102.177}{77.078}{2}%
\TextAt{8}{106.91}{80.034}{,}%
\TextAt{8}{114.059}{80.034}{which}%
\TextAt{8}{147.545}{80.034}{is}%
\TextAt{8}{159.309}{80.034}{not}%
\TextAt{8}{180.123}{80.034}{the}%
\TextAt{8}{200.28}{80.034}{case.}%
\TextAt{8}{229.602}{80.034}{Therefore}%
\TextAt{8}{282.599}{80.034}{we}%
\TextAt{8}{299.826}{80.034}{have}%
\TextAt{9}{326.821}{80.034}{\Character{40}}%
\TextAt{16}{331.373}{80.034}{\Character{0}}%
\TextAt{9}{340.672}{80.034}{1}%
\TextAt{13}{346.525}{80.034}{\Character{59}}%
\TextAt{16}{351.769}{80.034}{\Character{0}}%
\TextAt{9}{361.068}{80.034}{1\Character{41}}%
\TextAt{9}{374.769}{80.034}{\Character{61}}%
\TextAt{16}{387.219}{80.034}{\Character{0}}%
\TextAt{9}{396.518}{80.034}{1}%
\TextAt{8}{402.371}{80.034}{.}%
\TextAt{8}{297.211}{43.033}{309}%
\Destination{78.23}{809.815}{nav.page.309}%
\Destination{79.36}{776.155}{nav.print.337}%
\Destination{79.36}{521.955}{nav.number.5.3.6}%
\Hotlink{441.47}{275.155}{25.078}{11.047}{nav.number.5.3.4}%
\Hotlink{427.803}{274.401}{18.04}{11.047}{nav.number.5.3.4}%
}

% Page 322 (body page 310)
\Page{%
\TextAt{2}{79.226}{796.943}{V.}%
\TextAt{2}{94.073}{796.943}{Local}%
\TextAt{2}{123.615}{796.943}{Class}%
\TextAt{2}{152.208}{796.943}{Field}%
\TextAt{2}{179.939}{796.943}{Theory}%
\Vector{q 0 G 0 g 0.299 w 7.226 18.729 m 460.769 18.729 l S Q}%
\TextAt{8}{97.159}{761.124}{We}%
\TextAt{8}{118.997}{761.124}{have}%
\TextAt{9}{147.686}{761.124}{\Character{40}2}%
\TextAt{13}{158.091}{761.124}{\Character{59}}%
\TextAt{9}{163.336}{761.124}{2\Character{41}}%
\TextAt{9}{179.93}{761.124}{\Character{61}}%
\TextAt{9}{195.233}{761.124}{\Character{40}2}%
\TextAt{13}{205.666}{761.124}{\Character{59}}%
\TextAt{16}{210.91}{761.124}{\Character{0}}%
\TextAt{9}{220.209}{761.124}{1\Character{41}}%
\TextAt{9}{236.803}{761.124}{\Character{61}}%
\TextAt{9}{252.106}{761.124}{1}%
\TextAt{8}{263.578}{761.124}{and}%
\TextAt{9}{288.035}{761.124}{\Character{40}5}%
\TextAt{13}{298.44}{761.124}{\Character{59}}%
\TextAt{9}{303.684}{761.124}{5\Character{41}}%
\TextAt{9}{320.278}{761.124}{\Character{61}}%
\TextAt{9}{335.581}{761.124}{\Character{40}5}%
\TextAt{13}{346.015}{761.124}{\Character{59}}%
\TextAt{16}{351.259}{761.124}{\Character{0}}%
\TextAt{9}{360.557}{761.124}{1\Character{41}}%
\TextAt{9}{377.151}{761.124}{\Character{61}}%
\TextAt{9}{392.454}{761.124}{1}%
\TextAt{8}{398.336}{761.124}{.}%
\TextAt{8}{411.866}{761.124}{It}%
\TextAt{8}{426.232}{761.124}{remains}%
\TextAt{8}{471.563}{761.124}{therefore}%
\TextAt{8}{522.357}{761.124}{to}%
\TextAt{8}{79.226}{744.078}{determine}%
\TextAt{9}{134.151}{744.078}{\Character{40}2}%
\TextAt{13}{144.556}{744.078}{\Character{59}}%
\TextAt{9}{149.801}{744.078}{5\Character{41}}%
\TextAt{8}{160.206}{744.078}{.}%
\TextAt{9}{169.56}{744.078}{\Character{40}2}%
\TextAt{13}{179.965}{744.078}{\Character{59}}%
\TextAt{9}{185.209}{744.078}{5\Character{41}}%
\TextAt{9}{199.424}{744.078}{\Character{61}}%
\TextAt{9}{212.359}{744.078}{1}%
\TextAt{8}{222.442}{744.078}{would}%
\TextAt{8}{256.888}{744.078}{imply}%
\TextAt{9}{290.031}{744.078}{\Character{40}2}%
\TextAt{13}{300.436}{744.078}{\Character{59}}%
\TextAt{13}{305.66}{744.078}{x}%
\TextAt{9}{312.332}{744.078}{\Character{41}}%
\TextAt{9}{320.701}{744.078}{\Character{61}}%
\TextAt{9}{333.636}{744.078}{1}%
\TextAt{8}{343.712}{744.078}{for}%
\TextAt{8}{361.899}{744.078}{all}%
\TextAt{13}{378.46}{744.078}{x}%
\TextAt{8}{385.112}{744.078}{,}%
\TextAt{8}{392.638}{744.078}{i.e.,}%
\TextAt{9}{415.121}{744.078}{2}%
\TextAt{8}{425.173}{744.078}{would}%
\TextAt{8}{459.619}{744.078}{be}%
\TextAt{8}{475.845}{744.078}{a}%
\TextAt{8}{485.902}{744.078}{square}%
\TextAt{8}{523.014}{744.078}{in}%
\TextAt{18}{79.226}{727.032}{Q}%
\TextAt{3}{88.525}{731.37}{\Character{3}}%
\TextAt{17}{88.525}{724.076}{2}%
\TextAt{8}{93.257}{727.032}{,}%
\TextAt{8}{101.445}{727.032}{which}%
\TextAt{8}{135.763}{727.032}{is}%
\TextAt{8}{148.359}{727.032}{not}%
\TextAt{8}{170.004}{727.032}{the}%
\TextAt{8}{190.992}{727.032}{case.}%
\TextAt{8}{222.808}{727.032}{Hence}%
\TextAt{9}{258.417}{727.032}{\Character{40}2}%
\TextAt{13}{268.822}{727.032}{\Character{59}}%
\TextAt{9}{274.067}{727.032}{5\Character{41}}%
\TextAt{9}{289.19}{727.032}{\Character{61}}%
\TextAt{16}{303.052}{727.032}{\Character{0}}%
\TextAt{9}{312.35}{727.032}{1}%
\TextAt{8}{318.203}{727.032}{.}%
\TextAt{8}{329.146}{727.032}{By}%
\TextAt{8}{348.34}{727.032}{direct}%
\TextAt{8}{382.335}{727.032}{verification}%
\TextAt{8}{443.66}{727.032}{one}%
\TextAt{8}{465.951}{727.032}{sees}%
\TextAt{8}{490.31}{727.032}{that}%
\TextAt{8}{516.51}{727.032}{the}%
\TextAt{8}{79.226}{709.986}{values}%
\TextAt{8}{114.594}{709.986}{we}%
\TextAt{8}{132.327}{709.986}{just}%
\TextAt{8}{155.979}{709.986}{found}%
\TextAt{8}{189.326}{709.986}{coincide}%
\TextAt{8}{234.7}{709.986}{with}%
\TextAt{8}{261.866}{709.986}{those}%
\TextAt{8}{293.001}{709.986}{of}%
\TextAt{9}{306.838}{709.986}{\Character{40}}%
\TextAt{16}{311.39}{709.986}{\Character{0}}%
\TextAt{9}{320.688}{709.986}{1\Character{41}}%
\TextAt{21}{331.094}{714.324}{\Character{17}}%
\TextAt{17}{335.629}{714.324}{\Character{40}}%
\TextAt{21}{338.922}{714.324}{a}%
\TextAt{17}{343.42}{714.324}{\Character{41}}%
\TextAt{8}{347.212}{709.986}{,}%
\TextAt{8}{354.994}{709.986}{resp.}%
\TextAt{9}{385.834}{709.986}{\Character{40}}%
\TextAt{16}{390.387}{709.986}{\Character{0}}%
\TextAt{9}{399.685}{709.986}{1\Character{41}}%
\TextAt{21}{410.091}{714.324}{\Character{34}}%
\TextAt{17}{414.001}{714.324}{\Character{40}}%
\TextAt{21}{417.294}{714.324}{a}%
\TextAt{17}{421.792}{714.324}{\Character{41}}%
\TextAt{21}{425.085}{714.324}{\Character{34}}%
\TextAt{17}{428.996}{714.324}{\Character{40}}%
\TextAt{21}{432.289}{714.324}{b}%
\TextAt{17}{435.912}{714.324}{\Character{41}}%
\TextAt{8}{439.703}{709.986}{,}%
\TextAt{8}{447.485}{709.986}{in}%
\TextAt{8}{461.644}{709.986}{the}%
\TextAt{8}{482.307}{709.986}{respective}%
\TextAt{8}{79.226}{692.94}{cases.}%
\TextAt{8}{97.159}{672.905}{It}%
\TextAt{8}{110.426}{672.905}{remains}%
\TextAt{8}{154.657}{672.905}{to}%
\TextAt{8}{169.563}{672.905}{show}%
\TextAt{8}{199.161}{672.905}{that}%
\TextAt{13}{225.124}{672.905}{U}%
\TextAt{25}{233.075}{671.112}{Q}%
\TextAt{24}{239.274}{670.005}{2}%
\TextAt{9}{243.923}{672.905}{\Character{47}}%
\TextAt{13}{249.776}{672.905}{U}%
\TextAt{17}{258.975}{677.243}{2}%
\TextAt{25}{257.726}{669.84}{Q}%
\TextAt{24}{263.925}{668.733}{2}%
\TextAt{8}{273.067}{672.905}{is}%
\TextAt{8}{285.425}{672.905}{generated}%
\TextAt{8}{339.352}{672.905}{by}%
\TextAt{16}{356.206}{672.905}{\Character{102}}%
\TextAt{9}{362.184}{672.905}{5}%
\TextAt{13}{368.037}{672.905}{\Character{59}}%
\TextAt{16}{373.281}{672.905}{\Character{0}}%
\TextAt{9}{382.579}{672.905}{1}%
\TextAt{16}{388.432}{672.905}{\Character{103}}%
\TextAt{8}{394.41}{672.905}{.}%
\TextAt{8}{404.643}{672.905}{We}%
\TextAt{8}{425.382}{672.905}{set}%
\TextAt{13}{444.244}{672.905}{U}%
\TextAt{9}{457.778}{672.905}{\Character{61}}%
\TextAt{13}{471.217}{672.905}{U}%
\TextAt{25}{479.168}{671.112}{Q}%
\TextAt{24}{485.367}{670.005}{2}%
\TextAt{8}{490.016}{672.905}{,}%
\TextAt{13}{497.908}{672.905}{U}%
\TextAt{17}{507.108}{677.243}{\Character{40}}%
\TextAt{21}{510.401}{677.243}{n}%
\TextAt{17}{515.539}{677.243}{\Character{41}}%
\TextAt{9}{523.665}{672.905}{\Character{61}}%
\TextAt{13}{79.226}{655.859}{U}%
\TextAt{17}{88.425}{661.969}{\Character{40}}%
\TextAt{21}{91.718}{661.969}{n}%
\TextAt{17}{96.857}{661.969}{\Character{41}}%
\TextAt{25}{87.176}{652.574}{Q}%
\TextAt{24}{93.375}{651.467}{2}%
\TextAt{8}{100.648}{655.859}{.}%
\TextAt{8}{110.172}{655.859}{By}%
\TextAt{8}{128.894}{655.859}{chap.}%
\TextAt{8}{162.161}{655.859}{II,}%
\TextAt{8}{178.432}{655.859}{(5.5),}%
\TextAt{8}{210.107}{655.859}{exp}%
\TextAt{9}{231.924}{655.859}{\Character{58}}%
\TextAt{9}{239.085}{655.859}{2}%
\TextAt{21}{244.961}{660.197}{n}%
\TextAt{18}{250.597}{655.859}{Z}%
\TextAt{17}{258.567}{654.066}{2}%
\NativeMathAt{267.231}{655.859}{11.955}{\rightarrow}%
\TextAt{13}{283.119}{655.859}{U}%
\TextAt{17}{292.318}{660.197}{\Character{40}}%
\TextAt{21}{295.611}{660.198}{n}%
\TextAt{17}{300.749}{660.197}{\Character{41}}%
\TextAt{8}{308.796}{655.859}{is}%
\TextAt{8}{320.919}{655.859}{an}%
\TextAt{8}{337.537}{655.859}{isomorphism}%
\TextAt{8}{406.315}{655.859}{for}%
\TextAt{13}{424.558}{655.859}{n}%
\TextAt{13}{435.461}{655.859}{\Character{62}}%
\TextAt{9}{448.514}{655.859}{1}%
\TextAt{8}{454.367}{655.859}{.}%
\TextAt{8}{463.892}{655.859}{Since}%
\TextAt{13}{494.808}{655.859}{a}%
\NativeArrowAt{504.884}{655.859}{11.955}{11.955}{\mapsto}%
\TextAt{9}{520.772}{655.859}{2}%
\TextAt{13}{526.625}{655.859}{a}%
\TextAt{8}{79.226}{638.813}{defines}%
\TextAt{8}{118.471}{638.813}{an}%
\TextAt{8}{135.551}{638.813}{isomorphism}%
\TextAt{9}{204.791}{638.813}{2}%
\TextAt{17}{210.644}{643.152}{2}%
\TextAt{18}{215.377}{638.813}{Z}%
\TextAt{17}{223.347}{637.02}{2}%
\NativeMathAt{232.799}{638.813}{11.955}{\rightarrow}%
\TextAt{9}{249.474}{638.813}{2}%
\TextAt{17}{255.327}{643.152}{3}%
\TextAt{18}{260.059}{638.813}{Z}%
\TextAt{17}{268.03}{637.02}{2}%
\TextAt{8}{272.762}{638.813}{,}%
\TextAt{13}{280.937}{638.813}{x}%
\NativeArrowAt{292.309}{638.813}{11.955}{11.955}{\mapsto}%
\TextAt{13}{308.984}{638.813}{x}%
\TextAt{17}{315.636}{643.152}{2}%
\TextAt{8}{325.087}{638.813}{defines}%
\TextAt{8}{364.332}{638.813}{an}%
\TextAt{8}{381.412}{638.813}{isomorphism}%
\TextAt{13}{450.652}{638.813}{U}%
\TextAt{17}{459.852}{643.152}{\Character{40}2\Character{41}}%
\NativeMathAt{475.89}{638.813}{11.955}{\rightarrow}%
\TextAt{13}{492.565}{638.813}{U}%
\TextAt{17}{501.764}{643.152}{\Character{40}3\Character{41}}%
\TextAt{8}{513.083}{638.813}{.}%
\TextAt{8}{523.994}{638.813}{It}%
\TextAt{8}{79.226}{621.767}{follows}%
\TextAt{8}{118.95}{621.767}{that}%
\TextAt{13}{145.606}{621.767}{U}%
\TextAt{17}{154.805}{626.106}{\Character{40}3\Character{41}}%
\NativeMathAt{171.639}{621.767}{11.955}{\subseteq}%
\TextAt{13}{186.453}{621.767}{U}%
\TextAt{17}{195.652}{626.106}{2}%
\TextAt{8}{200.385}{621.767}{.}%
\TextAt{8}{212.696}{621.767}{Since}%
\TextAt{16}{244.541}{621.767}{\Character{102}}%
\TextAt{9}{250.519}{621.767}{1}%
\TextAt{13}{256.372}{621.767}{\Character{59}}%
\TextAt{16}{261.616}{621.767}{\Character{0}}%
\TextAt{9}{270.915}{621.767}{1}%
\TextAt{13}{276.767}{621.767}{\Character{59}}%
\TextAt{9}{282.012}{621.767}{5}%
\TextAt{13}{287.865}{621.767}{\Character{59}}%
\TextAt{16}{293.109}{621.767}{\Character{0}}%
\TextAt{9}{302.407}{621.767}{5}%
\TextAt{16}{308.26}{621.767}{\Character{103}}%
\TextAt{8}{319.423}{621.767}{is}%
\TextAt{8}{332.475}{621.767}{a}%
\TextAt{8}{343.518}{621.767}{system}%
\TextAt{8}{383.624}{621.767}{of}%
\TextAt{8}{398.242}{621.767}{representatives}%
\TextAt{8}{479.33}{621.767}{of}%
\TextAt{13}{493.948}{621.767}{U}%
\TextAt{9}{503.147}{621.767}{\Character{47}}%
\TextAt{13}{509}{621.767}{U}%
\TextAt{17}{518.199}{626.106}{\Character{40}3\Character{41}}%
\TextAt{8}{529.518}{621.767}{,}%
\TextAt{13}{79.226}{604.721}{U}%
\TextAt{9}{88.425}{604.721}{\Character{47}}%
\TextAt{13}{94.278}{604.721}{U}%
\TextAt{17}{103.477}{609.06}{2}%
\TextAt{8}{112.107}{604.721}{is}%
\TextAt{8}{123.871}{604.721}{generated}%
\TextAt{8}{177.203}{604.721}{by}%
\TextAt{16}{193.462}{604.721}{\Character{0}}%
\TextAt{9}{202.76}{604.721}{1}%
\TextAt{8}{212.511}{604.721}{and}%
\TextAt{9}{235.273}{604.721}{5}%
\TextAt{8}{241.126}{604.721}{.}%
\Vector{q 0 G 0 g 0.399 w 452.334 -169.279 m 452.334 -161.209 l S Q}%
\Vector{q 0 G 0 g 0.399 w 452.534 -161.408 m 459.707 -161.408 l S Q}%
\Vector{q 0 G 0 g 0.399 w 452.534 -169.08 m 459.707 -169.08 l S Q}%
\Vector{q 0 G 0 g 0.399 w 459.906 -169.279 m 459.906 -161.209 l S Q}%
\TextAt{8}{97.159}{575.72}{It}%
\TextAt{8}{109.845}{575.72}{is}%
\TextAt{8}{121.623}{575.72}{much}%
\TextAt{8}{152.852}{575.72}{more}%
\TextAt{8}{182.132}{575.72}{difficult}%
\TextAt{8}{225.065}{575.72}{to}%
\TextAt{8}{239.39}{575.72}{determine}%
\TextAt{8}{294.027}{575.72}{the}%
\TextAt{13}{314.197}{575.72}{n}%
\TextAt{8}{321.185}{575.72}{-th}%
\TextAt{8}{340.052}{575.72}{Hilbert}%
\TextAt{8}{380.379}{575.72}{symbol}%
\TextAt{8}{420.455}{575.72}{in}%
\TextAt{8}{434.121}{575.72}{the}%
\TextAt{8}{454.292}{575.72}{general}%
\TextAt{8}{494.63}{575.72}{case.}%
\TextAt{8}{523.994}{575.72}{It}%
\TextAt{8}{79.226}{558.674}{was}%
\TextAt{8}{101.209}{558.674}{discovered}%
\TextAt{8}{157.013}{558.674}{only}%
\TextAt{8}{182.188}{558.674}{in}%
\TextAt{8}{195.324}{558.674}{1964}%
\TextAt{8}{222.137}{558.674}{by}%
\TextAt{8}{237.88}{558.674}{the}%
\TextAt{8}{257.52}{558.674}{mathematician}%
\TextAt{26}{337.007}{558.674}{Helmut}%
\TextAt{26}{386.787}{558.674}{Brückner}%
\TextAt{8}{446.132}{558.674}{.}%
\TextAt{8}{454.408}{558.674}{Since}%
\TextAt{8}{484.449}{558.674}{the}%
\TextAt{8}{504.089}{558.674}{result}%
\TextAt{8}{79.226}{541.628}{has}%
\TextAt{8}{100.431}{541.628}{not}%
\TextAt{8}{121.577}{541.628}{previously}%
\TextAt{8}{177.906}{541.628}{been}%
\TextAt{8}{205.866}{541.628}{published}%
\TextAt{8}{258.931}{541.628}{in}%
\TextAt{8}{272.916}{541.628}{an}%
\TextAt{8}{289.506}{541.628}{easily}%
\TextAt{8}{322.093}{541.628}{accessible}%
\TextAt{8}{375.218}{541.628}{place,}%
\TextAt{8}{408.796}{541.628}{we}%
\TextAt{8}{426.355}{541.628}{state}%
\TextAt{8}{455.366}{541.628}{it}%
\TextAt{8}{467.402}{541.628}{here}%
\TextAt{8}{493.09}{541.628}{without}%
\TextAt{8}{79.226}{524.582}{proof}%
\TextAt{8}{109.796}{524.582}{for}%
\TextAt{8}{127.68}{524.582}{the}%
\TextAt{8}{147.837}{524.582}{case}%
\TextAt{13}{172.608}{524.582}{n}%
\TextAt{9}{182.916}{524.582}{\Character{61}}%
\TextAt{13}{195.342}{524.582}{p}%
\TextAt{21}{201.217}{528.921}{\Character{23}}%
\TextAt{8}{210.296}{524.582}{of}%
\TextAt{8}{223.626}{524.582}{odd}%
\TextAt{8}{246.712}{524.582}{residue}%
\TextAt{8}{286.439}{524.582}{characteristic}%
\TextAt{13}{358.373}{524.582}{p}%
\TextAt{8}{368.146}{524.582}{of}%
\TextAt{13}{381.476}{524.582}{K}%
\TextAt{8}{392.293}{524.582}{.}%
\TextAt{8}{97.159}{504.547}{So}%
\TextAt{8}{113.605}{504.547}{let}%
\TextAt{13}{130.696}{504.547}{\Character{22}}%
\TextAt{21}{137.739}{502.754}{p}%
\TextAt{22}{142.002}{505.023}{\Character{23}}%
\NativeMathAt{150.761}{504.547}{11.955}{\subseteq}%
\TextAt{13}{163.698}{504.547}{K}%
\TextAt{8}{174.516}{504.547}{,}%
\TextAt{8}{181.899}{504.547}{choose}%
\TextAt{8}{219.218}{504.547}{a}%
\TextAt{8}{229.161}{504.547}{prime}%
\TextAt{8}{262.511}{504.547}{element}%
\TextAt{13}{305.94}{504.547}{\Character{25}}%
\TextAt{8}{317.093}{504.547}{of}%
\TextAt{13}{330.61}{504.547}{K}%
\TextAt{8}{341.428}{504.547}{,}%
\TextAt{8}{348.81}{504.547}{and}%
\TextAt{8}{371.76}{504.547}{let}%
\TextAt{13}{388.851}{504.547}{W}%
\TextAt{8}{405.613}{504.547}{be}%
\TextAt{8}{421.724}{504.547}{the}%
\TextAt{8}{442.067}{504.547}{ring}%
\TextAt{8}{466.32}{504.547}{of}%
\TextAt{8}{479.836}{504.547}{integers}%
\TextAt{8}{523.337}{504.547}{of}%
\TextAt{8}{79.226}{487.501}{the}%
\TextAt{8}{99.802}{487.501}{maximal}%
\TextAt{8}{148.031}{487.501}{unramified}%
\TextAt{8}{206.983}{487.501}{subextension}%
\TextAt{13}{277.114}{487.501}{T}%
\TextAt{8}{289.918}{487.501}{of}%
\TextAt{13}{303.668}{487.501}{K}%
\TextAt{16}{314.486}{487.501}{\Character{106}}%
\TextAt{18}{317.806}{487.501}{Q}%
\TextAt{21}{327.105}{485.708}{p}%
\TextAt{8}{336.183}{487.501}{(i.e.,}%
\TextAt{8}{363.368}{487.501}{the}%
\TextAt{8}{383.944}{487.501}{ring}%
\TextAt{8}{408.43}{487.501}{of}%
\TextAt{8}{422.18}{487.501}{Witt}%
\TextAt{8}{450.886}{487.501}{vectors}%
\TextAt{8}{491.044}{487.501}{over}%
\TextAt{8}{516.51}{487.501}{the}%
\TextAt{8}{79.226}{470.455}{residue}%
\TextAt{8}{118.953}{470.455}{class}%
\TextAt{8}{146.39}{470.455}{field}%
\TextAt{8}{171.747}{470.455}{of}%
\TextAt{13}{185.077}{470.455}{K}%
\TextAt{8}{195.895}{470.455}{).}%
\TextAt{8}{208.898}{470.455}{Then}%
\TextAt{8}{239.456}{470.455}{every}%
\TextAt{8}{270.348}{470.455}{element}%
\TextAt{13}{313.59}{470.455}{x}%
\NativeMathAt{323.562}{470.455}{11.955}{\in}%
\TextAt{13}{334.853}{470.455}{K}%
\TextAt{8}{349.569}{470.455}{can}%
\TextAt{8}{371.028}{470.455}{be}%
\TextAt{8}{386.952}{470.455}{written}%
\TextAt{8}{427.923}{470.455}{in}%
\TextAt{8}{441.576}{470.455}{the}%
\TextAt{8}{461.732}{470.455}{form}%
\TextAt{13}{281.563}{441.454}{x}%
\TextAt{9}{291.536}{441.454}{\Character{61}}%
\TextAt{13}{303.961}{441.454}{f}%
\TextAt{9}{311.007}{441.454}{\Character{40}}%
\TextAt{13}{315.56}{441.454}{\Character{25}}%
\TextAt{9}{322.629}{441.454}{\Character{41}}%
\TextAt{13}{327.181}{441.454}{\Character{59}}%
\TextAt{8}{79.226}{412.453}{with}%
\TextAt{8}{105.886}{412.453}{a}%
\TextAt{8}{115.642}{412.453}{Laurent}%
\TextAt{8}{159.708}{412.453}{series}%
\TextAt{13}{191.043}{412.453}{f}%
\TextAt{9}{198.089}{412.453}{\Character{40}}%
\TextAt{13}{202.642}{412.453}{X}%
\TextAt{9}{213.297}{412.453}{\Character{41}}%
\NativeMathAt{221.17}{412.453}{11.955}{\in}%
\TextAt{13}{232.461}{412.453}{W}%
\TextAt{9}{245.138}{412.453}{\Character{40}\Character{40}}%
\TextAt{13}{254.243}{412.453}{X}%
\TextAt{9}{264.898}{412.453}{\Character{41}\Character{41}}%
\TextAt{8}{274.003}{412.453}{.}%
\TextAt{8}{97.159}{392.418}{For}%
\TextAt{8}{117.731}{392.418}{an}%
\TextAt{8}{133.594}{392.418}{arbitrary}%
\TextAt{8}{182.966}{392.418}{Laurent}%
\TextAt{8}{226.636}{392.418}{series}%
\TextAt{13}{257.574}{392.418}{f}%
\TextAt{9}{264.62}{392.418}{\Character{40}}%
\TextAt{13}{269.173}{392.418}{X}%
\TextAt{9}{279.828}{392.418}{\Character{41}}%
\TextAt{9}{287.683}{392.418}{\Character{61}}%
\TextAt{20}{300.126}{401.385}{\Character{80}}%
\TextAt{21}{312.746}{388.931}{i}%
\NativeMathAt{315.629}{388.931}{7.97}{\geq}%
\TextAt{3}{322.212}{388.931}{\Character{0}}%
\TextAt{21}{328.802}{388.931}{m}%
\TextAt{13}{338.783}{392.418}{a}%
\TextAt{21}{344.928}{390.625}{i}%
\TextAt{13}{348.309}{392.418}{X}%
\TextAt{21}{358.964}{396.757}{i}%
\NativeMathAt{365.667}{392.418}{11.955}{\in}%
\TextAt{13}{376.958}{392.418}{W}%
\TextAt{9}{389.635}{392.418}{\Character{40}\Character{40}}%
\TextAt{13}{398.74}{392.418}{X}%
\TextAt{9}{409.395}{392.418}{\Character{41}\Character{41}}%
\TextAt{8}{418.499}{392.418}{,}%
\TextAt{8}{425.331}{392.418}{let}%
\TextAt{13}{441.839}{392.418}{f}%
\TextAt{21}{448.885}{396.757}{P}%
\TextAt{9}{455.928}{392.418}{\Character{40}}%
\TextAt{13}{460.481}{392.418}{X}%
\TextAt{9}{471.136}{392.418}{\Character{41}}%
\TextAt{8}{479.189}{392.418}{denote}%
\TextAt{8}{516.511}{392.418}{the}%
\TextAt{8}{79.226}{375.372}{series}%
\TextAt{13}{254.713}{358.326}{f}%
\TextAt{21}{261.76}{363.263}{P}%
\TextAt{9}{268.803}{358.326}{\Character{40}}%
\TextAt{13}{273.355}{358.326}{X}%
\TextAt{9}{284.01}{358.326}{\Character{41}}%
\TextAt{9}{291.865}{358.326}{\Character{61}}%
\TextAt{20}{304.309}{369.684}{\Character{88}}%
\TextAt{21}{311.502}{344.49}{i}%
\TextAt{13}{323.57}{358.326}{a}%
\TextAt{21}{329.715}{363.994}{\Character{39}}%
\TextAt{21}{329.715}{355.263}{i}%
\TextAt{13}{335.731}{358.326}{X}%
\TextAt{21}{346.387}{363.263}{ip}%
\TextAt{13}{354.031}{358.326}{\Character{59}}%
\TextAt{8}{79.226}{326.856}{where}%
\TextAt{13}{113.567}{326.856}{\Character{39}}%
\TextAt{8}{125.664}{326.856}{is}%
\TextAt{8}{137.96}{326.856}{the}%
\TextAt{8}{158.648}{326.856}{Frobenius}%
\TextAt{8}{213.049}{326.856}{automorphism}%
\TextAt{8}{291.051}{326.856}{of}%
\TextAt{13}{304.912}{326.856}{W}%
\TextAt{8}{317.59}{326.856}{.}%
\TextAt{8}{327.633}{326.856}{Further,}%
\TextAt{8}{373.978}{326.856}{let}%
\TextAt{8}{391.414}{326.856}{Res}%
\TextAt{9}{409.849}{326.856}{\Character{40}}%
\TextAt{13}{414.401}{326.856}{f}%
\TextAt{13}{423.427}{326.856}{dX}%
\TextAt{9}{440.178}{326.856}{\Character{41}}%
\NativeMathAt{448.957}{326.856}{11.955}{\in}%
\TextAt{13}{461.154}{326.856}{W}%
\TextAt{8}{478.26}{326.856}{denote}%
\TextAt{8}{516.511}{326.856}{the}%
\TextAt{8}{79.226}{309.81}{residue}%
\TextAt{8}{118.953}{309.81}{of}%
\TextAt{8}{132.283}{309.81}{the}%
\TextAt{8}{152.44}{309.81}{differential}%
\TextAt{13}{210.972}{309.81}{f}%
\TextAt{13}{219.998}{309.81}{dX}%
\TextAt{8}{236.749}{309.81}{,}%
\TextAt{13}{263.676}{277.68}{d}%
\TextAt{8}{271.751}{277.68}{log}%
\TextAt{13}{288.711}{277.68}{f}%
\TextAt{9}{299.079}{277.68}{\Character{58}\Character{61}}%
\TextAt{13}{315.951}{285.767}{f}%
\TextAt{3}{322.998}{290.106}{\Character{48}}%
\Vector{q 0 G 0 g 0.478 w 243.951 -493.332 m 253.793 -493.332 l S Q}%
\TextAt{13}{317.349}{269.479}{f}%
\TextAt{13}{328.981}{277.68}{dX}%
\TextAt{13}{345.049}{277.68}{\Character{59}}%
\TextAt{8}{79.226}{245.906}{and}%
\TextAt{8}{233.812}{225.219}{log}%
\TextAt{13}{250.773}{225.219}{f}%
\TextAt{9}{261.14}{225.219}{\Character{58}\Character{61}}%
\NativeMathAt{281.217}{240.163}{7.97}{\infty}%
\TextAt{20}{276.817}{236.577}{\Character{88}}%
\TextAt{21}{278.599}{211.382}{i}%
\TextAt{17}{281.482}{211.382}{\Character{61}1}%
\TextAt{9}{294.086}{225.219}{\Character{40}}%
\TextAt{16}{298.638}{225.219}{\Character{0}}%
\TextAt{9}{307.936}{225.219}{1\Character{41}}%
\TextAt{21}{318.342}{230.155}{i}%
\TextAt{17}{321.225}{230.155}{\Character{43}1}%
\TextAt{9}{333.739}{233.307}{\Character{40}}%
\TextAt{13}{338.291}{233.307}{f}%
\TextAt{16}{347.995}{233.307}{\Character{0}}%
\TextAt{9}{359.95}{233.307}{1\Character{41}}%
\TextAt{21}{370.355}{237.645}{i}%
\Vector{q 0 G 0 g 0.478 w 261.739 -545.792 m 301.736 -545.792 l S Q}%
\TextAt{13}{351.741}{217.018}{i}%
\TextAt{13}{374.932}{225.219}{\Character{59}}%
\TextAt{8}{79.226}{193.748}{if}%
\TextAt{13}{89.95}{193.748}{f}%
\NativeMathAt{100.317}{193.748}{11.955}{\in}%
\TextAt{9}{111.608}{193.748}{1}%
\TextAt{9}{120.108}{193.748}{\Character{43}}%
\TextAt{13}{131.879}{193.748}{pW}%
\TextAt{9}{150.432}{193.748}{\Character{91}\Character{91}}%
\TextAt{13}{156.935}{193.748}{X}%
\TextAt{9}{167.59}{193.748}{\Character{93}\Character{93}}%
\TextAt{8}{174.093}{193.748}{.}%
\TextAt{8}{97.159}{173.714}{Now}%
\TextAt{8}{123.786}{173.714}{let}%
\TextAt{13}{140.657}{173.714}{\Character{16}}%
\TextAt{8}{150.537}{173.714}{be}%
\TextAt{8}{166.429}{173.714}{a}%
\TextAt{8}{176.151}{173.714}{primitive}%
\TextAt{13}{226.198}{173.714}{p}%
\TextAt{21}{232.073}{178.052}{\Character{23}}%
\TextAt{8}{237.255}{173.714}{-th}%
\TextAt{8}{256.076}{173.714}{root}%
\TextAt{8}{281.089}{173.714}{of}%
\TextAt{8}{294.386}{173.714}{unity.}%
\TextAt{8}{328.515}{173.714}{Then}%
\TextAt{9}{359.039}{173.714}{1}%
\TextAt{16}{367.481}{173.714}{\Character{0}}%
\TextAt{13}{379.369}{173.714}{\Character{16}}%
\TextAt{8}{389.248}{173.714}{is}%
\TextAt{8}{400.979}{173.714}{a}%
\TextAt{8}{410.702}{173.714}{prime}%
\TextAt{8}{443.832}{173.714}{element}%
\TextAt{8}{487.041}{173.714}{of}%
\TextAt{18}{500.338}{173.714}{Q}%
\TextAt{21}{509.637}{171.92}{p}%
\TextAt{9}{514.397}{173.714}{\Character{40}}%
\TextAt{13}{518.95}{173.714}{\Character{16}}%
\TextAt{9}{524.965}{173.714}{\Character{41}}%
\TextAt{8}{529.518}{173.714}{,}%
\TextAt{8}{79.226}{156.668}{and}%
\TextAt{8}{101.989}{156.668}{thus}%
\TextAt{9}{274.772}{139.622}{1}%
\TextAt{16}{283.281}{139.622}{\Character{0}}%
\TextAt{13}{295.237}{139.622}{\Character{16}}%
\TextAt{9}{304.573}{139.622}{\Character{61}}%
\TextAt{13}{316.998}{139.622}{\Character{25}}%
\TextAt{21}{324.068}{144.558}{e}%
\TextAt{13}{328.494}{139.622}{\Character{34}\Character{59}}%
\TextAt{8}{79.226}{116.1}{for}%
\TextAt{8}{97.417}{116.1}{some}%
\TextAt{8}{127.05}{116.1}{unit}%
\TextAt{13}{152.068}{116.1}{\Character{34}}%
\TextAt{8}{161.749}{116.1}{of}%
\TextAt{13}{175.386}{116.1}{K}%
\TextAt{8}{186.204}{116.1}{,}%
\TextAt{8}{193.736}{116.1}{where}%
\TextAt{13}{227.851}{116.1}{e}%
\TextAt{8}{237.481}{116.1}{is}%
\TextAt{8}{249.551}{116.1}{the}%
\TextAt{8}{270.014}{116.1}{ramification}%
\TextAt{8}{335.368}{116.1}{index}%
\TextAt{8}{367.212}{116.1}{of}%
\TextAt{13}{380.849}{116.1}{K}%
\TextAt{16}{391.666}{116.1}{\Character{106}}%
\TextAt{18}{394.987}{116.1}{Q}%
\TextAt{21}{404.286}{114.307}{p}%
\TextAt{9}{409.047}{116.1}{\Character{40}}%
\TextAt{13}{413.599}{116.1}{\Character{16}}%
\TextAt{9}{419.615}{116.1}{\Character{41}}%
\TextAt{8}{424.167}{116.1}{.}%
\TextAt{8}{433.534}{116.1}{Let}%
\TextAt{13}{454.81}{116.1}{\Character{17}}%
\TextAt{9}{461.015}{116.1}{\Character{40}}%
\TextAt{13}{465.567}{116.1}{X}%
\TextAt{9}{476.222}{116.1}{\Character{41}}%
\NativeMathAt{484.617}{116.1}{11.955}{\in}%
\TextAt{13}{496.43}{116.1}{W}%
\TextAt{9}{509.108}{116.1}{\Character{91}\Character{91}}%
\TextAt{13}{515.611}{116.1}{X}%
\TextAt{9}{526.266}{116.1}{\Character{93}\Character{93}}%
\TextAt{8}{79.226}{99.054}{be}%
\TextAt{8}{95.15}{99.054}{a}%
\TextAt{8}{104.906}{99.054}{power}%
\TextAt{8}{139.05}{99.054}{series}%
\TextAt{8}{170.384}{99.054}{such}%
\TextAt{8}{196.781}{99.054}{that}%
\TextAt{13}{282.571}{82.008}{\Character{34}}%
\TextAt{9}{291.369}{82.008}{\Character{61}}%
\TextAt{13}{303.795}{82.008}{\Character{17}}%
\TextAt{9}{309.999}{82.008}{\Character{40}}%
\TextAt{13}{314.552}{82.008}{\Character{25}}%
\TextAt{9}{321.621}{82.008}{\Character{41}}%
\TextAt{13}{326.173}{82.008}{\Character{59}}%
\TextAt{8}{297.211}{43.033}{310}%
\Destination{78.23}{809.815}{nav.page.310}%
\Destination{79.36}{759.055}{nav.print.338}%
\Hotlink{161.661}{653.655}{9.763}{11.047}{nav.chapter.2}%
\Hotlink{177.932}{653.655}{25.078}{11.047}{nav.number.2.5.5}%
}

% Page 323 (body page 311)
\Page{%
\TextAt{2}{79.226}{796.943}{V.}%
\TextAt{2}{94.073}{796.943}{Local}%
\TextAt{2}{123.615}{796.943}{Class}%
\TextAt{2}{152.208}{796.943}{Field}%
\TextAt{2}{179.939}{796.943}{Theory}%
\Vector{q 0 G 0 g 0.299 w 7.226 18.729 m 460.769 18.729 l S Q}%
\TextAt{8}{79.226}{761.124}{and}%
\TextAt{8}{101.989}{761.124}{let}%
\TextAt{13}{118.893}{761.124}{h}%
\TextAt{9}{125.632}{761.124}{\Character{40}}%
\TextAt{13}{130.184}{761.124}{X}%
\TextAt{9}{140.839}{761.124}{\Character{41}}%
\TextAt{8}{149.289}{761.124}{be}%
\TextAt{8}{165.213}{761.124}{the}%
\TextAt{8}{185.37}{761.124}{series}%
\TextAt{13}{124.417}{726.526}{h}%
\TextAt{9}{131.155}{726.526}{\Character{40}}%
\TextAt{13}{135.708}{726.526}{X}%
\TextAt{9}{146.363}{726.526}{\Character{41}}%
\TextAt{9}{154.218}{726.526}{\Character{61}}%
\TextAt{9}{167.857}{734.614}{1}%
\Vector{q 0 G 0 g 0.478 w 95.857 -44.485 m 101.71 -44.485 l S Q}%
\TextAt{9}{167.857}{718.325}{2}%
\TextAt{9}{176.198}{734.614}{1}%
\TextAt{9}{184.698}{734.614}{\Character{43}}%
\TextAt{9}{196.45}{734.614}{\Character{40}1}%
\TextAt{16}{209.531}{734.614}{\Character{0}}%
\TextAt{13}{221.486}{734.614}{X}%
\TextAt{21}{232.141}{738.952}{e}%
\TextAt{13}{236.567}{734.614}{\Character{17}}%
\TextAt{9}{242.772}{734.614}{\Character{40}}%
\TextAt{13}{247.324}{734.614}{X}%
\TextAt{9}{257.979}{734.614}{\Character{41}\Character{41}}%
\TextAt{21}{267.084}{738.952}{p}%
\TextAt{22}{271.347}{741.765}{\Character{23}}%
\Vector{q 0 G 0 g 0.478 w 104.101 -44.485 m 204.563 -44.485 l S Q}%
\TextAt{9}{176.101}{718.325}{1}%
\TextAt{16}{184.61}{718.325}{\Character{0}}%
\TextAt{9}{196.566}{718.325}{\Character{40}1}%
\TextAt{16}{209.628}{718.325}{\Character{0}}%
\TextAt{13}{221.583}{718.325}{X}%
\TextAt{21}{232.238}{721.779}{e}%
\TextAt{13}{236.664}{718.325}{\Character{17}}%
\TextAt{9}{242.869}{718.325}{\Character{40}}%
\TextAt{13}{247.421}{718.325}{X}%
\TextAt{9}{258.076}{718.325}{\Character{41}\Character{41}}%
\TextAt{21}{267.181}{721.779}{p}%
\TextAt{22}{271.444}{724.048}{\Character{23}}%
\TextAt{9}{281.08}{726.526}{\Character{61}}%
\NativeMathAt{301.533}{741.47}{7.97}{\infty}%
\TextAt{20}{297.133}{737.883}{\Character{88}}%
\TextAt{21}{293.505}{712.689}{i}%
\TextAt{17}{296.388}{712.689}{\Character{61}}%
\TextAt{3}{302.975}{712.689}{\Character{0}}%
\NativeMathAt{309.558}{712.689}{7.97}{\infty}%
\TextAt{13}{320.022}{726.526}{a}%
\TextAt{21}{326.167}{724.733}{i}%
\TextAt{13}{329.548}{726.526}{X}%
\TextAt{21}{340.203}{731.462}{i}%
\TextAt{13}{343.585}{726.526}{\Character{59}}%
\TextAt{13}{372.72}{726.526}{a}%
\TextAt{21}{378.884}{724.733}{i}%
\NativeMathAt{385.586}{726.526}{11.955}{\in}%
\TextAt{13}{396.877}{726.526}{W}%
\TextAt{13}{407.601}{726.526}{\Character{59}}%
\TextAt{8}{429.877}{726.526}{lim}%
\TextAt{21}{424.803}{719.264}{i}%
\NativeMathAt{427.686}{719.264}{7.97}{\rightarrow}%
\TextAt{3}{436.15}{719.264}{\Character{0}}%
\NativeMathAt{442.733}{719.264}{7.97}{\infty}%
\TextAt{13}{453.202}{726.526}{a}%
\TextAt{21}{459.347}{724.733}{i}%
\TextAt{9}{466.049}{726.526}{\Character{61}}%
\TextAt{9}{478.459}{726.526}{0}%
\TextAt{13}{484.327}{726.526}{\Character{58}}%
\TextAt{8}{79.226}{689.874}{With}%
\TextAt{8}{109.724}{689.874}{this}%
\TextAt{8}{132.809}{689.874}{notation}%
\TextAt{8}{179.913}{689.874}{we}%
\TextAt{8}{197.403}{689.874}{can}%
\TextAt{8}{219.126}{689.874}{now}%
\TextAt{8}{243.778}{689.874}{state}%
\TextAt{26}{272.721}{689.874}{Brückner}%
\TextAt{8}{332.067}{689.874}{’s}%
\TextAt{8}{344.094}{689.874}{formula}%
\TextAt{8}{387.288}{689.874}{for}%
\TextAt{8}{405.436}{689.874}{the}%
\TextAt{13}{425.856}{689.874}{p}%
\TextAt{21}{431.731}{694.213}{\Character{23}}%
\TextAt{8}{436.913}{689.874}{-th}%
\TextAt{8}{456.029}{689.874}{Hilbert}%
\TextAt{8}{496.605}{689.874}{symbol}%
\TextAt{20}{79.226}{682.512}{\Character{0}}%
\TextAt{21}{85.901}{678.084}{x\Character{59}y}%
\Vector{q 0 G 0 g 0.478 w 13.901 -98.183 m 25.492 -98.183 l S Q}%
\TextAt{27}{89.702}{668.706}{p}%
\TextAt{20}{98.688}{682.512}{\Character{1}}%
\TextAt{8}{104.167}{672.828}{,}%
\TextAt{13}{111.316}{672.828}{p}%
\TextAt{9}{120.512}{672.828}{\Character{61}}%
\TextAt{8}{132.938}{672.828}{char}%
\TextAt{9}{155.055}{672.828}{\Character{40}}%
\TextAt{13}{159.607}{672.828}{\Character{20}}%
\TextAt{9}{166.346}{672.828}{\Character{41}}%
\NativeMathAt{174.219}{672.828}{11.955}{\neq}%
\TextAt{9}{186.629}{672.828}{2}%
\TextAt{8}{192.497}{672.828}{.}%
\TextAt{12}{79.226}{647.077}{Theorem}%
\TextAt{12}{135.627}{647.077}{(3.7).}%
\TextAt{8}{170.409}{647.077}{If}%
\TextAt{13}{181.568}{647.077}{x\Character{59}}%
\TextAt{13}{193.44}{647.077}{y}%
\NativeMathAt{202.922}{647.077}{11.955}{\in}%
\TextAt{13}{214.213}{647.077}{K}%
\TextAt{3}{225.031}{651.415}{\Character{3}}%
\TextAt{8}{233.127}{647.077}{and}%
\TextAt{13}{255.357}{647.077}{f}%
\TextAt{13}{261.741}{647.077}{\Character{59}}%
\TextAt{13}{266.978}{647.077}{g}%
\NativeMathAt{276.352}{647.077}{11.955}{\in}%
\TextAt{13}{287.643}{647.077}{W}%
\TextAt{9}{300.32}{647.077}{\Character{40}\Character{40}}%
\TextAt{13}{309.425}{647.077}{X}%
\TextAt{9}{320.08}{647.077}{\Character{41}\Character{41}}%
\TextAt{3}{329.185}{651.415}{\Character{3}}%
\TextAt{8}{337.281}{647.077}{such}%
\TextAt{8}{363.145}{647.077}{that}%
\TextAt{13}{387.981}{647.077}{f}%
\TextAt{9}{395.027}{647.077}{\Character{40}}%
\TextAt{13}{399.579}{647.077}{\Character{25}}%
\TextAt{9}{406.649}{647.077}{\Character{41}}%
\TextAt{9}{414.504}{647.077}{\Character{61}}%
\TextAt{13}{426.947}{647.077}{x}%
\TextAt{8}{436.964}{647.077}{and}%
\TextAt{13}{459.193}{647.077}{g}%
\TextAt{9}{465.227}{647.077}{\Character{40}}%
\TextAt{13}{469.78}{647.077}{\Character{25}}%
\TextAt{9}{476.849}{647.077}{\Character{41}}%
\TextAt{9}{484.704}{647.077}{\Character{61}}%
\TextAt{13}{497.148}{647.077}{y}%
\TextAt{8}{503.284}{647.077}{,}%
\TextAt{8}{510.007}{647.077}{then}%
\TextAt{20}{263.621}{631.56}{\Character{18}}%
\TextAt{13}{273.617}{622.791}{x\Character{59}}%
\TextAt{13}{285.488}{622.791}{y}%
\Vector{q 0 G 0 g 0.478 w 201.617 -156.308 m 219.65 -156.308 l S Q}%
\TextAt{14}{279.642}{606.502}{p}%
\TextAt{20}{292.845}{631.56}{\Character{19}}%
\TextAt{9}{304.967}{614.703}{\Character{61}}%
\TextAt{13}{317.392}{614.703}{\Character{16}}%
\TextAt{21}{323.408}{619.639}{w}%
\TextAt{17}{329.699}{619.639}{\Character{40}}%
\TextAt{21}{332.992}{619.639}{x\Character{59}y}%
\TextAt{17}{344.583}{619.639}{\Character{41}}%
\TextAt{8}{79.226}{583.011}{where}%
\TextAt{13}{121.92}{553.518}{w}%
\TextAt{9}{130.61}{553.518}{\Character{40}}%
\TextAt{13}{135.162}{553.518}{x\Character{59}}%
\TextAt{13}{147.034}{553.518}{y}%
\TextAt{9}{153.195}{553.518}{\Character{41}}%
\TextAt{9}{161.05}{553.518}{\Character{61}}%
\TextAt{8}{173.494}{553.518}{Tr}%
\TextAt{21}{186.501}{551.659}{W}%
\TextAt{3}{195.667}{551.659}{\Character{106}}%
\TextAt{25}{198.019}{551.659}{Z}%
\TextAt{22}{203.333}{550.662}{p}%
\TextAt{8}{210.162}{553.518}{Res}%
\TextAt{13}{230.589}{553.518}{h}%
\TextAt{16}{239.984}{553.518}{\Character{1}}%
\TextAt{20}{245.962}{570.375}{\Character{18}}%
\TextAt{9}{255.969}{561.606}{1}%
\Vector{q 0 G 0 g 0.478 w 183.958 -217.493 m 189.833 -217.493 l S Q}%
\TextAt{13}{255.958}{545.318}{p}%
\TextAt{8}{265.021}{553.518}{log}%
\TextAt{13}{284.318}{561.606}{f}%
\TextAt{21}{291.364}{565.944}{p}%
\Vector{q 0 G 0 g 0.478 w 211.177 -217.493 m 225.266 -217.493 l S Q}%
\TextAt{13}{283.177}{545.318}{f}%
\TextAt{21}{290.223}{548.771}{P}%
\TextAt{13}{300.454}{553.518}{d}%
\TextAt{8}{308.53}{553.518}{log}%
\TextAt{13}{325.49}{553.518}{g}%
\TextAt{16}{334.181}{553.518}{\Character{0}}%
\TextAt{9}{347.343}{561.606}{1}%
\Vector{q 0 G 0 g 0.478 w 275.332 -217.493 m 281.207 -217.493 l S Q}%
\TextAt{13}{347.332}{545.318}{p}%
\TextAt{8}{356.395}{553.518}{log}%
\TextAt{13}{375.692}{561.606}{g}%
\TextAt{21}{381.726}{565.944}{p}%
\Vector{q 0 G 0 g 0.478 w 302.551 -217.493 m 315.628 -217.493 l S Q}%
\TextAt{13}{374.551}{545.318}{g}%
\TextAt{21}{380.585}{548.771}{P}%
\TextAt{9}{392.023}{561.606}{1}%
\Vector{q 0 G 0 g 0.478 w 320.012 -217.493 m 325.887 -217.493 l S Q}%
\TextAt{13}{392.012}{545.318}{p}%
\TextAt{13}{399.082}{553.518}{d}%
\TextAt{8}{407.157}{553.518}{log}%
\TextAt{13}{424.118}{553.518}{f}%
\TextAt{21}{431.164}{558.454}{P}%
\TextAt{20}{438.207}{570.375}{\Character{19}}%
\TextAt{8}{450.328}{553.518}{mod}%
\TextAt{13}{475.766}{553.518}{p}%
\TextAt{21}{481.642}{558.454}{\Character{23}}%
\TextAt{13}{486.824}{553.518}{\Character{58}}%
\TextAt{8}{97.159}{512.451}{For}%
\TextAt{8}{117.227}{512.451}{the}%
\TextAt{8}{136.483}{512.451}{proof}%
\TextAt{8}{166.152}{512.451}{of}%
\TextAt{8}{178.581}{512.451}{this}%
\TextAt{8}{200.502}{512.451}{theorem,}%
\TextAt{8}{248.559}{512.451}{we}%
\TextAt{8}{264.885}{512.451}{have}%
\TextAt{8}{290.979}{512.451}{to}%
\TextAt{8}{304.389}{512.451}{refer}%
\TextAt{8}{330.471}{512.451}{to}%
\TextAt{8}{343.88}{512.451}{[20]}%
\TextAt{8}{365.096}{512.451}{(see}%
\TextAt{8}{387.664}{512.451}{also}%
\TextAt{8}{410.243}{512.451}{[69]}%
\TextAt{8}{431.459}{512.451}{and}%
\TextAt{8}{453.321}{512.451}{[135]).}%
\TextAt{26}{490.101}{512.451}{Brück-}%
\TextAt{26}{79.226}{495.405}{ner}%
\TextAt{8}{104.536}{495.405}{has}%
\TextAt{8}{125.626}{495.405}{also}%
\TextAt{8}{149.322}{495.405}{deduced}%
\TextAt{8}{195.052}{495.405}{an}%
\TextAt{8}{211.527}{495.405}{explicit}%
\TextAt{8}{253.036}{495.405}{formula}%
\TextAt{8}{296.183}{495.405}{for}%
\TextAt{8}{314.284}{495.405}{the}%
\TextAt{8}{334.657}{495.405}{case}%
\TextAt{13}{359.644}{495.405}{n}%
\TextAt{9}{370.321}{495.405}{\Character{61}}%
\TextAt{9}{383.101}{495.405}{2}%
\TextAt{21}{388.967}{499.744}{\Character{23}}%
\TextAt{8}{394.149}{495.405}{,}%
\TextAt{8}{401.569}{495.405}{but}%
\TextAt{8}{423.244}{495.405}{it}%
\TextAt{8}{435.165}{495.405}{is}%
\TextAt{8}{447.145}{495.405}{much}%
\TextAt{8}{478.576}{495.405}{more}%
\TextAt{8}{508.058}{495.405}{com-}%
\TextAt{8}{79.226}{478.359}{plicated.}%
\TextAt{8}{129.225}{478.359}{A}%
\TextAt{8}{142.306}{478.359}{more}%
\TextAt{8}{171.981}{478.359}{recent}%
\TextAt{8}{207.179}{478.359}{treatment}%
\TextAt{8}{261.9}{478.359}{of}%
\TextAt{8}{275.638}{478.359}{the}%
\TextAt{8}{296.204}{478.359}{theorem,}%
\TextAt{8}{345.491}{478.359}{which}%
\TextAt{8}{379.386}{478.359}{also}%
\TextAt{8}{403.275}{478.359}{includes}%
\TextAt{8}{448.611}{478.359}{the}%
\TextAt{8}{469.176}{478.359}{case}%
\TextAt{13}{494.356}{478.359}{n}%
\TextAt{9}{505.361}{478.359}{\Character{61}}%
\TextAt{9}{518.476}{478.359}{2}%
\TextAt{21}{524.336}{482.698}{\Character{23}}%
\TextAt{8}{529.518}{478.359}{,}%
\TextAt{8}{79.226}{461.313}{has}%
\TextAt{8}{100.1}{461.313}{been}%
\TextAt{8}{127.728}{461.313}{given}%
\TextAt{8}{158.298}{461.313}{by}%
\TextAt{26}{174.557}{461.313}{G.}%
\TextAt{26}{193.063}{461.313}{Henniart}%
\TextAt{8}{252.421}{461.313}{[69].}%
\TextAt{8}{97.159}{441.278}{It}%
\TextAt{8}{109.799}{441.278}{would}%
\TextAt{8}{143.911}{441.278}{be}%
\TextAt{8}{159.803}{441.278}{interesting}%
\TextAt{8}{217.395}{441.278}{to}%
\TextAt{8}{231.673}{441.278}{deduce}%
\TextAt{8}{270.651}{441.278}{from}%
\TextAt{8}{298.259}{441.278}{these}%
\TextAt{8}{328.198}{441.278}{formulae}%
\TextAt{8}{376.297}{441.278}{the}%
\TextAt{8}{396.422}{441.278}{following}%
\TextAt{8}{445.824}{441.278}{classical}%
\TextAt{8}{490.791}{441.278}{result}%
\TextAt{8}{523.337}{441.278}{of}%
\TextAt{26}{79.226}{424.232}{Iwasawa}%
\TextAt{8}{132.235}{424.232}{[80],}%
\TextAt{26}{157.604}{424.232}{Artin}%
\TextAt{8}{195.634}{424.232}{and}%
\TextAt{26}{218.396}{424.232}{Hasse}%
\TextAt{8}{256.976}{424.232}{(see}%
\TextAt{8}{280.444}{424.232}{[9])}%
\TextAt{8}{301.258}{424.232}{relative}%
\TextAt{8}{342.885}{424.232}{to}%
\TextAt{8}{357.196}{424.232}{the}%
\TextAt{8}{377.352}{424.232}{field}%
\TextAt{9}{275.091}{396.145}{\Character{8}}%
\TextAt{21}{283.545}{394.352}{\Character{23}}%
\TextAt{9}{292.048}{396.145}{\Character{61}}%
\TextAt{18}{304.473}{396.145}{Q}%
\TextAt{21}{313.772}{394.352}{p}%
\TextAt{9}{318.533}{396.145}{\Character{40}}%
\TextAt{13}{323.085}{396.145}{\Character{16}}%
\TextAt{9}{329.101}{396.145}{\Character{41}}%
\TextAt{13}{333.653}{396.145}{\Character{59}}%
\TextAt{8}{79.226}{368.057}{where}%
\TextAt{13}{113.155}{368.057}{\Character{16}}%
\TextAt{8}{123.187}{368.057}{is}%
\TextAt{8}{135.07}{368.057}{a}%
\TextAt{8}{144.945}{368.057}{primitive}%
\TextAt{13}{195.145}{368.057}{p}%
\TextAt{21}{201.02}{372.396}{\Character{23}}%
\TextAt{8}{206.202}{368.057}{-th}%
\TextAt{8}{225.174}{368.057}{root}%
\TextAt{8}{250.34}{368.057}{of}%
\TextAt{8}{263.789}{368.057}{unity}%
\TextAt{8}{294.478}{368.057}{(}%
\TextAt{13}{299.033}{368.057}{p}%
\NativeMathAt{308.432}{368.057}{11.955}{\neq}%
\TextAt{9}{321.045}{368.057}{2}%
\TextAt{8}{326.914}{368.057}{).}%
\TextAt{8}{340.275}{368.057}{Putting}%
\TextAt{13}{383.481}{368.057}{\Character{25}}%
\TextAt{9}{394.074}{368.057}{\Character{61}}%
\TextAt{9}{406.687}{368.057}{1}%
\TextAt{16}{415.294}{368.057}{\Character{0}}%
\TextAt{13}{427.331}{368.057}{\Character{16}}%
\TextAt{8}{437.363}{368.057}{and}%
\TextAt{8}{460.245}{368.057}{denoting}%
\TextAt{8}{508.496}{368.057}{by}%
\TextAt{13}{524.874}{368.057}{S}%
\TextAt{8}{79.226}{351.011}{the}%
\TextAt{8}{99.052}{351.011}{trace}%
\TextAt{8}{127.987}{351.011}{map}%
\TextAt{8}{153.671}{351.011}{from}%
\TextAt{9}{180.981}{351.011}{\Character{8}}%
\TextAt{21}{189.435}{349.218}{\Character{23}}%
\TextAt{8}{198.184}{351.011}{to}%
\TextAt{18}{212.164}{351.011}{Q}%
\TextAt{21}{221.462}{349.218}{p}%
\TextAt{8}{226.223}{351.011}{,}%
\TextAt{8}{233.108}{351.011}{we}%
\TextAt{8}{250.004}{351.011}{obtain}%
\TextAt{8}{286.101}{351.011}{for}%
\TextAt{8}{303.655}{351.011}{the}%
\TextAt{13}{323.481}{351.011}{p}%
\TextAt{21}{329.356}{355.35}{\Character{23}}%
\TextAt{8}{334.538}{351.011}{-th}%
\TextAt{8}{353.061}{351.011}{Hilbert}%
\TextAt{8}{393.043}{351.011}{symbol}%
\TextAt{20}{432.774}{360.695}{\Character{0}}%
\TextAt{21}{439.449}{356.267}{x\Character{59}y}%
\Vector{q 0 G 0 g 0.478 w 367.449 -420 m 379.04 -420 l S Q}%
\TextAt{27}{443.25}{346.889}{p}%
\TextAt{20}{452.236}{360.695}{\Character{1}}%
\TextAt{8}{461.282}{351.011}{of}%
\TextAt{8}{474.281}{351.011}{the}%
\TextAt{8}{494.107}{351.011}{field}%
\TextAt{9}{519.133}{351.011}{\Character{8}}%
\TextAt{21}{527.587}{349.218}{\Character{23}}%
\TextAt{8}{79.226}{333.965}{the}%
\TextAt{12}{79.226}{308.214}{Proposition}%
\TextAt{12}{151.937}{308.214}{(3.8).}%
\TextAt{8}{187.253}{308.214}{For}%
\TextAt{13}{208.222}{308.214}{a}%
\NativeMathAt{217.688}{308.214}{11.955}{\in}%
\TextAt{13}{228.979}{308.214}{U}%
\TextAt{17}{238.178}{314.324}{\Character{40}2}%
\TextAt{21}{245.706}{314.324}{p}%
\TextAt{22}{249.968}{317.137}{\Character{23}}%
\TextAt{23}{254.092}{317.137}{\Character{0}}%
\TextAt{24}{259.848}{317.137}{1}%
\TextAt{17}{263.999}{314.324}{\Character{41}}%
\TextAt{17}{236.929}{304.973}{\Character{8}}%
\TextAt{22}{243.045}{303.977}{\Character{23}}%
\TextAt{8}{271.688}{308.214}{and}%
\TextAt{13}{294.45}{308.214}{b}%
\NativeMathAt{302.748}{308.214}{11.955}{\in}%
\TextAt{9}{314.039}{308.214}{\Character{8}}%
\TextAt{3}{322.494}{312.552}{\Character{3}}%
\TextAt{21}{322.494}{305.259}{\Character{23}}%
\TextAt{8}{331.573}{308.214}{one}%
\TextAt{8}{353.033}{308.214}{has}%
\TextAt{20}{238.283}{291.461}{\Character{18}}%
\TextAt{13}{248.279}{282.691}{a\Character{59}}%
\TextAt{13}{259.648}{282.691}{b}%
\Vector{q 0 G 0 g 0.478 w 176.279 -496.408 m 192.645 -496.408 l S Q}%
\TextAt{14}{253.47}{266.403}{p}%
\TextAt{20}{265.84}{291.461}{\Character{19}}%
\TextAt{9}{277.962}{274.604}{\Character{61}}%
\TextAt{13}{290.387}{274.604}{\Character{16}}%
\TextAt{21}{296.403}{279.54}{S}%
\TextAt{17}{301.998}{279.54}{\Character{40}}%
\TextAt{21}{305.292}{279.54}{\Character{16}}%
\TextAt{7}{310.972}{279.54}{log}%
\TextAt{21}{323.199}{279.54}{a}%
\TextAt{21}{329.105}{279.54}{D}%
\TextAt{7}{337.702}{279.54}{log}%
\TextAt{21}{349.929}{279.54}{b}%
\TextAt{17}{353.551}{279.54}{\Character{41}\Character{47}}%
\TextAt{21}{361.079}{279.54}{p}%
\TextAt{22}{365.342}{282.353}{\Character{23}}%
\TextAt{13}{370.461}{274.604}{\Character{59}}%
\TextAt{8}{517.802}{274.604}{(1)}%
\ErrataAt{79.19}{242.792}{455.81}{where $D\log b=f^{\prime}(\pi)/f(\pi)$ for an arbitrary representation $f(\pi)=b$ with}%
\ErrataAt{79.19}{225.746}{455.81}{$f(X)\in\mathbb Z_p((X))$, a Laurent series with coefficients in $\mathbb Z_p$.}%
\TextAt{8}{97.159}{205.711}{For}%
\TextAt{13}{118.128}{205.711}{a}%
\NativeMathAt{127.594}{205.711}{11.955}{\in}%
\TextAt{13}{138.885}{205.711}{U}%
\TextAt{17}{148.084}{211.821}{\Character{40}1\Character{41}}%
\TextAt{17}{146.835}{202.47}{\Character{8}}%
\TextAt{22}{152.951}{201.474}{\Character{23}}%
\TextAt{8}{163.3}{205.711}{one}%
\TextAt{8}{184.76}{205.711}{has}%
\TextAt{8}{205.634}{205.711}{furthermore}%
\TextAt{8}{270.347}{205.711}{the}%
\TextAt{8}{290.503}{205.711}{two}%
\TextAt{8}{312.62}{205.711}{supplementary}%
\TextAt{8}{391.381}{205.711}{theorems}%
\TextAt{20}{253.531}{188.958}{\Character{18}}%
\TextAt{13}{263.527}{180.188}{\Character{16}}%
\TextAt{13}{269.529}{180.188}{\Character{59}}%
\TextAt{13}{274.765}{180.188}{a}%
\Vector{q 0 G 0 g 0.478 w 191.527 -598.911 m 208.932 -598.911 l S Q}%
\TextAt{14}{269.238}{163.9}{p}%
\TextAt{20}{282.127}{188.958}{\Character{19}}%
\TextAt{9}{294.248}{172.101}{\Character{61}}%
\TextAt{13}{306.674}{172.101}{\Character{16}}%
\TextAt{21}{312.689}{177.037}{S}%
\TextAt{17}{318.285}{177.037}{\Character{40}}%
\TextAt{7}{321.578}{177.037}{log}%
\TextAt{21}{333.805}{177.037}{a}%
\TextAt{17}{338.303}{177.037}{\Character{41}\Character{47}}%
\TextAt{21}{345.831}{177.037}{p}%
\TextAt{22}{350.093}{179.85}{\Character{23}}%
\TextAt{13}{355.213}{172.101}{\Character{59}}%
\TextAt{8}{517.802}{172.101}{(2)}%
\TextAt{20}{242.179}{148.227}{\Character{18}}%
\TextAt{13}{252.175}{139.458}{a\Character{59}}%
\TextAt{13}{263.544}{139.458}{\Character{25}}%
\Vector{q 0 G 0 g 0.478 w 180.175 -639.641 m 198.634 -639.641 l S Q}%
\TextAt{14}{258.413}{123.169}{p}%
\TextAt{20}{271.829}{148.227}{\Character{19}}%
\TextAt{9}{283.95}{131.37}{\Character{61}}%
\TextAt{13}{296.376}{131.37}{\Character{16}}%
\TextAt{21}{302.391}{136.306}{S}%
\TextAt{17}{307.987}{136.306}{\Character{40}\Character{40}}%
\TextAt{21}{314.574}{136.306}{\Character{16}}%
\TextAt{17}{318.842}{136.306}{\Character{47}}%
\TextAt{21}{323.077}{136.306}{\Character{25}}%
\TextAt{17}{328.225}{136.306}{\Character{41}}%
\TextAt{7}{332.93}{136.306}{log}%
\TextAt{21}{345.157}{136.306}{a}%
\TextAt{17}{349.655}{136.306}{\Character{41}\Character{47}}%
\TextAt{21}{357.182}{136.306}{p}%
\TextAt{22}{361.445}{139.119}{\Character{23}}%
\TextAt{13}{366.565}{131.37}{\Character{58}}%
\TextAt{8}{517.802}{131.37}{(3)}%
\TextAt{8}{297.211}{43.033}{311}%
\Destination{522.36}{291.645}{nav.AMS.59}%
\Destination{522.36}{189.145}{nav.AMS.60}%
\Destination{522.36}{148.415}{nav.AMS.61}%
\Destination{78.23}{809.815}{nav.page.311}%
\Destination{79.36}{741.555}{nav.print.339}%
\Destination{79.36}{661.055}{nav.number.5.3.7}%
\Destination{79.36}{322.255}{nav.number.5.3.8}%
\Hotlink{343.38}{510.255}{19.22}{11.047}{nav.cite.N020}%
\Hotlink{409.743}{510.255}{19.22}{11.047}{nav.cite.N069}%
\Hotlink{452.821}{510.255}{25.078}{11.047}{nav.cite.N135}%
\Hotlink{251.921}{459.155}{19.22}{11.047}{nav.cite.N069}%
\Hotlink{131.735}{422.055}{19.22}{11.047}{nav.cite.N080}%
\Hotlink{279.944}{422.055}{13.362}{11.047}{nav.cite.N009}%
\Destination{79.226}{259.792}{errata.118}%
\Destination{79.226}{289.604}{nav.local.5.3.1}%
\Destination{79.226}{187.101}{nav.local.5.3.2}%
\Destination{79.226}{146.37}{nav.local.5.3.3}%
}

% Page 324 (body page 312)
\Page{%
\TextAt{2}{79.226}{796.943}{V.}%
\TextAt{2}{94.073}{796.943}{Local}%
\TextAt{2}{123.615}{796.943}{Class}%
\TextAt{2}{152.208}{796.943}{Field}%
\TextAt{2}{179.939}{796.943}{Theory}%
\Vector{q 0 G 0 g 0.299 w 7.226 18.729 m 460.769 18.729 l S Q}%
\TextAt{8}{97.159}{761.124}{The}%
\TextAt{8}{122.852}{761.124}{supplementary}%
\TextAt{8}{203.251}{761.124}{theorems}%
\TextAt{8}{255.03}{761.124}{(2)}%
\TextAt{8}{275.534}{761.124}{and}%
\TextAt{8}{299.936}{761.124}{(3)}%
\TextAt{8}{320.44}{761.124}{go}%
\TextAt{8}{337.692}{761.124}{back}%
\TextAt{8}{366.649}{761.124}{to}%
\TextAt{26}{382.598}{761.124}{Artin}%
\TextAt{8}{422.266}{761.124}{and}%
\TextAt{26}{446.668}{761.124}{Hasse}%
\TextAt{8}{486.886}{761.124}{[9].}%
\TextAt{8}{512.613}{761.124}{The}%
\TextAt{8}{79.226}{744.078}{formula}%
\TextAt{8}{122.479}{744.078}{(1)}%
\TextAt{8}{141.666}{744.078}{was}%
\TextAt{8}{164.487}{744.078}{proved}%
\TextAt{8}{202.862}{744.078}{independently}%
\TextAt{8}{278.944}{744.078}{by}%
\TextAt{26}{295.525}{744.078}{Artin}%
\TextAt{8}{333.876}{744.078}{[10]}%
\TextAt{8}{356.314}{744.078}{and}%
\TextAt{26}{379.399}{744.078}{Hasse}%
\TextAt{8}{418.3}{744.078}{[61]}%
\TextAt{8}{440.739}{744.078}{in}%
\TextAt{8}{454.713}{744.078}{the}%
\TextAt{8}{475.191}{744.078}{case}%
\TextAt{13}{500.284}{744.078}{\Character{23}}%
\TextAt{9}{510.691}{744.078}{\Character{61}}%
\TextAt{9}{523.65}{744.078}{1}%
\TextAt{8}{529.518}{744.078}{,}%
\TextAt{8}{79.226}{727.032}{and}%
\TextAt{8}{102.314}{727.032}{by}%
\TextAt{26}{118.898}{727.032}{Iwasawa}%
\TextAt{8}{172.233}{727.032}{[80]}%
\TextAt{8}{194.675}{727.032}{in}%
\TextAt{8}{208.653}{727.032}{general.}%
\TextAt{8}{254.504}{727.032}{In}%
\TextAt{8}{269.451}{727.032}{the}%
\TextAt{8}{289.932}{727.032}{case}%
\TextAt{13}{315.029}{727.032}{\Character{23}}%
\TextAt{9}{325.442}{727.032}{\Character{61}}%
\TextAt{9}{338.413}{727.032}{1}%
\TextAt{8}{344.274}{727.032}{,}%
\TextAt{8}{351.83}{727.032}{for}%
\TextAt{8}{370.04}{727.032}{instance,}%
\TextAt{8}{419.284}{727.032}{one}%
\TextAt{8}{441.068}{727.032}{can}%
\TextAt{8}{462.853}{727.032}{indeed}%
\TextAt{8}{500.239}{727.032}{obtain}%
\ErrataAt{79.19}{709.986}{455.81}{the formulae from Brückner\textquoteright{}s theorem \hyperlink{nav.number.5.3.7}{(3.7)}. Since, for integers $i\ge1$,}%
\TextAt{9}{152.752}{683.218}{1}%
\Vector{q 0 G 0 g 0.478 w 80.741 -95.881 m 86.616 -95.881 l S Q}%
\TextAt{13}{152.741}{666.93}{p}%
\TextAt{13}{159.812}{675.131}{S}%
\TextAt{9}{167.707}{675.131}{\Character{40}}%
\TextAt{13}{172.259}{675.131}{\Character{16}}%
\TextAt{13}{178.261}{675.131}{\Character{25}}%
\TextAt{21}{185.344}{680.067}{i}%
\TextAt{9}{188.725}{675.131}{\Character{41}}%
\NativeMathAt{196.599}{675.131}{11.955}{\equiv}%
\TextAt{20}{209.218}{699.639}{\Character{56}}%
\TextAt{20}{209.218}{688.879}{\Character{60}}%
\TextAt{20}{209.218}{667.36}{\Character{58}}%
\TextAt{9}{219.845}{684.256}{1}%
\TextAt{8}{229.019}{684.256}{mod}%
\TextAt{13}{254.456}{684.256}{p\Character{59}}%
\TextAt{13}{275.521}{684.256}{i}%
\TextAt{9}{282.853}{684.256}{\Character{61}}%
\TextAt{13}{295.278}{684.256}{p}%
\TextAt{16}{303.81}{684.256}{\Character{0}}%
\TextAt{9}{315.765}{684.256}{1}%
\TextAt{13}{321.618}{684.256}{\Character{59}}%
\TextAt{9}{219.845}{663.801}{0}%
\TextAt{8}{229.019}{663.801}{mod}%
\TextAt{13}{254.456}{663.801}{p\Character{59}}%
\TextAt{13}{275.521}{663.801}{i}%
\NativeMathAt{282.853}{663.801}{11.955}{\neq}%
\TextAt{13}{295.278}{663.801}{p}%
\TextAt{16}{303.81}{663.801}{\Character{0}}%
\TextAt{9}{315.765}{663.801}{1}%
\TextAt{13}{321.618}{663.801}{\Character{59}}%
\TextAt{8}{340.013}{675.131}{and}%
\TextAt{8}{372.826}{675.131}{log}%
\TextAt{13}{389.786}{675.13}{a}%
\NativeMathAt{399.252}{675.131}{11.955}{\equiv}%
\TextAt{9}{411.871}{675.13}{0}%
\TextAt{8}{421.045}{675.131}{mod}%
\TextAt{14}{446.483}{675.13}{p}%
\TextAt{17}{452.466}{680.067}{2}%
\TextAt{13}{457.199}{675.13}{\Character{59}}%
\TextAt{8}{79.226}{640.287}{one}%
\TextAt{8}{101.765}{640.287}{may}%
\TextAt{8}{128.214}{640.287}{also}%
\TextAt{8}{152.773}{640.287}{interpret}%
\TextAt{8}{202.307}{640.287}{the}%
\TextAt{13}{223.543}{640.287}{\Character{16}}%
\TextAt{8}{229.559}{640.287}{-exponent}%
\TextAt{8}{284.939}{640.287}{in}%
\TextAt{8}{299.671}{640.287}{the}%
\TextAt{8}{320.907}{640.287}{formulae}%
\TextAt{8}{370.118}{640.287}{(1)–(3)}%
\TextAt{8}{410.889}{640.287}{as}%
\TextAt{8}{426.339}{640.287}{the}%
\TextAt{9}{447.574}{640.287}{\Character{40}}%
\TextAt{13}{452.127}{640.287}{p}%
\TextAt{16}{461.395}{640.287}{\Character{0}}%
\TextAt{9}{474.086}{640.287}{1\Character{41}}%
\TextAt{8}{484.491}{640.287}{-st}%
\TextAt{8}{502.535}{640.287}{coeffi-}%
\TextAt{8}{79.226}{623.241}{cient}%
\TextAt{8}{108.592}{623.241}{of}%
\TextAt{8}{123.002}{623.241}{a}%
\TextAt{13}{133.838}{623.241}{\Character{25}}%
\TextAt{8}{140.907}{623.241}{-adic}%
\TextAt{8}{170.596}{623.241}{expansion}%
\TextAt{8}{226.048}{623.241}{of}%
\TextAt{8}{240.458}{623.241}{log}%
\TextAt{13}{257.419}{623.241}{a}%
\TextAt{13}{265.548}{623.241}{D}%
\TextAt{8}{277.593}{623.241}{log}%
\TextAt{13}{294.553}{623.241}{b}%
\TextAt{8}{299.53}{623.241}{.}%
\TextAt{8}{311.219}{623.241}{In}%
\TextAt{8}{326.92}{623.241}{this}%
\TextAt{8}{350.823}{623.241}{way}%
\TextAt{8}{375.646}{623.241}{it}%
\TextAt{8}{388.43}{623.241}{appears}%
\TextAt{8}{432.824}{623.241}{as}%
\TextAt{8}{448.274}{623.241}{a}%
\TextAt{8}{459.11}{623.241}{formal}%
\TextAt{8}{496.94}{623.241}{residue}%
\TextAt{8}{79.226}{606.195}{Res}%
\TextAt{21}{97.661}{604.402}{\Character{25}}%
\TextAt{17}{109.123}{610.902}{1}%
\Vector{q 0 G 0 g 0.478 w 34.496 -164.816 m 43.983 -164.816 l S Q}%
\TextAt{21}{106.496}{602.073}{\Character{25}}%
\TextAt{22}{111.645}{604.342}{p}%
\TextAt{8}{119.171}{606.195}{log}%
\TextAt{13}{136.132}{606.195}{a}%
\TextAt{13}{144.262}{606.195}{D}%
\TextAt{8}{156.306}{606.195}{log}%
\TextAt{13}{173.267}{606.195}{b}%
\TextAt{8}{178.244}{606.195}{.}%
\TextAt{8}{186.852}{606.195}{As}%
\TextAt{8}{204.193}{606.195}{to}%
\TextAt{8}{218.557}{606.195}{the}%
\TextAt{8}{238.767}{606.195}{supplementary}%
\TextAt{8}{317.581}{606.195}{theorems,}%
\TextAt{8}{371.04}{606.195}{one}%
\TextAt{8}{392.553}{606.195}{has}%
\TextAt{8}{413.48}{606.195}{to}%
\TextAt{8}{427.844}{606.195}{define}%
\TextAt{8}{461.706}{606.195}{also}%
\TextAt{13}{485.24}{606.195}{D}%
\TextAt{8}{497.277}{606.195}{log}%
\TextAt{13}{514.237}{606.195}{\Character{16}}%
\TextAt{9}{523.665}{606.195}{\Character{61}}%
\TextAt{16}{79.226}{589.149}{\Character{0}}%
\TextAt{13}{88.525}{589.149}{\Character{16}}%
\TextAt{3}{94.54}{593.488}{\Character{0}}%
\TextAt{17}{101.127}{593.488}{1}%
\TextAt{8}{105.859}{589.149}{,}%
\TextAt{13}{113.008}{589.149}{D}%
\TextAt{8}{125.045}{589.149}{log}%
\TextAt{13}{142.006}{589.149}{\Character{25}}%
\TextAt{9}{152.396}{589.149}{\Character{61}}%
\TextAt{13}{164.821}{589.149}{\Character{25}}%
\TextAt{3}{171.89}{593.488}{\Character{0}}%
\TextAt{17}{178.477}{593.488}{1}%
\TextAt{8}{183.209}{589.149}{.}%
\TextAt{12}{79.226}{564.738}{Exercise}%
\TextAt{12}{132.427}{564.738}{V.3.1.}%
\TextAt{8}{171.173}{564.738}{For}%
\TextAt{13}{192.143}{564.738}{n}%
\TextAt{9}{202.451}{564.738}{\Character{61}}%
\TextAt{9}{214.861}{564.738}{2}%
\TextAt{8}{224.627}{564.738}{the}%
\TextAt{8}{244.783}{564.738}{Hilbert}%
\TextAt{8}{285.096}{564.738}{symbol}%
\TextAt{8}{325.158}{564.738}{has}%
\TextAt{8}{346.032}{564.738}{the}%
\TextAt{8}{366.188}{564.738}{following}%
\TextAt{8}{415.623}{564.738}{concrete}%
\TextAt{8}{461.794}{564.738}{meaning:}%
\ErrataAt{116}{537.623}{410}{$\displaystyle\left(\frac{a,b}{2}\right)=1\quad\Longleftrightarrow\quad ax^2+by^2-z^2=0\text{ has a nontrivial solution in }K.$}%
\TextAt{12}{79.226}{498.195}{Exercise}%
\TextAt{12}{132.427}{498.195}{V.3.2.}%
\TextAt{8}{171.173}{498.195}{Deduce}%
\TextAt{8}{212.622}{498.195}{proposition}%
\TextAt{8}{274.155}{498.195}{(3.8)}%
\TextAt{8}{302.13}{498.195}{from}%
\TextAt{8}{329.771}{498.195}{theorem}%
\TextAt{8}{375.296}{498.195}{(3.7).}%
\TextAt{12}{79.226}{469.408}{Exercise}%
\TextAt{12}{132.081}{469.408}{V.3.3.}%
\TextAt{8}{170.527}{469.408}{Let}%
\TextAt{13}{191.196}{469.408}{K}%
\TextAt{8}{205.61}{469.408}{be}%
\TextAt{8}{221.234}{469.408}{a}%
\TextAt{8}{230.688}{469.408}{local}%
\TextAt{8}{258.028}{469.408}{field}%
\TextAt{8}{283.084}{469.408}{of}%
\TextAt{8}{296.114}{469.408}{characteristic}%
\TextAt{13}{367.747}{469.408}{p}%
\TextAt{8}{373.622}{469.408}{,}%
\TextAt{8}{380.531}{469.408}{let}%
\TextAt{9}{400.268}{472.43}{\Character{22}}%
\TextAt{13}{397.135}{469.408}{K}%
\TextAt{8}{411.549}{469.408}{be}%
\TextAt{8}{427.173}{469.408}{its}%
\TextAt{8}{443.191}{469.408}{separable}%
\TextAt{8}{494.334}{469.408}{closure,}%
\TextAt{8}{79.226}{452.362}{and}%
\TextAt{8}{102.255}{452.362}{let}%
\TextAt{13}{119.426}{452.362}{W}%
\TextAt{21}{130.477}{450.569}{n}%
\TextAt{9}{136.113}{452.362}{\Character{40}}%
\TextAt{9}{143.798}{455.384}{\Character{22}}%
\TextAt{13}{140.666}{452.362}{K}%
\TextAt{9}{151.483}{452.362}{\Character{41}}%
\TextAt{8}{160.199}{452.362}{be}%
\TextAt{8}{176.39}{452.362}{the}%
\TextAt{8}{196.812}{452.362}{ring}%
\TextAt{8}{221.144}{452.362}{of}%
\TextAt{8}{234.74}{452.362}{Witt}%
\TextAt{8}{263.292}{452.362}{vectors}%
\TextAt{8}{303.297}{452.362}{of}%
\TextAt{8}{316.894}{452.362}{length}%
\TextAt{13}{352.93}{452.362}{n}%
\TextAt{8}{359.917}{452.362}{,}%
\TextAt{8}{367.399}{452.362}{with}%
\TextAt{8}{394.325}{452.362}{the}%
\TextAt{8}{414.748}{452.362}{operator}%
\TextAt{13}{462.177}{452.362}{\Character{125}}%
\TextAt{9}{473.404}{452.362}{\Character{58}}%
\TextAt{13}{480.43}{452.362}{W}%
\TextAt{21}{491.481}{450.569}{n}%
\TextAt{9}{497.118}{452.362}{\Character{40}}%
\TextAt{9}{504.803}{455.384}{\Character{22}}%
\TextAt{13}{501.67}{452.362}{K}%
\TextAt{9}{512.488}{452.362}{\Character{41}}%
\NativeMathAt{520.814}{452.362}{11.955}{\rightarrow}%
\TextAt{13}{79.226}{435.316}{W}%
\TextAt{21}{90.278}{433.523}{n}%
\TextAt{9}{95.914}{435.316}{\Character{40}}%
\TextAt{9}{103.599}{438.338}{\Character{22}}%
\TextAt{13}{100.466}{435.316}{K}%
\TextAt{9}{111.284}{435.316}{\Character{41}}%
\TextAt{8}{115.836}{435.316}{,}%
\TextAt{13}{123.143}{435.316}{\Character{125}a}%
\TextAt{9}{140.275}{435.316}{\Character{61}}%
\TextAt{13}{152.915}{435.316}{F}%
\TextAt{13}{162.097}{435.316}{a}%
\TextAt{16}{171.006}{435.316}{\Character{0}}%
\TextAt{13}{183.046}{435.316}{a}%
\TextAt{8}{193.214}{435.316}{(see}%
\TextAt{8}{216.808}{435.316}{chap.}%
\TextAt{8}{249.376}{435.316}{IV,}%
\TextAt{8}{269.647}{435.316}{§3,}%
\TextAt{8}{288.49}{435.316}{exercises}%
\TextAt{8}{336.532}{435.316}{2}%
\TextAt{8}{346.413}{435.316}{and}%
\TextAt{8}{369.301}{435.316}{3).}%
\TextAt{8}{388.539}{435.316}{Show}%
\TextAt{8}{419.557}{435.316}{that}%
\TextAt{8}{445.052}{435.316}{one}%
\TextAt{8}{466.637}{435.316}{has}%
\TextAt{8}{487.636}{435.316}{ker}%
\TextAt{9}{503.573}{435.316}{\Character{40}}%
\TextAt{13}{508.125}{435.316}{\Character{125}}%
\TextAt{9}{515.578}{435.316}{\Character{41}}%
\TextAt{9}{523.648}{435.316}{\Character{61}}%
\TextAt{13}{79.226}{418.27}{W}%
\TextAt{21}{90.278}{416.477}{n}%
\TextAt{9}{95.914}{418.27}{\Character{40}}%
\TextAt{18}{100.466}{418.27}{F}%
\TextAt{21}{107.772}{416.477}{p}%
\TextAt{9}{112.533}{418.27}{\Character{41}}%
\TextAt{8}{117.086}{418.27}{.}%
\TextAt{12}{79.226}{389.483}{Exercise}%
\TextAt{12}{134.464}{389.483}{V.3.4.}%
\TextAt{8}{174.982}{389.483}{Abstract}%
\TextAt{8}{225.267}{389.483}{Kummer}%
\TextAt{8}{275.804}{389.483}{theory}%
\TextAt{8}{314.325}{389.483}{(chap.}%
\TextAt{8}{356.385}{389.483}{IV,}%
\TextAt{8}{378.301}{389.483}{(3.3))}%
\TextAt{8}{412.602}{389.483}{yields}%
\TextAt{8}{447.274}{389.483}{for}%
\TextAt{8}{466.93}{389.483}{the}%
\TextAt{8}{488.858}{389.483}{maximal}%
\ErrataAt{79.22}{372.455}{453.64}{abelian extension $L\mid K$ of exponent $p^n$, a surjective homomorphism}%
\TextAt{13}{177.442}{349.041}{W}%
\TextAt{21}{188.494}{347.248}{n}%
\TextAt{9}{194.13}{349.041}{\Character{40}}%
\TextAt{13}{198.683}{349.041}{K}%
\TextAt{9}{209.5}{349.041}{\Character{41}}%
\NativeArrowAt{217.373}{349.041}{11.955}{19.26}{\longrightarrow}%
\TextAt{8}{239.955}{349.041}{Hom}%
\TextAt{9}{264.344}{349.041}{\Character{40}}%
\TextAt{13}{268.896}{349.041}{G}%
\TextAt{9}{278.13}{349.041}{\Character{40}}%
\TextAt{13}{282.682}{349.041}{L}%
\TextAt{16}{290.647}{349.041}{\Character{106}}%
\TextAt{13}{293.968}{349.041}{K}%
\TextAt{9}{304.785}{349.041}{\Character{41}}%
\TextAt{13}{309.338}{349.041}{\Character{59}}%
\TextAt{13}{314.562}{349.041}{W}%
\TextAt{21}{325.633}{347.248}{n}%
\TextAt{9}{331.27}{349.041}{\Character{40}}%
\TextAt{18}{335.822}{349.041}{F}%
\TextAt{21}{343.128}{347.248}{p}%
\TextAt{9}{347.889}{349.041}{\Character{41}\Character{41}}%
\TextAt{13}{356.994}{349.041}{\Character{59}}%
\TextAt{13}{386.129}{349.041}{x}%
\NativeArrowAt{396.121}{349.041}{11.955}{19.26}{\longmapsto}%
\TextAt{13}{418.703}{349.041}{\Character{31}}%
\TextAt{21}{426.037}{347.248}{x}%
\TextAt{13}{431.302}{349.041}{\Character{59}}%
\TextAt{8}{79.226}{325.645}{where}%
\TextAt{8}{112.655}{325.645}{one}%
\TextAt{8}{133.734}{325.645}{has}%
\TextAt{13}{154.227}{325.645}{\Character{31}}%
\TextAt{21}{161.561}{323.852}{x}%
\TextAt{9}{166.826}{325.645}{\Character{40}}%
\TextAt{13}{171.378}{325.645}{\Character{27}}%
\TextAt{9}{178.46}{325.645}{\Character{41}}%
\TextAt{9}{186.315}{325.645}{\Character{61}}%
\TextAt{13}{198.759}{325.645}{\Character{27}}%
\TextAt{13}{205.825}{325.645}{\Character{24}}%
\TextAt{16}{213.411}{325.645}{\Character{0}}%
\TextAt{13}{224.588}{325.645}{\Character{24}}%
\TextAt{8}{233.795}{325.645}{for}%
\TextAt{8}{251.299}{325.645}{all}%
\TextAt{13}{267.178}{325.645}{\Character{27}}%
\NativeMathAt{277.581}{325.645}{11.955}{\in}%
\TextAt{13}{288.872}{325.645}{G}%
\TextAt{9}{298.106}{325.645}{\Character{40}}%
\TextAt{13}{302.658}{325.645}{L}%
\TextAt{16}{310.623}{325.645}{\Character{106}}%
\TextAt{13}{313.944}{325.645}{K}%
\TextAt{9}{324.761}{325.645}{\Character{41}}%
\TextAt{8}{329.314}{325.645}{,}%
\TextAt{8}{336.158}{325.645}{with}%
\TextAt{8}{362.438}{325.645}{an}%
\TextAt{8}{378.316}{325.645}{arbitrary}%
\TextAt{13}{427.705}{325.645}{\Character{24}}%
\NativeMathAt{436.717}{325.645}{11.955}{\in}%
\TextAt{13}{448.008}{325.645}{W}%
\TextAt{21}{459.059}{323.852}{n}%
\TextAt{9}{464.695}{325.645}{\Character{40}}%
\TextAt{13}{469.248}{325.645}{L}%
\TextAt{9}{477.212}{325.645}{\Character{41}}%
\TextAt{8}{485.282}{325.645}{such}%
\TextAt{8}{511.298}{325.645}{that}%
\TextAt{13}{79.226}{308.599}{\Character{125}\Character{24}}%
\TextAt{9}{95.69}{308.599}{\Character{61}}%
\TextAt{13}{108.115}{308.599}{x}%
\TextAt{8}{114.768}{308.599}{.}%
\TextAt{8}{123.216}{308.599}{Show}%
\TextAt{8}{154.108}{308.599}{that}%
\TextAt{13}{179.477}{308.599}{x}%
\NativeArrowAt{189.45}{308.599}{11.955}{11.955}{\mapsto}%
\TextAt{13}{204.726}{308.599}{\Character{31}}%
\TextAt{21}{212.059}{306.806}{x}%
\TextAt{8}{221.222}{308.599}{has}%
\TextAt{8}{242.096}{308.599}{kernel}%
\TextAt{13}{276.562}{308.599}{\Character{125}W}%
\TextAt{21}{295.066}{306.806}{n}%
\TextAt{9}{300.703}{308.599}{\Character{40}}%
\TextAt{13}{305.255}{308.599}{K}%
\TextAt{9}{316.073}{308.599}{\Character{41}}%
\TextAt{8}{320.625}{308.599}{.}%
\TextAt{12}{79.226}{279.812}{Exercise}%
\TextAt{12}{132.427}{279.812}{V.3.5.}%
\TextAt{8}{171.173}{279.812}{Define,}%
\TextAt{8}{210.673}{279.812}{for}%
\TextAt{13}{228.558}{279.812}{x}%
\NativeMathAt{238.531}{279.812}{11.955}{\in}%
\TextAt{13}{249.822}{279.812}{W}%
\TextAt{21}{260.874}{278.019}{n}%
\TextAt{9}{266.51}{279.812}{\Character{40}}%
\TextAt{13}{271.062}{279.812}{K}%
\TextAt{9}{281.88}{279.812}{\Character{41}}%
\TextAt{8}{290.33}{279.812}{and}%
\TextAt{13}{313.092}{279.812}{a}%
\NativeMathAt{322.558}{279.812}{11.955}{\in}%
\TextAt{13}{333.849}{279.812}{K}%
\TextAt{3}{344.667}{284.151}{\Character{3}}%
\TextAt{8}{349.399}{279.812}{,}%
\TextAt{8}{356.548}{279.812}{the}%
\TextAt{8}{376.705}{279.812}{symbol}%
\TextAt{9}{416.766}{279.812}{\Character{91}}%
\TextAt{13}{420.018}{279.812}{x\Character{59}}%
\TextAt{13}{431.889}{279.812}{a}%
\TextAt{9}{438.059}{279.812}{\Character{41}}%
\NativeMathAt{445.932}{279.812}{11.955}{\in}%
\TextAt{13}{457.223}{279.812}{W}%
\TextAt{21}{468.275}{278.019}{n}%
\TextAt{9}{473.911}{279.812}{\Character{40}}%
\TextAt{18}{478.464}{279.812}{F}%
\TextAt{21}{485.77}{278.019}{p}%
\TextAt{9}{490.53}{279.812}{\Character{41}}%
\TextAt{8}{498.98}{279.812}{by}%
\TextAt{9}{251.426}{256.417}{\Character{91}}%
\TextAt{13}{254.678}{256.417}{x\Character{59}}%
\TextAt{13}{266.55}{256.417}{a}%
\TextAt{9}{272.719}{256.417}{\Character{41}}%
\TextAt{9}{280.574}{256.417}{\Character{61}}%
\TextAt{13}{293.018}{256.417}{\Character{31}}%
\TextAt{21}{300.351}{254.623}{x}%
\TextAt{9}{305.616}{256.417}{\Character{40}\Character{40}}%
\TextAt{13}{314.721}{256.417}{a\Character{59}}%
\TextAt{13}{326.09}{256.417}{L}%
\TextAt{16}{334.075}{256.417}{\Character{106}}%
\TextAt{13}{337.395}{256.417}{K}%
\TextAt{9}{348.213}{256.417}{\Character{41}\Character{41}}%
\TextAt{13}{357.318}{256.417}{\Character{59}}%
\TextAt{8}{79.226}{233.021}{where}%
\TextAt{9}{113.035}{233.021}{\Character{40}}%
\TextAt{13}{121.485}{233.021}{\Character{59}}%
\TextAt{13}{126.709}{233.021}{L}%
\TextAt{16}{134.694}{233.021}{\Character{106}}%
\TextAt{13}{138.015}{233.021}{K}%
\TextAt{9}{148.832}{233.021}{\Character{41}}%
\TextAt{8}{157.282}{233.021}{is}%
\TextAt{8}{169.046}{233.021}{the}%
\TextAt{8}{189.202}{233.021}{norm}%
\TextAt{8}{219.772}{233.021}{residue}%
\TextAt{8}{259.499}{233.021}{symbol.}%
\TextAt{8}{304.111}{233.021}{Show:}%
\TextAt{8}{297.211}{43.033}{312}%
\Destination{78.23}{809.815}{nav.page.312}%
\Destination{79.226}{578.755}{nav.ex.5.3.1}%
\Destination{79.226}{512.155}{nav.ex.5.3.2}%
\Destination{79.226}{483.455}{nav.ex.5.3.3}%
\Destination{79.226}{403.455}{nav.ex.5.3.4}%
\Destination{79.226}{293.855}{nav.ex.5.3.5}%
\Destination{79.36}{776.155}{nav.print.340}%
\Hotlink{486.386}{758.955}{13.362}{11.047}{nav.cite.N009}%
\Hotlink{333.376}{741.855}{19.22}{11.047}{nav.cite.N010}%
\Hotlink{417.8}{741.855}{19.22}{11.047}{nav.cite.N061}%
\Hotlink{171.733}{724.855}{19.22}{11.047}{nav.cite.N080}%
\Hotlink{336.032}{433.155}{6.858}{11.047}{nav.ex.4.3.2}%
\Hotlink{248.876}{433.155}{13.995}{11.047}{nav.chapter.4}%
\Hotlink{274.826}{433.155}{6.858}{11.047}{nav.section.4.3}%
\Hotlink{355.885}{387.255}{13.995}{11.047}{nav.chapter.4}%
\Hotlink{144.02}{127.155}{5.981}{9.572}{nav.ex.5.3.4}%
\Hotlink{273.655}{495.955}{25.078}{11.047}{nav.number.5.3.8}%
\Hotlink{374.796}{495.955}{25.078}{11.047}{nav.number.5.3.7}%
\Hotlink{377.801}{387.255}{25.078}{11.047}{nav.number.4.3.3}%
\Destination{79.226}{386.455}{errata.26}%
\Destination{79.226}{726.986}{errata.119}%
\Destination{79.226}{554.623}{errata.120}%
\Hotlink{357.4}{433.15}{6.9}{11.047}{nav.ex.4.3.3}%
\Hotlink{255.03}{758.924}{18}{11.047}{nav.local.5.3.2}%
\Hotlink{299.936}{758.924}{18}{11.047}{nav.local.5.3.3}%
\Hotlink{122.479}{741.878}{18}{11.047}{nav.local.5.3.1}%
\Hotlink{370.118}{638.087}{18}{11.047}{nav.local.5.3.1}%
\Hotlink{394.25}{638.087}{18}{11.047}{nav.local.5.3.3}%
}

% Page 325 (body page 313)
\Page{%
\TextAt{2}{79.226}{796.943}{V.}%
\TextAt{2}{94.073}{796.943}{Local}%
\TextAt{2}{123.615}{796.943}{Class}%
\TextAt{2}{152.208}{796.943}{Field}%
\TextAt{2}{179.939}{796.943}{Theory}%
\Vector{q 0 G 0 g 0.299 w 7.226 18.729 m 460.769 18.729 l S Q}%
\TextAt{8}{89.579}{761.124}{(i)}%
\TextAt{9}{107.918}{761.124}{\Character{91}}%
\TextAt{13}{111.17}{761.124}{x\Character{59}}%
\TextAt{13}{123.042}{761.124}{a}%
\TextAt{9}{129.211}{761.124}{\Character{41}}%
\TextAt{9}{137.066}{761.124}{\Character{61}}%
\TextAt{9}{149.487}{761.124}{\Character{40}}%
\TextAt{13}{154.062}{761.124}{a\Character{59}}%
\TextAt{13}{165.431}{761.124}{K}%
\TextAt{9}{176.269}{761.124}{\Character{40}}%
\TextAt{13}{180.821}{761.124}{\Character{24}}%
\TextAt{9}{186.512}{761.124}{\Character{41}}%
\TextAt{16}{191.064}{761.124}{\Character{106}}%
\TextAt{13}{194.385}{761.124}{K}%
\TextAt{9}{205.203}{761.124}{\Character{41}}%
\TextAt{13}{209.755}{761.124}{\Character{24}}%
\TextAt{16}{218.102}{761.124}{\Character{0}}%
\TextAt{13}{230.057}{761.124}{\Character{24}}%
\TextAt{8}{235.748}{761.124}{,}%
\TextAt{8}{242.897}{761.124}{if}%
\TextAt{13}{253.621}{761.124}{\Character{24}}%
\NativeMathAt{262.632}{761.124}{11.955}{\in}%
\TextAt{13}{273.923}{761.124}{W}%
\TextAt{21}{284.974}{759.33}{n}%
\TextAt{9}{290.611}{761.124}{\Character{40}}%
\TextAt{9}{298.296}{764.146}{\Character{22}}%
\TextAt{13}{295.163}{761.124}{K}%
\TextAt{9}{305.981}{761.124}{\Character{41}}%
\TextAt{8}{314.431}{761.124}{with}%
\TextAt{13}{341.091}{761.124}{\Character{125}\Character{24}}%
\TextAt{9}{357.554}{761.124}{\Character{61}}%
\TextAt{13}{369.98}{761.124}{x}%
\TextAt{8}{376.632}{761.124}{.}%
\TextAt{8}{86.327}{744.078}{(ii)}%
\TextAt{9}{107.918}{744.078}{\Character{91}}%
\TextAt{13}{111.17}{744.078}{x}%
\TextAt{9}{120.479}{744.078}{\Character{43}}%
\TextAt{13}{132.24}{744.078}{y}%
\TextAt{13}{138.361}{744.078}{\Character{59}}%
\TextAt{13}{143.597}{744.078}{a}%
\TextAt{9}{149.766}{744.078}{\Character{41}}%
\TextAt{9}{157.621}{744.078}{\Character{61}}%
\TextAt{9}{170.042}{744.078}{\Character{91}}%
\TextAt{13}{173.316}{744.078}{x\Character{59}}%
\TextAt{13}{185.188}{744.078}{a}%
\TextAt{9}{191.357}{744.078}{\Character{41}}%
\TextAt{9}{198.554}{744.078}{\Character{43}}%
\TextAt{9}{210.306}{744.078}{\Character{91}}%
\TextAt{13}{213.579}{744.078}{y}%
\TextAt{13}{219.7}{744.078}{\Character{59}}%
\TextAt{13}{224.936}{744.078}{a}%
\TextAt{9}{231.105}{744.078}{\Character{41}}%
\TextAt{8}{235.657}{744.078}{.}%
\TextAt{8}{83.076}{727.032}{(iii)}%
\TextAt{9}{107.918}{727.032}{\Character{91}}%
\TextAt{13}{111.17}{727.032}{x\Character{59}}%
\TextAt{13}{123.042}{727.032}{ab}%
\TextAt{9}{134.188}{727.032}{\Character{41}}%
\TextAt{9}{142.043}{727.032}{\Character{61}}%
\TextAt{9}{154.464}{727.032}{\Character{91}}%
\TextAt{13}{157.739}{727.032}{x\Character{59}}%
\TextAt{13}{169.611}{727.032}{a}%
\TextAt{9}{175.78}{727.032}{\Character{41}}%
\TextAt{9}{182.977}{727.032}{\Character{43}}%
\TextAt{9}{194.729}{727.032}{\Character{91}}%
\TextAt{13}{198.002}{727.032}{x\Character{59}}%
\TextAt{13}{209.874}{727.032}{b}%
\TextAt{9}{214.875}{727.032}{\Character{41}}%
\TextAt{8}{219.428}{727.032}{.}%
\TextAt{8}{83.398}{709.986}{(iv)}%
\TextAt{9}{107.918}{709.986}{\Character{91}}%
\TextAt{13}{111.17}{709.986}{x\Character{59}}%
\TextAt{13}{123.042}{709.986}{a}%
\TextAt{9}{129.211}{709.986}{\Character{41}}%
\TextAt{9}{138.273}{709.986}{\Character{61}}%
\TextAt{9}{151.878}{709.986}{0}%
\NativeArrowAt{174.233}{709.986}{11.955}{21.926}{\Longleftrightarrow}%
\TextAt{13}{212.625}{709.986}{a}%
\NativeMathAt{223.288}{709.986}{11.955}{\in}%
\TextAt{13}{235.777}{709.986}{N}%
\TextAt{21}{245.151}{708.126}{K}%
\TextAt{17}{252.856}{708.126}{\Character{40}}%
\TextAt{21}{256.149}{708.126}{\Character{24}}%
\TextAt{17}{260.182}{708.126}{\Character{41}}%
\TextAt{3}{263.476}{708.126}{\Character{106}}%
\TextAt{21}{265.828}{708.126}{K}%
\TextAt{13}{274.031}{709.986}{K}%
\TextAt{9}{284.848}{709.986}{\Character{40}}%
\TextAt{13}{289.401}{709.986}{\Character{24}}%
\TextAt{9}{295.091}{709.986}{\Character{41}}%
\TextAt{3}{299.643}{714.324}{\Character{3}}%
\TextAt{8}{304.376}{709.986}{,}%
\TextAt{8}{312.404}{709.986}{where}%
\TextAt{13}{346.916}{709.986}{\Character{24}}%
\NativeMathAt{357.125}{709.986}{11.955}{\in}%
\TextAt{13}{369.614}{709.986}{W}%
\TextAt{21}{380.665}{708.193}{n}%
\TextAt{9}{386.302}{709.986}{\Character{40}}%
\TextAt{9}{393.987}{713.008}{\Character{22}}%
\TextAt{13}{390.854}{709.986}{K}%
\TextAt{9}{401.672}{709.986}{\Character{41}}%
\TextAt{8}{410.824}{709.986}{is}%
\TextAt{8}{423.291}{709.986}{an}%
\TextAt{8}{440.253}{709.986}{element}%
\TextAt{8}{484.198}{709.986}{such}%
\TextAt{8}{511.298}{709.986}{that}%
\TextAt{13}{107.918}{692.94}{\Character{125}\Character{24}}%
\TextAt{9}{124.382}{692.94}{\Character{61}}%
\TextAt{13}{136.808}{692.94}{x}%
\TextAt{8}{143.46}{692.94}{.}%
\TextAt{8}{86.65}{675.894}{(v)}%
\TextAt{9}{107.918}{675.894}{\Character{91}}%
\TextAt{13}{111.17}{675.894}{x\Character{59}}%
\TextAt{13}{123.042}{675.894}{a}%
\TextAt{9}{129.211}{675.894}{\Character{41}}%
\TextAt{9}{137.066}{675.894}{\Character{61}}%
\TextAt{9}{149.487}{675.894}{0}%
\TextAt{8}{159.26}{675.894}{for}%
\TextAt{8}{177.145}{675.894}{all}%
\TextAt{13}{193.404}{675.894}{a}%
\NativeMathAt{202.87}{675.894}{11.955}{\in}%
\TextAt{13}{214.161}{675.894}{K}%
\TextAt{3}{224.979}{680.232}{\Character{3}}%
\NativeArrowAt{244.987}{675.894}{11.955}{21.926}{\Longleftrightarrow}%
\TextAt{13}{282.181}{675.894}{x}%
\NativeMathAt{292.154}{675.894}{11.955}{\in}%
\TextAt{13}{303.445}{675.894}{\Character{125}W}%
\TextAt{21}{321.949}{674.101}{n}%
\TextAt{9}{327.585}{675.894}{\Character{40}}%
\TextAt{13}{332.138}{675.894}{K}%
\TextAt{9}{342.955}{675.894}{\Character{41}}%
\TextAt{8}{347.508}{675.894}{.}%
\TextAt{8}{83.398}{658.848}{(vi)}%
\TextAt{9}{107.918}{658.848}{\Character{91}}%
\TextAt{13}{111.17}{658.848}{x\Character{59}}%
\TextAt{13}{123.042}{658.848}{a}%
\TextAt{9}{129.211}{658.848}{\Character{41}}%
\TextAt{9}{137.066}{658.848}{\Character{61}}%
\TextAt{9}{149.487}{658.848}{0}%
\TextAt{8}{159.26}{658.848}{for}%
\TextAt{8}{177.145}{658.848}{all}%
\TextAt{13}{193.404}{658.848}{x}%
\NativeMathAt{203.377}{658.848}{11.955}{\in}%
\TextAt{13}{214.668}{658.848}{W}%
\TextAt{21}{225.72}{657.055}{n}%
\TextAt{9}{231.356}{658.848}{\Character{40}}%
\TextAt{13}{235.908}{658.848}{K}%
\TextAt{9}{246.726}{658.848}{\Character{41}}%
\NativeArrowAt{266.554}{658.848}{11.955}{21.926}{\Longleftrightarrow}%
\TextAt{13}{303.748}{658.848}{a}%
\NativeMathAt{313.214}{658.848}{11.955}{\in}%
\TextAt{13}{324.505}{658.848}{K}%
\TextAt{3}{335.323}{663.186}{\Character{3}}%
\TextAt{21}{339.557}{663.186}{p}%
\TextAt{22}{343.82}{665.999}{n}%
\TextAt{8}{349.417}{658.848}{.}%
\TextAt{12}{79.226}{626.858}{Exercise}%
\TextAt{12}{132.801}{626.858}{V.3.6.}%
\TextAt{8}{171.874}{626.858}{Let}%
\TextAt{13}{193.169}{626.858}{\Character{20}}%
\TextAt{8}{204.13}{626.858}{be}%
\TextAt{8}{220.38}{626.858}{the}%
\TextAt{8}{240.862}{626.858}{residue}%
\TextAt{8}{280.915}{626.858}{class}%
\TextAt{8}{308.678}{626.858}{field}%
\TextAt{8}{334.36}{626.858}{of}%
\TextAt{13}{348.016}{626.858}{K}%
\TextAt{8}{363.057}{626.858}{and}%
\TextAt{13}{386.145}{626.858}{\Character{25}}%
\TextAt{8}{397.437}{626.858}{a}%
\TextAt{8}{407.518}{626.858}{prime}%
\TextAt{8}{441.008}{626.858}{element}%
\TextAt{8}{484.575}{626.858}{such}%
\TextAt{8}{511.298}{626.858}{that}%
\TextAt{13}{79.226}{609.812}{K}%
\TextAt{9}{93.365}{609.812}{\Character{61}}%
\TextAt{13}{105.79}{609.812}{\Character{20}}%
\TextAt{9}{112.529}{609.812}{\Character{40}\Character{40}}%
\TextAt{13}{121.633}{609.812}{\Character{25}}%
\TextAt{9}{128.703}{609.812}{\Character{41}\Character{41}}%
\TextAt{8}{137.807}{609.812}{.}%
\TextAt{8}{146.256}{609.812}{Let}%
\TextAt{13}{229.763}{592.766}{d}%
\TextAt{9}{239.167}{592.766}{\Character{58}}%
\TextAt{13}{245.739}{592.766}{K}%
\NativeArrowAt{259.878}{592.766}{11.955}{19.26}{\longrightarrow}%
\TextAt{9}{282.46}{592.766}{\Character{10}}%
\TextAt{17}{290.914}{597.702}{1}%
\TextAt{21}{290.914}{589.81}{K}%
\TextAt{3}{298.619}{589.81}{\Character{106}}%
\TextAt{21}{300.971}{589.81}{\Character{20}}%
\TextAt{13}{306.349}{592.766}{\Character{59}}%
\TextAt{13}{335.484}{592.766}{f}%
\NativeArrowAt{345.871}{592.766}{11.955}{19.26}{\longmapsto}%
\TextAt{13}{368.453}{592.766}{d}%
\TextAt{13}{372.577}{592.766}{f}%
\TextAt{13}{378.962}{592.766}{\Character{59}}%
\TextAt{8}{79.226}{569.244}{be}%
\TextAt{8}{95.612}{569.244}{the}%
\TextAt{8}{116.23}{569.244}{canonical}%
\TextAt{8}{168.074}{569.244}{map}%
\TextAt{8}{194.55}{569.244}{to}%
\TextAt{8}{209.322}{569.244}{the}%
\TextAt{8}{229.94}{569.244}{differential}%
\TextAt{8}{288.934}{569.244}{module}%
\TextAt{8}{330.688}{569.244}{of}%
\TextAt{13}{344.48}{569.244}{K}%
\TextAt{16}{355.297}{569.244}{\Character{106}}%
\TextAt{13}{358.618}{569.244}{\Character{20}}%
\TextAt{8}{369.716}{569.244}{(see}%
\TextAt{8}{393.645}{569.244}{chap.}%
\TextAt{8}{427.22}{569.244}{III,}%
\TextAt{8}{448.137}{569.244}{§2,}%
\TextAt{8}{467.4}{569.244}{p.}%
\TextAt{8}{483.736}{569.244}{200).}%
\TextAt{8}{515.697}{569.244}{For}%
\TextAt{8}{79.226}{552.198}{every}%
\TextAt{13}{110.118}{552.198}{f}%
\NativeMathAt{120.486}{552.198}{11.955}{\in}%
\TextAt{13}{131.777}{552.198}{K}%
\TextAt{8}{146.492}{552.198}{one}%
\TextAt{8}{167.951}{552.198}{has}%
\TextAt{13}{277.629}{535.152}{d}%
\TextAt{13}{281.754}{535.152}{f}%
\TextAt{9}{292.128}{535.152}{\Character{61}}%
\TextAt{13}{304.553}{535.152}{f}%
\TextAt{3}{311.6}{540.089}{\Character{48}}%
\TextAt{21}{310.323}{532.197}{\Character{25}}%
\TextAt{13}{317.963}{535.152}{d\Character{25}}%
\TextAt{13}{331.09}{535.152}{\Character{59}}%
\TextAt{8}{79.226}{511.631}{where}%
\TextAt{13}{114.086}{511.631}{f}%
\TextAt{3}{121.133}{515.969}{\Character{48}}%
\TextAt{21}{119.856}{508.675}{\Character{25}}%
\TextAt{8}{130.451}{511.631}{is}%
\TextAt{8}{143.266}{511.631}{the}%
\TextAt{8}{164.473}{511.631}{formal}%
\TextAt{8}{202.274}{511.631}{derivative}%
\TextAt{8}{256.992}{511.631}{of}%
\TextAt{13}{271.372}{511.631}{f}%
\TextAt{8}{283.367}{511.631}{in}%
\TextAt{8}{298.071}{511.631}{the}%
\TextAt{8}{319.278}{511.631}{expansion}%
\TextAt{8}{374.701}{511.631}{according}%
\TextAt{8}{428.438}{511.631}{to}%
\TextAt{8}{443.799}{511.631}{powers}%
\TextAt{8}{483.609}{511.631}{of}%
\TextAt{13}{497.989}{511.631}{\Character{25}}%
\TextAt{8}{510.007}{511.631}{with}%
\TextAt{8}{79.226}{494.585}{coefficients}%
\TextAt{8}{138.464}{494.585}{in}%
\TextAt{13}{152.117}{494.585}{\Character{20}}%
\TextAt{8}{158.855}{494.585}{.}%
\TextAt{8}{167.304}{494.585}{Show}%
\TextAt{8}{198.196}{494.585}{that}%
\TextAt{8}{223.565}{494.585}{for}%
\TextAt{13}{241.45}{494.585}{\Character{33}}%
\TextAt{9}{252.489}{494.585}{\Character{61}}%
\TextAt{9}{264.899}{494.585}{\Character{40}}%
\TextAt{20}{269.467}{503.551}{\Character{80}}%
\TextAt{21}{282.086}{491.098}{i\Character{62}}%
\TextAt{3}{291.556}{491.098}{\Character{0}}%
\NativeMathAt{298.139}{491.098}{7.97}{\infty}%
\TextAt{13}{309.101}{494.585}{a}%
\TextAt{21}{315.246}{492.791}{i}%
\TextAt{13}{318.627}{494.585}{\Character{25}}%
\TextAt{21}{325.697}{498.923}{i}%
\TextAt{9}{329.078}{494.585}{\Character{41}}%
\TextAt{13}{333.63}{494.585}{d\Character{25}}%
\TextAt{8}{346.782}{494.585}{,}%
\TextAt{8}{353.931}{494.585}{the}%
\TextAt{8}{374.088}{494.585}{residue}%
\TextAt{8}{273.694}{465.584}{Res}%
\TextAt{13}{294.122}{465.584}{\Character{33}}%
\TextAt{9}{305.161}{465.584}{\Character{58}\Character{61}}%
\TextAt{13}{320.838}{465.584}{a}%
\TextAt{3}{326.983}{463.79}{\Character{0}}%
\TextAt{17}{333.569}{463.79}{1}%
\TextAt{8}{79.226}{436.582}{does}%
\TextAt{8}{105.623}{436.582}{not}%
\TextAt{8}{126.437}{436.582}{depend}%
\TextAt{8}{167.073}{436.582}{on}%
\TextAt{8}{183.332}{436.582}{the}%
\TextAt{8}{203.488}{436.582}{choice}%
\TextAt{8}{238.278}{436.582}{of}%
\TextAt{8}{251.608}{436.582}{the}%
\TextAt{8}{271.764}{436.582}{prime}%
\TextAt{8}{304.928}{436.582}{element}%
\TextAt{13}{348.17}{436.582}{\Character{25}}%
\TextAt{8}{355.239}{436.582}{.}%
\TextAt{12}{79.226}{404.592}{Exercise}%
\TextAt{12}{132.427}{404.592}{V.3.7.}%
\TextAt{8}{171.173}{404.592}{Show}%
\TextAt{8}{202.066}{404.592}{that}%
\TextAt{8}{227.434}{404.592}{in}%
\TextAt{8}{241.087}{404.592}{the}%
\TextAt{8}{261.244}{404.592}{case}%
\TextAt{13}{286.015}{404.592}{n}%
\TextAt{9}{296.323}{404.592}{\Character{61}}%
\TextAt{9}{308.733}{404.592}{1}%
\TextAt{8}{318.499}{404.592}{the}%
\TextAt{8}{338.655}{404.592}{symbol}%
\TextAt{9}{378.717}{404.592}{\Character{91}}%
\TextAt{13}{381.969}{404.592}{x\Character{59}}%
\TextAt{13}{393.841}{404.592}{a}%
\TextAt{9}{400.01}{404.592}{\Character{41}}%
\TextAt{8}{408.46}{404.592}{is}%
\TextAt{8}{420.224}{404.592}{given}%
\TextAt{8}{450.793}{404.592}{by}%
\TextAt{9}{235.965}{371.317}{\Character{91}}%
\TextAt{13}{239.216}{371.317}{x\Character{59}}%
\TextAt{13}{251.088}{371.317}{a}%
\TextAt{9}{257.257}{371.317}{\Character{41}}%
\TextAt{9}{265.112}{371.317}{\Character{61}}%
\TextAt{8}{277.556}{371.317}{Tr}%
\TextAt{21}{290.563}{369.457}{\Character{20}}%
\TextAt{3}{295.443}{369.457}{\Character{106}}%
\TextAt{25}{297.796}{369.457}{F}%
\TextAt{22}{302.666}{368.461}{p}%
\TextAt{8}{309.495}{371.317}{Res}%
\TextAt{20}{329.923}{388.174}{\Character{18}}%
\TextAt{13}{338.723}{371.317}{x}%
\TextAt{13}{348.563}{379.405}{da}%
\Vector{q 0 G 0 g 0.478 w 276.563 -399.694 m 288.791 -399.694 l S Q}%
\TextAt{13}{351.605}{363.116}{a}%
\TextAt{20}{361.986}{388.174}{\Character{19}}%
\TextAt{13}{372.779}{371.317}{\Character{58}}%
\TextAt{12}{79.226}{326.087}{Remark.}%
\TextAt{8}{133.383}{326.087}{Such}%
\TextAt{8}{161.669}{326.087}{a}%
\TextAt{8}{171.424}{326.087}{formula}%
\TextAt{8}{214.355}{326.087}{can}%
\TextAt{8}{235.815}{326.087}{also}%
\TextAt{8}{259.295}{326.087}{be}%
\TextAt{8}{275.219}{326.087}{given}%
\TextAt{8}{305.789}{326.087}{for}%
\TextAt{13}{323.674}{326.087}{n}%
\NativeMathAt{333.982}{326.087}{11.955}{\geq}%
\TextAt{9}{346.601}{326.087}{1}%
\TextAt{8}{356.352}{326.087}{(}%
\TextAt{26}{360.907}{326.087}{P.}%
\TextAt{26}{378.11}{326.087}{Kölcze}%
\TextAt{8}{425.213}{326.087}{[88]).}%
\TextAt{15}{79.226}{276.212}{§}%
\TextAt{15}{93.785}{276.212}{4}%
\TextAt{15}{120.693}{276.212}{Formal}%
\TextAt{15}{186.267}{276.212}{Groups}%
\TextAt{8}{79.226}{244.326}{The}%
\TextAt{8}{103.513}{244.326}{most}%
\TextAt{8}{132.427}{244.326}{explicit}%
\TextAt{8}{173.953}{244.326}{realization}%
\TextAt{8}{231.427}{244.326}{of}%
\TextAt{8}{244.991}{244.326}{local}%
\TextAt{8}{272.864}{244.326}{class}%
\TextAt{8}{300.534}{244.326}{field}%
\TextAt{8}{326.124}{244.326}{theory}%
\TextAt{8}{363.108}{244.326}{we}%
\TextAt{8}{380.568}{244.326}{have}%
\TextAt{8}{407.796}{244.326}{encountered}%
\TextAt{8}{473.388}{244.326}{for}%
\TextAt{8}{491.506}{244.326}{the}%
\TextAt{8}{511.896}{244.326}{case}%
\TextAt{8}{79.226}{227.28}{of}%
\TextAt{8}{93.249}{227.28}{cyclotomic}%
\TextAt{8}{152.151}{227.28}{fields}%
\TextAt{8}{182.815}{227.28}{over}%
\TextAt{8}{208.554}{227.28}{the}%
\TextAt{8}{229.402}{227.28}{field}%
\TextAt{18}{255.452}{227.28}{Q}%
\TextAt{21}{264.75}{225.487}{p}%
\TextAt{8}{269.511}{227.28}{,}%
\TextAt{8}{277.526}{227.28}{i.e.,}%
\TextAt{8}{300.497}{227.28}{with}%
\TextAt{8}{327.849}{227.28}{the}%
\TextAt{8}{348.698}{227.28}{extensions}%
\TextAt{18}{405.771}{227.28}{Q}%
\TextAt{21}{415.07}{225.487}{p}%
\TextAt{9}{419.831}{227.28}{\Character{40}}%
\TextAt{13}{424.383}{227.28}{\Character{16}}%
\TextAt{9}{430.399}{227.28}{\Character{41}}%
\TextAt{16}{439.452}{227.28}{\Character{106}}%
\TextAt{18}{447.274}{227.28}{Q}%
\TextAt{21}{456.572}{225.487}{p}%
\TextAt{8}{461.333}{227.28}{,}%
\TextAt{8}{469.348}{227.28}{where}%
\TextAt{13}{503.85}{227.28}{\Character{16}}%
\TextAt{8}{514.455}{227.28}{is}%
\TextAt{8}{526.911}{227.28}{a}%
\TextAt{13}{79.226}{210.234}{p}%
\TextAt{21}{85.101}{214.572}{n}%
\TextAt{8}{90.738}{210.234}{-th}%
\TextAt{8}{110.252}{210.234}{root}%
\TextAt{8}{135.96}{210.234}{of}%
\TextAt{8}{149.951}{210.234}{unity.}%
\TextAt{8}{186.075}{210.234}{The}%
\TextAt{8}{210.79}{210.234}{notion}%
\TextAt{8}{247.879}{210.234}{of}%
\TextAt{8}{261.871}{210.234}{formal}%
\TextAt{8}{299.282}{210.234}{group}%
\TextAt{8}{333.119}{210.234}{allows}%
\TextAt{8}{368.642}{210.234}{us}%
\TextAt{8}{384.319}{210.234}{to}%
\TextAt{8}{399.29}{210.234}{construct}%
\TextAt{8}{451.395}{210.234}{such}%
\TextAt{8}{478.453}{210.234}{an}%
\TextAt{8}{495.374}{210.234}{explicit}%
\TextAt{8}{79.226}{193.188}{cyclotomic}%
\TextAt{8}{138.004}{193.188}{theory}%
\TextAt{8}{175.323}{193.188}{over}%
\TextAt{8}{200.937}{193.188}{an}%
\TextAt{8}{217.765}{193.188}{arbitrary}%
\TextAt{8}{268.102}{193.188}{local}%
\TextAt{8}{296.311}{193.188}{field}%
\TextAt{13}{322.237}{193.188}{K}%
\TextAt{8}{337.52}{193.188}{by}%
\TextAt{8}{354.347}{193.188}{introducing}%
\TextAt{8}{417.358}{193.188}{a}%
\TextAt{8}{427.681}{193.188}{new}%
\TextAt{8}{452.304}{193.188}{kind}%
\TextAt{8}{479.209}{193.188}{of}%
\TextAt{8}{493.108}{193.188}{roots}%
\TextAt{8}{523.337}{193.188}{of}%
\TextAt{8}{79.226}{176.142}{unity}%
\TextAt{8}{110.22}{176.142}{which}%
\TextAt{8}{144.132}{176.142}{are}%
\TextAt{8}{164.068}{176.142}{“division}%
\TextAt{8}{213.329}{176.142}{points”}%
\TextAt{8}{254.462}{176.142}{that}%
\TextAt{8}{280.256}{176.142}{do}%
\TextAt{8}{296.94}{176.142}{the}%
\TextAt{8}{317.521}{176.142}{same}%
\TextAt{8}{347.272}{176.142}{for}%
\TextAt{8}{365.582}{176.142}{the}%
\TextAt{8}{386.163}{176.142}{field}%
\TextAt{13}{411.945}{176.142}{K}%
\TextAt{8}{427.085}{176.142}{as}%
\TextAt{8}{441.88}{176.142}{the}%
\TextAt{13}{462.461}{176.142}{p}%
\TextAt{21}{468.337}{180.48}{n}%
\TextAt{8}{473.973}{176.142}{-th}%
\TextAt{8}{493.251}{176.142}{roots}%
\TextAt{8}{523.337}{176.142}{of}%
\TextAt{8}{79.226}{159.096}{unity}%
\TextAt{8}{109.796}{159.096}{do}%
\TextAt{8}{126.055}{159.096}{for}%
\TextAt{8}{143.939}{159.096}{the}%
\TextAt{8}{164.096}{159.096}{field}%
\TextAt{18}{189.453}{159.096}{Q}%
\TextAt{21}{198.751}{157.303}{p}%
\TextAt{8}{203.512}{159.096}{.}%
\TextAt{12}{79.226}{133.084}{Definition}%
\TextAt{12}{143.84}{133.084}{(4.1).}%
\TextAt{8}{180.743}{133.084}{A}%
\TextAt{8}{195.001}{133.084}{(1-dimensional,}%
\TextAt{8}{278.998}{133.084}{commutative)}%
\TextAt{12}{355.066}{133.084}{formal}%
\TextAt{12}{399.213}{133.084}{group}%
\TextAt{8}{438.602}{133.084}{over}%
\TextAt{8}{465.236}{133.084}{a}%
\TextAt{8}{476.578}{133.084}{ring}%
\TextAt{14}{502.231}{133.084}{o}%
\TextAt{8}{513.56}{133.084}{is}%
\TextAt{8}{526.911}{133.084}{a}%
\TextAt{8}{79.226}{116.038}{formal}%
\TextAt{8}{115.976}{116.038}{power}%
\TextAt{8}{150.12}{116.038}{series}%
\TextAt{13}{181.455}{116.038}{F}%
\TextAt{9}{190.658}{116.038}{\Character{40}}%
\TextAt{13}{195.211}{116.038}{X}%
\TextAt{13}{205.206}{116.038}{\Character{59}}%
\TextAt{13}{210.43}{116.038}{Y}%
\TextAt{9}{219.853}{116.038}{\Character{41}}%
\NativeMathAt{227.726}{116.038}{11.955}{\in}%
\TextAt{14}{239.017}{116.038}{o}%
\TextAt{9}{244.861}{116.038}{\Character{91}\Character{91}}%
\TextAt{13}{251.365}{116.038}{X}%
\TextAt{13}{261.36}{116.038}{\Character{59}}%
\TextAt{13}{266.584}{116.038}{Y}%
\TextAt{9}{276.007}{116.038}{\Character{93}\Character{93}}%
\TextAt{8}{286.407}{116.038}{with}%
\TextAt{8}{313.067}{116.038}{the}%
\TextAt{8}{333.224}{116.038}{following}%
\TextAt{8}{382.658}{116.038}{properties:}%
\TextAt{8}{89.579}{96.003}{(i)}%
\TextAt{13}{107.918}{96.003}{F}%
\TextAt{9}{117.122}{96.003}{\Character{40}}%
\TextAt{13}{121.674}{96.003}{X}%
\TextAt{13}{131.668}{96.003}{\Character{59}}%
\TextAt{13}{136.893}{96.003}{Y}%
\TextAt{9}{146.317}{96.003}{\Character{41}}%
\TextAt{9}{154.172}{96.003}{\Character{61}}%
\TextAt{13}{166.615}{96.003}{X}%
\TextAt{9}{179.927}{96.003}{\Character{43}}%
\TextAt{13}{191.688}{96.003}{Y}%
\TextAt{8}{204.402}{96.003}{mod}%
\TextAt{8}{229.831}{96.003}{deg}%
\TextAt{9}{249.395}{96.003}{2}%
\TextAt{8}{255.248}{96.003}{,}%
\TextAt{8}{86.327}{78.957}{(ii)}%
\TextAt{13}{107.918}{78.957}{F}%
\TextAt{9}{117.122}{78.957}{\Character{40}}%
\TextAt{13}{121.674}{78.957}{X}%
\TextAt{13}{131.668}{78.957}{\Character{59}}%
\TextAt{13}{136.893}{78.957}{Y}%
\TextAt{9}{146.317}{78.957}{\Character{41}}%
\TextAt{9}{154.172}{78.957}{\Character{61}}%
\TextAt{13}{166.615}{78.957}{F}%
\TextAt{9}{175.819}{78.957}{\Character{40}}%
\TextAt{13}{180.371}{78.957}{Y}%
\TextAt{13}{187.807}{78.957}{\Character{59}}%
\TextAt{13}{193.044}{78.957}{X}%
\TextAt{9}{203.713}{78.957}{\Character{41}}%
\TextAt{8}{443.656}{78.957}{“commutativity”,}%
\TextAt{8}{297.211}{43.033}{313}%
\Destination{79.23}{773.075}{nav.Item.105}%
\Destination{79.23}{758.125}{nav.Item.106}%
\Destination{79.23}{741.075}{nav.Item.107}%
\Destination{79.23}{724.035}{nav.Item.108}%
\Destination{79.23}{690.615}{nav.Item.109}%
\Destination{79.23}{672.895}{nav.Item.110}%
\Destination{79.23}{110.055}{nav.Item.111}%
\Destination{79.23}{93.005}{nav.Item.112}%
\Destination{78.23}{809.815}{nav.page.313}%
\Destination{79.23}{302.065}{nav.section.5.4}%
\Destination{79.226}{640.855}{nav.ex.5.3.6}%
\Destination{79.226}{418.555}{nav.ex.5.3.7}%
\Destination{79.36}{641.855}{nav.print.341}%
\Destination{79.36}{147.055}{nav.number.5.4.1}%
\Destination{79.23}{302.065}{outline.52}%
\Hotlink{426.72}{567.055}{14.306}{11.047}{nav.chapter.3}%
\Hotlink{453.316}{567.055}{6.858}{11.047}{nav.section.3.2}%
\Hotlink{483.236}{567.055}{18.574}{11.047}{nav.print.200}%
\Hotlink{424.713}{323.855}{19.22}{11.047}{nav.cite.N088}%
}

% Page 326 (body page 314)
\Page{%
\TextAt{2}{79.226}{796.943}{V.}%
\TextAt{2}{94.073}{796.943}{Local}%
\TextAt{2}{123.615}{796.943}{Class}%
\TextAt{2}{152.208}{796.943}{Field}%
\TextAt{2}{179.939}{796.943}{Theory}%
\Vector{q 0 G 0 g 0.299 w 7.226 18.729 m 460.769 18.729 l S Q}%
\TextAt{8}{83.076}{761.124}{(iii)}%
\TextAt{13}{107.918}{761.124}{F}%
\TextAt{9}{117.122}{761.124}{\Character{40}}%
\TextAt{13}{121.674}{761.124}{X}%
\TextAt{13}{131.668}{761.124}{\Character{59}}%
\TextAt{13}{136.893}{761.124}{F}%
\TextAt{9}{146.127}{761.124}{\Character{40}}%
\TextAt{13}{150.679}{761.124}{Y}%
\TextAt{13}{158.115}{761.124}{\Character{59}}%
\TextAt{13}{163.352}{761.124}{Z}%
\TextAt{9}{172.236}{761.124}{\Character{41}\Character{41}}%
\TextAt{9}{184.634}{761.124}{\Character{61}}%
\TextAt{13}{197.087}{761.124}{F}%
\TextAt{9}{206.291}{761.124}{\Character{40}}%
\TextAt{13}{210.843}{761.124}{F}%
\TextAt{9}{220.047}{761.124}{\Character{40}}%
\TextAt{13}{224.599}{761.124}{X}%
\TextAt{13}{234.594}{761.124}{\Character{59}}%
\TextAt{13}{239.818}{761.124}{Y}%
\TextAt{9}{249.241}{761.124}{\Character{41}}%
\TextAt{13}{253.794}{761.124}{\Character{59}}%
\TextAt{13}{259.018}{761.124}{Z}%
\TextAt{9}{267.909}{761.124}{\Character{41}}%
\TextAt{8}{456.866}{761.124}{“associativity”.}%
\TextAt{8}{97.159}{736.215}{From}%
\TextAt{8}{127.576}{736.215}{a}%
\TextAt{8}{137.024}{736.215}{formal}%
\TextAt{8}{173.466}{736.215}{group}%
\TextAt{8}{206.334}{736.215}{one}%
\TextAt{8}{227.486}{736.215}{gets}%
\TextAt{8}{251.304}{736.215}{an}%
\TextAt{8}{267.255}{736.215}{ordinary}%
\TextAt{8}{314.111}{736.215}{group}%
\TextAt{8}{346.979}{736.215}{by}%
\TextAt{8}{362.93}{736.215}{evaluating}%
\TextAt{8}{418.896}{736.215}{in}%
\TextAt{8}{432.241}{736.215}{a}%
\TextAt{8}{441.689}{736.215}{domain}%
\TextAt{8}{483.009}{736.215}{where}%
\TextAt{8}{516.51}{736.215}{the}%
\TextAt{8}{79.226}{719.169}{power}%
\TextAt{8}{114.067}{719.169}{series}%
\TextAt{8}{146.098}{719.169}{converge.}%
\TextAt{8}{200.549}{719.169}{If}%
\TextAt{8}{212.938}{719.169}{for}%
\TextAt{8}{231.52}{719.169}{instance}%
\TextAt{14}{277.802}{719.169}{o}%
\TextAt{8}{288.24}{719.169}{is}%
\TextAt{8}{300.701}{719.169}{a}%
\TextAt{8}{311.154}{719.169}{complete}%
\TextAt{8}{361.273}{719.169}{valuation}%
\TextAt{8}{413.043}{719.169}{ring}%
\TextAt{8}{437.805}{719.169}{and}%
\TextAt{14}{461.265}{719.169}{p}%
\TextAt{8}{471.842}{719.169}{its}%
\TextAt{8}{488.858}{719.169}{maximal}%
\TextAt{8}{79.226}{702.124}{ideal,}%
\TextAt{8}{110.441}{702.124}{then}%
\TextAt{8}{137.101}{702.124}{the}%
\TextAt{8}{157.258}{702.124}{operation}%
\TextAt{13}{261.208}{685.078}{x}%
\TextAt{9}{270.517}{685.078}{\Character{43}}%
\TextAt{21}{279.622}{683.284}{F}%
\TextAt{13}{289.312}{685.078}{y}%
\TextAt{9}{298.769}{685.078}{\Character{58}\Character{61}}%
\TextAt{13}{314.446}{685.078}{F}%
\TextAt{9}{323.65}{685.078}{\Character{40}}%
\TextAt{13}{328.202}{685.078}{x\Character{59}}%
\TextAt{13}{340.074}{685.078}{y}%
\TextAt{9}{346.235}{685.078}{\Character{41}}%
\TextAt{8}{79.226}{663.212}{defines}%
\TextAt{8}{117.65}{663.212}{a}%
\TextAt{8}{127.405}{663.212}{new}%
\TextAt{8}{151.459}{663.212}{structure}%
\TextAt{8}{201.599}{663.212}{of}%
\TextAt{8}{214.929}{663.212}{abelian}%
\TextAt{8}{255.577}{663.212}{group}%
\TextAt{8}{288.752}{663.212}{on}%
\TextAt{8}{305.011}{663.212}{the}%
\TextAt{8}{325.168}{663.212}{set}%
\TextAt{14}{343.435}{663.212}{p}%
\TextAt{8}{349.418}{663.212}{.}%
\TextAt{12}{79.226}{643.177}{Examples:}%
\TextAt{8}{92.831}{623.143}{1.}%
\TextAt{18}{107.918}{623.143}{G}%
\TextAt{21}{117.217}{621.349}{a}%
\TextAt{9}{122.213}{623.143}{\Character{40}}%
\TextAt{13}{126.765}{623.143}{X}%
\TextAt{13}{136.76}{623.143}{\Character{59}}%
\TextAt{13}{141.984}{623.143}{Y}%
\TextAt{9}{151.408}{623.143}{\Character{41}}%
\TextAt{9}{159.263}{623.143}{\Character{61}}%
\TextAt{13}{171.706}{623.143}{X}%
\TextAt{9}{185.018}{623.143}{\Character{43}}%
\TextAt{13}{196.779}{623.143}{Y}%
\TextAt{8}{210.07}{623.143}{(the}%
\TextAt{8}{234.781}{623.143}{formal}%
\TextAt{8}{271.531}{623.143}{additive}%
\TextAt{8}{316.411}{623.143}{group).}%
\TextAt{8}{92.831}{606.097}{2.}%
\TextAt{18}{107.918}{606.097}{G}%
\TextAt{21}{117.217}{604.303}{m}%
\TextAt{9}{125.206}{606.097}{\Character{40}}%
\TextAt{13}{129.758}{606.097}{X}%
\TextAt{13}{139.753}{606.097}{\Character{59}}%
\TextAt{13}{144.977}{606.097}{Y}%
\TextAt{9}{154.4}{606.097}{\Character{41}}%
\TextAt{9}{162.255}{606.097}{\Character{61}}%
\TextAt{13}{174.699}{606.097}{X}%
\TextAt{9}{187.611}{606.097}{\Character{43}}%
\TextAt{13}{198.972}{606.097}{Y}%
\TextAt{9}{210.622}{606.097}{\Character{43}}%
\TextAt{13}{221.984}{606.097}{X}%
\TextAt{13}{232.624}{606.097}{Y}%
\TextAt{8}{245.734}{606.097}{(the}%
\TextAt{8}{270.25}{606.097}{formal}%
\TextAt{8}{306.805}{606.097}{multiplicative}%
\TextAt{8}{380.432}{606.097}{group).}%
\TextAt{8}{422.649}{606.097}{Since}%
\TextAt{13}{453.011}{606.097}{X}%
\TextAt{9}{465.923}{606.097}{\Character{43}}%
\TextAt{13}{477.284}{606.097}{Y}%
\TextAt{9}{488.934}{606.097}{\Character{43}}%
\TextAt{13}{500.296}{606.097}{X}%
\TextAt{13}{510.936}{606.097}{Y}%
\TextAt{9}{523.665}{606.097}{\Character{61}}%
\TextAt{9}{107.918}{589.051}{\Character{40}1}%
\TextAt{9}{120.961}{589.051}{\Character{43}}%
\TextAt{13}{132.742}{589.051}{X}%
\TextAt{9}{143.397}{589.051}{\Character{41}\Character{40}1}%
\TextAt{9}{160.983}{589.051}{\Character{43}}%
\TextAt{13}{172.772}{589.051}{Y}%
\TextAt{9}{182.166}{589.051}{\Character{41}}%
\TextAt{16}{189.375}{589.051}{\Character{0}}%
\TextAt{9}{201.33}{589.051}{1}%
\TextAt{8}{207.183}{589.051}{,}%
\TextAt{8}{214.332}{589.051}{we}%
\TextAt{8}{231.559}{589.051}{have}%
\TextAt{9}{235.6}{563.914}{\Character{40}}%
\TextAt{13}{240.152}{563.914}{x}%
\TextAt{9}{249.461}{563.914}{\Character{43}}%
\TextAt{25}{258.566}{562.121}{G}%
\TextAt{22}{264.765}{561.124}{m}%
\TextAt{13}{274.956}{563.914}{y}%
\TextAt{9}{281.092}{563.914}{\Character{41}}%
\TextAt{9}{288.289}{563.914}{\Character{43}}%
\TextAt{9}{300.041}{563.914}{1}%
\TextAt{9}{309.211}{563.914}{\Character{61}}%
\TextAt{9}{321.632}{563.914}{\Character{40}}%
\TextAt{13}{326.214}{563.914}{x}%
\TextAt{9}{335.523}{563.914}{\Character{43}}%
\TextAt{9}{347.275}{563.914}{1\Character{41}}%
\TextAt{16}{360.346}{563.914}{\Character{1}}%
\TextAt{9}{366.324}{563.914}{\Character{40}}%
\TextAt{13}{370.876}{563.914}{y}%
\TextAt{9}{379.67}{563.914}{\Character{43}}%
\TextAt{9}{391.422}{563.914}{1\Character{41}}%
\TextAt{13}{401.836}{563.914}{\Character{58}}%
\TextAt{8}{107.918}{538.777}{So}%
\TextAt{8}{124.33}{538.777}{the}%
\TextAt{8}{144.638}{538.777}{new}%
\TextAt{8}{168.844}{538.777}{operation}%
\TextAt{9}{221.36}{538.777}{\Character{43}}%
\TextAt{25}{230.465}{536.984}{G}%
\TextAt{22}{236.664}{535.988}{m}%
\TextAt{8}{248.248}{538.777}{is}%
\TextAt{8}{260.164}{538.777}{obtained}%
\TextAt{8}{308.448}{538.777}{from}%
\TextAt{8}{336.24}{538.777}{multiplication}%
\TextAt{8}{411.519}{538.777}{via}%
\TextAt{8}{430.859}{538.777}{the}%
\TextAt{8}{451.168}{538.777}{translation}%
\TextAt{13}{510.582}{538.777}{x}%
\NativeArrowAt{520.814}{538.777}{11.955}{11.955}{\mapsto}%
\TextAt{13}{107.918}{521.731}{x}%
\TextAt{9}{117.227}{521.731}{\Character{43}}%
\TextAt{9}{128.979}{521.731}{1}%
\TextAt{8}{134.842}{521.731}{.}%
\TextAt{8}{92.831}{504.685}{3.}%
\TextAt{8}{107.918}{504.685}{A}%
\TextAt{8}{120.325}{504.685}{power}%
\TextAt{8}{154.203}{504.685}{series}%
\TextAt{13}{185.272}{504.685}{f}%
\TextAt{9}{192.318}{504.685}{\Character{40}}%
\TextAt{13}{196.871}{504.685}{X}%
\TextAt{9}{207.526}{504.685}{\Character{41}}%
\TextAt{9}{215.381}{504.685}{\Character{61}}%
\TextAt{13}{227.824}{504.685}{a}%
\TextAt{17}{233.969}{502.892}{1}%
\TextAt{13}{238.702}{504.685}{X}%
\TextAt{9}{251.47}{504.685}{\Character{43}}%
\TextAt{13}{262.688}{504.685}{a}%
\TextAt{17}{268.833}{502.892}{2}%
\TextAt{13}{273.565}{504.685}{X}%
\TextAt{17}{284.22}{509.024}{2}%
\TextAt{9}{291.066}{504.685}{\Character{43}}%
\TextAt{16}{302.284}{504.685}{\Character{1}}%
\TextAt{16}{307.58}{504.685}{\Character{1}}%
\TextAt{16}{312.888}{504.685}{\Character{1}}%
\NativeMathAt{319.523}{504.685}{11.955}{\in}%
\TextAt{14}{330.843}{504.685}{o}%
\TextAt{9}{336.687}{504.685}{\Character{91}\Character{91}}%
\TextAt{13}{343.191}{504.685}{X}%
\TextAt{9}{353.846}{504.685}{\Character{93}\Character{93}}%
\TextAt{8}{363.981}{504.685}{whose}%
\TextAt{8}{398.242}{504.685}{first}%
\TextAt{8}{422.101}{504.685}{coefficient}%
\TextAt{13}{476.459}{504.685}{a}%
\TextAt{17}{482.604}{502.892}{1}%
\TextAt{8}{490.968}{504.685}{is}%
\TextAt{8}{502.466}{504.685}{a}%
\TextAt{8}{511.956}{504.685}{unit}%
\TextAt{8}{107.918}{487.639}{admits}%
\TextAt{8}{146.339}{487.639}{an}%
\TextAt{8}{162.583}{487.639}{“inverse”,}%
\TextAt{8}{215.628}{487.639}{i.e.,}%
\TextAt{8}{237.721}{487.639}{there}%
\TextAt{8}{267.618}{487.639}{exists}%
\TextAt{8}{299.917}{487.639}{a}%
\TextAt{8}{309.658}{487.639}{power}%
\TextAt{8}{343.787}{487.639}{series}%
\TextAt{13}{375.106}{487.639}{f}%
\TextAt{3}{382.152}{491.978}{\Character{0}}%
\TextAt{17}{388.739}{491.978}{1}%
\TextAt{9}{393.471}{487.639}{\Character{40}}%
\TextAt{13}{398.024}{487.639}{X}%
\TextAt{9}{408.679}{487.639}{\Character{41}}%
\TextAt{9}{416.534}{487.639}{\Character{61}}%
\TextAt{13}{428.977}{487.639}{a}%
\TextAt{3}{435.122}{492.587}{\Character{0}}%
\TextAt{17}{441.709}{492.587}{1}%
\TextAt{17}{435.122}{484.708}{1}%
\TextAt{13}{446.441}{487.639}{X}%
\TextAt{9}{459.722}{487.639}{\Character{43}}%
\TextAt{16}{471.452}{487.639}{\Character{1}}%
\TextAt{16}{476.748}{487.639}{\Character{1}}%
\TextAt{16}{482.056}{487.639}{\Character{1}}%
\NativeMathAt{488.691}{487.639}{11.955}{\in}%
\TextAt{14}{500.012}{487.639}{o}%
\TextAt{9}{505.856}{487.639}{\Character{91}\Character{91}}%
\TextAt{13}{512.359}{487.639}{X}%
\TextAt{9}{523.014}{487.639}{\Character{93}\Character{93}}%
\TextAt{8}{529.518}{487.639}{,}%
\TextAt{8}{107.918}{470.593}{such}%
\TextAt{8}{134.315}{470.593}{that}%
\TextAt{13}{159.684}{470.593}{f}%
\TextAt{3}{166.731}{474.932}{\Character{0}}%
\TextAt{17}{173.317}{474.932}{1}%
\TextAt{9}{178.05}{470.593}{\Character{40}}%
\TextAt{13}{182.602}{470.593}{f}%
\TextAt{9}{189.648}{470.593}{\Character{40}}%
\TextAt{13}{194.201}{470.593}{X}%
\TextAt{9}{204.856}{470.593}{\Character{41}\Character{41}}%
\TextAt{9}{217.254}{470.593}{\Character{61}}%
\TextAt{13}{229.707}{470.593}{f}%
\TextAt{9}{236.753}{470.593}{\Character{40}}%
\TextAt{13}{241.306}{470.593}{f}%
\TextAt{3}{248.352}{474.932}{\Character{0}}%
\TextAt{17}{254.938}{474.932}{1}%
\TextAt{9}{259.671}{470.593}{\Character{40}}%
\TextAt{13}{264.223}{470.593}{X}%
\TextAt{9}{274.878}{470.593}{\Character{41}\Character{41}}%
\TextAt{9}{287.276}{470.593}{\Character{61}}%
\TextAt{13}{299.729}{470.593}{X}%
\TextAt{8}{310.384}{470.593}{.}%
\TextAt{8}{318.833}{470.593}{For}%
\TextAt{8}{339.802}{470.593}{every}%
\TextAt{8}{370.694}{470.593}{such}%
\TextAt{8}{397.091}{470.593}{power}%
\TextAt{8}{431.235}{470.593}{series,}%
\TextAt{13}{243.876}{445.457}{F}%
\TextAt{9}{253.08}{445.457}{\Character{40}}%
\TextAt{13}{257.632}{445.457}{X}%
\TextAt{13}{267.627}{445.457}{\Character{59}}%
\TextAt{13}{272.851}{445.457}{Y}%
\TextAt{9}{282.275}{445.457}{\Character{41}}%
\TextAt{9}{290.13}{445.457}{\Character{61}}%
\TextAt{13}{302.573}{445.457}{f}%
\TextAt{3}{309.62}{450.393}{\Character{0}}%
\TextAt{17}{316.206}{450.393}{1}%
\TextAt{9}{320.938}{445.457}{\Character{40}}%
\TextAt{13}{325.491}{445.457}{f}%
\TextAt{9}{332.537}{445.457}{\Character{40}}%
\TextAt{13}{337.09}{445.457}{X}%
\TextAt{9}{347.745}{445.457}{\Character{41}}%
\TextAt{9}{354.942}{445.457}{\Character{43}}%
\TextAt{13}{366.715}{445.457}{f}%
\TextAt{9}{373.761}{445.457}{\Character{40}}%
\TextAt{13}{378.314}{445.457}{Y}%
\TextAt{9}{387.707}{445.457}{\Character{41}\Character{41}}%
\TextAt{8}{107.918}{420.32}{is}%
\TextAt{8}{119.682}{420.32}{a}%
\TextAt{8}{129.438}{420.32}{formal}%
\TextAt{8}{166.188}{420.32}{group.}%
\TextAt{12}{79.226}{395.412}{Definition}%
\TextAt{12}{143.603}{395.412}{(4.2).}%
\TextAt{8}{180.299}{395.412}{A}%
\TextAt{12}{194.351}{395.412}{homomorphism}%
\TextAt{13}{290.428}{395.412}{f}%
\TextAt{9}{303.148}{395.412}{\Character{58}}%
\TextAt{13}{312.073}{395.412}{F}%
\NativeArrowAt{326.95}{395.412}{11.955}{19.26}{\longrightarrow}%
\TextAt{13}{351.885}{395.412}{G}%
\TextAt{8}{366.397}{395.412}{between}%
\TextAt{8}{412.968}{395.412}{two}%
\TextAt{8}{436.465}{395.412}{formal}%
\TextAt{8}{474.596}{395.412}{groups}%
\TextAt{8}{513.767}{395.412}{is}%
\TextAt{8}{526.911}{395.412}{a}%
\TextAt{8}{79.226}{378.366}{power}%
\TextAt{8}{113.37}{378.366}{series}%
\TextAt{13}{144.705}{378.366}{f}%
\TextAt{9}{151.751}{378.366}{\Character{40}}%
\TextAt{13}{156.303}{378.366}{X}%
\TextAt{9}{166.958}{378.366}{\Character{41}}%
\TextAt{9}{174.813}{378.366}{\Character{61}}%
\TextAt{13}{187.257}{378.366}{a}%
\TextAt{17}{193.402}{376.572}{1}%
\TextAt{13}{198.134}{378.366}{X}%
\TextAt{9}{211.446}{378.366}{\Character{43}}%
\TextAt{13}{223.207}{378.366}{a}%
\TextAt{17}{229.352}{376.572}{2}%
\TextAt{13}{234.085}{378.366}{X}%
\TextAt{17}{244.74}{382.704}{2}%
\TextAt{9}{252.129}{378.366}{\Character{43}}%
\TextAt{16}{263.89}{378.366}{\Character{1}}%
\TextAt{16}{269.186}{378.366}{\Character{1}}%
\TextAt{16}{274.494}{378.366}{\Character{1}}%
\NativeMathAt{281.129}{378.366}{11.955}{\in}%
\TextAt{14}{292.45}{378.366}{o}%
\TextAt{9}{298.294}{378.366}{\Character{91}\Character{91}}%
\TextAt{13}{304.797}{378.366}{X}%
\TextAt{9}{315.452}{378.366}{\Character{93}\Character{93}}%
\TextAt{8}{325.853}{378.366}{such}%
\TextAt{8}{352.25}{378.366}{that}%
\TextAt{13}{228.982}{353.229}{f}%
\TextAt{9}{236.028}{353.229}{\Character{40}}%
\TextAt{13}{240.58}{353.229}{F}%
\TextAt{9}{249.784}{353.229}{\Character{40}}%
\TextAt{13}{254.336}{353.229}{X}%
\TextAt{13}{264.331}{353.229}{\Character{59}}%
\TextAt{13}{269.555}{353.229}{Y}%
\TextAt{9}{278.978}{353.229}{\Character{41}\Character{41}}%
\TextAt{9}{291.375}{353.229}{\Character{61}}%
\TextAt{13}{303.829}{353.229}{G}%
\TextAt{9}{313.063}{353.229}{\Character{40}}%
\TextAt{13}{317.615}{353.229}{f}%
\TextAt{9}{324.662}{353.229}{\Character{40}}%
\TextAt{13}{329.214}{353.229}{X}%
\TextAt{9}{339.869}{353.229}{\Character{41}}%
\TextAt{13}{344.422}{353.229}{\Character{59}}%
\TextAt{13}{349.646}{353.229}{f}%
\TextAt{9}{356.712}{353.229}{\Character{40}}%
\TextAt{13}{361.264}{353.229}{Y}%
\TextAt{9}{370.658}{353.229}{\Character{41}\Character{41}}%
\TextAt{13}{379.762}{353.229}{\Character{58}}%
\TextAt{8}{79.226}{320.23}{In}%
\TextAt{8}{93.845}{320.23}{example}%
\TextAt{8}{139.69}{320.23}{3,}%
\TextAt{8}{152.696}{320.23}{for}%
\TextAt{8}{170.578}{320.23}{instance,}%
\TextAt{8}{219.413}{320.23}{the}%
\TextAt{8}{239.567}{320.23}{power}%
\TextAt{8}{273.708}{320.23}{series}%
\TextAt{13}{305.04}{320.23}{f}%
\TextAt{8}{315.982}{320.23}{is}%
\TextAt{8}{327.743}{320.23}{a}%
\TextAt{8}{337.496}{320.23}{homomorphism}%
\TextAt{8}{420.164}{320.23}{of}%
\TextAt{8}{433.491}{320.23}{the}%
\TextAt{8}{453.645}{320.23}{formal}%
\TextAt{8}{490.393}{320.23}{group}%
\TextAt{13}{523.566}{320.23}{F}%
\TextAt{8}{79.226}{303.184}{to}%
\TextAt{8}{93.536}{303.184}{the}%
\TextAt{8}{113.693}{303.184}{additive}%
\TextAt{8}{158.573}{303.184}{group}%
\TextAt{18}{191.748}{303.184}{G}%
\TextAt{21}{201.047}{301.391}{a}%
\TextAt{8}{206.043}{303.184}{.}%
\TextAt{8}{214.491}{303.184}{It}%
\TextAt{8}{227.164}{303.184}{is}%
\TextAt{8}{238.928}{303.184}{called}%
\TextAt{8}{272.091}{303.184}{the}%
\TextAt{12}{292.248}{303.184}{logarithm}%
\TextAt{8}{353.052}{303.184}{of}%
\TextAt{13}{366.382}{303.184}{F}%
\TextAt{8}{375.585}{303.184}{.}%
\TextAt{8}{97.159}{283.149}{A}%
\TextAt{8}{109.384}{283.149}{homomorphism}%
\TextAt{13}{191.607}{283.149}{f}%
\TextAt{9}{201.974}{283.149}{\Character{58}}%
\TextAt{13}{208.547}{283.149}{F}%
\NativeArrowAt{221.071}{283.149}{11.955}{19.26}{\longrightarrow}%
\TextAt{13}{243.653}{283.149}{G}%
\TextAt{8}{256.337}{283.149}{is}%
\TextAt{8}{267.654}{283.149}{an}%
\TextAt{12}{283.466}{283.149}{isomorphism}%
\TextAt{8}{361.35}{283.149}{if}%
\TextAt{13}{371.626}{283.149}{a}%
\TextAt{17}{377.771}{281.356}{1}%
\TextAt{9}{385.825}{283.149}{\Character{61}}%
\TextAt{13}{398.25}{283.149}{f}%
\TextAt{3}{405.296}{287.488}{\Character{48}}%
\TextAt{9}{408.092}{283.149}{\Character{40}0\Character{41}}%
\TextAt{8}{426.5}{283.149}{is}%
\TextAt{8}{437.817}{283.149}{a}%
\TextAt{8}{447.125}{283.149}{unit,}%
\TextAt{8}{474.731}{283.149}{i.e.,}%
\TextAt{8}{496.478}{283.149}{if}%
\TextAt{8}{506.755}{283.149}{there}%
\TextAt{8}{79.226}{266.103}{is}%
\TextAt{8}{90.99}{266.103}{a}%
\TextAt{8}{100.745}{266.103}{homomorphism}%
\TextAt{13}{183.415}{266.103}{g}%
\TextAt{9}{192.771}{266.103}{\Character{61}}%
\TextAt{13}{205.196}{266.103}{f}%
\TextAt{3}{212.242}{270.442}{\Character{0}}%
\TextAt{17}{218.829}{270.442}{1}%
\TextAt{9}{226.882}{266.103}{\Character{58}}%
\TextAt{13}{233.455}{266.103}{G}%
\NativeArrowAt{246.009}{266.103}{11.955}{19.26}{\longrightarrow}%
\TextAt{13}{268.591}{266.103}{F}%
\TextAt{8}{281.692}{266.103}{such}%
\TextAt{8}{308.089}{266.103}{that}%
\TextAt{13}{333.458}{266.103}{f}%
\TextAt{9}{340.504}{266.103}{\Character{40}}%
\TextAt{13}{345.057}{266.103}{g}%
\TextAt{9}{351.091}{266.103}{\Character{40}}%
\TextAt{13}{355.643}{266.103}{X}%
\TextAt{9}{366.298}{266.103}{\Character{41}\Character{41}}%
\TextAt{9}{378.696}{266.103}{\Character{61}}%
\TextAt{13}{391.149}{266.103}{g}%
\TextAt{9}{397.184}{266.103}{\Character{40}}%
\TextAt{13}{401.736}{266.103}{f}%
\TextAt{9}{408.782}{266.103}{\Character{40}}%
\TextAt{13}{413.335}{266.103}{X}%
\TextAt{9}{423.99}{266.103}{\Character{41}\Character{41}}%
\TextAt{9}{436.387}{266.103}{\Character{61}}%
\TextAt{13}{448.841}{266.103}{X}%
\TextAt{8}{459.496}{266.103}{.}%
\TextAt{8}{97.159}{246.069}{If}%
\TextAt{8}{109.903}{246.069}{the}%
\TextAt{8}{131.111}{246.069}{power}%
\TextAt{8}{166.307}{246.069}{series}%
\TextAt{13}{198.693}{246.069}{f}%
\TextAt{9}{205.739}{246.069}{\Character{40}}%
\TextAt{13}{210.292}{246.069}{X}%
\TextAt{9}{220.947}{246.069}{\Character{41}}%
\TextAt{9}{230.595}{246.069}{\Character{61}}%
\TextAt{13}{244.83}{246.069}{a}%
\TextAt{17}{250.975}{244.275}{1}%
\TextAt{13}{255.707}{246.069}{X}%
\TextAt{9}{269.736}{246.069}{\Character{43}}%
\TextAt{13}{282.214}{246.069}{a}%
\TextAt{17}{288.359}{244.275}{2}%
\TextAt{13}{293.091}{246.069}{X}%
\TextAt{17}{303.746}{250.407}{2}%
\TextAt{9}{311.852}{246.069}{\Character{43}}%
\TextAt{16}{324.33}{246.069}{\Character{1}}%
\TextAt{16}{329.626}{246.069}{\Character{1}}%
\TextAt{16}{334.934}{246.069}{\Character{1}}%
\TextAt{8}{345.22}{246.069}{satisfies}%
\TextAt{8}{389.382}{246.069}{the}%
\TextAt{8}{410.59}{246.069}{equation}%
\TextAt{13}{459.45}{246.069}{f}%
\TextAt{9}{466.497}{246.069}{\Character{40}}%
\TextAt{13}{471.049}{246.069}{F}%
\TextAt{9}{480.253}{246.069}{\Character{40}}%
\TextAt{13}{484.805}{246.069}{X}%
\TextAt{13}{494.8}{246.069}{\Character{59}}%
\TextAt{13}{500.024}{246.069}{Y}%
\TextAt{9}{509.447}{246.069}{\Character{41}\Character{41}}%
\TextAt{9}{523.638}{246.069}{\Character{61}}%
\TextAt{13}{79.226}{229.023}{G}%
\TextAt{9}{88.46}{229.023}{\Character{40}}%
\TextAt{13}{93.012}{229.023}{f}%
\TextAt{9}{100.059}{229.023}{\Character{40}}%
\TextAt{13}{104.611}{229.023}{X}%
\TextAt{9}{115.266}{229.023}{\Character{41}}%
\TextAt{13}{119.818}{229.023}{\Character{59}}%
\TextAt{13}{125.042}{229.023}{f}%
\TextAt{9}{132.109}{229.023}{\Character{40}}%
\TextAt{13}{136.661}{229.023}{Y}%
\TextAt{9}{146.054}{229.023}{\Character{41}\Character{41}}%
\TextAt{8}{155.159}{229.023}{,}%
\TextAt{8}{164.018}{229.023}{but}%
\TextAt{8}{186.845}{229.023}{its}%
\TextAt{8}{204.532}{229.023}{coefficients}%
\TextAt{8}{265.137}{229.023}{belong}%
\TextAt{8}{303.901}{229.023}{to}%
\TextAt{8}{319.579}{229.023}{an}%
\TextAt{8}{337.206}{229.023}{extension}%
\TextAt{8}{390.339}{229.023}{ring}%
\TextAt{14}{415.773}{229.023}{o}%
\TextAt{3}{421.617}{233.361}{\Character{48}}%
\TextAt{8}{424.412}{229.023}{,}%
\TextAt{8}{433.271}{229.023}{then}%
\TextAt{8}{461.299}{229.023}{we}%
\TextAt{8}{479.894}{229.023}{call}%
\TextAt{8}{502.721}{229.023}{this}%
\TextAt{8}{526.911}{229.023}{a}%
\TextAt{8}{79.226}{211.977}{homomorphism}%
\TextAt{12}{161.896}{211.977}{defined}%
\TextAt{12}{208.533}{211.977}{over}%
\TextAt{14}{237.74}{211.977}{o}%
\TextAt{3}{243.584}{216.315}{\Character{48}}%
\TextAt{8}{246.379}{211.977}{.}%
\TextAt{8}{254.827}{211.977}{The}%
\TextAt{8}{278.881}{211.977}{following}%
\TextAt{8}{328.316}{211.977}{proposition}%
\TextAt{8}{389.849}{211.977}{is}%
\TextAt{8}{401.613}{211.977}{immediately}%
\TextAt{8}{468.275}{211.977}{evident.}%
\TextAt{12}{79.226}{187.068}{Proposition}%
\TextAt{12}{152.204}{187.068}{(4.3).}%
\TextAt{8}{187.751}{187.068}{The}%
\TextAt{8}{212.036}{187.068}{homomorphisms}%
\TextAt{13}{299.552}{187.068}{f}%
\TextAt{9}{310.313}{187.068}{\Character{58}}%
\TextAt{13}{317.28}{187.068}{F}%
\NativeArrowAt{330.199}{187.068}{11.955}{19.26}{\longrightarrow}%
\TextAt{13}{353.175}{187.068}{F}%
\TextAt{8}{366.507}{187.068}{of}%
\TextAt{8}{380.069}{187.068}{a}%
\TextAt{8}{390.055}{187.068}{formal}%
\TextAt{8}{427.037}{187.068}{group}%
\TextAt{13}{460.444}{187.068}{F}%
\TextAt{8}{473.776}{187.068}{over}%
\TextAt{14}{499.054}{187.068}{o}%
\TextAt{8}{509.027}{187.068}{form}%
\TextAt{8}{79.226}{170.022}{a}%
\TextAt{8}{88.982}{170.022}{ring}%
\TextAt{8}{113.047}{170.022}{End}%
\TextAt{27}{134.017}{168.229}{o}%
\TextAt{9}{138.411}{170.022}{\Character{40}}%
\TextAt{13}{142.963}{170.022}{F}%
\TextAt{9}{152.167}{170.022}{\Character{41}}%
\TextAt{8}{160.617}{170.022}{in}%
\TextAt{8}{174.269}{170.022}{which}%
\TextAt{8}{207.756}{170.022}{addition}%
\TextAt{8}{253.939}{170.022}{and}%
\TextAt{8}{276.701}{170.022}{multiplication}%
\TextAt{8}{351.828}{170.022}{are}%
\TextAt{8}{371.338}{170.022}{defined}%
\TextAt{8}{411.651}{170.022}{by}%
\TextAt{9}{156.683}{144.886}{\Character{40}}%
\TextAt{13}{161.236}{144.886}{f}%
\TextAt{9}{170.939}{144.886}{\Character{43}}%
\TextAt{21}{180.043}{143.092}{F}%
\TextAt{13}{189.733}{144.886}{g}%
\TextAt{9}{195.768}{144.886}{\Character{41}\Character{40}}%
\TextAt{13}{204.872}{144.886}{X}%
\TextAt{9}{215.528}{144.886}{\Character{41}}%
\TextAt{9}{223.383}{144.886}{\Character{61}}%
\TextAt{13}{235.826}{144.886}{F}%
\TextAt{9}{245.03}{144.886}{\Character{40}}%
\TextAt{13}{249.582}{144.886}{f}%
\TextAt{9}{256.629}{144.886}{\Character{40}}%
\TextAt{13}{261.181}{144.886}{X}%
\TextAt{9}{271.836}{144.886}{\Character{41}}%
\TextAt{13}{276.388}{144.886}{\Character{59}}%
\TextAt{13}{281.612}{144.886}{g}%
\TextAt{9}{287.667}{144.886}{\Character{40}}%
\TextAt{13}{292.219}{144.886}{X}%
\TextAt{9}{302.874}{144.886}{\Character{41}\Character{41}}%
\TextAt{13}{311.979}{144.886}{\Character{59}}%
\TextAt{9}{341.133}{144.886}{\Character{40}}%
\TextAt{13}{345.686}{144.886}{f}%
\TextAt{16}{355.389}{144.886}{\Character{14}}%
\TextAt{13}{364.023}{144.886}{g}%
\TextAt{9}{370.057}{144.886}{\Character{41}\Character{40}}%
\TextAt{13}{379.162}{144.886}{X}%
\TextAt{9}{389.817}{144.886}{\Character{41}}%
\TextAt{9}{397.672}{144.886}{\Character{61}}%
\TextAt{13}{410.116}{144.886}{f}%
\TextAt{9}{417.162}{144.886}{\Character{40}}%
\TextAt{13}{421.714}{144.886}{g}%
\TextAt{9}{427.749}{144.886}{\Character{40}}%
\TextAt{13}{432.301}{144.886}{X}%
\TextAt{9}{442.956}{144.886}{\Character{41}\Character{41}}%
\TextAt{13}{452.061}{144.886}{\Character{58}}%
\TextAt{12}{79.226}{107.013}{Definition}%
\TextAt{12}{142.089}{107.013}{(4.4).}%
\TextAt{8}{177.469}{107.013}{A}%
\TextAt{12}{190.206}{107.013}{formal}%
\TextAt{14}{232.601}{107.013}{o}%
\TextAt{12}{238.445}{107.013}{-module}%
\TextAt{8}{290.025}{107.013}{is}%
\TextAt{8}{301.853}{107.013}{a}%
\TextAt{8}{311.673}{107.013}{formal}%
\TextAt{8}{348.487}{107.013}{group}%
\TextAt{13}{381.727}{107.013}{F}%
\TextAt{8}{394.893}{107.013}{over}%
\TextAt{14}{420.003}{107.013}{o}%
\TextAt{8}{429.809}{107.013}{together}%
\TextAt{8}{476.057}{107.013}{with}%
\TextAt{8}{502.781}{107.013}{a}%
\TextAt{8}{512.601}{107.013}{ring}%
\TextAt{8}{79.226}{89.967}{homomorphism}%
\TextAt{14}{216.341}{72.921}{o}%
\NativeArrowAt{225.506}{72.921}{11.955}{19.26}{\longrightarrow}%
\TextAt{8}{248.088}{72.921}{End}%
\TextAt{27}{269.057}{71.128}{o}%
\TextAt{9}{273.451}{72.921}{\Character{40}}%
\TextAt{13}{278.004}{72.921}{F}%
\TextAt{9}{287.207}{72.921}{\Character{41}}%
\TextAt{13}{291.76}{72.921}{\Character{59}}%
\TextAt{13}{320.895}{72.921}{a}%
\NativeArrowAt{330.38}{72.921}{11.955}{19.26}{\longmapsto}%
\TextAt{9}{352.962}{72.921}{\Character{91}}%
\TextAt{13}{356.214}{72.921}{a}%
\TextAt{9}{362.358}{72.921}{\Character{93}}%
\TextAt{21}{365.61}{71.128}{F}%
\TextAt{9}{372.644}{72.921}{\Character{40}}%
\TextAt{13}{377.196}{72.921}{X}%
\TextAt{9}{387.851}{72.921}{\Character{41}}%
\TextAt{13}{392.403}{72.921}{\Character{59}}%
\TextAt{8}{297.211}{43.033}{314}%
\Destination{79.23}{773.075}{nav.Item.113}%
\Destination{79.23}{637.865}{nav.Item.114}%
\Destination{79.23}{620.145}{nav.Item.115}%
\Destination{79.23}{520.815}{nav.Item.116}%
\Destination{78.23}{809.815}{nav.page.314}%
\Destination{79.36}{751.255}{nav.print.342}%
\Destination{79.36}{335.255}{nav.print.343}%
\Destination{79.36}{409.455}{nav.number.5.4.2}%
\Destination{79.36}{201.055}{nav.number.5.4.3}%
\Destination{79.36}{121.055}{nav.number.5.4.4}%
}

% Page 327 (body page 315)
\Page{%
\TextAt{2}{79.226}{796.943}{V.}%
\TextAt{2}{94.073}{796.943}{Local}%
\TextAt{2}{123.615}{796.943}{Class}%
\TextAt{2}{152.208}{796.943}{Field}%
\TextAt{2}{179.939}{796.943}{Theory}%
\Vector{q 0 G 0 g 0.299 w 7.226 18.729 m 460.769 18.729 l S Q}%
\TextAt{8}{79.226}{761.124}{such}%
\TextAt{8}{105.623}{761.124}{that}%
\TextAt{9}{130.992}{761.124}{\Character{91}}%
\TextAt{13}{134.244}{761.124}{a}%
\TextAt{9}{140.389}{761.124}{\Character{93}}%
\TextAt{21}{143.64}{759.33}{F}%
\TextAt{9}{150.674}{761.124}{\Character{40}}%
\TextAt{13}{155.226}{761.124}{X}%
\TextAt{9}{165.881}{761.124}{\Character{41}}%
\NativeMathAt{173.754}{761.124}{11.955}{\equiv}%
\TextAt{13}{186.374}{761.124}{aX}%
\TextAt{8}{206.495}{761.124}{mod}%
\TextAt{8}{231.924}{761.124}{deg}%
\TextAt{9}{251.487}{761.124}{2}%
\TextAt{8}{257.34}{761.124}{.}%
\TextAt{8}{97.159}{741.089}{A}%
\TextAt{12}{109.391}{741.089}{homomorphism}%
\TextAt{8}{203.649}{741.089}{(over}%
\TextAt{14}{232.81}{741.089}{o}%
\TextAt{3}{238.654}{745.427}{\Character{48}}%
\NativeMathAt{244.77}{741.089}{11.955}{\supseteq}%
\TextAt{14}{257.39}{741.089}{o}%
\TextAt{8}{263.234}{741.089}{)}%
\TextAt{8}{271.247}{741.089}{between}%
\TextAt{8}{315.998}{741.089}{formal}%
\TextAt{14}{352.308}{741.089}{o}%
\TextAt{8}{358.152}{741.089}{-modules}%
\TextAt{13}{407.518}{741.089}{F}%
\TextAt{13}{415.408}{741.089}{\Character{59}}%
\TextAt{13}{420.633}{741.089}{G}%
\TextAt{8}{433.357}{741.089}{is}%
\TextAt{8}{444.681}{741.089}{a}%
\TextAt{8}{453.997}{741.089}{homomorphism}%
\TextAt{13}{79.226}{724.043}{f}%
\TextAt{9}{89.593}{724.043}{\Character{58}}%
\TextAt{13}{96.166}{724.043}{F}%
\NativeArrowAt{108.69}{724.043}{11.955}{19.26}{\longrightarrow}%
\TextAt{13}{131.272}{724.043}{G}%
\TextAt{8}{144.403}{724.043}{of}%
\TextAt{8}{157.733}{724.043}{formal}%
\TextAt{8}{194.484}{724.043}{groups}%
\TextAt{8}{232.274}{724.043}{(over}%
\TextAt{14}{261.875}{724.043}{o}%
\TextAt{3}{267.719}{728.381}{\Character{48}}%
\TextAt{8}{270.514}{724.043}{)}%
\TextAt{8}{278.966}{724.043}{in}%
\TextAt{8}{292.619}{724.043}{the}%
\TextAt{8}{312.776}{724.043}{sense}%
\TextAt{8}{342.807}{724.043}{of}%
\TextAt{8}{356.137}{724.043}{(4.2)}%
\TextAt{8}{384.112}{724.043}{such}%
\TextAt{8}{410.509}{724.043}{that}%
\TextAt{13}{198.547}{695.042}{f}%
\TextAt{9}{205.593}{695.042}{\Character{40}\Character{91}}%
\TextAt{13}{213.397}{695.042}{a}%
\TextAt{9}{219.542}{695.042}{\Character{93}}%
\TextAt{21}{222.794}{693.249}{F}%
\TextAt{9}{229.827}{695.042}{\Character{40}}%
\TextAt{13}{234.38}{695.042}{X}%
\TextAt{9}{245.035}{695.042}{\Character{41}\Character{41}}%
\TextAt{9}{257.433}{695.042}{\Character{61}}%
\TextAt{9}{269.854}{695.042}{\Character{91}}%
\TextAt{13}{273.138}{695.042}{a}%
\TextAt{9}{279.282}{695.042}{\Character{93}}%
\TextAt{21}{282.534}{693.249}{G}%
\TextAt{9}{289.631}{695.042}{\Character{40}}%
\TextAt{13}{294.183}{695.042}{f}%
\TextAt{9}{301.23}{695.042}{\Character{40}}%
\TextAt{13}{305.782}{695.042}{X}%
\TextAt{9}{316.437}{695.042}{\Character{41}\Character{41}}%
\TextAt{8}{349.452}{695.042}{for}%
\TextAt{8}{367.337}{695.042}{all}%
\TextAt{13}{383.596}{695.042}{a}%
\NativeMathAt{393.062}{695.042}{11.955}{\in}%
\TextAt{14}{404.353}{695.042}{o}%
\TextAt{13}{410.197}{695.042}{\Character{58}}%
\TextAt{8}{97.159}{657.074}{Now}%
\TextAt{8}{123.865}{657.074}{let}%
\TextAt{14}{140.816}{657.074}{o}%
\TextAt{9}{150.061}{657.074}{\Character{61}}%
\TextAt{14}{162.565}{657.074}{o}%
\TextAt{21}{168.41}{655.281}{K}%
\TextAt{8}{180.556}{657.074}{be}%
\TextAt{8}{196.527}{657.074}{the}%
\TextAt{8}{216.73}{657.074}{valuation}%
\TextAt{8}{267.849}{657.074}{ring}%
\TextAt{8}{291.961}{657.074}{of}%
\TextAt{8}{305.337}{657.074}{a}%
\TextAt{8}{315.139}{657.074}{local}%
\TextAt{8}{342.826}{657.074}{field}%
\TextAt{13}{368.23}{657.074}{K}%
\TextAt{8}{379.047}{657.074}{,}%
\TextAt{8}{386.254}{657.074}{and}%
\TextAt{8}{409.064}{657.074}{write}%
\TextAt{13}{439.022}{657.074}{q}%
\TextAt{9}{448.041}{657.074}{\Character{61}}%
\TextAt{9}{460.534}{657.074}{\Character{40}}%
\TextAt{14}{465.098}{657.074}{o}%
\TextAt{21}{470.942}{655.281}{K}%
\TextAt{9}{482.545}{657.074}{\Character{58}}%
\TextAt{14}{489.197}{657.074}{p}%
\TextAt{21}{495.18}{655.281}{K}%
\TextAt{9}{503.382}{657.074}{\Character{41}}%
\TextAt{8}{507.935}{657.074}{.}%
\TextAt{8}{516.522}{657.074}{We}%
\TextAt{8}{79.226}{640.028}{consider}%
\TextAt{8}{124.811}{640.028}{the}%
\TextAt{8}{144.968}{640.028}{following}%
\TextAt{8}{194.402}{640.028}{special}%
\TextAt{8}{232.503}{640.028}{formal}%
\TextAt{14}{269.254}{640.028}{o}%
\TextAt{21}{275.098}{638.235}{K}%
\TextAt{8}{283.3}{640.028}{-modules.}%
\TextAt{12}{79.226}{614.016}{Definition}%
\TextAt{12}{142.125}{614.016}{(4.5).}%
\TextAt{8}{177.536}{614.016}{A}%
\TextAt{12}{190.305}{614.016}{Lubin–Tate}%
\TextAt{12}{262.003}{614.016}{module}%
\TextAt{8}{309.13}{614.016}{over}%
\TextAt{14}{334.272}{614.016}{o}%
\TextAt{21}{340.116}{612.223}{K}%
\TextAt{8}{352.312}{614.016}{for}%
\TextAt{8}{370.293}{614.016}{the}%
\TextAt{8}{390.546}{614.016}{prime}%
\TextAt{8}{423.805}{614.016}{element}%
\TextAt{13}{467.143}{614.016}{\Character{25}}%
\TextAt{8}{478.206}{614.016}{is}%
\TextAt{8}{490.065}{614.016}{a}%
\TextAt{8}{499.917}{614.016}{formal}%
\TextAt{14}{79.226}{596.97}{o}%
\TextAt{21}{85.07}{595.177}{K}%
\TextAt{8}{93.273}{596.97}{-module}%
\TextAt{13}{138.464}{596.97}{F}%
\TextAt{8}{151.565}{596.97}{such}%
\TextAt{8}{177.962}{596.97}{that}%
\TextAt{9}{250.695}{579.924}{\Character{91}}%
\TextAt{13}{253.947}{579.924}{\Character{25}}%
\TextAt{9}{261.016}{579.924}{\Character{93}}%
\TextAt{21}{264.268}{578.131}{F}%
\TextAt{9}{271.301}{579.924}{\Character{40}}%
\TextAt{13}{275.853}{579.924}{X}%
\TextAt{9}{286.508}{579.924}{\Character{41}}%
\NativeMathAt{294.382}{579.924}{11.955}{\equiv}%
\TextAt{13}{307.001}{579.924}{X}%
\TextAt{21}{317.656}{584.86}{q}%
\TextAt{8}{325.542}{579.924}{mod}%
\TextAt{13}{350.98}{579.924}{\Character{25}}%
\TextAt{13}{358.034}{579.924}{\Character{58}}%
\TextAt{8}{97.159}{547.436}{This}%
\TextAt{8}{125.15}{547.436}{definition}%
\TextAt{8}{178.451}{547.436}{reflects}%
\TextAt{8}{219.45}{547.436}{once}%
\TextAt{8}{247.381}{547.436}{more}%
\TextAt{8}{277.919}{547.436}{the}%
\TextAt{8}{299.347}{547.436}{dominating}%
\TextAt{8}{362.415}{547.436}{principle}%
\TextAt{8}{411.806}{547.436}{of}%
\TextAt{8}{426.408}{547.436}{class}%
\TextAt{8}{455.117}{547.436}{field}%
\TextAt{8}{481.745}{547.436}{theory,}%
\TextAt{8}{522.357}{547.436}{to}%
\TextAt{8}{79.226}{530.39}{the}%
\TextAt{8}{100.293}{530.39}{effect}%
\TextAt{8}{132.084}{530.39}{that}%
\TextAt{8}{158.363}{530.39}{prime}%
\TextAt{8}{192.437}{530.39}{elements}%
\TextAt{8}{241.204}{530.39}{correspond}%
\TextAt{8}{301.687}{530.39}{to}%
\TextAt{8}{316.908}{530.39}{Frobenius}%
\TextAt{8}{371.688}{530.39}{elements.}%
\TextAt{8}{426.827}{530.39}{In}%
\TextAt{8}{442.359}{530.39}{fact,}%
\TextAt{8}{469.834}{530.39}{if}%
\TextAt{8}{481.468}{530.39}{we}%
\TextAt{8}{499.606}{530.39}{reduce}%
\TextAt{8}{79.226}{513.344}{the}%
\TextAt{8}{99.249}{513.344}{coefficients}%
\TextAt{8}{158.353}{513.344}{of}%
\TextAt{8}{171.549}{513.344}{some}%
\TextAt{8}{200.741}{513.344}{formal}%
\TextAt{14}{237.357}{513.344}{o}%
\TextAt{8}{243.201}{513.344}{-module}%
\TextAt{13}{288.258}{513.344}{F}%
\TextAt{8}{301.225}{513.344}{modulo}%
\TextAt{13}{343.042}{513.344}{\Character{25}}%
\TextAt{8}{350.111}{513.344}{,}%
\TextAt{8}{357.153}{513.344}{we}%
\TextAt{8}{374.246}{513.344}{obtain}%
\TextAt{8}{410.54}{513.344}{a}%
\TextAt{8}{420.161}{513.344}{formal}%
\TextAt{8}{456.777}{513.344}{group}%
\Vector{q 0 G 0 g 0.478 w 417.819 -250.813 m 427.023 -250.813 l S Q}%
\TextAt{13}{489.819}{513.344}{F}%
\TextAt{9}{499.023}{513.344}{\Character{40}}%
\TextAt{13}{503.575}{513.344}{X}%
\TextAt{13}{513.57}{513.344}{\Character{59}}%
\TextAt{13}{518.794}{513.344}{Y}%
\TextAt{9}{528.217}{513.344}{\Character{41}}%
\TextAt{8}{79.226}{496.298}{over}%
\TextAt{8}{104.761}{496.298}{the}%
\TextAt{8}{125.406}{496.298}{residue}%
\TextAt{8}{165.621}{496.298}{class}%
\TextAt{8}{193.547}{496.298}{field}%
\TextAt{18}{219.392}{496.298}{F}%
\TextAt{21}{226.698}{494.505}{q}%
\TextAt{8}{231.263}{496.298}{.}%
\TextAt{8}{241.177}{496.298}{The}%
\TextAt{8}{265.719}{496.298}{reduction}%
\TextAt{8}{318.237}{496.298}{mod}%
\TextAt{13}{345.062}{496.298}{\Character{25}}%
\TextAt{8}{356.518}{496.298}{of}%
\TextAt{9}{370.336}{496.298}{\Character{91}}%
\TextAt{13}{373.588}{496.298}{\Character{25}}%
\TextAt{9}{380.657}{496.298}{\Character{93}}%
\TextAt{21}{383.909}{494.505}{F}%
\TextAt{9}{390.942}{496.298}{\Character{40}}%
\TextAt{13}{395.494}{496.298}{X}%
\TextAt{9}{406.15}{496.298}{\Character{41}}%
\TextAt{8}{415.088}{496.298}{is}%
\TextAt{8}{427.34}{496.298}{an}%
\TextAt{8}{444.088}{496.298}{endomorphism}%
\TextAt{8}{523.337}{496.298}{of}%
\Vector{q 0 G 0 g 0.478 w 7.226 -284.905 m 16.43 -284.905 l S Q}%
\TextAt{13}{79.226}{479.252}{F}%
\TextAt{8}{88.43}{479.252}{.}%
\TextAt{8}{97.446}{479.252}{But}%
\TextAt{8}{120.876}{479.252}{on}%
\TextAt{8}{137.324}{479.252}{the}%
\TextAt{8}{157.67}{479.252}{other}%
\TextAt{8}{188.429}{479.252}{hand,}%
\TextAt{13}{221.184}{479.252}{f}%
\TextAt{9}{228.23}{479.252}{\Character{40}}%
\TextAt{13}{232.782}{479.252}{X}%
\TextAt{9}{243.437}{479.252}{\Character{41}}%
\TextAt{9}{251.614}{479.252}{\Character{61}}%
\TextAt{13}{264.381}{479.252}{X}%
\TextAt{21}{275.036}{483.591}{q}%
\TextAt{8}{283.688}{479.252}{is}%
\TextAt{8}{295.641}{479.252}{clearly}%
\TextAt{8}{333.226}{479.252}{an}%
\TextAt{8}{349.674}{479.252}{endomorphism}%
\TextAt{8}{428.624}{479.252}{of}%
\Vector{q 0 G 0 g 0.478 w 370.144 -284.905 m 379.347 -284.905 l S Q}%
\TextAt{13}{442.144}{479.252}{F}%
\TextAt{8}{451.347}{479.252}{,}%
\TextAt{8}{458.733}{479.252}{its}%
\TextAt{12}{475.241}{479.252}{Frobenius}%
\TextAt{12}{79.226}{462.206}{endomorphism}%
\TextAt{8}{165.686}{462.206}{.}%
\TextAt{8}{178.009}{462.206}{Thus}%
\TextAt{13}{208.949}{462.206}{F}%
\TextAt{8}{223.342}{462.206}{is}%
\TextAt{8}{236.397}{462.206}{a}%
\TextAt{8}{247.444}{462.206}{Lubin–Tate}%
\TextAt{8}{311.656}{462.206}{module}%
\TextAt{8}{354.241}{462.206}{if}%
\TextAt{8}{366.256}{462.206}{the}%
\TextAt{8}{387.704}{462.206}{endomorphism}%
\TextAt{8}{467.756}{462.206}{defined}%
\TextAt{8}{509.361}{462.206}{by}%
\TextAt{8}{526.911}{462.206}{a}%
\TextAt{8}{79.226}{445.16}{prime}%
\TextAt{8}{112.39}{445.16}{element}%
\TextAt{13}{155.632}{445.16}{\Character{25}}%
\TextAt{8}{166.598}{445.16}{gives}%
\TextAt{8}{195.279}{445.16}{via}%
\TextAt{8}{214.467}{445.16}{reduction}%
\TextAt{8}{266.496}{445.16}{the}%
\TextAt{8}{286.652}{445.16}{Frobenius}%
\TextAt{8}{340.522}{445.16}{endomorphism}%
\TextAt{8}{419.283}{445.16}{of}%
\Vector{q 0 G 0 g 0.478 w 360.613 -318.997 m 369.816 -318.997 l S Q}%
\TextAt{13}{432.613}{445.16}{F}%
\TextAt{8}{441.816}{445.16}{.}%
\TextAt{12}{79.226}{425.125}{Example:}%
\TextAt{8}{140.925}{425.125}{The}%
\TextAt{8}{165.846}{425.125}{formal}%
\TextAt{8}{203.464}{425.125}{multiplicative}%
\TextAt{8}{278.155}{425.125}{group}%
\TextAt{18}{312.198}{425.125}{G}%
\TextAt{21}{321.496}{423.332}{m}%
\TextAt{8}{334.25}{425.125}{is}%
\TextAt{8}{346.881}{425.125}{a}%
\TextAt{8}{357.504}{425.125}{formal}%
\TextAt{18}{395.122}{425.125}{Z}%
\TextAt{21}{403.092}{423.332}{p}%
\TextAt{8}{407.853}{425.125}{-module}%
\TextAt{8}{453.911}{425.125}{with}%
\TextAt{8}{481.439}{425.125}{respect}%
\TextAt{8}{522.357}{425.125}{to}%
\TextAt{8}{79.226}{408.079}{the}%
\TextAt{8}{99.383}{408.079}{mapping}%
\TextAt{18}{124.517}{373.106}{Z}%
\TextAt{21}{132.487}{371.313}{p}%
\NativeArrowAt{140.569}{373.106}{11.955}{19.26}{\longrightarrow}%
\TextAt{8}{163.151}{373.106}{End}%
\TextAt{25}{184.12}{371.313}{Z}%
\TextAt{22}{189.434}{370.317}{p}%
\TextAt{9}{194.27}{373.106}{\Character{40}}%
\TextAt{18}{198.823}{373.106}{G}%
\TextAt{21}{208.121}{371.313}{m}%
\TextAt{9}{216.11}{373.106}{\Character{41}}%
\TextAt{13}{220.662}{373.106}{\Character{59}}%
\TextAt{13}{249.797}{373.106}{a}%
\NativeArrowAt{259.283}{373.106}{11.955}{19.26}{\longmapsto}%
\TextAt{9}{281.865}{373.106}{\Character{91}}%
\TextAt{13}{285.116}{373.106}{a}%
\TextAt{9}{291.261}{373.106}{\Character{93}}%
\TextAt{25}{294.513}{371.313}{G}%
\TextAt{22}{300.712}{370.317}{m}%
\TextAt{9}{308.246}{373.106}{\Character{40}}%
\TextAt{13}{312.798}{373.106}{X}%
\TextAt{9}{323.453}{373.106}{\Character{41}}%
\TextAt{9}{331.308}{373.106}{\Character{61}}%
\TextAt{9}{343.729}{373.106}{\Character{40}1}%
\TextAt{9}{356.772}{373.106}{\Character{43}}%
\TextAt{13}{368.575}{373.106}{X}%
\TextAt{9}{379.23}{373.106}{\Character{41}}%
\TextAt{21}{383.783}{378.042}{a}%
\TextAt{16}{391.436}{373.106}{\Character{0}}%
\TextAt{9}{403.391}{373.106}{1}%
\TextAt{9}{412.549}{373.106}{\Character{61}}%
\NativeMathAt{429.39}{388.05}{7.97}{\infty}%
\TextAt{20}{424.99}{384.464}{\Character{88}}%
\TextAt{21}{425.872}{359.358}{\Character{23}}%
\TextAt{17}{430.556}{359.358}{\Character{61}1}%
\TextAt{20}{444.251}{389.963}{\Character{18}}%
\TextAt{13}{453.248}{381.194}{a}%
\TextAt{13}{453.052}{364.905}{\Character{23}}%
\TextAt{20}{459.589}{389.963}{\Character{19}}%
\TextAt{13}{468.39}{373.106}{X}%
\TextAt{21}{479.045}{378.042}{\Character{23}}%
\TextAt{13}{484.227}{373.106}{\Character{58}}%
\TextAt{18}{79.226}{336.794}{G}%
\TextAt{21}{88.525}{335.001}{m}%
\TextAt{8}{100.411}{336.794}{is}%
\TextAt{8}{112.175}{336.794}{a}%
\TextAt{8}{121.93}{336.794}{Lubin–Tate}%
\TextAt{8}{184.85}{336.794}{module}%
\TextAt{8}{226.143}{336.794}{for}%
\TextAt{8}{244.028}{336.794}{the}%
\TextAt{8}{264.185}{336.794}{prime}%
\TextAt{8}{297.348}{336.794}{element}%
\TextAt{13}{340.59}{336.794}{p}%
\TextAt{8}{350.363}{336.794}{because}%
\TextAt{9}{207.94}{307.793}{\Character{91}}%
\TextAt{13}{211.192}{307.793}{p}%
\TextAt{9}{217.067}{307.793}{\Character{93}}%
\TextAt{25}{220.318}{306}{G}%
\TextAt{22}{226.517}{305.003}{m}%
\TextAt{9}{234.052}{307.793}{\Character{40}}%
\TextAt{13}{238.604}{307.793}{X}%
\TextAt{9}{249.259}{307.793}{\Character{41}}%
\TextAt{9}{257.114}{307.793}{\Character{61}}%
\TextAt{9}{269.535}{307.793}{\Character{40}1}%
\TextAt{9}{282.578}{307.793}{\Character{43}}%
\TextAt{13}{294.381}{307.793}{X}%
\TextAt{9}{305.036}{307.793}{\Character{41}}%
\TextAt{21}{309.588}{312.729}{p}%
\TextAt{16}{317.006}{307.793}{\Character{0}}%
\TextAt{9}{328.961}{307.793}{1}%
\NativeMathAt{338.135}{307.793}{11.955}{\equiv}%
\TextAt{13}{350.754}{307.793}{X}%
\TextAt{21}{361.409}{312.729}{p}%
\TextAt{8}{369.491}{307.793}{mod}%
\TextAt{13}{394.929}{307.793}{p\Character{58}}%
\TextAt{8}{97.159}{275.803}{The}%
\TextAt{8}{122.423}{275.803}{following}%
\TextAt{8}{173.068}{275.803}{theorem}%
\TextAt{8}{219.803}{275.803}{gives}%
\TextAt{8}{249.694}{275.803}{a}%
\TextAt{8}{260.66}{275.803}{complete}%
\TextAt{8}{311.293}{275.803}{and}%
\TextAt{8}{335.265}{275.803}{explicit}%
\TextAt{8}{377.769}{275.803}{overall}%
\TextAt{8}{416.387}{275.803}{view}%
\TextAt{8}{444.58}{275.803}{of}%
\TextAt{8}{459.12}{275.803}{the}%
\TextAt{8}{480.487}{275.803}{totality}%
\TextAt{8}{523.337}{275.803}{of}%
\TextAt{8}{79.226}{258.757}{all}%
\TextAt{8}{96.369}{258.757}{Lubin–Tate}%
\TextAt{8}{160.173}{258.757}{modules.}%
\TextAt{8}{213.283}{258.757}{Let}%
\TextAt{13}{235.136}{258.757}{e}%
\TextAt{9}{240.561}{258.757}{\Character{40}}%
\TextAt{13}{245.114}{258.757}{X}%
\TextAt{9}{255.769}{258.757}{\Character{41}}%
\TextAt{13}{260.321}{258.757}{\Character{59}}%
\Vector{q 0 G 0 g 0.478 w 193.565 -508.422 m 198.991 -508.422 l S Q}%
\TextAt{13}{265.565}{258.757}{e}%
\TextAt{9}{270.991}{258.757}{\Character{40}}%
\TextAt{13}{275.543}{258.757}{X}%
\TextAt{9}{286.198}{258.757}{\Character{41}}%
\NativeMathAt{295.577}{258.757}{11.955}{\in}%
\TextAt{14}{308.374}{258.757}{o}%
\TextAt{21}{314.219}{256.964}{K}%
\TextAt{9}{322.421}{258.757}{\Character{91}\Character{91}}%
\TextAt{13}{328.925}{258.757}{X}%
\TextAt{9}{339.58}{258.757}{\Character{93}\Character{93}}%
\TextAt{8}{350.864}{258.757}{be}%
\TextAt{8}{367.672}{258.757}{Lubin–Tate}%
\TextAt{8}{431.476}{258.757}{series}%
\TextAt{8}{463.694}{258.757}{for}%
\TextAt{8}{482.463}{258.757}{the}%
\TextAt{8}{503.503}{258.757}{prime}%
\TextAt{8}{79.226}{241.711}{element}%
\TextAt{13}{122.468}{241.711}{\Character{25}}%
\TextAt{8}{133.435}{241.711}{of}%
\TextAt{13}{146.765}{241.711}{K}%
\TextAt{8}{157.582}{241.711}{,}%
\TextAt{8}{164.732}{241.711}{and}%
\TextAt{8}{187.494}{241.711}{let}%
\TextAt{13}{180.863}{212.71}{F}%
\TextAt{21}{188.441}{210.917}{e}%
\TextAt{9}{192.867}{212.71}{\Character{40}}%
\TextAt{13}{197.42}{212.71}{X}%
\TextAt{13}{207.415}{212.71}{\Character{59}}%
\TextAt{13}{212.639}{212.71}{Y}%
\TextAt{9}{222.062}{212.71}{\Character{41}}%
\NativeMathAt{229.935}{212.71}{11.955}{\in}%
\TextAt{14}{241.226}{212.71}{o}%
\TextAt{21}{247.07}{210.917}{K}%
\TextAt{9}{255.273}{212.71}{\Character{91}\Character{91}}%
\TextAt{13}{261.776}{212.71}{X}%
\TextAt{13}{271.771}{212.71}{\Character{59}}%
\TextAt{13}{276.995}{212.71}{Y}%
\TextAt{9}{286.418}{212.71}{\Character{93}\Character{93}}%
\TextAt{8}{304.877}{212.71}{and}%
\TextAt{9}{335.697}{212.71}{\Character{91}}%
\TextAt{13}{338.949}{212.71}{a}%
\TextAt{9}{345.094}{212.71}{\Character{93}}%
\TextAt{21}{348.345}{210.917}{e\Character{59}}%
\Vector{q 0 G 0 g 0.359 w 282.626 -558.396 m 286.554 -558.396 l S Q}%
\TextAt{21}{354.626}{210.917}{e}%
\TextAt{9}{359.053}{212.71}{\Character{40}}%
\TextAt{13}{363.605}{212.71}{X}%
\TextAt{9}{374.26}{212.71}{\Character{41}}%
\NativeMathAt{382.133}{212.71}{11.955}{\in}%
\TextAt{14}{393.424}{212.71}{o}%
\TextAt{21}{399.268}{210.917}{K}%
\TextAt{9}{407.471}{212.71}{\Character{91}\Character{91}}%
\TextAt{13}{413.974}{212.71}{X}%
\TextAt{9}{424.629}{212.71}{\Character{93}\Character{93}}%
\TextAt{8}{79.226}{183.709}{(}%
\TextAt{13}{83.781}{183.709}{a}%
\NativeMathAt{93.247}{183.709}{11.955}{\in}%
\TextAt{14}{104.538}{183.709}{o}%
\TextAt{21}{110.382}{181.916}{K}%
\TextAt{8}{118.585}{183.709}{)}%
\TextAt{8}{127.037}{183.709}{be}%
\TextAt{8}{142.961}{183.709}{the}%
\TextAt{8}{163.118}{183.709}{power}%
\TextAt{8}{197.262}{183.709}{series}%
\TextAt{8}{228.596}{183.709}{(uniquely}%
\TextAt{8}{280.625}{183.709}{determined}%
\TextAt{8}{341.752}{183.709}{according}%
\TextAt{8}{394.438}{183.709}{to}%
\TextAt{8}{408.749}{183.709}{(2.2))}%
\TextAt{8}{441.279}{183.709}{such}%
\TextAt{8}{467.676}{183.709}{that}%
\TextAt{13}{127.635}{154.708}{F}%
\TextAt{21}{135.213}{152.914}{e}%
\TextAt{9}{139.64}{154.708}{\Character{40}}%
\TextAt{13}{144.192}{154.708}{X}%
\TextAt{13}{154.187}{154.708}{\Character{59}}%
\TextAt{13}{159.411}{154.708}{Y}%
\TextAt{9}{168.834}{154.708}{\Character{41}}%
\NativeMathAt{176.707}{154.708}{11.955}{\equiv}%
\TextAt{13}{189.326}{154.708}{X}%
\TextAt{9}{202.638}{154.708}{\Character{43}}%
\TextAt{13}{214.399}{154.708}{Y}%
\TextAt{8}{227.113}{154.708}{mod}%
\TextAt{8}{252.542}{154.708}{deg}%
\TextAt{9}{272.106}{154.708}{2}%
\TextAt{13}{277.959}{154.708}{\Character{59}}%
\TextAt{13}{329.62}{154.708}{e}%
\TextAt{9}{335.045}{154.708}{\Character{40}}%
\TextAt{13}{339.597}{154.708}{F}%
\TextAt{21}{347.175}{152.914}{e}%
\TextAt{9}{351.602}{154.708}{\Character{40}}%
\TextAt{13}{356.154}{154.708}{X}%
\TextAt{13}{366.149}{154.708}{\Character{59}}%
\TextAt{13}{371.373}{154.708}{Y}%
\TextAt{9}{380.796}{154.708}{\Character{41}\Character{41}}%
\TextAt{9}{393.194}{154.708}{\Character{61}}%
\TextAt{13}{405.647}{154.708}{F}%
\TextAt{21}{413.225}{152.914}{e}%
\TextAt{9}{417.651}{154.708}{\Character{40}}%
\TextAt{13}{422.204}{154.708}{e}%
\TextAt{9}{427.629}{154.708}{\Character{40}}%
\TextAt{13}{432.182}{154.708}{X}%
\TextAt{9}{442.837}{154.708}{\Character{41}}%
\TextAt{13}{447.389}{154.708}{\Character{59}}%
\TextAt{13}{452.613}{154.708}{e}%
\TextAt{9}{458.059}{154.708}{\Character{40}}%
\TextAt{13}{462.611}{154.708}{Y}%
\TextAt{9}{472.004}{154.708}{\Character{41}\Character{41}}%
\TextAt{13}{481.109}{154.708}{\Character{59}}%
\TextAt{9}{130.271}{134.673}{\Character{91}}%
\TextAt{13}{133.523}{134.673}{a}%
\TextAt{9}{139.668}{134.673}{\Character{93}}%
\TextAt{21}{142.919}{132.88}{e\Character{59}}%
\Vector{q 0 G 0 g 0.359 w 77.2 -636.433 m 81.128 -636.433 l S Q}%
\TextAt{21}{149.2}{132.88}{e}%
\TextAt{9}{153.627}{134.673}{\Character{40}}%
\TextAt{13}{158.179}{134.673}{X}%
\TextAt{9}{168.834}{134.673}{\Character{41}}%
\NativeMathAt{176.707}{134.673}{11.955}{\equiv}%
\TextAt{13}{189.326}{134.673}{aX}%
\TextAt{8}{209.447}{134.673}{mod}%
\TextAt{8}{234.876}{134.673}{deg}%
\TextAt{9}{254.44}{134.673}{2}%
\TextAt{13}{260.293}{134.673}{\Character{59}}%
\TextAt{13}{332.256}{134.673}{e}%
\TextAt{9}{337.681}{134.673}{\Character{40}\Character{91}}%
\TextAt{13}{345.485}{134.673}{a}%
\TextAt{9}{351.63}{134.673}{\Character{93}}%
\TextAt{21}{354.882}{132.88}{e\Character{59}}%
\Vector{q 0 G 0 g 0.359 w 289.162 -636.433 m 293.091 -636.433 l S Q}%
\TextAt{21}{361.162}{132.88}{e}%
\TextAt{9}{365.589}{134.673}{\Character{40}}%
\TextAt{13}{370.141}{134.673}{X}%
\TextAt{9}{380.796}{134.673}{\Character{41}\Character{41}}%
\TextAt{9}{393.194}{134.673}{\Character{61}}%
\TextAt{9}{405.615}{134.673}{\Character{91}}%
\TextAt{13}{408.899}{134.673}{a}%
\TextAt{9}{415.044}{134.673}{\Character{93}}%
\TextAt{21}{418.295}{132.88}{e\Character{59}}%
\Vector{q 0 G 0 g 0.359 w 352.576 -636.433 m 356.505 -636.433 l S Q}%
\TextAt{21}{424.576}{132.88}{e}%
\TextAt{9}{429.003}{134.673}{\Character{40}}%
\Vector{q 0 G 0 g 0.478 w 361.555 -632.506 m 366.98 -632.506 l S Q}%
\TextAt{13}{433.555}{134.673}{e}%
\TextAt{9}{438.98}{134.673}{\Character{40}}%
\TextAt{13}{443.533}{134.673}{X}%
\TextAt{9}{454.188}{134.673}{\Character{41}\Character{41}}%
\TextAt{13}{463.293}{134.673}{\Character{58}}%
\TextAt{8}{79.226}{105.672}{If}%
\TextAt{13}{90.918}{105.672}{e}%
\TextAt{9}{96.344}{105.672}{\Character{40}}%
\TextAt{13}{100.896}{105.672}{X}%
\TextAt{9}{111.551}{105.672}{\Character{41}}%
\TextAt{9}{119.406}{105.672}{\Character{61}}%
\Vector{q 0 G 0 g 0.478 w 59.85 -661.507 m 65.275 -661.507 l S Q}%
\TextAt{13}{131.85}{105.672}{e}%
\TextAt{9}{137.275}{105.672}{\Character{40}}%
\TextAt{13}{141.828}{105.672}{X}%
\TextAt{9}{152.483}{105.672}{\Character{41}}%
\TextAt{8}{160.932}{105.672}{we}%
\TextAt{8}{178.16}{105.672}{simply}%
\TextAt{8}{215.615}{105.672}{write}%
\TextAt{9}{245.527}{105.672}{\Character{91}}%
\TextAt{13}{248.779}{105.672}{a}%
\TextAt{9}{254.924}{105.672}{\Character{93}}%
\TextAt{21}{258.175}{103.879}{e\Character{59}}%
\Vector{q 0 G 0 g 0.359 w 192.456 -665.435 m 196.385 -665.435 l S Q}%
\TextAt{21}{264.456}{103.879}{e}%
\TextAt{9}{268.883}{105.672}{\Character{40}}%
\TextAt{13}{273.435}{105.672}{X}%
\TextAt{9}{284.09}{105.672}{\Character{41}}%
\TextAt{9}{291.945}{105.672}{\Character{61}}%
\TextAt{9}{304.366}{105.672}{\Character{91}}%
\TextAt{13}{307.64}{105.672}{a}%
\TextAt{9}{313.785}{105.672}{\Character{93}}%
\TextAt{21}{317.037}{103.879}{e}%
\TextAt{9}{321.464}{105.672}{\Character{40}}%
\TextAt{13}{326.016}{105.672}{X}%
\TextAt{9}{336.671}{105.672}{\Character{41}}%
\TextAt{8}{341.223}{105.672}{.}%
\TextAt{12}{79.226}{79.659}{Theorem}%
\TextAt{12}{136.24}{79.659}{(4.6).}%
\TextAt{8}{297.211}{43.033}{315}%
\Destination{78.23}{809.815}{nav.page.315}%
\Destination{79.36}{562.455}{nav.print.344}%
\Destination{79.36}{628.055}{nav.number.5.4.5}%
\Destination{79.36}{93.655}{nav.number.5.4.6}%
\Hotlink{355.637}{721.855}{25.078}{11.047}{nav.number.5.4.2}%
\Hotlink{408.249}{181.555}{25.078}{11.047}{nav.number.5.2.2}%
}

% Page 328 (body page 316)
\Page{%
\TextAt{2}{79.226}{796.943}{V.}%
\TextAt{2}{94.073}{796.943}{Local}%
\TextAt{2}{123.615}{796.943}{Class}%
\TextAt{2}{152.208}{796.943}{Field}%
\TextAt{2}{179.939}{796.943}{Theory}%
\Vector{q 0 G 0 g 0.299 w 7.226 18.729 m 460.769 18.729 l S Q}%
\TextAt{8}{89.579}{761.124}{(i)}%
\TextAt{8}{107.918}{761.124}{The}%
\TextAt{8}{132.927}{761.124}{Lubin–Tate}%
\TextAt{8}{196.801}{761.124}{modules}%
\TextAt{8}{243.663}{761.124}{for}%
\TextAt{13}{262.502}{761.124}{\Character{25}}%
\TextAt{8}{274.423}{761.124}{are}%
\TextAt{8}{294.888}{761.124}{precisely}%
\TextAt{8}{343.699}{761.124}{the}%
\TextAt{8}{364.81}{761.124}{series}%
\TextAt{13}{397.099}{761.124}{F}%
\TextAt{21}{404.676}{759.33}{e}%
\TextAt{9}{409.103}{761.124}{\Character{40}}%
\TextAt{13}{413.655}{761.124}{X}%
\TextAt{13}{423.649}{761.124}{\Character{59}}%
\TextAt{13}{428.874}{761.124}{Y}%
\TextAt{9}{438.297}{761.124}{\Character{41}}%
\TextAt{8}{442.85}{761.124}{,}%
\TextAt{8}{451.192}{761.124}{with}%
\TextAt{8}{478.806}{761.124}{the}%
\TextAt{8}{499.917}{761.124}{formal}%
\TextAt{14}{107.918}{744.078}{o}%
\TextAt{21}{113.763}{742.284}{K}%
\TextAt{8}{121.965}{744.078}{-module}%
\TextAt{8}{167.156}{744.078}{structure}%
\TextAt{8}{217.296}{744.078}{given}%
\TextAt{8}{247.865}{744.078}{by}%
\TextAt{14}{222.967}{715.077}{o}%
\TextAt{21}{228.811}{713.283}{K}%
\NativeArrowAt{240.335}{715.077}{11.955}{19.26}{\longrightarrow}%
\TextAt{8}{262.917}{715.077}{End}%
\TextAt{27}{283.886}{713.283}{o}%
\TextAt{22}{287.782}{711.944}{K}%
\TextAt{9}{295.323}{715.077}{\Character{40}}%
\TextAt{13}{299.876}{715.077}{F}%
\TextAt{21}{307.454}{713.283}{e}%
\TextAt{9}{311.88}{715.077}{\Character{41}}%
\TextAt{13}{316.432}{715.077}{\Character{59}}%
\TextAt{13}{345.567}{715.077}{a}%
\NativeArrowAt{355.053}{715.077}{11.955}{19.26}{\longmapsto}%
\TextAt{9}{377.635}{715.077}{\Character{91}}%
\TextAt{13}{380.886}{715.077}{a}%
\TextAt{9}{387.031}{715.077}{\Character{93}}%
\TextAt{21}{390.283}{713.283}{e}%
\TextAt{9}{394.709}{715.077}{\Character{40}}%
\TextAt{13}{399.262}{715.077}{X}%
\TextAt{9}{409.917}{715.077}{\Character{41}}%
\TextAt{13}{414.469}{715.077}{\Character{58}}%
\TextAt{8}{86.327}{686.075}{(ii)}%
\TextAt{8}{107.918}{686.075}{For}%
\TextAt{8}{128.933}{686.075}{every}%
\TextAt{13}{159.871}{686.075}{a}%
\NativeMathAt{169.414}{686.075}{11.955}{\in}%
\TextAt{14}{180.782}{686.075}{o}%
\TextAt{21}{186.626}{684.282}{K}%
\TextAt{8}{198.771}{686.075}{the}%
\TextAt{8}{218.973}{686.075}{power}%
\TextAt{8}{253.162}{686.075}{series}%
\TextAt{9}{284.542}{686.075}{\Character{91}}%
\TextAt{13}{287.794}{686.075}{a}%
\TextAt{9}{293.939}{686.075}{\Character{93}}%
\TextAt{21}{297.19}{684.282}{e\Character{59}}%
\Vector{q 0 G 0 g 0.359 w 231.471 -85.031 m 235.399 -85.031 l S Q}%
\TextAt{21}{303.471}{684.282}{e}%
\TextAt{9}{307.898}{686.075}{\Character{40}}%
\TextAt{13}{312.45}{686.075}{X}%
\TextAt{9}{323.105}{686.075}{\Character{41}}%
\TextAt{8}{331.6}{686.075}{is}%
\TextAt{8}{343.409}{686.075}{a}%
\TextAt{8}{353.21}{686.075}{homomorphism}%
\TextAt{9}{435.925}{686.075}{\Character{91}}%
\TextAt{13}{439.177}{686.075}{a}%
\TextAt{9}{445.322}{686.075}{\Character{93}}%
\TextAt{21}{448.573}{684.282}{e\Character{59}}%
\Vector{q 0 G 0 g 0.359 w 382.854 -85.031 m 386.783 -85.031 l S Q}%
\TextAt{21}{454.854}{684.282}{e}%
\TextAt{9}{462.679}{686.075}{\Character{58}}%
\TextAt{13}{469.328}{686.075}{F}%
\Vector{q 0 G 0 g 0.359 w 404.906 -85.031 m 409.333 -85.031 l S Q}%
\TextAt{21}{476.906}{684.282}{e}%
\NativeArrowAt{484.731}{686.075}{11.955}{19.26}{\longrightarrow}%
\TextAt{13}{507.39}{686.075}{F}%
\TextAt{21}{514.968}{684.282}{e}%
\TextAt{8}{523.337}{686.075}{of}%
\TextAt{8}{107.918}{669.029}{formal}%
\TextAt{14}{144.669}{669.029}{o}%
\TextAt{8}{150.513}{669.029}{-modules,}%
\TextAt{8}{203.57}{669.029}{and}%
\TextAt{8}{226.333}{669.029}{it}%
\TextAt{8}{238.037}{669.029}{is}%
\TextAt{8}{249.801}{669.029}{an}%
\TextAt{8}{266.06}{669.029}{isomorphism}%
\TextAt{8}{334.479}{669.029}{if}%
\TextAt{13}{345.203}{669.029}{a}%
\TextAt{8}{355.245}{669.029}{is}%
\TextAt{8}{367.009}{669.029}{a}%
\TextAt{8}{376.764}{669.029}{unit.}%
\TextAt{19}{79.226}{634.051}{Proof.}%
\TextAt{8}{116}{634.051}{If}%
\TextAt{13}{127.982}{634.051}{F}%
\TextAt{8}{141.373}{634.051}{is}%
\TextAt{8}{153.427}{634.051}{any}%
\TextAt{8}{175.834}{634.051}{Lubin–Tate}%
\TextAt{8}{239.044}{634.051}{module,}%
\TextAt{8}{283.952}{634.051}{then}%
\TextAt{13}{310.902}{634.051}{e}%
\TextAt{9}{316.327}{634.051}{\Character{40}}%
\TextAt{13}{320.879}{634.051}{X}%
\TextAt{9}{331.534}{634.051}{\Character{41}}%
\TextAt{9}{339.891}{634.051}{\Character{58}\Character{61}}%
\TextAt{9}{356.042}{634.051}{\Character{91}}%
\TextAt{13}{359.325}{634.051}{\Character{25}}%
\TextAt{9}{366.394}{634.051}{\Character{93}}%
\TextAt{21}{369.646}{632.257}{F}%
\TextAt{9}{376.679}{634.051}{\Character{40}}%
\TextAt{13}{381.231}{634.051}{X}%
\TextAt{9}{391.886}{634.051}{\Character{41}}%
\NativeMathAt{400.254}{634.051}{11.955}{\in}%
\TextAt{16}{412.03}{634.051}{\Character{69}}%
\TextAt{21}{418.349}{632.257}{\Character{25}}%
\TextAt{8}{428.183}{634.051}{and}%
\TextAt{13}{451.235}{634.051}{F}%
\TextAt{9}{464.254}{634.051}{\Character{61}}%
\TextAt{13}{477.174}{634.051}{F}%
\TextAt{21}{484.751}{632.257}{e}%
\TextAt{8}{493.365}{634.051}{because}%
\TextAt{13}{79.226}{617.005}{e}%
\TextAt{9}{84.652}{617.005}{\Character{40}}%
\TextAt{13}{89.204}{617.005}{F}%
\TextAt{9}{98.408}{617.005}{\Character{40}}%
\TextAt{13}{102.96}{617.005}{X}%
\TextAt{13}{112.955}{617.005}{\Character{59}}%
\TextAt{13}{118.179}{617.005}{Y}%
\TextAt{9}{127.602}{617.005}{\Character{41}\Character{41}}%
\TextAt{9}{140.49}{617.005}{\Character{61}}%
\TextAt{13}{153.438}{617.005}{F}%
\TextAt{9}{162.641}{617.005}{\Character{40}}%
\TextAt{13}{167.194}{617.005}{e}%
\TextAt{9}{172.619}{617.005}{\Character{40}}%
\TextAt{13}{177.171}{617.005}{X}%
\TextAt{9}{187.827}{617.005}{\Character{41}}%
\TextAt{13}{192.379}{617.005}{\Character{59}}%
\TextAt{13}{197.603}{617.005}{e}%
\TextAt{9}{203.048}{617.005}{\Character{40}}%
\TextAt{13}{207.601}{617.005}{Y}%
\TextAt{9}{216.994}{617.005}{\Character{41}\Character{41}}%
\TextAt{8}{226.099}{617.005}{,}%
\TextAt{8}{233.609}{617.005}{and}%
\TextAt{8}{256.661}{617.005}{because}%
\TextAt{8}{300.251}{617.005}{of}%
\TextAt{8}{313.87}{617.005}{the}%
\TextAt{8}{334.315}{617.005}{uniqueness}%
\TextAt{8}{393.579}{617.005}{statement}%
\TextAt{8}{448.24}{617.005}{of}%
\TextAt{8}{461.859}{617.005}{(2.2).}%
\TextAt{8}{495.252}{617.005}{For}%
\TextAt{8}{516.51}{617.005}{the}%
\TextAt{8}{79.226}{599.959}{other}%
\TextAt{8}{109.796}{599.959}{claims}%
\TextAt{8}{145.625}{599.959}{of}%
\TextAt{8}{158.955}{599.959}{the}%
\TextAt{8}{179.112}{599.959}{theorem}%
\TextAt{8}{224.637}{599.959}{one}%
\TextAt{8}{246.096}{599.959}{has}%
\TextAt{8}{266.97}{599.959}{to}%
\TextAt{8}{281.28}{599.959}{show}%
\TextAt{8}{310.284}{599.959}{the}%
\TextAt{8}{330.44}{599.959}{following}%
\TextAt{8}{379.875}{599.959}{formulae.}%
\TextAt{13}{234.552}{570.958}{F}%
\TextAt{21}{242.13}{569.164}{e}%
\TextAt{9}{246.556}{570.958}{\Character{40}}%
\TextAt{13}{251.109}{570.958}{X}%
\TextAt{13}{261.104}{570.958}{\Character{59}}%
\TextAt{13}{266.328}{570.958}{Y}%
\TextAt{9}{275.751}{570.958}{\Character{41}}%
\TextAt{9}{283.606}{570.958}{\Character{61}}%
\TextAt{13}{296.049}{570.958}{F}%
\TextAt{21}{303.627}{569.164}{e}%
\TextAt{9}{308.054}{570.958}{\Character{40}}%
\TextAt{13}{312.606}{570.958}{Y}%
\TextAt{13}{320.042}{570.958}{\Character{59}}%
\TextAt{13}{325.279}{570.958}{X}%
\TextAt{9}{335.947}{570.958}{\Character{41}}%
\TextAt{13}{340.5}{570.958}{\Character{59}}%
\TextAt{8}{517.802}{570.958}{(1)}%
\TextAt{13}{201.279}{550.923}{F}%
\TextAt{21}{208.857}{549.13}{e}%
\TextAt{9}{213.283}{550.923}{\Character{40}}%
\TextAt{13}{217.836}{550.923}{X}%
\TextAt{13}{227.831}{550.923}{\Character{59}}%
\TextAt{13}{233.055}{550.923}{F}%
\TextAt{21}{240.662}{549.13}{e}%
\TextAt{9}{245.089}{550.923}{\Character{40}}%
\TextAt{13}{249.641}{550.923}{Y}%
\TextAt{13}{257.077}{550.923}{\Character{59}}%
\TextAt{13}{262.313}{550.923}{Z}%
\TextAt{9}{271.198}{550.923}{\Character{41}\Character{41}}%
\TextAt{9}{283.596}{550.923}{\Character{61}}%
\TextAt{13}{296.049}{550.923}{F}%
\TextAt{21}{303.627}{549.13}{e}%
\TextAt{9}{308.054}{550.923}{\Character{40}}%
\TextAt{13}{312.606}{550.923}{F}%
\TextAt{21}{320.184}{549.13}{e}%
\TextAt{9}{324.61}{550.923}{\Character{40}}%
\TextAt{13}{329.162}{550.923}{X}%
\TextAt{13}{339.156}{550.923}{\Character{59}}%
\TextAt{13}{344.381}{550.923}{Y}%
\TextAt{9}{353.805}{550.923}{\Character{41}}%
\TextAt{13}{358.357}{550.923}{\Character{59}}%
\TextAt{13}{363.581}{550.923}{Z}%
\TextAt{9}{372.472}{550.923}{\Character{41}}%
\TextAt{13}{377.024}{550.923}{\Character{59}}%
\TextAt{8}{517.802}{550.923}{(2)}%
\TextAt{9}{202.092}{530.888}{\Character{91}}%
\TextAt{13}{205.343}{530.888}{a}%
\TextAt{9}{211.488}{530.888}{\Character{93}}%
\TextAt{21}{214.74}{529.095}{e\Character{59}}%
\Vector{q 0 G 0 g 0.359 w 149.021 -240.218 m 152.949 -240.218 l S Q}%
\TextAt{21}{221.021}{529.095}{e}%
\TextAt{9}{225.447}{530.888}{\Character{40}}%
\TextAt{13}{230}{530.888}{F}%
\Vector{q 0 G 0 g 0.359 w 165.577 -240.218 m 170.004 -240.218 l S Q}%
\TextAt{21}{237.577}{529.095}{e}%
\TextAt{9}{242.004}{530.888}{\Character{40}}%
\TextAt{13}{246.556}{530.888}{X}%
\TextAt{13}{256.55}{530.888}{\Character{59}}%
\TextAt{13}{261.775}{530.888}{Y}%
\TextAt{9}{271.198}{530.888}{\Character{41}\Character{41}}%
\TextAt{9}{283.596}{530.888}{\Character{61}}%
\TextAt{13}{296.049}{530.888}{F}%
\TextAt{21}{303.627}{529.095}{e}%
\TextAt{9}{308.054}{530.888}{\Character{40}\Character{91}}%
\TextAt{13}{315.858}{530.888}{a}%
\TextAt{9}{322.002}{530.888}{\Character{93}}%
\TextAt{21}{325.254}{529.095}{e\Character{59}}%
\Vector{q 0 G 0 g 0.359 w 259.535 -240.218 m 263.463 -240.218 l S Q}%
\TextAt{21}{331.535}{529.095}{e}%
\TextAt{9}{335.961}{530.888}{\Character{40}}%
\TextAt{13}{340.514}{530.888}{X}%
\TextAt{9}{351.169}{530.888}{\Character{41}}%
\TextAt{13}{355.721}{530.888}{\Character{59}}%
\TextAt{9}{360.965}{530.888}{\Character{91}}%
\TextAt{13}{364.217}{530.888}{a}%
\TextAt{9}{370.362}{530.888}{\Character{93}}%
\TextAt{21}{373.614}{529.095}{e\Character{59}}%
\Vector{q 0 G 0 g 0.359 w 307.894 -240.218 m 311.823 -240.218 l S Q}%
\TextAt{21}{379.894}{529.095}{e}%
\TextAt{9}{384.321}{530.888}{\Character{40}}%
\TextAt{13}{388.873}{530.888}{Y}%
\TextAt{9}{398.266}{530.888}{\Character{41}\Character{41}}%
\TextAt{13}{407.371}{530.888}{\Character{59}}%
\TextAt{8}{517.802}{530.888}{(3)}%
\TextAt{9}{217.793}{510.853}{\Character{91}}%
\TextAt{13}{221.044}{510.853}{a}%
\TextAt{9}{229.846}{510.853}{\Character{43}}%
\TextAt{13}{241.607}{510.853}{b}%
\TextAt{9}{246.584}{510.853}{\Character{93}}%
\TextAt{21}{249.836}{509.06}{e\Character{59}}%
\Vector{q 0 G 0 g 0.359 w 184.117 -260.253 m 188.045 -260.253 l S Q}%
\TextAt{21}{256.117}{509.06}{e}%
\TextAt{9}{260.543}{510.853}{\Character{40}}%
\TextAt{13}{265.096}{510.853}{X}%
\TextAt{9}{275.751}{510.853}{\Character{41}}%
\TextAt{9}{283.606}{510.853}{\Character{61}}%
\TextAt{13}{296.049}{510.853}{F}%
\TextAt{21}{303.627}{509.06}{e}%
\TextAt{9}{308.054}{510.853}{\Character{40}\Character{91}}%
\TextAt{13}{315.858}{510.853}{a}%
\TextAt{9}{322.002}{510.853}{\Character{93}}%
\TextAt{21}{325.254}{509.06}{e\Character{59}}%
\Vector{q 0 G 0 g 0.359 w 259.535 -260.253 m 263.463 -260.253 l S Q}%
\TextAt{21}{331.535}{509.06}{e}%
\TextAt{9}{335.961}{510.853}{\Character{40}}%
\TextAt{13}{340.514}{510.853}{X}%
\TextAt{9}{351.169}{510.853}{\Character{41}}%
\TextAt{13}{355.721}{510.853}{\Character{59}}%
\TextAt{9}{360.965}{510.853}{\Character{91}}%
\TextAt{13}{364.217}{510.853}{b}%
\TextAt{9}{369.194}{510.853}{\Character{93}}%
\TextAt{21}{372.446}{509.06}{e\Character{59}}%
\Vector{q 0 G 0 g 0.359 w 306.726 -260.253 m 310.655 -260.253 l S Q}%
\TextAt{21}{378.726}{509.06}{e}%
\TextAt{9}{383.153}{510.853}{\Character{40}}%
\TextAt{13}{387.705}{510.853}{X}%
\TextAt{9}{398.36}{510.853}{\Character{41}\Character{41}}%
\TextAt{13}{407.465}{510.853}{\Character{59}}%
\TextAt{8}{517.802}{510.853}{(4)}%
\TextAt{9}{232.211}{490.819}{\Character{91}}%
\TextAt{13}{235.462}{490.819}{ab}%
\TextAt{9}{246.584}{490.819}{\Character{93}}%
\TextAt{21}{249.836}{489.025}{e\Character{59}}%
\Vector{q 0 G 0 g 0.359 w 184.117 -280.288 m 188.045 -280.288 l S Q}%
\TextAt{21}{256.117}{489.025}{e}%
\TextAt{9}{260.543}{490.819}{\Character{40}}%
\TextAt{13}{265.096}{490.819}{X}%
\TextAt{9}{275.751}{490.819}{\Character{41}}%
\TextAt{9}{283.606}{490.819}{\Character{61}}%
\TextAt{9}{296.027}{490.819}{\Character{91}}%
\TextAt{13}{299.301}{490.819}{a}%
\TextAt{9}{305.446}{490.819}{\Character{93}}%
\TextAt{21}{308.698}{489.025}{e\Character{59}}%
\Vector{q 0 G 0 g 0.359 w 242.978 -280.288 m 246.907 -280.288 l S Q}%
\TextAt{21}{314.978}{489.025}{e}%
\TextAt{9}{319.405}{490.819}{\Character{40}\Character{91}}%
\TextAt{13}{327.209}{490.819}{b}%
\TextAt{9}{332.186}{490.819}{\Character{93}}%
\Vector{q 0 G 0 g 0.359 w 263.437 -280.288 m 267.864 -280.288 l S Q}%
\TextAt{21}{335.437}{489.025}{e}%
\TextAt{9}{339.864}{490.819}{\Character{40}}%
\TextAt{13}{344.416}{490.819}{X}%
\TextAt{9}{355.071}{490.819}{\Character{41}\Character{41}}%
\TextAt{13}{364.176}{490.819}{\Character{59}}%
\TextAt{8}{517.802}{490.819}{(5)}%
\TextAt{9}{242.544}{470.784}{\Character{91}}%
\TextAt{13}{245.796}{470.784}{\Character{25}}%
\TextAt{9}{252.865}{470.784}{\Character{93}}%
\TextAt{21}{256.117}{468.991}{e}%
\TextAt{9}{260.543}{470.784}{\Character{40}}%
\TextAt{13}{265.096}{470.784}{X}%
\TextAt{9}{275.751}{470.784}{\Character{41}}%
\TextAt{9}{283.606}{470.784}{\Character{61}}%
\TextAt{13}{296.049}{470.784}{e}%
\TextAt{9}{301.475}{470.784}{\Character{40}}%
\TextAt{13}{306.027}{470.784}{X}%
\TextAt{9}{316.682}{470.784}{\Character{41}}%
\TextAt{13}{321.234}{470.784}{\Character{58}}%
\TextAt{8}{517.802}{470.784}{(6)}%
\TextAt{8}{79.226}{441.783}{(1)}%
\TextAt{8}{98.091}{441.783}{and}%
\TextAt{8}{120.854}{441.783}{(2)}%
\TextAt{8}{139.719}{441.783}{show}%
\TextAt{8}{168.723}{441.783}{that}%
\TextAt{13}{194.091}{441.783}{F}%
\TextAt{21}{201.669}{439.989}{e}%
\TextAt{8}{209.993}{441.783}{is}%
\TextAt{8}{221.757}{441.783}{a}%
\TextAt{8}{231.512}{441.783}{formal}%
\TextAt{8}{268.263}{441.783}{group.}%
\TextAt{8}{305.989}{441.783}{(3),}%
\TextAt{8}{328.106}{441.783}{(4),}%
\TextAt{8}{350.223}{441.783}{and}%
\TextAt{8}{372.986}{441.783}{(5)}%
\TextAt{8}{391.851}{441.783}{show}%
\TextAt{8}{420.854}{441.783}{that}%
\TextAt{14}{218.501}{412.782}{o}%
\TextAt{21}{224.345}{410.988}{K}%
\NativeArrowAt{235.868}{412.782}{11.955}{19.26}{\longrightarrow}%
\TextAt{8}{258.451}{412.782}{End}%
\TextAt{27}{279.42}{410.988}{o}%
\TextAt{22}{283.316}{409.649}{K}%
\TextAt{9}{290.857}{412.782}{\Character{40}}%
\TextAt{13}{295.41}{412.782}{F}%
\TextAt{21}{302.987}{410.988}{e}%
\TextAt{9}{307.414}{412.782}{\Character{41}}%
\TextAt{13}{311.966}{412.782}{\Character{59}}%
\TextAt{13}{341.101}{412.782}{a}%
\NativeArrowAt{350.586}{412.782}{11.955}{19.26}{\longmapsto}%
\TextAt{9}{373.168}{412.782}{\Character{91}}%
\TextAt{13}{376.42}{412.782}{a}%
\TextAt{9}{382.565}{412.782}{\Character{93}}%
\TextAt{21}{385.817}{410.988}{e}%
\TextAt{13}{390.243}{412.782}{\Character{59}}%
\TextAt{8}{79.226}{383.78}{is}%
\TextAt{8}{92.599}{383.78}{a}%
\TextAt{8}{103.963}{383.78}{homomorphism}%
\TextAt{8}{188.243}{383.78}{of}%
\TextAt{8}{203.182}{383.78}{rings,}%
\TextAt{8}{237.125}{383.78}{i.e.,}%
\TextAt{8}{261.241}{383.78}{that}%
\TextAt{13}{288.219}{383.78}{F}%
\TextAt{21}{295.797}{381.987}{e}%
\TextAt{8}{305.73}{383.78}{is}%
\TextAt{8}{319.103}{383.78}{a}%
\TextAt{8}{330.467}{383.78}{formal}%
\TextAt{14}{368.827}{383.78}{o}%
\TextAt{21}{374.671}{381.987}{K}%
\TextAt{8}{382.873}{383.78}{-module,}%
\TextAt{8}{433.327}{383.78}{and}%
\TextAt{8}{457.699}{383.78}{that}%
\TextAt{9}{484.677}{383.78}{\Character{91}}%
\TextAt{13}{487.928}{383.78}{a}%
\TextAt{9}{494.073}{383.78}{\Character{93}}%
\TextAt{21}{497.325}{381.987}{e\Character{59}}%
\Vector{q 0 G 0 g 0.359 w 431.606 -387.326 m 435.534 -387.326 l S Q}%
\TextAt{21}{503.606}{381.987}{e}%
\TextAt{8}{513.539}{383.78}{is}%
\TextAt{8}{526.911}{383.78}{a}%
\TextAt{8}{79.226}{366.734}{homomorphism}%
\TextAt{8}{163.444}{366.734}{of}%
\TextAt{8}{178.322}{366.734}{formal}%
\TextAt{14}{216.62}{366.734}{o}%
\TextAt{21}{222.464}{364.941}{K}%
\TextAt{8}{230.667}{366.734}{-modules}%
\TextAt{8}{282.02}{366.734}{from}%
\TextAt{13}{311.209}{366.734}{F}%
\Vector{q 0 G 0 g 0.359 w 246.786 -404.372 m 251.213 -404.372 l S Q}%
\TextAt{21}{318.786}{364.941}{e}%
\TextAt{8}{328.658}{366.734}{to}%
\TextAt{13}{344.517}{366.734}{F}%
\TextAt{21}{352.094}{364.941}{e}%
\TextAt{8}{356.521}{366.734}{.}%
\TextAt{8}{369.613}{366.734}{Finally,}%
\TextAt{8}{413.654}{366.734}{(6)}%
\TextAt{8}{434.067}{366.734}{shows}%
\TextAt{8}{469.233}{366.734}{that}%
\TextAt{13}{496.15}{366.734}{F}%
\TextAt{21}{503.728}{364.941}{e}%
\TextAt{8}{513.6}{366.734}{is}%
\TextAt{8}{526.911}{366.734}{a}%
\TextAt{8}{79.226}{349.688}{Lubin–Tate}%
\TextAt{8}{142.146}{349.688}{module.}%
\TextAt{8}{97.159}{329.654}{The}%
\TextAt{8}{120.362}{329.654}{proofs}%
\TextAt{8}{154.696}{329.654}{of}%
\TextAt{8}{167.175}{329.654}{these}%
\TextAt{8}{196.296}{329.654}{formulae}%
\TextAt{8}{243.577}{329.654}{all}%
\TextAt{8}{258.986}{329.654}{follow}%
\TextAt{8}{291.957}{329.654}{the}%
\TextAt{8}{311.263}{329.654}{same}%
\TextAt{8}{339.738}{329.654}{pattern.}%
\TextAt{8}{385.634}{329.654}{One}%
\TextAt{8}{409.494}{329.654}{checks}%
\TextAt{8}{444.796}{329.654}{that}%
\TextAt{8}{469.315}{329.654}{both}%
\TextAt{8}{496.105}{329.654}{sides}%
\TextAt{8}{523.337}{329.654}{of}%
\TextAt{8}{79.226}{312.608}{each}%
\TextAt{8}{105.654}{312.608}{formula}%
\TextAt{8}{148.676}{312.608}{are}%
\TextAt{8}{168.277}{312.608}{solutions}%
\TextAt{8}{217.276}{312.608}{of}%
\TextAt{8}{230.697}{312.608}{the}%
\TextAt{8}{250.944}{312.608}{same}%
\TextAt{8}{280.361}{312.608}{problem}%
\TextAt{8}{325.976}{312.608}{of}%
\TextAt{8}{339.397}{312.608}{(2.2),}%
\TextAt{8}{370.737}{312.608}{and}%
\TextAt{8}{393.591}{312.608}{then}%
\TextAt{8}{420.341}{312.608}{deduces}%
\TextAt{8}{464.056}{312.608}{their}%
\TextAt{8}{492.11}{312.608}{equality}%
\TextAt{8}{79.226}{295.562}{from}%
\TextAt{8}{107.932}{295.562}{the}%
\TextAt{8}{129.154}{295.562}{uniqueness}%
\TextAt{8}{189.195}{295.562}{statement.}%
\TextAt{8}{251.315}{295.562}{In}%
\TextAt{8}{267.002}{295.562}{(6)}%
\TextAt{8}{286.933}{295.562}{for}%
\TextAt{8}{305.884}{295.562}{instance,}%
\TextAt{8}{356.053}{295.562}{both}%
\TextAt{8}{384.759}{295.562}{power}%
\TextAt{8}{419.969}{295.562}{series}%
\TextAt{8}{452.369}{295.562}{commence}%
\TextAt{8}{510.007}{295.562}{with}%
\TextAt{8}{79.226}{278.516}{the}%
\TextAt{8}{100.237}{278.516}{linear}%
\TextAt{8}{133.609}{278.516}{form}%
\TextAt{13}{162.103}{278.516}{\Character{25}}%
\TextAt{13}{169.157}{278.516}{X}%
\TextAt{8}{184.579}{278.516}{and}%
\TextAt{8}{208.196}{278.516}{satisfy}%
\TextAt{8}{245.597}{278.516}{the}%
\TextAt{8}{266.607}{278.516}{condition}%
\TextAt{13}{318.845}{278.516}{e}%
\TextAt{9}{324.27}{278.516}{\Character{40}\Character{91}}%
\TextAt{13}{332.074}{278.516}{\Character{25}}%
\TextAt{9}{339.143}{278.516}{\Character{93}}%
\TextAt{21}{342.395}{276.723}{e}%
\TextAt{9}{346.821}{278.516}{\Character{40}}%
\TextAt{13}{351.374}{278.516}{X}%
\TextAt{9}{362.029}{278.516}{\Character{41}\Character{41}}%
\TextAt{9}{375.885}{278.516}{\Character{61}}%
\TextAt{9}{389.765}{278.516}{\Character{91}}%
\TextAt{13}{393.042}{278.516}{\Character{25}}%
\TextAt{9}{400.112}{278.516}{\Character{93}}%
\TextAt{21}{403.363}{276.723}{e}%
\TextAt{9}{407.79}{278.516}{\Character{40}}%
\TextAt{13}{412.342}{278.516}{e}%
\TextAt{9}{417.767}{278.516}{\Character{40}}%
\TextAt{13}{422.32}{278.516}{X}%
\TextAt{9}{432.975}{278.516}{\Character{41}\Character{41}}%
\TextAt{8}{442.08}{278.516}{,}%
\TextAt{8}{450.296}{278.516}{resp.}%
\TextAt{13}{479.173}{278.516}{e}%
\TextAt{9}{484.599}{278.516}{\Character{40}}%
\TextAt{13}{489.151}{278.516}{e}%
\TextAt{9}{494.577}{278.516}{\Character{40}}%
\TextAt{13}{499.129}{278.516}{X}%
\TextAt{9}{509.784}{278.516}{\Character{41}\Character{41}}%
\TextAt{9}{523.64}{278.516}{\Character{61}}%
\TextAt{13}{79.226}{261.47}{e}%
\TextAt{9}{84.652}{261.47}{\Character{40}}%
\TextAt{13}{89.204}{261.47}{e}%
\TextAt{9}{94.629}{261.47}{\Character{40}}%
\TextAt{13}{99.182}{261.47}{X}%
\TextAt{9}{109.837}{261.47}{\Character{41}\Character{41}}%
\TextAt{8}{118.941}{261.47}{.}%
\Vector{q 0 G 0 g 0.399 w 452.334 -512.53 m 452.334 -504.46 l S Q}%
\Vector{q 0 G 0 g 0.399 w 452.534 -504.66 m 459.707 -504.66 l S Q}%
\Vector{q 0 G 0 g 0.399 w 452.534 -512.331 m 459.707 -512.331 l S Q}%
\Vector{q 0 G 0 g 0.399 w 459.906 -512.53 m 459.906 -504.46 l S Q}%
\TextAt{12}{79.226}{226.491}{Exercise}%
\TextAt{12}{132.427}{226.491}{V.4.1.}%
\ErrataAt{169}{226.491}{366}{If $\mathfrak o$ has characteristic zero, $\operatorname{End}_{\mathfrak o}(G_a)$ consists of all $aX$ with $a\in\mathfrak o$.}%
\TextAt{12}{79.226}{194.501}{Exercise}%
\TextAt{12}{134.705}{194.501}{V.4.2.}%
\TextAt{8}{175.432}{194.501}{Let}%
\TextAt{13}{198.382}{194.501}{R}%
\TextAt{8}{213.268}{194.501}{be}%
\TextAt{8}{231.173}{194.501}{a}%
\TextAt{8}{242.908}{194.501}{commutative}%
\TextAt{18}{314.815}{194.501}{Q}%
\TextAt{8}{324.113}{194.501}{-algebra.}%
\TextAt{8}{379.485}{194.501}{Then}%
\TextAt{8}{412.023}{194.501}{for}%
\TextAt{8}{431.888}{194.501}{every}%
\TextAt{8}{464.761}{194.501}{formal}%
\TextAt{8}{503.491}{194.501}{group}%
\TextAt{13}{79.226}{177.455}{F}%
\TextAt{9}{88.43}{177.455}{\Character{40}}%
\TextAt{13}{92.982}{177.455}{X}%
\TextAt{13}{102.977}{177.455}{\Character{59}}%
\TextAt{13}{108.201}{177.455}{Y}%
\TextAt{9}{117.624}{177.455}{\Character{41}}%
\TextAt{8}{125.362}{177.455}{over}%
\TextAt{13}{149.697}{177.455}{R}%
\TextAt{8}{158.705}{177.455}{,}%
\TextAt{8}{165.285}{177.455}{there}%
\TextAt{8}{194.486}{177.455}{exists}%
\TextAt{8}{226.089}{177.455}{a}%
\TextAt{8}{235.133}{177.455}{unique}%
\TextAt{8}{272.463}{177.455}{isomorphism}%
\TextAt{8}{340.171}{177.455}{log}%
\TextAt{21}{355.138}{174.506}{F}%
\TextAt{9}{365.493}{177.455}{\Character{58}}%
\TextAt{13}{372.065}{177.455}{F}%
\NativeMathAt{390.927}{184.232}{7.97}{\sim}%
\NativeArrowAt{384.59}{177.455}{11.955}{19.26}{\longrightarrow}%
\TextAt{18}{407.172}{177.455}{G}%
\TextAt{21}{416.47}{175.662}{a}%
\TextAt{8}{421.466}{177.455}{,}%
\TextAt{8}{428.046}{177.455}{such}%
\TextAt{8}{453.732}{177.455}{that}%
\TextAt{8}{478.389}{177.455}{log}%
\TextAt{21}{493.357}{174.506}{F}%
\TextAt{9}{500.39}{177.455}{\Character{40}}%
\TextAt{13}{504.943}{177.455}{X}%
\TextAt{9}{515.598}{177.455}{\Character{41}}%
\NativeMathAt{523.471}{177.455}{11.955}{\equiv}%
\TextAt{13}{79.226}{160.409}{X}%
\TextAt{8}{93.202}{160.409}{mod}%
\TextAt{8}{118.631}{160.409}{deg}%
\TextAt{9}{138.195}{160.409}{2}%
\TextAt{8}{144.048}{160.409}{,}%
\TextAt{8}{151.197}{160.409}{the}%
\TextAt{12}{171.353}{160.409}{logarithm}%
\TextAt{8}{232.157}{160.409}{of}%
\TextAt{13}{245.487}{160.409}{F}%
\TextAt{8}{254.691}{160.409}{.}%
\ErrataAt{79.19}{140.352}{455.81}{\textbf{Hint:} Put $F_1=\partial_Y F$. Differentiating associativity gives $F_1(F(X,Y),0)=F_1(X,Y)F_1(Y,0)$.}%
\ErrataAt{79.19}{123.306}{455.81}{Let $\psi(X)=1/F_1(X,0)=1+\sum_{n\ge1}a_nX^n$. Then $\psi(F(X,Y))F_1(X,Y)=\psi(Y)$.}%
\ErrataAt{79.19}{106.26}{455.81}{The antiderivative $\log_F(X)=X+\sum_{n\ge1}\dfrac{a_n}{n+1}X^{n+1}$ has the required property.}%
\TextAt{12}{79.226}{74.293}{Exercise}%
\TextAt{12}{132.427}{74.293}{V.4.3.}%
\TextAt{8}{171.173}{74.293}{log}%
\TextAt{25}{186.141}{71.344}{G}%
\TextAt{22}{192.34}{70.347}{m}%
\TextAt{9}{199.874}{74.293}{\Character{40}}%
\TextAt{13}{204.427}{74.293}{X}%
\TextAt{9}{215.082}{74.293}{\Character{41}}%
\TextAt{9}{222.937}{74.293}{\Character{61}}%
\TextAt{20}{235.381}{83.259}{\Character{80}}%
\NativeMathAt{248}{80.104}{7.97}{\infty}%
\TextAt{21}{248}{70.806}{n}%
\TextAt{17}{253.138}{70.806}{\Character{61}1}%
\TextAt{9}{264.457}{74.293}{\Character{40}}%
\TextAt{16}{269.009}{74.293}{\Character{0}}%
\TextAt{9}{278.308}{74.293}{1\Character{41}}%
\TextAt{21}{288.713}{78.631}{n}%
\TextAt{17}{293.851}{78.631}{\Character{43}1}%
\TextAt{13}{305.17}{74.293}{X}%
\TextAt{21}{315.825}{78.631}{n}%
\TextAt{9}{321.462}{74.293}{\Character{47}}%
\TextAt{13}{327.315}{74.293}{n}%
\TextAt{9}{337.623}{74.293}{\Character{61}}%
\TextAt{8}{350.048}{74.293}{log}%
\TextAt{9}{365.016}{74.293}{\Character{40}1}%
\TextAt{9}{378.059}{74.293}{\Character{43}}%
\TextAt{13}{389.84}{74.293}{X}%
\TextAt{9}{400.495}{74.293}{\Character{41}}%
\TextAt{8}{405.047}{74.293}{.}%
\TextAt{8}{297.211}{43.033}{316}%
\Destination{522.36}{590.995}{nav.AMS.68}%
\Destination{522.36}{570.955}{nav.AMS.69}%
\Destination{522.36}{550.925}{nav.AMS.70}%
\Destination{522.36}{530.885}{nav.AMS.71}%
\Destination{522.36}{510.855}{nav.AMS.72}%
\Destination{522.36}{490.815}{nav.AMS.73}%
\Destination{79.23}{773.075}{nav.Item.117}%
\Destination{79.23}{699.985}{nav.Item.118}%
\Destination{78.23}{809.815}{nav.page.316}%
\Destination{79.226}{240.455}{nav.ex.5.4.1}%
\Destination{79.226}{208.455}{nav.ex.5.4.2}%
\Destination{79.226}{88.255}{nav.ex.5.4.3}%
\Destination{79.36}{649.055}{nav.print.345}%
\Hotlink{461.359}{614.855}{25.078}{11.047}{nav.number.5.2.2}%
\Hotlink{338.897}{310.455}{25.078}{11.047}{nav.number.5.2.2}%
\Destination{79.226}{243.491}{errata.127}%
\Destination{79.226}{157.352}{errata.128}%
}

% Page 329 (body page 317)
\Page{%
\TextAt{2}{79.226}{796.943}{V.}%
\TextAt{2}{94.073}{796.943}{Local}%
\TextAt{2}{123.615}{796.943}{Class}%
\TextAt{2}{152.208}{796.943}{Field}%
\TextAt{2}{179.939}{796.943}{Theory}%
\Vector{q 0 G 0 g 0.299 w 7.226 18.729 m 460.769 18.729 l S Q}%
\TextAt{12}{79.226}{761.124}{Exercise}%
\TextAt{12}{133.518}{761.124}{V.4.4.}%
\TextAt{8}{173.213}{761.124}{Let}%
\TextAt{13}{195.131}{761.124}{\Character{25}}%
\TextAt{8}{207.046}{761.124}{be}%
\TextAt{8}{223.919}{761.124}{a}%
\TextAt{8}{234.623}{761.124}{prime}%
\TextAt{8}{268.735}{761.124}{element}%
\TextAt{8}{312.926}{761.124}{of}%
\TextAt{8}{327.204}{761.124}{the}%
\TextAt{8}{348.309}{761.124}{local}%
\TextAt{8}{376.898}{761.124}{field}%
\TextAt{13}{403.204}{761.124}{K}%
\TextAt{8}{414.022}{761.124}{,}%
\TextAt{8}{422.356}{761.124}{and}%
\TextAt{8}{446.068}{761.124}{let}%
\TextAt{13}{463.921}{761.124}{f}%
\TextAt{9}{470.967}{761.124}{\Character{40}}%
\TextAt{13}{475.52}{761.124}{X}%
\TextAt{9}{486.175}{761.124}{\Character{41}}%
\TextAt{9}{495.644}{761.124}{\Character{61}}%
\TextAt{13}{509.706}{761.124}{X}%
\TextAt{9}{523.665}{761.124}{\Character{43}}%
\TextAt{13}{79.226}{744.078}{\Character{25}}%
\TextAt{3}{86.295}{748.416}{\Character{0}}%
\TextAt{17}{92.882}{748.416}{1}%
\TextAt{13}{97.614}{744.078}{X}%
\TextAt{21}{108.269}{748.416}{q}%
\TextAt{9}{115.491}{744.078}{\Character{43}}%
\TextAt{13}{127.252}{744.078}{\Character{25}}%
\TextAt{3}{134.321}{748.416}{\Character{0}}%
\TextAt{17}{140.908}{748.416}{2}%
\TextAt{13}{145.64}{744.078}{X}%
\TextAt{21}{156.295}{748.416}{q}%
\TextAt{24}{160.362}{751.229}{2}%
\TextAt{9}{167.668}{744.078}{\Character{43}}%
\TextAt{16}{179.429}{744.078}{\Character{1}}%
\TextAt{16}{184.725}{744.078}{\Character{1}}%
\TextAt{16}{190.033}{744.078}{\Character{1}}%
\TextAt{8}{195.369}{744.078}{.}%
\TextAt{8}{203.818}{744.078}{Then}%
\TextAt{13}{129.521}{716.815}{F}%
\TextAt{9}{138.725}{716.815}{\Character{40}}%
\TextAt{13}{143.277}{716.815}{X}%
\TextAt{13}{153.272}{716.815}{\Character{59}}%
\TextAt{13}{158.496}{716.815}{Y}%
\TextAt{9}{167.919}{716.815}{\Character{41}}%
\TextAt{9}{175.774}{716.815}{\Character{61}}%
\TextAt{13}{188.218}{716.815}{f}%
\TextAt{3}{195.264}{721.751}{\Character{0}}%
\TextAt{17}{201.851}{721.751}{1}%
\TextAt{9}{206.583}{716.815}{\Character{40}}%
\TextAt{13}{211.136}{716.815}{f}%
\TextAt{9}{218.182}{716.815}{\Character{40}}%
\TextAt{13}{222.734}{716.815}{X}%
\TextAt{9}{233.39}{716.815}{\Character{41}}%
\TextAt{9}{240.587}{716.815}{\Character{43}}%
\TextAt{13}{252.36}{716.815}{f}%
\TextAt{9}{259.406}{716.815}{\Character{40}}%
\TextAt{13}{263.959}{716.815}{Y}%
\TextAt{9}{273.352}{716.815}{\Character{41}\Character{41}}%
\TextAt{13}{282.456}{716.815}{\Character{59}}%
\TextAt{9}{311.611}{716.815}{\Character{91}}%
\TextAt{13}{314.863}{716.815}{a}%
\TextAt{9}{321.008}{716.815}{\Character{93}}%
\TextAt{21}{324.259}{715.022}{F}%
\TextAt{9}{331.293}{716.815}{\Character{40}}%
\TextAt{13}{335.845}{716.815}{X}%
\TextAt{9}{346.5}{716.815}{\Character{41}}%
\TextAt{9}{354.355}{716.815}{\Character{61}}%
\TextAt{13}{366.799}{716.815}{f}%
\TextAt{3}{373.845}{721.751}{\Character{0}}%
\TextAt{17}{380.432}{721.751}{1}%
\TextAt{9}{385.164}{716.815}{\Character{40}}%
\TextAt{13}{389.716}{716.815}{af}%
\TextAt{9}{402.908}{716.815}{\Character{40}}%
\TextAt{13}{407.46}{716.815}{X}%
\TextAt{9}{418.115}{716.815}{\Character{41}\Character{41}}%
\TextAt{13}{427.22}{716.815}{\Character{59}}%
\TextAt{13}{444.4}{716.815}{a}%
\NativeMathAt{453.885}{716.815}{11.955}{\in}%
\TextAt{14}{465.176}{716.815}{o}%
\TextAt{21}{471.02}{715.022}{K}%
\TextAt{13}{479.223}{716.815}{\Character{59}}%
\TextAt{8}{79.226}{689.552}{defines}%
\TextAt{8}{117.65}{689.552}{a}%
\TextAt{8}{127.405}{689.552}{Lubin–Tate}%
\TextAt{8}{190.326}{689.552}{module}%
\TextAt{8}{231.619}{689.552}{with}%
\TextAt{8}{258.279}{689.552}{logarithm}%
\TextAt{8}{311.623}{689.552}{log}%
\TextAt{21}{326.591}{686.603}{F}%
\TextAt{9}{336.945}{689.552}{\Character{61}}%
\TextAt{13}{349.37}{689.552}{f}%
\TextAt{8}{356.417}{689.552}{.}%
\TextAt{12}{79.226}{658.556}{Exercise}%
\TextAt{12}{132.343}{658.556}{V.4.5.}%
\TextAt{8}{171.017}{658.556}{Two}%
\TextAt{8}{196.302}{658.556}{Lubin–Tate}%
\TextAt{8}{259.149}{658.556}{modules}%
\TextAt{8}{304.985}{658.556}{over}%
\TextAt{8}{329.958}{658.556}{the}%
\TextAt{8}{350.042}{658.556}{valuation}%
\TextAt{8}{401.042}{658.556}{ring}%
\TextAt{14}{425.035}{658.556}{o}%
\TextAt{21}{430.88}{656.763}{K}%
\TextAt{8}{442.907}{658.556}{of}%
\TextAt{8}{456.165}{658.556}{a}%
\TextAt{8}{465.848}{658.556}{local}%
\TextAt{8}{493.416}{658.556}{field}%
\TextAt{13}{518.7}{658.556}{K}%
\TextAt{8}{529.518}{658.556}{,}%
\TextAt{8}{79.226}{641.51}{but}%
\TextAt{8}{100.686}{641.51}{for}%
\TextAt{8}{118.571}{641.51}{different}%
\TextAt{8}{164.741}{641.51}{prime}%
\TextAt{8}{197.905}{641.51}{elements}%
\TextAt{13}{245.762}{641.51}{\Character{25}}%
\TextAt{8}{256.728}{641.51}{and}%
\Vector{q 0 G 0 g 0.478 w 207.491 -125.669 m 214.56 -125.669 l S Q}%
\TextAt{13}{279.491}{641.51}{\Character{25}}%
\TextAt{8}{286.56}{641.51}{,}%
\TextAt{8}{293.709}{641.51}{are}%
\TextAt{8}{313.22}{641.51}{never}%
\TextAt{8}{344.435}{641.51}{isomorphic.}%
\TextAt{12}{79.226}{610.513}{Exercise}%
\TextAt{12}{133.747}{610.513}{V.4.6.}%
\TextAt{8}{173.641}{610.513}{Two}%
\TextAt{8}{200.145}{610.513}{Lubin–Tate}%
\TextAt{8}{264.213}{610.513}{modules}%
\TextAt{13}{311.268}{610.513}{F}%
\TextAt{21}{318.846}{608.72}{e}%
\TextAt{8}{328.317}{610.513}{and}%
\TextAt{13}{352.227}{610.513}{F}%
\Vector{q 0 G 0 g 0.359 w 287.805 -160.593 m 292.232 -160.593 l S Q}%
\TextAt{21}{359.805}{608.72}{e}%
\TextAt{8}{369.276}{610.513}{for}%
\TextAt{8}{388.309}{610.513}{prime}%
\TextAt{8}{422.62}{610.513}{elements}%
\TextAt{13}{471.624}{610.513}{\Character{25}}%
\TextAt{8}{483.738}{610.513}{and}%
\Vector{q 0 G 0 g 0.478 w 435.648 -156.666 m 442.717 -156.666 l S Q}%
\TextAt{13}{507.648}{610.513}{\Character{25}}%
\TextAt{8}{519.762}{610.513}{al-}%
\TextAt{8}{79.226}{593.467}{ways}%
\TextAt{8}{107.7}{593.467}{become}%
\TextAt{8}{149.756}{593.467}{isomorphic}%
\TextAt{8}{209.122}{593.467}{over}%
\TextAt{14}{234.285}{593.467}{o}%
\TextAt{30}{242.099}{591.829}{\Character{98}}%
\TextAt{21}{240.129}{589.814}{K}%
\TextAt{8}{248.332}{593.467}{,}%
\TextAt{8}{255.626}{593.467}{where}%
\TextAt{20}{292.291}{596.489}{\Character{98}}%
\TextAt{13}{289.552}{593.467}{K}%
\TextAt{8}{304.384}{593.467}{is}%
\TextAt{8}{316.264}{593.467}{the}%
\TextAt{8}{336.537}{593.467}{completion}%
\TextAt{8}{396.49}{593.467}{of}%
\TextAt{8}{409.936}{593.467}{the}%
\TextAt{8}{430.209}{593.467}{maximal}%
\TextAt{8}{478.134}{593.467}{unramified}%
\TextAt{8}{79.226}{576.421}{extension}%
\TextAt{20}{133.73}{579.443}{\Character{101}}%
\TextAt{13}{130.992}{576.421}{K}%
\TextAt{16}{145.131}{576.421}{\Character{106}}%
\TextAt{13}{151.772}{576.421}{K}%
\TextAt{8}{162.59}{576.421}{.}%
\TextAt{12}{79.226}{556.387}{Hint:}%
\TextAt{8}{114.765}{556.387}{The}%
\TextAt{8}{138.819}{556.387}{power}%
\TextAt{8}{172.963}{556.387}{series}%
\TextAt{13}{204.297}{556.387}{\Character{18}}%
\TextAt{8}{213.975}{556.387}{of}%
\TextAt{8}{227.305}{556.387}{(2.3)}%
\TextAt{8}{255.28}{556.387}{yields}%
\TextAt{8}{288.181}{556.387}{an}%
\TextAt{8}{304.44}{556.387}{isomorphism}%
\TextAt{13}{372.859}{556.387}{\Character{18}}%
\TextAt{9}{381.96}{556.387}{\Character{58}}%
\TextAt{13}{388.533}{556.387}{F}%
\Vector{q 0 G 0 g 0.359 w 324.111 -214.72 m 328.537 -214.72 l S Q}%
\TextAt{21}{396.111}{554.593}{e}%
\NativeArrowAt{403.858}{556.387}{11.955}{19.26}{\longrightarrow}%
\TextAt{13}{426.44}{556.387}{F}%
\TextAt{21}{434.018}{554.593}{e}%
\TextAt{8}{438.444}{556.387}{.}%
\TextAt{15}{79.226}{504.277}{§}%
\TextAt{15}{93.785}{504.277}{5}%
\TextAt{15}{120.693}{504.277}{Generalized}%
\TextAt{15}{227.102}{504.277}{Cyclotomic}%
\TextAt{15}{328.845}{504.277}{Theory}%
\TextAt{8}{79.226}{472.391}{Formal}%
\TextAt{8}{119.369}{472.391}{groups}%
\TextAt{8}{157.468}{472.391}{are}%
\TextAt{8}{177.287}{472.391}{relevant}%
\TextAt{8}{221.83}{472.391}{for}%
\TextAt{8}{240.023}{472.391}{local}%
\TextAt{8}{267.972}{472.391}{class}%
\TextAt{8}{295.718}{472.391}{field}%
\TextAt{8}{321.383}{472.391}{theory}%
\TextAt{8}{358.442}{472.391}{in}%
\TextAt{8}{372.403}{472.391}{that}%
\TextAt{8}{398.08}{472.391}{they}%
\TextAt{8}{424.726}{472.391}{allow}%
\TextAt{8}{455.281}{472.391}{us}%
\TextAt{8}{470.605}{472.391}{to}%
\TextAt{8}{485.224}{472.391}{construct}%
\TextAt{8}{79.226}{455.345}{a}%
\TextAt{8}{90.288}{455.345}{perfect}%
\TextAt{8}{130.604}{455.345}{analogue}%
\TextAt{8}{180.699}{455.345}{of}%
\TextAt{8}{195.336}{455.345}{the}%
\TextAt{8}{216.798}{455.345}{theory}%
\TextAt{8}{254.855}{455.345}{of}%
\TextAt{8}{269.491}{455.345}{the}%
\TextAt{13}{290.954}{455.345}{p}%
\TextAt{21}{296.829}{459.683}{n}%
\TextAt{8}{302.465}{455.345}{-th}%
\TextAt{8}{322.625}{455.345}{cyclotomic}%
\TextAt{8}{382.141}{455.345}{field}%
\TextAt{18}{408.804}{455.345}{Q}%
\TextAt{21}{418.103}{453.551}{p}%
\TextAt{9}{422.864}{455.345}{\Character{40}}%
\TextAt{13}{427.416}{455.345}{\Character{16}}%
\TextAt{9}{433.432}{455.345}{\Character{41}}%
\TextAt{8}{443.188}{455.345}{over}%
\TextAt{18}{469.54}{455.345}{Q}%
\TextAt{21}{478.839}{453.551}{p}%
\TextAt{8}{483.6}{455.345}{,}%
\TextAt{8}{492.382}{455.345}{with}%
\TextAt{8}{520.348}{455.345}{its}%
\TextAt{8}{79.226}{438.299}{fundamental}%
\TextAt{8}{146.868}{438.299}{isomorphism}%
\TextAt{13}{233.706}{411.036}{G}%
\TextAt{9}{242.939}{411.036}{\Character{40}}%
\TextAt{18}{247.492}{411.036}{Q}%
\TextAt{21}{256.79}{409.243}{p}%
\TextAt{9}{261.551}{411.036}{\Character{40}}%
\TextAt{13}{266.103}{411.036}{\Character{16}}%
\TextAt{9}{272.119}{411.036}{\Character{41}}%
\TextAt{16}{279.992}{411.036}{\Character{106}}%
\TextAt{18}{286.634}{411.036}{Q}%
\TextAt{21}{295.932}{409.243}{p}%
\TextAt{9}{300.693}{411.036}{\Character{41}}%
\NativeMathAt{314.904}{417.813}{7.97}{\sim}%
\NativeArrowAt{308.566}{411.036}{11.955}{19.26}{\longrightarrow}%
\TextAt{9}{331.148}{411.036}{\Character{40}}%
\TextAt{18}{335.701}{411.036}{Z}%
\TextAt{9}{343.671}{411.036}{\Character{47}}%
\TextAt{13}{349.524}{411.036}{p}%
\TextAt{21}{355.399}{415.972}{n}%
\TextAt{18}{361.035}{411.036}{Z}%
\TextAt{9}{369.005}{411.036}{\Character{41}}%
\TextAt{3}{373.558}{415.972}{\Character{3}}%
\ErrataAt{79.22}{383.755}{453.64}{(see chap. II, \hyperlink{nav.number.2.7.13}{(7.13)}), replacing $\mathbb Q_p$ by an arbitrary local ground field $K$. The formal}%
\TextAt{8}{79.226}{366.727}{groups}%
\TextAt{8}{116.268}{366.727}{furnish}%
\TextAt{8}{154.923}{366.727}{a}%
\TextAt{8}{163.93}{366.727}{generalization}%
\TextAt{8}{237.985}{366.727}{of}%
\TextAt{8}{250.566}{366.727}{the}%
\TextAt{8}{269.974}{366.727}{notion}%
\TextAt{8}{305.652}{366.727}{of}%
\TextAt{13}{318.234}{366.727}{p}%
\TextAt{21}{324.109}{371.066}{n}%
\TextAt{8}{329.745}{366.727}{-th}%
\TextAt{8}{347.85}{366.727}{root}%
\TextAt{8}{372.147}{366.727}{of}%
\TextAt{8}{384.728}{366.727}{unity,}%
\TextAt{8}{416.97}{366.727}{and}%
\TextAt{8}{438.984}{366.727}{provide}%
\TextAt{8}{479.863}{366.727}{an}%
\TextAt{8}{495.374}{366.727}{explicit}%
\TextAt{8}{79.226}{349.682}{version}%
\TextAt{8}{118.965}{349.682}{of}%
\TextAt{8}{132.295}{349.682}{the}%
\TextAt{8}{152.452}{349.682}{local}%
\TextAt{8}{180.092}{349.682}{reciprocity}%
\TextAt{8}{238.302}{349.682}{law}%
\TextAt{8}{259.438}{349.682}{in}%
\TextAt{8}{273.091}{349.682}{the}%
\TextAt{8}{293.248}{349.682}{corresponding}%
\TextAt{8}{368.434}{349.682}{extensions.}%
\TextAt{8}{97.159}{329.647}{A}%
\TextAt{8}{110.122}{329.647}{formal}%
\TextAt{14}{147.164}{329.647}{o}%
\TextAt{21}{153.008}{327.853}{K}%
\TextAt{8}{161.211}{329.647}{-module}%
\TextAt{8}{206.693}{329.647}{gives}%
\TextAt{8}{235.664}{329.647}{rise}%
\TextAt{8}{257.475}{329.647}{to}%
\TextAt{8}{272.076}{329.647}{an}%
\TextAt{8}{288.627}{329.647}{ordinary}%
\TextAt{14}{336.081}{329.647}{o}%
\TextAt{21}{341.925}{327.853}{K}%
\TextAt{8}{350.128}{329.647}{-module}%
\TextAt{8}{395.61}{329.647}{if}%
\TextAt{8}{406.625}{329.647}{we}%
\TextAt{8}{424.144}{329.647}{read}%
\TextAt{8}{450.449}{329.647}{the}%
\TextAt{8}{470.897}{329.647}{power}%
\TextAt{8}{505.332}{329.647}{series}%
\TextAt{8}{79.226}{312.601}{over}%
\TextAt{8}{103.875}{312.601}{a}%
\TextAt{8}{113.233}{312.601}{domain}%
\TextAt{8}{154.463}{312.601}{in}%
\TextAt{8}{167.718}{312.601}{which}%
\TextAt{8}{200.807}{312.601}{they}%
\TextAt{8}{226.747}{312.601}{converge.}%
\TextAt{8}{278.974}{312.601}{We}%
\TextAt{8}{298.721}{312.601}{now}%
\TextAt{8}{322.712}{312.601}{choose}%
\TextAt{8}{359.447}{312.601}{for}%
\TextAt{8}{376.934}{312.601}{this}%
\TextAt{8}{399.359}{312.601}{the}%
\TextAt{8}{419.118}{312.601}{maximal}%
\TextAt{8}{466.529}{312.601}{ideal}%
\Vector{q 0 G 0 g 0.478 w 422.095 -453.462 m 428.078 -453.462 l S Q}%
\TextAt{14}{494.095}{312.601}{p}%
\TextAt{8}{503.578}{312.601}{of}%
\TextAt{8}{516.51}{312.601}{the}%
\TextAt{8}{79.226}{295.555}{valuation}%
\TextAt{8}{130.752}{295.555}{ring}%
\TextAt{8}{155.271}{295.555}{of}%
\TextAt{8}{169.055}{295.555}{the}%
\TextAt{8}{189.664}{295.555}{algebraic}%
\TextAt{8}{239.552}{295.555}{closure}%
\Vector{q 0 G 0 g 0.478 w 207.087 -468.602 m 217.905 -468.602 l S Q}%
\TextAt{13}{279.087}{295.555}{K}%
\TextAt{8}{294.256}{295.555}{of}%
\TextAt{8}{308.039}{295.555}{the}%
\TextAt{8}{328.649}{295.555}{given}%
\TextAt{8}{359.672}{295.555}{local}%
\TextAt{8}{387.766}{295.555}{field}%
\TextAt{13}{413.576}{295.555}{K}%
\TextAt{8}{424.394}{295.555}{.}%
\TextAt{8}{434.202}{295.555}{If}%
\TextAt{13}{446.348}{295.555}{G}%
\TextAt{9}{455.581}{295.555}{\Character{40}}%
\TextAt{13}{460.134}{295.555}{X}%
\TextAt{17}{469.849}{293.762}{1}%
\TextAt{13}{474.581}{295.555}{\Character{59}}%
\TextAt{13}{479.805}{295.555}{\Character{58}}%
\TextAt{13}{485.042}{295.555}{\Character{58}}%
\TextAt{13}{490.278}{295.555}{\Character{58}}%
\TextAt{13}{495.503}{295.555}{\Character{59}}%
\TextAt{13}{500.739}{295.555}{X}%
\TextAt{21}{510.517}{293.762}{n}%
\TextAt{9}{516.154}{295.555}{\Character{41}}%
\NativeMathAt{524.799}{295.555}{11.955}{\in}%
\TextAt{14}{79.226}{278.509}{o}%
\TextAt{21}{85.07}{276.716}{K}%
\TextAt{9}{93.273}{278.509}{\Character{91}\Character{91}}%
\TextAt{13}{99.776}{278.509}{X}%
\TextAt{17}{109.492}{276.716}{1}%
\TextAt{13}{114.224}{278.509}{\Character{59}}%
\TextAt{13}{119.448}{278.509}{\Character{58}}%
\TextAt{13}{124.685}{278.509}{\Character{58}}%
\TextAt{13}{129.909}{278.509}{\Character{58}}%
\TextAt{13}{135.146}{278.509}{\Character{59}}%
\TextAt{13}{140.382}{278.509}{X}%
\TextAt{21}{150.16}{276.716}{n}%
\TextAt{9}{155.796}{278.509}{\Character{93}\Character{93}}%
\TextAt{8}{166.46}{278.509}{is}%
\TextAt{8}{178.486}{278.509}{a}%
\TextAt{8}{188.504}{278.509}{power}%
\TextAt{8}{222.911}{278.509}{series}%
\TextAt{8}{254.508}{278.509}{with}%
\TextAt{8}{281.431}{278.509}{constant}%
\TextAt{8}{328.917}{278.509}{coefficient}%
\TextAt{9}{383.803}{278.509}{0}%
\TextAt{8}{389.656}{278.509}{,}%
\TextAt{8}{397.133}{278.509}{and}%
\TextAt{8}{420.159}{278.509}{if}%
\TextAt{13}{431.145}{278.509}{x}%
\TextAt{17}{437.797}{276.716}{1}%
\TextAt{13}{442.53}{278.509}{\Character{59}}%
\TextAt{13}{447.754}{278.509}{\Character{58}}%
\TextAt{13}{452.991}{278.509}{\Character{58}}%
\TextAt{13}{458.215}{278.509}{\Character{58}}%
\TextAt{13}{463.452}{278.509}{\Character{59}}%
\TextAt{13}{468.688}{278.509}{x}%
\TextAt{21}{475.403}{276.716}{n}%
\NativeMathAt{484.807}{278.509}{11.955}{\in}%
\Vector{q 0 G 0 g 0.478 w 424.546 -487.554 m 430.529 -487.554 l S Q}%
\TextAt{14}{496.546}{278.509}{p}%
\TextAt{8}{502.529}{278.509}{,}%
\TextAt{8}{510.007}{278.509}{then}%
\TextAt{8}{79.226}{261.463}{the}%
\TextAt{8}{100.352}{261.463}{series}%
\TextAt{13}{132.656}{261.463}{G}%
\TextAt{9}{141.889}{261.463}{\Character{40}}%
\TextAt{13}{146.442}{261.463}{x}%
\TextAt{17}{153.094}{259.67}{1}%
\TextAt{13}{157.826}{261.463}{\Character{59}}%
\TextAt{13}{163.05}{261.463}{\Character{58}}%
\TextAt{13}{168.287}{261.463}{\Character{58}}%
\TextAt{13}{173.511}{261.463}{\Character{58}}%
\TextAt{13}{178.748}{261.463}{\Character{59}}%
\TextAt{13}{183.984}{261.463}{x}%
\TextAt{21}{190.699}{259.67}{n}%
\TextAt{9}{196.335}{261.463}{\Character{41}}%
\TextAt{8}{205.754}{261.463}{converges}%
\TextAt{8}{259.147}{261.463}{in}%
\TextAt{8}{273.769}{261.463}{the}%
\TextAt{8}{294.894}{261.463}{complete}%
\TextAt{8}{345.286}{261.463}{field}%
\TextAt{13}{371.613}{261.463}{K}%
\TextAt{9}{382.43}{261.463}{\Character{40}}%
\TextAt{13}{386.983}{261.463}{x}%
\TextAt{17}{393.635}{259.67}{1}%
\TextAt{13}{398.367}{261.463}{\Character{59}}%
\TextAt{13}{403.591}{261.463}{\Character{58}}%
\TextAt{13}{408.828}{261.463}{\Character{58}}%
\TextAt{13}{414.052}{261.463}{\Character{58}}%
\TextAt{13}{419.289}{261.463}{\Character{59}}%
\TextAt{13}{424.525}{261.463}{x}%
\TextAt{21}{431.24}{259.67}{n}%
\TextAt{9}{436.876}{261.463}{\Character{41}}%
\TextAt{8}{446.295}{261.463}{to}%
\TextAt{8}{461.575}{261.463}{an}%
\TextAt{8}{478.803}{261.463}{element}%
\TextAt{8}{523.014}{261.463}{in}%
\Vector{q 0 G 0 g 0.478 w 7.226 -521.645 m 13.209 -521.645 l S Q}%
\TextAt{14}{79.226}{244.417}{p}%
\TextAt{8}{85.209}{244.417}{.}%
\TextAt{8}{95.774}{244.417}{From}%
\TextAt{8}{127.204}{244.417}{the}%
\TextAt{8}{148.066}{244.417}{definition}%
\TextAt{8}{200.8}{244.417}{of}%
\TextAt{8}{214.835}{244.417}{the}%
\TextAt{8}{235.697}{244.417}{formal}%
\TextAt{14}{273.153}{244.417}{o}%
\TextAt{8}{278.997}{244.417}{-modules}%
\TextAt{8}{329.508}{244.417}{and}%
\TextAt{8}{352.976}{244.417}{their}%
\TextAt{8}{381.644}{244.417}{homomorphisms}%
\TextAt{8}{469.634}{244.417}{we}%
\TextAt{8}{487.567}{244.417}{therefore}%
\TextAt{8}{79.226}{227.371}{obtain}%
\TextAt{8}{115.654}{227.371}{immediately}%
\TextAt{8}{182.316}{227.371}{the}%
\TextAt{12}{79.226}{201.855}{Proposition}%
\TextAt{12}{151.937}{201.855}{(5.1).}%
\TextAt{8}{187.253}{201.855}{Let}%
\TextAt{13}{208.222}{201.855}{F}%
\TextAt{8}{221.323}{201.855}{be}%
\TextAt{8}{237.248}{201.855}{a}%
\TextAt{8}{247.003}{201.855}{formal}%
\TextAt{14}{283.753}{201.855}{o}%
\TextAt{21}{289.598}{200.062}{K}%
\TextAt{8}{297.8}{201.855}{-module.}%
\TextAt{8}{347.542}{201.855}{Then}%
\TextAt{8}{378.099}{201.855}{the}%
\TextAt{8}{398.256}{201.855}{set}%
\Vector{q 0 G 0 g 0.478 w 344.523 -564.207 m 350.506 -564.207 l S Q}%
\TextAt{14}{416.523}{201.855}{p}%
\TextAt{8}{426.403}{201.855}{with}%
\TextAt{8}{453.064}{201.855}{the}%
\TextAt{8}{473.22}{201.855}{operations}%
\TextAt{13}{203.512}{174.593}{x}%
\TextAt{9}{212.821}{174.593}{\Character{43}}%
\TextAt{21}{221.926}{172.799}{F}%
\TextAt{13}{231.616}{174.593}{y}%
\TextAt{9}{241.073}{174.593}{\Character{61}}%
\TextAt{13}{253.499}{174.593}{F}%
\TextAt{9}{262.702}{174.593}{\Character{40}}%
\TextAt{13}{267.255}{174.593}{x\Character{59}}%
\TextAt{13}{279.127}{174.593}{y}%
\TextAt{9}{285.288}{174.593}{\Character{41}}%
\TextAt{8}{301.795}{174.593}{and}%
\TextAt{13}{332.616}{174.593}{a}%
\TextAt{16}{341.417}{174.593}{\Character{1}}%
\TextAt{13}{347.395}{174.593}{x}%
\TextAt{9}{357.368}{174.593}{\Character{61}}%
\TextAt{9}{369.778}{174.593}{\Character{91}}%
\TextAt{13}{373.045}{174.593}{a}%
\TextAt{9}{379.19}{174.593}{\Character{93}}%
\TextAt{21}{382.441}{172.799}{F}%
\TextAt{9}{389.475}{174.593}{\Character{40}}%
\TextAt{13}{394.027}{174.593}{x}%
\TextAt{9}{400.679}{174.593}{\Character{41}}%
\TextAt{13}{405.232}{174.593}{\Character{59}}%
\TextAt{13}{79.226}{147.33}{x\Character{59}}%
\TextAt{13}{91.098}{147.33}{y}%
\NativeMathAt{100.58}{147.33}{11.955}{\in}%
\Vector{q 0 G 0 g 0.478 w 39.871 -618.732 m 45.854 -618.732 l S Q}%
\TextAt{14}{111.871}{147.33}{p}%
\TextAt{8}{117.854}{147.33}{,}%
\TextAt{13}{125.003}{147.33}{a}%
\NativeMathAt{134.469}{147.33}{11.955}{\in}%
\TextAt{14}{145.76}{147.33}{o}%
\TextAt{21}{151.604}{145.537}{K}%
\TextAt{8}{159.807}{147.33}{,}%
\TextAt{8}{166.956}{147.33}{is}%
\TextAt{8}{178.72}{147.33}{an}%
\TextAt{14}{194.979}{147.33}{o}%
\TextAt{21}{200.823}{145.537}{K}%
\TextAt{8}{209.026}{147.33}{-module}%
\TextAt{8}{254.216}{147.33}{in}%
\TextAt{8}{267.869}{147.33}{the}%
\TextAt{8}{288.026}{147.33}{usual}%
\TextAt{8}{318.655}{147.33}{sense.}%
\TextAt{8}{353.237}{147.33}{We}%
\TextAt{8}{373.381}{147.33}{denote}%
\TextAt{8}{411.1}{147.33}{it}%
\TextAt{8}{422.804}{147.33}{by}%
\Vector{q 0 G 0 g 0.478 w 367.063 -618.732 m 373.046 -618.732 l S Q}%
\TextAt{14}{439.063}{147.33}{p}%
\TextAt{21}{445.046}{144.572}{F}%
\TextAt{8}{452.08}{147.33}{.}%
\TextAt{8}{97.159}{121.814}{If}%
\TextAt{13}{108.851}{121.814}{f}%
\TextAt{9}{119.218}{121.814}{\Character{58}}%
\TextAt{13}{125.791}{121.814}{F}%
\NativeArrowAt{138.315}{121.814}{11.955}{19.26}{\longrightarrow}%
\TextAt{13}{160.897}{121.814}{G}%
\TextAt{8}{174.028}{121.814}{is}%
\TextAt{8}{185.792}{121.814}{a}%
\TextAt{8}{195.548}{121.814}{homomorphism}%
\TextAt{8}{278.218}{121.814}{(isomorphism)}%
\TextAt{8}{355.747}{121.814}{of}%
\TextAt{8}{369.077}{121.814}{formal}%
\TextAt{14}{405.827}{121.814}{o}%
\TextAt{21}{411.671}{120.021}{K}%
\TextAt{8}{419.874}{121.814}{-modules,}%
\TextAt{8}{472.931}{121.814}{then}%
\TextAt{13}{227.646}{94.552}{f}%
\TextAt{9}{238.013}{94.552}{\Character{58}}%
\Vector{q 0 G 0 g 0.478 w 172.586 -671.511 m 178.569 -671.511 l S Q}%
\TextAt{14}{244.586}{94.552}{p}%
\TextAt{21}{250.569}{91.793}{F}%
\NativeArrowAt{260.923}{94.552}{11.955}{19.26}{\longrightarrow}%
\Vector{q 0 G 0 g 0.478 w 211.505 -671.511 m 217.488 -671.511 l S Q}%
\TextAt{14}{283.505}{94.552}{p}%
\TextAt{21}{289.488}{91.793}{G}%
\TextAt{13}{296.585}{94.552}{\Character{59}}%
\TextAt{13}{325.72}{94.552}{x}%
\NativeArrowAt{335.713}{94.552}{11.955}{19.26}{\longmapsto}%
\TextAt{13}{358.295}{94.552}{f}%
\TextAt{9}{365.341}{94.552}{\Character{40}}%
\TextAt{13}{369.893}{94.552}{x}%
\TextAt{9}{376.545}{94.552}{\Character{41}}%
\TextAt{13}{381.098}{94.552}{\Character{59}}%
\TextAt{8}{297.211}{43.033}{317}%
\Destination{78.23}{809.815}{nav.page.317}%
\Destination{79.23}{530.125}{nav.section.5.5}%
\Destination{79.226}{775.155}{nav.ex.5.4.4}%
\Destination{79.226}{672.555}{nav.ex.5.4.5}%
\Destination{79.226}{624.555}{nav.ex.5.4.6}%
\Destination{79.36}{673.555}{nav.print.346}%
\Destination{79.36}{215.855}{nav.number.5.5.1}%
\Destination{79.23}{530.125}{outline.53}%
%
\Hotlink{226.805}{554.155}{25.078}{11.047}{nav.number.5.2.3}%
%
\Destination{79.226}{397.755}{errata.41}%
}

% Page 330 (body page 318)
\Page{%
\TextAt{2}{79.226}{796.943}{V.}%
\TextAt{2}{94.073}{796.943}{Local}%
\TextAt{2}{123.615}{796.943}{Class}%
\TextAt{2}{152.208}{796.943}{Field}%
\TextAt{2}{179.939}{796.943}{Theory}%
\Vector{q 0 G 0 g 0.299 w 7.226 18.729 m 460.769 18.729 l S Q}%
\TextAt{8}{79.226}{761.124}{is}%
\TextAt{8}{90.99}{761.124}{a}%
\TextAt{8}{100.745}{761.124}{homomorphism}%
\TextAt{8}{183.415}{761.124}{(isomorphism)}%
\TextAt{8}{260.945}{761.124}{of}%
\TextAt{8}{274.275}{761.124}{ordinary}%
\TextAt{14}{321.438}{761.124}{o}%
\TextAt{21}{327.282}{759.33}{K}%
\TextAt{8}{335.485}{761.124}{-modules.}%
\TextAt{8}{97.159}{741.089}{The}%
\TextAt{8}{121.737}{741.089}{operations}%
\TextAt{8}{179.24}{741.089}{in}%
\Vector{q 0 G 0 g 0.478 w 121.418 -24.973 m 127.401 -24.973 l S Q}%
\TextAt{14}{193.418}{741.089}{p}%
\TextAt{21}{199.401}{738.331}{F}%
\TextAt{8}{206.434}{741.089}{,}%
\TextAt{8}{214.239}{741.089}{and}%
\TextAt{8}{237.527}{741.089}{particularly}%
\TextAt{8}{301.473}{741.089}{scalar}%
\TextAt{8}{335.234}{741.089}{multiplication}%
\TextAt{13}{410.884}{741.089}{a}%
\TextAt{16}{420.044}{741.089}{\Character{1}}%
\TextAt{13}{426.379}{741.089}{x}%
\TextAt{9}{437.246}{741.089}{\Character{61}}%
\TextAt{9}{450.552}{741.089}{\Character{91}}%
\TextAt{13}{453.817}{741.089}{a}%
\TextAt{9}{459.962}{741.089}{\Character{93}}%
\TextAt{21}{463.214}{739.296}{F}%
\TextAt{9}{470.247}{741.089}{\Character{40}}%
\TextAt{13}{474.8}{741.089}{x}%
\TextAt{9}{481.452}{741.089}{\Character{41}}%
\TextAt{8}{486.004}{741.089}{,}%
\TextAt{8}{493.809}{741.089}{must}%
\TextAt{8}{523.337}{741.089}{of}%
\TextAt{8}{79.226}{724.043}{course}%
\TextAt{8}{115.056}{724.043}{not}%
\TextAt{8}{135.87}{724.043}{be}%
\TextAt{8}{151.794}{724.043}{confused}%
\TextAt{8}{199.651}{724.043}{with}%
\TextAt{8}{226.311}{724.043}{the}%
\TextAt{8}{246.467}{724.043}{usual}%
\TextAt{8}{277.096}{724.043}{operations}%
\TextAt{8}{334.075}{724.043}{in}%
\TextAt{8}{347.727}{724.043}{the}%
\TextAt{8}{367.884}{724.043}{field}%
\Vector{q 0 G 0 g 0.478 w 321.241 -40.114 m 332.058 -40.114 l S Q}%
\TextAt{13}{393.241}{724.043}{K}%
\TextAt{8}{404.058}{724.043}{.}%
\TextAt{8}{97.159}{704.008}{We}%
\TextAt{8}{117.723}{704.008}{now}%
\TextAt{8}{142.53}{704.008}{consider}%
\TextAt{8}{188.535}{704.008}{a}%
\TextAt{8}{198.71}{704.008}{Lubin–Tate}%
\TextAt{8}{262.049}{704.008}{module}%
\TextAt{13}{303.761}{704.008}{F}%
\TextAt{8}{317.282}{704.008}{for}%
\TextAt{8}{335.586}{704.008}{the}%
\TextAt{8}{356.162}{704.008}{prime}%
\TextAt{8}{389.744}{704.008}{element}%
\TextAt{13}{433.406}{704.008}{\Character{25}}%
\TextAt{8}{444.792}{704.008}{of}%
\TextAt{14}{458.541}{704.008}{o}%
\TextAt{21}{464.385}{702.215}{K}%
\TextAt{8}{472.588}{704.008}{.}%
\TextAt{8}{482.294}{704.008}{We}%
\TextAt{8}{502.858}{704.008}{define}%
\TextAt{8}{79.226}{686.962}{the}%
\TextAt{8}{99.383}{686.962}{group}%
\TextAt{8}{132.558}{686.962}{of}%
\TextAt{13}{146.402}{686.962}{\Character{25}}%
\TextAt{21}{153.472}{691.301}{n}%
\TextAt{12}{159.108}{686.962}{-division}%
\TextAt{12}{213.384}{686.962}{points}%
\TextAt{8}{253.243}{686.962}{by}%
\TextAt{13}{197.265}{658.218}{F}%
\TextAt{9}{206.469}{658.218}{\Character{40}}%
\TextAt{13}{211.021}{658.218}{n}%
\TextAt{9}{218.009}{658.218}{\Character{41}}%
\TextAt{9}{225.864}{658.218}{\Character{61}}%
\TextAt{16}{238.307}{658.218}{\Character{102}}%
\TextAt{13}{244.285}{658.218}{\Character{21}}%
\NativeMathAt{254.434}{658.218}{11.955}{\in}%
\Vector{q 0 G 0 g 0.478 w 193.725 -107.844 m 199.708 -107.844 l S Q}%
\TextAt{14}{265.725}{658.218}{p}%
\TextAt{21}{271.708}{655.46}{F}%
\TextAt{16}{282.062}{658.218}{\Character{106}}%
\TextAt{13}{288.704}{658.218}{\Character{25}}%
\TextAt{21}{295.773}{663.154}{n}%
\TextAt{16}{304.066}{658.218}{\Character{1}}%
\TextAt{13}{310.044}{658.218}{\Character{21}}%
\TextAt{9}{320.193}{658.218}{\Character{61}}%
\TextAt{9}{332.603}{658.218}{0}%
\TextAt{16}{338.472}{658.218}{\Character{103}}%
\TextAt{9}{347.77}{658.218}{\Character{61}}%
\TextAt{8}{360.196}{658.218}{ker}%
\TextAt{9}{376.132}{658.218}{\Character{40}\Character{91}}%
\TextAt{13}{383.936}{658.218}{\Character{25}}%
\TextAt{21}{391.005}{663.154}{n}%
\TextAt{9}{396.642}{658.218}{\Character{93}}%
\TextAt{21}{399.893}{656.425}{F}%
\TextAt{9}{406.927}{658.218}{\Character{41}}%
\TextAt{13}{411.479}{658.218}{\Character{58}}%
\TextAt{8}{79.226}{629.474}{This}%
\TextAt{8}{105.946}{629.474}{is}%
\TextAt{8}{117.71}{629.474}{an}%
\TextAt{14}{133.969}{629.474}{o}%
\TextAt{21}{139.813}{627.68}{K}%
\TextAt{8}{148.016}{629.474}{-module,}%
\TextAt{8}{196.458}{629.474}{and}%
\TextAt{8}{219.221}{629.474}{an}%
\TextAt{14}{235.48}{629.474}{o}%
\TextAt{21}{241.324}{627.68}{K}%
\TextAt{9}{249.527}{629.474}{\Character{47}}%
\TextAt{13}{255.38}{629.474}{\Character{25}}%
\TextAt{21}{262.449}{633.812}{n}%
\TextAt{14}{268.085}{629.474}{o}%
\TextAt{21}{273.93}{627.68}{K}%
\TextAt{8}{282.132}{629.474}{-module}%
\TextAt{8}{327.323}{629.474}{because}%
\TextAt{8}{370.624}{629.474}{it}%
\TextAt{8}{382.328}{629.474}{is}%
\TextAt{8}{394.092}{629.474}{killed}%
\TextAt{8}{425.63}{629.474}{by}%
\TextAt{13}{441.889}{629.474}{\Character{25}}%
\TextAt{21}{448.958}{633.812}{n}%
\TextAt{14}{454.595}{629.474}{o}%
\TextAt{21}{460.439}{627.68}{K}%
\TextAt{8}{468.642}{629.474}{.}%
\TextAt{12}{79.226}{603.535}{Proposition}%
\TextAt{12}{151.937}{603.535}{(5.2).}%
\TextAt{13}{187.253}{603.535}{F}%
\TextAt{9}{196.457}{603.535}{\Character{40}}%
\TextAt{13}{201.009}{603.535}{n}%
\TextAt{9}{207.997}{603.535}{\Character{41}}%
\TextAt{8}{216.446}{603.535}{is}%
\TextAt{8}{228.21}{603.535}{a}%
\TextAt{8}{237.966}{603.535}{free}%
\TextAt{14}{260.394}{603.535}{o}%
\TextAt{21}{266.238}{601.741}{K}%
\TextAt{9}{274.44}{603.535}{\Character{47}}%
\TextAt{13}{280.293}{603.535}{\Character{25}}%
\TextAt{21}{287.363}{607.873}{n}%
\TextAt{14}{292.999}{603.535}{o}%
\TextAt{21}{298.843}{601.741}{K}%
\TextAt{8}{307.046}{603.535}{-module}%
\TextAt{8}{352.236}{603.535}{of}%
\TextAt{8}{365.566}{603.535}{rank}%
\TextAt{9}{392.561}{603.535}{1}%
\TextAt{8}{398.414}{603.535}{.}%
\TextAt{19}{79.226}{568.703}{Proof.}%
\TextAt{8}{116}{568.703}{An}%
\TextAt{8}{135.891}{568.703}{isomorphism}%
\TextAt{13}{205.026}{568.703}{f}%
\TextAt{9}{216.612}{568.703}{\Character{58}}%
\TextAt{13}{224.403}{568.703}{F}%
\NativeArrowAt{238.146}{568.703}{11.955}{19.26}{\longrightarrow}%
\TextAt{13}{261.946}{568.703}{G}%
\TextAt{8}{275.792}{568.703}{of}%
\TextAt{8}{289.837}{568.703}{Lubin–Tate}%
\TextAt{8}{353.473}{568.703}{modules}%
\TextAt{8}{400.096}{568.703}{obviously}%
\TextAt{8}{452.588}{568.703}{induces}%
\TextAt{8}{494.979}{568.703}{isomor-}%
\TextAt{8}{79.226}{551.657}{phisms}%
\TextAt{13}{119.669}{551.657}{f}%
\TextAt{9}{132.255}{551.657}{\Character{58}}%
\Vector{q 0 G 0 g 0.478 w 69.047 -214.406 m 75.03 -214.406 l S Q}%
\TextAt{14}{141.047}{551.657}{p}%
\TextAt{21}{147.03}{548.898}{F}%
\NativeArrowAt{159.603}{551.657}{11.955}{19.26}{\longrightarrow}%
\Vector{q 0 G 0 g 0.478 w 112.404 -214.406 m 118.387 -214.406 l S Q}%
\TextAt{14}{184.404}{551.657}{p}%
\TextAt{21}{190.387}{548.898}{G}%
\TextAt{8}{202.683}{551.657}{and}%
\TextAt{13}{226.747}{551.657}{f}%
\TextAt{9}{239.333}{551.657}{\Character{58}}%
\TextAt{13}{248.125}{551.657}{F}%
\TextAt{9}{257.328}{551.657}{\Character{40}}%
\TextAt{13}{261.881}{551.657}{n}%
\TextAt{9}{268.868}{551.657}{\Character{41}}%
\NativeArrowAt{278.96}{551.657}{11.955}{19.26}{\longrightarrow}%
\TextAt{13}{303.761}{551.657}{G}%
\TextAt{9}{312.995}{551.657}{\Character{40}}%
\TextAt{13}{317.547}{551.657}{n}%
\TextAt{9}{324.535}{551.657}{\Character{41}}%
\TextAt{8}{334.286}{551.657}{of}%
\TextAt{14}{348.918}{551.657}{o}%
\TextAt{21}{354.762}{549.863}{K}%
\TextAt{8}{362.965}{551.657}{-modules.}%
\TextAt{8}{421.227}{551.657}{By}%
\TextAt{8}{440.892}{551.657}{(4.6),}%
\TextAt{8}{473.747}{551.657}{Lubin–Tate}%
\TextAt{8}{79.226}{534.611}{modules}%
\TextAt{8}{124.894}{534.611}{for}%
\TextAt{8}{142.538}{534.611}{the}%
\TextAt{8}{162.455}{534.611}{same}%
\TextAt{8}{191.54}{534.611}{prime}%
\TextAt{8}{224.464}{534.611}{element}%
\TextAt{13}{267.465}{534.611}{\Character{25}}%
\TextAt{8}{278.192}{534.611}{are}%
\TextAt{8}{297.462}{534.611}{all}%
\TextAt{8}{313.481}{534.611}{isomorphic.}%
\TextAt{8}{377.202}{534.611}{We}%
\TextAt{8}{397.106}{534.611}{may}%
\TextAt{8}{422.234}{534.611}{therefore}%
\TextAt{8}{471.094}{534.611}{assume}%
\TextAt{8}{511.298}{534.611}{that}%
\TextAt{13}{79.226}{517.565}{F}%
\TextAt{9}{91.751}{517.565}{\Character{61}}%
\TextAt{13}{104.176}{517.565}{F}%
\TextAt{21}{111.754}{515.771}{e}%
\TextAt{8}{116.18}{517.565}{,}%
\TextAt{8}{122.937}{517.565}{with}%
\TextAt{13}{149.106}{517.565}{e}%
\TextAt{9}{154.531}{517.565}{\Character{40}}%
\TextAt{13}{159.084}{517.565}{X}%
\TextAt{9}{169.739}{517.565}{\Character{41}}%
\TextAt{9}{177.594}{517.565}{\Character{61}}%
\TextAt{13}{190.038}{517.565}{X}%
\TextAt{21}{200.693}{521.903}{q}%
\TextAt{9}{206.91}{517.565}{\Character{43}}%
\TextAt{13}{217.668}{517.565}{\Character{25}}%
\TextAt{13}{224.722}{517.565}{X}%
\TextAt{9}{238.713}{517.565}{\Character{61}}%
\TextAt{9}{251.123}{517.565}{\Character{91}}%
\TextAt{13}{254.39}{517.565}{\Character{25}}%
\TextAt{9}{261.46}{517.565}{\Character{93}}%
\TextAt{21}{264.711}{515.771}{F}%
\TextAt{9}{271.745}{517.565}{\Character{40}}%
\TextAt{13}{276.297}{517.565}{X}%
\TextAt{9}{286.952}{517.565}{\Character{41}}%
\TextAt{8}{291.504}{517.565}{.}%
\TextAt{13}{299.789}{517.565}{F}%
\TextAt{9}{308.993}{517.565}{\Character{40}}%
\TextAt{13}{313.545}{517.565}{n}%
\TextAt{9}{320.533}{517.565}{\Character{41}}%
\TextAt{8}{328.491}{517.565}{then}%
\TextAt{8}{354.661}{517.565}{consists}%
\TextAt{8}{397.28}{517.565}{of}%
\TextAt{8}{410.119}{517.565}{the}%
\TextAt{13}{429.785}{517.565}{q}%
\TextAt{21}{435.404}{521.903}{n}%
\TextAt{8}{444.447}{517.565}{zeroes}%
\TextAt{8}{478.805}{517.565}{of}%
\TextAt{8}{491.644}{517.565}{the}%
\TextAt{8}{511.31}{517.565}{iter-}%
\TextAt{8}{79.226}{500.519}{ated}%
\TextAt{8}{104.903}{500.519}{polynomial}%
\TextAt{13}{165.058}{500.519}{e}%
\TextAt{21}{170.484}{504.857}{n}%
\TextAt{9}{176.12}{500.519}{\Character{40}}%
\TextAt{13}{180.672}{500.519}{X}%
\TextAt{9}{191.328}{500.519}{\Character{41}}%
\TextAt{9}{199.183}{500.519}{\Character{61}}%
\TextAt{9}{211.604}{500.519}{\Character{40}}%
\TextAt{13}{216.178}{500.519}{e}%
\TextAt{16}{223.57}{500.519}{\Character{14}}%
\TextAt{16}{231.508}{500.519}{\Character{1}}%
\TextAt{16}{236.816}{500.519}{\Character{1}}%
\TextAt{16}{242.113}{500.519}{\Character{1}}%
\TextAt{16}{247.397}{500.519}{\Character{14}}%
\TextAt{13}{255.371}{500.519}{e}%
\TextAt{9}{260.796}{500.519}{\Character{41}\Character{40}}%
\TextAt{13}{269.901}{500.519}{X}%
\TextAt{9}{280.556}{500.519}{\Character{41}}%
\TextAt{9}{288.411}{500.519}{\Character{61}}%
\TextAt{9}{300.832}{500.519}{\Character{91}}%
\TextAt{13}{304.106}{500.519}{\Character{25}}%
\TextAt{21}{311.176}{504.857}{n}%
\TextAt{9}{316.812}{500.519}{\Character{93}}%
\TextAt{21}{320.064}{498.725}{F}%
\TextAt{9}{327.097}{500.519}{\Character{40}}%
\TextAt{13}{331.65}{500.519}{X}%
\TextAt{9}{342.305}{500.519}{\Character{41}}%
\TextAt{8}{346.857}{500.519}{,}%
\TextAt{8}{353.736}{500.519}{which}%
\TextAt{8}{386.885}{500.519}{is}%
\TextAt{8}{398.311}{500.519}{easily}%
\TextAt{8}{430.228}{500.519}{shown,}%
\TextAt{8}{468.717}{500.519}{by}%
\TextAt{8}{484.638}{500.519}{induction}%
\TextAt{8}{79.226}{483.473}{on}%
\TextAt{13}{95.485}{483.473}{n}%
\TextAt{8}{102.473}{483.473}{,}%
\TextAt{8}{109.622}{483.473}{to}%
\TextAt{8}{123.932}{483.473}{be}%
\TextAt{8}{139.857}{483.473}{separable.}%
\TextAt{8}{195.851}{483.473}{Now}%
\TextAt{8}{222.511}{483.473}{if}%
\TextAt{13}{233.234}{483.473}{\Character{21}}%
\TextAt{21}{240.063}{481.679}{n}%
\NativeMathAt{249.02}{483.473}{11.955}{\in}%
\TextAt{13}{260.311}{483.473}{F}%
\TextAt{9}{269.515}{483.473}{\Character{40}}%
\TextAt{13}{274.067}{483.473}{n}%
\TextAt{9}{281.055}{483.473}{\Character{41}}%
\TextAt{16}{288.264}{483.473}{\Character{110}}%
\TextAt{13}{296.898}{483.473}{F}%
\TextAt{9}{306.102}{483.473}{\Character{40}}%
\TextAt{13}{310.654}{483.473}{n}%
\TextAt{16}{320.298}{483.473}{\Character{0}}%
\TextAt{9}{332.253}{483.473}{1\Character{41}}%
\TextAt{8}{342.659}{483.473}{,}%
\TextAt{8}{349.808}{483.473}{then}%
\TextAt{14}{227.526}{454.728}{o}%
\TextAt{21}{233.37}{452.935}{K}%
\NativeArrowAt{244.894}{454.728}{11.955}{19.26}{\longrightarrow}%
\TextAt{13}{267.476}{454.728}{F}%
\TextAt{9}{276.679}{454.728}{\Character{40}}%
\TextAt{13}{281.232}{454.728}{n}%
\TextAt{9}{288.219}{454.728}{\Character{41}}%
\TextAt{13}{292.772}{454.728}{\Character{59}}%
\TextAt{13}{321.907}{454.728}{a}%
\NativeArrowAt{331.392}{454.728}{11.955}{19.26}{\longmapsto}%
\TextAt{13}{353.974}{454.728}{a}%
\TextAt{16}{362.776}{454.728}{\Character{1}}%
\TextAt{13}{368.753}{454.728}{\Character{21}}%
\TextAt{21}{375.582}{452.935}{n}%
\TextAt{13}{381.218}{454.728}{\Character{59}}%
\TextAt{8}{79.226}{425.984}{is}%
\TextAt{8}{89.797}{425.984}{a}%
\TextAt{8}{98.359}{425.984}{homomorphism}%
\TextAt{8}{179.835}{425.984}{of}%
\TextAt{14}{191.972}{425.984}{o}%
\TextAt{21}{197.816}{424.191}{K}%
\TextAt{8}{206.019}{425.984}{-modules}%
\TextAt{8}{254.63}{425.984}{with}%
\TextAt{8}{280.097}{425.984}{kernel}%
\TextAt{13}{313.37}{425.984}{\Character{25}}%
\TextAt{21}{320.44}{430.322}{n}%
\TextAt{14}{326.076}{425.984}{o}%
\TextAt{21}{331.92}{424.191}{K}%
\TextAt{8}{340.123}{425.984}{.}%
\TextAt{8}{348.173}{425.984}{It}%
\TextAt{8}{359.652}{425.984}{induces}%
\TextAt{8}{400.135}{425.984}{a}%
\TextAt{8}{408.697}{425.984}{bijective}%
\TextAt{8}{453.997}{425.984}{homomorphism}%
\TextAt{14}{79.226}{408.938}{o}%
\TextAt{21}{85.07}{407.145}{K}%
\TextAt{9}{93.273}{408.938}{\Character{47}}%
\TextAt{13}{99.126}{408.938}{\Character{25}}%
\TextAt{21}{106.195}{413.276}{n}%
\TextAt{14}{111.832}{408.938}{o}%
\TextAt{21}{117.676}{407.145}{K}%
\NativeArrowAt{129.199}{408.938}{11.955}{19.26}{\longrightarrow}%
\TextAt{13}{151.781}{408.938}{F}%
\TextAt{9}{160.985}{408.938}{\Character{40}}%
\TextAt{13}{165.537}{408.938}{n}%
\TextAt{9}{172.525}{408.938}{\Character{41}}%
\TextAt{8}{180.975}{408.938}{because}%
\TextAt{8}{224.276}{408.938}{both}%
\TextAt{8}{251.917}{408.938}{sides}%
\TextAt{8}{279.999}{408.938}{are}%
\TextAt{8}{299.51}{408.938}{of}%
\TextAt{8}{312.84}{408.938}{order}%
\TextAt{13}{343.41}{408.938}{q}%
\TextAt{21}{349.029}{413.276}{n}%
\TextAt{8}{354.665}{408.938}{.}%
\Vector{q 0 G 0 g 0.399 w 452.334 -365.062 m 452.334 -356.992 l S Q}%
\Vector{q 0 G 0 g 0.399 w 452.534 -357.191 m 459.707 -357.191 l S Q}%
\Vector{q 0 G 0 g 0.399 w 452.534 -364.863 m 459.707 -364.863 l S Q}%
\Vector{q 0 G 0 g 0.399 w 459.906 -365.062 m 459.906 -356.992 l S Q}%
\TextAt{12}{79.226}{374.106}{Corollary}%
\TextAt{12}{139.014}{374.106}{(5.3).}%
\TextAt{8}{174.33}{374.106}{Associating}%
\TextAt{13}{236.891}{374.106}{a}%
\NativeArrowAt{246.357}{374.106}{11.955}{11.955}{\mapsto}%
\TextAt{9}{261.633}{374.106}{\Character{91}}%
\TextAt{13}{264.884}{374.106}{a}%
\TextAt{9}{271.029}{374.106}{\Character{93}}%
\TextAt{21}{274.281}{372.313}{F}%
\TextAt{8}{285.212}{374.106}{we}%
\TextAt{8}{302.439}{374.106}{obtain}%
\TextAt{8}{338.867}{374.106}{canonical}%
\TextAt{8}{390.25}{374.106}{isomorphisms}%
\TextAt{14}{145.802}{345.362}{o}%
\TextAt{21}{151.646}{343.568}{K}%
\TextAt{9}{159.849}{345.362}{\Character{47}}%
\TextAt{13}{165.702}{345.362}{\Character{25}}%
\TextAt{21}{172.771}{350.298}{n}%
\TextAt{14}{178.407}{345.362}{o}%
\TextAt{21}{184.252}{343.568}{K}%
\NativeArrowAt{195.775}{345.362}{11.955}{19.26}{\longrightarrow}%
\TextAt{8}{218.357}{345.362}{End}%
\TextAt{27}{239.327}{343.568}{o}%
\TextAt{22}{243.223}{342.229}{K}%
\TextAt{9}{250.764}{345.362}{\Character{40}}%
\TextAt{13}{255.316}{345.362}{F}%
\TextAt{9}{264.52}{345.362}{\Character{40}}%
\TextAt{13}{269.072}{345.362}{n}%
\TextAt{9}{276.06}{345.362}{\Character{41}\Character{41}}%
\TextAt{8}{297.12}{345.362}{and}%
\TextAt{13}{327.94}{345.362}{U}%
\TextAt{21}{335.89}{343.568}{K}%
\TextAt{9}{344.093}{345.362}{\Character{47}}%
\TextAt{13}{349.946}{345.362}{U}%
\TextAt{17}{359.145}{351.472}{\Character{40}}%
\TextAt{21}{362.438}{351.472}{n}%
\TextAt{17}{367.576}{351.472}{\Character{41}}%
\TextAt{21}{357.896}{342.121}{K}%
\NativeArrowAt{374.689}{345.362}{11.955}{19.26}{\longrightarrow}%
\TextAt{8}{397.271}{345.362}{Aut}%
\TextAt{27}{417.104}{343.568}{o}%
\TextAt{22}{421}{342.229}{K}%
\TextAt{9}{428.542}{345.362}{\Character{40}}%
\TextAt{13}{433.094}{345.362}{F}%
\TextAt{9}{442.297}{345.362}{\Character{40}}%
\TextAt{13}{446.85}{345.362}{n}%
\TextAt{9}{453.837}{345.362}{\Character{41}\Character{41}}%
\TextAt{13}{462.942}{345.362}{\Character{58}}%
\TextAt{19}{79.226}{298.831}{Proof.}%
\TextAt{8}{116}{298.831}{The}%
\TextAt{8}{140.054}{298.831}{map}%
\TextAt{8}{166.068}{298.831}{on}%
\TextAt{8}{182.328}{298.831}{the}%
\TextAt{8}{202.484}{298.831}{left}%
\TextAt{8}{222.963}{298.831}{is}%
\TextAt{8}{234.727}{298.831}{an}%
\TextAt{8}{250.986}{298.831}{isomorphism}%
\TextAt{8}{319.406}{298.831}{since}%
\TextAt{14}{348.074}{298.831}{o}%
\TextAt{21}{353.918}{297.038}{K}%
\TextAt{9}{362.121}{298.831}{\Character{47}}%
\TextAt{13}{367.974}{298.831}{\Character{25}}%
\TextAt{21}{375.043}{303.17}{n}%
\TextAt{14}{380.679}{298.831}{o}%
\TextAt{21}{386.524}{297.038}{K}%
\NativeMathAt{398.047}{298.831}{11.955}{\cong}%
\TextAt{13}{410.667}{298.831}{F}%
\TextAt{9}{419.87}{298.831}{\Character{40}}%
\TextAt{13}{424.422}{298.831}{n}%
\TextAt{9}{431.41}{298.831}{\Character{41}}%
\TextAt{8}{439.86}{298.831}{and}%
\TextAt{8}{228.994}{270.087}{End}%
\TextAt{27}{249.963}{268.294}{o}%
\TextAt{22}{253.859}{266.954}{K}%
\TextAt{9}{261.4}{270.087}{\Character{40}}%
\TextAt{14}{265.953}{270.087}{o}%
\TextAt{21}{271.797}{268.294}{K}%
\TextAt{9}{280}{270.087}{\Character{47}}%
\TextAt{13}{285.853}{270.087}{\Character{25}}%
\TextAt{21}{292.922}{275.023}{n}%
\TextAt{14}{298.558}{270.087}{o}%
\TextAt{21}{304.403}{268.294}{K}%
\TextAt{9}{312.605}{270.087}{\Character{41}}%
\NativeMathAt{320.478}{270.087}{11.955}{\cong}%
\TextAt{14}{333.098}{270.087}{o}%
\TextAt{21}{338.942}{268.294}{K}%
\TextAt{9}{347.145}{270.087}{\Character{47}}%
\TextAt{13}{352.998}{270.087}{\Character{25}}%
\TextAt{21}{360.067}{275.023}{n}%
\TextAt{14}{365.703}{270.087}{o}%
\TextAt{21}{371.547}{268.294}{K}%
\TextAt{13}{379.75}{270.087}{\Character{58}}%
\TextAt{8}{79.226}{241.343}{The}%
\TextAt{8}{103.28}{241.343}{one}%
\TextAt{8}{124.739}{241.343}{on}%
\TextAt{8}{140.999}{241.343}{the}%
\TextAt{8}{161.155}{241.343}{right}%
\TextAt{8}{189.453}{241.343}{is}%
\TextAt{8}{201.217}{241.343}{obtained}%
\TextAt{8}{249.348}{241.343}{by}%
\TextAt{8}{265.607}{241.343}{taking}%
\TextAt{8}{301.712}{241.343}{the}%
\TextAt{8}{321.868}{241.343}{unit}%
\TextAt{8}{346.58}{241.343}{groups}%
\TextAt{8}{384.37}{241.343}{of}%
\TextAt{8}{397.7}{241.343}{these}%
\TextAt{8}{427.672}{241.343}{rings.}%
\Vector{q 0 G 0 g 0.399 w 452.334 -532.657 m 452.334 -524.588 l S Q}%
\Vector{q 0 G 0 g 0.399 w 452.534 -524.787 m 459.707 -524.787 l S Q}%
\Vector{q 0 G 0 g 0.399 w 452.534 -532.458 m 459.707 -532.458 l S Q}%
\Vector{q 0 G 0 g 0.399 w 459.906 -532.657 m 459.906 -524.588 l S Q}%
\TextAt{8}{97.159}{212.415}{Given}%
\TextAt{8}{131.718}{212.415}{a}%
\TextAt{8}{142.139}{212.415}{Lubin–Tate}%
\TextAt{8}{205.725}{212.415}{module}%
\TextAt{13}{247.684}{212.415}{F}%
\TextAt{8}{261.451}{212.415}{for}%
\TextAt{8}{280.002}{212.415}{the}%
\TextAt{8}{300.824}{212.415}{prime}%
\TextAt{8}{334.654}{212.415}{element}%
\TextAt{13}{378.562}{212.415}{\Character{25}}%
\TextAt{8}{385.631}{212.415}{,}%
\TextAt{8}{393.613}{212.415}{we}%
\TextAt{8}{411.506}{212.415}{now}%
\TextAt{8}{436.56}{212.415}{define}%
\TextAt{8}{471.036}{212.415}{the}%
\TextAt{12}{491.858}{212.415}{field}%
\TextAt{12}{521.926}{212.415}{of}%
\TextAt{13}{79.226}{195.369}{\Character{25}}%
\TextAt{21}{86.295}{199.707}{n}%
\TextAt{12}{91.932}{195.369}{-division}%
\TextAt{12}{146.208}{195.369}{points}%
\TextAt{8}{186.067}{195.369}{by}%
\TextAt{13}{202.326}{195.369}{L}%
\TextAt{21}{210.29}{193.576}{n}%
\TextAt{9}{219.247}{195.369}{\Character{61}}%
\TextAt{13}{231.673}{195.369}{K}%
\TextAt{9}{242.491}{195.369}{\Character{40}}%
\TextAt{13}{247.043}{195.369}{F}%
\TextAt{9}{256.247}{195.369}{\Character{40}}%
\TextAt{13}{260.799}{195.369}{n}%
\TextAt{9}{267.786}{195.369}{\Character{41}\Character{41}}%
\TextAt{8}{276.891}{195.369}{.}%
\TextAt{8}{285.339}{195.369}{Since}%
\TextAt{13}{315.897}{195.369}{F}%
\TextAt{9}{325.1}{195.369}{\Character{40}}%
\TextAt{13}{329.653}{195.369}{n}%
\TextAt{9}{336.64}{195.369}{\Character{41}}%
\NativeMathAt{344.514}{195.369}{11.955}{\subseteq}%
\TextAt{13}{357.133}{195.369}{F}%
\TextAt{9}{366.337}{195.369}{\Character{40}}%
\TextAt{13}{370.889}{195.369}{n}%
\TextAt{9}{380.533}{195.369}{\Character{43}}%
\TextAt{9}{392.285}{195.369}{1\Character{41}}%
\TextAt{8}{406.597}{195.369}{we}%
\TextAt{8}{423.825}{195.369}{get}%
\TextAt{8}{443.335}{195.369}{a}%
\TextAt{8}{453.091}{195.369}{tower}%
\TextAt{8}{484.963}{195.369}{of}%
\TextAt{8}{498.293}{195.369}{fields}%
\TextAt{13}{215.331}{160.114}{K}%
\NativeMathAt{229.47}{160.114}{11.955}{\subseteq}%
\TextAt{13}{242.089}{160.114}{L}%
\TextAt{17}{250.054}{158.321}{1}%
\NativeMathAt{258.107}{160.114}{11.955}{\subseteq}%
\TextAt{13}{270.726}{160.114}{L}%
\TextAt{17}{278.691}{158.321}{2}%
\NativeMathAt{286.744}{160.114}{11.955}{\subseteq}%
\TextAt{16}{299.345}{160.114}{\Character{1}}%
\TextAt{16}{304.653}{160.114}{\Character{1}}%
\TextAt{16}{309.961}{160.114}{\Character{1}}%
\NativeMathAt{316.584}{160.114}{11.955}{\subseteq}%
\TextAt{13}{329.251}{160.114}{L}%
\TextAt{21}{337.215}{158.321}{\Character{25}}%
\TextAt{9}{346.183}{160.114}{\Character{58}\Character{61}}%
\NativeMathAt{365.606}{175.058}{7.97}{\infty}%
\TextAt{20}{363.198}{171.472}{\Character{91}}%
\TextAt{21}{361.86}{146.366}{n}%
\TextAt{17}{366.999}{146.366}{\Character{61}1}%
\TextAt{13}{379.812}{160.114}{L}%
\TextAt{21}{387.776}{158.321}{n}%
\TextAt{13}{393.413}{160.114}{\Character{58}}%
\TextAt{8}{79.226}{124.059}{These}%
\TextAt{8}{114.766}{124.059}{fields}%
\TextAt{8}{146.408}{124.059}{are}%
\TextAt{8}{167.589}{124.059}{also}%
\TextAt{8}{192.74}{124.059}{called}%
\TextAt{8}{227.574}{124.059}{the}%
\TextAt{12}{249.401}{124.059}{Lubin–Tate}%
\TextAt{12}{322.91}{124.059}{extensions}%
\TextAt{8}{383.547}{124.059}{.}%
\TextAt{8}{397.007}{124.059}{They}%
\TextAt{8}{428.912}{124.059}{only}%
\TextAt{8}{456.275}{124.059}{depend}%
\TextAt{8}{498.581}{124.059}{on}%
\TextAt{8}{516.51}{124.059}{the}%
\TextAt{8}{79.226}{107.013}{prime}%
\TextAt{8}{113.946}{107.013}{element}%
\TextAt{13}{158.743}{107.013}{\Character{25}}%
\TextAt{8}{165.813}{107.013}{,}%
\TextAt{8}{174.907}{107.013}{not}%
\TextAt{8}{197.277}{107.013}{on}%
\TextAt{8}{215.091}{107.013}{the}%
\TextAt{8}{236.804}{107.013}{Lubin–Tate}%
\TextAt{8}{301.28}{107.013}{module}%
\TextAt{13}{344.129}{107.013}{F}%
\TextAt{8}{353.332}{107.013}{.}%
\TextAt{8}{366.448}{107.013}{For}%
\TextAt{8}{388.974}{107.013}{if}%
\TextAt{13}{401.253}{107.013}{G}%
\TextAt{8}{415.94}{107.013}{is}%
\TextAt{8}{429.26}{107.013}{another}%
\TextAt{8}{473.747}{107.013}{Lubin–Tate}%
\TextAt{8}{79.226}{89.967}{module}%
\TextAt{8}{120.523}{89.967}{for}%
\TextAt{13}{138.411}{89.967}{\Character{25}}%
\TextAt{8}{145.481}{89.967}{,}%
\TextAt{8}{152.634}{89.967}{then}%
\TextAt{8}{179.298}{89.967}{by}%
\TextAt{8}{195.56}{89.967}{(4.6),}%
\TextAt{8}{226.792}{89.967}{there}%
\TextAt{8}{256.707}{89.967}{is}%
\TextAt{8}{268.474}{89.967}{an}%
\TextAt{8}{284.737}{89.967}{isomorphism}%
\TextAt{13}{353.16}{89.967}{f}%
\TextAt{9}{363.533}{89.967}{\Character{58}}%
\TextAt{13}{370.112}{89.967}{F}%
\NativeArrowAt{382.642}{89.967}{11.955}{19.26}{\longrightarrow}%
\TextAt{13}{405.23}{89.967}{G}%
\TextAt{8}{414.464}{89.967}{,}%
\TextAt{13}{421.618}{89.967}{f}%
\NativeMathAt{431.991}{89.967}{11.955}{\in}%
\TextAt{14}{443.288}{89.967}{o}%
\TextAt{21}{449.132}{88.174}{K}%
\TextAt{9}{457.335}{89.967}{\Character{91}\Character{91}}%
\TextAt{13}{463.838}{89.967}{X}%
\TextAt{9}{474.493}{89.967}{\Character{93}\Character{93}}%
\TextAt{8}{484.897}{89.967}{such}%
\TextAt{8}{511.298}{89.967}{that}%
\TextAt{13}{79.226}{72.921}{G}%
\TextAt{9}{88.46}{72.921}{\Character{40}}%
\TextAt{13}{93.012}{72.921}{n}%
\TextAt{9}{100}{72.921}{\Character{41}}%
\TextAt{9}{109.277}{72.921}{\Character{61}}%
\TextAt{13}{123.137}{72.921}{f}%
\TextAt{9}{130.183}{72.921}{\Character{40}}%
\TextAt{13}{134.735}{72.921}{F}%
\TextAt{9}{143.939}{72.921}{\Character{40}}%
\TextAt{13}{148.491}{72.921}{n}%
\TextAt{9}{155.479}{72.921}{\Character{41}\Character{41}}%
\NativeMathAt{169.324}{72.921}{11.955}{\subseteq}%
\TextAt{13}{183.362}{72.921}{K}%
\TextAt{9}{194.18}{72.921}{\Character{40}}%
\TextAt{13}{198.732}{72.921}{F}%
\TextAt{9}{207.936}{72.921}{\Character{40}}%
\TextAt{13}{212.488}{72.921}{n}%
\TextAt{9}{219.476}{72.921}{\Character{41}\Character{41}}%
\TextAt{8}{228.58}{72.921}{,}%
\TextAt{8}{236.77}{72.921}{and}%
\TextAt{8}{260.366}{72.921}{hence}%
\TextAt{13}{293.705}{72.921}{K}%
\TextAt{9}{304.522}{72.921}{\Character{40}}%
\TextAt{13}{309.075}{72.921}{G}%
\TextAt{9}{318.308}{72.921}{\Character{40}}%
\TextAt{13}{322.861}{72.921}{n}%
\TextAt{9}{329.848}{72.921}{\Character{41}\Character{41}}%
\TextAt{9}{343.668}{72.921}{\Character{61}}%
\TextAt{13}{357.537}{72.921}{K}%
\TextAt{9}{368.355}{72.921}{\Character{40}}%
\TextAt{13}{372.907}{72.921}{F}%
\TextAt{9}{382.111}{72.921}{\Character{40}}%
\TextAt{13}{386.663}{72.921}{n}%
\TextAt{9}{393.651}{72.921}{\Character{41}\Character{41}}%
\TextAt{8}{402.756}{72.921}{.}%
\TextAt{8}{413.702}{72.921}{If}%
\TextAt{13}{426.227}{72.921}{F}%
\TextAt{8}{440.161}{72.921}{is}%
\TextAt{8}{452.758}{72.921}{the}%
\TextAt{8}{473.747}{72.921}{Lubin–Tate}%
\TextAt{8}{297.211}{43.033}{318}%
\Destination{78.23}{809.815}{nav.page.318}%
\Destination{79.36}{756.055}{nav.print.347}%
\Destination{79.36}{719.055}{nav.print.348}%
\Destination{79.36}{617.555}{nav.number.5.5.2}%
\Destination{79.36}{388.155}{nav.number.5.5.3}%
\Hotlink{440.392}{549.455}{25.078}{11.047}{nav.number.5.4.6}%
\Hotlink{195.06}{87.755}{25.078}{11.047}{nav.number.5.4.6}%
}

% Page 331 (body page 319)
\Page{%
\TextAt{2}{79.226}{796.943}{V.}%
\TextAt{2}{94.073}{796.943}{Local}%
\TextAt{2}{123.615}{796.943}{Class}%
\TextAt{2}{152.208}{796.943}{Field}%
\TextAt{2}{179.939}{796.943}{Theory}%
\Vector{q 0 G 0 g 0.299 w 7.226 18.729 m 460.769 18.729 l S Q}%
\TextAt{8}{79.226}{761.124}{module}%
\TextAt{13}{120.885}{761.124}{F}%
\TextAt{21}{128.463}{759.33}{e}%
\TextAt{8}{137.152}{761.124}{belonging}%
\TextAt{8}{190.527}{761.124}{to}%
\TextAt{8}{205.202}{761.124}{a}%
\TextAt{8}{215.323}{761.124}{Lubin–Tate}%
\TextAt{8}{278.609}{761.124}{polynomial}%
\TextAt{13}{339.468}{761.124}{e}%
\TextAt{9}{344.893}{761.124}{\Character{40}}%
\TextAt{13}{349.445}{761.124}{X}%
\TextAt{9}{360.1}{761.124}{\Character{41}}%
\NativeMathAt{368.596}{761.124}{11.955}{\in}%
\TextAt{16}{380.491}{761.124}{\Character{69}}%
\TextAt{21}{386.82}{759.33}{\Character{25}}%
\TextAt{8}{392.467}{761.124}{,}%
\TextAt{8}{400.073}{761.124}{then}%
\TextAt{13}{427.098}{761.124}{e}%
\TextAt{9}{432.524}{761.124}{\Character{40}}%
\TextAt{13}{437.076}{761.124}{X}%
\TextAt{9}{447.731}{761.124}{\Character{41}}%
\TextAt{9}{456.207}{761.124}{\Character{61}}%
\TextAt{9}{469.25}{761.124}{\Character{91}}%
\TextAt{13}{472.527}{761.124}{\Character{25}}%
\TextAt{9}{479.597}{761.124}{\Character{93}}%
\TextAt{21}{482.848}{759.33}{F}%
\TextAt{9}{489.882}{761.124}{\Character{40}}%
\TextAt{13}{494.434}{761.124}{X}%
\TextAt{9}{505.089}{761.124}{\Character{41}}%
\TextAt{8}{513.904}{761.124}{and}%
\TextAt{13}{79.226}{744.078}{L}%
\TextAt{21}{87.191}{742.284}{n}%
\TextAt{16}{96.148}{744.078}{\Character{106}}%
\TextAt{13}{102.79}{744.078}{K}%
\TextAt{8}{117.505}{744.078}{is}%
\TextAt{8}{129.268}{744.078}{the}%
\TextAt{8}{149.425}{744.078}{splitting}%
\TextAt{8}{195.667}{744.078}{field}%
\TextAt{8}{221.024}{744.078}{of}%
\TextAt{8}{234.354}{744.078}{the}%
\TextAt{13}{254.511}{744.078}{n}%
\TextAt{8}{261.498}{744.078}{-fold}%
\TextAt{8}{288.481}{744.078}{iteration}%
\TextAt{13}{212.091}{715.077}{e}%
\TextAt{21}{217.517}{720.013}{n}%
\TextAt{9}{223.153}{715.077}{\Character{40}}%
\TextAt{13}{227.706}{715.077}{X}%
\TextAt{9}{238.361}{715.077}{\Character{41}}%
\TextAt{9}{246.216}{715.077}{\Character{61}}%
\TextAt{9}{258.637}{715.077}{\Character{40}}%
\TextAt{13}{263.212}{715.077}{e}%
\TextAt{16}{271.294}{715.077}{\Character{14}}%
\TextAt{16}{279.926}{715.077}{\Character{1}}%
\TextAt{16}{285.222}{715.077}{\Character{1}}%
\TextAt{16}{290.53}{715.077}{\Character{1}}%
\TextAt{16}{296.496}{715.077}{\Character{14}}%
\TextAt{13}{305.167}{715.077}{e}%
\TextAt{9}{310.592}{715.077}{\Character{41}\Character{40}}%
\TextAt{13}{319.697}{715.077}{X}%
\TextAt{9}{330.352}{715.077}{\Character{41}}%
\TextAt{9}{338.207}{715.077}{\Character{61}}%
\TextAt{9}{350.628}{715.077}{\Character{91}}%
\TextAt{13}{353.902}{715.077}{\Character{25}}%
\TextAt{21}{360.971}{720.013}{n}%
\TextAt{9}{366.608}{715.077}{\Character{93}}%
\TextAt{21}{369.859}{713.283}{F}%
\TextAt{9}{376.893}{715.077}{\Character{40}}%
\TextAt{13}{381.445}{715.077}{X}%
\TextAt{9}{392.1}{715.077}{\Character{41}}%
\TextAt{13}{396.653}{715.077}{\Character{58}}%
\TextAt{12}{79.226}{683.087}{Example:}%
\TextAt{8}{139.37}{683.087}{If}%
\TextAt{14}{151.152}{683.087}{o}%
\TextAt{21}{156.996}{681.293}{K}%
\TextAt{9}{168.673}{683.087}{\Character{61}}%
\TextAt{18}{181.252}{683.087}{Z}%
\TextAt{21}{189.222}{681.293}{p}%
\TextAt{8}{197.971}{683.087}{and}%
\TextAt{13}{220.823}{683.087}{F}%
\TextAt{8}{234.015}{683.087}{is}%
\TextAt{8}{245.869}{683.087}{the}%
\TextAt{8}{266.115}{683.087}{Lubin–Tate}%
\TextAt{8}{329.125}{683.087}{module}%
\TextAt{18}{370.509}{683.087}{G}%
\TextAt{21}{379.807}{681.293}{m}%
\TextAt{8}{387.796}{683.087}{,}%
\TextAt{8}{395.058}{683.087}{then}%
\TextAt{13}{421.808}{683.087}{e}%
\TextAt{21}{427.233}{687.425}{n}%
\TextAt{9}{432.87}{683.087}{\Character{40}}%
\TextAt{13}{437.422}{683.087}{X}%
\TextAt{9}{448.077}{683.087}{\Character{41}}%
\TextAt{9}{456.087}{683.087}{\Character{61}}%
\TextAt{9}{468.664}{683.087}{\Character{91}}%
\TextAt{13}{471.934}{683.087}{p}%
\TextAt{21}{477.81}{687.425}{n}%
\TextAt{9}{483.446}{683.087}{\Character{93}}%
\TextAt{25}{486.697}{681.293}{G}%
\TextAt{22}{492.896}{680.297}{m}%
\TextAt{9}{500.431}{683.087}{\Character{40}}%
\TextAt{13}{504.983}{683.087}{X}%
\TextAt{9}{515.638}{683.087}{\Character{41}}%
\TextAt{9}{523.648}{683.087}{\Character{61}}%
\TextAt{9}{79.226}{666.041}{\Character{40}1}%
\TextAt{9}{91.074}{666.041}{\Character{43}}%
\TextAt{13}{101.659}{666.041}{X}%
\TextAt{9}{112.315}{666.041}{\Character{41}}%
\TextAt{21}{116.867}{670.379}{p}%
\TextAt{22}{121.13}{673.192}{n}%
\TextAt{16}{128.188}{666.041}{\Character{0}}%
\TextAt{9}{138.949}{666.041}{1}%
\TextAt{8}{144.802}{666.041}{.}%
\TextAt{8}{153.055}{666.041}{So}%
\TextAt{18}{168.73}{666.041}{G}%
\TextAt{21}{178.028}{664.247}{m}%
\TextAt{9}{186.017}{666.041}{\Character{40}}%
\TextAt{13}{190.569}{666.041}{n}%
\TextAt{9}{197.557}{666.041}{\Character{41}}%
\TextAt{8}{205.422}{666.041}{consists}%
\TextAt{8}{247.948}{666.041}{of}%
\TextAt{8}{260.694}{666.041}{the}%
\TextAt{8}{280.266}{666.041}{elements}%
\TextAt{13}{327.538}{666.041}{\Character{16}}%
\TextAt{16}{335.015}{666.041}{\Character{0}}%
\TextAt{9}{345.776}{666.041}{1}%
\TextAt{8}{351.629}{666.041}{,}%
\TextAt{8}{358.31}{666.041}{where}%
\TextAt{13}{391.535}{666.041}{\Character{16}}%
\TextAt{8}{400.864}{666.041}{varies}%
\TextAt{8}{433.192}{666.041}{over}%
\TextAt{8}{457.654}{666.041}{the}%
\TextAt{13}{477.226}{666.041}{p}%
\TextAt{21}{483.101}{670.379}{n}%
\TextAt{8}{488.737}{666.041}{-th}%
\TextAt{8}{507.006}{666.041}{roots}%
\TextAt{8}{79.226}{648.995}{of}%
\TextAt{8}{92.754}{648.995}{unity.}%
\TextAt{13}{127.486}{648.995}{L}%
\TextAt{21}{135.45}{647.201}{n}%
\TextAt{16}{144.744}{648.995}{\Character{106}}%
\TextAt{13}{151.722}{648.995}{K}%
\TextAt{8}{166.634}{648.995}{is}%
\TextAt{8}{178.595}{648.995}{therefore}%
\TextAt{8}{227.893}{648.995}{the}%
\TextAt{13}{248.247}{648.995}{p}%
\TextAt{21}{254.122}{653.333}{n}%
\TextAt{8}{259.758}{648.995}{-th}%
\TextAt{8}{278.809}{648.995}{cyclotomic}%
\TextAt{8}{337.216}{648.995}{extension}%
\TextAt{18}{389.179}{648.995}{Q}%
\TextAt{21}{398.478}{647.201}{p}%
\TextAt{9}{403.238}{648.995}{\Character{40}}%
\TextAt{13}{407.791}{648.995}{\Character{22}}%
\TextAt{21}{414.834}{647.201}{p}%
\TextAt{22}{419.097}{649.471}{n}%
\TextAt{9}{424.694}{648.995}{\Character{41}}%
\TextAt{16}{432.903}{648.995}{\Character{106}}%
\TextAt{18}{439.881}{648.995}{Q}%
\TextAt{21}{449.18}{647.201}{p}%
\TextAt{8}{453.941}{648.995}{.}%
\TextAt{8}{462.981}{648.995}{The}%
\TextAt{8}{487.232}{648.995}{following}%
\TextAt{8}{79.226}{631.949}{theorem}%
\TextAt{8}{124.751}{631.949}{shows}%
\TextAt{8}{158.369}{631.949}{the}%
\TextAt{8}{178.526}{631.949}{complete}%
\TextAt{8}{227.948}{631.949}{analogy}%
\TextAt{8}{271.214}{631.949}{of}%
\TextAt{8}{284.544}{631.949}{Lubin–Tate}%
\TextAt{8}{347.464}{631.949}{extensions}%
\TextAt{8}{403.845}{631.949}{with}%
\TextAt{8}{430.505}{631.949}{cyclotomic}%
\TextAt{8}{488.715}{631.949}{fields.}%
\TextAt{12}{79.226}{605.936}{Theorem}%
\TextAt{12}{135.702}{605.936}{(5.4).}%
\TextAt{13}{170.549}{605.936}{L}%
\TextAt{21}{178.514}{604.143}{n}%
\TextAt{16}{187.471}{605.936}{\Character{106}}%
\TextAt{13}{194.113}{605.936}{K}%
\TextAt{8}{208.36}{605.936}{is}%
\TextAt{8}{219.656}{605.936}{a}%
\TextAt{8}{228.943}{605.936}{totally}%
\TextAt{8}{265.883}{605.936}{ramified}%
\TextAt{8}{310.94}{605.936}{abelian}%
\TextAt{8}{351.12}{605.936}{extension}%
\TextAt{8}{402.417}{605.936}{of}%
\TextAt{8}{415.279}{605.936}{degree}%
\TextAt{13}{451.227}{605.936}{q}%
\TextAt{21}{456.846}{610.275}{n}%
\TextAt{3}{461.984}{610.275}{\Character{0}}%
\TextAt{17}{468.571}{610.275}{1}%
\TextAt{9}{473.303}{605.936}{\Character{40}}%
\TextAt{13}{477.855}{605.936}{q}%
\TextAt{16}{485.174}{605.936}{\Character{0}}%
\TextAt{9}{496.172}{605.936}{1\Character{41}}%
\TextAt{8}{510.007}{605.936}{with}%
\TextAt{8}{79.226}{588.89}{Galois}%
\TextAt{8}{115.139}{588.89}{group}%
\TextAt{13}{207.523}{571.844}{G}%
\TextAt{9}{216.756}{571.844}{\Character{40}}%
\TextAt{13}{221.308}{571.844}{L}%
\TextAt{21}{229.273}{570.051}{n}%
\TextAt{16}{238.23}{571.844}{\Character{106}}%
\TextAt{13}{244.872}{571.844}{K}%
\TextAt{9}{255.69}{571.844}{\Character{41}}%
\NativeMathAt{263.563}{571.844}{11.955}{\cong}%
\TextAt{8}{276.182}{571.844}{Aut}%
\TextAt{27}{296.016}{570.051}{o}%
\TextAt{22}{299.912}{568.712}{K}%
\TextAt{9}{307.453}{571.844}{\Character{40}}%
\TextAt{13}{312.005}{571.844}{F}%
\TextAt{9}{321.209}{571.844}{\Character{40}}%
\TextAt{13}{325.761}{571.844}{n}%
\TextAt{9}{332.749}{571.844}{\Character{41}\Character{41}}%
\NativeMathAt{345.174}{571.844}{11.955}{\cong}%
\TextAt{13}{357.794}{571.844}{U}%
\TextAt{21}{365.744}{570.051}{K}%
\TextAt{9}{373.947}{571.844}{\Character{47}}%
\TextAt{13}{379.8}{571.844}{U}%
\TextAt{17}{388.999}{577.955}{\Character{40}}%
\TextAt{21}{392.292}{577.955}{n}%
\TextAt{17}{397.43}{577.955}{\Character{41}}%
\TextAt{21}{387.75}{568.603}{K}%
\TextAt{13}{401.221}{571.844}{\Character{59}}%
\TextAt{8}{79.226}{548.323}{i.e.,}%
\TextAt{8}{101.331}{548.323}{for}%
\TextAt{8}{119.216}{548.323}{every}%
\TextAt{13}{150.108}{548.323}{\Character{27}}%
\NativeMathAt{160.512}{548.323}{11.955}{\in}%
\TextAt{13}{171.803}{548.323}{G}%
\TextAt{9}{181.036}{548.323}{\Character{40}}%
\TextAt{13}{185.589}{548.323}{L}%
\TextAt{21}{193.553}{546.529}{n}%
\TextAt{16}{202.51}{548.323}{\Character{106}}%
\TextAt{13}{209.152}{548.323}{K}%
\TextAt{9}{219.97}{548.323}{\Character{41}}%
\TextAt{8}{228.419}{548.323}{there}%
\TextAt{8}{258.331}{548.323}{is}%
\TextAt{8}{270.095}{548.323}{a}%
\TextAt{8}{279.85}{548.323}{unique}%
\TextAt{8}{317.892}{548.323}{class}%
\TextAt{13}{345.329}{548.323}{u}%
\TextAt{8}{355.312}{548.323}{mod}%
\TextAt{13}{380.75}{548.323}{U}%
\TextAt{17}{389.949}{554.433}{\Character{40}}%
\TextAt{21}{393.242}{554.433}{n}%
\TextAt{17}{398.381}{554.433}{\Character{41}}%
\TextAt{21}{388.7}{545.082}{K}%
\TextAt{8}{402.172}{548.323}{,}%
\TextAt{8}{409.321}{548.323}{with}%
\TextAt{13}{435.981}{548.323}{u}%
\NativeMathAt{445.964}{548.323}{11.955}{\in}%
\TextAt{13}{457.255}{548.323}{U}%
\TextAt{21}{465.206}{546.529}{K}%
\TextAt{8}{473.408}{548.323}{,}%
\TextAt{8}{480.558}{548.323}{such}%
\TextAt{8}{506.955}{548.323}{that}%
\TextAt{13}{227.941}{519.322}{\Character{21}}%
\TextAt{21}{234.769}{524.258}{\Character{27}}%
\TextAt{9}{243.714}{519.322}{\Character{61}}%
\TextAt{9}{256.124}{519.322}{\Character{91}}%
\TextAt{13}{259.391}{519.322}{u}%
\TextAt{9}{266.053}{519.322}{\Character{93}}%
\TextAt{21}{269.305}{517.528}{F}%
\TextAt{9}{276.338}{519.322}{\Character{40}}%
\TextAt{13}{280.891}{519.322}{\Character{21}}%
\TextAt{9}{287.719}{519.322}{\Character{41}}%
\TextAt{8}{316.182}{519.322}{for}%
\TextAt{13}{334.067}{519.322}{\Character{21}}%
\NativeMathAt{344.216}{519.322}{11.955}{\in}%
\TextAt{13}{355.507}{519.322}{F}%
\TextAt{9}{364.711}{519.322}{\Character{40}}%
\TextAt{13}{369.263}{519.322}{n}%
\TextAt{9}{376.251}{519.322}{\Character{41}}%
\TextAt{13}{380.803}{519.322}{\Character{58}}%
\TextAt{8}{79.226}{490.32}{Furthermore}%
\TextAt{8}{147.199}{490.32}{the}%
\TextAt{8}{167.53}{490.32}{following}%
\TextAt{8}{217.14}{490.32}{is}%
\TextAt{8}{229.078}{490.32}{true:}%
\TextAt{8}{258.69}{490.32}{let}%
\TextAt{13}{275.77}{490.32}{F}%
\TextAt{8}{289.046}{490.32}{be}%
\TextAt{8}{305.145}{490.32}{the}%
\TextAt{8}{325.476}{490.32}{Lubin–Tate}%
\TextAt{8}{388.571}{490.32}{module}%
\TextAt{13}{430.039}{490.32}{F}%
\TextAt{21}{437.617}{488.527}{e}%
\TextAt{8}{446.115}{490.32}{associated}%
\TextAt{8}{502.025}{490.32}{to}%
\TextAt{8}{516.51}{490.32}{the}%
\ErrataAt{79.19}{473.275}{455.81}{\emph{monic polynomial $e(X)\in\mathcal E_\pi$, and let $\lambda_n\in F(n)\setminus F(n-1)$. Then $\lambda_n$ is a prime element of $L_n$,}}%
\TextAt{8}{79.226}{456.229}{i.e.,}%
\TextAt{13}{101.331}{456.229}{L}%
\TextAt{21}{109.296}{454.435}{n}%
\TextAt{9}{118.253}{456.229}{\Character{61}}%
\TextAt{13}{130.678}{456.229}{K}%
\TextAt{9}{141.496}{456.229}{\Character{40}}%
\TextAt{13}{146.048}{456.229}{\Character{21}}%
\TextAt{21}{152.877}{454.435}{n}%
\TextAt{9}{158.513}{456.229}{\Character{41}}%
\TextAt{8}{163.066}{456.229}{,}%
\TextAt{8}{170.215}{456.229}{and}%
\TextAt{13}{178.637}{423.234}{\Character{30}}%
\TextAt{21}{185.565}{421.441}{n}%
\TextAt{9}{191.201}{423.234}{\Character{40}}%
\TextAt{13}{195.754}{423.234}{X}%
\TextAt{9}{206.409}{423.234}{\Character{41}}%
\TextAt{9}{214.264}{423.234}{\Character{61}}%
\TextAt{13}{233.313}{431.322}{e}%
\TextAt{21}{238.739}{435.66}{n}%
\TextAt{9}{244.375}{431.322}{\Character{40}}%
\TextAt{13}{248.927}{431.322}{X}%
\TextAt{9}{259.582}{431.322}{\Character{41}}%
\Vector{q 0 G 0 g 0.478 w 155.903 -347.777 m 197.545 -347.777 l S Q}%
\TextAt{13}{227.903}{415.034}{e}%
\TextAt{21}{233.328}{418.487}{n}%
\TextAt{3}{238.467}{418.487}{\Character{0}}%
\TextAt{17}{245.053}{418.487}{1}%
\TextAt{9}{249.785}{415.034}{\Character{40}}%
\TextAt{13}{254.338}{415.034}{X}%
\TextAt{9}{264.993}{415.034}{\Character{41}}%
\TextAt{9}{274.061}{423.234}{\Character{61}}%
\TextAt{13}{286.487}{423.234}{X}%
\TextAt{21}{297.142}{428.17}{q}%
\TextAt{22}{301.209}{430.983}{n}%
\TextAt{23}{305.81}{430.983}{\Character{0}}%
\TextAt{24}{311.566}{430.983}{1}%
\TextAt{17}{315.717}{428.17}{\Character{40}}%
\TextAt{21}{319.01}{428.17}{q}%
\TextAt{3}{323.077}{428.17}{\Character{0}}%
\TextAt{17}{329.663}{428.17}{1\Character{41}}%
\TextAt{9}{340.345}{423.234}{\Character{43}}%
\TextAt{16}{352.107}{423.234}{\Character{1}}%
\TextAt{16}{357.403}{423.234}{\Character{1}}%
\TextAt{16}{362.711}{423.234}{\Character{1}}%
\TextAt{9}{368.711}{423.234}{\Character{43}}%
\TextAt{13}{380.472}{423.234}{\Character{25}}%
\NativeMathAt{390.863}{423.234}{11.955}{\in}%
\TextAt{14}{402.154}{423.234}{o}%
\TextAt{21}{407.998}{421.441}{K}%
\TextAt{9}{416.2}{423.234}{\Character{91}}%
\TextAt{13}{419.452}{423.234}{X}%
\TextAt{9}{430.107}{423.234}{\Character{93}}%
\TextAt{8}{79.226}{390.127}{is}%
\TextAt{8}{90.99}{390.127}{its}%
\TextAt{8}{107.309}{390.127}{minimal}%
\TextAt{8}{152.834}{390.127}{polynomial.}%
\TextAt{8}{217.878}{390.127}{In}%
\TextAt{8}{232.499}{390.127}{particular}%
\TextAt{8}{286.489}{390.127}{one}%
\TextAt{8}{307.948}{390.127}{has}%
\TextAt{13}{328.822}{390.127}{N}%
\TextAt{21}{338.196}{388.267}{L}%
\TextAt{22}{343.955}{387.271}{n}%
\TextAt{3}{349.054}{388.267}{\Character{106}}%
\TextAt{21}{351.406}{388.267}{K}%
\TextAt{9}{359.609}{390.127}{\Character{40}}%
\TextAt{16}{364.161}{390.127}{\Character{0}}%
\TextAt{13}{373.46}{390.127}{\Character{21}}%
\TextAt{21}{380.288}{388.334}{n}%
\TextAt{9}{385.925}{390.127}{\Character{41}}%
\TextAt{9}{393.78}{390.127}{\Character{61}}%
\TextAt{13}{406.223}{390.127}{\Character{25}}%
\TextAt{8}{413.293}{390.127}{.}%
\TextAt{19}{79.226}{355.148}{Proof.}%
\TextAt{8}{116}{355.148}{If}%
\TextAt{13}{127.454}{355.148}{e}%
\TextAt{9}{132.88}{355.148}{\Character{40}}%
\TextAt{13}{137.432}{355.148}{X}%
\TextAt{9}{148.087}{355.148}{\Character{41}}%
\TextAt{9}{155.942}{355.148}{\Character{61}}%
\TextAt{13}{168.386}{355.148}{X}%
\TextAt{21}{179.041}{359.487}{q}%
\TextAt{9}{185.776}{355.148}{\Character{43}}%
\TextAt{13}{197.051}{355.148}{\Character{25}}%
\TextAt{9}{204.12}{355.148}{\Character{40}}%
\TextAt{13}{208.672}{355.148}{a}%
\TextAt{21}{214.817}{353.355}{q}%
\TextAt{3}{218.884}{353.355}{\Character{0}}%
\TextAt{17}{225.471}{353.355}{1}%
\TextAt{13}{230.203}{355.148}{X}%
\TextAt{21}{240.858}{359.487}{q}%
\TextAt{3}{244.925}{359.487}{\Character{0}}%
\TextAt{17}{251.511}{359.487}{1}%
\TextAt{9}{258.414}{355.148}{\Character{43}}%
\TextAt{16}{269.688}{355.148}{\Character{1}}%
\TextAt{16}{274.984}{355.148}{\Character{1}}%
\TextAt{16}{280.292}{355.148}{\Character{1}}%
\TextAt{9}{285.806}{355.148}{\Character{43}}%
\TextAt{13}{297.081}{355.148}{a}%
\TextAt{17}{303.226}{353.355}{2}%
\TextAt{13}{307.958}{355.148}{X}%
\TextAt{17}{318.613}{359.487}{2}%
\TextAt{9}{323.346}{355.148}{\Character{41}}%
\TextAt{9}{330.053}{355.148}{\Character{43}}%
\TextAt{13}{341.343}{355.148}{\Character{25}}%
\TextAt{13}{348.397}{355.148}{X}%
\TextAt{8}{362.727}{355.148}{is}%
\TextAt{8}{374.253}{355.148}{a}%
\TextAt{8}{383.77}{355.148}{Lubin–Tate}%
\TextAt{8}{446.452}{355.148}{polynomial,}%
\TextAt{8}{510.007}{355.148}{then}%
\TextAt{13}{140.478}{321.8}{\Character{30}}%
\TextAt{21}{147.407}{320.006}{n}%
\TextAt{9}{153.043}{321.8}{\Character{40}}%
\TextAt{13}{157.595}{321.8}{X}%
\TextAt{9}{168.25}{321.8}{\Character{41}}%
\TextAt{9}{176.105}{321.8}{\Character{61}}%
\TextAt{13}{195.155}{329.887}{e}%
\TextAt{21}{200.58}{334.226}{n}%
\TextAt{9}{206.217}{329.887}{\Character{40}}%
\TextAt{13}{210.769}{329.887}{X}%
\TextAt{9}{221.424}{329.887}{\Character{41}}%
\Vector{q 0 G 0 g 0.478 w 117.744 -449.212 m 159.387 -449.212 l S Q}%
\TextAt{13}{189.744}{313.599}{e}%
\TextAt{21}{195.17}{317.053}{n}%
\TextAt{3}{200.308}{317.053}{\Character{0}}%
\TextAt{17}{206.895}{317.053}{1}%
\TextAt{9}{211.627}{313.599}{\Character{40}}%
\TextAt{13}{216.179}{313.599}{X}%
\TextAt{9}{226.834}{313.599}{\Character{41}}%
\TextAt{9}{176.123}{294.693}{\Character{61}}%
\TextAt{13}{188.549}{294.693}{e}%
\TextAt{21}{193.974}{299.629}{n}%
\TextAt{3}{199.113}{299.629}{\Character{0}}%
\TextAt{17}{205.699}{299.629}{1}%
\TextAt{9}{210.431}{294.693}{\Character{40}}%
\TextAt{13}{214.984}{294.693}{X}%
\TextAt{9}{225.639}{294.693}{\Character{41}}%
\TextAt{21}{230.191}{299.629}{q}%
\TextAt{3}{234.258}{299.629}{\Character{0}}%
\TextAt{17}{240.844}{299.629}{1}%
\TextAt{9}{248.233}{294.693}{\Character{43}}%
\TextAt{13}{259.995}{294.693}{\Character{25}}%
\TextAt{9}{267.064}{294.693}{\Character{40}}%
\TextAt{13}{271.616}{294.693}{a}%
\TextAt{21}{277.761}{292.9}{q}%
\TextAt{3}{281.828}{292.9}{\Character{0}}%
\TextAt{17}{288.414}{292.9}{1}%
\TextAt{13}{293.147}{294.693}{e}%
\TextAt{21}{298.572}{299.629}{n}%
\TextAt{3}{303.71}{299.629}{\Character{0}}%
\TextAt{17}{310.297}{299.629}{1}%
\TextAt{9}{315.029}{294.693}{\Character{40}}%
\TextAt{13}{319.582}{294.693}{X}%
\TextAt{9}{330.237}{294.693}{\Character{41}}%
\TextAt{21}{334.789}{299.629}{q}%
\TextAt{3}{338.856}{299.629}{\Character{0}}%
\TextAt{17}{345.442}{299.629}{2}%
\TextAt{9}{352.831}{294.693}{\Character{43}}%
\TextAt{16}{364.593}{294.693}{\Character{1}}%
\TextAt{16}{369.889}{294.693}{\Character{1}}%
\TextAt{16}{375.197}{294.693}{\Character{1}}%
\TextAt{9}{381.197}{294.693}{\Character{43}}%
\TextAt{13}{392.958}{294.693}{a}%
\TextAt{17}{399.103}{292.9}{2}%
\TextAt{13}{403.835}{294.693}{e}%
\TextAt{21}{409.261}{299.629}{n}%
\TextAt{3}{414.399}{299.629}{\Character{0}}%
\TextAt{17}{420.986}{299.629}{1}%
\TextAt{9}{425.718}{294.693}{\Character{40}}%
\TextAt{13}{430.27}{294.693}{X}%
\TextAt{9}{440.925}{294.693}{\Character{41}\Character{41}}%
\TextAt{9}{452.665}{294.693}{\Character{43}}%
\TextAt{13}{464.448}{294.693}{\Character{25}}%
\TextAt{8}{79.226}{265.692}{is}%
\TextAt{8}{89.848}{265.692}{an}%
\TextAt{8}{104.966}{265.692}{Eisenstein}%
\TextAt{8}{159.379}{265.692}{polynomial}%
\TextAt{8}{218.731}{265.692}{of}%
\TextAt{8}{230.919}{265.692}{degree}%
\TextAt{13}{266.193}{265.692}{q}%
\TextAt{21}{271.812}{270.03}{n}%
\TextAt{3}{276.95}{270.03}{\Character{0}}%
\TextAt{17}{283.537}{270.03}{1}%
\TextAt{9}{288.269}{265.692}{\Character{40}}%
\TextAt{13}{292.821}{265.692}{q}%
\TextAt{16}{298.762}{265.692}{\Character{0}}%
\TextAt{9}{308.383}{265.692}{1\Character{41}}%
\TextAt{8}{318.788}{265.692}{.}%
\TextAt{8}{326.856}{265.692}{If}%
\TextAt{13}{337.406}{265.692}{F}%
\TextAt{8}{349.365}{265.692}{is}%
\TextAt{8}{359.988}{265.692}{the}%
\TextAt{8}{379.002}{265.692}{Lubin–Tate}%
\TextAt{8}{440.78}{265.692}{module}%
\TextAt{8}{480.932}{265.692}{associated}%
\TextAt{8}{79.226}{248.646}{to}%
\TextAt{13}{94.001}{248.646}{e}%
\TextAt{8}{99.426}{248.646}{,}%
\TextAt{8}{107.156}{248.646}{and}%
\TextAt{13}{130.383}{248.646}{\Character{21}}%
\TextAt{21}{137.211}{246.853}{n}%
\NativeMathAt{146.96}{248.646}{11.955}{\in}%
\TextAt{13}{159.042}{248.646}{F}%
\TextAt{9}{168.246}{248.646}{\Character{40}}%
\TextAt{13}{172.798}{248.646}{n}%
\TextAt{9}{179.786}{248.646}{\Character{41}}%
\TextAt{16}{187.311}{248.646}{\Character{110}}%
\TextAt{13}{196.262}{248.646}{F}%
\TextAt{9}{205.466}{248.646}{\Character{40}}%
\TextAt{13}{210.018}{248.646}{n}%
\TextAt{16}{219.979}{248.646}{\Character{0}}%
\TextAt{9}{232.251}{248.646}{1\Character{41}}%
\TextAt{8}{242.656}{248.646}{,}%
\TextAt{8}{250.386}{248.646}{then}%
\TextAt{13}{277.51}{248.646}{\Character{21}}%
\TextAt{21}{284.339}{246.853}{n}%
\TextAt{8}{294.337}{248.646}{is}%
\TextAt{8}{306.565}{248.646}{clearly}%
\TextAt{8}{344.425}{248.646}{a}%
\TextAt{8}{354.645}{248.646}{zero}%
\TextAt{8}{379.821}{248.646}{of}%
\TextAt{8}{393.615}{248.646}{this}%
\TextAt{8}{416.902}{248.646}{Eisenstein}%
\TextAt{8}{472.922}{248.646}{polynomial,}%
\TextAt{8}{79.226}{231.6}{and}%
\TextAt{8}{102.76}{231.6}{is}%
\TextAt{8}{115.296}{231.6}{therefore}%
\TextAt{8}{165.167}{231.6}{a}%
\TextAt{8}{175.694}{231.6}{prime}%
\TextAt{8}{209.629}{231.6}{element}%
\TextAt{8}{253.642}{231.6}{of}%
\TextAt{8}{267.743}{231.6}{the}%
\TextAt{8}{288.671}{231.6}{totally}%
\TextAt{8}{326.851}{231.6}{ramified}%
\TextAt{8}{373.147}{231.6}{extension}%
\TextAt{13}{425.685}{231.6}{K}%
\TextAt{9}{436.502}{231.6}{\Character{40}}%
\TextAt{13}{441.055}{231.6}{\Character{21}}%
\TextAt{21}{447.883}{229.807}{n}%
\TextAt{9}{453.519}{231.6}{\Character{41}}%
\TextAt{16}{462.707}{231.6}{\Character{106}}%
\TextAt{13}{470.663}{231.6}{K}%
\TextAt{8}{486.15}{231.6}{of}%
\TextAt{8}{500.251}{231.6}{degree}%
\TextAt{13}{79.226}{214.554}{q}%
\TextAt{21}{84.845}{218.892}{n}%
\TextAt{3}{89.984}{218.892}{\Character{0}}%
\TextAt{17}{96.57}{218.892}{1}%
\TextAt{9}{101.302}{214.554}{\Character{40}}%
\TextAt{13}{105.855}{214.554}{q}%
\TextAt{16}{113.962}{214.554}{\Character{0}}%
\TextAt{9}{125.749}{214.554}{1\Character{41}}%
\TextAt{8}{136.155}{214.554}{.}%
\TextAt{8}{144.575}{214.554}{Each}%
\TextAt{13}{173.592}{214.554}{\Character{27}}%
\NativeMathAt{183.995}{214.554}{11.955}{\in}%
\TextAt{13}{195.286}{214.554}{G}%
\TextAt{9}{204.52}{214.554}{\Character{40}}%
\Vector{q 0 G 0 g 0.478 w 137.072 -549.603 m 147.89 -549.603 l S Q}%
\TextAt{13}{209.072}{214.554}{K}%
\TextAt{16}{223.211}{214.554}{\Character{106}}%
\TextAt{13}{229.852}{214.554}{K}%
\TextAt{9}{240.67}{214.554}{\Character{41}}%
\TextAt{8}{249.038}{214.554}{induces}%
\TextAt{8}{290.631}{214.554}{an}%
\TextAt{8}{306.808}{214.554}{automorphism}%
\TextAt{8}{384.195}{214.554}{of}%
\TextAt{13}{397.443}{214.554}{F}%
\TextAt{9}{406.646}{214.554}{\Character{40}}%
\TextAt{13}{411.199}{214.554}{n}%
\TextAt{9}{418.186}{214.554}{\Character{41}}%
\TextAt{8}{422.739}{214.554}{.}%
\TextAt{8}{431.16}{214.554}{We}%
\TextAt{8}{451.222}{214.554}{therefore}%
\TextAt{8}{500.239}{214.554}{obtain}%
\TextAt{8}{79.226}{197.508}{a}%
\TextAt{8}{88.982}{197.508}{homomorphism}%
\TextAt{13}{202.541}{180.462}{G}%
\TextAt{9}{211.775}{180.462}{\Character{40}}%
\TextAt{13}{216.327}{180.462}{L}%
\TextAt{21}{224.292}{178.669}{n}%
\TextAt{16}{233.249}{180.462}{\Character{106}}%
\TextAt{13}{239.891}{180.462}{K}%
\TextAt{9}{250.708}{180.462}{\Character{41}}%
\NativeArrowAt{258.581}{180.462}{11.955}{19.26}{\longrightarrow}%
\TextAt{8}{281.163}{180.462}{Aut}%
\TextAt{27}{300.997}{178.669}{o}%
\TextAt{22}{304.893}{177.329}{K}%
\TextAt{9}{312.434}{180.462}{\Character{40}}%
\TextAt{13}{316.987}{180.462}{F}%
\TextAt{9}{326.19}{180.462}{\Character{40}}%
\TextAt{13}{330.743}{180.462}{n}%
\TextAt{9}{337.73}{180.462}{\Character{41}\Character{41}}%
\NativeMathAt{350.156}{180.462}{11.955}{\cong}%
\TextAt{13}{362.775}{180.462}{U}%
\TextAt{21}{370.725}{178.669}{K}%
\TextAt{9}{378.928}{180.462}{\Character{47}}%
\TextAt{13}{384.781}{180.462}{U}%
\TextAt{17}{393.98}{186.572}{\Character{40}}%
\TextAt{21}{397.273}{186.572}{n}%
\TextAt{17}{402.411}{186.572}{\Character{41}}%
\TextAt{21}{392.731}{177.221}{K}%
\TextAt{13}{406.203}{180.462}{\Character{58}}%
\TextAt{8}{79.226}{156.94}{It}%
\TextAt{8}{91.899}{156.94}{is}%
\TextAt{8}{103.663}{156.94}{injective}%
\TextAt{8}{150.156}{156.94}{because}%
\TextAt{13}{193.458}{156.94}{L}%
\TextAt{21}{201.422}{155.147}{n}%
\TextAt{8}{210.956}{156.94}{is}%
\TextAt{8}{222.72}{156.94}{generated}%
\TextAt{8}{276.052}{156.94}{by}%
\TextAt{13}{292.311}{156.94}{F}%
\TextAt{9}{301.515}{156.94}{\Character{40}}%
\TextAt{13}{306.067}{156.94}{n}%
\TextAt{9}{313.054}{156.94}{\Character{41}}%
\TextAt{8}{317.607}{156.94}{,}%
\TextAt{8}{324.756}{156.94}{and}%
\TextAt{8}{347.519}{156.94}{it}%
\TextAt{8}{359.223}{156.94}{is}%
\TextAt{8}{370.987}{156.94}{surjective}%
\TextAt{8}{423.398}{156.94}{because}%
\TextAt{8}{164.393}{127.939}{\char35\relax }%
\TextAt{13}{174.149}{127.939}{G}%
\TextAt{9}{183.383}{127.939}{\Character{40}}%
\TextAt{13}{187.935}{127.939}{L}%
\TextAt{21}{195.899}{126.146}{n}%
\TextAt{16}{204.856}{127.939}{\Character{106}}%
\TextAt{13}{211.498}{127.939}{K}%
\TextAt{9}{222.316}{127.939}{\Character{41}}%
\NativeMathAt{230.189}{127.939}{11.955}{\geq}%
\TextAt{9}{242.808}{127.939}{\Character{91}}%
\TextAt{13}{246.06}{127.939}{K}%
\TextAt{9}{256.878}{127.939}{\Character{40}}%
\TextAt{13}{261.43}{127.939}{\Character{21}}%
\TextAt{21}{268.259}{126.146}{n}%
\TextAt{9}{273.895}{127.939}{\Character{41}}%
\TextAt{9}{281.75}{127.939}{\Character{58}}%
\TextAt{13}{288.341}{127.939}{K}%
\TextAt{9}{299.158}{127.939}{\Character{93}}%
\TextAt{9}{305.709}{127.939}{\Character{61}}%
\TextAt{13}{318.156}{127.939}{q}%
\TextAt{21}{323.775}{132.875}{n}%
\TextAt{3}{328.914}{132.875}{\Character{0}}%
\TextAt{17}{335.5}{132.875}{1}%
\TextAt{9}{340.232}{127.939}{\Character{40}}%
\TextAt{13}{344.785}{127.939}{q}%
\TextAt{16}{353.061}{127.939}{\Character{0}}%
\TextAt{9}{365.016}{127.939}{1\Character{41}}%
\TextAt{9}{378.717}{127.939}{\Character{61}}%
\TextAt{8}{391.167}{127.939}{\char35\relax }%
\TextAt{13}{400.923}{127.939}{U}%
\TextAt{21}{408.873}{126.146}{K}%
\TextAt{9}{417.076}{127.939}{\Character{47}}%
\TextAt{13}{422.929}{127.939}{U}%
\TextAt{17}{432.128}{134.05}{\Character{40}}%
\TextAt{21}{435.421}{134.05}{n}%
\TextAt{17}{440.559}{134.05}{\Character{41}}%
\TextAt{21}{430.879}{124.698}{K}%
\TextAt{13}{444.351}{127.939}{\Character{58}}%
\TextAt{8}{79.226}{98.938}{This}%
\TextAt{8}{105.946}{98.938}{proves}%
\TextAt{8}{142.11}{98.938}{the}%
\TextAt{8}{162.267}{98.938}{theorem.}%
\Vector{q 0 G 0 g 0.399 w 452.334 -675.062 m 452.334 -666.992 l S Q}%
\Vector{q 0 G 0 g 0.399 w 452.534 -667.191 m 459.707 -667.191 l S Q}%
\Vector{q 0 G 0 g 0.399 w 452.534 -674.863 m 459.707 -674.863 l S Q}%
\Vector{q 0 G 0 g 0.399 w 459.906 -675.062 m 459.906 -666.992 l S Q}%
\TextAt{8}{297.211}{43.033}{319}%
\Destination{78.23}{809.815}{nav.page.319}%
\Destination{79.36}{776.155}{nav.print.349}%
\Destination{79.36}{619.955}{nav.number.5.5.4}%
\Destination{79.226}{498}{errata.196}%
}

% Page 332 (body page 320)
\Page{%
\TextAt{2}{79.226}{796.943}{V.}%
\TextAt{2}{94.073}{796.943}{Local}%
\TextAt{2}{123.615}{796.943}{Class}%
\TextAt{2}{152.208}{796.943}{Field}%
\TextAt{2}{179.939}{796.943}{Theory}%
\Vector{q 0 G 0 g 0.299 w 7.226 18.729 m 460.769 18.729 l S Q}%
\TextAt{8}{97.159}{761.124}{Generalizing}%
\TextAt{8}{165.644}{761.124}{the}%
\TextAt{8}{186.571}{761.124}{explicit}%
\TextAt{8}{228.634}{761.124}{norm}%
\TextAt{8}{259.974}{761.124}{residue}%
\TextAt{8}{300.472}{761.124}{symbol}%
\TextAt{8}{341.304}{761.124}{of}%
\TextAt{8}{355.405}{761.124}{the}%
\TextAt{8}{376.332}{761.124}{cyclotomic}%
\TextAt{8}{435.312}{761.124}{fields}%
\TextAt{18}{466.054}{761.124}{Q}%
\TextAt{21}{475.353}{759.33}{p}%
\TextAt{9}{480.114}{761.124}{\Character{40}}%
\TextAt{13}{484.666}{761.124}{\Character{22}}%
\TextAt{21}{491.709}{759.33}{p}%
\TextAt{22}{495.972}{761.6}{n}%
\TextAt{9}{501.569}{761.124}{\Character{41}}%
\TextAt{16}{510.755}{761.124}{\Character{106}}%
\TextAt{18}{518.71}{761.124}{Q}%
\TextAt{21}{528.009}{759.33}{p}%
\TextAt{8}{79.226}{744.078}{(see}%
\TextAt{8}{103.314}{744.078}{(2.4)),}%
\TextAt{8}{139.871}{744.078}{we}%
\TextAt{8}{157.718}{744.078}{obtain}%
\TextAt{8}{194.766}{744.078}{the}%
\TextAt{8}{215.542}{744.078}{following}%
\TextAt{8}{265.597}{744.078}{explicit}%
\TextAt{8}{307.51}{744.078}{formula}%
\TextAt{8}{351.061}{744.078}{for}%
\TextAt{8}{369.566}{744.078}{the}%
\TextAt{8}{390.342}{744.078}{norm}%
\TextAt{8}{421.532}{744.078}{residue}%
\TextAt{8}{461.879}{744.078}{symbol}%
\TextAt{8}{502.56}{744.078}{of}%
\TextAt{8}{516.51}{744.078}{the}%
\TextAt{8}{79.226}{727.032}{Lubin–Tate}%
\TextAt{8}{142.146}{727.032}{extensions.}%
\TextAt{12}{79.226}{701.019}{Theorem}%
\TextAt{12}{137.079}{701.019}{(5.5).}%
\TextAt{8}{173.124}{701.019}{For}%
\TextAt{8}{194.823}{701.019}{the}%
\TextAt{8}{215.708}{701.019}{field}%
\TextAt{13}{241.795}{701.019}{L}%
\TextAt{21}{249.759}{699.226}{n}%
\TextAt{16}{259.959}{701.019}{\Character{106}}%
\TextAt{13}{267.844}{701.019}{K}%
\TextAt{8}{283.288}{701.019}{of}%
\TextAt{13}{297.347}{701.019}{\Character{25}}%
\TextAt{21}{304.417}{705.358}{n}%
\TextAt{8}{310.053}{701.019}{-division}%
\TextAt{8}{357.993}{701.019}{points}%
\TextAt{8}{393.906}{701.019}{and}%
\TextAt{8}{417.398}{701.019}{for}%
\TextAt{13}{436.013}{701.019}{a}%
\TextAt{9}{446.721}{701.019}{\Character{61}}%
\TextAt{13}{460.389}{701.019}{u\Character{25}}%
\TextAt{21}{474.121}{705.358}{v}%
\TextAt{22}{478.245}{704.018}{K}%
\TextAt{17}{485.288}{705.358}{\Character{40}}%
\TextAt{21}{488.581}{705.358}{a}%
\TextAt{17}{493.079}{705.358}{\Character{41}}%
\NativeMathAt{501.434}{701.019}{11.955}{\in}%
\TextAt{13}{513.968}{701.019}{K}%
\TextAt{3}{524.785}{705.358}{\Character{3}}%
\TextAt{8}{529.518}{701.019}{,}%
\TextAt{13}{79.226}{683.973}{u}%
\NativeMathAt{89.209}{683.973}{11.955}{\in}%
\TextAt{13}{100.5}{683.973}{U}%
\TextAt{21}{108.451}{682.18}{K}%
\TextAt{8}{116.653}{683.973}{,}%
\TextAt{8}{123.803}{683.973}{one}%
\TextAt{8}{145.262}{683.973}{has}%
\TextAt{9}{203.976}{666.927}{\Character{40}}%
\TextAt{13}{208.529}{666.927}{a\Character{59}}%
\TextAt{13}{219.898}{666.927}{L}%
\TextAt{21}{227.882}{665.134}{n}%
\TextAt{16}{236.839}{666.927}{\Character{106}}%
\TextAt{13}{243.481}{666.927}{K}%
\TextAt{9}{254.299}{666.927}{\Character{41}}%
\TextAt{13}{258.851}{666.927}{\Character{21}}%
\TextAt{9}{269}{666.927}{\Character{61}}%
\TextAt{9}{281.41}{666.927}{\Character{91}}%
\TextAt{13}{284.678}{666.927}{u}%
\TextAt{3}{291.34}{671.864}{\Character{0}}%
\TextAt{17}{297.927}{671.864}{1}%
\TextAt{9}{302.659}{666.927}{\Character{93}}%
\TextAt{21}{305.911}{665.134}{F}%
\TextAt{9}{312.944}{666.927}{\Character{40}}%
\TextAt{13}{317.496}{666.927}{\Character{21}}%
\TextAt{9}{324.325}{666.927}{\Character{41}}%
\TextAt{13}{328.877}{666.927}{\Character{59}}%
\TextAt{13}{358.012}{666.927}{\Character{21}}%
\NativeMathAt{368.181}{666.927}{11.955}{\in}%
\TextAt{13}{379.472}{666.927}{F}%
\TextAt{9}{388.675}{666.927}{\Character{40}}%
\TextAt{13}{393.228}{666.927}{n}%
\TextAt{9}{400.215}{666.927}{\Character{41}}%
\TextAt{13}{404.768}{666.927}{\Character{58}}%
\TextAt{19}{79.226}{625.473}{Proof.}%
\TextAt{8}{116}{625.473}{The}%
\TextAt{8}{140.346}{625.473}{proof}%
\TextAt{8}{171.208}{625.473}{is}%
\TextAt{8}{183.263}{625.473}{the}%
\TextAt{8}{203.712}{625.473}{same}%
\TextAt{8}{233.33}{625.473}{as}%
\TextAt{8}{247.992}{625.473}{that}%
\TextAt{8}{273.653}{625.473}{of}%
\TextAt{8}{287.275}{625.473}{(2.4).}%
\TextAt{8}{320.677}{625.473}{Let}%
\TextAt{13}{341.939}{625.473}{\Character{27}}%
\NativeMathAt{352.84}{625.473}{11.955}{\in}%
\TextAt{13}{364.628}{625.473}{G}%
\TextAt{9}{373.862}{625.473}{\Character{40}}%
\TextAt{13}{378.414}{625.473}{L}%
\TextAt{21}{386.379}{623.68}{n}%
\TextAt{16}{395.834}{625.473}{\Character{106}}%
\TextAt{13}{402.973}{625.473}{K}%
\TextAt{9}{413.791}{625.473}{\Character{41}}%
\TextAt{8}{422.533}{625.473}{be}%
\TextAt{8}{438.749}{625.473}{the}%
\TextAt{8}{459.197}{625.473}{automorphism}%
\TextAt{8}{79.226}{608.427}{such}%
\TextAt{8}{105.623}{608.427}{that}%
\TextAt{13}{130.992}{608.427}{\Character{21}}%
\TextAt{21}{137.821}{612.765}{\Character{27}}%
\TextAt{9}{146.765}{608.427}{\Character{61}}%
\TextAt{9}{159.175}{608.427}{\Character{91}}%
\TextAt{13}{162.442}{608.427}{u}%
\TextAt{3}{169.104}{612.765}{\Character{0}}%
\TextAt{17}{175.691}{612.765}{1}%
\TextAt{9}{180.423}{608.427}{\Character{93}}%
\TextAt{21}{183.675}{606.634}{F}%
\TextAt{9}{190.708}{608.427}{\Character{40}}%
\TextAt{13}{195.261}{608.427}{\Character{21}}%
\TextAt{9}{202.089}{608.427}{\Character{41}}%
\TextAt{8}{206.641}{608.427}{,}%
\TextAt{13}{213.791}{608.427}{\Character{21}}%
\NativeMathAt{223.94}{608.427}{11.955}{\in}%
\TextAt{13}{235.231}{608.427}{F}%
\TextAt{9}{244.435}{608.427}{\Character{40}}%
\TextAt{13}{248.987}{608.427}{n}%
\TextAt{9}{255.974}{608.427}{\Character{41}}%
\TextAt{8}{260.527}{608.427}{.}%
\ErrataAt{97.1}{588.392}{437.9}{Let $\widetilde\sigma\in\operatorname{Frob}(\widetilde L_n\mid K)$ satisfy $\sigma=\widetilde\sigma|_{L_n}$ and $d_K(\widetilde\sigma)=1$. We view}%
\TextAt{20}{79.446}{571.346}{\Character{101}}%
\TextAt{13}{79.226}{571.346}{\Character{27}}%
\TextAt{8}{90.936}{571.346}{as}%
\TextAt{8}{106.036}{571.346}{an}%
\TextAt{8}{123.024}{571.346}{automorphism}%
\TextAt{8}{201.224}{571.346}{of}%
\TextAt{8}{215.283}{571.346}{the}%
\TextAt{8}{236.17}{571.346}{completion}%
\TextAt{20}{297.722}{574.368}{\Character{98}}%
\TextAt{13}{296.735}{571.346}{L}%
\TextAt{21}{304.7}{569.553}{n}%
\TextAt{9}{314.9}{571.346}{\Character{61}}%
\TextAt{13}{328.57}{571.346}{L}%
\TextAt{21}{336.534}{569.553}{n}%
\TextAt{20}{344.909}{574.368}{\Character{98}}%
\TextAt{13}{342.17}{571.346}{K}%
\TextAt{8}{357.615}{571.346}{of}%
\TextAt{20}{372.662}{574.368}{\Character{101}}%
\TextAt{13}{371.675}{571.346}{L}%
\TextAt{21}{379.64}{569.553}{n}%
\TextAt{8}{385.276}{571.346}{.}%
\TextAt{8}{395.914}{571.346}{Let}%
\TextAt{9}{417.613}{571.346}{\Character{6}}%
\TextAt{8}{430.694}{571.346}{be}%
\TextAt{8}{447.348}{571.346}{the}%
\TextAt{8}{468.234}{571.346}{fixed}%
\TextAt{8}{497.25}{571.346}{field}%
\TextAt{8}{523.337}{571.346}{of}%
\TextAt{20}{79.446}{554.3}{\Character{101}}%
\TextAt{13}{79.226}{554.3}{\Character{27}}%
\TextAt{8}{86.309}{554.3}{.}%
\TextAt{8}{95.555}{554.3}{Since}%
\TextAt{13}{126.378}{554.3}{f}%
\TextAt{17}{132.148}{552.441}{\Character{6}}%
\TextAt{3}{138.264}{552.441}{\Character{106}}%
\TextAt{21}{140.617}{552.441}{K}%
\TextAt{9}{152.593}{554.3}{\Character{61}}%
\TextAt{13}{165.472}{554.3}{d}%
\TextAt{21}{171.555}{552.507}{K}%
\TextAt{9}{179.758}{554.3}{\Character{40}}%
\TextAt{20}{184.53}{554.3}{\Character{101}}%
\TextAt{13}{184.31}{554.3}{\Character{27}}%
\TextAt{9}{191.392}{554.3}{\Character{41}}%
\TextAt{9}{199.701}{554.3}{\Character{61}}%
\TextAt{9}{212.577}{554.3}{1}%
\TextAt{8}{218.451}{554.3}{,}%
\TextAt{9}{225.932}{554.3}{\Character{6}}%
\TextAt{16}{238.161}{554.3}{\Character{106}}%
\TextAt{13}{245.256}{554.3}{K}%
\TextAt{8}{260.237}{554.3}{is}%
\TextAt{8}{272.267}{554.3}{totally}%
\TextAt{8}{309.94}{554.3}{ramified.}%
\TextAt{8}{360.815}{554.3}{It}%
\TextAt{8}{373.753}{554.3}{has}%
\TextAt{8}{394.893}{554.3}{degree}%
\TextAt{13}{431.574}{554.3}{q}%
\TextAt{21}{437.194}{558.639}{n}%
\TextAt{3}{442.332}{558.639}{\Character{0}}%
\TextAt{17}{448.918}{558.639}{1}%
\TextAt{9}{453.651}{554.3}{\Character{40}}%
\TextAt{13}{458.203}{554.3}{q}%
\TextAt{16}{466.66}{554.3}{\Character{0}}%
\TextAt{9}{478.797}{554.3}{1\Character{41}}%
\TextAt{8}{493.365}{554.3}{because}%
\TextAt{9}{79.226}{537.254}{\Character{6}}%
\TextAt{16}{90.337}{537.254}{\Character{92}}%
\TextAt{20}{103.702}{540.276}{\Character{98}}%
\TextAt{13}{100.964}{537.254}{K}%
\TextAt{9}{115.102}{537.254}{\Character{61}}%
\TextAt{13}{127.528}{537.254}{K}%
\TextAt{8}{142.243}{537.254}{and}%
\TextAt{20}{165.912}{540.276}{\Character{98}}%
\TextAt{9}{165.006}{537.254}{\Character{6}}%
\TextAt{9}{176.77}{537.254}{\Character{61}}%
\TextAt{9}{189.191}{537.254}{\Character{6}}%
\TextAt{20}{200.399}{540.276}{\Character{98}}%
\TextAt{13}{197.661}{537.254}{K}%
\TextAt{9}{211.799}{537.254}{\Character{61}}%
\TextAt{20}{225.211}{540.276}{\Character{98}}%
\TextAt{13}{224.225}{537.254}{L}%
\TextAt{21}{232.189}{535.461}{n}%
\TextAt{8}{237.825}{537.254}{.}%
\TextAt{8}{246.274}{537.254}{Consequently}%
\TextAt{9}{318.854}{537.254}{\Character{91}\Character{6}}%
\TextAt{9}{333.858}{537.254}{\Character{58}}%
\TextAt{13}{340.453}{537.254}{K}%
\TextAt{9}{351.271}{537.254}{\Character{93}}%
\TextAt{9}{357.822}{537.254}{\Character{61}}%
\TextAt{9}{370.244}{537.254}{\Character{91}}%
\TextAt{20}{374.507}{540.276}{\Character{98}}%
\TextAt{13}{373.52}{537.254}{L}%
\TextAt{21}{381.485}{535.461}{n}%
\TextAt{9}{390.442}{537.254}{\Character{58}}%
\TextAt{20}{399.753}{540.276}{\Character{98}}%
\TextAt{13}{397.014}{537.254}{K}%
\TextAt{9}{407.832}{537.254}{\Character{93}}%
\TextAt{9}{414.383}{537.254}{\Character{61}}%
\TextAt{9}{426.805}{537.254}{\Character{91}}%
\TextAt{13}{430.082}{537.254}{L}%
\TextAt{21}{438.046}{535.461}{n}%
\TextAt{9}{447.003}{537.254}{\Character{58}}%
\TextAt{13}{453.576}{537.254}{K}%
\TextAt{9}{464.393}{537.254}{\Character{93}}%
\TextAt{8}{467.645}{537.254}{.}%
\ErrataAt{97.1}{517.233}{437.9}{Write $\bar\pi$ for the prime element of the statement, and set $\pi=u\bar\pi$. By \hyperlink{nav.number.5.4.6}{(4.6)}, choose}%
\ErrataAt{79.19}{500.187}{455.81}{monic Lubin--Tate polynomials $e\in\mathcal E_\pi$ and $\bar e\in\mathcal E_{\bar\pi}$, and put $\bar F=F_{\bar e}$. By \hyperlink{nav.number.5.2.3}{(2.3)}, there exists}%
\ErrataAt{79.19}{487.7}{455.81}{$\theta(X)=\epsilon X+\cdots\in\mathfrak o_{\widehat K}[[X]]$ with $\epsilon\in U_{\widehat K}$, such that}%
\TextAt{13}{181.049}{471.172}{\Character{18}}%
\TextAt{21}{186.829}{476.109}{\Character{39}}%
\TextAt{9}{196.167}{471.172}{\Character{61}}%
\TextAt{13}{208.592}{471.172}{\Character{18}}%
\TextAt{16}{217.029}{471.172}{\Character{14}}%
\TextAt{9}{225.663}{471.172}{\Character{91}}%
\TextAt{13}{228.915}{471.172}{u}%
\TextAt{9}{235.577}{471.172}{\Character{93}}%
\TextAt{21}{238.829}{469.379}{F}%
\TextAt{8}{257.818}{471.172}{and}%
\TextAt{13}{288.638}{471.172}{\Character{18}}%
\TextAt{21}{294.418}{476.109}{\Character{39}}%
\TextAt{16}{303.092}{471.172}{\Character{14}}%
\Vector{q 0 G 0 g 0.478 w 239.726 -296.006 m 245.151 -296.006 l S Q}%
\TextAt{13}{311.726}{471.172}{e}%
\TextAt{9}{320.472}{471.172}{\Character{61}}%
\TextAt{13}{332.898}{471.172}{e}%
\TextAt{16}{340.98}{471.172}{\Character{14}}%
\TextAt{13}{349.614}{471.172}{\Character{18}}%
\TextAt{9}{379.305}{471.172}{\Character{40}}%
\TextAt{13}{383.857}{471.172}{\Character{39}}%
\TextAt{9}{394.846}{471.172}{\Character{61}}%
\TextAt{13}{407.272}{471.172}{\Character{39}}%
\TextAt{21}{414.94}{469.379}{K}%
\TextAt{9}{423.143}{471.172}{\Character{41}}%
\TextAt{13}{427.695}{471.172}{\Character{58}}%
\TextAt{8}{79.226}{442.171}{Let}%
\TextAt{13}{99.693}{442.171}{\Character{21}}%
\TextAt{21}{106.522}{440.378}{n}%
\NativeMathAt{115.479}{442.171}{11.955}{\in}%
\TextAt{13}{126.77}{442.171}{F}%
\TextAt{9}{135.974}{442.171}{\Character{40}}%
\TextAt{13}{140.526}{442.171}{n}%
\TextAt{9}{147.513}{442.171}{\Character{41}}%
\TextAt{16}{153.695}{442.171}{\Character{110}}%
\TextAt{13}{161.303}{442.171}{F}%
\TextAt{9}{170.506}{442.171}{\Character{40}}%
\TextAt{13}{175.059}{442.171}{n}%
\TextAt{16}{183.676}{442.171}{\Character{0}}%
\TextAt{9}{194.604}{442.171}{1\Character{41}}%
\TextAt{8}{205.009}{442.171}{.}%
\TextAt{13}{213.29}{442.171}{\Character{21}}%
\TextAt{21}{220.119}{440.378}{n}%
\TextAt{8}{229.15}{442.171}{is}%
\TextAt{8}{240.412}{442.171}{a}%
\TextAt{8}{249.665}{442.171}{prime}%
\TextAt{8}{282.327}{442.171}{element}%
\TextAt{8}{325.066}{442.171}{of}%
\TextAt{13}{337.894}{442.171}{L}%
\TextAt{21}{345.859}{440.378}{n}%
\TextAt{8}{351.495}{442.171}{,}%
\TextAt{8}{358.242}{442.171}{and}%
\TextAt{13}{380.503}{442.171}{\Character{25}}%
\TextAt{17}{387.143}{440.378}{\Character{6}}%
\TextAt{9}{397.078}{442.171}{\Character{61}}%
\TextAt{13}{409.504}{442.171}{\Character{18}}%
\TextAt{9}{415.284}{442.171}{\Character{40}}%
\TextAt{13}{419.836}{442.171}{\Character{21}}%
\TextAt{21}{426.665}{440.378}{n}%
\TextAt{9}{432.301}{442.171}{\Character{41}}%
\TextAt{8}{440.249}{442.171}{is}%
\TextAt{8}{451.51}{442.171}{a}%
\TextAt{8}{460.764}{442.171}{prime}%
\TextAt{8}{493.425}{442.171}{element}%
\TextAt{8}{79.226}{425.125}{of}%
\TextAt{9}{92.556}{425.125}{\Character{6}}%
\TextAt{8}{104.908}{425.125}{because}%
\TextAt{13}{192.272}{408.079}{\Character{25}}%
\TextAt{30}{199.552}{413.016}{\Character{101}}%
\TextAt{21}{199.342}{413.016}{\Character{27}}%
\TextAt{17}{198.913}{405.124}{\Character{6}}%
\TextAt{9}{208.848}{408.079}{\Character{61}}%
\TextAt{13}{221.273}{408.079}{\Character{18}}%
\TextAt{21}{227.054}{413.016}{\Character{39}}%
\TextAt{9}{233.07}{408.079}{\Character{40}}%
\TextAt{13}{237.622}{408.079}{\Character{21}}%
\TextAt{21}{244.451}{413.016}{\Character{27}}%
\TextAt{21}{244.451}{405.124}{n}%
\TextAt{9}{250.087}{408.079}{\Character{41}}%
\TextAt{9}{257.942}{408.079}{\Character{61}}%
\TextAt{13}{270.386}{408.079}{\Character{18}}%
\TextAt{21}{276.166}{413.016}{\Character{39}}%
\TextAt{9}{282.182}{408.079}{\Character{40}\Character{91}}%
\TextAt{13}{289.986}{408.079}{u}%
\TextAt{3}{296.649}{413.016}{\Character{0}}%
\TextAt{17}{303.235}{413.016}{1}%
\TextAt{9}{307.968}{408.079}{\Character{93}}%
\TextAt{21}{311.219}{406.286}{F}%
\TextAt{9}{318.253}{408.079}{\Character{40}}%
\TextAt{13}{322.805}{408.079}{\Character{21}}%
\TextAt{21}{329.634}{406.286}{n}%
\TextAt{9}{335.27}{408.079}{\Character{41}\Character{41}}%
\TextAt{9}{347.668}{408.079}{\Character{61}}%
\TextAt{13}{360.121}{408.079}{\Character{18}}%
\TextAt{9}{365.901}{408.079}{\Character{40}}%
\TextAt{13}{370.454}{408.079}{\Character{21}}%
\TextAt{21}{377.282}{406.286}{n}%
\TextAt{9}{382.918}{408.079}{\Character{41}}%
\TextAt{9}{390.773}{408.079}{\Character{61}}%
\TextAt{13}{403.217}{408.079}{\Character{25}}%
\TextAt{17}{409.857}{406.286}{\Character{6}}%
\TextAt{13}{416.472}{408.079}{\Character{58}}%
\TextAt{8}{79.226}{384.558}{Since}%
\TextAt{13}{109.921}{384.558}{e}%
\TextAt{21}{115.346}{388.896}{i}%
\TextAt{9}{118.728}{384.558}{\Character{40}}%
\TextAt{13}{123.28}{384.558}{\Character{18}}%
\TextAt{9}{129.06}{384.558}{\Character{40}}%
\TextAt{13}{133.613}{384.558}{\Character{21}}%
\TextAt{21}{140.441}{382.764}{n}%
\TextAt{9}{146.077}{384.558}{\Character{41}\Character{41}}%
\TextAt{9}{158.714}{384.558}{\Character{61}}%
\TextAt{13}{171.396}{384.558}{\Character{18}}%
\TextAt{21}{177.177}{388.896}{\Character{39}}%
\TextAt{22}{182.695}{391.709}{i}%
\TextAt{9}{186.355}{384.558}{\Character{40}}%
\Vector{q 0 G 0 g 0.478 w 118.907 -382.621 m 124.333 -382.621 l S Q}%
\TextAt{13}{190.907}{384.558}{e}%
\TextAt{21}{196.333}{388.941}{i}%
\TextAt{9}{199.714}{384.558}{\Character{40}}%
\TextAt{13}{204.266}{384.558}{\Character{21}}%
\TextAt{21}{211.095}{382.764}{n}%
\TextAt{9}{216.731}{384.558}{\Character{41}\Character{41}}%
\TextAt{9}{229.368}{384.558}{\Character{61}}%
\TextAt{9}{242.016}{384.558}{0}%
\TextAt{8}{251.938}{384.558}{for}%
\TextAt{13}{269.96}{384.558}{i}%
\TextAt{9}{277.508}{384.558}{\Character{61}}%
\TextAt{13}{290.168}{384.558}{n}%
\TextAt{8}{297.155}{384.558}{,}%
\TextAt{8}{304.476}{384.558}{and}%
\NativeMathAt{327.376}{384.558}{11.955}{\neq}%
\TextAt{9}{340.025}{384.558}{0}%
\TextAt{8}{349.923}{384.558}{for}%
\TextAt{13}{367.946}{384.558}{i}%
\TextAt{9}{375.494}{384.558}{\Character{61}}%
\TextAt{13}{388.153}{384.558}{n}%
\TextAt{16}{397.891}{384.558}{\Character{0}}%
\TextAt{9}{409.94}{384.558}{1}%
\TextAt{8}{415.793}{384.558}{,}%
\TextAt{8}{423.114}{384.558}{we}%
\TextAt{8}{440.479}{384.558}{have}%
\TextAt{13}{467.611}{384.558}{\Character{25}}%
\TextAt{17}{474.251}{382.764}{\Character{6}}%
\NativeMathAt{484.42}{384.558}{11.955}{\in}%
\TextAt{13}{495.945}{384.558}{F}%
\TextAt{21}{503.523}{382.764}{e}%
\TextAt{9}{507.949}{384.558}{\Character{40}}%
\TextAt{13}{512.502}{384.558}{n}%
\TextAt{9}{519.489}{384.558}{\Character{41}}%
\TextAt{16}{526.792}{384.558}{\Character{110}}%
\TextAt{13}{79.226}{367.512}{F}%
\TextAt{21}{86.804}{365.718}{e}%
\TextAt{9}{91.23}{367.512}{\Character{40}}%
\TextAt{13}{95.783}{367.512}{n}%
\TextAt{16}{105.736}{367.512}{\Character{0}}%
\TextAt{9}{118}{367.512}{1\Character{41}}%
\TextAt{8}{128.406}{367.512}{.}%
\TextAt{8}{138.215}{367.512}{Hence}%
\TextAt{9}{173.446}{367.512}{\Character{6}}%
\TextAt{9}{185.987}{367.512}{\Character{61}}%
\TextAt{13}{199.192}{367.512}{K}%
\TextAt{9}{210.01}{367.512}{\Character{40}}%
\TextAt{13}{214.562}{367.512}{\Character{25}}%
\TextAt{17}{221.202}{365.718}{\Character{6}}%
\TextAt{9}{227.817}{367.512}{\Character{41}}%
\TextAt{8}{236.72}{367.512}{is}%
\TextAt{8}{248.937}{367.512}{the}%
\TextAt{8}{269.547}{367.512}{field}%
\TextAt{8}{295.357}{367.512}{of}%
\TextAt{13}{309.141}{367.512}{\Character{25}}%
\TextAt{21}{316.21}{371.85}{n}%
\TextAt{8}{321.847}{367.512}{-division}%
\TextAt{8}{369.511}{367.512}{points}%
\TextAt{8}{405.149}{367.512}{of}%
\TextAt{13}{418.932}{367.512}{F}%
\TextAt{21}{426.51}{365.718}{e}%
\TextAt{8}{430.936}{367.512}{,}%
\TextAt{8}{438.653}{367.512}{and}%
\TextAt{13}{461.869}{367.512}{N}%
\TextAt{17}{471.243}{365.652}{\Character{6}}%
\TextAt{3}{477.359}{365.652}{\Character{106}}%
\TextAt{21}{479.711}{365.652}{K}%
\TextAt{9}{487.914}{367.512}{\Character{40}}%
\TextAt{16}{492.466}{367.512}{\Character{0}}%
\TextAt{13}{501.764}{367.512}{\Character{25}}%
\TextAt{17}{508.405}{365.718}{\Character{6}}%
\TextAt{9}{515.019}{367.512}{\Character{41}}%
\TextAt{9}{523.651}{367.512}{\Character{61}}%
\TextAt{13}{79.226}{350.466}{\Character{25}}%
\TextAt{9}{89.616}{350.466}{\Character{61}}%
\TextAt{13}{102.042}{350.466}{u}%
\Vector{q 0 G 0 g 0.478 w 36.704 -416.713 m 43.773 -416.713 l S Q}%
\TextAt{13}{108.704}{350.466}{\Character{25}}%
\TextAt{8}{119.671}{350.466}{by}%
\TextAt{8}{135.93}{350.466}{(5.4).}%
\TextAt{8}{168.456}{350.466}{Since}%
\Vector{q 0 G 0 g 0.478 w 127.013 -416.713 m 134.083 -416.713 l S Q}%
\TextAt{13}{199.013}{350.466}{\Character{25}}%
\TextAt{9}{209.403}{350.466}{\Character{61}}%
\TextAt{13}{221.829}{350.466}{N}%
\TextAt{21}{231.203}{348.606}{L}%
\TextAt{22}{236.962}{347.61}{n}%
\TextAt{3}{242.061}{348.606}{\Character{106}}%
\TextAt{21}{244.413}{348.606}{K}%
\TextAt{9}{252.616}{350.466}{\Character{40}}%
\TextAt{16}{257.168}{350.466}{\Character{0}}%
\TextAt{13}{266.466}{350.466}{\Character{21}}%
\TextAt{21}{273.295}{348.673}{n}%
\TextAt{9}{278.931}{350.466}{\Character{41}}%
\NativeMathAt{286.804}{350.466}{11.955}{\in}%
\TextAt{13}{298.095}{350.466}{N}%
\TextAt{21}{307.469}{348.606}{L}%
\TextAt{22}{313.228}{347.61}{n}%
\TextAt{3}{318.327}{348.606}{\Character{106}}%
\TextAt{21}{320.68}{348.606}{K}%
\TextAt{13}{328.882}{350.466}{L}%
\TextAt{3}{336.847}{354.804}{\Character{3}}%
\TextAt{21}{336.847}{347.51}{n}%
\TextAt{8}{342.483}{350.466}{,}%
\TextAt{8}{349.632}{350.466}{we}%
\TextAt{8}{366.86}{350.466}{get}%
\TextAt{13}{186.928}{321.465}{r}%
\TextAt{21}{192.203}{319.605}{L}%
\TextAt{22}{197.962}{318.609}{n}%
\TextAt{3}{203.061}{319.605}{\Character{106}}%
\TextAt{21}{205.413}{319.605}{K}%
\TextAt{9}{213.616}{321.465}{\Character{40}}%
\TextAt{13}{218.168}{321.465}{\Character{27}}%
\TextAt{9}{225.251}{321.465}{\Character{41}}%
\TextAt{9}{233.106}{321.465}{\Character{61}}%
\TextAt{13}{245.549}{321.465}{N}%
\TextAt{17}{254.923}{319.605}{\Character{6}}%
\TextAt{3}{261.039}{319.605}{\Character{106}}%
\TextAt{21}{263.391}{319.605}{K}%
\TextAt{9}{271.594}{321.465}{\Character{40}}%
\TextAt{16}{276.146}{321.465}{\Character{0}}%
\TextAt{13}{285.445}{321.465}{\Character{25}}%
\TextAt{17}{292.085}{319.671}{\Character{6}}%
\TextAt{9}{298.699}{321.465}{\Character{41}}%
\TextAt{9}{306.554}{321.465}{\Character{61}}%
\TextAt{13}{318.998}{321.465}{\Character{25}}%
\NativeMathAt{329.388}{321.465}{11.955}{\equiv}%
\TextAt{13}{342.008}{321.465}{u}%
\TextAt{8}{351.991}{321.465}{mod}%
\TextAt{13}{377.429}{321.465}{N}%
\TextAt{21}{386.803}{319.605}{L}%
\TextAt{22}{392.562}{318.609}{n}%
\TextAt{3}{397.66}{319.605}{\Character{106}}%
\TextAt{21}{400.013}{319.605}{K}%
\TextAt{13}{408.215}{321.465}{L}%
\TextAt{3}{416.18}{326.401}{\Character{3}}%
\TextAt{21}{416.18}{318.509}{n}%
\TextAt{13}{421.816}{321.465}{\Character{59}}%
\TextAt{8}{79.226}{292.463}{and}%
\TextAt{8}{101.989}{292.463}{thus}%
\TextAt{9}{155.433}{275.418}{\Character{40}}%
\TextAt{13}{159.985}{275.418}{a\Character{59}}%
\TextAt{13}{171.354}{275.418}{L}%
\TextAt{21}{179.339}{273.624}{n}%
\TextAt{16}{188.296}{275.418}{\Character{106}}%
\TextAt{13}{194.937}{275.418}{K}%
\TextAt{9}{205.755}{275.418}{\Character{41}}%
\TextAt{9}{213.61}{275.418}{\Character{61}}%
\TextAt{9}{226.031}{275.418}{\Character{40}}%
\Vector{q 0 G 0 g 0.478 w 158.606 -491.761 m 165.675 -491.761 l S Q}%
\TextAt{13}{230.606}{275.418}{\Character{25}}%
\TextAt{21}{237.675}{280.354}{v}%
\TextAt{22}{241.799}{279.014}{K}%
\TextAt{17}{248.842}{280.354}{\Character{40}}%
\TextAt{21}{252.135}{280.354}{a}%
\TextAt{17}{256.633}{280.354}{\Character{41}}%
\TextAt{13}{260.424}{275.418}{\Character{59}}%
\TextAt{13}{265.648}{275.418}{L}%
\TextAt{21}{273.633}{273.624}{n}%
\TextAt{16}{282.59}{275.418}{\Character{106}}%
\TextAt{13}{289.232}{275.418}{K}%
\TextAt{9}{300.05}{275.418}{\Character{41}\Character{40}}%
\TextAt{13}{309.154}{275.418}{u\Character{59}}%
\TextAt{13}{321.038}{275.418}{L}%
\TextAt{21}{329.025}{273.624}{n}%
\TextAt{16}{337.983}{275.418}{\Character{106}}%
\TextAt{13}{344.624}{275.418}{K}%
\TextAt{9}{355.442}{275.418}{\Character{41}}%
\TextAt{9}{363.297}{275.418}{\Character{61}}%
\TextAt{9}{375.718}{275.418}{\Character{40}}%
\TextAt{13}{380.293}{275.418}{u\Character{59}}%
\TextAt{13}{392.177}{275.418}{L}%
\TextAt{21}{400.164}{273.624}{n}%
\TextAt{16}{409.121}{275.418}{\Character{106}}%
\TextAt{13}{415.763}{275.418}{K}%
\TextAt{9}{426.581}{275.418}{\Character{41}}%
\TextAt{9}{434.436}{275.418}{\Character{61}}%
\TextAt{13}{446.879}{275.418}{\Character{27}}%
\TextAt{13}{453.311}{275.418}{\Character{58}}%
\Vector{q 0 G 0 g 0.399 w 452.334 -522.104 m 452.334 -514.034 l S Q}%
\Vector{q 0 G 0 g 0.399 w 452.534 -514.234 m 459.707 -514.234 l S Q}%
\Vector{q 0 G 0 g 0.399 w 452.534 -521.905 m 459.707 -521.905 l S Q}%
\Vector{q 0 G 0 g 0.399 w 459.906 -522.104 m 459.906 -514.034 l S Q}%
\TextAt{12}{79.226}{196.882}{Corollary}%
\TextAt{12}{139.368}{196.882}{(5.6).}%
\TextAt{8}{174.992}{196.882}{The}%
\TextAt{8}{199.354}{196.882}{field}%
\TextAt{13}{225.019}{196.882}{L}%
\TextAt{21}{232.983}{195.089}{n}%
\TextAt{16}{242.465}{196.882}{\Character{106}}%
\TextAt{13}{249.632}{196.882}{K}%
\TextAt{8}{264.655}{196.882}{of}%
\TextAt{13}{278.293}{196.882}{\Character{25}}%
\TextAt{21}{285.362}{201.221}{n}%
\TextAt{8}{290.999}{196.882}{-division}%
\TextAt{8}{338.518}{196.882}{points}%
\TextAt{8}{374.01}{196.882}{is}%
\TextAt{8}{386.082}{196.882}{the}%
\TextAt{8}{406.546}{196.882}{class}%
\TextAt{8}{434.291}{196.882}{field}%
\TextAt{8}{459.956}{196.882}{relative}%
\TextAt{8}{501.892}{196.882}{to}%
\TextAt{8}{516.51}{196.882}{the}%
\TextAt{8}{79.226}{179.836}{group}%
\TextAt{9}{112.402}{179.836}{\Character{40}}%
\TextAt{13}{116.954}{179.836}{\Character{25}}%
\TextAt{9}{124.023}{179.836}{\Character{41}}%
\TextAt{16}{131.232}{179.836}{\Character{2}}%
\TextAt{13}{143.187}{179.836}{U}%
\TextAt{17}{152.387}{185.947}{\Character{40}}%
\TextAt{21}{155.68}{185.947}{n}%
\TextAt{17}{160.818}{185.947}{\Character{41}}%
\TextAt{21}{151.138}{176.595}{K}%
\NativeMathAt{167.93}{179.836}{11.955}{\subseteq}%
\TextAt{13}{180.55}{179.836}{K}%
\TextAt{3}{191.367}{184.175}{\Character{3}}%
\TextAt{8}{196.1}{179.836}{.}%
\TextAt{8}{297.211}{43.033}{320}%
\Destination{78.23}{809.815}{nav.page.320}%
\Destination{79.36}{716.055}{nav.print.350}%
\Destination{79.36}{715.055}{nav.number.5.5.5}%
\Destination{79.36}{210.855}{nav.number.5.5.6}%
\Hotlink{102.814}{741.855}{25.078}{11.047}{nav.number.5.2.4}%
\Hotlink{286.775}{623.255}{25.078}{11.047}{nav.number.5.2.4}%
\Hotlink{135.43}{348.255}{25.078}{11.047}{nav.number.5.5.4}%
\Destination{79.226}{605.392}{errata.121}%
\Destination{79.226}{534.233}{errata.122}%
}

% Page 333 (body page 321)
\Page{%
\TextAt{2}{79.226}{796.943}{V.}%
\TextAt{2}{94.073}{796.943}{Local}%
\TextAt{2}{123.615}{796.943}{Class}%
\TextAt{2}{152.208}{796.943}{Field}%
\TextAt{2}{179.939}{796.943}{Theory}%
\Vector{q 0 G 0 g 0.299 w 7.226 18.729 m 460.769 18.729 l S Q}%
\TextAt{19}{79.226}{761.124}{Proof.}%
\TextAt{8}{116}{761.124}{For}%
\TextAt{13}{136.97}{761.124}{a}%
\TextAt{9}{146.435}{761.124}{\Character{61}}%
\TextAt{13}{158.861}{761.124}{u\Character{25}}%
\TextAt{21}{172.593}{765.462}{v}%
\TextAt{22}{176.716}{764.123}{K}%
\TextAt{17}{183.759}{765.462}{\Character{40}}%
\TextAt{21}{187.052}{765.462}{a}%
\TextAt{17}{191.55}{765.462}{\Character{41}}%
\TextAt{8}{199.239}{761.124}{we}%
\TextAt{8}{216.466}{761.124}{have}%
\TextAt{13}{177.715}{732.123}{a}%
\NativeMathAt{187.181}{732.123}{11.955}{\in}%
\TextAt{13}{198.472}{732.123}{N}%
\TextAt{21}{207.846}{730.263}{L}%
\TextAt{22}{213.605}{729.267}{n}%
\TextAt{3}{218.704}{730.263}{\Character{106}}%
\TextAt{21}{221.056}{730.263}{K}%
\TextAt{13}{229.259}{732.123}{L}%
\TextAt{3}{237.223}{737.059}{\Character{3}}%
\TextAt{21}{237.223}{729.167}{n}%
\NativeArrowAt{246.18}{732.123}{11.955}{21.926}{\Longleftrightarrow}%
\TextAt{9}{271.419}{732.123}{\Character{40}}%
\TextAt{13}{275.971}{732.123}{a\Character{59}}%
\TextAt{13}{287.34}{732.123}{L}%
\TextAt{21}{295.325}{730.329}{n}%
\TextAt{16}{304.282}{732.123}{\Character{106}}%
\TextAt{13}{310.924}{732.123}{K}%
\TextAt{9}{321.742}{732.123}{\Character{41}}%
\TextAt{9}{329.597}{732.123}{\Character{61}}%
\TextAt{9}{342.018}{732.123}{1}%
\NativeArrowAt{246.18}{712.088}{11.955}{21.926}{\Longleftrightarrow}%
\TextAt{9}{271.419}{712.088}{\Character{91}}%
\TextAt{13}{274.671}{712.088}{u}%
\TextAt{3}{281.333}{717.024}{\Character{0}}%
\TextAt{17}{287.92}{717.024}{1}%
\TextAt{9}{292.652}{712.088}{\Character{93}}%
\TextAt{21}{295.904}{710.295}{F}%
\TextAt{9}{302.937}{712.088}{\Character{40}}%
\TextAt{13}{307.489}{712.088}{\Character{21}}%
\TextAt{9}{314.318}{712.088}{\Character{41}}%
\TextAt{9}{322.173}{712.088}{\Character{61}}%
\TextAt{13}{334.617}{712.088}{\Character{21}}%
\TextAt{8}{353.4}{712.088}{for}%
\TextAt{8}{371.285}{712.088}{all}%
\TextAt{13}{387.544}{712.088}{\Character{21}}%
\NativeMathAt{397.694}{712.088}{11.955}{\in}%
\TextAt{13}{408.985}{712.088}{F}%
\TextAt{9}{418.188}{712.088}{\Character{40}}%
\TextAt{13}{422.74}{712.088}{n}%
\TextAt{9}{429.728}{712.088}{\Character{41}}%
\NativeArrowAt{246.18}{690.901}{11.955}{21.926}{\Longleftrightarrow}%
\TextAt{9}{271.419}{690.901}{\Character{91}}%
\TextAt{13}{274.671}{690.901}{u}%
\TextAt{3}{281.333}{695.837}{\Character{0}}%
\TextAt{17}{287.92}{695.837}{1}%
\TextAt{9}{292.652}{690.901}{\Character{93}}%
\TextAt{21}{295.904}{689.108}{F}%
\TextAt{9}{306.258}{690.901}{\Character{61}}%
\TextAt{8}{318.683}{690.901}{id}%
\TextAt{21}{328.439}{689.041}{F}%
\TextAt{17}{334.974}{689.041}{\Character{40}}%
\TextAt{21}{338.267}{689.041}{n}%
\TextAt{17}{343.406}{689.041}{\Character{41}}%
\NativeArrowAt{350.518}{690.901}{11.955}{21.926}{\Longleftrightarrow}%
\TextAt{13}{375.757}{690.901}{u}%
\TextAt{3}{382.419}{695.837}{\Character{0}}%
\TextAt{17}{389.006}{695.837}{1}%
\NativeMathAt{397.059}{690.901}{11.955}{\in}%
\TextAt{13}{408.35}{690.901}{U}%
\TextAt{17}{417.549}{697.011}{\Character{40}}%
\TextAt{21}{420.842}{697.011}{n}%
\TextAt{17}{425.98}{697.011}{\Character{41}}%
\TextAt{21}{416.3}{687.66}{K}%
\NativeArrowAt{246.18}{669.714}{11.955}{21.926}{\Longleftrightarrow}%
\TextAt{13}{271.419}{669.714}{a}%
\NativeMathAt{280.885}{669.714}{11.955}{\in}%
\TextAt{9}{292.176}{669.714}{\Character{40}}%
\TextAt{13}{296.728}{669.714}{\Character{25}}%
\TextAt{9}{303.797}{669.714}{\Character{41}}%
\TextAt{16}{311.006}{669.714}{\Character{2}}%
\TextAt{13}{322.962}{669.714}{U}%
\TextAt{17}{332.161}{675.824}{\Character{40}}%
\TextAt{21}{335.454}{675.824}{n}%
\TextAt{17}{340.592}{675.824}{\Character{41}}%
\TextAt{21}{330.912}{666.473}{K}%
\TextAt{13}{344.384}{669.714}{\Character{58}}%
\Vector{q 0 G 0 g 0.399 w 452.334 -133.287 m 452.334 -125.217 l S Q}%
\Vector{q 0 G 0 g 0.399 w 452.534 -125.417 m 459.707 -125.417 l S Q}%
\Vector{q 0 G 0 g 0.399 w 452.534 -133.088 m 459.707 -133.088 l S Q}%
\Vector{q 0 G 0 g 0.399 w 459.906 -133.287 m 459.906 -125.217 l S Q}%
\TextAt{8}{97.159}{611.712}{For}%
\TextAt{8}{118.235}{611.712}{the}%
\TextAt{8}{138.498}{611.712}{maximal}%
\TextAt{8}{186.414}{611.712}{abelian}%
\TextAt{8}{227.169}{611.712}{extension}%
\TextAt{13}{279.041}{611.712}{K}%
\TextAt{21}{289.859}{616.05}{ab}%
\TextAt{16}{301.981}{611.712}{\Character{106}}%
\TextAt{13}{308.805}{611.712}{K}%
\TextAt{8}{319.622}{611.712}{,}%
\TextAt{8}{326.905}{611.712}{this}%
\TextAt{8}{349.834}{611.712}{gives}%
\TextAt{8}{378.622}{611.712}{the}%
\TextAt{8}{398.885}{611.712}{following}%
\TextAt{8}{448.426}{611.712}{generalization}%
\TextAt{8}{523.337}{611.712}{of}%
\TextAt{8}{79.226}{594.666}{the}%
\TextAt{8}{99.383}{594.666}{local}%
\TextAt{8}{127.023}{594.666}{Kronecker–Weber}%
\TextAt{8}{221.313}{594.666}{theorem}%
\TextAt{8}{266.839}{594.666}{(1.9):}%
\TextAt{12}{79.226}{568.653}{Corollary}%
\TextAt{12}{139.014}{568.653}{(5.7).}%
\TextAt{8}{174.33}{568.653}{The}%
\TextAt{8}{198.383}{568.653}{maximal}%
\TextAt{8}{246.192}{568.653}{abelian}%
\TextAt{8}{286.84}{568.653}{extension}%
\TextAt{8}{338.605}{568.653}{of}%
\TextAt{13}{351.936}{568.653}{K}%
\TextAt{8}{366.651}{568.653}{is}%
\TextAt{8}{378.414}{568.653}{the}%
\TextAt{8}{398.571}{568.653}{composite}%
\TextAt{13}{274.566}{539.652}{K}%
\TextAt{21}{285.384}{544.588}{ab}%
\TextAt{9}{297.323}{539.652}{\Character{61}}%
\TextAt{20}{312.487}{542.674}{\Character{101}}%
\TextAt{13}{309.749}{539.652}{K}%
\TextAt{13}{320.557}{539.652}{L}%
\TextAt{21}{328.531}{537.859}{\Character{25}}%
\TextAt{13}{334.178}{539.652}{\Character{59}}%
\TextAt{8}{79.226}{510.651}{where}%
\TextAt{13}{113.035}{510.651}{L}%
\TextAt{21}{121}{508.858}{\Character{25}}%
\TextAt{8}{130.544}{510.651}{is}%
\TextAt{8}{142.308}{510.651}{the}%
\TextAt{8}{162.465}{510.651}{union}%
\TextAt{20}{194.983}{519.618}{\Character{83}}%
\NativeMathAt{204.945}{516.463}{7.97}{\infty}%
\TextAt{21}{204.945}{507.164}{n}%
\TextAt{17}{210.083}{507.164}{\Character{61}1}%
\TextAt{13}{223.395}{510.651}{L}%
\TextAt{21}{231.359}{508.858}{n}%
\TextAt{8}{240.893}{510.651}{of}%
\TextAt{8}{254.223}{510.651}{the}%
\TextAt{8}{274.379}{510.651}{fields}%
\TextAt{13}{304.351}{510.651}{L}%
\TextAt{21}{312.316}{508.858}{n}%
\TextAt{8}{321.849}{510.651}{of}%
\TextAt{13}{335.179}{510.651}{\Character{25}}%
\TextAt{21}{342.249}{514.99}{n}%
\TextAt{8}{347.885}{510.651}{-division}%
\TextAt{8}{395.096}{510.651}{points.}%
\TextAt{19}{79.226}{475.672}{Proof.}%
\TextAt{8}{116}{475.672}{Let}%
\TextAt{13}{137.392}{475.672}{L}%
\TextAt{16}{149.397}{475.672}{\Character{106}}%
\TextAt{13}{156.758}{475.672}{K}%
\TextAt{8}{171.896}{475.672}{be}%
\TextAt{8}{188.242}{475.672}{a}%
\TextAt{8}{198.42}{475.672}{finite}%
\TextAt{8}{228.754}{475.672}{abelian}%
\TextAt{8}{269.824}{475.672}{extension.}%
\TextAt{8}{327.407}{475.672}{Then}%
\TextAt{8}{358.387}{475.672}{we}%
\TextAt{8}{376.037}{475.672}{have}%
\TextAt{13}{403.454}{475.672}{\Character{25}}%
\TextAt{21}{410.523}{480.011}{f}%
\NativeMathAt{420.008}{475.672}{11.955}{\in}%
\TextAt{13}{432.019}{475.672}{N}%
\TextAt{21}{441.393}{473.813}{L}%
\TextAt{3}{447.152}{473.813}{\Character{106}}%
\TextAt{21}{449.504}{473.813}{K}%
\TextAt{13}{457.707}{475.672}{L}%
\TextAt{3}{465.671}{480.011}{\Character{3}}%
\TextAt{8}{474.723}{475.672}{for}%
\TextAt{8}{493.031}{475.672}{suitable}%
\TextAt{13}{79.226}{458.626}{f}%
\TextAt{8}{86.273}{458.626}{.}%
\TextAt{8}{95.56}{458.626}{Since}%
\TextAt{13}{126.396}{458.626}{N}%
\TextAt{21}{135.77}{456.767}{L}%
\TextAt{3}{141.529}{456.767}{\Character{106}}%
\TextAt{21}{143.882}{456.767}{K}%
\TextAt{13}{152.084}{458.626}{L}%
\TextAt{3}{160.049}{462.965}{\Character{3}}%
\TextAt{8}{168.958}{458.626}{is}%
\TextAt{8}{181.001}{458.626}{open}%
\TextAt{8}{209.567}{458.626}{in}%
\TextAt{13}{223.499}{458.626}{K}%
\TextAt{3}{234.317}{462.965}{\Character{3}}%
\TextAt{8}{239.049}{458.626}{,}%
\TextAt{8}{246.548}{458.626}{and}%
\TextAt{8}{269.59}{458.626}{since}%
\TextAt{8}{298.538}{458.626}{the}%
\TextAt{13}{318.974}{458.626}{U}%
\TextAt{17}{328.173}{464.737}{\Character{40}}%
\TextAt{21}{331.466}{464.737}{n}%
\TextAt{17}{336.604}{464.737}{\Character{41}}%
\TextAt{21}{326.924}{455.385}{K}%
\TextAt{8}{344.573}{458.626}{form}%
\TextAt{8}{372.493}{458.626}{a}%
\TextAt{8}{382.528}{458.626}{basis}%
\TextAt{8}{411.547}{458.626}{of}%
\TextAt{8}{425.157}{458.626}{neighbourhoods}%
\TextAt{8}{510.055}{458.626}{of}%
\TextAt{9}{523.665}{458.626}{1}%
\TextAt{8}{529.518}{458.626}{,}%
\TextAt{8}{79.226}{441.58}{we}%
\TextAt{8}{96.415}{441.58}{have}%
\TextAt{9}{123.371}{441.58}{\Character{40}}%
\TextAt{13}{127.923}{441.58}{\Character{25}}%
\TextAt{21}{134.992}{445.919}{f}%
\TextAt{9}{140.437}{441.58}{\Character{41}}%
\TextAt{16}{147.567}{441.58}{\Character{2}}%
\TextAt{13}{159.443}{441.58}{U}%
\TextAt{17}{168.642}{447.691}{\Character{40}}%
\TextAt{21}{171.935}{447.691}{n}%
\TextAt{17}{177.073}{447.691}{\Character{41}}%
\TextAt{21}{167.393}{438.339}{K}%
\NativeMathAt{184.185}{441.58}{11.955}{\subseteq}%
\TextAt{13}{196.805}{441.58}{N}%
\TextAt{21}{206.179}{439.721}{L}%
\TextAt{3}{211.938}{439.721}{\Character{106}}%
\TextAt{21}{214.29}{439.721}{K}%
\TextAt{13}{222.493}{441.58}{L}%
\TextAt{3}{230.457}{445.919}{\Character{3}}%
\TextAt{8}{239.048}{441.58}{for}%
\TextAt{8}{256.894}{441.58}{a}%
\TextAt{8}{266.611}{441.58}{suitable}%
\TextAt{13}{310.208}{441.58}{n}%
\TextAt{8}{317.196}{441.58}{.}%
\TextAt{8}{325.631}{441.58}{Hence}%
\TextAt{13}{360.37}{441.58}{L}%
\TextAt{8}{372.193}{441.58}{is}%
\TextAt{8}{383.918}{441.58}{contained}%
\TextAt{8}{436.889}{441.58}{in}%
\TextAt{8}{450.503}{441.58}{the}%
\TextAt{8}{470.62}{441.58}{class}%
\TextAt{8}{498.019}{441.58}{field}%
\TextAt{8}{523.337}{441.58}{of}%
\TextAt{8}{79.226}{424.535}{the}%
\TextAt{8}{99.383}{424.535}{group}%
\TextAt{9}{199.631}{407.489}{\Character{40}}%
\TextAt{13}{204.183}{407.489}{\Character{25}}%
\TextAt{21}{211.252}{412.425}{f}%
\TextAt{9}{216.697}{407.489}{\Character{41}}%
\TextAt{16}{223.906}{407.489}{\Character{2}}%
\TextAt{13}{235.861}{407.489}{U}%
\TextAt{17}{245.06}{413.599}{\Character{40}}%
\TextAt{21}{248.354}{413.599}{n}%
\TextAt{17}{253.492}{413.599}{\Character{41}}%
\TextAt{21}{243.812}{404.247}{K}%
\TextAt{9}{260.604}{407.489}{\Character{61}}%
\TextAt{9}{273.014}{407.489}{\Character{40}\Character{40}}%
\TextAt{13}{282.134}{407.489}{\Character{25}}%
\TextAt{9}{289.203}{407.489}{\Character{41}}%
\TextAt{16}{296.412}{407.489}{\Character{2}}%
\TextAt{13}{308.368}{407.489}{U}%
\TextAt{17}{317.567}{413.599}{\Character{40}}%
\TextAt{21}{320.86}{413.599}{n}%
\TextAt{17}{325.998}{413.599}{\Character{41}}%
\TextAt{21}{316.318}{404.247}{K}%
\TextAt{9}{329.789}{407.489}{\Character{41}}%
\TextAt{16}{336.998}{407.489}{\Character{92}}%
\TextAt{9}{347.625}{407.489}{\Character{40}\Character{40}}%
\TextAt{13}{356.73}{407.489}{\Character{25}}%
\TextAt{21}{363.799}{412.425}{f}%
\TextAt{9}{369.244}{407.489}{\Character{41}}%
\TextAt{16}{376.453}{407.489}{\Character{2}}%
\TextAt{13}{388.408}{407.489}{U}%
\TextAt{21}{396.358}{405.695}{K}%
\TextAt{9}{404.561}{407.489}{\Character{41}}%
\TextAt{13}{409.113}{407.489}{\Character{58}}%
\TextAt{8}{79.226}{383.967}{The}%
\TextAt{8}{103.253}{383.967}{class}%
\TextAt{8}{130.663}{383.967}{field}%
\TextAt{8}{155.993}{383.967}{of}%
\TextAt{9}{169.296}{383.967}{\Character{40}}%
\TextAt{13}{173.849}{383.967}{\Character{25}}%
\TextAt{9}{180.918}{383.967}{\Character{41}}%
\TextAt{16}{188.072}{383.967}{\Character{2}}%
\TextAt{13}{199.972}{383.967}{U}%
\TextAt{17}{209.171}{390.077}{\Character{40}}%
\TextAt{21}{212.464}{390.077}{n}%
\TextAt{17}{217.602}{390.077}{\Character{41}}%
\TextAt{21}{207.922}{380.726}{K}%
\TextAt{8}{225.264}{383.967}{is}%
\TextAt{13}{237.001}{383.967}{L}%
\TextAt{21}{244.966}{382.174}{n}%
\TextAt{8}{250.602}{383.967}{,}%
\TextAt{8}{257.73}{383.967}{and}%
\TextAt{8}{280.466}{383.967}{that}%
\TextAt{8}{305.807}{383.967}{of}%
\TextAt{9}{319.111}{383.967}{\Character{40}}%
\TextAt{13}{323.663}{383.967}{\Character{25}}%
\TextAt{21}{330.732}{388.305}{f}%
\TextAt{9}{336.177}{383.967}{\Character{41}}%
\TextAt{16}{343.331}{383.967}{\Character{2}}%
\TextAt{13}{355.231}{383.967}{U}%
\TextAt{21}{363.181}{382.174}{K}%
\TextAt{8}{375.255}{383.967}{is}%
\TextAt{8}{386.991}{383.967}{the}%
\TextAt{8}{407.121}{383.967}{unramified}%
\TextAt{8}{465.627}{383.967}{extension}%
\TextAt{13}{517.366}{383.967}{K}%
\TextAt{21}{527.325}{382.174}{f}%
\TextAt{8}{79.226}{366.921}{of}%
\TextAt{8}{92.556}{366.921}{degree}%
\TextAt{13}{128.972}{366.921}{f}%
\TextAt{8}{136.018}{366.921}{.}%
\TextAt{8}{144.466}{366.921}{It}%
\TextAt{8}{157.139}{366.921}{follows}%
\TextAt{8}{195.575}{366.921}{that}%
\TextAt{13}{220.944}{366.921}{L}%
\NativeMathAt{232.229}{366.921}{11.955}{\subseteq}%
\TextAt{13}{244.848}{366.921}{K}%
\TextAt{21}{254.807}{365.128}{f}%
\TextAt{13}{260.252}{366.921}{L}%
\TextAt{21}{268.217}{365.128}{n}%
\NativeMathAt{277.174}{366.921}{11.955}{\subseteq}%
\TextAt{20}{292.532}{369.943}{\Character{101}}%
\TextAt{13}{289.793}{366.921}{K}%
\TextAt{13}{300.601}{366.921}{L}%
\TextAt{21}{308.575}{365.128}{\Character{25}}%
\TextAt{9}{317.543}{366.921}{\Character{61}}%
\TextAt{13}{329.969}{366.921}{K}%
\TextAt{21}{340.786}{371.259}{ab}%
\TextAt{8}{349.405}{366.921}{.}%
\Vector{q 0 G 0 g 0.399 w 452.334 -407.079 m 452.334 -399.009 l S Q}%
\Vector{q 0 G 0 g 0.399 w 452.534 -399.209 m 459.707 -399.209 l S Q}%
\Vector{q 0 G 0 g 0.399 w 452.534 -406.88 m 459.707 -406.88 l S Q}%
\Vector{q 0 G 0 g 0.399 w 459.906 -407.079 m 459.906 -399.009 l S Q}%
\TextAt{12}{79.226}{331.942}{Exercise}%
\TextAt{12}{131.664}{331.942}{V.5.1.}%
\TextAt{8}{169.747}{331.942}{Let}%
\TextAt{13}{190.053}{331.942}{F}%
\TextAt{9}{202.577}{331.942}{\Character{61}}%
\TextAt{13}{215.003}{331.942}{F}%
\TextAt{21}{222.581}{330.149}{e}%
\TextAt{8}{230.241}{331.942}{be}%
\TextAt{8}{245.502}{331.942}{the}%
\TextAt{8}{264.995}{331.942}{Lubin–Tate}%
\TextAt{8}{327.252}{331.942}{module}%
\TextAt{8}{367.881}{331.942}{for}%
\TextAt{8}{385.103}{331.942}{the}%
\TextAt{8}{404.596}{331.942}{Lubin–Tate}%
\TextAt{8}{466.853}{331.942}{series}%
\TextAt{13}{497.524}{331.942}{e}%
\NativeMathAt{506.27}{331.942}{11.955}{\in}%
\TextAt{16}{517.544}{331.942}{\Character{69}}%
\TextAt{21}{523.871}{330.149}{\Character{25}}%
\TextAt{8}{529.518}{331.942}{,}%
\TextAt{8}{79.226}{314.896}{with}%
\TextAt{8}{106.79}{314.896}{the}%
\TextAt{8}{127.85}{314.896}{endomorphisms}%
\TextAt{9}{212.129}{314.896}{\Character{91}}%
\TextAt{13}{215.381}{314.896}{a}%
\TextAt{9}{221.526}{314.896}{\Character{93}}%
\TextAt{9}{229.62}{314.896}{\Character{61}}%
\TextAt{9}{243.583}{314.896}{\Character{91}}%
\TextAt{13}{246.855}{314.896}{a}%
\TextAt{9}{253}{314.896}{\Character{93}}%
\TextAt{21}{256.252}{313.103}{e}%
\TextAt{8}{260.678}{314.896}{.}%
\TextAt{8}{271.838}{314.896}{Let}%
\TextAt{13}{293.711}{314.896}{S}%
\TextAt{9}{306.467}{314.896}{\Character{61}}%
\TextAt{14}{320.432}{314.896}{o}%
\TextAt{21}{326.276}{313.103}{K}%
\TextAt{9}{334.479}{314.896}{\Character{91}\Character{91}}%
\TextAt{13}{340.982}{314.896}{X}%
\TextAt{9}{351.638}{314.896}{\Character{93}\Character{93}}%
\TextAt{8}{362.942}{314.896}{and}%
\TextAt{13}{386.608}{314.896}{S}%
\TextAt{3}{394.504}{319.235}{\Character{3}}%
\TextAt{9}{404.097}{314.896}{\Character{61}}%
\TextAt{16}{418.062}{314.896}{\Character{102}}%
\TextAt{13}{424.04}{314.896}{g}%
\NativeMathAt{434.935}{314.896}{11.955}{\in}%
\TextAt{13}{447.766}{314.896}{S}%
\TextAt{16}{460.522}{314.896}{\Character{106}}%
\TextAt{13}{468.704}{314.896}{g}%
\TextAt{9}{474.738}{314.896}{\Character{40}0\Character{41}}%
\NativeMathAt{494.556}{314.896}{11.955}{\in}%
\TextAt{13}{507.387}{314.896}{U}%
\TextAt{21}{515.337}{313.103}{K}%
\TextAt{16}{523.54}{314.896}{\Character{103}}%
\TextAt{8}{529.518}{314.896}{.}%
\TextAt{8}{79.226}{297.85}{Show:}%
\TextAt{8}{89.579}{277.815}{(i)}%
\TextAt{8}{107.918}{277.815}{If}%
\TextAt{13}{120.707}{277.815}{g}%
\NativeMathAt{131.929}{277.815}{11.955}{\in}%
\TextAt{13}{145.088}{277.815}{S}%
\TextAt{8}{157.976}{277.815}{is}%
\TextAt{8}{170.836}{277.815}{a}%
\TextAt{8}{181.687}{277.815}{power}%
\TextAt{8}{216.927}{277.815}{series}%
\TextAt{8}{249.357}{277.815}{such}%
\TextAt{8}{276.85}{277.815}{that}%
\TextAt{13}{303.315}{277.815}{g}%
\TextAt{9}{309.349}{277.815}{\Character{40}}%
\TextAt{13}{313.902}{277.815}{F}%
\TextAt{9}{323.105}{277.815}{\Character{40}1\Character{41}\Character{41}}%
\TextAt{9}{347.757}{277.815}{\Character{61}}%
\TextAt{9}{362.043}{277.815}{0}%
\TextAt{8}{367.949}{277.815}{,}%
\TextAt{8}{376.468}{277.815}{then}%
\TextAt{13}{404.224}{277.815}{g}%
\TextAt{8}{415.252}{277.815}{is}%
\TextAt{8}{428.112}{277.815}{divisible}%
\TextAt{8}{475.115}{277.815}{by}%
\TextAt{9}{492.47}{277.815}{\Character{91}}%
\TextAt{13}{495.722}{277.815}{\Character{25}}%
\TextAt{9}{502.791}{277.815}{\Character{93}}%
\TextAt{8}{506.043}{277.815}{,}%
\TextAt{8}{514.562}{277.815}{i.e.,}%
\TextAt{13}{107.918}{260.769}{g}%
\TextAt{9}{113.953}{260.769}{\Character{40}}%
\TextAt{13}{118.505}{260.769}{X}%
\TextAt{9}{129.16}{260.769}{\Character{41}}%
\TextAt{9}{137.015}{260.769}{\Character{61}}%
\TextAt{9}{149.436}{260.769}{\Character{91}}%
\TextAt{13}{152.71}{260.769}{\Character{25}}%
\TextAt{9}{159.78}{260.769}{\Character{93}\Character{40}}%
\TextAt{13}{167.584}{260.769}{X}%
\TextAt{9}{178.239}{260.769}{\Character{41}}%
\TextAt{13}{182.791}{260.769}{h}%
\TextAt{9}{189.53}{260.769}{\Character{40}}%
\TextAt{13}{194.082}{260.769}{X}%
\TextAt{9}{204.737}{260.769}{\Character{41}}%
\TextAt{8}{209.289}{260.769}{,}%
\TextAt{13}{216.439}{260.769}{h}%
\TextAt{9}{223.177}{260.769}{\Character{40}}%
\TextAt{13}{227.73}{260.769}{X}%
\TextAt{9}{238.385}{260.769}{\Character{41}}%
\NativeMathAt{246.258}{260.769}{11.955}{\in}%
\TextAt{13}{257.549}{260.769}{S}%
\TextAt{8}{265.444}{260.769}{.}%
\TextAt{8}{86.327}{243.724}{(ii)}%
\TextAt{8}{107.918}{243.724}{Let}%
\TextAt{13}{128.59}{243.724}{g}%
\NativeMathAt{137.945}{243.724}{11.955}{\in}%
\TextAt{13}{149.236}{243.724}{S}%
\TextAt{8}{160.731}{243.724}{be}%
\TextAt{8}{176.358}{243.724}{a}%
\TextAt{8}{185.816}{243.724}{power}%
\TextAt{8}{219.662}{243.724}{series}%
\TextAt{8}{250.699}{243.724}{such}%
\TextAt{8}{276.798}{243.724}{that}%
\TextAt{13}{301.87}{243.724}{g}%
\TextAt{9}{307.904}{243.724}{\Character{40}}%
\TextAt{13}{312.456}{243.724}{X}%
\TextAt{9}{325.159}{243.724}{\Character{43}}%
\TextAt{21}{334.264}{241.93}{F}%
\TextAt{13}{343.346}{243.724}{\Character{21}}%
\TextAt{9}{350.174}{243.724}{\Character{41}}%
\TextAt{9}{358.029}{243.724}{\Character{61}}%
\TextAt{13}{370.473}{243.724}{g}%
\TextAt{9}{376.507}{243.724}{\Character{40}}%
\TextAt{13}{381.059}{243.724}{X}%
\TextAt{9}{391.714}{243.724}{\Character{41}}%
\TextAt{8}{399.867}{243.724}{for}%
\TextAt{8}{417.454}{243.724}{all}%
\TextAt{13}{433.415}{243.724}{\Character{21}}%
\NativeMathAt{443.565}{243.724}{11.955}{\in}%
\TextAt{13}{454.856}{243.724}{F}%
\TextAt{9}{464.059}{243.724}{\Character{40}1\Character{41}}%
\TextAt{8}{479.017}{243.724}{,}%
\TextAt{8}{485.928}{243.724}{where}%
\TextAt{8}{519.44}{243.724}{we}%
\TextAt{8}{107.918}{226.678}{write}%
\TextAt{13}{138.204}{226.678}{X}%
\TextAt{9}{151.771}{226.678}{\Character{43}}%
\TextAt{21}{160.876}{224.884}{F}%
\TextAt{13}{170.821}{226.678}{\Character{21}}%
\TextAt{9}{181.608}{226.678}{\Character{61}}%
\TextAt{13}{194.671}{226.678}{F}%
\TextAt{9}{203.875}{226.678}{\Character{40}}%
\TextAt{13}{208.427}{226.678}{X}%
\TextAt{13}{218.422}{226.678}{\Character{59}}%
\TextAt{13}{223.646}{226.678}{\Character{21}}%
\TextAt{9}{230.505}{226.678}{\Character{41}}%
\TextAt{8}{235.057}{226.678}{.}%
\TextAt{8}{244.628}{226.678}{Then}%
\TextAt{8}{275.559}{226.678}{there}%
\TextAt{8}{305.845}{226.678}{exists}%
\TextAt{8}{338.534}{226.678}{a}%
\TextAt{8}{348.664}{226.678}{unique}%
\TextAt{8}{387.079}{226.678}{power}%
\TextAt{8}{421.597}{226.678}{series}%
\TextAt{13}{453.306}{226.678}{h}%
\TextAt{9}{460.045}{226.678}{\Character{40}}%
\TextAt{13}{464.597}{226.678}{X}%
\TextAt{9}{475.252}{226.678}{\Character{41}}%
\TextAt{8}{484.076}{226.678}{in}%
\TextAt{13}{498.103}{226.678}{S}%
\TextAt{8}{510.27}{226.678}{such}%
\TextAt{8}{107.918}{209.632}{that}%
\TextAt{13}{133.287}{209.632}{g}%
\TextAt{9}{142.642}{209.632}{\Character{61}}%
\TextAt{13}{155.068}{209.632}{h}%
\TextAt{16}{164.463}{209.632}{\Character{14}}%
\TextAt{13}{173.097}{209.632}{\Character{25}}%
\TextAt{8}{180.167}{209.632}{.}%
\TextAt{12}{79.226}{177.642}{Exercise}%
\TextAt{12}{132.427}{177.642}{V.5.2.}%
\TextAt{8}{171.173}{177.642}{If}%
\TextAt{13}{182.866}{177.642}{h}%
\TextAt{9}{189.604}{177.642}{\Character{40}}%
\TextAt{13}{194.156}{177.642}{X}%
\TextAt{9}{204.812}{177.642}{\Character{41}}%
\TextAt{8}{213.261}{177.642}{is}%
\TextAt{8}{225.025}{177.642}{a}%
\TextAt{8}{234.781}{177.642}{power}%
\TextAt{8}{268.924}{177.642}{series}%
\TextAt{8}{300.259}{177.642}{in}%
\TextAt{13}{313.912}{177.642}{S}%
\TextAt{8}{321.807}{177.642}{,}%
\TextAt{8}{328.956}{177.642}{then}%
\TextAt{8}{355.616}{177.642}{the}%
\TextAt{8}{375.773}{177.642}{power}%
\TextAt{8}{409.917}{177.642}{series}%
\TextAt{13}{240.153}{148.641}{h}%
\TextAt{17}{246.892}{146.847}{1}%
\TextAt{9}{251.624}{148.64}{\Character{40}}%
\TextAt{13}{256.176}{148.64}{X}%
\TextAt{9}{266.832}{148.64}{\Character{41}}%
\TextAt{9}{274.687}{148.64}{\Character{61}}%
\TextAt{20}{293.463}{159.998}{\Character{89}}%
\TextAt{21}{287.13}{134.095}{\Character{21}}%
\NativeMathAt{292.07}{134.095}{7.97}{\in}%
\TextAt{21}{297.716}{134.095}{F}%
\TextAt{17}{304.251}{134.095}{\Character{40}1\Character{41}}%
\TextAt{13}{317.064}{148.64}{h}%
\TextAt{9}{323.803}{148.64}{\Character{40}}%
\TextAt{13}{328.355}{148.64}{X}%
\TextAt{9}{341.667}{148.64}{\Character{43}}%
\TextAt{21}{350.771}{146.847}{F}%
\TextAt{13}{360.462}{148.64}{\Character{21}}%
\TextAt{9}{367.29}{148.64}{\Character{41}}%
\TextAt{8}{79.226}{108.989}{also}%
\TextAt{8}{102.706}{108.989}{belongs}%
\TextAt{8}{144.717}{108.989}{to}%
\TextAt{13}{159.027}{108.989}{S}%
\TextAt{8}{166.922}{108.989}{,}%
\TextAt{8}{174.071}{108.989}{and}%
\TextAt{8}{196.834}{108.989}{one}%
\TextAt{8}{218.294}{108.989}{has}%
\TextAt{13}{239.167}{108.989}{h}%
\TextAt{17}{245.906}{107.196}{1}%
\TextAt{9}{250.638}{108.989}{\Character{40}}%
\TextAt{13}{255.191}{108.989}{X}%
\TextAt{9}{268.502}{108.989}{\Character{43}}%
\TextAt{21}{277.607}{107.196}{F}%
\TextAt{13}{287.297}{108.989}{\Character{21}}%
\TextAt{9}{294.126}{108.989}{\Character{41}}%
\TextAt{9}{301.981}{108.989}{\Character{61}}%
\TextAt{13}{314.424}{108.989}{h}%
\TextAt{17}{321.163}{107.196}{1}%
\TextAt{9}{325.895}{108.989}{\Character{40}}%
\TextAt{13}{330.447}{108.989}{X}%
\TextAt{9}{341.102}{108.989}{\Character{41}}%
\TextAt{8}{349.552}{108.989}{for}%
\TextAt{8}{367.437}{108.989}{all}%
\TextAt{13}{383.696}{108.989}{\Character{21}}%
\NativeMathAt{393.845}{108.989}{11.955}{\in}%
\TextAt{13}{405.136}{108.989}{F}%
\TextAt{9}{414.34}{108.989}{\Character{40}1\Character{41}}%
\TextAt{8}{429.298}{108.989}{.}%
\TextAt{8}{297.211}{43.033}{321}%
\Destination{79.23}{294.615}{nav.Item.119}%
\Destination{79.23}{257.785}{nav.Item.120}%
\Destination{78.23}{809.815}{nav.page.321}%
\Destination{79.226}{345.955}{nav.ex.5.5.1}%
\Destination{79.226}{191.655}{nav.ex.5.5.2}%
\Destination{79.36}{473.655}{nav.print.351}%
\Destination{79.36}{582.655}{nav.number.5.5.7}%
\Hotlink{266.339}{592.455}{25.078}{11.047}{nav.number.5.1.9}%
}

% Page 334 (body page 322)
\Page{%
\TextAt{2}{79.226}{796.943}{V.}%
\TextAt{2}{94.073}{796.943}{Local}%
\TextAt{2}{123.615}{796.943}{Class}%
\TextAt{2}{152.208}{796.943}{Field}%
\TextAt{2}{179.939}{796.943}{Theory}%
\Vector{q 0 G 0 g 0.299 w 7.226 18.729 m 460.769 18.729 l S Q}%
\TextAt{12}{79.226}{761.124}{Exercise}%
\TextAt{12}{132.99}{761.124}{V.5.3.}%
\TextAt{8}{172.227}{761.124}{Let}%
\TextAt{13}{193.686}{761.124}{N}%
\TextAt{9}{204.309}{761.124}{\Character{40}}%
\TextAt{13}{208.861}{761.124}{h}%
\TextAt{9}{215.6}{761.124}{\Character{41}}%
\NativeMathAt{224.308}{761.124}{11.955}{\in}%
\TextAt{13}{236.434}{761.124}{S}%
\TextAt{8}{248.716}{761.124}{be}%
\TextAt{8}{265.131}{761.124}{the}%
\TextAt{8}{285.777}{761.124}{power}%
\TextAt{8}{320.411}{761.124}{series}%
\TextAt{8}{352.235}{761.124}{(uniquely}%
\TextAt{8}{404.754}{761.124}{determined}%
\TextAt{8}{466.371}{761.124}{by}%
\TextAt{8}{483.12}{761.124}{exercise}%
\TextAt{8}{526.911}{761.124}{1}%
\TextAt{8}{79.226}{744.078}{and}%
\TextAt{8}{101.989}{744.078}{exercise}%
\TextAt{8}{145.29}{744.078}{2)}%
\TextAt{8}{159.601}{744.078}{such}%
\TextAt{8}{185.998}{744.078}{that}%
\TextAt{13}{228.478}{715.077}{N}%
\TextAt{9}{239.101}{715.077}{\Character{40}}%
\TextAt{13}{243.653}{715.077}{h}%
\TextAt{9}{250.392}{715.077}{\Character{41}}%
\TextAt{16}{257.601}{715.077}{\Character{14}}%
\TextAt{9}{266.235}{715.077}{\Character{91}}%
\TextAt{13}{269.486}{715.077}{\Character{25}}%
\TextAt{9}{276.556}{715.077}{\Character{93}}%
\TextAt{9}{283.107}{715.077}{\Character{61}}%
\TextAt{20}{301.886}{726.434}{\Character{89}}%
\TextAt{21}{295.554}{700.531}{\Character{21}}%
\NativeMathAt{300.494}{700.531}{7.97}{\in}%
\TextAt{21}{306.139}{700.531}{F}%
\TextAt{17}{312.674}{700.531}{\Character{40}1\Character{41}}%
\TextAt{13}{325.488}{715.077}{h}%
\TextAt{9}{332.226}{715.077}{\Character{40}}%
\TextAt{13}{336.779}{715.077}{X}%
\TextAt{9}{350.09}{715.077}{\Character{43}}%
\TextAt{21}{359.195}{713.283}{F}%
\TextAt{13}{368.885}{715.077}{\Character{21}}%
\TextAt{9}{375.714}{715.077}{\Character{41}}%
\TextAt{13}{380.266}{715.077}{\Character{58}}%
\TextAt{8}{79.226}{675.975}{The}%
\TextAt{8}{103.28}{675.975}{mapping}%
\TextAt{13}{151.411}{675.975}{N}%
\TextAt{9}{165.355}{675.975}{\Character{58}}%
\TextAt{13}{171.927}{675.975}{S}%
\NativeArrowAt{183.144}{675.975}{11.955}{19.26}{\longrightarrow}%
\TextAt{13}{205.726}{675.975}{S}%
\TextAt{8}{217.518}{675.975}{is}%
\TextAt{8}{229.282}{675.975}{called}%
\TextAt{12}{262.446}{675.975}{Coleman’s}%
\TextAt{12}{327.53}{675.975}{norm}%
\TextAt{12}{362.929}{675.975}{operator}%
\TextAt{8}{413.141}{675.975}{.}%
\TextAt{8}{421.589}{675.975}{Show:}%
\TextAt{8}{89.579}{655.94}{(i)}%
\TextAt{13}{107.918}{655.94}{N}%
\TextAt{9}{118.541}{655.94}{\Character{40}}%
\TextAt{13}{123.093}{655.94}{h}%
\TextAt{17}{129.832}{654.147}{1}%
\TextAt{13}{134.564}{655.94}{h}%
\TextAt{17}{141.303}{654.147}{2}%
\TextAt{9}{146.035}{655.94}{\Character{41}}%
\TextAt{9}{153.89}{655.94}{\Character{61}}%
\TextAt{13}{166.334}{655.94}{N}%
\TextAt{9}{176.956}{655.94}{\Character{40}}%
\TextAt{13}{181.509}{655.94}{h}%
\TextAt{17}{188.247}{654.147}{1}%
\TextAt{9}{192.98}{655.94}{\Character{41}}%
\TextAt{13}{197.532}{655.94}{N}%
\TextAt{9}{208.155}{655.94}{\Character{40}}%
\TextAt{13}{212.707}{655.94}{h}%
\TextAt{17}{219.445}{654.147}{2}%
\TextAt{9}{224.178}{655.94}{\Character{41}}%
\TextAt{8}{228.73}{655.94}{.}%
\TextAt{8}{86.327}{638.894}{(ii)}%
\TextAt{13}{107.918}{638.894}{N}%
\TextAt{9}{118.541}{638.894}{\Character{40}}%
\TextAt{13}{123.093}{638.894}{h}%
\TextAt{9}{129.832}{638.894}{\Character{41}}%
\NativeMathAt{137.705}{638.894}{11.955}{\equiv}%
\TextAt{13}{150.324}{638.894}{h}%
\TextAt{8}{160.384}{638.894}{mod}%
\TextAt{14}{185.822}{638.894}{p}%
\TextAt{8}{191.805}{638.894}{.}%
\TextAt{8}{83.076}{621.848}{(iii)}%
\TextAt{13}{107.918}{621.848}{h}%
\NativeMathAt{117.978}{621.848}{11.955}{\in}%
\TextAt{13}{129.269}{621.848}{X}%
\TextAt{21}{139.924}{626.187}{i}%
\TextAt{13}{143.305}{621.848}{S}%
\TextAt{3}{151.201}{626.187}{\Character{3}}%
\TextAt{8}{159.83}{621.848}{for}%
\TextAt{13}{177.715}{621.848}{i}%
\NativeMathAt{185.029}{621.848}{11.955}{\geq}%
\TextAt{9}{197.649}{621.848}{0}%
\NativeArrowAt{206.807}{621.848}{11.955}{19.083}{\Longrightarrow}%
\TextAt{13}{229.211}{621.848}{N}%
\TextAt{9}{239.833}{621.848}{\Character{40}}%
\TextAt{13}{244.386}{621.848}{h}%
\TextAt{9}{251.124}{621.848}{\Character{41}}%
\NativeMathAt{258.997}{621.848}{11.955}{\in}%
\TextAt{13}{270.288}{621.848}{X}%
\TextAt{21}{280.943}{626.187}{i}%
\TextAt{13}{284.325}{621.848}{S}%
\TextAt{3}{292.22}{626.187}{\Character{3}}%
\TextAt{8}{296.952}{621.848}{.}%
\TextAt{8}{83.398}{604.802}{(iv)}%
\TextAt{13}{107.918}{604.802}{h}%
\NativeMathAt{117.978}{604.802}{11.955}{\equiv}%
\TextAt{9}{130.597}{604.802}{1}%
\TextAt{8}{139.771}{604.802}{mod}%
\TextAt{14}{165.209}{604.802}{p}%
\TextAt{21}{171.192}{609.141}{i}%
\TextAt{8}{178.471}{604.802}{for}%
\TextAt{13}{196.356}{604.802}{i}%
\NativeMathAt{203.67}{604.802}{11.955}{\geq}%
\TextAt{9}{216.289}{604.802}{1}%
\NativeArrowAt{225.447}{604.802}{11.955}{19.083}{\Longrightarrow}%
\TextAt{13}{247.851}{604.802}{N}%
\TextAt{9}{258.474}{604.802}{\Character{40}}%
\TextAt{13}{263.026}{604.802}{h}%
\TextAt{9}{269.765}{604.802}{\Character{41}}%
\NativeMathAt{277.638}{604.802}{11.955}{\equiv}%
\TextAt{9}{290.257}{604.802}{1}%
\TextAt{8}{299.431}{604.802}{mod}%
\TextAt{14}{324.869}{604.802}{p}%
\TextAt{21}{330.852}{609.141}{i}%
\TextAt{17}{333.735}{609.141}{\Character{43}1}%
\TextAt{8}{345.054}{604.802}{.}%
\TextAt{8}{86.65}{587.756}{(v)}%
\TextAt{8}{107.918}{587.756}{For}%
\TextAt{8}{128.888}{587.756}{the}%
\TextAt{8}{149.044}{587.756}{operators}%
\TextAt{13}{200.822}{587.756}{N}%
\TextAt{17}{211.445}{592.095}{0}%
\TextAt{9}{216.177}{587.756}{\Character{40}}%
\TextAt{13}{220.729}{587.756}{h}%
\TextAt{9}{227.468}{587.756}{\Character{41}}%
\TextAt{9}{235.323}{587.756}{\Character{61}}%
\TextAt{13}{247.767}{587.756}{h}%
\TextAt{8}{254.505}{587.756}{,}%
\TextAt{13}{261.654}{587.756}{N}%
\TextAt{21}{272.277}{592.095}{n}%
\TextAt{9}{277.913}{587.756}{\Character{40}}%
\TextAt{13}{282.466}{587.756}{h}%
\TextAt{9}{289.204}{587.756}{\Character{41}}%
\TextAt{9}{297.059}{587.756}{\Character{61}}%
\TextAt{13}{309.503}{587.756}{N}%
\TextAt{9}{320.125}{587.756}{\Character{40}}%
\TextAt{13}{324.678}{587.756}{N}%
\TextAt{21}{335.3}{592.095}{n}%
\TextAt{3}{340.439}{592.095}{\Character{0}}%
\TextAt{17}{347.025}{592.095}{1}%
\TextAt{9}{351.757}{587.756}{\Character{40}}%
\TextAt{13}{356.31}{587.756}{h}%
\TextAt{9}{363.048}{587.756}{\Character{41}\Character{41}}%
\TextAt{8}{372.153}{587.756}{,}%
\TextAt{8}{379.302}{587.756}{one}%
\TextAt{8}{400.762}{587.756}{has}%
\TextAt{13}{207.768}{558.755}{N}%
\TextAt{21}{218.391}{563.692}{n}%
\TextAt{9}{224.027}{558.755}{\Character{40}}%
\TextAt{13}{228.58}{558.755}{h}%
\TextAt{9}{235.318}{558.755}{\Character{41}}%
\TextAt{16}{242.527}{558.755}{\Character{14}}%
\TextAt{9}{251.161}{558.755}{\Character{91}}%
\TextAt{13}{254.413}{558.755}{\Character{25}}%
\TextAt{21}{261.482}{563.692}{n}%
\TextAt{9}{267.119}{558.755}{\Character{93}}%
\TextAt{9}{273.67}{558.755}{\Character{61}}%
\TextAt{20}{292.901}{570.113}{\Character{89}}%
\TextAt{21}{286.117}{544.21}{\Character{21}}%
\NativeMathAt{291.057}{544.21}{7.97}{\in}%
\TextAt{21}{296.702}{544.21}{F}%
\TextAt{17}{303.237}{544.21}{\Character{40}}%
\TextAt{21}{306.531}{544.21}{n}%
\TextAt{17}{311.669}{544.21}{\Character{41}}%
\TextAt{13}{316.955}{558.755}{h}%
\TextAt{9}{323.693}{558.755}{\Character{40}}%
\TextAt{13}{328.245}{558.755}{X}%
\TextAt{9}{341.557}{558.755}{\Character{43}}%
\TextAt{21}{350.662}{556.962}{F}%
\TextAt{13}{360.352}{558.755}{\Character{21}}%
\TextAt{9}{367.18}{558.755}{\Character{41}}%
\TextAt{13}{371.733}{558.755}{\Character{59}}%
\TextAt{13}{400.868}{558.755}{n}%
\NativeMathAt{411.196}{558.755}{11.955}{\geq}%
\TextAt{9}{423.815}{558.755}{0}%
\TextAt{13}{429.668}{558.755}{\Character{58}}%
\TextAt{8}{83.398}{518.463}{(vi)}%
\TextAt{8}{107.918}{518.463}{If}%
\TextAt{13}{119.611}{518.463}{h}%
\NativeMathAt{129.67}{518.463}{11.955}{\in}%
\TextAt{13}{140.961}{518.463}{X}%
\TextAt{21}{151.616}{522.801}{i}%
\TextAt{13}{154.997}{518.463}{S}%
\TextAt{3}{162.893}{522.801}{\Character{3}}%
\TextAt{8}{167.625}{518.463}{,}%
\TextAt{13}{174.774}{518.463}{i}%
\NativeMathAt{182.088}{518.463}{11.955}{\geq}%
\TextAt{9}{194.708}{518.463}{0}%
\TextAt{8}{200.561}{518.463}{,}%
\TextAt{8}{207.71}{518.463}{then}%
\TextAt{13}{234.37}{518.463}{N}%
\TextAt{21}{244.993}{522.801}{n}%
\TextAt{17}{250.131}{522.801}{\Character{43}1}%
\TextAt{9}{261.45}{518.463}{\Character{40}}%
\TextAt{13}{266.002}{518.463}{h}%
\TextAt{9}{272.741}{518.463}{\Character{41}\Character{47}}%
\TextAt{13}{283.146}{518.463}{N}%
\TextAt{21}{293.768}{522.801}{n}%
\TextAt{9}{299.405}{518.463}{\Character{40}}%
\TextAt{13}{303.957}{518.463}{h}%
\TextAt{9}{310.696}{518.463}{\Character{41}}%
\NativeMathAt{318.569}{518.463}{11.955}{\in}%
\TextAt{13}{329.86}{518.463}{S}%
\TextAt{3}{337.755}{522.801}{\Character{3}}%
\TextAt{8}{346.385}{518.463}{and}%
\TextAt{13}{218.669}{489.462}{N}%
\TextAt{21}{229.291}{494.398}{n}%
\TextAt{17}{234.429}{494.398}{\Character{43}1}%
\TextAt{9}{245.748}{489.462}{\Character{40}}%
\TextAt{13}{250.3}{489.462}{h}%
\TextAt{9}{257.039}{489.462}{\Character{41}}%
\NativeMathAt{264.912}{489.462}{11.955}{\equiv}%
\TextAt{13}{277.532}{489.462}{N}%
\TextAt{21}{288.154}{494.398}{n}%
\TextAt{9}{293.79}{489.462}{\Character{40}}%
\TextAt{13}{298.343}{489.462}{h}%
\TextAt{9}{305.081}{489.462}{\Character{41}}%
\TextAt{8}{312.955}{489.462}{mod}%
\TextAt{14}{338.392}{489.462}{p}%
\TextAt{21}{344.375}{494.398}{n}%
\TextAt{17}{349.514}{494.398}{\Character{43}1}%
\TextAt{13}{360.833}{489.462}{\Character{59}}%
\TextAt{13}{389.968}{489.462}{n}%
\NativeMathAt{400.295}{489.462}{11.955}{\geq}%
\TextAt{9}{412.915}{489.462}{0}%
\TextAt{13}{418.768}{489.462}{\Character{58}}%
\TextAt{12}{79.226}{448.506}{Exercise}%
\TextAt{12}{133.585}{448.506}{V.5.4.}%
\TextAt{8}{173.339}{448.506}{Let}%
\TextAt{13}{195.315}{448.506}{\Character{21}}%
\NativeMathAt{207.181}{448.506}{11.955}{\in}%
\TextAt{13}{220.188}{448.506}{F}%
\TextAt{9}{229.392}{448.506}{\Character{40}}%
\TextAt{13}{233.944}{448.506}{n}%
\TextAt{9}{244.275}{448.506}{\Character{43}}%
\TextAt{9}{256.708}{448.506}{1\Character{41}}%
\TextAt{16}{270.471}{448.506}{\Character{110}}%
\TextAt{13}{279.792}{448.506}{F}%
\TextAt{9}{288.996}{448.506}{\Character{40}}%
\TextAt{13}{293.548}{448.506}{n}%
\TextAt{9}{300.536}{448.506}{\Character{41}}%
\TextAt{8}{305.088}{448.506}{,}%
\TextAt{13}{313.496}{448.506}{n}%
\NativeMathAt{325.521}{448.506}{11.955}{\geq}%
\TextAt{9}{339.856}{448.506}{0}%
\TextAt{8}{345.709}{448.506}{,}%
\TextAt{8}{354.117}{448.506}{and}%
\TextAt{13}{377.887}{448.506}{\Character{21}}%
\TextAt{21}{384.716}{446.712}{i}%
\TextAt{9}{393.134}{448.506}{\Character{61}}%
\TextAt{9}{407.265}{448.506}{\Character{91}}%
\TextAt{13}{410.527}{448.506}{\Character{25}}%
\TextAt{21}{417.597}{452.844}{n}%
\TextAt{3}{422.735}{452.844}{\Character{0}}%
\TextAt{21}{429.321}{452.844}{i}%
\TextAt{9}{432.703}{448.506}{\Character{93}\Character{40}}%
\TextAt{13}{440.507}{448.506}{\Character{21}}%
\TextAt{9}{447.335}{448.506}{\Character{41}}%
\NativeMathAt{456.925}{448.506}{11.955}{\in}%
\TextAt{13}{469.932}{448.506}{F}%
\TextAt{9}{479.135}{448.506}{\Character{40}}%
\TextAt{13}{483.688}{448.506}{i}%
\TextAt{9}{491.024}{448.506}{\Character{43}}%
\TextAt{9}{503.457}{448.506}{1\Character{41}}%
\TextAt{8}{518.782}{448.506}{for}%
\TextAt{9}{79.226}{431.46}{0}%
\NativeMathAt{88.536}{431.46}{11.955}{\leq}%
\TextAt{13}{101.291}{431.46}{i}%
\NativeMathAt{108.741}{431.46}{11.955}{\leq}%
\TextAt{13}{121.496}{431.46}{n}%
\TextAt{8}{128.484}{431.46}{.}%
\TextAt{8}{137.171}{431.46}{Then}%
\TextAt{13}{167.808}{431.46}{\Character{21}}%
\TextAt{21}{174.637}{429.666}{i}%
\TextAt{8}{181.995}{431.46}{is}%
\TextAt{8}{193.839}{431.46}{a}%
\TextAt{8}{203.674}{431.46}{prime}%
\TextAt{8}{236.917}{431.46}{element}%
\TextAt{8}{280.239}{431.46}{of}%
\TextAt{8}{293.648}{431.46}{the}%
\TextAt{8}{313.884}{431.46}{Lubin–Tate}%
\TextAt{8}{376.884}{431.46}{extension}%
\TextAt{13}{428.73}{431.46}{L}%
\TextAt{21}{436.694}{429.666}{i}%
\TextAt{17}{439.577}{429.666}{\Character{43}1}%
\TextAt{9}{454.353}{431.46}{\Character{61}}%
\TextAt{13}{466.914}{431.46}{K}%
\TextAt{9}{477.732}{431.46}{\Character{40}}%
\TextAt{13}{482.284}{431.46}{F}%
\TextAt{9}{491.488}{431.46}{\Character{40}}%
\TextAt{13}{496.04}{431.46}{i}%
\TextAt{9}{502.744}{431.46}{\Character{43}}%
\TextAt{9}{514.544}{431.46}{1\Character{41}\Character{41}}%
\TextAt{8}{529.518}{431.46}{,}%
\TextAt{8}{79.226}{414.414}{and}%
\TextAt{14}{101.989}{414.414}{o}%
\TextAt{21}{107.833}{412.62}{i}%
\TextAt{17}{110.716}{412.62}{\Character{43}1}%
\TextAt{9}{125.356}{414.414}{\Character{61}}%
\TextAt{14}{137.781}{414.414}{o}%
\TextAt{21}{143.626}{412.62}{K}%
\TextAt{9}{151.828}{414.414}{\Character{91}}%
\TextAt{13}{155.08}{414.414}{\Character{21}}%
\TextAt{21}{161.908}{412.62}{i}%
\TextAt{9}{165.29}{414.414}{\Character{93}}%
\TextAt{8}{172.439}{414.414}{is}%
\TextAt{8}{184.203}{414.414}{the}%
\TextAt{8}{204.359}{414.414}{valuation}%
\TextAt{8}{255.431}{414.414}{ring}%
\TextAt{8}{279.497}{414.414}{of}%
\TextAt{13}{292.827}{414.414}{L}%
\TextAt{21}{300.792}{412.62}{i}%
\TextAt{17}{303.675}{412.62}{\Character{43}1}%
\TextAt{8}{314.994}{414.414}{,}%
\TextAt{8}{322.143}{414.414}{with}%
\TextAt{8}{348.803}{414.414}{maximal}%
\TextAt{8}{396.612}{414.414}{ideal}%
\TextAt{14}{424.575}{414.414}{p}%
\TextAt{21}{430.558}{412.62}{i}%
\TextAt{17}{433.441}{412.62}{\Character{43}1}%
\TextAt{9}{448.081}{414.414}{\Character{61}}%
\TextAt{13}{460.506}{414.414}{\Character{21}}%
\TextAt{21}{467.335}{412.62}{i}%
\TextAt{14}{470.716}{414.414}{o}%
\TextAt{21}{476.56}{412.62}{i}%
\TextAt{17}{479.443}{412.62}{\Character{43}1}%
\TextAt{8}{490.762}{414.414}{.}%
\TextAt{8}{499.21}{414.414}{Show:}%
\TextAt{8}{97.159}{394.379}{let}%
\TextAt{13}{114.364}{394.379}{\Character{12}}%
\TextAt{21}{120.972}{392.586}{i}%
\NativeMathAt{128.185}{394.379}{11.955}{\in}%
\TextAt{13}{139.988}{394.379}{\Character{25}}%
\TextAt{21}{147.057}{398.717}{n}%
\TextAt{3}{152.195}{398.717}{\Character{0}}%
\TextAt{21}{158.782}{398.717}{i}%
\TextAt{14}{162.163}{394.379}{p}%
\TextAt{17}{168.146}{392.586}{1}%
\TextAt{14}{172.878}{394.379}{o}%
\TextAt{21}{178.723}{392.586}{i}%
\TextAt{17}{181.606}{392.586}{\Character{43}1}%
\TextAt{8}{192.925}{394.379}{,}%
\TextAt{9}{200.449}{394.379}{0}%
\NativeMathAt{210.134}{394.379}{11.955}{\leq}%
\TextAt{13}{223.265}{394.379}{i}%
\NativeMathAt{231.09}{394.379}{11.955}{\leq}%
\TextAt{13}{244.221}{394.379}{n}%
\TextAt{8}{251.209}{394.379}{.}%
\TextAt{8}{260.557}{394.379}{Then}%
\TextAt{8}{291.415}{394.379}{there}%
\TextAt{8}{321.627}{394.379}{exists}%
\TextAt{8}{354.241}{394.379}{a}%
\TextAt{8}{364.297}{394.379}{power}%
\TextAt{8}{398.741}{394.379}{series}%
\TextAt{13}{430.375}{394.379}{h}%
\TextAt{9}{437.114}{394.379}{\Character{40}}%
\TextAt{13}{441.666}{394.379}{X}%
\TextAt{9}{452.321}{394.379}{\Character{41}}%
\NativeMathAt{460.706}{394.379}{11.955}{\in}%
\TextAt{13}{472.508}{394.379}{S}%
\TextAt{8}{484.601}{394.379}{such}%
\TextAt{8}{511.298}{394.379}{that}%
\TextAt{13}{79.226}{377.333}{h}%
\TextAt{9}{85.965}{377.333}{\Character{40}}%
\TextAt{13}{90.517}{377.333}{\Character{21}}%
\TextAt{21}{97.346}{375.54}{i}%
\TextAt{9}{100.727}{377.333}{\Character{41}}%
\TextAt{9}{108.582}{377.333}{\Character{61}}%
\TextAt{13}{121.025}{377.333}{\Character{12}}%
\TextAt{21}{127.634}{375.54}{i}%
\TextAt{8}{134.913}{377.333}{for}%
\TextAt{9}{152.797}{377.333}{0}%
\NativeMathAt{161.971}{377.333}{11.955}{\leq}%
\TextAt{13}{174.591}{377.333}{i}%
\NativeMathAt{181.905}{377.333}{11.955}{\leq}%
\TextAt{13}{194.524}{377.333}{n}%
\TextAt{8}{201.512}{377.333}{.}%
\TextAt{12}{79.226}{357.298}{Hint:}%
\TextAt{8}{114.765}{357.298}{Write}%
\TextAt{13}{147.271}{357.298}{\Character{12}}%
\TextAt{21}{153.879}{355.505}{i}%
\TextAt{9}{160.582}{357.298}{\Character{61}}%
\TextAt{13}{173.007}{357.298}{\Character{25}}%
\TextAt{21}{180.076}{361.636}{n}%
\TextAt{3}{185.215}{361.637}{\Character{0}}%
\TextAt{21}{191.801}{361.636}{i}%
\TextAt{13}{195.182}{357.298}{\Character{21}}%
\TextAt{17}{202.011}{355.505}{0}%
\TextAt{13}{206.743}{357.298}{h}%
\TextAt{21}{213.482}{355.505}{i}%
\TextAt{9}{216.863}{357.298}{\Character{40}}%
\TextAt{13}{221.415}{357.298}{\Character{21}}%
\TextAt{21}{228.244}{355.505}{i}%
\TextAt{9}{231.625}{357.298}{\Character{41}}%
\TextAt{8}{236.177}{357.298}{,}%
\TextAt{8}{243.327}{357.298}{with}%
\TextAt{13}{269.987}{357.298}{h}%
\TextAt{21}{276.725}{355.505}{i}%
\TextAt{9}{280.106}{357.298}{\Character{40}}%
\TextAt{13}{284.659}{357.298}{X}%
\TextAt{9}{295.314}{357.298}{\Character{41}}%
\NativeMathAt{303.187}{357.298}{11.955}{\in}%
\TextAt{14}{314.478}{357.298}{o}%
\TextAt{9}{320.322}{357.298}{\Character{91}}%
\TextAt{13}{323.574}{357.298}{X}%
\TextAt{9}{334.229}{357.298}{\Character{93}}%
\TextAt{8}{341.378}{357.298}{and}%
\TextAt{8}{364.141}{357.298}{put,}%
\TextAt{8}{388.852}{357.298}{for}%
\TextAt{9}{406.737}{357.298}{0}%
\NativeMathAt{415.911}{357.298}{11.955}{\leq}%
\TextAt{13}{428.53}{357.298}{i}%
\NativeMathAt{435.844}{357.298}{11.955}{\leq}%
\TextAt{13}{448.464}{357.298}{n}%
\TextAt{8}{455.451}{357.298}{:}%
\TextAt{13}{257.438}{323.663}{g}%
\TextAt{21}{263.044}{321.869}{i}%
\TextAt{9}{266.425}{323.663}{\Character{40}}%
\TextAt{13}{270.977}{323.663}{X}%
\TextAt{9}{281.632}{323.663}{\Character{41}}%
\TextAt{9}{289.487}{323.663}{\Character{61}}%
\TextAt{9}{303.127}{331.75}{\Character{91}}%
\TextAt{13}{306.378}{331.75}{\Character{25}}%
\TextAt{21}{313.447}{336.089}{n}%
\TextAt{17}{318.586}{336.089}{\Character{43}1}%
\TextAt{9}{329.905}{331.75}{\Character{93}\Character{91}}%
\TextAt{13}{336.408}{331.75}{\Character{25}}%
\TextAt{21}{343.477}{336.089}{i}%
\TextAt{9}{346.858}{331.75}{\Character{93}}%
\Vector{q 0 G 0 g 0.478 w 231.127 -447.348 m 278.11 -447.348 l S Q}%
\TextAt{9}{312.731}{315.462}{\Character{91}}%
\TextAt{13}{315.983}{315.462}{\Character{25}}%
\TextAt{21}{323.052}{318.916}{i}%
\TextAt{17}{325.935}{318.916}{\Character{43}1}%
\TextAt{9}{337.254}{315.462}{\Character{93}}%
\TextAt{13}{351.306}{323.663}{\Character{58}}%
\TextAt{8}{79.226}{290.279}{Then}%
\TextAt{13}{109.784}{290.279}{h}%
\TextAt{9}{119.843}{290.279}{\Character{61}}%
\TextAt{20}{132.268}{299.245}{\Character{80}}%
\TextAt{21}{144.888}{296.09}{n}%
\TextAt{21}{144.888}{286.792}{i}%
\TextAt{17}{147.771}{286.792}{\Character{61}0}%
\TextAt{13}{161.082}{290.279}{h}%
\TextAt{21}{167.821}{288.485}{i}%
\TextAt{13}{171.202}{290.279}{g}%
\TextAt{21}{176.807}{288.485}{i}%
\TextAt{8}{184.086}{290.279}{is}%
\TextAt{8}{195.85}{290.279}{a}%
\TextAt{8}{205.605}{290.279}{solution.}%
\TextAt{12}{79.226}{258.289}{Exercise}%
\TextAt{12}{133.61}{258.289}{V.5.5.}%
\TextAt{8}{173.385}{258.289}{Let}%
\TextAt{13}{195.384}{258.289}{\Character{21}}%
\NativeMathAt{207.286}{258.289}{11.955}{\in}%
\TextAt{13}{220.33}{258.289}{F}%
\TextAt{9}{229.534}{258.289}{\Character{40}}%
\TextAt{13}{234.086}{258.289}{n}%
\TextAt{9}{244.432}{258.289}{\Character{43}}%
\TextAt{9}{256.877}{258.289}{1\Character{41}}%
\TextAt{16}{270.657}{258.289}{\Character{110}}%
\TextAt{13}{279.993}{258.289}{F}%
\TextAt{9}{289.196}{258.289}{\Character{40}}%
\TextAt{13}{293.749}{258.289}{n}%
\TextAt{9}{300.736}{258.289}{\Character{41}}%
\TextAt{8}{310.215}{258.289}{and}%
\TextAt{13}{334.006}{258.289}{\Character{21}}%
\TextAt{21}{340.835}{256.495}{i}%
\TextAt{9}{349.29}{258.289}{\Character{61}}%
\TextAt{9}{363.457}{258.289}{\Character{91}}%
\TextAt{13}{366.72}{258.289}{\Character{25}}%
\TextAt{21}{373.789}{262.627}{n}%
\TextAt{3}{378.928}{262.627}{\Character{0}}%
\TextAt{21}{385.514}{262.627}{i}%
\TextAt{9}{388.895}{258.289}{\Character{93}\Character{40}}%
\TextAt{13}{396.699}{258.289}{\Character{21}}%
\TextAt{9}{403.528}{258.289}{\Character{41}}%
\TextAt{8}{408.08}{258.289}{,}%
\TextAt{9}{416.515}{258.289}{0}%
\NativeMathAt{427.442}{258.289}{11.955}{\leq}%
\TextAt{13}{441.815}{258.289}{i}%
\NativeMathAt{450.882}{258.289}{11.955}{\leq}%
\TextAt{13}{465.255}{258.289}{n}%
\TextAt{8}{472.242}{258.289}{.}%
\TextAt{8}{483.777}{258.289}{For}%
\TextAt{8}{505.775}{258.289}{every}%
\TextAt{13}{79.226}{241.243}{u}%
\NativeMathAt{89.209}{241.243}{11.955}{\in}%
\TextAt{13}{100.5}{241.243}{U}%
\TextAt{21}{108.451}{239.449}{L}%
\TextAt{22}{114.21}{238.343}{n}%
\TextAt{24}{118.81}{238.343}{\Character{43}1}%
\TextAt{8}{129.05}{241.243}{,}%
\TextAt{8}{136.199}{241.243}{there}%
\TextAt{8}{166.111}{241.243}{exists}%
\TextAt{8}{198.426}{241.243}{a}%
\TextAt{8}{208.181}{241.243}{power}%
\TextAt{8}{242.325}{241.243}{series}%
\TextAt{13}{273.659}{241.243}{h}%
\TextAt{9}{280.398}{241.243}{\Character{40}}%
\TextAt{13}{284.95}{241.243}{X}%
\TextAt{9}{295.605}{241.243}{\Character{41}}%
\NativeMathAt{303.479}{241.243}{11.955}{\in}%
\TextAt{14}{314.77}{241.243}{o}%
\TextAt{9}{320.614}{241.243}{\Character{91}\Character{91}}%
\TextAt{13}{327.117}{241.243}{X}%
\TextAt{9}{337.772}{241.243}{\Character{93}\Character{93}}%
\TextAt{8}{348.173}{241.243}{such}%
\TextAt{8}{374.57}{241.243}{that}%
\TextAt{13}{220.211}{212.242}{N}%
\TextAt{21}{229.585}{210.448}{n\Character{59}i}%
\TextAt{9}{240.457}{212.242}{\Character{40}}%
\TextAt{13}{245.009}{212.242}{u}%
\TextAt{9}{251.672}{212.242}{\Character{41}}%
\TextAt{9}{259.527}{212.242}{\Character{61}}%
\TextAt{13}{271.97}{212.242}{h}%
\TextAt{9}{278.709}{212.242}{\Character{40}}%
\TextAt{13}{283.261}{212.242}{\Character{21}}%
\TextAt{21}{290.09}{210.448}{i}%
\TextAt{9}{293.471}{212.242}{\Character{41}}%
\TextAt{8}{321.934}{212.242}{for}%
\TextAt{9}{339.818}{212.242}{0}%
\NativeMathAt{348.992}{212.242}{11.955}{\leq}%
\TextAt{13}{361.612}{212.242}{i}%
\NativeMathAt{368.926}{212.242}{11.955}{\leq}%
\TextAt{13}{381.545}{212.242}{n\Character{59}}%
\TextAt{8}{79.226}{183.24}{where}%
\TextAt{13}{113.035}{183.24}{N}%
\TextAt{21}{122.409}{181.447}{n\Character{59}i}%
\TextAt{8}{137.178}{183.24}{is}%
\TextAt{8}{148.942}{183.24}{the}%
\TextAt{8}{169.099}{183.24}{norm}%
\TextAt{8}{199.668}{183.24}{from}%
\TextAt{13}{227.308}{183.24}{L}%
\TextAt{21}{235.273}{181.447}{n}%
\TextAt{17}{240.411}{181.447}{\Character{43}1}%
\TextAt{8}{255.627}{183.24}{to}%
\TextAt{13}{269.938}{183.24}{L}%
\TextAt{21}{277.902}{181.447}{i}%
\TextAt{17}{280.785}{181.447}{\Character{43}1}%
\TextAt{8}{292.104}{183.24}{.}%
\TextAt{12}{79.226}{163.206}{Hint:}%
\TextAt{8}{116.913}{163.206}{Write}%
\TextAt{13}{150.494}{163.206}{u}%
\TextAt{9}{162.308}{163.206}{\Character{61}}%
\TextAt{13}{176.564}{163.206}{h}%
\TextAt{17}{183.302}{161.412}{1}%
\TextAt{9}{188.035}{163.206}{\Character{40}}%
\TextAt{13}{192.587}{163.206}{\Character{21}}%
\TextAt{9}{199.416}{163.206}{\Character{41}}%
\TextAt{8}{203.968}{163.206}{,}%
\TextAt{13}{212.46}{163.206}{h}%
\TextAt{17}{219.198}{161.412}{1}%
\TextAt{9}{223.931}{163.206}{\Character{40}}%
\TextAt{13}{228.483}{163.206}{X}%
\TextAt{9}{239.138}{163.206}{\Character{41}}%
\NativeMathAt{248.842}{163.206}{11.955}{\in}%
\TextAt{14}{261.964}{163.206}{o}%
\TextAt{9}{267.808}{163.206}{\Character{91}}%
\TextAt{13}{271.06}{163.206}{X}%
\TextAt{9}{281.715}{163.206}{\Character{93}}%
\TextAt{8}{284.966}{163.206}{,}%
\TextAt{8}{293.458}{163.206}{and}%
\TextAt{8}{317.295}{163.206}{put}%
\TextAt{13}{339.829}{163.206}{h}%
\TextAt{17}{346.568}{161.412}{2}%
\TextAt{9}{356.452}{163.206}{\Character{61}}%
\TextAt{13}{370.708}{163.206}{N}%
\TextAt{21}{381.33}{167.544}{n}%
\TextAt{9}{386.967}{163.206}{\Character{40}}%
\TextAt{13}{391.519}{163.206}{h}%
\TextAt{17}{398.257}{161.412}{1}%
\TextAt{9}{402.99}{163.206}{\Character{41}}%
\NativeMathAt{412.694}{163.206}{11.955}{\in}%
\TextAt{13}{425.815}{163.206}{S}%
\TextAt{3}{433.711}{167.544}{\Character{3}}%
\TextAt{8}{438.443}{163.206}{.}%
\TextAt{8}{450.114}{163.206}{Show}%
\TextAt{8}{482.08}{163.206}{that}%
\TextAt{13}{508.524}{163.206}{\Character{12}}%
\TextAt{21}{515.132}{161.412}{i}%
\TextAt{9}{523.665}{163.206}{\Character{61}}%
\TextAt{13}{79.226}{146.16}{N}%
\TextAt{21}{88.6}{144.366}{n\Character{59}i}%
\TextAt{9}{99.472}{146.16}{\Character{40}}%
\TextAt{13}{104.024}{146.16}{u}%
\TextAt{9}{110.687}{146.16}{\Character{41}}%
\TextAt{16}{118.031}{146.16}{\Character{0}}%
\TextAt{13}{130.121}{146.16}{h}%
\TextAt{17}{136.86}{144.366}{2}%
\TextAt{9}{141.592}{146.16}{\Character{40}}%
\TextAt{13}{146.144}{146.16}{\Character{21}}%
\TextAt{21}{152.973}{144.366}{i}%
\TextAt{9}{156.354}{146.16}{\Character{41}}%
\NativeMathAt{164.566}{146.16}{11.955}{\in}%
\TextAt{13}{176.195}{146.16}{\Character{25}}%
\TextAt{21}{183.264}{150.498}{n}%
\TextAt{3}{188.402}{150.498}{\Character{0}}%
\TextAt{21}{194.989}{150.498}{i}%
\TextAt{14}{198.37}{146.16}{p}%
\TextAt{17}{204.353}{144.366}{1}%
\TextAt{14}{209.085}{146.16}{o}%
\TextAt{21}{214.93}{144.366}{i}%
\TextAt{17}{217.813}{144.366}{\Character{43}1}%
\TextAt{8}{229.132}{146.16}{.}%
\TextAt{8}{238.175}{146.16}{Then}%
\TextAt{8}{268.931}{146.16}{by}%
\TextAt{8}{285.389}{146.16}{exercise}%
\TextAt{8}{328.889}{146.16}{4}%
\TextAt{8}{338.843}{146.16}{there}%
\TextAt{8}{368.953}{146.16}{is}%
\TextAt{8}{380.915}{146.16}{a}%
\TextAt{8}{390.869}{146.16}{power}%
\TextAt{8}{425.212}{146.16}{series}%
\TextAt{13}{456.745}{146.16}{h}%
\TextAt{17}{463.483}{144.366}{3}%
\TextAt{9}{468.216}{146.16}{\Character{40}}%
\TextAt{13}{472.768}{146.16}{X}%
\TextAt{9}{483.423}{146.16}{\Character{41}}%
\NativeMathAt{491.634}{146.16}{11.955}{\in}%
\TextAt{14}{503.264}{146.16}{o}%
\TextAt{9}{509.108}{146.16}{\Character{91}\Character{91}}%
\TextAt{13}{515.611}{146.16}{X}%
\TextAt{9}{526.266}{146.16}{\Character{93}\Character{93}}%
\TextAt{8}{79.226}{129.114}{such}%
\TextAt{8}{105.623}{129.114}{that}%
\TextAt{13}{130.992}{129.114}{\Character{12}}%
\TextAt{21}{137.601}{127.32}{i}%
\TextAt{9}{144.303}{129.114}{\Character{61}}%
\TextAt{13}{156.728}{129.114}{h}%
\TextAt{17}{163.467}{127.32}{3}%
\TextAt{9}{168.199}{129.114}{\Character{40}}%
\TextAt{13}{172.751}{129.114}{\Character{21}}%
\TextAt{21}{179.58}{127.32}{i}%
\TextAt{9}{182.961}{129.114}{\Character{41}}%
\TextAt{8}{187.513}{129.114}{,}%
\TextAt{9}{194.663}{129.114}{0}%
\NativeMathAt{203.836}{129.114}{11.955}{\leq}%
\TextAt{13}{216.456}{129.114}{i}%
\NativeMathAt{223.77}{129.114}{11.955}{\leq}%
\TextAt{13}{236.389}{129.114}{n}%
\TextAt{8}{243.377}{129.114}{.}%
\TextAt{8}{251.825}{129.114}{Show}%
\TextAt{8}{282.717}{129.114}{that}%
\TextAt{13}{308.086}{129.114}{h}%
\TextAt{9}{318.146}{129.114}{\Character{61}}%
\TextAt{13}{330.571}{129.114}{h}%
\TextAt{17}{337.31}{127.32}{2}%
\TextAt{9}{344.699}{129.114}{\Character{43}}%
\TextAt{13}{356.46}{129.114}{h}%
\TextAt{17}{363.199}{127.32}{3}%
\TextAt{8}{371.828}{129.114}{works.}%
\TextAt{12}{79.226}{100.113}{Remark.}%
\TextAt{8}{133.383}{100.113}{The}%
\TextAt{8}{157.437}{100.113}{solutions}%
\TextAt{8}{206.345}{100.113}{of}%
\TextAt{8}{219.675}{100.113}{these}%
\TextAt{8}{249.647}{100.113}{exercises}%
\TextAt{8}{297.563}{100.113}{are}%
\TextAt{8}{317.074}{100.113}{discussed}%
\TextAt{8}{367.979}{100.113}{in}%
\TextAt{8}{381.632}{100.113}{detail}%
\TextAt{8}{414.15}{100.113}{in}%
\TextAt{8}{427.803}{100.113}{[79],}%
\TextAt{8}{453.172}{100.113}{5.2.}%
\TextAt{8}{297.211}{43.033}{322}%
\Destination{79.23}{670.535}{nav.Item.121}%
\Destination{79.23}{652.935}{nav.Item.122}%
\Destination{79.23}{635.895}{nav.Item.123}%
\Destination{79.23}{618.845}{nav.Item.124}%
\Destination{79.23}{601.805}{nav.Item.125}%
\Destination{79.23}{529.065}{nav.Item.126}%
\Destination{78.23}{809.815}{nav.page.322}%
\Destination{79.226}{775.155}{nav.ex.5.5.3}%
\Destination{79.226}{462.555}{nav.ex.5.5.4}%
\Destination{79.226}{272.255}{nav.ex.5.5.5}%
\Destination{79.36}{446.455}{nav.print.352}%
\Hotlink{526.411}{758.955}{6.858}{11.047}{nav.ex.5.5.1}%
\Hotlink{144.79}{741.855}{6.858}{11.047}{nav.ex.5.5.2}%
\Hotlink{328.389}{143.955}{6.858}{11.047}{nav.ex.5.5.4}%
\Hotlink{427.303}{97.955}{19.22}{11.047}{nav.cite.N079}%
}

% Page 335 (body page 323)
\Page{%
\TextAt{2}{79.226}{796.943}{V.}%
\TextAt{2}{94.073}{796.943}{Local}%
\TextAt{2}{123.615}{796.943}{Class}%
\TextAt{2}{152.208}{796.943}{Field}%
\TextAt{2}{179.939}{796.943}{Theory}%
\Vector{q 0 G 0 g 0.299 w 7.226 18.729 m 460.769 18.729 l S Q}%
\TextAt{15}{79.226}{760.804}{§}%
\TextAt{15}{93.785}{760.804}{6}%
\TextAt{15}{120.693}{760.804}{Higher}%
\TextAt{15}{184.889}{760.804}{Ramification}%
\TextAt{15}{298.907}{760.804}{Groups}%
\TextAt{8}{79.226}{728.918}{Considering}%
\TextAt{8}{143.676}{728.918}{the}%
\TextAt{8}{163.833}{728.918}{homomorphism}%
\TextAt{9}{233.288}{701.162}{\Character{40}}%
\TextAt{13}{241.737}{701.162}{\Character{59}}%
\TextAt{13}{246.961}{701.162}{L}%
\TextAt{16}{258.267}{701.162}{\Character{106}}%
\TextAt{13}{264.909}{701.162}{K}%
\TextAt{9}{275.726}{701.162}{\Character{41}}%
\TextAt{9}{283.581}{701.162}{\Character{58}}%
\TextAt{13}{290.172}{701.162}{K}%
\TextAt{3}{300.99}{706.098}{\Character{3}}%
\NativeArrowAt{309.043}{701.162}{11.955}{19.26}{\longrightarrow}%
\TextAt{13}{331.625}{701.162}{G}%
\TextAt{9}{340.858}{701.162}{\Character{40}}%
\TextAt{13}{345.411}{701.162}{L}%
\TextAt{16}{356.696}{701.162}{\Character{106}}%
\TextAt{13}{363.338}{701.162}{K}%
\TextAt{9}{374.156}{701.162}{\Character{41}}%
\TextAt{8}{79.226}{673.405}{defined}%
\TextAt{8}{120.262}{673.405}{for}%
\TextAt{8}{138.87}{673.405}{an}%
\TextAt{8}{155.852}{673.405}{abelian}%
\TextAt{8}{197.223}{673.405}{extension}%
\TextAt{13}{249.712}{673.405}{L}%
\TextAt{16}{262.23}{673.405}{\Character{106}}%
\TextAt{13}{270.104}{673.405}{K}%
\TextAt{8}{285.542}{673.405}{of}%
\TextAt{8}{299.595}{673.405}{local}%
\TextAt{8}{327.959}{673.405}{fields}%
\TextAt{8}{358.654}{673.405}{by}%
\TextAt{8}{375.636}{673.405}{the}%
\TextAt{8}{396.515}{673.405}{norm}%
\TextAt{8}{427.808}{673.405}{residue}%
\TextAt{8}{468.258}{673.405}{symbol,}%
\TextAt{8}{512.476}{673.405}{it}%
\TextAt{8}{524.903}{673.405}{is}%
\TextAt{8}{79.226}{656.359}{striking}%
\TextAt{8}{122.598}{656.359}{that}%
\TextAt{8}{148.672}{656.359}{both}%
\TextAt{8}{177.016}{656.359}{groups}%
\TextAt{8}{215.511}{656.359}{are}%
\TextAt{8}{235.726}{656.359}{equipped}%
\TextAt{8}{286.498}{656.359}{with}%
\TextAt{8}{313.863}{656.359}{a}%
\TextAt{8}{324.322}{656.359}{canonical}%
\TextAt{8}{376.41}{656.359}{filtration:}%
\TextAt{8}{431.159}{656.359}{in}%
\TextAt{8}{445.516}{656.359}{the}%
\TextAt{8}{466.376}{656.359}{group}%
\TextAt{13}{500.256}{656.359}{K}%
\TextAt{3}{511.074}{660.698}{\Character{3}}%
\TextAt{8}{520.408}{656.359}{on}%
\TextAt{8}{79.226}{639.313}{the}%
\TextAt{8}{99.383}{639.313}{left}%
\TextAt{8}{119.862}{639.313}{we}%
\TextAt{8}{137.089}{639.313}{have}%
\TextAt{8}{164.084}{639.313}{the}%
\TextAt{8}{184.24}{639.313}{descending}%
\TextAt{8}{243.478}{639.313}{chain}%
\TextAt{13}{212.642}{611.557}{K}%
\TextAt{3}{223.46}{616.493}{\Character{3}}%
\NativeMathAt{231.513}{611.557}{11.955}{\supseteq}%
\TextAt{13}{244.132}{611.557}{U}%
\TextAt{21}{252.083}{609.763}{K}%
\TextAt{9}{263.606}{611.557}{\Character{61}}%
\TextAt{13}{276.032}{611.557}{U}%
\TextAt{17}{285.231}{617.667}{\Character{40}0\Character{41}}%
\TextAt{21}{283.982}{608.316}{K}%
\NativeMathAt{299.87}{611.557}{11.955}{\supseteq}%
\TextAt{13}{312.49}{611.557}{U}%
\TextAt{17}{321.689}{617.667}{\Character{40}1\Character{41}}%
\TextAt{21}{320.44}{608.316}{K}%
\NativeMathAt{336.328}{611.557}{11.955}{\supseteq}%
\TextAt{13}{348.948}{611.557}{U}%
\TextAt{17}{358.147}{617.667}{\Character{40}2\Character{41}}%
\TextAt{21}{356.898}{608.316}{K}%
\NativeMathAt{372.786}{611.557}{11.955}{\supseteq}%
\TextAt{16}{385.387}{611.557}{\Character{1}}%
\TextAt{16}{390.695}{611.557}{\Character{1}}%
\TextAt{16}{396.003}{611.557}{\Character{1}}%
\TextAt{8}{517.802}{611.557}{(*)}%
\TextAt{8}{79.226}{583.8}{of}%
\TextAt{12}{93.07}{583.8}{higher}%
\TextAt{12}{134.602}{583.8}{unit}%
\TextAt{12}{163.008}{583.8}{groups}%
\TextAt{13}{206.118}{583.8}{U}%
\TextAt{17}{215.317}{589.911}{\Character{40}}%
\TextAt{21}{218.611}{589.911}{n}%
\TextAt{17}{223.749}{589.911}{\Character{41}}%
\TextAt{21}{214.069}{580.559}{K}%
\TextAt{8}{227.54}{583.8}{,}%
\TextAt{8}{234.689}{583.8}{and}%
\TextAt{8}{257.452}{583.8}{on}%
\TextAt{8}{273.711}{583.8}{the}%
\TextAt{8}{293.868}{583.8}{right}%
\TextAt{8}{322.165}{583.8}{there}%
\TextAt{8}{352.077}{583.8}{is}%
\TextAt{8}{363.841}{583.8}{the}%
\TextAt{8}{383.998}{583.8}{descending}%
\TextAt{8}{443.235}{583.8}{chain}%
\TextAt{13}{165.879}{556.044}{G}%
\TextAt{9}{175.113}{556.044}{\Character{40}}%
\TextAt{13}{179.665}{556.044}{L}%
\TextAt{16}{190.95}{556.044}{\Character{106}}%
\TextAt{13}{197.592}{556.044}{K}%
\TextAt{9}{208.41}{556.044}{\Character{41}}%
\NativeMathAt{216.283}{556.044}{11.955}{\supseteq}%
\TextAt{13}{228.902}{556.044}{G}%
\TextAt{17}{238.136}{560.98}{0}%
\TextAt{9}{242.868}{556.044}{\Character{40}}%
\TextAt{13}{247.421}{556.044}{L}%
\TextAt{16}{258.706}{556.044}{\Character{106}}%
\TextAt{13}{265.348}{556.044}{K}%
\TextAt{9}{276.165}{556.044}{\Character{41}}%
\NativeMathAt{284.039}{556.044}{11.955}{\supseteq}%
\TextAt{13}{296.658}{556.044}{G}%
\TextAt{17}{305.891}{560.98}{1}%
\TextAt{9}{310.624}{556.044}{\Character{40}}%
\TextAt{13}{315.176}{556.044}{L}%
\TextAt{16}{326.461}{556.044}{\Character{106}}%
\TextAt{13}{333.103}{556.044}{K}%
\TextAt{9}{343.921}{556.044}{\Character{41}}%
\NativeMathAt{351.794}{556.044}{11.955}{\supseteq}%
\TextAt{13}{364.413}{556.044}{G}%
\TextAt{17}{373.647}{560.98}{2}%
\TextAt{9}{378.379}{556.044}{\Character{40}}%
\TextAt{13}{382.932}{556.044}{L}%
\TextAt{16}{394.217}{556.044}{\Character{106}}%
\TextAt{13}{400.859}{556.044}{K}%
\TextAt{9}{411.676}{556.044}{\Character{41}}%
\NativeMathAt{419.55}{556.044}{11.955}{\supseteq}%
\TextAt{16}{432.151}{556.044}{\Character{1}}%
\TextAt{16}{437.459}{556.044}{\Character{1}}%
\TextAt{16}{442.767}{556.044}{\Character{1}}%
\TextAt{8}{511.944}{556.044}{(**)}%
\TextAt{8}{79.226}{528.287}{of}%
\TextAt{12}{94.09}{528.287}{ramification}%
\TextAt{12}{169.911}{528.287}{groups}%
\TextAt{13}{214.042}{528.287}{G}%
\TextAt{21}{223.275}{532.625}{i}%
\TextAt{9}{226.657}{528.287}{\Character{40}}%
\TextAt{13}{231.209}{528.287}{L}%
\TextAt{16}{244.232}{528.287}{\Character{106}}%
\TextAt{13}{252.612}{528.287}{K}%
\TextAt{9}{263.43}{528.287}{\Character{41}}%
\TextAt{8}{272.899}{528.287}{in}%
\TextAt{8}{287.572}{528.287}{the}%
\TextAt{8}{308.748}{528.287}{upper}%
\TextAt{8}{343.255}{528.287}{numbering}%
\TextAt{8}{402.484}{528.287}{(see}%
\TextAt{8}{426.972}{528.287}{chap.}%
\TextAt{8}{462.223}{528.287}{II,}%
\TextAt{8}{479.155}{528.287}{§10).}%
\TextAt{8}{512.613}{528.287}{The}%
\TextAt{8}{79.226}{511.241}{latter}%
\TextAt{8}{111.099}{511.241}{arose}%
\TextAt{8}{141.082}{511.241}{from}%
\TextAt{8}{168.723}{511.241}{the}%
\TextAt{8}{188.879}{511.241}{ramification}%
\TextAt{8}{253.927}{511.241}{groups}%
\TextAt{8}{291.717}{511.241}{in}%
\TextAt{8}{305.37}{511.241}{the}%
\TextAt{8}{325.527}{511.241}{lower}%
\TextAt{8}{356.096}{511.241}{numbering}%
\TextAt{13}{145.22}{483.484}{G}%
\TextAt{21}{154.454}{481.691}{i}%
\TextAt{9}{157.835}{483.484}{\Character{40}}%
\TextAt{13}{162.387}{483.484}{L}%
\TextAt{16}{173.672}{483.484}{\Character{106}}%
\TextAt{13}{180.314}{483.484}{K}%
\TextAt{9}{191.132}{483.484}{\Character{41}}%
\TextAt{9}{198.987}{483.484}{\Character{61}}%
\TextAt{16}{211.431}{483.484}{\Character{102}}%
\TextAt{13}{217.408}{483.484}{\Character{27}}%
\NativeMathAt{227.811}{483.484}{11.955}{\in}%
\TextAt{13}{239.102}{483.484}{G}%
\TextAt{9}{248.336}{483.484}{\Character{40}}%
\TextAt{13}{252.888}{483.484}{L}%
\TextAt{16}{264.174}{483.484}{\Character{106}}%
\TextAt{13}{270.815}{483.484}{K}%
\TextAt{9}{281.633}{483.484}{\Character{41}}%
\TextAt{16}{289.506}{483.484}{\Character{106}}%
\TextAt{13}{296.148}{483.484}{v}%
\TextAt{21}{301.807}{481.691}{L}%
\TextAt{9}{308.064}{483.484}{\Character{40}}%
\TextAt{13}{312.617}{483.484}{a}%
\TextAt{21}{318.762}{488.421}{\Character{27}}%
\TextAt{16}{327.042}{483.484}{\Character{0}}%
\TextAt{13}{338.997}{483.484}{a}%
\TextAt{9}{345.142}{483.484}{\Character{41}}%
\NativeMathAt{353.015}{483.484}{11.955}{\geq}%
\TextAt{13}{365.634}{483.484}{i}%
\TextAt{9}{372.284}{483.484}{\Character{43}}%
\TextAt{9}{384.036}{483.484}{1}%
\TextAt{8}{393.796}{483.484}{for}%
\TextAt{8}{411.681}{483.484}{all}%
\TextAt{13}{427.94}{483.484}{a}%
\NativeMathAt{437.406}{483.484}{11.955}{\in}%
\TextAt{14}{448.697}{483.484}{o}%
\TextAt{21}{454.541}{481.691}{L}%
\TextAt{16}{460.798}{483.484}{\Character{103}}%
\TextAt{8}{79.226}{455.728}{via}%
\TextAt{8}{98.414}{455.728}{the}%
\TextAt{8}{118.571}{455.728}{strictly}%
\TextAt{8}{158.632}{455.728}{increasing}%
\TextAt{8}{213.327}{455.728}{function}%
\TextAt{13}{245.663}{425.021}{\Character{17}}%
\TextAt{21}{251.438}{423.161}{L}%
\TextAt{3}{257.197}{423.161}{\Character{106}}%
\TextAt{21}{259.55}{423.161}{K}%
\TextAt{9}{267.752}{425.021}{\Character{40}}%
\TextAt{13}{272.305}{425.021}{s}%
\TextAt{9}{277.819}{425.021}{\Character{41}}%
\TextAt{9}{285.674}{425.021}{\Character{61}}%
\TextAt{20}{298.117}{441.294}{\Character{90}}%
\TextAt{21}{310.073}{438.139}{s}%
\TextAt{17}{304.759}{414.228}{0}%
\TextAt{8}{334.828}{433.109}{d}%
\TextAt{13}{341.332}{433.109}{x}%
\Vector{q 0 G 0 g 0.478 w 245.675 -345.99 m 293.137 -345.99 l S Q}%
\TextAt{9}{317.675}{416.82}{\Character{40}}%
\TextAt{13}{322.227}{416.82}{G}%
\TextAt{17}{331.461}{415.027}{0}%
\TextAt{9}{339.514}{416.82}{\Character{58}}%
\TextAt{13}{346.086}{416.82}{G}%
\TextAt{21}{355.32}{415.027}{x}%
\TextAt{9}{360.585}{416.82}{\Character{41}}%
\TextAt{8}{79.226}{392.517}{by}%
\TextAt{8}{94.755}{392.517}{the}%
\TextAt{8}{114.181}{392.517}{rule}%
\TextAt{13}{136.858}{392.517}{G}%
\TextAt{21}{146.092}{396.856}{i}%
\TextAt{9}{149.473}{392.517}{\Character{40}}%
\TextAt{13}{154.026}{392.517}{L}%
\TextAt{16}{165.311}{392.517}{\Character{106}}%
\TextAt{13}{171.953}{392.517}{K}%
\TextAt{9}{182.77}{392.517}{\Character{41}}%
\TextAt{9}{190.625}{392.517}{\Character{61}}%
\TextAt{13}{203.069}{392.517}{G}%
\TextAt{21}{212.303}{390.658}{\Character{32}}%
\TextAt{22}{217.838}{388.92}{L}%
\TextAt{23}{222.896}{388.92}{\Character{106}}%
\TextAt{22}{225.166}{388.92}{K}%
\TextAt{17}{232.209}{390.658}{\Character{40}}%
\TextAt{21}{235.502}{390.658}{i}%
\TextAt{17}{238.385}{390.658}{\Character{41}}%
\TextAt{9}{242.176}{392.517}{\Character{40}}%
\TextAt{13}{246.729}{392.517}{L}%
\TextAt{16}{258.014}{392.517}{\Character{106}}%
\TextAt{13}{264.656}{392.517}{K}%
\TextAt{9}{275.473}{392.517}{\Character{41}}%
\TextAt{8}{280.026}{392.517}{,}%
\TextAt{8}{286.591}{392.517}{where}%
\TextAt{13}{319.669}{392.517}{\Character{32}}%
\TextAt{8}{330.876}{392.517}{is}%
\TextAt{8}{341.909}{392.517}{the}%
\TextAt{8}{361.335}{392.517}{inverse}%
\TextAt{8}{399.363}{392.517}{function}%
\TextAt{8}{444.481}{392.517}{of}%
\TextAt{13}{457.08}{392.517}{\Character{17}}%
\TextAt{8}{463.285}{392.517}{.}%
\TextAt{8}{471.49}{392.517}{We}%
\TextAt{8}{490.904}{392.517}{will}%
\TextAt{8}{512.278}{392.517}{now}%
\TextAt{8}{79.226}{375.471}{prove}%
\TextAt{8}{110.292}{375.471}{the}%
\TextAt{8}{129.964}{375.471}{remarkable}%
\TextAt{8}{190.295}{375.471}{arithmetic}%
\TextAt{8}{246.395}{375.471}{fact}%
\TextAt{8}{268.996}{375.471}{that}%
\TextAt{8}{293.881}{375.471}{the}%
\TextAt{8}{313.553}{375.471}{norm}%
\TextAt{8}{343.638}{375.471}{residue}%
\TextAt{8}{382.881}{375.471}{symbol}%
\TextAt{9}{422.458}{375.471}{\Character{40}}%
\TextAt{13}{430.424}{375.471}{\Character{59}}%
\TextAt{13}{435.648}{375.471}{L}%
\TextAt{16}{446.953}{375.471}{\Character{106}}%
\TextAt{13}{453.595}{375.471}{K}%
\TextAt{9}{464.413}{375.471}{\Character{41}}%
\TextAt{8}{472.378}{375.471}{relates}%
\TextAt{8}{509.027}{375.471}{both}%
\TextAt{8}{79.226}{358.425}{filtrations}%
\TextAt{8}{132.446}{358.425}{(*)}%
\TextAt{8}{151.127}{358.425}{and}%
\TextAt{8}{173.706}{358.425}{(**)}%
\TextAt{8}{198.245}{358.425}{in}%
\TextAt{8}{211.714}{358.425}{a}%
\TextAt{8}{221.285}{358.425}{precise}%
\TextAt{8}{259.525}{358.425}{way.}%
\TextAt{8}{286.778}{358.425}{To}%
\TextAt{8}{303.821}{358.425}{this}%
\TextAt{8}{326.459}{358.425}{end}%
\TextAt{8}{348.38}{358.425}{we}%
\TextAt{8}{365.424}{358.425}{determine}%
\TextAt{8}{419.863}{358.425}{(generalizing}%
\TextAt{8}{488.625}{358.425}{chap.}%
\TextAt{8}{520.755}{358.425}{II,}%
\TextAt{8}{79.226}{341.379}{§10,}%
\TextAt{8}{103.77}{341.379}{exercise}%
\TextAt{8}{147.072}{341.379}{1)}%
\TextAt{8}{161.382}{341.379}{the}%
\TextAt{8}{181.538}{341.379}{higher}%
\TextAt{8}{217.308}{341.379}{ramification}%
\TextAt{8}{282.356}{341.379}{groups}%
\TextAt{8}{320.147}{341.379}{of}%
\TextAt{8}{333.477}{341.379}{the}%
\TextAt{8}{353.633}{341.379}{Lubin–Tate}%
\TextAt{8}{416.553}{341.379}{extensions.}%
\TextAt{12}{79.226}{315.723}{Proposition}%
\TextAt{12}{151.325}{315.723}{(6.1).}%
\TextAt{8}{186.108}{315.723}{Let}%
\TextAt{13}{206.545}{315.723}{L}%
\TextAt{21}{214.51}{313.929}{n}%
\TextAt{16}{223.467}{315.723}{\Character{106}}%
\TextAt{13}{230.109}{315.723}{K}%
\TextAt{8}{244.291}{315.723}{be}%
\TextAt{8}{259.683}{315.723}{the}%
\TextAt{8}{279.307}{315.723}{field}%
\TextAt{8}{304.132}{315.723}{of}%
\TextAt{13}{316.929}{315.723}{\Character{25}}%
\TextAt{21}{323.999}{320.061}{n}%
\TextAt{8}{329.635}{315.723}{-division}%
\TextAt{8}{376.314}{315.723}{points}%
\TextAt{8}{410.965}{315.723}{of}%
\TextAt{8}{423.763}{315.723}{a}%
\TextAt{8}{432.986}{315.723}{Lubin–Tate}%
\TextAt{8}{495.374}{315.723}{module}%
\ErrataAt{79.19}{298.677}{455.81}{\emph{for the prime element $\pi$. For integers $i$ and $1\le k\le n$ one has}}%
\TextAt{13}{175.739}{270.92}{G}%
\TextAt{21}{184.973}{269.127}{i}%
\TextAt{9}{188.354}{270.92}{\Character{40}}%
\TextAt{13}{192.907}{270.92}{L}%
\TextAt{21}{200.871}{269.127}{n}%
\TextAt{16}{209.828}{270.92}{\Character{106}}%
\TextAt{13}{216.47}{270.92}{K}%
\TextAt{9}{227.288}{270.92}{\Character{41}}%
\TextAt{9}{235.143}{270.92}{\Character{61}}%
\TextAt{13}{247.586}{270.92}{G}%
\TextAt{9}{256.82}{270.92}{\Character{40}}%
\TextAt{13}{261.372}{270.92}{L}%
\TextAt{21}{269.337}{269.127}{n}%
\TextAt{16}{278.294}{270.92}{\Character{106}}%
\TextAt{13}{284.936}{270.92}{L}%
\TextAt{21}{292.9}{269.127}{k}%
\TextAt{9}{298.02}{270.92}{\Character{41}}%
\TextAt{8}{326.483}{270.92}{for}%
\TextAt{13}{344.368}{270.92}{q}%
\TextAt{21}{349.987}{275.856}{k}%
\TextAt{3}{354.608}{275.856}{\Character{0}}%
\TextAt{17}{361.195}{275.856}{1}%
\NativeMathAt{369.248}{270.92}{11.955}{\leq}%
\TextAt{13}{381.867}{270.92}{i}%
\NativeMathAt{389.182}{270.92}{11.955}{\leq}%
\TextAt{13}{401.801}{270.92}{q}%
\TextAt{21}{407.42}{275.856}{k}%
\TextAt{16}{415.196}{270.92}{\Character{0}}%
\TextAt{9}{427.152}{270.92}{1}%
\TextAt{13}{433.005}{270.92}{\Character{58}}%
\TextAt{19}{79.226}{225.942}{Proof.}%
\TextAt{8}{116}{225.942}{By}%
\TextAt{8}{134.91}{225.942}{(5.4)}%
\TextAt{8}{163.431}{225.942}{and}%
\TextAt{8}{186.74}{225.942}{(5.5),}%
\TextAt{8}{218.649}{225.942}{the}%
\TextAt{8}{239.352}{225.942}{norm}%
\TextAt{8}{270.467}{225.942}{residue}%
\TextAt{8}{310.74}{225.942}{symbol}%
\TextAt{8}{351.348}{225.942}{gives}%
\TextAt{8}{380.575}{225.942}{an}%
\TextAt{8}{397.38}{225.942}{isomorphism}%
\TextAt{13}{466.346}{225.942}{U}%
\TextAt{21}{474.296}{224.149}{K}%
\TextAt{9}{482.499}{225.942}{\Character{47}}%
\TextAt{13}{488.351}{225.942}{U}%
\TextAt{17}{497.551}{232.052}{\Character{40}}%
\TextAt{21}{500.844}{232.052}{k}%
\TextAt{17}{505.465}{232.052}{\Character{41}}%
\TextAt{21}{496.302}{222.701}{K}%
\NativeArrowAt{513.508}{225.942}{11.955}{19.26}{\longrightarrow}%
\TextAt{13}{79.226}{208.896}{G}%
\TextAt{9}{88.46}{208.896}{\Character{40}}%
\TextAt{13}{93.012}{208.896}{L}%
\TextAt{21}{100.977}{207.103}{k}%
\TextAt{16}{110.31}{208.896}{\Character{106}}%
\TextAt{13}{117.845}{208.896}{K}%
\TextAt{9}{128.662}{208.896}{\Character{41}}%
\TextAt{8}{137.636}{208.896}{for}%
\TextAt{8}{156.045}{208.896}{every}%
\TextAt{13}{187.461}{208.896}{k}%
\TextAt{8}{193.951}{208.896}{.}%
\TextAt{8}{203.971}{208.896}{Hence}%
\TextAt{13}{239.273}{208.896}{G}%
\TextAt{9}{248.506}{208.896}{\Character{40}}%
\TextAt{13}{253.059}{208.896}{L}%
\TextAt{21}{261.023}{207.103}{n}%
\TextAt{16}{270.873}{208.896}{\Character{106}}%
\TextAt{13}{278.408}{208.896}{L}%
\TextAt{21}{286.372}{207.103}{k}%
\TextAt{9}{291.492}{208.896}{\Character{41}}%
\TextAt{9}{300.243}{208.896}{\Character{61}}%
\TextAt{9}{313.549}{208.896}{\Character{40}}%
\TextAt{13}{318.129}{208.896}{U}%
\TextAt{17}{327.328}{215.006}{\Character{40}}%
\TextAt{21}{330.621}{215.006}{k}%
\TextAt{17}{335.243}{215.006}{\Character{41}}%
\TextAt{21}{326.079}{205.655}{K}%
\TextAt{13}{339.034}{208.896}{\Character{59}}%
\TextAt{13}{344.258}{208.896}{L}%
\TextAt{21}{352.243}{207.103}{n}%
\TextAt{16}{362.093}{208.896}{\Character{106}}%
\TextAt{13}{369.628}{208.896}{K}%
\TextAt{9}{380.446}{208.896}{\Character{41}}%
\TextAt{8}{384.998}{208.896}{.}%
\TextAt{8}{395.018}{208.896}{We}%
\TextAt{8}{415.687}{208.896}{therefore}%
\TextAt{8}{465.311}{208.896}{have}%
\TextAt{8}{492.829}{208.896}{to}%
\TextAt{8}{507.664}{208.896}{show}%
\TextAt{8}{79.226}{191.85}{that}%
\TextAt{13}{168.415}{174.804}{G}%
\TextAt{21}{177.648}{173.011}{i}%
\TextAt{9}{181.03}{174.804}{\Character{40}}%
\TextAt{13}{185.582}{174.804}{L}%
\TextAt{21}{193.547}{173.011}{n}%
\TextAt{16}{202.504}{174.804}{\Character{106}}%
\TextAt{13}{209.145}{174.804}{K}%
\TextAt{9}{219.963}{174.804}{\Character{41}}%
\TextAt{9}{227.818}{174.804}{\Character{61}}%
\TextAt{9}{240.239}{174.804}{\Character{40}}%
\TextAt{13}{244.814}{174.804}{U}%
\TextAt{17}{254.013}{180.915}{\Character{40}}%
\TextAt{21}{257.306}{180.915}{k}%
\TextAt{17}{261.928}{180.915}{\Character{41}}%
\TextAt{21}{252.764}{171.563}{K}%
\TextAt{13}{265.719}{174.804}{\Character{59}}%
\TextAt{13}{270.943}{174.804}{L}%
\TextAt{21}{278.928}{173.011}{n}%
\TextAt{16}{287.885}{174.804}{\Character{106}}%
\TextAt{13}{294.527}{174.804}{K}%
\TextAt{9}{305.345}{174.804}{\Character{41}}%
\TextAt{8}{333.807}{174.804}{for}%
\TextAt{13}{351.692}{174.804}{q}%
\TextAt{21}{357.311}{179.74}{k}%
\TextAt{3}{361.933}{179.74}{\Character{0}}%
\TextAt{17}{368.519}{179.74}{1}%
\NativeMathAt{376.573}{174.804}{11.955}{\leq}%
\TextAt{13}{389.192}{174.804}{i}%
\NativeMathAt{396.506}{174.804}{11.955}{\leq}%
\TextAt{13}{409.126}{174.804}{q}%
\TextAt{21}{414.745}{179.74}{k}%
\TextAt{16}{422.521}{174.804}{\Character{0}}%
\TextAt{9}{434.476}{174.804}{1}%
\TextAt{13}{440.329}{174.804}{\Character{58}}%
\TextAt{8}{79.226}{151.816}{Let}%
\TextAt{13}{100.573}{151.816}{\Character{27}}%
\NativeMathAt{111.619}{151.816}{11.955}{\in}%
\TextAt{13}{123.553}{151.816}{G}%
\TextAt{17}{132.786}{150.023}{1}%
\TextAt{9}{137.518}{151.816}{\Character{40}}%
\TextAt{13}{142.071}{151.816}{L}%
\TextAt{21}{150.035}{150.023}{n}%
\TextAt{16}{159.635}{151.816}{\Character{106}}%
\TextAt{13}{166.92}{151.816}{K}%
\TextAt{9}{177.737}{151.816}{\Character{41}}%
\TextAt{8}{186.564}{151.816}{and}%
\TextAt{13}{209.704}{151.816}{\Character{27}}%
\TextAt{9}{220.75}{151.816}{\Character{61}}%
\TextAt{9}{233.805}{151.816}{\Character{40}}%
\TextAt{13}{238.371}{151.816}{u}%
\TextAt{3}{245.033}{156.154}{\Character{0}}%
\TextAt{17}{251.62}{156.154}{1}%
\TextAt{13}{256.352}{151.816}{\Character{59}}%
\TextAt{13}{261.576}{151.816}{L}%
\TextAt{21}{269.561}{150.023}{n}%
\TextAt{16}{279.161}{151.816}{\Character{106}}%
\TextAt{13}{286.445}{151.816}{K}%
\TextAt{9}{297.263}{151.816}{\Character{41}}%
\TextAt{8}{301.815}{151.816}{.}%
\TextAt{8}{311.395}{151.816}{Then}%
\TextAt{8}{342.33}{151.816}{we}%
\TextAt{8}{359.934}{151.816}{have}%
\TextAt{8}{387.306}{151.816}{necessarily}%
\TextAt{13}{446.013}{151.816}{u}%
\NativeMathAt{456.639}{151.816}{11.955}{\in}%
\TextAt{13}{468.573}{151.816}{U}%
\TextAt{17}{477.772}{157.926}{\Character{40}1\Character{41}}%
\TextAt{21}{476.523}{148.575}{K}%
\TextAt{8}{493.365}{151.816}{because}%
\TextAt{9}{79.226}{134.77}{\Character{40}}%
\TextAt{13}{88.859}{134.77}{\Character{59}}%
\TextAt{13}{94.083}{134.77}{L}%
\TextAt{21}{102.068}{132.977}{n}%
\TextAt{16}{113.042}{134.77}{\Character{106}}%
\TextAt{13}{121.7}{134.77}{K}%
\TextAt{9}{132.518}{134.77}{\Character{41}}%
\TextAt{9}{142.393}{134.77}{\Character{58}}%
\TextAt{13}{150.997}{134.77}{U}%
\TextAt{21}{158.947}{132.977}{K}%
\TextAt{9}{167.15}{134.77}{\Character{47}}%
\TextAt{13}{173.003}{134.77}{U}%
\TextAt{17}{182.202}{140.88}{\Character{40}}%
\TextAt{21}{185.495}{140.88}{n}%
\TextAt{17}{190.633}{140.88}{\Character{41}}%
\TextAt{21}{180.953}{131.529}{K}%
\NativeMathAt{206.099}{141.547}{7.97}{\sim}%
\NativeArrowAt{199.762}{134.77}{11.955}{19.26}{\longrightarrow}%
\TextAt{13}{224.361}{134.77}{G}%
\TextAt{9}{233.594}{134.77}{\Character{40}}%
\TextAt{13}{238.147}{134.77}{L}%
\TextAt{21}{246.111}{132.977}{n}%
\TextAt{16}{257.085}{134.77}{\Character{106}}%
\TextAt{13}{265.743}{134.77}{K}%
\TextAt{9}{276.561}{134.77}{\Character{41}}%
\TextAt{8}{286.194}{134.77}{maps}%
\TextAt{8}{318.007}{134.77}{the}%
\TextAt{13}{339.346}{134.77}{p}%
\TextAt{8}{345.222}{134.77}{-Sylow}%
\TextAt{8}{384.123}{134.77}{subgroup}%
\TextAt{13}{436.104}{134.77}{U}%
\TextAt{17}{445.303}{140.88}{\Character{40}1\Character{41}}%
\TextAt{21}{444.055}{131.529}{K}%
\TextAt{9}{456.622}{134.77}{\Character{47}}%
\TextAt{13}{462.475}{134.77}{U}%
\TextAt{17}{471.674}{140.88}{\Character{40}}%
\TextAt{21}{474.968}{140.88}{n}%
\TextAt{17}{480.106}{140.88}{\Character{41}}%
\TextAt{21}{470.426}{131.529}{K}%
\TextAt{8}{488.978}{134.77}{onto}%
\TextAt{8}{516.51}{134.77}{the}%
\TextAt{13}{79.226}{117.724}{p}%
\TextAt{8}{85.101}{117.724}{-Sylow}%
\TextAt{8}{124.511}{117.724}{subgroup}%
\TextAt{13}{177}{117.724}{G}%
\TextAt{17}{186.234}{115.931}{1}%
\TextAt{9}{190.966}{117.724}{\Character{40}}%
\TextAt{13}{195.518}{117.724}{L}%
\TextAt{21}{203.483}{115.931}{n}%
\TextAt{16}{215.322}{117.724}{\Character{106}}%
\TextAt{13}{224.846}{117.724}{K}%
\TextAt{9}{235.664}{117.724}{\Character{41}}%
\TextAt{8}{245.805}{117.724}{of}%
\TextAt{13}{260.826}{117.724}{G}%
\TextAt{9}{270.06}{117.724}{\Character{40}}%
\TextAt{13}{274.612}{117.724}{L}%
\TextAt{21}{282.577}{115.931}{n}%
\TextAt{16}{294.416}{117.724}{\Character{106}}%
\TextAt{13}{303.94}{117.724}{K}%
\TextAt{9}{314.758}{117.724}{\Character{41}}%
\TextAt{8}{319.31}{117.724}{.}%
\TextAt{8}{332.833}{117.724}{Let}%
\TextAt{13}{355.494}{117.724}{u}%
\TextAt{9}{368.359}{117.724}{\Character{61}}%
\TextAt{9}{383.65}{117.724}{1}%
\TextAt{9}{393.31}{117.724}{\Character{43}}%
\TextAt{13}{406.244}{117.724}{\Character{34}\Character{25}}%
\TextAt{21}{418.791}{122.062}{m}%
\TextAt{8}{426.78}{117.724}{,}%
\TextAt{13}{436.043}{117.724}{\Character{34}}%
\NativeMathAt{447.724}{117.724}{11.955}{\in}%
\TextAt{13}{461.897}{117.724}{U}%
\TextAt{21}{469.848}{115.931}{K}%
\TextAt{8}{478.05}{117.724}{,}%
\TextAt{8}{487.314}{117.724}{and}%
\TextAt{13}{511.768}{117.724}{\Character{21}}%
\NativeMathAt{524.799}{117.724}{11.955}{\in}%
\TextAt{13}{79.226}{100.678}{F}%
\TextAt{9}{88.43}{100.678}{\Character{40}}%
\TextAt{13}{92.982}{100.678}{n}%
\TextAt{9}{99.97}{100.678}{\Character{41}}%
\TextAt{16}{107.179}{100.678}{\Character{110}}%
\TextAt{13}{115.813}{100.678}{F}%
\TextAt{9}{125.017}{100.678}{\Character{40}}%
\TextAt{13}{129.569}{100.678}{n}%
\TextAt{16}{139.213}{100.678}{\Character{0}}%
\TextAt{9}{151.168}{100.678}{1\Character{41}}%
\TextAt{8}{161.574}{100.678}{.}%
\TextAt{8}{170.022}{100.678}{Then}%
\TextAt{13}{200.579}{100.678}{\Character{21}}%
\TextAt{8}{211.305}{100.678}{is}%
\TextAt{8}{223.069}{100.678}{a}%
\TextAt{8}{232.825}{100.678}{prime}%
\TextAt{8}{265.988}{100.678}{element}%
\TextAt{8}{309.23}{100.678}{of}%
\TextAt{13}{322.56}{100.678}{L}%
\TextAt{21}{330.525}{98.885}{n}%
\TextAt{8}{340.058}{100.678}{and}%
\TextAt{8}{362.821}{100.678}{from}%
\TextAt{8}{390.461}{100.678}{(5.4)}%
\TextAt{8}{418.436}{100.678}{we}%
\TextAt{8}{435.664}{100.678}{get}%
\TextAt{8}{455.175}{100.678}{that}%
\TextAt{13}{224.14}{72.921}{\Character{21}}%
\TextAt{21}{230.969}{77.858}{\Character{27}}%
\TextAt{9}{239.913}{72.921}{\Character{61}}%
\TextAt{9}{252.323}{72.921}{\Character{91}}%
\TextAt{13}{255.59}{72.921}{u}%
\TextAt{9}{262.253}{72.921}{\Character{93}}%
\TextAt{21}{265.504}{71.128}{F}%
\TextAt{9}{272.538}{72.921}{\Character{40}}%
\TextAt{13}{277.09}{72.921}{\Character{21}}%
\TextAt{9}{283.918}{72.921}{\Character{41}}%
\TextAt{9}{291.773}{72.921}{\Character{61}}%
\TextAt{13}{304.217}{72.921}{F}%
\TextAt{9}{313.421}{72.921}{\Character{40}}%
\TextAt{13}{317.973}{72.921}{\Character{21}\Character{59}}%
\TextAt{9}{330.046}{72.921}{\Character{91}}%
\TextAt{13}{333.297}{72.921}{\Character{34}\Character{25}}%
\TextAt{21}{345.845}{77.858}{m}%
\TextAt{9}{353.833}{72.921}{\Character{93}}%
\TextAt{21}{357.085}{71.128}{F}%
\TextAt{9}{364.118}{72.921}{\Character{40}}%
\TextAt{13}{368.671}{72.921}{\Character{21}}%
\TextAt{9}{375.499}{72.921}{\Character{41}\Character{41}}%
\TextAt{13}{384.604}{72.921}{\Character{58}}%
\TextAt{8}{297.211}{43.033}{323}%
\Destination{522.36}{628.605}{nav.AMS.74}%
\Destination{516.5}{573.085}{nav.AMS.75}%
\Destination{78.23}{809.815}{nav.page.323}%
\Destination{79.23}{773.075}{nav.section.5.6}%
\Destination{79.36}{407.555}{nav.print.353}%
\Destination{79.36}{329.755}{nav.number.5.6.1}%
\Destination{79.23}{773.075}{outline.54}%
\Hotlink{461.723}{526.055}{9.763}{11.047}{nav.chapter.2}%
\Hotlink{484.334}{526.055}{12.716}{11.047}{nav.section.2.10}%
\Hotlink{520.255}{356.255}{9.751}{11.047}{nav.chapter.2}%
\Hotlink{134.41}{223.755}{25.078}{11.047}{nav.number.5.5.4}%
\Hotlink{186.24}{223.755}{25.078}{11.047}{nav.number.5.5.5}%
\Hotlink{389.961}{98.455}{25.078}{11.047}{nav.number.5.5.4}%
\Destination{79.226}{316}{errata.197}%
\Hotlink{132.446}{356.225}{16}{11.047}{nav.AMS.74}%
\Hotlink{173.706}{356.225}{22}{11.047}{nav.AMS.75}%
\Hotlink{79.226}{339.179}{19.8}{11.047}{nav.section.2.10}%
\Hotlink{147.072}{339.179}{6.5}{11.047}{nav.ex.2.10.1}%
}

% Page 336 (body page 324)
\Page{%
\TextAt{2}{79.226}{796.943}{V.}%
\TextAt{2}{94.073}{796.943}{Local}%
\TextAt{2}{123.615}{796.943}{Class}%
\TextAt{2}{152.208}{796.943}{Field}%
\TextAt{2}{179.939}{796.943}{Theory}%
\Vector{q 0 G 0 g 0.299 w 7.226 18.729 m 460.769 18.729 l S Q}%
\TextAt{8}{79.226}{761.124}{If}%
\TextAt{13}{92.133}{761.124}{m}%
\NativeMathAt{107.764}{761.124}{11.955}{\geq}%
\TextAt{13}{122.454}{761.124}{n}%
\TextAt{8}{129.441}{761.124}{,}%
\TextAt{8}{138.109}{761.124}{then}%
\TextAt{13}{165.984}{761.124}{\Character{27}}%
\TextAt{9}{178.458}{761.124}{\Character{61}}%
\TextAt{9}{192.936}{761.124}{1}%
\TextAt{8}{198.807}{761.124}{,}%
\TextAt{8}{207.475}{761.124}{so}%
\TextAt{8}{223.06}{761.124}{that}%
\TextAt{13}{249.644}{761.124}{v}%
\TextAt{21}{255.303}{759.33}{L}%
\TextAt{22}{261.062}{758.334}{n}%
\TextAt{9}{266.659}{761.124}{\Character{40}}%
\TextAt{13}{271.212}{761.124}{\Character{21}}%
\TextAt{21}{278.04}{765.462}{\Character{27}}%
\TextAt{16}{287.148}{761.124}{\Character{0}}%
\TextAt{13}{299.932}{761.124}{\Character{21}}%
\TextAt{9}{306.76}{761.124}{\Character{41}}%
\TextAt{9}{316.683}{761.124}{\Character{61}}%
\NativeMathAt{331.2}{761.124}{11.955}{\infty}%
\TextAt{8}{343.155}{761.124}{.}%
\TextAt{8}{355.248}{761.124}{If}%
\TextAt{13}{368.155}{761.124}{m}%
\TextAt{13}{383.769}{761.124}{\Character{60}}%
\TextAt{13}{398.258}{761.124}{n}%
\TextAt{8}{405.27}{761.124}{,}%
\TextAt{8}{413.938}{761.124}{then}%
\TextAt{13}{441.813}{761.124}{\Character{21}}%
\TextAt{21}{448.641}{759.33}{n}%
\TextAt{3}{453.779}{759.33}{\Character{0}}%
\TextAt{21}{460.366}{759.33}{m}%
\TextAt{9}{473.746}{761.124}{\Character{61}}%
\TextAt{9}{488.224}{761.124}{\Character{91}}%
\TextAt{13}{491.493}{761.124}{\Character{25}}%
\TextAt{21}{498.563}{765.462}{m}%
\TextAt{9}{506.551}{761.124}{\Character{93}}%
\TextAt{21}{509.803}{759.33}{F}%
\TextAt{9}{516.836}{761.124}{\Character{40}}%
\TextAt{13}{521.389}{761.124}{\Character{21}}%
\TextAt{9}{528.217}{761.124}{\Character{41}}%
\TextAt{8}{79.226}{744.078}{is}%
\TextAt{8}{91.897}{744.078}{a}%
\TextAt{8}{102.56}{744.078}{prime}%
\TextAt{8}{136.631}{744.078}{element}%
\TextAt{8}{180.78}{744.078}{of}%
\TextAt{13}{195.018}{744.078}{L}%
\TextAt{21}{202.982}{742.284}{n}%
\TextAt{3}{208.12}{742.284}{\Character{0}}%
\TextAt{21}{214.707}{742.284}{m}%
\TextAt{8}{227.5}{744.078}{and}%
\TextAt{8}{251.17}{744.078}{therefore}%
\TextAt{8}{301.178}{744.078}{also}%
\TextAt{9}{325.565}{744.078}{\Character{91}}%
\TextAt{13}{328.817}{744.078}{\Character{34}\Character{25}}%
\TextAt{21}{341.364}{748.416}{m}%
\TextAt{9}{349.352}{744.078}{\Character{93}}%
\TextAt{21}{352.604}{742.284}{F}%
\TextAt{9}{359.638}{744.078}{\Character{40}}%
\TextAt{13}{364.19}{744.078}{\Character{21}}%
\TextAt{9}{371.018}{744.078}{\Character{41}}%
\TextAt{9}{380.427}{744.078}{\Character{61}}%
\TextAt{9}{394.39}{744.078}{\Character{91}}%
\TextAt{13}{397.661}{744.078}{\Character{34}}%
\TextAt{9}{403.139}{744.078}{\Character{93}}%
\TextAt{21}{406.391}{742.284}{F}%
\TextAt{9}{413.424}{744.078}{\Character{40}}%
\TextAt{13}{417.977}{744.078}{\Character{21}}%
\TextAt{21}{424.805}{742.284}{n}%
\TextAt{3}{429.943}{742.284}{\Character{0}}%
\TextAt{21}{436.53}{742.284}{m}%
\TextAt{9}{444.518}{744.078}{\Character{41}}%
\TextAt{8}{449.071}{744.078}{.}%
\TextAt{8}{460.241}{744.078}{As}%
\TextAt{13}{478.436}{744.078}{L}%
\TextAt{21}{486.4}{742.284}{n}%
\TextAt{16}{496.904}{744.078}{\Character{106}}%
\TextAt{13}{505.092}{744.078}{L}%
\TextAt{21}{513.056}{742.284}{n}%
\TextAt{3}{518.194}{742.284}{\Character{0}}%
\TextAt{21}{524.781}{742.284}{m}%
\TextAt{8}{79.226}{727.032}{is}%
\TextAt{8}{92.383}{727.032}{totally}%
\TextAt{8}{131.183}{727.032}{ramified}%
\TextAt{8}{178.101}{727.032}{of}%
\TextAt{8}{192.824}{727.032}{degree}%
\TextAt{13}{230.632}{727.032}{q}%
\TextAt{21}{236.251}{731.37}{m}%
\TextAt{8}{249.53}{727.032}{we}%
\TextAt{8}{268.15}{727.032}{may}%
\TextAt{8}{294.912}{727.032}{write}%
\TextAt{9}{326.217}{727.032}{\Character{91}}%
\TextAt{13}{329.468}{727.032}{\Character{34}\Character{25}}%
\TextAt{21}{342.016}{731.37}{m}%
\TextAt{9}{350.004}{727.032}{\Character{93}}%
\TextAt{21}{353.256}{725.238}{F}%
\TextAt{9}{360.289}{727.032}{\Character{40}}%
\TextAt{13}{364.842}{727.032}{\Character{21}}%
\TextAt{9}{371.67}{727.032}{\Character{41}}%
\TextAt{9}{381.904}{727.032}{\Character{61}}%
\TextAt{13}{396.716}{727.032}{\Character{34}}%
\TextAt{17}{402.194}{725.238}{0}%
\TextAt{13}{406.926}{727.032}{\Character{21}}%
\TextAt{21}{413.755}{731.37}{q}%
\TextAt{22}{417.821}{734.183}{m}%
\TextAt{8}{430.646}{727.032}{for}%
\TextAt{8}{449.923}{727.032}{some}%
\TextAt{13}{480.642}{727.032}{\Character{34}}%
\TextAt{17}{486.12}{725.238}{0}%
\NativeMathAt{496.547}{727.032}{11.955}{\in}%
\TextAt{13}{510.211}{727.032}{U}%
\TextAt{21}{518.162}{725.238}{L}%
\TextAt{22}{523.921}{724.242}{n}%
\TextAt{8}{529.518}{727.032}{.}%
\TextAt{8}{79.226}{709.986}{Since}%
\TextAt{13}{111.709}{709.986}{F}%
\TextAt{9}{120.912}{709.986}{\Character{40}}%
\TextAt{13}{125.465}{709.986}{X}%
\TextAt{13}{135.46}{709.986}{\Character{59}}%
\TextAt{9}{140.714}{709.986}{0\Character{41}}%
\TextAt{9}{157.702}{709.986}{\Character{61}}%
\TextAt{13}{173.426}{709.986}{X}%
\TextAt{8}{184.082}{709.986}{,}%
\TextAt{13}{193.637}{709.986}{F}%
\TextAt{9}{202.841}{709.986}{\Character{40}0}%
\TextAt{13}{213.246}{709.986}{\Character{59}}%
\TextAt{13}{218.47}{709.986}{Y}%
\TextAt{9}{227.883}{709.986}{\Character{41}}%
\TextAt{9}{239.025}{709.986}{\Character{61}}%
\TextAt{13}{254.743}{709.986}{Y}%
\TextAt{8}{264.137}{709.986}{,}%
\TextAt{8}{273.692}{709.986}{we}%
\TextAt{8}{292.845}{709.986}{have}%
\TextAt{13}{321.765}{709.986}{F}%
\TextAt{9}{330.968}{709.986}{\Character{40}}%
\TextAt{13}{335.52}{709.986}{X}%
\TextAt{13}{345.514}{709.986}{\Character{59}}%
\TextAt{13}{350.739}{709.986}{Y}%
\TextAt{9}{360.163}{709.986}{\Character{41}}%
\TextAt{9}{371.305}{709.986}{\Character{61}}%
\TextAt{13}{387.023}{709.986}{X}%
\TextAt{9}{401.647}{709.986}{\Character{43}}%
\TextAt{13}{414.72}{709.986}{Y}%
\TextAt{9}{428.082}{709.986}{\Character{43}}%
\TextAt{13}{441.156}{709.986}{X}%
\TextAt{13}{451.796}{709.986}{Y}%
\TextAt{13}{461.193}{709.986}{G}%
\TextAt{9}{470.438}{709.986}{\Character{40}}%
\TextAt{13}{474.99}{709.986}{X}%
\TextAt{13}{484.985}{709.986}{\Character{59}}%
\TextAt{13}{490.209}{709.986}{Y}%
\TextAt{9}{499.632}{709.986}{\Character{41}}%
\TextAt{8}{510.007}{709.986}{with}%
\TextAt{13}{79.226}{692.94}{G}%
\TextAt{9}{88.46}{692.94}{\Character{40}}%
\TextAt{13}{93.012}{692.94}{X}%
\TextAt{13}{103.007}{692.94}{\Character{59}}%
\TextAt{13}{108.231}{692.94}{Y}%
\TextAt{9}{117.654}{692.94}{\Character{41}}%
\NativeMathAt{125.527}{692.94}{11.955}{\in}%
\TextAt{14}{136.818}{692.94}{o}%
\TextAt{21}{142.663}{691.147}{K}%
\TextAt{9}{150.865}{692.94}{\Character{91}\Character{91}}%
\TextAt{13}{157.369}{692.94}{X}%
\TextAt{13}{167.364}{692.94}{\Character{59}}%
\TextAt{13}{172.588}{692.94}{Y}%
\TextAt{9}{182.011}{692.94}{\Character{93}\Character{93}}%
\TextAt{8}{188.514}{692.94}{.}%
\TextAt{8}{196.962}{692.94}{Thus}%
\TextAt{13}{158.667}{663.939}{\Character{21}}%
\TextAt{21}{165.495}{668.875}{\Character{27}}%
\TextAt{16}{173.775}{663.939}{\Character{0}}%
\TextAt{13}{185.731}{663.939}{\Character{21}}%
\TextAt{9}{195.88}{663.939}{\Character{61}}%
\TextAt{13}{208.305}{663.939}{F}%
\TextAt{9}{217.509}{663.939}{\Character{40}}%
\TextAt{13}{222.061}{663.939}{\Character{21}\Character{59}}%
\TextAt{13}{234.112}{663.939}{\Character{34}}%
\TextAt{17}{239.612}{662.145}{0}%
\TextAt{13}{244.344}{663.939}{\Character{21}}%
\TextAt{21}{251.173}{668.875}{q}%
\TextAt{22}{255.24}{671.688}{m}%
\TextAt{9}{262.774}{663.939}{\Character{41}}%
\TextAt{16}{269.983}{663.939}{\Character{0}}%
\TextAt{13}{281.938}{663.939}{\Character{21}}%
\TextAt{9}{292.087}{663.939}{\Character{61}}%
\TextAt{13}{304.513}{663.939}{\Character{34}}%
\TextAt{17}{309.991}{662.145}{0}%
\TextAt{13}{314.723}{663.939}{\Character{21}}%
\TextAt{21}{321.552}{668.875}{q}%
\TextAt{22}{325.618}{671.688}{m}%
\TextAt{9}{335.809}{663.939}{\Character{43}}%
\TextAt{13}{347.571}{663.939}{a\Character{21}}%
\TextAt{21}{360.544}{668.875}{q}%
\TextAt{22}{364.611}{671.688}{m}%
\TextAt{17}{371.647}{668.875}{\Character{43}1}%
\TextAt{13}{382.966}{663.939}{\Character{59}}%
\TextAt{13}{412.101}{663.939}{a}%
\NativeMathAt{421.586}{663.939}{11.955}{\in}%
\TextAt{14}{432.877}{663.939}{o}%
\TextAt{21}{438.721}{662.145}{L}%
\TextAt{22}{444.48}{661.149}{n}%
\TextAt{13}{450.077}{663.939}{\Character{59}}%
\TextAt{8}{79.226}{634.938}{i.e.,}%
\TextAt{13}{193.933}{607.126}{i}%
\TextAt{21}{197.927}{605.266}{L}%
\TextAt{22}{203.686}{604.27}{n}%
\TextAt{3}{208.784}{605.266}{\Character{106}}%
\TextAt{21}{211.137}{605.266}{K}%
\TextAt{9}{219.339}{607.126}{\Character{40}}%
\TextAt{13}{223.892}{607.126}{\Character{27}}%
\TextAt{9}{230.974}{607.126}{\Character{41}}%
\TextAt{9}{238.829}{607.126}{\Character{58}\Character{61}}%
\TextAt{13}{254.525}{607.126}{v}%
\TextAt{21}{260.184}{605.332}{L}%
\TextAt{22}{265.943}{604.336}{n}%
\TextAt{9}{271.54}{607.126}{\Character{40}}%
\TextAt{13}{276.092}{607.126}{\Character{21}}%
\TextAt{21}{282.921}{612.062}{\Character{27}}%
\TextAt{16}{291.201}{607.126}{\Character{0}}%
\TextAt{13}{303.156}{607.126}{\Character{21}}%
\TextAt{9}{309.984}{607.126}{\Character{41}}%
\TextAt{9}{317.839}{607.126}{\Character{61}}%
\TextAt{20}{330.283}{631.634}{\Character{56}}%
\TextAt{20}{330.283}{620.874}{\Character{60}}%
\TextAt{20}{330.283}{599.355}{\Character{58}}%
\TextAt{13}{340.91}{616.251}{q}%
\TextAt{21}{346.529}{620.589}{m}%
\TextAt{13}{354.518}{616.251}{\Character{59}}%
\TextAt{8}{369.725}{616.251}{if}%
\TextAt{13}{380.448}{616.251}{m}%
\TextAt{13}{393.993}{616.251}{\Character{60}}%
\TextAt{13}{406.415}{616.251}{n\Character{59}}%
\NativeMathAt{340.91}{595.796}{11.955}{\infty}%
\TextAt{13}{352.865}{595.796}{\Character{59}}%
\TextAt{8}{369.725}{595.796}{if}%
\TextAt{13}{380.448}{595.796}{m}%
\NativeMathAt{394.008}{595.796}{11.955}{\geq}%
\TextAt{13}{406.628}{595.796}{n\Character{58}}%
\TextAt{8}{79.226}{572.157}{By}%
\TextAt{8}{98.712}{572.157}{chap.}%
\TextAt{8}{134.272}{572.157}{II,}%
\TextAt{8}{151.307}{572.157}{§10,}%
\TextAt{8}{177.254}{572.157}{we}%
\TextAt{8}{195.605}{572.157}{have}%
\TextAt{13}{223.722}{572.157}{G}%
\TextAt{21}{232.956}{570.363}{i}%
\TextAt{9}{236.337}{572.157}{\Character{40}}%
\TextAt{13}{240.889}{572.157}{L}%
\TextAt{21}{248.854}{570.363}{n}%
\TextAt{16}{259.724}{572.157}{\Character{106}}%
\TextAt{13}{268.28}{572.157}{K}%
\TextAt{9}{279.097}{572.157}{\Character{41}}%
\TextAt{9}{288.864}{572.157}{\Character{61}}%
\TextAt{16}{303.223}{572.157}{\Character{102}}%
\TextAt{13}{309.2}{572.157}{\Character{27}}%
\NativeMathAt{321.517}{572.157}{11.955}{\in}%
\TextAt{13}{334.721}{572.157}{G}%
\TextAt{9}{343.955}{572.157}{\Character{40}}%
\TextAt{13}{348.507}{572.157}{L}%
\TextAt{21}{356.472}{570.363}{n}%
\TextAt{16}{367.342}{572.157}{\Character{106}}%
\TextAt{13}{375.898}{572.157}{K}%
\TextAt{9}{386.715}{572.157}{\Character{41}}%
\TextAt{16}{396.502}{572.157}{\Character{106}}%
\TextAt{13}{405.057}{572.157}{i}%
\TextAt{21}{409.05}{570.297}{L}%
\TextAt{22}{414.809}{569.301}{n}%
\TextAt{3}{419.908}{570.297}{\Character{106}}%
\TextAt{21}{422.261}{570.297}{K}%
\TextAt{9}{430.463}{572.157}{\Character{40}}%
\TextAt{13}{435.016}{572.157}{\Character{27}}%
\TextAt{9}{442.098}{572.157}{\Character{41}}%
\NativeMathAt{451.885}{572.157}{11.955}{\geq}%
\TextAt{13}{466.417}{572.157}{i}%
\TextAt{9}{473.833}{572.157}{\Character{43}}%
\TextAt{9}{486.35}{572.157}{1}%
\TextAt{16}{492.212}{572.157}{\Character{103}}%
\TextAt{8}{498.19}{572.157}{.}%
\TextAt{8}{510.007}{572.157}{Now}%
\TextAt{8}{79.226}{555.111}{let}%
\TextAt{13}{97.299}{555.111}{q}%
\TextAt{21}{102.918}{559.449}{k}%
\TextAt{3}{107.54}{559.449}{\Character{0}}%
\TextAt{17}{114.126}{559.449}{1}%
\NativeMathAt{124.17}{555.111}{11.955}{\leq}%
\TextAt{13}{138.78}{555.111}{i}%
\NativeMathAt{148.084}{555.111}{11.955}{\leq}%
\TextAt{13}{162.694}{555.111}{q}%
\TextAt{21}{168.313}{559.449}{k}%
\TextAt{16}{176.886}{555.111}{\Character{0}}%
\TextAt{9}{189.637}{555.111}{1}%
\TextAt{8}{195.49}{555.111}{.}%
\TextAt{8}{207.443}{555.111}{If}%
\TextAt{13}{220.303}{555.111}{u}%
\NativeMathAt{232.277}{555.111}{11.955}{\in}%
\TextAt{13}{245.558}{555.111}{U}%
\TextAt{17}{254.757}{561.221}{\Character{40}}%
\TextAt{21}{258.051}{561.221}{k}%
\TextAt{17}{262.672}{561.221}{\Character{41}}%
\TextAt{21}{253.509}{551.869}{K}%
\TextAt{8}{266.464}{555.111}{,}%
\TextAt{8}{275.073}{555.111}{then}%
\TextAt{13}{302.901}{555.111}{m}%
\NativeMathAt{318.452}{555.111}{11.955}{\geq}%
\TextAt{13}{333.061}{555.111}{k}%
\TextAt{8}{339.551}{555.111}{,}%
\TextAt{8}{348.16}{555.111}{i.e.,}%
\TextAt{13}{371.725}{555.111}{i}%
\TextAt{21}{375.719}{553.251}{L}%
\TextAt{22}{381.478}{552.255}{n}%
\TextAt{3}{386.577}{553.251}{\Character{106}}%
\TextAt{21}{388.929}{553.251}{K}%
\TextAt{9}{397.132}{555.111}{\Character{40}}%
\TextAt{13}{401.684}{555.111}{\Character{27}}%
\TextAt{9}{408.766}{555.111}{\Character{41}}%
\NativeMathAt{418.63}{555.111}{11.955}{\geq}%
\TextAt{13}{433.24}{555.111}{q}%
\TextAt{21}{438.859}{559.449}{k}%
\NativeMathAt{449.29}{555.111}{11.955}{\geq}%
\TextAt{13}{463.9}{555.111}{i}%
\TextAt{9}{471.346}{555.111}{\Character{43}}%
\TextAt{9}{483.887}{555.111}{1}%
\TextAt{8}{489.757}{555.111}{,}%
\TextAt{8}{498.366}{555.111}{and}%
\TextAt{8}{522.297}{555.111}{so}%
\TextAt{13}{79.226}{538.065}{\Character{27}}%
\NativeMathAt{90.936}{538.065}{11.955}{\in}%
\TextAt{13}{103.534}{538.065}{G}%
\TextAt{21}{112.768}{536.271}{i}%
\TextAt{9}{116.149}{538.065}{\Character{40}}%
\TextAt{13}{120.701}{538.065}{L}%
\TextAt{21}{128.666}{536.271}{n}%
\TextAt{16}{138.93}{538.065}{\Character{106}}%
\TextAt{13}{146.879}{538.065}{K}%
\TextAt{9}{157.696}{538.065}{\Character{41}}%
\TextAt{8}{162.249}{538.065}{.}%
\TextAt{8}{172.998}{538.065}{This}%
\TextAt{8}{200.485}{538.065}{proves}%
\TextAt{8}{237.416}{538.065}{the}%
\TextAt{8}{258.339}{538.065}{inclusion}%
\TextAt{9}{307.943}{538.065}{\Character{40}}%
\TextAt{13}{312.495}{538.065}{U}%
\TextAt{17}{321.694}{544.175}{\Character{40}}%
\TextAt{21}{324.988}{544.175}{k}%
\TextAt{17}{329.609}{544.175}{\Character{41}}%
\TextAt{21}{320.446}{534.823}{K}%
\TextAt{13}{333.401}{538.065}{\Character{59}}%
\TextAt{13}{338.625}{538.065}{L}%
\TextAt{21}{346.609}{536.271}{n}%
\TextAt{16}{356.873}{538.065}{\Character{106}}%
\TextAt{13}{364.822}{538.065}{K}%
\TextAt{9}{375.64}{538.065}{\Character{41}}%
\NativeMathAt{384.82}{538.065}{11.955}{\subseteq}%
\TextAt{13}{398.746}{538.065}{G}%
\TextAt{21}{407.98}{536.271}{i}%
\TextAt{9}{411.361}{538.065}{\Character{40}}%
\TextAt{13}{415.913}{538.065}{L}%
\TextAt{21}{423.878}{536.271}{n}%
\TextAt{16}{434.142}{538.065}{\Character{106}}%
\TextAt{13}{442.091}{538.065}{K}%
\TextAt{9}{452.908}{538.065}{\Character{41}}%
\TextAt{8}{457.461}{538.065}{.}%
\TextAt{8}{468.21}{538.065}{If}%
\TextAt{8}{480.669}{538.065}{conversely}%
\TextAt{13}{79.226}{521.019}{\Character{27}}%
\NativeMathAt{90.611}{521.019}{11.955}{\in}%
\TextAt{13}{102.883}{521.019}{G}%
\TextAt{21}{112.117}{519.225}{i}%
\TextAt{9}{115.498}{521.019}{\Character{40}}%
\TextAt{13}{120.051}{521.019}{L}%
\TextAt{21}{128.015}{519.225}{n}%
\TextAt{16}{137.954}{521.019}{\Character{106}}%
\TextAt{13}{145.577}{521.019}{K}%
\TextAt{9}{156.395}{521.019}{\Character{41}}%
\TextAt{8}{165.42}{521.019}{and}%
\TextAt{13}{188.759}{521.019}{\Character{27}}%
\NativeMathAt{200.144}{521.019}{11.955}{\neq}%
\TextAt{9}{213.534}{521.019}{1}%
\TextAt{8}{219.404}{521.019}{,}%
\TextAt{8}{227.273}{521.019}{then}%
\TextAt{13}{254.509}{521.019}{i}%
\TextAt{21}{258.502}{519.159}{L}%
\TextAt{22}{264.261}{518.163}{n}%
\TextAt{3}{269.36}{519.159}{\Character{106}}%
\TextAt{21}{271.712}{519.159}{K}%
\TextAt{9}{279.915}{521.019}{\Character{40}}%
\TextAt{13}{284.467}{521.019}{\Character{27}}%
\TextAt{9}{291.55}{521.019}{\Character{41}}%
\TextAt{9}{300.385}{521.019}{\Character{61}}%
\TextAt{13}{313.811}{521.019}{q}%
\TextAt{21}{319.431}{525.357}{m}%
\TextAt{13}{331.722}{521.019}{\Character{62}}%
\TextAt{13}{345.112}{521.019}{i}%
\NativeMathAt{353.424}{521.019}{11.955}{\geq}%
\TextAt{13}{367.025}{521.019}{q}%
\TextAt{21}{372.644}{525.357}{k}%
\TextAt{3}{377.266}{525.357}{\Character{0}}%
\TextAt{17}{383.852}{525.357}{1}%
\TextAt{8}{388.585}{521.019}{,}%
\TextAt{8}{396.454}{521.019}{i.e.,}%
\TextAt{13}{419.279}{521.019}{m}%
\NativeMathAt{433.821}{521.019}{11.955}{\geq}%
\TextAt{13}{447.421}{521.019}{k}%
\TextAt{8}{453.911}{521.019}{.}%
\TextAt{8}{464.087}{521.019}{Consequently}%
\TextAt{13}{79.226}{503.973}{u}%
\NativeMathAt{89.209}{503.973}{11.955}{\in}%
\TextAt{13}{100.5}{503.973}{U}%
\TextAt{17}{109.699}{510.083}{\Character{40}}%
\TextAt{21}{112.993}{510.083}{k}%
\TextAt{17}{117.614}{510.083}{\Character{41}}%
\TextAt{21}{108.451}{500.732}{K}%
\TextAt{8}{121.406}{503.973}{,}%
\TextAt{8}{128.555}{503.973}{and}%
\TextAt{8}{151.318}{503.973}{this}%
\TextAt{8}{174.14}{503.973}{shows}%
\TextAt{8}{207.758}{503.973}{the}%
\TextAt{8}{227.914}{503.973}{inclusion}%
\TextAt{13}{276.751}{503.973}{G}%
\TextAt{21}{285.985}{502.179}{i}%
\TextAt{9}{289.366}{503.973}{\Character{40}}%
\TextAt{13}{293.918}{503.973}{L}%
\TextAt{21}{301.883}{502.179}{n}%
\TextAt{16}{310.84}{503.973}{\Character{106}}%
\TextAt{13}{317.482}{503.973}{K}%
\TextAt{9}{328.299}{503.973}{\Character{41}}%
\NativeMathAt{336.173}{503.973}{11.955}{\subseteq}%
\TextAt{9}{348.792}{503.973}{\Character{40}}%
\TextAt{13}{353.344}{503.973}{U}%
\TextAt{17}{362.543}{510.083}{\Character{40}}%
\TextAt{21}{365.837}{510.083}{k}%
\TextAt{17}{370.458}{510.083}{\Character{41}}%
\TextAt{21}{361.295}{500.732}{K}%
\TextAt{13}{374.25}{503.973}{\Character{59}}%
\TextAt{13}{379.474}{503.973}{L}%
\TextAt{21}{387.458}{502.179}{n}%
\TextAt{16}{396.415}{503.973}{\Character{106}}%
\TextAt{13}{403.057}{503.973}{K}%
\TextAt{9}{413.875}{503.973}{\Character{41}}%
\TextAt{8}{418.427}{503.973}{.}%
\Vector{q 0 G 0 g 0.399 w 452.334 -270.027 m 452.334 -261.958 l S Q}%
\Vector{q 0 G 0 g 0.399 w 452.534 -262.157 m 459.707 -262.157 l S Q}%
\Vector{q 0 G 0 g 0.399 w 452.534 -269.828 m 459.707 -269.828 l S Q}%
\Vector{q 0 G 0 g 0.399 w 459.906 -270.027 m 459.906 -261.958 l S Q}%
\TextAt{8}{97.159}{474.972}{From}%
\TextAt{8}{127.663}{474.972}{this}%
\TextAt{8}{150.265}{474.972}{proposition}%
\TextAt{8}{211.578}{474.972}{we}%
\TextAt{8}{228.585}{474.972}{get}%
\TextAt{8}{247.875}{474.972}{the}%
\TextAt{8}{267.811}{474.972}{following}%
\TextAt{8}{317.026}{474.972}{result,}%
\TextAt{8}{352.679}{474.972}{which}%
\TextAt{8}{385.945}{474.972}{may}%
\TextAt{8}{411.093}{474.972}{be}%
\TextAt{8}{426.797}{474.972}{considered}%
\TextAt{8}{483.866}{474.972}{the}%
\TextAt{12}{503.802}{474.972}{main}%
\TextAt{12}{79.226}{457.926}{theorem}%
\TextAt{12}{132.128}{457.926}{of}%
\TextAt{12}{147.454}{457.926}{higher}%
\TextAt{12}{188.987}{457.926}{ramification}%
\TextAt{12}{263.635}{457.926}{theory}%
\TextAt{8}{301.807}{457.926}{.}%
\TextAt{12}{79.226}{431.913}{Theorem}%
\TextAt{12}{136.24}{431.913}{(6.2).}%
\TextAt{8}{171.556}{431.913}{If}%
\TextAt{13}{183.248}{431.913}{L}%
\TextAt{16}{194.533}{431.913}{\Character{106}}%
\TextAt{13}{201.175}{431.913}{K}%
\TextAt{8}{215.89}{431.913}{is}%
\TextAt{8}{227.654}{431.913}{a}%
\TextAt{8}{237.409}{431.913}{finite}%
\TextAt{8}{267.321}{431.913}{abelian}%
\TextAt{8}{307.969}{431.913}{extension,}%
\TextAt{8}{362.987}{431.913}{then}%
\TextAt{8}{389.647}{431.913}{the}%
\TextAt{8}{409.803}{431.913}{norm}%
\TextAt{8}{440.372}{431.913}{residue}%
\TextAt{8}{480.099}{431.913}{symbol}%
\TextAt{9}{233.288}{402.912}{\Character{40}}%
\TextAt{13}{241.737}{402.912}{\Character{59}}%
\TextAt{13}{246.961}{402.912}{L}%
\TextAt{16}{258.267}{402.912}{\Character{106}}%
\TextAt{13}{264.909}{402.912}{K}%
\TextAt{9}{275.726}{402.912}{\Character{41}}%
\TextAt{9}{283.581}{402.912}{\Character{58}}%
\TextAt{13}{290.172}{402.912}{K}%
\TextAt{3}{300.99}{407.848}{\Character{3}}%
\NativeArrowAt{309.043}{402.912}{11.955}{19.26}{\longrightarrow}%
\TextAt{13}{331.625}{402.912}{G}%
\TextAt{9}{340.858}{402.912}{\Character{40}}%
\TextAt{13}{345.411}{402.912}{L}%
\TextAt{16}{356.696}{402.912}{\Character{106}}%
\TextAt{13}{363.338}{402.912}{K}%
\TextAt{9}{374.156}{402.912}{\Character{41}}%
\TextAt{8}{79.226}{373.911}{maps}%
\TextAt{8}{109.855}{373.911}{the}%
\TextAt{8}{130.012}{373.911}{group}%
\TextAt{13}{163.187}{373.911}{U}%
\TextAt{17}{172.386}{380.021}{\Character{40}}%
\TextAt{21}{175.68}{380.021}{n}%
\TextAt{17}{180.818}{380.021}{\Character{41}}%
\TextAt{21}{171.138}{370.67}{K}%
\TextAt{8}{188.507}{373.911}{onto}%
\TextAt{8}{214.856}{373.911}{the}%
\TextAt{8}{235.012}{373.911}{group}%
\TextAt{13}{268.188}{373.911}{G}%
\TextAt{21}{277.421}{378.249}{n}%
\TextAt{9}{283.058}{373.911}{\Character{40}}%
\TextAt{13}{287.61}{373.911}{L}%
\TextAt{16}{298.895}{373.911}{\Character{106}}%
\TextAt{13}{305.537}{373.911}{K}%
\TextAt{9}{316.355}{373.911}{\Character{41}}%
\TextAt{8}{320.907}{373.911}{,}%
\TextAt{8}{328.056}{373.911}{for}%
\TextAt{13}{345.941}{373.911}{n}%
\NativeMathAt{356.25}{373.911}{11.955}{\geq}%
\TextAt{9}{368.869}{373.911}{0}%
\TextAt{8}{374.722}{373.911}{.}%
\TextAt{19}{79.226}{338.932}{Proof.}%
\TextAt{8}{116}{338.932}{We}%
\TextAt{8}{137.458}{338.932}{may}%
\TextAt{8}{164.14}{338.932}{assume}%
\TextAt{8}{205.898}{338.932}{that}%
\TextAt{13}{232.58}{338.932}{L}%
\TextAt{16}{246.104}{338.932}{\Character{106}}%
\TextAt{13}{254.984}{338.932}{K}%
\TextAt{8}{271.012}{338.932}{is}%
\TextAt{8}{284.089}{338.932}{totally}%
\TextAt{8}{322.811}{338.932}{ramified.}%
\TextAt{8}{376.827}{338.932}{For}%
\TextAt{8}{399.11}{338.932}{if}%
\TextAt{13}{411.147}{338.932}{L}%
\TextAt{17}{419.111}{343.271}{0}%
\TextAt{16}{429.403}{338.932}{\Character{106}}%
\TextAt{13}{438.283}{338.932}{K}%
\TextAt{8}{454.311}{338.932}{is}%
\TextAt{8}{467.388}{338.932}{the}%
\TextAt{8}{488.858}{338.932}{maximal}%
\TextAt{8}{79.226}{321.886}{unramified}%
\TextAt{8}{137.295}{321.886}{subextension}%
\TextAt{8}{206.541}{321.886}{of}%
\TextAt{13}{219.407}{321.886}{L}%
\TextAt{16}{230.693}{321.886}{\Character{106}}%
\TextAt{13}{237.334}{321.886}{K}%
\TextAt{8}{248.152}{321.886}{,}%
\TextAt{8}{254.93}{321.886}{then}%
\TextAt{8}{281.126}{321.886}{we}%
\TextAt{8}{297.889}{321.886}{have}%
\TextAt{8}{324.42}{321.886}{on}%
\TextAt{8}{340.215}{321.886}{the}%
\TextAt{8}{359.908}{321.886}{one}%
\TextAt{8}{380.903}{321.886}{hand}%
\TextAt{13}{409.705}{321.886}{G}%
\TextAt{21}{418.939}{326.225}{n}%
\TextAt{9}{424.575}{321.886}{\Character{40}}%
\TextAt{13}{429.128}{321.886}{L}%
\TextAt{16}{440.413}{321.886}{\Character{106}}%
\TextAt{13}{447.055}{321.886}{K}%
\TextAt{9}{457.872}{321.886}{\Character{41}}%
\TextAt{9}{465.727}{321.886}{\Character{61}}%
\TextAt{13}{478.171}{321.886}{G}%
\TextAt{21}{487.405}{326.225}{n}%
\TextAt{9}{493.041}{321.886}{\Character{40}}%
\TextAt{13}{497.593}{321.886}{L}%
\TextAt{16}{508.879}{321.886}{\Character{106}}%
\TextAt{13}{515.52}{321.886}{L}%
\TextAt{17}{523.485}{326.225}{0}%
\TextAt{9}{528.217}{321.886}{\Character{41}}%
\TextAt{8}{79.226}{304.84}{because}%
\TextAt{13}{122.062}{304.84}{\Character{32}}%
\TextAt{21}{129.672}{302.837}{L}%
\TextAt{24}{135.431}{305.106}{0}%
\TextAt{3}{139.582}{302.837}{\Character{106}}%
\TextAt{21}{141.934}{302.837}{K}%
\TextAt{9}{150.137}{304.84}{\Character{40}}%
\TextAt{13}{154.689}{304.84}{s}%
\TextAt{9}{160.203}{304.84}{\Character{41}}%
\TextAt{9}{168.058}{304.84}{\Character{61}}%
\TextAt{13}{180.502}{304.84}{s}%
\TextAt{8}{189.447}{304.84}{and}%
\TextAt{13}{211.744}{304.84}{\Character{32}}%
\TextAt{21}{219.354}{302.981}{L}%
\TextAt{3}{225.113}{302.981}{\Character{106}}%
\TextAt{21}{227.466}{302.981}{K}%
\TextAt{9}{235.668}{304.84}{\Character{40}}%
\TextAt{13}{240.221}{304.84}{s}%
\TextAt{9}{245.735}{304.84}{\Character{41}}%
\TextAt{9}{253.59}{304.84}{\Character{61}}%
\TextAt{13}{266.033}{304.84}{\Character{32}}%
\TextAt{21}{273.644}{302.837}{L}%
\TextAt{3}{279.403}{302.837}{\Character{106}}%
\TextAt{21}{281.755}{302.837}{L}%
\TextAt{24}{287.514}{305.106}{0}%
\TextAt{9}{292.163}{304.84}{\Character{40}}%
\TextAt{13}{296.716}{304.84}{\Character{32}}%
\TextAt{21}{304.326}{302.837}{L}%
\TextAt{24}{310.085}{305.106}{0}%
\TextAt{3}{314.236}{302.837}{\Character{106}}%
\TextAt{21}{316.588}{302.837}{K}%
\TextAt{9}{324.791}{304.84}{\Character{40}}%
\TextAt{13}{329.343}{304.84}{s}%
\TextAt{9}{334.857}{304.84}{\Character{41}\Character{41}}%
\TextAt{9}{347.255}{304.84}{\Character{61}}%
\TextAt{13}{359.708}{304.84}{\Character{32}}%
\TextAt{21}{367.319}{302.837}{L}%
\TextAt{3}{373.078}{302.837}{\Character{106}}%
\TextAt{21}{375.43}{302.837}{L}%
\TextAt{24}{381.189}{305.106}{0}%
\TextAt{9}{385.838}{304.84}{\Character{40}}%
\TextAt{13}{390.39}{304.84}{s}%
\TextAt{9}{395.904}{304.84}{\Character{41}}%
\TextAt{8}{403.888}{304.84}{(see}%
\TextAt{8}{426.89}{304.84}{chap.}%
\TextAt{8}{458.926}{304.84}{II,}%
\TextAt{8}{474.372}{304.84}{(10.8)).}%
\TextAt{8}{517.156}{304.84}{On}%
\TextAt{8}{79.226}{287.794}{the}%
\TextAt{8}{99.383}{287.794}{other}%
\TextAt{8}{129.952}{287.794}{hand,}%
\TextAt{8}{162.47}{287.794}{by}%
\TextAt{8}{178.729}{287.794}{chap.}%
\TextAt{8}{210.92}{287.794}{IV,}%
\TextAt{8}{231.065}{287.794}{(6.4),}%
\TextAt{8}{262.292}{287.794}{and}%
\TextAt{8}{285.054}{287.794}{chap.}%
\TextAt{8}{317.246}{287.794}{V,}%
\TextAt{8}{333.17}{287.794}{(1.2),}%
\TextAt{8}{364.397}{287.794}{we}%
\TextAt{8}{381.624}{287.794}{have}%
\TextAt{9}{175.994}{258.793}{\Character{40}}%
\TextAt{13}{180.546}{258.793}{U}%
\TextAt{17}{189.745}{264.904}{\Character{40}}%
\TextAt{21}{193.038}{264.904}{n}%
\TextAt{17}{198.176}{264.904}{\Character{41}}%
\TextAt{21}{188.496}{254.877}{L}%
\TextAt{24}{194.255}{257.146}{0}%
\TextAt{13}{201.968}{258.793}{\Character{59}}%
\TextAt{13}{207.192}{258.793}{L}%
\TextAt{16}{218.497}{258.793}{\Character{106}}%
\TextAt{13}{225.139}{258.793}{L}%
\TextAt{17}{233.103}{263.729}{0}%
\TextAt{9}{237.836}{258.793}{\Character{41}}%
\TextAt{9}{245.691}{258.793}{\Character{61}}%
\TextAt{9}{258.112}{258.793}{\Character{40}}%
\TextAt{13}{262.687}{258.793}{N}%
\TextAt{21}{272.061}{256.79}{L}%
\TextAt{24}{277.82}{259.059}{0}%
\TextAt{3}{281.971}{256.79}{\Character{106}}%
\TextAt{21}{284.323}{256.79}{K}%
\TextAt{13}{292.526}{258.793}{U}%
\TextAt{17}{301.725}{264.904}{\Character{40}}%
\TextAt{21}{305.018}{264.904}{n}%
\TextAt{17}{310.156}{264.904}{\Character{41}}%
\TextAt{21}{300.476}{254.877}{L}%
\TextAt{24}{306.235}{257.146}{0}%
\TextAt{13}{313.948}{258.793}{\Character{59}}%
\TextAt{13}{319.172}{258.793}{L}%
\TextAt{16}{330.477}{258.793}{\Character{106}}%
\TextAt{13}{337.119}{258.793}{K}%
\TextAt{9}{347.936}{258.793}{\Character{41}}%
\TextAt{9}{355.791}{258.793}{\Character{61}}%
\TextAt{9}{368.212}{258.793}{\Character{40}}%
\TextAt{13}{372.787}{258.793}{U}%
\TextAt{17}{381.986}{264.904}{\Character{40}}%
\TextAt{21}{385.28}{264.904}{n}%
\TextAt{17}{390.418}{264.904}{\Character{41}}%
\TextAt{21}{380.738}{255.552}{K}%
\TextAt{13}{394.209}{258.793}{\Character{59}}%
\TextAt{13}{399.433}{258.793}{L}%
\TextAt{16}{410.739}{258.793}{\Character{106}}%
\TextAt{13}{417.38}{258.793}{K}%
\TextAt{9}{428.198}{258.793}{\Character{41}}%
\TextAt{13}{432.75}{258.793}{\Character{59}}%
\TextAt{8}{79.226}{229.792}{so}%
\TextAt{8}{93.596}{229.792}{we}%
\TextAt{8}{110.824}{229.792}{may}%
\TextAt{8}{136.193}{229.792}{replace}%
\TextAt{13}{175.86}{229.792}{L}%
\TextAt{16}{187.145}{229.792}{\Character{106}}%
\TextAt{13}{193.787}{229.792}{K}%
\TextAt{8}{208.502}{229.792}{by}%
\TextAt{13}{224.761}{229.792}{L}%
\TextAt{16}{236.046}{229.792}{\Character{106}}%
\TextAt{13}{242.688}{229.792}{L}%
\TextAt{17}{250.653}{234.13}{0}%
\TextAt{8}{255.385}{229.792}{.}%
\TextAt{8}{97.159}{209.757}{If}%
\TextAt{8}{108.682}{209.757}{now}%
\TextAt{13}{132.902}{209.757}{L}%
\TextAt{16}{144.188}{209.757}{\Character{106}}%
\TextAt{13}{150.829}{209.757}{K}%
\TextAt{8}{165.376}{209.757}{is}%
\TextAt{8}{176.971}{209.757}{totally}%
\TextAt{8}{214.21}{209.757}{ramified}%
\TextAt{8}{259.567}{209.757}{and}%
\TextAt{13}{282.161}{209.757}{\Character{25}}%
\TextAt{21}{288.801}{207.964}{L}%
\TextAt{8}{298.787}{209.757}{is}%
\TextAt{8}{310.382}{209.757}{a}%
\TextAt{8}{319.969}{209.757}{prime}%
\TextAt{8}{352.964}{209.757}{element}%
\TextAt{8}{396.037}{209.757}{of}%
\TextAt{13}{409.198}{209.757}{L}%
\TextAt{8}{417.163}{209.757}{,}%
\TextAt{8}{424.177}{209.757}{then}%
\TextAt{13}{450.669}{209.757}{\Character{25}}%
\TextAt{9}{461.059}{209.757}{\Character{61}}%
\TextAt{13}{473.484}{209.757}{N}%
\TextAt{21}{482.858}{207.898}{L}%
\TextAt{3}{488.617}{207.898}{\Character{106}}%
\TextAt{21}{490.969}{207.898}{K}%
\TextAt{9}{499.172}{209.757}{\Character{40}}%
\TextAt{13}{503.724}{209.757}{\Character{25}}%
\TextAt{21}{510.365}{207.964}{L}%
\TextAt{9}{516.622}{209.757}{\Character{41}}%
\TextAt{8}{524.903}{209.757}{is}%
\TextAt{8}{79.226}{192.711}{a}%
\TextAt{8}{89.314}{192.711}{prime}%
\TextAt{8}{122.81}{192.711}{element}%
\TextAt{8}{166.385}{192.711}{of}%
\TextAt{13}{180.047}{192.711}{K}%
\TextAt{8}{195.095}{192.711}{and}%
\TextAt{9}{218.19}{192.711}{\Character{40}}%
\TextAt{13}{222.742}{192.711}{\Character{25}}%
\TextAt{9}{229.811}{192.711}{\Character{41}}%
\TextAt{16}{237.247}{192.711}{\Character{2}}%
\TextAt{13}{249.429}{192.711}{U}%
\TextAt{17}{258.628}{198.822}{\Character{40}}%
\TextAt{21}{261.921}{198.822}{m}%
\TextAt{17}{269.412}{198.822}{\Character{41}}%
\TextAt{21}{257.379}{189.47}{K}%
\NativeMathAt{277.091}{192.711}{11.955}{\subseteq}%
\TextAt{13}{290.277}{192.711}{N}%
\TextAt{21}{299.65}{190.852}{L}%
\TextAt{3}{305.409}{190.852}{\Character{106}}%
\TextAt{21}{307.762}{190.852}{K}%
\TextAt{13}{315.964}{192.711}{L}%
\TextAt{3}{323.929}{197.05}{\Character{3}}%
\TextAt{8}{332.891}{192.711}{for}%
\TextAt{13}{351.109}{192.711}{m}%
\TextAt{8}{365.578}{192.711}{sufficiently}%
\TextAt{8}{424.502}{192.711}{big.}%
\TextAt{8}{449.562}{192.711}{Therefore}%
\TextAt{13}{502.892}{192.711}{L}%
\TextAt{16}{514.743}{192.711}{\Character{106}}%
\TextAt{13}{521.952}{192.711}{K}%
\TextAt{8}{79.226}{175.665}{is}%
\TextAt{8}{91.573}{175.665}{contained}%
\TextAt{8}{145.165}{175.665}{in}%
\TextAt{8}{159.401}{175.665}{the}%
\TextAt{8}{180.14}{175.665}{class}%
\TextAt{8}{208.16}{175.665}{field}%
\TextAt{8}{234.1}{175.665}{of}%
\TextAt{9}{248.013}{175.665}{\Character{40}}%
\TextAt{13}{252.565}{175.665}{\Character{25}}%
\TextAt{9}{259.634}{175.665}{\Character{41}}%
\TextAt{16}{267.241}{175.665}{\Character{2}}%
\TextAt{13}{279.593}{175.665}{U}%
\TextAt{17}{288.792}{181.776}{\Character{40}}%
\TextAt{21}{292.085}{181.776}{m}%
\TextAt{17}{299.576}{181.776}{\Character{41}}%
\TextAt{21}{287.543}{172.424}{K}%
\TextAt{8}{303.367}{175.665}{,}%
\TextAt{8}{311.245}{175.665}{which,}%
\TextAt{8}{348.712}{175.665}{by}%
\TextAt{8}{365.554}{175.665}{(5.6),}%
\TextAt{8}{397.509}{175.665}{is}%
\TextAt{8}{409.856}{175.665}{equal}%
\TextAt{8}{441.331}{175.665}{to}%
\TextAt{8}{456.224}{175.665}{the}%
\TextAt{8}{476.964}{175.665}{field}%
\TextAt{13}{502.903}{175.665}{L}%
\TextAt{21}{510.868}{173.872}{m}%
\TextAt{8}{523.337}{175.665}{of}%
\TextAt{13}{79.226}{158.619}{\Character{25}}%
\TextAt{21}{86.295}{162.958}{m}%
\TextAt{8}{94.284}{158.619}{-division}%
\TextAt{8}{142.417}{158.619}{points}%
\TextAt{8}{178.524}{158.619}{of}%
\TextAt{8}{192.776}{158.619}{some}%
\TextAt{8}{223.024}{158.619}{Lubin–Tate}%
\TextAt{8}{286.866}{158.619}{module}%
\TextAt{8}{329.082}{158.619}{for}%
\TextAt{13}{347.889}{158.619}{\Character{25}}%
\TextAt{8}{354.958}{158.619}{.}%
\TextAt{8}{366.173}{158.619}{In}%
\TextAt{8}{381.717}{158.619}{view}%
\TextAt{8}{409.622}{158.619}{of}%
\TextAt{8}{423.874}{158.619}{chap.}%
\TextAt{8}{458.832}{158.619}{II,}%
\TextAt{8}{475.666}{158.619}{(10.9),}%
\TextAt{8}{513.904}{158.619}{and}%
\TextAt{8}{79.226}{141.573}{chap.}%
\TextAt{8}{112.868}{141.573}{IV,}%
\TextAt{8}{133.496}{141.573}{(6.4),}%
\TextAt{8}{165.327}{141.573}{we}%
\TextAt{8}{183.038}{141.573}{may}%
\TextAt{8}{208.89}{141.573}{even}%
\TextAt{8}{236.034}{141.573}{assume}%
\TextAt{8}{276.961}{141.573}{that}%
\TextAt{13}{302.814}{141.573}{L}%
\TextAt{9}{314.923}{141.573}{\Character{61}}%
\TextAt{13}{328.172}{141.573}{L}%
\TextAt{21}{336.137}{139.78}{m}%
\TextAt{8}{344.125}{141.573}{.}%
\TextAt{8}{354.024}{141.573}{By}%
\TextAt{8}{372.871}{141.573}{(6.1),}%
\TextAt{8}{404.702}{141.573}{the}%
\TextAt{8}{425.342}{141.573}{norm}%
\TextAt{8}{456.395}{141.573}{residue}%
\TextAt{8}{496.605}{141.573}{symbol}%
\ErrataAt{79.19}{124.514}{455.81}{maps $U_K$ onto $G^0(L\mid K)=G(L\mid K)$. For $1\le n\le m$, it maps $U_K^{(n)}$ onto}%
\TextAt{13}{172.612}{95.526}{G}%
\TextAt{9}{181.846}{95.526}{\Character{40}}%
\TextAt{13}{186.398}{95.526}{L}%
\TextAt{21}{194.363}{93.733}{m}%
\TextAt{16}{205.672}{95.526}{\Character{106}}%
\TextAt{13}{212.314}{95.526}{L}%
\TextAt{21}{220.278}{93.733}{n}%
\TextAt{9}{225.915}{95.526}{\Character{41}}%
\TextAt{9}{233.77}{95.526}{\Character{61}}%
\TextAt{13}{246.213}{95.526}{G}%
\TextAt{21}{255.447}{93.733}{i}%
\TextAt{9}{258.828}{95.526}{\Character{40}}%
\TextAt{13}{263.381}{95.526}{L}%
\TextAt{21}{271.345}{93.733}{m}%
\TextAt{16}{282.655}{95.526}{\Character{106}}%
\TextAt{13}{289.296}{95.526}{K}%
\TextAt{9}{300.114}{95.526}{\Character{41}}%
\TextAt{8}{328.577}{95.526}{for}%
\TextAt{13}{346.462}{95.526}{q}%
\TextAt{21}{352.081}{100.462}{n}%
\TextAt{3}{357.219}{100.462}{\Character{0}}%
\TextAt{17}{363.805}{100.462}{1}%
\NativeMathAt{371.859}{95.526}{11.955}{\leq}%
\TextAt{13}{384.478}{95.526}{i}%
\NativeMathAt{391.792}{95.526}{11.955}{\leq}%
\TextAt{13}{404.411}{95.526}{q}%
\TextAt{21}{410.031}{100.462}{n}%
\TextAt{16}{418.324}{95.526}{\Character{0}}%
\TextAt{9}{430.279}{95.526}{1}%
\TextAt{13}{436.132}{95.526}{\Character{58}}%
\TextAt{8}{297.211}{43.033}{324}%
\Destination{78.23}{809.815}{nav.page.324}%
\Destination{79.36}{472.955}{nav.print.354}%
\Destination{79.36}{445.955}{nav.number.5.6.2}%
\Hotlink{133.772}{569.955}{9.751}{11.047}{nav.chapter.2}%
\Hotlink{156.486}{569.955}{12.716}{11.047}{nav.section.2.10}%
\Hotlink{458.426}{302.655}{9.763}{11.047}{nav.chapter.2}%
\Hotlink{210.42}{285.555}{13.995}{11.047}{nav.chapter.4}%
\Hotlink{316.746}{285.555}{9.775}{11.047}{nav.chapter.5}%
\Hotlink{458.332}{156.455}{9.751}{11.047}{nav.chapter.2}%
\Hotlink{112.368}{139.355}{13.995}{11.047}{nav.chapter.4}%
\Hotlink{473.872}{302.655}{30.936}{11.047}{nav.number.2.10.8}%
\Hotlink{230.565}{285.555}{25.078}{11.047}{nav.number.4.6.4}%
\Hotlink{332.67}{285.555}{25.078}{11.047}{nav.number.5.1.2}%
\Hotlink{365.054}{173.455}{25.078}{11.047}{nav.number.5.5.6}%
\Hotlink{475.166}{156.455}{30.936}{11.047}{nav.number.2.10.9}%
\Hotlink{132.996}{139.355}{25.078}{11.047}{nav.number.4.6.4}%
\Destination{79.226}{141.514}{errata.123}%
\Hotlink{374.032}{138.582}{21.722}{11.047}{nav.number.5.6.1}%
}

% Page 337 (body page 325)
\Page{%
\TextAt{2}{79.226}{796.943}{V.}%
\TextAt{2}{94.073}{796.943}{Local}%
\TextAt{2}{123.615}{796.943}{Class}%
\TextAt{2}{152.208}{796.943}{Field}%
\TextAt{2}{179.939}{796.943}{Theory}%
\Vector{q 0 G 0 g 0.299 w 7.226 18.729 m 460.769 18.729 l S Q}%
\TextAt{8}{79.226}{761.124}{But}%
\TextAt{8}{102.467}{761.124}{we}%
\TextAt{8}{119.694}{761.124}{have}%
\TextAt{8}{146.689}{761.124}{(see}%
\TextAt{8}{170.157}{761.124}{chap.}%
\TextAt{8}{202.348}{761.124}{II,}%
\TextAt{8}{218.261}{761.124}{§10)}%
\TextAt{13}{216.145}{729.379}{\Character{17}}%
\TextAt{21}{221.921}{727.52}{L}%
\TextAt{3}{227.68}{727.52}{\Character{106}}%
\TextAt{21}{230.032}{727.52}{K}%
\TextAt{9}{238.235}{729.379}{\Character{40}}%
\TextAt{13}{242.787}{729.379}{q}%
\TextAt{21}{248.406}{734.316}{n}%
\TextAt{16}{256.699}{729.379}{\Character{0}}%
\TextAt{9}{268.654}{729.379}{1\Character{41}}%
\TextAt{9}{282.355}{729.379}{\Character{61}}%
\TextAt{9}{298.244}{737.467}{1}%
\Vector{q 0 G 0 g 0.478 w 224.002 -41.632 m 234.339 -41.632 l S Q}%
\TextAt{13}{296.002}{721.179}{g}%
\TextAt{17}{301.607}{719.385}{0}%
\TextAt{9}{307.535}{729.379}{\Character{40}}%
\TextAt{13}{312.087}{729.379}{g}%
\TextAt{17}{317.692}{727.586}{1}%
\TextAt{9}{325.081}{729.379}{\Character{43}}%
\TextAt{16}{336.843}{729.379}{\Character{1}}%
\TextAt{16}{342.139}{729.379}{\Character{1}}%
\TextAt{16}{347.447}{729.379}{\Character{1}}%
\TextAt{9}{353.447}{729.379}{\Character{43}}%
\TextAt{13}{365.208}{729.379}{g}%
\TextAt{21}{370.814}{727.586}{q}%
\TextAt{22}{374.88}{729.855}{n}%
\TextAt{3}{379.979}{727.586}{\Character{0}}%
\TextAt{17}{386.566}{727.586}{1}%
\TextAt{9}{391.298}{729.379}{\Character{41}}%
\ErrataAt{79.19}{696.428}{455.81}{Here $g_0=q^{m-1}(q-1)$ and $g_i=\#G(L_m\mid L_j)=q^{m-j}$ for $q^{j-1}\le i\le q^j-1$ and $1\le j\le m$.}%
\ErrataAt{79.19}{679.382}{455.81}{Each such interval contributes $1$ to $\eta$. Thus $\eta(q^n-1)=n$ and $(U_K^{(n)},L\mid K)=G^n(L\mid K)$.}%
\ErrataAt{79.19}{662.6}{455.81}{For $n\ge m$ both groups are trivial.}%
\ErrataAt{524.35}{662.6}{8.65}{$\square$}%
\TextAt{8}{97.159}{630.346}{Higher}%
\TextAt{8}{136.889}{630.346}{ramification}%
\TextAt{8}{203.625}{630.346}{groups}%
\TextAt{13}{243.103}{630.346}{G}%
\TextAt{21}{252.337}{634.684}{t}%
\TextAt{9}{255.893}{630.346}{\Character{40}}%
\TextAt{13}{260.446}{630.346}{L}%
\TextAt{16}{274.608}{630.346}{\Character{106}}%
\TextAt{13}{284.127}{630.346}{K}%
\TextAt{9}{294.944}{630.346}{\Character{41}}%
\TextAt{8}{305.082}{630.346}{were}%
\TextAt{8}{333.753}{630.346}{introduced}%
\TextAt{8}{393.974}{630.346}{for}%
\TextAt{8}{413.547}{630.346}{arbitrary}%
\TextAt{8}{465.005}{630.346}{real}%
\TextAt{8}{489.456}{630.346}{numbers}%
\TextAt{13}{79.226}{613.3}{t}%
\NativeMathAt{87.87}{613.3}{11.955}{\geq}%
\TextAt{16}{101.571}{613.3}{\Character{0}}%
\TextAt{9}{110.883}{613.3}{1}%
\TextAt{8}{116.736}{613.3}{.}%
\TextAt{8}{127.114}{613.3}{Thus}%
\TextAt{8}{157.406}{613.3}{we}%
\TextAt{8}{175.276}{613.3}{may}%
\TextAt{8}{201.288}{613.3}{ask}%
\TextAt{8}{222.482}{613.3}{for}%
\TextAt{8}{241.01}{613.3}{which}%
\TextAt{8}{275.139}{613.3}{numbers}%
\TextAt{8}{322.993}{613.3}{they}%
\TextAt{8}{349.973}{613.3}{change.}%
\TextAt{8}{395.152}{613.3}{We}%
\TextAt{8}{415.939}{613.3}{call}%
\TextAt{8}{438.042}{613.3}{these}%
\TextAt{8}{468.657}{613.3}{numbers}%
\TextAt{8}{516.51}{613.3}{the}%
\TextAt{12}{79.226}{596.254}{jumps}%
\TextAt{8}{119.071}{596.254}{of}%
\TextAt{8}{132.771}{596.254}{the}%
\TextAt{8}{153.296}{596.254}{filtration}%
\TextAt{16}{202.454}{596.254}{\Character{102}}%
\TextAt{13}{208.432}{596.254}{G}%
\TextAt{21}{217.666}{600.593}{t}%
\TextAt{9}{221.222}{596.254}{\Character{40}}%
\TextAt{13}{225.774}{596.254}{L}%
\TextAt{16}{237.689}{596.254}{\Character{106}}%
\TextAt{13}{244.959}{596.254}{K}%
\TextAt{9}{255.777}{596.254}{\Character{41}}%
\TextAt{16}{260.329}{596.254}{\Character{103}}%
\TextAt{21}{266.307}{594.461}{t}%
\NativeMathAt{269.365}{594.461}{7.97}{\geq}%
\TextAt{3}{275.948}{594.461}{\Character{0}}%
\TextAt{17}{282.538}{594.461}{1}%
\TextAt{8}{291.537}{596.254}{of}%
\TextAt{13}{305.236}{596.254}{G}%
\TextAt{9}{314.47}{596.254}{\Character{40}}%
\TextAt{13}{319.022}{596.254}{L}%
\TextAt{16}{330.937}{596.254}{\Character{106}}%
\TextAt{13}{338.207}{596.254}{K}%
\TextAt{9}{349.025}{596.254}{\Character{41}}%
\TextAt{8}{353.577}{596.254}{.}%
\TextAt{8}{363.133}{596.254}{In}%
\TextAt{8}{378.124}{596.254}{other}%
\TextAt{8}{409.062}{596.254}{words,}%
\TextAt{13}{446.334}{596.254}{t}%
\TextAt{8}{454.827}{596.254}{is}%
\TextAt{8}{466.961}{596.254}{a}%
\TextAt{8}{477.085}{596.254}{jump}%
\TextAt{8}{507.689}{596.254}{if}%
\TextAt{8}{518.782}{596.254}{for}%
\TextAt{8}{79.226}{579.208}{all}%
\TextAt{13}{95.485}{579.208}{\Character{34}}%
\TextAt{13}{104.272}{579.208}{\Character{62}}%
\TextAt{9}{116.71}{579.208}{0}%
\TextAt{8}{122.563}{579.208}{,}%
\TextAt{8}{129.712}{579.208}{one}%
\TextAt{8}{151.171}{579.208}{has}%
\TextAt{13}{240.611}{562.162}{G}%
\TextAt{21}{249.845}{567.098}{t}%
\TextAt{9}{253.401}{562.162}{\Character{40}}%
\TextAt{13}{257.953}{562.162}{L}%
\TextAt{16}{269.239}{562.162}{\Character{106}}%
\TextAt{13}{275.88}{562.162}{K}%
\TextAt{9}{286.698}{562.162}{\Character{41}}%
\NativeMathAt{294.571}{562.162}{11.955}{\neq}%
\TextAt{13}{306.997}{562.162}{G}%
\TextAt{21}{316.23}{567.098}{t}%
\TextAt{17}{319.288}{567.098}{\Character{43}}%
\TextAt{21}{325.875}{567.098}{\Character{34}}%
\TextAt{9}{330.283}{562.162}{\Character{40}}%
\TextAt{13}{334.836}{562.162}{L}%
\TextAt{16}{346.121}{562.162}{\Character{106}}%
\TextAt{13}{352.763}{562.162}{K}%
\TextAt{9}{363.58}{562.162}{\Character{41}}%
\TextAt{13}{368.133}{562.162}{\Character{58}}%
\TextAt{12}{79.226}{529.674}{Proposition}%
\TextAt{12}{152.902}{529.674}{(6.3)}%
\TextAt{12}{186.027}{529.674}{(}%
\TextAt{26}{191.263}{529.674}{Hasse–Arf}%
\TextAt{12}{255.893}{529.674}{).}%
\TextAt{8}{269.607}{529.674}{For}%
\TextAt{8}{291.416}{529.674}{a}%
\TextAt{8}{302.01}{529.674}{finite}%
\TextAt{8}{332.761}{529.674}{abelian}%
\TextAt{8}{374.247}{529.674}{extension}%
\TextAt{13}{426.852}{529.674}{L}%
\TextAt{16}{439.567}{529.674}{\Character{106}}%
\TextAt{13}{447.638}{529.674}{K}%
\TextAt{8}{458.456}{529.674}{,}%
\TextAt{8}{466.653}{529.674}{the}%
\TextAt{8}{487.649}{529.674}{jumps}%
\TextAt{8}{523.337}{529.674}{of}%
\TextAt{8}{79.226}{512.628}{the}%
\TextAt{8}{99.383}{512.628}{filtration}%
\TextAt{16}{148.172}{512.628}{\Character{102}}%
\TextAt{13}{154.149}{512.628}{G}%
\TextAt{21}{163.383}{516.967}{t}%
\TextAt{9}{166.939}{512.628}{\Character{40}}%
\TextAt{13}{171.491}{512.628}{L}%
\TextAt{16}{182.777}{512.628}{\Character{106}}%
\TextAt{13}{189.418}{512.628}{K}%
\TextAt{9}{200.236}{512.628}{\Character{41}}%
\TextAt{16}{204.788}{512.628}{\Character{103}}%
\TextAt{21}{210.766}{510.835}{t}%
\NativeMathAt{213.824}{510.835}{7.97}{\geq}%
\TextAt{3}{220.407}{510.835}{\Character{0}}%
\TextAt{17}{226.997}{510.835}{1}%
\TextAt{8}{235.627}{512.628}{of}%
\TextAt{13}{248.957}{512.628}{G}%
\TextAt{9}{258.19}{512.628}{\Character{40}}%
\TextAt{13}{262.743}{512.628}{L}%
\TextAt{16}{274.028}{512.628}{\Character{106}}%
\TextAt{13}{280.67}{512.628}{K}%
\TextAt{9}{291.487}{512.628}{\Character{41}}%
\TextAt{8}{299.937}{512.628}{are}%
\TextAt{8}{319.448}{512.628}{rational}%
\TextAt{8}{363.037}{512.628}{integers.}%
\ErrataAt{79.19}{477.61}{455.81}{\textbf{Proof.} The possible jump at $-1$ is integral. For $t\ge0$ one has $G^t(L\mid K)=G^t(L\mid L^0)$, so}%
\ErrataAt{79.19}{460.564}{455.81}{as in \hyperlink{nav.number.5.6.2}{(6.2)} we may assume $L\mid K$ is totally ramified and contained in $L_m\mid K$. If $t$}%
\TextAt{8}{79.226}{443.557}{is}%
\TextAt{8}{91.912}{443.557}{a}%
\TextAt{8}{102.589}{443.557}{jump}%
\TextAt{8}{133.746}{443.557}{of}%
\TextAt{16}{147.997}{443.557}{\Character{102}}%
\TextAt{13}{153.975}{443.557}{G}%
\TextAt{21}{163.209}{447.896}{t}%
\TextAt{9}{166.765}{443.557}{\Character{40}}%
\TextAt{13}{171.317}{443.557}{L}%
\TextAt{16}{184.173}{443.557}{\Character{106}}%
\TextAt{13}{192.386}{443.557}{K}%
\TextAt{9}{203.204}{443.557}{\Character{41}}%
\TextAt{16}{207.756}{443.557}{\Character{103}}%
\TextAt{8}{213.734}{443.557}{,}%
\TextAt{8}{222.035}{443.557}{then}%
\TextAt{8}{249.617}{443.557}{by}%
\TextAt{8}{266.798}{443.557}{chap.}%
\TextAt{8}{301.755}{443.557}{II}%
\TextAt{8}{315.336}{443.557}{(10.9),}%
\TextAt{13}{353.573}{443.557}{t}%
\TextAt{8}{362.62}{443.557}{is}%
\TextAt{8}{375.305}{443.557}{also}%
\TextAt{8}{399.707}{443.557}{a}%
\TextAt{8}{410.384}{443.557}{jump}%
\TextAt{8}{441.541}{443.557}{of}%
\TextAt{16}{455.793}{443.557}{\Character{102}}%
\TextAt{13}{461.77}{443.557}{G}%
\TextAt{21}{471.004}{447.896}{t}%
\TextAt{9}{474.56}{443.557}{\Character{40}}%
\TextAt{13}{479.113}{443.557}{L}%
\TextAt{21}{487.077}{441.764}{m}%
\TextAt{16}{499.957}{443.557}{\Character{106}}%
\TextAt{13}{508.17}{443.557}{K}%
\TextAt{9}{518.988}{443.557}{\Character{41}}%
\TextAt{16}{523.54}{443.557}{\Character{103}}%
\TextAt{8}{529.518}{443.557}{.}%
\TextAt{8}{79.226}{426.511}{Since}%
\TextAt{8}{109.905}{426.511}{by}%
\TextAt{8}{126.286}{426.511}{(6.1),}%
\TextAt{8}{157.665}{426.511}{the}%
\TextAt{8}{177.943}{426.511}{jumps}%
\TextAt{8}{212.914}{426.511}{of}%
\TextAt{16}{226.366}{426.511}{\Character{102}}%
\TextAt{13}{232.343}{426.511}{G}%
\TextAt{21}{241.577}{424.718}{s}%
\TextAt{9}{245.991}{426.511}{\Character{40}}%
\TextAt{13}{250.543}{426.511}{L}%
\TextAt{21}{258.508}{424.718}{m}%
\TextAt{16}{270.024}{426.511}{\Character{106}}%
\TextAt{13}{276.873}{426.511}{K}%
\TextAt{9}{287.691}{426.511}{\Character{41}}%
\TextAt{16}{292.243}{426.511}{\Character{103}}%
\TextAt{8}{302.24}{426.511}{are}%
\TextAt{8}{321.873}{426.511}{the}%
\TextAt{8}{342.151}{426.511}{numbers}%
\TextAt{13}{389.483}{426.511}{q}%
\TextAt{21}{395.102}{430.85}{n}%
\TextAt{16}{403.478}{426.511}{\Character{0}}%
\TextAt{9}{415.516}{426.511}{1}%
\TextAt{8}{421.369}{426.511}{,}%
\TextAt{8}{428.671}{426.511}{for}%
\TextAt{13}{446.677}{426.511}{n}%
\TextAt{9}{457.193}{426.511}{\Character{61}}%
\TextAt{9}{469.818}{426.511}{0}%
\TextAt{13}{475.679}{426.511}{\Character{59}}%
\TextAt{13}{480.903}{426.511}{\Character{58}}%
\TextAt{13}{486.14}{426.511}{\Character{58}}%
\TextAt{13}{491.364}{426.511}{\Character{58}}%
\TextAt{13}{496.601}{426.511}{\Character{59}}%
\TextAt{13}{501.837}{426.511}{m}%
\TextAt{16}{514.878}{426.511}{\Character{0}}%
\TextAt{9}{526.917}{426.511}{1}%
\TextAt{8}{79.226}{409.466}{(}%
\TextAt{13}{83.781}{409.466}{q}%
\TextAt{9}{93.979}{409.466}{\Character{61}}%
\TextAt{9}{107.644}{409.466}{2}%
\TextAt{8}{118.151}{409.466}{is}%
\TextAt{8}{130.653}{409.466}{an}%
\TextAt{8}{147.65}{409.466}{exception:}%
\TextAt{9}{206.028}{409.466}{0}%
\TextAt{8}{216.517}{409.466}{is}%
\TextAt{8}{229.019}{409.466}{not}%
\TextAt{8}{250.571}{409.466}{a}%
\TextAt{8}{261.064}{409.466}{jump),}%
\TextAt{8}{300.028}{409.466}{the}%
\TextAt{8}{320.923}{409.466}{jumps}%
\TextAt{8}{356.51}{409.466}{of}%
\TextAt{16}{370.578}{409.466}{\Character{102}}%
\TextAt{13}{376.556}{409.466}{G}%
\TextAt{21}{385.79}{413.804}{t}%
\TextAt{9}{389.346}{409.466}{\Character{40}}%
\TextAt{13}{393.898}{409.466}{L}%
\TextAt{21}{401.863}{407.672}{m}%
\TextAt{16}{414.43}{409.466}{\Character{106}}%
\TextAt{13}{422.329}{409.466}{K}%
\TextAt{9}{433.147}{409.466}{\Character{41}}%
\TextAt{16}{437.699}{409.466}{\Character{103}}%
\TextAt{8}{448.312}{409.466}{are}%
\TextAt{8}{468.561}{409.466}{the}%
\TextAt{8}{489.456}{409.466}{numbers}%
\TextAt{13}{79.226}{392.42}{\Character{17}}%
\TextAt{21}{85.002}{390.56}{L}%
\TextAt{22}{90.761}{389.564}{m}%
\TextAt{3}{97.797}{390.56}{\Character{106}}%
\TextAt{21}{100.149}{390.56}{K}%
\TextAt{9}{108.352}{392.42}{\Character{40}}%
\TextAt{13}{112.904}{392.42}{q}%
\TextAt{21}{118.523}{396.758}{n}%
\TextAt{16}{126.816}{392.42}{\Character{0}}%
\TextAt{9}{138.771}{392.42}{1\Character{41}}%
\TextAt{9}{152.472}{392.42}{\Character{61}}%
\TextAt{13}{164.923}{392.42}{n}%
\TextAt{8}{171.911}{392.42}{,}%
\TextAt{8}{179.06}{392.42}{for}%
\TextAt{13}{196.945}{392.42}{n}%
\TextAt{9}{207.253}{392.42}{\Character{61}}%
\TextAt{9}{219.663}{392.42}{0}%
\TextAt{13}{225.532}{392.42}{\Character{59}}%
\TextAt{13}{230.756}{392.42}{\Character{58}}%
\TextAt{13}{235.993}{392.42}{\Character{58}}%
\TextAt{13}{241.217}{392.42}{\Character{58}}%
\TextAt{13}{246.454}{392.42}{\Character{59}}%
\TextAt{13}{251.69}{392.42}{m}%
\TextAt{16}{264.648}{392.42}{\Character{0}}%
\TextAt{9}{276.604}{392.42}{1}%
\TextAt{8}{282.457}{392.42}{.}%
\Vector{q 0 G 0 g 0.399 w 452.334 -381.58 m 452.334 -373.511 l S Q}%
\Vector{q 0 G 0 g 0.399 w 452.534 -373.71 m 459.707 -373.71 l S Q}%
\Vector{q 0 G 0 g 0.399 w 452.534 -381.381 m 459.707 -381.381 l S Q}%
\Vector{q 0 G 0 g 0.399 w 459.906 -381.58 m 459.906 -373.511 l S Q}%
\TextAt{8}{97.159}{363.418}{The}%
\TextAt{8}{122}{363.418}{theorem}%
\TextAt{8}{168.312}{363.418}{of}%
\TextAt{26}{182.944}{363.418}{Hasse–Arf}%
\TextAt{8}{252.258}{363.418}{has}%
\TextAt{8}{273.919}{363.418}{an}%
\TextAt{8}{290.965}{363.418}{important}%
\TextAt{8}{347.045}{363.418}{application}%
\TextAt{8}{408.326}{363.418}{to}%
\TextAt{8}{423.423}{363.418}{Artin}%
\TextAt{13}{455.748}{363.418}{L}%
\TextAt{8}{463.713}{363.418}{-series,}%
\TextAt{8}{503.18}{363.418}{which}%
\TextAt{8}{79.226}{346.372}{we}%
\TextAt{8}{96.454}{346.372}{will}%
\TextAt{8}{118.559}{346.372}{study}%
\TextAt{8}{150.814}{346.372}{in}%
\TextAt{8}{164.467}{346.372}{chap.}%
\TextAt{8}{196.658}{346.372}{VII}%
\TextAt{8}{218.093}{346.372}{(see}%
\TextAt{8}{241.561}{346.372}{chap.}%
\TextAt{8}{273.753}{346.372}{VII,}%
\TextAt{8}{298.44}{346.372}{(11.4)).}%
\TextAt{8}{297.211}{43.033}{325}%
\Destination{78.23}{809.815}{nav.page.325}%
\Destination{79.36}{611.255}{nav.print.355}%
\Destination{79.36}{543.655}{nav.number.5.6.3}%
\Hotlink{201.848}{758.955}{9.763}{11.047}{nav.chapter.2}%
\Hotlink{223.44}{758.955}{12.716}{11.047}{nav.section.2.10}%
\Hotlink{301.255}{441.355}{9.751}{11.047}{nav.chapter.2}%
\Hotlink{196.158}{344.155}{18.526}{11.047}{nav.chapter.7}%
\Hotlink{273.253}{344.155}{18.526}{11.047}{nav.chapter.7}%
\Hotlink{216.423}{475.455}{25.078}{11.047}{nav.number.5.6.2}%
\Hotlink{314.836}{441.355}{30.936}{11.047}{nav.number.2.10.9}%
\Hotlink{125.786}{424.355}{25.078}{11.047}{nav.number.5.6.1}%
\Hotlink{297.94}{344.155}{30.936}{11.047}{nav.number.7.11.4}%
\Hotlink{248.226}{70.755}{21.483}{9.572}{nav.number.5.1.5}%
\Destination{79.226}{713.428}{errata.124}%
\Destination{79.226}{494.61}{errata.125}%
}


% Chapter 6
% Page 338 (body page 326)
\Page{%
\TextAt{1}{79.226}{673.054}{Chapter}%
\TextAt{1}{188.19}{673.054}{VI}%
\TextAt{1}{79.226}{614.873}{Global}%
\TextAt{1}{168.955}{614.873}{Class}%
\TextAt{1}{241.73}{614.873}{Field}%
\TextAt{1}{312.274}{614.873}{Theory}%
\TextAt{15}{79.226}{546.17}{§}%
\TextAt{15}{93.785}{546.17}{1}%
\TextAt{15}{120.693}{546.17}{Idèles}%
\TextAt{15}{175.817}{546.17}{and}%
\TextAt{15}{213.209}{546.17}{Idèle}%
\TextAt{15}{260.689}{546.17}{Classes}%
\TextAt{8}{79.226}{514.284}{The}%
\TextAt{8}{102.248}{514.284}{role}%
\TextAt{8}{123.979}{514.284}{held}%
\TextAt{8}{148.305}{514.284}{in}%
\TextAt{8}{160.926}{514.284}{local}%
\TextAt{8}{187.534}{514.284}{class}%
\TextAt{8}{213.94}{514.284}{field}%
\TextAt{8}{238.265}{514.284}{theory}%
\TextAt{8}{273.984}{514.284}{by}%
\TextAt{8}{289.211}{514.284}{the}%
\TextAt{8}{308.336}{514.284}{multiplicative}%
\TextAt{8}{381.127}{514.284}{group}%
\TextAt{8}{413.271}{514.284}{of}%
\TextAt{8}{425.57}{514.284}{the}%
\TextAt{8}{444.694}{514.284}{base}%
\TextAt{8}{469.737}{514.284}{field}%
\TextAt{8}{494.062}{514.284}{is}%
\TextAt{8}{504.794}{514.284}{taken}%
\TextAt{8}{79.226}{497.238}{in}%
\TextAt{8}{92.119}{497.238}{global}%
\TextAt{8}{125.839}{497.238}{class}%
\TextAt{8}{152.516}{497.238}{field}%
\TextAt{8}{177.114}{497.238}{theory}%
\TextAt{8}{213.104}{497.238}{by}%
\TextAt{8}{228.604}{497.238}{the}%
\TextAt{8}{248.001}{497.238}{idèle}%
\TextAt{8}{274.547}{497.238}{class}%
\TextAt{8}{301.224}{497.238}{group.}%
\TextAt{8}{338.698}{497.238}{The}%
\TextAt{8}{361.992}{497.238}{notion}%
\TextAt{8}{397.66}{497.238}{of}%
\TextAt{8}{410.23}{497.238}{idèle}%
\TextAt{8}{436.777}{497.238}{is}%
\TextAt{8}{447.781}{497.238}{a}%
\TextAt{8}{456.777}{497.238}{modification}%
\TextAt{8}{523.337}{497.238}{of}%
\TextAt{8}{79.226}{480.192}{the}%
\TextAt{8}{98.597}{480.192}{notion}%
\TextAt{8}{134.238}{480.192}{of}%
\TextAt{8}{146.782}{480.192}{ideal.}%
\TextAt{8}{179.034}{480.192}{It}%
\TextAt{8}{190.92}{480.192}{was}%
\TextAt{8}{212.634}{480.192}{introduced}%
\TextAt{8}{270.38}{480.192}{by}%
\TextAt{8}{285.853}{480.192}{the}%
\TextAt{8}{305.224}{480.192}{French}%
\TextAt{8}{342.634}{480.192}{mathematician}%
\TextAt{26}{421.852}{480.192}{Claude}%
\TextAt{26}{469.515}{480.192}{Chevalley}%
\TextAt{8}{79.226}{463.147}{(1909–1984)}%
\TextAt{8}{145.795}{463.147}{with}%
\TextAt{8}{173.293}{463.147}{a}%
\TextAt{8}{183.888}{463.147}{view}%
\TextAt{8}{211.709}{463.147}{to}%
\TextAt{8}{226.859}{463.147}{providing}%
\TextAt{8}{279.738}{463.147}{a}%
\TextAt{8}{290.333}{463.147}{suitable}%
\TextAt{8}{334.808}{463.147}{basis}%
\TextAt{8}{364.387}{463.147}{for}%
\TextAt{8}{383.111}{463.147}{the}%
\TextAt{8}{404.106}{463.147}{important}%
\TextAt{8}{460.237}{463.147}{local-to-global}%
\TextAt{8}{79.226}{446.101}{principle,}%
\TextAt{8}{131.267}{446.101}{i.e.,}%
\TextAt{8}{154.042}{446.101}{for}%
\TextAt{8}{172.463}{446.101}{the}%
\TextAt{8}{193.156}{446.101}{principle}%
\TextAt{8}{241.811}{446.101}{which}%
\TextAt{8}{275.833}{446.101}{reduces}%
\TextAt{8}{318.045}{446.101}{problems}%
\TextAt{8}{368.721}{446.101}{concerning}%
\TextAt{8}{427.789}{446.101}{a}%
\TextAt{8}{438.081}{446.101}{number}%
\TextAt{8}{481.213}{446.101}{field}%
\TextAt{13}{507.106}{446.101}{K}%
\TextAt{8}{522.357}{446.101}{to}%
\TextAt{8}{79.226}{429.055}{analogous}%
\TextAt{8}{134.234}{429.055}{problems}%
\TextAt{8}{185.32}{429.055}{for}%
\TextAt{8}{204.151}{429.055}{the}%
\TextAt{8}{225.254}{429.055}{various}%
\TextAt{8}{266.274}{429.055}{completions}%
\TextAt{13}{331.671}{429.055}{K}%
\TextAt{27}{341.63}{427.261}{p}%
\TextAt{8}{346.116}{429.055}{.}%
\TextAt{26}{357.404}{429.055}{Chevalley}%
\TextAt{8}{425.502}{429.055}{used}%
\TextAt{8}{453.168}{429.055}{the}%
\TextAt{8}{474.271}{429.055}{term}%
\TextAt{8}{503.18}{429.055}{“ideal}%
\TextAt{8}{79.226}{412.009}{element”,}%
\TextAt{8}{131.243}{412.009}{which}%
\TextAt{8}{164.73}{412.009}{was}%
\TextAt{8}{187.229}{412.009}{abbreviated}%
\TextAt{8}{251.297}{412.009}{as}%
\TextAt{8}{265.667}{412.009}{id.}%
\TextAt{8}{283.871}{412.009}{el.}%
\TextAt{8}{97.159}{391.974}{An}%
\TextAt{8}{116.766}{391.974}{adèle}%
\TextAt{8}{147.108}{391.974}{of}%
\TextAt{13}{160.868}{391.974}{K}%
\TextAt{8}{171.686}{391.974}{—this}%
\TextAt{8}{206.643}{391.974}{curious}%
\TextAt{8}{247.458}{391.974}{expression,}%
\TextAt{8}{307.629}{391.974}{which}%
\TextAt{8}{341.546}{391.974}{has}%
\TextAt{8}{362.85}{391.974}{the}%
\TextAt{8}{383.437}{391.974}{stress}%
\TextAt{8}{415.919}{391.974}{on}%
\TextAt{8}{432.609}{391.974}{the}%
\TextAt{8}{453.196}{391.974}{second}%
\TextAt{8}{491.405}{391.974}{syllable,}%
\TextAt{8}{79.226}{374.928}{is}%
\TextAt{8}{90.99}{374.928}{derived}%
\TextAt{8}{131.96}{374.928}{from}%
\TextAt{8}{159.601}{374.928}{the}%
\TextAt{8}{179.757}{374.928}{original}%
\TextAt{8}{222.043}{374.928}{term}%
\TextAt{8}{250.006}{374.928}{“additive}%
\TextAt{8}{300.409}{374.928}{idèle”—is}%
\TextAt{8}{352.808}{374.928}{a}%
\TextAt{8}{362.564}{374.928}{family}%
\TextAt{13}{283.845}{347.674}{\Character{11}}%
\TextAt{9}{294.688}{347.674}{\Character{61}}%
\TextAt{9}{307.098}{347.674}{\Character{40}}%
\TextAt{13}{311.666}{347.674}{\Character{11}}%
\TextAt{27}{319.111}{345.88}{p}%
\TextAt{9}{323.598}{347.674}{\Character{41}}%
\TextAt{8}{79.226}{320.419}{of}%
\TextAt{8}{93.001}{320.419}{elements}%
\TextAt{13}{141.302}{320.419}{\Character{11}}%
\TextAt{27}{148.748}{318.626}{p}%
\NativeMathAt{157.314}{320.419}{11.955}{\in}%
\TextAt{13}{169.363}{320.419}{K}%
\TextAt{27}{179.322}{318.626}{p}%
\TextAt{8}{188.151}{320.419}{where}%
\TextAt{14}{222.405}{320.419}{p}%
\TextAt{8}{232.73}{320.419}{runs}%
\TextAt{8}{259.249}{320.419}{through}%
\TextAt{8}{303.928}{320.419}{all}%
\TextAt{8}{320.632}{320.419}{primes}%
\TextAt{8}{358.855}{320.419}{of}%
\TextAt{13}{372.63}{320.419}{K}%
\TextAt{8}{383.447}{320.419}{,}%
\TextAt{8}{391.152}{320.419}{and}%
\TextAt{13}{414.36}{320.419}{\Character{11}}%
\TextAt{27}{421.805}{318.626}{p}%
\TextAt{8}{430.635}{320.419}{is}%
\TextAt{8}{442.843}{320.419}{integral}%
\TextAt{8}{485.896}{320.419}{in}%
\TextAt{13}{499.994}{320.419}{K}%
\TextAt{27}{509.953}{318.626}{p}%
\TextAt{8}{518.782}{320.419}{for}%
\TextAt{8}{79.226}{303.373}{almost}%
\TextAt{8}{117.016}{303.373}{all}%
\TextAt{14}{133.275}{303.373}{p}%
\TextAt{8}{139.259}{303.373}{.}%
\TextAt{8}{147.707}{303.373}{The}%
\TextAt{8}{171.761}{303.373}{adèles}%
\TextAt{8}{206.287}{303.373}{form}%
\TextAt{8}{233.928}{303.373}{a}%
\TextAt{8}{243.683}{303.373}{ring,}%
\TextAt{8}{271.001}{303.373}{which}%
\TextAt{8}{304.487}{303.373}{is}%
\TextAt{8}{316.251}{303.373}{denoted}%
\TextAt{8}{360.473}{303.373}{by}%
\TextAt{18}{272.223}{269.15}{A}%
\TextAt{21}{280.857}{267.357}{K}%
\TextAt{9}{292.381}{269.15}{\Character{61}}%
\TextAt{3}{311.296}{284.094}{\Character{48}}%
\TextAt{20}{304.806}{280.508}{\Character{89}}%
\TextAt{27}{310.45}{255.402}{p}%
\TextAt{13}{322.075}{269.15}{K}%
\TextAt{27}{332.034}{267.357}{p}%
\TextAt{13}{336.521}{269.15}{\Character{58}}%
\TextAt{8}{79.226}{232.959}{Addition}%
\TextAt{8}{128.784}{232.959}{and}%
\TextAt{8}{152.328}{232.959}{multiplication}%
\TextAt{8}{228.235}{232.959}{are}%
\TextAt{8}{248.527}{232.959}{defined}%
\TextAt{8}{289.621}{232.959}{componentwise.}%
\TextAt{8}{377.548}{232.959}{This}%
\TextAt{8}{405.048}{232.959}{kind}%
\TextAt{8}{432.167}{232.959}{of}%
\TextAt{8}{446.278}{232.959}{product}%
\TextAt{8}{490.958}{232.959}{is}%
\TextAt{8}{503.503}{232.959}{called}%
\TextAt{8}{79.226}{215.913}{the}%
\TextAt{8}{99.383}{215.913}{“restricted}%
\TextAt{8}{156.995}{215.913}{product”}%
\TextAt{8}{206.417}{215.913}{of}%
\TextAt{8}{219.747}{215.913}{the}%
\TextAt{13}{239.904}{215.913}{K}%
\TextAt{27}{249.863}{214.12}{p}%
\TextAt{8}{258.247}{215.913}{with}%
\TextAt{8}{284.907}{215.913}{respect}%
\TextAt{8}{324.957}{215.913}{to}%
\TextAt{8}{339.267}{215.913}{the}%
\TextAt{8}{359.424}{215.913}{subrings}%
\TextAt{16}{405.726}{215.913}{\Character{79}}%
\TextAt{27}{415.244}{214.12}{p}%
\NativeMathAt{423.051}{215.913}{11.955}{\subseteq}%
\TextAt{13}{435.671}{215.913}{K}%
\TextAt{27}{445.63}{214.12}{p}%
\TextAt{8}{450.117}{215.913}{.}%
\TextAt{8}{97.159}{195.878}{The}%
\TextAt{8}{121.213}{195.878}{idèle}%
\TextAt{8}{148.518}{195.878}{group}%
\TextAt{8}{181.694}{195.878}{of}%
\TextAt{13}{195.024}{195.878}{K}%
\TextAt{8}{209.739}{195.878}{is}%
\TextAt{8}{221.503}{195.878}{defined}%
\TextAt{8}{261.816}{195.878}{to}%
\TextAt{8}{276.126}{195.878}{be}%
\TextAt{8}{292.05}{195.878}{the}%
\TextAt{8}{312.207}{195.878}{unit}%
\TextAt{8}{336.918}{195.878}{group}%
\TextAt{13}{281.398}{168.624}{I}%
\TextAt{21}{286.56}{166.831}{K}%
\TextAt{9}{298.084}{168.624}{\Character{61}}%
\TextAt{18}{310.509}{168.624}{A}%
\TextAt{3}{319.144}{173.56}{\Character{3}}%
\TextAt{21}{319.144}{165.668}{K}%
\TextAt{13}{327.346}{168.624}{\Character{58}}%
\TextAt{8}{79.226}{141.37}{Thus}%
\TextAt{8}{108.87}{141.37}{an}%
\TextAt{8}{125.124}{141.37}{idèle}%
\TextAt{8}{152.425}{141.37}{is}%
\TextAt{8}{164.184}{141.37}{a}%
\TextAt{8}{173.934}{141.37}{family}%
\TextAt{13}{209.699}{141.37}{\Character{11}}%
\TextAt{9}{220.542}{141.369}{\Character{61}}%
\TextAt{9}{232.952}{141.369}{\Character{40}}%
\TextAt{13}{237.519}{141.37}{\Character{11}}%
\TextAt{27}{244.965}{139.576}{p}%
\TextAt{9}{249.452}{141.369}{\Character{41}}%
\TextAt{8}{257.897}{141.37}{of}%
\TextAt{8}{271.222}{141.37}{elements}%
\TextAt{13}{319.073}{141.37}{\Character{11}}%
\TextAt{27}{326.519}{139.576}{p}%
\NativeMathAt{334.327}{141.369}{11.955}{\in}%
\TextAt{13}{345.618}{141.369}{K}%
\TextAt{3}{356.435}{145.708}{\Character{3}}%
\TextAt{27}{355.577}{138.414}{p}%
\TextAt{8}{365.06}{141.37}{where}%
\TextAt{13}{398.864}{141.37}{\Character{11}}%
\TextAt{27}{406.31}{139.576}{p}%
\TextAt{8}{414.689}{141.37}{is}%
\TextAt{8}{426.448}{141.37}{a}%
\TextAt{8}{436.198}{141.37}{unit}%
\TextAt{8}{460.905}{141.37}{in}%
\TextAt{8}{474.553}{141.37}{the}%
\TextAt{8}{494.704}{141.37}{ring}%
\TextAt{16}{518.765}{141.37}{\Character{79}}%
\TextAt{27}{528.283}{139.576}{p}%
\TextAt{8}{79.226}{124.324}{of}%
\TextAt{8}{92.979}{124.324}{integers}%
\TextAt{8}{136.716}{124.324}{of}%
\TextAt{13}{150.468}{124.324}{K}%
\TextAt{27}{160.428}{122.53}{p}%
\TextAt{8}{164.914}{124.324}{,}%
\TextAt{8}{172.592}{124.324}{for}%
\TextAt{8}{190.9}{124.324}{almost}%
\TextAt{8}{229.113}{124.324}{all}%
\TextAt{14}{245.795}{124.324}{p}%
\TextAt{8}{251.778}{124.324}{.}%
\TextAt{8}{261.495}{124.324}{In}%
\TextAt{8}{276.54}{124.324}{analogy}%
\TextAt{8}{320.228}{124.324}{with}%
\TextAt{18}{347.311}{124.324}{A}%
\TextAt{21}{355.945}{122.53}{K}%
\TextAt{8}{364.148}{124.324}{,}%
\TextAt{8}{371.826}{124.324}{we}%
\TextAt{8}{389.476}{124.324}{write}%
\TextAt{8}{419.811}{124.324}{the}%
\TextAt{8}{440.39}{124.324}{idèle}%
\TextAt{8}{468.119}{124.324}{group}%
\TextAt{8}{501.717}{124.324}{as}%
\TextAt{8}{516.51}{124.324}{the}%
\TextAt{8}{79.226}{107.278}{restricted}%
\TextAt{8}{131.315}{107.278}{product}%
\TextAt{13}{275.033}{83.395}{I}%
\TextAt{21}{280.195}{81.601}{K}%
\TextAt{9}{291.719}{83.395}{\Character{61}}%
\TextAt{3}{310.634}{98.339}{\Character{48}}%
\TextAt{20}{304.144}{94.752}{\Character{89}}%
\TextAt{27}{309.788}{69.646}{p}%
\TextAt{13}{321.413}{83.395}{K}%
\TextAt{3}{332.231}{88.331}{\Character{3}}%
\TextAt{27}{331.372}{80.439}{p}%
\TextAt{8}{297.211}{43.033}{326}%
\Destination{79.23}{773.075}{nav.chapter.6}%
\Destination{78.23}{809.815}{nav.page.326}%
\Destination{79.23}{570.065}{nav.section.6.1}%
\Destination{79.36}{629.855}{nav.print.357}%
\Destination{79.23}{773.075}{outline.55}%
\Destination{79.23}{570.065}{outline.56}%
}

% Page 339 (body page 327)
\Page{%
\TextAt{2}{79.226}{796.943}{VI.}%
\TextAt{2}{98.012}{796.943}{Global}%
\TextAt{2}{133.248}{796.943}{Class}%
\TextAt{2}{161.841}{796.943}{Field}%
\TextAt{2}{189.572}{796.943}{Theory}%
\Vector{q 0 G 0 g 0.299 w 7.226 18.729 m 460.769 18.729 l S Q}%
\TextAt{8}{79.226}{761.124}{with}%
\TextAt{8}{107.048}{761.124}{respect}%
\TextAt{8}{148.259}{761.124}{to}%
\TextAt{8}{163.731}{761.124}{the}%
\TextAt{8}{185.049}{761.124}{unit}%
\TextAt{8}{210.922}{761.124}{groups}%
\TextAt{16}{249.874}{761.124}{\Character{79}}%
\TextAt{3}{259.724}{765.462}{\Character{3}}%
\TextAt{27}{259.392}{758.168}{p}%
\TextAt{8}{264.456}{761.124}{.}%
\TextAt{8}{276.389}{761.124}{For}%
\TextAt{8}{298.52}{761.124}{every}%
\TextAt{8}{330.574}{761.124}{finite}%
\TextAt{8}{361.647}{761.124}{set}%
\TextAt{8}{381.076}{761.124}{of}%
\TextAt{8}{395.568}{761.124}{primes}%
\TextAt{13}{434.508}{761.124}{S}%
\TextAt{8}{442.403}{761.124}{,}%
\TextAt{13}{451.004}{761.124}{I}%
\TextAt{21}{456.167}{759.33}{K}%
\TextAt{8}{469.429}{761.124}{contains}%
\TextAt{8}{516.51}{761.124}{the}%
\TextAt{8}{79.226}{744.078}{subgroup}%
\TextAt{13}{252.874}{727.032}{I}%
\TextAt{21}{258.977}{731.968}{S}%
\TextAt{21}{258.037}{724.076}{K}%
\TextAt{9}{269.56}{727.032}{\Character{61}}%
\TextAt{20}{281.986}{738.389}{\Character{89}}%
\TextAt{27}{282.009}{713.018}{p}%
\NativeMathAt{285.997}{713.018}{7.97}{\in}%
\TextAt{21}{291.643}{713.018}{S}%
\TextAt{13}{299.254}{727.032}{K}%
\TextAt{3}{310.072}{731.968}{\Character{3}}%
\TextAt{27}{309.213}{724.076}{p}%
\TextAt{16}{317.461}{727.032}{\Character{2}}%
\TextAt{20}{329.416}{738.389}{\Character{89}}%
\TextAt{27}{329.439}{712.486}{p}%
\TextAt{17}{334.369}{712.486}{\Character{47}}%
\NativeMathAt{333.428}{712.486}{7.97}{\in}%
\TextAt{21}{339.073}{712.486}{S}%
\TextAt{13}{346.685}{727.032}{U}%
\TextAt{27}{354.635}{725.238}{p}%
\TextAt{8}{79.226}{693.895}{of}%
\TextAt{13}{92.949}{693.895}{S}%
\TextAt{8}{100.844}{693.895}{-idèles,}%
\TextAt{8}{140.404}{693.895}{where}%
\TextAt{13}{174.606}{693.895}{U}%
\TextAt{27}{182.556}{692.102}{p}%
\TextAt{9}{191.033}{693.895}{\Character{61}}%
\TextAt{13}{204.127}{693.895}{K}%
\TextAt{3}{214.945}{698.234}{\Character{3}}%
\TextAt{27}{214.087}{690.94}{p}%
\TextAt{8}{223.967}{693.895}{for}%
\TextAt{14}{242.245}{693.895}{p}%
\TextAt{8}{252.518}{693.895}{infinite}%
\TextAt{8}{292.578}{693.895}{complex,}%
\TextAt{8}{342.168}{693.895}{and}%
\TextAt{13}{365.324}{693.895}{U}%
\TextAt{27}{373.274}{692.102}{p}%
\TextAt{9}{381.751}{693.895}{\Character{61}}%
\TextAt{18}{394.845}{693.895}{R}%
\TextAt{3}{403.479}{698.234}{\Character{3}}%
\TextAt{17}{403.479}{690.94}{\Character{43}}%
\TextAt{8}{414.854}{693.895}{for}%
\TextAt{14}{433.131}{693.895}{p}%
\TextAt{8}{443.404}{693.895}{infinite}%
\TextAt{8}{483.464}{693.895}{real.}%
\TextAt{8}{511.956}{693.895}{One}%
\TextAt{8}{79.226}{676.85}{clearly}%
\TextAt{8}{116.622}{676.85}{has}%
\TextAt{13}{275.495}{659.804}{I}%
\TextAt{21}{280.658}{658.01}{K}%
\TextAt{9}{292.182}{659.804}{\Character{61}}%
\TextAt{20}{304.607}{671.161}{\Character{91}}%
\TextAt{21}{308.451}{645.789}{S}%
\TextAt{13}{319.883}{659.804}{I}%
\TextAt{21}{325.986}{664.74}{S}%
\TextAt{21}{325.046}{656.848}{K}%
\TextAt{13}{333.249}{659.804}{\Character{59}}%
\TextAt{8}{79.226}{628.705}{if}%
\TextAt{13}{89.95}{628.705}{S}%
\TextAt{8}{101.743}{628.705}{varies}%
\TextAt{8}{134.655}{628.705}{over}%
\TextAt{8}{159.701}{628.705}{all}%
\TextAt{8}{175.96}{628.705}{finite}%
\TextAt{8}{205.872}{628.705}{sets}%
\TextAt{8}{228.754}{628.705}{of}%
\TextAt{8}{242.084}{628.705}{primes}%
\TextAt{8}{279.863}{628.705}{of}%
\TextAt{13}{293.193}{628.705}{K}%
\TextAt{8}{304.01}{628.705}{.}%
\TextAt{8}{97.159}{608.671}{The}%
\TextAt{8}{121.213}{608.671}{inclusions}%
\TextAt{13}{174.664}{608.671}{K}%
\NativeMathAt{188.803}{608.671}{11.955}{\subseteq}%
\TextAt{13}{201.422}{608.671}{K}%
\TextAt{27}{211.381}{606.877}{p}%
\TextAt{8}{219.765}{608.671}{allow}%
\TextAt{8}{250.012}{608.671}{us}%
\TextAt{8}{265.028}{608.671}{to}%
\TextAt{8}{279.338}{608.671}{define}%
\TextAt{8}{313.147}{608.671}{the}%
\TextAt{8}{333.304}{608.671}{diagonal}%
\TextAt{8}{380.144}{608.671}{embedding}%
\TextAt{13}{276.963}{579.67}{K}%
\TextAt{3}{287.78}{584.606}{\Character{3}}%
\NativeArrowAt{295.834}{579.67}{11.955}{19.26}{\longrightarrow}%
\TextAt{13}{318.416}{579.67}{I}%
\TextAt{21}{323.579}{577.876}{K}%
\TextAt{13}{331.781}{579.67}{\Character{59}}%
\TextAt{8}{79.226}{550.668}{which}%
\TextAt{8}{113.226}{550.668}{associates}%
\TextAt{8}{167.585}{550.668}{to}%
\TextAt{13}{182.408}{550.668}{a}%
\NativeMathAt{192.748}{550.668}{11.955}{\in}%
\TextAt{13}{204.913}{550.668}{K}%
\TextAt{3}{215.731}{555.007}{\Character{3}}%
\TextAt{8}{224.874}{550.668}{the}%
\TextAt{8}{245.543}{550.668}{idèle}%
\TextAt{13}{273.362}{550.668}{\Character{11}}%
\NativeMathAt{285.079}{550.668}{11.955}{\in}%
\TextAt{13}{297.244}{550.668}{I}%
\TextAt{21}{302.407}{548.875}{K}%
\TextAt{8}{315.02}{550.668}{whose}%
\TextAt{14}{350.059}{550.668}{p}%
\TextAt{8}{356.043}{550.668}{-th}%
\TextAt{8}{375.409}{550.668}{component}%
\TextAt{8}{435.758}{550.668}{is}%
\TextAt{8}{448.034}{550.668}{the}%
\TextAt{8}{468.704}{550.668}{element}%
\TextAt{13}{512.459}{550.668}{a}%
\TextAt{8}{523.014}{550.668}{in}%
\TextAt{13}{79.226}{533.622}{K}%
\TextAt{27}{89.185}{531.829}{p}%
\TextAt{8}{93.672}{533.622}{.}%
\TextAt{8}{101.902}{533.622}{We}%
\TextAt{8}{121.392}{533.622}{thus}%
\TextAt{8}{146.488}{533.622}{view}%
\TextAt{13}{172.816}{533.622}{K}%
\TextAt{3}{183.634}{537.961}{\Character{3}}%
\TextAt{8}{191.608}{533.622}{as}%
\TextAt{8}{205.324}{533.622}{a}%
\TextAt{8}{214.424}{533.622}{subgroup}%
\TextAt{8}{264.567}{533.622}{of}%
\TextAt{13}{277.242}{533.622}{I}%
\TextAt{21}{282.405}{531.829}{K}%
\TextAt{8}{293.85}{533.622}{and}%
\TextAt{8}{315.958}{533.622}{we}%
\TextAt{8}{332.53}{533.622}{call}%
\TextAt{8}{353.335}{533.622}{the}%
\TextAt{8}{372.836}{533.622}{elements}%
\TextAt{8}{420.038}{533.622}{of}%
\TextAt{13}{432.713}{533.622}{K}%
\TextAt{3}{443.531}{537.961}{\Character{3}}%
\TextAt{8}{451.505}{533.622}{in}%
\TextAt{13}{464.503}{533.622}{I}%
\TextAt{21}{469.666}{531.829}{K}%
\TextAt{12}{481.111}{533.622}{principal}%
\TextAt{12}{79.226}{516.576}{idèles}%
\TextAt{8}{111.756}{516.576}{.}%
\TextAt{8}{120.204}{516.576}{The}%
\TextAt{8}{144.258}{516.576}{intersection}%
\TextAt{13}{268.569}{499.53}{K}%
\TextAt{21}{279.387}{504.467}{S}%
\TextAt{9}{288.802}{499.53}{\Character{61}}%
\TextAt{13}{301.227}{499.53}{K}%
\TextAt{3}{312.045}{504.467}{\Character{3}}%
\TextAt{16}{319.434}{499.53}{\Character{92}}%
\TextAt{13}{330.061}{499.53}{I}%
\TextAt{21}{336.164}{504.467}{S}%
\TextAt{21}{335.224}{496.575}{K}%
\TextAt{8}{79.226}{476.009}{consists}%
\TextAt{8}{122.976}{476.009}{of}%
\TextAt{8}{136.946}{476.009}{the}%
\TextAt{8}{157.742}{476.009}{numbers}%
\TextAt{13}{205.593}{476.009}{a}%
\NativeMathAt{216.149}{476.009}{11.955}{\in}%
\TextAt{13}{228.531}{476.009}{K}%
\TextAt{3}{239.348}{480.347}{\Character{3}}%
\TextAt{8}{248.618}{476.009}{which}%
\TextAt{8}{282.744}{476.009}{are}%
\TextAt{8}{302.895}{476.009}{units}%
\TextAt{8}{332.861}{476.009}{at}%
\TextAt{8}{347.811}{476.009}{all}%
\TextAt{8}{364.71}{476.009}{primes}%
\TextAt{14}{403.128}{476.009}{p}%
\NativeMathAt{413.522}{476.009}{11.955}{\notin}%
\TextAt{13}{425.903}{476.009}{S}%
\TextAt{8}{433.799}{476.009}{,}%
\TextAt{14}{441.748}{476.009}{p}%
\TextAt{18}{452.142}{476.009}{\Character{45}}%
\NativeMathAt{459.874}{476.009}{11.955}{\infty}%
\TextAt{8}{471.829}{476.009}{,}%
\TextAt{8}{479.778}{476.009}{and}%
\TextAt{8}{503.18}{476.009}{which}%
\TextAt{8}{79.226}{458.963}{are}%
\TextAt{8}{99.648}{458.963}{positive}%
\TextAt{8}{143.873}{458.963}{in}%
\TextAt{13}{158.437}{458.963}{K}%
\TextAt{27}{168.396}{457.17}{p}%
\TextAt{9}{177.757}{458.963}{\Character{61}}%
\TextAt{18}{191.735}{458.963}{R}%
\TextAt{8}{205.178}{458.963}{for}%
\TextAt{8}{223.974}{458.963}{all}%
\TextAt{8}{241.144}{458.963}{real}%
\TextAt{8}{264.818}{458.963}{infinite}%
\TextAt{8}{305.397}{458.963}{places}%
\TextAt{14}{340.834}{458.963}{p}%
\NativeMathAt{351.691}{458.963}{11.955}{\notin}%
\TextAt{13}{364.535}{458.963}{S}%
\TextAt{8}{372.43}{458.963}{.}%
\TextAt{8}{383.612}{458.963}{They}%
\TextAt{8}{414.758}{458.963}{are}%
\TextAt{8}{435.18}{458.963}{called}%
\TextAt{13}{469.255}{458.963}{S}%
\TextAt{12}{477.15}{458.963}{-units}%
\TextAt{8}{510.864}{458.963}{.}%
\TextAt{8}{522.046}{458.963}{In}%
\TextAt{8}{79.226}{441.917}{particular,}%
\TextAt{8}{136.191}{441.917}{for}%
\TextAt{8}{153.73}{441.917}{the}%
\TextAt{8}{173.541}{441.917}{set}%
\TextAt{13}{191.463}{441.917}{S}%
\NativeMathAt{198.663}{440.124}{7.97}{\infty}%
\TextAt{8}{211.181}{441.917}{of}%
\TextAt{8}{224.165}{441.917}{infinite}%
\TextAt{8}{263.487}{441.917}{places,}%
\TextAt{13}{300.989}{441.917}{K}%
\TextAt{21}{311.807}{446.255}{S}%
\NativeMathAt{316.957}{445.259}{5.978}{\infty}%
\TextAt{8}{328.811}{441.917}{is}%
\TextAt{8}{340.229}{441.917}{the}%
\TextAt{8}{360.04}{441.917}{unit}%
\TextAt{8}{384.406}{441.917}{group}%
\TextAt{16}{417.236}{441.917}{\Character{79}}%
\TextAt{3}{427.086}{446.255}{\Character{3}}%
\TextAt{21}{426.754}{438.896}{K}%
\TextAt{8}{438.508}{441.917}{of}%
\TextAt{16}{451.493}{441.917}{\Character{79}}%
\TextAt{21}{461.01}{440.124}{K}%
\TextAt{8}{469.213}{441.917}{.}%
\TextAt{8}{477.546}{441.917}{We}%
\TextAt{8}{497.345}{441.917}{get}%
\TextAt{8}{516.51}{441.917}{the}%
\TextAt{8}{79.226}{424.871}{following}%
\TextAt{8}{128.661}{424.871}{generalization}%
\TextAt{8}{203.464}{424.871}{of}%
\TextAt{8}{216.794}{424.871}{Dirichlet’s}%
\TextAt{8}{272.948}{424.871}{unit}%
\TextAt{8}{297.659}{424.871}{theorem.}%
\TextAt{12}{79.226}{398.859}{Proposition}%
\TextAt{12}{151.937}{398.859}{(1.1).}%
\TextAt{8}{187.253}{398.859}{If}%
\TextAt{13}{198.945}{398.859}{S}%
\TextAt{8}{210.738}{398.859}{contains}%
\TextAt{8}{256.658}{398.859}{all}%
\TextAt{8}{272.917}{398.859}{infinite}%
\TextAt{8}{312.584}{398.859}{places,}%
\TextAt{8}{350.362}{398.859}{then}%
\TextAt{8}{377.022}{398.859}{the}%
\TextAt{8}{397.179}{398.859}{homomorphism}%
\TextAt{13}{198.63}{369.857}{\Character{21}}%
\TextAt{9}{208.779}{369.857}{\Character{58}}%
\TextAt{13}{215.351}{369.857}{K}%
\TextAt{21}{226.169}{374.794}{S}%
\NativeArrowAt{235.584}{369.857}{11.955}{19.26}{\longrightarrow}%
\TextAt{20}{258.166}{381.215}{\Character{89}}%
\TextAt{27}{258.189}{355.843}{p}%
\NativeMathAt{262.178}{355.843}{7.97}{\in}%
\TextAt{21}{267.823}{355.843}{S}%
\TextAt{18}{275.434}{369.857}{R}%
\TextAt{13}{284.069}{369.857}{\Character{59}}%
\TextAt{13}{313.204}{369.857}{\Character{21}}%
\TextAt{9}{320.052}{369.857}{\Character{40}}%
\TextAt{13}{324.604}{369.857}{a}%
\TextAt{9}{330.749}{369.857}{\Character{41}}%
\TextAt{9}{338.604}{369.857}{\Character{61}}%
\TextAt{9}{351.025}{369.857}{\Character{40}}%
\TextAt{8}{355.6}{369.857}{log}%
\TextAt{16}{372.56}{369.857}{\Character{106}}%
\TextAt{13}{375.881}{369.857}{a}%
\TextAt{16}{382.026}{369.857}{\Character{106}}%
\TextAt{27}{385.347}{368.064}{p}%
\TextAt{9}{389.834}{369.857}{\Character{41}}%
\TextAt{27}{394.386}{368.064}{p}%
\NativeMathAt{398.375}{368.064}{7.97}{\in}%
\TextAt{21}{404.02}{368.064}{S}%
\TextAt{13}{410.114}{369.857}{\Character{59}}%
\TextAt{8}{79.226}{331.223}{has}%
\TextAt{8}{99.931}{331.223}{kernel}%
\TextAt{13}{134.23}{331.223}{\Character{22}}%
\TextAt{9}{141.273}{331.223}{\Character{40}}%
\TextAt{13}{145.825}{331.223}{K}%
\TextAt{9}{156.643}{331.223}{\Character{41}}%
\TextAt{8}{161.195}{331.223}{,}%
\TextAt{8}{168.209}{331.223}{and}%
\TextAt{8}{190.803}{331.223}{its}%
\TextAt{8}{206.954}{331.223}{image}%
\TextAt{8}{240.606}{331.223}{is}%
\TextAt{8}{252.202}{331.223}{a}%
\TextAt{8}{261.788}{331.223}{complete}%
\TextAt{8}{311.043}{331.223}{lattice}%
\TextAt{8}{346.644}{331.223}{in}%
\TextAt{8}{360.128}{331.223}{the}%
\TextAt{9}{380.116}{331.223}{\Character{40}}%
\TextAt{13}{384.668}{331.223}{s}%
\TextAt{16}{392.494}{331.223}{\Character{0}}%
\TextAt{9}{404.105}{331.223}{1\Character{41}}%
\TextAt{8}{414.51}{331.223}{-dimensional}%
\TextAt{8}{482.689}{331.223}{trace-zero}%
\TextAt{8}{79.226}{314.177}{space}%
\TextAt{13}{190.142}{297.131}{H}%
\TextAt{9}{204.086}{297.131}{\Character{61}}%
\TextAt{16}{216.511}{297.131}{\Character{102}}%
\TextAt{9}{222.489}{297.131}{\Character{40}}%
\TextAt{13}{227.041}{297.131}{x}%
\TextAt{27}{233.693}{295.338}{p}%
\TextAt{9}{238.18}{297.131}{\Character{41}}%
\NativeMathAt{246.053}{297.131}{11.955}{\in}%
\TextAt{20}{257.344}{308.489}{\Character{89}}%
\TextAt{27}{257.367}{283.117}{p}%
\NativeMathAt{261.356}{283.117}{7.97}{\in}%
\TextAt{21}{267.001}{283.117}{S}%
\TextAt{18}{274.613}{297.131}{R}%
\TextAt{16}{286.568}{297.131}{\Character{106}}%
\TextAt{20}{293.21}{308.489}{\Character{88}}%
\TextAt{27}{294.229}{283.117}{p}%
\NativeMathAt{298.218}{283.117}{7.97}{\in}%
\TextAt{21}{303.863}{283.117}{S}%
\TextAt{13}{312.471}{297.131}{x}%
\TextAt{27}{319.123}{295.338}{p}%
\TextAt{9}{326.93}{297.131}{\Character{61}}%
\TextAt{9}{339.339}{297.131}{0}%
\TextAt{16}{345.209}{297.131}{\Character{103}}%
\TextAt{13}{351.187}{297.131}{\Character{59}}%
\TextAt{13}{380.322}{297.131}{s}%
\TextAt{9}{389.176}{297.131}{\Character{61}}%
\TextAt{8}{401.601}{297.131}{\char35\relax }%
\TextAt{13}{411.357}{297.131}{S\Character{58}}%
\TextAt{19}{79.226}{246.044}{Proof.}%
\TextAt{8}{116}{246.044}{For}%
\TextAt{8}{138.555}{246.044}{the}%
\TextAt{8}{160.297}{246.044}{set}%
\TextAt{13}{180.15}{246.044}{S}%
\NativeMathAt{187.35}{244.25}{7.97}{\infty}%
\TextAt{9}{202.339}{246.043}{\Character{61}}%
\TextAt{16}{217.466}{246.043}{\Character{102}}%
\TextAt{14}{223.444}{246.043}{p}%
\TextAt{16}{235.45}{246.043}{\Character{106}}%
\NativeMathAt{244.775}{246.043}{11.955}{\infty}%
\TextAt{16}{256.73}{246.043}{\Character{103}}%
\TextAt{8}{262.727}{246.044}{,}%
\TextAt{8}{271.858}{246.044}{this}%
\TextAt{8}{296.266}{246.044}{is}%
\TextAt{8}{309.615}{246.044}{the}%
\TextAt{8}{331.357}{246.044}{claim}%
\TextAt{8}{364.158}{246.044}{of}%
\TextAt{8}{379.073}{246.044}{chap.}%
\TextAt{8}{411.551}{246.044}{I,}%
\TextAt{8}{424.506}{246.044}{(7.1)}%
\TextAt{8}{454.067}{246.044}{and}%
\TextAt{8}{478.415}{246.044}{(7.3).}%
\TextAt{8}{515.698}{246.044}{Let}%
\TextAt{13}{79.226}{228.998}{S}%
\TextAt{21}{86.425}{227.204}{f}%
\TextAt{9}{96.689}{228.998}{\Character{61}}%
\TextAt{13}{110.613}{228.998}{S}%
\TextAt{16}{121.764}{228.998}{\Character{110}}%
\TextAt{13}{130.998}{228.998}{S}%
\NativeMathAt{138.197}{227.204}{7.97}{\infty}%
\TextAt{8}{147.163}{228.998}{,}%
\TextAt{8}{155.411}{228.998}{and}%
\TextAt{8}{179.053}{228.998}{let}%
\TextAt{13}{196.837}{228.998}{J}%
\TextAt{9}{204.414}{228.998}{\Character{40}}%
\TextAt{13}{208.966}{228.998}{S}%
\TextAt{21}{216.166}{227.204}{f}%
\TextAt{9}{221.61}{228.998}{\Character{41}}%
\TextAt{8}{230.939}{228.998}{be}%
\TextAt{8}{247.743}{228.998}{the}%
\TextAt{8}{268.778}{228.998}{subgroup}%
\TextAt{8}{320.455}{228.998}{of}%
\TextAt{13}{334.664}{228.998}{J}%
\TextAt{21}{341.134}{227.204}{K}%
\TextAt{8}{354.113}{228.998}{generated}%
\TextAt{8}{408.324}{228.998}{by}%
\TextAt{8}{425.462}{228.998}{the}%
\TextAt{8}{446.498}{228.998}{prime}%
\TextAt{8}{480.54}{228.998}{ideals}%
\TextAt{14}{513.997}{228.998}{p}%
\NativeMathAt{524.799}{228.998}{11.955}{\in}%
\TextAt{13}{79.226}{211.952}{S}%
\TextAt{21}{86.425}{210.158}{f}%
\TextAt{8}{91.87}{211.952}{.}%
\TextAt{8}{103.397}{211.952}{Associating}%
\TextAt{8}{166.985}{211.952}{to}%
\TextAt{8}{182.321}{211.952}{every}%
\TextAt{13}{214.24}{211.952}{a}%
\NativeMathAt{225.454}{211.952}{11.955}{\in}%
\TextAt{13}{238.494}{211.952}{K}%
\TextAt{21}{249.312}{216.29}{S}%
\TextAt{8}{260.329}{211.952}{the}%
\TextAt{8}{281.512}{211.952}{principal}%
\TextAt{8}{331.315}{211.952}{ideal}%
\TextAt{13}{360.304}{211.952}{ia}%
\TextAt{9}{375.512}{211.952}{\Character{61}}%
\TextAt{9}{389.679}{211.952}{\Character{40}}%
\TextAt{13}{394.239}{211.952}{a}%
\TextAt{9}{400.384}{211.952}{\Character{41}}%
\NativeMathAt{410.006}{211.952}{11.955}{\in}%
\TextAt{13}{423.045}{211.952}{J}%
\TextAt{9}{430.622}{211.952}{\Character{40}}%
\TextAt{13}{435.175}{211.952}{S}%
\TextAt{21}{442.374}{210.158}{f}%
\TextAt{9}{447.819}{211.952}{\Character{41}}%
\TextAt{8}{452.371}{211.952}{,}%
\TextAt{8}{460.803}{211.952}{we}%
\TextAt{8}{479.057}{211.952}{obtain}%
\TextAt{8}{516.51}{211.952}{the}%
\TextAt{8}{79.226}{194.906}{commutative}%
\TextAt{8}{149.152}{194.906}{diagram}%
\TextAt{9}{178.401}{165.339}{1}%
\TextAt{16}{245.346}{165.339}{\Character{79}}%
\TextAt{3}{255.196}{169.677}{\Character{3}}%
\TextAt{21}{254.864}{162.318}{K}%
\TextAt{13}{328.653}{165.339}{K}%
\TextAt{21}{339.471}{169.677}{S}%
\TextAt{13}{403.834}{165.339}{J}%
\TextAt{9}{411.411}{165.339}{\Character{40}}%
\TextAt{13}{415.964}{165.339}{S}%
\TextAt{21}{423.163}{163.546}{f}%
\TextAt{9}{428.608}{165.339}{\Character{41}}%
\TextAt{9}{178.401}{119.164}{0}%
\TextAt{20}{241.254}{130.521}{\Character{89}}%
\TextAt{27}{237.598}{105.15}{p}%
\NativeMathAt{241.587}{105.15}{7.97}{\in}%
\TextAt{21}{247.232}{105.15}{S}%
\NativeMathAt{252.383}{104.153}{5.978}{\infty}%
\TextAt{18}{262.179}{119.164}{R}%
\TextAt{20}{324.157}{130.521}{\Character{89}}%
\TextAt{27}{324.18}{105.15}{p}%
\NativeMathAt{328.169}{105.15}{7.97}{\in}%
\TextAt{21}{333.815}{105.15}{S}%
\TextAt{18}{341.426}{119.164}{R}%
\TextAt{20}{405.545}{130.521}{\Character{89}}%
\TextAt{27}{403.404}{105.15}{p}%
\NativeMathAt{407.393}{105.15}{7.97}{\in}%
\TextAt{21}{413.039}{105.15}{S}%
\TextAt{22}{418.189}{103.744}{f}%
\TextAt{18}{424.955}{119.164}{R}%
\NativeDiagram{%
\path[draw=black,line width=0.478bp] (117.607,-605.672) -- (167.517,-605.672);%
\path[draw=black,line width=0.478bp,line cap=round,line join=round] (165.27,-602.803) .. controls (165.724,-604.524) and (166.737,-605.337) .. (167.757,-605.672) .. controls (166.737,-606.007) and (165.724,-606.819) .. (165.27,-608.541);%
\path[draw=black,line width=0.478bp] (196.42,-605.672) -- (250.826,-605.672);%
\path[draw=black,line width=0.478bp,line cap=round,line join=round] (248.579,-602.803) .. controls (249.033,-604.524) and (250.045,-605.337) .. (251.065,-605.672) .. controls (250.045,-606.007) and (249.033,-606.819) .. (248.579,-608.541);%
\path[draw=black,line width=0.478bp] (182.208,-616.26) -- (182.208,-637.226);%
\path[draw=black,line width=0.478bp,line cap=round,line join=round] (185.077,-634.978) .. controls (183.356,-635.432) and (182.543,-636.445) .. (182.208,-637.465) .. controls (181.873,-636.445) and (181.06,-635.432) .. (179.339,-634.978);%
}%
\TextAt{21}{257.023}{143.951}{\Character{21}}%
\TextAt{23}{261.963}{146.764}{\Character{48}}%
\NativeDiagram{%
\path[draw=black,line width=0.478bp] (278.92,-605.672) -- (326.008,-605.672);%
\path[draw=black,line width=0.478bp,line cap=round,line join=round] (323.76,-602.803) .. controls (324.214,-604.524) and (325.227,-605.337) .. (326.247,-605.672) .. controls (325.227,-606.007) and (324.214,-606.819) .. (323.76,-608.541);%
}%
\TextAt{21}{373.261}{171.144}{i}%
\NativeDiagram{%
\path[draw=black,line width=0.478bp] (265.112,-613.24) -- (265.112,-637.226);%
\path[draw=black,line width=0.478bp,line cap=round,line join=round] (267.981,-634.978) .. controls (266.26,-635.432) and (265.447,-636.445) .. (265.112,-637.465) .. controls (264.777,-636.445) and (263.964,-635.432) .. (262.243,-634.978);%
}%
\TextAt{21}{339.927}{145.761}{\Character{21}}%
\NativeDiagram{%
\path[draw=black,line width=0.478bp] (346.501,-616.583) -- (346.501,-637.226);%
\path[draw=black,line width=0.478bp,line cap=round,line join=round] (349.37,-634.978) .. controls (347.649,-635.432) and (346.836,-636.445) .. (346.501,-637.465) .. controls (346.166,-636.445) and (345.353,-635.432) .. (343.632,-634.978);%
}%
\TextAt{21}{421.317}{143.79}{\Character{21}}%
\TextAt{23}{426.257}{146.603}{\Character{48}\Character{48}}%
\NativeDiagram{%
\path[draw=black,line width=0.478bp] (117.607,-651.847) -- (159.77,-651.847);%
\path[draw=black,line width=0.478bp,line cap=round,line join=round] (157.522,-648.978) .. controls (157.977,-650.7) and (158.989,-651.513) .. (160.009,-651.847) .. controls (158.989,-652.182) and (157.977,-652.995) .. (157.522,-654.716);%
\path[draw=black,line width=0.478bp] (204.168,-651.847) -- (246.33,-651.847);%
\path[draw=black,line width=0.478bp,line cap=round,line join=round] (244.083,-648.978) .. controls (244.537,-650.7) and (245.549,-651.513) .. (246.569,-651.847) .. controls (245.549,-652.182) and (244.537,-652.995) .. (244.083,-654.716);%
\path[draw=black,line width=0.478bp] (283.416,-651.847) -- (325.578,-651.847);%
\path[draw=black,line width=0.478bp,line cap=round,line join=round] (323.331,-648.978) .. controls (323.785,-650.7) and (324.797,-651.513) .. (325.817,-651.847) .. controls (324.797,-652.182) and (323.785,-652.995) .. (323.331,-654.716);%
}%
\TextAt{21}{375.294}{124.968}{i}%
\TextAt{8}{297.211}{43.033}{327}%
\Destination{78.23}{809.815}{nav.page.327}%
\Destination{79.36}{776.155}{nav.print.358}%
\Destination{79.36}{412.855}{nav.number.6.1.1}%
\Hotlink{411.051}{243.855}{5.22}{11.047}{nav.chapter.1}%
\Hotlink{424.006}{243.855}{25.078}{11.047}{nav.number.1.7.1}%
\Hotlink{477.915}{243.855}{25.078}{11.047}{nav.number.1.7.3}%
}

% Page 340 (body page 328)
\Page{%
\TextAt{2}{79.226}{796.943}{VI.}%
\TextAt{2}{98.012}{796.943}{Global}%
\TextAt{2}{133.248}{796.943}{Class}%
\TextAt{2}{161.841}{796.943}{Field}%
\TextAt{2}{189.572}{796.943}{Theory}%
\Vector{q 0 G 0 g 0.299 w 7.226 18.729 m 460.769 18.729 l S Q}%
\TextAt{8}{79.226}{761.124}{with}%
\TextAt{8}{105.886}{761.124}{exact}%
\TextAt{8}{136.778}{761.124}{rows.}%
\TextAt{8}{168.384}{761.124}{The}%
\TextAt{8}{192.438}{761.124}{map}%
\TextAt{13}{218.452}{761.124}{\Character{21}}%
\TextAt{3}{225.281}{765.462}{\Character{48}\Character{48}}%
\TextAt{8}{234.27}{761.124}{on}%
\TextAt{8}{250.529}{761.124}{the}%
\TextAt{8}{270.685}{761.124}{right}%
\TextAt{8}{298.983}{761.124}{is}%
\TextAt{8}{310.747}{761.124}{given}%
\TextAt{8}{341.317}{761.124}{by}%
\TextAt{13}{216.291}{722.778}{\Character{21}}%
\TextAt{3}{223.12}{727.714}{\Character{48}\Character{48}}%
\TextAt{20}{230.204}{746.808}{\Character{48}}%
\TextAt{20}{230.204}{725.289}{\Character{64}}%
\TextAt{20}{242.806}{734.135}{\Character{89}}%
\TextAt{27}{240.665}{708.764}{p}%
\NativeMathAt{244.654}{708.764}{7.97}{\in}%
\TextAt{21}{250.299}{708.764}{S}%
\TextAt{22}{255.449}{707.358}{f}%
\TextAt{14}{262.216}{722.778}{p}%
\TextAt{21}{268.199}{727.714}{v}%
\TextAt{28}{272.322}{726.718}{p}%
\TextAt{20}{276.903}{746.808}{\Character{49}}%
\TextAt{20}{276.903}{725.289}{\Character{65}}%
\TextAt{9}{290.685}{722.778}{\Character{61}}%
\TextAt{9}{303.095}{722.778}{\Character{40}}%
\TextAt{16}{307.663}{722.778}{\Character{0}}%
\TextAt{13}{316.961}{722.778}{v}%
\TextAt{27}{322.62}{720.985}{p}%
\TextAt{8}{329.1}{722.778}{log}%
\TextAt{14}{346.06}{722.778}{N}%
\TextAt{9}{356.008}{722.778}{\Character{40}}%
\TextAt{14}{360.56}{722.778}{p}%
\TextAt{9}{366.544}{722.778}{\Character{41}\Character{41}}%
\TextAt{27}{375.648}{720.985}{p}%
\NativeMathAt{379.637}{720.985}{7.97}{\in}%
\TextAt{21}{385.283}{720.985}{S}%
\TextAt{22}{390.433}{719.579}{f}%
\TextAt{8}{79.226}{682.545}{(observe}%
\TextAt{8}{124.814}{682.545}{that}%
\TextAt{16}{149.529}{682.545}{\Character{106}}%
\TextAt{13}{152.85}{682.545}{a}%
\TextAt{16}{158.995}{682.545}{\Character{106}}%
\TextAt{27}{162.316}{680.751}{p}%
\TextAt{9}{170.124}{682.545}{\Character{61}}%
\TextAt{14}{182.549}{682.545}{N}%
\TextAt{9}{192.497}{682.545}{\Character{40}}%
\TextAt{14}{197.049}{682.545}{p}%
\TextAt{9}{203.033}{682.545}{\Character{41}}%
\TextAt{3}{207.585}{686.883}{\Character{0}}%
\TextAt{21}{214.171}{686.883}{v}%
\TextAt{28}{218.295}{685.887}{p}%
\TextAt{17}{222.378}{686.883}{\Character{40}}%
\TextAt{21}{225.671}{686.883}{a}%
\TextAt{17}{230.169}{686.883}{\Character{41}}%
\TextAt{8}{233.96}{682.545}{),}%
\TextAt{8}{245.141}{682.545}{and}%
\TextAt{8}{267.249}{682.545}{maps}%
\TextAt{13}{297.224}{682.545}{J}%
\TextAt{9}{304.801}{682.545}{\Character{40}}%
\TextAt{13}{309.354}{682.545}{S}%
\TextAt{21}{316.553}{680.751}{f}%
\TextAt{9}{321.998}{682.545}{\Character{41}}%
\TextAt{8}{329.793}{682.545}{isomorphically}%
\TextAt{8}{406.931}{682.545}{onto}%
\TextAt{8}{432.626}{682.545}{the}%
\TextAt{8}{452.129}{682.545}{complete}%
\TextAt{8}{500.897}{682.545}{lattice}%
\TextAt{8}{79.226}{665.499}{spanned}%
\TextAt{8}{124.811}{665.499}{by}%
\TextAt{8}{141.07}{665.499}{the}%
\TextAt{8}{161.227}{665.499}{vectors}%
\TextAt{13}{222.821}{648.453}{e}%
\TextAt{27}{228.247}{646.659}{p}%
\TextAt{9}{236.054}{648.453}{\Character{61}}%
\TextAt{9}{248.464}{648.453}{\Character{40}0}%
\TextAt{13}{258.885}{648.453}{\Character{59}}%
\TextAt{13}{264.109}{648.453}{\Character{58}}%
\TextAt{13}{269.346}{648.453}{\Character{58}}%
\TextAt{13}{274.582}{648.453}{\Character{58}}%
\TextAt{13}{279.807}{648.453}{\Character{59}}%
\TextAt{9}{285.106}{648.453}{0}%
\TextAt{13}{290.959}{648.453}{\Character{59}}%
\TextAt{8}{296.203}{648.453}{log}%
\TextAt{14}{313.164}{648.453}{N}%
\TextAt{9}{323.112}{648.453}{\Character{40}}%
\TextAt{14}{327.664}{648.453}{p}%
\TextAt{9}{333.647}{648.453}{\Character{41}}%
\TextAt{13}{338.199}{648.453}{\Character{59}}%
\TextAt{9}{343.444}{648.453}{0}%
\TextAt{13}{349.297}{648.453}{\Character{59}}%
\TextAt{13}{354.521}{648.453}{\Character{58}}%
\TextAt{13}{359.758}{648.453}{\Character{58}}%
\TextAt{13}{364.982}{648.453}{\Character{58}}%
\TextAt{13}{370.219}{648.453}{\Character{59}}%
\TextAt{9}{375.517}{648.453}{0\Character{41}}%
\TextAt{13}{385.923}{648.453}{\Character{59}}%
\TextAt{8}{79.226}{625.602}{for}%
\TextAt{14}{97.111}{625.602}{p}%
\NativeMathAt{106.415}{625.602}{11.955}{\in}%
\TextAt{13}{117.706}{625.602}{S}%
\TextAt{21}{124.905}{623.808}{f}%
\TextAt{8}{130.35}{625.602}{.}%
\TextAt{8}{138.799}{625.602}{It}%
\TextAt{8}{151.471}{625.602}{follows}%
\TextAt{8}{189.907}{625.602}{that}%
\TextAt{8}{215.276}{625.602}{ker}%
\TextAt{9}{231.212}{625.602}{\Character{40}}%
\TextAt{13}{235.764}{625.602}{\Character{21}}%
\TextAt{9}{242.593}{625.602}{\Character{41}}%
\TextAt{9}{250.448}{625.602}{\Character{61}}%
\TextAt{8}{262.891}{625.602}{ker}%
\TextAt{9}{278.828}{625.602}{\Character{40}}%
\TextAt{13}{283.38}{625.602}{\Character{21}}%
\TextAt{3}{290.208}{629.94}{\Character{48}}%
\TextAt{9}{293.004}{625.602}{\Character{41}}%
\TextAt{9}{300.859}{625.602}{\Character{61}}%
\TextAt{13}{313.302}{625.602}{\Character{22}}%
\TextAt{9}{320.345}{625.602}{\Character{40}}%
\TextAt{13}{324.897}{625.602}{K}%
\TextAt{9}{335.715}{625.602}{\Character{41}}%
\TextAt{8}{340.268}{625.602}{,}%
\TextAt{8}{347.417}{625.602}{and}%
\TextAt{8}{370.179}{625.602}{we}%
\TextAt{8}{387.407}{625.602}{obtain}%
\TextAt{8}{423.834}{625.602}{the}%
\TextAt{8}{443.991}{625.602}{exact}%
\TextAt{8}{474.883}{625.602}{sequence}%
\TextAt{9}{218.89}{598.165}{0}%
\NativeArrowAt{228.064}{598.165}{11.955}{19.26}{\longrightarrow}%
\TextAt{8}{250.646}{598.165}{im}%
\TextAt{9}{263.653}{598.165}{\Character{40}}%
\TextAt{13}{268.205}{598.165}{\Character{21}}%
\TextAt{3}{275.034}{603.102}{\Character{48}}%
\TextAt{9}{277.829}{598.165}{\Character{41}}%
\NativeArrowAt{285.702}{598.165}{11.955}{19.26}{\longrightarrow}%
\TextAt{8}{308.284}{598.165}{im}%
\TextAt{9}{321.291}{598.165}{\Character{40}}%
\TextAt{13}{325.844}{598.165}{\Character{21}}%
\TextAt{9}{332.672}{598.165}{\Character{41}}%
\TextAt{21}{344.376}{604.942}{i}%
\NativeArrowAt{340.545}{598.165}{11.955}{11.955}{\longrightarrow}%
\TextAt{8}{355.822}{598.165}{im}%
\TextAt{9}{368.829}{598.165}{\Character{40}}%
\TextAt{13}{373.381}{598.165}{\Character{21}}%
\TextAt{3}{380.21}{603.102}{\Character{48}\Character{48}}%
\TextAt{9}{385.302}{598.165}{\Character{41}}%
\TextAt{13}{389.854}{598.165}{\Character{59}}%
\TextAt{8}{517.802}{598.165}{(*)}%
\TextAt{8}{79.226}{570.729}{where}%
\TextAt{8}{113.686}{570.729}{the}%
\TextAt{8}{134.494}{570.729}{groups}%
\TextAt{8}{172.935}{570.729}{on}%
\TextAt{8}{189.846}{570.729}{the}%
\TextAt{8}{210.653}{570.729}{left}%
\TextAt{8}{231.783}{570.729}{and}%
\TextAt{8}{255.197}{570.729}{on}%
\TextAt{8}{272.107}{570.729}{the}%
\TextAt{8}{292.915}{570.729}{right}%
\TextAt{8}{321.864}{570.729}{are}%
\TextAt{8}{342.026}{570.729}{lattices.}%
\TextAt{8}{388.915}{570.729}{This}%
\TextAt{8}{416.286}{570.729}{implies}%
\TextAt{8}{456.664}{570.729}{that}%
\TextAt{8}{482.684}{570.729}{the}%
\TextAt{8}{503.491}{570.729}{group}%
\TextAt{8}{79.226}{553.683}{in}%
\TextAt{8}{93.014}{553.683}{the}%
\TextAt{8}{113.305}{553.683}{middle}%
\TextAt{8}{151.804}{553.683}{is}%
\TextAt{8}{163.703}{553.683}{also}%
\TextAt{8}{187.317}{553.683}{a}%
\TextAt{8}{197.208}{553.683}{lattice.}%
\TextAt{8}{237.933}{553.683}{For}%
\TextAt{8}{259.037}{553.683}{if}%
\TextAt{13}{269.896}{553.683}{x}%
\NativeMathAt{280.098}{553.683}{11.955}{\in}%
\TextAt{8}{291.619}{553.683}{im}%
\TextAt{9}{304.626}{553.683}{\Character{40}}%
\TextAt{13}{309.178}{553.683}{\Character{21}}%
\TextAt{9}{316.007}{553.683}{\Character{41}}%
\TextAt{8}{320.559}{553.683}{,}%
\TextAt{8}{327.877}{553.683}{and}%
\TextAt{13}{350.774}{553.683}{U}%
\TextAt{8}{364.006}{553.683}{is}%
\TextAt{8}{375.904}{553.683}{a}%
\TextAt{8}{385.795}{553.683}{neighbourhood}%
\TextAt{8}{465.933}{553.683}{of}%
\TextAt{13}{479.398}{553.683}{i}%
\TextAt{9}{483.392}{553.683}{\Character{40}}%
\TextAt{13}{487.944}{553.683}{x}%
\TextAt{9}{494.596}{553.683}{\Character{41}}%
\TextAt{8}{503.18}{553.683}{which}%
\TextAt{8}{79.226}{536.637}{contains}%
\TextAt{8}{124.425}{536.637}{no}%
\TextAt{8}{139.962}{536.637}{other}%
\TextAt{8}{169.81}{536.637}{point}%
\TextAt{8}{199.658}{536.637}{of}%
\TextAt{8}{212.267}{536.637}{im}%
\TextAt{9}{225.274}{536.637}{\Character{40}}%
\TextAt{13}{229.826}{536.637}{\Character{21}}%
\TextAt{3}{236.655}{540.976}{\Character{48}\Character{48}}%
\TextAt{9}{241.747}{536.637}{\Character{41}}%
\TextAt{8}{246.299}{536.637}{,}%
\TextAt{8}{252.871}{536.637}{then}%
\TextAt{13}{278.81}{536.637}{i}%
\TextAt{3}{282.803}{540.976}{\Character{0}}%
\TextAt{17}{289.39}{540.976}{1}%
\TextAt{9}{294.122}{536.637}{\Character{40}}%
\TextAt{13}{298.674}{536.637}{U}%
\TextAt{9}{307.873}{536.637}{\Character{41}}%
\TextAt{8}{315.602}{536.637}{contains}%
\TextAt{8}{360.8}{536.637}{the}%
\TextAt{8}{380.235}{536.637}{coset}%
\TextAt{13}{408.84}{536.637}{x}%
\TextAt{9}{416.673}{536.637}{\Character{43}}%
\TextAt{8}{426.959}{536.637}{im}%
\TextAt{9}{439.967}{536.637}{\Character{40}}%
\TextAt{13}{444.519}{536.637}{\Character{21}}%
\TextAt{3}{451.347}{540.976}{\Character{48}}%
\TextAt{9}{454.142}{536.637}{\Character{41}}%
\TextAt{8}{458.695}{536.637}{,}%
\TextAt{8}{465.267}{536.637}{and}%
\TextAt{8}{487.308}{536.637}{no}%
\TextAt{8}{502.846}{536.637}{other.}%
\TextAt{8}{79.226}{519.591}{It}%
\TextAt{8}{91.899}{519.591}{is}%
\TextAt{8}{103.663}{519.591}{discrete}%
\TextAt{8}{146.641}{519.591}{since}%
\TextAt{8}{175.31}{519.591}{im}%
\TextAt{9}{188.317}{519.591}{\Character{40}}%
\TextAt{13}{192.869}{519.591}{\Character{21}}%
\TextAt{3}{199.698}{523.93}{\Character{48}}%
\TextAt{9}{202.493}{519.591}{\Character{41}}%
\TextAt{8}{210.943}{519.591}{is}%
\TextAt{8}{222.707}{519.591}{discrete.}%
\TextAt{8}{97.159}{499.556}{For}%
\TextAt{8}{118.128}{499.556}{every}%
\TextAt{14}{149.02}{499.556}{p}%
\NativeMathAt{158.324}{499.556}{11.955}{\in}%
\TextAt{13}{169.615}{499.556}{S}%
\TextAt{21}{176.815}{497.763}{f}%
\TextAt{8}{182.26}{499.556}{,}%
\TextAt{8}{189.409}{499.556}{if}%
\TextAt{13}{200.133}{499.556}{h}%
\TextAt{8}{210.768}{499.556}{is}%
\TextAt{8}{222.532}{499.556}{the}%
\TextAt{8}{242.689}{499.556}{class}%
\TextAt{8}{270.126}{499.556}{number}%
\TextAt{8}{312.722}{499.556}{of}%
\TextAt{13}{326.052}{499.556}{K}%
\TextAt{8}{336.87}{499.556}{,}%
\TextAt{8}{344.019}{499.556}{then}%
\TextAt{14}{370.679}{499.556}{p}%
\TextAt{21}{376.662}{503.895}{h}%
\TextAt{8}{385.938}{499.556}{belongs}%
\TextAt{8}{427.948}{499.556}{to}%
\TextAt{13}{442.258}{499.556}{i}%
\TextAt{9}{446.252}{499.556}{\Character{40}}%
\TextAt{13}{450.804}{499.556}{K}%
\TextAt{21}{461.622}{503.895}{S}%
\TextAt{9}{467.716}{499.556}{\Character{41}}%
\TextAt{8}{472.268}{499.556}{,}%
\TextAt{8}{479.417}{499.556}{i.e.,}%
\TextAt{13}{241.412}{472.12}{J}%
\TextAt{9}{248.989}{472.12}{\Character{40}}%
\TextAt{13}{253.541}{472.12}{S}%
\TextAt{21}{260.741}{470.327}{f}%
\TextAt{9}{266.186}{472.12}{\Character{41}}%
\TextAt{21}{270.738}{477.056}{h}%
\NativeMathAt{279.437}{472.12}{11.955}{\subseteq}%
\TextAt{13}{292.056}{472.12}{i}%
\TextAt{9}{296.05}{472.12}{\Character{40}}%
\TextAt{13}{300.602}{472.12}{K}%
\TextAt{21}{311.42}{477.056}{S}%
\TextAt{9}{317.513}{472.12}{\Character{41}}%
\NativeMathAt{325.387}{472.12}{11.955}{\subseteq}%
\TextAt{13}{338.006}{472.12}{J}%
\TextAt{9}{345.583}{472.12}{\Character{40}}%
\TextAt{13}{350.135}{472.12}{S}%
\TextAt{21}{357.335}{470.327}{f}%
\TextAt{9}{362.78}{472.12}{\Character{41}}%
\TextAt{13}{367.332}{472.12}{\Character{58}}%
\TextAt{8}{79.226}{444.684}{The}%
\TextAt{8}{104.38}{444.684}{groups}%
\TextAt{8}{143.271}{444.684}{on}%
\TextAt{8}{160.63}{444.684}{the}%
\TextAt{8}{181.886}{444.684}{left}%
\TextAt{8}{203.466}{444.684}{and}%
\TextAt{8}{227.329}{444.684}{on}%
\TextAt{8}{244.688}{444.684}{the}%
\TextAt{8}{265.945}{444.684}{right}%
\TextAt{8}{295.343}{444.684}{have}%
\TextAt{8}{323.438}{444.684}{rank}%
\TextAt{8}{351.533}{444.684}{\char35\relax }%
\TextAt{13}{361.288}{444.684}{S}%
\TextAt{21}{368.487}{442.891}{f}%
\TextAt{8}{373.932}{444.684}{,}%
\TextAt{8}{382.457}{444.684}{hence}%
\TextAt{8}{416.063}{444.684}{so}%
\TextAt{8}{431.533}{444.684}{does}%
\TextAt{13}{459.031}{444.684}{i}%
\TextAt{9}{463.024}{444.684}{\Character{40}}%
\TextAt{13}{467.576}{444.684}{K}%
\TextAt{21}{478.394}{449.022}{S}%
\TextAt{9}{484.488}{444.684}{\Character{41}}%
\TextAt{8}{489.04}{444.684}{.}%
\TextAt{8}{500.789}{444.684}{In}%
\TextAt{8}{516.51}{444.684}{the}%
\TextAt{8}{79.226}{427.638}{sequence}%
\TextAt{8}{128.103}{427.638}{(*),}%
\TextAt{8}{150.688}{427.638}{the}%
\TextAt{8}{171.219}{427.638}{image}%
\TextAt{8}{205.415}{427.638}{of}%
\TextAt{13}{219.119}{427.638}{i}%
\TextAt{8}{227.385}{427.638}{therefore}%
\TextAt{8}{276.859}{427.638}{has}%
\TextAt{8}{298.107}{427.638}{rank}%
\TextAt{8}{325.477}{427.638}{\char35\relax }%
\TextAt{13}{335.232}{427.638}{S}%
\TextAt{21}{342.432}{425.845}{f}%
\TextAt{8}{347.876}{427.638}{,}%
\TextAt{8}{355.494}{427.638}{and}%
\TextAt{8}{378.631}{427.638}{the}%
\TextAt{8}{399.162}{427.638}{kernel}%
\TextAt{8}{434.003}{427.638}{has}%
\TextAt{8}{455.252}{427.638}{rank}%
\TextAt{8}{482.621}{427.638}{\char35\relax }%
\TextAt{13}{492.376}{427.638}{S}%
\NativeMathAt{499.576}{425.845}{7.97}{\infty}%
\TextAt{16}{511.454}{427.638}{\Character{0}}%
\TextAt{9}{523.665}{427.638}{1}%
\TextAt{8}{529.518}{427.638}{.}%
\TextAt{8}{79.226}{410.592}{Hence}%
\TextAt{8}{115.097}{410.592}{im}%
\TextAt{9}{128.104}{410.592}{\Character{40}}%
\TextAt{13}{132.656}{410.592}{\Character{21}}%
\TextAt{9}{139.485}{410.592}{\Character{41}}%
\TextAt{8}{149.027}{410.592}{is}%
\TextAt{8}{161.884}{410.592}{a}%
\TextAt{8}{172.732}{410.592}{lattice}%
\TextAt{8}{209.595}{410.592}{of}%
\TextAt{8}{224.018}{410.592}{rank}%
\TextAt{8}{252.106}{410.592}{\char35\relax }%
\TextAt{13}{261.861}{410.592}{S}%
\NativeMathAt{269.06}{408.799}{7.97}{\infty}%
\TextAt{16}{281.428}{410.592}{\Character{0}}%
\TextAt{9}{294.129}{410.592}{1}%
\TextAt{9}{303.37}{410.592}{\Character{43}}%
\TextAt{8}{315.889}{410.592}{\char35\relax }%
\TextAt{13}{325.645}{410.592}{S}%
\TextAt{21}{332.844}{408.799}{f}%
\TextAt{9}{343.472}{410.592}{\Character{61}}%
\TextAt{8}{357.76}{410.592}{\char35\relax }%
\TextAt{13}{367.516}{410.592}{S}%
\TextAt{16}{378.813}{410.592}{\Character{0}}%
\TextAt{9}{391.513}{410.592}{1}%
\TextAt{8}{397.366}{410.592}{.}%
\TextAt{8}{409.093}{410.592}{It}%
\TextAt{8}{422.858}{410.592}{lies}%
\TextAt{8}{444.167}{410.592}{in}%
\TextAt{8}{458.913}{410.592}{the}%
\TextAt{9}{480.162}{410.592}{\Character{40}}%
\TextAt{8}{484.714}{410.592}{\char35\relax }%
\TextAt{13}{494.47}{410.592}{S}%
\TextAt{16}{505.767}{410.592}{\Character{0}}%
\TextAt{9}{518.467}{410.592}{1\Character{41}}%
\TextAt{8}{528.872}{410.592}{-}%
\TextAt{8}{79.226}{393.546}{dimensional}%
\TextAt{8}{143.676}{393.546}{trace-zero}%
\TextAt{8}{197.654}{393.546}{space}%
\TextAt{13}{228.929}{393.546}{H}%
\TextAt{8}{239.551}{393.546}{,}%
\TextAt{8}{246.701}{393.546}{since}%
\TextAt{20}{275.369}{402.512}{\Character{81}}%
\TextAt{27}{286.66}{390.059}{p}%
\NativeMathAt{290.649}{390.059}{7.97}{\in}%
\TextAt{21}{296.294}{390.059}{S}%
\TextAt{16}{304.381}{393.546}{\Character{106}}%
\TextAt{13}{307.702}{393.546}{a}%
\TextAt{16}{313.847}{393.546}{\Character{106}}%
\TextAt{27}{317.167}{391.753}{p}%
\TextAt{9}{324.975}{393.546}{\Character{61}}%
\TextAt{20}{337.401}{402.512}{\Character{81}}%
\TextAt{27}{348.692}{390.059}{p}%
\TextAt{16}{355.171}{393.546}{\Character{106}}%
\TextAt{13}{358.492}{393.546}{a}%
\TextAt{16}{364.637}{393.546}{\Character{106}}%
\TextAt{27}{367.958}{391.753}{p}%
\TextAt{9}{375.765}{393.546}{\Character{61}}%
\TextAt{9}{388.175}{393.546}{1}%
\TextAt{8}{397.941}{393.546}{for}%
\TextAt{13}{415.826}{393.546}{a}%
\NativeMathAt{425.292}{393.546}{11.955}{\in}%
\TextAt{13}{436.583}{393.546}{K}%
\TextAt{21}{447.401}{397.884}{S}%
\TextAt{8}{453.494}{393.546}{.}%
\Vector{q 0 G 0 g 0.399 w 452.334 -380.454 m 452.334 -372.384 l S Q}%
\Vector{q 0 G 0 g 0.399 w 452.534 -372.584 m 459.707 -372.584 l S Q}%
\Vector{q 0 G 0 g 0.399 w 452.534 -380.255 m 459.707 -380.255 l S Q}%
\Vector{q 0 G 0 g 0.399 w 459.906 -380.454 m 459.906 -372.384 l S Q}%
\TextAt{12}{79.226}{359.461}{Definition}%
\TextAt{12}{142.694}{359.461}{(1.2).}%
\TextAt{8}{178.601}{359.461}{The}%
\TextAt{8}{203.245}{359.461}{elements}%
\TextAt{8}{251.693}{359.461}{of}%
\TextAt{8}{265.614}{359.461}{the}%
\TextAt{8}{286.361}{359.461}{subgroup}%
\TextAt{13}{337.749}{359.461}{K}%
\TextAt{3}{348.567}{363.8}{\Character{3}}%
\TextAt{8}{357.787}{359.461}{of}%
\TextAt{13}{371.708}{359.461}{I}%
\TextAt{21}{376.871}{357.668}{K}%
\TextAt{8}{389.562}{359.461}{are}%
\TextAt{8}{409.664}{359.461}{called}%
\TextAt{12}{443.418}{359.461}{principal}%
\TextAt{12}{500.239}{359.461}{idèles}%
\TextAt{8}{79.226}{342.415}{and}%
\TextAt{8}{101.989}{342.415}{the}%
\TextAt{8}{122.145}{342.415}{quotient}%
\TextAt{8}{168.328}{342.415}{group}%
\TextAt{13}{272.452}{325.369}{C}%
\TextAt{21}{280.827}{323.576}{K}%
\TextAt{9}{292.35}{325.369}{\Character{61}}%
\TextAt{13}{304.776}{325.369}{I}%
\TextAt{21}{309.938}{323.576}{K}%
\TextAt{9}{318.141}{325.369}{\Character{47}}%
\TextAt{13}{323.994}{325.369}{K}%
\TextAt{3}{334.812}{330.306}{\Character{3}}%
\TextAt{8}{79.226}{302.518}{is}%
\TextAt{8}{90.99}{302.518}{called}%
\TextAt{8}{124.154}{302.518}{the}%
\TextAt{12}{144.31}{302.518}{idèle}%
\TextAt{12}{176.015}{302.518}{class}%
\TextAt{12}{207.374}{302.518}{group}%
\TextAt{8}{245.176}{302.518}{of}%
\TextAt{13}{258.506}{302.518}{K}%
\TextAt{8}{269.324}{302.518}{.}%
\TextAt{8}{97.159}{276.953}{The}%
\TextAt{8}{122.302}{276.953}{relation}%
\TextAt{8}{166.321}{276.953}{between}%
\TextAt{8}{212.601}{276.953}{the}%
\TextAt{8}{233.846}{276.953}{ideal}%
\TextAt{8}{262.898}{276.953}{class}%
\TextAt{8}{291.424}{276.953}{group}%
\TextAt{13}{325.689}{276.953}{C}%
\TextAt{13}{334.906}{276.953}{l}%
\TextAt{21}{338.409}{275.16}{K}%
\TextAt{8}{351.598}{276.953}{and}%
\TextAt{8}{375.45}{276.953}{the}%
\TextAt{8}{396.695}{276.953}{idèle}%
\TextAt{8}{425.09}{276.953}{class}%
\TextAt{8}{453.616}{276.953}{group}%
\TextAt{13}{487.88}{276.953}{C}%
\TextAt{21}{496.255}{275.16}{K}%
\TextAt{8}{509.444}{276.953}{is}%
\TextAt{8}{522.297}{276.953}{as}%
\TextAt{8}{79.226}{259.907}{follows.}%
\TextAt{8}{122.213}{259.907}{There}%
\TextAt{8}{156.022}{259.907}{is}%
\TextAt{8}{167.786}{259.907}{a}%
\TextAt{8}{177.541}{259.907}{surjective}%
\TextAt{8}{229.953}{259.907}{homomorphism}%
\TextAt{9}{194.446}{232.471}{\Character{40}}%
\TextAt{9}{202.874}{232.471}{\Character{41}}%
\TextAt{9}{210.741}{232.471}{\Character{58}}%
\TextAt{13}{217.341}{232.471}{I}%
\TextAt{21}{222.504}{230.678}{K}%
\NativeArrowAt{234.027}{232.471}{11.955}{19.26}{\longrightarrow}%
\TextAt{13}{256.61}{232.471}{J}%
\TextAt{21}{263.08}{230.678}{K}%
\TextAt{13}{271.282}{232.471}{\Character{59}}%
\TextAt{13}{300.417}{232.471}{\Character{11}}%
\NativeArrowAt{311.279}{232.471}{11.955}{19.26}{\longmapsto}%
\TextAt{9}{333.861}{232.471}{\Character{40}}%
\TextAt{13}{338.414}{232.471}{\Character{11}}%
\TextAt{9}{345.935}{232.471}{\Character{41}}%
\TextAt{9}{353.79}{232.471}{\Character{61}}%
\TextAt{20}{366.234}{243.828}{\Character{89}}%
\TextAt{27}{366.537}{217.996}{p}%
\TextAt{25}{370.525}{217.996}{\Character{45}}%
\NativeMathAt{372.739}{217.996}{7.97}{\infty}%
\TextAt{14}{383.503}{232.471}{p}%
\TextAt{21}{389.486}{237.407}{v}%
\TextAt{28}{393.609}{236.411}{p}%
\TextAt{17}{397.692}{237.407}{\Character{40}}%
\TextAt{21}{400.985}{237.407}{\Character{11}}%
\TextAt{28}{406.424}{236.411}{p}%
\TextAt{17}{410.507}{237.407}{\Character{41}}%
\TextAt{13}{414.298}{232.471}{\Character{59}}%
\TextAt{8}{79.226}{194.994}{from}%
\TextAt{8}{106.866}{194.994}{the}%
\TextAt{8}{127.023}{194.994}{idèle}%
\TextAt{8}{154.329}{194.994}{group}%
\TextAt{13}{187.504}{194.994}{I}%
\TextAt{21}{192.667}{193.2}{K}%
\TextAt{8}{204.767}{194.994}{to}%
\TextAt{8}{219.077}{194.994}{the}%
\TextAt{8}{239.234}{194.994}{ideal}%
\TextAt{8}{267.197}{194.994}{group}%
\TextAt{13}{300.373}{194.994}{J}%
\TextAt{21}{306.843}{193.2}{K}%
\TextAt{8}{315.045}{194.994}{.}%
\TextAt{8}{323.494}{194.994}{Its}%
\TextAt{8}{340.781}{194.994}{kernel}%
\TextAt{8}{375.248}{194.994}{is}%
\TextAt{13}{248.153}{167.557}{I}%
\TextAt{21}{254.256}{172.672}{S}%
\NativeMathAt{259.406}{171.675}{5.978}{\infty}%
\TextAt{21}{253.316}{164.316}{K}%
\TextAt{9}{271.029}{167.557}{\Character{61}}%
\TextAt{20}{283.455}{178.915}{\Character{89}}%
\TextAt{27}{283.688}{153.012}{p}%
\TextAt{3}{287.677}{153.012}{\Character{106}}%
\NativeMathAt{290.028}{153.012}{7.97}{\infty}%
\TextAt{13}{300.723}{167.557}{K}%
\TextAt{3}{311.541}{172.494}{\Character{3}}%
\TextAt{27}{310.682}{164.602}{p}%
\TextAt{16}{318.93}{167.557}{\Character{2}}%
\TextAt{20}{330.885}{178.915}{\Character{89}}%
\TextAt{27}{331.188}{153.083}{p}%
\TextAt{25}{335.176}{153.083}{\Character{45}}%
\NativeMathAt{337.39}{153.083}{7.97}{\infty}%
\TextAt{13}{348.154}{167.557}{U}%
\TextAt{27}{356.104}{165.764}{p}%
\TextAt{13}{360.591}{167.557}{\Character{58}}%
\TextAt{8}{79.226}{127.877}{It}%
\TextAt{8}{91.5}{127.877}{induces}%
\TextAt{8}{132.777}{127.877}{a}%
\TextAt{8}{142.133}{127.877}{surjective}%
\TextAt{8}{194.146}{127.877}{homomorphism}%
\TextAt{13}{276.417}{127.877}{C}%
\TextAt{21}{284.792}{126.084}{K}%
\NativeArrowAt{296.315}{127.877}{11.955}{19.26}{\longrightarrow}%
\TextAt{13}{318.897}{127.877}{C}%
\TextAt{13}{328.115}{127.877}{l}%
\TextAt{21}{331.618}{126.084}{K}%
\TextAt{8}{343.319}{127.877}{with}%
\TextAt{8}{369.58}{127.877}{kernel}%
\TextAt{13}{403.648}{127.877}{I}%
\TextAt{21}{409.751}{132.991}{S}%
\NativeMathAt{414.901}{131.995}{5.978}{\infty}%
\TextAt{21}{408.811}{124.636}{K}%
\TextAt{13}{423.203}{127.877}{K}%
\TextAt{3}{434.021}{132.215}{\Character{3}}%
\TextAt{9}{438.753}{127.877}{\Character{47}}%
\TextAt{13}{444.606}{127.877}{K}%
\TextAt{3}{455.424}{132.215}{\Character{3}}%
\TextAt{8}{460.156}{127.877}{.}%
\TextAt{8}{468.471}{127.877}{We}%
\TextAt{8}{488.217}{127.877}{may}%
\TextAt{8}{513.187}{127.877}{also}%
\TextAt{8}{79.226}{110.831}{consider}%
\TextAt{8}{124.811}{110.831}{the}%
\TextAt{8}{144.968}{110.831}{surjective}%
\TextAt{8}{197.379}{110.831}{homomorphism}%
\TextAt{13}{216.15}{83.395}{I}%
\TextAt{21}{221.313}{81.601}{K}%
\NativeArrowAt{232.837}{83.395}{11.955}{19.26}{\longrightarrow}%
\TextAt{13}{255.419}{83.395}{J}%
\TextAt{9}{262.996}{83.395}{\Character{40}}%
\Vector{q 0 G 0 g 0.478 w 195.548 -680.762 m 205.398 -680.762 l S Q}%
\TextAt{16}{267.548}{83.395}{\Character{79}}%
\TextAt{9}{277.398}{83.395}{\Character{41}}%
\TextAt{13}{281.95}{83.395}{\Character{59}}%
\TextAt{13}{311.085}{83.395}{\Character{11}}%
\NativeArrowAt{321.947}{83.395}{11.955}{19.26}{\longmapsto}%
\TextAt{20}{344.529}{94.752}{\Character{89}}%
\TextAt{27}{350.173}{69.646}{p}%
\TextAt{14}{361.798}{83.395}{p}%
\TextAt{21}{367.781}{88.331}{v}%
\TextAt{28}{371.905}{87.335}{p}%
\TextAt{17}{375.987}{88.331}{\Character{40}}%
\TextAt{21}{379.281}{88.331}{\Character{11}}%
\TextAt{28}{384.719}{87.335}{p}%
\TextAt{17}{388.802}{88.331}{\Character{41}}%
\TextAt{13}{392.594}{83.395}{\Character{59}}%
\TextAt{8}{297.211}{43.033}{328}%
\Destination{522.36}{615.215}{nav.AMS.76}%
\Destination{78.23}{809.815}{nav.page.328}%
\Destination{79.36}{776.155}{nav.print.359}%
\Destination{79.36}{373.455}{nav.number.6.1.2}%
}

% Page 341 (body page 329)
\Page{%
\TextAt{2}{79.226}{796.943}{VI.}%
\TextAt{2}{98.012}{796.943}{Global}%
\TextAt{2}{133.248}{796.943}{Class}%
\TextAt{2}{161.841}{796.943}{Field}%
\TextAt{2}{189.572}{796.943}{Theory}%
\Vector{q 0 G 0 g 0.299 w 7.226 18.729 m 460.769 18.729 l S Q}%
\TextAt{8}{79.226}{761.124}{onto}%
\TextAt{8}{105.575}{761.124}{the}%
\TextAt{8}{125.732}{761.124}{replete}%
\TextAt{8}{164.096}{761.124}{ideal}%
\TextAt{8}{192.059}{761.124}{group}%
\TextAt{13}{225.235}{761.124}{J}%
\TextAt{9}{232.812}{761.124}{\Character{40}}%
\Vector{q 0 G 0 g 0.478 w 165.364 -3.033 m 175.214 -3.033 l S Q}%
\TextAt{16}{237.364}{761.124}{\Character{79}}%
\TextAt{9}{247.214}{761.124}{\Character{41}}%
\TextAt{8}{251.766}{761.124}{.}%
\TextAt{8}{260.215}{761.124}{Its}%
\TextAt{8}{277.502}{761.124}{kernel}%
\TextAt{8}{311.968}{761.124}{is}%
\TextAt{13}{211.633}{732.123}{I}%
\TextAt{17}{217.736}{737.059}{0}%
\TextAt{21}{216.796}{729.167}{K}%
\TextAt{9}{228.32}{732.123}{\Character{61}}%
\TextAt{16}{240.745}{732.123}{\Character{102}}%
\TextAt{9}{246.723}{732.123}{\Character{40}}%
\TextAt{13}{251.275}{732.123}{\Character{11}}%
\TextAt{27}{258.721}{730.329}{p}%
\TextAt{9}{263.208}{732.123}{\Character{41}}%
\NativeMathAt{271.081}{732.123}{11.955}{\in}%
\TextAt{13}{282.372}{732.123}{I}%
\TextAt{21}{287.535}{730.329}{K}%
\TextAt{16}{299.058}{732.123}{\Character{106}}%
\TextAt{16}{305.681}{732.123}{\Character{106}}%
\TextAt{13}{309.021}{732.123}{\Character{11}}%
\TextAt{27}{316.466}{730.329}{p}%
\TextAt{16}{320.953}{732.123}{\Character{106}}%
\TextAt{27}{324.274}{730.329}{p}%
\TextAt{9}{332.082}{732.123}{\Character{61}}%
\TextAt{9}{344.492}{732.123}{1}%
\TextAt{8}{354.258}{732.123}{for}%
\TextAt{8}{372.143}{732.123}{all}%
\TextAt{14}{388.402}{732.123}{p}%
\TextAt{16}{394.385}{732.123}{\Character{103}}%
\TextAt{8}{79.226}{703.121}{(see}%
\TextAt{8}{104.351}{703.121}{chap.}%
\TextAt{8}{136.9}{703.121}{III,}%
\TextAt{8}{159.012}{703.121}{§1).}%
\TextAt{8}{188.522}{703.121}{It}%
\TextAt{8}{202.852}{703.121}{takes}%
\TextAt{8}{234.492}{703.121}{principal}%
\TextAt{8}{284.926}{703.121}{idèles}%
\TextAt{8}{318.503}{703.121}{to}%
\TextAt{8}{334.47}{703.121}{replete}%
\TextAt{8}{374.491}{703.121}{principal}%
\TextAt{8}{424.925}{703.121}{ideals}%
\TextAt{8}{459.159}{703.121}{and}%
\TextAt{8}{483.579}{703.121}{induces}%
\TextAt{8}{526.911}{703.121}{a}%
\TextAt{8}{79.226}{686.075}{surjective}%
\TextAt{8}{131.711}{686.075}{homomorphism}%
\TextAt{13}{214.454}{686.075}{C}%
\TextAt{21}{222.829}{684.282}{K}%
\NativeArrowAt{234.477}{686.075}{11.955}{19.26}{\longrightarrow}%
\TextAt{13}{257.184}{686.075}{P}%
\TextAt{13}{266.354}{686.075}{ic}%
\TextAt{9}{275.387}{686.075}{\Character{40}}%
\Vector{q 0 G 0 g 0.478 w 207.939 -78.082 m 217.789 -78.082 l S Q}%
\TextAt{16}{279.939}{686.075}{\Character{79}}%
\TextAt{9}{289.789}{686.075}{\Character{41}}%
\TextAt{8}{298.312}{686.075}{onto}%
\TextAt{8}{324.734}{686.075}{the}%
\TextAt{8}{344.964}{686.075}{replete}%
\TextAt{8}{383.401}{686.075}{ideal}%
\TextAt{8}{411.438}{686.075}{class}%
\TextAt{8}{438.948}{686.075}{group,}%
\TextAt{8}{475.467}{686.075}{with}%
\TextAt{8}{502.2}{686.075}{kernel}%
\TextAt{13}{79.226}{669.029}{I}%
\TextAt{17}{85.329}{673.368}{0}%
\TextAt{21}{84.389}{666.009}{K}%
\TextAt{13}{92.592}{669.029}{K}%
\TextAt{3}{103.409}{673.368}{\Character{3}}%
\TextAt{9}{108.142}{669.029}{\Character{47}}%
\TextAt{13}{113.995}{669.029}{K}%
\TextAt{3}{124.812}{673.368}{\Character{3}}%
\TextAt{8}{129.545}{669.029}{.}%
\TextAt{8}{137.993}{669.029}{We}%
\TextAt{8}{158.137}{669.029}{therefore}%
\TextAt{8}{207.237}{669.029}{have}%
\TextAt{8}{234.232}{669.029}{the}%
\TextAt{12}{79.226}{643.017}{Proposition}%
\TextAt{12}{151.937}{643.017}{(1.3).}%
\TextAt{13}{187.253}{643.017}{C}%
\TextAt{13}{196.471}{643.017}{l}%
\TextAt{21}{199.974}{641.224}{K}%
\NativeMathAt{211.497}{643.017}{11.955}{\cong}%
\TextAt{13}{224.116}{643.017}{I}%
\TextAt{21}{229.279}{641.224}{K}%
\TextAt{9}{237.482}{643.017}{\Character{47}}%
\TextAt{13}{243.335}{643.017}{I}%
\TextAt{21}{249.438}{648.131}{S}%
\NativeMathAt{254.588}{647.135}{5.978}{\infty}%
\TextAt{21}{248.498}{639.776}{K}%
\TextAt{13}{262.89}{643.017}{K}%
\TextAt{3}{273.708}{647.356}{\Character{3}}%
\TextAt{8}{278.44}{643.017}{,}%
\TextAt{8}{285.589}{643.017}{and}%
\TextAt{13}{308.352}{643.017}{P}%
\TextAt{13}{317.51}{643.017}{ic}%
\TextAt{9}{326.554}{643.017}{\Character{40}}%
\Vector{q 0 G 0 g 0.478 w 259.107 -121.14 m 268.957 -121.14 l S Q}%
\TextAt{16}{331.107}{643.017}{\Character{79}}%
\TextAt{9}{340.957}{643.017}{\Character{41}}%
\NativeMathAt{348.83}{643.017}{11.955}{\cong}%
\TextAt{13}{361.449}{643.017}{I}%
\TextAt{21}{366.612}{641.224}{K}%
\TextAt{9}{374.815}{643.017}{\Character{47}}%
\TextAt{13}{380.668}{643.017}{I}%
\TextAt{17}{386.77}{647.356}{0}%
\TextAt{21}{385.83}{639.997}{K}%
\TextAt{13}{394.033}{643.017}{K}%
\TextAt{3}{404.851}{647.356}{\Character{3}}%
\TextAt{8}{409.583}{643.017}{.}%
\TextAt{8}{97.159}{617.005}{In}%
\TextAt{8}{110.72}{617.005}{contrast}%
\TextAt{8}{154.935}{617.005}{to}%
\TextAt{8}{168.185}{617.005}{the}%
\TextAt{8}{187.281}{617.005}{ideal}%
\TextAt{8}{214.185}{617.005}{class}%
\TextAt{8}{240.562}{617.005}{group,}%
\TextAt{8}{276.141}{617.005}{the}%
\TextAt{8}{295.238}{617.005}{idèle}%
\TextAt{8}{321.484}{617.005}{class}%
\TextAt{8}{347.861}{617.005}{group}%
\TextAt{8}{379.977}{617.005}{is}%
\TextAt{8}{390.681}{617.005}{not}%
\TextAt{8}{410.435}{617.005}{finite.}%
\TextAt{8}{444.544}{617.005}{But}%
\TextAt{8}{466.725}{617.005}{the}%
\TextAt{8}{485.822}{617.005}{finiteness}%
\TextAt{8}{79.226}{599.959}{of}%
\TextAt{8}{92.556}{599.959}{the}%
\TextAt{8}{112.713}{599.959}{former}%
\TextAt{8}{150.108}{599.959}{is}%
\TextAt{8}{161.872}{599.959}{reflected}%
\TextAt{8}{208.689}{599.959}{in}%
\TextAt{8}{222.342}{599.959}{terms}%
\TextAt{8}{254.919}{599.959}{of}%
\TextAt{8}{268.249}{599.959}{the}%
\TextAt{8}{288.406}{599.959}{latter}%
\TextAt{8}{320.278}{599.959}{as}%
\TextAt{8}{334.648}{599.959}{follows.}%
\TextAt{12}{79.226}{573.946}{Proposition}%
\TextAt{12}{152.055}{573.946}{(1.4).}%
\TextAt{13}{187.474}{573.946}{I}%
\TextAt{21}{192.636}{572.153}{K}%
\TextAt{9}{204.335}{573.946}{\Character{61}}%
\TextAt{13}{216.935}{573.946}{I}%
\TextAt{21}{223.038}{578.285}{S}%
\TextAt{21}{222.098}{570.926}{K}%
\TextAt{13}{230.301}{573.946}{K}%
\TextAt{3}{241.118}{578.285}{\Character{3}}%
\TextAt{8}{245.851}{573.946}{,}%
\TextAt{8}{253.128}{573.946}{i.e.,}%
\TextAt{13}{275.361}{573.946}{C}%
\TextAt{21}{283.736}{572.153}{K}%
\TextAt{9}{295.435}{573.946}{\Character{61}}%
\TextAt{13}{308.035}{573.946}{I}%
\TextAt{21}{314.138}{578.285}{S}%
\TextAt{21}{313.198}{570.926}{K}%
\TextAt{13}{321.4}{573.946}{K}%
\TextAt{3}{332.218}{578.285}{\Character{3}}%
\TextAt{9}{336.95}{573.946}{\Character{47}}%
\TextAt{13}{342.803}{573.946}{K}%
\TextAt{3}{353.621}{578.285}{\Character{3}}%
\TextAt{8}{358.353}{573.946}{,}%
\TextAt{8}{365.631}{573.946}{if}%
\TextAt{13}{376.457}{573.946}{S}%
\TextAt{8}{388.352}{573.946}{is}%
\TextAt{8}{400.219}{573.946}{a}%
\TextAt{8}{410.077}{573.946}{sufficiently}%
\TextAt{8}{468.772}{573.946}{big}%
\TextAt{8}{488.385}{573.946}{finite}%
\TextAt{8}{518.399}{573.946}{set}%
\TextAt{8}{79.226}{556.9}{of}%
\TextAt{8}{92.556}{556.9}{places}%
\TextAt{8}{127.083}{556.9}{of}%
\TextAt{13}{140.413}{556.9}{K}%
\TextAt{8}{151.23}{556.9}{.}%
\TextAt{19}{79.226}{521.922}{Proof.}%
\TextAt{8}{116}{521.922}{Let}%
\TextAt{14}{137.21}{521.922}{a}%
\TextAt{17}{143.193}{520.128}{1}%
\TextAt{13}{147.926}{521.922}{\Character{59}}%
\TextAt{13}{153.15}{521.922}{\Character{58}}%
\TextAt{13}{158.387}{521.922}{\Character{58}}%
\TextAt{13}{163.611}{521.922}{\Character{58}}%
\TextAt{13}{168.848}{521.922}{\Character{59}}%
\TextAt{14}{174.146}{521.922}{a}%
\TextAt{21}{180.129}{520.128}{h}%
\TextAt{8}{189.645}{521.922}{be}%
\TextAt{8}{205.81}{521.922}{ideals}%
\TextAt{8}{238.629}{521.922}{representing}%
\TextAt{8}{304.945}{521.922}{the}%
\TextAt{13}{325.342}{521.922}{h}%
\TextAt{8}{336.219}{521.922}{classes}%
\TextAt{8}{373.712}{521.922}{of}%
\TextAt{13}{387.282}{521.922}{J}%
\TextAt{21}{393.752}{520.128}{K}%
\TextAt{9}{401.955}{521.922}{\Character{47}}%
\TextAt{13}{407.808}{521.922}{P}%
\TextAt{21}{415.353}{520.128}{K}%
\TextAt{8}{423.556}{521.922}{.}%
\TextAt{8}{432.726}{521.922}{They}%
\TextAt{8}{463.201}{521.922}{are}%
\TextAt{8}{482.952}{521.922}{composed}%
\TextAt{8}{79.226}{504.876}{of}%
\TextAt{8}{91.995}{504.876}{a}%
\TextAt{8}{101.19}{504.876}{finite}%
\TextAt{8}{130.541}{504.876}{number}%
\TextAt{8}{172.577}{504.876}{of}%
\TextAt{8}{185.346}{504.876}{prime}%
\TextAt{8}{217.949}{504.876}{ideals}%
\TextAt{14}{249.966}{504.876}{p}%
\TextAt{17}{255.949}{503.083}{1}%
\TextAt{13}{260.681}{504.876}{\Character{59}}%
\TextAt{13}{265.905}{504.876}{\Character{58}}%
\TextAt{13}{271.142}{504.876}{\Character{58}}%
\TextAt{13}{276.378}{504.876}{\Character{58}}%
\TextAt{13}{281.603}{504.876}{\Character{59}}%
\TextAt{14}{286.902}{504.876}{p}%
\TextAt{21}{292.885}{503.083}{r}%
\TextAt{8}{297.44}{504.876}{.}%
\TextAt{8}{305.702}{504.876}{Now}%
\TextAt{8}{331.801}{504.876}{if}%
\TextAt{13}{341.964}{504.876}{S}%
\TextAt{8}{353.196}{504.876}{is}%
\TextAt{8}{364.399}{504.876}{any}%
\TextAt{8}{385.955}{504.876}{finite}%
\TextAt{8}{415.306}{504.876}{set}%
\TextAt{8}{433.013}{504.876}{of}%
\TextAt{8}{445.783}{504.876}{places}%
\TextAt{8}{479.748}{504.876}{containing}%
\TextAt{8}{79.226}{487.83}{these}%
\TextAt{8}{109.198}{487.83}{primes}%
\TextAt{8}{146.976}{487.83}{and}%
\TextAt{8}{169.739}{487.83}{the}%
\TextAt{8}{189.895}{487.83}{places}%
\TextAt{8}{224.422}{487.83}{at}%
\TextAt{8}{238.732}{487.83}{infinity,}%
\TextAt{8}{281.328}{487.83}{then}%
\TextAt{8}{307.988}{487.83}{one}%
\TextAt{8}{329.448}{487.83}{has}%
\TextAt{13}{350.322}{487.83}{I}%
\TextAt{21}{355.485}{486.037}{K}%
\TextAt{9}{367.008}{487.83}{\Character{61}}%
\TextAt{13}{379.433}{487.83}{I}%
\TextAt{21}{385.536}{492.168}{S}%
\TextAt{21}{384.596}{484.809}{K}%
\TextAt{13}{392.799}{487.83}{K}%
\TextAt{3}{403.617}{492.168}{\Character{3}}%
\TextAt{8}{408.349}{487.83}{.}%
\TextAt{8}{97.159}{467.795}{In}%
\TextAt{8}{113.125}{467.795}{order}%
\TextAt{8}{145.039}{467.795}{to}%
\TextAt{8}{160.694}{467.795}{see}%
\TextAt{8}{180.952}{467.795}{this,}%
\TextAt{8}{208.707}{467.795}{we}%
\TextAt{8}{227.279}{467.795}{use}%
\TextAt{8}{248.84}{467.795}{the}%
\TextAt{8}{270.341}{467.795}{isomorphism}%
\TextAt{13}{340.105}{467.795}{I}%
\TextAt{21}{345.268}{466.002}{K}%
\TextAt{9}{353.47}{467.795}{\Character{47}}%
\TextAt{13}{359.323}{467.795}{I}%
\TextAt{21}{365.426}{472.909}{S}%
\NativeMathAt{370.576}{471.913}{5.978}{\infty}%
\TextAt{21}{364.486}{464.554}{K}%
\NativeMathAt{384.491}{467.795}{11.955}{\cong}%
\TextAt{13}{399.402}{467.795}{J}%
\TextAt{21}{405.872}{466.002}{K}%
\TextAt{8}{414.074}{467.795}{.}%
\TextAt{8}{426.557}{467.795}{If}%
\TextAt{13}{439.594}{467.795}{\Character{11}}%
\NativeMathAt{452.728}{467.795}{11.955}{\in}%
\TextAt{13}{466.31}{467.795}{I}%
\TextAt{21}{471.473}{466.002}{K}%
\TextAt{8}{479.676}{467.795}{,}%
\TextAt{8}{488.506}{467.795}{then}%
\TextAt{8}{516.51}{467.795}{the}%
\TextAt{8}{79.226}{450.749}{corresponding}%
\TextAt{8}{155.137}{450.749}{ideal}%
\TextAt{9}{183.825}{450.749}{\Character{40}}%
\TextAt{13}{188.378}{450.749}{\Character{11}}%
\TextAt{9}{195.899}{450.749}{\Character{41}}%
\TextAt{9}{204.997}{450.749}{\Character{61}}%
\TextAt{20}{218.669}{459.716}{\Character{81}}%
\TextAt{27}{229.96}{447.262}{p}%
\TextAt{25}{233.949}{447.262}{\Character{45}}%
\NativeMathAt{236.162}{447.262}{7.97}{\infty}%
\TextAt{14}{247.121}{450.749}{p}%
\TextAt{21}{253.105}{455.088}{v}%
\TextAt{28}{257.228}{454.091}{p}%
\TextAt{17}{261.311}{455.088}{\Character{40}}%
\TextAt{21}{264.604}{455.088}{\Character{11}}%
\TextAt{28}{270.043}{454.091}{p}%
\TextAt{17}{274.126}{455.088}{\Character{41}}%
\TextAt{8}{282.539}{450.749}{belongs}%
\TextAt{8}{325.275}{450.749}{to}%
\TextAt{8}{340.31}{450.749}{some}%
\TextAt{8}{370.361}{450.749}{class}%
\TextAt{14}{398.523}{450.749}{a}%
\TextAt{21}{404.506}{448.956}{i}%
\TextAt{13}{407.888}{450.749}{P}%
\TextAt{21}{415.433}{448.956}{K}%
\TextAt{8}{423.635}{450.749}{,}%
\TextAt{8}{431.691}{450.749}{i.e.,}%
\TextAt{9}{454.702}{450.749}{\Character{40}}%
\TextAt{13}{459.254}{450.749}{\Character{11}}%
\TextAt{9}{466.776}{450.749}{\Character{41}}%
\TextAt{9}{475.874}{450.749}{\Character{61}}%
\TextAt{14}{489.546}{450.749}{a}%
\TextAt{21}{495.529}{448.956}{i}%
\TextAt{9}{498.91}{450.749}{\Character{40}}%
\TextAt{13}{503.462}{450.749}{a}%
\TextAt{9}{509.607}{450.749}{\Character{41}}%
\TextAt{8}{518.782}{450.749}{for}%
\TextAt{8}{79.226}{433.703}{some}%
\TextAt{8}{109.86}{433.703}{principal}%
\TextAt{8}{159.945}{433.703}{ideal}%
\TextAt{9}{189.216}{433.703}{\Character{40}}%
\TextAt{13}{193.769}{433.703}{a}%
\TextAt{9}{199.914}{433.703}{\Character{41}}%
\TextAt{8}{204.466}{433.703}{.}%
\TextAt{8}{216.838}{433.703}{The}%
\TextAt{8}{242.2}{433.703}{idèle}%
\TextAt{13}{270.813}{433.703}{\Character{11}}%
\TextAt{3}{278.335}{438.042}{\Character{48}}%
\TextAt{9}{286.68}{433.703}{\Character{61}}%
\TextAt{13}{301.334}{433.703}{\Character{11}a}%
\TextAt{3}{315.001}{438.042}{\Character{0}}%
\TextAt{17}{321.587}{438.042}{1}%
\TextAt{8}{331.525}{433.703}{is}%
\TextAt{8}{344.597}{433.703}{mapped}%
\TextAt{8}{390.45}{433.703}{by}%
\TextAt{13}{408.017}{433.703}{I}%
\TextAt{21}{413.18}{431.91}{K}%
\NativeArrowAt{426.932}{433.703}{11.955}{19.26}{\longrightarrow}%
\TextAt{13}{451.743}{433.703}{J}%
\TextAt{21}{458.213}{431.91}{K}%
\TextAt{8}{471.621}{433.703}{to}%
\TextAt{8}{487.239}{433.703}{the}%
\TextAt{8}{508.704}{433.703}{ideal}%
\TextAt{14}{79.226}{416.657}{a}%
\TextAt{21}{85.209}{414.864}{i}%
\TextAt{9}{91.911}{416.657}{\Character{61}}%
\TextAt{20}{104.337}{425.624}{\Character{81}}%
\TextAt{27}{115.628}{413.17}{p}%
\TextAt{25}{119.617}{413.17}{\Character{45}}%
\NativeMathAt{121.831}{413.17}{7.97}{\infty}%
\TextAt{14}{132.789}{416.657}{p}%
\TextAt{21}{138.773}{421.068}{v}%
\TextAt{28}{142.896}{420.071}{p}%
\TextAt{17}{146.979}{421.068}{\Character{40}}%
\TextAt{21}{150.272}{421.068}{\Character{11}}%
\TextAt{23}{155.711}{423.88}{\Character{48}}%
\TextAt{28}{155.711}{419.075}{p}%
\TextAt{17}{159.794}{421.068}{\Character{41}}%
\TextAt{8}{163.585}{416.657}{.}%
\TextAt{8}{172.008}{416.657}{Since}%
\TextAt{8}{202.488}{416.657}{the}%
\TextAt{8}{222.567}{416.657}{prime}%
\TextAt{8}{255.653}{416.657}{ideals}%
\TextAt{8}{288.154}{416.657}{occurring}%
\TextAt{8}{339.782}{416.657}{in}%
\TextAt{14}{353.358}{416.657}{a}%
\TextAt{21}{359.341}{414.864}{i}%
\TextAt{8}{366.542}{416.657}{lie}%
\TextAt{8}{382.066}{416.657}{in}%
\TextAt{13}{395.642}{416.657}{S}%
\TextAt{8}{403.537}{416.657}{,}%
\TextAt{8}{410.624}{416.657}{we}%
\TextAt{8}{427.775}{416.657}{have}%
\TextAt{13}{454.692}{416.657}{v}%
\TextAt{27}{460.351}{414.864}{p}%
\TextAt{9}{464.838}{416.657}{\Character{40}}%
\TextAt{13}{469.39}{416.657}{\Character{11}}%
\TextAt{3}{476.912}{420.996}{\Character{48}}%
\TextAt{27}{476.836}{413.702}{p}%
\TextAt{9}{481.323}{416.657}{\Character{41}}%
\TextAt{9}{489.178}{416.657}{\Character{61}}%
\TextAt{9}{501.599}{416.657}{0}%
\TextAt{8}{507.474}{416.657}{,}%
\TextAt{8}{514.562}{416.657}{i.e.,}%
\TextAt{13}{79.226}{399.611}{\Character{11}}%
\TextAt{3}{86.748}{403.95}{\Character{48}}%
\TextAt{27}{86.672}{396.656}{p}%
\NativeMathAt{94.479}{399.611}{11.955}{\in}%
\TextAt{13}{105.77}{399.611}{U}%
\TextAt{27}{113.721}{397.818}{p}%
\TextAt{8}{122.105}{399.611}{for}%
\TextAt{8}{139.99}{399.611}{all}%
\TextAt{14}{156.249}{399.611}{p}%
\NativeMathAt{165.553}{399.611}{11.955}{\notin}%
\TextAt{13}{176.844}{399.611}{S}%
\TextAt{8}{184.739}{399.611}{.}%
\TextAt{8}{193.187}{399.611}{Hence}%
\TextAt{13}{227.965}{399.611}{\Character{11}}%
\TextAt{3}{235.487}{403.95}{\Character{48}}%
\TextAt{9}{241.603}{399.611}{\Character{61}}%
\TextAt{13}{254.028}{399.611}{\Character{11}a}%
\TextAt{3}{267.695}{403.95}{\Character{0}}%
\TextAt{17}{274.281}{403.95}{1}%
\NativeMathAt{282.334}{399.611}{11.955}{\in}%
\TextAt{13}{293.625}{399.611}{I}%
\TextAt{21}{299.728}{403.95}{S}%
\TextAt{21}{298.788}{396.591}{K}%
\TextAt{8}{306.991}{399.611}{,}%
\TextAt{8}{314.14}{399.611}{and}%
\TextAt{8}{336.903}{399.611}{thus}%
\TextAt{13}{362.654}{399.611}{\Character{11}}%
\NativeMathAt{373.497}{399.611}{11.955}{\in}%
\TextAt{13}{384.788}{399.611}{I}%
\TextAt{21}{390.891}{403.95}{S}%
\TextAt{21}{389.951}{396.591}{K}%
\TextAt{13}{398.153}{399.611}{K}%
\TextAt{3}{408.971}{403.95}{\Character{3}}%
\TextAt{8}{413.703}{399.611}{.}%
\Vector{q 0 G 0 g 0.399 w 452.334 -374.389 m 452.334 -366.319 l S Q}%
\Vector{q 0 G 0 g 0.399 w 452.534 -366.518 m 459.707 -366.518 l S Q}%
\Vector{q 0 G 0 g 0.399 w 452.534 -374.19 m 459.707 -374.19 l S Q}%
\Vector{q 0 G 0 g 0.399 w 459.906 -374.389 m 459.906 -366.319 l S Q}%
\TextAt{8}{97.159}{370.61}{The}%
\TextAt{8}{121.32}{370.61}{idèle}%
\TextAt{8}{148.732}{370.61}{group}%
\TextAt{8}{182.015}{370.61}{comes}%
\TextAt{8}{216.649}{370.61}{equipped}%
\TextAt{8}{266.824}{370.61}{with}%
\TextAt{8}{293.591}{370.61}{a}%
\TextAt{8}{303.454}{370.61}{canonical}%
\TextAt{8}{354.944}{370.61}{topology.}%
\TextAt{8}{406.979}{370.61}{A}%
\TextAt{8}{419.759}{370.61}{basic}%
\TextAt{8}{449.192}{370.61}{system}%
\TextAt{8}{488.117}{370.61}{of}%
\TextAt{8}{501.555}{370.61}{neigh-}%
\TextAt{8}{79.226}{353.564}{bourhoods}%
\TextAt{8}{136.85}{353.564}{of}%
\TextAt{9}{150.18}{353.564}{1}%
\NativeMathAt{159.354}{353.564}{11.955}{\in}%
\TextAt{13}{170.645}{353.564}{I}%
\TextAt{21}{175.808}{351.771}{K}%
\TextAt{8}{187.908}{353.564}{is}%
\TextAt{8}{199.672}{353.564}{given}%
\TextAt{8}{230.241}{353.564}{by}%
\TextAt{8}{246.5}{353.564}{the}%
\TextAt{8}{266.657}{353.564}{sets}%
\TextAt{20}{251.157}{335.92}{\Character{89}}%
\TextAt{27}{251.18}{310.549}{p}%
\NativeMathAt{255.169}{310.549}{7.97}{\in}%
\TextAt{21}{260.814}{310.549}{S}%
\TextAt{13}{268.426}{324.563}{W}%
\TextAt{27}{279.477}{322.77}{p}%
\TextAt{16}{286.621}{324.563}{\Character{2}}%
\TextAt{20}{298.576}{335.92}{\Character{89}}%
\TextAt{27}{298.599}{310.017}{p}%
\TextAt{17}{303.528}{310.017}{\Character{47}}%
\NativeMathAt{302.587}{310.017}{7.97}{\in}%
\TextAt{21}{308.233}{310.017}{S}%
\TextAt{13}{315.844}{324.563}{U}%
\TextAt{27}{323.795}{322.77}{p}%
\NativeMathAt{331.602}{324.563}{11.955}{\subseteq}%
\TextAt{13}{344.222}{324.563}{I}%
\TextAt{21}{349.384}{322.77}{K}%
\TextAt{13}{357.587}{324.563}{\Character{59}}%
\TextAt{8}{79.226}{285.397}{where}%
\TextAt{13}{112.585}{285.397}{S}%
\TextAt{8}{123.928}{285.397}{runs}%
\TextAt{8}{149.552}{285.397}{through}%
\TextAt{8}{193.336}{285.397}{the}%
\TextAt{8}{213.043}{285.397}{finite}%
\TextAt{8}{242.504}{285.397}{sets}%
\TextAt{8}{264.937}{285.397}{of}%
\TextAt{8}{277.816}{285.397}{places}%
\TextAt{8}{311.893}{285.397}{of}%
\TextAt{13}{324.773}{285.397}{K}%
\TextAt{8}{339.038}{285.397}{which}%
\TextAt{8}{372.074}{285.397}{contain}%
\TextAt{8}{412.929}{285.397}{all}%
\TextAt{14}{428.738}{285.397}{p}%
\TextAt{16}{438.042}{285.397}{\Character{106}}%
\NativeMathAt{444.665}{285.397}{11.955}{\infty}%
\TextAt{8}{456.639}{285.397}{,}%
\TextAt{8}{463.428}{285.397}{and}%
\TextAt{13}{485.741}{285.397}{W}%
\TextAt{27}{496.792}{283.604}{p}%
\NativeMathAt{504.6}{285.397}{11.955}{\subseteq}%
\TextAt{13}{517.22}{285.397}{K}%
\TextAt{3}{528.037}{289.736}{\Character{3}}%
\TextAt{27}{527.179}{282.442}{p}%
\TextAt{8}{79.226}{268.351}{is}%
\TextAt{8}{91.603}{268.351}{a}%
\TextAt{8}{101.972}{268.351}{basic}%
\TextAt{8}{131.911}{268.351}{system}%
\TextAt{8}{171.343}{268.351}{of}%
\TextAt{8}{185.286}{268.351}{neighbourhoods}%
\TextAt{8}{270.518}{268.351}{of}%
\TextAt{9}{284.462}{268.351}{1}%
\NativeMathAt{294.681}{268.351}{11.955}{\in}%
\TextAt{13}{307.017}{268.351}{K}%
\TextAt{3}{317.834}{272.69}{\Character{3}}%
\TextAt{27}{316.976}{265.396}{p}%
\TextAt{8}{322.567}{268.351}{.}%
\TextAt{8}{332.855}{268.351}{The}%
\TextAt{8}{357.522}{268.351}{groups}%
\TextAt{13}{395.925}{268.351}{U}%
\TextAt{27}{403.876}{266.558}{p}%
\TextAt{8}{412.873}{268.351}{are}%
\TextAt{8}{432.997}{268.351}{compact}%
\TextAt{8}{480.439}{268.351}{for}%
\TextAt{14}{498.937}{268.351}{p}%
\NativeMathAt{509.286}{268.351}{11.955}{\notin}%
\TextAt{13}{521.622}{268.351}{S}%
\TextAt{8}{529.518}{268.351}{.}%
\TextAt{8}{79.226}{251.305}{Therefore}%
\TextAt{8}{132.036}{251.305}{the}%
\TextAt{8}{152.005}{251.305}{same}%
\TextAt{8}{181.143}{251.305}{is}%
\TextAt{8}{192.719}{251.305}{true}%
\TextAt{8}{217.243}{251.305}{of}%
\TextAt{8}{230.386}{251.305}{the}%
\TextAt{8}{250.354}{251.305}{group}%
\TextAt{20}{283.342}{260.272}{\Character{81}}%
\TextAt{27}{294.633}{247.818}{p}%
\TextAt{17}{299.563}{247.818}{\Character{47}}%
\NativeMathAt{298.622}{247.818}{7.97}{\in}%
\TextAt{21}{304.268}{247.818}{S}%
\TextAt{13}{312.354}{251.305}{U}%
\TextAt{27}{320.304}{249.512}{p}%
\TextAt{8}{324.791}{251.305}{.}%
\TextAt{8}{333.177}{251.305}{If}%
\TextAt{8}{344.682}{251.305}{the}%
\TextAt{13}{364.65}{251.305}{W}%
\TextAt{27}{375.702}{249.512}{p}%
\TextAt{8}{380.189}{251.305}{,}%
\TextAt{8}{387.188}{251.305}{for}%
\TextAt{14}{404.885}{251.305}{p}%
\TextAt{16}{414.189}{251.305}{\Character{106}}%
\NativeMathAt{420.812}{251.305}{11.955}{\infty}%
\TextAt{8}{432.786}{251.305}{,}%
\TextAt{8}{439.785}{251.305}{are}%
\TextAt{8}{459.108}{251.305}{bounded,}%
\TextAt{8}{510.007}{251.305}{then}%
\TextAt{20}{79.226}{243.226}{\Character{81}}%
\TextAt{27}{90.517}{230.772}{p}%
\NativeMathAt{94.506}{230.772}{7.97}{\in}%
\TextAt{21}{100.151}{230.772}{S}%
\TextAt{13}{108.238}{234.259}{W}%
\TextAt{27}{119.289}{232.466}{p}%
\TextAt{16}{126.773}{234.259}{\Character{2}}%
\TextAt{20}{139.068}{243.226}{\Character{81}}%
\TextAt{27}{150.359}{230.772}{p}%
\TextAt{17}{155.289}{230.772}{\Character{47}}%
\NativeMathAt{154.348}{230.772}{7.97}{\in}%
\TextAt{21}{159.994}{230.772}{S}%
\TextAt{13}{168.08}{234.259}{U}%
\TextAt{27}{176.03}{232.466}{p}%
\TextAt{8}{184.913}{234.259}{is}%
\TextAt{8}{197.176}{234.259}{a}%
\TextAt{8}{207.43}{234.259}{neighbourhood}%
\TextAt{8}{287.933}{234.259}{of}%
\TextAt{9}{301.762}{234.259}{1}%
\TextAt{8}{312.011}{234.259}{in}%
\TextAt{13}{326.163}{234.259}{I}%
\TextAt{21}{331.326}{232.466}{K}%
\TextAt{8}{343.925}{234.259}{whose}%
\TextAt{8}{378.95}{234.259}{closure}%
\TextAt{8}{418.531}{234.259}{is}%
\TextAt{8}{430.794}{234.259}{compact.}%
\TextAt{8}{483.67}{234.259}{Therefore}%
\TextAt{13}{79.226}{217.213}{I}%
\TextAt{21}{84.389}{215.42}{K}%
\TextAt{8}{96.489}{217.213}{is}%
\TextAt{8}{108.253}{217.213}{a}%
\TextAt{12}{118.008}{217.213}{locally}%
\TextAt{12}{160.437}{217.213}{compact}%
\TextAt{12}{214.068}{217.213}{topological}%
\TextAt{12}{282.296}{217.213}{group}%
\TextAt{8}{316.201}{217.213}{.}%
\TextAt{12}{79.226}{191.201}{Proposition}%
\TextAt{12}{151.937}{191.201}{(1.5).}%
\TextAt{13}{187.253}{191.201}{K}%
\TextAt{3}{198.071}{195.539}{\Character{3}}%
\TextAt{8}{206.7}{191.201}{is}%
\TextAt{8}{218.464}{191.201}{a}%
\TextAt{8}{228.22}{191.201}{discrete,}%
\TextAt{8}{274.45}{191.201}{and}%
\TextAt{8}{297.213}{191.201}{therefore}%
\TextAt{8}{346.313}{191.201}{closed,}%
\TextAt{8}{384.091}{191.201}{subgroup}%
\TextAt{8}{434.889}{191.201}{of}%
\TextAt{13}{448.219}{191.201}{I}%
\TextAt{21}{453.382}{189.408}{K}%
\TextAt{8}{461.584}{191.201}{.}%
\TextAt{19}{79.226}{156.222}{Proof.}%
\TextAt{8}{116}{156.222}{It}%
\TextAt{8}{129.442}{156.222}{is}%
\TextAt{8}{141.976}{156.222}{enough}%
\TextAt{8}{183.07}{156.222}{to}%
\TextAt{8}{198.15}{156.222}{show}%
\TextAt{8}{227.923}{156.222}{that}%
\TextAt{9}{254.061}{156.222}{1}%
\NativeMathAt{264.547}{156.222}{11.955}{\in}%
\TextAt{13}{277.149}{156.222}{I}%
\TextAt{21}{282.312}{154.429}{K}%
\TextAt{8}{295.182}{156.222}{has}%
\TextAt{8}{316.825}{156.222}{a}%
\TextAt{8}{327.35}{156.222}{neighbourhood}%
\TextAt{8}{408.124}{156.222}{which}%
\TextAt{8}{442.38}{156.222}{contains}%
\TextAt{8}{489.069}{156.222}{no}%
\TextAt{8}{506.098}{156.222}{other}%
\TextAt{8}{79.226}{139.176}{principal}%
\TextAt{8}{128.003}{139.176}{idèle}%
\TextAt{8}{155.309}{139.176}{besides}%
\TextAt{9}{195.418}{139.176}{1}%
\TextAt{8}{201.271}{139.176}{.}%
\TextAt{16}{151.906}{110.175}{\Character{85}}%
\TextAt{9}{163.896}{110.175}{\Character{61}}%
\TextAt{16}{176.322}{110.175}{\Character{102}}%
\TextAt{13}{182.299}{110.175}{\Character{11}}%
\NativeMathAt{193.142}{110.175}{11.955}{\in}%
\TextAt{13}{204.433}{110.175}{I}%
\TextAt{21}{209.596}{108.382}{K}%
\TextAt{16}{221.119}{110.175}{\Character{106}}%
\TextAt{16}{227.742}{110.175}{\Character{106}}%
\TextAt{13}{231.082}{110.175}{\Character{11}}%
\TextAt{27}{238.527}{108.382}{p}%
\TextAt{16}{243.014}{110.175}{\Character{106}}%
\TextAt{27}{246.335}{108.382}{p}%
\TextAt{9}{254.143}{110.175}{\Character{61}}%
\TextAt{9}{266.553}{110.175}{1}%
\TextAt{8}{276.319}{110.175}{for}%
\TextAt{14}{294.204}{110.175}{p}%
\TextAt{18}{303.507}{110.175}{\Character{45}}%
\NativeMathAt{310.149}{110.175}{11.955}{\infty}%
\TextAt{13}{322.104}{110.175}{\Character{59}}%
\TextAt{16}{339.304}{110.175}{\Character{106}}%
\TextAt{13}{342.625}{110.175}{\Character{11}}%
\TextAt{27}{350.07}{108.382}{p}%
\TextAt{16}{357.214}{110.175}{\Character{0}}%
\TextAt{9}{369.169}{110.175}{1}%
\TextAt{16}{375.022}{110.175}{\Character{106}}%
\TextAt{27}{378.343}{108.382}{p}%
\TextAt{13}{386.15}{110.175}{\Character{60}}%
\TextAt{9}{398.576}{110.175}{1}%
\TextAt{8}{408.326}{110.175}{for}%
\TextAt{14}{426.211}{110.175}{p}%
\TextAt{16}{435.515}{110.175}{\Character{106}}%
\NativeMathAt{442.138}{110.175}{11.955}{\infty}%
\TextAt{16}{454.093}{110.175}{\Character{103}}%
\TextAt{8}{297.211}{43.033}{329}%
\Destination{78.23}{809.815}{nav.page.329}%
\Destination{79.36}{776.155}{nav.print.360}%
\Destination{79.36}{385.655}{nav.print.361}%
\Destination{79.36}{657.055}{nav.number.6.1.3}%
\Destination{79.36}{587.955}{nav.number.6.1.4}%
\Destination{79.36}{205.155}{nav.number.6.1.5}%
\Hotlink{136.4}{700.955}{14.294}{11.047}{nav.chapter.3}%
\Hotlink{164.191}{700.955}{6.858}{11.047}{nav.section.3.1}%
}

% Page 342 (body page 330)
\Page{%
\TextAt{2}{79.226}{796.943}{VI.}%
\TextAt{2}{98.012}{796.943}{Global}%
\TextAt{2}{133.248}{796.943}{Class}%
\TextAt{2}{161.841}{796.943}{Field}%
\TextAt{2}{189.572}{796.943}{Theory}%
\Vector{q 0 G 0 g 0.299 w 7.226 18.729 m 460.769 18.729 l S Q}%
\TextAt{8}{79.226}{761.124}{is}%
\TextAt{8}{91.125}{761.124}{such}%
\TextAt{8}{117.656}{761.124}{a}%
\TextAt{8}{127.546}{761.124}{neighbourhood.}%
\TextAt{8}{212.505}{761.124}{For}%
\TextAt{8}{233.61}{761.124}{if}%
\TextAt{8}{244.468}{761.124}{we}%
\TextAt{8}{261.83}{761.124}{had}%
\TextAt{8}{284.727}{761.124}{a}%
\TextAt{8}{294.617}{761.124}{principal}%
\TextAt{8}{343.529}{761.124}{idèle}%
\TextAt{13}{370.97}{761.124}{x}%
\NativeMathAt{381.172}{761.124}{11.955}{\in}%
\TextAt{16}{392.673}{761.124}{\Character{85}}%
\TextAt{8}{405.394}{761.124}{different}%
\TextAt{8}{451.699}{761.124}{from}%
\TextAt{9}{479.474}{761.124}{1}%
\TextAt{8}{485.327}{761.124}{,}%
\TextAt{8}{492.645}{761.124}{then}%
\TextAt{8}{519.439}{761.124}{we}%
\TextAt{8}{79.226}{744.078}{get}%
\TextAt{8}{98.737}{744.078}{the}%
\TextAt{8}{118.893}{744.078}{contradiction}%
\TextAt{9}{199.737}{718.768}{1}%
\TextAt{9}{208.895}{718.768}{\Character{61}}%
\TextAt{20}{221.336}{730.126}{\Character{89}}%
\TextAt{27}{226.98}{705.02}{p}%
\TextAt{16}{238.605}{718.768}{\Character{106}}%
\TextAt{13}{241.926}{718.768}{x}%
\TextAt{16}{251.234}{718.768}{\Character{0}}%
\TextAt{9}{263.189}{718.768}{1}%
\TextAt{16}{269.042}{718.768}{\Character{106}}%
\TextAt{27}{272.363}{716.975}{p}%
\TextAt{9}{280.171}{718.768}{\Character{61}}%
\TextAt{20}{292.596}{730.126}{\Character{89}}%
\TextAt{27}{292.899}{704.293}{p}%
\TextAt{25}{296.888}{704.293}{\Character{45}}%
\NativeMathAt{299.102}{704.293}{7.97}{\infty}%
\TextAt{16}{309.865}{718.768}{\Character{106}}%
\TextAt{13}{313.186}{718.768}{x}%
\TextAt{16}{322.495}{718.768}{\Character{0}}%
\TextAt{9}{334.45}{718.768}{1}%
\TextAt{16}{340.303}{718.768}{\Character{106}}%
\TextAt{27}{343.624}{716.975}{p}%
\TextAt{16}{350.767}{718.768}{\Character{1}}%
\TextAt{20}{356.745}{730.126}{\Character{89}}%
\TextAt{27}{356.978}{704.223}{p}%
\TextAt{3}{360.967}{704.223}{\Character{106}}%
\NativeMathAt{363.318}{704.223}{7.97}{\infty}%
\TextAt{16}{374.013}{718.768}{\Character{106}}%
\TextAt{13}{377.334}{718.768}{x}%
\TextAt{16}{386.643}{718.768}{\Character{0}}%
\TextAt{9}{398.598}{718.768}{1}%
\TextAt{16}{404.451}{718.768}{\Character{106}}%
\TextAt{27}{407.772}{716.975}{p}%
\TextAt{13}{208.911}{684.496}{\Character{60}}%
\TextAt{20}{221.336}{695.854}{\Character{89}}%
\TextAt{27}{221.639}{670.022}{p}%
\TextAt{25}{225.627}{670.022}{\Character{45}}%
\NativeMathAt{227.841}{670.022}{7.97}{\infty}%
\TextAt{16}{238.605}{684.496}{\Character{106}}%
\TextAt{13}{241.926}{684.496}{x}%
\TextAt{16}{251.234}{684.496}{\Character{0}}%
\TextAt{9}{263.189}{684.496}{1}%
\TextAt{16}{269.042}{684.496}{\Character{106}}%
\TextAt{27}{272.363}{682.703}{p}%
\NativeMathAt{280.171}{684.496}{11.955}{\leq}%
\TextAt{20}{292.79}{695.854}{\Character{89}}%
\TextAt{27}{293.093}{670.022}{p}%
\TextAt{25}{297.082}{670.022}{\Character{45}}%
\NativeMathAt{299.296}{670.022}{7.97}{\infty}%
\TextAt{8}{310.059}{684.496}{max}%
\TextAt{16}{331.853}{684.496}{\Character{102}\Character{106}}%
\TextAt{13}{341.152}{684.496}{x}%
\TextAt{16}{347.804}{684.496}{\Character{106}}%
\TextAt{27}{351.125}{682.703}{p}%
\TextAt{13}{355.612}{684.496}{\Character{59}}%
\TextAt{9}{360.856}{684.496}{1}%
\TextAt{16}{366.709}{684.496}{\Character{103}}%
\TextAt{9}{376.007}{684.496}{\Character{61}}%
\TextAt{9}{388.417}{684.496}{1}%
\TextAt{13}{394.286}{684.496}{\Character{58}}%
\TextAt{8}{79.226}{648.088}{That}%
\TextAt{8}{108.628}{648.088}{the}%
\TextAt{8}{128.92}{648.088}{subgroup}%
\TextAt{8}{179.854}{648.088}{is}%
\TextAt{8}{191.753}{648.088}{closed}%
\TextAt{8}{226.416}{648.088}{follows}%
\TextAt{8}{264.987}{648.088}{for}%
\TextAt{8}{283.008}{648.088}{a}%
\TextAt{8}{292.899}{648.088}{completely}%
\TextAt{8}{351.89}{648.088}{general}%
\TextAt{8}{392.351}{648.088}{reason:}%
\TextAt{8}{433.66}{648.088}{since}%
\TextAt{9}{462.465}{648.088}{\Character{40}}%
\TextAt{13}{467.017}{648.088}{x\Character{59}}%
\TextAt{13}{478.889}{648.088}{y}%
\TextAt{9}{485.05}{648.088}{\Character{41}}%
\NativeArrowAt{493.155}{648.088}{11.955}{11.955}{\mapsto}%
\TextAt{13}{508.662}{648.088}{xy}%
\TextAt{3}{521.451}{652.426}{\Character{0}}%
\TextAt{17}{528.037}{652.426}{1}%
\TextAt{8}{79.226}{631.042}{is}%
\TextAt{8}{91.331}{631.042}{continuous,}%
\TextAt{8}{153.613}{631.042}{there}%
\TextAt{8}{183.866}{631.042}{is}%
\TextAt{8}{195.971}{631.042}{a}%
\TextAt{8}{206.067}{631.042}{neighbourhood}%
\TextAt{13}{286.412}{631.042}{V}%
\TextAt{8}{300.076}{631.042}{of}%
\TextAt{9}{313.747}{631.042}{1}%
\TextAt{8}{323.838}{631.042}{such}%
\TextAt{8}{350.576}{631.042}{that}%
\TextAt{13}{376.286}{631.042}{V}%
\TextAt{13}{385.695}{631.042}{V}%
\TextAt{3}{395.137}{635.38}{\Character{0}}%
\TextAt{17}{401.724}{635.38}{1}%
\NativeMathAt{410.358}{631.042}{11.955}{\subseteq}%
\TextAt{16}{423.545}{631.042}{\Character{85}}%
\TextAt{8}{432.227}{631.042}{.}%
\TextAt{8}{441.698}{631.042}{For}%
\TextAt{8}{463.009}{631.042}{every}%
\TextAt{13}{494.242}{631.042}{y}%
\NativeMathAt{504.28}{631.042}{11.955}{\in}%
\TextAt{13}{516.152}{631.042}{I}%
\TextAt{21}{521.315}{629.248}{K}%
\TextAt{8}{529.518}{631.042}{,}%
\TextAt{8}{79.226}{613.996}{the}%
\TextAt{8}{99.842}{613.996}{neighbourhood}%
\TextAt{13}{180.305}{613.996}{y}%
\TextAt{13}{186.426}{613.996}{V}%
\TextAt{8}{200.224}{613.996}{then}%
\TextAt{8}{227.343}{613.996}{contains}%
\TextAt{8}{273.722}{613.996}{at}%
\TextAt{8}{288.491}{613.996}{most}%
\TextAt{8}{317.631}{613.996}{one}%
\TextAt{13}{339.55}{613.996}{x}%
\NativeMathAt{350.305}{613.996}{11.955}{\in}%
\TextAt{13}{362.379}{613.996}{K}%
\TextAt{3}{373.196}{618.334}{\Character{3}}%
\TextAt{8}{377.929}{613.996}{.}%
\TextAt{8}{387.754}{613.996}{Indeed,}%
\TextAt{8}{429.61}{613.996}{from}%
\TextAt{13}{457.709}{613.996}{x}%
\TextAt{17}{464.361}{612.202}{1}%
\TextAt{9}{473.197}{613.996}{\Character{61}}%
\TextAt{13}{486.405}{613.996}{y}%
\TextAt{13}{492.526}{613.996}{v}%
\TextAt{17}{498.201}{612.202}{1}%
\TextAt{13}{502.933}{613.996}{\Character{59}}%
\TextAt{13}{508.157}{613.996}{x}%
\TextAt{17}{514.829}{612.202}{2}%
\TextAt{9}{523.665}{613.996}{\Character{61}}%
\TextAt{13}{79.226}{596.95}{y}%
\TextAt{13}{85.347}{596.95}{v}%
\TextAt{17}{91.022}{595.156}{2}%
\NativeMathAt{99.075}{596.95}{11.955}{\in}%
\TextAt{13}{110.366}{596.95}{K}%
\TextAt{3}{121.184}{601.288}{\Character{3}}%
\TextAt{8}{125.916}{596.95}{,}%
\TextAt{8}{133.065}{596.95}{with}%
\TextAt{13}{159.725}{596.95}{x}%
\TextAt{17}{166.378}{595.156}{1}%
\NativeMathAt{174.431}{596.95}{11.955}{\neq}%
\TextAt{13}{186.856}{596.95}{x}%
\TextAt{17}{193.508}{595.156}{2}%
\TextAt{8}{198.241}{596.95}{,}%
\TextAt{8}{205.39}{596.95}{one}%
\TextAt{8}{226.849}{596.95}{deduces}%
\TextAt{13}{270.474}{596.95}{x}%
\TextAt{17}{277.126}{595.156}{1}%
\TextAt{13}{281.858}{596.95}{x}%
\TextAt{3}{288.51}{601.898}{\Character{0}}%
\TextAt{17}{295.097}{601.898}{1}%
\TextAt{17}{288.51}{594.019}{2}%
\TextAt{9}{303.15}{596.95}{\Character{61}}%
\TextAt{13}{315.575}{596.95}{v}%
\TextAt{17}{321.235}{595.156}{1}%
\TextAt{13}{325.967}{596.95}{v}%
\TextAt{3}{332.055}{601.898}{\Character{0}}%
\TextAt{17}{338.642}{601.898}{1}%
\TextAt{17}{331.626}{594.019}{2}%
\NativeMathAt{346.695}{596.95}{11.955}{\in}%
\TextAt{16}{357.969}{596.95}{\Character{85}}%
\TextAt{8}{366.655}{596.95}{,}%
\TextAt{8}{373.804}{596.95}{a}%
\TextAt{8}{383.56}{596.95}{contradiction.}%
\Vector{q 0 G 0 g 0.399 w 452.334 -177.05 m 452.334 -168.98 l S Q}%
\Vector{q 0 G 0 g 0.399 w 452.534 -169.18 m 459.707 -169.18 l S Q}%
\Vector{q 0 G 0 g 0.399 w 452.534 -176.851 m 459.707 -176.851 l S Q}%
\Vector{q 0 G 0 g 0.399 w 459.906 -177.05 m 459.906 -168.98 l S Q}%
\TextAt{8}{97.159}{569.003}{As}%
\TextAt{13}{115.633}{569.003}{K}%
\TextAt{3}{126.451}{573.342}{\Character{3}}%
\TextAt{8}{136.268}{569.003}{is}%
\TextAt{8}{149.219}{569.003}{closed}%
\TextAt{8}{184.933}{569.003}{in}%
\TextAt{13}{199.773}{569.003}{I}%
\TextAt{21}{204.935}{567.21}{K}%
\TextAt{8}{213.138}{569.003}{,}%
\TextAt{8}{221.771}{569.003}{the}%
\TextAt{8}{243.115}{569.003}{fact}%
\TextAt{8}{267.388}{569.003}{that}%
\TextAt{13}{293.944}{569.003}{I}%
\TextAt{21}{299.106}{567.21}{K}%
\TextAt{8}{312.394}{569.003}{is}%
\TextAt{8}{325.345}{569.003}{a}%
\TextAt{8}{336.287}{569.003}{locally}%
\TextAt{8}{374.548}{569.003}{compact}%
\TextAt{8}{422.563}{569.003}{Hausdorff}%
\TextAt{8}{477.142}{569.003}{topological}%
\TextAt{8}{79.226}{551.957}{group}%
\TextAt{8}{112.994}{551.957}{carries}%
\TextAt{8}{150.718}{551.957}{over}%
\TextAt{8}{176.356}{551.957}{to}%
\TextAt{8}{191.258}{551.957}{the}%
\TextAt{8}{212.006}{551.957}{idèle}%
\TextAt{8}{239.904}{551.957}{class}%
\TextAt{8}{267.933}{551.957}{group}%
\TextAt{13}{301.7}{551.957}{C}%
\TextAt{21}{310.075}{550.164}{K}%
\TextAt{9}{322.607}{551.957}{\Character{61}}%
\TextAt{13}{336.041}{551.957}{I}%
\TextAt{21}{341.204}{550.164}{K}%
\TextAt{9}{349.407}{551.957}{\Character{47}}%
\TextAt{13}{355.26}{551.957}{K}%
\TextAt{3}{366.077}{556.296}{\Character{3}}%
\TextAt{8}{370.81}{551.957}{.}%
\TextAt{8}{381.034}{551.957}{For}%
\TextAt{8}{402.595}{551.957}{any}%
\TextAt{8}{425.304}{551.957}{idèle}%
\TextAt{13}{453.201}{551.957}{\Character{11}}%
\TextAt{9}{465.052}{551.957}{\Character{61}}%
\TextAt{9}{478.478}{551.957}{\Character{40}}%
\TextAt{13}{483.038}{551.957}{\Character{11}}%
\TextAt{27}{490.484}{550.164}{p}%
\TextAt{9}{494.971}{551.957}{\Character{41}}%
\NativeMathAt{503.853}{551.957}{11.955}{\in}%
\TextAt{13}{516.152}{551.957}{I}%
\TextAt{21}{521.315}{550.164}{K}%
\TextAt{8}{529.518}{551.957}{,}%
\TextAt{8}{79.226}{534.911}{its}%
\TextAt{8}{95.471}{534.911}{class}%
\TextAt{8}{122.834}{534.911}{in}%
\TextAt{13}{136.412}{534.911}{C}%
\TextAt{21}{144.788}{533.118}{K}%
\TextAt{8}{156.813}{534.911}{will}%
\TextAt{8}{178.844}{534.911}{be}%
\TextAt{8}{194.695}{534.911}{denoted}%
\TextAt{8}{238.843}{534.911}{by}%
\TextAt{9}{255.028}{534.911}{\Character{91}}%
\TextAt{13}{258.279}{534.911}{\Character{11}}%
\TextAt{9}{265.801}{534.911}{\Character{93}}%
\TextAt{8}{269.053}{534.911}{.}%
\TextAt{8}{277.476}{534.911}{We}%
\TextAt{8}{297.546}{534.911}{define}%
\TextAt{8}{331.282}{534.911}{the}%
\TextAt{12}{351.364}{534.911}{absolute}%
\TextAt{12}{404.395}{534.911}{norm}%
\TextAt{8}{439.135}{534.911}{of}%
\TextAt{13}{452.391}{534.911}{\Character{11}}%
\TextAt{8}{463.736}{534.911}{to}%
\TextAt{8}{477.972}{534.911}{be}%
\TextAt{8}{493.822}{534.911}{the}%
\TextAt{8}{513.904}{534.911}{real}%
\TextAt{8}{79.226}{517.865}{number}%
\TextAt{14}{218.199}{500.819}{N}%
\TextAt{9}{228.147}{500.819}{\Character{40}}%
\TextAt{13}{232.7}{500.819}{\Character{11}}%
\TextAt{9}{240.221}{500.819}{\Character{41}}%
\TextAt{9}{248.076}{500.819}{\Character{61}}%
\TextAt{20}{260.52}{512.177}{\Character{89}}%
\TextAt{27}{266.164}{487.071}{p}%
\TextAt{14}{277.789}{500.819}{N}%
\TextAt{9}{287.737}{500.819}{\Character{40}}%
\TextAt{14}{292.289}{500.819}{p}%
\TextAt{9}{298.272}{500.819}{\Character{41}}%
\TextAt{21}{302.824}{505.756}{v}%
\TextAt{28}{306.948}{504.759}{p}%
\TextAt{17}{311.031}{505.756}{\Character{40}}%
\TextAt{21}{314.324}{505.756}{\Character{11}}%
\TextAt{28}{319.763}{504.759}{p}%
\TextAt{17}{323.845}{505.756}{\Character{41}}%
\TextAt{9}{330.958}{500.819}{\Character{61}}%
\TextAt{20}{343.383}{512.177}{\Character{89}}%
\TextAt{27}{349.027}{487.071}{p}%
\TextAt{16}{360.652}{500.819}{\Character{106}}%
\TextAt{13}{363.973}{500.819}{\Character{11}}%
\TextAt{27}{371.418}{499.026}{p}%
\TextAt{16}{375.905}{500.819}{\Character{106}}%
\TextAt{3}{379.226}{505.756}{\Character{0}}%
\TextAt{17}{385.813}{505.756}{1}%
\TextAt{27}{379.226}{497.864}{p}%
\TextAt{13}{390.545}{500.819}{\Character{58}}%
\TextAt{8}{79.226}{469.004}{If}%
\TextAt{13}{90.922}{469.004}{x}%
\NativeMathAt{100.902}{469.004}{11.955}{\in}%
\TextAt{13}{112.201}{469.004}{K}%
\TextAt{3}{123.018}{473.342}{\Character{3}}%
\TextAt{8}{131.652}{469.004}{is}%
\TextAt{8}{143.42}{469.004}{a}%
\TextAt{8}{153.18}{469.004}{principal}%
\TextAt{8}{201.961}{469.004}{idèle,}%
\TextAt{8}{232.524}{469.004}{then}%
\TextAt{8}{259.188}{469.004}{we}%
\TextAt{8}{276.419}{469.004}{find}%
\TextAt{8}{299.832}{469.004}{by}%
\TextAt{8}{316.095}{469.004}{chap.}%
\TextAt{8}{346.991}{469.004}{III,}%
\TextAt{8}{367.451}{469.004}{(1.3),}%
\TextAt{8}{398.683}{469.004}{that}%
\TextAt{14}{424.056}{469.004}{N}%
\TextAt{9}{434.004}{469.004}{\Character{40}}%
\TextAt{13}{438.556}{469.004}{x}%
\TextAt{9}{445.208}{469.004}{\Character{41}}%
\TextAt{9}{453.075}{469.004}{\Character{61}}%
\TextAt{20}{465.521}{477.97}{\Character{81}}%
\TextAt{27}{476.812}{465.517}{p}%
\TextAt{16}{483.291}{469.004}{\Character{106}}%
\TextAt{13}{486.612}{469.004}{x}%
\TextAt{16}{493.264}{469.004}{\Character{106}}%
\TextAt{3}{496.585}{473.342}{\Character{0}}%
\TextAt{17}{503.172}{473.342}{1}%
\TextAt{27}{496.585}{466.048}{p}%
\TextAt{9}{511.232}{469.004}{\Character{61}}%
\TextAt{9}{523.654}{469.004}{1}%
\TextAt{8}{529.518}{469.004}{.}%
\TextAt{8}{79.226}{451.958}{We}%
\TextAt{8}{99.371}{451.958}{thus}%
\TextAt{8}{125.122}{451.958}{have}%
\TextAt{8}{152.117}{451.958}{a}%
\TextAt{8}{161.872}{451.958}{continuous}%
\TextAt{8}{220.476}{451.958}{homomorphism}%
\TextAt{14}{265.352}{426.648}{N}%
\TextAt{9}{278.62}{426.648}{\Character{58}}%
\TextAt{13}{285.193}{426.648}{C}%
\TextAt{21}{293.568}{424.855}{K}%
\NativeArrowAt{305.091}{426.648}{11.955}{19.26}{\longrightarrow}%
\TextAt{18}{327.674}{426.648}{R}%
\TextAt{3}{336.308}{431.584}{\Character{3}}%
\TextAt{17}{336.308}{423.693}{\Character{43}}%
\TextAt{13}{343.392}{426.648}{\Character{58}}%
\TextAt{8}{79.226}{401.339}{It}%
\TextAt{8}{91.045}{401.339}{is}%
\TextAt{8}{101.955}{401.339}{related}%
\TextAt{8}{140.123}{401.339}{to}%
\TextAt{8}{153.58}{401.339}{the}%
\TextAt{8}{172.883}{401.339}{absolute}%
\TextAt{8}{218.272}{401.339}{norm}%
\TextAt{8}{247.988}{401.339}{on}%
\TextAt{8}{263.393}{401.339}{the}%
\TextAt{8}{282.696}{401.339}{replete}%
\TextAt{8}{320.206}{401.339}{Picard}%
\TextAt{8}{356.581}{401.339}{group}%
\TextAt{13}{388.903}{401.339}{P}%
\TextAt{13}{398.073}{401.339}{ic}%
\TextAt{9}{407.106}{401.339}{\Character{40}}%
\Vector{q 0 G 0 g 0.478 w 339.658 -362.818 m 349.508 -362.818 l S Q}%
\TextAt{16}{411.658}{401.339}{\Character{79}}%
\TextAt{9}{421.508}{401.339}{\Character{41}}%
\TextAt{8}{429.104}{401.339}{via}%
\TextAt{8}{447.438}{401.339}{the}%
\TextAt{8}{466.741}{401.339}{commutative}%
\TextAt{8}{79.226}{384.293}{diagram}%
\TextAt{13}{270.35}{366.954}{C}%
\TextAt{21}{278.725}{365.161}{K}%
\TextAt{18}{336.215}{366.954}{R}%
\TextAt{3}{344.849}{371.292}{\Character{3}}%
\TextAt{17}{344.849}{363.998}{\Character{43}}%
\TextAt{13}{260.06}{322.329}{P}%
\TextAt{13}{269.218}{322.329}{ic}%
\TextAt{9}{278.262}{322.329}{\Character{40}}%
\NativeDiagram{%
\path[draw=black,line width=0.478bp] (210.815,-441.828) -- (220.665,-441.828);%
}%
\TextAt{16}{282.815}{322.329}{\Character{79}}%
\TextAt{9}{292.665}{322.329}{\Character{41}}%
\TextAt{18}{336.215}{322.329}{R}%
\TextAt{3}{344.849}{326.667}{\Character{3}}%
\TextAt{17}{344.849}{319.373}{\Character{43}}%
\NativeDiagram{%
\path[draw=black,line width=0.478bp] (220.28,-404.056) -- (258.386,-404.056);%
\path[draw=black,line width=0.478bp,line cap=round,line join=round] (256.139,-401.187) .. controls (256.593,-402.909) and (257.605,-403.722) .. (258.625,-404.056) .. controls (257.605,-404.391) and (256.593,-405.204) .. (256.139,-406.926);%
}%
\TextAt{27}{308.256}{372.759}{N}%
\NativeDiagram{%
\path[draw=black,line width=0.478bp] (206.639,-413.418) -- (206.639,-436.053);%
\path[draw=black,line width=0.478bp,line cap=round,line join=round] (209.508,-433.806) .. controls (207.787,-434.26) and (206.974,-435.272) .. (206.639,-436.292) .. controls (206.305,-435.272) and (205.492,-434.26) .. (203.77,-433.806);%
\path[draw=black,line width=2.797bp] (272.076,-415.41) -- (272.076,-438.856);%
\path[draw=black!0,line width=1.841bp] (272.076,-415.41) -- (272.076,-438.856);%
\path[draw=black,line width=0.478bp] (230.57,-448.682) -- (258.386,-448.682);%
\path[draw=black,line width=0.478bp,line cap=round,line join=round] (256.139,-445.813) .. controls (256.593,-447.534) and (257.605,-448.347) .. (258.625,-448.682) .. controls (257.605,-449.017) and (256.593,-449.83) .. (256.139,-451.551);%
}%
\TextAt{27}{313.401}{328.134}{N}%
\TextAt{8}{79.226}{297.714}{Here}%
\TextAt{8}{106.855}{297.714}{the}%
\TextAt{8}{127.011}{297.714}{arrow}%
\TextAt{13}{159.864}{297.714}{C}%
\TextAt{21}{168.239}{295.92}{K}%
\NativeArrowAt{179.762}{297.714}{11.955}{19.26}{\longrightarrow}%
\TextAt{13}{202.344}{297.714}{P}%
\TextAt{13}{211.514}{297.714}{ic}%
\TextAt{9}{220.547}{297.714}{\Character{40}}%
\Vector{q 0 G 0 g 0.478 w 153.099 -466.443 m 162.949 -466.443 l S Q}%
\TextAt{16}{225.099}{297.714}{\Character{79}}%
\TextAt{9}{234.949}{297.714}{\Character{41}}%
\TextAt{8}{243.399}{297.714}{is}%
\TextAt{8}{255.163}{297.714}{induced}%
\TextAt{8}{298.727}{297.714}{by}%
\TextAt{8}{314.986}{297.714}{the}%
\TextAt{8}{335.143}{297.714}{continuous}%
\TextAt{8}{393.747}{297.714}{surjective}%
\TextAt{8}{446.158}{297.714}{homomorphism}%
\TextAt{13}{209.393}{272.404}{I}%
\TextAt{21}{214.556}{270.611}{K}%
\NativeArrowAt{226.079}{272.404}{11.955}{19.26}{\longrightarrow}%
\TextAt{13}{248.661}{272.404}{J}%
\TextAt{9}{256.238}{272.404}{\Character{40}}%
\Vector{q 0 G 0 g 0.478 w 188.791 -491.753 m 198.64 -491.753 l S Q}%
\TextAt{16}{260.791}{272.404}{\Character{79}}%
\TextAt{9}{270.64}{272.404}{\Character{41}}%
\TextAt{13}{275.193}{272.404}{\Character{59}}%
\TextAt{9}{304.347}{272.404}{\Character{40}}%
\TextAt{13}{308.899}{272.404}{\Character{11}}%
\TextAt{27}{316.345}{270.611}{p}%
\TextAt{9}{320.832}{272.404}{\Character{41}}%
\NativeArrowAt{328.705}{272.404}{11.955}{19.26}{\longmapsto}%
\TextAt{20}{351.287}{283.762}{\Character{89}}%
\TextAt{27}{356.931}{258.656}{p}%
\TextAt{14}{368.556}{272.404}{p}%
\TextAt{21}{374.539}{277.34}{v}%
\TextAt{28}{378.662}{276.344}{p}%
\TextAt{17}{382.745}{277.34}{\Character{40}}%
\TextAt{21}{386.038}{277.34}{\Character{11}}%
\TextAt{28}{391.477}{276.344}{p}%
\TextAt{17}{395.56}{277.34}{\Character{41}}%
\TextAt{13}{399.351}{272.404}{\Character{59}}%
\TextAt{8}{79.226}{237.219}{with}%
\TextAt{8}{105.157}{237.219}{kernel}%
\TextAt{13}{138.894}{237.219}{I}%
\TextAt{17}{144.997}{241.557}{0}%
\TextAt{21}{144.057}{234.198}{K}%
\TextAt{9}{155.58}{237.219}{\Character{61}}%
\TextAt{16}{168.006}{237.218}{\Character{102}}%
\TextAt{9}{173.983}{237.219}{\Character{40}}%
\TextAt{13}{178.536}{237.218}{\Character{11}}%
\TextAt{27}{185.981}{235.425}{p}%
\TextAt{9}{190.468}{237.219}{\Character{41}}%
\NativeMathAt{198.341}{237.219}{11.955}{\in}%
\TextAt{13}{209.632}{237.219}{I}%
\TextAt{21}{214.795}{235.425}{K}%
\TextAt{16}{226.319}{237.219}{\Character{106}}%
\TextAt{16}{232.942}{237.219}{\Character{106}}%
\TextAt{13}{236.281}{237.219}{\Character{11}}%
\TextAt{27}{243.727}{235.425}{p}%
\TextAt{16}{248.214}{237.219}{\Character{106}}%
\TextAt{27}{251.535}{235.425}{p}%
\TextAt{9}{259.342}{237.219}{\Character{61}}%
\TextAt{9}{271.752}{237.219}{1}%
\TextAt{8}{281.518}{237.219}{for}%
\TextAt{8}{299.403}{237.219}{all}%
\TextAt{14}{315.662}{237.219}{p}%
\TextAt{16}{321.645}{237.218}{\Character{103}}%
\TextAt{8}{327.623}{237.219}{.}%
\TextAt{8}{335.828}{237.219}{As}%
\TextAt{8}{352.386}{237.219}{to}%
\TextAt{8}{365.967}{237.219}{the}%
\TextAt{8}{385.393}{237.219}{kernel}%
\TextAt{13}{419.131}{237.219}{C}%
\TextAt{17}{428.364}{241.557}{0}%
\TextAt{21}{427.506}{234.198}{K}%
\TextAt{8}{438.876}{237.219}{of}%
\TextAt{14}{451.477}{237.219}{N}%
\TextAt{9}{464.746}{237.218}{\Character{58}}%
\TextAt{13}{471.318}{237.219}{C}%
\TextAt{21}{479.693}{235.425}{K}%
\NativeArrowAt{491.217}{237.218}{11.955}{19.26}{\longrightarrow}%
\TextAt{18}{513.799}{237.218}{R}%
\TextAt{17}{522.433}{235.425}{\Character{43}}%
\TextAt{8}{529.518}{237.219}{,}%
\TextAt{8}{79.226}{220.173}{we}%
\TextAt{8}{96.529}{220.173}{obtain,}%
\TextAt{8}{136.302}{220.173}{in}%
\TextAt{8}{150.03}{220.173}{analogy}%
\TextAt{8}{193.37}{220.173}{with}%
\TextAt{8}{220.106}{220.173}{chap.}%
\TextAt{8}{251.073}{220.173}{III,}%
\TextAt{8}{271.603}{220.173}{(1.14),}%
\TextAt{8}{308.782}{220.173}{the}%
\TextAt{8}{329.013}{220.173}{following}%
\TextAt{8}{378.523}{220.173}{important}%
\TextAt{8}{433.891}{220.173}{theorem.}%
\TextAt{8}{484.192}{220.173}{It}%
\TextAt{8}{496.94}{220.173}{reflects}%
\TextAt{8}{79.226}{203.127}{the}%
\TextAt{8}{99.383}{203.127}{finiteness}%
\TextAt{8}{150.228}{203.127}{of}%
\TextAt{8}{163.558}{203.127}{the}%
\TextAt{8}{183.714}{203.127}{unit}%
\TextAt{8}{208.426}{203.127}{rank}%
\TextAt{8}{235.42}{203.127}{of}%
\TextAt{13}{248.75}{203.127}{K}%
\TextAt{8}{263.466}{203.127}{as}%
\TextAt{8}{277.836}{203.127}{well}%
\TextAt{8}{301.567}{203.127}{as}%
\TextAt{8}{315.937}{203.127}{the}%
\TextAt{8}{336.093}{203.127}{finiteness}%
\TextAt{8}{386.939}{203.127}{of}%
\TextAt{8}{400.269}{203.127}{the}%
\TextAt{8}{420.425}{203.127}{class}%
\TextAt{8}{447.862}{203.127}{number.}%
\TextAt{12}{79.226}{178.169}{Theorem}%
\TextAt{12}{136.24}{178.169}{(1.6).}%
\TextAt{8}{171.556}{178.169}{The}%
\TextAt{8}{195.61}{178.169}{group}%
\TextAt{13}{228.785}{178.169}{C}%
\TextAt{17}{238.019}{182.507}{0}%
\TextAt{21}{237.16}{175.148}{K}%
\TextAt{9}{248.684}{178.169}{\Character{61}}%
\TextAt{16}{261.109}{178.169}{\Character{102}}%
\TextAt{9}{267.087}{178.169}{\Character{91}}%
\TextAt{13}{270.339}{178.169}{\Character{11}}%
\TextAt{9}{277.86}{178.169}{\Character{93}}%
\NativeMathAt{284.433}{178.169}{11.955}{\in}%
\TextAt{13}{295.724}{178.169}{C}%
\TextAt{21}{304.099}{176.376}{K}%
\TextAt{16}{315.622}{178.169}{\Character{106}}%
\TextAt{14}{322.264}{178.169}{N}%
\TextAt{9}{332.212}{178.169}{\Character{40}\Character{91}}%
\TextAt{13}{340.016}{178.169}{\Character{11}}%
\TextAt{9}{347.538}{178.169}{\Character{93}\Character{41}}%
\TextAt{9}{358.632}{178.169}{\Character{61}}%
\TextAt{9}{371.054}{178.169}{1}%
\TextAt{16}{376.941}{178.169}{\Character{103}}%
\TextAt{8}{386.816}{178.169}{is}%
\TextAt{8}{398.58}{178.169}{compact.}%
\TextAt{19}{79.226}{145.3}{Proof.}%
\TextAt{8}{116}{145.3}{The}%
\TextAt{8}{140.054}{145.3}{claim}%
\TextAt{8}{171.269}{145.3}{concerning}%
\TextAt{8}{229.801}{145.3}{the}%
\TextAt{8}{249.958}{145.3}{commutative}%
\TextAt{8}{319.884}{145.3}{exact}%
\TextAt{8}{350.776}{145.3}{diagram}%
\TextAt{9}{146.072}{119.734}{1}%
\TextAt{13}{217.925}{119.734}{C}%
\TextAt{17}{227.159}{124.073}{0}%
\TextAt{21}{226.3}{116.714}{K}%
\TextAt{13}{310.792}{119.734}{C}%
\TextAt{21}{319.167}{117.941}{K}%
\TextAt{18}{391.003}{119.734}{R}%
\TextAt{3}{399.638}{124.073}{\Character{3}}%
\TextAt{17}{399.638}{116.779}{\Character{43}}%
\TextAt{9}{460.066}{119.734}{1}%
\TextAt{9}{146.072}{75.109}{1}%
\TextAt{13}{205.269}{75.109}{P}%
\TextAt{13}{214.439}{75.109}{ic}%
\TextAt{9}{223.472}{75.109}{\Character{40}}%
\NativeDiagram{%
\path[draw=black,line width=0.478bp] (156.024,-689.048) -- (165.874,-689.048);%
}%
\TextAt{16}{228.024}{75.109}{\Character{79}}%
\TextAt{9}{237.874}{75.109}{\Character{41}}%
\TextAt{17}{242.426}{79.448}{0}%
\TextAt{13}{300.502}{75.109}{P}%
\TextAt{13}{309.672}{75.109}{ic}%
\TextAt{9}{318.705}{75.109}{\Character{40}}%
\NativeDiagram{%
\path[draw=black,line width=0.478bp] (251.257,-689.048) -- (261.107,-689.048);%
}%
\TextAt{16}{323.257}{75.109}{\Character{79}}%
\TextAt{9}{333.107}{75.109}{\Character{41}}%
\TextAt{18}{391.003}{75.109}{R}%
\TextAt{3}{399.638}{79.448}{\Character{3}}%
\TextAt{17}{399.638}{72.154}{\Character{43}}%
\TextAt{9}{460.066}{75.109}{1}%
\NativeDiagram{%
\path[draw=black,line width=0.478bp] (85.277,-651.277) -- (140.096,-651.277);%
\path[draw=black,line width=0.478bp,line cap=round,line join=round] (137.848,-648.408) .. controls (138.303,-650.129) and (139.315,-650.942) .. (140.335,-651.277) .. controls (139.315,-651.612) and (138.303,-652.425) .. (137.848,-654.146);%
\path[draw=black,line width=0.478bp] (167.856,-651.277) -- (232.964,-651.277);%
\path[draw=black,line width=0.478bp,line cap=round,line join=round] (230.717,-648.408) .. controls (231.171,-650.129) and (232.183,-650.942) .. (233.203,-651.277) .. controls (232.183,-651.612) and (231.171,-652.425) .. (230.717,-654.146);%
\path[draw=black,line width=0.478bp] (154.215,-661.865) -- (154.215,-683.274);%
\path[draw=black,line width=0.478bp,line cap=round,line join=round] (157.084,-681.026) .. controls (155.363,-681.48) and (154.55,-682.493) .. (154.215,-683.513) .. controls (153.88,-682.493) and (153.067,-681.48) .. (151.346,-681.026);%
\path[draw=black,line width=0.478bp] (260.724,-651.277) -- (313.177,-651.277);%
\path[draw=black,line width=0.478bp,line cap=round,line join=round] (310.929,-648.408) .. controls (311.383,-650.129) and (312.396,-650.942) .. (313.416,-651.277) .. controls (312.396,-651.612) and (311.383,-652.425) .. (310.929,-654.146);%
\path[draw=black,line width=0.478bp] (247.083,-660.638) -- (247.083,-683.274);%
\path[draw=black,line width=0.478bp,line cap=round,line join=round] (249.952,-681.026) .. controls (248.231,-681.48) and (247.418,-682.493) .. (247.083,-683.513) .. controls (246.749,-682.493) and (245.936,-681.48) .. (244.214,-681.026);%
\path[draw=black,line width=0.478bp] (340.078,-651.277) -- (382.24,-651.277);%
\path[draw=black,line width=0.478bp,line cap=round,line join=round] (379.993,-648.408) .. controls (380.447,-650.129) and (381.46,-650.942) .. (382.48,-651.277) .. controls (381.46,-651.612) and (380.447,-652.425) .. (379.993,-654.146);%
\path[draw=black,line width=2.797bp] (326.866,-662.631) -- (326.866,-686.076);%
\path[draw=black!0,line width=1.841bp] (326.866,-662.631) -- (326.866,-686.076);%
\path[draw=black,line width=0.478bp] (85.277,-695.902) -- (127.44,-695.902);%
\path[draw=black,line width=0.478bp,line cap=round,line join=round] (125.192,-693.033) .. controls (125.647,-694.755) and (126.659,-695.568) .. (127.679,-695.903) .. controls (126.659,-696.237) and (125.647,-697.05) .. (125.192,-698.772);%
\path[draw=black,line width=0.478bp] (180.512,-695.902) -- (222.674,-695.902);%
\path[draw=black,line width=0.478bp,line cap=round,line join=round] (220.427,-693.033) .. controls (220.881,-694.755) and (221.894,-695.568) .. (222.913,-695.903) .. controls (221.894,-696.237) and (220.881,-697.05) .. (220.427,-698.772);%
\path[draw=black,line width=0.478bp] (271.014,-695.902) -- (313.177,-695.902);%
\path[draw=black,line width=0.478bp,line cap=round,line join=round] (310.929,-693.033) .. controls (311.383,-694.755) and (312.396,-695.568) .. (313.416,-695.903) .. controls (312.396,-696.237) and (311.383,-697.05) .. (310.929,-698.772);%
\path[draw=black,line width=0.478bp] (340.078,-695.902) -- (382.24,-695.902);%
\path[draw=black,line width=0.478bp,line cap=round,line join=round] (379.993,-693.033) .. controls (380.447,-694.755) and (381.46,-695.568) .. (382.48,-695.903) .. controls (381.46,-696.237) and (380.447,-697.05) .. (379.993,-698.772);%
}%
\TextAt{8}{297.211}{43.033}{330}%
\Destination{78.23}{809.815}{nav.page.330}%
\Destination{79.36}{483.955}{nav.print.362}%
\Destination{79.36}{192.155}{nav.number.6.1.6}%
\Hotlink{346.491}{466.755}{14.306}{11.047}{nav.chapter.3}%
\Hotlink{250.573}{217.955}{14.294}{11.047}{nav.chapter.3}%
\Hotlink{366.951}{466.755}{25.078}{11.047}{nav.number.3.1.3}%
\Hotlink{271.103}{217.955}{30.936}{11.047}{nav.number.3.1.14}%
}

% Page 343 (body page 331)
\Page{%
\TextAt{2}{79.226}{796.943}{VI.}%
\TextAt{2}{98.012}{796.943}{Global}%
\TextAt{2}{133.248}{796.943}{Class}%
\TextAt{2}{161.841}{796.943}{Field}%
\TextAt{2}{189.572}{796.943}{Theory}%
\Vector{q 0 G 0 g 0.299 w 7.226 18.729 m 460.769 18.729 l S Q}%
\TextAt{8}{79.226}{761.124}{will}%
\TextAt{8}{101.255}{761.124}{be}%
\TextAt{8}{117.104}{761.124}{reduced}%
\TextAt{8}{160.593}{761.124}{to}%
\TextAt{8}{174.827}{761.124}{the}%
\TextAt{8}{194.908}{761.124}{compactness}%
\TextAt{8}{262.594}{761.124}{of}%
\TextAt{8}{275.848}{761.124}{the}%
\TextAt{8}{295.929}{761.124}{group}%
\TextAt{13}{329.029}{761.124}{P}%
\TextAt{13}{338.187}{761.124}{ic}%
\TextAt{9}{347.231}{761.124}{\Character{40}}%
\Vector{q 0 G 0 g 0.478 w 279.783 -3.033 m 289.633 -3.033 l S Q}%
\TextAt{16}{351.783}{761.124}{\Character{79}}%
\TextAt{9}{361.633}{761.124}{\Character{41}}%
\TextAt{17}{366.185}{765.462}{0}%
\TextAt{8}{370.918}{761.124}{,}%
\TextAt{8}{378.006}{761.124}{which}%
\TextAt{8}{411.417}{761.124}{was}%
\TextAt{8}{433.841}{761.124}{proved}%
\TextAt{8}{471.818}{761.124}{in}%
\TextAt{8}{485.395}{761.124}{chap.}%
\TextAt{8}{516.212}{761.124}{III,}%
\TextAt{8}{79.226}{744.078}{(1.14).}%
\TextAt{8}{117.554}{744.078}{The}%
\TextAt{8}{141.439}{744.078}{kernel}%
\TextAt{8}{175.737}{744.078}{of}%
\TextAt{8}{188.898}{744.078}{the}%
\TextAt{8}{208.886}{744.078}{vertical}%
\TextAt{8}{250.345}{744.078}{arrow}%
\TextAt{8}{283.029}{744.078}{in}%
\TextAt{8}{296.513}{744.078}{the}%
\TextAt{8}{316.501}{744.078}{middle}%
\TextAt{8}{354.697}{744.078}{is}%
\TextAt{8}{366.292}{744.078}{the}%
\TextAt{8}{386.28}{744.078}{group}%
\TextAt{13}{419.286}{744.078}{I}%
\TextAt{17}{425.389}{748.416}{0}%
\TextAt{21}{424.449}{741.057}{K}%
\TextAt{13}{432.652}{744.078}{K}%
\TextAt{3}{443.47}{748.416}{\Character{3}}%
\TextAt{9}{448.202}{744.078}{\Character{47}}%
\TextAt{13}{454.055}{744.078}{K}%
\TextAt{3}{464.873}{748.416}{\Character{3}}%
\TextAt{9}{472.926}{744.078}{\Character{61}}%
\TextAt{13}{485.351}{744.078}{I}%
\TextAt{17}{491.454}{748.416}{0}%
\TextAt{21}{490.514}{741.057}{K}%
\TextAt{9}{498.717}{744.078}{\Character{47}\Character{40}}%
\TextAt{13}{509.122}{744.078}{I}%
\TextAt{17}{515.225}{748.416}{0}%
\TextAt{21}{514.285}{741.057}{K}%
\TextAt{16}{524.799}{744.078}{\Character{92}}%
\TextAt{13}{79.226}{727.032}{K}%
\TextAt{3}{90.044}{731.37}{\Character{3}}%
\TextAt{9}{94.776}{727.032}{\Character{41}}%
\TextAt{8}{99.328}{727.032}{,}%
\TextAt{8}{107.221}{727.032}{where}%
\TextAt{8}{141.624}{727.032}{we}%
\TextAt{8}{159.446}{727.032}{have}%
\TextAt{13}{187.035}{727.032}{I}%
\TextAt{17}{193.137}{731.37}{0}%
\TextAt{21}{192.198}{724.011}{K}%
\TextAt{9}{204.734}{727.032}{\Character{61}}%
\TextAt{20}{218.172}{735.998}{\Character{81}}%
\TextAt{27}{229.463}{723.545}{p}%
\TextAt{13}{235.942}{727.032}{I}%
\TextAt{17}{242.045}{731.37}{0}%
\TextAt{27}{241.105}{724.076}{p}%
\TextAt{8}{246.777}{727.032}{,}%
\TextAt{13}{254.67}{727.032}{I}%
\TextAt{17}{260.772}{731.37}{0}%
\TextAt{27}{259.832}{724.076}{p}%
\TextAt{9}{269.838}{727.032}{\Character{61}}%
\TextAt{16}{283.276}{727.032}{\Character{102}}%
\TextAt{13}{289.254}{727.032}{\Character{11}}%
\TextAt{27}{296.7}{725.238}{p}%
\NativeMathAt{305.52}{727.032}{11.955}{\in}%
\TextAt{13}{317.824}{727.032}{K}%
\TextAt{27}{327.783}{725.238}{p}%
\TextAt{16}{336.603}{727.032}{\Character{106}}%
\TextAt{16}{344.242}{727.032}{\Character{106}}%
\TextAt{13}{347.579}{727.032}{\Character{11}}%
\TextAt{27}{355.024}{725.238}{p}%
\TextAt{16}{359.511}{727.032}{\Character{106}}%
\TextAt{27}{362.832}{725.238}{p}%
\TextAt{9}{371.652}{727.032}{\Character{61}}%
\TextAt{9}{385.078}{727.032}{1}%
\TextAt{16}{390.943}{727.032}{\Character{103}}%
\TextAt{8}{396.921}{727.032}{,}%
\TextAt{8}{404.813}{727.032}{and}%
\TextAt{13}{428.17}{727.032}{I}%
\TextAt{17}{434.273}{731.37}{0}%
\TextAt{21}{433.333}{724.011}{K}%
\TextAt{16}{444.597}{727.032}{\Character{92}}%
\TextAt{13}{455.629}{727.032}{K}%
\TextAt{3}{466.447}{731.37}{\Character{3}}%
\TextAt{9}{475.513}{727.032}{\Character{61}}%
\TextAt{13}{488.951}{727.032}{\Character{22}}%
\TextAt{9}{495.994}{727.032}{\Character{40}}%
\TextAt{13}{500.546}{727.032}{K}%
\TextAt{9}{511.364}{727.032}{\Character{41}}%
\TextAt{8}{520.408}{727.032}{by}%
\TextAt{8}{79.226}{709.986}{chap.}%
\TextAt{8}{110.118}{709.986}{III,}%
\TextAt{8}{130.574}{709.986}{(1.9).}%
\TextAt{8}{163.1}{709.986}{This}%
\TextAt{8}{189.819}{709.986}{kernel}%
\TextAt{8}{224.286}{709.986}{is}%
\TextAt{8}{236.05}{709.986}{clearly}%
\TextAt{8}{273.446}{709.986}{compact.}%
\TextAt{8}{324.825}{709.986}{We}%
\TextAt{8}{344.97}{709.986}{obtain}%
\TextAt{8}{381.397}{709.986}{an}%
\TextAt{8}{397.656}{709.986}{exact}%
\TextAt{8}{428.548}{709.986}{sequence}%
\TextAt{9}{193.946}{681.947}{1}%
\NativeArrowAt{203.12}{681.947}{11.955}{19.26}{\longrightarrow}%
\TextAt{13}{225.702}{681.947}{I}%
\TextAt{17}{231.805}{686.884}{0}%
\TextAt{21}{230.865}{678.992}{K}%
\TextAt{13}{239.068}{681.947}{K}%
\TextAt{3}{249.885}{686.884}{\Character{3}}%
\TextAt{9}{254.618}{681.947}{\Character{47}}%
\TextAt{13}{260.471}{681.947}{K}%
\TextAt{3}{271.288}{686.884}{\Character{3}}%
\NativeArrowAt{279.342}{681.947}{11.955}{19.26}{\longrightarrow}%
\TextAt{13}{301.924}{681.947}{C}%
\TextAt{17}{311.157}{686.884}{0}%
\TextAt{21}{310.299}{678.992}{K}%
\NativeArrowAt{321.822}{681.947}{11.955}{19.26}{\longrightarrow}%
\TextAt{13}{344.404}{681.947}{P}%
\TextAt{13}{353.574}{681.947}{ic}%
\TextAt{9}{362.607}{681.947}{\Character{40}}%
\Vector{q 0 G 0 g 0.478 w 295.159 -82.21 m 305.009 -82.21 l S Q}%
\TextAt{16}{367.159}{681.947}{\Character{79}}%
\TextAt{9}{377.009}{681.947}{\Character{41}}%
\TextAt{17}{381.561}{686.884}{0}%
\NativeArrowAt{389.614}{681.947}{11.955}{19.26}{\longrightarrow}%
\TextAt{9}{412.196}{681.947}{1}%
\TextAt{8}{79.226}{653.909}{of}%
\TextAt{8}{93.128}{653.909}{continuous}%
\TextAt{8}{152.304}{653.909}{homomorphisms.}%
\TextAt{8}{245.855}{653.909}{Since}%
\TextAt{13}{276.985}{653.909}{P}%
\TextAt{13}{286.143}{653.909}{ic}%
\TextAt{9}{295.187}{653.909}{\Character{40}}%
\Vector{q 0 G 0 g 0.478 w 227.74 -110.248 m 237.589 -110.248 l S Q}%
\TextAt{16}{299.74}{653.909}{\Character{79}}%
\TextAt{9}{309.589}{653.909}{\Character{41}}%
\TextAt{17}{314.142}{658.247}{0}%
\TextAt{8}{323.343}{653.909}{is}%
\TextAt{8}{335.679}{653.909}{compact,}%
\TextAt{8}{386.474}{653.909}{and}%
\TextAt{8}{409.808}{653.909}{the}%
\TextAt{8}{430.537}{653.909}{same}%
\TextAt{8}{460.435}{653.909}{is}%
\TextAt{8}{472.77}{653.909}{true}%
\TextAt{8}{498.054}{653.909}{for}%
\TextAt{8}{516.51}{653.909}{the}%
\TextAt{8}{79.226}{636.863}{fibres}%
\TextAt{8}{110.259}{636.863}{of}%
\TextAt{8}{123.348}{636.863}{the}%
\TextAt{8}{143.263}{636.863}{mapping}%
\TextAt{13}{191.153}{636.863}{C}%
\TextAt{17}{200.387}{641.201}{0}%
\TextAt{21}{199.528}{633.842}{K}%
\NativeArrowAt{211.052}{636.863}{11.955}{19.26}{\longrightarrow}%
\TextAt{13}{233.634}{636.863}{P}%
\TextAt{13}{242.792}{636.863}{ic}%
\TextAt{9}{251.836}{636.863}{\Character{40}}%
\Vector{q 0 G 0 g 0.478 w 184.388 -127.294 m 194.238 -127.294 l S Q}%
\TextAt{16}{256.388}{636.863}{\Character{79}}%
\TextAt{9}{266.238}{636.863}{\Character{41}}%
\TextAt{17}{270.791}{641.201}{0}%
\TextAt{8}{279.179}{636.863}{(they}%
\TextAt{8}{309.83}{636.863}{are}%
\TextAt{8}{329.099}{636.863}{cosets,}%
\TextAt{8}{366.098}{636.863}{all}%
\TextAt{8}{382.116}{636.863}{homeomorphic}%
\TextAt{8}{460.575}{636.863}{to}%
\TextAt{13}{474.644}{636.863}{I}%
\TextAt{17}{480.747}{641.201}{0}%
\TextAt{21}{479.807}{633.842}{K}%
\TextAt{13}{488.01}{636.863}{K}%
\TextAt{3}{498.827}{641.201}{\Character{3}}%
\TextAt{9}{503.56}{636.863}{\Character{47}}%
\TextAt{13}{509.413}{636.863}{K}%
\TextAt{3}{520.23}{641.201}{\Character{3}}%
\TextAt{8}{524.963}{636.863}{),}%
\TextAt{8}{79.226}{619.817}{hence}%
\TextAt{8}{111.732}{619.817}{so}%
\TextAt{8}{126.102}{619.817}{is}%
\TextAt{13}{137.866}{619.817}{C}%
\TextAt{17}{147.1}{624.155}{0}%
\TextAt{21}{146.241}{616.796}{K}%
\TextAt{8}{154.444}{619.817}{.}%
\Vector{q 0 G 0 g 0.399 w 452.334 -154.183 m 452.334 -146.113 l S Q}%
\Vector{q 0 G 0 g 0.399 w 452.534 -146.313 m 459.707 -146.313 l S Q}%
\Vector{q 0 G 0 g 0.399 w 452.534 -153.984 m 459.707 -153.984 l S Q}%
\Vector{q 0 G 0 g 0.399 w 459.906 -154.183 m 459.906 -146.113 l S Q}%
\TextAt{8}{97.159}{591.091}{The}%
\TextAt{8}{121.303}{591.091}{idèle}%
\TextAt{8}{148.7}{591.091}{class}%
\TextAt{8}{176.228}{591.091}{group}%
\TextAt{13}{209.494}{591.091}{C}%
\TextAt{21}{217.869}{589.298}{K}%
\TextAt{8}{230.06}{591.091}{plays}%
\TextAt{8}{260.134}{591.091}{a}%
\TextAt{8}{269.981}{591.091}{similar}%
\TextAt{8}{308.507}{591.091}{role}%
\TextAt{8}{331.361}{591.091}{for}%
\TextAt{8}{349.336}{591.091}{the}%
\TextAt{8}{369.584}{591.091}{algebraic}%
\TextAt{8}{419.109}{591.091}{number}%
\TextAt{8}{461.796}{591.091}{field}%
\TextAt{13}{487.244}{591.091}{K}%
\TextAt{8}{502.05}{591.091}{as}%
\TextAt{8}{516.51}{591.091}{the}%
\TextAt{8}{79.226}{574.045}{multiplicative}%
\TextAt{8}{154.171}{574.045}{group}%
\TextAt{13}{188.468}{574.045}{K}%
\TextAt{3}{199.285}{578.383}{\Character{3}}%
\TextAt{27}{198.427}{571.089}{p}%
\TextAt{8}{209.037}{574.045}{does}%
\TextAt{8}{236.555}{574.045}{for}%
\TextAt{8}{255.561}{574.045}{a}%
\TextAt{13}{266.438}{574.045}{p}%
\TextAt{8}{272.313}{574.045}{-adic}%
\TextAt{8}{302.044}{574.045}{number}%
\TextAt{8}{345.761}{574.045}{field}%
\TextAt{13}{372.24}{574.045}{K}%
\TextAt{27}{382.199}{572.252}{p}%
\TextAt{8}{386.686}{574.045}{.}%
\TextAt{8}{398.498}{574.045}{It}%
\TextAt{8}{412.292}{574.045}{comes}%
\TextAt{8}{447.94}{574.045}{equipped}%
\TextAt{8}{499.13}{574.045}{with}%
\TextAt{8}{526.911}{574.045}{a}%
\TextAt{8}{79.226}{556.999}{collection}%
\TextAt{8}{131.394}{556.999}{of}%
\TextAt{8}{144.863}{556.999}{canonical}%
\TextAt{8}{196.386}{556.999}{subgroups}%
\TextAt{8}{251.937}{556.999}{which}%
\TextAt{8}{285.563}{556.999}{are}%
\TextAt{8}{305.213}{556.999}{to}%
\TextAt{8}{319.663}{556.999}{be}%
\TextAt{8}{335.726}{556.999}{viewed}%
\TextAt{8}{374.229}{556.999}{as}%
\TextAt{8}{388.739}{556.999}{analogues}%
\TextAt{8}{442.282}{556.999}{of}%
\TextAt{8}{455.751}{556.999}{the}%
\TextAt{8}{476.046}{556.999}{higher}%
\TextAt{8}{511.956}{556.999}{unit}%
\TextAt{8}{79.226}{539.953}{groups}%
\TextAt{13}{117.849}{539.953}{U}%
\TextAt{17}{127.048}{546.063}{\Character{40}}%
\TextAt{21}{130.341}{546.063}{n}%
\TextAt{17}{135.48}{546.063}{\Character{41}}%
\TextAt{27}{125.799}{537.982}{p}%
\TextAt{9}{144.011}{539.953}{\Character{61}}%
\TextAt{9}{157.843}{539.953}{1}%
\TextAt{9}{166.917}{539.953}{\Character{43}}%
\TextAt{14}{179.262}{539.953}{p}%
\TextAt{21}{185.245}{544.291}{n}%
\TextAt{8}{195.611}{539.953}{of}%
\TextAt{8}{209.774}{539.953}{a}%
\TextAt{13}{220.362}{539.953}{p}%
\TextAt{8}{226.237}{539.953}{-adic}%
\TextAt{8}{255.678}{539.953}{number}%
\TextAt{8}{299.107}{539.953}{field}%
\TextAt{13}{325.297}{539.953}{K}%
\TextAt{27}{335.256}{538.16}{p}%
\TextAt{8}{339.743}{539.953}{.}%
\TextAt{8}{350.689}{539.953}{Instead}%
\TextAt{8}{392.875}{539.953}{of}%
\TextAt{14}{407.038}{539.953}{p}%
\TextAt{21}{413.021}{544.291}{n}%
\TextAt{8}{418.657}{539.953}{,}%
\TextAt{8}{426.847}{539.953}{we}%
\TextAt{8}{444.907}{539.953}{take}%
\TextAt{8}{471.109}{539.953}{any}%
\TextAt{8}{494.059}{539.953}{integral}%
\TextAt{8}{79.226}{522.907}{ideal}%
\TextAt{14}{107.189}{522.907}{m}%
\TextAt{9}{119.677}{522.907}{\Character{61}}%
\TextAt{20}{132.102}{531.873}{\Character{81}}%
\TextAt{27}{143.393}{519.42}{p}%
\TextAt{25}{147.382}{519.42}{\Character{45}}%
\NativeMathAt{149.596}{519.42}{7.97}{\infty}%
\TextAt{14}{160.555}{522.907}{p}%
\TextAt{21}{166.538}{527.245}{n}%
\TextAt{28}{171.676}{526.249}{p}%
\TextAt{8}{176.257}{522.907}{.}%
\TextAt{8}{184.705}{522.907}{We}%
\TextAt{8}{204.85}{522.907}{may}%
\TextAt{8}{230.219}{522.907}{also}%
\TextAt{8}{253.699}{522.907}{write}%
\TextAt{8}{283.611}{522.907}{it}%
\TextAt{8}{295.315}{522.907}{as}%
\TextAt{8}{309.685}{522.907}{a}%
\TextAt{8}{319.44}{522.907}{replete}%
\TextAt{8}{357.804}{522.907}{ideal}%
\TextAt{14}{277.056}{492.874}{m}%
\TextAt{9}{289.544}{492.874}{\Character{61}}%
\TextAt{20}{301.969}{504.232}{\Character{89}}%
\TextAt{27}{307.613}{479.126}{p}%
\TextAt{14}{319.238}{492.874}{p}%
\TextAt{21}{325.221}{497.811}{n}%
\TextAt{28}{330.359}{496.814}{p}%
\TextAt{8}{79.226}{455.468}{with}%
\TextAt{13}{105.517}{455.468}{n}%
\TextAt{27}{112.505}{453.675}{p}%
\TextAt{9}{120.312}{455.468}{\Character{61}}%
\TextAt{9}{132.722}{455.468}{0}%
\TextAt{8}{142.119}{455.468}{for}%
\TextAt{14}{159.635}{455.468}{p}%
\TextAt{16}{168.939}{455.468}{\Character{106}}%
\NativeMathAt{175.562}{455.468}{11.955}{\infty}%
\TextAt{8}{187.536}{455.468}{,}%
\TextAt{8}{194.39}{455.468}{and}%
\TextAt{8}{216.784}{455.468}{we}%
\TextAt{8}{233.642}{455.468}{treat}%
\TextAt{8}{261.894}{455.468}{it}%
\TextAt{8}{273.229}{455.468}{in}%
\TextAt{8}{286.513}{455.468}{what}%
\TextAt{8}{315.41}{455.468}{follows}%
\TextAt{8}{353.477}{455.468}{as}%
\TextAt{8}{367.478}{455.468}{a}%
\TextAt{12}{376.865}{455.468}{module}%
\TextAt{8}{423.527}{455.468}{of}%
\TextAt{13}{436.488}{455.468}{K}%
\TextAt{8}{447.306}{455.468}{.}%
\TextAt{8}{455.631}{455.468}{For}%
\TextAt{8}{476.232}{455.468}{every}%
\TextAt{8}{506.755}{455.468}{place}%
\TextAt{14}{79.226}{438.422}{p}%
\TextAt{8}{89.107}{438.422}{of}%
\TextAt{13}{102.437}{438.422}{K}%
\TextAt{8}{117.152}{438.422}{we}%
\TextAt{8}{134.379}{438.422}{put}%
\TextAt{13}{155.839}{438.422}{U}%
\TextAt{17}{165.038}{444.533}{\Character{40}0\Character{41}}%
\TextAt{27}{163.789}{436.452}{p}%
\TextAt{9}{179.677}{438.422}{\Character{61}}%
\TextAt{13}{192.103}{438.422}{U}%
\TextAt{27}{200.053}{436.629}{p}%
\TextAt{8}{204.54}{438.422}{,}%
\TextAt{8}{211.689}{438.422}{and}%
\TextAt{13}{207.842}{389.285}{U}%
\TextAt{17}{217.041}{395.53}{\Character{40}}%
\TextAt{21}{220.334}{395.53}{n}%
\TextAt{28}{225.473}{394.534}{p}%
\TextAt{17}{229.555}{395.53}{\Character{41}}%
\TextAt{27}{215.792}{387.314}{p}%
\TextAt{9}{236.668}{389.285}{\Character{58}\Character{61}}%
\TextAt{20}{252.345}{420.966}{\Character{56}}%
\TextAt{20}{252.345}{410.206}{\Character{62}}%
\TextAt{20}{252.345}{406.62}{\Character{62}}%
\TextAt{20}{252.345}{403.033}{\Character{60}}%
\TextAt{20}{252.345}{381.514}{\Character{62}}%
\TextAt{20}{252.345}{377.927}{\Character{62}}%
\TextAt{20}{252.345}{374.34}{\Character{58}}%
\TextAt{9}{262.972}{408.638}{1}%
\TextAt{9}{271.472}{408.638}{\Character{43}}%
\TextAt{14}{283.243}{408.638}{p}%
\TextAt{21}{289.226}{412.976}{n}%
\TextAt{28}{294.364}{411.98}{p}%
\TextAt{13}{298.945}{408.638}{\Character{59}}%
\TextAt{8}{325.388}{408.638}{if}%
\TextAt{14}{336.111}{408.638}{p}%
\TextAt{18}{345.415}{408.638}{\Character{45}}%
\NativeMathAt{352.057}{408.638}{11.955}{\infty}%
\TextAt{13}{364.012}{408.638}{\Character{59}}%
\TextAt{18}{262.972}{388.182}{R}%
\TextAt{3}{271.606}{392.521}{\Character{3}}%
\TextAt{17}{271.606}{385.227}{\Character{43}}%
\NativeMathAt{282.011}{388.182}{11.955}{\subseteq}%
\TextAt{13}{294.631}{388.182}{K}%
\TextAt{3}{305.448}{392.521}{\Character{3}}%
\TextAt{27}{304.59}{385.227}{p}%
\TextAt{13}{310.181}{388.182}{\Character{59}}%
\TextAt{8}{325.388}{388.182}{if}%
\TextAt{14}{336.111}{388.182}{p}%
\TextAt{8}{345.992}{388.182}{is}%
\TextAt{8}{357.756}{388.182}{real}%
\TextAt{13}{376.621}{388.182}{\Character{59}}%
\TextAt{18}{262.972}{367.727}{C}%
\TextAt{3}{271.606}{372.066}{\Character{3}}%
\TextAt{9}{279.659}{367.727}{\Character{61}}%
\TextAt{13}{292.084}{367.727}{K}%
\TextAt{3}{302.902}{372.066}{\Character{3}}%
\TextAt{27}{302.044}{364.772}{p}%
\TextAt{13}{307.635}{367.727}{\Character{59}}%
\TextAt{8}{325.388}{367.727}{if}%
\TextAt{14}{336.111}{367.727}{p}%
\TextAt{8}{345.992}{367.727}{is}%
\TextAt{8}{357.756}{367.727}{complex}%
\TextAt{13}{399.706}{367.727}{\Character{59}}%
\TextAt{8}{79.226}{341.186}{for}%
\TextAt{13}{97.111}{341.186}{n}%
\TextAt{27}{104.099}{339.392}{p}%
\TextAt{13}{111.906}{341.186}{\Character{62}}%
\TextAt{9}{124.332}{341.186}{0}%
\TextAt{8}{130.185}{341.186}{.}%
\TextAt{8}{138.633}{341.186}{Given}%
\TextAt{13}{172.526}{341.186}{\Character{11}}%
\TextAt{27}{179.972}{339.392}{p}%
\NativeMathAt{187.779}{341.186}{11.955}{\in}%
\TextAt{13}{199.07}{341.186}{K}%
\TextAt{3}{209.888}{345.524}{\Character{3}}%
\TextAt{27}{209.029}{338.23}{p}%
\TextAt{8}{218.518}{341.186}{we}%
\TextAt{8}{235.745}{341.186}{write}%
\TextAt{13}{213.019}{313.147}{\Character{11}}%
\TextAt{27}{220.465}{311.354}{p}%
\NativeMathAt{228.273}{313.147}{11.955}{\equiv}%
\TextAt{9}{240.892}{313.147}{1}%
\TextAt{8}{250.066}{313.147}{mod}%
\TextAt{14}{275.504}{313.147}{p}%
\TextAt{21}{281.487}{318.083}{n}%
\TextAt{28}{286.625}{317.087}{p}%
\NativeArrowAt{306.482}{313.147}{11.955}{21.926}{\Longleftrightarrow}%
\TextAt{13}{343.676}{313.147}{\Character{11}}%
\TextAt{27}{351.121}{311.354}{p}%
\NativeMathAt{358.929}{313.147}{11.955}{\in}%
\TextAt{13}{370.22}{313.147}{U}%
\TextAt{17}{379.419}{319.392}{\Character{40}}%
\TextAt{21}{382.712}{319.392}{n}%
\TextAt{28}{387.851}{318.396}{p}%
\TextAt{17}{391.933}{319.392}{\Character{41}}%
\TextAt{27}{378.17}{311.176}{p}%
\TextAt{13}{395.725}{313.147}{\Character{58}}%
\TextAt{8}{79.226}{285.109}{For}%
\TextAt{8}{101.465}{285.109}{a}%
\TextAt{8}{112.489}{285.109}{finite}%
\TextAt{8}{143.67}{285.109}{prime}%
\TextAt{14}{178.103}{285.109}{p}%
\TextAt{8}{189.253}{285.109}{and}%
\TextAt{13}{213.285}{285.109}{n}%
\TextAt{27}{220.272}{283.315}{p}%
\TextAt{13}{230.243}{285.109}{\Character{62}}%
\TextAt{9}{244.831}{285.109}{0}%
\TextAt{8}{255.851}{285.109}{this}%
\TextAt{8}{279.942}{285.109}{means}%
\TextAt{8}{317.041}{285.109}{the}%
\TextAt{8}{338.467}{285.109}{usual}%
\TextAt{8}{370.365}{285.109}{congruence;}%
\TextAt{8}{436.002}{285.109}{for}%
\TextAt{8}{455.156}{285.109}{a}%
\TextAt{8}{466.181}{285.109}{real}%
\TextAt{8}{490.213}{285.109}{place,}%
\TextAt{8}{524.963}{285.109}{it}%
\TextAt{8}{79.226}{268.063}{symbolizes}%
\TextAt{8}{137.555}{268.063}{positivity,}%
\TextAt{8}{191.928}{268.063}{and}%
\TextAt{8}{214.69}{268.063}{for}%
\TextAt{8}{232.575}{268.063}{a}%
\TextAt{8}{242.331}{268.063}{complex}%
\TextAt{8}{288.179}{268.063}{place}%
\TextAt{8}{318.09}{268.063}{it}%
\TextAt{8}{329.795}{268.063}{is}%
\TextAt{8}{341.558}{268.063}{the}%
\TextAt{8}{361.715}{268.063}{empty}%
\TextAt{8}{397.485}{268.063}{condition.}%
\TextAt{12}{79.226}{242.325}{Definition}%
\TextAt{12}{142.015}{242.325}{(1.7).}%
\TextAt{8}{177.33}{242.325}{The}%
\TextAt{8}{201.384}{242.325}{group}%
\TextAt{13}{263.051}{225.28}{C}%
\TextAt{27}{272.284}{230.216}{m}%
\TextAt{21}{271.426}{222.324}{K}%
\TextAt{9}{282.949}{225.28}{\Character{61}}%
\TextAt{13}{295.375}{225.28}{I}%
\TextAt{27}{301.478}{230.216}{m}%
\TextAt{21}{300.538}{222.324}{K}%
\TextAt{13}{308.74}{225.28}{K}%
\TextAt{3}{319.558}{230.216}{\Character{3}}%
\TextAt{9}{324.29}{225.28}{\Character{47}}%
\TextAt{13}{330.143}{225.28}{K}%
\TextAt{3}{340.961}{230.216}{\Character{3}}%
\TextAt{13}{345.693}{225.28}{\Character{59}}%
\TextAt{8}{79.226}{202.17}{formed}%
\TextAt{8}{118.571}{202.17}{from}%
\TextAt{8}{146.211}{202.17}{the}%
\TextAt{8}{166.367}{202.17}{idèle}%
\TextAt{8}{193.673}{202.17}{group}%
\TextAt{13}{268.429}{185.124}{I}%
\TextAt{27}{274.532}{190.061}{m}%
\TextAt{21}{273.592}{182.169}{K}%
\TextAt{9}{285.116}{185.124}{\Character{61}}%
\TextAt{20}{297.541}{196.482}{\Character{89}}%
\TextAt{27}{303.185}{171.376}{p}%
\TextAt{13}{314.81}{185.124}{U}%
\TextAt{17}{324.009}{191.37}{\Character{40}}%
\TextAt{21}{327.302}{191.37}{n}%
\TextAt{28}{332.44}{190.373}{p}%
\TextAt{17}{336.523}{191.37}{\Character{41}}%
\TextAt{27}{322.76}{183.154}{p}%
\TextAt{13}{340.315}{185.124}{\Character{59}}%
\TextAt{8}{79.226}{152.648}{is}%
\TextAt{8}{91.845}{152.648}{called}%
\TextAt{8}{125.864}{152.648}{the}%
\TextAt{12}{146.875}{152.648}{congruence}%
\TextAt{12}{217.94}{152.648}{subgroup}%
\TextAt{8}{276.849}{152.648}{mod}%
\TextAt{14}{304.041}{152.648}{m}%
\TextAt{8}{313.208}{152.648}{,}%
\TextAt{8}{321.426}{152.648}{and}%
\TextAt{8}{345.044}{152.648}{the}%
\TextAt{8}{366.055}{152.648}{quotient}%
\TextAt{8}{413.093}{152.648}{group}%
\TextAt{13}{447.124}{152.648}{C}%
\TextAt{21}{455.499}{150.855}{K}%
\TextAt{9}{463.701}{152.648}{\Character{47}}%
\TextAt{13}{469.554}{152.648}{C}%
\TextAt{27}{478.788}{156.986}{m}%
\TextAt{21}{477.929}{149.627}{K}%
\TextAt{8}{490.884}{152.648}{is}%
\TextAt{8}{503.503}{152.648}{called}%
\TextAt{8}{79.226}{135.602}{the}%
\TextAt{12}{99.383}{135.602}{ray}%
\TextAt{12}{122.635}{135.602}{class}%
\TextAt{12}{153.994}{135.602}{group}%
\TextAt{8}{191.796}{135.602}{mod}%
\TextAt{14}{218.133}{135.602}{m}%
\TextAt{8}{227.3}{135.602}{.}%
\TextAt{12}{79.226}{107.013}{Remark.}%
\TextAt{8}{133.787}{107.013}{This}%
\TextAt{8}{160.91}{107.013}{definition}%
\TextAt{8}{213.342}{107.013}{of}%
\TextAt{8}{227.076}{107.013}{the}%
\TextAt{8}{247.636}{107.013}{ray}%
\TextAt{8}{268.208}{107.013}{class}%
\TextAt{8}{296.048}{107.013}{group}%
\TextAt{8}{329.628}{107.013}{does}%
\TextAt{8}{356.428}{107.013}{correspond}%
\TextAt{8}{416.404}{107.013}{to}%
\TextAt{8}{431.118}{107.013}{the}%
\TextAt{8}{451.678}{107.013}{classical}%
\TextAt{8}{497.081}{107.013}{one,}%
\TextAt{8}{522.297}{107.013}{as}%
\TextAt{8}{79.226}{89.967}{given}%
\TextAt{8}{109.951}{89.967}{(in}%
\TextAt{8}{128.315}{89.967}{the}%
\TextAt{8}{148.627}{89.967}{ideal-theoretic}%
\TextAt{8}{225.523}{89.967}{version)}%
\TextAt{8}{269.972}{89.967}{for}%
\TextAt{8}{288.013}{89.967}{instance}%
\TextAt{8}{333.754}{89.967}{in}%
\TextAt{8}{347.562}{89.967}{Hasse’s}%
\TextAt{8}{388.545}{89.967}{“Zahlbericht”}%
\TextAt{8}{461.854}{89.967}{[53].}%
\TextAt{8}{488.989}{89.967}{It}%
\TextAt{8}{501.818}{89.967}{differs}%
\TextAt{8}{79.226}{72.921}{from}%
\TextAt{8}{107.124}{72.921}{those}%
\TextAt{8}{138.011}{72.921}{found}%
\TextAt{8}{171.11}{72.921}{in}%
\TextAt{8}{185.021}{72.921}{modern}%
\TextAt{8}{227.875}{72.921}{textbooks,}%
\TextAt{8}{285.499}{72.921}{and}%
\TextAt{8}{308.519}{72.921}{also}%
\TextAt{8}{332.257}{72.921}{from}%
\TextAt{8}{360.155}{72.921}{that}%
\TextAt{8}{385.782}{72.921}{given}%
\TextAt{8}{416.61}{72.921}{in}%
\TextAt{8}{430.52}{72.921}{[107]}%
\TextAt{8}{458.753}{72.921}{by}%
\TextAt{8}{475.27}{72.921}{the}%
\TextAt{8}{495.685}{72.921}{author:}%
\TextAt{8}{297.211}{43.033}{331}%
\Destination{78.23}{809.815}{nav.page.331}%
\Destination{79.36}{606.055}{nav.print.363}%
\Destination{79.36}{256.355}{nav.number.6.1.7}%
\Hotlink{515.712}{758.955}{14.294}{11.047}{nav.chapter.3}%
\Hotlink{109.618}{707.755}{14.306}{11.047}{nav.chapter.3}%
\Hotlink{461.354}{87.755}{19.22}{11.047}{nav.cite.N053}%
\Hotlink{430.02}{70.755}{25.078}{11.047}{nav.cite.N107}%
\Hotlink{130.074}{707.755}{25.078}{11.047}{nav.number.3.1.9}%
}

% Page 344 (body page 332)
\Page{%
\TextAt{2}{79.226}{796.943}{VI.}%
\TextAt{2}{98.012}{796.943}{Global}%
\TextAt{2}{133.248}{796.943}{Class}%
\TextAt{2}{161.841}{796.943}{Field}%
\TextAt{2}{189.572}{796.943}{Theory}%
\Vector{q 0 G 0 g 0.299 w 7.226 18.729 m 460.769 18.729 l S Q}%
\TextAt{8}{79.226}{761.124}{in}%
\TextAt{8}{93.406}{761.124}{the}%
\TextAt{8}{114.088}{761.124}{present}%
\TextAt{8}{155.322}{761.124}{book,}%
\TextAt{8}{188.176}{761.124}{the}%
\TextAt{8}{208.859}{761.124}{components}%
\TextAt{13}{273.836}{761.124}{\Character{11}}%
\TextAt{27}{281.281}{759.33}{p}%
\TextAt{8}{290.192}{761.124}{of}%
\TextAt{8}{304.049}{761.124}{idèles}%
\TextAt{13}{336.496}{761.124}{\Character{11}}%
\TextAt{8}{348.441}{761.124}{in}%
\TextAt{13}{362.621}{761.124}{I}%
\TextAt{27}{368.723}{765.462}{m}%
\TextAt{21}{367.784}{758.103}{K}%
\TextAt{8}{380.41}{761.124}{are}%
\TextAt{8}{400.448}{761.124}{always}%
\TextAt{8}{438.442}{761.124}{positive}%
\TextAt{8}{482.282}{761.124}{at}%
\TextAt{8}{497.119}{761.124}{all}%
\TextAt{8}{513.904}{761.124}{real}%
\TextAt{8}{79.226}{744.078}{places}%
\TextAt{14}{113.231}{744.078}{p}%
\TextAt{8}{119.214}{744.078}{,}%
\TextAt{8}{125.945}{744.078}{so}%
\TextAt{8}{139.793}{744.078}{we}%
\TextAt{8}{156.499}{744.078}{have}%
\TextAt{8}{182.972}{744.078}{here}%
\TextAt{8}{207.806}{744.078}{fewer}%
\TextAt{8}{237.842}{744.078}{congruence}%
\TextAt{8}{297.801}{744.078}{subgroups}%
\TextAt{8}{352.691}{744.078}{than}%
\TextAt{8}{379.487}{744.078}{in}%
\TextAt{8}{392.618}{744.078}{the}%
\TextAt{8}{412.252}{744.078}{other}%
\TextAt{8}{442.299}{744.078}{texts.}%
\TextAt{8}{475.68}{744.078}{This}%
\TextAt{8}{501.877}{744.078}{choice}%
\TextAt{8}{79.226}{727.032}{does}%
\TextAt{8}{106.789}{727.032}{not}%
\TextAt{8}{128.768}{727.032}{only}%
\TextAt{8}{155.626}{727.032}{simplify}%
\TextAt{8}{201.073}{727.032}{matters.}%
\TextAt{8}{252.112}{727.032}{Most}%
\TextAt{8}{282.926}{727.032}{of}%
\TextAt{8}{297.422}{727.032}{all,}%
\TextAt{8}{318.39}{727.032}{it}%
\TextAt{8}{331.259}{727.032}{was}%
\TextAt{8}{354.924}{727.032}{made}%
\TextAt{8}{387.305}{727.032}{substantially}%
\TextAt{8}{457.548}{727.032}{because}%
\TextAt{8}{502.015}{727.032}{of}%
\TextAt{8}{516.51}{727.032}{the}%
\TextAt{8}{79.226}{709.986}{choice}%
\TextAt{8}{113.902}{709.986}{of}%
\TextAt{8}{127.118}{709.986}{the}%
\TextAt{8}{147.161}{709.986}{canonical}%
\TextAt{8}{198.43}{709.986}{metric}%
\TextAt{9}{234.732}{709.986}{\Character{40}}%
\TextAt{13}{243.067}{709.986}{\Character{59}}%
\TextAt{9}{252.095}{709.986}{\Character{41}}%
\TextAt{8}{260.431}{709.986}{on}%
\TextAt{8}{276.576}{709.986}{the}%
\TextAt{8}{296.619}{709.986}{Minkowski}%
\TextAt{8}{354.774}{709.986}{space}%
\TextAt{13}{385.935}{709.986}{K}%
\TextAt{25}{395.894}{708.193}{R}%
\TextAt{8}{405.932}{709.986}{(see}%
\TextAt{8}{429.286}{709.986}{chap.}%
\TextAt{8}{460.065}{709.986}{I,}%
\TextAt{8}{471.32}{709.986}{§5).}%
\TextAt{8}{495.822}{709.986}{In}%
\TextAt{8}{510.33}{709.986}{fact,}%
\TextAt{8}{79.226}{692.94}{we}%
\TextAt{8}{96.529}{692.94}{saw}%
\TextAt{8}{119.104}{692.94}{in}%
\TextAt{8}{132.832}{692.94}{chap.}%
\TextAt{8}{163.8}{692.94}{III,}%
\TextAt{8}{184.33}{692.94}{§3,}%
\TextAt{8}{203.111}{692.94}{that}%
\TextAt{8}{228.555}{692.94}{this}%
\TextAt{8}{251.453}{692.94}{choice}%
\TextAt{8}{286.317}{692.94}{forces}%
\TextAt{8}{319.293}{692.94}{the}%
\TextAt{8}{339.525}{692.94}{extension}%
\TextAt{18}{391.367}{692.94}{C}%
\TextAt{16}{403.45}{692.94}{\Character{106}}%
\TextAt{18}{410.22}{692.94}{R}%
\TextAt{8}{422.827}{692.94}{to}%
\TextAt{8}{437.213}{692.94}{be}%
\TextAt{8}{453.213}{692.94}{unramified.}%
\TextAt{8}{516.522}{692.94}{We}%
\TextAt{8}{79.226}{675.894}{will}%
\TextAt{8}{101.47}{675.894}{explain}%
\TextAt{8}{142.257}{675.894}{in}%
\TextAt{8}{156.049}{675.894}{§6}%
\TextAt{8}{171.622}{675.894}{below}%
\TextAt{8}{204.925}{675.894}{how}%
\TextAt{8}{229.452}{675.894}{to}%
\TextAt{8}{243.902}{675.894}{interpret}%
\TextAt{8}{292.495}{675.894}{this}%
\TextAt{8}{315.456}{675.894}{situation,}%
\TextAt{8}{367.731}{675.894}{and}%
\TextAt{8}{390.633}{675.894}{how}%
\TextAt{8}{415.16}{675.894}{to}%
\TextAt{8}{429.61}{675.894}{reconcile}%
\TextAt{8}{477.868}{675.894}{it}%
\TextAt{8}{489.711}{675.894}{with}%
\TextAt{8}{516.51}{675.894}{the}%
\TextAt{8}{79.226}{658.848}{definition}%
\TextAt{8}{131.255}{658.848}{of}%
\TextAt{8}{144.585}{658.848}{ray}%
\TextAt{8}{164.753}{658.848}{classes}%
\TextAt{8}{202.006}{658.848}{in}%
\TextAt{8}{215.659}{658.848}{other}%
\TextAt{8}{246.228}{658.848}{texts.}%
\TextAt{8}{97.159}{635.824}{The}%
\TextAt{8}{122.266}{635.824}{significance}%
\TextAt{8}{185.163}{635.824}{of}%
\TextAt{8}{199.546}{635.824}{the}%
\TextAt{8}{220.755}{635.824}{congruence}%
\TextAt{8}{282.29}{635.824}{subgroups}%
\TextAt{8}{338.755}{635.824}{lies}%
\TextAt{8}{360.024}{635.824}{in}%
\TextAt{8}{374.73}{635.824}{that}%
\TextAt{8}{401.152}{635.824}{they}%
\TextAt{8}{428.542}{635.824}{provide}%
\TextAt{8}{471.223}{635.824}{an}%
\TextAt{8}{488.535}{635.824}{overview}%
\TextAt{8}{79.226}{618.778}{over}%
\TextAt{8}{104.272}{618.778}{all}%
\TextAt{8}{120.531}{618.778}{closed}%
\TextAt{8}{155.058}{618.778}{subgroups}%
\TextAt{8}{210.47}{618.778}{of}%
\TextAt{8}{223.8}{618.778}{finite}%
\TextAt{8}{253.712}{618.778}{index}%
\TextAt{8}{285.25}{618.778}{in}%
\TextAt{13}{298.902}{618.778}{C}%
\TextAt{21}{307.277}{616.985}{K}%
\TextAt{8}{315.48}{618.778}{.}%
\TextAt{8}{323.928}{618.778}{More}%
\TextAt{8}{354.163}{618.778}{precisely,}%
\TextAt{8}{404.291}{618.778}{we}%
\TextAt{8}{421.519}{618.778}{have}%
\TextAt{8}{448.513}{618.778}{the}%
\TextAt{12}{79.226}{592.766}{Proposition}%
\TextAt{12}{152.237}{592.766}{(1.8).}%
\TextAt{8}{187.813}{592.766}{The}%
\TextAt{8}{212.127}{592.766}{closed}%
\TextAt{8}{246.914}{592.766}{subgroups}%
\TextAt{8}{302.586}{592.766}{of}%
\TextAt{8}{316.177}{592.766}{finite}%
\TextAt{8}{346.349}{592.766}{index}%
\TextAt{8}{378.147}{592.766}{of}%
\TextAt{13}{391.737}{592.766}{C}%
\TextAt{21}{400.112}{590.973}{K}%
\TextAt{8}{412.473}{592.766}{are}%
\TextAt{8}{432.244}{592.766}{precisely}%
\TextAt{8}{480.361}{592.766}{those}%
\TextAt{8}{511.25}{592.766}{sub-}%
\TextAt{8}{79.226}{575.72}{groups}%
\TextAt{8}{117.016}{575.72}{that}%
\TextAt{8}{142.385}{575.72}{contain}%
\TextAt{8}{183.69}{575.72}{a}%
\TextAt{8}{193.446}{575.72}{congruence}%
\TextAt{8}{253.927}{575.72}{subgroup}%
\TextAt{13}{304.725}{575.72}{C}%
\TextAt{27}{313.958}{580.058}{m}%
\TextAt{21}{313.1}{572.699}{K}%
\TextAt{8}{321.302}{575.72}{.}%
\TextAt{19}{79.226}{540.741}{Proof.}%
\TextAt{13}{116}{540.741}{C}%
\TextAt{27}{125.234}{545.08}{m}%
\TextAt{21}{124.375}{537.721}{K}%
\TextAt{8}{137.214}{540.741}{is}%
\TextAt{8}{149.716}{540.741}{open}%
\TextAt{8}{178.74}{540.741}{in}%
\TextAt{13}{193.131}{540.741}{C}%
\TextAt{21}{201.506}{538.948}{K}%
\TextAt{8}{214.344}{540.741}{because}%
\TextAt{13}{258.384}{540.741}{I}%
\TextAt{27}{264.487}{545.08}{m}%
\TextAt{21}{263.547}{537.721}{K}%
\TextAt{9}{276.329}{540.741}{\Character{61}}%
\TextAt{20}{290.012}{549.708}{\Character{81}}%
\TextAt{27}{301.303}{537.254}{p}%
\TextAt{13}{307.782}{540.741}{U}%
\TextAt{17}{316.982}{546.986}{\Character{40}}%
\TextAt{21}{320.275}{546.986}{n}%
\TextAt{28}{325.413}{545.99}{p}%
\TextAt{17}{329.496}{546.986}{\Character{41}}%
\TextAt{27}{315.733}{538.77}{p}%
\TextAt{8}{337.923}{540.741}{is}%
\TextAt{8}{350.425}{540.741}{open}%
\TextAt{8}{379.449}{540.741}{in}%
\TextAt{13}{393.84}{540.741}{I}%
\TextAt{21}{399.003}{538.948}{K}%
\TextAt{8}{407.206}{540.741}{.}%
\TextAt{13}{417.869}{540.741}{I}%
\TextAt{27}{423.971}{545.08}{m}%
\TextAt{21}{423.032}{537.721}{K}%
\TextAt{8}{435.87}{540.741}{is}%
\TextAt{8}{448.372}{540.741}{contained}%
\TextAt{8}{502.119}{540.741}{in}%
\TextAt{8}{516.51}{540.741}{the}%
\TextAt{8}{79.226}{523.695}{group}%
\TextAt{13}{113.018}{523.695}{I}%
\TextAt{21}{119.121}{528.809}{S}%
\NativeMathAt{124.271}{527.813}{5.978}{\infty}%
\TextAt{21}{118.181}{520.454}{K}%
\TextAt{9}{136.945}{523.695}{\Character{61}}%
\TextAt{20}{150.421}{532.662}{\Character{81}}%
\TextAt{27}{161.712}{520.208}{p}%
\TextAt{3}{165.701}{520.208}{\Character{106}}%
\NativeMathAt{168.052}{520.208}{7.97}{\infty}%
\TextAt{13}{179.012}{523.695}{K}%
\TextAt{3}{189.83}{528.034}{\Character{3}}%
\TextAt{27}{188.972}{520.74}{p}%
\TextAt{16}{197.639}{523.695}{\Character{2}}%
\TextAt{20}{210.015}{532.662}{\Character{81}}%
\TextAt{27}{221.306}{520.208}{p}%
\TextAt{25}{225.295}{520.208}{\Character{45}}%
\NativeMathAt{227.509}{520.208}{7.97}{\infty}%
\TextAt{13}{238.468}{523.695}{U}%
\TextAt{27}{246.418}{521.902}{p}%
\TextAt{8}{250.905}{523.695}{,}%
\TextAt{8}{258.825}{523.695}{and}%
\TextAt{8}{282.204}{523.695}{since}%
\TextAt{9}{311.489}{523.695}{\Character{40}}%
\TextAt{13}{316.041}{523.695}{C}%
\TextAt{21}{324.416}{521.902}{K}%
\TextAt{9}{336.991}{523.695}{\Character{58}}%
\TextAt{13}{344.614}{523.695}{I}%
\TextAt{21}{350.717}{528.809}{S}%
\NativeMathAt{355.867}{527.813}{5.978}{\infty}%
\TextAt{21}{349.777}{520.454}{K}%
\TextAt{13}{364.169}{523.695}{K}%
\TextAt{3}{374.987}{528.034}{\Character{3}}%
\TextAt{9}{379.719}{523.695}{\Character{47}}%
\TextAt{13}{385.572}{523.695}{K}%
\TextAt{3}{396.39}{528.034}{\Character{3}}%
\TextAt{9}{401.122}{523.695}{\Character{41}}%
\TextAt{9}{410.029}{523.695}{\Character{61}}%
\TextAt{8}{423.522}{523.695}{\char35\relax }%
\TextAt{13}{433.278}{523.695}{C}%
\TextAt{13}{442.495}{523.695}{l}%
\TextAt{21}{445.998}{521.902}{K}%
\TextAt{9}{458.573}{523.695}{\Character{61}}%
\TextAt{13}{472.049}{523.695}{h}%
\TextAt{13}{483.143}{523.695}{\Character{60}}%
\NativeMathAt{496.635}{523.695}{11.955}{\infty}%
\TextAt{8}{508.591}{523.695}{,}%
\TextAt{8}{516.51}{523.695}{the}%
\TextAt{8}{79.226}{506.649}{index}%
\TextAt{9}{189.871}{477.648}{\Character{40}}%
\TextAt{13}{194.423}{477.648}{C}%
\TextAt{21}{202.798}{475.855}{K}%
\TextAt{9}{214.322}{477.648}{\Character{58}}%
\TextAt{13}{220.894}{477.648}{C}%
\TextAt{27}{230.128}{482.584}{m}%
\TextAt{21}{229.269}{474.693}{K}%
\TextAt{9}{237.472}{477.648}{\Character{41}}%
\TextAt{9}{245.327}{477.648}{\Character{61}}%
\TextAt{13}{257.77}{477.648}{h}%
\TextAt{9}{264.509}{477.648}{\Character{40}}%
\TextAt{13}{269.061}{477.648}{I}%
\TextAt{21}{275.164}{482.762}{S}%
\NativeMathAt{280.314}{481.766}{5.978}{\infty}%
\TextAt{21}{274.224}{474.407}{K}%
\TextAt{13}{288.616}{477.648}{K}%
\TextAt{3}{299.434}{482.584}{\Character{3}}%
\TextAt{9}{307.487}{477.648}{\Character{58}}%
\TextAt{13}{314.06}{477.648}{I}%
\TextAt{27}{320.162}{482.584}{m}%
\TextAt{21}{319.223}{474.693}{K}%
\TextAt{13}{327.425}{477.648}{K}%
\TextAt{3}{338.243}{482.584}{\Character{3}}%
\TextAt{9}{342.975}{477.648}{\Character{41}}%
\NativeMathAt{350.848}{477.648}{11.955}{\leq}%
\TextAt{13}{363.468}{477.648}{h}%
\TextAt{9}{370.206}{477.648}{\Character{40}}%
\TextAt{13}{374.759}{477.648}{I}%
\TextAt{21}{380.861}{482.762}{S}%
\NativeMathAt{386.012}{481.766}{5.978}{\infty}%
\TextAt{21}{379.922}{474.407}{K}%
\TextAt{9}{397.635}{477.648}{\Character{58}}%
\TextAt{13}{404.207}{477.648}{I}%
\TextAt{27}{410.31}{482.584}{m}%
\TextAt{21}{409.37}{474.693}{K}%
\TextAt{9}{417.573}{477.648}{\Character{41}}%
\TextAt{9}{245.345}{455.996}{\Character{61}}%
\TextAt{13}{257.77}{455.996}{h}%
\TextAt{20}{266.501}{467.354}{\Character{89}}%
\TextAt{27}{266.804}{441.522}{p}%
\TextAt{25}{270.793}{441.522}{\Character{45}}%
\NativeMathAt{273.007}{441.522}{7.97}{\infty}%
\TextAt{9}{281.778}{455.996}{\Character{40}}%
\TextAt{13}{286.33}{455.996}{U}%
\TextAt{27}{294.28}{454.203}{p}%
\TextAt{9}{302.088}{455.996}{\Character{58}}%
\TextAt{13}{308.66}{455.996}{U}%
\TextAt{17}{317.859}{462.241}{\Character{40}}%
\TextAt{21}{321.153}{462.241}{n}%
\TextAt{28}{326.291}{461.245}{p}%
\TextAt{17}{330.374}{462.241}{\Character{41}}%
\TextAt{27}{316.611}{454.025}{p}%
\TextAt{9}{334.165}{455.996}{\Character{41}}%
\TextAt{20}{340.71}{467.354}{\Character{89}}%
\TextAt{27}{340.943}{441.451}{p}%
\TextAt{3}{344.932}{441.451}{\Character{106}}%
\NativeMathAt{347.283}{441.451}{7.97}{\infty}%
\TextAt{9}{355.986}{455.996}{\Character{40}}%
\TextAt{13}{360.538}{455.996}{K}%
\TextAt{3}{371.356}{460.933}{\Character{3}}%
\TextAt{27}{370.497}{453.041}{p}%
\TextAt{9}{379.409}{455.996}{\Character{58}}%
\TextAt{13}{385.981}{455.996}{U}%
\TextAt{17}{395.18}{462.241}{\Character{40}}%
\TextAt{21}{398.474}{462.241}{n}%
\TextAt{28}{403.612}{461.245}{p}%
\TextAt{17}{407.695}{462.241}{\Character{41}}%
\TextAt{27}{393.932}{454.025}{p}%
\TextAt{9}{411.486}{455.996}{\Character{41}}%
\TextAt{8}{79.226}{416.895}{is}%
\TextAt{8}{90.836}{416.895}{finite.}%
\TextAt{8}{125.248}{416.895}{Being}%
\TextAt{8}{158.091}{416.895}{the}%
\TextAt{8}{178.094}{416.895}{complement}%
\TextAt{8}{243.299}{416.895}{of}%
\TextAt{8}{256.476}{416.895}{the}%
\TextAt{8}{276.478}{416.895}{nontrivial}%
\TextAt{8}{329.669}{416.895}{open}%
\TextAt{8}{357.801}{416.895}{cosets,}%
\TextAt{8}{394.871}{416.895}{which}%
\TextAt{8}{428.204}{416.895}{are}%
\TextAt{8}{447.561}{416.895}{finite}%
\TextAt{8}{477.32}{416.895}{in}%
\TextAt{8}{490.819}{416.895}{number,}%
\TextAt{13}{79.226}{399.849}{C}%
\TextAt{27}{88.46}{404.187}{m}%
\TextAt{21}{87.601}{396.828}{K}%
\TextAt{8}{100.435}{399.849}{is}%
\TextAt{8}{112.934}{399.849}{closed}%
\TextAt{8}{148.194}{399.849}{of}%
\TextAt{8}{162.259}{399.849}{finite}%
\TextAt{8}{192.905}{399.849}{index.}%
\TextAt{8}{231.196}{399.849}{Consequently,}%
\TextAt{8}{306.965}{399.849}{every}%
\TextAt{8}{338.592}{399.849}{group}%
\TextAt{8}{372.502}{399.849}{containing}%
\TextAt{13}{430.154}{399.849}{C}%
\TextAt{27}{439.388}{404.187}{m}%
\TextAt{21}{438.529}{396.828}{K}%
\TextAt{8}{451.364}{399.849}{is}%
\TextAt{8}{463.862}{399.849}{also}%
\TextAt{8}{488.076}{399.849}{closed}%
\TextAt{8}{523.337}{399.849}{of}%
\TextAt{8}{79.226}{382.803}{finite}%
\TextAt{8}{109.138}{382.803}{index,}%
\TextAt{8}{143.928}{382.803}{for}%
\TextAt{8}{161.812}{382.803}{it}%
\TextAt{8}{173.517}{382.803}{is}%
\TextAt{8}{185.28}{382.803}{the}%
\TextAt{8}{205.437}{382.803}{union}%
\TextAt{8}{237.955}{382.803}{of}%
\TextAt{8}{251.285}{382.803}{finitely}%
\TextAt{8}{290.629}{382.803}{many}%
\TextAt{8}{322.502}{382.803}{cosets}%
\TextAt{8}{356.443}{382.803}{of}%
\TextAt{13}{369.773}{382.803}{C}%
\TextAt{27}{379.006}{387.141}{m}%
\TextAt{21}{378.148}{379.782}{K}%
\TextAt{8}{386.35}{382.803}{.}%
\TextAt{8}{97.159}{362.768}{Conversely,}%
\TextAt{8}{160.045}{362.768}{let}%
\TextAt{16}{178.041}{362.768}{\Character{78}}%
\TextAt{8}{194.6}{362.768}{be}%
\TextAt{8}{211.616}{362.768}{an}%
\TextAt{8}{228.967}{362.768}{arbitrary}%
\TextAt{8}{279.828}{362.768}{closed}%
\TextAt{8}{315.446}{362.768}{subgroup}%
\TextAt{8}{367.335}{362.768}{of}%
\TextAt{8}{381.756}{362.768}{finite}%
\TextAt{8}{412.76}{362.768}{index.}%
\TextAt{8}{452.123}{362.768}{Then}%
\TextAt{16}{483.772}{362.768}{\Character{78}}%
\TextAt{8}{500.331}{362.768}{is}%
\TextAt{8}{513.187}{362.768}{also}%
\TextAt{8}{79.226}{345.722}{open,}%
\TextAt{8}{110.925}{345.722}{being}%
\TextAt{8}{142.591}{345.722}{the}%
\TextAt{8}{162.876}{345.722}{complement}%
\TextAt{8}{228.363}{345.722}{of}%
\TextAt{8}{241.822}{345.722}{a}%
\TextAt{8}{251.706}{345.722}{finite}%
\TextAt{8}{281.747}{345.722}{number}%
\TextAt{8}{324.471}{345.722}{of}%
\TextAt{8}{337.93}{345.722}{closed}%
\TextAt{8}{372.585}{345.722}{cosets.}%
\TextAt{8}{411.463}{345.722}{Thus}%
\TextAt{8}{441.24}{345.722}{the}%
\TextAt{8}{461.525}{345.722}{preimage}%
\TextAt{13}{511.734}{345.722}{J}%
\TextAt{8}{523.337}{345.722}{of}%
\TextAt{16}{79.226}{328.676}{\Character{78}}%
\TextAt{8}{94.694}{328.676}{in}%
\TextAt{13}{108.347}{328.676}{I}%
\TextAt{21}{113.51}{326.883}{K}%
\TextAt{8}{125.61}{328.676}{is}%
\TextAt{8}{137.374}{328.676}{also}%
\TextAt{8}{160.854}{328.676}{open,}%
\TextAt{8}{192.392}{328.676}{and}%
\TextAt{8}{215.154}{328.676}{it}%
\TextAt{8}{226.858}{328.676}{thus}%
\TextAt{8}{252.61}{328.676}{contains}%
\TextAt{8}{298.53}{328.676}{a}%
\TextAt{8}{308.285}{328.676}{subset}%
\TextAt{8}{344.175}{328.676}{of}%
\TextAt{8}{357.505}{328.676}{the}%
\TextAt{8}{377.661}{328.676}{form}%
\TextAt{13}{251.598}{299.675}{W}%
\TextAt{9}{267.596}{299.675}{\Character{61}}%
\TextAt{20}{280.022}{311.033}{\Character{89}}%
\TextAt{27}{280.045}{285.661}{p}%
\NativeMathAt{284.033}{285.661}{7.97}{\in}%
\TextAt{21}{289.679}{285.661}{S}%
\TextAt{13}{297.29}{299.675}{W}%
\TextAt{27}{308.342}{297.882}{p}%
\TextAt{16}{315.485}{299.675}{\Character{2}}%
\TextAt{20}{327.44}{311.033}{\Character{89}}%
\TextAt{27}{327.463}{285.13}{p}%
\TextAt{17}{332.393}{285.13}{\Character{47}}%
\NativeMathAt{331.452}{285.13}{7.97}{\in}%
\TextAt{21}{337.098}{285.13}{S}%
\TextAt{13}{344.709}{299.675}{U}%
\TextAt{27}{352.659}{297.882}{p}%
\TextAt{13}{357.146}{299.675}{\Character{59}}%
\TextAt{8}{79.226}{261.059}{where}%
\TextAt{13}{114.222}{261.059}{S}%
\TextAt{8}{127.202}{261.059}{is}%
\TextAt{8}{140.153}{261.059}{a}%
\TextAt{8}{151.095}{261.059}{finite}%
\TextAt{8}{182.194}{261.059}{set}%
\TextAt{8}{201.649}{261.059}{of}%
\TextAt{8}{216.166}{261.059}{places}%
\TextAt{8}{251.879}{261.059}{of}%
\TextAt{13}{266.396}{261.059}{K}%
\TextAt{8}{282.298}{261.059}{containing}%
\TextAt{8}{340.404}{261.059}{the}%
\TextAt{8}{361.748}{261.059}{infinite}%
\TextAt{8}{402.602}{261.059}{ones,}%
\TextAt{8}{433.412}{261.059}{and}%
\TextAt{13}{457.361}{261.059}{W}%
\TextAt{27}{468.413}{259.266}{p}%
\TextAt{8}{477.984}{261.059}{is}%
\TextAt{8}{490.935}{261.059}{an}%
\TextAt{8}{508.381}{261.059}{open}%
\TextAt{8}{79.226}{244.013}{neighbourhood}%
\TextAt{8}{159.386}{244.013}{of}%
\TextAt{9}{172.871}{244.013}{1}%
\NativeMathAt{182.31}{244.013}{11.955}{\in}%
\TextAt{13}{193.866}{244.013}{K}%
\TextAt{3}{204.684}{248.352}{\Character{3}}%
\TextAt{27}{203.825}{241.058}{p}%
\TextAt{8}{209.416}{244.013}{.}%
\TextAt{8}{218.331}{244.013}{If}%
\TextAt{14}{230.178}{244.013}{p}%
\NativeMathAt{239.747}{244.013}{11.955}{\in}%
\TextAt{13}{251.303}{244.013}{S}%
\TextAt{8}{263.251}{244.013}{is}%
\TextAt{8}{275.171}{244.013}{finite,}%
\TextAt{8}{308.529}{244.013}{we}%
\TextAt{8}{325.912}{244.013}{are}%
\TextAt{8}{345.578}{244.013}{liable}%
\TextAt{8}{376.948}{244.013}{to}%
\TextAt{8}{391.414}{244.013}{choose}%
\TextAt{13}{428.702}{244.013}{W}%
\TextAt{27}{439.754}{242.22}{p}%
\TextAt{9}{447.827}{244.013}{\Character{61}}%
\TextAt{13}{460.517}{244.013}{U}%
\TextAt{17}{469.716}{250.259}{\Character{40}}%
\TextAt{21}{473.009}{250.259}{n}%
\TextAt{28}{478.148}{249.262}{p}%
\TextAt{17}{482.23}{250.259}{\Character{41}}%
\TextAt{27}{468.467}{242.043}{p}%
\TextAt{8}{486.022}{244.013}{,}%
\TextAt{8}{493.365}{244.013}{because}%
\TextAt{8}{79.226}{226.967}{the}%
\TextAt{8}{99.042}{226.967}{groups}%
\TextAt{13}{136.492}{226.967}{U}%
\TextAt{17}{145.691}{233.213}{\Character{40}}%
\TextAt{21}{148.984}{233.213}{n}%
\TextAt{28}{154.123}{232.216}{p}%
\TextAt{17}{158.205}{233.213}{\Character{41}}%
\TextAt{27}{144.442}{224.997}{p}%
\NativeMathAt{165.318}{226.967}{11.955}{\subseteq}%
\TextAt{13}{177.937}{226.967}{K}%
\TextAt{3}{188.755}{231.306}{\Character{3}}%
\TextAt{27}{187.896}{224.012}{p}%
\TextAt{8}{197.044}{226.967}{form}%
\TextAt{8}{224.344}{226.967}{a}%
\TextAt{8}{233.759}{226.967}{basic}%
\TextAt{8}{262.745}{226.967}{system}%
\TextAt{8}{301.223}{226.967}{of}%
\TextAt{8}{314.212}{226.967}{neighbourhoods}%
\TextAt{8}{398.491}{226.967}{of}%
\TextAt{9}{411.48}{226.967}{1}%
\NativeMathAt{420.654}{226.967}{11.955}{\in}%
\TextAt{13}{431.945}{226.967}{K}%
\TextAt{3}{442.763}{231.306}{\Character{3}}%
\TextAt{27}{441.904}{224.012}{p}%
\TextAt{8}{447.495}{226.967}{.}%
\TextAt{8}{455.83}{226.967}{If}%
\TextAt{14}{467.182}{226.967}{p}%
\NativeMathAt{476.486}{226.967}{11.955}{\in}%
\TextAt{13}{487.777}{226.967}{S}%
\TextAt{8}{499.229}{226.967}{is}%
\TextAt{8}{510.652}{226.967}{real,}%
\TextAt{8}{79.226}{209.921}{we}%
\TextAt{8}{96.959}{209.921}{may}%
\TextAt{8}{122.834}{209.921}{choose}%
\TextAt{13}{160.472}{209.921}{W}%
\TextAt{27}{171.524}{208.128}{p}%
\NativeMathAt{180.194}{209.921}{11.955}{\subseteq}%
\TextAt{18}{193.675}{209.921}{R}%
\TextAt{3}{202.309}{214.26}{\Character{3}}%
\TextAt{17}{202.309}{206.966}{\Character{43}}%
\TextAt{8}{209.394}{209.921}{.}%
\TextAt{8}{219.359}{209.921}{The}%
\TextAt{8}{243.919}{209.921}{open}%
\TextAt{8}{272.711}{209.921}{set}%
\TextAt{13}{291.484}{209.921}{W}%
\TextAt{27}{302.535}{208.128}{p}%
\TextAt{8}{311.425}{209.921}{will}%
\TextAt{8}{334.036}{209.921}{then}%
\TextAt{8}{361.202}{209.921}{generate}%
\TextAt{8}{408.536}{209.921}{the}%
\TextAt{8}{429.198}{209.921}{group}%
\TextAt{18}{462.88}{209.921}{R}%
\TextAt{3}{471.514}{214.26}{\Character{3}}%
\TextAt{17}{471.514}{206.966}{\Character{43}}%
\TextAt{8}{478.599}{209.921}{,}%
\TextAt{8}{486.38}{209.921}{resp.}%
\TextAt{13}{517.22}{209.921}{K}%
\TextAt{3}{528.037}{214.26}{\Character{3}}%
\TextAt{27}{527.179}{206.966}{p}%
\TextAt{8}{79.226}{192.875}{in}%
\TextAt{8}{93.15}{192.875}{the}%
\TextAt{8}{113.577}{192.875}{case}%
\TextAt{8}{138.618}{192.875}{of}%
\TextAt{8}{152.219}{192.875}{a}%
\TextAt{8}{162.245}{192.875}{complex}%
\TextAt{8}{208.364}{192.875}{place}%
\TextAt{14}{238.546}{192.875}{p}%
\TextAt{8}{244.529}{192.875}{.}%
\TextAt{8}{253.789}{192.875}{The}%
\TextAt{8}{278.114}{192.875}{subgroup}%
\TextAt{8}{329.182}{192.875}{of}%
\TextAt{13}{342.782}{192.875}{J}%
\TextAt{8}{354.527}{192.875}{generated}%
\TextAt{8}{408.13}{192.875}{by}%
\TextAt{13}{424.66}{192.875}{W}%
\TextAt{8}{441.505}{192.875}{is}%
\TextAt{8}{453.539}{192.875}{therefore}%
\TextAt{8}{502.91}{192.875}{of}%
\TextAt{8}{516.51}{192.875}{the}%
\TextAt{8}{79.226}{175.83}{form}%
\TextAt{13}{106.866}{175.83}{I}%
\TextAt{27}{112.969}{180.168}{m}%
\TextAt{21}{112.029}{172.809}{K}%
\TextAt{8}{120.232}{175.83}{,}%
\TextAt{8}{127.381}{175.83}{so}%
\TextAt{16}{141.751}{175.83}{\Character{78}}%
\TextAt{8}{157.22}{175.83}{contains}%
\TextAt{8}{203.139}{175.83}{the}%
\TextAt{8}{223.296}{175.83}{congruence}%
\TextAt{8}{283.777}{175.83}{subgroup}%
\TextAt{13}{334.575}{175.83}{C}%
\TextAt{27}{343.808}{180.168}{m}%
\TextAt{21}{342.95}{172.809}{K}%
\TextAt{8}{351.152}{175.83}{.}%
\Vector{q 0 G 0 g 0.399 w 452.334 -598.17 m 452.334 -590.101 l S Q}%
\Vector{q 0 G 0 g 0.399 w 452.534 -590.3 m 459.707 -590.3 l S Q}%
\Vector{q 0 G 0 g 0.399 w 452.534 -597.971 m 459.707 -597.971 l S Q}%
\Vector{q 0 G 0 g 0.399 w 459.906 -598.17 m 459.906 -590.101 l S Q}%
\TextAt{8}{97.159}{146.828}{The}%
\TextAt{8}{120.707}{146.828}{ray}%
\TextAt{8}{140.369}{146.828}{class}%
\TextAt{8}{167.3}{146.828}{groups}%
\TextAt{8}{204.585}{146.828}{can}%
\TextAt{8}{225.538}{146.828}{be}%
\TextAt{8}{240.956}{146.828}{given}%
\TextAt{8}{271.02}{146.828}{the}%
\TextAt{8}{290.67}{146.828}{following}%
\TextAt{8}{339.599}{146.828}{purely}%
\TextAt{8}{375.186}{146.828}{ideal-theoretic}%
\TextAt{8}{451.42}{146.828}{description.}%
\TextAt{8}{515.697}{146.828}{Let}%
\TextAt{13}{79.226}{129.782}{J}%
\TextAt{27}{86.803}{134.121}{m}%
\TextAt{21}{85.696}{126.762}{K}%
\TextAt{8}{97.701}{129.782}{be}%
\TextAt{8}{113.53}{129.782}{the}%
\TextAt{8}{133.591}{129.782}{group}%
\TextAt{8}{166.671}{129.782}{of}%
\TextAt{8}{179.906}{129.782}{all}%
\TextAt{8}{196.069}{129.782}{fractional}%
\TextAt{8}{248.338}{129.782}{ideals}%
\TextAt{8}{280.82}{129.782}{relatively}%
\TextAt{8}{331.785}{129.782}{prime}%
\TextAt{8}{364.853}{129.782}{to}%
\TextAt{14}{379.068}{129.782}{m}%
\TextAt{8}{388.235}{129.782}{,}%
\TextAt{8}{395.308}{129.782}{and}%
\TextAt{8}{417.975}{129.782}{let}%
\TextAt{13}{434.784}{129.782}{P}%
\TextAt{27}{443.955}{134.121}{m}%
\TextAt{21}{442.33}{126.762}{K}%
\TextAt{8}{454.367}{129.782}{be}%
\TextAt{8}{470.196}{129.782}{the}%
\TextAt{8}{490.257}{129.782}{group}%
\TextAt{8}{523.337}{129.782}{of}%
\TextAt{8}{79.226}{112.736}{all}%
\TextAt{8}{95.485}{112.736}{principal}%
\TextAt{8}{144.262}{112.736}{ideals}%
\TextAt{9}{176.84}{112.736}{\Character{40}}%
\TextAt{13}{181.392}{112.736}{a}%
\TextAt{9}{187.537}{112.736}{\Character{41}}%
\NativeMathAt{195.411}{112.736}{11.955}{\in}%
\TextAt{13}{206.701}{112.736}{P}%
\TextAt{21}{214.247}{110.943}{K}%
\TextAt{8}{226.347}{112.736}{such}%
\TextAt{8}{252.744}{112.736}{that}%
\TextAt{13}{194.664}{83.735}{a}%
\NativeMathAt{204.13}{83.735}{11.955}{\equiv}%
\TextAt{9}{216.749}{83.735}{1}%
\TextAt{8}{225.923}{83.735}{mod}%
\TextAt{14}{251.361}{83.735}{m}%
\TextAt{8}{284.438}{83.735}{and}%
\TextAt{13}{327.214}{83.735}{a}%
\TextAt{8}{337.256}{83.735}{totally}%
\TextAt{8}{374.664}{83.735}{positive}%
\TextAt{13}{414.08}{83.735}{\Character{58}}%
\TextAt{8}{297.211}{43.033}{332}%
\Destination{78.23}{809.815}{nav.page.332}%
\Destination{79.36}{724.955}{nav.print.364}%
\Destination{79.36}{161.855}{nav.print.365}%
\Destination{79.36}{606.755}{nav.number.6.1.8}%
\Hotlink{459.565}{707.755}{5.22}{11.047}{nav.chapter.1}%
\Hotlink{476.499}{707.755}{6.858}{11.047}{nav.section.1.5}%
\Hotlink{163.3}{690.755}{14.294}{11.047}{nav.chapter.3}%
\Hotlink{189.509}{690.755}{6.858}{11.047}{nav.section.3.3}%
\Hotlink{161.228}{673.655}{6.858}{11.047}{nav.section.6.6}%
}

% Page 345 (body page 333)
\Page{%
\TextAt{2}{79.226}{796.943}{VI.}%
\TextAt{2}{98.012}{796.943}{Global}%
\TextAt{2}{133.248}{796.943}{Class}%
\TextAt{2}{161.841}{796.943}{Field}%
\TextAt{2}{189.572}{796.943}{Theory}%
\Vector{q 0 G 0 g 0.299 w 7.226 18.729 m 460.769 18.729 l S Q}%
\TextAt{8}{79.226}{761.124}{The}%
\TextAt{8}{104.857}{761.124}{latter}%
\TextAt{8}{138.306}{761.124}{condition}%
\TextAt{8}{191.267}{761.124}{means}%
\TextAt{8}{228.673}{761.124}{that,}%
\TextAt{8}{259.265}{761.124}{for}%
\TextAt{8}{278.727}{761.124}{every}%
\TextAt{8}{311.196}{761.124}{real}%
\TextAt{8}{335.535}{761.124}{embedding}%
\TextAt{13}{396.29}{761.124}{K}%
\NativeArrowAt{413.116}{761.124}{11.955}{19.26}{\longrightarrow}%
\TextAt{18}{438.386}{761.124}{R}%
\TextAt{8}{447.02}{761.124}{,}%
\TextAt{13}{456.14}{761.124}{a}%
\TextAt{8}{467.76}{761.124}{turns}%
\TextAt{8}{499.966}{761.124}{out}%
\TextAt{8}{522.357}{761.124}{to}%
\TextAt{8}{79.226}{744.078}{be}%
\TextAt{8}{97.041}{744.078}{positive.}%
\TextAt{8}{150.577}{744.078}{The}%
\TextAt{8}{176.521}{744.078}{congruence}%
\TextAt{13}{238.893}{744.078}{a}%
\NativeMathAt{251.58}{744.078}{11.955}{\equiv}%
\TextAt{9}{267.421}{744.078}{1}%
\TextAt{8}{276.595}{744.078}{mod}%
\TextAt{14}{302.032}{744.078}{m}%
\TextAt{8}{316.987}{744.078}{means}%
\TextAt{8}{354.707}{744.078}{that}%
\TextAt{13}{381.966}{744.078}{a}%
\TextAt{8}{393.899}{744.078}{is}%
\TextAt{8}{407.553}{744.078}{the}%
\TextAt{8}{429.6}{744.078}{quotient}%
\TextAt{13}{477.674}{744.078}{b}%
\TextAt{9}{482.651}{744.078}{\Character{47}}%
\TextAt{13}{488.504}{744.078}{c}%
\TextAt{8}{499.329}{744.078}{of}%
\TextAt{8}{514.55}{744.078}{two}%
\TextAt{8}{79.226}{727.032}{integers}%
\TextAt{8}{122.576}{727.032}{relatively}%
\TextAt{8}{173.673}{727.032}{prime}%
\TextAt{8}{206.873}{727.032}{to}%
\TextAt{14}{221.22}{727.032}{m}%
\TextAt{8}{234.321}{727.032}{such}%
\TextAt{8}{260.754}{727.032}{that}%
\TextAt{13}{286.16}{727.032}{b}%
\NativeMathAt{294.52}{727.032}{11.955}{\equiv}%
\TextAt{13}{307.202}{727.032}{c}%
\TextAt{8}{315.56}{727.032}{mod}%
\TextAt{14}{340.998}{727.032}{m}%
\TextAt{8}{350.165}{727.032}{.}%
\TextAt{8}{358.723}{727.032}{This}%
\TextAt{8}{385.479}{727.032}{is}%
\TextAt{8}{397.28}{727.032}{tantamount}%
\TextAt{8}{461.073}{727.032}{to}%
\TextAt{8}{475.42}{727.032}{saying}%
\TextAt{8}{511.298}{727.032}{that}%
\TextAt{13}{79.226}{709.986}{a}%
\NativeMathAt{88.841}{709.986}{11.955}{\equiv}%
\TextAt{9}{101.61}{709.986}{1}%
\TextAt{8}{110.784}{709.986}{mod}%
\TextAt{14}{136.222}{709.986}{p}%
\TextAt{21}{142.205}{714.324}{n}%
\TextAt{28}{147.343}{713.328}{p}%
\TextAt{8}{155.909}{709.986}{in}%
\TextAt{13}{169.65}{709.986}{K}%
\TextAt{27}{179.609}{708.193}{p}%
\TextAt{8}{184.096}{709.986}{,}%
\TextAt{8}{191.355}{709.986}{i.e.,}%
\TextAt{13}{213.569}{709.986}{a}%
\NativeMathAt{223.185}{709.986}{11.955}{\in}%
\TextAt{13}{234.625}{709.986}{U}%
\TextAt{17}{243.824}{716.231}{\Character{40}}%
\TextAt{21}{247.117}{716.231}{n}%
\TextAt{28}{252.256}{715.235}{p}%
\TextAt{17}{256.338}{716.231}{\Character{41}}%
\TextAt{27}{242.575}{708.015}{p}%
\TextAt{8}{264.115}{709.986}{for}%
\TextAt{8}{282.087}{709.986}{all}%
\TextAt{14}{298.434}{709.986}{p}%
\TextAt{16}{307.887}{709.986}{\Character{106}}%
\TextAt{14}{314.679}{709.986}{m}%
\TextAt{9}{327.316}{709.986}{\Character{61}}%
\TextAt{20}{339.891}{718.952}{\Character{81}}%
\TextAt{27}{351.182}{706.499}{p}%
\TextAt{25}{355.17}{706.499}{\Character{45}}%
\NativeMathAt{357.384}{706.499}{7.97}{\infty}%
\TextAt{14}{368.343}{709.986}{p}%
\TextAt{21}{374.326}{714.324}{n}%
\TextAt{28}{379.465}{713.328}{p}%
\TextAt{8}{384.045}{709.986}{.}%
\TextAt{8}{392.757}{709.986}{We}%
\TextAt{8}{412.989}{709.986}{put}%
\TextAt{13}{434.536}{709.986}{C}%
\TextAt{13}{443.754}{709.986}{l}%
\TextAt{27}{447.52}{714.324}{m}%
\TextAt{21}{447.257}{706.965}{K}%
\TextAt{9}{458.93}{709.986}{\Character{61}}%
\TextAt{13}{471.505}{709.986}{J}%
\TextAt{27}{479.082}{714.324}{m}%
\TextAt{21}{477.975}{706.965}{K}%
\TextAt{9}{486.178}{709.986}{\Character{47}}%
\TextAt{13}{492.031}{709.986}{P}%
\TextAt{27}{501.202}{714.324}{m}%
\TextAt{21}{499.576}{706.965}{K}%
\TextAt{8}{507.811}{709.986}{.}%
\TextAt{8}{516.523}{709.986}{We}%
\TextAt{8}{79.226}{692.94}{then}%
\TextAt{8}{105.886}{692.94}{have}%
\TextAt{8}{132.881}{692.94}{the}%
\TextAt{12}{79.226}{667.31}{Proposition}%
\TextAt{12}{151.937}{667.31}{(1.9).}%
\TextAt{8}{187.253}{667.31}{The}%
\TextAt{8}{211.307}{667.31}{homomorphism}%
\TextAt{9}{194.446}{639.646}{\Character{40}}%
\TextAt{9}{202.874}{639.646}{\Character{41}}%
\TextAt{9}{210.741}{639.646}{\Character{58}}%
\TextAt{13}{217.341}{639.646}{I}%
\TextAt{21}{222.504}{637.853}{K}%
\NativeArrowAt{234.027}{639.646}{11.955}{19.26}{\longrightarrow}%
\TextAt{13}{256.61}{639.646}{J}%
\TextAt{21}{263.08}{637.853}{K}%
\TextAt{13}{271.282}{639.646}{\Character{59}}%
\TextAt{13}{300.417}{639.646}{\Character{11}}%
\NativeArrowAt{311.279}{639.646}{11.955}{19.26}{\longmapsto}%
\TextAt{9}{333.861}{639.646}{\Character{40}}%
\TextAt{13}{338.414}{639.646}{\Character{11}}%
\TextAt{9}{345.935}{639.646}{\Character{41}}%
\TextAt{9}{353.79}{639.646}{\Character{61}}%
\TextAt{20}{366.234}{651.003}{\Character{89}}%
\TextAt{27}{366.537}{625.171}{p}%
\TextAt{25}{370.525}{625.171}{\Character{45}}%
\NativeMathAt{372.739}{625.171}{7.97}{\infty}%
\TextAt{14}{383.503}{639.646}{p}%
\TextAt{21}{389.486}{644.582}{v}%
\TextAt{28}{393.609}{643.586}{p}%
\TextAt{17}{397.692}{644.582}{\Character{40}}%
\TextAt{21}{400.985}{644.582}{\Character{11}}%
\TextAt{28}{406.424}{643.586}{p}%
\TextAt{17}{410.507}{644.582}{\Character{41}}%
\TextAt{13}{414.298}{639.646}{\Character{59}}%
\TextAt{8}{79.226}{601.392}{induces}%
\TextAt{8}{120.902}{601.392}{an}%
\TextAt{8}{137.161}{601.392}{isomorphism}%
\TextAt{13}{205.58}{601.392}{C}%
\TextAt{21}{213.955}{599.598}{K}%
\TextAt{9}{222.158}{601.392}{\Character{47}}%
\TextAt{13}{228.011}{601.392}{C}%
\TextAt{27}{237.245}{605.73}{m}%
\TextAt{21}{236.386}{598.371}{K}%
\NativeMathAt{247.91}{601.392}{11.955}{\cong}%
\TextAt{13}{260.529}{601.392}{C}%
\TextAt{13}{269.747}{601.392}{l}%
\TextAt{27}{273.512}{605.73}{m}%
\TextAt{21}{273.249}{598.371}{K}%
\TextAt{8}{281.452}{601.392}{.}%
\TextAt{19}{79.226}{567.177}{Proof.}%
\TextAt{8}{116}{567.177}{Let}%
\TextAt{14}{136.97}{567.177}{m}%
\TextAt{9}{149.457}{567.177}{\Character{61}}%
\TextAt{20}{161.883}{576.144}{\Character{81}}%
\TextAt{27}{173.174}{563.69}{p}%
\TextAt{14}{179.653}{567.177}{p}%
\TextAt{21}{185.636}{571.516}{n}%
\TextAt{28}{190.774}{570.519}{p}%
\TextAt{8}{195.355}{567.177}{,}%
\TextAt{8}{202.504}{567.177}{and}%
\TextAt{8}{225.267}{567.177}{let}%
\TextAt{13}{202.692}{538.347}{I}%
\TextAt{17}{208.795}{544.457}{\Character{40}}%
\TextAt{27}{212.088}{544.457}{m}%
\TextAt{17}{218.199}{544.457}{\Character{41}}%
\TextAt{21}{207.855}{535.106}{K}%
\TextAt{9}{225.311}{538.347}{\Character{61}}%
\TextAt{16}{237.737}{538.347}{\Character{102}}%
\TextAt{13}{243.714}{538.347}{\Character{11}}%
\NativeMathAt{254.557}{538.347}{11.955}{\in}%
\TextAt{13}{265.848}{538.347}{I}%
\TextAt{21}{271.011}{536.553}{K}%
\TextAt{16}{282.534}{538.347}{\Character{106}}%
\TextAt{13}{289.176}{538.347}{\Character{11}}%
\TextAt{27}{296.622}{536.553}{p}%
\NativeMathAt{304.429}{538.347}{11.955}{\in}%
\TextAt{13}{315.72}{538.347}{U}%
\TextAt{17}{324.919}{544.592}{\Character{40}}%
\TextAt{21}{328.213}{544.592}{n}%
\TextAt{28}{333.351}{543.596}{p}%
\TextAt{17}{337.433}{544.592}{\Character{41}}%
\TextAt{27}{323.67}{536.376}{p}%
\TextAt{8}{345.122}{538.347}{for}%
\TextAt{14}{363.007}{538.347}{p}%
\TextAt{16}{372.311}{538.347}{\Character{106}}%
\TextAt{14}{378.953}{538.347}{m}%
\NativeMathAt{388.12}{538.347}{11.955}{\infty}%
\TextAt{16}{400.075}{538.347}{\Character{103}}%
\TextAt{13}{406.052}{538.347}{\Character{58}}%
\TextAt{8}{79.226}{510.683}{Then}%
\TextAt{13}{109.357}{510.683}{I}%
\TextAt{21}{114.52}{508.89}{K}%
\TextAt{9}{126.043}{510.683}{\Character{61}}%
\TextAt{13}{138.469}{510.683}{I}%
\TextAt{17}{144.571}{516.794}{\Character{40}}%
\TextAt{27}{147.865}{516.794}{m}%
\TextAt{17}{153.976}{516.794}{\Character{41}}%
\TextAt{21}{143.631}{507.442}{K}%
\TextAt{13}{157.767}{510.683}{K}%
\TextAt{3}{168.585}{515.022}{\Character{3}}%
\TextAt{8}{173.317}{510.683}{,}%
\TextAt{8}{180.125}{510.683}{because}%
\TextAt{8}{223}{510.683}{for}%
\TextAt{8}{240.458}{510.683}{every}%
\TextAt{13}{270.923}{510.683}{\Character{11}}%
\NativeMathAt{281.765}{510.683}{11.955}{\in}%
\TextAt{13}{293.056}{510.683}{I}%
\TextAt{21}{298.219}{508.89}{K}%
\TextAt{8}{306.422}{510.683}{,}%
\TextAt{8}{313.23}{510.683}{by}%
\TextAt{8}{329.062}{510.683}{the}%
\TextAt{8}{348.791}{510.683}{approximation}%
\TextAt{8}{426.432}{510.683}{theorem,}%
\TextAt{8}{474.867}{510.683}{there}%
\TextAt{8}{504.352}{510.683}{exists}%
\TextAt{8}{79.226}{493.637}{an}%
\TextAt{13}{96.668}{493.637}{a}%
\NativeMathAt{108.149}{493.637}{11.955}{\in}%
\TextAt{13}{121.456}{493.637}{K}%
\TextAt{3}{132.274}{497.976}{\Character{3}}%
\TextAt{8}{142.086}{493.637}{such}%
\TextAt{8}{169.666}{493.637}{that}%
\TextAt{13}{196.217}{493.637}{\Character{11}}%
\TextAt{27}{203.663}{491.844}{p}%
\TextAt{13}{208.15}{493.637}{a}%
\NativeMathAt{219.631}{493.637}{11.955}{\equiv}%
\TextAt{9}{234.266}{493.637}{1}%
\TextAt{8}{243.44}{493.637}{mod}%
\TextAt{14}{268.878}{493.637}{p}%
\TextAt{21}{274.861}{497.976}{n}%
\TextAt{28}{279.999}{496.979}{p}%
\TextAt{8}{289.66}{493.637}{for}%
\TextAt{14}{308.728}{493.637}{p}%
\TextAt{16}{320.048}{493.637}{\Character{106}}%
\TextAt{14}{328.705}{493.637}{m}%
\TextAt{8}{337.872}{493.637}{,}%
\TextAt{8}{346.499}{493.637}{and}%
\TextAt{13}{370.445}{493.637}{\Character{11}}%
\TextAt{27}{377.89}{491.844}{p}%
\TextAt{13}{382.377}{493.637}{a}%
\TextAt{13}{393.854}{493.637}{\Character{62}}%
\TextAt{9}{408.3}{493.637}{0}%
\TextAt{8}{419.233}{493.637}{for}%
\TextAt{14}{438.3}{493.637}{p}%
\TextAt{8}{449.364}{493.637}{real.}%
\TextAt{8}{480.226}{493.637}{Thus}%
\TextAt{13}{511.057}{493.637}{\Character{12}}%
\TextAt{9}{523.665}{493.637}{\Character{61}}%
\TextAt{9}{79.226}{476.591}{\Character{40}}%
\TextAt{13}{83.778}{476.591}{\Character{11}}%
\TextAt{27}{91.224}{474.798}{p}%
\TextAt{13}{95.711}{476.591}{a}%
\TextAt{9}{101.856}{476.591}{\Character{41}}%
\NativeMathAt{109.785}{476.591}{11.955}{\in}%
\TextAt{13}{121.131}{476.591}{I}%
\TextAt{17}{127.234}{482.702}{\Character{40}}%
\TextAt{27}{130.527}{482.702}{m}%
\TextAt{17}{136.638}{482.702}{\Character{41}}%
\TextAt{21}{126.294}{473.35}{K}%
\TextAt{8}{140.43}{476.591}{,}%
\TextAt{8}{147.62}{476.591}{so}%
\TextAt{8}{162.022}{476.591}{that}%
\TextAt{13}{187.424}{476.591}{\Character{11}}%
\TextAt{9}{198.322}{476.591}{\Character{61}}%
\TextAt{13}{210.803}{476.591}{\Character{12}}%
\TextAt{13}{218.072}{476.591}{a}%
\TextAt{3}{224.219}{480.93}{\Character{0}}%
\TextAt{17}{230.806}{480.93}{1}%
\NativeMathAt{238.914}{476.591}{11.955}{\in}%
\TextAt{13}{250.261}{476.591}{I}%
\TextAt{17}{256.364}{482.702}{\Character{40}}%
\TextAt{27}{259.657}{482.702}{m}%
\TextAt{17}{265.768}{482.702}{\Character{41}}%
\TextAt{21}{255.424}{473.35}{K}%
\TextAt{13}{269.559}{476.591}{K}%
\TextAt{3}{280.377}{480.93}{\Character{3}}%
\TextAt{8}{285.109}{476.591}{.}%
\TextAt{8}{293.655}{476.591}{The}%
\TextAt{8}{317.742}{476.591}{elements}%
\TextAt{13}{365.631}{476.591}{a}%
\NativeMathAt{375.152}{476.591}{11.955}{\in}%
\TextAt{13}{386.499}{476.591}{I}%
\TextAt{17}{392.602}{482.702}{\Character{40}}%
\TextAt{27}{395.895}{482.702}{m}%
\TextAt{17}{402.006}{482.702}{\Character{41}}%
\TextAt{21}{391.662}{473.35}{K}%
\TextAt{16}{408.476}{476.591}{\Character{92}}%
\TextAt{13}{419.125}{476.591}{K}%
\TextAt{3}{429.943}{480.93}{\Character{3}}%
\TextAt{8}{438.605}{476.591}{are}%
\TextAt{8}{458.149}{476.591}{precisely}%
\TextAt{8}{506.038}{476.591}{those}%
\TextAt{8}{79.226}{459.545}{generating}%
\TextAt{8}{136.735}{459.545}{principal}%
\TextAt{8}{185.779}{459.545}{ideals}%
\TextAt{8}{218.624}{459.545}{in}%
\TextAt{13}{232.544}{459.545}{P}%
\TextAt{27}{241.715}{463.884}{m}%
\TextAt{21}{240.09}{456.525}{K}%
\TextAt{8}{248.325}{459.545}{.}%
\TextAt{8}{257.575}{459.545}{Therefore}%
\TextAt{8}{310.839}{459.545}{the}%
\TextAt{8}{331.263}{459.545}{correspondence}%
\TextAt{13}{413.208}{459.545}{\Character{11}}%
\NativeArrowAt{424.506}{459.545}{11.955}{11.955}{\mapsto}%
\TextAt{9}{440.238}{459.545}{\Character{40}}%
\TextAt{13}{444.79}{459.545}{\Character{11}}%
\TextAt{9}{452.312}{459.545}{\Character{41}}%
\TextAt{9}{460.621}{459.545}{\Character{61}}%
\TextAt{20}{473.521}{468.512}{\Character{81}}%
\TextAt{27}{484.812}{456.058}{p}%
\TextAt{25}{488.801}{456.058}{\Character{45}}%
\NativeMathAt{491.015}{456.058}{7.97}{\infty}%
\TextAt{14}{501.974}{459.545}{p}%
\TextAt{21}{507.957}{463.884}{v}%
\TextAt{28}{512.08}{462.887}{p}%
\TextAt{17}{516.163}{463.884}{\Character{40}}%
\TextAt{21}{519.456}{463.884}{\Character{11}}%
\TextAt{28}{524.895}{462.887}{p}%
\TextAt{17}{528.978}{463.884}{\Character{41}}%
\TextAt{8}{79.226}{442.499}{defines}%
\TextAt{8}{117.65}{442.499}{a}%
\TextAt{8}{127.405}{442.499}{surjective}%
\TextAt{8}{179.817}{442.499}{homomorphism}%
\TextAt{13}{179.912}{414.836}{C}%
\TextAt{21}{188.287}{413.042}{K}%
\TextAt{9}{199.811}{414.836}{\Character{61}}%
\TextAt{13}{212.236}{414.836}{I}%
\TextAt{17}{218.339}{420.946}{\Character{40}}%
\TextAt{27}{221.632}{420.946}{m}%
\TextAt{17}{227.744}{420.946}{\Character{41}}%
\TextAt{21}{217.399}{411.594}{K}%
\TextAt{13}{231.535}{414.836}{K}%
\TextAt{3}{242.353}{419.772}{\Character{3}}%
\TextAt{9}{247.085}{414.836}{\Character{47}}%
\TextAt{13}{252.938}{414.836}{K}%
\TextAt{3}{263.756}{419.772}{\Character{3}}%
\TextAt{9}{271.809}{414.836}{\Character{61}}%
\TextAt{13}{284.234}{414.836}{I}%
\TextAt{17}{290.337}{420.946}{\Character{40}}%
\TextAt{27}{293.63}{420.946}{m}%
\TextAt{17}{299.742}{420.946}{\Character{41}}%
\TextAt{21}{289.397}{411.594}{K}%
\TextAt{9}{303.533}{414.836}{\Character{47}\Character{40}}%
\TextAt{13}{313.938}{414.836}{I}%
\TextAt{17}{320.041}{420.946}{\Character{40}}%
\TextAt{27}{323.334}{420.946}{m}%
\TextAt{17}{329.445}{420.946}{\Character{41}}%
\TextAt{21}{319.101}{411.594}{K}%
\TextAt{16}{335.893}{414.836}{\Character{92}}%
\TextAt{13}{346.52}{414.836}{K}%
\TextAt{3}{357.338}{419.772}{\Character{3}}%
\TextAt{9}{362.07}{414.836}{\Character{41}}%
\NativeArrowAt{369.943}{414.836}{11.955}{19.26}{\longrightarrow}%
\TextAt{13}{392.525}{414.836}{J}%
\TextAt{27}{400.103}{419.772}{m}%
\TextAt{21}{398.996}{411.88}{K}%
\TextAt{9}{407.198}{414.836}{\Character{47}}%
\TextAt{13}{413.051}{414.836}{P}%
\TextAt{27}{422.222}{419.772}{m}%
\TextAt{21}{420.597}{411.88}{K}%
\TextAt{13}{428.832}{414.836}{\Character{58}}%
\TextAt{8}{79.226}{387.172}{Since}%
\TextAt{9}{109.75}{387.172}{\Character{40}}%
\TextAt{13}{114.302}{387.172}{\Character{11}}%
\TextAt{9}{121.824}{387.172}{\Character{41}}%
\TextAt{9}{129.679}{387.172}{\Character{61}}%
\TextAt{9}{142.1}{387.172}{1}%
\TextAt{8}{151.839}{387.172}{for}%
\TextAt{13}{169.69}{387.172}{\Character{11}}%
\NativeMathAt{180.532}{387.172}{11.955}{\in}%
\TextAt{13}{191.823}{387.172}{I}%
\TextAt{27}{197.926}{391.51}{m}%
\TextAt{21}{196.986}{384.151}{K}%
\TextAt{8}{205.189}{387.172}{,}%
\TextAt{8}{212.311}{387.172}{the}%
\TextAt{8}{232.434}{387.172}{group}%
\TextAt{13}{265.575}{387.172}{C}%
\TextAt{27}{274.809}{391.51}{m}%
\TextAt{21}{273.95}{384.151}{K}%
\TextAt{9}{285.474}{387.172}{\Character{61}}%
\TextAt{13}{297.899}{387.172}{I}%
\TextAt{27}{304.002}{391.51}{m}%
\TextAt{21}{303.062}{384.151}{K}%
\TextAt{13}{311.265}{387.172}{K}%
\TextAt{3}{322.082}{391.51}{\Character{3}}%
\TextAt{9}{326.815}{387.172}{\Character{47}}%
\TextAt{13}{332.668}{387.172}{K}%
\TextAt{3}{343.485}{391.51}{\Character{3}}%
\TextAt{8}{352.081}{387.172}{is}%
\TextAt{8}{363.811}{387.172}{certainly}%
\TextAt{8}{412.232}{387.172}{contained}%
\TextAt{8}{465.207}{387.172}{in}%
\TextAt{8}{478.826}{387.172}{the}%
\TextAt{8}{498.948}{387.172}{kernel.}%
\TextAt{8}{79.226}{370.126}{Conversely,}%
\TextAt{8}{140.59}{370.126}{if}%
\TextAt{8}{151.118}{370.126}{the}%
\TextAt{8}{171.078}{370.126}{class}%
\TextAt{9}{198.319}{370.126}{\Character{91}}%
\TextAt{13}{201.57}{370.126}{\Character{11}}%
\TextAt{9}{209.092}{370.126}{\Character{93}}%
\TextAt{8}{216.045}{370.126}{represented}%
\TextAt{8}{278.015}{370.126}{by}%
\TextAt{13}{294.078}{370.126}{\Character{11}}%
\NativeMathAt{304.921}{370.126}{11.955}{\in}%
\TextAt{13}{316.212}{370.126}{I}%
\TextAt{17}{322.314}{376.236}{\Character{40}}%
\TextAt{27}{325.608}{376.236}{m}%
\TextAt{17}{331.719}{376.236}{\Character{41}}%
\TextAt{21}{321.375}{366.885}{K}%
\TextAt{8}{339.211}{370.126}{belongs}%
\TextAt{8}{381.025}{370.126}{to}%
\TextAt{8}{395.139}{370.126}{the}%
\TextAt{8}{415.1}{370.126}{kernel,}%
\TextAt{8}{452.661}{370.126}{then}%
\TextAt{8}{479.125}{370.126}{there}%
\TextAt{8}{508.84}{370.126}{is}%
\TextAt{8}{520.408}{370.126}{an}%
\TextAt{9}{79.226}{353.08}{\Character{40}}%
\TextAt{13}{83.778}{353.08}{a}%
\TextAt{9}{89.923}{353.08}{\Character{41}}%
\NativeMathAt{97.797}{353.08}{11.955}{\in}%
\TextAt{13}{109.088}{353.08}{P}%
\TextAt{27}{118.259}{357.418}{m}%
\TextAt{21}{116.633}{350.059}{K}%
\TextAt{8}{124.868}{353.08}{,}%
\TextAt{8}{131.613}{353.08}{with}%
\TextAt{13}{157.768}{353.08}{a}%
\NativeMathAt{167.234}{353.08}{11.955}{\in}%
\TextAt{13}{178.525}{353.08}{I}%
\TextAt{17}{184.628}{359.19}{\Character{40}}%
\TextAt{27}{187.921}{359.19}{m}%
\TextAt{17}{194.032}{359.19}{\Character{41}}%
\TextAt{21}{183.688}{349.839}{K}%
\TextAt{16}{199.448}{353.08}{\Character{92}}%
\TextAt{13}{209.042}{353.08}{K}%
\TextAt{3}{219.859}{357.418}{\Character{3}}%
\TextAt{8}{224.592}{353.08}{,}%
\TextAt{8}{231.337}{353.08}{such}%
\TextAt{8}{257.229}{353.08}{that}%
\TextAt{9}{282.093}{353.08}{\Character{40}}%
\TextAt{13}{286.645}{353.08}{\Character{11}}%
\TextAt{9}{294.167}{353.08}{\Character{41}}%
\TextAt{9}{302.022}{353.08}{\Character{61}}%
\TextAt{9}{314.443}{353.08}{\Character{40}}%
\TextAt{13}{319.018}{353.08}{a}%
\TextAt{9}{325.163}{353.08}{\Character{41}}%
\TextAt{8}{329.715}{353.08}{.}%
\TextAt{8}{337.995}{353.08}{The}%
\TextAt{8}{361.544}{353.08}{components}%
\TextAt{8}{425.489}{353.08}{of}%
\TextAt{8}{438.314}{353.08}{the}%
\TextAt{8}{457.966}{353.08}{idèle}%
\TextAt{13}{484.766}{353.08}{\Character{12}}%
\TextAt{9}{495.359}{353.08}{\Character{61}}%
\TextAt{13}{507.784}{353.08}{\Character{11}a}%
\TextAt{3}{521.451}{357.418}{\Character{0}}%
\TextAt{17}{528.037}{357.418}{1}%
\TextAt{8}{79.226}{336.034}{satisfy}%
\TextAt{13}{115.616}{336.034}{\Character{12}}%
\TextAt{27}{122.224}{334.241}{p}%
\NativeMathAt{130.032}{336.034}{11.955}{\in}%
\TextAt{13}{141.323}{336.034}{U}%
\TextAt{27}{149.273}{334.241}{p}%
\TextAt{8}{157.5}{336.034}{for}%
\TextAt{14}{175.227}{336.034}{p}%
\TextAt{18}{184.531}{336.034}{\Character{45}}%
\TextAt{14}{191.173}{336.034}{m}%
\NativeMathAt{200.34}{336.034}{11.955}{\infty}%
\TextAt{8}{212.295}{336.034}{,}%
\TextAt{8}{219.318}{336.034}{and}%
\TextAt{13}{241.923}{336.034}{\Character{12}}%
\TextAt{27}{248.532}{334.241}{p}%
\NativeMathAt{256.339}{336.034}{11.955}{\in}%
\TextAt{13}{267.63}{336.034}{U}%
\TextAt{17}{276.829}{342.279}{\Character{40}}%
\TextAt{21}{280.122}{342.279}{n}%
\TextAt{28}{285.261}{341.283}{p}%
\TextAt{17}{289.343}{342.279}{\Character{41}}%
\TextAt{27}{275.58}{334.063}{p}%
\TextAt{8}{296.875}{336.034}{for}%
\TextAt{14}{314.602}{336.034}{p}%
\TextAt{16}{323.906}{336.034}{\Character{106}}%
\TextAt{14}{330.548}{336.034}{m}%
\NativeMathAt{339.714}{336.034}{11.955}{\infty}%
\TextAt{8}{351.67}{336.034}{,}%
\TextAt{8}{358.693}{336.034}{in}%
\TextAt{8}{372.188}{336.034}{other}%
\TextAt{8}{402.6}{336.034}{words,}%
\TextAt{13}{439.284}{336.034}{\Character{12}}%
\NativeMathAt{449.876}{336.034}{11.955}{\in}%
\TextAt{13}{461.167}{336.034}{I}%
\TextAt{27}{467.27}{340.372}{m}%
\TextAt{21}{466.33}{333.014}{K}%
\TextAt{8}{474.533}{336.034}{,}%
\TextAt{8}{481.556}{336.034}{and}%
\TextAt{8}{504.161}{336.034}{hence}%
\TextAt{9}{79.226}{318.988}{\Character{91}}%
\TextAt{13}{82.478}{318.988}{\Character{11}}%
\TextAt{9}{90}{318.988}{\Character{93}}%
\TextAt{9}{97.089}{318.988}{\Character{61}}%
\TextAt{9}{110.037}{318.988}{\Character{91}}%
\TextAt{13}{113.319}{318.988}{\Character{12}}%
\TextAt{9}{120.59}{318.988}{\Character{93}}%
\NativeMathAt{127.698}{318.988}{11.955}{\in}%
\TextAt{13}{139.523}{318.988}{I}%
\TextAt{27}{145.626}{323.327}{m}%
\TextAt{21}{144.686}{315.968}{K}%
\TextAt{13}{152.889}{318.988}{K}%
\TextAt{3}{163.707}{323.327}{\Character{3}}%
\TextAt{9}{168.439}{318.988}{\Character{47}}%
\TextAt{13}{174.292}{318.988}{K}%
\TextAt{3}{185.11}{323.327}{\Character{3}}%
\TextAt{9}{193.698}{318.988}{\Character{61}}%
\TextAt{13}{206.658}{318.988}{C}%
\TextAt{27}{215.892}{323.327}{m}%
\TextAt{21}{215.033}{315.968}{K}%
\TextAt{8}{223.236}{318.988}{.}%
\TextAt{8}{232.626}{318.988}{Therefore}%
\TextAt{13}{285.937}{318.988}{C}%
\TextAt{27}{295.171}{323.327}{m}%
\TextAt{21}{294.312}{315.968}{K}%
\TextAt{8}{306.726}{318.988}{is}%
\TextAt{8}{318.804}{318.988}{the}%
\TextAt{8}{339.274}{318.988}{kernel}%
\TextAt{8}{374.055}{318.988}{of}%
\TextAt{8}{387.698}{318.988}{the}%
\TextAt{8}{408.169}{318.988}{above}%
\TextAt{8}{441.658}{318.988}{mapping,}%
\TextAt{8}{493.434}{318.988}{and}%
\TextAt{8}{516.51}{318.988}{the}%
\TextAt{8}{79.226}{301.942}{proposition}%
\TextAt{8}{140.759}{301.942}{is}%
\TextAt{8}{152.523}{301.942}{proved.}%
\Vector{q 0 G 0 g 0.399 w 452.334 -472.058 m 452.334 -463.988 l S Q}%
\Vector{q 0 G 0 g 0.399 w 452.534 -464.187 m 459.707 -464.187 l S Q}%
\Vector{q 0 G 0 g 0.399 w 452.534 -471.859 m 459.707 -471.859 l S Q}%
\Vector{q 0 G 0 g 0.399 w 459.906 -472.058 m 459.906 -463.988 l S Q}%
\TextAt{8}{97.159}{273.323}{The}%
\TextAt{8}{121.291}{273.323}{ray}%
\TextAt{8}{141.537}{273.323}{class}%
\TextAt{8}{169.052}{273.323}{groups}%
\TextAt{8}{206.92}{273.323}{in}%
\TextAt{8}{220.651}{273.323}{the}%
\TextAt{8}{240.885}{273.323}{ideal-theoretic}%
\TextAt{8}{317.703}{273.323}{version}%
\TextAt{13}{357.52}{273.323}{C}%
\TextAt{13}{366.738}{273.323}{l}%
\TextAt{27}{370.504}{277.662}{m}%
\TextAt{21}{370.241}{270.303}{K}%
\TextAt{9}{381.897}{273.323}{\Character{61}}%
\TextAt{13}{394.455}{273.323}{J}%
\TextAt{27}{402.032}{277.662}{m}%
\TextAt{21}{400.925}{270.303}{K}%
\TextAt{9}{409.128}{273.323}{\Character{47}}%
\TextAt{13}{414.981}{273.323}{P}%
\TextAt{27}{424.152}{277.662}{m}%
\TextAt{21}{422.526}{270.303}{K}%
\TextAt{8}{434.737}{273.323}{were}%
\TextAt{8}{461.797}{273.323}{introduced}%
\TextAt{8}{520.408}{273.323}{by}%
\TextAt{26}{79.226}{256.277}{Heinrich}%
\TextAt{26}{136.418}{256.277}{Weber}%
\TextAt{8}{181.123}{256.277}{(1842–1913)}%
\TextAt{8}{246.965}{256.277}{as}%
\TextAt{8}{261.447}{256.277}{a}%
\TextAt{8}{271.314}{256.277}{common}%
\TextAt{8}{318.255}{256.277}{generalization}%
\TextAt{8}{393.17}{256.277}{of}%
\TextAt{8}{406.612}{256.277}{ideal}%
\TextAt{8}{434.688}{256.277}{class}%
\TextAt{8}{462.237}{256.277}{groups}%
\TextAt{8}{500.139}{256.277}{on}%
\TextAt{8}{516.51}{256.277}{the}%
\TextAt{8}{79.226}{239.231}{one}%
\TextAt{8}{100.911}{239.231}{hand,}%
\TextAt{8}{133.711}{239.231}{and}%
\TextAt{8}{156.7}{239.231}{the}%
\TextAt{8}{177.082}{239.231}{groups}%
\TextAt{9}{215.098}{239.231}{\Character{40}}%
\TextAt{18}{219.65}{239.231}{Z}%
\TextAt{9}{227.62}{239.231}{\Character{47}}%
\TextAt{13}{233.473}{239.231}{m}%
\TextAt{18}{243.712}{239.231}{Z}%
\TextAt{9}{251.682}{239.231}{\Character{41}}%
\TextAt{3}{256.235}{243.57}{\Character{3}}%
\TextAt{8}{265.09}{239.231}{on}%
\TextAt{8}{281.575}{239.231}{the}%
\TextAt{8}{301.957}{239.231}{other.}%
\TextAt{8}{337.754}{239.231}{These}%
\TextAt{8}{371.849}{239.231}{latter}%
\TextAt{8}{403.947}{239.231}{groups}%
\TextAt{8}{441.963}{239.231}{may}%
\TextAt{8}{467.557}{239.231}{be}%
\TextAt{8}{483.707}{239.231}{viewed}%
\TextAt{8}{522.297}{239.231}{as}%
\TextAt{8}{79.226}{222.185}{the}%
\TextAt{8}{99.383}{222.185}{ray}%
\TextAt{8}{119.551}{222.185}{class}%
\TextAt{8}{146.988}{222.185}{groups}%
\TextAt{8}{184.778}{222.185}{of}%
\TextAt{8}{198.108}{222.185}{the}%
\TextAt{8}{218.265}{222.185}{field}%
\TextAt{18}{243.622}{222.185}{Q}%
\TextAt{8}{252.92}{222.185}{:}%
\TextAt{12}{79.226}{196.555}{Proposition}%
\TextAt{12}{151.937}{196.555}{(1.10).}%
\TextAt{8}{193.984}{196.555}{For}%
\TextAt{8}{214.953}{196.555}{any}%
\TextAt{8}{237.07}{196.555}{module}%
\TextAt{14}{278.363}{196.555}{m}%
\TextAt{9}{290.851}{196.555}{\Character{61}}%
\TextAt{9}{303.261}{196.555}{\Character{40}}%
\TextAt{13}{307.829}{196.555}{m}%
\TextAt{9}{318.068}{196.555}{\Character{41}}%
\TextAt{8}{326.518}{196.555}{of}%
\TextAt{8}{339.848}{196.555}{the}%
\TextAt{8}{360.004}{196.555}{field}%
\TextAt{18}{385.361}{196.555}{Q}%
\TextAt{8}{394.66}{196.555}{,}%
\TextAt{8}{401.809}{196.555}{one}%
\TextAt{8}{423.268}{196.555}{has}%
\TextAt{13}{237.317}{168.891}{C}%
\TextAt{25}{245.692}{167.098}{Q}%
\TextAt{9}{252.389}{168.891}{\Character{47}}%
\TextAt{13}{258.242}{168.891}{C}%
\TextAt{27}{267.475}{173.828}{m}%
\TextAt{25}{266.617}{165.936}{Q}%
\NativeMathAt{277.406}{168.891}{11.955}{\cong}%
\TextAt{13}{290.025}{168.891}{C}%
\TextAt{13}{299.243}{168.891}{l}%
\TextAt{27}{303.008}{173.828}{m}%
\TextAt{25}{302.745}{165.936}{Q}%
\NativeMathAt{312.938}{168.891}{11.955}{\cong}%
\TextAt{9}{325.558}{168.891}{\Character{40}}%
\TextAt{18}{330.11}{168.891}{Z}%
\TextAt{9}{338.08}{168.891}{\Character{47}}%
\TextAt{13}{343.933}{168.891}{m}%
\TextAt{18}{354.172}{168.891}{Z}%
\TextAt{9}{362.143}{168.891}{\Character{41}}%
\TextAt{3}{366.695}{173.828}{\Character{3}}%
\TextAt{13}{371.427}{168.891}{\Character{58}}%
\TextAt{19}{79.226}{124.059}{Proof.}%
\TextAt{8}{116}{124.059}{Every}%
\TextAt{8}{148.477}{124.059}{ideal}%
\TextAt{9}{175.264}{124.059}{\Character{40}}%
\TextAt{13}{179.816}{124.059}{a}%
\TextAt{9}{185.961}{124.059}{\Character{41}}%
\NativeMathAt{193.834}{124.059}{11.955}{\in}%
\TextAt{13}{205.125}{124.059}{J}%
\TextAt{27}{212.702}{128.398}{m}%
\TextAt{25}{211.595}{120.994}{Q}%
\TextAt{8}{222.032}{124.059}{has}%
\TextAt{8}{241.729}{124.059}{two}%
\TextAt{8}{262.669}{124.059}{generators,}%
\TextAt{13}{321.635}{124.059}{a}%
\TextAt{8}{330.501}{124.059}{and}%
\TextAt{16}{352.087}{124.059}{\Character{0}}%
\TextAt{13}{361.385}{124.059}{a}%
\TextAt{8}{367.53}{124.059}{.}%
\TextAt{8}{375.586}{124.059}{Mapping}%
\TextAt{8}{423.51}{124.059}{the}%
\TextAt{8}{442.489}{124.059}{positive}%
\TextAt{8}{484.626}{124.059}{generator}%
\TextAt{8}{79.226}{107.013}{onto}%
\TextAt{8}{105.165}{107.013}{the}%
\TextAt{8}{124.911}{107.013}{residue}%
\TextAt{8}{164.228}{107.013}{class}%
\TextAt{8}{191.255}{107.013}{mod}%
\TextAt{14}{217.182}{107.013}{m}%
\TextAt{8}{226.349}{107.013}{,}%
\TextAt{8}{233.17}{107.013}{we}%
\TextAt{8}{249.987}{107.013}{get}%
\TextAt{8}{269.088}{107.013}{a}%
\TextAt{8}{278.433}{107.013}{surjective}%
\TextAt{8}{330.435}{107.013}{homomorphism}%
\TextAt{13}{412.694}{107.013}{J}%
\TextAt{27}{420.272}{111.352}{m}%
\TextAt{25}{419.165}{103.948}{Q}%
\NativeArrowAt{430.202}{107.013}{11.955}{19.26}{\longrightarrow}%
\TextAt{9}{452.784}{107.013}{\Character{40}}%
\TextAt{18}{457.336}{107.013}{Z}%
\TextAt{9}{465.306}{107.013}{\Character{47}}%
\TextAt{13}{471.159}{107.013}{m}%
\TextAt{18}{481.398}{107.013}{Z}%
\TextAt{9}{489.369}{107.013}{\Character{41}}%
\TextAt{3}{493.921}{111.352}{\Character{3}}%
\TextAt{8}{502.14}{107.013}{whose}%
\TextAt{8}{79.226}{89.967}{kernel}%
\TextAt{8}{113.347}{89.967}{consists}%
\TextAt{8}{156.112}{89.967}{of}%
\TextAt{8}{169.096}{89.967}{all}%
\TextAt{8}{185.009}{89.967}{ideals}%
\TextAt{9}{217.241}{89.967}{\Character{40}}%
\TextAt{13}{221.793}{89.967}{a}%
\TextAt{9}{227.938}{89.967}{\Character{41}}%
\TextAt{8}{236.042}{89.967}{which}%
\TextAt{8}{269.183}{89.967}{have}%
\TextAt{8}{295.832}{89.967}{a}%
\TextAt{8}{305.242}{89.967}{positive}%
\TextAt{8}{348.209}{89.967}{generator}%
\NativeMathAt{399.904}{89.967}{11.955}{\equiv}%
\TextAt{9}{412.524}{89.967}{1}%
\TextAt{8}{421.697}{89.967}{mod}%
\TextAt{14}{447.135}{89.967}{m}%
\TextAt{8}{456.302}{89.967}{.}%
\TextAt{8}{464.635}{89.967}{But}%
\TextAt{8}{487.53}{89.967}{these}%
\TextAt{8}{517.156}{89.967}{are}%
\TextAt{8}{79.226}{72.921}{precisely}%
\TextAt{8}{127.083}{72.921}{the}%
\TextAt{8}{147.239}{72.921}{ideals}%
\TextAt{9}{179.817}{72.921}{\Character{40}}%
\TextAt{13}{184.369}{72.921}{a}%
\TextAt{9}{190.514}{72.921}{\Character{41}}%
\TextAt{8}{198.964}{72.921}{such}%
\TextAt{8}{225.361}{72.921}{that}%
\TextAt{13}{250.73}{72.921}{a}%
\NativeMathAt{260.196}{72.921}{11.955}{\equiv}%
\TextAt{9}{272.815}{72.921}{1}%
\TextAt{8}{281.989}{72.921}{mod}%
\TextAt{14}{307.427}{72.921}{p}%
\TextAt{21}{313.41}{77.26}{n}%
\TextAt{28}{318.548}{76.264}{p}%
\TextAt{8}{327.026}{72.921}{for}%
\TextAt{14}{344.911}{72.921}{p}%
\TextAt{16}{354.215}{72.921}{\Character{106}}%
\TextAt{14}{360.857}{72.921}{m}%
\NativeMathAt{370.023}{72.921}{11.955}{\infty}%
\TextAt{8}{381.979}{72.921}{,}%
\TextAt{8}{389.128}{72.921}{i.e.,}%
\TextAt{8}{411.233}{72.921}{the}%
\TextAt{8}{431.389}{72.921}{kernel}%
\TextAt{8}{465.856}{72.921}{of}%
\TextAt{13}{479.186}{72.921}{P}%
\TextAt{27}{488.357}{77.26}{m}%
\TextAt{25}{486.731}{69.857}{Q}%
\TextAt{8}{494.967}{72.921}{.}%
\Vector{q 0 G 0 g 0.399 w 452.334 -701.079 m 452.334 -693.009 l S Q}%
\Vector{q 0 G 0 g 0.399 w 452.534 -693.208 m 459.707 -693.208 l S Q}%
\Vector{q 0 G 0 g 0.399 w 452.534 -700.879 m 459.707 -700.879 l S Q}%
\Vector{q 0 G 0 g 0.399 w 459.906 -701.079 m 459.906 -693.009 l S Q}%
\TextAt{8}{297.211}{43.033}{333}%
\Destination{78.23}{809.815}{nav.page.333}%
\Destination{79.36}{288.355}{nav.print.366}%
\Destination{79.36}{681.355}{nav.number.6.1.9}%
\Destination{79.36}{210.555}{nav.number.6.1.10}%
}

% Page 346 (body page 334)
\Page{%
\TextAt{2}{79.226}{796.943}{VI.}%
\TextAt{2}{98.012}{796.943}{Global}%
\TextAt{2}{133.248}{796.943}{Class}%
\TextAt{2}{161.841}{796.943}{Field}%
\TextAt{2}{189.572}{796.943}{Theory}%
\Vector{q 0 G 0 g 0.299 w 7.226 18.729 m 460.769 18.729 l S Q}%
\TextAt{8}{97.159}{761.124}{The}%
\TextAt{8}{121.828}{761.124}{group}%
\TextAt{9}{155.619}{761.124}{\Character{40}}%
\TextAt{18}{160.171}{761.124}{Z}%
\TextAt{9}{168.141}{761.124}{\Character{47}}%
\TextAt{13}{173.994}{761.124}{m}%
\TextAt{18}{184.234}{761.124}{Z}%
\TextAt{9}{192.204}{761.124}{\Character{41}}%
\TextAt{3}{196.756}{765.462}{\Character{3}}%
\TextAt{8}{206.001}{761.124}{is}%
\TextAt{8}{218.38}{761.124}{canonically}%
\TextAt{8}{279.811}{761.124}{isomorphic}%
\TextAt{8}{339.676}{761.124}{to}%
\TextAt{8}{354.602}{761.124}{the}%
\TextAt{8}{375.374}{761.124}{Galois}%
\TextAt{8}{411.902}{761.124}{group}%
\TextAt{13}{445.693}{761.124}{G}%
\TextAt{9}{454.927}{761.124}{\Character{40}}%
\TextAt{18}{459.479}{761.124}{Q}%
\TextAt{9}{468.778}{761.124}{\Character{40}}%
\TextAt{13}{473.33}{761.124}{\Character{22}}%
\TextAt{21}{480.373}{759.33}{m}%
\TextAt{9}{488.362}{761.124}{\Character{41}}%
\TextAt{16}{497.283}{761.124}{\Character{106}}%
\TextAt{18}{504.973}{761.124}{Q}%
\TextAt{9}{514.272}{761.124}{\Character{41}}%
\TextAt{8}{523.337}{761.124}{of}%
\TextAt{8}{79.226}{744.078}{the}%
\TextAt{13}{99.383}{744.078}{m}%
\TextAt{8}{109.622}{744.078}{-th}%
\TextAt{8}{128.475}{744.078}{cyclotomic}%
\TextAt{8}{186.685}{744.078}{field}%
\TextAt{18}{212.042}{744.078}{Q}%
\TextAt{9}{221.34}{744.078}{\Character{40}}%
\TextAt{13}{225.893}{744.078}{\Character{22}}%
\TextAt{21}{232.936}{742.284}{m}%
\TextAt{9}{240.924}{744.078}{\Character{41}}%
\TextAt{8}{245.477}{744.078}{.}%
\TextAt{8}{253.925}{744.078}{We}%
\TextAt{8}{274.069}{744.078}{therefore}%
\TextAt{8}{323.169}{744.078}{obtain}%
\TextAt{8}{359.597}{744.078}{a}%
\TextAt{8}{369.352}{744.078}{canonical}%
\TextAt{8}{420.735}{744.078}{isomorphism}%
\TextAt{13}{242.501}{715.119}{G}%
\TextAt{9}{251.734}{715.119}{\Character{40}}%
\TextAt{18}{256.287}{715.119}{Q}%
\TextAt{9}{265.585}{715.119}{\Character{40}}%
\TextAt{13}{270.138}{715.119}{\Character{22}}%
\TextAt{21}{277.181}{713.326}{m}%
\TextAt{9}{285.169}{715.119}{\Character{41}}%
\TextAt{16}{293.042}{715.119}{\Character{106}}%
\TextAt{18}{299.684}{715.119}{Q}%
\TextAt{9}{308.983}{715.119}{\Character{41}}%
\NativeMathAt{316.856}{715.119}{11.955}{\cong}%
\TextAt{13}{329.475}{715.119}{C}%
\TextAt{25}{337.85}{713.326}{Q}%
\TextAt{9}{344.547}{715.119}{\Character{47}}%
\TextAt{13}{350.4}{715.119}{C}%
\TextAt{27}{359.634}{720.055}{m}%
\TextAt{25}{358.775}{712.164}{Q}%
\TextAt{13}{366.243}{715.119}{\Character{58}}%
\TextAt{8}{79.226}{686.161}{It}%
\TextAt{8}{91.487}{686.161}{is}%
\TextAt{8}{102.839}{686.161}{class}%
\TextAt{8}{129.864}{686.161}{field}%
\TextAt{8}{154.809}{686.161}{theory,}%
\TextAt{8}{193.501}{686.161}{which}%
\TextAt{8}{226.576}{686.161}{provides}%
\TextAt{8}{272.406}{686.161}{a}%
\TextAt{8}{281.75}{686.161}{far-reaching}%
\TextAt{8}{345.729}{686.161}{generalization}%
\TextAt{8}{420.12}{686.161}{of}%
\TextAt{8}{433.038}{686.161}{this}%
\TextAt{8}{455.449}{686.161}{important}%
\TextAt{8}{510.33}{686.161}{fact.}%
\TextAt{8}{79.226}{669.115}{For}%
\TextAt{8}{99.765}{669.115}{all}%
\TextAt{8}{115.593}{669.115}{modules}%
\TextAt{14}{161.07}{669.115}{m}%
\TextAt{8}{173.703}{669.115}{of}%
\TextAt{8}{186.602}{669.115}{an}%
\TextAt{8}{202.43}{669.115}{arbitrary}%
\TextAt{8}{251.768}{669.115}{number}%
\TextAt{8}{293.934}{669.115}{field}%
\TextAt{13}{318.86}{669.115}{K}%
\TextAt{8}{329.677}{669.115}{,}%
\TextAt{8}{336.482}{669.115}{there}%
\TextAt{8}{365.963}{669.115}{will}%
\TextAt{8}{387.637}{669.115}{be}%
\TextAt{8}{403.13}{669.115}{Galois}%
\TextAt{8}{438.613}{669.115}{extensions}%
\TextAt{13}{494.562}{669.115}{K}%
\TextAt{27}{505.38}{673.453}{m}%
\TextAt{16}{515.31}{669.115}{\Character{106}}%
\TextAt{13}{521.952}{669.115}{K}%
\TextAt{8}{79.226}{652.069}{generalizing}%
\TextAt{8}{142.648}{652.069}{the}%
\TextAt{8}{161.836}{652.069}{cyclotomic}%
\TextAt{8}{219.078}{652.069}{fields:}%
\TextAt{8}{253.116}{652.069}{the}%
\TextAt{8}{272.304}{652.069}{so-called}%
\TextAt{12}{318.869}{652.069}{ray}%
\TextAt{12}{341.008}{652.069}{class}%
\TextAt{12}{371.253}{652.069}{fields}%
\TextAt{8}{401.38}{652.069}{,}%
\TextAt{8}{407.754}{652.069}{which}%
\TextAt{8}{440.272}{652.069}{satisfy}%
\TextAt{8}{475.851}{652.069}{canonically}%
\TextAt{13}{250.251}{623.11}{G}%
\TextAt{9}{259.484}{623.11}{\Character{40}}%
\TextAt{13}{264.037}{623.11}{K}%
\TextAt{27}{274.854}{628.046}{m}%
\TextAt{16}{284.785}{623.11}{\Character{106}}%
\TextAt{13}{291.426}{623.11}{K}%
\TextAt{9}{302.244}{623.11}{\Character{41}}%
\NativeMathAt{310.117}{623.11}{11.955}{\cong}%
\TextAt{13}{322.736}{623.11}{C}%
\TextAt{21}{331.111}{621.317}{K}%
\TextAt{9}{339.314}{623.11}{\Character{47}}%
\TextAt{13}{345.167}{623.11}{C}%
\TextAt{27}{354.401}{628.046}{m}%
\TextAt{21}{353.542}{620.155}{K}%
\TextAt{8}{79.226}{594.152}{(see}%
\TextAt{8}{103.627}{594.152}{§6).}%
\TextAt{8}{130.964}{594.152}{The}%
\TextAt{8}{155.951}{594.152}{ray}%
\TextAt{8}{177.052}{594.152}{class}%
\TextAt{8}{205.421}{594.152}{group}%
\TextAt{8}{239.529}{594.152}{mod}%
\TextAt{9}{266.799}{594.152}{1}%
\TextAt{8}{277.482}{594.152}{is}%
\TextAt{8}{290.179}{594.152}{of}%
\TextAt{8}{304.441}{594.152}{particular}%
\TextAt{8}{359.363}{594.152}{interest}%
\TextAt{8}{402.306}{594.152}{here.}%
\TextAt{8}{435.012}{594.152}{It}%
\TextAt{8}{448.617}{594.152}{is}%
\TextAt{8}{461.313}{594.152}{related}%
\TextAt{8}{501.268}{594.152}{to}%
\TextAt{8}{516.51}{594.152}{the}%
\TextAt{8}{79.226}{577.106}{ideal}%
\TextAt{8}{106.54}{577.106}{class}%
\TextAt{8}{133.327}{577.106}{group}%
\TextAt{13}{165.853}{577.106}{C}%
\TextAt{13}{175.071}{577.106}{l}%
\TextAt{21}{178.574}{575.312}{K}%
\TextAt{8}{186.776}{577.106}{—which}%
\TextAt{8}{231.317}{577.106}{according}%
\TextAt{8}{283.354}{577.106}{to}%
\TextAt{8}{297.015}{577.106}{our}%
\TextAt{8}{317.179}{577.106}{definition}%
\TextAt{8}{368.559}{577.106}{here,}%
\TextAt{8}{396.648}{577.106}{is}%
\TextAt{8}{407.762}{577.106}{in}%
\TextAt{8}{420.765}{577.106}{general}%
\TextAt{8}{460.441}{577.106}{not}%
\TextAt{8}{480.605}{577.106}{a}%
\TextAt{8}{489.711}{577.106}{ray}%
\TextAt{8}{509.23}{577.106}{class}%
\TextAt{8}{79.226}{560.06}{group—as}%
\TextAt{8}{134.579}{560.06}{follows.}%
\TextAt{12}{79.226}{534.06}{Proposition}%
\TextAt{12}{151.937}{534.06}{(1.11).}%
\TextAt{8}{193.984}{534.06}{There}%
\TextAt{8}{227.793}{534.06}{is}%
\TextAt{8}{239.557}{534.06}{an}%
\TextAt{8}{255.816}{534.06}{exact}%
\TextAt{8}{286.708}{534.06}{sequence}%
\TextAt{9}{165.627}{505.101}{1}%
\NativeArrowAt{174.801}{505.101}{11.955}{19.26}{\longrightarrow}%
\TextAt{16}{197.372}{505.101}{\Character{79}}%
\TextAt{3}{207.233}{510.037}{\Character{3}}%
\TextAt{9}{211.965}{505.101}{\Character{47}}%
\TextAt{16}{217.818}{505.101}{\Character{79}}%
\TextAt{3}{227.668}{510.037}{\Character{3}}%
\TextAt{17}{227.336}{502.146}{\Character{43}}%
\NativeArrowAt{237.741}{505.101}{11.955}{19.26}{\longrightarrow}%
\TextAt{20}{262.913}{516.458}{\Character{89}}%
\TextAt{27}{260.323}{491.002}{p}%
\TextAt{7}{267.133}{491.002}{real}%
\TextAt{18}{282.771}{505.101}{R}%
\TextAt{3}{291.405}{510.037}{\Character{3}}%
\TextAt{9}{296.137}{505.101}{\Character{47}}%
\TextAt{18}{301.99}{505.101}{R}%
\TextAt{3}{310.624}{510.037}{\Character{3}}%
\TextAt{17}{310.624}{502.146}{\Character{43}}%
\NativeArrowAt{321.03}{505.101}{11.955}{19.26}{\longrightarrow}%
\TextAt{13}{343.612}{505.101}{C}%
\TextAt{13}{352.829}{505.101}{l}%
\TextAt{17}{356.595}{510.037}{1}%
\TextAt{21}{356.332}{502.146}{K}%
\NativeArrowAt{367.856}{505.101}{11.955}{19.26}{\longrightarrow}%
\TextAt{13}{390.438}{505.101}{C}%
\TextAt{13}{399.656}{505.101}{l}%
\TextAt{21}{403.158}{503.308}{K}%
\NativeArrowAt{414.682}{505.101}{11.955}{19.26}{\longrightarrow}%
\TextAt{9}{437.264}{505.101}{1}%
\TextAt{13}{443.117}{505.101}{\Character{59}}%
\TextAt{8}{79.226}{466.974}{where}%
\TextAt{16}{113.035}{466.974}{\Character{79}}%
\TextAt{3}{122.885}{471.313}{\Character{3}}%
\TextAt{17}{122.553}{464.019}{\Character{43}}%
\TextAt{8}{133.535}{466.974}{is}%
\TextAt{8}{145.299}{466.974}{the}%
\TextAt{8}{165.455}{466.974}{group}%
\TextAt{8}{198.631}{466.974}{of}%
\TextAt{8}{211.961}{466.974}{totally}%
\TextAt{8}{249.369}{466.974}{positive}%
\TextAt{8}{292.682}{466.974}{units}%
\TextAt{8}{322.008}{466.974}{of}%
\TextAt{13}{335.338}{466.974}{K}%
\TextAt{8}{346.156}{466.974}{.}%
\TextAt{19}{79.226}{432.02}{Proof.}%
\TextAt{8}{116}{432.02}{One}%
\TextAt{8}{142.146}{432.02}{has}%
\TextAt{13}{164.454}{432.02}{C}%
\TextAt{13}{173.672}{432.02}{l}%
\TextAt{17}{177.438}{436.358}{1}%
\TextAt{21}{177.175}{428.999}{K}%
\NativeMathAt{191.143}{432.02}{11.955}{\cong}%
\TextAt{13}{206.207}{432.02}{C}%
\TextAt{21}{214.582}{430.227}{K}%
\TextAt{9}{222.785}{432.02}{\Character{47}}%
\TextAt{13}{228.638}{432.02}{C}%
\TextAt{17}{237.872}{436.358}{1}%
\TextAt{21}{237.013}{428.999}{K}%
\TextAt{9}{250.981}{432.02}{\Character{61}}%
\TextAt{13}{265.851}{432.02}{I}%
\TextAt{21}{271.014}{430.227}{K}%
\TextAt{9}{279.217}{432.02}{\Character{47}}%
\TextAt{13}{285.07}{432.02}{I}%
\TextAt{17}{291.173}{436.358}{1}%
\TextAt{21}{290.233}{428.999}{K}%
\TextAt{13}{298.435}{432.02}{K}%
\TextAt{3}{309.253}{436.358}{\Character{3}}%
\TextAt{8}{319.317}{432.02}{and,}%
\TextAt{8}{347.125}{432.02}{by}%
\TextAt{8}{364.819}{432.02}{(1.3),}%
\TextAt{13}{397.839}{432.02}{C}%
\TextAt{13}{407.057}{432.02}{l}%
\TextAt{21}{410.559}{430.227}{K}%
\NativeMathAt{424.528}{432.02}{11.955}{\cong}%
\TextAt{13}{439.592}{432.02}{I}%
\TextAt{21}{444.754}{430.227}{K}%
\TextAt{9}{452.957}{432.02}{\Character{47}}%
\TextAt{13}{458.81}{432.02}{I}%
\TextAt{21}{464.913}{437.134}{S}%
\NativeMathAt{470.063}{436.138}{5.978}{\infty}%
\TextAt{21}{463.973}{428.779}{K}%
\TextAt{13}{478.365}{432.02}{K}%
\TextAt{3}{489.183}{436.358}{\Character{3}}%
\TextAt{8}{493.915}{432.02}{,}%
\TextAt{8}{502.858}{432.02}{where}%
\TextAt{13}{79.226}{414.974}{I}%
\TextAt{17}{85.329}{419.312}{1}%
\TextAt{21}{84.389}{411.953}{K}%
\TextAt{9}{95.913}{414.974}{\Character{61}}%
\TextAt{20}{108.338}{423.94}{\Character{81}}%
\TextAt{27}{119.629}{411.487}{p}%
\TextAt{13}{126.108}{414.974}{U}%
\TextAt{27}{134.059}{413.181}{p}%
\TextAt{8}{142.443}{414.974}{and}%
\TextAt{13}{165.206}{414.974}{I}%
\TextAt{21}{171.308}{420.088}{S}%
\NativeMathAt{176.459}{419.092}{5.978}{\infty}%
\TextAt{21}{170.368}{411.733}{K}%
\TextAt{9}{188.082}{414.974}{\Character{61}}%
\TextAt{20}{200.507}{423.94}{\Character{81}}%
\TextAt{27}{211.798}{411.487}{p}%
\TextAt{25}{215.787}{411.487}{\Character{45}}%
\NativeMathAt{218.001}{411.487}{7.97}{\infty}%
\TextAt{13}{228.96}{414.974}{U}%
\TextAt{27}{236.91}{413.181}{p}%
\TextAt{16}{244.053}{414.974}{\Character{2}}%
\TextAt{20}{256.009}{423.94}{\Character{81}}%
\TextAt{27}{267.3}{411.487}{p}%
\TextAt{3}{271.288}{411.487}{\Character{106}}%
\NativeMathAt{273.639}{411.487}{7.97}{\infty}%
\TextAt{13}{284.6}{414.974}{K}%
\TextAt{3}{295.417}{419.312}{\Character{3}}%
\TextAt{27}{294.559}{412.018}{p}%
\TextAt{8}{300.15}{414.974}{.}%
\TextAt{8}{308.598}{414.974}{We}%
\TextAt{8}{328.742}{414.974}{therefore}%
\TextAt{8}{377.842}{414.974}{obtain}%
\TextAt{8}{414.27}{414.974}{an}%
\TextAt{8}{430.529}{414.974}{exact}%
\TextAt{8}{461.421}{414.974}{sequence}%
\TextAt{9}{181.811}{386.015}{1}%
\NativeArrowAt{190.985}{386.015}{11.955}{19.26}{\longrightarrow}%
\TextAt{13}{213.567}{386.015}{I}%
\TextAt{21}{219.669}{391.13}{S}%
\NativeMathAt{224.82}{390.133}{5.978}{\infty}%
\TextAt{21}{218.729}{382.774}{K}%
\TextAt{13}{233.122}{386.015}{K}%
\TextAt{3}{243.939}{390.952}{\Character{3}}%
\TextAt{9}{248.672}{386.015}{\Character{47}}%
\TextAt{13}{254.525}{386.015}{I}%
\TextAt{17}{260.627}{390.952}{1}%
\TextAt{21}{259.688}{383.06}{K}%
\TextAt{13}{267.89}{386.015}{K}%
\TextAt{3}{278.708}{390.952}{\Character{3}}%
\NativeArrowAt{286.761}{386.015}{11.955}{19.26}{\longrightarrow}%
\TextAt{13}{309.343}{386.015}{C}%
\TextAt{21}{317.718}{384.222}{K}%
\TextAt{9}{325.921}{386.015}{\Character{47}}%
\TextAt{13}{331.774}{386.015}{C}%
\TextAt{17}{341.007}{390.952}{1}%
\TextAt{21}{340.149}{383.06}{K}%
\NativeArrowAt{351.672}{386.015}{11.955}{19.26}{\longrightarrow}%
\TextAt{13}{374.254}{386.015}{C}%
\TextAt{13}{383.471}{386.015}{l}%
\TextAt{21}{386.975}{384.222}{K}%
\NativeArrowAt{398.498}{386.015}{11.955}{19.26}{\longrightarrow}%
\TextAt{9}{421.08}{386.015}{1}%
\TextAt{13}{426.933}{386.015}{\Character{58}}%
\TextAt{8}{79.226}{357.057}{For}%
\TextAt{8}{100.196}{357.057}{the}%
\TextAt{8}{120.352}{357.057}{group}%
\TextAt{8}{153.528}{357.057}{on}%
\TextAt{8}{169.787}{357.057}{the}%
\TextAt{8}{189.943}{357.057}{left}%
\TextAt{8}{210.422}{357.057}{we}%
\TextAt{8}{227.65}{357.057}{have}%
\TextAt{8}{254.644}{357.057}{the}%
\TextAt{8}{274.801}{357.057}{exact}%
\TextAt{8}{305.693}{357.057}{sequence}%
\TextAt{9}{167}{323.049}{1}%
\NativeArrowAt{176.174}{323.049}{11.955}{19.26}{\longrightarrow}%
\TextAt{13}{199.951}{331.137}{I}%
\TextAt{21}{206.054}{336.251}{S}%
\NativeMathAt{211.204}{335.255}{5.978}{\infty}%
\TextAt{21}{205.114}{327.896}{K}%
\TextAt{16}{222.163}{331.137}{\Character{92}}%
\TextAt{13}{232.79}{331.137}{K}%
\TextAt{3}{243.607}{335.475}{\Character{3}}%
\Vector{q 0 G 0 g 0.478 w 127.951 -447.962 m 176.34 -447.962 l S Q}%
\TextAt{13}{203.046}{314.848}{I}%
\TextAt{17}{209.149}{318.966}{1}%
\TextAt{21}{208.209}{311.607}{K}%
\TextAt{16}{219.068}{314.848}{\Character{92}}%
\TextAt{13}{229.695}{314.848}{K}%
\TextAt{3}{240.513}{318.302}{\Character{3}}%
\NativeArrowAt{252.856}{323.049}{11.955}{19.26}{\longrightarrow}%
\TextAt{13}{275.438}{323.049}{I}%
\TextAt{21}{281.541}{328.163}{S}%
\NativeMathAt{286.691}{327.167}{5.978}{\infty}%
\TextAt{21}{280.601}{319.808}{K}%
\TextAt{9}{294.993}{323.049}{\Character{47}}%
\TextAt{13}{300.846}{323.049}{I}%
\TextAt{17}{306.949}{327.985}{1}%
\TextAt{21}{306.009}{320.094}{K}%
\NativeArrowAt{317.533}{323.049}{11.955}{19.26}{\longrightarrow}%
\TextAt{13}{340.115}{323.049}{I}%
\TextAt{21}{346.217}{328.163}{S}%
\NativeMathAt{351.368}{327.167}{5.978}{\infty}%
\TextAt{21}{345.278}{319.808}{K}%
\TextAt{13}{359.67}{323.049}{K}%
\TextAt{3}{370.487}{327.985}{\Character{3}}%
\TextAt{9}{375.22}{323.049}{\Character{47}}%
\TextAt{13}{381.073}{323.049}{I}%
\TextAt{17}{387.176}{327.985}{1}%
\TextAt{21}{386.236}{320.094}{K}%
\TextAt{13}{394.438}{323.049}{K}%
\TextAt{3}{405.256}{327.985}{\Character{3}}%
\NativeArrowAt{413.309}{323.049}{11.955}{19.26}{\longrightarrow}%
\TextAt{9}{435.891}{323.049}{1}%
\TextAt{13}{441.744}{323.049}{\Character{58}}%
\TextAt{8}{79.226}{288.138}{But}%
\TextAt{13}{102.467}{288.138}{I}%
\TextAt{21}{108.57}{293.252}{S}%
\NativeMathAt{113.72}{292.256}{5.978}{\infty}%
\TextAt{21}{107.63}{284.897}{K}%
\TextAt{16}{124.679}{288.138}{\Character{92}}%
\TextAt{13}{135.306}{288.138}{K}%
\TextAt{3}{146.123}{292.477}{\Character{3}}%
\TextAt{9}{154.176}{288.138}{\Character{61}}%
\TextAt{16}{166.602}{288.138}{\Character{79}}%
\TextAt{3}{176.452}{292.477}{\Character{3}}%
\TextAt{8}{181.184}{288.138}{,}%
\TextAt{13}{188.333}{288.138}{I}%
\TextAt{17}{194.436}{292.477}{1}%
\TextAt{21}{193.496}{285.118}{K}%
\TextAt{16}{204.355}{288.138}{\Character{92}}%
\TextAt{13}{214.982}{288.138}{K}%
\TextAt{3}{225.8}{292.477}{\Character{3}}%
\TextAt{9}{233.853}{288.138}{\Character{61}}%
\TextAt{16}{246.279}{288.138}{\Character{79}}%
\TextAt{3}{256.128}{292.477}{\Character{3}}%
\TextAt{17}{255.796}{285.183}{\Character{43}}%
\TextAt{8}{262.881}{288.138}{,}%
\TextAt{8}{270.03}{288.138}{and}%
\TextAt{13}{292.793}{288.138}{I}%
\TextAt{21}{298.895}{293.252}{S}%
\NativeMathAt{304.046}{292.256}{5.978}{\infty}%
\TextAt{21}{297.956}{284.897}{K}%
\TextAt{9}{312.348}{288.138}{\Character{47}}%
\TextAt{13}{318.201}{288.138}{I}%
\TextAt{17}{324.304}{292.477}{1}%
\TextAt{21}{323.364}{285.118}{K}%
\TextAt{9}{334.887}{288.138}{\Character{61}}%
\TextAt{20}{347.313}{297.105}{\Character{81}}%
\TextAt{27}{358.604}{284.651}{p}%
\TextAt{3}{362.592}{284.651}{\Character{106}}%
\NativeMathAt{364.943}{284.651}{7.97}{\infty}%
\TextAt{13}{375.904}{288.138}{K}%
\TextAt{3}{386.721}{292.477}{\Character{3}}%
\TextAt{27}{385.863}{285.183}{p}%
\TextAt{9}{391.454}{288.138}{\Character{47}}%
\TextAt{13}{397.307}{288.138}{U}%
\TextAt{27}{405.257}{286.345}{p}%
\TextAt{9}{413.065}{288.138}{\Character{61}}%
\TextAt{20}{425.49}{297.105}{\Character{81}}%
\TextAt{27}{436.781}{284.651}{p}%
\TextAt{7}{443.591}{284.651}{real}%
\TextAt{18}{459.727}{288.138}{R}%
\TextAt{3}{468.361}{292.477}{\Character{3}}%
\TextAt{9}{473.093}{288.138}{\Character{47}}%
\TextAt{18}{478.946}{288.138}{R}%
\TextAt{3}{487.581}{292.477}{\Character{3}}%
\TextAt{17}{487.581}{285.183}{\Character{43}}%
\TextAt{8}{494.665}{288.138}{.}%
\Vector{q 0 G 0 g 0.399 w 452.334 -485.862 m 452.334 -477.792 l S Q}%
\Vector{q 0 G 0 g 0.399 w 452.534 -477.991 m 459.707 -477.991 l S Q}%
\Vector{q 0 G 0 g 0.399 w 452.534 -485.663 m 459.707 -485.663 l S Q}%
\Vector{q 0 G 0 g 0.399 w 459.906 -485.862 m 459.906 -477.792 l S Q}%
\ErrataAt{79.22}{252.055}{453.64}{\textbf{Exercise VI.1.1.} (i) $A_{\mathbb Q}=(\widehat{\mathbb Z}\otimes_{\mathbb Z}\mathbb Q)\times\mathbb R$. (ii) The quotient $A_{\mathbb Q}/\mathbb Q$ is}%
\ErrataAt{79.22}{234.955}{453.64}{compact and connected. (iii) $A_{\mathbb Q}/\mathbb Q$ is uniquely divisible: $nx=y$ has a unique}%
\ErrataAt{79.22}{217.955}{453.64}{solution for each $n\ge1$ and $y\in A_{\mathbb Q}/\mathbb Q$.}%
\TextAt{12}{79.226}{185.974}{Exercise}%
\TextAt{12}{133.26}{185.974}{VI.1.2.}%
\TextAt{8}{177.74}{185.974}{Let}%
\TextAt{13}{199.434}{185.974}{K}%
\TextAt{8}{214.874}{185.974}{be}%
\TextAt{8}{231.523}{185.974}{a}%
\TextAt{8}{242.003}{185.974}{number}%
\TextAt{8}{285.323}{185.974}{field,}%
\TextAt{13}{314.838}{185.974}{m}%
\TextAt{9}{329.632}{185.974}{\Character{61}}%
\TextAt{9}{343.285}{185.974}{2}%
\TextAt{21}{349.146}{190.312}{\Character{23}}%
\TextAt{13}{354.328}{185.974}{m}%
\TextAt{3}{364.567}{190.312}{\Character{48}}%
\TextAt{8}{371.984}{185.974}{(}%
\TextAt{13}{376.539}{185.974}{m}%
\TextAt{3}{386.778}{190.312}{\Character{48}}%
\TextAt{8}{394.195}{185.974}{odd),}%
\TextAt{8}{425.993}{185.974}{and}%
\TextAt{8}{449.48}{185.974}{let}%
\TextAt{13}{467.109}{185.974}{S}%
\TextAt{8}{479.626}{185.974}{be}%
\TextAt{8}{496.275}{185.974}{a}%
\TextAt{8}{506.755}{185.974}{finite}%
\TextAt{8}{79.226}{168.928}{set}%
\TextAt{8}{97.726}{168.928}{of}%
\TextAt{8}{111.289}{168.928}{primes.}%
\TextAt{8}{154.317}{168.928}{Let}%
\TextAt{13}{175.519}{168.928}{a}%
\NativeMathAt{185.382}{168.928}{11.955}{\in}%
\TextAt{13}{197.07}{168.928}{K}%
\TextAt{3}{207.887}{173.266}{\Character{3}}%
\TextAt{8}{216.75}{168.928}{and}%
\TextAt{13}{239.745}{168.928}{a}%
\NativeMathAt{249.608}{168.928}{11.955}{\in}%
\TextAt{13}{261.296}{168.928}{K}%
\TextAt{3}{272.113}{173.266}{\Character{3}}%
\TextAt{21}{276.347}{173.266}{m}%
\TextAt{27}{271.255}{165.972}{p}%
\TextAt{8}{284.336}{168.928}{,}%
\TextAt{8}{291.776}{168.928}{for}%
\TextAt{8}{309.894}{168.928}{all}%
\TextAt{14}{326.386}{168.928}{p}%
\NativeMathAt{336.087}{168.928}{11.955}{\notin}%
\TextAt{13}{347.774}{168.928}{S}%
\TextAt{8}{355.67}{168.928}{.}%
\TextAt{8}{364.817}{168.928}{Show:}%
\TextAt{8}{400.725}{168.928}{(i)}%
\TextAt{8}{417.217}{168.928}{If}%
\TextAt{13}{429.142}{168.928}{K}%
\TextAt{9}{439.96}{168.928}{\Character{40}}%
\TextAt{13}{444.512}{168.928}{\Character{16}}%
\TextAt{17}{449.625}{167.134}{2}%
\TextAt{22}{453.859}{169.404}{\Character{23}}%
\TextAt{9}{458.979}{168.928}{\Character{41}}%
\TextAt{16}{467.249}{168.928}{\Character{106}}%
\TextAt{13}{474.287}{168.928}{K}%
\TextAt{8}{489.235}{168.928}{is}%
\TextAt{8}{501.232}{168.928}{cyclic,}%
\TextAt{8}{79.226}{151.882}{where}%
\TextAt{13}{112.691}{151.882}{\Character{16}}%
\TextAt{17}{117.803}{150.088}{2}%
\TextAt{22}{122.038}{152.358}{\Character{23}}%
\TextAt{8}{130.71}{151.882}{is}%
\TextAt{8}{142.129}{151.882}{a}%
\TextAt{8}{151.54}{151.882}{primitive}%
\TextAt{9}{201.276}{151.882}{2}%
\TextAt{21}{207.129}{156.22}{\Character{23}}%
\TextAt{8}{212.311}{151.882}{-th}%
\TextAt{8}{230.819}{151.882}{root}%
\TextAt{8}{255.521}{151.882}{of}%
\TextAt{8}{268.506}{151.882}{unity,}%
\TextAt{8}{301.072}{151.882}{then}%
\TextAt{13}{327.387}{151.882}{a}%
\NativeMathAt{336.853}{151.882}{11.955}{\in}%
\TextAt{13}{348.144}{151.882}{K}%
\TextAt{3}{358.961}{156.22}{\Character{3}}%
\TextAt{21}{363.196}{156.22}{m}%
\TextAt{8}{371.184}{151.882}{.}%
\TextAt{8}{379.518}{151.882}{(ii)}%
\TextAt{8}{398.684}{151.882}{Otherwise}%
\TextAt{8}{453.68}{151.882}{one}%
\TextAt{8}{474.795}{151.882}{has}%
\TextAt{8}{495.324}{151.882}{at}%
\TextAt{8}{509.29}{151.882}{least}%
\TextAt{8}{79.226}{134.836}{that}%
\TextAt{13}{104.595}{134.836}{a}%
\NativeMathAt{114.061}{134.836}{11.955}{\in}%
\TextAt{13}{125.352}{134.836}{K}%
\TextAt{3}{136.169}{139.174}{\Character{3}}%
\TextAt{21}{140.404}{139.174}{m}%
\TextAt{17}{147.894}{139.174}{\Character{47}2}%
\TextAt{8}{156.861}{134.836}{.}%
\TextAt{8}{297.211}{43.033}{334}%
\Destination{78.23}{809.815}{nav.page.334}%
\Destination{79.226}{266.055}{nav.ex.6.1.1}%
\Destination{79.226}{199.955}{nav.ex.6.1.2}%
\Destination{79.36}{372.055}{nav.print.367}%
\Destination{79.36}{548.055}{nav.number.6.1.11}%
\Hotlink{108.806}{591.955}{6.858}{11.047}{nav.section.6.6}%
\Hotlink{364.319}{429.855}{25.078}{11.047}{nav.number.6.1.3}%
\Destination{79.226}{266.055}{errata.27}%
}

% Page 347 (body page 335)
\Page{%
\TextAt{2}{79.226}{796.943}{VI.}%
\TextAt{2}{98.012}{796.943}{Global}%
\TextAt{2}{133.248}{796.943}{Class}%
\TextAt{2}{161.841}{796.943}{Field}%
\TextAt{2}{189.572}{796.943}{Theory}%
\Vector{q 0 G 0 g 0.299 w 7.226 18.729 m 460.769 18.729 l S Q}%
\TextAt{12}{97.159}{761.124}{Hint:}%
\TextAt{8}{133.435}{761.124}{Use}%
\TextAt{8}{156.292}{761.124}{the}%
\TextAt{8}{176.817}{761.124}{following}%
\TextAt{8}{226.62}{761.124}{fact,}%
\TextAt{8}{253.419}{761.124}{proved}%
\TextAt{8}{291.841}{761.124}{in}%
\TextAt{8}{305.862}{761.124}{(3.8):}%
\TextAt{8}{339.126}{761.124}{if}%
\TextAt{13}{350.219}{761.124}{L}%
\TextAt{16}{362.132}{761.124}{\Character{106}}%
\TextAt{13}{369.403}{761.124}{K}%
\TextAt{8}{384.487}{761.124}{is}%
\TextAt{8}{396.619}{761.124}{a}%
\TextAt{8}{406.743}{761.124}{finite}%
\TextAt{8}{437.024}{761.124}{extension}%
\TextAt{8}{489.159}{761.124}{in}%
\TextAt{8}{503.18}{761.124}{which}%
\TextAt{8}{79.226}{744.078}{almost}%
\TextAt{8}{117.016}{744.078}{all}%
\TextAt{8}{133.275}{744.078}{prime}%
\TextAt{8}{166.439}{744.078}{ideals}%
\TextAt{8}{199.017}{744.078}{split}%
\TextAt{8}{225.091}{744.078}{completely,}%
\TextAt{8}{286.218}{744.078}{then}%
\TextAt{13}{312.878}{744.078}{L}%
\TextAt{9}{324.163}{744.078}{\Character{61}}%
\TextAt{13}{336.589}{744.078}{K}%
\TextAt{8}{347.406}{744.078}{.}%
\TextAt{12}{79.226}{712.088}{Exercise}%
\TextAt{12}{133.215}{712.088}{VI.1.3.}%
\TextAt{8}{177.657}{712.088}{Write}%
\TextAt{13}{210.849}{712.088}{I}%
\TextAt{17}{216.951}{716.426}{1}%
\TextAt{21}{216.012}{709.067}{K}%
\TextAt{9}{228.703}{712.088}{\Character{61}}%
\TextAt{13}{242.297}{712.088}{I}%
\TextAt{17}{248.4}{716.426}{1}%
\TextAt{21}{247.46}{708.979}{f}%
\TextAt{16}{256.256}{712.088}{\Character{2}}%
\TextAt{13}{268.679}{712.088}{I}%
\TextAt{17}{274.782}{716.426}{1}%
\NativeMathAt{273.842}{709.132}{7.97}{\infty}%
\TextAt{8}{282.808}{712.088}{,}%
\TextAt{8}{290.815}{712.088}{with}%
\TextAt{13}{318.16}{712.088}{I}%
\TextAt{17}{324.263}{716.426}{1}%
\TextAt{21}{323.323}{708.979}{f}%
\TextAt{9}{333.485}{712.088}{\Character{61}}%
\TextAt{20}{347.078}{721.054}{\Character{81}}%
\TextAt{27}{358.369}{708.601}{p}%
\TextAt{25}{362.358}{708.601}{\Character{45}}%
\NativeMathAt{364.572}{708.601}{7.97}{\infty}%
\TextAt{13}{375.531}{712.088}{U}%
\TextAt{27}{383.481}{710.295}{p}%
\TextAt{8}{387.968}{712.088}{,}%
\TextAt{13}{395.974}{712.088}{I}%
\TextAt{17}{402.077}{716.426}{1}%
\NativeMathAt{401.137}{709.132}{7.97}{\infty}%
\TextAt{9}{414.593}{712.088}{\Character{61}}%
\TextAt{20}{428.187}{721.054}{\Character{81}}%
\TextAt{27}{439.478}{708.601}{p}%
\TextAt{3}{443.467}{708.601}{\Character{106}}%
\NativeMathAt{445.818}{708.601}{7.97}{\infty}%
\TextAt{13}{456.778}{712.088}{U}%
\TextAt{27}{464.728}{710.295}{p}%
\TextAt{8}{469.215}{712.088}{.}%
\TextAt{8}{479.72}{712.088}{Show}%
\TextAt{8}{511.298}{712.088}{that}%
\TextAt{8}{79.226}{695.042}{taking}%
\TextAt{8}{115.994}{695.042}{integer}%
\TextAt{8}{155.356}{695.042}{powers}%
\TextAt{8}{194.778}{695.042}{of}%
\TextAt{8}{208.771}{695.042}{idèles}%
\TextAt{13}{241.355}{695.042}{\Character{11}}%
\NativeMathAt{253.327}{695.042}{11.955}{\in}%
\TextAt{13}{265.749}{695.042}{I}%
\TextAt{17}{271.851}{699.38}{1}%
\TextAt{21}{270.911}{691.933}{f}%
\TextAt{8}{281.144}{695.042}{extends}%
\TextAt{8}{324.463}{695.042}{by}%
\TextAt{8}{341.386}{695.042}{continuity}%
\TextAt{8}{397.342}{695.042}{to}%
\TextAt{8}{412.315}{695.042}{exponentiation}%
\TextAt{13}{492.66}{695.042}{\Character{11}}%
\TextAt{21}{500.181}{699.38}{x}%
\TextAt{8}{510.007}{695.042}{with}%
\TextAt{13}{79.226}{677.996}{x}%
\NativeMathAt{89.199}{677.996}{11.955}{\in}%
\TextAt{20}{101.154}{681.084}{\Character{98}}%
\TextAt{18}{100.49}{677.996}{Z}%
\TextAt{8}{108.46}{677.996}{.}%
\TextAt{12}{79.226}{646.006}{Exercise}%
\TextAt{12}{132.528}{646.006}{VI.1.4.}%
\TextAt{8}{176.371}{646.006}{Let}%
\TextAt{13}{197.429}{646.006}{\Character{34}}%
\TextAt{17}{202.907}{644.213}{1}%
\TextAt{13}{207.639}{646.006}{\Character{59}}%
\TextAt{13}{212.863}{646.006}{\Character{58}}%
\TextAt{13}{218.1}{646.006}{\Character{58}}%
\TextAt{13}{223.324}{646.006}{\Character{58}}%
\TextAt{13}{228.561}{646.006}{\Character{59}}%
\TextAt{13}{233.797}{646.006}{\Character{34}}%
\TextAt{21}{239.338}{644.213}{t}%
\NativeMathAt{246.364}{646.006}{11.955}{\in}%
\TextAt{16}{257.793}{646.006}{\Character{79}}%
\TextAt{3}{267.655}{650.344}{\Character{3}}%
\TextAt{21}{267.323}{642.985}{K}%
\TextAt{8}{279.51}{646.006}{be}%
\TextAt{8}{295.523}{646.006}{independent}%
\TextAt{8}{361.938}{646.006}{units.}%
\TextAt{8}{396.078}{646.006}{The}%
\TextAt{8}{420.22}{646.006}{images}%
\Vector{q 0 G 0 g 0.478 w 386.743 -121.173 m 392.221 -121.173 l S Q}%
\TextAt{13}{458.743}{646.006}{\Character{34}}%
\TextAt{17}{464.221}{644.213}{1}%
\TextAt{13}{468.954}{646.006}{\Character{59}}%
\TextAt{13}{474.178}{646.006}{\Character{58}}%
\TextAt{13}{479.415}{646.006}{\Character{58}}%
\TextAt{13}{484.639}{646.006}{\Character{58}}%
\TextAt{13}{489.876}{646.006}{\Character{59}}%
\Vector{q 0 G 0 g 0.478 w 423.174 -121.173 m 428.652 -121.173 l S Q}%
\TextAt{13}{495.174}{646.006}{\Character{34}}%
\TextAt{21}{500.652}{644.213}{t}%
\TextAt{8}{508.194}{646.006}{in}%
\TextAt{13}{521.934}{646.006}{I}%
\TextAt{17}{528.037}{650.344}{1}%
\TextAt{21}{527.097}{642.897}{f}%
\TextAt{8}{79.226}{628.96}{are}%
\TextAt{8}{99.468}{628.96}{then}%
\TextAt{8}{126.859}{628.96}{independent}%
\TextAt{8}{193.918}{628.96}{units}%
\TextAt{8}{223.975}{628.96}{with}%
\TextAt{8}{251.366}{628.96}{respect}%
\TextAt{8}{292.147}{628.96}{to}%
\TextAt{8}{307.189}{628.96}{the}%
\TextAt{8}{328.076}{628.96}{exponentiation}%
\TextAt{8}{408.488}{628.96}{with}%
\TextAt{8}{435.88}{628.96}{elements}%
\TextAt{8}{484.467}{628.96}{of}%
\TextAt{20}{499.193}{632.048}{\Character{98}}%
\TextAt{18}{498.528}{628.96}{Z}%
\TextAt{8}{506.499}{628.96}{,}%
\TextAt{8}{514.562}{628.96}{i.e.,}%
\TextAt{8}{79.226}{611.914}{any}%
\TextAt{8}{101.343}{611.914}{relation}%
\Vector{q 0 G 0 g 0.478 w 167.228 -172.311 m 172.706 -172.311 l S Q}%
\TextAt{13}{239.228}{594.868}{\Character{34}}%
\TextAt{21}{244.706}{600.093}{x}%
\TextAt{24}{249.473}{598.986}{1}%
\TextAt{17}{244.706}{591.937}{1}%
\TextAt{16}{256.115}{594.868}{\Character{1}}%
\TextAt{16}{261.411}{594.868}{\Character{1}}%
\TextAt{16}{266.719}{594.868}{\Character{1}}%
\Vector{q 0 G 0 g 0.478 w 200.055 -172.311 m 205.533 -172.311 l S Q}%
\TextAt{13}{272.055}{594.868}{\Character{34}}%
\TextAt{21}{277.533}{599.804}{x}%
\TextAt{22}{282.3}{598.808}{t}%
\TextAt{21}{277.533}{591.912}{t}%
\TextAt{9}{289.467}{594.868}{\Character{61}}%
\TextAt{9}{301.877}{594.868}{1}%
\TextAt{13}{307.746}{594.868}{\Character{59}}%
\TextAt{13}{336.881}{594.868}{x}%
\TextAt{21}{343.553}{593.075}{i}%
\NativeMathAt{350.255}{594.868}{11.955}{\in}%
\TextAt{20}{362.21}{597.956}{\Character{98}}%
\TextAt{18}{361.546}{594.868}{Z}%
\TextAt{13}{369.516}{594.868}{\Character{59}}%
\TextAt{8}{79.226}{571.346}{implies}%
\TextAt{13}{118.953}{571.346}{x}%
\TextAt{21}{125.605}{569.553}{i}%
\TextAt{9}{132.307}{571.346}{\Character{61}}%
\TextAt{9}{144.717}{571.346}{0}%
\TextAt{8}{150.586}{571.346}{,}%
\TextAt{13}{157.735}{571.346}{i}%
\TextAt{9}{165.049}{571.346}{\Character{61}}%
\TextAt{9}{177.459}{571.346}{1}%
\TextAt{13}{183.328}{571.346}{\Character{59}}%
\TextAt{13}{188.552}{571.346}{\Character{58}}%
\TextAt{13}{193.789}{571.346}{\Character{58}}%
\TextAt{13}{199.013}{571.346}{\Character{58}}%
\TextAt{13}{204.25}{571.346}{\Character{59}}%
\TextAt{13}{209.486}{571.346}{t}%
\TextAt{8}{213.776}{571.346}{.}%
\TextAt{12}{79.226}{539.356}{Exercise}%
\TextAt{12}{131.658}{539.356}{VI.1.5.}%
\TextAt{8}{174.745}{539.356}{Let}%
\TextAt{13}{195.046}{539.356}{\Character{34}}%
\NativeMathAt{203.845}{539.356}{11.955}{\in}%
\TextAt{16}{215.119}{539.356}{\Character{79}}%
\TextAt{3}{224.986}{543.695}{\Character{3}}%
\TextAt{21}{224.654}{536.336}{K}%
\TextAt{8}{236.086}{539.356}{be}%
\TextAt{8}{251.342}{539.356}{totally}%
\TextAt{8}{288.081}{539.356}{positive,}%
\TextAt{8}{334.112}{539.356}{i.e.,}%
\TextAt{13}{355.682}{539.356}{\Character{34}}%
\NativeMathAt{364.481}{539.356}{11.955}{\in}%
\TextAt{13}{375.772}{539.356}{I}%
\TextAt{17}{381.875}{543.695}{1}%
\TextAt{21}{380.935}{536.336}{K}%
\TextAt{8}{389.137}{539.356}{.}%
\TextAt{8}{397.363}{539.356}{Extend}%
\TextAt{8}{437.498}{539.356}{the}%
\TextAt{8}{456.986}{539.356}{exponentiation}%
\TextAt{18}{79.226}{522.31}{Z}%
\NativeArrowAt{90.574}{522.31}{11.955}{19.26}{\longrightarrow}%
\TextAt{13}{113.213}{522.31}{I}%
\TextAt{17}{119.316}{526.649}{1}%
\TextAt{21}{118.376}{519.29}{K}%
\TextAt{8}{126.579}{522.31}{,}%
\TextAt{13}{133.77}{522.31}{n}%
\NativeArrowAt{144.136}{522.31}{11.955}{11.955}{\mapsto}%
\TextAt{13}{159.469}{522.31}{\Character{34}}%
\TextAt{21}{164.947}{526.649}{n}%
\TextAt{8}{170.583}{522.31}{,}%
\TextAt{8}{177.774}{522.31}{by}%
\TextAt{8}{194.067}{522.31}{continuity}%
\TextAt{8}{249.393}{522.31}{to}%
\TextAt{8}{263.737}{522.31}{an}%
\TextAt{8}{280.029}{522.31}{exponentiation}%
\TextAt{20}{360.408}{525.399}{\Character{98}}%
\TextAt{18}{359.744}{522.31}{Z}%
\TextAt{16}{370.394}{522.31}{\Character{2}}%
\TextAt{18}{382.372}{522.31}{R}%
\NativeArrowAt{394.384}{522.31}{11.955}{19.26}{\longrightarrow}%
\TextAt{13}{417.023}{522.31}{I}%
\TextAt{17}{423.126}{526.649}{1}%
\TextAt{21}{422.186}{519.29}{K}%
\TextAt{9}{433.766}{522.31}{\Character{61}}%
\TextAt{13}{446.249}{522.31}{I}%
\TextAt{17}{452.352}{526.649}{1}%
\TextAt{21}{451.412}{519.201}{f}%
\TextAt{16}{459.764}{522.31}{\Character{2}}%
\TextAt{13}{471.742}{522.31}{I}%
\TextAt{17}{477.844}{526.649}{1}%
\NativeMathAt{476.905}{519.355}{7.97}{\infty}%
\TextAt{8}{485.871}{522.31}{,}%
\TextAt{13}{493.062}{522.31}{\Character{21}}%
\NativeArrowAt{503.269}{522.31}{11.955}{11.955}{\mapsto}%
\TextAt{13}{518.602}{522.31}{\Character{34}}%
\TextAt{21}{524.08}{526.649}{\Character{21}}%
\TextAt{8}{529.518}{522.31}{,}%
\TextAt{8}{79.226}{505.264}{in}%
\TextAt{8}{92.879}{505.264}{such}%
\TextAt{8}{119.276}{505.264}{a}%
\TextAt{8}{129.031}{505.264}{way}%
\TextAt{8}{152.774}{505.264}{that}%
\TextAt{14}{178.143}{505.264}{N}%
\TextAt{9}{188.091}{505.264}{\Character{40}}%
\TextAt{13}{192.644}{505.264}{\Character{34}}%
\TextAt{21}{198.122}{509.603}{\Character{21}}%
\TextAt{9}{203.56}{505.264}{\Character{41}}%
\TextAt{9}{211.415}{505.264}{\Character{61}}%
\TextAt{9}{223.836}{505.264}{1}%
\TextAt{8}{229.711}{505.264}{.}%
\TextAt{12}{79.226}{473.275}{Exercise}%
\TextAt{12}{131.968}{473.275}{VI.1.6.}%
\TextAt{8}{175.326}{473.275}{Let}%
\TextAt{14}{195.897}{473.275}{p}%
\TextAt{17}{201.88}{471.481}{1}%
\TextAt{13}{206.612}{473.274}{\Character{59}}%
\TextAt{13}{211.836}{473.274}{\Character{58}}%
\TextAt{13}{217.073}{473.274}{\Character{58}}%
\TextAt{13}{222.297}{473.274}{\Character{58}}%
\TextAt{13}{227.534}{473.274}{\Character{59}}%
\TextAt{14}{232.833}{473.274}{p}%
\TextAt{21}{238.816}{471.481}{s}%
\TextAt{8}{246.729}{473.275}{be}%
\TextAt{8}{262.255}{473.275}{the}%
\TextAt{8}{282.012}{473.275}{complex}%
\TextAt{8}{327.462}{473.275}{primes}%
\TextAt{8}{364.842}{473.275}{of}%
\TextAt{13}{377.773}{473.275}{K}%
\TextAt{8}{388.591}{473.275}{.}%
\TextAt{8}{396.907}{473.275}{For}%
\TextAt{13}{417.477}{473.275}{y}%
\NativeMathAt{426.935}{473.274}{11.955}{\in}%
\TextAt{18}{438.226}{473.275}{R}%
\TextAt{8}{446.86}{473.275}{,}%
\TextAt{8}{453.691}{473.275}{let}%
\TextAt{13}{470.197}{473.275}{\Character{30}}%
\TextAt{21}{477.125}{471.481}{k}%
\TextAt{9}{482.245}{473.274}{\Character{40}}%
\TextAt{13}{486.797}{473.274}{y}%
\TextAt{9}{492.934}{473.274}{\Character{41}}%
\TextAt{8}{500.985}{473.275}{be}%
\TextAt{8}{516.511}{473.275}{the}%
\TextAt{8}{79.226}{456.229}{idèle}%
\TextAt{8}{106.897}{456.229}{having}%
\TextAt{8}{144.993}{456.229}{component}%
\TextAt{13}{205.194}{456.229}{e}%
\TextAt{17}{210.619}{460.567}{2}%
\TextAt{21}{214.853}{460.567}{\Character{25}}%
\TextAt{21}{219.994}{460.567}{iy}%
\TextAt{8}{232.118}{456.229}{at}%
\TextAt{14}{246.794}{456.229}{p}%
\TextAt{21}{252.777}{454.435}{k}%
\TextAt{8}{257.896}{456.229}{,}%
\TextAt{8}{265.502}{456.229}{and}%
\TextAt{8}{288.63}{456.229}{components}%
\TextAt{9}{353.446}{456.229}{1}%
\TextAt{8}{363.561}{456.229}{at}%
\TextAt{8}{378.237}{456.229}{all}%
\TextAt{8}{394.861}{456.229}{other}%
\TextAt{8}{425.796}{456.229}{places.}%
\TextAt{8}{465.97}{456.229}{Let}%
\TextAt{13}{487.304}{456.229}{\Character{34}}%
\TextAt{17}{492.782}{454.435}{1}%
\TextAt{13}{497.515}{456.229}{\Character{59}}%
\TextAt{13}{502.739}{456.229}{\Character{58}}%
\TextAt{13}{507.976}{456.229}{\Character{58}}%
\TextAt{13}{513.2}{456.229}{\Character{58}}%
\TextAt{13}{518.437}{456.229}{\Character{59}}%
\TextAt{13}{523.673}{456.229}{\Character{34}}%
\TextAt{21}{529.213}{454.435}{t}%
\ErrataAt{79.19}{439.183}{455.81}{have images forming a $\mathbb Z$-basis of $\mathcal O_{K,+}^{\times}/\mu_+$, where $\mu_+$ is the torsion subgroup. (i) The idèles of the form}%
\TextAt{13}{161.889}{410.181}{\Character{11}}%
\TextAt{9}{172.732}{410.181}{\Character{61}}%
\TextAt{13}{185.157}{410.181}{\Character{34}}%
\TextAt{21}{190.635}{415.406}{\Character{21}}%
\TextAt{24}{195.575}{414.299}{1}%
\TextAt{17}{190.635}{407.25}{1}%
\TextAt{16}{202.217}{410.181}{\Character{1}}%
\TextAt{16}{207.513}{410.181}{\Character{1}}%
\TextAt{16}{212.821}{410.181}{\Character{1}}%
\TextAt{13}{218.157}{410.181}{\Character{34}}%
\TextAt{21}{223.635}{415.118}{\Character{21}}%
\TextAt{22}{228.575}{414.121}{t}%
\TextAt{21}{223.635}{407.226}{t}%
\TextAt{13}{232.422}{410.181}{\Character{30}}%
\TextAt{17}{239.35}{408.388}{1}%
\TextAt{9}{244.082}{410.181}{\Character{40}}%
\TextAt{13}{248.634}{410.181}{y}%
\TextAt{17}{254.342}{408.388}{1}%
\TextAt{9}{259.074}{410.181}{\Character{41}}%
\TextAt{16}{265.619}{410.181}{\Character{1}}%
\TextAt{16}{270.915}{410.181}{\Character{1}}%
\TextAt{16}{276.223}{410.181}{\Character{1}}%
\TextAt{13}{281.559}{410.181}{\Character{30}}%
\TextAt{21}{288.488}{408.388}{s}%
\TextAt{9}{292.902}{410.181}{\Character{40}}%
\TextAt{13}{297.454}{410.181}{y}%
\TextAt{21}{303.162}{408.388}{s}%
\TextAt{9}{307.576}{410.181}{\Character{41}}%
\TextAt{13}{312.128}{410.181}{\Character{59}}%
\TextAt{13}{341.263}{410.181}{\Character{21}}%
\TextAt{21}{348.111}{408.388}{i}%
\NativeMathAt{354.813}{410.181}{11.955}{\in}%
\TextAt{20}{366.768}{413.27}{\Character{98}}%
\TextAt{18}{366.104}{410.181}{Z}%
\TextAt{16}{376.731}{410.181}{\Character{2}}%
\TextAt{18}{388.686}{410.181}{R}%
\TextAt{13}{397.32}{410.181}{\Character{59}}%
\TextAt{13}{414.5}{410.181}{y}%
\TextAt{21}{420.227}{408.388}{i}%
\NativeMathAt{426.929}{410.181}{11.955}{\in}%
\TextAt{18}{438.22}{410.181}{R}%
\TextAt{13}{446.855}{410.181}{\Character{59}}%
\ErrataAt{79.19}{381.18}{455.81}{form a group and have absolute norm $N(\alpha)=1$. (ii) The idèle $\alpha$ is principal if and only if}%
\TextAt{13}{79.226}{364.134}{\Character{21}}%
\TextAt{21}{86.055}{362.341}{i}%
\NativeMathAt{92.757}{364.134}{11.955}{\in}%
\TextAt{18}{104.048}{364.134}{Z}%
\NativeMathAt{115.339}{364.134}{11.955}{\subseteq}%
\TextAt{20}{128.622}{367.223}{\Character{98}}%
\TextAt{18}{127.958}{364.134}{Z}%
\TextAt{16}{138.585}{364.134}{\Character{2}}%
\TextAt{18}{150.54}{364.134}{R}%
\TextAt{8}{163.072}{364.134}{and}%
\TextAt{13}{185.834}{364.134}{y}%
\TextAt{21}{191.542}{362.341}{i}%
\NativeMathAt{198.244}{364.134}{11.955}{\in}%
\TextAt{18}{209.535}{364.134}{Z}%
\NativeMathAt{220.826}{364.134}{11.955}{\subseteq}%
\TextAt{18}{233.445}{364.134}{R}%
\TextAt{8}{242.08}{364.134}{.}%
\TextAt{12}{79.226}{332.144}{Exercise}%
\TextAt{12}{132.427}{332.144}{VI.1.7.}%
\TextAt{8}{176.183}{332.144}{Sending}%
\TextAt{9}{174.912}{303.143}{\Character{40}}%
\TextAt{13}{179.465}{303.143}{\Character{21}}%
\TextAt{17}{186.293}{301.35}{1}%
\TextAt{13}{191.025}{303.143}{\Character{59}}%
\TextAt{13}{196.249}{303.143}{\Character{58}}%
\TextAt{13}{201.486}{303.143}{\Character{58}}%
\TextAt{13}{206.722}{303.143}{\Character{58}}%
\TextAt{13}{211.947}{303.143}{\Character{59}}%
\TextAt{13}{217.183}{303.143}{\Character{21}}%
\TextAt{21}{224.075}{301.35}{t}%
\TextAt{13}{227.631}{303.143}{\Character{59}}%
\TextAt{13}{232.855}{303.143}{y}%
\TextAt{17}{238.583}{301.35}{1}%
\TextAt{13}{243.315}{303.143}{\Character{59}}%
\TextAt{13}{248.539}{303.143}{\Character{58}}%
\TextAt{13}{253.776}{303.143}{\Character{58}}%
\TextAt{13}{259}{303.143}{\Character{58}}%
\TextAt{13}{264.237}{303.143}{\Character{59}}%
\TextAt{13}{269.473}{303.143}{y}%
\TextAt{21}{275.244}{301.35}{s}%
\TextAt{9}{279.658}{303.143}{\Character{41}}%
\NativeArrowAt{287.531}{303.143}{11.955}{19.26}{\longmapsto}%
\TextAt{13}{310.113}{303.143}{\Character{34}}%
\TextAt{21}{315.591}{308.368}{\Character{21}}%
\TextAt{24}{320.531}{307.261}{1}%
\TextAt{17}{315.591}{300.212}{1}%
\TextAt{16}{327.172}{303.143}{\Character{1}}%
\TextAt{16}{332.468}{303.143}{\Character{1}}%
\TextAt{16}{337.776}{303.143}{\Character{1}}%
\TextAt{13}{343.113}{303.143}{\Character{34}}%
\TextAt{21}{348.591}{308.079}{\Character{21}}%
\TextAt{22}{353.53}{307.083}{t}%
\TextAt{21}{348.591}{300.188}{t}%
\TextAt{13}{357.377}{303.143}{\Character{30}}%
\TextAt{17}{364.305}{301.35}{1}%
\TextAt{9}{369.037}{303.143}{\Character{40}}%
\TextAt{13}{373.59}{303.143}{y}%
\TextAt{17}{379.298}{301.35}{1}%
\TextAt{9}{384.03}{303.143}{\Character{41}}%
\TextAt{16}{390.575}{303.143}{\Character{1}}%
\TextAt{16}{395.871}{303.143}{\Character{1}}%
\TextAt{16}{401.179}{303.143}{\Character{1}}%
\TextAt{13}{406.515}{303.143}{\Character{30}}%
\TextAt{21}{413.443}{301.35}{s}%
\TextAt{9}{417.857}{303.143}{\Character{40}}%
\TextAt{13}{422.409}{303.143}{y}%
\TextAt{21}{428.117}{301.35}{s}%
\TextAt{9}{432.531}{303.143}{\Character{41}}%
\TextAt{8}{79.226}{274.142}{defines}%
\TextAt{8}{117.65}{274.142}{a}%
\TextAt{8}{127.405}{274.142}{continuous}%
\TextAt{8}{186.01}{274.142}{homomorphism}%
\TextAt{13}{240.519}{245.141}{f}%
\TextAt{9}{250.886}{245.141}{\Character{58}}%
\TextAt{9}{257.437}{245.141}{\Character{40}}%
\TextAt{20}{262.675}{248.229}{\Character{98}}%
\TextAt{18}{262.011}{245.141}{Z}%
\TextAt{16}{272.638}{245.141}{\Character{2}}%
\TextAt{18}{284.593}{245.141}{R}%
\TextAt{9}{293.227}{245.141}{\Character{41}}%
\TextAt{21}{297.78}{250.077}{t}%
\TextAt{16}{303.993}{245.141}{\Character{2}}%
\TextAt{18}{315.948}{245.141}{R}%
\TextAt{21}{324.582}{250.077}{s}%
\NativeArrowAt{332.317}{245.141}{11.955}{19.26}{\longrightarrow}%
\TextAt{13}{354.899}{245.141}{C}%
\TextAt{17}{364.132}{250.077}{0}%
\TextAt{21}{363.274}{242.185}{K}%
\TextAt{8}{79.226}{216.14}{into}%
\TextAt{8}{102.969}{216.14}{the}%
\TextAt{8}{123.126}{216.14}{group}%
\TextAt{13}{156.301}{216.14}{C}%
\TextAt{17}{165.535}{220.478}{0}%
\TextAt{21}{164.676}{213.119}{K}%
\TextAt{9}{176.2}{216.14}{\Character{61}}%
\TextAt{16}{188.625}{216.14}{\Character{102}}%
\TextAt{9}{194.603}{216.14}{\Character{91}}%
\TextAt{13}{197.854}{216.14}{\Character{11}}%
\TextAt{9}{205.376}{216.14}{\Character{93}}%
\NativeMathAt{211.949}{216.14}{11.955}{\in}%
\TextAt{13}{223.24}{216.14}{C}%
\TextAt{21}{231.615}{214.347}{K}%
\TextAt{16}{243.138}{216.14}{\Character{106}}%
\TextAt{14}{249.78}{216.14}{N}%
\TextAt{9}{259.728}{216.14}{\Character{40}\Character{91}}%
\TextAt{13}{267.532}{216.14}{\Character{11}}%
\TextAt{9}{275.054}{216.14}{\Character{93}\Character{41}}%
\TextAt{9}{286.148}{216.14}{\Character{61}}%
\TextAt{9}{298.57}{216.14}{1}%
\TextAt{16}{304.457}{216.14}{\Character{103}}%
\TextAt{8}{310.435}{216.14}{,}%
\TextAt{8}{317.584}{216.14}{with}%
\TextAt{8}{344.244}{216.14}{kernel}%
\TextAt{18}{378.711}{216.14}{Z}%
\TextAt{21}{386.681}{220.478}{t}%
\TextAt{16}{392.894}{216.14}{\Character{2}}%
\TextAt{18}{404.849}{216.14}{Z}%
\TextAt{21}{412.819}{220.478}{s}%
\TextAt{8}{417.233}{216.14}{.}%
\TextAt{12}{79.226}{184.15}{Exercise}%
\TextAt{12}{132.035}{184.15}{VI.1.8.}%
\TextAt{8}{175.451}{184.15}{(i)}%
\TextAt{8}{191.37}{184.15}{The}%
\TextAt{8}{215.084}{184.15}{image}%
\TextAt{13}{248.565}{184.15}{D}%
\TextAt{17}{258.609}{188.488}{0}%
\TextAt{21}{258.284}{181.129}{K}%
\TextAt{8}{270.044}{184.15}{of}%
\TextAt{13}{283.034}{184.15}{f}%
\TextAt{8}{293.637}{184.15}{is}%
\TextAt{8}{305.061}{184.15}{compact,}%
\TextAt{8}{354.869}{184.15}{connected}%
\TextAt{8}{409.152}{184.15}{and}%
\TextAt{8}{431.574}{184.15}{arbitrarily}%
\TextAt{8}{487.507}{184.15}{divisible.}%
\TextAt{8}{79.226}{167.104}{(ii)}%
\TextAt{13}{98.737}{167.104}{f}%
\TextAt{8}{109.681}{167.104}{yields}%
\TextAt{8}{142.581}{167.104}{a}%
\TextAt{8}{152.337}{167.104}{topological}%
\TextAt{8}{211.862}{167.104}{isomorphism}%
\TextAt{13}{217.41}{138.103}{f}%
\TextAt{9}{227.778}{138.103}{\Character{58}}%
\TextAt{20}{234.35}{147.787}{\Character{0}}%
\TextAt{9}{239.829}{138.103}{\Character{40}}%
\TextAt{20}{245.046}{141.191}{\Character{98}}%
\TextAt{18}{244.382}{138.103}{Z}%
\TextAt{16}{255.009}{138.103}{\Character{2}}%
\TextAt{18}{266.964}{138.103}{R}%
\TextAt{9}{275.598}{138.103}{\Character{41}\Character{47}}%
\TextAt{18}{286.003}{138.103}{Z}%
\TextAt{20}{293.974}{147.787}{\Character{1}}%
\TextAt{21}{299.453}{145.11}{t}%
\TextAt{16}{305.666}{138.103}{\Character{2}}%
\TextAt{9}{317.621}{138.103}{\Character{40}}%
\TextAt{18}{322.173}{138.103}{R}%
\TextAt{9}{330.808}{138.103}{\Character{47}}%
\TextAt{18}{336.661}{138.103}{Z}%
\TextAt{9}{344.631}{138.103}{\Character{41}}%
\TextAt{21}{349.183}{143.039}{s}%
\NativeMathAt{359.506}{144.88}{7.97}{\sim}%
\NativeArrowAt{356.918}{138.103}{11.955}{13.175}{\longrightarrow}%
\TextAt{13}{373.412}{138.103}{D}%
\TextAt{17}{383.456}{143.039}{0}%
\TextAt{21}{383.131}{135.147}{K}%
\TextAt{13}{391.334}{138.103}{\Character{58}}%
\TextAt{8}{297.211}{43.033}{335}%
\Destination{78.23}{809.815}{nav.page.335}%
\Destination{79.226}{726.055}{nav.ex.6.1.3}%
\Destination{79.226}{660.055}{nav.ex.6.1.4}%
\Destination{79.226}{553.355}{nav.ex.6.1.5}%
\Destination{79.226}{487.255}{nav.ex.6.1.6}%
\Destination{79.226}{346.155}{nav.ex.6.1.7}%
\Destination{79.226}{198.155}{nav.ex.6.1.8}%
\Hotlink{305.362}{758.955}{25.078}{11.047}{nav.number.6.3.8}%
\Destination{79.226}{456.183}{errata.132}%
\Destination{79.226}{398.18}{errata.133}%
}

% Page 348 (body page 336)
\Page{%
\TextAt{2}{79.226}{796.943}{VI.}%
\TextAt{2}{98.012}{796.943}{Global}%
\TextAt{2}{133.248}{796.943}{Class}%
\TextAt{2}{161.841}{796.943}{Field}%
\TextAt{2}{189.572}{796.943}{Theory}%
\Vector{q 0 G 0 g 0.299 w 7.226 18.729 m 460.769 18.729 l S Q}%
\TextAt{12}{79.226}{761.124}{Exercise}%
\TextAt{12}{131.807}{761.124}{VI.1.9.}%
\TextAt{8}{175.025}{761.124}{The}%
\TextAt{8}{198.541}{761.124}{group}%
\TextAt{13}{231.178}{761.124}{D}%
\TextAt{17}{241.222}{765.462}{0}%
\TextAt{21}{240.897}{758.103}{K}%
\TextAt{8}{252.459}{761.124}{is}%
\TextAt{8}{263.685}{761.124}{the}%
\TextAt{8}{283.303}{761.124}{intersection}%
\TextAt{8}{345.589}{761.124}{of}%
\TextAt{8}{358.381}{761.124}{all}%
\TextAt{8}{374.102}{761.124}{closed}%
\TextAt{8}{408.09}{761.124}{subgroups}%
\TextAt{8}{462.964}{761.124}{of}%
\TextAt{8}{475.756}{761.124}{finite}%
\TextAt{8}{505.129}{761.124}{index}%
\TextAt{8}{79.226}{744.078}{in}%
\TextAt{13}{92.879}{744.078}{C}%
\TextAt{17}{102.113}{748.416}{0}%
\TextAt{21}{101.254}{741.057}{K}%
\TextAt{8}{109.457}{744.078}{,}%
\TextAt{8}{116.606}{744.078}{and}%
\TextAt{8}{139.369}{744.078}{it}%
\TextAt{8}{151.073}{744.078}{is}%
\TextAt{8}{162.837}{744.078}{the}%
\TextAt{8}{182.993}{744.078}{connected}%
\TextAt{8}{237.616}{744.078}{component}%
\TextAt{8}{297.452}{744.078}{of}%
\TextAt{9}{310.782}{744.078}{1}%
\TextAt{8}{320.532}{744.078}{in}%
\TextAt{13}{334.185}{744.078}{C}%
\TextAt{17}{343.419}{748.416}{0}%
\TextAt{21}{342.56}{741.057}{K}%
\TextAt{8}{350.763}{744.078}{.}%
\TextAt{12}{79.226}{712.151}{Exercise}%
\TextAt{12}{132.14}{712.151}{VI.1.10.}%
\TextAt{8}{182.377}{712.151}{The}%
\TextAt{8}{206.181}{712.151}{connected}%
\TextAt{8}{260.554}{712.151}{component}%
\TextAt{13}{320.14}{712.151}{D}%
\TextAt{21}{329.86}{710.358}{K}%
\TextAt{8}{341.71}{712.151}{of}%
\TextAt{9}{354.79}{712.151}{1}%
\TextAt{8}{364.291}{712.151}{in}%
\TextAt{8}{377.694}{712.151}{the}%
\TextAt{8}{397.601}{712.151}{idèle}%
\TextAt{8}{424.657}{712.151}{class}%
\TextAt{8}{451.845}{712.151}{group}%
\TextAt{13}{484.771}{712.151}{C}%
\TextAt{21}{493.146}{710.358}{K}%
\TextAt{8}{504.996}{712.151}{is}%
\TextAt{8}{516.51}{712.151}{the}%
\TextAt{8}{79.226}{695.105}{direct}%
\TextAt{8}{112.39}{695.105}{product}%
\TextAt{8}{156.289}{695.105}{of}%
\TextAt{13}{169.619}{695.105}{t}%
\TextAt{8}{177.744}{695.105}{copies}%
\TextAt{8}{212.27}{695.105}{of}%
\TextAt{8}{225.6}{695.105}{the}%
\TextAt{8}{245.757}{695.105}{“solenoid”}%
\TextAt{9}{301.743}{695.105}{\Character{40}}%
\TextAt{20}{306.959}{698.193}{\Character{98}}%
\TextAt{18}{306.295}{695.105}{Z}%
\TextAt{16}{316.922}{695.105}{\Character{2}}%
\TextAt{18}{328.877}{695.105}{R}%
\TextAt{9}{337.511}{695.105}{\Character{41}\Character{47}}%
\TextAt{18}{347.917}{695.105}{Z}%
\TextAt{8}{355.887}{695.105}{,}%
\TextAt{13}{363.036}{695.105}{s}%
\TextAt{8}{372.447}{695.105}{circles}%
\TextAt{18}{407.62}{695.105}{R}%
\TextAt{9}{416.254}{695.105}{\Character{47}}%
\TextAt{18}{422.107}{695.105}{Z}%
\TextAt{8}{430.077}{695.105}{,}%
\TextAt{8}{437.226}{695.105}{and}%
\TextAt{8}{459.989}{695.105}{a}%
\TextAt{8}{469.744}{695.105}{real}%
\TextAt{8}{492.507}{695.105}{line.}%
\TextAt{12}{79.226}{663.178}{Exercise}%
\TextAt{12}{131.936}{663.178}{VI.1.11.}%
\TextAt{8}{181.997}{663.178}{Every}%
\TextAt{8}{215.224}{663.178}{ideal}%
\TextAt{8}{242.761}{663.178}{class}%
\TextAt{8}{269.772}{663.178}{of}%
\TextAt{8}{282.676}{663.178}{the}%
\TextAt{8}{302.406}{663.178}{ray}%
\TextAt{8}{322.148}{663.178}{class}%
\TextAt{8}{349.159}{663.178}{group}%
\TextAt{13}{381.909}{663.178}{C}%
\TextAt{13}{391.127}{663.178}{l}%
\TextAt{27}{394.892}{667.517}{m}%
\TextAt{21}{394.629}{660.158}{K}%
\TextAt{8}{406.303}{663.178}{can}%
\TextAt{8}{427.336}{663.178}{be}%
\TextAt{8}{442.834}{663.178}{represented}%
\TextAt{8}{504.575}{663.178}{by}%
\TextAt{8}{520.408}{663.178}{an}%
\TextAt{8}{79.226}{646.132}{integral}%
\TextAt{8}{121.834}{646.132}{ideal}%
\TextAt{8}{149.798}{646.132}{which}%
\TextAt{8}{183.284}{646.132}{is}%
\TextAt{8}{195.048}{646.132}{prime}%
\TextAt{8}{228.211}{646.132}{to}%
\TextAt{8}{242.522}{646.132}{an}%
\TextAt{8}{258.781}{646.132}{arbitrary}%
\TextAt{8}{308.55}{646.132}{fixed}%
\TextAt{8}{336.836}{646.132}{ideal.}%
\TextAt{12}{79.226}{614.205}{Exercise}%
\TextAt{12}{132.47}{614.205}{VI.1.12.}%
\TextAt{8}{182.994}{614.205}{Let}%
\TextAt{16}{204.001}{614.205}{\Character{79}}%
\TextAt{9}{217.236}{614.205}{\Character{61}}%
\TextAt{16}{229.726}{614.205}{\Character{79}}%
\TextAt{21}{239.243}{612.412}{K}%
\TextAt{8}{247.446}{614.205}{.}%
\TextAt{8}{256.007}{614.205}{Every}%
\TextAt{8}{289.699}{614.205}{class}%
\TextAt{8}{317.174}{614.205}{in}%
\TextAt{9}{330.864}{614.205}{\Character{40}}%
\TextAt{16}{335.417}{614.205}{\Character{79}}%
\TextAt{9}{345.266}{614.205}{\Character{47}}%
\TextAt{14}{351.119}{614.205}{m}%
\TextAt{9}{360.286}{614.205}{\Character{41}}%
\TextAt{3}{364.838}{618.544}{\Character{3}}%
\TextAt{8}{373.506}{614.205}{can}%
\TextAt{8}{395.003}{614.205}{be}%
\TextAt{8}{410.965}{614.205}{represented}%
\TextAt{8}{473.169}{614.205}{by}%
\TextAt{8}{489.466}{614.205}{a}%
\TextAt{8}{499.259}{614.205}{totally}%
\TextAt{8}{79.226}{597.159}{positive}%
\TextAt{8}{122.54}{597.159}{number}%
\TextAt{8}{165.136}{597.159}{in}%
\TextAt{16}{178.789}{597.159}{\Character{79}}%
\TextAt{8}{192.536}{597.159}{which}%
\TextAt{8}{226.022}{597.159}{is}%
\TextAt{8}{237.786}{597.159}{prime}%
\TextAt{8}{270.95}{597.159}{to}%
\TextAt{8}{285.26}{597.159}{an}%
\TextAt{8}{301.519}{597.159}{arbitrary}%
\TextAt{8}{351.289}{597.159}{fixed}%
\TextAt{8}{379.575}{597.159}{ideal.}%
\TextAt{12}{79.226}{565.233}{Exercise}%
\TextAt{12}{132.427}{565.233}{VI.1.13.}%
\TextAt{8}{182.913}{565.233}{For}%
\TextAt{8}{203.883}{565.233}{every}%
\TextAt{8}{234.775}{565.233}{module}%
\TextAt{14}{276.068}{565.233}{m}%
\TextAt{8}{285.235}{565.233}{,}%
\TextAt{8}{292.384}{565.233}{one}%
\TextAt{8}{313.843}{565.233}{has}%
\TextAt{8}{334.717}{565.233}{an}%
\TextAt{8}{350.976}{565.233}{exact}%
\TextAt{8}{381.868}{565.233}{sequence}%
\TextAt{9}{173.957}{536.342}{1}%
\NativeArrowAt{183.131}{536.342}{11.955}{19.26}{\longrightarrow}%
\TextAt{16}{205.702}{536.342}{\Character{79}}%
\TextAt{3}{215.562}{541.278}{\Character{3}}%
\TextAt{17}{215.23}{533.387}{\Character{43}}%
\TextAt{9}{222.315}{536.342}{\Character{47}}%
\TextAt{16}{228.168}{536.342}{\Character{79}}%
\TextAt{27}{238.018}{541.278}{m}%
\TextAt{17}{237.686}{533.387}{\Character{43}}%
\NativeArrowAt{248.091}{536.342}{11.955}{19.26}{\longrightarrow}%
\TextAt{9}{270.673}{536.342}{\Character{40}}%
\TextAt{16}{275.225}{536.342}{\Character{79}}%
\TextAt{9}{285.075}{536.342}{\Character{47}}%
\TextAt{14}{290.928}{536.342}{m}%
\TextAt{9}{300.095}{536.342}{\Character{41}}%
\TextAt{3}{304.647}{541.278}{\Character{3}}%
\NativeArrowAt{312.7}{536.342}{11.955}{19.26}{\longrightarrow}%
\TextAt{13}{335.282}{536.342}{C}%
\TextAt{13}{344.5}{536.342}{l}%
\TextAt{27}{348.266}{541.278}{m}%
\TextAt{21}{348.003}{533.387}{K}%
\NativeArrowAt{359.526}{536.342}{11.955}{19.26}{\longrightarrow}%
\TextAt{13}{382.108}{536.342}{C}%
\TextAt{13}{391.326}{536.342}{l}%
\TextAt{17}{395.092}{541.278}{1}%
\TextAt{21}{394.829}{533.387}{K}%
\NativeArrowAt{406.352}{536.342}{11.955}{19.26}{\longrightarrow}%
\TextAt{9}{428.934}{536.342}{1}%
\TextAt{13}{434.787}{536.342}{\Character{59}}%
\TextAt{8}{79.226}{507.451}{where}%
\TextAt{16}{113.259}{507.451}{\Character{79}}%
\TextAt{3}{123.108}{511.79}{\Character{3}}%
\TextAt{17}{122.776}{504.496}{\Character{43}}%
\TextAt{8}{129.861}{507.451}{,}%
\TextAt{8}{137.289}{507.451}{resp.}%
\TextAt{16}{167.281}{507.451}{\Character{79}}%
\TextAt{27}{177.13}{511.79}{m}%
\TextAt{17}{176.798}{504.496}{\Character{43}}%
\TextAt{8}{183.883}{507.451}{,}%
\TextAt{8}{191.311}{507.451}{is}%
\TextAt{8}{203.298}{507.451}{the}%
\TextAt{8}{223.678}{507.451}{group}%
\TextAt{8}{257.076}{507.451}{of}%
\TextAt{8}{270.63}{507.451}{totally}%
\TextAt{8}{308.261}{507.451}{positive}%
\TextAt{8}{351.797}{507.451}{units}%
\TextAt{8}{381.347}{507.451}{of}%
\TextAt{16}{394.9}{507.451}{\Character{79}}%
\TextAt{8}{404.75}{507.451}{,}%
\TextAt{8}{412.178}{507.451}{resp.}%
\TextAt{8}{442.169}{507.451}{of}%
\TextAt{8}{455.722}{507.451}{totally}%
\TextAt{8}{493.353}{507.451}{positive}%
\TextAt{8}{79.226}{490.405}{units}%
\NativeMathAt{108.552}{490.405}{11.955}{\equiv}%
\TextAt{9}{121.172}{490.405}{1}%
\TextAt{8}{130.345}{490.405}{mod}%
\TextAt{14}{155.783}{490.405}{m}%
\TextAt{8}{164.95}{490.405}{.}%
\TextAt{12}{79.226}{458.479}{Exercise}%
\TextAt{12}{132.427}{458.479}{VI.1.14.}%
\TextAt{8}{182.913}{458.479}{Compute}%
\TextAt{8}{233.639}{458.479}{the}%
\TextAt{8}{253.796}{458.479}{kernels}%
\TextAt{8}{292.877}{458.479}{of}%
\TextAt{13}{306.207}{458.479}{C}%
\TextAt{13}{315.425}{458.479}{l}%
\TextAt{27}{319.19}{462.817}{m}%
\TextAt{21}{318.928}{455.458}{K}%
\NativeArrowAt{330.451}{458.479}{11.955}{19.26}{\longrightarrow}%
\TextAt{13}{353.033}{458.479}{C}%
\TextAt{13}{362.25}{458.479}{l}%
\TextAt{21}{365.754}{456.685}{K}%
\TextAt{8}{377.854}{458.479}{and}%
\TextAt{13}{400.616}{458.479}{C}%
\TextAt{13}{409.834}{458.479}{l}%
\TextAt{27}{413.6}{462.817}{m}%
\TextAt{21}{413.337}{455.458}{K}%
\NativeArrowAt{424.86}{458.479}{11.955}{19.26}{\longrightarrow}%
\TextAt{13}{447.442}{458.479}{C}%
\TextAt{13}{456.66}{458.479}{l}%
\TextAt{27}{460.426}{462.817}{m}%
\TextAt{23}{466.537}{465.63}{\Character{48}}%
\TextAt{21}{460.163}{455.458}{K}%
\TextAt{8}{473.617}{458.479}{for}%
\TextAt{14}{491.502}{458.479}{m}%
\TextAt{3}{500.668}{462.817}{\Character{48}}%
\TextAt{16}{506.784}{458.479}{\Character{106}}%
\TextAt{14}{513.426}{458.479}{m}%
\TextAt{8}{522.593}{458.479}{.}%
\TextAt{15}{79.226}{405.663}{§}%
\TextAt{15}{93.785}{405.663}{2}%
\TextAt{15}{120.693}{405.663}{Idèles}%
\TextAt{15}{175.817}{405.663}{in}%
\TextAt{15}{198.421}{405.663}{Field}%
\TextAt{15}{247.416}{405.663}{Extensions}%
\TextAt{8}{79.226}{373.777}{We}%
\TextAt{8}{99.106}{373.777}{shall}%
\TextAt{8}{126.218}{373.777}{now}%
\TextAt{8}{150.341}{373.777}{study}%
\TextAt{8}{182.331}{373.777}{the}%
\TextAt{8}{202.223}{373.777}{behaviour}%
\TextAt{8}{256.27}{373.777}{of}%
\TextAt{8}{269.335}{373.777}{idèles}%
\TextAt{8}{300.99}{373.777}{and}%
\TextAt{8}{323.488}{373.777}{idèle}%
\TextAt{8}{350.529}{373.777}{classes}%
\TextAt{8}{387.516}{373.777}{when}%
\TextAt{8}{417.808}{373.777}{we}%
\TextAt{8}{434.771}{373.777}{pass}%
\TextAt{8}{459.994}{373.777}{from}%
\TextAt{8}{487.369}{373.777}{a}%
\TextAt{8}{496.86}{373.777}{field}%
\TextAt{13}{521.952}{373.777}{K}%
\TextAt{8}{79.226}{356.731}{to}%
\TextAt{8}{92.748}{356.731}{an}%
\TextAt{8}{108.219}{356.731}{extension}%
\TextAt{13}{159.197}{356.731}{L}%
\TextAt{8}{167.161}{356.731}{.}%
\TextAt{8}{175.347}{356.731}{So}%
\TextAt{8}{190.818}{356.731}{let}%
\TextAt{13}{206.934}{356.731}{L}%
\TextAt{16}{218.22}{356.731}{\Character{106}}%
\TextAt{13}{224.861}{356.731}{K}%
\TextAt{8}{238.788}{356.731}{be}%
\TextAt{8}{253.925}{356.731}{a}%
\TextAt{8}{262.892}{356.731}{finite}%
\TextAt{8}{292.015}{356.731}{extension}%
\TextAt{8}{342.993}{356.731}{of}%
\TextAt{8}{355.535}{356.731}{algebraic}%
\TextAt{8}{404.182}{356.731}{number}%
\TextAt{8}{445.99}{356.731}{fields.}%
\TextAt{8}{480.25}{356.731}{We}%
\TextAt{8}{499.606}{356.731}{embed}%
\TextAt{8}{79.226}{339.685}{the}%
\TextAt{8}{99.73}{339.685}{idèle}%
\TextAt{8}{127.383}{339.685}{group}%
\TextAt{13}{160.905}{339.685}{I}%
\TextAt{21}{166.068}{337.892}{K}%
\TextAt{8}{178.515}{339.685}{of}%
\TextAt{13}{192.193}{339.685}{K}%
\TextAt{8}{207.255}{339.685}{into}%
\TextAt{8}{231.345}{339.685}{the}%
\TextAt{8}{251.849}{339.685}{idèle}%
\TextAt{8}{279.501}{339.685}{group}%
\TextAt{13}{313.024}{339.685}{I}%
\TextAt{21}{318.187}{337.892}{L}%
\TextAt{8}{328.689}{339.685}{of}%
\TextAt{13}{342.366}{339.685}{L}%
\TextAt{8}{354.575}{339.685}{by}%
\TextAt{8}{371.181}{339.685}{sending}%
\TextAt{8}{413.862}{339.685}{an}%
\TextAt{8}{430.468}{339.685}{idèle}%
\TextAt{13}{458.121}{339.685}{\Character{11}}%
\TextAt{9}{469.555}{339.685}{\Character{61}}%
\TextAt{9}{482.562}{339.685}{\Character{40}}%
\TextAt{13}{487.124}{339.685}{\Character{11}}%
\TextAt{27}{494.57}{337.892}{p}%
\TextAt{9}{499.057}{339.685}{\Character{41}}%
\NativeMathAt{507.521}{339.685}{11.955}{\in}%
\TextAt{13}{519.404}{339.685}{I}%
\TextAt{21}{524.567}{337.892}{K}%
\TextAt{8}{79.226}{322.639}{to}%
\TextAt{8}{93.536}{322.639}{the}%
\TextAt{8}{113.693}{322.639}{idèle}%
\TextAt{13}{140.999}{322.639}{\Character{11}}%
\TextAt{3}{148.52}{326.978}{\Character{48}}%
\TextAt{9}{154.636}{322.639}{\Character{61}}%
\TextAt{9}{167.046}{322.639}{\Character{40}}%
\TextAt{13}{171.614}{322.639}{\Character{11}}%
\TextAt{3}{179.136}{326.978}{\Character{48}}%
\TextAt{27}{179.06}{319.554}{P}%
\TextAt{9}{186.153}{322.639}{\Character{41}}%
\NativeMathAt{194.026}{322.639}{11.955}{\in}%
\TextAt{13}{205.317}{322.639}{I}%
\TextAt{21}{210.48}{320.846}{L}%
\TextAt{8}{220.635}{322.639}{whose}%
\TextAt{8}{255.161}{322.639}{components}%
\TextAt{13}{319.612}{322.639}{\Character{11}}%
\TextAt{3}{327.133}{326.978}{\Character{48}}%
\TextAt{27}{327.057}{319.554}{P}%
\TextAt{8}{338.048}{322.639}{are}%
\TextAt{8}{357.559}{322.639}{given}%
\TextAt{8}{388.128}{322.639}{by}%
\TextAt{13}{218.866}{293.749}{\Character{11}}%
\TextAt{3}{226.388}{298.685}{\Character{48}}%
\TextAt{27}{226.312}{290.793}{P}%
\TextAt{9}{236.726}{293.749}{\Character{61}}%
\TextAt{13}{249.151}{293.749}{\Character{11}}%
\TextAt{27}{256.597}{291.955}{p}%
\NativeMathAt{264.405}{293.749}{11.955}{\in}%
\TextAt{13}{275.696}{293.749}{K}%
\TextAt{3}{286.513}{298.685}{\Character{3}}%
\TextAt{27}{285.655}{290.793}{p}%
\NativeMathAt{294.567}{293.749}{11.955}{\subseteq}%
\TextAt{13}{307.186}{293.749}{L}%
\TextAt{3}{315.15}{298.685}{\Character{3}}%
\TextAt{27}{315.15}{290.793}{P}%
\TextAt{8}{346.154}{293.749}{for}%
\TextAt{14}{364.039}{293.749}{P}%
\TextAt{16}{377.253}{293.749}{\Character{106}}%
\TextAt{14}{383.895}{293.749}{p}%
\TextAt{13}{389.878}{293.749}{\Character{58}}%
\TextAt{8}{79.226}{264.858}{In}%
\TextAt{8}{93.348}{264.858}{this}%
\TextAt{8}{115.671}{264.858}{way}%
\TextAt{8}{138.914}{264.858}{we}%
\TextAt{8}{155.642}{264.858}{obtain}%
\TextAt{8}{191.57}{264.858}{an}%
\TextAt{8}{207.329}{264.858}{injective}%
\TextAt{8}{253.323}{264.858}{homomorphism}%
\TextAt{13}{335.494}{264.858}{I}%
\TextAt{21}{340.657}{263.065}{K}%
\NativeArrowAt{352.18}{264.858}{11.955}{19.26}{\longrightarrow}%
\TextAt{13}{374.762}{264.858}{I}%
\TextAt{21}{379.925}{263.065}{L}%
\TextAt{8}{386.182}{264.858}{,}%
\TextAt{8}{392.932}{264.858}{which}%
\TextAt{8}{425.919}{264.858}{will}%
\TextAt{8}{447.524}{264.858}{always}%
\TextAt{8}{484.492}{264.858}{be}%
\TextAt{8}{499.917}{264.858}{tacitly}%
\TextAt{8}{79.226}{247.812}{used}%
\TextAt{8}{106.637}{247.812}{to}%
\TextAt{8}{121.638}{247.812}{consider}%
\TextAt{13}{167.914}{247.812}{I}%
\TextAt{21}{173.077}{246.019}{K}%
\TextAt{8}{185.868}{247.812}{as}%
\TextAt{8}{200.929}{247.812}{a}%
\TextAt{8}{211.375}{247.812}{subgroup}%
\TextAt{8}{262.863}{247.812}{of}%
\TextAt{13}{276.884}{247.812}{I}%
\TextAt{21}{282.047}{246.019}{L}%
\TextAt{8}{288.304}{247.812}{.}%
\TextAt{8}{298.825}{247.812}{An}%
\TextAt{8}{318.692}{247.812}{element}%
\TextAt{13}{362.625}{247.812}{\Character{11}}%
\TextAt{9}{374.644}{247.812}{\Character{61}}%
\TextAt{9}{388.237}{247.812}{\Character{40}}%
\TextAt{13}{392.8}{247.812}{\Character{11}}%
\TextAt{27}{400.245}{246.019}{P}%
\TextAt{9}{407.339}{247.812}{\Character{41}}%
\NativeMathAt{416.389}{247.812}{11.955}{\in}%
\TextAt{13}{428.857}{247.812}{I}%
\TextAt{21}{434.02}{246.019}{L}%
\TextAt{8}{444.866}{247.812}{therefore}%
\TextAt{8}{494.656}{247.812}{belongs}%
\TextAt{8}{79.226}{230.766}{to}%
\TextAt{8}{94.25}{230.766}{the}%
\TextAt{8}{115.121}{230.766}{group}%
\TextAt{13}{149.01}{230.766}{I}%
\TextAt{21}{154.173}{228.973}{K}%
\TextAt{8}{166.987}{230.766}{if}%
\TextAt{8}{178.424}{230.766}{and}%
\TextAt{8}{201.901}{230.766}{only}%
\TextAt{8}{228.307}{230.766}{if}%
\TextAt{8}{239.744}{230.766}{its}%
\TextAt{8}{256.777}{230.766}{components}%
\TextAt{13}{321.941}{230.766}{\Character{11}}%
\TextAt{27}{329.387}{228.973}{P}%
\TextAt{8}{341.091}{230.766}{belong}%
\TextAt{8}{379.201}{230.766}{to}%
\TextAt{13}{394.225}{230.766}{K}%
\TextAt{27}{404.184}{228.973}{p}%
\TextAt{8}{413.282}{230.766}{(}%
\TextAt{14}{417.837}{230.766}{P}%
\TextAt{16}{432.268}{230.766}{\Character{106}}%
\TextAt{14}{440.126}{230.766}{p}%
\TextAt{8}{446.109}{230.766}{),}%
\TextAt{8}{458.706}{230.766}{and}%
\TextAt{8}{482.182}{230.766}{if}%
\TextAt{8}{493.62}{230.766}{one}%
\TextAt{8}{515.793}{230.766}{has}%
\TextAt{8}{79.226}{213.72}{furthermore}%
\TextAt{13}{143.939}{213.72}{\Character{11}}%
\TextAt{27}{151.385}{211.927}{P}%
\TextAt{9}{161.8}{213.72}{\Character{61}}%
\TextAt{13}{174.225}{213.72}{\Character{11}}%
\TextAt{27}{181.671}{211.927}{P}%
\TextAt{23}{188.266}{214.196}{\Character{48}}%
\TextAt{8}{195.346}{213.72}{whenever}%
\TextAt{14}{246.717}{213.72}{P}%
\TextAt{8}{260.508}{213.72}{and}%
\TextAt{14}{283.271}{213.72}{P}%
\TextAt{3}{293.164}{218.058}{\Character{48}}%
\TextAt{8}{299.856}{213.72}{lie}%
\TextAt{8}{315.458}{213.72}{above}%
\TextAt{8}{348.633}{213.72}{the}%
\TextAt{8}{368.79}{213.72}{same}%
\TextAt{8}{398.116}{213.72}{place}%
\TextAt{14}{428.028}{213.72}{p}%
\TextAt{8}{437.908}{213.72}{of}%
\TextAt{13}{451.238}{213.72}{K}%
\TextAt{8}{462.056}{213.72}{.}%
\TextAt{8}{97.159}{193.685}{Every}%
\TextAt{8}{130.423}{193.685}{isomorphism}%
\TextAt{13}{198.452}{193.685}{\Character{27}}%
\TextAt{9}{208.855}{193.685}{\Character{58}}%
\TextAt{13}{215.428}{193.685}{L}%
\NativeArrowAt{226.713}{193.685}{11.955}{19.26}{\longrightarrow}%
\TextAt{13}{249.295}{193.685}{\Character{27}}%
\TextAt{13}{256.361}{193.685}{L}%
\TextAt{8}{267.849}{193.685}{induces}%
\TextAt{8}{309.135}{193.685}{an}%
\TextAt{8}{325.004}{193.685}{isomorphism}%
\TextAt{13}{393.033}{193.685}{\Character{27}}%
\TextAt{9}{403.437}{193.685}{\Character{58}}%
\TextAt{13}{410.009}{193.685}{I}%
\TextAt{21}{415.172}{191.892}{L}%
\NativeArrowAt{424.75}{193.685}{11.955}{19.26}{\longrightarrow}%
\TextAt{13}{447.332}{193.685}{I}%
\TextAt{21}{452.495}{191.892}{\Character{27}}%
\TextAt{21}{457.612}{191.892}{L}%
\TextAt{8}{467.385}{193.685}{like}%
\TextAt{8}{488.454}{193.685}{this.}%
\TextAt{8}{515.698}{193.685}{For}%
\TextAt{8}{79.226}{176.639}{each}%
\TextAt{8}{105.563}{176.639}{place}%
\TextAt{14}{135.475}{176.639}{P}%
\TextAt{8}{149.266}{176.639}{of}%
\TextAt{13}{162.596}{176.639}{L}%
\TextAt{8}{170.56}{176.639}{,}%
\TextAt{13}{177.71}{176.639}{\Character{27}}%
\TextAt{8}{188.689}{176.639}{induces}%
\TextAt{8}{230.365}{176.639}{an}%
\TextAt{8}{246.624}{176.639}{isomorphism}%
\TextAt{13}{257.218}{147.749}{\Character{27}}%
\TextAt{9}{267.622}{147.749}{\Character{58}}%
\TextAt{13}{274.194}{147.749}{L}%
\TextAt{27}{282.159}{145.955}{P}%
\NativeArrowAt{292.573}{147.749}{11.955}{19.26}{\longrightarrow}%
\TextAt{9}{315.155}{147.749}{\Character{40}}%
\TextAt{13}{319.707}{147.749}{\Character{27}}%
\TextAt{13}{326.773}{147.749}{L}%
\TextAt{9}{334.754}{147.749}{\Character{41}}%
\TextAt{21}{339.307}{145.955}{\Character{27}}%
\TextAt{27}{344.432}{145.955}{P}%
\TextAt{13}{351.526}{147.749}{\Character{58}}%
\TextAt{8}{79.226}{118.858}{For}%
\TextAt{8}{101.216}{118.858}{if}%
\TextAt{8}{112.961}{118.858}{we}%
\TextAt{8}{131.209}{118.858}{have}%
\TextAt{13}{159.225}{118.858}{\Character{11}}%
\TextAt{9}{171.807}{118.858}{\Character{61}}%
\TextAt{14}{185.972}{118.858}{P}%
\TextAt{8}{195.865}{118.858}{-}%
\TextAt{8}{201.755}{118.858}{lim}%
\TextAt{13}{220.006}{118.858}{\Character{11}}%
\TextAt{21}{227.452}{117.065}{i}%
\TextAt{8}{230.833}{118.858}{,}%
\TextAt{8}{239.259}{118.858}{for}%
\TextAt{8}{258.164}{118.858}{some}%
\TextAt{8}{288.511}{118.858}{sequence}%
\TextAt{13}{338.034}{118.858}{\Character{11}}%
\TextAt{21}{345.48}{117.065}{i}%
\NativeMathAt{353.921}{118.858}{11.955}{\in}%
\TextAt{13}{366.952}{118.858}{L}%
\TextAt{8}{374.916}{118.858}{,}%
\TextAt{8}{383.341}{118.858}{then}%
\TextAt{8}{411.022}{118.858}{the}%
\TextAt{8}{432.199}{118.858}{sequence}%
\TextAt{13}{481.722}{118.858}{\Character{27}}%
\TextAt{13}{488.788}{118.858}{\Character{11}}%
\TextAt{21}{496.25}{117.065}{i}%
\NativeMathAt{504.692}{118.858}{11.955}{\in}%
\TextAt{13}{517.723}{118.858}{\Character{27}}%
\TextAt{13}{524.788}{118.858}{L}%
\TextAt{8}{79.226}{101.812}{converges}%
\TextAt{8}{131.65}{101.812}{with}%
\TextAt{8}{158.31}{101.812}{respect}%
\TextAt{8}{198.359}{101.812}{to}%
\TextAt{16}{212.67}{101.812}{\Character{106}}%
\TextAt{16}{219.867}{101.812}{\Character{106}}%
\TextAt{21}{223.209}{100.019}{\Character{27}}%
\TextAt{27}{228.334}{100.019}{P}%
\TextAt{8}{239.325}{101.812}{in}%
\TextAt{9}{252.978}{101.812}{\Character{40}}%
\TextAt{13}{257.53}{101.812}{\Character{27}}%
\TextAt{13}{264.596}{101.812}{L}%
\TextAt{9}{272.577}{101.812}{\Character{41}}%
\TextAt{21}{277.13}{100.019}{\Character{27}}%
\TextAt{27}{282.255}{100.019}{P}%
\TextAt{8}{289.348}{101.812}{,}%
\TextAt{8}{296.498}{101.812}{and}%
\TextAt{8}{319.26}{101.812}{the}%
\TextAt{8}{339.417}{101.812}{isomorphism}%
\TextAt{8}{407.836}{101.812}{is}%
\TextAt{8}{419.6}{101.812}{given}%
\TextAt{8}{450.169}{101.812}{by}%
\TextAt{13}{212.667}{72.921}{\Character{11}}%
\TextAt{9}{223.51}{72.921}{\Character{61}}%
\TextAt{14}{235.935}{72.921}{P}%
\TextAt{8}{245.829}{72.921}{-}%
\TextAt{8}{251.719}{72.921}{lim}%
\TextAt{13}{269.97}{72.921}{\Character{11}}%
\TextAt{21}{277.416}{71.128}{i}%
\NativeArrowAt{284.118}{72.921}{11.955}{19.26}{\longmapsto}%
\TextAt{13}{306.7}{72.921}{\Character{27}}%
\TextAt{13}{313.766}{72.921}{\Character{11}}%
\TextAt{9}{324.625}{72.921}{\Character{61}}%
\TextAt{13}{337.05}{72.921}{\Character{27}}%
\TextAt{14}{344.133}{72.921}{P}%
\TextAt{8}{354.026}{72.921}{-}%
\TextAt{8}{359.916}{72.921}{lim}%
\TextAt{13}{378.167}{72.921}{\Character{27}}%
\TextAt{13}{385.233}{72.921}{\Character{11}}%
\TextAt{21}{392.695}{71.128}{i}%
\TextAt{13}{396.077}{72.921}{\Character{58}}%
\TextAt{8}{297.211}{43.033}{336}%
\Destination{78.23}{809.815}{nav.page.336}%
\Destination{79.23}{431.495}{nav.section.6.2}%
\Destination{79.226}{775.155}{nav.ex.6.1.9}%
\Destination{79.226}{726.155}{nav.ex.6.1.10}%
\Destination{79.226}{677.155}{nav.ex.6.1.11}%
\Destination{79.226}{628.155}{nav.ex.6.1.12}%
\Destination{79.226}{579.255}{nav.ex.6.1.13}%
\Destination{79.226}{472.455}{nav.ex.6.1.14}%
\Destination{79.36}{776.155}{nav.print.368}%
\Destination{79.23}{431.495}{outline.57}%
}

% Page 349 (body page 337)
\Page{%
\TextAt{2}{79.226}{796.943}{VI.}%
\TextAt{2}{98.012}{796.943}{Global}%
\TextAt{2}{133.248}{796.943}{Class}%
\TextAt{2}{161.841}{796.943}{Field}%
\TextAt{2}{189.572}{796.943}{Theory}%
\Vector{q 0 G 0 g 0.299 w 7.226 18.729 m 460.769 18.729 l S Q}%
\TextAt{8}{79.226}{761.124}{For}%
\TextAt{8}{100.196}{761.124}{an}%
\TextAt{8}{116.455}{761.124}{idèle}%
\TextAt{13}{143.76}{761.124}{\Character{11}}%
\NativeMathAt{154.603}{761.124}{11.955}{\in}%
\TextAt{13}{165.894}{761.124}{I}%
\TextAt{21}{171.057}{759.33}{L}%
\TextAt{8}{177.314}{761.124}{,}%
\TextAt{8}{184.463}{761.124}{we}%
\TextAt{8}{201.69}{761.124}{then}%
\TextAt{8}{228.35}{761.124}{define}%
\TextAt{13}{262.16}{761.124}{\Character{27}}%
\TextAt{13}{269.226}{761.124}{\Character{11}}%
\NativeMathAt{280.085}{761.124}{11.955}{\in}%
\TextAt{13}{291.376}{761.124}{I}%
\TextAt{21}{296.538}{759.33}{\Character{27}}%
\TextAt{21}{301.655}{759.33}{L}%
\TextAt{8}{311.818}{761.124}{to}%
\TextAt{8}{326.129}{761.124}{be}%
\TextAt{8}{342.053}{761.124}{the}%
\TextAt{8}{362.209}{761.124}{idèle}%
\TextAt{8}{389.515}{761.124}{with}%
\TextAt{8}{416.175}{761.124}{components}%
\TextAt{9}{242.233}{732.123}{\Character{40}}%
\TextAt{13}{246.785}{732.123}{\Character{27}}%
\TextAt{13}{253.851}{732.123}{\Character{11}}%
\TextAt{9}{261.389}{732.123}{\Character{41}}%
\TextAt{21}{265.942}{730.329}{\Character{27}}%
\TextAt{27}{271.067}{730.329}{P}%
\TextAt{9}{281.482}{732.123}{\Character{61}}%
\TextAt{13}{293.907}{732.123}{\Character{27}}%
\TextAt{13}{300.973}{732.123}{\Character{11}}%
\TextAt{27}{308.435}{730.329}{P}%
\NativeMathAt{318.849}{732.123}{11.955}{\in}%
\TextAt{9}{330.14}{732.123}{\Character{40}}%
\TextAt{13}{334.693}{732.123}{\Character{27}}%
\TextAt{13}{341.759}{732.123}{L}%
\TextAt{9}{349.74}{732.123}{\Character{41}}%
\TextAt{21}{354.292}{730.329}{\Character{27}}%
\TextAt{27}{359.417}{730.329}{P}%
\TextAt{13}{366.511}{732.123}{\Character{58}}%
\TextAt{8}{79.226}{703.121}{If}%
\TextAt{13}{91.055}{703.121}{L}%
\TextAt{16}{102.573}{703.121}{\Character{106}}%
\TextAt{13}{109.447}{703.121}{K}%
\TextAt{8}{124.299}{703.121}{is}%
\TextAt{8}{136.199}{703.121}{a}%
\TextAt{8}{146.091}{703.121}{Galois}%
\TextAt{8}{182.141}{703.121}{extension}%
\TextAt{8}{234.043}{703.121}{with}%
\TextAt{8}{260.839}{703.121}{Galois}%
\TextAt{8}{296.889}{703.121}{group}%
\TextAt{13}{330.201}{703.121}{G}%
\TextAt{9}{342.988}{703.121}{\Character{61}}%
\TextAt{13}{355.646}{703.121}{G}%
\TextAt{9}{364.88}{703.121}{\Character{40}}%
\TextAt{13}{369.432}{703.121}{L}%
\TextAt{16}{380.95}{703.121}{\Character{106}}%
\TextAt{13}{387.824}{703.121}{K}%
\TextAt{9}{398.642}{703.121}{\Character{41}}%
\TextAt{8}{403.194}{703.121}{,}%
\TextAt{8}{410.514}{703.121}{then}%
\TextAt{8}{437.311}{703.121}{every}%
\TextAt{13}{468.339}{703.121}{\Character{27}}%
\NativeMathAt{478.975}{703.121}{11.955}{\in}%
\TextAt{13}{490.499}{703.121}{G}%
\TextAt{8}{503.766}{703.121}{yields}%
\TextAt{8}{79.226}{686.075}{an}%
\TextAt{8}{94.722}{686.075}{automorphism}%
\TextAt{13}{171.429}{686.075}{\Character{27}}%
\TextAt{9}{181.832}{686.075}{\Character{58}}%
\TextAt{13}{188.405}{686.075}{I}%
\TextAt{21}{193.567}{684.282}{L}%
\NativeArrowAt{203.145}{686.075}{11.955}{19.26}{\longrightarrow}%
\TextAt{13}{225.727}{686.075}{I}%
\TextAt{21}{230.89}{684.282}{L}%
\TextAt{8}{237.148}{686.075}{,}%
\TextAt{8}{243.686}{686.075}{i.e.,}%
\TextAt{13}{265.181}{686.075}{I}%
\TextAt{21}{270.344}{684.282}{L}%
\TextAt{8}{279.736}{686.075}{is}%
\TextAt{8}{290.737}{686.075}{turned}%
\TextAt{8}{327.692}{686.075}{into}%
\TextAt{8}{350.672}{686.075}{a}%
\TextAt{13}{359.665}{686.075}{G}%
\TextAt{8}{368.898}{686.075}{-module.}%
\TextAt{8}{418.386}{686.075}{As}%
\TextAt{8}{434.91}{686.075}{to}%
\TextAt{8}{448.457}{686.075}{the}%
\TextAt{8}{467.851}{686.075}{fixed}%
\TextAt{8}{495.374}{686.075}{module}%
\TextAt{13}{79.226}{669.029}{I}%
\TextAt{21}{85.329}{673.368}{G}%
\TextAt{21}{84.389}{666.009}{L}%
\TextAt{9}{95.747}{669.029}{\Character{61}}%
\TextAt{16}{108.172}{669.029}{\Character{102}}%
\TextAt{13}{114.15}{669.029}{\Character{11}}%
\NativeMathAt{124.992}{669.029}{11.955}{\in}%
\TextAt{13}{136.283}{669.029}{I}%
\TextAt{21}{141.446}{667.236}{L}%
\TextAt{16}{151.024}{669.029}{\Character{106}}%
\TextAt{13}{157.666}{669.029}{\Character{27}}%
\TextAt{13}{164.732}{669.029}{\Character{11}}%
\TextAt{9}{175.591}{669.029}{\Character{61}}%
\TextAt{13}{188.016}{669.029}{\Character{11}}%
\TextAt{8}{199.435}{669.029}{for}%
\TextAt{8}{217.32}{669.029}{all}%
\TextAt{13}{233.579}{669.029}{\Character{27}}%
\NativeMathAt{243.983}{669.029}{11.955}{\in}%
\TextAt{13}{255.274}{669.029}{G}%
\TextAt{16}{264.507}{669.029}{\Character{103}}%
\TextAt{8}{270.485}{669.029}{,}%
\TextAt{8}{277.634}{669.029}{we}%
\TextAt{8}{294.861}{669.029}{have}%
\TextAt{8}{321.856}{669.029}{the}%
\TextAt{12}{79.226}{643.017}{Proposition}%
\TextAt{12}{151.937}{643.017}{(2.1).}%
\TextAt{8}{187.253}{643.017}{If}%
\TextAt{13}{198.945}{643.017}{L}%
\TextAt{16}{210.231}{643.017}{\Character{106}}%
\TextAt{13}{216.872}{643.017}{K}%
\TextAt{8}{231.587}{643.017}{is}%
\TextAt{8}{243.351}{643.017}{a}%
\TextAt{8}{253.107}{643.017}{Galois}%
\TextAt{8}{289.02}{643.017}{extension}%
\TextAt{8}{340.786}{643.017}{with}%
\TextAt{8}{367.446}{643.017}{Galois}%
\TextAt{8}{403.359}{643.017}{group}%
\TextAt{13}{436.535}{643.017}{G}%
\TextAt{8}{445.768}{643.017}{,}%
\TextAt{8}{452.918}{643.017}{then}%
\TextAt{13}{479.578}{643.017}{I}%
\TextAt{21}{485.68}{647.356}{G}%
\TextAt{21}{484.741}{639.997}{L}%
\TextAt{9}{496.098}{643.017}{\Character{61}}%
\TextAt{13}{508.524}{643.017}{I}%
\TextAt{21}{513.686}{641.224}{K}%
\TextAt{8}{521.889}{643.017}{.}%
\TextAt{19}{79.226}{608.038}{Proof.}%
\TextAt{8}{116}{608.038}{Let}%
\TextAt{13}{138.807}{608.038}{\Character{11}}%
\NativeMathAt{152.782}{608.038}{11.955}{\in}%
\TextAt{13}{167.204}{608.038}{I}%
\TextAt{21}{172.367}{606.245}{K}%
\NativeMathAt{187.023}{608.038}{11.955}{\subseteq}%
\TextAt{13}{202.774}{608.038}{I}%
\TextAt{21}{207.937}{606.245}{L}%
\TextAt{8}{214.194}{608.038}{.}%
\TextAt{8}{228.155}{608.038}{For}%
\TextAt{13}{250.963}{608.038}{\Character{27}}%
\NativeMathAt{264.498}{608.038}{11.955}{\in}%
\TextAt{13}{278.92}{608.038}{G}%
\TextAt{8}{288.154}{608.038}{,}%
\TextAt{8}{297.6}{608.038}{the}%
\TextAt{8}{319.594}{608.038}{induced}%
\TextAt{8}{364.997}{608.038}{map}%
\TextAt{13}{392.849}{608.038}{\Character{27}}%
\TextAt{9}{406.384}{608.038}{\Character{58}}%
\TextAt{13}{416.088}{608.038}{L}%
\TextAt{27}{424.053}{606.245}{P}%
\NativeArrowAt{437.599}{608.038}{11.955}{19.26}{\longrightarrow}%
\TextAt{13}{463.313}{608.038}{L}%
\TextAt{21}{471.277}{606.245}{\Character{27}}%
\TextAt{27}{476.403}{606.245}{P}%
\TextAt{8}{489.231}{608.038}{is}%
\TextAt{8}{502.833}{608.038}{a}%
\TextAt{13}{514.426}{608.038}{K}%
\TextAt{27}{524.385}{606.245}{p}%
\TextAt{8}{528.872}{608.038}{-}%
\TextAt{8}{79.226}{590.992}{isomorphism,}%
\TextAt{8}{150.897}{590.992}{if}%
\TextAt{14}{161.621}{590.992}{P}%
\TextAt{16}{174.835}{590.992}{\Character{106}}%
\TextAt{14}{181.477}{590.992}{p}%
\TextAt{8}{187.46}{590.992}{.}%
\TextAt{8}{195.908}{590.992}{Therefore}%
\TextAt{9}{234.876}{561.991}{\Character{40}}%
\TextAt{13}{239.428}{561.991}{\Character{27}}%
\TextAt{13}{246.494}{561.991}{\Character{11}}%
\TextAt{9}{254.032}{561.991}{\Character{41}}%
\TextAt{21}{258.585}{560.198}{\Character{27}}%
\TextAt{27}{263.71}{560.198}{P}%
\TextAt{9}{274.125}{561.991}{\Character{61}}%
\TextAt{13}{286.55}{561.991}{\Character{27}}%
\TextAt{13}{293.615}{561.991}{\Character{11}}%
\TextAt{27}{301.078}{560.198}{P}%
\TextAt{9}{311.492}{561.991}{\Character{61}}%
\TextAt{13}{323.918}{561.991}{\Character{11}}%
\TextAt{27}{331.364}{560.198}{P}%
\TextAt{9}{341.778}{561.991}{\Character{61}}%
\TextAt{13}{354.203}{561.991}{\Character{11}}%
\TextAt{21}{361.649}{560.198}{\Character{27}}%
\TextAt{27}{366.774}{560.198}{P}%
\TextAt{13}{373.868}{561.991}{\Character{59}}%
\TextAt{8}{79.226}{532.99}{so}%
\TextAt{8}{94.622}{532.99}{that}%
\TextAt{13}{121.017}{532.99}{\Character{27}}%
\TextAt{13}{128.083}{532.99}{\Character{11}}%
\TextAt{9}{140.691}{532.99}{\Character{61}}%
\TextAt{13}{154.865}{532.99}{\Character{11}}%
\TextAt{8}{162.387}{532.99}{,}%
\TextAt{8}{170.818}{532.99}{and}%
\TextAt{8}{194.607}{532.99}{therefore}%
\TextAt{13}{244.733}{532.99}{\Character{11}}%
\NativeMathAt{257.324}{532.99}{11.955}{\in}%
\TextAt{13}{270.364}{532.99}{I}%
\TextAt{21}{276.467}{537.329}{G}%
\TextAt{21}{275.527}{529.97}{L}%
\TextAt{8}{283.563}{532.99}{.}%
\TextAt{8}{295.09}{532.99}{If}%
\TextAt{8}{307.808}{532.99}{conversely}%
\TextAt{13}{364.832}{532.99}{\Character{11}}%
\TextAt{9}{377.423}{532.99}{\Character{61}}%
\TextAt{9}{391.59}{532.99}{\Character{40}}%
\TextAt{13}{396.15}{532.99}{\Character{11}}%
\TextAt{27}{403.595}{531.197}{P}%
\TextAt{9}{410.689}{532.99}{\Character{41}}%
\NativeMathAt{420.311}{532.99}{11.955}{\in}%
\TextAt{13}{433.35}{532.99}{I}%
\TextAt{21}{439.453}{537.329}{G}%
\TextAt{21}{438.513}{529.97}{L}%
\TextAt{8}{446.55}{532.99}{,}%
\TextAt{8}{454.982}{532.99}{then}%
\TextAt{9}{482.668}{532.99}{\Character{40}}%
\TextAt{13}{487.22}{532.99}{\Character{27}}%
\TextAt{13}{494.286}{532.99}{\Character{11}}%
\TextAt{9}{501.824}{532.99}{\Character{41}}%
\TextAt{21}{506.377}{531.197}{\Character{27}}%
\TextAt{27}{511.502}{531.197}{P}%
\TextAt{9}{523.665}{532.99}{\Character{61}}%
\TextAt{13}{79.226}{515.944}{\Character{27}}%
\TextAt{13}{86.292}{515.944}{\Character{11}}%
\TextAt{27}{93.754}{514.151}{P}%
\TextAt{9}{104.885}{515.944}{\Character{61}}%
\TextAt{13}{118.026}{515.944}{\Character{11}}%
\TextAt{21}{125.472}{514.151}{\Character{27}}%
\TextAt{27}{130.597}{514.151}{P}%
\TextAt{8}{142.008}{515.944}{for}%
\TextAt{8}{160.313}{515.944}{all}%
\TextAt{13}{176.992}{515.944}{\Character{27}}%
\NativeMathAt{188.112}{515.944}{11.955}{\in}%
\TextAt{13}{200.118}{515.944}{G}%
\TextAt{8}{209.352}{515.944}{.}%
\TextAt{8}{219.061}{515.944}{In}%
\TextAt{8}{234.102}{515.944}{particular,}%
\TextAt{8}{291.869}{515.944}{if}%
\TextAt{13}{303.013}{515.944}{\Character{27}}%
\TextAt{8}{314.413}{515.944}{belongs}%
\TextAt{8}{356.843}{515.944}{to}%
\TextAt{8}{371.574}{515.944}{the}%
\TextAt{8}{392.15}{515.944}{decomposition}%
\TextAt{8}{469.705}{515.944}{group}%
\TextAt{13}{503.301}{515.944}{G}%
\TextAt{27}{512.534}{514.151}{P}%
\TextAt{9}{523.665}{515.944}{\Character{61}}%
\TextAt{13}{79.226}{498.898}{G}%
\TextAt{9}{88.46}{498.898}{\Character{40}}%
\TextAt{13}{93.012}{498.898}{L}%
\TextAt{27}{100.977}{497.105}{P}%
\TextAt{16}{111.979}{498.898}{\Character{106}}%
\TextAt{13}{119.208}{498.898}{K}%
\TextAt{27}{129.167}{497.105}{p}%
\TextAt{9}{133.654}{498.898}{\Character{41}}%
\TextAt{8}{138.206}{498.898}{,}%
\TextAt{8}{145.787}{498.898}{then}%
\TextAt{13}{172.791}{498.898}{\Character{27}}%
\TextAt{14}{179.874}{498.898}{P}%
\TextAt{9}{193.676}{498.898}{\Character{61}}%
\TextAt{14}{206.689}{498.898}{P}%
\TextAt{8}{220.824}{498.898}{and}%
\TextAt{13}{243.932}{498.898}{\Character{27}}%
\TextAt{13}{250.998}{498.898}{\Character{11}}%
\TextAt{27}{258.46}{497.105}{P}%
\TextAt{9}{269.462}{498.898}{\Character{61}}%
\TextAt{13}{282.475}{498.898}{\Character{11}}%
\TextAt{27}{289.92}{497.105}{P}%
\TextAt{8}{301.256}{498.898}{so}%
\TextAt{8}{315.971}{498.898}{that}%
\TextAt{13}{341.685}{498.898}{\Character{11}}%
\TextAt{27}{349.13}{497.105}{P}%
\NativeMathAt{360.133}{498.898}{11.955}{\in}%
\TextAt{13}{372.011}{498.898}{K}%
\TextAt{3}{382.829}{503.237}{\Character{3}}%
\TextAt{27}{381.97}{495.943}{p}%
\TextAt{8}{387.561}{498.898}{.}%
\TextAt{8}{397.044}{498.898}{If}%
\TextAt{13}{409.081}{498.898}{\Character{27}}%
\NativeMathAt{420.072}{498.898}{11.955}{\in}%
\TextAt{13}{431.95}{498.898}{G}%
\TextAt{8}{445.426}{498.898}{is}%
\TextAt{8}{457.535}{498.898}{arbitrary,}%
\TextAt{8}{510.007}{498.898}{then}%
\TextAt{13}{79.226}{481.852}{\Character{27}}%
\TextAt{9}{90.923}{481.852}{\Character{58}}%
\TextAt{13}{98.789}{481.852}{L}%
\TextAt{27}{106.754}{480.059}{P}%
\NativeArrowAt{118.462}{481.852}{11.955}{19.26}{\longrightarrow}%
\TextAt{13}{142.337}{481.852}{L}%
\TextAt{21}{150.302}{480.059}{\Character{27}}%
\TextAt{27}{155.427}{480.059}{P}%
\TextAt{8}{167.177}{481.852}{induces}%
\TextAt{8}{209.612}{481.852}{the}%
\TextAt{8}{230.527}{481.852}{identity}%
\TextAt{8}{274.54}{481.852}{on}%
\TextAt{13}{291.559}{481.852}{K}%
\TextAt{27}{301.518}{480.059}{p}%
\TextAt{8}{306.004}{481.852}{,}%
\TextAt{8}{314.103}{481.852}{and}%
\TextAt{8}{337.624}{481.852}{we}%
\TextAt{8}{355.611}{481.852}{get}%
\TextAt{13}{375.881}{481.852}{\Character{11}}%
\TextAt{27}{383.326}{480.059}{P}%
\TextAt{9}{395.034}{481.852}{\Character{61}}%
\TextAt{13}{408.753}{481.852}{\Character{27}}%
\TextAt{13}{415.819}{481.852}{\Character{11}}%
\TextAt{27}{423.281}{480.059}{P}%
\TextAt{9}{434.99}{481.852}{\Character{61}}%
\TextAt{13}{448.709}{481.852}{\Character{11}}%
\TextAt{21}{456.154}{480.059}{\Character{27}}%
\TextAt{27}{461.28}{480.059}{P}%
\TextAt{8}{473.03}{481.852}{for}%
\TextAt{8}{491.674}{481.852}{any}%
\TextAt{8}{514.55}{481.852}{two}%
\TextAt{8}{79.226}{464.806}{places}%
\TextAt{14}{113.753}{464.806}{P}%
\TextAt{8}{127.543}{464.806}{and}%
\TextAt{13}{150.306}{464.806}{\Character{27}}%
\TextAt{14}{157.388}{464.806}{P}%
\TextAt{8}{171.179}{464.806}{above}%
\TextAt{14}{204.355}{464.806}{p}%
\TextAt{8}{210.338}{464.806}{.}%
\TextAt{8}{218.786}{464.806}{This}%
\TextAt{8}{245.506}{464.806}{shows}%
\TextAt{8}{279.124}{464.806}{that}%
\TextAt{13}{304.493}{464.806}{\Character{11}}%
\NativeMathAt{315.335}{464.806}{11.955}{\in}%
\TextAt{13}{326.626}{464.806}{I}%
\TextAt{21}{331.789}{463.013}{K}%
\TextAt{8}{339.992}{464.806}{.}%
\Vector{q 0 G 0 g 0.399 w 452.334 -309.194 m 452.334 -301.124 l S Q}%
\Vector{q 0 G 0 g 0.399 w 452.534 -301.323 m 459.707 -301.323 l S Q}%
\Vector{q 0 G 0 g 0.399 w 452.534 -308.994 m 459.707 -308.994 l S Q}%
\Vector{q 0 G 0 g 0.399 w 459.906 -309.194 m 459.906 -301.124 l S Q}%
\TextAt{8}{97.159}{435.805}{The}%
\TextAt{8}{121.173}{435.805}{idèle}%
\TextAt{8}{148.439}{435.805}{group}%
\TextAt{13}{181.575}{435.805}{I}%
\TextAt{21}{186.738}{434.012}{L}%
\TextAt{8}{196.853}{435.805}{is}%
\TextAt{8}{208.577}{435.805}{the}%
\TextAt{8}{228.694}{435.805}{unit}%
\TextAt{8}{253.366}{435.805}{group}%
\TextAt{8}{286.502}{435.805}{of}%
\TextAt{8}{299.792}{435.805}{the}%
\TextAt{8}{319.909}{435.805}{ring}%
\TextAt{8}{343.935}{435.805}{of}%
\TextAt{8}{357.226}{435.805}{adèles}%
\TextAt{18}{391.713}{435.805}{A}%
\TextAt{21}{400.347}{434.012}{L}%
\TextAt{8}{410.462}{435.805}{of}%
\TextAt{13}{423.752}{435.805}{L}%
\TextAt{8}{431.717}{435.805}{.}%
\TextAt{8}{440.152}{435.805}{It}%
\TextAt{8}{452.785}{435.805}{is}%
\TextAt{8}{464.509}{435.805}{convenient}%
\TextAt{8}{522.357}{435.805}{to}%
\TextAt{8}{79.226}{418.759}{write}%
\TextAt{8}{109.138}{418.759}{this}%
\TextAt{8}{131.96}{418.759}{ring}%
\TextAt{8}{156.026}{418.759}{as}%
\TextAt{18}{229.354}{394.745}{A}%
\TextAt{21}{237.988}{392.951}{L}%
\TextAt{9}{247.566}{394.745}{\Character{61}}%
\TextAt{3}{266.481}{409.689}{\Character{48}}%
\TextAt{20}{259.991}{406.102}{\Character{89}}%
\TextAt{27}{265.635}{380.996}{p}%
\TextAt{13}{277.26}{394.745}{L}%
\TextAt{27}{285.225}{392.951}{p}%
\TextAt{13}{289.711}{394.745}{\Character{59}}%
\TextAt{13}{318.846}{394.745}{L}%
\TextAt{27}{326.83}{392.951}{p}%
\TextAt{9}{334.638}{394.745}{\Character{61}}%
\TextAt{20}{347.064}{406.102}{\Character{89}}%
\TextAt{27}{348.233}{380.199}{P}%
\TextAt{3}{354.829}{380.199}{\Character{106}}%
\TextAt{27}{357.181}{380.199}{p}%
\TextAt{13}{364.332}{394.745}{L}%
\TextAt{27}{372.297}{392.951}{P}%
\TextAt{13}{379.39}{394.745}{\Character{58}}%
\TextAt{8}{79.226}{359.3}{The}%
\TextAt{8}{103.715}{359.3}{restricted}%
\TextAt{8}{156.239}{359.3}{product}%
\TextAt{20}{200.574}{368.266}{\Character{81}}%
\TextAt{3}{211.865}{365.111}{\Character{48}}%
\TextAt{27}{211.865}{355.813}{p}%
\TextAt{13}{218.345}{359.3}{L}%
\TextAt{27}{226.309}{357.506}{p}%
\TextAt{8}{235.129}{359.3}{consists}%
\TextAt{8}{278.675}{359.3}{of}%
\TextAt{8}{292.44}{359.3}{all}%
\TextAt{8}{309.134}{359.3}{families}%
\TextAt{9}{352.226}{359.3}{\Character{40}}%
\TextAt{13}{356.778}{359.3}{\Character{11}}%
\TextAt{27}{364.224}{357.506}{p}%
\TextAt{9}{368.711}{359.3}{\Character{41}}%
\TextAt{8}{377.596}{359.3}{of}%
\TextAt{8}{391.361}{359.3}{elements}%
\TextAt{13}{439.653}{359.3}{\Character{11}}%
\TextAt{27}{447.099}{357.506}{p}%
\NativeMathAt{455.648}{359.3}{11.955}{\in}%
\TextAt{13}{467.681}{359.3}{L}%
\TextAt{27}{475.646}{357.506}{p}%
\TextAt{8}{484.466}{359.3}{such}%
\TextAt{8}{511.298}{359.3}{that}%
\TextAt{13}{79.226}{342.254}{\Character{11}}%
\TextAt{27}{86.672}{340.46}{p}%
\NativeMathAt{94.914}{342.254}{11.955}{\in}%
\TextAt{16}{106.63}{342.254}{\Character{79}}%
\TextAt{27}{116.158}{340.46}{p}%
\TextAt{9}{124.401}{342.254}{\Character{61}}%
\TextAt{20}{137.261}{351.22}{\Character{81}}%
\TextAt{27}{148.552}{338.767}{P}%
\TextAt{3}{155.148}{338.767}{\Character{106}}%
\TextAt{27}{157.5}{338.767}{p}%
\TextAt{16}{163.979}{342.254}{\Character{79}}%
\TextAt{27}{173.497}{340.46}{P}%
\TextAt{8}{184.743}{342.254}{for}%
\TextAt{8}{202.883}{342.254}{almost}%
\TextAt{8}{240.929}{342.254}{all}%
\TextAt{14}{257.443}{342.254}{p}%
\TextAt{8}{263.426}{342.254}{.}%
\TextAt{8}{272.64}{342.254}{Via}%
\TextAt{8}{294.678}{342.254}{the}%
\TextAt{8}{315.09}{342.254}{diagonal}%
\TextAt{8}{362.185}{342.254}{embedding}%
\TextAt{13}{421.619}{342.254}{K}%
\TextAt{27}{431.578}{340.46}{p}%
\NativeArrowAt{439.82}{342.254}{11.955}{19.26}{\longrightarrow}%
\TextAt{13}{462.837}{342.254}{L}%
\TextAt{27}{470.802}{340.46}{p}%
\TextAt{8}{475.289}{342.254}{,}%
\TextAt{8}{482.757}{342.254}{the}%
\TextAt{8}{503.169}{342.254}{factor}%
\TextAt{13}{79.226}{325.208}{L}%
\TextAt{27}{87.191}{323.414}{p}%
\TextAt{8}{96.121}{325.208}{is}%
\TextAt{8}{108.43}{325.208}{a}%
\TextAt{8}{118.732}{325.208}{commutative}%
\TextAt{13}{189.203}{325.208}{K}%
\TextAt{27}{199.163}{323.414}{p}%
\TextAt{8}{203.649}{325.208}{-algebra}%
\TextAt{8}{249.075}{325.208}{of}%
\TextAt{8}{262.951}{325.208}{degree}%
\TextAt{20}{299.912}{334.174}{\Character{80}}%
\TextAt{27}{312.532}{321.721}{P}%
\TextAt{3}{319.127}{321.721}{\Character{106}}%
\TextAt{27}{321.479}{321.721}{p}%
\TextAt{9}{325.966}{325.208}{\Character{91}}%
\TextAt{13}{329.218}{325.208}{L}%
\TextAt{27}{337.182}{323.414}{P}%
\TextAt{9}{348.527}{325.208}{\Character{58}}%
\TextAt{13}{356.03}{325.208}{K}%
\TextAt{27}{365.989}{323.414}{p}%
\TextAt{9}{370.476}{325.208}{\Character{93}}%
\TextAt{9}{377.96}{325.208}{\Character{61}}%
\TextAt{9}{391.314}{325.208}{\Character{91}}%
\TextAt{13}{394.586}{325.208}{L}%
\TextAt{9}{406.801}{325.208}{\Character{58}}%
\TextAt{13}{414.304}{325.208}{K}%
\TextAt{9}{425.122}{325.208}{\Character{93}}%
\TextAt{8}{428.373}{325.208}{.}%
\TextAt{8}{438.459}{325.208}{These}%
\TextAt{8}{472.874}{325.208}{embeddings}%
\TextAt{8}{79.226}{308.162}{yield}%
\TextAt{8}{107.512}{308.162}{the}%
\TextAt{8}{127.668}{308.162}{embedding}%
\TextAt{18}{186.847}{308.162}{A}%
\TextAt{21}{195.481}{306.368}{K}%
\NativeArrowAt{207.004}{308.162}{11.955}{19.26}{\longrightarrow}%
\TextAt{18}{229.586}{308.162}{A}%
\TextAt{21}{238.221}{306.368}{L}%
\TextAt{8}{244.478}{308.162}{,}%
\TextAt{8}{251.627}{308.162}{whose}%
\TextAt{8}{286.154}{308.162}{restriction}%
\TextAt{13}{249.043}{279.161}{I}%
\TextAt{21}{254.206}{277.367}{K}%
\TextAt{9}{265.729}{279.161}{\Character{61}}%
\TextAt{18}{278.155}{279.161}{A}%
\TextAt{3}{286.789}{284.097}{\Character{3}}%
\TextAt{21}{286.789}{276.205}{K}%
\NativeArrowAt{298.313}{279.161}{11.955}{19.26}{\longrightarrow}%
\TextAt{18}{320.895}{279.161}{A}%
\TextAt{3}{329.529}{284.097}{\Character{3}}%
\TextAt{21}{329.529}{276.205}{L}%
\TextAt{9}{339.107}{279.161}{\Character{61}}%
\TextAt{13}{351.533}{279.161}{I}%
\TextAt{21}{356.695}{277.367}{L}%
\TextAt{8}{79.226}{250.159}{turns}%
\TextAt{8}{109.855}{250.159}{out}%
\TextAt{8}{130.669}{250.159}{to}%
\TextAt{8}{144.98}{250.159}{be}%
\TextAt{8}{160.904}{250.159}{the}%
\TextAt{8}{181.06}{250.159}{inclusion}%
\TextAt{8}{229.897}{250.159}{considered}%
\TextAt{8}{287.186}{250.159}{above.}%
\TextAt{8}{97.159}{230.125}{Every}%
\TextAt{13}{130.813}{230.125}{\Character{11}}%
\TextAt{27}{138.258}{228.331}{p}%
\NativeMathAt{146.066}{230.125}{11.955}{\in}%
\TextAt{13}{157.357}{230.125}{L}%
\TextAt{3}{165.321}{234.463}{\Character{3}}%
\TextAt{27}{165.321}{227.169}{p}%
\TextAt{8}{173.951}{230.125}{defines}%
\TextAt{8}{212.375}{230.125}{an}%
\TextAt{8}{228.634}{230.125}{automorphism}%
\TextAt{20}{228.674}{201.124}{\Character{101}}%
\TextAt{13}{227.909}{201.124}{\Character{11}}%
\TextAt{27}{235.355}{199.33}{p}%
\TextAt{9}{243.163}{201.124}{\Character{58}}%
\TextAt{13}{249.735}{201.124}{L}%
\TextAt{27}{257.7}{199.33}{p}%
\NativeArrowAt{265.507}{201.124}{11.955}{19.26}{\longrightarrow}%
\TextAt{13}{288.089}{201.124}{L}%
\TextAt{27}{296.054}{199.33}{p}%
\TextAt{13}{300.541}{201.124}{\Character{59}}%
\TextAt{13}{329.676}{201.124}{x}%
\NativeArrowAt{339.668}{201.124}{11.955}{19.26}{\longmapsto}%
\TextAt{13}{362.25}{201.124}{\Character{11}}%
\TextAt{27}{369.696}{199.33}{p}%
\TextAt{13}{374.183}{201.124}{x\Character{59}}%
\TextAt{8}{79.226}{172.122}{of}%
\TextAt{8}{92.845}{172.122}{the}%
\TextAt{13}{113.291}{172.122}{K}%
\TextAt{27}{123.25}{170.329}{p}%
\TextAt{8}{127.736}{172.122}{-vector}%
\TextAt{8}{167.047}{172.122}{space}%
\TextAt{13}{198.611}{172.122}{L}%
\TextAt{27}{206.575}{170.329}{p}%
\TextAt{8}{211.062}{172.122}{,}%
\TextAt{8}{218.573}{172.122}{and}%
\TextAt{8}{241.624}{172.122}{as}%
\TextAt{8}{256.283}{172.122}{in}%
\TextAt{8}{270.225}{172.122}{the}%
\TextAt{8}{290.67}{172.122}{case}%
\TextAt{8}{315.73}{172.122}{of}%
\TextAt{8}{329.349}{172.122}{a}%
\TextAt{8}{339.394}{172.122}{field}%
\TextAt{8}{365.04}{172.122}{extension,}%
\TextAt{8}{420.419}{172.122}{we}%
\TextAt{8}{437.935}{172.122}{define}%
\TextAt{8}{472.033}{172.122}{the}%
\TextAt{8}{492.479}{172.122}{norm}%
\TextAt{8}{523.337}{172.122}{of}%
\TextAt{13}{79.226}{155.076}{\Character{11}}%
\TextAt{27}{86.672}{153.283}{p}%
\TextAt{8}{95.056}{155.076}{by}%
\TextAt{13}{250.687}{138.03}{N}%
\TextAt{21}{260.061}{136.171}{L}%
\TextAt{28}{265.82}{135.175}{p}%
\TextAt{3}{269.903}{136.171}{\Character{106}}%
\TextAt{21}{272.255}{136.171}{K}%
\TextAt{28}{279.396}{135.175}{p}%
\TextAt{9}{283.977}{138.03}{\Character{40}}%
\TextAt{13}{288.53}{138.03}{\Character{11}}%
\TextAt{27}{295.975}{136.237}{p}%
\TextAt{9}{300.462}{138.03}{\Character{41}}%
\TextAt{9}{308.317}{138.03}{\Character{61}}%
\TextAt{8}{320.761}{138.03}{det}%
\TextAt{9}{337.02}{138.03}{\Character{40}}%
\TextAt{20}{342.337}{138.031}{\Character{101}}%
\TextAt{13}{341.572}{138.03}{\Character{11}}%
\TextAt{27}{349.018}{136.237}{p}%
\TextAt{9}{353.504}{138.03}{\Character{41}}%
\TextAt{13}{358.057}{138.03}{\Character{58}}%
\TextAt{8}{79.226}{114.509}{In}%
\TextAt{8}{94.517}{114.509}{this}%
\TextAt{8}{118.01}{114.509}{way}%
\TextAt{8}{142.423}{114.509}{we}%
\TextAt{8}{160.321}{114.509}{obtain}%
\TextAt{8}{197.418}{114.509}{a}%
\TextAt{8}{207.844}{114.509}{homomorphism}%
\TextAt{13}{291.184}{114.509}{N}%
\TextAt{21}{300.558}{112.649}{L}%
\TextAt{28}{306.317}{111.653}{p}%
\TextAt{3}{310.399}{112.649}{\Character{106}}%
\TextAt{21}{312.752}{112.649}{K}%
\TextAt{28}{319.893}{111.653}{p}%
\TextAt{9}{328.937}{114.509}{\Character{58}}%
\TextAt{13}{336.651}{114.509}{L}%
\TextAt{3}{344.616}{118.847}{\Character{3}}%
\TextAt{27}{344.616}{111.553}{p}%
\NativeArrowAt{353.811}{114.509}{11.955}{19.26}{\longrightarrow}%
\TextAt{13}{377.535}{114.509}{K}%
\TextAt{3}{388.353}{118.847}{\Character{3}}%
\TextAt{27}{387.494}{111.553}{p}%
\TextAt{8}{393.085}{114.509}{.}%
\TextAt{8}{403.544}{114.509}{It}%
\TextAt{8}{416.886}{114.509}{induces}%
\TextAt{8}{459.232}{114.509}{a}%
\TextAt{8}{469.657}{114.509}{norm}%
\TextAt{8}{500.897}{114.509}{homo-}%
\TextAt{8}{79.226}{97.463}{morphism}%
\TextAt{13}{262.863}{80.417}{N}%
\TextAt{21}{272.237}{78.557}{L}%
\TextAt{3}{277.996}{78.557}{\Character{106}}%
\TextAt{21}{280.348}{78.557}{K}%
\TextAt{9}{291.872}{80.417}{\Character{58}}%
\TextAt{13}{298.444}{80.417}{I}%
\TextAt{21}{303.607}{78.624}{L}%
\NativeArrowAt{313.185}{80.417}{11.955}{19.26}{\longrightarrow}%
\TextAt{13}{335.767}{80.417}{I}%
\TextAt{21}{340.93}{78.624}{K}%
\TextAt{8}{297.211}{43.033}{337}%
\Destination{78.23}{809.815}{nav.page.337}%
\Destination{79.36}{776.155}{nav.print.369}%
\Destination{79.36}{357.255}{nav.print.370}%
\Destination{79.36}{657.055}{nav.number.6.2.1}%
}

% Page 350 (body page 338)
\Page{%
\TextAt{2}{79.226}{796.943}{VI.}%
\TextAt{2}{98.012}{796.943}{Global}%
\TextAt{2}{133.248}{796.943}{Class}%
\TextAt{2}{161.841}{796.943}{Field}%
\TextAt{2}{189.572}{796.943}{Theory}%
\Vector{q 0 G 0 g 0.299 w 7.226 18.729 m 460.769 18.729 l S Q}%
\TextAt{8}{79.226}{761.124}{between}%
\TextAt{8}{124.332}{761.124}{the}%
\TextAt{8}{144.403}{761.124}{idèle}%
\TextAt{8}{171.624}{761.124}{groups}%
\TextAt{13}{209.329}{761.124}{I}%
\TextAt{21}{214.492}{759.33}{L}%
\TextAt{9}{224.07}{761.124}{\Character{61}}%
\TextAt{20}{236.495}{770.09}{\Character{81}}%
\TextAt{3}{247.787}{766.935}{\Character{48}}%
\TextAt{27}{247.786}{757.637}{p}%
\TextAt{13}{254.266}{761.124}{L}%
\TextAt{3}{262.23}{765.462}{\Character{3}}%
\TextAt{27}{262.23}{758.168}{p}%
\TextAt{8}{270.775}{761.124}{and}%
\TextAt{13}{293.453}{761.124}{I}%
\TextAt{21}{298.616}{759.33}{K}%
\TextAt{9}{310.139}{761.124}{\Character{61}}%
\TextAt{20}{322.565}{770.09}{\Character{81}}%
\TextAt{3}{333.856}{766.935}{\Character{48}}%
\TextAt{27}{333.856}{757.637}{p}%
\TextAt{13}{340.335}{761.124}{K}%
\TextAt{3}{351.153}{765.462}{\Character{3}}%
\TextAt{27}{350.294}{758.168}{p}%
\TextAt{8}{355.885}{761.124}{.}%
\TextAt{8}{364.305}{761.124}{Explicitly}%
\TextAt{8}{417.707}{761.124}{the}%
\TextAt{8}{437.779}{761.124}{norm}%
\TextAt{8}{468.263}{761.124}{of}%
\TextAt{8}{481.508}{761.124}{an}%
\TextAt{8}{497.682}{761.124}{idèle}%
\TextAt{8}{524.903}{761.124}{is}%
\TextAt{8}{79.226}{744.078}{given}%
\TextAt{8}{109.796}{744.078}{by}%
\TextAt{8}{126.055}{744.078}{the}%
\TextAt{8}{146.211}{744.078}{following}%
\TextAt{8}{195.646}{744.078}{proposition.}%
\TextAt{12}{79.226}{718.065}{Proposition}%
\TextAt{12}{150.665}{718.065}{(2.2).}%
\TextAt{8}{184.874}{718.065}{If}%
\TextAt{13}{195.46}{718.065}{L}%
\TextAt{16}{206.745}{718.065}{\Character{106}}%
\TextAt{13}{213.387}{718.065}{K}%
\TextAt{8}{226.996}{718.065}{is}%
\TextAt{8}{237.654}{718.065}{a}%
\TextAt{8}{246.303}{718.065}{finite}%
\TextAt{8}{275.108}{718.065}{extension}%
\TextAt{8}{325.768}{718.065}{and}%
\TextAt{13}{347.424}{718.065}{\Character{11}}%
\TextAt{9}{358.267}{718.065}{\Character{61}}%
\TextAt{9}{370.677}{718.065}{\Character{40}}%
\TextAt{13}{375.245}{718.065}{\Character{11}}%
\TextAt{27}{382.69}{716.272}{P}%
\TextAt{9}{389.784}{718.065}{\Character{41}}%
\NativeMathAt{397.657}{718.065}{11.955}{\in}%
\TextAt{13}{408.948}{718.065}{I}%
\TextAt{21}{414.111}{716.272}{L}%
\TextAt{8}{420.368}{718.065}{,}%
\TextAt{8}{426.632}{718.065}{the}%
\TextAt{8}{445.682}{718.065}{local}%
\TextAt{8}{472.217}{718.065}{components}%
\TextAt{8}{79.226}{701.019}{of}%
\TextAt{8}{92.556}{701.019}{the}%
\TextAt{8}{112.713}{701.019}{idèle}%
\TextAt{13}{140.018}{701.019}{N}%
\TextAt{21}{149.392}{699.16}{L}%
\TextAt{3}{155.151}{699.16}{\Character{106}}%
\TextAt{21}{157.503}{699.16}{K}%
\TextAt{9}{165.706}{701.019}{\Character{40}}%
\TextAt{13}{170.258}{701.019}{\Character{11}}%
\TextAt{9}{177.78}{701.019}{\Character{41}}%
\TextAt{8}{186.23}{701.019}{are}%
\TextAt{8}{205.741}{701.019}{given}%
\TextAt{8}{236.31}{701.019}{by}%
\TextAt{13}{234.94}{672.018}{N}%
\TextAt{21}{244.314}{670.159}{L}%
\TextAt{3}{250.073}{670.159}{\Character{106}}%
\TextAt{21}{252.425}{670.159}{K}%
\TextAt{9}{260.628}{672.018}{\Character{40}}%
\TextAt{13}{265.18}{672.018}{\Character{11}}%
\TextAt{9}{272.702}{672.018}{\Character{41}}%
\TextAt{27}{277.254}{670.225}{p}%
\TextAt{9}{285.062}{672.018}{\Character{61}}%
\TextAt{20}{297.487}{683.376}{\Character{89}}%
\TextAt{27}{298.657}{657.473}{P}%
\TextAt{3}{305.252}{657.473}{\Character{106}}%
\TextAt{27}{307.605}{657.473}{p}%
\TextAt{13}{314.756}{672.018}{N}%
\TextAt{21}{324.13}{670.159}{L}%
\TextAt{28}{329.889}{668.768}{P}%
\TextAt{3}{336.086}{670.159}{\Character{106}}%
\TextAt{21}{338.438}{670.159}{K}%
\TextAt{28}{345.579}{669.162}{p}%
\TextAt{9}{350.16}{672.018}{\Character{40}}%
\TextAt{13}{354.713}{672.018}{\Character{11}}%
\TextAt{27}{362.158}{670.225}{P}%
\TextAt{9}{369.252}{672.018}{\Character{41}}%
\TextAt{13}{373.804}{672.018}{\Character{58}}%
\TextAt{19}{79.226}{614.434}{Proof.}%
\TextAt{8}{116}{614.434}{Putting}%
\TextAt{13}{160.082}{614.434}{\Character{11}}%
\TextAt{27}{167.527}{612.641}{p}%
\TextAt{9}{177.031}{614.434}{\Character{61}}%
\TextAt{9}{191.138}{614.434}{\Character{40}}%
\TextAt{13}{195.705}{614.434}{\Character{11}}%
\TextAt{27}{203.151}{612.641}{P}%
\TextAt{9}{210.244}{614.434}{\Character{41}}%
\TextAt{27}{214.796}{612.574}{P}%
\TextAt{3}{221.392}{612.574}{\Character{106}}%
\TextAt{27}{223.744}{612.574}{p}%
\NativeMathAt{233.248}{614.434}{11.955}{\in}%
\TextAt{13}{246.235}{614.434}{L}%
\TextAt{3}{254.199}{618.773}{\Character{3}}%
\TextAt{27}{254.199}{611.479}{p}%
\TextAt{8}{258.932}{614.434}{,}%
\TextAt{8}{267.325}{614.434}{the}%
\TextAt{13}{288.477}{614.434}{K}%
\TextAt{27}{298.436}{612.641}{p}%
\TextAt{8}{302.923}{614.434}{-automorphism}%
\TextAt{20}{386.05}{614.434}{\Character{101}}%
\TextAt{13}{385.285}{614.434}{\Character{11}}%
\TextAt{27}{392.73}{612.641}{p}%
\TextAt{9}{402.234}{614.434}{\Character{58}}%
\TextAt{13}{410.502}{614.434}{L}%
\TextAt{27}{418.467}{612.641}{p}%
\NativeArrowAt{427.971}{614.434}{11.955}{19.26}{\longrightarrow}%
\TextAt{13}{452.249}{614.434}{L}%
\TextAt{27}{460.213}{612.641}{p}%
\TextAt{8}{469.592}{614.434}{is}%
\TextAt{8}{482.352}{614.434}{the}%
\TextAt{8}{503.503}{614.434}{direct}%
\TextAt{8}{79.226}{597.388}{product}%
\TextAt{8}{123.126}{597.388}{of}%
\TextAt{8}{136.456}{597.388}{the}%
\TextAt{13}{156.612}{597.388}{K}%
\TextAt{27}{166.571}{595.595}{p}%
\TextAt{8}{171.058}{597.388}{-automorphisms}%
\TextAt{20}{257.805}{597.388}{\Character{101}}%
\TextAt{13}{257.04}{597.388}{\Character{11}}%
\TextAt{27}{264.485}{595.595}{P}%
\TextAt{9}{274.9}{597.388}{\Character{58}}%
\TextAt{13}{281.472}{597.388}{L}%
\TextAt{27}{289.437}{595.595}{P}%
\NativeArrowAt{299.851}{597.388}{11.955}{19.26}{\longrightarrow}%
\TextAt{13}{322.433}{597.388}{L}%
\TextAt{27}{330.398}{595.595}{P}%
\TextAt{8}{337.491}{597.388}{.}%
\TextAt{8}{345.939}{597.388}{Therefore}%
\TextAt{13}{168.197}{568.387}{N}%
\TextAt{21}{177.57}{566.527}{L}%
\TextAt{28}{183.329}{565.531}{p}%
\TextAt{3}{187.412}{566.527}{\Character{106}}%
\TextAt{21}{189.764}{566.527}{K}%
\TextAt{28}{196.906}{565.531}{p}%
\TextAt{9}{201.487}{568.387}{\Character{40}}%
\TextAt{13}{206.039}{568.387}{\Character{11}}%
\TextAt{27}{213.485}{566.594}{p}%
\TextAt{9}{217.971}{568.387}{\Character{41}}%
\TextAt{9}{225.826}{568.387}{\Character{61}}%
\TextAt{8}{238.27}{568.387}{det}%
\TextAt{9}{254.529}{568.387}{\Character{40}}%
\TextAt{20}{259.847}{568.387}{\Character{101}}%
\TextAt{13}{259.081}{568.387}{\Character{11}}%
\TextAt{27}{266.527}{566.594}{p}%
\TextAt{9}{271.014}{568.387}{\Character{41}}%
\TextAt{9}{278.869}{568.387}{\Character{61}}%
\TextAt{20}{291.313}{579.744}{\Character{89}}%
\TextAt{27}{292.482}{553.841}{P}%
\TextAt{3}{299.078}{553.841}{\Character{106}}%
\TextAt{27}{301.43}{553.841}{p}%
\TextAt{8}{308.581}{568.387}{det}%
\TextAt{9}{324.84}{568.387}{\Character{40}}%
\TextAt{20}{330.158}{568.387}{\Character{101}}%
\TextAt{13}{329.393}{568.387}{\Character{11}}%
\TextAt{27}{336.838}{566.594}{P}%
\TextAt{9}{343.932}{568.387}{\Character{41}}%
\TextAt{9}{351.787}{568.387}{\Character{61}}%
\TextAt{20}{364.23}{579.744}{\Character{89}}%
\TextAt{27}{365.4}{553.842}{P}%
\TextAt{3}{371.996}{553.841}{\Character{106}}%
\TextAt{27}{374.348}{553.842}{p}%
\TextAt{13}{381.499}{568.387}{N}%
\TextAt{21}{390.873}{566.527}{L}%
\TextAt{28}{396.632}{565.136}{P}%
\TextAt{3}{402.829}{566.527}{\Character{106}}%
\TextAt{21}{405.181}{566.527}{K}%
\TextAt{28}{412.323}{565.531}{p}%
\TextAt{9}{416.904}{568.387}{\Character{40}}%
\TextAt{13}{421.456}{568.387}{\Character{11}}%
\TextAt{27}{428.901}{566.594}{P}%
\TextAt{9}{435.995}{568.387}{\Character{41}}%
\TextAt{13}{440.547}{568.387}{\Character{58}}%
\Vector{q 0 G 0 g 0.399 w 452.334 -244.368 m 452.334 -236.298 l S Q}%
\Vector{q 0 G 0 g 0.399 w 452.534 -236.497 m 459.707 -236.497 l S Q}%
\Vector{q 0 G 0 g 0.399 w 452.534 -244.169 m 459.707 -244.169 l S Q}%
\Vector{q 0 G 0 g 0.399 w 459.906 -244.368 m 459.906 -236.298 l S Q}%
\TextAt{8}{97.159}{500.631}{The}%
\TextAt{8}{121.213}{500.631}{idèle}%
\TextAt{8}{148.518}{500.631}{norm}%
\TextAt{8}{179.088}{500.631}{enjoys}%
\TextAt{8}{214.595}{500.631}{the}%
\TextAt{8}{234.751}{500.631}{following}%
\TextAt{8}{284.186}{500.631}{properties.}%
\TextAt{12}{79.226}{474.619}{Proposition}%
\TextAt{12}{151.937}{474.619}{(2.3).}%
\TextAt{8}{187.253}{474.619}{(i)}%
\TextAt{8}{203.512}{474.619}{For}%
\TextAt{8}{224.481}{474.619}{a}%
\TextAt{8}{234.237}{474.619}{tower}%
\TextAt{8}{266.109}{474.619}{of}%
\TextAt{8}{279.439}{474.619}{fields}%
\TextAt{13}{309.411}{474.619}{K}%
\NativeMathAt{323.549}{474.619}{11.955}{\subseteq}%
\TextAt{13}{336.169}{474.619}{L}%
\NativeMathAt{347.454}{474.619}{11.955}{\subseteq}%
\TextAt{13}{360.073}{474.619}{M}%
\TextAt{8}{376.544}{474.619}{we}%
\TextAt{8}{393.772}{474.619}{have}%
\TextAt{13}{250.072}{445.618}{N}%
\TextAt{21}{259.445}{443.758}{M}%
\TextAt{3}{268.427}{443.758}{\Character{106}}%
\TextAt{21}{270.78}{443.758}{K}%
\TextAt{9}{282.303}{445.618}{\Character{61}}%
\TextAt{13}{294.729}{445.618}{N}%
\TextAt{21}{304.102}{443.758}{L}%
\TextAt{3}{309.861}{443.758}{\Character{106}}%
\TextAt{21}{312.214}{443.758}{K}%
\TextAt{16}{323.073}{445.618}{\Character{14}}%
\TextAt{13}{331.707}{445.618}{N}%
\TextAt{21}{341.081}{443.758}{M}%
\TextAt{3}{350.063}{443.758}{\Character{106}}%
\TextAt{21}{352.415}{443.758}{L}%
\TextAt{13}{358.672}{445.618}{\Character{58}}%
\TextAt{8}{79.226}{416.616}{(ii)}%
\TextAt{8}{99.522}{416.616}{If}%
\TextAt{13}{112}{416.616}{L}%
\TextAt{16}{124.623}{416.616}{\Character{106}}%
\TextAt{13}{132.603}{416.616}{K}%
\TextAt{8}{148.104}{416.616}{is}%
\TextAt{8}{160.653}{416.616}{embedded}%
\TextAt{8}{216.707}{416.616}{into}%
\TextAt{8}{241.235}{416.616}{the}%
\TextAt{8}{262.177}{416.616}{Galois}%
\TextAt{8}{298.875}{416.616}{extension}%
\TextAt{13}{351.427}{416.616}{M}%
\TextAt{16}{368.659}{416.616}{\Character{106}}%
\TextAt{13}{376.639}{416.616}{K}%
\TextAt{8}{392.14}{416.616}{and}%
\TextAt{8}{415.687}{416.616}{if}%
\TextAt{13}{427.197}{416.616}{G}%
\TextAt{9}{441.089}{416.616}{\Character{61}}%
\TextAt{13}{454.853}{416.616}{G}%
\TextAt{9}{464.087}{416.616}{\Character{40}}%
\TextAt{13}{468.639}{416.616}{M}%
\TextAt{16}{485.872}{416.616}{\Character{106}}%
\TextAt{13}{493.852}{416.616}{K}%
\TextAt{9}{504.669}{416.616}{\Character{41}}%
\TextAt{8}{513.904}{416.616}{and}%
\TextAt{13}{79.226}{399.57}{H}%
\TextAt{9}{93.17}{399.57}{\Character{61}}%
\TextAt{13}{105.595}{399.57}{G}%
\TextAt{9}{114.829}{399.57}{\Character{40}}%
\TextAt{13}{119.381}{399.57}{M}%
\TextAt{16}{135.275}{399.57}{\Character{106}}%
\TextAt{13}{141.917}{399.57}{L}%
\TextAt{9}{149.882}{399.57}{\Character{41}}%
\TextAt{8}{154.434}{399.57}{,}%
\TextAt{8}{161.583}{399.57}{then}%
\TextAt{8}{188.243}{399.57}{one}%
\TextAt{8}{209.703}{399.57}{has}%
\TextAt{8}{230.577}{399.57}{for}%
\TextAt{13}{248.461}{399.57}{\Character{11}}%
\NativeMathAt{259.304}{399.57}{11.955}{\in}%
\TextAt{13}{270.595}{399.57}{I}%
\TextAt{21}{275.758}{397.777}{L}%
\TextAt{8}{282.015}{399.57}{:}%
\TextAt{13}{252.456}{370.569}{N}%
\TextAt{21}{261.83}{368.71}{L}%
\TextAt{3}{267.589}{368.71}{\Character{106}}%
\TextAt{21}{269.942}{368.71}{K}%
\TextAt{9}{278.144}{370.569}{\Character{40}}%
\TextAt{13}{282.697}{370.569}{\Character{11}}%
\TextAt{9}{290.218}{370.569}{\Character{41}}%
\TextAt{9}{298.073}{370.569}{\Character{61}}%
\TextAt{20}{317.466}{381.927}{\Character{89}}%
\TextAt{21}{310.517}{356.024}{\Character{27}}%
\NativeMathAt{315.642}{356.024}{7.97}{\in}%
\TextAt{21}{321.288}{356.024}{G}%
\TextAt{17}{327.886}{356.024}{\Character{47}}%
\TextAt{21}{332.121}{356.024}{H}%
\TextAt{13}{341.683}{370.569}{\Character{27}}%
\TextAt{13}{348.749}{370.569}{\Character{11}\Character{58}}%
\TextAt{8}{79.226}{329.568}{(iii)}%
\TextAt{13}{103.017}{329.568}{N}%
\TextAt{21}{112.391}{327.709}{L}%
\TextAt{3}{118.15}{327.709}{\Character{106}}%
\TextAt{21}{120.502}{327.709}{K}%
\TextAt{9}{128.705}{329.568}{\Character{40}}%
\TextAt{13}{133.257}{329.568}{\Character{11}}%
\TextAt{9}{140.779}{329.568}{\Character{41}}%
\TextAt{9}{150.391}{329.568}{\Character{61}}%
\TextAt{13}{164.582}{329.568}{\Character{11}}%
\TextAt{17}{172.104}{333.907}{\Character{91}}%
\TextAt{21}{174.456}{333.907}{L}%
\TextAt{17}{180.215}{333.907}{\Character{58}}%
\TextAt{21}{182.567}{333.907}{K}%
\TextAt{17}{190.272}{333.907}{\Character{93}}%
\TextAt{8}{198.048}{329.568}{for}%
\TextAt{13}{216.961}{329.568}{\Character{11}}%
\NativeMathAt{229.556}{329.568}{11.955}{\in}%
\TextAt{13}{242.599}{329.568}{I}%
\TextAt{21}{247.762}{327.775}{K}%
\TextAt{8}{255.965}{329.568}{.}%
\TextAt{8}{267.498}{329.568}{(iv)}%
\TextAt{8}{290.966}{329.568}{The}%
\TextAt{8}{316.048}{329.568}{norm}%
\TextAt{8}{347.646}{329.568}{of}%
\TextAt{8}{362.004}{329.568}{the}%
\TextAt{8}{383.188}{329.568}{principal}%
\TextAt{8}{432.994}{329.568}{idèle}%
\TextAt{13}{461.328}{329.568}{x}%
\NativeMathAt{473.053}{329.568}{11.955}{\in}%
\TextAt{13}{486.096}{329.568}{L}%
\TextAt{3}{494.06}{333.907}{\Character{3}}%
\TextAt{8}{503.718}{329.568}{is}%
\TextAt{8}{516.51}{329.568}{the}%
\TextAt{8}{79.226}{312.522}{principal}%
\TextAt{8}{128.003}{312.522}{idèle}%
\TextAt{8}{155.309}{312.522}{of}%
\TextAt{13}{168.639}{312.522}{K}%
\TextAt{8}{183.354}{312.522}{defined}%
\TextAt{8}{223.667}{312.522}{by}%
\TextAt{8}{239.926}{312.522}{the}%
\TextAt{8}{260.082}{312.522}{usual}%
\TextAt{8}{290.711}{312.522}{norm}%
\TextAt{13}{321.281}{312.522}{N}%
\TextAt{21}{330.655}{310.663}{L}%
\TextAt{3}{336.414}{310.663}{\Character{106}}%
\TextAt{21}{338.766}{310.663}{K}%
\TextAt{9}{346.969}{312.522}{\Character{40}}%
\TextAt{13}{351.521}{312.522}{x}%
\TextAt{9}{358.173}{312.522}{\Character{41}}%
\TextAt{8}{362.725}{312.522}{.}%
\TextAt{8}{97.159}{286.51}{The}%
\TextAt{8}{120.228}{286.51}{proofs}%
\TextAt{8}{154.428}{286.51}{of}%
\TextAt{8}{166.774}{286.51}{(i),}%
\TextAt{8}{185.497}{286.51}{(ii),}%
\TextAt{8}{207.473}{286.51}{(iii)}%
\TextAt{8}{229.251}{286.51}{are}%
\TextAt{8}{247.778}{286.51}{literally}%
\TextAt{8}{290.047}{286.51}{the}%
\TextAt{8}{309.219}{286.51}{same}%
\TextAt{8}{337.561}{286.51}{as}%
\TextAt{8}{350.947}{286.51}{for}%
\TextAt{8}{367.848}{286.51}{the}%
\TextAt{8}{387.02}{286.51}{norm}%
\TextAt{8}{416.605}{286.51}{in}%
\TextAt{8}{429.273}{286.51}{a}%
\TextAt{8}{438.045}{286.51}{field}%
\TextAt{8}{462.417}{286.51}{extension}%
\TextAt{8}{513.199}{286.51}{(see}%
\TextAt{8}{79.226}{269.464}{chap.}%
\TextAt{8}{109.909}{269.464}{I,}%
\TextAt{8}{121.069}{269.464}{§2).}%
\TextAt{8}{145.539}{269.464}{(iv)}%
\TextAt{8}{167.77}{269.464}{follows}%
\TextAt{8}{205.996}{269.464}{from}%
\TextAt{8}{233.427}{269.464}{the}%
\TextAt{8}{253.374}{269.464}{fact}%
\TextAt{8}{276.25}{269.464}{that,}%
\TextAt{8}{304.703}{269.464}{once}%
\TextAt{8}{331.154}{269.464}{we}%
\TextAt{8}{348.172}{269.464}{identify}%
\TextAt{13}{390.559}{269.464}{L}%
\TextAt{27}{398.523}{267.671}{p}%
\TextAt{9}{406.331}{269.464}{\Character{61}}%
\TextAt{13}{418.757}{269.464}{L}%
\TextAt{16}{428.95}{269.464}{\Character{10}}%
\TextAt{21}{438.248}{267.671}{K}%
\TextAt{13}{448.679}{269.464}{K}%
\TextAt{27}{458.638}{267.671}{p}%
\TextAt{8}{466.813}{269.464}{(see}%
\TextAt{8}{490.072}{269.464}{chap.}%
\TextAt{8}{520.755}{269.464}{II,}%
\TextAt{8}{79.226}{252.418}{(8.3)),}%
\TextAt{8}{115.966}{252.418}{the}%
\TextAt{13}{136.889}{252.418}{K}%
\TextAt{27}{146.848}{250.625}{p}%
\TextAt{8}{151.335}{252.418}{-automorphism}%
\TextAt{20}{233.799}{252.418}{\Character{101}}%
\TextAt{13}{233.468}{252.418}{x}%
\TextAt{9}{244.748}{252.418}{\Character{58}}%
\TextAt{13}{252.626}{252.418}{L}%
\TextAt{27}{260.591}{250.625}{p}%
\NativeArrowAt{269.705}{252.418}{11.955}{19.26}{\longrightarrow}%
\TextAt{13}{293.593}{252.418}{L}%
\TextAt{27}{301.558}{250.625}{p}%
\TextAt{8}{306.044}{252.418}{,}%
\TextAt{13}{314.152}{252.418}{y}%
\NativeArrowAt{324.916}{252.418}{11.955}{11.955}{\mapsto}%
\TextAt{13}{341.498}{252.418}{xy}%
\TextAt{8}{354.287}{252.418}{,}%
\TextAt{8}{362.394}{252.418}{arises}%
\TextAt{8}{395.152}{252.418}{from}%
\TextAt{8}{423.559}{252.418}{the}%
\TextAt{13}{444.482}{252.418}{K}%
\TextAt{8}{455.3}{252.418}{-automorphism}%
\TextAt{13}{79.226}{235.372}{x}%
\TextAt{9}{89.199}{235.372}{\Character{58}}%
\TextAt{13}{95.772}{235.372}{L}%
\NativeArrowAt{107.057}{235.372}{11.955}{19.26}{\longrightarrow}%
\TextAt{13}{129.639}{235.372}{L}%
\TextAt{8}{141.501}{235.372}{by}%
\TextAt{8}{157.76}{235.372}{tensoring}%
\TextAt{8}{208.557}{235.372}{with}%
\TextAt{13}{235.217}{235.372}{K}%
\TextAt{27}{245.176}{233.579}{p}%
\TextAt{8}{249.663}{235.372}{.}%
\TextAt{8}{258.112}{235.372}{Hence}%
\TextAt{8}{292.889}{235.372}{det}%
\TextAt{9}{309.148}{235.372}{\Character{40}}%
\TextAt{20}{314.031}{235.372}{\Character{101}}%
\TextAt{13}{313.701}{235.372}{x}%
\TextAt{9}{320.353}{235.372}{\Character{41}}%
\TextAt{9}{328.208}{235.372}{\Character{61}}%
\TextAt{8}{340.651}{235.372}{det}%
\TextAt{9}{356.91}{235.372}{\Character{40}}%
\TextAt{13}{361.463}{235.372}{x}%
\TextAt{9}{368.115}{235.372}{\Character{41}}%
\TextAt{8}{372.667}{235.372}{.}%
\TextAt{12}{79.226}{212.349}{Remark.}%
\TextAt{8}{132.975}{212.349}{For}%
\TextAt{8}{153.536}{212.349}{fundamental}%
\TextAt{8}{220.771}{212.349}{as}%
\TextAt{8}{234.733}{212.349}{well}%
\TextAt{8}{258.056}{212.349}{as}%
\TextAt{8}{272.018}{212.349}{practical}%
\TextAt{8}{319.741}{212.349}{reasons,}%
\TextAt{8}{363.769}{212.349}{it}%
\TextAt{8}{375.065}{212.349}{is}%
\TextAt{8}{386.421}{212.349}{convenient}%
\TextAt{8}{443.899}{212.349}{to}%
\TextAt{8}{457.802}{212.349}{adopt}%
\TextAt{8}{490.569}{212.349}{a}%
\TextAt{8}{499.917}{212.349}{formal}%
\TextAt{8}{79.226}{195.303}{point}%
\TextAt{8}{109.334}{195.303}{of}%
\TextAt{8}{122.203}{195.303}{view}%
\TextAt{8}{148.724}{195.303}{for}%
\TextAt{8}{166.147}{195.303}{the}%
\TextAt{8}{185.842}{195.303}{above}%
\TextAt{8}{218.557}{195.303}{considerations}%
\TextAt{8}{294.321}{195.303}{which}%
\TextAt{8}{327.346}{195.303}{allows}%
\TextAt{8}{361.746}{195.303}{us}%
\TextAt{8}{376.3}{195.303}{to}%
\TextAt{8}{390.149}{195.303}{avoid}%
\TextAt{8}{420.592}{195.303}{the}%
\TextAt{8}{440.287}{195.303}{constant}%
\TextAt{8}{487.048}{195.303}{back}%
\TextAt{8}{513.904}{195.303}{and}%
\TextAt{8}{79.226}{178.257}{forth}%
\TextAt{8}{107.801}{178.257}{between}%
\TextAt{8}{152.624}{178.257}{idèles}%
\TextAt{8}{184.176}{178.257}{and}%
\TextAt{8}{206.57}{178.257}{their}%
\TextAt{8}{234.165}{178.257}{components.}%
\TextAt{8}{303.044}{178.257}{This}%
\TextAt{8}{329.395}{178.257}{point}%
\TextAt{8}{359.596}{178.257}{of}%
\TextAt{8}{372.558}{178.257}{view}%
\TextAt{8}{399.173}{178.257}{is}%
\TextAt{8}{410.568}{178.257}{based}%
\TextAt{8}{442.778}{178.257}{on}%
\TextAt{8}{458.669}{178.257}{identifying}%
\TextAt{8}{516.51}{178.257}{the}%
\TextAt{8}{79.226}{161.211}{ring}%
\TextAt{8}{103.292}{161.211}{of}%
\TextAt{8}{116.622}{161.211}{adèles}%
\TextAt{18}{151.148}{161.211}{A}%
\TextAt{21}{159.783}{159.417}{L}%
\TextAt{8}{169.937}{161.211}{of}%
\TextAt{13}{183.267}{161.211}{L}%
\TextAt{8}{195.129}{161.211}{as}%
\TextAt{18}{265.245}{144.165}{A}%
\TextAt{21}{273.879}{142.371}{L}%
\TextAt{9}{283.457}{144.165}{\Character{61}}%
\TextAt{18}{295.883}{144.165}{A}%
\TextAt{21}{304.517}{142.371}{K}%
\TextAt{16}{315.377}{144.165}{\Character{10}}%
\TextAt{21}{324.675}{142.371}{K}%
\TextAt{13}{335.534}{144.165}{L\Character{59}}%
\TextAt{8}{79.226}{120.643}{which}%
\TextAt{8}{112.713}{120.643}{results}%
\TextAt{8}{149.905}{120.643}{from}%
\TextAt{8}{177.545}{120.643}{the}%
\TextAt{8}{197.702}{120.643}{canonical}%
\TextAt{8}{249.085}{120.643}{isomorphisms}%
\TextAt{8}{322.119}{120.643}{(see}%
\TextAt{8}{345.587}{120.643}{chap.}%
\TextAt{8}{376.479}{120.643}{II,}%
\TextAt{8}{392.392}{120.643}{(8.3))}%
\TextAt{13}{179.974}{91.642}{K}%
\TextAt{27}{189.933}{89.849}{p}%
\TextAt{16}{197.077}{91.642}{\Character{10}}%
\TextAt{21}{206.375}{89.849}{K}%
\TextAt{13}{217.235}{91.642}{L}%
\NativeMathAt{228.52}{91.642}{11.955}{\cong}%
\TextAt{13}{241.139}{91.642}{L}%
\TextAt{27}{249.104}{89.849}{p}%
\TextAt{9}{256.912}{91.642}{\Character{61}}%
\TextAt{20}{269.337}{102.999}{\Character{89}}%
\TextAt{27}{270.507}{77.096}{P}%
\TextAt{3}{277.102}{77.096}{\Character{106}}%
\TextAt{27}{279.455}{77.096}{p}%
\TextAt{13}{286.606}{91.642}{L}%
\TextAt{27}{294.57}{89.849}{P}%
\TextAt{13}{301.664}{91.642}{\Character{59}}%
\TextAt{13}{330.799}{91.642}{\Character{11}}%
\TextAt{27}{338.264}{89.849}{p}%
\TextAt{16}{345.407}{91.642}{\Character{10}}%
\TextAt{13}{357.363}{91.642}{a}%
\NativeArrowAt{366.828}{91.642}{11.955}{19.26}{\longmapsto}%
\TextAt{13}{389.41}{91.642}{\Character{11}}%
\TextAt{27}{396.856}{89.849}{p}%
\TextAt{9}{401.343}{91.642}{\Character{40}}%
\TextAt{13}{405.895}{91.642}{\Character{28}}%
\TextAt{27}{410.979}{89.849}{P}%
\TextAt{13}{418.072}{91.642}{a}%
\TextAt{9}{424.217}{91.642}{\Character{41}}%
\TextAt{13}{428.77}{91.642}{\Character{58}}%
\TextAt{8}{297.211}{43.033}{338}%
\Destination{78.23}{809.815}{nav.page.338}%
\Destination{79.36}{344.555}{nav.print.371}%
\Destination{79.36}{732.055}{nav.number.6.2.2}%
\Destination{79.36}{488.655}{nav.number.6.2.3}%
\Hotlink{109.409}{267.255}{5.22}{11.047}{nav.chapter.1}%
\Hotlink{520.255}{267.255}{9.751}{11.047}{nav.chapter.2}%
\Hotlink{126.248}{267.255}{6.858}{11.047}{nav.section.1.2}%
\Hotlink{375.979}{118.455}{9.763}{11.047}{nav.chapter.2}%
\Hotlink{78.726}{250.255}{25.078}{11.047}{nav.number.6.8.3}%
\Hotlink{391.892}{118.455}{25.078}{11.047}{nav.number.2.8.3}%
}

% Page 351 (body page 339)
\Page{%
\TextAt{2}{79.226}{796.943}{VI.}%
\TextAt{2}{98.012}{796.943}{Global}%
\TextAt{2}{133.248}{796.943}{Class}%
\TextAt{2}{161.841}{796.943}{Field}%
\TextAt{2}{189.572}{796.943}{Theory}%
\Vector{q 0 G 0 g 0.299 w 7.226 18.729 m 460.769 18.729 l S Q}%
\TextAt{8}{79.226}{761.124}{Here}%
\TextAt{13}{106.855}{761.124}{\Character{28}}%
\TextAt{27}{111.938}{759.33}{P}%
\TextAt{8}{122.929}{761.124}{denotes}%
\TextAt{8}{165.262}{761.124}{the}%
\TextAt{8}{185.419}{761.124}{canonical}%
\TextAt{8}{236.802}{761.124}{embedding}%
\TextAt{13}{295.98}{761.124}{\Character{28}}%
\TextAt{27}{301.064}{759.33}{P}%
\TextAt{9}{311.478}{761.124}{\Character{58}}%
\TextAt{13}{318.051}{761.124}{L}%
\NativeArrowAt{329.336}{761.124}{11.955}{19.26}{\longrightarrow}%
\TextAt{13}{351.918}{761.124}{L}%
\TextAt{27}{359.883}{759.33}{P}%
\TextAt{8}{366.976}{761.124}{.}%
\TextAt{8}{97.159}{741.089}{In}%
\TextAt{8}{112.305}{741.089}{this}%
\TextAt{8}{135.652}{741.089}{way}%
\TextAt{8}{159.92}{741.089}{the}%
\TextAt{8}{180.602}{741.089}{inclusion}%
\TextAt{8}{229.963}{741.089}{by}%
\TextAt{8}{246.747}{741.089}{components}%
\TextAt{13}{311.723}{741.089}{I}%
\TextAt{21}{316.886}{739.296}{K}%
\NativeMathAt{329.304}{741.089}{11.955}{\subseteq}%
\TextAt{13}{342.817}{741.089}{I}%
\TextAt{21}{347.98}{739.296}{L}%
\TextAt{8}{358.66}{741.089}{is}%
\TextAt{8}{370.949}{741.089}{simply}%
\TextAt{8}{408.929}{741.089}{given}%
\TextAt{8}{440.023}{741.089}{by}%
\TextAt{8}{456.807}{741.089}{the}%
\TextAt{8}{477.489}{741.089}{embedding}%
\TextAt{18}{79.226}{724.043}{A}%
\TextAt{21}{87.86}{722.25}{K}%
\NativeArrowAt{100.037}{724.043}{11.955}{19.26}{\longrightarrow}%
\TextAt{18}{123.273}{724.043}{A}%
\TextAt{21}{131.907}{722.25}{L}%
\TextAt{8}{138.164}{724.043}{,}%
\TextAt{13}{145.793}{724.043}{\Character{11}}%
\NativeArrowAt{157.289}{724.043}{11.955}{11.955}{\mapsto}%
\TextAt{13}{173.219}{724.043}{\Character{11}}%
\TextAt{16}{183.658}{724.043}{\Character{10}}%
\TextAt{9}{195.875}{724.043}{1}%
\TextAt{8}{201.728}{724.043}{,}%
\TextAt{8}{209.357}{724.043}{induced}%
\TextAt{8}{253.305}{724.043}{by}%
\TextAt{13}{269.947}{724.043}{K}%
\NativeMathAt{284.739}{724.043}{11.955}{\subseteq}%
\TextAt{13}{298.012}{724.043}{L}%
\TextAt{8}{305.977}{724.043}{.}%
\ErrataAt{315.55}{724.043}{219.45}{A $K$-isomorphism $\sigma:L\to\sigma L$ then yields}%
\TextAt{8}{79.226}{706.997}{the}%
\TextAt{8}{99.383}{706.997}{isomorphism}%
\TextAt{13}{200.201}{689.951}{\Character{27}}%
\TextAt{9}{210.604}{689.951}{\Character{58}}%
\TextAt{18}{217.177}{689.951}{A}%
\TextAt{21}{225.811}{688.158}{L}%
\TextAt{9}{235.389}{689.951}{\Character{61}}%
\TextAt{18}{247.814}{689.951}{A}%
\TextAt{21}{256.449}{688.158}{K}%
\TextAt{16}{267.308}{689.951}{\Character{10}}%
\TextAt{21}{276.607}{688.158}{K}%
\TextAt{13}{287.466}{689.951}{L}%
\NativeArrowAt{298.751}{689.951}{11.955}{19.26}{\longrightarrow}%
\TextAt{18}{321.333}{689.951}{A}%
\TextAt{21}{329.968}{688.158}{K}%
\TextAt{16}{340.827}{689.951}{\Character{10}}%
\TextAt{21}{350.125}{688.158}{K}%
\TextAt{13}{360.985}{689.951}{\Character{27}}%
\TextAt{13}{368.05}{689.951}{L}%
\TextAt{9}{379.352}{689.951}{\Character{61}}%
\TextAt{18}{391.778}{689.951}{A}%
\TextAt{21}{400.412}{688.158}{\Character{27}}%
\TextAt{21}{405.529}{688.158}{L}%
\TextAt{8}{79.226}{666.82}{via}%
\TextAt{13}{98.414}{666.82}{\Character{27}}%
\TextAt{9}{105.497}{666.82}{\Character{40}}%
\TextAt{13}{110.049}{666.82}{\Character{11}}%
\TextAt{16}{120.227}{666.82}{\Character{10}}%
\TextAt{13}{132.182}{666.82}{a}%
\TextAt{9}{138.327}{666.82}{\Character{41}}%
\TextAt{9}{146.182}{666.82}{\Character{61}}%
\TextAt{13}{158.626}{666.82}{\Character{11}}%
\TextAt{16}{168.804}{666.82}{\Character{10}}%
\TextAt{13}{180.76}{666.82}{\Character{27}}%
\TextAt{13}{187.826}{666.82}{a}%
\TextAt{8}{193.987}{666.82}{,}%
\TextAt{8}{201.136}{666.82}{and}%
\TextAt{8}{223.899}{666.82}{the}%
\TextAt{8}{244.055}{666.82}{norm}%
\TextAt{8}{274.625}{666.82}{of}%
\TextAt{8}{287.955}{666.82}{an}%
\TextAt{13}{304.214}{666.82}{L}%
\TextAt{8}{312.178}{666.82}{-idèle}%
\TextAt{13}{343.381}{666.82}{\Character{11}}%
\NativeMathAt{354.224}{666.82}{11.955}{\in}%
\TextAt{18}{365.515}{666.82}{A}%
\TextAt{3}{374.149}{671.159}{\Character{3}}%
\TextAt{21}{374.149}{663.8}{L}%
\TextAt{8}{384.304}{666.82}{is}%
\TextAt{8}{396.067}{666.82}{simply}%
\TextAt{8}{433.523}{666.82}{the}%
\TextAt{8}{453.679}{666.82}{determinant}%
\TextAt{13}{260.525}{638.732}{N}%
\TextAt{21}{269.899}{636.872}{L}%
\TextAt{3}{275.658}{636.872}{\Character{106}}%
\TextAt{21}{278.01}{636.872}{K}%
\TextAt{9}{286.213}{638.732}{\Character{40}}%
\TextAt{13}{290.765}{638.732}{\Character{11}}%
\TextAt{9}{298.287}{638.732}{\Character{41}}%
\TextAt{9}{306.142}{638.732}{\Character{61}}%
\TextAt{8}{318.585}{638.732}{det}%
\TextAt{25}{320.315}{631.129}{A}%
\TextAt{22}{326.072}{629.79}{K}%
\TextAt{9}{334.844}{638.732}{\Character{40}}%
\TextAt{13}{339.397}{638.732}{\Character{11}}%
\TextAt{9}{346.918}{638.732}{\Character{41}}%
\TextAt{8}{79.226}{608.139}{of}%
\TextAt{8}{93.701}{608.139}{the}%
\TextAt{8}{115.003}{608.139}{endomorphism}%
\TextAt{13}{194.909}{608.139}{\Character{11}}%
\TextAt{9}{207.703}{608.139}{\Character{58}}%
\TextAt{18}{216.227}{608.139}{A}%
\TextAt{21}{224.862}{606.346}{L}%
\NativeArrowAt{236.391}{608.139}{11.955}{19.26}{\longrightarrow}%
\TextAt{18}{260.925}{608.139}{A}%
\TextAt{21}{269.559}{606.346}{L}%
\TextAt{8}{280.859}{608.139}{which}%
\TextAt{13}{315.491}{608.139}{\Character{11}}%
\TextAt{8}{328.055}{608.139}{induces}%
\TextAt{8}{370.876}{608.139}{on}%
\TextAt{8}{388.28}{608.139}{the}%
\TextAt{8}{409.582}{608.139}{finite}%
\TextAt{18}{440.639}{608.139}{A}%
\TextAt{21}{449.273}{606.346}{K}%
\TextAt{8}{457.476}{608.139}{-algebra}%
\TextAt{18}{503.501}{608.139}{A}%
\TextAt{21}{512.135}{606.346}{L}%
\TextAt{9}{523.665}{608.139}{\Character{61}}%
\TextAt{18}{79.226}{591.093}{A}%
\TextAt{21}{87.86}{589.3}{K}%
\TextAt{16}{98.72}{591.093}{\Character{10}}%
\TextAt{21}{108.018}{589.3}{K}%
\TextAt{13}{118.878}{591.093}{L}%
\TextAt{8}{126.842}{591.093}{.}%
\TextAt{8}{97.159}{568.2}{Here}%
\TextAt{8}{124.787}{568.2}{are}%
\TextAt{8}{144.298}{568.2}{consequences}%
\TextAt{8}{214.977}{568.2}{of}%
\TextAt{8}{228.307}{568.2}{the}%
\TextAt{8}{248.464}{568.2}{preceding}%
\TextAt{8}{301.138}{568.2}{investigations}%
\TextAt{8}{374.447}{568.2}{for}%
\TextAt{8}{392.332}{568.2}{the}%
\TextAt{8}{412.488}{568.2}{idèle}%
\TextAt{8}{439.794}{568.2}{class}%
\TextAt{8}{467.231}{568.2}{groups.}%
\TextAt{12}{79.226}{542.448}{Proposition}%
\TextAt{12}{152.857}{542.448}{(2.4).}%
\TextAt{8}{188.972}{542.448}{If}%
\TextAt{13}{201.463}{542.448}{L}%
\TextAt{16}{214.11}{542.448}{\Character{106}}%
\TextAt{13}{222.114}{542.448}{K}%
\TextAt{8}{237.628}{542.448}{is}%
\TextAt{8}{250.191}{542.448}{a}%
\TextAt{8}{260.746}{542.448}{finite}%
\TextAt{8}{291.457}{542.448}{extension,}%
\TextAt{8}{347.473}{542.448}{then}%
\TextAt{8}{374.932}{542.448}{the}%
\TextAt{8}{395.888}{542.448}{homomorphism}%
\TextAt{13}{479.357}{542.448}{I}%
\TextAt{21}{484.52}{540.655}{K}%
\NativeArrowAt{497.406}{542.448}{11.955}{19.26}{\longrightarrow}%
\TextAt{13}{521.349}{542.448}{I}%
\TextAt{21}{526.512}{540.655}{L}%
\TextAt{8}{79.226}{525.403}{induces}%
\TextAt{8}{120.902}{525.403}{an}%
\TextAt{8}{137.161}{525.403}{injection}%
\TextAt{8}{184.958}{525.403}{of}%
\TextAt{8}{198.288}{525.403}{idèle}%
\TextAt{8}{225.593}{525.403}{class}%
\TextAt{8}{253.03}{525.403}{groups}%
\TextAt{13}{226.642}{497.314}{C}%
\TextAt{21}{235.017}{495.521}{K}%
\NativeArrowAt{246.54}{497.314}{11.955}{19.26}{\longrightarrow}%
\TextAt{13}{269.122}{497.314}{C}%
\TextAt{21}{277.497}{495.521}{L}%
\TextAt{13}{283.755}{497.314}{\Character{59}}%
\TextAt{13}{312.89}{497.314}{\Character{11}K}%
\TextAt{3}{331.248}{502.25}{\Character{3}}%
\NativeArrowAt{339.302}{497.314}{11.955}{19.26}{\longmapsto}%
\TextAt{13}{361.884}{497.314}{\Character{11}L}%
\TextAt{3}{377.37}{502.25}{\Character{3}}%
\TextAt{13}{382.102}{497.314}{\Character{58}}%
\TextAt{19}{79.226}{451.814}{Proof.}%
\TextAt{8}{116}{451.814}{The}%
\TextAt{8}{140.304}{451.814}{injection}%
\TextAt{13}{188.35}{451.814}{I}%
\TextAt{21}{193.513}{450.021}{K}%
\NativeArrowAt{205.462}{451.814}{11.955}{19.26}{\longrightarrow}%
\TextAt{13}{228.47}{451.814}{I}%
\TextAt{21}{233.633}{450.021}{L}%
\TextAt{8}{244.037}{451.814}{clearly}%
\TextAt{8}{281.682}{451.814}{maps}%
\TextAt{13}{312.561}{451.814}{K}%
\TextAt{3}{323.379}{456.153}{\Character{3}}%
\TextAt{8}{332.258}{451.814}{into}%
\TextAt{13}{356.251}{451.814}{L}%
\TextAt{3}{364.216}{456.153}{\Character{3}}%
\TextAt{8}{368.948}{451.814}{.}%
\TextAt{8}{378.145}{451.814}{For}%
\TextAt{8}{399.364}{451.814}{the}%
\TextAt{8}{419.771}{451.814}{injectivity,}%
\TextAt{8}{477.635}{451.814}{we}%
\TextAt{8}{495.112}{451.814}{have}%
\TextAt{8}{522.357}{451.814}{to}%
\TextAt{8}{79.226}{434.768}{show}%
\TextAt{8}{108.816}{434.768}{that}%
\TextAt{13}{134.772}{434.768}{I}%
\TextAt{21}{139.935}{432.975}{K}%
\TextAt{16}{151.195}{434.768}{\Character{92}}%
\TextAt{13}{162.222}{434.768}{L}%
\TextAt{3}{170.186}{439.107}{\Character{3}}%
\TextAt{9}{179.24}{434.768}{\Character{61}}%
\TextAt{13}{192.665}{434.768}{K}%
\TextAt{3}{203.483}{439.107}{\Character{3}}%
\TextAt{8}{208.215}{434.768}{.}%
\TextAt{8}{218.425}{434.768}{Let}%
\TextAt{13}{239.981}{434.768}{M}%
\TextAt{16}{256.876}{434.768}{\Character{106}}%
\TextAt{13}{264.518}{434.768}{K}%
\TextAt{8}{279.82}{434.768}{be}%
\TextAt{8}{296.331}{434.768}{a}%
\TextAt{8}{306.674}{434.768}{finite}%
\TextAt{8}{337.173}{434.768}{Galois}%
\TextAt{8}{373.673}{434.768}{extension}%
\TextAt{8}{426.026}{434.768}{with}%
\TextAt{8}{453.273}{434.768}{Galois}%
\TextAt{8}{489.773}{434.768}{group}%
\TextAt{13}{523.536}{434.768}{G}%
\TextAt{8}{79.226}{417.722}{containing}%
\TextAt{13}{136.145}{417.722}{L}%
\TextAt{8}{144.109}{417.722}{.}%
\TextAt{8}{152.558}{417.722}{Then}%
\TextAt{8}{183.115}{417.722}{we}%
\TextAt{8}{200.342}{417.722}{have}%
\TextAt{13}{227.337}{417.722}{I}%
\TextAt{21}{232.5}{415.929}{K}%
\NativeMathAt{244.024}{417.722}{11.955}{\subseteq}%
\TextAt{13}{256.643}{417.722}{I}%
\TextAt{21}{261.806}{415.929}{L}%
\NativeMathAt{271.384}{417.722}{11.955}{\subseteq}%
\TextAt{13}{284.003}{417.722}{I}%
\TextAt{21}{289.166}{415.929}{M}%
\TextAt{8}{298.646}{417.722}{,}%
\TextAt{8}{305.795}{417.722}{and}%
\TextAt{13}{138.932}{389.634}{I}%
\TextAt{21}{144.095}{387.841}{K}%
\TextAt{16}{154.955}{389.634}{\Character{92}}%
\TextAt{13}{165.582}{389.634}{L}%
\TextAt{3}{173.546}{394.57}{\Character{3}}%
\NativeMathAt{181.599}{389.634}{11.955}{\subseteq}%
\TextAt{13}{194.219}{389.634}{I}%
\TextAt{21}{199.381}{387.841}{K}%
\TextAt{16}{210.241}{389.634}{\Character{92}}%
\TextAt{13}{220.868}{389.634}{M}%
\TextAt{3}{233.441}{394.57}{\Character{3}}%
\NativeMathAt{241.494}{389.634}{11.955}{\subseteq}%
\TextAt{9}{254.114}{389.634}{\Character{40}}%
\TextAt{13}{258.666}{389.634}{I}%
\TextAt{21}{263.829}{387.841}{K}%
\TextAt{16}{274.688}{389.634}{\Character{92}}%
\TextAt{13}{285.315}{389.634}{M}%
\TextAt{3}{297.889}{394.57}{\Character{3}}%
\TextAt{9}{302.621}{389.634}{\Character{41}}%
\TextAt{21}{307.173}{394.57}{G}%
\TextAt{9}{317.591}{389.634}{\Character{61}}%
\TextAt{13}{330.016}{389.634}{I}%
\TextAt{21}{335.179}{387.841}{K}%
\TextAt{16}{346.039}{389.634}{\Character{92}}%
\TextAt{13}{356.665}{389.634}{M}%
\TextAt{3}{369.239}{394.57}{\Character{3}}%
\TextAt{21}{373.473}{394.57}{G}%
\TextAt{9}{383.891}{389.634}{\Character{61}}%
\TextAt{13}{396.316}{389.634}{I}%
\TextAt{21}{401.479}{387.841}{K}%
\TextAt{16}{412.339}{389.634}{\Character{92}}%
\TextAt{13}{422.965}{389.634}{K}%
\TextAt{3}{433.783}{394.57}{\Character{3}}%
\TextAt{9}{441.836}{389.634}{\Character{61}}%
\TextAt{13}{454.262}{389.634}{K}%
\TextAt{3}{465.079}{394.57}{\Character{3}}%
\TextAt{13}{469.812}{389.634}{\Character{58}}%
\Vector{q 0 G 0 g 0.399 w 452.334 -412.455 m 452.334 -404.385 l S Q}%
\Vector{q 0 G 0 g 0.399 w 452.534 -404.584 m 459.707 -404.584 l S Q}%
\Vector{q 0 G 0 g 0.399 w 452.534 -412.255 m 459.707 -412.255 l S Q}%
\Vector{q 0 G 0 g 0.399 w 459.906 -412.455 m 459.906 -404.385 l S Q}%
\TextAt{8}{97.159}{332.805}{Via}%
\TextAt{8}{119.143}{332.805}{the}%
\TextAt{8}{139.501}{332.805}{embedding}%
\TextAt{13}{198.881}{332.805}{C}%
\TextAt{21}{207.256}{331.012}{K}%
\NativeArrowAt{219.124}{332.805}{11.955}{19.26}{\longrightarrow}%
\TextAt{13}{242.049}{332.805}{C}%
\TextAt{21}{250.424}{331.012}{L}%
\TextAt{8}{256.682}{332.805}{,}%
\TextAt{8}{264.083}{332.805}{the}%
\TextAt{8}{284.441}{332.805}{idèle}%
\TextAt{8}{311.949}{332.805}{class}%
\TextAt{8}{339.588}{332.805}{group}%
\TextAt{13}{372.965}{332.805}{C}%
\TextAt{21}{381.34}{331.012}{K}%
\TextAt{8}{393.642}{332.805}{becomes}%
\TextAt{8}{440.397}{332.805}{a}%
\TextAt{8}{450.354}{332.805}{subgroup}%
\TextAt{8}{501.354}{332.805}{of}%
\TextAt{13}{514.885}{332.805}{C}%
\TextAt{21}{523.26}{331.012}{L}%
\TextAt{8}{529.518}{332.805}{:}%
\TextAt{8}{79.226}{315.759}{an}%
\TextAt{8}{95.251}{315.759}{element}%
\TextAt{13}{138.258}{315.759}{\Character{11}L}%
\TextAt{3}{153.744}{320.097}{\Character{3}}%
\NativeMathAt{161.797}{315.759}{11.955}{\in}%
\TextAt{13}{173.088}{315.759}{C}%
\TextAt{21}{181.463}{313.966}{L}%
\TextAt{8}{191.384}{315.759}{(}%
\TextAt{13}{195.938}{315.759}{\Character{11}}%
\NativeMathAt{206.781}{315.759}{11.955}{\in}%
\TextAt{13}{218.072}{315.759}{I}%
\TextAt{21}{223.235}{313.966}{L}%
\TextAt{8}{229.492}{315.759}{)}%
\TextAt{8}{237.71}{315.759}{lies}%
\TextAt{8}{257.692}{315.759}{in}%
\TextAt{13}{271.11}{315.759}{C}%
\TextAt{21}{279.485}{313.966}{K}%
\TextAt{8}{291.351}{315.759}{if}%
\TextAt{8}{301.84}{315.759}{and}%
\TextAt{8}{324.368}{315.759}{only}%
\TextAt{8}{349.825}{315.759}{if}%
\TextAt{8}{360.315}{315.759}{the}%
\TextAt{8}{380.237}{315.759}{class}%
\TextAt{13}{407.439}{315.759}{\Character{11}L}%
\TextAt{3}{422.925}{320.097}{\Character{3}}%
\TextAt{8}{431.321}{315.759}{has}%
\TextAt{8}{451.96}{315.759}{a}%
\TextAt{8}{461.481}{315.759}{representative}%
\TextAt{13}{79.226}{298.713}{\Character{11}}%
\TextAt{3}{86.748}{303.051}{\Character{48}}%
\TextAt{8}{93.44}{298.713}{in}%
\TextAt{13}{107.093}{298.713}{I}%
\TextAt{21}{112.256}{296.92}{K}%
\TextAt{8}{120.459}{298.713}{.}%
\TextAt{8}{128.907}{298.713}{It}%
\TextAt{8}{141.58}{298.713}{is}%
\TextAt{8}{153.343}{298.713}{important}%
\TextAt{8}{208.636}{298.713}{to}%
\TextAt{8}{222.946}{298.713}{know}%
\TextAt{8}{253.516}{298.713}{that}%
\TextAt{8}{278.885}{298.713}{we}%
\TextAt{8}{296.112}{298.713}{have}%
\TextAt{8}{323.107}{298.713}{Galois}%
\TextAt{8}{359.02}{298.713}{descent}%
\TextAt{8}{400.373}{298.713}{for}%
\TextAt{8}{418.258}{298.713}{the}%
\TextAt{8}{438.414}{298.713}{idèle}%
\TextAt{8}{465.72}{298.713}{class}%
\TextAt{8}{493.157}{298.713}{group:}%
\TextAt{12}{79.226}{272.961}{Proposition}%
\TextAt{12}{153.816}{272.961}{(2.5).}%
\TextAt{8}{190.765}{272.961}{If}%
\TextAt{13}{204.09}{272.961}{L}%
\TextAt{16}{218.159}{272.961}{\Character{106}}%
\TextAt{13}{227.584}{272.961}{K}%
\TextAt{8}{243.932}{272.961}{is}%
\TextAt{8}{257.33}{272.961}{a}%
\TextAt{8}{268.718}{272.961}{Galois}%
\TextAt{8}{306.265}{272.961}{extension}%
\TextAt{8}{359.664}{272.961}{and}%
\TextAt{13}{384.06}{272.961}{G}%
\TextAt{9}{399.398}{272.961}{\Character{61}}%
\TextAt{13}{414.606}{272.961}{G}%
\TextAt{9}{423.84}{272.961}{\Character{40}}%
\TextAt{13}{428.392}{272.961}{L}%
\TextAt{16}{442.461}{272.961}{\Character{106}}%
\TextAt{13}{451.886}{272.961}{K}%
\TextAt{9}{462.704}{272.961}{\Character{41}}%
\TextAt{8}{467.256}{272.961}{,}%
\TextAt{8}{476.447}{272.961}{then}%
\TextAt{13}{504.74}{272.961}{C}%
\TextAt{21}{513.115}{271.168}{L}%
\TextAt{8}{524.903}{272.961}{is}%
\TextAt{8}{79.226}{255.915}{canonically}%
\TextAt{8}{140.042}{255.915}{a}%
\TextAt{13}{149.798}{255.915}{G}%
\TextAt{8}{159.031}{255.915}{-module}%
\TextAt{8}{204.222}{255.915}{and}%
\TextAt{13}{226.984}{255.915}{C}%
\TextAt{21}{236.218}{260.254}{G}%
\TextAt{21}{235.359}{252.895}{L}%
\TextAt{9}{246.636}{255.916}{\Character{61}}%
\TextAt{13}{259.061}{255.915}{C}%
\TextAt{21}{267.436}{254.122}{K}%
\TextAt{8}{275.639}{255.915}{.}%
\TextAt{19}{79.226}{221.458}{Proof.}%
\TextAt{8}{116}{221.458}{The}%
\TextAt{13}{140.353}{221.458}{G}%
\TextAt{8}{149.587}{221.458}{-module}%
\TextAt{13}{195.076}{221.458}{I}%
\TextAt{21}{200.239}{219.665}{L}%
\TextAt{8}{210.693}{221.458}{contains}%
\TextAt{13}{256.911}{221.458}{L}%
\TextAt{3}{264.876}{225.797}{\Character{3}}%
\TextAt{8}{273.805}{221.458}{as}%
\TextAt{8}{288.474}{221.458}{a}%
\TextAt{13}{298.528}{221.458}{G}%
\TextAt{8}{307.762}{221.458}{-submodule.}%
\TextAt{8}{376.022}{221.458}{Hence}%
\TextAt{8}{411.099}{221.458}{every}%
\TextAt{13}{442.29}{221.458}{\Character{27}}%
\NativeMathAt{453.203}{221.458}{11.955}{\in}%
\TextAt{13}{465.003}{221.458}{G}%
\TextAt{8}{478.433}{221.458}{induces}%
\TextAt{8}{520.408}{221.458}{an}%
\TextAt{8}{79.226}{204.412}{automorphism}%
\TextAt{13}{212.46}{187.366}{\Character{27}}%
\TextAt{9}{222.863}{187.366}{\Character{58}}%
\TextAt{13}{229.435}{187.366}{C}%
\TextAt{21}{237.81}{185.573}{L}%
\NativeArrowAt{247.389}{187.366}{11.955}{19.26}{\longrightarrow}%
\TextAt{13}{269.971}{187.366}{C}%
\TextAt{21}{278.346}{185.573}{L}%
\TextAt{13}{284.603}{187.366}{\Character{59}}%
\TextAt{13}{313.738}{187.366}{\Character{11}L}%
\TextAt{3}{329.243}{192.303}{\Character{3}}%
\NativeArrowAt{337.297}{187.366}{11.955}{19.26}{\longmapsto}%
\TextAt{9}{359.879}{187.366}{\Character{40}}%
\TextAt{13}{364.431}{187.366}{\Character{27}}%
\TextAt{13}{371.497}{187.366}{\Character{11}}%
\TextAt{9}{379.035}{187.366}{\Character{41}}%
\TextAt{13}{383.587}{187.366}{L}%
\TextAt{3}{391.552}{192.303}{\Character{3}}%
\TextAt{13}{396.284}{187.366}{\Character{58}}%
\TextAt{8}{79.226}{164.236}{This}%
\TextAt{8}{106.146}{164.236}{gives}%
\TextAt{8}{135.027}{164.236}{us}%
\TextAt{8}{150.242}{164.236}{an}%
\TextAt{8}{166.701}{164.236}{exact}%
\TextAt{8}{197.794}{164.236}{sequence}%
\TextAt{8}{246.496}{164.236}{of}%
\TextAt{13}{260.026}{164.236}{G}%
\TextAt{8}{269.26}{164.236}{-modules}%
\TextAt{9}{319.265}{164.236}{1}%
\NativeArrowAt{328.78}{164.236}{11.955}{19.26}{\longrightarrow}%
\TextAt{13}{351.703}{164.236}{L}%
\TextAt{3}{359.667}{168.574}{\Character{3}}%
\NativeArrowAt{368.061}{164.236}{11.955}{19.26}{\longrightarrow}%
\TextAt{13}{390.984}{164.236}{I}%
\TextAt{21}{396.147}{162.443}{L}%
\NativeArrowAt{406.066}{164.236}{11.955}{19.26}{\longrightarrow}%
\TextAt{13}{428.989}{164.236}{C}%
\TextAt{21}{437.364}{162.443}{L}%
\NativeArrowAt{447.283}{164.236}{11.955}{19.26}{\longrightarrow}%
\TextAt{9}{470.206}{164.236}{1}%
\TextAt{8}{476.059}{164.236}{.}%
\TextAt{8}{485.107}{164.236}{We}%
\TextAt{8}{505.452}{164.236}{claim}%
\TextAt{8}{79.226}{147.19}{that}%
\TextAt{8}{104.595}{147.19}{the}%
\TextAt{8}{124.751}{147.19}{sequence}%
\TextAt{9}{223.926}{130.144}{1}%
\NativeArrowAt{233.1}{130.144}{11.955}{19.26}{\longrightarrow}%
\TextAt{13}{255.682}{130.144}{L}%
\TextAt{3}{263.647}{135.08}{\Character{3}}%
\TextAt{21}{267.881}{135.08}{G}%
\NativeArrowAt{278.299}{130.144}{11.955}{19.26}{\longrightarrow}%
\TextAt{13}{300.881}{130.144}{I}%
\TextAt{21}{306.983}{135.08}{G}%
\TextAt{21}{306.043}{127.188}{L}%
\NativeArrowAt{317.401}{130.144}{11.955}{19.26}{\longrightarrow}%
\TextAt{13}{339.983}{130.144}{C}%
\TextAt{21}{349.217}{135.08}{G}%
\TextAt{21}{348.358}{127.188}{L}%
\NativeArrowAt{359.634}{130.144}{11.955}{19.26}{\longrightarrow}%
\TextAt{9}{382.216}{130.144}{1}%
\TextAt{8}{79.226}{107.013}{deduced}%
\TextAt{8}{124.991}{107.013}{from}%
\TextAt{8}{152.882}{107.013}{the}%
\TextAt{8}{173.29}{107.013}{first}%
\TextAt{8}{197.667}{107.013}{is}%
\TextAt{8}{209.682}{107.013}{still}%
\TextAt{8}{232.756}{107.013}{exact.}%
\TextAt{8}{268.953}{107.013}{The}%
\TextAt{8}{293.258}{107.013}{injectivity}%
\TextAt{8}{348.79}{107.013}{of}%
\TextAt{13}{362.371}{107.013}{L}%
\TextAt{3}{370.335}{111.352}{\Character{3}}%
\TextAt{21}{374.569}{111.352}{G}%
\NativeArrowAt{385.415}{107.013}{11.955}{19.26}{\longrightarrow}%
\TextAt{13}{408.425}{107.013}{I}%
\TextAt{21}{414.528}{111.352}{G}%
\TextAt{21}{413.588}{103.993}{L}%
\TextAt{8}{425.774}{107.013}{is}%
\TextAt{8}{437.789}{107.013}{trivial.}%
\TextAt{8}{477.895}{107.013}{The}%
\TextAt{8}{502.2}{107.013}{kernel}%
\TextAt{8}{79.226}{89.967}{of}%
\TextAt{13}{92.994}{89.967}{I}%
\TextAt{21}{99.097}{94.306}{G}%
\TextAt{21}{98.157}{86.947}{L}%
\NativeArrowAt{110.26}{89.967}{11.955}{19.26}{\longrightarrow}%
\TextAt{13}{133.588}{89.967}{C}%
\TextAt{21}{142.822}{94.306}{G}%
\TextAt{21}{141.963}{86.947}{L}%
\TextAt{8}{154.254}{89.967}{is}%
\TextAt{13}{166.455}{89.967}{I}%
\TextAt{21}{172.558}{94.306}{G}%
\TextAt{21}{171.618}{86.947}{L}%
\TextAt{16}{182.61}{89.967}{\Character{92}}%
\TextAt{13}{193.535}{89.967}{L}%
\TextAt{3}{201.5}{94.306}{\Character{3}}%
\TextAt{9}{210.299}{89.967}{\Character{61}}%
\TextAt{13}{223.47}{89.967}{I}%
\TextAt{21}{228.633}{88.174}{K}%
\TextAt{16}{239.791}{89.967}{\Character{92}}%
\TextAt{13}{250.716}{89.967}{L}%
\TextAt{3}{258.68}{94.306}{\Character{3}}%
\TextAt{9}{267.479}{89.967}{\Character{61}}%
\TextAt{13}{280.651}{89.967}{K}%
\TextAt{3}{291.469}{94.306}{\Character{3}}%
\TextAt{9}{300.268}{89.967}{\Character{61}}%
\TextAt{13}{313.439}{89.967}{L}%
\TextAt{3}{321.404}{94.306}{\Character{3}}%
\TextAt{21}{325.638}{94.306}{G}%
\TextAt{8}{332.735}{89.967}{.}%
\TextAt{8}{342.496}{89.967}{The}%
\TextAt{8}{366.988}{89.967}{surjectivity}%
\TextAt{8}{428.624}{89.967}{of}%
\TextAt{13}{442.392}{89.967}{I}%
\TextAt{21}{448.494}{94.306}{G}%
\TextAt{21}{447.554}{86.947}{L}%
\NativeArrowAt{459.658}{89.967}{11.955}{19.26}{\longrightarrow}%
\TextAt{13}{482.986}{89.967}{C}%
\TextAt{21}{492.22}{94.306}{G}%
\TextAt{21}{491.361}{86.947}{L}%
\TextAt{8}{503.651}{89.967}{is}%
\TextAt{8}{515.853}{89.967}{not}%
\TextAt{8}{79.226}{72.921}{altogether}%
\TextAt{8}{135.208}{72.921}{straightforward.}%
\TextAt{8}{224.184}{72.921}{To}%
\TextAt{8}{242.1}{72.921}{prove}%
\TextAt{8}{274.339}{72.921}{it,}%
\TextAt{8}{290.155}{72.921}{let}%
\TextAt{13}{307.749}{72.921}{\Character{11}L}%
\TextAt{3}{323.235}{77.26}{\Character{3}}%
\NativeMathAt{332.462}{72.921}{11.955}{\in}%
\TextAt{13}{344.927}{72.921}{C}%
\TextAt{21}{354.16}{77.26}{G}%
\TextAt{21}{353.302}{69.901}{L}%
\TextAt{8}{361.257}{72.921}{.}%
\TextAt{8}{371.772}{72.921}{For}%
\TextAt{8}{393.43}{72.921}{every}%
\TextAt{13}{425.011}{72.921}{\Character{27}}%
\NativeMathAt{436.588}{72.921}{11.955}{\in}%
\TextAt{13}{449.052}{72.921}{G}%
\TextAt{8}{458.286}{72.921}{,}%
\TextAt{8}{466.296}{72.921}{one}%
\TextAt{8}{488.444}{72.921}{then}%
\TextAt{8}{515.793}{72.921}{has}%
\TextAt{8}{297.211}{43.033}{339}%
\Destination{78.23}{809.815}{nav.page.339}%
\Destination{79.36}{557.455}{nav.print.372}%
\Destination{79.36}{556.455}{nav.number.6.2.4}%
\Destination{79.36}{286.955}{nav.number.6.2.5}%
\Destination{79.226}{741.043}{errata.142}%
}

% Page 352 (body page 340)
\Page{%
\TextAt{2}{79.226}{796.943}{VI.}%
\TextAt{2}{98.012}{796.943}{Global}%
\TextAt{2}{133.248}{796.943}{Class}%
\TextAt{2}{161.841}{796.943}{Field}%
\TextAt{2}{189.572}{796.943}{Theory}%
\Vector{q 0 G 0 g 0.299 w 7.226 18.729 m 460.769 18.729 l S Q}%
\TextAt{13}{79.226}{761.124}{\Character{27}}%
\TextAt{9}{86.309}{761.124}{\Character{40}}%
\TextAt{13}{90.861}{761.124}{\Character{11}L}%
\TextAt{3}{106.347}{765.462}{\Character{3}}%
\TextAt{9}{111.079}{761.124}{\Character{41}}%
\TextAt{9}{119.28}{761.124}{\Character{61}}%
\TextAt{13}{132.054}{761.124}{\Character{11}L}%
\TextAt{3}{147.54}{765.462}{\Character{3}}%
\TextAt{8}{152.272}{761.124}{,}%
\TextAt{8}{159.669}{761.124}{i.e.,}%
\TextAt{13}{182.022}{761.124}{\Character{27}}%
\TextAt{13}{189.088}{761.124}{\Character{11}}%
\TextAt{9}{200.285}{761.124}{\Character{61}}%
\TextAt{13}{213.048}{761.124}{\Character{11}x}%
\TextAt{21}{227.222}{759.33}{\Character{27}}%
\TextAt{8}{236.941}{761.124}{for}%
\TextAt{8}{255.024}{761.124}{some}%
\TextAt{13}{284.548}{761.124}{x}%
\TextAt{21}{291.2}{759.33}{\Character{27}}%
\NativeMathAt{300.482}{761.124}{11.955}{\in}%
\TextAt{13}{312.111}{761.124}{L}%
\TextAt{3}{320.076}{765.462}{\Character{3}}%
\TextAt{8}{324.808}{761.124}{.}%
\TextAt{8}{333.851}{761.124}{This}%
\TextAt{13}{360.769}{761.124}{x}%
\TextAt{21}{367.421}{759.33}{\Character{27}}%
\TextAt{8}{377.14}{761.124}{is}%
\TextAt{8}{389.102}{761.124}{a}%
\TextAt{8}{399.056}{761.124}{“crossed}%
\TextAt{8}{445.222}{761.124}{homomorphism”,}%
\TextAt{8}{79.226}{744.078}{i.e.,}%
\TextAt{8}{101.331}{744.078}{we}%
\TextAt{8}{118.559}{744.078}{have}%
\TextAt{13}{271.432}{727.032}{x}%
\TextAt{21}{278.084}{725.238}{\Character{27}}%
\TextAt{21}{283.201}{725.238}{\Character{28}}%
\TextAt{9}{291.707}{727.032}{\Character{61}}%
\TextAt{13}{304.133}{727.032}{x}%
\TextAt{21}{310.785}{725.238}{\Character{27}}%
\TextAt{13}{318.401}{727.032}{\Character{27}}%
\TextAt{13}{325.466}{727.032}{x}%
\TextAt{21}{332.135}{725.238}{\Character{28}}%
\TextAt{13}{337.312}{727.032}{\Character{58}}%
\TextAt{8}{79.226}{705.806}{Indeed,}%
\TextAt{13}{122.351}{705.806}{x}%
\TextAt{21}{129.003}{704.013}{\Character{27}}%
\TextAt{21}{134.12}{704.013}{\Character{28}}%
\TextAt{9}{145.139}{705.806}{\Character{61}}%
\TextAt{13}{160.078}{705.806}{\Character{27}}%
\TextAt{13}{167.144}{705.806}{\Character{28}}%
\TextAt{13}{173.564}{705.806}{\Character{11}}%
\TextAt{9}{181.101}{705.806}{\Character{47}}%
\TextAt{13}{186.954}{705.806}{\Character{11}}%
\TextAt{9}{200.31}{705.806}{\Character{61}}%
\TextAt{9}{215.23}{705.806}{\Character{40}}%
\TextAt{13}{219.801}{705.806}{\Character{27}}%
\TextAt{13}{226.867}{705.806}{\Character{28}}%
\TextAt{13}{233.287}{705.806}{\Character{11}}%
\TextAt{9}{240.824}{705.806}{\Character{47}}%
\TextAt{13}{246.677}{705.806}{\Character{27}}%
\TextAt{13}{253.743}{705.806}{\Character{11}}%
\TextAt{9}{261.281}{705.806}{\Character{41}\Character{40}}%
\TextAt{13}{270.386}{705.806}{\Character{27}}%
\TextAt{13}{277.452}{705.806}{\Character{11}}%
\TextAt{9}{284.99}{705.806}{\Character{47}}%
\TextAt{13}{290.843}{705.806}{\Character{11}}%
\TextAt{9}{298.364}{705.806}{\Character{41}}%
\TextAt{9}{308.741}{705.806}{\Character{61}}%
\TextAt{13}{323.69}{705.806}{\Character{27}}%
\TextAt{9}{330.772}{705.806}{\Character{40}}%
\TextAt{13}{335.325}{705.806}{\Character{28}}%
\TextAt{13}{341.733}{705.806}{\Character{11}}%
\TextAt{9}{349.265}{705.806}{\Character{47}}%
\TextAt{13}{355.118}{705.806}{\Character{11}}%
\TextAt{9}{362.64}{705.806}{\Character{41}}%
\TextAt{13}{367.192}{705.806}{x}%
\TextAt{21}{373.844}{704.013}{\Character{27}}%
\TextAt{9}{385.302}{705.806}{\Character{61}}%
\TextAt{13}{400.241}{705.806}{\Character{27}}%
\TextAt{13}{407.307}{705.806}{x}%
\TextAt{21}{413.975}{704.013}{\Character{28}}%
\TextAt{13}{419.152}{705.806}{x}%
\TextAt{21}{425.804}{704.013}{\Character{27}}%
\TextAt{8}{431.427}{705.806}{.}%
\TextAt{8}{444.3}{705.806}{By}%
\TextAt{8}{464.138}{705.806}{Hilbert}%
\TextAt{8}{505.926}{705.806}{90}%
\TextAt{8}{523.014}{705.806}{in}%
\TextAt{8}{79.226}{688.76}{Noether’s}%
\TextAt{8}{131.759}{688.76}{version}%
\TextAt{8}{171.296}{688.76}{(see}%
\TextAt{8}{194.562}{688.76}{chap.}%
\TextAt{8}{225.253}{688.76}{IV,}%
\TextAt{8}{245.196}{688.76}{(3.8))}%
\TextAt{8}{277.524}{688.76}{such}%
\TextAt{8}{303.72}{688.76}{a}%
\TextAt{8}{313.273}{688.76}{crossed}%
\TextAt{8}{353.516}{688.76}{homomorphism}%
\TextAt{8}{435.984}{688.76}{is}%
\TextAt{8}{447.547}{688.76}{of}%
\TextAt{8}{460.675}{688.76}{the}%
\TextAt{8}{480.63}{688.76}{form}%
\TextAt{13}{508.068}{688.76}{x}%
\TextAt{21}{514.721}{686.967}{\Character{27}}%
\TextAt{9}{523.665}{688.76}{\Character{61}}%
\TextAt{13}{79.226}{671.714}{\Character{27}}%
\TextAt{13}{86.292}{671.714}{y}%
\TextAt{9}{92.445}{671.714}{\Character{47}}%
\TextAt{13}{98.298}{671.714}{y}%
\TextAt{8}{109.55}{671.714}{for}%
\TextAt{8}{128.652}{671.714}{some}%
\TextAt{13}{159.196}{671.714}{y}%
\NativeMathAt{170.728}{671.714}{11.955}{\in}%
\TextAt{13}{184.093}{671.714}{L}%
\TextAt{3}{192.058}{676.053}{\Character{3}}%
\TextAt{8}{196.79}{671.714}{.}%
\TextAt{8}{208.891}{671.714}{Putting}%
\TextAt{13}{253.195}{671.714}{\Character{11}}%
\TextAt{3}{260.716}{676.053}{\Character{48}}%
\TextAt{9}{268.907}{671.714}{\Character{61}}%
\TextAt{13}{283.407}{671.714}{\Character{11}y}%
\TextAt{3}{297.066}{676.053}{\Character{0}}%
\TextAt{17}{303.652}{676.053}{1}%
\TextAt{8}{313.499}{671.714}{yields}%
\TextAt{13}{347.617}{671.714}{\Character{11}}%
\TextAt{3}{355.139}{676.053}{\Character{48}}%
\TextAt{13}{357.934}{671.714}{L}%
\TextAt{3}{365.899}{676.053}{\Character{3}}%
\TextAt{9}{376.026}{671.714}{\Character{61}}%
\TextAt{13}{390.527}{671.714}{\Character{11}L}%
\TextAt{3}{406.013}{676.053}{\Character{3}}%
\TextAt{8}{415.86}{671.714}{and}%
\TextAt{13}{439.84}{671.714}{\Character{27}}%
\TextAt{13}{446.906}{671.714}{\Character{11}}%
\TextAt{3}{454.444}{676.053}{\Character{48}}%
\TextAt{9}{462.635}{671.714}{\Character{61}}%
\TextAt{13}{477.135}{671.714}{\Character{27}}%
\TextAt{13}{484.201}{671.714}{\Character{11}}%
\TextAt{13}{493.705}{671.714}{\Character{27}}%
\TextAt{13}{500.782}{671.714}{y}%
\TextAt{3}{506.951}{676.053}{\Character{0}}%
\TextAt{17}{513.537}{676.053}{1}%
\TextAt{9}{523.665}{671.714}{\Character{61}}%
\TextAt{13}{79.226}{654.668}{\Character{11}x}%
\TextAt{21}{93.4}{652.875}{\Character{27}}%
\TextAt{13}{101.016}{654.668}{\Character{27}}%
\TextAt{13}{108.082}{654.668}{y}%
\TextAt{3}{114.235}{659.007}{\Character{0}}%
\TextAt{17}{120.821}{659.007}{1}%
\TextAt{9}{128.875}{654.668}{\Character{61}}%
\TextAt{13}{141.3}{654.668}{\Character{11}y}%
\TextAt{3}{154.958}{659.007}{\Character{0}}%
\TextAt{17}{161.545}{659.007}{1}%
\TextAt{9}{169.598}{654.668}{\Character{61}}%
\TextAt{13}{182.024}{654.668}{\Character{11}}%
\TextAt{3}{189.545}{659.007}{\Character{48}}%
\TextAt{8}{192.34}{654.668}{,}%
\TextAt{8}{199.49}{654.668}{hence}%
\TextAt{13}{231.996}{654.668}{\Character{11}}%
\TextAt{3}{239.517}{659.007}{\Character{48}}%
\NativeMathAt{245.633}{654.668}{11.955}{\in}%
\TextAt{13}{256.924}{654.668}{I}%
\TextAt{21}{262.087}{652.875}{K}%
\TextAt{8}{270.29}{654.668}{.}%
\TextAt{8}{278.738}{654.668}{This}%
\TextAt{8}{305.458}{654.668}{proves}%
\TextAt{8}{341.622}{654.668}{surjectivity.}%
\Vector{q 0 G 0 g 0.399 w 452.334 -119.332 m 452.334 -111.262 l S Q}%
\Vector{q 0 G 0 g 0.399 w 452.534 -111.461 m 459.707 -111.461 l S Q}%
\Vector{q 0 G 0 g 0.399 w 452.534 -119.132 m 459.707 -119.132 l S Q}%
\Vector{q 0 G 0 g 0.399 w 459.906 -119.332 m 459.906 -111.262 l S Q}%
\TextAt{8}{97.159}{627.198}{The}%
\TextAt{8}{122.008}{627.198}{norm}%
\TextAt{8}{153.374}{627.198}{map}%
\TextAt{13}{180.184}{627.198}{N}%
\TextAt{21}{189.558}{625.338}{L}%
\TextAt{3}{195.317}{625.338}{\Character{106}}%
\TextAt{21}{197.669}{625.338}{K}%
\TextAt{9}{210.549}{627.198}{\Character{58}}%
\TextAt{13}{218.477}{627.198}{I}%
\TextAt{21}{223.64}{625.405}{L}%
\NativeArrowAt{234.574}{627.198}{11.955}{19.26}{\longrightarrow}%
\TextAt{13}{258.512}{627.198}{I}%
\TextAt{21}{263.675}{625.405}{K}%
\TextAt{8}{276.571}{627.198}{sends}%
\TextAt{8}{308.701}{627.198}{principal}%
\TextAt{8}{358.274}{627.198}{idèles}%
\TextAt{8}{390.99}{627.198}{to}%
\TextAt{8}{406.096}{627.198}{principal}%
\TextAt{8}{455.669}{627.198}{idèles}%
\TextAt{8}{488.385}{627.198}{by}%
\TextAt{8}{505.44}{627.198}{(2.3).}%
\TextAt{8}{79.226}{610.152}{Hence}%
\TextAt{8}{114.004}{610.152}{we}%
\TextAt{8}{131.231}{610.152}{get}%
\TextAt{8}{150.742}{610.152}{a}%
\TextAt{8}{160.497}{610.152}{norm}%
\TextAt{8}{191.067}{610.152}{map}%
\TextAt{8}{217.081}{610.152}{also}%
\TextAt{8}{240.561}{610.152}{for}%
\TextAt{8}{258.446}{610.152}{the}%
\TextAt{8}{278.603}{610.152}{idèle}%
\TextAt{8}{305.908}{610.152}{class}%
\TextAt{8}{333.345}{610.152}{group}%
\TextAt{13}{366.521}{610.152}{C}%
\TextAt{21}{374.896}{608.359}{L}%
\TextAt{8}{381.153}{610.152}{,}%
\TextAt{13}{258.025}{586.509}{N}%
\TextAt{21}{267.399}{584.649}{L}%
\TextAt{3}{273.158}{584.649}{\Character{106}}%
\TextAt{21}{275.51}{584.649}{K}%
\TextAt{9}{287.034}{586.508}{\Character{58}}%
\TextAt{13}{293.606}{586.509}{C}%
\TextAt{21}{301.981}{584.715}{L}%
\NativeArrowAt{311.559}{586.508}{11.955}{19.26}{\longrightarrow}%
\TextAt{13}{334.141}{586.508}{C}%
\TextAt{21}{342.516}{584.715}{K}%
\TextAt{13}{350.719}{586.508}{\Character{58}}%
\TextAt{8}{79.226}{562.865}{It}%
\TextAt{8}{91.899}{562.865}{enjoys}%
\TextAt{8}{127.405}{562.865}{the}%
\TextAt{8}{147.562}{562.865}{same}%
\TextAt{8}{176.888}{562.865}{properties}%
\TextAt{8}{231.906}{562.865}{(2.3),}%
\TextAt{8}{263.133}{562.865}{(i),}%
\TextAt{8}{282.643}{562.865}{(ii),}%
\TextAt{8}{305.406}{562.865}{(iii),}%
\TextAt{8}{331.42}{562.865}{as}%
\TextAt{8}{345.791}{562.865}{the}%
\TextAt{8}{365.947}{562.865}{norm}%
\TextAt{8}{396.516}{562.865}{map}%
\TextAt{8}{422.531}{562.865}{on}%
\TextAt{8}{438.79}{562.865}{the}%
\TextAt{8}{458.946}{562.865}{idèle}%
\TextAt{8}{486.252}{562.865}{group.}%
\TextAt{12}{79.226}{538.383}{Exercise}%
\TextAt{12}{132.326}{538.383}{VI.2.1.}%
\TextAt{8}{175.994}{538.383}{Let}%
\TextAt{13}{196.875}{538.383}{\Character{33}}%
\TextAt{17}{204.164}{536.59}{1}%
\TextAt{13}{208.897}{538.383}{\Character{59}}%
\TextAt{13}{214.121}{538.383}{\Character{58}}%
\TextAt{13}{219.358}{538.383}{\Character{58}}%
\TextAt{13}{224.582}{538.383}{\Character{58}}%
\TextAt{13}{229.819}{538.383}{\Character{59}}%
\TextAt{13}{235.055}{538.383}{\Character{33}}%
\TextAt{21}{242.407}{536.59}{n}%
\TextAt{8}{251.853}{538.383}{be}%
\TextAt{8}{267.689}{538.383}{a}%
\TextAt{8}{277.357}{538.383}{basis}%
\TextAt{8}{306.009}{538.383}{of}%
\TextAt{13}{319.251}{538.383}{L}%
\TextAt{16}{330.536}{538.383}{\Character{106}}%
\TextAt{13}{337.178}{538.383}{K}%
\TextAt{8}{347.996}{538.383}{.}%
\TextAt{8}{356.415}{538.383}{Then}%
\TextAt{8}{386.884}{538.383}{the}%
\TextAt{8}{406.953}{538.383}{isomorphism}%
\TextAt{13}{475.285}{538.383}{L}%
\TextAt{16}{485.726}{538.383}{\Character{10}}%
\TextAt{21}{495.025}{536.59}{K}%
\TextAt{13}{505.704}{538.383}{K}%
\TextAt{27}{515.663}{536.59}{p}%
\NativeMathAt{523.471}{538.383}{11.955}{\cong}%
\TextAt{20}{79.226}{530.304}{\Character{81}}%
\TextAt{27}{90.517}{517.851}{P}%
\TextAt{3}{97.113}{517.851}{\Character{106}}%
\TextAt{27}{99.465}{517.851}{p}%
\TextAt{13}{105.944}{521.338}{L}%
\TextAt{27}{113.909}{519.544}{P}%
\TextAt{8}{124.9}{521.338}{induces,}%
\TextAt{8}{169.827}{521.338}{for}%
\TextAt{8}{187.712}{521.338}{almost}%
\TextAt{8}{225.503}{521.338}{all}%
\TextAt{8}{241.762}{521.338}{prime}%
\TextAt{8}{274.925}{521.338}{ideals}%
\TextAt{14}{307.503}{521.338}{p}%
\TextAt{8}{317.384}{521.338}{of}%
\TextAt{13}{330.714}{521.338}{K}%
\TextAt{8}{341.531}{521.338}{,}%
\TextAt{8}{348.68}{521.338}{an}%
\TextAt{8}{364.939}{521.338}{isomorphism}%
\TextAt{13}{231.398}{495.711}{\Character{33}}%
\TextAt{17}{238.687}{493.918}{1}%
\TextAt{16}{243.42}{495.711}{\Character{79}}%
\TextAt{27}{252.938}{493.918}{p}%
\TextAt{16}{260.081}{495.711}{\Character{8}}%
\TextAt{16}{272.024}{495.711}{\Character{1}}%
\TextAt{16}{277.32}{495.711}{\Character{1}}%
\TextAt{16}{282.629}{495.711}{\Character{1}}%
\TextAt{16}{288.594}{495.711}{\Character{8}}%
\TextAt{13}{300.596}{495.711}{\Character{33}}%
\TextAt{21}{307.885}{493.918}{n}%
\TextAt{16}{313.521}{495.711}{\Character{79}}%
\TextAt{27}{323.039}{493.918}{p}%
\NativeMathAt{330.847}{495.711}{11.955}{\cong}%
\TextAt{20}{343.466}{507.069}{\Character{89}}%
\TextAt{27}{344.636}{481.166}{P}%
\TextAt{3}{351.231}{481.166}{\Character{106}}%
\TextAt{27}{353.584}{481.166}{p}%
\TextAt{16}{360.735}{495.711}{\Character{79}}%
\TextAt{27}{370.252}{493.918}{P}%
\TextAt{13}{377.346}{495.711}{\Character{59}}%
\TextAt{8}{79.226}{461.967}{where}%
\TextAt{16}{113.035}{461.967}{\Character{79}}%
\TextAt{27}{122.553}{460.174}{p}%
\TextAt{8}{127.04}{461.967}{,}%
\TextAt{8}{134.189}{461.967}{resp.}%
\TextAt{16}{163.511}{461.967}{\Character{79}}%
\TextAt{27}{173.029}{460.174}{P}%
\TextAt{8}{180.122}{461.967}{,}%
\TextAt{8}{187.272}{461.967}{is}%
\TextAt{8}{199.036}{461.967}{the}%
\TextAt{8}{219.192}{461.967}{valuation}%
\TextAt{8}{270.264}{461.967}{ring}%
\TextAt{8}{294.33}{461.967}{of}%
\TextAt{13}{307.66}{461.967}{K}%
\TextAt{27}{317.619}{460.174}{p}%
\TextAt{8}{322.106}{461.967}{,}%
\TextAt{8}{329.255}{461.967}{resp.}%
\TextAt{13}{358.577}{461.967}{L}%
\TextAt{27}{366.542}{460.174}{P}%
\TextAt{8}{373.636}{461.967}{.}%
\TextAt{12}{79.226}{433.039}{Exercise}%
\TextAt{12}{132.634}{433.039}{VI.2.2.}%
\TextAt{8}{176.57}{433.039}{Let}%
\TextAt{13}{197.72}{433.039}{L}%
\TextAt{16}{209.312}{433.039}{\Character{106}}%
\TextAt{13}{216.261}{433.039}{K}%
\TextAt{8}{231.157}{433.039}{be}%
\TextAt{8}{247.261}{433.039}{a}%
\TextAt{8}{257.197}{433.039}{finite}%
\TextAt{8}{287.289}{433.039}{extension.}%
\TextAt{8}{344.147}{433.039}{The}%
\TextAt{8}{368.381}{433.039}{absolute}%
\TextAt{8}{414.803}{433.039}{norm}%
\TextAt{14}{445.553}{433.039}{N}%
\TextAt{8}{459.579}{433.039}{of}%
\TextAt{8}{473.089}{433.039}{idèles}%
\TextAt{8}{505.19}{433.039}{of}%
\TextAt{13}{518.7}{433.039}{K}%
\TextAt{8}{529.518}{433.039}{,}%
\TextAt{8}{79.226}{415.993}{resp.}%
\TextAt{13}{110.241}{415.993}{L}%
\TextAt{8}{118.205}{415.993}{,}%
\TextAt{8}{126.06}{415.993}{behaves}%
\TextAt{8}{170.26}{415.993}{as}%
\TextAt{8}{185.195}{415.993}{follows}%
\TextAt{8}{224.195}{415.993}{under}%
\TextAt{8}{257.923}{415.993}{the}%
\TextAt{8}{278.643}{415.993}{inclusion}%
\TextAt{13}{328.044}{415.993}{i}%
\TextAt{21}{332.038}{414.133}{L}%
\TextAt{3}{337.797}{414.133}{\Character{106}}%
\TextAt{21}{340.149}{414.133}{K}%
\TextAt{9}{352.634}{415.993}{\Character{58}}%
\TextAt{13}{360.168}{415.993}{I}%
\TextAt{21}{365.331}{414.2}{K}%
\NativeArrowAt{377.816}{415.993}{11.955}{19.26}{\longrightarrow}%
\TextAt{13}{401.36}{415.993}{I}%
\TextAt{21}{406.523}{414.2}{L}%
\TextAt{8}{412.78}{415.993}{,}%
\TextAt{8}{420.634}{415.993}{resp.}%
\TextAt{8}{451.649}{415.993}{under}%
\TextAt{8}{485.377}{415.993}{the}%
\TextAt{8}{506.098}{415.993}{norm}%
\TextAt{13}{79.226}{398.947}{N}%
\TextAt{21}{88.6}{397.087}{L}%
\TextAt{3}{94.359}{397.087}{\Character{106}}%
\TextAt{21}{96.711}{397.087}{K}%
\TextAt{9}{108.235}{398.947}{\Character{58}}%
\TextAt{13}{114.807}{398.947}{I}%
\TextAt{21}{119.97}{397.154}{L}%
\NativeArrowAt{129.548}{398.947}{11.955}{19.26}{\longrightarrow}%
\TextAt{13}{152.13}{398.947}{I}%
\TextAt{21}{157.293}{397.154}{K}%
\TextAt{8}{165.496}{398.947}{:}%
\TextAt{14}{113.921}{375.304}{N}%
\TextAt{9}{123.87}{375.304}{\Character{40}}%
\TextAt{13}{128.422}{375.304}{i}%
\TextAt{21}{132.415}{373.444}{L}%
\TextAt{3}{138.174}{373.444}{\Character{106}}%
\TextAt{21}{140.527}{373.444}{K}%
\TextAt{9}{148.729}{375.304}{\Character{40}}%
\TextAt{13}{153.282}{375.304}{\Character{11}}%
\TextAt{9}{160.803}{375.304}{\Character{41}\Character{41}}%
\TextAt{9}{173.201}{375.304}{\Character{61}}%
\TextAt{14}{185.654}{375.304}{N}%
\TextAt{9}{195.602}{375.304}{\Character{40}}%
\TextAt{13}{200.155}{375.304}{\Character{11}}%
\TextAt{9}{207.676}{375.304}{\Character{41}}%
\TextAt{17}{212.229}{380.24}{\Character{91}}%
\TextAt{21}{214.581}{380.24}{L}%
\TextAt{17}{220.34}{380.24}{\Character{58}}%
\TextAt{21}{222.692}{380.24}{K}%
\TextAt{17}{230.397}{380.24}{\Character{93}}%
\TextAt{8}{245.203}{375.304}{for}%
\TextAt{13}{263.088}{375.304}{\Character{11}}%
\NativeMathAt{273.93}{375.304}{11.955}{\in}%
\TextAt{13}{285.221}{375.304}{I}%
\TextAt{21}{290.384}{373.51}{K}%
\TextAt{13}{298.587}{375.304}{\Character{59}}%
\TextAt{14}{327.741}{375.304}{N}%
\TextAt{9}{337.689}{375.304}{\Character{40}}%
\TextAt{13}{342.242}{375.304}{N}%
\TextAt{21}{351.615}{373.444}{L}%
\TextAt{3}{357.374}{373.444}{\Character{106}}%
\TextAt{21}{359.727}{373.444}{K}%
\TextAt{9}{367.929}{375.304}{\Character{40}}%
\TextAt{13}{372.482}{375.304}{\Character{11}}%
\TextAt{9}{380.004}{375.304}{\Character{41}\Character{41}}%
\TextAt{9}{392.401}{375.304}{\Character{61}}%
\TextAt{14}{404.854}{375.304}{N}%
\TextAt{9}{414.803}{375.304}{\Character{40}}%
\TextAt{13}{419.355}{375.304}{\Character{11}}%
\TextAt{9}{426.877}{375.304}{\Character{41}}%
\TextAt{8}{443.384}{375.304}{for}%
\TextAt{13}{461.269}{375.304}{\Character{11}}%
\NativeMathAt{472.112}{375.304}{11.955}{\in}%
\TextAt{13}{483.403}{375.304}{I}%
\TextAt{21}{488.565}{373.51}{L}%
\TextAt{13}{494.823}{375.304}{\Character{58}}%
\TextAt{12}{79.226}{339.778}{Exercise}%
\TextAt{12}{132.908}{339.778}{VI.2.3.}%
\TextAt{8}{177.083}{339.778}{The}%
\TextAt{8}{201.556}{339.778}{correspondence}%
\TextAt{8}{283.653}{339.778}{between}%
\TextAt{8}{329.262}{339.778}{idèles}%
\TextAt{8}{361.601}{339.778}{and}%
\TextAt{8}{384.783}{339.778}{ideals,}%
\TextAt{13}{421.136}{339.778}{\Character{11}}%
\NativeArrowAt{432.692}{339.778}{11.955}{11.955}{\mapsto}%
\TextAt{9}{448.682}{339.778}{\Character{40}}%
\TextAt{13}{453.234}{339.778}{\Character{11}}%
\TextAt{9}{460.756}{339.778}{\Character{41}}%
\TextAt{8}{465.308}{339.778}{,}%
\TextAt{8}{472.981}{339.778}{satisfies}%
\TextAt{8}{516.51}{339.778}{the}%
\TextAt{8}{79.226}{322.732}{following}%
\TextAt{8}{128.661}{322.732}{rule,}%
\TextAt{8}{155.321}{322.732}{in}%
\TextAt{8}{168.974}{322.732}{the}%
\TextAt{8}{189.13}{322.732}{case}%
\TextAt{8}{213.901}{322.732}{of}%
\TextAt{8}{227.231}{322.732}{a}%
\TextAt{8}{236.987}{322.732}{Galois}%
\TextAt{8}{272.9}{322.732}{extension}%
\TextAt{13}{324.666}{322.732}{L}%
\TextAt{16}{335.951}{322.732}{\Character{106}}%
\TextAt{13}{342.593}{322.732}{K}%
\TextAt{8}{353.411}{322.732}{,}%
\TextAt{9}{245.08}{299.088}{\Character{40}}%
\TextAt{13}{249.632}{299.088}{N}%
\TextAt{21}{259.006}{297.229}{L}%
\TextAt{3}{264.765}{297.229}{\Character{106}}%
\TextAt{21}{267.117}{297.229}{K}%
\TextAt{9}{275.32}{299.088}{\Character{40}}%
\TextAt{13}{279.872}{299.088}{\Character{11}}%
\TextAt{9}{287.394}{299.088}{\Character{41}\Character{41}}%
\TextAt{9}{299.792}{299.088}{\Character{61}}%
\TextAt{13}{312.245}{299.088}{N}%
\TextAt{21}{321.619}{297.229}{L}%
\TextAt{3}{327.378}{297.229}{\Character{106}}%
\TextAt{21}{329.73}{297.229}{K}%
\TextAt{9}{337.933}{299.088}{\Character{40}\Character{40}}%
\TextAt{13}{347.038}{299.088}{\Character{11}}%
\TextAt{9}{354.559}{299.088}{\Character{41}\Character{41}}%
\TextAt{13}{363.664}{299.088}{\Character{58}}%
\TextAt{8}{79.226}{275.445}{(For}%
\TextAt{8}{104.75}{275.445}{the}%
\TextAt{8}{124.907}{275.445}{norm}%
\TextAt{8}{155.476}{275.445}{on}%
\TextAt{8}{171.735}{275.445}{ideals,}%
\TextAt{8}{207.565}{275.445}{see}%
\TextAt{8}{226.478}{275.445}{chap.}%
\TextAt{8}{257.37}{275.445}{III,}%
\TextAt{8}{277.825}{275.445}{§1.)}%
\TextAt{12}{79.226}{246.517}{Exercise}%
\TextAt{12}{134.5}{246.517}{VI.2.4.}%
\TextAt{8}{180.058}{246.517}{The}%
\TextAt{8}{205.915}{246.517}{ideal}%
\TextAt{8}{235.68}{246.517}{class}%
\TextAt{8}{264.92}{246.517}{group,}%
\TextAt{8}{303.6}{246.517}{unlike}%
\TextAt{8}{339.869}{246.517}{the}%
\TextAt{8}{361.828}{246.517}{idèle}%
\TextAt{8}{390.936}{246.517}{class}%
\TextAt{8}{420.176}{246.517}{group,}%
\TextAt{8}{458.856}{246.517}{does}%
\TextAt{8}{487.056}{246.517}{not}%
\TextAt{8}{509.672}{246.517}{have}%
\TextAt{8}{79.226}{229.471}{Galois}%
\TextAt{8}{114.517}{229.471}{descent.}%
\TextAt{8}{160.214}{229.471}{More}%
\TextAt{8}{189.826}{229.471}{precisely,}%
\TextAt{8}{239.456}{229.471}{for}%
\TextAt{8}{256.719}{229.471}{a}%
\TextAt{8}{265.852}{229.471}{Galois}%
\TextAt{8}{301.143}{229.471}{extension}%
\TextAt{13}{352.286}{229.471}{L}%
\TextAt{16}{363.572}{229.471}{\Character{106}}%
\TextAt{13}{370.213}{229.471}{K}%
\TextAt{8}{381.031}{229.471}{,}%
\TextAt{8}{387.682}{229.471}{the}%
\TextAt{8}{407.217}{229.471}{homomorphism}%
\TextAt{13}{489.264}{229.471}{C}%
\TextAt{13}{498.482}{229.471}{l}%
\TextAt{21}{501.985}{227.677}{K}%
\NativeArrowAt{513.508}{229.471}{11.955}{19.26}{\longrightarrow}%
\TextAt{13}{79.226}{212.425}{C}%
\TextAt{13}{88.444}{212.425}{l}%
\TextAt{21}{92.21}{218.535}{G}%
\TextAt{17}{98.808}{218.535}{\Character{40}}%
\TextAt{21}{102.101}{218.535}{L}%
\TextAt{3}{107.861}{218.535}{\Character{106}}%
\TextAt{21}{110.213}{218.535}{K}%
\TextAt{17}{117.917}{218.535}{\Character{41}}%
\TextAt{21}{91.947}{209.184}{L}%
\TextAt{8}{125.606}{212.425}{is}%
\TextAt{8}{137.37}{212.425}{in}%
\TextAt{8}{151.023}{212.425}{general}%
\TextAt{8}{191.348}{212.425}{neither}%
\TextAt{8}{231.015}{212.425}{injective}%
\TextAt{8}{277.509}{212.425}{nor}%
\TextAt{8}{298.322}{212.425}{surjective.}%
\TextAt{12}{79.226}{183.496}{Exercise}%
\TextAt{12}{132.367}{183.496}{VI.2.5.}%
\TextAt{8}{176.07}{183.496}{Define}%
\TextAt{8}{212.266}{183.496}{the}%
\TextAt{8}{232.37}{183.496}{trace}%
\TextAt{13}{261.584}{183.496}{T}%
\TextAt{13}{270.06}{183.496}{r}%
\TextAt{21}{275.346}{181.637}{L}%
\TextAt{3}{281.106}{181.637}{\Character{106}}%
\TextAt{21}{283.458}{181.637}{K}%
\TextAt{9}{294.981}{183.496}{\Character{58}}%
\TextAt{18}{301.554}{183.496}{A}%
\TextAt{21}{310.188}{181.703}{L}%
\NativeArrowAt{319.766}{183.496}{11.955}{19.26}{\longrightarrow}%
\TextAt{18}{342.348}{183.496}{A}%
\TextAt{21}{350.982}{181.703}{K}%
\TextAt{8}{363.03}{183.496}{by}%
\TextAt{13}{379.237}{183.496}{T}%
\TextAt{13}{387.713}{183.496}{r}%
\TextAt{21}{392.999}{181.637}{L}%
\TextAt{3}{398.758}{181.637}{\Character{106}}%
\TextAt{21}{401.111}{181.637}{K}%
\TextAt{9}{409.313}{183.496}{\Character{40}}%
\TextAt{13}{413.866}{183.496}{\Character{11}}%
\TextAt{9}{421.387}{183.496}{\Character{41}}%
\TextAt{9}{429.242}{183.496}{\Character{61}}%
\TextAt{8}{442.21}{183.496}{trace}%
\TextAt{8}{471.424}{183.496}{of}%
\TextAt{8}{484.702}{183.496}{the}%
\TextAt{8}{504.806}{183.496}{endo-}%
\TextAt{8}{79.226}{166.45}{morphism}%
\TextAt{13}{134.158}{166.45}{x}%
\NativeArrowAt{144.534}{166.45}{11.955}{11.955}{\mapsto}%
\TextAt{13}{160.213}{166.45}{\Character{11}x}%
\TextAt{8}{178.521}{166.45}{of}%
\TextAt{8}{192.087}{166.45}{the}%
\TextAt{18}{212.48}{166.45}{A}%
\TextAt{21}{221.115}{164.657}{K}%
\TextAt{8}{229.317}{166.45}{-algebra}%
\TextAt{18}{274.434}{166.45}{A}%
\TextAt{21}{283.068}{164.657}{L}%
\TextAt{8}{289.325}{166.45}{,}%
\TextAt{8}{296.77}{166.45}{and}%
\TextAt{8}{319.769}{166.45}{show:}%
\TextAt{8}{353.797}{166.45}{(i)}%
\TextAt{13}{370.292}{166.45}{T}%
\TextAt{13}{378.768}{166.45}{r}%
\TextAt{21}{384.054}{164.591}{L}%
\TextAt{3}{389.813}{164.591}{\Character{106}}%
\TextAt{21}{392.166}{164.591}{K}%
\TextAt{9}{400.368}{166.45}{\Character{40}}%
\TextAt{13}{404.921}{166.45}{\Character{11}}%
\TextAt{9}{412.443}{166.45}{\Character{41}}%
\TextAt{27}{416.995}{164.657}{p}%
\TextAt{9}{425.206}{166.45}{\Character{61}}%
\TextAt{20}{438.034}{175.417}{\Character{80}}%
\TextAt{27}{450.654}{162.963}{P}%
\TextAt{3}{457.249}{162.963}{\Character{106}}%
\TextAt{27}{459.602}{162.963}{p}%
\TextAt{13}{466.081}{166.45}{T}%
\TextAt{13}{474.545}{166.45}{r}%
\TextAt{21}{479.843}{164.591}{L}%
\TextAt{28}{485.602}{163.2}{P}%
\TextAt{3}{491.799}{164.591}{\Character{106}}%
\TextAt{21}{494.152}{164.591}{K}%
\TextAt{28}{501.293}{163.594}{p}%
\TextAt{9}{505.874}{166.45}{\Character{40}}%
\TextAt{13}{510.426}{166.45}{\Character{11}}%
\TextAt{27}{517.872}{164.657}{P}%
\TextAt{9}{524.965}{166.45}{\Character{41}}%
\TextAt{8}{529.518}{166.45}{.}%
\TextAt{8}{79.226}{149.404}{(ii)}%
\TextAt{8}{99.075}{149.404}{For}%
\TextAt{8}{120.383}{149.404}{a}%
\TextAt{8}{130.477}{149.404}{tower}%
\TextAt{8}{162.688}{149.404}{of}%
\TextAt{8}{176.357}{149.404}{fields}%
\TextAt{13}{206.667}{149.404}{K}%
\NativeMathAt{221.382}{149.404}{11.955}{\subseteq}%
\TextAt{13}{234.578}{149.404}{L}%
\NativeMathAt{246.44}{149.404}{11.955}{\subseteq}%
\TextAt{13}{259.636}{149.404}{M}%
\TextAt{8}{272.21}{149.404}{,}%
\TextAt{8}{279.782}{149.404}{one}%
\TextAt{8}{301.58}{149.404}{has}%
\TextAt{13}{322.793}{149.404}{T}%
\TextAt{13}{331.257}{149.404}{r}%
\TextAt{21}{336.555}{147.545}{M}%
\TextAt{3}{345.536}{147.545}{\Character{106}}%
\TextAt{21}{347.889}{147.545}{K}%
\TextAt{9}{359.989}{149.404}{\Character{61}}%
\TextAt{13}{372.991}{149.404}{T}%
\TextAt{13}{381.467}{149.404}{r}%
\TextAt{21}{386.754}{147.545}{L}%
\TextAt{3}{392.513}{147.545}{\Character{106}}%
\TextAt{21}{394.865}{147.545}{K}%
\TextAt{16}{405.955}{149.404}{\Character{14}}%
\TextAt{13}{414.82}{149.404}{T}%
\TextAt{13}{423.284}{149.404}{r}%
\TextAt{21}{428.582}{147.545}{M}%
\TextAt{3}{437.564}{147.545}{\Character{106}}%
\TextAt{21}{439.916}{147.545}{L}%
\TextAt{8}{446.173}{149.404}{.}%
\TextAt{8}{455.637}{149.404}{(iii)}%
\TextAt{8}{478.738}{149.404}{If}%
\TextAt{13}{490.769}{149.404}{L}%
\TextAt{16}{502.631}{149.404}{\Character{106}}%
\TextAt{13}{509.849}{149.404}{K}%
\TextAt{8}{524.903}{149.404}{is}%
\TextAt{8}{79.226}{132.358}{embedded}%
\TextAt{8}{134.895}{132.358}{into}%
\TextAt{8}{159.038}{132.358}{the}%
\TextAt{8}{179.594}{132.358}{Galois}%
\TextAt{8}{215.907}{132.358}{extension}%
\TextAt{13}{268.073}{132.358}{M}%
\TextAt{16}{284.649}{132.358}{\Character{106}}%
\TextAt{13}{291.972}{132.358}{K}%
\TextAt{8}{302.79}{132.358}{,}%
\TextAt{8}{310.439}{132.358}{and}%
\TextAt{8}{333.602}{132.358}{if}%
\TextAt{13}{344.726}{132.358}{G}%
\TextAt{9}{357.961}{132.358}{\Character{61}}%
\TextAt{13}{371.069}{132.358}{G}%
\TextAt{9}{380.302}{132.358}{\Character{40}}%
\TextAt{13}{384.854}{132.358}{M}%
\TextAt{16}{401.43}{132.358}{\Character{106}}%
\TextAt{13}{408.754}{132.358}{K}%
\TextAt{9}{419.571}{132.358}{\Character{41}}%
\TextAt{8}{428.421}{132.358}{and}%
\TextAt{13}{451.584}{132.358}{H}%
\TextAt{9}{466.209}{132.358}{\Character{61}}%
\TextAt{13}{479.316}{132.358}{G}%
\TextAt{9}{488.549}{132.358}{\Character{40}}%
\TextAt{13}{493.102}{132.358}{M}%
\TextAt{16}{509.678}{132.358}{\Character{106}}%
\TextAt{13}{517.001}{132.358}{L}%
\TextAt{9}{524.965}{132.358}{\Character{41}}%
\TextAt{8}{529.518}{132.358}{,}%
\TextAt{8}{79.226}{115.312}{then}%
\TextAt{8}{105.925}{115.312}{one}%
\TextAt{8}{127.424}{115.312}{has}%
\TextAt{8}{148.337}{115.312}{for}%
\TextAt{13}{166.262}{115.312}{\Character{11}}%
\NativeMathAt{177.171}{115.312}{11.955}{\in}%
\TextAt{18}{188.529}{115.312}{A}%
\TextAt{21}{197.163}{113.519}{L}%
\TextAt{8}{203.42}{115.312}{,}%
\TextAt{13}{210.619}{115.312}{T}%
\TextAt{13}{219.083}{115.312}{r}%
\TextAt{21}{224.381}{113.453}{L}%
\TextAt{3}{230.14}{113.453}{\Character{106}}%
\TextAt{21}{232.492}{113.453}{K}%
\TextAt{9}{240.695}{115.312}{\Character{40}}%
\TextAt{13}{245.247}{115.312}{\Character{11}}%
\TextAt{9}{252.769}{115.312}{\Character{41}}%
\TextAt{9}{260.695}{115.312}{\Character{61}}%
\TextAt{20}{273.202}{124.279}{\Character{80}}%
\TextAt{21}{285.821}{111.825}{\Character{27}}%
\NativeMathAt{290.946}{111.825}{7.97}{\in}%
\TextAt{21}{296.592}{111.825}{G}%
\TextAt{17}{303.19}{111.825}{\Character{47}}%
\TextAt{21}{307.425}{111.825}{H}%
\TextAt{13}{317.486}{115.312}{\Character{27}}%
\TextAt{13}{324.552}{115.312}{\Character{11}}%
\TextAt{8}{332.09}{115.312}{.}%
\TextAt{8}{340.656}{115.312}{(iv)}%
\TextAt{13}{363.135}{115.312}{T}%
\TextAt{13}{371.611}{115.312}{r}%
\TextAt{21}{376.897}{113.453}{L}%
\TextAt{3}{382.656}{113.453}{\Character{106}}%
\TextAt{21}{385.009}{113.453}{K}%
\TextAt{9}{393.211}{115.312}{\Character{40}}%
\TextAt{13}{397.764}{115.312}{\Character{11}}%
\TextAt{9}{405.285}{115.312}{\Character{41}}%
\TextAt{9}{413.211}{115.312}{\Character{61}}%
\TextAt{9}{425.692}{115.312}{\Character{91}}%
\TextAt{13}{428.97}{115.312}{L}%
\TextAt{9}{440.322}{115.312}{\Character{58}}%
\TextAt{13}{446.961}{115.312}{K}%
\TextAt{9}{457.779}{115.312}{\Character{93}}%
\TextAt{13}{461.031}{115.312}{\Character{11}}%
\TextAt{8}{472.489}{115.312}{for}%
\TextAt{13}{490.413}{115.312}{\Character{11}}%
\NativeMathAt{501.323}{115.312}{11.955}{\in}%
\TextAt{18}{512.681}{115.312}{A}%
\TextAt{21}{521.315}{113.519}{K}%
\TextAt{8}{529.518}{115.312}{.}%
\TextAt{8}{79.226}{98.266}{(v)}%
\TextAt{8}{98.278}{98.266}{The}%
\TextAt{8}{122.195}{98.266}{trace}%
\TextAt{8}{151.325}{98.266}{of}%
\TextAt{8}{164.518}{98.266}{a}%
\TextAt{8}{174.137}{98.266}{principal}%
\TextAt{8}{222.777}{98.266}{adèle}%
\TextAt{13}{252.553}{98.266}{x}%
\NativeMathAt{262.526}{98.266}{11.955}{\in}%
\TextAt{13}{273.817}{98.266}{L}%
\TextAt{8}{285.542}{98.266}{is}%
\TextAt{8}{297.169}{98.266}{the}%
\TextAt{8}{317.189}{98.266}{principal}%
\TextAt{8}{365.83}{98.266}{adèle}%
\TextAt{8}{395.605}{98.266}{in}%
\TextAt{18}{409.121}{98.266}{A}%
\TextAt{21}{417.756}{96.473}{K}%
\TextAt{8}{429.719}{98.266}{defined}%
\TextAt{8}{469.895}{98.266}{by}%
\TextAt{8}{486.018}{98.266}{the}%
\TextAt{8}{506.038}{98.266}{usual}%
\TextAt{8}{79.226}{81.22}{trace}%
\TextAt{13}{108.492}{81.22}{T}%
\TextAt{13}{116.968}{81.22}{r}%
\TextAt{21}{122.255}{79.361}{L}%
\TextAt{3}{128.014}{79.361}{\Character{106}}%
\TextAt{21}{130.366}{79.361}{K}%
\TextAt{9}{138.569}{81.22}{\Character{40}}%
\TextAt{13}{143.121}{81.22}{x}%
\TextAt{9}{149.773}{81.22}{\Character{41}}%
\TextAt{8}{154.325}{81.22}{.}%
\TextAt{8}{297.211}{43.033}{340}%
\Destination{78.23}{809.815}{nav.page.340}%
\Destination{79.226}{552.355}{nav.ex.6.2.1}%
\Destination{79.226}{447.055}{nav.ex.6.2.2}%
\Destination{79.226}{353.755}{nav.ex.6.2.3}%
\Destination{79.226}{260.555}{nav.ex.6.2.4}%
\Destination{79.226}{197.455}{nav.ex.6.2.5}%
\Destination{79.36}{642.155}{nav.print.373}%
\Hotlink{224.753}{686.555}{13.995}{11.047}{nav.chapter.4}%
\Hotlink{256.87}{273.255}{14.306}{11.047}{nav.chapter.3}%
\Hotlink{283.004}{273.255}{6.858}{11.047}{nav.section.3.1}%
\Hotlink{244.696}{686.555}{25.078}{11.047}{nav.number.4.3.8}%
\Hotlink{504.94}{624.955}{25.078}{11.047}{nav.number.6.2.3}%
\Hotlink{231.406}{560.655}{25.078}{11.047}{nav.number.6.2.3}%
}

% Page 353 (body page 341)
\Page{%
\TextAt{2}{79.226}{796.943}{VI.}%
\TextAt{2}{98.012}{796.943}{Global}%
\TextAt{2}{133.248}{796.943}{Class}%
\TextAt{2}{161.841}{796.943}{Field}%
\TextAt{2}{189.572}{796.943}{Theory}%
\Vector{q 0 G 0 g 0.299 w 7.226 18.729 m 460.769 18.729 l S Q}%
\TextAt{15}{79.226}{760.804}{§}%
\TextAt{15}{93.785}{760.804}{3}%
\TextAt{15}{120.693}{760.804}{The}%
\TextAt{15}{159.571}{760.804}{Herbrand}%
\TextAt{15}{246.911}{760.804}{Quotient}%
\TextAt{15}{327.262}{760.804}{of}%
\TextAt{15}{348.701}{760.804}{the}%
\TextAt{15}{381.658}{760.804}{Idèle}%
\TextAt{15}{428.507}{760.804}{Class}%
\TextAt{15}{478.42}{760.804}{Group}%
\TextAt{8}{79.226}{728.918}{Our}%
\TextAt{8}{102.209}{728.918}{goal}%
\TextAt{8}{125.85}{728.918}{now}%
\TextAt{8}{149.156}{728.918}{is}%
\TextAt{8}{159.837}{728.918}{to}%
\TextAt{8}{173.065}{728.918}{show}%
\TextAt{8}{200.986}{728.918}{that}%
\TextAt{8}{225.272}{728.918}{the}%
\TextAt{8}{244.346}{728.918}{idèle}%
\TextAt{8}{270.569}{728.918}{class}%
\TextAt{8}{296.923}{728.918}{group}%
\TextAt{8}{329.016}{728.918}{satisfies}%
\TextAt{8}{371.044}{728.918}{the}%
\TextAt{8}{390.118}{728.918}{class}%
\TextAt{8}{416.472}{728.918}{field}%
\TextAt{8}{440.747}{728.918}{axiom}%
\TextAt{8}{474.465}{728.918}{of}%
\TextAt{8}{486.713}{728.918}{chap.}%
\TextAt{8}{516.522}{728.918}{IV,}%
\TextAt{8}{79.226}{711.872}{(6.1).}%
\TextAt{8}{111.633}{711.872}{To}%
\TextAt{8}{128.503}{711.872}{do}%
\TextAt{8}{144.405}{711.872}{this}%
\TextAt{8}{166.87}{711.872}{we}%
\TextAt{8}{183.74}{711.872}{will}%
\TextAt{8}{205.488}{711.872}{first}%
\TextAt{8}{229.256}{711.872}{compute}%
\TextAt{8}{276.373}{711.872}{its}%
\TextAt{8}{292.334}{711.872}{Herbrand}%
\TextAt{8}{344.329}{711.872}{quotient.}%
\TextAt{8}{394.943}{711.872}{It}%
\TextAt{8}{407.259}{711.872}{is}%
\TextAt{8}{418.665}{711.872}{constituted}%
\TextAt{8}{479.506}{711.872}{on}%
\TextAt{8}{495.408}{711.872}{the}%
\TextAt{8}{515.207}{711.872}{one}%
\TextAt{8}{79.226}{694.826}{hand}%
\TextAt{8}{108.714}{694.826}{by}%
\TextAt{8}{125.195}{694.826}{the}%
\TextAt{8}{145.573}{694.826}{Herbrand}%
\TextAt{8}{198.147}{694.826}{quotient}%
\TextAt{8}{244.551}{694.826}{of}%
\TextAt{8}{258.103}{694.826}{the}%
\TextAt{8}{278.481}{694.826}{idèle}%
\TextAt{8}{306.009}{694.826}{group,}%
\TextAt{8}{342.713}{694.826}{and}%
\TextAt{8}{365.698}{694.826}{by}%
\TextAt{8}{382.178}{694.826}{that}%
\TextAt{8}{407.769}{694.826}{of}%
\TextAt{8}{421.321}{694.826}{the}%
\TextAt{8}{441.699}{694.826}{unit}%
\TextAt{8}{466.632}{694.826}{group}%
\TextAt{8}{500.03}{694.826}{on}%
\TextAt{8}{516.51}{694.826}{the}%
\TextAt{8}{79.226}{677.78}{other.}%
\TextAt{8}{114.346}{677.78}{We}%
\TextAt{8}{134.491}{677.78}{study}%
\TextAt{8}{166.746}{677.78}{the}%
\TextAt{8}{186.902}{677.78}{idèle}%
\TextAt{8}{214.208}{677.78}{group}%
\TextAt{8}{247.384}{677.78}{first.}%
\TextAt{8}{97.159}{657.746}{Let}%
\TextAt{13}{117.717}{657.746}{L}%
\TextAt{16}{129.002}{657.746}{\Character{106}}%
\TextAt{13}{135.644}{657.746}{K}%
\TextAt{8}{149.948}{657.746}{be}%
\TextAt{8}{165.461}{657.746}{a}%
\TextAt{8}{174.805}{657.746}{finite}%
\TextAt{8}{204.306}{657.746}{Galois}%
\TextAt{8}{239.808}{657.746}{extension}%
\TextAt{8}{291.163}{657.746}{with}%
\TextAt{8}{317.412}{657.746}{Galois}%
\TextAt{8}{352.914}{657.746}{group}%
\TextAt{13}{385.678}{657.746}{G}%
\TextAt{8}{394.912}{657.746}{.}%
\TextAt{8}{403.223}{657.746}{The}%
\TextAt{13}{426.866}{657.746}{G}%
\TextAt{8}{436.099}{657.746}{-module}%
\TextAt{13}{480.879}{657.746}{I}%
\TextAt{21}{486.042}{655.952}{L}%
\TextAt{8}{495.785}{657.746}{may}%
\TextAt{8}{520.743}{657.746}{be}%
\TextAt{8}{79.226}{640.7}{described}%
\TextAt{8}{131.668}{640.7}{in}%
\TextAt{8}{146.009}{640.7}{the}%
\TextAt{8}{166.853}{640.7}{following}%
\TextAt{8}{216.976}{640.7}{simple}%
\TextAt{8}{254.139}{640.7}{manner,}%
\TextAt{8}{300.525}{640.7}{which}%
\TextAt{8}{334.699}{640.7}{immediately}%
\TextAt{8}{402.049}{640.7}{reduces}%
\TextAt{8}{444.413}{640.7}{us}%
\TextAt{8}{460.117}{640.7}{to}%
\TextAt{8}{475.115}{640.7}{local}%
\TextAt{8}{503.443}{640.7}{fields.}%
\TextAt{8}{79.226}{623.654}{For}%
\TextAt{8}{100.196}{623.654}{every}%
\TextAt{8}{131.088}{623.654}{place}%
\TextAt{14}{161}{623.654}{p}%
\TextAt{8}{170.88}{623.654}{of}%
\TextAt{13}{184.21}{623.654}{K}%
\TextAt{8}{198.925}{623.654}{we}%
\TextAt{8}{216.152}{623.654}{put}%
\TextAt{13}{218.296}{594.653}{L}%
\TextAt{3}{226.26}{599.589}{\Character{3}}%
\TextAt{27}{226.26}{591.697}{p}%
\TextAt{9}{234.313}{594.653}{\Character{61}}%
\TextAt{20}{246.739}{606.01}{\Character{89}}%
\TextAt{27}{247.909}{580.107}{P}%
\TextAt{3}{254.504}{580.107}{\Character{106}}%
\TextAt{27}{256.856}{580.107}{p}%
\TextAt{13}{264.007}{594.653}{L}%
\TextAt{3}{271.972}{599.589}{\Character{3}}%
\TextAt{27}{271.972}{591.697}{P}%
\TextAt{8}{291.021}{594.653}{and}%
\TextAt{13}{321.841}{594.653}{U}%
\TextAt{21}{329.791}{592.859}{L\Character{59}}%
\TextAt{27}{337.903}{592.859}{p}%
\TextAt{9}{345.71}{594.653}{\Character{61}}%
\TextAt{20}{358.136}{606.01}{\Character{89}}%
\TextAt{27}{359.306}{580.107}{P}%
\TextAt{3}{365.901}{580.107}{\Character{106}}%
\TextAt{27}{368.253}{580.107}{p}%
\TextAt{13}{375.404}{594.653}{U}%
\TextAt{27}{383.355}{592.859}{P}%
\TextAt{13}{390.448}{594.653}{\Character{58}}%
\TextAt{8}{79.226}{555.551}{Since}%
\TextAt{8}{109.494}{555.551}{the}%
\TextAt{8}{129.361}{555.551}{automorphisms}%
\TextAt{13}{211.155}{555.551}{\Character{27}}%
\NativeMathAt{221.559}{555.551}{11.955}{\in}%
\TextAt{13}{232.85}{555.551}{G}%
\TextAt{8}{245.691}{555.551}{permute}%
\TextAt{8}{291.572}{555.551}{the}%
\TextAt{8}{311.439}{555.551}{places}%
\TextAt{8}{345.676}{555.551}{of}%
\TextAt{13}{358.717}{555.551}{L}%
\TextAt{8}{370.289}{555.551}{above}%
\TextAt{14}{403.175}{555.551}{p}%
\TextAt{8}{409.158}{555.551}{,}%
\TextAt{8}{416.076}{555.551}{the}%
\TextAt{8}{435.943}{555.551}{groups}%
\TextAt{13}{473.443}{555.551}{L}%
\TextAt{3}{481.408}{559.89}{\Character{3}}%
\TextAt{27}{481.408}{552.596}{p}%
\TextAt{8}{489.748}{555.551}{and}%
\TextAt{13}{512.221}{555.551}{U}%
\TextAt{21}{520.171}{553.758}{L\Character{59}}%
\TextAt{27}{528.283}{553.758}{p}%
\TextAt{8}{79.226}{538.505}{are}%
\TextAt{13}{99.051}{538.505}{G}%
\TextAt{8}{108.284}{538.505}{-modules,}%
\TextAt{8}{161.733}{538.505}{and}%
\TextAt{8}{184.809}{538.505}{we}%
\TextAt{8}{202.35}{538.505}{have}%
\TextAt{8}{229.659}{538.505}{for}%
\TextAt{8}{247.857}{538.505}{the}%
\TextAt{13}{268.327}{538.505}{G}%
\TextAt{8}{277.561}{538.505}{-module}%
\TextAt{13}{323.065}{538.505}{I}%
\TextAt{21}{328.228}{536.712}{L}%
\TextAt{8}{338.696}{538.505}{the}%
\TextAt{8}{359.166}{538.505}{decomposition}%
\TextAt{13}{436.614}{538.505}{I}%
\TextAt{21}{441.777}{536.712}{L}%
\TextAt{9}{451.889}{538.505}{\Character{61}}%
\TextAt{20}{464.849}{547.472}{\Character{81}}%
\TextAt{3}{476.14}{544.317}{\Character{48}}%
\TextAt{27}{476.14}{535.018}{p}%
\TextAt{13}{482.62}{538.505}{L}%
\TextAt{3}{490.584}{542.844}{\Character{3}}%
\TextAt{27}{490.584}{535.55}{p}%
\TextAt{8}{495.317}{538.505}{,}%
\TextAt{8}{502.858}{538.505}{where}%
\TextAt{8}{79.226}{521.459}{the}%
\TextAt{8}{99.549}{521.459}{restricted}%
\TextAt{8}{151.804}{521.459}{product}%
\TextAt{8}{195.87}{521.459}{is}%
\TextAt{8}{207.8}{521.459}{taken}%
\TextAt{8}{239.839}{521.459}{with}%
\TextAt{8}{266.665}{521.459}{respect}%
\TextAt{8}{306.881}{521.459}{to}%
\TextAt{8}{321.358}{521.459}{the}%
\TextAt{8}{341.681}{521.459}{subgroups}%
\TextAt{13}{397.26}{521.459}{U}%
\TextAt{21}{405.21}{519.666}{L\Character{59}}%
\TextAt{27}{413.321}{519.666}{p}%
\NativeMathAt{421.412}{521.459}{11.955}{\subseteq}%
\TextAt{13}{434.315}{521.459}{L}%
\TextAt{3}{442.28}{525.798}{\Character{3}}%
\TextAt{27}{442.28}{518.504}{p}%
\TextAt{8}{447.012}{521.459}{.}%
\TextAt{8}{455.96}{521.459}{Choose}%
\TextAt{8}{496.833}{521.459}{a}%
\TextAt{8}{506.755}{521.459}{place}%
\TextAt{14}{79.226}{504.413}{P}%
\TextAt{8}{93.215}{504.413}{of}%
\TextAt{13}{106.743}{504.413}{L}%
\TextAt{8}{118.804}{504.413}{above}%
\TextAt{14}{152.177}{504.413}{p}%
\TextAt{8}{158.16}{504.413}{,}%
\TextAt{8}{165.558}{504.413}{and}%
\TextAt{8}{188.518}{504.413}{let}%
\TextAt{13}{205.621}{504.413}{G}%
\TextAt{27}{214.855}{502.62}{P}%
\TextAt{9}{225.607}{504.413}{\Character{61}}%
\TextAt{13}{238.371}{504.413}{G}%
\TextAt{9}{247.604}{504.413}{\Character{40}}%
\TextAt{13}{252.157}{504.413}{L}%
\TextAt{27}{260.121}{502.62}{P}%
\TextAt{16}{270.874}{504.413}{\Character{106}}%
\TextAt{13}{277.853}{504.413}{K}%
\TextAt{27}{287.812}{502.62}{p}%
\TextAt{9}{292.299}{504.413}{\Character{41}}%
\NativeMathAt{300.51}{504.413}{11.955}{\subseteq}%
\TextAt{13}{313.467}{504.413}{G}%
\TextAt{8}{326.797}{504.413}{be}%
\TextAt{8}{342.919}{504.413}{its}%
\TextAt{8}{359.436}{504.413}{decomposition}%
\TextAt{8}{436.769}{504.413}{group.}%
\TextAt{8}{475.091}{504.413}{As}%
\TextAt{13}{492.576}{504.413}{\Character{27}}%
\TextAt{8}{503.754}{504.413}{varies}%
\TextAt{8}{79.226}{487.367}{over}%
\TextAt{8}{103.9}{487.367}{a}%
\TextAt{8}{113.283}{487.367}{system}%
\TextAt{8}{151.729}{487.367}{of}%
\TextAt{8}{164.687}{487.367}{representatives}%
\TextAt{8}{244.116}{487.367}{of}%
\TextAt{13}{257.074}{487.367}{G}%
\TextAt{9}{266.307}{487.367}{\Character{47}}%
\TextAt{13}{272.16}{487.367}{G}%
\TextAt{27}{281.394}{485.574}{P}%
\TextAt{8}{288.487}{487.367}{,}%
\TextAt{13}{295.339}{487.367}{\Character{27}}%
\TextAt{14}{302.421}{487.367}{P}%
\TextAt{8}{315.84}{487.367}{runs}%
\TextAt{8}{341.542}{487.367}{through}%
\TextAt{8}{385.404}{487.367}{the}%
\TextAt{8}{405.188}{487.367}{various}%
\TextAt{8}{444.889}{487.367}{places}%
\TextAt{8}{479.044}{487.367}{of}%
\TextAt{13}{492.002}{487.367}{L}%
\TextAt{8}{503.491}{487.367}{above}%
\TextAt{14}{79.226}{470.321}{p}%
\TextAt{8}{85.209}{470.321}{,}%
\TextAt{8}{92.358}{470.321}{and}%
\TextAt{8}{115.121}{470.321}{we}%
\TextAt{8}{132.348}{470.321}{get}%
\TextAt{13}{155.728}{441.32}{L}%
\TextAt{3}{163.692}{446.256}{\Character{3}}%
\TextAt{27}{163.692}{438.365}{p}%
\TextAt{9}{171.746}{441.32}{\Character{61}}%
\TextAt{20}{184.171}{452.678}{\Character{89}}%
\TextAt{21}{189.247}{427.572}{\Character{27}}%
\TextAt{13}{201.44}{441.32}{L}%
\TextAt{3}{209.404}{446.256}{\Character{3}}%
\TextAt{21}{209.404}{438.365}{\Character{27}}%
\TextAt{27}{214.53}{438.365}{P}%
\TextAt{9}{224.944}{441.32}{\Character{61}}%
\TextAt{20}{237.369}{452.678}{\Character{89}}%
\TextAt{21}{242.445}{427.572}{\Character{27}}%
\TextAt{13}{254.638}{441.32}{\Character{27}}%
\TextAt{9}{261.72}{441.32}{\Character{40}}%
\TextAt{13}{266.273}{441.32}{L}%
\TextAt{3}{274.237}{446.256}{\Character{3}}%
\TextAt{27}{274.237}{438.365}{P}%
\TextAt{9}{281.331}{441.32}{\Character{41}}%
\TextAt{13}{285.883}{441.32}{\Character{59}}%
\TextAt{13}{315.018}{441.32}{U}%
\TextAt{21}{322.988}{439.527}{L\Character{59}}%
\TextAt{27}{331.099}{439.527}{p}%
\TextAt{9}{338.907}{441.32}{\Character{61}}%
\TextAt{20}{351.332}{452.678}{\Character{89}}%
\TextAt{21}{356.408}{427.572}{\Character{27}}%
\TextAt{13}{368.601}{441.32}{U}%
\TextAt{21}{376.551}{439.527}{\Character{27}}%
\TextAt{27}{381.677}{439.527}{P}%
\TextAt{9}{392.091}{441.32}{\Character{61}}%
\TextAt{20}{404.517}{452.678}{\Character{89}}%
\TextAt{21}{409.592}{427.572}{\Character{27}}%
\TextAt{13}{421.785}{441.32}{\Character{27}}%
\TextAt{9}{428.868}{441.32}{\Character{40}}%
\TextAt{13}{433.42}{441.32}{U}%
\TextAt{27}{441.37}{439.527}{P}%
\TextAt{9}{448.464}{441.32}{\Character{41}}%
\TextAt{13}{453.016}{441.32}{\Character{58}}%
\TextAt{8}{79.226}{404.901}{In}%
\TextAt{8}{94.93}{404.901}{terms}%
\TextAt{8}{128.59}{404.901}{of}%
\TextAt{8}{143.002}{404.901}{the}%
\TextAt{8}{164.24}{404.901}{notion}%
\TextAt{8}{201.75}{404.901}{of}%
\TextAt{8}{216.162}{404.901}{induced}%
\TextAt{8}{260.809}{404.901}{module}%
\TextAt{8}{303.185}{404.901}{introduced}%
\TextAt{8}{362.799}{404.901}{in}%
\TextAt{8}{377.534}{404.901}{chap.}%
\TextAt{8}{409.509}{404.901}{IV,}%
\TextAt{8}{430.735}{404.901}{§7,}%
\TextAt{8}{450.774}{404.901}{we}%
\TextAt{8}{469.084}{404.901}{thus}%
\TextAt{8}{495.917}{404.901}{get}%
\TextAt{8}{516.51}{404.901}{the}%
\TextAt{8}{79.226}{387.855}{following}%
\TextAt{12}{79.226}{361.842}{Proposition}%
\TextAt{12}{151.937}{361.842}{(3.1).}%
\TextAt{13}{187.253}{361.842}{L}%
\TextAt{3}{195.218}{366.181}{\Character{3}}%
\TextAt{27}{195.218}{358.887}{p}%
\TextAt{8}{203.847}{361.842}{and}%
\TextAt{13}{226.61}{361.842}{U}%
\TextAt{21}{234.56}{360.049}{L\Character{59}}%
\TextAt{27}{242.672}{360.049}{p}%
\TextAt{8}{251.056}{361.842}{are}%
\TextAt{8}{270.567}{361.842}{the}%
\TextAt{8}{290.723}{361.842}{induced}%
\TextAt{13}{334.288}{361.842}{G}%
\TextAt{8}{343.521}{361.842}{-modules}%
\TextAt{13}{202.749}{332.841}{L}%
\TextAt{3}{210.713}{337.777}{\Character{3}}%
\TextAt{27}{210.713}{329.886}{p}%
\TextAt{9}{218.766}{332.841}{\Character{61}}%
\TextAt{8}{231.192}{332.841}{Ind}%
\TextAt{21}{248.419}{339.481}{G}%
\TextAt{28}{255.018}{338.09}{P}%
\TextAt{21}{248.419}{329.6}{G}%
\TextAt{9}{261.713}{332.841}{\Character{40}}%
\TextAt{13}{266.266}{332.841}{L}%
\TextAt{3}{274.23}{337.777}{\Character{3}}%
\TextAt{27}{274.23}{329.886}{P}%
\TextAt{9}{281.324}{332.841}{\Character{41}}%
\TextAt{13}{285.876}{332.841}{\Character{59}}%
\TextAt{13}{315.011}{332.841}{U}%
\TextAt{21}{322.981}{331.048}{L\Character{59}}%
\TextAt{27}{331.092}{331.048}{p}%
\TextAt{9}{338.9}{332.841}{\Character{61}}%
\TextAt{8}{351.325}{332.841}{Ind}%
\TextAt{21}{368.553}{339.481}{G}%
\TextAt{28}{375.151}{338.09}{P}%
\TextAt{21}{368.553}{329.6}{G}%
\TextAt{9}{381.847}{332.841}{\Character{40}}%
\TextAt{13}{386.399}{332.841}{U}%
\TextAt{27}{394.349}{331.048}{P}%
\TextAt{9}{401.443}{332.841}{\Character{41}}%
\TextAt{13}{405.995}{332.841}{\Character{58}}%
\TextAt{8}{97.159}{294.874}{Now}%
\TextAt{8}{123.694}{294.874}{let}%
\TextAt{13}{140.473}{294.874}{S}%
\TextAt{8}{152.141}{294.874}{be}%
\TextAt{8}{167.94}{294.874}{a}%
\TextAt{8}{177.57}{294.874}{finite}%
\TextAt{8}{207.357}{294.874}{set}%
\TextAt{8}{225.5}{294.874}{of}%
\TextAt{8}{238.704}{294.874}{places}%
\TextAt{8}{273.106}{294.874}{of}%
\TextAt{13}{286.311}{294.874}{K}%
\TextAt{8}{300.901}{294.874}{containing}%
\TextAt{8}{357.694}{294.874}{the}%
\TextAt{8}{377.725}{294.874}{infinite}%
\TextAt{8}{417.268}{294.874}{places.}%
\TextAt{8}{456.303}{294.874}{We}%
\TextAt{8}{476.323}{294.874}{then}%
\TextAt{8}{502.858}{294.874}{define}%
\TextAt{13}{79.226}{277.828}{I}%
\TextAt{21}{85.329}{282.166}{S}%
\TextAt{21}{84.389}{274.807}{L}%
\TextAt{9}{95.24}{277.828}{\Character{61}}%
\TextAt{13}{108.163}{277.828}{I}%
\Vector{q 0 G 0 g 0.359 w 42.265 -485.132 m 48.359 -485.132 l S Q}%
\TextAt{21}{114.265}{282.166}{S}%
\TextAt{21}{113.325}{274.807}{L}%
\TextAt{8}{120.359}{277.828}{,}%
\TextAt{8}{127.873}{277.828}{where}%
\Vector{q 0 G 0 g 0.478 w 89.973 -486.329 m 97.869 -486.329 l S Q}%
\TextAt{13}{161.973}{277.828}{S}%
\TextAt{8}{174.058}{277.828}{denotes}%
\TextAt{8}{216.682}{277.828}{the}%
\TextAt{8}{237.13}{277.828}{set}%
\TextAt{8}{255.689}{277.828}{of}%
\TextAt{8}{269.311}{277.828}{all}%
\TextAt{8}{285.861}{277.828}{places}%
\TextAt{8}{320.679}{277.828}{of}%
\TextAt{13}{334.301}{277.828}{L}%
\TextAt{8}{346.454}{277.828}{which}%
\TextAt{8}{380.232}{277.828}{lie}%
\TextAt{8}{396.125}{277.828}{above}%
\TextAt{8}{429.592}{277.828}{the}%
\TextAt{8}{450.04}{277.828}{places}%
\TextAt{8}{484.858}{277.828}{of}%
\TextAt{13}{498.479}{277.828}{S}%
\TextAt{8}{506.375}{277.828}{.}%
\TextAt{8}{515.698}{277.828}{For}%
\TextAt{13}{79.226}{260.782}{I}%
\TextAt{21}{85.329}{265.12}{S}%
\TextAt{21}{84.389}{257.761}{L}%
\TextAt{8}{95.32}{260.782}{we}%
\TextAt{8}{112.548}{260.782}{have}%
\TextAt{8}{139.542}{260.782}{the}%
\TextAt{13}{159.699}{260.782}{G}%
\TextAt{8}{168.932}{260.782}{-module}%
\TextAt{8}{214.123}{260.782}{decomposition}%
\TextAt{13}{249.203}{231.781}{I}%
\TextAt{21}{255.306}{236.717}{S}%
\TextAt{21}{254.366}{228.825}{L}%
\TextAt{9}{264.721}{231.781}{\Character{61}}%
\TextAt{20}{277.146}{243.138}{\Character{89}}%
\TextAt{27}{277.169}{217.766}{p}%
\NativeMathAt{281.158}{217.766}{7.97}{\in}%
\TextAt{21}{286.804}{217.766}{S}%
\TextAt{13}{294.415}{231.781}{L}%
\TextAt{3}{302.379}{236.717}{\Character{3}}%
\TextAt{27}{302.379}{228.825}{p}%
\TextAt{16}{309.768}{231.781}{\Character{2}}%
\TextAt{20}{321.724}{243.138}{\Character{89}}%
\TextAt{27}{321.747}{217.235}{p}%
\TextAt{17}{326.676}{217.235}{\Character{47}}%
\NativeMathAt{325.735}{217.235}{7.97}{\in}%
\TextAt{21}{331.381}{217.235}{S}%
\TextAt{13}{338.992}{231.781}{U}%
\TextAt{21}{346.942}{229.987}{L\Character{59}}%
\TextAt{27}{355.054}{229.987}{p}%
\TextAt{13}{359.541}{231.781}{\Character{59}}%
\TextAt{8}{79.226}{192.615}{and}%
\TextAt{8}{101.989}{192.615}{(3.1)}%
\TextAt{8}{129.964}{192.615}{gives}%
\TextAt{8}{158.644}{192.615}{the}%
\TextAt{12}{79.226}{166.603}{Proposition}%
\TextAt{12}{152.299}{166.603}{(3.2).}%
\TextAt{8}{187.929}{166.603}{If}%
\TextAt{13}{199.935}{166.603}{L}%
\TextAt{16}{211.756}{166.603}{\Character{106}}%
\TextAt{13}{218.934}{166.603}{K}%
\TextAt{8}{233.963}{166.603}{is}%
\TextAt{8}{246.041}{166.603}{a}%
\TextAt{8}{256.111}{166.603}{cyclic}%
\TextAt{8}{288.609}{166.603}{extension,}%
\TextAt{8}{344.019}{166.603}{and}%
\TextAt{8}{367.096}{166.603}{if}%
\TextAt{13}{378.134}{166.603}{S}%
\TextAt{8}{390.241}{166.603}{contains}%
\TextAt{8}{436.476}{166.603}{all}%
\TextAt{8}{453.049}{166.603}{primes}%
\TextAt{8}{491.142}{166.603}{ramified}%
\TextAt{8}{79.226}{149.557}{in}%
\TextAt{13}{92.879}{149.557}{L}%
\TextAt{8}{100.843}{149.557}{,}%
\TextAt{8}{107.993}{149.557}{then}%
\TextAt{8}{134.653}{149.557}{we}%
\TextAt{8}{151.88}{149.557}{have}%
\TextAt{8}{178.875}{149.557}{for}%
\TextAt{13}{196.76}{149.557}{i}%
\TextAt{9}{204.074}{149.557}{\Character{61}}%
\TextAt{9}{216.484}{149.557}{0}%
\TextAt{13}{222.353}{149.557}{\Character{59}}%
\TextAt{16}{227.597}{149.557}{\Character{0}}%
\TextAt{9}{236.895}{149.557}{1}%
\TextAt{8}{246.646}{149.557}{that}%
\TextAt{13}{137.853}{120.555}{H}%
\TextAt{21}{148.476}{125.492}{i}%
\TextAt{9}{151.857}{120.555}{\Character{40}}%
\TextAt{13}{156.409}{120.555}{G\Character{59}}%
\TextAt{13}{170.863}{120.555}{I}%
\TextAt{21}{176.99}{125.492}{S}%
\TextAt{21}{176.05}{117.6}{L}%
\TextAt{9}{183.084}{120.555}{\Character{41}}%
\NativeMathAt{190.957}{120.555}{11.955}{\cong}%
\TextAt{20}{203.576}{131.913}{\Character{77}}%
\TextAt{27}{204.994}{106.541}{p}%
\NativeMathAt{208.983}{106.541}{7.97}{\in}%
\TextAt{21}{214.628}{106.541}{S}%
\TextAt{13}{223.634}{120.555}{H}%
\TextAt{21}{234.257}{125.492}{i}%
\TextAt{9}{237.638}{120.555}{\Character{40}}%
\TextAt{13}{242.191}{120.555}{G}%
\TextAt{27}{251.424}{118.762}{P}%
\TextAt{13}{258.518}{120.555}{\Character{59}}%
\TextAt{13}{263.742}{120.555}{L}%
\TextAt{3}{271.727}{125.492}{\Character{3}}%
\TextAt{27}{271.727}{117.6}{P}%
\TextAt{9}{278.82}{120.555}{\Character{41}}%
\TextAt{8}{295.328}{120.555}{and}%
\TextAt{13}{326.148}{120.555}{H}%
\TextAt{21}{336.771}{125.492}{i}%
\TextAt{9}{340.152}{120.555}{\Character{40}}%
\TextAt{13}{344.704}{120.555}{G\Character{59}}%
\TextAt{13}{359.158}{120.555}{I}%
\TextAt{21}{364.345}{118.762}{L}%
\TextAt{9}{370.602}{120.555}{\Character{41}}%
\NativeMathAt{378.475}{120.555}{11.955}{\cong}%
\TextAt{20}{391.095}{131.913}{\Character{77}}%
\TextAt{27}{398.133}{106.807}{p}%
\TextAt{13}{411.153}{120.555}{H}%
\TextAt{21}{421.775}{125.492}{i}%
\TextAt{9}{425.157}{120.555}{\Character{40}}%
\TextAt{13}{429.709}{120.555}{G}%
\TextAt{27}{438.943}{118.762}{P}%
\TextAt{13}{446.036}{120.555}{\Character{59}}%
\TextAt{13}{451.26}{120.555}{L}%
\TextAt{3}{459.245}{125.492}{\Character{3}}%
\TextAt{27}{459.245}{117.6}{P}%
\TextAt{9}{466.338}{120.555}{\Character{41}}%
\TextAt{13}{470.891}{120.555}{\Character{59}}%
\TextAt{8}{79.226}{82.471}{where}%
\TextAt{8}{113.035}{82.471}{for}%
\TextAt{8}{130.92}{82.471}{each}%
\TextAt{14}{157.258}{82.471}{p}%
\TextAt{8}{163.241}{82.471}{,}%
\TextAt{14}{170.39}{82.471}{P}%
\TextAt{8}{184.18}{82.471}{is}%
\TextAt{8}{195.944}{82.471}{a}%
\TextAt{8}{205.7}{82.471}{chosen}%
\TextAt{8}{243.155}{82.471}{prime}%
\TextAt{8}{276.319}{82.471}{of}%
\TextAt{13}{289.649}{82.471}{L}%
\TextAt{8}{301.511}{82.471}{above}%
\TextAt{14}{334.686}{82.471}{p}%
\TextAt{8}{340.67}{82.471}{.}%
\TextAt{8}{297.211}{43.033}{341}%
\Destination{78.23}{809.815}{nav.page.341}%
\Destination{79.23}{773.075}{nav.section.6.3}%
\Destination{79.36}{775.755}{nav.print.374}%
\Destination{79.36}{375.855}{nav.number.6.3.1}%
\Destination{79.36}{180.555}{nav.number.6.3.2}%
\Destination{79.23}{773.075}{outline.58}%
\Hotlink{516.022}{726.755}{13.995}{11.047}{nav.chapter.4}%
\Hotlink{409.009}{402.655}{13.995}{11.047}{nav.chapter.4}%
\Hotlink{435.914}{402.655}{6.858}{11.047}{nav.section.4.7}%
\Hotlink{78.726}{709.655}{25.078}{11.047}{nav.number.6.6.1}%
\Hotlink{101.489}{190.455}{25.078}{11.047}{nav.number.6.3.1}%
}

% Page 354 (body page 342)
\Page{%
\TextAt{2}{79.226}{796.943}{VI.}%
\TextAt{2}{98.012}{796.943}{Global}%
\TextAt{2}{133.248}{796.943}{Class}%
\TextAt{2}{161.841}{796.943}{Field}%
\TextAt{2}{189.572}{796.943}{Theory}%
\Vector{q 0 G 0 g 0.299 w 7.226 18.729 m 460.769 18.729 l S Q}%
\TextAt{19}{79.226}{761.124}{Proof.}%
\TextAt{8}{116}{761.124}{The}%
\TextAt{8}{139.8}{761.124}{decomposition}%
\TextAt{13}{216.681}{761.124}{I}%
\TextAt{21}{222.784}{765.462}{S}%
\TextAt{21}{221.844}{758.103}{L}%
\TextAt{9}{232.199}{761.124}{\Character{61}}%
\TextAt{9}{244.609}{761.124}{\Character{40}}%
\TextAt{20}{249.176}{770.09}{\Character{76}}%
\TextAt{27}{262.46}{757.637}{p}%
\NativeMathAt{266.449}{757.637}{7.97}{\in}%
\TextAt{21}{272.094}{757.637}{S}%
\TextAt{13}{280.181}{761.124}{L}%
\TextAt{3}{288.145}{765.462}{\Character{3}}%
\TextAt{27}{288.145}{758.168}{p}%
\TextAt{9}{292.878}{761.124}{\Character{41}}%
\TextAt{16}{299.567}{761.124}{\Character{8}}%
\TextAt{13}{311.004}{761.124}{V}%
\TextAt{8}{320.429}{761.124}{,}%
\TextAt{13}{327.375}{761.124}{V}%
\TextAt{9}{340.122}{761.124}{\Character{61}}%
\TextAt{20}{352.547}{770.09}{\Character{81}}%
\TextAt{27}{363.838}{757.637}{p}%
\TextAt{17}{368.768}{757.637}{\Character{47}}%
\NativeMathAt{367.827}{757.637}{7.97}{\in}%
\TextAt{21}{373.473}{757.637}{S}%
\TextAt{13}{381.559}{761.124}{U}%
\TextAt{21}{389.509}{759.33}{L\Character{59}}%
\TextAt{27}{397.621}{759.33}{p}%
\TextAt{8}{402.108}{761.124}{,}%
\TextAt{8}{409.054}{761.124}{gives}%
\TextAt{8}{437.48}{761.124}{us}%
\TextAt{8}{452.242}{761.124}{an}%
\TextAt{8}{468.247}{761.124}{isomorphism}%
\TextAt{13}{205.528}{730.625}{H}%
\TextAt{21}{216.151}{735.562}{i}%
\TextAt{9}{219.532}{730.625}{\Character{40}}%
\TextAt{13}{224.084}{730.625}{G\Character{59}}%
\TextAt{13}{238.538}{730.625}{I}%
\TextAt{21}{244.665}{735.562}{S}%
\TextAt{21}{243.725}{727.67}{L}%
\TextAt{9}{250.759}{730.625}{\Character{41}}%
\NativeMathAt{258.632}{730.625}{11.955}{\cong}%
\TextAt{20}{271.251}{741.983}{\Character{77}}%
\TextAt{27}{272.669}{716.611}{p}%
\NativeMathAt{276.658}{716.611}{7.97}{\in}%
\TextAt{21}{282.303}{716.611}{S}%
\TextAt{13}{291.309}{730.625}{H}%
\TextAt{21}{301.932}{735.562}{i}%
\TextAt{9}{305.313}{730.625}{\Character{40}}%
\TextAt{13}{309.865}{730.625}{G\Character{59}}%
\TextAt{13}{324.319}{730.625}{L}%
\TextAt{3}{332.308}{735.562}{\Character{3}}%
\TextAt{27}{332.308}{727.67}{p}%
\TextAt{9}{337.04}{730.625}{\Character{41}}%
\TextAt{16}{344.249}{730.625}{\Character{8}}%
\TextAt{13}{356.204}{730.625}{H}%
\TextAt{21}{366.827}{735.562}{i}%
\TextAt{9}{370.208}{730.625}{\Character{40}}%
\TextAt{13}{374.76}{730.625}{G\Character{59}}%
\TextAt{13}{389.214}{730.625}{V}%
\TextAt{9}{398.664}{730.625}{\Character{41}}%
\TextAt{13}{403.216}{730.625}{\Character{59}}%
\TextAt{8}{79.226}{691.35}{and}%
\TextAt{8}{102.381}{691.35}{an}%
\TextAt{8}{119.031}{691.35}{injection}%
\TextAt{13}{167.22}{691.35}{H}%
\TextAt{21}{177.842}{695.689}{i}%
\TextAt{9}{181.224}{691.35}{\Character{40}}%
\TextAt{13}{185.776}{691.35}{G\Character{59}}%
\TextAt{13}{200.23}{691.35}{V}%
\TextAt{9}{209.68}{691.35}{\Character{41}}%
\NativeArrowAt{218.22}{691.35}{11.955}{19.26}{\longrightarrow}%
\TextAt{20}{241.47}{700.317}{\Character{81}}%
\TextAt{27}{252.761}{687.863}{p}%
\TextAt{17}{257.691}{687.863}{\Character{47}}%
\NativeMathAt{256.75}{687.863}{7.97}{\in}%
\TextAt{21}{262.395}{687.863}{S}%
\TextAt{13}{270.482}{691.35}{H}%
\TextAt{21}{281.104}{695.689}{i}%
\TextAt{9}{284.485}{691.35}{\Character{40}}%
\TextAt{13}{289.038}{691.35}{G\Character{59}}%
\TextAt{13}{303.492}{691.35}{U}%
\TextAt{21}{311.466}{689.557}{L\Character{59}}%
\TextAt{27}{319.577}{689.557}{p}%
\TextAt{9}{324.064}{691.35}{\Character{41}}%
\TextAt{8}{328.616}{691.35}{.}%
\TextAt{8}{338.24}{691.35}{By}%
\TextAt{8}{356.995}{691.35}{(3.1)}%
\TextAt{8}{385.362}{691.35}{and}%
\TextAt{8}{408.516}{691.35}{chap.}%
\TextAt{8}{439.8}{691.35}{IV,}%
\TextAt{8}{460.336}{691.35}{(7.4),}%
\TextAt{8}{492.053}{691.35}{we}%
\TextAt{8}{509.672}{691.35}{have}%
\TextAt{8}{79.226}{674.304}{the}%
\TextAt{8}{100.083}{674.304}{isomorphisms}%
\TextAt{13}{173.818}{674.304}{H}%
\TextAt{21}{184.441}{678.643}{i}%
\TextAt{9}{187.822}{674.304}{\Character{40}}%
\TextAt{13}{192.375}{674.304}{G\Character{59}}%
\TextAt{13}{206.829}{674.304}{L}%
\TextAt{3}{214.817}{678.643}{\Character{3}}%
\TextAt{27}{214.817}{671.349}{p}%
\TextAt{9}{219.549}{674.304}{\Character{41}}%
\NativeMathAt{228.617}{674.304}{11.955}{\cong}%
\TextAt{13}{242.43}{674.304}{H}%
\TextAt{21}{253.053}{678.643}{i}%
\TextAt{9}{256.434}{674.304}{\Character{40}}%
\TextAt{13}{260.987}{674.304}{G}%
\TextAt{27}{270.22}{672.511}{P}%
\TextAt{13}{277.314}{674.304}{\Character{59}}%
\TextAt{13}{282.538}{674.304}{L}%
\TextAt{3}{290.523}{678.643}{\Character{3}}%
\TextAt{27}{290.523}{671.219}{P}%
\TextAt{9}{297.616}{674.304}{\Character{41}}%
\TextAt{8}{306.767}{674.304}{and}%
\TextAt{13}{330.23}{674.304}{H}%
\TextAt{21}{340.853}{678.643}{i}%
\TextAt{9}{344.234}{674.304}{\Character{40}}%
\TextAt{13}{348.786}{674.304}{G\Character{59}}%
\TextAt{13}{363.24}{674.304}{U}%
\TextAt{21}{371.214}{672.511}{L\Character{59}}%
\TextAt{27}{379.326}{672.511}{p}%
\TextAt{9}{383.813}{674.304}{\Character{41}}%
\NativeMathAt{392.88}{674.304}{11.955}{\cong}%
\TextAt{13}{406.694}{674.304}{H}%
\TextAt{21}{417.316}{678.643}{i}%
\TextAt{9}{420.698}{674.304}{\Character{40}}%
\TextAt{13}{425.25}{674.304}{G}%
\TextAt{27}{434.484}{672.511}{P}%
\TextAt{13}{441.577}{674.304}{\Character{59}}%
\TextAt{13}{446.801}{674.304}{U}%
\TextAt{27}{454.772}{672.511}{P}%
\TextAt{9}{461.865}{674.304}{\Character{41}}%
\TextAt{8}{466.418}{674.304}{.}%
\TextAt{8}{476.969}{674.304}{For}%
\TextAt{14}{498.639}{674.304}{p}%
\NativeMathAt{509.137}{674.304}{11.955}{\notin}%
\TextAt{13}{521.622}{674.304}{S}%
\TextAt{8}{529.518}{674.304}{,}%
\TextAt{13}{79.226}{657.258}{L}%
\TextAt{27}{87.191}{655.465}{P}%
\TextAt{16}{98.769}{657.258}{\Character{106}}%
\TextAt{13}{106.575}{657.258}{K}%
\TextAt{27}{116.534}{655.465}{p}%
\TextAt{8}{125.602}{657.258}{is}%
\TextAt{8}{138.049}{657.258}{unramified.}%
\TextAt{8}{203.182}{657.258}{Hence}%
\TextAt{13}{238.642}{657.258}{H}%
\TextAt{21}{249.265}{661.597}{i}%
\TextAt{9}{252.646}{657.258}{\Character{40}}%
\TextAt{13}{257.199}{657.258}{G}%
\TextAt{27}{266.432}{655.465}{P}%
\TextAt{13}{273.526}{657.258}{\Character{59}}%
\TextAt{13}{278.75}{657.258}{U}%
\TextAt{27}{286.72}{655.465}{P}%
\TextAt{9}{293.814}{657.258}{\Character{41}}%
\TextAt{9}{302.84}{657.258}{\Character{61}}%
\TextAt{9}{316.421}{657.258}{1}%
\TextAt{8}{322.294}{657.258}{,}%
\TextAt{8}{330.297}{657.258}{by}%
\TextAt{8}{347.239}{657.258}{chap.}%
\TextAt{8}{378.815}{657.258}{V,}%
\TextAt{8}{395.422}{657.258}{(1.2).}%
\TextAt{8}{429.998}{657.258}{This}%
\TextAt{8}{457.401}{657.258}{shows}%
\TextAt{8}{491.702}{657.258}{the}%
\TextAt{8}{512.541}{657.258}{first}%
\TextAt{8}{79.226}{640.212}{claim}%
\TextAt{8}{110.441}{640.212}{of}%
\TextAt{8}{123.771}{640.212}{the}%
\TextAt{8}{143.928}{640.212}{proposition.}%
\TextAt{8}{210.012}{640.212}{The}%
\TextAt{8}{234.066}{640.212}{second}%
\TextAt{8}{271.844}{640.212}{is}%
\TextAt{8}{283.608}{640.212}{an}%
\TextAt{8}{299.867}{640.212}{immediate}%
\TextAt{8}{357.096}{640.212}{consequence:}%
\TextAt{13}{133.31}{611.211}{H}%
\TextAt{21}{143.933}{616.147}{i}%
\TextAt{9}{147.314}{611.211}{\Character{40}}%
\TextAt{13}{151.866}{611.211}{G\Character{59}}%
\TextAt{13}{166.32}{611.211}{I}%
\TextAt{21}{171.507}{609.418}{L}%
\TextAt{9}{177.764}{611.211}{\Character{41}}%
\TextAt{9}{185.619}{611.211}{\Character{61}}%
\NativeMathAt{198.063}{611.211}{11.955}{\varinjlim}%
\TextAt{21}{203.394}{598.831}{S}%
\TextAt{13}{216.314}{611.211}{H}%
\TextAt{21}{226.937}{616.147}{i}%
\TextAt{9}{230.318}{611.211}{\Character{40}}%
\TextAt{13}{234.871}{611.211}{G\Character{59}}%
\TextAt{13}{249.325}{611.211}{I}%
\TextAt{21}{255.451}{616.147}{S}%
\TextAt{21}{254.511}{608.256}{L}%
\TextAt{9}{261.545}{611.211}{\Character{41}}%
\TextAt{9}{269.4}{611.211}{\Character{61}}%
\NativeMathAt{281.844}{611.211}{11.955}{\varinjlim}%
\TextAt{21}{287.175}{598.831}{S}%
\TextAt{20}{300.095}{622.569}{\Character{77}}%
\TextAt{27}{301.513}{597.197}{p}%
\NativeMathAt{305.502}{597.197}{7.97}{\in}%
\TextAt{21}{311.147}{597.197}{S}%
\TextAt{13}{320.153}{611.211}{H}%
\TextAt{21}{330.776}{616.147}{i}%
\TextAt{9}{334.157}{611.211}{\Character{40}}%
\TextAt{13}{338.71}{611.211}{G}%
\TextAt{27}{347.943}{609.418}{P}%
\TextAt{13}{355.037}{611.211}{\Character{59}}%
\TextAt{13}{360.261}{611.211}{L}%
\TextAt{3}{368.245}{616.147}{\Character{3}}%
\TextAt{27}{368.245}{608.256}{P}%
\TextAt{9}{375.339}{611.211}{\Character{41}}%
\TextAt{9}{383.194}{611.211}{\Character{61}}%
\TextAt{20}{395.638}{622.569}{\Character{77}}%
\TextAt{27}{402.676}{597.463}{p}%
\TextAt{13}{415.696}{611.211}{H}%
\TextAt{21}{426.318}{616.147}{i}%
\TextAt{9}{429.7}{611.211}{\Character{40}}%
\TextAt{13}{434.252}{611.211}{G}%
\TextAt{27}{443.486}{609.418}{P}%
\TextAt{13}{450.579}{611.211}{\Character{59}}%
\TextAt{13}{455.803}{611.211}{L}%
\TextAt{3}{463.788}{616.147}{\Character{3}}%
\TextAt{27}{463.788}{608.256}{P}%
\TextAt{9}{470.882}{611.211}{\Character{41}}%
\TextAt{13}{475.434}{611.211}{\Character{58}}%
\Vector{q 0 G 0 g 0.399 w 452.334 -200.527 m 452.334 -192.457 l S Q}%
\Vector{q 0 G 0 g 0.399 w 452.534 -192.656 m 459.707 -192.656 l S Q}%
\Vector{q 0 G 0 g 0.399 w 452.534 -200.327 m 459.707 -200.327 l S Q}%
\Vector{q 0 G 0 g 0.399 w 459.906 -200.527 m 459.906 -192.457 l S Q}%
\TextAt{8}{97.159}{544.472}{The}%
\TextAt{8}{121.104}{544.472}{proposition}%
\TextAt{8}{182.528}{544.472}{says}%
\TextAt{8}{207.262}{544.472}{that}%
\TextAt{8}{232.521}{544.472}{one}%
\TextAt{8}{253.872}{544.472}{has}%
\TextAt{13}{274.636}{544.472}{H}%
\TextAt{3}{285.259}{548.811}{\Character{0}}%
\TextAt{17}{291.846}{548.811}{1}%
\TextAt{9}{296.578}{544.472}{\Character{40}}%
\TextAt{13}{301.13}{544.472}{G\Character{59}}%
\TextAt{13}{315.584}{544.472}{I}%
\TextAt{21}{320.771}{542.679}{L}%
\TextAt{9}{327.028}{544.472}{\Character{41}}%
\TextAt{9}{334.883}{544.472}{\Character{61}}%
\TextAt{16}{347.327}{544.472}{\Character{102}}%
\TextAt{9}{353.304}{544.472}{1}%
\TextAt{16}{359.157}{544.472}{\Character{103}}%
\TextAt{8}{365.135}{544.472}{,}%
\TextAt{8}{372.197}{544.472}{because}%
\TextAt{13}{415.389}{544.472}{H}%
\TextAt{3}{426.012}{548.811}{\Character{0}}%
\TextAt{17}{432.599}{548.811}{1}%
\TextAt{9}{437.331}{544.472}{\Character{40}}%
\TextAt{13}{441.883}{544.472}{G}%
\TextAt{27}{451.117}{542.679}{P}%
\TextAt{13}{458.21}{544.472}{\Character{59}}%
\TextAt{13}{463.434}{544.472}{L}%
\TextAt{3}{471.419}{548.811}{\Character{3}}%
\TextAt{27}{471.419}{541.387}{P}%
\TextAt{9}{478.513}{544.472}{\Character{41}}%
\TextAt{9}{486.368}{544.472}{\Character{61}}%
\TextAt{16}{498.811}{544.472}{\Character{102}}%
\TextAt{9}{504.789}{544.472}{1}%
\TextAt{16}{510.642}{544.472}{\Character{103}}%
\TextAt{8}{520.408}{544.472}{by}%
\TextAt{8}{79.226}{527.426}{Hilbert}%
\TextAt{8}{119.539}{527.426}{90.}%
\TextAt{8}{139.703}{527.426}{Further}%
\TextAt{8}{182.132}{527.426}{it}%
\TextAt{8}{193.836}{527.426}{says}%
\TextAt{8}{218.679}{527.426}{that}%
\TextAt{13}{222.277}{498.425}{I}%
\TextAt{21}{227.44}{496.632}{K}%
\TextAt{9}{235.642}{498.425}{\Character{47}}%
\TextAt{13}{241.495}{498.425}{N}%
\TextAt{21}{250.869}{496.565}{L}%
\TextAt{3}{256.628}{496.565}{\Character{106}}%
\TextAt{21}{258.981}{496.565}{K}%
\TextAt{13}{267.183}{498.425}{I}%
\TextAt{21}{272.346}{496.632}{L}%
\NativeMathAt{281.924}{498.425}{11.955}{\cong}%
\TextAt{20}{294.543}{509.783}{\Character{77}}%
\TextAt{27}{301.582}{484.677}{p}%
\TextAt{13}{314.602}{498.425}{K}%
\TextAt{3}{325.419}{503.361}{\Character{3}}%
\TextAt{27}{324.561}{495.47}{p}%
\TextAt{9}{330.152}{498.425}{\Character{47}}%
\TextAt{13}{336.005}{498.425}{N}%
\TextAt{21}{345.378}{496.565}{L}%
\TextAt{28}{351.137}{495.174}{P}%
\TextAt{3}{357.335}{496.565}{\Character{106}}%
\TextAt{21}{359.687}{496.565}{K}%
\TextAt{28}{366.828}{495.569}{p}%
\TextAt{13}{371.409}{498.425}{L}%
\TextAt{3}{379.374}{503.361}{\Character{3}}%
\TextAt{27}{379.374}{495.47}{P}%
\TextAt{13}{386.467}{498.425}{\Character{59}}%
\TextAt{8}{79.226}{460.726}{where}%
\TextAt{14}{113.035}{460.726}{P}%
\TextAt{8}{126.826}{460.726}{is}%
\TextAt{8}{138.59}{460.726}{a}%
\TextAt{8}{148.345}{460.726}{chosen}%
\TextAt{8}{185.801}{460.726}{place}%
\TextAt{8}{215.713}{460.726}{above}%
\TextAt{14}{248.888}{460.726}{p}%
\TextAt{8}{254.871}{460.726}{.}%
\TextAt{8}{263.32}{460.726}{In}%
\TextAt{8}{277.941}{460.726}{other}%
\TextAt{8}{308.51}{460.726}{words:}%
\TextAt{19}{109.114}{430.729}{An}%
\TextAt{19}{128.047}{430.729}{idèle}%
\TextAt{13}{155.624}{430.729}{\Character{11}}%
\NativeMathAt{166.594}{430.729}{11.955}{\in}%
\TextAt{13}{178.012}{430.729}{I}%
\TextAt{21}{183.175}{428.935}{K}%
\TextAt{19}{195.642}{430.729}{is}%
\TextAt{19}{208.275}{430.729}{a}%
\TextAt{19}{218.517}{430.729}{norm}%
\TextAt{19}{249.836}{430.729}{of}%
\TextAt{19}{263.664}{430.729}{an}%
\TextAt{19}{280.481}{430.729}{idèle}%
\TextAt{19}{308.058}{430.729}{of}%
\TextAt{13}{321.887}{430.729}{L}%
\TextAt{19}{334.116}{430.729}{if}%
\TextAt{19}{345.553}{430.729}{and}%
\TextAt{19}{368.348}{430.729}{only}%
\TextAt{19}{393.833}{430.729}{if}%
\TextAt{19}{405.27}{430.729}{it}%
\TextAt{19}{417.007}{430.729}{is}%
\TextAt{19}{429.64}{430.729}{a}%
\TextAt{19}{439.881}{430.729}{norm}%
\TextAt{19}{471.2}{430.729}{locally}%
\TextAt{19}{109.114}{413.683}{everywhere,}%
\TextAt{19}{172.488}{413.683}{i.e.,}%
\TextAt{19}{196.399}{413.683}{if}%
\TextAt{19}{207.756}{413.683}{every}%
\TextAt{19}{238.696}{413.683}{component}%
\TextAt{13}{296.978}{413.683}{\Character{11}}%
\TextAt{27}{304.423}{411.889}{p}%
\TextAt{19}{313.094}{413.683}{is}%
\TextAt{19}{325.647}{413.683}{the}%
\TextAt{19}{345.074}{413.683}{norm}%
\TextAt{19}{376.313}{413.683}{of}%
\TextAt{19}{390.062}{413.683}{an}%
\TextAt{19}{406.799}{413.683}{element}%
\TextAt{19}{450.136}{413.683}{of}%
\TextAt{13}{463.885}{413.683}{L}%
\TextAt{3}{471.849}{418.021}{\Character{3}}%
\TextAt{27}{471.849}{410.598}{P}%
\TextAt{19}{478.943}{413.683}{.}%
\TextAt{8}{79.226}{383.685}{As}%
\TextAt{8}{96.513}{383.685}{for}%
\TextAt{8}{114.398}{383.685}{the}%
\TextAt{8}{134.555}{383.685}{Herbrand}%
\TextAt{8}{186.906}{383.685}{quotient}%
\TextAt{13}{233.089}{383.685}{h}%
\TextAt{9}{239.828}{383.685}{\Character{40}}%
\TextAt{13}{244.38}{383.685}{G\Character{59}}%
\TextAt{13}{258.834}{383.685}{I}%
\TextAt{21}{264.961}{388.024}{S}%
\TextAt{21}{264.021}{380.665}{L}%
\TextAt{9}{271.055}{383.685}{\Character{41}}%
\TextAt{8}{279.504}{383.685}{we}%
\TextAt{8}{296.732}{383.685}{obtain}%
\TextAt{8}{333.159}{383.685}{the}%
\TextAt{8}{353.315}{383.685}{result:}%
\TextAt{12}{79.226}{357.673}{Proposition}%
\TextAt{12}{152.304}{357.673}{(3.3).}%
\TextAt{8}{187.939}{357.673}{If}%
\TextAt{13}{199.95}{357.673}{L}%
\TextAt{16}{211.778}{357.673}{\Character{106}}%
\TextAt{13}{218.963}{357.673}{K}%
\TextAt{8}{233.997}{357.673}{is}%
\TextAt{8}{246.08}{357.673}{a}%
\TextAt{8}{256.154}{357.673}{cyclic}%
\TextAt{8}{288.656}{357.673}{extension}%
\TextAt{8}{340.741}{357.673}{and}%
\TextAt{8}{363.822}{357.673}{if}%
\TextAt{13}{374.865}{357.673}{S}%
\TextAt{8}{386.976}{357.673}{contains}%
\TextAt{8}{433.215}{357.673}{all}%
\TextAt{8}{449.793}{357.673}{ramified}%
\TextAt{8}{495.637}{357.673}{primes,}%
\TextAt{8}{79.226}{340.627}{then}%
\TextAt{13}{260.868}{323.581}{h}%
\TextAt{9}{267.607}{323.581}{\Character{40}}%
\TextAt{13}{272.159}{323.581}{G\Character{59}}%
\TextAt{13}{286.613}{323.581}{I}%
\TextAt{21}{292.74}{328.517}{S}%
\TextAt{21}{291.8}{320.625}{L}%
\TextAt{9}{298.834}{323.581}{\Character{41}}%
\TextAt{9}{306.689}{323.581}{\Character{61}}%
\TextAt{20}{319.132}{334.938}{\Character{89}}%
\TextAt{27}{319.155}{309.567}{p}%
\NativeMathAt{323.144}{309.567}{7.97}{\in}%
\TextAt{21}{328.79}{309.567}{S}%
\TextAt{13}{336.401}{323.581}{n}%
\TextAt{27}{343.389}{321.788}{p}%
\TextAt{13}{347.876}{323.581}{\Character{59}}%
\TextAt{8}{79.226}{290.426}{where}%
\TextAt{13}{113.035}{290.426}{n}%
\TextAt{27}{120.023}{288.633}{p}%
\TextAt{9}{127.831}{290.426}{\Character{61}}%
\TextAt{9}{140.241}{290.426}{\Character{91}}%
\TextAt{13}{143.508}{290.426}{L}%
\TextAt{27}{151.472}{288.633}{P}%
\TextAt{9}{161.887}{290.426}{\Character{58}}%
\TextAt{13}{168.459}{290.426}{K}%
\TextAt{27}{178.418}{288.633}{p}%
\TextAt{9}{182.905}{290.426}{\Character{93}}%
\TextAt{8}{186.157}{290.426}{.}%
\TextAt{19}{79.226}{255.447}{Proof.}%
\TextAt{8}{116}{255.447}{We}%
\TextAt{8}{136.145}{255.447}{have}%
\TextAt{13}{163.139}{255.447}{H}%
\TextAt{3}{173.762}{259.786}{\Character{0}}%
\TextAt{17}{180.349}{259.786}{1}%
\TextAt{9}{185.081}{255.447}{\Character{40}}%
\TextAt{13}{189.633}{255.447}{G\Character{59}}%
\TextAt{13}{204.087}{255.447}{I}%
\TextAt{21}{210.214}{259.786}{S}%
\TextAt{21}{209.274}{252.427}{L}%
\TextAt{9}{216.308}{255.447}{\Character{41}}%
\TextAt{9}{224.163}{255.447}{\Character{61}}%
\TextAt{20}{236.606}{264.414}{\Character{81}}%
\TextAt{27}{247.897}{251.96}{p}%
\NativeMathAt{251.886}{251.96}{7.97}{\in}%
\TextAt{21}{257.532}{251.96}{S}%
\TextAt{13}{265.618}{255.447}{H}%
\TextAt{3}{276.241}{259.786}{\Character{0}}%
\TextAt{17}{282.827}{259.786}{1}%
\TextAt{9}{287.559}{255.447}{\Character{40}}%
\TextAt{13}{292.112}{255.447}{G}%
\TextAt{27}{301.345}{253.654}{P}%
\TextAt{13}{308.439}{255.447}{\Character{59}}%
\TextAt{13}{313.663}{255.447}{L}%
\TextAt{3}{321.648}{259.786}{\Character{3}}%
\TextAt{27}{321.648}{252.362}{P}%
\TextAt{9}{328.741}{255.447}{\Character{41}}%
\TextAt{9}{336.596}{255.447}{\Character{61}}%
\TextAt{9}{349.017}{255.447}{1}%
\TextAt{8}{358.79}{255.447}{and}%
\TextAt{13}{231.753}{224.949}{H}%
\TextAt{17}{242.376}{229.885}{0}%
\TextAt{9}{247.108}{224.949}{\Character{40}}%
\TextAt{13}{251.66}{224.949}{G\Character{59}}%
\TextAt{13}{266.114}{224.949}{I}%
\TextAt{21}{272.241}{229.885}{S}%
\TextAt{21}{271.301}{221.994}{L}%
\TextAt{9}{278.335}{224.949}{\Character{41}}%
\TextAt{9}{286.19}{224.949}{\Character{61}}%
\TextAt{20}{298.633}{236.307}{\Character{89}}%
\TextAt{27}{298.656}{210.935}{p}%
\NativeMathAt{302.645}{210.935}{7.97}{\in}%
\TextAt{21}{308.291}{210.935}{S}%
\TextAt{13}{315.902}{224.949}{H}%
\TextAt{17}{326.525}{229.885}{0}%
\TextAt{9}{331.257}{224.949}{\Character{40}}%
\TextAt{13}{335.809}{224.949}{G}%
\TextAt{27}{345.043}{223.156}{P}%
\TextAt{13}{352.136}{224.949}{\Character{59}}%
\TextAt{13}{357.36}{224.949}{L}%
\TextAt{3}{365.345}{229.885}{\Character{3}}%
\TextAt{27}{365.345}{221.994}{P}%
\TextAt{9}{372.439}{224.949}{\Character{41}}%
\TextAt{13}{376.991}{224.949}{\Character{58}}%
\TextAt{8}{79.226}{185.806}{By}%
\TextAt{8}{97.589}{185.806}{local}%
\TextAt{8}{125.23}{185.806}{class}%
\TextAt{8}{152.667}{185.806}{field}%
\TextAt{8}{178.024}{185.806}{theory,}%
\TextAt{8}{217.045}{185.806}{we}%
\TextAt{8}{234.273}{185.806}{find}%
\TextAt{8}{257.681}{185.806}{\char35\relax }%
\TextAt{13}{267.436}{185.806}{H}%
\TextAt{17}{278.059}{190.145}{0}%
\TextAt{9}{282.791}{185.806}{\Character{40}}%
\TextAt{13}{287.344}{185.806}{G}%
\TextAt{27}{296.577}{184.013}{P}%
\TextAt{13}{303.671}{185.806}{\Character{59}}%
\TextAt{13}{308.895}{185.806}{L}%
\TextAt{3}{316.88}{190.145}{\Character{3}}%
\TextAt{27}{316.88}{182.722}{P}%
\TextAt{9}{323.973}{185.806}{\Character{41}}%
\TextAt{9}{331.828}{185.806}{\Character{61}}%
\TextAt{9}{344.249}{185.806}{\Character{40}}%
\TextAt{13}{348.824}{185.806}{K}%
\TextAt{3}{359.642}{190.145}{\Character{3}}%
\TextAt{27}{358.783}{182.851}{p}%
\TextAt{9}{367.695}{185.806}{\Character{58}}%
\TextAt{13}{374.267}{185.806}{N}%
\TextAt{21}{383.641}{183.947}{L}%
\TextAt{28}{389.4}{182.556}{P}%
\TextAt{3}{395.597}{183.947}{\Character{106}}%
\TextAt{21}{397.95}{183.947}{K}%
\TextAt{28}{405.091}{182.95}{p}%
\TextAt{13}{409.672}{185.806}{L}%
\TextAt{3}{417.637}{190.145}{\Character{3}}%
\TextAt{27}{417.637}{182.722}{P}%
\TextAt{9}{424.73}{185.806}{\Character{41}}%
\TextAt{9}{432.585}{185.806}{\Character{61}}%
\TextAt{13}{445.029}{185.806}{n}%
\TextAt{27}{452.016}{184.013}{p}%
\TextAt{8}{456.503}{185.806}{.}%
\TextAt{8}{464.952}{185.806}{Hence}%
\TextAt{13}{218.062}{150.391}{h}%
\TextAt{9}{224.8}{150.391}{\Character{40}}%
\TextAt{13}{229.353}{150.391}{G\Character{59}}%
\TextAt{13}{243.807}{150.391}{I}%
\TextAt{21}{249.933}{155.327}{S}%
\TextAt{21}{248.993}{147.435}{L}%
\TextAt{9}{256.027}{150.391}{\Character{41}}%
\TextAt{9}{263.882}{150.391}{\Character{61}}%
\TextAt{8}{280.815}{158.479}{\char35\relax }%
\TextAt{13}{290.57}{158.479}{H}%
\TextAt{17}{301.193}{162.817}{0}%
\TextAt{9}{305.925}{158.479}{\Character{40}}%
\TextAt{13}{310.477}{158.479}{G\Character{59}}%
\TextAt{13}{324.931}{158.479}{I}%
\TextAt{21}{331.058}{162.817}{S}%
\TextAt{21}{330.118}{155.458}{L}%
\TextAt{9}{337.152}{158.479}{\Character{41}}%
\Vector{q 0 G 0 g 0.478 w 205.521 -620.62 m 272.997 -620.62 l S Q}%
\TextAt{8}{277.521}{142.142}{\char35\relax }%
\TextAt{13}{287.277}{142.142}{H}%
\TextAt{3}{297.899}{145.596}{\Character{0}}%
\TextAt{17}{304.486}{145.596}{1}%
\TextAt{9}{309.218}{142.142}{\Character{40}}%
\TextAt{13}{313.771}{142.142}{G\Character{59}}%
\TextAt{13}{328.225}{142.142}{I}%
\TextAt{21}{334.351}{146.26}{S}%
\TextAt{21}{333.411}{138.901}{L}%
\TextAt{9}{340.445}{142.142}{\Character{41}}%
\TextAt{9}{349.514}{150.391}{\Character{61}}%
\TextAt{20}{361.939}{161.748}{\Character{89}}%
\TextAt{27}{361.962}{136.377}{p}%
\NativeMathAt{365.951}{136.377}{7.97}{\in}%
\TextAt{21}{371.596}{136.377}{S}%
\TextAt{13}{379.208}{150.391}{n}%
\TextAt{27}{386.195}{148.598}{p}%
\TextAt{13}{390.682}{150.391}{\Character{58}}%
\Vector{q 0 G 0 g 0.399 w 452.334 -661.347 m 452.334 -653.277 l S Q}%
\Vector{q 0 G 0 g 0.399 w 452.534 -653.476 m 459.707 -653.476 l S Q}%
\Vector{q 0 G 0 g 0.399 w 452.534 -661.148 m 459.707 -661.148 l S Q}%
\Vector{q 0 G 0 g 0.399 w 459.906 -661.347 m 459.906 -653.277 l S Q}%
\TextAt{8}{297.211}{43.033}{342}%
\Destination{78.23}{809.815}{nav.page.342}%
\Destination{79.36}{689.255}{nav.print.375}%
\Destination{79.36}{371.655}{nav.number.6.3.3}%
\Hotlink{439.3}{689.155}{13.995}{11.047}{nav.chapter.4}%
\Hotlink{378.315}{655.055}{9.775}{11.047}{nav.chapter.5}%
\Hotlink{356.495}{689.155}{25.078}{11.047}{nav.number.6.3.1}%
\Hotlink{459.836}{689.155}{25.078}{11.047}{nav.number.4.7.4}%
\Hotlink{394.922}{655.055}{25.078}{11.047}{nav.number.5.1.2}%
}

% Page 355 (body page 343)
\Page{%
\TextAt{2}{79.226}{796.943}{VI.}%
\TextAt{2}{98.012}{796.943}{Global}%
\TextAt{2}{133.248}{796.943}{Class}%
\TextAt{2}{161.841}{796.943}{Field}%
\TextAt{2}{189.572}{796.943}{Theory}%
\Vector{q 0 G 0 g 0.299 w 7.226 18.729 m 460.769 18.729 l S Q}%
\TextAt{8}{97.159}{761.124}{Next}%
\TextAt{8}{125.75}{761.124}{we}%
\TextAt{8}{142.96}{761.124}{determine}%
\TextAt{8}{197.566}{761.124}{the}%
\TextAt{8}{217.705}{761.124}{Herbrand}%
\TextAt{8}{270.039}{761.124}{quotient}%
\TextAt{8}{316.204}{761.124}{of}%
\TextAt{8}{329.517}{761.124}{the}%
\TextAt{13}{349.656}{761.124}{G}%
\TextAt{8}{358.889}{761.124}{-module}%
\TextAt{13}{404.062}{761.124}{L}%
\TextAt{21}{412.027}{765.462}{S}%
\TextAt{9}{421.442}{761.124}{\Character{61}}%
\TextAt{13}{433.867}{761.124}{L}%
\TextAt{16}{444.452}{761.124}{\Character{92}}%
\TextAt{13}{455.044}{761.124}{I}%
\TextAt{21}{461.146}{765.462}{S}%
\TextAt{21}{460.206}{758.103}{L}%
\TextAt{8}{467.24}{761.124}{.}%
\TextAt{8}{475.683}{761.124}{For}%
\TextAt{8}{496.635}{761.124}{this}%
\TextAt{8}{519.439}{761.124}{we}%
\TextAt{8}{79.226}{744.078}{need}%
\TextAt{8}{106.532}{744.078}{the}%
\TextAt{8}{126.688}{744.078}{following}%
\TextAt{8}{176.123}{744.078}{general}%
\TextAt{12}{79.226}{719.918}{Lemma}%
\TextAt{12}{126.678}{719.918}{(3.4).}%
\TextAt{8}{161.809}{719.918}{Let}%
\TextAt{13}{182.593}{719.918}{V}%
\TextAt{8}{195.73}{719.918}{be}%
\TextAt{8}{211.469}{719.918}{an}%
\TextAt{13}{227.543}{719.918}{s}%
\TextAt{8}{233.057}{719.918}{-dimensional}%
\TextAt{18}{301.22}{719.918}{R}%
\TextAt{8}{309.854}{719.918}{-vector}%
\TextAt{8}{348.69}{719.918}{space,}%
\TextAt{8}{383.069}{719.918}{and}%
\TextAt{8}{405.646}{719.918}{let}%
\TextAt{13}{422.365}{719.918}{G}%
\TextAt{8}{435.311}{719.918}{be}%
\TextAt{8}{451.05}{719.918}{a}%
\TextAt{8}{460.62}{719.918}{finite}%
\TextAt{8}{490.347}{719.918}{group}%
\TextAt{8}{523.337}{719.918}{of}%
\TextAt{8}{79.226}{702.872}{automorphisms}%
\TextAt{8}{161.9}{702.872}{of}%
\TextAt{13}{175.82}{702.872}{V}%
\TextAt{8}{189.733}{702.872}{which}%
\TextAt{8}{223.809}{702.872}{operates}%
\TextAt{8}{270.964}{702.872}{as}%
\TextAt{8}{285.924}{702.872}{a}%
\TextAt{8}{296.269}{702.872}{permutation}%
\TextAt{8}{363.855}{702.872}{group}%
\TextAt{8}{397.621}{702.872}{on}%
\TextAt{8}{414.469}{702.872}{the}%
\TextAt{8}{435.215}{702.872}{elements}%
\TextAt{8}{483.662}{702.872}{of}%
\TextAt{8}{497.581}{702.872}{a}%
\TextAt{8}{507.927}{702.872}{basis}%
\TextAt{13}{79.226}{685.826}{v}%
\TextAt{17}{84.885}{684.033}{1}%
\TextAt{13}{89.618}{685.826}{\Character{59}}%
\TextAt{13}{94.842}{685.826}{\Character{58}}%
\TextAt{13}{100.079}{685.826}{\Character{58}}%
\TextAt{13}{105.303}{685.826}{\Character{58}}%
\TextAt{13}{110.54}{685.826}{\Character{59}}%
\TextAt{13}{115.776}{685.826}{v}%
\TextAt{21}{121.498}{684.033}{s}%
\TextAt{8}{125.912}{685.826}{:}%
\TextAt{13}{134.922}{685.826}{\Character{27}}%
\TextAt{13}{141.988}{685.826}{v}%
\TextAt{21}{147.664}{684.033}{i}%
\TextAt{9}{154.844}{685.826}{\Character{61}}%
\TextAt{13}{167.748}{685.826}{v}%
\TextAt{21}{173.408}{683.966}{\Character{27}}%
\TextAt{17}{178.533}{683.966}{\Character{40}}%
\TextAt{21}{181.826}{683.966}{i}%
\TextAt{17}{184.709}{683.966}{\Character{41}}%
\TextAt{8}{188.501}{685.826}{.}%
\TextAt{8}{197.791}{685.826}{If}%
\TextAt{9}{209.764}{685.826}{\Character{0}}%
\TextAt{8}{221.257}{685.826}{is}%
\TextAt{8}{233.302}{685.826}{a}%
\TextAt{13}{243.338}{685.826}{G}%
\TextAt{8}{252.571}{685.826}{-invariant}%
\TextAt{8}{305.873}{685.826}{complete}%
\TextAt{8}{355.577}{685.826}{lattice}%
\TextAt{8}{391.628}{685.826}{in}%
\TextAt{13}{405.561}{685.826}{V}%
\TextAt{8}{414.987}{685.826}{,}%
\TextAt{8}{422.487}{685.826}{i.e.,}%
\TextAt{13}{444.943}{685.826}{\Character{27}}%
\TextAt{9}{452.026}{685.826}{\Character{0}}%
\NativeMathAt{463.139}{685.826}{11.955}{\subseteq}%
\TextAt{9}{476.237}{685.826}{\Character{0}}%
\TextAt{8}{487.73}{685.826}{for}%
\TextAt{8}{505.895}{685.826}{all}%
\TextAt{13}{522.435}{685.826}{\Character{27}}%
\TextAt{8}{529.518}{685.826}{,}%
\TextAt{8}{79.226}{668.78}{then}%
\TextAt{8}{105.886}{668.78}{there}%
\TextAt{8}{135.798}{668.78}{exists}%
\TextAt{8}{168.113}{668.78}{a}%
\TextAt{8}{177.868}{668.78}{complete}%
\TextAt{8}{227.291}{668.78}{sublattice}%
\TextAt{8}{280.683}{668.78}{in}%
\TextAt{9}{294.336}{668.78}{\Character{0}}%
\TextAt{8}{301.65}{668.78}{,}%
\TextAt{9}{249.141}{646.262}{\Character{0}}%
\TextAt{3}{256.455}{651.198}{\Character{48}}%
\TextAt{9}{262.571}{646.262}{\Character{61}}%
\TextAt{18}{274.997}{646.262}{Z}%
\TextAt{13}{282.967}{646.262}{w}%
\TextAt{17}{291.335}{644.469}{1}%
\TextAt{9}{298.724}{646.262}{\Character{43}}%
\TextAt{16}{310.485}{646.262}{\Character{1}}%
\TextAt{16}{315.781}{646.262}{\Character{1}}%
\TextAt{16}{321.089}{646.262}{\Character{1}}%
\TextAt{9}{327.089}{646.262}{\Character{43}}%
\TextAt{18}{338.851}{646.262}{Z}%
\TextAt{13}{346.821}{646.262}{w}%
\TextAt{21}{355.189}{644.469}{s}%
\TextAt{13}{359.603}{646.262}{\Character{59}}%
\TextAt{8}{79.226}{623.744}{such}%
\TextAt{8}{105.623}{623.744}{that}%
\TextAt{13}{130.992}{623.744}{\Character{27}}%
\TextAt{13}{138.058}{623.744}{w}%
\TextAt{21}{146.442}{621.951}{i}%
\TextAt{9}{153.144}{623.744}{\Character{61}}%
\TextAt{13}{165.57}{623.744}{w}%
\TextAt{21}{173.938}{621.884}{\Character{27}}%
\TextAt{17}{179.063}{621.884}{\Character{40}}%
\TextAt{21}{182.356}{621.884}{i}%
\TextAt{17}{185.239}{621.884}{\Character{41}}%
\TextAt{8}{192.928}{623.744}{for}%
\TextAt{8}{210.813}{623.744}{all}%
\TextAt{13}{227.072}{623.744}{\Character{27}}%
\NativeMathAt{237.475}{623.744}{11.955}{\in}%
\TextAt{13}{248.766}{623.744}{G}%
\TextAt{8}{258}{623.744}{.}%
\TextAt{19}{79.226}{592.47}{Proof.}%
\TextAt{8}{116}{592.47}{Let}%
\TextAt{16}{138.232}{592.47}{\Character{106}}%
\TextAt{16}{146.696}{592.47}{\Character{106}}%
\TextAt{8}{155.194}{592.47}{be}%
\TextAt{8}{172.38}{592.47}{the}%
\TextAt{8}{193.799}{592.47}{sup-norm}%
\TextAt{8}{247.15}{592.47}{with}%
\TextAt{8}{275.073}{592.47}{respect}%
\TextAt{8}{316.385}{592.47}{to}%
\TextAt{8}{331.958}{592.47}{the}%
\TextAt{8}{353.377}{592.47}{coordinates}%
\TextAt{8}{416.818}{592.47}{of}%
\TextAt{8}{431.41}{592.47}{the}%
\TextAt{8}{452.829}{592.47}{basis}%
\TextAt{13}{482.832}{592.47}{v}%
\TextAt{17}{488.491}{590.677}{1}%
\TextAt{13}{493.224}{592.47}{\Character{59}}%
\TextAt{13}{498.448}{592.47}{\Character{58}}%
\TextAt{13}{503.685}{592.47}{\Character{58}}%
\TextAt{13}{508.909}{592.47}{\Character{58}}%
\TextAt{13}{514.146}{592.47}{\Character{59}}%
\TextAt{13}{519.382}{592.47}{v}%
\TextAt{21}{525.104}{590.677}{s}%
\TextAt{8}{529.518}{592.47}{.}%
\TextAt{8}{79.226}{575.424}{Since}%
\TextAt{9}{110.152}{575.424}{\Character{0}}%
\TextAt{8}{121.731}{575.424}{is}%
\TextAt{8}{133.863}{575.424}{a}%
\TextAt{8}{143.987}{575.424}{lattice,}%
\TextAt{8}{183.469}{575.424}{there}%
\TextAt{8}{213.749}{575.424}{exists}%
\TextAt{8}{246.432}{575.424}{a}%
\TextAt{8}{256.555}{575.424}{number}%
\TextAt{13}{299.52}{575.424}{b}%
\TextAt{8}{308.762}{575.424}{such}%
\TextAt{8}{335.527}{575.424}{that}%
\TextAt{8}{361.264}{575.424}{for}%
\TextAt{8}{379.517}{575.424}{every}%
\TextAt{13}{410.778}{575.424}{x}%
\NativeMathAt{421.378}{575.424}{11.955}{\in}%
\TextAt{13}{433.296}{575.424}{V}%
\TextAt{8}{442.722}{575.424}{,}%
\TextAt{8}{450.331}{575.424}{there}%
\TextAt{8}{480.611}{575.424}{is}%
\TextAt{8}{492.743}{575.424}{a}%
\TextAt{13}{502.867}{575.424}{\Character{13}}%
\NativeMathAt{513.537}{575.424}{11.955}{\in}%
\TextAt{9}{525.455}{575.424}{\Character{0}}%
\ErrataAt{79.19}{558.378}{455.81}{satisfying $|x-\gamma|<b$. Choose $t>sb$ and $\gamma_i\in\Gamma$ with $|tv_i-\gamma_i|<b$ for $i=1,\ldots,s$.}%
\ErrataAt{79.19}{541.332}{455.81}{For each basis index define}%
\ErrataAt{216}{524.286}{236}{$\displaystyle w_i=\sum_{\sigma\in G}\sigma\gamma_{\sigma^{-1}(i)},\qquad i=1,\ldots,s.$}%
\ErrataAt{79.19}{492.901}{455.81}{For every $\tau\in G$, change the summation variable from $\sigma$ to $\rho=\tau\sigma$. Then}%
\ErrataAt{192}{470.398}{263}{$\displaystyle\tau w_i=\sum_{\sigma\in G}\tau\sigma\gamma_{\sigma^{-1}(i)}=\sum_{\rho\in G}\rho\gamma_{\rho^{-1}(\tau(i))}=w_{\tau(i)}.$}%
\TextAt{8}{79.226}{437.757}{It}%
\TextAt{8}{94.357}{437.757}{is}%
\TextAt{8}{108.58}{437.757}{therefore}%
\TextAt{8}{160.138}{437.757}{enough}%
\TextAt{8}{202.922}{437.757}{to}%
\TextAt{8}{219.691}{437.757}{check}%
\TextAt{8}{253.687}{437.757}{the}%
\TextAt{8}{276.302}{437.757}{linear}%
\TextAt{8}{311.279}{437.757}{independence}%
\TextAt{8}{386.234}{437.757}{of}%
\TextAt{8}{402.022}{437.757}{the}%
\TextAt{13}{424.638}{437.757}{w}%
\TextAt{21}{433.005}{435.963}{i}%
\TextAt{8}{436.387}{437.757}{.}%
\TextAt{8}{452.211}{437.757}{To}%
\TextAt{8}{471.897}{437.757}{do}%
\TextAt{8}{490.615}{437.757}{this,}%
\TextAt{8}{519.762}{437.757}{let}%
\TextAt{20}{79.226}{429.677}{\Character{80}}%
\TextAt{21}{91.846}{426.522}{s}%
\TextAt{21}{91.846}{417.224}{i}%
\TextAt{17}{94.729}{417.224}{\Character{61}1}%
\TextAt{13}{108.04}{420.711}{c}%
\TextAt{21}{113.078}{418.917}{i}%
\TextAt{13}{116.459}{420.711}{w}%
\TextAt{21}{124.827}{418.917}{i}%
\TextAt{9}{134.276}{420.711}{\Character{61}}%
\TextAt{9}{149.435}{420.711}{0}%
\TextAt{8}{155.302}{420.711}{,}%
\TextAt{13}{164.466}{420.711}{c}%
\TextAt{21}{169.504}{418.917}{i}%
\NativeMathAt{178.953}{420.711}{11.955}{\in}%
\TextAt{18}{192.991}{420.711}{R}%
\TextAt{8}{201.625}{420.711}{.}%
\TextAt{8}{214.91}{420.711}{If}%
\TextAt{8}{228.214}{420.711}{not}%
\TextAt{8}{250.64}{420.711}{all}%
\TextAt{8}{268.511}{420.711}{of}%
\TextAt{8}{283.453}{420.711}{the}%
\TextAt{13}{305.221}{420.711}{c}%
\TextAt{21}{310.259}{418.917}{i}%
\TextAt{9}{319.708}{420.711}{\Character{61}}%
\TextAt{9}{334.867}{420.711}{0}%
\TextAt{8}{340.734}{420.711}{,}%
\TextAt{8}{349.898}{420.711}{then}%
\TextAt{8}{378.17}{420.711}{we}%
\TextAt{8}{397.009}{420.711}{may}%
\TextAt{8}{423.99}{420.711}{assume}%
\TextAt{16}{466.047}{420.711}{\Character{106}}%
\TextAt{13}{469.368}{420.711}{c}%
\TextAt{21}{474.406}{418.917}{i}%
\TextAt{16}{477.787}{420.711}{\Character{106}}%
\NativeMathAt{487.16}{420.711}{11.955}{\leq}%
\TextAt{9}{502.542}{420.711}{1}%
\TextAt{8}{513.904}{420.711}{and}%
\ErrataAt{79.19}{403.736}{455.81}{$c_j=1$ for some $j$. Write $\gamma_i=tv_i-y_i$ with $|y_i|<b$, and put $g=\#G$. Then}%
\ErrataAt{184}{381.076}{276}{$\displaystyle w_i=\sum_{\sigma\in G}\sigma\gamma_{\sigma^{-1}(i)}=gtv_i-z_i,\qquad z_i=\sum_{\sigma\in G}\sigma y_{\sigma^{-1}(i)}.$}%
\ErrataAt{79.19}{348.026}{455.81}{Since $G$ permutes the basis, $|z_i|\le gb$. The assumed dependence gives}%
\ErrataAt{202}{318.896}{268}{$\displaystyle 0=\sum_{i=1}^s c_iw_i=gt\sum_{i=1}^s c_i v_i-z,\qquad z=\sum_{i=1}^s c_i z_i.$}%
\ErrataAt{79.19}{288.492}{455.81}{Here $|z|\le sgb$ since $|c_i|\le1$. The $j$-th coordinate of $z$ is $gt$, so}%
\ErrataAt{232}{271.446}{152}{$gt\le sgb.$}%
\ErrataAt{79.19}{241.48}{455.81}{This contradicts $t>sb$. Thus all the $c_i$ vanish, and the linear independence of the $w_i$}%
\ErrataAt{79.19}{224.434}{455.81}{proves the lemma.}%
\Vector{q 0 G 0 g 0.399 w 452.334 -549.566 m 452.334 -541.496 l S Q}%
\Vector{q 0 G 0 g 0.399 w 452.534 -541.696 m 459.707 -541.696 l S Q}%
\Vector{q 0 G 0 g 0.399 w 452.534 -549.367 m 459.707 -549.367 l S Q}%
\Vector{q 0 G 0 g 0.399 w 459.906 -549.566 m 459.906 -541.496 l S Q}%
\TextAt{8}{116.179}{197.285}{Now}%
\TextAt{8}{142.767}{197.285}{let}%
\TextAt{13}{159.599}{197.285}{L}%
\TextAt{16}{170.884}{197.285}{\Character{106}}%
\TextAt{13}{177.526}{197.285}{K}%
\TextAt{8}{192.169}{197.285}{be}%
\TextAt{8}{208.021}{197.285}{a}%
\TextAt{8}{217.704}{197.285}{cyclic}%
\TextAt{8}{249.815}{197.285}{extension}%
\TextAt{8}{301.509}{197.285}{of}%
\TextAt{8}{314.767}{197.285}{degree}%
\TextAt{13}{351.11}{197.285}{n}%
\TextAt{8}{361.923}{197.285}{with}%
\TextAt{8}{388.51}{197.285}{Galois}%
\TextAt{8}{424.351}{197.285}{group}%
\TextAt{13}{457.455}{197.285}{G}%
\TextAt{9}{470.009}{197.285}{\Character{61}}%
\TextAt{13}{482.435}{197.285}{G}%
\TextAt{9}{491.668}{197.285}{\Character{40}}%
\TextAt{13}{496.221}{197.285}{L}%
\TextAt{16}{507.506}{197.285}{\Character{106}}%
\TextAt{13}{514.148}{197.285}{K}%
\TextAt{9}{524.965}{197.285}{\Character{41}}%
\TextAt{8}{529.518}{197.285}{,}%
\TextAt{8}{79.226}{180.239}{let}%
\TextAt{13}{96.069}{180.239}{S}%
\TextAt{8}{107.8}{180.239}{be}%
\TextAt{8}{123.663}{180.239}{a}%
\TextAt{8}{133.357}{180.239}{finite}%
\TextAt{8}{163.207}{180.239}{set}%
\TextAt{8}{181.413}{180.239}{of}%
\TextAt{8}{194.682}{180.239}{places}%
\TextAt{8}{229.147}{180.239}{containing}%
\TextAt{8}{286.003}{180.239}{the}%
\TextAt{8}{306.098}{180.239}{infinite}%
\TextAt{8}{345.704}{180.239}{places,}%
\TextAt{8}{383.433}{180.239}{and}%
\TextAt{8}{406.134}{180.239}{let}%
\Vector{q 0 G 0 g 0.478 w 350.977 -583.918 m 358.873 -583.918 l S Q}%
\TextAt{13}{422.977}{180.239}{S}%
\TextAt{8}{434.708}{180.239}{be}%
\TextAt{8}{450.571}{180.239}{the}%
\TextAt{8}{470.666}{180.239}{set}%
\TextAt{8}{488.872}{180.239}{of}%
\TextAt{8}{502.14}{180.239}{places}%
\TextAt{8}{79.226}{163.193}{of}%
\TextAt{13}{92.556}{163.193}{L}%
\TextAt{8}{104.418}{163.193}{that}%
\TextAt{8}{129.787}{163.193}{lie}%
\TextAt{8}{145.388}{163.193}{above}%
\TextAt{8}{178.564}{163.193}{the}%
\TextAt{8}{198.72}{163.193}{places}%
\TextAt{8}{233.247}{163.193}{of}%
\TextAt{13}{246.577}{163.193}{S}%
\TextAt{8}{254.472}{163.193}{.}%
\TextAt{8}{262.921}{163.193}{We}%
\TextAt{8}{283.065}{163.193}{denote}%
\TextAt{8}{320.784}{163.193}{the}%
\TextAt{8}{340.94}{163.193}{group}%
\TextAt{13}{374.116}{163.193}{L}%
\Vector{q 0 G 0 g 0.359 w 310.08 -599.767 m 316.174 -599.767 l S Q}%
\TextAt{21}{382.08}{167.531}{S}%
\TextAt{8}{392.071}{163.193}{of}%
\Vector{q 0 G 0 g 0.478 w 333.401 -600.964 m 341.297 -600.964 l S Q}%
\TextAt{13}{405.401}{163.193}{S}%
\TextAt{8}{413.297}{163.193}{-units}%
\TextAt{8}{446.52}{163.193}{simply}%
\TextAt{8}{483.976}{163.193}{by}%
\TextAt{13}{500.235}{163.193}{L}%
\TextAt{21}{508.199}{167.531}{S}%
\TextAt{8}{514.293}{163.193}{.}%
\TextAt{8}{297.211}{43.033}{343}%
\Destination{78.23}{809.815}{nav.page.343}%
\Destination{79.36}{734.955}{nav.print.376}%
\Destination{79.36}{776.155}{nav.print.377}%
\Destination{79.36}{733.955}{nav.number.6.3.4}%
\Destination{79.226}{575.378}{errata.134}%
}

% Page 356 (body page 344)
\Page{%
\TextAt{2}{79.226}{796.943}{VI.}%
\TextAt{2}{98.012}{796.943}{Global}%
\TextAt{2}{133.248}{796.943}{Class}%
\TextAt{2}{161.841}{796.943}{Field}%
\TextAt{2}{189.572}{796.943}{Theory}%
\Vector{q 0 G 0 g 0.299 w 7.226 18.729 m 460.769 18.729 l S Q}%
\TextAt{12}{79.226}{761.124}{Proposition}%
\TextAt{12}{151.937}{761.124}{(3.5).}%
\TextAt{8}{187.253}{761.124}{The}%
\TextAt{8}{211.307}{761.124}{Herbrand}%
\TextAt{8}{263.659}{761.124}{quotient}%
\TextAt{8}{309.841}{761.124}{of}%
\TextAt{8}{323.171}{761.124}{the}%
\TextAt{13}{343.328}{761.124}{G}%
\TextAt{8}{352.561}{761.124}{-module}%
\TextAt{13}{397.752}{761.124}{L}%
\TextAt{21}{405.716}{765.462}{S}%
\TextAt{8}{415.708}{761.124}{satisfies}%
\TextAt{13}{254.252}{732.26}{h}%
\TextAt{9}{260.991}{732.26}{\Character{40}}%
\TextAt{13}{265.543}{732.26}{G\Character{59}}%
\TextAt{13}{279.997}{732.26}{L}%
\TextAt{21}{287.985}{737.196}{S}%
\TextAt{9}{294.079}{732.26}{\Character{41}}%
\TextAt{9}{301.934}{732.26}{\Character{61}}%
\TextAt{9}{316.141}{740.347}{1}%
\Vector{q 0 G 0 g 0.478 w 243.573 -38.752 m 250.561 -38.752 l S Q}%
\TextAt{13}{315.573}{724.059}{n}%
\TextAt{20}{325.749}{743.617}{\Character{89}}%
\TextAt{27}{325.772}{718.245}{p}%
\NativeMathAt{329.761}{718.245}{7.97}{\in}%
\TextAt{21}{335.406}{718.245}{S}%
\TextAt{13}{343.017}{732.26}{n}%
\TextAt{27}{350.005}{730.466}{p}%
\TextAt{13}{354.492}{732.26}{\Character{59}}%
\TextAt{8}{79.226}{696.506}{where}%
\TextAt{13}{113.035}{696.506}{n}%
\TextAt{27}{120.023}{694.712}{p}%
\TextAt{9}{127.831}{696.506}{\Character{61}}%
\TextAt{9}{140.241}{696.506}{\Character{91}}%
\TextAt{13}{143.508}{696.506}{L}%
\TextAt{27}{151.472}{694.712}{P}%
\TextAt{9}{161.887}{696.506}{\Character{58}}%
\TextAt{13}{168.459}{696.506}{K}%
\TextAt{27}{178.418}{694.712}{p}%
\TextAt{9}{182.905}{696.506}{\Character{93}}%
\TextAt{8}{186.157}{696.506}{.}%
\TextAt{19}{79.226}{663.173}{Proof.}%
\TextAt{8}{116}{663.173}{Let}%
\TextAt{16}{137.967}{663.173}{\Character{102}}%
\TextAt{13}{143.945}{663.173}{e}%
\TextAt{27}{149.37}{661.379}{P}%
\TextAt{16}{161.485}{663.173}{\Character{106}}%
\TextAt{14}{169.827}{663.173}{P}%
\NativeMathAt{184.741}{663.173}{11.955}{\in}%
\Vector{q 0 G 0 g 0.478 w 125.732 -100.984 m 133.628 -100.984 l S Q}%
\TextAt{13}{197.732}{663.173}{S}%
\TextAt{16}{205.628}{663.173}{\Character{103}}%
\TextAt{8}{216.5}{663.173}{be}%
\TextAt{8}{233.422}{663.173}{the}%
\TextAt{8}{254.577}{663.173}{standard}%
\TextAt{8}{304.423}{663.173}{basis}%
\TextAt{8}{334.161}{663.173}{of}%
\TextAt{8}{348.489}{663.173}{the}%
\TextAt{8}{369.643}{663.173}{vector}%
\TextAt{8}{405.765}{663.173}{space}%
\TextAt{13}{438.037}{663.173}{V}%
\TextAt{9}{452.484}{663.173}{\Character{61}}%
\TextAt{20}{466.61}{672.139}{\Character{81}}%
\TextAt{27}{477.901}{659.686}{P}%
\NativeMathAt{484.496}{659.686}{7.97}{\in}%
\Vector{q 0 G 0 g 0.359 w 418.142 -107.613 m 423.737 -107.613 l S Q}%
\TextAt{21}{490.142}{659.686}{S}%
\TextAt{18}{498.228}{663.173}{R}%
\TextAt{8}{506.862}{663.173}{.}%
\TextAt{8}{518.304}{663.173}{By}%
\TextAt{8}{79.226}{646.127}{(1.1),}%
\TextAt{8}{110.453}{646.127}{the}%
\TextAt{8}{130.609}{646.127}{homomorphism}%
\TextAt{13}{204.209}{620.006}{\Character{21}}%
\TextAt{9}{214.358}{620.006}{\Character{58}}%
\TextAt{13}{220.931}{620.006}{L}%
\TextAt{21}{228.895}{624.942}{S}%
\NativeArrowAt{238.31}{620.006}{11.955}{19.26}{\longrightarrow}%
\TextAt{13}{260.892}{620.006}{V}%
\TextAt{13}{268.352}{620.006}{\Character{59}}%
\TextAt{13}{297.499}{620.006}{\Character{21}}%
\TextAt{9}{304.35}{620.006}{\Character{40}}%
\TextAt{13}{308.902}{620.006}{a}%
\TextAt{9}{315.047}{620.006}{\Character{41}}%
\TextAt{9}{322.902}{620.006}{\Character{61}}%
\TextAt{20}{335.63}{631.363}{\Character{88}}%
\TextAt{27}{335.346}{604.198}{P}%
\NativeMathAt{341.941}{604.198}{7.97}{\in}%
\Vector{q 0 G 0 g 0.359 w 275.587 -163.1 m 281.183 -163.1 l S Q}%
\TextAt{21}{347.587}{604.198}{S}%
\TextAt{8}{355.175}{620.006}{log}%
\TextAt{16}{372.135}{620.006}{\Character{106}}%
\TextAt{13}{375.456}{620.006}{a}%
\TextAt{16}{381.601}{620.006}{\Character{106}}%
\TextAt{27}{384.922}{618.213}{P}%
\TextAt{13}{392.016}{620.006}{e}%
\TextAt{27}{397.441}{618.213}{P}%
\TextAt{13}{404.535}{620.006}{\Character{59}}%
\TextAt{8}{79.226}{580.865}{has}%
\TextAt{8}{101.074}{580.865}{kernel}%
\TextAt{13}{136.515}{580.865}{\Character{22}}%
\TextAt{9}{143.558}{580.865}{\Character{40}}%
\TextAt{13}{148.111}{580.865}{L}%
\TextAt{9}{156.075}{580.865}{\Character{41}}%
\TextAt{8}{165.499}{580.865}{and}%
\TextAt{8}{189.236}{580.865}{its}%
\TextAt{8}{206.529}{580.865}{image}%
\TextAt{8}{241.325}{580.865}{is}%
\TextAt{8}{254.063}{580.865}{an}%
\TextAt{9}{271.297}{580.865}{\Character{40}}%
\TextAt{13}{275.849}{580.865}{s}%
\TextAt{16}{284.684}{580.865}{\Character{0}}%
\TextAt{9}{297.303}{580.865}{1\Character{41}}%
\TextAt{8}{307.708}{580.865}{-dimensional}%
\TextAt{8}{377.03}{580.865}{lattice,}%
\TextAt{13}{417.27}{580.865}{s}%
\TextAt{9}{427.765}{580.865}{\Character{61}}%
\TextAt{8}{441.851}{580.865}{\char35\relax }%
\Vector{q 0 G 0 g 0.478 w 379.607 -183.292 m 387.502 -183.292 l S Q}%
\TextAt{13}{451.607}{580.865}{S}%
\TextAt{8}{459.502}{580.865}{.}%
\TextAt{8}{470.873}{580.865}{We}%
\TextAt{8}{491.992}{580.865}{make}%
\TextAt{13}{523.536}{580.865}{G}%
\TextAt{8}{79.226}{563.819}{operate}%
\TextAt{8}{121.172}{563.819}{on}%
\TextAt{13}{137.426}{563.819}{V}%
\TextAt{8}{150.745}{563.819}{via}%
\TextAt{13}{169.928}{563.819}{\Character{27}}%
\TextAt{13}{176.994}{563.819}{e}%
\TextAt{27}{182.436}{562.025}{P}%
\TextAt{9}{192.85}{563.819}{\Character{61}}%
\TextAt{13}{205.276}{563.819}{e}%
\TextAt{21}{210.701}{562.025}{\Character{27}}%
\TextAt{27}{215.827}{562.025}{P}%
\TextAt{8}{222.92}{563.819}{.}%
\TextAt{8}{231.367}{563.819}{Then}%
\TextAt{13}{261.92}{563.819}{\Character{21}}%
\TextAt{8}{272.641}{563.819}{is}%
\TextAt{8}{284.4}{563.819}{a}%
\TextAt{13}{294.151}{563.819}{G}%
\TextAt{8}{303.384}{563.819}{-homomorphism}%
\TextAt{8}{389.947}{563.819}{because}%
\TextAt{8}{433.244}{563.819}{we}%
\TextAt{8}{450.467}{563.819}{have,}%
\TextAt{8}{480.71}{563.819}{for}%
\TextAt{13}{498.59}{563.819}{\Character{27}}%
\NativeMathAt{508.993}{563.819}{11.955}{\in}%
\TextAt{13}{520.284}{563.819}{G}%
\TextAt{8}{529.518}{563.819}{,}%
\TextAt{13}{183.939}{537.698}{\Character{21}}%
\TextAt{9}{190.768}{537.698}{\Character{40}}%
\TextAt{13}{195.32}{537.698}{\Character{27}}%
\TextAt{13}{202.386}{537.698}{a}%
\TextAt{9}{208.547}{537.698}{\Character{41}}%
\TextAt{9}{216.402}{537.698}{\Character{61}}%
\TextAt{20}{228.846}{549.055}{\Character{88}}%
\TextAt{27}{234.183}{523.619}{P}%
\TextAt{8}{248.107}{537.698}{log}%
\TextAt{16}{265.068}{537.698}{\Character{106}}%
\TextAt{13}{268.388}{537.698}{\Character{27}}%
\TextAt{13}{275.454}{537.698}{a}%
\TextAt{16}{281.616}{537.698}{\Character{106}}%
\TextAt{27}{284.937}{535.905}{P}%
\TextAt{13}{292.03}{537.698}{e}%
\TextAt{27}{297.456}{535.905}{P}%
\TextAt{9}{307.87}{537.698}{\Character{61}}%
\TextAt{20}{320.296}{549.055}{\Character{88}}%
\TextAt{27}{325.632}{523.619}{P}%
\TextAt{8}{339.557}{537.698}{log}%
\TextAt{16}{356.517}{537.698}{\Character{106}}%
\TextAt{13}{359.838}{537.698}{a}%
\TextAt{16}{365.983}{537.698}{\Character{106}}%
\TextAt{21}{369.304}{535.694}{\Character{27}}%
\TextAt{23}{374.429}{537.964}{\Character{0}}%
\TextAt{24}{380.185}{537.964}{1}%
\TextAt{27}{384.336}{535.694}{P}%
\TextAt{13}{393.422}{537.698}{\Character{27}}%
\TextAt{13}{400.488}{537.698}{e}%
\TextAt{21}{405.93}{535.694}{\Character{27}}%
\TextAt{23}{411.056}{537.964}{\Character{0}}%
\TextAt{24}{416.812}{537.964}{1}%
\TextAt{27}{420.963}{535.694}{P}%
\TextAt{9}{216.421}{496.01}{\Character{61}}%
\TextAt{13}{228.846}{496.01}{\Character{27}}%
\TextAt{20}{237.921}{516.454}{\Character{32}}%
\TextAt{20}{247.386}{507.368}{\Character{88}}%
\TextAt{27}{252.722}{481.932}{P}%
\TextAt{8}{266.647}{496.01}{log}%
\TextAt{16}{283.607}{496.01}{\Character{106}}%
\TextAt{13}{286.928}{496.01}{a}%
\TextAt{16}{293.073}{496.01}{\Character{106}}%
\TextAt{21}{296.394}{494.007}{\Character{27}}%
\TextAt{23}{301.519}{496.276}{\Character{0}}%
\TextAt{24}{307.275}{496.276}{1}%
\TextAt{27}{311.426}{494.007}{P}%
\TextAt{13}{318.52}{496.01}{e}%
\TextAt{21}{323.945}{494.007}{\Character{27}}%
\TextAt{23}{329.071}{496.276}{\Character{0}}%
\TextAt{24}{334.827}{496.276}{1}%
\TextAt{27}{338.978}{494.007}{P}%
\TextAt{20}{346.071}{516.454}{\Character{33}}%
\TextAt{9}{358.857}{496.01}{\Character{61}}%
\TextAt{13}{371.282}{496.01}{\Character{27}}%
\TextAt{13}{378.348}{496.01}{\Character{21}}%
\TextAt{9}{385.193}{496.01}{\Character{40}}%
\TextAt{13}{389.745}{496.01}{a}%
\TextAt{9}{395.89}{496.01}{\Character{41}}%
\TextAt{13}{400.443}{496.01}{\Character{58}}%
\TextAt{8}{79.226}{459.374}{Therefore}%
\TextAt{13}{131.8}{459.374}{e}%
\TextAt{17}{137.225}{457.58}{0}%
\TextAt{9}{145.278}{459.374}{\Character{61}}%
\TextAt{20}{157.704}{468.34}{\Character{80}}%
\TextAt{27}{170.323}{455.887}{P}%
\NativeMathAt{176.919}{455.887}{7.97}{\in}%
\Vector{q 0 G 0 g 0.359 w 110.564 -311.412 m 116.16 -311.412 l S Q}%
\TextAt{21}{182.564}{455.887}{S}%
\TextAt{13}{190.651}{459.374}{e}%
\TextAt{27}{196.076}{457.58}{P}%
\TextAt{8}{206.643}{459.374}{and}%
\TextAt{13}{228.982}{459.374}{\Character{21}}%
\TextAt{9}{235.811}{459.374}{\Character{40}}%
\TextAt{13}{240.363}{459.374}{L}%
\TextAt{21}{248.327}{463.712}{S}%
\TextAt{9}{254.421}{459.374}{\Character{41}}%
\TextAt{8}{262.447}{459.374}{generate}%
\TextAt{8}{308.852}{459.374}{a}%
\TextAt{13}{318.184}{459.374}{G}%
\TextAt{8}{327.417}{459.374}{-invariant}%
\TextAt{8}{380.015}{459.374}{complete}%
\TextAt{8}{429.014}{459.374}{lattice}%
\TextAt{9}{464.36}{459.374}{\Character{0}}%
\TextAt{8}{475.148}{459.374}{in}%
\TextAt{13}{488.377}{459.374}{V}%
\TextAt{8}{497.802}{459.374}{.}%
\TextAt{8}{506.109}{459.374}{Since}%
\TextAt{18}{79.226}{442.328}{Z}%
\TextAt{13}{87.196}{442.328}{e}%
\TextAt{17}{92.622}{440.534}{0}%
\TextAt{8}{101.251}{442.328}{is}%
\TextAt{13}{113.015}{442.328}{G}%
\TextAt{8}{122.249}{442.328}{-isomorphic}%
\TextAt{8}{185.396}{442.328}{to}%
\TextAt{18}{199.706}{442.328}{Z}%
\TextAt{8}{207.677}{442.328}{,}%
\TextAt{8}{214.826}{442.328}{the}%
\TextAt{8}{234.982}{442.328}{exact}%
\TextAt{8}{265.874}{442.328}{sequence}%
\TextAt{9}{218.345}{416.207}{0}%
\NativeArrowAt{227.519}{416.207}{11.955}{19.26}{\longrightarrow}%
\TextAt{18}{250.101}{416.207}{Z}%
\TextAt{13}{258.071}{416.207}{e}%
\TextAt{17}{263.496}{414.413}{0}%
\NativeArrowAt{271.549}{416.207}{11.955}{19.26}{\longrightarrow}%
\TextAt{9}{294.131}{416.207}{\Character{0}}%
\NativeArrowAt{304.766}{416.207}{11.955}{19.26}{\longrightarrow}%
\TextAt{9}{327.348}{416.207}{\Character{0}\Character{47}}%
\TextAt{18}{340.515}{416.207}{Z}%
\TextAt{13}{348.486}{416.207}{e}%
\TextAt{17}{353.911}{414.413}{0}%
\NativeArrowAt{361.964}{416.207}{11.955}{19.26}{\longrightarrow}%
\TextAt{9}{384.546}{416.207}{0}%
\TextAt{13}{390.399}{416.207}{\Character{59}}%
\TextAt{8}{79.226}{390.086}{together}%
\TextAt{8}{125.409}{390.086}{with}%
\TextAt{8}{152.069}{390.086}{the}%
\TextAt{8}{172.225}{390.086}{fact}%
\TextAt{8}{195.311}{390.086}{that}%
\TextAt{9}{220.68}{390.086}{\Character{0}\Character{47}}%
\TextAt{18}{233.847}{390.086}{Z}%
\TextAt{13}{241.817}{390.086}{e}%
\TextAt{17}{247.242}{388.293}{0}%
\NativeMathAt{255.296}{390.086}{11.955}{\cong}%
\TextAt{13}{267.915}{390.086}{\Character{21}}%
\TextAt{9}{274.743}{390.086}{\Character{40}}%
\TextAt{13}{279.296}{390.086}{L}%
\TextAt{21}{287.26}{394.424}{S}%
\TextAt{9}{293.354}{390.086}{\Character{41}}%
\TextAt{8}{297.906}{390.086}{,}%
\TextAt{8}{305.056}{390.086}{yields}%
\TextAt{8}{337.956}{390.086}{the}%
\TextAt{8}{358.113}{390.086}{identities}%
\TextAt{13}{161.277}{361.234}{h}%
\TextAt{9}{168.016}{361.234}{\Character{40}}%
\TextAt{13}{172.568}{361.234}{G\Character{59}}%
\TextAt{13}{187.022}{361.234}{L}%
\TextAt{21}{195.01}{366.17}{S}%
\TextAt{9}{201.104}{361.234}{\Character{41}}%
\TextAt{9}{208.959}{361.234}{\Character{61}}%
\TextAt{13}{221.403}{361.234}{h}%
\TextAt{9}{228.141}{361.234}{\Character{40}}%
\TextAt{13}{232.694}{361.234}{G\Character{59}}%
\TextAt{13}{247.148}{361.234}{\Character{21}}%
\TextAt{9}{254}{361.234}{\Character{40}}%
\TextAt{13}{258.552}{361.234}{L}%
\TextAt{21}{266.517}{366.17}{S}%
\TextAt{9}{272.611}{361.234}{\Character{41}\Character{41}}%
\TextAt{9}{285.009}{361.234}{\Character{61}}%
\TextAt{13}{297.462}{361.234}{h}%
\TextAt{9}{304.2}{361.234}{\Character{40}}%
\TextAt{13}{308.753}{361.234}{G\Character{59}}%
\TextAt{18}{323.23}{361.234}{Z}%
\TextAt{9}{331.2}{361.234}{\Character{41}}%
\TextAt{3}{335.753}{366.17}{\Character{0}}%
\TextAt{17}{342.339}{366.17}{1}%
\TextAt{13}{347.072}{361.234}{h}%
\TextAt{9}{353.81}{361.234}{\Character{40}}%
\TextAt{13}{358.362}{361.234}{G\Character{59}}%
\TextAt{9}{372.84}{361.234}{\Character{0}\Character{41}}%
\TextAt{9}{387.999}{361.234}{\Character{61}}%
\TextAt{9}{402.216}{369.321}{1}%
\Vector{q 0 G 0 g 0.478 w 329.649 -409.777 m 336.636 -409.777 l S Q}%
\TextAt{13}{401.649}{353.033}{n}%
\TextAt{13}{409.832}{361.234}{h}%
\TextAt{9}{416.57}{361.234}{\Character{40}}%
\TextAt{13}{421.123}{361.234}{G\Character{59}}%
\TextAt{9}{435.6}{361.234}{\Character{0}\Character{41}}%
\TextAt{13}{447.467}{361.234}{\Character{58}}%
\TextAt{8}{79.226}{333.996}{We}%
\TextAt{8}{99.371}{333.996}{now}%
\TextAt{8}{123.759}{333.996}{choose}%
\TextAt{8}{160.892}{333.996}{in}%
\TextAt{9}{174.545}{333.996}{\Character{0}}%
\TextAt{8}{185.756}{333.996}{a}%
\TextAt{8}{195.512}{333.996}{sublattice}%
\TextAt{9}{248.903}{333.996}{\Character{0}}%
\TextAt{3}{256.218}{338.334}{\Character{48}}%
\TextAt{8}{259.013}{333.996}{,}%
\TextAt{8}{266.162}{333.996}{in}%
\TextAt{8}{279.815}{333.996}{accordance}%
\TextAt{8}{339.65}{333.996}{with}%
\TextAt{8}{366.31}{333.996}{lemma}%
\TextAt{8}{404.029}{333.996}{(3.4).}%
\TextAt{8}{436.555}{333.996}{Then}%
\TextAt{8}{467.112}{333.996}{we}%
\TextAt{8}{484.34}{333.996}{have}%
\TextAt{9}{207.875}{307.875}{\Character{0}}%
\TextAt{3}{215.19}{312.811}{\Character{48}}%
\TextAt{9}{221.305}{307.875}{\Character{61}}%
\TextAt{20}{233.731}{319.232}{\Character{77}}%
\TextAt{27}{239.466}{293.796}{P}%
\TextAt{18}{253.789}{307.875}{Z}%
\TextAt{13}{261.759}{307.875}{w}%
\TextAt{27}{270.127}{306.081}{P}%
\TextAt{9}{280.541}{307.875}{\Character{61}}%
\TextAt{20}{292.967}{319.232}{\Character{77}}%
\TextAt{27}{294.385}{293.86}{p}%
\NativeMathAt{298.373}{293.861}{7.97}{\in}%
\TextAt{21}{304.019}{293.86}{S}%
\TextAt{20}{313.025}{319.232}{\Character{77}}%
\TextAt{27}{315.59}{293.329}{P}%
\TextAt{3}{322.185}{293.329}{\Character{106}}%
\TextAt{27}{324.537}{293.329}{p}%
\TextAt{18}{333.083}{307.875}{Z}%
\TextAt{13}{341.053}{307.875}{w}%
\TextAt{27}{349.421}{306.081}{P}%
\TextAt{9}{359.836}{307.875}{\Character{61}}%
\TextAt{20}{372.261}{319.232}{\Character{77}}%
\TextAt{27}{373.679}{293.861}{p}%
\NativeMathAt{377.668}{293.861}{7.97}{\in}%
\TextAt{21}{383.313}{293.861}{S}%
\TextAt{9}{392.319}{307.875}{\Character{0}}%
\TextAt{3}{399.633}{312.811}{\Character{48}}%
\TextAt{27}{399.633}{304.919}{p}%
\TextAt{8}{79.226}{271.304}{and}%
\TextAt{13}{101.989}{271.304}{\Character{27}}%
\TextAt{13}{109.055}{271.304}{w}%
\TextAt{27}{117.439}{269.51}{P}%
\TextAt{9}{127.853}{271.304}{\Character{61}}%
\TextAt{13}{140.279}{271.304}{w}%
\TextAt{21}{148.647}{269.51}{\Character{27}}%
\TextAt{27}{153.772}{269.51}{P}%
\TextAt{8}{160.866}{271.304}{.}%
\TextAt{8}{169.314}{271.304}{This}%
\TextAt{8}{196.034}{271.304}{identifies}%
\TextAt{9}{245.193}{271.304}{\Character{0}}%
\TextAt{3}{252.508}{275.642}{\Character{48}}%
\TextAt{27}{252.508}{268.348}{p}%
\TextAt{8}{260.892}{271.304}{as}%
\TextAt{8}{275.262}{271.304}{the}%
\TextAt{8}{295.418}{271.304}{induced}%
\TextAt{13}{338.983}{271.304}{G}%
\TextAt{8}{348.217}{271.304}{-module}%
\TextAt{9}{175.271}{244.022}{\Character{0}}%
\TextAt{3}{182.585}{248.958}{\Character{48}}%
\TextAt{27}{182.585}{241.066}{p}%
\TextAt{9}{190.393}{244.022}{\Character{61}}%
\TextAt{20}{202.818}{255.379}{\Character{77}}%
\TextAt{27}{205.383}{229.476}{P}%
\TextAt{3}{211.978}{229.476}{\Character{106}}%
\TextAt{27}{214.331}{229.476}{p}%
\TextAt{18}{222.876}{244.022}{Z}%
\TextAt{13}{230.847}{244.022}{w}%
\TextAt{27}{239.214}{242.229}{P}%
\TextAt{9}{249.629}{244.022}{\Character{61}}%
\TextAt{20}{272.297}{255.379}{\Character{77}}%
\TextAt{21}{262.054}{229.476}{\Character{27}}%
\NativeMathAt{267.18}{229.476}{7.97}{\in}%
\TextAt{21}{272.825}{229.476}{G}%
\TextAt{17}{279.424}{229.476}{\Character{47}}%
\TextAt{21}{283.658}{229.476}{G}%
\TextAt{28}{290.257}{228.085}{P}%
\TextAt{24}{295.956}{226.292}{0}%
\TextAt{13}{302.597}{244.022}{\Character{27}}%
\TextAt{9}{309.68}{244.022}{\Character{40}}%
\TextAt{18}{314.232}{244.022}{Z}%
\TextAt{13}{322.202}{244.022}{w}%
\TextAt{27}{330.57}{242.229}{P}%
\TextAt{24}{337.166}{241.122}{0}%
\TextAt{9}{341.815}{244.022}{\Character{41}}%
\TextAt{9}{349.67}{244.022}{\Character{61}}%
\TextAt{8}{362.113}{244.022}{Ind}%
\TextAt{21}{379.341}{251.324}{G}%
\TextAt{28}{385.939}{249.933}{P}%
\TextAt{24}{391.639}{248.14}{0}%
\TextAt{21}{379.341}{240.781}{G}%
\TextAt{9}{396.786}{244.022}{\Character{40}}%
\TextAt{18}{401.338}{244.022}{Z}%
\TextAt{13}{409.308}{244.022}{w}%
\TextAt{27}{417.676}{242.229}{P}%
\TextAt{24}{424.272}{241.122}{0}%
\TextAt{9}{428.921}{244.022}{\Character{41}}%
\TextAt{13}{433.473}{244.022}{\Character{59}}%
\TextAt{8}{79.226}{206.259}{where}%
\TextAt{14}{113.364}{206.259}{P}%
\TextAt{17}{123.257}{204.466}{0}%
\TextAt{8}{132.216}{206.259}{is}%
\TextAt{8}{144.309}{206.259}{a}%
\TextAt{8}{154.393}{206.259}{chosen}%
\TextAt{8}{192.177}{206.259}{place}%
\TextAt{8}{222.418}{206.259}{above}%
\TextAt{14}{255.922}{206.259}{p}%
\TextAt{8}{261.906}{206.259}{,}%
\TextAt{8}{269.466}{206.259}{and}%
\TextAt{13}{292.557}{206.259}{G}%
\TextAt{27}{301.791}{204.466}{P}%
\TextAt{24}{308.386}{203.359}{0}%
\TextAt{8}{317.262}{206.259}{is}%
\TextAt{8}{329.355}{206.259}{its}%
\TextAt{8}{346.002}{206.259}{decomposition}%
\TextAt{8}{423.466}{206.259}{group.}%
\TextAt{8}{462.179}{206.259}{The}%
\TextAt{8}{486.562}{206.259}{lattice}%
\TextAt{9}{522.66}{206.259}{\Character{0}}%
\TextAt{3}{529.974}{210.597}{\Character{48}}%
\TextAt{8}{79.226}{189.213}{has}%
\TextAt{8}{101.154}{189.213}{the}%
\TextAt{8}{122.365}{189.213}{same}%
\TextAt{8}{152.745}{189.213}{rank}%
\TextAt{8}{180.794}{189.213}{as}%
\TextAt{9}{196.218}{189.213}{\Character{0}}%
\TextAt{8}{203.533}{189.213}{,}%
\TextAt{8}{212}{189.213}{so}%
\TextAt{8}{227.424}{189.213}{is}%
\TextAt{8}{240.242}{189.213}{therefore}%
\TextAt{8}{290.396}{189.213}{of}%
\TextAt{8}{304.78}{189.213}{finite}%
\TextAt{8}{335.746}{189.213}{index}%
\TextAt{8}{368.338}{189.213}{in}%
\TextAt{9}{383.045}{189.213}{\Character{0}}%
\TextAt{8}{390.36}{189.213}{.}%
\TextAt{8}{401.971}{189.213}{From}%
\TextAt{8}{433.75}{189.213}{chap.}%
\TextAt{8}{465.696}{189.213}{IV,}%
\TextAt{8}{486.895}{189.213}{(7.4),}%
\TextAt{8}{519.439}{189.213}{we}%
\TextAt{8}{79.226}{172.167}{conclude}%
\TextAt{8}{127.346}{172.167}{that}%
\TextAt{13}{182.605}{146.046}{h}%
\TextAt{9}{189.344}{146.046}{\Character{40}}%
\TextAt{13}{193.896}{146.046}{G\Character{59}}%
\TextAt{13}{208.35}{146.046}{L}%
\TextAt{21}{216.338}{150.982}{S}%
\TextAt{9}{222.432}{146.046}{\Character{41}}%
\TextAt{9}{230.287}{146.046}{\Character{61}}%
\TextAt{9}{244.494}{154.134}{1}%
\Vector{q 0 G 0 g 0.478 w 171.926 -624.965 m 178.914 -624.965 l S Q}%
\TextAt{13}{243.926}{137.846}{n}%
\TextAt{13}{252.11}{146.046}{h}%
\TextAt{9}{258.848}{146.046}{\Character{40}}%
\TextAt{13}{263.4}{146.046}{G\Character{59}}%
\TextAt{9}{277.878}{146.046}{\Character{0}}%
\TextAt{3}{285.192}{150.982}{\Character{48}}%
\TextAt{9}{287.988}{146.046}{\Character{41}}%
\TextAt{9}{295.843}{146.046}{\Character{61}}%
\TextAt{9}{310.049}{154.134}{1}%
\Vector{q 0 G 0 g 0.478 w 237.482 -624.965 m 244.469 -624.965 l S Q}%
\TextAt{13}{309.482}{137.846}{n}%
\TextAt{20}{319.657}{157.404}{\Character{89}}%
\TextAt{27}{319.68}{132.032}{p}%
\NativeMathAt{323.669}{132.032}{7.97}{\in}%
\TextAt{21}{329.315}{132.032}{S}%
\TextAt{13}{336.926}{146.046}{h}%
\TextAt{9}{343.664}{146.046}{\Character{40}}%
\TextAt{13}{348.217}{146.046}{G\Character{59}}%
\TextAt{9}{362.695}{146.046}{\Character{0}}%
\TextAt{3}{370.009}{150.982}{\Character{48}}%
\TextAt{27}{370.009}{143.091}{p}%
\TextAt{9}{374.496}{146.046}{\Character{41}}%
\TextAt{9}{230.306}{109.552}{\Character{61}}%
\TextAt{9}{244.494}{117.64}{1}%
\Vector{q 0 G 0 g 0.478 w 171.926 -661.459 m 178.914 -661.459 l S Q}%
\TextAt{13}{243.926}{101.352}{n}%
\TextAt{20}{254.102}{120.91}{\Character{89}}%
\TextAt{27}{254.125}{95.538}{p}%
\NativeMathAt{258.114}{95.538}{7.97}{\in}%
\TextAt{21}{263.759}{95.538}{S}%
\TextAt{13}{271.371}{109.552}{h}%
\TextAt{9}{278.109}{109.552}{\Character{40}}%
\TextAt{13}{282.662}{109.552}{G}%
\TextAt{27}{291.895}{107.759}{P}%
\TextAt{13}{298.989}{109.552}{\Character{59}}%
\TextAt{18}{304.233}{109.552}{Z}%
\TextAt{13}{312.203}{109.552}{w}%
\TextAt{27}{320.571}{107.759}{P}%
\TextAt{24}{327.166}{106.652}{0}%
\TextAt{9}{331.816}{109.552}{\Character{41}}%
\TextAt{9}{339.671}{109.552}{\Character{61}}%
\TextAt{9}{353.877}{117.64}{1}%
\Vector{q 0 G 0 g 0.478 w 281.31 -661.459 m 288.297 -661.459 l S Q}%
\TextAt{13}{353.31}{101.352}{n}%
\TextAt{20}{363.485}{120.91}{\Character{89}}%
\TextAt{27}{363.508}{95.538}{p}%
\NativeMathAt{367.497}{95.538}{7.97}{\in}%
\TextAt{21}{373.143}{95.538}{S}%
\TextAt{13}{380.754}{109.552}{h}%
\TextAt{9}{387.493}{109.552}{\Character{40}}%
\TextAt{13}{392.045}{109.552}{G}%
\TextAt{27}{401.278}{107.759}{P}%
\TextAt{13}{408.372}{109.552}{\Character{59}}%
\TextAt{18}{413.616}{109.552}{Z}%
\TextAt{9}{421.586}{109.552}{\Character{41}}%
\TextAt{13}{426.139}{109.552}{\Character{58}}%
\TextAt{8}{79.226}{72.921}{Thus}%
\TextAt{8}{108.875}{72.921}{we}%
\TextAt{8}{126.102}{72.921}{do}%
\TextAt{8}{142.361}{72.921}{find}%
\TextAt{8}{165.77}{72.921}{that}%
\TextAt{13}{191.139}{72.921}{h}%
\TextAt{9}{197.877}{72.921}{\Character{40}}%
\TextAt{13}{202.429}{72.921}{G\Character{59}}%
\TextAt{13}{216.883}{72.921}{L}%
\TextAt{21}{224.872}{77.26}{S}%
\TextAt{9}{230.966}{72.921}{\Character{41}}%
\TextAt{9}{238.821}{72.921}{\Character{61}}%
\TextAt{17}{252.912}{77.628}{1}%
\Vector{q 0 G 0 g 0.478 w 180.46 -698.09 m 185.598 -698.09 l S Q}%
\TextAt{21}{252.46}{68.799}{n}%
\TextAt{20}{260.786}{81.888}{\Character{81}}%
\TextAt{27}{272.077}{69.434}{p}%
\NativeMathAt{276.066}{69.434}{7.97}{\in}%
\TextAt{21}{281.711}{69.434}{S}%
\TextAt{13}{289.798}{72.921}{n}%
\TextAt{27}{296.785}{71.128}{p}%
\TextAt{8}{301.272}{72.921}{,}%
\TextAt{8}{308.421}{72.921}{where}%
\TextAt{13}{342.23}{72.921}{n}%
\TextAt{27}{349.218}{71.128}{p}%
\TextAt{9}{357.026}{72.921}{\Character{61}}%
\TextAt{8}{369.451}{72.921}{\char35\relax }%
\TextAt{13}{379.207}{72.921}{G}%
\TextAt{27}{388.44}{71.128}{P}%
\TextAt{9}{398.855}{72.921}{\Character{61}}%
\TextAt{9}{411.264}{72.921}{\Character{91}}%
\TextAt{13}{414.532}{72.921}{L}%
\TextAt{27}{422.496}{71.128}{P}%
\TextAt{9}{432.911}{72.921}{\Character{58}}%
\TextAt{13}{439.483}{72.921}{K}%
\TextAt{27}{449.442}{71.128}{p}%
\TextAt{9}{453.929}{72.921}{\Character{93}}%
\TextAt{8}{457.181}{72.921}{.}%
\Vector{q 0 G 0 g 0.399 w 452.334 -701.079 m 452.334 -693.009 l S Q}%
\Vector{q 0 G 0 g 0.399 w 452.534 -693.208 m 459.707 -693.208 l S Q}%
\Vector{q 0 G 0 g 0.399 w 452.534 -700.879 m 459.707 -700.879 l S Q}%
\Vector{q 0 G 0 g 0.399 w 459.906 -701.079 m 459.906 -693.009 l S Q}%
\TextAt{8}{297.211}{43.033}{344}%
\Destination{78.23}{809.815}{nav.page.344}%
\Destination{79.36}{775.155}{nav.number.6.3.5}%
\Hotlink{465.196}{187.055}{13.995}{11.047}{nav.chapter.4}%
\Hotlink{78.726}{643.955}{25.078}{11.047}{nav.number.6.1.1}%
\Hotlink{403.529}{331.755}{25.078}{11.047}{nav.number.6.3.4}%
\Hotlink{486.395}{187.055}{25.078}{11.047}{nav.number.4.7.4}%
}

% Page 357 (body page 345)
\Page{%
\TextAt{2}{79.226}{796.943}{VI.}%
\TextAt{2}{98.012}{796.943}{Global}%
\TextAt{2}{133.248}{796.943}{Class}%
\TextAt{2}{161.841}{796.943}{Field}%
\TextAt{2}{189.572}{796.943}{Theory}%
\Vector{q 0 G 0 g 0.299 w 7.226 18.729 m 460.769 18.729 l S Q}%
\TextAt{8}{97.159}{761.124}{From}%
\TextAt{8}{127.719}{761.124}{the}%
\TextAt{8}{147.711}{761.124}{Herbrand}%
\TextAt{8}{199.898}{761.124}{quotient}%
\TextAt{8}{245.916}{761.124}{of}%
\TextAt{13}{259.081}{761.124}{I}%
\TextAt{21}{265.184}{765.462}{S}%
\TextAt{21}{264.244}{758.103}{L}%
\TextAt{8}{275.01}{761.124}{and}%
\TextAt{13}{297.608}{761.124}{L}%
\TextAt{21}{305.573}{765.462}{S}%
\TextAt{8}{315.399}{761.124}{we}%
\TextAt{8}{332.462}{761.124}{immediately}%
\TextAt{8}{398.959}{761.124}{get}%
\TextAt{8}{418.305}{761.124}{the}%
\TextAt{8}{438.297}{761.124}{Herbrand}%
\TextAt{8}{490.484}{761.124}{quotient}%
\TextAt{8}{79.226}{744.078}{of}%
\TextAt{8}{92.566}{744.078}{the}%
\TextAt{8}{112.733}{744.078}{idèle}%
\TextAt{8}{140.049}{744.078}{class}%
\TextAt{8}{167.496}{744.078}{group}%
\TextAt{13}{200.682}{744.078}{C}%
\TextAt{21}{209.057}{742.284}{L}%
\TextAt{8}{215.314}{744.078}{.}%
\TextAt{8}{223.793}{744.078}{To}%
\TextAt{8}{241.03}{744.078}{do}%
\TextAt{8}{257.3}{744.078}{it}%
\TextAt{8}{269.014}{744.078}{choose}%
\TextAt{8}{306.157}{744.078}{a}%
\TextAt{8}{315.922}{744.078}{finite}%
\TextAt{8}{345.844}{744.078}{set}%
\TextAt{8}{364.122}{744.078}{of}%
\TextAt{8}{377.462}{744.078}{places}%
\TextAt{13}{411.999}{744.078}{S}%
\TextAt{8}{423.802}{744.078}{containing}%
\TextAt{8}{480.73}{744.078}{all}%
\TextAt{8}{497}{744.078}{infinite}%
\TextAt{8}{79.226}{727.032}{ones}%
\TextAt{8}{105.291}{727.032}{and}%
\TextAt{8}{128.045}{727.032}{all}%
\TextAt{8}{144.294}{727.032}{primes}%
\TextAt{8}{182.064}{727.032}{ramified}%
\TextAt{8}{227.58}{727.032}{in}%
\TextAt{13}{241.223}{727.032}{L}%
\TextAt{8}{249.188}{727.032}{,}%
\TextAt{8}{256.33}{727.032}{such}%
\TextAt{8}{282.718}{727.032}{that}%
\TextAt{13}{308.077}{727.032}{I}%
\TextAt{21}{313.24}{725.238}{L}%
\TextAt{9}{322.818}{727.032}{\Character{61}}%
\TextAt{13}{335.244}{727.032}{I}%
\TextAt{21}{341.346}{731.37}{S}%
\TextAt{21}{340.407}{724.011}{L}%
\TextAt{13}{347.44}{727.032}{L}%
\TextAt{3}{355.405}{731.37}{\Character{3}}%
\TextAt{8}{360.137}{727.032}{.}%
\TextAt{8}{368.582}{727.032}{Such}%
\TextAt{8}{396.859}{727.032}{a}%
\TextAt{8}{406.605}{727.032}{set}%
\TextAt{8}{424.864}{727.032}{exists}%
\TextAt{8}{457.169}{727.032}{by}%
\TextAt{8}{473.419}{727.032}{(1.4).}%
\TextAt{8}{505.942}{727.032}{From}%
\TextAt{8}{79.226}{709.986}{the}%
\TextAt{8}{99.383}{709.986}{exact}%
\TextAt{8}{130.275}{709.986}{sequence}%
\TextAt{9}{213.49}{692.94}{1}%
\NativeArrowAt{222.664}{692.94}{11.955}{19.26}{\longrightarrow}%
\TextAt{13}{245.246}{692.94}{L}%
\TextAt{21}{253.21}{697.876}{S}%
\NativeArrowAt{262.625}{692.94}{11.955}{19.26}{\longrightarrow}%
\TextAt{13}{285.207}{692.94}{I}%
\TextAt{21}{291.31}{697.876}{S}%
\TextAt{21}{290.37}{689.984}{L}%
\NativeArrowAt{300.725}{692.94}{11.955}{19.26}{\longrightarrow}%
\TextAt{13}{323.307}{692.94}{I}%
\TextAt{21}{329.409}{697.876}{S}%
\TextAt{21}{328.469}{689.984}{L}%
\TextAt{13}{335.503}{692.94}{L}%
\TextAt{3}{343.468}{697.876}{\Character{3}}%
\TextAt{9}{348.2}{692.94}{\Character{47}}%
\TextAt{13}{354.053}{692.94}{L}%
\TextAt{3}{362.018}{697.876}{\Character{3}}%
\NativeArrowAt{370.071}{692.94}{11.955}{19.26}{\longrightarrow}%
\TextAt{9}{392.653}{692.94}{1}%
\TextAt{8}{79.226}{670.794}{arises}%
\TextAt{8}{110.574}{670.794}{the}%
\TextAt{8}{130.087}{670.794}{identity}%
\TextAt{13}{172.697}{670.794}{h}%
\TextAt{9}{179.436}{670.794}{\Character{40}}%
\TextAt{13}{183.988}{670.794}{G\Character{59}}%
\TextAt{13}{198.442}{670.794}{C}%
\TextAt{21}{206.841}{669.001}{L}%
\TextAt{9}{213.098}{670.794}{\Character{41}}%
\TextAt{9}{220.953}{670.794}{\Character{61}}%
\TextAt{13}{233.396}{670.794}{h}%
\TextAt{9}{240.135}{670.794}{\Character{40}}%
\TextAt{13}{244.687}{670.794}{G\Character{59}}%
\TextAt{13}{259.141}{670.794}{I}%
\TextAt{21}{265.268}{675.133}{S}%
\TextAt{21}{264.328}{667.774}{L}%
\TextAt{9}{271.362}{670.794}{\Character{41}}%
\TextAt{13}{275.914}{670.794}{h}%
\TextAt{9}{282.653}{670.794}{\Character{40}}%
\TextAt{13}{287.205}{670.794}{G\Character{59}}%
\TextAt{13}{301.659}{670.794}{L}%
\TextAt{21}{309.647}{675.133}{S}%
\TextAt{9}{315.741}{670.794}{\Character{41}}%
\TextAt{3}{320.294}{675.133}{\Character{0}}%
\TextAt{17}{326.88}{675.133}{1}%
\TextAt{8}{331.612}{670.794}{,}%
\TextAt{8}{338.246}{670.794}{and}%
\TextAt{8}{360.365}{670.794}{from}%
\TextAt{8}{387.362}{670.794}{(3.3)}%
\TextAt{8}{414.693}{670.794}{and}%
\TextAt{8}{436.812}{670.794}{(3.5)}%
\TextAt{8}{464.143}{670.794}{we}%
\TextAt{8}{480.727}{670.794}{obtain}%
\TextAt{8}{516.51}{670.794}{the}%
\TextAt{12}{79.226}{645.699}{Theorem}%
\TextAt{12}{136.325}{645.699}{(3.6).}%
\TextAt{8}{171.713}{645.699}{If}%
\TextAt{13}{183.479}{645.699}{L}%
\TextAt{16}{194.889}{645.699}{\Character{106}}%
\TextAt{13}{201.655}{645.699}{K}%
\TextAt{8}{216.443}{645.699}{is}%
\TextAt{8}{228.28}{645.699}{a}%
\TextAt{8}{238.109}{645.699}{cyclic}%
\TextAt{8}{270.366}{645.699}{extension}%
\TextAt{8}{322.205}{645.699}{of}%
\TextAt{8}{335.608}{645.699}{degree}%
\TextAt{13}{372.096}{645.699}{n}%
\TextAt{8}{383.055}{645.699}{with}%
\TextAt{8}{409.788}{645.699}{Galois}%
\TextAt{8}{445.774}{645.699}{group}%
\TextAt{13}{479.023}{645.699}{G}%
\TextAt{9}{491.702}{645.699}{\Character{61}}%
\TextAt{13}{504.253}{645.699}{G}%
\TextAt{9}{513.486}{645.699}{\Character{40}}%
\TextAt{13}{518.039}{645.699}{L}%
\TextAt{16}{529.449}{645.699}{\Character{106}}%
\TextAt{13}{79.226}{628.653}{K}%
\TextAt{9}{90.044}{628.653}{\Character{41}}%
\TextAt{8}{94.596}{628.653}{,}%
\TextAt{8}{101.745}{628.653}{then}%
\TextAt{13}{226.504}{607.106}{h}%
\TextAt{9}{233.243}{607.106}{\Character{40}}%
\TextAt{13}{237.795}{607.106}{G\Character{59}}%
\TextAt{13}{252.249}{607.106}{C}%
\TextAt{21}{260.648}{605.312}{L}%
\TextAt{9}{266.905}{607.106}{\Character{41}}%
\TextAt{9}{274.76}{607.106}{\Character{61}}%
\TextAt{8}{291.692}{615.193}{\char35\relax }%
\TextAt{13}{301.448}{615.193}{H}%
\TextAt{17}{312.07}{619.532}{0}%
\TextAt{9}{316.803}{615.193}{\Character{40}}%
\TextAt{13}{321.355}{615.193}{G\Character{59}}%
\TextAt{13}{335.809}{615.193}{C}%
\TextAt{21}{344.208}{613.4}{L}%
\TextAt{9}{350.465}{615.193}{\Character{41}}%
\Vector{q 0 G 0 g 0.478 w 216.399 -163.906 m 286.311 -163.906 l S Q}%
\TextAt{8}{288.399}{598.905}{\char35\relax }%
\TextAt{13}{298.154}{598.905}{H}%
\TextAt{3}{308.777}{602.359}{\Character{0}}%
\TextAt{17}{315.364}{602.359}{1}%
\TextAt{9}{320.096}{598.905}{\Character{40}}%
\TextAt{13}{324.648}{598.905}{G\Character{59}}%
\TextAt{13}{339.102}{598.905}{C}%
\TextAt{21}{347.501}{597.112}{L}%
\TextAt{9}{353.758}{598.905}{\Character{41}}%
\TextAt{9}{362.827}{607.106}{\Character{61}}%
\TextAt{13}{375.252}{607.106}{n\Character{58}}%
\TextAt{8}{79.226}{580.854}{In}%
\TextAt{8}{93.847}{580.854}{particular}%
\TextAt{9}{147.837}{580.854}{\Character{40}}%
\TextAt{13}{152.389}{580.854}{C}%
\TextAt{21}{160.764}{579.061}{K}%
\TextAt{9}{172.288}{580.854}{\Character{58}}%
\TextAt{13}{178.86}{580.854}{N}%
\TextAt{21}{188.234}{578.994}{L}%
\TextAt{3}{193.993}{578.994}{\Character{106}}%
\TextAt{21}{196.345}{578.994}{K}%
\TextAt{13}{204.548}{580.854}{C}%
\TextAt{21}{212.923}{579.061}{L}%
\TextAt{9}{219.18}{580.854}{\Character{41}}%
\NativeMathAt{227.053}{580.854}{11.955}{\geq}%
\TextAt{13}{239.673}{580.854}{n}%
\TextAt{8}{246.66}{580.854}{.}%
\TextAt{8}{97.159}{555.759}{From}%
\TextAt{8}{127.884}{555.759}{this}%
\TextAt{8}{150.706}{555.759}{result}%
\TextAt{8}{183.284}{555.759}{we}%
\TextAt{8}{200.511}{555.759}{deduce}%
\TextAt{8}{239.521}{555.759}{the}%
\TextAt{8}{259.677}{555.759}{following}%
\TextAt{8}{309.112}{555.759}{interesting}%
\TextAt{8}{366.736}{555.759}{consequence.}%
\TextAt{12}{79.226}{530.664}{Corollary}%
\TextAt{12}{138.918}{530.664}{(3.7).}%
\TextAt{8}{174.15}{530.664}{If}%
\TextAt{13}{185.759}{530.664}{L}%
\TextAt{16}{197.044}{530.664}{\Character{106}}%
\TextAt{13}{203.686}{530.664}{K}%
\TextAt{8}{218.317}{530.664}{is}%
\TextAt{8}{229.998}{530.664}{cyclic}%
\TextAt{8}{262.098}{530.664}{of}%
\TextAt{8}{275.344}{530.664}{prime}%
\TextAt{8}{308.424}{530.664}{power}%
\TextAt{8}{342.485}{530.664}{degree}%
\TextAt{13}{378.817}{530.664}{n}%
\TextAt{9}{389.125}{530.664}{\Character{61}}%
\TextAt{13}{401.551}{530.664}{p}%
\TextAt{21}{407.426}{535.002}{\Character{23}}%
\TextAt{8}{416.422}{530.664}{(}%
\TextAt{13}{420.977}{530.664}{\Character{23}}%
\TextAt{13}{430.816}{530.664}{\Character{62}}%
\TextAt{9}{443.261}{530.664}{0}%
\TextAt{8}{449.114}{530.664}{),}%
\TextAt{8}{460.751}{530.664}{then}%
\TextAt{8}{487.328}{530.664}{there}%
\TextAt{8}{517.156}{530.664}{are}%
\TextAt{8}{79.226}{513.618}{infinitely}%
\TextAt{8}{128.326}{513.618}{many}%
\TextAt{8}{160.199}{513.618}{places}%
\TextAt{8}{194.725}{513.618}{of}%
\TextAt{13}{208.055}{513.618}{K}%
\TextAt{8}{222.77}{513.618}{which}%
\TextAt{8}{256.257}{513.618}{do}%
\TextAt{8}{272.516}{513.618}{not}%
\TextAt{8}{293.33}{513.618}{split}%
\TextAt{8}{319.404}{513.618}{in}%
\TextAt{13}{333.057}{513.618}{L}%
\TextAt{8}{341.021}{513.618}{.}%
\TextAt{19}{79.226}{480.474}{Proof.}%
\TextAt{8}{116}{480.474}{Assume}%
\TextAt{8}{160.021}{480.474}{that}%
\TextAt{8}{186.05}{480.474}{the}%
\TextAt{8}{206.866}{480.474}{set}%
\TextAt{13}{225.793}{480.474}{S}%
\TextAt{8}{238.246}{480.474}{of}%
\TextAt{8}{252.236}{480.474}{nonsplit}%
\TextAt{8}{297.835}{480.474}{primes}%
\TextAt{8}{336.273}{480.474}{were}%
\TextAt{8}{363.915}{480.474}{finite.}%
\TextAt{8}{400.357}{480.474}{Let}%
\TextAt{13}{421.986}{480.474}{M}%
\TextAt{16}{439.005}{480.474}{\Character{106}}%
\TextAt{13}{446.771}{480.474}{K}%
\TextAt{8}{462.146}{480.474}{be}%
\TextAt{8}{478.73}{480.474}{the}%
\TextAt{8}{499.546}{480.474}{subex-}%
\TextAt{8}{79.226}{463.428}{tension}%
\TextAt{8}{120.014}{463.428}{of}%
\TextAt{13}{133.748}{463.428}{L}%
\TextAt{16}{145.721}{463.428}{\Character{106}}%
\TextAt{13}{153.05}{463.428}{K}%
\TextAt{8}{168.169}{463.428}{of}%
\TextAt{8}{181.902}{463.428}{degree}%
\TextAt{13}{218.721}{463.428}{p}%
\TextAt{8}{224.596}{463.428}{.}%
\TextAt{8}{234.255}{463.428}{For}%
\TextAt{8}{255.628}{463.428}{every}%
\TextAt{14}{286.924}{463.428}{p}%
\NativeMathAt{296.915}{463.428}{11.955}{\notin}%
\TextAt{13}{308.894}{463.428}{S}%
\TextAt{8}{316.789}{463.428}{,}%
\TextAt{8}{324.443}{463.428}{the}%
\TextAt{8}{345.003}{463.428}{decomposition}%
\TextAt{8}{422.541}{463.428}{group}%
\TextAt{13}{456.12}{463.428}{G}%
\TextAt{27}{465.354}{461.635}{P}%
\TextAt{8}{476.748}{463.428}{of}%
\TextAt{13}{490.482}{463.428}{L}%
\TextAt{16}{502.455}{463.428}{\Character{106}}%
\TextAt{13}{509.784}{463.428}{K}%
\TextAt{8}{524.903}{463.428}{is}%
\TextAt{8}{79.226}{446.382}{different}%
\TextAt{8}{125.129}{446.382}{from}%
\TextAt{13}{152.501}{446.382}{G}%
\TextAt{9}{161.735}{446.382}{\Character{40}}%
\TextAt{13}{166.287}{446.382}{L}%
\TextAt{16}{177.572}{446.382}{\Character{106}}%
\TextAt{13}{184.214}{446.382}{K}%
\TextAt{9}{195.032}{446.382}{\Character{41}}%
\TextAt{8}{199.584}{446.382}{.}%
\TextAt{8}{207.943}{446.382}{Hence}%
\TextAt{13}{242.453}{446.382}{G}%
\TextAt{27}{251.686}{444.589}{P}%
\NativeMathAt{262.101}{446.382}{11.955}{\subseteq}%
\TextAt{13}{274.72}{446.382}{G}%
\TextAt{9}{283.954}{446.382}{\Character{40}}%
\TextAt{13}{288.506}{446.382}{L}%
\TextAt{16}{299.791}{446.382}{\Character{106}}%
\TextAt{13}{306.433}{446.382}{M}%
\TextAt{9}{319.007}{446.382}{\Character{41}}%
\TextAt{8}{323.559}{446.382}{.}%
\TextAt{8}{331.918}{446.382}{Therefore}%
\TextAt{8}{384.647}{446.382}{every}%
\TextAt{14}{415.271}{446.382}{p}%
\NativeMathAt{424.575}{446.382}{11.955}{\notin}%
\TextAt{13}{435.866}{446.382}{S}%
\TextAt{8}{447.391}{446.382}{splits}%
\TextAt{8}{477.812}{446.382}{completely}%
\TextAt{8}{79.226}{429.336}{in}%
\TextAt{13}{92.879}{429.336}{M}%
\TextAt{8}{105.453}{429.336}{.}%
\TextAt{8}{113.901}{429.336}{We}%
\TextAt{8}{134.045}{429.336}{deduce}%
\TextAt{8}{173.055}{429.336}{from}%
\TextAt{8}{200.695}{429.336}{this}%
\TextAt{8}{223.518}{429.336}{that}%
\TextAt{13}{248.887}{429.336}{N}%
\TextAt{21}{258.261}{427.476}{M}%
\TextAt{3}{267.242}{427.476}{\Character{106}}%
\TextAt{21}{269.595}{427.476}{K}%
\TextAt{13}{277.797}{429.336}{C}%
\TextAt{21}{286.172}{427.543}{M}%
\TextAt{9}{298.973}{429.336}{\Character{61}}%
\TextAt{13}{311.398}{429.336}{C}%
\TextAt{21}{319.773}{427.543}{K}%
\TextAt{8}{327.976}{429.336}{,}%
\TextAt{8}{335.125}{429.336}{thus}%
\TextAt{8}{360.877}{429.336}{contradicting}%
\TextAt{8}{432.106}{429.336}{(3.6).}%
\TextAt{8}{97.159}{409.301}{Indeed,}%
\TextAt{8}{138.798}{409.301}{let}%
\TextAt{13}{155.989}{409.301}{\Character{11}}%
\NativeMathAt{167.32}{409.301}{11.955}{\in}%
\TextAt{13}{179.099}{409.301}{I}%
\TextAt{21}{184.262}{407.508}{K}%
\TextAt{8}{192.464}{409.301}{.}%
\TextAt{8}{201.772}{409.301}{By}%
\TextAt{8}{220.422}{409.301}{the}%
\TextAt{8}{240.864}{409.301}{approximation}%
\TextAt{8}{319.218}{409.301}{theorem}%
\TextAt{8}{365.03}{409.301}{of}%
\TextAt{8}{378.646}{409.301}{chap.}%
\TextAt{8}{409.825}{409.301}{II,}%
\TextAt{8}{426.023}{409.301}{(3.4),}%
\TextAt{8}{457.608}{409.301}{there}%
\TextAt{8}{487.807}{409.301}{exists}%
\TextAt{8}{520.408}{409.301}{an}%
\TextAt{13}{79.226}{392.255}{a}%
\NativeMathAt{89.312}{392.255}{11.955}{\in}%
\TextAt{13}{101.223}{392.255}{K}%
\TextAt{3}{112.04}{396.594}{\Character{3}}%
\TextAt{8}{121.034}{392.255}{such}%
\TextAt{8}{147.795}{392.255}{that}%
\TextAt{13}{173.527}{392.255}{\Character{11}}%
\TextAt{27}{180.973}{390.462}{p}%
\TextAt{13}{185.46}{392.255}{a}%
\TextAt{3}{191.605}{396.594}{\Character{0}}%
\TextAt{17}{198.191}{396.594}{1}%
\TextAt{8}{207.185}{392.255}{is}%
\TextAt{8}{219.312}{392.255}{contained}%
\TextAt{8}{272.685}{392.255}{in}%
\TextAt{8}{286.702}{392.255}{the}%
\TextAt{8}{307.222}{392.255}{open}%
\TextAt{8}{335.872}{392.255}{subgroup}%
\TextAt{13}{387.033}{392.255}{N}%
\TextAt{21}{396.407}{390.396}{M}%
\TextAt{28}{404.54}{389.005}{P}%
\TextAt{3}{410.737}{390.396}{\Character{106}}%
\TextAt{21}{413.089}{390.396}{K}%
\TextAt{28}{420.23}{389.399}{p}%
\TextAt{13}{424.811}{392.255}{M}%
\TextAt{3}{437.385}{396.594}{\Character{3}}%
\TextAt{27}{436.136}{389.17}{P}%
\TextAt{8}{443.23}{392.255}{,}%
\TextAt{8}{450.834}{392.255}{for}%
\TextAt{8}{469.082}{392.255}{all}%
\TextAt{14}{485.705}{392.255}{p}%
\NativeMathAt{495.629}{392.255}{11.955}{\in}%
\TextAt{13}{507.54}{392.255}{S}%
\TextAt{8}{515.435}{392.255}{.}%
\TextAt{8}{524.975}{392.255}{If}%
\TextAt{14}{79.226}{375.209}{p}%
\NativeMathAt{88.53}{375.209}{11.955}{\notin}%
\TextAt{13}{99.821}{375.209}{S}%
\TextAt{8}{107.716}{375.209}{,}%
\TextAt{8}{114.866}{375.209}{then}%
\TextAt{13}{141.526}{375.209}{\Character{11}}%
\TextAt{27}{148.971}{373.416}{p}%
\TextAt{13}{153.458}{375.209}{a}%
\TextAt{3}{159.603}{379.548}{\Character{0}}%
\TextAt{17}{166.189}{379.548}{1}%
\TextAt{8}{174.819}{375.209}{is}%
\TextAt{8}{186.583}{375.209}{automatically}%
\TextAt{8}{260.418}{375.209}{contained}%
\TextAt{8}{313.427}{375.209}{in}%
\TextAt{13}{327.08}{375.209}{N}%
\TextAt{21}{336.454}{373.35}{M}%
\TextAt{28}{344.587}{371.959}{P}%
\TextAt{3}{350.784}{373.35}{\Character{106}}%
\TextAt{21}{353.136}{373.35}{K}%
\TextAt{28}{360.278}{372.353}{p}%
\TextAt{13}{364.859}{375.209}{M}%
\TextAt{3}{377.432}{379.548}{\Character{3}}%
\TextAt{27}{376.183}{372.124}{P}%
\TextAt{8}{387.174}{375.209}{because}%
\TextAt{13}{430.476}{375.209}{M}%
\TextAt{27}{441.801}{373.416}{P}%
\TextAt{9}{452.215}{375.209}{\Character{61}}%
\TextAt{13}{464.641}{375.209}{K}%
\TextAt{27}{474.6}{373.416}{p}%
\TextAt{8}{479.087}{375.209}{.}%
\TextAt{8}{487.535}{375.209}{Since}%
\TextAt{13}{216.187}{348.324}{I}%
\TextAt{21}{221.35}{346.531}{K}%
\TextAt{9}{229.553}{348.324}{\Character{47}}%
\TextAt{13}{235.406}{348.324}{N}%
\TextAt{21}{244.779}{346.464}{M}%
\TextAt{3}{253.761}{346.464}{\Character{106}}%
\TextAt{21}{256.113}{346.464}{K}%
\TextAt{13}{264.316}{348.324}{I}%
\TextAt{21}{269.479}{346.531}{M}%
\NativeMathAt{282.28}{348.324}{11.955}{\cong}%
\TextAt{20}{294.899}{359.682}{\Character{77}}%
\TextAt{27}{301.938}{334.576}{p}%
\TextAt{13}{314.957}{348.324}{K}%
\TextAt{3}{325.775}{353.26}{\Character{3}}%
\TextAt{27}{324.916}{345.369}{p}%
\TextAt{9}{330.507}{348.324}{\Character{47}}%
\TextAt{13}{336.36}{348.324}{N}%
\TextAt{21}{345.734}{346.464}{M}%
\TextAt{28}{353.867}{345.073}{P}%
\TextAt{3}{360.064}{346.464}{\Character{106}}%
\TextAt{21}{362.416}{346.464}{K}%
\TextAt{28}{369.558}{345.468}{p}%
\TextAt{13}{374.138}{348.324}{M}%
\TextAt{3}{386.712}{353.26}{\Character{3}}%
\TextAt{27}{385.463}{345.369}{P}%
\TextAt{13}{392.557}{348.324}{\Character{59}}%
\TextAt{8}{79.226}{312.658}{the}%
\TextAt{8}{99.686}{312.658}{idèle}%
\TextAt{13}{127.294}{312.658}{\Character{11}a}%
\TextAt{3}{140.961}{316.996}{\Character{0}}%
\TextAt{17}{147.548}{316.996}{1}%
\TextAt{8}{156.48}{312.658}{is}%
\TextAt{8}{168.547}{312.658}{a}%
\TextAt{8}{178.606}{312.658}{norm}%
\TextAt{8}{209.478}{312.658}{of}%
\TextAt{8}{223.111}{312.658}{some}%
\TextAt{8}{252.741}{312.658}{idèle}%
\TextAt{13}{280.349}{312.658}{\Character{12}}%
\TextAt{8}{291.821}{312.658}{of}%
\TextAt{13}{305.454}{312.658}{I}%
\TextAt{21}{310.617}{310.865}{M}%
\TextAt{8}{320.097}{312.658}{,}%
\TextAt{8}{327.625}{312.658}{i.e.,}%
\TextAt{13}{350.109}{312.658}{\Character{11}}%
\TextAt{9}{361.468}{312.658}{\Character{61}}%
\TextAt{9}{374.392}{312.658}{\Character{40}}%
\TextAt{13}{378.962}{312.658}{N}%
\TextAt{21}{388.336}{310.798}{M}%
\TextAt{3}{397.318}{310.798}{\Character{106}}%
\TextAt{21}{399.67}{310.798}{K}%
\TextAt{13}{407.873}{312.658}{\Character{12}}%
\TextAt{9}{415.144}{312.658}{\Character{41}}%
\TextAt{13}{419.696}{312.658}{a}%
\NativeMathAt{429.679}{312.658}{11.955}{\in}%
\TextAt{13}{441.486}{312.658}{N}%
\TextAt{21}{450.86}{310.798}{M}%
\TextAt{3}{459.842}{310.798}{\Character{106}}%
\TextAt{21}{462.194}{310.798}{K}%
\TextAt{13}{470.397}{312.658}{I}%
\TextAt{21}{475.56}{310.865}{M}%
\TextAt{13}{485.039}{312.658}{K}%
\TextAt{3}{495.857}{316.996}{\Character{3}}%
\TextAt{8}{500.589}{312.658}{.}%
\TextAt{8}{509.947}{312.658}{This}%
\TextAt{8}{79.226}{295.612}{shows}%
\TextAt{8}{112.844}{295.612}{that}%
\TextAt{8}{138.213}{295.612}{the}%
\TextAt{8}{158.369}{295.612}{class}%
\TextAt{8}{185.806}{295.612}{of}%
\TextAt{13}{199.137}{295.612}{\Character{11}}%
\TextAt{8}{210.556}{295.612}{belongs}%
\TextAt{8}{252.566}{295.612}{to}%
\TextAt{13}{266.876}{295.612}{N}%
\TextAt{21}{276.25}{293.752}{M}%
\TextAt{3}{285.232}{293.752}{\Character{106}}%
\TextAt{21}{287.584}{293.752}{K}%
\TextAt{13}{295.787}{295.612}{C}%
\TextAt{21}{304.162}{293.819}{M}%
\TextAt{8}{313.642}{295.612}{,}%
\TextAt{8}{320.791}{295.612}{so}%
\TextAt{8}{335.161}{295.612}{that}%
\TextAt{13}{360.53}{295.612}{C}%
\TextAt{21}{368.905}{293.819}{K}%
\TextAt{9}{380.429}{295.612}{\Character{61}}%
\TextAt{13}{392.854}{295.612}{N}%
\TextAt{21}{402.228}{293.752}{M}%
\TextAt{3}{411.21}{293.752}{\Character{106}}%
\TextAt{21}{413.562}{293.752}{K}%
\TextAt{13}{421.765}{295.612}{C}%
\TextAt{21}{430.14}{293.819}{M}%
\TextAt{8}{439.619}{295.612}{.}%
\Vector{q 0 G 0 g 0.399 w 452.334 -478.388 m 452.334 -470.318 l S Q}%
\Vector{q 0 G 0 g 0.399 w 452.534 -470.518 m 459.707 -470.518 l S Q}%
\Vector{q 0 G 0 g 0.399 w 452.534 -478.189 m 459.707 -478.189 l S Q}%
\Vector{q 0 G 0 g 0.399 w 459.906 -478.388 m 459.906 -470.318 l S Q}%
\TextAt{12}{79.226}{262.468}{Corollary}%
\TextAt{12}{138.651}{262.468}{(3.8).}%
\TextAt{8}{173.651}{262.468}{Let}%
\TextAt{13}{194.304}{262.468}{L}%
\TextAt{16}{205.59}{262.468}{\Character{106}}%
\TextAt{13}{212.231}{262.468}{K}%
\TextAt{8}{226.631}{262.468}{be}%
\TextAt{8}{242.239}{262.468}{a}%
\TextAt{8}{251.679}{262.468}{finite}%
\TextAt{8}{281.275}{262.468}{extension}%
\TextAt{8}{332.726}{262.468}{of}%
\TextAt{8}{345.74}{262.468}{algebraic}%
\TextAt{8}{394.859}{262.468}{number}%
\TextAt{8}{437.139}{262.468}{fields.}%
\TextAt{8}{471.557}{262.468}{If}%
\TextAt{8}{482.933}{262.468}{almost}%
\TextAt{8}{520.408}{262.468}{all}%
\TextAt{8}{79.226}{245.422}{primes}%
\TextAt{8}{117.004}{245.422}{of}%
\TextAt{13}{130.334}{245.422}{K}%
\TextAt{8}{145.05}{245.422}{split}%
\TextAt{8}{171.124}{245.422}{completely}%
\TextAt{8}{229.979}{245.422}{in}%
\TextAt{13}{243.632}{245.422}{L}%
\TextAt{8}{251.596}{245.422}{,}%
\TextAt{8}{258.746}{245.422}{then}%
\TextAt{13}{285.406}{245.422}{L}%
\TextAt{9}{296.691}{245.422}{\Character{61}}%
\TextAt{13}{309.116}{245.422}{K}%
\TextAt{8}{319.934}{245.422}{.}%
\TextAt{19}{79.226}{212.278}{Proof.}%
\TextAt{8}{116}{212.278}{We}%
\TextAt{8}{135.857}{212.278}{may}%
\TextAt{8}{160.938}{212.278}{assume}%
\TextAt{8}{201.095}{212.278}{without}%
\TextAt{8}{244.384}{212.278}{loss}%
\TextAt{8}{266.333}{212.278}{of}%
\TextAt{8}{279.376}{212.278}{generality}%
\TextAt{8}{333.077}{212.278}{that}%
\TextAt{13}{358.159}{212.278}{L}%
\TextAt{16}{369.444}{212.278}{\Character{106}}%
\TextAt{13}{376.086}{212.278}{K}%
\TextAt{8}{390.513}{212.278}{is}%
\TextAt{8}{401.989}{212.278}{Galois.}%
\TextAt{8}{442.358}{212.278}{In}%
\TextAt{8}{456.691}{212.278}{fact,}%
\TextAt{8}{482.799}{212.278}{let}%
\TextAt{13}{499.416}{212.278}{M}%
\TextAt{16}{515.31}{212.278}{\Character{106}}%
\TextAt{13}{521.952}{212.278}{K}%
\TextAt{8}{79.226}{195.232}{be}%
\TextAt{8}{95.167}{195.232}{the}%
\TextAt{8}{115.34}{195.232}{normal}%
\TextAt{8}{155.035}{195.232}{closure}%
\TextAt{8}{194.133}{195.232}{of}%
\TextAt{13}{207.48}{195.232}{L}%
\TextAt{16}{218.793}{195.232}{\Character{106}}%
\TextAt{13}{225.463}{195.232}{K}%
\TextAt{8}{236.28}{195.232}{,}%
\TextAt{8}{243.45}{195.232}{and}%
\TextAt{8}{266.229}{195.232}{write}%
\TextAt{13}{296.157}{195.232}{G}%
\TextAt{9}{308.74}{195.232}{\Character{61}}%
\TextAt{13}{321.193}{195.232}{G}%
\TextAt{9}{330.427}{195.232}{\Character{40}}%
\TextAt{13}{334.979}{195.232}{M}%
\TextAt{16}{350.902}{195.232}{\Character{106}}%
\TextAt{13}{357.572}{195.232}{K}%
\TextAt{9}{368.389}{195.232}{\Character{41}}%
\TextAt{8}{376.855}{195.232}{and}%
\TextAt{13}{399.634}{195.232}{H}%
\TextAt{9}{413.606}{195.232}{\Character{61}}%
\TextAt{13}{426.059}{195.232}{G}%
\TextAt{9}{435.293}{195.232}{\Character{40}}%
\TextAt{13}{439.845}{195.232}{M}%
\TextAt{16}{455.768}{195.232}{\Character{106}}%
\TextAt{13}{462.437}{195.232}{L}%
\TextAt{9}{470.402}{195.232}{\Character{41}}%
\TextAt{8}{474.954}{195.232}{.}%
\TextAt{8}{483.452}{195.232}{Also}%
\TextAt{8}{509.865}{195.232}{let}%
\TextAt{14}{526.786}{195.232}{p}%
\TextAt{8}{79.226}{178.186}{be}%
\TextAt{8}{95.31}{178.186}{a}%
\TextAt{8}{105.225}{178.186}{place}%
\TextAt{8}{135.297}{178.186}{of}%
\TextAt{13}{148.786}{178.186}{K}%
\TextAt{8}{159.604}{178.186}{,}%
\TextAt{14}{166.953}{178.186}{P}%
\TextAt{8}{180.903}{178.186}{a}%
\TextAt{8}{190.818}{178.186}{place}%
\TextAt{8}{220.889}{178.186}{of}%
\TextAt{13}{234.379}{178.186}{M}%
\TextAt{8}{251.01}{178.186}{above}%
\TextAt{14}{284.345}{178.186}{p}%
\TextAt{8}{290.328}{178.186}{,}%
\TextAt{8}{297.677}{178.186}{and}%
\TextAt{8}{320.599}{178.186}{let}%
\TextAt{13}{337.663}{178.186}{G}%
\TextAt{27}{346.897}{176.393}{P}%
\TextAt{8}{358.047}{178.186}{be}%
\TextAt{8}{374.131}{178.186}{its}%
\TextAt{8}{390.61}{178.186}{decomposition}%
\TextAt{8}{467.904}{178.186}{group.}%
\TextAt{8}{506.109}{178.186}{Then}%
\TextAt{8}{79.226}{161.14}{the}%
\TextAt{8}{99.098}{161.14}{number}%
\TextAt{8}{141.411}{161.14}{of}%
\TextAt{8}{154.456}{161.14}{places}%
\TextAt{8}{188.699}{161.14}{of}%
\TextAt{13}{201.745}{161.14}{L}%
\TextAt{8}{213.323}{161.14}{above}%
\TextAt{14}{246.214}{161.14}{p}%
\TextAt{8}{255.81}{161.14}{equals}%
\TextAt{8}{291.033}{161.14}{the}%
\TextAt{8}{310.905}{161.14}{number}%
\TextAt{8}{353.218}{161.14}{\char35\relax }%
\TextAt{13}{362.973}{161.14}{H}%
\TextAt{16}{373.596}{161.14}{\Character{110}}%
\TextAt{13}{379.573}{161.14}{G}%
\TextAt{9}{388.807}{161.14}{\Character{47}}%
\TextAt{13}{394.66}{161.14}{G}%
\TextAt{27}{403.893}{159.347}{P}%
\TextAt{8}{414.6}{161.14}{of}%
\TextAt{8}{427.646}{161.14}{double}%
\TextAt{8}{465.081}{161.14}{cosets}%
\TextAt{13}{498.737}{161.14}{H}%
\TextAt{13}{509.353}{161.14}{\Character{27}}%
\TextAt{13}{516.431}{161.14}{G}%
\TextAt{27}{525.676}{159.347}{P}%
\TextAt{8}{79.226}{144.094}{in}%
\TextAt{13}{92.989}{144.094}{G}%
\TextAt{8}{106.229}{144.094}{(see}%
\TextAt{8}{129.807}{144.094}{chap.}%
\TextAt{8}{160.809}{144.094}{I,}%
\TextAt{8}{172.288}{144.094}{§9).}%
\ErrataAt{197.15}{144.094}{337.85}{Complete splitting in $L$ makes $G_{\mathfrak P}$ act trivially on $H\backslash G$.}%
\ErrataAt{79.19}{127.048}{455.81}{Since $M$ is the normal closure, $\bigcap_{g\in G}g^{-1}Hg=1$, so $G_{\mathfrak P}=1$ and $\mathfrak p$ splits completely in $M$.}%
\TextAt{8}{97.159}{107.013}{So}%
\TextAt{8}{114.276}{107.013}{assume}%
\TextAt{13}{155.579}{107.013}{L}%
\TextAt{16}{168.328}{107.013}{\Character{106}}%
\TextAt{13}{176.433}{107.013}{K}%
\TextAt{8}{192.006}{107.013}{is}%
\TextAt{8}{204.629}{107.013}{Galois,}%
\TextAt{13}{244.867}{107.013}{L}%
\NativeMathAt{257.615}{107.013}{11.955}{\neq}%
\TextAt{13}{271.504}{107.013}{K}%
\TextAt{8}{282.322}{107.013}{,}%
\TextAt{8}{290.544}{107.013}{and}%
\TextAt{8}{314.165}{107.013}{let}%
\TextAt{13}{331.928}{107.013}{\Character{27}}%
\NativeMathAt{343.795}{107.013}{11.955}{\in}%
\TextAt{13}{356.549}{107.013}{G}%
\TextAt{9}{365.782}{107.013}{\Character{40}}%
\TextAt{13}{370.335}{107.013}{L}%
\TextAt{16}{383.083}{107.013}{\Character{106}}%
\TextAt{13}{391.188}{107.013}{K}%
\TextAt{9}{402.006}{107.013}{\Character{41}}%
\TextAt{8}{411.314}{107.013}{be}%
\TextAt{8}{428.097}{107.013}{an}%
\TextAt{8}{445.214}{107.013}{element}%
\TextAt{8}{489.315}{107.013}{of}%
\TextAt{8}{503.503}{107.013}{prime}%
\TextAt{8}{79.226}{89.967}{order,}%
\TextAt{8}{113.847}{89.967}{with}%
\TextAt{8}{141.146}{89.967}{fixed}%
\TextAt{8}{170.072}{89.967}{field}%
\TextAt{13}{196.068}{89.967}{K}%
\TextAt{3}{206.886}{94.306}{\Character{48}}%
\TextAt{8}{209.681}{89.967}{.}%
\TextAt{8}{220.048}{89.967}{If}%
\TextAt{8}{232.38}{89.967}{almost}%
\TextAt{8}{270.81}{89.967}{all}%
\TextAt{8}{287.708}{89.967}{primes}%
\TextAt{14}{326.126}{89.967}{p}%
\TextAt{8}{336.646}{89.967}{of}%
\TextAt{13}{350.616}{89.967}{K}%
\TextAt{8}{365.97}{89.967}{were}%
\TextAt{8}{393.593}{89.967}{completely}%
\TextAt{8}{453.088}{89.967}{split}%
\TextAt{8}{479.801}{89.967}{in}%
\TextAt{13}{494.094}{89.967}{L}%
\TextAt{8}{502.058}{89.967}{,}%
\TextAt{8}{510.007}{89.967}{then}%
\TextAt{8}{79.226}{72.921}{the}%
\TextAt{8}{99.383}{72.921}{same}%
\TextAt{8}{128.709}{72.921}{would}%
\TextAt{8}{162.853}{72.921}{hold}%
\TextAt{8}{188.867}{72.921}{for}%
\TextAt{8}{206.752}{72.921}{the}%
\TextAt{8}{226.908}{72.921}{primes}%
\TextAt{14}{264.687}{72.921}{p}%
\TextAt{3}{270.67}{77.26}{\Character{48}}%
\TextAt{8}{277.362}{72.921}{of}%
\TextAt{13}{290.692}{72.921}{K}%
\TextAt{3}{301.51}{77.26}{\Character{48}}%
\TextAt{8}{304.305}{72.921}{.}%
\TextAt{8}{312.753}{72.921}{This}%
\TextAt{8}{339.473}{72.921}{contradicts}%
\TextAt{8}{399.703}{72.921}{(3.7).}%
\Vector{q 0 G 0 g 0.399 w 452.334 -701.079 m 452.334 -693.009 l S Q}%
\Vector{q 0 G 0 g 0.399 w 452.534 -693.208 m 459.707 -693.208 l S Q}%
\Vector{q 0 G 0 g 0.399 w 452.534 -700.879 m 459.707 -700.879 l S Q}%
\Vector{q 0 G 0 g 0.399 w 459.906 -701.079 m 459.906 -693.009 l S Q}%
\TextAt{8}{297.211}{43.033}{345}%
\Destination{78.23}{809.815}{nav.page.345}%
\Destination{79.36}{776.155}{nav.print.378}%
\Destination{79.36}{424.255}{nav.print.379}%
\Destination{79.36}{659.655}{nav.number.6.3.6}%
\Destination{79.36}{544.655}{nav.number.6.3.7}%
\Destination{79.36}{276.455}{nav.number.6.3.8}%
\Hotlink{409.325}{407.055}{9.751}{11.047}{nav.chapter.2}%
\Hotlink{160.309}{141.855}{5.22}{11.047}{nav.chapter.1}%
\Hotlink{177.467}{141.855}{6.858}{11.047}{nav.section.1.9}%
\Hotlink{472.919}{724.855}{25.078}{11.047}{nav.number.6.1.4}%
\Hotlink{386.862}{668.555}{25.078}{11.047}{nav.number.6.3.3}%
\Hotlink{436.312}{668.555}{25.078}{11.047}{nav.number.6.3.5}%
\Hotlink{431.606}{427.155}{25.078}{11.047}{nav.number.6.3.6}%
\Hotlink{399.203}{70.755}{25.078}{11.047}{nav.number.6.3.7}%
\Destination{79.226}{161.094}{errata.143}%
\Hotlink{427.187}{406.307}{21.722}{11.047}{nav.number.2.3.4}%
}

% Page 358 (body page 346)
\Page{%
\TextAt{2}{79.226}{796.943}{VI.}%
\TextAt{2}{98.012}{796.943}{Global}%
\TextAt{2}{133.248}{796.943}{Class}%
\TextAt{2}{161.841}{796.943}{Field}%
\TextAt{2}{189.572}{796.943}{Theory}%
\Vector{q 0 G 0 g 0.299 w 7.226 18.729 m 460.769 18.729 l S Q}%
\TextAt{12}{79.226}{761.124}{Exercise}%
\TextAt{12}{133.798}{761.124}{VI.3.1.}%
\TextAt{8}{178.746}{761.124}{If}%
\TextAt{8}{191.63}{761.124}{the}%
\TextAt{8}{212.978}{761.124}{Galois}%
\TextAt{8}{250.084}{761.124}{extension}%
\TextAt{13}{303.042}{761.124}{L}%
\TextAt{16}{316.358}{761.124}{\Character{106}}%
\TextAt{13}{325.031}{761.124}{K}%
\TextAt{8}{340.938}{761.124}{is}%
\TextAt{8}{353.894}{761.124}{not}%
\TextAt{8}{375.9}{761.124}{cyclic,}%
\TextAt{8}{412.825}{761.124}{then}%
\TextAt{8}{440.677}{761.124}{there}%
\TextAt{8}{471.781}{761.124}{are}%
\TextAt{8}{492.484}{761.124}{at}%
\TextAt{8}{507.986}{761.124}{most}%
\TextAt{8}{79.226}{744.078}{finitely}%
\TextAt{8}{118.571}{744.078}{many}%
\TextAt{8}{150.443}{744.078}{primes}%
\TextAt{8}{188.221}{744.078}{of}%
\TextAt{13}{201.551}{744.078}{K}%
\TextAt{8}{216.267}{744.078}{which}%
\TextAt{8}{249.753}{744.078}{do}%
\TextAt{8}{266.012}{744.078}{not}%
\TextAt{8}{286.826}{744.078}{split}%
\TextAt{8}{312.9}{744.078}{in}%
\TextAt{13}{326.553}{744.078}{L}%
\TextAt{8}{334.517}{744.078}{.}%
\TextAt{12}{79.226}{712.088}{Exercise}%
\TextAt{12}{132.819}{712.088}{VI.3.2.}%
\TextAt{8}{176.917}{712.088}{If}%
\TextAt{13}{188.95}{712.088}{L}%
\TextAt{16}{200.817}{712.088}{\Character{106}}%
\TextAt{13}{208.041}{712.088}{K}%
\TextAt{8}{223.097}{712.088}{is}%
\TextAt{8}{235.203}{712.088}{a}%
\TextAt{8}{245.299}{712.088}{finite}%
\TextAt{8}{275.553}{712.088}{Galois}%
\TextAt{8}{311.807}{712.088}{extension,}%
\TextAt{8}{367.252}{712.088}{then}%
\TextAt{8}{394.253}{712.088}{the}%
\TextAt{8}{414.751}{712.088}{Galois}%
\TextAt{8}{451.006}{712.088}{group}%
\TextAt{13}{484.523}{712.088}{G}%
\TextAt{9}{493.757}{712.088}{\Character{40}}%
\TextAt{13}{498.309}{712.088}{L}%
\TextAt{16}{510.176}{712.088}{\Character{106}}%
\TextAt{13}{517.399}{712.088}{K}%
\TextAt{9}{528.217}{712.088}{\Character{41}}%
\TextAt{8}{79.226}{695.042}{is}%
\TextAt{8}{91.924}{695.042}{generated}%
\TextAt{8}{146.189}{695.042}{by}%
\TextAt{8}{163.382}{695.042}{the}%
\TextAt{8}{184.472}{695.042}{Frobenius}%
\TextAt{8}{239.276}{695.042}{automorphisms}%
\TextAt{13}{322.294}{695.042}{\Character{39}}%
\TextAt{27}{329.962}{693.249}{P}%
\TextAt{8}{341.887}{695.042}{of}%
\TextAt{8}{356.151}{695.042}{all}%
\TextAt{8}{373.343}{695.042}{prime}%
\TextAt{8}{407.441}{695.042}{ideals}%
\TextAt{14}{440.952}{695.042}{P}%
\TextAt{8}{455.677}{695.042}{of}%
\TextAt{13}{469.94}{695.042}{L}%
\TextAt{8}{482.736}{695.042}{which}%
\TextAt{8}{517.156}{695.042}{are}%
\TextAt{8}{79.226}{677.996}{unramified}%
\TextAt{8}{137.759}{677.996}{over}%
\TextAt{13}{162.805}{677.996}{K}%
\TextAt{8}{173.622}{677.996}{.}%
\TextAt{12}{79.226}{646.006}{Exercise}%
\TextAt{12}{132.061}{646.006}{VI.3.3.}%
\TextAt{8}{175.5}{646.006}{Let}%
\TextAt{13}{196.152}{646.006}{L}%
\TextAt{16}{207.437}{646.006}{\Character{106}}%
\TextAt{13}{214.079}{646.006}{K}%
\TextAt{8}{228.477}{646.006}{be}%
\TextAt{8}{244.083}{646.006}{a}%
\TextAt{8}{253.521}{646.006}{finite}%
\TextAt{8}{283.116}{646.006}{abelian}%
\TextAt{8}{323.446}{646.006}{extension,}%
\TextAt{8}{378.21}{646.006}{and}%
\TextAt{8}{400.655}{646.006}{let}%
\TextAt{13}{417.242}{646.006}{D}%
\TextAt{8}{430.866}{646.006}{be}%
\TextAt{8}{446.473}{646.006}{a}%
\TextAt{8}{455.911}{646.006}{subgroup}%
\TextAt{8}{506.391}{646.006}{of}%
\TextAt{13}{519.404}{646.006}{I}%
\TextAt{21}{524.567}{644.213}{K}%
\ErrataAt{79.22}{628.955}{453.64}{such that $K^*D$ is dense in $I_K$ and $D\subseteq N_{L\mid K}I_L$. Then $L=K$.}%
\TextAt{12}{79.226}{596.97}{Exercise}%
\TextAt{12}{133.496}{596.97}{VI.3.4.}%
\TextAt{8}{178.182}{596.97}{Let}%
\TextAt{13}{200.081}{596.97}{L}%
\TextAt{17}{208.045}{595.177}{1}%
\TextAt{13}{212.778}{596.97}{\Character{59}}%
\TextAt{13}{218.002}{596.97}{\Character{58}}%
\TextAt{13}{223.239}{596.97}{\Character{58}}%
\TextAt{13}{228.463}{596.97}{\Character{58}}%
\TextAt{13}{233.7}{596.97}{\Character{59}}%
\TextAt{13}{238.936}{596.97}{L}%
\TextAt{21}{246.963}{595.177}{r}%
\TextAt{16}{256.423}{596.97}{\Character{106}}%
\TextAt{13}{264.649}{596.97}{K}%
\TextAt{8}{280.294}{596.97}{be}%
\TextAt{8}{297.148}{596.97}{cyclic}%
\TextAt{8}{330.261}{596.97}{extensions}%
\TextAt{8}{387.571}{596.97}{of}%
\TextAt{8}{401.831}{596.97}{prime}%
\TextAt{8}{435.924}{596.97}{degree}%
\TextAt{13}{473.269}{596.97}{p}%
\TextAt{8}{483.971}{596.97}{such}%
\TextAt{8}{511.298}{596.97}{that}%
\ErrataAt{79.22}{579.955}{453.64}{$L_1\not\subseteq L_2\cdots L_r$. Then there are infinitely many primes of $K$ which split completely}%
\TextAt{8}{79.226}{562.878}{in}%
\TextAt{13}{92.879}{562.878}{L}%
\TextAt{21}{100.843}{561.085}{i}%
\TextAt{8}{104.225}{562.878}{,}%
\TextAt{8}{111.374}{562.878}{for}%
\TextAt{13}{129.259}{562.878}{i}%
\TextAt{13}{136.564}{562.878}{\Character{62}}%
\TextAt{9}{148.999}{562.878}{1}%
\TextAt{8}{154.852}{562.878}{,}%
\TextAt{8}{162.001}{562.878}{but}%
\TextAt{8}{183.46}{562.878}{which}%
\TextAt{8}{216.947}{562.878}{are}%
\TextAt{8}{236.458}{562.878}{nonsplit}%
\TextAt{8}{281.397}{562.878}{in}%
\TextAt{13}{295.05}{562.878}{L}%
\TextAt{17}{303.014}{561.085}{1}%
\TextAt{8}{307.747}{562.878}{.}%
\TextAt{15}{79.226}{510.014}{§}%
\TextAt{15}{93.785}{510.014}{4}%
\TextAt{15}{120.693}{510.014}{The}%
\TextAt{15}{160.202}{510.014}{Class}%
\TextAt{15}{210.747}{510.014}{Field}%
\TextAt{15}{259.742}{510.014}{Axiom}%
\TextAt{8}{79.226}{478.129}{Having}%
\TextAt{8}{120.734}{478.129}{determined}%
\TextAt{8}{183.366}{478.129}{the}%
\TextAt{8}{205.028}{478.129}{Herbrand}%
\TextAt{8}{258.885}{478.129}{quotient}%
\TextAt{13}{306.574}{478.129}{h}%
\TextAt{9}{313.312}{478.129}{\Character{40}}%
\TextAt{13}{317.865}{478.129}{G\Character{59}}%
\TextAt{13}{332.319}{478.129}{C}%
\TextAt{21}{340.717}{476.335}{L}%
\TextAt{9}{346.975}{478.129}{\Character{41}}%
\TextAt{8}{356.93}{478.129}{to}%
\TextAt{8}{372.746}{478.129}{be}%
\TextAt{8}{390.176}{478.129}{the}%
\TextAt{8}{411.837}{478.129}{degree}%
\TextAt{13}{449.758}{478.129}{n}%
\TextAt{9}{462.633}{478.129}{\Character{61}}%
\TextAt{9}{477.613}{478.129}{\Character{91}}%
\TextAt{13}{480.875}{478.129}{L}%
\TextAt{9}{494.726}{478.129}{\Character{58}}%
\TextAt{13}{503.865}{478.129}{K}%
\TextAt{9}{514.682}{478.129}{\Character{93}}%
\TextAt{8}{523.337}{478.129}{of}%
\TextAt{8}{79.226}{461.083}{the}%
\TextAt{8}{100.706}{461.083}{cyclic}%
\TextAt{8}{134.212}{461.083}{extension}%
\TextAt{13}{187.302}{461.083}{L}%
\TextAt{16}{200.842}{461.083}{\Character{106}}%
\TextAt{13}{209.739}{461.083}{K}%
\TextAt{8}{220.557}{461.083}{,}%
\TextAt{8}{229.36}{461.083}{it}%
\TextAt{8}{242.387}{461.083}{will}%
\TextAt{8}{265.816}{461.083}{now}%
\TextAt{8}{291.528}{461.083}{be}%
\TextAt{8}{308.775}{461.083}{enough}%
\TextAt{8}{350.423}{461.083}{to}%
\TextAt{8}{366.057}{461.083}{show}%
\TextAt{8}{396.383}{461.083}{either}%
\TextAt{13}{430.87}{461.083}{H}%
\TextAt{3}{441.493}{465.421}{\Character{0}}%
\TextAt{17}{448.079}{465.421}{1}%
\TextAt{9}{452.812}{461.083}{\Character{40}}%
\TextAt{13}{457.364}{461.083}{G\Character{59}}%
\TextAt{13}{471.818}{461.083}{C}%
\TextAt{21}{480.217}{459.289}{L}%
\TextAt{9}{486.474}{461.083}{\Character{41}}%
\TextAt{9}{496.588}{461.083}{\Character{61}}%
\TextAt{9}{511.257}{461.083}{1}%
\TextAt{8}{522.357}{461.083}{or}%
\TextAt{13}{79.226}{444.037}{H}%
\TextAt{17}{89.849}{448.375}{0}%
\TextAt{9}{94.581}{444.037}{\Character{40}}%
\TextAt{13}{99.133}{444.037}{G\Character{59}}%
\TextAt{13}{113.587}{444.037}{C}%
\TextAt{21}{121.986}{442.243}{L}%
\TextAt{9}{128.243}{444.037}{\Character{41}}%
\TextAt{9}{137.078}{444.037}{\Character{61}}%
\TextAt{9}{150.468}{444.037}{\Character{40}}%
\TextAt{13}{155.046}{444.037}{C}%
\TextAt{21}{163.421}{442.243}{K}%
\TextAt{9}{175.92}{444.037}{\Character{58}}%
\TextAt{13}{183.468}{444.037}{N}%
\TextAt{21}{192.842}{442.177}{L}%
\TextAt{3}{198.601}{442.177}{\Character{106}}%
\TextAt{21}{200.953}{442.177}{K}%
\TextAt{13}{209.156}{444.037}{C}%
\TextAt{21}{217.531}{442.243}{L}%
\TextAt{9}{223.788}{444.037}{\Character{41}}%
\TextAt{9}{232.623}{444.037}{\Character{61}}%
\TextAt{13}{246.038}{444.037}{n}%
\TextAt{8}{253.025}{444.037}{.}%
\TextAt{8}{263.191}{444.037}{The}%
\TextAt{8}{287.818}{444.037}{first}%
\TextAt{8}{312.516}{444.037}{identity}%
\TextAt{8}{356.342}{444.037}{is}%
\TextAt{8}{368.678}{444.037}{curiously}%
\TextAt{8}{419.068}{444.037}{inaccessible}%
\TextAt{8}{482.19}{444.037}{by}%
\TextAt{8}{499.021}{444.037}{way}%
\TextAt{8}{523.337}{444.037}{of}%
\TextAt{8}{79.226}{426.991}{direct}%
\TextAt{8}{112.017}{426.991}{attack.}%
\TextAt{8}{152.226}{426.991}{We}%
\TextAt{8}{171.997}{426.991}{are}%
\TextAt{8}{191.135}{426.991}{thus}%
\TextAt{8}{216.514}{426.991}{stuck}%
\TextAt{8}{246.77}{426.991}{with}%
\TextAt{8}{273.058}{426.991}{the}%
\TextAt{8}{292.841}{426.991}{second.}%
\TextAt{8}{335.046}{426.991}{We}%
\TextAt{8}{354.818}{426.991}{will}%
\TextAt{8}{376.55}{426.991}{reduce}%
\TextAt{8}{413.238}{426.991}{the}%
\TextAt{8}{433.022}{426.991}{problem}%
\TextAt{8}{478.175}{426.991}{to}%
\TextAt{8}{492.112}{426.991}{the}%
\TextAt{8}{511.896}{426.991}{case}%
\TextAt{8}{79.226}{409.945}{of}%
\TextAt{8}{93.068}{409.945}{a}%
\TextAt{19}{103.335}{409.945}{Kummer}%
\TextAt{19}{152.183}{409.945}{extension}%
\TextAt{8}{199.704}{409.945}{.}%
\TextAt{8}{209.687}{409.945}{For}%
\TextAt{8}{231.168}{409.945}{such}%
\TextAt{8}{258.077}{409.945}{an}%
\TextAt{8}{274.847}{409.945}{extension}%
\TextAt{8}{327.125}{409.945}{the}%
\TextAt{8}{347.793}{409.945}{norm}%
\TextAt{8}{378.874}{409.945}{group}%
\TextAt{13}{412.561}{409.945}{N}%
\TextAt{21}{421.934}{408.085}{L}%
\TextAt{3}{427.694}{408.085}{\Character{106}}%
\TextAt{21}{430.046}{408.085}{K}%
\TextAt{13}{438.249}{409.945}{C}%
\TextAt{21}{446.624}{408.151}{L}%
\TextAt{8}{457.29}{409.945}{can}%
\TextAt{8}{479.261}{409.945}{be}%
\TextAt{8}{495.696}{409.945}{written}%
\TextAt{8}{79.226}{392.899}{down}%
\TextAt{8}{110.118}{392.899}{explicitly,}%
\TextAt{8}{163.116}{392.899}{and}%
\TextAt{8}{185.878}{392.899}{this}%
\TextAt{8}{208.701}{392.899}{allows}%
\TextAt{8}{243.562}{392.899}{us}%
\TextAt{8}{258.578}{392.899}{to}%
\TextAt{8}{272.888}{392.899}{compute}%
\TextAt{8}{320.362}{392.899}{the}%
\TextAt{8}{340.518}{392.899}{index}%
\TextAt{9}{372.056}{392.899}{\Character{40}}%
\TextAt{13}{376.608}{392.899}{C}%
\TextAt{21}{384.983}{391.105}{K}%
\TextAt{9}{396.507}{392.899}{\Character{58}}%
\TextAt{13}{403.079}{392.899}{N}%
\TextAt{21}{412.453}{391.039}{L}%
\TextAt{3}{418.212}{391.039}{\Character{106}}%
\TextAt{21}{420.565}{391.039}{K}%
\TextAt{13}{428.767}{392.899}{C}%
\TextAt{21}{437.142}{391.105}{L}%
\TextAt{9}{443.4}{392.899}{\Character{41}}%
\TextAt{8}{447.952}{392.899}{.}%
\TextAt{8}{97.159}{372.864}{So}%
\TextAt{8}{113.987}{372.864}{let}%
\TextAt{13}{131.46}{372.864}{K}%
\TextAt{8}{146.744}{372.864}{be}%
\TextAt{8}{163.237}{372.864}{a}%
\TextAt{8}{173.561}{372.864}{number}%
\TextAt{8}{216.726}{372.864}{field}%
\TextAt{8}{242.652}{372.864}{that}%
\TextAt{8}{268.589}{372.864}{contains}%
\TextAt{8}{315.078}{372.864}{the}%
\TextAt{13}{335.803}{372.864}{n}%
\TextAt{8}{342.79}{372.864}{-th}%
\TextAt{8}{362.213}{372.864}{roots}%
\TextAt{8}{392.442}{372.864}{of}%
\TextAt{8}{406.341}{372.864}{unity,}%
\TextAt{8}{439.893}{372.864}{where}%
\TextAt{13}{474.27}{372.864}{n}%
\TextAt{8}{485.724}{372.864}{is}%
\TextAt{8}{498.057}{372.864}{a}%
\TextAt{8}{508.381}{372.864}{fixed}%
\TextAt{8}{79.226}{355.818}{prime}%
\TextAt{8}{112.39}{355.818}{power,}%
\TextAt{8}{149.786}{355.818}{and}%
\TextAt{8}{172.548}{355.818}{let}%
\TextAt{13}{189.453}{355.818}{L}%
\TextAt{16}{200.738}{355.818}{\Character{106}}%
\TextAt{13}{207.38}{355.818}{K}%
\TextAt{8}{222.095}{355.818}{be}%
\TextAt{8}{238.019}{355.818}{a}%
\TextAt{8}{247.775}{355.818}{Galois}%
\TextAt{8}{283.688}{355.818}{extension}%
\TextAt{8}{335.454}{355.818}{with}%
\TextAt{8}{362.114}{355.818}{a}%
\TextAt{8}{371.869}{355.818}{Galois}%
\TextAt{8}{407.783}{355.818}{group}%
\TextAt{8}{440.958}{355.818}{of}%
\TextAt{8}{454.288}{355.818}{the}%
\TextAt{8}{474.445}{355.818}{form}%
\TextAt{13}{251.64}{326.817}{G}%
\TextAt{9}{260.874}{326.817}{\Character{40}}%
\TextAt{13}{265.426}{326.817}{L}%
\TextAt{16}{276.712}{326.817}{\Character{106}}%
\TextAt{13}{283.353}{326.817}{K}%
\TextAt{9}{294.171}{326.817}{\Character{41}}%
\NativeMathAt{302.044}{326.817}{11.955}{\cong}%
\TextAt{9}{314.663}{326.817}{\Character{40}}%
\TextAt{18}{319.216}{326.817}{Z}%
\TextAt{9}{327.186}{326.817}{\Character{47}}%
\TextAt{13}{333.039}{326.817}{n}%
\TextAt{18}{340.026}{326.817}{Z}%
\TextAt{9}{347.997}{326.817}{\Character{41}}%
\TextAt{21}{352.549}{331.753}{r}%
\TextAt{13}{357.104}{326.817}{\Character{58}}%
\TextAt{8}{79.226}{297.816}{We}%
\TextAt{8}{99.019}{297.816}{choose}%
\TextAt{8}{135.8}{297.816}{a}%
\TextAt{8}{145.204}{297.816}{finite}%
\TextAt{8}{174.764}{297.816}{set}%
\TextAt{8}{192.68}{297.816}{of}%
\TextAt{8}{205.658}{297.816}{places}%
\TextAt{13}{239.833}{297.816}{S}%
\TextAt{8}{251.274}{297.816}{containing}%
\TextAt{8}{307.841}{297.816}{the}%
\TextAt{8}{327.646}{297.816}{ramified}%
\TextAt{8}{372.819}{297.816}{places,}%
\TextAt{8}{410.316}{297.816}{those}%
\TextAt{8}{440.594}{297.816}{that}%
\TextAt{8}{465.611}{297.816}{divide}%
\TextAt{13}{500.049}{297.816}{n}%
\TextAt{8}{507.036}{297.816}{,}%
\TextAt{8}{513.904}{297.816}{and}%
\TextAt{8}{79.226}{280.77}{the}%
\TextAt{8}{99.558}{280.77}{infinite}%
\TextAt{8}{139.401}{280.77}{ones,}%
\TextAt{8}{168.947}{280.77}{and}%
\TextAt{8}{191.885}{280.77}{which}%
\TextAt{8}{225.547}{280.77}{is}%
\TextAt{8}{237.486}{280.77}{such}%
\TextAt{8}{264.059}{280.77}{that}%
\TextAt{13}{289.603}{280.77}{I}%
\TextAt{21}{294.766}{278.976}{K}%
\TextAt{9}{306.589}{280.77}{\Character{61}}%
\TextAt{13}{319.314}{280.77}{I}%
\TextAt{21}{325.416}{285.108}{S}%
\TextAt{21}{324.477}{277.749}{K}%
\TextAt{13}{332.679}{280.77}{K}%
\TextAt{3}{343.497}{285.108}{\Character{3}}%
\TextAt{8}{348.229}{280.77}{.}%
\TextAt{8}{357.204}{280.77}{We}%
\TextAt{8}{377.524}{280.77}{write}%
\TextAt{8}{407.612}{280.77}{again}%
\TextAt{13}{439.014}{280.77}{K}%
\TextAt{21}{449.832}{285.108}{S}%
\TextAt{9}{459.546}{280.77}{\Character{61}}%
\TextAt{13}{472.271}{280.77}{I}%
\TextAt{21}{478.373}{285.108}{S}%
\TextAt{21}{477.434}{277.749}{K}%
\TextAt{16}{488.413}{280.77}{\Character{92}}%
\TextAt{13}{499.159}{280.77}{K}%
\TextAt{3}{509.977}{285.108}{\Character{3}}%
\TextAt{8}{518.782}{280.77}{for}%
\TextAt{8}{79.226}{263.724}{the}%
\TextAt{8}{99.383}{263.724}{group}%
\TextAt{8}{132.558}{263.724}{of}%
\TextAt{13}{145.888}{263.724}{S}%
\TextAt{8}{153.783}{263.724}{-units,}%
\TextAt{8}{190.259}{263.724}{and}%
\TextAt{8}{213.021}{263.724}{we}%
\TextAt{8}{230.249}{263.724}{put}%
\TextAt{13}{251.708}{263.724}{s}%
\TextAt{9}{260.543}{263.724}{\Character{61}}%
\TextAt{8}{272.969}{263.724}{\char35\relax }%
\TextAt{13}{282.724}{263.724}{S}%
\TextAt{8}{290.619}{263.724}{.}%
\TextAt{12}{79.226}{237.711}{Proposition}%
\TextAt{12}{151.6}{237.711}{(4.1).}%
\TextAt{8}{186.622}{237.711}{One}%
\TextAt{8}{211.04}{237.711}{has}%
\TextAt{13}{231.62}{237.711}{s}%
\NativeMathAt{240.455}{237.711}{11.955}{\geq}%
\TextAt{13}{253.074}{237.711}{r}%
\TextAt{8}{258.674}{237.711}{,}%
\TextAt{8}{265.589}{237.711}{and}%
\TextAt{8}{288.058}{237.711}{there}%
\TextAt{8}{317.676}{237.711}{exists}%
\TextAt{8}{349.697}{237.711}{a}%
\TextAt{8}{359.159}{237.711}{set}%
\TextAt{13}{377.133}{237.711}{T}%
\TextAt{8}{389.224}{237.711}{of}%
\TextAt{13}{402.26}{237.711}{s}%
\TextAt{16}{409.831}{237.711}{\Character{0}}%
\TextAt{13}{421.185}{237.711}{r}%
\TextAt{8}{430.39}{237.711}{primes}%
\TextAt{8}{467.875}{237.711}{of}%
\TextAt{13}{480.911}{237.711}{K}%
\TextAt{8}{495.333}{237.711}{that}%
\TextAt{8}{520.408}{237.711}{do}%
\TextAt{8}{79.226}{220.665}{not}%
\TextAt{8}{100.04}{220.665}{belong}%
\TextAt{8}{137.436}{220.665}{to}%
\TextAt{13}{151.746}{220.665}{S}%
\TextAt{8}{163.539}{220.665}{such}%
\TextAt{8}{189.936}{220.665}{that}%
\TextAt{13}{272.057}{203.619}{L}%
\TextAt{9}{283.342}{203.619}{\Character{61}}%
\TextAt{13}{295.767}{203.619}{K}%
\TextAt{9}{306.585}{203.619}{\Character{40}}%
\TextAt{22}{314.458}{210.154}{n}%
\TextAt{16}{312.417}{214.325}{\Character{112}}%
\Vector{q 0 G 0 g 0.478 w 250.38 -559.436 m 260.135 -559.436 l S Q}%
\TextAt{9}{322.38}{203.619}{\Character{1}\Character{41}}%
\TextAt{13}{336.687}{203.619}{\Character{59}}%
\TextAt{8}{79.226}{180.098}{where}%
\TextAt{9}{113.035}{180.098}{\Character{1}}%
\TextAt{8}{126.688}{180.098}{is}%
\TextAt{8}{138.452}{180.098}{the}%
\TextAt{8}{158.608}{180.098}{kernel}%
\TextAt{8}{193.075}{180.098}{of}%
\TextAt{8}{206.405}{180.098}{the}%
\TextAt{8}{226.561}{180.098}{map}%
\TextAt{13}{252.576}{180.098}{K}%
\TextAt{21}{263.393}{184.436}{S}%
\NativeArrowAt{272.808}{180.098}{11.955}{19.26}{\longrightarrow}%
\TextAt{20}{295.39}{189.064}{\Character{81}}%
\TextAt{27}{306.681}{176.611}{p}%
\NativeMathAt{310.67}{176.611}{7.97}{\in}%
\TextAt{21}{316.315}{176.611}{T}%
\TextAt{13}{324.912}{180.098}{K}%
\TextAt{3}{335.73}{184.436}{\Character{3}}%
\TextAt{27}{334.872}{177.142}{p}%
\TextAt{9}{340.462}{180.098}{\Character{47}}%
\TextAt{13}{346.315}{180.098}{K}%
\TextAt{3}{357.133}{184.436}{\Character{3}}%
\TextAt{21}{361.367}{184.436}{n}%
\TextAt{27}{356.274}{177.142}{p}%
\TextAt{8}{367.004}{180.098}{.}%
\TextAt{8}{297.211}{43.033}{346}%
\Destination{78.23}{809.815}{nav.page.346}%
\Destination{79.23}{536.545}{nav.section.6.4}%
\Destination{79.226}{775.155}{nav.ex.6.3.1}%
\Destination{79.226}{726.055}{nav.ex.6.3.2}%
\Destination{79.226}{660.055}{nav.ex.6.3.3}%
\Destination{79.226}{610.955}{nav.ex.6.3.4}%
\Destination{79.36}{493.155}{nav.print.380}%
\Destination{79.36}{251.755}{nav.number.6.4.1}%
\Destination{79.23}{536.545}{outline.59}%
\Destination{79.226}{642.955}{errata.28}%
\Destination{79.226}{593.955}{errata.29}%
}

% Page 359 (body page 347)
\Page{%
\TextAt{2}{79.226}{796.943}{VI.}%
\TextAt{2}{98.012}{796.943}{Global}%
\TextAt{2}{133.248}{796.943}{Class}%
\TextAt{2}{161.841}{796.943}{Field}%
\TextAt{2}{189.572}{796.943}{Theory}%
\Vector{q 0 G 0 g 0.299 w 7.226 18.729 m 460.769 18.729 l S Q}%
\TextAt{19}{79.226}{761.124}{Proof.}%
\TextAt{8}{116}{761.124}{We}%
\TextAt{8}{136.966}{761.124}{show}%
\TextAt{8}{166.791}{761.124}{first}%
\TextAt{8}{191.738}{761.124}{that}%
\TextAt{13}{217.928}{761.124}{L}%
\TextAt{9}{230.613}{761.124}{\Character{61}}%
\TextAt{13}{244.438}{761.124}{K}%
\TextAt{9}{255.256}{761.124}{\Character{40}}%
\TextAt{22}{263.129}{766.958}{n}%
\TextAt{16}{261.088}{771.246}{\Character{112}}%
\Vector{q 0 G 0 g 0.478 w 199.051 -2.515 m 208.806 -2.515 l S Q}%
\TextAt{9}{271.051}{761.124}{\Character{1}\Character{41}}%
\TextAt{8}{290.077}{761.124}{if}%
\TextAt{9}{301.622}{761.124}{\Character{1}}%
\TextAt{9}{316.076}{761.124}{\Character{61}}%
\TextAt{13}{329.923}{761.124}{L}%
\TextAt{3}{337.887}{765.462}{\Character{3}}%
\TextAt{21}{342.122}{765.462}{n}%
\TextAt{16}{350.975}{761.124}{\Character{92}}%
\TextAt{13}{362.161}{761.124}{K}%
\TextAt{21}{372.979}{765.462}{S}%
\TextAt{8}{379.073}{761.124}{,}%
\TextAt{8}{387.249}{761.124}{and}%
\TextAt{8}{410.833}{761.124}{then}%
\TextAt{8}{438.314}{761.124}{that}%
\TextAt{9}{464.504}{761.124}{\Character{1}}%
\TextAt{8}{478.978}{761.124}{is}%
\TextAt{8}{491.564}{761.124}{the}%
\TextAt{8}{512.541}{761.124}{said}%
\TextAt{8}{79.226}{744.078}{kernel.}%
\TextAt{8}{120.258}{744.078}{By}%
\TextAt{8}{139.293}{744.078}{chap.}%
\TextAt{8}{170.857}{744.078}{IV,}%
\TextAt{8}{191.673}{744.078}{(3.6),}%
\TextAt{8}{223.739}{744.078}{we}%
\TextAt{8}{241.638}{744.078}{certainly}%
\TextAt{8}{290.764}{744.078}{have}%
\TextAt{8}{318.43}{744.078}{that}%
\TextAt{13}{344.471}{744.078}{L}%
\TextAt{9}{356.901}{744.078}{\Character{61}}%
\TextAt{13}{370.47}{744.078}{K}%
\TextAt{9}{381.288}{744.078}{\Character{40}}%
\TextAt{22}{389.161}{749.912}{n}%
\TextAt{16}{387.12}{754.2}{\Character{112}}%
\Vector{q 0 G 0 g 0.478 w 325.083 -19.561 m 335.128 -19.561 l S Q}%
\TextAt{13}{397.083}{744.078}{D}%
\TextAt{9}{407.128}{744.078}{\Character{41}}%
\TextAt{8}{416.249}{744.078}{with}%
\TextAt{13}{443.58}{744.078}{D}%
\TextAt{9}{458.09}{744.078}{\Character{61}}%
\TextAt{13}{471.66}{744.078}{L}%
\TextAt{3}{479.625}{748.416}{\Character{3}}%
\TextAt{21}{483.859}{748.416}{n}%
\TextAt{16}{492.609}{744.078}{\Character{92}}%
\TextAt{13}{503.694}{744.078}{K}%
\TextAt{8}{514.512}{744.078}{.}%
\TextAt{8}{524.975}{744.078}{If}%
\TextAt{13}{79.226}{727.032}{x}%
\NativeMathAt{91.477}{727.032}{11.955}{\in}%
\TextAt{13}{105.047}{727.032}{D}%
\TextAt{8}{115.091}{727.032}{,}%
\TextAt{8}{123.911}{727.032}{then}%
\TextAt{13}{151.908}{727.032}{K}%
\TextAt{27}{161.867}{725.238}{p}%
\TextAt{9}{166.354}{727.032}{\Character{40}}%
\TextAt{22}{174.227}{731.053}{n}%
\TextAt{16}{172.187}{735.643}{\Character{112}}%
\Vector{q 0 G 0 g 0.478 w 110.149 -38.118 m 116.801 -38.118 l S Q}%
\TextAt{13}{182.149}{727.032}{x}%
\TextAt{9}{188.801}{727.032}{\Character{41}}%
\TextAt{16}{198.953}{727.032}{\Character{106}}%
\TextAt{13}{207.873}{727.032}{K}%
\TextAt{27}{217.832}{725.238}{p}%
\TextAt{8}{227.553}{727.032}{is}%
\TextAt{8}{240.654}{727.032}{unramified}%
\TextAt{8}{300.523}{727.032}{for}%
\TextAt{8}{319.745}{727.032}{all}%
\TextAt{14}{337.341}{727.032}{p}%
\NativeMathAt{348.923}{727.032}{11.955}{\notin}%
\TextAt{13}{362.493}{727.032}{S}%
\TextAt{8}{375.622}{727.032}{because}%
\TextAt{13}{420.261}{727.032}{S}%
\TextAt{8}{433.39}{727.032}{contains}%
\TextAt{8}{480.647}{727.032}{the}%
\TextAt{8}{502.14}{727.032}{places}%
\TextAt{8}{79.226}{709.986}{ramified}%
\TextAt{8}{124.112}{709.986}{in}%
\TextAt{13}{137.125}{709.986}{L}%
\TextAt{8}{145.089}{709.986}{.}%
\TextAt{8}{153.324}{709.986}{By}%
\TextAt{8}{171.047}{709.986}{chap.}%
\TextAt{8}{201.3}{709.986}{V,}%
\TextAt{8}{216.584}{709.986}{(3.3),}%
\TextAt{8}{247.299}{709.986}{we}%
\TextAt{8}{263.887}{709.986}{may}%
\TextAt{8}{288.616}{709.986}{therefore}%
\TextAt{8}{337.076}{709.986}{write}%
\TextAt{13}{366.348}{709.986}{x}%
\TextAt{9}{376.321}{709.986}{\Character{61}}%
\TextAt{13}{388.746}{709.986}{u}%
\TextAt{27}{395.409}{708.193}{p}%
\TextAt{13}{399.895}{709.986}{y}%
\TextAt{21}{406.032}{714.324}{n}%
\TextAt{27}{405.603}{707.03}{p}%
\TextAt{8}{411.668}{709.986}{,}%
\TextAt{8}{418.306}{709.986}{with}%
\TextAt{13}{444.326}{709.986}{u}%
\TextAt{27}{450.988}{708.193}{p}%
\NativeMathAt{458.796}{709.986}{11.955}{\in}%
\TextAt{13}{470.087}{709.986}{U}%
\TextAt{27}{478.037}{708.193}{p}%
\TextAt{8}{482.524}{709.986}{,}%
\TextAt{13}{489.161}{709.986}{y}%
\TextAt{27}{494.869}{708.193}{p}%
\NativeMathAt{502.677}{709.986}{11.955}{\in}%
\TextAt{13}{513.968}{709.986}{K}%
\TextAt{3}{524.785}{714.324}{\Character{3}}%
\TextAt{27}{523.927}{707.03}{p}%
\TextAt{8}{529.518}{709.986}{.}%
\TextAt{8}{79.226}{692.94}{Putting}%
\TextAt{13}{123.093}{692.94}{y}%
\TextAt{27}{128.801}{691.147}{p}%
\TextAt{9}{137.939}{692.94}{\Character{61}}%
\TextAt{9}{151.688}{692.94}{1}%
\TextAt{8}{162.227}{692.94}{for}%
\TextAt{14}{180.892}{692.94}{p}%
\NativeMathAt{191.527}{692.94}{11.955}{\in}%
\TextAt{13}{204.148}{692.94}{S}%
\TextAt{8}{212.044}{692.94}{,}%
\TextAt{8}{220.169}{692.94}{we}%
\TextAt{8}{238.177}{692.94}{get}%
\TextAt{8}{258.469}{692.94}{an}%
\TextAt{8}{275.509}{692.94}{idèle}%
\TextAt{13}{303.595}{692.94}{y}%
\TextAt{9}{314.383}{692.94}{\Character{61}}%
\TextAt{9}{328.132}{692.94}{\Character{40}}%
\TextAt{13}{332.691}{692.94}{y}%
\TextAt{27}{338.399}{691.147}{p}%
\TextAt{9}{342.886}{692.94}{\Character{41}}%
\TextAt{8}{352.117}{692.94}{which}%
\TextAt{8}{386.384}{692.94}{can}%
\TextAt{8}{408.624}{692.94}{be}%
\TextAt{8}{425.329}{692.94}{written}%
\TextAt{8}{467.08}{692.94}{as}%
\TextAt{8}{482.231}{692.94}{a}%
\TextAt{8}{492.768}{692.94}{product}%
\TextAt{13}{79.226}{675.894}{y}%
\TextAt{9}{88.684}{675.894}{\Character{61}}%
\TextAt{13}{101.109}{675.894}{\Character{11}z}%
\TextAt{8}{118.499}{675.894}{with}%
\TextAt{13}{145.159}{675.894}{\Character{11}}%
\NativeMathAt{156.001}{675.894}{11.955}{\in}%
\TextAt{13}{167.292}{675.894}{I}%
\TextAt{21}{173.395}{680.232}{S}%
\TextAt{21}{172.455}{672.873}{K}%
\TextAt{8}{180.658}{675.894}{,}%
\TextAt{13}{187.807}{675.894}{z}%
\NativeMathAt{197.099}{675.894}{11.955}{\in}%
\TextAt{13}{208.39}{675.894}{K}%
\TextAt{3}{219.207}{680.232}{\Character{3}}%
\TextAt{8}{223.94}{675.894}{.}%
\TextAt{8}{232.388}{675.894}{Then}%
\TextAt{13}{262.945}{675.894}{xz}%
\TextAt{3}{275.568}{680.232}{\Character{0}}%
\TextAt{21}{282.154}{680.232}{n}%
\TextAt{9}{291.112}{675.894}{\Character{61}}%
\TextAt{13}{303.537}{675.894}{u}%
\TextAt{27}{310.2}{674.101}{p}%
\TextAt{13}{314.686}{675.894}{\Character{11}}%
\TextAt{21}{322.208}{680.232}{n}%
\TextAt{27}{322.132}{672.938}{p}%
\NativeMathAt{331.165}{675.894}{11.955}{\in}%
\TextAt{13}{342.456}{675.894}{U}%
\TextAt{27}{350.407}{674.101}{p}%
\TextAt{8}{358.791}{675.894}{for}%
\TextAt{8}{376.676}{675.894}{all}%
\TextAt{14}{392.935}{675.894}{p}%
\NativeMathAt{402.239}{675.894}{11.955}{\notin}%
\TextAt{13}{413.53}{675.894}{S}%
\TextAt{8}{421.425}{675.894}{,}%
\TextAt{8}{428.574}{675.894}{i.e.,}%
\TextAt{13}{247.215}{647.939}{xz}%
\TextAt{3}{259.838}{652.875}{\Character{0}}%
\TextAt{21}{266.424}{652.875}{n}%
\NativeMathAt{275.381}{647.939}{11.955}{\in}%
\TextAt{13}{286.672}{647.939}{I}%
\TextAt{21}{292.775}{652.875}{S}%
\TextAt{21}{291.835}{644.984}{K}%
\TextAt{16}{302.694}{647.939}{\Character{92}}%
\TextAt{13}{313.321}{647.939}{K}%
\TextAt{3}{324.139}{652.875}{\Character{3}}%
\TextAt{9}{332.192}{647.939}{\Character{61}}%
\TextAt{13}{344.618}{647.939}{K}%
\TextAt{21}{355.435}{652.875}{S}%
\TextAt{13}{361.529}{647.939}{\Character{59}}%
\TextAt{8}{79.226}{619.985}{so}%
\TextAt{8}{93.596}{619.985}{that}%
\TextAt{13}{118.965}{619.985}{xz}%
\TextAt{3}{131.588}{624.323}{\Character{0}}%
\TextAt{21}{138.174}{624.323}{n}%
\NativeMathAt{147.131}{619.985}{11.955}{\in}%
\TextAt{9}{158.422}{619.985}{\Character{1}}%
\TextAt{8}{168.177}{619.985}{.}%
\TextAt{8}{176.626}{619.985}{This}%
\TextAt{8}{203.346}{619.985}{shows}%
\TextAt{8}{236.964}{619.985}{that}%
\TextAt{13}{262.332}{619.985}{D}%
\TextAt{9}{275.698}{619.985}{\Character{61}}%
\TextAt{9}{288.108}{619.985}{\Character{1}}%
\TextAt{13}{297.878}{619.985}{K}%
\TextAt{3}{308.696}{624.323}{\Character{3}}%
\TextAt{21}{312.93}{624.323}{n}%
\TextAt{8}{318.566}{619.985}{,}%
\TextAt{8}{325.716}{619.985}{and}%
\TextAt{8}{348.478}{619.985}{thus}%
\TextAt{13}{374.23}{619.985}{L}%
\TextAt{9}{385.515}{619.985}{\Character{61}}%
\TextAt{13}{397.94}{619.985}{K}%
\TextAt{9}{408.758}{619.985}{\Character{40}}%
\TextAt{22}{416.631}{625.819}{n}%
\TextAt{16}{414.59}{630.107}{\Character{112}}%
\Vector{q 0 G 0 g 0.478 w 352.553 -143.654 m 362.308 -143.654 l S Q}%
\TextAt{9}{424.553}{619.985}{\Character{1}\Character{41}}%
\TextAt{8}{438.86}{619.985}{.}%
\TextAt{8}{97.159}{599.95}{The}%
\TextAt{8}{122.628}{599.95}{field}%
\TextAt{13}{149.4}{599.95}{N}%
\TextAt{9}{165.755}{599.95}{\Character{61}}%
\TextAt{13}{180.591}{599.95}{K}%
\TextAt{9}{191.409}{599.95}{\Character{40}}%
\TextAt{22}{199.282}{606.222}{n}%
\TextAt{16}{197.241}{610.437}{\Character{112}}%
\Vector{q 0 G 0 g 0.478 w 135.204 -163.324 m 152.116 -163.324 l S Q}%
\TextAt{13}{207.204}{599.95}{K}%
\TextAt{21}{218.022}{603.404}{S}%
\TextAt{9}{224.116}{599.95}{\Character{41}}%
\TextAt{8}{233.98}{599.95}{contains}%
\TextAt{8}{281.315}{599.95}{the}%
\TextAt{8}{302.887}{599.95}{field}%
\TextAt{13}{329.659}{599.95}{L}%
\TextAt{8}{342.936}{599.95}{because}%
\TextAt{9}{387.652}{599.95}{\Character{1}}%
\TextAt{9}{403.122}{599.95}{\Character{61}}%
\TextAt{13}{417.976}{599.95}{L}%
\TextAt{3}{425.941}{604.288}{\Character{3}}%
\TextAt{21}{430.175}{604.288}{n}%
\TextAt{16}{439.433}{599.95}{\Character{92}}%
\TextAt{13}{451.024}{599.95}{K}%
\TextAt{21}{461.842}{604.288}{S}%
\NativeMathAt{473.668}{599.95}{11.955}{\subseteq}%
\TextAt{13}{488.699}{599.95}{K}%
\TextAt{21}{499.516}{604.288}{S}%
\TextAt{8}{505.61}{599.95}{.}%
\TextAt{8}{518.304}{599.95}{By}%
\TextAt{8}{79.226}{582.904}{Kummer}%
\TextAt{8}{127.991}{582.904}{theory,}%
\TextAt{8}{167.013}{582.904}{chap.}%
\TextAt{8}{197.905}{582.904}{IV,}%
\TextAt{8}{218.05}{582.904}{(3.6),}%
\TextAt{8}{249.276}{582.904}{we}%
\TextAt{8}{266.504}{582.904}{have}%
\TextAt{13}{210.564}{554.949}{G}%
\TextAt{9}{219.797}{554.949}{\Character{40}}%
\TextAt{13}{224.35}{554.949}{N}%
\TextAt{16}{238.293}{554.949}{\Character{106}}%
\TextAt{13}{244.935}{554.949}{K}%
\TextAt{9}{255.752}{554.949}{\Character{41}}%
\NativeMathAt{263.626}{554.949}{11.955}{\cong}%
\TextAt{8}{276.245}{554.949}{Hom}%
\TextAt{9}{300.634}{554.949}{\Character{40}}%
\TextAt{13}{305.186}{554.949}{K}%
\TextAt{21}{316.004}{559.885}{S}%
\TextAt{9}{322.097}{554.949}{\Character{47}\Character{40}}%
\TextAt{13}{332.503}{554.949}{K}%
\TextAt{21}{343.32}{559.885}{S}%
\TextAt{9}{349.414}{554.949}{\Character{41}}%
\TextAt{21}{353.967}{559.885}{n}%
\TextAt{13}{359.603}{554.949}{\Character{59}}%
\TextAt{18}{364.847}{554.949}{Z}%
\TextAt{9}{372.817}{554.949}{\Character{47}}%
\TextAt{13}{378.67}{554.949}{n}%
\TextAt{18}{385.658}{554.949}{Z}%
\TextAt{9}{393.628}{554.949}{\Character{41}}%
\TextAt{13}{398.18}{554.949}{\Character{58}}%
\TextAt{8}{79.226}{526.995}{By}%
\TextAt{8}{98.092}{526.995}{(1.1),}%
\TextAt{13}{129.947}{526.995}{K}%
\TextAt{21}{140.765}{531.333}{S}%
\TextAt{8}{151.259}{526.995}{is}%
\TextAt{8}{163.526}{526.995}{the}%
\TextAt{8}{184.185}{526.995}{product}%
\TextAt{8}{228.587}{526.995}{of}%
\TextAt{8}{242.42}{526.995}{a}%
\TextAt{8}{252.678}{526.995}{free}%
\TextAt{8}{275.609}{526.995}{group}%
\TextAt{8}{309.287}{526.995}{of}%
\TextAt{8}{323.12}{526.995}{rank}%
\TextAt{13}{350.618}{526.995}{s}%
\TextAt{16}{359.131}{526.995}{\Character{0}}%
\TextAt{9}{371.429}{526.995}{1}%
\TextAt{8}{381.682}{526.995}{and}%
\TextAt{8}{404.948}{526.995}{of}%
\TextAt{8}{418.78}{526.995}{the}%
\TextAt{8}{439.44}{526.995}{cyclic}%
\TextAt{8}{472.126}{526.995}{group}%
\TextAt{13}{505.804}{526.995}{\Character{22}}%
\TextAt{9}{512.847}{526.995}{\Character{40}}%
\TextAt{13}{517.399}{526.995}{K}%
\TextAt{9}{528.217}{526.995}{\Character{41}}%
\TextAt{8}{79.226}{509.949}{whose}%
\TextAt{8}{114.481}{509.949}{order}%
\TextAt{8}{145.779}{509.949}{is}%
\TextAt{8}{158.271}{509.949}{divisible}%
\TextAt{8}{204.908}{509.949}{by}%
\TextAt{13}{221.895}{509.949}{n}%
\TextAt{8}{228.883}{509.949}{.}%
\TextAt{8}{239.516}{509.949}{Therefore}%
\TextAt{13}{293.242}{509.949}{K}%
\TextAt{21}{304.06}{514.287}{S}%
\TextAt{9}{310.154}{509.949}{\Character{47}\Character{40}}%
\TextAt{13}{320.559}{509.949}{K}%
\TextAt{21}{331.377}{514.287}{S}%
\TextAt{9}{337.47}{509.949}{\Character{41}}%
\TextAt{21}{342.023}{514.287}{n}%
\TextAt{8}{352.285}{509.949}{is}%
\TextAt{8}{364.777}{509.949}{a}%
\TextAt{8}{375.261}{509.949}{free}%
\TextAt{9}{398.417}{509.949}{\Character{40}}%
\TextAt{18}{402.97}{509.949}{Z}%
\TextAt{9}{410.94}{509.949}{\Character{47}}%
\TextAt{13}{416.793}{509.949}{n}%
\TextAt{18}{423.781}{509.949}{Z}%
\TextAt{9}{431.751}{509.949}{\Character{41}}%
\TextAt{8}{436.303}{509.949}{-module}%
\TextAt{8}{482.222}{509.949}{of}%
\TextAt{8}{496.28}{509.949}{rank}%
\TextAt{13}{524.004}{509.949}{s}%
\TextAt{8}{529.518}{509.949}{,}%
\TextAt{8}{79.226}{492.903}{and}%
\TextAt{8}{103.552}{492.903}{so}%
\TextAt{8}{119.485}{492.903}{is}%
\TextAt{13}{132.811}{492.903}{G}%
\TextAt{9}{142.045}{492.903}{\Character{40}}%
\TextAt{13}{146.597}{492.903}{N}%
\TextAt{16}{163.204}{492.903}{\Character{106}}%
\TextAt{13}{172.509}{492.903}{K}%
\TextAt{9}{183.327}{492.903}{\Character{41}}%
\TextAt{8}{187.879}{492.903}{.}%
\TextAt{8}{201.016}{492.903}{Moreover,}%
\TextAt{13}{257.605}{492.903}{G}%
\TextAt{9}{266.839}{492.903}{\Character{40}}%
\TextAt{13}{271.391}{492.903}{N}%
\TextAt{16}{287.998}{492.903}{\Character{106}}%
\TextAt{13}{297.303}{492.903}{K}%
\TextAt{9}{308.121}{492.903}{\Character{41}\Character{47}}%
\TextAt{13}{318.526}{492.903}{G}%
\TextAt{9}{327.76}{492.903}{\Character{40}}%
\TextAt{13}{332.312}{492.903}{N}%
\TextAt{16}{348.919}{492.903}{\Character{106}}%
\TextAt{13}{358.224}{492.903}{L}%
\TextAt{9}{366.188}{492.903}{\Character{41}}%
\NativeMathAt{376.725}{492.903}{11.955}{\cong}%
\TextAt{13}{392.007}{492.903}{G}%
\TextAt{9}{401.241}{492.903}{\Character{40}}%
\TextAt{13}{405.793}{492.903}{L}%
\TextAt{16}{419.742}{492.903}{\Character{106}}%
\TextAt{13}{429.047}{492.903}{K}%
\TextAt{9}{439.865}{492.903}{\Character{41}}%
\NativeMathAt{450.401}{492.903}{11.955}{\cong}%
\TextAt{9}{465.684}{492.903}{\Character{40}}%
\TextAt{18}{470.236}{492.903}{Z}%
\TextAt{9}{478.207}{492.903}{\Character{47}}%
\TextAt{13}{484.06}{492.903}{n}%
\TextAt{18}{491.047}{492.903}{Z}%
\TextAt{9}{499.017}{492.903}{\Character{41}}%
\TextAt{21}{503.57}{497.241}{r}%
\TextAt{8}{513.585}{492.903}{is}%
\TextAt{8}{526.911}{492.903}{a}%
\TextAt{8}{79.226}{475.857}{free}%
\TextAt{9}{103.029}{475.857}{\Character{40}}%
\TextAt{18}{107.582}{475.857}{Z}%
\TextAt{9}{115.552}{475.857}{\Character{47}}%
\TextAt{13}{121.405}{475.857}{n}%
\TextAt{18}{128.393}{475.857}{Z}%
\TextAt{9}{136.363}{475.857}{\Character{41}}%
\TextAt{8}{140.915}{475.857}{-module}%
\TextAt{8}{187.481}{475.857}{of}%
\TextAt{8}{202.186}{475.857}{rank}%
\TextAt{13}{230.557}{475.857}{r}%
\TextAt{8}{241.43}{475.857}{so}%
\TextAt{8}{257.175}{475.857}{that}%
\TextAt{13}{283.92}{475.857}{r}%
\NativeMathAt{295.185}{475.857}{11.955}{\leq}%
\TextAt{13}{310.148}{475.857}{s}%
\TextAt{8}{315.662}{475.857}{,}%
\TextAt{8}{324.531}{475.857}{and}%
\TextAt{13}{348.669}{475.857}{G}%
\TextAt{9}{357.902}{475.857}{\Character{40}}%
\TextAt{13}{362.455}{475.857}{N}%
\TextAt{16}{378.742}{475.857}{\Character{106}}%
\TextAt{13}{387.728}{475.857}{L}%
\TextAt{9}{395.692}{475.857}{\Character{41}}%
\TextAt{8}{405.517}{475.857}{is}%
\TextAt{8}{418.657}{475.857}{a}%
\TextAt{8}{429.788}{475.857}{free}%
\TextAt{9}{453.591}{475.857}{\Character{40}}%
\TextAt{18}{458.143}{475.857}{Z}%
\TextAt{9}{466.113}{475.857}{\Character{47}}%
\TextAt{13}{471.966}{475.857}{n}%
\TextAt{18}{478.954}{475.857}{Z}%
\TextAt{9}{486.924}{475.857}{\Character{41}}%
\TextAt{8}{491.476}{475.857}{-module}%
\TextAt{8}{79.226}{458.811}{of}%
\TextAt{8}{93.628}{458.811}{rank}%
\TextAt{13}{121.695}{458.811}{s}%
\TextAt{16}{130.596}{458.811}{\Character{0}}%
\TextAt{13}{143.282}{458.811}{r}%
\TextAt{8}{148.882}{458.811}{.}%
\TextAt{8}{160.547}{458.811}{Let}%
\TextAt{13}{182.588}{458.811}{\Character{27}}%
\TextAt{17}{189.241}{457.017}{1}%
\TextAt{13}{193.974}{458.811}{\Character{59}}%
\TextAt{13}{199.198}{458.811}{\Character{58}}%
\TextAt{13}{204.435}{458.811}{\Character{58}}%
\TextAt{13}{209.659}{458.811}{\Character{58}}%
\TextAt{13}{214.896}{458.811}{\Character{59}}%
\TextAt{13}{220.132}{458.811}{\Character{27}}%
\TextAt{21}{226.848}{457.017}{s}%
\TextAt{3}{230.764}{457.017}{\Character{0}}%
\TextAt{21}{237.35}{457.017}{r}%
\TextAt{8}{246.874}{458.811}{be}%
\TextAt{8}{263.871}{458.811}{a}%
\TextAt{18}{274.698}{458.811}{Z}%
\TextAt{9}{282.668}{458.811}{\Character{47}}%
\TextAt{13}{288.521}{458.811}{n}%
\TextAt{18}{295.509}{458.811}{Z}%
\TextAt{8}{303.479}{458.811}{-basis}%
\TextAt{8}{337.188}{458.811}{of}%
\TextAt{13}{351.59}{458.811}{G}%
\TextAt{9}{360.824}{458.811}{\Character{40}}%
\TextAt{13}{365.376}{458.811}{N}%
\TextAt{16}{381.146}{458.811}{\Character{106}}%
\TextAt{13}{389.615}{458.811}{L}%
\TextAt{9}{397.579}{458.811}{\Character{41}}%
\TextAt{8}{402.132}{458.811}{,}%
\TextAt{8}{410.621}{458.811}{and}%
\TextAt{8}{434.455}{458.811}{let}%
\TextAt{13}{452.432}{458.811}{N}%
\TextAt{21}{461.806}{457.017}{i}%
\TextAt{8}{470.156}{458.811}{be}%
\TextAt{8}{487.153}{458.811}{the}%
\TextAt{8}{508.381}{458.811}{fixed}%
\TextAt{8}{79.226}{441.765}{field}%
\TextAt{8}{105.139}{441.765}{of}%
\TextAt{13}{119.025}{441.765}{\Character{27}}%
\TextAt{21}{125.679}{439.971}{i}%
\TextAt{8}{129.06}{441.765}{,}%
\TextAt{13}{136.904}{441.765}{i}%
\TextAt{9}{145.166}{441.765}{\Character{61}}%
\TextAt{9}{158.532}{441.765}{1}%
\TextAt{13}{164.393}{441.765}{\Character{59}}%
\TextAt{13}{169.617}{441.765}{\Character{58}}%
\TextAt{13}{174.854}{441.765}{\Character{58}}%
\TextAt{13}{180.078}{441.765}{\Character{58}}%
\TextAt{13}{185.315}{441.765}{\Character{59}}%
\TextAt{13}{190.551}{441.765}{s}%
\TextAt{16}{199.163}{441.765}{\Character{0}}%
\TextAt{13}{211.497}{441.765}{r}%
\TextAt{8}{217.098}{441.765}{.}%
\TextAt{8}{227.215}{441.765}{Then}%
\TextAt{13}{258.328}{441.765}{L}%
\TextAt{9}{270.561}{441.765}{\Character{61}}%
\TextAt{20}{283.934}{450.731}{\Character{84}}%
\TextAt{21}{293.897}{447.576}{s}%
\TextAt{3}{297.813}{447.576}{\Character{0}}%
\TextAt{21}{304.399}{447.576}{r}%
\TextAt{21}{293.897}{438.278}{i}%
\TextAt{17}{296.78}{438.278}{\Character{61}1}%
\TextAt{13}{310.947}{441.765}{N}%
\TextAt{21}{320.32}{439.971}{i}%
\TextAt{8}{323.702}{441.765}{.}%
\TextAt{8}{333.818}{441.765}{For}%
\TextAt{8}{355.344}{441.765}{every}%
\TextAt{13}{386.792}{441.765}{i}%
\TextAt{9}{395.054}{441.765}{\Character{61}}%
\TextAt{9}{408.42}{441.765}{1}%
\TextAt{13}{414.28}{441.765}{\Character{59}}%
\TextAt{13}{419.504}{441.765}{\Character{58}}%
\TextAt{13}{424.741}{441.765}{\Character{58}}%
\TextAt{13}{429.977}{441.765}{\Character{58}}%
\TextAt{13}{435.202}{441.765}{\Character{59}}%
\TextAt{13}{440.438}{441.765}{s}%
\TextAt{16}{449.051}{441.765}{\Character{0}}%
\TextAt{13}{461.385}{441.765}{r}%
\TextAt{8}{471.439}{441.765}{we}%
\TextAt{8}{489.223}{441.765}{choose}%
\TextAt{8}{526.911}{441.765}{a}%
\TextAt{8}{79.226}{424.719}{prime}%
\TextAt{14}{112.215}{424.719}{P}%
\TextAt{21}{122.109}{422.925}{i}%
\TextAt{8}{129.213}{424.719}{of}%
\TextAt{13}{142.368}{424.719}{N}%
\TextAt{21}{151.742}{422.925}{i}%
\TextAt{8}{158.846}{424.719}{which}%
\TextAt{8}{192.158}{424.719}{is}%
\TextAt{8}{203.747}{424.719}{nonsplit}%
\TextAt{8}{248.512}{424.719}{in}%
\TextAt{13}{261.991}{424.719}{N}%
\TextAt{8}{276.336}{424.719}{such}%
\TextAt{8}{302.559}{424.719}{that}%
\TextAt{8}{327.753}{424.719}{the}%
\TextAt{8}{347.735}{424.719}{primes}%
\TextAt{14}{385.339}{424.719}{p}%
\TextAt{17}{391.322}{422.925}{1}%
\TextAt{13}{396.054}{424.719}{\Character{59}}%
\TextAt{13}{401.278}{424.719}{\Character{58}}%
\TextAt{13}{406.515}{424.719}{\Character{58}}%
\TextAt{13}{411.751}{424.719}{\Character{58}}%
\TextAt{13}{416.976}{424.719}{\Character{59}}%
\TextAt{14}{422.275}{424.719}{p}%
\TextAt{21}{428.258}{422.925}{s}%
\TextAt{3}{432.174}{422.925}{\Character{0}}%
\TextAt{21}{438.761}{422.925}{r}%
\TextAt{8}{447.038}{424.719}{of}%
\TextAt{13}{460.194}{424.719}{K}%
\TextAt{8}{474.734}{424.719}{lying}%
\TextAt{8}{503.503}{424.719}{below}%
\TextAt{14}{79.226}{407.673}{P}%
\TextAt{17}{89.119}{405.88}{1}%
\TextAt{13}{93.852}{407.673}{\Character{59}}%
\TextAt{13}{99.076}{407.673}{\Character{58}}%
\TextAt{13}{104.313}{407.673}{\Character{58}}%
\TextAt{13}{109.537}{407.673}{\Character{58}}%
\TextAt{13}{114.774}{407.673}{\Character{59}}%
\TextAt{14}{120.072}{407.673}{P}%
\TextAt{21}{129.966}{405.88}{s}%
\TextAt{3}{133.882}{405.88}{\Character{0}}%
\TextAt{21}{140.468}{405.88}{r}%
\TextAt{8}{149.053}{407.673}{are}%
\TextAt{8}{168.696}{407.673}{all}%
\TextAt{8}{185.088}{407.673}{distinct,}%
\TextAt{8}{230.839}{407.673}{and}%
\TextAt{8}{253.734}{407.673}{do}%
\TextAt{8}{270.126}{407.673}{not}%
\TextAt{8}{291.072}{407.673}{belong}%
\TextAt{8}{328.601}{407.673}{to}%
\TextAt{13}{343.044}{407.673}{S}%
\TextAt{8}{350.939}{407.673}{.}%
\TextAt{8}{359.785}{407.673}{This}%
\TextAt{8}{386.638}{407.673}{is}%
\TextAt{8}{398.534}{407.673}{possible}%
\TextAt{8}{442.686}{407.673}{by}%
\TextAt{8}{459.077}{407.673}{(3.7).}%
\TextAt{8}{492.001}{407.673}{We}%
\TextAt{8}{512.278}{407.673}{now}%
\TextAt{8}{79.226}{390.627}{show}%
\TextAt{8}{108.847}{390.627}{that}%
\TextAt{8}{134.834}{390.627}{the}%
\TextAt{8}{155.608}{390.627}{set}%
\TextAt{13}{174.494}{390.627}{T}%
\TextAt{9}{187.354}{390.627}{\Character{61}}%
\TextAt{16}{200.833}{390.627}{\Character{102}}%
\TextAt{14}{206.81}{390.627}{p}%
\TextAt{17}{212.793}{388.834}{1}%
\TextAt{13}{217.526}{390.627}{\Character{59}}%
\TextAt{13}{222.75}{390.627}{\Character{58}}%
\TextAt{13}{227.987}{390.627}{\Character{58}}%
\TextAt{13}{233.211}{390.627}{\Character{58}}%
\TextAt{13}{238.448}{390.627}{\Character{59}}%
\TextAt{14}{243.746}{390.627}{p}%
\TextAt{21}{249.73}{388.834}{s}%
\TextAt{3}{253.645}{388.834}{\Character{0}}%
\TextAt{21}{260.232}{388.834}{r}%
\TextAt{16}{264.787}{390.627}{\Character{103}}%
\TextAt{8}{275.28}{390.627}{realizes}%
\TextAt{8}{316.928}{390.627}{the}%
\TextAt{8}{337.702}{390.627}{group}%
\TextAt{9}{371.495}{390.627}{\Character{1}}%
\TextAt{9}{385.602}{390.627}{\Character{61}}%
\TextAt{13}{399.102}{390.627}{L}%
\TextAt{3}{407.067}{394.965}{\Character{3}}%
\TextAt{21}{411.301}{394.965}{n}%
\TextAt{16}{420.015}{390.627}{\Character{92}}%
\TextAt{13}{431.063}{390.627}{K}%
\TextAt{21}{441.881}{394.965}{S}%
\TextAt{8}{452.49}{390.627}{as}%
\TextAt{8}{467.478}{390.627}{the}%
\TextAt{8}{488.252}{390.627}{kernel}%
\TextAt{8}{523.337}{390.627}{of}%
\TextAt{13}{79.226}{373.581}{K}%
\TextAt{21}{90.044}{377.919}{S}%
\NativeArrowAt{99.459}{373.581}{11.955}{19.26}{\longrightarrow}%
\TextAt{20}{122.041}{382.547}{\Character{81}}%
\TextAt{27}{133.332}{370.094}{p}%
\NativeMathAt{137.32}{370.094}{7.97}{\in}%
\TextAt{21}{142.966}{370.094}{T}%
\TextAt{13}{151.563}{373.581}{K}%
\TextAt{3}{162.381}{377.919}{\Character{3}}%
\TextAt{27}{161.522}{370.625}{p}%
\TextAt{9}{167.113}{373.581}{\Character{47}}%
\TextAt{13}{172.966}{373.581}{K}%
\TextAt{3}{183.784}{377.919}{\Character{3}}%
\TextAt{21}{188.018}{377.919}{n}%
\TextAt{27}{182.925}{370.625}{p}%
\TextAt{8}{193.654}{373.581}{.}%
\TextAt{13}{97.159}{353.546}{N}%
\TextAt{21}{106.533}{351.753}{i}%
\TextAt{8}{115.047}{353.546}{is}%
\TextAt{8}{128.047}{353.546}{the}%
\TextAt{8}{149.44}{353.546}{decomposition}%
\TextAt{8}{227.811}{353.546}{field}%
\TextAt{8}{254.404}{353.546}{of}%
\TextAt{13}{268.97}{353.546}{N}%
\TextAt{16}{285.02}{353.546}{\Character{106}}%
\TextAt{13}{293.768}{353.546}{K}%
\TextAt{8}{309.719}{353.546}{at}%
\TextAt{8}{325.265}{353.546}{the}%
\TextAt{8}{346.658}{353.546}{unique}%
\TextAt{8}{385.935}{353.546}{prime}%
\TextAt{14}{420.335}{353.546}{P}%
\TextAt{3}{430.228}{357.885}{\Character{48}}%
\TextAt{21}{430.228}{350.591}{i}%
\TextAt{8}{438.743}{353.546}{above}%
\TextAt{14}{473.154}{353.546}{P}%
\TextAt{21}{483.048}{351.753}{i}%
\TextAt{8}{486.429}{353.546}{,}%
\TextAt{8}{495.123}{353.546}{for}%
\TextAt{13}{514.244}{353.546}{i}%
\TextAt{9}{523.665}{353.546}{\Character{61}}%
\TextAt{9}{79.226}{336.5}{1}%
\TextAt{13}{85.079}{336.5}{\Character{59}}%
\TextAt{13}{90.303}{336.5}{\Character{58}}%
\TextAt{13}{95.54}{336.5}{\Character{58}}%
\TextAt{13}{100.776}{336.5}{\Character{58}}%
\TextAt{13}{106.001}{336.5}{\Character{59}}%
\TextAt{13}{111.237}{336.5}{s}%
\TextAt{16}{119.099}{336.5}{\Character{0}}%
\TextAt{13}{130.682}{336.5}{r}%
\TextAt{8}{136.283}{336.5}{.}%
\TextAt{8}{144.67}{336.5}{Indeed,}%
\TextAt{8}{185.806}{336.5}{this}%
\TextAt{8}{208.447}{336.5}{decomposition}%
\TextAt{8}{285.4}{336.5}{field}%
\TextAt{13}{310.575}{336.5}{Z}%
\TextAt{21}{318.587}{334.707}{i}%
\TextAt{8}{325.684}{336.5}{is}%
\TextAt{8}{337.266}{336.5}{contained}%
\TextAt{8}{390.094}{336.5}{in}%
\TextAt{13}{403.565}{336.5}{N}%
\TextAt{21}{412.938}{334.707}{i}%
\TextAt{8}{420.035}{336.5}{because}%
\TextAt{14}{463.155}{336.5}{P}%
\TextAt{21}{473.048}{334.707}{i}%
\TextAt{8}{480.145}{336.5}{is}%
\TextAt{8}{491.727}{336.5}{nonsplit}%
\TextAt{8}{79.226}{319.454}{in}%
\TextAt{13}{93.089}{319.454}{N}%
\TextAt{8}{103.711}{319.454}{.}%
\TextAt{8}{112.789}{319.454}{On}%
\TextAt{8}{132.509}{319.454}{the}%
\TextAt{8}{152.875}{319.454}{other}%
\TextAt{8}{183.654}{319.454}{hand,}%
\TextAt{8}{216.435}{319.454}{the}%
\TextAt{8}{236.801}{319.454}{prime}%
\TextAt{14}{270.174}{319.454}{p}%
\TextAt{21}{276.157}{317.661}{i}%
\TextAt{8}{283.646}{319.454}{is}%
\TextAt{8}{295.619}{319.454}{unramified}%
\TextAt{8}{354.361}{319.454}{in}%
\TextAt{13}{368.224}{319.454}{N}%
\TextAt{8}{378.846}{319.454}{,}%
\TextAt{8}{386.258}{319.454}{because}%
\TextAt{8}{429.769}{319.454}{by}%
\TextAt{8}{446.238}{319.454}{chap.}%
\TextAt{8}{477.34}{319.454}{V,}%
\TextAt{8}{493.474}{319.454}{(3.3),}%
\TextAt{8}{524.963}{319.454}{it}%
\TextAt{8}{79.226}{302.408}{is}%
\TextAt{8}{90.774}{302.408}{unramified}%
\TextAt{8}{149.09}{302.408}{in}%
\TextAt{8}{162.527}{302.408}{every}%
\TextAt{8}{193.203}{302.408}{extension}%
\TextAt{13}{244.752}{302.408}{K}%
\TextAt{9}{255.57}{302.408}{\Character{40}}%
\TextAt{22}{263.443}{306.429}{n}%
\TextAt{16}{261.402}{311.019}{\Character{112}}%
\Vector{q 0 G 0 g 0.478 w 199.365 -462.742 m 206.027 -462.742 l S Q}%
\TextAt{13}{271.365}{302.408}{u}%
\TextAt{9}{278.027}{302.408}{\Character{41}}%
\TextAt{8}{282.58}{302.408}{,}%
\TextAt{13}{289.556}{302.408}{u}%
\NativeMathAt{299.539}{302.408}{11.955}{\in}%
\TextAt{13}{310.83}{302.408}{K}%
\TextAt{21}{321.648}{306.747}{S}%
\TextAt{8}{327.742}{302.408}{.}%
\TextAt{8}{336.118}{302.408}{The}%
\TextAt{8}{359.955}{302.408}{decomposition}%
\TextAt{8}{436.874}{302.408}{group}%
\TextAt{13}{469.833}{302.408}{G}%
\TextAt{9}{479.067}{302.408}{\Character{40}}%
\TextAt{13}{483.619}{302.408}{N}%
\TextAt{16}{497.563}{302.408}{\Character{106}}%
\TextAt{13}{504.204}{302.408}{Z}%
\TextAt{21}{512.217}{300.615}{i}%
\TextAt{9}{515.598}{302.408}{\Character{41}}%
\NativeMathAt{523.471}{302.408}{11.955}{\supseteq}%
\TextAt{13}{79.226}{285.362}{G}%
\TextAt{9}{88.46}{285.362}{\Character{40}}%
\TextAt{13}{93.012}{285.362}{N}%
\TextAt{16}{107.335}{285.362}{\Character{106}}%
\TextAt{13}{114.357}{285.362}{N}%
\TextAt{21}{123.731}{283.569}{i}%
\TextAt{9}{127.112}{285.362}{\Character{41}}%
\TextAt{8}{135.785}{285.362}{is}%
\TextAt{8}{147.772}{285.362}{therefore}%
\TextAt{8}{197.094}{285.362}{cyclic,}%
\TextAt{8}{232.808}{285.362}{and}%
\TextAt{8}{255.794}{285.362}{necessarily}%
\TextAt{8}{314.346}{285.362}{of}%
\TextAt{8}{327.899}{285.362}{order}%
\TextAt{13}{358.691}{285.362}{n}%
\TextAt{8}{369.799}{285.362}{since}%
\TextAt{8}{398.691}{285.362}{each}%
\TextAt{8}{425.251}{285.362}{element}%
\TextAt{8}{468.716}{285.362}{of}%
\TextAt{13}{482.269}{285.362}{G}%
\TextAt{9}{491.502}{285.362}{\Character{40}}%
\TextAt{13}{496.054}{285.362}{N}%
\TextAt{16}{510.378}{285.362}{\Character{106}}%
\TextAt{13}{517.399}{285.362}{K}%
\TextAt{9}{528.217}{285.362}{\Character{41}}%
\TextAt{8}{79.226}{268.316}{has}%
\TextAt{8}{100.1}{268.316}{order}%
\TextAt{8}{130.669}{268.316}{dividing}%
\TextAt{13}{175.872}{268.316}{n}%
\TextAt{8}{182.859}{268.316}{.}%
\TextAt{8}{191.308}{268.316}{This}%
\TextAt{8}{218.027}{268.316}{shows}%
\TextAt{8}{251.645}{268.316}{that}%
\TextAt{13}{277.014}{268.316}{N}%
\TextAt{21}{286.388}{266.523}{i}%
\TextAt{9}{293.09}{268.316}{\Character{61}}%
\TextAt{13}{305.516}{268.316}{Z}%
\TextAt{21}{313.528}{266.523}{i}%
\TextAt{8}{316.909}{268.316}{.}%
\TextAt{8}{97.159}{248.281}{From}%
\TextAt{13}{128.853}{248.281}{L}%
\TextAt{9}{141.789}{248.281}{\Character{61}}%
\TextAt{20}{155.866}{257.248}{\Character{84}}%
\TextAt{21}{165.828}{254.093}{s}%
\TextAt{3}{169.744}{254.093}{\Character{0}}%
\TextAt{21}{176.331}{254.093}{r}%
\TextAt{21}{165.828}{244.794}{i}%
\TextAt{17}{168.712}{244.794}{\Character{61}1}%
\TextAt{13}{182.878}{248.281}{N}%
\TextAt{21}{192.252}{246.488}{i}%
\TextAt{8}{200.499}{248.281}{it}%
\TextAt{8}{213.172}{248.281}{follows}%
\TextAt{8}{252.577}{248.281}{that}%
\TextAt{13}{278.915}{248.281}{L}%
\TextAt{16}{291.852}{248.281}{\Character{106}}%
\TextAt{13}{300.144}{248.281}{K}%
\TextAt{8}{315.828}{248.281}{is}%
\TextAt{8}{328.561}{248.281}{the}%
\TextAt{8}{349.687}{248.281}{maximal}%
\TextAt{8}{398.464}{248.281}{subextension}%
\TextAt{8}{469.144}{248.281}{of}%
\TextAt{13}{483.443}{248.281}{N}%
\TextAt{16}{499.037}{248.281}{\Character{106}}%
\TextAt{13}{507.33}{248.281}{K}%
\TextAt{8}{523.014}{248.281}{in}%
\TextAt{8}{79.226}{231.235}{which}%
\TextAt{8}{112.713}{231.235}{the}%
\TextAt{8}{132.869}{231.235}{primes}%
\TextAt{14}{170.647}{231.235}{p}%
\TextAt{17}{176.63}{229.442}{1}%
\TextAt{13}{181.363}{231.235}{\Character{59}}%
\TextAt{13}{186.587}{231.235}{\Character{58}}%
\TextAt{13}{191.824}{231.235}{\Character{58}}%
\TextAt{13}{197.048}{231.235}{\Character{58}}%
\TextAt{13}{202.285}{231.235}{\Character{59}}%
\TextAt{14}{207.584}{231.235}{p}%
\TextAt{21}{213.567}{229.442}{s}%
\TextAt{3}{217.483}{229.442}{\Character{0}}%
\TextAt{21}{224.069}{229.442}{r}%
\TextAt{8}{232.521}{231.235}{split}%
\TextAt{8}{258.595}{231.235}{completely.}%
\TextAt{8}{321.021}{231.235}{For}%
\TextAt{13}{341.991}{231.235}{x}%
\NativeMathAt{351.964}{231.235}{11.955}{\in}%
\TextAt{13}{363.255}{231.235}{K}%
\TextAt{21}{374.072}{235.574}{S}%
\TextAt{8}{384.063}{231.235}{we}%
\TextAt{8}{401.291}{231.235}{therefore}%
\TextAt{8}{450.391}{231.235}{have}%
\TextAt{13}{188.558}{203.281}{x}%
\NativeMathAt{198.531}{203.281}{11.955}{\in}%
\TextAt{9}{209.822}{203.281}{\Character{1}}%
\NativeArrowAt{222.897}{203.281}{11.955}{21.926}{\Longleftrightarrow}%
\TextAt{13}{248.136}{203.281}{K}%
\TextAt{9}{258.954}{203.281}{\Character{40}}%
\TextAt{22}{266.827}{208.002}{n}%
\TextAt{16}{264.786}{212.476}{\Character{112}}%
\Vector{q 0 G 0 g 0.478 w 202.749 -561.285 m 209.401 -561.285 l S Q}%
\TextAt{13}{274.749}{203.281}{x}%
\TextAt{9}{281.401}{203.281}{\Character{41}}%
\NativeMathAt{289.274}{203.281}{11.955}{\subseteq}%
\TextAt{13}{301.893}{203.281}{L}%
\NativeArrowAt{222.897}{183.246}{11.955}{21.926}{\Longleftrightarrow}%
\TextAt{13}{248.136}{183.246}{K}%
\TextAt{27}{258.095}{181.453}{p}%
\TextAt{22}{262.084}{180.238}{i}%
\TextAt{9}{265.744}{183.246}{\Character{40}}%
\TextAt{22}{273.617}{187.967}{n}%
\TextAt{16}{271.576}{192.441}{\Character{112}}%
\Vector{q 0 G 0 g 0.478 w 209.539 -581.32 m 216.191 -581.32 l S Q}%
\TextAt{13}{281.539}{183.246}{x}%
\TextAt{9}{288.191}{183.246}{\Character{41}}%
\TextAt{9}{296.046}{183.246}{\Character{61}}%
\TextAt{13}{308.49}{183.246}{K}%
\TextAt{27}{318.449}{181.453}{p}%
\TextAt{22}{322.437}{180.238}{i}%
\TextAt{13}{326.097}{183.246}{\Character{59}}%
\TextAt{13}{343.277}{183.246}{i}%
\TextAt{9}{350.611}{183.246}{\Character{61}}%
\TextAt{9}{363.021}{183.246}{1}%
\TextAt{13}{368.889}{183.246}{\Character{59}}%
\TextAt{13}{374.113}{183.246}{\Character{58}}%
\TextAt{13}{379.35}{183.246}{\Character{58}}%
\TextAt{13}{384.586}{183.246}{\Character{58}}%
\TextAt{13}{389.811}{183.246}{\Character{59}}%
\TextAt{13}{395.047}{183.246}{s}%
\TextAt{16}{403.281}{183.246}{\Character{0}}%
\TextAt{13}{415.236}{183.246}{r}%
\TextAt{13}{420.185}{183.246}{\Character{59}}%
\NativeArrowAt{222.897}{163.211}{11.955}{21.926}{\Longleftrightarrow}%
\TextAt{13}{248.136}{163.211}{x}%
\NativeMathAt{258.109}{163.211}{11.955}{\in}%
\TextAt{13}{269.4}{163.211}{K}%
\TextAt{3}{280.218}{168.148}{\Character{3}}%
\TextAt{21}{284.452}{168.148}{n}%
\TextAt{27}{279.359}{160.256}{p}%
\TextAt{22}{283.348}{159.041}{i}%
\TextAt{13}{290.088}{163.211}{\Character{59}}%
\TextAt{13}{307.268}{163.211}{i}%
\TextAt{9}{314.602}{163.211}{\Character{61}}%
\TextAt{9}{327.012}{163.211}{1}%
\TextAt{13}{332.88}{163.211}{\Character{59}}%
\TextAt{13}{338.104}{163.211}{\Character{58}}%
\TextAt{13}{343.341}{163.211}{\Character{58}}%
\TextAt{13}{348.577}{163.211}{\Character{58}}%
\TextAt{13}{353.802}{163.211}{\Character{59}}%
\TextAt{13}{359.038}{163.211}{s}%
\TextAt{16}{367.272}{163.211}{\Character{0}}%
\TextAt{13}{379.227}{163.211}{r}%
\TextAt{13}{384.177}{163.211}{\Character{58}}%
\TextAt{8}{79.226}{135.257}{This}%
\TextAt{8}{105.946}{135.257}{shows}%
\TextAt{8}{139.564}{135.257}{that}%
\TextAt{9}{164.933}{135.257}{\Character{1}}%
\TextAt{8}{178.585}{135.257}{is}%
\TextAt{8}{190.349}{135.257}{the}%
\TextAt{8}{210.505}{135.257}{kernel}%
\TextAt{8}{244.972}{135.257}{of}%
\TextAt{8}{258.302}{135.257}{the}%
\TextAt{8}{278.459}{135.257}{map}%
\TextAt{13}{304.473}{135.257}{K}%
\TextAt{21}{315.291}{139.595}{S}%
\NativeArrowAt{324.705}{135.257}{11.955}{19.26}{\longrightarrow}%
\TextAt{20}{347.287}{144.223}{\Character{81}}%
\TextAt{21}{358.579}{141.068}{s}%
\TextAt{3}{362.494}{141.068}{\Character{0}}%
\TextAt{21}{369.081}{141.068}{r}%
\TextAt{21}{358.578}{131.77}{i}%
\TextAt{17}{361.462}{131.77}{\Character{61}1}%
\TextAt{13}{375.628}{135.257}{K}%
\TextAt{3}{386.446}{139.595}{\Character{3}}%
\TextAt{27}{385.587}{132.301}{p}%
\TextAt{22}{389.576}{131.086}{i}%
\TextAt{9}{393.236}{135.257}{\Character{47}}%
\TextAt{13}{399.089}{135.257}{K}%
\TextAt{3}{409.907}{139.595}{\Character{3}}%
\TextAt{21}{414.141}{139.595}{n}%
\TextAt{27}{409.048}{132.301}{p}%
\TextAt{22}{413.037}{131.086}{i}%
\TextAt{8}{419.777}{135.257}{.}%
\Vector{q 0 G 0 g 0.399 w 452.334 -638.743 m 452.334 -630.674 l S Q}%
\Vector{q 0 G 0 g 0.399 w 452.534 -630.873 m 459.707 -630.873 l S Q}%
\Vector{q 0 G 0 g 0.399 w 452.534 -638.544 m 459.707 -638.544 l S Q}%
\Vector{q 0 G 0 g 0.399 w 459.906 -638.743 m 459.906 -630.674 l S Q}%
\TextAt{12}{79.226}{100.876}{Theorem}%
\TextAt{12}{136.24}{100.876}{(4.2).}%
\TextAt{8}{171.556}{100.876}{Let}%
\TextAt{13}{192.525}{100.876}{T}%
\TextAt{8}{204.909}{100.876}{be}%
\TextAt{8}{220.834}{100.876}{a}%
\TextAt{8}{230.589}{100.876}{set}%
\TextAt{8}{248.857}{100.876}{of}%
\TextAt{8}{262.187}{100.876}{places}%
\TextAt{8}{296.713}{100.876}{as}%
\TextAt{8}{311.083}{100.876}{in}%
\TextAt{8}{324.736}{100.876}{(4.1),}%
\TextAt{8}{355.963}{100.876}{and}%
\TextAt{8}{378.726}{100.876}{let}%
\TextAt{13}{232.97}{72.921}{C}%
\TextAt{21}{241.345}{71.128}{K}%
\TextAt{9}{249.548}{72.921}{\Character{40}}%
\TextAt{13}{254.1}{72.921}{S\Character{59}}%
\TextAt{13}{266.569}{72.921}{T}%
\TextAt{9}{275.076}{72.921}{\Character{41}}%
\TextAt{9}{282.931}{72.921}{\Character{61}}%
\TextAt{13}{295.375}{72.921}{I}%
\TextAt{21}{300.538}{71.128}{K}%
\TextAt{9}{308.74}{72.921}{\Character{40}}%
\TextAt{13}{313.293}{72.921}{S\Character{59}}%
\TextAt{13}{325.762}{72.921}{T}%
\TextAt{9}{334.269}{72.921}{\Character{41}}%
\TextAt{13}{338.821}{72.921}{K}%
\TextAt{3}{349.639}{77.858}{\Character{3}}%
\TextAt{9}{354.371}{72.921}{\Character{47}}%
\TextAt{13}{360.224}{72.921}{K}%
\TextAt{3}{371.042}{77.858}{\Character{3}}%
\TextAt{13}{375.774}{72.921}{\Character{59}}%
\TextAt{8}{297.211}{43.033}{347}%
\Destination{78.23}{809.815}{nav.page.347}%
\Destination{79.36}{473.855}{nav.print.381}%
\Destination{79.36}{114.855}{nav.number.6.4.2}%
\Hotlink{170.357}{741.855}{13.995}{11.047}{nav.chapter.4}%
\Hotlink{200.8}{707.755}{9.775}{11.047}{nav.chapter.5}%
\Hotlink{197.405}{580.655}{13.995}{11.047}{nav.chapter.4}%
\Hotlink{476.84}{317.255}{9.775}{11.047}{nav.chapter.5}%
\Hotlink{191.173}{741.855}{25.078}{11.047}{nav.number.4.3.6}%
\Hotlink{216.084}{707.755}{25.078}{11.047}{nav.number.5.3.3}%
\Hotlink{217.55}{580.655}{25.078}{11.047}{nav.number.4.3.6}%
\Hotlink{97.592}{524.755}{25.078}{11.047}{nav.number.6.1.1}%
\Hotlink{458.577}{405.455}{25.078}{11.047}{nav.number.6.3.7}%
\Hotlink{492.974}{317.255}{25.078}{11.047}{nav.number.5.3.3}%
\Hotlink{324.236}{98.655}{25.078}{11.047}{nav.number.6.4.1}%
}

% Page 360 (body page 348)
\Page{%
\TextAt{2}{79.226}{796.943}{VI.}%
\TextAt{2}{98.012}{796.943}{Global}%
\TextAt{2}{133.248}{796.943}{Class}%
\TextAt{2}{161.841}{796.943}{Field}%
\TextAt{2}{189.572}{796.943}{Theory}%
\Vector{q 0 G 0 g 0.299 w 7.226 18.729 m 460.769 18.729 l S Q}%
\TextAt{8}{79.226}{761.124}{where}%
\TextAt{13}{203.838}{744.078}{I}%
\TextAt{21}{209.001}{742.284}{K}%
\TextAt{9}{217.204}{744.078}{\Character{40}}%
\TextAt{13}{221.756}{744.078}{S\Character{59}}%
\TextAt{13}{234.225}{744.078}{T}%
\TextAt{9}{242.732}{744.078}{\Character{41}}%
\TextAt{9}{250.587}{744.078}{\Character{61}}%
\TextAt{20}{263.031}{755.435}{\Character{89}}%
\TextAt{27}{263.054}{730.064}{p}%
\NativeMathAt{267.042}{730.064}{7.97}{\in}%
\TextAt{21}{272.688}{730.064}{S}%
\TextAt{13}{280.299}{744.078}{K}%
\TextAt{3}{291.117}{749.014}{\Character{3}}%
\TextAt{21}{295.351}{749.014}{n}%
\TextAt{27}{290.258}{741.122}{p}%
\TextAt{16}{303.644}{744.078}{\Character{2}}%
\TextAt{20}{315.832}{755.435}{\Character{89}}%
\TextAt{27}{315.599}{730.064}{p}%
\NativeMathAt{319.588}{730.064}{7.97}{\in}%
\TextAt{21}{325.234}{730.064}{T}%
\TextAt{13}{333.332}{744.078}{K}%
\TextAt{3}{344.15}{749.014}{\Character{3}}%
\TextAt{27}{343.292}{741.122}{p}%
\TextAt{16}{351.539}{744.078}{\Character{2}}%
\TextAt{20}{369.347}{755.435}{\Character{89}}%
\TextAt{27}{363.494}{729.532}{p}%
\TextAt{17}{368.424}{729.532}{\Character{47}}%
\NativeMathAt{367.483}{729.532}{7.97}{\in}%
\TextAt{21}{373.129}{729.532}{S}%
\TextAt{3}{378.724}{729.532}{\Character{91}}%
\TextAt{21}{384.37}{729.532}{T}%
\TextAt{13}{392.469}{744.078}{U}%
\TextAt{27}{400.419}{742.284}{p}%
\TextAt{13}{404.906}{744.078}{\Character{58}}%
\TextAt{8}{79.226}{712.703}{Then}%
\TextAt{8}{109.784}{712.703}{one}%
\TextAt{8}{131.243}{712.703}{has}%
\TextAt{13}{152.99}{687.533}{N}%
\TextAt{21}{162.364}{685.673}{L}%
\TextAt{3}{168.123}{685.673}{\Character{106}}%
\TextAt{21}{170.475}{685.673}{K}%
\TextAt{13}{178.678}{687.533}{C}%
\TextAt{21}{187.053}{685.74}{L}%
\NativeMathAt{196.631}{687.533}{11.955}{\supseteq}%
\TextAt{13}{209.25}{687.533}{C}%
\TextAt{21}{217.625}{685.74}{K}%
\TextAt{9}{225.828}{687.533}{\Character{40}}%
\TextAt{13}{230.38}{687.533}{S\Character{59}}%
\TextAt{13}{242.849}{687.533}{T}%
\TextAt{9}{251.356}{687.533}{\Character{41}}%
\TextAt{8}{279.819}{687.533}{and}%
\TextAt{9}{322.595}{687.533}{\Character{40}}%
\TextAt{13}{327.147}{687.533}{C}%
\TextAt{21}{335.522}{685.74}{K}%
\TextAt{9}{347.046}{687.533}{\Character{58}}%
\TextAt{13}{353.618}{687.533}{C}%
\TextAt{21}{361.993}{685.74}{K}%
\TextAt{9}{370.196}{687.533}{\Character{40}}%
\TextAt{13}{374.748}{687.533}{S\Character{59}}%
\TextAt{13}{387.217}{687.533}{T}%
\TextAt{9}{395.724}{687.533}{\Character{41}\Character{41}}%
\TextAt{9}{408.122}{687.533}{\Character{61}}%
\TextAt{9}{420.543}{687.533}{\Character{91}}%
\TextAt{13}{423.827}{687.533}{L}%
\TextAt{9}{435.112}{687.533}{\Character{58}}%
\TextAt{13}{441.684}{687.533}{K}%
\TextAt{9}{452.502}{687.533}{\Character{93}}%
\TextAt{13}{455.754}{687.533}{\Character{58}}%
\TextAt{8}{79.226}{662.363}{In}%
\TextAt{8}{93.847}{662.363}{particular,}%
\TextAt{8}{151.089}{662.363}{if}%
\TextAt{13}{161.812}{662.363}{L}%
\TextAt{16}{173.098}{662.363}{\Character{106}}%
\TextAt{13}{179.74}{662.363}{K}%
\TextAt{8}{194.455}{662.363}{is}%
\TextAt{8}{206.218}{662.363}{cyclic,}%
\TextAt{8}{241.654}{662.363}{then}%
\TextAt{13}{268.314}{662.363}{N}%
\TextAt{21}{277.687}{660.503}{L}%
\TextAt{3}{283.446}{660.503}{\Character{106}}%
\TextAt{21}{285.799}{660.503}{K}%
\TextAt{13}{294.002}{662.363}{C}%
\TextAt{21}{302.377}{660.57}{L}%
\TextAt{9}{311.955}{662.363}{\Character{61}}%
\TextAt{13}{324.38}{662.363}{C}%
\TextAt{21}{332.755}{660.57}{K}%
\TextAt{9}{340.958}{662.363}{\Character{40}}%
\TextAt{13}{345.51}{662.363}{S\Character{59}}%
\TextAt{13}{357.979}{662.363}{T}%
\TextAt{9}{366.486}{662.363}{\Character{41}}%
\TextAt{8}{371.038}{662.363}{.}%
\TextAt{12}{79.226}{635.004}{Remark.}%
\TextAt{8}{133.383}{635.004}{It}%
\TextAt{8}{146.056}{635.004}{will}%
\TextAt{8}{168.161}{635.004}{follow}%
\TextAt{8}{201.982}{635.004}{from}%
\TextAt{8}{229.622}{635.004}{(5.5)}%
\TextAt{8}{257.597}{635.004}{that}%
\TextAt{13}{282.966}{635.004}{N}%
\TextAt{21}{292.34}{633.144}{L}%
\TextAt{3}{298.099}{633.144}{\Character{106}}%
\TextAt{21}{300.451}{633.144}{K}%
\TextAt{13}{308.654}{635.004}{C}%
\TextAt{21}{317.029}{633.211}{L}%
\TextAt{9}{326.607}{635.004}{\Character{61}}%
\TextAt{13}{339.033}{635.004}{C}%
\TextAt{21}{347.408}{633.211}{K}%
\TextAt{9}{355.61}{635.004}{\Character{40}}%
\TextAt{13}{360.163}{635.004}{S\Character{59}}%
\TextAt{13}{372.632}{635.004}{T}%
\TextAt{9}{381.139}{635.004}{\Character{41}}%
\TextAt{8}{389.588}{635.004}{also}%
\TextAt{8}{413.068}{635.004}{holds}%
\TextAt{8}{443.697}{635.004}{in}%
\TextAt{8}{457.35}{635.004}{general.}%
\TextAt{8}{97.159}{612.528}{For}%
\TextAt{8}{118.128}{612.528}{the}%
\TextAt{8}{138.285}{612.528}{proof}%
\TextAt{8}{168.854}{612.528}{of}%
\TextAt{8}{182.184}{612.528}{the}%
\TextAt{8}{202.341}{612.528}{theorem}%
\TextAt{8}{247.866}{612.528}{we}%
\TextAt{8}{265.093}{612.528}{need}%
\TextAt{8}{292.399}{612.528}{the}%
\TextAt{8}{312.555}{612.528}{following}%
\TextAt{12}{79.226}{587.61}{Lemma}%
\TextAt{12}{126.891}{587.61}{(4.3).}%
\TextAt{13}{162.207}{587.61}{I}%
\TextAt{21}{167.37}{585.817}{K}%
\TextAt{9}{175.573}{587.61}{\Character{40}}%
\TextAt{13}{180.125}{587.61}{S\Character{59}}%
\TextAt{13}{192.594}{587.61}{T}%
\TextAt{9}{201.101}{587.61}{\Character{41}}%
\TextAt{16}{208.31}{587.61}{\Character{92}}%
\TextAt{13}{218.937}{587.61}{K}%
\TextAt{3}{229.754}{591.948}{\Character{3}}%
\TextAt{9}{237.807}{587.61}{\Character{61}}%
\TextAt{9}{250.217}{587.61}{\Character{40}}%
\TextAt{13}{254.785}{587.61}{K}%
\TextAt{21}{265.603}{591.948}{S}%
\TextAt{3}{271.199}{591.948}{\Character{91}}%
\TextAt{21}{276.844}{591.948}{T}%
\TextAt{9}{283.449}{587.61}{\Character{41}}%
\TextAt{21}{288.001}{591.948}{n}%
\TextAt{8}{293.637}{587.61}{.}%
\TextAt{19}{79.226}{554.82}{Proof.}%
\TextAt{8}{116}{554.82}{The}%
\TextAt{8}{140.971}{554.82}{inclusion}%
\TextAt{9}{190.724}{554.82}{\Character{40}}%
\TextAt{13}{195.277}{554.82}{K}%
\TextAt{21}{206.094}{559.159}{S}%
\TextAt{3}{211.69}{559.159}{\Character{91}}%
\TextAt{21}{217.336}{559.159}{T}%
\TextAt{9}{223.94}{554.82}{\Character{41}}%
\TextAt{21}{228.493}{559.159}{n}%
\NativeMathAt{239.012}{554.82}{11.955}{\subseteq}%
\TextAt{13}{253.194}{554.82}{I}%
\TextAt{21}{258.357}{553.027}{K}%
\TextAt{9}{266.559}{554.82}{\Character{40}}%
\TextAt{13}{271.112}{554.82}{S\Character{59}}%
\TextAt{13}{283.581}{554.82}{T}%
\TextAt{9}{292.087}{554.82}{\Character{41}}%
\TextAt{16}{299.921}{554.82}{\Character{92}}%
\TextAt{13}{311.173}{554.82}{K}%
\TextAt{3}{321.991}{559.159}{\Character{3}}%
\TextAt{8}{331.537}{554.82}{is}%
\TextAt{8}{344.218}{554.82}{trivial.}%
\TextAt{8}{386.321}{554.82}{Let}%
\TextAt{13}{408.207}{554.82}{y}%
\NativeMathAt{419.226}{554.82}{11.955}{\in}%
\TextAt{13}{432.08}{554.82}{I}%
\TextAt{21}{437.243}{553.027}{K}%
\TextAt{9}{445.445}{554.82}{\Character{40}}%
\TextAt{13}{449.998}{554.82}{S\Character{59}}%
\TextAt{13}{462.467}{554.82}{T}%
\TextAt{9}{470.974}{554.82}{\Character{41}}%
\TextAt{16}{478.807}{554.82}{\Character{92}}%
\TextAt{13}{490.059}{554.82}{K}%
\TextAt{3}{500.877}{559.159}{\Character{3}}%
\TextAt{8}{505.609}{554.82}{,}%
\TextAt{8}{513.904}{554.82}{and}%
\TextAt{13}{79.226}{537.774}{M}%
\TextAt{9}{96.388}{537.774}{\Character{61}}%
\TextAt{13}{110.081}{537.774}{K}%
\TextAt{9}{120.898}{537.774}{\Character{40}}%
\TextAt{22}{128.771}{540.401}{n}%
\TextAt{16}{126.731}{545.223}{\Character{112}}%
\Vector{q 0 G 0 g 0.478 w 64.693 -228.538 m 70.83 -228.538 l S Q}%
\TextAt{13}{136.693}{537.774}{y}%
\TextAt{9}{142.83}{537.774}{\Character{41}}%
\TextAt{8}{147.382}{537.774}{.}%
\TextAt{8}{158.061}{537.774}{It}%
\TextAt{8}{171.477}{537.774}{suffices}%
\TextAt{8}{212.008}{537.774}{to}%
\TextAt{8}{227.062}{537.774}{show}%
\TextAt{8}{256.809}{537.774}{that}%
\TextAt{13}{282.921}{537.774}{N}%
\TextAt{21}{292.295}{535.915}{M}%
\TextAt{3}{301.277}{535.915}{\Character{106}}%
\TextAt{21}{303.629}{535.915}{K}%
\TextAt{13}{311.832}{537.774}{C}%
\TextAt{21}{320.207}{535.981}{M}%
\TextAt{9}{334.275}{537.774}{\Character{61}}%
\TextAt{13}{347.967}{537.774}{C}%
\TextAt{21}{356.343}{535.981}{K}%
\TextAt{8}{364.545}{537.774}{,}%
\TextAt{8}{372.624}{537.774}{for}%
\TextAt{8}{391.252}{537.774}{then}%
\TextAt{8}{418.656}{537.774}{(3.6)}%
\TextAt{8}{447.375}{537.774}{implies}%
\TextAt{13}{487.846}{537.774}{M}%
\TextAt{9}{505.007}{537.774}{\Character{61}}%
\TextAt{13}{518.7}{537.774}{K}%
\TextAt{8}{529.518}{537.774}{,}%
\TextAt{8}{79.226}{520.728}{hence}%
\TextAt{13}{112.25}{520.728}{y}%
\NativeMathAt{122.59}{520.728}{11.955}{\in}%
\TextAt{13}{134.763}{520.728}{K}%
\TextAt{3}{145.58}{525.067}{\Character{3}}%
\TextAt{21}{149.815}{525.067}{n}%
\TextAt{16}{158.461}{520.728}{\Character{92}}%
\TextAt{13}{169.44}{520.728}{I}%
\TextAt{21}{174.603}{518.935}{K}%
\TextAt{9}{182.806}{520.728}{\Character{40}}%
\TextAt{13}{187.358}{520.728}{S\Character{59}}%
\TextAt{13}{199.827}{520.728}{T}%
\TextAt{9}{208.334}{520.728}{\Character{41}}%
\NativeMathAt{217.089}{520.728}{11.955}{\subseteq}%
\TextAt{9}{230.591}{520.728}{\Character{40}}%
\TextAt{13}{235.143}{520.728}{K}%
\TextAt{21}{245.961}{525.067}{S}%
\TextAt{3}{251.557}{525.067}{\Character{91}}%
\TextAt{21}{257.202}{525.067}{T}%
\TextAt{9}{263.807}{520.728}{\Character{41}}%
\TextAt{21}{268.359}{525.067}{n}%
\TextAt{8}{273.995}{520.728}{.}%
\TextAt{8}{283.997}{520.728}{Let}%
\TextAt{9}{305.484}{520.728}{\Character{91}}%
\TextAt{13}{308.736}{520.728}{\Character{11}}%
\TextAt{9}{316.257}{520.728}{\Character{93}}%
\NativeMathAt{323.712}{520.728}{11.955}{\in}%
\TextAt{13}{335.885}{520.728}{C}%
\TextAt{21}{344.26}{518.935}{K}%
\TextAt{9}{356.666}{520.728}{\Character{61}}%
\TextAt{13}{369.974}{520.728}{I}%
\TextAt{21}{376.076}{525.067}{S}%
\TextAt{21}{375.136}{517.708}{K}%
\TextAt{13}{383.339}{520.728}{K}%
\TextAt{3}{394.157}{525.067}{\Character{3}}%
\TextAt{9}{398.889}{520.728}{\Character{47}}%
\TextAt{13}{404.742}{520.728}{K}%
\TextAt{3}{415.56}{525.067}{\Character{3}}%
\TextAt{8}{420.292}{520.728}{,}%
\TextAt{8}{428.088}{520.728}{and}%
\TextAt{8}{451.369}{520.728}{let}%
\TextAt{13}{468.791}{520.728}{\Character{11}}%
\NativeMathAt{480.516}{520.728}{11.955}{\in}%
\TextAt{13}{492.689}{520.728}{I}%
\TextAt{21}{498.792}{525.067}{S}%
\TextAt{21}{497.852}{517.708}{K}%
\TextAt{8}{510.47}{520.728}{be}%
\TextAt{8}{526.912}{520.728}{a}%
\TextAt{8}{79.226}{503.682}{representative}%
\TextAt{8}{154.412}{503.682}{of}%
\TextAt{8}{167.742}{503.682}{the}%
\TextAt{8}{187.899}{503.682}{class}%
\TextAt{9}{215.336}{503.682}{\Character{91}}%
\TextAt{13}{218.587}{503.682}{\Character{11}}%
\TextAt{9}{226.109}{503.682}{\Character{93}}%
\TextAt{8}{229.361}{503.682}{.}%
\TextAt{8}{237.809}{503.682}{The}%
\TextAt{8}{261.863}{503.682}{map}%
\TextAt{13}{259.161}{478.513}{K}%
\TextAt{21}{269.979}{483.449}{S}%
\NativeArrowAt{279.394}{478.512}{11.955}{19.26}{\longrightarrow}%
\TextAt{20}{302.208}{489.87}{\Character{89}}%
\TextAt{27}{301.976}{464.498}{p}%
\NativeMathAt{305.964}{464.498}{7.97}{\in}%
\TextAt{21}{311.61}{464.498}{T}%
\TextAt{13}{319.709}{478.513}{U}%
\TextAt{27}{327.659}{476.719}{p}%
\TextAt{9}{332.146}{478.513}{\Character{47}}%
\TextAt{13}{337.999}{478.513}{U}%
\TextAt{21}{347.198}{483.449}{n}%
\TextAt{27}{345.949}{475.557}{p}%
\TextAt{8}{79.226}{442.891}{is}%
\TextAt{8}{93.591}{442.891}{surjective.}%
\TextAt{8}{158.356}{442.891}{For}%
\TextAt{8}{181.926}{442.891}{if}%
\TextAt{9}{195.251}{442.891}{\Character{1}}%
\TextAt{8}{211.504}{442.891}{denotes}%
\TextAt{8}{256.439}{442.891}{its}%
\TextAt{8}{275.358}{442.891}{kernel,}%
\TextAt{8}{316.328}{442.891}{then}%
\TextAt{8}{345.589}{442.891}{obviously}%
\TextAt{13}{399.968}{442.891}{K}%
\TextAt{3}{410.785}{447.23}{\Character{3}}%
\TextAt{21}{415.02}{447.23}{n}%
\TextAt{16}{425.086}{442.891}{\Character{92}}%
\TextAt{9}{437.485}{442.891}{\Character{1}}%
\TextAt{9}{454.976}{442.891}{\Character{61}}%
\TextAt{9}{471.832}{442.891}{\Character{40}}%
\TextAt{13}{476.404}{442.891}{K}%
\TextAt{21}{487.221}{447.23}{S}%
\TextAt{9}{493.315}{442.891}{\Character{41}}%
\TextAt{21}{497.868}{447.23}{n}%
\TextAt{8}{503.504}{442.891}{,}%
\TextAt{8}{513.904}{442.891}{and}%
\TextAt{9}{79.226}{425.845}{\Character{1}}%
\TextAt{13}{88.981}{425.845}{K}%
\TextAt{3}{99.799}{430.184}{\Character{3}}%
\TextAt{21}{104.033}{430.184}{n}%
\TextAt{9}{109.669}{425.845}{\Character{47}}%
\TextAt{13}{115.522}{425.845}{K}%
\TextAt{3}{126.34}{430.184}{\Character{3}}%
\TextAt{21}{130.574}{430.184}{n}%
\NativeMathAt{139.531}{425.845}{11.955}{\cong}%
\TextAt{9}{152.151}{425.845}{\Character{1}\Character{47}\Character{40}}%
\TextAt{13}{172.311}{425.845}{K}%
\TextAt{21}{183.129}{430.184}{S}%
\TextAt{9}{189.223}{425.845}{\Character{41}}%
\TextAt{21}{193.775}{430.184}{n}%
\TextAt{8}{199.411}{425.845}{.}%
\TextAt{8}{207.86}{425.845}{From}%
\TextAt{8}{238.584}{425.845}{(1.1)}%
\TextAt{8}{266.559}{425.845}{and}%
\TextAt{8}{289.322}{425.845}{Kummer}%
\TextAt{8}{338.087}{425.845}{theory,}%
\TextAt{8}{377.109}{425.845}{we}%
\TextAt{8}{394.336}{425.845}{therefore}%
\TextAt{8}{443.436}{425.845}{get}%
\TextAt{8}{181.144}{395.852}{\char35\relax }%
\TextAt{9}{190.899}{395.852}{\Character{40}}%
\TextAt{13}{195.452}{395.852}{K}%
\TextAt{21}{206.269}{400.788}{S}%
\TextAt{9}{212.363}{395.852}{\Character{47}\Character{1}\Character{41}}%
\TextAt{9}{235.807}{395.852}{\Character{61}}%
\TextAt{8}{249.465}{403.939}{\char35\relax }%
\TextAt{13}{259.221}{403.939}{K}%
\TextAt{21}{270.038}{408.278}{S}%
\TextAt{9}{276.132}{403.939}{\Character{47}\Character{40}}%
\TextAt{13}{286.538}{403.939}{K}%
\TextAt{21}{297.355}{408.278}{S}%
\TextAt{9}{303.449}{403.939}{\Character{41}}%
\TextAt{21}{308.001}{408.278}{n}%
\Vector{q 0 G 0 g 0.478 w 177.465 -375.16 m 241.638 -375.16 l S Q}%
\TextAt{8}{253.044}{387.651}{\char35\relax }%
\TextAt{9}{262.799}{387.651}{\Character{1}\Character{47}\Character{40}}%
\TextAt{13}{282.959}{387.651}{K}%
\TextAt{21}{293.777}{391.105}{S}%
\TextAt{9}{299.871}{387.651}{\Character{41}}%
\TextAt{21}{304.423}{391.105}{n}%
\TextAt{9}{318.154}{395.852}{\Character{61}}%
\TextAt{13}{354.494}{403.939}{n}%
\TextAt{21}{361.481}{408.278}{s}%
\Vector{q 0 G 0 g 0.478 w 259.775 -375.16 m 316.614 -375.16 l S Q}%
\TextAt{8}{331.775}{387.651}{\char35\relax }%
\TextAt{13}{341.531}{387.651}{G}%
\TextAt{9}{350.764}{387.651}{\Character{40}}%
\TextAt{13}{355.316}{387.651}{L}%
\TextAt{16}{366.602}{387.651}{\Character{106}}%
\TextAt{13}{373.244}{387.651}{K}%
\TextAt{9}{384.061}{387.651}{\Character{41}}%
\TextAt{9}{393.13}{395.852}{\Character{61}}%
\TextAt{13}{405.555}{395.852}{n}%
\TextAt{21}{412.543}{400.788}{s}%
\TextAt{3}{416.459}{400.788}{\Character{0}}%
\TextAt{21}{423.045}{400.788}{r}%
\TextAt{13}{427.6}{395.852}{\Character{58}}%
\TextAt{8}{79.226}{366.576}{This}%
\TextAt{8}{106.806}{366.576}{is}%
\TextAt{8}{119.429}{366.576}{also}%
\TextAt{8}{143.769}{366.576}{the}%
\TextAt{8}{164.785}{366.576}{order}%
\TextAt{8}{196.214}{366.576}{of}%
\TextAt{8}{210.403}{366.576}{the}%
\TextAt{8}{231.419}{366.576}{product}%
\TextAt{8}{276.178}{366.576}{because}%
\TextAt{8}{320.34}{366.576}{by}%
\TextAt{8}{337.458}{366.576}{chap.}%
\TextAt{8}{369.21}{366.576}{II,}%
\TextAt{8}{385.982}{366.576}{(5.8),}%
\TextAt{8}{418.283}{366.576}{we}%
\TextAt{8}{436.37}{366.576}{have}%
\TextAt{8}{464.225}{366.576}{\char35\relax }%
\TextAt{13}{473.98}{366.576}{U}%
\TextAt{27}{481.931}{364.782}{p}%
\TextAt{9}{486.417}{366.576}{\Character{47}}%
\TextAt{13}{492.27}{366.576}{U}%
\TextAt{21}{501.469}{370.914}{n}%
\TextAt{27}{500.221}{363.62}{p}%
\TextAt{9}{511.892}{366.576}{\Character{61}}%
\TextAt{13}{525.782}{366.576}{n}%
\TextAt{8}{79.226}{349.53}{since}%
\TextAt{14}{108.555}{349.53}{p}%
\TextAt{18}{118.986}{349.53}{\Character{45}}%
\TextAt{13}{126.753}{349.53}{n}%
\TextAt{8}{133.741}{349.53}{.}%
\TextAt{8}{144.172}{349.53}{We}%
\TextAt{8}{164.977}{349.53}{thus}%
\TextAt{8}{191.39}{349.53}{find}%
\TextAt{8}{215.459}{349.53}{an}%
\TextAt{8}{232.379}{349.53}{element}%
\TextAt{13}{276.281}{349.53}{x}%
\NativeMathAt{287.38}{349.53}{11.955}{\in}%
\TextAt{13}{299.797}{349.53}{K}%
\TextAt{21}{310.615}{353.868}{S}%
\TextAt{8}{321.267}{349.53}{such}%
\TextAt{8}{348.325}{349.53}{that}%
\TextAt{13}{374.355}{349.53}{\Character{11}}%
\TextAt{27}{381.8}{347.736}{p}%
\TextAt{9}{390.734}{349.53}{\Character{61}}%
\TextAt{13}{404.286}{349.53}{xu}%
\TextAt{21}{417.6}{353.868}{n}%
\TextAt{27}{417.6}{346.574}{p}%
\TextAt{8}{423.237}{349.53}{,}%
\TextAt{13}{431.212}{349.53}{u}%
\TextAt{27}{437.875}{347.736}{p}%
\NativeMathAt{446.808}{349.53}{11.955}{\in}%
\TextAt{13}{459.226}{349.53}{U}%
\TextAt{27}{467.176}{347.736}{p}%
\TextAt{8}{471.663}{349.53}{,}%
\TextAt{8}{479.638}{349.53}{for}%
\TextAt{14}{498.184}{349.53}{p}%
\NativeMathAt{508.614}{349.53}{11.955}{\in}%
\TextAt{13}{521.031}{349.53}{T}%
\TextAt{8}{529.518}{349.53}{.}%
\TextAt{8}{79.226}{332.484}{The}%
\TextAt{8}{103.187}{332.484}{idèle}%
\TextAt{13}{130.4}{332.484}{\Character{11}}%
\TextAt{3}{137.922}{336.822}{\Character{48}}%
\TextAt{9}{144.038}{332.484}{\Character{61}}%
\TextAt{13}{156.463}{332.484}{\Character{11}x}%
\TextAt{3}{170.637}{336.822}{\Character{0}}%
\TextAt{17}{177.224}{336.822}{1}%
\TextAt{8}{185.761}{332.484}{belongs}%
\TextAt{8}{227.679}{332.484}{to}%
\TextAt{8}{241.896}{332.484}{the}%
\TextAt{8}{261.96}{332.484}{same}%
\TextAt{8}{291.194}{332.484}{class}%
\TextAt{8}{318.538}{332.484}{as}%
\TextAt{13}{332.816}{332.484}{\Character{11}}%
\TextAt{8}{340.338}{332.484}{,}%
\TextAt{8}{347.413}{332.484}{and}%
\TextAt{8}{370.083}{332.484}{we}%
\TextAt{8}{387.218}{332.484}{show}%
\TextAt{8}{416.128}{332.484}{that}%
\TextAt{13}{441.404}{332.484}{\Character{11}}%
\TextAt{3}{448.926}{336.822}{\Character{48}}%
\NativeMathAt{455.042}{332.484}{11.955}{\in}%
\TextAt{13}{466.333}{332.484}{N}%
\TextAt{21}{475.707}{330.624}{M}%
\TextAt{3}{484.689}{330.624}{\Character{106}}%
\TextAt{21}{487.041}{330.624}{K}%
\TextAt{13}{495.244}{332.484}{I}%
\TextAt{21}{500.406}{330.69}{M}%
\TextAt{8}{509.886}{332.484}{.}%
\TextAt{8}{518.304}{332.484}{By}%
\TextAt{8}{79.226}{315.438}{(3.2),}%
\TextAt{8}{111.14}{315.438}{this}%
\TextAt{8}{134.512}{315.438}{amounts}%
\TextAt{8}{182.285}{315.438}{to}%
\TextAt{8}{197.145}{315.438}{checking}%
\TextAt{8}{244.846}{315.438}{that}%
\TextAt{8}{270.765}{315.438}{every}%
\TextAt{8}{302.207}{315.438}{component}%
\TextAt{13}{362.592}{315.438}{\Character{11}}%
\TextAt{3}{370.114}{319.776}{\Character{48}}%
\TextAt{27}{370.037}{312.482}{p}%
\TextAt{8}{378.971}{315.438}{is}%
\TextAt{8}{391.285}{315.438}{a}%
\TextAt{8}{401.59}{315.438}{norm}%
\TextAt{8}{432.709}{315.438}{from}%
\TextAt{13}{460.899}{315.438}{M}%
\TextAt{27}{472.224}{313.644}{P}%
\TextAt{16}{483.575}{315.438}{\Character{106}}%
\TextAt{13}{491.154}{315.438}{K}%
\TextAt{27}{501.113}{313.644}{p}%
\TextAt{8}{505.6}{315.438}{.}%
\TextAt{8}{515.697}{315.438}{For}%
\TextAt{14}{79.226}{298.392}{p}%
\NativeMathAt{89.348}{298.392}{11.955}{\in}%
\TextAt{13}{101.457}{298.392}{S}%
\TextAt{8}{113.73}{298.392}{this}%
\TextAt{8}{137.033}{298.392}{holds}%
\TextAt{8}{168.142}{298.392}{because}%
\TextAt{13}{211.924}{298.392}{y}%
\NativeMathAt{222.2}{298.392}{11.955}{\in}%
\TextAt{13}{234.309}{298.392}{K}%
\TextAt{3}{245.126}{302.73}{\Character{3}}%
\TextAt{21}{249.361}{302.73}{n}%
\TextAt{27}{244.268}{295.436}{p}%
\TextAt{8}{254.997}{298.392}{.}%
\TextAt{8}{264.886}{298.392}{Hence}%
\TextAt{8}{300.143}{298.392}{we}%
\TextAt{8}{317.851}{298.392}{have}%
\TextAt{13}{345.326}{298.392}{M}%
\TextAt{27}{356.651}{296.598}{P}%
\TextAt{9}{367.883}{298.392}{\Character{61}}%
\TextAt{13}{381.127}{298.392}{K}%
\TextAt{27}{391.086}{296.598}{p}%
\TextAt{8}{395.573}{298.392}{.}%
\TextAt{8}{405.462}{298.392}{For}%
\TextAt{14}{426.911}{298.392}{p}%
\NativeMathAt{437.033}{298.392}{11.955}{\in}%
\TextAt{13}{449.143}{298.392}{T}%
\TextAt{8}{462.007}{298.392}{since}%
\TextAt{13}{491.156}{298.392}{\Character{11}}%
\TextAt{3}{498.677}{302.73}{\Character{48}}%
\TextAt{27}{498.601}{295.436}{p}%
\TextAt{9}{507.227}{298.392}{\Character{61}}%
\TextAt{13}{520.471}{298.392}{u}%
\TextAt{21}{527.133}{302.73}{n}%
\TextAt{27}{527.133}{295.436}{p}%
\TextAt{8}{79.226}{281.346}{is}%
\TextAt{8}{91.077}{281.346}{an}%
\TextAt{13}{107.423}{281.346}{n}%
\TextAt{8}{114.411}{281.346}{-th}%
\TextAt{8}{133.351}{281.346}{power.}%
\TextAt{8}{172.307}{281.346}{For}%
\TextAt{14}{193.363}{281.346}{p}%
\NativeMathAt{202.815}{281.346}{11.955}{\notin}%
\TextAt{13}{214.254}{281.346}{S}%
\TextAt{16}{224.865}{281.346}{\Character{91}}%
\TextAt{13}{235.551}{281.346}{T}%
\TextAt{8}{248.022}{281.346}{it}%
\TextAt{8}{259.814}{281.346}{holds}%
\TextAt{8}{290.53}{281.346}{because}%
\TextAt{13}{333.918}{281.346}{\Character{11}}%
\TextAt{3}{341.44}{285.684}{\Character{48}}%
\TextAt{27}{341.364}{278.39}{p}%
\TextAt{8}{349.835}{281.346}{is}%
\TextAt{8}{361.686}{281.346}{a}%
\TextAt{8}{371.528}{281.346}{unit}%
\TextAt{8}{396.326}{281.346}{and}%
\TextAt{13}{419.176}{281.346}{M}%
\TextAt{27}{430.501}{279.553}{P}%
\TextAt{16}{441.063}{281.346}{\Character{106}}%
\TextAt{13}{447.853}{281.346}{K}%
\TextAt{27}{457.812}{279.553}{p}%
\TextAt{8}{466.284}{281.346}{is}%
\TextAt{8}{478.134}{281.346}{unramified}%
\TextAt{8}{79.226}{264.3}{(see}%
\TextAt{8}{102.694}{264.3}{chap.}%
\TextAt{8}{133.586}{264.3}{V,}%
\TextAt{8}{149.511}{264.3}{(3.3)).}%
\TextAt{8}{186.592}{264.3}{This}%
\TextAt{8}{213.311}{264.3}{is}%
\TextAt{8}{225.075}{264.3}{why}%
\TextAt{9}{249.787}{264.3}{\Character{91}}%
\TextAt{13}{253.038}{264.3}{\Character{11}}%
\TextAt{9}{260.56}{264.3}{\Character{93}}%
\NativeMathAt{267.132}{264.3}{11.955}{\in}%
\TextAt{13}{278.423}{264.3}{N}%
\TextAt{21}{287.797}{262.44}{M}%
\TextAt{3}{296.779}{262.44}{\Character{106}}%
\TextAt{21}{299.131}{262.44}{K}%
\TextAt{13}{307.334}{264.3}{C}%
\TextAt{21}{315.709}{262.507}{M}%
\TextAt{8}{325.189}{264.3}{,}%
\TextAt{8}{332.338}{264.3}{q.e.d.}%
\Vector{q 0 G 0 g 0.399 w 452.334 -509.7 m 452.334 -501.63 l S Q}%
\Vector{q 0 G 0 g 0.399 w 452.534 -501.83 m 459.707 -501.83 l S Q}%
\Vector{q 0 G 0 g 0.399 w 452.534 -509.501 m 459.707 -509.501 l S Q}%
\Vector{q 0 G 0 g 0.399 w 459.906 -509.7 m 459.906 -501.63 l S Q}%
\TextAt{12}{79.226}{236.393}{Proof}%
\TextAt{12}{116.347}{236.393}{of}%
\TextAt{12}{131.673}{236.393}{theorem}%
\TextAt{12}{184.575}{236.393}{(4.2):}%
\TextAt{19}{79.226}{208.487}{Proof.}%
\TextAt{8}{116}{208.487}{The}%
\TextAt{8}{140.054}{208.487}{identity}%
\TextAt{9}{183.308}{208.487}{\Character{40}}%
\TextAt{13}{187.86}{208.487}{C}%
\TextAt{21}{196.235}{206.694}{K}%
\TextAt{9}{207.759}{208.487}{\Character{58}}%
\TextAt{13}{214.331}{208.487}{C}%
\TextAt{21}{222.706}{206.694}{K}%
\TextAt{9}{230.909}{208.487}{\Character{40}}%
\TextAt{13}{235.461}{208.487}{S\Character{59}}%
\TextAt{13}{247.93}{208.487}{T}%
\TextAt{9}{256.437}{208.487}{\Character{41}\Character{41}}%
\TextAt{9}{268.834}{208.487}{\Character{61}}%
\TextAt{9}{281.256}{208.487}{\Character{91}}%
\TextAt{13}{284.54}{208.487}{L}%
\TextAt{9}{295.825}{208.487}{\Character{58}}%
\TextAt{13}{302.398}{208.487}{K}%
\TextAt{9}{313.215}{208.487}{\Character{93}}%
\TextAt{8}{320.364}{208.487}{follows}%
\TextAt{8}{358.8}{208.487}{from}%
\TextAt{8}{386.441}{208.487}{the}%
\TextAt{8}{406.597}{208.487}{exact}%
\TextAt{8}{437.489}{208.487}{sequence}%
\TextAt{9}{89.189}{175.517}{1}%
\NativeArrowAt{98.363}{175.517}{11.955}{19.26}{\longrightarrow}%
\TextAt{13}{131.889}{183.604}{I}%
\TextAt{21}{137.992}{187.943}{S}%
\TextAt{3}{143.587}{187.943}{\Character{91}}%
\TextAt{21}{149.233}{187.943}{T}%
\TextAt{21}{137.052}{180.584}{K}%
\TextAt{16}{158.494}{183.604}{\Character{92}}%
\TextAt{13}{169.121}{183.604}{K}%
\TextAt{3}{179.939}{187.943}{\Character{3}}%
\Vector{q 0 G 0 g 0.478 w 50.14 -595.495 m 122.42 -595.495 l S Q}%
\TextAt{13}{122.14}{167.316}{I}%
\TextAt{21}{127.303}{165.523}{K}%
\TextAt{9}{135.506}{167.316}{\Character{40}}%
\TextAt{13}{140.058}{167.316}{S\Character{59}}%
\TextAt{13}{152.527}{167.316}{T}%
\TextAt{9}{161.034}{167.316}{\Character{41}}%
\TextAt{16}{168.243}{167.316}{\Character{92}}%
\TextAt{13}{178.87}{167.316}{K}%
\TextAt{3}{189.687}{170.77}{\Character{3}}%
\NativeArrowAt{198.936}{175.517}{11.955}{19.26}{\longrightarrow}%
\TextAt{13}{221.518}{175.517}{I}%
\TextAt{21}{227.621}{180.453}{S}%
\TextAt{3}{233.217}{180.453}{\Character{91}}%
\TextAt{21}{238.862}{180.453}{T}%
\TextAt{21}{226.681}{172.561}{K}%
\TextAt{9}{245.467}{175.517}{\Character{47}}%
\TextAt{13}{251.32}{175.517}{I}%
\TextAt{21}{256.483}{173.723}{K}%
\TextAt{9}{264.685}{175.517}{\Character{40}}%
\TextAt{13}{269.238}{175.517}{S\Character{59}}%
\TextAt{13}{281.707}{175.517}{T}%
\TextAt{9}{290.214}{175.517}{\Character{41}}%
\NativeArrowAt{360.87}{148.41}{11.955}{19.26}{\longrightarrow}%
\TextAt{13}{383.452}{148.41}{I}%
\TextAt{21}{389.554}{153.346}{S}%
\TextAt{3}{395.15}{153.346}{\Character{91}}%
\TextAt{21}{400.796}{153.346}{T}%
\TextAt{21}{388.614}{145.454}{K}%
\TextAt{13}{407.4}{148.41}{K}%
\TextAt{3}{418.218}{153.346}{\Character{3}}%
\TextAt{9}{422.95}{148.41}{\Character{47}}%
\TextAt{13}{428.803}{148.41}{I}%
\TextAt{21}{433.966}{146.617}{K}%
\TextAt{9}{442.169}{148.41}{\Character{40}}%
\TextAt{13}{446.721}{148.41}{S\Character{59}}%
\TextAt{13}{459.19}{148.41}{T}%
\TextAt{9}{467.697}{148.41}{\Character{41}}%
\TextAt{13}{472.249}{148.41}{K}%
\TextAt{3}{483.067}{153.346}{\Character{3}}%
\NativeArrowAt{491.12}{148.41}{11.955}{19.26}{\longrightarrow}%
\TextAt{9}{513.702}{148.41}{1}%
\TextAt{13}{519.555}{148.41}{\Character{58}}%
\TextAt{8}{79.226}{123.24}{Since}%
\TextAt{13}{109.784}{123.24}{I}%
\TextAt{21}{115.886}{127.578}{S}%
\TextAt{3}{121.482}{127.578}{\Character{91}}%
\TextAt{21}{127.128}{127.578}{T}%
\TextAt{21}{114.946}{120.219}{K}%
\TextAt{13}{133.732}{123.24}{K}%
\TextAt{3}{144.55}{127.578}{\Character{3}}%
\TextAt{9}{152.603}{123.24}{\Character{61}}%
\TextAt{13}{165.028}{123.24}{I}%
\TextAt{21}{170.191}{121.447}{K}%
\TextAt{8}{178.394}{123.24}{,}%
\TextAt{8}{185.543}{123.24}{the}%
\TextAt{8}{205.7}{123.24}{order}%
\TextAt{8}{236.269}{123.24}{of}%
\TextAt{8}{249.599}{123.24}{the}%
\TextAt{8}{269.755}{123.24}{group}%
\TextAt{8}{302.931}{123.24}{on}%
\TextAt{8}{319.19}{123.24}{the}%
\TextAt{8}{339.347}{123.24}{right}%
\TextAt{8}{367.644}{123.24}{is}%
\TextAt{9}{164.52}{98.07}{\Character{40}}%
\TextAt{13}{169.073}{98.07}{I}%
\TextAt{21}{175.176}{103.006}{S}%
\TextAt{3}{180.771}{103.006}{\Character{91}}%
\TextAt{21}{186.417}{103.006}{T}%
\TextAt{21}{174.236}{95.114}{K}%
\TextAt{13}{193.021}{98.07}{K}%
\TextAt{3}{203.839}{103.006}{\Character{3}}%
\TextAt{9}{211.892}{98.07}{\Character{58}}%
\TextAt{13}{218.465}{98.07}{I}%
\TextAt{21}{223.628}{96.277}{K}%
\TextAt{9}{231.83}{98.07}{\Character{40}}%
\TextAt{13}{236.383}{98.07}{S\Character{59}}%
\TextAt{13}{248.852}{98.07}{T}%
\TextAt{9}{257.359}{98.07}{\Character{41}}%
\TextAt{13}{261.911}{98.07}{K}%
\TextAt{3}{272.729}{103.006}{\Character{3}}%
\TextAt{9}{277.461}{98.07}{\Character{41}}%
\TextAt{9}{285.316}{98.07}{\Character{61}}%
\TextAt{9}{297.737}{98.07}{\Character{40}}%
\TextAt{13}{302.312}{98.07}{I}%
\TextAt{21}{307.475}{96.277}{K}%
\TextAt{13}{315.677}{98.07}{K}%
\TextAt{3}{326.495}{103.006}{\Character{3}}%
\TextAt{9}{331.227}{98.07}{\Character{47}}%
\TextAt{13}{337.08}{98.07}{K}%
\TextAt{3}{347.898}{103.006}{\Character{3}}%
\TextAt{9}{355.951}{98.07}{\Character{58}}%
\TextAt{13}{362.524}{98.07}{I}%
\TextAt{21}{367.687}{96.277}{K}%
\TextAt{9}{375.889}{98.07}{\Character{40}}%
\TextAt{13}{380.442}{98.07}{S\Character{59}}%
\TextAt{13}{392.911}{98.07}{T}%
\TextAt{9}{401.418}{98.07}{\Character{41}}%
\TextAt{13}{405.97}{98.07}{K}%
\TextAt{3}{416.788}{103.006}{\Character{3}}%
\TextAt{9}{421.52}{98.07}{\Character{47}}%
\TextAt{13}{427.373}{98.07}{K}%
\TextAt{3}{438.191}{103.006}{\Character{3}}%
\TextAt{9}{442.923}{98.07}{\Character{41}}%
\TextAt{9}{285.334}{78.035}{\Character{61}}%
\TextAt{9}{297.744}{78.035}{\Character{40}}%
\TextAt{13}{302.312}{78.035}{C}%
\TextAt{21}{310.687}{76.242}{K}%
\TextAt{9}{322.21}{78.035}{\Character{58}}%
\TextAt{13}{328.783}{78.035}{C}%
\TextAt{21}{337.158}{76.242}{K}%
\TextAt{9}{345.361}{78.035}{\Character{40}}%
\TextAt{13}{349.913}{78.035}{S\Character{59}}%
\TextAt{13}{362.382}{78.035}{T}%
\TextAt{9}{370.889}{78.035}{\Character{41}\Character{41}}%
\TextAt{13}{379.994}{78.035}{\Character{58}}%
\TextAt{8}{297.211}{43.033}{348}%
\Destination{78.23}{809.815}{nav.page.348}%
\Destination{79.36}{552.755}{nav.print.382}%
\Destination{79.36}{601.655}{nav.number.6.4.3}%
\Hotlink{368.71}{364.355}{9.763}{11.047}{nav.chapter.2}%
\Hotlink{133.086}{262.055}{9.775}{11.047}{nav.chapter.5}%
\Hotlink{229.122}{632.755}{25.078}{11.047}{nav.number.6.5.5}%
\Hotlink{418.156}{535.555}{25.078}{11.047}{nav.number.6.3.6}%
\Hotlink{238.084}{423.655}{25.078}{11.047}{nav.number.6.1.1}%
\Hotlink{385.482}{364.355}{25.078}{11.047}{nav.number.2.5.8}%
\Hotlink{78.726}{313.255}{25.078}{11.047}{nav.number.6.3.2}%
\Hotlink{149.011}{262.055}{25.078}{11.047}{nav.number.5.3.3}%
\Hotlink{184.075}{234.155}{28.676}{11.047}{nav.number.6.4.2}%
}

% Page 361 (body page 349)
\Page{%
\TextAt{2}{79.226}{796.943}{VI.}%
\TextAt{2}{98.012}{796.943}{Global}%
\TextAt{2}{133.248}{796.943}{Class}%
\TextAt{2}{161.841}{796.943}{Field}%
\TextAt{2}{189.572}{796.943}{Theory}%
\Vector{q 0 G 0 g 0.299 w 7.226 18.729 m 460.769 18.729 l S Q}%
\TextAt{8}{79.226}{761.124}{The}%
\TextAt{8}{103.28}{761.124}{order}%
\TextAt{8}{133.849}{761.124}{of}%
\TextAt{8}{147.179}{761.124}{the}%
\TextAt{8}{167.336}{761.124}{group}%
\TextAt{8}{200.511}{761.124}{on}%
\TextAt{8}{216.77}{761.124}{the}%
\TextAt{8}{236.927}{761.124}{left}%
\TextAt{8}{257.406}{761.124}{is}%
\TextAt{9}{159.549}{732.123}{\Character{40}}%
\TextAt{13}{164.102}{732.123}{I}%
\TextAt{21}{170.204}{737.059}{S}%
\TextAt{3}{175.8}{737.059}{\Character{91}}%
\TextAt{21}{181.446}{737.059}{T}%
\TextAt{21}{169.264}{729.167}{K}%
\TextAt{16}{190.707}{732.123}{\Character{92}}%
\TextAt{13}{201.334}{732.123}{K}%
\TextAt{3}{212.151}{737.059}{\Character{3}}%
\TextAt{9}{220.204}{732.123}{\Character{58}}%
\TextAt{13}{226.777}{732.123}{I}%
\TextAt{21}{231.94}{730.329}{K}%
\TextAt{9}{240.143}{732.123}{\Character{40}}%
\TextAt{13}{244.695}{732.123}{S\Character{59}}%
\TextAt{13}{257.164}{732.123}{T}%
\TextAt{9}{265.671}{732.123}{\Character{41}}%
\TextAt{16}{272.88}{732.123}{\Character{92}}%
\TextAt{13}{283.507}{732.123}{K}%
\TextAt{3}{294.324}{737.059}{\Character{3}}%
\TextAt{9}{299.057}{732.123}{\Character{41}}%
\TextAt{9}{306.912}{732.123}{\Character{61}}%
\TextAt{9}{319.333}{732.123}{\Character{40}}%
\TextAt{13}{323.908}{732.123}{K}%
\TextAt{21}{334.725}{737.059}{S}%
\TextAt{3}{340.321}{737.059}{\Character{91}}%
\TextAt{21}{345.967}{737.059}{T}%
\TextAt{9}{355.892}{732.123}{\Character{58}}%
\TextAt{9}{362.443}{732.123}{\Character{40}}%
\TextAt{13}{367.017}{732.123}{K}%
\TextAt{21}{377.834}{737.059}{S}%
\TextAt{3}{383.43}{737.059}{\Character{91}}%
\TextAt{21}{389.076}{737.059}{T}%
\TextAt{9}{395.68}{732.123}{\Character{41}}%
\TextAt{21}{400.232}{737.059}{n}%
\TextAt{9}{405.869}{732.123}{\Character{41}}%
\TextAt{9}{413.724}{732.123}{\Character{61}}%
\TextAt{13}{426.167}{732.123}{n}%
\TextAt{17}{433.155}{737.059}{2}%
\TextAt{21}{437.389}{737.059}{s}%
\TextAt{3}{441.305}{737.059}{\Character{0}}%
\TextAt{21}{447.892}{737.059}{r}%
\TextAt{8}{79.226}{703.121}{because}%
\TextAt{8}{122.746}{703.121}{\char35\relax }%
\TextAt{9}{132.501}{703.121}{\Character{40}}%
\TextAt{13}{137.054}{703.121}{S}%
\TextAt{16}{147.755}{703.121}{\Character{91}}%
\TextAt{13}{158.53}{703.121}{T}%
\TextAt{9}{167.017}{703.121}{\Character{41}}%
\TextAt{9}{175.242}{703.121}{\Character{61}}%
\TextAt{9}{188.034}{703.121}{2}%
\TextAt{13}{193.913}{703.121}{s}%
\TextAt{16}{202.232}{703.121}{\Character{0}}%
\TextAt{13}{214.336}{703.121}{r}%
\TextAt{8}{219.937}{703.121}{,}%
\TextAt{8}{227.359}{703.121}{and}%
\TextAt{13}{250.34}{703.121}{\Character{22}}%
\TextAt{21}{257.383}{701.328}{n}%
\NativeMathAt{266.712}{703.121}{11.955}{\subseteq}%
\TextAt{13}{279.703}{703.121}{K}%
\TextAt{21}{290.521}{707.46}{S}%
\TextAt{3}{296.117}{707.46}{\Character{91}}%
\TextAt{21}{301.762}{707.46}{T}%
\TextAt{8}{308.367}{703.121}{.}%
\TextAt{8}{317.47}{703.121}{In}%
\TextAt{8}{332.309}{703.121}{view}%
\TextAt{8}{359.511}{703.121}{of}%
\TextAt{8}{373.059}{703.121}{chap.}%
\TextAt{8}{404.169}{703.121}{II,}%
\TextAt{8}{420.3}{703.121}{(5.8),}%
\TextAt{8}{451.8}{703.121}{the}%
\TextAt{8}{472.174}{703.121}{order}%
\TextAt{8}{502.962}{703.121}{of}%
\TextAt{8}{516.51}{703.121}{the}%
\TextAt{8}{79.226}{686.075}{group}%
\TextAt{8}{112.402}{686.075}{in}%
\TextAt{8}{126.055}{686.075}{the}%
\TextAt{8}{146.211}{686.075}{middle}%
\TextAt{8}{184.575}{686.075}{is}%
\TextAt{9}{137.798}{653.111}{\Character{40}}%
\TextAt{13}{142.351}{653.111}{I}%
\TextAt{21}{148.454}{658.047}{S}%
\TextAt{3}{154.049}{658.047}{\Character{91}}%
\TextAt{21}{159.695}{658.047}{T}%
\TextAt{21}{147.514}{650.155}{K}%
\TextAt{9}{169.62}{653.111}{\Character{58}}%
\TextAt{13}{176.193}{653.111}{I}%
\TextAt{21}{181.356}{651.317}{K}%
\TextAt{9}{189.558}{653.111}{\Character{40}}%
\TextAt{13}{194.111}{653.111}{S\Character{59}}%
\TextAt{13}{206.58}{653.111}{T}%
\TextAt{9}{215.087}{653.111}{\Character{41}\Character{41}}%
\TextAt{9}{227.485}{653.111}{\Character{61}}%
\TextAt{20}{239.938}{664.468}{\Character{89}}%
\TextAt{27}{239.961}{639.097}{p}%
\NativeMathAt{243.949}{639.097}{7.97}{\in}%
\TextAt{21}{249.595}{639.097}{S}%
\TextAt{9}{255.214}{653.111}{\Character{40}}%
\TextAt{13}{259.766}{653.111}{K}%
\TextAt{3}{270.584}{658.047}{\Character{3}}%
\TextAt{27}{269.725}{650.155}{p}%
\TextAt{9}{278.637}{653.111}{\Character{58}}%
\TextAt{13}{285.209}{653.111}{K}%
\TextAt{3}{296.027}{658.047}{\Character{3}}%
\TextAt{21}{300.261}{658.047}{n}%
\TextAt{27}{295.168}{650.155}{p}%
\TextAt{9}{305.897}{653.111}{\Character{41}}%
\TextAt{9}{313.752}{653.111}{\Character{61}}%
\TextAt{20}{326.196}{664.468}{\Character{89}}%
\TextAt{27}{326.219}{639.097}{p}%
\NativeMathAt{330.208}{639.097}{7.97}{\in}%
\TextAt{21}{335.853}{639.097}{S}%
\TextAt{13}{347.858}{661.198}{n}%
\TextAt{17}{354.846}{665.537}{2}%
\Vector{q 0 G 0 g 0.478 w 272.66 -117.901 m 290.776 -117.901 l S Q}%
\TextAt{16}{344.66}{644.91}{\Character{106}}%
\TextAt{13}{347.981}{644.91}{n}%
\TextAt{16}{354.969}{644.91}{\Character{106}}%
\TextAt{27}{358.29}{643.117}{p}%
\TextAt{9}{367.293}{653.111}{\Character{61}}%
\TextAt{13}{379.718}{653.111}{n}%
\TextAt{17}{386.706}{658.047}{2}%
\TextAt{21}{390.94}{658.047}{s}%
\TextAt{20}{397.347}{664.468}{\Character{89}}%
\TextAt{27}{402.99}{639.362}{p}%
\TextAt{16}{414.615}{653.111}{\Character{106}}%
\TextAt{13}{417.936}{653.111}{n}%
\TextAt{16}{424.924}{653.111}{\Character{106}}%
\TextAt{3}{428.245}{658.047}{\Character{0}}%
\TextAt{17}{434.831}{658.047}{1}%
\TextAt{27}{428.245}{650.155}{p}%
\TextAt{9}{442.884}{653.111}{\Character{61}}%
\TextAt{13}{455.31}{653.111}{n}%
\TextAt{17}{462.297}{658.047}{2}%
\TextAt{21}{466.532}{658.047}{s}%
\TextAt{13}{470.946}{653.111}{\Character{58}}%
\TextAt{8}{79.226}{614.907}{Altogether}%
\TextAt{8}{137.436}{614.907}{this}%
\TextAt{8}{160.258}{614.907}{gives}%
\TextAt{9}{201.94}{581.942}{\Character{40}}%
\TextAt{13}{206.492}{581.942}{C}%
\TextAt{21}{214.867}{580.149}{K}%
\TextAt{9}{226.391}{581.942}{\Character{58}}%
\TextAt{13}{232.963}{581.942}{C}%
\TextAt{21}{241.338}{580.149}{K}%
\TextAt{9}{249.541}{581.942}{\Character{40}}%
\TextAt{13}{254.093}{581.942}{S\Character{59}}%
\TextAt{13}{266.562}{581.942}{T}%
\TextAt{9}{275.069}{581.942}{\Character{41}\Character{41}}%
\TextAt{9}{287.466}{581.942}{\Character{61}}%
\TextAt{13}{306.437}{590.03}{n}%
\TextAt{17}{313.425}{594.368}{2}%
\TextAt{21}{317.659}{594.368}{s}%
\Vector{q 0 G 0 g 0.478 w 229.116 -189.069 m 255.395 -189.069 l S Q}%
\TextAt{13}{301.116}{573.741}{n}%
\TextAt{17}{308.103}{577.195}{2}%
\TextAt{21}{312.338}{577.195}{s}%
\TextAt{3}{316.253}{577.195}{\Character{0}}%
\TextAt{21}{322.84}{577.195}{r}%
\TextAt{9}{331.911}{581.942}{\Character{61}}%
\TextAt{13}{344.337}{581.942}{n}%
\TextAt{21}{351.324}{586.878}{r}%
\TextAt{9}{359.2}{581.942}{\Character{61}}%
\TextAt{9}{371.61}{581.942}{\Character{91}}%
\TextAt{13}{374.877}{581.942}{L}%
\TextAt{9}{386.162}{581.942}{\Character{58}}%
\TextAt{13}{392.735}{581.942}{K}%
\TextAt{9}{403.552}{581.942}{\Character{93}}%
\TextAt{13}{406.804}{581.942}{\Character{58}}%
\TextAt{8}{79.226}{551.824}{We}%
\TextAt{8}{100.289}{551.824}{now}%
\TextAt{8}{125.597}{551.824}{show}%
\TextAt{8}{155.518}{551.824}{the}%
\TextAt{8}{176.594}{551.824}{inclusion}%
\TextAt{13}{226.349}{551.824}{C}%
\TextAt{21}{234.724}{550.03}{K}%
\TextAt{9}{242.927}{551.824}{\Character{40}}%
\TextAt{13}{247.479}{551.824}{S\Character{59}}%
\TextAt{13}{259.948}{551.824}{T}%
\TextAt{9}{268.455}{551.824}{\Character{41}}%
\NativeMathAt{277.894}{551.824}{11.955}{\subseteq}%
\TextAt{13}{292.079}{551.824}{N}%
\TextAt{21}{301.453}{549.964}{L}%
\TextAt{3}{307.212}{549.964}{\Character{106}}%
\TextAt{21}{309.564}{549.964}{K}%
\TextAt{13}{317.767}{551.824}{C}%
\TextAt{21}{326.142}{550.03}{L}%
\TextAt{8}{332.399}{551.824}{.}%
\TextAt{8}{343.603}{551.824}{Let}%
\TextAt{13}{365.491}{551.824}{\Character{11}}%
\NativeMathAt{377.9}{551.824}{11.955}{\in}%
\TextAt{13}{390.756}{551.824}{I}%
\TextAt{21}{395.919}{550.03}{K}%
\TextAt{9}{404.122}{551.824}{\Character{40}}%
\TextAt{13}{408.674}{551.824}{S\Character{59}}%
\TextAt{13}{421.143}{551.824}{T}%
\TextAt{9}{429.65}{551.824}{\Character{41}}%
\TextAt{8}{434.202}{551.824}{.}%
\TextAt{8}{445.407}{551.824}{In}%
\TextAt{8}{460.947}{551.824}{order}%
\TextAt{8}{492.435}{551.824}{to}%
\TextAt{8}{507.664}{551.824}{show}%
\TextAt{8}{79.226}{534.778}{that}%
\TextAt{13}{105.37}{534.778}{\Character{11}}%
\NativeMathAt{117.533}{534.778}{11.955}{\in}%
\TextAt{13}{130.145}{534.778}{N}%
\TextAt{21}{139.519}{532.918}{L}%
\TextAt{3}{145.278}{532.918}{\Character{106}}%
\TextAt{21}{147.63}{532.918}{K}%
\TextAt{13}{155.833}{534.778}{I}%
\TextAt{21}{160.996}{532.984}{L}%
\TextAt{8}{171.925}{534.778}{all}%
\TextAt{8}{188.959}{534.778}{we}%
\TextAt{8}{206.962}{534.778}{have}%
\TextAt{8}{234.732}{534.778}{to}%
\TextAt{8}{249.817}{534.778}{check,}%
\TextAt{8}{285.575}{534.778}{by}%
\TextAt{8}{302.609}{534.778}{(3.2),}%
\TextAt{8}{334.805}{534.778}{is}%
\TextAt{8}{347.344}{534.778}{again}%
\TextAt{8}{379.346}{534.778}{that}%
\TextAt{8}{405.49}{534.778}{every}%
\TextAt{8}{437.157}{534.778}{component}%
\TextAt{13}{497.768}{534.778}{\Character{11}}%
\TextAt{27}{505.213}{532.984}{p}%
\TextAt{8}{514.373}{534.778}{is}%
\TextAt{8}{526.911}{534.778}{a}%
\TextAt{8}{79.226}{517.732}{norm}%
\TextAt{8}{110.185}{517.732}{from}%
\TextAt{13}{138.215}{517.732}{L}%
\TextAt{27}{146.18}{515.938}{P}%
\TextAt{16}{157.259}{517.732}{\Character{106}}%
\TextAt{13}{164.565}{517.732}{K}%
\TextAt{27}{174.524}{515.938}{p}%
\TextAt{8}{179.01}{517.732}{.}%
\TextAt{8}{188.628}{517.732}{For}%
\TextAt{14}{209.987}{517.732}{p}%
\NativeMathAt{219.955}{517.732}{11.955}{\in}%
\TextAt{13}{231.911}{517.732}{S}%
\TextAt{8}{244.093}{517.732}{this}%
\TextAt{8}{267.305}{517.732}{is}%
\TextAt{8}{279.459}{517.732}{true}%
\TextAt{8}{304.56}{517.732}{because}%
\TextAt{13}{348.251}{517.732}{\Character{11}}%
\TextAt{27}{355.697}{515.938}{p}%
\NativeMathAt{364.169}{517.732}{11.955}{\in}%
\TextAt{13}{376.124}{517.732}{K}%
\TextAt{3}{386.942}{522.07}{\Character{3}}%
\TextAt{21}{391.176}{522.07}{n}%
\TextAt{27}{386.083}{514.776}{p}%
\TextAt{8}{401.099}{517.732}{is}%
\TextAt{8}{413.253}{517.732}{an}%
\TextAt{13}{429.902}{517.732}{n}%
\TextAt{8}{436.89}{517.732}{-th}%
\TextAt{8}{456.133}{517.732}{power,}%
\TextAt{8}{494.016}{517.732}{hence}%
\TextAt{8}{526.911}{517.732}{a}%
\TextAt{8}{79.226}{500.686}{norm}%
\TextAt{8}{109.879}{500.686}{from}%
\TextAt{13}{137.604}{500.686}{K}%
\TextAt{27}{147.563}{498.892}{p}%
\TextAt{9}{152.05}{500.686}{\Character{40}}%
\TextAt{22}{159.923}{503.842}{n}%
\TextAt{20}{157.882}{509.772}{\Character{112}}%
\Vector{q 0 G 0 g 0.478 w 97.837 -263.989 m 113.387 -263.989 l S Q}%
\TextAt{13}{169.837}{500.686}{K}%
\TextAt{3}{180.655}{504.139}{\Character{3}}%
\TextAt{27}{179.796}{497.73}{p}%
\TextAt{9}{185.387}{500.686}{\Character{41}}%
\TextAt{8}{193.921}{500.686}{(see}%
\TextAt{8}{217.473}{500.686}{chap.}%
\TextAt{8}{248.449}{500.686}{V,}%
\TextAt{8}{264.457}{500.686}{(1.5)),}%
\TextAt{8}{300.344}{500.686}{so}%
\TextAt{8}{314.798}{500.686}{in}%
\TextAt{8}{328.534}{500.686}{particular}%
\TextAt{8}{382.608}{500.686}{also}%
\TextAt{8}{406.172}{500.686}{from}%
\TextAt{13}{433.896}{500.686}{L}%
\TextAt{27}{441.86}{498.892}{P}%
\TextAt{16}{452.418}{500.686}{\Character{106}}%
\TextAt{13}{459.203}{500.686}{K}%
\TextAt{27}{469.162}{498.892}{p}%
\TextAt{8}{473.648}{500.686}{.}%
\TextAt{8}{482.349}{500.686}{For}%
\TextAt{14}{503.402}{500.686}{p}%
\NativeMathAt{512.849}{500.686}{11.955}{\in}%
\TextAt{13}{524.283}{500.686}{T}%
\TextAt{8}{79.226}{483.64}{it}%
\TextAt{8}{91.038}{483.64}{holds}%
\TextAt{8}{121.775}{483.64}{because}%
\TextAt{8}{165.185}{483.64}{(4.1)}%
\TextAt{8}{193.268}{483.64}{gives}%
\TextAt{9}{222.056}{483.64}{\Character{1}}%
\NativeMathAt{235.316}{483.64}{11.955}{\subseteq}%
\TextAt{13}{248.119}{483.64}{K}%
\TextAt{3}{258.937}{487.978}{\Character{3}}%
\TextAt{21}{263.171}{487.978}{n}%
\TextAt{27}{258.078}{480.684}{p}%
\TextAt{8}{268.807}{483.64}{,}%
\TextAt{8}{276.091}{483.64}{and}%
\TextAt{8}{298.962}{483.64}{thus}%
\TextAt{13}{324.821}{483.64}{L}%
\TextAt{27}{332.786}{481.846}{P}%
\TextAt{9}{343.384}{483.64}{\Character{61}}%
\TextAt{13}{355.993}{483.64}{K}%
\TextAt{27}{365.952}{481.846}{p}%
\TextAt{8}{370.439}{483.64}{.}%
\TextAt{8}{379.211}{483.64}{Finally,}%
\TextAt{8}{421.452}{483.64}{it}%
\TextAt{8}{433.264}{483.64}{holds}%
\TextAt{8}{464.001}{483.64}{for}%
\TextAt{14}{481.994}{483.64}{p}%
\NativeMathAt{491.482}{483.64}{11.955}{\notin}%
\TextAt{13}{502.957}{483.64}{S}%
\TextAt{16}{513.582}{483.64}{\Character{91}}%
\TextAt{13}{524.283}{483.64}{T}%
\TextAt{8}{79.226}{466.594}{since}%
\TextAt{13}{109.202}{466.594}{\Character{11}}%
\TextAt{27}{116.648}{464.8}{p}%
\TextAt{8}{126.34}{466.594}{is}%
\TextAt{8}{139.411}{466.594}{a}%
\TextAt{8}{150.474}{466.594}{unit}%
\TextAt{8}{176.493}{466.594}{and}%
\TextAt{13}{200.563}{466.594}{L}%
\TextAt{27}{208.527}{464.8}{P}%
\TextAt{16}{221.17}{466.594}{\Character{106}}%
\TextAt{13}{230.04}{466.594}{K}%
\TextAt{27}{239.999}{464.8}{p}%
\TextAt{8}{249.691}{466.594}{is}%
\TextAt{8}{262.762}{466.594}{unramified}%
\TextAt{8}{322.602}{466.594}{(see}%
\TextAt{8}{347.378}{466.594}{chap.}%
\TextAt{8}{379.577}{466.594}{V,}%
\TextAt{8}{396.809}{466.594}{(3.3)).}%
\TextAt{8}{437.813}{466.594}{We}%
\TextAt{8}{459.265}{466.594}{therefore}%
\TextAt{8}{509.672}{466.594}{have}%
\TextAt{13}{79.226}{449.548}{I}%
\TextAt{21}{84.389}{447.754}{K}%
\TextAt{9}{92.592}{449.548}{\Character{40}}%
\TextAt{13}{97.144}{449.548}{S\Character{59}}%
\TextAt{13}{109.613}{449.548}{T}%
\TextAt{9}{118.12}{449.548}{\Character{41}}%
\NativeMathAt{126.559}{449.548}{11.955}{\subseteq}%
\TextAt{13}{139.745}{449.548}{N}%
\TextAt{21}{149.119}{447.688}{L}%
\TextAt{3}{154.878}{447.688}{\Character{106}}%
\TextAt{21}{157.23}{447.688}{K}%
\TextAt{13}{165.433}{449.548}{I}%
\TextAt{21}{170.595}{447.754}{L}%
\TextAt{8}{176.853}{449.548}{,}%
\TextAt{8}{184.417}{449.548}{i.e.,}%
\TextAt{13}{206.937}{449.548}{C}%
\TextAt{21}{215.312}{447.754}{K}%
\TextAt{9}{223.515}{449.548}{\Character{40}}%
\TextAt{13}{228.067}{449.548}{S\Character{59}}%
\TextAt{13}{240.536}{449.548}{T}%
\TextAt{9}{249.043}{449.548}{\Character{41}}%
\NativeMathAt{257.483}{449.548}{11.955}{\subseteq}%
\TextAt{13}{270.668}{449.548}{N}%
\TextAt{21}{280.042}{447.688}{L}%
\TextAt{3}{285.801}{447.688}{\Character{106}}%
\TextAt{21}{288.153}{447.688}{K}%
\TextAt{13}{296.356}{449.548}{C}%
\TextAt{21}{304.731}{447.754}{L}%
\TextAt{8}{310.988}{449.548}{.}%
\TextAt{8}{320.433}{449.548}{Now}%
\TextAt{8}{347.425}{449.548}{if}%
\TextAt{13}{358.481}{449.548}{L}%
\TextAt{16}{370.333}{449.548}{\Character{106}}%
\TextAt{13}{377.54}{449.548}{K}%
\TextAt{8}{392.588}{449.548}{is}%
\TextAt{8}{404.684}{449.548}{cyclic,}%
\TextAt{8}{440.534}{449.548}{i.e.,}%
\TextAt{8}{463.055}{449.548}{if}%
\TextAt{13}{474.11}{449.548}{r}%
\TextAt{9}{483.598}{449.548}{\Character{61}}%
\TextAt{9}{496.581}{449.548}{1}%
\TextAt{8}{502.442}{449.548}{,}%
\TextAt{8}{510.007}{449.548}{then}%
\TextAt{8}{79.226}{432.502}{from}%
\TextAt{8}{106.866}{432.502}{(3.6),}%
\TextAt{9}{166.315}{415.456}{\Character{91}}%
\TextAt{13}{169.567}{415.456}{L}%
\TextAt{9}{180.852}{415.456}{\Character{58}}%
\TextAt{13}{187.424}{415.456}{K}%
\TextAt{9}{198.242}{415.456}{\Character{93}}%
\NativeMathAt{204.815}{415.456}{11.955}{\leq}%
\TextAt{9}{217.434}{415.456}{\Character{40}}%
\TextAt{13}{221.986}{415.456}{C}%
\TextAt{21}{230.361}{413.663}{K}%
\TextAt{9}{241.885}{415.456}{\Character{58}}%
\TextAt{13}{248.457}{415.456}{N}%
\TextAt{21}{257.831}{413.596}{L}%
\TextAt{3}{263.59}{413.596}{\Character{106}}%
\TextAt{21}{265.942}{413.596}{K}%
\TextAt{13}{274.145}{415.456}{C}%
\TextAt{21}{282.52}{413.663}{L}%
\TextAt{9}{288.777}{415.456}{\Character{41}}%
\NativeMathAt{296.651}{415.456}{11.955}{\leq}%
\TextAt{9}{309.27}{415.456}{\Character{40}}%
\TextAt{13}{313.822}{415.456}{C}%
\TextAt{21}{322.197}{413.663}{K}%
\TextAt{9}{333.721}{415.456}{\Character{58}}%
\TextAt{13}{340.293}{415.456}{C}%
\TextAt{21}{348.668}{413.663}{K}%
\TextAt{9}{356.871}{415.456}{\Character{40}}%
\TextAt{13}{361.423}{415.456}{S\Character{59}}%
\TextAt{13}{373.892}{415.456}{T}%
\TextAt{9}{382.399}{415.456}{\Character{41}\Character{41}}%
\TextAt{9}{394.797}{415.456}{\Character{61}}%
\TextAt{9}{407.218}{415.456}{\Character{91}}%
\TextAt{13}{410.502}{415.456}{L}%
\TextAt{9}{421.787}{415.456}{\Character{58}}%
\TextAt{13}{428.36}{415.456}{K}%
\TextAt{9}{439.177}{415.456}{\Character{93}}%
\TextAt{13}{442.429}{415.456}{\Character{59}}%
\TextAt{8}{79.226}{391.934}{hence}%
\TextAt{13}{111.732}{391.934}{N}%
\TextAt{21}{121.106}{390.074}{L}%
\TextAt{3}{126.865}{390.074}{\Character{106}}%
\TextAt{21}{129.217}{390.074}{K}%
\TextAt{13}{137.42}{391.934}{C}%
\TextAt{21}{145.795}{390.141}{L}%
\TextAt{9}{155.373}{391.934}{\Character{61}}%
\TextAt{13}{167.799}{391.934}{C}%
\TextAt{21}{176.174}{390.141}{K}%
\TextAt{9}{184.376}{391.934}{\Character{40}}%
\TextAt{13}{188.929}{391.934}{S\Character{59}}%
\TextAt{13}{201.398}{391.934}{T}%
\TextAt{9}{209.905}{391.934}{\Character{41}}%
\TextAt{8}{214.457}{391.934}{.}%
\Vector{q 0 G 0 g 0.399 w 452.334 -382.066 m 452.334 -373.996 l S Q}%
\Vector{q 0 G 0 g 0.399 w 452.534 -374.195 m 459.707 -374.195 l S Q}%
\Vector{q 0 G 0 g 0.399 w 452.534 -381.867 m 459.707 -381.867 l S Q}%
\Vector{q 0 G 0 g 0.399 w 459.906 -382.066 m 459.906 -373.996 l S Q}%
\TextAt{8}{97.159}{362.933}{Now}%
\TextAt{8}{123.129}{362.933}{that}%
\TextAt{8}{147.809}{362.933}{we}%
\TextAt{8}{164.347}{362.933}{have}%
\TextAt{8}{190.652}{362.933}{an}%
\TextAt{8}{206.221}{362.933}{explicit}%
\TextAt{8}{246.825}{362.933}{picture}%
\TextAt{8}{285.803}{362.933}{in}%
\TextAt{8}{298.766}{362.933}{the}%
\TextAt{8}{318.233}{362.933}{case}%
\TextAt{8}{342.314}{362.933}{of}%
\TextAt{8}{354.955}{362.933}{a}%
\TextAt{8}{364.021}{362.933}{Kummer}%
\TextAt{8}{412.096}{362.933}{field,}%
\TextAt{8}{440.153}{362.933}{the}%
\TextAt{8}{459.62}{362.933}{result}%
\TextAt{8}{491.508}{362.933}{we}%
\TextAt{8}{508.046}{362.933}{want}%
\TextAt{8}{79.226}{345.887}{follows}%
\TextAt{8}{117.662}{345.887}{also}%
\TextAt{8}{141.142}{345.887}{in}%
\TextAt{8}{154.795}{345.887}{complete}%
\TextAt{8}{204.217}{345.887}{generality:}%
\TextAt{12}{79.226}{319.875}{Theorem}%
\TextAt{12}{135.089}{319.875}{(4.4)}%
\TextAt{12}{166.097}{319.875}{(Global}%
\TextAt{12}{213.46}{319.875}{Class}%
\TextAt{12}{247.409}{319.875}{Field}%
\TextAt{12}{280.281}{319.875}{Axiom).}%
\TextAt{8}{331.106}{319.875}{If}%
\TextAt{13}{341.797}{319.875}{L}%
\TextAt{16}{353.083}{319.875}{\Character{106}}%
\TextAt{13}{359.724}{319.875}{K}%
\TextAt{8}{373.439}{319.875}{is}%
\TextAt{8}{384.201}{319.875}{a}%
\TextAt{8}{392.956}{319.875}{cyclic}%
\TextAt{8}{424.138}{319.875}{extension}%
\TextAt{8}{474.903}{319.875}{of}%
\TextAt{8}{487.232}{319.875}{algebraic}%
\TextAt{8}{79.226}{302.829}{number}%
\TextAt{8}{121.822}{302.829}{fields,}%
\TextAt{8}{155.046}{302.829}{then}%
\TextAt{8}{190.721}{263.062}{\char35\relax }%
\TextAt{13}{200.476}{263.062}{H}%
\TextAt{21}{211.099}{267.998}{i}%
\TextAt{9}{214.48}{263.062}{\Character{40}}%
\TextAt{13}{219.033}{263.062}{G}%
\TextAt{9}{228.266}{263.062}{\Character{40}}%
\TextAt{13}{232.818}{263.062}{L}%
\TextAt{16}{244.104}{263.062}{\Character{106}}%
\TextAt{13}{250.746}{263.062}{K}%
\TextAt{9}{261.563}{263.062}{\Character{41}}%
\TextAt{13}{266.116}{263.062}{\Character{59}}%
\TextAt{13}{271.34}{263.062}{C}%
\TextAt{21}{279.735}{261.268}{L}%
\TextAt{9}{285.992}{263.062}{\Character{41}}%
\TextAt{9}{293.847}{263.062}{\Character{61}}%
\TextAt{20}{306.291}{287.57}{\Character{56}}%
\TextAt{20}{306.291}{276.81}{\Character{60}}%
\TextAt{20}{306.291}{255.291}{\Character{58}}%
\TextAt{9}{316.917}{272.187}{\Character{91}}%
\TextAt{13}{320.169}{272.187}{L}%
\TextAt{9}{331.454}{272.187}{\Character{58}}%
\TextAt{13}{338.027}{272.187}{K}%
\TextAt{9}{348.845}{272.187}{\Character{93}}%
\TextAt{8}{364.051}{272.187}{for}%
\TextAt{13}{381.936}{272.187}{i}%
\TextAt{9}{389.251}{272.187}{\Character{61}}%
\TextAt{9}{401.661}{272.187}{0}%
\TextAt{13}{407.529}{272.187}{\Character{59}}%
\TextAt{9}{316.917}{251.732}{1}%
\TextAt{8}{364.051}{251.732}{for}%
\TextAt{13}{381.936}{251.732}{i}%
\TextAt{9}{389.251}{251.732}{\Character{61}}%
\TextAt{16}{401.676}{251.732}{\Character{0}}%
\TextAt{9}{410.975}{251.732}{1}%
\TextAt{13}{416.828}{251.732}{\Character{58}}%
\TextAt{19}{79.226}{204.172}{Proof.}%
\TextAt{8}{116}{204.172}{Since}%
\TextAt{13}{147.38}{204.172}{h}%
\TextAt{9}{154.118}{204.172}{\Character{40}}%
\TextAt{13}{158.671}{204.172}{G}%
\TextAt{9}{167.904}{204.172}{\Character{40}}%
\TextAt{13}{172.456}{204.172}{L}%
\TextAt{16}{185.143}{204.172}{\Character{106}}%
\TextAt{13}{193.185}{204.172}{K}%
\TextAt{9}{204.003}{204.172}{\Character{41}}%
\TextAt{13}{208.555}{204.172}{\Character{59}}%
\TextAt{13}{213.779}{204.172}{C}%
\TextAt{21}{222.174}{202.379}{L}%
\TextAt{9}{228.431}{204.172}{\Character{41}}%
\TextAt{9}{237.684}{204.172}{\Character{61}}%
\TextAt{9}{251.505}{204.172}{\Character{91}}%
\TextAt{13}{254.783}{204.172}{L}%
\TextAt{9}{267.47}{204.172}{\Character{58}}%
\TextAt{13}{275.443}{204.172}{K}%
\TextAt{9}{286.26}{204.172}{\Character{93}}%
\TextAt{8}{289.512}{204.172}{,}%
\TextAt{8}{297.689}{204.172}{it}%
\TextAt{8}{310.215}{204.172}{is}%
\TextAt{8}{322.801}{204.172}{clearly}%
\TextAt{8}{361.019}{204.172}{enough}%
\TextAt{8}{402.165}{204.172}{to}%
\TextAt{8}{417.298}{204.172}{show}%
\TextAt{8}{447.123}{204.172}{that}%
\TextAt{8}{473.314}{204.172}{\char35\relax }%
\TextAt{13}{483.069}{204.172}{H}%
\TextAt{17}{493.692}{208.51}{0}%
\TextAt{9}{498.424}{204.172}{\Character{40}}%
\TextAt{13}{502.976}{204.172}{G}%
\TextAt{9}{512.21}{204.172}{\Character{40}}%
\TextAt{13}{516.762}{204.172}{L}%
\TextAt{16}{529.449}{204.172}{\Character{106}}%
\TextAt{13}{79.226}{187.126}{K}%
\TextAt{9}{90.044}{187.126}{\Character{41}}%
\TextAt{13}{94.596}{187.126}{\Character{59}}%
\TextAt{13}{99.82}{187.126}{C}%
\TextAt{21}{108.215}{185.333}{L}%
\TextAt{9}{114.472}{187.126}{\Character{41}}%
\TextAt{16}{122.915}{187.126}{\Character{106}}%
\TextAt{9}{130.125}{187.126}{\Character{91}}%
\TextAt{13}{133.377}{187.126}{L}%
\TextAt{9}{145.231}{187.126}{\Character{58}}%
\TextAt{13}{152.373}{187.126}{K}%
\TextAt{9}{163.191}{187.126}{\Character{93}}%
\TextAt{8}{166.442}{187.126}{.}%
\TextAt{8}{175.892}{187.126}{We}%
\TextAt{8}{196.371}{187.126}{will}%
\TextAt{8}{218.81}{187.126}{prove}%
\TextAt{8}{250.693}{187.126}{this}%
\TextAt{8}{273.85}{187.126}{by}%
\TextAt{8}{290.442}{187.126}{induction}%
\TextAt{8}{342.805}{187.126}{on}%
\TextAt{8}{359.398}{187.126}{the}%
\TextAt{8}{379.889}{187.126}{degree}%
\TextAt{13}{416.638}{187.126}{n}%
\TextAt{9}{427.515}{187.126}{\Character{61}}%
\TextAt{9}{440.498}{187.126}{\Character{91}}%
\TextAt{13}{443.761}{187.126}{L}%
\TextAt{9}{455.616}{187.126}{\Character{58}}%
\TextAt{13}{462.757}{187.126}{K}%
\TextAt{9}{473.575}{187.126}{\Character{93}}%
\TextAt{8}{476.827}{187.126}{.}%
\TextAt{8}{486.277}{187.126}{We}%
\TextAt{8}{506.755}{187.126}{write}%
\TextAt{8}{79.226}{170.08}{for}%
\TextAt{8}{97.004}{170.08}{short}%
\TextAt{13}{126.88}{170.08}{H}%
\TextAt{17}{137.503}{174.418}{0}%
\TextAt{9}{142.235}{170.08}{\Character{40}}%
\TextAt{13}{146.787}{170.08}{L}%
\TextAt{16}{158.073}{170.08}{\Character{106}}%
\TextAt{13}{164.714}{170.08}{K}%
\TextAt{9}{175.532}{170.08}{\Character{41}}%
\TextAt{8}{183.874}{170.08}{instead}%
\TextAt{8}{224.152}{170.08}{of}%
\TextAt{13}{237.374}{170.08}{H}%
\TextAt{17}{247.997}{174.418}{0}%
\TextAt{9}{252.729}{170.08}{\Character{40}}%
\TextAt{13}{257.282}{170.08}{G}%
\TextAt{9}{266.515}{170.08}{\Character{40}}%
\TextAt{13}{271.068}{170.08}{L}%
\TextAt{16}{282.353}{170.08}{\Character{106}}%
\TextAt{13}{288.995}{170.08}{K}%
\TextAt{9}{299.812}{170.08}{\Character{41}}%
\TextAt{13}{304.365}{170.08}{\Character{59}}%
\TextAt{13}{309.589}{170.08}{C}%
\TextAt{21}{317.984}{168.287}{L}%
\TextAt{9}{324.241}{170.08}{\Character{41}}%
\TextAt{8}{328.793}{170.08}{.}%
\TextAt{8}{337.206}{170.08}{Let}%
\TextAt{13}{358.068}{170.08}{M}%
\TextAt{16}{373.963}{170.08}{\Character{106}}%
\TextAt{13}{380.604}{170.08}{K}%
\TextAt{8}{395.212}{170.08}{be}%
\TextAt{8}{411.029}{170.08}{a}%
\TextAt{8}{420.677}{170.08}{subextension}%
\TextAt{8}{490.28}{170.08}{of}%
\TextAt{8}{503.503}{170.08}{prime}%
\TextAt{8}{79.226}{153.034}{degree}%
\TextAt{13}{115.642}{153.034}{p}%
\TextAt{8}{121.517}{153.034}{.}%
\TextAt{8}{129.965}{153.034}{We}%
\TextAt{8}{150.11}{153.034}{consider}%
\TextAt{8}{195.695}{153.034}{the}%
\TextAt{8}{215.851}{153.034}{exact}%
\TextAt{8}{246.743}{153.034}{sequence}%
\TextAt{13}{158.492}{122.043}{C}%
\TextAt{21}{166.867}{120.25}{M}%
\TextAt{9}{176.347}{122.043}{\Character{47}}%
\TextAt{13}{182.2}{122.043}{N}%
\TextAt{21}{191.573}{120.184}{L}%
\TextAt{3}{197.332}{120.184}{\Character{106}}%
\TextAt{21}{199.685}{120.184}{M}%
\TextAt{13}{209.165}{122.043}{C}%
\TextAt{21}{217.54}{120.25}{L}%
\TextAt{21}{229.705}{130.99}{N}%
\TextAt{22}{236.427}{129.252}{M}%
\TextAt{23}{244.015}{129.252}{\Character{106}}%
\TextAt{22}{246.284}{129.252}{K}%
\NativeArrowAt{227.118}{122.043}{11.955}{30.187}{\longrightarrow}%
\TextAt{13}{260.647}{122.043}{C}%
\TextAt{21}{269.022}{120.25}{K}%
\TextAt{9}{277.224}{122.043}{\Character{47}}%
\TextAt{13}{283.077}{122.043}{N}%
\TextAt{21}{292.451}{120.184}{L}%
\TextAt{3}{298.21}{120.184}{\Character{106}}%
\TextAt{21}{300.563}{120.184}{K}%
\TextAt{13}{308.765}{122.043}{C}%
\TextAt{21}{317.14}{120.25}{L}%
\NativeArrowAt{326.718}{122.043}{11.955}{19.26}{\longrightarrow}%
\TextAt{13}{349.3}{122.043}{C}%
\TextAt{21}{357.675}{120.25}{K}%
\TextAt{9}{365.878}{122.043}{\Character{47}}%
\TextAt{13}{371.731}{122.043}{N}%
\TextAt{21}{381.105}{120.184}{M}%
\TextAt{3}{390.087}{120.184}{\Character{106}}%
\TextAt{21}{392.439}{120.184}{K}%
\TextAt{13}{400.642}{122.043}{C}%
\TextAt{21}{409.017}{120.25}{M}%
\NativeArrowAt{421.817}{122.043}{11.955}{19.26}{\longrightarrow}%
\TextAt{9}{444.399}{122.043}{1}%
\TextAt{13}{450.252}{122.043}{\Character{59}}%
\TextAt{8}{79.226}{93.042}{i.e.,}%
\TextAt{8}{102.513}{93.042}{the}%
\TextAt{8}{123.614}{93.042}{exact}%
\TextAt{8}{155.452}{93.042}{sequence}%
\TextAt{13}{204.899}{93.042}{H}%
\TextAt{17}{215.521}{97.381}{0}%
\TextAt{9}{220.254}{93.042}{\Character{40}}%
\TextAt{13}{224.806}{93.042}{L}%
\TextAt{16}{237.702}{93.042}{\Character{106}}%
\TextAt{13}{245.954}{93.042}{M}%
\TextAt{9}{258.528}{93.042}{\Character{41}}%
\NativeArrowAt{268.012}{93.042}{11.955}{19.26}{\longrightarrow}%
\TextAt{13}{292.205}{93.042}{H}%
\TextAt{17}{302.827}{97.381}{0}%
\TextAt{9}{307.559}{93.042}{\Character{40}}%
\TextAt{13}{312.112}{93.042}{L}%
\TextAt{16}{325.008}{93.042}{\Character{106}}%
\TextAt{13}{333.26}{93.042}{K}%
\TextAt{9}{344.078}{93.042}{\Character{41}}%
\NativeArrowAt{353.562}{93.042}{11.955}{19.26}{\longrightarrow}%
\TextAt{13}{377.754}{93.042}{H}%
\TextAt{17}{388.377}{97.381}{0}%
\TextAt{9}{393.109}{93.042}{\Character{40}}%
\TextAt{13}{397.662}{93.042}{M}%
\TextAt{16}{415.167}{93.042}{\Character{106}}%
\TextAt{13}{423.419}{93.042}{K}%
\TextAt{9}{434.237}{93.042}{\Character{41}}%
\NativeArrowAt{443.721}{93.042}{11.955}{19.26}{\longrightarrow}%
\TextAt{9}{467.913}{93.042}{1}%
\TextAt{8}{473.766}{93.042}{.}%
\TextAt{8}{485.05}{93.042}{If}%
\TextAt{13}{497.687}{93.042}{p}%
\TextAt{13}{508.483}{93.042}{\Character{60}}%
\TextAt{13}{522.518}{93.042}{n}%
\TextAt{8}{529.518}{93.042}{,}%
\TextAt{8}{79.226}{75.996}{then}%
\TextAt{8}{106.205}{75.996}{\char35\relax }%
\TextAt{13}{115.961}{75.996}{H}%
\TextAt{17}{126.583}{80.335}{0}%
\TextAt{9}{131.316}{75.996}{\Character{40}}%
\TextAt{13}{135.868}{75.996}{L}%
\TextAt{16}{147.697}{75.996}{\Character{106}}%
\TextAt{13}{154.882}{75.996}{M}%
\TextAt{9}{167.456}{75.996}{\Character{41}}%
\TextAt{16}{175.872}{75.996}{\Character{106}}%
\TextAt{9}{183.058}{75.996}{\Character{91}}%
\TextAt{13}{186.309}{75.996}{L}%
\TextAt{9}{198.138}{75.996}{\Character{58}}%
\TextAt{13}{205.255}{75.996}{M}%
\TextAt{9}{217.828}{75.996}{\Character{93}}%
\TextAt{8}{221.08}{75.996}{,}%
\TextAt{8}{228.628}{75.996}{\char35\relax }%
\TextAt{13}{238.383}{75.996}{H}%
\TextAt{17}{249.006}{80.335}{0}%
\TextAt{9}{253.738}{75.996}{\Character{40}}%
\TextAt{13}{258.29}{75.996}{M}%
\TextAt{16}{274.728}{75.996}{\Character{106}}%
\TextAt{13}{281.914}{75.996}{K}%
\TextAt{9}{292.731}{75.996}{\Character{41}}%
\TextAt{16}{301.148}{75.996}{\Character{106}}%
\TextAt{9}{308.334}{75.996}{\Character{91}}%
\TextAt{13}{311.585}{75.996}{M}%
\TextAt{9}{328.023}{75.996}{\Character{58}}%
\TextAt{13}{335.139}{75.996}{K}%
\TextAt{9}{345.957}{75.996}{\Character{93}}%
\TextAt{8}{353.425}{75.996}{by}%
\TextAt{8}{370.003}{75.996}{the}%
\TextAt{8}{390.479}{75.996}{induction}%
\TextAt{8}{442.827}{75.996}{hypothesis,}%
\TextAt{8}{504.161}{75.996}{hence}%
\TextAt{8}{297.211}{43.033}{349}%
\Destination{78.23}{809.815}{nav.page.349}%
\Destination{79.36}{566.855}{nav.print.383}%
\Destination{79.36}{333.855}{nav.number.6.4.4}%
\Hotlink{403.669}{700.955}{9.763}{11.047}{nav.chapter.2}%
\Hotlink{247.949}{498.455}{9.775}{11.047}{nav.chapter.5}%
\Hotlink{379.077}{464.355}{9.775}{11.047}{nav.chapter.5}%
\Hotlink{419.8}{700.955}{25.078}{11.047}{nav.number.2.5.8}%
\Hotlink{302.109}{532.555}{25.078}{11.047}{nav.number.6.3.2}%
\Hotlink{263.957}{498.455}{25.078}{11.047}{nav.number.5.1.5}%
\Hotlink{164.685}{481.455}{25.078}{11.047}{nav.number.6.4.1}%
\Hotlink{396.309}{464.355}{25.078}{11.047}{nav.number.5.3.3}%
\Hotlink{106.366}{430.255}{25.078}{11.047}{nav.number.6.3.6}%
}

% Page 362 (body page 350)
\Page{%
\TextAt{2}{79.226}{796.943}{VI.}%
\TextAt{2}{98.012}{796.943}{Global}%
\TextAt{2}{133.248}{796.943}{Class}%
\TextAt{2}{161.841}{796.943}{Field}%
\TextAt{2}{189.572}{796.943}{Theory}%
\Vector{q 0 G 0 g 0.299 w 7.226 18.729 m 460.769 18.729 l S Q}%
\TextAt{8}{79.226}{761.124}{\char35\relax }%
\TextAt{13}{88.982}{761.124}{H}%
\TextAt{17}{99.604}{765.462}{0}%
\TextAt{9}{104.336}{761.124}{\Character{40}}%
\TextAt{13}{108.889}{761.124}{L}%
\TextAt{16}{121.887}{761.124}{\Character{106}}%
\TextAt{13}{130.241}{761.124}{K}%
\TextAt{9}{141.058}{761.124}{\Character{41}}%
\TextAt{16}{150.644}{761.124}{\Character{106}}%
\TextAt{9}{158.998}{761.124}{\Character{91}}%
\TextAt{13}{162.25}{761.124}{L}%
\TextAt{9}{175.247}{761.124}{\Character{58}}%
\TextAt{13}{183.532}{761.124}{M}%
\TextAt{9}{196.106}{761.124}{\Character{93}\Character{91}}%
\TextAt{13}{202.609}{761.124}{M}%
\TextAt{9}{220.216}{761.124}{\Character{58}}%
\TextAt{13}{228.501}{761.124}{K}%
\TextAt{9}{239.318}{761.124}{\Character{93}}%
\TextAt{9}{247.591}{761.124}{\Character{61}}%
\TextAt{9}{261.722}{761.124}{\Character{91}}%
\TextAt{13}{264.993}{761.124}{L}%
\TextAt{9}{277.99}{761.124}{\Character{58}}%
\TextAt{13}{286.275}{761.124}{K}%
\TextAt{9}{297.093}{761.124}{\Character{93}}%
\TextAt{8}{300.345}{761.124}{.}%
\TextAt{8}{311.807}{761.124}{Now}%
\TextAt{8}{339.472}{761.124}{let}%
\TextAt{13}{357.382}{761.124}{p}%
\TextAt{9}{368.29}{761.124}{\Character{61}}%
\TextAt{13}{382.428}{761.124}{n}%
\TextAt{8}{389.416}{761.124}{.}%
\TextAt{8}{400.878}{761.124}{We}%
\TextAt{8}{422.028}{761.124}{put}%
\TextAt{13}{444.492}{761.124}{K}%
\TextAt{3}{455.31}{765.462}{\Character{48}}%
\TextAt{9}{463.138}{761.124}{\Character{61}}%
\TextAt{13}{477.276}{761.124}{K}%
\TextAt{9}{488.093}{761.124}{\Character{40}}%
\TextAt{13}{492.646}{761.124}{\Character{22}}%
\TextAt{21}{499.689}{759.33}{p}%
\TextAt{9}{504.45}{761.124}{\Character{41}}%
\TextAt{8}{513.904}{761.124}{and}%
\TextAt{13}{79.226}{744.078}{L}%
\TextAt{3}{87.191}{748.416}{\Character{48}}%
\TextAt{9}{94.508}{744.078}{\Character{61}}%
\TextAt{13}{108.136}{744.078}{L}%
\TextAt{9}{116.1}{744.078}{\Character{40}}%
\TextAt{13}{120.653}{744.078}{\Character{22}}%
\TextAt{21}{127.695}{742.284}{p}%
\TextAt{9}{132.456}{744.078}{\Character{41}}%
\TextAt{8}{137.009}{744.078}{.}%
\TextAt{8}{147.573}{744.078}{Since}%
\TextAt{13}{178.835}{744.078}{d}%
\TextAt{9}{189.441}{744.078}{\Character{61}}%
\TextAt{9}{203.058}{744.078}{\Character{91}}%
\TextAt{13}{206.32}{744.078}{K}%
\TextAt{3}{217.137}{748.416}{\Character{48}}%
\TextAt{9}{224.455}{744.078}{\Character{58}}%
\TextAt{13}{232.229}{744.078}{K}%
\TextAt{9}{243.047}{744.078}{\Character{93}}%
\TextAt{16}{250.821}{744.078}{\Character{106}}%
\TextAt{13}{258.665}{744.078}{p}%
\TextAt{16}{267.678}{744.078}{\Character{0}}%
\TextAt{9}{280.113}{744.078}{1}%
\TextAt{8}{285.966}{744.078}{,}%
\TextAt{8}{293.997}{744.078}{we}%
\TextAt{8}{311.93}{744.078}{have}%
\TextAt{13}{339.63}{744.078}{G}%
\TextAt{9}{348.863}{744.078}{\Character{40}}%
\TextAt{13}{353.416}{744.078}{L}%
\TextAt{16}{365.903}{744.078}{\Character{106}}%
\TextAt{13}{373.746}{744.078}{K}%
\TextAt{9}{384.564}{744.078}{\Character{41}}%
\NativeMathAt{393.639}{744.078}{11.955}{\cong}%
\TextAt{13}{407.46}{744.078}{G}%
\TextAt{9}{416.694}{744.078}{\Character{40}}%
\TextAt{13}{421.246}{744.078}{L}%
\TextAt{3}{429.211}{748.416}{\Character{48}}%
\TextAt{16}{436.528}{744.078}{\Character{106}}%
\TextAt{13}{444.372}{744.078}{K}%
\TextAt{3}{455.19}{748.416}{\Character{48}}%
\TextAt{9}{457.985}{744.078}{\Character{41}}%
\TextAt{8}{462.537}{744.078}{.}%
\TextAt{13}{473.101}{744.078}{L}%
\TextAt{3}{481.066}{748.416}{\Character{48}}%
\TextAt{16}{488.383}{744.078}{\Character{106}}%
\TextAt{13}{496.227}{744.078}{K}%
\TextAt{3}{507.045}{748.416}{\Character{48}}%
\TextAt{8}{514.442}{744.078}{is}%
\TextAt{8}{526.911}{744.078}{a}%
\TextAt{8}{79.226}{727.032}{cyclic}%
\TextAt{8}{111.526}{727.032}{Kummer}%
\TextAt{8}{160.407}{727.032}{extension,}%
\TextAt{8}{215.57}{727.032}{so}%
\TextAt{8}{230.057}{727.032}{by}%
\TextAt{8}{246.432}{727.032}{(4.2),}%
\TextAt{8}{277.805}{727.032}{\char35\relax }%
\TextAt{13}{287.56}{727.032}{H}%
\TextAt{17}{298.183}{731.37}{0}%
\TextAt{9}{302.915}{727.032}{\Character{40}}%
\TextAt{13}{307.467}{727.032}{L}%
\TextAt{3}{315.432}{731.37}{\Character{48}}%
\TextAt{16}{321.746}{727.032}{\Character{106}}%
\TextAt{13}{328.586}{727.032}{K}%
\TextAt{3}{339.404}{731.37}{\Character{48}}%
\TextAt{9}{342.199}{727.032}{\Character{41}}%
\TextAt{9}{350.257}{727.032}{\Character{61}}%
\TextAt{9}{362.87}{727.032}{\Character{91}}%
\TextAt{13}{366.146}{727.032}{L}%
\TextAt{3}{374.11}{731.37}{\Character{48}}%
\TextAt{9}{380.424}{727.032}{\Character{58}}%
\TextAt{13}{387.195}{727.032}{K}%
\TextAt{3}{398.013}{731.37}{\Character{48}}%
\TextAt{9}{400.808}{727.032}{\Character{93}}%
\TextAt{9}{407.563}{727.032}{\Character{61}}%
\TextAt{13}{420.202}{727.032}{p}%
\TextAt{8}{426.078}{727.032}{.}%
\TextAt{8}{434.875}{727.032}{It}%
\TextAt{8}{447.664}{727.032}{therefore}%
\TextAt{8}{496.88}{727.032}{suffices}%
\TextAt{8}{79.226}{709.986}{to}%
\TextAt{8}{93.536}{709.986}{show}%
\TextAt{8}{122.54}{709.986}{that}%
\TextAt{8}{147.909}{709.986}{the}%
\TextAt{8}{168.065}{709.986}{homomorphism}%
\TextAt{13}{237.047}{684.858}{H}%
\TextAt{17}{247.67}{689.794}{0}%
\TextAt{9}{252.402}{684.858}{\Character{40}}%
\TextAt{13}{256.954}{684.858}{L}%
\TextAt{16}{268.24}{684.858}{\Character{106}}%
\TextAt{13}{274.881}{684.858}{K}%
\TextAt{9}{285.699}{684.858}{\Character{41}}%
\NativeArrowAt{293.572}{684.858}{11.955}{19.26}{\longrightarrow}%
\TextAt{13}{316.154}{684.858}{H}%
\TextAt{17}{326.777}{689.794}{0}%
\TextAt{9}{331.509}{684.858}{\Character{40}}%
\TextAt{13}{336.061}{684.858}{L}%
\TextAt{3}{344.026}{689.794}{\Character{48}}%
\TextAt{16}{350.142}{684.858}{\Character{106}}%
\TextAt{13}{356.784}{684.858}{K}%
\TextAt{3}{367.601}{689.794}{\Character{48}}%
\TextAt{9}{370.396}{684.858}{\Character{41}}%
\TextAt{8}{517.802}{684.858}{(*)}%
\TextAt{8}{79.226}{659.731}{induced}%
\TextAt{8}{122.861}{659.731}{by}%
\TextAt{8}{139.19}{659.731}{the}%
\TextAt{8}{159.417}{659.731}{inclusion}%
\TextAt{13}{208.324}{659.731}{C}%
\TextAt{21}{216.699}{657.938}{L}%
\NativeArrowAt{226.396}{659.731}{11.955}{19.26}{\longrightarrow}%
\TextAt{13}{249.098}{659.731}{C}%
\TextAt{21}{257.473}{657.938}{L}%
\TextAt{23}{263.232}{660.207}{\Character{48}}%
\TextAt{8}{270.382}{659.731}{is}%
\TextAt{8}{282.216}{659.731}{injective.}%
\TextAt{13}{333.471}{659.731}{H}%
\TextAt{17}{344.093}{664.069}{0}%
\TextAt{9}{348.825}{659.731}{\Character{40}}%
\TextAt{13}{353.378}{659.731}{L}%
\TextAt{16}{364.783}{659.731}{\Character{106}}%
\TextAt{13}{371.544}{659.731}{K}%
\TextAt{9}{382.361}{659.731}{\Character{41}}%
\TextAt{8}{390.881}{659.731}{has}%
\TextAt{8}{411.825}{659.731}{exponent}%
\TextAt{13}{462.298}{659.731}{p}%
\TextAt{8}{468.173}{659.731}{,}%
\TextAt{8}{475.41}{659.731}{because}%
\TextAt{8}{518.782}{659.731}{for}%
\TextAt{13}{79.226}{642.685}{x}%
\NativeMathAt{90.676}{642.685}{11.955}{\in}%
\TextAt{13}{103.445}{642.685}{C}%
\TextAt{21}{111.82}{640.892}{K}%
\TextAt{8}{124.787}{642.685}{we}%
\TextAt{8}{142.881}{642.685}{always}%
\TextAt{8}{181.215}{642.685}{have}%
\TextAt{13}{209.077}{642.685}{x}%
\TextAt{21}{215.729}{647.023}{p}%
\TextAt{9}{225.288}{642.685}{\Character{61}}%
\TextAt{13}{239.191}{642.685}{N}%
\TextAt{21}{248.565}{640.825}{L}%
\TextAt{3}{254.324}{640.825}{\Character{106}}%
\TextAt{21}{256.676}{640.825}{K}%
\TextAt{9}{264.879}{642.685}{\Character{40}}%
\TextAt{13}{269.431}{642.685}{x}%
\TextAt{9}{276.083}{642.685}{\Character{41}}%
\TextAt{8}{280.636}{642.685}{.}%
\TextAt{8}{291.685}{642.685}{Taking}%
\TextAt{13}{331.573}{642.685}{d}%
\TextAt{9}{342.454}{642.685}{\Character{61}}%
\TextAt{9}{356.346}{642.685}{\Character{91}}%
\TextAt{13}{359.609}{642.685}{K}%
\TextAt{3}{370.426}{647.023}{\Character{48}}%
\TextAt{9}{378.02}{642.685}{\Character{58}}%
\TextAt{13}{386.069}{642.685}{K}%
\TextAt{9}{396.887}{642.685}{\Character{93}}%
\TextAt{8}{400.139}{642.685}{-th}%
\TextAt{8}{419.859}{642.685}{powers}%
\TextAt{8}{459.485}{642.685}{on}%
\TextAt{13}{476.61}{642.685}{H}%
\TextAt{17}{487.233}{647.023}{0}%
\TextAt{9}{491.965}{642.685}{\Character{40}}%
\TextAt{13}{496.518}{642.685}{L}%
\TextAt{16}{509.28}{642.685}{\Character{106}}%
\TextAt{13}{517.399}{642.685}{K}%
\TextAt{9}{528.217}{642.685}{\Character{41}}%
\TextAt{8}{79.226}{625.639}{is}%
\TextAt{8}{92.138}{625.639}{therefore}%
\TextAt{8}{142.385}{625.639}{an}%
\TextAt{8}{159.792}{625.639}{isomorphism.}%
\TextAt{8}{236.206}{625.639}{Now}%
\TextAt{8}{264.013}{625.639}{let}%
\Vector{q 0 G 0 g 0.478 w 210.066 -141.54 m 216.718 -141.54 l S Q}%
\TextAt{13}{282.066}{625.639}{x}%
\TextAt{9}{293.995}{625.639}{\Character{61}}%
\TextAt{13}{308.376}{625.639}{x}%
\TextAt{8}{318.349}{625.639}{mod}%
\TextAt{13}{343.787}{625.639}{N}%
\TextAt{21}{353.161}{623.779}{L}%
\TextAt{3}{358.92}{623.779}{\Character{106}}%
\TextAt{21}{361.272}{623.779}{K}%
\TextAt{13}{369.475}{625.639}{C}%
\TextAt{21}{377.85}{623.846}{L}%
\TextAt{8}{389.152}{625.639}{belong}%
\TextAt{8}{427.695}{625.639}{to}%
\TextAt{8}{443.153}{625.639}{the}%
\TextAt{8}{464.458}{625.639}{kernel}%
\TextAt{8}{500.072}{625.639}{of}%
\TextAt{8}{514.55}{625.639}{(*).}%
\TextAt{8}{79.226}{608.593}{We}%
\TextAt{8}{99.882}{608.593}{write}%
\Vector{q 0 G 0 g 0.478 w 58.304 -158.586 m 64.956 -158.586 l S Q}%
\TextAt{13}{130.304}{608.593}{x}%
\TextAt{9}{141.148}{608.593}{\Character{61}}%
\Vector{q 0 G 0 g 0.478 w 82.444 -158.586 m 88.581 -158.586 l S Q}%
\TextAt{13}{154.444}{608.593}{y}%
\TextAt{21}{160.581}{612.976}{d}%
\TextAt{8}{165.436}{608.593}{,}%
\TextAt{8}{173.224}{608.593}{for}%
\TextAt{8}{191.62}{608.593}{some}%
\Vector{q 0 G 0 g 0.478 w 149.457 -158.586 m 155.593 -158.586 l S Q}%
\TextAt{13}{221.457}{608.593}{y}%
\TextAt{9}{231.785}{608.593}{\Character{61}}%
\TextAt{13}{245.081}{608.593}{y}%
\TextAt{8}{254.538}{608.593}{mod}%
\TextAt{13}{279.976}{608.593}{N}%
\TextAt{21}{289.35}{606.733}{L}%
\TextAt{3}{295.109}{606.733}{\Character{106}}%
\TextAt{21}{297.462}{606.733}{K}%
\TextAt{13}{305.664}{608.593}{C}%
\TextAt{21}{314.039}{606.8}{L}%
\TextAt{8}{320.296}{608.593}{.}%
\TextAt{8}{330.277}{608.593}{Then}%
\Vector{q 0 G 0 g 0.478 w 289.346 -158.586 m 295.482 -158.586 l S Q}%
\TextAt{13}{361.346}{608.593}{y}%
\TextAt{8}{371.891}{608.593}{also}%
\TextAt{8}{395.882}{608.593}{is}%
\TextAt{8}{408.156}{608.593}{in}%
\TextAt{8}{422.32}{608.593}{the}%
\TextAt{8}{442.987}{608.593}{kernel}%
\TextAt{8}{477.965}{608.593}{of}%
\TextAt{8}{491.806}{608.593}{(*),}%
\TextAt{8}{514.562}{608.593}{i.e.,}%
\TextAt{13}{79.226}{591.547}{y}%
\TextAt{9}{88.684}{591.547}{\Character{61}}%
\TextAt{13}{101.109}{591.547}{N}%
\TextAt{21}{110.483}{589.687}{L}%
\TextAt{23}{116.242}{591.956}{\Character{48}}%
\TextAt{3}{118.926}{589.687}{\Character{106}}%
\TextAt{21}{121.279}{589.687}{K}%
\TextAt{23}{128.983}{591.956}{\Character{48}}%
\TextAt{9}{132.166}{591.547}{\Character{40}}%
\TextAt{13}{136.718}{591.547}{z}%
\TextAt{3}{142.689}{595.885}{\Character{48}}%
\TextAt{9}{145.484}{591.547}{\Character{41}}%
\TextAt{8}{150.036}{591.547}{,}%
\TextAt{13}{157.185}{591.547}{z}%
\TextAt{3}{163.156}{595.885}{\Character{48}}%
\NativeMathAt{169.272}{591.547}{11.955}{\in}%
\TextAt{13}{180.563}{591.547}{C}%
\TextAt{21}{188.938}{589.754}{L}%
\TextAt{23}{194.697}{592.023}{\Character{48}}%
\TextAt{8}{197.879}{591.547}{,}%
\TextAt{8}{205.029}{591.547}{and}%
\TextAt{8}{227.791}{591.547}{we}%
\TextAt{8}{245.019}{591.547}{find:}%
\TextAt{13}{162.345}{566.419}{y}%
\TextAt{21}{168.482}{571.356}{d}%
\TextAt{9}{176.658}{566.419}{\Character{61}}%
\TextAt{13}{189.083}{566.419}{N}%
\TextAt{21}{198.457}{564.56}{K}%
\TextAt{23}{206.162}{566.829}{\Character{48}}%
\TextAt{3}{208.846}{564.56}{\Character{106}}%
\TextAt{21}{211.198}{564.56}{K}%
\TextAt{9}{219.401}{566.419}{\Character{40}}%
\TextAt{13}{223.953}{566.419}{y}%
\TextAt{9}{230.09}{566.419}{\Character{41}}%
\TextAt{9}{237.945}{566.419}{\Character{61}}%
\TextAt{13}{250.389}{566.419}{N}%
\TextAt{21}{259.763}{564.56}{L}%
\TextAt{23}{265.522}{566.829}{\Character{48}}%
\TextAt{3}{268.206}{564.56}{\Character{106}}%
\TextAt{21}{270.558}{564.56}{K}%
\TextAt{9}{278.761}{566.419}{\Character{40}}%
\TextAt{13}{283.313}{566.419}{z}%
\TextAt{3}{289.284}{571.356}{\Character{48}}%
\TextAt{9}{292.079}{566.419}{\Character{41}}%
\TextAt{9}{299.934}{566.419}{\Character{61}}%
\TextAt{13}{312.378}{566.419}{N}%
\TextAt{21}{321.751}{564.56}{L}%
\TextAt{3}{327.51}{564.56}{\Character{106}}%
\TextAt{21}{329.863}{564.56}{K}%
\TextAt{9}{338.065}{566.419}{\Character{40}}%
\TextAt{13}{342.618}{566.419}{N}%
\TextAt{21}{351.992}{564.56}{L}%
\TextAt{23}{357.751}{566.829}{\Character{48}}%
\TextAt{3}{360.435}{564.56}{\Character{106}}%
\TextAt{21}{362.787}{564.56}{L}%
\TextAt{9}{369.045}{566.419}{\Character{40}}%
\TextAt{13}{373.597}{566.419}{z}%
\TextAt{3}{379.567}{571.356}{\Character{48}}%
\TextAt{9}{382.363}{566.419}{\Character{41}\Character{41}}%
\NativeMathAt{394.788}{566.419}{11.955}{\in}%
\TextAt{13}{406.079}{566.419}{N}%
\TextAt{21}{415.453}{564.56}{L}%
\TextAt{3}{421.212}{564.56}{\Character{106}}%
\TextAt{21}{423.564}{564.56}{K}%
\TextAt{13}{431.767}{566.419}{C}%
\TextAt{21}{440.142}{564.626}{L}%
\TextAt{13}{446.399}{566.419}{\Character{58}}%
\TextAt{8}{79.226}{541.292}{Hence}%
\Vector{q 0 G 0 g 0.478 w 42.004 -225.887 m 48.656 -225.887 l S Q}%
\TextAt{13}{114.004}{541.292}{x}%
\TextAt{9}{123.977}{541.292}{\Character{61}}%
\Vector{q 0 G 0 g 0.478 w 64.402 -225.887 m 70.539 -225.887 l S Q}%
\TextAt{13}{136.402}{541.292}{y}%
\TextAt{21}{142.539}{545.675}{d}%
\TextAt{9}{150.715}{541.292}{\Character{61}}%
\TextAt{9}{163.125}{541.292}{1}%
\TextAt{8}{168.994}{541.292}{.}%
\Vector{q 0 G 0 g 0.399 w 452.334 -232.708 m 452.334 -224.638 l S Q}%
\Vector{q 0 G 0 g 0.399 w 452.534 -224.838 m 459.707 -224.838 l S Q}%
\Vector{q 0 G 0 g 0.399 w 452.534 -232.509 m 459.707 -232.509 l S Q}%
\Vector{q 0 G 0 g 0.399 w 459.906 -232.708 m 459.906 -224.638 l S Q}%
\TextAt{8}{97.159}{513.398}{An}%
\TextAt{8}{116.924}{513.398}{immediate}%
\TextAt{8}{174.743}{513.398}{consequence}%
\TextAt{8}{241.397}{513.398}{of}%
\TextAt{8}{255.317}{513.398}{the}%
\TextAt{8}{276.063}{513.398}{theorem}%
\TextAt{8}{322.177}{513.398}{we}%
\TextAt{8}{339.994}{513.398}{have}%
\TextAt{8}{367.578}{513.398}{just}%
\TextAt{8}{391.313}{513.398}{proved}%
\TextAt{8}{429.956}{513.398}{is}%
\TextAt{8}{442.309}{513.398}{the}%
\TextAt{8}{463.055}{513.398}{famous}%
\TextAt{8}{503.706}{513.398}{Hasse}%
\TextAt{8}{79.226}{496.352}{Norm}%
\TextAt{8}{112.067}{496.352}{Theorem:}%
\TextAt{12}{79.226}{471.446}{Corollary}%
\TextAt{12}{139.294}{471.446}{(4.5).}%
\TextAt{8}{174.854}{471.446}{Let}%
\TextAt{13}{196.067}{471.446}{L}%
\TextAt{16}{207.767}{471.446}{\Character{106}}%
\TextAt{13}{214.825}{471.446}{K}%
\TextAt{8}{229.783}{471.446}{be}%
\TextAt{8}{245.951}{471.446}{a}%
\TextAt{8}{255.951}{471.446}{cyclic}%
\TextAt{8}{288.378}{471.446}{extension.}%
\TextAt{8}{345.426}{471.446}{An}%
\TextAt{8}{364.845}{471.446}{element}%
\TextAt{13}{408.331}{471.446}{x}%
\NativeMathAt{418.719}{471.446}{11.955}{\in}%
\TextAt{13}{430.426}{471.446}{K}%
\TextAt{3}{441.243}{475.784}{\Character{3}}%
\TextAt{8}{450.117}{471.446}{is}%
\TextAt{8}{462.124}{471.446}{a}%
\TextAt{8}{472.124}{471.446}{norm}%
\TextAt{8}{502.937}{471.446}{if}%
\TextAt{8}{513.904}{471.446}{and}%
\TextAt{8}{79.226}{454.4}{only}%
\TextAt{8}{104.918}{454.4}{if}%
\TextAt{8}{115.642}{454.4}{it}%
\TextAt{8}{127.346}{454.4}{is}%
\TextAt{8}{139.11}{454.4}{a}%
\TextAt{8}{148.865}{454.4}{norm}%
\TextAt{8}{179.434}{454.4}{locally}%
\TextAt{8}{216.507}{454.4}{everywhere,}%
\TextAt{8}{280.563}{454.4}{i.e.,}%
\TextAt{8}{302.668}{454.4}{a}%
\TextAt{8}{312.424}{454.4}{norm}%
\TextAt{8}{342.993}{454.4}{in}%
\TextAt{8}{356.646}{454.4}{every}%
\TextAt{8}{387.538}{454.4}{completion}%
\TextAt{13}{447.374}{454.4}{L}%
\TextAt{27}{455.338}{452.607}{P}%
\TextAt{16}{465.753}{454.4}{\Character{106}}%
\TextAt{13}{472.394}{454.4}{K}%
\TextAt{27}{482.353}{452.607}{p}%
\TextAt{8}{490.738}{454.4}{(}%
\TextAt{14}{495.293}{454.4}{P}%
\TextAt{16}{508.507}{454.4}{\Character{106}}%
\TextAt{14}{515.148}{454.4}{p}%
\TextAt{8}{521.131}{454.4}{).}%
\TextAt{19}{79.226}{421.635}{Proof.}%
\TextAt{8}{116}{421.635}{Let}%
\TextAt{13}{137.988}{421.635}{G}%
\TextAt{9}{152.279}{421.635}{\Character{61}}%
\TextAt{13}{166.441}{421.635}{G}%
\TextAt{9}{175.674}{421.635}{\Character{40}}%
\TextAt{13}{180.227}{421.635}{L}%
\TextAt{16}{193.248}{421.635}{\Character{106}}%
\TextAt{13}{201.626}{421.635}{K}%
\TextAt{9}{212.444}{421.635}{\Character{41}}%
\TextAt{8}{221.912}{421.635}{and}%
\TextAt{13}{245.694}{421.635}{G}%
\TextAt{27}{254.927}{419.842}{P}%
\TextAt{9}{267.078}{421.635}{\Character{61}}%
\TextAt{13}{281.24}{421.635}{G}%
\TextAt{9}{290.473}{421.635}{\Character{40}}%
\TextAt{13}{295.026}{421.635}{L}%
\TextAt{27}{302.99}{419.842}{P}%
\TextAt{16}{315.141}{421.635}{\Character{106}}%
\TextAt{13}{323.519}{421.635}{K}%
\TextAt{27}{333.478}{419.842}{p}%
\TextAt{9}{337.965}{421.635}{\Character{41}}%
\TextAt{8}{342.517}{421.635}{.}%
\TextAt{8}{354.022}{421.635}{The}%
\TextAt{8}{379.094}{421.635}{exact}%
\TextAt{8}{411.005}{421.635}{sequence}%
\TextAt{9}{460.526}{421.635}{1}%
\NativeArrowAt{471.436}{421.635}{11.955}{19.26}{\longrightarrow}%
\TextAt{13}{495.754}{421.635}{L}%
\TextAt{3}{503.719}{425.973}{\Character{3}}%
\NativeArrowAt{513.508}{421.635}{11.955}{19.26}{\longrightarrow}%
\TextAt{13}{79.226}{404.589}{I}%
\TextAt{21}{84.389}{402.796}{L}%
\NativeArrowAt{93.967}{404.589}{11.955}{19.26}{\longrightarrow}%
\TextAt{13}{116.549}{404.589}{C}%
\TextAt{21}{124.924}{402.796}{L}%
\NativeArrowAt{134.502}{404.589}{11.955}{19.26}{\longrightarrow}%
\TextAt{9}{157.084}{404.589}{1}%
\TextAt{8}{166.835}{404.589}{of}%
\TextAt{13}{180.165}{404.589}{G}%
\TextAt{8}{189.398}{404.589}{-modules}%
\TextAt{8}{239.203}{404.589}{gives,}%
\TextAt{8}{271.136}{404.589}{by}%
\TextAt{8}{287.395}{404.589}{chap.}%
\TextAt{8}{318.287}{404.589}{IV,}%
\TextAt{8}{338.431}{404.589}{(7.1),}%
\TextAt{8}{369.658}{404.589}{an}%
\TextAt{8}{385.917}{404.589}{exact}%
\TextAt{8}{416.809}{404.589}{sequence}%
\TextAt{13}{197.395}{379.461}{H}%
\TextAt{3}{208.018}{384.398}{\Character{0}}%
\TextAt{17}{214.604}{384.398}{1}%
\TextAt{9}{219.337}{379.461}{\Character{40}}%
\TextAt{13}{223.889}{379.461}{G\Character{59}}%
\TextAt{13}{238.343}{379.461}{C}%
\TextAt{21}{246.742}{377.668}{L}%
\TextAt{9}{252.999}{379.461}{\Character{41}}%
\NativeArrowAt{260.872}{379.461}{11.955}{19.26}{\longrightarrow}%
\TextAt{13}{283.454}{379.461}{H}%
\TextAt{17}{294.077}{384.398}{0}%
\TextAt{9}{298.809}{379.461}{\Character{40}}%
\TextAt{13}{303.361}{379.461}{G\Character{59}}%
\TextAt{13}{317.815}{379.461}{L}%
\TextAt{3}{325.804}{384.398}{\Character{3}}%
\TextAt{9}{330.536}{379.461}{\Character{41}}%
\NativeArrowAt{338.409}{379.461}{11.955}{19.26}{\longrightarrow}%
\TextAt{13}{360.991}{379.461}{H}%
\TextAt{17}{371.614}{384.398}{0}%
\TextAt{9}{376.346}{379.461}{\Character{40}}%
\TextAt{13}{380.898}{379.461}{G\Character{59}}%
\TextAt{13}{395.352}{379.461}{I}%
\TextAt{21}{400.539}{377.668}{L}%
\TextAt{9}{406.796}{379.461}{\Character{41}}%
\TextAt{13}{411.349}{379.461}{\Character{58}}%
\TextAt{8}{79.226}{354.334}{By}%
\TextAt{8}{100.849}{354.334}{(4.4),}%
\TextAt{8}{136.151}{354.334}{we}%
\TextAt{8}{156.639}{354.334}{have}%
\TextAt{13}{186.893}{354.334}{H}%
\TextAt{3}{197.516}{358.672}{\Character{0}}%
\TextAt{17}{204.103}{358.672}{1}%
\TextAt{9}{208.835}{354.334}{\Character{40}}%
\TextAt{13}{213.387}{354.334}{G\Character{59}}%
\TextAt{13}{227.841}{354.334}{C}%
\TextAt{21}{236.24}{352.541}{L}%
\TextAt{9}{242.497}{354.334}{\Character{41}}%
\TextAt{9}{255.911}{354.334}{\Character{61}}%
\TextAt{9}{273.879}{354.334}{1}%
\TextAt{8}{279.76}{354.334}{,}%
\TextAt{8}{290.984}{354.334}{and}%
\TextAt{8}{317.007}{354.334}{from}%
\TextAt{8}{347.907}{354.334}{(3.2)}%
\TextAt{8}{379.142}{354.334}{it}%
\TextAt{8}{394.106}{354.334}{follows}%
\TextAt{8}{435.802}{354.334}{that}%
\TextAt{13}{464.431}{354.334}{H}%
\TextAt{17}{475.054}{358.672}{0}%
\TextAt{9}{479.786}{354.334}{\Character{40}}%
\TextAt{13}{484.338}{354.334}{G\Character{59}}%
\TextAt{13}{498.792}{354.334}{I}%
\TextAt{21}{503.979}{352.541}{L}%
\TextAt{9}{510.236}{354.334}{\Character{41}}%
\TextAt{9}{523.65}{354.334}{\Character{61}}%
\TextAt{20}{79.226}{346.254}{\Character{76}}%
\TextAt{27}{92.51}{333.801}{p}%
\TextAt{13}{98.989}{337.288}{H}%
\TextAt{17}{109.612}{341.626}{0}%
\TextAt{9}{114.344}{337.288}{\Character{40}}%
\TextAt{13}{118.896}{337.288}{G}%
\TextAt{27}{128.13}{335.495}{P}%
\TextAt{13}{135.224}{337.288}{\Character{59}}%
\TextAt{13}{140.448}{337.288}{L}%
\TextAt{3}{148.432}{341.626}{\Character{3}}%
\TextAt{27}{148.432}{334.203}{P}%
\TextAt{9}{155.526}{337.288}{\Character{41}}%
\TextAt{8}{160.078}{337.288}{.}%
\TextAt{8}{168.526}{337.288}{Therefore}%
\TextAt{8}{221.524}{337.288}{the}%
\TextAt{8}{241.68}{337.288}{homomorphism}%
\TextAt{13}{217.191}{310.663}{K}%
\TextAt{3}{228.008}{315.6}{\Character{3}}%
\TextAt{9}{232.741}{310.663}{\Character{47}}%
\TextAt{13}{238.594}{310.663}{N}%
\TextAt{21}{247.967}{308.804}{L}%
\TextAt{3}{253.726}{308.804}{\Character{106}}%
\TextAt{21}{256.079}{308.804}{K}%
\TextAt{13}{264.282}{310.663}{L}%
\TextAt{3}{272.246}{315.6}{\Character{3}}%
\NativeArrowAt{280.299}{310.663}{11.955}{19.26}{\longrightarrow}%
\TextAt{20}{302.881}{322.021}{\Character{77}}%
\TextAt{27}{309.92}{296.915}{p}%
\TextAt{13}{322.939}{310.663}{K}%
\TextAt{3}{333.757}{315.6}{\Character{3}}%
\TextAt{27}{332.898}{307.708}{p}%
\TextAt{9}{338.489}{310.663}{\Character{47}}%
\TextAt{13}{344.342}{310.663}{N}%
\TextAt{21}{353.716}{308.804}{L}%
\TextAt{28}{359.475}{307.413}{P}%
\TextAt{3}{365.672}{308.804}{\Character{106}}%
\TextAt{21}{368.025}{308.804}{K}%
\TextAt{28}{375.166}{307.807}{p}%
\TextAt{13}{379.747}{310.663}{L}%
\TextAt{3}{387.711}{315.599}{\Character{3}}%
\TextAt{27}{387.711}{307.708}{P}%
\TextAt{8}{79.226}{276.718}{is}%
\TextAt{8}{90.99}{276.718}{injective.}%
\TextAt{8}{142.035}{276.718}{But}%
\TextAt{8}{165.275}{276.718}{this}%
\TextAt{8}{188.098}{276.718}{is}%
\TextAt{8}{199.862}{276.718}{the}%
\TextAt{8}{220.018}{276.718}{claim}%
\TextAt{8}{251.233}{276.718}{of}%
\TextAt{8}{264.563}{276.718}{the}%
\TextAt{8}{284.72}{276.718}{corollary.}%
\Vector{q 0 G 0 g 0.399 w 452.334 -497.282 m 452.334 -489.212 l S Q}%
\Vector{q 0 G 0 g 0.399 w 452.534 -489.411 m 459.707 -489.411 l S Q}%
\Vector{q 0 G 0 g 0.399 w 452.534 -497.082 m 459.707 -497.082 l S Q}%
\Vector{q 0 G 0 g 0.399 w 459.906 -497.282 m 459.906 -489.212 l S Q}%
\TextAt{8}{97.159}{248.824}{It}%
\TextAt{8}{109.524}{248.824}{should}%
\TextAt{8}{146.349}{248.824}{be}%
\TextAt{8}{161.966}{248.824}{noted}%
\TextAt{8}{194.177}{248.824}{that}%
\TextAt{8}{219.238}{248.824}{cyclicity}%
\TextAt{8}{264.779}{248.824}{is}%
\TextAt{8}{276.236}{248.824}{crucial}%
\TextAt{8}{313.647}{248.824}{for}%
\TextAt{8}{331.224}{248.824}{Hasse’s}%
\TextAt{8}{371.744}{248.824}{norm}%
\TextAt{8}{402.006}{248.824}{theorem.}%
\TextAt{8}{451.98}{248.824}{In}%
\TextAt{8}{466.294}{248.824}{fact,}%
\TextAt{8}{492.385}{248.824}{whereas}%
\TextAt{8}{79.226}{231.778}{it}%
\TextAt{8}{91.109}{231.778}{is}%
\TextAt{8}{103.051}{231.778}{true}%
\TextAt{8}{127.941}{231.778}{by}%
\TextAt{8}{144.379}{231.778}{(3.2)}%
\TextAt{8}{172.533}{231.778}{that}%
\TextAt{8}{198.08}{231.778}{an}%
\TextAt{8}{214.518}{231.778}{element}%
\TextAt{13}{257.938}{231.778}{x}%
\NativeMathAt{268.215}{231.778}{11.955}{\in}%
\TextAt{13}{279.811}{231.778}{K}%
\TextAt{3}{290.628}{236.116}{\Character{3}}%
\TextAt{8}{299.437}{231.778}{which}%
\TextAt{8}{333.102}{231.778}{is}%
\TextAt{8}{345.044}{231.778}{everywhere}%
\TextAt{8}{406.027}{231.778}{locally}%
\TextAt{8}{443.278}{231.778}{a}%
\TextAt{8}{453.212}{231.778}{norm,}%
\TextAt{8}{487.257}{231.778}{is}%
\TextAt{8}{499.199}{231.778}{always}%
\TextAt{8}{79.226}{214.732}{the}%
\TextAt{8}{99.342}{214.732}{norm}%
\TextAt{8}{129.872}{214.732}{of}%
\TextAt{8}{143.162}{214.732}{some}%
\TextAt{8}{172.448}{214.732}{idèle}%
\TextAt{13}{199.713}{214.732}{\Character{11}}%
\TextAt{8}{211.092}{214.732}{of}%
\TextAt{13}{224.382}{214.732}{L}%
\TextAt{8}{232.347}{214.732}{,}%
\TextAt{8}{239.464}{214.732}{this}%
\TextAt{8}{262.246}{214.732}{need}%
\TextAt{8}{289.512}{214.732}{not}%
\TextAt{8}{310.286}{214.732}{be}%
\TextAt{8}{326.17}{214.732}{by}%
\TextAt{8}{342.389}{214.732}{any}%
\TextAt{8}{364.466}{214.732}{means}%
\TextAt{8}{400.255}{214.732}{a}%
\TextAt{8}{409.971}{214.732}{principal}%
\TextAt{8}{458.708}{214.732}{idèle,}%
\TextAt{8}{489.233}{214.732}{not}%
\TextAt{8}{510.007}{214.732}{even}%
\TextAt{8}{79.226}{197.686}{in}%
\TextAt{8}{92.879}{197.686}{the}%
\TextAt{8}{113.035}{197.686}{case}%
\TextAt{8}{137.806}{197.686}{of}%
\TextAt{8}{151.137}{197.686}{arbitrary}%
\TextAt{8}{200.906}{197.686}{abelian}%
\TextAt{8}{241.553}{197.686}{extensions.}%
\TextAt{12}{79.226}{172.78}{Exercise}%
\TextAt{12}{131.484}{172.78}{VI.4.1.}%
\TextAt{8}{174.42}{172.78}{Determine}%
\TextAt{8}{230.662}{172.78}{the}%
\TextAt{8}{249.999}{172.78}{norm}%
\TextAt{8}{279.748}{172.78}{group}%
\TextAt{13}{312.104}{172.78}{N}%
\TextAt{21}{321.478}{170.921}{L}%
\TextAt{3}{327.237}{170.921}{\Character{106}}%
\TextAt{21}{329.589}{170.921}{K}%
\TextAt{13}{337.792}{172.78}{C}%
\TextAt{21}{346.167}{170.987}{L}%
\TextAt{8}{355.502}{172.78}{for}%
\TextAt{8}{372.567}{172.78}{an}%
\TextAt{8}{388.006}{172.78}{arbitrary}%
\TextAt{8}{436.956}{172.78}{Kummer}%
\TextAt{8}{484.901}{172.78}{extension}%
\TextAt{8}{79.226}{155.734}{in}%
\TextAt{8}{92.879}{155.734}{a}%
\TextAt{8}{102.634}{155.734}{way}%
\TextAt{8}{126.377}{155.734}{analogous}%
\TextAt{8}{180.439}{155.734}{to}%
\TextAt{8}{194.749}{155.734}{the}%
\TextAt{8}{214.905}{155.734}{case}%
\TextAt{8}{239.676}{155.734}{treated}%
\TextAt{8}{280.001}{155.734}{in}%
\TextAt{8}{293.654}{155.734}{(4.2)}%
\TextAt{8}{321.629}{155.734}{where}%
\TextAt{13}{355.438}{155.734}{G}%
\TextAt{9}{364.672}{155.734}{\Character{40}}%
\TextAt{13}{369.224}{155.734}{L}%
\TextAt{16}{380.51}{155.734}{\Character{106}}%
\TextAt{13}{387.151}{155.734}{K}%
\TextAt{9}{397.969}{155.734}{\Character{41}}%
\NativeMathAt{405.842}{155.734}{11.955}{\cong}%
\TextAt{9}{418.462}{155.734}{\Character{40}}%
\TextAt{18}{423.014}{155.734}{Z}%
\TextAt{9}{430.984}{155.734}{\Character{47}}%
\TextAt{13}{436.837}{155.734}{p}%
\TextAt{21}{442.712}{160.073}{\Character{23}}%
\TextAt{18}{447.894}{155.734}{Z}%
\TextAt{9}{455.864}{155.734}{\Character{41}}%
\TextAt{21}{460.417}{160.073}{r}%
\TextAt{8}{464.971}{155.734}{.}%
\TextAt{12}{79.226}{125.958}{Exercise}%
\TextAt{12}{133.236}{125.958}{VI.4.2.}%
\TextAt{8}{177.696}{125.958}{Let}%
\TextAt{13}{199.369}{125.958}{\Character{16}}%
\TextAt{8}{209.986}{125.958}{be}%
\TextAt{8}{226.614}{125.958}{a}%
\TextAt{8}{237.073}{125.958}{primitive}%
\TextAt{13}{287.857}{125.958}{m}%
\TextAt{8}{298.097}{125.958}{-th}%
\TextAt{8}{317.654}{125.958}{root}%
\TextAt{8}{343.404}{125.958}{of}%
\TextAt{8}{357.438}{125.958}{unity.}%
\TextAt{8}{393.689}{125.958}{Show}%
\TextAt{8}{425.285}{125.958}{that}%
\TextAt{8}{451.358}{125.958}{the}%
\TextAt{8}{472.218}{125.958}{norm}%
\TextAt{8}{503.491}{125.958}{group}%
\ErrataAt{79.22}{108.955}{453.64}{$N_{\mathbb Q(\zeta)\mid\mathbb Q}C_{\mathbb Q(\zeta)}$ is the congruence subgroup $C_{\mathbb Q}^{\mathfrak m}$ for $\mathfrak m=(m)$.}%
\TextAt{8}{297.211}{43.033}{350}%
\Destination{522.36}{701.905}{nav.AMS.77}%
\Destination{78.23}{809.815}{nav.page.350}%
\Destination{79.226}{186.755}{nav.ex.6.4.1}%
\Destination{79.226}{139.955}{nav.ex.6.4.2}%
\Destination{79.36}{724.955}{nav.print.384}%
\Destination{79.36}{485.455}{nav.number.6.4.5}%
\Hotlink{317.787}{402.355}{13.995}{11.047}{nav.chapter.4}%
\Hotlink{245.932}{724.855}{25.078}{11.047}{nav.number.6.4.2}%
\Hotlink{337.931}{402.355}{25.078}{11.047}{nav.number.4.7.1}%
\Hotlink{100.349}{352.155}{25.078}{11.047}{nav.number.6.4.4}%
\Hotlink{347.407}{352.155}{25.078}{11.047}{nav.number.6.3.2}%
\Hotlink{143.879}{229.555}{25.078}{11.047}{nav.number.6.3.2}%
\Hotlink{293.154}{153.555}{25.078}{11.047}{nav.number.6.4.2}%
\Destination{79.226}{122.955}{errata.42}%
}

% Page 363 (body page 351)
\Page{%
\TextAt{2}{79.226}{796.943}{VI.}%
\TextAt{2}{98.012}{796.943}{Global}%
\TextAt{2}{133.248}{796.943}{Class}%
\TextAt{2}{161.841}{796.943}{Field}%
\TextAt{2}{189.572}{796.943}{Theory}%
\Vector{q 0 G 0 g 0.299 w 7.226 18.729 m 460.769 18.729 l S Q}%
\TextAt{12}{79.226}{761.124}{Exercise}%
\TextAt{12}{132.768}{761.124}{VI.4.3.}%
\TextAt{8}{176.82}{761.124}{An}%
\TextAt{8}{196.292}{761.124}{equation}%
\TextAt{13}{244.398}{761.124}{x}%
\TextAt{17}{251.05}{765.462}{2}%
\TextAt{16}{258.641}{761.124}{\Character{0}}%
\TextAt{13}{270.798}{761.124}{ay}%
\TextAt{17}{283.079}{765.462}{2}%
\TextAt{9}{291.638}{761.124}{\Character{61}}%
\TextAt{13}{304.568}{761.124}{b}%
\TextAt{8}{309.545}{761.124}{,}%
\TextAt{13}{317.065}{761.124}{a\Character{59}}%
\TextAt{13}{328.434}{761.124}{b}%
\NativeMathAt{337.257}{761.124}{11.955}{\in}%
\TextAt{13}{349.053}{761.124}{K}%
\TextAt{3}{359.871}{765.462}{\Character{3}}%
\TextAt{8}{364.603}{761.124}{,}%
\TextAt{8}{372.123}{761.124}{has}%
\TextAt{8}{393.293}{761.124}{a}%
\TextAt{8}{403.345}{761.124}{solution}%
\TextAt{8}{447.935}{761.124}{in}%
\TextAt{13}{461.885}{761.124}{K}%
\TextAt{8}{476.896}{761.124}{if}%
\TextAt{8}{487.916}{761.124}{and}%
\TextAt{8}{510.975}{761.124}{only}%
\TextAt{8}{79.226}{744.078}{if}%
\TextAt{8}{89.95}{744.078}{it}%
\TextAt{8}{101.654}{744.078}{is}%
\TextAt{8}{113.418}{744.078}{solvable}%
\TextAt{8}{157.389}{744.078}{everywhere}%
\TextAt{8}{218.193}{744.078}{locally,}%
\TextAt{8}{257.537}{744.078}{i.e.,}%
\TextAt{8}{279.643}{744.078}{in}%
\TextAt{8}{293.295}{744.078}{each}%
\TextAt{8}{319.633}{744.078}{completion}%
\TextAt{13}{379.468}{744.078}{K}%
\TextAt{27}{389.427}{742.284}{p}%
\TextAt{8}{393.914}{744.078}{.}%
\ErrataAt{97.1}{724.02}{437.9}{\textbf{Hint:} Use the norm from $K(\sqrt a)$ if $a\notin K^{*2}$. If $a=c^2$, take $x=(b+1)/2$ and $y=(b-1)/(2c)$.}%
\TextAt{12}{79.226}{692.053}{Exercise}%
\TextAt{12}{133.305}{692.053}{VI.4.4.}%
\ErrataAt{178}{692.008}{357}{Let $K$ have more than five elements and $\operatorname{char}K\ne2$. Suppose $a_i\in K^*$.}%
\ErrataAt{79.19}{674.962}{455.81}{If $a_1x_1^2+\cdots+a_nx_n^2=0$ has a nontrivial solution in $K$,}%
\TextAt{8}{79.226}{657.961}{then}%
\TextAt{8}{105.886}{657.961}{there}%
\TextAt{8}{135.798}{657.961}{is}%
\TextAt{8}{147.562}{657.961}{a}%
\TextAt{8}{157.317}{657.961}{representation}%
\TextAt{8}{233.806}{657.961}{of}%
\TextAt{8}{247.137}{657.961}{zero}%
\TextAt{8}{271.848}{657.961}{in}%
\TextAt{8}{285.501}{657.961}{which}%
\TextAt{8}{318.987}{657.961}{all}%
\TextAt{13}{335.246}{657.961}{x}%
\TextAt{21}{341.898}{656.168}{i}%
\NativeMathAt{348.6}{657.961}{11.955}{\neq}%
\TextAt{9}{361.01}{657.961}{0}%
\TextAt{8}{366.879}{657.961}{.}%
\TextAt{12}{97.159}{637.926}{Hint:}%
\TextAt{8}{134.488}{637.926}{If}%
\TextAt{13}{147.075}{637.926}{a\Character{24}}%
\TextAt{17}{158.91}{642.265}{2}%
\TextAt{9}{168.489}{637.926}{\Character{61}}%
\TextAt{13}{182.44}{637.926}{\Character{21}}%
\NativeMathAt{194.114}{637.926}{11.955}{\neq}%
\TextAt{9}{208.054}{637.926}{0}%
\TextAt{8}{213.918}{637.926}{,}%
\TextAt{13}{222.186}{637.926}{b}%
\NativeMathAt{232.01}{637.926}{11.955}{\neq}%
\TextAt{9}{245.95}{637.926}{0}%
\TextAt{8}{251.813}{637.926}{,}%
\TextAt{8}{260.082}{637.926}{then}%
\TextAt{8}{287.637}{637.926}{there}%
\TextAt{8}{318.444}{637.926}{are}%
\TextAt{8}{338.849}{637.926}{non-zero}%
\TextAt{8}{387.219}{637.926}{elements}%
\TextAt{13}{435.97}{637.926}{\Character{11}}%
\TextAt{8}{448.284}{637.926}{and}%
\TextAt{13}{471.942}{637.926}{\Character{12}}%
\TextAt{8}{484.006}{637.926}{such}%
\TextAt{8}{511.298}{637.926}{that}%
\TextAt{13}{79.226}{620.88}{a\Character{11}}%
\TextAt{17}{92.893}{625.219}{2}%
\TextAt{9}{100.282}{620.88}{\Character{43}}%
\TextAt{13}{112.043}{620.88}{b\Character{12}}%
\TextAt{17}{124.292}{625.219}{2}%
\TextAt{9}{132.345}{620.88}{\Character{61}}%
\TextAt{13}{144.77}{620.88}{\Character{21}}%
\TextAt{8}{151.599}{620.88}{.}%
\TextAt{8}{160.047}{620.88}{To}%
\TextAt{8}{177.274}{620.88}{prove}%
\TextAt{8}{208.824}{620.88}{this,}%
\TextAt{8}{234.898}{620.88}{multiply}%
\TextAt{8}{281.727}{620.88}{the}%
\TextAt{8}{301.883}{620.88}{identity}%
\TextAt{9}{248.362}{596.015}{\Character{40}}%
\TextAt{13}{252.914}{596.015}{t}%
\TextAt{16}{259.798}{596.015}{\Character{0}}%
\TextAt{9}{271.753}{596.015}{1\Character{41}}%
\TextAt{17}{282.158}{600.354}{2}%
\Vector{q 0 G 0 g 0.478 w 176.362 -183.084 m 214.891 -183.084 l S Q}%
\TextAt{9}{248.459}{579.727}{\Character{40}}%
\TextAt{13}{253.011}{579.727}{t}%
\TextAt{9}{259.895}{579.727}{\Character{43}}%
\TextAt{9}{271.647}{579.727}{1\Character{41}}%
\TextAt{17}{282.061}{583.181}{2}%
\TextAt{9}{290.743}{587.928}{\Character{43}}%
\TextAt{9}{317.827}{596.015}{4}%
\TextAt{13}{323.68}{596.015}{t}%
\Vector{q 0 G 0 g 0.478 w 231.7 -183.084 m 270.035 -183.084 l S Q}%
\TextAt{9}{303.7}{579.727}{\Character{40}}%
\TextAt{13}{308.252}{579.727}{t}%
\TextAt{9}{315.136}{579.727}{\Character{43}}%
\TextAt{9}{326.888}{579.727}{1\Character{41}}%
\TextAt{17}{337.302}{583.181}{2}%
\TextAt{9}{346.551}{587.928}{\Character{61}}%
\TextAt{9}{358.961}{587.928}{1}%
\TextAt{8}{79.226}{554.312}{by}%
\TextAt{13}{95.807}{554.312}{a\Character{24}}%
\TextAt{17}{107.642}{558.65}{2}%
\TextAt{9}{116.244}{554.312}{\Character{61}}%
\TextAt{13}{129.218}{554.312}{\Character{21}}%
\TextAt{8}{140.266}{554.312}{and}%
\TextAt{8}{163.35}{554.312}{insert}%
\TextAt{13}{196.25}{554.312}{t}%
\TextAt{9}{204.346}{554.312}{\Character{61}}%
\TextAt{13}{217.32}{554.312}{b\Character{13}}%
\TextAt{17}{229.02}{558.65}{2}%
\TextAt{9}{233.752}{554.312}{\Character{47}}%
\TextAt{13}{239.605}{554.312}{a}%
\TextAt{8}{245.75}{554.312}{,}%
\TextAt{8}{253.301}{554.312}{for}%
\TextAt{8}{271.508}{554.312}{some}%
\TextAt{8}{301.156}{554.312}{element}%
\TextAt{13}{344.72}{554.312}{\Character{13}}%
\NativeMathAt{355.311}{554.312}{11.955}{\neq}%
\TextAt{9}{368.27}{554.312}{0}%
\TextAt{8}{378.357}{554.312}{such}%
\TextAt{8}{405.076}{554.312}{that}%
\TextAt{13}{430.767}{554.312}{t}%
\NativeMathAt{438.863}{554.312}{11.955}{\neq}%
\NativeMathAt{451.837}{554.312}{11.955}{\pm}%
\TextAt{9}{461.136}{554.312}{1}%
\TextAt{8}{466.989}{554.312}{.}%
\TextAt{8}{476.403}{554.312}{Use}%
\TextAt{8}{499.212}{554.312}{this}%
\TextAt{8}{522.357}{554.312}{to}%
\TextAt{8}{79.226}{537.266}{prove}%
\TextAt{8}{110.776}{537.266}{the}%
\TextAt{8}{130.932}{537.266}{claim}%
\TextAt{8}{162.147}{537.266}{by}%
\TextAt{8}{178.406}{537.266}{induction.}%
\TextAt{12}{79.226}{505.276}{Exercise}%
\TextAt{12}{132.802}{505.276}{VI.4.5.}%
\TextAt{8}{176.884}{505.276}{A}%
\TextAt{8}{189.882}{505.276}{quadratic}%
\TextAt{8}{242.571}{505.276}{form}%
\TextAt{13}{270.538}{505.276}{ax}%
\TextAt{17}{283.335}{509.614}{2}%
\TextAt{9}{290.946}{505.276}{\Character{43}}%
\TextAt{13}{302.93}{505.276}{by}%
\TextAt{17}{314.043}{509.614}{2}%
\TextAt{9}{321.654}{505.276}{\Character{43}}%
\TextAt{13}{333.638}{505.276}{cz}%
\TextAt{17}{344.647}{509.614}{2}%
\TextAt{8}{349.379}{505.276}{,}%
\TextAt{13}{356.936}{505.276}{a\Character{59}}%
\TextAt{13}{368.305}{505.276}{b\Character{59}}%
\TextAt{13}{378.515}{505.276}{c}%
\NativeMathAt{387.46}{505.276}{11.955}{\in}%
\TextAt{13}{399.307}{505.276}{K}%
\TextAt{3}{410.125}{509.614}{\Character{3}}%
\TextAt{8}{414.857}{505.276}{,}%
\TextAt{8}{422.414}{505.276}{represents}%
\TextAt{8}{477.817}{505.276}{zero}%
\TextAt{8}{502.854}{505.276}{if}%
\TextAt{8}{513.904}{505.276}{and}%
\TextAt{8}{79.226}{488.23}{only}%
\TextAt{8}{104.918}{488.23}{if}%
\TextAt{8}{115.642}{488.23}{it}%
\TextAt{8}{127.346}{488.23}{represents}%
\TextAt{8}{182.423}{488.23}{zero}%
\TextAt{8}{207.135}{488.23}{everywhere}%
\TextAt{8}{267.939}{488.23}{locally.}%
\TextAt{12}{97.159}{468.195}{Remark:}%
\TextAt{8}{152.615}{468.195}{In}%
\TextAt{8}{167.236}{468.195}{complete}%
\TextAt{8}{216.659}{468.195}{generality,}%
\TextAt{8}{272.92}{468.195}{one}%
\TextAt{8}{294.379}{468.195}{has}%
\TextAt{8}{315.253}{468.195}{the}%
\TextAt{8}{335.41}{468.195}{following}%
\TextAt{8}{384.844}{468.195}{“local-to-global}%
\TextAt{8}{466.797}{468.195}{principle”:}%
\TextAt{12}{97.159}{448.161}{Theorem}%
\TextAt{12}{153.058}{448.161}{of}%
\TextAt{12}{167.27}{448.161}{Minkowski–Hasse:}%
\TextAt{8}{279.196}{448.161}{A}%
\TextAt{8}{290.899}{448.161}{quadratic}%
\TextAt{8}{342.293}{448.161}{form}%
\TextAt{8}{368.965}{448.161}{over}%
\TextAt{8}{393.042}{448.161}{a}%
\TextAt{8}{401.828}{448.161}{number}%
\TextAt{8}{443.455}{448.161}{field}%
\TextAt{13}{467.843}{448.161}{K}%
\TextAt{8}{481.589}{448.161}{represents}%
\TextAt{8}{79.226}{431.115}{zero}%
\TextAt{8}{104.641}{431.115}{if}%
\TextAt{8}{116.069}{431.115}{and}%
\TextAt{8}{139.536}{431.115}{only}%
\TextAt{8}{165.931}{431.115}{if}%
\TextAt{8}{177.359}{431.115}{it}%
\TextAt{8}{189.767}{431.115}{represents}%
\TextAt{8}{245.549}{431.115}{zero}%
\TextAt{8}{270.964}{431.115}{over}%
\TextAt{8}{296.714}{431.115}{every}%
\TextAt{8}{328.31}{431.115}{completion}%
\TextAt{13}{388.85}{431.115}{K}%
\TextAt{27}{398.809}{429.321}{p}%
\TextAt{8}{403.295}{431.115}{.}%
\TextAt{8}{413.856}{431.115}{The}%
\TextAt{8}{438.613}{431.115}{proof}%
\TextAt{8}{469.887}{431.115}{follows}%
\TextAt{8}{509.027}{431.115}{from}%
\TextAt{8}{79.226}{414.069}{the}%
\TextAt{8}{99.383}{414.069}{result}%
\TextAt{8}{131.96}{414.069}{stated}%
\TextAt{8}{167.144}{414.069}{in}%
\TextAt{8}{180.797}{414.069}{exercise}%
\TextAt{8}{224.099}{414.069}{5}%
\TextAt{8}{233.854}{414.069}{by}%
\TextAt{8}{250.113}{414.069}{pure}%
\TextAt{8}{276.773}{414.069}{algebra}%
\TextAt{8}{317.756}{414.069}{(see}%
\TextAt{8}{341.224}{414.069}{[113]).}%
\TextAt{15}{79.226}{361.205}{§}%
\TextAt{15}{93.785}{361.205}{5}%
\TextAt{15}{120.693}{361.205}{The}%
\TextAt{15}{160.202}{361.205}{Global}%
\TextAt{15}{222.522}{361.205}{Reciprocity}%
\TextAt{15}{326.383}{361.205}{Law}%
\TextAt{8}{79.226}{329.319}{Now}%
\TextAt{8}{105.587}{329.319}{that}%
\TextAt{8}{130.657}{329.319}{we}%
\TextAt{8}{147.586}{329.319}{know}%
\TextAt{8}{177.856}{329.319}{that}%
\TextAt{8}{202.926}{329.319}{the}%
\TextAt{8}{222.784}{329.319}{idèle}%
\TextAt{8}{249.791}{329.319}{class}%
\TextAt{8}{276.929}{329.319}{group}%
\TextAt{8}{309.806}{329.319}{satisfies}%
\TextAt{8}{352.617}{329.319}{the}%
\TextAt{8}{372.475}{329.319}{class}%
\TextAt{8}{399.613}{329.319}{field}%
\TextAt{8}{424.671}{329.319}{axiom,}%
\TextAt{8}{462.485}{329.319}{we}%
\TextAt{8}{479.414}{329.319}{proceed}%
\TextAt{8}{522.357}{329.319}{to}%
\TextAt{8}{79.226}{312.273}{determine}%
\TextAt{8}{133.849}{312.273}{a}%
\TextAt{8}{143.605}{312.273}{pair}%
\TextAt{8}{167.67}{312.273}{of}%
\TextAt{8}{181.001}{312.273}{homomorphisms}%
\TextAt{9}{254.806}{283.272}{\Character{40}}%
\TextAt{13}{259.358}{283.272}{G}%
\TextAt{25}{268.592}{281.479}{Q}%
\TextAt{21}{281.703}{290.049}{d}%
\NativeArrowAt{278.61}{283.272}{11.955}{11.955}{\longrightarrow}%
\TextAt{20}{294.55}{286.36}{\Character{98}}%
\TextAt{18}{293.886}{283.272}{Z}%
\TextAt{13}{301.856}{283.272}{\Character{59}}%
\TextAt{13}{310.978}{283.272}{C}%
\TextAt{25}{319.373}{281.479}{Q}%
\TextAt{21}{332.458}{290.049}{v}%
\NativeArrowAt{329.391}{283.272}{11.955}{11.955}{\longrightarrow}%
\TextAt{20}{345.331}{286.36}{\Character{98}}%
\TextAt{18}{344.667}{283.272}{Z}%
\TextAt{9}{352.637}{283.272}{\Character{41}}%
\TextAt{8}{79.226}{254.271}{obeying}%
\TextAt{8}{124.141}{254.271}{the}%
\TextAt{8}{145.635}{254.271}{rules}%
\TextAt{8}{174.996}{254.271}{of}%
\TextAt{8}{189.663}{254.271}{abstract}%
\TextAt{8}{236.598}{254.271}{class}%
\TextAt{8}{265.373}{254.271}{field}%
\TextAt{8}{292.068}{254.271}{theory}%
\TextAt{8}{330.156}{254.271}{as}%
\TextAt{8}{345.864}{254.271}{developed}%
\TextAt{8}{401.502}{254.271}{in}%
\TextAt{8}{416.493}{254.271}{chap.}%
\TextAt{8}{448.723}{254.271}{IV,}%
\TextAt{8}{470.205}{254.271}{§4.}%
\TextAt{8}{494.203}{254.271}{For}%
\TextAt{8}{516.51}{254.271}{the}%
\TextAt{20}{79.89}{240.313}{\Character{98}}%
\TextAt{18}{79.226}{237.225}{Z}%
\TextAt{8}{87.196}{237.225}{-extension}%
\TextAt{8}{142.86}{237.225}{of}%
\TextAt{18}{156.19}{237.225}{Q}%
\TextAt{8}{169.385}{237.225}{given}%
\TextAt{8}{199.955}{237.225}{by}%
\TextAt{13}{216.214}{237.225}{d}%
\TextAt{8}{222.297}{237.225}{,}%
\TextAt{8}{229.446}{237.225}{we}%
\TextAt{8}{246.673}{237.225}{have}%
\TextAt{8}{273.668}{237.225}{only}%
\TextAt{8}{299.36}{237.225}{one}%
\TextAt{8}{320.819}{237.225}{choice.}%
\TextAt{8}{360.16}{237.225}{It}%
\TextAt{8}{372.832}{237.225}{is}%
\TextAt{8}{384.596}{237.225}{described}%
\TextAt{8}{436.35}{237.225}{in}%
\TextAt{8}{450.003}{237.225}{the}%
\TextAt{8}{470.159}{237.225}{following}%
\TextAt{12}{79.226}{211.213}{Proposition}%
\TextAt{12}{152.266}{211.213}{(5.1).}%
\TextAt{8}{187.868}{211.213}{Let}%
\TextAt{9}{209.123}{211.213}{\Character{10}}%
\TextAt{16}{221.385}{211.213}{\Character{106}}%
\TextAt{18}{228.514}{211.213}{Q}%
\TextAt{8}{241.995}{211.213}{be}%
\TextAt{8}{258.206}{211.213}{the}%
\TextAt{8}{278.648}{211.213}{field}%
\TextAt{8}{304.29}{211.213}{obtained}%
\TextAt{8}{352.708}{211.213}{by}%
\TextAt{8}{369.253}{211.213}{adjoining}%
\TextAt{8}{420.599}{211.213}{all}%
\TextAt{8}{437.144}{211.213}{roots}%
\TextAt{8}{467.09}{211.213}{of}%
\TextAt{8}{480.706}{211.213}{unity,}%
\TextAt{8}{513.904}{211.213}{and}%
\ErrataAt{79.22}{194.155}{453.64}{let $T$ be the torsion subgroup of $G(\Omega\mid\mathbb Q)$ (the group of all elements of finite order).}%
\TextAt{8}{79.226}{177.121}{Then}%
\TextAt{8}{109.784}{177.121}{the}%
\TextAt{8}{129.94}{177.121}{fixed}%
\TextAt{8}{158.226}{177.121}{field}%
\TextAt{20}{184.911}{180.209}{\Character{101}}%
\TextAt{18}{183.583}{177.121}{Q}%
\TextAt{16}{196.202}{177.121}{\Character{106}}%
\TextAt{18}{202.844}{177.121}{Q}%
\TextAt{8}{216.04}{177.121}{of}%
\TextAt{13}{229.37}{177.121}{T}%
\TextAt{8}{241.754}{177.121}{is}%
\TextAt{8}{253.518}{177.121}{a}%
\TextAt{20}{263.937}{180.209}{\Character{98}}%
\TextAt{18}{263.273}{177.121}{Z}%
\TextAt{8}{271.243}{177.121}{-extension.}%
\TextAt{8}{297.211}{43.033}{351}%
\Destination{78.23}{809.815}{nav.page.351}%
\Destination{79.23}{387.055}{nav.section.6.5}%
\Destination{79.226}{775.155}{nav.ex.6.4.3}%
\Destination{79.226}{706.055}{nav.ex.6.4.4}%
\Destination{79.226}{519.255}{nav.ex.6.4.5}%
\Destination{79.36}{776.155}{nav.print.385}%
\Destination{79.36}{225.255}{nav.number.6.5.1}%
\Destination{79.23}{387.055}{outline.60}%
\Hotlink{223.599}{411.855}{6.858}{11.047}{nav.ex.6.4.5}%
\Hotlink{340.724}{411.855}{25.078}{11.047}{nav.cite.N113}%
\Hotlink{448.223}{252.055}{13.995}{11.047}{nav.chapter.4}%
\Hotlink{475.384}{252.055}{6.858}{11.047}{nav.section.4.4}%
\Destination{79.226}{208.155}{errata.30}%
\Destination{79.226}{741.02}{errata.135}%
\Destination{79.226}{709.008}{errata.136}%
}

% Page 364 (body page 352)
\Page{%
\TextAt{2}{79.226}{796.943}{VI.}%
\TextAt{2}{98.012}{796.943}{Global}%
\TextAt{2}{133.248}{796.943}{Class}%
\TextAt{2}{161.841}{796.943}{Field}%
\TextAt{2}{189.572}{796.943}{Theory}%
\Vector{q 0 G 0 g 0.299 w 7.226 18.729 m 460.769 18.729 l S Q}%
\TextAt{19}{79.226}{761.124}{Proof.}%
\TextAt{8}{116}{761.124}{Since}%
\TextAt{9}{146.558}{761.124}{\Character{10}}%
\TextAt{9}{158.322}{761.124}{\Character{61}}%
\TextAt{20}{170.758}{770.09}{\Character{83}}%
\TextAt{21}{180.721}{757.637}{n}%
\NativeMathAt{185.859}{757.637}{7.97}{\geq}%
\TextAt{17}{192.446}{757.637}{1}%
\TextAt{18}{199.17}{761.124}{Q}%
\TextAt{9}{208.469}{761.124}{\Character{40}}%
\TextAt{13}{213.021}{761.124}{\Character{22}}%
\TextAt{21}{220.064}{759.33}{n}%
\TextAt{9}{225.701}{761.124}{\Character{41}}%
\TextAt{8}{230.253}{761.124}{,}%
\TextAt{8}{237.402}{761.124}{we}%
\TextAt{8}{254.63}{761.124}{find}%
\TextAt{13}{178.469}{732.123}{G}%
\TextAt{9}{187.703}{732.123}{\Character{40}\Character{10}}%
\TextAt{16}{204.03}{732.123}{\Character{106}}%
\TextAt{18}{210.672}{732.123}{Q}%
\TextAt{9}{219.971}{732.123}{\Character{41}}%
\TextAt{9}{227.826}{732.123}{\Character{61}}%
\NativeMathAt{240.269}{732.123}{11.955}{\varprojlim}%
\TextAt{21}{245.83}{720.008}{n}%
\TextAt{13}{258.521}{732.123}{G}%
\TextAt{9}{267.755}{732.123}{\Character{40}}%
\TextAt{18}{272.307}{732.123}{Q}%
\TextAt{9}{281.605}{732.123}{\Character{40}}%
\TextAt{13}{286.158}{732.123}{\Character{22}}%
\TextAt{21}{293.201}{730.329}{n}%
\TextAt{9}{298.837}{732.123}{\Character{41}}%
\TextAt{16}{306.71}{732.123}{\Character{106}}%
\TextAt{18}{313.352}{732.123}{Q}%
\TextAt{9}{322.65}{732.123}{\Character{41}}%
\TextAt{9}{330.505}{732.123}{\Character{61}}%
\NativeMathAt{342.949}{732.123}{11.955}{\varprojlim}%
\TextAt{21}{348.509}{720.008}{n}%
\TextAt{9}{359.208}{732.123}{\Character{40}}%
\TextAt{18}{363.76}{732.123}{Z}%
\TextAt{9}{371.73}{732.123}{\Character{47}}%
\TextAt{13}{377.583}{732.123}{n}%
\TextAt{18}{384.571}{732.123}{Z}%
\TextAt{9}{392.541}{732.123}{\Character{41}}%
\TextAt{3}{397.094}{737.059}{\Character{3}}%
\TextAt{9}{405.147}{732.123}{\Character{61}}%
\TextAt{20}{418.236}{735.211}{\Character{98}}%
\TextAt{18}{417.572}{732.123}{Z}%
\TextAt{3}{425.542}{737.059}{\Character{3}}%
\TextAt{13}{430.275}{732.123}{\Character{58}}%
\TextAt{8}{79.226}{694.138}{But}%
\TextAt{20}{102.808}{697.227}{\Character{98}}%
\TextAt{18}{102.144}{694.138}{Z}%
\TextAt{9}{113.435}{694.138}{\Character{61}}%
\TextAt{20}{125.86}{703.105}{\Character{81}}%
\TextAt{21}{137.151}{690.651}{p}%
\TextAt{18}{143.904}{694.138}{Z}%
\TextAt{21}{151.875}{692.345}{p}%
\TextAt{8}{156.635}{694.138}{,}%
\TextAt{8}{163.526}{694.138}{and}%
\TextAt{18}{185.965}{694.138}{Z}%
\TextAt{3}{193.935}{698.477}{\Character{3}}%
\TextAt{21}{193.935}{691.183}{p}%
\NativeMathAt{202.017}{694.138}{11.955}{\cong}%
\TextAt{18}{214.636}{694.138}{Z}%
\TextAt{21}{222.607}{692.345}{p}%
\TextAt{16}{229.363}{694.138}{\Character{2}}%
\TextAt{18}{240.657}{694.138}{Z}%
\TextAt{9}{248.627}{694.138}{\Character{47}\Character{40}}%
\TextAt{13}{259.032}{694.138}{p}%
\TextAt{16}{266.902}{694.138}{\Character{0}}%
\TextAt{9}{278.196}{694.138}{1\Character{41}}%
\TextAt{18}{288.602}{694.138}{Z}%
\TextAt{8}{300.146}{694.138}{for}%
\TextAt{13}{317.707}{694.138}{p}%
\NativeMathAt{326.903}{694.138}{11.955}{\neq}%
\TextAt{9}{339.312}{694.138}{2}%
\TextAt{8}{348.756}{694.138}{and}%
\TextAt{18}{371.195}{694.138}{Z}%
\TextAt{3}{379.165}{698.477}{\Character{3}}%
\TextAt{17}{379.165}{691.183}{2}%
\NativeMathAt{387.218}{694.138}{11.955}{\cong}%
\TextAt{18}{399.837}{694.138}{Z}%
\TextAt{17}{407.808}{692.345}{2}%
\TextAt{16}{414.535}{694.138}{\Character{2}}%
\TextAt{18}{425.829}{694.138}{Z}%
\TextAt{9}{433.799}{694.138}{\Character{47}2}%
\TextAt{18}{445.505}{694.138}{Z}%
\TextAt{8}{453.475}{694.138}{.}%
\TextAt{8}{461.816}{694.138}{Consequently,}%
\TextAt{13}{158.443}{663.597}{G}%
\TextAt{9}{167.677}{663.597}{\Character{40}\Character{10}}%
\TextAt{16}{184.004}{663.597}{\Character{106}}%
\TextAt{18}{190.646}{663.597}{Q}%
\TextAt{9}{199.945}{663.597}{\Character{41}}%
\NativeMathAt{207.818}{663.597}{11.955}{\cong}%
\TextAt{20}{221.101}{666.686}{\Character{98}}%
\TextAt{18}{220.437}{663.597}{Z}%
\TextAt{3}{228.407}{668.533}{\Character{3}}%
\NativeMathAt{236.46}{663.597}{11.955}{\cong}%
\TextAt{20}{249.744}{666.686}{\Character{98}}%
\TextAt{18}{249.08}{663.597}{Z}%
\TextAt{16}{259.707}{663.597}{\Character{2}}%
\TextAt{20}{273.56}{666.619}{\Character{98}}%
\TextAt{13}{271.662}{663.597}{T}%
\TextAt{13}{280.126}{663.597}{\Character{59}}%
\TextAt{20}{311.201}{666.619}{\Character{98}}%
\TextAt{13}{309.303}{663.597}{T}%
\TextAt{9}{321.111}{663.597}{\Character{61}}%
\TextAt{20}{333.536}{674.955}{\Character{89}}%
\TextAt{21}{333.632}{649.495}{p}%
\NativeMathAt{337.895}{649.495}{7.97}{\neq}%
\TextAt{17}{344.478}{649.495}{2}%
\TextAt{18}{350.805}{663.597}{Z}%
\TextAt{9}{358.775}{663.597}{\Character{47}\Character{40}}%
\TextAt{13}{369.18}{663.597}{p}%
\TextAt{16}{377.712}{663.597}{\Character{0}}%
\TextAt{9}{389.667}{663.597}{1\Character{41}}%
\TextAt{18}{400.072}{663.597}{Z}%
\TextAt{16}{410.699}{663.597}{\Character{2}}%
\TextAt{18}{422.654}{663.597}{Z}%
\TextAt{9}{430.625}{663.597}{\Character{47}2}%
\TextAt{18}{442.331}{663.597}{Z}%
\TextAt{13}{450.301}{663.597}{\Character{58}}%
\TextAt{8}{79.226}{624.831}{This}%
\TextAt{8}{105.278}{624.831}{shows}%
\TextAt{8}{138.228}{624.831}{that}%
\TextAt{8}{162.929}{624.831}{the}%
\TextAt{8}{182.418}{624.831}{torsion}%
\TextAt{8}{220.843}{624.831}{subgroup}%
\TextAt{13}{270.973}{624.831}{T}%
\TextAt{8}{282.689}{624.831}{of}%
\TextAt{13}{295.352}{624.831}{G}%
\TextAt{9}{304.585}{624.831}{\Character{40}\Character{10}}%
\TextAt{16}{320.913}{624.831}{\Character{106}}%
\TextAt{18}{327.554}{624.831}{Q}%
\TextAt{9}{336.853}{624.831}{\Character{41}}%
\TextAt{8}{344.635}{624.831}{is}%
\TextAt{8}{355.731}{624.831}{isomorphic}%
\TextAt{8}{414.313}{624.831}{to}%
\TextAt{8}{427.955}{624.831}{the}%
\TextAt{8}{447.444}{624.831}{torsion}%
\TextAt{8}{485.869}{624.831}{subgroup}%
\TextAt{8}{79.226}{607.785}{of}%
\TextAt{20}{95.996}{610.808}{\Character{98}}%
\TextAt{13}{94.098}{607.785}{T}%
\TextAt{8}{102.584}{607.785}{.}%
\TextAt{8}{115.657}{607.785}{Since}%
\TextAt{8}{147.756}{607.785}{the}%
\TextAt{8}{169.454}{607.785}{latter}%
\TextAt{8}{202.868}{607.785}{contains}%
\TextAt{8}{250.329}{607.785}{the}%
\TextAt{8}{272.027}{607.785}{group}%
\TextAt{20}{306.744}{616.752}{\Character{76}}%
\TextAt{21}{320.028}{604.299}{p}%
\NativeMathAt{324.29}{604.298}{7.97}{\neq}%
\TextAt{17}{330.873}{604.299}{2}%
\TextAt{18}{337.602}{607.785}{Z}%
\TextAt{9}{345.572}{607.786}{\Character{47}\Character{40}}%
\TextAt{13}{355.977}{607.785}{p}%
\TextAt{16}{365.56}{607.786}{\Character{0}}%
\TextAt{9}{378.566}{607.785}{1\Character{41}}%
\TextAt{18}{388.971}{607.786}{Z}%
\TextAt{16}{400.649}{607.785}{\Character{8}}%
\TextAt{18}{413.654}{607.786}{Z}%
\TextAt{9}{421.625}{607.785}{\Character{47}2}%
\TextAt{18}{433.331}{607.786}{Z}%
\TextAt{8}{441.301}{607.785}{,}%
\TextAt{8}{450.377}{607.785}{we}%
\TextAt{8}{469.146}{607.785}{see}%
\TextAt{8}{489.6}{607.785}{that}%
\TextAt{8}{516.511}{607.785}{the}%
\TextAt{8}{79.226}{590.74}{closure}%
\Vector{q 0 G 0 g 0.478 w 47.318 -173.417 m 55.805 -173.417 l S Q}%
\TextAt{13}{119.318}{590.74}{T}%
\TextAt{8}{132.713}{590.74}{of}%
\TextAt{13}{147.053}{590.74}{T}%
\TextAt{8}{160.447}{590.74}{is}%
\TextAt{8}{173.222}{590.74}{isomorphic}%
\TextAt{8}{233.482}{590.74}{to}%
\TextAt{20}{250.701}{593.762}{\Character{98}}%
\TextAt{13}{248.803}{590.74}{T}%
\TextAt{8}{257.289}{590.74}{.}%
\TextAt{8}{268.769}{590.74}{Now,}%
\TextAt{8}{299.944}{590.74}{if}%
\TextAt{20}{313.006}{593.828}{\Character{101}}%
\TextAt{18}{311.678}{590.74}{Q}%
\TextAt{8}{325.884}{590.74}{is}%
\TextAt{8}{338.658}{590.74}{the}%
\TextAt{8}{359.825}{590.74}{fixed}%
\TextAt{8}{389.121}{590.74}{field}%
\TextAt{8}{415.488}{590.74}{of}%
\TextAt{13}{429.829}{590.74}{T}%
\TextAt{8}{438.316}{590.74}{,}%
\TextAt{8}{446.728}{590.74}{this}%
\TextAt{8}{470.561}{590.74}{implies}%
\TextAt{8}{511.298}{590.74}{that}%
\TextAt{13}{79.226}{573.694}{G}%
\TextAt{9}{88.46}{573.694}{\Character{40}}%
\TextAt{20}{94.34}{576.782}{\Character{101}}%
\TextAt{18}{93.012}{573.694}{Q}%
\TextAt{16}{105.631}{573.694}{\Character{106}}%
\TextAt{18}{112.273}{573.694}{Q}%
\TextAt{9}{121.572}{573.694}{\Character{41}}%
\TextAt{9}{129.427}{573.694}{\Character{61}}%
\TextAt{13}{141.87}{573.694}{G}%
\TextAt{9}{151.104}{573.694}{\Character{40}\Character{10}}%
\TextAt{16}{167.431}{573.694}{\Character{106}}%
\TextAt{18}{174.073}{573.694}{Q}%
\TextAt{9}{183.372}{573.694}{\Character{41}\Character{47}}%
\Vector{q 0 G 0 g 0.478 w 121.777 -190.463 m 130.264 -190.463 l S Q}%
\TextAt{13}{193.777}{573.694}{T}%
\NativeMathAt{205.585}{573.693}{11.955}{\cong}%
\TextAt{20}{218.868}{576.782}{\Character{98}}%
\TextAt{18}{218.204}{573.694}{Z}%
\TextAt{8}{226.174}{573.694}{.}%
\Vector{q 0 G 0 g 0.399 w 452.334 -200.306 m 452.334 -192.237 l S Q}%
\Vector{q 0 G 0 g 0.399 w 452.534 -192.436 m 459.707 -192.436 l S Q}%
\Vector{q 0 G 0 g 0.399 w 452.534 -200.107 m 459.707 -200.107 l S Q}%
\Vector{q 0 G 0 g 0.399 w 459.906 -200.306 m 459.906 -192.237 l S Q}%
\TextAt{8}{97.159}{544.692}{Another}%
\TextAt{8}{142.412}{544.692}{description}%
\TextAt{8}{201.712}{544.692}{of}%
\TextAt{8}{214.447}{544.692}{the}%
\TextAt{20}{234.672}{547.781}{\Character{98}}%
\TextAt{18}{234.008}{544.692}{Z}%
\TextAt{8}{241.978}{544.692}{-extension}%
\TextAt{20}{298.375}{547.781}{\Character{101}}%
\TextAt{18}{297.046}{544.692}{Q}%
\TextAt{16}{309.666}{544.692}{\Character{106}}%
\TextAt{18}{316.307}{544.692}{Q}%
\TextAt{8}{328.908}{544.692}{is}%
\TextAt{8}{340.077}{544.692}{obtained}%
\TextAt{8}{387.613}{544.692}{in}%
\TextAt{8}{400.671}{544.692}{the}%
\TextAt{8}{420.232}{544.692}{following}%
\TextAt{8}{469.071}{544.692}{manner.}%
\TextAt{8}{515.698}{544.692}{For}%
\TextAt{8}{79.226}{527.646}{every}%
\TextAt{8}{110.345}{527.646}{prime}%
\TextAt{8}{143.734}{527.646}{number}%
\TextAt{13}{186.557}{527.646}{p}%
\TextAt{8}{192.432}{527.646}{,}%
\TextAt{8}{199.864}{527.646}{let}%
\TextAt{9}{216.995}{527.646}{\Character{10}}%
\TextAt{21}{225.449}{525.853}{p}%
\TextAt{16}{233.917}{527.646}{\Character{106}}%
\TextAt{18}{240.944}{527.646}{Q}%
\TextAt{8}{254.366}{527.646}{be}%
\TextAt{8}{270.517}{527.646}{the}%
\TextAt{8}{290.899}{527.646}{field}%
\TextAt{8}{316.483}{527.646}{obtained}%
\TextAt{8}{364.84}{527.646}{by}%
\TextAt{8}{381.326}{527.646}{adjoining}%
\TextAt{8}{432.613}{527.646}{all}%
\TextAt{8}{449.098}{527.646}{roots}%
\TextAt{8}{478.985}{527.646}{of}%
\TextAt{8}{492.541}{527.646}{unity}%
\TextAt{8}{523.337}{527.646}{of}%
\TextAt{13}{79.226}{510.6}{p}%
\TextAt{8}{85.101}{510.6}{-power}%
\TextAt{8}{123.143}{510.6}{order.}%
\TextAt{8}{158.263}{510.6}{Then}%
\TextAt{13}{172.167}{481.599}{G}%
\TextAt{9}{181.4}{481.599}{\Character{40}\Character{10}}%
\TextAt{21}{194.407}{479.806}{p}%
\TextAt{16}{202.489}{481.599}{\Character{106}}%
\TextAt{18}{209.131}{481.599}{Q}%
\TextAt{9}{218.429}{481.599}{\Character{41}}%
\TextAt{9}{226.284}{481.599}{\Character{61}}%
\NativeMathAt{238.728}{481.599}{11.955}{\varprojlim}%
\TextAt{21}{244.515}{469.485}{\Character{23}}%
\TextAt{13}{256.979}{481.599}{G}%
\TextAt{9}{266.213}{481.599}{\Character{40}}%
\TextAt{18}{270.765}{481.599}{Q}%
\TextAt{9}{280.064}{481.599}{\Character{40}}%
\TextAt{13}{284.616}{481.599}{\Character{22}}%
\TextAt{21}{291.659}{479.806}{p}%
\TextAt{22}{295.922}{482.075}{\Character{23}}%
\TextAt{9}{301.041}{481.599}{\Character{41}}%
\TextAt{16}{308.915}{481.599}{\Character{106}}%
\TextAt{18}{315.556}{481.599}{Q}%
\TextAt{9}{324.855}{481.599}{\Character{41}}%
\TextAt{9}{332.71}{481.599}{\Character{61}}%
\NativeMathAt{345.153}{481.599}{11.955}{\varprojlim}%
\TextAt{21}{350.941}{469.485}{\Character{23}}%
\TextAt{9}{361.412}{481.599}{\Character{40}}%
\TextAt{18}{365.965}{481.599}{Z}%
\TextAt{9}{373.935}{481.599}{\Character{47}}%
\TextAt{13}{379.788}{481.599}{p}%
\TextAt{21}{385.663}{486.536}{\Character{23}}%
\TextAt{18}{390.845}{481.599}{Z}%
\TextAt{9}{398.815}{481.599}{\Character{41}}%
\TextAt{3}{403.368}{486.536}{\Character{3}}%
\TextAt{9}{411.421}{481.599}{\Character{61}}%
\TextAt{18}{423.846}{481.599}{Z}%
\TextAt{3}{431.816}{486.536}{\Character{3}}%
\TextAt{21}{431.816}{478.644}{p}%
\TextAt{13}{436.577}{481.599}{\Character{59}}%
\TextAt{8}{79.226}{446.372}{and}%
\TextAt{18}{103.023}{446.372}{Z}%
\TextAt{3}{110.993}{450.71}{\Character{3}}%
\TextAt{21}{110.993}{443.416}{p}%
\NativeMathAt{120.838}{446.371}{11.955}{\cong}%
\TextAt{18}{135.22}{446.372}{Z}%
\TextAt{21}{143.19}{444.578}{p}%
\TextAt{16}{151.312}{446.372}{\Character{2}}%
\TextAt{18}{163.972}{446.372}{Z}%
\TextAt{9}{171.943}{446.372}{\Character{47}\Character{40}}%
\TextAt{13}{182.348}{446.372}{p}%
\TextAt{16}{191.585}{446.372}{\Character{0}}%
\TextAt{9}{204.245}{446.372}{1\Character{41}}%
\TextAt{18}{214.65}{446.372}{Z}%
\TextAt{8}{227.552}{446.372}{for}%
\TextAt{13}{246.471}{446.372}{p}%
\NativeMathAt{257.43}{446.372}{11.955}{\neq}%
\TextAt{9}{271.609}{446.372}{2}%
\TextAt{8}{282.403}{446.372}{and}%
\TextAt{18}{306.2}{446.372}{Z}%
\TextAt{3}{314.17}{450.71}{\Character{3}}%
\TextAt{17}{314.17}{443.416}{2}%
\NativeMathAt{323.986}{446.371}{11.955}{\cong}%
\TextAt{18}{338.367}{446.372}{Z}%
\TextAt{17}{346.338}{444.578}{2}%
\TextAt{16}{354.432}{446.372}{\Character{2}}%
\TextAt{18}{367.092}{446.372}{Z}%
\TextAt{9}{375.062}{446.372}{\Character{47}2}%
\TextAt{18}{386.768}{446.372}{Z}%
\TextAt{8}{394.738}{446.372}{.}%
\TextAt{8}{406.289}{446.372}{The}%
\TextAt{8}{431.377}{446.372}{torsion}%
\TextAt{8}{471.505}{446.372}{subgroup}%
\TextAt{8}{523.337}{446.372}{of}%
\TextAt{18}{79.226}{429.326}{Z}%
\TextAt{3}{87.196}{433.664}{\Character{3}}%
\TextAt{21}{87.196}{426.37}{p}%
\TextAt{8}{95.936}{429.326}{is}%
\TextAt{8}{107.781}{429.326}{isomorphic}%
\TextAt{8}{167.112}{429.326}{to}%
\TextAt{18}{181.503}{429.326}{Z}%
\TextAt{9}{189.473}{429.326}{\Character{47}\Character{40}}%
\TextAt{13}{199.879}{429.326}{p}%
\TextAt{16}{208.466}{429.326}{\Character{0}}%
\TextAt{9}{220.476}{429.326}{1\Character{41}}%
\TextAt{18}{230.882}{429.326}{Z}%
\TextAt{8}{238.852}{429.326}{,}%
\TextAt{8}{246.102}{429.326}{resp.}%
\TextAt{18}{275.668}{429.326}{Z}%
\TextAt{9}{283.638}{429.326}{\Character{47}2}%
\TextAt{18}{295.344}{429.326}{Z}%
\TextAt{8}{303.314}{429.326}{,}%
\TextAt{8}{310.565}{429.326}{and}%
\TextAt{8}{333.408}{429.326}{taking}%
\TextAt{8}{369.594}{429.326}{its}%
\TextAt{8}{385.994}{429.326}{fixed}%
\TextAt{8}{414.361}{429.326}{field}%
\TextAt{8}{439.799}{429.326}{gives}%
\TextAt{8}{468.561}{429.326}{an}%
\TextAt{8}{484.901}{429.326}{extension}%
\TextAt{20}{80.555}{415.368}{\Character{101}}%
\TextAt{18}{79.226}{412.28}{Q}%
\TextAt{17}{88.525}{416.618}{\Character{40}}%
\TextAt{21}{91.818}{416.618}{p}%
\TextAt{17}{96.081}{416.618}{\Character{41}}%
\TextAt{16}{103.193}{412.28}{\Character{106}}%
\TextAt{18}{109.835}{412.28}{Q}%
\TextAt{8}{123.03}{412.28}{with}%
\TextAt{8}{149.691}{412.28}{Galois}%
\TextAt{8}{185.604}{412.28}{group}%
\TextAt{13}{260.914}{395.234}{G}%
\TextAt{9}{270.147}{395.234}{\Character{40}}%
\TextAt{20}{276.028}{398.322}{\Character{101}}%
\TextAt{18}{274.7}{395.234}{Q}%
\TextAt{17}{283.998}{400.17}{\Character{40}}%
\TextAt{21}{287.291}{400.17}{p}%
\TextAt{17}{291.554}{400.17}{\Character{41}}%
\TextAt{16}{298.666}{395.234}{\Character{106}}%
\TextAt{18}{305.308}{395.234}{Q}%
\TextAt{9}{314.607}{395.234}{\Character{41}}%
\NativeMathAt{322.48}{395.234}{11.955}{\cong}%
\TextAt{18}{335.099}{395.234}{Z}%
\TextAt{21}{343.069}{393.44}{p}%
\TextAt{13}{347.83}{395.234}{\Character{58}}%
\TextAt{8}{79.226}{371.712}{The}%
\TextAt{20}{103.944}{374.8}{\Character{98}}%
\TextAt{18}{103.28}{371.712}{Z}%
\TextAt{8}{111.25}{371.712}{-extension}%
\TextAt{20}{168.242}{374.8}{\Character{101}}%
\TextAt{18}{166.913}{371.712}{Q}%
\TextAt{16}{179.533}{371.712}{\Character{106}}%
\TextAt{18}{186.174}{371.712}{Q}%
\TextAt{8}{199.37}{371.712}{is}%
\TextAt{8}{211.134}{371.712}{then}%
\TextAt{8}{237.794}{371.712}{the}%
\TextAt{8}{257.951}{371.712}{composite}%
\TextAt{20}{314.297}{374.8}{\Character{101}}%
\TextAt{18}{312.968}{371.712}{Q}%
\TextAt{9}{325.588}{371.712}{\Character{61}}%
\TextAt{20}{338.013}{380.678}{\Character{81}}%
\TextAt{21}{349.304}{368.225}{p}%
\TextAt{20}{357.386}{374.8}{\Character{101}}%
\TextAt{18}{356.058}{371.712}{Q}%
\TextAt{17}{365.356}{376.05}{\Character{40}}%
\TextAt{21}{368.649}{376.05}{p}%
\TextAt{17}{372.912}{376.05}{\Character{41}}%
\TextAt{8}{376.703}{371.712}{.}%
\TextAt{8}{97.159}{351.677}{We}%
\TextAt{8}{117.503}{351.677}{fix}%
\TextAt{8}{134.285}{351.677}{an}%
\TextAt{8}{150.744}{351.677}{isomorphism}%
\TextAt{13}{219.363}{351.677}{G}%
\TextAt{9}{228.597}{351.677}{\Character{40}}%
\TextAt{20}{234.478}{354.766}{\Character{101}}%
\TextAt{18}{233.149}{351.677}{Q}%
\TextAt{16}{246.109}{351.677}{\Character{106}}%
\TextAt{18}{253.092}{351.677}{Q}%
\TextAt{9}{262.39}{351.677}{\Character{41}}%
\NativeMathAt{270.604}{351.677}{11.955}{\cong}%
\TextAt{20}{284.228}{354.766}{\Character{98}}%
\TextAt{18}{283.564}{351.677}{Z}%
\TextAt{8}{291.534}{351.677}{.}%
\TextAt{8}{300.582}{351.677}{There}%
\TextAt{8}{334.591}{351.677}{is}%
\TextAt{8}{346.555}{351.677}{no}%
\TextAt{8}{363.014}{351.677}{canonical}%
\TextAt{8}{414.597}{351.677}{choice}%
\TextAt{8}{449.587}{351.677}{as}%
\TextAt{8}{464.157}{351.677}{in}%
\TextAt{8}{478.01}{351.677}{the}%
\TextAt{8}{498.366}{351.677}{case}%
\TextAt{8}{523.337}{351.677}{of}%
\TextAt{8}{79.226}{334.631}{local}%
\TextAt{8}{105.674}{334.631}{fields.}%
\TextAt{8}{139.799}{334.631}{However,}%
\TextAt{8}{189.247}{334.631}{the}%
\TextAt{8}{208.211}{334.631}{reciprocity}%
\TextAt{8}{265.228}{334.631}{law}%
\TextAt{8}{285.172}{334.631}{will}%
\TextAt{8}{306.084}{334.631}{not}%
\TextAt{8}{325.705}{334.631}{depend}%
\TextAt{8}{365.148}{334.631}{on}%
\TextAt{8}{380.214}{334.631}{the}%
\TextAt{8}{399.177}{334.631}{choice.}%
\TextAt{8}{438.12}{334.631}{Via}%
\TextAt{13}{458.71}{334.631}{G}%
\TextAt{9}{467.943}{334.631}{\Character{40}}%
\TextAt{20}{473.824}{337.72}{\Character{101}}%
\TextAt{18}{472.496}{334.631}{Q}%
\TextAt{16}{485.115}{334.631}{\Character{106}}%
\TextAt{18}{491.757}{334.631}{Q}%
\TextAt{9}{501.055}{334.631}{\Character{41}}%
\NativeMathAt{508.928}{334.631}{11.955}{\cong}%
\TextAt{20}{522.212}{337.72}{\Character{98}}%
\TextAt{18}{521.548}{334.631}{Z}%
\TextAt{8}{529.518}{334.631}{,}%
\TextAt{8}{79.226}{317.585}{we}%
\TextAt{8}{96.454}{317.585}{obtain}%
\TextAt{8}{132.881}{317.585}{a}%
\TextAt{8}{142.636}{317.585}{continuous}%
\TextAt{8}{201.241}{317.585}{surjective}%
\TextAt{8}{253.652}{317.585}{homomorphism}%
\TextAt{13}{273.108}{288.584}{d}%
\TextAt{9}{282.511}{288.584}{\Character{58}}%
\TextAt{13}{289.084}{288.584}{G}%
\TextAt{25}{298.318}{286.791}{Q}%
\NativeArrowAt{308.336}{288.584}{11.955}{19.26}{\longrightarrow}%
\TextAt{20}{331.582}{291.673}{\Character{98}}%
\TextAt{18}{330.918}{288.584}{Z}%
\TextAt{8}{79.226}{259.583}{of}%
\TextAt{8}{93.722}{259.583}{the}%
\TextAt{8}{115.044}{259.583}{absolute}%
\TextAt{8}{162.452}{259.583}{Galois}%
\TextAt{8}{199.53}{259.583}{group}%
\TextAt{13}{233.872}{259.583}{G}%
\TextAt{25}{243.105}{257.79}{Q}%
\TextAt{9}{255.109}{259.583}{\Character{61}}%
\TextAt{13}{269.521}{259.583}{G}%
\TextAt{9}{278.755}{259.583}{\Character{40}}%
\Vector{q 0 G 0 g 0.478 w 211.307 -504.508 m 220.605 -504.508 l S Q}%
\TextAt{18}{283.307}{259.583}{Q}%
\TextAt{16}{297.912}{259.583}{\Character{106}}%
\TextAt{18}{306.54}{259.583}{Q}%
\TextAt{9}{315.839}{259.583}{\Character{41}}%
\TextAt{8}{320.391}{259.583}{.}%
\TextAt{8}{332.336}{259.583}{With}%
\TextAt{8}{363.736}{259.583}{this}%
\TextAt{8}{387.724}{259.583}{we}%
\TextAt{8}{406.117}{259.583}{continue}%
\TextAt{8}{454.111}{259.583}{as}%
\TextAt{8}{469.646}{259.583}{in}%
\TextAt{8}{484.465}{259.583}{chap.}%
\TextAt{8}{516.522}{259.583}{IV,}%
\TextAt{8}{79.226}{242.537}{§4,}%
\TextAt{8}{100.336}{242.537}{choosing}%
\TextAt{13}{149.82}{242.537}{k}%
\TextAt{9}{162.935}{242.537}{\Character{61}}%
\TextAt{18}{178.664}{242.537}{Q}%
\TextAt{8}{193.799}{242.537}{as}%
\TextAt{8}{210.108}{242.537}{our}%
\TextAt{8}{232.861}{242.537}{base}%
\TextAt{8}{260.874}{242.537}{field.}%
\TextAt{8}{296.599}{242.537}{If}%
\TextAt{13}{310.23}{242.537}{K}%
\TextAt{16}{327.672}{242.537}{\Character{106}}%
\TextAt{18}{337.618}{242.537}{Q}%
\TextAt{8}{352.753}{242.537}{is}%
\TextAt{8}{366.456}{242.537}{a}%
\TextAt{8}{378.15}{242.537}{finite}%
\TextAt{8}{410.001}{242.537}{extension,}%
\TextAt{8}{467.442}{242.537}{then}%
\TextAt{8}{496.041}{242.537}{we}%
\TextAt{8}{515.207}{242.537}{put}%
\TextAt{13}{79.226}{225.491}{f}%
\TextAt{21}{84.996}{223.698}{K}%
\TextAt{9}{96.52}{225.491}{\Character{61}}%
\TextAt{9}{108.93}{225.491}{\Character{91}}%
\TextAt{13}{112.197}{225.491}{K}%
\TextAt{16}{125.671}{225.491}{\Character{92}}%
\TextAt{20}{137.626}{228.579}{\Character{101}}%
\TextAt{18}{136.298}{225.491}{Q}%
\TextAt{9}{148.917}{225.491}{\Character{58}}%
\TextAt{18}{155.49}{225.491}{Q}%
\TextAt{9}{164.788}{225.491}{\Character{93}}%
\TextAt{8}{171.937}{225.491}{and}%
\TextAt{8}{194.7}{225.491}{get}%
\TextAt{8}{214.211}{225.491}{a}%
\TextAt{8}{223.966}{225.491}{surjective}%
\TextAt{8}{276.378}{225.491}{homomorphism}%
\TextAt{13}{247.532}{193.759}{d}%
\TextAt{21}{253.614}{191.965}{K}%
\TextAt{9}{265.138}{193.759}{\Character{61}}%
\TextAt{9}{282.819}{201.846}{1}%
\Vector{q 0 G 0 g 0.478 w 206.759 -577.253 m 220.731 -577.253 l S Q}%
\TextAt{13}{278.759}{185.558}{f}%
\TextAt{21}{284.529}{183.765}{K}%
\TextAt{13}{293.927}{193.759}{d}%
\TextAt{9}{303.331}{193.759}{\Character{58}}%
\TextAt{13}{309.903}{193.759}{G}%
\TextAt{21}{319.137}{191.965}{K}%
\NativeArrowAt{330.66}{193.759}{11.955}{19.26}{\longrightarrow}%
\TextAt{20}{353.906}{196.847}{\Character{98}}%
\TextAt{18}{353.242}{193.759}{Z}%
\TextAt{13}{361.212}{193.759}{\Character{59}}%
\TextAt{8}{79.226}{158.559}{which}%
\TextAt{8}{114.369}{158.559}{defines}%
\TextAt{8}{154.449}{158.559}{the}%
\TextAt{20}{176.925}{161.648}{\Character{98}}%
\TextAt{18}{176.261}{158.559}{Z}%
\TextAt{8}{184.231}{158.559}{-extension}%
\TextAt{20}{244.289}{161.581}{\Character{101}}%
\TextAt{13}{241.551}{158.559}{K}%
\TextAt{9}{258.512}{158.559}{\Character{61}}%
\TextAt{13}{273.759}{158.559}{K}%
\TextAt{20}{285.905}{161.648}{\Character{101}}%
\TextAt{18}{284.577}{158.559}{Q}%
\TextAt{8}{299.429}{158.559}{of}%
\TextAt{13}{314.415}{158.559}{K}%
\TextAt{8}{325.233}{158.559}{.}%
\TextAt{20}{341.388}{161.581}{\Character{101}}%
\TextAt{13}{338.65}{158.559}{K}%
\TextAt{16}{355.61}{158.559}{\Character{106}}%
\TextAt{13}{365.074}{158.559}{K}%
\TextAt{8}{381.446}{158.559}{is}%
\TextAt{8}{394.866}{158.559}{called}%
\TextAt{8}{429.685}{158.559}{the}%
\TextAt{12}{451.498}{158.559}{cyclotomic}%
\TextAt{20}{520.98}{161.648}{\Character{98}}%
\TextAt{18}{520.316}{158.559}{Z}%
\TextAt{12}{528.286}{158.559}{-}%
\TextAt{12}{79.226}{141.513}{extension}%
\TextAt{8}{139.49}{141.513}{of}%
\TextAt{13}{153.858}{141.513}{K}%
\TextAt{8}{164.676}{141.513}{.}%
\TextAt{8}{176.238}{141.513}{We}%
\TextAt{8}{197.42}{141.513}{denote}%
\TextAt{8}{236.177}{141.513}{again}%
\TextAt{8}{268.441}{141.513}{by}%
\TextAt{13}{285.738}{141.513}{\Character{39}}%
\TextAt{21}{293.407}{139.72}{K}%
\TextAt{8}{306.545}{141.513}{the}%
\TextAt{8}{327.739}{141.513}{element}%
\TextAt{8}{372.019}{141.513}{of}%
\TextAt{13}{386.387}{141.513}{G}%
\TextAt{9}{395.62}{141.513}{\Character{40}}%
\TextAt{20}{402.911}{144.535}{\Character{101}}%
\TextAt{13}{400.173}{141.513}{K}%
\TextAt{16}{416.08}{141.513}{\Character{106}}%
\TextAt{13}{424.49}{141.513}{K}%
\TextAt{9}{435.308}{141.513}{\Character{41}}%
\TextAt{8}{444.796}{141.513}{which}%
\TextAt{8}{479.32}{141.513}{is}%
\TextAt{8}{492.122}{141.513}{mapped}%
\TextAt{8}{79.226}{124.467}{to}%
\TextAt{9}{93.939}{124.467}{1}%
\TextAt{8}{104.093}{124.467}{by}%
\TextAt{8}{120.754}{124.467}{the}%
\TextAt{8}{141.314}{124.467}{isomorphism}%
\TextAt{13}{210.136}{124.467}{G}%
\TextAt{9}{219.369}{124.467}{\Character{40}}%
\TextAt{20}{226.66}{127.489}{\Character{101}}%
\TextAt{13}{223.922}{124.467}{K}%
\TextAt{16}{238.747}{124.467}{\Character{106}}%
\TextAt{13}{246.075}{124.467}{K}%
\TextAt{9}{256.893}{124.467}{\Character{41}}%
\NativeMathAt{265.452}{124.467}{11.955}{\cong}%
\TextAt{20}{279.422}{127.556}{\Character{98}}%
\TextAt{18}{278.758}{124.467}{Z}%
\TextAt{8}{286.728}{124.467}{,}%
\TextAt{8}{294.381}{124.467}{and}%
\TextAt{8}{317.546}{124.467}{by}%
\TextAt{13}{334.208}{124.467}{\Character{39}}%
\TextAt{21}{341.877}{122.608}{L}%
\TextAt{3}{347.636}{122.608}{\Character{106}}%
\TextAt{21}{349.988}{122.608}{K}%
\TextAt{8}{362.491}{124.467}{the}%
\TextAt{8}{383.05}{124.467}{restriction}%
\TextAt{13}{439.451}{124.467}{\Character{39}}%
\TextAt{21}{447.119}{122.674}{K}%
\TextAt{16}{455.322}{124.467}{\Character{106}}%
\TextAt{21}{458.643}{122.674}{L}%
\TextAt{8}{469.2}{124.467}{if}%
\TextAt{13}{480.327}{124.467}{L}%
\TextAt{16}{492.299}{124.467}{\Character{106}}%
\TextAt{13}{499.627}{124.467}{K}%
\TextAt{8}{514.745}{124.467}{is}%
\TextAt{8}{526.911}{124.467}{a}%
\TextAt{8}{79.226}{107.421}{subextension}%
\TextAt{8}{148.113}{107.421}{of}%
\TextAt{20}{163.358}{110.443}{\Character{101}}%
\TextAt{13}{160.62}{107.421}{K}%
\TextAt{16}{174.758}{107.421}{\Character{106}}%
\TextAt{13}{181.4}{107.421}{K}%
\TextAt{8}{192.218}{107.421}{.}%
\TextAt{8}{200.391}{107.421}{The}%
\TextAt{8}{223.622}{107.421}{automorphism}%
\TextAt{13}{300.268}{107.421}{\Character{39}}%
\TextAt{21}{307.936}{105.562}{L}%
\TextAt{3}{313.695}{105.562}{\Character{106}}%
\TextAt{21}{316.047}{105.562}{K}%
\TextAt{8}{327.324}{107.421}{must}%
\TextAt{8}{355.503}{107.421}{not}%
\TextAt{8}{375.494}{107.421}{be}%
\TextAt{8}{390.595}{107.421}{confused}%
\TextAt{8}{437.628}{107.421}{with}%
\TextAt{8}{463.464}{107.421}{the}%
\TextAt{8}{482.797}{107.421}{Frobenius}%
\TextAt{8}{79.226}{90.375}{automorphism}%
\TextAt{8}{156.696}{90.375}{corresponding}%
\TextAt{8}{231.882}{90.375}{to}%
\TextAt{8}{246.192}{90.375}{a}%
\TextAt{8}{255.947}{90.375}{prime}%
\TextAt{8}{289.111}{90.375}{ideal}%
\TextAt{8}{317.074}{90.375}{of}%
\TextAt{13}{330.404}{90.375}{L}%
\TextAt{8}{342.266}{90.375}{(see}%
\TextAt{8}{365.734}{90.375}{§7).}%
\TextAt{8}{297.211}{43.033}{352}%
\Destination{78.23}{809.815}{nav.page.352}%
\Destination{79.36}{639.855}{nav.print.386}%
\Hotlink{516.022}{257.355}{13.995}{11.047}{nav.chapter.4}%
\Hotlink{84.405}{240.355}{6.858}{11.047}{nav.section.4.4}%
\Hotlink{370.913}{88.155}{6.858}{11.047}{nav.section.6.7}%
}

% Page 365 (body page 353)
\Page{%
\TextAt{2}{79.226}{796.943}{VI.}%
\TextAt{2}{98.012}{796.943}{Global}%
\TextAt{2}{133.248}{796.943}{Class}%
\TextAt{2}{161.841}{796.943}{Field}%
\TextAt{2}{189.572}{796.943}{Theory}%
\Vector{q 0 G 0 g 0.299 w 7.226 18.729 m 460.769 18.729 l S Q}%
\TextAt{8}{97.159}{761.124}{For}%
\TextAt{8}{118.639}{761.124}{the}%
\TextAt{13}{139.307}{761.124}{G}%
\TextAt{25}{148.541}{759.33}{Q}%
\TextAt{8}{155.238}{761.124}{-module}%
\TextAt{13}{200.939}{761.124}{A}%
\TextAt{8}{209.715}{761.124}{,}%
\TextAt{8}{217.503}{761.124}{we}%
\TextAt{8}{235.241}{761.124}{choose}%
\TextAt{8}{272.885}{761.124}{the}%
\TextAt{8}{293.553}{761.124}{union}%
\TextAt{8}{326.582}{761.124}{of}%
\TextAt{8}{340.423}{761.124}{the}%
\TextAt{8}{361.091}{761.124}{idèle}%
\TextAt{8}{388.908}{761.124}{class}%
\TextAt{8}{416.856}{761.124}{groups}%
\TextAt{13}{455.157}{761.124}{C}%
\TextAt{21}{463.532}{759.33}{K}%
\TextAt{8}{476.144}{761.124}{of}%
\TextAt{8}{489.985}{761.124}{all}%
\TextAt{8}{506.755}{761.124}{finite}%
\TextAt{8}{79.226}{744.078}{extensions}%
\TextAt{13}{134.942}{744.078}{K}%
\TextAt{16}{149.08}{744.078}{\Character{106}}%
\TextAt{18}{155.722}{744.078}{Q}%
\TextAt{8}{165.021}{744.078}{.}%
\TextAt{8}{173.247}{744.078}{Thus}%
\TextAt{13}{202.231}{744.078}{A}%
\TextAt{21}{211.006}{742.284}{K}%
\TextAt{9}{222.53}{744.078}{\Character{61}}%
\TextAt{13}{234.955}{744.078}{C}%
\TextAt{21}{243.33}{742.284}{K}%
\TextAt{8}{251.533}{744.078}{.}%
\TextAt{8}{259.76}{744.078}{The}%
\TextAt{8}{283.149}{744.078}{henselian}%
\TextAt{8}{333.269}{744.078}{valuation}%
\TextAt{13}{383.677}{744.078}{v}%
\TextAt{9}{393.086}{744.078}{\Character{58}}%
\TextAt{13}{399.658}{744.078}{C}%
\TextAt{25}{408.033}{742.284}{Q}%
\NativeArrowAt{418.051}{744.078}{11.955}{19.26}{\longrightarrow}%
\TextAt{20}{441.297}{747.166}{\Character{98}}%
\TextAt{18}{440.633}{744.078}{Z}%
\TextAt{8}{451.836}{744.078}{will}%
\TextAt{8}{473.276}{744.078}{be}%
\TextAt{8}{488.535}{744.078}{obtained}%
\TextAt{8}{79.226}{727.032}{as}%
\TextAt{8}{93.596}{727.032}{the}%
\TextAt{8}{113.753}{727.032}{composite}%
\TextAt{13}{240.637}{706.999}{C}%
\TextAt{25}{249.012}{705.205}{Q}%
\TextAt{17}{261.618}{714.706}{\Character{91}}%
\TextAt{21}{267.867}{714.706}{\Character{59}}%
\TextAt{30}{270.967}{716.764}{\Character{101}}%
\TextAt{25}{270.22}{714.706}{Q}%
\TextAt{3}{276.419}{714.706}{\Character{106}}%
\TextAt{25}{278.771}{714.706}{Q}%
\TextAt{17}{284.97}{714.706}{\Character{93}}%
\NativeArrowAt{259.03}{706.999}{11.955}{32.255}{\longrightarrow}%
\TextAt{13}{294.642}{706.999}{G}%
\TextAt{9}{303.876}{706.999}{\Character{40}}%
\TextAt{20}{309.756}{710.087}{\Character{101}}%
\TextAt{18}{308.428}{706.999}{Q}%
\TextAt{16}{321.047}{706.999}{\Character{106}}%
\TextAt{18}{327.689}{706.999}{Q}%
\TextAt{9}{336.988}{706.999}{\Character{41}}%
\TextAt{21}{347.954}{713.776}{d}%
\NativeArrowAt{344.861}{706.999}{11.955}{11.955}{\longrightarrow}%
\TextAt{20}{360.801}{710.087}{\Character{98}}%
\TextAt{18}{360.137}{706.999}{Z}%
\TextAt{13}{368.107}{706.999}{\Character{59}}%
\TextAt{8}{79.226}{683.699}{where}%
\TextAt{8}{113.007}{683.699}{the}%
\TextAt{8}{133.134}{683.699}{mapping}%
\TextAt{9}{181.237}{683.699}{\Character{91}}%
\TextAt{13}{188.357}{683.699}{\Character{59}}%
\TextAt{20}{194.929}{686.787}{\Character{101}}%
\TextAt{18}{193.601}{683.699}{Q}%
\TextAt{16}{206.22}{683.699}{\Character{106}}%
\TextAt{18}{212.862}{683.699}{Q}%
\TextAt{9}{222.161}{683.699}{\Character{93}}%
\TextAt{8}{229.281}{683.699}{will}%
\TextAt{8}{251.357}{683.699}{later}%
\TextAt{8}{278.646}{683.699}{turn}%
\TextAt{8}{304.631}{683.699}{out}%
\TextAt{8}{325.417}{683.699}{to}%
\TextAt{8}{339.698}{683.699}{be}%
\TextAt{8}{355.593}{683.699}{the}%
\TextAt{8}{375.721}{683.699}{norm}%
\TextAt{8}{406.262}{683.699}{residue}%
\TextAt{8}{445.96}{683.699}{symbol}%
\TextAt{9}{485.993}{683.699}{\Character{40}}%
\TextAt{13}{494.413}{683.699}{\Character{59}}%
\TextAt{20}{500.986}{686.787}{\Character{101}}%
\TextAt{18}{499.658}{683.699}{Q}%
\TextAt{16}{512.277}{683.699}{\Character{106}}%
\TextAt{18}{518.919}{683.699}{Q}%
\TextAt{9}{528.217}{683.699}{\Character{41}}%
\TextAt{8}{79.226}{666.653}{of}%
\TextAt{8}{92.556}{666.653}{global}%
\TextAt{8}{127.035}{666.653}{class}%
\TextAt{8}{154.472}{666.653}{field}%
\TextAt{8}{179.829}{666.653}{theory}%
\TextAt{8}{216.579}{666.653}{(see}%
\TextAt{8}{240.047}{666.653}{(5.7)).}%
\TextAt{8}{277.128}{666.653}{For}%
\TextAt{8}{298.097}{666.653}{the}%
\TextAt{8}{318.254}{666.653}{moment}%
\TextAt{8}{363.456}{666.653}{we}%
\TextAt{8}{380.684}{666.653}{merely}%
\TextAt{8}{418.725}{666.653}{define}%
\TextAt{8}{452.534}{666.653}{it}%
\TextAt{8}{464.238}{666.653}{as}%
\TextAt{8}{478.609}{666.653}{follows.}%
\TextAt{8}{97.159}{646.618}{For}%
\TextAt{8}{118.128}{646.618}{an}%
\TextAt{8}{134.387}{646.618}{arbitrary}%
\TextAt{8}{184.157}{646.618}{finite}%
\TextAt{8}{214.069}{646.618}{abelian}%
\TextAt{8}{254.716}{646.618}{extension}%
\TextAt{13}{306.482}{646.618}{L}%
\TextAt{16}{317.767}{646.618}{\Character{106}}%
\TextAt{13}{324.409}{646.618}{K}%
\TextAt{8}{335.227}{646.618}{,}%
\TextAt{8}{342.376}{646.618}{we}%
\TextAt{8}{359.603}{646.618}{define}%
\TextAt{8}{393.413}{646.618}{the}%
\TextAt{8}{413.569}{646.618}{homomorphism}%
\TextAt{9}{235.681}{618.135}{\Character{91}}%
\TextAt{13}{242.83}{618.135}{\Character{59}}%
\TextAt{13}{248.054}{618.135}{L}%
\TextAt{16}{259.359}{618.135}{\Character{106}}%
\TextAt{13}{266.001}{618.135}{K}%
\TextAt{9}{276.819}{618.135}{\Character{93}}%
\TextAt{9}{283.37}{618.135}{\Character{58}}%
\TextAt{13}{289.964}{618.135}{I}%
\TextAt{21}{295.126}{616.342}{K}%
\NativeArrowAt{306.65}{618.135}{11.955}{19.26}{\longrightarrow}%
\TextAt{13}{329.232}{618.135}{G}%
\TextAt{9}{338.466}{618.135}{\Character{40}}%
\TextAt{13}{343.018}{618.135}{L}%
\TextAt{16}{354.303}{618.135}{\Character{106}}%
\TextAt{13}{360.945}{618.135}{K}%
\TextAt{9}{371.763}{618.135}{\Character{41}}%
\TextAt{8}{79.226}{589.652}{by}%
\TextAt{9}{233.283}{572.606}{\Character{91}}%
\TextAt{13}{236.535}{572.606}{\Character{11}\Character{59}}%
\TextAt{13}{249.279}{572.606}{L}%
\TextAt{16}{260.586}{572.606}{\Character{106}}%
\TextAt{13}{267.228}{572.606}{K}%
\TextAt{9}{278.046}{572.606}{\Character{93}}%
\TextAt{9}{284.597}{572.606}{\Character{61}}%
\TextAt{20}{297.044}{583.963}{\Character{89}}%
\TextAt{27}{302.687}{558.857}{p}%
\TextAt{9}{312.32}{572.606}{\Character{40}}%
\TextAt{13}{316.872}{572.606}{\Character{11}}%
\TextAt{27}{324.318}{570.813}{p}%
\TextAt{13}{328.804}{572.606}{\Character{59}}%
\TextAt{13}{334.028}{572.606}{L}%
\TextAt{27}{342.013}{570.813}{p}%
\TextAt{16}{349.821}{572.606}{\Character{106}}%
\TextAt{13}{356.462}{572.606}{K}%
\TextAt{27}{366.422}{570.813}{p}%
\TextAt{9}{370.908}{572.606}{\Character{41}}%
\TextAt{13}{375.461}{572.606}{\Character{59}}%
\TextAt{8}{79.226}{539.939}{where}%
\TextAt{13}{113.24}{539.939}{L}%
\TextAt{27}{121.205}{538.146}{p}%
\TextAt{8}{129.794}{539.939}{denotes}%
\TextAt{8}{172.333}{539.939}{the}%
\TextAt{8}{192.694}{539.939}{completion}%
\TextAt{8}{252.735}{539.939}{of}%
\TextAt{13}{266.27}{539.939}{L}%
\TextAt{8}{278.337}{539.939}{with}%
\TextAt{8}{305.202}{539.939}{respect}%
\TextAt{8}{345.457}{539.939}{to}%
\TextAt{8}{359.972}{539.939}{a}%
\TextAt{8}{369.933}{539.939}{place}%
\TextAt{14}{400.05}{539.939}{P}%
\TextAt{16}{413.613}{539.939}{\Character{106}}%
\TextAt{14}{420.604}{539.939}{p}%
\TextAt{8}{426.587}{539.939}{,}%
\TextAt{8}{433.993}{539.939}{and}%
\TextAt{9}{456.961}{539.939}{\Character{40}}%
\TextAt{13}{461.513}{539.939}{\Character{11}}%
\TextAt{27}{468.958}{538.146}{p}%
\TextAt{13}{473.445}{539.939}{\Character{59}}%
\TextAt{13}{478.669}{539.939}{L}%
\TextAt{27}{486.654}{538.146}{p}%
\TextAt{16}{494.811}{539.939}{\Character{106}}%
\TextAt{13}{501.802}{539.939}{K}%
\TextAt{27}{511.761}{538.146}{p}%
\TextAt{9}{516.248}{539.939}{\Character{41}}%
\TextAt{8}{524.903}{539.939}{is}%
\TextAt{8}{79.226}{522.893}{the}%
\TextAt{8}{100.16}{522.893}{norm}%
\TextAt{8}{131.507}{522.893}{residue}%
\TextAt{8}{172.011}{522.893}{symbol}%
\TextAt{8}{212.85}{522.893}{of}%
\TextAt{8}{226.957}{522.893}{local}%
\TextAt{8}{255.375}{522.893}{class}%
\TextAt{8}{283.589}{522.893}{field}%
\TextAt{8}{309.724}{522.893}{theory.}%
\TextAt{8}{352.376}{522.893}{Note}%
\TextAt{8}{381.44}{522.893}{that}%
\TextAt{8}{407.586}{522.893}{almost}%
\TextAt{8}{446.154}{522.893}{all}%
\TextAt{8}{463.19}{522.893}{factors}%
\TextAt{8}{502.08}{522.893}{in}%
\TextAt{8}{516.51}{522.893}{the}%
\TextAt{8}{79.226}{505.847}{product}%
\TextAt{8}{123.097}{505.847}{are}%
\TextAt{9}{142.579}{505.847}{1}%
\TextAt{8}{152.3}{505.847}{because}%
\TextAt{8}{195.573}{505.847}{almost}%
\TextAt{8}{233.334}{505.847}{all}%
\TextAt{8}{249.564}{505.847}{extensions}%
\TextAt{13}{305.916}{505.847}{L}%
\TextAt{27}{313.88}{504.054}{p}%
\TextAt{16}{321.688}{505.847}{\Character{106}}%
\TextAt{13}{328.33}{505.847}{K}%
\TextAt{27}{338.289}{504.054}{p}%
\TextAt{8}{346.644}{505.847}{are}%
\TextAt{8}{366.126}{505.847}{unramified}%
\TextAt{8}{424.63}{505.847}{and}%
\TextAt{8}{447.364}{505.847}{almost}%
\TextAt{8}{485.125}{505.847}{all}%
\TextAt{13}{501.355}{505.847}{\Character{11}}%
\TextAt{27}{508.801}{504.054}{p}%
\TextAt{8}{517.156}{505.847}{are}%
\TextAt{8}{79.226}{488.801}{units.}%
\TextAt{12}{79.226}{462.937}{Proposition}%
\TextAt{12}{152.711}{462.937}{(5.2).}%
\TextAt{8}{188.699}{462.937}{If}%
\TextAt{13}{201.063}{462.937}{L}%
\TextAt{16}{213.494}{462.937}{\Character{106}}%
\TextAt{13}{221.282}{462.937}{K}%
\TextAt{8}{236.669}{462.937}{and}%
\TextAt{13}{260.104}{462.937}{L}%
\TextAt{3}{268.069}{467.275}{\Character{48}}%
\TextAt{16}{275.33}{462.937}{\Character{106}}%
\TextAt{13}{283.118}{462.937}{K}%
\TextAt{3}{293.936}{467.275}{\Character{48}}%
\TextAt{8}{301.3}{462.937}{are}%
\TextAt{8}{321.483}{462.937}{two}%
\TextAt{8}{344.273}{462.937}{abelian}%
\TextAt{8}{385.593}{462.937}{extensions}%
\TextAt{8}{442.646}{462.937}{of}%
\TextAt{8}{456.648}{462.937}{finite}%
\TextAt{8}{487.232}{462.937}{algebraic}%
\TextAt{8}{79.226}{445.891}{number}%
\TextAt{8}{121.822}{445.891}{fields}%
\TextAt{8}{151.794}{445.891}{such}%
\TextAt{8}{178.191}{445.891}{that}%
\TextAt{13}{203.56}{445.891}{K}%
\NativeMathAt{217.698}{445.891}{11.955}{\subseteq}%
\TextAt{13}{230.318}{445.891}{K}%
\TextAt{3}{241.135}{450.229}{\Character{48}}%
\TextAt{8}{247.828}{445.891}{and}%
\TextAt{13}{270.591}{445.891}{L}%
\NativeMathAt{281.876}{445.891}{11.955}{\subseteq}%
\TextAt{13}{294.495}{445.891}{L}%
\TextAt{3}{302.46}{450.229}{\Character{48}}%
\TextAt{8}{305.255}{445.891}{,}%
\TextAt{8}{312.404}{445.891}{then}%
\TextAt{8}{339.064}{445.891}{we}%
\TextAt{8}{356.291}{445.891}{have}%
\TextAt{8}{383.286}{445.891}{the}%
\TextAt{8}{403.443}{445.891}{commutative}%
\TextAt{8}{473.368}{445.891}{diagram}%
\TextAt{13}{253.702}{414.499}{I}%
\TextAt{21}{258.865}{412.706}{K}%
\TextAt{23}{266.569}{414.975}{\Character{48}}%
\TextAt{13}{323.096}{414.499}{G}%
\TextAt{9}{332.329}{414.499}{\Character{40}}%
\TextAt{13}{336.881}{414.499}{L}%
\TextAt{3}{344.846}{418.838}{\Character{48}}%
\TextAt{16}{350.962}{414.499}{\Character{106}}%
\TextAt{13}{357.604}{414.499}{K}%
\TextAt{3}{368.421}{418.838}{\Character{48}}%
\TextAt{9}{371.216}{414.499}{\Character{41}}%
\TextAt{13}{255.044}{372.265}{I}%
\TextAt{21}{260.207}{370.472}{K}%
\TextAt{13}{324.265}{372.265}{G}%
\TextAt{9}{333.498}{372.265}{\Character{40}}%
\TextAt{13}{338.05}{372.265}{L}%
\TextAt{16}{349.336}{372.265}{\Character{106}}%
\TextAt{13}{355.978}{372.265}{K}%
\TextAt{9}{366.795}{372.265}{\Character{41}}%
\TextAt{13}{371.348}{372.265}{\Character{58}}%
\NativeDiagram{%
\path[draw=black,line width=0.478bp] (203.104,-356.512) -- (245.266,-356.512);%
\path[draw=black,line width=0.478bp,line cap=round,line join=round] (243.019,-353.643) .. controls (243.473,-355.364) and (244.485,-356.177) .. (245.505,-356.512) .. controls (244.485,-356.846) and (243.473,-357.659) .. (243.019,-359.381);%
}%
\TextAt{17}{280.354}{422.296}{\Character{91}}%
\TextAt{21}{286.604}{422.296}{\Character{59}L}%
\TextAt{23}{294.715}{425.109}{\Character{48}}%
\TextAt{3}{297.4}{422.296}{\Character{106}}%
\TextAt{21}{299.752}{422.296}{K}%
\TextAt{23}{307.457}{425.109}{\Character{48}}%
\TextAt{17}{310.141}{422.296}{\Character{93}}%
\NativeDiagram{%
\path[draw=black,line width=0.478bp] (189.727,-365.873) -- (189.727,-388.508);%
\path[draw=black,line width=0.478bp,line cap=round,line join=round] (192.596,-386.261) .. controls (190.874,-386.715) and (190.061,-387.728) .. (189.727,-388.748) .. controls (189.392,-387.728) and (188.579,-386.715) .. (186.858,-386.261);%
}%
\TextAt{21}{233.648}{395.74}{N}%
\TextAt{22}{240.37}{393.449}{K}%
\TextAt{23}{246.915}{395.165}{\Character{48}}%
\TextAt{23}{249.599}{393.449}{\Character{106}}%
\TextAt{22}{251.868}{393.449}{K}%
\NativeDiagram{%
\path[draw=black,line width=0.478bp] (277.433,-367.068) -- (277.433,-387.711);%
\path[draw=black,line width=0.478bp,line cap=round,line join=round] (280.303,-385.464) .. controls (278.581,-385.918) and (277.768,-386.931) .. (277.433,-387.95) .. controls (277.099,-386.931) and (276.286,-385.918) .. (274.564,-385.464);%
\path[draw=black,line width=0.478bp] (201.762,-398.746) -- (246.436,-398.746);%
\path[draw=black,line width=0.478bp,line cap=round,line join=round] (244.188,-395.877) .. controls (244.642,-397.599) and (245.655,-398.412) .. (246.675,-398.746) .. controls (245.655,-399.081) and (244.642,-399.894) .. (244.188,-401.615);%
}%
\TextAt{17}{282.952}{380.062}{\Character{91}}%
\TextAt{21}{289.202}{380.062}{\Character{59}L}%
\TextAt{3}{297.313}{380.062}{\Character{106}}%
\TextAt{21}{299.666}{380.062}{K}%
\TextAt{17}{307.37}{380.062}{\Character{93}}%
\TextAt{19}{79.226}{325.86}{Proof.}%
\TextAt{8}{116}{325.86}{For}%
\TextAt{8}{136.97}{325.86}{an}%
\TextAt{8}{153.229}{325.86}{idèle}%
\TextAt{13}{180.534}{325.86}{\Character{11}}%
\TextAt{9}{191.377}{325.86}{\Character{61}}%
\TextAt{9}{203.787}{325.86}{\Character{40}}%
\TextAt{13}{208.355}{325.86}{\Character{11}}%
\TextAt{27}{215.8}{324.067}{P}%
\TextAt{9}{222.894}{325.86}{\Character{41}}%
\NativeMathAt{230.767}{325.86}{11.955}{\in}%
\TextAt{13}{242.058}{325.86}{I}%
\TextAt{21}{247.221}{324.067}{K}%
\TextAt{23}{254.925}{326.336}{\Character{48}}%
\TextAt{8}{262.005}{325.86}{of}%
\TextAt{13}{275.335}{325.86}{K}%
\TextAt{3}{286.153}{330.199}{\Character{48}}%
\TextAt{8}{288.948}{325.86}{,}%
\TextAt{8}{296.097}{325.86}{we}%
\TextAt{8}{313.325}{325.86}{find}%
\TextAt{8}{336.733}{325.86}{by}%
\TextAt{8}{352.992}{325.86}{chap.}%
\TextAt{8}{383.884}{325.86}{IV,}%
\TextAt{8}{404.029}{325.86}{(6.4),}%
\TextAt{8}{435.255}{325.86}{that}%
\TextAt{9}{170.871}{297.377}{\Character{40}}%
\TextAt{13}{175.424}{297.377}{\Character{11}}%
\TextAt{27}{182.869}{295.584}{P}%
\TextAt{13}{189.963}{297.377}{\Character{59}}%
\TextAt{13}{195.187}{297.377}{L}%
\TextAt{3}{203.172}{302.313}{\Character{48}}%
\TextAt{27}{203.172}{294.422}{P}%
\TextAt{16}{213.586}{297.377}{\Character{106}}%
\TextAt{13}{220.228}{297.377}{K}%
\TextAt{3}{231.045}{302.313}{\Character{48}}%
\TextAt{27}{230.187}{294.422}{P}%
\TextAt{9}{237.28}{297.377}{\Character{41}}%
\TextAt{16}{241.833}{297.377}{\Character{106}}%
\TextAt{21}{245.154}{295.584}{L}%
\TextAt{28}{250.913}{294.588}{p}%
\TextAt{9}{258.814}{297.377}{\Character{61}}%
\TextAt{9}{271.224}{297.377}{\Character{40}}%
\TextAt{13}{275.792}{297.377}{N}%
\TextAt{21}{285.166}{295.429}{K}%
\TextAt{23}{292.871}{298.174}{\Character{48}}%
\TextAt{28}{292.307}{292.603}{P}%
\TextAt{3}{298.505}{295.429}{\Character{106}}%
\TextAt{21}{300.857}{295.429}{K}%
\TextAt{28}{307.998}{294.433}{p}%
\TextAt{9}{312.579}{297.377}{\Character{40}}%
\TextAt{13}{317.131}{297.377}{\Character{11}}%
\TextAt{27}{324.577}{295.584}{P}%
\TextAt{9}{331.671}{297.377}{\Character{41}}%
\TextAt{13}{336.223}{297.377}{\Character{59}}%
\TextAt{13}{341.447}{297.377}{L}%
\TextAt{27}{349.432}{295.584}{p}%
\TextAt{16}{357.239}{297.377}{\Character{106}}%
\TextAt{13}{363.881}{297.377}{K}%
\TextAt{27}{373.84}{295.584}{p}%
\TextAt{9}{378.327}{297.377}{\Character{41}}%
\TextAt{13}{382.879}{297.377}{\Character{59}}%
\TextAt{14}{412.034}{297.377}{P}%
\TextAt{16}{425.248}{297.377}{\Character{106}}%
\TextAt{14}{431.89}{297.377}{p}%
\TextAt{13}{437.873}{297.377}{\Character{59}}%
\TextAt{8}{79.226}{268.894}{and}%
\TextAt{8}{101.989}{268.894}{(2.2)}%
\TextAt{8}{129.964}{268.894}{implies}%
\TextAt{9}{165.409}{240.411}{\Character{91}}%
\TextAt{13}{168.66}{240.411}{N}%
\TextAt{21}{178.034}{238.551}{K}%
\TextAt{23}{185.739}{240.821}{\Character{48}}%
\TextAt{3}{188.423}{238.551}{\Character{106}}%
\TextAt{21}{190.775}{238.551}{K}%
\TextAt{9}{198.978}{240.411}{\Character{40}}%
\TextAt{13}{203.53}{240.411}{\Character{11}}%
\TextAt{9}{211.052}{240.411}{\Character{41}}%
\TextAt{13}{215.604}{240.411}{\Character{59}}%
\TextAt{13}{220.828}{240.411}{L}%
\TextAt{16}{232.134}{240.411}{\Character{106}}%
\TextAt{13}{238.776}{240.411}{K}%
\TextAt{9}{249.593}{240.411}{\Character{93}}%
\TextAt{9}{256.144}{240.411}{\Character{61}}%
\TextAt{20}{268.591}{251.768}{\Character{89}}%
\TextAt{27}{274.235}{226.662}{p}%
\TextAt{9}{283.867}{240.411}{\Character{40}}%
\TextAt{13}{288.42}{240.411}{N}%
\TextAt{21}{297.794}{238.551}{K}%
\TextAt{23}{305.498}{240.821}{\Character{48}}%
\TextAt{3}{308.182}{238.551}{\Character{106}}%
\TextAt{21}{310.535}{238.551}{K}%
\TextAt{9}{318.737}{240.411}{\Character{40}}%
\TextAt{13}{323.29}{240.411}{\Character{11}}%
\TextAt{9}{330.811}{240.411}{\Character{41}}%
\TextAt{27}{335.364}{238.618}{p}%
\TextAt{13}{339.851}{240.411}{\Character{59}}%
\TextAt{13}{345.075}{240.411}{L}%
\TextAt{27}{353.059}{238.618}{p}%
\TextAt{16}{360.867}{240.411}{\Character{106}}%
\TextAt{13}{367.509}{240.411}{K}%
\TextAt{27}{377.468}{238.618}{p}%
\TextAt{9}{381.955}{240.411}{\Character{41}}%
\TextAt{9}{256.166}{207.422}{\Character{61}}%
\TextAt{20}{268.591}{218.78}{\Character{89}}%
\TextAt{27}{274.235}{193.674}{p}%
\TextAt{20}{285.86}{218.78}{\Character{89}}%
\TextAt{27}{287.03}{192.877}{P}%
\TextAt{3}{293.625}{192.877}{\Character{106}}%
\TextAt{27}{295.977}{192.877}{p}%
\TextAt{9}{301.136}{207.422}{\Character{40}}%
\TextAt{13}{305.688}{207.422}{N}%
\TextAt{21}{315.062}{205.474}{K}%
\TextAt{23}{322.767}{208.219}{\Character{48}}%
\TextAt{28}{322.203}{202.648}{P}%
\TextAt{3}{328.401}{205.474}{\Character{106}}%
\TextAt{21}{330.753}{205.474}{K}%
\TextAt{28}{337.894}{204.478}{p}%
\TextAt{9}{342.475}{207.422}{\Character{40}}%
\TextAt{13}{347.028}{207.422}{\Character{11}}%
\TextAt{27}{354.473}{205.629}{P}%
\TextAt{9}{361.567}{207.422}{\Character{41}}%
\TextAt{13}{366.119}{207.422}{\Character{59}}%
\TextAt{13}{371.343}{207.422}{L}%
\TextAt{27}{379.328}{205.629}{p}%
\TextAt{16}{387.135}{207.422}{\Character{106}}%
\TextAt{13}{393.777}{207.422}{K}%
\TextAt{27}{403.736}{205.629}{p}%
\TextAt{9}{408.223}{207.422}{\Character{41}}%
\TextAt{9}{256.166}{173.15}{\Character{61}}%
\TextAt{20}{268.591}{184.508}{\Character{89}}%
\TextAt{27}{272.932}{159.072}{P}%
\TextAt{9}{283.867}{173.15}{\Character{40}}%
\TextAt{13}{288.42}{173.15}{\Character{11}}%
\TextAt{27}{295.865}{171.357}{P}%
\TextAt{13}{302.959}{173.15}{\Character{59}}%
\TextAt{13}{308.183}{173.15}{L}%
\TextAt{3}{316.168}{178.087}{\Character{48}}%
\TextAt{27}{316.168}{170.195}{P}%
\TextAt{16}{326.582}{173.15}{\Character{106}}%
\TextAt{13}{333.224}{173.15}{K}%
\TextAt{3}{344.041}{178.087}{\Character{48}}%
\TextAt{27}{343.183}{170.195}{P}%
\TextAt{9}{350.276}{173.15}{\Character{41}}%
\TextAt{16}{354.829}{173.15}{\Character{106}}%
\TextAt{21}{358.15}{171.357}{L}%
\TextAt{9}{367.728}{173.15}{\Character{61}}%
\TextAt{9}{380.137}{173.15}{\Character{91}}%
\TextAt{13}{383.405}{173.15}{\Character{11}\Character{59}}%
\TextAt{13}{396.149}{173.15}{L}%
\TextAt{3}{404.135}{178.087}{\Character{48}}%
\TextAt{16}{410.251}{173.15}{\Character{106}}%
\TextAt{13}{416.893}{173.15}{K}%
\TextAt{3}{427.711}{178.087}{\Character{48}}%
\TextAt{9}{430.506}{173.15}{\Character{93}}%
\TextAt{16}{433.757}{173.15}{\Character{106}}%
\TextAt{21}{437.078}{171.357}{L}%
\TextAt{13}{443.335}{173.15}{\Character{58}}%
\Vector{q 0 G 0 g 0.399 w 452.334 -638.134 m 452.334 -630.064 l S Q}%
\Vector{q 0 G 0 g 0.399 w 452.534 -630.263 m 459.707 -630.263 l S Q}%
\Vector{q 0 G 0 g 0.399 w 452.534 -637.934 m 459.707 -637.934 l S Q}%
\Vector{q 0 G 0 g 0.399 w 459.906 -638.134 m 459.906 -630.064 l S Q}%
\TextAt{8}{97.159}{107.013}{If}%
\TextAt{13}{108.944}{107.013}{L}%
\TextAt{16}{120.388}{107.013}{\Character{106}}%
\TextAt{13}{127.188}{107.013}{K}%
\TextAt{8}{141.996}{107.013}{is}%
\TextAt{8}{153.853}{107.013}{an}%
\TextAt{8}{170.205}{107.013}{abelian}%
\TextAt{8}{210.946}{107.013}{extension}%
\TextAt{8}{262.805}{107.013}{of}%
\TextAt{8}{276.228}{107.013}{infinite}%
\TextAt{8}{315.988}{107.013}{degree,}%
\TextAt{8}{355.772}{107.013}{then}%
\TextAt{8}{382.525}{107.013}{we}%
\TextAt{8}{399.845}{107.013}{define}%
\TextAt{8}{433.747}{107.013}{the}%
\TextAt{8}{453.997}{107.013}{homomorphism}%
\TextAt{9}{79.226}{89.967}{\Character{91}}%
\TextAt{13}{87.22}{89.967}{\Character{59}}%
\TextAt{13}{92.444}{89.967}{L}%
\TextAt{16}{105.19}{89.967}{\Character{106}}%
\TextAt{13}{113.271}{89.967}{K}%
\TextAt{9}{124.089}{89.967}{\Character{93}}%
\TextAt{9}{132.087}{89.967}{\Character{58}}%
\TextAt{13}{140.114}{89.967}{I}%
\TextAt{21}{145.276}{88.174}{K}%
\NativeArrowAt{158.24}{89.967}{11.955}{19.26}{\longrightarrow}%
\TextAt{13}{182.262}{89.967}{G}%
\TextAt{9}{191.495}{89.967}{\Character{40}}%
\TextAt{13}{196.048}{89.967}{L}%
\TextAt{16}{208.773}{89.967}{\Character{106}}%
\TextAt{13}{216.855}{89.967}{K}%
\TextAt{9}{227.672}{89.967}{\Character{41}}%
\TextAt{8}{236.967}{89.967}{by}%
\TextAt{8}{254.071}{89.967}{its}%
\TextAt{8}{271.235}{89.967}{restrictions}%
\TextAt{9}{332.692}{89.967}{\Character{91}}%
\TextAt{13}{340.686}{89.967}{\Character{59}}%
\TextAt{13}{345.91}{89.967}{L}%
\TextAt{16}{358.656}{89.967}{\Character{106}}%
\TextAt{13}{366.737}{89.967}{K}%
\TextAt{9}{377.555}{89.967}{\Character{93}}%
\TextAt{16}{380.807}{89.967}{\Character{106}}%
\TextAt{21}{384.127}{88.174}{L}%
\TextAt{23}{389.886}{90.443}{\Character{48}}%
\TextAt{9}{397.83}{89.967}{\Character{58}\Character{61}}%
\TextAt{9}{414.926}{89.967}{\Character{91}}%
\TextAt{13}{422.941}{89.967}{\Character{59}}%
\TextAt{13}{428.165}{89.967}{L}%
\TextAt{3}{436.149}{94.306}{\Character{48}}%
\TextAt{16}{443.705}{89.967}{\Character{106}}%
\TextAt{13}{451.787}{89.967}{K}%
\TextAt{9}{462.604}{89.967}{\Character{93}}%
\TextAt{8}{470.598}{89.967}{to}%
\TextAt{8}{485.754}{89.967}{the}%
\TextAt{8}{506.755}{89.967}{finite}%
\TextAt{8}{79.226}{72.921}{subextensions}%
\TextAt{13}{153.226}{72.921}{L}%
\TextAt{3}{161.191}{77.26}{\Character{48}}%
\TextAt{8}{167.559}{72.921}{of}%
\TextAt{13}{180.564}{72.921}{L}%
\TextAt{16}{191.849}{72.921}{\Character{106}}%
\TextAt{13}{198.491}{72.921}{K}%
\TextAt{8}{209.308}{72.921}{.}%
\TextAt{8}{217.648}{72.921}{In}%
\TextAt{8}{231.945}{72.921}{other}%
\TextAt{8}{262.189}{72.921}{words,}%
\TextAt{8}{298.739}{72.921}{if}%
\TextAt{13}{309.138}{72.921}{\Character{11}}%
\NativeMathAt{319.98}{72.921}{11.955}{\in}%
\TextAt{13}{331.271}{72.921}{I}%
\TextAt{21}{336.434}{71.128}{K}%
\TextAt{8}{344.637}{72.921}{,}%
\TextAt{8}{351.526}{72.921}{then}%
\TextAt{8}{377.861}{72.921}{the}%
\TextAt{8}{397.693}{72.921}{elements}%
\TextAt{9}{445.224}{72.921}{\Character{91}}%
\TextAt{13}{448.476}{72.921}{\Character{11}\Character{59}}%
\TextAt{13}{461.22}{72.921}{L}%
\TextAt{3}{469.206}{77.26}{\Character{48}}%
\TextAt{16}{475.322}{72.921}{\Character{106}}%
\TextAt{13}{481.964}{72.921}{K}%
\TextAt{9}{492.782}{72.921}{\Character{93}}%
\TextAt{8}{499.606}{72.921}{define,}%
\TextAt{8}{297.211}{43.033}{353}%
\Destination{78.23}{809.815}{nav.page.353}%
\Destination{79.36}{776.155}{nav.print.387}%
\Destination{79.36}{476.955}{nav.number.6.5.2}%
\Hotlink{383.384}{323.655}{13.995}{11.047}{nav.chapter.4}%
\Hotlink{239.547}{664.455}{25.078}{11.047}{nav.number.6.5.7}%
\Hotlink{403.529}{323.655}{25.078}{11.047}{nav.number.4.6.4}%
\Hotlink{101.489}{266.655}{25.078}{11.047}{nav.number.6.2.2}%
}

% Page 366 (body page 354)
\Page{%
\TextAt{2}{79.226}{796.943}{VI.}%
\TextAt{2}{98.012}{796.943}{Global}%
\TextAt{2}{133.248}{796.943}{Class}%
\TextAt{2}{161.841}{796.943}{Field}%
\TextAt{2}{189.572}{796.943}{Theory}%
\Vector{q 0 G 0 g 0.299 w 7.226 18.729 m 460.769 18.729 l S Q}%
\TextAt{8}{79.226}{761.124}{by}%
\TextAt{8}{94.984}{761.124}{(5.2),}%
\TextAt{8}{125.811}{761.124}{an}%
\TextAt{8}{141.569}{761.124}{element}%
\TextAt{8}{184.31}{761.124}{of}%
\TextAt{8}{197.139}{761.124}{the}%
\TextAt{8}{216.795}{761.124}{projective}%
\TextAt{8}{270.271}{761.124}{limit}%
\NativeMathAt{297.734}{761.124}{11.955}{\varprojlim}%
\TextAt{21}{313.993}{755.684}{L}%
\TextAt{23}{319.752}{757.953}{\Character{48}}%
\TextAt{13}{324.927}{761.124}{G}%
\TextAt{9}{334.16}{761.124}{\Character{40}}%
\TextAt{13}{338.713}{761.124}{L}%
\TextAt{3}{346.677}{765.462}{\Character{48}}%
\TextAt{16}{352.793}{761.124}{\Character{106}}%
\TextAt{13}{359.435}{761.124}{K}%
\TextAt{9}{370.253}{761.124}{\Character{41}}%
\TextAt{8}{374.805}{761.124}{,}%
\TextAt{8}{381.553}{761.124}{and}%
\TextAt{9}{403.815}{761.124}{\Character{91}}%
\TextAt{13}{407.067}{761.124}{\Character{11}\Character{59}}%
\TextAt{13}{419.811}{761.124}{L}%
\TextAt{16}{431.118}{761.124}{\Character{106}}%
\TextAt{13}{437.76}{761.124}{K}%
\TextAt{9}{448.577}{761.124}{\Character{93}}%
\TextAt{8}{455.226}{761.124}{is}%
\TextAt{8}{466.489}{761.124}{precisely}%
\TextAt{8}{513.844}{761.124}{this}%
\TextAt{8}{79.226}{744.078}{element,}%
\TextAt{8}{125.72}{744.078}{once}%
\TextAt{8}{152.38}{744.078}{we}%
\TextAt{8}{169.607}{744.078}{identify}%
\TextAt{13}{212.204}{744.078}{G}%
\TextAt{9}{221.437}{744.078}{\Character{40}}%
\TextAt{13}{225.989}{744.078}{L}%
\TextAt{16}{237.275}{744.078}{\Character{106}}%
\TextAt{13}{243.917}{744.078}{K}%
\TextAt{9}{254.734}{744.078}{\Character{41}}%
\TextAt{9}{262.589}{744.078}{\Character{61}}%
\NativeMathAt{275.033}{744.078}{11.955}{\varprojlim}%
\TextAt{21}{291.292}{738.638}{L}%
\TextAt{23}{297.051}{740.908}{\Character{48}}%
\TextAt{13}{302.226}{744.078}{G}%
\TextAt{9}{311.46}{744.078}{\Character{40}}%
\TextAt{13}{316.012}{744.078}{L}%
\TextAt{3}{323.976}{748.416}{\Character{48}}%
\TextAt{16}{330.092}{744.078}{\Character{106}}%
\TextAt{13}{336.734}{744.078}{K}%
\TextAt{9}{347.552}{744.078}{\Character{41}}%
\TextAt{8}{352.104}{744.078}{.}%
\TextAt{8}{360.552}{744.078}{Again}%
\TextAt{8}{394.696}{744.078}{one}%
\TextAt{8}{416.156}{744.078}{has}%
\TextAt{8}{437.03}{744.078}{the}%
\TextAt{8}{457.186}{744.078}{equation}%
\TextAt{9}{233.283}{713.164}{\Character{91}}%
\TextAt{13}{236.535}{713.164}{\Character{11}\Character{59}}%
\TextAt{13}{249.279}{713.164}{L}%
\TextAt{16}{260.586}{713.164}{\Character{106}}%
\TextAt{13}{267.228}{713.164}{K}%
\TextAt{9}{278.046}{713.164}{\Character{93}}%
\TextAt{9}{284.597}{713.164}{\Character{61}}%
\TextAt{20}{297.044}{724.521}{\Character{89}}%
\TextAt{27}{302.687}{699.415}{p}%
\TextAt{9}{312.32}{713.164}{\Character{40}}%
\TextAt{13}{316.872}{713.164}{\Character{11}}%
\TextAt{27}{324.318}{711.371}{p}%
\TextAt{13}{328.804}{713.164}{\Character{59}}%
\TextAt{13}{334.028}{713.164}{L}%
\TextAt{27}{342.013}{711.371}{p}%
\TextAt{16}{349.821}{713.164}{\Character{106}}%
\TextAt{13}{356.462}{713.164}{K}%
\TextAt{27}{366.422}{711.371}{p}%
\TextAt{9}{370.908}{713.164}{\Character{41}}%
\TextAt{13}{375.461}{713.164}{\Character{59}}%
\TextAt{8}{79.226}{675.375}{where}%
\TextAt{13}{112.654}{675.375}{L}%
\TextAt{27}{120.618}{673.582}{p}%
\TextAt{8}{128.621}{675.375}{does}%
\TextAt{8}{154.637}{675.375}{not}%
\TextAt{8}{175.069}{675.375}{denote}%
\TextAt{8}{212.406}{675.375}{the}%
\TextAt{8}{232.181}{675.375}{completion}%
\TextAt{8}{291.635}{675.375}{of}%
\TextAt{13}{304.583}{675.375}{L}%
\TextAt{8}{316.064}{675.375}{with}%
\TextAt{8}{342.342}{675.375}{respect}%
\TextAt{8}{382.011}{675.375}{to}%
\TextAt{8}{395.939}{675.375}{a}%
\TextAt{8}{405.313}{675.375}{place}%
\TextAt{8}{434.843}{675.375}{above}%
\TextAt{14}{467.638}{675.375}{p}%
\TextAt{8}{473.621}{675.375}{,}%
\TextAt{8}{480.465}{675.375}{but}%
\TextAt{8}{501.543}{675.375}{rather}%
\TextAt{8}{79.226}{658.329}{the}%
\TextAt{19}{101.039}{658.329}{localization}%
\TextAt{8}{157.527}{658.329}{,}%
\TextAt{8}{166.746}{658.329}{i.e.,}%
\TextAt{8}{190.922}{658.329}{the}%
\TextAt{8}{212.734}{658.329}{union}%
\TextAt{8}{246.909}{658.329}{of}%
\TextAt{8}{261.895}{658.329}{the}%
\TextAt{8}{283.707}{658.329}{completions}%
\TextAt{13}{349.814}{658.329}{L}%
\TextAt{3}{357.778}{662.668}{\Character{48}}%
\TextAt{27}{357.778}{655.374}{p}%
\TextAt{16}{368.408}{658.329}{\Character{106}}%
\TextAt{13}{377.873}{658.329}{K}%
\TextAt{27}{387.832}{656.536}{p}%
\TextAt{8}{397.872}{658.329}{of}%
\TextAt{8}{412.858}{658.329}{all}%
\TextAt{8}{430.774}{658.329}{finite}%
\TextAt{8}{462.342}{658.329}{subextensions}%
\TextAt{8}{79.226}{641.283}{(see}%
\TextAt{8}{103.733}{641.283}{chap.}%
\TextAt{8}{135.665}{641.283}{II,}%
\TextAt{8}{152.616}{641.283}{§8).}%
\TextAt{8}{180.274}{641.283}{Then}%
\TextAt{13}{211.87}{641.283}{L}%
\TextAt{27}{219.835}{639.49}{p}%
\TextAt{16}{229.413}{641.283}{\Character{106}}%
\TextAt{13}{237.826}{641.283}{K}%
\TextAt{27}{247.785}{639.49}{p}%
\TextAt{8}{257.208}{641.283}{is}%
\TextAt{8}{270.011}{641.283}{Galois,}%
\TextAt{13}{310.475}{641.283}{G}%
\TextAt{9}{319.709}{641.283}{\Character{40}}%
\TextAt{13}{324.261}{641.283}{L}%
\TextAt{27}{332.226}{639.49}{p}%
\TextAt{16}{341.804}{641.283}{\Character{106}}%
\TextAt{13}{350.217}{641.283}{K}%
\TextAt{27}{360.176}{639.49}{p}%
\TextAt{9}{364.663}{641.283}{\Character{41}}%
\NativeMathAt{374.307}{641.283}{11.955}{\subseteq}%
\TextAt{13}{388.697}{641.283}{G}%
\TextAt{9}{397.931}{641.283}{\Character{40}}%
\TextAt{13}{402.483}{641.283}{L}%
\TextAt{16}{415.539}{641.283}{\Character{106}}%
\TextAt{13}{423.952}{641.283}{K}%
\TextAt{9}{434.77}{641.283}{\Character{41}}%
\TextAt{8}{439.322}{641.283}{,}%
\TextAt{8}{447.77}{641.283}{and}%
\TextAt{8}{471.572}{641.283}{the}%
\TextAt{8}{492.767}{641.283}{product}%
\TextAt{20}{79.226}{633.204}{\Character{81}}%
\TextAt{27}{90.517}{620.75}{p}%
\TextAt{9}{95.004}{624.237}{\Character{40}}%
\TextAt{13}{99.556}{624.237}{\Character{11}}%
\TextAt{27}{107.002}{622.444}{p}%
\TextAt{13}{111.489}{624.237}{\Character{59}}%
\TextAt{13}{116.713}{624.237}{L}%
\TextAt{27}{124.697}{622.444}{p}%
\TextAt{16}{132.505}{624.237}{\Character{106}}%
\TextAt{13}{139.147}{624.237}{K}%
\TextAt{27}{149.106}{622.444}{p}%
\TextAt{9}{153.593}{624.237}{\Character{41}}%
\TextAt{8}{160.855}{624.237}{converges}%
\TextAt{8}{212.09}{624.237}{in}%
\TextAt{8}{224.555}{624.237}{the}%
\TextAt{8}{243.524}{624.237}{profinite}%
\TextAt{8}{289.164}{624.237}{group}%
\TextAt{8}{321.152}{624.237}{to}%
\TextAt{8}{334.274}{624.237}{the}%
\TextAt{8}{353.243}{624.237}{element}%
\TextAt{9}{395.297}{624.237}{\Character{91}}%
\TextAt{13}{398.548}{624.237}{\Character{11}\Character{59}}%
\TextAt{13}{411.292}{624.237}{L}%
\TextAt{16}{422.599}{624.237}{\Character{106}}%
\TextAt{13}{429.241}{624.237}{K}%
\TextAt{9}{440.059}{624.237}{\Character{93}}%
\TextAt{8}{443.311}{624.237}{.}%
\TextAt{8}{451.363}{624.237}{Indeed,}%
\TextAt{8}{491.694}{624.237}{if}%
\TextAt{13}{501.23}{624.237}{L}%
\TextAt{3}{509.194}{628.576}{\Character{48}}%
\TextAt{16}{515.31}{624.237}{\Character{106}}%
\TextAt{13}{521.952}{624.237}{K}%
\TextAt{8}{79.226}{607.191}{varies}%
\TextAt{8}{112.422}{607.191}{over}%
\TextAt{8}{137.752}{607.191}{the}%
\TextAt{8}{158.192}{607.191}{finite}%
\TextAt{8}{188.387}{607.191}{subextensions}%
\TextAt{8}{262.996}{607.191}{of}%
\TextAt{13}{276.61}{607.191}{L}%
\TextAt{16}{288.379}{607.191}{\Character{106}}%
\TextAt{13}{295.504}{607.191}{K}%
\TextAt{8}{306.321}{607.191}{,}%
\TextAt{8}{313.825}{607.191}{then}%
\TextAt{8}{340.769}{607.191}{the}%
\TextAt{8}{361.209}{607.191}{sets}%
\TextAt{13}{384.375}{607.191}{S}%
\TextAt{21}{391.574}{605.398}{L}%
\TextAt{23}{397.333}{607.667}{\Character{48}}%
\TextAt{9}{404.32}{607.191}{\Character{61}}%
\TextAt{16}{417.228}{607.191}{\Character{102}}%
\TextAt{14}{423.206}{607.191}{p}%
\TextAt{16}{432.993}{607.191}{\Character{106}}%
\TextAt{9}{440.118}{607.191}{\Character{40}}%
\TextAt{13}{444.671}{607.191}{\Character{11}}%
\TextAt{27}{452.116}{605.398}{p}%
\TextAt{13}{456.603}{607.191}{\Character{59}}%
\TextAt{13}{461.827}{607.191}{L}%
\TextAt{3}{469.812}{611.53}{\Character{48}}%
\TextAt{27}{469.812}{604.236}{p}%
\TextAt{16}{478.103}{607.191}{\Character{106}}%
\TextAt{13}{485.228}{607.191}{K}%
\TextAt{27}{495.187}{605.398}{p}%
\TextAt{9}{499.674}{607.191}{\Character{41}}%
\NativeMathAt{508.03}{607.191}{11.955}{\neq}%
\TextAt{9}{520.93}{607.191}{1}%
\TextAt{16}{526.792}{607.191}{\Character{103}}%
\TextAt{8}{79.226}{590.145}{are}%
\TextAt{8}{98.737}{590.145}{all}%
\TextAt{8}{114.996}{590.145}{finite,}%
\TextAt{8}{148.16}{590.145}{so}%
\TextAt{8}{162.53}{590.145}{that}%
\TextAt{8}{187.899}{590.145}{we}%
\TextAt{8}{205.126}{590.145}{may}%
\TextAt{8}{230.495}{590.145}{write}%
\TextAt{8}{260.407}{590.145}{down}%
\TextAt{8}{291.299}{590.145}{the}%
\TextAt{8}{311.455}{590.145}{finite}%
\TextAt{8}{341.367}{590.145}{products}%
\TextAt{13}{214.77}{561.174}{\Character{27}}%
\TextAt{21}{221.424}{559.381}{L}%
\TextAt{23}{227.183}{561.65}{\Character{48}}%
\TextAt{9}{233.686}{561.174}{\Character{61}}%
\TextAt{20}{249.987}{572.532}{\Character{89}}%
\TextAt{27}{246.112}{547.16}{p}%
\NativeMathAt{250.101}{547.16}{7.97}{\in}%
\TextAt{21}{255.746}{547.16}{S}%
\TextAt{22}{260.896}{544.869}{L}%
\TextAt{23}{265.955}{546.584}{\Character{48}}%
\TextAt{9}{269.138}{561.174}{\Character{40}}%
\TextAt{13}{273.69}{561.174}{\Character{11}}%
\TextAt{27}{281.135}{559.381}{p}%
\TextAt{13}{285.622}{561.174}{\Character{59}}%
\TextAt{13}{290.846}{561.174}{L}%
\TextAt{27}{298.831}{559.381}{p}%
\TextAt{16}{306.639}{561.174}{\Character{106}}%
\TextAt{13}{313.28}{561.174}{K}%
\TextAt{27}{323.24}{559.381}{p}%
\TextAt{9}{327.726}{561.174}{\Character{41}}%
\NativeMathAt{335.6}{561.174}{11.955}{\in}%
\TextAt{13}{346.891}{561.174}{G}%
\TextAt{9}{356.124}{561.174}{\Character{40}}%
\TextAt{13}{360.676}{561.174}{L}%
\TextAt{16}{371.962}{561.174}{\Character{106}}%
\TextAt{13}{378.603}{561.174}{K}%
\TextAt{9}{389.421}{561.174}{\Character{41}}%
\TextAt{13}{393.973}{561.174}{\Character{58}}%
\TextAt{8}{79.226}{521.785}{They}%
\TextAt{8}{110.041}{521.785}{converge}%
\TextAt{8}{158.431}{521.785}{to}%
\TextAt{9}{173.322}{521.785}{\Character{91}}%
\TextAt{13}{176.574}{521.785}{\Character{11}\Character{59}}%
\TextAt{13}{189.306}{521.785}{L}%
\TextAt{16}{201.614}{521.785}{\Character{106}}%
\TextAt{13}{209.246}{521.785}{K}%
\TextAt{9}{220.063}{521.785}{\Character{93}}%
\TextAt{8}{223.315}{521.785}{,}%
\TextAt{8}{231.19}{521.785}{for}%
\TextAt{8}{249.656}{521.785}{if}%
\TextAt{9}{260.96}{521.785}{\Character{91}}%
\TextAt{13}{264.212}{521.785}{\Character{11}\Character{59}}%
\TextAt{13}{276.956}{521.785}{L}%
\TextAt{16}{289.253}{521.785}{\Character{106}}%
\TextAt{13}{296.884}{521.785}{K}%
\TextAt{9}{307.702}{521.785}{\Character{93}}%
\TextAt{13}{310.953}{521.785}{G}%
\TextAt{9}{320.187}{521.785}{\Character{40}}%
\TextAt{13}{324.739}{521.785}{L}%
\TextAt{16}{337.014}{521.785}{\Character{106}}%
\TextAt{13}{344.646}{521.785}{N}%
\TextAt{9}{355.268}{521.785}{\Character{41}}%
\TextAt{8}{364.299}{521.785}{is}%
\TextAt{8}{376.643}{521.785}{one}%
\TextAt{8}{398.684}{521.785}{of}%
\TextAt{8}{412.594}{521.785}{the}%
\TextAt{8}{433.331}{521.785}{fundamental}%
\TextAt{8}{501.555}{521.785}{neigh-}%
\TextAt{8}{79.226}{504.739}{bourhoods}%
\TextAt{8}{136.85}{504.739}{(i.e.,}%
\TextAt{13}{163.51}{504.739}{N}%
\TextAt{16}{177.454}{504.739}{\Character{106}}%
\TextAt{13}{184.095}{504.739}{K}%
\TextAt{8}{198.81}{504.739}{is}%
\TextAt{8}{210.574}{504.739}{one}%
\TextAt{8}{232.034}{504.739}{of}%
\TextAt{8}{245.364}{504.739}{the}%
\TextAt{8}{265.52}{504.739}{finite}%
\TextAt{8}{295.432}{504.739}{subextensions}%
\TextAt{8}{369.757}{504.739}{of}%
\TextAt{13}{383.087}{504.739}{L}%
\TextAt{16}{394.373}{504.739}{\Character{106}}%
\TextAt{13}{401.014}{504.739}{K}%
\TextAt{8}{411.832}{504.739}{),}%
\TextAt{8}{423.536}{504.739}{then}%
\TextAt{13}{243.443}{475.768}{\Character{27}}%
\TextAt{21}{250.097}{473.975}{L}%
\TextAt{23}{255.856}{476.244}{\Character{48}}%
\NativeMathAt{262.359}{475.768}{11.955}{\in}%
\TextAt{9}{273.65}{475.768}{\Character{91}}%
\TextAt{13}{276.902}{475.768}{\Character{11}\Character{59}}%
\TextAt{13}{289.646}{475.768}{L}%
\TextAt{16}{300.953}{475.768}{\Character{106}}%
\TextAt{13}{307.595}{475.768}{K}%
\TextAt{9}{318.413}{475.768}{\Character{93}}%
\TextAt{13}{321.664}{475.768}{G}%
\TextAt{9}{330.898}{475.768}{\Character{40}}%
\TextAt{13}{335.45}{475.768}{L}%
\TextAt{16}{346.735}{475.768}{\Character{106}}%
\TextAt{13}{353.377}{475.768}{N}%
\TextAt{9}{364}{475.768}{\Character{41}}%
\TextAt{8}{79.226}{446.797}{for}%
\TextAt{8}{97.111}{446.797}{all}%
\TextAt{13}{113.37}{446.797}{L}%
\TextAt{3}{121.335}{451.135}{\Character{48}}%
\NativeMathAt{127.451}{446.797}{11.955}{\supseteq}%
\TextAt{13}{140.07}{446.797}{N}%
\TextAt{8}{154.59}{446.797}{because}%
\TextAt{13}{172.309}{417.826}{\Character{27}}%
\TextAt{21}{178.963}{416.033}{L}%
\TextAt{23}{184.722}{418.302}{\Character{48}}%
\TextAt{16}{187.904}{417.826}{\Character{106}}%
\TextAt{21}{191.225}{416.033}{N}%
\TextAt{9}{202.615}{417.826}{\Character{61}}%
\TextAt{20}{215.04}{429.183}{\Character{89}}%
\TextAt{27}{220.684}{404.077}{p}%
\TextAt{9}{230.316}{417.826}{\Character{40}}%
\TextAt{13}{234.868}{417.826}{\Character{11}}%
\TextAt{27}{242.314}{416.033}{p}%
\TextAt{13}{246.801}{417.826}{\Character{59}}%
\TextAt{13}{252.025}{417.826}{N}%
\TextAt{27}{261.419}{416.033}{p}%
\TextAt{16}{269.227}{417.826}{\Character{106}}%
\TextAt{13}{275.868}{417.826}{K}%
\TextAt{27}{285.827}{416.033}{p}%
\TextAt{9}{290.314}{417.826}{\Character{41}}%
\TextAt{9}{298.169}{417.826}{\Character{61}}%
\TextAt{9}{310.59}{417.826}{\Character{91}}%
\TextAt{13}{313.865}{417.826}{\Character{11}\Character{59}}%
\TextAt{13}{326.597}{417.826}{N}%
\TextAt{16}{340.574}{417.826}{\Character{106}}%
\TextAt{13}{347.216}{417.826}{K}%
\TextAt{9}{358.033}{417.826}{\Character{93}}%
\TextAt{9}{364.584}{417.826}{\Character{61}}%
\TextAt{9}{377.006}{417.826}{\Character{91}}%
\TextAt{13}{380.283}{417.826}{\Character{11}\Character{59}}%
\TextAt{13}{393.027}{417.826}{L}%
\TextAt{16}{404.334}{417.826}{\Character{106}}%
\TextAt{13}{410.976}{417.826}{K}%
\TextAt{9}{421.794}{417.826}{\Character{93}}%
\TextAt{16}{425.045}{417.826}{\Character{106}}%
\TextAt{21}{428.366}{416.033}{N}%
\TextAt{13}{436.435}{417.826}{\Character{58}}%
\TextAt{8}{79.226}{379.487}{This}%
\TextAt{8}{106.049}{379.487}{shows}%
\TextAt{8}{139.771}{379.487}{that}%
\TextAt{9}{165.243}{379.487}{\Character{91}}%
\TextAt{13}{168.495}{379.487}{\Character{11}\Character{59}}%
\TextAt{13}{181.239}{379.487}{L}%
\TextAt{16}{192.723}{379.487}{\Character{106}}%
\TextAt{13}{199.541}{379.487}{K}%
\TextAt{9}{210.358}{379.487}{\Character{93}}%
\TextAt{8}{217.611}{379.487}{is}%
\TextAt{8}{229.478}{379.487}{the}%
\TextAt{8}{249.738}{379.487}{only}%
\TextAt{8}{275.533}{379.487}{accumulation}%
\TextAt{8}{347.511}{379.487}{point}%
\TextAt{8}{378.184}{379.487}{of}%
\TextAt{8}{391.618}{379.487}{the}%
\TextAt{8}{411.878}{379.487}{family}%
\TextAt{16}{447.751}{379.487}{\Character{102}}%
\TextAt{13}{453.729}{379.487}{\Character{27}}%
\TextAt{21}{460.382}{377.694}{L}%
\TextAt{23}{466.141}{379.963}{\Character{48}}%
\TextAt{16}{469.324}{379.487}{\Character{103}}%
\TextAt{8}{475.301}{379.487}{.}%
\TextAt{8}{484.06}{379.487}{It}%
\TextAt{8}{496.836}{379.487}{is}%
\TextAt{8}{508.704}{379.487}{clear}%
\TextAt{8}{79.226}{362.441}{that}%
\TextAt{8}{106.005}{362.441}{proposition}%
\TextAt{8}{168.948}{362.441}{(5.2)}%
\TextAt{8}{198.332}{362.441}{remains}%
\TextAt{8}{243.378}{362.441}{true}%
\TextAt{8}{269.499}{362.441}{for}%
\TextAt{8}{288.794}{362.441}{infinite}%
\TextAt{8}{329.871}{362.441}{extensions}%
\TextAt{13}{387.661}{362.441}{L}%
\TextAt{8}{400.932}{362.441}{and}%
\TextAt{13}{425.105}{362.441}{L}%
\TextAt{3}{433.069}{366.78}{\Character{48}}%
\TextAt{8}{441.171}{362.441}{of}%
\TextAt{8}{455.911}{362.441}{finite}%
\TextAt{8}{487.232}{362.441}{algebraic}%
\TextAt{8}{79.226}{345.395}{number}%
\TextAt{8}{121.822}{345.395}{fields}%
\TextAt{13}{151.794}{345.395}{K}%
\TextAt{8}{166.509}{345.395}{and}%
\TextAt{13}{189.272}{345.395}{K}%
\TextAt{3}{200.089}{349.734}{\Character{48}}%
\TextAt{8}{202.884}{345.395}{.}%
\TextAt{12}{79.226}{319.391}{Proposition}%
\TextAt{12}{152.358}{319.391}{(5.3).}%
\TextAt{8}{188.039}{319.391}{For}%
\TextAt{8}{209.374}{319.391}{every}%
\TextAt{8}{240.632}{319.391}{root}%
\TextAt{8}{266.043}{319.391}{of}%
\TextAt{8}{279.739}{319.391}{unity}%
\TextAt{13}{310.674}{319.391}{\Character{16}}%
\TextAt{8}{320.953}{319.391}{and}%
\TextAt{8}{344.081}{319.391}{every}%
\TextAt{8}{375.339}{319.391}{principal}%
\TextAt{8}{424.481}{319.391}{idèle}%
\TextAt{13}{452.152}{319.391}{a}%
\NativeMathAt{462.241}{319.391}{11.955}{\in}%
\TextAt{13}{474.155}{319.391}{K}%
\TextAt{3}{484.973}{323.73}{\Character{3}}%
\TextAt{8}{493.968}{319.391}{one}%
\TextAt{8}{515.793}{319.391}{has}%
\TextAt{9}{79.226}{302.345}{\Character{91}}%
\TextAt{13}{82.478}{302.345}{a\Character{59}}%
\TextAt{13}{93.847}{302.345}{K}%
\TextAt{9}{104.685}{302.345}{\Character{40}}%
\TextAt{13}{109.237}{302.345}{\Character{16}}%
\TextAt{9}{115.253}{302.345}{\Character{41}}%
\TextAt{16}{123.126}{302.345}{\Character{106}}%
\TextAt{13}{129.767}{302.345}{K}%
\TextAt{9}{140.585}{302.345}{\Character{93}}%
\TextAt{9}{147.136}{302.345}{\Character{61}}%
\TextAt{9}{159.558}{302.345}{1}%
\TextAt{8}{165.436}{302.345}{.}%
\TextAt{19}{79.226}{267.384}{Proof.}%
\TextAt{8}{116}{267.384}{By}%
\TextAt{8}{134.488}{267.384}{(5.2),}%
\TextAt{8}{165.871}{267.384}{we}%
\TextAt{8}{183.223}{267.384}{have}%
\TextAt{9}{210.342}{267.384}{\Character{91}}%
\TextAt{13}{213.594}{267.384}{N}%
\TextAt{21}{222.968}{265.524}{K}%
\TextAt{3}{230.672}{265.524}{\Character{106}}%
\TextAt{25}{233.025}{265.524}{Q}%
\TextAt{9}{239.722}{267.384}{\Character{40}}%
\TextAt{13}{244.274}{267.384}{a}%
\TextAt{9}{250.419}{267.384}{\Character{41}}%
\TextAt{13}{254.971}{267.384}{\Character{59}}%
\TextAt{18}{260.216}{267.384}{Q}%
\TextAt{9}{269.514}{267.384}{\Character{40}}%
\TextAt{13}{274.066}{267.384}{\Character{16}}%
\TextAt{9}{280.082}{267.384}{\Character{41}}%
\TextAt{16}{288.168}{267.384}{\Character{106}}%
\TextAt{18}{295.022}{267.384}{Q}%
\TextAt{9}{304.32}{267.384}{\Character{93}}%
\TextAt{9}{311.087}{267.384}{\Character{61}}%
\TextAt{9}{323.723}{267.384}{\Character{91}}%
\TextAt{13}{326.995}{267.384}{a\Character{59}}%
\TextAt{13}{338.364}{267.384}{K}%
\TextAt{9}{349.202}{267.384}{\Character{40}}%
\TextAt{13}{353.754}{267.384}{\Character{16}}%
\TextAt{9}{359.77}{267.384}{\Character{41}}%
\TextAt{16}{367.855}{267.384}{\Character{106}}%
\TextAt{13}{374.71}{267.384}{K}%
\TextAt{9}{385.527}{267.384}{\Character{93}}%
\TextAt{16}{388.779}{267.384}{\Character{106}}%
\TextAt{25}{392.1}{265.524}{Q}%
\TextAt{17}{398.299}{265.524}{\Character{40}}%
\TextAt{21}{401.592}{265.524}{\Character{16}}%
\TextAt{17}{405.861}{265.524}{\Character{41}}%
\TextAt{8}{409.652}{267.384}{.}%
\TextAt{8}{418.475}{267.384}{Hence}%
\TextAt{8}{453.377}{267.384}{we}%
\TextAt{8}{470.729}{267.384}{may}%
\TextAt{8}{496.223}{267.384}{assume}%
\TextAt{8}{79.226}{250.338}{that}%
\TextAt{13}{106.024}{250.338}{K}%
\TextAt{9}{122.598}{250.338}{\Character{61}}%
\TextAt{18}{137.459}{250.338}{Q}%
\TextAt{8}{146.757}{250.338}{.}%
\TextAt{8}{159.493}{250.338}{Likewise}%
\TextAt{8}{207.966}{250.338}{we}%
\TextAt{8}{226.622}{250.338}{may}%
\TextAt{8}{253.42}{250.338}{assume}%
\TextAt{8}{295.293}{250.338}{that}%
\TextAt{13}{322.091}{250.338}{\Character{16}}%
\TextAt{8}{333.433}{250.338}{has}%
\TextAt{8}{355.736}{250.338}{prime}%
\TextAt{8}{390.329}{250.338}{power}%
\TextAt{8}{425.902}{250.338}{order}%
\TextAt{13}{457.9}{250.338}{\Character{96}}%
\TextAt{21}{462.813}{254.676}{m}%
\NativeMathAt{476.557}{250.338}{11.955}{\neq}%
\TextAt{9}{491.405}{250.338}{2}%
\TextAt{8}{497.271}{250.338}{.}%
\TextAt{8}{510.007}{250.338}{Now}%
\TextAt{8}{79.226}{233.292}{let}%
\TextAt{13}{97.226}{233.292}{a}%
\NativeMathAt{108.557}{233.292}{11.955}{\in}%
\TextAt{18}{121.714}{233.292}{Q}%
\TextAt{3}{131.012}{237.63}{\Character{3}}%
\TextAt{8}{135.745}{233.292}{,}%
\TextAt{8}{144.262}{233.292}{let}%
\TextAt{13}{162.262}{233.292}{v}%
\TextAt{21}{167.921}{231.499}{p}%
\TextAt{8}{177.674}{233.292}{be}%
\TextAt{8}{194.693}{233.292}{the}%
\TextAt{8}{215.945}{233.292}{normalized}%
\TextAt{8}{276.875}{233.292}{exponential}%
\TextAt{8}{340.735}{233.292}{valuation}%
\TextAt{8}{392.902}{233.292}{of}%
\TextAt{18}{407.327}{233.292}{Q}%
\TextAt{8}{421.617}{233.292}{for}%
\TextAt{13}{440.597}{233.292}{p}%
\NativeMathAt{451.659}{233.292}{11.955}{\neq}%
\NativeMathAt{465.95}{233.292}{11.955}{\infty}%
\TextAt{8}{482.898}{233.292}{and}%
\TextAt{8}{506.755}{233.292}{write}%
\TextAt{13}{79.226}{216.246}{a}%
\TextAt{9}{88.997}{216.246}{\Character{61}}%
\TextAt{13}{101.728}{216.246}{u}%
\TextAt{21}{108.39}{214.453}{p}%
\TextAt{13}{113.151}{216.246}{p}%
\TextAt{21}{119.026}{220.584}{v}%
\TextAt{22}{123.15}{219.588}{p}%
\TextAt{17}{127.488}{220.584}{\Character{40}}%
\TextAt{21}{130.782}{220.584}{a}%
\TextAt{17}{135.28}{220.584}{\Character{41}}%
\TextAt{8}{139.071}{216.246}{.}%
\TextAt{8}{148.057}{216.246}{For}%
\TextAt{13}{169.205}{216.246}{p}%
\NativeMathAt{178.706}{216.246}{11.955}{\neq}%
\TextAt{13}{191.437}{216.246}{\Character{96}\Character{59}}%
\NativeMathAt{201.593}{216.246}{11.955}{\infty}%
\TextAt{8}{213.548}{216.246}{,}%
\TextAt{18}{220.921}{216.246}{Q}%
\TextAt{21}{230.22}{214.453}{p}%
\TextAt{9}{234.981}{216.246}{\Character{40}}%
\TextAt{13}{239.533}{216.246}{\Character{16}}%
\TextAt{9}{245.549}{216.246}{\Character{41}}%
\TextAt{16}{253.727}{216.246}{\Character{106}}%
\TextAt{18}{260.674}{216.246}{Q}%
\TextAt{21}{269.972}{214.453}{p}%
\TextAt{8}{278.81}{216.246}{is}%
\TextAt{8}{290.753}{216.246}{unramified}%
\TextAt{8}{349.464}{216.246}{and}%
\TextAt{9}{372.406}{216.246}{\Character{40}}%
\TextAt{13}{376.958}{216.246}{p\Character{59}}%
\TextAt{18}{388.078}{216.246}{Q}%
\TextAt{21}{397.376}{214.453}{p}%
\TextAt{9}{402.137}{216.246}{\Character{40}}%
\TextAt{13}{406.689}{216.246}{\Character{16}}%
\TextAt{9}{412.705}{216.246}{\Character{41}}%
\TextAt{16}{420.883}{216.246}{\Character{106}}%
\TextAt{18}{427.83}{216.246}{Q}%
\TextAt{21}{437.129}{214.453}{p}%
\TextAt{9}{441.89}{216.246}{\Character{41}}%
\TextAt{8}{450.518}{216.246}{is}%
\TextAt{8}{462.461}{216.246}{the}%
\TextAt{8}{482.797}{216.246}{Frobenius}%
\TextAt{8}{79.226}{199.2}{automorphism}%
\TextAt{13}{156.696}{199.2}{\Character{39}}%
\TextAt{21}{164.364}{197.407}{p}%
\TextAt{9}{172.446}{199.2}{\Character{58}}%
\TextAt{13}{179.018}{199.2}{\Character{16}}%
\NativeArrowAt{188.355}{199.2}{11.955}{11.955}{\mapsto}%
\TextAt{13}{203.631}{199.2}{\Character{16}}%
\TextAt{21}{209.646}{203.538}{p}%
\TextAt{8}{214.407}{199.2}{.}%
\TextAt{8}{222.856}{199.2}{From}%
\TextAt{8}{253.58}{199.2}{chap.}%
\TextAt{8}{284.472}{199.2}{V,}%
\TextAt{8}{300.397}{199.2}{(2.4),}%
\TextAt{8}{331.624}{199.2}{we}%
\TextAt{8}{348.851}{199.2}{thus}%
\TextAt{8}{374.603}{199.2}{get}%
\TextAt{9}{150.805}{149.264}{\Character{40}}%
\TextAt{13}{155.357}{149.264}{a\Character{59}}%
\TextAt{18}{166.746}{149.264}{Q}%
\TextAt{21}{176.045}{147.471}{p}%
\TextAt{9}{180.806}{149.264}{\Character{40}}%
\TextAt{13}{185.358}{149.264}{\Character{16}}%
\TextAt{9}{191.374}{149.264}{\Character{41}}%
\TextAt{16}{199.247}{149.264}{\Character{106}}%
\TextAt{18}{205.888}{149.264}{Q}%
\TextAt{21}{215.187}{147.471}{p}%
\TextAt{9}{219.948}{149.264}{\Character{41}}%
\TextAt{13}{224.5}{149.264}{\Character{16}}%
\TextAt{9}{233.837}{149.264}{\Character{61}}%
\TextAt{13}{246.262}{149.264}{\Character{16}}%
\TextAt{21}{252.278}{154.2}{n}%
\TextAt{22}{257.416}{153.204}{p}%
\TextAt{8}{286.163}{149.264}{with}%
\TextAt{13}{312.823}{149.264}{n}%
\TextAt{21}{319.811}{147.471}{p}%
\TextAt{9}{327.892}{149.264}{\Character{61}}%
\TextAt{20}{340.318}{180.946}{\Character{56}}%
\TextAt{20}{340.318}{170.186}{\Character{62}}%
\TextAt{20}{340.318}{166.599}{\Character{62}}%
\TextAt{20}{340.318}{163.013}{\Character{60}}%
\TextAt{20}{340.318}{141.493}{\Character{62}}%
\TextAt{20}{340.318}{137.907}{\Character{62}}%
\TextAt{20}{340.318}{134.32}{\Character{58}}%
\TextAt{13}{350.945}{168.617}{p}%
\TextAt{21}{356.82}{172.955}{v}%
\TextAt{22}{360.943}{171.959}{p}%
\TextAt{17}{365.282}{172.956}{\Character{40}}%
\TextAt{21}{368.575}{172.956}{a}%
\TextAt{17}{373.073}{172.956}{\Character{41}}%
\TextAt{8}{395.126}{168.617}{for}%
\TextAt{13}{413.011}{168.617}{p}%
\NativeMathAt{422.207}{168.617}{11.955}{\neq}%
\TextAt{13}{434.632}{168.617}{\Character{96}\Character{59}}%
\NativeMathAt{444.788}{168.617}{11.955}{\infty}%
\TextAt{13}{456.744}{168.617}{\Character{59}}%
\TextAt{13}{350.945}{148.162}{u}%
\TextAt{3}{357.607}{152.5}{\Character{0}}%
\TextAt{17}{364.194}{152.5}{1}%
\TextAt{21}{357.607}{145.207}{p}%
\TextAt{8}{395.126}{148.162}{for}%
\TextAt{13}{413.011}{148.162}{p}%
\TextAt{9}{422.207}{148.162}{\Character{61}}%
\TextAt{13}{434.632}{148.162}{\Character{96}\Character{59}}%
\TextAt{8}{350.945}{127.707}{sgn}%
\TextAt{9}{367.921}{127.707}{\Character{40}}%
\TextAt{13}{372.473}{127.707}{a}%
\TextAt{9}{378.618}{127.707}{\Character{41}}%
\TextAt{8}{395.126}{127.707}{for}%
\TextAt{13}{413.011}{127.707}{p}%
\TextAt{9}{422.207}{127.707}{\Character{61}}%
\NativeMathAt{434.632}{127.707}{11.955}{\infty}%
\TextAt{13}{446.587}{127.707}{\Character{58}}%
\TextAt{8}{79.226}{100.484}{Hence}%
\TextAt{9}{202.988}{83.438}{\Character{91}}%
\TextAt{13}{206.239}{83.438}{a\Character{59}}%
\TextAt{18}{217.628}{83.438}{Q}%
\TextAt{9}{226.927}{83.438}{\Character{40}}%
\TextAt{13}{231.479}{83.438}{\Character{16}}%
\TextAt{9}{237.495}{83.438}{\Character{41}}%
\TextAt{16}{245.368}{83.438}{\Character{106}}%
\TextAt{18}{252.01}{83.438}{Q}%
\TextAt{9}{261.308}{83.438}{\Character{93}}%
\TextAt{13}{264.56}{83.438}{\Character{16}}%
\TextAt{9}{273.896}{83.438}{\Character{61}}%
\TextAt{20}{286.322}{94.795}{\Character{89}}%
\TextAt{21}{291.829}{69.689}{p}%
\TextAt{9}{301.598}{83.438}{\Character{40}}%
\TextAt{13}{306.15}{83.438}{a\Character{59}}%
\TextAt{18}{317.539}{83.438}{Q}%
\TextAt{21}{326.838}{81.644}{p}%
\TextAt{9}{331.599}{83.438}{\Character{40}}%
\TextAt{13}{336.151}{83.438}{\Character{16}}%
\TextAt{9}{342.167}{83.438}{\Character{41}}%
\TextAt{16}{350.04}{83.438}{\Character{106}}%
\TextAt{18}{356.682}{83.438}{Q}%
\TextAt{21}{365.98}{81.644}{p}%
\TextAt{9}{370.741}{83.438}{\Character{41}}%
\TextAt{13}{375.293}{83.438}{\Character{16}}%
\TextAt{9}{384.63}{83.438}{\Character{61}}%
\TextAt{13}{397.055}{83.438}{\Character{16}}%
\TextAt{21}{403.071}{88.374}{\Character{11}}%
\TextAt{8}{297.211}{43.033}{354}%
\Destination{78.23}{809.815}{nav.page.354}%
\Destination{79.36}{776.155}{nav.print.388}%
\Destination{79.36}{334.355}{nav.print.389}%
\Destination{79.36}{333.355}{nav.number.6.5.3}%
\Hotlink{135.165}{639.055}{9.751}{11.047}{nav.chapter.2}%
\Hotlink{157.795}{639.055}{6.858}{11.047}{nav.section.2.8}%
\Hotlink{283.972}{196.955}{9.775}{11.047}{nav.chapter.5}%
\Hotlink{94.484}{758.955}{25.078}{11.047}{nav.number.6.5.2}%
\Hotlink{168.448}{360.255}{25.078}{11.047}{nav.number.6.5.2}%
\Hotlink{133.988}{265.155}{25.078}{11.047}{nav.number.6.5.2}%
\Hotlink{299.897}{196.955}{25.078}{11.047}{nav.number.5.2.4}%
}

% Page 367 (body page 355)
\Page{%
\TextAt{2}{79.226}{796.943}{VI.}%
\TextAt{2}{98.012}{796.943}{Global}%
\TextAt{2}{133.248}{796.943}{Class}%
\TextAt{2}{161.841}{796.943}{Field}%
\TextAt{2}{189.572}{796.943}{Theory}%
\Vector{q 0 G 0 g 0.299 w 7.226 18.729 m 460.769 18.729 l S Q}%
\TextAt{8}{79.226}{761.124}{where}%
\TextAt{13}{149.607}{744.078}{\Character{11}}%
\TextAt{9}{160.449}{744.078}{\Character{61}}%
\TextAt{20}{172.875}{755.435}{\Character{89}}%
\TextAt{21}{178.381}{730.329}{p}%
\TextAt{13}{190.143}{744.078}{n}%
\TextAt{21}{197.131}{742.284}{p}%
\TextAt{9}{205.213}{744.078}{\Character{61}}%
\TextAt{8}{217.638}{744.078}{sgn}%
\TextAt{9}{234.614}{744.078}{\Character{40}}%
\TextAt{13}{239.167}{744.078}{a}%
\TextAt{9}{245.312}{744.078}{\Character{41}}%
\TextAt{20}{256.783}{755.435}{\Character{89}}%
\TextAt{21}{251.857}{729.975}{p}%
\NativeMathAt{256.119}{729.975}{7.97}{\neq}%
\TextAt{21}{262.706}{729.975}{\Character{96}\Character{59}}%
\NativeMathAt{268.517}{729.975}{7.97}{\infty}%
\TextAt{13}{278.978}{744.078}{p}%
\TextAt{21}{284.853}{749.014}{v}%
\TextAt{22}{288.977}{748.018}{p}%
\TextAt{17}{293.316}{749.014}{\Character{40}}%
\TextAt{21}{296.609}{749.014}{a}%
\TextAt{17}{301.107}{749.014}{\Character{41}}%
\TextAt{13}{304.898}{744.078}{u}%
\TextAt{3}{311.561}{749.026}{\Character{0}}%
\TextAt{17}{318.147}{749.026}{1}%
\TextAt{21}{311.561}{740.748}{\Character{96}}%
\TextAt{9}{326.2}{744.078}{\Character{61}}%
\TextAt{8}{338.626}{744.078}{sgn}%
\TextAt{9}{355.602}{744.078}{\Character{40}}%
\TextAt{13}{360.154}{744.078}{a}%
\TextAt{9}{366.299}{744.078}{\Character{41}}%
\TextAt{20}{374.865}{755.435}{\Character{89}}%
\TextAt{21}{372.844}{729.975}{p}%
\NativeMathAt{377.107}{729.975}{7.97}{\neq}%
\NativeMathAt{383.693}{729.975}{7.97}{\infty}%
\TextAt{13}{394.154}{744.078}{p}%
\TextAt{21}{400.029}{749.014}{v}%
\TextAt{22}{404.153}{748.018}{p}%
\TextAt{17}{408.492}{749.014}{\Character{40}}%
\TextAt{21}{411.785}{749.014}{a}%
\TextAt{17}{416.283}{749.014}{\Character{41}}%
\TextAt{13}{420.074}{744.078}{a}%
\TextAt{3}{426.219}{749.014}{\Character{0}}%
\TextAt{17}{432.806}{749.014}{1}%
\TextAt{9}{440.859}{744.078}{\Character{61}}%
\TextAt{9}{453.269}{744.078}{1}%
\TextAt{13}{459.137}{744.078}{\Character{58}}%
\Vector{q 0 G 0 g 0.399 w 452.334 -62.312 m 452.334 -54.242 l S Q}%
\Vector{q 0 G 0 g 0.399 w 452.534 -54.442 m 459.707 -54.442 l S Q}%
\Vector{q 0 G 0 g 0.399 w 452.534 -62.113 m 459.707 -62.113 l S Q}%
\Vector{q 0 G 0 g 0.399 w 459.906 -62.312 m 459.906 -54.242 l S Q}%
\TextAt{8}{97.159}{682.687}{Since}%
\TextAt{8}{128.29}{682.687}{the}%
\TextAt{8}{149.02}{682.687}{extension}%
\TextAt{20}{204.098}{685.709}{\Character{101}}%
\TextAt{13}{201.359}{682.687}{K}%
\TextAt{16}{216.475}{682.687}{\Character{106}}%
\TextAt{13}{224.094}{682.687}{K}%
\TextAt{8}{239.383}{682.687}{is}%
\TextAt{8}{251.72}{682.687}{contained}%
\TextAt{8}{305.303}{682.687}{in}%
\TextAt{8}{319.53}{682.687}{the}%
\TextAt{8}{340.26}{682.687}{field}%
\TextAt{8}{366.19}{682.687}{of}%
\TextAt{8}{380.094}{682.687}{all}%
\TextAt{8}{396.926}{682.687}{roots}%
\TextAt{8}{427.16}{682.687}{of}%
\TextAt{8}{441.064}{682.687}{unity}%
\TextAt{8}{472.207}{682.687}{over}%
\TextAt{13}{497.827}{682.687}{K}%
\TextAt{8}{508.644}{682.687}{,}%
\TextAt{8}{516.51}{682.687}{the}%
\TextAt{8}{79.226}{665.641}{proposition}%
\TextAt{8}{140.59}{665.641}{implies}%
\TextAt{9}{180.148}{665.641}{\Character{91}}%
\TextAt{13}{183.4}{665.641}{a\Character{59}}%
\TextAt{20}{197.527}{668.663}{\Character{101}}%
\TextAt{13}{194.789}{665.641}{K}%
\TextAt{16}{208.927}{665.641}{\Character{106}}%
\TextAt{13}{215.569}{665.641}{K}%
\TextAt{9}{226.387}{665.641}{\Character{93}}%
\TextAt{9}{232.938}{665.641}{\Character{61}}%
\TextAt{9}{245.36}{665.641}{1}%
\TextAt{8}{254.966}{665.641}{for}%
\TextAt{8}{272.681}{665.641}{all}%
\TextAt{13}{288.771}{665.641}{a}%
\NativeMathAt{298.237}{665.641}{11.955}{\in}%
\TextAt{13}{309.528}{665.641}{K}%
\TextAt{3}{320.346}{669.979}{\Character{3}}%
\TextAt{8}{325.078}{665.641}{.}%
\TextAt{8}{333.47}{665.641}{The}%
\TextAt{8}{357.354}{665.641}{homomorphism}%
\TextAt{9}{439.855}{665.641}{\Character{91}}%
\TextAt{13}{446.835}{665.641}{\Character{59}}%
\TextAt{20}{454.817}{668.663}{\Character{101}}%
\TextAt{13}{452.079}{665.641}{K}%
\TextAt{16}{466.218}{665.641}{\Character{106}}%
\TextAt{13}{472.859}{665.641}{K}%
\TextAt{9}{483.677}{665.641}{\Character{93}}%
\TextAt{9}{490.228}{665.641}{\Character{58}}%
\TextAt{13}{496.822}{665.641}{I}%
\TextAt{21}{501.985}{663.848}{K}%
\NativeArrowAt{513.508}{665.641}{11.955}{19.26}{\longrightarrow}%
\TextAt{13}{79.226}{648.595}{G}%
\TextAt{9}{88.46}{648.595}{\Character{40}}%
\TextAt{20}{95.75}{651.617}{\Character{101}}%
\TextAt{13}{93.012}{648.595}{K}%
\TextAt{16}{107.151}{648.595}{\Character{106}}%
\TextAt{13}{113.792}{648.595}{K}%
\TextAt{9}{124.61}{648.595}{\Character{41}}%
\TextAt{8}{133.06}{648.595}{therefore}%
\TextAt{8}{182.16}{648.595}{induces}%
\TextAt{8}{223.835}{648.595}{a}%
\TextAt{8}{233.591}{648.595}{homomorphism}%
\TextAt{9}{229.596}{619.594}{\Character{91}}%
\TextAt{13}{236.745}{619.594}{\Character{59}}%
\TextAt{20}{244.727}{622.616}{\Character{101}}%
\TextAt{13}{241.989}{619.594}{K}%
\TextAt{16}{256.127}{619.594}{\Character{106}}%
\TextAt{13}{262.769}{619.594}{K}%
\TextAt{9}{273.587}{619.594}{\Character{93}}%
\TextAt{9}{280.138}{619.594}{\Character{58}}%
\TextAt{13}{286.732}{619.594}{C}%
\TextAt{21}{295.107}{617.8}{K}%
\NativeArrowAt{306.63}{619.594}{11.955}{19.26}{\longrightarrow}%
\TextAt{13}{329.212}{619.594}{G}%
\TextAt{9}{338.446}{619.594}{\Character{40}}%
\TextAt{20}{345.736}{622.616}{\Character{101}}%
\TextAt{13}{342.998}{619.594}{K}%
\TextAt{16}{357.137}{619.594}{\Character{106}}%
\TextAt{13}{363.778}{619.594}{K}%
\TextAt{9}{374.596}{619.594}{\Character{41}}%
\TextAt{13}{379.148}{619.594}{\Character{59}}%
\TextAt{8}{79.226}{590.593}{and}%
\TextAt{8}{103.102}{590.593}{we}%
\TextAt{8}{121.442}{590.593}{consider}%
\TextAt{8}{168.14}{590.593}{its}%
\TextAt{8}{185.572}{590.593}{composite}%
\TextAt{13}{241.703}{590.593}{v}%
\TextAt{21}{247.362}{588.799}{K}%
\TextAt{9}{260.782}{590.593}{\Character{58}}%
\TextAt{13}{269.251}{590.593}{C}%
\TextAt{21}{277.626}{588.799}{K}%
\NativeArrowAt{291.047}{590.593}{11.955}{19.26}{\longrightarrow}%
\TextAt{20}{316.19}{593.681}{\Character{98}}%
\TextAt{18}{315.525}{590.593}{Z}%
\TextAt{8}{328.506}{590.593}{with}%
\TextAt{13}{356.279}{590.593}{d}%
\TextAt{21}{362.362}{588.799}{K}%
\TextAt{9}{375.782}{590.593}{\Character{58}}%
\TextAt{13}{384.251}{590.593}{G}%
\TextAt{9}{393.485}{590.593}{\Character{40}}%
\TextAt{20}{400.775}{593.615}{\Character{101}}%
\TextAt{13}{398.037}{590.593}{K}%
\TextAt{16}{414.072}{590.593}{\Character{106}}%
\TextAt{13}{422.611}{590.593}{K}%
\TextAt{9}{433.428}{590.593}{\Character{41}}%
\NativeArrowAt{443.198}{590.593}{11.955}{19.26}{\longrightarrow}%
\TextAt{20}{468.341}{593.681}{\Character{98}}%
\TextAt{18}{467.677}{590.593}{Z}%
\TextAt{8}{475.647}{590.593}{.}%
\TextAt{8}{487.434}{590.593}{The}%
\TextAt{8}{512.601}{590.593}{pair}%
\TextAt{9}{79.226}{573.547}{\Character{40}}%
\TextAt{13}{83.778}{573.547}{d}%
\TextAt{21}{89.861}{571.753}{K}%
\TextAt{13}{98.064}{573.547}{\Character{59}}%
\TextAt{13}{103.288}{573.547}{v}%
\TextAt{21}{108.967}{571.753}{K}%
\TextAt{9}{117.17}{573.547}{\Character{41}}%
\TextAt{8}{125.62}{573.547}{is}%
\TextAt{8}{137.384}{573.547}{then}%
\TextAt{8}{164.044}{573.547}{a}%
\TextAt{8}{173.799}{573.547}{class}%
\TextAt{8}{201.236}{573.547}{field}%
\TextAt{8}{226.593}{573.547}{theory,}%
\TextAt{8}{265.615}{573.547}{for}%
\TextAt{8}{283.5}{573.547}{we}%
\TextAt{8}{300.727}{573.547}{have}%
\TextAt{8}{327.722}{573.547}{the}%
\TextAt{12}{79.226}{547.534}{Proposition}%
\TextAt{12}{152.637}{547.534}{(5.4).}%
\TextAt{8}{188.561}{547.534}{The}%
\TextAt{8}{213.223}{547.534}{map}%
\TextAt{13}{239.845}{547.534}{v}%
\TextAt{21}{245.505}{545.741}{K}%
\TextAt{9}{258.065}{547.534}{\Character{58}}%
\TextAt{13}{265.673}{547.534}{C}%
\TextAt{21}{274.048}{545.741}{K}%
\NativeArrowAt{286.608}{547.534}{11.955}{19.26}{\longrightarrow}%
\TextAt{20}{310.891}{550.623}{\Character{98}}%
\TextAt{18}{310.227}{547.534}{Z}%
\TextAt{8}{322.703}{547.534}{is}%
\TextAt{8}{335.075}{547.534}{surjective}%
\TextAt{8}{388.094}{547.534}{and}%
\TextAt{8}{411.465}{547.534}{is}%
\TextAt{8}{423.837}{547.534}{a}%
\TextAt{8}{434.201}{547.534}{henselian}%
\TextAt{8}{485.594}{547.534}{valuation}%
\TextAt{8}{79.226}{530.488}{with}%
\TextAt{8}{105.886}{530.488}{respect}%
\TextAt{8}{145.936}{530.488}{to}%
\TextAt{13}{160.246}{530.488}{d}%
\TextAt{21}{166.329}{528.695}{K}%
\TextAt{8}{174.532}{530.488}{.}%
\TextAt{19}{79.226}{495.51}{Proof.}%
\TextAt{8}{116}{495.51}{We}%
\TextAt{8}{135.554}{495.51}{first}%
\TextAt{8}{159.089}{495.51}{show}%
\TextAt{8}{187.501}{495.51}{surjectivity.}%
\TextAt{8}{252.073}{495.51}{If}%
\TextAt{13}{263.174}{495.51}{L}%
\TextAt{16}{274.46}{495.51}{\Character{106}}%
\TextAt{13}{281.101}{495.51}{K}%
\TextAt{8}{295.226}{495.51}{is}%
\TextAt{8}{306.399}{495.51}{a}%
\TextAt{8}{315.563}{495.51}{finite}%
\TextAt{8}{344.884}{495.51}{subextension}%
\TextAt{8}{414.004}{495.51}{of}%
\TextAt{20}{429.481}{498.532}{\Character{101}}%
\TextAt{13}{426.743}{495.51}{K}%
\TextAt{16}{440.882}{495.51}{\Character{106}}%
\TextAt{13}{447.523}{495.51}{K}%
\TextAt{8}{458.341}{495.51}{,}%
\TextAt{8}{465.018}{495.51}{then}%
\TextAt{8}{491.087}{495.51}{the}%
\TextAt{8}{510.652}{495.51}{map}%
\TextAt{9}{192.616}{466.508}{\Character{91}}%
\TextAt{13}{199.765}{466.508}{\Character{59}}%
\TextAt{13}{204.989}{466.508}{L}%
\TextAt{16}{216.295}{466.508}{\Character{106}}%
\TextAt{13}{222.937}{466.508}{K}%
\TextAt{9}{233.754}{466.508}{\Character{93}}%
\TextAt{9}{240.305}{466.508}{\Character{61}}%
\TextAt{20}{252.752}{477.866}{\Character{89}}%
\TextAt{27}{258.396}{452.76}{p}%
\TextAt{9}{268.028}{466.508}{\Character{40}}%
\TextAt{13}{276.478}{466.508}{\Character{59}}%
\TextAt{13}{281.702}{466.508}{L}%
\TextAt{27}{289.687}{464.715}{p}%
\TextAt{16}{297.494}{466.508}{\Character{106}}%
\TextAt{13}{304.136}{466.508}{K}%
\TextAt{27}{314.095}{464.715}{p}%
\TextAt{9}{318.582}{466.508}{\Character{41}}%
\TextAt{9}{326.437}{466.508}{\Character{58}}%
\TextAt{13}{333.028}{466.508}{I}%
\TextAt{21}{338.191}{464.715}{K}%
\NativeArrowAt{349.714}{466.508}{11.955}{19.26}{\longrightarrow}%
\TextAt{13}{372.296}{466.508}{G}%
\TextAt{9}{381.53}{466.508}{\Character{40}}%
\TextAt{13}{386.082}{466.508}{L}%
\TextAt{16}{397.368}{466.508}{\Character{106}}%
\TextAt{13}{404.009}{466.508}{K}%
\TextAt{9}{414.827}{466.508}{\Character{41}}%
\TextAt{8}{79.226}{428.14}{is}%
\TextAt{8}{92.008}{428.14}{surjective.}%
\TextAt{8}{152.024}{428.14}{Indeed,}%
\TextAt{8}{194.578}{428.14}{since}%
\TextAt{9}{224.264}{428.14}{\Character{40}}%
\TextAt{13}{233.732}{428.14}{\Character{59}}%
\TextAt{13}{238.956}{428.14}{L}%
\TextAt{27}{246.94}{426.347}{p}%
\TextAt{16}{256.483}{428.14}{\Character{106}}%
\TextAt{13}{264.859}{428.14}{K}%
\TextAt{27}{274.818}{426.347}{p}%
\TextAt{9}{279.305}{428.14}{\Character{41}}%
\TextAt{9}{288.893}{428.14}{\Character{58}}%
\TextAt{13}{297.22}{428.14}{K}%
\TextAt{3}{308.038}{432.478}{\Character{3}}%
\TextAt{27}{307.179}{425.184}{p}%
\NativeArrowAt{317.825}{428.14}{11.955}{19.26}{\longrightarrow}%
\TextAt{13}{342.142}{428.14}{G}%
\TextAt{9}{351.376}{428.14}{\Character{40}}%
\TextAt{13}{355.928}{428.14}{L}%
\TextAt{27}{363.893}{426.347}{p}%
\TextAt{16}{373.435}{428.14}{\Character{106}}%
\TextAt{13}{381.811}{428.14}{K}%
\TextAt{27}{391.77}{426.347}{p}%
\TextAt{9}{396.257}{428.14}{\Character{41}}%
\TextAt{8}{405.725}{428.14}{is}%
\TextAt{8}{418.507}{428.14}{surjective,}%
\TextAt{9}{475.442}{428.14}{\Character{91}}%
\TextAt{13}{478.694}{428.14}{I}%
\TextAt{21}{483.857}{426.347}{K}%
\TextAt{13}{492.06}{428.14}{\Character{59}}%
\TextAt{13}{497.284}{428.14}{L}%
\TextAt{16}{510.324}{428.14}{\Character{106}}%
\TextAt{13}{518.7}{428.14}{K}%
\TextAt{9}{529.518}{428.14}{\Character{93}}%
\TextAt{8}{79.226}{411.094}{contains}%
\TextAt{8}{125.943}{411.094}{all}%
\TextAt{8}{142.999}{411.094}{decomposition}%
\TextAt{8}{220.931}{411.094}{groups}%
\TextAt{13}{259.519}{411.094}{G}%
\TextAt{9}{268.752}{411.094}{\Character{40}}%
\TextAt{13}{273.305}{411.094}{L}%
\TextAt{27}{281.269}{409.301}{p}%
\TextAt{16}{290.435}{411.094}{\Character{106}}%
\TextAt{13}{298.436}{411.094}{K}%
\TextAt{27}{308.395}{409.301}{p}%
\TextAt{9}{312.882}{411.094}{\Character{41}}%
\TextAt{8}{317.434}{411.094}{.}%
\TextAt{8}{328.274}{411.094}{Thus}%
\TextAt{8}{358.72}{411.094}{all}%
\TextAt{14}{375.776}{411.094}{p}%
\TextAt{8}{386.454}{411.094}{split}%
\TextAt{8}{413.325}{411.094}{completely}%
\TextAt{8}{472.977}{411.094}{in}%
\TextAt{8}{487.427}{411.094}{the}%
\TextAt{8}{508.381}{411.094}{fixed}%
\TextAt{8}{79.226}{394.048}{field}%
\TextAt{13}{104.145}{394.048}{M}%
\TextAt{8}{120.178}{394.048}{of}%
\TextAt{9}{133.07}{394.048}{\Character{91}}%
\TextAt{13}{136.322}{394.048}{I}%
\TextAt{21}{141.485}{392.255}{K}%
\TextAt{13}{149.688}{394.048}{\Character{59}}%
\TextAt{13}{154.912}{394.048}{L}%
\TextAt{16}{166.217}{394.048}{\Character{106}}%
\TextAt{13}{172.859}{394.048}{K}%
\TextAt{9}{183.677}{394.048}{\Character{93}}%
\TextAt{8}{186.928}{394.048}{.}%
\TextAt{8}{195.231}{394.048}{By}%
\TextAt{8}{213.156}{394.048}{(3.8),}%
\TextAt{8}{244.032}{394.048}{this}%
\TextAt{8}{266.417}{394.048}{implies}%
\TextAt{8}{305.706}{394.048}{that}%
\TextAt{13}{330.637}{394.048}{M}%
\TextAt{9}{346.532}{394.048}{\Character{61}}%
\TextAt{13}{358.957}{394.048}{K}%
\TextAt{8}{369.775}{394.048}{,}%
\TextAt{8}{376.574}{394.048}{and}%
\TextAt{8}{398.898}{394.048}{so}%
\TextAt{9}{412.831}{394.048}{\Character{91}}%
\TextAt{13}{416.082}{394.048}{I}%
\TextAt{21}{421.245}{392.255}{K}%
\TextAt{13}{429.448}{394.048}{\Character{59}}%
\TextAt{13}{434.672}{394.048}{L}%
\TextAt{16}{445.977}{394.048}{\Character{106}}%
\TextAt{13}{452.619}{394.048}{K}%
\TextAt{9}{463.437}{394.048}{\Character{93}}%
\TextAt{9}{469.988}{394.048}{\Character{61}}%
\TextAt{13}{482.435}{394.048}{G}%
\TextAt{9}{491.668}{394.048}{\Character{40}}%
\TextAt{13}{496.221}{394.048}{L}%
\TextAt{16}{507.506}{394.048}{\Character{106}}%
\TextAt{13}{514.148}{394.048}{K}%
\TextAt{9}{524.965}{394.048}{\Character{41}}%
\TextAt{8}{529.518}{394.048}{.}%
\TextAt{8}{79.226}{377.002}{This}%
\TextAt{8}{106.764}{377.002}{yields}%
\TextAt{8}{140.482}{377.002}{furthermore}%
\TextAt{8}{206.013}{377.002}{that}%
\TextAt{9}{232.2}{377.002}{\Character{91}}%
\TextAt{13}{235.451}{377.002}{I}%
\TextAt{21}{240.614}{375.209}{K}%
\TextAt{13}{248.817}{377.002}{\Character{59}}%
\TextAt{20}{256.799}{380.024}{\Character{101}}%
\TextAt{13}{254.061}{377.002}{K}%
\TextAt{16}{269.593}{377.002}{\Character{106}}%
\TextAt{13}{277.628}{377.002}{K}%
\TextAt{9}{288.446}{377.002}{\Character{93}}%
\TextAt{9}{296.396}{377.002}{\Character{61}}%
\TextAt{9}{310.204}{377.002}{\Character{91}}%
\TextAt{13}{313.483}{377.002}{C}%
\TextAt{21}{321.858}{375.209}{K}%
\TextAt{13}{330.061}{377.002}{\Character{59}}%
\TextAt{20}{338.043}{380.024}{\Character{101}}%
\TextAt{13}{335.305}{377.002}{K}%
\TextAt{16}{350.837}{377.002}{\Character{106}}%
\TextAt{13}{358.872}{377.002}{K}%
\TextAt{9}{369.69}{377.002}{\Character{93}}%
\TextAt{8}{377.656}{377.002}{is}%
\TextAt{8}{390.238}{377.002}{dense}%
\TextAt{8}{422.976}{377.002}{in}%
\TextAt{13}{437.447}{377.002}{G}%
\TextAt{9}{446.68}{377.002}{\Character{40}}%
\TextAt{20}{453.971}{380.024}{\Character{101}}%
\TextAt{13}{451.233}{377.002}{K}%
\TextAt{16}{466.765}{377.002}{\Character{106}}%
\TextAt{13}{474.8}{377.002}{K}%
\TextAt{9}{485.618}{377.002}{\Character{41}}%
\TextAt{8}{490.17}{377.002}{.}%
\TextAt{8}{501.072}{377.002}{In}%
\TextAt{8}{516.51}{377.002}{the}%
\TextAt{8}{79.226}{359.956}{exact}%
\TextAt{8}{110.118}{359.956}{sequence}%
\TextAt{9}{226.923}{342.91}{1}%
\NativeArrowAt{236.097}{342.91}{11.955}{19.26}{\longrightarrow}%
\TextAt{13}{258.679}{342.91}{C}%
\TextAt{17}{267.913}{347.846}{0}%
\TextAt{21}{267.054}{339.954}{K}%
\NativeArrowAt{278.578}{342.91}{11.955}{19.26}{\longrightarrow}%
\TextAt{13}{301.16}{342.91}{C}%
\TextAt{21}{309.535}{341.117}{K}%
\TextAt{27}{323.646}{349.687}{N}%
\NativeArrowAt{321.058}{342.91}{11.955}{13.211}{\longrightarrow}%
\TextAt{18}{337.598}{342.91}{R}%
\TextAt{3}{346.232}{347.846}{\Character{3}}%
\TextAt{17}{346.232}{339.954}{\Character{43}}%
\NativeArrowAt{356.637}{342.91}{11.955}{19.26}{\longrightarrow}%
\TextAt{9}{379.219}{342.91}{1}%
\TextAt{8}{79.226}{319.388}{(see}%
\TextAt{8}{104.384}{319.388}{§1)}%
\TextAt{8}{126.063}{319.388}{the}%
\TextAt{8}{147.909}{319.388}{group}%
\TextAt{13}{182.775}{319.388}{C}%
\TextAt{17}{192.008}{323.727}{0}%
\TextAt{21}{191.15}{316.368}{K}%
\TextAt{8}{204.94}{319.388}{is}%
\TextAt{8}{218.394}{319.388}{compact}%
\TextAt{8}{266.912}{319.388}{by}%
\TextAt{8}{284.861}{319.388}{(1.6),}%
\TextAt{8}{318.2}{319.388}{and}%
\TextAt{8}{342.653}{319.388}{we}%
\TextAt{8}{361.57}{319.388}{obtain}%
\TextAt{8}{399.687}{319.388}{a}%
\TextAt{8}{411.133}{319.388}{splitting,}%
\TextAt{8}{462.74}{319.388}{if}%
\TextAt{8}{475.153}{319.388}{we}%
\TextAt{8}{494.071}{319.388}{identify}%
\TextAt{18}{79.226}{302.342}{R}%
\TextAt{3}{87.86}{306.681}{\Character{3}}%
\TextAt{17}{87.86}{299.387}{\Character{43}}%
\TextAt{8}{99.435}{302.342}{with}%
\TextAt{8}{126.689}{302.342}{the}%
\TextAt{8}{147.438}{302.342}{group}%
\TextAt{8}{181.207}{302.342}{of}%
\TextAt{8}{195.13}{302.342}{positive}%
\TextAt{8}{239.036}{302.342}{real}%
\TextAt{8}{262.392}{302.342}{numbers}%
\TextAt{8}{310.196}{302.342}{in}%
\TextAt{8}{324.442}{302.342}{any}%
\TextAt{8}{347.152}{302.342}{infinite}%
\TextAt{8}{387.412}{302.342}{completion}%
\TextAt{13}{447.841}{302.342}{K}%
\TextAt{27}{457.8}{300.549}{p}%
\TextAt{8}{462.287}{302.342}{.}%
\TextAt{8}{472.514}{302.342}{Thus}%
\TextAt{13}{502.756}{302.342}{C}%
\TextAt{21}{511.131}{300.549}{K}%
\TextAt{9}{523.665}{302.342}{\Character{61}}%
\TextAt{13}{79.226}{285.296}{C}%
\TextAt{17}{88.46}{289.635}{0}%
\TextAt{21}{87.601}{282.276}{K}%
\TextAt{16}{99.182}{285.296}{\Character{2}}%
\TextAt{18}{111.858}{285.296}{R}%
\TextAt{3}{120.492}{289.635}{\Character{3}}%
\TextAt{17}{120.492}{282.341}{\Character{43}}%
\TextAt{8}{127.577}{285.296}{.}%
\TextAt{8}{139.199}{285.296}{Now,}%
\TextAt{9}{170.433}{285.296}{\Character{91}}%
\TextAt{18}{173.685}{285.296}{R}%
\TextAt{3}{182.319}{289.635}{\Character{3}}%
\TextAt{17}{182.319}{282.341}{\Character{43}}%
\TextAt{13}{189.404}{285.296}{\Character{59}}%
\TextAt{20}{197.386}{288.318}{\Character{101}}%
\TextAt{13}{194.648}{285.296}{K}%
\TextAt{16}{210.589}{285.296}{\Character{106}}%
\TextAt{13}{219.034}{285.296}{K}%
\TextAt{9}{229.851}{285.296}{\Character{93}}%
\TextAt{9}{238.208}{285.296}{\Character{61}}%
\TextAt{9}{252.434}{285.296}{1}%
\TextAt{8}{258.308}{285.296}{,}%
\TextAt{8}{266.779}{285.296}{for}%
\TextAt{8}{285.722}{285.296}{if}%
\TextAt{13}{297.504}{285.296}{x}%
\NativeMathAt{309.28}{285.296}{11.955}{\in}%
\TextAt{18}{322.373}{285.296}{R}%
\TextAt{3}{331.008}{289.635}{\Character{3}}%
\TextAt{17}{331.008}{282.341}{\Character{43}}%
\TextAt{8}{338.092}{285.296}{,}%
\TextAt{8}{346.564}{285.296}{then}%
\TextAt{9}{374.282}{285.296}{\Character{91}}%
\TextAt{13}{377.533}{285.296}{x\Character{59}}%
\TextAt{20}{392.168}{288.318}{\Character{101}}%
\TextAt{13}{389.43}{285.296}{K}%
\TextAt{16}{405.371}{285.296}{\Character{106}}%
\TextAt{13}{413.815}{285.296}{K}%
\TextAt{9}{424.633}{285.296}{\Character{93}}%
\TextAt{16}{427.885}{285.296}{\Character{106}}%
\TextAt{21}{431.206}{283.503}{L}%
\TextAt{9}{442.586}{285.296}{\Character{61}}%
\TextAt{9}{456.801}{285.296}{\Character{91}}%
\TextAt{13}{460.066}{285.296}{x\Character{59}}%
\TextAt{13}{471.938}{285.296}{L}%
\TextAt{16}{485.051}{285.296}{\Character{106}}%
\TextAt{13}{493.495}{285.296}{K}%
\TextAt{9}{504.313}{285.296}{\Character{93}}%
\TextAt{9}{512.67}{285.296}{\Character{61}}%
\TextAt{9}{526.896}{285.296}{1}%
\TextAt{8}{79.226}{268.25}{for}%
\TextAt{8}{97.38}{268.25}{every}%
\TextAt{8}{128.541}{268.25}{finite}%
\TextAt{8}{158.722}{268.25}{subextension}%
\TextAt{13}{228.702}{268.25}{L}%
\TextAt{16}{240.446}{268.25}{\Character{106}}%
\TextAt{13}{247.546}{268.25}{K}%
\TextAt{8}{262.53}{268.25}{of}%
\TextAt{20}{278.867}{271.272}{\Character{101}}%
\TextAt{13}{276.129}{268.25}{K}%
\TextAt{16}{290.726}{268.25}{\Character{106}}%
\TextAt{13}{297.826}{268.25}{K}%
\TextAt{8}{308.644}{268.25}{,}%
\TextAt{8}{316.13}{268.25}{because}%
\TextAt{8}{359.7}{268.25}{we}%
\TextAt{8}{377.197}{268.25}{may}%
\TextAt{8}{402.835}{268.25}{always}%
\TextAt{8}{440.571}{268.25}{write}%
\TextAt{13}{470.752}{268.25}{x}%
\TextAt{9}{481.183}{268.25}{\Character{61}}%
\TextAt{13}{494.067}{268.25}{y}%
\TextAt{21}{500.204}{272.589}{n}%
\TextAt{8}{510.007}{268.25}{with}%
\TextAt{13}{79.226}{251.204}{y}%
\NativeMathAt{91.14}{251.204}{11.955}{\in}%
\TextAt{18}{104.887}{251.204}{R}%
\TextAt{3}{113.521}{255.543}{\Character{3}}%
\TextAt{17}{113.521}{248.249}{\Character{43}}%
\TextAt{8}{125.945}{251.204}{and}%
\TextAt{13}{150.148}{251.204}{n}%
\TextAt{9}{162.913}{251.204}{\Character{61}}%
\TextAt{9}{177.785}{251.204}{\Character{91}}%
\TextAt{13}{181.046}{251.204}{L}%
\TextAt{9}{194.788}{251.204}{\Character{58}}%
\TextAt{13}{203.817}{251.204}{K}%
\TextAt{9}{214.634}{251.204}{\Character{93}}%
\TextAt{8}{217.886}{251.204}{.}%
\TextAt{8}{230.658}{251.204}{Therefore}%
\TextAt{9}{285.097}{251.204}{\Character{91}}%
\TextAt{13}{288.348}{251.204}{C}%
\TextAt{21}{296.723}{249.411}{K}%
\TextAt{13}{304.926}{251.204}{\Character{59}}%
\TextAt{20}{312.908}{254.226}{\Character{101}}%
\TextAt{13}{310.17}{251.204}{K}%
\TextAt{16}{326.765}{251.204}{\Character{106}}%
\TextAt{13}{335.863}{251.204}{K}%
\TextAt{9}{346.68}{251.204}{\Character{93}}%
\TextAt{9}{355.694}{251.204}{\Character{61}}%
\TextAt{9}{370.567}{251.204}{\Character{91}}%
\TextAt{13}{373.842}{251.204}{C}%
\TextAt{17}{383.076}{255.543}{0}%
\TextAt{21}{382.217}{248.184}{K}%
\TextAt{13}{390.42}{251.204}{\Character{59}}%
\TextAt{20}{398.402}{254.226}{\Character{101}}%
\TextAt{13}{395.664}{251.204}{K}%
\TextAt{16}{412.259}{251.204}{\Character{106}}%
\TextAt{13}{421.357}{251.204}{K}%
\TextAt{9}{432.174}{251.204}{\Character{93}}%
\TextAt{8}{440.765}{251.204}{is}%
\TextAt{8}{453.97}{251.204}{a}%
\TextAt{8}{465.167}{251.204}{closed,}%
\TextAt{8}{504.747}{251.204}{dense}%
\TextAt{8}{79.226}{234.158}{subgroup}%
\TextAt{8}{131.022}{234.158}{of}%
\TextAt{13}{145.351}{234.158}{G}%
\TextAt{9}{154.585}{234.158}{\Character{40}}%
\TextAt{20}{161.875}{237.18}{\Character{101}}%
\TextAt{13}{159.137}{234.158}{K}%
\TextAt{16}{174.977}{234.158}{\Character{106}}%
\TextAt{13}{183.321}{234.158}{K}%
\TextAt{9}{194.139}{234.158}{\Character{41}}%
\TextAt{8}{203.587}{234.158}{and}%
\TextAt{8}{227.349}{234.158}{therefore}%
\TextAt{8}{277.447}{234.158}{equal}%
\TextAt{8}{309.338}{234.158}{to}%
\TextAt{13}{324.647}{234.158}{G}%
\TextAt{9}{333.881}{234.158}{\Character{40}}%
\TextAt{20}{341.171}{237.18}{\Character{101}}%
\TextAt{13}{338.433}{234.158}{K}%
\TextAt{16}{354.273}{234.158}{\Character{106}}%
\TextAt{13}{362.617}{234.158}{K}%
\TextAt{9}{373.435}{234.158}{\Character{41}}%
\TextAt{8}{377.987}{234.158}{.}%
\TextAt{8}{389.432}{234.158}{This}%
\TextAt{8}{417.15}{234.158}{proves}%
\TextAt{8}{454.313}{234.158}{the}%
\TextAt{8}{475.468}{234.158}{surjectivity}%
\TextAt{8}{79.226}{217.112}{of}%
\TextAt{13}{93.688}{217.112}{v}%
\TextAt{21}{99.347}{215.319}{K}%
\TextAt{9}{112.798}{217.112}{\Character{61}}%
\TextAt{13}{127.152}{217.112}{d}%
\TextAt{21}{133.235}{215.319}{K}%
\TextAt{16}{144.865}{217.112}{\Character{14}}%
\TextAt{9}{154.271}{217.112}{\Character{91}}%
\TextAt{13}{162.551}{217.112}{\Character{59}}%
\TextAt{20}{170.534}{220.134}{\Character{101}}%
\TextAt{13}{167.795}{217.112}{K}%
\TextAt{16}{183.862}{217.112}{\Character{106}}%
\TextAt{13}{192.432}{217.112}{K}%
\TextAt{9}{203.25}{217.112}{\Character{93}}%
\TextAt{8}{206.501}{217.112}{.}%
\TextAt{8}{218.344}{217.112}{In}%
\TextAt{8}{234.096}{217.112}{the}%
\TextAt{8}{255.384}{217.112}{definition}%
\TextAt{8}{308.545}{217.112}{of}%
\TextAt{8}{323.006}{217.112}{a}%
\TextAt{8}{333.893}{217.112}{henselian}%
\TextAt{8}{385.81}{217.112}{valuation}%
\TextAt{8}{438.014}{217.112}{given}%
\TextAt{8}{469.715}{217.112}{in}%
\TextAt{8}{484.499}{217.112}{chap.}%
\TextAt{8}{516.522}{217.112}{IV,}%
\TextAt{8}{79.226}{200.066}{(4.6),}%
\TextAt{8}{110.484}{200.066}{condition}%
\TextAt{8}{161.891}{200.066}{(i)}%
\TextAt{8}{178.175}{200.066}{is}%
\TextAt{8}{189.963}{200.066}{satisfied}%
\TextAt{8}{234.987}{200.066}{because}%
\TextAt{13}{278.313}{200.066}{v}%
\TextAt{21}{283.972}{198.273}{K}%
\TextAt{9}{292.175}{200.066}{\Character{40}}%
\TextAt{13}{296.727}{200.066}{C}%
\TextAt{21}{305.102}{198.273}{K}%
\TextAt{9}{313.305}{200.066}{\Character{41}}%
\TextAt{9}{321.207}{200.066}{\Character{61}}%
\TextAt{20}{334.351}{203.155}{\Character{98}}%
\TextAt{18}{333.687}{200.066}{Z}%
\TextAt{8}{341.657}{200.066}{,}%
\TextAt{8}{348.837}{200.066}{and}%
\TextAt{8}{371.624}{200.066}{condition}%
\TextAt{8}{423.031}{200.066}{(ii)}%
\TextAt{8}{442.567}{200.066}{follows}%
\TextAt{8}{481.027}{200.066}{from}%
\TextAt{8}{508.692}{200.066}{(5.2)}%
\TextAt{8}{79.226}{183.02}{because}%
\TextAt{8}{122.528}{183.02}{for}%
\TextAt{8}{140.413}{183.02}{every}%
\TextAt{8}{171.305}{183.02}{finite}%
\TextAt{8}{201.217}{183.02}{extension}%
\TextAt{13}{252.983}{183.02}{L}%
\TextAt{16}{264.268}{183.02}{\Character{106}}%
\TextAt{13}{270.91}{183.02}{K}%
\TextAt{8}{285.625}{183.02}{we}%
\TextAt{8}{302.852}{183.02}{have}%
\TextAt{8}{329.847}{183.02}{the}%
\TextAt{8}{350.003}{183.02}{identity}%
\TextAt{13}{163.472}{154.019}{v}%
\TextAt{21}{169.131}{152.226}{K}%
\TextAt{9}{177.334}{154.019}{\Character{40}}%
\TextAt{13}{181.886}{154.019}{N}%
\TextAt{21}{191.26}{152.16}{L}%
\TextAt{3}{197.019}{152.16}{\Character{106}}%
\TextAt{21}{199.371}{152.16}{K}%
\TextAt{13}{207.574}{154.019}{C}%
\TextAt{21}{215.949}{152.226}{L}%
\TextAt{9}{222.206}{154.019}{\Character{41}}%
\TextAt{9}{230.061}{154.019}{\Character{61}}%
\TextAt{13}{242.505}{154.019}{v}%
\TextAt{21}{248.164}{152.226}{K}%
\TextAt{9}{256.367}{154.019}{\Character{40}}%
\TextAt{13}{260.919}{154.019}{N}%
\TextAt{21}{270.293}{152.16}{L}%
\TextAt{3}{276.052}{152.16}{\Character{106}}%
\TextAt{21}{278.404}{152.16}{K}%
\TextAt{13}{286.607}{154.019}{I}%
\TextAt{21}{291.77}{152.226}{L}%
\TextAt{9}{298.027}{154.019}{\Character{41}}%
\TextAt{9}{305.882}{154.019}{\Character{61}}%
\TextAt{13}{318.326}{154.019}{d}%
\TextAt{21}{324.408}{152.226}{K}%
\TextAt{9}{332.611}{154.019}{\Character{91}}%
\TextAt{13}{335.863}{154.019}{N}%
\TextAt{21}{345.237}{152.16}{L}%
\TextAt{3}{350.996}{152.16}{\Character{106}}%
\TextAt{21}{353.348}{152.16}{K}%
\TextAt{13}{361.551}{154.019}{I}%
\TextAt{21}{366.713}{152.226}{L}%
\TextAt{13}{372.971}{154.019}{\Character{59}}%
\TextAt{20}{380.953}{157.041}{\Character{101}}%
\TextAt{13}{378.215}{154.019}{K}%
\TextAt{16}{392.353}{154.019}{\Character{106}}%
\TextAt{13}{398.995}{154.019}{K}%
\TextAt{9}{409.813}{154.019}{\Character{93}}%
\TextAt{9}{230.079}{133.985}{\Character{61}}%
\TextAt{13}{242.505}{133.985}{f}%
\TextAt{21}{248.275}{132.125}{L}%
\TextAt{3}{254.034}{132.125}{\Character{106}}%
\TextAt{21}{256.386}{132.125}{K}%
\TextAt{13}{264.589}{133.985}{d}%
\TextAt{21}{270.672}{132.191}{L}%
\TextAt{9}{276.929}{133.985}{\Character{91}}%
\TextAt{13}{280.18}{133.985}{I}%
\TextAt{21}{285.343}{132.191}{L}%
\TextAt{13}{291.601}{133.985}{\Character{59}}%
\TextAt{20}{297.831}{137.007}{\Character{101}}%
\TextAt{13}{296.845}{133.985}{L}%
\TextAt{16}{308.13}{133.985}{\Character{106}}%
\TextAt{13}{314.772}{133.985}{L}%
\TextAt{9}{322.736}{133.985}{\Character{93}}%
\TextAt{9}{329.287}{133.985}{\Character{61}}%
\TextAt{13}{341.734}{133.985}{f}%
\TextAt{21}{347.504}{132.125}{L}%
\TextAt{3}{353.263}{132.125}{\Character{106}}%
\TextAt{21}{355.616}{132.125}{K}%
\TextAt{13}{363.818}{133.985}{v}%
\TextAt{21}{369.478}{132.191}{L}%
\TextAt{9}{375.735}{133.985}{\Character{40}}%
\TextAt{13}{380.287}{133.985}{C}%
\TextAt{21}{388.662}{132.191}{L}%
\TextAt{9}{394.919}{133.985}{\Character{41}}%
\TextAt{9}{402.774}{133.985}{\Character{61}}%
\TextAt{13}{415.218}{133.985}{f}%
\TextAt{21}{420.988}{132.125}{L}%
\TextAt{3}{426.747}{132.125}{\Character{106}}%
\TextAt{21}{429.099}{132.125}{K}%
\TextAt{20}{437.966}{137.073}{\Character{98}}%
\TextAt{18}{437.302}{133.985}{Z}%
\TextAt{13}{445.272}{133.985}{\Character{58}}%
\Vector{q 0 G 0 g 0.399 w 452.334 -669.017 m 452.334 -660.947 l S Q}%
\Vector{q 0 G 0 g 0.399 w 452.534 -661.146 m 459.707 -661.146 l S Q}%
\Vector{q 0 G 0 g 0.399 w 452.534 -668.817 m 459.707 -668.817 l S Q}%
\Vector{q 0 G 0 g 0.399 w 459.906 -669.017 m 459.906 -660.947 l S Q}%
\TextAt{8}{297.211}{43.033}{355}%
\Destination{78.23}{809.815}{nav.page.355}%
\Destination{79.36}{510.555}{nav.print.390}%
\Destination{79.36}{561.555}{nav.number.6.5.4}%
\Hotlink{109.563}{317.155}{6.858}{11.047}{nav.section.6.1}%
\Hotlink{516.022}{214.955}{13.995}{11.047}{nav.chapter.4}%
\Hotlink{212.656}{391.855}{25.078}{11.047}{nav.number.6.3.8}%
\Hotlink{284.361}{317.155}{25.078}{11.047}{nav.number.6.1.6}%
\Hotlink{508.192}{197.855}{25.078}{11.047}{nav.number.6.5.2}%
}

% Page 368 (body page 356)
\Page{%
\TextAt{2}{79.226}{796.943}{VI.}%
\TextAt{2}{98.012}{796.943}{Global}%
\TextAt{2}{133.248}{796.943}{Class}%
\TextAt{2}{161.841}{796.943}{Field}%
\TextAt{2}{189.572}{796.943}{Theory}%
\Vector{q 0 G 0 g 0.299 w 7.226 18.729 m 460.769 18.729 l S Q}%
\TextAt{8}{97.159}{761.124}{In}%
\TextAt{8}{112.589}{761.124}{view}%
\TextAt{8}{140.38}{761.124}{of}%
\TextAt{8}{154.519}{761.124}{the}%
\TextAt{8}{175.484}{761.124}{fact}%
\TextAt{8}{199.379}{761.124}{that}%
\TextAt{8}{225.556}{761.124}{the}%
\TextAt{8}{246.522}{761.124}{idèle}%
\TextAt{8}{274.636}{761.124}{class}%
\TextAt{8}{302.882}{761.124}{group}%
\TextAt{13}{336.866}{761.124}{C}%
\TextAt{21}{345.241}{759.33}{K}%
\TextAt{8}{358.15}{761.124}{satisfies}%
\TextAt{8}{402.069}{761.124}{the}%
\TextAt{8}{423.035}{761.124}{class}%
\TextAt{8}{451.28}{761.124}{field}%
\TextAt{8}{477.446}{761.124}{axiom,}%
\TextAt{8}{516.51}{761.124}{the}%
\TextAt{8}{79.226}{744.078}{pair}%
\TextAt{9}{102.711}{744.078}{\Character{40}}%
\TextAt{13}{107.263}{744.078}{d}%
\TextAt{25}{113.346}{742.284}{Q}%
\TextAt{9}{123.364}{744.078}{\Character{58}}%
\TextAt{13}{129.936}{744.078}{G}%
\TextAt{25}{139.17}{742.284}{Q}%
\NativeArrowAt{149.188}{744.078}{11.955}{19.26}{\longrightarrow}%
\TextAt{20}{172.434}{747.166}{\Character{98}}%
\TextAt{18}{171.77}{744.078}{Z}%
\TextAt{13}{179.74}{744.078}{\Character{59}}%
\TextAt{13}{188.288}{744.078}{v}%
\TextAt{25}{193.96}{742.284}{Q}%
\TextAt{9}{203.978}{744.078}{\Character{58}}%
\TextAt{13}{210.551}{744.078}{C}%
\TextAt{25}{218.926}{742.284}{Q}%
\NativeArrowAt{228.944}{744.078}{11.955}{19.26}{\longrightarrow}%
\TextAt{20}{252.19}{747.166}{\Character{98}}%
\TextAt{18}{251.526}{744.078}{Z}%
\TextAt{9}{259.496}{744.078}{\Character{41}}%
\TextAt{8}{267.365}{744.078}{constitutes}%
\TextAt{8}{326.093}{744.078}{a}%
\TextAt{8}{335.268}{744.078}{class}%
\TextAt{8}{362.124}{744.078}{field}%
\TextAt{8}{386.9}{744.078}{theory,}%
\TextAt{8}{425.457}{744.078}{the}%
\TextAt{8}{445.033}{744.078}{“global}%
\TextAt{8}{484.454}{744.078}{class}%
\TextAt{8}{511.31}{744.078}{field}%
\TextAt{8}{79.226}{727.032}{theory”.}%
\TextAt{8}{124.373}{727.032}{The}%
\TextAt{8}{148.27}{727.032}{above}%
\TextAt{8}{181.289}{727.032}{homomorphism}%
\TextAt{13}{263.803}{727.032}{v}%
\TextAt{21}{269.462}{725.238}{K}%
\TextAt{9}{280.986}{727.032}{\Character{61}}%
\TextAt{13}{293.411}{727.032}{d}%
\TextAt{21}{299.494}{725.238}{K}%
\TextAt{16}{310.033}{727.032}{\Character{14}}%
\TextAt{9}{318.348}{727.032}{\Character{91}}%
\TextAt{13}{325.34}{727.032}{\Character{59}}%
\TextAt{20}{333.323}{730.054}{\Character{101}}%
\TextAt{13}{330.585}{727.032}{K}%
\TextAt{16}{344.723}{727.032}{\Character{106}}%
\TextAt{13}{351.365}{727.032}{K}%
\TextAt{9}{362.183}{727.032}{\Character{93}}%
\TextAt{9}{368.734}{727.032}{\Character{58}}%
\TextAt{13}{375.328}{727.032}{C}%
\TextAt{21}{383.703}{725.238}{K}%
\NativeArrowAt{395.226}{727.032}{11.955}{19.26}{\longrightarrow}%
\TextAt{20}{418.472}{730.12}{\Character{98}}%
\TextAt{18}{417.808}{727.032}{Z}%
\TextAt{8}{425.778}{727.032}{,}%
\TextAt{8}{432.802}{727.032}{for}%
\TextAt{8}{450.531}{727.032}{finite}%
\TextAt{8}{480.286}{727.032}{extensions}%
\TextAt{13}{79.226}{709.986}{K}%
\TextAt{16}{93.365}{709.986}{\Character{106}}%
\TextAt{18}{100.006}{709.986}{Q}%
\TextAt{8}{109.305}{709.986}{,}%
\TextAt{8}{116.454}{709.986}{satisfies}%
\TextAt{8}{159.564}{709.986}{the}%
\TextAt{8}{179.721}{709.986}{formula}%
\TextAt{13}{189.222}{681.471}{v}%
\TextAt{21}{194.882}{679.678}{K}%
\TextAt{9}{206.405}{681.471}{\Character{61}}%
\TextAt{9}{224.086}{689.559}{1}%
\Vector{q 0 G 0 g 0.478 w 148.026 -89.54 m 161.999 -89.54 l S Q}%
\TextAt{13}{220.026}{673.271}{f}%
\TextAt{21}{225.796}{671.477}{K}%
\TextAt{13}{235.194}{681.471}{d}%
\TextAt{25}{241.277}{679.678}{Q}%
\TextAt{16}{250.631}{681.471}{\Character{14}}%
\TextAt{9}{259.265}{681.471}{\Character{91}}%
\TextAt{13}{266.414}{681.471}{\Character{59}}%
\TextAt{20}{272.987}{684.56}{\Character{101}}%
\TextAt{18}{271.658}{681.471}{Q}%
\TextAt{16}{284.278}{681.471}{\Character{106}}%
\TextAt{18}{290.919}{681.471}{Q}%
\TextAt{9}{300.218}{681.471}{\Character{93}}%
\TextAt{16}{306.126}{681.471}{\Character{14}}%
\TextAt{13}{314.76}{681.471}{N}%
\TextAt{21}{324.134}{679.611}{K}%
\TextAt{3}{331.839}{679.611}{\Character{106}}%
\TextAt{25}{334.191}{679.611}{Q}%
\TextAt{9}{344.209}{681.471}{\Character{61}}%
\TextAt{9}{361.89}{689.559}{1}%
\Vector{q 0 G 0 g 0.478 w 285.83 -89.54 m 299.803 -89.54 l S Q}%
\TextAt{13}{357.83}{673.271}{f}%
\TextAt{21}{363.6}{671.477}{K}%
\TextAt{13}{372.998}{681.471}{v}%
\TextAt{25}{378.657}{679.678}{Q}%
\TextAt{16}{388.011}{681.471}{\Character{14}}%
\TextAt{13}{396.646}{681.471}{N}%
\TextAt{21}{406.019}{679.611}{K}%
\TextAt{3}{413.724}{679.611}{\Character{106}}%
\TextAt{25}{416.076}{679.611}{Q}%
\TextAt{8}{79.226}{652.796}{and}%
\TextAt{8}{101.981}{652.796}{is}%
\TextAt{8}{113.736}{652.796}{therefore}%
\TextAt{8}{162.828}{652.796}{precisely}%
\TextAt{8}{210.676}{652.796}{the}%
\TextAt{8}{230.824}{652.796}{induced}%
\TextAt{8}{274.381}{652.796}{homomorphism}%
\TextAt{8}{357.042}{652.796}{in}%
\TextAt{8}{370.687}{652.796}{the}%
\TextAt{8}{390.835}{652.796}{sense}%
\TextAt{8}{420.858}{652.796}{of}%
\TextAt{8}{434.18}{652.796}{the}%
\TextAt{8}{454.328}{652.796}{abstract}%
\TextAt{8}{499.917}{652.796}{theory}%
\TextAt{8}{79.226}{635.75}{in}%
\TextAt{8}{92.879}{635.75}{chap.}%
\TextAt{8}{123.771}{635.75}{IV,}%
\TextAt{8}{143.916}{635.75}{(4.7).}%
\TextAt{8}{97.159}{615.715}{As}%
\TextAt{8}{113.815}{615.715}{the}%
\TextAt{8}{133.34}{615.715}{main}%
\TextAt{8}{161.975}{615.715}{result}%
\TextAt{8}{193.922}{615.715}{of}%
\TextAt{8}{206.62}{615.715}{global}%
\TextAt{8}{240.468}{615.715}{class}%
\TextAt{8}{267.274}{615.715}{field}%
\TextAt{8}{291.999}{615.715}{theory}%
\TextAt{8}{328.118}{615.715}{we}%
\TextAt{8}{344.715}{615.715}{now}%
\TextAt{8}{368.472}{615.715}{obtain}%
\TextAt{8}{404.268}{615.715}{the}%
\TextAt{8}{423.793}{615.715}{Artin}%
\TextAt{8}{454.7}{615.715}{reciprocity}%
\TextAt{8}{512.278}{615.715}{law:}%
\TextAt{12}{79.226}{590.622}{Theorem}%
\TextAt{12}{137.013}{590.622}{(5.5).}%
\TextAt{8}{173.001}{590.622}{For}%
\TextAt{8}{194.642}{590.622}{every}%
\TextAt{8}{226.207}{590.622}{Galois}%
\TextAt{8}{262.792}{590.622}{extension}%
\TextAt{13}{315.23}{590.622}{L}%
\TextAt{16}{327.661}{590.622}{\Character{106}}%
\TextAt{13}{335.448}{590.622}{K}%
\TextAt{8}{350.835}{590.622}{of}%
\TextAt{8}{364.837}{590.622}{finite}%
\TextAt{8}{395.421}{590.622}{algebraic}%
\TextAt{8}{445.527}{590.622}{number}%
\TextAt{8}{488.796}{590.622}{fields}%
\TextAt{8}{519.439}{590.622}{we}%
\TextAt{8}{79.226}{573.576}{have}%
\TextAt{8}{106.221}{573.576}{a}%
\TextAt{8}{115.976}{573.576}{canonical}%
\TextAt{8}{167.36}{573.576}{isomorphism}%
\TextAt{13}{219.497}{547.793}{r}%
\TextAt{21}{224.772}{545.933}{L}%
\TextAt{3}{230.531}{545.933}{\Character{106}}%
\TextAt{21}{232.884}{545.933}{K}%
\TextAt{9}{244.407}{547.793}{\Character{58}}%
\TextAt{13}{250.98}{547.793}{G}%
\TextAt{9}{260.213}{547.793}{\Character{40}}%
\TextAt{13}{264.766}{547.793}{L}%
\TextAt{16}{276.051}{547.793}{\Character{106}}%
\TextAt{13}{282.693}{547.793}{K}%
\TextAt{9}{293.51}{547.793}{\Character{41}}%
\TextAt{21}{298.063}{552.729}{ab}%
\NativeMathAt{312.59}{554.57}{7.97}{\sim}%
\NativeArrowAt{310.002}{547.793}{11.955}{13.175}{\longrightarrow}%
\TextAt{13}{326.496}{547.793}{C}%
\TextAt{21}{334.871}{546}{K}%
\TextAt{9}{343.074}{547.793}{\Character{47}}%
\TextAt{13}{348.927}{547.793}{N}%
\TextAt{21}{358.301}{545.933}{L}%
\TextAt{3}{364.06}{545.933}{\Character{106}}%
\TextAt{21}{366.412}{545.933}{K}%
\TextAt{13}{374.615}{547.793}{C}%
\TextAt{21}{382.99}{546}{L}%
\TextAt{13}{389.247}{547.793}{\Character{58}}%
\TextAt{8}{97.159}{513.963}{The}%
\TextAt{8}{121.213}{513.963}{inverse}%
\TextAt{8}{159.971}{513.963}{map}%
\TextAt{8}{185.986}{513.963}{of}%
\TextAt{13}{199.316}{513.963}{r}%
\TextAt{21}{204.591}{512.103}{L}%
\TextAt{3}{210.35}{512.103}{\Character{106}}%
\TextAt{21}{212.702}{512.103}{K}%
\TextAt{8}{224.803}{513.963}{yields}%
\TextAt{8}{257.703}{513.963}{a}%
\TextAt{8}{267.459}{513.963}{surjective}%
\TextAt{8}{319.87}{513.963}{homomorphism}%
\TextAt{9}{228.465}{488.179}{\Character{40}}%
\TextAt{13}{236.914}{488.179}{\Character{59}}%
\TextAt{13}{242.138}{488.179}{L}%
\TextAt{16}{253.444}{488.179}{\Character{106}}%
\TextAt{13}{260.085}{488.179}{K}%
\TextAt{9}{270.903}{488.179}{\Character{41}}%
\TextAt{9}{278.758}{488.179}{\Character{58}}%
\TextAt{13}{285.349}{488.179}{C}%
\TextAt{21}{293.724}{486.386}{K}%
\NativeArrowAt{305.247}{488.179}{11.955}{19.26}{\longrightarrow}%
\TextAt{13}{327.829}{488.179}{G}%
\TextAt{9}{337.063}{488.179}{\Character{40}}%
\TextAt{13}{341.615}{488.179}{L}%
\TextAt{16}{352.901}{488.179}{\Character{106}}%
\TextAt{13}{359.542}{488.179}{K}%
\TextAt{9}{370.36}{488.179}{\Character{41}}%
\TextAt{21}{374.912}{493.116}{ab}%
\TextAt{8}{79.226}{462.396}{with}%
\TextAt{8}{106.507}{462.396}{kernel}%
\TextAt{13}{141.594}{462.396}{N}%
\TextAt{21}{150.968}{460.536}{L}%
\TextAt{3}{156.727}{460.536}{\Character{106}}%
\TextAt{21}{159.08}{460.536}{K}%
\TextAt{13}{167.282}{462.396}{C}%
\TextAt{21}{175.657}{460.603}{L}%
\TextAt{8}{181.914}{462.396}{.}%
\TextAt{8}{192.225}{462.396}{The}%
\TextAt{8}{216.899}{462.396}{map}%
\TextAt{9}{243.534}{462.396}{\Character{40}}%
\TextAt{13}{252.605}{462.396}{\Character{59}}%
\TextAt{13}{257.829}{462.396}{L}%
\TextAt{16}{270.192}{462.396}{\Character{106}}%
\TextAt{13}{277.892}{462.396}{K}%
\TextAt{9}{288.709}{462.396}{\Character{41}}%
\TextAt{8}{297.78}{462.396}{is}%
\TextAt{8}{310.164}{462.396}{called}%
\TextAt{8}{343.949}{462.396}{the}%
\TextAt{12}{364.726}{462.396}{global}%
\TextAt{12}{404.88}{462.396}{norm}%
\TextAt{12}{440.993}{462.396}{residue}%
\TextAt{12}{487.95}{462.396}{symbol}%
\TextAt{8}{529.518}{462.396}{.}%
\TextAt{8}{79.226}{445.35}{We}%
\TextAt{8}{99.211}{445.35}{view}%
\TextAt{8}{126.034}{445.35}{it}%
\TextAt{8}{137.578}{445.35}{also}%
\TextAt{8}{160.899}{445.35}{as}%
\TextAt{8}{175.109}{445.35}{a}%
\TextAt{8}{184.705}{445.35}{homomorphism}%
\TextAt{13}{267.215}{445.35}{I}%
\TextAt{21}{272.378}{443.557}{K}%
\NativeArrowAt{283.901}{445.35}{11.955}{19.26}{\longrightarrow}%
\TextAt{13}{306.483}{445.35}{G}%
\TextAt{9}{315.717}{445.35}{\Character{40}}%
\TextAt{13}{320.269}{445.35}{L}%
\TextAt{16}{331.555}{445.35}{\Character{106}}%
\TextAt{13}{338.196}{445.35}{K}%
\TextAt{9}{349.014}{445.35}{\Character{41}}%
\TextAt{21}{353.566}{449.689}{ab}%
\TextAt{8}{362.185}{445.35}{.}%
\TextAt{8}{370.58}{445.35}{For}%
\TextAt{8}{391.39}{445.35}{every}%
\TextAt{8}{422.122}{445.35}{place}%
\TextAt{14}{451.874}{445.35}{p}%
\TextAt{8}{461.595}{445.35}{of}%
\TextAt{13}{474.765}{445.35}{K}%
\TextAt{8}{485.583}{445.35}{,}%
\TextAt{8}{492.604}{445.35}{we}%
\TextAt{8}{509.672}{445.35}{have}%
\TextAt{8}{79.226}{428.304}{on}%
\TextAt{8}{95.503}{428.304}{the}%
\TextAt{8}{115.676}{428.304}{one}%
\TextAt{8}{137.153}{428.304}{hand}%
\TextAt{8}{166.437}{428.304}{the}%
\TextAt{8}{186.611}{428.304}{embedding}%
\TextAt{13}{245.806}{428.304}{G}%
\TextAt{9}{255.04}{428.304}{\Character{40}}%
\TextAt{13}{259.592}{428.304}{L}%
\TextAt{27}{267.557}{426.511}{p}%
\TextAt{16}{275.394}{428.304}{\Character{106}}%
\TextAt{13}{282.066}{428.304}{K}%
\TextAt{27}{292.025}{426.511}{p}%
\TextAt{9}{296.512}{428.304}{\Character{41}}%
\NativeArrowAt{304.414}{428.304}{11.955}{13.215}{\hookrightarrow}%
\TextAt{13}{320.979}{428.304}{G}%
\TextAt{9}{330.213}{428.304}{\Character{40}}%
\TextAt{13}{334.765}{428.304}{L}%
\TextAt{16}{346.08}{428.304}{\Character{106}}%
\TextAt{13}{352.752}{428.304}{K}%
\TextAt{9}{363.569}{428.304}{\Character{41}}%
\TextAt{8}{368.122}{428.304}{,}%
\TextAt{8}{375.293}{428.304}{and}%
\TextAt{8}{398.073}{428.304}{on}%
\TextAt{8}{414.349}{428.304}{the}%
\TextAt{8}{434.523}{428.304}{other}%
\TextAt{8}{465.11}{428.304}{the}%
\TextAt{8}{485.284}{428.304}{canonical}%
\TextAt{8}{79.226}{411.258}{injection}%
\TextAt{16}{263.812}{394.212}{\Character{104}}%
\TextAt{16}{272.348}{394.212}{\Character{105}}%
\TextAt{9}{280.329}{394.212}{\Character{58}}%
\TextAt{13}{286.901}{394.212}{K}%
\TextAt{3}{297.719}{399.148}{\Character{3}}%
\TextAt{27}{296.86}{391.257}{p}%
\NativeArrowAt{305.772}{394.212}{11.955}{19.26}{\longrightarrow}%
\TextAt{13}{328.354}{394.212}{C}%
\TextAt{21}{336.729}{392.419}{K}%
\TextAt{13}{344.932}{394.212}{\Character{59}}%
\TextAt{8}{79.226}{372.07}{which}%
\TextAt{8}{113.826}{372.07}{sends}%
\TextAt{13}{146.275}{372.07}{a}%
\TextAt{27}{152.42}{370.276}{p}%
\NativeMathAt{162.126}{372.07}{11.955}{\in}%
\TextAt{13}{175.315}{372.07}{K}%
\TextAt{3}{186.132}{376.408}{\Character{3}}%
\TextAt{27}{185.274}{369.114}{p}%
\TextAt{8}{195.876}{372.07}{to}%
\TextAt{8}{211.3}{372.07}{the}%
\TextAt{8}{232.57}{372.07}{class}%
\TextAt{8}{261.121}{372.07}{of}%
\TextAt{8}{275.565}{372.07}{the}%
\TextAt{8}{296.836}{372.07}{idèle}%
\TextAt{16}{325.255}{372.07}{\Character{104}}%
\TextAt{13}{329.904}{372.07}{a}%
\TextAt{27}{336.049}{370.276}{p}%
\TextAt{16}{340.536}{372.07}{\Character{105}}%
\TextAt{9}{350.404}{372.07}{\Character{61}}%
\TextAt{9}{364.714}{372.07}{\Character{40}}%
\TextAt{13}{369.28}{372.07}{\Character{58}}%
\TextAt{13}{374.504}{372.07}{\Character{58}}%
\TextAt{13}{379.741}{372.07}{\Character{58}}%
\TextAt{13}{384.977}{372.07}{\Character{59}}%
\TextAt{9}{390.257}{372.07}{1}%
\TextAt{13}{396.11}{372.07}{\Character{59}}%
\TextAt{9}{401.354}{372.07}{1}%
\TextAt{13}{407.207}{372.07}{\Character{59}}%
\TextAt{9}{412.451}{372.07}{1}%
\TextAt{13}{418.304}{372.07}{\Character{59}}%
\TextAt{13}{423.528}{372.07}{a}%
\TextAt{27}{429.693}{370.276}{p}%
\TextAt{13}{434.18}{372.07}{\Character{59}}%
\TextAt{9}{439.424}{372.07}{1}%
\TextAt{13}{445.277}{372.07}{\Character{59}}%
\TextAt{9}{450.521}{372.07}{1}%
\TextAt{13}{456.374}{372.07}{\Character{59}}%
\TextAt{9}{461.619}{372.07}{1}%
\TextAt{13}{467.472}{372.07}{\Character{59}}%
\TextAt{13}{472.696}{372.07}{\Character{58}}%
\TextAt{13}{477.933}{372.07}{\Character{58}}%
\TextAt{13}{483.157}{372.07}{\Character{58}}%
\TextAt{9}{486.456}{372.07}{\Character{41}}%
\TextAt{8}{491.008}{372.07}{.}%
\TextAt{8}{502.798}{372.07}{These}%
\TextAt{8}{79.226}{355.024}{homomorphisms}%
\TextAt{8}{165.799}{355.024}{express}%
\TextAt{8}{205.853}{355.024}{the}%
\TextAt{8}{225.298}{355.024}{compatibility}%
\TextAt{8}{296.137}{355.024}{of}%
\TextAt{8}{308.755}{355.024}{local}%
\TextAt{8}{335.683}{355.024}{and}%
\TextAt{8}{357.733}{355.024}{global}%
\TextAt{8}{391.5}{355.024}{class}%
\TextAt{8}{418.225}{355.024}{field}%
\TextAt{8}{442.869}{355.024}{theory,}%
\TextAt{8}{481.321}{355.024}{as}%
\TextAt{8}{494.979}{355.024}{follows.}%
\TextAt{12}{79.226}{329.931}{Proposition}%
\TextAt{12}{153.242}{329.931}{(5.6).}%
\TextAt{8}{189.691}{329.931}{If}%
\TextAt{13}{202.516}{329.931}{L}%
\TextAt{16}{215.734}{329.931}{\Character{106}}%
\TextAt{13}{224.307}{329.931}{K}%
\TextAt{8}{240.156}{329.931}{is}%
\TextAt{8}{253.053}{329.931}{an}%
\TextAt{8}{270.446}{329.931}{abelian}%
\TextAt{8}{312.227}{329.931}{extension}%
\TextAt{8}{365.127}{329.931}{and}%
\TextAt{14}{389.023}{329.931}{p}%
\TextAt{8}{400.037}{329.931}{is}%
\TextAt{8}{412.935}{329.931}{a}%
\TextAt{8}{423.824}{329.931}{place}%
\TextAt{8}{454.869}{329.931}{of}%
\TextAt{13}{469.333}{329.931}{K}%
\TextAt{8}{480.151}{329.931}{,}%
\TextAt{8}{488.717}{329.931}{then}%
\TextAt{8}{516.51}{329.931}{the}%
\TextAt{8}{79.226}{312.885}{diagram}%
\TextAt{13}{245.168}{292.952}{K}%
\TextAt{3}{255.986}{297.29}{\Character{3}}%
\TextAt{27}{255.128}{289.996}{p}%
\TextAt{13}{314.576}{292.952}{G}%
\TextAt{9}{323.81}{292.952}{\Character{40}}%
\TextAt{13}{328.362}{292.952}{L}%
\TextAt{27}{336.327}{291.159}{p}%
\TextAt{16}{344.134}{292.952}{\Character{106}}%
\TextAt{13}{350.776}{292.952}{K}%
\TextAt{27}{360.735}{291.159}{p}%
\TextAt{9}{365.222}{292.952}{\Character{41}}%
\TextAt{13}{244.654}{249.244}{C}%
\TextAt{21}{253.029}{247.451}{K}%
\TextAt{13}{318.633}{249.244}{G}%
\TextAt{9}{327.867}{249.244}{\Character{40}}%
\TextAt{13}{332.419}{249.244}{L}%
\TextAt{16}{343.705}{249.244}{\Character{106}}%
\TextAt{13}{350.346}{249.244}{K}%
\TextAt{9}{361.164}{249.244}{\Character{41}}%
\NativeDiagram{%
\path[draw=black,line width=0.478bp] (194.07,-478.059) -- (236.747,-478.059);%
\path[draw=black,line width=0.478bp,line cap=round,line join=round] (234.499,-475.19) .. controls (234.954,-476.911) and (235.966,-477.724) .. (236.986,-478.059) .. controls (235.966,-478.394) and (234.954,-479.207) .. (234.499,-480.928);%
}%
\TextAt{17}{269.52}{300.884}{\Character{40}}%
\TextAt{21}{276.711}{300.884}{\Character{59}L}%
\TextAt{28}{284.822}{299.888}{p}%
\TextAt{3}{288.905}{300.884}{\Character{106}}%
\TextAt{21}{291.257}{300.884}{K}%
\TextAt{28}{298.398}{299.888}{p}%
\TextAt{17}{302.481}{300.884}{\Character{41}}%
\NativeDiagram{%
\path[draw=black,line width=0.478bp] (180.943,-490.089) -- (180.943,-511.529);%
\path[draw=black,line width=0.478bp,line cap=round,line join=round] (183.812,-509.282) .. controls (182.091,-509.736) and (181.278,-510.748) .. (180.943,-511.768) .. controls (180.608,-510.748) and (179.796,-509.736) .. (178.074,-509.282);%
}%
\TextAt{3}{239.644}{270.959}{\Character{104}}%
\TextAt{3}{246.825}{270.959}{\Character{105}}%
\NativeDiagram{%
\path[draw=black,line width=0.478bp] (270.176,-488.927) -- (270.176,-510.732);%
\path[draw=black,line width=0.478bp,line cap=round,line join=round] (273.045,-508.485) .. controls (271.324,-508.939) and (270.511,-509.951) .. (270.176,-510.971) .. controls (269.842,-509.951) and (269.029,-508.939) .. (267.307,-508.485);%
\path[draw=black,line width=0.478bp] (194.584,-521.767) -- (240.804,-521.767);%
\path[draw=black,line width=0.478bp,line cap=round,line join=round] (238.557,-518.898) .. controls (239.011,-520.619) and (240.023,-521.432) .. (241.043,-521.767) .. controls (240.023,-522.102) and (239.011,-522.915) .. (238.557,-524.636);%
}%
\TextAt{17}{275.607}{257.041}{\Character{40}}%
\TextAt{21}{282.798}{257.041}{\Character{59}L}%
\TextAt{3}{290.909}{257.041}{\Character{106}}%
\TextAt{21}{293.261}{257.041}{K}%
\TextAt{17}{300.966}{257.041}{\Character{41}}%
\TextAt{8}{79.226}{227.976}{is}%
\TextAt{8}{90.99}{227.976}{commutative.}%
\TextAt{19}{79.226}{194.836}{Proof.}%
\TextAt{8}{116}{194.836}{We}%
\TextAt{8}{136.559}{194.836}{first}%
\TextAt{8}{161.098}{194.836}{show}%
\TextAt{8}{190.516}{194.836}{that}%
\TextAt{8}{216.299}{194.836}{the}%
\TextAt{8}{236.869}{194.836}{proposition}%
\TextAt{8}{298.816}{194.836}{holds}%
\TextAt{8}{329.859}{194.836}{if}%
\TextAt{13}{340.997}{194.836}{L}%
\TextAt{16}{352.988}{194.836}{\Character{106}}%
\TextAt{13}{360.336}{194.836}{K}%
\TextAt{8}{375.465}{194.836}{is}%
\TextAt{8}{387.643}{194.836}{a}%
\TextAt{8}{397.812}{194.836}{subextension}%
\TextAt{8}{467.937}{194.836}{of}%
\TextAt{20}{484.419}{197.858}{\Character{101}}%
\TextAt{13}{481.681}{194.836}{K}%
\TextAt{16}{496.525}{194.836}{\Character{106}}%
\TextAt{13}{503.872}{194.836}{K}%
\TextAt{8}{514.69}{194.836}{,}%
\TextAt{8}{522.357}{194.836}{or}%
\TextAt{8}{79.226}{177.79}{if}%
\TextAt{13}{90.291}{177.79}{L}%
\TextAt{9}{102.157}{177.79}{\Character{61}}%
\TextAt{13}{115.164}{177.79}{K}%
\TextAt{9}{125.981}{177.79}{\Character{40}}%
\TextAt{13}{130.534}{177.79}{i}%
\TextAt{9}{134.527}{177.79}{\Character{41}}%
\TextAt{8}{139.08}{177.79}{,}%
\TextAt{13}{146.655}{177.79}{i}%
\TextAt{9}{154.55}{177.79}{\Character{61}}%
\TextAt{16}{167.557}{187.182}{\Character{112}}%
\Vector{q 0 G 0 g 0.478 w 105.519 -586.579 m 120.671 -586.579 l S Q}%
\TextAt{16}{177.519}{177.79}{\Character{0}}%
\TextAt{9}{186.818}{177.79}{1}%
\TextAt{8}{192.671}{177.79}{,}%
\TextAt{8}{200.246}{177.79}{and}%
\TextAt{14}{223.35}{177.79}{p}%
\TextAt{16}{233.235}{177.79}{\Character{106}}%
\NativeMathAt{240.444}{177.79}{11.955}{\infty}%
\TextAt{8}{252.413}{177.79}{.}%
\TextAt{8}{261.884}{177.79}{Indeed,}%
\TextAt{8}{303.592}{177.79}{the}%
\TextAt{8}{324.089}{177.79}{two}%
\TextAt{8}{346.547}{177.79}{maps}%
\TextAt{9}{377.517}{177.79}{\Character{91}}%
\TextAt{13}{385.007}{177.79}{\Character{59}}%
\TextAt{20}{392.99}{180.812}{\Character{101}}%
\TextAt{13}{390.251}{177.79}{K}%
\TextAt{16}{404.971}{177.79}{\Character{106}}%
\TextAt{13}{412.194}{177.79}{K}%
\TextAt{9}{423.011}{177.79}{\Character{93}}%
\TextAt{8}{426.263}{177.79}{,}%
\TextAt{9}{433.838}{177.79}{\Character{40}}%
\TextAt{13}{442.629}{177.79}{\Character{59}}%
\TextAt{20}{450.611}{180.812}{\Character{101}}%
\TextAt{13}{447.873}{177.79}{K}%
\TextAt{16}{462.593}{177.79}{\Character{106}}%
\TextAt{13}{469.815}{177.79}{K}%
\TextAt{9}{480.633}{177.79}{\Character{41}}%
\TextAt{9}{489.073}{177.79}{\Character{58}}%
\TextAt{13}{496.241}{177.79}{I}%
\TextAt{21}{501.404}{175.997}{K}%
\NativeArrowAt{513.508}{177.79}{11.955}{19.26}{\longrightarrow}%
\TextAt{13}{79.226}{160.744}{G}%
\TextAt{9}{88.46}{160.745}{\Character{40}}%
\TextAt{20}{95.75}{163.767}{\Character{101}}%
\TextAt{13}{93.012}{160.745}{K}%
\TextAt{16}{107.151}{160.744}{\Character{106}}%
\TextAt{13}{113.792}{160.745}{K}%
\TextAt{9}{124.61}{160.744}{\Character{41}}%
\TextAt{8}{133.06}{160.744}{agree}%
\TextAt{8}{163.629}{160.744}{because}%
\TextAt{8}{206.931}{160.744}{from}%
\TextAt{8}{234.571}{160.744}{chap.}%
\TextAt{8}{265.463}{160.744}{IV,}%
\TextAt{8}{285.608}{160.744}{(6.5),}%
\TextAt{8}{316.835}{160.744}{we}%
\TextAt{8}{334.062}{160.744}{have}%
\TextAt{13}{207.575}{134.961}{d}%
\TextAt{21}{213.658}{133.168}{K}%
\TextAt{16}{224.517}{134.961}{\Character{14}}%
\TextAt{9}{233.151}{134.961}{\Character{40}}%
\TextAt{13}{241.601}{134.961}{\Character{59}}%
\TextAt{20}{249.583}{137.983}{\Character{101}}%
\TextAt{13}{246.845}{134.961}{K}%
\TextAt{16}{260.984}{134.961}{\Character{106}}%
\TextAt{13}{267.625}{134.961}{K}%
\TextAt{9}{278.443}{134.961}{\Character{41}}%
\TextAt{9}{286.298}{134.961}{\Character{61}}%
\TextAt{13}{298.742}{134.961}{v}%
\TextAt{21}{304.401}{133.168}{K}%
\TextAt{9}{315.924}{134.961}{\Character{61}}%
\TextAt{13}{328.35}{134.961}{d}%
\TextAt{21}{334.433}{133.168}{K}%
\TextAt{16}{345.292}{134.961}{\Character{14}}%
\TextAt{9}{353.926}{134.961}{\Character{91}}%
\TextAt{13}{361.075}{134.961}{\Character{59}}%
\TextAt{20}{369.058}{137.983}{\Character{101}}%
\TextAt{13}{366.319}{134.961}{K}%
\TextAt{16}{380.458}{134.961}{\Character{106}}%
\TextAt{13}{387.1}{134.961}{K}%
\TextAt{9}{397.917}{134.961}{\Character{93}}%
\TextAt{13}{401.169}{134.961}{\Character{58}}%
\TextAt{8}{79.226}{109.178}{Thus,}%
\TextAt{8}{112.127}{109.178}{if}%
\TextAt{13}{122.851}{109.178}{L}%
\TextAt{16}{134.136}{109.178}{\Character{106}}%
\TextAt{13}{140.778}{109.178}{K}%
\TextAt{8}{155.493}{109.178}{is}%
\TextAt{8}{167.257}{109.178}{a}%
\TextAt{8}{177.012}{109.178}{subextension}%
\TextAt{8}{246.723}{109.178}{of}%
\TextAt{20}{262.791}{112.2}{\Character{101}}%
\TextAt{13}{260.053}{109.178}{K}%
\TextAt{16}{274.191}{109.178}{\Character{106}}%
\TextAt{13}{280.833}{109.178}{K}%
\TextAt{8}{295.548}{109.178}{and}%
\TextAt{13}{318.311}{109.178}{\Character{11}}%
\TextAt{9}{329.153}{109.178}{\Character{61}}%
\TextAt{9}{341.562}{109.178}{\Character{40}}%
\TextAt{13}{346.131}{109.178}{\Character{11}}%
\TextAt{27}{353.576}{107.385}{p}%
\TextAt{9}{358.063}{109.178}{\Character{41}}%
\NativeMathAt{365.936}{109.178}{11.955}{\in}%
\TextAt{13}{377.227}{109.178}{I}%
\TextAt{21}{382.39}{107.385}{K}%
\TextAt{8}{390.593}{109.178}{,}%
\TextAt{8}{397.742}{109.178}{then}%
\TextAt{9}{200.102}{83.395}{\Character{40}}%
\TextAt{13}{204.655}{83.395}{\Character{11}\Character{59}}%
\TextAt{13}{217.399}{83.395}{L}%
\TextAt{16}{228.706}{83.395}{\Character{106}}%
\TextAt{13}{235.348}{83.395}{K}%
\TextAt{9}{246.165}{83.395}{\Character{41}}%
\TextAt{9}{254.02}{83.395}{\Character{61}}%
\TextAt{9}{266.441}{83.395}{\Character{91}}%
\TextAt{13}{269.716}{83.395}{\Character{11}\Character{59}}%
\TextAt{13}{282.46}{83.395}{L}%
\TextAt{16}{293.767}{83.395}{\Character{106}}%
\TextAt{13}{300.409}{83.395}{K}%
\TextAt{9}{311.226}{83.395}{\Character{93}}%
\TextAt{9}{317.777}{83.395}{\Character{61}}%
\TextAt{20}{330.224}{94.752}{\Character{89}}%
\TextAt{27}{335.868}{69.646}{p}%
\TextAt{9}{345.5}{83.395}{\Character{40}}%
\TextAt{13}{350.053}{83.395}{\Character{11}}%
\TextAt{27}{357.498}{81.601}{p}%
\TextAt{13}{361.985}{83.395}{\Character{59}}%
\TextAt{13}{367.209}{83.395}{L}%
\TextAt{27}{375.194}{81.601}{p}%
\TextAt{16}{383.002}{83.395}{\Character{106}}%
\TextAt{13}{389.643}{83.395}{K}%
\TextAt{27}{399.602}{81.601}{p}%
\TextAt{9}{404.089}{83.395}{\Character{41}}%
\TextAt{13}{408.642}{83.395}{\Character{58}}%
\TextAt{8}{297.211}{43.033}{356}%
\Destination{78.23}{809.815}{nav.page.356}%
\Destination{79.36}{605.655}{nav.print.391}%
\Destination{79.36}{604.655}{nav.number.6.5.5}%
\Destination{79.36}{343.955}{nav.number.6.5.6}%
\Hotlink{123.271}{633.555}{13.995}{11.047}{nav.chapter.4}%
\Hotlink{264.963}{158.555}{13.995}{11.047}{nav.chapter.4}%
\Hotlink{143.416}{633.555}{25.078}{11.047}{nav.number.4.4.7}%
\Hotlink{285.108}{158.555}{25.078}{11.047}{nav.number.4.6.5}%
}

% Page 369 (body page 357)
\Page{%
\TextAt{2}{79.226}{796.943}{VI.}%
\TextAt{2}{98.012}{796.943}{Global}%
\TextAt{2}{133.248}{796.943}{Class}%
\TextAt{2}{161.841}{796.943}{Field}%
\TextAt{2}{189.572}{796.943}{Theory}%
\Vector{q 0 G 0 g 0.299 w 7.226 18.729 m 460.769 18.729 l S Q}%
\TextAt{8}{79.226}{761.124}{In}%
\TextAt{8}{93.494}{761.124}{particular,}%
\TextAt{8}{150.452}{761.124}{for}%
\TextAt{13}{167.983}{761.124}{\Character{11}}%
\TextAt{27}{175.429}{759.33}{p}%
\NativeMathAt{183.236}{761.124}{11.955}{\in}%
\TextAt{13}{194.527}{761.124}{K}%
\TextAt{3}{205.345}{765.462}{\Character{3}}%
\TextAt{27}{204.486}{758.168}{p}%
\TextAt{8}{213.621}{761.124}{we}%
\TextAt{8}{230.495}{761.124}{have}%
\TextAt{8}{257.136}{761.124}{the}%
\TextAt{8}{276.938}{761.124}{identity}%
\TextAt{9}{319.838}{761.124}{\Character{40}}%
\TextAt{16}{324.391}{761.124}{\Character{104}}%
\TextAt{13}{329.04}{761.124}{\Character{11}}%
\TextAt{27}{336.485}{759.33}{p}%
\TextAt{16}{340.972}{761.124}{\Character{105}}%
\TextAt{13}{345.621}{761.124}{\Character{59}}%
\TextAt{13}{350.845}{761.124}{L}%
\TextAt{16}{362.151}{761.124}{\Character{106}}%
\TextAt{13}{368.793}{761.124}{K}%
\TextAt{9}{379.61}{761.124}{\Character{41}}%
\TextAt{9}{387.465}{761.124}{\Character{61}}%
\TextAt{9}{399.886}{761.124}{\Character{40}}%
\TextAt{13}{404.461}{761.124}{\Character{11}}%
\TextAt{27}{411.907}{759.33}{p}%
\TextAt{13}{416.394}{761.124}{\Character{59}}%
\TextAt{13}{421.618}{761.124}{L}%
\TextAt{27}{429.602}{759.33}{p}%
\TextAt{16}{437.41}{761.124}{\Character{106}}%
\TextAt{13}{444.052}{761.124}{K}%
\TextAt{27}{454.011}{759.33}{p}%
\TextAt{9}{458.498}{761.124}{\Character{41}}%
\TextAt{8}{463.05}{761.124}{,}%
\TextAt{8}{469.916}{761.124}{which}%
\TextAt{8}{503.049}{761.124}{shows}%
\TextAt{8}{79.226}{744.078}{that}%
\TextAt{8}{104.595}{744.078}{the}%
\TextAt{8}{124.751}{744.078}{diagram}%
\TextAt{8}{170.289}{744.078}{is}%
\TextAt{8}{182.053}{744.078}{commutative}%
\TextAt{8}{251.978}{744.078}{when}%
\TextAt{8}{282.536}{744.078}{restricted}%
\TextAt{8}{334.624}{744.078}{to}%
\TextAt{8}{348.935}{744.078}{the}%
\TextAt{8}{369.091}{744.078}{finite}%
\TextAt{8}{399.003}{744.078}{subextensions}%
\TextAt{8}{473.328}{744.078}{of}%
\TextAt{20}{489.397}{747.1}{\Character{101}}%
\TextAt{13}{486.658}{744.078}{K}%
\TextAt{16}{500.797}{744.078}{\Character{106}}%
\TextAt{13}{507.439}{744.078}{K}%
\TextAt{8}{518.256}{744.078}{.}%
\TextAt{8}{97.159}{724.043}{On}%
\TextAt{8}{117.421}{724.043}{the}%
\TextAt{8}{138.329}{724.043}{other}%
\TextAt{8}{169.649}{724.043}{hand,}%
\TextAt{8}{203.106}{724.043}{let}%
\TextAt{13}{220.762}{724.043}{L}%
\TextAt{9}{233.328}{724.043}{\Character{61}}%
\TextAt{13}{247.034}{724.043}{K}%
\TextAt{9}{257.851}{724.043}{\Character{40}}%
\TextAt{13}{262.404}{724.043}{i}%
\TextAt{9}{266.397}{724.043}{\Character{41}}%
\TextAt{8}{270.949}{724.043}{,}%
\TextAt{14}{279.038}{724.043}{p}%
\TextAt{16}{289.622}{724.043}{\Character{106}}%
\NativeMathAt{297.524}{724.043}{11.955}{\infty}%
\TextAt{8}{309.499}{724.043}{,}%
\TextAt{8}{317.587}{724.043}{and}%
\TextAt{13}{341.101}{724.043}{L}%
\TextAt{27}{349.066}{722.25}{p}%
\NativeMathAt{358.154}{724.043}{11.955}{\neq}%
\TextAt{13}{371.859}{724.043}{K}%
\TextAt{27}{381.818}{722.25}{p}%
\TextAt{8}{386.305}{724.043}{.}%
\TextAt{8}{397.007}{724.043}{Then}%
\TextAt{13}{428.316}{724.043}{K}%
\TextAt{3}{439.134}{728.381}{\Character{3}}%
\TextAt{27}{438.275}{721.087}{p}%
\TextAt{9}{448.467}{724.043}{\Character{61}}%
\TextAt{18}{462.173}{724.043}{R}%
\TextAt{3}{470.807}{728.381}{\Character{3}}%
\TextAt{8}{475.539}{724.043}{,}%
\TextAt{18}{483.628}{724.043}{R}%
\TextAt{3}{492.262}{728.381}{\Character{3}}%
\TextAt{17}{492.262}{721.087}{\Character{43}}%
\TextAt{8}{503.995}{724.043}{is}%
\TextAt{8}{516.51}{724.043}{the}%
\TextAt{8}{79.226}{706.997}{kernel}%
\TextAt{8}{113.287}{706.997}{of}%
\TextAt{9}{126.21}{706.997}{\Character{40}}%
\TextAt{13}{134.254}{706.997}{\Character{59}}%
\TextAt{13}{139.478}{706.997}{L}%
\TextAt{27}{147.462}{705.204}{p}%
\TextAt{16}{155.27}{706.997}{\Character{106}}%
\TextAt{13}{161.912}{706.997}{K}%
\TextAt{27}{171.871}{705.204}{p}%
\TextAt{9}{176.358}{706.997}{\Character{41}}%
\TextAt{8}{180.91}{706.997}{,}%
\TextAt{8}{187.734}{706.997}{and}%
\TextAt{9}{210.091}{706.997}{\Character{40}}%
\TextAt{16}{214.643}{706.997}{\Character{0}}%
\TextAt{9}{223.942}{706.997}{1}%
\TextAt{13}{229.794}{706.997}{\Character{59}}%
\TextAt{13}{235.018}{706.997}{L}%
\TextAt{27}{243.003}{705.204}{p}%
\TextAt{16}{250.811}{706.997}{\Character{106}}%
\TextAt{13}{257.453}{706.997}{K}%
\TextAt{27}{267.412}{705.204}{p}%
\TextAt{9}{271.899}{706.997}{\Character{41}}%
\TextAt{8}{279.942}{706.997}{is}%
\TextAt{8}{291.3}{706.997}{complex}%
\TextAt{8}{336.741}{706.997}{conjugation}%
\TextAt{8}{399.757}{706.997}{in}%
\TextAt{13}{413.004}{706.997}{G}%
\TextAt{9}{422.237}{706.997}{\Character{40}}%
\TextAt{13}{426.79}{706.997}{L}%
\TextAt{27}{434.754}{705.204}{p}%
\TextAt{16}{442.562}{706.997}{\Character{106}}%
\TextAt{13}{449.204}{706.997}{K}%
\TextAt{27}{459.163}{705.204}{p}%
\TextAt{9}{463.65}{706.997}{\Character{41}}%
\TextAt{9}{471.505}{706.997}{\Character{61}}%
\TextAt{13}{483.948}{706.997}{G}%
\TextAt{9}{493.182}{706.997}{\Character{40}}%
\TextAt{18}{497.734}{706.997}{C}%
\TextAt{16}{509.689}{706.997}{\Character{106}}%
\TextAt{18}{516.331}{706.997}{R}%
\TextAt{9}{524.965}{706.997}{\Character{41}}%
\TextAt{8}{529.518}{706.997}{.}%
\TextAt{8}{79.226}{689.951}{Thus,}%
\TextAt{8}{113.832}{689.951}{all}%
\TextAt{8}{131.456}{689.951}{we}%
\TextAt{8}{150.048}{689.951}{have}%
\TextAt{8}{178.407}{689.951}{to}%
\TextAt{8}{194.082}{689.951}{show}%
\TextAt{8}{224.45}{689.951}{is}%
\TextAt{8}{237.578}{689.951}{that}%
\TextAt{9}{264.312}{689.951}{\Character{40}}%
\TextAt{16}{268.864}{689.951}{\Character{104}\Character{0}}%
\TextAt{9}{282.812}{689.951}{1}%
\TextAt{16}{288.665}{689.951}{\Character{105}}%
\TextAt{13}{293.314}{689.951}{\Character{59}}%
\TextAt{13}{298.538}{689.951}{L}%
\TextAt{16}{312.169}{689.951}{\Character{106}}%
\TextAt{13}{321.136}{689.951}{K}%
\TextAt{9}{331.954}{689.951}{\Character{41}}%
\NativeMathAt{342.152}{689.951}{11.955}{\neq}%
\TextAt{9}{356.893}{689.951}{1}%
\TextAt{8}{362.756}{689.951}{.}%
\TextAt{8}{375.298}{689.951}{If}%
\TextAt{8}{388.354}{689.951}{we}%
\TextAt{8}{406.946}{689.951}{had}%
\TextAt{9}{431.074}{689.951}{\Character{40}}%
\TextAt{16}{435.626}{689.951}{\Character{104}\Character{0}}%
\TextAt{9}{449.574}{689.951}{1}%
\TextAt{16}{455.427}{689.951}{\Character{105}}%
\TextAt{13}{460.076}{689.951}{\Character{59}}%
\TextAt{13}{465.3}{689.951}{L}%
\TextAt{16}{478.931}{689.951}{\Character{106}}%
\TextAt{13}{487.898}{689.951}{K}%
\TextAt{9}{498.715}{689.951}{\Character{41}}%
\TextAt{9}{508.901}{689.951}{\Character{61}}%
\TextAt{9}{523.642}{689.951}{1}%
\TextAt{8}{529.518}{689.951}{,}%
\TextAt{8}{79.226}{672.905}{then}%
\TextAt{8}{107.018}{672.905}{the}%
\TextAt{8}{128.307}{672.905}{class}%
\TextAt{8}{156.876}{672.905}{of}%
\TextAt{16}{171.338}{672.905}{\Character{104}\Character{0}}%
\TextAt{9}{185.286}{672.905}{1}%
\TextAt{16}{191.139}{672.905}{\Character{105}}%
\TextAt{8}{200.818}{672.905}{would}%
\TextAt{8}{236.094}{672.905}{be}%
\TextAt{8}{253.15}{672.905}{the}%
\TextAt{8}{274.439}{672.905}{norm}%
\TextAt{8}{306.14}{672.905}{of}%
\TextAt{8}{320.603}{672.905}{a}%
\TextAt{8}{331.49}{672.905}{class}%
\TextAt{8}{360.059}{672.905}{of}%
\TextAt{13}{374.522}{672.905}{C}%
\TextAt{21}{382.897}{671.112}{L}%
\TextAt{8}{389.154}{672.905}{,}%
\TextAt{8}{397.718}{672.905}{i.e.,}%
\TextAt{16}{421.238}{672.905}{\Character{104}\Character{0}}%
\TextAt{9}{435.186}{672.905}{1}%
\TextAt{16}{441.039}{672.905}{\Character{105}}%
\TextAt{13}{445.688}{672.905}{a}%
\TextAt{9}{457.083}{672.905}{\Character{61}}%
\TextAt{13}{471.438}{672.905}{N}%
\TextAt{21}{480.812}{671.045}{L}%
\TextAt{3}{486.571}{671.045}{\Character{106}}%
\TextAt{21}{488.923}{671.045}{K}%
\TextAt{9}{497.126}{672.905}{\Character{40}}%
\TextAt{13}{501.678}{672.905}{\Character{11}}%
\TextAt{9}{509.2}{672.905}{\Character{41}}%
\TextAt{8}{518.782}{672.905}{for}%
\TextAt{8}{79.226}{655.859}{some}%
\TextAt{13}{108.906}{655.859}{a}%
\NativeMathAt{118.975}{655.859}{11.955}{\in}%
\TextAt{13}{130.869}{655.859}{K}%
\TextAt{3}{141.687}{660.197}{\Character{3}}%
\TextAt{8}{150.67}{655.859}{and}%
\TextAt{8}{173.787}{655.859}{an}%
\TextAt{8}{190.4}{655.859}{idèle}%
\TextAt{13}{218.059}{655.859}{\Character{11}}%
\NativeMathAt{229.505}{655.859}{11.955}{\in}%
\TextAt{13}{241.399}{655.859}{I}%
\TextAt{21}{246.562}{654.066}{L}%
\TextAt{8}{252.819}{655.859}{.}%
\TextAt{8}{262.329}{655.859}{This}%
\TextAt{8}{289.403}{655.859}{would}%
\TextAt{8}{323.901}{655.859}{mean}%
\TextAt{8}{355.469}{655.859}{that}%
\TextAt{13}{381.192}{655.859}{a}%
\TextAt{9}{391.261}{655.859}{\Character{61}}%
\TextAt{13}{404.29}{655.859}{N}%
\TextAt{21}{413.663}{653.999}{L}%
\TextAt{28}{419.423}{653.003}{q}%
\TextAt{3}{423.432}{653.999}{\Character{106}}%
\TextAt{21}{425.784}{653.999}{K}%
\TextAt{28}{432.926}{653.003}{q}%
\TextAt{9}{437.433}{655.859}{\Character{40}}%
\TextAt{13}{441.985}{655.859}{\Character{11}}%
\TextAt{27}{449.431}{654.066}{q}%
\TextAt{9}{453.828}{655.859}{\Character{41}}%
\TextAt{8}{462.631}{655.859}{for}%
\TextAt{14}{480.87}{655.859}{q}%
\NativeMathAt{490.641}{655.859}{11.955}{\neq}%
\TextAt{14}{503.67}{655.859}{p}%
\TextAt{8}{513.904}{655.859}{and}%
\TextAt{16}{79.226}{638.813}{\Character{0}}%
\TextAt{13}{88.525}{638.813}{a}%
\TextAt{9}{99.28}{638.813}{\Character{61}}%
\TextAt{13}{112.994}{638.813}{N}%
\TextAt{21}{122.368}{636.953}{L}%
\TextAt{28}{128.127}{635.957}{p}%
\TextAt{3}{132.21}{636.953}{\Character{106}}%
\TextAt{21}{134.562}{636.953}{K}%
\TextAt{28}{141.704}{635.957}{p}%
\TextAt{9}{146.285}{638.813}{\Character{40}}%
\TextAt{13}{150.837}{638.813}{\Character{11}}%
\TextAt{27}{158.283}{637.02}{p}%
\TextAt{9}{162.769}{638.813}{\Character{41}}%
\TextAt{8}{167.322}{638.813}{,}%
\TextAt{8}{175.417}{638.813}{i.e.,}%
\TextAt{9}{198.467}{638.813}{\Character{40}}%
\TextAt{13}{203.02}{638.813}{a\Character{59}}%
\TextAt{13}{214.389}{638.813}{L}%
\TextAt{27}{222.373}{637.02}{q}%
\TextAt{16}{231.38}{638.813}{\Character{106}}%
\TextAt{13}{239.311}{638.813}{K}%
\TextAt{27}{249.27}{637.02}{q}%
\TextAt{9}{253.666}{638.813}{\Character{41}}%
\TextAt{9}{262.812}{638.813}{\Character{61}}%
\TextAt{9}{276.524}{638.813}{1}%
\TextAt{8}{287.051}{638.813}{for}%
\TextAt{14}{305.692}{638.813}{q}%
\NativeMathAt{316.15}{638.813}{11.955}{\neq}%
\TextAt{14}{329.864}{638.813}{p}%
\TextAt{8}{340.501}{638.813}{and}%
\TextAt{9}{364.021}{638.813}{\Character{40}}%
\TextAt{16}{368.573}{638.813}{\Character{0}}%
\TextAt{13}{377.871}{638.813}{a\Character{59}}%
\TextAt{13}{389.24}{638.813}{L}%
\TextAt{27}{397.225}{637.02}{p}%
\TextAt{16}{406.322}{638.813}{\Character{106}}%
\TextAt{13}{414.253}{638.813}{K}%
\TextAt{27}{424.212}{637.02}{p}%
\TextAt{9}{428.699}{638.813}{\Character{41}}%
\TextAt{9}{437.845}{638.813}{\Character{61}}%
\TextAt{9}{451.557}{638.813}{1}%
\TextAt{8}{457.429}{638.813}{.}%
\TextAt{8}{468.147}{638.813}{By}%
\TextAt{8}{487.267}{638.813}{(5.3),}%
\TextAt{8}{519.439}{638.813}{we}%
\TextAt{8}{79.226}{621.767}{would}%
\TextAt{8}{113.417}{621.767}{have}%
\TextAt{9}{140.46}{621.767}{1}%
\TextAt{9}{149.701}{621.767}{\Character{61}}%
\TextAt{9}{162.207}{621.767}{\Character{91}}%
\TextAt{13}{165.472}{621.767}{a\Character{59}}%
\TextAt{13}{176.841}{621.767}{L}%
\TextAt{16}{188.227}{621.767}{\Character{106}}%
\TextAt{13}{194.95}{621.767}{K}%
\TextAt{9}{205.767}{621.767}{\Character{93}}%
\TextAt{9}{212.402}{621.767}{\Character{61}}%
\TextAt{20}{224.927}{630.734}{\Character{81}}%
\TextAt{27}{236.218}{618.28}{q}%
\TextAt{9}{240.614}{621.767}{\Character{40}}%
\TextAt{13}{245.167}{621.767}{a\Character{59}}%
\TextAt{13}{256.536}{621.767}{L}%
\TextAt{27}{264.52}{619.974}{q}%
\TextAt{16}{272.318}{621.767}{\Character{106}}%
\TextAt{13}{279.041}{621.767}{K}%
\TextAt{27}{289}{619.974}{q}%
\TextAt{9}{293.396}{621.767}{\Character{41}}%
\TextAt{9}{301.334}{621.767}{\Character{61}}%
\TextAt{9}{313.839}{621.767}{\Character{40}}%
\TextAt{13}{318.409}{621.767}{a\Character{59}}%
\TextAt{13}{329.778}{621.767}{L}%
\TextAt{27}{337.762}{619.974}{p}%
\TextAt{16}{345.651}{621.767}{\Character{106}}%
\TextAt{13}{352.373}{621.767}{K}%
\TextAt{27}{362.332}{619.974}{p}%
\TextAt{9}{366.819}{621.767}{\Character{41}}%
\TextAt{8}{371.371}{621.767}{,}%
\TextAt{8}{378.58}{621.767}{so}%
\TextAt{8}{392.997}{621.767}{that}%
\TextAt{9}{418.413}{621.767}{\Character{40}}%
\TextAt{16}{422.966}{621.767}{\Character{0}}%
\TextAt{9}{432.264}{621.767}{1}%
\TextAt{13}{438.117}{621.767}{\Character{59}}%
\TextAt{13}{443.341}{621.767}{L}%
\TextAt{27}{451.326}{619.974}{p}%
\TextAt{16}{459.214}{621.767}{\Character{106}}%
\TextAt{13}{465.937}{621.767}{K}%
\TextAt{27}{475.896}{619.974}{p}%
\TextAt{9}{480.383}{621.767}{\Character{41}}%
\TextAt{9}{488.321}{621.767}{\Character{61}}%
\TextAt{9}{500.826}{621.767}{1}%
\TextAt{8}{506.696}{621.767}{,}%
\TextAt{8}{513.904}{621.767}{and}%
\TextAt{8}{79.226}{604.721}{therefore}%
\TextAt{16}{128.326}{604.721}{\Character{0}}%
\TextAt{9}{137.625}{604.721}{1}%
\NativeMathAt{146.798}{604.721}{11.955}{\in}%
\TextAt{13}{158.089}{604.721}{N}%
\TextAt{21}{167.463}{602.861}{L}%
\TextAt{28}{173.222}{601.865}{p}%
\TextAt{3}{177.305}{602.861}{\Character{106}}%
\TextAt{21}{179.657}{602.861}{K}%
\TextAt{28}{186.799}{601.865}{p}%
\TextAt{9}{191.379}{604.721}{\Character{40}}%
\TextAt{13}{195.932}{604.721}{L}%
\TextAt{3}{203.896}{609.06}{\Character{3}}%
\TextAt{27}{203.896}{601.766}{p}%
\TextAt{9}{208.629}{604.721}{\Character{41}}%
\TextAt{9}{216.484}{604.721}{\Character{61}}%
\TextAt{13}{228.927}{604.721}{N}%
\TextAt{25}{238.301}{602.861}{C}%
\TextAt{3}{244.057}{602.861}{\Character{106}}%
\TextAt{25}{246.41}{602.861}{R}%
\TextAt{18}{252.664}{604.721}{C}%
\TextAt{3}{261.298}{609.06}{\Character{3}}%
\TextAt{9}{269.351}{604.721}{\Character{61}}%
\TextAt{18}{281.777}{604.721}{R}%
\TextAt{3}{290.411}{609.06}{\Character{3}}%
\TextAt{17}{290.411}{601.766}{\Character{43}}%
\TextAt{8}{297.496}{604.721}{,}%
\TextAt{8}{304.645}{604.721}{a}%
\TextAt{8}{314.4}{604.721}{contradiction.}%
\TextAt{8}{97.159}{584.686}{We}%
\TextAt{8}{117.899}{584.686}{now}%
\TextAt{8}{142.883}{584.686}{reduce}%
\TextAt{8}{180.539}{584.686}{the}%
\TextAt{8}{201.291}{584.686}{general}%
\TextAt{8}{242.211}{584.686}{case}%
\TextAt{8}{267.578}{584.686}{to}%
\TextAt{8}{282.483}{584.686}{these}%
\TextAt{8}{313.05}{584.686}{special}%
\TextAt{8}{351.747}{584.686}{cases}%
\TextAt{8}{381.728}{584.686}{as}%
\TextAt{8}{396.693}{584.686}{follows.}%
\TextAt{8}{441.466}{584.686}{Let}%
\TextAt{13}{463.031}{584.686}{L}%
\TextAt{3}{470.996}{589.025}{\Character{48}}%
\TextAt{16}{478.126}{584.686}{\Character{106}}%
\TextAt{13}{485.783}{584.686}{K}%
\TextAt{3}{496.6}{589.025}{\Character{48}}%
\TextAt{8}{503.888}{584.686}{be}%
\TextAt{8}{520.408}{584.686}{an}%
\TextAt{8}{79.226}{567.64}{abelian}%
\TextAt{8}{119.874}{567.64}{extension,}%
\TextAt{8}{174.891}{567.64}{so}%
\TextAt{8}{189.262}{567.64}{that}%
\TextAt{13}{214.63}{567.64}{K}%
\NativeMathAt{228.769}{567.64}{11.955}{\subseteq}%
\TextAt{13}{241.388}{567.64}{K}%
\TextAt{3}{252.206}{571.979}{\Character{48}}%
\TextAt{8}{255.001}{567.64}{,}%
\TextAt{13}{262.15}{567.64}{L}%
\NativeMathAt{273.436}{567.64}{11.955}{\subseteq}%
\TextAt{13}{286.055}{567.64}{L}%
\TextAt{3}{294.019}{571.979}{\Character{48}}%
\TextAt{8}{296.814}{567.64}{.}%
\TextAt{8}{305.263}{567.64}{We}%
\TextAt{8}{325.407}{567.64}{then}%
\TextAt{8}{352.067}{567.64}{consider}%
\TextAt{8}{397.652}{567.64}{the}%
\TextAt{8}{417.809}{567.64}{diagram}%
\NativeDiagram{%
\path[fill=black!0] (95.853,-263.592) -- (151.655,-263.592) -- (151.655,-248.282) -- (95.853,-248.282) -- cycle;%
}%
\TextAt{47}{169.846}{515.544}{G}%
\TextAt{48}{178.423}{515.544}{\Character{40}}%
\TextAt{47}{182.666}{515.544}{L}%
\TextAt{27}{190.09}{513.908}{p}%
\TextAt{49}{197.607}{515.544}{\Character{106}}%
\TextAt{47}{203.668}{515.544}{K}%
\TextAt{27}{212.933}{513.908}{p}%
\TextAt{48}{217.42}{515.544}{\Character{41}}%
\NativeDiagram{%
\path[fill=black!0] (259.307,-264.194) -- (328.363,-264.194) -- (328.363,-247.68) -- (259.307,-247.68) -- cycle;%
}%
\TextAt{47}{333.299}{516.146}{K}%
\TextAt{3}{343.345}{520.105}{\Character{3}}%
\TextAt{27}{342.564}{513.449}{p}%
\TextAt{48}{348.077}{516.146}{\Character{47}}%
\TextAt{47}{353.531}{516.146}{N}%
\TextAt{21}{362.297}{513.926}{L}%
\TextAt{28}{368.056}{512.93}{p}%
\TextAt{3}{372.138}{513.926}{\Character{106}}%
\TextAt{21}{374.491}{513.926}{K}%
\TextAt{28}{381.632}{512.93}{p}%
\TextAt{47}{386.213}{516.146}{L}%
\TextAt{3}{393.637}{520.105}{\Character{3}}%
\TextAt{27}{393.637}{513.449}{p}%
\NativeDiagram{%
\path[fill=black!0] (99.95,-325.747) -- (147.558,-325.747) -- (147.558,-310.853) -- (99.95,-310.853) -- cycle;%
}%
\TextAt{47}{173.942}{452.973}{G}%
\TextAt{48}{182.519}{452.973}{\Character{40}}%
\TextAt{47}{186.762}{452.973}{L}%
\TextAt{49}{197.216}{452.973}{\Character{106}}%
\TextAt{47}{203.277}{452.973}{K}%
\TextAt{48}{213.322}{452.973}{\Character{41}}%
\NativeDiagram{%
\path[fill=black!0] (261.548,-326.49) -- (326.122,-326.49) -- (326.122,-310.11) -- (261.548,-310.11) -- cycle;%
}%
\TextAt{47}{335.541}{453.715}{C}%
\TextAt{21}{343.338}{452.026}{K}%
\TextAt{48}{351.541}{453.715}{\Character{47}}%
\TextAt{47}{356.995}{453.715}{N}%
\TextAt{21}{365.76}{451.495}{L}%
\TextAt{3}{371.52}{451.495}{\Character{106}}%
\TextAt{21}{373.872}{451.495}{K}%
\TextAt{47}{382.075}{453.715}{C}%
\TextAt{21}{389.871}{452.026}{L}%
\NativeDiagram{%
\path[fill=black!0] (135.538,-235.878) -- (191.34,-235.878) -- (191.34,-219.302) -- (135.538,-219.302) -- cycle;%
}%
\TextAt{47}{209.531}{544.318}{G}%
\TextAt{48}{218.108}{544.318}{\Character{40}}%
\TextAt{47}{222.351}{544.318}{L}%
\TextAt{3}{229.775}{548.277}{\Character{48}}%
\TextAt{27}{229.775}{541.621}{p}%
\TextAt{49}{237.292}{544.318}{\Character{106}}%
\TextAt{47}{243.353}{544.318}{K}%
\TextAt{3}{253.398}{548.277}{\Character{48}}%
\TextAt{27}{252.618}{541.621}{p}%
\TextAt{48}{257.105}{544.318}{\Character{41}}%
\NativeDiagram{%
\path[fill=black!0] (294.895,-236.408) -- (372.145,-236.408) -- (372.145,-218.773) -- (294.895,-218.773) -- cycle;%
}%
\TextAt{47}{368.888}{544.848}{K}%
\TextAt{3}{378.934}{548.807}{\Character{48}}%
\TextAt{27}{378.153}{542.151}{p}%
\TextAt{3}{382.64}{548.807}{\Character{3}}%
\TextAt{48}{387.372}{544.848}{\Character{47}}%
\TextAt{47}{392.827}{544.848}{N}%
\TextAt{21}{401.592}{542.539}{L}%
\TextAt{23}{407.351}{545.285}{\Character{48}}%
\TextAt{28}{407.351}{540.715}{p}%
\TextAt{3}{411.434}{542.539}{\Character{106}}%
\TextAt{21}{413.786}{542.539}{K}%
\TextAt{23}{421.491}{545.285}{\Character{48}}%
\TextAt{28}{420.928}{540.715}{p}%
\TextAt{47}{425.509}{544.848}{L}%
\TextAt{3}{432.933}{548.807}{\Character{48}}%
\TextAt{27}{432.933}{542.151}{p}%
\TextAt{3}{437.42}{548.807}{\Character{3}}%
\NativeDiagram{%
\path[fill=black!0] (136.84,-297.503) -- (190.039,-297.503) -- (190.039,-282.403) -- (136.84,-282.403) -- cycle;%
}%
\TextAt{47}{210.833}{481.217}{G}%
\TextAt{48}{219.41}{481.217}{\Character{40}}%
\TextAt{47}{223.653}{481.217}{L}%
\TextAt{3}{231.077}{485.176}{\Character{48}}%
\TextAt{49}{236.902}{481.217}{\Character{106}}%
\TextAt{47}{242.963}{481.217}{K}%
\TextAt{3}{253.008}{485.176}{\Character{48}}%
\TextAt{48}{255.804}{481.217}{\Character{41}}%
\NativeDiagram{%
\path[fill=black!0] (295.865,-298.143) -- (371.176,-298.143) -- (371.176,-281.763) -- (295.865,-281.763) -- cycle;%
}%
\TextAt{47}{369.857}{482.062}{C}%
\TextAt{21}{377.654}{480.229}{K}%
\TextAt{23}{385.359}{482.499}{\Character{48}}%
\TextAt{48}{388.541}{482.062}{\Character{47}}%
\TextAt{47}{393.996}{482.062}{N}%
\TextAt{21}{402.761}{479.842}{L}%
\TextAt{23}{408.52}{482.111}{\Character{48}}%
\TextAt{3}{411.204}{479.842}{\Character{106}}%
\TextAt{21}{413.557}{479.842}{K}%
\TextAt{23}{421.261}{482.111}{\Character{48}}%
\TextAt{47}{424.444}{482.062}{C}%
\TextAt{21}{432.241}{480.229}{L}%
\TextAt{23}{438}{482.499}{\Character{48}}%
\NativeDiagram{%
\path[draw=black,line width=0.399bp] (191.54,-227.59) -- (291.632,-227.59);%
\path[draw=black,line width=0.399bp,fill=black] (294.129,-227.59) -- (290.557,-226.244) -- (291.732,-227.59) -- (290.557,-228.936) -- cycle;%
\path[draw=black,line width=0.399bp] (151.557,-236.077) -- (137.241,-246.302);%
\path[draw=black,line width=0.399bp,fill=black] (135.21,-247.753) -- (138.898,-246.772) -- (137.16,-246.36) -- (137.334,-244.582) -- cycle;%
\path[draw=black,line width=0.399bp] (320.898,-236.607) -- (308.165,-245.7);%
\path[draw=black,line width=0.399bp,fill=black] (306.134,-247.151) -- (309.822,-246.171) -- (308.084,-245.758) -- (308.258,-243.98) -- cycle;%
\path[draw=black,line width=0.399bp] (151.854,-255.937) -- (256.044,-255.937);%
\path[draw=black,line width=0.399bp,fill=black] (258.54,-255.937) -- (254.969,-254.591) -- (256.143,-255.937) -- (254.969,-257.283) -- cycle;%
\path[draw=black,line width=0.399bp] (147.757,-318.3) -- (258.285,-318.3);%
\path[draw=black,line width=0.399bp,fill=black] (260.781,-318.3) -- (257.21,-316.954) -- (258.385,-318.3) -- (257.21,-319.646) -- cycle;%
\path[draw=black,line width=0.399bp] (190.238,-289.953) -- (292.602,-289.953);%
\path[draw=black,line width=0.399bp,fill=black] (295.098,-289.953) -- (291.527,-288.607) -- (292.701,-289.953) -- (291.527,-291.299) -- cycle;%
\path[draw=black,line width=0.399bp] (163.439,-236.077) -- (163.439,-279.14);%
\path[draw=black,line width=0.399bp,fill=black] (163.439,-281.637) -- (164.785,-278.066) -- (163.439,-279.24) -- (162.094,-278.066) -- cycle;%
\path[draw=black,line width=0.399bp] (333.52,-236.607) -- (333.52,-278.5);%
\path[draw=black,line width=0.399bp,fill=black] (333.52,-280.997) -- (334.866,-277.425) -- (333.52,-278.6) -- (332.174,-277.425) -- cycle;%
\path[draw=black,line width=0.399bp] (123.754,-263.791) -- (123.754,-307.59);%
\path[draw=black,line width=0.399bp,fill=black] (123.754,-310.086) -- (125.1,-306.515) -- (123.754,-307.689) -- (122.408,-306.515) -- cycle;%
\path[draw=black,line width=0.399bp] (293.835,-264.393) -- (293.835,-306.847);%
\path[draw=black,line width=0.399bp,fill=black] (293.835,-309.343) -- (295.181,-305.772) -- (293.835,-306.946) -- (292.489,-305.772) -- cycle;%
\path[draw=black,line width=0.399bp] (152.591,-297.702) -- (136.95,-308.873);%
\path[draw=black,line width=0.399bp,fill=black] (134.919,-310.324) -- (138.607,-309.343) -- (136.869,-308.931) -- (137.043,-307.153) -- cycle;%
\path[draw=black,line width=0.399bp] (321.777,-298.342) -- (308.071,-308.13);%
\path[draw=black,line width=0.399bp,fill=black] (306.04,-309.581) -- (309.727,-308.601) -- (307.99,-308.188) -- (308.163,-306.411) -- cycle;%
}%
\TextAt{8}{79.226}{421.499}{where}%
\TextAt{13}{114.197}{421.499}{L}%
\TextAt{27}{122.161}{419.706}{p}%
\TextAt{9}{131.948}{421.499}{\Character{61}}%
\TextAt{13}{146.352}{421.499}{K}%
\TextAt{27}{156.311}{419.706}{p}%
\TextAt{13}{160.798}{421.499}{L}%
\TextAt{8}{168.763}{421.499}{,}%
\TextAt{13}{177.363}{421.499}{K}%
\TextAt{3}{188.181}{425.837}{\Character{48}}%
\TextAt{27}{187.322}{418.543}{p}%
\TextAt{9}{197.109}{421.499}{\Character{61}}%
\TextAt{13}{211.514}{421.499}{K}%
\TextAt{27}{221.473}{419.706}{p}%
\TextAt{13}{225.96}{421.499}{K}%
\TextAt{3}{236.777}{425.837}{\Character{48}}%
\TextAt{8}{239.572}{421.499}{,}%
\TextAt{13}{248.173}{421.499}{L}%
\TextAt{3}{256.138}{425.837}{\Character{48}}%
\TextAt{27}{256.138}{418.543}{p}%
\TextAt{9}{265.924}{421.499}{\Character{61}}%
\TextAt{13}{280.329}{421.499}{K}%
\TextAt{27}{290.288}{419.706}{p}%
\TextAt{13}{294.775}{421.499}{L}%
\TextAt{3}{302.739}{425.837}{\Character{48}}%
\TextAt{8}{305.534}{421.499}{.}%
\TextAt{8}{317.466}{421.499}{In}%
\TextAt{8}{333.249}{421.499}{this}%
\TextAt{8}{357.232}{421.499}{diagram,}%
\TextAt{8}{407.473}{421.499}{the}%
\TextAt{8}{428.791}{421.499}{top}%
\TextAt{8}{450.766}{421.499}{and}%
\TextAt{8}{474.69}{421.499}{bottom}%
\TextAt{8}{517.156}{421.499}{are}%
\TextAt{8}{79.226}{404.453}{commutative}%
\TextAt{8}{148.705}{404.453}{by}%
\TextAt{8}{164.518}{404.453}{chap.}%
\TextAt{8}{194.963}{404.453}{IV,}%
\TextAt{8}{214.661}{404.453}{(6.4),}%
\TextAt{8}{245.531}{404.453}{and}%
\TextAt{8}{267.847}{404.453}{the}%
\TextAt{8}{287.556}{404.453}{sides}%
\TextAt{8}{315.192}{404.453}{are}%
\TextAt{8}{334.257}{404.453}{commutative}%
\TextAt{8}{403.736}{404.453}{for}%
\TextAt{8}{421.174}{404.453}{trivial}%
\TextAt{8}{455.529}{404.453}{reasons.}%
\TextAt{8}{501.033}{404.453}{If}%
\TextAt{8}{512.278}{404.453}{now}%
\TextAt{13}{79.226}{387.407}{L}%
\TextAt{3}{87.191}{391.745}{\Character{48}}%
\TextAt{16}{94.208}{387.407}{\Character{106}}%
\TextAt{13}{101.751}{387.407}{K}%
\TextAt{3}{112.569}{391.745}{\Character{48}}%
\TextAt{8}{119.79}{387.407}{is}%
\TextAt{8}{132.083}{387.407}{one}%
\TextAt{8}{154.072}{387.407}{of}%
\TextAt{8}{167.931}{387.407}{the}%
\TextAt{8}{188.616}{387.407}{special}%
\TextAt{8}{227.246}{387.407}{extensions}%
\TextAt{8}{284.156}{387.407}{for}%
\TextAt{8}{302.57}{387.407}{which}%
\TextAt{8}{336.585}{387.407}{the}%
\TextAt{8}{357.271}{387.407}{proposition}%
\TextAt{8}{419.333}{387.407}{is}%
\TextAt{8}{431.626}{387.407}{already}%
\TextAt{8}{473.46}{387.407}{established,}%
\TextAt{8}{79.226}{370.361}{then}%
\TextAt{8}{106.226}{370.361}{the}%
\TextAt{8}{126.721}{370.361}{back}%
\TextAt{8}{154.378}{370.361}{diagram}%
\TextAt{8}{200.255}{370.361}{is}%
\TextAt{8}{212.358}{370.361}{commutative,}%
\TextAt{8}{285.96}{370.361}{and}%
\TextAt{8}{309.062}{370.361}{hence}%
\TextAt{8}{341.907}{370.361}{also}%
\TextAt{8}{365.727}{370.361}{the}%
\TextAt{8}{386.223}{370.361}{front}%
\TextAt{8}{415.183}{370.361}{one,}%
\TextAt{8}{440.318}{370.361}{for}%
\TextAt{8}{458.543}{370.361}{all}%
\TextAt{8}{475.141}{370.361}{elements}%
\TextAt{8}{523.337}{370.361}{of}%
\TextAt{13}{79.226}{353.315}{G}%
\TextAt{9}{88.46}{353.315}{\Character{40}}%
\TextAt{13}{93.012}{353.315}{L}%
\TextAt{27}{100.977}{351.522}{p}%
\TextAt{16}{110.155}{353.315}{\Character{106}}%
\TextAt{13}{118.168}{353.315}{K}%
\TextAt{27}{128.127}{351.522}{p}%
\TextAt{9}{132.614}{353.315}{\Character{41}}%
\TextAt{8}{141.868}{353.315}{in}%
\TextAt{8}{156.325}{353.315}{the}%
\TextAt{8}{177.286}{353.315}{image}%
\TextAt{8}{211.912}{353.315}{of}%
\TextAt{13}{226.046}{353.315}{G}%
\TextAt{9}{235.28}{353.315}{\Character{40}}%
\TextAt{13}{239.832}{353.315}{L}%
\TextAt{3}{247.797}{357.653}{\Character{48}}%
\TextAt{27}{247.797}{350.359}{p}%
\TextAt{16}{256.975}{353.315}{\Character{106}}%
\TextAt{13}{264.988}{353.315}{K}%
\TextAt{3}{275.805}{357.653}{\Character{48}}%
\TextAt{27}{274.947}{350.359}{p}%
\TextAt{9}{279.434}{353.315}{\Character{41}}%
\NativeArrowAt{288.678}{353.315}{11.955}{19.26}{\longrightarrow}%
\TextAt{13}{312.631}{353.315}{G}%
\TextAt{9}{321.864}{353.315}{\Character{40}}%
\TextAt{13}{326.417}{353.315}{L}%
\TextAt{27}{334.381}{351.522}{p}%
\TextAt{16}{343.56}{353.315}{\Character{106}}%
\TextAt{13}{351.572}{353.315}{K}%
\TextAt{27}{361.531}{351.522}{p}%
\TextAt{9}{366.018}{353.315}{\Character{41}}%
\TextAt{8}{370.571}{353.315}{.}%
\TextAt{8}{381.432}{353.315}{This}%
\TextAt{8}{408.956}{353.315}{makes}%
\TextAt{8}{444.945}{353.315}{it}%
\TextAt{8}{457.454}{353.315}{clear}%
\TextAt{8}{486.221}{353.315}{that}%
\TextAt{8}{512.394}{353.315}{it}%
\TextAt{8}{524.903}{353.315}{is}%
\TextAt{8}{79.226}{336.269}{enough}%
\TextAt{8}{119.675}{336.269}{to}%
\TextAt{8}{134.109}{336.269}{find,}%
\TextAt{8}{160.924}{336.269}{for}%
\TextAt{8}{178.933}{336.269}{every}%
\TextAt{13}{209.949}{336.269}{\Character{27}}%
\NativeMathAt{220.563}{336.269}{11.955}{\in}%
\TextAt{13}{232.065}{336.269}{G}%
\TextAt{9}{241.299}{336.269}{\Character{40}}%
\TextAt{13}{245.851}{336.269}{L}%
\TextAt{27}{253.816}{334.476}{p}%
\TextAt{16}{261.835}{336.269}{\Character{106}}%
\TextAt{13}{268.688}{336.269}{K}%
\TextAt{27}{278.647}{334.476}{p}%
\TextAt{9}{283.134}{336.269}{\Character{41}}%
\TextAt{8}{287.686}{336.269}{,}%
\TextAt{8}{294.99}{336.269}{some}%
\TextAt{8}{324.44}{336.269}{special}%
\TextAt{8}{362.665}{336.269}{extension}%
\TextAt{13}{414.555}{336.269}{L}%
\TextAt{3}{422.519}{340.607}{\Character{48}}%
\TextAt{16}{428.846}{336.269}{\Character{106}}%
\TextAt{13}{435.699}{336.269}{K}%
\TextAt{3}{446.517}{340.607}{\Character{48}}%
\TextAt{8}{453.333}{336.269}{such}%
\TextAt{8}{479.854}{336.269}{that}%
\TextAt{13}{505.347}{336.269}{\Character{27}}%
\TextAt{8}{516.451}{336.269}{lies}%
\TextAt{8}{79.226}{319.223}{in}%
\TextAt{8}{93.209}{319.223}{the}%
\TextAt{8}{113.696}{319.223}{image}%
\TextAt{8}{147.848}{319.223}{of}%
\TextAt{13}{161.508}{319.223}{G}%
\TextAt{9}{170.742}{319.223}{\Character{40}}%
\TextAt{13}{175.294}{319.223}{L}%
\TextAt{3}{183.258}{323.561}{\Character{48}}%
\TextAt{27}{183.258}{316.267}{p}%
\TextAt{16}{191.629}{319.223}{\Character{106}}%
\TextAt{13}{198.834}{319.223}{K}%
\TextAt{3}{209.651}{323.561}{\Character{48}}%
\TextAt{27}{208.793}{316.267}{p}%
\TextAt{9}{213.28}{319.223}{\Character{41}}%
\TextAt{8}{217.832}{319.223}{.}%
\TextAt{8}{227.271}{319.223}{It}%
\TextAt{8}{240.274}{319.223}{is}%
\TextAt{8}{252.368}{319.223}{even}%
\TextAt{8}{279.359}{319.223}{sufficient}%
\TextAt{8}{328.849}{319.223}{to}%
\TextAt{8}{343.49}{319.223}{do}%
\TextAt{8}{360.079}{319.223}{this}%
\TextAt{8}{383.232}{319.223}{only}%
\TextAt{8}{409.254}{319.223}{for}%
\TextAt{8}{427.469}{319.223}{all}%
\TextAt{13}{444.058}{319.223}{\Character{27}}%
\TextAt{8}{455.369}{319.223}{of}%
\TextAt{8}{469.029}{319.223}{prime}%
\TextAt{8}{502.523}{319.223}{power}%
\TextAt{8}{79.226}{302.177}{order,}%
\TextAt{8}{113.886}{302.177}{because}%
\TextAt{8}{157.858}{302.177}{they}%
\TextAt{8}{184.866}{302.177}{generate}%
\TextAt{8}{232.366}{302.177}{the}%
\TextAt{8}{253.193}{302.177}{group.}%
\TextAt{8}{292.932}{302.177}{Passing}%
\TextAt{8}{335.84}{302.177}{to}%
\TextAt{8}{350.821}{302.177}{the}%
\TextAt{8}{371.649}{302.177}{fixed}%
\TextAt{8}{400.605}{302.177}{field}%
\TextAt{8}{426.633}{302.177}{of}%
\TextAt{13}{440.634}{302.177}{\Character{27}}%
\TextAt{8}{452.285}{302.177}{we}%
\TextAt{8}{470.183}{302.177}{may}%
\TextAt{8}{496.223}{302.177}{assume}%
\TextAt{8}{79.226}{285.131}{moreover}%
\TextAt{8}{129.641}{285.131}{that}%
\TextAt{13}{155.01}{285.131}{G}%
\TextAt{9}{164.244}{285.131}{\Character{40}}%
\TextAt{13}{168.796}{285.131}{L}%
\TextAt{16}{180.081}{285.131}{\Character{106}}%
\TextAt{13}{186.723}{285.131}{K}%
\TextAt{9}{197.541}{285.131}{\Character{41}}%
\TextAt{8}{205.99}{285.131}{is}%
\TextAt{8}{217.754}{285.131}{generated}%
\TextAt{8}{271.086}{285.131}{by}%
\TextAt{13}{287.345}{285.131}{\Character{27}}%
\TextAt{8}{294.428}{285.131}{.}%
\TextAt{8}{97.159}{265.096}{When}%
\TextAt{14}{131.63}{265.096}{p}%
\TextAt{16}{141.512}{265.096}{\Character{106}}%
\NativeMathAt{148.721}{265.096}{11.955}{\infty}%
\TextAt{8}{164.923}{265.096}{and}%
\TextAt{13}{188.024}{265.096}{L}%
\TextAt{27}{195.989}{263.303}{p}%
\NativeMathAt{204.374}{265.096}{11.955}{\neq}%
\TextAt{13}{217.377}{265.096}{K}%
\TextAt{27}{227.337}{263.303}{p}%
\TextAt{8}{231.823}{265.096}{,}%
\TextAt{8}{239.396}{265.096}{i.e.,}%
\TextAt{13}{261.925}{265.096}{K}%
\TextAt{27}{271.884}{263.303}{p}%
\TextAt{9}{280.27}{265.096}{\Character{61}}%
\TextAt{18}{293.273}{265.096}{R}%
\TextAt{8}{301.907}{265.096}{,}%
\TextAt{13}{309.48}{265.096}{L}%
\TextAt{27}{317.445}{263.303}{p}%
\TextAt{9}{325.83}{265.096}{\Character{61}}%
\TextAt{18}{338.833}{265.096}{C}%
\TextAt{8}{347.468}{265.096}{,}%
\TextAt{8}{355.041}{265.096}{we}%
\TextAt{8}{372.607}{265.096}{put}%
\TextAt{13}{394.406}{265.096}{L}%
\TextAt{3}{402.37}{269.435}{\Character{48}}%
\TextAt{9}{409.064}{265.096}{\Character{61}}%
\TextAt{13}{422.067}{265.096}{L}%
\TextAt{9}{430.031}{265.096}{\Character{40}}%
\TextAt{13}{434.584}{265.096}{i}%
\TextAt{9}{438.577}{265.096}{\Character{41}}%
\NativeMathAt{447.028}{265.096}{11.955}{\subseteq}%
\TextAt{18}{460.225}{265.096}{C}%
\TextAt{8}{468.86}{265.096}{,}%
\TextAt{8}{476.432}{265.096}{and}%
\TextAt{8}{499.534}{265.096}{choose}%
\TextAt{8}{79.226}{248.05}{for}%
\TextAt{13}{97.882}{248.05}{K}%
\TextAt{3}{108.7}{252.389}{\Character{48}}%
\TextAt{8}{116.164}{248.05}{the}%
\TextAt{8}{137.092}{248.05}{fixed}%
\TextAt{8}{166.149}{248.05}{field}%
\TextAt{8}{192.278}{248.05}{of}%
\TextAt{8}{206.379}{248.05}{the}%
\TextAt{8}{227.307}{248.05}{restriction}%
\TextAt{8}{284.076}{248.05}{of}%
\TextAt{8}{298.178}{248.05}{complex}%
\TextAt{8}{344.797}{248.05}{conjugation}%
\TextAt{8}{408.991}{248.05}{to}%
\TextAt{13}{424.073}{248.05}{L}%
\TextAt{3}{432.037}{252.389}{\Character{48}}%
\TextAt{8}{434.832}{248.05}{.}%
\TextAt{8}{445.595}{248.05}{Then}%
\TextAt{13}{476.924}{248.05}{L}%
\TextAt{3}{484.888}{252.389}{\Character{48}}%
\TextAt{9}{492.319}{248.05}{\Character{61}}%
\TextAt{13}{506.059}{248.05}{K}%
\TextAt{3}{516.876}{252.389}{\Character{48}}%
\TextAt{9}{519.671}{248.05}{\Character{40}}%
\TextAt{13}{524.224}{248.05}{i}%
\TextAt{9}{528.217}{248.05}{\Character{41}}%
\TextAt{8}{79.226}{231.004}{and}%
\TextAt{13}{101.989}{231.004}{K}%
\TextAt{3}{112.806}{235.343}{\Character{48}}%
\TextAt{27}{111.948}{228.049}{p}%
\TextAt{9}{119.756}{231.004}{\Character{61}}%
\TextAt{18}{132.181}{231.004}{R}%
\TextAt{8}{140.815}{231.004}{,}%
\TextAt{13}{147.965}{231.004}{L}%
\TextAt{3}{155.929}{235.343}{\Character{48}}%
\TextAt{27}{155.929}{228.049}{p}%
\TextAt{9}{163.737}{231.004}{\Character{61}}%
\TextAt{18}{176.162}{231.004}{C}%
\TextAt{8}{184.797}{231.004}{,}%
\TextAt{8}{191.946}{231.004}{so}%
\TextAt{8}{206.316}{231.004}{the}%
\TextAt{8}{226.472}{231.004}{mapping}%
\TextAt{13}{274.604}{231.004}{G}%
\TextAt{9}{283.837}{231.004}{\Character{40}}%
\TextAt{13}{288.39}{231.004}{L}%
\TextAt{3}{296.354}{235.343}{\Character{48}}%
\TextAt{27}{296.354}{228.049}{p}%
\TextAt{16}{304.162}{231.004}{\Character{106}}%
\TextAt{13}{310.804}{231.004}{K}%
\TextAt{3}{321.621}{235.343}{\Character{48}}%
\TextAt{27}{320.763}{228.049}{p}%
\TextAt{9}{325.25}{231.004}{\Character{41}}%
\NativeArrowAt{333.123}{231.004}{11.955}{19.26}{\longrightarrow}%
\TextAt{13}{355.705}{231.004}{G}%
\TextAt{9}{364.938}{231.004}{\Character{40}}%
\TextAt{13}{369.491}{231.004}{L}%
\TextAt{27}{377.455}{229.211}{p}%
\TextAt{16}{385.263}{231.004}{\Character{106}}%
\TextAt{13}{391.905}{231.004}{K}%
\TextAt{27}{401.864}{229.211}{p}%
\TextAt{9}{406.351}{231.004}{\Character{41}}%
\TextAt{8}{414.8}{231.004}{is}%
\TextAt{8}{426.564}{231.004}{surjective.}%
\TextAt{8}{97.159}{210.97}{When}%
\TextAt{14}{131.36}{210.97}{p}%
\TextAt{18}{140.781}{210.97}{\Character{45}}%
\NativeMathAt{147.54}{210.97}{11.955}{\infty}%
\TextAt{8}{159.495}{210.97}{,}%
\TextAt{8}{166.73}{210.97}{we}%
\TextAt{8}{184.027}{210.97}{find}%
\TextAt{8}{207.504}{210.97}{the}%
\TextAt{8}{227.729}{210.97}{extension}%
\TextAt{13}{279.564}{210.97}{L}%
\TextAt{3}{287.528}{215.308}{\Character{48}}%
\TextAt{16}{293.761}{210.97}{\Character{106}}%
\TextAt{13}{300.52}{210.97}{K}%
\TextAt{3}{311.338}{215.308}{\Character{48}}%
\TextAt{8}{318.099}{210.97}{as}%
\TextAt{8}{332.538}{210.97}{follows.}%
\TextAt{8}{375.732}{210.97}{Let}%
\TextAt{13}{396.77}{210.97}{\Character{27}}%
\TextAt{8}{407.818}{210.97}{be}%
\TextAt{8}{423.811}{210.97}{of}%
\TextAt{13}{437.21}{210.97}{p}%
\TextAt{8}{443.085}{210.97}{-power}%
\TextAt{8}{481.196}{210.97}{order.}%
\TextAt{8}{516.522}{210.97}{We}%
\TextAt{8}{79.226}{193.924}{denote}%
\TextAt{8}{117.704}{193.924}{by}%
\TextAt{20}{137.461}{196.946}{\Character{98}}%
\TextAt{13}{134.723}{193.924}{K}%
\TextAt{16}{150.156}{193.924}{\Character{106}}%
\TextAt{13}{158.092}{193.924}{K}%
\TextAt{8}{168.91}{193.924}{,}%
\TextAt{8}{177.008}{193.924}{resp.}%
\TextAt{20}{209.596}{196.946}{\Character{98}}%
\TextAt{13}{208.609}{193.924}{L}%
\TextAt{16}{221.189}{193.924}{\Character{106}}%
\TextAt{13}{229.125}{193.924}{L}%
\TextAt{8}{237.089}{193.924}{,}%
\TextAt{8}{245.188}{193.924}{the}%
\TextAt{18}{266.104}{193.924}{Z}%
\TextAt{21}{274.074}{192.13}{p}%
\TextAt{8}{278.835}{193.924}{-extension}%
\TextAt{8}{335.258}{193.924}{contained}%
\TextAt{8}{389.027}{193.924}{in}%
\TextAt{20}{406.177}{196.946}{\Character{101}}%
\TextAt{13}{403.439}{193.924}{K}%
\TextAt{16}{418.872}{193.924}{\Character{106}}%
\TextAt{13}{426.808}{193.924}{K}%
\TextAt{8}{437.626}{193.924}{,}%
\TextAt{8}{445.724}{193.924}{resp.}%
\TextAt{20}{478.312}{196.946}{\Character{101}}%
\TextAt{13}{477.325}{193.924}{L}%
\TextAt{16}{489.905}{193.924}{\Character{106}}%
\TextAt{13}{497.841}{193.924}{L}%
\TextAt{8}{505.806}{193.924}{,}%
\TextAt{8}{513.904}{193.924}{and}%
\TextAt{8}{79.226}{176.878}{consider}%
\TextAt{8}{124.811}{176.878}{the}%
\TextAt{8}{144.968}{176.878}{field}%
\TextAt{8}{170.325}{176.878}{diagram}%
\TextAt{8}{297.211}{43.033}{357}%
\Destination{78.23}{809.815}{nav.page.357}%
\Destination{79.36}{739.055}{nav.print.392}%
\Hotlink{194.463}{402.255}{13.995}{11.047}{nav.chapter.4}%
\Hotlink{486.767}{636.655}{25.078}{11.047}{nav.number.6.5.3}%
\Hotlink{214.161}{402.255}{25.078}{11.047}{nav.number.4.6.4}%
}

% Page 370 (body page 358)
\Page{%
\TextAt{2}{79.226}{796.943}{VI.}%
\TextAt{2}{98.012}{796.943}{Global}%
\TextAt{2}{133.248}{796.943}{Class}%
\TextAt{2}{161.841}{796.943}{Field}%
\TextAt{2}{189.572}{796.943}{Theory}%
\Vector{q 0 G 0 g 0.299 w 7.226 18.729 m 460.769 18.729 l S Q}%
\NativeDiagram{%
\path[fill=black!0] (153.415,-44.815) -- (165.365,-44.815) -- (165.365,-32.661) -- (153.415,-32.661) -- cycle;%
}%
\TextAt{13}{227.408}{731.177}{L}%
\NativeDiagram{%
\path[fill=black!0] (266.802,-46.559) -- (278.752,-46.559) -- (278.752,-30.917) -- (266.802,-30.917) -- cycle;%
}%
\TextAt{20}{341.782}{732.456}{\Character{98}}%
\TextAt{13}{340.795}{729.434}{L}%
\NativeDiagram{%
\path[fill=black!0] (151.989,-101.509) -- (166.791,-101.509) -- (166.791,-89.354) -- (151.989,-89.354) -- cycle;%
}%
\TextAt{13}{225.981}{674.484}{K}%
\NativeDiagram{%
\path[fill=black!0] (265.376,-103.252) -- (280.179,-103.252) -- (280.179,-87.611) -- (265.376,-87.611) -- cycle;%
}%
\TextAt{20}{342.106}{675.762}{\Character{98}}%
\TextAt{13}{339.368}{672.74}{K}%
\NativeDiagram{%
\path[fill=black!0] (185.188,-18.119) -- (201.624,-18.119) -- (201.624,-2.664) -- (185.188,-2.664) -- cycle;%
}%
\TextAt{13}{259.18}{761.174}{L}%
\TextAt{27}{267.145}{759.38}{p}%
\NativeDiagram{%
\path[fill=black!0] (298.575,-19.862) -- (315.012,-19.862) -- (315.012,-0.92) -- (298.575,-0.92) -- cycle;%
}%
\TextAt{20}{373.555}{762.453}{\Character{98}}%
\TextAt{13}{372.568}{759.431}{L}%
\TextAt{27}{380.533}{757.637}{p}%
\NativeDiagram{%
\path[fill=black!0] (184.19,-74.812) -- (202.622,-74.812) -- (202.622,-59.357) -- (184.19,-59.357) -- cycle;%
}%
\TextAt{13}{258.183}{704.481}{K}%
\TextAt{27}{268.142}{702.687}{p}%
\NativeDiagram{%
\path[fill=black!0] (297.578,-76.556) -- (316.009,-76.556) -- (316.009,-57.614) -- (297.578,-57.614) -- cycle;%
}%
\TextAt{20}{374.308}{705.759}{\Character{98}}%
\TextAt{13}{371.57}{702.737}{K}%
\TextAt{27}{381.529}{700.943}{p}%
\NativeDiagram{%
\path[draw=black,line width=0.399bp] (165.564,-38.738) -- (266.603,-38.738) (272.777,-46.758) -- (272.777,-87.412) (265.177,-95.432) -- (166.991,-95.432) (159.39,-89.155) -- (159.39,-45.015);%
\path[draw=black,line width=0.399bp] (201.824,-10.391) -- (298.376,-10.391) (306.793,-20.061) -- (306.793,-57.415) (297.378,-67.085) -- (202.821,-67.085) (193.406,-59.158) -- (193.406,-18.318);%
\path[draw=black,line width=0.399bp] (165.564,-33.594) -- (184.989,-17.404);%
\path[draw=black,line width=0.399bp] (278.951,-33.594) -- (298.376,-17.404);%
\path[draw=black,line width=0.399bp] (166.92,-89.155) -- (183.991,-74.929);%
\path[draw=black,line width=0.399bp] (280.378,-89.097) -- (297.378,-74.929);%
}%
\TextAt{8}{79.226}{644.736}{with}%
\TextAt{8}{106.998}{644.736}{localizations}%
\TextAt{20}{177.582}{647.758}{\Character{98}}%
\TextAt{13}{174.843}{644.736}{K}%
\TextAt{27}{184.803}{642.943}{p}%
\TextAt{9}{194.505}{644.736}{\Character{61}}%
\TextAt{13}{208.825}{644.736}{K}%
\TextAt{27}{218.784}{642.943}{p}%
\TextAt{20}{226.009}{647.758}{\Character{98}}%
\TextAt{13}{223.271}{644.736}{K}%
\TextAt{8}{234.088}{644.736}{,}%
\TextAt{20}{243.614}{647.758}{\Character{98}}%
\TextAt{13}{242.627}{644.736}{L}%
\TextAt{27}{250.592}{642.943}{p}%
\TextAt{9}{260.294}{644.736}{\Character{61}}%
\TextAt{13}{274.614}{644.736}{L}%
\TextAt{27}{282.579}{642.943}{p}%
\TextAt{20}{288.052}{647.758}{\Character{98}}%
\TextAt{13}{287.065}{644.736}{L}%
\TextAt{8}{300.039}{644.736}{(all}%
\TextAt{8}{321.965}{644.736}{fields}%
\TextAt{8}{353.048}{644.736}{are}%
\TextAt{8}{373.671}{644.736}{considered}%
\TextAt{8}{432.071}{644.736}{to}%
\TextAt{8}{447.494}{644.736}{lie}%
\TextAt{8}{464.207}{644.736}{in}%
\TextAt{8}{478.971}{644.736}{a}%
\TextAt{8}{489.838}{644.736}{common}%
\TextAt{8}{79.226}{627.69}{bigger}%
\TextAt{8}{114.608}{627.69}{field).}%
\TextAt{8}{149.842}{627.69}{We}%
\TextAt{8}{170.243}{627.69}{may}%
\TextAt{8}{195.869}{627.69}{now}%
\TextAt{8}{220.515}{627.69}{lift}%
\TextAt{13}{239.302}{627.69}{\Character{27}}%
\NativeMathAt{250.144}{627.69}{11.955}{\in}%
\TextAt{13}{261.873}{627.69}{G}%
\TextAt{9}{271.107}{627.69}{\Character{40}}%
\TextAt{13}{275.659}{627.69}{L}%
\TextAt{27}{283.623}{625.897}{p}%
\TextAt{16}{291.869}{627.69}{\Character{106}}%
\TextAt{13}{298.949}{627.69}{K}%
\TextAt{27}{308.908}{625.897}{p}%
\TextAt{9}{313.395}{627.69}{\Character{41}}%
\TextAt{9}{321.692}{627.69}{\Character{61}}%
\TextAt{13}{334.57}{627.69}{G}%
\TextAt{9}{343.804}{627.69}{\Character{40}}%
\TextAt{13}{348.356}{627.69}{L}%
\TextAt{16}{360.079}{627.69}{\Character{106}}%
\TextAt{13}{367.159}{627.69}{K}%
\TextAt{9}{377.977}{627.69}{\Character{41}}%
\TextAt{8}{386.684}{627.69}{to}%
\TextAt{8}{401.251}{627.69}{an}%
\TextAt{8}{417.767}{627.69}{automorphism}%
\TextAt{20}{495.714}{627.69}{\Character{98}}%
\TextAt{13}{495.494}{627.69}{\Character{27}}%
\TextAt{8}{506.731}{627.69}{of}%
\TextAt{20}{521.305}{630.712}{\Character{98}}%
\TextAt{13}{520.318}{627.69}{L}%
\TextAt{27}{528.283}{625.897}{p}%
\TextAt{8}{79.226}{610.644}{such}%
\TextAt{8}{105.623}{610.644}{that}%
\TextAt{8}{86.973}{590.609}{(1)}%
\TextAt{20}{108.139}{590.609}{\Character{98}}%
\TextAt{13}{107.918}{590.609}{\Character{27}}%
\NativeMathAt{118.322}{590.609}{11.955}{\in}%
\TextAt{13}{129.613}{590.609}{G}%
\TextAt{9}{138.846}{590.609}{\Character{40}}%
\TextAt{20}{144.385}{593.631}{\Character{98}}%
\TextAt{13}{143.399}{590.609}{L}%
\TextAt{27}{151.363}{588.816}{p}%
\TextAt{16}{159.171}{590.609}{\Character{106}}%
\TextAt{13}{165.813}{590.609}{K}%
\TextAt{27}{175.772}{588.816}{p}%
\TextAt{9}{180.259}{590.609}{\Character{41}}%
\TextAt{8}{184.811}{590.609}{,}%
\TextAt{8}{86.973}{573.563}{(2)}%
\TextAt{20}{108.139}{573.563}{\Character{98}}%
\TextAt{13}{107.918}{573.563}{\Character{27}}%
\TextAt{16}{115.001}{573.563}{\Character{106}}%
\TextAt{30}{120.292}{571.925}{\Character{98}}%
\TextAt{21}{118.322}{569.91}{K}%
\TextAt{9}{129.845}{573.563}{\Character{61}}%
\TextAt{13}{142.271}{573.563}{\Character{39}}%
\TextAt{21}{151.351}{577.902}{n}%
\TextAt{30}{151.91}{570.233}{\Character{98}}%
\TextAt{21}{149.939}{568.218}{K}%
\TextAt{3}{157.644}{568.218}{\Character{106}}%
\TextAt{21}{159.996}{568.218}{K}%
\TextAt{8}{172.096}{573.563}{for}%
\TextAt{8}{189.981}{573.563}{some}%
\TextAt{13}{219.307}{573.563}{n}%
\NativeMathAt{229.615}{573.563}{11.955}{\in}%
\TextAt{18}{240.906}{573.563}{N}%
\TextAt{8}{249.541}{573.563}{.}%
\TextAt{8}{79.226}{550.584}{Indeed,}%
\TextAt{8}{121.51}{550.584}{since}%
\TextAt{20}{153.719}{553.606}{\Character{98}}%
\TextAt{13}{150.981}{550.584}{K}%
\TextAt{27}{160.94}{548.791}{p}%
\TextAt{9}{170.115}{550.584}{\Character{61}}%
\TextAt{13}{183.907}{550.584}{K}%
\TextAt{27}{193.866}{548.791}{p}%
\TextAt{20}{201.092}{553.606}{\Character{98}}%
\TextAt{13}{198.353}{550.584}{K}%
\NativeMathAt{213.859}{550.584}{11.955}{\neq}%
\TextAt{13}{227.651}{550.584}{K}%
\TextAt{27}{237.611}{548.791}{p}%
\TextAt{8}{242.097}{550.584}{,}%
\TextAt{8}{250.249}{550.584}{the}%
\TextAt{8}{271.208}{550.584}{group}%
\TextAt{13}{305.186}{550.584}{G}%
\TextAt{9}{314.419}{550.584}{\Character{40}}%
\TextAt{20}{321.71}{553.606}{\Character{98}}%
\TextAt{13}{318.972}{550.584}{K}%
\TextAt{27}{328.931}{548.791}{p}%
\TextAt{16}{338.106}{550.584}{\Character{106}}%
\TextAt{13}{346.114}{550.584}{K}%
\TextAt{27}{356.074}{548.791}{p}%
\TextAt{9}{360.56}{550.584}{\Character{41}}%
\NativeMathAt{369.801}{550.584}{11.955}{\neq}%
\TextAt{9}{383.585}{550.584}{1}%
\TextAt{8}{389.446}{550.584}{,}%
\TextAt{8}{397.598}{550.584}{and}%
\TextAt{8}{421.163}{550.584}{thus}%
\TextAt{8}{447.717}{550.584}{is}%
\TextAt{8}{460.283}{550.584}{of}%
\TextAt{8}{474.415}{550.584}{finite}%
\TextAt{8}{505.129}{550.584}{index}%
\TextAt{8}{79.226}{533.538}{if}%
\TextAt{8}{90.628}{533.538}{viewed}%
\TextAt{8}{129.67}{533.538}{as}%
\TextAt{8}{144.718}{533.538}{subgroup}%
\TextAt{8}{196.193}{533.538}{of}%
\TextAt{13}{210.201}{533.538}{G}%
\TextAt{9}{219.435}{533.538}{\Character{40}}%
\TextAt{20}{226.726}{536.56}{\Character{98}}%
\TextAt{13}{223.987}{533.538}{K}%
\TextAt{16}{239.281}{533.538}{\Character{106}}%
\TextAt{13}{247.078}{533.538}{K}%
\TextAt{9}{257.896}{533.538}{\Character{41}}%
\NativeMathAt{266.924}{533.538}{11.955}{\cong}%
\TextAt{18}{280.699}{533.538}{Z}%
\TextAt{21}{288.669}{531.745}{p}%
\TextAt{8}{293.43}{533.538}{.}%
\TextAt{8}{303.912}{533.538}{It}%
\TextAt{8}{317.263}{533.538}{is}%
\TextAt{8}{329.705}{533.538}{therefore}%
\TextAt{8}{379.482}{533.538}{generated}%
\TextAt{8}{433.492}{533.538}{by}%
\TextAt{8}{450.429}{533.538}{a}%
\TextAt{8}{460.863}{533.538}{natural}%
\TextAt{8}{502.523}{533.538}{power}%
\TextAt{13}{79.226}{516.492}{\Character{32}}%
\TextAt{9}{91.78}{516.492}{\Character{61}}%
\TextAt{13}{105.4}{516.492}{\Character{39}}%
\TextAt{21}{114.48}{520.831}{k}%
\TextAt{30}{115.039}{513.162}{\Character{98}}%
\TextAt{21}{113.068}{511.147}{K}%
\TextAt{3}{120.773}{511.147}{\Character{106}}%
\TextAt{21}{123.125}{511.147}{K}%
\TextAt{8}{135.926}{516.492}{of}%
\TextAt{8}{149.957}{516.492}{Frobenius}%
\TextAt{13}{204.527}{516.492}{\Character{39}}%
\TextAt{30}{214.166}{514.854}{\Character{98}}%
\TextAt{21}{212.196}{512.839}{K}%
\TextAt{3}{219.9}{512.839}{\Character{106}}%
\TextAt{21}{222.252}{512.839}{K}%
\NativeMathAt{234.97}{516.492}{11.955}{\in}%
\TextAt{13}{247.455}{516.492}{G}%
\TextAt{9}{256.689}{516.492}{\Character{40}}%
\TextAt{20}{263.979}{519.514}{\Character{98}}%
\TextAt{13}{261.241}{516.492}{K}%
\TextAt{16}{276.574}{516.492}{\Character{106}}%
\TextAt{13}{284.409}{516.492}{K}%
\TextAt{9}{295.227}{516.492}{\Character{41}}%
\TextAt{8}{299.779}{516.492}{.}%
\TextAt{8}{310.33}{516.492}{As}%
\TextAt{8}{328.318}{516.492}{in}%
\TextAt{8}{342.671}{516.492}{the}%
\TextAt{8}{363.528}{516.492}{proof}%
\TextAt{8}{394.798}{516.492}{of}%
\TextAt{8}{408.829}{516.492}{chap.}%
\TextAt{8}{440.422}{516.492}{IV,}%
\TextAt{8}{461.267}{516.492}{(4.4),}%
\TextAt{8}{493.37}{516.492}{we}%
\TextAt{8}{511.298}{516.492}{may}%
\TextAt{8}{79.226}{496.502}{then}%
\TextAt{8}{105.886}{496.502}{lift}%
\TextAt{13}{124.417}{496.502}{\Character{27}}%
\TextAt{8}{135.396}{496.502}{to}%
\TextAt{8}{149.707}{496.502}{a}%
\TextAt{20}{159.683}{496.502}{\Character{98}}%
\TextAt{13}{159.462}{496.502}{\Character{27}}%
\NativeMathAt{169.865}{496.502}{11.955}{\in}%
\TextAt{13}{181.156}{496.502}{G}%
\TextAt{9}{190.39}{496.502}{\Character{40}}%
\TextAt{20}{195.929}{499.524}{\Character{98}}%
\TextAt{13}{194.942}{496.502}{L}%
\TextAt{27}{202.907}{494.709}{p}%
\TextAt{16}{210.715}{496.502}{\Character{106}}%
\TextAt{13}{217.356}{496.502}{K}%
\TextAt{27}{227.315}{494.709}{p}%
\TextAt{9}{231.802}{496.502}{\Character{41}}%
\TextAt{8}{240.252}{496.502}{such}%
\TextAt{8}{266.649}{496.502}{that}%
\TextAt{20}{292.238}{496.502}{\Character{98}}%
\TextAt{13}{292.018}{496.502}{\Character{27}}%
\TextAt{16}{299.1}{496.502}{\Character{106}}%
\TextAt{30}{304.392}{494.864}{\Character{98}}%
\TextAt{21}{302.421}{492.849}{K}%
\TextAt{28}{309.563}{491.853}{p}%
\TextAt{9}{317.464}{496.502}{\Character{61}}%
\TextAt{13}{329.89}{496.502}{\Character{32}}%
\TextAt{21}{337.929}{500.84}{m}%
\TextAt{8}{345.918}{496.502}{,}%
\TextAt{13}{353.067}{496.502}{m}%
\NativeMathAt{366.627}{496.502}{11.955}{\in}%
\TextAt{18}{377.918}{496.502}{N}%
\TextAt{8}{386.552}{496.502}{,}%
\TextAt{8}{393.701}{496.502}{so}%
\TextAt{8}{408.071}{496.502}{that}%
\TextAt{20}{433.661}{496.502}{\Character{98}}%
\TextAt{13}{433.44}{496.502}{\Character{27}}%
\TextAt{16}{440.523}{496.502}{\Character{106}}%
\TextAt{30}{445.814}{494.864}{\Character{98}}%
\TextAt{21}{443.844}{492.849}{K}%
\TextAt{9}{455.367}{496.502}{\Character{61}}%
\TextAt{13}{467.793}{496.502}{\Character{39}}%
\TextAt{21}{476.872}{500.84}{k}%
\TextAt{21}{481.487}{500.84}{m}%
\TextAt{30}{477.431}{493.171}{\Character{98}}%
\TextAt{21}{475.461}{491.157}{K}%
\TextAt{3}{483.166}{491.157}{\Character{106}}%
\TextAt{21}{485.518}{491.157}{K}%
\TextAt{8}{493.721}{496.502}{.}%
\TextAt{8}{97.159}{476.467}{We}%
\TextAt{8}{117.671}{476.467}{now}%
\TextAt{8}{142.427}{476.467}{take}%
\TextAt{8}{168.164}{476.467}{the}%
\TextAt{8}{188.688}{476.467}{fixed}%
\TextAt{8}{217.342}{476.467}{field}%
\TextAt{13}{243.066}{476.467}{K}%
\TextAt{3}{253.884}{480.806}{\Character{48}}%
\TextAt{8}{260.944}{476.467}{of}%
\TextAt{20}{274.862}{476.467}{\Character{98}}%
\TextAt{13}{274.642}{476.467}{\Character{27}}%
\TextAt{16}{281.724}{476.467}{\Character{106}}%
\TextAt{30}{285.808}{474.829}{\Character{98}}%
\TextAt{21}{285.045}{472.814}{L}%
\TextAt{8}{291.302}{476.467}{,}%
\TextAt{8}{298.911}{476.467}{and}%
\TextAt{8}{322.042}{476.467}{the}%
\TextAt{8}{342.566}{476.467}{extension}%
\TextAt{13}{394.699}{476.467}{L}%
\TextAt{3}{402.664}{480.806}{\Character{48}}%
\TextAt{9}{409.406}{476.467}{\Character{61}}%
\TextAt{13}{422.458}{476.467}{K}%
\TextAt{3}{433.276}{480.806}{\Character{48}}%
\TextAt{13}{436.071}{476.467}{L}%
\TextAt{8}{444.036}{476.467}{.}%
\TextAt{8}{453.587}{476.467}{As}%
\TextAt{8}{471.242}{476.467}{in}%
\TextAt{8}{485.263}{476.467}{chap.}%
\TextAt{8}{516.522}{476.467}{IV,}%
\TextAt{8}{79.226}{459.421}{(4.5),}%
\TextAt{8}{111.367}{459.421}{conditions}%
\TextAt{8}{168.096}{459.421}{(ii)}%
\TextAt{8}{188.338}{459.421}{and}%
\TextAt{8}{211.831}{459.421}{(iii),}%
\TextAt{8}{238.76}{459.421}{it}%
\TextAt{8}{251.195}{459.421}{then}%
\TextAt{8}{278.586}{459.421}{follows}%
\TextAt{8}{317.753}{459.421}{that}%
\TextAt{9}{343.853}{459.421}{\Character{91}}%
\TextAt{13}{347.104}{459.421}{K}%
\TextAt{3}{357.922}{463.76}{\Character{48}}%
\TextAt{9}{365.284}{459.421}{\Character{58}}%
\TextAt{13}{373.102}{459.421}{K}%
\TextAt{9}{383.92}{459.421}{\Character{93}}%
\TextAt{13}{391.738}{459.421}{\Character{60}}%
\NativeMathAt{405.409}{459.421}{11.955}{\infty}%
\TextAt{8}{421.993}{459.421}{and}%
\TextAt{20}{448.225}{462.443}{\Character{98}}%
\TextAt{13}{445.486}{459.421}{K}%
\TextAt{3}{456.304}{463.76}{\Character{48}}%
\TextAt{9}{463.666}{459.421}{\Character{61}}%
\TextAt{20}{478.324}{462.443}{\Character{98}}%
\TextAt{13}{477.337}{459.421}{L}%
\TextAt{8}{485.302}{459.421}{.}%
\TextAt{13}{495.943}{459.421}{L}%
\TextAt{3}{503.908}{463.76}{\Character{48}}%
\TextAt{16}{511.269}{459.421}{\Character{106}}%
\TextAt{13}{519.157}{459.421}{K}%
\TextAt{3}{529.974}{463.76}{\Character{48}}%
\TextAt{8}{79.226}{442.375}{is}%
\TextAt{8}{91.304}{442.375}{therefore}%
\TextAt{8}{140.718}{442.375}{a}%
\TextAt{8}{150.787}{442.375}{subextension}%
\TextAt{8}{220.812}{442.375}{of}%
\TextAt{20}{237.194}{445.397}{\Character{98}}%
\TextAt{13}{234.455}{442.375}{K}%
\TextAt{3}{245.273}{446.714}{\Character{48}}%
\TextAt{16}{251.924}{442.375}{\Character{106}}%
\TextAt{13}{259.101}{442.375}{K}%
\TextAt{3}{269.918}{446.714}{\Character{48}}%
\TextAt{8}{272.713}{442.375}{,}%
\TextAt{8}{280.255}{442.375}{and}%
\TextAt{13}{303.332}{442.375}{\Character{27}}%
\TextAt{8}{314.625}{442.375}{is}%
\TextAt{8}{326.703}{442.375}{the}%
\TextAt{8}{347.173}{442.375}{image}%
\TextAt{8}{381.308}{442.375}{of}%
\TextAt{20}{395.173}{442.375}{\Character{98}}%
\TextAt{13}{394.952}{442.375}{\Character{27}}%
\TextAt{16}{402.035}{442.375}{\Character{106}}%
\TextAt{21}{405.356}{440.427}{L}%
\TextAt{23}{411.115}{443.172}{\Character{48}}%
\TextAt{28}{411.115}{438.603}{p}%
\TextAt{8}{419.907}{442.375}{under}%
\TextAt{13}{453.384}{442.375}{G}%
\TextAt{9}{462.618}{442.375}{\Character{40}}%
\TextAt{13}{467.17}{442.375}{L}%
\TextAt{3}{475.135}{446.714}{\Character{48}}%
\TextAt{27}{475.135}{439.42}{p}%
\TextAt{16}{483.478}{442.375}{\Character{106}}%
\TextAt{13}{490.654}{442.375}{K}%
\TextAt{3}{501.472}{446.714}{\Character{48}}%
\TextAt{27}{500.613}{439.42}{p}%
\TextAt{9}{505.1}{442.375}{\Character{41}}%
\NativeArrowAt{513.508}{442.375}{11.955}{19.26}{\longrightarrow}%
\TextAt{13}{79.226}{425.329}{G}%
\TextAt{9}{88.46}{425.329}{\Character{40}}%
\TextAt{13}{93.012}{425.329}{L}%
\TextAt{27}{100.977}{423.536}{p}%
\TextAt{16}{108.784}{425.329}{\Character{106}}%
\TextAt{13}{115.426}{425.329}{K}%
\TextAt{27}{125.385}{423.536}{p}%
\TextAt{9}{129.872}{425.329}{\Character{41}}%
\TextAt{8}{134.424}{425.329}{.}%
\TextAt{8}{142.873}{425.329}{This}%
\TextAt{8}{169.592}{425.329}{finishes}%
\TextAt{8}{210.682}{425.329}{the}%
\TextAt{8}{230.839}{425.329}{proof.}%
\Vector{q 0 G 0 g 0.399 w 452.334 -348.671 m 452.334 -340.601 l S Q}%
\Vector{q 0 G 0 g 0.399 w 452.534 -340.8 m 459.707 -340.8 l S Q}%
\Vector{q 0 G 0 g 0.399 w 452.534 -348.472 m 459.707 -348.472 l S Q}%
\Vector{q 0 G 0 g 0.399 w 459.906 -348.671 m 459.906 -340.601 l S Q}%
\TextAt{12}{79.226}{390.35}{Corollary}%
\TextAt{12}{139.014}{390.35}{(5.7).}%
\TextAt{8}{174.33}{390.35}{If}%
\TextAt{13}{186.022}{390.35}{L}%
\TextAt{16}{197.307}{390.35}{\Character{106}}%
\TextAt{13}{203.949}{390.35}{K}%
\TextAt{8}{218.664}{390.35}{is}%
\TextAt{8}{230.428}{390.35}{an}%
\TextAt{8}{246.687}{390.35}{abelian}%
\TextAt{8}{287.334}{390.35}{extension}%
\TextAt{8}{339.1}{390.35}{and}%
\TextAt{13}{361.863}{390.35}{\Character{11}}%
\TextAt{9}{372.705}{390.35}{\Character{61}}%
\TextAt{9}{385.115}{390.35}{\Character{40}}%
\TextAt{13}{389.683}{390.35}{\Character{11}}%
\TextAt{27}{397.129}{388.557}{p}%
\TextAt{9}{401.616}{390.35}{\Character{41}}%
\NativeMathAt{409.489}{390.35}{11.955}{\in}%
\TextAt{13}{420.78}{390.35}{I}%
\TextAt{21}{425.943}{388.557}{K}%
\TextAt{8}{434.145}{390.35}{,}%
\TextAt{8}{441.295}{390.35}{then}%
\TextAt{9}{231.983}{361.349}{\Character{40}}%
\TextAt{13}{236.535}{361.349}{\Character{11}\Character{59}}%
\TextAt{13}{249.279}{361.349}{L}%
\TextAt{16}{260.586}{361.349}{\Character{106}}%
\TextAt{13}{267.228}{361.349}{K}%
\TextAt{9}{278.046}{361.349}{\Character{41}}%
\TextAt{9}{285.901}{361.349}{\Character{61}}%
\TextAt{20}{298.344}{372.707}{\Character{89}}%
\TextAt{27}{303.988}{347.601}{p}%
\TextAt{9}{313.62}{361.349}{\Character{40}}%
\TextAt{13}{318.173}{361.349}{\Character{11}}%
\TextAt{27}{325.618}{359.556}{p}%
\TextAt{13}{330.105}{361.349}{\Character{59}}%
\TextAt{13}{335.329}{361.349}{L}%
\TextAt{27}{343.314}{359.556}{p}%
\TextAt{16}{351.121}{361.349}{\Character{106}}%
\TextAt{13}{357.763}{361.349}{K}%
\TextAt{27}{367.722}{359.556}{p}%
\TextAt{9}{372.209}{361.349}{\Character{41}}%
\TextAt{13}{376.761}{361.349}{\Character{58}}%
\TextAt{8}{79.226}{322.981}{In}%
\TextAt{8}{92.878}{322.981}{particular,}%
\TextAt{8}{149.344}{322.981}{for}%
\TextAt{8}{166.26}{322.981}{a}%
\TextAt{8}{175.046}{322.981}{principal}%
\TextAt{8}{222.853}{322.981}{idèle}%
\TextAt{13}{249.19}{322.981}{a}%
\NativeMathAt{258.655}{322.981}{11.955}{\in}%
\TextAt{13}{269.946}{322.981}{K}%
\TextAt{3}{280.764}{327.319}{\Character{3}}%
\TextAt{8}{288.424}{322.981}{we}%
\TextAt{8}{304.683}{322.981}{have}%
\TextAt{8}{330.708}{322.981}{the}%
\TextAt{8}{349.895}{322.981}{product}%
\TextAt{8}{392.825}{322.981}{formula}%
\TextAt{20}{434.787}{331.947}{\Character{81}}%
\TextAt{27}{446.078}{319.494}{p}%
\TextAt{9}{450.565}{322.981}{\Character{40}}%
\TextAt{13}{455.117}{322.981}{a\Character{59}}%
\TextAt{13}{466.486}{322.981}{L}%
\TextAt{27}{474.471}{321.187}{p}%
\TextAt{16}{482.278}{322.981}{\Character{106}}%
\TextAt{13}{488.92}{322.981}{K}%
\TextAt{27}{498.879}{321.187}{p}%
\TextAt{9}{503.366}{322.981}{\Character{41}}%
\TextAt{9}{511.221}{322.981}{\Character{61}}%
\TextAt{9}{523.642}{322.981}{1}%
\TextAt{8}{529.518}{322.981}{.}%
\TextAt{19}{79.226}{288.002}{Proof.}%
\TextAt{8}{116}{288.002}{Since}%
\TextAt{13}{146.299}{288.002}{I}%
\TextAt{21}{151.461}{286.209}{K}%
\TextAt{8}{163.302}{288.002}{is}%
\TextAt{8}{174.807}{288.002}{topologically}%
\TextAt{8}{243.505}{288.002}{generated}%
\TextAt{8}{296.578}{288.002}{by}%
\TextAt{8}{312.578}{288.002}{the}%
\TextAt{8}{332.475}{288.002}{idèles}%
\TextAt{8}{364.136}{288.002}{of}%
\TextAt{8}{377.207}{288.002}{the}%
\TextAt{8}{397.105}{288.002}{form}%
\TextAt{13}{424.486}{288.002}{\Character{11}}%
\TextAt{9}{435.328}{288.002}{\Character{61}}%
\TextAt{16}{447.754}{288.002}{\Character{104}}%
\TextAt{13}{452.403}{288.002}{\Character{11}}%
\TextAt{27}{459.849}{286.209}{p}%
\TextAt{16}{464.336}{288.002}{\Character{105}}%
\TextAt{8}{468.985}{288.002}{,}%
\TextAt{13}{475.927}{288.002}{\Character{11}}%
\TextAt{27}{483.372}{286.209}{p}%
\NativeMathAt{491.18}{288.002}{11.955}{\in}%
\TextAt{13}{502.471}{288.002}{K}%
\TextAt{3}{513.289}{292.34}{\Character{3}}%
\TextAt{27}{512.43}{285.046}{p}%
\TextAt{8}{518.021}{288.002}{,}%
\TextAt{8}{524.963}{288.002}{it}%
\TextAt{8}{79.226}{270.956}{is}%
\TextAt{8}{91.098}{270.956}{enough}%
\TextAt{8}{131.53}{270.956}{to}%
\TextAt{8}{145.949}{270.956}{prove}%
\TextAt{8}{177.606}{270.956}{the}%
\TextAt{8}{197.87}{270.956}{first}%
\TextAt{8}{222.104}{270.956}{formula}%
\TextAt{8}{265.143}{270.956}{for}%
\TextAt{8}{283.135}{270.956}{these}%
\TextAt{8}{313.215}{270.956}{idèles.}%
\TextAt{8}{350.009}{270.956}{But}%
\TextAt{8}{373.358}{270.956}{this}%
\TextAt{8}{396.288}{270.956}{is}%
\TextAt{8}{408.16}{270.956}{exactly}%
\TextAt{8}{448.593}{270.956}{the}%
\TextAt{8}{468.857}{270.956}{statement}%
\TextAt{8}{523.337}{270.956}{of}%
\TextAt{8}{79.226}{253.91}{(5.6):}%
\TextAt{9}{152.639}{236.864}{\Character{40}}%
\TextAt{13}{157.191}{236.864}{\Character{11}\Character{59}}%
\TextAt{13}{169.935}{236.864}{L}%
\TextAt{16}{181.243}{236.864}{\Character{106}}%
\TextAt{13}{187.884}{236.864}{K}%
\TextAt{9}{198.702}{236.864}{\Character{41}}%
\TextAt{9}{206.557}{236.864}{\Character{61}}%
\TextAt{9}{218.978}{236.864}{\Character{40}}%
\TextAt{16}{223.553}{236.864}{\Character{104}}%
\TextAt{13}{228.202}{236.864}{\Character{11}}%
\TextAt{27}{235.648}{235.071}{p}%
\TextAt{16}{240.135}{236.864}{\Character{105}}%
\TextAt{13}{244.784}{236.864}{\Character{59}}%
\TextAt{13}{250.008}{236.864}{L}%
\TextAt{16}{261.313}{236.864}{\Character{106}}%
\TextAt{13}{267.955}{236.864}{K}%
\TextAt{9}{278.773}{236.864}{\Character{41}}%
\TextAt{9}{286.628}{236.864}{\Character{61}}%
\TextAt{9}{299.049}{236.864}{\Character{40}}%
\TextAt{13}{303.624}{236.864}{\Character{11}}%
\TextAt{27}{311.069}{235.071}{p}%
\TextAt{13}{315.556}{236.864}{\Character{59}}%
\TextAt{13}{320.78}{236.864}{L}%
\TextAt{27}{328.765}{235.071}{p}%
\TextAt{16}{336.573}{236.864}{\Character{106}}%
\TextAt{13}{343.214}{236.864}{K}%
\TextAt{27}{353.173}{235.071}{p}%
\TextAt{9}{357.66}{236.864}{\Character{41}}%
\TextAt{9}{365.515}{236.864}{\Character{61}}%
\TextAt{20}{377.959}{248.221}{\Character{89}}%
\TextAt{27}{383.648}{223.115}{q}%
\TextAt{9}{393.235}{236.864}{\Character{40}}%
\TextAt{13}{397.787}{236.864}{\Character{11}}%
\TextAt{27}{405.233}{235.071}{q}%
\TextAt{13}{409.629}{236.864}{\Character{59}}%
\TextAt{13}{414.853}{236.864}{L}%
\TextAt{27}{422.838}{235.071}{q}%
\TextAt{16}{430.555}{236.864}{\Character{106}}%
\TextAt{13}{437.197}{236.864}{K}%
\TextAt{27}{447.156}{235.071}{q}%
\TextAt{9}{451.553}{236.864}{\Character{41}}%
\TextAt{13}{456.105}{236.864}{\Character{58}}%
\TextAt{8}{79.226}{203.975}{The}%
\TextAt{8}{104.151}{203.975}{product}%
\TextAt{8}{148.921}{203.975}{formula}%
\TextAt{8}{192.723}{203.975}{is}%
\TextAt{8}{205.358}{203.975}{a}%
\TextAt{8}{215.985}{203.975}{consequence}%
\TextAt{8}{282.92}{203.975}{of}%
\TextAt{8}{297.121}{203.975}{the}%
\TextAt{8}{318.148}{203.975}{fact}%
\TextAt{8}{342.105}{203.975}{that}%
\TextAt{9}{368.344}{203.975}{\Character{40}}%
\TextAt{13}{372.897}{203.975}{\Character{11}\Character{59}}%
\TextAt{13}{385.641}{203.975}{L}%
\TextAt{16}{398.432}{203.975}{\Character{106}}%
\TextAt{13}{406.558}{203.975}{K}%
\TextAt{9}{417.376}{203.975}{\Character{41}}%
\TextAt{8}{426.697}{203.975}{depends}%
\TextAt{8}{472.818}{203.975}{only}%
\TextAt{8}{499.38}{203.975}{on}%
\TextAt{8}{516.51}{203.975}{the}%
\TextAt{8}{79.226}{186.929}{idèle}%
\TextAt{8}{106.532}{186.929}{class}%
\TextAt{13}{133.969}{186.929}{\Character{11}}%
\TextAt{8}{144.811}{186.929}{mod}%
\TextAt{13}{170.249}{186.929}{K}%
\TextAt{3}{181.067}{191.267}{\Character{3}}%
\TextAt{8}{185.799}{186.929}{.}%
\Vector{q 0 G 0 g 0.399 w 452.334 -587.071 m 452.334 -579.001 l S Q}%
\Vector{q 0 G 0 g 0.399 w 452.534 -579.201 m 459.707 -579.201 l S Q}%
\Vector{q 0 G 0 g 0.399 w 452.534 -586.872 m 459.707 -586.872 l S Q}%
\Vector{q 0 G 0 g 0.399 w 459.906 -587.071 m 459.906 -579.001 l S Q}%
\TextAt{8}{97.159}{157.928}{Identifying}%
\TextAt{13}{155.69}{157.928}{K}%
\TextAt{3}{166.507}{162.266}{\Character{3}}%
\TextAt{27}{165.649}{154.972}{p}%
\TextAt{8}{174.49}{157.928}{with}%
\TextAt{8}{200.502}{157.928}{its}%
\TextAt{8}{216.174}{157.928}{image}%
\TextAt{8}{249.347}{157.928}{in}%
\TextAt{13}{262.353}{157.928}{C}%
\TextAt{21}{270.728}{156.134}{K}%
\TextAt{8}{282.181}{157.928}{under}%
\TextAt{8}{314.697}{157.928}{the}%
\TextAt{8}{334.206}{157.928}{map}%
\TextAt{13}{359.573}{157.928}{a}%
\TextAt{27}{365.718}{156.134}{p}%
\NativeArrowAt{373.526}{157.928}{11.955}{11.955}{\mapsto}%
\TextAt{16}{388.793}{157.928}{\Character{104}}%
\TextAt{13}{393.451}{157.928}{a}%
\TextAt{27}{399.596}{156.134}{p}%
\TextAt{16}{404.083}{157.928}{\Character{105}}%
\TextAt{8}{408.732}{157.928}{,}%
\TextAt{8}{415.363}{157.928}{we}%
\TextAt{8}{431.943}{157.928}{obtain}%
\TextAt{8}{467.723}{157.928}{the}%
\TextAt{8}{487.232}{157.928}{following}%
\TextAt{8}{79.226}{140.882}{further}%
\TextAt{8}{118.571}{140.882}{corollary,}%
\TextAt{8}{169.308}{140.882}{where}%
\TextAt{8}{203.118}{140.882}{we}%
\TextAt{8}{220.345}{140.882}{use}%
\TextAt{8}{240.561}{140.882}{the}%
\TextAt{8}{260.718}{140.882}{abbreviations}%
\TextAt{13}{333.309}{140.882}{N}%
\TextAt{9}{347.253}{140.882}{\Character{61}}%
\TextAt{13}{359.678}{140.882}{N}%
\TextAt{21}{369.052}{139.022}{L}%
\TextAt{3}{374.811}{139.022}{\Character{106}}%
\TextAt{21}{377.164}{139.022}{K}%
\TextAt{8}{389.264}{140.882}{and}%
\TextAt{13}{412.026}{140.882}{N}%
\TextAt{27}{421.4}{139.088}{p}%
\TextAt{9}{429.208}{140.882}{\Character{61}}%
\TextAt{13}{441.633}{140.882}{N}%
\TextAt{21}{451.007}{139.022}{L}%
\TextAt{28}{456.766}{138.026}{p}%
\TextAt{3}{460.849}{139.022}{\Character{106}}%
\TextAt{21}{463.201}{139.022}{K}%
\TextAt{28}{470.342}{138.026}{p}%
\TextAt{8}{474.923}{140.882}{.}%
\TextAt{12}{79.226}{114.869}{Corollary}%
\TextAt{12}{139.014}{114.869}{(5.8).}%
\TextAt{8}{174.33}{114.869}{For}%
\TextAt{8}{195.299}{114.869}{every}%
\TextAt{8}{226.191}{114.869}{finite}%
\TextAt{8}{256.103}{114.869}{abelian}%
\TextAt{8}{296.75}{114.869}{extension}%
\TextAt{8}{348.516}{114.869}{one}%
\TextAt{8}{369.976}{114.869}{has}%
\TextAt{13}{390.85}{114.869}{N}%
\TextAt{13}{401.466}{114.869}{C}%
\TextAt{21}{409.847}{113.076}{L}%
\TextAt{16}{418.761}{114.869}{\Character{92}}%
\TextAt{13}{429.388}{114.869}{K}%
\TextAt{3}{440.206}{119.208}{\Character{3}}%
\TextAt{27}{439.347}{111.914}{p}%
\TextAt{9}{448.259}{114.869}{\Character{61}}%
\TextAt{13}{460.684}{114.869}{N}%
\TextAt{27}{470.058}{113.076}{p}%
\TextAt{13}{474.545}{114.869}{L}%
\TextAt{3}{482.509}{119.208}{\Character{3}}%
\TextAt{27}{482.509}{111.914}{p}%
\TextAt{8}{487.242}{114.869}{.}%
\TextAt{8}{297.211}{43.033}{358}%
\Destination{79.23}{607.535}{nav.Item.127}%
\Destination{79.23}{587.305}{nav.Item.128}%
\Destination{78.23}{809.815}{nav.page.358}%
\Destination{79.36}{642.655}{nav.print.393}%
\Destination{79.36}{404.355}{nav.number.6.5.7}%
\Destination{79.36}{128.855}{nav.number.6.5.8}%
\Hotlink{439.922}{514.255}{13.995}{11.047}{nav.chapter.4}%
\Hotlink{516.022}{474.255}{13.995}{11.047}{nav.chapter.4}%
\Hotlink{460.767}{514.255}{25.078}{11.047}{nav.number.4.4.4}%
\Hotlink{78.726}{457.255}{25.078}{11.047}{nav.number.6.4.5}%
\Hotlink{78.726}{251.755}{25.078}{11.047}{nav.number.6.5.6}%
}

% Page 371 (body page 359)
\Page{%
\TextAt{2}{79.226}{796.943}{VI.}%
\TextAt{2}{98.012}{796.943}{Global}%
\TextAt{2}{133.248}{796.943}{Class}%
\TextAt{2}{161.841}{796.943}{Field}%
\TextAt{2}{189.572}{796.943}{Theory}%
\Vector{q 0 G 0 g 0.299 w 7.226 18.729 m 460.769 18.729 l S Q}%
\TextAt{19}{79.226}{761.124}{Proof.}%
\TextAt{8}{116}{761.124}{For}%
\TextAt{13}{137.144}{761.124}{x}%
\TextAt{27}{143.796}{759.33}{p}%
\NativeMathAt{151.901}{761.124}{11.955}{\in}%
\TextAt{13}{163.489}{761.124}{N}%
\TextAt{27}{172.863}{759.33}{p}%
\TextAt{13}{177.35}{761.124}{L}%
\TextAt{3}{185.314}{765.462}{\Character{3}}%
\TextAt{27}{185.314}{758.168}{p}%
\TextAt{8}{194.118}{761.124}{we}%
\TextAt{8}{211.52}{761.124}{see}%
\TextAt{8}{230.607}{761.124}{from}%
\TextAt{8}{258.422}{761.124}{(5.6)}%
\TextAt{8}{286.572}{761.124}{that}%
\TextAt{9}{312.115}{761.124}{\Character{40}}%
\TextAt{16}{316.667}{761.124}{\Character{104}}%
\TextAt{13}{321.316}{761.124}{x}%
\TextAt{27}{327.968}{759.33}{p}%
\TextAt{16}{332.455}{761.124}{\Character{105}}%
\TextAt{13}{337.105}{761.124}{\Character{59}}%
\TextAt{13}{342.329}{761.124}{L}%
\TextAt{16}{353.931}{761.124}{\Character{106}}%
\TextAt{13}{360.87}{761.124}{K}%
\TextAt{9}{371.688}{761.124}{\Character{41}}%
\TextAt{9}{379.841}{761.124}{\Character{61}}%
\TextAt{9}{392.562}{761.124}{\Character{40}}%
\TextAt{13}{397.133}{761.124}{x}%
\TextAt{27}{403.785}{759.33}{p}%
\TextAt{13}{408.272}{761.124}{\Character{59}}%
\TextAt{13}{413.496}{761.124}{L}%
\TextAt{27}{421.48}{759.33}{p}%
\TextAt{16}{429.585}{761.124}{\Character{106}}%
\TextAt{13}{436.524}{761.124}{K}%
\TextAt{27}{446.483}{759.33}{p}%
\TextAt{9}{450.97}{761.124}{\Character{41}}%
\TextAt{9}{459.123}{761.124}{\Character{61}}%
\TextAt{9}{471.844}{761.124}{1}%
\TextAt{8}{477.716}{761.124}{.}%
\TextAt{8}{486.687}{761.124}{Thus}%
\TextAt{8}{516.51}{761.124}{the}%
\TextAt{8}{79.226}{744.078}{class}%
\TextAt{8}{105.868}{744.078}{of}%
\TextAt{16}{118.402}{744.078}{\Character{104}}%
\TextAt{13}{123.051}{744.078}{x}%
\TextAt{27}{129.703}{742.284}{p}%
\TextAt{16}{134.19}{744.078}{\Character{105}}%
\TextAt{8}{141.941}{744.078}{is}%
\TextAt{8}{152.91}{744.078}{contained}%
\TextAt{8}{205.123}{744.078}{in}%
\TextAt{13}{217.981}{744.078}{N}%
\TextAt{13}{228.597}{744.078}{C}%
\TextAt{21}{236.978}{742.284}{L}%
\TextAt{8}{243.235}{744.078}{.}%
\TextAt{8}{251.419}{744.078}{Therefore}%
\TextAt{13}{303.62}{744.078}{N}%
\TextAt{27}{312.994}{742.284}{p}%
\TextAt{13}{317.481}{744.078}{L}%
\TextAt{3}{325.445}{748.416}{\Character{3}}%
\TextAt{27}{325.445}{741.122}{p}%
\NativeMathAt{333.499}{744.078}{11.955}{\subseteq}%
\TextAt{13}{346.118}{744.078}{N}%
\TextAt{13}{356.734}{744.078}{C}%
\TextAt{21}{365.116}{742.284}{L}%
\TextAt{8}{371.373}{744.078}{.}%
\TextAt{8}{379.556}{744.078}{Conversely,}%
\TextAt{8}{440.441}{744.078}{let}%
\Vector{q 0 G 0 g 0.478 w 384.55 -23.101 m 392.071 -23.101 l S Q}%
\TextAt{13}{456.55}{744.078}{\Character{11}}%
\NativeMathAt{467.392}{744.078}{11.955}{\in}%
\TextAt{13}{478.683}{744.078}{N}%
\TextAt{13}{489.299}{744.078}{C}%
\TextAt{21}{497.681}{742.284}{L}%
\TextAt{16}{504.968}{744.078}{\Character{92}}%
\TextAt{13}{513.968}{744.078}{K}%
\TextAt{3}{524.785}{748.416}{\Character{3}}%
\TextAt{27}{523.927}{741.122}{p}%
\TextAt{8}{529.518}{744.078}{.}%
\TextAt{8}{79.226}{727.032}{Then}%
\Vector{q 0 G 0 g 0.478 w 38.751 -40.147 m 46.273 -40.147 l S Q}%
\TextAt{13}{110.751}{727.032}{\Character{11}}%
\TextAt{8}{123.138}{727.032}{is}%
\TextAt{8}{135.869}{727.032}{represented}%
\TextAt{8}{199.004}{727.032}{on}%
\TextAt{8}{216.23}{727.032}{the}%
\TextAt{8}{237.355}{727.032}{one}%
\TextAt{8}{259.782}{727.032}{hand}%
\TextAt{8}{290.016}{727.032}{by}%
\TextAt{8}{307.242}{727.032}{a}%
\TextAt{8}{317.965}{727.032}{norm}%
\TextAt{8}{349.502}{727.032}{idèle}%
\TextAt{13}{377.775}{727.032}{\Character{11}}%
\TextAt{9}{390.267}{727.032}{\Character{61}}%
\TextAt{13}{404.341}{727.032}{N}%
\TextAt{13}{414.957}{727.032}{\Character{12}}%
\TextAt{8}{422.235}{727.032}{,}%
\TextAt{13}{430.594}{727.032}{\Character{12}}%
\NativeMathAt{442.835}{727.032}{11.955}{\in}%
\TextAt{13}{455.775}{727.032}{I}%
\TextAt{21}{460.938}{725.238}{L}%
\TextAt{8}{467.195}{727.032}{,}%
\TextAt{8}{475.554}{727.032}{and}%
\TextAt{8}{499.284}{727.032}{on}%
\TextAt{8}{516.51}{727.032}{the}%
\TextAt{8}{79.226}{709.986}{other}%
\TextAt{8}{110.267}{709.986}{hand}%
\TextAt{8}{140.005}{709.986}{by}%
\TextAt{8}{156.736}{709.986}{an}%
\TextAt{8}{173.467}{709.986}{idèle}%
\TextAt{16}{201.244}{709.986}{\Character{104}}%
\TextAt{13}{205.894}{709.986}{x}%
\TextAt{27}{212.546}{708.193}{p}%
\TextAt{16}{217.033}{709.986}{\Character{105}}%
\TextAt{8}{221.682}{709.986}{,}%
\TextAt{13}{229.421}{709.986}{x}%
\TextAt{27}{236.073}{708.193}{p}%
\NativeMathAt{244.685}{709.986}{11.955}{\in}%
\TextAt{13}{256.78}{709.986}{K}%
\TextAt{3}{267.597}{714.324}{\Character{3}}%
\TextAt{27}{266.739}{707.03}{p}%
\TextAt{8}{272.33}{709.986}{.}%
\TextAt{8}{282.193}{709.986}{This}%
\TextAt{8}{309.385}{709.986}{gives}%
\TextAt{16}{338.537}{709.986}{\Character{104}}%
\TextAt{13}{343.186}{709.986}{x}%
\TextAt{27}{349.838}{708.193}{p}%
\TextAt{16}{354.325}{709.986}{\Character{105}}%
\TextAt{13}{358.975}{709.986}{a}%
\TextAt{9}{369.244}{709.986}{\Character{61}}%
\TextAt{13}{382.474}{709.986}{N}%
\TextAt{13}{393.09}{709.986}{\Character{12}}%
\TextAt{8}{404.737}{709.986}{with}%
\TextAt{13}{431.869}{709.986}{a}%
\NativeMathAt{442.139}{709.986}{11.955}{\in}%
\TextAt{13}{454.233}{709.986}{K}%
\TextAt{3}{465.051}{714.324}{\Character{3}}%
\TextAt{8}{469.783}{709.986}{.}%
\TextAt{8}{479.647}{709.986}{Passing}%
\TextAt{8}{522.357}{709.986}{to}%
\TextAt{8}{79.226}{692.94}{components}%
\TextAt{8}{143.028}{692.94}{shows}%
\TextAt{8}{175.998}{692.94}{that}%
\TextAt{13}{200.718}{692.94}{a}%
\TextAt{8}{210.112}{692.94}{is}%
\TextAt{8}{221.227}{692.94}{a}%
\TextAt{8}{230.334}{692.94}{norm}%
\TextAt{8}{260.255}{692.94}{from}%
\TextAt{13}{287.247}{692.94}{L}%
\TextAt{27}{295.212}{691.147}{q}%
\TextAt{16}{302.929}{692.94}{\Character{106}}%
\TextAt{13}{309.571}{692.94}{K}%
\TextAt{27}{319.53}{691.147}{q}%
\TextAt{8}{327.175}{692.94}{for}%
\TextAt{8}{344.412}{692.94}{every}%
\TextAt{14}{374.655}{692.94}{q}%
\NativeMathAt{383.824}{692.94}{11.955}{\neq}%
\TextAt{14}{396.249}{692.94}{p}%
\TextAt{8}{402.232}{692.94}{,}%
\TextAt{8}{408.863}{692.94}{and}%
\TextAt{8}{430.977}{692.94}{the}%
\TextAt{8}{450.485}{692.94}{product}%
\TextAt{8}{493.736}{692.94}{formula}%
\TextAt{8}{79.226}{675.894}{(5.7)}%
\TextAt{8}{107.029}{675.894}{shows}%
\TextAt{8}{140.474}{675.894}{that}%
\TextAt{13}{165.67}{675.894}{a}%
\TextAt{8}{175.539}{675.894}{is}%
\TextAt{8}{187.131}{675.894}{also}%
\TextAt{8}{210.438}{675.894}{a}%
\TextAt{8}{220.021}{675.894}{norm}%
\TextAt{8}{250.417}{675.894}{from}%
\TextAt{13}{277.885}{675.894}{L}%
\TextAt{27}{285.849}{674.101}{p}%
\TextAt{16}{293.657}{675.894}{\Character{106}}%
\TextAt{13}{300.299}{675.894}{K}%
\TextAt{27}{310.258}{674.101}{p}%
\TextAt{8}{314.745}{675.894}{.}%
\TextAt{8}{323.135}{675.894}{Therefore}%
\TextAt{13}{375.96}{675.894}{x}%
\TextAt{27}{382.612}{674.101}{p}%
\NativeMathAt{390.42}{675.894}{11.955}{\in}%
\TextAt{13}{401.711}{675.894}{N}%
\TextAt{27}{411.084}{674.101}{p}%
\TextAt{13}{415.571}{675.894}{L}%
\TextAt{3}{423.536}{680.232}{\Character{3}}%
\TextAt{27}{423.536}{672.938}{p}%
\TextAt{8}{428.268}{675.894}{,}%
\TextAt{8}{435.279}{675.894}{and}%
\TextAt{8}{457.869}{675.894}{this}%
\TextAt{8}{480.519}{675.894}{proves}%
\TextAt{8}{516.51}{675.894}{the}%
\TextAt{8}{79.226}{658.848}{inclusion}%
\TextAt{13}{128.063}{658.848}{N}%
\TextAt{13}{138.679}{658.848}{C}%
\TextAt{21}{147.061}{657.055}{L}%
\TextAt{16}{155.975}{658.848}{\Character{92}}%
\TextAt{13}{166.601}{658.848}{K}%
\TextAt{3}{177.419}{663.186}{\Character{3}}%
\TextAt{27}{176.56}{655.892}{p}%
\NativeMathAt{185.472}{658.848}{11.955}{\subseteq}%
\TextAt{13}{198.091}{658.848}{N}%
\TextAt{27}{207.465}{657.055}{p}%
\TextAt{13}{211.952}{658.848}{L}%
\TextAt{3}{219.917}{663.186}{\Character{3}}%
\TextAt{27}{219.917}{655.892}{p}%
\TextAt{8}{224.649}{658.848}{.}%
\Vector{q 0 G 0 g 0.399 w 452.334 -115.152 m 452.334 -107.082 l S Q}%
\Vector{q 0 G 0 g 0.399 w 452.534 -107.282 m 459.707 -107.282 l S Q}%
\Vector{q 0 G 0 g 0.399 w 452.534 -114.953 m 459.707 -114.953 l S Q}%
\Vector{q 0 G 0 g 0.399 w 459.906 -115.152 m 459.906 -107.082 l S Q}%
\TextAt{12}{79.226}{627.07}{Exercise}%
\TextAt{12}{133.07}{627.07}{VI.5.1.}%
\TextAt{8}{177.385}{627.07}{If}%
\TextAt{13}{189.636}{627.07}{D}%
\TextAt{21}{199.355}{625.276}{K}%
\TextAt{8}{212.015}{627.07}{is}%
\TextAt{8}{224.338}{627.07}{the}%
\TextAt{8}{245.053}{627.07}{connected}%
\TextAt{8}{300.235}{627.07}{component}%
\TextAt{8}{360.63}{627.07}{of}%
\TextAt{8}{374.519}{627.07}{the}%
\TextAt{8}{395.235}{627.07}{unit}%
\TextAt{8}{420.505}{627.07}{element}%
\TextAt{8}{464.306}{627.07}{of}%
\TextAt{13}{478.195}{627.07}{C}%
\TextAt{21}{486.571}{625.276}{K}%
\TextAt{8}{494.773}{627.07}{,}%
\TextAt{8}{502.621}{627.07}{and}%
\TextAt{8}{525.943}{627.07}{if}%
\TextAt{13}{79.226}{610.024}{K}%
\TextAt{21}{90.044}{614.362}{ab}%
\TextAt{16}{101.983}{610.024}{\Character{106}}%
\TextAt{13}{108.625}{610.024}{K}%
\TextAt{8}{123.34}{610.024}{is}%
\TextAt{8}{135.104}{610.024}{the}%
\TextAt{8}{155.26}{610.024}{maximal}%
\TextAt{8}{203.069}{610.024}{abelian}%
\TextAt{8}{243.717}{610.024}{extension}%
\TextAt{8}{295.483}{610.024}{of}%
\TextAt{13}{308.813}{610.024}{K}%
\TextAt{8}{319.63}{610.024}{,}%
\TextAt{8}{326.78}{610.024}{then}%
\TextAt{13}{353.44}{610.024}{C}%
\TextAt{21}{361.815}{608.23}{K}%
\TextAt{9}{370.017}{610.024}{\Character{47}}%
\TextAt{13}{375.87}{610.024}{D}%
\TextAt{21}{385.59}{608.23}{K}%
\NativeMathAt{397.113}{610.024}{11.955}{\cong}%
\TextAt{13}{409.732}{610.024}{G}%
\TextAt{9}{418.966}{610.024}{\Character{40}}%
\TextAt{13}{423.518}{610.024}{K}%
\TextAt{21}{434.336}{614.362}{ab}%
\TextAt{16}{446.276}{610.024}{\Character{106}}%
\TextAt{13}{452.917}{610.024}{K}%
\TextAt{9}{463.735}{610.024}{\Character{41}}%
\TextAt{8}{468.287}{610.024}{.}%
\TextAt{12}{79.226}{581.234}{Exercise}%
\TextAt{12}{132.427}{581.234}{VI.5.2.}%
\TextAt{8}{176.183}{581.234}{For}%
\TextAt{8}{197.152}{581.234}{every}%
\TextAt{8}{228.044}{581.234}{place}%
\TextAt{14}{257.956}{581.234}{p}%
\TextAt{8}{267.836}{581.234}{of}%
\TextAt{13}{281.166}{581.234}{K}%
\TextAt{8}{295.881}{581.234}{one}%
\TextAt{8}{317.341}{581.234}{has}%
\TextAt{13}{338.215}{581.234}{K}%
\TextAt{21}{349.032}{585.573}{ab}%
\TextAt{27}{348.174}{578.279}{p}%
\TextAt{9}{360.972}{581.234}{\Character{61}}%
\TextAt{13}{373.397}{581.234}{K}%
\TextAt{21}{384.215}{585.573}{ab}%
\TextAt{13}{392.834}{581.234}{K}%
\TextAt{27}{402.793}{579.441}{p}%
\TextAt{8}{407.28}{581.234}{.}%
\TextAt{12}{97.159}{561.199}{Hint:}%
\TextAt{8}{132.698}{561.199}{Use}%
\TextAt{8}{155.185}{561.199}{(5.6)}%
\TextAt{8}{183.16}{561.199}{and}%
\TextAt{8}{205.923}{561.199}{(5.8).}%
\TextAt{12}{79.226}{532.41}{Exercise}%
\TextAt{12}{133.724}{532.41}{VI.5.3.}%
\TextAt{8}{178.608}{532.41}{Let}%
\TextAt{13}{200.705}{532.41}{p}%
\TextAt{8}{211.605}{532.41}{be}%
\TextAt{8}{228.657}{532.41}{a}%
\TextAt{8}{239.54}{532.41}{prime}%
\TextAt{8}{273.831}{532.41}{number,}%
\TextAt{8}{321.089}{532.41}{and}%
\TextAt{8}{344.98}{532.41}{let}%
\TextAt{13}{363.012}{532.41}{M}%
\TextAt{21}{374.337}{530.617}{p}%
\TextAt{16}{384.34}{532.41}{\Character{106}}%
\TextAt{13}{392.904}{532.41}{K}%
\TextAt{8}{408.747}{532.41}{be}%
\TextAt{8}{425.799}{532.41}{the}%
\TextAt{8}{447.083}{532.41}{maximal}%
\TextAt{8}{496.019}{532.41}{abelian}%
\TextAt{13}{79.226}{515.364}{p}%
\TextAt{8}{85.101}{515.364}{-extension}%
\TextAt{8}{140.821}{515.364}{unramified}%
\TextAt{8}{199.409}{515.364}{outside}%
\TextAt{8}{239.85}{515.364}{of}%
\TextAt{16}{253.236}{515.364}{\Character{102}}%
\TextAt{14}{259.213}{515.364}{p}%
\TextAt{16}{268.613}{515.364}{\Character{106}}%
\TextAt{13}{275.35}{515.364}{p}%
\TextAt{16}{281.225}{515.364}{\Character{103}}%
\TextAt{8}{287.202}{515.364}{.}%
\TextAt{8}{295.819}{515.364}{Further,}%
\TextAt{8}{341.569}{515.364}{let}%
\TextAt{13}{358.53}{515.364}{H}%
\TextAt{16}{372.569}{515.364}{\Character{106}}%
\TextAt{13}{379.306}{515.364}{K}%
\TextAt{8}{394.077}{515.364}{be}%
\TextAt{8}{410.057}{515.364}{the}%
\TextAt{8}{430.27}{515.364}{maximal}%
\TextAt{8}{478.134}{515.364}{unramified}%
\TextAt{8}{79.226}{498.318}{subextension}%
\TextAt{8}{149.76}{498.318}{of}%
\TextAt{13}{163.913}{498.318}{M}%
\TextAt{21}{175.238}{496.525}{p}%
\TextAt{16}{184.723}{498.318}{\Character{106}}%
\TextAt{13}{192.768}{498.318}{K}%
\TextAt{8}{208.306}{498.318}{in}%
\TextAt{8}{222.782}{498.318}{which}%
\TextAt{8}{257.092}{498.318}{the}%
\TextAt{8}{278.071}{498.318}{infinite}%
\TextAt{8}{318.562}{498.318}{places}%
\TextAt{8}{353.912}{498.318}{split}%
\TextAt{8}{380.809}{498.318}{completely.}%
\TextAt{8}{445.705}{498.318}{Then}%
\TextAt{8}{477.086}{498.318}{there}%
\TextAt{8}{507.821}{498.318}{is}%
\TextAt{8}{520.408}{498.318}{an}%
\TextAt{8}{79.226}{481.272}{exact}%
\TextAt{8}{110.118}{481.272}{sequence}%
\TextAt{9}{173.655}{457.872}{1}%
\NativeArrowAt{182.829}{457.872}{11.955}{19.26}{\longrightarrow}%
\TextAt{13}{205.411}{457.872}{G}%
\TextAt{9}{214.645}{457.872}{\Character{40}}%
\TextAt{13}{219.197}{457.872}{M}%
\TextAt{21}{230.522}{456.079}{p}%
\TextAt{16}{238.603}{457.872}{\Character{106}}%
\TextAt{13}{245.245}{457.872}{H}%
\TextAt{9}{255.868}{457.872}{\Character{41}}%
\NativeArrowAt{263.741}{457.872}{11.955}{19.26}{\longrightarrow}%
\TextAt{13}{286.323}{457.872}{G}%
\TextAt{9}{295.557}{457.872}{\Character{40}}%
\TextAt{13}{300.109}{457.872}{M}%
\TextAt{21}{311.434}{456.079}{p}%
\TextAt{16}{319.515}{457.872}{\Character{106}}%
\TextAt{13}{326.157}{457.872}{K}%
\TextAt{9}{336.975}{457.872}{\Character{41}}%
\NativeArrowAt{344.848}{457.872}{11.955}{19.26}{\longrightarrow}%
\TextAt{13}{367.43}{457.872}{C}%
\TextAt{13}{376.648}{457.872}{l}%
\TextAt{21}{380.151}{456.079}{K}%
\TextAt{9}{388.353}{457.872}{\Character{40}}%
\TextAt{13}{392.906}{457.872}{p}%
\TextAt{9}{398.781}{457.872}{\Character{41}}%
\NativeArrowAt{406.654}{457.872}{11.955}{19.26}{\longrightarrow}%
\TextAt{9}{429.236}{457.872}{1}%
\TextAt{13}{435.089}{457.872}{\Character{59}}%
\TextAt{8}{79.226}{434.472}{where}%
\TextAt{13}{111.884}{434.472}{C}%
\TextAt{13}{121.102}{434.472}{l}%
\TextAt{21}{124.604}{432.678}{K}%
\TextAt{9}{132.807}{434.472}{\Character{40}}%
\TextAt{13}{137.359}{434.472}{p}%
\TextAt{9}{143.235}{434.472}{\Character{41}}%
\TextAt{8}{150.533}{434.472}{is}%
\TextAt{8}{161.145}{434.472}{the}%
\TextAt{13}{180.15}{434.472}{p}%
\TextAt{8}{186.026}{434.472}{-Sylow}%
\TextAt{8}{222.593}{434.472}{subgroup}%
\TextAt{8}{272.239}{434.472}{of}%
\TextAt{8}{284.417}{434.472}{the}%
\TextAt{8}{303.422}{434.472}{ideal}%
\TextAt{8}{330.234}{434.472}{class}%
\TextAt{8}{356.52}{434.472}{group}%
\TextAt{13}{388.544}{434.472}{C}%
\TextAt{13}{397.762}{434.472}{l}%
\TextAt{21}{401.265}{432.678}{K}%
\TextAt{8}{409.467}{434.472}{,}%
\TextAt{8}{415.695}{434.472}{and}%
\TextAt{8}{437.307}{434.472}{there}%
\TextAt{8}{466.067}{434.472}{is}%
\TextAt{8}{476.68}{434.472}{a}%
\TextAt{8}{485.284}{434.472}{canonical}%
\TextAt{8}{79.226}{417.426}{isomorphism}%
\ErrataAt{170}{382.5}{315}{$\displaystyle G(M_p\mid H)\cong\frac{P\times S_p}{(P\times S_p)\cap\overline E^{\,\mathrm{sgn}}},\qquad P=\prod_{\mathfrak p\mid p}U_{\mathfrak p}^{(1)}.$}%
\ErrataAt{79.19}{352.5}{455.81}{Here $S_2=\{\pm1\}^r$ for the $r$ real places and $S_p=1$ for odd $p$.}%
\ErrataAt{79.19}{334.583}{455.81}{Here $\overline E^{\,\mathrm{sgn}}$ is the closure of the units in $\bigl(\prod_{\mathfrak p\mid p}U_{\mathfrak p}\bigr)\times S_p$, including the real signs.}%
\Vector{q 0 G 0 g 0.478 w 41.035 -429.599 m 50.402 -429.599 l S Q}%
\TextAt{12}{79.226}{304.251}{Exercise}%
\TextAt{12}{132.427}{304.251}{VI.5.4.}%
\TextAt{8}{176.183}{304.251}{The}%
\TextAt{8}{200.236}{304.251}{group}%
\Vector{q 0 G 0 g 0.478 w 161.412 -459.906 m 170.778 -459.906 l S Q}%
\TextAt{13}{233.412}{304.251}{E}%
\TextAt{9}{242.778}{304.251}{\Character{40}}%
\TextAt{13}{247.33}{304.251}{p}%
\TextAt{9}{253.206}{304.251}{\Character{41}}%
\TextAt{9}{261.061}{304.251}{\Character{58}\Character{61}}%
\Vector{q 0 G 0 g 0.478 w 204.756 -459.906 m 214.122 -459.906 l S Q}%
\TextAt{13}{276.756}{304.251}{E}%
\TextAt{16}{288.779}{304.251}{\Character{92}}%
\TextAt{20}{299.406}{313.217}{\Character{81}}%
\TextAt{27}{310.697}{300.764}{p}%
\TextAt{3}{314.685}{300.764}{\Character{106}}%
\TextAt{21}{317.038}{300.764}{p}%
\TextAt{13}{323.791}{304.251}{U}%
\TextAt{17}{332.99}{310.361}{\Character{40}1\Character{41}}%
\TextAt{27}{331.741}{302.28}{p}%
\TextAt{8}{348.206}{304.251}{is}%
\TextAt{8}{359.97}{304.251}{a}%
\TextAt{18}{369.726}{304.251}{Z}%
\TextAt{21}{377.696}{302.457}{p}%
\TextAt{8}{382.457}{304.251}{-module}%
\TextAt{8}{427.647}{304.251}{of}%
\TextAt{8}{440.977}{304.251}{rank}%
\TextAt{13}{168.613}{280.85}{r}%
\TextAt{21}{173.888}{279.057}{p}%
\TextAt{9}{178.649}{280.85}{\Character{40}}%
\TextAt{13}{183.201}{280.85}{E}%
\TextAt{9}{192.568}{280.85}{\Character{41}}%
\TextAt{9}{200.423}{280.85}{\Character{58}\Character{61}}%
\TextAt{8}{216.118}{280.85}{rank}%
\TextAt{25}{239.215}{279.057}{Z}%
\TextAt{22}{244.529}{278.061}{p}%
\TextAt{9}{249.365}{280.85}{\Character{40}}%
\Vector{q 0 G 0 g 0.478 w 181.918 -483.307 m 191.284 -483.307 l S Q}%
\TextAt{13}{253.918}{280.85}{E}%
\TextAt{9}{263.284}{280.85}{\Character{40}}%
\TextAt{13}{267.836}{280.85}{p}%
\TextAt{9}{273.711}{280.85}{\Character{41}\Character{41}}%
\TextAt{9}{286.108}{280.85}{\Character{61}}%
\TextAt{9}{298.53}{280.85}{\Character{91}}%
\TextAt{13}{301.814}{280.85}{K}%
\TextAt{9}{315.952}{280.85}{\Character{58}}%
\TextAt{18}{322.525}{280.85}{Q}%
\TextAt{9}{331.823}{280.85}{\Character{93}}%
\TextAt{16}{337.732}{280.85}{\Character{0}}%
\TextAt{8}{349.687}{280.85}{rank}%
\TextAt{25}{372.784}{279.057}{Z}%
\TextAt{22}{378.098}{278.061}{p}%
\TextAt{13}{384.927}{280.85}{G}%
\TextAt{9}{394.161}{280.85}{\Character{40}}%
\TextAt{13}{398.713}{280.85}{M}%
\TextAt{21}{410.038}{279.057}{p}%
\TextAt{16}{418.119}{280.85}{\Character{106}}%
\TextAt{13}{424.761}{280.85}{K}%
\TextAt{9}{435.579}{280.85}{\Character{41}}%
\TextAt{13}{440.131}{280.85}{\Character{58}}%
\TextAt{13}{79.226}{257.45}{r}%
\TextAt{21}{84.501}{255.657}{p}%
\TextAt{9}{89.262}{257.45}{\Character{40}}%
\TextAt{13}{93.815}{257.45}{E}%
\TextAt{9}{103.181}{257.45}{\Character{41}}%
\TextAt{8}{111.631}{257.45}{is}%
\TextAt{8}{123.394}{257.45}{called}%
\TextAt{8}{156.558}{257.45}{the}%
\TextAt{13}{176.715}{257.45}{p}%
\TextAt{12}{182.59}{257.45}{-adic}%
\TextAt{12}{215.287}{257.45}{unit}%
\TextAt{12}{243.693}{257.45}{rank}%
\TextAt{8}{270.305}{257.45}{.}%
\TextAt{12}{97.159}{237.415}{Problem:}%
\TextAt{8}{156.082}{237.415}{For}%
\TextAt{8}{177.051}{237.415}{the}%
\TextAt{13}{197.208}{237.415}{p}%
\TextAt{8}{203.083}{237.415}{-adic}%
\TextAt{8}{231.692}{237.415}{unit}%
\TextAt{8}{256.403}{237.415}{rank,}%
\TextAt{8}{286.65}{237.415}{one}%
\TextAt{8}{308.109}{237.415}{has}%
\TextAt{8}{328.983}{237.415}{the}%
\TextAt{8}{349.139}{237.415}{famous}%
\TextAt{8}{389.201}{237.415}{Leopoldt}%
\TextAt{8}{438.468}{237.415}{conjecture:}%
\TextAt{13}{259.247}{214.015}{r}%
\TextAt{21}{264.522}{212.222}{p}%
\TextAt{9}{269.283}{214.015}{\Character{40}}%
\TextAt{13}{273.835}{214.015}{E}%
\TextAt{9}{283.201}{214.015}{\Character{41}}%
\TextAt{9}{291.056}{214.015}{\Character{61}}%
\TextAt{13}{303.5}{214.015}{r}%
\TextAt{9}{311.757}{214.015}{\Character{43}}%
\TextAt{13}{323.518}{214.015}{s}%
\TextAt{16}{331.689}{214.015}{\Character{0}}%
\TextAt{9}{343.644}{214.015}{1}%
\TextAt{13}{349.497}{214.015}{\Character{59}}%
\TextAt{8}{79.226}{190.615}{where}%
\TextAt{13}{115.062}{190.615}{r}%
\TextAt{8}{120.662}{190.615}{,}%
\TextAt{8}{130.345}{190.615}{resp.}%
\TextAt{13}{165.746}{190.615}{s}%
\TextAt{8}{171.26}{190.615}{,}%
\TextAt{8}{180.943}{190.615}{is}%
\TextAt{8}{194.733}{190.615}{the}%
\TextAt{8}{216.916}{190.615}{number}%
\TextAt{8}{261.539}{190.615}{of}%
\TextAt{8}{276.896}{190.615}{all}%
\TextAt{8}{295.181}{190.615}{real,}%
\TextAt{8}{323.729}{190.615}{resp.}%
\TextAt{8}{359.13}{190.615}{complex,}%
\TextAt{8}{410.763}{190.615}{places;}%
\TextAt{8}{451.582}{190.615}{in}%
\TextAt{8}{467.261}{190.615}{other}%
\TextAt{8}{499.857}{190.615}{words,}%
\TextAt{8}{79.226}{173.569}{rank}%
\TextAt{25}{102.324}{171.776}{Z}%
\TextAt{22}{107.637}{170.779}{p}%
\TextAt{13}{114.466}{173.569}{G}%
\TextAt{9}{123.7}{173.569}{\Character{40}}%
\TextAt{13}{128.252}{173.569}{M}%
\TextAt{21}{139.577}{171.776}{p}%
\TextAt{16}{150.074}{173.569}{\Character{106}}%
\TextAt{13}{159.131}{173.569}{K}%
\TextAt{9}{169.948}{173.569}{\Character{41}}%
\TextAt{9}{180.218}{173.569}{\Character{61}}%
\TextAt{13}{195.077}{173.569}{s}%
\TextAt{9}{204.214}{173.569}{\Character{43}}%
\TextAt{9}{216.934}{173.569}{1}%
\TextAt{8}{222.794}{173.569}{.}%
\TextAt{8}{235.494}{173.569}{The}%
\TextAt{8}{260.965}{173.569}{Leopoldt}%
\TextAt{8}{311.65}{173.569}{conjecture}%
\TextAt{8}{369.316}{173.569}{was}%
\TextAt{8}{393.233}{173.569}{proved}%
\TextAt{8}{432.704}{173.569}{for}%
\TextAt{8}{452.006}{173.569}{abelian}%
\TextAt{8}{494.071}{173.569}{number}%
\TextAt{8}{79.226}{156.523}{fields}%
\TextAt{13}{109.535}{156.523}{K}%
\TextAt{16}{124.248}{156.523}{\Character{106}}%
\TextAt{18}{131.465}{156.523}{Q}%
\TextAt{8}{144.998}{156.523}{by}%
\TextAt{8}{161.594}{156.523}{the}%
\TextAt{8}{182.088}{156.523}{American}%
\TextAt{8}{235.422}{156.523}{mathematician}%
\TextAt{26}{315.764}{156.523}{Armand}%
\TextAt{26}{368.65}{156.523}{Brumer}%
\TextAt{8}{419.163}{156.523}{[22].}%
\TextAt{8}{446.843}{156.523}{The}%
\TextAt{8}{471.234}{156.523}{general}%
\TextAt{8}{511.896}{156.523}{case}%
\TextAt{8}{79.226}{139.477}{is}%
\TextAt{8}{90.99}{139.477}{still}%
\TextAt{8}{113.812}{139.477}{open}%
\TextAt{8}{142.098}{139.477}{to}%
\TextAt{8}{156.409}{139.477}{date.}%
\TextAt{8}{297.211}{43.033}{359}%
\Destination{78.23}{809.815}{nav.page.359}%
\Destination{79.226}{641.055}{nav.ex.6.5.1}%
\Destination{79.226}{595.255}{nav.ex.6.5.2}%
\Destination{79.226}{546.455}{nav.ex.6.5.3}%
\Destination{79.226}{318.255}{nav.ex.6.5.4}%
\Destination{79.36}{776.155}{nav.print.394}%
\Hotlink{418.663}{154.355}{19.22}{11.047}{nav.cite.N022}%
\Hotlink{257.922}{758.955}{25.078}{11.047}{nav.number.6.5.6}%
\Hotlink{78.726}{673.655}{25.078}{11.047}{nav.number.6.5.7}%
\Hotlink{154.685}{558.955}{25.078}{11.047}{nav.number.6.5.6}%
\Hotlink{205.423}{558.955}{25.078}{11.047}{nav.number.6.5.8}%
\Destination{79.226}{399.5}{errata.137}%
}

% Page 372 (body page 360)
\Page{%
\TextAt{2}{79.226}{796.943}{VI.}%
\TextAt{2}{98.012}{796.943}{Global}%
\TextAt{2}{133.248}{796.943}{Class}%
\TextAt{2}{161.841}{796.943}{Field}%
\TextAt{2}{189.572}{796.943}{Theory}%
\Vector{q 0 G 0 g 0.299 w 7.226 18.729 m 460.769 18.729 l S Q}%
\TextAt{15}{79.226}{760.804}{§}%
\TextAt{15}{93.785}{760.804}{6}%
\TextAt{15}{120.693}{760.804}{Global}%
\TextAt{15}{183.013}{760.804}{Class}%
\TextAt{15}{233.557}{760.804}{Fields}%
\TextAt{8}{79.226}{728.918}{As}%
\TextAt{8}{96.523}{728.918}{in}%
\TextAt{8}{110.185}{728.918}{local}%
\TextAt{8}{137.835}{728.918}{class}%
\TextAt{8}{165.281}{728.918}{field}%
\TextAt{8}{190.648}{728.918}{theory,}%
\TextAt{8}{229.681}{728.918}{the}%
\TextAt{8}{249.847}{728.918}{reciprocity}%
\TextAt{8}{308.066}{728.918}{law}%
\TextAt{8}{329.212}{728.918}{provides}%
\TextAt{8}{375.464}{728.918}{also}%
\TextAt{8}{398.953}{728.918}{in}%
\TextAt{8}{412.616}{728.918}{global}%
\TextAt{8}{447.104}{728.918}{class}%
\TextAt{8}{474.55}{728.918}{field}%
\TextAt{8}{499.917}{728.918}{theory}%
\TextAt{8}{79.226}{711.872}{a}%
\TextAt{8}{88.427}{711.872}{complete}%
\TextAt{8}{137.295}{711.872}{classification}%
\TextAt{8}{205.16}{711.872}{of}%
\TextAt{8}{217.935}{711.872}{all}%
\TextAt{8}{233.64}{711.872}{abelian}%
\TextAt{8}{273.733}{711.872}{extensions}%
\TextAt{8}{329.559}{711.872}{of}%
\TextAt{8}{342.334}{711.872}{a}%
\TextAt{8}{351.535}{711.872}{finite}%
\TextAt{8}{380.892}{711.872}{algebraic}%
\TextAt{8}{429.772}{711.872}{number}%
\TextAt{8}{471.814}{711.872}{field}%
\TextAt{13}{496.616}{711.872}{K}%
\TextAt{8}{507.434}{711.872}{.}%
\TextAt{8}{515.697}{711.872}{For}%
\TextAt{8}{79.226}{694.826}{this}%
\TextAt{8}{101.585}{694.826}{it}%
\TextAt{8}{112.826}{694.826}{is}%
\TextAt{8}{124.126}{694.826}{necessary}%
\TextAt{8}{175.488}{694.826}{to}%
\TextAt{8}{189.335}{694.826}{view}%
\TextAt{8}{215.854}{694.826}{the}%
\TextAt{8}{235.547}{694.826}{idèle}%
\TextAt{8}{262.389}{694.826}{class}%
\TextAt{8}{289.362}{694.826}{group}%
\TextAt{13}{322.074}{694.826}{C}%
\TextAt{21}{330.449}{693.033}{K}%
\TextAt{8}{342.086}{694.826}{as}%
\TextAt{8}{355.993}{694.826}{a}%
\TextAt{8}{365.284}{694.826}{topological}%
\TextAt{8}{424.346}{694.826}{group,}%
\TextAt{8}{460.402}{694.826}{equipped}%
\TextAt{8}{510.007}{694.826}{with}%
\TextAt{8}{79.226}{677.78}{its}%
\TextAt{8}{95.747}{677.78}{natural}%
\TextAt{8}{136.931}{677.78}{topology}%
\TextAt{8}{185.277}{677.78}{which}%
\TextAt{8}{218.965}{677.78}{the}%
\TextAt{8}{239.324}{677.78}{valuations}%
\TextAt{8}{295.213}{677.78}{of}%
\TextAt{8}{308.745}{677.78}{the}%
\TextAt{8}{329.103}{677.78}{various}%
\TextAt{8}{369.379}{677.78}{completions}%
\TextAt{13}{434.031}{677.78}{K}%
\TextAt{27}{443.99}{675.987}{p}%
\TextAt{8}{452.577}{677.78}{impress}%
\TextAt{8}{495.172}{677.78}{upon}%
\TextAt{8}{524.963}{677.78}{it}%
\TextAt{8}{79.226}{660.734}{(see}%
\TextAt{8}{102.694}{660.734}{§1).}%
\TextAt{12}{79.226}{634.88}{Theorem}%
\TextAt{12}{136.24}{634.88}{(6.1).}%
\TextAt{8}{171.556}{634.88}{The}%
\TextAt{8}{195.61}{634.88}{map}%
\TextAt{13}{252.998}{617.834}{L}%
\NativeArrowAt{264.283}{617.834}{11.955}{19.26}{\longmapsto}%
\TextAt{16}{286.854}{617.834}{\Character{78}}%
\TextAt{21}{296.674}{616.041}{L}%
\TextAt{9}{306.252}{617.834}{\Character{61}}%
\TextAt{13}{318.678}{617.834}{N}%
\TextAt{21}{328.052}{615.974}{L}%
\TextAt{3}{333.811}{615.974}{\Character{106}}%
\TextAt{21}{336.163}{615.974}{K}%
\TextAt{13}{344.366}{617.834}{C}%
\TextAt{21}{352.741}{616.041}{L}%
\TextAt{8}{79.226}{594.549}{is}%
\TextAt{8}{91.441}{594.549}{a}%
\TextAt{9}{101.648}{594.549}{1}%
\TextAt{8}{107.501}{594.549}{-}%
\TextAt{9}{111.398}{594.549}{1}%
\TextAt{8}{117.251}{594.549}{-correspondence}%
\TextAt{8}{203.277}{594.549}{between}%
\TextAt{8}{248.919}{594.549}{the}%
\TextAt{8}{269.527}{594.549}{finite}%
\TextAt{8}{299.89}{594.549}{abelian}%
\TextAt{8}{340.988}{594.549}{extensions}%
\TextAt{13}{397.82}{594.549}{L}%
\TextAt{16}{409.874}{594.549}{\Character{106}}%
\TextAt{13}{417.285}{594.549}{K}%
\TextAt{8}{432.451}{594.549}{and}%
\TextAt{8}{455.665}{594.549}{the}%
\TextAt{8}{476.272}{594.549}{closed}%
\TextAt{8}{511.25}{594.549}{sub-}%
\TextAt{8}{79.226}{577.503}{groups}%
\TextAt{8}{117.016}{577.503}{of}%
\TextAt{8}{130.346}{577.503}{finite}%
\TextAt{8}{160.258}{577.503}{index}%
\TextAt{8}{191.796}{577.503}{in}%
\TextAt{13}{205.449}{577.503}{C}%
\TextAt{21}{213.824}{575.71}{K}%
\TextAt{8}{222.027}{577.503}{.}%
\TextAt{8}{230.475}{577.503}{Moreover}%
\TextAt{8}{281.858}{577.503}{one}%
\TextAt{8}{303.318}{577.503}{has:}%
\TextAt{13}{116.311}{549.055}{L}%
\TextAt{17}{124.275}{547.261}{1}%
\NativeMathAt{132.328}{549.055}{11.955}{\subseteq}%
\TextAt{13}{144.948}{549.055}{L}%
\TextAt{17}{152.912}{547.261}{2}%
\NativeArrowAt{160.965}{549.055}{11.955}{21.926}{\Longleftrightarrow}%
\TextAt{16}{186.202}{549.055}{\Character{78}}%
\TextAt{21}{196.013}{547.261}{L}%
\TextAt{24}{201.772}{546.154}{1}%
\NativeMathAt{209.742}{549.055}{11.955}{\supseteq}%
\TextAt{16}{222.343}{549.055}{\Character{78}}%
\TextAt{21}{232.17}{547.261}{L}%
\TextAt{24}{237.93}{546.154}{2}%
\TextAt{13}{242.579}{549.055}{\Character{59}}%
\TextAt{16}{271.733}{549.055}{\Character{78}}%
\TextAt{21}{281.542}{547.261}{L}%
\TextAt{24}{287.301}{546.154}{1}%
\TextAt{21}{291.452}{547.261}{L}%
\TextAt{24}{297.211}{546.154}{2}%
\TextAt{9}{305.181}{549.055}{\Character{61}}%
\TextAt{16}{317.607}{549.055}{\Character{78}}%
\TextAt{21}{327.416}{547.261}{L}%
\TextAt{24}{333.175}{546.154}{1}%
\TextAt{16}{340.481}{549.055}{\Character{92}}%
\TextAt{16}{351.097}{549.055}{\Character{78}}%
\TextAt{21}{360.917}{547.261}{L}%
\TextAt{24}{366.676}{546.154}{2}%
\TextAt{13}{371.325}{549.055}{\Character{59}}%
\TextAt{16}{400.48}{549.055}{\Character{78}}%
\TextAt{21}{410.289}{547.261}{L}%
\TextAt{24}{416.048}{546.154}{1}%
\TextAt{3}{420.199}{547.261}{\Character{92}}%
\TextAt{21}{425.844}{547.261}{L}%
\TextAt{24}{431.603}{546.154}{2}%
\TextAt{9}{439.573}{549.055}{\Character{61}}%
\TextAt{16}{451.999}{549.055}{\Character{78}}%
\TextAt{21}{461.808}{547.261}{L}%
\TextAt{24}{467.567}{546.154}{1}%
\TextAt{16}{472.216}{549.055}{\Character{78}}%
\TextAt{21}{482.025}{547.261}{L}%
\TextAt{24}{487.784}{546.154}{2}%
\TextAt{13}{492.434}{549.055}{\Character{58}}%
\TextAt{8}{79.226}{520.606}{The}%
\TextAt{8}{103.712}{520.606}{field}%
\TextAt{13}{129.501}{520.606}{L}%
\TextAt{16}{141.524}{520.606}{\Character{106}}%
\TextAt{13}{148.902}{520.606}{K}%
\TextAt{8}{164.049}{520.606}{corresponding}%
\TextAt{8}{239.668}{520.606}{to}%
\TextAt{8}{254.41}{520.606}{the}%
\TextAt{8}{274.999}{520.606}{subgroup}%
\TextAt{16}{326.229}{520.606}{\Character{78}}%
\TextAt{8}{342.13}{520.606}{of}%
\TextAt{13}{355.892}{520.606}{C}%
\TextAt{21}{364.267}{518.813}{K}%
\TextAt{8}{376.799}{520.606}{is}%
\TextAt{8}{388.995}{520.606}{called}%
\TextAt{8}{422.591}{520.606}{the}%
\TextAt{12}{443.18}{520.606}{class}%
\TextAt{12}{475.036}{520.606}{field}%
\TextAt{8}{504.185}{520.606}{of}%
\TextAt{16}{517.947}{520.606}{\Character{78}}%
\TextAt{8}{529.518}{520.606}{.}%
\TextAt{8}{79.226}{503.56}{It}%
\TextAt{8}{91.899}{503.56}{satisfies}%
\TextAt{13}{255.86}{486.514}{G}%
\TextAt{9}{265.093}{486.514}{\Character{40}}%
\TextAt{13}{269.646}{486.514}{L}%
\TextAt{16}{280.931}{486.514}{\Character{106}}%
\TextAt{13}{287.573}{486.514}{K}%
\TextAt{9}{298.39}{486.514}{\Character{41}}%
\NativeMathAt{306.264}{486.514}{11.955}{\cong}%
\TextAt{13}{318.883}{486.514}{C}%
\TextAt{21}{327.258}{484.721}{K}%
\TextAt{9}{335.461}{486.514}{\Character{47}}%
\TextAt{16}{341.314}{486.514}{\Character{78}}%
\TextAt{13}{352.884}{486.514}{\Character{58}}%
\TextAt{19}{79.226}{445.612}{Proof.}%
\TextAt{8}{116}{445.612}{By}%
\TextAt{8}{135.642}{445.612}{chap.}%
\TextAt{8}{167.814}{445.612}{IV,}%
\TextAt{8}{189.237}{445.612}{(6.7),}%
\TextAt{8}{222.063}{445.612}{all}%
\TextAt{8}{239.601}{445.612}{we}%
\TextAt{8}{258.107}{445.612}{have}%
\TextAt{8}{286.381}{445.612}{to}%
\TextAt{8}{301.97}{445.612}{show}%
\TextAt{8}{332.253}{445.612}{is}%
\TextAt{8}{345.296}{445.612}{that}%
\TextAt{8}{371.943}{445.612}{the}%
\TextAt{8}{393.379}{445.612}{subgroups}%
\TextAt{16}{450.07}{445.612}{\Character{78}}%
\TextAt{8}{466.817}{445.612}{of}%
\TextAt{13}{481.426}{445.612}{C}%
\TextAt{21}{489.801}{443.819}{K}%
\TextAt{8}{503.18}{445.612}{which}%
\TextAt{8}{79.226}{428.566}{are}%
\TextAt{8}{99.193}{428.566}{open}%
\TextAt{8}{127.935}{428.566}{in}%
\TextAt{8}{142.043}{428.566}{the}%
\TextAt{8}{162.656}{428.566}{norm}%
\TextAt{8}{193.681}{428.566}{topology}%
\TextAt{8}{242.28}{428.566}{are}%
\TextAt{8}{262.247}{428.566}{precisely}%
\TextAt{8}{310.559}{428.566}{the}%
\TextAt{8}{331.172}{428.566}{closed}%
\TextAt{8}{366.154}{428.566}{subgroups}%
\TextAt{8}{422.022}{428.566}{of}%
\TextAt{8}{435.808}{428.566}{finite}%
\TextAt{8}{466.176}{428.566}{index}%
\TextAt{8}{498.17}{428.566}{for}%
\TextAt{8}{516.51}{428.566}{the}%
\TextAt{8}{79.226}{411.52}{natural}%
\TextAt{8}{120.208}{411.52}{topology.}%
\TextAt{8}{97.159}{391.486}{If}%
\TextAt{8}{107.873}{391.486}{the}%
\TextAt{8}{127.052}{391.486}{subgroup}%
\TextAt{16}{176.871}{391.486}{\Character{78}}%
\TextAt{8}{191.361}{391.486}{is}%
\TextAt{8}{202.147}{391.486}{open}%
\TextAt{8}{229.455}{391.486}{in}%
\TextAt{8}{242.13}{391.486}{the}%
\TextAt{8}{261.309}{391.486}{norm}%
\TextAt{8}{290.9}{391.486}{topology,}%
\TextAt{8}{340.533}{391.486}{then}%
\TextAt{8}{366.215}{391.486}{it}%
\TextAt{8}{376.941}{391.486}{contains}%
\TextAt{8}{421.883}{391.486}{a}%
\TextAt{8}{430.66}{391.486}{norm}%
\TextAt{8}{460.252}{391.486}{group}%
\TextAt{13}{492.449}{391.486}{N}%
\TextAt{21}{501.823}{389.626}{L}%
\TextAt{3}{507.582}{389.626}{\Character{106}}%
\TextAt{21}{509.935}{389.626}{K}%
\TextAt{13}{518.137}{391.486}{C}%
\TextAt{21}{526.512}{389.692}{L}%
\TextAt{8}{79.226}{374.44}{and}%
\TextAt{8}{102.466}{374.44}{is}%
\TextAt{8}{114.707}{374.44}{therefore}%
\TextAt{8}{164.283}{374.44}{of}%
\TextAt{8}{178.09}{374.44}{finite}%
\TextAt{8}{208.479}{374.44}{index,}%
\TextAt{8}{243.865}{374.44}{because}%
\TextAt{8}{287.644}{374.44}{from}%
\TextAt{8}{315.761}{374.44}{(5.5),}%
\TextAt{9}{347.584}{374.44}{\Character{40}}%
\TextAt{13}{352.136}{374.44}{C}%
\TextAt{21}{360.511}{372.646}{K}%
\TextAt{9}{372.848}{374.44}{\Character{58}}%
\TextAt{13}{380.233}{374.44}{N}%
\TextAt{21}{389.607}{372.58}{L}%
\TextAt{3}{395.366}{372.58}{\Character{106}}%
\TextAt{21}{397.718}{372.58}{K}%
\TextAt{13}{405.921}{374.44}{C}%
\TextAt{21}{414.296}{372.646}{L}%
\TextAt{9}{420.553}{374.44}{\Character{41}}%
\TextAt{9}{429.221}{374.44}{\Character{61}}%
\TextAt{8}{442.477}{374.44}{\char35\relax }%
\TextAt{13}{452.233}{374.44}{G}%
\TextAt{9}{461.466}{374.44}{\Character{40}}%
\TextAt{13}{466.019}{374.44}{L}%
\TextAt{16}{478.117}{374.44}{\Character{106}}%
\TextAt{13}{485.571}{374.44}{K}%
\TextAt{9}{496.389}{374.44}{\Character{41}}%
\TextAt{21}{500.941}{378.778}{ab}%
\TextAt{8}{509.56}{374.44}{.}%
\TextAt{8}{519.439}{374.44}{To}%
\TextAt{8}{79.226}{357.394}{show}%
\TextAt{8}{107.702}{357.394}{that}%
\TextAt{16}{132.544}{357.394}{\Character{78}}%
\TextAt{8}{147.486}{357.394}{is}%
\TextAt{8}{158.723}{357.394}{closed}%
\TextAt{8}{192.722}{357.394}{it}%
\TextAt{8}{203.9}{357.394}{is}%
\TextAt{8}{215.137}{357.394}{enough}%
\TextAt{8}{254.934}{357.394}{to}%
\TextAt{8}{268.718}{357.394}{show}%
\TextAt{8}{297.194}{357.394}{that}%
\TextAt{13}{322.036}{357.394}{N}%
\TextAt{21}{331.41}{355.534}{L}%
\TextAt{3}{337.169}{355.534}{\Character{106}}%
\TextAt{21}{339.521}{355.534}{K}%
\TextAt{13}{347.724}{357.394}{C}%
\TextAt{21}{356.099}{355.6}{L}%
\TextAt{8}{365.727}{357.394}{is.}%
\TextAt{8}{381.866}{357.394}{For}%
\TextAt{8}{402.308}{357.394}{this,}%
\TextAt{8}{427.961}{357.394}{we}%
\TextAt{8}{444.662}{357.394}{choose}%
\TextAt{8}{481.267}{357.394}{an}%
\TextAt{8}{497}{357.394}{infinite}%
\TextAt{8}{79.226}{340.348}{place}%
\TextAt{14}{110.015}{340.348}{p}%
\TextAt{8}{120.773}{340.348}{of}%
\TextAt{13}{134.981}{340.348}{K}%
\TextAt{8}{150.573}{340.348}{and}%
\TextAt{8}{174.214}{340.348}{denote}%
\TextAt{8}{212.81}{340.348}{by}%
\TextAt{9}{229.946}{340.348}{\Character{0}}%
\TextAt{21}{237.26}{338.554}{K}%
\TextAt{8}{250.238}{340.348}{the}%
\TextAt{8}{271.272}{340.348}{image}%
\TextAt{8}{305.97}{340.348}{of}%
\TextAt{8}{320.178}{340.348}{the}%
\TextAt{8}{341.212}{340.348}{subgroup}%
\TextAt{8}{392.887}{340.348}{of}%
\TextAt{8}{407.094}{340.348}{positive}%
\TextAt{8}{451.285}{340.348}{real}%
\TextAt{8}{474.926}{340.348}{numbers}%
\TextAt{8}{523.014}{340.348}{in}%
\TextAt{13}{79.226}{323.302}{K}%
\TextAt{3}{90.044}{327.64}{\Character{3}}%
\TextAt{27}{89.185}{320.346}{p}%
\TextAt{8}{99.162}{323.302}{under}%
\TextAt{8}{132.813}{323.302}{the}%
\TextAt{8}{153.458}{323.302}{mapping}%
\TextAt{16}{202.077}{323.302}{\Character{104}}%
\TextAt{16}{211.091}{323.302}{\Character{105}}%
\TextAt{9}{219.914}{323.302}{\Character{58}}%
\TextAt{13}{227.318}{323.302}{K}%
\TextAt{3}{238.135}{327.64}{\Character{3}}%
\TextAt{27}{237.277}{320.346}{p}%
\NativeArrowAt{247.02}{323.302}{11.955}{19.26}{\longrightarrow}%
\TextAt{13}{270.434}{323.302}{C}%
\TextAt{21}{278.809}{321.508}{K}%
\TextAt{8}{287.012}{323.302}{.}%
\TextAt{8}{296.924}{323.302}{Then}%
\TextAt{9}{327.969}{323.302}{\Character{0}}%
\TextAt{21}{335.284}{321.508}{K}%
\TextAt{8}{347.872}{323.302}{is}%
\TextAt{8}{360.124}{323.302}{a}%
\TextAt{8}{370.367}{323.302}{group}%
\TextAt{8}{404.031}{323.302}{of}%
\TextAt{8}{417.849}{323.302}{representatives}%
\TextAt{8}{498.138}{323.302}{for}%
\TextAt{8}{516.51}{323.302}{the}%
\TextAt{8}{79.226}{306.256}{homomorphism}%
\TextAt{14}{162.578}{306.256}{N}%
\TextAt{9}{177.008}{306.256}{\Character{58}}%
\TextAt{13}{184.742}{306.256}{C}%
\TextAt{21}{193.117}{304.462}{K}%
\NativeArrowAt{205.802}{306.256}{11.955}{19.26}{\longrightarrow}%
\TextAt{18}{229.546}{306.256}{R}%
\TextAt{3}{238.18}{310.594}{\Character{3}}%
\TextAt{17}{238.18}{303.3}{\Character{43}}%
\TextAt{8}{249.844}{306.256}{with}%
\TextAt{8}{277.185}{306.256}{kernel}%
\TextAt{13}{312.334}{306.256}{C}%
\TextAt{17}{321.567}{310.594}{0}%
\TextAt{21}{320.709}{303.235}{K}%
\TextAt{8}{333.49}{306.256}{(see}%
\TextAt{8}{357.64}{306.256}{§1),}%
\TextAt{8}{381.733}{306.256}{i.e.,}%
\TextAt{13}{404.69}{306.256}{C}%
\TextAt{21}{413.065}{304.462}{K}%
\TextAt{9}{425.75}{306.256}{\Character{61}}%
\TextAt{13}{439.337}{306.256}{C}%
\TextAt{17}{448.571}{310.594}{0}%
\TextAt{21}{447.712}{303.235}{K}%
\TextAt{16}{459.036}{306.256}{\Character{2}}%
\TextAt{9}{471.456}{306.256}{\Character{0}}%
\TextAt{21}{478.77}{304.462}{K}%
\TextAt{8}{486.973}{306.256}{.}%
\TextAt{8}{497.466}{306.256}{By}%
\TextAt{8}{516.51}{306.256}{the}%
\TextAt{8}{79.226}{289.21}{same}%
\TextAt{8}{108.569}{289.21}{token,}%
\TextAt{9}{143.714}{289.21}{\Character{0}}%
\TextAt{21}{151.028}{287.417}{K}%
\TextAt{8}{163.145}{289.21}{is}%
\TextAt{8}{174.925}{289.21}{a}%
\TextAt{8}{184.697}{289.21}{group}%
\TextAt{8}{217.89}{289.21}{of}%
\TextAt{8}{231.236}{289.21}{representatives}%
\TextAt{8}{311.054}{289.21}{for}%
\TextAt{8}{328.955}{289.21}{the}%
\TextAt{8}{349.128}{289.21}{homomorphism}%
\TextAt{14}{431.815}{289.21}{N}%
\TextAt{9}{445.112}{289.21}{\Character{58}}%
\TextAt{13}{451.713}{289.21}{C}%
\TextAt{21}{460.088}{287.417}{L}%
\NativeArrowAt{469.695}{289.21}{11.955}{19.26}{\longrightarrow}%
\TextAt{18}{492.305}{289.21}{R}%
\TextAt{3}{500.939}{293.548}{\Character{3}}%
\TextAt{17}{500.939}{286.254}{\Character{43}}%
\TextAt{8}{508.024}{289.21}{.}%
\TextAt{8}{516.522}{289.21}{We}%
\TextAt{8}{79.226}{272.164}{therefore}%
\TextAt{8}{128.326}{272.164}{get}%
\TextAt{13}{142.075}{243.715}{N}%
\TextAt{21}{151.449}{241.856}{L}%
\TextAt{3}{157.208}{241.856}{\Character{106}}%
\TextAt{21}{159.561}{241.856}{K}%
\TextAt{13}{167.763}{243.715}{C}%
\TextAt{21}{176.138}{241.922}{L}%
\TextAt{9}{185.716}{243.715}{\Character{61}}%
\TextAt{13}{198.142}{243.715}{N}%
\TextAt{21}{207.516}{241.856}{L}%
\TextAt{3}{213.275}{241.856}{\Character{106}}%
\TextAt{21}{215.627}{241.856}{K}%
\TextAt{13}{223.83}{243.715}{C}%
\TextAt{17}{233.063}{248.651}{0}%
\TextAt{21}{232.205}{240.76}{L}%
\TextAt{16}{241.119}{243.715}{\Character{2}}%
\TextAt{13}{253.074}{243.715}{N}%
\TextAt{21}{262.448}{241.856}{L}%
\TextAt{3}{268.207}{241.856}{\Character{106}}%
\TextAt{21}{270.559}{241.856}{K}%
\TextAt{9}{278.762}{243.715}{\Character{0}}%
\TextAt{21}{286.076}{241.922}{K}%
\TextAt{9}{297.599}{243.715}{\Character{61}}%
\TextAt{13}{310.025}{243.715}{N}%
\TextAt{21}{319.399}{241.856}{L}%
\TextAt{3}{325.158}{241.856}{\Character{106}}%
\TextAt{21}{327.51}{241.856}{K}%
\TextAt{13}{335.713}{243.715}{C}%
\TextAt{17}{344.946}{248.651}{0}%
\TextAt{21}{344.088}{240.76}{L}%
\TextAt{16}{353.001}{243.715}{\Character{2}}%
\TextAt{9}{364.957}{243.715}{\Character{0}}%
\TextAt{21}{372.271}{248.651}{n}%
\TextAt{21}{372.271}{240.76}{K}%
\TextAt{9}{383.794}{243.715}{\Character{61}}%
\TextAt{13}{396.22}{243.715}{N}%
\TextAt{21}{405.594}{241.856}{L}%
\TextAt{3}{411.353}{241.856}{\Character{106}}%
\TextAt{21}{413.705}{241.856}{K}%
\TextAt{13}{421.908}{243.715}{C}%
\TextAt{17}{431.141}{248.651}{0}%
\TextAt{21}{430.283}{240.76}{L}%
\TextAt{16}{439.197}{243.715}{\Character{2}}%
\TextAt{9}{451.152}{243.715}{\Character{0}}%
\TextAt{21}{458.466}{241.922}{K}%
\TextAt{13}{466.669}{243.715}{\Character{58}}%
\TextAt{8}{79.226}{215.267}{The}%
\TextAt{8}{102.843}{215.267}{norm}%
\TextAt{8}{132.976}{215.267}{map}%
\TextAt{8}{158.554}{215.267}{is}%
\TextAt{8}{169.881}{215.267}{continuous,}%
\TextAt{8}{231.388}{215.267}{and}%
\TextAt{13}{253.714}{215.267}{C}%
\TextAt{17}{262.947}{219.605}{0}%
\TextAt{21}{262.089}{212.246}{L}%
\TextAt{8}{271.807}{215.267}{is}%
\TextAt{8}{283.134}{215.267}{compact}%
\TextAt{8}{329.526}{215.267}{by}%
\TextAt{8}{345.348}{215.267}{(1.6).}%
\TextAt{8}{377.729}{215.267}{Hence}%
\TextAt{13}{412.07}{215.267}{N}%
\TextAt{21}{421.443}{213.407}{L}%
\TextAt{3}{427.202}{213.407}{\Character{106}}%
\TextAt{21}{429.555}{213.407}{K}%
\TextAt{13}{437.757}{215.267}{C}%
\TextAt{17}{446.991}{219.605}{0}%
\TextAt{21}{446.132}{212.246}{L}%
\TextAt{8}{455.85}{215.267}{is}%
\TextAt{8}{467.178}{215.267}{closed.}%
\TextAt{8}{506.109}{215.267}{Since}%
\TextAt{9}{79.226}{198.221}{\Character{0}}%
\TextAt{21}{86.54}{196.427}{K}%
\TextAt{8}{98.64}{198.221}{is}%
\TextAt{8}{110.404}{198.221}{clearly}%
\TextAt{8}{147.8}{198.221}{also}%
\TextAt{8}{171.28}{198.221}{closed}%
\TextAt{8}{205.807}{198.221}{in}%
\TextAt{13}{219.459}{198.221}{C}%
\TextAt{21}{227.834}{196.427}{K}%
\TextAt{8}{236.037}{198.221}{,}%
\TextAt{8}{243.186}{198.221}{we}%
\TextAt{8}{260.414}{198.221}{get}%
\TextAt{8}{279.924}{198.221}{that}%
\TextAt{13}{305.293}{198.221}{N}%
\TextAt{21}{314.667}{196.361}{L}%
\TextAt{3}{320.426}{196.361}{\Character{106}}%
\TextAt{21}{322.779}{196.361}{K}%
\TextAt{13}{330.981}{198.221}{C}%
\TextAt{21}{339.356}{196.427}{L}%
\TextAt{8}{349.511}{198.221}{is}%
\TextAt{8}{361.275}{198.221}{closed.}%
\TextAt{8}{97.159}{178.186}{Conversely}%
\TextAt{8}{156.849}{178.186}{let}%
\TextAt{16}{174.194}{178.186}{\Character{78}}%
\TextAt{8}{190.102}{178.186}{be}%
\TextAt{8}{206.467}{178.186}{a}%
\TextAt{8}{216.662}{178.186}{closed}%
\TextAt{8}{251.629}{178.186}{subgroup}%
\TextAt{8}{302.867}{178.186}{of}%
\TextAt{13}{316.637}{178.186}{C}%
\TextAt{21}{325.012}{176.393}{K}%
\TextAt{8}{337.553}{178.186}{of}%
\TextAt{8}{351.323}{178.186}{finite}%
\TextAt{8}{381.675}{178.186}{index.}%
\TextAt{8}{419.084}{178.186}{We}%
\TextAt{8}{439.669}{178.186}{have}%
\TextAt{8}{467.104}{178.186}{to}%
\TextAt{8}{481.855}{178.186}{show}%
\TextAt{8}{511.298}{178.186}{that}%
\TextAt{16}{79.226}{161.14}{\Character{78}}%
\TextAt{8}{95.259}{161.14}{is}%
\TextAt{8}{107.588}{161.14}{open}%
\TextAt{8}{136.438}{161.14}{in}%
\TextAt{8}{150.656}{161.14}{the}%
\TextAt{8}{171.377}{161.14}{norm}%
\TextAt{8}{202.511}{161.14}{topology,}%
\TextAt{8}{253.632}{161.14}{i.e.,}%
\TextAt{8}{276.443}{161.14}{contains}%
\TextAt{8}{322.927}{161.14}{a}%
\TextAt{8}{333.248}{161.14}{norm}%
\TextAt{8}{364.382}{161.14}{group}%
\TextAt{13}{398.122}{161.14}{N}%
\TextAt{21}{407.496}{159.28}{L}%
\TextAt{3}{413.255}{159.28}{\Character{106}}%
\TextAt{21}{415.607}{159.28}{K}%
\TextAt{13}{423.81}{161.14}{C}%
\TextAt{21}{432.185}{159.347}{L}%
\TextAt{8}{438.442}{161.14}{.}%
\TextAt{8}{448.585}{161.14}{For}%
\TextAt{8}{470.119}{161.14}{this}%
\TextAt{8}{493.506}{161.14}{we}%
\TextAt{8}{511.298}{161.14}{may}%
\TextAt{8}{79.226}{144.094}{assume}%
\TextAt{8}{120.605}{144.094}{that}%
\TextAt{8}{146.909}{144.094}{the}%
\TextAt{8}{168}{144.094}{index}%
\TextAt{13}{200.472}{144.094}{n}%
\TextAt{8}{212.291}{144.094}{is}%
\TextAt{8}{224.99}{144.094}{a}%
\TextAt{8}{235.68}{144.094}{prime}%
\TextAt{8}{269.778}{144.094}{power.}%
\TextAt{8}{311.277}{144.094}{For}%
\TextAt{8}{333.181}{144.094}{if}%
\TextAt{13}{344.839}{144.094}{n}%
\TextAt{9}{356.74}{144.094}{\Character{61}}%
\TextAt{13}{370.758}{144.094}{p}%
\TextAt{21}{376.633}{149.319}{\Character{23}}%
\TextAt{24}{380.796}{148.212}{1}%
\TextAt{17}{376.633}{141.163}{1}%
\TextAt{16}{387.438}{144.094}{\Character{1}}%
\TextAt{16}{392.734}{144.094}{\Character{1}}%
\TextAt{16}{398.042}{144.094}{\Character{1}}%
\TextAt{13}{403.378}{144.094}{p}%
\TextAt{21}{409.253}{148.432}{\Character{23}}%
\TextAt{22}{413.416}{147.436}{r}%
\TextAt{21}{409.253}{141.139}{r}%
\TextAt{8}{418.072}{144.094}{,}%
\TextAt{8}{426.389}{144.094}{and}%
\TextAt{16}{450.086}{144.094}{\Character{78}}%
\TextAt{21}{459.896}{142.301}{i}%
\NativeMathAt{468.19}{144.094}{11.955}{\subseteq}%
\TextAt{13}{482.402}{144.094}{C}%
\TextAt{21}{490.777}{142.301}{K}%
\TextAt{8}{503.812}{144.094}{is}%
\TextAt{8}{516.51}{144.094}{the}%
\TextAt{8}{79.226}{127.048}{group}%
\TextAt{8}{112.542}{127.048}{containing}%
\TextAt{16}{169.6}{127.048}{\Character{78}}%
\TextAt{8}{185.208}{127.048}{of}%
\TextAt{8}{198.678}{127.048}{index}%
\TextAt{13}{230.355}{127.048}{p}%
\TextAt{21}{236.231}{132.381}{\Character{23}}%
\TextAt{22}{240.393}{131.166}{i}%
\TextAt{21}{236.231}{123.984}{i}%
\TextAt{8}{244.053}{127.048}{,}%
\TextAt{8}{251.377}{127.048}{then}%
\TextAt{16}{278.177}{127.048}{\Character{78}}%
\TextAt{9}{293.307}{127.048}{\Character{61}}%
\TextAt{20}{305.971}{136.015}{\Character{84}}%
\TextAt{21}{315.933}{132.86}{r}%
\TextAt{21}{315.933}{123.561}{i}%
\TextAt{17}{318.817}{123.561}{\Character{61}1}%
\TextAt{16}{332.128}{127.048}{\Character{78}}%
\TextAt{21}{341.937}{125.255}{i}%
\TextAt{8}{345.318}{127.048}{,}%
\TextAt{8}{352.642}{127.048}{and}%
\TextAt{8}{375.545}{127.048}{if}%
\TextAt{8}{386.408}{127.048}{the}%
\TextAt{16}{406.705}{127.048}{\Character{78}}%
\TextAt{21}{416.514}{125.255}{i}%
\TextAt{8}{423.932}{127.048}{are}%
\TextAt{8}{443.583}{127.048}{open}%
\TextAt{8}{472.009}{127.048}{in}%
\TextAt{8}{485.801}{127.048}{the}%
\TextAt{8}{506.097}{127.048}{norm}%
\TextAt{8}{79.226}{110.002}{topology,}%
\TextAt{8}{129.641}{110.002}{then}%
\TextAt{8}{156.301}{110.002}{so}%
\TextAt{8}{170.671}{110.002}{is}%
\TextAt{16}{182.435}{110.002}{\Character{78}}%
\TextAt{8}{194.006}{110.002}{.}%
\TextAt{8}{97.159}{89.967}{Now}%
\TextAt{8}{124.149}{89.967}{let}%
\TextAt{13}{141.383}{89.967}{J}%
\TextAt{8}{153.187}{89.967}{be}%
\TextAt{8}{169.441}{89.967}{the}%
\TextAt{8}{189.927}{89.967}{preimage}%
\TextAt{8}{240.336}{89.967}{of}%
\TextAt{16}{253.996}{89.967}{\Character{78}}%
\TextAt{8}{269.794}{89.967}{with}%
\TextAt{8}{296.783}{89.967}{respect}%
\TextAt{8}{337.163}{89.967}{to}%
\TextAt{8}{351.803}{89.967}{the}%
\TextAt{8}{372.289}{89.967}{projection}%
\TextAt{13}{427.899}{89.967}{I}%
\TextAt{21}{433.062}{88.174}{K}%
\NativeArrowAt{445.147}{89.967}{11.955}{19.26}{\longrightarrow}%
\TextAt{13}{468.291}{89.967}{C}%
\TextAt{21}{476.666}{88.174}{K}%
\TextAt{8}{484.868}{89.967}{.}%
\TextAt{8}{494.305}{89.967}{Then}%
\TextAt{13}{525.192}{89.967}{J}%
\TextAt{8}{79.226}{72.921}{is}%
\TextAt{8}{90.766}{72.921}{open}%
\TextAt{8}{118.828}{72.921}{in}%
\TextAt{13}{132.257}{72.921}{I}%
\TextAt{21}{137.42}{71.128}{K}%
\TextAt{8}{149.296}{72.921}{because}%
\TextAt{16}{192.374}{72.921}{\Character{78}}%
\TextAt{8}{207.618}{72.921}{is}%
\TextAt{8}{219.158}{72.921}{open}%
\TextAt{8}{247.22}{72.921}{in}%
\TextAt{13}{260.649}{72.921}{C}%
\TextAt{21}{269.024}{71.128}{K}%
\TextAt{8}{280.9}{72.921}{(with}%
\TextAt{8}{311.892}{72.921}{respect}%
\TextAt{8}{351.717}{72.921}{to}%
\TextAt{8}{365.804}{72.921}{the}%
\TextAt{8}{385.736}{72.921}{natural}%
\TextAt{8}{426.495}{72.921}{topology).}%
\TextAt{8}{483.67}{72.921}{Therefore}%
\TextAt{8}{297.211}{43.033}{360}%
\Destination{78.23}{809.815}{nav.page.360}%
\Destination{79.23}{773.075}{nav.section.6.6}%
\Destination{79.36}{743.955}{nav.print.395}%
\Destination{79.36}{648.855}{nav.number.6.6.1}%
\Destination{79.23}{773.075}{outline.61}%
\Hotlink{107.873}{658.555}{6.858}{11.047}{nav.section.6.1}%
\Hotlink{167.314}{443.455}{13.995}{11.047}{nav.chapter.4}%
\Hotlink{362.819}{304.055}{6.858}{11.047}{nav.section.6.1}%
\Hotlink{188.737}{443.455}{25.078}{11.047}{nav.number.4.6.7}%
\Hotlink{315.261}{372.255}{25.078}{11.047}{nav.number.6.5.5}%
\Hotlink{344.848}{213.055}{25.078}{11.047}{nav.number.6.1.6}%
}

% Page 373 (body page 361)
\Page{%
\TextAt{2}{79.226}{796.943}{VI.}%
\TextAt{2}{98.012}{796.943}{Global}%
\TextAt{2}{133.248}{796.943}{Class}%
\TextAt{2}{161.841}{796.943}{Field}%
\TextAt{2}{189.572}{796.943}{Theory}%
\Vector{q 0 G 0 g 0.299 w 7.226 18.729 m 460.769 18.729 l S Q}%
\TextAt{13}{79.226}{761.124}{J}%
\TextAt{8}{90.701}{761.124}{contains}%
\TextAt{8}{136.62}{761.124}{a}%
\TextAt{8}{146.376}{761.124}{group}%
\TextAt{13}{249.721}{744.078}{U}%
\TextAt{21}{258.921}{749.014}{S}%
\TextAt{21}{257.672}{741.122}{K}%
\TextAt{9}{269.195}{744.078}{\Character{61}}%
\TextAt{20}{281.621}{755.435}{\Character{89}}%
\TextAt{27}{281.644}{730.064}{p}%
\NativeMathAt{285.632}{730.064}{7.97}{\in}%
\TextAt{21}{291.278}{730.064}{S}%
\TextAt{16}{296.897}{744.078}{\Character{102}}%
\TextAt{9}{302.874}{744.078}{1}%
\TextAt{16}{308.727}{744.078}{\Character{103}}%
\TextAt{16}{317.359}{744.078}{\Character{2}}%
\TextAt{20}{329.317}{755.435}{\Character{89}}%
\TextAt{27}{329.34}{729.532}{p}%
\TextAt{17}{334.27}{729.532}{\Character{47}}%
\NativeMathAt{333.329}{729.532}{7.97}{\in}%
\TextAt{21}{338.974}{729.532}{S}%
\TextAt{13}{346.585}{744.078}{U}%
\TextAt{27}{354.536}{742.284}{p}%
\TextAt{13}{359.023}{744.078}{\Character{59}}%
\TextAt{8}{79.226}{710.941}{where}%
\TextAt{13}{112.56}{710.941}{S}%
\TextAt{8}{123.878}{710.941}{is}%
\TextAt{8}{135.167}{710.941}{a}%
\TextAt{8}{144.448}{710.941}{sufficiently}%
\TextAt{8}{202.565}{710.941}{big}%
\TextAt{8}{221.601}{710.941}{finite}%
\TextAt{8}{251.038}{710.941}{set}%
\TextAt{8}{268.831}{710.941}{of}%
\TextAt{8}{281.686}{710.941}{places}%
\TextAt{8}{315.738}{710.941}{of}%
\TextAt{13}{328.593}{710.941}{K}%
\TextAt{8}{342.833}{710.941}{containing}%
\TextAt{8}{399.277}{710.941}{the}%
\TextAt{8}{418.958}{710.941}{infinite}%
\TextAt{8}{458.151}{710.941}{ones}%
\TextAt{8}{483.75}{710.941}{and}%
\TextAt{8}{506.038}{710.941}{those}%
\TextAt{8}{79.226}{693.895}{primes}%
\TextAt{8}{116.379}{693.895}{that}%
\TextAt{8}{141.123}{693.895}{divide}%
\TextAt{13}{175.287}{693.895}{n}%
\TextAt{8}{182.275}{693.895}{,}%
\TextAt{8}{188.924}{693.895}{such}%
\TextAt{8}{214.695}{693.895}{that}%
\TextAt{13}{239.439}{693.895}{I}%
\TextAt{21}{244.602}{692.102}{K}%
\TextAt{9}{256.125}{693.895}{\Character{61}}%
\TextAt{13}{268.551}{693.895}{I}%
\TextAt{21}{274.653}{698.234}{S}%
\TextAt{21}{273.714}{690.875}{K}%
\TextAt{13}{281.916}{693.895}{K}%
\TextAt{3}{292.734}{698.234}{\Character{3}}%
\TextAt{8}{297.466}{693.895}{.}%
\TextAt{8}{305.706}{693.895}{Since}%
\TextAt{9}{335.638}{693.895}{\Character{40}}%
\TextAt{13}{340.191}{693.895}{I}%
\TextAt{21}{345.354}{692.102}{K}%
\TextAt{9}{356.877}{693.895}{\Character{58}}%
\TextAt{13}{363.45}{693.895}{J}%
\TextAt{9}{371.027}{693.895}{\Character{41}}%
\TextAt{9}{378.882}{693.895}{\Character{61}}%
\TextAt{13}{391.325}{693.895}{n}%
\TextAt{8}{398.313}{693.895}{,}%
\TextAt{13}{404.962}{693.895}{J}%
\TextAt{8}{415.811}{693.895}{also}%
\TextAt{8}{438.666}{693.895}{contains}%
\TextAt{8}{483.96}{693.895}{the}%
\TextAt{8}{503.491}{693.895}{group}%
\TextAt{20}{79.226}{685.816}{\Character{81}}%
\TextAt{27}{90.517}{673.363}{p}%
\NativeMathAt{94.506}{673.362}{7.97}{\in}%
\TextAt{21}{100.151}{673.363}{S}%
\TextAt{13}{108.238}{676.849}{K}%
\TextAt{3}{119.056}{681.188}{\Character{3}}%
\TextAt{21}{123.29}{681.188}{n}%
\TextAt{27}{118.197}{673.894}{p}%
\TextAt{16}{131.583}{676.85}{\Character{2}}%
\TextAt{20}{143.538}{685.816}{\Character{81}}%
\TextAt{27}{154.829}{673.363}{p}%
\TextAt{17}{159.759}{673.362}{\Character{47}}%
\NativeMathAt{158.818}{673.363}{7.97}{\in}%
\TextAt{21}{164.463}{673.362}{S}%
\TextAt{16}{170.557}{676.85}{\Character{102}}%
\TextAt{9}{176.535}{676.849}{1}%
\TextAt{16}{182.388}{676.85}{\Character{103}}%
\TextAt{8}{188.365}{676.85}{,}%
\TextAt{8}{195.514}{676.85}{and}%
\TextAt{8}{218.277}{676.85}{hence}%
\TextAt{8}{250.783}{676.85}{the}%
\TextAt{8}{270.94}{676.85}{group}%
\TextAt{13}{240.179}{646.351}{I}%
\TextAt{21}{245.342}{644.558}{K}%
\TextAt{9}{253.545}{646.351}{\Character{40}}%
\TextAt{13}{258.097}{646.351}{S}%
\TextAt{9}{265.992}{646.351}{\Character{41}}%
\TextAt{9}{273.847}{646.351}{\Character{61}}%
\TextAt{20}{286.291}{657.709}{\Character{89}}%
\TextAt{27}{286.314}{632.337}{p}%
\NativeMathAt{290.303}{632.337}{7.97}{\in}%
\TextAt{21}{295.948}{632.337}{S}%
\TextAt{13}{303.559}{646.351}{K}%
\TextAt{3}{314.377}{651.287}{\Character{3}}%
\TextAt{21}{318.611}{651.287}{n}%
\TextAt{27}{313.518}{643.396}{p}%
\TextAt{16}{326.904}{646.351}{\Character{2}}%
\TextAt{20}{338.859}{657.709}{\Character{89}}%
\TextAt{27}{338.882}{631.806}{p}%
\TextAt{17}{343.812}{631.806}{\Character{47}}%
\NativeMathAt{342.871}{631.806}{7.97}{\in}%
\TextAt{21}{348.517}{631.806}{S}%
\TextAt{13}{356.128}{646.351}{U}%
\TextAt{27}{364.078}{644.558}{p}%
\TextAt{13}{368.565}{646.351}{\Character{58}}%
\TextAt{8}{79.226}{607.186}{Thus}%
\TextAt{8}{109.432}{607.186}{it}%
\TextAt{8}{121.694}{607.186}{is}%
\TextAt{8}{134.015}{607.186}{enough}%
\TextAt{8}{174.897}{607.186}{to}%
\TextAt{8}{189.765}{607.186}{show}%
\TextAt{8}{219.326}{607.186}{that}%
\TextAt{13}{245.252}{607.186}{C}%
\TextAt{21}{253.627}{605.392}{K}%
\TextAt{9}{261.829}{607.186}{\Character{40}}%
\TextAt{13}{266.382}{607.186}{S}%
\TextAt{9}{274.277}{607.186}{\Character{41}}%
\TextAt{9}{283.088}{607.186}{\Character{61}}%
\TextAt{13}{296.475}{607.186}{I}%
\TextAt{21}{301.638}{605.392}{K}%
\TextAt{9}{309.841}{607.186}{\Character{40}}%
\TextAt{13}{314.393}{607.186}{S}%
\TextAt{9}{322.289}{607.186}{\Character{41}}%
\TextAt{13}{326.841}{607.186}{K}%
\TextAt{3}{337.659}{611.524}{\Character{3}}%
\TextAt{9}{342.391}{607.186}{\Character{47}}%
\TextAt{13}{348.244}{607.186}{K}%
\TextAt{3}{359.062}{611.524}{\Character{3}}%
\NativeMathAt{368.065}{607.186}{11.955}{\subseteq}%
\TextAt{16}{381.622}{607.186}{\Character{78}}%
\TextAt{8}{397.659}{607.186}{contains}%
\TextAt{8}{444.137}{607.186}{a}%
\TextAt{8}{454.449}{607.186}{norm}%
\TextAt{8}{485.576}{607.186}{group.}%
\TextAt{8}{524.975}{607.186}{If}%
\TextAt{8}{79.226}{590.14}{the}%
\TextAt{13}{99.306}{590.14}{n}%
\TextAt{8}{106.293}{590.14}{-th}%
\TextAt{8}{125.07}{590.14}{roots}%
\TextAt{8}{154.654}{590.14}{of}%
\TextAt{8}{167.908}{590.14}{unity}%
\TextAt{8}{198.4}{590.14}{belong}%
\TextAt{8}{235.719}{590.14}{to}%
\TextAt{13}{249.953}{590.14}{K}%
\TextAt{8}{260.771}{590.14}{,}%
\TextAt{8}{267.859}{590.14}{then}%
\TextAt{13}{294.442}{590.14}{C}%
\TextAt{21}{302.817}{588.346}{K}%
\TextAt{9}{311.02}{590.14}{\Character{40}}%
\TextAt{13}{315.572}{590.14}{S}%
\TextAt{9}{323.467}{590.14}{\Character{41}}%
\TextAt{9}{331.322}{590.14}{\Character{61}}%
\TextAt{13}{343.766}{590.14}{N}%
\TextAt{21}{353.14}{588.28}{L}%
\TextAt{3}{358.899}{588.28}{\Character{106}}%
\TextAt{21}{361.251}{588.28}{K}%
\TextAt{13}{369.454}{590.14}{C}%
\TextAt{21}{377.829}{588.346}{L}%
\TextAt{8}{387.907}{590.14}{with}%
\TextAt{13}{414.49}{590.14}{L}%
\TextAt{9}{425.775}{590.14}{\Character{61}}%
\TextAt{13}{438.201}{590.14}{K}%
\TextAt{9}{449.019}{590.14}{\Character{40}}%
\TextAt{22}{456.892}{596.412}{n}%
\TextAt{16}{454.851}{600.627}{\Character{112}}%
\Vector{q 0 G 0 g 0.478 w 392.813 -173.134 m 409.725 -173.134 l S Q}%
\TextAt{13}{464.813}{590.14}{K}%
\TextAt{21}{475.631}{593.593}{S}%
\TextAt{9}{481.725}{590.14}{\Character{41}}%
\TextAt{8}{486.277}{590.14}{,}%
\TextAt{8}{493.365}{590.14}{because}%
\TextAt{8}{79.226}{573.094}{of}%
\TextAt{8}{93.937}{573.094}{the}%
\TextAt{8}{115.475}{573.094}{remark}%
\TextAt{8}{156.858}{573.094}{following}%
\TextAt{8}{207.674}{573.094}{(4.2).}%
\TextAt{8}{244.344}{573.094}{If}%
\TextAt{8}{257.417}{573.094}{they}%
\TextAt{8}{285.136}{573.094}{do}%
\TextAt{8}{302.776}{573.094}{not}%
\TextAt{8}{324.971}{573.094}{belong}%
\TextAt{8}{363.748}{573.094}{to}%
\TextAt{13}{379.44}{573.094}{K}%
\TextAt{8}{390.257}{573.094}{,}%
\TextAt{8}{399.133}{573.094}{then}%
\TextAt{8}{427.174}{573.094}{we}%
\TextAt{8}{445.783}{573.094}{adjoin}%
\TextAt{8}{482.611}{573.094}{them}%
\TextAt{8}{513.904}{573.094}{and}%
\TextAt{8}{79.226}{556.048}{obtain}%
\TextAt{8}{116.809}{556.048}{an}%
\TextAt{8}{134.224}{556.048}{extension}%
\TextAt{13}{187.146}{556.048}{K}%
\TextAt{3}{197.963}{560.386}{\Character{48}}%
\TextAt{16}{206.049}{556.048}{\Character{106}}%
\TextAt{13}{214.66}{556.048}{K}%
\TextAt{8}{225.478}{556.048}{.}%
\TextAt{8}{237.394}{556.048}{Let}%
\TextAt{13}{259.519}{556.048}{S}%
\TextAt{3}{267.414}{560.386}{\Character{48}}%
\TextAt{8}{275.262}{556.048}{be}%
\TextAt{8}{292.342}{556.048}{the}%
\TextAt{8}{313.655}{556.048}{set}%
\TextAt{8}{333.078}{556.048}{of}%
\TextAt{8}{347.564}{556.048}{primes}%
\TextAt{8}{386.498}{556.048}{of}%
\TextAt{13}{400.983}{556.048}{K}%
\TextAt{3}{411.801}{560.386}{\Character{48}}%
\TextAt{8}{419.649}{556.048}{lying}%
\TextAt{8}{449.749}{556.048}{above}%
\TextAt{8}{484.08}{556.048}{primes}%
\TextAt{8}{523.014}{556.048}{in}%
\TextAt{13}{79.226}{539.002}{S}%
\TextAt{8}{87.121}{539.002}{.}%
\TextAt{8}{97.223}{539.002}{If}%
\TextAt{13}{109.466}{539.002}{S}%
\TextAt{8}{121.809}{539.002}{was}%
\TextAt{8}{144.86}{539.002}{chosen}%
\TextAt{8}{182.866}{539.002}{sufficiently}%
\TextAt{8}{242.01}{539.002}{large,}%
\TextAt{8}{274.571}{539.002}{then}%
\TextAt{13}{301.782}{539.002}{I}%
\TextAt{21}{306.945}{537.208}{K}%
\TextAt{23}{314.649}{539.478}{\Character{48}}%
\TextAt{9}{322.091}{539.002}{\Character{61}}%
\TextAt{13}{335.456}{539.002}{I}%
\TextAt{21}{341.559}{543.34}{S}%
\TextAt{23}{347.154}{546.153}{\Character{48}}%
\TextAt{21}{340.619}{535.837}{K}%
\TextAt{23}{348.323}{538.107}{\Character{48}}%
\TextAt{13}{351.506}{539.002}{K}%
\TextAt{3}{362.323}{543.34}{\Character{48}\Character{3}}%
\TextAt{8}{373.801}{539.002}{and}%
\TextAt{13}{397.115}{539.002}{C}%
\TextAt{21}{405.49}{537.208}{K}%
\TextAt{23}{413.194}{539.478}{\Character{48}}%
\TextAt{9}{416.377}{539.002}{\Character{40}}%
\TextAt{13}{420.929}{539.002}{S}%
\TextAt{3}{428.824}{543.34}{\Character{48}}%
\TextAt{9}{431.619}{539.002}{\Character{41}}%
\TextAt{9}{440.418}{539.002}{\Character{61}}%
\TextAt{13}{453.796}{539.002}{N}%
\TextAt{21}{463.17}{537.142}{L}%
\TextAt{23}{468.929}{539.411}{\Character{48}}%
\TextAt{3}{471.613}{537.142}{\Character{106}}%
\TextAt{21}{473.965}{537.142}{K}%
\TextAt{23}{481.67}{539.411}{\Character{48}}%
\TextAt{13}{484.852}{539.002}{C}%
\TextAt{21}{493.227}{537.208}{L}%
\TextAt{23}{498.986}{539.478}{\Character{48}}%
\TextAt{8}{502.169}{539.002}{,}%
\TextAt{8}{510.007}{539.002}{with}%
\TextAt{13}{79.226}{521.956}{L}%
\TextAt{3}{87.191}{526.294}{\Character{48}}%
\TextAt{9}{94.388}{521.956}{\Character{61}}%
\TextAt{13}{107.894}{521.956}{K}%
\TextAt{3}{118.712}{526.294}{\Character{48}}%
\TextAt{9}{121.507}{521.956}{\Character{40}}%
\TextAt{22}{129.38}{528.315}{n}%
\TextAt{16}{127.339}{532.515}{\Character{112}}%
\Vector{q 0 G 0 g 0.478 w 65.302 -241.246 m 87.195 -241.246 l S Q}%
\TextAt{13}{137.302}{521.956}{K}%
\TextAt{3}{148.12}{525.409}{\Character{48}}%
\TextAt{21}{150.417}{525.409}{S}%
\TextAt{23}{156.012}{527.679}{\Character{48}}%
\TextAt{9}{159.195}{521.956}{\Character{41}}%
\TextAt{8}{163.747}{521.956}{,}%
\TextAt{8}{171.689}{521.956}{by}%
\TextAt{8}{188.583}{521.956}{the}%
\TextAt{8}{209.374}{521.956}{above}%
\TextAt{8}{243.184}{521.956}{argument.}%
\TextAt{8}{302.001}{521.956}{Using}%
\TextAt{8}{335.536}{521.956}{chap.}%
\TextAt{8}{367.063}{521.956}{V,}%
\TextAt{8}{383.622}{521.956}{(1.5),}%
\TextAt{8}{415.642}{521.956}{this}%
\TextAt{8}{439.098}{521.956}{gives}%
\TextAt{8}{468.413}{521.956}{on}%
\TextAt{8}{485.307}{521.956}{the}%
\TextAt{8}{506.097}{521.956}{other}%
\TextAt{8}{79.226}{504.91}{hand}%
\TextAt{8}{108.492}{504.91}{that}%
\TextAt{13}{133.861}{504.91}{N}%
\TextAt{21}{143.235}{503.05}{K}%
\TextAt{23}{150.94}{505.319}{\Character{48}}%
\TextAt{3}{153.624}{503.05}{\Character{106}}%
\TextAt{21}{155.976}{503.05}{K}%
\TextAt{9}{164.179}{504.91}{\Character{40}}%
\TextAt{13}{168.731}{504.91}{I}%
\TextAt{21}{173.894}{503.117}{K}%
\TextAt{23}{181.599}{505.386}{\Character{48}}%
\TextAt{9}{184.781}{504.91}{\Character{40}}%
\TextAt{13}{189.334}{504.91}{S}%
\TextAt{3}{197.229}{509.248}{\Character{48}}%
\TextAt{9}{200.024}{504.91}{\Character{41}\Character{41}}%
\NativeMathAt{212.449}{504.91}{11.955}{\subseteq}%
\TextAt{13}{225.069}{504.91}{I}%
\TextAt{21}{230.232}{503.117}{K}%
\TextAt{9}{238.434}{504.91}{\Character{40}}%
\TextAt{13}{242.987}{504.91}{S}%
\TextAt{9}{250.882}{504.91}{\Character{41}}%
\TextAt{8}{255.434}{504.91}{,}%
\TextAt{8}{262.584}{504.91}{so}%
\TextAt{8}{276.954}{504.91}{that}%
\TextAt{13}{157.885}{475.909}{N}%
\TextAt{21}{167.259}{474.049}{L}%
\TextAt{23}{173.018}{476.318}{\Character{48}}%
\TextAt{3}{175.702}{474.049}{\Character{106}}%
\TextAt{21}{178.054}{474.049}{K}%
\TextAt{13}{186.257}{475.909}{C}%
\TextAt{21}{194.632}{474.115}{L}%
\TextAt{23}{200.391}{476.385}{\Character{48}}%
\TextAt{9}{206.894}{475.909}{\Character{61}}%
\TextAt{13}{219.32}{475.909}{N}%
\TextAt{21}{228.694}{474.049}{K}%
\TextAt{23}{236.398}{476.318}{\Character{48}}%
\TextAt{3}{239.083}{474.049}{\Character{106}}%
\TextAt{21}{241.435}{474.049}{K}%
\TextAt{9}{249.638}{475.909}{\Character{40}}%
\TextAt{13}{254.19}{475.909}{N}%
\TextAt{21}{263.564}{474.049}{L}%
\TextAt{23}{269.323}{476.318}{\Character{48}}%
\TextAt{3}{272.007}{474.049}{\Character{106}}%
\TextAt{21}{274.36}{474.049}{K}%
\TextAt{23}{282.064}{476.318}{\Character{48}}%
\TextAt{13}{285.247}{475.909}{C}%
\TextAt{21}{293.622}{474.115}{L}%
\TextAt{23}{299.381}{476.385}{\Character{48}}%
\TextAt{9}{302.563}{475.909}{\Character{41}}%
\TextAt{9}{310.418}{475.909}{\Character{61}}%
\TextAt{13}{322.862}{475.909}{N}%
\TextAt{21}{332.236}{474.049}{K}%
\TextAt{23}{339.94}{476.318}{\Character{48}}%
\TextAt{3}{342.625}{474.049}{\Character{106}}%
\TextAt{21}{344.977}{474.049}{K}%
\TextAt{9}{353.18}{475.909}{\Character{40}}%
\TextAt{13}{357.732}{475.909}{C}%
\TextAt{21}{366.107}{474.115}{K}%
\TextAt{23}{373.811}{476.385}{\Character{48}}%
\TextAt{9}{376.994}{475.909}{\Character{40}}%
\TextAt{13}{381.546}{475.909}{S}%
\TextAt{3}{389.442}{480.845}{\Character{48}}%
\TextAt{9}{392.237}{475.909}{\Character{41}\Character{41}}%
\NativeMathAt{404.662}{475.909}{11.955}{\subseteq}%
\TextAt{13}{417.281}{475.909}{C}%
\TextAt{21}{425.657}{474.115}{K}%
\TextAt{9}{433.859}{475.909}{\Character{40}}%
\TextAt{13}{438.412}{475.909}{S}%
\TextAt{9}{446.307}{475.909}{\Character{41}}%
\TextAt{13}{450.859}{475.909}{\Character{58}}%
\TextAt{8}{79.226}{446.908}{This}%
\TextAt{8}{105.946}{446.908}{finishes}%
\TextAt{8}{147.036}{446.908}{the}%
\TextAt{8}{167.192}{446.908}{proof.}%
\Vector{q 0 G 0 g 0.399 w 452.334 -327.092 m 452.334 -319.023 l S Q}%
\Vector{q 0 G 0 g 0.399 w 452.534 -319.222 m 459.707 -319.222 l S Q}%
\Vector{q 0 G 0 g 0.399 w 452.534 -326.893 m 459.707 -326.893 l S Q}%
\Vector{q 0 G 0 g 0.399 w 459.906 -327.092 m 459.906 -319.023 l S Q}%
\TextAt{8}{97.159}{417.906}{The}%
\TextAt{8}{120.031}{417.906}{above}%
\TextAt{8}{152.025}{417.906}{theorem}%
\TextAt{8}{196.368}{417.906}{is}%
\TextAt{8}{206.951}{417.906}{called}%
\TextAt{8}{238.932}{417.906}{the}%
\TextAt{8}{257.907}{417.906}{“existence}%
\TextAt{8}{312.054}{417.906}{theorem”}%
\TextAt{8}{361.921}{417.906}{of}%
\TextAt{8}{374.069}{417.906}{global}%
\TextAt{8}{407.366}{417.906}{class}%
\TextAt{8}{433.622}{417.906}{field}%
\TextAt{8}{457.797}{417.906}{theory}%
\TextAt{8}{493.365}{417.906}{because}%
\TextAt{8}{79.226}{400.86}{its}%
\TextAt{8}{95.193}{400.86}{main}%
\TextAt{8}{124.108}{400.86}{assertion}%
\TextAt{8}{172.665}{400.86}{is}%
\TextAt{8}{184.078}{400.86}{the}%
\TextAt{8}{203.883}{400.86}{existence,}%
\TextAt{8}{256.659}{400.86}{for}%
\TextAt{8}{274.192}{400.86}{any}%
\TextAt{8}{295.958}{400.86}{given}%
\TextAt{8}{326.176}{400.86}{closed}%
\TextAt{8}{360.351}{400.86}{subgroup}%
\TextAt{16}{410.797}{400.86}{\Character{78}}%
\TextAt{8}{425.913}{400.86}{of}%
\TextAt{8}{438.892}{400.86}{finite}%
\TextAt{8}{468.452}{400.86}{index}%
\TextAt{8}{499.639}{400.86}{in}%
\TextAt{13}{512.94}{400.86}{C}%
\TextAt{21}{521.315}{399.067}{K}%
\TextAt{8}{529.518}{400.86}{,}%
\TextAt{8}{79.226}{383.814}{of}%
\TextAt{8}{93.713}{383.814}{an}%
\TextAt{8}{111.129}{383.814}{abelian}%
\TextAt{8}{152.933}{383.814}{extension}%
\TextAt{13}{205.856}{383.814}{L}%
\TextAt{16}{219.113}{383.814}{\Character{106}}%
\TextAt{13}{227.726}{383.814}{K}%
\TextAt{8}{243.598}{383.814}{such}%
\TextAt{8}{271.152}{383.814}{that}%
\TextAt{13}{297.678}{383.814}{N}%
\TextAt{21}{307.052}{381.955}{L}%
\TextAt{3}{312.811}{381.955}{\Character{106}}%
\TextAt{21}{315.163}{381.955}{K}%
\TextAt{13}{323.366}{383.814}{C}%
\TextAt{21}{331.741}{382.021}{L}%
\TextAt{9}{343.29}{383.814}{\Character{61}}%
\TextAt{16}{357.687}{383.814}{\Character{78}}%
\TextAt{8}{369.258}{383.814}{.}%
\TextAt{8}{381.177}{383.814}{This}%
\TextAt{8}{409.054}{383.814}{extension}%
\TextAt{13}{461.977}{383.814}{L}%
\TextAt{8}{474.996}{383.814}{is}%
\TextAt{8}{487.916}{383.814}{the}%
\TextAt{8}{509.23}{383.814}{class}%
\TextAt{8}{79.226}{366.768}{field}%
\TextAt{8}{104.593}{366.768}{for}%
\TextAt{16}{122.487}{366.768}{\Character{78}}%
\TextAt{8}{134.058}{366.768}{.}%
\TextAt{8}{142.535}{366.768}{The}%
\TextAt{8}{166.598}{366.768}{existence}%
\TextAt{8}{216.413}{366.768}{theorem}%
\TextAt{8}{261.947}{366.768}{gives}%
\TextAt{8}{290.637}{366.768}{a}%
\TextAt{8}{300.402}{366.768}{clear}%
\TextAt{8}{328.375}{366.768}{overview}%
\TextAt{8}{376.516}{366.768}{of}%
\TextAt{8}{389.855}{366.768}{all}%
\TextAt{8}{406.124}{366.768}{the}%
\TextAt{8}{426.29}{366.768}{abelian}%
\TextAt{8}{466.947}{366.768}{extensions}%
\TextAt{8}{523.337}{366.768}{of}%
\TextAt{13}{79.226}{349.722}{K}%
\TextAt{8}{94.564}{349.722}{once}%
\TextAt{8}{121.848}{349.722}{we}%
\TextAt{8}{139.698}{349.722}{bring}%
\TextAt{8}{170.891}{349.722}{in}%
\TextAt{8}{185.167}{349.722}{the}%
\TextAt{8}{205.947}{349.722}{congruence}%
\TextAt{8}{267.051}{349.722}{subgroups}%
\TextAt{13}{323.087}{349.722}{C}%
\TextAt{27}{332.32}{354.061}{m}%
\TextAt{21}{331.462}{346.702}{K}%
\TextAt{8}{344.185}{349.722}{of}%
\TextAt{13}{358.138}{349.722}{C}%
\TextAt{21}{366.513}{347.929}{K}%
\TextAt{8}{379.236}{349.722}{corresponding}%
\TextAt{8}{455.046}{349.722}{to}%
\TextAt{8}{469.979}{349.722}{the}%
\TextAt{8}{490.759}{349.722}{modules}%
\TextAt{14}{79.226}{332.677}{m}%
\TextAt{9}{92.278}{332.677}{\Character{61}}%
\TextAt{20}{105.267}{341.643}{\Character{81}}%
\TextAt{27}{116.558}{329.19}{p}%
\TextAt{25}{120.547}{329.19}{\Character{45}}%
\NativeMathAt{122.761}{329.19}{7.97}{\infty}%
\TextAt{14}{133.72}{332.677}{p}%
\TextAt{21}{139.703}{337.015}{n}%
\TextAt{28}{144.841}{336.019}{p}%
\TextAt{8}{153.65}{332.677}{(see}%
\TextAt{8}{177.449}{332.677}{(1.7)).}%
\TextAt{8}{215.523}{332.677}{They}%
\TextAt{8}{246.088}{332.677}{are}%
\TextAt{8}{265.93}{332.677}{closed}%
\TextAt{8}{300.787}{332.677}{of}%
\TextAt{8}{314.448}{332.677}{finite}%
\TextAt{8}{344.691}{332.677}{index}%
\TextAt{8}{376.56}{332.677}{by}%
\TextAt{8}{393.149}{332.677}{(1.8),}%
\TextAt{8}{424.79}{332.677}{and}%
\TextAt{8}{447.884}{332.677}{they}%
\TextAt{8}{474.552}{332.677}{prompt}%
\TextAt{8}{516.51}{332.677}{the}%
\TextAt{8}{79.226}{315.631}{following}%
\TextAt{8}{128.661}{315.631}{definition.}%
\TextAt{12}{79.226}{289.618}{Definition}%
\TextAt{12}{142.397}{289.618}{(6.2).}%
\TextAt{8}{178.045}{289.618}{The}%
\TextAt{8}{202.431}{289.618}{class}%
\TextAt{8}{230.201}{289.618}{field}%
\TextAt{13}{255.89}{289.618}{K}%
\TextAt{27}{266.708}{293.957}{m}%
\TextAt{16}{277.205}{289.618}{\Character{106}}%
\TextAt{13}{284.413}{289.618}{K}%
\TextAt{8}{299.461}{289.618}{for}%
\TextAt{8}{317.678}{289.618}{the}%
\TextAt{8}{338.167}{289.618}{congruence}%
\TextAt{8}{398.98}{289.618}{subgroup}%
\TextAt{13}{450.11}{289.618}{C}%
\TextAt{27}{459.344}{293.957}{m}%
\TextAt{21}{458.485}{286.598}{K}%
\TextAt{8}{470.918}{289.618}{is}%
\TextAt{8}{483.014}{289.618}{called}%
\TextAt{8}{516.51}{289.618}{the}%
\TextAt{12}{79.226}{272.572}{ray}%
\TextAt{12}{102.479}{272.572}{class}%
\TextAt{12}{133.837}{272.572}{field}%
\TextAt{12}{163.139}{272.572}{mod}%
\TextAt{14}{193.41}{272.572}{m}%
\TextAt{8}{202.577}{272.572}{.}%
\TextAt{8}{97.159}{246.56}{The}%
\TextAt{8}{120.735}{246.56}{Galois}%
\TextAt{8}{156.171}{246.56}{group}%
\TextAt{8}{188.869}{246.56}{of}%
\TextAt{8}{201.722}{246.56}{the}%
\TextAt{8}{221.401}{246.56}{ray}%
\TextAt{8}{241.092}{246.56}{class}%
\TextAt{8}{268.052}{246.56}{field}%
\TextAt{8}{292.931}{246.56}{is}%
\TextAt{8}{304.218}{246.56}{canonically}%
\TextAt{8}{364.556}{246.56}{isomorphic}%
\TextAt{8}{423.329}{246.56}{to}%
\TextAt{8}{437.162}{246.56}{the}%
\TextAt{8}{456.841}{246.56}{ray}%
\TextAt{8}{476.532}{246.56}{class}%
\TextAt{8}{503.491}{246.56}{group}%
\TextAt{8}{79.226}{229.514}{mod}%
\TextAt{14}{105.563}{229.514}{m}%
\TextAt{8}{114.73}{229.514}{:}%
\TextAt{13}{248.625}{212.468}{G}%
\TextAt{9}{257.859}{212.468}{\Character{40}}%
\TextAt{13}{262.411}{212.468}{K}%
\TextAt{27}{273.229}{217.404}{m}%
\TextAt{16}{283.159}{212.468}{\Character{106}}%
\TextAt{13}{289.8}{212.468}{K}%
\TextAt{9}{300.618}{212.468}{\Character{41}}%
\NativeMathAt{308.491}{212.468}{11.955}{\cong}%
\TextAt{13}{321.111}{212.468}{C}%
\TextAt{21}{329.486}{210.675}{K}%
\TextAt{9}{337.688}{212.468}{\Character{47}}%
\TextAt{13}{343.541}{212.468}{C}%
\TextAt{27}{352.775}{217.404}{m}%
\TextAt{21}{351.916}{209.512}{K}%
\TextAt{13}{360.119}{212.468}{\Character{58}}%
\TextAt{8}{79.226}{188.946}{One}%
\TextAt{8}{103.937}{188.946}{has}%
\TextAt{14}{237.278}{171.9}{m}%
\TextAt{16}{249.765}{171.9}{\Character{106}}%
\TextAt{14}{256.407}{171.9}{m}%
\TextAt{3}{265.574}{176.836}{\Character{48}}%
\NativeArrowAt{283.645}{171.9}{11.955}{19.067}{\Longrightarrow}%
\TextAt{13}{317.988}{171.9}{K}%
\TextAt{27}{328.806}{176.836}{m}%
\NativeMathAt{338.736}{171.9}{11.955}{\subseteq}%
\TextAt{13}{351.355}{171.9}{K}%
\TextAt{27}{362.173}{176.836}{m}%
\TextAt{23}{368.284}{179.649}{\Character{48}}%
\TextAt{13}{371.466}{171.9}{\Character{59}}%
\TextAt{8}{79.226}{148.379}{because}%
\TextAt{8}{122.731}{148.379}{clearly}%
\TextAt{13}{160.33}{148.379}{C}%
\TextAt{27}{169.564}{152.717}{m}%
\TextAt{21}{168.705}{145.358}{K}%
\NativeMathAt{180.575}{148.379}{11.955}{\supseteq}%
\TextAt{13}{193.541}{148.379}{C}%
\TextAt{27}{202.774}{152.717}{m}%
\TextAt{23}{208.885}{155.53}{\Character{48}}%
\TextAt{21}{201.916}{145.358}{K}%
\TextAt{8}{212.068}{148.379}{.}%
\TextAt{8}{221.126}{148.379}{Since}%
\TextAt{8}{251.887}{148.379}{the}%
\TextAt{8}{272.246}{148.379}{closed}%
\TextAt{8}{306.976}{148.379}{subgroups}%
\TextAt{8}{362.592}{148.379}{of}%
\TextAt{8}{376.125}{148.379}{finite}%
\TextAt{8}{406.24}{148.379}{index}%
\TextAt{8}{437.981}{148.379}{in}%
\TextAt{13}{451.837}{148.379}{C}%
\TextAt{21}{460.212}{146.585}{K}%
\TextAt{8}{472.515}{148.379}{are}%
\TextAt{8}{492.229}{148.379}{by}%
\TextAt{8}{508.692}{148.379}{(1.8)}%
\TextAt{8}{79.226}{131.333}{precisely}%
\TextAt{8}{127.083}{131.333}{those}%
\TextAt{8}{157.712}{131.333}{subgroups}%
\TextAt{8}{213.124}{131.333}{containing}%
\TextAt{8}{270.043}{131.333}{a}%
\TextAt{8}{279.798}{131.333}{congruence}%
\TextAt{8}{340.279}{131.333}{subgroup}%
\TextAt{13}{391.077}{131.333}{C}%
\TextAt{27}{400.31}{135.671}{m}%
\TextAt{21}{399.452}{128.312}{K}%
\TextAt{8}{407.654}{131.333}{,}%
\TextAt{8}{414.804}{131.333}{we}%
\TextAt{8}{432.031}{131.333}{get}%
\TextAt{8}{451.542}{131.333}{from}%
\TextAt{8}{479.182}{131.333}{(6.1)}%
\TextAt{8}{507.157}{131.333}{the}%
\TextAt{12}{79.226}{105.32}{Corollary}%
\TextAt{12}{139.887}{105.32}{(6.3).}%
\TextAt{8}{175.962}{105.32}{Every}%
\TextAt{8}{210.375}{105.32}{finite}%
\TextAt{8}{241.046}{105.32}{abelian}%
\TextAt{8}{282.453}{105.32}{extension}%
\TextAt{13}{334.978}{105.32}{L}%
\TextAt{16}{347.558}{105.32}{\Character{106}}%
\TextAt{13}{355.493}{105.32}{K}%
\TextAt{8}{370.968}{105.32}{is}%
\TextAt{8}{383.491}{105.32}{contained}%
\TextAt{8}{437.259}{105.32}{in}%
\TextAt{8}{451.671}{105.32}{a}%
\TextAt{8}{462.186}{105.32}{ray}%
\TextAt{8}{483.114}{105.32}{class}%
\TextAt{8}{511.31}{105.32}{field}%
\TextAt{13}{79.226}{88.274}{K}%
\TextAt{27}{90.044}{92.613}{m}%
\TextAt{16}{99.974}{88.274}{\Character{106}}%
\TextAt{13}{106.616}{88.274}{K}%
\TextAt{8}{117.433}{88.274}{.}%
\TextAt{8}{297.211}{43.033}{361}%
\Destination{78.23}{809.815}{nav.page.361}%
\Destination{79.36}{725.955}{nav.print.396}%
\Destination{79.36}{303.655}{nav.number.6.6.2}%
\Destination{79.36}{119.355}{nav.number.6.6.3}%
\Hotlink{366.563}{519.755}{9.775}{11.047}{nav.chapter.5}%
\Hotlink{207.174}{570.855}{25.078}{11.047}{nav.number.6.4.2}%
\Hotlink{383.122}{519.755}{25.078}{11.047}{nav.number.5.1.5}%
\Hotlink{176.949}{330.455}{25.078}{11.047}{nav.number.6.1.7}%
\Hotlink{392.649}{330.455}{25.078}{11.047}{nav.number.6.1.8}%
\Hotlink{508.192}{146.155}{25.078}{11.047}{nav.number.6.1.8}%
\Hotlink{478.682}{129.155}{25.078}{11.047}{nav.number.6.6.1}%
}

% Page 374 (body page 362)
\Page{%
\TextAt{2}{79.226}{796.943}{VI.}%
\TextAt{2}{98.012}{796.943}{Global}%
\TextAt{2}{133.248}{796.943}{Class}%
\TextAt{2}{161.841}{796.943}{Field}%
\TextAt{2}{189.572}{796.943}{Theory}%
\Vector{q 0 G 0 g 0.299 w 7.226 18.729 m 460.769 18.729 l S Q}%
\TextAt{12}{79.226}{761.124}{Definition}%
\TextAt{12}{142.566}{761.124}{(6.4).}%
\TextAt{8}{178.361}{761.124}{Let}%
\TextAt{13}{199.809}{761.124}{L}%
\TextAt{16}{211.911}{761.124}{\Character{106}}%
\TextAt{13}{219.369}{761.124}{K}%
\TextAt{8}{234.564}{761.124}{be}%
\TextAt{8}{250.967}{761.124}{a}%
\TextAt{8}{261.202}{761.124}{finite}%
\TextAt{8}{291.593}{761.124}{abelian}%
\TextAt{8}{332.719}{761.124}{extension,}%
\TextAt{8}{388.336}{761.124}{and}%
\TextAt{8}{411.578}{761.124}{let}%
\TextAt{16}{428.961}{761.124}{\Character{78}}%
\TextAt{21}{438.771}{759.33}{L}%
\TextAt{9}{449.165}{761.124}{\Character{61}}%
\TextAt{13}{462.407}{761.124}{N}%
\TextAt{21}{471.781}{759.264}{L}%
\TextAt{3}{477.54}{759.264}{\Character{106}}%
\TextAt{21}{479.892}{759.264}{K}%
\TextAt{13}{488.095}{761.124}{C}%
\TextAt{21}{496.47}{759.33}{L}%
\TextAt{8}{502.727}{761.124}{.}%
\TextAt{8}{512.613}{761.124}{The}%
\TextAt{12}{79.226}{744.078}{conductor}%
\TextAt{14}{142.485}{744.078}{f}%
\TextAt{8}{151.077}{744.078}{of}%
\TextAt{13}{165.2}{744.078}{L}%
\TextAt{16}{177.837}{744.078}{\Character{106}}%
\TextAt{13}{185.831}{744.078}{K}%
\TextAt{8}{201.339}{744.078}{(or}%
\TextAt{8}{220.998}{744.078}{of}%
\TextAt{16}{235.121}{744.078}{\Character{78}}%
\TextAt{21}{244.931}{742.284}{L}%
\TextAt{8}{251.188}{744.078}{)}%
\TextAt{8}{260.433}{744.078}{is}%
\TextAt{8}{272.991}{744.078}{the}%
\TextAt{8}{293.94}{744.078}{gcd}%
\TextAt{8}{316.193}{744.078}{of}%
\TextAt{8}{330.317}{744.078}{all}%
\TextAt{8}{347.369}{744.078}{modules}%
\TextAt{14}{394.07}{744.078}{m}%
\TextAt{8}{407.928}{744.078}{such}%
\TextAt{8}{435.118}{744.078}{that}%
\TextAt{13}{461.28}{744.078}{L}%
\NativeMathAt{473.918}{744.078}{11.955}{\subseteq}%
\TextAt{13}{487.889}{744.078}{K}%
\TextAt{27}{498.707}{748.416}{m}%
\TextAt{8}{510.007}{744.078}{(i.e.,}%
\TextAt{13}{79.226}{727.032}{C}%
\TextAt{27}{88.46}{731.37}{m}%
\TextAt{21}{87.601}{724.011}{K}%
\NativeMathAt{99.125}{727.032}{11.955}{\subseteq}%
\TextAt{16}{111.726}{727.032}{\Character{78}}%
\TextAt{21}{121.553}{725.238}{L}%
\TextAt{8}{127.81}{727.032}{).}%
\TextAt{13}{97.159}{701.137}{K}%
\TextAt{27}{107.977}{705.475}{f}%
\TextAt{16}{114.534}{701.137}{\Character{106}}%
\TextAt{13}{121.314}{701.137}{K}%
\TextAt{8}{136.111}{701.137}{is}%
\TextAt{8}{147.956}{701.137}{therefore}%
\TextAt{8}{197.137}{701.137}{the}%
\TextAt{8}{217.374}{701.137}{smallest}%
\TextAt{8}{262.455}{701.137}{ray}%
\TextAt{8}{282.704}{701.137}{class}%
\TextAt{8}{310.223}{701.137}{field}%
\TextAt{8}{335.661}{701.137}{containing}%
\TextAt{13}{392.661}{701.137}{L}%
\TextAt{16}{404.084}{701.137}{\Character{106}}%
\TextAt{13}{410.864}{701.137}{K}%
\TextAt{8}{421.682}{701.137}{.}%
\TextAt{8}{430.374}{701.137}{But}%
\TextAt{8}{453.696}{701.137}{it}%
\TextAt{8}{465.481}{701.137}{is}%
\TextAt{8}{477.326}{701.137}{not}%
\TextAt{8}{498.222}{701.137}{true}%
\TextAt{8}{523.014}{701.137}{in}%
\TextAt{8}{79.226}{684.091}{general}%
\TextAt{8}{119.897}{684.091}{that}%
\TextAt{14}{145.612}{684.091}{m}%
\TextAt{8}{159.022}{684.091}{is}%
\TextAt{8}{171.131}{684.091}{the}%
\TextAt{8}{191.634}{684.091}{conductor}%
\TextAt{8}{246.615}{684.091}{of}%
\TextAt{13}{260.291}{684.091}{K}%
\TextAt{27}{271.108}{688.429}{m}%
\TextAt{16}{281.628}{684.091}{\Character{106}}%
\TextAt{13}{288.859}{684.091}{K}%
\TextAt{8}{299.677}{684.091}{.}%
\TextAt{8}{309.163}{684.091}{In}%
\TextAt{8}{324.13}{684.091}{chap.}%
\TextAt{8}{355.368}{684.091}{V,}%
\TextAt{8}{371.638}{684.091}{(1.6),}%
\TextAt{8}{403.297}{684.091}{we}%
\TextAt{8}{420.871}{684.091}{defined}%
\TextAt{8}{461.529}{684.091}{the}%
\TextAt{8}{482.032}{684.091}{conductor}%
\TextAt{14}{79.226}{667.045}{f}%
\TextAt{27}{83.127}{665.251}{p}%
\TextAt{8}{91.241}{667.045}{of}%
\TextAt{8}{104.302}{667.045}{a}%
\TextAt{14}{113.788}{667.045}{p}%
\TextAt{8}{119.771}{667.045}{-adic}%
\TextAt{8}{148.11}{667.045}{extension}%
\TextAt{13}{199.606}{667.045}{L}%
\TextAt{27}{207.571}{665.251}{p}%
\TextAt{16}{215.379}{667.045}{\Character{106}}%
\TextAt{13}{222.02}{667.045}{K}%
\TextAt{27}{231.979}{665.251}{p}%
\TextAt{8}{236.466}{667.045}{,}%
\TextAt{8}{243.4}{667.045}{for}%
\TextAt{8}{261.015}{667.045}{a}%
\TextAt{8}{270.501}{667.045}{finite}%
\TextAt{8}{300.143}{667.045}{place}%
\TextAt{14}{329.786}{667.045}{p}%
\TextAt{8}{335.769}{667.045}{,}%
\TextAt{8}{342.702}{667.045}{to}%
\TextAt{8}{356.743}{667.045}{be}%
\TextAt{8}{372.398}{667.045}{the}%
\TextAt{8}{392.285}{667.045}{smallest}%
\TextAt{8}{437.015}{667.045}{power}%
\TextAt{14}{470.889}{667.045}{f}%
\TextAt{27}{474.789}{665.251}{p}%
\TextAt{9}{482.597}{667.045}{\Character{61}}%
\TextAt{14}{495.023}{667.045}{p}%
\TextAt{21}{501.006}{671.383}{n}%
\TextAt{8}{510.27}{667.045}{such}%
\TextAt{8}{79.226}{649.999}{that}%
\TextAt{13}{105.254}{649.999}{U}%
\TextAt{17}{114.453}{656.109}{\Character{40}}%
\TextAt{21}{117.747}{656.109}{n}%
\TextAt{17}{122.885}{656.109}{\Character{41}}%
\TextAt{21}{113.205}{646.758}{K}%
\NativeMathAt{131.121}{649.999}{11.955}{\subseteq}%
\TextAt{13}{144.864}{649.999}{N}%
\TextAt{21}{154.238}{648.139}{L}%
\TextAt{28}{159.997}{647.143}{p}%
\TextAt{3}{164.079}{648.139}{\Character{106}}%
\TextAt{21}{166.432}{648.139}{K}%
\TextAt{28}{173.573}{647.143}{p}%
\TextAt{13}{178.154}{649.999}{L}%
\TextAt{3}{186.118}{654.337}{\Character{3}}%
\TextAt{27}{186.118}{647.043}{p}%
\TextAt{8}{190.851}{649.999}{.}%
\TextAt{8}{201.277}{649.999}{For}%
\TextAt{8}{222.906}{649.999}{an}%
\TextAt{8}{239.824}{649.999}{infinite}%
\TextAt{8}{280.151}{649.999}{place}%
\TextAt{14}{310.722}{649.999}{p}%
\TextAt{8}{321.262}{649.999}{we}%
\TextAt{8}{339.149}{649.999}{define}%
\TextAt{14}{373.618}{649.999}{f}%
\TextAt{27}{377.518}{648.205}{p}%
\TextAt{9}{386.449}{649.999}{\Character{61}}%
\TextAt{9}{399.982}{649.999}{1}%
\TextAt{8}{405.852}{649.999}{.}%
\TextAt{8}{416.278}{649.999}{Then}%
\TextAt{8}{447.495}{649.999}{we}%
\TextAt{8}{465.382}{649.999}{view}%
\TextAt{14}{493.024}{649.999}{f}%
\TextAt{8}{501.481}{649.999}{as}%
\TextAt{8}{516.511}{649.999}{the}%
\TextAt{8}{79.226}{632.953}{replete}%
\TextAt{8}{117.59}{632.953}{ideal}%
\TextAt{14}{145.553}{632.953}{f}%
\TextAt{20}{151.446}{641.919}{\Character{81}}%
\TextAt{27}{162.737}{629.466}{p}%
\TextAt{3}{166.726}{629.466}{\Character{106}}%
\NativeMathAt{169.077}{629.466}{7.97}{\infty}%
\TextAt{14}{180.037}{632.953}{p}%
\TextAt{17}{186.021}{637.291}{0}%
\TextAt{8}{194.65}{632.953}{and}%
\TextAt{8}{217.413}{632.953}{obtain}%
\TextAt{8}{253.84}{632.953}{the}%
\TextAt{12}{79.226}{607.057}{Proposition}%
\TextAt{12}{152.831}{607.057}{(6.5).}%
\TextAt{8}{188.923}{607.057}{If}%
\TextAt{14}{201.392}{607.057}{f}%
\TextAt{8}{209.966}{607.057}{is}%
\TextAt{8}{222.507}{607.057}{the}%
\TextAt{8}{243.44}{607.057}{conductor}%
\TextAt{8}{298.851}{607.057}{of}%
\TextAt{8}{312.958}{607.057}{the}%
\TextAt{8}{333.891}{607.057}{abelian}%
\TextAt{8}{375.315}{607.057}{extension}%
\TextAt{13}{427.858}{607.057}{L}%
\TextAt{16}{440.467}{607.058}{\Character{106}}%
\TextAt{13}{448.432}{607.057}{K}%
\TextAt{8}{459.249}{607.057}{,}%
\TextAt{8}{467.369}{607.057}{and}%
\TextAt{14}{490.909}{607.057}{f}%
\TextAt{27}{494.809}{605.264}{p}%
\TextAt{8}{503.97}{607.057}{is}%
\TextAt{8}{516.51}{607.057}{the}%
\TextAt{8}{79.226}{590.012}{conductor}%
\TextAt{8}{133.861}{590.012}{of}%
\TextAt{8}{147.191}{590.012}{the}%
\TextAt{8}{167.348}{590.012}{local}%
\TextAt{8}{194.988}{590.012}{extension}%
\TextAt{13}{246.754}{590.012}{L}%
\TextAt{27}{254.718}{588.218}{p}%
\TextAt{16}{262.526}{590.012}{\Character{106}}%
\TextAt{13}{269.168}{590.012}{K}%
\TextAt{27}{279.127}{588.218}{p}%
\TextAt{8}{283.614}{590.012}{,}%
\TextAt{8}{290.763}{590.012}{then}%
\TextAt{14}{281.721}{561.42}{f}%
\TextAt{9}{288.942}{561.42}{\Character{61}}%
\TextAt{20}{301.367}{572.778}{\Character{89}}%
\TextAt{27}{307.011}{547.672}{p}%
\TextAt{14}{318.636}{561.42}{f}%
\TextAt{27}{322.537}{559.627}{p}%
\TextAt{13}{327.023}{561.42}{\Character{58}}%
\TextAt{19}{79.226}{505.763}{Proof.}%
\TextAt{8}{116}{505.763}{Let}%
\TextAt{16}{137.285}{505.763}{\Character{78}}%
\TextAt{9}{152.714}{505.763}{\Character{61}}%
\TextAt{13}{165.676}{505.763}{N}%
\TextAt{21}{175.05}{503.904}{L}%
\TextAt{3}{180.809}{503.904}{\Character{106}}%
\TextAt{21}{183.161}{503.904}{K}%
\TextAt{13}{191.364}{505.763}{C}%
\TextAt{21}{199.739}{503.97}{L}%
\TextAt{8}{205.996}{505.763}{,}%
\TextAt{8}{213.539}{505.763}{and}%
\TextAt{8}{236.617}{505.763}{let}%
\TextAt{14}{253.837}{505.763}{m}%
\TextAt{9}{266.862}{505.763}{\Character{61}}%
\TextAt{20}{279.824}{514.73}{\Character{81}}%
\TextAt{14}{293.108}{505.763}{p}%
\TextAt{21}{299.091}{510.102}{n}%
\TextAt{28}{304.229}{509.106}{p}%
\TextAt{8}{313.023}{505.763}{be}%
\TextAt{8}{329.262}{505.763}{a}%
\TextAt{8}{339.333}{505.763}{module}%
\TextAt{8}{380.941}{505.763}{(}%
\TextAt{13}{385.496}{505.763}{n}%
\TextAt{27}{392.483}{503.97}{p}%
\TextAt{9}{400.828}{505.763}{\Character{61}}%
\TextAt{9}{413.776}{505.763}{0}%
\TextAt{8}{423.856}{505.763}{for}%
\TextAt{14}{442.057}{505.763}{p}%
\TextAt{16}{451.898}{505.763}{\Character{106}}%
\NativeMathAt{459.059}{505.763}{11.955}{\infty}%
\TextAt{8}{471.032}{505.763}{).}%
\TextAt{8}{484.98}{505.763}{One}%
\TextAt{8}{510.007}{505.763}{then}%
\TextAt{8}{79.226}{488.717}{has}%
\TextAt{13}{154.409}{471.671}{C}%
\TextAt{27}{163.643}{476.608}{m}%
\TextAt{21}{162.785}{468.716}{K}%
\NativeMathAt{174.308}{471.671}{11.955}{\subseteq}%
\TextAt{16}{186.909}{471.671}{\Character{78}}%
\NativeArrowAt{201.793}{471.671}{11.955}{21.938}{\Longleftrightarrow}%
\TextAt{14}{227.058}{471.671}{f}%
\TextAt{16}{234.279}{471.671}{\Character{106}}%
\TextAt{14}{240.921}{471.671}{m}%
\TextAt{8}{262.043}{471.671}{and}%
\TextAt{20}{294.856}{483.029}{\Character{89}}%
\TextAt{27}{300.499}{457.923}{p}%
\TextAt{14}{312.124}{471.671}{f}%
\TextAt{27}{316.025}{469.878}{p}%
\TextAt{16}{323.832}{471.671}{\Character{106}}%
\TextAt{14}{330.474}{471.671}{m}%
\NativeArrowAt{342.962}{471.671}{11.955}{21.926}{\Longleftrightarrow}%
\TextAt{14}{368.2}{471.671}{f}%
\TextAt{27}{372.101}{469.878}{p}%
\TextAt{16}{379.908}{471.671}{\Character{106}}%
\TextAt{14}{386.55}{471.671}{p}%
\TextAt{21}{392.533}{476.608}{n}%
\TextAt{28}{397.672}{475.611}{p}%
\TextAt{8}{414.208}{471.671}{for}%
\TextAt{8}{432.092}{471.671}{all}%
\TextAt{14}{448.351}{471.671}{p}%
\TextAt{13}{454.335}{471.671}{\Character{58}}%
\TextAt{8}{79.226}{438.958}{So}%
\TextAt{8}{95.485}{438.958}{to}%
\TextAt{8}{109.796}{438.958}{prove}%
\TextAt{14}{141.345}{438.958}{f}%
\TextAt{9}{148.567}{438.958}{\Character{61}}%
\TextAt{20}{160.992}{447.924}{\Character{81}}%
\TextAt{27}{172.283}{435.471}{p}%
\TextAt{14}{178.762}{438.958}{f}%
\TextAt{27}{182.663}{437.165}{p}%
\TextAt{8}{187.15}{438.958}{,}%
\TextAt{8}{194.299}{438.958}{we}%
\TextAt{8}{211.526}{438.958}{have}%
\TextAt{8}{238.521}{438.958}{to}%
\TextAt{8}{252.831}{438.958}{show}%
\TextAt{8}{281.835}{438.958}{the}%
\TextAt{8}{301.991}{438.958}{equivalence}%
\TextAt{13}{224.981}{410.367}{C}%
\TextAt{27}{234.214}{415.303}{m}%
\TextAt{21}{233.356}{407.411}{K}%
\NativeMathAt{244.879}{410.367}{11.955}{\subseteq}%
\TextAt{16}{257.48}{410.367}{\Character{78}}%
\NativeArrowAt{272.364}{410.367}{11.955}{21.938}{\Longleftrightarrow}%
\TextAt{14}{297.629}{410.367}{f}%
\TextAt{27}{301.529}{408.574}{p}%
\TextAt{16}{309.337}{410.367}{\Character{106}}%
\TextAt{14}{315.979}{410.367}{p}%
\TextAt{21}{321.962}{415.303}{n}%
\TextAt{28}{327.1}{414.307}{p}%
\TextAt{8}{343.636}{410.367}{for}%
\TextAt{8}{361.521}{410.367}{all}%
\TextAt{14}{377.78}{410.367}{p}%
\TextAt{13}{383.763}{410.367}{\Character{58}}%
\TextAt{8}{79.226}{381.776}{It}%
\TextAt{8}{91.899}{381.776}{follows}%
\TextAt{8}{130.334}{381.776}{from}%
\TextAt{8}{157.975}{381.776}{the}%
\TextAt{8}{178.131}{381.776}{identity}%
\TextAt{16}{221.385}{381.776}{\Character{78}}%
\TextAt{16}{235.6}{381.776}{\Character{92}}%
\TextAt{13}{246.239}{381.776}{K}%
\TextAt{3}{257.057}{386.114}{\Character{3}}%
\TextAt{27}{256.198}{378.82}{p}%
\TextAt{9}{265.11}{381.776}{\Character{61}}%
\TextAt{13}{277.536}{381.776}{N}%
\TextAt{27}{286.909}{379.983}{p}%
\TextAt{13}{291.396}{381.776}{L}%
\TextAt{3}{299.361}{386.114}{\Character{3}}%
\TextAt{27}{299.361}{378.82}{p}%
\TextAt{8}{307.991}{381.776}{(see}%
\TextAt{8}{331.459}{381.776}{(5.8)):}%
\TextAt{13}{131.413}{353.185}{C}%
\TextAt{27}{140.647}{358.121}{m}%
\TextAt{21}{139.788}{350.229}{K}%
\NativeMathAt{151.312}{353.185}{11.955}{\subseteq}%
\TextAt{16}{163.913}{353.185}{\Character{78}}%
\NativeArrowAt{178.797}{353.185}{11.955}{21.938}{\Longleftrightarrow}%
\TextAt{9}{204.061}{353.185}{\Character{40}}%
\TextAt{13}{208.614}{353.185}{\Character{11}}%
\NativeMathAt{219.456}{353.185}{11.955}{\in}%
\TextAt{13}{230.747}{353.185}{I}%
\TextAt{27}{236.85}{358.121}{m}%
\TextAt{21}{235.91}{350.229}{K}%
\NativeMathAt{247.434}{353.185}{11.955}{\Rightarrow}%
\Vector{q 0 G 0 g 0.478 w 190.71 -413.994 m 198.231 -413.994 l S Q}%
\TextAt{13}{262.71}{353.185}{\Character{11}}%
\NativeMathAt{273.552}{353.185}{11.955}{\in}%
\TextAt{16}{284.826}{353.185}{\Character{78}}%
\TextAt{9}{296.414}{353.185}{\Character{41}}%
\TextAt{8}{312.922}{353.185}{for}%
\TextAt{13}{330.806}{353.185}{\Character{11}}%
\NativeMathAt{341.649}{353.185}{11.955}{\in}%
\TextAt{13}{352.94}{353.185}{I}%
\TextAt{21}{358.103}{351.392}{K}%
\NativeArrowAt{178.823}{333.15}{11.955}{21.926}{\Longleftrightarrow}%
\TextAt{9}{204.061}{333.15}{\Character{40}}%
\TextAt{13}{208.614}{333.15}{\Character{11}}%
\TextAt{27}{216.059}{331.357}{p}%
\NativeMathAt{223.867}{333.15}{11.955}{\equiv}%
\TextAt{9}{236.486}{333.15}{1}%
\TextAt{8}{245.66}{333.15}{mod}%
\TextAt{14}{271.098}{333.15}{p}%
\TextAt{21}{277.081}{338.086}{n}%
\TextAt{28}{282.219}{337.09}{p}%
\NativeMathAt{290.121}{333.15}{11.955}{\Rightarrow}%
\TextAt{16}{305.388}{333.15}{\Character{104}}%
\TextAt{13}{310.046}{333.15}{\Character{11}}%
\TextAt{27}{317.492}{331.357}{p}%
\TextAt{16}{321.979}{333.15}{\Character{105}}%
\NativeMathAt{329.929}{333.15}{11.955}{\in}%
\TextAt{16}{341.215}{333.15}{\Character{78}}%
\TextAt{16}{355.442}{333.15}{\Character{92}}%
\TextAt{13}{366.094}{333.15}{K}%
\TextAt{3}{376.912}{338.086}{\Character{3}}%
\TextAt{27}{376.053}{330.195}{p}%
\TextAt{9}{384.965}{333.15}{\Character{61}}%
\TextAt{13}{397.39}{333.15}{N}%
\TextAt{27}{406.764}{331.357}{p}%
\TextAt{13}{411.251}{333.15}{L}%
\TextAt{3}{419.216}{338.086}{\Character{3}}%
\TextAt{27}{419.216}{330.195}{p}%
\TextAt{9}{423.948}{333.15}{\Character{41}}%
\TextAt{8}{440.455}{333.15}{for}%
\TextAt{8}{458.34}{333.15}{all}%
\TextAt{14}{474.599}{333.15}{p}%
\NativeArrowAt{178.823}{311.828}{11.955}{21.926}{\Longleftrightarrow}%
\TextAt{9}{204.061}{311.828}{\Character{40}}%
\TextAt{13}{208.614}{311.828}{\Character{11}}%
\TextAt{27}{216.059}{310.035}{p}%
\NativeMathAt{223.867}{311.828}{11.955}{\in}%
\TextAt{13}{235.158}{311.828}{U}%
\TextAt{17}{244.357}{318.074}{\Character{40}}%
\TextAt{21}{247.65}{318.074}{n}%
\TextAt{28}{252.789}{317.077}{p}%
\TextAt{17}{256.871}{318.074}{\Character{41}}%
\TextAt{27}{243.108}{309.858}{p}%
\NativeMathAt{263.984}{311.828}{11.955}{\Rightarrow}%
\TextAt{13}{279.26}{311.828}{\Character{11}}%
\TextAt{27}{286.705}{310.035}{p}%
\NativeMathAt{294.513}{311.828}{11.955}{\in}%
\TextAt{13}{305.804}{311.828}{N}%
\TextAt{27}{315.178}{310.035}{p}%
\TextAt{13}{319.664}{311.828}{L}%
\TextAt{3}{327.629}{316.765}{\Character{3}}%
\TextAt{27}{327.629}{308.873}{p}%
\TextAt{9}{332.361}{311.828}{\Character{41}}%
\NativeArrowAt{178.823}{290.507}{11.955}{21.926}{\Longleftrightarrow}%
\TextAt{13}{204.061}{290.507}{U}%
\TextAt{17}{213.261}{296.752}{\Character{40}}%
\TextAt{21}{216.554}{296.752}{n}%
\TextAt{28}{221.692}{295.756}{p}%
\TextAt{17}{225.775}{296.752}{\Character{41}}%
\TextAt{27}{212.012}{288.536}{p}%
\NativeMathAt{232.887}{290.507}{11.955}{\subseteq}%
\TextAt{13}{245.506}{290.507}{N}%
\TextAt{27}{254.88}{288.714}{p}%
\TextAt{13}{259.367}{290.507}{L}%
\TextAt{3}{267.331}{295.443}{\Character{3}}%
\TextAt{27}{267.331}{287.551}{p}%
\NativeArrowAt{275.385}{290.507}{11.955}{21.926}{\Longleftrightarrow}%
\TextAt{14}{300.623}{290.507}{f}%
\TextAt{27}{304.524}{288.714}{p}%
\TextAt{16}{312.331}{290.507}{\Character{106}}%
\TextAt{14}{318.973}{290.507}{p}%
\TextAt{21}{324.956}{295.443}{n}%
\TextAt{28}{330.094}{294.447}{p}%
\TextAt{13}{334.675}{290.507}{\Character{58}}%
\Vector{q 0 G 0 g 0.399 w 452.334 -512.084 m 452.334 -504.015 l S Q}%
\Vector{q 0 G 0 g 0.399 w 452.534 -504.214 m 459.707 -504.214 l S Q}%
\Vector{q 0 G 0 g 0.399 w 452.534 -511.885 m 459.707 -511.885 l S Q}%
\Vector{q 0 G 0 g 0.399 w 459.906 -512.084 m 459.906 -504.015 l S Q}%
\TextAt{8}{97.159}{233.032}{By}%
\TextAt{8}{115.765}{233.032}{chap.}%
\TextAt{8}{146.901}{233.032}{V,}%
\TextAt{8}{163.069}{233.032}{(1.7),}%
\TextAt{8}{194.6}{233.032}{the}%
\TextAt{8}{215}{233.032}{local}%
\TextAt{8}{242.883}{233.032}{extension}%
\TextAt{13}{294.893}{233.032}{L}%
\TextAt{27}{302.857}{231.238}{p}%
\TextAt{16}{311.079}{233.032}{\Character{106}}%
\TextAt{13}{318.136}{233.032}{K}%
\TextAt{27}{328.095}{231.238}{p}%
\TextAt{8}{332.582}{233.032}{,}%
\TextAt{8}{340.035}{233.032}{for}%
\TextAt{8}{358.164}{233.032}{a}%
\TextAt{8}{368.162}{233.032}{finite}%
\TextAt{8}{398.318}{233.032}{prime}%
\TextAt{14}{431.725}{233.032}{p}%
\TextAt{8}{437.708}{233.032}{,}%
\TextAt{8}{445.161}{233.032}{is}%
\TextAt{8}{457.168}{233.032}{ramified}%
\TextAt{8}{502.937}{233.032}{if}%
\TextAt{8}{513.904}{233.032}{and}%
\TextAt{8}{79.226}{215.986}{only}%
\TextAt{8}{104.443}{215.986}{if}%
\TextAt{8}{114.692}{215.986}{its}%
\TextAt{8}{130.536}{215.986}{conductor}%
\TextAt{14}{184.697}{215.986}{f}%
\TextAt{27}{188.597}{214.193}{p}%
\TextAt{8}{196.507}{215.986}{is}%
\NativeMathAt{207.796}{215.986}{11.955}{\neq}%
\TextAt{9}{220.206}{215.986}{1}%
\TextAt{8}{226.074}{215.986}{.}%
\TextAt{8}{234.364}{215.986}{This}%
\TextAt{8}{260.61}{215.986}{continues}%
\TextAt{8}{311.578}{215.986}{to}%
\TextAt{8}{325.414}{215.986}{hold}%
\TextAt{8}{350.953}{215.986}{also}%
\TextAt{8}{373.959}{215.986}{for}%
\TextAt{8}{391.369}{215.986}{an}%
\TextAt{8}{407.153}{215.986}{infinite}%
\TextAt{8}{446.346}{215.986}{place}%
\TextAt{14}{475.783}{215.986}{p}%
\TextAt{8}{481.766}{215.986}{,}%
\TextAt{8}{488.535}{215.986}{provided}%
\TextAt{8}{79.226}{198.94}{we}%
\TextAt{8}{96.868}{198.94}{call}%
\TextAt{8}{118.741}{198.94}{the}%
\TextAt{8}{139.312}{198.94}{extension}%
\TextAt{13}{191.492}{198.94}{L}%
\TextAt{27}{199.456}{197.147}{p}%
\TextAt{16}{207.969}{198.94}{\Character{106}}%
\TextAt{13}{215.317}{198.94}{K}%
\TextAt{27}{225.276}{197.147}{p}%
\TextAt{19}{234.074}{198.94}{unramified}%
\TextAt{8}{292.638}{198.94}{in}%
\TextAt{8}{306.705}{198.94}{this}%
\TextAt{8}{329.941}{198.94}{case,}%
\TextAt{8}{358.482}{198.94}{as}%
\TextAt{8}{373.266}{198.94}{we}%
\TextAt{8}{390.907}{198.94}{did}%
\TextAt{8}{411.478}{198.94}{in}%
\TextAt{8}{425.545}{198.94}{chap.}%
\TextAt{8}{456.851}{198.94}{III.}%
\TextAt{8}{477.72}{198.94}{Then}%
\TextAt{8}{508.692}{198.94}{(6.5)}%
\TextAt{8}{79.226}{181.894}{yields}%
\TextAt{8}{112.127}{181.894}{the}%
\TextAt{12}{79.226}{155.999}{Corollary}%
\TextAt{12}{139.014}{155.999}{(6.6).}%
\TextAt{8}{174.33}{155.999}{Let}%
\TextAt{13}{195.299}{155.999}{L}%
\TextAt{16}{206.584}{155.999}{\Character{106}}%
\TextAt{13}{213.226}{155.999}{K}%
\TextAt{8}{227.941}{155.999}{be}%
\TextAt{8}{243.865}{155.999}{a}%
\TextAt{8}{253.621}{155.999}{finite}%
\TextAt{8}{283.533}{155.999}{abelian}%
\TextAt{8}{324.18}{155.999}{extension}%
\TextAt{8}{375.946}{155.999}{and}%
\TextAt{14}{398.709}{155.999}{f}%
\TextAt{8}{406.507}{155.999}{its}%
\TextAt{8}{422.825}{155.999}{conductor.}%
\TextAt{8}{482.011}{155.999}{Then:}%
\TextAt{14}{223.821}{127.408}{p}%
\TextAt{8}{233.701}{127.408}{is}%
\TextAt{8}{245.465}{127.408}{ramified}%
\TextAt{8}{290.99}{127.408}{in}%
\TextAt{13}{304.643}{127.408}{L}%
\NativeArrowAt{327.883}{127.408}{11.955}{21.926}{\Longleftrightarrow}%
\TextAt{14}{365.077}{127.408}{p}%
\TextAt{16}{374.381}{127.408}{\Character{106}}%
\TextAt{14}{381.023}{127.408}{f}%
\TextAt{13}{384.924}{127.408}{\Character{58}}%
\TextAt{8}{97.159}{89.967}{In}%
\TextAt{8}{111.784}{89.967}{the}%
\TextAt{8}{131.944}{89.967}{case}%
\TextAt{8}{156.718}{89.967}{of}%
\TextAt{8}{170.052}{89.967}{the}%
\TextAt{8}{190.212}{89.967}{base}%
\TextAt{8}{216.29}{89.967}{field}%
\TextAt{18}{241.651}{89.967}{Q}%
\TextAt{8}{250.949}{89.967}{,}%
\TextAt{8}{258.103}{89.967}{the}%
\TextAt{8}{278.263}{89.967}{ray}%
\TextAt{8}{298.435}{89.967}{class}%
\TextAt{8}{325.876}{89.967}{fields}%
\TextAt{8}{355.851}{89.967}{are}%
\TextAt{8}{375.365}{89.967}{nothing}%
\TextAt{8}{418.3}{89.967}{but}%
\TextAt{8}{439.763}{89.967}{the}%
\TextAt{8}{459.923}{89.967}{familiar}%
\TextAt{8}{503.18}{89.967}{cyclo-}%
\TextAt{8}{79.226}{72.921}{tomic}%
\TextAt{8}{111.744}{72.921}{fields:}%
\TextAt{8}{297.211}{43.033}{362}%
\Destination{78.23}{809.815}{nav.page.362}%
\Destination{79.36}{776.155}{nav.print.397}%
\Destination{79.36}{776.155}{nav.print.398}%
\Destination{79.36}{775.155}{nav.number.6.6.4}%
\Destination{79.36}{621.055}{nav.number.6.6.5}%
\Destination{79.36}{169.955}{nav.number.6.6.6}%
\Hotlink{354.868}{681.855}{9.775}{11.047}{nav.chapter.5}%
\Hotlink{146.401}{230.855}{9.775}{11.047}{nav.chapter.5}%
\Hotlink{456.351}{196.755}{14.294}{11.047}{nav.chapter.3}%
\Hotlink{371.138}{681.855}{25.078}{11.047}{nav.number.5.1.6}%
\Hotlink{330.959}{379.555}{25.078}{11.047}{nav.number.6.5.8}%
\Hotlink{162.569}{230.855}{25.078}{11.047}{nav.number.5.1.7}%
\Hotlink{508.192}{196.755}{25.078}{11.047}{nav.number.3.6.5}%
}

% Page 375 (body page 363)
\Page{%
\TextAt{2}{79.226}{796.943}{VI.}%
\TextAt{2}{98.012}{796.943}{Global}%
\TextAt{2}{133.248}{796.943}{Class}%
\TextAt{2}{161.841}{796.943}{Field}%
\TextAt{2}{189.572}{796.943}{Theory}%
\Vector{q 0 G 0 g 0.299 w 7.226 18.729 m 460.769 18.729 l S Q}%
\TextAt{12}{79.226}{761.124}{Proposition}%
\TextAt{12}{152.285}{761.124}{(6.7).}%
\TextAt{8}{187.902}{761.124}{Let}%
\TextAt{13}{209.174}{761.124}{m}%
\TextAt{8}{223.612}{761.124}{be}%
\TextAt{8}{239.838}{761.124}{a}%
\TextAt{8}{249.896}{761.124}{natural}%
\TextAt{8}{291.18}{761.124}{number}%
\TextAt{8}{334.078}{761.124}{and}%
\TextAt{14}{357.143}{761.124}{m}%
\TextAt{9}{370.145}{761.124}{\Character{61}}%
\TextAt{9}{383.069}{761.124}{\Character{40}}%
\TextAt{13}{387.637}{761.124}{m}%
\TextAt{9}{397.876}{761.124}{\Character{41}}%
\TextAt{8}{402.429}{761.124}{.}%
\TextAt{8}{411.783}{761.124}{Then}%
\TextAt{8}{442.642}{761.124}{the}%
\TextAt{8}{463.101}{761.124}{ray}%
\TextAt{8}{483.571}{761.124}{class}%
\TextAt{8}{511.31}{761.124}{field}%
\TextAt{8}{79.226}{744.078}{mod}%
\TextAt{14}{105.563}{744.078}{m}%
\TextAt{8}{118.628}{744.078}{of}%
\TextAt{18}{131.958}{744.078}{Q}%
\TextAt{8}{145.153}{744.078}{is}%
\TextAt{8}{156.917}{744.078}{the}%
\TextAt{8}{177.074}{744.078}{field}%
\TextAt{18}{273.453}{727.032}{Q}%
\TextAt{27}{282.752}{731.968}{m}%
\TextAt{9}{292.682}{727.032}{\Character{61}}%
\TextAt{18}{305.107}{727.032}{Q}%
\TextAt{9}{314.406}{727.032}{\Character{40}}%
\TextAt{13}{318.958}{727.032}{\Character{22}}%
\TextAt{21}{326.001}{725.238}{m}%
\TextAt{9}{333.99}{727.032}{\Character{41}}%
\TextAt{8}{79.226}{703.826}{of}%
\TextAt{13}{92.556}{703.826}{m}%
\TextAt{8}{102.795}{703.826}{-th}%
\TextAt{8}{121.649}{703.826}{roots}%
\TextAt{8}{151.31}{703.826}{of}%
\TextAt{8}{164.64}{703.826}{unity.}%
\TextAt{19}{79.226}{669.269}{Proof.}%
\TextAt{8}{116}{669.269}{Let}%
\TextAt{13}{136.593}{669.269}{m}%
\TextAt{9}{150.153}{669.269}{\Character{61}}%
\TextAt{20}{162.578}{678.235}{\Character{81}}%
\TextAt{21}{173.869}{665.782}{p}%
\NativeMathAt{178.132}{665.782}{7.97}{\neq}%
\TextAt{21}{184.718}{665.782}{p}%
\NativeMathAt{188.981}{664.786}{5.978}{\infty}%
\TextAt{13}{199.276}{669.269}{p}%
\TextAt{21}{205.151}{673.607}{n}%
\TextAt{22}{210.289}{672.611}{p}%
\TextAt{8}{215.126}{669.269}{.}%
\TextAt{8}{223.449}{669.269}{Then}%
\TextAt{13}{253.629}{669.269}{I}%
\TextAt{27}{259.732}{673.607}{m}%
\TextAt{25}{258.792}{666.204}{Q}%
\TextAt{9}{269.662}{669.269}{\Character{61}}%
\TextAt{20}{282.087}{678.235}{\Character{81}}%
\TextAt{21}{293.378}{665.782}{p}%
\NativeMathAt{297.641}{665.782}{7.97}{\neq}%
\TextAt{21}{304.228}{665.782}{p}%
\NativeMathAt{308.49}{664.786}{5.978}{\infty}%
\TextAt{13}{318.785}{669.269}{U}%
\TextAt{17}{327.984}{675.545}{\Character{40}}%
\TextAt{21}{331.277}{675.545}{n}%
\TextAt{22}{336.416}{674.549}{p}%
\TextAt{17}{340.754}{675.545}{\Character{41}}%
\TextAt{21}{326.735}{668.042}{p}%
\TextAt{16}{346.431}{669.269}{\Character{2}}%
\TextAt{18}{357.615}{669.269}{R}%
\TextAt{3}{366.25}{673.607}{\Character{3}}%
\TextAt{17}{366.25}{666.313}{\Character{43}}%
\TextAt{8}{373.334}{669.269}{.}%
\TextAt{8}{381.657}{669.269}{Let}%
\TextAt{13}{402.249}{669.269}{m}%
\TextAt{9}{415.809}{669.269}{\Character{61}}%
\TextAt{13}{428.235}{669.269}{m}%
\TextAt{3}{438.474}{673.607}{\Character{48}}%
\TextAt{13}{441.269}{669.269}{p}%
\TextAt{21}{447.144}{673.607}{n}%
\TextAt{22}{452.282}{672.611}{p}%
\TextAt{8}{457.119}{669.269}{.}%
\TextAt{8}{465.442}{669.269}{Then}%
\TextAt{13}{495.622}{669.269}{U}%
\TextAt{17}{504.821}{675.545}{\Character{40}}%
\TextAt{21}{508.115}{675.545}{n}%
\TextAt{22}{513.253}{674.549}{p}%
\TextAt{17}{517.591}{675.545}{\Character{41}}%
\TextAt{21}{503.572}{668.042}{p}%
\TextAt{8}{524.903}{669.269}{is}%
\TextAt{8}{79.226}{652.223}{certainly}%
\TextAt{8}{127.481}{652.223}{contained}%
\TextAt{8}{180.291}{652.223}{in}%
\TextAt{8}{193.744}{652.223}{the}%
\TextAt{8}{213.701}{652.223}{norm}%
\TextAt{8}{244.071}{652.223}{group}%
\TextAt{8}{277.048}{652.223}{of}%
\TextAt{8}{290.178}{652.223}{the}%
\TextAt{8}{310.135}{652.223}{unramified}%
\TextAt{8}{368.468}{652.223}{extension}%
\TextAt{18}{420.035}{652.223}{Q}%
\TextAt{21}{429.334}{650.43}{p}%
\TextAt{9}{434.094}{652.223}{\Character{40}}%
\TextAt{13}{438.647}{652.223}{\Character{22}}%
\TextAt{21}{445.69}{650.43}{m}%
\TextAt{23}{453.18}{652.699}{\Character{48}}%
\TextAt{9}{456.363}{652.223}{\Character{41}}%
\TextAt{16}{464.236}{652.223}{\Character{106}}%
\TextAt{18}{470.878}{652.223}{Q}%
\TextAt{21}{480.176}{650.43}{p}%
\TextAt{8}{484.937}{652.223}{,}%
\TextAt{8}{491.927}{652.223}{but}%
\TextAt{8}{513.187}{652.223}{also}%
\TextAt{8}{79.226}{635.177}{in}%
\TextAt{8}{92.772}{635.177}{the}%
\TextAt{8}{112.821}{635.177}{norm}%
\TextAt{8}{143.283}{635.177}{group}%
\TextAt{8}{176.351}{635.177}{of}%
\TextAt{18}{189.574}{635.177}{Q}%
\TextAt{21}{198.872}{633.384}{p}%
\TextAt{9}{203.633}{635.177}{\Character{40}}%
\TextAt{13}{208.185}{635.177}{\Character{22}}%
\TextAt{21}{215.228}{633.384}{p}%
\TextAt{22}{219.491}{636.4}{n}%
\TextAt{22}{224.092}{635.404}{p}%
\TextAt{9}{229.427}{635.177}{\Character{41}}%
\TextAt{16}{237.3}{635.177}{\Character{106}}%
\TextAt{18}{243.942}{635.177}{Q}%
\TextAt{21}{253.24}{633.384}{p}%
\TextAt{8}{258.001}{635.177}{,}%
\TextAt{8}{265.064}{635.177}{according}%
\TextAt{8}{317.643}{635.177}{to}%
\TextAt{8}{331.846}{635.177}{chap.}%
\TextAt{8}{362.631}{635.177}{V,}%
\TextAt{8}{378.448}{635.177}{(1.8).}%
\TextAt{8}{410.938}{635.177}{This}%
\TextAt{8}{437.551}{635.177}{means,}%
\TextAt{8}{476.546}{635.177}{by}%
\TextAt{8}{492.698}{635.177}{§3,}%
\TextAt{8}{511.298}{635.177}{that}%
\TextAt{8}{79.226}{618.131}{every}%
\TextAt{8}{110.425}{618.131}{idèle}%
\TextAt{8}{138.038}{618.131}{in}%
\TextAt{13}{151.998}{618.131}{I}%
\TextAt{27}{158.101}{622.469}{m}%
\TextAt{25}{157.161}{615.066}{Q}%
\TextAt{8}{168.915}{618.131}{is}%
\TextAt{8}{180.986}{618.131}{a}%
\TextAt{8}{191.048}{618.131}{norm}%
\TextAt{8}{221.925}{618.131}{of}%
\TextAt{8}{235.562}{618.131}{some}%
\TextAt{8}{265.195}{618.131}{idèle}%
\TextAt{8}{292.808}{618.131}{of}%
\TextAt{18}{306.445}{618.131}{Q}%
\TextAt{9}{315.743}{618.131}{\Character{40}}%
\TextAt{13}{320.295}{618.131}{\Character{22}}%
\TextAt{21}{327.338}{616.338}{m}%
\TextAt{9}{335.327}{618.131}{\Character{41}}%
\TextAt{8}{339.879}{618.131}{.}%
\TextAt{8}{349.249}{618.131}{Thus}%
\TextAt{13}{379.205}{618.131}{C}%
\TextAt{27}{388.439}{622.469}{m}%
\TextAt{25}{387.58}{615.066}{Q}%
\NativeMathAt{398.892}{618.131}{11.955}{\subseteq}%
\TextAt{13}{412.035}{618.131}{N}%
\TextAt{13}{422.651}{618.131}{C}%
\TextAt{25}{431.032}{616.271}{Q}%
\TextAt{17}{437.231}{616.271}{\Character{40}}%
\TextAt{21}{440.525}{616.271}{\Character{22}}%
\TextAt{22}{445.619}{615.275}{m}%
\TextAt{17}{452.655}{616.271}{\Character{41}}%
\TextAt{8}{456.446}{618.131}{.}%
\TextAt{8}{465.816}{618.131}{On}%
\TextAt{8}{485.634}{618.131}{the}%
\TextAt{8}{506.098}{618.131}{other}%
\TextAt{8}{79.226}{601.085}{hand,}%
\TextAt{13}{111.744}{601.085}{C}%
\TextAt{25}{120.119}{599.292}{Q}%
\TextAt{9}{126.816}{601.085}{\Character{47}}%
\TextAt{13}{132.669}{601.085}{C}%
\TextAt{27}{141.903}{605.423}{m}%
\TextAt{25}{141.044}{598.02}{Q}%
\NativeMathAt{151.833}{601.085}{11.955}{\cong}%
\TextAt{9}{164.452}{601.085}{\Character{40}}%
\TextAt{18}{169.005}{601.085}{Z}%
\TextAt{9}{176.975}{601.085}{\Character{47}}%
\TextAt{13}{182.828}{601.085}{m}%
\TextAt{18}{193.067}{601.085}{Z}%
\TextAt{9}{201.037}{601.085}{\Character{41}}%
\TextAt{3}{205.59}{605.423}{\Character{3}}%
\TextAt{8}{214.219}{601.085}{by}%
\TextAt{8}{230.478}{601.085}{(1.10),}%
\TextAt{8}{267.563}{601.085}{and}%
\TextAt{8}{290.326}{601.085}{therefore}%
\TextAt{9}{194.863}{572.821}{\Character{40}}%
\TextAt{13}{199.416}{572.821}{C}%
\TextAt{25}{207.791}{571.028}{Q}%
\TextAt{9}{217.809}{572.821}{\Character{58}}%
\TextAt{13}{224.381}{572.821}{C}%
\TextAt{27}{233.615}{577.757}{m}%
\TextAt{25}{232.756}{569.866}{Q}%
\TextAt{9}{240.224}{572.821}{\Character{41}}%
\TextAt{9}{248.079}{572.821}{\Character{61}}%
\TextAt{9}{260.5}{572.821}{\Character{91}}%
\TextAt{18}{263.774}{572.821}{Q}%
\TextAt{9}{273.073}{572.821}{\Character{40}}%
\TextAt{13}{277.625}{572.821}{\Character{22}}%
\TextAt{21}{284.668}{571.028}{m}%
\TextAt{9}{292.657}{572.821}{\Character{41}}%
\TextAt{9}{300.512}{572.821}{\Character{58}}%
\TextAt{18}{307.102}{572.821}{Q}%
\TextAt{9}{316.401}{572.821}{\Character{93}}%
\TextAt{9}{322.952}{572.821}{\Character{61}}%
\TextAt{9}{335.374}{572.821}{\Character{40}}%
\TextAt{13}{339.951}{572.821}{C}%
\TextAt{25}{348.326}{571.028}{Q}%
\TextAt{9}{358.344}{572.821}{\Character{58}}%
\TextAt{13}{364.917}{572.821}{N}%
\TextAt{13}{375.533}{572.821}{C}%
\TextAt{25}{383.914}{570.962}{Q}%
\TextAt{17}{390.113}{570.962}{\Character{40}}%
\TextAt{21}{393.407}{570.962}{\Character{22}}%
\TextAt{22}{398.501}{569.965}{m}%
\TextAt{17}{405.537}{570.962}{\Character{41}}%
\TextAt{9}{409.328}{572.821}{\Character{41}}%
\TextAt{13}{413.881}{572.821}{\Character{59}}%
\TextAt{8}{79.226}{544.558}{so}%
\TextAt{8}{93.596}{544.558}{that}%
\TextAt{13}{118.965}{544.558}{C}%
\TextAt{27}{128.199}{548.896}{m}%
\TextAt{25}{127.34}{541.493}{Q}%
\TextAt{9}{138.129}{544.558}{\Character{61}}%
\TextAt{13}{150.554}{544.558}{N}%
\TextAt{13}{161.17}{544.558}{C}%
\TextAt{25}{169.552}{542.698}{Q}%
\TextAt{17}{175.751}{542.698}{\Character{40}}%
\TextAt{21}{179.044}{542.698}{\Character{22}}%
\TextAt{22}{184.139}{541.702}{m}%
\TextAt{17}{191.175}{542.698}{\Character{41}}%
\TextAt{8}{194.966}{544.558}{,}%
\TextAt{8}{202.115}{544.558}{and}%
\TextAt{8}{224.878}{544.558}{this}%
\TextAt{8}{247.7}{544.558}{proves}%
\TextAt{8}{283.865}{544.558}{the}%
\TextAt{8}{304.021}{544.558}{claim.}%
\Vector{q 0 G 0 g 0.399 w 452.334 -229.442 m 452.334 -221.373 l S Q}%
\Vector{q 0 G 0 g 0.399 w 452.534 -221.572 m 459.707 -221.572 l S Q}%
\Vector{q 0 G 0 g 0.399 w 452.534 -229.243 m 459.707 -229.243 l S Q}%
\Vector{q 0 G 0 g 0.399 w 459.906 -229.442 m 459.906 -221.373 l S Q}%
\TextAt{8}{97.159}{515.767}{According}%
\TextAt{8}{153.278}{515.767}{to}%
\TextAt{8}{168.427}{515.767}{this}%
\TextAt{8}{192.087}{515.767}{proposition,}%
\TextAt{8}{257.92}{515.767}{one}%
\TextAt{8}{280.218}{515.767}{may}%
\TextAt{8}{306.425}{515.767}{view}%
\TextAt{8}{334.247}{515.767}{the}%
\TextAt{8}{355.241}{515.767}{general}%
\TextAt{8}{396.404}{515.767}{ray}%
\TextAt{8}{417.411}{515.767}{class}%
\TextAt{8}{445.687}{515.767}{fields}%
\TextAt{13}{476.497}{515.767}{K}%
\TextAt{27}{487.314}{520.106}{m}%
\TextAt{16}{498.673}{515.767}{\Character{106}}%
\TextAt{13}{506.743}{515.767}{K}%
\TextAt{8}{522.297}{515.767}{as}%
\TextAt{8}{79.226}{498.721}{analogues}%
\TextAt{8}{133.55}{498.721}{of}%
\TextAt{8}{147.799}{498.721}{the}%
\TextAt{8}{168.876}{498.721}{cyclotomic}%
\TextAt{8}{228.005}{498.721}{fields}%
\TextAt{18}{258.897}{498.721}{Q}%
\TextAt{9}{268.195}{498.721}{\Character{40}}%
\TextAt{13}{272.747}{498.721}{\Character{22}}%
\TextAt{21}{279.79}{496.928}{m}%
\TextAt{9}{287.779}{498.721}{\Character{41}}%
\TextAt{16}{297.22}{498.721}{\Character{106}}%
\TextAt{18}{305.429}{498.721}{Q}%
\TextAt{8}{314.727}{498.721}{.}%
\TextAt{8}{325.935}{498.721}{Nonetheless,}%
\TextAt{8}{394.512}{498.721}{they}%
\TextAt{8}{421.769}{498.721}{are}%
\TextAt{8}{442.199}{498.721}{not}%
\TextAt{8}{463.933}{498.721}{made}%
\TextAt{8}{496.068}{498.721}{to}%
\TextAt{8}{511.298}{498.721}{take}%
\TextAt{8}{79.226}{481.675}{over}%
\TextAt{8}{104.499}{481.675}{the}%
\TextAt{8}{124.883}{481.675}{important}%
\TextAt{8}{180.403}{481.675}{role}%
\TextAt{8}{203.393}{481.675}{of}%
\TextAt{8}{216.95}{481.675}{the}%
\TextAt{8}{237.333}{481.675}{latter}%
\TextAt{8}{269.433}{481.675}{because}%
\TextAt{8}{312.962}{481.675}{all}%
\TextAt{8}{329.448}{481.675}{we}%
\TextAt{8}{346.902}{481.675}{know}%
\TextAt{8}{377.699}{481.675}{about}%
\TextAt{8}{411.424}{481.675}{them}%
\TextAt{8}{441.563}{481.675}{is}%
\TextAt{8}{453.554}{481.675}{that}%
\TextAt{8}{479.151}{481.675}{they}%
\TextAt{8}{505.715}{481.675}{exist,}%
\TextAt{8}{79.226}{464.629}{but}%
\TextAt{8}{100.538}{464.629}{not}%
\TextAt{8}{121.204}{464.629}{how}%
\TextAt{8}{145.445}{464.629}{to}%
\TextAt{8}{159.607}{464.629}{generate}%
\TextAt{8}{206.288}{464.629}{them.}%
\TextAt{8}{240.701}{464.629}{In}%
\TextAt{8}{255.175}{464.629}{the}%
\TextAt{8}{275.183}{464.629}{case}%
\TextAt{8}{299.807}{464.629}{of}%
\TextAt{8}{312.989}{464.629}{local}%
\TextAt{8}{340.481}{464.629}{fields}%
\TextAt{8}{370.305}{464.629}{things}%
\TextAt{8}{405.341}{464.629}{were}%
\TextAt{8}{432.176}{464.629}{different.}%
\TextAt{8}{482.849}{464.629}{There}%
\TextAt{8}{516.51}{464.629}{the}%
\TextAt{8}{79.226}{447.583}{analogues}%
\TextAt{8}{132.059}{447.583}{of}%
\TextAt{8}{144.818}{447.583}{the}%
\TextAt{8}{164.403}{447.583}{ray}%
\TextAt{8}{184}{447.583}{class}%
\TextAt{8}{210.866}{447.583}{fields}%
\TextAt{8}{240.267}{447.583}{were}%
\TextAt{8}{266.679}{447.583}{the}%
\TextAt{8}{286.264}{447.583}{Lubin–Tate}%
\TextAt{8}{348.613}{447.583}{extensions}%
\TextAt{8}{404.422}{447.583}{which}%
\TextAt{8}{437.338}{447.583}{could}%
\TextAt{8}{467.982}{447.583}{be}%
\TextAt{8}{483.335}{447.583}{generated}%
\TextAt{8}{79.226}{430.537}{by}%
\TextAt{8}{95.003}{430.537}{the}%
\TextAt{8}{114.677}{430.537}{division}%
\TextAt{8}{157.509}{430.537}{points}%
\TextAt{8}{192.21}{430.537}{of}%
\TextAt{8}{205.058}{430.537}{formal}%
\TextAt{8}{241.326}{430.537}{groups}%
\TextAt{8}{278.634}{430.537}{—}%
\TextAt{8}{293.754}{430.537}{a}%
\TextAt{8}{303.027}{430.537}{fact}%
\TextAt{8}{325.63}{430.537}{that}%
\TextAt{8}{350.517}{430.537}{carries}%
\TextAt{8}{387.167}{430.537}{a}%
\TextAt{8}{396.441}{430.537}{long}%
\TextAt{8}{421.327}{430.537}{way}%
\TextAt{8}{444.588}{430.537}{(see}%
\TextAt{8}{467.574}{430.537}{chap.}%
\TextAt{8}{497.984}{430.537}{V,}%
\TextAt{8}{513.426}{430.537}{§5).}%
\TextAt{8}{79.226}{413.491}{This}%
\TextAt{8}{105.938}{413.491}{local}%
\TextAt{8}{133.57}{413.491}{discovery}%
\TextAt{8}{184.359}{413.491}{does,}%
\TextAt{8}{214.001}{413.491}{however,}%
\TextAt{8}{262.126}{413.491}{originate}%
\TextAt{8}{310.907}{413.491}{from}%
\TextAt{8}{338.539}{413.491}{the}%
\TextAt{8}{358.687}{413.491}{problem}%
\TextAt{8}{404.204}{413.491}{of}%
\TextAt{8}{417.526}{413.491}{generating}%
\TextAt{8}{474.759}{413.491}{global}%
\TextAt{8}{509.23}{413.491}{class}%
\TextAt{8}{79.226}{396.445}{fields,}%
\TextAt{8}{112.45}{396.445}{which}%
\TextAt{8}{145.936}{396.445}{will}%
\TextAt{8}{168.041}{396.445}{be}%
\TextAt{8}{183.965}{396.445}{discussed}%
\TextAt{8}{234.871}{396.445}{at}%
\TextAt{8}{249.181}{396.445}{the}%
\TextAt{8}{269.337}{396.445}{end}%
\TextAt{8}{291.442}{396.445}{of}%
\TextAt{8}{304.772}{396.445}{this}%
\TextAt{8}{327.595}{396.445}{section.}%
\TextAt{8}{97.159}{376.411}{Note}%
\TextAt{8}{127.115}{376.411}{in}%
\TextAt{8}{142.438}{376.411}{passing}%
\TextAt{8}{185.21}{376.411}{that}%
\TextAt{8}{212.249}{376.411}{the}%
\TextAt{8}{234.075}{376.411}{above}%
\TextAt{8}{268.921}{376.411}{proposition}%
\TextAt{8}{332.124}{376.411}{gives}%
\TextAt{8}{362.474}{376.411}{another}%
\TextAt{8}{407.076}{376.411}{proof}%
\TextAt{8}{439.315}{376.411}{of}%
\TextAt{8}{454.315}{376.411}{the}%
\TextAt{8}{476.142}{376.411}{theorem}%
\TextAt{8}{523.337}{376.411}{of}%
\TextAt{8}{79.226}{359.365}{Kronecker}%
\TextAt{8}{134.097}{359.365}{and}%
\TextAt{8}{156.127}{359.365}{Weber}%
\TextAt{8}{192.121}{359.365}{(see}%
\TextAt{8}{214.856}{359.365}{chap.}%
\TextAt{8}{245.015}{359.365}{V,}%
\TextAt{8}{260.207}{359.365}{(1.10))}%
\TextAt{8}{297.862}{359.365}{to}%
\TextAt{8}{311.44}{359.365}{the}%
\TextAt{8}{330.864}{359.365}{effect}%
\TextAt{8}{361.011}{359.365}{that}%
\TextAt{8}{385.647}{359.365}{every}%
\TextAt{8}{415.807}{359.365}{finite}%
\TextAt{8}{444.986}{359.365}{abelian}%
\TextAt{8}{484.901}{359.365}{extension}%
\ErrataAt{79.22}{342.355}{453.64}{$L\mid\mathbb Q$ is contained in a field $\mathbb Q(\mu_m)\mid\mathbb Q$, because by \hyperlink{nav.number.6.1.8}{(1.8)} the norm group $N_{L\mid\mathbb Q}C_L$}%
\ErrataAt{79.22}{325.255}{453.64}{contains a congruence subgroup $C_{\mathbb Q}^{\mathfrak m}$, $\mathfrak m=(m)$, so that $L\subseteq\mathbb Q(\mu_m)$.}%
\TextAt{8}{97.159}{305.238}{Among}%
\TextAt{8}{137.637}{305.238}{all}%
\TextAt{8}{153.727}{305.238}{abelian}%
\TextAt{8}{194.206}{305.238}{extensions}%
\TextAt{8}{250.417}{305.238}{of}%
\TextAt{13}{263.578}{305.238}{K}%
\TextAt{8}{274.395}{305.238}{,}%
\TextAt{8}{281.409}{305.238}{the}%
\TextAt{8}{301.397}{305.238}{ray}%
\TextAt{8}{321.396}{305.238}{class}%
\TextAt{8}{348.664}{305.238}{field}%
\TextAt{8}{373.851}{305.238}{mod}%
\TextAt{9}{400.02}{305.238}{1}%
\TextAt{8}{409.601}{305.238}{occupies}%
\TextAt{8}{455.985}{305.238}{a}%
\TextAt{8}{465.571}{305.238}{special}%
\TextAt{8}{503.503}{305.238}{place.}%
\TextAt{8}{79.226}{288.192}{It}%
\TextAt{8}{91.899}{288.192}{is}%
\TextAt{8}{103.663}{288.192}{called}%
\TextAt{8}{136.826}{288.192}{the}%
\TextAt{12}{156.983}{288.192}{big}%
\TextAt{12}{179.41}{288.192}{Hilbert}%
\TextAt{12}{226.609}{288.192}{class}%
\TextAt{12}{257.968}{288.192}{field}%
\TextAt{8}{286.684}{288.192}{and}%
\TextAt{8}{309.447}{288.192}{has}%
\TextAt{8}{330.321}{288.192}{Galois}%
\TextAt{8}{366.234}{288.192}{group}%
\TextAt{13}{258.606}{259.928}{G}%
\TextAt{9}{267.84}{259.928}{\Character{40}}%
\TextAt{13}{272.392}{259.928}{K}%
\TextAt{17}{283.21}{264.864}{1}%
\TextAt{16}{291.263}{259.928}{\Character{106}}%
\TextAt{13}{297.905}{259.928}{K}%
\TextAt{9}{308.722}{259.928}{\Character{41}}%
\NativeMathAt{316.595}{259.928}{11.955}{\cong}%
\TextAt{13}{329.215}{259.928}{C}%
\TextAt{13}{338.433}{259.928}{l}%
\TextAt{17}{342.198}{264.864}{1}%
\TextAt{21}{341.935}{256.973}{K}%
\TextAt{13}{350.138}{259.928}{\Character{58}}%
\TextAt{8}{79.226}{231.665}{By}%
\TextAt{8}{97.057}{231.665}{(1.11),}%
\TextAt{8}{133.717}{231.665}{the}%
\TextAt{8}{153.341}{231.665}{group}%
\TextAt{13}{185.985}{231.665}{C}%
\TextAt{13}{195.203}{231.665}{l}%
\TextAt{17}{198.968}{236.003}{1}%
\TextAt{21}{198.705}{228.644}{K}%
\TextAt{8}{210.273}{231.665}{is}%
\TextAt{8}{221.505}{231.665}{linked}%
\TextAt{8}{255.44}{231.665}{to}%
\TextAt{8}{269.218}{231.665}{the}%
\TextAt{8}{288.842}{231.665}{ordinary}%
\TextAt{8}{335.473}{231.665}{ideal}%
\TextAt{8}{362.905}{231.665}{class}%
\TextAt{8}{389.81}{231.665}{group}%
\TextAt{8}{422.453}{231.665}{by}%
\TextAt{8}{438.18}{231.665}{the}%
\TextAt{8}{457.805}{231.665}{exact}%
\TextAt{8}{488.165}{231.665}{sequence}%
\TextAt{9}{165.627}{203.401}{1}%
\NativeArrowAt{174.801}{203.401}{11.955}{19.26}{\longrightarrow}%
\TextAt{16}{197.372}{203.401}{\Character{79}}%
\TextAt{3}{207.233}{208.337}{\Character{3}}%
\TextAt{9}{211.965}{203.401}{\Character{47}}%
\TextAt{16}{217.818}{203.401}{\Character{79}}%
\TextAt{3}{227.668}{208.337}{\Character{3}}%
\TextAt{17}{227.336}{200.445}{\Character{43}}%
\NativeArrowAt{237.741}{203.401}{11.955}{19.26}{\longrightarrow}%
\TextAt{20}{262.913}{214.758}{\Character{89}}%
\TextAt{27}{260.323}{189.302}{p}%
\TextAt{7}{267.133}{189.302}{real}%
\TextAt{18}{282.771}{203.401}{R}%
\TextAt{3}{291.405}{208.337}{\Character{3}}%
\TextAt{9}{296.137}{203.401}{\Character{47}}%
\TextAt{18}{301.99}{203.401}{R}%
\TextAt{3}{310.624}{208.337}{\Character{3}}%
\TextAt{17}{310.624}{200.445}{\Character{43}}%
\NativeArrowAt{321.03}{203.401}{11.955}{19.26}{\longrightarrow}%
\TextAt{13}{343.612}{203.401}{C}%
\TextAt{13}{352.829}{203.401}{l}%
\TextAt{17}{356.595}{208.337}{1}%
\TextAt{21}{356.332}{200.445}{K}%
\NativeArrowAt{367.856}{203.401}{11.955}{19.26}{\longrightarrow}%
\TextAt{13}{390.438}{203.401}{C}%
\TextAt{13}{399.656}{203.401}{l}%
\TextAt{21}{403.158}{201.608}{K}%
\NativeArrowAt{414.682}{203.401}{11.955}{19.26}{\longrightarrow}%
\TextAt{9}{437.264}{203.401}{1}%
\TextAt{13}{443.117}{203.401}{\Character{58}}%
\TextAt{8}{79.226}{165.969}{The}%
\TextAt{8}{104.208}{165.969}{big}%
\TextAt{8}{124.647}{165.969}{Hilbert}%
\TextAt{8}{165.888}{165.969}{class}%
\TextAt{8}{194.254}{165.969}{field}%
\TextAt{8}{220.539}{165.969}{has}%
\TextAt{8}{242.341}{165.969}{conductor}%
\TextAt{14}{297.904}{165.969}{f}%
\TextAt{9}{306.708}{165.969}{\Character{61}}%
\TextAt{9}{320.708}{165.969}{1}%
\TextAt{8}{331.394}{165.969}{and}%
\TextAt{8}{355.085}{165.969}{may}%
\TextAt{8}{381.382}{165.969}{therefore}%
\TextAt{8}{431.41}{165.969}{be}%
\TextAt{8}{448.263}{165.969}{characterized}%
\TextAt{8}{520.408}{165.969}{by}%
\TextAt{8}{79.226}{148.923}{(6.6)}%
\TextAt{8}{107.201}{148.923}{in}%
\TextAt{8}{120.854}{148.923}{the}%
\TextAt{8}{141.01}{148.923}{following}%
\TextAt{8}{190.445}{148.923}{way.}%
\TextAt{12}{79.226}{123.121}{Proposition}%
\TextAt{12}{151.786}{123.121}{(6.8).}%
\TextAt{8}{186.97}{123.121}{The}%
\TextAt{8}{210.892}{123.121}{big}%
\TextAt{8}{230.27}{123.121}{Hilbert}%
\TextAt{8}{270.451}{123.121}{class}%
\TextAt{8}{297.757}{123.121}{field}%
\TextAt{8}{322.982}{123.121}{is}%
\TextAt{8}{334.614}{123.121}{the}%
\TextAt{8}{354.638}{123.121}{maximal}%
\TextAt{8}{402.315}{123.121}{unramified}%
\TextAt{8}{460.716}{123.121}{abelian}%
\TextAt{8}{501.232}{123.121}{exten-}%
\TextAt{8}{79.226}{106.075}{sion}%
\TextAt{8}{103.352}{106.075}{of}%
\TextAt{13}{116.682}{106.075}{K}%
\TextAt{8}{127.499}{106.075}{.}%
\TextAt{8}{297.211}{43.033}{363}%
\Destination{78.23}{809.815}{nav.page.363}%
\Destination{79.36}{340.255}{nav.print.399}%
\Destination{79.36}{775.155}{nav.number.6.6.7}%
\Destination{79.36}{137.155}{nav.number.6.6.8}%
\Hotlink{362.131}{632.955}{9.775}{11.047}{nav.chapter.5}%
\Hotlink{497.877}{632.955}{6.858}{11.047}{nav.section.5.3}%
\Hotlink{497.484}{428.355}{9.775}{11.047}{nav.chapter.5}%
\Hotlink{518.605}{428.355}{6.858}{11.047}{nav.section.5.5}%
\Hotlink{244.515}{357.155}{9.775}{11.047}{nav.chapter.5}%
\Hotlink{377.948}{632.955}{25.078}{11.047}{nav.number.5.1.8}%
\Hotlink{229.978}{598.855}{30.936}{11.047}{nav.number.6.1.10}%
\Hotlink{259.707}{357.155}{30.936}{11.047}{nav.number.5.1.10}%
\Hotlink{96.557}{229.455}{30.936}{11.047}{nav.number.6.1.11}%
\Hotlink{78.726}{146.755}{25.078}{11.047}{nav.number.6.6.6}%
%
\Destination{79.226}{356.355}{errata.31}%
}

% Page 376 (body page 364)
\Page{%
\TextAt{2}{79.226}{796.943}{VI.}%
\TextAt{2}{98.012}{796.943}{Global}%
\TextAt{2}{133.248}{796.943}{Class}%
\TextAt{2}{161.841}{796.943}{Field}%
\TextAt{2}{189.572}{796.943}{Theory}%
\Vector{q 0 G 0 g 0.299 w 7.226 18.729 m 460.769 18.729 l S Q}%
\TextAt{8}{97.159}{761.124}{Since}%
\TextAt{8}{128.494}{761.124}{the}%
\TextAt{8}{149.428}{761.124}{infinite}%
\TextAt{8}{189.873}{761.124}{places}%
\TextAt{8}{225.178}{761.124}{are}%
\TextAt{8}{245.466}{761.124}{always}%
\TextAt{8}{283.711}{761.124}{unramified,}%
\TextAt{8}{346.468}{761.124}{this}%
\TextAt{8}{370.068}{761.124}{means}%
\TextAt{8}{406.676}{761.124}{that}%
\TextAt{8}{432.822}{761.124}{all}%
\TextAt{8}{449.859}{761.124}{prime}%
\TextAt{8}{483.8}{761.124}{ideals}%
\TextAt{8}{517.156}{761.124}{are}%
\TextAt{8}{79.226}{744.078}{unramified.}%
\TextAt{8}{142.202}{744.078}{The}%
\TextAt{12}{165.933}{744.078}{Hilbert}%
\TextAt{12}{212.761}{744.078}{class}%
\TextAt{12}{243.749}{744.078}{field}%
\TextAt{8}{268.568}{744.078}{,}%
\TextAt{8}{275.459}{744.078}{or}%
\TextAt{8}{289.447}{744.078}{more}%
\TextAt{8}{318.39}{744.078}{precisely,}%
\TextAt{8}{368.26}{744.078}{the}%
\TextAt{8}{388.094}{744.078}{“small}%
\TextAt{8}{423.924}{744.078}{Hilbert}%
\TextAt{8}{463.914}{744.078}{class}%
\TextAt{8}{491.029}{744.078}{field”,}%
\TextAt{8}{524.903}{744.078}{is}%
\TextAt{8}{79.226}{727.032}{defined}%
\TextAt{8}{119.161}{727.032}{to}%
\TextAt{8}{133.094}{727.032}{be}%
\TextAt{8}{148.64}{727.032}{the}%
\TextAt{8}{168.419}{727.032}{maximal}%
\TextAt{8}{215.85}{727.032}{unramified}%
\TextAt{8}{274.005}{727.032}{abelian}%
\TextAt{8}{314.275}{727.032}{extension}%
\TextAt{13}{365.663}{727.032}{H}%
\TextAt{16}{379.607}{727.032}{\Character{106}}%
\TextAt{13}{386.248}{727.032}{K}%
\TextAt{8}{400.586}{727.032}{in}%
\TextAt{8}{413.861}{727.032}{which}%
\TextAt{8}{446.969}{727.032}{all}%
\TextAt{8}{462.851}{727.032}{infinite}%
\TextAt{8}{502.14}{727.032}{places}%
\TextAt{8}{79.226}{709.986}{split}%
\TextAt{8}{105.3}{709.986}{completely,}%
\TextAt{8}{166.427}{709.986}{i.e.,}%
\TextAt{8}{188.532}{709.986}{the}%
\TextAt{8}{208.689}{709.986}{real}%
\TextAt{8}{231.451}{709.986}{places}%
\TextAt{8}{265.978}{709.986}{stay}%
\TextAt{8}{290.761}{709.986}{real.}%
\TextAt{8}{318.074}{709.986}{It}%
\TextAt{8}{330.747}{709.986}{satisfies}%
\TextAt{8}{373.857}{709.986}{the}%
\TextAt{12}{79.226}{683.973}{Proposition}%
\TextAt{12}{151.165}{683.973}{(6.9).}%
\TextAt{8}{185.809}{683.973}{The}%
\TextAt{8}{209.191}{683.973}{Galois}%
\TextAt{8}{244.433}{683.973}{group}%
\TextAt{8}{276.937}{683.973}{of}%
\TextAt{8}{289.595}{683.973}{the}%
\TextAt{8}{309.08}{683.973}{small}%
\TextAt{8}{339.038}{683.973}{Hilbert}%
\TextAt{8}{378.679}{683.973}{class}%
\TextAt{8}{405.445}{683.973}{field}%
\TextAt{13}{430.13}{683.973}{H}%
\TextAt{16}{444.073}{683.973}{\Character{106}}%
\TextAt{13}{450.715}{683.973}{K}%
\TextAt{8}{464.759}{683.973}{is}%
\TextAt{8}{475.851}{683.973}{canonically}%
\TextAt{8}{79.226}{666.927}{isomorphic}%
\TextAt{8}{138.476}{666.927}{to}%
\TextAt{8}{152.786}{666.927}{the}%
\TextAt{8}{172.943}{666.927}{ideal}%
\TextAt{8}{200.906}{666.927}{class}%
\TextAt{8}{228.343}{666.927}{group:}%
\TextAt{13}{261.07}{637.926}{G}%
\TextAt{9}{270.303}{637.926}{\Character{40}}%
\TextAt{13}{274.856}{637.926}{H}%
\TextAt{16}{288.799}{637.926}{\Character{106}}%
\TextAt{13}{295.441}{637.926}{K}%
\TextAt{9}{306.259}{637.926}{\Character{41}}%
\NativeMathAt{314.132}{637.926}{11.955}{\cong}%
\TextAt{13}{326.751}{637.926}{C}%
\TextAt{13}{335.969}{637.926}{l}%
\TextAt{21}{339.472}{636.133}{K}%
\TextAt{13}{347.674}{637.926}{\Character{58}}%
\TextAt{8}{79.226}{608.925}{In}%
\TextAt{8}{93.847}{608.925}{particular,}%
\TextAt{8}{151.089}{608.925}{the}%
\TextAt{8}{171.245}{608.925}{degree}%
\TextAt{9}{207.661}{608.925}{\Character{91}}%
\TextAt{13}{210.912}{608.925}{H}%
\TextAt{9}{224.856}{608.925}{\Character{58}}%
\TextAt{13}{231.428}{608.925}{K}%
\TextAt{9}{242.246}{608.925}{\Character{93}}%
\TextAt{8}{249.395}{608.925}{is}%
\TextAt{8}{261.159}{608.925}{the}%
\TextAt{8}{281.315}{608.925}{class}%
\TextAt{8}{308.752}{608.925}{number}%
\TextAt{13}{351.349}{608.925}{h}%
\TextAt{21}{358.087}{607.132}{K}%
\TextAt{8}{370.187}{608.925}{of}%
\TextAt{13}{383.517}{608.925}{K}%
\TextAt{8}{394.335}{608.925}{.}%
\TextAt{19}{79.226}{573.946}{Proof.}%
\TextAt{8}{116}{573.946}{We}%
\TextAt{8}{136.383}{573.946}{consider}%
\TextAt{8}{182.207}{573.946}{the}%
\TextAt{8}{202.603}{573.946}{big}%
\TextAt{8}{222.352}{573.946}{Hilbert}%
\TextAt{8}{262.904}{573.946}{class}%
\TextAt{8}{290.58}{573.946}{field}%
\TextAt{13}{316.175}{573.946}{K}%
\TextAt{17}{326.993}{578.285}{1}%
\TextAt{16}{335.453}{573.946}{\Character{106}}%
\TextAt{13}{342.502}{573.946}{K}%
\TextAt{8}{357.456}{573.946}{and,}%
\TextAt{8}{383.768}{573.946}{for}%
\TextAt{8}{401.892}{573.946}{every}%
\TextAt{8}{433.023}{573.946}{infinite}%
\TextAt{8}{472.929}{573.946}{place}%
\TextAt{14}{503.08}{573.946}{p}%
\TextAt{8}{509.063}{573.946}{,}%
\TextAt{8}{516.51}{573.946}{the}%
\TextAt{8}{79.226}{556.9}{commutative}%
\TextAt{8}{149.152}{556.9}{diagram}%
\TextAt{8}{194.689}{556.9}{(see}%
\TextAt{8}{218.157}{556.9}{(5.6))}%
\TextAt{13}{242.402}{522.779}{K}%
\TextAt{3}{253.219}{527.118}{\Character{3}}%
\TextAt{27}{252.361}{519.824}{p}%
\TextAt{13}{327.587}{522.779}{G}%
\TextAt{9}{336.821}{522.779}{\Character{40}}%
\TextAt{13}{341.373}{522.779}{K}%
\TextAt{17}{352.191}{527.118}{1}%
\TextAt{27}{351.332}{519.824}{p}%
\TextAt{16}{360.244}{522.779}{\Character{106}}%
\TextAt{13}{366.885}{522.779}{K}%
\TextAt{27}{376.845}{520.986}{p}%
\TextAt{9}{381.331}{522.779}{\Character{41}}%
\TextAt{13}{226.11}{478.563}{I}%
\TextAt{21}{231.272}{476.77}{K}%
\TextAt{9}{239.475}{478.563}{\Character{47}}%
\TextAt{13}{245.328}{478.563}{I}%
\TextAt{17}{251.431}{482.902}{1}%
\TextAt{21}{250.491}{475.543}{K}%
\TextAt{13}{258.694}{478.563}{K}%
\TextAt{3}{269.511}{482.902}{\Character{3}}%
\TextAt{13}{327.776}{478.563}{G}%
\TextAt{9}{337.01}{478.563}{\Character{40}}%
\TextAt{13}{341.562}{478.563}{K}%
\TextAt{17}{352.38}{482.902}{1}%
\TextAt{16}{360.433}{478.563}{\Character{106}}%
\TextAt{13}{367.074}{478.563}{K}%
\TextAt{9}{377.892}{478.563}{\Character{41}}%
\TextAt{13}{382.444}{478.563}{\Character{58}}%
\NativeDiagram{%
\path[draw=black,line width=0.478bp] (191.304,-248.232) -- (249.759,-248.232);%
\path[draw=black,line width=0.478bp,line cap=round,line join=round] (247.511,-245.363) .. controls (247.966,-247.084) and (248.978,-247.897) .. (249.998,-248.232) .. controls (248.978,-248.567) and (247.966,-249.38) .. (247.511,-251.101);%
}%
\TextAt{17}{273.636}{531.707}{\Character{40}}%
\TextAt{21}{280.826}{531.707}{\Character{59}K}%
\TextAt{24}{290.883}{534.52}{1}%
\TextAt{28}{290.32}{529.715}{p}%
\TextAt{3}{295.034}{531.707}{\Character{106}}%
\TextAt{21}{297.387}{531.707}{K}%
\TextAt{28}{304.528}{530.711}{p}%
\TextAt{17}{308.611}{531.707}{\Character{41}}%
\NativeDiagram{%
\path[draw=black,line width=0.478bp] (178.177,-260.262) -- (178.177,-280.905);%
\path[draw=black,line width=0.478bp,line cap=round,line join=round] (181.046,-278.658) .. controls (179.325,-279.112) and (178.512,-280.124) .. (178.177,-281.144) .. controls (177.842,-280.124) and (177.029,-279.112) .. (175.308,-278.658);%
}%
\TextAt{3}{236.877}{501.185}{\Character{104}}%
\TextAt{3}{244.058}{501.185}{\Character{105}}%
\NativeDiagram{%
\path[draw=black,line width=0.478bp] (284.738,-260.262) -- (284.738,-280.905);%
\path[draw=black,line width=0.478bp,line cap=round,line join=round] (287.607,-278.658) .. controls (285.885,-279.112) and (285.072,-280.124) .. (284.738,-281.144) .. controls (284.403,-280.124) and (283.59,-279.112) .. (281.869,-278.658);%
\path[draw=black,line width=0.478bp] (207.596,-292.449) -- (249.947,-292.449);%
\path[draw=black,line width=0.478bp,line cap=round,line join=round] (247.7,-289.579) .. controls (248.154,-291.301) and (249.166,-292.114) .. (250.186,-292.449) .. controls (249.166,-292.783) and (248.154,-293.596) .. (247.7,-295.318);%
}%
\TextAt{17}{283.636}{472.071}{\Character{40}}%
\TextAt{21}{290.826}{472.071}{\Character{59}K}%
\TextAt{24}{300.883}{474.884}{1}%
\TextAt{3}{305.034}{472.071}{\Character{106}}%
\TextAt{21}{307.387}{472.071}{K}%
\TextAt{17}{315.091}{472.071}{\Character{41}}%
\TextAt{8}{79.226}{445.585}{The}%
\TextAt{8}{103.065}{445.585}{small}%
\TextAt{8}{133.479}{445.585}{Hilbert}%
\TextAt{8}{173.576}{445.585}{class}%
\TextAt{8}{200.798}{445.585}{field}%
\TextAt{13}{225.94}{445.585}{H}%
\TextAt{16}{239.883}{445.585}{\Character{106}}%
\TextAt{13}{246.525}{445.585}{K}%
\TextAt{8}{261.025}{445.585}{is}%
\TextAt{8}{272.573}{445.585}{the}%
\TextAt{8}{292.514}{445.585}{fixed}%
\TextAt{8}{320.585}{445.585}{field}%
\TextAt{8}{345.727}{445.585}{of}%
\TextAt{8}{358.842}{445.585}{the}%
\TextAt{8}{378.783}{445.585}{subgroup}%
\TextAt{13}{429.365}{445.585}{G}%
\NativeMathAt{438.599}{443.792}{7.97}{\infty}%
\TextAt{8}{451.247}{445.585}{generated}%
\TextAt{8}{504.364}{445.585}{by}%
\TextAt{8}{520.408}{445.585}{all}%
\TextAt{13}{79.226}{428.539}{G}%
\TextAt{9}{88.46}{428.539}{\Character{40}}%
\TextAt{13}{93.012}{428.539}{K}%
\TextAt{17}{103.83}{432.877}{1}%
\TextAt{27}{102.971}{425.583}{p}%
\TextAt{16}{111.883}{428.539}{\Character{106}}%
\TextAt{13}{118.525}{428.539}{K}%
\TextAt{27}{128.484}{426.746}{p}%
\TextAt{9}{132.971}{428.539}{\Character{41}}%
\TextAt{8}{137.523}{428.539}{,}%
\TextAt{14}{144.672}{428.539}{p}%
\TextAt{16}{153.976}{428.539}{\Character{106}}%
\NativeMathAt{160.599}{428.539}{11.955}{\infty}%
\TextAt{8}{172.573}{428.539}{.}%
\TextAt{8}{181.021}{428.539}{Under}%
\TextAt{9}{216.456}{428.539}{\Character{40}}%
\TextAt{13}{224.906}{428.539}{\Character{59}}%
\TextAt{13}{230.13}{428.539}{K}%
\TextAt{17}{240.968}{432.877}{1}%
\TextAt{16}{249.021}{428.539}{\Character{106}}%
\TextAt{13}{255.663}{428.539}{K}%
\TextAt{9}{266.481}{428.539}{\Character{41}}%
\TextAt{8}{274.93}{428.539}{this}%
\TextAt{8}{297.753}{428.539}{is}%
\TextAt{8}{309.517}{428.539}{the}%
\TextAt{8}{329.673}{428.539}{image}%
\TextAt{8}{363.494}{428.539}{of}%
\TextAt{20}{201.854}{410.647}{\Character{48}}%
\TextAt{20}{201.854}{389.128}{\Character{64}}%
\TextAt{20}{212.314}{397.974}{\Character{89}}%
\TextAt{27}{212.548}{372.071}{p}%
\TextAt{3}{216.536}{372.071}{\Character{106}}%
\NativeMathAt{218.887}{372.071}{7.97}{\infty}%
\TextAt{13}{229.583}{386.617}{K}%
\TextAt{3}{240.401}{391.553}{\Character{3}}%
\TextAt{27}{239.542}{383.661}{p}%
\TextAt{20}{245.133}{410.647}{\Character{49}}%
\TextAt{20}{245.133}{389.128}{\Character{65}}%
\TextAt{13}{257.586}{386.617}{I}%
\TextAt{17}{263.689}{391.553}{1}%
\TextAt{21}{262.749}{383.661}{K}%
\TextAt{13}{270.952}{386.617}{K}%
\TextAt{3}{281.77}{391.553}{\Character{3}}%
\TextAt{9}{286.502}{386.617}{\Character{47}}%
\TextAt{13}{292.355}{386.617}{I}%
\TextAt{17}{298.458}{391.553}{1}%
\TextAt{21}{297.518}{383.661}{K}%
\TextAt{13}{305.72}{386.617}{K}%
\TextAt{3}{316.538}{391.553}{\Character{3}}%
\TextAt{9}{324.591}{386.617}{\Character{61}}%
\TextAt{13}{337.017}{386.617}{I}%
\TextAt{21}{343.12}{391.731}{S}%
\NativeMathAt{348.27}{390.735}{5.978}{\infty}%
\TextAt{21}{342.18}{383.376}{K}%
\TextAt{13}{356.572}{386.617}{K}%
\TextAt{3}{367.39}{391.553}{\Character{3}}%
\TextAt{9}{372.122}{386.617}{\Character{47}}%
\TextAt{13}{377.975}{386.617}{I}%
\TextAt{17}{384.078}{391.553}{1}%
\TextAt{21}{383.138}{383.661}{K}%
\TextAt{13}{391.34}{386.617}{K}%
\TextAt{3}{402.158}{391.553}{\Character{3}}%
\TextAt{13}{406.89}{386.617}{\Character{59}}%
\TextAt{8}{79.226}{344.574}{where}%
\TextAt{13}{113.035}{344.574}{I}%
\TextAt{21}{119.138}{349.689}{S}%
\NativeMathAt{124.288}{348.692}{5.978}{\infty}%
\TextAt{21}{118.198}{341.333}{K}%
\TextAt{9}{135.911}{344.574}{\Character{61}}%
\TextAt{20}{148.337}{353.541}{\Character{81}}%
\TextAt{27}{159.628}{341.087}{p}%
\TextAt{3}{163.617}{341.087}{\Character{106}}%
\NativeMathAt{165.968}{341.087}{7.97}{\infty}%
\TextAt{13}{176.928}{344.574}{K}%
\TextAt{3}{187.746}{348.913}{\Character{3}}%
\TextAt{27}{186.887}{341.619}{p}%
\TextAt{16}{195.135}{344.574}{\Character{2}}%
\TextAt{20}{207.09}{353.541}{\Character{81}}%
\TextAt{27}{218.381}{341.087}{p}%
\TextAt{25}{222.369}{341.087}{\Character{45}}%
\NativeMathAt{224.583}{341.087}{7.97}{\infty}%
\TextAt{13}{235.542}{344.574}{U}%
\TextAt{27}{243.493}{342.781}{p}%
\TextAt{8}{247.979}{344.574}{.}%
\TextAt{8}{256.428}{344.574}{Therefore}%
\TextAt{8}{309.425}{344.574}{by}%
\TextAt{8}{325.684}{344.574}{(1.3),}%
\TextAt{13}{178.704}{315.573}{G}%
\TextAt{9}{187.938}{315.573}{\Character{40}}%
\TextAt{13}{192.49}{315.573}{H}%
\TextAt{16}{206.433}{315.573}{\Character{106}}%
\TextAt{13}{213.075}{315.573}{K}%
\TextAt{9}{223.893}{315.573}{\Character{41}}%
\TextAt{9}{231.748}{315.573}{\Character{61}}%
\TextAt{13}{244.191}{315.573}{G}%
\TextAt{9}{253.425}{315.573}{\Character{40}}%
\TextAt{13}{257.977}{315.573}{K}%
\TextAt{17}{268.795}{320.509}{1}%
\TextAt{16}{276.848}{315.573}{\Character{106}}%
\TextAt{13}{283.49}{315.573}{K}%
\TextAt{9}{294.307}{315.573}{\Character{41}\Character{47}}%
\TextAt{13}{304.713}{315.573}{G}%
\NativeMathAt{313.946}{313.78}{7.97}{\infty}%
\NativeMathAt{326.234}{315.573}{11.955}{\cong}%
\TextAt{13}{338.853}{315.573}{I}%
\TextAt{21}{344.016}{313.78}{K}%
\TextAt{9}{352.219}{315.573}{\Character{47}}%
\TextAt{13}{358.072}{315.573}{I}%
\TextAt{21}{364.174}{320.687}{S}%
\NativeMathAt{369.325}{319.691}{5.978}{\infty}%
\TextAt{21}{363.235}{312.332}{K}%
\TextAt{13}{377.627}{315.573}{K}%
\TextAt{3}{388.444}{320.509}{\Character{3}}%
\NativeMathAt{396.498}{315.573}{11.955}{\cong}%
\TextAt{13}{409.117}{315.573}{C}%
\TextAt{13}{418.334}{315.573}{l}%
\TextAt{21}{421.837}{313.78}{K}%
\TextAt{13}{430.04}{315.573}{\Character{58}}%
\Vector{q 0 G 0 g 0.399 w 452.334 -487.428 m 452.334 -479.358 l S Q}%
\Vector{q 0 G 0 g 0.399 w 452.534 -479.557 m 459.707 -479.557 l S Q}%
\Vector{q 0 G 0 g 0.399 w 452.534 -487.229 m 459.707 -487.229 l S Q}%
\Vector{q 0 G 0 g 0.399 w 459.906 -487.428 m 459.906 -479.358 l S Q}%
\TextAt{12}{79.226}{254.582}{Remark.}%
\TextAt{8}{133.766}{254.582}{The}%
\TextAt{8}{158.202}{254.582}{small}%
\TextAt{8}{189.214}{254.582}{Hilbert}%
\TextAt{8}{229.909}{254.582}{class}%
\TextAt{8}{257.729}{254.582}{field}%
\TextAt{8}{283.468}{254.582}{is}%
\TextAt{8}{295.615}{254.582}{in}%
\TextAt{8}{309.65}{254.582}{general}%
\TextAt{8}{350.357}{254.582}{not}%
\TextAt{8}{371.554}{254.582}{a}%
\TextAt{8}{381.692}{254.582}{ray}%
\TextAt{8}{402.243}{254.582}{class}%
\TextAt{8}{430.063}{254.582}{field}%
\TextAt{8}{455.802}{254.582}{in}%
\TextAt{8}{469.837}{254.582}{terms}%
\TextAt{8}{502.798}{254.582}{of}%
\TextAt{8}{516.51}{254.582}{the}%
\TextAt{8}{79.226}{237.536}{theory}%
\TextAt{8}{115.801}{237.536}{developed}%
\TextAt{8}{169.927}{237.536}{here.}%
\TextAt{8}{199.777}{237.536}{But}%
\TextAt{8}{222.842}{237.536}{it}%
\TextAt{8}{234.372}{237.536}{is}%
\TextAt{8}{245.961}{237.536}{in}%
\TextAt{8}{259.439}{237.536}{many}%
\TextAt{8}{291.136}{237.536}{other}%
\TextAt{8}{321.531}{237.536}{textbooks}%
\TextAt{8}{375.405}{237.536}{where}%
\TextAt{8}{409.04}{237.536}{ray}%
\TextAt{8}{429.033}{237.536}{class}%
\TextAt{8}{456.295}{237.536}{groups}%
\TextAt{8}{493.911}{237.536}{and}%
\TextAt{8}{516.498}{237.536}{ray}%
\TextAt{8}{79.226}{220.49}{class}%
\TextAt{8}{107.107}{220.49}{fields}%
\TextAt{8}{137.522}{220.49}{are}%
\TextAt{8}{157.476}{220.49}{defined}%
\TextAt{8}{198.232}{220.49}{differently}%
\TextAt{8}{254.279}{220.49}{(see}%
\TextAt{8}{278.19}{220.49}{for}%
\TextAt{8}{296.519}{220.49}{instance}%
\TextAt{8}{342.547}{220.49}{[107]).}%
\TextAt{8}{380.958}{220.49}{This}%
\TextAt{8}{408.122}{220.49}{other}%
\TextAt{8}{439.134}{220.49}{theory}%
\TextAt{8}{476.328}{220.49}{is}%
\TextAt{8}{488.535}{220.49}{obtained}%
\TextAt{8}{79.226}{203.444}{by}%
\TextAt{8}{95.485}{203.444}{equipping}%
\TextAt{8}{149.14}{203.444}{all}%
\TextAt{8}{165.399}{203.444}{number}%
\TextAt{8}{207.995}{203.444}{fields}%
\TextAt{8}{237.967}{203.444}{with}%
\TextAt{8}{264.627}{203.444}{the}%
\TextAt{8}{284.783}{203.444}{Minkowski}%
\TextAt{8}{343.053}{203.444}{metric}%
\TextAt{16}{195.427}{174.443}{\Character{104}}%
\TextAt{13}{200.076}{174.443}{x\Character{59}}%
\TextAt{13}{211.948}{174.443}{y}%
\TextAt{16}{218.109}{174.443}{\Character{105}}%
\TextAt{21}{222.758}{172.65}{K}%
\TextAt{9}{234.282}{174.443}{\Character{61}}%
\TextAt{20}{246.707}{185.801}{\Character{88}}%
\TextAt{21}{253.003}{160.695}{\Character{28}}%
\TextAt{13}{265.969}{174.443}{\Character{11}}%
\TextAt{21}{273.414}{172.65}{\Character{28}}%
\TextAt{13}{278.591}{174.443}{x}%
\TextAt{21}{285.243}{172.65}{\Character{28}}%
\Vector{q 0 G 0 g 0.478 w 218.419 -592.736 m 224.556 -592.736 l S Q}%
\TextAt{13}{290.419}{174.443}{y}%
\TextAt{21}{296.556}{171.62}{\Character{28}}%
\TextAt{9}{325.643}{174.443}{\Character{40}}%
\TextAt{13}{330.195}{174.443}{\Character{28}}%
\NativeMathAt{339.935}{174.443}{11.955}{\in}%
\TextAt{8}{351.226}{174.443}{Hom}%
\TextAt{9}{375.614}{174.443}{\Character{40}}%
\TextAt{13}{380.167}{174.443}{K}%
\TextAt{13}{390.329}{174.443}{\Character{59}}%
\TextAt{18}{395.578}{174.443}{C}%
\TextAt{9}{404.212}{174.443}{\Character{41}\Character{41}}%
\TextAt{13}{413.317}{174.443}{\Character{59}}%
\TextAt{13}{79.226}{136.704}{\Character{11}}%
\TextAt{21}{86.672}{134.911}{\Character{28}}%
\TextAt{9}{95.557}{136.704}{\Character{61}}%
\TextAt{9}{108.361}{136.704}{1}%
\TextAt{8}{118.348}{136.704}{if}%
\TextAt{13}{129.299}{136.704}{\Character{28}}%
\TextAt{9}{139.427}{136.704}{\Character{61}}%
\Vector{q 0 G 0 g 0.478 w 80.24 -630.475 m 86.659 -630.475 l S Q}%
\TextAt{13}{152.24}{136.704}{\Character{28}}%
\TextAt{8}{158.659}{136.704}{,}%
\TextAt{13}{166.093}{136.704}{\Character{11}}%
\TextAt{21}{173.539}{134.911}{\Character{28}}%
\TextAt{9}{182.424}{136.704}{\Character{61}}%
\TextAt{17}{196.433}{141.411}{1}%
\Vector{q 0 G 0 g 0.478 w 124.433 -634.307 m 128.667 -634.307 l S Q}%
\TextAt{17}{196.433}{132.582}{2}%
\TextAt{8}{205.987}{136.704}{if}%
\TextAt{13}{216.939}{136.704}{\Character{28}}%
\NativeMathAt{227.066}{136.704}{11.955}{\neq}%
\Vector{q 0 G 0 g 0.478 w 167.879 -630.475 m 174.298 -630.475 l S Q}%
\TextAt{13}{239.879}{136.704}{\Character{28}}%
\TextAt{8}{246.298}{136.704}{.}%
\TextAt{8}{255.429}{136.704}{A}%
\TextAt{8}{268.329}{136.704}{ray}%
\TextAt{8}{288.725}{136.704}{class}%
\TextAt{8}{316.39}{136.704}{group}%
\TextAt{8}{349.793}{136.704}{can}%
\TextAt{8}{371.48}{136.704}{then}%
\TextAt{8}{398.368}{136.704}{be}%
\TextAt{8}{414.52}{136.704}{attached}%
\TextAt{8}{462.556}{136.704}{to}%
\TextAt{8}{477.094}{136.704}{any}%
\TextAt{19}{499.438}{136.704}{replete}%
\TextAt{8}{79.226}{119.658}{module}%
\TextAt{14}{275.43}{102.612}{m}%
\TextAt{9}{287.918}{102.612}{\Character{61}}%
\TextAt{20}{300.343}{113.97}{\Character{89}}%
\TextAt{27}{305.987}{88.864}{p}%
\TextAt{14}{317.612}{102.612}{p}%
\TextAt{21}{323.595}{107.549}{n}%
\TextAt{28}{328.733}{106.552}{p}%
\TextAt{13}{333.314}{102.612}{\Character{59}}%
\TextAt{8}{297.211}{43.033}{364}%
\Destination{78.23}{809.815}{nav.page.364}%
\Destination{79.36}{460.555}{nav.print.400}%
\Destination{79.36}{697.955}{nav.number.6.6.9}%
\Hotlink{342.047}{218.255}{25.078}{11.047}{nav.cite.N107}%
\Hotlink{217.657}{554.655}{25.078}{11.047}{nav.number.6.5.6}%
\Hotlink{325.184}{342.355}{25.078}{11.047}{nav.number.6.1.3}%
}

% Page 377 (body page 365)
\Page{%
\TextAt{2}{79.226}{796.943}{VI.}%
\TextAt{2}{98.012}{796.943}{Global}%
\TextAt{2}{133.248}{796.943}{Class}%
\TextAt{2}{161.841}{796.943}{Field}%
\TextAt{2}{189.572}{796.943}{Theory}%
\Vector{q 0 G 0 g 0.299 w 7.226 18.729 m 460.769 18.729 l S Q}%
\TextAt{8}{79.226}{760.856}{where}%
\TextAt{13}{113.827}{760.856}{n}%
\TextAt{27}{120.814}{759.063}{p}%
\NativeMathAt{129.971}{760.856}{11.955}{\in}%
\TextAt{18}{142.61}{760.856}{Z}%
\TextAt{8}{150.58}{760.856}{,}%
\TextAt{13}{158.719}{760.856}{n}%
\TextAt{27}{165.706}{759.063}{p}%
\NativeMathAt{174.863}{760.856}{11.955}{\geq}%
\TextAt{9}{188.83}{760.856}{0}%
\TextAt{8}{194.683}{760.856}{,}%
\TextAt{8}{202.822}{760.856}{and}%
\TextAt{13}{226.376}{760.856}{n}%
\TextAt{27}{233.363}{759.063}{p}%
\TextAt{9}{242.52}{760.856}{\Character{61}}%
\TextAt{9}{256.28}{760.856}{0}%
\TextAt{8}{266.836}{760.856}{or}%
\TextAt{9}{281.937}{760.856}{\Character{61}}%
\TextAt{9}{295.697}{760.856}{1}%
\TextAt{8}{306.253}{760.856}{if}%
\TextAt{14}{317.768}{760.856}{p}%
\TextAt{16}{328.421}{760.856}{\Character{106}}%
\NativeMathAt{336.395}{760.856}{11.955}{\infty}%
\TextAt{8}{348.366}{760.856}{.}%
\TextAt{8}{359.189}{760.856}{The}%
\TextAt{8}{384.034}{760.856}{groups}%
\TextAt{13}{422.615}{760.856}{U}%
\TextAt{17}{431.814}{767.101}{\Character{40}}%
\TextAt{21}{435.108}{767.101}{n}%
\TextAt{28}{440.246}{766.105}{p}%
\TextAt{17}{444.329}{767.101}{\Character{41}}%
\TextAt{27}{430.566}{758.885}{p}%
\TextAt{8}{452.809}{760.856}{attached}%
\TextAt{8}{501.409}{760.856}{to}%
\TextAt{8}{516.511}{760.856}{the}%
\TextAt{8}{79.226}{743.81}{metrized}%
\TextAt{8}{127.346}{743.81}{number}%
\TextAt{8}{169.942}{743.81}{field}%
\TextAt{9}{195.299}{743.81}{\Character{40}}%
\TextAt{13}{199.851}{743.81}{K}%
\TextAt{13}{210.013}{743.81}{\Character{59}}%
\TextAt{16}{215.263}{743.81}{\Character{104}}%
\TextAt{13}{223.809}{743.81}{\Character{59}}%
\TextAt{16}{232.951}{743.81}{\Character{105}}%
\TextAt{21}{237.6}{742.017}{K}%
\TextAt{9}{245.803}{743.81}{\Character{41}}%
\TextAt{8}{254.253}{743.81}{are}%
\TextAt{8}{273.763}{743.81}{defined}%
\TextAt{8}{314.076}{743.81}{by}%
\TextAt{13}{113.084}{684.443}{U}%
\TextAt{17}{122.283}{690.688}{\Character{40}}%
\TextAt{21}{125.576}{690.688}{n}%
\TextAt{28}{130.715}{689.692}{p}%
\TextAt{17}{134.797}{690.688}{\Character{41}}%
\TextAt{27}{121.034}{682.472}{p}%
\TextAt{9}{141.91}{684.443}{\Character{61}}%
\TextAt{20}{154.335}{726.884}{\Character{56}}%
\TextAt{20}{154.335}{716.124}{\Character{62}}%
\TextAt{20}{154.335}{712.538}{\Character{62}}%
\TextAt{20}{154.335}{708.951}{\Character{62}}%
\TextAt{20}{154.335}{705.365}{\Character{62}}%
\TextAt{20}{154.335}{701.778}{\Character{62}}%
\TextAt{20}{154.335}{698.191}{\Character{60}}%
\TextAt{20}{154.335}{676.672}{\Character{62}}%
\TextAt{20}{154.335}{673.085}{\Character{62}}%
\TextAt{20}{154.335}{669.499}{\Character{62}}%
\TextAt{20}{154.335}{665.912}{\Character{62}}%
\TextAt{20}{154.335}{662.326}{\Character{62}}%
\TextAt{20}{154.335}{658.739}{\Character{58}}%
\TextAt{9}{164.962}{714.023}{1}%
\TextAt{9}{173.462}{714.023}{\Character{43}}%
\TextAt{14}{185.233}{714.023}{p}%
\TextAt{21}{191.216}{718.362}{n}%
\TextAt{28}{196.354}{717.365}{p}%
\TextAt{13}{200.935}{714.023}{\Character{59}}%
\TextAt{8}{218.134}{714.023}{for}%
\TextAt{13}{236.019}{714.023}{n}%
\TextAt{27}{243.007}{712.23}{p}%
\TextAt{13}{250.815}{714.023}{\Character{62}}%
\TextAt{9}{263.24}{714.023}{0}%
\TextAt{13}{269.093}{714.023}{\Character{59}}%
\TextAt{8}{278.235}{714.023}{and}%
\TextAt{13}{300.997}{714.023}{U}%
\TextAt{27}{308.947}{712.23}{p}%
\TextAt{8}{317.332}{714.023}{for}%
\TextAt{13}{335.217}{714.023}{n}%
\TextAt{27}{342.204}{712.23}{p}%
\TextAt{9}{350.012}{714.023}{\Character{61}}%
\TextAt{9}{362.422}{714.023}{0}%
\TextAt{13}{368.29}{714.023}{\Character{59}}%
\TextAt{8}{383.497}{714.023}{if}%
\TextAt{14}{394.221}{714.023}{p}%
\TextAt{18}{403.525}{714.023}{\Character{45}}%
\NativeMathAt{410.167}{714.023}{11.955}{\infty}%
\TextAt{13}{422.122}{714.023}{\Character{59}}%
\TextAt{18}{164.962}{693.568}{R}%
\TextAt{3}{173.596}{697.907}{\Character{3}}%
\TextAt{13}{178.328}{693.568}{\Character{59}}%
\TextAt{8}{383.497}{693.568}{if}%
\TextAt{14}{394.221}{693.568}{p}%
\TextAt{8}{404.101}{693.568}{is}%
\TextAt{8}{415.865}{693.568}{real}%
\TextAt{8}{438.628}{693.568}{and}%
\TextAt{13}{461.391}{693.568}{n}%
\TextAt{27}{468.378}{691.775}{p}%
\TextAt{9}{476.186}{693.568}{\Character{61}}%
\TextAt{9}{488.596}{693.568}{0}%
\TextAt{13}{494.464}{693.568}{\Character{59}}%
\TextAt{18}{164.962}{673.113}{R}%
\TextAt{3}{173.596}{677.452}{\Character{3}}%
\TextAt{17}{173.596}{670.158}{\Character{43}}%
\TextAt{13}{180.681}{673.113}{\Character{59}}%
\TextAt{8}{383.497}{673.113}{if}%
\TextAt{14}{394.221}{673.113}{p}%
\TextAt{8}{404.101}{673.113}{is}%
\TextAt{8}{415.865}{673.113}{real}%
\TextAt{8}{438.628}{673.113}{and}%
\TextAt{13}{461.391}{673.113}{n}%
\TextAt{27}{468.378}{671.32}{p}%
\TextAt{9}{476.186}{673.113}{\Character{61}}%
\TextAt{9}{488.596}{673.113}{1}%
\TextAt{13}{494.464}{673.113}{\Character{59}}%
\TextAt{18}{164.962}{652.658}{C}%
\TextAt{3}{173.596}{656.997}{\Character{3}}%
\TextAt{9}{181.649}{652.658}{\Character{61}}%
\TextAt{13}{194.075}{652.658}{K}%
\TextAt{3}{204.892}{656.997}{\Character{3}}%
\TextAt{27}{204.034}{649.703}{p}%
\TextAt{13}{209.625}{652.658}{\Character{59}}%
\TextAt{8}{383.497}{652.658}{if}%
\TextAt{14}{394.221}{652.658}{p}%
\TextAt{8}{404.101}{652.658}{is}%
\TextAt{8}{415.865}{652.658}{complex}%
\TextAt{13}{457.816}{652.658}{\Character{58}}%
\TextAt{8}{79.226}{625.076}{The}%
\TextAt{19}{103.591}{625.076}{congruence}%
\TextAt{19}{164.155}{625.076}{subgroup}%
\TextAt{19}{212.464}{625.076}{mod}%
\TextAt{14}{237.904}{625.076}{m}%
\TextAt{8}{251.279}{625.076}{of}%
\TextAt{9}{264.92}{625.076}{\Character{40}}%
\TextAt{13}{269.472}{625.076}{K}%
\TextAt{13}{279.634}{625.076}{\Character{59}}%
\TextAt{16}{284.883}{625.076}{\Character{104}}%
\TextAt{13}{293.741}{625.076}{\Character{59}}%
\TextAt{16}{303.193}{625.076}{\Character{105}}%
\TextAt{21}{307.843}{623.282}{K}%
\TextAt{9}{316.045}{625.076}{\Character{41}}%
\TextAt{8}{324.806}{625.076}{is}%
\TextAt{8}{336.88}{625.076}{then}%
\TextAt{8}{363.851}{625.076}{the}%
\TextAt{8}{384.318}{625.076}{subgroup}%
\TextAt{13}{435.427}{625.076}{C}%
\TextAt{27}{444.66}{629.414}{m}%
\TextAt{21}{443.802}{622.055}{K}%
\TextAt{9}{455.855}{625.076}{\Character{61}}%
\TextAt{13}{468.81}{625.076}{I}%
\TextAt{27}{474.913}{629.414}{m}%
\TextAt{21}{473.973}{622.055}{K}%
\TextAt{13}{482.176}{625.076}{K}%
\TextAt{3}{492.993}{629.414}{\Character{3}}%
\TextAt{9}{497.726}{625.076}{\Character{47}}%
\TextAt{13}{503.579}{625.076}{K}%
\TextAt{3}{514.396}{629.414}{\Character{3}}%
\TextAt{8}{523.337}{625.076}{of}%
\TextAt{13}{79.226}{608.03}{C}%
\TextAt{21}{87.601}{606.236}{K}%
\TextAt{8}{99.701}{608.03}{formed}%
\TextAt{8}{139.046}{608.03}{with}%
\TextAt{8}{165.706}{608.03}{the}%
\TextAt{8}{185.862}{608.03}{group}%
\TextAt{13}{268.429}{590.984}{I}%
\TextAt{27}{274.532}{595.92}{m}%
\TextAt{21}{273.592}{588.028}{K}%
\TextAt{9}{285.116}{590.984}{\Character{61}}%
\TextAt{20}{297.541}{602.341}{\Character{89}}%
\TextAt{27}{303.185}{577.235}{p}%
\TextAt{13}{314.81}{590.984}{U}%
\TextAt{17}{324.009}{597.229}{\Character{40}}%
\TextAt{21}{327.302}{597.229}{n}%
\TextAt{28}{332.44}{596.233}{p}%
\TextAt{17}{336.523}{597.229}{\Character{41}}%
\TextAt{27}{322.76}{589.013}{p}%
\TextAt{13}{340.315}{590.984}{\Character{59}}%
\TextAt{8}{79.226}{558.297}{and}%
\TextAt{8}{102.687}{558.297}{the}%
\TextAt{8}{123.542}{558.297}{factor}%
\TextAt{8}{157.738}{558.297}{group}%
\TextAt{13}{191.612}{558.297}{C}%
\TextAt{21}{199.987}{556.503}{K}%
\TextAt{9}{208.19}{558.297}{\Character{47}}%
\TextAt{13}{214.043}{558.297}{C}%
\TextAt{27}{223.277}{562.635}{m}%
\TextAt{21}{222.418}{555.276}{K}%
\TextAt{8}{235.216}{558.297}{is}%
\TextAt{8}{247.678}{558.297}{the}%
\TextAt{19}{268.533}{558.297}{ray}%
\TextAt{19}{289.463}{558.297}{class}%
\TextAt{19}{318.307}{558.297}{group}%
\TextAt{19}{351.193}{558.297}{mod}%
\TextAt{14}{377.048}{558.297}{m}%
\TextAt{8}{386.215}{558.297}{.}%
\TextAt{8}{396.758}{558.297}{The}%
\TextAt{8}{421.51}{558.297}{ray}%
\TextAt{8}{442.377}{558.297}{class}%
\TextAt{8}{470.512}{558.297}{field}%
\TextAt{8}{496.567}{558.297}{mod}%
\TextAt{14}{523.603}{558.297}{m}%
\TextAt{8}{79.226}{541.251}{of}%
\TextAt{9}{93.654}{541.251}{\Character{40}}%
\TextAt{13}{98.207}{541.251}{K}%
\TextAt{13}{108.369}{541.251}{\Character{59}}%
\TextAt{16}{113.618}{541.251}{\Character{104}}%
\TextAt{13}{123.263}{541.251}{\Character{59}}%
\TextAt{16}{133.503}{541.251}{\Character{105}}%
\TextAt{21}{138.152}{539.458}{K}%
\TextAt{9}{146.355}{541.251}{\Character{41}}%
\TextAt{8}{155.903}{541.251}{is}%
\TextAt{8}{168.765}{541.251}{again}%
\TextAt{8}{201.09}{541.251}{the}%
\TextAt{8}{222.345}{541.251}{class}%
\TextAt{8}{250.88}{541.251}{field}%
\TextAt{8}{277.335}{541.251}{of}%
\TextAt{13}{291.763}{541.251}{K}%
\TextAt{8}{307.577}{541.251}{corresponding}%
\TextAt{8}{383.861}{541.251}{to}%
\TextAt{8}{399.27}{541.251}{the}%
\TextAt{8}{420.524}{541.251}{group}%
\TextAt{13}{454.798}{541.251}{C}%
\TextAt{27}{464.032}{545.589}{m}%
\TextAt{21}{463.173}{538.23}{K}%
\NativeMathAt{476.568}{541.251}{11.955}{\subseteq}%
\TextAt{13}{491.059}{541.251}{C}%
\TextAt{21}{499.434}{539.458}{K}%
\TextAt{8}{507.637}{541.251}{.}%
\TextAt{8}{519.38}{541.251}{As}%
\TextAt{8}{79.226}{524.205}{explained}%
\TextAt{8}{132.382}{524.205}{in}%
\TextAt{8}{146.84}{524.205}{chap.}%
\TextAt{8}{178.536}{524.205}{III,}%
\TextAt{8}{199.796}{524.205}{§3,}%
\TextAt{8}{219.488}{524.205}{the}%
\TextAt{8}{240.449}{524.205}{infinite}%
\TextAt{8}{280.92}{524.205}{places}%
\TextAt{14}{316.252}{524.205}{p}%
\TextAt{8}{326.937}{524.205}{have}%
\TextAt{8}{354.736}{524.205}{to}%
\TextAt{8}{369.851}{524.205}{be}%
\TextAt{8}{386.58}{524.205}{considered}%
\TextAt{8}{444.673}{524.205}{as}%
\TextAt{19}{459.848}{524.205}{ramified}%
\TextAt{8}{505.95}{524.205}{in}%
\TextAt{8}{520.408}{524.205}{an}%
\TextAt{8}{79.226}{507.159}{extension}%
\TextAt{13}{131.246}{507.159}{L}%
\TextAt{16}{142.965}{507.159}{\Character{106}}%
\TextAt{13}{150.039}{507.159}{K}%
\TextAt{8}{165.009}{507.159}{if}%
\TextAt{13}{175.987}{507.159}{L}%
\TextAt{27}{183.951}{505.366}{p}%
\NativeMathAt{192.192}{507.159}{11.955}{\neq}%
\TextAt{13}{205.051}{507.159}{K}%
\TextAt{27}{215.01}{505.366}{p}%
\TextAt{8}{219.496}{507.159}{.}%
\TextAt{8}{228.707}{507.159}{Likewise,}%
\TextAt{8}{279.32}{507.159}{the}%
\TextAt{19}{299.731}{507.159}{conductor}%
\TextAt{8}{354.465}{507.159}{of}%
\TextAt{8}{368.049}{507.159}{an}%
\TextAt{8}{384.562}{507.159}{abelian}%
\TextAt{8}{425.464}{507.159}{extension}%
\TextAt{13}{477.484}{507.159}{L}%
\TextAt{16}{489.202}{507.159}{\Character{106}}%
\TextAt{13}{496.277}{507.159}{K}%
\TextAt{8}{507.095}{507.159}{,}%
\TextAt{8}{514.562}{507.159}{i.e.,}%
\TextAt{8}{79.226}{490.113}{the}%
\TextAt{8}{99.383}{490.113}{gcd}%
\TextAt{8}{120.842}{490.113}{of}%
\TextAt{8}{134.172}{490.113}{all}%
\TextAt{8}{150.431}{490.113}{modules}%
\TextAt{14}{196.339}{490.113}{m}%
\TextAt{9}{208.827}{490.113}{\Character{61}}%
\TextAt{20}{221.252}{499.079}{\Character{81}}%
\TextAt{27}{232.543}{486.626}{p}%
\TextAt{14}{239.022}{490.113}{p}%
\TextAt{21}{245.005}{494.451}{n}%
\TextAt{28}{250.144}{493.455}{p}%
\TextAt{8}{258.622}{490.113}{such}%
\TextAt{8}{285.019}{490.113}{that}%
\TextAt{13}{310.388}{490.113}{C}%
\TextAt{27}{319.621}{494.451}{m}%
\TextAt{21}{318.763}{487.092}{K}%
\NativeMathAt{330.286}{490.113}{11.955}{\subseteq}%
\TextAt{13}{342.906}{490.113}{N}%
\TextAt{21}{352.28}{488.253}{L}%
\TextAt{3}{358.039}{488.253}{\Character{106}}%
\TextAt{21}{360.391}{488.253}{K}%
\TextAt{13}{368.594}{490.113}{C}%
\TextAt{21}{376.969}{488.32}{L}%
\TextAt{8}{383.226}{490.113}{,}%
\TextAt{8}{390.375}{490.113}{is}%
\TextAt{8}{402.139}{490.113}{the}%
\TextAt{8}{422.295}{490.113}{replete}%
\TextAt{8}{460.659}{490.113}{ideal}%
\TextAt{14}{281.721}{460.087}{f}%
\TextAt{9}{288.942}{460.087}{\Character{61}}%
\TextAt{20}{301.367}{471.444}{\Character{89}}%
\TextAt{27}{307.011}{446.338}{p}%
\TextAt{14}{318.636}{460.087}{f}%
\TextAt{27}{322.537}{458.293}{p}%
\TextAt{13}{327.023}{460.087}{\Character{59}}%
\TextAt{8}{79.226}{422.74}{where}%
\TextAt{8}{113.024}{422.74}{now}%
\TextAt{8}{137.402}{422.74}{for}%
\TextAt{8}{155.276}{422.74}{an}%
\TextAt{8}{171.523}{422.74}{infinite}%
\TextAt{8}{211.18}{422.74}{place}%
\TextAt{14}{241.08}{422.74}{p}%
\TextAt{8}{247.063}{422.74}{,}%
\TextAt{8}{254.204}{422.74}{we}%
\TextAt{8}{271.42}{422.74}{have}%
\TextAt{14}{298.404}{422.74}{f}%
\TextAt{27}{302.304}{420.946}{p}%
\TextAt{9}{310.112}{422.74}{\Character{61}}%
\TextAt{14}{322.537}{422.74}{p}%
\TextAt{21}{328.521}{427.078}{n}%
\TextAt{28}{333.659}{426.082}{p}%
\TextAt{8}{342.126}{422.74}{with}%
\TextAt{13}{368.775}{422.74}{n}%
\TextAt{27}{375.762}{420.946}{p}%
\TextAt{9}{383.57}{422.74}{\Character{61}}%
\TextAt{9}{395.98}{422.74}{0}%
\TextAt{8}{405.735}{422.74}{if}%
\TextAt{13}{416.448}{422.74}{L}%
\TextAt{27}{424.412}{420.946}{p}%
\TextAt{9}{432.22}{422.74}{\Character{61}}%
\TextAt{13}{444.645}{422.74}{K}%
\TextAt{27}{454.604}{420.946}{p}%
\TextAt{8}{459.091}{422.74}{,}%
\TextAt{8}{466.232}{422.74}{and}%
\TextAt{13}{488.983}{422.74}{n}%
\TextAt{27}{495.971}{420.946}{p}%
\TextAt{9}{503.778}{422.74}{\Character{61}}%
\TextAt{9}{516.188}{422.74}{1}%
\TextAt{8}{525.943}{422.74}{if}%
\TextAt{13}{79.226}{405.694}{L}%
\TextAt{27}{87.191}{403.901}{p}%
\NativeMathAt{95.102}{405.694}{11.955}{\neq}%
\TextAt{13}{107.631}{405.694}{K}%
\TextAt{27}{117.59}{403.901}{p}%
\TextAt{8}{122.077}{405.694}{.}%
\TextAt{8}{130.708}{405.694}{Corollary}%
\TextAt{8}{182.487}{405.694}{(6.6)}%
\TextAt{8}{210.522}{405.694}{then}%
\TextAt{8}{237.243}{405.694}{continues}%
\TextAt{8}{288.747}{405.694}{to}%
\TextAt{8}{303.118}{405.694}{hold:}%
\TextAt{8}{333.805}{405.694}{a}%
\TextAt{8}{343.622}{405.694}{place}%
\TextAt{14}{373.594}{405.694}{p}%
\TextAt{8}{383.535}{405.694}{is}%
\TextAt{8}{395.36}{405.694}{ramified}%
\TextAt{8}{440.946}{405.694}{in}%
\TextAt{13}{454.66}{405.694}{L}%
\TextAt{8}{466.583}{405.694}{if}%
\TextAt{8}{477.367}{405.694}{and}%
\TextAt{8}{500.191}{405.694}{only}%
\TextAt{8}{525.943}{405.694}{if}%
\TextAt{14}{79.226}{388.648}{p}%
\TextAt{8}{89.107}{388.648}{occurs}%
\TextAt{8}{125.259}{388.648}{in}%
\TextAt{8}{138.912}{388.648}{the}%
\TextAt{8}{159.068}{388.648}{conductor}%
\TextAt{14}{213.703}{388.648}{f}%
\TextAt{8}{217.604}{388.648}{.}%
\TextAt{8}{97.159}{368.613}{This}%
\TextAt{8}{122.869}{368.613}{entails}%
\TextAt{8}{158.669}{368.613}{the}%
\TextAt{8}{177.816}{368.613}{following}%
\TextAt{8}{226.241}{368.613}{modifications}%
\TextAt{8}{297.165}{368.613}{of}%
\TextAt{8}{309.486}{368.613}{the}%
\TextAt{8}{328.632}{368.613}{above}%
\TextAt{8}{360.798}{368.613}{theory,}%
\TextAt{8}{399.012}{368.613}{as}%
\TextAt{8}{412.373}{368.613}{far}%
\TextAt{8}{429.248}{368.613}{as}%
\TextAt{8}{442.608}{368.613}{ray}%
\TextAt{8}{461.767}{368.613}{class}%
\TextAt{8}{488.194}{368.613}{fields}%
\TextAt{8}{517.156}{368.613}{are}%
\TextAt{8}{79.226}{351.567}{concerned.}%
\TextAt{8}{138.126}{351.567}{The}%
\TextAt{8}{161.355}{351.567}{ray}%
\TextAt{8}{180.7}{351.567}{class}%
\TextAt{8}{207.313}{351.567}{field}%
\TextAt{8}{231.845}{351.567}{mod}%
\TextAt{9}{257.359}{351.567}{1}%
\TextAt{8}{266.285}{351.567}{is}%
\TextAt{8}{277.225}{351.567}{the}%
\TextAt{19}{296.557}{351.567}{small}%
\TextAt{8}{327.438}{351.567}{Hilbert}%
\TextAt{8}{366.927}{351.567}{class}%
\TextAt{8}{393.54}{351.567}{field.}%
\TextAt{8}{423.173}{351.567}{It}%
\TextAt{8}{435.021}{351.567}{is}%
\TextAt{8}{445.961}{351.567}{now}%
\TextAt{8}{469.526}{351.567}{the}%
\TextAt{8}{488.858}{351.567}{maximal}%
\TextAt{8}{79.226}{334.521}{abelian}%
\TextAt{8}{119.473}{334.521}{extension}%
\TextAt{8}{170.839}{334.521}{of}%
\TextAt{13}{183.768}{334.521}{K}%
\TextAt{8}{198.083}{334.521}{which}%
\TextAt{8}{231.169}{334.521}{is}%
\TextAt{8}{242.533}{334.521}{unramified}%
\TextAt{8}{300.665}{334.521}{at}%
\TextAt{19}{314.575}{334.521}{all}%
\TextAt{8}{331.533}{334.521}{places.}%
\TextAt{8}{370.477}{334.521}{The}%
\TextAt{8}{394.131}{334.521}{big}%
\TextAt{8}{413.241}{334.521}{Hilbert}%
\TextAt{8}{453.154}{334.521}{class}%
\TextAt{8}{480.19}{334.521}{field}%
\TextAt{8}{505.147}{334.521}{is}%
\TextAt{8}{516.51}{334.521}{the}%
\TextAt{8}{79.226}{317.475}{ray}%
\TextAt{8}{99.503}{317.475}{class}%
\TextAt{8}{127.049}{317.475}{field}%
\TextAt{8}{152.515}{317.475}{for}%
\TextAt{8}{170.509}{317.475}{the}%
\TextAt{8}{190.774}{317.475}{module}%
\TextAt{14}{232.176}{317.475}{m}%
\TextAt{9}{244.849}{317.475}{\Character{61}}%
\TextAt{20}{257.46}{326.442}{\Character{81}}%
\TextAt{27}{268.751}{313.988}{p}%
\TextAt{3}{272.739}{313.988}{\Character{106}}%
\NativeMathAt{275.09}{313.988}{7.97}{\infty}%
\TextAt{14}{286.051}{317.475}{p}%
\TextAt{8}{292.034}{317.475}{.}%
\TextAt{8}{300.808}{317.475}{In}%
\TextAt{8}{315.538}{317.475}{the}%
\TextAt{8}{335.804}{317.475}{case}%
\TextAt{8}{360.683}{317.475}{of}%
\TextAt{8}{374.122}{317.475}{the}%
\TextAt{8}{394.387}{317.475}{base}%
\TextAt{8}{420.57}{317.475}{field}%
\TextAt{18}{446.036}{317.475}{Q}%
\TextAt{8}{455.335}{317.475}{,}%
\TextAt{8}{462.62}{317.475}{the}%
\TextAt{8}{482.885}{317.475}{field}%
\TextAt{18}{508.351}{317.475}{Q}%
\TextAt{9}{517.649}{317.475}{\Character{40}}%
\TextAt{13}{522.202}{317.475}{\Character{16}}%
\TextAt{9}{528.217}{317.475}{\Character{41}}%
\TextAt{8}{79.226}{300.429}{of}%
\TextAt{13}{92.285}{300.429}{m}%
\TextAt{8}{102.524}{300.429}{-th}%
\TextAt{8}{121.107}{300.429}{roots}%
\TextAt{8}{150.496}{300.429}{of}%
\TextAt{8}{163.555}{300.429}{unity}%
\TextAt{8}{193.853}{300.429}{is}%
\TextAt{8}{205.346}{300.429}{the}%
\TextAt{8}{225.232}{300.429}{ray}%
\TextAt{8}{245.129}{300.429}{class}%
\TextAt{8}{272.295}{300.429}{field}%
\TextAt{8}{297.381}{300.429}{mod}%
\TextAt{13}{323.447}{300.429}{mp}%
\NativeMathAt{339.561}{298.636}{7.97}{\infty}%
\TextAt{8}{348.528}{300.429}{,}%
\TextAt{8}{355.46}{300.429}{where}%
\TextAt{13}{388.998}{300.429}{p}%
\NativeMathAt{394.874}{298.636}{7.97}{\infty}%
\TextAt{8}{407.466}{300.429}{is}%
\TextAt{8}{418.959}{300.429}{the}%
\TextAt{8}{438.844}{300.429}{infinite}%
\TextAt{8}{478.241}{300.429}{place.}%
\TextAt{8}{512.613}{300.429}{The}%
\TextAt{8}{79.226}{283.383}{ray}%
\TextAt{8}{98.637}{283.383}{class}%
\TextAt{8}{125.317}{283.383}{field}%
\TextAt{8}{149.917}{283.383}{for}%
\TextAt{8}{167.044}{283.383}{the}%
\TextAt{8}{186.443}{283.383}{module}%
\TextAt{13}{226.979}{283.383}{m}%
\TextAt{8}{240.358}{283.383}{becomes}%
\TextAt{8}{286.154}{283.383}{the}%
\TextAt{8}{305.553}{283.383}{maximal}%
\TextAt{8}{352.605}{283.383}{real}%
\TextAt{8}{374.61}{283.383}{subextension}%
\TextAt{18}{443.563}{283.383}{Q}%
\TextAt{9}{452.862}{283.383}{\Character{40}}%
\TextAt{13}{457.414}{283.383}{\Character{16}}%
\TextAt{9}{464.538}{283.383}{\Character{43}}%
\TextAt{13}{474.75}{283.383}{\Character{16}}%
\TextAt{3}{480.766}{287.722}{\Character{0}}%
\TextAt{17}{487.352}{287.722}{1}%
\TextAt{9}{492.085}{283.383}{\Character{41}}%
\TextAt{8}{496.637}{283.383}{,}%
\TextAt{8}{503.18}{283.383}{which}%
\TextAt{8}{79.226}{266.337}{was}%
\TextAt{8}{101.922}{266.337}{not}%
\TextAt{8}{122.933}{266.337}{a}%
\TextAt{8}{132.885}{266.337}{ray}%
\TextAt{8}{153.25}{266.337}{class}%
\TextAt{8}{180.883}{266.337}{field}%
\TextAt{8}{206.437}{266.337}{before.}%
\TextAt{8}{246.69}{266.337}{This}%
\TextAt{8}{273.606}{266.337}{is}%
\TextAt{8}{285.566}{266.337}{the}%
\TextAt{8}{305.919}{266.337}{theory}%
\TextAt{8}{342.866}{266.337}{one}%
\TextAt{8}{364.522}{266.337}{finds}%
\TextAt{8}{392.741}{266.337}{in}%
\TextAt{8}{406.591}{266.337}{the}%
\TextAt{8}{426.944}{266.337}{textbooks}%
\TextAt{8}{481.19}{266.337}{alluded}%
\TextAt{8}{522.357}{266.337}{to}%
\TextAt{8}{79.226}{249.291}{above.}%
\TextAt{8}{119.709}{249.291}{It}%
\TextAt{8}{133.301}{249.291}{corresponds}%
\TextAt{8}{198.407}{249.291}{to}%
\TextAt{8}{213.636}{249.291}{the}%
\TextAt{8}{234.712}{249.291}{number}%
\TextAt{8}{278.227}{249.291}{fields}%
\TextAt{8}{309.117}{249.291}{with}%
\TextAt{8}{336.696}{249.291}{Minkowski}%
\TextAt{8}{395.885}{249.291}{metric.}%
\TextAt{8}{439.608}{249.291}{The}%
\TextAt{8}{464.58}{249.291}{theory}%
\TextAt{8}{502.25}{249.291}{of}%
\TextAt{8}{516.498}{249.291}{ray}%
\TextAt{8}{79.226}{232.245}{class}%
\TextAt{8}{106.026}{232.245}{fields}%
\TextAt{8}{135.36}{232.245}{according}%
\TextAt{8}{187.409}{232.245}{to}%
\TextAt{8}{201.082}{232.245}{the}%
\TextAt{8}{220.601}{232.245}{treatment}%
\TextAt{8}{274.276}{232.245}{of}%
\TextAt{8}{286.969}{232.245}{this}%
\TextAt{8}{309.154}{232.245}{book}%
\TextAt{8}{337.46}{232.245}{is}%
\TextAt{8}{348.587}{232.245}{forced}%
\TextAt{8}{382.739}{232.245}{upon}%
\TextAt{8}{411.69}{232.245}{us}%
\TextAt{8}{426.069}{232.245}{already}%
\TextAt{8}{466.737}{232.245}{by}%
\TextAt{8}{482.358}{232.245}{the}%
\TextAt{8}{501.877}{232.245}{choice}%
\TextAt{8}{79.226}{215.199}{of}%
\TextAt{8}{92.117}{215.199}{the}%
\TextAt{8}{111.835}{215.199}{standard}%
\TextAt{8}{160.245}{215.199}{metric}%
\TextAt{16}{196.222}{215.199}{\Character{104}}%
\TextAt{13}{200.871}{215.199}{x\Character{59}}%
\TextAt{13}{212.743}{215.199}{y}%
\TextAt{16}{218.904}{215.199}{\Character{105}}%
\TextAt{9}{226.874}{215.199}{\Character{61}}%
\TextAt{20}{239.3}{224.166}{\Character{80}}%
\TextAt{21}{251.919}{211.712}{\Character{28}}%
\TextAt{13}{259.088}{215.199}{x}%
\TextAt{21}{265.74}{213.406}{\Character{28}}%
\Vector{q 0 G 0 g 0.478 w 198.916 -551.98 m 205.053 -551.98 l S Q}%
\TextAt{13}{270.916}{215.199}{y}%
\TextAt{21}{277.053}{212.377}{\Character{28}}%
\TextAt{8}{285.688}{215.199}{on}%
\TextAt{13}{301.508}{215.199}{K}%
\TextAt{25}{311.468}{213.406}{R}%
\TextAt{8}{321.181}{215.199}{taken}%
\TextAt{8}{352.614}{215.199}{in}%
\TextAt{8}{365.828}{215.199}{chap.}%
\TextAt{8}{396.282}{215.199}{I,}%
\TextAt{8}{407.212}{215.199}{§5.}%
\TextAt{8}{427.051}{215.199}{It}%
\TextAt{8}{439.285}{215.199}{is}%
\TextAt{8}{450.61}{215.199}{compatible}%
\TextAt{8}{510.007}{215.199}{with}%
\TextAt{8}{79.226}{198.153}{the}%
\TextAt{8}{99.383}{198.153}{Riemann–Roch}%
\TextAt{8}{181.012}{198.153}{theory}%
\TextAt{8}{217.763}{198.153}{of}%
\TextAt{8}{231.093}{198.153}{chap.}%
\TextAt{8}{261.985}{198.153}{III,}%
\TextAt{8}{282.44}{198.153}{and}%
\TextAt{8}{305.203}{198.153}{has}%
\TextAt{8}{326.077}{198.153}{the}%
\TextAt{8}{346.233}{198.153}{advantage}%
\TextAt{8}{401.538}{198.153}{of}%
\TextAt{8}{414.868}{198.153}{being}%
\TextAt{8}{446.405}{198.153}{simpler.}%
\TextAt{8}{97.159}{175.197}{Over}%
\TextAt{8}{125.635}{175.197}{the}%
\TextAt{8}{145.647}{175.197}{field}%
\TextAt{18}{170.86}{175.197}{Q}%
\TextAt{8}{180.159}{175.197}{,}%
\TextAt{8}{187.192}{175.197}{the}%
\TextAt{8}{207.205}{175.197}{ray}%
\TextAt{8}{227.229}{175.197}{class}%
\TextAt{8}{254.522}{175.197}{field}%
\TextAt{8}{279.734}{175.197}{mod}%
\TextAt{9}{305.927}{175.197}{\Character{40}}%
\TextAt{13}{310.48}{175.197}{m}%
\TextAt{9}{320.719}{175.197}{\Character{41}}%
\TextAt{8}{329.024}{175.197}{can}%
\TextAt{8}{350.34}{175.197}{be}%
\TextAt{8}{366.12}{175.197}{generated,}%
\TextAt{8}{422.588}{175.197}{according}%
\TextAt{8}{475.13}{175.197}{to}%
\TextAt{8}{489.296}{175.197}{(6.7),}%
\TextAt{8}{520.408}{175.197}{by}%
\TextAt{8}{79.226}{158.151}{the}%
\TextAt{13}{98.928}{158.151}{m}%
\TextAt{8}{109.168}{158.151}{-th}%
\TextAt{8}{127.567}{158.151}{roots}%
\TextAt{8}{156.774}{158.151}{of}%
\TextAt{8}{169.649}{158.151}{unity,}%
\TextAt{8}{202.127}{158.151}{i.e.,}%
\TextAt{8}{223.869}{158.151}{by}%
\TextAt{8}{239.674}{158.151}{special}%
\TextAt{8}{277.321}{158.151}{values}%
\TextAt{8}{311.728}{158.151}{of}%
\TextAt{8}{324.604}{158.151}{the}%
\TextAt{8}{344.306}{158.151}{exponential}%
\TextAt{8}{406.617}{158.151}{function}%
\TextAt{13}{452.01}{158.151}{e}%
\TextAt{17}{457.436}{162.49}{2}%
\TextAt{21}{461.67}{162.49}{\Character{25}}%
\TextAt{21}{466.811}{162.49}{iz}%
\TextAt{8}{474.476}{158.151}{.}%
\TextAt{8}{482.773}{158.151}{The}%
\TextAt{8}{506.372}{158.151}{ques-}%
\TextAt{8}{79.226}{141.105}{tion}%
\TextAt{8}{102.821}{141.105}{suggested}%
\TextAt{8}{155.156}{141.105}{by}%
\TextAt{8}{170.944}{141.105}{this}%
\TextAt{8}{193.296}{141.105}{observation}%
\TextAt{8}{255.016}{141.105}{is}%
\TextAt{8}{266.309}{141.105}{whether}%
\TextAt{8}{310.706}{141.105}{one}%
\TextAt{8}{331.695}{141.105}{may}%
\TextAt{8}{356.593}{141.105}{construct}%
\TextAt{8}{407.565}{141.105}{the}%
\TextAt{8}{427.25}{141.105}{abelian}%
\TextAt{8}{467.427}{141.105}{extensions}%
\TextAt{8}{523.337}{141.105}{of}%
\TextAt{8}{79.226}{124.059}{an}%
\TextAt{8}{95.512}{124.059}{arbitrary}%
\TextAt{8}{145.308}{124.059}{number}%
\TextAt{8}{187.931}{124.059}{field}%
\TextAt{8}{213.315}{124.059}{in}%
\TextAt{8}{226.995}{124.059}{a}%
\TextAt{8}{236.777}{124.059}{similarly}%
\TextAt{8}{284.672}{124.059}{concrete}%
\TextAt{8}{330.87}{124.059}{way,}%
\TextAt{8}{356.918}{124.059}{via}%
\TextAt{8}{376.133}{124.059}{special}%
\TextAt{8}{414.261}{124.059}{values}%
\TextAt{8}{449.149}{124.059}{of}%
\TextAt{8}{462.506}{124.059}{analytic}%
\TextAt{8}{507.09}{124.059}{func-}%
\TextAt{8}{79.226}{107.013}{tions.}%
\TextAt{8}{113.338}{107.013}{This}%
\TextAt{8}{140.351}{107.013}{was}%
\TextAt{8}{163.144}{107.013}{the}%
\TextAt{8}{183.594}{107.013}{historic}%
\TextAt{8}{225.575}{107.013}{origin}%
\TextAt{8}{259.044}{107.013}{of}%
\TextAt{8}{272.667}{107.013}{the}%
\TextAt{8}{293.117}{107.013}{notion}%
\TextAt{8}{329.837}{107.013}{of}%
\TextAt{8}{343.461}{107.013}{class}%
\TextAt{8}{371.191}{107.013}{field.}%
\TextAt{8}{401.979}{107.013}{A}%
\TextAt{8}{414.945}{107.013}{completely}%
\TextAt{8}{474.094}{107.013}{satisfactory}%
\TextAt{8}{79.226}{89.967}{answer}%
\TextAt{8}{117.447}{89.967}{to}%
\TextAt{8}{131.219}{89.967}{this}%
\TextAt{8}{153.504}{89.967}{question}%
\TextAt{8}{199.531}{89.967}{has}%
\TextAt{8}{219.867}{89.967}{been}%
\TextAt{8}{246.958}{89.967}{given}%
\TextAt{8}{276.989}{89.967}{only}%
\TextAt{8}{302.143}{89.967}{in}%
\TextAt{8}{315.258}{89.967}{the}%
\TextAt{8}{334.877}{89.967}{case}%
\TextAt{8}{359.11}{89.967}{of}%
\TextAt{8}{371.902}{89.967}{an}%
\TextAt{8}{387.623}{89.967}{imaginary}%
\TextAt{8}{442.055}{89.967}{quadratic}%
\TextAt{8}{493.881}{89.967}{field}%
\TextAt{13}{518.7}{89.967}{K}%
\TextAt{8}{529.518}{89.967}{.}%
\TextAt{8}{79.226}{72.921}{The}%
\TextAt{8}{103.508}{72.921}{results}%
\TextAt{8}{140.929}{72.921}{for}%
\TextAt{8}{159.042}{72.921}{this}%
\TextAt{8}{182.093}{72.921}{case}%
\TextAt{8}{207.092}{72.921}{are}%
\TextAt{8}{226.831}{72.921}{subsumed}%
\TextAt{8}{281.157}{72.921}{under}%
\TextAt{8}{314.549}{72.921}{the}%
\TextAt{8}{334.933}{72.921}{name}%
\TextAt{8}{366.377}{72.921}{of}%
\TextAt{12}{380.449}{72.921}{Kronecker’s}%
\TextAt{12}{454.583}{72.921}{Jugendtraum}%
\TextAt{8}{297.211}{43.033}{365}%
\Destination{78.23}{809.815}{nav.page.365}%
\Destination{79.36}{505.155}{nav.print.401}%
\Hotlink{178.036}{521.955}{14.306}{11.047}{nav.chapter.3}%
\Hotlink{204.975}{521.955}{6.858}{11.047}{nav.section.3.3}%
\Hotlink{395.782}{212.955}{5.22}{11.047}{nav.chapter.1}%
\Hotlink{412.391}{212.955}{6.858}{11.047}{nav.section.1.5}%
\Hotlink{261.485}{195.955}{14.294}{11.047}{nav.chapter.3}%
\Hotlink{181.987}{403.455}{25.078}{11.047}{nav.number.6.6.6}%
\Hotlink{488.796}{172.955}{25.078}{11.047}{nav.number.6.6.7}%
}

% Page 378 (body page 366)
\Page{%
\TextAt{2}{79.226}{796.943}{VI.}%
\TextAt{2}{98.012}{796.943}{Global}%
\TextAt{2}{133.248}{796.943}{Class}%
\TextAt{2}{161.841}{796.943}{Field}%
\TextAt{2}{189.572}{796.943}{Theory}%
\Vector{q 0 G 0 g 0.299 w 7.226 18.729 m 460.769 18.729 l S Q}%
\TextAt{8}{79.226}{761.124}{(Kronecker’s}%
\TextAt{8}{148.239}{761.124}{dream}%
\TextAt{8}{184.997}{761.124}{of}%
\TextAt{8}{199.314}{761.124}{his}%
\TextAt{8}{218.57}{761.124}{youth).}%
\TextAt{8}{263.815}{761.124}{We}%
\TextAt{8}{284.947}{761.124}{will}%
\TextAt{8}{308.04}{761.124}{briefly}%
\TextAt{8}{345.12}{761.124}{describe}%
\TextAt{8}{391.359}{761.124}{them}%
\TextAt{8}{422.258}{761.124}{here.}%
\TextAt{8}{455.13}{761.124}{For}%
\TextAt{8}{477.087}{761.124}{the}%
\TextAt{8}{498.231}{761.124}{proofs,}%
\TextAt{8}{79.226}{744.078}{which}%
\TextAt{8}{112.482}{744.078}{presuppose}%
\TextAt{8}{172.529}{744.078}{an}%
\TextAt{8}{188.557}{744.078}{in-depth}%
\TextAt{8}{235.142}{744.078}{knowledge}%
\TextAt{8}{291.495}{744.078}{of}%
\TextAt{8}{304.594}{744.078}{the}%
\TextAt{8}{324.52}{744.078}{theory}%
\TextAt{8}{361.039}{744.078}{of}%
\TextAt{19}{374.138}{744.078}{elliptic}%
\TextAt{19}{412.744}{744.078}{curves}%
\TextAt{8}{444.88}{744.078}{,}%
\TextAt{8}{451.844}{744.078}{we}%
\TextAt{8}{468.841}{744.078}{have}%
\TextAt{8}{495.605}{744.078}{to}%
\TextAt{8}{509.684}{744.078}{refer}%
\TextAt{8}{79.226}{727.032}{to}%
\TextAt{8}{93.536}{727.032}{[96]}%
\TextAt{8}{115.654}{727.032}{and}%
\TextAt{8}{138.416}{727.032}{[28].}%
\TextAt{8}{97.159}{706.997}{An}%
\TextAt{8}{117.14}{706.997}{elliptic}%
\TextAt{8}{156.309}{706.997}{curve}%
\TextAt{8}{188.329}{706.997}{is}%
\TextAt{8}{200.898}{706.997}{given}%
\TextAt{8}{232.273}{706.997}{as}%
\TextAt{8}{247.448}{706.997}{the}%
\TextAt{8}{268.41}{706.997}{quotient}%
\TextAt{13}{315.398}{706.997}{E}%
\TextAt{9}{329.457}{706.997}{\Character{61}}%
\TextAt{18}{343.254}{706.997}{C}%
\TextAt{9}{351.889}{706.997}{\Character{47}\Character{0}}%
\TextAt{8}{369.758}{706.997}{of}%
\TextAt{18}{383.893}{706.997}{C}%
\TextAt{8}{397.23}{706.997}{by}%
\TextAt{8}{414.294}{706.997}{a}%
\TextAt{8}{424.855}{706.997}{complete}%
\TextAt{8}{475.083}{706.997}{lattice}%
\TextAt{9}{511.658}{706.997}{\Character{0}}%
\TextAt{9}{523.649}{706.997}{\Character{61}}%
\TextAt{18}{79.226}{689.951}{Z}%
\TextAt{13}{87.196}{689.951}{\Character{33}}%
\TextAt{17}{94.486}{688.158}{1}%
\TextAt{9}{101.226}{689.951}{\Character{43}}%
\TextAt{18}{112.339}{689.951}{Z}%
\TextAt{13}{120.309}{689.951}{\Character{33}}%
\TextAt{17}{127.598}{688.158}{2}%
\TextAt{8}{135.91}{689.951}{in}%
\TextAt{18}{149.246}{689.951}{C}%
\TextAt{8}{157.88}{689.951}{.}%
\TextAt{8}{166.223}{689.951}{This}%
\TextAt{8}{192.626}{689.951}{is}%
\TextAt{8}{204.072}{689.951}{a}%
\TextAt{8}{213.511}{689.951}{torus}%
\TextAt{8}{243.177}{689.951}{which}%
\TextAt{8}{276.346}{689.951}{receives}%
\TextAt{8}{319.008}{689.951}{the}%
\TextAt{8}{338.847}{689.951}{structure}%
\TextAt{8}{388.67}{689.951}{of}%
\TextAt{8}{401.683}{689.951}{an}%
\TextAt{8}{417.624}{689.951}{algebraic}%
\TextAt{8}{466.742}{689.951}{curve}%
\TextAt{8}{497.64}{689.951}{via}%
\TextAt{8}{516.51}{689.951}{the}%
\TextAt{12}{79.226}{672.905}{Weierstrass}%
\TextAt{13}{151.184}{672.905}{\Character{125}}%
\TextAt{12}{158.637}{672.905}{-function}%
\TextAt{13}{194.57}{641.57}{\Character{125}}%
\TextAt{9}{202.023}{641.57}{\Character{40}}%
\TextAt{13}{206.575}{641.57}{z}%
\TextAt{9}{212.546}{641.57}{\Character{41}}%
\TextAt{9}{220.401}{641.57}{\Character{61}}%
\TextAt{13}{232.844}{641.57}{\Character{125}}%
\TextAt{17}{240.297}{639.777}{\Character{0}}%
\TextAt{9}{246.084}{641.57}{\Character{40}}%
\TextAt{13}{250.636}{641.57}{z}%
\TextAt{9}{256.607}{641.57}{\Character{41}}%
\TextAt{9}{264.462}{641.57}{\Character{61}}%
\TextAt{9}{280.526}{649.658}{1}%
\Vector{q 0 G 0 g 0.478 w 206.101 -129.441 m 216.804 -129.441 l S Q}%
\TextAt{13}{278.101}{633.369}{z}%
\TextAt{17}{284.071}{636.823}{2}%
\TextAt{9}{292.656}{641.57}{\Character{43}}%
\TextAt{20}{305.368}{652.927}{\Character{88}}%
\TextAt{21}{304.417}{627.412}{\Character{33}}%
\NativeMathAt{309.969}{627.412}{7.97}{\in}%
\TextAt{17}{315.614}{627.412}{\Character{0}}%
\TextAt{23}{320.903}{629.681}{\Character{48}}%
\TextAt{20}{325.58}{658.427}{\Character{20}}%
\TextAt{9}{351.227}{649.658}{1}%
\Vector{q 0 G 0 g 0.478 w 261.085 -129.441 m 303.223 -129.441 l S Q}%
\TextAt{9}{333.085}{633.369}{\Character{40}}%
\TextAt{13}{337.637}{633.369}{z}%
\TextAt{16}{346.265}{633.369}{\Character{0}}%
\TextAt{13}{358.22}{633.369}{\Character{33}}%
\TextAt{9}{365.938}{633.369}{\Character{41}}%
\TextAt{17}{370.49}{636.823}{2}%
\TextAt{16}{379.075}{641.57}{\Character{0}}%
\TextAt{9}{395.524}{649.658}{1}%
\Vector{q 0 G 0 g 0.478 w 320.225 -129.441 m 332.676 -129.441 l S Q}%
\TextAt{13}{392.225}{633.369}{\Character{33}}%
\TextAt{17}{399.944}{636.823}{2}%
\TextAt{20}{405.871}{658.427}{\Character{21}}%
\TextAt{13}{414.174}{641.57}{\Character{59}}%
\TextAt{8}{79.226}{605.071}{where}%
\TextAt{9}{113.035}{605.071}{\Character{0}}%
\TextAt{3}{120.35}{609.41}{\Character{48}}%
\TextAt{9}{126.465}{605.071}{\Character{61}}%
\TextAt{9}{138.875}{605.071}{\Character{0}}%
\TextAt{16}{148.862}{605.071}{\Character{110}}%
\TextAt{16}{157.494}{605.071}{\Character{102}}%
\TextAt{9}{163.474}{605.071}{0}%
\TextAt{16}{169.327}{605.071}{\Character{103}}%
\TextAt{8}{175.304}{605.071}{.}%
\TextAt{13}{183.753}{605.071}{\Character{125}}%
\TextAt{9}{191.205}{605.071}{\Character{40}}%
\TextAt{13}{195.757}{605.071}{z}%
\TextAt{9}{201.728}{605.071}{\Character{41}}%
\TextAt{8}{210.178}{605.071}{is}%
\TextAt{8}{221.942}{605.071}{a}%
\TextAt{8}{231.697}{605.071}{meromorphic}%
\TextAt{8}{302.591}{605.071}{doubly}%
\TextAt{8}{341.29}{605.071}{periodic}%
\TextAt{8}{386.158}{605.071}{function,}%
\TextAt{8}{435.258}{605.071}{i.e.,}%
\TextAt{13}{225.031}{577.346}{\Character{125}}%
\TextAt{9}{232.484}{577.346}{\Character{40}}%
\TextAt{13}{237.036}{577.346}{z}%
\TextAt{9}{245.663}{577.346}{\Character{43}}%
\TextAt{13}{257.425}{577.346}{\Character{33}}%
\TextAt{9}{265.143}{577.346}{\Character{41}}%
\TextAt{9}{272.998}{577.346}{\Character{61}}%
\TextAt{13}{285.442}{577.346}{\Character{125}}%
\TextAt{9}{292.894}{577.346}{\Character{40}}%
\TextAt{13}{297.446}{577.346}{z}%
\TextAt{9}{303.417}{577.346}{\Character{41}}%
\TextAt{8}{319.925}{577.346}{for}%
\TextAt{8}{337.809}{577.346}{all}%
\TextAt{13}{354.069}{577.346}{\Character{33}}%
\NativeMathAt{365.108}{577.346}{11.955}{\in}%
\TextAt{9}{376.399}{577.346}{\Character{0}}%
\TextAt{13}{383.713}{577.346}{\Character{59}}%
\TextAt{8}{79.226}{549.621}{and}%
\TextAt{8}{101.989}{549.621}{it}%
\TextAt{8}{113.693}{549.621}{satisfies,}%
\TextAt{8}{160.055}{549.621}{along}%
\TextAt{8}{191.282}{549.621}{with}%
\TextAt{8}{217.942}{549.621}{its}%
\TextAt{8}{234.261}{549.621}{derivative}%
\TextAt{13}{287.928}{549.621}{\Character{125}}%
\TextAt{3}{295.38}{553.959}{\Character{48}}%
\TextAt{9}{298.175}{549.621}{\Character{40}}%
\TextAt{13}{302.728}{549.621}{z}%
\TextAt{9}{308.698}{549.621}{\Character{41}}%
\TextAt{8}{313.25}{549.621}{,}%
\TextAt{8}{320.4}{549.621}{an}%
\TextAt{8}{336.659}{549.621}{identity}%
\TextAt{13}{228.701}{521.896}{\Character{125}}%
\TextAt{3}{236.154}{526.832}{\Character{48}}%
\TextAt{9}{238.949}{521.896}{\Character{40}}%
\TextAt{13}{243.501}{521.896}{z}%
\TextAt{9}{249.472}{521.896}{\Character{41}}%
\TextAt{17}{254.024}{526.832}{2}%
\TextAt{9}{262.077}{521.896}{\Character{61}}%
\TextAt{9}{274.487}{521.896}{4}%
\TextAt{13}{280.356}{521.896}{\Character{125}}%
\TextAt{9}{287.808}{521.896}{\Character{40}}%
\TextAt{13}{292.361}{521.896}{z}%
\TextAt{9}{298.331}{521.896}{\Character{41}}%
\TextAt{17}{302.884}{526.832}{3}%
\TextAt{16}{310.273}{521.896}{\Character{0}}%
\TextAt{13}{322.228}{521.896}{g}%
\TextAt{17}{327.833}{520.102}{2}%
\TextAt{13}{332.565}{521.896}{\Character{125}}%
\TextAt{9}{340.018}{521.896}{\Character{40}}%
\TextAt{13}{344.57}{521.896}{z}%
\TextAt{9}{350.541}{521.896}{\Character{41}}%
\TextAt{16}{357.75}{521.896}{\Character{0}}%
\TextAt{13}{369.705}{521.896}{g}%
\TextAt{17}{375.31}{520.102}{3}%
\TextAt{13}{380.043}{521.896}{\Character{58}}%
\TextAt{8}{79.226}{494.17}{The}%
\TextAt{8}{105.485}{494.17}{constants}%
\TextAt{13}{159.528}{494.17}{g}%
\TextAt{17}{165.133}{492.377}{2}%
\TextAt{13}{169.865}{494.17}{\Character{59}}%
\TextAt{13}{175.089}{494.17}{g}%
\TextAt{17}{180.715}{492.377}{3}%
\TextAt{8}{191.55}{494.17}{only}%
\TextAt{8}{219.447}{494.17}{depend}%
\TextAt{8}{262.287}{494.17}{on}%
\TextAt{8}{280.751}{494.17}{the}%
\TextAt{8}{303.113}{494.17}{lattice}%
\TextAt{9}{341.088}{494.17}{\Character{0}}%
\TextAt{8}{348.402}{494.17}{,}%
\TextAt{8}{358.308}{494.17}{and}%
\TextAt{8}{383.275}{494.17}{are}%
\TextAt{8}{404.991}{494.17}{given}%
\TextAt{8}{437.766}{494.17}{by}%
\TextAt{13}{456.23}{494.17}{g}%
\TextAt{17}{461.835}{492.377}{2}%
\TextAt{9}{473.646}{494.17}{\Character{61}}%
\TextAt{13}{489.83}{494.17}{g}%
\TextAt{17}{495.435}{492.377}{2}%
\TextAt{9}{500.167}{494.17}{\Character{40}\Character{0}\Character{41}}%
\TextAt{9}{523.635}{494.17}{\Character{61}}%
\TextAt{9}{79.226}{477.124}{60}%
\TextAt{20}{92.925}{486.091}{\Character{80}}%
\TextAt{21}{105.544}{473.637}{\Character{33}}%
\NativeMathAt{111.095}{473.637}{7.97}{\neq}%
\TextAt{17}{117.678}{473.637}{0}%
\TextAt{17}{128.336}{481.831}{1}%
\Vector{q 0 G 0 g 0.478 w 53.602 -293.887 m 63.305 -293.887 l S Q}%
\TextAt{21}{125.602}{473.002}{\Character{33}}%
\TextAt{24}{131.154}{475.271}{4}%
\TextAt{8}{136.5}{477.124}{,}%
\TextAt{13}{143.358}{477.124}{g}%
\TextAt{17}{148.963}{475.331}{3}%
\TextAt{9}{157.016}{477.124}{\Character{61}}%
\TextAt{13}{169.442}{477.124}{g}%
\TextAt{17}{175.047}{475.331}{3}%
\TextAt{9}{179.779}{477.124}{\Character{40}\Character{0}\Character{41}}%
\TextAt{9}{199.481}{477.124}{\Character{61}}%
\TextAt{9}{211.903}{477.124}{140}%
\TextAt{20}{231.496}{486.091}{\Character{80}}%
\TextAt{21}{244.115}{473.637}{\Character{33}}%
\NativeMathAt{249.667}{473.637}{7.97}{\neq}%
\TextAt{17}{256.25}{473.637}{0}%
\TextAt{17}{266.908}{481.831}{1}%
\Vector{q 0 G 0 g 0.478 w 192.173 -293.887 m 201.876 -293.887 l S Q}%
\TextAt{21}{264.173}{473.002}{\Character{33}}%
\TextAt{24}{269.725}{475.271}{6}%
\TextAt{8}{275.072}{477.124}{.}%
\TextAt{13}{283.398}{477.124}{\Character{125}}%
\TextAt{8}{294.383}{477.124}{and}%
\TextAt{13}{316.781}{477.124}{\Character{125}}%
\TextAt{3}{324.234}{481.463}{\Character{48}}%
\TextAt{8}{330.561}{477.124}{may}%
\TextAt{8}{355.565}{477.124}{thus}%
\TextAt{8}{380.952}{477.124}{be}%
\TextAt{8}{396.511}{477.124}{interpreted}%
\TextAt{8}{456.305}{477.124}{as}%
\TextAt{8}{470.31}{477.124}{functions}%
\TextAt{8}{520.408}{477.124}{on}%
\TextAt{18}{79.226}{460.078}{C}%
\TextAt{9}{87.86}{460.078}{\Character{47}\Character{0}}%
\TextAt{8}{101.028}{460.078}{.}%
\TextAt{8}{109.476}{460.078}{If}%
\TextAt{8}{121.168}{460.078}{one}%
\TextAt{8}{142.628}{460.078}{takes}%
\TextAt{8}{172.611}{460.078}{away}%
\TextAt{8}{201.889}{460.078}{the}%
\TextAt{8}{222.046}{460.078}{finite}%
\TextAt{8}{251.958}{460.078}{set}%
\TextAt{13}{270.225}{460.078}{S}%
\NativeMathAt{281.441}{460.078}{11.955}{\subseteq}%
\TextAt{18}{294.061}{460.078}{C}%
\TextAt{9}{302.695}{460.078}{\Character{47}\Character{0}}%
\TextAt{8}{319.76}{460.078}{of}%
\TextAt{8}{333.09}{460.078}{poles,}%
\TextAt{8}{365.99}{460.078}{one}%
\TextAt{8}{387.45}{460.078}{gets}%
\TextAt{8}{411.575}{460.078}{a}%
\TextAt{8}{421.331}{460.078}{bijection}%
\TextAt{18}{123.638}{432.353}{C}%
\TextAt{9}{132.273}{432.353}{\Character{47}\Character{0}}%
\TextAt{16}{148.097}{432.353}{\Character{110}}%
\TextAt{13}{156.731}{432.353}{S}%
\NativeMathAt{170.535}{439.13}{7.97}{\sim}%
\NativeArrowAt{167.947}{432.353}{11.955}{13.175}{\longrightarrow}%
\TextAt{16}{184.421}{432.353}{\Character{102}}%
\TextAt{9}{190.418}{432.353}{\Character{40}}%
\TextAt{13}{194.971}{432.353}{x\Character{59}}%
\TextAt{13}{206.843}{432.353}{y}%
\TextAt{9}{213.004}{432.353}{\Character{41}}%
\NativeMathAt{220.877}{432.353}{11.955}{\in}%
\TextAt{18}{232.168}{432.353}{C}%
\TextAt{17}{240.802}{437.289}{2}%
\TextAt{16}{248.855}{432.353}{\Character{106}}%
\TextAt{13}{255.497}{432.353}{y}%
\TextAt{17}{261.634}{437.289}{2}%
\TextAt{9}{269.687}{432.353}{\Character{61}}%
\TextAt{9}{282.096}{432.353}{4}%
\TextAt{13}{287.965}{432.353}{x}%
\TextAt{17}{294.617}{437.289}{3}%
\TextAt{16}{302.006}{432.353}{\Character{0}}%
\TextAt{13}{313.961}{432.353}{g}%
\TextAt{17}{319.567}{430.56}{2}%
\TextAt{13}{324.299}{432.353}{x}%
\TextAt{16}{333.608}{432.353}{\Character{0}}%
\TextAt{13}{345.563}{432.353}{g}%
\TextAt{17}{351.168}{430.56}{3}%
\TextAt{16}{355.901}{432.353}{\Character{103}}%
\TextAt{13}{361.878}{432.353}{\Character{59}}%
\TextAt{13}{391.013}{432.353}{z}%
\NativeArrowAt{400.324}{432.353}{11.955}{19.26}{\longmapsto}%
\TextAt{9}{422.906}{432.353}{\Character{40}}%
\TextAt{13}{427.458}{432.353}{\Character{125}}%
\TextAt{9}{434.911}{432.353}{\Character{40}}%
\TextAt{13}{439.463}{432.353}{z}%
\TextAt{9}{445.434}{432.353}{\Character{41}}%
\TextAt{13}{449.986}{432.353}{\Character{59}}%
\TextAt{13}{455.21}{432.353}{\Character{125}}%
\TextAt{3}{462.683}{437.289}{\Character{48}}%
\TextAt{9}{465.478}{432.353}{\Character{40}}%
\TextAt{13}{470.03}{432.353}{z}%
\TextAt{9}{476.001}{432.353}{\Character{41}\Character{41}}%
\TextAt{13}{485.106}{432.353}{\Character{59}}%
\TextAt{8}{79.226}{404.628}{onto}%
\TextAt{8}{105.044}{404.628}{the}%
\TextAt{8}{124.67}{404.628}{affine}%
\TextAt{8}{155.354}{404.628}{algebraic}%
\TextAt{8}{204.257}{404.628}{curve}%
\TextAt{8}{234.941}{404.628}{in}%
\TextAt{18}{248.063}{404.628}{C}%
\TextAt{17}{256.698}{408.966}{2}%
\TextAt{8}{264.796}{404.628}{given}%
\TextAt{8}{294.835}{404.628}{by}%
\TextAt{8}{310.563}{404.628}{the}%
\TextAt{8}{330.188}{404.628}{equation}%
\TextAt{13}{377.466}{404.628}{y}%
\TextAt{17}{383.603}{408.966}{2}%
\TextAt{9}{391.656}{404.628}{\Character{61}}%
\TextAt{9}{404.065}{404.628}{4}%
\TextAt{13}{409.934}{404.628}{x}%
\TextAt{17}{416.586}{408.966}{3}%
\TextAt{16}{422.889}{404.628}{\Character{0}}%
\TextAt{13}{433.759}{404.628}{g}%
\TextAt{17}{439.364}{402.835}{2}%
\TextAt{13}{444.096}{404.628}{x}%
\TextAt{16}{452.319}{404.628}{\Character{0}}%
\TextAt{13}{463.189}{404.628}{g}%
\TextAt{17}{468.794}{402.835}{3}%
\TextAt{8}{473.526}{404.628}{.}%
\TextAt{8}{481.798}{404.628}{This}%
\TextAt{8}{507.986}{404.628}{gives}%
\TextAt{8}{79.226}{387.582}{the}%
\TextAt{8}{98.961}{387.582}{torus}%
\TextAt{18}{128.522}{387.582}{C}%
\TextAt{9}{137.156}{387.582}{\Character{47}\Character{0}}%
\TextAt{8}{153.799}{387.582}{the}%
\TextAt{8}{173.534}{387.582}{structure}%
\TextAt{8}{223.252}{387.582}{of}%
\TextAt{8}{236.16}{387.582}{an}%
\TextAt{8}{251.997}{387.582}{algebraic}%
\TextAt{8}{301.009}{387.582}{curve}%
\TextAt{13}{331.802}{387.582}{E}%
\TextAt{8}{344.644}{387.582}{over}%
\TextAt{18}{369.268}{387.582}{C}%
\TextAt{8}{381.378}{387.582}{of}%
\TextAt{8}{394.286}{387.582}{genus}%
\TextAt{9}{426.119}{387.582}{1}%
\TextAt{8}{431.972}{387.582}{.}%
\TextAt{8}{440.279}{387.582}{An}%
\TextAt{8}{459.034}{387.582}{important}%
\TextAt{8}{513.904}{387.582}{role}%
\TextAt{8}{79.226}{370.536}{is}%
\TextAt{8}{90.99}{370.536}{played}%
\TextAt{8}{127.74}{370.536}{by}%
\TextAt{8}{143.999}{370.536}{the}%
\TextAt{13}{164.156}{370.536}{j}%
\TextAt{12}{169.671}{370.536}{-invariant}%
\TextAt{13}{182.511}{338.859}{j}%
\TextAt{9}{188.027}{338.859}{\Character{40}}%
\TextAt{13}{192.579}{338.859}{E}%
\TextAt{9}{201.945}{338.859}{\Character{41}}%
\TextAt{9}{209.8}{338.859}{\Character{61}}%
\TextAt{13}{222.244}{338.859}{j}%
\TextAt{9}{227.759}{338.859}{\Character{40}\Character{0}\Character{41}}%
\TextAt{9}{247.461}{338.859}{\Character{61}}%
\TextAt{9}{261.12}{346.947}{2}%
\TextAt{17}{266.973}{351.285}{6}%
\TextAt{9}{271.705}{346.947}{3}%
\TextAt{17}{277.558}{351.285}{3}%
\TextAt{13}{282.29}{346.947}{g}%
\TextAt{17}{288.325}{351.285}{3}%
\TextAt{17}{287.896}{343.991}{2}%
\Vector{q 0 G 0 g 0.478 w 189.12 -432.152 m 221.057 -432.152 l S Q}%
\TextAt{9}{272.211}{330.658}{\Character{1}}%
\TextAt{8}{318.163}{338.859}{with}%
\TextAt{9}{352.881}{338.859}{\Character{1}}%
\TextAt{9}{365.936}{338.859}{\Character{61}}%
\TextAt{13}{378.382}{338.859}{g}%
\TextAt{17}{384.416}{343.795}{3}%
\TextAt{17}{383.987}{335.904}{2}%
\TextAt{16}{391.805}{338.859}{\Character{0}}%
\TextAt{9}{403.76}{338.859}{27}%
\TextAt{13}{415.466}{338.859}{g}%
\TextAt{17}{421.501}{343.795}{2}%
\TextAt{17}{421.072}{335.904}{3}%
\TextAt{13}{426.233}{338.859}{\Character{58}}%
\TextAt{8}{79.226}{311.134}{It}%
\TextAt{8}{92.361}{311.134}{determines}%
\TextAt{8}{152.061}{311.134}{the}%
\TextAt{8}{172.679}{311.134}{elliptic}%
\TextAt{8}{211.505}{311.134}{curve}%
\TextAt{13}{243.182}{311.134}{E}%
\TextAt{8}{256.908}{311.134}{up}%
\TextAt{8}{274.274}{311.134}{to}%
\TextAt{8}{289.047}{311.134}{isomorphism.}%
\TextAt{8}{363.403}{311.134}{Writing}%
\TextAt{8}{406.784}{311.134}{generators}%
\TextAt{13}{463.902}{311.134}{\Character{33}}%
\TextAt{17}{471.191}{309.341}{1}%
\TextAt{13}{475.923}{311.134}{\Character{59}}%
\TextAt{13}{481.147}{311.134}{\Character{33}}%
\TextAt{17}{488.457}{309.341}{2}%
\TextAt{8}{497.548}{311.134}{of}%
\TextAt{9}{511.341}{311.134}{\Character{0}}%
\TextAt{8}{523.014}{311.134}{in}%
\TextAt{8}{79.226}{294.088}{such}%
\TextAt{8}{105.159}{294.088}{an}%
\TextAt{8}{120.953}{294.088}{order}%
\TextAt{8}{151.059}{294.088}{that}%
\TextAt{13}{175.963}{294.088}{\Character{28}}%
\TextAt{9}{185.703}{294.088}{\Character{61}}%
\TextAt{13}{198.128}{294.088}{\Character{33}}%
\TextAt{17}{205.418}{292.295}{1}%
\TextAt{9}{210.15}{294.088}{\Character{47}}%
\TextAt{13}{216.003}{294.088}{\Character{33}}%
\TextAt{17}{223.292}{292.295}{2}%
\TextAt{8}{231.457}{294.088}{lies}%
\TextAt{8}{251.209}{294.088}{in}%
\TextAt{8}{264.398}{294.088}{the}%
\TextAt{8}{284.09}{294.088}{upper}%
\TextAt{8}{317.112}{294.088}{half-plane}%
\TextAt{18}{370.948}{294.088}{H}%
\TextAt{8}{380.246}{294.088}{,}%
\TextAt{8}{387.024}{294.088}{then}%
\TextAt{13}{413.22}{294.088}{j}%
\TextAt{9}{418.735}{294.088}{\Character{40}}%
\TextAt{13}{423.288}{294.088}{E}%
\TextAt{9}{432.654}{294.088}{\Character{41}}%
\TextAt{8}{440.639}{294.088}{becomes}%
\TextAt{8}{486.728}{294.088}{the}%
\TextAt{8}{506.42}{294.088}{value}%
\TextAt{13}{79.226}{277.042}{j}%
\TextAt{9}{84.742}{277.042}{\Character{40}}%
\TextAt{13}{89.294}{277.042}{\Character{28}}%
\TextAt{9}{95.713}{277.042}{\Character{41}}%
\TextAt{8}{104.754}{277.042}{of}%
\TextAt{8}{118.675}{277.042}{a}%
\TextAt{12}{129.022}{277.042}{modular}%
\TextAt{12}{183.225}{277.042}{function}%
\TextAt{8}{231.441}{277.042}{,}%
\TextAt{8}{239.329}{277.042}{i.e.,}%
\TextAt{8}{262.173}{277.042}{of}%
\TextAt{8}{276.094}{277.042}{a}%
\TextAt{8}{286.441}{277.042}{holomorphic}%
\TextAt{8}{354.029}{277.042}{function}%
\TextAt{13}{400.469}{277.042}{j}%
\TextAt{8}{410.473}{277.042}{on}%
\TextAt{18}{427.323}{277.042}{H}%
\TextAt{8}{441.11}{277.042}{which}%
\TextAt{8}{475.188}{277.042}{is}%
\TextAt{8}{487.543}{277.042}{invariant}%
\TextAt{8}{79.226}{259.996}{under}%
\TextAt{8}{112.39}{259.996}{the}%
\TextAt{8}{132.546}{259.996}{substitution}%
\TextAt{13}{173.948}{227.279}{\Character{28}}%
\NativeArrowAt{183.688}{227.279}{11.955}{19.26}{\longmapsto}%
\TextAt{13}{207.465}{235.367}{a\Character{28}}%
\TextAt{9}{222.686}{235.367}{\Character{43}}%
\TextAt{13}{234.447}{235.367}{b}%
\Vector{q 0 G 0 g 0.478 w 135.465 -543.732 m 167.424 -543.732 l S Q}%
\TextAt{13}{207.466}{219.079}{c\Character{28}}%
\TextAt{9}{221.58}{219.079}{\Character{43}}%
\TextAt{13}{233.341}{219.079}{d}%
\TextAt{8}{252.575}{227.279}{for}%
\TextAt{8}{270.46}{227.279}{every}%
\TextAt{8}{301.352}{227.279}{matrix}%
\TextAt{20}{341.398}{247.723}{\Character{32}}%
\TextAt{13}{350.862}{235.382}{a}%
\TextAt{13}{367.516}{235.382}{b}%
\TextAt{13}{351.416}{218.336}{c}%
\TextAt{13}{366.958}{218.336}{d}%
\TextAt{20}{373.053}{247.723}{\Character{33}}%
\NativeMathAt{385.838}{227.279}{11.955}{\in}%
\TextAt{13}{397.129}{227.279}{S}%
\TextAt{13}{405.019}{227.279}{L}%
\TextAt{17}{412.989}{225.486}{2}%
\TextAt{9}{417.721}{227.279}{\Character{40}}%
\TextAt{18}{422.273}{227.279}{Z}%
\TextAt{9}{430.244}{227.279}{\Character{41}}%
\TextAt{13}{434.796}{227.279}{\Character{58}}%
\TextAt{8}{79.226}{192.243}{Now}%
\TextAt{8}{106.377}{192.243}{let}%
\TextAt{13}{123.772}{192.243}{K}%
\NativeMathAt{138.747}{192.243}{11.955}{\subseteq}%
\TextAt{18}{152.203}{192.243}{C}%
\TextAt{8}{165.225}{192.243}{be}%
\TextAt{8}{181.64}{192.243}{an}%
\TextAt{8}{198.39}{192.243}{imaginary}%
\TextAt{8}{253.851}{192.243}{quadratic}%
\TextAt{8}{306.705}{192.243}{number}%
\TextAt{8}{349.792}{192.243}{field.}%
\TextAt{8}{381.172}{192.243}{Then}%
\TextAt{8}{412.22}{192.243}{the}%
\TextAt{8}{432.868}{192.243}{ring}%
\TextAt{16}{457.424}{192.243}{\Character{79}}%
\TextAt{21}{466.942}{190.45}{K}%
\TextAt{8}{479.533}{192.243}{of}%
\TextAt{8}{493.353}{192.243}{integers}%
\TextAt{8}{79.226}{175.197}{forms}%
\TextAt{8}{112.05}{175.197}{a}%
\TextAt{8}{122.375}{175.197}{lattice}%
\TextAt{8}{158.714}{175.197}{in}%
\TextAt{18}{172.936}{175.197}{C}%
\TextAt{8}{181.571}{175.197}{,}%
\TextAt{8}{189.432}{175.197}{and}%
\TextAt{8}{212.763}{175.197}{more}%
\TextAt{8}{242.599}{175.197}{generally,}%
\TextAt{8}{295.339}{175.197}{any}%
\TextAt{8}{318.026}{175.197}{ideal}%
\TextAt{14}{346.558}{175.197}{a}%
\TextAt{8}{357.008}{175.197}{of}%
\TextAt{16}{370.907}{175.197}{\Character{79}}%
\TextAt{21}{380.425}{173.404}{K}%
\TextAt{8}{393.094}{175.197}{does}%
\TextAt{8}{420.061}{175.197}{as}%
\TextAt{8}{435}{175.197}{well.}%
\TextAt{8}{464.99}{175.197}{The}%
\TextAt{8}{489.613}{175.197}{tori}%
\TextAt{18}{512.299}{175.197}{C}%
\TextAt{9}{520.933}{175.197}{\Character{47}}%
\TextAt{14}{526.786}{175.197}{a}%
\TextAt{8}{79.226}{158.151}{constructed}%
\TextAt{8}{142.99}{158.151}{in}%
\TextAt{8}{157.26}{158.151}{this}%
\TextAt{8}{180.7}{158.151}{way}%
\TextAt{8}{205.06}{158.151}{are}%
\TextAt{8}{225.187}{158.151}{elliptic}%
\TextAt{8}{264.169}{158.151}{curves}%
\TextAt{8}{300.615}{158.151}{with}%
\TextAt{19}{327.892}{158.151}{complex}%
\TextAt{19}{372.789}{158.151}{multiplication}%
\TextAt{8}{442.428}{158.151}{.}%
\TextAt{8}{452.727}{158.151}{This}%
\TextAt{8}{480.064}{158.151}{means}%
\TextAt{8}{516.51}{158.151}{the}%
\TextAt{8}{79.226}{141.105}{following.}%
\TextAt{8}{134.496}{141.105}{An}%
\TextAt{8}{154.1}{141.105}{endomorphism}%
\TextAt{8}{233.288}{141.105}{of}%
\TextAt{8}{247.046}{141.105}{an}%
\TextAt{8}{263.733}{141.105}{elliptic}%
\TextAt{8}{302.525}{141.105}{curve}%
\TextAt{13}{334.168}{141.105}{E}%
\TextAt{9}{347.585}{141.105}{\Character{61}}%
\TextAt{18}{360.739}{141.105}{C}%
\TextAt{9}{369.374}{141.105}{\Character{47}\Character{0}}%
\TextAt{8}{386.866}{141.105}{is}%
\TextAt{8}{399.058}{141.105}{given}%
\TextAt{8}{430.056}{141.105}{as}%
\TextAt{8}{444.854}{141.105}{multiplication}%
\TextAt{8}{520.408}{141.105}{by}%
\TextAt{8}{79.226}{124.059}{a}%
\TextAt{8}{89.629}{124.059}{complex}%
\TextAt{8}{136.125}{124.059}{number}%
\TextAt{13}{179.369}{124.059}{z}%
\TextAt{8}{189.885}{124.059}{such}%
\TextAt{8}{216.93}{124.059}{that}%
\TextAt{13}{242.946}{124.059}{z}%
\TextAt{9}{248.917}{124.059}{\Character{0}}%
\NativeMathAt{260.656}{124.059}{11.955}{\subseteq}%
\TextAt{9}{274.379}{124.059}{\Character{0}}%
\TextAt{8}{281.693}{124.059}{.}%
\TextAt{8}{292.085}{124.059}{Generally,}%
\TextAt{8}{348.247}{124.059}{one}%
\TextAt{8}{370.354}{124.059}{has}%
\TextAt{8}{391.875}{124.059}{End}%
\TextAt{9}{412.845}{124.059}{\Character{40}}%
\TextAt{13}{417.397}{124.059}{E}%
\TextAt{9}{426.763}{124.059}{\Character{41}}%
\TextAt{9}{435.729}{124.059}{\Character{61}}%
\TextAt{18}{449.27}{124.059}{Z}%
\TextAt{8}{457.24}{124.059}{.}%
\TextAt{8}{467.631}{124.059}{If}%
\TextAt{8}{479.971}{124.059}{this}%
\TextAt{8}{503.441}{124.059}{is}%
\TextAt{8}{515.853}{124.059}{not}%
\TextAt{8}{79.226}{107.013}{the}%
\TextAt{8}{100.319}{107.013}{case,}%
\TextAt{8}{129.513}{107.013}{then}%
\TextAt{8}{157.11}{107.013}{End}%
\TextAt{9}{178.079}{107.013}{\Character{40}}%
\TextAt{13}{182.631}{107.013}{E}%
\TextAt{9}{191.998}{107.013}{\Character{41}}%
\TextAt{16}{199.845}{107.013}{\Character{10}}%
\TextAt{18}{212.439}{107.013}{Q}%
\TextAt{8}{226.571}{107.013}{is}%
\TextAt{8}{239.272}{107.013}{necessarily}%
\TextAt{8}{298.538}{107.013}{an}%
\TextAt{8}{315.733}{107.013}{imaginary}%
\TextAt{8}{371.64}{107.013}{quadratic}%
\TextAt{8}{424.94}{107.013}{number}%
\TextAt{8}{468.473}{107.013}{field}%
\TextAt{13}{494.767}{107.013}{K}%
\TextAt{8}{505.584}{107.013}{,}%
\TextAt{8}{513.904}{107.013}{and}%
\TextAt{8}{79.226}{89.967}{one}%
\TextAt{8}{101.273}{89.967}{says}%
\TextAt{8}{126.703}{89.967}{that}%
\TextAt{8}{152.659}{89.967}{this}%
\TextAt{8}{176.069}{89.967}{is}%
\TextAt{8}{188.42}{89.967}{an}%
\TextAt{8}{205.267}{89.967}{elliptic}%
\TextAt{8}{244.218}{89.967}{curve}%
\TextAt{8}{276.021}{89.967}{with}%
\TextAt{8}{303.268}{89.967}{complex}%
\TextAt{8}{349.703}{89.967}{multiplication.}%
\TextAt{8}{431.143}{89.967}{The}%
\TextAt{8}{455.784}{89.967}{curves}%
\TextAt{18}{492.201}{89.967}{C}%
\TextAt{9}{500.835}{89.967}{\Character{47}}%
\TextAt{14}{506.688}{89.967}{a}%
\TextAt{8}{517.156}{89.967}{are}%
\TextAt{8}{79.226}{72.921}{obviously}%
\TextAt{8}{131.004}{72.921}{of}%
\TextAt{8}{144.334}{72.921}{this}%
\TextAt{8}{167.156}{72.921}{kind.}%
\TextAt{8}{297.211}{43.033}{366}%
\Destination{78.23}{809.815}{nav.page.366}%
\Destination{79.36}{704.955}{nav.print.402}%
\Hotlink{93.036}{724.855}{19.22}{11.047}{nav.cite.N096}%
\Hotlink{137.916}{724.855}{19.22}{11.047}{nav.cite.N028}%
}

% Page 379 (body page 367)
\Page{%
\TextAt{2}{79.226}{796.943}{VI.}%
\TextAt{2}{98.012}{796.943}{Global}%
\TextAt{2}{133.248}{796.943}{Class}%
\TextAt{2}{161.841}{796.943}{Field}%
\TextAt{2}{189.572}{796.943}{Theory}%
\Vector{q 0 G 0 g 0.299 w 7.226 18.729 m 460.769 18.729 l S Q}%
\TextAt{8}{97.159}{761.124}{The}%
\TextAt{8}{121.112}{761.124}{consequences}%
\TextAt{8}{191.69}{761.124}{of}%
\TextAt{8}{204.92}{761.124}{these}%
\TextAt{8}{234.791}{761.124}{analytic}%
\TextAt{8}{279.247}{761.124}{investigations}%
\TextAt{8}{352.456}{761.124}{for}%
\TextAt{8}{370.24}{761.124}{class}%
\TextAt{8}{397.576}{761.124}{field}%
\TextAt{8}{422.833}{761.124}{theory}%
\TextAt{8}{459.482}{761.124}{are}%
\TextAt{8}{478.892}{761.124}{the}%
\TextAt{8}{498.948}{761.124}{follow-}%
\TextAt{8}{79.226}{744.078}{ing.}%
\TextAt{12}{79.226}{718.492}{Theorem}%
\TextAt{12}{136.253}{718.492}{(6.10).}%
\TextAt{8}{178.31}{718.492}{Let}%
\TextAt{13}{199.289}{718.492}{K}%
\TextAt{8}{214.015}{718.492}{be}%
\TextAt{8}{229.95}{718.492}{an}%
\TextAt{8}{246.22}{718.492}{imaginary}%
\TextAt{8}{301.2}{718.492}{quadratic}%
\TextAt{8}{353.574}{718.492}{number}%
\TextAt{8}{396.181}{718.492}{field}%
\TextAt{8}{421.549}{718.492}{and}%
\TextAt{14}{444.322}{718.492}{a}%
\TextAt{8}{454.213}{718.492}{an}%
\TextAt{8}{470.483}{718.492}{ideal}%
\TextAt{8}{498.457}{718.492}{of}%
\TextAt{16}{511.797}{718.492}{\Character{79}}%
\TextAt{21}{521.315}{716.699}{K}%
\TextAt{8}{529.518}{718.492}{.}%
\TextAt{8}{79.226}{701.446}{Then}%
\TextAt{8}{109.784}{701.446}{one}%
\TextAt{8}{131.243}{701.446}{has:}%
\TextAt{8}{89.579}{681.412}{(i)}%
\TextAt{8}{107.918}{681.412}{The}%
\TextAt{13}{131.901}{681.412}{j}%
\TextAt{8}{137.416}{681.412}{-invariant}%
\TextAt{13}{190.366}{681.412}{j}%
\TextAt{9}{195.881}{681.412}{\Character{40}}%
\TextAt{14}{200.433}{681.412}{a}%
\TextAt{9}{206.416}{681.412}{\Character{41}}%
\TextAt{8}{214.794}{681.412}{of}%
\TextAt{18}{228.053}{681.412}{C}%
\TextAt{9}{236.687}{681.412}{\Character{47}}%
\TextAt{14}{242.54}{681.412}{a}%
\TextAt{8}{252.349}{681.412}{is}%
\TextAt{8}{264.041}{681.412}{an}%
\TextAt{8}{280.228}{681.412}{algebraic}%
\TextAt{8}{329.591}{681.412}{integer}%
\TextAt{8}{368.218}{681.412}{which}%
\TextAt{8}{401.633}{681.412}{depends}%
\TextAt{8}{446.812}{681.412}{only}%
\TextAt{8}{472.432}{681.412}{on}%
\TextAt{8}{488.619}{681.412}{the}%
\TextAt{8}{508.704}{681.412}{ideal}%
\TextAt{8}{107.918}{664.366}{class}%
\TextAt{14}{135.356}{664.366}{K}%
\TextAt{8}{147.244}{664.366}{of}%
\TextAt{14}{160.574}{664.366}{a}%
\TextAt{8}{170.455}{664.366}{and}%
\TextAt{8}{193.218}{664.366}{will}%
\TextAt{8}{215.323}{664.366}{therefore}%
\TextAt{8}{264.423}{664.366}{be}%
\TextAt{8}{280.347}{664.366}{denoted}%
\TextAt{8}{324.569}{664.366}{by}%
\TextAt{13}{340.828}{664.366}{j}%
\TextAt{9}{346.344}{664.366}{\Character{40}}%
\TextAt{14}{350.896}{664.366}{K}%
\TextAt{9}{358.887}{664.366}{\Character{41}}%
\TextAt{8}{363.44}{664.366}{.}%
\TextAt{8}{86.327}{647.32}{(ii)}%
\TextAt{8}{107.918}{647.32}{Every}%
\TextAt{13}{141.572}{647.32}{j}%
\TextAt{9}{147.088}{647.32}{\Character{40}}%
\TextAt{14}{151.64}{647.32}{a}%
\TextAt{9}{157.623}{647.32}{\Character{41}}%
\TextAt{8}{166.073}{647.32}{generates}%
\TextAt{8}{217.516}{647.32}{the}%
\TextAt{8}{237.672}{647.32}{Hilbert}%
\TextAt{8}{277.985}{647.32}{class}%
\TextAt{8}{305.422}{647.32}{field}%
\TextAt{8}{330.779}{647.32}{over}%
\TextAt{13}{355.825}{647.32}{K}%
\TextAt{8}{366.643}{647.32}{.}%
\TextAt{8}{83.076}{630.274}{(iii)}%
\TextAt{8}{107.918}{630.274}{If}%
\TextAt{14}{119.103}{630.274}{a}%
\TextAt{17}{125.086}{628.48}{1}%
\TextAt{13}{129.819}{630.274}{\Character{59}}%
\TextAt{13}{135.043}{630.274}{\Character{58}}%
\TextAt{13}{140.28}{630.274}{\Character{58}}%
\TextAt{13}{145.504}{630.274}{\Character{58}}%
\TextAt{13}{150.741}{630.274}{\Character{59}}%
\TextAt{14}{156.039}{630.274}{a}%
\TextAt{21}{162.023}{628.48}{h}%
\TextAt{8}{170.791}{630.274}{are}%
\TextAt{8}{189.794}{630.274}{representatives}%
\TextAt{8}{269.088}{630.274}{of}%
\TextAt{8}{281.91}{630.274}{the}%
\TextAt{8}{301.559}{630.274}{ideal}%
\TextAt{8}{329.015}{630.274}{class}%
\TextAt{8}{355.945}{630.274}{group}%
\TextAt{13}{388.613}{630.274}{C}%
\TextAt{13}{397.831}{630.274}{l}%
\TextAt{21}{401.334}{628.48}{K}%
\TextAt{8}{409.536}{630.274}{,}%
\TextAt{8}{416.28}{630.274}{then}%
\TextAt{8}{442.432}{630.274}{the}%
\TextAt{8}{462.081}{630.274}{numbers}%
\TextAt{13}{508.785}{630.274}{j}%
\TextAt{9}{514.3}{630.274}{\Character{40}}%
\TextAt{14}{518.853}{630.274}{a}%
\TextAt{21}{524.836}{628.48}{i}%
\TextAt{9}{528.217}{630.274}{\Character{41}}%
\TextAt{8}{107.918}{613.228}{are}%
\TextAt{8}{127.429}{613.228}{conjugate}%
\TextAt{8}{180.439}{613.228}{to}%
\TextAt{8}{194.749}{613.228}{one}%
\TextAt{8}{216.208}{613.228}{another}%
\TextAt{8}{259.139}{613.228}{over}%
\TextAt{13}{284.185}{613.228}{K}%
\TextAt{8}{295.003}{613.228}{.}%
\TextAt{8}{83.398}{596.182}{(iv)}%
\TextAt{8}{107.918}{596.182}{For}%
\TextAt{8}{128.888}{596.182}{almost}%
\TextAt{8}{166.678}{596.182}{all}%
\TextAt{8}{182.937}{596.182}{prime}%
\TextAt{8}{216.101}{596.182}{ideals}%
\TextAt{14}{248.679}{596.182}{p}%
\TextAt{8}{258.559}{596.182}{of}%
\TextAt{13}{271.889}{596.182}{K}%
\TextAt{8}{286.604}{596.182}{one}%
\TextAt{8}{308.064}{596.182}{has}%
\TextAt{13}{275.513}{568.675}{\Character{39}}%
\TextAt{27}{283.182}{566.882}{p}%
\TextAt{13}{287.668}{568.675}{j}%
\TextAt{9}{293.184}{568.675}{\Character{40}}%
\TextAt{14}{297.736}{568.675}{a}%
\TextAt{9}{303.719}{568.675}{\Character{41}}%
\TextAt{9}{311.574}{568.675}{\Character{61}}%
\TextAt{13}{324.018}{568.675}{j}%
\TextAt{9}{329.533}{568.675}{\Character{40}}%
\TextAt{14}{334.086}{568.675}{p}%
\TextAt{3}{340.069}{573.611}{\Character{0}}%
\TextAt{17}{346.655}{573.611}{1}%
\TextAt{14}{351.388}{568.675}{a}%
\TextAt{9}{357.371}{568.675}{\Character{41}}%
\TextAt{13}{361.923}{568.675}{\Character{59}}%
\TextAt{8}{107.918}{541.168}{where}%
\TextAt{13}{142.291}{541.168}{\Character{39}}%
\TextAt{27}{149.959}{539.375}{p}%
\NativeMathAt{158.727}{541.168}{11.955}{\in}%
\TextAt{13}{170.978}{541.168}{G}%
\TextAt{9}{180.212}{541.168}{\Character{40}}%
\TextAt{13}{184.764}{541.168}{K}%
\TextAt{9}{195.582}{541.168}{\Character{40}}%
\TextAt{13}{200.134}{541.168}{j}%
\TextAt{9}{205.65}{541.168}{\Character{40}}%
\TextAt{14}{210.202}{541.168}{a}%
\TextAt{9}{216.185}{541.168}{\Character{41}\Character{41}}%
\TextAt{16}{229.571}{541.168}{\Character{106}}%
\TextAt{13}{237.173}{541.168}{K}%
\TextAt{9}{247.99}{541.168}{\Character{41}}%
\TextAt{8}{257.004}{541.168}{is}%
\TextAt{8}{269.331}{541.168}{the}%
\TextAt{8}{290.051}{541.168}{Frobenius}%
\TextAt{8}{344.484}{541.168}{automorphism}%
\TextAt{8}{422.517}{541.168}{of}%
\TextAt{8}{436.41}{541.168}{a}%
\TextAt{8}{446.729}{541.168}{prime}%
\TextAt{8}{480.456}{541.168}{ideal}%
\TextAt{14}{508.983}{541.168}{P}%
\TextAt{8}{523.337}{541.168}{of}%
\TextAt{13}{107.918}{524.122}{K}%
\TextAt{9}{118.736}{524.122}{\Character{40}}%
\TextAt{13}{123.288}{524.122}{j}%
\TextAt{9}{128.804}{524.122}{\Character{40}}%
\TextAt{14}{133.356}{524.122}{a}%
\TextAt{9}{139.339}{524.122}{\Character{41}\Character{41}}%
\TextAt{8}{152.341}{524.122}{above}%
\TextAt{14}{185.517}{524.122}{p}%
\TextAt{8}{191.5}{524.122}{.}%
\TextAt{8}{97.159}{498.537}{Note}%
\TextAt{8}{125.271}{498.537}{that}%
\TextAt{8}{150.466}{498.537}{for}%
\TextAt{8}{168.178}{498.537}{a}%
\TextAt{8}{177.76}{498.537}{totally}%
\TextAt{8}{214.994}{498.537}{imaginary}%
\TextAt{8}{269.79}{498.537}{field}%
\TextAt{13}{294.973}{498.537}{K}%
\TextAt{8}{309.515}{498.537}{there}%
\TextAt{8}{339.253}{498.537}{is}%
\TextAt{8}{350.843}{498.537}{no}%
\TextAt{8}{366.928}{498.537}{difference}%
\TextAt{8}{419.095}{498.537}{between}%
\TextAt{8}{464.111}{498.537}{big}%
\TextAt{8}{483.449}{498.537}{and}%
\TextAt{8}{506.038}{498.537}{small}%
\TextAt{8}{79.226}{481.491}{Hilbert}%
\TextAt{8}{120.203}{481.491}{class}%
\TextAt{8}{148.304}{481.491}{field.}%
\TextAt{8}{180.203}{481.491}{In}%
\TextAt{8}{195.488}{481.491}{order}%
\TextAt{8}{226.722}{481.491}{to}%
\TextAt{8}{241.696}{481.491}{go}%
\TextAt{8}{257.973}{481.491}{beyond}%
\TextAt{8}{299.285}{481.491}{the}%
\TextAt{8}{320.105}{481.491}{Hilbert}%
\TextAt{8}{361.082}{481.491}{class}%
\TextAt{8}{389.183}{481.491}{field,}%
\TextAt{8}{418.621}{481.491}{i.e.,}%
\TextAt{8}{441.556}{481.491}{the}%
\TextAt{8}{462.377}{481.491}{ray}%
\TextAt{8}{483.209}{481.491}{class}%
\TextAt{8}{511.31}{481.491}{field}%
\TextAt{8}{79.226}{464.445}{mod}%
\TextAt{9}{105.021}{464.445}{1}%
\TextAt{8}{110.874}{464.445}{,}%
\TextAt{8}{117.59}{464.445}{to}%
\TextAt{8}{131.358}{464.445}{the}%
\TextAt{8}{150.972}{464.445}{ray}%
\TextAt{8}{170.598}{464.445}{class}%
\TextAt{8}{197.493}{464.445}{fields}%
\TextAt{8}{226.923}{464.445}{for}%
\TextAt{8}{244.266}{464.445}{arbitrary}%
\TextAt{8}{293.493}{464.445}{modules}%
\TextAt{14}{338.858}{464.445}{m}%
\NativeMathAt{351.346}{464.445}{11.955}{\neq}%
\TextAt{9}{363.756}{464.445}{1}%
\TextAt{8}{369.624}{464.445}{,}%
\TextAt{8}{376.34}{464.445}{we}%
\TextAt{8}{393.025}{464.445}{form,}%
\TextAt{8}{423.484}{464.445}{for}%
\TextAt{8}{440.826}{464.445}{any}%
\TextAt{8}{462.401}{464.445}{lattice}%
\TextAt{9}{497.629}{464.445}{\Character{0}}%
\NativeMathAt{508.264}{464.445}{11.955}{\subseteq}%
\TextAt{18}{520.883}{464.445}{C}%
\TextAt{8}{529.518}{464.445}{,}%
\TextAt{8}{79.226}{447.399}{the}%
\TextAt{12}{99.383}{447.399}{Weber}%
\TextAt{12}{142.242}{447.399}{function}%
\TextAt{13}{202.587}{395.512}{\Character{28}}%
\TextAt{17}{207.671}{393.718}{\Character{0}}%
\TextAt{9}{213.458}{395.512}{\Character{40}}%
\TextAt{13}{218.01}{395.512}{z}%
\TextAt{9}{223.981}{395.512}{\Character{41}}%
\TextAt{9}{231.836}{395.512}{\Character{61}}%
\TextAt{20}{244.279}{434.366}{\Character{56}}%
\TextAt{20}{244.279}{423.607}{\Character{62}}%
\TextAt{20}{244.279}{420.02}{\Character{62}}%
\TextAt{20}{244.279}{416.433}{\Character{62}}%
\TextAt{20}{244.279}{412.847}{\Character{62}}%
\TextAt{20}{244.279}{409.26}{\Character{60}}%
\TextAt{20}{244.279}{387.741}{\Character{62}}%
\TextAt{20}{244.279}{384.154}{\Character{62}}%
\TextAt{20}{244.279}{380.568}{\Character{62}}%
\TextAt{20}{244.279}{376.981}{\Character{62}}%
\TextAt{20}{244.279}{373.394}{\Character{58}}%
\TextAt{16}{254.906}{421.504}{\Character{0}}%
\TextAt{9}{264.205}{421.504}{2}%
\TextAt{17}{270.058}{425.842}{7}%
\TextAt{9}{274.79}{421.504}{3}%
\TextAt{17}{280.643}{425.842}{5}%
\TextAt{13}{286.571}{429.592}{g}%
\TextAt{17}{292.176}{427.798}{2}%
\TextAt{13}{296.908}{429.591}{g}%
\TextAt{17}{302.514}{427.798}{3}%
\Vector{q 0 G 0 g 0.478 w 214.571 -349.507 m 235.246 -349.507 l S Q}%
\TextAt{9}{292.031}{413.303}{\Character{1}}%
\TextAt{13}{308.441}{421.504}{\Character{125}}%
\TextAt{17}{315.894}{419.711}{\Character{0}}%
\TextAt{9}{321.681}{421.504}{\Character{40}}%
\TextAt{13}{326.233}{421.504}{z}%
\TextAt{9}{332.204}{421.504}{\Character{41}}%
\TextAt{13}{336.756}{421.504}{\Character{59}}%
\TextAt{8}{351.963}{421.504}{if}%
\TextAt{13}{362.687}{421.504}{g}%
\TextAt{17}{368.292}{419.711}{2}%
\TextAt{13}{373.024}{421.504}{g}%
\TextAt{17}{378.63}{419.711}{3}%
\NativeMathAt{386.683}{421.504}{11.955}{\neq}%
\TextAt{9}{399.093}{421.504}{0}%
\TextAt{13}{404.961}{421.504}{\Character{59}}%
\TextAt{16}{254.906}{397.064}{\Character{0}}%
\TextAt{9}{264.205}{397.064}{2}%
\TextAt{17}{270.058}{401.402}{9}%
\TextAt{9}{274.79}{397.064}{3}%
\TextAt{17}{280.643}{401.402}{6}%
\TextAt{13}{286.571}{405.151}{g}%
\TextAt{17}{292.176}{403.358}{3}%
\Vector{q 0 G 0 g 0.478 w 214.571 -373.948 m 224.908 -373.948 l S Q}%
\TextAt{9}{286.862}{388.863}{\Character{1}}%
\TextAt{13}{298.104}{397.064}{\Character{125}}%
\TextAt{17}{305.556}{401.402}{3}%
\TextAt{17}{305.556}{394.043}{\Character{0}}%
\TextAt{9}{311.343}{397.064}{\Character{40}}%
\TextAt{13}{315.895}{397.064}{z}%
\TextAt{9}{321.866}{397.064}{\Character{41}}%
\TextAt{13}{326.418}{397.064}{\Character{59}}%
\TextAt{8}{351.963}{397.064}{if}%
\TextAt{13}{362.687}{397.064}{g}%
\TextAt{17}{368.292}{395.27}{2}%
\TextAt{9}{376.345}{397.064}{\Character{61}}%
\TextAt{9}{388.755}{397.064}{0}%
\TextAt{13}{394.624}{397.064}{\Character{59}}%
\TextAt{9}{254.906}{369.38}{2}%
\TextAt{17}{260.759}{373.718}{8}%
\TextAt{9}{265.491}{369.38}{3}%
\TextAt{17}{271.344}{373.718}{4}%
\TextAt{13}{277.272}{377.467}{g}%
\TextAt{17}{283.306}{381.806}{2}%
\TextAt{17}{282.878}{374.512}{2}%
\Vector{q 0 G 0 g 0.478 w 205.272 -401.632 m 216.039 -401.632 l S Q}%
\TextAt{9}{277.778}{361.179}{\Character{1}}%
\TextAt{13}{289.234}{369.379}{\Character{125}}%
\TextAt{17}{296.687}{373.718}{2}%
\TextAt{17}{296.687}{366.359}{\Character{0}}%
\TextAt{9}{302.474}{369.38}{\Character{40}}%
\TextAt{13}{307.026}{369.379}{z}%
\TextAt{9}{312.996}{369.38}{\Character{41}}%
\TextAt{13}{317.549}{369.379}{\Character{59}}%
\TextAt{8}{351.963}{369.38}{if}%
\TextAt{13}{362.687}{369.38}{g}%
\TextAt{17}{368.292}{367.586}{3}%
\TextAt{9}{376.345}{369.379}{\Character{61}}%
\TextAt{9}{388.755}{369.379}{0}%
\TextAt{13}{394.624}{369.379}{\Character{58}}%
\TextAt{8}{79.226}{341.305}{Let}%
\TextAt{14}{99.845}{341.305}{K}%
\NativeMathAt{111.158}{341.305}{11.955}{\in}%
\TextAt{13}{122.449}{341.305}{C}%
\TextAt{13}{131.667}{341.305}{l}%
\TextAt{21}{135.169}{339.512}{K}%
\TextAt{8}{146.919}{341.305}{be}%
\TextAt{8}{162.493}{341.305}{an}%
\TextAt{8}{178.402}{341.305}{ideal}%
\TextAt{8}{206.015}{341.305}{class}%
\TextAt{8}{233.102}{341.305}{chosen}%
\TextAt{8}{270.207}{341.305}{once}%
\TextAt{8}{296.517}{341.305}{and}%
\TextAt{8}{318.929}{341.305}{for}%
\TextAt{8}{336.464}{341.305}{all.}%
\TextAt{8}{357.157}{341.305}{We}%
\TextAt{8}{376.951}{341.305}{denote}%
\TextAt{8}{414.32}{341.305}{by}%
\TextAt{14}{430.229}{341.305}{K}%
\TextAt{3}{438.22}{345.644}{\Character{3}}%
\TextAt{8}{446.5}{341.305}{the}%
\TextAt{8}{466.306}{341.305}{classes}%
\TextAt{8}{503.208}{341.305}{in}%
\TextAt{8}{516.51}{341.305}{the}%
\TextAt{8}{79.226}{324.259}{ray}%
\TextAt{8}{99.395}{324.259}{class}%
\TextAt{8}{126.832}{324.259}{group}%
\TextAt{13}{160.007}{324.259}{C}%
\TextAt{13}{169.225}{324.259}{l}%
\TextAt{27}{172.991}{328.598}{m}%
\TextAt{21}{172.728}{321.239}{K}%
\TextAt{9}{184.251}{324.259}{\Character{61}}%
\TextAt{13}{196.677}{324.259}{J}%
\TextAt{27}{204.254}{328.598}{m}%
\TextAt{21}{203.147}{321.239}{K}%
\TextAt{9}{211.35}{324.259}{\Character{47}}%
\TextAt{13}{217.203}{324.259}{P}%
\TextAt{27}{226.374}{328.598}{m}%
\TextAt{21}{224.748}{321.239}{K}%
\TextAt{8}{236.88}{324.259}{which}%
\TextAt{8}{270.367}{324.259}{under}%
\TextAt{8}{303.53}{324.259}{the}%
\TextAt{8}{323.687}{324.259}{homomorphism}%
\TextAt{13}{272.123}{296.753}{C}%
\TextAt{13}{281.341}{296.753}{l}%
\TextAt{27}{285.107}{301.689}{m}%
\TextAt{21}{284.844}{293.797}{K}%
\NativeArrowAt{296.367}{296.753}{11.955}{19.26}{\longrightarrow}%
\TextAt{13}{318.949}{296.753}{C}%
\TextAt{13}{328.167}{296.753}{l}%
\TextAt{21}{331.67}{294.959}{K}%
\TextAt{8}{79.226}{269.246}{are}%
\TextAt{8}{98.977}{269.246}{sent}%
\TextAt{8}{123.665}{269.246}{to}%
\TextAt{8}{138.216}{269.246}{the}%
\TextAt{8}{158.612}{269.246}{ideal}%
\TextAt{8}{186.815}{269.246}{class}%
\TextAt{9}{214.493}{269.246}{\Character{40}}%
\TextAt{14}{219.045}{269.246}{m}%
\TextAt{9}{228.212}{269.246}{\Character{41}}%
\TextAt{14}{232.764}{269.246}{K}%
\TextAt{3}{240.755}{273.584}{\Character{0}}%
\TextAt{17}{247.342}{273.584}{1}%
\TextAt{8}{252.074}{269.246}{.}%
\TextAt{8}{261.243}{269.246}{Let}%
\TextAt{14}{282.452}{269.246}{a}%
\TextAt{8}{292.572}{269.246}{be}%
\TextAt{8}{308.737}{269.246}{an}%
\TextAt{8}{325.236}{269.246}{ideal}%
\TextAt{8}{353.439}{269.246}{in}%
\TextAt{14}{367.332}{269.246}{K}%
\TextAt{8}{375.323}{269.246}{,}%
\TextAt{8}{382.773}{269.246}{and}%
\TextAt{8}{405.775}{269.246}{let}%
\TextAt{14}{422.92}{269.246}{b}%
\TextAt{8}{433.192}{269.246}{be}%
\TextAt{8}{449.356}{269.246}{an}%
\TextAt{8}{465.855}{269.246}{integral}%
\TextAt{8}{508.704}{269.246}{ideal}%
\TextAt{8}{79.226}{252.2}{in}%
\TextAt{14}{93.055}{252.2}{K}%
\TextAt{3}{101.046}{256.538}{\Character{3}}%
\TextAt{8}{105.778}{252.2}{.}%
\TextAt{8}{114.754}{252.2}{Then}%
\TextAt{14}{145.487}{252.2}{abm}%
\TextAt{3}{166.772}{256.538}{\Character{0}}%
\TextAt{17}{173.358}{256.538}{1}%
\TextAt{9}{181.711}{252.2}{\Character{61}}%
\TextAt{9}{194.419}{252.2}{\Character{40}}%
\TextAt{13}{198.988}{252.2}{a}%
\TextAt{9}{205.133}{252.2}{\Character{41}}%
\TextAt{8}{213.759}{252.2}{is}%
\TextAt{8}{225.699}{252.2}{a}%
\TextAt{8}{235.63}{252.2}{principal}%
\TextAt{8}{284.583}{252.2}{ideal.}%
\TextAt{8}{317.624}{252.2}{The}%
\TextAt{8}{341.854}{252.2}{value}%
\TextAt{13}{372.276}{252.2}{\Character{28}}%
\TextAt{27}{377.36}{250.407}{a}%
\TextAt{9}{381.846}{252.2}{\Character{40}}%
\TextAt{13}{386.399}{252.2}{a}%
\TextAt{9}{392.544}{252.2}{\Character{41}}%
\TextAt{8}{401.169}{252.2}{only}%
\TextAt{8}{427.037}{252.2}{depends}%
\TextAt{8}{472.463}{252.2}{on}%
\TextAt{8}{488.898}{252.2}{the}%
\TextAt{8}{509.23}{252.2}{class}%
\TextAt{14}{79.226}{235.154}{K}%
\TextAt{3}{87.218}{239.492}{\Character{3}}%
\TextAt{8}{91.95}{235.154}{,}%
\TextAt{8}{99.099}{235.154}{not}%
\TextAt{8}{119.913}{235.154}{on}%
\TextAt{8}{136.172}{235.154}{the}%
\TextAt{8}{156.328}{235.154}{choice}%
\TextAt{8}{191.118}{235.154}{of}%
\TextAt{14}{204.448}{235.154}{a}%
\TextAt{13}{210.431}{235.154}{\Character{59}}%
\TextAt{14}{215.675}{235.154}{b}%
\TextAt{8}{225.708}{235.154}{and}%
\TextAt{13}{248.47}{235.154}{a}%
\TextAt{8}{254.615}{235.154}{.}%
\TextAt{8}{263.063}{235.154}{It}%
\TextAt{8}{275.736}{235.154}{will}%
\TextAt{8}{297.841}{235.154}{be}%
\TextAt{8}{313.765}{235.154}{denoted}%
\TextAt{8}{357.987}{235.154}{by}%
\TextAt{13}{269.965}{207.647}{\Character{28}}%
\TextAt{9}{276.384}{207.647}{\Character{40}}%
\TextAt{14}{280.936}{207.647}{K}%
\TextAt{3}{288.928}{212.583}{\Character{3}}%
\TextAt{9}{293.66}{207.647}{\Character{41}}%
\TextAt{9}{301.515}{207.647}{\Character{61}}%
\TextAt{13}{313.959}{207.647}{\Character{28}}%
\TextAt{27}{319.042}{205.854}{a}%
\TextAt{9}{323.529}{207.647}{\Character{40}}%
\TextAt{13}{328.082}{207.647}{a}%
\TextAt{9}{334.227}{207.647}{\Character{41}}%
\TextAt{13}{338.779}{207.647}{\Character{58}}%
\TextAt{8}{79.226}{180.14}{With}%
\TextAt{8}{109.461}{180.14}{these}%
\TextAt{8}{139.432}{180.14}{conventions}%
\TextAt{8}{202.592}{180.14}{we}%
\TextAt{8}{219.819}{180.14}{then}%
\TextAt{8}{246.479}{180.14}{have}%
\TextAt{8}{273.474}{180.14}{the}%
\TextAt{12}{79.226}{154.555}{Theorem}%
\TextAt{12}{136.24}{154.555}{(6.11).}%
\TextAt{8}{89.579}{134.52}{(i)}%
\TextAt{8}{107.918}{134.52}{The}%
\TextAt{8}{131.55}{134.52}{invariants}%
\TextAt{13}{184.867}{134.52}{\Character{28}}%
\TextAt{9}{191.285}{134.52}{\Character{40}}%
\TextAt{14}{195.838}{134.52}{K}%
\TextAt{3}{203.829}{138.859}{\Character{3}}%
\TextAt{17}{203.829}{131.565}{1}%
\TextAt{9}{208.562}{134.52}{\Character{41}}%
\TextAt{13}{213.114}{134.52}{\Character{59}}%
\TextAt{13}{218.338}{134.52}{\Character{28}}%
\TextAt{9}{224.777}{134.52}{\Character{40}}%
\TextAt{14}{229.329}{134.52}{K}%
\TextAt{3}{237.321}{138.859}{\Character{3}}%
\TextAt{17}{237.321}{131.565}{2}%
\TextAt{9}{242.053}{134.52}{\Character{41}}%
\TextAt{13}{246.605}{134.52}{\Character{59}}%
\TextAt{13}{251.829}{134.52}{\Character{58}}%
\TextAt{13}{257.066}{134.52}{\Character{58}}%
\TextAt{13}{262.302}{134.52}{\Character{58}}%
\TextAt{8}{265.589}{134.52}{,}%
\TextAt{8}{272.401}{134.52}{for}%
\TextAt{8}{289.864}{134.52}{a}%
\TextAt{8}{299.197}{134.52}{fixed}%
\TextAt{8}{327.061}{134.52}{ideal}%
\TextAt{8}{354.602}{134.52}{class}%
\TextAt{14}{381.617}{134.52}{K}%
\TextAt{8}{389.608}{134.52}{,}%
\TextAt{8}{396.42}{134.52}{are}%
\TextAt{8}{415.509}{134.52}{distinct}%
\TextAt{8}{457.42}{134.52}{algebraic}%
\TextAt{8}{506.432}{134.52}{num-}%
\TextAt{8}{107.918}{117.474}{bers}%
\TextAt{8}{133.012}{117.474}{which}%
\TextAt{8}{166.499}{117.474}{are}%
\TextAt{8}{186.01}{117.474}{conjugate}%
\TextAt{8}{239.019}{117.474}{over}%
\TextAt{8}{264.065}{117.474}{the}%
\TextAt{8}{284.221}{117.474}{Hilbert}%
\TextAt{8}{324.534}{117.474}{class}%
\TextAt{8}{351.971}{117.474}{field}%
\TextAt{13}{377.328}{117.474}{K}%
\TextAt{17}{388.146}{121.813}{1}%
\TextAt{9}{396.199}{117.474}{\Character{61}}%
\TextAt{13}{408.625}{117.474}{K}%
\TextAt{9}{419.442}{117.474}{\Character{40}}%
\TextAt{13}{423.995}{117.474}{j}%
\TextAt{9}{429.51}{117.474}{\Character{40}}%
\TextAt{14}{434.062}{117.474}{K}%
\TextAt{9}{442.054}{117.474}{\Character{41}\Character{41}}%
\TextAt{8}{451.158}{117.474}{.}%
\TextAt{8}{86.327}{100.428}{(ii)}%
\TextAt{8}{107.918}{100.428}{For}%
\TextAt{8}{128.888}{100.428}{an}%
\TextAt{8}{145.147}{100.428}{arbitrary}%
\TextAt{14}{194.916}{100.428}{K}%
\TextAt{3}{202.908}{104.767}{\Character{3}}%
\TextAt{8}{207.64}{100.428}{,}%
\TextAt{8}{214.789}{100.428}{the}%
\TextAt{8}{234.946}{100.428}{field}%
\TextAt{13}{260.303}{100.428}{K}%
\TextAt{9}{271.12}{100.428}{\Character{40}}%
\TextAt{13}{275.673}{100.428}{j}%
\TextAt{9}{281.188}{100.428}{\Character{40}}%
\TextAt{14}{285.74}{100.428}{K}%
\TextAt{9}{293.732}{100.428}{\Character{41}}%
\TextAt{13}{298.284}{100.428}{\Character{59}}%
\TextAt{13}{303.508}{100.428}{\Character{28}}%
\TextAt{9}{309.947}{100.428}{\Character{40}}%
\TextAt{14}{314.5}{100.428}{K}%
\TextAt{3}{322.491}{104.767}{\Character{3}}%
\TextAt{9}{327.223}{100.428}{\Character{41}\Character{41}}%
\TextAt{8}{340.225}{100.428}{is}%
\TextAt{8}{351.989}{100.428}{the}%
\TextAt{8}{372.146}{100.428}{ray}%
\TextAt{8}{392.314}{100.428}{class}%
\TextAt{8}{419.751}{100.428}{field}%
\TextAt{8}{445.108}{100.428}{mod}%
\TextAt{14}{471.445}{100.428}{m}%
\TextAt{8}{484.509}{100.428}{over}%
\TextAt{13}{509.555}{100.428}{K}%
\TextAt{8}{520.373}{100.428}{:}%
\TextAt{13}{264.119}{72.921}{K}%
\TextAt{27}{274.936}{77.858}{m}%
\TextAt{9}{284.867}{72.921}{\Character{61}}%
\TextAt{13}{297.292}{72.921}{K}%
\TextAt{9}{308.11}{72.921}{\Character{40}}%
\TextAt{13}{312.662}{72.921}{j}%
\TextAt{9}{318.178}{72.921}{\Character{40}}%
\TextAt{14}{322.73}{72.921}{K}%
\TextAt{9}{330.721}{72.921}{\Character{41}}%
\TextAt{13}{335.274}{72.921}{\Character{59}}%
\TextAt{13}{340.498}{72.921}{\Character{28}}%
\TextAt{9}{346.937}{72.921}{\Character{40}}%
\TextAt{14}{351.489}{72.921}{K}%
\TextAt{3}{359.481}{77.858}{\Character{3}}%
\TextAt{9}{364.213}{72.921}{\Character{41}\Character{41}}%
\TextAt{13}{373.317}{72.921}{\Character{58}}%
\TextAt{8}{297.211}{43.033}{367}%
\Destination{79.23}{698.335}{nav.Item.129}%
\Destination{79.23}{661.375}{nav.Item.130}%
\Destination{79.23}{644.315}{nav.Item.131}%
\Destination{79.23}{610.775}{nav.Item.132}%
\Destination{79.23}{148.565}{nav.Item.133}%
\Destination{79.23}{114.485}{nav.Item.134}%
\Destination{78.23}{809.815}{nav.page.367}%
\Destination{79.36}{645.255}{nav.print.403}%
\Destination{79.36}{732.455}{nav.number.6.6.10}%
\Destination{79.36}{168.555}{nav.number.6.6.11}%
}

% Page 380 (body page 368)
\Page{%
\TextAt{2}{79.226}{796.943}{VI.}%
\TextAt{2}{98.012}{796.943}{Global}%
\TextAt{2}{133.248}{796.943}{Class}%
\TextAt{2}{161.841}{796.943}{Field}%
\TextAt{2}{189.572}{796.943}{Theory}%
\Vector{q 0 G 0 g 0.299 w 7.226 18.729 m 460.769 18.729 l S Q}%
\TextAt{12}{79.226}{761.124}{Exercise}%
\TextAt{12}{132.594}{761.124}{VI.6.1.}%
\TextAt{8}{176.496}{761.124}{Let}%
\TextAt{13}{197.611}{761.124}{K}%
\TextAt{17}{208.428}{765.462}{1}%
\TextAt{16}{216.73}{761.124}{\Character{106}}%
\TextAt{13}{223.619}{761.124}{K}%
\TextAt{8}{238.48}{761.124}{be}%
\TextAt{8}{254.55}{761.124}{the}%
\TextAt{8}{274.852}{761.124}{big,}%
\TextAt{8}{297.797}{761.124}{and}%
\TextAt{13}{320.705}{761.124}{H}%
\TextAt{16}{334.896}{761.124}{\Character{106}}%
\TextAt{13}{341.786}{761.124}{K}%
\TextAt{8}{356.647}{761.124}{the}%
\TextAt{8}{376.949}{761.124}{small}%
\TextAt{8}{407.724}{761.124}{Hilbert}%
\TextAt{8}{448.182}{761.124}{class}%
\TextAt{8}{475.765}{761.124}{field.}%
\TextAt{8}{506.109}{761.124}{Then}%
\TextAt{13}{79.226}{744.078}{G}%
\TextAt{9}{88.46}{744.078}{\Character{40}}%
\TextAt{13}{93.012}{744.078}{K}%
\TextAt{17}{103.83}{748.416}{1}%
\TextAt{16}{111.883}{744.078}{\Character{106}}%
\TextAt{13}{118.525}{744.078}{H}%
\TextAt{9}{129.147}{744.078}{\Character{41}}%
\NativeMathAt{137.02}{744.078}{11.955}{\cong}%
\TextAt{9}{149.64}{744.078}{\Character{40}}%
\TextAt{18}{154.192}{744.078}{Z}%
\TextAt{9}{162.162}{744.078}{\Character{47}2}%
\TextAt{18}{173.868}{744.078}{Z}%
\TextAt{9}{181.838}{744.078}{\Character{41}}%
\TextAt{21}{186.391}{748.416}{r}%
\TextAt{3}{190.447}{748.416}{\Character{0}}%
\TextAt{21}{197.034}{748.416}{t}%
\TextAt{8}{200.59}{744.078}{,}%
\TextAt{8}{207.739}{744.078}{where}%
\TextAt{13}{241.548}{744.078}{r}%
\TextAt{8}{251.046}{744.078}{is}%
\TextAt{8}{262.81}{744.078}{the}%
\TextAt{8}{282.966}{744.078}{number}%
\TextAt{8}{325.563}{744.078}{of}%
\TextAt{8}{338.893}{744.078}{real}%
\TextAt{8}{361.655}{744.078}{places,}%
\TextAt{8}{399.434}{744.078}{and}%
\TextAt{9}{422.196}{744.078}{2}%
\TextAt{21}{428.049}{748.416}{t}%
\TextAt{9}{434.926}{744.078}{\Character{61}}%
\TextAt{9}{447.336}{744.078}{\Character{40}}%
\TextAt{16}{451.904}{744.078}{\Character{79}}%
\TextAt{3}{461.754}{748.416}{\Character{3}}%
\TextAt{9}{469.807}{744.078}{\Character{58}}%
\TextAt{16}{476.38}{744.078}{\Character{79}}%
\TextAt{3}{486.229}{748.416}{\Character{3}}%
\TextAt{17}{485.897}{741.122}{\Character{43}}%
\TextAt{9}{492.982}{744.078}{\Character{41}}%
\TextAt{8}{497.534}{744.078}{.}%
\ErrataAt{79.22}{712.055}{453.64}{\textbf{Exercise VI.6.2.} Let $d>1$ be squarefree and $K=\mathbb Q(\sqrt d)$. Let $\varepsilon$ be a}%
\ErrataAt{79.22}{695.055}{453.64}{fundamental unit of $K$. Then $[K^1:H]=1$ or $2$, according as $N_{K\mid\mathbb Q}(\varepsilon)=-1$}%
\ErrataAt{79.22}{677.955}{453.64}{or $+1$.}%
\TextAt{12}{79.226}{646.006}{Exercise}%
\TextAt{12}{131.461}{646.006}{VI.6.3.}%
\TextAt{8}{174.378}{646.006}{The}%
\TextAt{8}{197.592}{646.006}{group}%
\TextAt{9}{229.929}{646.006}{\Character{40}}%
\TextAt{13}{234.481}{646.006}{C}%
\TextAt{21}{242.856}{644.213}{K}%
\TextAt{9}{251.059}{646.006}{\Character{41}}%
\TextAt{21}{255.611}{650.344}{n}%
\TextAt{9}{264.568}{646.006}{\Character{61}}%
\TextAt{9}{276.978}{646.006}{\Character{40}}%
\TextAt{13}{281.546}{646.006}{I}%
\TextAt{21}{286.709}{644.213}{K}%
\TextAt{9}{294.912}{646.006}{\Character{41}}%
\TextAt{21}{299.464}{650.344}{n}%
\TextAt{13}{305.1}{646.006}{K}%
\TextAt{3}{315.918}{650.344}{\Character{3}}%
\TextAt{9}{320.65}{646.006}{\Character{47}}%
\TextAt{13}{326.503}{646.006}{K}%
\TextAt{3}{337.321}{650.344}{\Character{3}}%
\TextAt{8}{345.112}{646.006}{is}%
\TextAt{8}{356.036}{646.006}{the}%
\TextAt{8}{375.353}{646.006}{intersection}%
\TextAt{8}{437.338}{646.006}{of}%
\TextAt{8}{449.829}{646.006}{the}%
\TextAt{8}{469.146}{646.006}{norm}%
\TextAt{8}{498.877}{646.006}{groups}%
\TextAt{13}{79.226}{628.96}{N}%
\TextAt{21}{88.6}{627.1}{L}%
\TextAt{3}{94.359}{627.1}{\Character{106}}%
\TextAt{21}{96.711}{627.1}{K}%
\TextAt{13}{104.914}{628.96}{C}%
\TextAt{21}{113.289}{627.167}{L}%
\TextAt{8}{123.444}{628.96}{of}%
\TextAt{8}{136.774}{628.96}{all}%
\TextAt{8}{153.033}{628.96}{abelian}%
\TextAt{8}{193.68}{628.96}{extensions}%
\TextAt{13}{250.061}{628.96}{L}%
\TextAt{16}{261.346}{628.96}{\Character{106}}%
\TextAt{13}{267.988}{628.96}{K}%
\TextAt{8}{282.703}{628.96}{of}%
\TextAt{8}{296.033}{628.96}{exponent}%
\TextAt{13}{346.436}{628.96}{n}%
\TextAt{8}{353.424}{628.96}{.}%
\TextAt{12}{79.226}{596.97}{Exercise}%
\TextAt{12}{132.501}{596.97}{VI.6.4.}%
\TextAt{8}{176.321}{596.97}{(i)}%
\TextAt{8}{192.644}{596.97}{For}%
\TextAt{8}{213.678}{596.97}{a}%
\TextAt{8}{223.497}{596.97}{number}%
\TextAt{8}{266.158}{596.97}{field}%
\TextAt{13}{291.579}{596.97}{K}%
\TextAt{8}{302.397}{596.97}{,}%
\TextAt{8}{309.626}{596.97}{local}%
\TextAt{8}{337.331}{596.97}{Tate}%
\TextAt{8}{364.378}{596.97}{duality}%
\TextAt{8}{404.122}{596.97}{(see}%
\TextAt{8}{427.654}{596.97}{chap.}%
\TextAt{8}{458.61}{596.97}{V,}%
\TextAt{8}{474.599}{596.97}{§1,}%
\TextAt{8}{493.365}{596.97}{exercise}%
\TextAt{8}{79.226}{579.924}{2)}%
\TextAt{8}{93.536}{579.924}{yields}%
\TextAt{8}{126.437}{579.924}{a}%
\TextAt{8}{136.193}{579.924}{non-degenerate}%
\TextAt{8}{217.488}{579.924}{pairing}%
\TextAt{3}{195.691}{558.349}{\Character{48}}%
\TextAt{20}{189.202}{554.762}{\Character{89}}%
\TextAt{27}{194.845}{529.656}{p}%
\TextAt{13}{206.47}{543.404}{H}%
\TextAt{17}{217.093}{548.341}{1}%
\TextAt{9}{221.825}{543.404}{\Character{40}}%
\TextAt{13}{226.378}{543.404}{K}%
\TextAt{27}{236.337}{541.611}{p}%
\TextAt{13}{240.824}{543.404}{\Character{59}}%
\TextAt{18}{246.068}{543.404}{Z}%
\TextAt{9}{254.038}{543.404}{\Character{47}}%
\TextAt{13}{259.891}{543.404}{n}%
\TextAt{18}{266.878}{543.404}{Z}%
\TextAt{9}{274.849}{543.404}{\Character{41}}%
\TextAt{16}{282.058}{543.404}{\Character{2}}%
\TextAt{3}{300.502}{558.349}{\Character{48}}%
\TextAt{20}{294.013}{554.762}{\Character{89}}%
\TextAt{27}{299.656}{529.656}{p}%
\TextAt{13}{311.281}{543.404}{H}%
\TextAt{17}{321.904}{548.341}{1}%
\TextAt{9}{326.636}{543.404}{\Character{40}}%
\TextAt{13}{331.189}{543.404}{K}%
\TextAt{27}{341.148}{541.611}{p}%
\TextAt{13}{345.634}{543.404}{\Character{59}}%
\TextAt{13}{350.858}{543.404}{\Character{22}}%
\TextAt{21}{357.922}{541.611}{n}%
\TextAt{9}{363.558}{543.404}{\Character{41}}%
\NativeArrowAt{371.431}{543.404}{11.955}{19.26}{\longrightarrow}%
\TextAt{18}{394.013}{543.404}{Z}%
\TextAt{9}{401.983}{543.404}{\Character{47}}%
\TextAt{13}{407.836}{543.404}{n}%
\TextAt{18}{414.824}{543.404}{Z}%
\TextAt{8}{517.802}{543.404}{(*)}%
\TextAt{8}{79.226}{505.586}{of}%
\TextAt{8}{93.465}{505.586}{locally}%
\TextAt{8}{131.446}{505.586}{compact}%
\TextAt{8}{179.183}{505.586}{groups,}%
\TextAt{8}{221.361}{505.586}{where}%
\TextAt{8}{256.079}{505.586}{the}%
\TextAt{8}{277.144}{505.586}{restricted}%
\TextAt{8}{330.141}{505.586}{products}%
\TextAt{8}{379.564}{505.586}{are}%
\TextAt{8}{399.983}{505.586}{taken}%
\TextAt{8}{432.764}{505.586}{with}%
\TextAt{8}{460.333}{505.586}{respect}%
\TextAt{8}{501.292}{505.586}{to}%
\TextAt{8}{516.51}{505.586}{the}%
\TextAt{8}{79.226}{488.54}{subgroups}%
\TextAt{13}{134.994}{488.54}{H}%
\TextAt{17}{145.617}{492.878}{1}%
\TextAt{21}{144.693}{485.584}{nr}%
\TextAt{9}{154.386}{488.54}{\Character{40}}%
\TextAt{13}{158.938}{488.54}{K}%
\TextAt{27}{168.897}{486.747}{p}%
\TextAt{13}{173.384}{488.54}{\Character{59}}%
\TextAt{18}{178.628}{488.54}{Z}%
\TextAt{9}{186.599}{488.54}{\Character{47}}%
\TextAt{13}{192.452}{488.54}{n}%
\TextAt{18}{199.439}{488.54}{Z}%
\TextAt{9}{207.409}{488.54}{\Character{41}}%
\TextAt{8}{211.962}{488.54}{,}%
\TextAt{8}{219.556}{488.54}{resp.}%
\TextAt{13}{249.945}{488.54}{H}%
\TextAt{17}{260.568}{492.878}{1}%
\TextAt{21}{259.644}{485.584}{nr}%
\TextAt{9}{269.337}{488.54}{\Character{40}}%
\TextAt{13}{273.889}{488.54}{K}%
\TextAt{27}{283.848}{486.747}{p}%
\TextAt{13}{288.335}{488.54}{\Character{59}}%
\TextAt{13}{293.559}{488.54}{\Character{22}}%
\TextAt{21}{300.622}{486.747}{n}%
\TextAt{9}{306.259}{488.54}{\Character{41}}%
\TextAt{8}{310.811}{488.54}{.}%
\TextAt{8}{320.327}{488.54}{For}%
\TextAt{13}{341.652}{488.54}{\Character{31}}%
\TextAt{9}{352.913}{488.54}{\Character{61}}%
\TextAt{9}{365.932}{488.54}{\Character{40}}%
\TextAt{13}{370.497}{488.54}{\Character{31}}%
\TextAt{27}{377.83}{486.747}{p}%
\TextAt{9}{382.317}{488.54}{\Character{41}}%
\TextAt{8}{391.123}{488.54}{in}%
\TextAt{8}{405.131}{488.54}{the}%
\TextAt{8}{425.643}{488.54}{first}%
\TextAt{8}{450.125}{488.54}{and}%
\TextAt{13}{473.243}{488.54}{\Character{11}}%
\TextAt{9}{484.692}{488.54}{\Character{61}}%
\TextAt{9}{497.711}{488.54}{\Character{40}}%
\TextAt{13}{502.276}{488.54}{\Character{11}}%
\TextAt{27}{509.722}{486.747}{p}%
\TextAt{9}{514.209}{488.54}{\Character{41}}%
\TextAt{8}{523.014}{488.54}{in}%
\TextAt{8}{79.226}{471.494}{the}%
\TextAt{8}{99.383}{471.494}{second}%
\TextAt{8}{137.161}{471.494}{product,}%
\TextAt{8}{184.312}{471.494}{it}%
\TextAt{8}{196.016}{471.494}{is}%
\TextAt{8}{207.78}{471.494}{given}%
\TextAt{8}{238.349}{471.494}{by}%
\TextAt{16}{233.262}{442.493}{\Character{104}}%
\TextAt{13}{237.911}{442.493}{\Character{31}\Character{59}}%
\TextAt{13}{250.464}{442.493}{\Character{11}}%
\TextAt{16}{258.011}{442.493}{\Character{105}}%
\TextAt{9}{265.981}{442.493}{\Character{61}}%
\TextAt{20}{278.406}{453.85}{\Character{88}}%
\TextAt{27}{285.046}{428.744}{p}%
\TextAt{13}{297.667}{442.493}{\Character{31}}%
\TextAt{27}{305.001}{440.699}{p}%
\TextAt{9}{309.488}{442.493}{\Character{40}}%
\TextAt{13}{314.04}{442.493}{\Character{11}}%
\TextAt{27}{321.486}{440.699}{p}%
\TextAt{13}{325.973}{442.493}{\Character{59}}%
\Vector{q 0 G 0 g 0.478 w 259.217 -321.664 m 270.034 -321.664 l S Q}%
\TextAt{13}{331.217}{442.493}{K}%
\TextAt{27}{342.034}{440.699}{p}%
\TextAt{16}{349.842}{442.493}{\Character{106}}%
\TextAt{13}{356.484}{442.493}{K}%
\TextAt{27}{366.443}{440.699}{p}%
\TextAt{9}{370.93}{442.493}{\Character{41}}%
\TextAt{13}{375.482}{442.493}{\Character{58}}%
\TextAt{8}{79.226}{404.124}{(ii)}%
\TextAt{8}{98.737}{404.124}{If}%
\TextAt{13}{110.429}{404.124}{L}%
\TextAt{16}{121.714}{404.124}{\Character{106}}%
\TextAt{13}{128.356}{404.124}{K}%
\TextAt{8}{143.071}{404.124}{is}%
\TextAt{8}{154.835}{404.124}{a}%
\TextAt{8}{164.591}{404.124}{finite}%
\TextAt{8}{194.502}{404.124}{extension,}%
\TextAt{8}{249.52}{404.124}{then}%
\TextAt{8}{276.18}{404.124}{one}%
\TextAt{8}{297.64}{404.124}{has}%
\TextAt{8}{318.513}{404.124}{a}%
\TextAt{8}{328.269}{404.124}{commutative}%
\TextAt{8}{398.195}{404.124}{diagram}%
\TextAt{3}{189.057}{378.169}{\Character{48}}%
\TextAt{20}{182.567}{374.582}{\Character{89}}%
\TextAt{27}{186.907}{349.146}{P}%
\TextAt{13}{199.836}{363.225}{H}%
\TextAt{17}{210.458}{368.161}{1}%
\TextAt{9}{215.191}{363.225}{\Character{40}}%
\TextAt{13}{219.743}{363.225}{L}%
\TextAt{27}{227.707}{361.431}{P}%
\TextAt{13}{234.801}{363.225}{\Character{59}}%
\TextAt{18}{240.045}{363.225}{Z}%
\TextAt{9}{248.015}{363.225}{\Character{47}}%
\TextAt{13}{253.868}{363.225}{n}%
\TextAt{18}{260.856}{363.225}{Z}%
\TextAt{9}{268.826}{363.225}{\Character{41}}%
\TextAt{3}{304.52}{378.169}{\Character{48}}%
\TextAt{20}{298.03}{374.582}{\Character{89}}%
\TextAt{27}{302.37}{349.146}{P}%
\TextAt{13}{315.299}{363.225}{H}%
\TextAt{17}{325.921}{368.161}{1}%
\TextAt{9}{330.654}{363.225}{\Character{40}}%
\TextAt{13}{335.206}{363.225}{L}%
\TextAt{27}{343.17}{361.431}{P}%
\TextAt{13}{350.264}{363.225}{\Character{59}}%
\TextAt{13}{355.488}{363.225}{\Character{22}}%
\TextAt{21}{362.551}{361.431}{n}%
\TextAt{9}{368.188}{363.225}{\Character{41}}%
\TextAt{18}{399.017}{363.225}{Z}%
\TextAt{9}{406.987}{363.225}{\Character{47}}%
\TextAt{13}{412.84}{363.225}{n}%
\TextAt{18}{419.827}{363.225}{Z}%
\TextAt{3}{189.363}{310.541}{\Character{48}}%
\TextAt{20}{182.873}{306.955}{\Character{89}}%
\TextAt{27}{188.517}{281.849}{p}%
\TextAt{13}{200.142}{295.597}{H}%
\TextAt{17}{210.764}{300.534}{1}%
\TextAt{9}{215.497}{295.597}{\Character{40}}%
\TextAt{13}{220.049}{295.597}{K}%
\TextAt{27}{230.008}{293.804}{p}%
\TextAt{13}{234.495}{295.597}{\Character{59}}%
\TextAt{18}{239.739}{295.597}{Z}%
\TextAt{9}{247.709}{295.597}{\Character{47}}%
\TextAt{13}{253.562}{295.597}{n}%
\TextAt{18}{260.55}{295.597}{Z}%
\TextAt{9}{268.52}{295.597}{\Character{41}}%
\TextAt{3}{304.826}{310.541}{\Character{48}}%
\TextAt{20}{298.336}{306.955}{\Character{89}}%
\TextAt{27}{303.98}{281.849}{p}%
\TextAt{13}{315.605}{295.597}{H}%
\TextAt{17}{326.227}{300.534}{1}%
\TextAt{9}{330.96}{295.597}{\Character{40}}%
\TextAt{13}{335.512}{295.597}{K}%
\TextAt{27}{345.471}{293.804}{p}%
\TextAt{13}{349.958}{295.597}{\Character{59}}%
\TextAt{13}{355.182}{295.597}{\Character{22}}%
\TextAt{21}{362.245}{293.804}{n}%
\TextAt{9}{367.881}{295.597}{\Character{41}}%
\TextAt{18}{397.392}{295.597}{Z}%
\TextAt{9}{405.362}{295.597}{\Character{47}}%
\TextAt{13}{411.215}{295.597}{n}%
\TextAt{18}{418.202}{295.597}{Z}%
\TextAt{13}{426.173}{295.597}{\Character{58}}%
\TextAt{16}{281.057}{363.225}{\Character{2}}%
\NativeDiagram{%
\path[draw=black,line width=0.478bp] (306.095,-407.786) -- (321.191,-407.786);%
\path[draw=black,line width=0.478bp,line cap=round,line join=round] (318.943,-404.917) .. controls (319.397,-406.639) and (320.41,-407.452) .. (321.43,-407.786) .. controls (320.41,-408.121) and (319.397,-408.934) .. (318.943,-410.656);%
\path[draw=black,line width=0.478bp] (263.388,-432.135) -- (263.388,-452.778);%
\path[draw=black,line width=0.478bp,line cap=round,line join=round] (266.257,-450.531) .. controls (264.535,-450.985) and (263.723,-451.998) .. (263.388,-453.017) .. controls (263.053,-451.998) and (262.24,-450.985) .. (260.519,-450.531);%
}%
\TextAt{21}{338.203}{330.197}{N}%
\TextAt{22}{344.924}{328.459}{L}%
\TextAt{23}{349.983}{328.459}{\Character{106}}%
\TextAt{22}{352.252}{328.459}{K}%
\NativeDiagram{%
\path[draw=black,line width=2.797bp] (341.411,-418.343) -- (341.411,-464.858);%
\path[draw=black!0,line width=1.841bp] (341.411,-418.343) -- (341.411,-464.858);%
\path[draw=black,line width=0.478bp] (155.974,-453.257) -- (155.974,-432.614);%
\path[draw=black,line width=0.478bp,line cap=round,line join=round] (153.105,-434.861) .. controls (154.827,-434.407) and (155.64,-433.395) .. (155.974,-432.375) .. controls (156.309,-433.395) and (157.122,-434.407) .. (158.844,-434.861);%
}%
\TextAt{16}{281.057}{295.597}{\Character{2}}%
\NativeDiagram{%
\path[draw=black,line width=0.478bp] (305.789,-475.414) -- (319.565,-475.414);%
\path[draw=black,line width=0.478bp,line cap=round,line join=round] (317.317,-472.545) .. controls (317.772,-474.267) and (318.784,-475.08) .. (319.804,-475.414) .. controls (318.784,-475.749) and (317.772,-476.562) .. (317.317,-478.284);%
}%
\TextAt{8}{79.226}{252.849}{(iii)}%
\TextAt{8}{101.989}{252.849}{The}%
\TextAt{8}{126.043}{252.849}{images}%
\TextAt{8}{164.478}{252.849}{of}%
\TextAt{13}{213.621}{228.285}{H}%
\TextAt{17}{224.243}{233.221}{1}%
\TextAt{9}{228.976}{228.285}{\Character{40}}%
\TextAt{13}{233.528}{228.285}{K}%
\TextAt{13}{243.69}{228.285}{\Character{59}}%
\TextAt{18}{248.94}{228.285}{Z}%
\TextAt{9}{256.91}{228.285}{\Character{47}}%
\TextAt{13}{262.763}{228.285}{n}%
\TextAt{18}{269.75}{228.285}{Z}%
\TextAt{9}{277.72}{228.285}{\Character{41}}%
\NativeArrowAt{285.594}{228.285}{11.955}{19.26}{\longrightarrow}%
\TextAt{3}{314.665}{243.229}{\Character{48}}%
\TextAt{20}{308.176}{239.642}{\Character{89}}%
\TextAt{27}{313.819}{214.536}{p}%
\TextAt{13}{325.444}{228.285}{H}%
\TextAt{17}{336.067}{233.221}{1}%
\TextAt{9}{340.799}{228.285}{\Character{40}}%
\TextAt{13}{345.351}{228.285}{K}%
\TextAt{27}{355.311}{226.491}{p}%
\TextAt{13}{359.797}{228.285}{\Character{59}}%
\TextAt{18}{365.042}{228.285}{Z}%
\TextAt{9}{373.012}{228.285}{\Character{47}}%
\TextAt{13}{378.865}{228.285}{n}%
\TextAt{18}{385.852}{228.285}{Z}%
\TextAt{9}{393.822}{228.285}{\Character{41}}%
\TextAt{8}{79.226}{196.065}{and}%
\TextAt{13}{229.722}{174.382}{H}%
\TextAt{17}{240.345}{179.318}{1}%
\TextAt{9}{245.077}{174.382}{\Character{40}}%
\TextAt{13}{249.63}{174.382}{K}%
\TextAt{13}{259.792}{174.382}{\Character{59}}%
\TextAt{13}{265.028}{174.382}{\Character{22}}%
\TextAt{21}{272.084}{172.589}{n}%
\TextAt{9}{277.72}{174.382}{\Character{41}}%
\NativeArrowAt{285.594}{174.382}{11.955}{19.26}{\longrightarrow}%
\TextAt{3}{314.665}{189.326}{\Character{48}}%
\TextAt{20}{308.176}{185.739}{\Character{89}}%
\TextAt{27}{313.819}{160.633}{p}%
\TextAt{13}{325.444}{174.382}{H}%
\TextAt{17}{336.067}{179.318}{1}%
\TextAt{9}{340.799}{174.382}{\Character{40}}%
\TextAt{13}{345.351}{174.382}{K}%
\TextAt{27}{355.311}{172.589}{p}%
\TextAt{13}{359.797}{174.382}{\Character{59}}%
\TextAt{13}{365.021}{174.382}{\Character{22}}%
\TextAt{21}{372.085}{172.589}{n}%
\TextAt{9}{377.721}{174.382}{\Character{41}}%
\TextAt{8}{79.226}{141.493}{are}%
\TextAt{8}{98.737}{141.493}{mutual}%
\TextAt{8}{138.739}{141.493}{orthogonal}%
\TextAt{8}{197.295}{141.493}{complements}%
\TextAt{8}{267.269}{141.493}{with}%
\TextAt{8}{293.929}{141.493}{respect}%
\TextAt{8}{333.979}{141.493}{to}%
\TextAt{8}{348.289}{141.493}{the}%
\TextAt{8}{368.446}{141.493}{pairing}%
\TextAt{9}{408.125}{141.493}{\Character{40}}%
\TextAt{16}{412.677}{141.493}{\Character{3}}%
\TextAt{9}{418.655}{141.493}{\Character{41}}%
\TextAt{8}{423.207}{141.493}{.}%
\TextAt{8}{297.211}{43.033}{368}%
\Destination{522.36}{560.455}{nav.AMS.78}%
\Destination{78.23}{809.815}{nav.page.368}%
\Destination{79.226}{775.155}{nav.ex.6.6.1}%
\Destination{79.226}{726.055}{nav.ex.6.6.2}%
\Destination{79.226}{660.055}{nav.ex.6.6.3}%
\Destination{79.226}{610.955}{nav.ex.6.6.4}%
\Destination{79.36}{776.155}{nav.print.404}%
\Hotlink{458.11}{594.755}{9.775}{11.047}{nav.chapter.5}%
\Hotlink{479.778}{594.755}{6.858}{11.047}{nav.section.5.1}%
\Destination{79.226}{726.055}{errata.32}%
}

% Page 381 (body page 369)
\Page{%
\TextAt{2}{79.226}{796.943}{VI.}%
\TextAt{2}{98.012}{796.943}{Global}%
\TextAt{2}{133.248}{796.943}{Class}%
\TextAt{2}{161.841}{796.943}{Field}%
\TextAt{2}{189.572}{796.943}{Theory}%
\Vector{q 0 G 0 g 0.299 w 7.226 18.729 m 460.769 18.729 l S Q}%
\ErrataAt{97.06}{761.155}{435.8}{\textbf{Hint for (iii):} The cokernel of the second map is $C_K/(C_K)^n$, and one has}%
\TextAt{13}{79.226}{744.078}{H}%
\TextAt{17}{89.849}{748.416}{1}%
\TextAt{9}{94.581}{744.078}{\Character{40}}%
\TextAt{13}{99.133}{744.078}{K}%
\TextAt{13}{109.295}{744.078}{\Character{59}}%
\TextAt{18}{114.545}{744.078}{Z}%
\TextAt{9}{122.515}{744.078}{\Character{47}}%
\TextAt{13}{128.368}{744.078}{n}%
\TextAt{18}{135.356}{744.078}{Z}%
\TextAt{9}{143.326}{744.078}{\Character{41}}%
\TextAt{9}{151.181}{744.078}{\Character{61}}%
\TextAt{8}{163.624}{744.078}{Hom}%
\TextAt{9}{188.013}{744.078}{\Character{40}}%
\TextAt{13}{192.565}{744.078}{G}%
\TextAt{9}{201.799}{744.078}{\Character{40}}%
\TextAt{13}{206.351}{744.078}{L}%
\TextAt{16}{217.637}{744.078}{\Character{106}}%
\TextAt{13}{224.278}{744.078}{K}%
\TextAt{9}{235.096}{744.078}{\Character{41}}%
\TextAt{13}{239.648}{744.078}{\Character{59}}%
\TextAt{18}{244.892}{744.078}{Z}%
\TextAt{9}{252.863}{744.078}{\Character{47}}%
\TextAt{13}{258.716}{744.078}{n}%
\TextAt{18}{265.703}{744.078}{Z}%
\TextAt{9}{273.673}{744.078}{\Character{41}}%
\TextAt{8}{278.226}{744.078}{,}%
\TextAt{8}{285.31}{744.078}{where}%
\TextAt{13}{319.038}{744.078}{L}%
\TextAt{16}{330.324}{744.078}{\Character{106}}%
\TextAt{13}{336.965}{744.078}{K}%
\TextAt{8}{351.599}{744.078}{is}%
\TextAt{8}{363.282}{744.078}{the}%
\TextAt{8}{383.358}{744.078}{maximal}%
\TextAt{8}{431.085}{744.078}{abelian}%
\TextAt{8}{471.652}{744.078}{extension}%
\TextAt{8}{523.337}{744.078}{of}%
\TextAt{8}{79.226}{727.032}{exponent}%
\TextAt{13}{129.629}{727.032}{n}%
\TextAt{8}{136.617}{727.032}{.}%
\TextAt{12}{79.226}{695.042}{Exercise}%
\TextAt{12}{134.275}{695.042}{VI.6.5.}%
\TextAt{12}{179.639}{695.042}{(Global}%
\TextAt{12}{230.002}{695.042}{Tate}%
\TextAt{12}{262.468}{695.042}{Duality).}%
\TextAt{8}{325.232}{695.042}{Show}%
\TextAt{8}{357.731}{695.042}{that}%
\TextAt{8}{384.708}{695.042}{the}%
\TextAt{8}{406.471}{695.042}{statements}%
\TextAt{8}{467.065}{695.042}{of}%
\TextAt{8}{482.003}{695.042}{exercise}%
\TextAt{8}{526.911}{695.042}{4}%
\TextAt{8}{79.226}{677.996}{extend}%
\TextAt{8}{119.121}{677.996}{to}%
\TextAt{8}{135.284}{677.996}{an}%
\TextAt{8}{153.397}{677.996}{arbitrary}%
\TextAt{8}{205.019}{677.996}{finite}%
\TextAt{13}{236.784}{677.996}{G}%
\TextAt{21}{246.018}{676.203}{K}%
\TextAt{8}{254.221}{677.996}{-module}%
\TextAt{13}{301.264}{677.996}{A}%
\TextAt{8}{315.79}{677.996}{instead}%
\TextAt{8}{358.028}{677.996}{of}%
\TextAt{18}{373.211}{677.996}{Z}%
\TextAt{9}{381.181}{677.996}{\Character{47}}%
\TextAt{13}{387.034}{677.996}{n}%
\TextAt{18}{394.022}{677.996}{Z}%
\TextAt{8}{401.992}{677.996}{,}%
\TextAt{8}{411.458}{677.996}{and}%
\TextAt{13}{436.074}{677.996}{A}%
\TextAt{3}{444.849}{682.334}{\Character{48}}%
\TextAt{9}{454.123}{677.996}{\Character{61}}%
\TextAt{8}{469.707}{677.996}{Hom}%
\TextAt{9}{494.095}{677.996}{\Character{40}}%
\TextAt{13}{498.648}{677.996}{A\Character{59}}%
\Vector{q 0 G 0 g 0.478 w 440.667 -86.161 m 451.485 -86.161 l S Q}%
\TextAt{13}{512.667}{677.996}{K}%
\TextAt{3}{523.485}{685.401}{\Character{3}}%
\TextAt{9}{528.217}{677.996}{\Character{41}}%
\TextAt{8}{79.226}{660.95}{instead}%
\TextAt{8}{119.611}{660.95}{of}%
\TextAt{13}{132.941}{660.95}{\Character{22}}%
\TextAt{21}{139.984}{659.157}{n}%
\TextAt{8}{145.62}{660.95}{.}%
\TextAt{12}{97.159}{640.915}{Hint:}%
\TextAt{8}{132.698}{640.915}{Use}%
\TextAt{8}{155.185}{640.915}{exercises}%
\TextAt{8}{203.102}{640.915}{4–8}%
\TextAt{8}{224.573}{640.915}{of}%
\TextAt{8}{237.903}{640.915}{chap.}%
\TextAt{8}{268.795}{640.915}{IV,}%
\TextAt{8}{288.94}{640.915}{§3,}%
\TextAt{8}{307.626}{640.915}{and}%
\TextAt{8}{330.388}{640.915}{exercise}%
\TextAt{8}{373.69}{640.915}{4}%
\TextAt{8}{383.445}{640.915}{of}%
\TextAt{8}{396.775}{640.915}{chap.}%
\TextAt{8}{427.668}{640.915}{V,}%
\TextAt{8}{443.592}{640.915}{§1.}%
\TextAt{12}{79.226}{608.925}{Exercise}%
\TextAt{12}{132.427}{608.925}{VI.6.6.}%
\TextAt{8}{176.183}{608.925}{If}%
\TextAt{13}{187.875}{608.925}{S}%
\TextAt{8}{199.667}{608.925}{is}%
\TextAt{8}{211.431}{608.925}{a}%
\TextAt{8}{221.187}{608.925}{finite}%
\TextAt{8}{251.099}{608.925}{set}%
\TextAt{8}{269.366}{608.925}{of}%
\TextAt{8}{282.696}{608.925}{places}%
\TextAt{8}{317.223}{608.925}{of}%
\TextAt{13}{330.553}{608.925}{K}%
\TextAt{8}{341.37}{608.925}{,}%
\TextAt{8}{348.52}{608.925}{then}%
\TextAt{8}{375.18}{608.925}{the}%
\TextAt{8}{395.336}{608.925}{map}%
\TextAt{13}{213.621}{579.924}{H}%
\TextAt{17}{224.243}{584.86}{1}%
\TextAt{9}{228.976}{579.924}{\Character{40}}%
\TextAt{13}{233.528}{579.924}{K}%
\TextAt{13}{243.69}{579.924}{\Character{59}}%
\TextAt{18}{248.94}{579.924}{Z}%
\TextAt{9}{256.91}{579.924}{\Character{47}}%
\TextAt{13}{262.763}{579.924}{n}%
\TextAt{18}{269.75}{579.924}{Z}%
\TextAt{9}{277.72}{579.924}{\Character{41}}%
\NativeArrowAt{285.594}{579.924}{11.955}{19.26}{\longrightarrow}%
\TextAt{20}{308.176}{591.281}{\Character{89}}%
\TextAt{27}{308.199}{565.91}{p}%
\NativeMathAt{312.187}{565.91}{7.97}{\in}%
\TextAt{21}{317.833}{565.91}{S}%
\TextAt{13}{325.444}{579.924}{H}%
\TextAt{17}{336.067}{584.86}{1}%
\TextAt{9}{340.799}{579.924}{\Character{40}}%
\TextAt{13}{345.351}{579.924}{K}%
\TextAt{27}{355.311}{578.131}{p}%
\TextAt{13}{359.797}{579.924}{\Character{59}}%
\TextAt{18}{365.042}{579.924}{Z}%
\TextAt{9}{373.012}{579.924}{\Character{47}}%
\TextAt{13}{378.865}{579.924}{n}%
\TextAt{18}{385.852}{579.924}{Z}%
\TextAt{9}{393.822}{579.924}{\Character{41}}%
\ErrataAt{79.19}{541.764}{455.81}{is surjective if and only if the kernel of the map}%
\TextAt{13}{229.722}{512.839}{H}%
\TextAt{17}{240.345}{517.775}{1}%
\TextAt{9}{245.077}{512.839}{\Character{40}}%
\TextAt{13}{249.63}{512.839}{K}%
\TextAt{13}{259.792}{512.839}{\Character{59}}%
\TextAt{13}{265.028}{512.839}{\Character{22}}%
\TextAt{21}{272.084}{511.045}{n}%
\TextAt{9}{277.72}{512.839}{\Character{41}}%
\NativeArrowAt{285.594}{512.839}{11.955}{19.26}{\longrightarrow}%
\TextAt{20}{308.176}{524.196}{\Character{89}}%
\TextAt{27}{308.199}{498.293}{p}%
\TextAt{17}{313.128}{498.293}{\Character{47}}%
\NativeMathAt{312.187}{498.293}{7.97}{\in}%
\TextAt{21}{317.833}{498.293}{S}%
\TextAt{13}{325.444}{512.839}{H}%
\TextAt{17}{336.067}{517.775}{1}%
\TextAt{9}{340.799}{512.839}{\Character{40}}%
\TextAt{13}{345.351}{512.839}{K}%
\TextAt{27}{355.311}{511.045}{p}%
\TextAt{13}{359.797}{512.839}{\Character{59}}%
\TextAt{13}{365.021}{512.839}{\Character{22}}%
\TextAt{21}{372.085}{511.045}{n}%
\TextAt{9}{377.721}{512.839}{\Character{41}}%
\ErrataAt{79.19}{473.682}{455.81}{equals the kernel of localization at all places. In particular, injectivity of the displayed map}%
\ErrataAt{79.19}{456.636}{455.81}{is sufficient. Write $n=2^\nu m$ with $m$ odd. Injectivity holds if $K(\mu_{2^\nu})\mid K$ is cyclic,}%
\ErrataAt{79.19}{439.59}{455.81}{or if a place above $2$ outside $S$ has full decomposition group in $K(\mu_{2^\nu})\mid K$.}%
\TextAt{12}{79.226}{407.591}{Exercise}%
\TextAt{12}{134.61}{407.591}{VI.6.7.}%
\ErrataAt{181}{407.593}{354}{(Theorem of Grunwald). Given cyclic extensions $L_{\mathfrak p}\mid K_{\mathfrak p}$ for $\mathfrak p\in S$, put}%
\ErrataAt{79.19}{390.547}{455.81}{$n=\operatorname{lcm}_{\mathfrak p\in S}[L_{\mathfrak p}:K_{\mathfrak p}]$. If the character restriction of Exercise \hyperlink{nav.ex.6.6.6}{6} is surjective for $(K,n,S)$, then}%
\ErrataAt{79.19}{373.501}{455.81}{there exists a cyclic extension $L\mid K$ which has $L_{\mathfrak p}\mid K_{\mathfrak p}$ as a completion for $\mathfrak p\in S$, and}%
\TextAt{8}{79.226}{356.453}{which}%
\TextAt{8}{112.713}{356.453}{satisfies}%
\TextAt{8}{155.823}{356.453}{the}%
\TextAt{8}{175.979}{356.453}{identity}%
\TextAt{8}{219.233}{356.453}{of}%
\TextAt{8}{232.563}{356.453}{degrees}%
\TextAt{9}{243.125}{327.452}{\Character{91}}%
\TextAt{13}{246.377}{327.452}{L}%
\TextAt{9}{257.662}{327.452}{\Character{58}}%
\TextAt{13}{264.235}{327.452}{K}%
\TextAt{9}{275.053}{327.452}{\Character{93}}%
\TextAt{9}{281.604}{327.452}{\Character{61}}%
\TextAt{8}{294.05}{327.452}{scm}%
\TextAt{16}{313.621}{327.452}{\Character{102}}%
\TextAt{9}{319.599}{327.452}{\Character{91}}%
\TextAt{13}{322.85}{327.452}{L}%
\TextAt{27}{330.815}{325.659}{p}%
\TextAt{9}{338.623}{327.452}{\Character{58}}%
\TextAt{13}{345.195}{327.452}{K}%
\TextAt{27}{355.154}{325.659}{p}%
\TextAt{9}{359.641}{327.452}{\Character{93}}%
\TextAt{16}{362.893}{327.452}{\Character{103}}%
\TextAt{8}{79.226}{298.451}{(see}%
\TextAt{8}{102.694}{298.451}{also}%
\TextAt{8}{126.174}{298.451}{[10],}%
\TextAt{8}{151.543}{298.451}{chap.}%
\TextAt{8}{182.435}{298.451}{X,}%
\TextAt{8}{198.359}{298.451}{§2).}%
\TextAt{12}{79.226}{272.438}{Note:}%
\TextAt{8}{116.402}{272.438}{Let}%
\TextAt{13}{136.604}{272.438}{G}%
\TextAt{8}{148.967}{272.438}{be}%
\TextAt{8}{164.124}{272.438}{a}%
\TextAt{8}{173.112}{272.438}{finite}%
\TextAt{8}{202.257}{272.438}{group}%
\TextAt{8}{234.665}{272.438}{of}%
\TextAt{8}{247.228}{272.438}{order}%
\TextAt{8}{277.03}{272.438}{prime}%
\TextAt{8}{309.426}{272.438}{to}%
\TextAt{8}{322.969}{272.438}{\char35\relax }%
\TextAt{13}{332.724}{272.438}{\Character{22}}%
\TextAt{9}{339.767}{272.438}{\Character{40}}%
\TextAt{13}{344.32}{272.438}{K}%
\TextAt{9}{355.137}{272.438}{\Character{41}}%
\TextAt{8}{359.69}{272.438}{,}%
\TextAt{8}{366.225}{272.438}{let}%
\TextAt{13}{382.362}{272.438}{S}%
\TextAt{8}{393.388}{272.438}{be}%
\TextAt{8}{408.544}{272.438}{a}%
\TextAt{8}{417.533}{272.438}{finite}%
\TextAt{8}{446.677}{272.438}{set}%
\TextAt{8}{464.177}{272.438}{of}%
\TextAt{8}{476.74}{272.438}{places,}%
\TextAt{8}{513.904}{272.438}{and}%
\TextAt{8}{79.226}{255.392}{let}%
\TextAt{13}{96.308}{255.392}{L}%
\TextAt{27}{104.272}{253.599}{p}%
\TextAt{16}{112.382}{255.392}{\Character{106}}%
\TextAt{13}{119.326}{255.392}{K}%
\TextAt{27}{129.285}{253.599}{p}%
\TextAt{8}{133.772}{255.392}{,}%
\TextAt{14}{141.142}{255.392}{p}%
\NativeMathAt{150.748}{255.392}{11.955}{\in}%
\TextAt{13}{162.341}{255.392}{S}%
\TextAt{8}{170.236}{255.392}{,}%
\TextAt{8}{177.607}{255.392}{be}%
\TextAt{8}{193.708}{255.392}{given}%
\TextAt{8}{224.455}{255.392}{Galois}%
\TextAt{8}{260.546}{255.392}{extensions}%
\TextAt{8}{317.103}{255.392}{whose}%
\TextAt{8}{351.807}{255.392}{Galois}%
\TextAt{8}{387.897}{255.392}{groups}%
\TextAt{13}{425.865}{255.392}{G}%
\TextAt{27}{435.099}{253.599}{p}%
\TextAt{8}{443.66}{255.392}{can}%
\TextAt{8}{465.297}{255.392}{be}%
\TextAt{8}{481.398}{255.392}{embedded}%
\TextAt{8}{79.226}{238.347}{into}%
\TextAt{13}{103.74}{238.347}{G}%
\TextAt{8}{112.974}{238.347}{.}%
\TextAt{8}{123.736}{238.347}{Then}%
\TextAt{8}{155.065}{238.347}{there}%
\TextAt{8}{185.748}{238.347}{exists}%
\TextAt{8}{218.834}{238.347}{a}%
\TextAt{8}{229.361}{238.347}{Galois}%
\TextAt{8}{266.045}{238.347}{extension}%
\TextAt{13}{318.582}{238.347}{L}%
\TextAt{16}{331.182}{238.347}{\Character{106}}%
\TextAt{13}{339.138}{238.347}{K}%
\TextAt{8}{354.625}{238.347}{which}%
\TextAt{8}{388.882}{238.347}{on}%
\TextAt{8}{405.913}{238.347}{the}%
\TextAt{8}{426.84}{238.347}{one}%
\TextAt{8}{449.071}{238.347}{hand}%
\TextAt{8}{479.109}{238.347}{has}%
\TextAt{8}{500.754}{238.347}{Galois}%
\TextAt{8}{79.226}{221.301}{group}%
\TextAt{8}{112.315}{221.301}{isomorphic}%
\TextAt{8}{171.478}{221.301}{to}%
\TextAt{13}{185.701}{221.301}{G}%
\TextAt{8}{194.935}{221.301}{,}%
\TextAt{8}{202.015}{221.301}{and}%
\TextAt{8}{224.69}{221.301}{which}%
\TextAt{8}{258.09}{221.301}{on}%
\TextAt{8}{274.262}{221.301}{the}%
\TextAt{8}{294.332}{221.301}{other}%
\TextAt{8}{324.814}{221.301}{hand}%
\TextAt{8}{353.994}{221.301}{has}%
\TextAt{8}{374.781}{221.301}{the}%
\TextAt{8}{394.85}{221.301}{given}%
\TextAt{8}{425.333}{221.301}{extensions}%
\TextAt{13}{481.626}{221.301}{L}%
\TextAt{27}{489.591}{219.507}{p}%
\TextAt{16}{497.399}{221.301}{\Character{106}}%
\TextAt{13}{504.04}{221.301}{K}%
\TextAt{27}{513.999}{219.507}{p}%
\TextAt{8}{522.297}{221.301}{as}%
\TextAt{8}{79.226}{204.255}{completions}%
\TextAt{8}{143.676}{204.255}{(see}%
\TextAt{8}{167.144}{204.255}{[109]).}%
\TextAt{8}{297.211}{43.033}{369}%
\Destination{78.23}{809.815}{nav.page.369}%
\Destination{79.226}{709.055}{nav.ex.6.6.5}%
\Destination{79.226}{622.955}{nav.ex.6.6.6}%
\Destination{79.226}{421.555}{nav.ex.6.6.7}%
\Destination{79.36}{488.655}{nav.print.405}%
\Hotlink{526.411}{692.855}{6.858}{11.047}{nav.ex.6.6.4}%
\Hotlink{202.602}{638.755}{6.858}{11.047}{nav.ex.6.3.4}%
\Hotlink{373.19}{638.755}{6.858}{11.047}{nav.ex.5.1.4}%
\Hotlink{268.295}{638.755}{13.995}{11.047}{nav.chapter.4}%
\Hotlink{427.168}{638.755}{9.775}{11.047}{nav.chapter.5}%
\Hotlink{294.119}{638.755}{6.858}{11.047}{nav.section.4.3}%
\Hotlink{448.771}{638.755}{6.858}{11.047}{nav.section.5.1}%
\Hotlink{164.182}{437.355}{6.858}{11.047}{nav.ex.6.1.2}%
\Hotlink{512.749}{405.355}{6.858}{11.047}{nav.ex.6.6.6}%
\Hotlink{203.538}{296.255}{6.858}{11.047}{nav.section.6.2}%
\Hotlink{125.674}{296.255}{19.22}{11.047}{nav.cite.N010}%
\Hotlink{166.644}{202.055}{25.078}{11.047}{nav.cite.N109}%
\Destination{79.226}{775.155}{errata.33}%
\Destination{79.226}{558.764}{errata.138}%
\Destination{79.226}{424.593}{errata.139}%
\Destination{79.226}{487}{errata.141}%
}

% Page 382 (body page 370)
\Page{%
\TextAt{2}{79.226}{796.943}{VI.}%
\TextAt{2}{98.012}{796.943}{Global}%
\TextAt{2}{133.248}{796.943}{Class}%
\TextAt{2}{161.841}{796.943}{Field}%
\TextAt{2}{189.572}{796.943}{Theory}%
\Vector{q 0 G 0 g 0.299 w 7.226 18.729 m 460.769 18.729 l S Q}%
\TextAt{15}{79.226}{760.804}{§}%
\TextAt{15}{93.785}{760.804}{7}%
\TextAt{15}{120.693}{760.804}{The}%
\TextAt{15}{161.234}{760.804}{Ideal-Theoretic}%
\TextAt{15}{297.82}{760.804}{Version}%
\TextAt{15}{368.556}{760.804}{of}%
\TextAt{15}{391.658}{760.804}{Class}%
\TextAt{15}{443.234}{760.804}{Field}%
\TextAt{15}{493.26}{760.804}{The-}%
\TextAt{15}{120.693}{734.941}{ory}%
\TextAt{8}{79.226}{703.055}{Class}%
\TextAt{8}{110.791}{703.055}{field}%
\TextAt{8}{137.025}{703.055}{theory}%
\TextAt{8}{174.651}{703.055}{has}%
\TextAt{8}{196.402}{703.055}{found}%
\TextAt{8}{230.119}{703.055}{its}%
\TextAt{8}{247.314}{703.055}{idèle-theoretic}%
\TextAt{8}{324.273}{703.055}{formulation}%
\TextAt{8}{388.249}{703.055}{only}%
\TextAt{8}{414.817}{703.055}{after}%
\TextAt{8}{443.334}{703.055}{it}%
\TextAt{8}{455.914}{703.055}{had}%
\TextAt{8}{479.553}{703.055}{been}%
\TextAt{8}{508.058}{703.055}{com-}%
\TextAt{8}{79.226}{686.01}{pleted}%
\TextAt{8}{113.415}{686.01}{in}%
\TextAt{8}{126.144}{686.01}{the}%
\TextAt{8}{145.376}{686.01}{language}%
\TextAt{8}{193.242}{686.01}{of}%
\TextAt{8}{205.648}{686.01}{ideals.}%
\TextAt{8}{242.469}{686.01}{From}%
\TextAt{8}{272.27}{686.01}{the}%
\TextAt{8}{291.502}{686.01}{very}%
\TextAt{8}{316.27}{686.01}{start,}%
\TextAt{8}{346.818}{686.01}{it}%
\TextAt{8}{357.598}{686.01}{was}%
\TextAt{8}{379.174}{686.01}{guided}%
\TextAt{8}{415.969}{686.01}{by}%
\TextAt{8}{431.304}{686.01}{the}%
\TextAt{8}{450.536}{686.01}{desire}%
\TextAt{8}{482.836}{686.01}{to}%
\TextAt{8}{496.223}{686.01}{classify}%
\TextAt{8}{79.226}{668.964}{all}%
\TextAt{8}{95.876}{668.964}{abelian}%
\TextAt{8}{136.914}{668.964}{extensions}%
\TextAt{8}{193.686}{668.964}{of}%
\TextAt{8}{207.406}{668.964}{a}%
\TextAt{8}{217.552}{668.964}{number}%
\TextAt{8}{260.539}{668.964}{field}%
\TextAt{13}{286.287}{668.964}{K}%
\TextAt{8}{297.105}{668.964}{.}%
\TextAt{8}{306.725}{668.964}{But}%
\TextAt{8}{330.357}{668.964}{at}%
\TextAt{8}{345.058}{668.964}{first,}%
\TextAt{8}{372.924}{668.964}{instead}%
\TextAt{8}{413.699}{668.964}{of}%
\TextAt{8}{427.42}{668.964}{the}%
\TextAt{8}{447.967}{668.964}{idèle}%
\TextAt{8}{475.663}{668.964}{class}%
\TextAt{8}{503.491}{668.964}{group}%
\TextAt{13}{79.226}{651.918}{C}%
\TextAt{21}{87.601}{650.124}{K}%
\TextAt{8}{95.804}{651.918}{,}%
\TextAt{8}{102.421}{651.918}{there}%
\TextAt{8}{131.669}{651.918}{was}%
\TextAt{8}{153.504}{651.918}{only}%
\TextAt{8}{178.531}{651.918}{the}%
\TextAt{8}{198.023}{651.918}{ideal}%
\TextAt{8}{225.321}{651.918}{class}%
\TextAt{8}{252.094}{651.918}{group}%
\TextAt{13}{284.605}{651.918}{C}%
\TextAt{13}{293.822}{651.918}{l}%
\TextAt{21}{297.325}{650.124}{K}%
\TextAt{8}{308.761}{651.918}{at}%
\TextAt{8}{322.406}{651.918}{hand}%
\TextAt{8}{351.008}{651.918}{to}%
\TextAt{8}{364.654}{651.918}{do}%
\TextAt{8}{380.248}{651.918}{this,}%
\TextAt{8}{405.791}{651.918}{along}%
\TextAt{8}{436.353}{651.918}{with}%
\TextAt{8}{462.349}{651.918}{its}%
\TextAt{8}{478.003}{651.918}{subgroups.}%
\TextAt{8}{79.226}{634.872}{In}%
\TextAt{8}{93.694}{634.872}{terms}%
\TextAt{8}{126.118}{634.872}{of}%
\TextAt{8}{139.294}{634.872}{the}%
\TextAt{8}{159.297}{634.872}{insights}%
\TextAt{8}{201.871}{634.872}{that}%
\TextAt{8}{227.087}{634.872}{we}%
\TextAt{8}{244.16}{634.872}{have}%
\TextAt{8}{271.002}{634.872}{gained}%
\TextAt{8}{307.921}{634.872}{in}%
\TextAt{8}{321.42}{634.872}{the}%
\TextAt{8}{341.423}{634.872}{preceding}%
\TextAt{8}{393.944}{634.872}{sections,}%
\TextAt{8}{440.769}{634.872}{this}%
\TextAt{8}{463.438}{634.872}{means}%
\TextAt{8}{499.114}{634.872}{the}%
\TextAt{8}{519.117}{634.872}{re-}%
\TextAt{8}{79.226}{617.826}{striction}%
\TextAt{8}{124.553}{617.826}{to}%
\TextAt{8}{137.947}{617.826}{the}%
\TextAt{8}{157.188}{617.826}{subfields}%
\TextAt{8}{203.865}{617.826}{of}%
\TextAt{8}{216.279}{617.826}{the}%
\TextAt{8}{235.52}{617.826}{Hilbert}%
\TextAt{8}{274.916}{617.826}{class}%
\TextAt{8}{301.438}{617.826}{field,}%
\TextAt{8}{329.314}{617.826}{i.e.,}%
\TextAt{8}{350.686}{617.826}{to}%
\TextAt{8}{364.08}{617.826}{the}%
\TextAt{19}{383.321}{617.826}{unramified}%
\TextAt{8}{440.555}{617.826}{abelian}%
\TextAt{8}{480.286}{617.826}{extensions}%
\TextAt{8}{79.226}{600.78}{of}%
\TextAt{13}{92.812}{600.78}{K}%
\TextAt{8}{103.629}{600.78}{.}%
\TextAt{8}{112.844}{600.78}{If}%
\TextAt{8}{124.791}{600.78}{the}%
\TextAt{8}{145.203}{600.78}{base}%
\TextAt{8}{171.533}{600.78}{field}%
\TextAt{8}{197.145}{600.78}{is}%
\TextAt{18}{209.164}{600.78}{Q}%
\TextAt{8}{218.463}{600.78}{,}%
\TextAt{8}{225.931}{600.78}{this}%
\TextAt{8}{249.009}{600.78}{restriction}%
\TextAt{8}{305.262}{600.78}{is}%
\TextAt{8}{317.282}{600.78}{of}%
\TextAt{8}{330.867}{600.78}{course}%
\TextAt{8}{366.952}{600.78}{radical,}%
\TextAt{8}{408.899}{600.78}{for}%
\TextAt{18}{427.04}{600.78}{Q}%
\TextAt{8}{440.491}{600.78}{has}%
\TextAt{8}{461.62}{600.78}{no}%
\TextAt{8}{478.134}{600.78}{unramified}%
\TextAt{8}{79.226}{583.734}{extensions}%
\TextAt{8}{136.437}{583.734}{at}%
\TextAt{8}{151.577}{583.734}{all}%
\TextAt{8}{168.666}{583.734}{by}%
\TextAt{8}{185.755}{583.734}{Minkowski’s}%
\TextAt{8}{252.72}{583.734}{theorem.}%
\TextAt{8}{305.286}{583.734}{But}%
\TextAt{8}{329.357}{583.734}{over}%
\TextAt{18}{355.233}{583.734}{Q}%
\TextAt{8}{364.531}{583.734}{,}%
\TextAt{8}{372.718}{583.734}{we}%
\TextAt{8}{390.775}{583.734}{naturally}%
\TextAt{8}{442.02}{583.734}{encounter}%
\TextAt{8}{496.505}{583.734}{the}%
\TextAt{8}{517.491}{583.734}{cy-}%
\TextAt{8}{79.226}{566.688}{clotomic}%
\TextAt{8}{126.818}{566.688}{fields}%
\TextAt{18}{157.552}{566.688}{Q}%
\TextAt{9}{166.851}{566.688}{\Character{40}}%
\TextAt{13}{171.403}{566.688}{\Character{22}}%
\TextAt{21}{178.446}{564.894}{m}%
\TextAt{9}{186.435}{566.688}{\Character{41}}%
\TextAt{16}{195.608}{566.688}{\Character{106}}%
\TextAt{18}{203.55}{566.688}{Q}%
\TextAt{8}{217.509}{566.688}{with}%
\TextAt{8}{244.932}{566.688}{their}%
\TextAt{8}{273.659}{566.688}{familiar}%
\TextAt{8}{317.675}{566.688}{isomorphisms}%
\TextAt{13}{391.473}{566.688}{G}%
\TextAt{9}{400.706}{566.688}{\Character{40}}%
\TextAt{18}{405.259}{566.688}{Q}%
\TextAt{9}{414.557}{566.688}{\Character{40}}%
\TextAt{13}{419.109}{566.688}{\Character{22}}%
\TextAt{21}{426.152}{564.894}{m}%
\TextAt{9}{434.141}{566.688}{\Character{41}}%
\TextAt{16}{443.314}{566.688}{\Character{106}}%
\TextAt{18}{451.257}{566.688}{Q}%
\TextAt{9}{460.555}{566.688}{\Character{41}}%
\NativeMathAt{469.729}{566.688}{11.955}{\cong}%
\TextAt{9}{483.648}{566.688}{\Character{40}}%
\TextAt{18}{488.201}{566.688}{Z}%
\TextAt{9}{496.171}{566.688}{\Character{47}}%
\TextAt{13}{502.024}{566.688}{m}%
\TextAt{18}{512.263}{566.688}{Z}%
\TextAt{9}{520.233}{566.688}{\Character{41}}%
\TextAt{3}{524.785}{571.026}{\Character{3}}%
\TextAt{8}{529.518}{566.688}{.}%
\TextAt{26}{79.226}{549.642}{Heinrich}%
\TextAt{26}{135.363}{549.642}{Weber}%
\TextAt{8}{179.158}{549.642}{realized,}%
\TextAt{8}{224.69}{549.642}{as}%
\TextAt{8}{238.262}{549.642}{was}%
\TextAt{8}{259.964}{549.642}{already}%
\TextAt{8}{300.471}{549.642}{mentioned,}%
\TextAt{8}{359.991}{549.642}{that}%
\TextAt{8}{384.562}{549.642}{the}%
\TextAt{8}{403.92}{549.642}{groups}%
\TextAt{13}{440.913}{549.642}{C}%
\TextAt{13}{450.131}{549.642}{l}%
\TextAt{21}{453.633}{547.848}{K}%
\TextAt{8}{464.935}{549.642}{and}%
\TextAt{9}{486.9}{549.642}{\Character{40}}%
\TextAt{18}{491.452}{549.642}{Z}%
\TextAt{9}{499.422}{549.642}{\Character{47}}%
\TextAt{13}{505.275}{549.642}{m}%
\TextAt{18}{515.515}{549.642}{Z}%
\TextAt{9}{523.485}{549.642}{\Character{41}}%
\TextAt{3}{528.037}{553.98}{\Character{3}}%
\TextAt{8}{79.226}{532.596}{are}%
\TextAt{8}{98.631}{532.596}{—}%
\TextAt{8}{114.126}{532.596}{with}%
\TextAt{8}{140.68}{532.596}{a}%
\TextAt{8}{150.329}{532.596}{grain}%
\TextAt{8}{180.147}{532.596}{of}%
\TextAt{8}{193.371}{532.596}{salt}%
\TextAt{8}{215.442}{532.596}{—}%
\TextAt{8}{230.937}{532.596}{only}%
\TextAt{8}{256.523}{532.596}{different}%
\TextAt{8}{302.587}{532.596}{instances}%
\TextAt{8}{352.681}{532.596}{of}%
\TextAt{8}{365.905}{532.596}{a}%
\TextAt{8}{375.554}{532.596}{common}%
\TextAt{8}{422.277}{532.596}{concept,}%
\TextAt{8}{468.363}{532.596}{that}%
\TextAt{8}{493.625}{532.596}{of}%
\TextAt{8}{506.849}{532.596}{a}%
\TextAt{8}{516.498}{532.596}{ray}%
\TextAt{8}{79.226}{515.55}{class}%
\TextAt{8}{106.663}{515.55}{group,}%
\TextAt{8}{143.091}{515.55}{which}%
\TextAt{8}{176.577}{515.55}{he}%
\TextAt{8}{192.179}{515.55}{defined}%
\TextAt{8}{232.491}{515.55}{in}%
\TextAt{8}{246.144}{515.55}{an}%
\TextAt{8}{262.403}{515.55}{ideal-theoretic}%
\TextAt{8}{339.143}{515.55}{way}%
\TextAt{8}{362.886}{515.55}{as}%
\TextAt{8}{377.257}{515.55}{the}%
\TextAt{8}{397.413}{515.55}{quotient}%
\TextAt{8}{443.596}{515.55}{group}%
\TextAt{13}{269.51}{491.617}{C}%
\TextAt{13}{278.728}{491.617}{l}%
\TextAt{27}{282.493}{496.553}{m}%
\TextAt{21}{282.23}{488.661}{K}%
\TextAt{9}{293.754}{491.617}{\Character{61}}%
\TextAt{13}{306.179}{491.617}{J}%
\TextAt{27}{313.757}{496.553}{m}%
\TextAt{21}{312.65}{488.661}{K}%
\TextAt{9}{320.852}{491.617}{\Character{47}}%
\TextAt{13}{326.705}{491.617}{P}%
\TextAt{27}{335.876}{496.553}{m}%
\TextAt{21}{334.251}{488.661}{K}%
\TextAt{8}{79.226}{467.684}{of}%
\TextAt{8}{93.103}{467.684}{all}%
\TextAt{8}{109.908}{467.684}{ideals}%
\TextAt{8}{143.033}{467.684}{relatively}%
\TextAt{8}{194.64}{467.684}{prime}%
\TextAt{8}{228.35}{467.684}{to}%
\TextAt{8}{243.207}{467.684}{a}%
\TextAt{8}{253.509}{467.684}{given}%
\TextAt{8}{284.625}{467.684}{module}%
\TextAt{14}{326.464}{467.684}{m}%
\TextAt{8}{335.631}{467.684}{,}%
\TextAt{8}{343.464}{467.684}{by}%
\TextAt{8}{360.269}{467.684}{the}%
\TextAt{8}{380.972}{467.684}{principal}%
\TextAt{8}{430.296}{467.684}{ideals}%
\TextAt{9}{463.42}{467.684}{\Character{40}}%
\TextAt{13}{467.972}{467.684}{\Character{11}}%
\TextAt{9}{475.494}{467.684}{\Character{41}}%
\TextAt{8}{484.49}{467.684}{with}%
\TextAt{13}{511.697}{467.684}{\Character{11}}%
\NativeMathAt{523.471}{467.684}{11.955}{\equiv}%
\TextAt{9}{79.226}{450.638}{1}%
\TextAt{8}{88.4}{450.638}{mod}%
\TextAt{14}{113.838}{450.638}{m}%
\TextAt{8}{123.005}{450.638}{,}%
\TextAt{8}{131.992}{450.638}{and}%
\TextAt{13}{156.224}{450.638}{\Character{11}}%
\TextAt{8}{169.114}{450.638}{totally}%
\TextAt{8}{207.992}{450.638}{positive.}%
\TextAt{8}{260.267}{450.638}{He}%
\TextAt{8}{279.61}{450.638}{conjectured}%
\TextAt{8}{343.833}{450.638}{that}%
\TextAt{8}{370.672}{450.638}{this}%
\TextAt{8}{394.965}{450.638}{group}%
\TextAt{13}{429.611}{450.638}{C}%
\TextAt{13}{438.829}{450.638}{l}%
\TextAt{27}{442.594}{454.977}{m}%
\TextAt{21}{442.331}{447.618}{K}%
\TextAt{8}{450.534}{450.638}{,}%
\TextAt{8}{459.521}{450.638}{along}%
\TextAt{8}{492.218}{450.638}{with}%
\TextAt{8}{520.348}{450.638}{its}%
\TextAt{8}{79.226}{433.592}{subgroups,}%
\TextAt{8}{137.849}{433.592}{would}%
\TextAt{8}{171.941}{433.592}{do}%
\TextAt{8}{188.148}{433.592}{the}%
\TextAt{8}{208.252}{433.592}{same}%
\TextAt{8}{237.527}{433.592}{for}%
\TextAt{8}{255.36}{433.592}{the}%
\TextAt{8}{275.464}{433.592}{subextensions}%
\TextAt{8}{349.738}{433.592}{of}%
\TextAt{8}{363.016}{433.592}{a}%
\TextAt{8}{372.719}{433.592}{“ray}%
\TextAt{8}{398.359}{433.592}{class}%
\TextAt{8}{425.744}{433.592}{field”}%
\TextAt{13}{456.573}{433.592}{K}%
\TextAt{27}{467.39}{437.931}{m}%
\TextAt{16}{477.321}{433.592}{\Character{106}}%
\TextAt{13}{483.962}{433.592}{K}%
\TextAt{8}{498.626}{433.592}{(which}%
\TextAt{8}{79.226}{416.546}{at}%
\TextAt{8}{92.723}{416.546}{first}%
\TextAt{8}{116.035}{416.546}{was}%
\TextAt{8}{137.72}{416.546}{only}%
\TextAt{8}{162.598}{416.546}{postulated}%
\TextAt{8}{219.408}{416.546}{to}%
\TextAt{8}{232.905}{416.546}{exist)}%
\TextAt{8}{264.346}{416.546}{as}%
\TextAt{8}{277.903}{416.546}{the}%
\TextAt{8}{297.245}{416.546}{ideal}%
\TextAt{8}{324.395}{416.546}{class}%
\TextAt{8}{351.018}{416.546}{group}%
\TextAt{13}{383.38}{416.546}{C}%
\TextAt{13}{392.598}{416.546}{l}%
\TextAt{21}{396.101}{414.753}{K}%
\TextAt{8}{407.387}{416.546}{and}%
\TextAt{8}{429.336}{416.546}{its}%
\TextAt{8}{444.841}{416.546}{subgroups}%
\TextAt{8}{499.439}{416.546}{did}%
\TextAt{8}{518.782}{416.546}{for}%
\TextAt{8}{79.226}{399.5}{the}%
\TextAt{8}{99.757}{399.5}{subfields}%
\TextAt{8}{147.724}{399.5}{of}%
\TextAt{8}{161.428}{399.5}{the}%
\TextAt{8}{181.958}{399.5}{Hilbert}%
\TextAt{8}{222.645}{399.5}{class}%
\TextAt{8}{250.456}{399.5}{field.}%
\TextAt{8}{281.486}{399.5}{Over}%
\TextAt{8}{310.481}{399.5}{and}%
\TextAt{8}{333.617}{399.5}{above}%
\TextAt{8}{367.167}{399.5}{this,}%
\TextAt{8}{393.709}{399.5}{he}%
\TextAt{8}{409.684}{399.5}{hypothesized}%
\TextAt{8}{480.032}{399.5}{that}%
\TextAt{8}{505.775}{399.5}{every}%
\TextAt{8}{79.226}{382.454}{abelian}%
\TextAt{8}{119.775}{382.454}{extension}%
\TextAt{8}{171.443}{382.454}{ought}%
\TextAt{8}{204.197}{382.454}{to}%
\TextAt{8}{218.409}{382.454}{be}%
\TextAt{8}{234.235}{382.454}{captured}%
\TextAt{8}{282.914}{382.454}{by}%
\TextAt{8}{299.075}{382.454}{such}%
\TextAt{8}{325.373}{382.454}{a}%
\TextAt{8}{335.03}{382.454}{ray}%
\TextAt{8}{355.1}{382.454}{class}%
\TextAt{8}{382.439}{382.454}{field,}%
\TextAt{8}{410.969}{382.454}{as}%
\TextAt{8}{425.241}{382.454}{was}%
\TextAt{8}{447.642}{382.454}{suggested}%
\TextAt{8}{500.35}{382.454}{by}%
\TextAt{8}{516.51}{382.454}{the}%
\TextAt{8}{79.226}{365.408}{case}%
\TextAt{8}{104.479}{365.408}{where}%
\TextAt{8}{138.77}{365.408}{the}%
\TextAt{8}{159.408}{365.408}{base}%
\TextAt{8}{185.964}{365.408}{field}%
\TextAt{8}{211.803}{365.408}{is}%
\TextAt{18}{224.048}{365.408}{Q}%
\TextAt{8}{233.347}{365.408}{,}%
\TextAt{8}{241.098}{365.408}{whose}%
\TextAt{8}{276.107}{365.408}{abelian}%
\TextAt{8}{317.236}{365.408}{extensions}%
\TextAt{8}{374.098}{365.408}{are}%
\TextAt{8}{394.091}{365.408}{all}%
\TextAt{8}{410.832}{365.408}{contained}%
\TextAt{8}{464.323}{365.408}{in}%
\TextAt{8}{478.457}{365.408}{cyclotomic}%
\TextAt{8}{79.226}{348.362}{fields}%
\TextAt{18}{108.028}{348.362}{Q}%
\TextAt{9}{117.327}{348.362}{\Character{40}}%
\TextAt{13}{121.879}{348.362}{\Character{22}}%
\TextAt{21}{128.922}{346.569}{m}%
\TextAt{9}{136.911}{348.362}{\Character{41}}%
\TextAt{16}{144.784}{348.362}{\Character{106}}%
\TextAt{18}{151.426}{348.362}{Q}%
\TextAt{8}{163.452}{348.362}{by}%
\TextAt{8}{178.542}{348.362}{the}%
\TextAt{8}{197.529}{348.362}{Kronecker–Weber}%
\TextAt{8}{290.651}{348.362}{theorem.}%
\TextAt{8}{340.337}{348.362}{After}%
\TextAt{8}{369.725}{348.362}{the}%
\TextAt{8}{388.712}{348.362}{seminal}%
\TextAt{8}{429.876}{348.362}{work}%
\TextAt{8}{457.328}{348.362}{of}%
\TextAt{8}{469.488}{348.362}{the}%
\TextAt{8}{488.476}{348.362}{Austrian}%
\TextAt{8}{79.226}{331.316}{mathematician}%
\TextAt{26}{160.503}{331.316}{Philipp}%
\TextAt{26}{209.394}{331.316}{Furtwängler}%
\TextAt{8}{295.13}{331.316}{[44],}%
\TextAt{8}{322.09}{331.316}{these}%
\TextAt{8}{353.335}{331.316}{conjectures}%
\TextAt{8}{415.472}{331.316}{were}%
\TextAt{8}{443.728}{331.316}{confirmed}%
\TextAt{8}{498.978}{331.316}{by}%
\TextAt{8}{516.51}{331.316}{the}%
\TextAt{8}{79.226}{314.27}{Japanese}%
\TextAt{8}{130.507}{314.27}{arithmetician}%
\TextAt{26}{204.335}{314.27}{Teiji}%
\TextAt{26}{239.104}{314.27}{Takagi}%
\TextAt{8}{286.248}{314.27}{(1875–1960),}%
\TextAt{8}{357.268}{314.27}{and}%
\TextAt{8}{381.662}{314.27}{cast}%
\TextAt{8}{407.419}{314.27}{by}%
\TextAt{26}{425.309}{314.27}{Emil}%
\TextAt{26}{459.264}{314.27}{Artin}%
\TextAt{8}{498.924}{314.27}{(1898–}%
\TextAt{8}{79.226}{297.224}{1962)}%
\TextAt{8}{111.111}{297.224}{into}%
\TextAt{8}{134.854}{297.224}{a}%
\TextAt{8}{144.609}{297.224}{definite,}%
\TextAt{8}{189.477}{297.224}{canonical}%
\TextAt{8}{240.86}{297.224}{form.}%
\TextAt{8}{97.159}{277.19}{The}%
\TextAt{8}{122.752}{277.19}{idèle-theoretic}%
\TextAt{8}{200.374}{277.19}{language}%
\TextAt{8}{250.702}{277.19}{introduced}%
\TextAt{8}{310.774}{277.19}{by}%
\TextAt{26}{328.572}{277.19}{Chevalley}%
\TextAt{8}{397.264}{277.19}{brought}%
\TextAt{8}{442.714}{277.19}{the}%
\TextAt{8}{464.41}{277.19}{simplification}%
\TextAt{8}{79.226}{260.144}{that}%
\TextAt{8}{103.426}{260.144}{the}%
\TextAt{8}{122.414}{260.144}{idèle}%
\TextAt{8}{148.551}{260.144}{class}%
\TextAt{8}{174.819}{260.144}{group}%
\TextAt{13}{206.826}{260.144}{C}%
\TextAt{21}{215.201}{258.35}{K}%
\TextAt{8}{226.133}{260.144}{encapsulated}%
\TextAt{8}{294.615}{260.144}{all}%
\TextAt{8}{309.705}{260.144}{abelian}%
\TextAt{8}{349.184}{260.144}{extensions}%
\TextAt{8}{404.396}{260.144}{of}%
\TextAt{13}{416.557}{260.144}{L}%
\TextAt{16}{427.843}{260.144}{\Character{106}}%
\TextAt{13}{434.484}{260.144}{K}%
\TextAt{8}{448.031}{260.144}{at}%
\TextAt{8}{461.172}{260.144}{once,}%
\TextAt{8}{490.149}{260.144}{avoiding}%
\TextAt{8}{79.226}{243.098}{choosing}%
\TextAt{8}{127.132}{243.098}{a}%
\TextAt{8}{137.248}{243.098}{module}%
\TextAt{14}{178.901}{243.098}{m}%
\TextAt{8}{192.325}{243.098}{every}%
\TextAt{8}{223.577}{243.098}{time}%
\TextAt{8}{250.597}{243.098}{such}%
\TextAt{8}{277.355}{243.098}{an}%
\TextAt{8}{293.974}{243.098}{extension}%
\TextAt{8}{346.1}{243.098}{was}%
\TextAt{8}{368.959}{243.098}{given,}%
\TextAt{8}{403.231}{243.098}{in}%
\TextAt{8}{417.244}{243.098}{order}%
\TextAt{8}{448.173}{243.098}{to}%
\TextAt{8}{462.844}{243.098}{accommodate}%
\TextAt{8}{79.226}{226.052}{it}%
\TextAt{8}{91.898}{226.052}{into}%
\TextAt{8}{116.608}{226.052}{the}%
\TextAt{8}{137.732}{226.052}{ray}%
\TextAt{8}{158.868}{226.052}{class}%
\TextAt{8}{187.273}{226.052}{field}%
\TextAt{13}{213.598}{226.052}{K}%
\TextAt{27}{224.415}{230.39}{m}%
\TextAt{16}{235.994}{226.052}{\Character{106}}%
\TextAt{13}{244.285}{226.052}{K}%
\TextAt{8}{255.103}{226.052}{,}%
\TextAt{8}{263.461}{226.052}{and}%
\TextAt{8}{287.192}{226.052}{thereby}%
\TextAt{8}{330.433}{226.052}{make}%
\TextAt{8}{361.97}{226.052}{it}%
\TextAt{8}{374.641}{226.052}{amenable}%
\TextAt{8}{427.638}{226.052}{to}%
\TextAt{8}{442.916}{226.052}{class}%
\TextAt{8}{471.321}{226.052}{field}%
\TextAt{8}{497.645}{226.052}{theory.}%
\TextAt{8}{79.226}{209.006}{The}%
\TextAt{8}{104.095}{209.006}{classical}%
\TextAt{8}{149.91}{209.006}{point}%
\TextAt{8}{181.295}{209.006}{of}%
\TextAt{8}{195.441}{209.006}{view}%
\TextAt{8}{223.239}{209.006}{can}%
\TextAt{8}{245.514}{209.006}{be}%
\TextAt{8}{262.254}{209.006}{vindicated}%
\TextAt{8}{319.976}{209.006}{in}%
\TextAt{8}{334.444}{209.006}{terms}%
\TextAt{8}{367.838}{209.006}{of}%
\TextAt{8}{381.983}{209.006}{the}%
\TextAt{8}{402.955}{209.006}{idèle-theoretic}%
\TextAt{8}{479.853}{209.006}{version}%
\TextAt{8}{520.408}{209.006}{by}%
\TextAt{8}{79.226}{191.96}{looking}%
\TextAt{8}{121.29}{191.96}{at}%
\TextAt{8}{136.682}{191.96}{congruence}%
\TextAt{8}{198.244}{191.96}{subgroups}%
\TextAt{13}{254.738}{191.96}{C}%
\TextAt{27}{263.971}{196.298}{m}%
\TextAt{21}{263.113}{188.939}{K}%
\TextAt{8}{276.294}{191.96}{in}%
\TextAt{13}{291.029}{191.96}{C}%
\TextAt{21}{299.404}{190.166}{K}%
\TextAt{8}{307.606}{191.96}{,}%
\TextAt{8}{316.107}{191.96}{which}%
\TextAt{8}{350.675}{191.96}{define}%
\TextAt{8}{385.566}{191.96}{the}%
\TextAt{8}{406.803}{191.96}{ray}%
\TextAt{8}{428.053}{191.96}{class}%
\TextAt{8}{456.572}{191.96}{fields}%
\TextAt{13}{487.625}{191.96}{K}%
\TextAt{27}{498.442}{196.298}{m}%
\TextAt{16}{510.215}{191.96}{\Character{106}}%
\TextAt{13}{518.7}{191.96}{K}%
\TextAt{8}{529.518}{191.96}{.}%
\TextAt{8}{79.226}{174.914}{Their}%
\TextAt{8}{110.491}{174.914}{subfields}%
\TextAt{8}{157.489}{174.914}{correspond,}%
\TextAt{8}{219.837}{174.914}{according}%
\TextAt{8}{271.927}{174.914}{to}%
\TextAt{8}{285.642}{174.914}{the}%
\TextAt{8}{305.203}{174.914}{new}%
\TextAt{8}{328.661}{174.914}{point}%
\TextAt{8}{358.635}{174.914}{of}%
\TextAt{8}{371.369}{174.914}{view,}%
\TextAt{8}{401.127}{174.914}{to}%
\TextAt{8}{414.842}{174.914}{the}%
\TextAt{8}{434.402}{174.914}{groups}%
\TextAt{8}{471.597}{174.914}{between}%
\TextAt{13}{516.192}{174.914}{C}%
\TextAt{27}{525.425}{179.252}{m}%
\TextAt{21}{524.567}{171.893}{K}%
\TextAt{8}{79.226}{157.868}{and}%
\TextAt{13}{101.989}{157.868}{C}%
\TextAt{21}{110.364}{156.075}{K}%
\TextAt{8}{118.567}{157.868}{,}%
\TextAt{8}{125.716}{157.868}{and}%
\TextAt{8}{148.478}{157.868}{hence,}%
\TextAt{8}{184.236}{157.868}{in}%
\TextAt{8}{197.889}{157.868}{view}%
\TextAt{8}{224.872}{157.868}{of}%
\TextAt{8}{238.202}{157.868}{the}%
\TextAt{8}{258.358}{157.868}{isomorphism}%
\TextAt{13}{266.436}{133.935}{C}%
\TextAt{21}{274.811}{132.142}{K}%
\TextAt{9}{283.014}{133.935}{\Character{47}}%
\TextAt{13}{288.867}{133.935}{C}%
\TextAt{27}{298.1}{138.871}{m}%
\TextAt{21}{297.242}{130.979}{K}%
\NativeMathAt{308.765}{133.935}{11.955}{\cong}%
\TextAt{13}{321.385}{133.935}{C}%
\TextAt{13}{330.603}{133.935}{l}%
\TextAt{27}{334.368}{138.871}{m}%
\TextAt{21}{334.105}{130.979}{K}%
\TextAt{13}{342.308}{133.935}{\Character{59}}%
\TextAt{8}{79.226}{110.002}{to}%
\TextAt{8}{93.536}{110.002}{the}%
\TextAt{8}{113.693}{110.002}{subgroups}%
\TextAt{8}{169.105}{110.002}{of}%
\TextAt{8}{182.435}{110.002}{the}%
\TextAt{8}{202.592}{110.002}{ray}%
\TextAt{8}{222.76}{110.002}{class}%
\TextAt{8}{250.197}{110.002}{group}%
\TextAt{13}{283.373}{110.002}{C}%
\TextAt{13}{292.591}{110.002}{l}%
\TextAt{27}{296.356}{114.341}{m}%
\TextAt{21}{296.093}{106.982}{K}%
\TextAt{8}{304.296}{110.002}{.}%
\TextAt{8}{97.159}{89.967}{In}%
\TextAt{8}{111.354}{89.967}{what}%
\TextAt{8}{140.194}{89.967}{follows,}%
\TextAt{8}{181.541}{89.967}{we}%
\TextAt{8}{198.343}{89.967}{want}%
\TextAt{8}{226.537}{89.967}{to}%
\TextAt{8}{240.422}{89.967}{deduce}%
\TextAt{8}{279.005}{89.967}{the}%
\TextAt{8}{298.736}{89.967}{classical,}%
\TextAt{8}{346.646}{89.967}{ideal-theoretic}%
\TextAt{8}{422.96}{89.967}{version}%
\TextAt{8}{462.273}{89.967}{of}%
\TextAt{8}{475.177}{89.967}{global}%
\TextAt{8}{509.23}{89.967}{class}%
\TextAt{8}{79.226}{72.921}{field}%
\TextAt{8}{104.864}{72.921}{theory}%
\TextAt{8}{141.895}{72.921}{from}%
\TextAt{8}{169.816}{72.921}{the}%
\TextAt{8}{190.253}{72.921}{idèle-theoretic}%
\TextAt{8}{266.617}{72.921}{one.}%
\TextAt{8}{293.47}{72.921}{This}%
\TextAt{8}{320.47}{72.921}{is}%
\TextAt{8}{332.515}{72.921}{not}%
\TextAt{8}{353.61}{72.921}{only}%
\TextAt{8}{379.582}{72.921}{an}%
\TextAt{8}{396.122}{72.921}{obligation}%
\TextAt{8}{451.05}{72.921}{towards}%
\TextAt{8}{494.979}{72.921}{history,}%
\TextAt{8}{297.211}{43.033}{370}%
\Destination{78.23}{809.815}{nav.page.370}%
\Destination{79.23}{773.075}{nav.section.6.7}%
\Destination{79.36}{380.455}{nav.print.406}%
\Destination{79.23}{773.075}{outline.62}%
\Hotlink{294.63}{329.155}{19.22}{11.047}{nav.cite.N044}%
}

% Page 383 (body page 371)
\Page{%
\TextAt{2}{79.226}{796.943}{VI.}%
\TextAt{2}{98.012}{796.943}{Global}%
\TextAt{2}{133.248}{796.943}{Class}%
\TextAt{2}{161.841}{796.943}{Field}%
\TextAt{2}{189.572}{796.943}{Theory}%
\Vector{q 0 G 0 g 0.299 w 7.226 18.729 m 460.769 18.729 l S Q}%
\TextAt{8}{79.226}{761.124}{but}%
\TextAt{8}{101.003}{761.124}{a}%
\TextAt{8}{111.076}{761.124}{factual}%
\TextAt{8}{150.092}{761.124}{necessity}%
\TextAt{8}{199.305}{761.124}{that}%
\TextAt{8}{224.992}{761.124}{is}%
\TextAt{8}{237.073}{761.124}{forced}%
\TextAt{8}{272.179}{761.124}{upon}%
\TextAt{8}{302.086}{761.124}{us}%
\TextAt{8}{317.419}{761.124}{by}%
\TextAt{8}{333.995}{761.124}{the}%
\TextAt{8}{354.468}{761.124}{numerous}%
\TextAt{8}{407.855}{761.124}{applications}%
\TextAt{8}{473.28}{761.124}{of}%
\TextAt{8}{486.927}{761.124}{the}%
\TextAt{8}{507.401}{761.124}{more}%
\TextAt{8}{79.226}{744.078}{elementary}%
\TextAt{8}{139.062}{744.078}{and}%
\TextAt{8}{161.824}{744.078}{more}%
\TextAt{8}{191.091}{744.078}{immediately}%
\TextAt{8}{257.753}{744.078}{accessible}%
\TextAt{8}{310.547}{744.078}{ideal}%
\TextAt{8}{338.51}{744.078}{groups.}%
\TextAt{8}{97.159}{724.043}{Let}%
\TextAt{13}{118.426}{724.043}{L}%
\TextAt{16}{130.218}{724.043}{\Character{106}}%
\TextAt{13}{137.367}{724.043}{K}%
\TextAt{8}{152.38}{724.043}{be}%
\TextAt{8}{168.602}{724.043}{an}%
\TextAt{8}{185.158}{724.043}{abelian}%
\TextAt{8}{226.103}{724.043}{extension,}%
\TextAt{8}{281.493}{724.043}{and}%
\TextAt{8}{304.553}{724.043}{let}%
\TextAt{14}{321.755}{724.043}{p}%
\TextAt{8}{331.933}{724.043}{be}%
\TextAt{8}{348.155}{724.043}{an}%
\TextAt{8}{364.712}{724.043}{unramified}%
\TextAt{8}{423.542}{724.043}{prime}%
\TextAt{8}{457.003}{724.043}{ideal}%
\TextAt{8}{485.264}{724.043}{of}%
\TextAt{13}{498.892}{724.043}{K}%
\TextAt{8}{513.904}{724.043}{and}%
\TextAt{14}{79.226}{706.997}{P}%
\TextAt{8}{93.024}{706.997}{a}%
\TextAt{8}{102.786}{706.997}{prime}%
\TextAt{8}{135.956}{706.997}{ideal}%
\TextAt{8}{163.926}{706.997}{of}%
\TextAt{13}{177.263}{706.997}{L}%
\TextAt{8}{189.132}{706.997}{lying}%
\TextAt{8}{218.082}{706.997}{above}%
\TextAt{14}{251.265}{706.997}{p}%
\TextAt{8}{257.248}{706.997}{.}%
\TextAt{8}{265.716}{706.997}{The}%
\TextAt{8}{289.777}{706.997}{decomposition}%
\TextAt{8}{366.919}{706.997}{group}%
\TextAt{13}{400.101}{706.997}{G}%
\TextAt{9}{409.335}{706.997}{\Character{40}}%
\TextAt{13}{413.887}{706.997}{L}%
\TextAt{27}{421.851}{705.204}{P}%
\TextAt{16}{432.277}{706.997}{\Character{106}}%
\TextAt{13}{438.931}{706.997}{K}%
\TextAt{27}{448.89}{705.204}{p}%
\TextAt{9}{453.377}{706.997}{\Character{41}}%
\NativeMathAt{461.262}{706.997}{11.955}{\subseteq}%
\TextAt{13}{473.893}{706.997}{G}%
\TextAt{9}{483.126}{706.997}{\Character{40}}%
\TextAt{13}{487.678}{706.997}{L}%
\TextAt{16}{498.975}{706.997}{\Character{106}}%
\TextAt{13}{505.629}{706.997}{K}%
\TextAt{9}{516.446}{706.997}{\Character{41}}%
\TextAt{8}{524.903}{706.997}{is}%
\TextAt{8}{79.226}{689.951}{then}%
\TextAt{8}{105.886}{689.951}{generated}%
\TextAt{8}{159.218}{689.951}{by}%
\TextAt{8}{175.477}{689.951}{the}%
\TextAt{8}{195.634}{689.951}{classical}%
\TextAt{12}{240.633}{689.951}{Frobenius}%
\TextAt{12}{302.644}{689.951}{automorphism}%
\TextAt{13}{257.95}{666.949}{\Character{39}}%
\TextAt{27}{265.618}{665.155}{p}%
\TextAt{9}{273.426}{666.949}{\Character{61}}%
\TextAt{9}{285.836}{666.949}{\Character{40}}%
\TextAt{13}{290.404}{666.949}{\Character{25}}%
\TextAt{27}{297.044}{665.155}{p}%
\TextAt{13}{301.531}{666.949}{\Character{59}}%
\TextAt{13}{306.755}{666.949}{L}%
\TextAt{27}{314.74}{665.155}{P}%
\TextAt{16}{325.154}{666.949}{\Character{106}}%
\TextAt{13}{331.796}{666.949}{K}%
\TextAt{27}{341.755}{665.155}{p}%
\TextAt{9}{346.242}{666.949}{\Character{41}}%
\TextAt{13}{350.794}{666.949}{\Character{59}}%
\TextAt{8}{79.226}{643.946}{where}%
\TextAt{13}{111.99}{643.946}{\Character{25}}%
\TextAt{27}{118.631}{642.153}{p}%
\TextAt{8}{125.97}{643.946}{is}%
\TextAt{8}{136.689}{643.946}{a}%
\TextAt{8}{145.4}{643.946}{prime}%
\TextAt{8}{177.518}{643.946}{element}%
\TextAt{8}{219.715}{643.946}{of}%
\TextAt{13}{232}{643.946}{K}%
\TextAt{27}{241.96}{642.153}{p}%
\TextAt{8}{246.446}{643.946}{.}%
\TextAt{8}{254.546}{643.946}{As}%
\TextAt{8}{270.789}{643.946}{an}%
\TextAt{8}{286.003}{643.946}{automorphism}%
\TextAt{8}{362.427}{643.946}{of}%
\TextAt{13}{374.713}{643.946}{L}%
\TextAt{8}{382.677}{643.946}{,}%
\TextAt{13}{388.99}{643.946}{\Character{39}}%
\TextAt{27}{396.659}{642.153}{p}%
\TextAt{8}{403.998}{643.946}{is}%
\TextAt{8}{414.717}{643.946}{obviously}%
\TextAt{8}{465.45}{643.946}{characterized}%
\TextAt{8}{79.226}{626.9}{by}%
\TextAt{8}{95.485}{626.9}{the}%
\TextAt{8}{115.642}{626.9}{congruence}%
\TextAt{13}{222.881}{609.854}{\Character{39}}%
\TextAt{27}{230.55}{608.061}{p}%
\TextAt{13}{235.037}{609.854}{a}%
\NativeMathAt{244.502}{609.854}{11.955}{\equiv}%
\TextAt{13}{257.122}{609.854}{a}%
\TextAt{21}{263.267}{614.79}{q}%
\TextAt{8}{271.152}{609.854}{mod}%
\TextAt{14}{296.59}{609.854}{P}%
\TextAt{8}{318.439}{609.854}{for}%
\TextAt{8}{336.324}{609.854}{all}%
\TextAt{13}{352.583}{609.854}{a}%
\NativeMathAt{362.048}{609.854}{11.955}{\in}%
\TextAt{16}{373.322}{609.854}{\Character{79}}%
\TextAt{21}{382.857}{608.061}{L}%
\TextAt{8}{79.226}{588.903}{where}%
\TextAt{13}{113.035}{588.903}{q}%
\TextAt{8}{122.552}{588.903}{is}%
\TextAt{8}{134.316}{588.903}{the}%
\TextAt{8}{154.472}{588.903}{number}%
\TextAt{8}{197.068}{588.903}{of}%
\TextAt{8}{210.399}{588.903}{elements}%
\TextAt{8}{258.255}{588.903}{in}%
\TextAt{8}{271.908}{588.903}{the}%
\TextAt{8}{292.064}{588.903}{residue}%
\TextAt{8}{331.791}{588.903}{class}%
\TextAt{8}{359.228}{588.903}{field}%
\TextAt{8}{384.585}{588.903}{of}%
\TextAt{14}{397.915}{588.903}{p}%
\TextAt{8}{403.898}{588.903}{.}%
\TextAt{8}{412.347}{588.903}{We}%
\TextAt{8}{432.491}{588.903}{put}%
\TextAt{13}{263.431}{562.291}{\Character{39}}%
\TextAt{27}{271.099}{560.498}{p}%
\TextAt{9}{278.907}{562.291}{\Character{61}\Character{58}}%
\TextAt{20}{294.584}{579.148}{\Character{18}}%
\TextAt{13}{304.58}{570.379}{L}%
\TextAt{16}{315.865}{570.379}{\Character{106}}%
\TextAt{13}{322.507}{570.379}{K}%
\Vector{q 0 G 0 g 0.478 w 232.58 -208.72 m 261.325 -208.72 l S Q}%
\TextAt{14}{315.961}{554.09}{p}%
\TextAt{20}{334.52}{579.148}{\Character{19}}%
\TextAt{13}{345.313}{562.291}{\Character{58}}%
\TextAt{8}{79.226}{535.564}{Now}%
\TextAt{8}{104.955}{535.564}{let}%
\TextAt{14}{120.928}{535.564}{m}%
\TextAt{8}{133.06}{535.564}{be}%
\TextAt{8}{148.053}{535.564}{a}%
\TextAt{8}{156.877}{535.564}{module}%
\TextAt{8}{197.239}{535.564}{of}%
\TextAt{13}{209.637}{535.564}{K}%
\TextAt{8}{223.421}{535.564}{such}%
\TextAt{8}{248.886}{535.564}{that}%
\TextAt{13}{273.324}{535.564}{L}%
\TextAt{8}{284.254}{535.564}{lies}%
\TextAt{8}{303.539}{535.564}{in}%
\TextAt{8}{316.26}{535.564}{the}%
\TextAt{8}{335.485}{535.564}{ray}%
\TextAt{8}{354.722}{535.564}{class}%
\TextAt{8}{381.227}{535.564}{field}%
\TextAt{8}{405.653}{535.564}{mod}%
\TextAt{14}{431.058}{535.564}{m}%
\TextAt{8}{440.225}{535.564}{.}%
\TextAt{8}{448.363}{535.564}{Such}%
\TextAt{8}{475.717}{535.564}{a}%
\TextAt{8}{484.541}{535.564}{module}%
\TextAt{8}{524.903}{535.564}{is}%
\TextAt{8}{79.226}{518.518}{called}%
\TextAt{8}{111.696}{518.518}{an}%
\TextAt{12}{127.261}{518.518}{module}%
\TextAt{12}{174.08}{518.518}{of}%
\TextAt{12}{188.609}{518.518}{definition}%
\TextAt{8}{247.284}{518.518}{for}%
\TextAt{13}{264.475}{518.518}{L}%
\TextAt{8}{272.439}{518.518}{.}%
\TextAt{8}{280.657}{518.518}{Since}%
\TextAt{8}{310.52}{518.518}{by}%
\TextAt{8}{326.085}{518.518}{(6.6)}%
\TextAt{8}{353.366}{518.518}{each}%
\TextAt{8}{379.01}{518.518}{prime}%
\TextAt{8}{411.479}{518.518}{ideal}%
\TextAt{14}{438.749}{518.518}{p}%
\TextAt{18}{448.052}{518.518}{\Character{45}}%
\TextAt{14}{454.694}{518.518}{m}%
\TextAt{8}{467.064}{518.518}{is}%
\TextAt{8}{478.134}{518.518}{unramified}%
\TextAt{8}{79.226}{501.472}{in}%
\TextAt{13}{92.879}{501.472}{L}%
\TextAt{8}{100.843}{501.472}{,}%
\TextAt{8}{107.993}{501.472}{we}%
\TextAt{8}{125.22}{501.472}{get}%
\TextAt{8}{144.731}{501.472}{a}%
\TextAt{8}{154.486}{501.472}{canonical}%
\TextAt{8}{205.87}{501.472}{homomorphism}%
\TextAt{20}{232.854}{491.59}{\Character{18}}%
\TextAt{13}{242.849}{482.821}{L}%
\TextAt{16}{254.135}{482.821}{\Character{106}}%
\TextAt{13}{260.776}{482.821}{K}%
\Vector{q 0 G 0 g 0.478 w 170.849 -296.278 m 199.594 -296.278 l S Q}%
\TextAt{20}{272.79}{491.59}{\Character{19}}%
\TextAt{9}{284.911}{474.733}{\Character{58}}%
\TextAt{13}{291.483}{474.733}{J}%
\TextAt{27}{299.06}{479.67}{m}%
\TextAt{21}{297.954}{471.778}{K}%
\NativeArrowAt{309.477}{474.733}{11.955}{19.26}{\longrightarrow}%
\TextAt{13}{332.059}{474.733}{G}%
\TextAt{9}{341.293}{474.733}{\Character{40}}%
\TextAt{13}{345.845}{474.733}{L}%
\TextAt{16}{357.13}{474.733}{\Character{106}}%
\TextAt{13}{363.772}{474.733}{K}%
\TextAt{9}{374.59}{474.733}{\Character{41}}%
\TextAt{8}{79.226}{448.007}{from}%
\TextAt{8}{106.973}{448.007}{the}%
\TextAt{8}{127.237}{448.007}{group}%
\TextAt{13}{160.519}{448.007}{J}%
\TextAt{27}{168.096}{452.345}{m}%
\TextAt{21}{166.989}{444.986}{K}%
\TextAt{8}{179.196}{448.007}{of}%
\TextAt{8}{192.633}{448.007}{all}%
\TextAt{8}{208.999}{448.007}{ideals}%
\TextAt{8}{241.683}{448.007}{of}%
\TextAt{13}{255.12}{448.007}{K}%
\TextAt{8}{269.942}{448.007}{which}%
\TextAt{8}{303.535}{448.007}{are}%
\TextAt{8}{323.153}{448.007}{relatively}%
\TextAt{8}{374.32}{448.007}{prime}%
\TextAt{8}{407.591}{448.007}{to}%
\TextAt{14}{422.008}{448.007}{m}%
\TextAt{8}{435.179}{448.007}{by}%
\TextAt{8}{451.545}{448.007}{putting,}%
\TextAt{8}{496.558}{448.007}{for}%
\TextAt{8}{514.55}{448.007}{any}%
\TextAt{8}{79.226}{430.961}{ideal}%
\TextAt{14}{107.189}{430.961}{a}%
\TextAt{9}{116.493}{430.961}{\Character{61}}%
\TextAt{20}{128.919}{439.927}{\Character{81}}%
\TextAt{27}{140.21}{427.474}{p}%
\TextAt{14}{146.689}{430.961}{p}%
\TextAt{21}{152.672}{435.299}{\Character{23}}%
\TextAt{28}{156.835}{434.303}{p}%
\TextAt{8}{161.416}{430.961}{:}%
\TextAt{20}{233.76}{424.216}{\Character{18}}%
\TextAt{13}{243.756}{415.447}{L}%
\TextAt{16}{255.041}{415.447}{\Character{106}}%
\TextAt{13}{261.683}{415.447}{K}%
\Vector{q 0 G 0 g 0.478 w 171.756 -363.652 m 200.501 -363.652 l S Q}%
\TextAt{14}{255.137}{399.158}{a}%
\TextAt{20}{273.696}{424.216}{\Character{19}}%
\TextAt{9}{285.817}{407.359}{\Character{61}}%
\TextAt{20}{298.243}{418.716}{\Character{89}}%
\TextAt{27}{303.886}{393.61}{p}%
\TextAt{20}{315.511}{424.216}{\Character{18}}%
\TextAt{13}{325.507}{415.447}{L}%
\TextAt{16}{336.793}{415.447}{\Character{106}}%
\TextAt{13}{343.434}{415.447}{K}%
\Vector{q 0 G 0 g 0.478 w 253.507 -363.652 m 282.252 -363.652 l S Q}%
\TextAt{14}{336.888}{399.158}{p}%
\TextAt{20}{355.448}{424.216}{\Character{19}}%
\TextAt{21}{364.248}{421.539}{\Character{23}}%
\TextAt{28}{368.411}{420.543}{p}%
\TextAt{13}{374.984}{407.359}{\Character{58}}%
\TextAt{20}{79.226}{385.529}{\Character{16}}%
\TextAt{21}{87.562}{377.957}{L}%
\TextAt{3}{93.321}{377.957}{\Character{106}}%
\TextAt{21}{95.673}{377.957}{K}%
\Vector{q 0 G 0 g 0.478 w 15.562 -398.753 m 31.377 -398.753 l S Q}%
\TextAt{27}{93.475}{368.136}{a}%
\TextAt{20}{104.573}{385.529}{\Character{17}}%
\TextAt{8}{115.524}{372.258}{is}%
\TextAt{8}{127.202}{372.258}{called}%
\TextAt{8}{160.279}{372.258}{the}%
\TextAt{12}{180.35}{372.258}{Artin}%
\TextAt{12}{216.845}{372.258}{symbol}%
\TextAt{8}{258.413}{372.258}{.}%
\TextAt{8}{266.833}{372.258}{If}%
\TextAt{14}{278.439}{372.258}{p}%
\NativeMathAt{287.743}{372.258}{11.955}{\in}%
\TextAt{13}{299.034}{372.258}{J}%
\TextAt{27}{306.611}{376.597}{m}%
\TextAt{21}{305.504}{369.238}{K}%
\TextAt{8}{317.518}{372.258}{is}%
\TextAt{8}{329.196}{372.258}{a}%
\TextAt{8}{338.865}{372.258}{prime}%
\TextAt{8}{371.942}{372.258}{ideal}%
\TextAt{8}{399.819}{372.258}{and}%
\TextAt{13}{422.496}{372.258}{\Character{25}}%
\TextAt{27}{429.136}{370.465}{p}%
\TextAt{8}{437.434}{372.258}{a}%
\TextAt{8}{447.104}{372.258}{prime}%
\TextAt{8}{480.181}{372.258}{element}%
\TextAt{8}{523.337}{372.258}{of}%
\TextAt{13}{79.226}{355.212}{K}%
\TextAt{27}{89.185}{353.419}{p}%
\TextAt{8}{93.672}{355.212}{,}%
\TextAt{8}{100.821}{355.212}{then}%
\TextAt{8}{127.481}{355.212}{clearly}%
\TextAt{20}{240.371}{350.438}{\Character{18}}%
\TextAt{13}{250.367}{341.669}{L}%
\TextAt{16}{261.652}{341.669}{\Character{106}}%
\TextAt{13}{268.294}{341.669}{K}%
\Vector{q 0 G 0 g 0.478 w 178.367 -437.43 m 207.112 -437.43 l S Q}%
\TextAt{14}{261.748}{325.38}{p}%
\TextAt{20}{280.307}{350.438}{\Character{19}}%
\TextAt{9}{292.428}{333.581}{\Character{61}}%
\TextAt{9}{304.838}{333.581}{\Character{40}}%
\TextAt{16}{309.406}{333.581}{\Character{104}}%
\TextAt{13}{314.055}{333.581}{\Character{25}}%
\TextAt{27}{320.696}{331.788}{p}%
\TextAt{16}{325.183}{333.581}{\Character{105}}%
\TextAt{13}{329.832}{333.581}{\Character{59}}%
\TextAt{13}{335.056}{333.581}{L}%
\TextAt{16}{346.361}{333.581}{\Character{106}}%
\TextAt{13}{353.003}{333.581}{K}%
\TextAt{9}{363.821}{333.581}{\Character{41}}%
\TextAt{13}{368.373}{333.581}{\Character{59}}%
\TextAt{8}{79.226}{308.356}{if}%
\TextAt{16}{89.95}{308.356}{\Character{104}}%
\TextAt{13}{94.599}{308.356}{\Character{25}}%
\TextAt{27}{101.24}{306.563}{p}%
\TextAt{16}{105.726}{308.356}{\Character{105}}%
\NativeMathAt{113.676}{308.356}{11.955}{\in}%
\TextAt{13}{124.987}{308.356}{C}%
\TextAt{21}{133.362}{306.563}{K}%
\TextAt{8}{145.463}{308.356}{denotes}%
\TextAt{8}{187.796}{308.356}{the}%
\TextAt{8}{207.952}{308.356}{class}%
\TextAt{8}{235.389}{308.356}{of}%
\TextAt{8}{248.719}{308.356}{the}%
\TextAt{8}{268.876}{308.356}{idèle}%
\TextAt{9}{296.181}{308.356}{\Character{40}}%
\TextAt{13}{300.734}{308.356}{\Character{58}}%
\TextAt{13}{305.958}{308.356}{\Character{58}}%
\TextAt{13}{311.195}{308.356}{\Character{58}}%
\TextAt{13}{316.419}{308.356}{\Character{59}}%
\TextAt{9}{321.71}{308.356}{1}%
\TextAt{13}{327.563}{308.356}{\Character{59}}%
\TextAt{9}{332.807}{308.356}{1}%
\TextAt{13}{338.66}{308.356}{\Character{59}}%
\TextAt{13}{343.884}{308.356}{\Character{25}}%
\TextAt{27}{350.545}{306.563}{p}%
\TextAt{13}{355.032}{308.356}{\Character{59}}%
\TextAt{9}{360.276}{308.356}{1}%
\TextAt{13}{366.129}{308.356}{\Character{59}}%
\TextAt{9}{371.373}{308.356}{1}%
\TextAt{13}{377.226}{308.356}{\Character{59}}%
\TextAt{13}{382.45}{308.356}{\Character{58}}%
\TextAt{13}{387.687}{308.356}{\Character{58}}%
\TextAt{13}{392.911}{308.356}{\Character{58}}%
\TextAt{9}{396.21}{308.356}{\Character{41}}%
\TextAt{8}{400.763}{308.356}{.}%
\TextAt{8}{97.159}{288.321}{The}%
\TextAt{8}{122.698}{288.321}{relation}%
\TextAt{8}{167.113}{288.321}{between}%
\TextAt{8}{213.789}{288.321}{the}%
\TextAt{8}{235.43}{288.321}{idèle-theoretic}%
\TextAt{8}{312.998}{288.321}{and}%
\TextAt{8}{337.245}{288.321}{the}%
\TextAt{8}{358.886}{288.321}{ideal-theoretic}%
\TextAt{8}{437.111}{288.321}{formulation}%
\TextAt{8}{501.696}{288.321}{of}%
\TextAt{8}{516.51}{288.321}{the}%
\TextAt{8}{79.226}{271.275}{Artin}%
\TextAt{8}{110.764}{271.275}{reciprocity}%
\TextAt{8}{168.974}{271.275}{law}%
\TextAt{8}{190.11}{271.275}{is}%
\TextAt{8}{201.874}{271.275}{now}%
\TextAt{8}{226.263}{271.275}{provided}%
\TextAt{8}{274.394}{271.275}{by}%
\TextAt{8}{290.653}{271.275}{the}%
\TextAt{8}{310.81}{271.275}{following}%
\TextAt{8}{360.244}{271.275}{theorem.}%
\TextAt{12}{79.226}{246.977}{Theorem}%
\TextAt{12}{136.132}{246.977}{(7.1).}%
\TextAt{8}{171.353}{246.977}{Let}%
\TextAt{13}{192.228}{246.977}{L}%
\TextAt{16}{203.513}{246.977}{\Character{106}}%
\TextAt{13}{210.155}{246.977}{K}%
\TextAt{8}{224.775}{246.977}{be}%
\TextAt{8}{240.605}{246.977}{an}%
\TextAt{8}{256.77}{246.977}{abelian}%
\TextAt{8}{297.323}{246.977}{extension,}%
\TextAt{8}{352.265}{246.977}{and}%
\TextAt{8}{374.933}{246.977}{let}%
\TextAt{14}{391.743}{246.977}{m}%
\TextAt{8}{404.713}{246.977}{be}%
\TextAt{8}{420.543}{246.977}{a}%
\TextAt{8}{430.204}{246.977}{module}%
\TextAt{8}{471.402}{246.977}{of}%
\TextAt{8}{484.638}{246.977}{definition}%
\TextAt{8}{79.226}{229.931}{for}%
\TextAt{8}{97.111}{229.931}{it.}%
\TextAt{8}{113.366}{229.931}{Then}%
\TextAt{8}{143.924}{229.931}{the}%
\TextAt{8}{164.08}{229.931}{Artin}%
\TextAt{8}{195.618}{229.931}{symbol}%
\TextAt{8}{235.679}{229.931}{induces}%
\TextAt{8}{277.355}{229.931}{a}%
\TextAt{8}{287.111}{229.931}{surjective}%
\TextAt{8}{339.522}{229.931}{homomorphism}%
\TextAt{20}{229.728}{220.061}{\Character{18}}%
\TextAt{13}{239.724}{211.292}{L}%
\TextAt{16}{251.01}{211.292}{\Character{106}}%
\TextAt{13}{257.651}{211.292}{K}%
\Vector{q 0 G 0 g 0.478 w 167.724 -567.807 m 196.469 -567.807 l S Q}%
\TextAt{20}{269.665}{220.061}{\Character{19}}%
\TextAt{9}{281.786}{203.204}{\Character{58}}%
\TextAt{13}{288.358}{203.204}{C}%
\TextAt{13}{297.576}{203.204}{l}%
\TextAt{27}{301.342}{208.14}{m}%
\TextAt{21}{301.079}{200.248}{K}%
\NativeArrowAt{312.602}{203.204}{11.955}{19.26}{\longrightarrow}%
\TextAt{13}{335.184}{203.204}{G}%
\TextAt{9}{344.418}{203.204}{\Character{40}}%
\TextAt{13}{348.97}{203.204}{L}%
\TextAt{16}{360.255}{203.204}{\Character{106}}%
\TextAt{13}{366.897}{203.204}{K}%
\TextAt{9}{377.715}{203.204}{\Character{41}}%
\TextAt{8}{79.226}{175.927}{with}%
\TextAt{8}{106.163}{175.927}{kernel}%
\TextAt{13}{140.907}{175.927}{H}%
\TextAt{27}{151.53}{180.266}{m}%
\TextAt{9}{158.139}{175.927}{\Character{47}}%
\TextAt{13}{163.992}{175.927}{P}%
\TextAt{27}{173.163}{180.266}{m}%
\TextAt{21}{171.538}{172.907}{K}%
\TextAt{8}{179.773}{175.927}{,}%
\TextAt{8}{187.269}{175.927}{where}%
\TextAt{13}{221.355}{175.927}{H}%
\TextAt{27}{231.978}{180.266}{m}%
\TextAt{9}{242.38}{175.927}{\Character{61}}%
\TextAt{9}{255.268}{175.927}{\Character{40}}%
\TextAt{13}{259.831}{175.927}{N}%
\TextAt{21}{269.204}{174.068}{L}%
\TextAt{3}{274.963}{174.068}{\Character{106}}%
\TextAt{21}{277.316}{174.068}{K}%
\TextAt{13}{285.518}{175.927}{J}%
\TextAt{27}{293.095}{180.266}{m}%
\TextAt{21}{291.989}{172.907}{L}%
\TextAt{9}{299.705}{175.927}{\Character{41}}%
\TextAt{13}{304.257}{175.927}{P}%
\TextAt{27}{313.428}{180.266}{m}%
\TextAt{21}{311.802}{172.907}{K}%
\TextAt{8}{320.037}{175.927}{,}%
\TextAt{8}{327.533}{175.927}{and}%
\TextAt{8}{350.573}{175.927}{we}%
\TextAt{8}{368.078}{175.927}{have}%
\TextAt{8}{395.35}{175.927}{an}%
\TextAt{8}{411.886}{175.927}{exact}%
\TextAt{8}{443.056}{175.927}{commutative}%
\TextAt{8}{513.259}{175.927}{dia-}%
\TextAt{8}{79.226}{158.881}{gram}%
\TextAt{9}{195.051}{139.083}{1}%
\TextAt{13}{225.556}{139.083}{N}%
\TextAt{21}{234.93}{137.224}{L}%
\TextAt{3}{240.689}{137.224}{\Character{106}}%
\TextAt{21}{243.041}{137.224}{K}%
\TextAt{13}{251.244}{139.083}{C}%
\TextAt{21}{259.619}{137.29}{L}%
\TextAt{13}{292.7}{139.083}{C}%
\TextAt{21}{301.075}{137.29}{K}%
\TextAt{13}{336.102}{139.083}{G}%
\TextAt{9}{345.335}{139.083}{\Character{40}}%
\TextAt{13}{349.888}{139.083}{L}%
\TextAt{16}{361.173}{139.083}{\Character{106}}%
\TextAt{13}{367.815}{139.083}{K}%
\TextAt{9}{378.632}{139.083}{\Character{41}}%
\TextAt{9}{409.462}{139.083}{1}%
\TextAt{9}{195.051}{95.986}{1}%
\TextAt{13}{226.283}{95.986}{H}%
\TextAt{27}{236.906}{100.324}{m}%
\TextAt{9}{243.515}{95.986}{\Character{47}}%
\TextAt{13}{249.368}{95.986}{P}%
\TextAt{27}{258.539}{100.324}{m}%
\TextAt{21}{256.913}{92.965}{K}%
\TextAt{13}{290.527}{95.986}{C}%
\TextAt{13}{299.745}{95.986}{l}%
\TextAt{27}{303.511}{100.324}{m}%
\TextAt{21}{303.248}{92.965}{K}%
\TextAt{13}{336.102}{95.986}{G}%
\TextAt{9}{345.335}{95.986}{\Character{40}}%
\TextAt{13}{349.888}{95.986}{L}%
\TextAt{16}{361.173}{95.986}{\Character{106}}%
\TextAt{13}{367.815}{95.986}{K}%
\TextAt{9}{378.632}{95.986}{\Character{41}}%
\TextAt{9}{407.837}{95.986}{1}%
\TextAt{13}{413.69}{95.986}{\Character{58}}%
\NativeDiagram{%
\path[draw=black,line width=0.478bp] (134.257,-631.927) -- (147.726,-631.927);%
\path[draw=black,line width=0.478bp,line cap=round,line join=round] (145.479,-629.058) .. controls (145.933,-630.779) and (146.945,-631.592) .. (147.965,-631.927) .. controls (146.945,-632.262) and (145.933,-633.075) .. (145.479,-634.796);%
\path[draw=black,line width=0.478bp] (199.229,-631.927) -- (214.871,-631.927);%
\path[draw=black,line width=0.478bp,line cap=round,line join=round] (212.624,-629.058) .. controls (213.078,-630.779) and (214.091,-631.592) .. (215.111,-631.927) .. controls (214.091,-632.262) and (213.078,-633.075) .. (212.624,-634.796);%
\path[draw=black,line width=0.478bp] (173.717,-643.347) -- (173.717,-663.99);%
\path[draw=black,line width=0.478bp,line cap=round,line join=round] (176.586,-661.743) .. controls (174.864,-662.197) and (174.052,-663.209) .. (173.717,-664.229) .. controls (173.382,-663.209) and (172.569,-662.197) .. (170.848,-661.743);%
\path[draw=black,line width=0.478bp] (242.632,-631.927) -- (258.274,-631.927);%
\path[draw=black,line width=0.478bp,line cap=round,line join=round] (256.026,-629.058) .. controls (256.481,-630.779) and (257.493,-631.592) .. (258.513,-631.927) .. controls (257.493,-632.262) and (256.481,-633.075) .. (256.026,-634.796);%
}%
\TextAt{17}{308.366}{146.881}{\Character{40}}%
\TextAt{21}{315.556}{146.881}{\Character{59}L}%
\TextAt{3}{323.668}{146.881}{\Character{106}}%
\TextAt{21}{326.02}{146.881}{K}%
\TextAt{17}{333.725}{146.881}{\Character{41}}%
\NativeDiagram{%
\path[draw=black,line width=0.478bp] (228.991,-641.288) -- (228.991,-664.654);%
\path[draw=black,line width=0.478bp,line cap=round,line join=round] (231.86,-662.407) .. controls (230.138,-662.861) and (229.325,-663.874) .. (228.991,-664.894) .. controls (228.656,-663.874) and (227.843,-662.861) .. (226.122,-662.407);%
\path[draw=black,line width=0.478bp] (316.54,-631.927) -- (331.635,-631.927);%
\path[draw=black,line width=0.478bp,line cap=round,line join=round] (329.388,-629.058) .. controls (329.842,-630.779) and (330.854,-631.592) .. (331.874,-631.927) .. controls (330.854,-632.262) and (329.842,-633.075) .. (329.388,-634.796);%
\path[draw=black,line width=0.478bp] (287.646,-642.484) -- (287.646,-663.99);%
\path[draw=black,line width=0.478bp,line cap=round,line join=round] (290.515,-661.743) .. controls (288.794,-662.197) and (287.981,-663.209) .. (287.646,-664.229) .. controls (287.311,-663.209) and (286.498,-662.197) .. (284.777,-661.743);%
}%
\TextAt{7}{362.461}{117.798}{id}%
\NativeDiagram{%
\path[draw=black,line width=0.478bp] (134.257,-675.025) -- (148.454,-675.025);%
\path[draw=black,line width=0.478bp,line cap=round,line join=round] (146.206,-672.156) .. controls (146.661,-673.878) and (147.673,-674.69) .. (148.693,-675.025) .. controls (147.673,-675.36) and (146.661,-676.173) .. (146.206,-677.894);%
\path[draw=black,line width=0.478bp] (198.502,-675.025) -- (212.699,-675.025);%
\path[draw=black,line width=0.478bp,line cap=round,line join=round] (210.451,-672.156) .. controls (210.906,-673.878) and (211.918,-674.69) .. (212.938,-675.025) .. controls (211.918,-675.36) and (210.906,-676.173) .. (210.451,-677.894);%
\path[draw=black,line width=0.478bp] (244.804,-675.025) -- (258.274,-675.025);%
\path[draw=black,line width=0.478bp,line cap=round,line join=round] (256.026,-672.156) .. controls (256.481,-673.878) and (257.493,-674.69) .. (258.513,-675.025) .. controls (257.493,-675.36) and (256.481,-676.173) .. (256.026,-677.894);%
}%
\TextAt{9}{311.094}{86.654}{\Character{40}}%
\TextAt{22}{316.842}{91.676}{L}%
\TextAt{23}{321.9}{91.676}{\Character{106}}%
\TextAt{22}{324.17}{91.676}{K}%
\NativeDiagram{%
\path[draw=black,line width=0.359bp] (244.842,-684.357) -- (258.715,-684.357);%
}%
\TextAt{9}{331.91}{86.654}{\Character{41}}%
\NativeDiagram{%
\path[draw=black,line width=0.478bp] (316.54,-675.025) -- (330.009,-675.025);%
\path[draw=black,line width=0.478bp,line cap=round,line join=round] (327.762,-672.156) .. controls (328.216,-673.878) and (329.229,-674.69) .. (330.249,-675.025) .. controls (329.229,-675.36) and (328.216,-676.173) .. (327.762,-677.894);%
}%
\TextAt{8}{297.211}{43.033}{371}%
\Destination{78.23}{809.815}{nav.page.371}%
\Destination{79.36}{721.955}{nav.print.407}%
\Destination{79.36}{260.955}{nav.number.6.7.1}%
\Hotlink{325.585}{516.355}{25.078}{11.047}{nav.number.6.6.6}%
}

% Page 384 (body page 372)
\Page{%
\TextAt{2}{79.226}{796.943}{VI.}%
\TextAt{2}{98.012}{796.943}{Global}%
\TextAt{2}{133.248}{796.943}{Class}%
\TextAt{2}{161.841}{796.943}{Field}%
\TextAt{2}{189.572}{796.943}{Theory}%
\Vector{q 0 G 0 g 0.299 w 7.226 18.729 m 460.769 18.729 l S Q}%
\TextAt{19}{79.226}{761.124}{Proof.}%
\TextAt{8}{116}{761.124}{In}%
\TextAt{8}{131.143}{761.124}{§1,}%
\TextAt{8}{150.48}{761.124}{we}%
\TextAt{8}{168.229}{761.124}{obtained}%
\TextAt{8}{216.882}{761.124}{the}%
\TextAt{8}{237.56}{761.124}{isomorphism}%
\TextAt{9}{306.5}{761.124}{\Character{40}}%
\TextAt{9}{315.454}{761.124}{\Character{41}}%
\TextAt{9}{324.206}{761.124}{\Character{58}}%
\TextAt{13}{331.694}{761.124}{C}%
\TextAt{21}{340.069}{759.33}{K}%
\TextAt{9}{348.272}{761.124}{\Character{47}}%
\TextAt{13}{354.125}{761.124}{C}%
\TextAt{27}{363.358}{765.462}{m}%
\TextAt{21}{362.5}{758.103}{K}%
\NativeArrowAt{374.912}{761.124}{11.955}{19.26}{\longrightarrow}%
\TextAt{13}{398.382}{761.124}{C}%
\TextAt{13}{407.6}{761.124}{l}%
\TextAt{27}{411.365}{765.462}{m}%
\TextAt{21}{411.103}{758.103}{K}%
\TextAt{9}{423.514}{761.124}{\Character{61}}%
\TextAt{13}{436.828}{761.124}{J}%
\TextAt{27}{444.405}{765.462}{m}%
\TextAt{21}{443.298}{758.103}{K}%
\TextAt{9}{451.501}{761.124}{\Character{47}}%
\TextAt{13}{457.354}{761.124}{P}%
\TextAt{27}{466.525}{765.462}{m}%
\TextAt{21}{464.899}{758.103}{K}%
\TextAt{8}{477.553}{761.124}{by}%
\TextAt{8}{494.334}{761.124}{sending}%
\ErrataAt{79.19}{744.078}{455.81}{a representative $\alpha\in I_K^{(\mathfrak m)}$ to the ideal $(\alpha)=\prod_{\mathfrak p\nmid\infty}\mathfrak p^{v_{\mathfrak p}(\alpha_{\mathfrak p})}$. This isomorphism yields a commutative}%
\TextAt{8}{79.226}{727.032}{diagram}%
\TextAt{13}{234.654}{707.234}{C}%
\TextAt{21}{243.029}{705.441}{K}%
\TextAt{9}{251.231}{707.234}{\Character{47}}%
\TextAt{13}{257.084}{707.234}{C}%
\TextAt{27}{266.318}{711.572}{m}%
\TextAt{21}{265.459}{704.213}{K}%
\TextAt{13}{328.631}{707.234}{G}%
\TextAt{9}{337.865}{707.234}{\Character{40}}%
\TextAt{13}{342.417}{707.234}{L}%
\TextAt{16}{353.703}{707.234}{\Character{106}}%
\TextAt{13}{360.344}{707.234}{K}%
\TextAt{9}{371.162}{707.234}{\Character{41}}%
\TextAt{13}{243.696}{664.968}{C}%
\TextAt{13}{252.914}{664.968}{l}%
\TextAt{27}{256.679}{669.306}{m}%
\TextAt{21}{256.416}{661.947}{K}%
\TextAt{13}{327.005}{664.968}{G}%
\TextAt{9}{336.239}{664.968}{\Character{40}}%
\TextAt{13}{340.791}{664.968}{L}%
\TextAt{16}{352.077}{664.968}{\Character{106}}%
\TextAt{13}{358.718}{664.968}{K}%
\TextAt{9}{369.536}{664.968}{\Character{41}}%
\TextAt{13}{374.088}{664.968}{\Character{59}}%
\NativeDiagram{%
\path[draw=black,line width=0.478bp] (207.014,-63.777) -- (250.803,-63.777);%
\path[draw=black,line width=0.478bp,line cap=round,line join=round] (248.555,-60.908) .. controls (249.01,-62.629) and (250.022,-63.442) .. (251.042,-63.777) .. controls (250.022,-64.112) and (249.01,-64.925) .. (248.555,-66.646);%
}%
\TextAt{17}{286.821}{715.031}{\Character{40}}%
\TextAt{21}{294.011}{715.031}{\Character{59}L}%
\TextAt{3}{302.123}{715.031}{\Character{106}}%
\TextAt{21}{304.475}{715.031}{K}%
\TextAt{17}{312.18}{715.031}{\Character{41}}%
\NativeDiagram{%
\path[draw=black,line width=0.478bp] (182.158,-74.365) -- (182.158,-95.673);%
\path[draw=black,line width=0.478bp,line cap=round,line join=round] (185.027,-93.425) .. controls (183.306,-93.879) and (182.493,-94.892) .. (182.158,-95.912) .. controls (181.823,-94.892) and (181.01,-93.879) .. (179.289,-93.425);%
}%
\TextAt{17}{240.859}{686.749}{\Character{40}}%
\TextAt{17}{248.04}{686.749}{\Character{41}}%
%
\NativeDiagram{%
\path[draw=black,line width=0.478bp] (280.175,-74.334) -- (280.175,-95.008);%
\path[draw=black,line width=0.478bp,line cap=round,line join=round] (283.044,-92.761) .. controls (281.322,-93.215) and (280.51,-94.228) .. (280.175,-95.247) .. controls (279.84,-94.228) and (279.027,-93.215) .. (277.306,-92.761);%
}%
\TextAt{7}{354.99}{686.364}{id}%
\NativeDiagram{%
\path[draw=black,line width=0.478bp] (197.972,-106.043) -- (249.177,-106.043);%
\path[draw=black,line width=0.478bp,line cap=round,line join=round] (246.93,-103.174) .. controls (247.384,-104.896) and (248.396,-105.709) .. (249.416,-106.043) .. controls (248.396,-106.378) and (247.384,-107.191) .. (246.93,-108.912);%
}%
\TextAt{21}{293.34}{659.606}{f}%
\TextAt{8}{79.226}{639.735}{and}%
\TextAt{8}{101.989}{639.735}{we}%
\TextAt{8}{119.216}{639.735}{show}%
\TextAt{8}{148.219}{639.735}{that}%
\TextAt{13}{173.588}{639.735}{f}%
\TextAt{8}{184.532}{639.735}{is}%
\TextAt{8}{196.296}{639.735}{given}%
\TextAt{8}{226.865}{639.735}{by}%
\TextAt{8}{243.124}{639.735}{the}%
\TextAt{8}{263.281}{639.735}{Artin}%
\TextAt{8}{294.819}{639.735}{symbol.}%
\TextAt{8}{97.159}{619.7}{Let}%
\TextAt{14}{118.004}{619.7}{p}%
\TextAt{8}{127.76}{619.7}{be}%
\TextAt{8}{143.56}{619.7}{a}%
\TextAt{8}{153.191}{619.7}{prime}%
\TextAt{8}{186.23}{619.7}{ideal}%
\TextAt{8}{214.068}{619.7}{not}%
\TextAt{8}{234.758}{619.7}{dividing}%
\TextAt{14}{279.836}{619.7}{m}%
\TextAt{8}{289.003}{619.7}{,}%
\TextAt{13}{296.052}{619.7}{\Character{25}}%
\TextAt{27}{302.693}{617.907}{p}%
\TextAt{8}{310.952}{619.7}{a}%
\TextAt{8}{320.583}{619.7}{prime}%
\TextAt{8}{353.623}{619.7}{element}%
\TextAt{8}{396.74}{619.7}{of}%
\TextAt{13}{409.946}{619.7}{K}%
\TextAt{27}{419.905}{617.907}{p}%
\TextAt{8}{424.392}{619.7}{,}%
\TextAt{8}{431.441}{619.7}{and}%
\TextAt{13}{454.079}{619.7}{c}%
\NativeMathAt{462.438}{619.7}{11.955}{\in}%
\TextAt{13}{473.729}{619.7}{C}%
\TextAt{21}{482.104}{617.907}{K}%
\TextAt{9}{490.307}{619.7}{\Character{47}}%
\TextAt{13}{496.16}{619.7}{C}%
\TextAt{27}{505.393}{624.039}{m}%
\TextAt{21}{504.535}{616.68}{K}%
\TextAt{8}{516.51}{619.7}{the}%
\TextAt{8}{79.226}{602.654}{class}%
\TextAt{8}{106.663}{602.654}{of}%
\TextAt{8}{119.993}{602.654}{the}%
\TextAt{8}{140.15}{602.654}{idèle}%
\TextAt{16}{167.455}{602.654}{\Character{104}}%
\TextAt{13}{172.105}{602.654}{\Character{25}}%
\TextAt{27}{178.745}{600.861}{p}%
\TextAt{16}{183.232}{602.654}{\Character{105}}%
\TextAt{9}{191.202}{602.654}{\Character{61}}%
\TextAt{9}{203.612}{602.654}{\Character{40}}%
\TextAt{13}{208.18}{602.654}{\Character{58}}%
\TextAt{13}{213.404}{602.654}{\Character{58}}%
\TextAt{13}{218.641}{602.654}{\Character{58}}%
\TextAt{13}{223.865}{602.654}{\Character{59}}%
\TextAt{9}{229.156}{602.654}{1}%
\TextAt{13}{235.009}{602.654}{\Character{59}}%
\TextAt{9}{240.253}{602.654}{1}%
\TextAt{13}{246.106}{602.654}{\Character{59}}%
\TextAt{13}{251.33}{602.654}{\Character{25}}%
\TextAt{27}{257.991}{600.861}{p}%
\TextAt{13}{262.478}{602.654}{\Character{59}}%
\TextAt{9}{267.722}{602.654}{1}%
\TextAt{13}{273.575}{602.654}{\Character{59}}%
\TextAt{9}{278.819}{602.654}{1}%
\TextAt{13}{284.672}{602.654}{\Character{59}}%
\TextAt{13}{289.896}{602.654}{\Character{58}}%
\TextAt{13}{295.133}{602.654}{\Character{58}}%
\TextAt{13}{300.357}{602.654}{\Character{58}}%
\TextAt{9}{303.656}{602.654}{\Character{41}}%
\TextAt{8}{308.208}{602.654}{.}%
\TextAt{8}{316.657}{602.654}{Then}%
\TextAt{9}{347.214}{602.654}{\Character{40}}%
\TextAt{13}{351.767}{602.654}{c}%
\TextAt{9}{356.805}{602.654}{\Character{41}}%
\TextAt{9}{364.66}{602.654}{\Character{61}}%
\TextAt{14}{377.103}{602.654}{p}%
\TextAt{8}{386.407}{602.654}{mod}%
\TextAt{13}{411.845}{602.654}{P}%
\TextAt{27}{421.016}{606.993}{m}%
\TextAt{21}{419.39}{599.634}{K}%
\TextAt{8}{431.523}{602.654}{and}%
\TextAt{13}{184.416}{569.68}{f}%
\TextAt{9}{191.462}{569.68}{\Character{40}\Character{40}}%
\TextAt{13}{200.567}{569.68}{c}%
\TextAt{9}{205.605}{569.68}{\Character{41}\Character{41}}%
\TextAt{9}{218.003}{569.68}{\Character{61}}%
\TextAt{9}{230.424}{569.68}{\Character{40}}%
\TextAt{13}{235.008}{569.68}{c\Character{59}}%
\TextAt{13}{245.266}{569.68}{L}%
\TextAt{16}{256.576}{569.68}{\Character{106}}%
\TextAt{13}{263.217}{569.68}{K}%
\TextAt{9}{274.035}{569.68}{\Character{41}}%
\TextAt{9}{281.89}{569.68}{\Character{61}}%
\TextAt{9}{294.311}{569.68}{\Character{40}}%
\TextAt{16}{298.886}{569.68}{\Character{104}}%
\TextAt{13}{303.535}{569.68}{\Character{25}}%
\TextAt{27}{310.176}{567.887}{p}%
\TextAt{16}{314.662}{569.68}{\Character{105}}%
\TextAt{13}{319.312}{569.68}{\Character{59}}%
\TextAt{13}{324.536}{569.68}{L}%
\TextAt{16}{335.841}{569.68}{\Character{106}}%
\TextAt{13}{342.483}{569.68}{K}%
\TextAt{9}{353.301}{569.68}{\Character{41}}%
\TextAt{9}{361.156}{569.68}{\Character{61}}%
\TextAt{20}{373.599}{586.537}{\Character{18}}%
\TextAt{13}{383.595}{577.768}{L}%
\TextAt{16}{394.88}{577.768}{\Character{106}}%
\TextAt{13}{401.522}{577.768}{K}%
\Vector{q 0 G 0 g 0.478 w 311.595 -201.331 m 340.34 -201.331 l S Q}%
\TextAt{14}{394.976}{561.479}{p}%
\TextAt{20}{413.535}{586.537}{\Character{19}}%
\TextAt{13}{424.328}{569.68}{\Character{58}}%
\TextAt{8}{79.226}{532.235}{This}%
\TextAt{8}{105.419}{532.235}{shows}%
\TextAt{8}{138.51}{532.235}{that}%
\TextAt{13}{163.351}{532.235}{f}%
\TextAt{9}{173.719}{532.235}{\Character{58}}%
\TextAt{13}{180.291}{532.235}{J}%
\TextAt{27}{187.868}{536.573}{m}%
\TextAt{21}{186.761}{529.214}{K}%
\TextAt{9}{194.964}{532.235}{\Character{47}}%
\TextAt{13}{200.817}{532.235}{P}%
\TextAt{27}{209.988}{536.573}{m}%
\TextAt{21}{208.362}{529.214}{K}%
\NativeArrowAt{219.918}{532.235}{11.955}{19.26}{\longrightarrow}%
\TextAt{13}{242.5}{532.235}{G}%
\TextAt{9}{251.734}{532.235}{\Character{40}}%
\TextAt{13}{256.286}{532.235}{L}%
\TextAt{16}{267.572}{532.235}{\Character{106}}%
\TextAt{13}{274.213}{532.235}{K}%
\TextAt{9}{285.031}{532.235}{\Character{41}}%
\TextAt{8}{292.954}{532.235}{is}%
\TextAt{8}{304.19}{532.235}{induced}%
\TextAt{8}{347.228}{532.235}{by}%
\TextAt{8}{362.96}{532.235}{the}%
\TextAt{8}{382.589}{532.235}{Artin}%
\TextAt{8}{413.6}{532.235}{symbol}%
\TextAt{20}{453.135}{545.505}{\Character{16}}%
\TextAt{21}{461.47}{537.933}{L}%
\TextAt{3}{467.229}{537.933}{\Character{106}}%
\TextAt{21}{469.581}{537.933}{K}%
\Vector{q 0 G 0 g 0.478 w 389.47 -238.777 m 405.286 -238.777 l S Q}%
\TextAt{20}{478.481}{545.505}{\Character{17}}%
\TextAt{9}{488.942}{532.235}{\Character{58}}%
\TextAt{13}{495.515}{532.235}{J}%
\TextAt{27}{503.092}{536.573}{m}%
\TextAt{21}{501.985}{529.214}{K}%
\NativeArrowAt{513.508}{532.235}{11.955}{19.26}{\longrightarrow}%
\TextAt{13}{79.226}{515.189}{G}%
\TextAt{9}{88.46}{515.189}{\Character{40}}%
\TextAt{13}{93.012}{515.189}{L}%
\TextAt{16}{104.297}{515.189}{\Character{106}}%
\TextAt{13}{110.939}{515.189}{K}%
\TextAt{9}{121.757}{515.189}{\Character{41}}%
\TextAt{8}{126.309}{515.189}{,}%
\TextAt{8}{133.458}{515.189}{and}%
\TextAt{8}{156.221}{515.189}{that}%
\TextAt{8}{181.59}{515.189}{it}%
\TextAt{8}{193.294}{515.189}{is}%
\TextAt{8}{205.058}{515.189}{surjective.}%
\TextAt{8}{97.159}{495.154}{It}%
\TextAt{8}{110.135}{495.154}{remains}%
\TextAt{8}{154.074}{495.154}{to}%
\TextAt{8}{168.688}{495.154}{show}%
\TextAt{8}{197.995}{495.154}{that}%
\TextAt{8}{223.667}{495.154}{the}%
\TextAt{8}{244.126}{495.154}{image}%
\TextAt{8}{278.251}{495.154}{of}%
\TextAt{13}{291.884}{495.154}{N}%
\TextAt{21}{301.258}{493.294}{L}%
\TextAt{3}{307.017}{493.294}{\Character{106}}%
\TextAt{21}{309.369}{493.294}{K}%
\TextAt{13}{317.572}{495.154}{C}%
\TextAt{21}{325.947}{493.361}{L}%
\TextAt{8}{336.405}{495.154}{under}%
\TextAt{8}{369.872}{495.154}{the}%
\TextAt{8}{390.332}{495.154}{map}%
\TextAt{9}{416.649}{495.154}{\Character{40}}%
\TextAt{9}{425.388}{495.154}{\Character{41}}%
\TextAt{9}{433.769}{495.154}{\Character{58}}%
\TextAt{13}{440.882}{495.154}{C}%
\TextAt{21}{449.257}{493.361}{K}%
\NativeArrowAt{461.297}{495.154}{11.955}{19.26}{\longrightarrow}%
\TextAt{13}{484.396}{495.154}{J}%
\TextAt{27}{491.973}{499.492}{m}%
\TextAt{21}{490.866}{492.133}{K}%
\TextAt{9}{499.069}{495.154}{\Character{47}}%
\TextAt{13}{504.922}{495.154}{P}%
\TextAt{27}{514.093}{499.492}{m}%
\TextAt{21}{512.467}{492.133}{K}%
\TextAt{8}{524.903}{495.154}{is}%
\TextAt{8}{79.226}{478.108}{the}%
\TextAt{8}{99.548}{478.108}{group}%
\TextAt{13}{132.89}{478.108}{H}%
\TextAt{27}{143.512}{482.446}{m}%
\TextAt{9}{150.121}{478.108}{\Character{47}}%
\TextAt{13}{155.974}{478.108}{P}%
\TextAt{27}{165.146}{482.446}{m}%
\TextAt{21}{163.52}{475.087}{K}%
\TextAt{8}{171.755}{478.108}{.}%
\TextAt{8}{180.7}{478.108}{We}%
\TextAt{8}{201.01}{478.108}{view}%
\TextAt{8}{228.159}{478.108}{the}%
\TextAt{8}{248.481}{478.108}{module}%
\TextAt{14}{289.94}{478.108}{m}%
\TextAt{9}{302.71}{478.108}{\Character{61}}%
\TextAt{20}{315.418}{487.074}{\Character{81}}%
\TextAt{27}{326.709}{474.621}{p}%
\TextAt{25}{330.697}{474.621}{\Character{45}}%
\NativeMathAt{332.911}{474.621}{7.97}{\infty}%
\TextAt{14}{343.87}{478.108}{p}%
\TextAt{21}{349.853}{482.446}{n}%
\TextAt{28}{354.992}{481.45}{p}%
\TextAt{8}{363.635}{478.108}{as}%
\TextAt{8}{378.171}{478.108}{a}%
\TextAt{8}{388.092}{478.108}{module}%
\TextAt{8}{429.551}{478.108}{of}%
\TextAt{13}{443.047}{478.108}{L}%
\TextAt{8}{455.075}{478.108}{by}%
\TextAt{8}{471.499}{478.108}{substituting}%
\TextAt{8}{79.226}{461.062}{for}%
\TextAt{8}{96.909}{461.062}{each}%
\TextAt{8}{123.043}{461.062}{prime}%
\TextAt{8}{156.005}{461.062}{ideal}%
\TextAt{14}{183.765}{461.062}{p}%
\TextAt{8}{193.443}{461.062}{of}%
\TextAt{13}{206.571}{461.062}{K}%
\TextAt{8}{221.083}{461.062}{the}%
\TextAt{8}{241.037}{461.062}{product}%
\TextAt{14}{284.734}{461.062}{p}%
\TextAt{9}{294.038}{461.062}{\Character{61}}%
\TextAt{20}{306.464}{470.028}{\Character{81}}%
\TextAt{27}{317.755}{457.575}{P}%
\TextAt{3}{324.35}{457.575}{\Character{106}}%
\TextAt{27}{326.703}{457.575}{p}%
\TextAt{14}{333.182}{461.062}{P}%
\TextAt{21}{343.075}{465.581}{e}%
\TextAt{28}{347.004}{463.843}{P}%
\TextAt{23}{352.703}{463.843}{\Character{106}}%
\TextAt{28}{354.972}{463.843}{p}%
\TextAt{8}{359.553}{461.062}{.}%
\TextAt{8}{367.933}{461.062}{As}%
\TextAt{8}{385.018}{461.062}{in}%
\TextAt{8}{398.469}{461.062}{the}%
\TextAt{8}{418.423}{461.062}{proof}%
\TextAt{8}{448.789}{461.062}{of}%
\TextAt{8}{461.917}{461.062}{(1.9),}%
\TextAt{8}{492.982}{461.062}{we}%
\TextAt{8}{510.007}{461.062}{then}%
\TextAt{8}{79.226}{442.498}{get}%
\TextAt{13}{98.737}{442.498}{C}%
\TextAt{21}{107.112}{440.705}{L}%
\TextAt{9}{116.69}{442.498}{\Character{61}}%
\TextAt{13}{129.116}{442.498}{I}%
\TextAt{17}{135.218}{448.608}{\Character{40}}%
\TextAt{27}{138.512}{448.608}{m}%
\TextAt{17}{144.623}{448.608}{\Character{41}}%
\TextAt{21}{134.278}{439.257}{L}%
\TextAt{13}{148.414}{442.498}{L}%
\TextAt{3}{156.379}{446.837}{\Character{3}}%
\TextAt{9}{161.111}{442.498}{\Character{47}}%
\TextAt{13}{166.964}{442.498}{L}%
\TextAt{3}{174.928}{446.837}{\Character{3}}%
\TextAt{8}{179.661}{442.498}{,}%
\TextAt{8}{186.81}{442.498}{where}%
\TextAt{13}{192.416}{414.109}{I}%
\TextAt{17}{198.519}{420.22}{\Character{40}}%
\TextAt{27}{201.812}{420.22}{m}%
\TextAt{17}{207.924}{420.22}{\Character{41}}%
\TextAt{21}{197.579}{410.868}{L}%
\TextAt{9}{215.036}{414.109}{\Character{61}}%
\TextAt{16}{227.461}{414.109}{\Character{102}}%
\TextAt{13}{233.439}{414.109}{\Character{11}}%
\NativeMathAt{244.281}{414.109}{11.955}{\in}%
\TextAt{13}{255.572}{414.109}{I}%
\TextAt{21}{260.735}{412.316}{L}%
\TextAt{16}{270.313}{414.109}{\Character{106}}%
\TextAt{13}{276.955}{414.109}{\Character{11}}%
\TextAt{27}{284.401}{412.316}{P}%
\NativeMathAt{294.815}{414.109}{11.955}{\in}%
\TextAt{13}{306.106}{414.109}{U}%
\TextAt{17}{315.305}{421.459}{\Character{40}}%
\TextAt{21}{318.598}{421.459}{e}%
\TextAt{28}{322.527}{419.721}{P}%
\TextAt{23}{328.226}{419.721}{\Character{106}}%
\TextAt{28}{330.495}{419.721}{p}%
\TextAt{21}{334.578}{421.459}{n}%
\TextAt{28}{339.716}{420.463}{p}%
\TextAt{17}{343.799}{421.459}{\Character{41}}%
\TextAt{27}{314.056}{410.804}{P}%
\TextAt{8}{351.487}{414.109}{for}%
\TextAt{14}{369.372}{414.109}{P}%
\TextAt{16}{382.586}{414.109}{\Character{106}}%
\TextAt{14}{389.228}{414.109}{m}%
\NativeMathAt{398.395}{414.109}{11.955}{\infty}%
\TextAt{16}{410.35}{414.109}{\Character{103}}%
\TextAt{13}{416.328}{414.109}{\Character{58}}%
\TextAt{8}{79.226}{385.72}{The}%
\TextAt{8}{103.28}{385.72}{elements}%
\TextAt{8}{151.136}{385.72}{of}%
\TextAt{13}{232.443}{368.674}{N}%
\TextAt{21}{241.816}{366.815}{L}%
\TextAt{3}{247.575}{366.815}{\Character{106}}%
\TextAt{21}{249.928}{366.815}{K}%
\TextAt{13}{258.13}{368.674}{C}%
\TextAt{21}{266.506}{366.881}{L}%
\TextAt{9}{276.084}{368.674}{\Character{61}}%
\TextAt{13}{288.509}{368.674}{N}%
\TextAt{21}{297.883}{366.815}{L}%
\TextAt{3}{303.642}{366.815}{\Character{106}}%
\TextAt{21}{305.994}{366.815}{K}%
\TextAt{9}{314.197}{368.674}{\Character{40}}%
\TextAt{13}{318.749}{368.674}{I}%
\TextAt{17}{324.852}{374.785}{\Character{40}}%
\TextAt{27}{328.145}{374.785}{m}%
\TextAt{17}{334.256}{374.785}{\Character{41}}%
\TextAt{21}{323.912}{365.433}{L}%
\TextAt{9}{338.048}{368.674}{\Character{41}}%
\TextAt{13}{342.6}{368.674}{K}%
\TextAt{3}{353.418}{373.61}{\Character{3}}%
\TextAt{9}{358.15}{368.674}{\Character{47}}%
\TextAt{13}{364.003}{368.674}{K}%
\TextAt{3}{374.821}{373.61}{\Character{3}}%
\TextAt{8}{79.226}{345.415}{are}%
\TextAt{8}{98.737}{345.415}{the}%
\TextAt{8}{118.893}{345.415}{classes}%
\TextAt{8}{156.146}{345.415}{of}%
\TextAt{8}{169.476}{345.415}{norm}%
\TextAt{8}{200.045}{345.415}{idèles}%
\TextAt{13}{231.965}{345.415}{N}%
\TextAt{21}{241.339}{343.555}{L}%
\TextAt{3}{247.098}{343.555}{\Character{106}}%
\TextAt{21}{249.451}{343.555}{K}%
\TextAt{9}{257.653}{345.415}{\Character{40}}%
\TextAt{13}{262.206}{345.415}{\Character{11}}%
\TextAt{9}{269.727}{345.415}{\Character{41}}%
\TextAt{8}{274.28}{345.415}{,}%
\TextAt{8}{281.429}{345.415}{for}%
\TextAt{13}{299.314}{345.415}{\Character{11}}%
\NativeMathAt{310.156}{345.415}{11.955}{\in}%
\TextAt{13}{321.447}{345.415}{I}%
\TextAt{17}{327.55}{351.525}{\Character{40}}%
\TextAt{27}{330.843}{351.525}{m}%
\TextAt{17}{336.954}{351.525}{\Character{41}}%
\TextAt{21}{326.61}{342.174}{L}%
\TextAt{8}{340.746}{345.415}{.}%
\TextAt{8}{349.194}{345.415}{As}%
\TextAt{13}{236.566}{317.026}{N}%
\TextAt{21}{245.939}{315.166}{L}%
\TextAt{3}{251.698}{315.166}{\Character{106}}%
\TextAt{21}{254.051}{315.166}{K}%
\TextAt{9}{262.253}{317.026}{\Character{40}}%
\TextAt{13}{266.806}{317.026}{\Character{11}}%
\TextAt{9}{274.328}{317.026}{\Character{41}}%
\TextAt{27}{278.88}{315.233}{p}%
\TextAt{9}{286.688}{317.026}{\Character{61}}%
\TextAt{20}{299.113}{328.383}{\Character{89}}%
\TextAt{27}{300.283}{302.48}{P}%
\TextAt{3}{306.878}{302.48}{\Character{106}}%
\TextAt{27}{309.231}{302.48}{p}%
\TextAt{13}{316.382}{317.026}{N}%
\TextAt{21}{325.755}{315.166}{L}%
\TextAt{28}{331.514}{313.775}{P}%
\TextAt{3}{337.712}{315.166}{\Character{106}}%
\TextAt{21}{340.064}{315.166}{K}%
\TextAt{28}{347.205}{314.17}{p}%
\TextAt{9}{351.786}{317.026}{\Character{40}}%
\TextAt{13}{356.339}{317.026}{\Character{11}}%
\TextAt{27}{363.784}{315.233}{P}%
\TextAt{9}{370.878}{317.026}{\Character{41}}%
\TextAt{8}{79.226}{277.987}{(see}%
\TextAt{8}{101.651}{277.987}{(2.2)),}%
\TextAt{8}{136.599}{277.987}{and}%
\TextAt{8}{158.318}{277.987}{since}%
\TextAt{13}{185.944}{277.987}{v}%
\TextAt{27}{191.603}{276.194}{p}%
\TextAt{9}{196.09}{277.987}{\Character{40}}%
\TextAt{13}{200.642}{277.987}{N}%
\TextAt{21}{210.016}{276.127}{L}%
\TextAt{28}{215.775}{274.736}{P}%
\TextAt{3}{221.972}{276.127}{\Character{106}}%
\TextAt{21}{224.325}{276.127}{K}%
\TextAt{28}{231.466}{275.131}{p}%
\TextAt{9}{236.047}{277.987}{\Character{40}}%
\TextAt{13}{240.599}{277.987}{\Character{11}}%
\TextAt{27}{248.045}{276.194}{P}%
\TextAt{9}{255.138}{277.987}{\Character{41}\Character{41}}%
\TextAt{9}{267.536}{277.987}{\Character{61}}%
\TextAt{13}{279.989}{277.987}{f}%
\TextAt{27}{285.759}{276.127}{P}%
\TextAt{3}{292.355}{276.127}{\Character{106}}%
\TextAt{27}{294.707}{276.127}{p}%
\TextAt{13}{299.194}{277.987}{v}%
\TextAt{27}{304.853}{276.194}{P}%
\TextAt{9}{311.947}{277.987}{\Character{40}}%
\TextAt{13}{316.499}{277.987}{\Character{11}}%
\TextAt{27}{323.945}{276.194}{P}%
\TextAt{9}{331.038}{277.987}{\Character{41}}%
\TextAt{8}{338.445}{277.987}{(see}%
\TextAt{8}{360.87}{277.987}{chap.}%
\TextAt{8}{390.719}{277.987}{III,}%
\TextAt{8}{410.132}{277.987}{(1.2)),}%
\TextAt{8}{445.079}{277.987}{the}%
\TextAt{8}{464.193}{277.987}{idèle}%
\TextAt{13}{490.455}{277.987}{N}%
\TextAt{21}{499.829}{276.127}{L}%
\TextAt{3}{505.588}{276.127}{\Character{106}}%
\TextAt{21}{507.94}{276.127}{K}%
\TextAt{9}{516.143}{277.987}{\Character{40}}%
\TextAt{13}{520.695}{277.987}{\Character{11}}%
\TextAt{9}{528.217}{277.987}{\Character{41}}%
\TextAt{8}{79.226}{260.941}{is}%
\TextAt{8}{90.99}{260.941}{mapped}%
\TextAt{8}{135.535}{260.941}{by}%
\TextAt{9}{151.794}{260.941}{\Character{40}}%
\TextAt{9}{160.234}{260.941}{\Character{41}}%
\TextAt{8}{168.693}{260.941}{to}%
\TextAt{8}{183.004}{260.941}{the}%
\TextAt{8}{203.16}{260.941}{ideal}%
\TextAt{9}{165.676}{221.104}{\Character{40}}%
\TextAt{13}{170.228}{221.104}{N}%
\TextAt{21}{179.602}{219.245}{L}%
\TextAt{3}{185.361}{219.245}{\Character{106}}%
\TextAt{21}{187.714}{219.245}{K}%
\TextAt{9}{195.916}{221.104}{\Character{40}}%
\TextAt{13}{200.469}{221.104}{\Character{11}}%
\TextAt{9}{207.99}{221.104}{\Character{41}\Character{41}}%
\TextAt{9}{220.388}{221.104}{\Character{61}}%
\TextAt{20}{232.841}{232.462}{\Character{89}}%
\TextAt{27}{233.144}{206.63}{p}%
\TextAt{25}{237.133}{206.63}{\Character{45}}%
\NativeMathAt{239.347}{206.63}{7.97}{\infty}%
\TextAt{20}{250.11}{232.462}{\Character{89}}%
\TextAt{27}{251.28}{206.559}{P}%
\TextAt{3}{257.875}{206.559}{\Character{106}}%
\TextAt{27}{260.228}{206.559}{p}%
\TextAt{14}{267.379}{221.104}{p}%
\TextAt{21}{273.362}{226.041}{f}%
\TextAt{28}{277.43}{224.303}{P}%
\TextAt{23}{283.129}{224.303}{\Character{106}}%
\TextAt{28}{285.398}{224.303}{p}%
\TextAt{21}{289.481}{226.041}{v}%
\TextAt{28}{293.604}{224.65}{P}%
\TextAt{17}{299.801}{226.041}{\Character{40}}%
\TextAt{21}{303.095}{226.041}{\Character{11}}%
\TextAt{28}{308.534}{224.65}{P}%
\TextAt{17}{314.731}{226.041}{\Character{41}}%
\TextAt{9}{321.843}{221.104}{\Character{61}}%
\TextAt{13}{334.268}{221.104}{N}%
\TextAt{21}{343.642}{219.245}{L}%
\TextAt{3}{349.401}{219.245}{\Character{106}}%
\TextAt{21}{351.754}{219.245}{K}%
\TextAt{20}{361.949}{245.135}{\Character{48}}%
\TextAt{20}{361.949}{223.615}{\Character{64}}%
\TextAt{20}{373.41}{232.462}{\Character{89}}%
\TextAt{27}{372.41}{206.63}{P}%
\TextAt{25}{379.005}{206.63}{\Character{45}}%
\NativeMathAt{381.219}{206.63}{7.97}{\infty}%
\TextAt{14}{391.68}{221.104}{P}%
\TextAt{21}{401.573}{226.041}{v}%
\TextAt{28}{405.697}{224.65}{P}%
\TextAt{17}{411.894}{226.041}{\Character{40}}%
\TextAt{21}{415.187}{226.041}{\Character{11}}%
\TextAt{28}{420.626}{224.65}{P}%
\TextAt{17}{426.823}{226.041}{\Character{41}}%
\TextAt{20}{430.615}{245.135}{\Character{49}}%
\TextAt{20}{430.615}{223.615}{\Character{65}}%
\TextAt{13}{443.068}{221.104}{\Character{58}}%
\TextAt{8}{79.226}{181.268}{Therefore}%
\TextAt{8}{131.471}{181.268}{the}%
\TextAt{8}{150.876}{181.268}{image}%
\TextAt{8}{183.945}{181.268}{of}%
\TextAt{13}{196.523}{181.268}{N}%
\TextAt{21}{205.897}{179.408}{L}%
\TextAt{3}{211.656}{179.408}{\Character{106}}%
\TextAt{21}{214.008}{179.408}{K}%
\TextAt{13}{222.211}{181.268}{C}%
\TextAt{21}{230.586}{179.475}{L}%
\TextAt{8}{239.989}{181.268}{under}%
\TextAt{8}{272.4}{181.268}{the}%
\TextAt{8}{291.805}{181.268}{homomorphism}%
\TextAt{9}{373.723}{181.268}{\Character{40}}%
\TextAt{9}{381.41}{181.268}{\Character{41}}%
\TextAt{9}{389.265}{181.268}{\Character{58}}%
\TextAt{13}{395.866}{181.268}{C}%
\TextAt{21}{404.241}{179.475}{K}%
\NativeArrowAt{415.765}{181.268}{11.955}{19.26}{\longrightarrow}%
\TextAt{13}{438.347}{181.268}{J}%
\TextAt{27}{445.924}{185.607}{m}%
\TextAt{21}{444.817}{178.248}{K}%
\TextAt{9}{453.02}{181.268}{\Character{47}}%
\TextAt{13}{458.873}{181.268}{P}%
\TextAt{27}{468.044}{185.607}{m}%
\TextAt{21}{466.418}{178.248}{K}%
\TextAt{8}{477.798}{181.268}{is}%
\TextAt{8}{488.81}{181.268}{precisely}%
\TextAt{8}{79.226}{164.222}{the}%
\TextAt{8}{99.383}{164.222}{group}%
\TextAt{9}{132.558}{164.222}{\Character{40}}%
\TextAt{13}{137.11}{164.222}{N}%
\TextAt{21}{146.484}{162.362}{L}%
\TextAt{3}{152.243}{162.362}{\Character{106}}%
\TextAt{21}{154.596}{162.362}{K}%
\TextAt{13}{162.798}{164.222}{J}%
\TextAt{27}{170.375}{168.561}{m}%
\TextAt{21}{169.269}{161.202}{L}%
\TextAt{9}{176.985}{164.222}{\Character{41}}%
\TextAt{13}{181.537}{164.222}{P}%
\TextAt{27}{190.708}{168.561}{m}%
\TextAt{21}{189.082}{161.202}{K}%
\TextAt{9}{197.317}{164.222}{\Character{47}}%
\TextAt{13}{203.17}{164.222}{P}%
\TextAt{27}{212.342}{168.561}{m}%
\TextAt{21}{210.716}{161.202}{K}%
\TextAt{8}{218.951}{164.222}{,}%
\TextAt{8}{226.1}{164.222}{q.e.d.}%
\Vector{q 0 G 0 g 0.399 w 452.334 -609.778 m 452.334 -601.708 l S Q}%
\Vector{q 0 G 0 g 0.399 w 452.534 -601.907 m 459.707 -601.907 l S Q}%
\Vector{q 0 G 0 g 0.399 w 452.534 -609.579 m 459.707 -609.579 l S Q}%
\Vector{q 0 G 0 g 0.399 w 459.906 -609.778 m 459.906 -601.708 l S Q}%
\TextAt{12}{79.226}{128.042}{Corollary}%
\TextAt{12}{138.71}{128.042}{(7.2).}%
\TextAt{8}{173.762}{128.042}{The}%
\TextAt{8}{197.551}{128.042}{Artin}%
\TextAt{8}{228.825}{128.042}{symbol}%
\TextAt{20}{268.623}{141.313}{\Character{16}}%
\TextAt{21}{276.958}{133.741}{L}%
\TextAt{3}{282.717}{133.741}{\Character{106}}%
\TextAt{21}{285.07}{133.741}{K}%
\Vector{q 0 G 0 g 0.478 w 204.958 -642.969 m 220.774 -642.969 l S Q}%
\TextAt{27}{282.872}{123.92}{a}%
\TextAt{20}{293.97}{141.313}{\Character{17}}%
\TextAt{8}{301.11}{128.042}{,}%
\TextAt{8}{308.048}{128.042}{for}%
\TextAt{14}{325.669}{128.042}{a}%
\NativeMathAt{334.972}{128.042}{11.955}{\in}%
\TextAt{13}{346.263}{128.042}{J}%
\TextAt{27}{353.841}{132.381}{m}%
\TextAt{21}{352.734}{125.022}{K}%
\TextAt{8}{360.936}{128.042}{,}%
\TextAt{8}{367.874}{128.042}{only}%
\TextAt{8}{393.302}{128.042}{depends}%
\TextAt{8}{438.288}{128.042}{on}%
\TextAt{8}{454.283}{128.042}{the}%
\TextAt{8}{474.175}{128.042}{class}%
\TextAt{14}{501.349}{128.042}{a}%
\TextAt{8}{510.652}{128.042}{mod}%
\TextAt{13}{79.226}{110.996}{P}%
\TextAt{27}{88.397}{115.335}{m}%
\TextAt{21}{86.771}{107.976}{K}%
\TextAt{8}{95.007}{110.996}{.}%
\TextAt{8}{103.455}{110.996}{It}%
\TextAt{8}{116.127}{110.996}{defines}%
\TextAt{8}{154.551}{110.996}{an}%
\TextAt{8}{170.81}{110.996}{isomorphism}%
\TextAt{20}{222.729}{95.158}{\Character{18}}%
\TextAt{13}{232.725}{86.389}{L}%
\TextAt{16}{244.011}{86.389}{\Character{106}}%
\TextAt{13}{250.652}{86.389}{K}%
\Vector{q 0 G 0 g 0.478 w 160.725 -692.71 m 189.47 -692.71 l S Q}%
\TextAt{20}{262.666}{95.158}{\Character{19}}%
\TextAt{9}{274.787}{78.301}{\Character{58}}%
\TextAt{13}{281.359}{78.301}{J}%
\TextAt{27}{288.936}{83.238}{m}%
\TextAt{21}{287.829}{75.346}{K}%
\TextAt{9}{296.032}{78.301}{\Character{47}}%
\TextAt{13}{301.885}{78.301}{H}%
\TextAt{27}{312.508}{83.238}{m}%
\NativeMathAt{325.025}{85.078}{7.97}{\sim}%
\NativeArrowAt{322.438}{78.301}{11.955}{13.175}{\longrightarrow}%
\TextAt{13}{338.932}{78.301}{G}%
\TextAt{9}{348.165}{78.301}{\Character{40}}%
\TextAt{13}{352.717}{78.301}{L}%
\TextAt{16}{364.003}{78.301}{\Character{106}}%
\TextAt{13}{370.645}{78.301}{K}%
\TextAt{9}{381.462}{78.301}{\Character{41}}%
\TextAt{13}{386.015}{78.301}{\Character{58}}%
\TextAt{8}{297.211}{43.033}{372}%
\Destination{78.23}{809.815}{nav.page.372}%
\Destination{79.36}{634.655}{nav.print.408}%
\Destination{79.36}{142.055}{nav.number.6.7.2}%
\Hotlink{136.322}{758.955}{6.858}{11.047}{nav.section.6.1}%
\Hotlink{390.219}{275.755}{14.306}{11.047}{nav.chapter.3}%
\Hotlink{461.417}{458.855}{25.078}{11.047}{nav.number.6.1.9}%
\Hotlink{101.151}{275.755}{25.078}{11.047}{nav.number.6.2.2}%
\Hotlink{409.632}{275.755}{25.078}{11.047}{nav.number.3.1.2}%
\Destination{79.226}{761.078}{errata.144}%
\Destination{79.226}{735}{errata.198}%
}

% Page 385 (body page 373)
\Page{%
\TextAt{2}{79.226}{796.943}{VI.}%
\TextAt{2}{98.012}{796.943}{Global}%
\TextAt{2}{133.248}{796.943}{Class}%
\TextAt{2}{161.841}{796.943}{Field}%
\TextAt{2}{189.572}{796.943}{Theory}%
\Vector{q 0 G 0 g 0.299 w 7.226 18.729 m 460.769 18.729 l S Q}%
\TextAt{8}{97.159}{761.124}{The}%
\TextAt{8}{120.908}{761.124}{group}%
\TextAt{13}{153.778}{761.124}{H}%
\TextAt{27}{164.401}{765.462}{m}%
\TextAt{9}{174.331}{761.124}{\Character{61}}%
\TextAt{9}{186.741}{761.124}{\Character{40}}%
\TextAt{13}{191.309}{761.124}{N}%
\TextAt{21}{200.683}{759.264}{L}%
\TextAt{3}{206.442}{759.264}{\Character{106}}%
\TextAt{21}{208.794}{759.264}{K}%
\TextAt{13}{216.997}{761.124}{J}%
\TextAt{27}{224.574}{765.462}{m}%
\TextAt{21}{223.467}{758.103}{L}%
\TextAt{9}{231.183}{761.124}{\Character{41}}%
\TextAt{13}{235.735}{761.124}{P}%
\TextAt{27}{244.907}{765.462}{m}%
\TextAt{21}{243.281}{758.103}{K}%
\TextAt{8}{255.108}{761.124}{is}%
\TextAt{8}{266.567}{761.124}{called}%
\TextAt{8}{299.426}{761.124}{the}%
\TextAt{8}{319.277}{761.124}{“ideal}%
\TextAt{8}{352.459}{761.124}{group}%
\TextAt{8}{385.33}{761.124}{defined}%
\TextAt{8}{425.337}{761.124}{mod}%
\TextAt{14}{451.37}{761.124}{m}%
\TextAt{8}{460.537}{761.124}{”}%
\TextAt{8}{469.652}{761.124}{belonging}%
\TextAt{8}{522.357}{761.124}{to}%
\TextAt{8}{79.226}{744.078}{the}%
\TextAt{8}{99.81}{744.078}{extension}%
\TextAt{13}{152.003}{744.078}{L}%
\TextAt{16}{164.016}{744.078}{\Character{106}}%
\TextAt{13}{171.386}{744.078}{K}%
\TextAt{8}{182.203}{744.078}{.}%
\TextAt{8}{191.933}{744.078}{From}%
\TextAt{8}{223.085}{744.078}{the}%
\TextAt{8}{243.669}{744.078}{existence}%
\TextAt{8}{293.901}{744.078}{theorem}%
\TextAt{8}{339.854}{744.078}{(6.1),}%
\TextAt{8}{371.615}{744.078}{we}%
\TextAt{8}{389.269}{744.078}{see}%
\TextAt{8}{408.609}{744.078}{that}%
\TextAt{8}{434.406}{744.078}{the}%
\TextAt{8}{454.989}{744.078}{correspondence}%
\TextAt{13}{79.226}{727.032}{L}%
\NativeArrowAt{90.793}{727.032}{11.955}{11.955}{\mapsto}%
\TextAt{13}{106.351}{727.032}{H}%
\TextAt{27}{116.974}{731.37}{m}%
\TextAt{8}{127.646}{727.032}{is}%
\TextAt{9}{139.575}{727.032}{1}%
\TextAt{8}{145.428}{727.032}{-}%
\TextAt{9}{149.325}{727.032}{1}%
\TextAt{8}{159.241}{727.032}{between}%
\TextAt{8}{204.597}{727.032}{subextensions}%
\TextAt{8}{279.088}{727.032}{of}%
\TextAt{8}{292.583}{727.032}{the}%
\TextAt{8}{312.905}{727.032}{ray}%
\TextAt{8}{333.239}{727.032}{class}%
\TextAt{8}{360.841}{727.032}{field}%
\TextAt{8}{386.364}{727.032}{mod}%
\TextAt{14}{412.866}{727.032}{m}%
\TextAt{8}{426.096}{727.032}{and}%
\TextAt{8}{449.024}{727.032}{subgroups}%
\TextAt{8}{504.601}{727.032}{of}%
\TextAt{13}{518.097}{727.032}{J}%
\TextAt{27}{525.674}{731.37}{m}%
\TextAt{21}{524.567}{724.011}{K}%
\TextAt{8}{79.226}{709.986}{containing}%
\TextAt{13}{136.145}{709.986}{P}%
\TextAt{27}{145.316}{714.324}{m}%
\TextAt{21}{143.69}{706.965}{K}%
\TextAt{8}{151.925}{709.986}{.}%
\TextAt{8}{97.159}{689.951}{The}%
\TextAt{8}{121.317}{689.951}{most}%
\TextAt{8}{150.102}{689.951}{important}%
\TextAt{8}{205.498}{689.951}{consequence}%
\TextAt{8}{271.667}{689.951}{of}%
\TextAt{8}{285.101}{689.951}{theorem}%
\TextAt{8}{330.73}{689.951}{(7.1)}%
\TextAt{8}{358.81}{689.951}{is}%
\TextAt{8}{370.678}{689.951}{a}%
\TextAt{8}{380.537}{689.951}{precise}%
\TextAt{8}{419.066}{689.951}{analysis}%
\TextAt{8}{463.201}{689.951}{of}%
\TextAt{8}{476.635}{689.951}{the}%
\TextAt{8}{496.895}{689.951}{kind}%
\TextAt{8}{523.337}{689.951}{of}%
\TextAt{8}{79.226}{672.905}{decomposition}%
\TextAt{8}{156.657}{672.905}{of}%
\TextAt{8}{170.283}{672.905}{any}%
\TextAt{8}{192.696}{672.905}{unramified}%
\TextAt{8}{251.524}{672.905}{prime}%
\TextAt{8}{284.984}{672.905}{ideal}%
\TextAt{14}{313.243}{672.905}{p}%
\TextAt{8}{323.42}{672.905}{in}%
\TextAt{8}{337.368}{672.905}{an}%
\TextAt{8}{353.923}{672.905}{abelian}%
\TextAt{8}{394.867}{672.905}{extension}%
\TextAt{13}{446.929}{672.905}{L}%
\TextAt{16}{458.719}{672.905}{\Character{106}}%
\TextAt{13}{465.865}{672.905}{K}%
\TextAt{8}{476.682}{672.905}{.}%
\TextAt{8}{486.019}{672.905}{It}%
\TextAt{8}{498.987}{672.905}{can}%
\TextAt{8}{520.743}{672.905}{be}%
\TextAt{8}{79.226}{655.859}{immediately}%
\TextAt{8}{145.505}{655.859}{read}%
\TextAt{8}{171.135}{655.859}{off}%
\TextAt{8}{187.334}{655.859}{the}%
\TextAt{8}{207.106}{655.859}{ideal}%
\TextAt{8}{234.686}{655.859}{group}%
\TextAt{13}{267.478}{655.859}{H}%
\TextAt{27}{278.101}{660.197}{m}%
\NativeMathAt{288.031}{655.859}{11.955}{\subseteq}%
\TextAt{13}{300.65}{655.859}{J}%
\TextAt{27}{308.227}{660.197}{m}%
\TextAt{21}{307.12}{652.839}{K}%
\TextAt{8}{318.837}{655.859}{which}%
\TextAt{8}{351.939}{655.859}{determines}%
\TextAt{8}{410.794}{655.859}{the}%
\TextAt{8}{430.567}{655.859}{field}%
\TextAt{13}{455.54}{655.859}{L}%
\TextAt{8}{467.018}{655.859}{as}%
\TextAt{8}{481.005}{655.859}{class}%
\TextAt{8}{508.058}{655.859}{field.}%
\TextAt{12}{79.226}{630.537}{Theorem}%
\TextAt{12}{136.75}{630.537}{(7.3)}%
\TextAt{12}{169.419}{630.537}{(Decomposition}%
\TextAt{12}{266.813}{630.537}{Law).}%
\TextAt{8}{304.102}{630.537}{Let}%
\TextAt{13}{325.514}{630.537}{L}%
\TextAt{16}{337.555}{630.537}{\Character{106}}%
\TextAt{13}{344.951}{630.537}{K}%
\TextAt{8}{360.11}{630.537}{be}%
\TextAt{8}{376.477}{630.537}{an}%
\TextAt{8}{393.179}{630.537}{abelian}%
\TextAt{8}{434.269}{630.537}{extension}%
\TextAt{8}{486.478}{630.537}{of}%
\TextAt{8}{500.251}{630.537}{degree}%
\TextAt{13}{79.226}{613.491}{n}%
\TextAt{8}{86.214}{613.491}{,}%
\TextAt{8}{93.134}{613.491}{and}%
\TextAt{8}{115.611}{613.491}{let}%
\TextAt{14}{132.23}{613.491}{p}%
\TextAt{8}{141.824}{613.491}{be}%
\TextAt{8}{157.463}{613.491}{an}%
\TextAt{8}{173.436}{613.491}{unramified}%
\TextAt{8}{231.682}{613.491}{prime}%
\TextAt{8}{264.56}{613.491}{ideal.}%
\TextAt{8}{296.979}{613.491}{Let}%
\TextAt{14}{317.662}{613.491}{m}%
\TextAt{8}{330.441}{613.491}{be}%
\TextAt{8}{346.079}{613.491}{a}%
\TextAt{8}{355.548}{613.491}{module}%
\TextAt{8}{396.556}{613.491}{of}%
\TextAt{8}{409.6}{613.491}{definition}%
\TextAt{8}{461.343}{613.491}{for}%
\TextAt{13}{478.942}{613.491}{L}%
\TextAt{16}{490.227}{613.491}{\Character{106}}%
\TextAt{13}{496.869}{613.491}{K}%
\TextAt{8}{511.298}{613.491}{that}%
\TextAt{8}{79.226}{596.445}{is}%
\TextAt{8}{91.132}{596.445}{not}%
\TextAt{8}{112.088}{596.445}{divisible}%
\TextAt{8}{158.138}{596.445}{by}%
\TextAt{14}{174.539}{596.445}{p}%
\TextAt{8}{184.561}{596.445}{(for}%
\TextAt{8}{207.143}{596.445}{instance}%
\TextAt{8}{252.87}{596.445}{the}%
\TextAt{8}{273.169}{596.445}{conductor),}%
\TextAt{8}{335.788}{596.445}{and}%
\TextAt{8}{358.693}{596.445}{let}%
\TextAt{13}{375.74}{596.445}{H}%
\TextAt{27}{386.362}{600.783}{m}%
\TextAt{8}{397.011}{596.445}{be}%
\TextAt{8}{413.077}{596.445}{the}%
\TextAt{8}{433.376}{596.445}{corresponding}%
\TextAt{8}{508.704}{596.445}{ideal}%
\TextAt{8}{79.226}{579.399}{group.}%
\TextAt{8}{97.159}{559.364}{If}%
\TextAt{13}{109.837}{559.364}{f}%
\TextAt{8}{121.766}{559.364}{is}%
\TextAt{8}{134.516}{559.364}{the}%
\TextAt{8}{155.658}{559.364}{order}%
\TextAt{8}{187.213}{559.364}{of}%
\TextAt{14}{201.529}{559.364}{p}%
\TextAt{8}{210.833}{559.364}{mod}%
\TextAt{13}{236.27}{559.364}{H}%
\TextAt{27}{246.893}{563.702}{m}%
\TextAt{8}{258.385}{559.364}{in}%
\TextAt{8}{273.024}{559.364}{the}%
\TextAt{8}{294.166}{559.364}{class}%
\TextAt{8}{322.589}{559.364}{group}%
\TextAt{13}{356.75}{559.364}{J}%
\TextAt{27}{364.327}{563.702}{m}%
\TextAt{21}{363.22}{556.343}{K}%
\TextAt{9}{371.423}{559.364}{\Character{47}}%
\TextAt{13}{377.276}{559.364}{H}%
\TextAt{27}{387.898}{563.702}{m}%
\TextAt{8}{394.508}{559.364}{,}%
\TextAt{8}{402.889}{559.364}{i.e.,}%
\TextAt{8}{426.226}{559.364}{the}%
\TextAt{8}{447.368}{559.364}{smallest}%
\TextAt{8}{493.353}{559.364}{positive}%
\TextAt{8}{79.226}{542.318}{integer}%
\TextAt{8}{117.925}{542.318}{such}%
\TextAt{8}{144.322}{542.318}{that}%
\TextAt{14}{282.736}{525.272}{p}%
\TextAt{21}{288.719}{530.208}{f}%
\NativeMathAt{297.485}{525.272}{11.955}{\in}%
\TextAt{13}{308.776}{525.272}{H}%
\TextAt{27}{319.399}{530.208}{m}%
\TextAt{13}{326.008}{525.272}{\Character{59}}%
\TextAt{8}{79.226}{502.785}{then}%
\TextAt{14}{105.886}{502.785}{p}%
\TextAt{8}{115.767}{502.785}{decomposes}%
\TextAt{8}{179.296}{502.785}{in}%
\TextAt{13}{192.949}{502.785}{L}%
\TextAt{8}{204.811}{502.785}{into}%
\TextAt{8}{228.554}{502.785}{a}%
\TextAt{8}{238.31}{502.785}{product}%
\TextAt{14}{271.63}{476.199}{p}%
\TextAt{9}{280.934}{476.199}{\Character{61}}%
\TextAt{14}{293.359}{476.199}{P}%
\TextAt{17}{303.253}{474.406}{1}%
\TextAt{16}{309.978}{476.199}{\Character{1}}%
\TextAt{16}{315.274}{476.199}{\Character{1}}%
\TextAt{16}{320.582}{476.199}{\Character{1}}%
\TextAt{14}{325.918}{476.199}{P}%
\TextAt{21}{335.811}{474.406}{r}%
\TextAt{8}{79.226}{449.612}{of}%
\TextAt{13}{92.556}{449.612}{r}%
\TextAt{9}{101.477}{449.612}{\Character{61}}%
\TextAt{13}{113.903}{449.612}{n}%
\TextAt{9}{120.891}{449.612}{\Character{47}}%
\TextAt{13}{126.744}{449.612}{f}%
\TextAt{8}{137.687}{449.612}{distinct}%
\TextAt{8}{180.021}{449.612}{prime}%
\TextAt{8}{213.184}{449.612}{ideals}%
\TextAt{8}{245.762}{449.612}{of}%
\TextAt{8}{259.092}{449.612}{degree}%
\TextAt{13}{295.508}{449.612}{f}%
\TextAt{8}{306.451}{449.612}{over}%
\TextAt{14}{331.497}{449.612}{p}%
\TextAt{8}{337.481}{449.612}{.}%
\TextAt{19}{79.226}{416.014}{Proof.}%
\TextAt{8}{116}{416.014}{Let}%
\TextAt{14}{137.699}{416.014}{p}%
\TextAt{9}{148.245}{416.014}{\Character{61}}%
\TextAt{14}{161.913}{416.014}{P}%
\TextAt{17}{171.807}{414.22}{1}%
\TextAt{16}{178.532}{416.014}{\Character{1}}%
\TextAt{16}{183.828}{416.014}{\Character{1}}%
\TextAt{16}{189.136}{416.014}{\Character{1}}%
\TextAt{14}{194.472}{416.014}{P}%
\TextAt{21}{204.365}{414.22}{r}%
\TextAt{8}{213.546}{416.014}{be}%
\TextAt{8}{230.2}{416.014}{the}%
\TextAt{8}{251.085}{416.014}{prime}%
\TextAt{8}{284.978}{416.014}{decomposition}%
\TextAt{8}{362.842}{416.014}{of}%
\TextAt{14}{376.901}{416.014}{p}%
\TextAt{8}{387.511}{416.014}{in}%
\TextAt{13}{401.893}{416.014}{L}%
\TextAt{8}{409.857}{416.014}{.}%
\TextAt{8}{420.493}{416.014}{Since}%
\TextAt{14}{451.78}{416.014}{p}%
\TextAt{8}{462.39}{416.014}{is}%
\TextAt{8}{474.883}{416.014}{unramified,}%
\TextAt{8}{79.226}{398.968}{the}%
\TextAt{14}{101.054}{398.968}{P}%
\TextAt{21}{110.947}{397.174}{i}%
\TextAt{8}{119.896}{398.968}{are}%
\TextAt{8}{141.078}{398.968}{all}%
\TextAt{8}{159.008}{398.968}{distinct}%
\TextAt{8}{203.012}{398.968}{and}%
\TextAt{8}{227.446}{398.968}{have}%
\TextAt{8}{256.112}{398.968}{the}%
\TextAt{8}{277.939}{398.968}{same}%
\TextAt{8}{308.936}{398.968}{degree}%
\TextAt{13}{347.022}{398.968}{f}%
\TextAt{8}{354.069}{398.968}{.}%
\TextAt{8}{367.53}{398.968}{This}%
\TextAt{8}{395.921}{398.968}{degree}%
\TextAt{8}{434.007}{398.968}{is}%
\TextAt{8}{447.442}{398.968}{the}%
\TextAt{8}{469.269}{398.968}{order}%
\TextAt{8}{501.509}{398.968}{of}%
\TextAt{8}{516.51}{398.968}{the}%
\TextAt{8}{79.226}{381.922}{decomposition}%
\TextAt{8}{156.853}{381.922}{group}%
\TextAt{8}{190.521}{381.922}{of}%
\TextAt{14}{204.343}{381.922}{P}%
\TextAt{21}{214.236}{380.128}{i}%
\TextAt{8}{222.006}{381.922}{over}%
\TextAt{13}{247.545}{381.922}{K}%
\TextAt{8}{258.362}{381.922}{,}%
\TextAt{8}{266.126}{381.922}{i.e.,}%
\TextAt{8}{288.847}{381.922}{the}%
\TextAt{8}{309.495}{381.922}{order}%
\TextAt{8}{340.556}{381.922}{of}%
\TextAt{8}{354.378}{381.922}{the}%
\TextAt{8}{375.027}{381.922}{Frobenius}%
\TextAt{8}{429.389}{381.922}{automorphism}%
\TextAt{13}{507.35}{381.922}{\Character{39}}%
\TextAt{27}{515.019}{380.128}{p}%
\TextAt{9}{523.665}{381.922}{\Character{61}}%
\TextAt{20}{79.226}{377.147}{\Character{16}}%
\TextAt{21}{87.562}{369.575}{L}%
\TextAt{3}{93.321}{369.575}{\Character{106}}%
\TextAt{21}{95.673}{369.575}{K}%
\Vector{q 0 G 0 g 0.478 w 15.562 -407.134 m 31.377 -407.134 l S Q}%
\TextAt{27}{93.475}{359.754}{p}%
\TextAt{20}{104.573}{377.147}{\Character{17}}%
\TextAt{8}{111.713}{363.877}{.}%
\TextAt{8}{119.806}{363.877}{In}%
\TextAt{8}{133.362}{363.877}{view}%
\TextAt{8}{159.279}{363.877}{of}%
\TextAt{8}{171.544}{363.877}{the}%
\TextAt{8}{190.635}{363.877}{isomorphism}%
\TextAt{13}{257.989}{363.877}{J}%
\TextAt{27}{265.566}{368.215}{m}%
\TextAt{21}{264.459}{360.856}{K}%
\TextAt{9}{272.662}{363.877}{\Character{47}}%
\TextAt{13}{278.515}{363.877}{H}%
\TextAt{27}{289.137}{368.215}{m}%
\NativeMathAt{299.067}{363.877}{11.955}{\cong}%
\TextAt{13}{311.687}{363.877}{G}%
\TextAt{9}{320.92}{363.877}{\Character{40}}%
\TextAt{13}{325.472}{363.877}{L}%
\TextAt{16}{336.758}{363.877}{\Character{106}}%
\TextAt{13}{343.4}{363.877}{K}%
\TextAt{9}{354.217}{363.877}{\Character{41}}%
\TextAt{8}{358.77}{363.877}{,}%
\TextAt{8}{365.066}{363.877}{this}%
\TextAt{8}{386.823}{363.877}{is}%
\TextAt{8}{397.522}{363.877}{also}%
\TextAt{8}{419.936}{363.877}{the}%
\TextAt{8}{439.027}{363.877}{order}%
\TextAt{8}{468.531}{363.877}{of}%
\TextAt{14}{480.796}{363.877}{p}%
\TextAt{8}{490.1}{363.877}{mod}%
\TextAt{13}{515.538}{363.877}{H}%
\TextAt{27}{526.16}{368.215}{m}%
\TextAt{8}{79.226}{346.831}{in}%
\TextAt{13}{92.879}{346.831}{J}%
\TextAt{27}{100.456}{351.169}{m}%
\TextAt{21}{99.349}{343.81}{K}%
\TextAt{9}{107.552}{346.831}{\Character{47}}%
\TextAt{13}{113.405}{346.831}{H}%
\TextAt{27}{124.027}{351.169}{m}%
\TextAt{8}{130.637}{346.831}{.}%
\TextAt{8}{139.085}{346.831}{This}%
\TextAt{8}{165.805}{346.831}{finishes}%
\TextAt{8}{206.895}{346.831}{the}%
\TextAt{8}{227.051}{346.831}{proof.}%
\Vector{q 0 G 0 g 0.399 w 452.334 -427.169 m 452.334 -419.099 l S Q}%
\Vector{q 0 G 0 g 0.399 w 452.534 -419.299 m 459.707 -419.299 l S Q}%
\Vector{q 0 G 0 g 0.399 w 452.534 -426.97 m 459.707 -426.97 l S Q}%
\Vector{q 0 G 0 g 0.399 w 459.906 -427.169 m 459.906 -419.099 l S Q}%
\TextAt{8}{97.159}{318.52}{The}%
\TextAt{8}{122.836}{318.52}{theorem}%
\TextAt{8}{169.984}{318.52}{shows}%
\TextAt{8}{205.225}{318.52}{in}%
\TextAt{8}{220.5}{318.52}{particular}%
\TextAt{8}{276.113}{318.52}{that}%
\TextAt{8}{303.105}{318.52}{the}%
\TextAt{8}{324.884}{318.52}{prime}%
\TextAt{8}{359.671}{318.52}{ideals}%
\TextAt{8}{393.871}{318.52}{which}%
\TextAt{8}{428.981}{318.52}{split}%
\TextAt{8}{456.678}{318.52}{completely}%
\TextAt{8}{517.156}{318.52}{are}%
\TextAt{8}{79.226}{301.474}{precisely}%
\TextAt{8}{127.083}{301.474}{those}%
\TextAt{8}{157.712}{301.474}{contained}%
\TextAt{8}{210.721}{301.474}{in}%
\TextAt{8}{224.374}{301.474}{the}%
\TextAt{8}{244.53}{301.474}{ideal}%
\TextAt{8}{272.493}{301.474}{group}%
\TextAt{13}{305.669}{301.474}{H}%
\TextAt{27}{316.292}{305.812}{f}%
\TextAt{8}{319.39}{301.474}{,}%
\TextAt{8}{326.539}{301.474}{if}%
\TextAt{14}{337.263}{301.474}{f}%
\TextAt{8}{345.061}{301.474}{is}%
\TextAt{8}{356.825}{301.474}{the}%
\TextAt{8}{376.981}{301.474}{conductor}%
\TextAt{8}{431.616}{301.474}{of}%
\TextAt{13}{444.946}{301.474}{L}%
\TextAt{16}{456.232}{301.474}{\Character{106}}%
\TextAt{13}{462.873}{301.474}{K}%
\TextAt{8}{473.691}{301.474}{.}%
\TextAt{8}{97.159}{281.439}{Let}%
\TextAt{8}{119.325}{281.439}{us}%
\TextAt{8}{135.538}{281.439}{highlight}%
\TextAt{8}{185.847}{281.439}{two}%
\TextAt{8}{209.161}{281.439}{special}%
\TextAt{8}{248.459}{281.439}{cases.}%
\TextAt{8}{285.987}{281.439}{If}%
\TextAt{8}{298.876}{281.439}{the}%
\TextAt{8}{320.23}{281.439}{base}%
\TextAt{8}{347.501}{281.439}{field}%
\TextAt{8}{374.055}{281.439}{is}%
\TextAt{13}{387.016}{281.439}{K}%
\TextAt{9}{403.195}{281.439}{\Character{61}}%
\TextAt{18}{417.66}{281.439}{Q}%
\TextAt{8}{432.053}{281.439}{and}%
\TextAt{8}{456.013}{281.439}{we}%
\TextAt{8}{474.437}{281.439}{look}%
\TextAt{8}{501.003}{281.439}{at}%
\TextAt{8}{516.51}{281.439}{the}%
\ErrataAt{79.22}{264.355}{453.64}{cyclotomic field $\mathbb Q(\mu_m)\mid\mathbb Q$, a module of definition is $\mathfrak m=(m)$, and the ideal group}%
\TextAt{8}{79.226}{247.347}{corresponding}%
\TextAt{8}{154.862}{247.347}{to}%
\TextAt{18}{169.623}{247.347}{Q}%
\TextAt{9}{178.922}{247.347}{\Character{40}}%
\TextAt{13}{183.474}{247.347}{\Character{22}}%
\TextAt{21}{190.517}{245.554}{m}%
\TextAt{9}{198.506}{247.347}{\Character{41}}%
\TextAt{8}{207.406}{247.347}{in}%
\TextAt{13}{221.509}{247.347}{J}%
\TextAt{27}{229.086}{251.685}{m}%
\TextAt{25}{227.979}{244.282}{Q}%
\TextAt{8}{240.043}{247.347}{is}%
\TextAt{8}{252.257}{247.347}{the}%
\TextAt{8}{272.864}{247.347}{group}%
\TextAt{13}{306.49}{247.347}{P}%
\TextAt{27}{315.661}{251.685}{m}%
\TextAt{25}{314.035}{244.282}{Q}%
\TextAt{8}{322.27}{247.347}{.}%
\TextAt{8}{332.069}{247.347}{As}%
\TextAt{13}{349.807}{247.347}{J}%
\TextAt{27}{357.384}{251.685}{m}%
\TextAt{25}{356.277}{244.282}{Q}%
\TextAt{9}{363.993}{247.347}{\Character{47}}%
\TextAt{13}{369.846}{247.347}{P}%
\TextAt{27}{379.017}{251.685}{m}%
\TextAt{25}{377.391}{244.282}{Q}%
\NativeMathAt{389.715}{247.347}{11.955}{\cong}%
\TextAt{9}{403.101}{247.347}{\Character{40}}%
\TextAt{18}{407.654}{247.347}{Z}%
\TextAt{9}{415.624}{247.347}{\Character{47}}%
\TextAt{13}{421.477}{247.347}{m}%
\TextAt{18}{431.716}{247.347}{Z}%
\TextAt{9}{439.686}{247.347}{\Character{41}}%
\TextAt{3}{444.238}{251.685}{\Character{3}}%
\TextAt{8}{453.318}{247.347}{(see}%
\TextAt{8}{477.237}{247.347}{(1.10)),}%
\TextAt{8}{519.439}{247.347}{we}%
\TextAt{8}{79.226}{230.301}{obtain}%
\TextAt{8}{116.965}{230.301}{for}%
\TextAt{8}{136.161}{230.301}{the}%
\TextAt{8}{157.628}{230.301}{decomposition}%
\TextAt{8}{236.074}{230.301}{of}%
\TextAt{8}{250.715}{230.301}{rational}%
\TextAt{8}{295.614}{230.301}{primes}%
\TextAt{13}{334.704}{230.301}{p}%
\TextAt{18}{346.134}{230.301}{\Character{45}}%
\TextAt{13}{355.01}{230.301}{m}%
\TextAt{8}{365.249}{230.301}{,}%
\TextAt{8}{374.037}{230.301}{the}%
\TextAt{8}{395.504}{230.301}{law}%
\TextAt{8}{417.952}{230.301}{which}%
\TextAt{8}{452.75}{230.301}{we}%
\TextAt{8}{471.288}{230.301}{had}%
\TextAt{8}{495.362}{230.301}{already}%
\TextAt{8}{79.226}{213.255}{deduced}%
\TextAt{8}{124.727}{213.255}{in}%
\TextAt{8}{138.368}{213.255}{chap.}%
\TextAt{8}{169.248}{213.255}{I,}%
\TextAt{8}{180.605}{213.255}{(10.4),}%
\TextAt{8}{217.681}{213.255}{and}%
\TextAt{8}{240.431}{213.255}{in}%
\TextAt{8}{254.072}{213.255}{particular}%
\TextAt{8}{308.049}{213.255}{the}%
\TextAt{8}{328.194}{213.255}{fact}%
\TextAt{8}{351.267}{213.255}{that}%
\TextAt{8}{376.624}{213.255}{the}%
\TextAt{8}{396.768}{213.255}{prime}%
\TextAt{8}{429.919}{213.255}{numbers}%
\TextAt{8}{477.118}{213.255}{which}%
\TextAt{8}{510.593}{213.255}{split}%
\TextAt{8}{79.226}{196.209}{completely}%
\TextAt{8}{138.081}{196.209}{are}%
\TextAt{8}{157.592}{196.209}{characterized}%
\TextAt{8}{228.809}{196.209}{by}%
\TextAt{13}{271.039}{179.163}{p}%
\NativeMathAt{280.235}{179.163}{11.955}{\equiv}%
\TextAt{9}{292.854}{179.163}{1}%
\TextAt{8}{302.028}{179.163}{mod}%
\TextAt{13}{327.466}{179.163}{m\Character{58}}%
\TextAt{8}{79.226}{156.676}{In}%
\TextAt{8}{93.419}{156.676}{the}%
\TextAt{8}{113.146}{156.676}{case}%
\TextAt{8}{137.489}{156.676}{of}%
\TextAt{8}{150.39}{156.676}{the}%
\TextAt{8}{170.118}{156.676}{Hilbert}%
\TextAt{8}{210.002}{156.676}{class}%
\TextAt{8}{237.01}{156.676}{field}%
\TextAt{13}{261.939}{156.676}{L}%
\TextAt{16}{273.224}{156.676}{\Character{106}}%
\TextAt{13}{279.866}{156.676}{K}%
\TextAt{8}{290.683}{156.676}{,}%
\TextAt{8}{297.49}{156.676}{i.e.,}%
\TextAt{8}{319.252}{156.676}{of}%
\TextAt{8}{332.153}{156.676}{the}%
\TextAt{8}{351.881}{156.676}{field}%
\TextAt{8}{376.809}{156.676}{inside}%
\TextAt{8}{409.604}{156.676}{the}%
\TextAt{8}{429.332}{156.676}{ray}%
\TextAt{8}{449.071}{156.676}{class}%
\TextAt{8}{476.08}{156.676}{field}%
\TextAt{8}{501.008}{156.676}{mod}%
\TextAt{9}{526.916}{156.676}{1}%
\TextAt{8}{79.226}{139.63}{in}%
\TextAt{8}{92.575}{139.63}{which}%
\TextAt{8}{125.757}{139.63}{the}%
\TextAt{8}{145.61}{139.63}{infinite}%
\TextAt{8}{184.973}{139.63}{places}%
\TextAt{8}{219.196}{139.63}{split}%
\TextAt{8}{244.966}{139.63}{completely,}%
\TextAt{8}{305.849}{139.63}{the}%
\TextAt{8}{325.702}{139.63}{corresponding}%
\TextAt{8}{400.584}{139.63}{ideal}%
\TextAt{8}{428.243}{139.63}{group}%
\TextAt{13}{461.115}{139.63}{H}%
\NativeMathAt{475.058}{139.63}{11.955}{\subseteq}%
\TextAt{13}{487.678}{139.63}{J}%
\TextAt{17}{495.255}{143.969}{1}%
\TextAt{21}{494.148}{136.61}{K}%
\TextAt{9}{505.671}{139.63}{\Character{61}}%
\TextAt{13}{518.097}{139.63}{J}%
\TextAt{21}{524.567}{137.837}{K}%
\TextAt{8}{79.226}{122.584}{is}%
\TextAt{8}{90.99}{122.584}{the}%
\TextAt{8}{111.146}{122.584}{group}%
\TextAt{13}{144.322}{122.584}{P}%
\TextAt{21}{151.867}{120.791}{K}%
\TextAt{8}{163.967}{122.584}{of}%
\TextAt{8}{177.297}{122.584}{principal}%
\TextAt{8}{226.074}{122.584}{ideals}%
\TextAt{8}{258.652}{122.584}{(see}%
\TextAt{8}{282.12}{122.584}{(6.9)).}%
\TextAt{8}{319.201}{122.584}{This}%
\TextAt{8}{345.921}{122.584}{gives}%
\TextAt{8}{374.602}{122.584}{us}%
\TextAt{8}{389.617}{122.584}{the}%
\TextAt{8}{409.774}{122.584}{strikingly}%
\TextAt{8}{461.874}{122.584}{simple}%
\TextAt{8}{297.211}{43.033}{373}%
\Destination{78.23}{809.815}{nav.page.373}%
\Destination{79.36}{645.555}{nav.print.409}%
\Destination{79.36}{644.555}{nav.number.6.7.3}%
\Hotlink{168.748}{211.055}{5.22}{11.047}{nav.chapter.1}%
\Hotlink{339.354}{741.855}{25.078}{11.047}{nav.number.6.6.1}%
\Hotlink{330.23}{687.755}{25.078}{11.047}{nav.number.6.7.1}%
\Hotlink{476.737}{245.155}{30.936}{11.047}{nav.number.6.1.10}%
\Hotlink{180.105}{211.055}{30.936}{11.047}{nav.number.1.10.4}%
\Hotlink{281.62}{120.355}{25.078}{11.047}{nav.number.6.6.9}%
\Destination{79.226}{278.355}{errata.34}%
}

% Page 386 (body page 374)
\Page{%
\TextAt{2}{79.226}{796.943}{VI.}%
\TextAt{2}{98.012}{796.943}{Global}%
\TextAt{2}{133.248}{796.943}{Class}%
\TextAt{2}{161.841}{796.943}{Field}%
\TextAt{2}{189.572}{796.943}{Theory}%
\Vector{q 0 G 0 g 0.299 w 7.226 18.729 m 460.769 18.729 l S Q}%
\TextAt{12}{79.226}{761.124}{Corollary}%
\TextAt{12}{138.946}{761.124}{(7.4).}%
\TextAt{8}{174.203}{761.124}{The}%
\TextAt{8}{198.197}{761.124}{prime}%
\TextAt{8}{231.302}{761.124}{ideals}%
\TextAt{8}{263.821}{761.124}{of}%
\TextAt{13}{277.092}{761.124}{K}%
\TextAt{8}{291.748}{761.124}{which}%
\TextAt{8}{325.175}{761.124}{split}%
\TextAt{8}{351.191}{761.124}{completely}%
\TextAt{8}{409.987}{761.124}{in}%
\TextAt{8}{423.581}{761.124}{the}%
\TextAt{8}{443.678}{761.124}{Hilbert}%
\TextAt{8}{483.932}{761.124}{class}%
\TextAt{8}{511.31}{761.124}{field}%
\TextAt{8}{79.226}{744.078}{are}%
\TextAt{8}{98.737}{744.078}{precisely}%
\TextAt{8}{146.594}{744.078}{the}%
\TextAt{8}{166.75}{744.078}{principal}%
\TextAt{8}{215.527}{744.078}{prime}%
\TextAt{8}{248.691}{744.078}{ideals.}%
\TextAt{8}{97.159}{718.225}{Another}%
\TextAt{8}{144.869}{718.225}{highly}%
\TextAt{8}{182.178}{718.225}{remarkable}%
\TextAt{8}{244.856}{718.225}{property}%
\TextAt{8}{294.527}{718.225}{of}%
\TextAt{8}{309.719}{718.225}{the}%
\TextAt{8}{331.737}{718.225}{Hilbert}%
\TextAt{8}{373.912}{718.225}{class}%
\TextAt{8}{403.211}{718.225}{field}%
\TextAt{8}{430.43}{718.225}{is}%
\TextAt{8}{444.056}{718.225}{expressed}%
\TextAt{8}{498.389}{718.225}{by}%
\TextAt{8}{516.51}{718.225}{the}%
\TextAt{8}{79.226}{701.179}{following}%
\TextAt{8}{128.661}{701.179}{theorem,}%
\TextAt{8}{177.438}{701.179}{known}%
\TextAt{8}{214.511}{701.179}{as}%
\TextAt{8}{228.881}{701.179}{the}%
\TextAt{12}{249.037}{701.179}{principal}%
\TextAt{12}{305.179}{701.179}{ideal}%
\TextAt{12}{337.29}{701.179}{theorem}%
\TextAt{8}{385.709}{701.179}{.}%
\TextAt{12}{79.226}{675.326}{Theorem}%
\TextAt{12}{136.24}{675.326}{(7.5).}%
\TextAt{8}{171.556}{675.326}{In}%
\TextAt{8}{186.177}{675.326}{the}%
\TextAt{8}{206.334}{675.326}{Hilbert}%
\TextAt{8}{246.646}{675.326}{class}%
\TextAt{8}{274.083}{675.326}{field}%
\TextAt{8}{299.44}{675.326}{every}%
\TextAt{8}{330.333}{675.326}{ideal}%
\TextAt{14}{358.296}{675.326}{a}%
\TextAt{8}{368.176}{675.326}{of}%
\TextAt{13}{381.506}{675.326}{K}%
\TextAt{8}{396.221}{675.326}{becomes}%
\TextAt{8}{442.775}{675.326}{a}%
\TextAt{8}{452.53}{675.326}{principal}%
\TextAt{8}{501.307}{675.326}{ideal.}%
\TextAt{19}{79.226}{640.666}{Proof.}%
\TextAt{8}{116}{640.666}{Let}%
\TextAt{13}{137.719}{640.666}{K}%
\TextAt{17}{147.678}{638.873}{1}%
\TextAt{16}{157.009}{640.666}{\Character{106}}%
\TextAt{13}{164.928}{640.666}{K}%
\TextAt{8}{180.393}{640.666}{be}%
\TextAt{8}{197.067}{640.666}{the}%
\TextAt{8}{217.973}{640.666}{Hilbert}%
\TextAt{8}{259.035}{640.666}{class}%
\TextAt{8}{287.222}{640.666}{field}%
\TextAt{8}{313.329}{640.666}{of}%
\TextAt{13}{327.408}{640.666}{K}%
\TextAt{8}{342.873}{640.666}{and}%
\TextAt{8}{366.386}{640.666}{let}%
\TextAt{13}{384.04}{640.666}{K}%
\TextAt{17}{393.999}{638.873}{2}%
\TextAt{16}{403.33}{640.666}{\Character{106}}%
\TextAt{13}{411.249}{640.666}{K}%
\TextAt{17}{421.208}{638.873}{1}%
\TextAt{8}{430.587}{640.666}{be}%
\TextAt{8}{447.261}{640.666}{the}%
\TextAt{8}{468.167}{640.666}{Hilbert}%
\TextAt{8}{509.23}{640.666}{class}%
\TextAt{8}{79.226}{623.62}{field}%
\TextAt{8}{104.583}{623.62}{of}%
\TextAt{13}{117.913}{623.62}{K}%
\TextAt{17}{127.872}{621.827}{1}%
\TextAt{8}{132.605}{623.62}{.}%
\TextAt{8}{141.053}{623.62}{We}%
\TextAt{8}{161.197}{623.62}{have}%
\TextAt{8}{188.192}{623.62}{to}%
\TextAt{8}{202.502}{623.62}{show}%
\TextAt{8}{231.506}{623.62}{that}%
\TextAt{8}{256.875}{623.62}{the}%
\TextAt{8}{277.031}{623.62}{canonical}%
\TextAt{8}{328.414}{623.62}{homomorphism}%
\TextAt{13}{253.185}{595.177}{J}%
\TextAt{21}{259.655}{593.384}{K}%
\TextAt{9}{267.858}{595.177}{\Character{47}}%
\TextAt{13}{273.711}{595.177}{P}%
\TextAt{21}{281.256}{593.384}{K}%
\NativeArrowAt{292.779}{595.177}{11.955}{19.26}{\longrightarrow}%
\TextAt{13}{315.361}{595.177}{J}%
\TextAt{21}{321.832}{593.384}{K}%
\TextAt{24}{328.973}{592.277}{1}%
\TextAt{9}{333.622}{595.177}{\Character{47}}%
\TextAt{13}{339.475}{595.177}{P}%
\TextAt{21}{347.02}{593.384}{K}%
\TextAt{24}{354.162}{592.277}{1}%
\TextAt{8}{79.226}{566.734}{is}%
\TextAt{8}{90.99}{566.734}{trivial.}%
\TextAt{8}{130.342}{566.734}{By}%
\TextAt{8}{148.706}{566.734}{chap.}%
\TextAt{8}{179.598}{566.734}{IV,}%
\TextAt{8}{199.742}{566.734}{(5.9),}%
\TextAt{8}{230.969}{566.734}{we}%
\TextAt{8}{248.197}{566.734}{have}%
\TextAt{8}{275.191}{566.734}{a}%
\TextAt{8}{284.947}{566.734}{commutative}%
\TextAt{8}{354.873}{566.734}{diagram}%
\TextAt{13}{190.281}{537.994}{J}%
\TextAt{21}{196.751}{536.201}{K}%
\TextAt{24}{203.893}{535.094}{1}%
\TextAt{9}{208.542}{537.994}{\Character{47}}%
\TextAt{13}{214.395}{537.994}{P}%
\TextAt{21}{221.94}{536.201}{K}%
\TextAt{24}{229.082}{535.094}{1}%
\TextAt{13}{258.382}{537.994}{C}%
\TextAt{21}{266.757}{536.201}{K}%
\TextAt{24}{273.899}{535.094}{1}%
\TextAt{9}{278.548}{537.994}{\Character{47}}%
\TextAt{13}{284.401}{537.994}{N}%
\TextAt{21}{293.775}{536.134}{K}%
\TextAt{24}{300.916}{535.027}{2}%
\TextAt{3}{305.067}{536.134}{\Character{106}}%
\TextAt{21}{307.419}{536.134}{K}%
\TextAt{24}{314.561}{535.027}{1}%
\TextAt{13}{319.21}{537.994}{C}%
\TextAt{21}{327.585}{536.201}{K}%
\TextAt{24}{334.726}{535.094}{2}%
\TextAt{13}{364.027}{537.994}{G}%
\TextAt{9}{373.261}{537.994}{\Character{40}}%
\TextAt{13}{377.813}{537.994}{K}%
\TextAt{17}{387.772}{536.201}{2}%
\TextAt{16}{395.825}{537.994}{\Character{106}}%
\TextAt{13}{402.467}{537.994}{K}%
\TextAt{17}{412.426}{536.201}{1}%
\TextAt{9}{417.158}{537.994}{\Character{41}}%
\TextAt{13}{193.869}{494.896}{J}%
\TextAt{21}{200.339}{493.103}{K}%
\TextAt{9}{208.542}{494.896}{\Character{47}}%
\TextAt{13}{214.395}{494.896}{P}%
\TextAt{21}{221.94}{493.103}{K}%
\TextAt{13}{261.97}{494.896}{C}%
\TextAt{21}{270.345}{493.103}{K}%
\TextAt{9}{278.548}{494.896}{\Character{47}}%
\TextAt{13}{284.401}{494.896}{N}%
\TextAt{21}{293.775}{493.037}{K}%
\TextAt{24}{300.916}{491.93}{1}%
\TextAt{3}{305.067}{493.037}{\Character{106}}%
\TextAt{21}{307.419}{493.037}{K}%
\TextAt{13}{315.622}{494.896}{C}%
\TextAt{21}{323.997}{493.103}{K}%
\TextAt{24}{331.138}{491.996}{1}%
\TextAt{13}{364.338}{494.896}{G}%
\TextAt{9}{373.572}{494.896}{\Character{40}}%
\TextAt{13}{378.124}{494.896}{K}%
\TextAt{17}{388.083}{493.103}{1}%
\TextAt{16}{396.136}{494.896}{\Character{106}}%
\TextAt{13}{402.778}{494.896}{K}%
\TextAt{9}{413.596}{494.896}{\Character{41}}%
\TextAt{13}{418.148}{494.896}{\Character{59}}%
\NativeMathAt{241.408}{537.38}{11.955}{\cong}%
\NativeMathAt{347.055}{537.38}{11.955}{\cong}%
\NativeDiagram{%
\path[draw=black,line width=0.478bp] (140.007,-265.559) -- (140.007,-244.052);%
\path[draw=black,line width=0.478bp,line cap=round,line join=round] (137.138,-246.3) .. controls (138.859,-245.845) and (139.672,-244.833) .. (140.007,-243.813) .. controls (140.342,-244.833) and (141.154,-245.845) .. (142.876,-246.3);%
}%
\NativeMathAt{241.408}{494.281}{11.955}{\cong}%
\NativeDiagram{%
\path[draw=black,line width=0.478bp] (226.881,-265.559) -- (226.881,-244.916);%
\path[draw=black,line width=0.478bp,line cap=round,line join=round] (224.012,-247.163) .. controls (225.733,-246.709) and (226.546,-245.696) .. (226.881,-244.677) .. controls (227.216,-245.696) and (228.028,-246.709) .. (229.75,-247.163);%
}%
\TextAt{21}{293.182}{516.368}{i}%
\NativeMathAt{345.416}{494.281}{11.955}{\cong}%
\NativeDiagram{%
\path[draw=black,line width=0.478bp] (320.872,-265.559) -- (320.872,-244.052);%
\path[draw=black,line width=0.478bp,line cap=round,line join=round] (318.003,-246.3) .. controls (319.724,-245.845) and (320.537,-244.833) .. (320.872,-243.813) .. controls (321.207,-244.833) and (322.02,-245.845) .. (323.741,-246.3);%
}%
\TextAt{7}{395.687}{516.796}{Ver}%
\TextAt{8}{79.226}{465.855}{where}%
\TextAt{13}{113.59}{465.855}{i}%
\TextAt{8}{122.036}{465.855}{is}%
\TextAt{8}{134.355}{465.855}{induced}%
\TextAt{8}{178.475}{465.855}{by}%
\TextAt{8}{195.289}{465.855}{the}%
\TextAt{8}{216.001}{465.855}{inclusion}%
\TextAt{13}{265.393}{465.855}{C}%
\TextAt{21}{273.768}{464.062}{K}%
\NativeMathAt{286.237}{465.855}{11.955}{\subseteq}%
\TextAt{13}{299.802}{465.855}{C}%
\TextAt{21}{308.177}{464.062}{K}%
\TextAt{24}{315.319}{462.955}{1}%
\TextAt{8}{319.968}{465.855}{.}%
\TextAt{8}{330.081}{465.855}{It}%
\TextAt{8}{343.309}{465.855}{is}%
\TextAt{8}{355.628}{465.855}{therefore}%
\TextAt{8}{405.283}{465.855}{enough}%
\TextAt{8}{446.163}{465.855}{to}%
\TextAt{8}{461.028}{465.855}{show}%
\TextAt{8}{490.586}{465.855}{that}%
\TextAt{8}{516.51}{465.855}{the}%
\TextAt{8}{79.226}{448.809}{transfer}%
\TextAt{8}{223.088}{431.763}{Ver}%
\TextAt{9}{244.939}{431.763}{\Character{58}}%
\TextAt{13}{251.512}{431.763}{G}%
\TextAt{9}{260.745}{431.763}{\Character{40}}%
\TextAt{13}{265.297}{431.763}{K}%
\TextAt{17}{275.257}{429.97}{1}%
\TextAt{16}{283.31}{431.763}{\Character{106}}%
\TextAt{13}{289.951}{431.763}{K}%
\TextAt{9}{300.769}{431.763}{\Character{41}}%
\NativeArrowAt{308.642}{431.763}{11.955}{19.26}{\longrightarrow}%
\TextAt{13}{331.224}{431.763}{G}%
\TextAt{9}{340.458}{431.763}{\Character{40}}%
\TextAt{13}{345.01}{431.763}{K}%
\TextAt{17}{354.969}{429.97}{2}%
\TextAt{16}{363.022}{431.763}{\Character{106}}%
\TextAt{13}{369.664}{431.763}{K}%
\TextAt{17}{379.623}{429.97}{1}%
\TextAt{9}{384.356}{431.763}{\Character{41}}%
\TextAt{8}{79.226}{408.48}{is}%
\TextAt{8}{91.167}{408.48}{the}%
\TextAt{8}{111.5}{408.48}{trivial}%
\TextAt{8}{146.478}{408.48}{homomorphism.}%
\TextAt{8}{234.228}{408.48}{Since}%
\TextAt{13}{264.962}{408.48}{K}%
\TextAt{17}{274.921}{406.687}{1}%
\TextAt{16}{283.275}{408.48}{\Character{106}}%
\TextAt{13}{290.218}{408.48}{K}%
\TextAt{8}{305.109}{408.48}{is}%
\TextAt{8}{317.05}{408.48}{the}%
\TextAt{8}{337.383}{408.48}{maximal}%
\TextAt{8}{385.368}{408.48}{unramified}%
\TextAt{8}{444.077}{408.48}{abelian}%
\TextAt{8}{484.901}{408.48}{extension}%
\TextAt{8}{79.226}{391.434}{of}%
\TextAt{13}{92.577}{391.434}{K}%
\TextAt{8}{107.312}{391.434}{in}%
\TextAt{8}{120.986}{391.434}{which}%
\TextAt{8}{154.492}{391.434}{the}%
\TextAt{8}{174.669}{391.434}{infinite}%
\TextAt{8}{214.357}{391.434}{places}%
\TextAt{8}{248.904}{391.434}{split}%
\TextAt{8}{274.999}{391.434}{completely,}%
\TextAt{8}{336.151}{391.434}{i.e.,}%
\TextAt{8}{358.282}{391.434}{the}%
\TextAt{8}{378.459}{391.434}{maximal}%
\TextAt{8}{426.288}{391.434}{abelian}%
\TextAt{8}{466.956}{391.434}{subextension}%
\TextAt{8}{79.226}{374.388}{of}%
\TextAt{13}{92.51}{374.388}{K}%
\TextAt{17}{102.469}{372.595}{2}%
\TextAt{16}{110.522}{374.388}{\Character{106}}%
\TextAt{13}{117.164}{374.388}{K}%
\TextAt{8}{127.982}{374.388}{,}%
\TextAt{8}{135.094}{374.388}{we}%
\TextAt{8}{152.276}{374.388}{see}%
\TextAt{8}{171.143}{374.388}{that}%
\TextAt{13}{196.466}{374.388}{G}%
\TextAt{9}{205.699}{374.388}{\Character{40}}%
\TextAt{13}{210.251}{374.388}{K}%
\TextAt{17}{220.211}{372.595}{2}%
\TextAt{16}{228.264}{374.388}{\Character{106}}%
\TextAt{13}{234.905}{374.388}{K}%
\TextAt{17}{244.865}{372.595}{1}%
\TextAt{9}{249.597}{374.388}{\Character{41}}%
\TextAt{8}{258.001}{374.388}{is}%
\TextAt{8}{269.718}{374.388}{the}%
\TextAt{8}{289.829}{374.388}{commutator}%
\TextAt{8}{355.811}{374.388}{subgroup}%
\TextAt{8}{406.563}{374.388}{of}%
\TextAt{13}{419.847}{374.388}{G}%
\TextAt{9}{429.08}{374.388}{\Character{40}}%
\TextAt{13}{433.633}{374.388}{K}%
\TextAt{17}{443.592}{372.595}{2}%
\TextAt{16}{451.645}{374.388}{\Character{106}}%
\TextAt{13}{458.287}{374.388}{K}%
\TextAt{9}{469.104}{374.388}{\Character{41}}%
\TextAt{8}{473.657}{374.388}{.}%
\TextAt{8}{482.09}{374.388}{The}%
\TextAt{8}{506.097}{374.388}{proof}%
\TextAt{8}{79.226}{357.342}{of}%
\TextAt{8}{94.111}{357.342}{the}%
\TextAt{8}{115.823}{357.342}{principal}%
\TextAt{8}{166.155}{357.342}{ideal}%
\TextAt{8}{195.674}{357.342}{theorem}%
\TextAt{8}{242.754}{357.342}{is}%
\TextAt{8}{256.073}{357.342}{thus}%
\TextAt{8}{283.38}{357.342}{reduced}%
\TextAt{8}{328.499}{357.342}{to}%
\TextAt{8}{344.365}{357.342}{the}%
\TextAt{8}{366.077}{357.342}{following}%
\TextAt{8}{417.066}{357.342}{purely}%
\TextAt{8}{454.714}{357.342}{group-theoretic}%
\TextAt{8}{79.226}{340.297}{result.}%
\Vector{q 0 G 0 g 0.399 w 452.334 -433.703 m 452.334 -425.634 l S Q}%
\Vector{q 0 G 0 g 0.399 w 452.534 -425.833 m 459.707 -425.833 l S Q}%
\Vector{q 0 G 0 g 0.399 w 452.534 -433.504 m 459.707 -433.504 l S Q}%
\Vector{q 0 G 0 g 0.399 w 459.906 -433.703 m 459.906 -425.634 l S Q}%
\TextAt{12}{79.226}{305.637}{Theorem}%
\TextAt{12}{136.337}{305.637}{(7.6).}%
\TextAt{8}{171.736}{305.637}{Let}%
\TextAt{13}{192.789}{305.637}{G}%
\TextAt{8}{206.003}{305.637}{be}%
\TextAt{8}{222.011}{305.637}{a}%
\TextAt{8}{231.85}{305.637}{finitely}%
\TextAt{8}{271.279}{305.637}{generated}%
\TextAt{8}{324.694}{305.637}{group,}%
\TextAt{13}{361.226}{305.637}{G}%
\TextAt{3}{370.46}{309.975}{\Character{48}}%
\TextAt{8}{377.236}{305.637}{its}%
\TextAt{8}{393.638}{305.637}{commutator}%
\TextAt{8}{459.75}{305.637}{subgroup,}%
\TextAt{8}{513.904}{305.637}{and}%
\TextAt{13}{79.226}{288.591}{G}%
\TextAt{3}{88.46}{292.929}{\Character{48}\Character{48}}%
\TextAt{8}{97.449}{288.591}{the}%
\TextAt{8}{117.606}{288.591}{commutator}%
\TextAt{8}{183.634}{288.591}{subgroup}%
\TextAt{8}{234.432}{288.591}{of}%
\TextAt{13}{247.762}{288.591}{G}%
\TextAt{3}{256.995}{292.929}{\Character{48}}%
\TextAt{8}{259.79}{288.591}{.}%
\TextAt{8}{268.239}{288.591}{If}%
\TextAt{9}{279.931}{288.591}{\Character{40}}%
\TextAt{13}{284.483}{288.591}{G}%
\TextAt{9}{297.037}{288.591}{\Character{58}}%
\TextAt{13}{303.61}{288.591}{G}%
\TextAt{3}{312.844}{292.929}{\Character{48}}%
\TextAt{9}{315.639}{288.591}{\Character{41}}%
\TextAt{13}{323.512}{288.591}{\Character{60}}%
\NativeMathAt{335.937}{288.591}{11.955}{\infty}%
\TextAt{8}{347.893}{288.591}{,}%
\TextAt{8}{355.042}{288.591}{then}%
\TextAt{8}{381.702}{288.591}{the}%
\TextAt{8}{401.858}{288.591}{transfer}%
\TextAt{8}{249.173}{260.148}{Ver}%
\TextAt{9}{271.024}{260.148}{\Character{58}}%
\TextAt{13}{277.597}{260.148}{G}%
\TextAt{9}{286.831}{260.148}{\Character{47}}%
\TextAt{13}{292.684}{260.148}{G}%
\TextAt{3}{301.917}{265.084}{\Character{48}}%
\NativeArrowAt{308.033}{260.148}{11.955}{19.26}{\longrightarrow}%
\TextAt{13}{330.615}{260.148}{G}%
\TextAt{3}{339.849}{265.084}{\Character{48}}%
\TextAt{9}{342.644}{260.148}{\Character{47}}%
\TextAt{13}{348.497}{260.148}{G}%
\TextAt{3}{357.73}{265.084}{\Character{48}\Character{48}}%
\TextAt{8}{79.226}{231.705}{is}%
\TextAt{8}{90.99}{231.705}{the}%
\TextAt{8}{111.146}{231.705}{trivial}%
\TextAt{8}{145.948}{231.705}{homomorphism.}%
\TextAt{8}{97.159}{205.852}{We}%
\TextAt{8}{118.175}{205.852}{give}%
\TextAt{8}{143.113}{205.852}{a}%
\TextAt{8}{153.74}{205.852}{proof}%
\TextAt{8}{185.181}{205.852}{of}%
\TextAt{8}{199.383}{205.852}{this}%
\TextAt{8}{223.078}{205.852}{theorem}%
\TextAt{8}{269.475}{205.852}{which}%
\TextAt{8}{303.833}{205.852}{is}%
\TextAt{8}{316.469}{205.852}{due}%
\TextAt{8}{339.446}{205.852}{to}%
\TextAt{26}{354.628}{205.852}{Ernst}%
\TextAt{26}{396.095}{205.852}{Witt}%
\TextAt{8}{431.828}{205.852}{[141].}%
\TextAt{8}{466.97}{205.852}{In}%
\TextAt{8}{482.463}{205.852}{the}%
\TextAt{8}{503.491}{205.852}{group}%
\TextAt{8}{79.226}{188.806}{ring}%
\TextAt{18}{103.03}{188.806}{Z}%
\TextAt{9}{111}{188.806}{\Character{91}}%
\TextAt{13}{114.252}{188.806}{G}%
\TextAt{9}{123.485}{188.806}{\Character{93}}%
\TextAt{9}{130.036}{188.806}{\Character{61}}%
\TextAt{16}{142.483}{188.806}{\Character{102}}%
\TextAt{20}{148.461}{197.772}{\Character{80}}%
\TextAt{21}{161.08}{185.319}{\Character{27}}%
\NativeMathAt{166.206}{185.319}{7.97}{\in}%
\TextAt{21}{171.851}{185.319}{G}%
\TextAt{13}{180.94}{188.806}{n}%
\TextAt{21}{187.928}{187.013}{\Character{27}}%
\TextAt{13}{193.551}{188.806}{\Character{27}}%
\TextAt{16}{203.955}{188.806}{\Character{106}}%
\TextAt{13}{210.596}{188.806}{n}%
\TextAt{21}{217.584}{187.013}{\Character{27}}%
\NativeMathAt{226.528}{188.806}{11.955}{\in}%
\TextAt{18}{237.819}{188.806}{Z}%
\TextAt{16}{245.789}{188.806}{\Character{103}}%
\TextAt{8}{251.767}{188.806}{,}%
\TextAt{8}{258.707}{188.806}{we}%
\TextAt{8}{275.672}{188.806}{consider}%
\TextAt{8}{320.995}{188.806}{the}%
\TextAt{12}{340.889}{188.806}{augmentation}%
\TextAt{12}{425.219}{188.806}{ideal}%
\TextAt{13}{456.482}{188.806}{I}%
\TextAt{21}{461.645}{187.013}{G}%
\TextAt{8}{468.742}{188.806}{,}%
\TextAt{8}{475.682}{188.806}{which}%
\TextAt{8}{508.906}{188.806}{is}%
\TextAt{8}{520.408}{188.806}{by}%
\TextAt{8}{79.226}{171.76}{definition}%
\TextAt{8}{131.255}{171.76}{the}%
\TextAt{8}{151.411}{171.76}{kernel}%
\TextAt{8}{185.878}{171.76}{of}%
\TextAt{8}{199.208}{171.76}{the}%
\TextAt{8}{219.365}{171.76}{ring}%
\TextAt{8}{243.43}{171.76}{homomorphism}%
\TextAt{18}{212.64}{143.317}{Z}%
\TextAt{9}{220.61}{143.317}{\Character{91}}%
\TextAt{13}{223.862}{143.317}{G}%
\TextAt{9}{233.095}{143.317}{\Character{93}}%
\NativeArrowAt{239.668}{143.317}{11.955}{19.26}{\longrightarrow}%
\TextAt{18}{262.25}{143.317}{Z}%
\TextAt{13}{270.22}{143.317}{\Character{59}}%
\TextAt{20}{299.374}{154.674}{\Character{88}}%
\TextAt{21}{305.446}{129.568}{\Character{27}}%
\TextAt{13}{318.636}{143.317}{n}%
\TextAt{21}{325.623}{141.524}{\Character{27}}%
\TextAt{13}{331.247}{143.317}{\Character{27}}%
\NativeArrowAt{341.65}{143.317}{11.955}{19.26}{\longmapsto}%
\TextAt{20}{364.232}{154.674}{\Character{88}}%
\TextAt{21}{370.304}{129.568}{\Character{27}}%
\TextAt{13}{383.493}{143.317}{n}%
\TextAt{21}{390.481}{141.524}{\Character{27}}%
\TextAt{13}{396.104}{143.317}{\Character{58}}%
\TextAt{8}{79.226}{107.013}{For}%
\TextAt{8}{100.282}{107.013}{every}%
\TextAt{8}{131.261}{107.013}{subgroup}%
\TextAt{13}{182.146}{107.013}{H}%
\TextAt{8}{196.753}{107.013}{of}%
\TextAt{13}{210.17}{107.013}{G}%
\TextAt{8}{219.403}{107.013}{,}%
\TextAt{8}{226.661}{107.013}{we}%
\TextAt{8}{243.975}{107.013}{have}%
\TextAt{13}{271.057}{107.013}{I}%
\TextAt{21}{276.22}{105.22}{H}%
\NativeMathAt{287.757}{107.013}{11.955}{\subseteq}%
\TextAt{13}{300.524}{107.013}{I}%
\TextAt{21}{305.687}{105.22}{G}%
\TextAt{8}{312.784}{107.013}{,}%
\TextAt{8}{320.042}{107.013}{and}%
\TextAt{16}{342.891}{107.013}{\Character{102}}%
\TextAt{13}{348.869}{107.013}{\Character{28}}%
\TextAt{16}{358.004}{107.013}{\Character{0}}%
\TextAt{9}{370.018}{107.013}{1}%
\TextAt{16}{379.34}{107.013}{\Character{106}}%
\TextAt{13}{386.13}{107.013}{\Character{28}}%
\NativeMathAt{396.018}{107.013}{11.955}{\in}%
\TextAt{13}{407.457}{107.013}{H}%
\TextAt{13}{417.416}{107.013}{\Character{59}}%
\TextAt{13}{422.652}{107.013}{\Character{28}}%
\NativeMathAt{432.561}{107.013}{11.955}{\neq}%
\TextAt{9}{445.126}{107.013}{1}%
\TextAt{16}{450.987}{107.013}{\Character{103}}%
\TextAt{8}{460.949}{107.013}{is}%
\TextAt{8}{472.8}{107.013}{a}%
\TextAt{18}{482.642}{107.013}{Z}%
\TextAt{8}{490.612}{107.013}{-basis}%
\TextAt{8}{523.337}{107.013}{of}%
\TextAt{13}{79.226}{89.967}{I}%
\TextAt{21}{84.389}{88.174}{H}%
\TextAt{8}{92.457}{89.967}{.}%
\TextAt{8}{100.799}{89.967}{We}%
\TextAt{8}{120.621}{89.967}{first}%
\TextAt{8}{144.425}{89.967}{establish}%
\TextAt{8}{192.355}{89.967}{the}%
\TextAt{8}{212.189}{89.967}{following}%
\TextAt{8}{261.302}{89.967}{lemma,}%
\TextAt{8}{302.015}{89.967}{which}%
\TextAt{8}{335.18}{89.967}{also}%
\TextAt{8}{358.338}{89.967}{has}%
\TextAt{8}{378.89}{89.967}{independent}%
\TextAt{8}{444.896}{89.967}{interest}%
\TextAt{8}{486.584}{89.967}{in}%
\TextAt{8}{499.916}{89.967}{that}%
\TextAt{8}{524.963}{89.967}{it}%
\TextAt{8}{79.226}{72.921}{gives}%
\TextAt{8}{107.907}{72.921}{an}%
\TextAt{8}{124.166}{72.921}{additive}%
\TextAt{8}{169.045}{72.921}{interpretation}%
\TextAt{8}{243.526}{72.921}{of}%
\TextAt{8}{256.856}{72.921}{the}%
\TextAt{8}{277.012}{72.921}{transfer.}%
\TextAt{8}{297.211}{43.033}{374}%
\Destination{78.23}{809.815}{nav.page.374}%
\Destination{79.36}{733.255}{nav.print.410}%
\Destination{79.36}{775.155}{nav.number.6.7.4}%
\Destination{79.36}{689.355}{nav.number.6.7.5}%
\Destination{79.36}{319.655}{nav.number.6.7.6}%
\Hotlink{179.098}{564.555}{13.995}{11.047}{nav.chapter.4}%
\Hotlink{431.328}{203.655}{25.078}{11.047}{nav.cite.N141}%
\Hotlink{199.242}{564.555}{25.078}{11.047}{nav.number.4.5.9}%
}

% Page 387 (body page 375)
\Page{%
\TextAt{2}{79.226}{796.943}{VI.}%
\TextAt{2}{98.012}{796.943}{Global}%
\TextAt{2}{133.248}{796.943}{Class}%
\TextAt{2}{161.841}{796.943}{Field}%
\TextAt{2}{189.572}{796.943}{Theory}%
\Vector{q 0 G 0 g 0.299 w 7.226 18.729 m 460.769 18.729 l S Q}%
\TextAt{12}{79.226}{761.124}{Lemma}%
\TextAt{12}{125.731}{761.124}{(7.7).}%
\TextAt{8}{160.037}{761.124}{For}%
\TextAt{8}{179.997}{761.124}{every}%
\TextAt{8}{209.881}{761.124}{subgroup}%
\TextAt{13}{259.669}{761.124}{H}%
\TextAt{8}{273.18}{761.124}{of}%
\TextAt{8}{285.501}{761.124}{finite}%
\TextAt{8}{314.404}{761.124}{index}%
\TextAt{8}{344.932}{761.124}{in}%
\TextAt{13}{357.576}{761.124}{G}%
\TextAt{8}{366.81}{761.124}{,}%
\TextAt{8}{373.151}{761.124}{one}%
\TextAt{8}{393.602}{761.124}{has}%
\TextAt{8}{413.467}{761.124}{a}%
\TextAt{8}{422.213}{761.124}{commutative}%
\TextAt{8}{491.13}{761.124}{diagram}%
\TextAt{13}{217.347}{736.571}{G}%
\TextAt{9}{226.58}{736.571}{\Character{47}}%
\TextAt{13}{232.433}{736.571}{G}%
\TextAt{3}{241.667}{740.909}{\Character{48}}%
\TextAt{13}{332.908}{736.571}{H}%
\TextAt{9}{343.531}{736.571}{\Character{47}}%
\TextAt{13}{349.384}{736.571}{H}%
\TextAt{3}{360.006}{740.909}{\Character{48}}%
\TextAt{13}{215.719}{693.828}{I}%
\TextAt{21}{220.882}{692.035}{G}%
\TextAt{9}{227.978}{693.828}{\Character{47}}%
\TextAt{13}{233.831}{693.828}{I}%
\TextAt{17}{239.934}{698.166}{2}%
\TextAt{21}{238.994}{690.808}{G}%
\TextAt{9}{299.434}{693.828}{\Character{40}}%
\TextAt{13}{303.987}{693.828}{I}%
\TextAt{21}{309.149}{692.035}{H}%
\TextAt{9}{319.875}{693.828}{\Character{43}}%
\TextAt{13}{331.636}{693.828}{I}%
\TextAt{21}{336.799}{692.035}{G}%
\TextAt{13}{343.896}{693.828}{I}%
\TextAt{21}{349.058}{692.035}{H}%
\TextAt{9}{357.127}{693.828}{\Character{41}\Character{47}}%
\TextAt{13}{367.532}{693.828}{I}%
\TextAt{21}{372.695}{692.035}{G}%
\TextAt{13}{379.792}{693.828}{I}%
\TextAt{21}{384.955}{692.035}{H}%
\TextAt{13}{393.023}{693.828}{\Character{59}}%
\NativeDiagram{%
\path[draw=black,line width=0.478bp] (177.815,-34.441) -- (255.08,-34.441);%
\path[draw=black,line width=0.478bp,line cap=round,line join=round] (252.832,-31.571) .. controls (253.286,-33.293) and (254.299,-34.106) .. (255.319,-34.441) .. controls (254.299,-34.775) and (253.286,-35.588) .. (252.832,-37.31);%
}%
\TextAt{7}{281.983}{742.542}{Ver}%
\NativeDiagram{%
\path[draw=black,line width=0.478bp] (158.905,-44.997) -- (158.905,-65.64);%
\path[draw=black,line width=0.478bp,line cap=round,line join=round] (161.774,-63.393) .. controls (160.052,-63.847) and (159.239,-64.859) .. (158.905,-65.879) .. controls (158.57,-64.859) and (157.757,-63.847) .. (156.036,-63.393);%
}%
\TextAt{21}{224.043}{715.675}{\Character{14}}%
\NativeMathAt{233.72}{718.193}{7.97}{\sim}%
\TextAt{17}{233.72}{715.679}{\Character{61}}%
\NativeDiagram{%
\path[draw=black,line width=0.478bp] (275.857,-44.997) -- (275.857,-66.149);%
\path[draw=black,line width=0.478bp,line cap=round,line join=round] (278.726,-63.901) .. controls (277.004,-64.356) and (276.191,-65.368) .. (275.857,-66.388) .. controls (275.522,-65.368) and (274.709,-64.356) .. (272.988,-63.901);%
}%
\TextAt{21}{340.995}{715.42}{\Character{14}}%
\NativeMathAt{350.672}{717.938}{7.97}{\sim}%
\TextAt{17}{350.672}{715.424}{\Character{61}}%
\NativeDiagram{%
\path[draw=black,line width=0.478bp] (179.443,-77.184) -- (221.606,-77.184);%
\path[draw=black,line width=0.478bp,line cap=round,line join=round] (219.358,-74.314) .. controls (219.812,-76.036) and (220.825,-76.849) .. (221.845,-77.184) .. controls (220.825,-77.518) and (219.812,-78.331) .. (219.358,-80.053);%
}%
\TextAt{21}{269.966}{688.555}{S}%
\TextAt{8}{79.226}{670.115}{where}%
\TextAt{8}{113.092}{670.115}{the}%
\TextAt{8}{133.306}{670.115}{homomorphisms}%
\TextAt{13}{220.647}{670.115}{\Character{14}}%
\TextAt{8}{230.235}{670.115}{are}%
\TextAt{8}{249.802}{670.115}{induced}%
\TextAt{8}{293.424}{670.115}{by}%
\TextAt{13}{309.74}{670.115}{\Character{27}}%
\NativeArrowAt{320.24}{670.115}{11.955}{11.955}{\mapsto}%
\TextAt{13}{335.613}{670.115}{\Character{14}}%
\TextAt{13}{341.232}{670.115}{\Character{27}}%
\TextAt{9}{351.747}{670.115}{\Character{61}}%
\TextAt{13}{364.269}{670.115}{\Character{27}}%
\TextAt{16}{374.047}{670.115}{\Character{0}}%
\TextAt{9}{386.041}{670.115}{1}%
\TextAt{8}{391.894}{670.115}{,}%
\TextAt{8}{399.114}{670.115}{and}%
\TextAt{8}{421.934}{670.115}{the}%
\TextAt{8}{442.147}{670.115}{homomorphism}%
\TextAt{13}{524.874}{670.115}{S}%
\TextAt{8}{79.226}{653.069}{is}%
\TextAt{8}{90.99}{653.069}{given}%
\TextAt{8}{121.559}{653.069}{by}%
\TextAt{13}{220.067}{636.023}{S}%
\TextAt{9}{227.962}{636.023}{\Character{40}}%
\TextAt{13}{232.515}{636.023}{x}%
\TextAt{8}{242.488}{636.023}{mod}%
\TextAt{13}{267.925}{636.023}{I}%
\TextAt{17}{274.028}{640.96}{2}%
\TextAt{21}{273.088}{633.068}{G}%
\TextAt{9}{280.185}{636.023}{\Character{41}}%
\TextAt{9}{288.04}{636.023}{\Character{61}}%
\TextAt{13}{300.484}{636.023}{x}%
\TextAt{20}{309.128}{647.381}{\Character{88}}%
\TextAt{21}{309.558}{622.009}{\Character{26}}%
\NativeMathAt{313.895}{622.009}{7.97}{\in}%
\TextAt{21}{319.54}{622.009}{R}%
\TextAt{13}{328.39}{636.023}{\Character{26}}%
\TextAt{8}{337.748}{636.023}{mod}%
\TextAt{13}{363.186}{636.023}{I}%
\TextAt{21}{368.349}{634.23}{G}%
\TextAt{13}{375.446}{636.023}{I}%
\TextAt{21}{380.609}{634.23}{H}%
\TextAt{13}{388.677}{636.023}{\Character{59}}%
\TextAt{8}{79.226}{605.136}{for}%
\TextAt{8}{97.111}{605.136}{a}%
\TextAt{8}{106.866}{605.136}{system}%
\TextAt{8}{145.685}{605.136}{of}%
\TextAt{8}{159.015}{605.136}{representatives}%
\TextAt{8}{238.816}{605.136}{of}%
\TextAt{8}{252.146}{605.136}{the}%
\TextAt{8}{272.302}{605.136}{left}%
\TextAt{8}{292.781}{605.136}{cosets}%
\TextAt{13}{326.722}{605.136}{R}%
\NativeMathAt{339.051}{605.136}{11.955}{\ni}%
\TextAt{9}{350.342}{605.136}{1}%
\TextAt{8}{360.093}{605.136}{of}%
\TextAt{13}{373.423}{605.136}{G}%
\TextAt{9}{382.656}{605.136}{\Character{47}}%
\TextAt{13}{388.509}{605.136}{H}%
\TextAt{8}{399.132}{605.136}{.}%
\TextAt{19}{79.226}{573.237}{Proof.}%
\TextAt{8}{116}{573.237}{We}%
\TextAt{8}{136.145}{573.237}{first}%
\TextAt{8}{160.27}{573.237}{show}%
\TextAt{8}{189.273}{573.237}{that}%
\TextAt{8}{214.642}{573.237}{the}%
\TextAt{8}{234.799}{573.237}{homomorphism}%
\TextAt{13}{234.958}{549.627}{H}%
\TextAt{9}{245.581}{549.627}{\Character{47}}%
\TextAt{13}{251.434}{549.627}{H}%
\TextAt{3}{262.056}{554.563}{\Character{48}}%
\TextAt{21}{271.421}{556.404}{\Character{14}}%
\NativeArrowAt{268.172}{549.627}{11.955}{11.955}{\longrightarrow}%
\TextAt{9}{283.448}{549.627}{\Character{40}}%
\TextAt{13}{288.001}{549.627}{I}%
\TextAt{21}{293.164}{547.833}{H}%
\TextAt{9}{303.889}{549.627}{\Character{43}}%
\TextAt{13}{315.65}{549.627}{I}%
\TextAt{21}{320.813}{547.833}{G}%
\TextAt{13}{327.91}{549.627}{I}%
\TextAt{21}{333.073}{547.833}{H}%
\TextAt{9}{341.141}{549.627}{\Character{41}\Character{47}}%
\TextAt{13}{351.546}{549.627}{I}%
\TextAt{21}{356.709}{547.833}{G}%
\TextAt{13}{363.806}{549.627}{I}%
\TextAt{21}{368.969}{547.833}{H}%
\TextAt{8}{517.802}{549.627}{(*)}%
\TextAt{8}{79.226}{526.016}{induced}%
\TextAt{8}{122.655}{526.016}{by}%
\TextAt{13}{138.778}{526.016}{\Character{28}}%
\NativeArrowAt{148.518}{526.016}{11.955}{11.955}{\mapsto}%
\TextAt{13}{163.794}{526.016}{\Character{14}}%
\TextAt{13}{169.413}{526.016}{\Character{28}}%
\TextAt{9}{179.167}{526.016}{\Character{61}}%
\TextAt{13}{191.592}{526.016}{\Character{28}}%
\TextAt{16}{200.39}{526.016}{\Character{0}}%
\TextAt{9}{212.067}{526.016}{1}%
\TextAt{8}{221.682}{526.016}{has}%
\TextAt{8}{242.419}{526.016}{an}%
\TextAt{8}{258.543}{526.016}{inverse.}%
\TextAt{8}{301.807}{526.016}{The}%
\TextAt{8}{325.725}{526.016}{elements}%
\TextAt{13}{373.445}{526.016}{\Character{26}\Character{14}}%
\TextAt{13}{385.101}{526.016}{\Character{28}}%
\TextAt{8}{391.535}{526.016}{,}%
\TextAt{13}{398.575}{526.016}{\Character{28}}%
\NativeMathAt{408.315}{526.016}{11.955}{\in}%
\TextAt{13}{419.606}{526.016}{H}%
\TextAt{8}{430.229}{526.016}{,}%
\TextAt{13}{437.269}{526.016}{\Character{28}}%
\NativeMathAt{447.009}{526.016}{11.955}{\neq}%
\TextAt{9}{459.418}{526.016}{1}%
\TextAt{8}{465.288}{526.016}{,}%
\TextAt{13}{472.328}{526.016}{\Character{26}}%
\NativeMathAt{481.687}{526.016}{11.955}{\in}%
\TextAt{13}{492.978}{526.016}{R}%
\TextAt{8}{501.986}{526.016}{,}%
\TextAt{8}{509.027}{526.016}{form}%
\TextAt{8}{79.226}{508.97}{a}%
\TextAt{18}{88.982}{508.97}{Z}%
\TextAt{8}{96.952}{508.97}{-basis}%
\TextAt{8}{129.589}{508.97}{of}%
\TextAt{13}{142.919}{508.97}{I}%
\TextAt{21}{148.082}{507.176}{H}%
\TextAt{9}{158.807}{508.97}{\Character{43}}%
\TextAt{13}{170.569}{508.97}{I}%
\TextAt{21}{175.732}{507.176}{G}%
\TextAt{13}{182.828}{508.97}{I}%
\TextAt{21}{187.991}{507.176}{H}%
\TextAt{8}{196.06}{508.97}{.}%
\TextAt{8}{204.508}{508.97}{Indeed,}%
\TextAt{8}{245.789}{508.97}{it}%
\TextAt{8}{257.493}{508.97}{follows}%
\TextAt{8}{295.929}{508.97}{from}%
\TextAt{13}{263.984}{485.359}{\Character{26}\Character{14}}%
\TextAt{13}{275.64}{485.359}{\Character{28}}%
\TextAt{9}{285.394}{485.359}{\Character{61}}%
\TextAt{13}{297.82}{485.359}{\Character{14}}%
\TextAt{13}{303.439}{485.359}{\Character{28}}%
\TextAt{9}{312.528}{485.359}{\Character{43}}%
\TextAt{13}{324.289}{485.359}{\Character{14}}%
\TextAt{13}{329.908}{485.359}{\Character{26}\Character{14}}%
\TextAt{13}{341.576}{485.359}{\Character{28}}%
\TextAt{8}{79.226}{461.748}{that}%
\TextAt{8}{104.595}{461.748}{they}%
\TextAt{8}{130.932}{461.748}{generate}%
\TextAt{13}{177.761}{461.748}{I}%
\TextAt{21}{182.924}{459.955}{H}%
\TextAt{9}{193.649}{461.748}{\Character{43}}%
\TextAt{13}{205.41}{461.748}{I}%
\TextAt{21}{210.573}{459.955}{G}%
\TextAt{13}{217.67}{461.748}{I}%
\TextAt{21}{222.833}{459.955}{H}%
\TextAt{8}{230.901}{461.748}{,}%
\TextAt{8}{238.05}{461.748}{and}%
\TextAt{8}{260.813}{461.748}{if}%
\TextAt{9}{134.803}{430.815}{0}%
\TextAt{9}{143.961}{430.815}{\Character{61}}%
\TextAt{20}{156.402}{442.172}{\Character{88}}%
\TextAt{21}{159.353}{417.066}{\Character{26}\Character{59}\Character{28}}%
\TextAt{13}{175.664}{430.815}{n}%
\TextAt{21}{182.651}{429.021}{\Character{26}\Character{59}\Character{28}}%
\TextAt{13}{194.517}{430.815}{\Character{26}\Character{14}}%
\TextAt{13}{206.173}{430.815}{\Character{28}}%
\TextAt{9}{215.927}{430.815}{\Character{61}}%
\TextAt{20}{228.353}{442.172}{\Character{88}}%
\TextAt{21}{231.303}{417.066}{\Character{26}\Character{59}\Character{28}}%
\TextAt{13}{247.614}{430.815}{n}%
\TextAt{21}{254.601}{429.021}{\Character{26}\Character{59}\Character{28}}%
\TextAt{9}{266.467}{430.815}{\Character{40}}%
\TextAt{13}{271.019}{430.815}{\Character{26}\Character{28}}%
\TextAt{16}{286.132}{430.815}{\Character{0}}%
\TextAt{13}{298.087}{430.815}{\Character{26}}%
\TextAt{9}{304.125}{430.815}{\Character{41}}%
\TextAt{9}{311.98}{430.815}{\Character{61}}%
\TextAt{20}{324.424}{442.172}{\Character{88}}%
\TextAt{21}{327.374}{417.066}{\Character{26}\Character{59}\Character{28}}%
\TextAt{13}{343.685}{430.815}{n}%
\TextAt{21}{350.673}{429.021}{\Character{26}\Character{59}\Character{28}}%
\TextAt{13}{362.538}{430.815}{\Character{26}\Character{28}}%
\TextAt{16}{377.651}{430.815}{\Character{0}}%
\TextAt{20}{389.606}{442.172}{\Character{88}}%
\TextAt{21}{396.072}{417.066}{\Character{26}}%
\TextAt{20}{408.867}{451.258}{\Character{32}}%
\TextAt{20}{418.332}{442.172}{\Character{88}}%
\TextAt{21}{424.627}{417.066}{\Character{28}}%
\TextAt{13}{437.593}{430.815}{n}%
\TextAt{21}{444.581}{429.021}{\Character{26}\Character{59}\Character{28}}%
\TextAt{20}{456.446}{451.258}{\Character{33}}%
\TextAt{13}{467.903}{430.815}{\Character{26}\Character{59}}%
\TextAt{8}{79.226}{398.458}{then}%
\TextAt{8}{106.158}{398.458}{we}%
\TextAt{8}{123.657}{398.458}{conclude}%
\TextAt{8}{172.048}{398.458}{that}%
\TextAt{13}{197.688}{398.458}{n}%
\TextAt{21}{204.676}{396.664}{\Character{26}\Character{59}\Character{28}}%
\TextAt{9}{220.325}{398.458}{\Character{61}}%
\TextAt{9}{233.201}{398.458}{0}%
\TextAt{8}{243.235}{398.458}{because}%
\TextAt{8}{286.808}{398.458}{the}%
\TextAt{13}{307.236}{398.458}{\Character{26}\Character{28}}%
\TextAt{8}{319.693}{398.458}{,}%
\TextAt{13}{327.182}{398.458}{\Character{26}}%
\TextAt{8}{337.388}{398.458}{are}%
\TextAt{8}{357.171}{398.458}{pairwise}%
\TextAt{8}{403.028}{398.458}{distinct.}%
\TextAt{8}{450.726}{398.458}{Mapping}%
\TextAt{13}{500.098}{398.458}{\Character{26}\Character{14}}%
\TextAt{13}{511.754}{398.458}{\Character{28}}%
\TextAt{8}{522.357}{398.458}{to}%
\TextAt{13}{79.226}{381.412}{\Character{28}}%
\TextAt{8}{88.966}{381.412}{mod}%
\TextAt{13}{114.404}{381.412}{H}%
\TextAt{3}{125.026}{385.75}{\Character{48}}%
\TextAt{8}{127.821}{381.412}{,}%
\TextAt{8}{134.971}{381.412}{we}%
\TextAt{8}{152.198}{381.412}{now}%
\TextAt{8}{176.587}{381.412}{have}%
\TextAt{8}{203.581}{381.412}{a}%
\TextAt{8}{213.337}{381.412}{surjective}%
\TextAt{8}{265.748}{381.412}{homomorphism}%
\TextAt{13}{249.904}{357.801}{I}%
\TextAt{21}{255.067}{356.007}{H}%
\TextAt{9}{265.792}{357.801}{\Character{43}}%
\TextAt{13}{277.553}{357.801}{I}%
\TextAt{21}{282.716}{356.007}{G}%
\TextAt{13}{289.813}{357.801}{I}%
\TextAt{21}{294.976}{356.007}{H}%
\NativeArrowAt{306.365}{357.801}{11.955}{19.26}{\longrightarrow}%
\TextAt{13}{328.947}{357.801}{H}%
\TextAt{9}{339.57}{357.801}{\Character{47}}%
\TextAt{13}{345.423}{357.801}{H}%
\TextAt{3}{356.045}{362.737}{\Character{48}}%
\TextAt{13}{358.84}{357.801}{\Character{58}}%
\TextAt{8}{79.226}{334.19}{It}%
\TextAt{8}{90.841}{334.19}{sends}%
\TextAt{13}{121.118}{334.19}{\Character{14}}%
\TextAt{9}{126.751}{334.19}{\Character{40}}%
\TextAt{13}{131.303}{334.19}{\Character{26}\Character{28}}%
\TextAt{3}{143.76}{338.528}{\Character{48}}%
\TextAt{9}{146.555}{334.19}{\Character{41}}%
\TextAt{13}{151.107}{334.19}{\Character{14}}%
\TextAt{13}{156.726}{334.19}{\Character{28}}%
\NativeMathAt{166.48}{334.19}{11.955}{\in}%
\TextAt{13}{177.771}{334.19}{I}%
\TextAt{21}{182.934}{332.397}{G}%
\TextAt{13}{190.031}{334.19}{I}%
\TextAt{21}{195.193}{332.397}{H}%
\TextAt{8}{206.102}{334.19}{to}%
\TextAt{13}{219.354}{334.19}{\Character{28}}%
\TextAt{3}{225.773}{338.528}{\Character{48}}%
\TextAt{13}{228.568}{334.19}{\Character{28}}%
\TextAt{13}{234.976}{334.19}{\Character{28}}%
\TextAt{3}{241.406}{338.528}{\Character{48}\Character{0}}%
\TextAt{17}{250.29}{338.528}{1}%
\TextAt{13}{255.022}{334.19}{\Character{28}}%
\TextAt{3}{261.441}{338.528}{\Character{0}}%
\TextAt{17}{268.027}{338.528}{1}%
\NativeMathAt{276.081}{334.19}{11.955}{\equiv}%
\TextAt{9}{288.7}{334.19}{1}%
\TextAt{8}{297.874}{334.19}{mod}%
\TextAt{13}{323.312}{334.19}{H}%
\TextAt{3}{333.934}{338.528}{\Character{48}}%
\TextAt{8}{339.569}{334.19}{because}%
\TextAt{13}{381.813}{334.19}{\Character{14}}%
\TextAt{9}{387.446}{334.19}{\Character{40}}%
\TextAt{13}{391.998}{334.19}{\Character{26}\Character{28}}%
\TextAt{3}{404.455}{338.528}{\Character{48}}%
\TextAt{9}{407.25}{334.19}{\Character{41}}%
\TextAt{13}{411.802}{334.19}{\Character{14}}%
\TextAt{13}{417.421}{334.19}{\Character{28}}%
\TextAt{9}{427.175}{334.19}{\Character{61}}%
\TextAt{13}{439.601}{334.19}{\Character{26}\Character{14}}%
\TextAt{9}{451.271}{334.19}{\Character{40}}%
\TextAt{13}{455.824}{334.19}{\Character{28}}%
\TextAt{3}{462.243}{338.528}{\Character{48}}%
\TextAt{13}{465.038}{334.19}{\Character{28}}%
\TextAt{9}{471.457}{334.19}{\Character{41}}%
\TextAt{16}{476.503}{334.19}{\Character{0}}%
\TextAt{13}{486.295}{334.19}{\Character{26}\Character{14}}%
\TextAt{13}{497.951}{334.19}{\Character{28}}%
\TextAt{3}{504.385}{338.528}{\Character{48}}%
\TextAt{16}{507.674}{334.19}{\Character{0}}%
\TextAt{13}{517.466}{334.19}{\Character{14}}%
\TextAt{13}{523.085}{334.19}{\Character{28}}%
\TextAt{8}{529.518}{334.19}{.}%
\TextAt{8}{79.226}{317.144}{It}%
\TextAt{8}{90.764}{317.144}{thus}%
\TextAt{8}{115.38}{317.144}{induces}%
\TextAt{8}{155.921}{317.144}{a}%
\TextAt{8}{164.542}{317.144}{homomorphism}%
\TextAt{8}{246.077}{317.144}{which}%
\TextAt{8}{278.428}{317.144}{is}%
\TextAt{8}{289.057}{317.144}{inverse}%
\TextAt{8}{326.681}{317.144}{to}%
\TextAt{8}{339.857}{317.144}{(*).}%
\TextAt{8}{362.894}{317.144}{In}%
\TextAt{8}{376.381}{317.144}{particular,}%
\TextAt{8}{432.714}{317.144}{if}%
\TextAt{13}{442.303}{317.144}{H}%
\TextAt{9}{456.247}{317.144}{\Character{61}}%
\TextAt{13}{468.672}{317.144}{G}%
\TextAt{8}{477.906}{317.144}{,}%
\TextAt{8}{484.147}{317.144}{we}%
\TextAt{8}{500.239}{317.144}{obtain}%
\TextAt{8}{79.226}{300.098}{the}%
\TextAt{8}{99.383}{300.098}{isomorphism}%
\TextAt{13}{167.802}{300.098}{G}%
\TextAt{9}{177.036}{300.098}{\Character{47}}%
\TextAt{13}{182.889}{300.098}{G}%
\TextAt{3}{192.122}{304.436}{\Character{48}}%
\TextAt{21}{201.487}{306.875}{\Character{14}}%
\NativeArrowAt{198.238}{300.098}{11.955}{11.955}{\longrightarrow}%
\TextAt{13}{213.514}{300.098}{I}%
\TextAt{21}{218.677}{298.305}{G}%
\TextAt{9}{225.774}{300.098}{\Character{47}}%
\TextAt{13}{231.627}{300.098}{I}%
\TextAt{17}{237.73}{304.436}{2}%
\TextAt{21}{236.79}{297.077}{G}%
\TextAt{8}{243.887}{300.098}{.}%
\TextAt{8}{97.159}{280.063}{The}%
\TextAt{8}{121.213}{280.063}{transfer}%
\TextAt{8}{164.526}{280.063}{is}%
\TextAt{8}{176.29}{280.063}{now}%
\TextAt{8}{200.679}{280.063}{obtained}%
\TextAt{8}{248.81}{280.063}{as}%
\TextAt{8}{222.713}{256.452}{Ver}%
\TextAt{9}{241.244}{256.452}{\Character{40}}%
\TextAt{13}{245.796}{256.452}{\Character{27}}%
\TextAt{8}{256.199}{256.452}{mod}%
\TextAt{13}{281.637}{256.452}{G}%
\TextAt{3}{290.871}{261.389}{\Character{48}}%
\TextAt{9}{293.666}{256.452}{\Character{41}}%
\TextAt{9}{301.521}{256.452}{\Character{61}}%
\TextAt{20}{314.531}{267.81}{\Character{89}}%
\TextAt{21}{313.964}{242.438}{\Character{26}}%
\NativeMathAt{318.301}{242.438}{7.97}{\in}%
\TextAt{21}{323.946}{242.438}{R}%
\TextAt{13}{332.366}{256.452}{\Character{27}}%
\TextAt{21}{339.02}{254.659}{\Character{26}}%
\TextAt{8}{347.175}{256.452}{mod}%
\TextAt{13}{372.613}{256.452}{H}%
\TextAt{3}{383.236}{261.389}{\Character{48}}%
\TextAt{13}{386.031}{256.452}{\Character{59}}%
\TextAt{8}{79.226}{223.366}{where}%
\TextAt{13}{113.035}{223.366}{\Character{27}}%
\TextAt{21}{119.689}{221.572}{\Character{26}}%
\NativeMathAt{127.844}{223.366}{11.955}{\in}%
\TextAt{13}{139.135}{223.366}{H}%
\TextAt{8}{153.655}{223.366}{is}%
\TextAt{8}{165.419}{223.366}{defined}%
\TextAt{8}{205.732}{223.366}{by}%
\TextAt{13}{221.991}{223.366}{\Character{27}}%
\TextAt{13}{229.056}{223.366}{\Character{26}}%
\TextAt{9}{238.432}{223.366}{\Character{61}}%
\TextAt{13}{250.858}{223.366}{\Character{26}}%
\TextAt{3}{256.895}{227.704}{\Character{48}}%
\TextAt{13}{259.69}{223.366}{\Character{27}}%
\TextAt{21}{266.344}{221.572}{\Character{26}}%
\TextAt{8}{271.178}{223.366}{,}%
\TextAt{13}{278.328}{223.366}{\Character{26}}%
\TextAt{3}{284.365}{227.704}{\Character{48}}%
\NativeMathAt{290.481}{223.366}{11.955}{\in}%
\TextAt{13}{301.772}{223.366}{R}%
\TextAt{8}{310.781}{223.366}{.}%
\TextAt{8}{319.229}{223.366}{Ver}%
\TextAt{8}{341.657}{223.366}{thus}%
\TextAt{8}{367.408}{223.366}{induces}%
\TextAt{8}{409.084}{223.366}{the}%
\TextAt{8}{429.241}{223.366}{homomorphism}%
\TextAt{13}{222.171}{199.755}{S}%
\TextAt{9}{233.388}{199.755}{\Character{58}}%
\TextAt{13}{239.96}{199.755}{I}%
\TextAt{21}{245.123}{197.962}{G}%
\TextAt{9}{252.22}{199.755}{\Character{47}}%
\TextAt{13}{258.073}{199.755}{I}%
\TextAt{17}{264.176}{204.691}{2}%
\TextAt{21}{263.236}{196.799}{G}%
\NativeArrowAt{273.653}{199.755}{11.955}{19.26}{\longrightarrow}%
\TextAt{9}{296.235}{199.755}{\Character{40}}%
\TextAt{13}{300.788}{199.755}{I}%
\TextAt{21}{305.951}{197.962}{H}%
\TextAt{9}{316.676}{199.755}{\Character{43}}%
\TextAt{13}{328.437}{199.755}{I}%
\TextAt{21}{333.6}{197.962}{G}%
\TextAt{13}{340.697}{199.755}{I}%
\TextAt{21}{345.86}{197.962}{H}%
\TextAt{9}{353.928}{199.755}{\Character{41}\Character{47}}%
\TextAt{13}{364.333}{199.755}{I}%
\TextAt{21}{369.496}{197.962}{G}%
\TextAt{13}{376.593}{199.755}{I}%
\TextAt{21}{381.756}{197.962}{H}%
\TextAt{8}{79.226}{176.144}{given}%
\TextAt{8}{109.796}{176.144}{by}%
\TextAt{13}{126.055}{176.144}{S}%
\TextAt{9}{133.95}{176.144}{\Character{40}}%
\TextAt{13}{138.502}{176.144}{\Character{14}}%
\TextAt{13}{144.121}{176.144}{\Character{27}}%
\TextAt{8}{154.538}{176.144}{mod}%
\TextAt{13}{179.976}{176.144}{I}%
\TextAt{17}{186.079}{180.482}{2}%
\TextAt{21}{185.139}{173.124}{G}%
\TextAt{9}{192.236}{176.144}{\Character{41}}%
\TextAt{9}{200.091}{176.144}{\Character{61}}%
\TextAt{20}{212.535}{185.11}{\Character{80}}%
\TextAt{21}{225.154}{172.657}{\Character{26}}%
\NativeMathAt{229.491}{172.657}{7.97}{\in}%
\TextAt{21}{235.136}{172.657}{R}%
\TextAt{13}{244.054}{176.144}{\Character{14}}%
\TextAt{13}{249.673}{176.144}{\Character{27}}%
\TextAt{21}{256.341}{174.351}{\Character{26}}%
\TextAt{8}{264.496}{176.144}{mod}%
\TextAt{13}{289.934}{176.144}{I}%
\TextAt{21}{295.097}{174.351}{G}%
\TextAt{13}{302.194}{176.144}{I}%
\TextAt{21}{307.357}{174.351}{H}%
\TextAt{8}{315.425}{176.144}{.}%
\TextAt{8}{323.873}{176.144}{From}%
\TextAt{13}{354.598}{176.144}{\Character{27}}%
\TextAt{13}{361.663}{176.144}{\Character{26}}%
\TextAt{9}{371.039}{176.144}{\Character{61}}%
\TextAt{13}{383.465}{176.144}{\Character{26}}%
\TextAt{3}{389.502}{180.482}{\Character{48}}%
\TextAt{13}{392.297}{176.144}{\Character{27}}%
\TextAt{21}{398.951}{174.351}{\Character{26}}%
\TextAt{8}{407.683}{176.144}{follows}%
\TextAt{8}{446.119}{176.144}{the}%
\TextAt{8}{466.275}{176.144}{identity}%
\TextAt{13}{223.521}{152.533}{\Character{14}}%
\TextAt{13}{229.14}{152.533}{\Character{26}}%
\TextAt{9}{237.848}{152.533}{\Character{43}}%
\TextAt{9}{249.6}{152.533}{\Character{40}}%
\TextAt{13}{254.162}{152.533}{\Character{14}}%
\TextAt{13}{259.781}{152.533}{\Character{27}}%
\TextAt{9}{266.877}{152.533}{\Character{41}}%
\TextAt{13}{271.429}{152.533}{\Character{26}}%
\TextAt{9}{280.788}{152.533}{\Character{61}}%
\TextAt{13}{293.214}{152.533}{\Character{14}}%
\TextAt{13}{298.833}{152.533}{\Character{27}}%
\TextAt{21}{305.5}{150.74}{\Character{26}}%
\TextAt{9}{312.991}{152.533}{\Character{43}}%
\TextAt{13}{324.753}{152.533}{\Character{14}}%
\TextAt{13}{330.372}{152.533}{\Character{26}}%
\TextAt{3}{336.423}{157.469}{\Character{48}}%
\TextAt{9}{341.875}{152.533}{\Character{43}}%
\TextAt{13}{353.636}{152.533}{\Character{14}}%
\TextAt{13}{359.255}{152.533}{\Character{26}}%
\TextAt{3}{365.307}{157.469}{\Character{48}}%
\TextAt{13}{368.102}{152.533}{\Character{14}}%
\TextAt{13}{373.721}{152.533}{\Character{27}}%
\TextAt{21}{380.389}{150.74}{\Character{26}}%
\TextAt{13}{385.223}{152.533}{\Character{58}}%
\TextAt{8}{79.226}{128.922}{Since}%
\TextAt{13}{109.784}{128.922}{\Character{26}}%
\TextAt{3}{115.821}{133.261}{\Character{48}}%
\TextAt{8}{122.514}{128.922}{runs}%
\TextAt{8}{148.588}{128.922}{through}%
\TextAt{8}{192.822}{128.922}{the}%
\TextAt{8}{212.979}{128.922}{set}%
\TextAt{13}{231.246}{128.922}{R}%
\TextAt{8}{244.152}{128.922}{if}%
\TextAt{13}{254.876}{128.922}{\Character{26}}%
\TextAt{8}{264.811}{128.922}{does,}%
\TextAt{8}{294.46}{128.922}{we}%
\TextAt{8}{311.687}{128.922}{get}%
\TextAt{8}{331.198}{128.922}{as}%
\TextAt{8}{345.568}{128.922}{claimed}%
\TextAt{13}{157.735}{105.312}{S}%
\TextAt{9}{165.63}{105.312}{\Character{40}}%
\TextAt{13}{170.182}{105.312}{\Character{14}}%
\TextAt{13}{175.801}{105.312}{\Character{26}}%
\TextAt{8}{185.174}{105.312}{mod}%
\TextAt{13}{210.612}{105.312}{I}%
\TextAt{17}{216.715}{110.248}{2}%
\TextAt{21}{215.775}{102.356}{G}%
\TextAt{9}{222.871}{105.312}{\Character{41}}%
\NativeMathAt{230.745}{105.312}{11.955}{\equiv}%
\TextAt{20}{243.364}{116.669}{\Character{88}}%
\TextAt{21}{243.794}{91.297}{\Character{26}}%
\NativeMathAt{248.13}{91.298}{7.97}{\in}%
\TextAt{21}{253.776}{91.297}{R}%
\TextAt{13}{262.625}{105.312}{\Character{14}}%
\TextAt{13}{268.244}{105.312}{\Character{27}}%
\TextAt{21}{274.912}{103.518}{\Character{26}}%
\NativeMathAt{283.067}{105.312}{11.955}{\equiv}%
\TextAt{20}{295.686}{116.669}{\Character{88}}%
\TextAt{21}{296.116}{91.298}{\Character{26}}%
\NativeMathAt{300.453}{91.298}{7.97}{\in}%
\TextAt{21}{306.098}{91.298}{R}%
\TextAt{9}{312.955}{105.312}{\Character{40}}%
\TextAt{13}{317.507}{105.312}{\Character{14}}%
\TextAt{13}{323.126}{105.312}{\Character{27}}%
\TextAt{9}{330.223}{105.312}{\Character{41}}%
\TextAt{13}{334.775}{105.312}{\Character{26}}%
\NativeMathAt{344.134}{105.312}{11.955}{\equiv}%
\TextAt{13}{356.753}{105.312}{\Character{14}}%
\TextAt{13}{362.372}{105.312}{\Character{27}}%
\TextAt{20}{371.461}{116.669}{\Character{88}}%
\TextAt{21}{371.89}{91.298}{\Character{26}}%
\NativeMathAt{376.227}{91.298}{7.97}{\in}%
\TextAt{21}{381.872}{91.298}{R}%
\TextAt{13}{390.722}{105.312}{\Character{26}}%
\TextAt{8}{400.08}{105.312}{mod}%
\TextAt{13}{425.518}{105.312}{I}%
\TextAt{21}{430.681}{103.518}{G}%
\TextAt{13}{437.778}{105.312}{I}%
\TextAt{21}{442.941}{103.518}{H}%
\TextAt{13}{451.009}{105.312}{\Character{58}}%
\Vector{q 0 G 0 g 0.399 w 452.334 -701.079 m 452.334 -693.009 l S Q}%
\Vector{q 0 G 0 g 0.399 w 452.534 -693.208 m 459.707 -693.208 l S Q}%
\Vector{q 0 G 0 g 0.399 w 452.534 -700.879 m 459.707 -700.879 l S Q}%
\Vector{q 0 G 0 g 0.399 w 459.906 -701.079 m 459.906 -693.009 l S Q}%
\TextAt{8}{297.211}{43.033}{375}%
\Destination{522.36}{566.675}{nav.AMS.79}%
\Destination{78.23}{809.815}{nav.page.375}%
\Destination{79.36}{776.155}{nav.print.411}%
\Destination{79.36}{775.155}{nav.number.6.7.7}%
}

% Page 388 (body page 376)
\Page{%
\TextAt{2}{79.226}{796.943}{VI.}%
\TextAt{2}{98.012}{796.943}{Global}%
\TextAt{2}{133.248}{796.943}{Class}%
\TextAt{2}{161.841}{796.943}{Field}%
\TextAt{2}{189.572}{796.943}{Theory}%
\Vector{q 0 G 0 g 0.299 w 7.226 18.729 m 460.769 18.729 l S Q}%
\TextAt{12}{79.226}{761.124}{Proof}%
\TextAt{12}{116.347}{761.124}{of}%
\TextAt{12}{131.673}{761.124}{theorem}%
\TextAt{12}{184.575}{761.124}{(7.6):}%
\TextAt{19}{79.226}{732.422}{Proof.}%
\TextAt{8}{116}{732.422}{Replacing}%
\TextAt{13}{171.198}{732.422}{G}%
\TextAt{8}{185.382}{732.422}{by}%
\TextAt{13}{202.694}{732.422}{G}%
\TextAt{9}{211.928}{732.422}{\Character{47}}%
\TextAt{13}{217.781}{732.422}{G}%
\TextAt{3}{227.015}{736.76}{\Character{48}\Character{48}}%
\TextAt{8}{232.107}{732.422}{,}%
\TextAt{8}{240.572}{732.422}{we}%
\TextAt{8}{258.852}{732.422}{may}%
\TextAt{8}{285.274}{732.422}{assume}%
\TextAt{8}{326.772}{732.422}{that}%
\TextAt{13}{353.194}{732.422}{G}%
\TextAt{3}{362.427}{736.76}{\Character{48}\Character{48}}%
\TextAt{9}{372.635}{732.422}{\Character{61}}%
\TextAt{16}{386.855}{732.422}{\Character{102}}%
\TextAt{9}{392.832}{732.422}{1}%
\TextAt{16}{398.685}{732.422}{\Character{103}}%
\TextAt{8}{404.663}{732.422}{,}%
\TextAt{8}{413.128}{732.422}{i.e.,}%
\TextAt{8}{436.55}{732.422}{that}%
\TextAt{13}{462.971}{732.422}{G}%
\TextAt{3}{472.205}{736.76}{\Character{48}}%
\TextAt{8}{479.951}{732.422}{is}%
\TextAt{8}{492.767}{732.422}{abelian.}%
\TextAt{8}{79.226}{715.376}{Let}%
\TextAt{13}{101.307}{715.376}{R}%
\NativeMathAt{115.531}{715.376}{11.955}{\ni}%
\TextAt{9}{128.716}{715.376}{1}%
\TextAt{8}{139.578}{715.376}{be}%
\TextAt{8}{156.614}{715.376}{a}%
\TextAt{8}{167.482}{715.376}{system}%
\TextAt{8}{207.412}{715.376}{of}%
\TextAt{8}{221.853}{715.376}{representatives}%
\TextAt{8}{302.766}{715.376}{of}%
\TextAt{8}{317.207}{715.376}{left}%
\TextAt{8}{338.798}{715.376}{cosets}%
\TextAt{8}{373.851}{715.376}{of}%
\TextAt{13}{388.292}{715.376}{G}%
\TextAt{9}{397.526}{715.376}{\Character{47}}%
\TextAt{13}{403.379}{715.376}{G}%
\TextAt{3}{412.613}{719.714}{\Character{48}}%
\TextAt{8}{415.408}{715.376}{,}%
\TextAt{8}{423.947}{715.376}{and}%
\TextAt{8}{447.821}{715.376}{let}%
\TextAt{13}{465.837}{715.376}{\Character{27}}%
\TextAt{17}{472.491}{713.583}{1}%
\TextAt{13}{477.223}{715.376}{\Character{59}}%
\TextAt{13}{482.447}{715.376}{\Character{58}}%
\TextAt{13}{487.684}{715.376}{\Character{58}}%
\TextAt{13}{492.908}{715.376}{\Character{58}}%
\TextAt{13}{498.145}{715.376}{\Character{59}}%
\TextAt{13}{503.381}{715.376}{\Character{27}}%
\TextAt{21}{510.097}{713.583}{n}%
\TextAt{8}{520.743}{715.376}{be}%
\TextAt{8}{79.226}{698.33}{generators}%
\TextAt{8}{135.765}{698.33}{of}%
\TextAt{13}{148.979}{698.33}{G}%
\TextAt{8}{158.212}{698.33}{.}%
\TextAt{8}{166.622}{698.33}{Mapping}%
\TextAt{13}{215.605}{698.33}{e}%
\TextAt{21}{221.031}{696.537}{i}%
\TextAt{9}{227.733}{698.33}{\Character{61}}%
\TextAt{9}{240.143}{698.33}{\Character{40}0}%
\TextAt{13}{250.564}{698.33}{\Character{59}}%
\TextAt{13}{255.788}{698.33}{\Character{58}}%
\TextAt{13}{261.025}{698.33}{\Character{58}}%
\TextAt{13}{266.249}{698.33}{\Character{58}}%
\TextAt{13}{271.486}{698.33}{\Character{59}}%
\TextAt{9}{276.784}{698.33}{0}%
\TextAt{13}{282.637}{698.33}{\Character{59}}%
\TextAt{9}{287.881}{698.33}{1}%
\TextAt{13}{293.734}{698.33}{\Character{59}}%
\TextAt{9}{298.979}{698.33}{0}%
\TextAt{13}{304.832}{698.33}{\Character{59}}%
\TextAt{13}{310.056}{698.33}{\Character{58}}%
\TextAt{13}{315.293}{698.33}{\Character{58}}%
\TextAt{13}{320.517}{698.33}{\Character{58}}%
\TextAt{13}{325.754}{698.33}{\Character{59}}%
\TextAt{9}{331.052}{698.33}{0\Character{41}}%
\NativeMathAt{344.779}{698.33}{11.955}{\in}%
\TextAt{18}{356.069}{698.33}{Z}%
\TextAt{21}{364.04}{702.668}{n}%
\TextAt{8}{373.457}{698.33}{to}%
\TextAt{13}{387.651}{698.33}{\Character{27}}%
\TextAt{21}{394.304}{696.537}{i}%
\TextAt{8}{397.685}{698.33}{,}%
\TextAt{8}{404.741}{698.33}{we}%
\TextAt{8}{421.852}{698.33}{get}%
\TextAt{8}{441.247}{698.33}{an}%
\TextAt{8}{457.389}{698.33}{exact}%
\TextAt{8}{488.165}{698.33}{sequence}%
\TextAt{9}{223.202}{670.377}{0}%
\NativeArrowAt{232.376}{670.377}{11.955}{19.26}{\longrightarrow}%
\TextAt{18}{254.958}{670.377}{Z}%
\TextAt{21}{262.928}{675.313}{n}%
\TextAt{21}{274.684}{677.641}{f}%
\NativeArrowAt{271.885}{670.377}{11.955}{11.955}{\longrightarrow}%
\TextAt{18}{287.161}{670.377}{Z}%
\TextAt{21}{295.132}{675.313}{n}%
\NativeArrowAt{304.089}{670.377}{11.955}{19.26}{\longrightarrow}%
\TextAt{13}{326.671}{670.377}{G}%
\TextAt{9}{335.904}{670.377}{\Character{47}}%
\TextAt{13}{341.757}{670.377}{G}%
\TextAt{3}{350.991}{675.313}{\Character{48}}%
\NativeArrowAt{357.107}{670.377}{11.955}{19.26}{\longrightarrow}%
\TextAt{9}{379.689}{670.377}{1}%
\TextAt{13}{385.542}{670.377}{\Character{59}}%
\TextAt{8}{79.226}{642.423}{where}%
\TextAt{13}{113.035}{642.423}{f}%
\TextAt{8}{123.979}{642.423}{is}%
\TextAt{8}{135.743}{642.423}{given}%
\TextAt{8}{166.312}{642.423}{by}%
\TextAt{8}{182.571}{642.423}{an}%
\TextAt{13}{198.831}{642.423}{n}%
\TextAt{16}{208.475}{642.423}{\Character{2}}%
\TextAt{13}{220.43}{642.423}{n}%
\TextAt{8}{227.418}{642.423}{-matrix}%
\TextAt{9}{269.368}{642.423}{\Character{40}}%
\TextAt{13}{273.921}{642.423}{m}%
\TextAt{21}{284.16}{640.63}{ik}%
\TextAt{9}{292.163}{642.423}{\Character{41}}%
\TextAt{8}{300.612}{642.423}{with}%
\TextAt{8}{327.272}{642.423}{det}%
\TextAt{9}{343.532}{642.423}{\Character{40}}%
\TextAt{13}{348.084}{642.423}{m}%
\TextAt{21}{358.323}{640.63}{ik}%
\TextAt{9}{366.326}{642.423}{\Character{41}}%
\TextAt{9}{374.181}{642.423}{\Character{61}}%
\TextAt{9}{386.602}{642.423}{\Character{40}}%
\TextAt{13}{391.177}{642.423}{G}%
\TextAt{9}{403.731}{642.423}{\Character{58}}%
\TextAt{13}{410.304}{642.423}{G}%
\TextAt{3}{419.538}{646.762}{\Character{48}}%
\TextAt{9}{422.333}{642.423}{\Character{41}}%
\TextAt{8}{426.885}{642.423}{.}%
\TextAt{8}{435.333}{642.423}{Consequently,}%
\TextAt{21}{230.044}{622.904}{n}%
\TextAt{20}{224.975}{619.317}{\Character{89}}%
\TextAt{21}{225.761}{594.123}{i}%
\TextAt{17}{228.645}{594.123}{\Character{61}1}%
\TextAt{13}{242.244}{607.96}{\Character{27}}%
\TextAt{21}{249.326}{613.484}{m}%
\TextAt{22}{256.817}{612.078}{ik}%
\TextAt{21}{248.897}{604.896}{i}%
\TextAt{13}{264.552}{607.96}{\Character{28}}%
\TextAt{21}{269.635}{606.167}{k}%
\TextAt{9}{278.076}{607.96}{\Character{61}}%
\TextAt{9}{290.486}{607.96}{1}%
\TextAt{8}{320.265}{607.96}{with}%
\TextAt{13}{346.925}{607.96}{\Character{28}}%
\TextAt{21}{352.008}{606.167}{k}%
\NativeMathAt{360.449}{607.96}{11.955}{\in}%
\TextAt{13}{371.74}{607.96}{G}%
\TextAt{3}{380.974}{612.896}{\Character{48}}%
\TextAt{13}{383.769}{607.96}{\Character{58}}%
\TextAt{8}{79.226}{571.549}{The}%
\TextAt{8}{103.28}{571.549}{formulae}%
\TextAt{13}{151.411}{571.549}{\Character{14}}%
\TextAt{9}{157.044}{571.549}{\Character{40}}%
\TextAt{13}{161.597}{571.549}{xy}%
\TextAt{9}{174.385}{571.549}{\Character{41}}%
\TextAt{9}{182.24}{571.549}{\Character{61}}%
\TextAt{13}{194.684}{571.549}{\Character{14}}%
\TextAt{13}{200.303}{571.549}{x}%
\TextAt{9}{209.626}{571.549}{\Character{43}}%
\TextAt{13}{221.387}{571.549}{\Character{14}}%
\TextAt{13}{227.006}{571.549}{y}%
\TextAt{9}{235.813}{571.549}{\Character{43}}%
\TextAt{13}{247.575}{571.549}{\Character{14}}%
\TextAt{13}{253.194}{571.549}{x\Character{14}}%
\TextAt{13}{265.46}{571.549}{y}%
\TextAt{8}{271.629}{571.549}{,}%
\TextAt{13}{278.779}{571.549}{\Character{14}}%
\TextAt{9}{284.412}{571.549}{\Character{40}}%
\TextAt{13}{288.964}{571.549}{x}%
\TextAt{3}{295.616}{575.887}{\Character{0}}%
\TextAt{17}{302.202}{575.887}{1}%
\TextAt{9}{306.935}{571.549}{\Character{41}}%
\TextAt{9}{314.79}{571.549}{\Character{61}}%
\TextAt{16}{327.233}{571.549}{\Character{0}}%
\TextAt{9}{336.532}{571.549}{\Character{40}}%
\TextAt{13}{341.084}{571.549}{\Character{14}}%
\TextAt{13}{346.703}{571.549}{x}%
\TextAt{9}{353.369}{571.549}{\Character{41}}%
\TextAt{13}{357.922}{571.549}{x}%
\TextAt{3}{364.574}{575.887}{\Character{0}}%
\TextAt{17}{371.16}{575.887}{1}%
\TextAt{8}{379.79}{571.549}{yield}%
\TextAt{8}{408.076}{571.549}{by}%
\TextAt{8}{424.335}{571.549}{iteration}%
\TextAt{8}{471.821}{571.549}{that}%
\TextAt{13}{218.989}{535.735}{\Character{14}}%
\TextAt{20}{226.614}{556.178}{\Character{32}}%
\TextAt{21}{241.148}{550.679}{n}%
\TextAt{20}{236.079}{547.092}{\Character{89}}%
\TextAt{21}{236.865}{521.898}{i}%
\TextAt{17}{239.748}{521.898}{\Character{61}1}%
\TextAt{13}{253.347}{535.735}{\Character{27}}%
\TextAt{21}{260.43}{541.259}{m}%
\TextAt{22}{267.92}{539.853}{ik}%
\TextAt{21}{260.001}{532.671}{i}%
\TextAt{13}{275.655}{535.735}{\Character{28}}%
\TextAt{21}{280.739}{533.942}{k}%
\TextAt{20}{285.858}{556.178}{\Character{33}}%
\TextAt{9}{298.644}{535.735}{\Character{61}}%
\TextAt{21}{317.134}{550.679}{n}%
\TextAt{20}{311.069}{547.092}{\Character{88}}%
\TextAt{21}{312.852}{521.898}{i}%
\TextAt{17}{315.735}{521.898}{\Character{61}1}%
\TextAt{9}{328.338}{535.735}{\Character{40}}%
\TextAt{13}{332.89}{535.735}{\Character{14}}%
\TextAt{13}{338.509}{535.735}{\Character{27}}%
\TextAt{21}{345.177}{533.942}{i}%
\TextAt{9}{348.558}{535.735}{\Character{41}}%
\TextAt{13}{353.11}{535.735}{\Character{22}}%
\TextAt{21}{360.153}{533.942}{ik}%
\TextAt{9}{371.477}{535.735}{\Character{61}}%
\TextAt{9}{383.887}{535.735}{0}%
\TextAt{13}{389.755}{535.735}{\Character{59}}%
\TextAt{8}{79.226}{498.714}{where}%
\TextAt{13}{113.137}{498.714}{\Character{22}}%
\TextAt{21}{120.18}{496.921}{ik}%
\NativeMathAt{131.676}{498.714}{11.955}{\equiv}%
\TextAt{13}{144.468}{498.714}{m}%
\TextAt{21}{154.707}{496.921}{ik}%
\TextAt{8}{166.031}{498.714}{mod}%
\TextAt{13}{191.469}{498.714}{I}%
\TextAt{21}{196.632}{496.921}{G}%
\TextAt{8}{203.729}{498.714}{.}%
\TextAt{8}{212.481}{498.714}{In}%
\TextAt{8}{227.203}{498.714}{fact,}%
\TextAt{8}{253.667}{498.714}{the}%
\TextAt{13}{273.925}{498.714}{\Character{28}}%
\TextAt{21}{279.008}{496.921}{k}%
\TextAt{8}{288.127}{498.714}{are}%
\TextAt{8}{307.739}{498.714}{products}%
\TextAt{8}{356.354}{498.714}{of}%
\TextAt{8}{369.786}{498.714}{commutators}%
\TextAt{8}{440.53}{498.714}{of}%
\TextAt{8}{453.961}{498.714}{the}%
\TextAt{13}{474.219}{498.714}{\Character{27}}%
\TextAt{21}{480.873}{496.921}{i}%
\TextAt{8}{488.253}{498.714}{and}%
\TextAt{13}{511.116}{498.714}{\Character{27}}%
\TextAt{3}{518.199}{503.663}{\Character{0}}%
\TextAt{17}{524.785}{503.663}{1}%
\TextAt{21}{517.77}{495.651}{i}%
\TextAt{8}{529.518}{498.714}{.}%
\TextAt{8}{79.226}{481.668}{We}%
\TextAt{8}{99.371}{481.668}{view}%
\TextAt{9}{126.353}{481.668}{\Character{40}}%
\TextAt{13}{130.906}{481.668}{\Character{22}}%
\TextAt{21}{137.949}{479.875}{ik}%
\TextAt{9}{145.952}{481.668}{\Character{41}}%
\TextAt{8}{154.401}{481.668}{as}%
\TextAt{8}{168.771}{481.668}{a}%
\TextAt{8}{178.527}{481.668}{matrix}%
\TextAt{8}{216.58}{481.668}{over}%
\TextAt{8}{241.626}{481.668}{the}%
\TextAt{8}{261.783}{481.668}{commutative}%
\TextAt{8}{331.708}{481.668}{ring}%
\TextAt{18}{241.502}{453.715}{Z}%
\TextAt{9}{249.472}{453.715}{\Character{91}}%
\TextAt{13}{252.724}{453.715}{G}%
\TextAt{9}{261.957}{453.715}{\Character{47}}%
\TextAt{13}{267.81}{453.715}{G}%
\TextAt{3}{277.044}{458.651}{\Character{48}}%
\TextAt{9}{279.839}{453.715}{\Character{93}}%
\NativeMathAt{286.412}{453.715}{11.955}{\cong}%
\TextAt{18}{299.031}{453.715}{Z}%
\TextAt{9}{307.001}{453.715}{\Character{91}}%
\TextAt{13}{310.253}{453.715}{G}%
\TextAt{9}{319.486}{453.715}{\Character{93}\Character{47}}%
\TextAt{18}{328.591}{453.715}{Z}%
\TextAt{9}{336.561}{453.715}{\Character{91}}%
\TextAt{13}{339.813}{453.715}{G}%
\TextAt{9}{349.046}{453.715}{\Character{93}}%
\TextAt{13}{352.298}{453.715}{I}%
\TextAt{21}{357.461}{451.922}{G}%
\TextAt{23}{364.06}{454.191}{\Character{48}}%
\TextAt{13}{367.242}{453.715}{\Character{59}}%
\TextAt{8}{79.226}{425.762}{which}%
\TextAt{8}{113.948}{425.762}{gives}%
\TextAt{8}{143.863}{425.762}{a}%
\TextAt{8}{154.854}{425.762}{meaning}%
\TextAt{8}{202.917}{425.762}{to}%
\TextAt{8}{218.463}{425.762}{the}%
\TextAt{8}{239.854}{425.762}{determinant}%
\TextAt{13}{307.105}{425.762}{\Character{22}}%
\TextAt{9}{319.574}{425.762}{\Character{61}}%
\TextAt{8}{334.104}{425.762}{det}%
\TextAt{9}{350.363}{425.762}{\Character{40}}%
\TextAt{13}{354.916}{425.762}{\Character{22}}%
\TextAt{21}{361.959}{423.969}{ik}%
\TextAt{9}{369.961}{425.762}{\Character{41}}%
\NativeMathAt{379.939}{425.762}{11.955}{\in}%
\TextAt{18}{393.335}{425.762}{Z}%
\TextAt{9}{401.305}{425.762}{\Character{91}}%
\TextAt{13}{404.557}{425.762}{G}%
\TextAt{9}{413.79}{425.762}{\Character{47}}%
\TextAt{13}{419.643}{425.762}{G}%
\TextAt{3}{428.877}{430.1}{\Character{48}}%
\TextAt{9}{431.672}{425.762}{\Character{93}}%
\TextAt{8}{434.924}{425.762}{.}%
\TextAt{8}{447.077}{425.762}{Let}%
\TextAt{9}{469.282}{425.762}{\Character{40}}%
\TextAt{13}{473.834}{425.762}{\Character{21}}%
\TextAt{21}{480.663}{423.969}{k}%
\TextAt{21}{485.278}{423.969}{j}%
\TextAt{9}{489.666}{425.762}{\Character{41}}%
\TextAt{8}{499.351}{425.762}{be}%
\TextAt{8}{516.51}{425.762}{the}%
\TextAt{8}{79.226}{408.716}{adjoint}%
\TextAt{8}{118.905}{408.716}{matrix}%
\TextAt{8}{156.959}{408.716}{of}%
\TextAt{9}{170.289}{408.716}{\Character{40}}%
\TextAt{13}{174.841}{408.716}{\Character{22}}%
\TextAt{21}{181.884}{406.923}{ik}%
\TextAt{9}{189.887}{408.716}{\Character{41}}%
\TextAt{8}{194.439}{408.716}{.}%
\TextAt{8}{202.888}{408.716}{Then}%
\TextAt{9}{192.9}{380.763}{\Character{40}}%
\TextAt{13}{197.452}{380.763}{\Character{14}}%
\TextAt{13}{203.071}{380.763}{\Character{27}}%
\TextAt{21}{209.739}{378.969}{j}%
\TextAt{9}{214.121}{380.763}{\Character{41}}%
\TextAt{13}{218.673}{380.763}{\Character{22}}%
\TextAt{9}{229.037}{380.763}{\Character{61}}%
\TextAt{20}{241.462}{392.12}{\Character{88}}%
\TextAt{21}{245.168}{366.66}{i\Character{59}k}%
\TextAt{9}{258.731}{380.763}{\Character{40}}%
\TextAt{13}{263.283}{380.763}{\Character{14}}%
\TextAt{13}{268.902}{380.763}{\Character{27}}%
\TextAt{21}{275.57}{378.969}{i}%
\TextAt{9}{278.951}{380.763}{\Character{41}}%
\TextAt{13}{283.503}{380.763}{\Character{22}}%
\TextAt{21}{290.546}{378.969}{ik}%
\TextAt{13}{298.549}{380.763}{\Character{21}}%
\TextAt{21}{305.378}{378.969}{k}%
\TextAt{21}{309.993}{378.969}{j}%
\NativeMathAt{317.702}{380.763}{11.955}{\equiv}%
\TextAt{9}{330.322}{380.763}{0}%
\TextAt{8}{339.495}{380.763}{mod}%
\TextAt{13}{364.933}{380.763}{I}%
\TextAt{21}{370.096}{378.969}{G}%
\TextAt{18}{377.193}{380.763}{Z}%
\TextAt{9}{385.163}{380.763}{\Character{91}}%
\TextAt{13}{388.415}{380.763}{G}%
\TextAt{9}{397.648}{380.763}{\Character{93}}%
\TextAt{13}{400.9}{380.763}{I}%
\TextAt{21}{406.063}{378.969}{G}%
\TextAt{23}{412.662}{381.239}{\Character{48}}%
\TextAt{13}{415.844}{380.763}{\Character{59}}%
\TextAt{8}{79.226}{343.045}{so}%
\TextAt{8}{93.596}{343.045}{that}%
\TextAt{9}{118.965}{343.045}{\Character{40}}%
\TextAt{13}{123.517}{343.045}{\Character{14}}%
\TextAt{13}{129.136}{343.045}{\Character{27}}%
\TextAt{9}{136.233}{343.045}{\Character{41}}%
\TextAt{13}{140.785}{343.045}{\Character{22}}%
\NativeMathAt{151.149}{343.045}{11.955}{\equiv}%
\TextAt{9}{163.768}{343.045}{0}%
\TextAt{8}{172.942}{343.045}{mod}%
\TextAt{13}{198.38}{343.045}{I}%
\TextAt{21}{203.543}{341.251}{G}%
\TextAt{18}{210.64}{343.045}{Z}%
\TextAt{9}{218.61}{343.045}{\Character{91}}%
\TextAt{13}{221.861}{343.045}{G}%
\TextAt{9}{231.095}{343.045}{\Character{93}}%
\TextAt{13}{234.347}{343.045}{I}%
\TextAt{21}{239.51}{341.251}{G}%
\TextAt{23}{246.108}{343.521}{\Character{48}}%
\TextAt{9}{252.612}{343.045}{\Character{61}}%
\TextAt{13}{265.037}{343.044}{I}%
\TextAt{21}{270.2}{341.251}{G}%
\TextAt{13}{277.297}{343.045}{I}%
\TextAt{21}{282.46}{341.251}{G}%
\TextAt{23}{289.058}{343.521}{\Character{48}}%
\TextAt{8}{296.138}{343.045}{for}%
\TextAt{8}{314.023}{343.045}{all}%
\TextAt{13}{330.282}{343.045}{\Character{27}}%
\TextAt{8}{337.365}{343.045}{.}%
\TextAt{8}{345.813}{343.045}{This}%
\TextAt{8}{372.533}{343.045}{yields}%
\TextAt{13}{246.526}{315.091}{\Character{22}}%
\NativeMathAt{256.89}{315.091}{11.955}{\equiv}%
\TextAt{20}{269.509}{326.449}{\Character{88}}%
\TextAt{21}{269.939}{301.077}{\Character{26}}%
\NativeMathAt{274.275}{301.077}{7.97}{\in}%
\TextAt{21}{279.921}{301.077}{R}%
\TextAt{13}{288.77}{315.091}{\Character{26}}%
\TextAt{8}{298.129}{315.091}{mod}%
\TextAt{18}{323.567}{315.091}{Z}%
\TextAt{9}{331.537}{315.091}{\Character{91}}%
\TextAt{13}{334.789}{315.091}{G}%
\TextAt{9}{344.022}{315.091}{\Character{93}}%
\TextAt{13}{347.274}{315.091}{I}%
\TextAt{21}{352.437}{313.298}{G}%
\TextAt{23}{359.035}{315.567}{\Character{48}}%
\TextAt{13}{362.218}{315.091}{\Character{58}}%
\TextAt{8}{79.226}{277.462}{For}%
\TextAt{8}{100.196}{277.462}{if}%
\TextAt{8}{110.919}{277.462}{we}%
\TextAt{8}{128.147}{277.462}{put}%
\TextAt{13}{149.606}{277.462}{\Character{22}}%
\TextAt{9}{159.97}{277.462}{\Character{61}}%
\TextAt{20}{172.396}{286.428}{\Character{80}}%
\TextAt{21}{185.015}{273.975}{\Character{26}}%
\NativeMathAt{189.351}{273.975}{7.97}{\in}%
\TextAt{21}{194.997}{273.975}{R}%
\TextAt{13}{203.915}{277.462}{n}%
\TextAt{21}{210.903}{275.668}{\Character{26}}%
\Vector{q 0 G 0 g 0.478 w 143.737 -489.717 m 149.775 -489.717 l S Q}%
\TextAt{13}{215.737}{277.462}{\Character{26}}%
\TextAt{8}{221.775}{277.462}{,}%
\TextAt{8}{228.924}{277.462}{where}%
\Vector{q 0 G 0 g 0.478 w 190.733 -489.717 m 196.771 -489.717 l S Q}%
\TextAt{13}{262.733}{277.462}{\Character{26}}%
\TextAt{9}{272.092}{277.462}{\Character{61}}%
\TextAt{13}{284.518}{277.462}{\Character{26}}%
\TextAt{8}{293.876}{277.462}{mod}%
\TextAt{13}{319.314}{277.462}{G}%
\TextAt{3}{328.548}{281.8}{\Character{48}}%
\TextAt{8}{331.343}{277.462}{,}%
\TextAt{8}{338.492}{277.462}{then}%
\TextAt{8}{365.152}{277.462}{for}%
\TextAt{8}{383.037}{277.462}{all}%
\Vector{q 0 G 0 g 0.478 w 327.296 -489.717 m 334.378 -489.717 l S Q}%
\TextAt{13}{399.296}{277.462}{\Character{27}}%
\NativeMathAt{409.699}{277.462}{11.955}{\in}%
\TextAt{13}{420.99}{277.462}{G}%
\TextAt{9}{430.224}{277.462}{\Character{47}}%
\TextAt{13}{436.077}{277.462}{G}%
\TextAt{3}{445.31}{281.8}{\Character{48}}%
\TextAt{8}{448.105}{277.462}{,}%
\Vector{q 0 G 0 g 0.478 w 167.904 -519.211 m 174.987 -519.211 l S Q}%
\TextAt{13}{239.904}{247.968}{\Character{27}}%
\TextAt{13}{246.97}{247.968}{\Character{22}}%
\TextAt{9}{257.351}{247.968}{\Character{61}}%
\TextAt{20}{269.776}{259.326}{\Character{88}}%
\Vector{q 0 G 0 g 0.359 w 204.242 -535.138 m 208.579 -535.138 l S Q}%
\TextAt{21}{276.242}{234.176}{\Character{26}}%
\TextAt{13}{289.037}{247.968}{n}%
\TextAt{21}{296.025}{246.175}{\Character{26}}%
\Vector{q 0 G 0 g 0.478 w 228.859 -519.211 m 235.942 -519.211 l S Q}%
\TextAt{13}{300.859}{247.968}{\Character{27}}%
\Vector{q 0 G 0 g 0.478 w 237.934 -519.211 m 243.972 -519.211 l S Q}%
\TextAt{13}{309.934}{247.968}{\Character{26}}%
\TextAt{9}{319.293}{247.968}{\Character{61}}%
\TextAt{20}{331.718}{259.326}{\Character{88}}%
\Vector{q 0 G 0 g 0.359 w 266.184 -535.138 m 270.521 -535.138 l S Q}%
\TextAt{21}{338.184}{234.176}{\Character{26}}%
\TextAt{13}{350.98}{247.968}{n}%
\TextAt{21}{357.967}{246.175}{\Character{26}}%
\Vector{q 0 G 0 g 0.478 w 290.802 -519.211 m 296.84 -519.211 l S Q}%
\TextAt{13}{362.802}{247.968}{\Character{26}\Character{58}}%
\TextAt{8}{79.226}{210.56}{This}%
\TextAt{8}{105.946}{210.56}{implies}%
\TextAt{8}{145.673}{210.56}{that}%
\TextAt{8}{171.042}{210.56}{all}%
\TextAt{13}{187.301}{210.56}{n}%
\TextAt{21}{194.288}{208.767}{\Character{26}}%
\TextAt{8}{203.021}{210.56}{are}%
\TextAt{8}{222.531}{210.56}{equal,}%
\TextAt{8}{256.675}{210.56}{hence}%
\TextAt{13}{289.181}{210.56}{\Character{22}}%
\NativeMathAt{299.545}{210.56}{11.955}{\equiv}%
\TextAt{13}{312.165}{210.56}{m}%
\TextAt{20}{324.396}{219.527}{\Character{80}}%
\TextAt{21}{337.016}{207.073}{\Character{26}}%
\NativeMathAt{341.352}{207.073}{7.97}{\in}%
\TextAt{21}{346.998}{207.073}{R}%
\TextAt{13}{355.916}{210.56}{\Character{26}}%
\TextAt{8}{365.274}{210.56}{mod}%
\TextAt{18}{390.712}{210.56}{Z}%
\TextAt{9}{398.682}{210.56}{\Character{91}}%
\TextAt{13}{401.934}{210.56}{G}%
\TextAt{9}{411.168}{210.56}{\Character{93}}%
\TextAt{13}{414.419}{210.56}{I}%
\TextAt{21}{419.582}{208.767}{G}%
\TextAt{23}{426.181}{211.036}{\Character{48}}%
\TextAt{8}{429.364}{210.56}{,}%
\TextAt{8}{436.513}{210.56}{and}%
\TextAt{8}{459.275}{210.56}{as}%
\TextAt{13}{189.248}{182.607}{\Character{22}}%
\NativeMathAt{199.612}{182.607}{11.955}{\equiv}%
\TextAt{8}{212.231}{182.607}{det}%
\TextAt{9}{228.49}{182.607}{\Character{40}}%
\TextAt{13}{233.043}{182.607}{m}%
\TextAt{21}{243.282}{180.814}{ik}%
\TextAt{9}{251.285}{182.607}{\Character{41}}%
\NativeMathAt{259.158}{182.607}{11.955}{\equiv}%
\TextAt{9}{271.777}{182.607}{\Character{40}}%
\TextAt{13}{276.33}{182.607}{G}%
\TextAt{9}{288.884}{182.607}{\Character{58}}%
\TextAt{13}{295.457}{182.607}{G}%
\TextAt{3}{304.69}{187.543}{\Character{48}}%
\TextAt{9}{307.485}{182.607}{\Character{41}}%
\NativeMathAt{315.358}{182.607}{11.955}{\equiv}%
\TextAt{13}{327.978}{182.607}{m}%
\TextAt{9}{338.217}{182.607}{\Character{40}}%
\TextAt{13}{342.769}{182.607}{G}%
\TextAt{9}{355.324}{182.607}{\Character{58}}%
\TextAt{13}{361.896}{182.607}{G}%
\TextAt{3}{371.13}{187.543}{\Character{48}}%
\TextAt{9}{373.925}{182.607}{\Character{41}}%
\TextAt{8}{381.798}{182.607}{mod}%
\TextAt{13}{407.236}{182.607}{I}%
\TextAt{21}{412.399}{180.814}{G}%
\TextAt{13}{419.496}{182.607}{\Character{59}}%
\TextAt{8}{79.226}{154.654}{we}%
\TextAt{8}{96.962}{154.654}{even}%
\TextAt{8}{124.131}{154.654}{have}%
\TextAt{13}{151.635}{154.654}{m}%
\TextAt{9}{166.062}{154.653}{\Character{61}}%
\TextAt{9}{179.344}{154.653}{1}%
\TextAt{8}{185.207}{154.654}{.}%
\TextAt{8}{195.182}{154.654}{Applying}%
\TextAt{8}{246.417}{154.654}{now}%
\TextAt{8}{271.314}{154.654}{lemma}%
\TextAt{8}{309.541}{154.654}{(7.7),}%
\TextAt{8}{341.404}{154.654}{we}%
\TextAt{8}{359.14}{154.654}{see}%
\TextAt{8}{378.562}{154.654}{that}%
\TextAt{8}{404.44}{154.654}{the}%
\TextAt{8}{425.105}{154.654}{transfer}%
\TextAt{8}{468.927}{154.654}{is}%
\TextAt{8}{481.2}{154.654}{the}%
\TextAt{8}{501.865}{154.654}{trivial}%
\TextAt{8}{79.226}{137.608}{homomorphism}%
\TextAt{8}{161.896}{137.608}{since}%
\TextAt{13}{179.849}{109.654}{S}%
\TextAt{9}{187.745}{109.654}{\Character{40}}%
\TextAt{13}{192.297}{109.654}{\Character{14}}%
\TextAt{13}{197.916}{109.654}{\Character{27}}%
\TextAt{8}{208.333}{109.654}{mod}%
\TextAt{13}{233.771}{109.654}{I}%
\TextAt{17}{239.874}{114.59}{2}%
\TextAt{21}{238.934}{106.699}{G}%
\TextAt{9}{246.031}{109.654}{\Character{41}}%
\TextAt{9}{253.886}{109.654}{\Character{61}}%
\TextAt{13}{266.329}{109.654}{\Character{14}}%
\TextAt{13}{271.948}{109.654}{\Character{27}}%
\TextAt{20}{281.037}{121.012}{\Character{88}}%
\TextAt{21}{281.467}{95.64}{\Character{26}}%
\NativeMathAt{285.803}{95.64}{7.97}{\in}%
\TextAt{21}{291.449}{95.64}{R}%
\TextAt{13}{300.298}{109.654}{\Character{26}}%
\NativeMathAt{309.657}{109.654}{11.955}{\equiv}%
\TextAt{9}{322.276}{109.654}{\Character{40}}%
\TextAt{13}{326.829}{109.654}{\Character{14}}%
\TextAt{13}{332.448}{109.654}{\Character{27}}%
\TextAt{9}{339.544}{109.654}{\Character{41}}%
\TextAt{13}{344.096}{109.654}{\Character{22}}%
\NativeMathAt{354.46}{109.654}{11.955}{\equiv}%
\TextAt{9}{367.079}{109.654}{0}%
\TextAt{8}{376.253}{109.654}{mod}%
\TextAt{13}{401.691}{109.654}{I}%
\TextAt{21}{406.854}{107.861}{G}%
\TextAt{13}{413.951}{109.654}{I}%
\TextAt{21}{419.114}{107.861}{G}%
\TextAt{23}{425.712}{110.13}{\Character{48}}%
\TextAt{13}{428.895}{109.654}{\Character{58}}%
\Vector{q 0 G 0 g 0.399 w 452.334 -701.079 m 452.334 -693.009 l S Q}%
\Vector{q 0 G 0 g 0.399 w 452.534 -693.208 m 459.707 -693.208 l S Q}%
\Vector{q 0 G 0 g 0.399 w 452.534 -700.879 m 459.707 -700.879 l S Q}%
\Vector{q 0 G 0 g 0.399 w 459.906 -701.079 m 459.906 -693.009 l S Q}%
\TextAt{8}{297.211}{43.033}{376}%
\Destination{78.23}{809.815}{nav.page.376}%
\Destination{79.36}{747.455}{nav.print.412}%
\Hotlink{184.075}{758.955}{28.676}{11.047}{nav.number.6.7.6}%
\Hotlink{309.041}{152.455}{25.078}{11.047}{nav.number.6.7.7}%
}

% Page 389 (body page 377)
\Page{%
\TextAt{2}{79.226}{796.943}{VI.}%
\TextAt{2}{98.012}{796.943}{Global}%
\TextAt{2}{133.248}{796.943}{Class}%
\TextAt{2}{161.841}{796.943}{Field}%
\TextAt{2}{189.572}{796.943}{Theory}%
\Vector{q 0 G 0 g 0.299 w 7.226 18.729 m 460.769 18.729 l S Q}%
\TextAt{8}{97.159}{761.124}{A}%
\TextAt{8}{109.83}{761.124}{problem}%
\TextAt{8}{155.353}{761.124}{which}%
\TextAt{8}{188.838}{761.124}{is}%
\TextAt{8}{200.6}{761.124}{closely}%
\TextAt{8}{238.054}{761.124}{related}%
\TextAt{8}{277.074}{761.124}{to}%
\TextAt{8}{291.383}{761.124}{the}%
\TextAt{8}{311.537}{761.124}{principal}%
\TextAt{8}{360.313}{761.124}{ideal}%
\TextAt{8}{388.274}{761.124}{theorem}%
\TextAt{8}{433.798}{761.124}{and}%
\TextAt{8}{456.559}{761.124}{which}%
\TextAt{8}{490.043}{761.124}{was}%
\TextAt{8}{512.541}{761.124}{first}%
\TextAt{8}{79.226}{744.078}{put}%
\TextAt{8}{100.171}{744.078}{forward}%
\TextAt{8}{142.588}{744.078}{by}%
\TextAt{26}{158.333}{744.078}{Philipp}%
\TextAt{26}{205.151}{744.078}{Furtwängler}%
\TextAt{8}{289.1}{744.078}{is}%
\TextAt{8}{300.35}{744.078}{the}%
\TextAt{12}{319.992}{744.078}{problem}%
\TextAt{12}{372.147}{744.078}{of}%
\TextAt{12}{386.882}{744.078}{the}%
\TextAt{12}{409.615}{744.078}{class}%
\TextAt{12}{440.382}{744.078}{field}%
\TextAt{12}{469.092}{744.078}{tower}%
\TextAt{8}{501.67}{744.078}{.}%
\TextAt{8}{509.947}{744.078}{This}%
\TextAt{8}{79.226}{727.032}{is}%
\TextAt{8}{90.99}{727.032}{the}%
\TextAt{8}{111.146}{727.032}{question}%
\TextAt{8}{157.712}{727.032}{whether}%
\TextAt{8}{202.58}{727.032}{the}%
\TextAt{8}{222.736}{727.032}{class}%
\TextAt{8}{250.173}{727.032}{field}%
\TextAt{8}{275.53}{727.032}{tower}%
\TextAt{13}{220.86}{699.192}{K}%
\TextAt{9}{234.999}{699.192}{\Character{61}}%
\TextAt{13}{247.424}{699.192}{K}%
\TextAt{17}{257.384}{697.399}{0}%
\NativeMathAt{265.437}{699.192}{11.955}{\subseteq}%
\TextAt{13}{278.056}{699.192}{K}%
\TextAt{17}{288.015}{697.399}{1}%
\NativeMathAt{296.068}{699.192}{11.955}{\subseteq}%
\TextAt{13}{308.688}{699.192}{K}%
\TextAt{17}{318.647}{697.399}{2}%
\NativeMathAt{326.7}{699.192}{11.955}{\subseteq}%
\TextAt{13}{339.319}{699.192}{K}%
\TextAt{17}{349.278}{697.399}{3}%
\NativeMathAt{357.331}{699.192}{11.955}{\subseteq}%
\TextAt{16}{369.932}{699.192}{\Character{1}}%
\TextAt{16}{375.24}{699.192}{\Character{1}}%
\TextAt{16}{380.548}{699.192}{\Character{1}}%
\TextAt{13}{387.883}{699.192}{\Character{59}}%
\TextAt{8}{79.226}{671.352}{where}%
\TextAt{13}{112.634}{671.352}{K}%
\TextAt{21}{122.593}{669.559}{i}%
\TextAt{17}{125.476}{669.559}{\Character{43}1}%
\TextAt{8}{140.291}{671.352}{is}%
\TextAt{8}{151.654}{671.352}{the}%
\TextAt{8}{171.409}{671.352}{Hilbert}%
\TextAt{8}{211.321}{671.352}{class}%
\TextAt{8}{238.357}{671.352}{field}%
\TextAt{8}{263.312}{671.352}{of}%
\TextAt{13}{276.241}{671.352}{K}%
\TextAt{21}{286.2}{669.559}{i}%
\TextAt{8}{289.581}{671.352}{,}%
\TextAt{8}{296.41}{671.352}{stops}%
\TextAt{8}{326.052}{671.352}{after}%
\TextAt{8}{353.291}{671.352}{a}%
\TextAt{8}{362.645}{671.352}{finite}%
\TextAt{8}{392.156}{671.352}{number}%
\TextAt{8}{434.351}{671.352}{of}%
\TextAt{8}{447.279}{671.352}{steps.}%
\TextAt{8}{481.082}{671.352}{A}%
\TextAt{8}{493.353}{671.352}{positive}%
\TextAt{8}{79.226}{654.306}{answer}%
\TextAt{8}{118.982}{654.306}{would}%
\TextAt{8}{154.124}{654.306}{have}%
\TextAt{8}{182.116}{654.306}{the}%
\TextAt{8}{203.27}{654.306}{implication}%
\TextAt{8}{265.406}{654.306}{that}%
\TextAt{8}{291.772}{654.306}{the}%
\TextAt{8}{312.926}{654.306}{last}%
\TextAt{8}{336.1}{654.306}{field}%
\TextAt{8}{362.454}{654.306}{in}%
\TextAt{8}{377.105}{654.306}{the}%
\TextAt{8}{398.258}{654.306}{tower}%
\TextAt{8}{431.128}{654.306}{had}%
\TextAt{8}{454.888}{654.306}{class}%
\TextAt{8}{483.323}{654.306}{number}%
\TextAt{9}{526.916}{654.306}{1}%
\TextAt{8}{79.226}{637.26}{so}%
\TextAt{8}{94.409}{637.26}{that}%
\TextAt{8}{120.592}{637.26}{in}%
\TextAt{8}{135.058}{637.26}{it}%
\TextAt{8}{147.575}{637.26}{not}%
\TextAt{8}{169.202}{637.26}{only}%
\TextAt{8}{195.707}{637.26}{the}%
\TextAt{8}{216.677}{637.26}{ideals}%
\TextAt{8}{250.068}{637.26}{of}%
\TextAt{13}{264.211}{637.26}{K}%
\TextAt{8}{275.029}{637.26}{,}%
\TextAt{8}{283.194}{637.26}{but}%
\TextAt{8}{305.467}{637.26}{in}%
\TextAt{8}{319.933}{637.26}{fact}%
\TextAt{8}{343.832}{637.26}{all}%
\TextAt{8}{360.904}{637.26}{its}%
\TextAt{8}{378.036}{637.26}{ideals}%
\TextAt{8}{411.427}{637.26}{become}%
\TextAt{8}{454.179}{637.26}{principal.}%
\TextAt{8}{509.947}{637.26}{This}%
\TextAt{8}{79.226}{620.214}{perspective}%
\TextAt{8}{139.468}{620.214}{naturally}%
\TextAt{8}{188.938}{620.214}{generated}%
\TextAt{8}{241.326}{620.214}{the}%
\TextAt{8}{260.538}{620.214}{greatest}%
\TextAt{8}{303.887}{620.214}{interest.}%
\TextAt{8}{350.133}{620.214}{But}%
\TextAt{8}{372.43}{620.214}{the}%
\TextAt{8}{391.641}{620.214}{problem,}%
\TextAt{8}{439.663}{620.214}{after}%
\TextAt{8}{466.359}{620.214}{withstanding}%
\TextAt{8}{79.226}{603.168}{for}%
\TextAt{8}{96.749}{603.168}{a}%
\TextAt{8}{106.142}{603.168}{long}%
\TextAt{8}{131.149}{603.168}{time}%
\TextAt{8}{157.447}{603.168}{all}%
\TextAt{8}{173.343}{603.168}{attempts}%
\TextAt{8}{222.476}{603.168}{to}%
\TextAt{8}{236.424}{603.168}{solve}%
\TextAt{8}{264.742}{603.168}{it,}%
\TextAt{8}{279.408}{603.168}{was}%
\TextAt{8}{301.546}{603.168}{finally}%
\TextAt{8}{336.63}{603.168}{decided}%
\TextAt{8}{378.53}{603.168}{in}%
\TextAt{8}{391.82}{603.168}{the}%
\TextAt{8}{411.614}{603.168}{negative}%
\TextAt{8}{457.435}{603.168}{by}%
\TextAt{8}{473.332}{603.168}{the}%
\TextAt{8}{493.126}{603.168}{Russian}%
\TextAt{8}{79.226}{586.122}{mathematicians}%
\TextAt{26}{163.845}{586.122}{E.S.}%
\TextAt{26}{192.167}{586.122}{Golod}%
\TextAt{8}{235.038}{586.122}{and}%
\TextAt{26}{257.8}{586.122}{I.R.}%
\TextAt{26}{284.365}{586.122}{Šafarevič}%
\TextAt{8}{347.44}{586.122}{in}%
\TextAt{8}{361.093}{586.122}{1964}%
\TextAt{8}{388.423}{586.122}{(see}%
\TextAt{8}{411.891}{586.122}{[48],}%
\TextAt{8}{437.26}{586.122}{[24]).}%
\TextAt{12}{79.226}{560.442}{Exercise}%
\TextAt{12}{132.125}{560.442}{VI.7.1.}%
\TextAt{8}{175.619}{560.442}{The}%
\TextAt{8}{199.41}{560.442}{decomposition}%
\TextAt{8}{276.283}{560.442}{law}%
\TextAt{8}{297.157}{560.442}{for}%
\TextAt{8}{314.78}{560.442}{the}%
\TextAt{8}{334.674}{560.442}{prime}%
\TextAt{8}{367.575}{560.442}{ideals}%
\TextAt{14}{399.891}{560.442}{p}%
\TextAt{8}{409.509}{560.442}{which}%
\TextAt{8}{442.733}{560.442}{are}%
\TextAt{19}{461.981}{560.442}{ramified}%
\TextAt{8}{507.017}{560.442}{in}%
\TextAt{8}{520.408}{560.442}{an}%
\TextAt{8}{79.226}{543.396}{abelian}%
\TextAt{8}{120.675}{543.396}{extension}%
\TextAt{13}{173.241}{543.396}{L}%
\TextAt{16}{185.891}{543.396}{\Character{106}}%
\TextAt{13}{193.898}{543.396}{K}%
\TextAt{8}{209.413}{543.396}{can}%
\TextAt{8}{231.674}{543.396}{be}%
\TextAt{8}{248.399}{543.396}{formulated}%
\TextAt{8}{308.39}{543.396}{like}%
\TextAt{8}{330.65}{543.396}{this.}%
\TextAt{8}{360.426}{543.396}{Let}%
\TextAt{14}{382.196}{543.396}{f}%
\TextAt{8}{390.795}{543.396}{be}%
\TextAt{8}{407.52}{543.396}{the}%
\TextAt{8}{428.477}{543.396}{conductor}%
\TextAt{8}{483.913}{543.396}{of}%
\TextAt{13}{498.044}{543.396}{L}%
\TextAt{16}{510.694}{543.396}{\Character{106}}%
\TextAt{13}{518.7}{543.396}{K}%
\TextAt{8}{529.518}{543.396}{,}%
\TextAt{13}{79.226}{526.35}{H}%
\TextAt{27}{89.849}{530.688}{f}%
\NativeMathAt{96.268}{526.35}{11.955}{\subseteq}%
\TextAt{13}{108.887}{526.35}{J}%
\TextAt{27}{116.464}{531.975}{f}%
\TextAt{21}{115.358}{523.109}{K}%
\TextAt{8}{126.463}{526.35}{the}%
\TextAt{8}{145.624}{526.35}{ideal}%
\TextAt{8}{172.592}{526.35}{group}%
\TextAt{8}{204.772}{526.35}{for}%
\TextAt{13}{221.662}{526.35}{L}%
\TextAt{8}{229.627}{526.35}{,}%
\TextAt{8}{235.98}{526.35}{and}%
\TextAt{13}{257.748}{526.35}{H}%
\TextAt{27}{267.447}{524.557}{p}%
\TextAt{8}{274.836}{526.35}{the}%
\TextAt{8}{293.997}{526.35}{smallest}%
\TextAt{8}{338.001}{526.35}{ideal}%
\TextAt{8}{364.969}{526.35}{group}%
\TextAt{8}{397.15}{526.35}{containing}%
\TextAt{13}{453.073}{526.35}{H}%
\TextAt{27}{463.696}{530.688}{f}%
\TextAt{8}{469.697}{526.35}{of}%
\TextAt{8}{482.032}{526.35}{conductor}%
\ErrataAt{79.19}{509.29}{455.81}{prime to $\mathfrak p$. If $e=(H_{\mathfrak p}\cap J_K^{\mathfrak f}:H^{\mathfrak f})$ and $\mathfrak p^f$ is the smallest power of $\mathfrak p$ in $H_{\mathfrak p}$, then}%
\TextAt{14}{263.239}{481.464}{p}%
\TextAt{9}{272.543}{481.464}{\Character{61}}%
\TextAt{9}{284.952}{481.464}{\Character{40}}%
\TextAt{14}{289.52}{481.464}{P}%
\TextAt{17}{299.414}{479.671}{1}%
\TextAt{16}{306.138}{481.464}{\Character{1}}%
\TextAt{16}{311.434}{481.464}{\Character{1}}%
\TextAt{16}{316.742}{481.464}{\Character{1}}%
\TextAt{14}{322.079}{481.464}{P}%
\TextAt{21}{331.972}{479.671}{r}%
\TextAt{9}{336.527}{481.464}{\Character{41}}%
\TextAt{21}{341.079}{486.4}{e}%
\TextAt{13}{345.505}{481.464}{\Character{59}}%
\TextAt{8}{79.226}{453.624}{where}%
\TextAt{8}{113.035}{453.624}{the}%
\TextAt{14}{133.192}{453.624}{P}%
\TextAt{21}{143.085}{451.831}{i}%
\TextAt{8}{150.364}{453.624}{are}%
\TextAt{8}{169.875}{453.624}{of}%
\TextAt{8}{183.205}{453.624}{degree}%
\TextAt{13}{219.62}{453.624}{f}%
\TextAt{8}{230.564}{453.624}{over}%
\TextAt{13}{255.61}{453.624}{K}%
\TextAt{8}{266.428}{453.624}{,}%
\TextAt{8}{273.577}{453.624}{and}%
\TextAt{13}{296.339}{453.624}{r}%
\TextAt{9}{305.261}{453.624}{\Character{61}}%
\TextAt{13}{317.686}{453.624}{n}%
\TextAt{9}{324.674}{453.624}{\Character{47}\Character{40}}%
\TextAt{13}{335.079}{453.624}{ef}%
\TextAt{9}{347.551}{453.624}{\Character{41}}%
\TextAt{8}{352.103}{453.624}{,}%
\TextAt{13}{359.253}{453.624}{n}%
\TextAt{9}{369.561}{453.624}{\Character{61}}%
\TextAt{9}{381.971}{453.624}{\Character{91}}%
\TextAt{13}{385.238}{453.624}{L}%
\TextAt{9}{396.523}{453.624}{\Character{58}}%
\TextAt{13}{403.096}{453.624}{K}%
\TextAt{9}{413.914}{453.624}{\Character{93}}%
\TextAt{8}{417.165}{453.624}{.}%
\TextAt{12}{97.159}{433.59}{Hint:}%
\TextAt{8}{132.698}{433.59}{The}%
\TextAt{8}{156.751}{433.59}{class}%
\TextAt{8}{184.189}{433.59}{field}%
\TextAt{8}{209.545}{433.59}{for}%
\TextAt{13}{227.43}{433.59}{H}%
\TextAt{27}{237.129}{431.796}{p}%
\TextAt{8}{245.514}{433.59}{is}%
\TextAt{8}{257.277}{433.59}{the}%
\TextAt{8}{277.434}{433.59}{inertia}%
\TextAt{8}{314.507}{433.59}{field}%
\TextAt{8}{339.864}{433.59}{above}%
\TextAt{14}{373.039}{433.59}{p}%
\TextAt{8}{379.023}{433.59}{.}%
\TextAt{8}{97.159}{407.909}{The}%
\TextAt{8}{121.213}{407.909}{following}%
\TextAt{8}{170.647}{407.909}{exercises}%
\TextAt{8}{218.564}{407.909}{2–6}%
\TextAt{8}{240.035}{407.909}{concern}%
\TextAt{8}{282.954}{407.909}{a}%
\TextAt{8}{292.71}{407.909}{non-abelian}%
\TextAt{8}{356.12}{407.909}{example}%
\TextAt{8}{401.968}{407.909}{of}%
\TextAt{26}{415.812}{407.909}{E.}%
\TextAt{26}{433.015}{407.909}{Artin}%
\TextAt{8}{467.147}{407.909}{.}%
\TextAt{12}{79.226}{382.229}{Exercise}%
\TextAt{12}{132.634}{382.229}{VI.7.2.}%
\TextAt{8}{176.571}{382.229}{The}%
\TextAt{8}{200.805}{382.229}{polynomial}%
\TextAt{13}{261.479}{382.229}{f}%
\TextAt{9}{268.525}{382.228}{\Character{40}}%
\TextAt{13}{273.077}{382.229}{X}%
\TextAt{9}{283.732}{382.228}{\Character{41}}%
\TextAt{9}{291.897}{382.228}{\Character{61}}%
\TextAt{13}{304.646}{382.229}{X}%
\TextAt{17}{315.301}{386.567}{5}%
\TextAt{16}{322.813}{382.229}{\Character{0}}%
\TextAt{13}{334.892}{382.228}{X}%
\TextAt{9}{348.326}{382.229}{\Character{43}}%
\TextAt{9}{360.197}{382.229}{1}%
\TextAt{8}{370.142}{382.229}{is}%
\TextAt{8}{382.086}{382.229}{irreducible.}%
\TextAt{8}{445.053}{382.229}{The}%
\TextAt{8}{469.288}{382.229}{discriminant}%
\TextAt{8}{79.226}{365.183}{of}%
\TextAt{8}{92.556}{365.183}{a}%
\TextAt{8}{102.312}{365.183}{root}%
\TextAt{13}{127.358}{365.183}{\Character{11}}%
\TextAt{8}{138.777}{365.183}{(i.e.,}%
\TextAt{8}{165.437}{365.183}{the}%
\TextAt{8}{185.593}{365.183}{discriminant}%
\TextAt{8}{252.973}{365.183}{of}%
\TextAt{18}{266.303}{365.183}{Z}%
\TextAt{9}{274.273}{365.183}{\Character{91}}%
\TextAt{13}{277.524}{365.183}{\Character{11}}%
\TextAt{9}{285.046}{365.183}{\Character{93}}%
\TextAt{8}{288.298}{365.183}{)}%
\TextAt{8}{296.75}{365.183}{is}%
\TextAt{13}{308.514}{365.183}{d}%
\TextAt{9}{317.917}{365.183}{\Character{61}}%
\TextAt{9}{330.327}{365.183}{19}%
\TextAt{16}{344.706}{365.183}{\Character{1}}%
\TextAt{9}{350.683}{365.183}{151}%
\TextAt{8}{368.242}{365.183}{.}%
\TextAt{12}{97.159}{345.148}{Hint:}%
\TextAt{8}{132.698}{345.148}{The}%
\TextAt{8}{156.751}{345.148}{discriminant}%
\TextAt{8}{224.131}{345.148}{of}%
\TextAt{8}{237.461}{345.148}{a}%
\TextAt{8}{247.216}{345.148}{root}%
\TextAt{8}{272.262}{345.148}{of}%
\TextAt{13}{285.592}{345.148}{X}%
\TextAt{17}{296.247}{349.486}{5}%
\TextAt{9}{303.636}{345.148}{\Character{43}}%
\TextAt{13}{315.398}{345.148}{aX}%
\TextAt{9}{334.854}{345.148}{\Character{43}}%
\TextAt{13}{346.616}{345.148}{b}%
\TextAt{8}{355.49}{345.148}{is}%
\TextAt{9}{367.254}{345.148}{5}%
\TextAt{17}{373.107}{349.486}{5}%
\TextAt{13}{377.839}{345.148}{b}%
\TextAt{17}{382.817}{349.486}{4}%
\TextAt{9}{390.205}{345.148}{\Character{43}}%
\TextAt{9}{401.957}{345.148}{2}%
\TextAt{17}{407.82}{349.486}{8}%
\TextAt{13}{412.552}{345.148}{a}%
\TextAt{17}{418.697}{349.486}{5}%
\TextAt{8}{423.429}{345.148}{.}%
\TextAt{12}{79.226}{313.821}{Exercise}%
\TextAt{12}{132.427}{313.821}{VI.7.3.}%
\TextAt{8}{176.183}{313.821}{Let}%
\TextAt{13}{197.152}{313.821}{k}%
\TextAt{9}{206.962}{313.821}{\Character{61}}%
\TextAt{18}{219.388}{313.821}{Q}%
\TextAt{9}{228.686}{313.821}{\Character{40}}%
\TextAt{13}{233.239}{313.821}{\Character{11}}%
\TextAt{9}{240.76}{313.821}{\Character{41}}%
\TextAt{8}{245.313}{313.821}{.}%
\TextAt{8}{253.761}{313.821}{Then}%
\TextAt{18}{284.318}{313.821}{Z}%
\TextAt{9}{292.289}{313.821}{\Character{91}}%
\TextAt{13}{295.54}{313.821}{\Character{11}}%
\TextAt{9}{303.062}{313.821}{\Character{93}}%
\TextAt{8}{310.211}{313.821}{is}%
\TextAt{8}{321.975}{313.821}{the}%
\TextAt{8}{342.131}{313.821}{ring}%
\TextAt{16}{366.197}{313.821}{\Character{79}}%
\TextAt{21}{375.715}{312.028}{k}%
\TextAt{8}{384.732}{313.821}{of}%
\TextAt{8}{398.062}{313.821}{integers}%
\TextAt{8}{441.375}{313.821}{of}%
\TextAt{13}{454.705}{313.821}{k}%
\TextAt{8}{461.195}{313.821}{.}%
\TextAt{12}{97.159}{293.787}{Hint:}%
\TextAt{8}{135.159}{293.787}{The}%
\TextAt{8}{160.443}{293.787}{discriminant}%
\TextAt{8}{229.053}{293.787}{of}%
\TextAt{18}{243.614}{293.787}{Z}%
\TextAt{9}{251.584}{293.787}{\Character{91}}%
\TextAt{13}{254.836}{293.787}{\Character{11}}%
\TextAt{9}{262.358}{293.787}{\Character{93}}%
\TextAt{8}{270.737}{293.787}{equals}%
\TextAt{8}{307.475}{293.787}{the}%
\TextAt{8}{328.862}{293.787}{discriminant}%
\TextAt{8}{397.472}{293.787}{of}%
\TextAt{16}{412.033}{293.787}{\Character{79}}%
\TextAt{21}{421.55}{291.993}{k}%
\TextAt{8}{431.798}{293.787}{because}%
\TextAt{8}{476.331}{293.787}{on}%
\TextAt{8}{493.82}{293.787}{the}%
\TextAt{8}{515.207}{293.787}{one}%
\TextAt{8}{79.226}{276.741}{hand,}%
\TextAt{8}{111.321}{276.741}{both}%
\TextAt{8}{138.431}{276.741}{differ}%
\TextAt{8}{168.136}{276.741}{only}%
\TextAt{8}{193.298}{276.741}{by}%
\TextAt{8}{209.028}{276.741}{a}%
\TextAt{8}{218.254}{276.741}{square,}%
\TextAt{8}{257.892}{276.741}{and}%
\TextAt{8}{280.125}{276.741}{on}%
\TextAt{8}{295.854}{276.741}{the}%
\TextAt{8}{315.481}{276.741}{other}%
\TextAt{8}{345.521}{276.741}{hand,}%
\TextAt{8}{377.615}{276.741}{it}%
\TextAt{8}{388.79}{276.741}{is}%
\TextAt{8}{400.024}{276.741}{squarefree.}%
\TextAt{8}{459.739}{276.741}{The}%
\TextAt{8}{483.263}{276.741}{transition}%
\TextAt{8}{79.226}{259.695}{matrix}%
\TextAt{8}{116.708}{259.695}{from}%
\TextAt{9}{143.776}{259.695}{1}%
\TextAt{13}{149.629}{259.695}{\Character{59}}%
\TextAt{13}{154.853}{259.695}{\Character{11}\Character{59}}%
\TextAt{13}{167.598}{259.695}{\Character{58}}%
\TextAt{13}{172.834}{259.695}{\Character{58}}%
\TextAt{13}{178.07}{259.695}{\Character{58}}%
\TextAt{13}{183.295}{259.695}{\Character{59}}%
\TextAt{13}{188.531}{259.695}{\Character{11}}%
\TextAt{21}{196.138}{264.033}{n}%
\TextAt{3}{201.276}{264.033}{\Character{0}}%
\TextAt{17}{207.862}{264.033}{1}%
\TextAt{8}{215.92}{259.695}{to}%
\TextAt{8}{229.659}{259.695}{an}%
\TextAt{8}{245.346}{259.695}{integral}%
\TextAt{8}{287.383}{259.695}{basis}%
\TextAt{13}{315.551}{259.695}{\Character{33}}%
\TextAt{17}{322.84}{257.902}{1}%
\TextAt{13}{327.573}{259.695}{\Character{59}}%
\TextAt{13}{332.797}{259.695}{\Character{58}}%
\TextAt{13}{338.034}{259.695}{\Character{58}}%
\TextAt{13}{343.258}{259.695}{\Character{58}}%
\TextAt{13}{348.495}{259.695}{\Character{59}}%
\TextAt{13}{353.731}{259.695}{\Character{33}}%
\TextAt{21}{361.083}{257.902}{n}%
\TextAt{8}{370.045}{259.695}{of}%
\TextAt{16}{382.803}{259.695}{\Character{79}}%
\TextAt{21}{392.321}{257.902}{k}%
\TextAt{8}{400.766}{259.695}{is}%
\TextAt{8}{411.958}{259.695}{therefore}%
\TextAt{8}{460.486}{259.695}{invertible}%
\TextAt{8}{511.621}{259.695}{over}%
\TextAt{18}{79.226}{242.649}{Z}%
\TextAt{8}{87.196}{242.649}{.}%
\TextAt{12}{79.226}{211.322}{Exercise}%
\TextAt{12}{132.642}{211.322}{VI.7.4.}%
\TextAt{8}{176.584}{211.322}{The}%
\TextAt{8}{200.825}{211.322}{decomposition}%
\TextAt{8}{278.147}{211.322}{field}%
\TextAt{13}{303.69}{211.322}{K}%
\TextAt{16}{318.147}{211.323}{\Character{106}}%
\TextAt{18}{325.108}{211.322}{Q}%
\TextAt{8}{338.49}{211.322}{of}%
\TextAt{13}{352.007}{211.322}{f}%
\TextAt{9}{359.054}{211.323}{\Character{40}}%
\TextAt{13}{363.606}{211.322}{X}%
\TextAt{9}{374.261}{211.323}{\Character{41}}%
\TextAt{8}{382.898}{211.322}{has}%
\TextAt{8}{403.958}{211.322}{as}%
\TextAt{8}{418.515}{211.322}{Galois}%
\TextAt{8}{454.615}{211.322}{group}%
\TextAt{8}{487.978}{211.322}{the}%
\TextAt{8}{508.321}{211.322}{sym-}%
\TextAt{8}{79.226}{194.277}{metric}%
\TextAt{8}{115.642}{194.277}{group}%
\TextAt{14}{148.817}{194.277}{S}%
\TextAt{17}{158.723}{192.483}{5}%
\TextAt{8}{163.456}{194.277}{,}%
\TextAt{8}{170.605}{194.277}{i.e.,}%
\TextAt{8}{192.71}{194.277}{it}%
\TextAt{8}{204.414}{194.277}{is}%
\TextAt{8}{216.178}{194.277}{of}%
\TextAt{8}{229.508}{194.277}{degree}%
\TextAt{9}{265.923}{194.277}{120}%
\TextAt{8}{283.482}{194.277}{.}%
\TextAt{12}{79.226}{162.95}{Exercise}%
\TextAt{12}{132.427}{162.95}{VI.7.5.}%
\ErrataAt{176.17}{162.95}{358.83}{$k$ has class number $1$.}%
\ErrataAt{97.1}{142.92}{437.9}{\textbf{Hint:} Show, using chap. I, \hyperlink{nav.ex.1.6.3}{\S6, exercise 3}, that every ideal class of $k$ contains an ideal}%
\TextAt{14}{79.226}{125.869}{a}%
\TextAt{8}{89.183}{125.869}{with}%
\TextAt{14}{115.919}{125.869}{N}%
\TextAt{9}{125.867}{125.869}{\Character{40}}%
\TextAt{14}{130.419}{125.869}{a}%
\TextAt{9}{136.402}{125.869}{\Character{41}}%
\TextAt{13}{144.405}{125.869}{\Character{60}}%
\TextAt{9}{156.96}{125.869}{4}%
\TextAt{8}{162.813}{125.869}{.}%
\TextAt{8}{171.49}{125.869}{If}%
\TextAt{14}{183.258}{125.869}{N}%
\TextAt{9}{193.206}{125.869}{\Character{40}}%
\TextAt{14}{197.758}{125.869}{a}%
\TextAt{9}{203.741}{125.869}{\Character{41}}%
\NativeMathAt{211.744}{125.869}{11.955}{\neq}%
\TextAt{9}{224.285}{125.869}{1}%
\TextAt{8}{230.152}{125.869}{,}%
\TextAt{8}{237.396}{125.869}{then}%
\TextAt{14}{264.133}{125.869}{a}%
\TextAt{8}{274.089}{125.869}{has}%
\TextAt{8}{295.039}{125.869}{to}%
\TextAt{8}{309.425}{125.869}{be}%
\TextAt{8}{325.426}{125.869}{a}%
\TextAt{8}{335.257}{125.869}{prime}%
\TextAt{8}{368.497}{125.869}{ideal}%
\TextAt{14}{396.536}{125.869}{p}%
\TextAt{8}{406.493}{125.869}{such}%
\TextAt{8}{432.966}{125.869}{that}%
\TextAt{14}{458.411}{125.869}{N}%
\TextAt{9}{468.359}{125.869}{\Character{40}}%
\TextAt{14}{472.911}{125.869}{p}%
\TextAt{9}{478.894}{125.869}{\Character{41}}%
\TextAt{9}{486.88}{125.869}{\Character{61}}%
\TextAt{9}{499.433}{125.869}{2}%
\TextAt{8}{509.278}{125.869}{or}%
\TextAt{9}{523.665}{125.869}{3}%
\TextAt{8}{529.518}{125.869}{.}%
\TextAt{8}{79.226}{108.824}{Hence}%
\TextAt{16}{114.004}{108.824}{\Character{79}}%
\TextAt{21}{123.521}{107.03}{k}%
\TextAt{9}{128.641}{108.824}{\Character{47}}%
\TextAt{14}{134.494}{108.824}{p}%
\TextAt{9}{143.798}{108.824}{\Character{61}}%
\TextAt{18}{156.224}{108.824}{Z}%
\TextAt{9}{164.194}{108.824}{\Character{47}2}%
\TextAt{18}{175.9}{108.824}{Z}%
\TextAt{8}{187.767}{108.824}{or}%
\TextAt{9}{202.078}{108.824}{\Character{61}}%
\TextAt{18}{214.503}{108.824}{Z}%
\TextAt{9}{222.473}{108.824}{\Character{47}3}%
\TextAt{18}{234.179}{108.824}{Z}%
\TextAt{8}{242.149}{108.824}{,}%
\TextAt{8}{249.298}{108.824}{so}%
\TextAt{13}{263.669}{108.824}{f}%
\TextAt{8}{274.612}{108.824}{has}%
\TextAt{8}{295.486}{108.824}{a}%
\TextAt{8}{305.242}{108.824}{root}%
\TextAt{8}{330.288}{108.824}{mod}%
\TextAt{9}{356.625}{108.824}{2}%
\TextAt{8}{366.375}{108.824}{or}%
\TextAt{9}{380.686}{108.824}{3}%
\TextAt{8}{386.539}{108.824}{,}%
\TextAt{8}{393.688}{108.824}{which}%
\TextAt{8}{427.174}{108.824}{is}%
\TextAt{8}{438.938}{108.824}{not}%
\TextAt{8}{459.752}{108.824}{the}%
\TextAt{8}{479.908}{108.824}{case.}%
\TextAt{8}{297.211}{43.033}{377}%
\Destination{78.23}{809.815}{nav.page.377}%
\Destination{79.226}{574.455}{nav.ex.6.7.1}%
\Destination{79.226}{396.255}{nav.ex.6.7.2}%
\Destination{79.226}{327.855}{nav.ex.6.7.3}%
\Destination{79.226}{225.355}{nav.ex.6.7.4}%
\Destination{79.226}{176.955}{nav.ex.6.7.5}%
\Destination{79.36}{776.155}{nav.print.413}%
\Destination{79.36}{226.355}{nav.print.414}%
\Hotlink{411.391}{583.955}{19.22}{11.047}{nav.cite.N048}%
\Hotlink{436.76}{583.955}{19.22}{11.047}{nav.cite.N024}%
\Hotlink{218.064}{405.755}{6.858}{11.047}{nav.ex.6.7.2}%
\Hotlink{297.982}{140.755}{6.858}{11.047}{nav.ex.1.6.3}%
\Hotlink{226.111}{140.755}{5.22}{11.047}{nav.chapter.1}%
\Destination{79.226}{179.95}{errata.140}%
\Destination{79.226}{526.29}{errata.145}%
}

% Page 390 (body page 378)
\Page{%
\TextAt{2}{79.226}{796.943}{VI.}%
\TextAt{2}{98.012}{796.943}{Global}%
\TextAt{2}{133.248}{796.943}{Class}%
\TextAt{2}{161.841}{796.943}{Field}%
\TextAt{2}{189.572}{796.943}{Theory}%
\Vector{q 0 G 0 g 0.299 w 7.226 18.729 m 460.769 18.729 l S Q}%
\TextAt{12}{79.226}{761.124}{Exercise}%
\TextAt{12}{132.427}{761.124}{VI.7.6.}%
\TextAt{8}{176.183}{761.124}{Show}%
\TextAt{8}{207.075}{761.124}{that}%
\TextAt{13}{232.444}{761.124}{K}%
\TextAt{16}{246.582}{761.124}{\Character{106}}%
\TextAt{18}{253.224}{761.124}{Q}%
\TextAt{9}{262.522}{761.124}{\Character{40}}%
\TextAt{16}{267.075}{771.013}{\Character{112}}%
\Vector{q 0 G 0 g 0.478 w 205.037 -2.748 m 242.936 -2.748 l S Q}%
\TextAt{9}{277.037}{761.124}{19}%
\TextAt{16}{291.4}{761.124}{\Character{1}}%
\TextAt{9}{297.378}{761.124}{151\Character{41}}%
\TextAt{8}{323.386}{761.124}{is}%
\TextAt{8}{335.15}{761.124}{a}%
\TextAt{8}{344.906}{761.124}{(non-abelian!)}%
\TextAt{8}{421.977}{761.124}{unramified}%
\TextAt{8}{480.509}{761.124}{extension.}%
\TextAt{12}{79.226}{730.362}{Exercise}%
\TextAt{12}{131.454}{730.362}{VI.7.7.}%
\TextAt{8}{174.364}{730.362}{For}%
\TextAt{8}{194.488}{730.362}{every}%
\TextAt{8}{224.534}{730.362}{Galois}%
\TextAt{8}{259.602}{730.362}{extension}%
\TextAt{13}{310.522}{730.362}{L}%
\TextAt{16}{321.807}{730.362}{\Character{106}}%
\TextAt{13}{328.449}{730.362}{K}%
\TextAt{8}{342.318}{730.362}{of}%
\TextAt{8}{354.803}{730.362}{finite}%
\TextAt{8}{383.869}{730.362}{algebraic}%
\TextAt{8}{432.458}{730.362}{number}%
\TextAt{8}{474.208}{730.362}{fields,}%
\TextAt{8}{506.755}{730.362}{there}%
\TextAt{8}{79.226}{713.316}{exist}%
\TextAt{8}{106.978}{713.316}{infinitely}%
\TextAt{8}{156.129}{713.316}{many}%
\TextAt{8}{188.052}{713.316}{finite}%
\TextAt{8}{218.016}{713.316}{extensions}%
\TextAt{13}{274.447}{713.316}{K}%
\TextAt{3}{285.265}{717.654}{\Character{48}}%
\TextAt{8}{292.009}{713.316}{such}%
\TextAt{8}{318.457}{713.316}{that}%
\TextAt{13}{343.877}{713.316}{L}%
\TextAt{16}{354.533}{713.316}{\Character{92}}%
\TextAt{13}{365.195}{713.316}{K}%
\TextAt{3}{376.013}{717.654}{\Character{48}}%
\TextAt{9}{382.216}{713.316}{\Character{61}}%
\TextAt{13}{394.729}{713.316}{K}%
\TextAt{8}{405.547}{713.316}{,}%
\TextAt{8}{412.76}{713.316}{and}%
\TextAt{8}{435.574}{713.316}{such}%
\TextAt{8}{462.022}{713.316}{that}%
\TextAt{13}{487.442}{713.316}{LK}%
\TextAt{3}{506.224}{717.654}{\Character{48}}%
\TextAt{16}{512.428}{713.316}{\Character{106}}%
\TextAt{13}{519.157}{713.316}{K}%
\TextAt{3}{529.974}{717.654}{\Character{48}}%
\TextAt{8}{79.226}{696.27}{is}%
\TextAt{8}{90.99}{696.27}{unramified.}%
\TextAt{12}{97.159}{676.235}{Hint:}%
\TextAt{8}{134.582}{676.235}{Let}%
\TextAt{13}{156.493}{676.235}{S}%
\TextAt{8}{169.228}{676.235}{be}%
\TextAt{8}{186.094}{676.235}{the}%
\TextAt{8}{207.193}{676.235}{set}%
\TextAt{8}{226.402}{676.235}{of}%
\TextAt{8}{240.675}{676.235}{places}%
\TextAt{8}{276.143}{676.235}{ramified}%
\TextAt{8}{322.61}{676.235}{in}%
\TextAt{13}{337.205}{676.235}{L}%
\TextAt{16}{350.096}{676.235}{\Character{106}}%
\TextAt{13}{358.343}{676.235}{K}%
\TextAt{8}{369.161}{676.235}{,}%
\TextAt{8}{377.488}{676.235}{and}%
\TextAt{8}{401.192}{676.235}{let}%
\TextAt{13}{419.039}{676.235}{L}%
\TextAt{27}{427.004}{674.442}{p}%
\TextAt{9}{436.417}{676.235}{\Character{61}}%
\TextAt{13}{450.448}{676.235}{K}%
\TextAt{27}{460.407}{674.442}{p}%
\TextAt{9}{464.894}{676.235}{\Character{40}}%
\TextAt{13}{469.446}{676.235}{\Character{11}}%
\TextAt{27}{476.891}{674.442}{p}%
\TextAt{9}{481.378}{676.235}{\Character{41}}%
\TextAt{8}{485.931}{676.235}{.}%
\TextAt{8}{497.205}{676.235}{By}%
\TextAt{8}{516.51}{676.235}{the}%
\TextAt{8}{79.226}{659.189}{approximation}%
\TextAt{8}{158.086}{659.189}{theorem,}%
\TextAt{8}{207.855}{659.189}{choose}%
\TextAt{8}{245.78}{659.189}{an}%
\TextAt{8}{262.832}{659.189}{algebraic}%
\TextAt{8}{313.06}{659.189}{number}%
\TextAt{13}{356.449}{659.189}{\Character{11}}%
\TextAt{8}{368.661}{659.189}{which,}%
\TextAt{8}{406.391}{659.189}{for}%
\TextAt{8}{425.069}{659.189}{every}%
\TextAt{14}{456.754}{659.189}{p}%
\NativeMathAt{467.409}{659.189}{11.955}{\in}%
\TextAt{13}{480.051}{659.189}{S}%
\TextAt{8}{487.947}{659.189}{,}%
\TextAt{8}{496.087}{659.189}{is}%
\TextAt{8}{508.644}{659.189}{close}%
\TextAt{8}{79.226}{642.143}{to}%
\TextAt{13}{93.652}{642.143}{\Character{11}}%
\TextAt{27}{101.098}{640.35}{p}%
\TextAt{8}{109.598}{642.143}{when}%
\TextAt{8}{140.271}{642.143}{embedded}%
\TextAt{8}{195.656}{642.143}{into}%
\Vector{q 0 G 0 g 0.478 w 147.514 -122.014 m 158.332 -122.014 l S Q}%
\TextAt{13}{219.514}{642.143}{K}%
\TextAt{27}{230.332}{640.35}{p}%
\TextAt{8}{234.819}{642.143}{.}%
\TextAt{8}{243.614}{642.143}{Then}%
\TextAt{13}{274.288}{642.143}{K}%
\TextAt{27}{284.247}{640.35}{p}%
\TextAt{9}{288.734}{642.143}{\Character{40}}%
\TextAt{13}{293.286}{642.143}{\Character{11}}%
\TextAt{27}{300.731}{640.35}{p}%
\TextAt{9}{305.218}{642.143}{\Character{41}}%
\NativeMathAt{313.289}{642.143}{11.955}{\subseteq}%
\TextAt{13}{326.105}{642.143}{K}%
\TextAt{27}{336.064}{640.35}{p}%
\TextAt{9}{340.551}{642.143}{\Character{40}}%
\TextAt{13}{345.104}{642.143}{\Character{11}}%
\TextAt{9}{352.625}{642.143}{\Character{41}}%
\TextAt{8}{361.191}{642.143}{by}%
\TextAt{8}{377.566}{642.143}{Krasner’s}%
\TextAt{8}{429.83}{642.143}{lemma,}%
\TextAt{8}{470.945}{642.143}{chap.}%
\TextAt{8}{501.953}{642.143}{II,}%
\TextAt{8}{517.981}{642.143}{§6,}%
\TextAt{8}{79.226}{625.097}{exercise}%
\TextAt{8}{122.959}{625.097}{2.}%
\TextAt{8}{138.559}{625.097}{Put}%
\TextAt{13}{161.909}{625.097}{K}%
\TextAt{3}{172.726}{629.436}{\Character{48}}%
\TextAt{9}{179.577}{625.097}{\Character{61}}%
\TextAt{13}{192.738}{625.097}{K}%
\TextAt{9}{203.556}{625.097}{\Character{40}}%
\TextAt{13}{208.108}{625.097}{\Character{11}}%
\TextAt{9}{215.63}{625.097}{\Character{41}}%
\TextAt{8}{224.511}{625.097}{and}%
\TextAt{8}{247.705}{625.097}{show}%
\TextAt{8}{277.139}{625.097}{that}%
\TextAt{13}{302.939}{625.097}{LK}%
\TextAt{3}{321.722}{629.436}{\Character{48}}%
\TextAt{16}{328.573}{625.097}{\Character{106}}%
\TextAt{13}{335.949}{625.097}{K}%
\TextAt{3}{346.767}{629.436}{\Character{48}}%
\TextAt{8}{353.891}{625.097}{is}%
\TextAt{8}{366.086}{625.097}{unramified.}%
\TextAt{8}{430.463}{625.097}{To}%
\TextAt{8}{448.122}{625.097}{show}%
\TextAt{8}{477.557}{625.097}{that}%
\TextAt{13}{503.357}{625.097}{\Character{11}}%
\TextAt{8}{515.207}{625.097}{can}%
\TextAt{8}{79.226}{608.051}{be}%
\TextAt{8}{95.677}{608.051}{chosen}%
\TextAt{8}{133.66}{608.051}{such}%
\TextAt{8}{160.584}{608.051}{that}%
\TextAt{13}{186.48}{608.051}{L}%
\TextAt{16}{197.46}{608.051}{\Character{92}}%
\TextAt{13}{208.446}{608.051}{K}%
\TextAt{3}{219.264}{612.39}{\Character{48}}%
\TextAt{9}{226.278}{608.051}{\Character{61}}%
\TextAt{13}{239.602}{608.051}{K}%
\TextAt{8}{254.844}{608.051}{use}%
\TextAt{8}{275.587}{608.051}{(3.7),}%
\TextAt{8}{307.473}{608.051}{and}%
\TextAt{8}{330.762}{608.051}{the}%
\TextAt{8}{351.446}{608.051}{fact}%
\TextAt{8}{375.058}{608.051}{that}%
\TextAt{13}{400.954}{608.051}{G}%
\TextAt{9}{410.188}{608.051}{\Character{40}}%
\TextAt{13}{414.74}{608.051}{L}%
\TextAt{16}{426.924}{608.051}{\Character{106}}%
\TextAt{13}{434.463}{608.051}{K}%
\TextAt{9}{445.281}{608.051}{\Character{41}}%
\TextAt{8}{454.258}{608.051}{is}%
\TextAt{8}{466.549}{608.051}{generated}%
\TextAt{8}{520.408}{608.051}{by}%
\TextAt{8}{79.226}{591.005}{elements}%
\TextAt{8}{127.083}{591.005}{of}%
\TextAt{8}{140.413}{591.005}{prime}%
\TextAt{8}{173.576}{591.005}{power}%
\TextAt{8}{207.72}{591.005}{order.}%
\TextAt{15}{79.226}{539.073}{§}%
\TextAt{15}{93.785}{539.073}{8}%
\TextAt{15}{120.693}{539.073}{The}%
\TextAt{15}{160.202}{539.073}{Reciprocity}%
\TextAt{15}{264.063}{539.073}{Law}%
\TextAt{15}{305.036}{539.073}{of}%
\TextAt{15}{327.106}{539.073}{the}%
\TextAt{15}{360.693}{539.073}{Power}%
\TextAt{15}{419.226}{539.073}{Residues}%
\TextAt{8}{79.226}{507.187}{In}%
\TextAt{8}{94.674}{507.187}{class}%
\TextAt{8}{122.938}{507.187}{field}%
\TextAt{8}{149.122}{507.187}{theory}%
\TextAt{8}{186.699}{507.187}{Gauss’s}%
\TextAt{8}{230.062}{507.187}{reciprocity}%
\TextAt{8}{289.098}{507.187}{law}%
\TextAt{8}{311.062}{507.187}{meets}%
\TextAt{8}{345.112}{507.187}{its}%
\TextAt{8}{362.257}{507.187}{most}%
\TextAt{8}{391.765}{507.187}{general}%
\TextAt{8}{432.916}{507.187}{and}%
\TextAt{8}{456.506}{507.187}{definite}%
\TextAt{8}{498.948}{507.187}{formu-}%
\TextAt{8}{79.226}{490.141}{lation.}%
\TextAt{8}{119.385}{490.141}{Let}%
\TextAt{13}{141.166}{490.141}{n}%
\TextAt{8}{152.862}{490.141}{be}%
\TextAt{8}{169.597}{490.141}{a}%
\TextAt{8}{180.163}{490.141}{positive}%
\TextAt{8}{224.288}{490.141}{integer}%
\NativeMathAt{263.797}{490.141}{11.955}{\geq}%
\TextAt{9}{277.605}{490.141}{2}%
\TextAt{8}{288.166}{490.141}{and}%
\TextAt{13}{311.74}{490.141}{K}%
\TextAt{8}{327.266}{490.141}{a}%
\TextAt{8}{337.832}{490.141}{number}%
\TextAt{8}{381.239}{490.141}{field}%
\TextAt{8}{407.407}{490.141}{containing}%
\TextAt{8}{465.136}{490.141}{the}%
\TextAt{8}{486.104}{490.141}{group}%
\TextAt{13}{520.09}{490.141}{\Character{22}}%
\TextAt{21}{527.133}{488.348}{n}%
\TextAt{8}{79.226}{473.096}{of}%
\TextAt{13}{93.334}{473.096}{n}%
\TextAt{8}{100.322}{473.096}{-th}%
\TextAt{8}{119.953}{473.096}{roots}%
\TextAt{8}{150.391}{473.096}{of}%
\TextAt{8}{164.499}{473.096}{unity.}%
\TextAt{8}{200.973}{473.096}{In}%
\TextAt{8}{216.372}{473.096}{chap.}%
\TextAt{8}{248.042}{473.096}{V,}%
\TextAt{8}{264.744}{473.096}{§3,}%
\TextAt{8}{284.402}{473.096}{we}%
\TextAt{8}{302.407}{473.096}{introduced,}%
\TextAt{8}{365.164}{473.096}{for}%
\TextAt{8}{383.827}{473.096}{every}%
\TextAt{8}{415.497}{473.096}{place}%
\TextAt{14}{446.186}{473.096}{p}%
\TextAt{8}{456.845}{473.096}{of}%
\TextAt{13}{470.953}{473.096}{K}%
\TextAt{8}{481.77}{473.096}{,}%
\TextAt{8}{489.892}{473.096}{the}%
\TextAt{13}{510.826}{473.096}{n}%
\TextAt{8}{517.814}{473.096}{-th}%
\TextAt{8}{79.226}{456.05}{Hilbert}%
\TextAt{8}{119.539}{456.05}{symbol}%
\TextAt{20}{241.759}{452.136}{\Character{18}}%
\TextAt{13}{255.653}{443.367}{\Character{59}}%
\Vector{q 0 G 0 g 0.478 w 179.755 -335.732 m 190.802 -335.732 l S Q}%
\TextAt{14}{254.287}{427.079}{p}%
\TextAt{20}{263.997}{452.136}{\Character{19}}%
\TextAt{9}{276.118}{435.279}{\Character{58}}%
\TextAt{13}{282.691}{435.279}{K}%
\TextAt{3}{293.508}{440.215}{\Character{3}}%
\TextAt{27}{292.65}{432.324}{p}%
\TextAt{16}{300.897}{435.279}{\Character{2}}%
\TextAt{13}{312.853}{435.279}{K}%
\TextAt{3}{323.67}{440.215}{\Character{3}}%
\TextAt{27}{322.812}{432.324}{p}%
\NativeArrowAt{331.723}{435.279}{11.955}{19.26}{\longrightarrow}%
\TextAt{13}{354.305}{435.279}{\Character{22}}%
\TextAt{21}{361.348}{433.486}{n}%
\TextAt{13}{366.985}{435.279}{\Character{58}}%
\TextAt{8}{79.226}{409.074}{It}%
\TextAt{8}{91.899}{409.074}{is}%
\TextAt{8}{103.663}{409.074}{given}%
\TextAt{8}{134.232}{409.074}{via}%
\TextAt{8}{153.42}{409.074}{the}%
\TextAt{8}{173.576}{409.074}{norm}%
\TextAt{8}{204.146}{409.074}{residue}%
\TextAt{8}{243.873}{409.074}{symbol}%
\TextAt{8}{283.935}{409.074}{by}%
\TextAt{9}{218.768}{378.486}{\Character{40}}%
\TextAt{13}{223.32}{378.486}{a\Character{59}}%
\TextAt{13}{234.689}{378.486}{K}%
\TextAt{27}{244.668}{376.693}{p}%
\TextAt{9}{249.155}{378.486}{\Character{40}}%
\TextAt{22}{257.028}{385.1}{n}%
\TextAt{16}{254.987}{389.258}{\Character{112}}%
\Vector{q 0 G 0 g 0.478 w 192.95 -384.503 m 197.927 -384.503 l S Q}%
\TextAt{13}{264.95}{378.486}{b}%
\TextAt{9}{269.927}{378.486}{\Character{41}}%
\TextAt{16}{277.8}{378.486}{\Character{106}}%
\TextAt{13}{284.442}{378.486}{K}%
\TextAt{27}{294.401}{376.693}{p}%
\TextAt{9}{298.888}{378.486}{\Character{41}}%
\TextAt{22}{306.761}{385.1}{n}%
\TextAt{16}{304.72}{389.258}{\Character{112}}%
\Vector{q 0 G 0 g 0.478 w 242.683 -384.503 m 247.66 -384.503 l S Q}%
\TextAt{13}{314.683}{378.486}{b}%
\TextAt{9}{322.981}{378.486}{\Character{61}}%
\TextAt{20}{335.406}{395.343}{\Character{18}}%
\TextAt{13}{345.402}{386.574}{a\Character{59}}%
\TextAt{13}{356.771}{386.574}{b}%
\Vector{q 0 G 0 g 0.478 w 273.402 -392.525 m 289.768 -392.525 l S Q}%
\TextAt{14}{350.594}{370.285}{p}%
\TextAt{20}{362.964}{395.343}{\Character{19}}%
\TextAt{22}{377.078}{385.1}{n}%
\TextAt{16}{375.037}{389.258}{\Character{112}}%
\Vector{q 0 G 0 g 0.478 w 312.999 -384.503 m 317.976 -384.503 l S Q}%
\TextAt{13}{384.999}{378.486}{b\Character{58}}%
\TextAt{8}{79.226}{347.91}{These}%
\TextAt{8}{113.095}{347.91}{symbols}%
\TextAt{8}{157.772}{347.91}{all}%
\TextAt{8}{174.031}{347.91}{fit}%
\TextAt{8}{188.987}{347.91}{together}%
\TextAt{8}{235.169}{347.91}{in}%
\TextAt{8}{248.822}{347.91}{the}%
\TextAt{8}{268.979}{347.91}{following}%
\TextAt{19}{318.413}{347.91}{product}%
\TextAt{19}{359.814}{347.91}{formula}%
\TextAt{8}{399.123}{347.91}{.}%
\TextAt{12}{79.226}{322.511}{Theorem}%
\TextAt{12}{136.24}{322.511}{(8.1).}%
\TextAt{8}{171.556}{322.511}{For}%
\TextAt{13}{192.525}{322.511}{a\Character{59}}%
\TextAt{13}{203.894}{322.511}{b}%
\NativeMathAt{212.212}{322.511}{11.955}{\in}%
\TextAt{13}{223.503}{322.511}{K}%
\TextAt{3}{234.321}{326.85}{\Character{3}}%
\TextAt{8}{242.951}{322.511}{one}%
\TextAt{8}{264.41}{322.511}{has}%
\TextAt{20}{266.759}{302.731}{\Character{89}}%
\TextAt{27}{272.403}{277.625}{p}%
\TextAt{20}{284.028}{308.23}{\Character{18}}%
\TextAt{13}{294.024}{299.461}{a\Character{59}}%
\TextAt{13}{305.393}{299.461}{b}%
\Vector{q 0 G 0 g 0.478 w 222.024 -479.638 m 238.39 -479.638 l S Q}%
\TextAt{14}{299.215}{283.173}{p}%
\TextAt{20}{311.585}{308.23}{\Character{19}}%
\TextAt{9}{323.706}{291.373}{\Character{61}}%
\TextAt{9}{336.115}{291.373}{1}%
\TextAt{13}{341.985}{291.373}{\Character{58}}%
\TextAt{19}{79.226}{238.449}{Proof.}%
\TextAt{8}{116}{238.449}{From}%
\TextAt{8}{146.725}{238.449}{(5.7),}%
\TextAt{8}{177.952}{238.449}{we}%
\TextAt{8}{195.179}{238.449}{find}%
\TextAt{20}{189.177}{224.168}{\Character{34}}%
\TextAt{20}{196.151}{215.082}{\Character{89}}%
\TextAt{27}{201.795}{189.976}{p}%
\TextAt{20}{213.42}{220.581}{\Character{18}}%
\TextAt{13}{223.415}{211.812}{a\Character{59}}%
\TextAt{13}{234.784}{211.812}{b}%
\Vector{q 0 G 0 g 0.478 w 151.415 -567.287 m 167.782 -567.287 l S Q}%
\TextAt{14}{228.607}{195.524}{p}%
\TextAt{20}{240.977}{220.581}{\Character{19}}%
\TextAt{20}{249.777}{224.168}{\Character{35}}%
\TextAt{22}{262.065}{210.339}{n}%
\TextAt{16}{260.024}{214.497}{\Character{112}}%
\Vector{q 0 G 0 g 0.478 w 197.986 -559.264 m 202.964 -559.264 l S Q}%
\TextAt{13}{269.986}{203.725}{b}%
\TextAt{9}{278.284}{203.725}{\Character{61}}%
\TextAt{20}{290.71}{224.168}{\Character{34}}%
\TextAt{20}{297.684}{215.082}{\Character{89}}%
\TextAt{27}{303.327}{189.976}{p}%
\TextAt{9}{312.96}{203.725}{\Character{40}}%
\TextAt{13}{317.512}{203.725}{a\Character{59}}%
\TextAt{13}{328.881}{203.725}{K}%
\TextAt{27}{338.86}{201.931}{p}%
\TextAt{9}{343.347}{203.725}{\Character{40}}%
\TextAt{22}{351.22}{210.339}{n}%
\TextAt{16}{349.18}{214.497}{\Character{112}}%
\Vector{q 0 G 0 g 0.478 w 287.142 -559.264 m 292.119 -559.264 l S Q}%
\TextAt{13}{359.142}{203.725}{b}%
\TextAt{9}{364.119}{203.725}{\Character{41}}%
\TextAt{16}{371.992}{203.725}{\Character{106}}%
\TextAt{13}{378.634}{203.725}{K}%
\TextAt{27}{388.593}{201.931}{p}%
\TextAt{9}{393.08}{203.725}{\Character{41}}%
\TextAt{20}{397.632}{224.168}{\Character{35}}%
\TextAt{22}{409.92}{210.339}{n}%
\TextAt{16}{407.879}{214.497}{\Character{112}}%
\Vector{q 0 G 0 g 0.478 w 345.841 -559.264 m 350.819 -559.264 l S Q}%
\TextAt{13}{417.841}{203.725}{b}%
\TextAt{9}{278.284}{171.357}{\Character{61}}%
\TextAt{9}{290.694}{171.357}{\Character{40}}%
\TextAt{13}{295.262}{171.357}{a\Character{59}}%
\TextAt{13}{306.631}{171.357}{K}%
\TextAt{9}{317.469}{171.357}{\Character{40}}%
\TextAt{22}{325.342}{177.971}{n}%
\TextAt{16}{323.301}{182.129}{\Character{112}}%
\Vector{q 0 G 0 g 0.478 w 261.264 -591.632 m 266.241 -591.632 l S Q}%
\TextAt{13}{333.264}{171.357}{b}%
\TextAt{9}{338.241}{171.357}{\Character{41}}%
\TextAt{16}{346.114}{171.357}{\Character{106}}%
\TextAt{13}{352.756}{171.357}{K}%
\TextAt{9}{363.574}{171.357}{\Character{41}}%
\TextAt{22}{371.447}{177.971}{n}%
\TextAt{16}{369.406}{182.129}{\Character{112}}%
\Vector{q 0 G 0 g 0.478 w 307.369 -591.632 m 312.346 -591.632 l S Q}%
\TextAt{13}{379.369}{171.357}{b}%
\TextAt{9}{387.666}{171.357}{\Character{61}}%
\TextAt{22}{403.413}{177.971}{n}%
\TextAt{16}{401.372}{182.129}{\Character{112}}%
\Vector{q 0 G 0 g 0.478 w 339.335 -591.632 m 344.312 -591.632 l S Q}%
\TextAt{13}{411.335}{171.357}{b\Character{59}}%
\TextAt{8}{79.226}{144.505}{and}%
\TextAt{8}{101.989}{144.505}{hence}%
\TextAt{8}{134.495}{144.505}{the}%
\TextAt{8}{154.651}{144.505}{theorem.}%
\Vector{q 0 G 0 g 0.399 w 452.334 -629.495 m 452.334 -621.426 l S Q}%
\Vector{q 0 G 0 g 0.399 w 452.534 -621.625 m 459.707 -621.625 l S Q}%
\Vector{q 0 G 0 g 0.399 w 452.534 -629.296 m 459.707 -629.296 l S Q}%
\Vector{q 0 G 0 g 0.399 w 459.906 -629.495 m 459.906 -621.426 l S Q}%
\TextAt{8}{97.159}{116.118}{In}%
\TextAt{8}{113.215}{116.118}{chap.}%
\TextAt{8}{145.542}{116.118}{V,}%
\TextAt{8}{162.902}{116.118}{§3,}%
\TextAt{8}{183.381}{116.118}{we}%
\TextAt{8}{202.044}{116.118}{defined}%
\TextAt{8}{243.792}{116.118}{the}%
\TextAt{13}{265.383}{116.118}{n}%
\TextAt{8}{272.371}{116.118}{-th}%
\TextAt{8}{292.659}{116.118}{power}%
\TextAt{8}{328.238}{116.118}{residue}%
\TextAt{8}{369.4}{116.118}{symbol}%
\TextAt{8}{410.897}{116.118}{in}%
\TextAt{8}{425.985}{116.118}{terms}%
\TextAt{8}{459.998}{116.118}{of}%
\TextAt{8}{474.763}{116.118}{the}%
\TextAt{8}{496.354}{116.118}{Hilbert}%
\TextAt{8}{79.226}{99.072}{symbol:}%
\TextAt{20}{263.209}{95.158}{\Character{18}}%
\TextAt{13}{273.205}{86.389}{a}%
\Vector{q 0 G 0 g 0.478 w 201.205 -692.71 m 207.35 -692.71 l S Q}%
\TextAt{14}{273.286}{70.101}{p}%
\TextAt{20}{280.545}{95.158}{\Character{19}}%
\TextAt{9}{292.667}{78.301}{\Character{61}}%
\TextAt{20}{305.092}{95.158}{\Character{18}}%
\TextAt{13}{315.088}{86.389}{\Character{25}}%
\TextAt{13}{322.142}{86.389}{\Character{59}}%
\TextAt{13}{327.378}{86.389}{a}%
\Vector{q 0 G 0 g 0.478 w 243.088 -692.71 m 261.546 -692.71 l S Q}%
\TextAt{14}{321.326}{70.101}{p}%
\TextAt{20}{334.742}{95.158}{\Character{19}}%
\TextAt{13}{345.535}{78.301}{\Character{59}}%
\TextAt{8}{297.211}{43.033}{378}%
\Destination{78.23}{809.815}{nav.page.378}%
\Destination{79.23}{565.605}{nav.section.6.8}%
\Destination{79.226}{775.155}{nav.ex.6.7.6}%
\Destination{79.226}{744.355}{nav.ex.6.7.7}%
\Destination{79.36}{336.555}{nav.number.6.8.1}%
\Destination{79.23}{565.605}{outline.63}%
\Hotlink{501.453}{639.955}{9.751}{11.047}{nav.chapter.2}%
\Hotlink{523.16}{639.955}{6.858}{11.047}{nav.section.2.6}%
\Hotlink{247.542}{470.855}{9.775}{11.047}{nav.chapter.5}%
\Hotlink{269.923}{470.855}{6.858}{11.047}{nav.section.5.3}%
\Hotlink{145.042}{113.955}{9.775}{11.047}{nav.chapter.5}%
\Hotlink{168.081}{113.955}{6.858}{11.047}{nav.section.5.3}%
\Hotlink{275.087}{605.855}{25.078}{11.047}{nav.number.6.3.7}%
\Hotlink{146.225}{236.255}{25.078}{11.047}{nav.number.6.5.7}%
\Destination{79.226}{507.141}{errata.146}%
}

% Page 391 (body page 379)
\Page{%
\TextAt{2}{79.226}{796.943}{VI.}%
\TextAt{2}{98.012}{796.943}{Global}%
\TextAt{2}{133.248}{796.943}{Class}%
\TextAt{2}{161.841}{796.943}{Field}%
\TextAt{2}{189.572}{796.943}{Theory}%
\Vector{q 0 G 0 g 0.299 w 7.226 18.729 m 460.769 18.729 l S Q}%
\TextAt{8}{79.226}{761.124}{where}%
\TextAt{14}{112.747}{761.124}{p}%
\TextAt{8}{122.339}{761.124}{is}%
\TextAt{8}{133.815}{761.124}{a}%
\TextAt{8}{143.282}{761.124}{prime}%
\TextAt{8}{176.157}{761.124}{ideal}%
\TextAt{8}{203.832}{761.124}{of}%
\TextAt{13}{216.874}{761.124}{K}%
\TextAt{8}{231.3}{761.124}{not}%
\TextAt{8}{251.826}{761.124}{dividing}%
\TextAt{13}{296.74}{761.124}{n}%
\TextAt{8}{303.728}{761.124}{,}%
\TextAt{13}{310.646}{761.124}{a}%
\NativeMathAt{320.112}{761.124}{11.955}{\in}%
\TextAt{13}{331.403}{761.124}{U}%
\TextAt{27}{339.353}{759.33}{p}%
\TextAt{8}{343.84}{761.124}{,}%
\TextAt{8}{350.759}{761.124}{and}%
\TextAt{13}{373.233}{761.124}{\Character{25}}%
\TextAt{8}{383.911}{761.124}{is}%
\TextAt{8}{395.387}{761.124}{a}%
\TextAt{8}{404.854}{761.124}{prime}%
\TextAt{8}{437.729}{761.124}{element}%
\TextAt{8}{480.683}{761.124}{of}%
\TextAt{13}{493.724}{761.124}{K}%
\TextAt{27}{503.683}{759.33}{p}%
\TextAt{8}{508.17}{761.124}{.}%
\TextAt{8}{516.522}{761.124}{We}%
\TextAt{8}{79.226}{744.078}{have}%
\TextAt{8}{106.354}{744.078}{seen}%
\TextAt{8}{131.905}{744.078}{that}%
\TextAt{8}{157.407}{744.078}{this}%
\TextAt{8}{180.363}{744.078}{definition}%
\TextAt{8}{232.525}{744.078}{does}%
\TextAt{8}{259.056}{744.078}{not}%
\TextAt{8}{280.003}{744.078}{depend}%
\TextAt{8}{320.773}{744.078}{on}%
\TextAt{8}{337.165}{744.078}{the}%
\TextAt{8}{357.455}{744.078}{choice}%
\TextAt{8}{392.378}{744.078}{of}%
\TextAt{8}{405.842}{744.078}{the}%
\TextAt{8}{426.132}{744.078}{prime}%
\TextAt{8}{459.429}{744.078}{element}%
\TextAt{13}{502.804}{744.078}{\Character{25}}%
\TextAt{8}{513.904}{744.078}{and}%
\TextAt{8}{79.226}{727.032}{that}%
\TextAt{8}{104.595}{727.032}{one}%
\TextAt{8}{126.055}{727.032}{has}%
\TextAt{20}{219.277}{725.438}{\Character{18}}%
\TextAt{13}{229.272}{716.669}{a}%
\Vector{q 0 G 0 g 0.478 w 157.272 -62.43 m 163.417 -62.43 l S Q}%
\TextAt{14}{229.353}{700.38}{p}%
\TextAt{20}{236.613}{725.438}{\Character{19}}%
\TextAt{9}{248.734}{708.581}{\Character{61}}%
\TextAt{9}{261.144}{708.581}{1}%
\NativeArrowAt{282.289}{708.581}{11.955}{21.926}{\Longleftrightarrow}%
\TextAt{13}{319.482}{708.581}{a}%
\NativeMathAt{328.948}{708.581}{11.955}{\equiv}%
\TextAt{13}{341.568}{708.581}{\Character{11}}%
\TextAt{21}{349.089}{713.517}{n}%
\TextAt{8}{358.046}{708.581}{mod}%
\TextAt{14}{383.484}{708.581}{p}%
\TextAt{13}{389.467}{708.581}{\Character{59}}%
\TextAt{8}{79.226}{682.758}{and}%
\TextAt{8}{101.989}{682.758}{more}%
\TextAt{8}{131.255}{682.758}{generally}%
\TextAt{20}{209.44}{678.833}{\Character{18}}%
\TextAt{13}{219.436}{670.064}{a}%
\Vector{q 0 G 0 g 0.478 w 147.436 -109.035 m 153.581 -109.035 l S Q}%
\TextAt{14}{219.517}{653.775}{p}%
\TextAt{20}{226.776}{678.833}{\Character{19}}%
\NativeMathAt{238.898}{661.976}{11.955}{\equiv}%
\TextAt{13}{251.517}{661.976}{a}%
\TextAt{17}{257.662}{666.912}{\Character{40}}%
\TextAt{21}{260.955}{666.912}{q}%
\TextAt{3}{265.022}{666.912}{\Character{0}}%
\TextAt{17}{271.608}{666.912}{1\Character{41}\Character{47}}%
\TextAt{21}{283.37}{666.912}{n}%
\TextAt{8}{292.327}{661.976}{mod}%
\TextAt{14}{317.765}{661.976}{p}%
\TextAt{13}{323.748}{661.976}{\Character{59}}%
\TextAt{13}{352.883}{661.976}{q}%
\TextAt{9}{361.843}{661.976}{\Character{61}}%
\TextAt{14}{374.268}{661.976}{N}%
\TextAt{9}{384.216}{661.976}{\Character{40}}%
\TextAt{14}{388.769}{661.976}{p}%
\TextAt{9}{394.752}{661.976}{\Character{41}}%
\TextAt{13}{399.304}{661.976}{\Character{58}}%
\TextAt{12}{79.226}{627.387}{Definition}%
\TextAt{12}{141.892}{627.387}{(8.2).}%
\TextAt{8}{177.1}{627.387}{For}%
\TextAt{8}{197.962}{627.387}{every}%
\TextAt{8}{228.747}{627.387}{ideal}%
\TextAt{14}{256.603}{627.387}{b}%
\TextAt{9}{266.059}{627.387}{\Character{61}}%
\TextAt{20}{278.485}{636.353}{\Character{81}}%
\TextAt{27}{289.776}{623.9}{p}%
\TextAt{25}{293.764}{623.9}{\Character{45}}%
\TextAt{21}{295.978}{623.9}{n}%
\TextAt{14}{303.607}{627.387}{p}%
\TextAt{21}{309.59}{631.725}{\Character{23}}%
\TextAt{28}{313.753}{630.729}{p}%
\TextAt{8}{322.124}{627.387}{prime}%
\TextAt{8}{355.181}{627.387}{to}%
\TextAt{13}{369.384}{627.387}{n}%
\TextAt{8}{376.371}{627.387}{,}%
\TextAt{8}{383.435}{627.387}{and}%
\TextAt{8}{406.09}{627.387}{every}%
\TextAt{8}{436.876}{627.387}{number}%
\TextAt{13}{479.365}{627.387}{a}%
\TextAt{8}{489.3}{627.387}{prime}%
\TextAt{8}{522.357}{627.387}{to}%
\TextAt{14}{79.226}{610.341}{b}%
\TextAt{8}{85.361}{610.341}{,}%
\TextAt{8}{92.51}{610.341}{we}%
\TextAt{8}{109.738}{610.341}{define}%
\TextAt{8}{143.547}{610.341}{the}%
\TextAt{13}{163.703}{610.341}{n}%
\TextAt{12}{170.691}{610.341}{-th}%
\TextAt{12}{192.366}{610.341}{power}%
\TextAt{12}{232.033}{610.341}{residue}%
\TextAt{12}{278.275}{610.341}{symbol}%
\TextAt{8}{323.741}{610.341}{by}%
\TextAt{20}{258.02}{593.651}{\Character{16}}%
\TextAt{13}{266.356}{588.469}{a}%
\Vector{q 0 G 0 g 0.478 w 194.356 -190.63 m 200.501 -190.63 l S Q}%
\TextAt{14}{266.361}{572.18}{b}%
\TextAt{20}{273.696}{593.651}{\Character{17}}%
\TextAt{9}{284.157}{580.381}{\Character{61}}%
\TextAt{20}{296.582}{591.738}{\Character{89}}%
\TextAt{27}{298.55}{565.906}{p}%
\TextAt{25}{302.539}{565.906}{\Character{45}}%
\TextAt{21}{304.753}{565.906}{n}%
\TextAt{20}{313.851}{597.238}{\Character{18}}%
\TextAt{13}{323.847}{588.469}{a}%
\Vector{q 0 G 0 g 0.478 w 251.847 -190.63 m 257.992 -190.63 l S Q}%
\TextAt{14}{323.928}{572.18}{p}%
\TextAt{20}{331.187}{597.238}{\Character{19}}%
\TextAt{21}{339.988}{594.561}{\Character{23}}%
\TextAt{28}{344.15}{593.565}{p}%
\TextAt{13}{350.724}{580.381}{\Character{58}}%
\TextAt{8}{79.226}{538.773}{Here}%
\TextAt{20}{106.855}{552.043}{\Character{16}}%
\TextAt{21}{115.19}{543.48}{a}%
\Vector{q 0 G 0 g 0.478 w 43.19 -232.239 m 47.688 -232.239 l S Q}%
\TextAt{27}{115.445}{534.65}{p}%
\TextAt{20}{120.883}{552.043}{\Character{17}}%
\TextAt{21}{128.023}{549.366}{\Character{23}}%
\TextAt{28}{132.186}{548.37}{p}%
\TextAt{9}{140.088}{538.773}{\Character{61}}%
\TextAt{9}{152.498}{538.773}{1}%
\TextAt{8}{162.264}{538.773}{when}%
\TextAt{13}{192.821}{538.773}{\Character{23}}%
\TextAt{27}{198.605}{536.979}{p}%
\TextAt{9}{206.413}{538.773}{\Character{61}}%
\TextAt{9}{218.823}{538.773}{0}%
\TextAt{8}{224.691}{538.773}{.}%
\TextAt{8}{97.159}{511.746}{The}%
\TextAt{8}{121.619}{511.746}{power}%
\TextAt{8}{156.17}{511.746}{residue}%
\TextAt{8}{196.304}{511.746}{symbol}%
\TextAt{20}{236.772}{521.43}{\Character{0}}%
\TextAt{21}{243.447}{516.453}{a}%
\Vector{q 0 G 0 g 0.478 w 171.447 -259.265 m 175.945 -259.265 l S Q}%
\TextAt{27}{243.651}{507.624}{b}%
\TextAt{20}{249.141}{521.43}{\Character{1}}%
\TextAt{8}{258.924}{511.746}{is}%
\TextAt{8}{271.095}{511.746}{obviously}%
\TextAt{8}{323.279}{511.746}{multiplicative}%
\TextAt{8}{397.509}{511.746}{in}%
\TextAt{8}{411.569}{511.746}{both}%
\TextAt{8}{439.616}{511.746}{arguments.}%
\TextAt{8}{502.365}{511.746}{If}%
\TextAt{14}{514.464}{511.746}{b}%
\TextAt{8}{524.903}{511.746}{is}%
\TextAt{8}{79.226}{494.7}{a}%
\TextAt{8}{88.984}{494.7}{principal}%
\TextAt{8}{137.763}{494.7}{ideal}%
\TextAt{9}{165.728}{494.7}{\Character{40}}%
\TextAt{13}{170.281}{494.7}{b}%
\TextAt{9}{175.258}{494.7}{\Character{41}}%
\TextAt{8}{179.81}{494.7}{,}%
\TextAt{8}{186.962}{494.7}{we}%
\TextAt{8}{204.191}{494.7}{write}%
\TextAt{8}{234.105}{494.7}{for}%
\TextAt{8}{251.992}{494.7}{short}%
\TextAt{20}{281.978}{504.384}{\Character{0}}%
\TextAt{21}{288.653}{499.407}{a}%
\Vector{q 0 G 0 g 0.478 w 216.653 -276.311 m 221.151 -276.311 l S Q}%
\TextAt{27}{288.857}{490.578}{b}%
\TextAt{20}{294.347}{504.384}{\Character{1}}%
\TextAt{9}{303.151}{494.7}{\Character{61}}%
\TextAt{20}{315.58}{504.384}{\Character{0}}%
\TextAt{21}{322.255}{499.407}{a}%
\Vector{q 0 G 0 g 0.478 w 250.255 -276.311 m 254.753 -276.311 l S Q}%
\TextAt{21}{322.692}{490.578}{b}%
\TextAt{20}{327.948}{504.384}{\Character{1}}%
\TextAt{8}{333.428}{494.7}{.}%
\TextAt{8}{341.883}{494.7}{We}%
\TextAt{8}{362.029}{494.7}{now}%
\TextAt{8}{386.42}{494.7}{prove}%
\TextAt{8}{417.972}{494.7}{the}%
\TextAt{8}{438.13}{494.7}{general}%
\TextAt{8}{478.457}{494.7}{reciprocity}%
\TextAt{8}{79.226}{477.654}{law}%
\TextAt{8}{100.363}{477.654}{for}%
\TextAt{8}{118.248}{477.654}{the}%
\TextAt{13}{138.404}{477.654}{n}%
\TextAt{8}{145.392}{477.654}{-th}%
\TextAt{8}{164.245}{477.654}{power}%
\TextAt{8}{198.389}{477.654}{residues.}%
\TextAt{12}{79.226}{452.511}{Theorem}%
\TextAt{12}{136.24}{452.511}{(8.3).}%
\TextAt{8}{171.556}{452.511}{If}%
\TextAt{13}{183.248}{452.511}{a\Character{59}}%
\TextAt{13}{194.617}{452.511}{b}%
\NativeMathAt{202.935}{452.511}{11.955}{\in}%
\TextAt{13}{214.226}{452.511}{K}%
\TextAt{3}{225.044}{456.85}{\Character{3}}%
\TextAt{8}{233.673}{452.511}{are}%
\TextAt{8}{253.184}{452.511}{prime}%
\TextAt{8}{286.348}{452.511}{to}%
\TextAt{8}{300.658}{452.511}{each}%
\TextAt{8}{326.996}{452.511}{other}%
\TextAt{8}{357.565}{452.511}{and}%
\TextAt{8}{380.328}{452.511}{to}%
\TextAt{13}{394.638}{452.511}{n}%
\TextAt{8}{401.625}{452.511}{,}%
\TextAt{8}{408.775}{452.511}{then}%
\TextAt{20}{235.222}{433.555}{\Character{16}}%
\TextAt{13}{243.557}{428.373}{a}%
\Vector{q 0 G 0 g 0.478 w 171.557 -350.726 m 177.702 -350.726 l S Q}%
\TextAt{13}{244.141}{412.084}{b}%
\TextAt{20}{250.898}{433.555}{\Character{17}}%
\TextAt{20}{260.03}{437.142}{\Character{18}}%
\TextAt{13}{270.61}{428.373}{b}%
\Vector{q 0 G 0 g 0.478 w 198.026 -350.726 m 204.171 -350.726 l S Q}%
\TextAt{13}{270.026}{412.084}{a}%
\TextAt{20}{277.366}{437.142}{\Character{19}}%
\TextAt{3}{286.167}{434.465}{\Character{0}}%
\TextAt{17}{292.753}{434.465}{1}%
\TextAt{9}{300.806}{420.285}{\Character{61}}%
\TextAt{20}{315.568}{431.642}{\Character{89}}%
\TextAt{27}{313.232}{405.739}{p}%
\TextAt{3}{317.22}{405.739}{\Character{106}}%
\TextAt{21}{319.573}{405.739}{n}%
\NativeMathAt{324.711}{405.739}{7.97}{\infty}%
\TextAt{20}{335.172}{437.142}{\Character{18}}%
\TextAt{13}{345.168}{428.373}{a\Character{59}}%
\TextAt{13}{356.537}{428.373}{b}%
\Vector{q 0 G 0 g 0.478 w 273.168 -350.726 m 289.534 -350.726 l S Q}%
\TextAt{14}{350.359}{412.084}{p}%
\TextAt{20}{362.729}{437.142}{\Character{19}}%
\TextAt{13}{373.522}{420.285}{\Character{58}}%
\TextAt{19}{79.226}{368.031}{Proof.}%
\TextAt{8}{116}{368.031}{If}%
\TextAt{14}{127.692}{368.031}{p}%
\TextAt{8}{137.573}{368.031}{is}%
\TextAt{8}{149.337}{368.031}{prime}%
\TextAt{8}{182.5}{368.031}{to}%
\TextAt{13}{196.811}{368.031}{bn}%
\NativeMathAt{208.775}{368.031}{11.955}{\infty}%
\TextAt{8}{220.731}{368.031}{,}%
\TextAt{8}{227.88}{368.031}{then}%
\TextAt{8}{254.54}{368.031}{we}%
\TextAt{8}{271.767}{368.031}{have}%
\TextAt{20}{218.033}{352.383}{\Character{18}}%
\TextAt{13}{228.532}{343.613}{b}%
\Vector{q 0 G 0 g 0.478 w 156.029 -435.486 m 162.012 -435.486 l S Q}%
\TextAt{14}{228.029}{327.325}{p}%
\TextAt{20}{235.207}{352.383}{\Character{19}}%
\TextAt{21}{244.008}{349.706}{v}%
\TextAt{28}{248.131}{348.71}{p}%
\TextAt{17}{252.214}{349.706}{\Character{40}}%
\TextAt{21}{255.507}{349.706}{a}%
\TextAt{17}{260.005}{349.706}{\Character{41}}%
\TextAt{9}{267.118}{335.526}{\Character{61}}%
\TextAt{20}{279.543}{352.383}{\Character{18}}%
\TextAt{13}{289.539}{343.613}{\Character{25}}%
\TextAt{13}{296.593}{343.613}{\Character{59}}%
\TextAt{13}{301.829}{343.613}{b}%
\Vector{q 0 G 0 g 0.478 w 217.539 -435.486 m 234.829 -435.486 l S Q}%
\TextAt{14}{295.193}{327.325}{p}%
\TextAt{20}{308.025}{352.383}{\Character{19}}%
\TextAt{21}{316.825}{349.706}{v}%
\TextAt{28}{320.949}{348.71}{p}%
\TextAt{17}{325.032}{349.706}{\Character{40}}%
\TextAt{21}{328.325}{349.706}{a}%
\TextAt{17}{332.823}{349.706}{\Character{41}}%
\TextAt{9}{339.935}{335.526}{\Character{61}}%
\TextAt{20}{352.36}{352.383}{\Character{18}}%
\TextAt{13}{362.356}{343.613}{a\Character{59}}%
\TextAt{13}{373.725}{343.613}{b}%
\Vector{q 0 G 0 g 0.478 w 290.356 -435.486 m 306.723 -435.486 l S Q}%
\TextAt{14}{367.548}{327.325}{p}%
\TextAt{20}{379.918}{352.383}{\Character{19}}%
\TextAt{13}{390.711}{335.526}{\Character{59}}%
\TextAt{8}{79.226}{300.51}{where}%
\TextAt{13}{113.973}{300.51}{\Character{25}}%
\TextAt{8}{125.877}{300.51}{is}%
\TextAt{8}{138.579}{300.51}{a}%
\TextAt{8}{149.272}{300.51}{prime}%
\TextAt{8}{183.374}{300.51}{element}%
\TextAt{8}{227.553}{300.51}{of}%
\TextAt{13}{241.821}{300.51}{K}%
\TextAt{27}{251.78}{298.717}{p}%
\TextAt{8}{256.267}{300.51}{.}%
\TextAt{8}{267.528}{300.51}{For}%
\TextAt{8}{289.435}{300.51}{if}%
\TextAt{8}{301.097}{300.51}{we}%
\TextAt{8}{319.262}{300.51}{put}%
\TextAt{13}{341.659}{300.51}{a}%
\TextAt{9}{352.723}{300.51}{\Character{61}}%
\TextAt{13}{366.746}{300.51}{u\Character{25}}%
\TextAt{21}{380.478}{304.849}{v}%
\TextAt{28}{384.602}{303.853}{p}%
\TextAt{17}{388.684}{304.849}{\Character{40}}%
\TextAt{21}{391.978}{304.849}{a}%
\TextAt{17}{396.476}{304.849}{\Character{41}}%
\TextAt{8}{400.267}{300.51}{,}%
\TextAt{8}{408.588}{300.51}{then}%
\TextAt{20}{436.186}{313.781}{\Character{16}}%
\TextAt{21}{444.522}{305.766}{u\Character{59}b}%
\Vector{q 0 G 0 g 0.478 w 372.522 -470.501 m 383.399 -470.501 l S Q}%
\TextAt{27}{447.966}{296.388}{p}%
\TextAt{20}{456.595}{313.781}{\Character{17}}%
\TextAt{9}{468.654}{300.51}{\Character{61}}%
\TextAt{9}{482.666}{300.51}{1}%
\TextAt{8}{493.365}{300.51}{because}%
\TextAt{13}{79.226}{283.464}{u\Character{59}}%
\TextAt{13}{91.11}{283.464}{b}%
\NativeMathAt{99.431}{283.464}{11.955}{\in}%
\TextAt{13}{110.722}{283.464}{U}%
\TextAt{27}{118.672}{281.671}{p}%
\TextAt{8}{123.159}{283.464}{.}%
\TextAt{8}{131.607}{283.464}{For}%
\TextAt{8}{152.576}{283.464}{the}%
\TextAt{8}{172.733}{283.464}{same}%
\TextAt{8}{202.059}{283.464}{reason,}%
\TextAt{8}{241.798}{283.464}{we}%
\TextAt{8}{259.025}{283.464}{find}%
\TextAt{20}{216.764}{269.777}{\Character{18}}%
\TextAt{13}{226.76}{261.008}{a\Character{59}}%
\TextAt{13}{238.129}{261.008}{b}%
\Vector{q 0 G 0 g 0.478 w 154.76 -518.091 m 171.126 -518.091 l S Q}%
\TextAt{14}{231.951}{244.719}{p}%
\TextAt{20}{244.321}{269.777}{\Character{19}}%
\TextAt{9}{256.442}{252.92}{\Character{61}}%
\TextAt{9}{268.851}{252.92}{1}%
\TextAt{8}{286.676}{252.92}{for}%
\TextAt{14}{304.561}{252.92}{p}%
\TextAt{8}{314.441}{252.92}{prime}%
\TextAt{8}{347.605}{252.92}{to}%
\TextAt{13}{361.915}{252.92}{abn}%
\NativeMathAt{380.025}{252.92}{11.955}{\infty}%
\TextAt{13}{391.98}{252.92}{\Character{58}}%
\TextAt{8}{79.226}{222.687}{(8.1)}%
\TextAt{8}{107.201}{222.687}{then}%
\TextAt{8}{133.861}{222.687}{gives}%
\TextAt{20}{160.696}{202.89}{\Character{16}}%
\TextAt{13}{169.032}{197.707}{a}%
\Vector{q 0 G 0 g 0.478 w 97.032 -581.392 m 103.177 -581.392 l S Q}%
\TextAt{13}{169.616}{181.419}{b}%
\TextAt{20}{176.372}{202.89}{\Character{17}}%
\TextAt{20}{185.504}{206.476}{\Character{18}}%
\TextAt{13}{196.084}{197.707}{b}%
\Vector{q 0 G 0 g 0.478 w 123.5 -581.392 m 129.645 -581.392 l S Q}%
\TextAt{13}{195.5}{181.419}{a}%
\TextAt{20}{202.841}{206.476}{\Character{19}}%
\TextAt{3}{211.641}{203.8}{\Character{0}}%
\TextAt{17}{218.228}{203.8}{1}%
\TextAt{9}{226.281}{189.619}{\Character{61}}%
\TextAt{20}{239.343}{200.977}{\Character{89}}%
\TextAt{27}{238.706}{175.074}{p}%
\TextAt{3}{242.695}{175.074}{\Character{106}}%
\TextAt{17}{245.047}{175.074}{\Character{40}}%
\TextAt{21}{248.341}{175.074}{b}%
\TextAt{17}{251.963}{175.074}{\Character{41}}%
\TextAt{20}{257.249}{206.476}{\Character{18}}%
\TextAt{13}{267.245}{197.707}{a}%
\Vector{q 0 G 0 g 0.478 w 195.245 -581.392 m 201.39 -581.392 l S Q}%
\TextAt{14}{267.326}{181.419}{p}%
\TextAt{20}{274.585}{206.476}{\Character{19}}%
\TextAt{21}{283.386}{203.8}{v}%
\TextAt{28}{287.509}{202.803}{p}%
\TextAt{17}{291.592}{203.8}{\Character{40}}%
\TextAt{21}{294.885}{203.8}{b}%
\TextAt{17}{298.508}{203.8}{\Character{41}}%
\TextAt{20}{305.366}{200.977}{\Character{89}}%
\TextAt{27}{304.292}{175.074}{p}%
\TextAt{3}{308.28}{175.074}{\Character{106}}%
\TextAt{17}{310.633}{175.074}{\Character{40}}%
\TextAt{21}{313.926}{175.074}{a}%
\TextAt{17}{318.424}{175.074}{\Character{41}}%
\TextAt{20}{323.71}{206.476}{\Character{18}}%
\TextAt{13}{334.209}{197.707}{b}%
\Vector{q 0 G 0 g 0.478 w 261.706 -581.392 m 267.689 -581.392 l S Q}%
\TextAt{14}{333.706}{181.419}{p}%
\TextAt{20}{340.884}{206.476}{\Character{19}}%
\TextAt{3}{349.684}{203.8}{\Character{0}}%
\TextAt{21}{356.271}{203.8}{v}%
\TextAt{28}{360.394}{202.803}{p}%
\TextAt{17}{364.477}{203.8}{\Character{40}}%
\TextAt{21}{367.77}{203.8}{a}%
\TextAt{17}{372.268}{203.8}{\Character{41}}%
\TextAt{9}{226.281}{148.584}{\Character{61}}%
\TextAt{20}{239.343}{159.942}{\Character{89}}%
\TextAt{27}{238.706}{134.039}{p}%
\TextAt{3}{242.695}{134.039}{\Character{106}}%
\TextAt{17}{245.047}{134.039}{\Character{40}}%
\TextAt{21}{248.341}{134.039}{b}%
\TextAt{17}{251.963}{134.039}{\Character{41}}%
\TextAt{20}{257.249}{165.441}{\Character{18}}%
\TextAt{13}{267.245}{156.672}{b\Character{59}}%
\TextAt{13}{277.443}{156.672}{a}%
\Vector{q 0 G 0 g 0.478 w 195.245 -622.427 m 211.611 -622.427 l S Q}%
\TextAt{14}{272.436}{140.384}{p}%
\TextAt{20}{284.807}{165.441}{\Character{19}}%
\TextAt{20}{296.674}{159.942}{\Character{89}}%
\TextAt{27}{295.599}{134.039}{p}%
\TextAt{3}{299.588}{134.039}{\Character{106}}%
\TextAt{17}{301.94}{134.039}{\Character{40}}%
\TextAt{21}{305.234}{134.039}{a}%
\TextAt{17}{309.732}{134.039}{\Character{41}}%
\TextAt{20}{315.017}{165.441}{\Character{18}}%
\TextAt{13}{325.013}{156.672}{a\Character{59}}%
\TextAt{13}{336.382}{156.672}{b}%
\Vector{q 0 G 0 g 0.478 w 253.013 -622.427 m 269.38 -622.427 l S Q}%
\TextAt{14}{330.205}{140.384}{p}%
\TextAt{20}{342.575}{165.441}{\Character{19}}%
\TextAt{3}{351.375}{162.764}{\Character{0}}%
\TextAt{17}{357.962}{162.764}{1}%
\TextAt{9}{226.281}{109.53}{\Character{61}}%
\TextAt{20}{241.592}{120.888}{\Character{89}}%
\TextAt{27}{238.706}{94.985}{p}%
\TextAt{3}{242.695}{94.985}{\Character{106}}%
\TextAt{17}{245.047}{94.985}{\Character{40}}%
\TextAt{21}{248.341}{94.985}{ab}%
\TextAt{17}{256.461}{94.985}{\Character{41}}%
\TextAt{20}{261.747}{126.387}{\Character{18}}%
\TextAt{13}{271.743}{117.618}{b\Character{59}}%
\TextAt{13}{281.941}{117.618}{a}%
\Vector{q 0 G 0 g 0.478 w 199.743 -661.481 m 216.109 -661.481 l S Q}%
\TextAt{14}{276.934}{101.33}{p}%
\TextAt{20}{289.305}{126.387}{\Character{19}}%
\TextAt{9}{301.426}{109.53}{\Character{61}}%
\TextAt{20}{316.118}{120.888}{\Character{89}}%
\TextAt{27}{313.851}{95.056}{p}%
\TextAt{25}{317.84}{95.056}{\Character{45}}%
\TextAt{21}{320.054}{95.056}{n}%
\NativeMathAt{325.192}{95.056}{7.97}{\infty}%
\TextAt{20}{335.653}{126.387}{\Character{18}}%
\TextAt{13}{345.649}{117.618}{b\Character{59}}%
\TextAt{13}{355.847}{117.618}{a}%
\Vector{q 0 G 0 g 0.478 w 273.649 -661.481 m 290.015 -661.481 l S Q}%
\TextAt{14}{350.84}{101.33}{p}%
\TextAt{20}{363.211}{126.387}{\Character{19}}%
\TextAt{9}{375.332}{109.53}{\Character{61}}%
\TextAt{20}{390.093}{120.888}{\Character{89}}%
\TextAt{27}{387.757}{94.985}{p}%
\TextAt{3}{391.746}{94.985}{\Character{106}}%
\TextAt{21}{394.098}{94.985}{n}%
\NativeMathAt{399.236}{94.985}{7.97}{\infty}%
\TextAt{20}{409.697}{126.387}{\Character{18}}%
\TextAt{13}{419.693}{117.618}{a\Character{59}}%
\TextAt{13}{431.062}{117.618}{b}%
\Vector{q 0 G 0 g 0.478 w 347.693 -661.481 m 364.059 -661.481 l S Q}%
\TextAt{14}{424.885}{101.33}{p}%
\TextAt{20}{437.255}{126.387}{\Character{19}}%
\TextAt{13}{448.048}{109.53}{\Character{58}}%
\TextAt{8}{79.226}{72.921}{Here}%
\TextAt{14}{106.855}{72.921}{p}%
\TextAt{16}{116.158}{72.921}{\Character{106}}%
\TextAt{9}{122.8}{72.921}{\Character{40}}%
\TextAt{13}{127.353}{72.921}{b}%
\TextAt{9}{132.33}{72.921}{\Character{41}}%
\TextAt{8}{140.779}{72.921}{means}%
\TextAt{8}{176.609}{72.921}{that}%
\TextAt{14}{201.978}{72.921}{p}%
\TextAt{8}{211.858}{72.921}{occurs}%
\TextAt{8}{248.011}{72.921}{in}%
\TextAt{8}{261.664}{72.921}{the}%
\TextAt{8}{281.82}{72.921}{prime}%
\TextAt{8}{314.984}{72.921}{decomposition}%
\TextAt{8}{392.118}{72.921}{of}%
\TextAt{9}{405.448}{72.921}{\Character{40}}%
\TextAt{13}{410.001}{72.921}{b}%
\TextAt{9}{414.978}{72.921}{\Character{41}}%
\TextAt{8}{419.53}{72.921}{.}%
\Vector{q 0 G 0 g 0.399 w 452.334 -701.079 m 452.334 -693.009 l S Q}%
\Vector{q 0 G 0 g 0.399 w 452.534 -693.208 m 459.707 -693.208 l S Q}%
\Vector{q 0 G 0 g 0.399 w 452.534 -700.879 m 459.707 -700.879 l S Q}%
\Vector{q 0 G 0 g 0.399 w 459.906 -701.079 m 459.906 -693.009 l S Q}%
\TextAt{8}{297.211}{43.033}{379}%
\Destination{78.23}{809.815}{nav.page.379}%
\Destination{79.36}{776.155}{nav.print.415}%
\Destination{79.36}{641.355}{nav.number.6.8.2}%
\Destination{79.36}{466.555}{nav.number.6.8.3}%
\Hotlink{78.726}{220.455}{25.078}{11.047}{nav.number.6.8.1}%
}

% Page 392 (body page 380)
\Page{%
\TextAt{2}{79.226}{796.943}{VI.}%
\TextAt{2}{98.012}{796.943}{Global}%
\TextAt{2}{133.248}{796.943}{Class}%
\TextAt{2}{161.841}{796.943}{Field}%
\TextAt{2}{189.572}{796.943}{Theory}%
\Vector{q 0 G 0 g 0.299 w 7.226 18.729 m 460.769 18.729 l S Q}%
\TextAt{8}{97.159}{761.124}{Gauss’s}%
\TextAt{8}{140.593}{761.124}{reciprocity}%
\TextAt{8}{199.701}{761.124}{law,}%
\TextAt{8}{225.211}{761.124}{for}%
\TextAt{8}{243.994}{761.124}{which}%
\TextAt{8}{278.378}{761.124}{we}%
\TextAt{8}{296.503}{761.124}{gave}%
\TextAt{8}{323.75}{761.124}{an}%
\TextAt{8}{340.907}{761.124}{elementary}%
\TextAt{8}{401.641}{761.124}{proof}%
\TextAt{8}{433.108}{761.124}{using}%
\TextAt{8}{464.635}{761.124}{the}%
\TextAt{8}{485.689}{761.124}{theory}%
\TextAt{8}{523.337}{761.124}{of}%
\TextAt{8}{79.226}{744.078}{Gauss}%
\TextAt{8}{113.223}{744.078}{sums}%
\TextAt{8}{141.936}{744.078}{in}%
\TextAt{8}{154.916}{744.078}{chap.}%
\TextAt{8}{185.135}{744.078}{I,}%
\TextAt{8}{195.831}{744.078}{(8.6),}%
\TextAt{8}{226.52}{744.078}{in}%
\TextAt{8}{239.5}{744.078}{the}%
\TextAt{8}{258.983}{744.078}{case}%
\TextAt{8}{283.081}{744.078}{of}%
\TextAt{8}{295.738}{744.078}{two}%
\TextAt{8}{317.182}{744.078}{odd}%
\TextAt{8}{339.595}{744.078}{prime}%
\TextAt{8}{372.086}{744.078}{numbers}%
\TextAt{13}{418.624}{744.078}{p\Character{59}}%
\TextAt{13}{429.718}{744.078}{l}%
\TextAt{8}{433.493}{744.078}{,}%
\TextAt{8}{440.103}{744.078}{is}%
\TextAt{8}{451.194}{744.078}{contained}%
\TextAt{8}{503.531}{744.078}{in}%
\TextAt{8}{516.51}{744.078}{the}%
\TextAt{8}{79.226}{727.032}{general}%
\TextAt{8}{120.022}{727.032}{reciprocity}%
\TextAt{8}{178.702}{727.032}{law}%
\TextAt{8}{200.31}{727.032}{(8.3)}%
\TextAt{8}{228.756}{727.032}{as}%
\TextAt{8}{243.597}{727.032}{a}%
\TextAt{8}{253.823}{727.032}{special}%
\TextAt{8}{292.395}{727.032}{case.}%
\TextAt{8}{323.13}{727.032}{For}%
\TextAt{8}{344.57}{727.032}{if}%
\TextAt{8}{355.764}{727.032}{we}%
\TextAt{8}{373.463}{727.032}{substitute,}%
\TextAt{8}{432.058}{727.032}{in}%
\TextAt{8}{446.181}{727.032}{the}%
\TextAt{8}{466.809}{727.032}{case}%
\TextAt{13}{492.05}{727.032}{K}%
\TextAt{9}{506.991}{727.032}{\Character{61}}%
\TextAt{18}{520.219}{727.032}{Q}%
\TextAt{8}{529.518}{727.032}{,}%
\TextAt{13}{79.226}{709.986}{n}%
\TextAt{9}{90.348}{709.986}{\Character{61}}%
\TextAt{9}{103.571}{709.986}{2}%
\TextAt{8}{109.439}{709.986}{,}%
\TextAt{8}{117.185}{709.986}{into}%
\TextAt{8}{141.405}{709.986}{formula}%
\TextAt{8}{184.813}{709.986}{(8.3)}%
\TextAt{8}{213.265}{709.986}{the}%
\TextAt{8}{233.899}{709.986}{explicit}%
\TextAt{8}{275.669}{709.986}{description}%
\TextAt{8}{336.042}{709.986}{(chap.}%
\TextAt{8}{371.966}{709.986}{V,}%
\TextAt{8}{388.367}{709.986}{(3.6))}%
\TextAt{8}{421.374}{709.986}{of}%
\TextAt{8}{435.182}{709.986}{the}%
\TextAt{8}{455.815}{709.986}{Hilbert}%
\TextAt{8}{496.605}{709.986}{symbol}%
\TextAt{20}{79.226}{706.21}{\Character{16}}%
\TextAt{21}{87.562}{698.196}{a\Character{59}b}%
\Vector{q 0 G 0 g 0.478 w 15.562 -78.071 m 26.034 -78.071 l S Q}%
\TextAt{21}{90.667}{688.817}{p}%
\TextAt{20}{99.23}{706.21}{\Character{17}}%
\TextAt{8}{110.229}{692.94}{for}%
\TextAt{13}{128.075}{692.94}{p}%
\TextAt{9}{137.271}{692.94}{\Character{61}}%
\TextAt{9}{149.681}{692.94}{2}%
\TextAt{8}{159.408}{692.94}{and}%
\TextAt{13}{182.132}{692.94}{p}%
\TextAt{9}{191.328}{692.94}{\Character{61}}%
\NativeMathAt{203.753}{692.94}{11.955}{\infty}%
\TextAt{8}{215.708}{692.94}{,}%
\TextAt{8}{222.826}{692.94}{we}%
\TextAt{8}{240.015}{692.94}{obtain}%
\TextAt{8}{276.404}{692.94}{the}%
\TextAt{8}{296.521}{692.94}{following}%
\TextAt{8}{345.917}{692.94}{theorem,}%
\TextAt{8}{394.663}{692.94}{which}%
\TextAt{8}{428.111}{692.94}{is}%
\TextAt{8}{439.836}{692.94}{more}%
\TextAt{8}{469.063}{692.94}{general}%
\TextAt{8}{509.349}{692.94}{than}%
\TextAt{8}{79.226}{675.894}{chap.}%
\TextAt{8}{110.118}{675.894}{I,}%
\TextAt{8}{121.488}{675.894}{(8.6).}%
\TextAt{12}{79.226}{650.781}{Gauss’s}%
\TextAt{12}{127.825}{650.781}{Reciprocity}%
\TextAt{12}{199.82}{650.781}{Law}%
\TextAt{12}{228.143}{650.781}{(8.4).}%
\TextAt{8}{263.346}{650.781}{Let}%
\TextAt{13}{284.202}{650.781}{K}%
\TextAt{9}{298.34}{650.781}{\Character{61}}%
\TextAt{18}{310.766}{650.781}{Q}%
\TextAt{8}{320.064}{650.781}{,}%
\TextAt{13}{327.123}{650.781}{n}%
\TextAt{9}{337.431}{650.781}{\Character{61}}%
\TextAt{9}{349.841}{650.781}{2}%
\TextAt{8}{355.71}{650.781}{,}%
\TextAt{8}{362.768}{650.781}{and}%
\TextAt{8}{385.418}{650.781}{let}%
\TextAt{13}{402.209}{650.781}{a}%
\TextAt{8}{412.138}{650.781}{and}%
\TextAt{13}{434.787}{650.781}{b}%
\TextAt{8}{443.549}{650.781}{be}%
\TextAt{8}{459.36}{650.781}{odd,}%
\TextAt{8}{485.606}{650.781}{relatively}%
\TextAt{8}{79.226}{633.735}{prime}%
\TextAt{8}{112.39}{633.735}{integers.}%
\TextAt{8}{160.254}{633.735}{Then}%
\TextAt{8}{190.812}{633.735}{one}%
\TextAt{8}{212.271}{633.735}{has}%
\TextAt{20}{203.481}{617.415}{\Character{16}}%
\TextAt{13}{211.816}{612.233}{a}%
\Vector{q 0 G 0 g 0.478 w 139.816 -166.866 m 145.961 -166.866 l S Q}%
\TextAt{13}{212.4}{595.944}{b}%
\TextAt{20}{219.157}{617.415}{\Character{17}}%
\TextAt{20}{228.289}{621.002}{\Character{18}}%
\TextAt{13}{238.869}{612.233}{b}%
\Vector{q 0 G 0 g 0.478 w 166.285 -166.866 m 172.43 -166.866 l S Q}%
\TextAt{13}{238.285}{595.944}{a}%
\TextAt{20}{245.625}{621.002}{\Character{19}}%
\TextAt{9}{257.747}{604.145}{\Character{61}}%
\TextAt{9}{270.156}{604.145}{\Character{40}}%
\TextAt{16}{274.724}{604.145}{\Character{0}}%
\TextAt{9}{284.023}{604.145}{1\Character{41}}%
\TextAt{22}{295.624}{612.442}{a}%
\TextAt{23}{299.658}{612.442}{\Character{0}}%
\TextAt{24}{305.414}{612.442}{1}%
\Vector{q 0 G 0 g 0.359 w 223.624 -162.926 m 237.067 -162.926 l S Q}%
\TextAt{24}{300.519}{605.955}{2}%
\TextAt{22}{311.458}{612.442}{b}%
\TextAt{23}{314.732}{612.442}{\Character{0}}%
\TextAt{24}{320.488}{612.442}{1}%
\Vector{q 0 G 0 g 0.359 w 239.458 -162.926 m 252.141 -162.926 l S Q}%
\TextAt{24}{315.973}{605.955}{2}%
\TextAt{9}{325.834}{604.145}{\Character{40}}%
\TextAt{16}{330.387}{604.145}{\Character{0}}%
\TextAt{9}{339.685}{604.145}{1\Character{41}}%
\TextAt{40}{351.286}{612.837}{sgn}%
\TextAt{22}{363.169}{612.837}{a}%
\TextAt{23}{367.203}{612.837}{\Character{0}}%
\TextAt{24}{372.96}{612.837}{1}%
\Vector{q 0 G 0 g 0.359 w 279.286 -162.926 m 304.613 -162.926 l S Q}%
\TextAt{24}{362.123}{605.955}{2}%
\TextAt{40}{379.004}{612.837}{sgn}%
\TextAt{22}{390.887}{612.837}{b}%
\TextAt{23}{394.161}{612.837}{\Character{0}}%
\TextAt{24}{399.917}{612.837}{1}%
\Vector{q 0 G 0 g 0.359 w 307.004 -162.926 m 331.57 -162.926 l S Q}%
\TextAt{24}{389.46}{605.955}{2}%
\TextAt{13}{405.263}{604.145}{\Character{59}}%
\TextAt{8}{79.226}{574.567}{and}%
\TextAt{8}{101.989}{574.567}{for}%
\TextAt{8}{119.874}{574.567}{positive}%
\TextAt{8}{163.187}{574.567}{odd}%
\TextAt{8}{186.273}{574.567}{integers}%
\TextAt{13}{229.586}{574.567}{b}%
\TextAt{8}{234.563}{574.567}{,}%
\TextAt{8}{241.713}{574.567}{we}%
\TextAt{8}{258.94}{574.567}{have}%
\TextAt{8}{285.935}{574.567}{the}%
\TextAt{8}{306.091}{574.567}{two}%
\TextAt{8}{328.208}{574.567}{“supplementary}%
\TextAt{8}{412.492}{574.567}{theorems”}%
\TextAt{20}{201.651}{561.834}{\Character{18}}%
\TextAt{16}{211.647}{553.065}{\Character{0}}%
\TextAt{9}{220.945}{553.065}{1}%
\Vector{q 0 G 0 g 0.478 w 139.647 -226.034 m 154.798 -226.034 l S Q}%
\TextAt{13}{216.734}{536.776}{b}%
\TextAt{20}{227.994}{561.834}{\Character{19}}%
\TextAt{9}{240.115}{544.977}{\Character{61}}%
\TextAt{9}{252.525}{544.977}{\Character{40}}%
\TextAt{16}{257.093}{544.977}{\Character{0}}%
\TextAt{9}{266.391}{544.977}{1\Character{41}}%
\TextAt{22}{277.992}{553.274}{b}%
\TextAt{23}{281.266}{553.274}{\Character{0}}%
\TextAt{24}{287.022}{553.274}{1}%
\Vector{q 0 G 0 g 0.359 w 205.992 -222.094 m 218.675 -222.094 l S Q}%
\TextAt{24}{282.507}{546.788}{2}%
\TextAt{13}{292.368}{544.977}{\Character{59}}%
\TextAt{20}{321.523}{561.834}{\Character{18}}%
\TextAt{9}{331.519}{553.065}{2}%
\Vector{q 0 G 0 g 0.478 w 259.519 -226.034 m 265.372 -226.034 l S Q}%
\TextAt{13}{331.957}{536.776}{b}%
\TextAt{20}{338.567}{561.834}{\Character{19}}%
\TextAt{9}{350.689}{544.977}{\Character{61}}%
\TextAt{9}{363.099}{544.977}{\Character{40}}%
\TextAt{16}{367.666}{544.977}{\Character{0}}%
\TextAt{9}{376.965}{544.977}{1\Character{41}}%
\TextAt{22}{388.566}{553.274}{b}%
\TextAt{24}{391.84}{555.782}{2}%
\TextAt{23}{395.991}{553.274}{\Character{0}}%
\TextAt{24}{401.747}{553.274}{1}%
\Vector{q 0 G 0 g 0.359 w 316.566 -222.094 m 333.4 -222.094 l S Q}%
\TextAt{24}{395.156}{546.788}{8}%
\TextAt{13}{407.093}{544.977}{\Character{58}}%
\TextAt{8}{97.159}{507.332}{For}%
\TextAt{8}{118.128}{507.332}{the}%
\TextAt{8}{138.285}{507.332}{last}%
\TextAt{8}{160.462}{507.332}{equation}%
\TextAt{8}{208.27}{507.332}{we}%
\TextAt{8}{225.498}{507.332}{need}%
\TextAt{8}{252.803}{507.332}{again}%
\TextAt{8}{284.03}{507.332}{the}%
\TextAt{8}{304.187}{507.332}{product}%
\TextAt{8}{348.086}{507.332}{formula:}%
\TextAt{20}{203.333}{491.776}{\Character{18}}%
\TextAt{9}{213.329}{483.007}{2}%
\Vector{q 0 G 0 g 0.478 w 141.329 -296.092 m 147.182 -296.092 l S Q}%
\TextAt{13}{213.767}{466.719}{b}%
\TextAt{20}{220.377}{491.776}{\Character{19}}%
\TextAt{9}{232.498}{474.919}{\Character{61}}%
\TextAt{20}{250.238}{486.277}{\Character{89}}%
\TextAt{21}{244.924}{460.816}{p}%
\NativeMathAt{249.187}{460.816}{7.97}{\neq}%
\TextAt{17}{255.77}{460.816}{2}%
\TextAt{21}{260.007}{460.816}{\Character{59}}%
\NativeMathAt{262.36}{460.816}{7.97}{\infty}%
\TextAt{20}{272.821}{491.776}{\Character{18}}%
\TextAt{13}{282.816}{483.007}{p\Character{59}}%
\TextAt{9}{293.936}{483.007}{2}%
\Vector{q 0 G 0 g 0.478 w 210.816 -296.092 m 227.789 -296.092 l S Q}%
\TextAt{13}{288.365}{466.719}{p}%
\TextAt{20}{300.984}{491.776}{\Character{19}}%
\TextAt{21}{309.785}{489.099}{v}%
\TextAt{22}{313.908}{488.103}{p}%
\TextAt{17}{318.247}{489.099}{\Character{40}}%
\TextAt{21}{321.54}{489.099}{b}%
\TextAt{17}{325.162}{489.099}{\Character{41}}%
\TextAt{9}{332.275}{474.919}{\Character{61}}%
\TextAt{20}{350.014}{486.277}{\Character{89}}%
\TextAt{21}{344.7}{460.816}{p}%
\NativeMathAt{348.963}{460.816}{7.97}{\neq}%
\TextAt{17}{355.546}{460.816}{2}%
\TextAt{21}{359.784}{460.816}{\Character{59}}%
\NativeMathAt{362.136}{460.816}{7.97}{\infty}%
\TextAt{20}{372.597}{491.776}{\Character{18}}%
\TextAt{13}{382.593}{483.007}{b\Character{59}}%
\TextAt{9}{392.814}{483.007}{2}%
\Vector{q 0 G 0 g 0.478 w 310.593 -296.092 m 326.667 -296.092 l S Q}%
\TextAt{13}{387.692}{466.719}{p}%
\TextAt{20}{399.862}{491.776}{\Character{19}}%
\TextAt{9}{232.498}{436.751}{\Character{61}}%
\TextAt{20}{244.924}{453.608}{\Character{18}}%
\TextAt{9}{254.92}{444.839}{2}%
\TextAt{13}{260.773}{444.839}{\Character{59}}%
\TextAt{13}{265.997}{444.839}{b}%
\Vector{q 0 G 0 g 0.478 w 182.92 -334.26 m 198.994 -334.26 l S Q}%
\TextAt{9}{260.03}{428.55}{2}%
\TextAt{20}{272.189}{453.608}{\Character{19}}%
\TextAt{20}{282.973}{453.608}{\Character{18}}%
\TextAt{9}{292.978}{444.839}{2}%
\TextAt{13}{298.831}{444.839}{\Character{59}}%
\TextAt{13}{304.055}{444.839}{b}%
\Vector{q 0 G 0 g 0.478 w 220.978 -334.26 m 237.052 -334.26 l S Q}%
\NativeMathAt{295.038}{428.55}{11.955}{\infty}%
\TextAt{20}{310.248}{453.608}{\Character{19}}%
\TextAt{9}{322.369}{436.751}{\Character{61}}%
\TextAt{9}{334.779}{436.751}{\Character{40}}%
\TextAt{16}{339.347}{436.751}{\Character{0}}%
\TextAt{9}{348.646}{436.751}{1\Character{41}}%
\TextAt{22}{360.246}{445.048}{b}%
\TextAt{24}{363.52}{447.556}{2}%
\TextAt{23}{367.671}{445.048}{\Character{0}}%
\TextAt{24}{373.427}{445.048}{1}%
\Vector{q 0 G 0 g 0.359 w 288.246 -330.32 m 305.08 -330.32 l S Q}%
\ErrataAt{366.86}{438.555}{13.5}{$\scriptscriptstyle 8$}%
\TextAt{8}{97.412}{405.427}{The}%
\TextAt{8}{120.559}{405.427}{symbol}%
\TextAt{20}{159.714}{415.111}{\Character{0}}%
\TextAt{21}{166.389}{410.135}{a}%
\Vector{q 0 G 0 g 0.478 w 94.389 -365.584 m 98.887 -365.584 l S Q}%
\TextAt{21}{166.827}{401.305}{b}%
\TextAt{20}{172.082}{415.111}{\Character{1}}%
\TextAt{8}{180.552}{405.427}{is}%
\TextAt{8}{191.41}{405.427}{called}%
\TextAt{8}{223.666}{405.427}{the}%
\TextAt{12}{242.916}{405.427}{Jacobi}%
\TextAt{12}{283.764}{405.427}{symbol}%
\TextAt{8}{325.332}{405.427}{,}%
\TextAt{8}{331.756}{405.427}{or}%
\TextAt{8}{345.159}{405.427}{also}%
\TextAt{8}{367.733}{405.427}{the}%
\TextAt{12}{386.982}{405.427}{quadratic}%
\TextAt{12}{446.002}{405.427}{residue}%
\TextAt{12}{491.201}{405.427}{symbol}%
\TextAt{8}{79.226}{388.382}{(although,}%
\TextAt{8}{135.88}{388.382}{for}%
\TextAt{13}{153.812}{388.382}{b}%
\TextAt{8}{162.734}{388.382}{not}%
\TextAt{8}{183.594}{388.382}{a}%
\TextAt{8}{193.397}{388.382}{prime}%
\TextAt{8}{226.607}{388.382}{number,}%
\TextAt{8}{272.514}{388.382}{the}%
\TextAt{8}{292.717}{388.382}{condition}%
\TextAt{8}{344.147}{388.382}{that}%
\TextAt{8}{369.563}{388.382}{the}%
\TextAt{8}{389.766}{388.382}{symbol}%
\TextAt{20}{429.875}{398.065}{\Character{0}}%
\TextAt{21}{436.55}{393.089}{a}%
\Vector{q 0 G 0 g 0.478 w 364.55 -382.63 m 369.048 -382.63 l S Q}%
\TextAt{21}{436.988}{384.259}{b}%
\TextAt{20}{442.243}{398.065}{\Character{1}}%
\TextAt{9}{451.123}{388.382}{\Character{61}}%
\TextAt{9}{463.616}{388.382}{1}%
\TextAt{8}{473.426}{388.382}{is}%
\TextAt{8}{485.237}{388.382}{no}%
\TextAt{8}{501.543}{388.382}{longer}%
\TextAt{8}{79.226}{371.336}{equivalent}%
\TextAt{8}{134.842}{371.336}{to}%
\TextAt{8}{149.152}{371.336}{the}%
\TextAt{8}{169.308}{371.336}{condition}%
\TextAt{8}{220.692}{371.336}{that}%
\TextAt{13}{246.061}{371.336}{a}%
\TextAt{8}{256.103}{371.336}{is}%
\TextAt{8}{267.867}{371.336}{a}%
\TextAt{8}{277.622}{371.336}{quadratic}%
\TextAt{8}{329.986}{371.336}{residue}%
\TextAt{8}{369.713}{371.336}{modulo}%
\TextAt{13}{411.664}{371.336}{b}%
\TextAt{8}{416.641}{371.336}{).}%
\TextAt{8}{97.159}{351.301}{In}%
\TextAt{8}{111.221}{351.301}{the}%
\TextAt{8}{130.818}{351.301}{above}%
\TextAt{8}{163.434}{351.301}{formulation,}%
\TextAt{8}{229.337}{351.301}{the}%
\TextAt{8}{248.934}{351.301}{reciprocity}%
\TextAt{8}{306.585}{351.301}{law}%
\TextAt{8}{327.162}{351.301}{allows}%
\TextAt{8}{361.464}{351.301}{us}%
\TextAt{8}{375.92}{351.301}{to}%
\TextAt{8}{389.671}{351.301}{compute}%
\TextAt{8}{436.585}{351.301}{simply}%
\TextAt{8}{473.481}{351.301}{by}%
\TextAt{8}{489.181}{351.301}{iteration}%
\TextAt{8}{79.226}{334.255}{the}%
\TextAt{8}{99.383}{334.255}{quadratic}%
\TextAt{8}{151.746}{334.255}{residue}%
\TextAt{8}{191.473}{334.255}{symbol}%
\TextAt{20}{231.535}{343.939}{\Character{0}}%
\TextAt{21}{238.21}{338.962}{a}%
\Vector{q 0 G 0 g 0.478 w 166.21 -436.756 m 170.708 -436.756 l S Q}%
\TextAt{21}{238.648}{330.132}{b}%
\TextAt{20}{243.904}{343.939}{\Character{1}}%
\TextAt{8}{249.383}{334.255}{,}%
\TextAt{8}{256.532}{334.255}{as}%
\TextAt{8}{270.902}{334.255}{is}%
\TextAt{8}{282.666}{334.255}{shown}%
\TextAt{8}{318.173}{334.255}{in}%
\TextAt{8}{331.826}{334.255}{the}%
\TextAt{8}{351.982}{334.255}{following}%
\TextAt{8}{401.417}{334.255}{example:}%
\TextAt{20}{105.151}{319.561}{\Character{18}}%
\TextAt{9}{115.147}{310.792}{40077}%
\Vector{q 0 G 0 g 0.478 w 43.147 -468.307 m 72.412 -468.307 l S Q}%
\TextAt{9}{115.147}{294.503}{65537}%
\TextAt{20}{145.607}{319.561}{\Character{19}}%
\TextAt{9}{157.728}{302.704}{\Character{61}}%
\TextAt{20}{170.154}{319.561}{\Character{18}}%
\TextAt{9}{180.15}{310.792}{65537}%
\Vector{q 0 G 0 g 0.478 w 108.15 -468.307 m 137.415 -468.307 l S Q}%
\TextAt{9}{180.15}{294.503}{40077}%
\TextAt{20}{210.61}{319.561}{\Character{19}}%
\TextAt{9}{222.731}{302.704}{\Character{61}}%
\TextAt{20}{235.157}{319.561}{\Character{18}}%
\TextAt{9}{245.153}{310.792}{25460}%
\Vector{q 0 G 0 g 0.478 w 173.153 -468.307 m 202.418 -468.307 l S Q}%
\TextAt{9}{245.153}{294.503}{40077}%
\TextAt{20}{275.613}{319.561}{\Character{19}}%
\TextAt{9}{287.734}{302.704}{\Character{61}}%
\TextAt{20}{300.16}{319.561}{\Character{18}}%
\TextAt{9}{319.496}{310.792}{2}%
\TextAt{17}{325.349}{315.13}{2}%
\Vector{q 0 G 0 g 0.478 w 238.156 -468.307 m 267.421 -468.307 l S Q}%
\TextAt{9}{310.156}{294.503}{40077}%
\TextAt{20}{340.616}{319.561}{\Character{19}}%
\TextAt{20}{351.4}{319.561}{\Character{18}}%
\TextAt{9}{364.331}{310.792}{6365}%
\Vector{q 0 G 0 g 0.478 w 289.405 -468.307 m 318.67 -468.307 l S Q}%
\TextAt{9}{361.405}{294.503}{40077}%
\TextAt{20}{391.865}{319.561}{\Character{19}}%
\TextAt{9}{403.987}{302.704}{\Character{61}}%
\TextAt{20}{416.412}{319.561}{\Character{18}}%
\TextAt{9}{426.408}{310.792}{40077}%
\Vector{q 0 G 0 g 0.478 w 354.408 -468.307 m 383.673 -468.307 l S Q}%
\TextAt{9}{429.334}{294.503}{6365}%
\TextAt{20}{456.868}{319.561}{\Character{19}}%
\TextAt{9}{157.728}{269.799}{\Character{61}}%
\TextAt{20}{170.154}{286.656}{\Character{18}}%
\TextAt{9}{180.15}{277.887}{1887}%
\Vector{q 0 G 0 g 0.478 w 108.15 -501.212 m 131.562 -501.212 l S Q}%
\TextAt{9}{180.15}{261.598}{6365}%
\TextAt{20}{204.757}{286.656}{\Character{19}}%
\TextAt{9}{216.878}{269.799}{\Character{61}}%
\TextAt{20}{229.304}{286.656}{\Character{18}}%
\TextAt{9}{239.3}{277.887}{6365}%
\Vector{q 0 G 0 g 0.478 w 167.3 -501.212 m 190.712 -501.212 l S Q}%
\TextAt{9}{239.3}{261.598}{1887}%
\TextAt{20}{263.907}{286.656}{\Character{19}}%
\TextAt{9}{276.028}{269.799}{\Character{61}}%
\TextAt{20}{288.454}{286.656}{\Character{18}}%
\TextAt{9}{301.376}{277.887}{704}%
\Vector{q 0 G 0 g 0.478 w 226.45 -501.212 m 249.862 -501.212 l S Q}%
\TextAt{9}{298.45}{261.598}{1887}%
\TextAt{20}{323.057}{286.656}{\Character{19}}%
\TextAt{9}{335.178}{269.799}{\Character{61}}%
\TextAt{20}{347.604}{286.656}{\Character{18}}%
\TextAt{9}{364.013}{277.887}{4}%
\TextAt{17}{369.866}{282.225}{3}%
\Vector{q 0 G 0 g 0.478 w 285.6 -501.212 m 309.012 -501.212 l S Q}%
\TextAt{9}{357.6}{261.598}{1887}%
\TextAt{20}{382.207}{286.656}{\Character{19}}%
\TextAt{20}{392.991}{286.656}{\Character{18}}%
\TextAt{9}{408.849}{277.887}{11}%
\Vector{q 0 G 0 g 0.478 w 330.996 -501.212 m 354.408 -501.212 l S Q}%
\TextAt{9}{402.996}{261.598}{1887}%
\TextAt{20}{427.603}{286.656}{\Character{19}}%
\TextAt{9}{439.725}{269.799}{\Character{61}}%
\TextAt{16}{452.15}{269.799}{\Character{0}}%
\TextAt{20}{463.441}{286.656}{\Character{18}}%
\TextAt{9}{473.437}{277.887}{1887}%
\Vector{q 0 G 0 g 0.478 w 401.437 -501.212 m 424.849 -501.212 l S Q}%
\TextAt{9}{479.29}{261.598}{11}%
\TextAt{20}{498.045}{286.656}{\Character{19}}%
\TextAt{9}{157.728}{237.121}{\Character{61}}%
\TextAt{16}{170.154}{237.121}{\Character{0}}%
\TextAt{20}{181.445}{253.978}{\Character{18}}%
\TextAt{9}{194.367}{245.209}{6}%
\Vector{q 0 G 0 g 0.478 w 119.441 -533.89 m 131.147 -533.89 l S Q}%
\TextAt{9}{191.441}{228.92}{11}%
\TextAt{20}{204.342}{253.978}{\Character{19}}%
\TextAt{9}{216.463}{237.121}{\Character{61}}%
\TextAt{16}{228.889}{237.121}{\Character{0}}%
\TextAt{20}{240.18}{253.978}{\Character{18}}%
\TextAt{9}{253.102}{245.209}{2}%
\Vector{q 0 G 0 g 0.478 w 178.176 -533.89 m 189.882 -533.89 l S Q}%
\TextAt{9}{250.176}{228.92}{11}%
\TextAt{20}{263.077}{253.978}{\Character{19}}%
\TextAt{20}{273.861}{253.978}{\Character{18}}%
\TextAt{9}{286.792}{245.209}{3}%
\Vector{q 0 G 0 g 0.478 w 211.866 -533.89 m 223.572 -533.89 l S Q}%
\TextAt{9}{283.866}{228.92}{11}%
\TextAt{20}{296.767}{253.978}{\Character{19}}%
\TextAt{9}{308.889}{237.121}{\Character{61}}%
\TextAt{20}{321.314}{253.978}{\Character{18}}%
\TextAt{9}{334.237}{245.209}{3}%
\Vector{q 0 G 0 g 0.478 w 259.31 -533.89 m 271.016 -533.89 l S Q}%
\TextAt{9}{331.31}{228.92}{11}%
\TextAt{20}{344.212}{253.978}{\Character{19}}%
\TextAt{9}{356.333}{237.121}{\Character{61}}%
\TextAt{16}{368.758}{237.121}{\Character{0}}%
\TextAt{20}{380.049}{253.978}{\Character{18}}%
\TextAt{9}{390.045}{245.209}{11}%
\Vector{q 0 G 0 g 0.478 w 318.045 -533.89 m 329.751 -533.89 l S Q}%
\TextAt{9}{392.972}{228.92}{3}%
\TextAt{20}{402.947}{253.978}{\Character{19}}%
\TextAt{9}{415.068}{237.121}{\Character{61}}%
\TextAt{16}{427.493}{237.121}{\Character{0}}%
\TextAt{20}{438.784}{253.978}{\Character{18}}%
\TextAt{9}{448.78}{245.209}{2}%
\Vector{q 0 G 0 g 0.478 w 376.78 -533.89 m 382.633 -533.89 l S Q}%
\TextAt{9}{448.78}{228.92}{3}%
\TextAt{20}{455.829}{253.978}{\Character{19}}%
\TextAt{9}{467.95}{237.121}{\Character{61}}%
\TextAt{9}{480.36}{237.121}{1}%
\TextAt{13}{486.228}{237.121}{\Character{58}}%
\TextAt{8}{79.226}{207.543}{Class}%
\TextAt{8}{109.553}{207.543}{field}%
\TextAt{8}{134.548}{207.543}{theory}%
\TextAt{8}{170.936}{207.543}{originated}%
\TextAt{8}{225.867}{207.543}{from}%
\TextAt{8}{253.145}{207.543}{Gauss’s}%
\TextAt{8}{295.32}{207.543}{reciprocity}%
\TextAt{8}{353.168}{207.543}{law.}%
\TextAt{8}{378.735}{207.543}{The}%
\TextAt{8}{402.427}{207.543}{quest}%
\TextAt{8}{433.017}{207.543}{for}%
\TextAt{8}{450.54}{207.543}{a}%
\TextAt{8}{459.933}{207.543}{similar}%
\TextAt{8}{498.007}{207.543}{law}%
\TextAt{8}{518.782}{207.543}{for}%
\TextAt{8}{79.226}{190.497}{the}%
\TextAt{13}{99.177}{190.497}{n}%
\TextAt{8}{106.164}{190.497}{-th}%
\TextAt{8}{124.812}{190.497}{power}%
\TextAt{8}{158.75}{190.497}{residues}%
\TextAt{8}{202.885}{190.497}{dominated}%
\TextAt{8}{260.566}{190.497}{number}%
\TextAt{8}{302.957}{190.497}{theory}%
\TextAt{8}{339.501}{190.497}{for}%
\TextAt{8}{357.18}{190.497}{a}%
\TextAt{8}{366.73}{190.497}{long}%
\TextAt{8}{391.892}{190.497}{time,}%
\TextAt{8}{421.64}{190.497}{and}%
\TextAt{8}{444.196}{190.497}{the}%
\TextAt{8}{464.147}{190.497}{all-embracing}%
\TextAt{8}{79.226}{173.451}{answer}%
\TextAt{8}{117.453}{173.451}{was}%
\TextAt{8}{139.42}{173.451}{finally}%
\TextAt{8}{174.335}{173.451}{found}%
\TextAt{8}{206.643}{173.451}{in}%
\TextAt{8}{219.764}{173.451}{Artin’s}%
\TextAt{8}{258.636}{173.451}{reciprocity}%
\TextAt{8}{316.313}{173.451}{law.}%
\TextAt{8}{341.823}{173.451}{The}%
\TextAt{8}{365.345}{173.451}{above}%
\TextAt{8}{397.988}{173.451}{reciprocity}%
\TextAt{8}{455.665}{173.451}{law}%
\TextAt{8}{476.27}{173.451}{(8.3)}%
\TextAt{8}{503.713}{173.451}{of}%
\TextAt{8}{516.51}{173.451}{the}%
\TextAt{8}{79.226}{156.405}{power}%
\TextAt{8}{112.477}{156.405}{residues}%
\TextAt{8}{155.926}{156.405}{now}%
\TextAt{8}{179.422}{156.405}{appears}%
\TextAt{8}{221.842}{156.405}{as}%
\TextAt{8}{235.319}{156.405}{a}%
\TextAt{8}{244.182}{156.405}{simple}%
\TextAt{8}{279.764}{156.405}{and}%
\TextAt{8}{301.634}{156.405}{special}%
\TextAt{8}{338.842}{156.405}{consequence}%
\TextAt{8}{404.013}{156.405}{of}%
\TextAt{8}{416.45}{156.405}{Artin’s}%
\TextAt{8}{454.962}{156.405}{reciprocity}%
\TextAt{8}{512.278}{156.405}{law.}%
\TextAt{8}{79.226}{139.359}{But}%
\TextAt{8}{103.113}{139.359}{to}%
\TextAt{8}{118.069}{139.359}{really}%
\TextAt{8}{150.911}{139.359}{settle}%
\TextAt{8}{182.831}{139.359}{the}%
\TextAt{8}{203.634}{139.359}{original}%
\TextAt{8}{246.565}{139.359}{problem,}%
\TextAt{8}{296.15}{139.359}{class}%
\TextAt{8}{324.233}{139.359}{field}%
\TextAt{8}{350.236}{139.359}{theory}%
\TextAt{8}{387.632}{139.359}{was}%
\TextAt{8}{410.778}{139.359}{still}%
\TextAt{8}{434.246}{139.359}{lacking}%
\TextAt{8}{474.571}{139.359}{the}%
\TextAt{8}{495.374}{139.359}{explicit}%
\TextAt{8}{79.226}{122.313}{computation}%
\TextAt{8}{147.743}{122.313}{of}%
\TextAt{8}{161.291}{122.313}{the}%
\TextAt{8}{181.664}{122.313}{Hilbert}%
\TextAt{8}{222.195}{122.313}{symbols}%
\TextAt{20}{267.088}{135.584}{\Character{16}}%
\TextAt{21}{275.424}{127.569}{a\Character{59}b}%
\Vector{q 0 G 0 g 0.478 w 203.424 -648.698 m 213.897 -648.698 l S Q}%
\TextAt{27}{278.666}{118.191}{p}%
\TextAt{20}{287.092}{135.584}{\Character{17}}%
\TextAt{8}{298.347}{122.313}{for}%
\TextAt{14}{316.449}{122.313}{p}%
\TextAt{16}{326.123}{122.313}{\Character{106}}%
\TextAt{13}{333.135}{122.313}{n}%
\NativeMathAt{340.123}{122.313}{11.955}{\infty}%
\TextAt{8}{352.078}{122.313}{.}%
\TextAt{8}{361.178}{122.313}{This}%
\TextAt{8}{388.115}{122.313}{was}%
\TextAt{8}{410.832}{122.313}{finally}%
\TextAt{8}{446.497}{122.313}{completed}%
\TextAt{8}{502.64}{122.313}{in}%
\TextAt{8}{516.51}{122.313}{the}%
\TextAt{8}{79.226}{105.267}{1960s}%
\TextAt{8}{111.17}{105.267}{by}%
\TextAt{8}{127.429}{105.267}{the}%
\TextAt{8}{147.586}{105.267}{mathematician}%
\TextAt{26}{227.59}{105.267}{Helmut}%
\TextAt{26}{277.969}{105.267}{Brückner}%
\TextAt{8}{337.314}{105.267}{,}%
\TextAt{8}{344.464}{105.267}{see}%
\TextAt{8}{363.377}{105.267}{chap.}%
\TextAt{8}{394.269}{105.267}{V,}%
\TextAt{8}{410.193}{105.267}{(3.7).}%
\TextAt{8}{297.211}{43.033}{380}%
\Destination{78.23}{809.815}{nav.page.380}%
\Destination{79.36}{776.155}{nav.print.416}%
\Destination{79.36}{420.455}{nav.print.417}%
\Hotlink{184.635}{741.855}{5.22}{11.047}{nav.chapter.1}%
\Hotlink{371.466}{707.755}{9.775}{11.047}{nav.chapter.5}%
\Hotlink{109.618}{673.655}{5.22}{11.047}{nav.chapter.1}%
\Hotlink{393.769}{103.055}{9.775}{11.047}{nav.chapter.5}%
\Hotlink{195.331}{741.855}{25.078}{11.047}{nav.number.1.8.6}%
\Hotlink{199.81}{724.855}{25.078}{11.047}{nav.number.6.8.3}%
\Hotlink{184.313}{707.755}{25.078}{11.047}{nav.number.6.8.3}%
\Hotlink{387.867}{707.755}{25.078}{11.047}{nav.number.5.3.6}%
\Hotlink{120.988}{673.655}{25.078}{11.047}{nav.number.1.8.6}%
\Hotlink{475.77}{171.255}{25.078}{11.047}{nav.number.6.8.3}%
\Hotlink{409.693}{103.055}{25.078}{11.047}{nav.number.5.3.7}%
\Destination{79.226}{452.555}{errata.35}%
}


\input{Corrected/chapter7}
\input{Corrected/backmatter}
\end{document}
